\documentclass{book}
%\usepackage[letterpaper, portrait, margin=1cm]{geometry}
%\usepackage[letterpaper, bindingoffset=0.2in, left=1in,right=1in,top=.5in,bottom=.5in,footskip=.25in,marginparwidth=5em]{geometry}
\usepackage[letterpaper, left=1in,right=1in,top=.5in,bottom=.5in,footskip=.25in,marginparwidth=1cm]{geometry}
% ---------------------
% mini-table-of-content
% ---------------------
\usepackage{minitoc}
\setcounter{minitocdepth}{1}
\setlength{\mtcindent}{24pt}
\setcounter{secnumdepth}{-2}
%\renewcommand{\mtcfont}{\small\rm}
%\renewcommand{\mtcSfont}{\small\bf}
%\usepackage{setspace}
%\usepackage{tocloft}
%\setlength\cftparskip{-1.2pt}
%\setlength\cftbeforesecskip{1.3pt}
%\setlength\cftaftertoctitleskip{2pt}
%\renewcommand{\cftsecafterpnum}{\hspace*{02.0em}}
%\renewcommand{\cftsubsecafterpnum}{\hspace*{02.0em}}

% ---------------------------
% Chinese Characters Packages
% ---------------------------
\usepackage{fontspec} 
\usepackage{xeCJK}
\setmainfont{Times}
\setCJKmainfont{BiauKai}
\newfontfamily\sblgoodhebrew{SBL BibLit}[Script=Hebrew,Contextuals=Alternate]
\newfontfamily\sblgoodgreek{SBL BibLit}[Script=Greek,Contextuals=Alternate]

\usepackage{ifpdf,cite,algorithmic,url,tikz}
\usepackage[cmex10]{amsmath}

% ---------------------------
% Hebrew Characters Packages
% ---------------------------
\usepackage{polyglossia}
\setmainfont{Times New Roman}

% -------
% General
% -------
\usepackage{multicol}
\usepackage{multirow}
\usepackage{color,colortbl}
\usepackage{xparse}
\usepackage{pbox}
\usepackage{stackengine}
\usepackage{titlesec}% http://ctan.org/pkg/titlesec
\usepackage{tabularx}
\usepackage{xltabular}
\usepackage{titlesec}
\usepackage{makecell}
\newcommand{\sectionbreak}{\clearpage}

\author{
  Editor, Michael Chan\\
  \texttt{michaelchan\_wahyan@yahoo.com.hk}
}
\usepackage{tocloft}

\usepackage{hyperref}
\hypersetup{
    colorlinks=true, % set true if you want colored links
    linktoc   =all , % set to all if you want both sections and subsections linked
    linkcolor =blue, % choose some color if you want links to stand out
}

% ----------
% Afterword
% ----------
\usepackage{marginnote}
\usepackage{sectsty}
\usepackage{ragged2e}
\usepackage{lineno}
\usepackage{xcolor}
\usepackage{paracol}

\begin{document}

\clearpage
%% temporary titles
% command to provide stretchy vertical space in proportion
\newcommand\nbvspace[1][3]{\vspace*{\stretch{#1}}}
% allow some slack to avoid under/overfull boxes
\newcommand\nbstretchyspace{\spaceskip0.5em plus 0.25em minus 0.25em}
% To improve spacing on titlepages
\newcommand{\nbtitlestretch}{\spaceskip0.6em}
\pagestyle{empty}
\begin{center}
\bfseries
\nbvspace[1]
\Huge
{%\nbtitlestretch
\Large
\textbf{香港中文大學崇基神學院~~校牧主日崇拜講章~~2009-present \\
       Chaplaincy, Divinity School of Chung Chi College, CUHK
       }}

\nbvspace[1]

{\large
Editor: Michael\\
\texttt{michaelchan\_wahyan@yahoo.com.hk}
}

\nbvspace[1]

{\large
Revision: \texttt{v1.3}\\
Last Update: \today
}


\vfill
\begin{tikzpicture}
    %remove comment for OT cover%\node (0,0) [opacity=0.03]{\includegraphics[width=15cm]{../bible_out/ot_frontcover.png}} ;
    %remove comment for NT cover%\node (0,0) [opacity=0.03]{\includegraphics[width=15cm]{../bible_out/christ_on_cross.png}} ;
    %remove comment for Bible cover%\node (0,0) [xshift=0.8cm, yshift=+2cm, opacity=0.03]{\includegraphics[width=10cm]{./christ_on_cross.png}} ;
    %remove comment for Bible cover%\node (0,0) [              yshift=-2cm, opacity=0.03]{\includegraphics[width=14cm]{./ot_frontcover.png}} ;
\end{tikzpicture}
\vfill

\end{center}

\newpage

\setcounter{tocdepth}{0}
\dominitoc
\begin{multicols}{3}
\addtocontents{toc}{\protect\hypertarget{toc}{}}
\tableofcontents
\end{multicols}

\large
%\twocolumn

% the color definition syntax is as follow:
% \definecolor{name}{system}{definition}
% example: a mono-channel color can be defined as
%          \definecolor{Gray}{gray}{0.9}
% example: an rgb-3-channel color can be defined as
%          \definecolor{LightCyan}{rgb}{0.88,1,1}
%          \definecolor{pink}{rgb}{0.68,0,0.68}

\definecolor{CUV1LightRed}{rgb}{1,0.75,0.75}     % for CUV1
\definecolor{LZZVLightGray}{rgb}{0.9,0.9,0.9}    % for LZZ
\definecolor{KJVVLightGreen}{rgb}{0.75,1,0.85}   % for KJV
\definecolor{CUV2LightYellow}{rgb}{1,1,0.75}     % for CUV2
\definecolor{CNVVLightBrown}{rgb}{1,0.85,0.7}    % for CNV
\definecolor{NRSVLightBlue}{rgb}{0.75,1,1}       % for NRSV
\definecolor{WENLLightPurple}{rgb}{0.95,0.85,0.9}% for WENL
\definecolor{TCV19PaleGreen}{rgb}{0.85,1,0.95}   % for TCV19
\definecolor{MSGVLightWhite}{rgb}{0.98,0.98,0.98}% for MSGV
\definecolor{NETSLightRed}{rgb}{1,0.75,0.75}     % for NETS
\definecolor{JPS1917LightYellow}{rgb}{1,1,0.75}  % for JPS1917
\definecolor{SBLGNTPaleRed}{rgb}{1,0.85,0.80}    % for SBLGNT

\section{目錄}\label{sec:toc}
{ \scriptsize


\begin{xltabular}{\textwidth}{|p{0.15\textwidth}|p{0.15\textwidth}|p{0.4\textwidth}|p{0.3\textwidth}|}
\hline
日期 & 講員 & 經課 & 講題 \\ \hline \hline
\hline \hline
2009年6月14日 & 伍渭文牧師 & 詩篇1:3 & \hyperref[sec:0]{栽在溪水旁的樹} \\
2009年6月21日 & 伍渭文牧師 & 約翰福音4:5-29 & \hyperref[sec:1]{喜出望外的邂逅} \\
2009年6月28日 & 林豪恩先生 & 馬可福音4:35-41 & \hyperref[sec:2]{為甚麼膽怯呢} \\
\hline
2009年7月5日 & 伍渭文牧師 & 馬可福音16:1-8 & \hyperref[sec:3]{誰滾開那石頭} \\
2009年7月12日 & 鄧瑞強博士 & 馬可福音5 & \hyperref[sec:4]{靈機一觸} \\
2009年7月19日 & 孫寶玲牧師 & 馬可福音6:1-6 & \hyperref[sec:5]{知不知} \\
2009年7月26日 & 陳孟賢博士 & 雅各書1:9-11 & \hyperref[sec:6]{超越貧與富} \\
\hline
2009年8月2日 & 伍渭文牧師 & 約翰福音9:24-41 & \hyperref[sec:7]{好牧人與雇工} \\
2009年8月9日 & 賀志勇博士 & 路加福音10:38-42 & \hyperref[sec:8]{你不能缺少的一件事} \\
2009年8月16日 & 梁元生教授 &  & \hyperref[sec:9]{曠野行旅嗎哪} \\
2009年8月23日 & 鄧瑞強博士 & 約翰福音6:35,41-51 & \hyperref[sec:10]{餓孩子飽孩子} \\
2009年8月30日 & Prof. Dennis P. McCann &  & \hyperref[sec:11]{沒有群羊的世界} \\
\hline
2009年9月13日 & 盧龍光牧師 & 約翰福音6 & \hyperref[sec:12]{我的肉是可喫的} \\
2009年9月20日 & 賴品超教授 & 約翰福音6:69-71 & \hyperref[sec:13]{我們還跟從誰呢} \\
2009年9月27日 & 戴國璋牧師 & 馬可福音12:28-34 & \hyperref[sec:14]{律法遺傳與良心} \\
\hline
2009年10月4日 & 伍渭文牧師 & 馬可福音7:24-30 & \hyperref[sec:15]{從桌子掉下來的餅碎} \\
2009年10月11日 & 陳南芬教授 &  & \hyperref[sec:16]{外人在摩押地} \\
2009年10月18日 & 伍渭文牧師 & 馬可福音9:30-37 & \hyperref[sec:17]{領袖僕人和小孩} \\
2009年10月25日 & 林豪恩先生 & 馬可福音9:38-50 & \hyperref[sec:18]{拴在頸上的石磨} \\
\hline
2009年11月1日 & 伍渭文牧師 & 希伯來書2 & \hyperref[sec:19]{人算什麼你竟眷顧他} \\
2009年11月8日 & 賀志勇博士 & 路加福音14:25-33 & \hyperref[sec:20]{你總要付的代價} \\
2009年11月15日 & 林豪恩先生 & 馬可福音10:35-45 & \hyperref[sec:21]{你知道所求的是甚麼嗎} \\
2009年11月22日 & 鄧瑞強博士 & 馬可福音10:46-52 & \hyperref[sec:22]{我要能看見} \\
2009年11月29日 & 伍渭文牧師 & 約翰福音11:32-44 & \hyperref[sec:23]{你們把石頭挪開} \\
\hline
2009年12月6日 & 胡美儀 &  & \hyperref[sec:24]{我行過死蔭的幽谷} \\
2009年12月13日 & 陳龍斌教授 &  & \hyperref[sec:25]{石頭記} \\
2009年12月20日 & 梁元生教授 & 馬太福音2;路加福音2 & \hyperref[sec:26]{十二月的星空} \\
2009年12月27日 & 伍渭文牧師 & 路加福音2:12 & \hyperref[sec:27]{新生王的記號} \\
\hline \hline
2010年1月3日 & 鄧瑞強博士 & 路加福音1:46-53 & \hyperref[sec:28]{我心尊主為大} \\
2010年1月10日 & 鄧瑞強博士 & 約翰福音18:33-38;啟示錄1:4-8 & \hyperref[sec:29]{誰是老闆} \\
2010年1月17日 & 戴國璋牧師 &  & \hyperref[sec:30]{一件美事} \\
2010年1月24日 & 林豪恩先生 & 路加福音2:41-52 & \hyperref[sec:31]{問道的小孩} \\
2010年1月31日 & 陳龍斌教授 & 以弗所書3:14 & \hyperref[sec:32]{仰望星空} \\
\hline
2010年2月7日 & 戴國璋牧師 &  & \hyperref[sec:33]{上好的福份} \\
2010年2月14日 & 賀志勇博士 & 路加福音8:40-56 & \hyperref[sec:34]{隱藏有時顯明有時} \\
2010年2月21日 & 林豪恩先生 & 路加福音4:14-21 & \hyperref[sec:35]{綑縛與釋放} \\
2010年2月28日 & 鄧瑞強博士 & 路加福音4:16-30 & \hyperref[sec:36]{獨食難肥} \\
\hline
2010年3月7日 & 伍渭文牧師 & 以賽亞書6:1-8 & \hyperref[sec:37]{我可以差遣誰呢} \\
2010年3月14日 & 鄧瑞強博士 & 路加福音9:21-36 & \hyperref[sec:38]{光明在人間} \\
2010年3月21日 & 林豪恩先生 & 路加福音4:1-13 & \hyperref[sec:39]{再細看試探的臉龐} \\
2010年3月28日 & 賀志勇博士 & 路加福音22 & \hyperref[sec:40]{認信背棄和重建} \\
\hline
2010年4月4日 & 鄧瑞強博士 & 馬太福音5:9 & \hyperref[sec:41]{使人和睦的人有福了} \\
2010年4月11日 & 林豪恩先生 & 路加福音19:28-41 & \hyperref[sec:42]{在橄欖山與聖殿山之間} \\
2010年4月18日 & 伍渭文牧師 &  & \hyperref[sec:43]{為義受逼迫的人有福了} \\
2010年4月25日 & Dr. Kuruvilla Abraham &  & \hyperref[sec:44]{讓復活與我們走過每一天} \\
\hline
2010年5月2日 & 陳曉東博士 & 約翰福音20:19-31 & \hyperref[sec:45]{願你們平安} \\
2010年5月9日 & 伍渭文牧師 & 詩篇30:1-12 & \hyperref[sec:46]{轉眼的怒氣一生的恩典} \\
2010年5月16日 & 鄧瑞強博士 & 詩篇23:1-6 & \hyperref[sec:47]{恩惠慈愛隨著我} \\
2010年5月30日 & 梁元生教授 & 以西結書17:1-3,22-24 & \hyperref[sec:48]{尋找一棵香柏樹} \\
\hline
2010年6月6日 & 林豪恩先生 & 約翰福音14:15-31 & \hyperref[sec:49]{離開祇留下愛} \\
2010年6月13日 & 鄧瑞強博士 & 使徒行傳2:1-21 & \hyperref[sec:50]{我沒有醉} \\
2010年6月20日 & 賴品超教授 & 羅馬書5 & \hyperref[sec:51]{施行拯救的三一上帝} \\
2010年6月27日 & 伍渭文牧師 &  & \hyperref[sec:52]{漫漫朝聖路} \\
\hline
2010年7月4日 & 賀志勇博士 & 路加福音7:36-50 & \hyperref[sec:53]{你可以行多一步嗎} \\
2010年7月11日 & 林豪恩先生 & 路加福音8:26-39 & \hyperref[sec:54]{一個人值多少} \\
2010年7月18日 & 鄧瑞強博士 & 路加福音9:57-62 & \hyperref[sec:55]{重價的恩典} \\
2010年7月25日 & 伍渭文牧師 & 路加福音11:17-24 & \hyperref[sec:56]{沒有甚麼能傷害你們了} \\
\hline
2010年8月1日 & 曹偉彤博士 &  & \hyperref[sec:57]{使生命完滿的愛} \\
2010年8月8日 & 林豪恩先生 & 路加福音10:25-37 & \hyperref[sec:58]{如果愛是真的} \\
2010年8月15日 & 賀志勇博士 & 路加福音7:1-10 & \hyperref[sec:59]{信心與權柄} \\
2010年8月22日 & 伍渭文牧師 & 路12:13-21 & \hyperref[sec:60]{窮得祇剩下錢} \\
2010年8月29日 & 張美珍博士 & 路加福音12 & \hyperref[sec:61]{看得見的與看不見的} \\
\hline
2010年9月5日 & 賀志勇博士 & 路加福音12:49-53 & \hyperref[sec:62]{戰爭與和平} \\
2010年9月12日 & 王群女士 & 馬可福音6:30-44 & \hyperref[sec:63]{你們給他們吃吧} \\
2010年9月19日 & 林豪恩先生 & 路加福音14:1-24 & \hyperref[sec:64]{戰雲飯局} \\
2010年9月26日 & 賴品超教授 & 阿摩司書8:4-7;詩篇113;提摩太前書2:1-7;路加福音16:1-13 & \hyperref[sec:65]{個人理財與社會公義} \\
\hline
2010年10月3日 & 姚西依博士 & 使徒行傳10 & \hyperref[sec:66]{普世的福音} \\
2010年10月10日 & 戴國璋牧師 &  & \hyperref[sec:67]{一片燒魚和一塊蜜房} \\
2010年10月17日 & 林豪恩先生 & 路加福音16:19-31 & \hyperref[sec:68]{乞丐與財主} \\
2010年10月24日 & 鄧瑞強博士 & 詩篇37:1-40 & \hyperref[sec:69]{卑微者要擁有土地} \\
2010年10月31日 & 王群女士 & 哥林多後書5:11-21 & \hyperref[sec:70]{基督的特使} \\
\hline
2010年11月7日 & 林豪恩先生 & 路加福音12:16-21,16:19-31,18:18-30 & \hyperref[sec:71]{路加福音的最後財主} \\
2010年11月14日 & 陳曉東博士(普通話) & 哥林多後書5:11-21 & \hyperref[sec:72]{感恩的痲瘋病人} \\
2010年11月21日 & 沈祖堯教授 & 約書亞記1 & \hyperref[sec:73]{你的應許之地} \\
2010年11月28日 & 周小安牧師 & 約翰福音4:13-14 & \hyperref[sec:74]{給你喝一杯生命之水} \\
\hline
2010年12月5日 & 梁元生教授 & 使徒行傳8:26-40 & \hyperref[sec:75]{腓利的旅程} \\
2010年12月12日 & 王建宣先生 & 詩篇23:1-6 & \hyperref[sec:76]{耶和華是我的牧者} \\
2010年12月19日 & 羅勝強教授 &  & \hyperref[sec:77]{如此我信} \\
2010年12月26日 & 伍渭文牧師 & 路加福音1:26-45 & \hyperref[sec:78]{聖誕節第一個問安} \\
\hline \hline
2011年1月2日 & 林豪恩先生 & 路加福音3:1-6 & \hyperref[sec:79]{聖言哪裡去了} \\
2011年1月9日 & 鄧瑞強博士 & 路加福音3:7-18 & \hyperref[sec:80]{及時行善} \\
2011年1月16日 & 盧龍光牧師 & 路加福音2:41-52 & \hyperref[sec:81]{為上主的事著迷} \\
2011年1月23日 & 伍渭文牧師 & 提摩太後書4: & \hyperref[sec:82]{這個冬天不算冷} \\
2011年1月30日 & 賀志勇博士 & 馬太福音3:13-17 & \hyperref[sec:83]{盡諸般的義} \\
\hline
2011年2月6日 & 林豪恩先生 & 約翰福音1:29-42 & \hyperref[sec:84]{你住在哪裡} \\
2011年2月13日 & 鄧瑞強博士 & 馬太福音4:12-23 & \hyperref[sec:85]{無法抗拒的呼喚} \\
2011年2月20日 & 高志棻博士 & 彌迦書6:8 & \hyperref[sec:86]{與神同行奔天路} \\
2011年2月27日 & 陳榮超牧師 & 馬太福音5:9 & \hyperref[sec:87]{和睦是福} \\
\hline
2011年3月6日 & 林豪恩先生 & 利未記19:9-18 & \hyperref[sec:88]{西乃山下與獅子山下} \\
2011年3月13日 & 陳曉東博士 &  & \hyperref[sec:89]{失去味的鹽味道如何呢} \\
2011年3月20日 & 鄧瑞強博士 & 馬太福音6:24-34 & \hyperref[sec:90]{飛鳥野花神的國} \\
2011年3月27日 & 伍渭文牧師 & 馬太福音17:1-9 & \hyperref[sec:91]{登山變像命途} \\
\hline
2011年4月3日 & 朱建忠先生 & 馬太福音4:1-11 & \hyperref[sec:92]{人活著不是單靠食物} \\
2011年4月10日 & 林豪恩先生 & 約翰福音3:1-17 & \hyperref[sec:93]{起跑線上定輸贏} \\
2011年4月17日 & 鄧瑞強博士 & 約翰福音4 & \hyperref[sec:94]{井旁的邂逅} \\
2011年4月24日 & 鄧瑞強博士 & 馬太福音21:1-11 & \hyperref[sec:95]{棕樹枝革命} \\
\hline
2011年5月1日 & 林豪恩先生 & 約翰福音20:11-18 & \hyperref[sec:96]{抹大拉的馬利亞也復活} \\
2011年5月8日 & 伍渭文牧師 & 約翰福音9 & \hyperref[sec:97]{看見真相的代價} \\
2011年5月15日 & 賀志勇博士 & 希伯來書12:1-2 & \hyperref[sec:98]{人生的馬拉松} \\
2011年5月19日 & 鄧瑞強博士 & 路加福音24:13-35 & \hyperref[sec:99]{在途上} \\
2011年5月22日 & 伍渭文牧師 & 詩篇16 & \hyperref[sec:100]{虛假與代替} \\
\hline
2011年6月5日 & 王教佺先生 & 詩133:1-3 & \hyperref[sec:101]{黑門的甘露} \\
2011年6月12日 & 林豪恩先生 & 約翰福音14:1-14 & \hyperref[sec:102]{愁眉得展} \\
2011年6月19日 & 鄧瑞強博士 & 約翰福音17:1-11 & \hyperref[sec:103]{參悟死生} \\
2011年6月26日 & 梁元生教授 &  & \hyperref[sec:104]{都市牧羊} \\
\hline
2011年7月3日 & 伍渭文牧師 & 約翰福音15:1-5 & \hyperref[sec:105]{真葡萄樹與紅酒的故事} \\
2011年7月10日 & 戴國璋牧師 & 約翰福音14:21-24 & \hyperref[sec:106]{在愛與真理中} \\
2011年7月17日 & 陳以諾牧師 & 以弗所書4:15-16 & \hyperref[sec:107]{在愛中建立自己} \\
2011年7月24日 & 林豪恩先生 & 羅馬書8:12-25 & \hyperref[sec:108]{無債一身輕} \\
2011年7月31日 & 賴品超教授 & 約拿書4:1-11 & \hyperref[sec:109]{知與情的交戰} \\
\hline
2011年8月7日 & 戴國璋牧師 & 馬太福音13:10-16 & \hyperref[sec:110]{種子土壤與撒種的人} \\
2011年8月14日 & 周志豪牧師 & 羅馬書8:26-27 & \hyperref[sec:111]{禱告的能力} \\
2011年8月21日 & 賀志勇博士 & 創世紀28:10-22 & \hyperref[sec:112]{禱告初階} \\
2011年8月28日 & 林豪恩先生 & 列王記上19:1-18 & \hyperref[sec:113]{當你走到無力} \\
\hline
2011年9月4日 & 賀志勇博士 & 羅馬書12:1-2 & \hyperref[sec:114]{夏生命盛放的季節} \\
2011年9月11日 & 賴品超教授 &  & \hyperref[sec:115]{救恩的奧秘} \\
2011年9月18日 & 伍渭文牧師 & 創4:1-16 & \hyperref[sec:116]{伏在門前的野獸} \\
2011年9月25日 & 馮煒文先生 & 馬太福音16:21-28 & \hyperref[sec:117]{若有人要跟隨我} \\
\hline
2011年10月2日 & 林豪恩先生 & 馬太福音18:21-35 & \hyperref[sec:118]{七十個七次} \\
2011年10月9日 & 鄧瑞強博士 & 馬太福音21:23-27 & \hyperref[sec:119]{關於權威的討論} \\
2011年10月16日 & 梁元生教授 &  & \hyperref[sec:120]{基督徒的品牌} \\
2011年10月23日 & 陳曉東博士 &  & \hyperref[sec:121]{在後的必要在先} \\
2011年10月30日 & 鄧瑞強博士 & 馬太福音22:15-22 & \hyperref[sec:122]{因我們只有一個老闆} \\
\hline
2011年11月6日 & 張伯笠牧師 & 羅馬書1:16-17 & \hyperref[sec:123]{這福音是上帝的大能} \\
2011年11月13日 & 沈祖堯教授 & 詩篇46 & \hyperref[sec:124]{上主仍是我們911後的} \\
2011年11月20日 & 伍渭文牧師 &  & \hyperref[sec:125]{為福音而活的人生} \\
2011年11月27日 & 林豪恩先生 & 馬太福音25:31-46 & \hyperref[sec:126]{看不見耶穌的時候} \\
\hline
2011年12月4日 & 林豪恩先生 & 詩篇1 & \hyperref[sec:127]{果樹與糠秕} \\
2011年12月11日 & 鄧瑞強博士 & 詩篇80 & \hyperref[sec:128]{神啊復興我們} \\
2011年12月18日 & 李炳光牧師 &  & \hyperref[sec:129]{有的加給沒有奪去} \\
2011年12月25日 & 鄧瑞強博士 & 約翰福音1:1-14 & \hyperref[sec:130]{光明與黑暗} \\
\hline \hline
2012年1月1日 & 伍渭文牧師 & 約翰福音1:14-18 & \hyperref[sec:131]{恩典與真理} \\
2012年1月8日 & 伍偉亨教授 &  & \hyperref[sec:132]{義薄坤輿} \\
2012年1月15日 & 鄧瑞強博士 & 路加福音2:22-35 & \hyperref[sec:133]{救恩就在平凡處} \\
2012年1月22日 & 賀志勇博士 & 詩篇29:1-11 & \hyperref[sec:134]{禱告進階聆聽上主} \\
2012年1月29日 & 林豪恩先生 & 約翰福音1:43-51 & \hyperref[sec:135]{知遇之恩} \\
\hline
2012年2月5日 & 伍渭文牧師 & 創世記47 & \hyperref[sec:136]{如何看待富裕和貧窮} \\
2012年2月12日 & 伍渭文牧師 & 馬可福音1:21-28 & \hyperref[sec:137]{人鬼神} \\
2012年2月19日 & 李慶宏先生 & 詩篇121:1-8 & \hyperref[sec:138]{保障期永遠有效} \\
2012年2月26日 & 林豪恩先生 & 馬可福音1:40-45;列王記下5:1-14 & \hyperref[sec:139]{跨越界限} \\
\hline
2012年3月4日 & 賴品超教授 & 哥林多前書8:1-13 & \hyperref[sec:140]{吃還是不吃} \\
2012年3月11日 & 鄧瑞強博士 & 馬可福音1:9-13 & \hyperref[sec:141]{天上曠野人間} \\
2012年3月18日 & 陳曉東博士 &  & \hyperref[sec:142]{心靈誠實與祭祀} \\
2012年3月25日 & 伍渭文牧師 & 馬太福音22:1-14 & \hyperref[sec:143]{婚禮的禮服} \\
\hline
2012年4月29日 & 林豪恩先生 & 馬可福音11:1-11 & \hyperref[sec:144]{驢以載道} \\
\hline
2012年5月6日 & 陳榮超牧師 & 路加福音24:13-32 & \hyperref[sec:145]{撻著生命的動力} \\
2012年5月13日 & 伍渭文牧師 & 路加福音24:13-35 & \hyperref[sec:146]{你們為何愁煩} \\
2012年5月20日 & 鄧瑞強博士 & 約翰福音10:11-18 & \hyperref[sec:147]{狼來了怎麼辦} \\
2012年5月27日 & 陳曉東博士 & 約翰福音20:19-31 & \hyperref[sec:148]{願你們平安} \\
\hline
2012年6月10日 & 林豪恩先生 & 約翰福音15:1-8 & \hyperref[sec:149]{常在我心間} \\
2012年6月17日 & 伍渭文牧師 & 約翰福音4:19-26 & \hyperref[sec:150]{崇拜失樂園的逆轉} \\
2012年6月24日 & 鄧瑞強博士 & 約翰福音16:8-11 & \hyperref[sec:151]{我知錯了} \\
\hline
2012年7月1日 & 戴國璋牧師 &  & \hyperref[sec:152]{柏拉奇圖士} \\
2012年7月8日 & 鄧瑞強博士 & 馬可福音3:20-35 & \hyperref[sec:153]{他癲狂了} \\
2012年7月15日 & 鄭漢文博士 & 約翰福音3:1-17 & \hyperref[sec:154]{神的永在人的永生} \\
2012年7月22日 & 賀志勇博士 & 詩篇30:1-12 & \hyperref[sec:155]{哭泣與歡呼流淚的禱告} \\
2012年7月29日 & 林豪恩先生 & 哥林多後書12:1-10 & \hyperref[sec:156]{一根刺的恩典} \\
\hline
2012年8月5日 & 林豪恩先生 & 馬可福音4:35-41 & \hyperref[sec:157]{信仰路上的四句話} \\
2012年8月12日 & 賀志勇博士 & 以弗所書4:25-30 & \hyperref[sec:158]{別讓聖靈擔憂你} \\
2012年8月19日 & 梁燕城博士 & 約翰福音16:8-11 & \hyperref[sec:159]{單純與委身} \\
2012年8月26日 & 陳志強牧師 &  & \hyperref[sec:160]{豁免年} \\
\hline
2012年9月2日 & 梁元生教授 & 使徒行傳16:1-10 & \hyperref[sec:161]{躑躅邊城} \\
2012年9月9日 & 張美珍博士 &  & \hyperref[sec:162]{做個有心人} \\
2012年9月16日 & 伍渭文牧師 & 約翰福音6:51-58 & \hyperref[sec:163]{我是生命的糧} \\
2012年9月23日 & 林豪恩先生 & 約翰福音6:56-69 & \hyperref[sec:164]{你們也要離開嗎} \\
2012年9月30日 & 陳天祥博士 &  & \hyperref[sec:165]{神的豐盛充滿我們} \\
\hline
2012年10月7日 & 伍渭文牧師 & 雅各書2:1-17 & \hyperref[sec:166]{位置決定我們的腦袋} \\
2012年10月14日 & 鄧瑞強博士 & 馬可福音8:27-38 & \hyperref[sec:167]{天國的國民教育} \\
2012年10月21日 & 林豪恩先生 & 雅各書3:13-4:3 & \hyperref[sec:168]{智者必懷仁} \\
2012年10月28日 & 賀志勇博士 & 雅各書5:13-18 & \hyperref[sec:169]{家的盼望} \\
\hline
2012年11月4日 & 鄧瑞強博士 & 馬可福音10:2-9 & \hyperref[sec:170]{不要那樣問而要} \\
2012年11月11日 & 伍渭文牧師 & 以賽亞書52:13-53:12 & \hyperref[sec:171]{負傷的治療者} \\
2012年11月18日 & 張伯笠牧師 & 約翰一書4 & \hyperref[sec:172]{讓我們彼此相愛} \\
2012年11月25日 & 沈祖堯教授 & 利未記19:1-18 & \hyperref[sec:173]{守望你的鄰舍} \\
\hline
2012年12月2日 & 鄧瑞強博士 & 馬可福音12:38-44 & \hyperref[sec:174]{不正義制度下的窮寡婦} \\
2012年12月9日 & 林豪恩先生 & 約翰福音18:33-37;啟示錄1:4-8 & \hyperref[sec:175]{爾國爾民} \\
2012年12月16日 & 伍渭文牧師 & 詩25 & \hyperref[sec:176]{塗抹與記念} \\
2012年12月23日 & 張義南牧師 &  & \hyperref[sec:177]{豐盛生命} \\
2012年12月30日 & 鄧瑞強博士(香港神學院) & 路加福音1:39-45 & \hyperref[sec:178]{兩個不平凡的孕婦} \\
2012年12月30日 & 林豪恩先生 & 路加福音2:41-52 & \hyperref[sec:179]{孩子不見了} \\
\hline \hline
2013年1月6日 & 鄧瑞強博士 & 以弗所書3:1-12 & \hyperref[sec:180]{奧秘揭開了} \\
2013年1月13日 & 陳偉光教授 &  & \hyperref[sec:181]{主都知道} \\
2013年1月13日 & 陳榮超牧師 &  & \hyperref[sec:182]{愛的洗禮} \\
2013年1月20日 & 伍渭文牧師 & 詩篇36 & \hyperref[sec:183]{公義如高山判斷如深淵} \\
2013年1月20日 & 陳曉東博士 & 瑪拉基書3:1-5 & \hyperref[sec:184]{期待基督的喜悅} \\
2013年1月27日 & 鄧瑞強博士(香港神學院) & 以弗所書3:1-12 & \hyperref[sec:185]{奧秘揭開了} \\
2013年1月27日 & 陳恩明牧師 & 路加福音4:14-21 & \hyperref[sec:186]{狂人的福音} \\
\hline
2013年2月3日 & 伍渭文牧師 & 詩篇36 & \hyperref[sec:187]{公義如高山判斷} \\
2013年2月10日 & 梁燕城博士 &  & \hyperref[sec:188]{義無反顧的獻身} \\
2013年2月17日 & 陳榮超牧師 &  & \hyperref[sec:189]{愛的洗禮} \\
2013年2月17日 & 林豪恩先生 & 申命記26:1-11 & \hyperref[sec:190]{說一個感恩的故事} \\
2013年2月24日 & 林豪恩先生 & 申命記26:1-11 & \hyperref[sec:191]{說一個感恩的故事} \\
\hline
2013年3月3日 & 伍渭文牧師 & 詩篇63 & \hyperref[sec:192]{在曠野的一天} \\
2013年3月10日 & 盧龍光牧師 &  & \hyperref[sec:193]{誰是那浪子} \\
2013年3月17日 & 林豪恩先生 & 約翰福音12:1-8 & \hyperref[sec:194]{伯大尼的筵席} \\
2013年3月17日 & 陳天祥博士 &  & \hyperref[sec:195]{真知道神救恩中的三} \\
2013年3月24日 & 張闊先生 & 腓立比書2:5-11 & \hyperref[sec:196]{金色的麥穗是下垂的} \\
2013年3月31日 & 伍渭文牧師 & 詩篇118 & \hyperref[sec:197]{我必不至死仍要存活} \\
\hline
2013年4月7日 & 鄭漢文博士 & 約翰福音20:19-31 & \hyperref[sec:198]{信主之後} \\
2013年4月14日 & 陳耀鵬牧師 & 約翰福音21:1-19 & \hyperref[sec:199]{以愛相繫} \\
2013年4月21日 & 鄧瑞強博士 & 啟示錄7:9-17 & \hyperref[sec:200]{凱旋的群眾} \\
2013年4月28日 & 許開明牧師 & 約翰福音13:31-35 & \hyperref[sec:201]{使我成為別人的祝福} \\
\hline
2013年5月19日 & 伍渭文牧師 &  & \hyperref[sec:202]{從水和聖靈生的} \\
2013年5月26日 & 張伯笠牧師 & 約翰福音16:12-13 & \hyperref[sec:203]{引我們進入真理的聖靈} \\
\hline
2013年6月2日 & 陳曉東博士 & 加拉太書1:1-12 & \hyperref[sec:204]{自由的滋味} \\
2013年6月9日 & 謝喜鵬先生 & 列王記下4:1-7 & \hyperref[sec:205]{主的恩典夠用} \\
2013年6月23日 & 伍渭文牧師 & 詩篇第22 & \hyperref[sec:206]{敬拜與讚美} \\
2013年6月30日 & 鄧瑞強博士 & 路加福音9:51-62 & \hyperref[sec:207]{溫柔與嚴厲} \\
\hline
2013年7月7日 & 賀志勇博士 & 路加福音10:1-12 & \hyperref[sec:208]{機不可失} \\
2013年7月14日 & 林豪恩先生 & 路加福音10:25-37 & \hyperref[sec:209]{好撒馬利亞人是如何鍊成的} \\
2013年7月21日 & 許開明牧師 & 路加福音10:38-42 & \hyperref[sec:210]{重整生命焦點} \\
2013年7月28日 & 鄧瑞強博士 & 路加福音11:1-4 & \hyperref[sec:211]{人生於世何所禱也} \\
\hline
2013年8月4日 & 林豪恩先生 & 路加福音12:13-21 & \hyperref[sec:212]{泰然安心何處求} \\
2013年8月11日 & 賀志勇博士 & 路加福音12:35-48 & \hyperref[sec:213]{警醒警醒再警醒} \\
2013年8月21日 & 鄭漢文博士 & 路加福音12:49-56 & \hyperref[sec:214]{在和平與紛爭的時代作信仰解讀} \\
2013年8月25日 & 許開明牧師 & 路加福音13:10-17 & \hyperref[sec:215]{使人生轉彎為直} \\
\hline
2013年9月1日 & 伍渭文牧師 & 詩篇112 & \hyperref[sec:216]{不怕兇惡的信息} \\
2013年9月9日 & 許開明牧師 & 路加福音14:25-35 & \hyperref[sec:217]{不斷思量的門生} \\
2013年9月15日 & 梁元生教授 & 路加福音15 & \hyperref[sec:218]{人生中之圓與缺} \\
2013年9月22日 & 林豪恩先生 & 路加福音16:1-13 & \hyperref[sec:219]{應使則使還是應使不使} \\
2013年9月29日 & 鄧瑞強博士 & 詩篇146 & \hyperref[sec:220]{為受屈的伸冤} \\
\hline
2013年10月6日 & 沈祖堯教授 &  & \hyperref[sec:221]{計劃命運與人生} \\
2013年10月13日 & 許開明牧師 &  & \hyperref[sec:222]{做個感恩快樂人} \\
2013年10月20日 & 伍渭文牧師 & 詩篇121:1-8 & \hyperref[sec:223]{我要向山舉目} \\
2013年10月27日 & 鄧瑞強博士 & 路加福音18:9-14 & \hyperref[sec:224]{這是甚麼態度} \\
\hline
2013年11月3日 & 彭順強博士 & 撒母耳記下11:2-5,14-17,12:1-7 & \hyperref[sec:225]{拿單為公義發聲} \\
2013年11月10日 & 林豪恩先生 & 路加福音20:27-38 & \hyperref[sec:226]{她是誰的妻} \\
2013年11月17日 & 許開明牧師 & 路加福音21:5-19 & \hyperref[sec:227]{苦難是玩意或愛意} \\
2013年11月24日 & 鄧瑞強博士 & 路加福音23:33-43 & \hyperref[sec:228]{天堂與地獄} \\
\hline
2013年12月15日 & 鄧瑞強博士 & 馬太福音11:2-11 & \hyperref[sec:229]{眾裡尋他千百度} \\
2013年12月29日 & 鄭漢文博士 & 馬太福音2:13-23 & \hyperref[sec:230]{生命因何活得更美好} \\
\hline \hline
2014年1月5日 & 許開明牧師 & 約翰福音1:1-18 & \hyperref[sec:231]{年有道與祂共渡} \\
2014年1月12日 & 賀志勇博士 &  & \hyperref[sec:232]{祈禱進階2天意從來高難問} \\
2014年1月19日 & 鄧瑞強博士 & 約翰福音1:35-39 & \hyperref[sec:233]{成長的階梯} \\
\hline
2014年2月2日 & 伍渭文牧師 &  & \hyperref[sec:234]{誰能住你的聖山} \\
2014年2月9日 & 許開明牧師 & 馬太福音5:13-20 & \hyperref[sec:235]{活出鹽與光的生命} \\
2014年2月23日 & 鄧瑞強博士 & 詩篇119:33-40 & \hyperref[sec:236]{心懷天道人自正} \\
\hline
2014年3月16日 & 鄧瑞強博士 & 創世記12:1-4 & \hyperref[sec:237]{天降大任} \\
2014年3月23日 & 伍渭文牧師 & 詩篇95 & \hyperref[sec:238]{容易迷失的歸途} \\
2014年3月30日 & 鄭漢文博士 & 約翰福音9:1-41;詩篇23:1-6 & \hyperref[sec:239]{信仰讓生命不再一樣} \\
\hline
2014年4月13日 & 許開明牧師 &  & \hyperref[sec:240]{作個有用的人} \\
2014年4月20日 & 伍渭文牧師 & 路加福音24:13-24;約翰福音20:21-22 & \hyperref[sec:241]{以馬忤斯晚餐的光影} \\
2014年4月27日 & 鄧瑞強博士 & 約翰福音20:19-23 & \hyperref[sec:242]{當世界變得越來越得人驚的時候} \\
\hline
2014年5月4日 & 許開明牧師 & 約翰福音21 & \hyperref[sec:243]{晨光第一線復活主再現} \\
2014年5月11日 & 伍渭文牧師 & 詩篇23:1-6 & \hyperref[sec:244]{我必不致缺乏} \\
2014年5月18日 & 賀志勇博士 & 約翰福音14:1-11 & \hyperref[sec:245]{生路一條} \\
2014年5月25日 & 鄧瑞強博士 & 約翰福音14:15-24 & \hyperref[sec:246]{基督在我們裡面} \\
\hline
2014年6月1日 & 伍渭文牧師 & 使徒行傳1:6-14 & \hyperref[sec:247]{澆灌前夕的等候} \\
2014年6月8日 & 許開明牧師 &  & \hyperref[sec:248]{領受靈氣平安伴隨} \\
2014年6月15日 & 林豪恩先生 & 創世記1:2-4 & \hyperref[sec:249]{美好} \\
2014年6月22日 & 陳劍小姐 &  & \hyperref[sec:250]{時間的旅人盼望之旅} \\
2014年6月29日 & 鄭漢文博士 & 羅馬書6:12-23 & \hyperref[sec:251]{讓愛與和平佔領人間世} \\
\hline
2014年7月6日 & 伍渭文牧師 & 詩篇145 & \hyperref[sec:252]{國度國土國民} \\
2014年7月13日 & 許開明牧師 & 馬太福音13:1-13,18-23 & \hyperref[sec:253]{你聽見嗎} \\
2014年7月27日 & 林豪恩先生 & 哥林多前書13:4-8 & \hyperref[sec:254]{一生之久的愛} \\
\hline
2014年8月3日 & 許開明牧師 & 馬太福音14:13-21 & \hyperref[sec:255]{耶穌頭腦簡單麼} \\
2014年8月10日 & 張美珍博士 & 列王記上19:9-18;詩篇85:8-13;羅馬書10:5-15;馬太福音14:22-33 & \hyperref[sec:256]{堅強與軟弱} \\
2014年8月17日 & 賀志勇博士 & 馬太福音15:21-28 & \hyperref[sec:257]{驕傲與偏見} \\
2014年8月24日 & 林豪恩先生 & 出埃及記1:8-2:10 & \hyperref[sec:258]{有不認識約瑟的新王起來} \\
2014年8月31日 & 松谷曄介牧師 &  & \hyperref[sec:259]{我們在東亞處境下怎麼背起我們的十字架} \\
\hline
2014年9月14日 & 伍渭文牧師 & 創世記50:15-21;37:18-27 & \hyperref[sec:260]{饒恕別人就是釋放一個囚犯} \\
2014年9月21日 & 許開明牧師 & 馬太福音20:1-16 & \hyperref[sec:261]{厚此薄彼豈有此理} \\
2014年9月28日 & 鄧瑞強博士 & 馬太福音21:23-27 & \hyperref[sec:262]{你仗著什麼權柄做這些事} \\
\hline
2014年10月2日 & 伍渭文牧師 & 詩篇80 & \hyperref[sec:263]{當枯萎葡萄吐新葉} \\
2014年10月12日 & 鄧瑞強博士 & 馬太福音22:1-14 & \hyperref[sec:264]{生命是一場筵席} \\
2014年10月19日 & 陳偉光教授 &  & \hyperref[sec:265]{這個上帝不夠棒} \\
\hline
2014年11月2日 & 許開明牧師 &  & \hyperref[sec:266]{講得好好行得離譜} \\
2014年11月9日 & 鄧瑞強博士 & 馬太福音25:1-13 & \hyperref[sec:267]{等待不是靜看流水落花} \\
2014年11月16日 & 伍渭文牧師 & 詩篇90 & \hyperref[sec:268]{數算自己的日子} \\
\hline
2014年12月7日 & 鄧瑞強博士 & 詩篇85 & \hyperref[sec:269]{公義和平安互相親嘴} \\
2014年12月21日 & 伍渭文牧師 & 路加福音1:26-38;46-48 & \hyperref[sec:270]{天使的宣告} \\
2014年12月28日 & 林豪恩先生 & 路加福音2:22-40 & \hyperref[sec:271]{報到} \\
\hline \hline
2015年1月4日 & 許開明牧師 &  & \hyperref[sec:272]{洞悉人生視野真實現實超現實的世界} \\
2015年1月18日 & 梁明浩牧師 & 約翰福音20:24-29 & \hyperref[sec:273]{多馬的認信} \\
2015年1月25日 & 鄧瑞強博士 & 詩篇６２ & \hyperref[sec:274]{堅石與空氣} \\
\hline
2015年2月1日 & 伍渭文牧師 & 馬可福音1:21-28;詩篇110 & \hyperref[sec:275]{權柄與權力} \\
2015年2月8日 & 梁元生教授 & 路加福音9:21-56;10:13-16 & \hyperref[sec:276]{往聖城路上的城與鄉} \\
2015年2月22日 & 鄧瑞強博士 & 馬可福音1:9-15 & \hyperref[sec:277]{你是我的愛子} \\
\hline
2015年3月1日 & 伍渭文牧師 & 詩篇二十二 & \hyperref[sec:278]{彌賽亞的命運交響曲} \\
2015年3月15日 & 鄧瑞強博士 & 約翰福音3:14-21 & \hyperref[sec:279]{行真理者必來就光} \\
2015年3月22日 & 許開明牧師 & 約翰福音12:20-36 & \hyperref[sec:280]{人生關鍵時刻時候到了} \\
2015年3月29日 & 賀志勇博士 & 馬可福音11:1-11,15-18 & \hyperref[sec:281]{騎驢進京殺機四伏} \\
\hline
2015年4月12日 & 邢福增教授 & 約翰福音20:19-31 & \hyperref[sec:282]{願你們平安} \\
2015年4月19日 & 鄧瑞強博士 & 詩篇4 & \hyperref[sec:283]{床前明月光} \\
2015年4月26日 & 戴浩輝牧師 & 詩篇23:1-6;使徒行傳4:5-12;約翰一書3:16-24;約翰福音10:11-18 & \hyperref[sec:284]{好牧人} \\
\hline
2015年5月10日 & 許開明牧師 &  & \hyperref[sec:285]{繫連妙結} \\
2015年5月17日 & 伍渭文牧師 & 提摩太後書4:9-22 & \hyperref[sec:286]{天上的信息地上的使者} \\
2015年5月31日 & 鄧瑞強博士 & 羅馬書8:12-17 & \hyperref[sec:287]{起來不是奴隸的你} \\
\hline
2015年6月7日 & 鄧瑞強博士 & 創世記3:8-13 & \hyperref[sec:288]{捉迷藏} \\
2015年6月28日 & 伍渭文牧師 & 馬可福音5:21-43 & \hyperref[sec:289]{我的女兒病重了} \\
\hline
2015年7月5日 & 賀志勇博士 & 馬可福音6:7-13,30 & \hyperref[sec:290]{我們被差遣} \\
2015年7月12日 & 張美珍博士 &  & \hyperref[sec:291]{上帝你還在嗎} \\
\hline
2015年8月2日 & 松谷曄介牧師 & 申命記6:4-9/詩篇131:1-3/羅馬福音10:14-17/馬太福音19:13-15 & \hyperref[sec:292]{教會沒有畢業典禮的語言學校} \\
2015年8月9日 & 許開明牧師 & 約翰福音6:35,41-51 & \hyperref[sec:293]{語出驚人} \\
2015年8月16日 & 關瑞文教授 & 約翰福音6:51-58 & \hyperref[sec:294]{吃肉喝血的群體} \\
2015年8月23日 & 林豪恩先生 & 約書亞記24:1-18;詩篇34;以弗所書6:10-20;約翰福音6:56-69 & \hyperref[sec:295]{信仰路上的四個重要時刻} \\
\hline
2015年9月6日 & 伍渭文牧師 & 馬可福音7:24-30 & \hyperref[sec:296]{從餐桌掉下來的餅碎} \\
2015年9月20日 & 鄧瑞強博士 & 馬可福音9:30-37 & \hyperref[sec:297]{你們在路上議論的是什麼} \\
2015年9月27日 & 蒲錦昌牧師 & 提摩太後書2:1-7 & \hyperref[sec:298]{三個屬靈成長的秘訣} \\
\hline
2015年10月11日 & 許開明牧師 & 馬可福音10:17-37 & \hyperref[sec:299]{上帝與金錢是誰在掌權} \\
2015年10月18日 & 林豪恩先生 & 馬可福音10:35-45 & \hyperref[sec:300]{前來與被召} \\
2015年10月25日 & 鄧瑞強博士 & 詩篇126 & \hyperref[sec:301]{我們好像做夢的人} \\
\hline
2015年11月1日 & 許開明牧師 &  & \hyperref[sec:302]{誡與命愛與怕} \\
2015年11月8日 & 鄭漢文博士 & 馬可福音12:38-44;列王記上17:8-16;詩篇146:1-10 & \hyperref[sec:303]{信仰中的有餘多餘與不足充足} \\
2015年11月15日 & 林豪恩宣教師 & 希伯來書10:11-14,19-25;馬可福音13:1-8 & \hyperref[sec:304]{那日子臨近} \\
2015年11月29日 & 鄧瑞強博士 & 路加福音21:25-36 & \hyperref[sec:305]{當挺身昂首} \\
\hline
2015年12月13日 & 鄧瑞強博士 & 腓立比書4:4-7 & \hyperref[sec:306]{應當一無掛慮} \\
2015年12月27日 & 許開明牧師 & 路加福音2:41-52 & \hyperref[sec:307]{成年耶穌也問題多多麼} \\
\hline \hline
2016年1月3日 & 林豪恩先生 & 耶利米書31:7-14;詩篇147:12-20;約翰福音1:10-18 & \hyperref[sec:308]{戀戀山城} \\
2016年1月24日 & 鄧瑞強博士 & 哥林多前書12:12-27 & \hyperref[sec:309]{教會論之下的人觀} \\
2016年1月31日 & 許開明牧師 & 路加福音4:14-30 & \hyperref[sec:310]{漠視真話封口滅口} \\
\hline
2016年2月7日 & 林豪恩先生 & 路加福音9:28-36 & \hyperref[sec:311]{堅離地道} \\
2016年2月28日 & 鄧瑞強博士 & 路加福音13:1-9 & \hyperref[sec:312]{人禍意外天譴} \\
\hline
2016年3月6日 & 鄧瑞強博士 & 路加福音15:11-32 & \hyperref[sec:313]{歸去來兮吾歸何處} \\
2016年3月13日 & 陳偉光教授 &  & \hyperref[sec:314]{何處覓蓬萊} \\
2016年3月27日 & 許開明牧師 & 路加福音24:1-12 & \hyperref[sec:315]{復活節信息生命的意想不到} \\
\hline
2016年4月3日 & 鄧瑞強博士 & 啟示錄1:4-8 & \hyperref[sec:316]{從死裡首先復活的耶穌} \\
2016年4月10日 & 許開明牧師 & 約翰福音21:1-19 & \hyperref[sec:317]{民以食為先} \\
2016年4月17日 & 伍渭文牧師 & 啟示錄7:9-17 & \hyperref[sec:318]{被殺的羔羊與掌管萬有的主} \\
\hline
2016年5月1日 & 許開明牧師 & 約翰福音14:21-31 & \hyperref[sec:319]{平安主賜勇往直前} \\
2016年5月8日 & 林豪恩先生 & 使徒行傳16:16-39 & \hyperref[sec:320]{監獄風雲} \\
2016年5月15日 & 鄧瑞強博士 & 羅馬書8:12-17 & \hyperref[sec:321]{誰是我的父親} \\
2016年5月22日 & 朱建忠宣教師 & 創世記32:22-32 & \hyperref[sec:322]{與神摔跤破碎然而蒙福} \\
\hline
2016年6月5日 & 鄧瑞強博士 & 詩篇30:1-12 & \hyperref[sec:323]{一生之久的恩典} \\
2016年6月12日 & 許開明牧師 & 路加福音7:36-50 & \hyperref[sec:324]{獨排眾議舉動觸目} \\
2016年6月19日 & 林豪恩先生 & 路加福音8:26-39 & \hyperref[sec:325]{或去或留} \\
2016年6月26日 & 梁元生教授 & 使徒行傳21:37-40,22:1-5;路加福音8:11-15 & \hyperref[sec:326]{尋根歷史的叩問} \\
\hline
2016年7月10日 & 鄭漢文博士 & 路加福音10:25-37 & \hyperref[sec:327]{生命同行的鄰舍} \\
2016年7月17日 & 林豪恩先生 & 路加福音10:38-42 & \hyperref[sec:328]{馬大世界馬利亞心} \\
2016年7月24日 & 蘇成溢牧師 & 何西阿書1:2-10;歌羅西書2:6-15;路加福音11:1-13 & \hyperref[sec:329]{求主教導我們禱告} \\
2016年7月31日 & 范晉豪座堂牧師 &  & \hyperref[sec:330]{從差異與分歧到合一與共融} \\
\hline
2016年8月7日 & 賴品超教授 & 彼得前書2:9-17 & \hyperref[sec:331]{異鄉人在故鄉} \\
2016年8月14日 & 鄧瑞強博士 & 路加福音12:49-59 & \hyperref[sec:332]{看得見的抗爭看不見的和平} \\
2016年8月21日 & 林豪恩先生 & 路加福音13:10-17 & \hyperref[sec:333]{皆大歡喜} \\
\hline
2016年9月4日 & 林豪恩先生 & 詩篇1 & \hyperref[sec:334]{潤物細無聲} \\
2016年9月18日 & 鄧瑞強博士 & 阿摩司書8:1-12 & \hyperref[sec:335]{人飢餓非因無餅} \\
2016年9月25日 & 許開明牧師 &  & \hyperref[sec:336]{不做壞事無做好事} \\
\hline
2016年10月2日 & 鄧瑞強博士 & 提摩太後書1:1-14 & \hyperref[sec:337]{天行健君子以自強不息} \\
2016年10月9日 & 謝任生牧師 &  & \hyperref[sec:338]{後現代性與基督教信仰} \\
2016年10月16日 & 林崇智牧師 &  & \hyperref[sec:339]{宣教的行動} \\
2016年10月23日 & 邢福增教授 & 提摩太後書4:6-8 & \hyperref[sec:340]{終極勝利} \\
\hline
2016年11月6日 & 鄧瑞強博士 & 路加福音20:27-28 & \hyperref[sec:341]{有今生有來世} \\
2016年11月13日 & 許開明牧師 & 路加福音21:5-19 & \hyperref[sec:342]{末世凶兆如何自處} \\
2016年11月27日 & 林豪恩先生 & 以賽亞書2:1-5;詩篇122:1-9;羅馬書13:11-14;馬太福音24:36-44 & \hyperref[sec:343]{化干戈為玉帛} \\
\hline
2016年12月4日 & 鄧瑞強博士 & 以賽亞書11:1-10 & \hyperref[sec:344]{我們等待的那個新世界} \\
2016年12月11日 & 許開明牧師 &  & \hyperref[sec:345]{自覺我是誰} \\
2016年12月18日 & 許立中先生 &  & \hyperref[sec:346]{聖誕默想} \\
2016年12月26日 & 林豪恩先生 & 以賽亞書52:7-10;詩篇98;希伯來書1:1-4;約翰福音1:1-14 & \hyperref[sec:347]{葉公好龍話聖誕} \\
\hline \hline
2017年1月1日 & 鄧瑞強博士 & 希伯來書2:10-18 & \hyperref[sec:348]{出死入生的兄弟} \\
2017年1月15日 & 許開明牧師 & 約翰福音1:19-51 & \hyperref[sec:349]{你看見甚麼視覺與人生的探索} \\
2017年1月22日 & 李均熊博士 & 馬太福音4:12-23;以賽亞書9:1-4 & \hyperref[sec:350]{坐在黑暗裡的百姓看見了大光} \\
2017年1月29日 & 林豪恩先生 & 彌迦書6:1-8 & \hyperref[sec:351]{何以對蒼天} \\
\hline
2017年2月5日 & 許開明牧師 &  & \hyperref[sec:352]{活出麗亮人生} \\
2017年2月12日 & 鄧瑞強博士 & 馬太福音5:21-26 & \hyperref[sec:353]{君子務本本立而道生} \\
2017年2月19日 & 梁元生院長 & 利未記19:1-2,9-18;詩篇119:33-40;哥林多前書3:10-11,16-23;馬太福音5:38-48 & \hyperref[sec:354]{說之間} \\
2017年2月26日 & 戴浩輝牧師 & 馬太福音17:1-9;出埃及記24:12-18;詩篇2;彼德後書1:16-21 & \hyperref[sec:355]{靈性高峰} \\
\hline
2017年3月5日 & 許開明牧師 & 馬可福音1:12-13;路加福音4:1-13 & \hyperref[sec:356]{正邪交鋒愈戰愈勇風起雲湧} \\
2017年3月12日 & 林豪恩先生 & 撒母耳記上9 & \hyperref[sec:357]{尋找的就尋見} \\
2017年3月19日 & 鄭漢文博士 & 羅馬書5:1-11 & \hyperref[sec:358]{因信稱義的因恩欣} \\
2017年3月26日 & 何崇謙牧師 & 創世記2:1-17 & \hyperref[sec:359]{我對你錯窒住盲堵的陰霾} \\
\hline
2017年4月2日 & 林豪恩先生 & 以西結書37:1-14;詩篇130:1-8;羅馬書8:6-11;約翰福音11:1-45 & \hyperref[sec:360]{我們還剩下甚麼} \\
2017年4月9日 & 陳天祥博士 &  & \hyperref[sec:361]{住在主裡的重要性} \\
2017年4月16日 & 許開明牧師 & 約翰福音20:1-20 & \hyperref[sec:362]{復活的震撼生命存在於死亡之中} \\
2017年4月23日 & 朱建忠先生 & 以弗所書4:1 & \hyperref[sec:363]{既然蒙召} \\
2017年4月30日 & 伍渭文牧師 & 使徒行傳2:14,36-41;詩篇116:1-4,12-19;彼得前書1:17-23;路加福音24:13-35 & \hyperref[sec:364]{凡有血氣盡都如草} \\
\hline
2017年5月7日 & 林豪恩先生 & 詩篇23:1-6 & \hyperref[sec:365]{追追趕趕高高低低} \\
2017年5月14日 & 許開明牧師 & 約翰福音14:1-14 & \hyperref[sec:366]{時日不多了} \\
2017年5月21日 & 許立中先生 & 使徒行傳17:22-31 & \hyperref[sec:367]{當耶路撒冷遇上雅典} \\
2017年5月28日 & 關瑞文教授 &  & \hyperref[sec:368]{尾場的扣問} \\
\hline
2017年6月4日 & 許開明牧師 & 約翰福音20:19-29 & \hyperref[sec:369]{願你們平安} \\
2017年6月11日 & 蒲錦昌牧師 & 創世記1:1-2:4;詩篇8:1-9;哥林多後書13:11-13;馬太福音28:16-20 & \hyperref[sec:370]{上帝愛世界} \\
2017年6月18日 & 邢福增院長 & 羅馬書5:1-8 & \hyperref[sec:371]{因信稱義} \\
2017年6月25日 & 王福義長老 & 創世記9:8-17 & \hyperref[sec:372]{首屆地球峰會} \\
\hline
2017年7月2日 & 林豪恩先生 & 耶利米書28:5-9 & \hyperref[sec:373]{忠誠的異見者} \\
2017年7月9日 & 許開明牧師 & 創世記12:1-20 & \hyperref[sec:374]{惜別講章與神同行祝福人} \\
2017年7月30日 & 鄧瑞強博士 & 羅馬書8:26-39 & \hyperref[sec:375]{在泥淖中起舞} \\
\hline
2017年8月6日 & 謝任生牧師 & 路加福音1:28-55 & \hyperref[sec:376]{上帝的逆道與人間的常道} \\
2017年8月13日 & 林豪恩先生 & 列王記上19:9-18 & \hyperref[sec:377]{哀莫大於心死} \\
2017年8月20日 & 鄧瑞強博士 & 馬太福音15:21-28 & \hyperref[sec:378]{不要奉旨} \\
\hline
2017年9月10日 & 鄧瑞強博士 & 啟示錄2:1-7 & \hyperref[sec:379]{毋忘初心} \\
2017年9月13日 & 伍渭文牧師 & 哥林多前書11:23-24 & \hyperref[sec:380]{吃甚麼成了甚麼} \\
2017年9月17日 & 林豪恩先生 & 創世記50:15-21;詩篇103;羅馬書14:1-12;馬太福音18:21-35 & \hyperref[sec:381]{淚眼泯恩仇} \\
2017年9月24日 & 許立中先生 & 馬太福音11:16-19,25-30 & \hyperref[sec:382]{信仰的重軛} \\
\hline
2017年10月1日 & 白德培牧師 & 腓立比書2:1-13 & \hyperref[sec:383]{基督無能為力嗎} \\
2017年10月8日 & 鄧瑞強博士 & 啟示錄3:14-22 & \hyperref[sec:384]{富裕的貧窮} \\
2017年10月22日 & 關瑞文教授 & 以賽亞書45:1-7 & \hyperref[sec:385]{我施平安又降災禍} \\
2017年10月29日 & 賴品超教授 & 啟示錄3:1-6 & \hyperref[sec:386]{不斷改革的教會} \\
\hline
2017年11月5日 & 蒲錦昌牧師 & 哥林多前書10:17;以弗所書4:1-6 & \hyperref[sec:387]{從衝突到共融} \\
2017年11月12日 & 鄧瑞強博士 & 啟示錄4:1-11 & \hyperref[sec:388]{真相寶座敬拜} \\
2017年11月18日 & 鄭漢文博士 & 馬太福音25:14-30 & \hyperref[sec:389]{財富才幹與在地若天} \\
2017年11月26日 & 梁元生院長 & 馬太福音16 & \hyperref[sec:390]{我是誰} \\
\hline
2017年12月3日 & 伍渭文牧師 & 提摩太前書3:16;哥林多前書11:23-29 & \hyperref[sec:391]{無人不以為然的奧秘} \\
2017年12月10日 & 鄧瑞強博士 & 以賽亞書40:1-11 & \hyperref[sec:392]{花褪殘紅青杏小} \\
2017年12月17日 & 林豪恩先生 & 約翰福音1:6-8,19-28 & \hyperref[sec:393]{燕子飛時綠水人家繞} \\
2017年12月24日 & 許立中先生 & 提摩太前書3:16 & \hyperref[sec:394]{可道與不可道} \\
2017年12月31日 & 朱建忠博士 & 雅各書1:1-12 & \hyperref[sec:395]{以試煉為喜樂} \\
\hline \hline
2018年1月7日 & 鄧瑞強博士 & 啟示錄5:1-14 & \hyperref[sec:396]{誰配揭開歷史之謎} \\
2018年1月21日 & 林豪恩先生 & 馬太福音22:1-14 & \hyperref[sec:397]{喜事起事} \\
2018年1月28日 & 林豪恩先生 & 哥林多前書8:1-13 & \hyperref[sec:398]{高端信徒與低端信徒} \\
\hline
2018年2月4日 & 伍渭文牧師 & 約拿書3:1-5,10;詩篇62:5-12;哥林多前書7:29-31;馬可福音1:14-20 & \hyperref[sec:399]{不能承受生命的速} \\
2018年2月11日 & 鄧瑞強博士 & 啟示錄6:1-17 & \hyperref[sec:400]{預你一份} \\
2018年2月18日 & 鄭漢文博士 & 馬可福音1:9-15 & \hyperref[sec:401]{悔改靈洗信福音} \\
2018年2月25日 & 邢福增院長 & 馬可福音8:31-38;創世記17:1-7,15-16;羅馬書4:13-25 & \hyperref[sec:402]{捨與得} \\
\hline
2018年3月4日 & 許開明牧師 & 約翰福音2:13-22 & \hyperref[sec:403]{耶穌都有火} \\
2018年3月11日 & 鄧瑞強博士 & 啟示錄7:1-17 & \hyperref[sec:404]{數不盡的精兵} \\
2018年3月25日 & 李均熊博士 &  & \hyperref[sec:405]{棕枝主日的驢工} \\
\hline
2018年4月1日 & 戴浩輝牧師 &  & \hyperref[sec:406]{祂在我們以先} \\
2018年4月8日 & 鄧瑞強博士 & 啟示錄8:1-6 & \hyperref[sec:407]{天上靜然無聲半小時} \\
2018年4月29日 & 林豪恩先生 & 使徒行傳8:26-40 & \hyperref[sec:408]{向南走向北走} \\
\hline
2018年5月6日 & 伍渭文牧師 & 使徒行傳10:44-48;詩篇98:1-9;約翰一書5:1-6;約翰福音15:9-17 & \hyperref[sec:409]{流淌不住的水像神聖的洗禮} \\
2018年5月13日 & 鄧瑞強博士 & 啟示錄11:1-19 & \hyperref[sec:410]{重述福音} \\
2018年5月27日 & 戴國璋牧師 & 哥羅西書4:10-14;提摩太後書4:9-11 & \hyperref[sec:411]{保羅身旁的三個夥伴} \\
\hline
2018年6月3日 & 吳碧珊牧師 & 撒母耳記上3:1-10 & \hyperref[sec:412]{事奉耶和華} \\
2018年6月10日 & 鄧瑞強博士 & 啟示錄12:1-17 & \hyperref[sec:413]{屋頂上的提琴手} \\
2018年6月17日 & 林豪恩先生 & 以西結書17:22-24,31:2-11,37:5-10 & \hyperref[sec:414]{香柏與芥菜} \\
\hline
2018年7月1日 & 王家輝牧師 & 哥林多後書8:7-15 & \hyperref[sec:415]{返璞歸真} \\
2018年7月8日 & 戴國璋牧師 & 馬太福音6:1-4 & \hyperref[sec:416]{耶穌論施捨} \\
2018年7月15日 & 梁元生院長 & 列王記上18:17-22;馬可福音6:14-29 & \hyperref[sec:417]{從迦密山回望世界} \\
2018年7月29日 & 林豪恩先生 & 列王記下4:42-44;詩篇145:10-18;以弗所書3:14-21;約翰福音6:1-15 & \hyperref[sec:418]{不足夠} \\
\hline
2018年8月5日 & 戴國璋牧師 & 以弗所書5:1-7 & \hyperref[sec:419]{作上帝蒙慈愛的兒女} \\
2018年8月12日 & 周宇牧師 & 腓立比書3:7-14 & \hyperref[sec:420]{使徒保羅的使命人生} \\
2018年8月19日 & 林豪恩先生 & 箴言9:1-6;詩篇34:9-14;以弗所書5:15-20;約翰福音6:51-58 & \hyperref[sec:421]{生活} \\
2018年8月26日 & 鄭漢文博士 &  & \hyperref[sec:422]{難為正邪定分界也得裝身作靈戰} \\
\hline
2018年9月2日 & 賴品超牧師 & 箴言3:1-12 & \hyperref[sec:423]{智慧的傳承} \\
2018年9月9日 & 劉國偉先生 & 約翰福音11:1-16 & \hyperref[sec:424]{來跟從我} \\
2018年9月30日 & 許立中先生 & 彼得前書4:1-6 & \hyperref[sec:425]{異鄉人} \\
\hline
2018年10月7日 & 吳炳華牧師 & 詩篇42,43 & \hyperref[sec:426]{給我一個笑臉} \\
2018年10月14日 & 鄭漢文博士 & 馬可福音10:17-22 & \hyperref[sec:427]{為乜你可以去到幾盡} \\
2018年10月21日 & 鄧瑞強博士 & 啟示錄13:1-18 & \hyperref[sec:428]{立於亂世之中} \\
2018年10月28日 & 林豪恩先生 & 耶利米書31:7-9 & \hyperref[sec:429]{復興之道} \\
\hline
2018年11月4日 & 伍渭文牧師 & 以賽亞書25:6-9;詩篇24:1-10;啟示錄21:1-6;約翰福音11:32-44 & \hyperref[sec:430]{神要擦去他們一切的眼淚} \\
2018年11月11日 & 鄧瑞強博士 & 啟示錄14:1-13 & \hyperref[sec:431]{音樂起了} \\
2018年11月18日 & 劉國偉先生 & 約翰福音13:1-15 & \hyperref[sec:432]{主愛我們到底} \\
2018年11月25日 & 吳梓明教授 & 約翰一書4:7-8,12,16 & \hyperref[sec:433]{轉化心靈的教育} \\
\hline
2018年12月2日 & 袁天佑牧師 & 耶利米書33:14-16;帖撒羅尼迦前書3:9-13;路加福音21:25-36 & \hyperref[sec:434]{無力的時代有力的信仰} \\
2018年12月9日 & 林豪恩先生 & 瑪拉基書3 & \hyperref[sec:435]{傳而有信} \\
2018年12月24日 & 鄧瑞強博士 & 馬太福音2:2-12 & \hyperref[sec:436]{我們多麼需要星星} \\
2018年12月30日 & 許立中先生 & 路加福音2:41-52 & \hyperref[sec:437]{人之子} \\
\hline \hline
2019年1月6日 & 白德培牧師 & 馬太福音2:1-12 & \hyperref[sec:438]{當東方智慧遇上基督} \\
2019年1月13日 & 劉國偉先生 & 約翰福音11:17-44 & \hyperref[sec:439]{死亡眼淚與生命} \\
2019年1月20日 & 鄧瑞強博士 & 啟示錄15:1-4 & \hyperref[sec:440]{來唱歌吧} \\
2019年1月27日 & 林豪恩先生 &  & \hyperref[sec:441]{講定唔講做得到先好講} \\
\hline
2019年2月3日 & 戴浩輝牧師 & 路加福音4:21-30 & \hyperref[sec:442]{拆毀與建立} \\
2019年2月10日 & 劉國偉先生 & 約翰福音11:45-57,12:1-11 & \hyperref[sec:443]{你們甚麼都不懂我們懂得甚麼} \\
2019年2月17日 & 鄧瑞強博士 & 啟示錄18:1-20 & \hyperref[sec:444]{不義而富且貴於我如浮雲} \\
2019年2月24日 & 區祥江博士 & 創世記45:1-15 & \hyperref[sec:445]{饒恕的中途站} \\
\hline
2019年3月3日 & 伍渭文牧師 & 路加福音9:28-31 & \hyperref[sec:446]{幽黑愈下瞥見榮光} \\
2019年3月10日 & 劉國偉先生 & 約翰福音13:1-2,12-30 & \hyperref[sec:447]{耶穌為猶大洗腳} \\
2019年3月17日 & 鄧瑞強博士 & 啟示錄21:1-27 & \hyperref[sec:448]{上帝的家在人間了} \\
2019年3月24日 & 林豪恩先生 & 路加福音13:1-9 & \hyperref[sec:449]{五十步笑百步} \\
\hline
2019年4月14日 & 鄭漢文博士 &  & \hyperref[sec:450]{以何心為心為何所用} \\
2019年4月21日 & 鄧瑞強博士 & 路加福音24:1-12 & \hyperref[sec:451]{為甚麼在死人中找活人呢} \\
2019年4月28日 & 劉國偉先生 & 約翰福音12:12-50 & \hyperref[sec:452]{人子得榮耀的時候到了} \\
\hline
2019年5月12日 & 劉國偉先生 & 約翰福音12:12-50 & \hyperref[sec:453]{雞叫以先} \\
2019年5月19日 & 鄧瑞強博士 & 約翰福音13:31-35 & \hyperref[sec:454]{在黑暗與榮耀之間} \\
2019年5月26日 & 林豪恩先生 & 使徒行傳16:6-15 & \hyperref[sec:455]{何處是吾家} \\
\hline
2019年6月9日 & 林豪恩先生 & 使徒行傳2:1-21 & \hyperref[sec:456]{忽然從天上有響聲下來} \\
2019年6月16日 & 鄧瑞強博士 & 約翰福音16:12-15 & \hyperref[sec:457]{真理在途上} \\
2019年6月23日 & 劉國偉先生 & 約翰福音18:1-14 & \hyperref[sec:458]{你們找誰} \\
2019年6月30日 & 朱光華牧師 & 詩篇78:1- & \hyperref[sec:459]{陶造新世代} \\
\hline
2019年7月7日 & 劉國偉先生 & 約翰福音18:15-27 & \hyperref[sec:460]{你為甚麼問我呢} \\
2019年7月14日 & 林豪恩先生 & 路加福音10:25-37 & \hyperref[sec:461]{智者出口聖者出走仁者出手} \\
2019年7月21日 & 許立中先生 & 路加福音10:38-42 & \hyperref[sec:462]{不可少的只有一件} \\
2019年7月28日 & 劉國偉先生 & 約翰福音18:28-38 & \hyperref[sec:463]{真理是甚麼呢} \\
\hline
2019年8月4日 & 伍渭文牧師 & 以弗所書5:26-27,31-32 & \hyperref[sec:464]{作為基督新婦的教會} \\
2019年8月11日 & 劉國偉先生 & 約翰福音19:1-16 & \hyperref[sec:465]{釘死你們的王} \\
2019年8月18日 & 林豪恩先生 & 路加福音12:49-56 & \hyperref[sec:466]{耶穌係呢度} \\
2019年8月25日 & 區祥江博士 & 路加福音4:16-19 & \hyperref[sec:467]{合理解綁帶來神的榮耀} \\
\hline
2019年9月8日 & 劉國偉先生 & 約翰福音19:17-27 & \hyperref[sec:468]{就在這裡把祂釘十字架} \\
2019年9月15日 & 鄧瑞強博士 & 出埃及記32:1-14 & \hyperref[sec:469]{他們如此容易便墮落了} \\
2019年9月22日 & 林豪恩先生 & 阿摩司書8:4-7;詩篇113;提摩太前書2:1-7;路加福音16:1-13 & \hyperref[sec:470]{資源位置與社群共融} \\
2019年9月29日 & 鄭漢文博士 &  & \hyperref[sec:471]{前輩教後輩打美好的仗} \\
\hline
2019年10月6日 & 袁天佑牧師 & 路加福音17:5-10 & \hyperref[sec:472]{我們是無用的僕人} \\
2019年10月13日 & 劉國偉先生 & 約翰福音20:1-18 & \hyperref[sec:473]{我已經看見了主} \\
2019年10月20日 & 鄧瑞強博士 & 創世記32:6-33:11 & \hyperref[sec:474]{渡口} \\
2019年10月27日 & 林豪恩先生 & 耶利米書14:4-22 & \hyperref[sec:475]{這麼近那麼遠} \\
\hline
2019年11月3日 & 戴浩輝牧師 & 路加福音19:1-10 & \hyperref[sec:476]{今天救恩到了這家} \\
2019年11月10日 & 劉國偉先生 & 約翰福音20:19-29 & \hyperref[sec:477]{那沒有看見就信的有福了} \\
2019年11月24日 & 林豪恩先生 & 耶利米書23:3-6 & \hyperref[sec:478]{從相分之痛到相聚之盼} \\
\hline
2019年12月1日 & 陳衍昌法政牧師 & 哥林多前書13:1-13 & \hyperref[sec:479]{超越的愛} \\
2019年12月8日 & 劉國偉先生 & 約翰福音21:1-19 & \hyperref[sec:480]{跟從我吧} \\
2019年12月15日 & 鄧瑞強博士 & 馬太福音11:2-11 & \hyperref[sec:481]{那將要來的是你嗎} \\
2019年12月22日 & 許立中先生 & 以賽亞書40:27-31 & \hyperref[sec:482]{等候神} \\
2019年12月24日 & 劉國偉先生 & 路加福音2:8-14 & \hyperref[sec:483]{那就是記號了} \\
2019年12月29日 & 林豪恩先生 &  & \hyperref[sec:484]{寒夜數星星} \\
\hline \hline
2020年1月5日 & 蘇成溢牧師 & 女人弗所書1:1-14 & \hyperref[sec:485]{在基督裡的多重福氣} \\
2020年1月19日 & 鄧瑞強博士 & 詩篇40 & \hyperref[sec:486]{神啊求你不要耽延} \\
2020年1月26日 & 鄭漢文博士 & 哥林多前書1:10-18 & \hyperref[sec:487]{不可分黨一心相合} \\
\hline
2020年2月2日 & 光華牧師 & 詩篇15:1-5 & \hyperref[sec:488]{聖所門口的要理問答} \\
2020年2月9日 & 劉國偉先生 &  & \hyperref[sec:489]{願榮光歸上帝} \\
2020年2月16日 & 鄧瑞強博士 & 申命記30:15-20 & \hyperref[sec:490]{這樣的因這樣的果} \\
2020年2月23日 & 林豪恩先生 & 馬太福音17:1-9 & \hyperref[sec:491]{遇見神} \\
\hline
2020年3月1日 & 伍渭文牧師 & 詩篇32 & \hyperref[sec:492]{塗抹的手才能寫出真實} \\
2020年3月8日 & 劉國偉先生 & 詩篇121:1-8 & \hyperref[sec:493]{本於信以至於信} \\
2020年3月15日 & 鄧瑞強博士 & 出埃及記17:1-7 & \hyperref[sec:494]{上帝是不是在我們中間呢} \\
2020年3月22日 & 關瑞文教授 & 約翰福音9 & \hyperref[sec:495]{失明的才能看見} \\
2020年3月29日 & 鄭漢文博士 & 羅馬書8:6-11 & \hyperref[sec:496]{體因罪而死心靈因義而活} \\
\hline
2020年4月12日 & 劉國偉先生 &  & \hyperref[sec:497]{基督是我們的生命} \\
2020年4月18日 & 鄧瑞強博士 & 約翰福音20:19-31 & \hyperref[sec:498]{困局中的突破} \\
2020年4月26日 & 林豪恩先生 & 使徒行傳2:14,36-41 & \hyperref[sec:499]{踏上這無盡旅途} \\
\hline
2020年5月10日 & 劉國偉先生 &  & \hyperref[sec:500]{更大的事} \\
2020年5月17日 & 鄧瑞強博士 & 約翰福音14:15-21 & \hyperref[sec:501]{古來聖賢皆寂寞} \\
2020年5月24日 & 豪恩先生 & 使徒行傳1:6-14;詩篇68:1-10,32-35;彼得前書4:12-14,5:6-11;約翰福音17:1-11 & \hyperref[sec:502]{作我的見證} \\
2020年5月31日 & 鄭漢文博士 & 使徒行傳2:1-21 & \hyperref[sec:503]{靈動我心力行奇事言說基督見證上帝} \\
\hline
2020年6月7日 & 王家輝牧師 & 馬太福音28:16-20 & \hyperref[sec:504]{奉命行事} \\
2020年6月14日 & 劉國偉先生 & 出埃及記19:2-8 & \hyperref[sec:505]{背在鷹的翅膀上} \\
2020年6月21日 & 邢福增教授 & 耶利米書20:7-13;羅馬書6:1-11;馬太福音10:24-39 & \hyperref[sec:506]{生命中不能承受的輕} \\
2020年6月28日 & 鄭漢文博士 & 馬太福音10:40-42 & \hyperref[sec:507]{惡待善待恩待等等接待的信仰意義} \\
\hline
2020年7月5日 & 袁天佑牧師 & 馬太福音11:16-19,25-30 & \hyperref[sec:508]{荒謬時代中信仰的實踐} \\
2020年7月12日 & 劉國偉先生 & 馬太福音13 & \hyperref[sec:509]{道在心田} \\
2020年7月19日 & 林豪恩先生 & 以賽亞書44:6-8 & \hyperref[sec:510]{生命中能承受的重} \\
2020年7月26日 & 劉國偉先生 & 馬太福音13 & \hyperref[sec:511]{誰能使我們與基督的愛隔絕呢} \\
\hline
2020年8月2日 & 高國雄牧師 & 馬太福音14:13-21 & \hyperref[sec:512]{我們這裡只有五個餅兩條魚} \\
2020年8月9日 & 張美珍博士 &  & \hyperref[sec:513]{只剩下我一個人} \\
2020年8月16日 & 許立中先生 & 馬可福音3:13-19 & \hyperref[sec:514]{教會這個群體} \\
2020年8月23日 & 林豪恩先生 & 以賽亞書51:1-6 & \hyperref[sec:515]{給求義尋道者之言} \\
\hline
2020年9月6日 & 高國雄牧師 &  & \hyperref[sec:516]{最大的權柄復和} \\
2020年9月13日 & 林豪恩先生 & 詩篇103:8-13 & \hyperref[sec:517]{痛與怒愛與恕} \\
2020年9月20日 & 鄧瑞強博士 & 腓立比書1:20-30 & \hyperref[sec:518]{在承擔責任的生命中活出永生} \\
2020年9月27日 & 伍渭文牧師 & 馬太福音21:23-32;以西結書18:1-4,25-32;詩篇25:1-9;腓立比書2:1-13 & \hyperref[sec:519]{你仗著甚麼權柄} \\
\hline
2020年10月4日 & 高國雄牧師 & 馬太福音21:33-46 & \hyperref[sec:520]{我們看得到這希奇嗎} \\
2020年10月11日 & 林豪恩先生 & 以賽亞書25:1-9;詩篇23:1-6;腓立比書4:1-9;馬太福音22:1-14 & \hyperref[sec:521]{要飯聚不犯罪} \\
2020年10月18日 & 鄧瑞強博士 & 以賽亞書45:1-7 & \hyperref[sec:522]{神造光又造暗} \\
2020年10月25日 & 梁遠耀牧師 & 馬太福音22:34-46;利未記19:1-2,15-18;詩篇1:1-6;帖撒羅尼迦前書2:1-8 & \hyperref[sec:523]{請問} \\
\hline
2020年11月1日 & 高國雄牧師 & 馬太福音23:1-12 & \hyperref[sec:524]{信仰只說不行行嗎} \\
2020年11月8日 & 林豪恩先生 & 阿摩司書5:18-24;詩篇70:1-5;帖撒羅尼迦前書4:13-18;馬太福音25:1-13 & \hyperref[sec:525]{那日子} \\
2020年11月15日 & 鄧瑞強博士 & 詩篇90 & \hyperref[sec:526]{主啊你世世代代作我們的避難所} \\
2020年11月22日 & 王玉教授 & 馬太福音25:31-46;以西結書34:11-16,20-24 & \hyperref[sec:527]{切膚之痛} \\
2020年11月29日 & 伍渭文牧師 & 以賽亞書64:1-9;詩篇80:1-7,17-19;哥林多前書1:3-9;馬可福音13:24-37 & \hyperref[sec:528]{萬物朝終} \\
\hline
2020年12月6日 & 高國雄牧師 & 馬可福音1:1-8;以賽亞書40:1-11;詩篇85:1-2,8-13;彼得後書3:8-15 & \hyperref[sec:529]{各就各位預備} \\
2020年12月13日 & 王家輝牧師 & 以賽亞書61:1-4,8,11 & \hyperref[sec:530]{終極好消息} \\
2020年12月20日 & 鄧瑞強博士 & 路加福音1:46-55 & \hyperref[sec:531]{祂顧念窮人的卑微} \\
2020年12月27日 & 林豪恩先生 & 路加福音2:22-40;加拉太書4:4-7 & \hyperref[sec:532]{不凡的竟然如此平凡} \\
\hline \hline
2021年1月3日 & 高國雄牧師 & 約翰福音1:1-18 & \hyperref[sec:533]{道成了肉身} \\
2021年1月17日 & 鄧瑞強博士 & 約翰福音1:43-51 & \hyperref[sec:534]{你們將要看見天開了} \\
2021年1月24日 & 張亦農牧師 & 馬可福音1:14-20 & \hyperref[sec:535]{悔改相信} \\
2021年1月31日 & 林豪恩先生 &  & \hyperref[sec:536]{先知文士耶穌} \\
\hline
2021年2月7日 & 高國雄牧師 & 馬可福音1:29-39 & \hyperref[sec:537]{眾人都在找你} \\
2021年2月14日 & 伍渭文牧師 & 馬可福音9:2-9 & \hyperref[sec:538]{登山顯容} \\
2021年2月21日 & 鄧瑞強博士 & 馬可福音1:9-15 & \hyperref[sec:539]{身分考驗信息} \\
2021年2月28日 & 林豪恩先生 &  & \hyperref[sec:540]{我們與魔的距離} \\
\hline
2021年3月14日 & 高國雄牧師 & 約翰福音3:14-21 & \hyperref[sec:541]{來就光} \\
2021年3月21日 & 鄧瑞強博士 & 希伯來書5:5-10 & \hyperref[sec:542]{因所受的苦難學了順從} \\
2021年3月28日 & 鄧達榮牧師 &  & \hyperref[sec:543]{榮耀之主謙卑服侍} \\
\hline
2021年4月4日 & 高國雄牧師 & 馬可福音16:1-8 & \hyperref[sec:544]{他已經復活了不在這裡} \\
2021年4月11日 & 伍渭文牧師 &  & \hyperref[sec:545]{我主我神一位誠實存疑者的認信} \\
2021年4月18日 & 鄧瑞強博士 & 路加福音24:36-43 & \hyperref[sec:546]{你們看我的手我的腳} \\
2021年4月25日 & 林豪恩先生 & 使徒行傳4:5-13 & \hyperref[sec:547]{存在的勇氣} \\
\hline
2021年5月2日 & 王家輝牧師 & 使徒行傳8:26-40 & \hyperref[sec:548]{問道門徒} \\
2021年5月9日 & 高國雄牧師 &  & \hyperref[sec:549]{我愛} \\
2021年5月16日 & 鄧瑞強博士 & 約翰福音17:13-19 & \hyperref[sec:550]{沒有幻想不用絕望} \\
2021年5月23日 & 林豪恩先生 & 使徒行傳2:1-21 & \hyperref[sec:551]{從五旬節到聖靈降臨} \\
2021年5月30日 & 蒲錦昌牧師 & 以賽亞書6:1-8;詩篇29:1-11;羅馬書8:12-17;約翰福音3:1-17 & \hyperref[sec:552]{上主的臨在} \\
\hline
2021年6月6日 & 高國雄牧師 & 馬可福音3:20-35 & \hyperref[sec:553]{癲狂的耶穌} \\
2021年6月13日 & 馮天聰牧師 & 馬可福音4:26-34 & \hyperref[sec:554]{一切從一粒種子開始} \\
2021年6月20日 & 鄧瑞強博士 & 馬可福音4:35-41 & \hyperref[sec:555]{為何耶穌可以睡得這般香甜} \\
2021年6月27日 & 林豪恩先生 & 耶利米哀歌3:22-33 & \hyperref[sec:556]{靜默有時言語有時} \\
\hline
2021年7月4日 & 高國雄牧師 & 馬可福音6:1-13 & \hyperref[sec:557]{平凡信徒肩負不平凡任務} \\
2021年7月11日 & 林崇智牧師 &  & \hyperref[sec:558]{生命中不能承受的輕} \\
2021年7月18日 & 鄧瑞強博士 & 以弗所書2:11-22 & \hyperref[sec:559]{中間的牆沒有了} \\
2021年7月25日 & 林豪恩先生 & 約翰福音6 & \hyperref[sec:560]{耶穌比曹操更快} \\
\hline
2021年8月1日 & 高國雄牧師 & 約翰福音6:24-35 & \hyperref[sec:561]{鳥為食亡} \\
2021年8月8日 & 李均熊博士 & 約翰福音6 & \hyperref[sec:562]{我就是生命的糧} \\
2021年8月15日 & 林豪恩先生 & 箴言9:1-6 & \hyperref[sec:563]{人生的智慧} \\
2021年8月22日 & 張美珍博士 &  & \hyperref[sec:564]{你認識祂嗎} \\
2021年8月29日 & 關瑞文教授 & 馬可福音7:1-8,14-15,21-23 & \hyperref[sec:565]{不要讓傳統死去} \\
\hline
2021年9月5日 & 張亦農牧師 & 馬可福音7 & \hyperref[sec:566]{把孩子的餅給狗} \\
2021年9月12日 & 李均熊博士 & 馬可福音8:27-38 & \hyperref[sec:567]{你介意甚麼} \\
2021年9月19日 & 高國雄牧師 & 馬可福音9:30-37 & \hyperref[sec:568]{們在路上議論的是甚麼} \\
2021年9月26日 & 林豪恩先生 & 民數記11:4-6,10-16,24-29 & \hyperref[sec:569]{擔當不起} \\
\hline
2021年10月3日 & 高國雄牧師 & 馬可福音10:2-16 & \hyperref[sec:570]{人不可分開} \\
2021年10月10日 & 吳梓明教授 & 阿摩司書5:6-7;詩篇90:12-17;約翰一書4:7-12,18-19;馬可福音12:28-34 & \hyperref[sec:571]{近代中國三位傑出的教育家} \\
2021年10月17日 & 林豪恩先生 & 以賽亞書53:3-12 & \hyperref[sec:572]{人生或苦不亢不卑} \\
2021年10月31日 & 蘇成溢牧師 & 希伯來書9:11-14 & \hyperref[sec:573]{成聖的因與果} \\
\hline
2021年11月14日 & 林豪恩先生 & 馬太福音13:1-10 & \hyperref[sec:574]{活在山崩石裂時} \\
2021年11月21日 & 盧龍光牧師 & 約翰福音18:33-37 & \hyperref[sec:575]{你是猶太人的王嗎} \\
2021年11月28日 & 林津牧師 & 路加福音21:25-36 & \hyperref[sec:576]{積極的等待} \\
\hline
2021年12月12日 & 林豪恩先生 & 西番雅書3:14-20 & \hyperref[sec:577]{等一個人} \\
2021年12月19日 & 鄧瑞強博士 & 路加福音1:46-55 & \hyperref[sec:578]{絕望者的希望} \\
2021年12月24日 & 高國雄牧師 & 路加福音2:1-20 & \hyperref[sec:579]{這就是要給你們的記號} \\
2021年12月26日 & 高國雄牧師 & 路加福音2:41-52 & \hyperref[sec:580]{你們為甚麼找尋我呢} \\
\hline \hline
2022年1月2日 & 高國雄牧師 & 約翰福音1:10-18 & \hyperref[sec:581]{道成了住在我們中間的肉身} \\
2022年1月9日 & 李均熊博士 & 天上來的聲音 & \hyperref[sec:582]{天上來的眾聲} \\
2022年1月16日 & 鄧瑞強博士 & 創世記12:1-3 & \hyperref[sec:583]{獨上高樓望盡天涯路} \\
2022年1月23日 & 伍渭文牧師 & 詩篇19 & \hyperref[sec:584]{他在那裡他並不鍼默} \\
2022年1月30日 & 林豪恩先生 & 路加福音4:16-29 & \hyperref[sec:585]{一時的光輝未必是永恆} \\
\hline
2022年2月6日 & 高國雄牧師 & 路加福音5:1-11 & \hyperref[sec:586]{來跟從耶穌} \\
2022年2月13日 & 林豪恩先生 & 耶利米書17:5-10 & \hyperref[sec:587]{無形事物的有形記號} \\
2022年2月20日 & 鄧瑞強博士 & 出埃及記14:10-18 & \hyperref[sec:588]{浪花淘盡英雄} \\
\hline
2022年3月6日 & 高國雄牧師 & 路加福音4:1-13 & \hyperref[sec:589]{戰勝魔鬼之道} \\
2022年3月13日 & 李均熊博士 & 路加福音13:31-35 & \hyperref[sec:590]{雞舍狐狸} \\
2022年3月20日 & 鄧瑞強博士 & 利未記19:1-18 & \hyperref[sec:591]{我聖爾亦當聖} \\
2022年3月27日 & 伍渭文牧師 & 路加福音15 & \hyperref[sec:592]{吉甲的覺醒進入應許之地} \\
\hline
2022年4月3日 & 林豪恩先生 &  & \hyperref[sec:593]{十二門徒容易得知心一個也難求} \\
2022年4月10日 & 梁元生教授 & 路加福音19:28-40 & \hyperref[sec:594]{人生如驢驢如人生} \\
2022年4月17日 & 鄧瑞強博士 & 民數記14:1-10 & \hyperref[sec:595]{路漫漫其曲折兮} \\
2022年4月24日 & 高國雄牧師 &  & \hyperref[sec:596]{超越多馬的信心} \\
\hline
2022年5月1日 & 林豪恩先生 &  & \hyperref[sec:597]{生亦何歡死亦何苦} \\
2022年5月8日 & 李均熊牧師 &  & \hyperref[sec:598]{彌賽亞牧人} \\
2022年5月15日 & 鄧瑞強博士 & 申命記10:12-22及30:5-20 & \hyperref[sec:599]{生路如此美好何解選擇滅亡} \\
2022年5月22日 & 馮天聰牧師 & 詩篇67;約翰福音14:23-29;使徒行傳16:9-15 & \hyperref[sec:600]{開始與終結} \\
2022年5月29日 & 林豪恩先生 & 使徒行傳16:16-34 & \hyperref[sec:601]{無情世界有情天} \\
\hline
2022年6月5日 & 姚志豪牧師 &  & \hyperref[sec:602]{清楚嗎} \\
2022年6月12日 & 高國雄牧師 & 羅馬書5:1-5 & \hyperref[sec:603]{我們因信稱義嗎} \\
2022年6月19日 & 鄧瑞強博士 &  & \hyperref[sec:604]{你們好好選擇吧} \\
2022年6月26日 & 林豪恩先生 & 列王記下2:1-14;路加福音9:51-62;加拉太書5:1-26;詩篇16:1-11 & \hyperref[sec:605]{離開的和留下的} \\
\hline
2022年7月3日 & 林豪恩先生 & 列王記下5:1-14 & \hyperref[sec:606]{痊癒始於放下身份} \\
2022年7月10日 & 李均熊牧師 & 阿摩司書7:7-17 & \hyperref[sec:607]{好人當賊辦} \\
2022年7月17日 & 鄧瑞強博士 & 士師記19:1-30 & \hyperref[sec:608]{將樂園變成了地獄} \\
2022年7月24日 & 高國雄牧師 & 路加福音11:1-13 & \hyperref[sec:609]{讓我們一同禱告} \\
2022年7月31日 & 王玉教授 & 歌羅西書3:1-11 & \hyperref[sec:610]{煥然一新} \\
\hline
2022年8月7日 & 林豪恩先 & 以賽亞書1:10-20 & \hyperref[sec:611]{攝禮歸仁} \\
2022年8月14日 & 高國雄牧師 & 路加福音12:49-56 & \hyperref[sec:612]{帶來分裂的信仰} \\
2022年8月28日 & 高國雄牧師 &  & \hyperref[sec:613]{耶穌教導的餐桌禮儀} \\
\hline
2022年9月4日 & 林豪恩先生 & 腓利門書1:1-21 & \hyperref[sec:614]{千里共嬋娟} \\
2022年9月11日 & 李均熊牧師 & 耶利米書4:11-12,22-28;詩篇51:1-10;提摩太前書1:12-17;路加福音15:1-10 & \hyperref[sec:615]{混沌中的盼望} \\
2022年9月18日 & 鄧瑞強博士 & 路得記1:1、6-18;2:1-13;3:1-13;4:1-6、13-17 & \hyperref[sec:616]{比亞伯拉罕更亞伯拉罕} \\
2022年9月25日 & 伍渭文牧師 &  & \hyperref[sec:617]{那日之後} \\
\hline
2022年10月2日 & 林豪恩先生 &  & \hyperref[sec:618]{當你凝視深} \\
2022年10月16日 & 鄧瑞強博士 & 撒母耳記上13:5-14 & \hyperref[sec:619]{王者的隕落} \\
2022年10月23日 & 高國雄牧師 & 路加福音18:9-14 & \hyperref[sec:620]{上帝啊開恩可憐我這個罪人} \\
2022年10月30日 & 吳梓明教授 & 約拿書3:1-10;詩篇90:1-6,12-17;哥林多前書13:4-12;路加福音10:30-35 & \hyperref[sec:621]{何明華贈予崇基的三份禮物} \\
\hline
2022年11月6日 & 林豪恩先生 & 路加福音20:27-38 & \hyperref[sec:622]{女人永遠只是男人的某某} \\
2022年11月13日 & 李均熊牧師 & 以賽亞書65:17-25 & \hyperref[sec:623]{看哪我造新天新地} \\
2022年11月20日 & 鄧瑞強博士 & 撒母耳記下11:1-4,13:7,10-14,21-22,15:13-14,16:20-22 & \hyperref[sec:624]{一子錯} \\
2022年11月27日 & 葉菁華教授 & 馬太福音24:36-44 & \hyperref[sec:625]{黑夜到底多長} \\
\hline
2022年12月4日 & 關瑞文教授 & 馬太福音3:1-12 & \hyperref[sec:626]{那盡了責的見證人} \\
2022年12月11日 & 姚志豪牧師 & 詩146 & \hyperref[sec:627]{凡不因我跌倒的就有福了} \\
2022年12月18日 & 鄧瑞強博士 & 馬太福音1:18-25 & \hyperref[sec:628]{噢流星雨} \\
2022年12月25日 & 林豪恩先生 &  & \hyperref[sec:629]{在沒有聖景之中尋找聖誕意義} \\
\hline \hline
2023年1月1日 & 林豪恩先生 & 傳道書3:1-11 & \hyperref[sec:630]{迷霧中如何前行} \\
2023年1月8日 & 李均熊牧師 & 以賽亞書42 & \hyperref[sec:631]{看哪我的僕人} \\
2023年1月15日 & 鄧瑞強博士 & 申命記17:16-19;列王記上3:1,9:15-21,24,10:26-11:9 & \hyperref[sec:632]{如此風流人物} \\
2023年1月22日 & 高國雄牧師 & 馬太福音4:12-23 & \hyperref[sec:633]{天國近了悔改吧} \\
2023年1月29日 & 林豪恩先生 &  & \hyperref[sec:634]{天堂的位置為誰而留} \\
\hline
2023年2月5日 & 高國雄牧師 & 馬太福音5:13-48 & \hyperref[sec:635]{如何可以勝過文士及法利賽人的義} \\
2023年2月19日 & 鄧瑞強博士 & 列王記下5:1-15 & \hyperref[sec:636]{大人物與小薯仔} \\
2023年2月26日 & 陳智衡牧師 & 路加福音9:51-62 & \hyperref[sec:637]{走不易的道路} \\
\hline
2023年3月12日 & 李均熊牧師 & 出埃及記17:1-7 & \hyperref[sec:638]{無水之道} \\
2023年3月19日 & 鄧瑞強博士 & 歷代志上1:1,28,3:1-10,17-19,4:1,6:1,31-33,49-50,8:1-2,9:17-19,39-40 & \hyperref[sec:639]{家譜會說話} \\
2023年3月26日 & 林豪恩先生 & 羅馬書8:6-11;詩篇130;以西結書37:1-14;約翰福音11:1-45 & \hyperref[sec:640]{人在此身中} \\
\hline
2023年4月9日 & 李均熊牧師 & 約翰福音20:1-18 & \hyperref[sec:641]{主的復活} \\
2023年4月16日 & 鄧瑞強博士 & 歷代志下36:5-23 & \hyperref[sec:642]{是結束也是開始} \\
2023年4月23日 & 林豪恩先生 & 路加福音24:13-35 & \hyperref[sec:643]{如果耶穌顯現你認得出他嗎} \\
2023年4月30日 & 林豪恩先生 &  & \hyperref[sec:644]{同羊唔同命} \\
\hline
2023年5月7日 & 高國雄牧師 & 約翰福音14:1-14 & \hyperref[sec:645]{若不藉著我沒有人能到父那裡去} \\
2023年5月14日 & 李均熊牧師 & 使徒行傳17:22-31;詩篇66:8-20;彼得前書3:13-22;約翰福音14:15-21 & \hyperref[sec:646]{為那不可知的神作見證} \\
2023年5月21日 & 鄧瑞強博士 & 以斯拉記1:1-4,7,4:1-7,18-24,5:1-6,6:1-15,19-22;7:21-23 & \hyperref[sec:647]{奉誰的旨} \\
2023年5月28日 & 高國雄牧師 & 約20:21-23 & \hyperref[sec:648]{領受聖靈吧} \\
\hline
2023年6月4日 & 林豪恩先生 &  & \hyperref[sec:649]{三一萬能俠} \\
2023年6月11日 & 李均熊牧師 &  & \hyperref[sec:650]{lekh-lekhah你要去} \\
2023年6月18日 & 鄧瑞強博士 & 尼希米記13:10-31 & \hyperref[sec:651]{在淤泥上植一池蓮花} \\
2023年6月25日 & 高國雄牧師 & 馬太福音10:24-39;羅馬書6:3-11 & \hyperref[sec:652]{真是我的最佳選擇嗎} \\
\hline
2023年7月2日 & 蒲錦昌牧師 & 路得記1:1-22,4:13-22 & \hyperref[sec:653]{命運就算顛沛流離} \\
2023年7月9日 & 李均熊牧師 & 創世記24 & \hyperref[sec:654]{識見僕人柔和謙卑} \\
2023年7月16日 & 鄧瑞強博士 & 以斯帖記1:1,10-22,2:16-17,21-23,3:1-6,4:7-16,5:12-14,6:1-9,12-13;7:1-10 & \hyperref[sec:655]{荒謬而諧趣} \\
2023年7月23日 & 林豪恩先生 & 創世記28:10-14;馬太福音13:24-43 & \hyperref[sec:656]{人在旅途做夢時} \\
2023年7月30日 & 梁元生教授 &  & \hyperref[sec:657]{來到了尼波山} \\
\hline
2023年8月6日 & 高國雄牧師 & 馬太福音14:13-21 & \hyperref[sec:658]{你們給他們吃吧} \\
2023年8月13日 & 林豪恩先生 &  & \hyperref[sec:659]{時不利兮} \\
2023年8月20日 & 高國雄牧師 & 馬太福音15:21-30 & \hyperref[sec:660]{婦人你的信心很大} \\
2023年8月27日 & 李均熊牧師 &  & \hyperref[sec:661]{施弗拉和普亞} \\
\hline
2023年9月3日 & 伍渭文牧師 &  & \hyperref[sec:662]{焚而不毀的荊棘} \\
2023年9月10日 & 李均熊牧師 &  & \hyperref[sec:663]{逾越抑保護} \\
2023年9月17日 & 高國雄牧師 & 馬太福音18:21-35 & \hyperref[sec:664]{主啊如若弟兄得罪我} \\
\hline
2023年10月1日 & 林豪恩先生 &  & \hyperref[sec:665]{大能的摩西與無能的基督} \\
2023年10月8日 & 李均熊牧師 &  & \hyperref[sec:666]{對話中的十誡} \\
2023年10月15日 & 鄧瑞強博士 & 約伯記1:1-22 & \hyperref[sec:667]{天地間的一場豪賭} \\
2023年10月22日 & 吳梓明教授 &  & \hyperref[sec:668]{博雅教育與中國基督教大學的故事} \\
2023年10月29日 & 蘇成溢牧師 &  & \hyperref[sec:669]{為福音作見證} \\
\hline
2023年11月5日 & 林豪恩先生 &  & \hyperref[sec:670]{誰是叛教者} \\
2023年11月12日 & 李均熊牧師 &  & \hyperref[sec:671]{事奉神諗清楚好喎} \\
2023年11月19日 & 鄧瑞強博士 & 詩篇82:1-8 & \hyperref[sec:672]{啊上帝真忙} \\
2023年11月26日 & 高國雄牧師 & 馬太福音25:31-46 & \hyperref[sec:673]{主啊我甚麼時候見到你呢} \\
\hline
2023年12月3日 & 鄧瑞強博士 & 箴言26:1-12 & \hyperref[sec:674]{要不要照愚昧人的愚妄話回答他} \\
2023年12月10日 & 李均熊牧師 & 以賽亞書40:1-31 & \hyperref[sec:675]{黎緊咁即係未到} \\
2023年12月17日 & 姚志豪牧師 &  & \hyperref[sec:676]{猜猜我是誰} \\
2023年12月24日 & 高國雄牧師 & 路加福音1:26-38 & \hyperref[sec:677]{從上帝而來的沒有任何事情是不可能的} \\
2023年12月24日 & 李均熊牧師 & 馬太福音1:18-25;以賽亞書7:1-17 & \hyperref[sec:678]{應許與期望} \\
2023年12月31日 & 林豪恩先生 & 路加福音2:22-40 & \hyperref[sec:679]{至人只是常} \\
\hline \hline
2024年1月7日 & 高國雄牧師 &  & \hyperref[sec:680]{傳道者約翰} \\
2024年1月14日 & 李均熊牧師 &  & \hyperref[sec:681]{始料不及} \\
2024年1月21日 & 鄧瑞強博士 & 傳道書3:1-14 & \hyperref[sec:682]{命限恩典義} \\
2024年1月28日 & 林豪恩先生 &  & \hyperref[sec:683]{顯現有時隱藏有時} \\
\hline
2024年2月11日 & 李均熊牧師 & 馬可福音9:2-13 & \hyperref[sec:684]{以利亞已經來了} \\
2024年2月18日 & 鄧瑞強博士 & 雅歌4:1-7,12-16,5:6-8,10-16,8:5-7 & \hyperref[sec:685]{重返伊甸} \\
2024年2月25日 & 郭志榮牧師 & 馬太福音5:17-20 & \hyperref[sec:686]{止於至善} \\
\hline
2024年3月3日 & 李均熊牧師 & 哥林多前書1:18-25 & \hyperref[sec:687]{智慧抑兆頭} \\
2024年3月10日 & 高國雄牧師 & 約翰福音3:16 & \hyperref[sec:688]{上帝愛世人} \\
2024年3月17日 & 鄧瑞強博士 & 以賽亞書6:1-13 & \hyperref[sec:689]{我見主坐在高高的寶座上} \\
2024年3月24日 & 林豪恩先生 &  & \hyperref[sec:690]{揮舞棕枝的激情} \\
2024年3月31日 & 張漢強牧師 & 馬太福音28:1-15 & \hyperref[sec:691]{願你們平安} \\
\hline
2024年4月7日 & 高國雄牧師 & 約翰福音20:19-31 & \hyperref[sec:692]{領受聖靈就會看見} \\
2024年4月14日 & 李均熊牧師 & 路加福音24:36-48 & \hyperref[sec:693]{復活基督的顯現} \\
2024年4月21日 & 高國雄牧師 & 約翰福音10:11-18 & \hyperref[sec:694]{捨去與重拾} \\
2024年4月28日 & 林豪恩先生 &  & \hyperref[sec:695]{往加沙的路上} \\
\hline
2024年5月5日 & 高國雄牧師 &  & \hyperref[sec:696]{常在主愛中} \\
2024年5月12日 & 李均熊牧師 & 詩篇1:1-6;約翰福音17:6-19 & \hyperref[sec:697]{栽在溪水旁} \\
2024年5月19日 & 鄧瑞強博士 & 耶利米書28:1-17 & \hyperref[sec:698]{誰是真先知} \\
\hline
2024年6月2日 & 林豪恩先生 &  & \hyperref[sec:699]{當遵守律法與追求美善不可兼得的時候} \\
2024年6月9日 & 李均熊牧師 & 詩篇130:1-9 & \hyperref[sec:700]{深淵裡的拯救} \\
2024年6月16日 & 鄧瑞強博士 & 耶利米哀歌1:1-5,2:20-21,3:22-27,4:9-10,5:21-22 & \hyperref[sec:701]{無處話淒涼} \\
2024年6月23日 & 高國雄牧師 & 馬可福音4:35-41 & \hyperref[sec:702]{老師我們喪命你不顧嗎} \\
2024年6月30日 & 林豪恩先生 &  & \hyperref[sec:703]{在絕望孤單中偶然遇到你} \\
\hline
2024年7月7日 & 李均熊牧師 & 詩篇123:1-4 & \hyperref[sec:704]{上主啊求你憐憫我們} \\
2024年7月14日 & 林豪恩先生 &  & \hyperref[sec:705]{盤子上的頭和籃子上的餅} \\
2024年7月21日 & 鄧瑞強博士 &  & \hyperref[sec:706]{這些骸骨能復活嗎} \\
2024年7月28日 & 高國雄牧師 &  & \hyperref[sec:707]{你看見什麼} \\
\hline
2024年8月4日 & 陳永業傳道 & 約翰福音4:19-26 & \hyperref[sec:708]{心靈和誠實} \\
2024年8月11日 & 義人的好處 & 詩篇34:1-22 & \hyperref[sec:709]{義人的好處} \\
2024年8月18日 & 鄧瑞強博士 & 但以理書7:1-14 & \hyperref[sec:710]{千古興亡多少事悠悠} \\
2024年8月25日 & 林豪恩先生 &  & \hyperref[sec:711]{靈與肉} \\
\hline
2024年9月1日 & 李均熊牧師 & 詩篇15:1-5 & \hyperref[sec:712]{得居聖山者之品行} \\
2024年9月8日 & 姚志豪牧師 &  & \hyperref[sec:713]{睇} \\
2024年9月15日 & 高國雄牧師 & 馬可福音8:27-38 & \hyperref[sec:714]{性價比與等價交換} \\
2024年9月22日 & 鄧瑞強博士 & 何西阿書4:11-14及11:8-11 & \hyperref[sec:715]{苦苦單戀} \\
2024年9月29日 & 關瑞文教授 & 雅各書5:13-20 & \hyperref[sec:716]{信心這回事} \\
\hline
2024年10月6日 & 李均熊牧師 & 詩篇8:1-9 & \hyperref[sec:717]{人算甚麼} \\
2024年10月13日 & 林豪恩先生 & 馬可福音10:17-31 & \hyperref[sec:718]{錢財是永生的敵人嗎} \\
2024年10月20日 & 伍渭文牧師 & 馬可福音10:35-45 & \hyperref[sec:719]{因真理得自由以服務} \\
2024年10月27日 & 吳梓明教授 & 詩篇126篇1-6;約伯記42:1-6及10-17;使徒行傳2:1-21;馬可福音10:46-52 & \hyperref[sec:720]{聖靈引導教會進入真理} \\
\hline
2024年11月3日 & 高國雄牧師 &  & \hyperref[sec:721]{上帝的國離我有多遠} \\
2024年11月10日 & 李均熊牧師 &  & \hyperref[sec:722]{寡婦和信仰} \\
2024年11月17日 & 鄧瑞強博士 & 約珥書1:4-2:32 & \hyperref[sec:723]{耶和華的日子臨近了} \\
2024年11月24日 & 林豪恩先生 &  & \hyperref[sec:724]{在亂世中再思基督君王的意義} \\
\hline
2024年12月1日 & 李均熊牧師 &  & \hyperref[sec:725]{將臨期的意義} \\
2024年12月8日 & 高國雄牧師 & 路加福音3:1-6 & \hyperref[sec:726]{在盼望中邁向更美好的新一天} \\
2024年12月15日 & 鄧瑞強博士 & 阿摩司書4:6-5:24 & \hyperref[sec:727]{就這樣習以為常地活了一生} \\
2024年12月22日 & 林豪恩先生 &  & \hyperref[sec:728]{為這個世界帶來一個新生命} \\
2024年12月24日 & 高國雄牧師 &  & \hyperref[sec:729]{渴求的平安在地上屬於人的} \\
2024年12月29日 & 郭志榮牧師 &  & \hyperref[sec:730]{東方博士的結局正邪交戰的序幕} \\
\hline \hline
2025年1月5日 & 高國雄牧師 &  & \hyperref[sec:731]{在我們當中住下的道} \\
2025年1月12日 & 林豪恩先生 &  & \hyperref[sec:732]{從耶穌受洗說起} \\
2025年1月19日 & 鄧瑞強博士 & 俄巴底亞書1:1-21（和合本修訂版） & \hyperref[sec:733]{相煎何太急} \\
2025年1月26日 & 李均熊牧師 &  & \hyperref[sec:734]{耶穌講道} \\
\hline
2025年2月2日 & 梁遠耀牧師 & 路加福音4:21-30 & \hyperref[sec:735]{心想事成} \\
2025年2月9日 & 李均熊牧師 &  & \hyperref[sec:736]{呼召} \\
2025年2月16日 & 高國雄牧師 &  & \hyperref[sec:737]{愛到底} \\
2025年2月23日 & 鄧瑞強博士 & 約拿書3:1-4:11 & \hyperref[sec:738]{斬斷延續怨恨的鎖鏈} \\
\hline
2025年3月2日 & 李均熊牧師 &  & \hyperref[sec:739]{主的榮光} \\
2025年3月9日 & 高國雄牧師 &  & \hyperref[sec:740]{勝過人生三大試探} \\
2025年3月16日 & 林豪恩先生 &  & \hyperref[sec:741]{今天明天後天我必須向前走} \\
2025年3月30日 & 高國雄牧師 & 路加福音15:1-3及11b-32 & \hyperref[sec:742]{當歡喜快樂} \\
\hline
2025年4月13日 & 李均熊牧師 & 路加福音19:28-40 & \hyperref[sec:743]{石頭也要呼叫} \\
2025年4月20日 & 鄧瑞強博士 & 約翰福音20:11-18 & \hyperref[sec:744]{復活天人合一的故事} \\
2025年4月27日 & 高國雄牧師 & 約翰福音20:19-31 & \hyperref[sec:745]{不要疑惑總要信} \\
\hline
2025年5月4日 & 李均熊牧師 & 約翰福音21:1-19 & \hyperref[sec:746]{岸邊的陌生人} \\
2025年5月11日 & 林豪恩先生 & 使徒行傳9:36-43;詩篇23:1-6;啟示錄7:9-17;約翰福音10:22-30 & \hyperref[sec:747]{誰為你哭泣誰為你抹淚} \\
2025年5月18日 & 鄧瑞強博士 & 彌迦書4:1-4;6:1-13 & \hyperref[sec:748]{我們獻的神不喜悅神喜悅的我們願獻嗎} \\
2025年5月25日 & 郭志榮牧師 & 詩篇127:1-5 & \hyperref[sec:749]{詩篇一百二十七篇的隱藏信息} \\
\hline
2025年6月1日 & 林豪恩先生 & 使徒行傳16:16-34 & \hyperref[sec:750]{一個關於宗教信仰與公共安全的故事} \\
2025年6月8日 & 李均熊牧師 & 使徒行傳2:1-21 & \hyperref[sec:751]{五旬節到了} \\
2025年6月15日 & 鄧瑞強博士 & 那鴻書1:1-3;1:9-2:2;3:1-7,19 & \hyperref[sec:752]{天地有正氣} \\
2025年6月22日 & 高國雄牧師 & 路加福音8:26-39 & \hyperref[sec:753]{回到你的家} \\
2025年6月29日 & 鄭躍培牧師 & 路加福音10:25-37;雅各書1:27 & \hyperref[sec:754]{誰是我的鄰舍呢} \\
\hline
2025年7月6日 & 高國雄牧師 & 路加福音10:1-20 & \hyperref[sec:755]{成為狼群中的羊} \\
2025年7月13日 & 李均熊牧師 & 路加福音10:25-37 & \hyperref[sec:756]{好撒馬利亞人} \\
2025年7月27日 & 林豪恩先生 & 何西亞書1:2-10 & \hyperref[sec:757]{被歷史烙印的家庭} \\
\hline
2025年8月3日 & 梁遠耀牧師 & 路加福音12:13-21 & \hyperref[sec:758]{追啊追} \\
2025年8月10日 & 李均熊牧師 & 路加福音12:32-40 & \hyperref[sec:759]{警醒的僕人} \\
2025年8月17日 & 鄧瑞強博士 & 哈巴谷書1:1-3:19 & \hyperref[sec:760]{絕望中的希望} \\
2025年8月24日 & 葉菁華教授 & 路加福音13:10-17 & \hyperref[sec:761]{直起腰來} \\
2025年8月31日 & 林豪恩先生 & 耶利米書2:4-13 & \hyperref[sec:762]{破裂水池永不盈活水滋潤生命源} \\
\hline
2025年9月7日 & 李均熊牧師 & 路加福音14:25-33 & \hyperref[sec:763]{作門徒的代價} \\
2025年9月14日 & 高國雄牧師 & 路加福音15:1-10 & \hyperref[sec:764]{1比99更喜樂} \\
2025年9月21日 & 鄧瑞強牧師 & 西番雅書1:1-12;3:1-4;2:10-11;3:9-20 & \hyperref[sec:765]{若果世人不是那麼邪惡} \\
2025年9月28日 & 林豪恩先生 & 耶利米書32:1-15 & \hyperref[sec:766]{在世界崩塌之際買一塊地} \\
\hline
2025年10月5日 & 姚志豪牧師 & 耶利米哀歌1:1-6 & \hyperref[sec:767]{請增加我們的信心} \\
2025年10月12日 & 李均熊牧師 & 耶利米書29:1-7 & \hyperref[sec:768]{為這城求平安} \\
2025年10月19日 & 鄧瑞強牧師 & 哈該書1:1-2:23 & \hyperref[sec:769]{欲窮千里目更上一層樓} \\
2025年10月26日 & 林豪恩先生 & 提摩太後書4:6-18 & \hyperref[sec:770]{閉起雙眼前你最掛念誰} \\
\hline
2025年11月2日 & 李均熊牧師 & 哈巴谷書1:1-4;2:1-4 & \hyperref[sec:771]{這要到幾時} \\
2025年11月9日 & 高國雄牧師 & 路加福音20:27-38 & \hyperref[sec:772]{來認識活人的上帝} \\
2025年11月16日 & 關瑞文教授 & 路加福音21:5-19 & \hyperref[sec:773]{甚麼時候有這些事呢} \\
2025年11月23日 & 高國雄牧師 & 路加福音23:33-43 & \hyperref[sec:774]{我們的上帝要施行救贖} \\
2025年11月30日 & 郭志榮牧師 & 詩篇24:1-10 & \hyperref[sec:775]{榮耀的王將要進來} \\
\hline
2025年12月7日 & 李均熊牧師 & 以賽亞書11:1-10 & \hyperref[sec:776]{耶西的根} \\
2025年12月14日 & 林豪恩先生 & 以賽亞書35:1-10;馬太福音11:2-11 & \hyperref[sec:777]{有甚麼值得喜樂} \\
2025年12月21日 & 鄧瑞強牧師 & 撒迦利亞書4:1-14 & \hyperref[sec:778]{不是倚靠勢力不是倚靠才能乃是倚靠神的靈} \\
2025年12月24日 & 高國雄牧師 & 路加福音2:1-14 & \hyperref[sec:779]{你聽過嗎} \\
2025年12月28日 & 林豪恩先生 & 馬太福音2:13-23 & \hyperref[sec:780]{天國大歷史之嬰孩耶穌逃命記} \\
\hline \hline
2026年1月4日 & 李均熊牧師 & 耶利米書31:7-14 & \hyperref[sec:781]{因為我要使他們的悲哀變為歡喜} \\
2026年1月11日 & 高國雄牧師 & 馬太福音3:13-17 & \hyperref[sec:782]{盡諸般的義} \\
2026年1月18日 & 鄧瑞強牧師 & 瑪拉基書3:13-4:6 & \hyperref[sec:783]{你們用話頂撞我} \\
2026年1月25日 & 王家輝牧師 & 哥林多前書1:10-18 & \hyperref[sec:784]{十字架的道理} \\
\hline
2026年2月1日 & 李均熊牧師 & 彌迦書6:1-8 & \hyperref[sec:785]{信仰的要求} \\
2026年2月15日 & 鄧瑞強牧師 & 馬太福音5:38-48 & \hyperref[sec:786]{不切實際的耶穌} \\
2026年2月22日 & 林豪恩先生 & 創世記2:15-17,3:1-7;詩篇32:1-11;羅馬書5:12-19;馬太福音4:1-11 & \hyperref[sec:787]{三項試探一個中心} \\
\hline
2026年3月1日 & 高國雄牧師 & 約翰福音3:1-17 & \hyperref[sec:788]{上帝愛世界} \\
2026年3月8日 & 李均熊牧師 & 約翰福音4:5-42 & \hyperref[sec:789]{活水之道} \\
2026年3月22日 & 郭志榮牧師 & 詩篇130:1-8 & \hyperref[sec:790]{從深處的覺醒} \\
2026年3月29日 & 林豪恩先生 & 馬太福音21:1-11 & \hyperref[sec:791]{路不能亂走車不可亂坐} \\
\hline
2026年4月5日 & 高國雄牧師 & 約翰福音20:1-18 & \hyperref[sec:792]{基督今復活} \\
2026年4月12日 & 李均熊牧師 & 使徒行傳2:14,22-32;詩篇16:1-11;彼得前書1:3-9;約翰福音20:19-31 & \hyperref[sec:793]{我信身體復活} \\
2026年4月19日 & 鄧瑞強牧師 & 馬可福音8:29-9:1 & \hyperref[sec:794]{只有一位受苦的上帝能救我們} \\
\hline
2026年5月3日 & 李均熊牧師 & 彼得前書2:2-10 & \hyperref[sec:795]{活石} \\
2026年5月10日 & 林豪恩先生 & 彼得前書3:13-22;約翰福音14:15-21 & \hyperref[sec:796]{耶穌不在的時候} \\
2026年5月17日 & 鄧瑞強牧師 & 路加福音7:11-17 & \hyperref[sec:797]{夕陽西下斷腸人在天涯} \\
2026年5月24日 & 高國雄牧師 & 使徒行傳2:1-21;詩篇104:24-34;哥林多前書12:3-13;約翰福音20:19-23 & \hyperref[sec:798]{我們今天一同得這能力吧} \\
\hline
2026年6月7日 & 李均熊牧師 & 創世記12:1-9;詩篇33:1-12;羅馬書4:13-25;馬太福音9:9-13,18-26 & \hyperref[sec:799]{放心你的信治好你了} \\
2026年6月14日 & 高國雄牧師 & 創世記18:1-15;詩篇116篇1-2及12-19;羅馬書5:1-8;馬太福音9:35-10:8 & \hyperref[sec:800]{在主前蒙恩} \\
2026年6月21日 & 鄧瑞強牧師 & 約翰福音10:1-18 & \hyperref[sec:801]{從一扇門到心深處} \\
2026年6月28日 & 林豪恩先生 & 創世記22:1-14;詩篇13篇1-6;羅馬書6:12-23;馬太福音10:40-42 & \hyperref[sec:802]{要到幾時呢} \\
\hline
2026年7月5日 & 高國雄牧師 & 羅馬書7:15-25;馬太福音11:16-19 & \hyperref[sec:803]{好人難做} \\
2026年7月12日 & 李均熊牧師 & 創世記25:19-34 & \hyperref[sec:804]{我都就死這於我有甚麼益處} \\
2026年7月19日 & 鄧瑞強牧師 & 使徒行傳16:6-10 & \hyperref[sec:805]{直搗羅馬} \\
\end{xltabular}

}


\newpage

\section{栽在溪水旁的樹}
\label{sec:0}
\textbf{伍渭文牧師 - 栽在溪水旁的樹}
\newline
\newline
\hyperref[sec:0]{< < < PREV SERMON < < <}
~
\hyperref[sec:toc]{[返目錄]}
~
\hyperref[sec:1]{> > > NEXT SERMON > > >}


詩篇 1:3
\newline
\begin{longtable}{cl}
\hline
\hline
章節 & 經文 (和合本修訂版)\\
\hline
1:3 & \begin{tabularx}{0.7\textwidth}{X} 他要像一棵樹栽在溪水旁，按時候結果子，葉子也不枯乾。凡他所做的盡都順利。 \end{tabularx} \\ \\
[1ex]
\hline
\hline
\end{longtable}
日期:2009年6月14日

証道:伍渭文牧師

題目:栽在溪水旁的樹

經文:詩篇 1:3

栽在溪水旁的樹,按時候結果子,葉子也不枯乾...(詩 1:3)

聽到這段經文的時候,你聯想到甚麼呢?

我聯想到以前的香港電訊以及匯豐銀行的股份,對於很多本很退休的人來說,如果手持香港電訊或者匯豐銀行的股份的話,它就按時候結果子,葉子也不枯乾；甚至你擁有幾座物業來收租,你也可以說好像一棵樹,按時候結果子,葉子也不枯乾.特別曾經香港電訊或匯豐銀行的股票價格,不斷的高攀,但曾幾何時,香港電訊成為電訊盈科的時候,葉子就枯乾了.而這個星期,匯豐銀行告訴我們,你存一萬塊錢一年,只有一塊錢的利息.

我們感謝上主,詩篇第一篇所講的,有福的人,他的福氣就不是一棵的搖錢樹可以給你財富,那些的樹會茂盛,但是也會枯乾.在華爾街用貪婪所種植的樹,現在已經連根拔起.

詩篇第一篇所描述的圖畫,是講一個上帝所喜悅的人,義人他的生活的寫照:第一節提到,這些人是甚麼呢?是「不從惡人的計謀,不站罪人的道路,不坐褻慢人的座位」.開宗名義,這篇詩篇提到一個蒙福的人,他不是擁有甚麼,不是擁有財富,乃是擁有一種生命!

我們知道真正的福氣,不是你所「有」,乃是你所「是」.

擁有一張好的床墊,不擔保你有好的睡眠；擁有外型亮麗的身材,不擔保你有一個peace of mind,一個很祥和,平伏的思想,特別是你經過很努力的纖體工程,很亮麗,但是整天都寢食不安,深怕自己又吃多了,吃多了就又變型了,更加的沒有心裡頭的平安.

你有財富不擔保你有好的健康,聖經裡頭所描述有關蒙福的人,是他的生命.首先這個生命,「不從惡人的計謀,不站罪人的道路,不坐褻慢人的座位」,是遠離罪惡的,這裡提到的罪是越來越嚴重的!

「不從惡人的計謀」:是說到一個人要解決一個實際的 practical的問題,他要解決具體的課題,他可能有試探,去尋求不合神心意的計謀.這三個是一個禁戒:不從,不站罪人的道路,不坐褻慢人的座位；用這三個很明顯的禁戒,就是針對人性的傾向,早在十誡的裡頭,上帝很清楚的禁戒:「不可殺人,不可姦淫,不可作假見證,不可偷盜」,這些禁戒都是針對人性的本性,人很容易按自己的喜好去偷盜,做不應該做的事,所以這裡提到「不可」,不要「從惡人的計謀」.

開始的時候是從惡人的計謀,是要解決pratical的問題,慢慢這個意念進入腦海裡的時候,我們就會很容易成為一個例常的習慣,「不站罪人的道路」就成為一個 habitual behavior,將想念成為一個具體行為的時候,就成了具體的罪行了.

當在樂園的時候,蛇誘惑亞當,夏娃:你看見那個果子是很漂亮,上帝說「不可吃」,是否真的這麼說呢?神豈說,你吃的時候,眼睛就明亮,你像上帝一樣能夠分別善惡,你會像上帝的啦!這是一個計謀──「不從惡人的計謀」,但當夏娃吃了的時候,就成為一個具體的罪行,結果在創世記十一章的時候,在示拿平原,巴比倫人要建巴別塔,「我們要建造一座城,和一座塔,塔頂通天,為要傳揚我們的名,免得我們分散在全地上」.建塔就好像「神」,傳揚我們的名,取代上帝.

在樂園裡的計謀,已經成為了一個工程──做神的工程──人要抗拒上帝,取代上帝:我們建一個塔,傳揚我們的名.所以從第一個 Evil thought,慢慢地就成為一個habitual behavior──具體的,慣常的罪行,接著,更嚴重的是每況愈下,成為一個「坐褻慢人的座位」.

「褻慢人」的意思是去嗤笑一些不和他同一夥的人,照聖經的傳統,講道的人,講律法的人,坐在摩西的座位上,是坐著的,耶穌很清楚的說:當「文士和法利賽人,坐在摩西的位上．凡他們所吩咐你們的,你們都要謹守,遵行．但不要效法他們的行為.」(太23:2-3),所以「坐褻慢人的座位」的意思是,甚至將一些抗拒神的意念成為教導,教師去教導人,越來越嚴重.

當耶穌在約翰福音第九章醫好一個生來就瞎眼的人時,那天是安息日,文士和法利賽人竟然說,耶穌是罪人,因為祂在安息日醫病,那瞎眼的人說,我不知道祂是誰,但我知道,沒理由神聽罪人的說話把我給醫好了,所以成了神學的辯論,是不是呢?

所以看到這三節的經文,最初是為了解決practical具體的課題,可能有試探,聽從惡人的計謀,接著成為具體的罪行,一路在做,不自覺的時候,還去教導人去反對,去嗤笑一些不跟從自己錯誤的人.

崇基學院每個星期五,如果你是學生,你就知道要參加週會的,而參加週會要過半數的出席率,十三次要參加七次,如果不是你就不及格了,但不及格也沒甚麼大不了,只是在成績單寫著:Fail,但並不影響你畢業的,也不影響你的績點的,不過給人看見就不好看罷了,而同學們都很認真的,但有些就不夠堂數,那怎麼辦?因為同時間需要參加另外一個活動,他就找一些同學幫他刷卡,刷了卡就把你的名字給記下來,就是不是自己的名字,也不是自的容貌,因那卡有名字,有照片,幫人刷卡同學沒帶卡來,我就說,你簽名吧,然後我們檢查後發現有一些冒簽的,如是冒簽的時候,就被邀請來到我的辦公室談話,我很想跟他談話,我首先跟他說,這是紀律的問題,你一定沒有分數 Fail 了,但我關注的是... 還好你這次給我捉到,如果不讓我捉到的時候,你會覺得:咦... 還不錯哦!可以幫助人!或許是一個希望人幫助的人,就想:我不用去了,有人幫我呀!結果呢?解決了一個practical的問題,具體的問題,可以拿到「出席」記錄,慢慢成為習慣,偽造文書的習慣,這個很嚴重.

我用一個例子,在澳門,一個很高層的貪污官員,姓歐的,大家都知道.我們想像不到在查他的案件,去他家搜查的時候,發現很多很好的紅酒,很多珠寶,很多現金,不能想像之多,不可能的,就是說,這位先生他的貪污已成為一種習慣,成為一個obsession,成為一個不能擺脫的沉溺行為,那是很可怕的,根本他不需要那麼多錢,但他一開始了之後,越來越不能自拔,如果你不貪的話,就有人笑你啦,「褻慢人」說,為什麼不這麼做呢?最要緊的是在一開始之時,你要守住的心,不要從惡人計謀的試探,一失足成千古恨,你慢慢不自覺的.

我有一個具體的例子,是心裡面很不開心的,所以我更加關心同學他們的行為,有一個同學(報章已報導過的了),大學生來的,他有一次去澳門賭錢,竟然贏了很多錢,心就紅了,再去賭錢,後來全輸掉了,接著就向財務公司借錢再去拚博,之後要還很多的利息,怎麼辦呢?偷竊!成為偷竊的慣匪──偷手機,手提電腦,結果有一次人贓俱獲,就被判監,沒法完成學位,是一個很具體的事例,所以我心裡頭很擔憂,一旦同學的心不正的時候,就一步一步.... 你從惡人的計謀,慢慢成為慣性的罪行,甚至最後你會譏笑那些不跟從你一樣去行的人,所以聖經裡面,箴言很清楚說:「你要保守你心,勝過保守一切,因為一生的果效,是由心發出」(箴4:23).

心很重要,我們作父母的,要教導孩童,到他年老的時候也不會偏離.我們用神的道,神的律法,來教訓我們的孩童,使他的心正.

詩篇第一篇講到一個有福的人,第一個很重要的是要離開罪惡,從意念開始,遠離惡人的計謀,要警覺,因為罪是不斷蠶食我們的心靈,一上這個斜坡的時候,就每況愈下的滑下來,不能自制.

第二個有福的人的記號,除了離開罪惡之後,積極要怎麼樣呢?

第二節:「惟喜愛耶和華的律法,晝夜思想,這人便為有福.」這裡提到我們不從惡人的計謀,不從罪人的道路,不坐褻慢人的座位,我們需要有力量,有力量就是要默想神的說話──晝夜思想.

「喜愛」是甚麼意思呢?「喜愛」就是很自發的,甘心情願的,很積極的,主動的喜愛神的律法,我希望你來到崇拜是喜愛的,不是人家強迫你來的,你很喜愛星期天來到崇基禮拜堂親近上帝.這個心是一個好的土壤的心,為了這個土壤,我們開始思想神的律法.

思想這個字,在英文是:Meditation,是思考,是默想,默想是需要用力的!

有一次,有一個雕刻家,米蓋蘭基羅Michelangelo,一直對著一個石頭敲擊,用力的敲擊,有時就輕輕的雕琢,有人問他,為什麼你那麼大力,勤勞,又那麼細心地敲擊這石頭呢?他說:我敲擊這石頭,因為我要將裡面的人──釋放他.很有意思!我看到裡面有一個人,一個人形,我看到了,我現在要用力,用心,去把他給釋放出來!

我相信這個說法,都可以應用在我們對神的律法的心,上帝的話語很奇妙,裡面的meaning意義,需要我們用心,用力去推敲,透過上文下理來思考它,如果你很懶惰,就水過鴨背,看完就算,你不會找到上帝話語的奇妙,唯有那些人,他們晝夜思想,很想明白神話語的奇妙,很想知道神的話語對我有甚麼關係呢?他默想,默想不單單用腦去思考,亦是對自己心靈來說話,所以在思考神的話語之先,首先,第一節,離開罪惡,以致能夠親近神的話語,你親近神的話語,又可以幫助你離開罪惡,是兩個為因果的,如果你不肯離開罪惡,你是從惡人的計謀,甚至你自己有慣常的罪行,甚至嘲弄不跟從你的人,而成為褻慢人,你怎可明白神的話語呢?

神的話語是刀,兩刃的利劍,能夠剖開我們的靈與魂,我們裡面隱而未現的罪,神的話語會光照我們,所以不肯離開罪的人,是沒有可能完全明白神話語的奇妙,你只不過明白字面上的意思,它不會觸動你心靈裡的深處,我們要讓神的話語來光照我們的靈魂,我們要很謙卑,我們要省查自己的罪行,晝夜思想,我們希望得到裡面的意思,這就一個有福氣之人的記號,就是他會去思考神的話語,他以神的話語為滿足,又透過神話語的力量,可以逃避惡人計謀的誘惑.

第三,這個有福的人,是按時候結果子的人──「他要像一棵樹栽在溪水旁,按時候結果子,葉子也不枯乾．凡他所作的,盡都順利.」

一,他是上帝所栽種的,他不是一棵野的樹,不知種在那裡的,乃是上帝很清楚,有計劃將他栽種於某一個地方,他是尋求神旨意的人.每一個人生在那一個家庭,都不是他的選擇,他做甚麼工作,好像是他自己找的,但其實是上帝冥冥之中的引導,他是一棵被上主栽種的樹,是上主在他生命路程裡面有一個很清楚的引導和計劃,而他得到的栽種,雨露,以致使他能夠得到成長,結果子,都不是他自己的,一棵樹的成長之後,他的果子,是為人而生產的,他的樹蔭是給經過的旅客遮擋太陽的,在他的蔭下歇息,而他所需要的水,陽光是上天給他的,所以這個果子不是說你去跟隨上帝,你有很多的財富,神給你有健康,不是你自己享用的,乃是為其他人來享用的,一個人他蒙福,蒙福最要緊是甚麼呢?是使別人得福,這是亞伯拉罕的原則,我們幾個月裡在舊約看亞伯拉罕,以撒,雅各,對不?

亞伯拉罕被上帝呼召的時候說:我要使你得福氣,但你要使別人得福,所以樹被栽種在某一個地方,一個位置,不是偶然的,是上主透過你的果子,使其他的人得到飽足的,是你的樹蔭,使其他的人得到蔭庛的,這是上帝為他人而栽種你的,人最好的生命,最蒙福的生命是使別人得福氣的生命.

講到這個蒙福的生命,就使我想起八福.八福是有福的人生,如同詩篇第一篇講到有福的人生.八福一開頭,有三個福氣,是在上主面前──「虛心的人」－Poor in the Spirit,面對上主,我一無所有,以致可以來到神面前,謙卑尋求神的恩典.但因這個緣故,他看自己為有罪的──是「哀慟的人」,看上主是偉大的,看自己是一個罪人,接著推出去與人相處的時候──是「哀慟的人」.虛心對上主,哀慟對自己,溫柔對人,接著第四福是甚麼呢?「飢渴慕義的人」,他會渴想公義,又追求神的旨意,做甚麼?以至可以接著以下五,六,七,八──「使人和睦,憐恤人,甚至為義受逼迫」,使人和睦要付代價的,憐恤人要付出的,甚至為義受逼迫要犧牲自己的生命,當為他人的緣故,一個有福的人,他從虛心,從哀慟,從哀慟,從飢渴慕義之後開出去,是幫助有需要的人,使人和睦,憐恤人,是為他人而結的果子,所以在這裡說到,他的果子是按他的時候.「按時候結果子」原文是指"His Time"是「他的時候」──"His Season".我們知道做事情都有時有候的,翻雲覆雨有時,天朗氣清有時,對不?股票的牛市有時,熊市有時,不曉得甚麼時候再回到牛市啦!聚集有時,分離有時.

一棵普通的樹,它會遇到乾旱的時候,落葉有時,遇到風季的時候,它的樹枝會被折斷,枝幹折斷有時.我們的財富──富裕有時,貧窮有時.加薪有時,減薪一定有時.我們是受外在,不同的季節影響著我們.

但這個人,這個有福的人,他是掌控自己的時間,就是說,無論得時,不得時,他都在人的面前結出果子,都是那麼地真實的,他蒙神祝福他,他會感謝上帝,一個Grateful 的態度,很開心.但是當處貧窮的時候,他亦感到知足,敬虔加上知足便是大利了(提前6:6).當他在受苦或患病的時候,他會有信心,他有忍耐,盼望上帝的拯救,雖然他會面對死亡,但殺身體不能殺靈魂的,他不怕(太10:28).因為他與上帝的愛,關係不可以藉死亡將他分開的(羅8:35),在某個時刻,任何的時刻,這個人得時,不得時,都可以結果子,他的season,他的時候,是他自己的自由,不是外在的環境影響他的,所以原文很清楚的,他要結果子是按他的時候,這是一個蒙福的基督徒,蒙福的人何等自由,他無論處於任何的環境,他都可以結出果子,而這些果子是可以有果效的,凡他所作的盡都順利.

那裡不是說,你離開罪惡,離開惡人的計謀,你多默想神的話,那你作甚麼都順利的啦!你申請做研究生,一定取錄你,你參加外面的大賽,一定可拿冠軍,你去旅行,不會怕H1N1啦,神會幫助你──凡事順利嘛!不是的!

如果看上文下理,講的是那個果子的果效,就是說,一個義人,一個離開罪惡,敬畏上帝,默想上帝話語的人,他任何時間,都可以結出令人得幫助的果子,而這個果子是一路在成長的maturing,盡都順利... 下文很清楚的說,對比就是那些果效,就像糠秕被風吹散(詩1:4),他們不能結任何的果子,他們的果效就像糠秕被風吹散,所以是在說一個人,無論得時或不得時,他自己是神祝福的人,他不受環境的限制,他生命的果子可以存到永遠,是有果效的.

我想起最近的一個人,在前幾個月,一位我們很尊重,很景仰的長輩的離世,就是戴紹曾牧師的離世,大家有沒有印象呢?他是戴得生牧師的曾孫,他的名字叫紹曾,承繼他的曾祖父,戴得生是內地會的創辦人 -- OMF,前身是Inland Mission,後來成為OMF,他曾經講過他的曾祖父:我有一千英鎊,我也要全然的送給中國人,如果我有一千條生命,我全部都送給中國人.當時在英國,教會做宣教,就不是在內地的,因為內地沒有甚麼保障,所以都在沿岸一些大城市,五口通商那一類,但因神的呼召,他很想... 他不倚靠差會,憑信心入內地傳福音,因此叫「內地會」,戴紹曾牧師1929年在中國河南出生,他曾經在日本統治中國,就是中日戰爭的時候,在淪陷的地方,坐了幾年的牢,被日本人判去集中營,然後在1955年,就舉家去了台灣,生了三個小孩,1970年他是台灣一間中華福音神學院的創辦院長,他講的國語很流利,1980年他被邀請成為OMF基督海外使團(內地會離開之後成為OMF)國際主任,總部設立在新加坡,1993年他自己為中國的醫療服務,成立一個國際專業服務中心Medical Services International,負責將海外的醫生,護士推介到許多醫院內服務,特別是四川的華西醫院,所以在2007年4月4日,四川省是公開讚揚他對華人,中國人的心,頒了一個四川省榮譽公民的身份給他,在他的追思會裡(在九龍城浸信會),來了很多人,充滿了整個會場來懷念他,他是比中國人更中國人的西教士,他一生獻上了他自己的生命,只為了中國人-福音的好處,他的兒子寫了一篇悼文,悼文很短,但題目很有意思-``Living to Make a Difference",他爸爸就是活出一個不一樣的生命.

詩篇第一篇就講一個不一樣的生命,「不從惡人的計謀,不站罪人的道路,不坐褻慢人的座位．惟喜愛耶和華的律法,晝夜思想,這人便為有福」.我們下個星期就是五旬節聖靈降臨主日,聖靈降臨的時候,每一個人,你和我已經擁有聖靈；每一個人,你和我都有聖靈結出的果子,是已經有的了,已經送了給我們,不過可能我們很瘦弱,不是很壯大,這些果子是甚麼呢?就是仁愛,喜樂,和平,忍耐,恩慈,良善,信實,溫柔,節制(加拉太書 5:22-23),這些果子,是上主透過聖靈在我們心中已經結出來的.

這條溪水是不會枯乾的,這溪水是耶穌基督給我們的活水,我們的活水江河在我們心靈深處來滋潤我們的生命,可以讓我們得時,不得時,都能結出聖靈的果子.

各位親愛的弟兄姊妹,你知不知道,你是上帝所栽種的樹呢?你不是一棵野樹,你是上帝有計劃,在某地點,某時間內要栽種的樹,這棵樹的栽種不是為你自己的,是為很多經過的人成為樹蔭,成為歇息的地方,你結出的果子,為他們成為幫助,你的果子不是為自己的,是給其他人的,這些果子是大有功效的,「凡他所作的,盡都順利」.是有永遠的效果!

當我們在喪禮的時候,我們自己的安息禮拜的時候,我們會不會讓人懷念我們的生命所結出的果子呢?

我們禱告:

天父上帝我們感謝你!因為你栽種我們,你賜聖靈給我們,其實我們內心可以結出諸般聖靈的果子,成為很多人的幫助,主啊!不要讓我們的生命最後像糠秕被風吹散,讓我們生命結出的果子,永遠有功效!

我們感謝你!雖然我們人生的生命很短,但是你讓我們可以結果子,為你的國度而結果子,為人的需要而結果子,而我們本身的需要,就好像樹所需要的陽光和水份,是你所供應的!

我們感謝你!我們當中如果有人正面對著許多抉擇,面對著許多難題要去解決,求你叫他們不要依從惡人計謀,主啊!叫他們遠離這些的計謀,遠離試探,求主幫助他們能夠倚靠上主,你要幫助他,因為你說過,凡是尊重你的人,你也會尊重他的.

奉耶穌的名禱告,阿們!

“Living to Make a Difference”James Hudson Taylor IIIAugust 12, 1929 ~ March 20, 2009

Our dear father was called home to be with the Lord. Following a five year battle with liver cancer, Dad has shed his earthly tent and now is clothed in the white robe specially reserved for him. Joining the saints of the past, he stands beside Samuel  and amp; Maria Dyer, Hudson  and amp; Maria Taylor and countless others, offering up eternal praise to the Creator God and Jesus Christ. Death has indeed been swallowed up in victory!

1 Corinthians 15:54

We would like to take this opportunity to express our deepest gratitude to so many who have journeyed with our father and our family on this pilgrimage. Your prayers, calls, emails, visits, and support have been an indescribable blessing and encouragement to us. We have been conscious each and every day of the Lord’s effectual, sufficient, and sustaining grace.

The following hymn,

Finally Home

, contains words our father often quoted to encourage others facing the frailty and brevity of life. The lyrics truly speak of the imperishable hope each of us have in Jesus Christ.

When engulfed by the terror of tempestuous sea,Unknown waves before you roll.At the end of doubt and peril is eternity,Though fear and conflict seize your soul;

(Refrain)

But just think of stepping on shore and finding it heaven!Of touching a hand and finding it God’s!Of breathing new air and finding it celestial!Of waking up in glory and finding it home!When surrounded by the blackness of the darkest night,Oh how lonely death can be;At the end of the long tunnel is a shining light,For death is swallowed up in victory!

(Refrain)

Oh, that will be glory for me; Glory for me (2X);When by His grace I shall look on His face,That will be glory, be glory for me.

Dad, you have fought well, raced to the finish, and kept the faith unwaveringly

(2 Timothy 4:7)

~ now enter into the joy of your Master!

(Matthew 25:21)

A memorial service will be held in Hong Kong on April 4, 2009, at Kowloon City Baptist Church at 10:00 a.m. The address of the church is 206 Argyle Street, Kowloon, HONG KONG. As requested by our father, a “Tai-le* Educational Fund” will be established with MSI Professional Services, the organization started by our father along with two Chinese Christian brothers in 1994. MSI is located at \#19 Mongkok Road, 3rd Floor, Mongkok, HONG KONG Contribution in his memory to this fund will support continuing education of Christian professionals serving in Creative Access Nations. Separate memorial services are also planned for Taiwan and Singapore.

Jamie TaylorOn behalf ofLeone Taylor, momAmy Jackson, sisterSigne Alsin, sister



\newpage

\section{喜出望外的邂逅}
\label{sec:1}
\textbf{伍渭文牧師 - 喜出望外的邂逅}
\newline
\newline
\hyperref[sec:0]{< < < PREV SERMON < < <}
~
\hyperref[sec:toc]{[返目錄]}
~
\hyperref[sec:2]{> > > NEXT SERMON > > >}


約翰福音 4:5-29
\newline
\begin{longtable}{cl}
\hline
\hline
章節 & 經文 (和合本修訂版)\\
\hline
4:5 & \begin{tabularx}{0.7\textwidth}{X} 於是到了撒瑪利亞的一座城，名叫敘加，靠近雅各給他兒子約瑟的那塊地。 \end{tabularx} \\ \\ \relax
4:6 & \begin{tabularx}{0.7\textwidth}{X} 雅各井就在那裡；耶穌因旅途疲乏，坐在井旁。那時約是正午。 \end{tabularx} \\ \\ \relax
4:7 & \begin{tabularx}{0.7\textwidth}{X} 有一個撒瑪利亞婦人來打水。耶穌對她說：「請給我水喝。」 \end{tabularx} \\ \\ \relax
4:8 & \begin{tabularx}{0.7\textwidth}{X} 因為那時門徒進城買食物去了。 \end{tabularx} \\ \\ \relax
4:9 & \begin{tabularx}{0.7\textwidth}{X} 撒瑪利亞婦人對他說：「你是猶太人，怎麼向我一個撒瑪利亞女人要水喝呢？」因為猶太人和撒瑪利亞人沒有來往。 \end{tabularx} \\ \\ \relax
4:10 & \begin{tabularx}{0.7\textwidth}{X} 耶穌回答她說：「你若知道神的恩賜，和對你說『請給我水喝』的是誰，你早就會求他，他也早就會給了你活水。」 \end{tabularx} \\ \\ \relax
4:11 & \begin{tabularx}{0.7\textwidth}{X} 婦人對耶穌說：「先生，你沒有打水的器具，井又深，哪裡去取活水呢？ \end{tabularx} \\ \\ \relax
4:12 & \begin{tabularx}{0.7\textwidth}{X} 我們的祖宗雅各把這井留給我們，他自己和兒女以及牲畜都喝這井裡的水，難道你比他還大嗎？」 \end{tabularx} \\ \\ \relax
4:13 & \begin{tabularx}{0.7\textwidth}{X} 耶穌回答，對她說：「凡喝這水的，還要再渴； \end{tabularx} \\ \\ \relax
4:14 & \begin{tabularx}{0.7\textwidth}{X} 誰喝我所賜的水，就永遠不渴。我所賜的水要在他裡面成為泉源，直湧到永生。」 \end{tabularx} \\ \\ \relax
4:15 & \begin{tabularx}{0.7\textwidth}{X} 婦人對他說：「先生，請把這水賜給我，使我不渴，也不用到這裡來打水。」 \end{tabularx} \\ \\ \relax
4:16 & \begin{tabularx}{0.7\textwidth}{X} 耶穌對她說：「你去，叫你的丈夫，再到這裡來。」 \end{tabularx} \\ \\ \relax
4:17 & \begin{tabularx}{0.7\textwidth}{X} 婦人回答，對耶穌說：「我沒有丈夫。」耶穌說：「你說沒有丈夫是對的。 \end{tabularx} \\ \\ \relax
4:18 & \begin{tabularx}{0.7\textwidth}{X} 你已經有過五個丈夫，你現在有的並不是你的丈夫。你這話是真的。」 \end{tabularx} \\ \\ \relax
4:19 & \begin{tabularx}{0.7\textwidth}{X} 婦人對他說：「先生，我看你是一位先知。 \end{tabularx} \\ \\ \relax
4:20 & \begin{tabularx}{0.7\textwidth}{X} 我們的祖宗在這山上敬拜神，你們倒說，應當敬拜的地方是在耶路撒冷。」 \end{tabularx} \\ \\ \relax
4:21 & \begin{tabularx}{0.7\textwidth}{X} 耶穌對她說：「婦人，你要信我。時候將到，你們敬拜父，既不在這山上，也不在耶路撒冷。 \end{tabularx} \\ \\ \relax
4:22 & \begin{tabularx}{0.7\textwidth}{X} 你們所敬拜的，你們不知道；我們所敬拜的，我們知道，因為救恩是從猶太人出來的。 \end{tabularx} \\ \\ \relax
4:23 & \begin{tabularx}{0.7\textwidth}{X} 時候將到，現在就是了，那真正敬拜父的，要用心靈和誠實敬拜他，因為父要這樣的人敬拜他。 \end{tabularx} \\ \\ \relax
4:24 & \begin{tabularx}{0.7\textwidth}{X} 神是靈，所以敬拜他的必須用心靈和誠實敬拜他。」 \end{tabularx} \\ \\ \relax
4:25 & \begin{tabularx}{0.7\textwidth}{X} 婦人對他說：「我知道彌賽亞—就是那稱為基督的—要來；他來了，會把一切的事都告訴我們。」 \end{tabularx} \\ \\ \relax
4:26 & \begin{tabularx}{0.7\textwidth}{X} 耶穌對她說：「我就是，正在跟你說話呢！」 \end{tabularx} \\ \\ \relax
4:27 & \begin{tabularx}{0.7\textwidth}{X} 正在這時，門徒回來了。他們對耶穌正在和一個婦人說話感到驚訝，可是沒有人說：「你要甚麼？」或說：「你為甚麼和她說話？」 \end{tabularx} \\ \\ \relax
4:28 & \begin{tabularx}{0.7\textwidth}{X} 那婦人留下水罐，往城裡去，對眾人說： \end{tabularx} \\ \\ \relax
4:29 & \begin{tabularx}{0.7\textwidth}{X} 「你們來看！有一個人把我素來所做的一切事都說了出來，難道這個人就是基督嗎？」 \end{tabularx} \\ \\
[1ex]
\hline
\hline
\end{longtable}
日期:2009年6月21日

証道:伍渭文牧師

題目:喜出望外的邂逅

Surprised By Joy

經文:約翰福音 4:5-29

今天唸的經文,特別在福音書的記載(約翰福音 4:5-29),說到一個撒瑪利亞的婦人,她在中午的時候,來到一個地方打水,因為她想避開免得遇見其他人,在她背後竊竊私語.一般打水都是在傍晚日落才去的,中午烈日當空沒有人會去打水的.但這婦人就怕傍晚聚在井旁的人,因為很多人知道她的背景,她有五個丈夫,而現在有的並不是她的丈夫,所以她很怕.正午在那裡,不需要面對任何人.但這次在那裡,她不期而遇遇見耶穌基督,耶穌基督就和她談到活水的問題,接著就談到崇拜的問題,最後這位婦人發現她遇見了彌賽亞.

在這故事的後半部份的28節:「那婦人就留下水罐子,往城裡去,對眾人說,你們來看,有一個人將我素來所行的一切事,都給我說出來了,莫非這就是基督麼.眾人就出城往耶穌那裡去.」一個完全相反的反應:和耶穌相遇之先,她很怕見人；和耶穌相遇之後,她整個生命被扭轉了!她要到城裡面宣告,有一個將我素來所發生的事,所作的事都講清講楚,莫非這就是基督嗎?

今天經課好幾段經文──舊約以賽亞書6:1-8,約翰福音4章,詩篇29篇,羅馬書8:12-17都有一個很清晰的共同主題,就是我們的敬拜.有關主日崇拜是甚麼?為什麼進行這種的崇拜程序?我很久沒講了,我查閱過上一次講是兩年三個月之前的事了,我想我們今天再思崇拜的主題, 特別這兩年來了不少新朋友.讓我們先看看約翰福音書第四章這段經文,來思考有關崇拜的真理.

一. 上帝主動尋求人敬拜祂

第一,在耶穌和婦人的對話中,指出崇拜是上帝主動尋求人的.第23節:「時候將到,如今就是了,那真正拜父的,要用心靈和誠實拜他,因為父要這樣的人拜他.」特別留意「父要這樣的人拜他」,另外一個翻譯就是:父,上主,是尋求(seek)人去拜祂的.「邂逅」,如果你看說文解字,有兩個意思,一個是「不期而遇」,就如同這個婦人,在想不到的時候見到了耶穌；另一個的解釋就是,邂逅的「邂」和解除的「解」有關係的,就是「解除分離的痛苦而相會」.

對婦人來說,那一個的「邂逅」是「不期而遇」,但對耶穌來說,那個時刻正等候尋求和她見面.祂知道這個婦人,她不能透過肉慾的歡愉而得到真正的滿足.「凡喝這水的,還要再渴．人若喝我所賜的水就永遠不渴．我所賜的水,要在他裡頭成為泉源,直湧到永生.」(約 4:14)主耶穌基督尋求與這婦人相遇,帶領她認識祂就是彌賽亞,邂逅是解除分離的痛苦而與她相遇,表明上帝很想與人再一次和好,以致祂的兒子耶穌基督成為肉身,死在十字架上來親近人.

宗教或者類宗教的哲學,是人尋求上主,但我們的信仰,是上主主動尋求與人和好.在中文大學裡有兩個很有名的地標,都蠻有意思的,一個叫「天人合一」,位於新亞書院那裡,我想很多人都去過.另一個就是「崇基禮拜堂」,我認為都是很重要的地標.兩個都代表著不同的精神,對「救贖」的看法很不一樣的.新亞書院的校訓是「誠明」,「天人合一」可以說具體的將它的精神表露出來.當你看到「天人合一」這個建築傑作的時候,你見到一泓很淺的池水,當你望過去,這池水好像連接著吐露港,而吐露港又好像連接著整個天空.池水很寧靜,當你看到池水的時候,就好像看到海與天.

甚麼叫「誠明」呢?,王陽明說:「心明便是天理」,「夫萬事萬物之理不外於吾心」,「萬物一體」,這是「天人合一」的基本意思.所以儒家或者東方很多的宗教都追求甚麼?追求解脫甚麼呢?就是追求「除人慾,存天理」.當你去到「天人合一」的池水面前,你可以很安靜,很個人,很單獨的去默想,靈修,或者去體驗一個崇拜經驗,屬靈追求的經驗.你可將你裡頭的雜念滌淨,池水可以幫你淨化內心,當你望過去,好像可以與海,天一齊連結,這就可以有一個屬靈的經歷了!你是可以單獨去進行的,是你個人的問題,所有的問題是你自己的「人慾」,你除去這個「人慾」就可以存「天理」,那你就「天人一物,內外一理」.

我們和上帝根本就沒有分別的,我們裡面有上帝的種子,是一種傾向泛神的思想,一般東方的宗教都有這種的傾向.究竟是蝴蝶夢莊周,還是莊周夢蝴蝶呢?我們怎麼可以保存一滴水而不會讓它消失呢?回歸大海吧!如果你避免一滴水不消失的話,唯一的方法是甚麼?讓它回歸大海!這是佛家的講法.不過,你會發現回歸大海的水,就不再是以前的水了.所以你是作夢看見蝴蝶,或者蝴蝶作夢看見你,都沒關係,萬物為一體了,已經消解了你的位格(或情格,personhood)了.

但是崇基禮拜堂這個地標,我們往聖壇望,有一個很清楚的十字架,垂在玻璃幕牆,這個地方不是給個人靈修的,這個地方是給人集體來敬拜的.敬拜是一個集體的行為,因為我們認為「罪」不是個人的蒙蔽,被自己的人慾拂亂；罪是對一個很清楚的位格──上帝的拒絕,「罪」很清楚對另一個位格──人的情格的拒絕,甚至是藐視,將他者的人性眨抑,這就是罪.

愛上帝和愛人是連結的,這十字架很清楚給我們看到,橫的是人與人的關係,直的是人與上帝的關係,是位格之間的破裂 (alienation),拒絕.罪不是蔽障,人所需要的不是知道更多,因為知道了仍做不到.人所需要的是拯救──愛的拯救.這個十字架中間的那點,就好像一個中心,下面很長,好像召喚我們──Come──來到基督的面前；那個焦點就是基督,基督使人與上帝和好,又使我們彼此和好.

剛才所唸的以賽亞先書第六章,提到先知以賽亞覺得自己是嘴唇不潔的,當他不潔的時候,不能夠靠自已滅人慾,除不潔,人是不行的；仍是上帝使天使將一塊炭潔淨他的嘴唇.潔淨的能力,來自以賽亞之外的「他力」,這是一個「拯救」.這個思想是我們崇拜的基本.崇拜是上主尋求我們歸回祂.

二. 上主以真理和聖靈引導人敬拜祂

上主用甚麼而方法尋找人歸向祂呢?上主按啟示的真理與聖靈的光照引導人敬拜祂.所以在約翰福音第四章也講到我們要用「心靈和誠實」來敬拜.這「心靈誠實」不是個人主觀的虔敬.很多時候我們誤解敬拜是純粹個人主觀心靈的活動,但「誠實」原文是「真理」(in truth),在新譯本譯作按真理來敬拜,而「真理」就是耶穌基督,我們要在耶穌基督裡來敬拜.

細看我們的崇拜程序,實在蘊涵很多真理,很多聖經的話語,剛才唸了四段經課,每個禮拜我們都誦讀四段經課；我們唱的詩歌內容,都是充滿著聖經真理,這些真理不是我們想出來的,是上帝很清楚啟示給我們的.這個婦人的問題不是沒有虔敬,她也去敬拜的,不過她的敬拜乃是錯誤的敬拜.她與耶穌爭論到底在那裡敬拜上帝呢?你們在耶路撒冷,我們的祖宗教我們在撒瑪利亞城的一個叫基利心山,那裡有一個邱壇.耶穌怎麼解釋呢?耶穌解釋說,真的應該是在耶路撒冷,「因為救恩是從猶太人出來的」.

為何聖經如此清楚說要到耶路撒冷的聖殿去敬拜呢?因為那裡是一個預表,那裡是上帝與人相遇的地方.指定在那個地方就能帶出唯一的預表.聖殿是上帝與人相遇的地方,而耶穌就是那聖殿,但「時候將到,如今就是了,那真正拜父的,要按真理和靠聖靈來敬拜.」不再說要往某處敬拜了,祇要就按著聖經的真理,隨著聖靈的感動,在那一處都可以了!(因為實體來了,預表就成為過去,毋需墨守某一敬拜地點了.)

敬拜又按著「心靈」,是指「聖靈」.這位婦人她除了是錯誤敬拜之外,另外就是她的敬拜,虔敬是和她的道德,私人的生活割裂的,她可以很虔敬地去到撒瑪利亞城的邱壇敬拜上帝,但平時的生活卻很糟糕──有五個丈夫,現在有的並不是她的(這個是第六個),所以她不單是錯誤的虔敬,也是沒有被聖靈光照的虔敬.我有一串禮拜堂的鑰匙,其中只有一把是用來打開我辦公室的門,其他的都很相似的,我很虔敬地用任何一把鑰匙都未能打開的,你再怎麼虔敬也無法打開,只有那一把才能打開.所以敬拜是按照上帝的啟示──耶穌基督是神的救恩,耶穌基督就是「真理」──「我就是道路,真理,生命．若不藉著我,沒有人能到父那裡去.」(約14:6)因此,敬拜是很清楚地按照上帝的啟示.

同時敬拜也是在聖靈.「心靈」一辭在約翰福音裡的翻譯應該是「聖靈」,像約翰福音十四章,spirit指聖靈.耶穌離開我們之後,祂要差遣聖靈下來,使我們明白真理.因為真理需要聖靈的光照我們才能明白,沒有聖靈的光照和感動,我們不能完全明白真理.像剛才唸羅馬書那段經文,「聖靈與我們的心同證(Knowing together)我們是上帝的兒女．」(羅8:16)如非聖靈感動我們,我們沒可能明白耶穌是基督.所以「心靈誠實」,不是談及到主觀的,個人的虔誠,乃是我們按照聖經的教導來敬拜,在敬拜中聖靈將聖經的教導來感動我們的心,好叫我們能夠明白神的話語,所以這個十字架的符號就很清楚了──敬拜就是來到耶穌的面前,並接受祂的恩典.

三. 敬拜是上帝給我們的服事

另外,敬拜是上帝給我們的服事,上帝透過敬拜好叫我們能服事祂──我們唱歌讚美祂,歌頌祂──這是我們給上帝的服事,同時,也是上帝透過詩歌,禱告,講道,讀經給我們的服事.委實,上帝也在每一次崇拜的時候在我們心中工作.所以,每一次崇拜,我希望我們每一個人都能夠很安靜等候,期待上帝:求你對我說話；上帝阿!求你幫助我,叫我能明白你的旨意.

以賽亞先知不是第一次去聖殿,我們不知道這是第幾次,但這一次的經歷很特別,第一節提到那個時候是「烏西雅王崩的那年」,適時他的情緒很低落,這位王曾經是一位很好的王帝,帶來整個民族,國家,宗教的中興,但是因為他驕傲,擅自入聖殿燒香,因而受罰長了大痳瘋就死了.當他死的時候,舉國很低沉,先知覺得很絕望,好像沒有王的帶領了,但是上帝給他看見──你不用懼怕,君王有崩逝,政黨有輪替,但是上主仍坐著為王.這是針對他個人的需要,給他知道他不需要懼怕,上帝是掌控一切的──君王的心好像隴溝的水,在神的手中隨意流轉(箴21:1).上主針對他的需要給他知道,他不需要因王的崩逝而沮喪,上帝仍然坐著為王.

同時,在崇拜透過天使唱的詩歌──聖哉,聖哉,聖哉,神的榮耀充滿整個聖殿,使以賽亞覺得自己是嘴唇不潔的人,不配做先知,然後神用炭火潔淨他的嘴唇,而且不單潔淨他,還差遣他:「我可以差遣誰呢,誰肯為我們去呢」.這位先知被主潔淨之後,他願意被主使用,這也是我們崇拜中很重要的一個高峰.如果說東方的宗教或類哲學的高峰是「渾然忘我」,忘記自己.到底是蝴蝶夢見我呢?還是我夢見蝴蝶呢?這滴水去了大海就消解了,因此你不用太過執著,所有的事都是你太執著,你要有覺悟,達到渾然忘我的境界,沒有自己.但我們崇拜的高峰是「安身立命」.以賽亞先知找到自己的人生方向,上帝呼召他在民中──雖然民是很悖逆,他們耳朵發沉,不聽先知所說的話,但神說你要去.他找到自己人生的目標,這個就是崇拜很重要的一個高峰,就是來到上帝面前──我應該做甚麼呢?我在這裡,請差遣我.

四. 崇拜是失樂園的逆轉

最後我們看看這位婦人的經驗,這位婦人最後怎樣呢?經過與主相遇後,她很開心,她能面對她懼怕見的人.「有一個人將我素來所行的一切事,都給我說出來了」.將一個人素來所行的事說出來,是不簡單的,如果你不是為了(像龔如心遺囑官司)爭取一千億的財產,你不會將你素來所行的說出來,因為裡面隱藏著許多羞恥的事.能夠將素來所行的都講出來而不避諱的,一定有一種很大的力量勝過你的羞恥心,一定有很大的力量能夠使你不用懼怕,這一種是喜樂的力量──Surprised by joy的邂逅,是一個不期而遇找到基督耶穌──祂可以給我活水；有了這活水的滿足,我不再需要追逐肉體的歡愉.這位婦人從此放棄過去那一種生活,她不怕羞恥告訴人家:祂將我過去所講,所做的都說出來,我有五個丈夫,現在有的都不是我的.有哪一個婦人敢這樣子講呢?一定是她本身給一種很大的力量來駕馭她,抬升她,以致她無所懼怕!這個力量就是在崇拜中,我們與主相遇!

我們感謝上帝,崇拜是上帝呼召我們,然後我們回到上帝那裡,所以在我們崇拜的程序表一開始為什麼有「宣召」的項目呢?而最後有祝福呢?為什麼要有這兩項呢?「宣召」是上帝主動呼喚我們,呼喚我們敬拜上帝,今天我們宣召的經文就是詩篇95篇6-7節:「來阿,我們要屈身敬拜,在造我們的耶和華面前跪下.因為他是我們的　神．我們是他草場的羊,是他手下的民.惟願你們今天聽他的話」.但是你留意,在墮落的時候,上帝來到樂園裡頭尋找人,始祖亞當,夏娃躲藏起來,他們不敢見上帝,但我們唱一首詩歌「聖哉三一」來歌頌上帝.中間的彼此問安是表示集體的關係,我們整個被贖的群體一起來敬拜上帝,不是單獨個人的靈修.當然我們可以有個人靈修,但是來到禮拜堂這個地方,是集體敬拜的敬拜.我們有很多人坐在這裡,是一群被救贖的子民,來到父面前唱詩,讚美上帝,所以我們是回應上帝的召喚.

最後「祝福」是甚麼意思呢?祝福就是當我們敬拜完畢的時候,我們好像一個孩童,很 Child-like地來接受上帝美善的話語,上帝祝福的話是美善的,祂是美好的話語,我們要接受祂.但是在樂園裡面,人的墮落就是因為懷疑神的話語不是美善的.蛇如何誘惑始祖呢?上帝明明吩咐亞當,夏娃不要吃這個果子,你吃的日子你會死的,神的禁戒很清楚,這個「禁果」是表示甚麼呢?是分別善惡的果子,善惡的果子又是甚麼呢?在希伯來人的文化背景,分別善惡表示「自主獨立」.人不是早上被創造,晚上就墮落的.上帝已有很多次來到樂園裡,人知道上帝創造樂園整個環境,是給他們享受,是給他們守護,所有的食物都可以吃,只是那棵樹的果子你不可吃!這個命令不是永遠的,若是永遠的命令,總會有一天會失敗的,對不?而在樂園有另外一棵「生命樹」,人若經過一段時間順服上帝的話語不吃這個禁果,就可以分享「生命樹」.做老師出題目是不想學生failed(不及格)的,乃希望他能成功的,如果只有禁果而沒有生命樹的話,人就是在「無限」中等候墮落的機會,但不是的,上帝透過這棵樹表明與人的關係.那棵樹就學習信賴,順服上帝.那時候始祖很清楚知道他是上帝所創造的,上帝也有很多次來到樂園裡,他與上帝有一個很親密的關係,但是他竟然聽從了撒旦的說話──懷疑上帝.

撒旦說:你豈會死呢?不一定會亡的!你吃了之後,就會如同上帝一樣了.吃禁果不是偶然不慎,因為所有樹的果子都可以吃,除了那一棵以外.是人很清楚蓄意地向上帝說不,向上帝說:我拒絕你,我不信任你,是與上帝撕裂一個很長久的愛的關係.並不是偶然的犯罪,是人很明顯的拒絕,撕裂一個很長久的愛的關係,這就是墮落,絕對不是早上創造,晚上墮落的.

所以上帝的心很痛,因為人拒絕上帝,拒絕祂的愛.因此祂要耶穌基督來到世界上尋找人,拯救人.整個崇拜的設計,就說明我們得回那個失落的樂園,當宣召經文一唸的時候,我們立即回應詩歌,是很開心的.到結束時,我們就如同來到樂園裡的孩童般信賴上帝,接受祂的祝福,面對明天.

我們感謝上帝,因為是上帝主動的來尋求我們,拯救我們.我們看看自己過去信主的經歷,你是怎樣信主呢?是否你主動地尋求上帝呢?還是上帝透過很多的同學,你的朋友,甚至一些你不認識的人,向你主動分享福音?你知不知道,你信耶穌是有很多人在為你禱告,以至你能夠被聖靈光照,得著耶穌的救恩,這全都是上帝的恩典,我們感謝上帝!

讓我們一齊來禱告,感謝上帝的拯救:天父上帝,我們求你光照我們,不是我們尋求你,是你在耶穌裡面拯救我們,而且藉著這個的救贖,我們可以得著活水的泉源,我們不再懼怕去面對明天,如同那位婦人不再懼怕去面對以前所懼怕的人,而且她走到人群當中,宣揚這個好的福音,我們感謝你,奉耶穌的名禱告,阿們!



\newpage

\section{為甚麼膽怯呢}
\label{sec:2}
\textbf{林豪恩先生 - 為甚麼膽怯呢? Why You Are So Fearful?}
\newline
\newline
\hyperref[sec:1]{< < < PREV SERMON < < <}
~
\hyperref[sec:toc]{[返目錄]}
~
\hyperref[sec:3]{> > > NEXT SERMON > > >}


馬可福音 4:35-41
\newline
\begin{longtable}{cl}
\hline
\hline
章節 & 經文 (和合本修訂版)\\
\hline
4:35 & \begin{tabularx}{0.7\textwidth}{X} 那天晚上，耶穌對門徒說：「我們渡到對岸去吧。」 \end{tabularx} \\ \\ \relax
4:36 & \begin{tabularx}{0.7\textwidth}{X} 門徒離開眾人，耶穌已在船上，他們就請他一同去；也有別的船和他同行。 \end{tabularx} \\ \\ \relax
4:37 & \begin{tabularx}{0.7\textwidth}{X} 忽然狂風大作，波浪打入船內，以致船灌滿了水。 \end{tabularx} \\ \\ \relax
4:38 & \begin{tabularx}{0.7\textwidth}{X} 耶穌在船尾上，枕著枕頭睡覺。門徒叫醒他，說：「老師！我們快沒命了，你不管嗎？」 \end{tabularx} \\ \\ \relax
4:39 & \begin{tabularx}{0.7\textwidth}{X} 耶穌醒了，斥責那風，向海說：「住了吧！靜了吧！」風就止住，大大平靜了。 \end{tabularx} \\ \\ \relax
4:40 & \begin{tabularx}{0.7\textwidth}{X} 耶穌對他們說：「為甚麼膽怯？你們還沒有信心嗎？」 \end{tabularx} \\ \\ \relax
4:41 & \begin{tabularx}{0.7\textwidth}{X} 他們就非常懼怕，彼此說：「這到底是誰？連風和海都聽從他。」 \end{tabularx} \\ \\
[1ex]
\hline
\hline
\end{longtable}
日期:2009年6月28日

証道:林豪恩先生

題目:為甚麼膽怯呢? Why You Are So Fearful?

經文:馬可福音4:35-41節

一．    經文

35　當那天晚上,耶穌對門徒說:「我們渡到那邊去吧.」36　門徒離開眾人,耶穌仍在船上,他們就把他一同帶去；也有別的船和他同行.37　忽然起了暴風,波浪打入船內,甚至船要滿了水.38　耶穌在船尾上,枕著枕頭睡覺.門徒叫醒了他,說:「夫子!我們喪命,你不顧嗎?」39　耶穌醒了,斥責風,向海說:「住了吧!靜了吧!」風就止住,大大地平靜了.40　耶穌對他們說:「為甚麼膽怯?你們還沒有信心嗎?」41　他們就大大地懼怕,彼此說:「這到底是誰,連風和海也聽從他了?」

二．    經文分析

記載在馬可福音第4 章35-41節的這個「耶穌平靜風和海」的事件,雖然簡短,但故事相當完整:有時間:當那天晚上；有主要人物:耶穌,門徒；有地點:海上；有緊張的情節:渡海的時候遇上風浪,門徒快被淹死的危險時刻,耶穌平靜了風浪,解除了危機；有結束解說(debriefing):耶穌的反問及門徒的反省.

除此之外,作者用了引人注意的對比來描述這個精簡的故事:37節「忽然起了暴風」的「暴風」直譯是大大的風；耶穌行動後,39節描述「風就止住,大大地平靜了」；41節故事結束時作者描述門徒對耶穌的反應是「大大地懼怕」.作者三次使用「大大地」,對我們閱讀這個故事有甚麼啟發呢?讓我們暫且保留這個問題,先回到今天的講道

題目:為甚麼膽怯呢?

三．    釋經與應用

1．    「為甚麼膽怯呢?」作為一個答辯的邀請膽怯,是人類共同的情緒之一,差不多所有人都曾經有過膽怯的經驗,所以,也能夠明白及體會膽怯是怎樣一回事.甚麼事情能夠令你膽怯呢?你在怎樣的狀況下會有膽怯的反應呢?是甚麼原因令你膽怯呢?

有些情況下,膽怯可能因為無知,出於高估了事情的威脅.幾個星期之前,崇基團契有三十位同學前往柬埔寨服務,其中一群服務對象是帶有HIV病毒的孤兒.有些同學的印象中,接觸愛滋病患者要穿上保護衣物,戴上手套.對於要赤手空拳去服務這些孤兒,這些同學難免感到膽怯.後來,有關人士向同學解釋了愛滋病病毒的傳染方法和途徑,於是,起初感到膽怯的同學都放心了.到達孤兒院看見小孩子時,同學們毫無顧慮地拖著小孩子的手,與他們玩遊戲,甚至抱在懷中,每天道別時都難捨難離.

令我們膽怯的情況之中,會否有些是因為我們的誤解,因為我們高估了事情的威脅而自己嚇自己呢?

如果我們把這個問題問門徒:「你們膽怯,是否由於無知呢?你們會否高估了風浪的威脅呢」.門徒可能如此回答:「我們膽怯,不是無知,反而是因為我們知道得十分清楚這種風暴的威脅有多厲害.」持平地看,我們其實沒有資格去挑戰門徒的判斷,因為他們當中有不少人本來是當地的漁夫,例如:西門彼得,安得列,雅各和約翰,他們對於當地的天氣狀況及風浪的威力有親身的體會.能夠令這些飽歷風浪的漁夫膽怯的,一定是威力極強的暴風.事實上,由於加利利湖處於盆地,四周是高山,高山之外是地中海.當地中海冷空氣的氣流越過高山,與湖面的熱空氣的氣流交匯時,便會產生狂風巨浪,這就是連漁夫也無法預測,以及令他們感到膽怯的風浪.

讓我們問一個似乎是太過簡單的問題:暴風巨浪為甚麼令門徒膽怯呢?也同樣問自己類似的問題:那些令我們膽怯的事情為甚麼能夠令我們膽怯呢?對於門徒來說,海上搖船遇上暴風巨浪,他們推想到的後果是:「我們喪命」.當我們面對令我們膽怯的事情的時候,我們推想到甚麼後果呢?是事情本身令我們膽怯,還是我們的推想令我們膽怯呢?

有一位朋友非常害怕被裁員,他推想:沒有工作,就沒有收入；沒有收入,就沒有錢給家人；沒有錢給家人,就成為家庭的負累；成為家庭的負累,就會被趕走；被趕走,無處藏身；病了,沒有錢醫病,也沒有人照顧,．．．因此,想到裁員,便膽怯.有時,令我們膽怯的,可能是類似的過度推想,恐嚇得我們膽戰心驚.然而,令門徒膽怯的,又是否如此呢?讓我們重申:不要懷疑門徒的判斷,因為他們是經驗豐富的當地漁夫,而且作者明言風暴翻起的波浪已經湧入船艙,船裡差不多入滿了水,這正是水浸眼眉,生死邊緣的危險時刻.

「喪命」,是風暴令門徒膽怯的原因.「喪命」,即是失去生命,本來擁有的東西被奪走.讓我們重提剛才的問題:「那些令我們膽怯的事情為甚麼能夠令我們膽怯呢?」那些事情對我們的威嚇,會否也是把我們本來擁有的東西強行奪走,令我們失去所擁有的?金融風暴和股市下跌令人膽怯的地方是它奪走人們的財富；疾病令人膽怯的地方是它奪走我們的健康；分離令人膽怯的地方是它奪走我們的關係；衝突令人膽怯的地方是它奪走我們的安全感；極權令人膽怯的地方是它奪走我們的自由；失敗令人膽怯的地方是它奪走我們的成功感；誹謗令人膽怯的地方是它奪走我們的聲譽.有多少事情可以成為我們的威嚇,可以令我們膽怯,可能與我們以為自己擁有多少,以及如何看這些擁有有關.

讓我們來思考一個較為爭議性的例子:如果有人偷了你與配偶的親密照片,然後向你勒索,聲言若你不就犯,就把照片公開.你會受威嚇嗎?你會感到膽怯嗎?若竊賊把相片公開,為何令你膽怯呢?

據台灣《聯合報》報道,台灣苗栗縣劉姓老農與妻子追隨時尚自拍性愛照欣賞,不料相機遭竊並被竊賊以公開裸照勒索180萬元(新台幣,下同),老夫婦表示"又不是外遇照,怕什麼!"除報警外還回嗆竊賊說:"給誰看都沒關係."

雖然有些事情會奪走我們的擁有,令我們失去一些東西；不過,會否令我們感到膽怯,還在乎我們如何看待這些擁有,以及如何面對失去.

「為甚麼膽怯呢?」如果這是一個需要回答的問題,門徒的答案很簡單:「我們要死了」.生命,已經是一個人最後可以失去的東西,怎可能不膽怯呢?「你又為甚麼膽怯呢?」如果這是一個需要你回答的問題,你一定能夠找出令你膽怯的理由,合理化你的膽怯.不過,如果我們只停留在回答問題,合理化自己的膽怯,那麼,對於解決膽怯的困擾並沒有太大的幫助.

2．    「為甚麼膽怯呢?」作為一個體會關係的邀請第40節耶穌對門徒說:「為甚麼膽怯?你們還沒有信心嗎?」耶穌對門徒說的話,似乎把膽怯與信心連上關係.信心與膽怯有甚麼關係呢?有信心就不會膽怯嗎?為甚麼有信心就不需要膽怯?信心如何令人不膽怯?在這裡,「信心」所指的是甚麼?是在暴風忽然來臨時,有信心耶穌必然能夠平靜風浪?是在生命危在旦夕時,有信心耶穌必然會出手拯救?是當我們遇到膽怯的事情的時候,有信心耶穌必然會為我們擊退令我們膽怯的事情?耶穌平靜風和浪的事件,好像在表達這樣的信息；然而,若我們仔細閱讀經文,又不是如此.

讓我們換一個角度再看一看經文,特別留意經文的交叉平衡及對比.

35當那天晚上,耶穌對門徒說:「我們渡到那邊去吧.」    41他們就大大地懼怕,彼此說:「這到底是誰,連風和海也聽從他了?」36門徒離開眾人,耶穌仍在船上,他們就把他一同帶去；也有別的船和他同行.    40耶穌對他們說:「為甚麼膽怯?你們還沒有信心嗎?」37忽然起了暴風,波浪打入船內,甚至船要滿了水.    39耶穌醒了,斥責風,向海說:「住了吧!靜了吧!」風就止住,大大地平靜了.38耶穌在船尾上,枕著枕頭睡覺.38門徒叫醒了他,說:「夫子!我們喪命,你不顧嗎?」

第一項平衡及對比:(35)故事以耶穌的說話作為開始,記錄了耶穌對門徒作出一項指示；(41)故事以門徒的說話作為結束,記錄了門徒的自我發問.

第二項平衡及對比:(37)暴風突然來臨,波浪翻騰,海水湧入船；(39)耶穌斥責風,吩咐海,風馬上消失,海立即平靜.

第三項平衡及對比:(38上)耶穌在風浪中安睡；(38下)門徒在風浪中喊救命.

當門徒在面對風浪沒頂的威脅的時候,耶穌仍在安睡,風浪仍未平靜,門徒尚不知道耶穌會否出手相救,也不知道耶穌有沒有能力出手相救.如果有信心的意思是門徒相信他們必然被耶穌所救,那麼,按當時的狀況,門徒膽怯是合情合理的反應,因為威脅仍未解除.

「為甚麼膽怯?你們還沒有信心嗎?」耶穌對門徒說這句話到底有甚麼意思呢?所指何事呢?讓我們閱讀這故事的第四項平衡及對比:(40)「耶穌對他們說:「為甚麼膽怯?你們還沒有信心嗎?」；(36)「門徒離開眾人,耶穌仍在船上,他們就把他一同帶去；也有別的船和他同行.」按這對平衡提供的線索,第40節耶穌的問題的回應在36節所描述的狀況,就是耶穌與門徒同在一條船上.按前文後理,耶穌彷彿如此對門徒說:「我與你們同坐一條船上,你們還沒有信心嗎?為甚麼膽怯呢?」從故事的敘述可見,耶穌所說勝過膽怯的信心,並不是因為耶穌必然解除令我們膽怯的事情,而只是因為有耶穌同在而已.耶穌也和我們說:「為甚麼膽怯呢?」耶穌並不期望我們向他解釋膽怯的理由,而是邀請我們與他同行,與他同坐一條船.

在一列歐洲的長途火車上,有一小孩玩耍,但是不見大人照顧.旅客都覺詫異,終於有一旅客忍不住問:「孩子,你一個人怕不怕?」小孩回答說「不怕,開火車頭的是我爸爸.」(網絡故事)

有一首兒童詩歌,名為「有主在我船上」是這樣說的:

有主在我船上,我就不怕風浪,不怕風浪,不怕風浪,有主在我船上,我就不怕風浪,直到安抵天家.搖啊!搖回家,搖啊!搖回家,有主在我船上,我就不怕風浪,直到安抵天家.

3．「為甚麼膽怯呢?」作為一個讓心靈震撼的邀請按經文的記載,門徒沒有直接回答「為甚麼膽怯?你們還沒有信心嗎?」這個問題.故事結束的時候,門徒仍然驚魂未定,「大大地懼怕」(另有翻譯為「大起了敬畏的心」.從門徒接著的彼此對問:「這到底是誰,連風和海也聽從他了?」可見,門徒現在的懼怕,與之前的懼怕並不一樣.之前懼怕的是「大大的風浪」；看見耶穌令風浪「大大地平靜」後,現在門徒「大大的懼怕」的,是他們跟隨的耶穌.「大大的風浪」固然令人懼怕,但能夠令風浪「大大地平靜」的耶穌,更是令人敬畏.事件震撼著門徒的心靈,他們不懂回應,也未及理解,只有彼此驚嘆:「這到底是誰,連風和海也聽從他了?」

對著大自然的鬼斧神工,壯麗山河,絢爛晚霞,又或者地動山搖的地震,瞬雷不及掩耳的海嘯,很多人的心靈都被震撼,敬畏的感覺油然而生.

(網絡文章)我由於工作關係,長年住在合歡山農場.有一個冬天早晨起來,突感一種與往昔不同的光亮,清新而柔和,開門一看－下雪了!

雪光雖然刺眼,但心中的驚異喜悅,使整個人興奮起來.大地一片潔白寂靜,昨日的一切景觀都改變了,這一切非人所能作為或掌控的事實就呈現在眼前.

平時我們常強調「人定勝天」,這只證明人類的驕傲,及自以為是,井底之蛙的看法.老天在一夜之間無聲無息默默地將幾千萬噸的水堆在你面前,不用引擎,不用電力,就這樣一片一片輕似鵝毛的雪花成就了他的傑作.偉哉!自然的力量何其大!從這聯想到九二一大地震不也在幾秒鐘瞬間摧毀了幾萬戶家園嗎?

這兩者給我的啟示是:人不可以一天忙到晚將自己的寶貴生命在混沌中與分秒間流失,要有一些時間靜下來回到心靈中尋求一些答案.試想,上天都要藉著下雪,強迫大地沈寂而得著養息,等待來春萬象更新,何況我們人呢?冬雪融化滋潤孕育大地能量以供養萬物生長,當然這恩典給好人也給壞人,這是老天爺的公義.我們人為萬物之靈不亦應向老天學習?

老農如我,那日站在這雪窗下遙望天地蒼茫,內心更加充滿對大自然的謙卑與感恩.

大自然給我們的震撼,我們對大自然的敬畏,有沒有衝擊我們的心靈呢?有沒有令我們有所啟發呢?如果大自然本身已經足以令我們震撼,足以令我們敬畏,足以令我們謙卑,足以令我們感恩,那麼,大自然的創造者豈不更令我們大大地震撼,大起敬畏的心,更加謙卑,更加感恩嗎?對於門徒來說,暴風巨浪令他們懼怕,平靜暴風巨浪的耶穌豈不更值得他們敬畏呢?

「為甚麼膽怯呢?」這個問題的重要性不在於答案,而在於震撼我們的心靈,在於刺激我們去體會所信的是誰,在於引發我們對創造者的敬畏.對一個更偉大者的敬畏或懼怕,能夠解除或減輕對其他威嚇的懼怕嗎?故事的結局是開放的,耶穌問了門徒問題:「為甚麼膽怯?你們還沒有信心嗎?」沒有回答；門徒彼此互問:「這到底是誰,連風和海也聽從他了?」也被有回答.心靈的震撼,對創造者的敬畏,往往不是言語所能描述的.雖是言語不能描述,對於當事人來說,卻比言語能描述的更真實和深刻.故事以一連串的問號作結,留下大量空間給讀者去思考及體會.

四．    總結

最後,我相當肯定這篇講章並沒有解除作為聽眾的你的膽怯.你聽完之後,仍有不少事情足以令你膽怯.這不是講員的失敗,也不是聽眾的失敗,這是事實.大家有沒有留意按馬可福音的記載,門徒經歷這件膽怯的事情之前正在做甚麼呢?他們在聽耶穌講道!門徒被耶穌親自教導了一整天,面對風浪的威脅時,仍然膽怯,何況是我們呢?門徒的經驗告訴我們,信仰有知識的部份,所以我們要聆聽,閱讀,研習；同時,信仰也有實踐的部份,而且唯有在實踐之中,才能真實地經歷信仰.各位,信仰的精彩之處不在講章裡面,卻在你的經歷,在你膽怯之處發現上帝的同在.放心向前行吧!



\newpage

\section{誰滾開那石頭}
\label{sec:3}
\textbf{伍渭文牧師 - 逆轉乾坤--誰滾開那石頭?The Great Reversal--Who Moved the Stone? 馬可福音 16: 1-8}
\newline
\newline
\hyperref[sec:2]{< < < PREV SERMON < < <}
~
\hyperref[sec:toc]{[返目錄]}
~
\hyperref[sec:4]{> > > NEXT SERMON > > >}


馬可福音 16:1-8
\newline
\begin{longtable}{cl}
\hline
\hline
章節 & 經文 (和合本修訂版)\\
\hline
16:1 & \begin{tabularx}{0.7\textwidth}{X} 過了安息日，抹大拉的馬利亞、雅各的母親馬利亞，和撒羅米，買了香料，要去膏耶穌的身體。 \end{tabularx} \\ \\ \relax
16:2 & \begin{tabularx}{0.7\textwidth}{X} 七日的第一日清早，太陽出來後，她們來到墳墓那裡， \end{tabularx} \\ \\ \relax
16:3 & \begin{tabularx}{0.7\textwidth}{X} 彼此說：「誰要替我們把石頭從墓門滾開呢？」 \end{tabularx} \\ \\ \relax
16:4 & \begin{tabularx}{0.7\textwidth}{X} 她們抬頭一看，看見石頭已經滾開了，原來那石頭很大。 \end{tabularx} \\ \\ \relax
16:5 & \begin{tabularx}{0.7\textwidth}{X} 她們進了墳墓，看見一個年輕人坐在右邊，穿著白袍，就很驚奇。 \end{tabularx} \\ \\ \relax
16:6 & \begin{tabularx}{0.7\textwidth}{X} 那年輕人對她們說：「不要驚慌！你們尋找那釘十字架的拿撒勒人耶穌，他已經復活了，不在這裡。來看安放他的地方。 \end{tabularx} \\ \\ \relax
16:7 & \begin{tabularx}{0.7\textwidth}{X} 你們去，對他的門徒和彼得說：『他要比你們先到加利利去，在那裡你們會看見他，正如他從前所告訴你們的。』」 \end{tabularx} \\ \\ \relax
16:8 & \begin{tabularx}{0.7\textwidth}{X} 於是她們出來，從墳墓那裡逃走，又發抖又驚訝，甚麼也沒有告訴人，因為她們害怕。 \end{tabularx} \\ \\
[1ex]
\hline
\hline
\end{longtable}
日期:2009年7月5日

証道:伍渭文牧師

題目:逆轉乾坤--誰滾開那石頭?The Great Reversal--Who Moved the Stone? 馬可福音 16: 1-8

經文:馬可福音 16:1-8

誰滾開那石頭?是今天題目的副題.Who Moved the Stone?也是一本書的名字,中文繙作墓石懸謎.1930年英國一位著名律師和專欄作家Frank Morison所寫的護教名著.作者本意是把耶穌受宙開始,受死和復活前後七天的經過,用律師的精密思考,收集證据,模擬在法庭提出,希望推翻耶穌確從死裡復活的說法.但當作者進入認真的資料搜集後,所有的証据剛剛相反,都指向耶穌從死裡復活的合理推論.諸多証据中,空的墳墓最為有力.墓石滾開,墳墓空了,屍体不見了.羅馬兵丁豈不是在墓旁看守,怎能說睡著的時候被人前來偷走呢?誰人有這麼大的氣力把石頭滾開?因為封墓就容易,推下去就可以,石頭推下去就落在陷下去的坑,鎖緊了.輥開陷在坑裡的石頭就不是那麼簡單了,儼然是一項大工程,如果兵丁睡了,也會被弄醒.而耶穌墳墓空了,他復活了,是非常有顛覆性的傳說,可以動搖羅馬政權.要是可以把屍體拿出來,很快就把危機平伏下來.但耶穌己從死裡復活的信息,不斷傳開,猶太人或羅馬政權一直找不到屍體.這宗教運動愈來愈壯大,從耶路撒冷,撒瑪利亞,席捲歐洲.

誰給我們把石頭從墓門滾開呢?

Who will roll away the stone for us from the entrance to the tomb?

這是三位婦女抹大拉的馬利亞,和雅各的母親馬利亞並撒羅米的問題.這問題和Morison的書名誰滾開那石頭?多了三個字:給我們.誰滾開那石頭?是純粹客現的推敲,與自已沒有個人的關係.就算找不到答案,對很多人來說衹不過是眾多歷史懸案中的一個個案.當然,對Morison來說,這是耶穌復活合理推論的明證.但誰給我們把石頭從墓門滾開呢?是非常實際迫切的個人關懷了.三位婦女帶著香膏一清早就到耶穌的墓園,就是希望膏抹耶穌的屍首,完成入殮的禮儀,因耶穌從十字架取下來時快安息日,設有時間作這儀式.我們若看過最近的奧斯卡最佳外語片禮儀師之鳴奏曲,就明白膏抹離世者是何等莊嚴神聖的儀禮了.但她們其實早已知道(可15: 46, 47),石頭己封了墳墓,帶著患得患失的心情,來到耶穌的墓前試試看.這班婦女因著對耶穌濃烈的情愛,不是那麼客觀分析實際的情況,要是墳墓被石頭封死了,三人合力也不可能滾開石頭哩!有時愛就是不太理性.她們希望石頭滾開,能進入墳內膏抹耶穌.香膏己帶來了,是她們最珍貴,最喜愛的香膏.

這三位婦女的問題:誰給我們把石頭從墓門滾開呢?正是復活節的信息.基督的復活沒有除去墳墓,但移開了緊封墳墓的石頭.

1. 復活使我們無懼罪惡和死亡的勢力

復活並沒有取消罪惡和死亡,乃叫我們不再懼怕罪惡和死亡.對潘霍華來說,復活就是把封密世界打開的風,這風滾開了封閉墓門的石頭,並把墳墓內的死寂絕望的氣氛吹走,這嶄新淨化一切的風,己吹進現今的世界.現今的世界對當時的潘霍華來說,是納粹軍國主義當道,是猶太人遭到種族清洗:正沖涼的時候,毒氣就開動,然後屍首就搬進火爐焚化.現今的世界對潘霍華來說,是沒有日光的窄小囚室,是他自願為信仰為真理而被囚的.1939年,當希特勒的恐怖統治在德國肆虐,戰爭已經在歐洲開始之時,他正在美國巡游講學,並公開抨擊納粹主義.他本可以留在美國,但是,一邊享受異國的自由和安全,一邊隔著遼闊海洋譴責納粹,對他來說無疑於靈魂犯罪.他在給友人的信中說:"我來美國實在是一個錯誤.假如此時不分擔同胞的苦難,我將無權參加戰後的重建."這嶄新淨化一切的風吹進他的囚室,藉鼻孔進入他身體每一個細胞.這風使他在獄中所寫的書簡,鼓舞著無數因為敢於說真理而失去自由的人.一位作家寫下獄中書簡讀後感說:學會敬畏,學會謙卑,學會在危險和恐懼中坦然面對厄運,並對未來保持樂觀的信心,學會以發自內心的坦誠與人相處,這就是潘霍華聖徒人格的啟示.......潘霍華始終平靜地對自己微笑,彷彿他是一個從一出生就只會笑的怪物.......他之所以能夠保持絕境中的希望,發出開朗而令人驚奇的笑,就在於他始終信仰著.他的文字不只是用來表達悲哀罷了,更是表達歡樂的,其份量沉得比任何絕望都豐富.正如他在臨刑前向獄友告別時所說:"這,就是終點.對我來說,是生命的開端."寫這段文字的作者,就是劉曉波先生,零八宣言的起草人和發起者,現在被當局囚禁在不知名的地方.

復活如風靜靜吹襲,逆轉乾坤

耶穌降生之時,有異星橫空,吸引遠方的賢士千里迢迢來到耶路撒冷,尋找新生王.耶穌誕生當晚,伯利恆上景星明亮似晝,天使歌聲響澈雲霄:在至高之處榮耀歸於上主,在地上平安歸於他所喜悅的人.耶穌死的時候,從午正至申的三個小時,遍地為黑喑所吞噬,日頭變黑；氣斷時,聖殿分隔神人的幔子裂開,地大震動,磐石也崩裂,墳墓開了.但復活的清晨,溫度,氣壓如常,墓園的草如常沾滿了露水,雀鳥鳴叫,好像甚麼事都沒有發生過,祗是墓石滾開了,若探頭往墓裡頭看看,才見到裹屍布像人形放在石棺上,耶穌屍體不見了.

這嶄新淨化一切的風,雖然不動聲息吹進現今的世界,但逆轉乾坤.神學家巴爾特(Karl Barth)說得好:復活的信息告訴我們,我們的仇敵,罪惡,詛咒和死亡已被打敗.他們已徹徹底底沒有作為.他們也許詐作未知鹿死誰手,勝負尚未分曉,我們也不應輕敵,但已不再懼怕它們了(The Easter message tells us that our enemies, sin, the curse and death, are beaten. Ultimately they can no longer start mischief. They still behave as though the game were not decided, the battle not fought; we must still reckon with them, but fundamentally we must cease to fear them any more).

2. 復活使我們成為門徒,宣講使人得釋放的福音

復活的信息,召喚我們去宣揚福音,使被擄的得釋放,被壓制的得自由.德國一位神學家莫特曼在其被釘死的上帝(The Crucified God)一書中說:基督藉受苦把世上的苦難歷史,連結到永生的上主,並他創造生命的聖潔的公義和正義上；也讓上主認同被暴力所壓制的人(Through his passion Christ brings the passion history of this world the eternal fellowship of God, and the divine justice and righteousness that create life; and he identifies God with the victims of violence).

環顧今天的世界,罪惡依然猖獗,不公義的事到處可見,但因為耶穌已從死裡復活,我們不會氣餒.最近北京大學有一位孫教授指上訪者有精神病的言論就是一例.引起的民憤越鬧越大.接連幾天有數百民眾在北大門前聚集抗議,要求與孫及北大校長見面,並要他們就孫的言論作解釋.据報北京每年有數百萬來自各地的上訪者.2003年,國家信訪局局長公開承認上訪者「80\%以上有道理,或有一定實際困難和問題應予解決」.正如有評論者指出,由於國家信訪局的權力有限,只能把問題轉交各地處理,而各地諸侯的權力層官官相衛,官商相衛,就使上訪者的問題鮮有獲得解決.其實,上訪者到北京告御狀時,各地掌權者己派人截訪,毆打.上訪者若掙扎到得以面見信訪局人員,問題會由信訪局交回各地,也仍然沒有解決.腐敗體制就是這樣進行著惡性循環.我想孫教授也許是對的,上訪民眾可能百分之九十以上精神有問題,因為他們的冤情得不到公義的申訴,長期鬱結而精神受了刺激,但罪不在他們哩.復活的信息,召喚我們站在這些被邊沿化,被人遺忘,受壓制的人的一邊,並堅信公義最終獲得超雪,因基督已從死裡復活.一股嶄新純化一切的風,己吹進現今的世界,我們不會放棄.

3. 復活使我們回到加利利,從新上路

``

那少年人(天使)對她們說:不要驚恐!你們尋找那釘十字架的拿撒勒人耶穌,他已經後活了,不在這裡.請看安放他的地方.你們可以去告訴他的門徒和彼得說:他在你們以先住加利利去.在那裡你們要見他``(可16: 6, 7).

若我們仔細查看整卷福音書的記載,便發現不是門徒編造復活的故事,乃復活的事實使他們成為門徒.加利利是耶穌與門徒三年多一起生活的地方.在加利利的海邊,耶穌呼召彼得,雅各,約翰,安得烈為門徒.在加利利的迦拿的婚宴,他在門徒面前行了第一件神蹟,使水變酒.在加利利的黑門山,耶穌在彼得,雅各,約翰面前變像,顯露他的榮耀,使他們極其懼怕,伏在地上.相約在加利利,是使他們記起主昔日的教訓,再踏上門徒之路.

今天你帶著甚麼來問這問題?

今天早上,也許你感到挫折,沮喪,絕望,心靈被困,像被死亡的力量綑縛住,覺得沒有出路,看不到半點亮光.誰給我們把石頭從墓門滾開呢?也許我們不會帶著香膏問這問題,我們可能帶著對上主的疑問,埋怨,甚至憤怒來問這問題.為何上主你讓這些不幸的事臨到我呢?為何我一直信任,提攜的人可以這樣在背後刺我一刀呢?為何主讓我失去這麼好的職業呢?為何我會把一生買了雷曼債卷,血本無歸?你是否感到絕望,像被困在墳墓中,看不到光明?

誰給我們把石頭從墓門滾開呢,讓我們不被死亡的氣味所窒息? 誰給我們把石頭從墓門滾開呢,讓我們看到墓外的陽光? 誰給我們把石頭從墓門滾開,讓我們可以直視死亡而不懼怕?

不過當你問:誰給我們把石頭從墓門滾開呢的時候,請你再看清楚,"石頭已從墓門滾開了,那石頭原來很大,她們抬頭一看,卻見石頭已經滾開了."(可16: 4)

讓我們回到我們的加利利,回顧昔日主與我們與親近的日子,他怎樣改變我們的生命,呼召我們跟隨他,我們也經歷過類似登山變像的屬靈高峰經歷.讓我們再一次回到加利利,再次立志作主的門徒.

讓我們今早深深呼吸一下,再次接受那靈風,那第一個復活主日刮起的風,這風滾開了封閉墓門的石頭,帶走墓裡的幽暗.這風在五旬節的時候,成為大響,落在門徒身上:〝忽然,從天上有響聲下來,好像一陣大風吹過.〞(使2: 2)使他們得著能力,為福音作見証.阿們.



\newpage

\section{靈機一觸}
\label{sec:4}
\textbf{鄧瑞強博士 - 靈機一觸}
\newline
\newline
\hyperref[sec:3]{< < < PREV SERMON < < <}
~
\hyperref[sec:toc]{[返目錄]}
~
\hyperref[sec:5]{> > > NEXT SERMON > > >}


日期:2009年7月12日

証道:鄧瑞強博士

題目:靈機一觸

經文:馬可福音5:21-43

各位弟兄姊妹,主內平安.

前幾日,吃了一味夏天常吃的sung 菜:「苦瓜釀豬肉」.

單吃苦瓜,可能「齋」了一點；單吃豬肉,可能「漏」了一點.兩者合璧,剛好,不齋不漏.再者,苦瓜令豬肉多了甘味,豬肉使苦瓜多了肉味,兩者的味道都豐富了.今日講道的經文,是馬可福音5:21-43,這段經文的特別之處,是一個故事「釀」在另一個故事裡,一個故事的味道滲入到另一個故事的味道裡,使兩者帶出一更豐富的味道來.

馬可福音5:21-43節的故事,是這樣開始的.

5:21    耶穌坐船又渡到那邊去,就有許多人到他那裡聚集；他正在海邊上.5:22    有一個管會堂的人,名叫睚魯,來見耶穌,就俯伏在他腳前,5:23    再三地求他,說:「我的小女兒快要死了,求你去按手在她身上,使她痊愈,得以活了.」5:24    耶穌就和他同去.有許多人跟隨擁擠他.

一個管會堂的人,是宗教上的敬虔人,通常是被社群尊重的人,有一點地位,處身在社會的中心位置.這個人叫睚魯,他遇到一件常人解決不了的事情.他的小女兒快死了.醫生對他女兒的病束手無策.他求耶穌去醫治這小女孩.睚魯求耶穌使她的女兒「痊愈」.聖經這裡的用字有點不尋常,「痊愈」這個字的希臘文常用來指「獲得救贖」.這是一個病得醫治的故事,也是一個人獲得救贖的故事.

耶穌回應這個社會上層人物的呼求,耶穌和他一起回家.

這個故事發展下去時,有另一個故事發生了.一個故事「釀」在另一個故事中.兩個故事互相呼應,去說明神的大愛何等豐富.

另一個故事是這樣的:

5:25    有一個女人,患了十二年的血漏,5:26    在好些醫生手裡受了許多的苦,又花盡了她所有的,一點也不見好,病勢反倒更重了.5:27    她聽見耶穌的事,就從後頭來,雜在眾人中間,摸耶穌的衣裳,5:28    意思說:「我只摸他的衣裳,就必痊愈.」5:29    於是她血漏的源頭立刻乾了；她便覺得身上的災病好了.

這裡有一個患了十二年血漏病的女人.這是一個可憐人.經文說,她花盡一切的錢財,在醫生手中受了許多苦,人間的財富及醫學的知識都不能幫到她.患這種病的人,在宗教禮儀上是不潔淨的人,受到社會排斥,被人孤立,看來與神無份了.她與睚魯處於社會的兩極,她活在社會的邊緣.

若睚魯的小女兒正經歷肉體的死亡,這女人則經歷社群生活的死亡.

她還有希望嗎?

對於我們來說,心靈裡不斷流血,不是不可能的.人是很脆弱的,受到傷害,人間的幫助未必能醫好我們.自己知道自己不住流血,但卻沒有辦法止住它,這是悲哀的.這女人的故事也可以是我們的故事.

這種生命還有希望嗎?

這一切的不幸,沒有叫這個女人失去存在下去的勇氣.她要活下去.剛才睚魯求耶穌,叫他的女兒「痊愈,得以活了.」這個女人也渴望「痊愈」,渴望好好地活下去.

她聽過耶穌醫病的事,她擠在人群中,決意要摸耶穌的衣裳.

這個女人的行動,有拼死一試的決心.她的行動在宗教上及社會上也是不容的.不潔的人竟膽敢去碰高貴的耶穌.宗教上及社會上的習慣叫她在可憐的光景中繼續下去,宗教上及社會上的習慣阻止她改變,但是,她渴望新的生命,她寧願打破種種人為的限制去尋求新生.

人群重重圍著耶穌,這些人群可說代表了世俗的藩籬,他們阻止這女人接觸耶穌,但這女人卻堅定地突破人群的阻隔,她要接觸耶穌.這一種決心,讓她真的接觸到耶穌,她的生命停止淌血了,她的生命有了新的開始.

這個社會低層的女人,她的不潔被潔淨了,這如何可能?

是她接觸了耶穌.或者說,是耶穌讓她接觸.潔淨的耶穌並沒有逃避不潔者的接觸.慈愛的神並沒有拒絕不潔者.只有當聖潔者甘願被不潔者接觸,不潔者才被支持,才被接納到聖潔的領域中.

近來,很多人討論「正生書院」事件.正生書院教的學生,是過去濫藥的青少年.他們被認為是不應接觸者,被梅窩居民拒絕進入他們的社區.當「潔淨者」不肯讓他們圈外的人接觸時,則那些被認為「不潔」的人只能停留在原地,心中淌血,很難開展新生.耶穌讓那女人接觸,那女人才能開始新生.這一接觸,屬靈的生機便開始了,這就是今日講道主題「靈機一觸」的意思.

耶穌此時是趕時間去救一個小女孩,他可以繼續他的腳步.這女人暗中作的事,讓它永遠留在暗中,也沒有什麼大不了.但是,耶穌停下來.天地間好像一切都停頓了,那女人觸摸耶穌的那一刻,像凝固了下來.

5:30    耶穌頓時心裡覺得有能力從自己身上出去,就在眾人中間轉過來,說:「誰摸我的衣裳?」5:31    門徒對他說:「你看眾人擁擠你,還說『誰摸我』嗎?」5:32    耶穌周圍觀看,要見做這事的女人.5:33    那女人知道在自己身上所成的事,就恐懼戰兢,來俯伏在耶穌跟前,將實情全告訴他.5:34    耶穌對她說:「女兒,你的信救了你,平平安安地回去吧!你的災病痊愈了.」

耶穌停下來,問:「誰摸我的衣裳?」

對問徒來說,這是一頗荒謬的問題,他們說:「你看眾人擁擠你,還說『誰摸我』嗎?」

當然,在一個擠擁的處境,很多人也可能摸過耶穌的衣裳.耶穌當然不是針對這種外在的接觸,耶穌針對一種內在的接觸.耶穌不是處理一種無心的接觸,耶穌處理一種在絕望中作出的,具有拼死一試的信心的接觸.耶穌要將這種接觸的含意表露出來.

人希望神為自己做一些「事情」,但神渴望與一個活生生的「人」相遇.人希望得到某些東西,但神渴望得到某個人.人渴望生活的幸福,這是好的,但神更希望我們渴望神自己.耶穌要與這女人相遇.

那被社會遺棄的女人走出來,帶著恐懼,俯伏在耶穌跟前,將實情全告訴他.

耶穌對她說:「女兒,你的信救了你,平平安安地回去吧!你的災病痊愈了.」

耶穌的回答對當時具有宗教意識的群眾是震撼性的.

首先,耶穌稱呼這女人為「女兒」.這是最緊密的關係的表達.這個女人,原先為社會唾棄,失去了一切人際關係.如今,耶穌稱她為女兒.這種關係比她的病得到醫治更重要.若耶穌不叫她從人群中出來,她只是得到醫治,但如今,她的生命與神有了關係.

你想單單得到醫治,抑更珍惜你與神的關係?

當時,耶穌正在趕往睚魯的家去醫治睚魯的女兒.醫治睚魯的女兒,在當時的人看來,是理所當然的.但是,耶穌卻不會因為要醫治一個尊貴人的女兒而忘記一個被社會遺棄的女人.這個女人不是任何重要人物的女兒.耶穌停下來,為要叫這女人明白自己也能成為神的女兒.對睚魯這個尊貴人而言,耶穌這行動是富有啟發性的,他的女兒不是唯一耶穌著緊的人,耶穌也愛邊緣人.睚魯日後,會否仍然不許不潔者進入他的會堂?

其次,耶穌指出,這女人的信救了她.信心,即依靠神的心,是重要的.這個女人,信而不怕,她信靠耶穌而不怕作出違反社會規範的行動.等一會,我們便見到,睚魯怕而不信,他在恐懼中不敢相信神的作為.

你信而不怕,抑怕而不信?

第三,耶穌對那女人說:「你的災病痊愈了.」耶穌宣告這女人從今以後脫離她過去的疾病.這句話也叫當時的宗教人士困惑.按利未記15:28-30的記載,女人患這種病,要痊愈七天後才能被視為潔淨,並在第八天,由祭司為她獻贖罪祭,這事情才了斷.耶穌在此取代了這一切宗教程序.耶穌宣告,人和神關係的復和,才是至關重要的.

這個女人得醫治的故事,是插在睚魯的女兒等待耶穌去醫治這個大故事中的.以這個大故事為背景,我們看到耶穌會為了一個小人物,會放下大人物的事情.大人物與小人物,在耶穌看來,都是一樣.這個小人物有信心,是睚魯這個大人物未必有的.而睚魯的女兒得到醫治的事,卻為這個不潔的女人得醫治的故事,提供一幅更完美的圖畫.

我們回到包著小故事的大故事.

5:35    還說話的時候,有人從管會堂的家裡來,說:「你的女兒死了,何必還勞動先生呢?」5:36    耶穌聽見所說的話,就對管會堂的說:「不要怕,只要信!」5:37    於是帶著彼得,雅各,和雅各的兄弟約翰同去,不許別人跟隨他.5:38    他們來到管會堂的家裡；耶穌看見那裡亂嚷,並有人大大地哭泣哀號,5:39    進到裡面,就對他們說:「為甚麼亂嚷哭泣呢?孩子不是死了,是睡著了.」5:40    他們就嗤笑耶穌.耶穌把他們都攆出去,就帶著孩子的父母,和跟隨的人進了孩子所在的地方,5:41    就拉著孩子的手,對她說:「大利大,古米!」(翻出來就是說:「閨女,我吩咐你起來!」)5:42    那閨女立時起來走.他們就大大地驚奇；閨女已經十二歲了.5:43    耶穌切切的囑咐他們,不要叫人知道這事,又吩咐給她東西吃.

有人報告,睚魯的女兒死了.耶穌剛剛醫好了一個女兒,如今,這個女兒卻死了.死,消滅一切希望.人面對死,只是懼怕,很難相信.耶穌卻說:「不要怕,只要信!」剛才那女人表達了那種不怕的信心,要有這種信心,人才能在神那裡看到生命的可能性.耶穌對眾人說:「為甚麼亂嚷哭泣呢?孩子不是死了,是睡著了.」

聖經在別的地方,也有形容「死亡」為「睡覺」的.死亡是一切的終點,但睡覺,是睡在神的懷裡,靜待神喚醒.這女孩現在等待神喚醒她.

耶穌拉著這女孩的手.

接觸死人,會帶來不潔.耶穌就是敢於走出自己的聖潔領域,去接觸不潔,這一觸,帶來了生機.正是「靈機一觸」.

耶穌這行動,解釋了先前的故事那女人得潔淨之所以可能的理由.那女人得以接觸耶穌,是因為耶穌願意接觸不潔者.神愛罪人,是罪人得救的根本原因.

耶穌吩咐女孩起來,「起來」這個字,在聖經其他地方,大部分譯為「復活」.耶穌叫死人復活.這個大故事,告訴我們,神醫治人,不單單是醫這個病或那個病,神希望人整個生命更新,神希望人復活,神希望人活一個新的生命.這是神終極的心意.

今日的經文,是一個小故事連著一個大故事.小故事提示我們,要懷著信心接近神,但人能接近神,不是因為人的信心必然會帶來這個作用,而是因為神願意主動接近罪人.這是大故事要告訴我們的.大故事告訴我們,神的終極心意,不是處理人這個或那個個別問題,神渴望人的整個生命在他裡面全然更新.

你如何看待一個社會上的邊緣人?你如何看待富人和窮人?你是否願意放下身段去幫助他人?這一切問題的答案,會告訴我們,我們的生命是否已復活了.

但願榮耀歸給聖父,聖子,聖靈.阿們.



\newpage

\section{知不知}
\label{sec:5}
\textbf{孫寶玲牧師 - 知．不知}
\newline
\newline
\hyperref[sec:4]{< < < PREV SERMON < < <}
~
\hyperref[sec:toc]{[返目錄]}
~
\hyperref[sec:6]{> > > NEXT SERMON > > >}


馬可福音 6:1-6
\newline
\begin{longtable}{cl}
\hline
\hline
章節 & 經文 (和合本修訂版)\\
\hline
6:1 & \begin{tabularx}{0.7\textwidth}{X} 耶穌離開那裡，來到自己的家鄉；門徒也跟從他。 \end{tabularx} \\ \\ \relax
6:2 & \begin{tabularx}{0.7\textwidth}{X} 到了安息日，他在會堂裡教導人。眾人聽見，就很驚奇，說：「這人哪來這本事呢？所賜給他的是甚麼智慧？他手所做的是何等的異能呢？ \end{tabularx} \\ \\ \relax
6:3 & \begin{tabularx}{0.7\textwidth}{X} 這不是那木匠嗎？不是馬利亞的兒子雅各、約西、猶大、西門的長兄嗎？他姊妹們不也是在我們這裡嗎？」他們就厭棄他。 \end{tabularx} \\ \\ \relax
6:4 & \begin{tabularx}{0.7\textwidth}{X} 耶穌對他們說：「先知除了在本鄉、本族和自己的家之外，沒有不被尊敬的。」 \end{tabularx} \\ \\ \relax
6:5 & \begin{tabularx}{0.7\textwidth}{X} 耶穌在那裡不能行甚麼異能，不過為幾個病人按手，治好他們。 \end{tabularx} \\ \\ \relax
6:6 & \begin{tabularx}{0.7\textwidth}{X} 他也詫異他們不信。 \end{tabularx} \\ \\
[1ex]
\hline
\hline
\end{longtable}
日期:2009年7月19日

証道:孫寶玲牧師

題目:知．不知

經文:馬可福音6:1-6

6:1

耶穌離開那裡,來到自己的家鄉．門徒也跟從他.

6:2

到了安息日,他在會堂裡教訓人．眾人聽見,就甚希奇,說,這人從那裡有這些事呢,所賜給他的是甚麼智慧,他手所作的是何等的異能呢.

6:3

這不是那木匠麼．不是馬利亞的兒子,雅各約西猶大西門的長兄麼．他妹妹們不也是在我們這裡麼．他們就厭棄他.〔厭棄他原文作因他跌倒〕

6:4

耶穌對他們說,大凡先知,除了本地親屬本家之外,沒有不被人尊敬的.

6:5

耶穌就在那裡不得行甚麼異能,不過按手在幾個病人身上,治好他們.

6:6

他也詫異他們不信,就往周圍鄉村教訓人去了.

?resize=250\%2C173 and \#038;ssl=1" alt="" data-recalc-dims="1" />今天早上的經文跟我們華人所熟悉的一些講話蠻相似的,比方說,我們有一句這樣子的話說:「外來的和尚會唸經」；香港人會講:「本地薑不辣」或是「隔離阿婆飯香」(隔壁的姥姥作的菜比較香),還有「外國的月亮比較圓」.

我們的講話和今天的經文 and \#8212;-我們所看見的就是有一些人他所看見的一些知識,他以為自己知道一些事情,講法,這些講法就是說,好的東西不會是本地出產的,出色的人不會在本地可以找到的,當然我們也認識一些的講法是跟剛剛我們所說的這些... 比方說甚麼「本地薑不辣」等等的講法是完全相反的.

比方說有一句這麼的話說:「猛虎不及地頭蟲」,還有一句是「茶是故鄉濃」.根據這樣的講法就是說 and \#8212;-地的東西是最濃,最好,最強的,究竟那一個講法比較對呢?那一個講法比較接近這個真相呢?

我想持著不同觀點的人,大概有辦法去證明任何一個觀點,它的合理性和它正確的地方,不過其實真正重要是我們心裡面怎麼樣去想 and \#8212;-我們以為我們看見的是怎麼樣的一種情況.其實真正的情況跟我們在那裡沒有關係,跟一個從那裡來沒有關係,跟這個地理的距離完全沒有關係,真正的問題在我們的心裡面,不是一個在那裡,從那裡來,是我們的心怎麼樣去想他,怎麼樣去看他.我說就是沒有認識以前,我們以為我們認識他,若果說我們以為認識他,已經完全認識他了,我們就說這是一些「先入為主」的一些想法,知識.

快20多年前,我去美國讀書,我在美國認識好些朋友,很好的朋友,人很溫和也很樂意的幫助我們,是很好,很好的朋友,在很好的朋友裡面有少數,大概一個,兩個這樣子,他們對於我們從東方來的人有一些先入為主的想法,他們以為知道我們,有一些先入為主的知識,他們以為我們東方來的就一定是來自「第三世界」,然後他們對於第三世界又有一種特定的想法,有一種特別的知識,他們就以為第三世界的生活是比較簡單,純樸,設備比較沒有那麼先進等等.

其中有一個朋友,就是我們住的地方的鄰舍,有一天晚上有另外一個朋友的家庭請我們家和鄰舍一起去吃飯,那一個晚上實在是非常的愜意的一個晚上,我們聊天,吃東,很愉快的,可是到了某一點,可能是我們在吃牛排或豬排的時間,這位鄰舍就很隨便不過是很確定的吐了一句話,他說我們現在吃的這個牛排或豬排做的方法跟Joe 和Salina (就是我和我太太的英文名字)在家裡面做的很不一樣,請我們吃飯的朋友就說:「怎麼會不一樣呢?他們在家裡面怎麼做肉的呢?」我還不曉得我們的鄰舍會講甚麼,他接著很認真,很有把握的繼續講下去...

他就說:「他們在家裡面做這個肉,平常是切得很細,切絲絲或肉丁這樣子,切得很細很細的,他們吃肉就是這個習慣,跟我們是一大塊一大塊不一樣.」

我們的朋友就問:「為什麼這樣子?」

他就接著說:「你不知道?讓我解釋一下!」他就很權威,很了解地講...「因為他們住的地方沒有冰箱,所以他們就買了很少很少的肉,買了很少的肉,當然就切絲,肉丁啦,每一次把它吃完,不會留下來,因為沒有冰箱嘛!」

那個情況是蠻不好意思,很窘的,我也不便多講甚麼,不過我也得承認在那個晚上我學會了一些東西,一些我自己也不知道的東西,就是原來我們中國人煮菜是因為沒有冰箱的緣故,不過這個對我來講是一個新的知識,我從來不曉得,當然我的朋友是不是那麼準確,我也不能講出一個講法,不過我這位朋友在認識我們以前是沒有中國人的朋友,他也沒有去過中國,也沒有去過台灣,香港的地方,他也不懂看中文,他也不知道怎樣做中國的菜,他看到中國人可能就在電影,電視裡面看到的那些的中國人,不過有一點我很確定的知道,他對中國人的看法就是說,中國人就是屬於第三世界的,然後他對第三世界是很清楚的把握 and \#8212;-很簡單,很純樸,生活的條件是很基本,但是當我們在聽的時候,我們想說也許這些先入為主的概念,甚至我們可以說是偏見的東西,是屬於世界某一個角落或是某一個國家吧!也不見得!

幾年前在互聯網流行一個故事或講法,我不知道是不是真的,反正就是說,很多人讀過這樣的講法 and \#8212;-聯合國發表了一個問題或思想,希望世界各地的小孩子能夠去討論一下,互動一下,題目是這樣:請照你自己個人的想法談一談其他國家的糧食充足的問題 ``PLEASE SHARE YOUR OWN THOUGHT CONCERNING THE ADEQUACY OF FOOD SUPPLY IN OTHER NATIONS."  那個問題就傳出去了,收回來的結果好像不太滿意,怎麼講呢?

他們就說非洲的孩子不懂得甚麼叫「充足的糧食」,美國的小孩子不知道「其他的國家」,拉丁美洲的國家不懂得「請」,亞洲的小孩子不懂得「有自己的想法」,所以收回來的結果不怎麼滿意,我們聽起來好像覺得蠻好笑的,不管怎麼想,這裡面有沒有一些的真理,真相,我說的真理,真相不是說每一個亞洲的小孩子不會自己思想,不是這個意思；我說的真相,不是說美國的每一個小孩子都是這樣子的自我,狂傲；我說的真相,不是說每一個拉丁的小孩子都是這樣子的自私的把別人的東西拿來.

我說的真相,跟上面的都沒有關係,我說的真相是:不論是你是怎麼樣的年紀,不管你是男的,女的,不管你長得怎麼樣子 and \#8212;-高的,矮的,胖的,瘦的,富有還是貧窮,成績好或是一般,每一個人都有一些先入為主的想法,每一個人都有一些東西他知道,也有許多東西他不知道,每一個人都有一些東西他知道 and \#8212;-他以為這個知道就可以知道所有的東西,這個問題就是,我們每一個人心裡面都有一種先入為主的「知道」或是「想法」,沒有一個人可以免於先入為主的「想法」,沒有一個人可以免於先入為主的「偏見」,沒有一個人可以!可是這也不是一個很大的問題,我說我們彼此生活的時候,也必需有一些先入為主的概念才可以,要不然我們沒有辦法溝通,像我現在所講的東西,如果我們不懂中文的話呢... 我們就不能明白了；如果沒有一些英文的想法的話呢... 那些現在透過傳譯來明白的話,也不能夠知道我在講甚麼,如果我們沒有一些共同先入為主的一些文法的規則,語言的法規的話呢... 我們實在沒有辦法溝通,我現在所講的,都是一些聲音.

沒有一些現存的想法概念,我們沒有辦法一齊崇拜,這個問題不是先入為主的想法,問題不是我們是不是先知道一些東西,問題是我們知道的把它「絕對化」,我們以為我們所經歷的,我們把它「絕對化」,那一個絕對化的經驗不允許任何人的知道進入我們的心裡面,不允許任何的經驗進入我們的心裡面,不允許我們伸開我們的手去接觸別人,不允許我們張開手懷抱別人,不允許任何改變成長的空間.生活,生命就好像是計算一樣,一加一等於二,一加二一定等於三,一加三一定等於四,這個法規是很清楚的,不在這個法規以內的我們就不要接受,拒絕,甚至是一種很卑微的一種講法.

我們剛才所讀的經文就是那麼一回事:這不是約瑟的孩子嗎?我們從小就知道他,他叫耶穌,我們就知道,他的弟弟們,雅各,猶大,和他的妹妹都在我們中間,我們都很清楚,他以為他自己是誰?我們都很清楚!這是拿撒勒家鄉的人對耶穌的拒絕,因為他們知道所有的事情,他們會覺得這孩子很不錯,很行,他確實很好,做得也不錯,不過一個木匠的兒子,將來就是一個木匠,不可能是別的!我們就知道,他現在如此,將來也必如此!拿撒勒的家鄉是很清楚,很肯定的下了一個判斷!在他們的計算裡面沒有信心,在他們沒有信心的計算裡面,沒有神奇,他們的生活就是連串的「我知道就是這樣子」,「我知道就是這樣子」.經文告訢我們說,耶穌詫異他們的不信,於是也沒在他們中間多做甚麼,然後就走了.

我們剛過的這個星期的禮拜二,有三萬八千的學生拿到他們的A Level(高級程度會考)的成績,當然為那些考取好成績的一齊來歡樂,也要鼓勵那些沒有考好的年青人,不過現在在想,無論是考好還是不好,有沒有覺得:我就知道他會考得這樣子!我就知道他會拿這個成績,我就知道他將來會怎麼樣... 我就想想,任何一個考生的成績,會不會是一些看扁他們的一些人的眼睛突然發亮呢?我不知道!我知道有一些人他是考6A的學生,他一定是前程似錦,一定有一個很光明的將來,我知道一定會!另外有的人會講:他很失敗啊!他將來就會失敗!他當然失敗,因他從小就失敗,出自一個失敗的家庭,他的父親,他的母親,一塌糊塗,他就會失敗,將來也會失敗,我就會知道!

我們會不會有這樣的講法?我就知道我自己是這樣子的,我長得這個樣子,我考得這個樣子,我的背景這個樣子...我就知道,我會很好..,. 不是,我就知道,我永遠都是如此的糟糕,也許... 也許...

不過最近有一些年青人,他們的生活,生命,證明許多人的知道是錯誤的,這些年青人的生活,他們的改變,證明了很多的,很多人的知識,觀念是完全錯的,這些孩子都有不同的背景,他們有幾樣東西是一樣的:第一,他們都吸了毒品,在不同的環境,但有一個相似的地方,就是沾染了毒品；第二,他們都送到一個叫正生書院的地方裡面學習.我想在過去一個月裡面,大家已很熟悉這方面的新聞,我就不多講了,我想大家也知道他們將來的校舍可能在梅窩這樣的一個問題,已引起了許多許多的討論,我在這裡就不多講了.

我要講第三個共同的地方,第一個是他們吸毒,第二個他們都送到正生書院,第三個共同就是:他們的生活,生命,活出一個完全令人匪夷所思的一個情況,超過一切人的計算,超過許多人的預測,超過了很多的邏輯,甚至我們可以說,他們的生活,生命所展現出來的,就是一個的「神蹟」.這個說明了許多人以為有的先前的知識是錯誤的,他們的知道其實是不知道,他們曾經被人家看扁,可能現在還有人家把他們看扁,不過有些人能看見他們,不僅僅看見他們,還看見他們背後的神 and \#8212;-改變他們背後的上帝,我在這裡讀一篇很短的文章,是一位作者在一份報章 and \#8212;-明報的副刊裡面寫一段不太長的文字,請聽我讀:

【明報專訊】

你們作了一個很好的見證

我不是教徒.

我也沒有特別的信仰.

這天,在電視鏡頭的洶湧中,忽然覺得,這個神,實在了不起.

正生書院的一班孩子,去到哪裡都是那麼整齊,那般沉默,那樣克制,面對敵意的群情,他們抬起頭,抿著嘴,即使落淚,也光明正大.

他們曾經是社會上的地底泥,是什麼令他們今日站得如此挺直?

面對挑釁式的謾罵,他們居然沉得住氣；面對搜奇式鏡頭,他們光明磊落.

我想起每逢電視出現排隊買票,打蛇餅爭籌,歌迷影迷追星,待業者找工作,動漫展霸頭位等新聞時,記者鏡頭一掃,人人拿手袋報紙掩面之窘.

不過排排隊,已怕世人見笑,更何況是曾入歧途的隱君子?!

羞人,只是心魔,讓不可告人的過去成為努力動力,說來容易,真的能做到,卻需要無限勇氣.

這群勇敢的孩子,讓我看到他們所信奉那神的力量.

梅窩居民找了一班小孩拿著示威牌匾圍著正生校長校監團團轉,兩個大男人一直笑盈盈,誰利用小孩子?誰才真心為下一代?一眼看破.

多年來,正生出來的孩子,像品牌,一班班壞孩子,走進這孤島,出來個個昂首闊步,重踏正軌.

如今年輕人毒患為禍,像正生書院這樣的學校,應該若七‧十一,每區一幢,梗有一間在附近,有什麼對付毒品的方法,機構,制度,成功得過正生書院?!

當地反對者說,外人隔岸觀火,自然講得口響.對不起,我絕不贊同.劉皇發建議正生遷往其他空置學校,其中一所在我家區內.如果屬實,我第一個支持,能成為正生鄰居,實在是我的榮幸.

孩子們,你們已為上帝作了最好的見證.

作者: 屈穎妍　(2009年6月19日) 等一下我們就會在十字架的下面領聖餐,當你領這個餅和這個杯的時候,你記得這位釘在十字架上的耶穌,曾經被人家看扁了,祂不過是一個木匠,祂不過是一個猶太人,祂不過是釘在十字架上面死了的那個囚犯,祂是一個失敗者,可是祂的復活告訴我們說,生命是超過我們的計算,生命是超過我們的知識,超過我們先入為主的知識和經驗,在祂裡面,藉著祂的復活,我們可以有新的生活,我們可以明白經歷甚麼叫赦免,甚麼叫喜樂,甚麼叫奇蹟,神蹟,因為祂的復活的緣故,可以把我們從不知,偏見,先入為主的綑綁當中釋放出來,只是綑綁可能綑綁我們自己,綑綁我們跟別人的關係,綑綁我們跟家人的關係,因著祂的復活,我們得到釋放.

我們都馬上要聚集在這個十字架的下面,我們領這個餅和這個杯的時候,願主釋放我們,阿們!



\newpage

\section{超越貧與富}
\label{sec:6}
\textbf{陳孟賢博士 - 超越貧與富}
\newline
\newline
\hyperref[sec:5]{< < < PREV SERMON < < <}
~
\hyperref[sec:toc]{[返目錄]}
~
\hyperref[sec:7]{> > > NEXT SERMON > > >}


雅各書 1:9-11
\newline
\begin{longtable}{cl}
\hline
\hline
章節 & 經文 (和合本修訂版)\\
\hline
1:9 & \begin{tabularx}{0.7\textwidth}{X} 卑微的弟兄要因高升而誇耀， \end{tabularx} \\ \\ \relax
1:10 & \begin{tabularx}{0.7\textwidth}{X} 富足的卻要因被降卑而誇耀，因為富足的人要消逝，如同草上的花一樣。 \end{tabularx} \\ \\ \relax
1:11 & \begin{tabularx}{0.7\textwidth}{X} 太陽出來，熱風颳起，草就枯乾，花也凋謝，它美麗的樣子就消失了；那富足的人在他一生的奔波中也要這樣衰殘。 \end{tabularx} \\ \\
[1ex]
\hline
\hline
\end{longtable}
日期:2009年7月26日

証道:陳孟賢博士

題目:超越貧與富

讀經:雅各書1:9-11

我們一齊看聖經之前,我們再一次禱告:

天父感謝你,給我們一群的弟兄姊妹能夠再一次地向你敬拜,也一齊地來讀你的話語,我們恭敬地將這段時間交在你的手裡,求聖靈帶領我們進入真理,我們知道唯有聖靈能夠帶領我們進入真理,求你幫助我們每一位弟兄姊妹都在這段時間裡頭蒙恩,得到你的幫助,祈禱交託奉主耶穌基督的名,阿們!

很感謝神今天可以和各位弟兄姊妹一齊來敬拜以及一齊來看聖經.

2009年──大家都知道是金融海嘯,也是經濟不穩定的一年,這個金融海嘯帶來的後果如何呢?其實仍然還未明朗,所以今天我們一齊來看這個題目──超越貧與富 .

我們今天要看的聖經主要在雅各書第一章9-11節.我們不要看輕這三節的經文,雖然只有三節的經文,而且其中的用詞都很「大眾化」,沒有甚麼艱深的字眼,但正是因為這段經文的用詞及結構都是很普通,它就能產生眾多的聯想,可能的解釋是相當多,就反而不容易確定經文本身的原意,其實在聖經裡頭難解釋的經文往往是這類型的,就是用詞和結構都是很普通的,聯想繁多而不容易找出經文根本的信息,這類經文是難以去解釋的,而在經文中有一些很特殊,艱深字眼的經文,卻是較容易去解釋,你只要找到那些特殊詞語既定的定義,詮釋起來就容易多了!

我們看看這三節的經文,雅各書第一章9-11節:

1:9 卑微的弟兄升高,就該喜樂．1:10 富足的降卑,也該如此．因為他必要過去,如同草上的花一樣．1:11 太陽出來,熱風颳起,草就枯乾,花也凋謝,美容就消沒了．那富足的人,在他所行的事上,也要這樣衰殘.

在解釋聖經的歷史裡頭,這三節的經文,起碼有三種很不同而且主要的詮釋方向,今天可能會稍微的提一提第一種和第二種的取向,我自己就較為傾向第三種解釋的取向.

一,心靈性解釋

第一種解釋這段經文的方向,就是一個所謂「心靈性」的解釋,意思就是說,用一個心靈感受的角度來解釋這三節的經文.

根據這個取向,認為這段經文的重點,是在第9,10節的上半節:「卑微的弟兄升高,就該喜樂．富足的降卑,也該如此．」而且解釋的時候,是用一種心靈性的取向,意思就是說,認為這裡所說的「升高」及「降卑」都是指著「心靈」的經驗,因此這裡的「卑微的弟兄升高」或者「富足的降卑」,這個「升高」和「降卑」就被解釋為心靈裡頭的滿足或謙卑──當貧窮的人信了主後,他的心靈就富足了(心靈裡頭的感受就富足了)；富有的人又怎麼樣呢?當他認識上帝信了主後,在上帝的面前就謙卑了,後果也是心靈的滿足了,所以都是很快樂的.像在第9,10節的上半節所說:「卑微的... 升高,就該喜樂．富足的降卑,也該如此．」,「也該如此」就是指著上一節的「喜樂」,是一個很「心靈性」的解釋.

以這種取向來解釋這幾節的經文,有沒有其他聖經的根據呢?都有!甚至在雅各書本身都有這種的經文,譬如第二章5節:「我親愛的弟兄們請聽,　神豈不是揀選了世上的貧窮人,叫他們在信上富足,並承受他所應許給那些愛他之人的國麼.」是一種很宗教性,心靈性的銓釋,貧窮的人信主了,他在生活仍然貧窮,不過「在信上富足」,在心靈上滿足快樂.

以這種「心靈性」的解釋來說雅各書第一章這裡,動機是很善良的,聽起來也很屬靈的,所以就很討好基督徒,而且我們也很喜歡這種的解釋,因為聽起來很清高.生活上如何是不打緊的,生活上或實際的經驗是貧窮或富足都不打緊的,最重要是「心靈」是否滿足?但是縱使動機是那麼的良好,用這樣的解釋來看這段經文是有問題,有弱點的!

從釋經的原則來說,有一個問題:究竟用這種「心靈化」方式來解釋這段經文的內容,在釋經的原則上是不是夠紮實呢?有沒有問題呢?因為如果我們強行以「心靈性」來解釋,就有一個困難,就是在釋經的原則上所用的字眼,特別是「升高」和「降卑」兩個字,我們把它最正常,最直接的解釋癈去,將之改變,只要「心靈性」的解釋.如果你問問當時雅各書年代的那些初期教會,當他們收到這封信的時候,他們看到這兩個字眼「升高」和「降卑」,他們會如何去理解呢?其實對當時候的信徒來說,這字眼是很簡單,很直接及很正常的,它是指著「生活上」的經驗,是指著在實際經濟生活上你是貧窮還是富貴,是簡單不過了!不是談心靈的感受的,是指著實際生活的經驗,所以如果我們不採取這個直接正常的解釋,那在釋經的原則上是有問題的,且是有點牽強的,甚至會將我們自己的,先入為主的意思... 縱使動機是很善良,卻是強行地解釋進去,從釋經的原則來看,這是不尊重聖經,所以並不是可取的!

而且以「心靈性」去解釋這段經文,還有一個問題,就是很容易將「屬靈」和「屬世」對立起來,形成背後另一個含意:如果上帝祝福我們,只有在「心靈」上的祝福,那就毋需提及有關「生活」層面的東西,如此一來,在實際生活上我們是貧窮或富足都已不打緊,不重要了,甚至不會加以理會了,只單單的能在信仰上得到上帝的祝福就是心靈的祝福,滿足及喜樂,這樣子的話,背地裡就是去把「屬靈」和「屬世」這兩方面給對立起來(上帝的祝福只會呈現在內心裡,而不是在生活層面裡頭).

但我們需不需要以此來看待基督徒的信仰呢?其實是不需要的!基督徒的信仰是不需要將生活的經驗與信仰的祝福給對立起來.所以無論如何這是第一種對雅各書這段經文的取向性解釋.

二,無情的咒詛

第二種取向的解釋是剛剛好是第一種取向的相反,是蠻有趣的!第一種取向的優點,是第二種取向的弱點；而第一種取向的弱點,反而又是第二種取向的優點.從第二種的取向來看雅各書這裡,是有一種很強烈的釋經基礎,但確不容易和基督徒信仰中那些一貫的信念互相協調.

根據第二種的取向,核心經文被認定是屬於第10節上半節:「富足的降卑,也該如此」這句,認為第9節「卑微的弟兄升高,就該喜樂」只是一個過度語,一個引子,「富足的降卑,也該如此(喜樂)」才是真正的重點!

但這重點是被認定為一種的諷刺語或反語!信息是咒詛有錢的人,這種所謂無情的咒詛,像「富足的降卑,也該如此」是第二種解釋性的取向.以「風水輪流轉」,「三十年河東,三十年河西」或「你們這群有錢人,以前就意氣風發,很了不起,很行,今天你們貧窮了」等等... 這類的咒詛.「富足的降卑,也該如此」意思是:你也應該喜樂呀!這是一種諷刺的話.

用這種的取向來解釋雅各書第一章這裡,有沒有根據呢?有!

在猶太人文獻傳統裡面,有一個很強的基礎及根據,就是在他們的傳統裡確實有一種叫「無情的咒詛」,這乃是咒詛一些有錢的人,富足的人,在這種文獻的風格底下,甚至在雅各書其他的章節內也有出現,剛才大家一齊唸的第四章13節至第五章6節就是屬於這類的經文,你唸的時候或許也會留意到,其實那一整段的經文都是在罵那些有錢的人: 第四章講論到:你們這些有錢人,因著有錢就誇口,計劃要怎樣去發展你們的業務,但神卻說:「明天如何,你們還不知道」； 第五章就說他們為富不仁:工人為你們作工,你們不但不付工資,還剋扣工人的工資.所以雅各書第四章13節至第五章6節整段經文是一幅所謂「土豪劣紳」的圖畫,有錢的人不過是一些為富不仁的人.另外在第二章6-7節也是描述那些的有錢人,罵那些富足的人,雖然有錢,富足,但對人刻薄,對神不虔不義:「你們反倒羞辱貧窮人.那富足人豈不是欺壓你們,拉你們到公堂去麼.他們不是褻瀆你們所敬奉的尊名麼.」

所以這種的解釋取向,是有一個很強的文獻傳統在背後支持著.但問題就如剛才所提及的,就是它和基督徒信仰中那些一貫的信念不協調.

基督徒信仰中的一貫匡架就是「憐憫」,「挽回」,對犯錯,犯罪的人,不是一味的諷刺,無情的咒詛,而雅各書第一章9-11節,這一段經文是基督徒的貧,富觀,是一段很基礎性的經文,如果加上「憐憫」和「挽回」這些元素的話,那就不一樣了,這就不是一味的咒詛,指責,而是希望這些富足的人能夠有一個「當頭棒喝」的教訓,可以回頭去做一個善良,敬虔的富足人,並力勸他們不要倚靠自己的財勢,對神狂妄,對人刻薄,正如同提摩太前書第六章一段經文,那段正是教導有錢人應該如何當一個敬虔的有錢人:

6:6-10 然而敬虔加上知足的心便是大利了．因為我們沒有帶甚麼到世上來,也不能帶甚麼去．只要有衣有食,就當知足.但那些想要發財的人,就陷在迷惑,落在網羅,和許多無知有害的私慾裡,叫人沉在敗壞和滅亡中.貪財是萬惡之根．有人貪戀錢財,就被引誘離了真道,用許多愁苦把自己刺透了.

6:17-19 你要囑咐那些今世富足的人,不要自高,也不要倚靠無定的錢財．只要倚靠那厚賜百物給我們享受的　神.又要囑咐他們行善,在好事上富足,甘心施捨,樂意供給人,為自己積成美好的根基,預備將來,叫他們持定那真正的生命.

有了「挽回」這個元素,就成為雅各書第一章的第三種的解釋取向了!根據這個取向來看雅各書一章9-11節,你就發現有幾個主要重點的信息.

三,得固欣然　失亦無礙

3.1,得固欣然這段經文我們說它的主題是基督徒如何看成敗得失,富貴貧窮,人生的升降起跌.可是雅各的信息到底是想說些甚麼呢?在這個問題上,你會發現原來雅各為我們勾畫出基督徒的富貴觀,成就觀:對富貴,成就是「得固欣然」,這是第一個信息.

第一章第9節:「卑微的弟兄升高,就該喜樂」,現代中文譯本翻譯成:「貧窮的弟兄蒙上帝提升,應該高興」,新譯本翻譯為「卑微的弟兄應該以高升為榮」.講的是生活上的經驗(你實際上是富貴或貧窮),所以當你說貧窮的人能夠成為富足,你是應該喜樂,高興的,這便就是「得固欣然」.

為什麼高興啊?很正常的,一個人之常情之事.但重要的是,在基督徒的信仰裡頭,富足本身不是罪惡.雖然我們知道在人類的歷史當中,有一種的主義的確是把有錢看為有罪,而且這種主義在現實的政治世界裡,曾經擁有一個很強大的力量.但基督徒的信仰並非如此,基督徒不認為有錢就有罪,因為富足本身不是罪惡,人如果能夠享用現世當中的富足,在生活上是一種好事,所以人才會喜歡生活富足,而且不單單希望自己的生活富足,也希望別人的生活也富足,這也是一種的慈悲心,亦唯有如此,賬災,救濟的工作才會變得有意義(基督徒的確做了很多的賬濟工作,哪裡有災難,我們就去).如果我們說富足本身是罪惡,有錢就有罪的話,那就根本不需要去理會那些救濟的工作,因為越窮越光榮嘛!假使越窮越光榮的話,那救濟的工作就變得沒意義,由他貧窮就好了!所以基督徒的信仰不應從這角度來看!

富足──上帝給你得到富足,你是應該高興的.當你能夠享用富足,又去分享你的富足,好讓其他人也能經驗到富足,也是一件好事!因為富足有它的「現世價值」──在現世的生活當中,人在富足當中,是一件美好的事──所以雅各說基督徒的富貴觀是「得固欣然」.3.2,失亦無礙當然最精彩的是下一句,當雅各他說,我們看富貴,成就為「得固欣然」的時候,馬上在第10節就說:「富足的降卑,也該如此」,那這句就不簡單了.在雅各的思想裡頭,基督徒的富貴觀是「得固欣然」,但也「失亦無礙」──失去也無所謂,那這就不簡單了.所以現代中文譯本翻譯第10節上半節為:富有的人被扁低,也該這樣；新譯本的翻譯就是:富足的也不該以降卑為辱.

那為何「失亦無礙」呢?雅各是否在唱高調,發神經呢?基督徒的信仰不是一個唱高調的信仰,雅各是有很多的原因,使他說出這句話來的,為什麼可以「失亦無礙」?3.2.1,沒有永遠的意義第一雅各說,因為富貴沒有「永遠的意義」.第一章第10,11節下半節「因為他必要過去,如同草上的花一樣... 草就枯乾,花也凋謝,美容就消沒了．那富足的人,在他所行的事上,也要這樣衰殘」.意思是說,無論我們願不願意,承不承認,富貴是沒有永遠的意義──現代中文譯本把這裡翻譯成:富有的人要像野花一樣凋謝,富有的人在事業上忙忙碌碌的時候,也會像花草一樣枯乾凋謝；新譯本的翻譯就說:因為他如早上的花必要過去,富足的人也必在他的奔波經營中這樣衰落.這一種的思想就是講明了「人無法帶走地上的一分一毫」.聖經上也多次提及這一個的思想,這是聖經一貫的信息,譬如,舊約約伯記十四章1-2節,詩篇九十篇3-6節,詩篇一百零三篇15-16節,以賽亞書第四十章6-8節,這些都在強調一個相同的信息──富貴是沒有永遠的意義,但有現世的價值,所以「得固欣然」(當我們享用現世當中的富足,理當高興,可是因著它沒有永遠的意義,所以失去了也不需要太介意).而且這些經文從舊約到新約都有一個共通點,就是都喜愛用「花」和「草」來形容,比喻人生的脆弱,短暫,所以雅各給的第一個原因就是為什麼我們可以「失亦無礙」.3.2.1,生活中難免的事情雅各繼續說第二個原因,為什麼基督徒的富貴觀可以「得固欣然,失亦無礙」呢?因為「得」與「失」是生活裡頭難免的事情,好像曾經有一齣很流行的電視劇主題曲,歌詞裡頭說:得與失是每個人都要經歷到的事,是一條成長之路,是每一個人都會遇見的,無論我們是否願意,承認與否,「得」與「失」都會發生在你我的身上,它不會在來臨前之際,先敲敲你的房門,徵詢你的意見說:「我來探訪你,好嗎?」,它絕對不會如此做的,它要來就來!

「得」與「失」是難免的事,就好像雅各在第一章11節上半節所說,「太陽出來,熱風颳起,草就枯乾,花也凋謝」,「得」「失」好像太陽,熱風一般,它要來就來,不容你來選擇──「太陽出來,熱風颳起」.所以一個生命成熟的人,他會明白這事,他知道人生的鍛鍊(「得」與「失」,「成」與「敗」,「升」與「降」)是必然會發生,所以重要的是無論處在得失,升降中,都要保存一份「生命的尊嚴」,不要因著成敗得失,貧富升降就失去了生命的尊嚴,亦不要因為這些事,而看不起其他的人,不要以別人的貧富,貴賤而去看輕或看重別人,亦不要因為這些東西而看不起你自己或者過份看重你自己.

金融海嘯使很多人的財產少了一大截,我也不敢問那些投資顧問,他今年虧了多少?有些人可能會失業,會失去一些東西,但最可怕的不是失去這些的東西,最可怕的是當你失去這些東西的時候,同時也失去了生命的尊嚴──對自己失去了信心,失去了生命的尊嚴.

英文的聖經譯本把這裡甚至翻譯成「驕傲」這個字,譬如,新國際的英文標準聖經(New International Version -- UK)譯本(第一章第9節至10節上半節):"The brother in humble circumstances ought to take pride in his high position. But the one who is rich should take pride in his low position..."；而英文的新美國標準聖經(English Standard Version)就用"boast"這個字"Let the lowly brother boast in his exaltation...".採用「驕傲」這個字,當然不是鼓勵我們去驕傲,其實它是指著「尊嚴」這個字(生命的尊嚴),這種生命的自尊,自重,不會因為成敗,貧富,升降而有所改變,所以取了一個很強烈的字眼「驕傲」來說明這個意思──不要因為「貧富」,「升降」左右了做人那種基本的尊嚴.3.2.1,上帝掌管一切雅各說基督徒的富貴觀是「得固欣然,失亦無礙」,為何「失亦無礙」呀?第一,因為富貴在現世的價值卻沒有永遠的意義；第二,因為得與失是生活中難免的事情；第三,因為上帝掌管一切,畢竟是神掌管富貴,賞罰,施行公義,而我們強求也沒用,這就是雅各提出來的第三個原因為何「失亦無礙」.

這個信息在聖經中也多次提到,從舊約到新約,一貫以來都強調著這個信息:上帝的主權管富貴,賞罰,施行公義,人如何強求也於事無補.詩篇七十五篇4-7及10節,撒母耳記上二章4-10節,約伯記五章8-13節,路加福音一章52-53節皆強調這方面的信息.我們記得在撒母耳記上二章,是先知撒母耳母親的祈禱文,路加福音一章52-53節是耶穌的母親馬利亞的祈禱文(很出名的尊主頌),而有趣的是,它們的內容都很相近,就是講到「上帝的主權」.

四,富貴觀的信仰信念

4.1,「永恆」的信念所以總括來說,今天我們看雅各書第一章9-11節這一段的經文,我們說這是很重要的,因為是基督徒建立如何看待富貴,貧窮,成敗,得失的基礎性經文,基督徒「得固欣然,失亦無礙」的富貴觀,其實是基於幾個很重要的信仰信念,第一個信念是永恆的信念:富貴有它現世的價值但沒有永恆的意義──有現世的價值所以「得固欣然」,沒有永遠的意義所以「失亦無礙」；而且得與失是人生難免的事,縱使我們一生能夠保住我們所擁有的東西而沒有失去,難到真能帶走一分一毫而離開世界嗎?這是從永恆的觀念去建立基督徒的富貴觀.4.2,「上帝的主權」的信念第二個信念去建立基督徒的富貴觀就是上帝的主權,上帝掌管貧富,賞罰,強求也沒用.4.3,「人的價值」的信念第三個信念去建立基督徒的富貴觀就是人的價值的信念,這是十分重要的.在上帝的眼中,人的價值是崇高的,每一個人都有一個很高的價值,只要他是一個人,他就有一個很高的生命價值和尊嚴,為什麼?因為只要他是一個人,他就是按照上帝的形象被創造的(創世記第一章),所以人就有一個很高的價值.

「得固欣然,失亦無礙」這個講法,表面上看來與「人的價值」相比較,好像有點矛盾,但其實不單止一點也不矛盾,也能相互統一.在「人的價值」的信念底下,因著人的價值是很高的,所以富足就能夠使人的日子好過一點,生活好過一點,是應當高興的,上帝也喜歡看到人的生活是幸福,國泰民安的,這是一件好事!但正因為「人的價值」很高,所以人就有高價值,高尊嚴,這樣貧窮根本就不會減少人生命的基本尊嚴,人不會因為貧窮而改變事實,因為人是按照上帝的形象被創造,或貧或富,人的基本的尊嚴與價值都仍是存在的,所以「失亦無礙」.在人的價值的信念底下就可以統一「得固欣然,失亦無礙」這個的看法.

五,總結

雅各書第一章的信息,其實是一個很實際的信息,是基督徒該怎麼去看升降,得失,成敗,貧富.

2009年下半年在經濟上,社會的穩定上會怎麼發展,其實我們仍然未知道,所以我們可能還會再面對一些的損失,一些的降卑(像第一章10節那裡所說的),但基督徒的富貴觀提醒我們,若它們真的發生在我們身上,我們該曉得怎樣去面對.

讓我們一齊禱告:

天父!謝謝你給我們今天能夠有一群弟兄姊妹聚在一起來讀你的話語,我們懇求的是你將你的話語放在我們心裡面,好幫助我們,使我們得力量,我恭敬地將這裡所有的弟兄姊妹都放在你的手裡,你知道我們的弟兄姊妹各自生活在不同的處境裡,或面對著人生不一樣的挑戰,生活的環境,但我們深信聖靈能夠在我們心中及生命中作工,聖靈能夠將你的話語在我們心中轉化,成為適合在我們每一個人在不同的處境當中所需要的信息,以至我們的生命得到幫助,求你恩待我們這裡的每一位弟兄姊妹,讓我們的生命更蒙你的喜悅,我們整體的教會亦更蒙你的喜悅,祈禱交託仰望是奉主耶穌基督的名,阿們!



\newpage

\section{好牧人與雇工}
\label{sec:7}
\textbf{伍渭文牧師 - 好牧人與雇工}
\newline
\newline
\hyperref[sec:6]{< < < PREV SERMON < < <}
~
\hyperref[sec:toc]{[返目錄]}
~
\hyperref[sec:8]{> > > NEXT SERMON > > >}


約翰福音 9:24-41
\newline
\begin{longtable}{cl}
\hline
\hline
章節 & 經文 (和合本修訂版)\\
\hline
9:24 & \begin{tabularx}{0.7\textwidth}{X} 於是法利賽人第二次叫了那以前失明的人來，對他說：「你要將榮耀歸給神，我們知道這人是個罪人。」 \end{tabularx} \\ \\ \relax
9:25 & \begin{tabularx}{0.7\textwidth}{X} 那人就回答：「他是不是個罪人，我不知道；有一件事我知道，我本來是失明的，現在我看見了。」 \end{tabularx} \\ \\ \relax
9:26 & \begin{tabularx}{0.7\textwidth}{X} 他們就問他：「他給你做了甚麼？是怎麼開了你的眼睛？」 \end{tabularx} \\ \\ \relax
9:27 & \begin{tabularx}{0.7\textwidth}{X} 他回答他們：「我已經告訴你們了，你們不聽，為甚麼又要聽呢？難道你們也要作他的門徒嗎？」 \end{tabularx} \\ \\ \relax
9:28 & \begin{tabularx}{0.7\textwidth}{X} 他們就罵他：「你是他的門徒，而我們是摩西的門徒。 \end{tabularx} \\ \\ \relax
9:29 & \begin{tabularx}{0.7\textwidth}{X} 神對摩西說話是我們知道的，可是這個人，我們不知道他從哪裡來。」 \end{tabularx} \\ \\ \relax
9:30 & \begin{tabularx}{0.7\textwidth}{X} 那人回答，對他們說：「他開了我的眼睛，你們竟不知道他從哪裡來，這真是奇怪！ \end{tabularx} \\ \\ \relax
9:31 & \begin{tabularx}{0.7\textwidth}{X} 我們知道神不聽罪人，惟有敬奉神、遵行他旨意的，神才聽他。 \end{tabularx} \\ \\ \relax
9:32 & \begin{tabularx}{0.7\textwidth}{X} 從創世以來，未曾聽見有人開了生來就失明的人的眼睛。 \end{tabularx} \\ \\ \relax
9:33 & \begin{tabularx}{0.7\textwidth}{X} 這人若不是從神來的，甚麼也不能做。」 \end{tabularx} \\ \\ \relax
9:34 & \begin{tabularx}{0.7\textwidth}{X} 他們回答他說：「你完全是生在罪中的，還要來教訓我們嗎？」於是他們把他趕出去了。 \end{tabularx} \\ \\ \relax
9:35 & \begin{tabularx}{0.7\textwidth}{X} 耶穌聽說他們把他趕出去，就找到他，說：「你信人子嗎？」 \end{tabularx} \\ \\ \relax
9:36 & \begin{tabularx}{0.7\textwidth}{X} 那人回答說：「主啊，人子是誰？告訴我，好讓我信他。」 \end{tabularx} \\ \\ \relax
9:37 & \begin{tabularx}{0.7\textwidth}{X} 耶穌對他說：「你已經看見他，現在和你說話的就是他。」 \end{tabularx} \\ \\ \relax
9:38 & \begin{tabularx}{0.7\textwidth}{X} 他說：「主啊，我信！」他就拜耶穌。 \end{tabularx} \\ \\ \relax
9:39 & \begin{tabularx}{0.7\textwidth}{X} 耶穌說：「我為審判到這世上來，使不能看見的看見，能看見的反而失明。」 \end{tabularx} \\ \\ \relax
9:40 & \begin{tabularx}{0.7\textwidth}{X} 同他在那裡的法利賽人聽見這些話，就對他說：「難道我們也失明了嗎？」 \end{tabularx} \\ \\ \relax
9:41 & \begin{tabularx}{0.7\textwidth}{X} 耶穌對他們說：「你們若是失明的，就沒有罪了；但現在你們說『我們能看見』，你們的罪還在。」 \end{tabularx} \\ \\
[1ex]
\hline
\hline
\end{longtable}
日期:2009年8月2日

証道:伍渭文牧師

題目:好牧人與雇工

經文:約翰福音9:24-41

有工會的領袖說,在金融海嘯之下,很多老板都把經濟的損失轉嫁到工人的身上；有些的公司縱使有盈利,但依然裁員,減薪,很少公司願意與員工一齊共渡時艱.

這些的遊行,示威,申訴,使我想起大概約一,兩個月之前的新聞(是令我很受感動的),在美國某一個城市裡的一間快餐店裡,老板因為金融海嘯的緣故而導致資金周轉不靈,就要將這間快餐店給關閉,不再營業,要遣散那些工作多年的員工... 那些員工因此深感惶惶不安... 但有一對退了休的夫婦,常在這間快餐店內用餐的,當他們知道這件事後,他們捨不得這群員工失去工作,失去了工作之後,他們的家人怎麼辦呢?看到他們的苦況之後,願意將自己畢生的積蓄,養老的退休金等等... 買下這間的快餐店,而且所有的員工還加薪,他說,我要好好待我的員工,就如同我的親人一樣,善待他們.這間的快餐店,透過電視媒體的報導之後,我心中覺得很開心,竟然有人看見這些員工面臨失業,自己本來都退休了,所有的積蓄是留待養老之用,竟付上自己養老金,以保員工們的飯碗(工作).我相信這家的快餐店一定會賺錢的... 我沒再去追蹤這件新聞,因為這所快餐店的名字在美國已在頭條新聞中大肆被報導了,很多人就會慕名而來,對吧?!而且員工加薪了之後,工作起來特別地有士氣,要去為老板賺錢呀!

與這群「幸運」員工相反,約翰福音第九章記載到一位生下來就瞎了眼的瞎子,是被耶穌在安息日中醫好雙眼的瞎子,竟然這群管理猶太人會堂的領袖們說:你不應該被這個人醫治好的,他不是從上帝來的,因為從上帝而來的是不會破壞安息日的.這位瞎子就說:上帝不會聽罪人所說的話,他的確是把我給醫好了,從來沒有生下來是瞎眼的人,能夠被醫治好的,他應該是來自上帝的.他就和這群猶太人的領袖展開了一場神學的辯論,結果他們說,你要是定要接受他是從上帝而來的話,那你就得離開!因為安息日是不能夠醫病的!最後這位被醫好的瞎子就被趕出會堂.

對這位生下來就瞎眼的人來說,會堂是很重要的,每一個禮拜的那一天的安息日,都會是人生的一個高潮.就像很多的菲傭一樣,一個星期最開心的就是星期日那天,因為那天不需要工作,穿自己喜歡的衣服,聯群結隊地去到禮拜堂或中環的皇后像廣場... 很高興的.所以會堂對瞎子來說是很重要的,在那裡... 做為神的選民,他們一齊禱告,頌詩,聽道,等候彌賽亞的降臨,在會堂裡,在敬拜上帝的會中,他知道上帝不是看人的外貌,乃是看人的內心.生下來就是瞎眼的瞎子,特別講到上帝的憐愛,在會堂當中,他特別講到自己生命的尊貴,會堂對他來說,本是一個公義,仁愛,盼望的地方,這是一個神聖的空間.但是現在這個被醫好的瞎子,竟然面臨被逐出這個神聖的空間,全因他堅持信仰,而被隔切離開熟悉的群族.

為什麼眼睛被醫治好了,會被逐出會堂呢?因為這位瞎眼的瞎子,他相信這個神蹟是來自上主的,因此耶穌應該是來自上主的,這是一個很清楚的邏輯,但是反對耶穌的猶太人領袖及法利賽人,他們是硬著心,他們不願接受耶穌基督是來自上帝的彌賽亞,他們要盤問這位瞎子... 第一次盤問不得要領,進而盤問他的父母親,希望得到瞎子是生下來不是瞎眼的證詞,他的父母親因為害怕猶太人,因為他們已經商議定了,凡相信耶穌是基督的,是要被趕出會堂的,因此都不敢支持他們的兒子,只是將問題推回給他說:「

是誰開了他的眼睛,我們也不知道．他已經成了人,你們問他罷．」

我們看看約翰福音第九章24節,整個事件的背景:

9:24  所以法利賽人第二次叫了那從前瞎眼的人來,對他說,你該將榮耀歸給　神．我們知道這人是個罪人.

9:25  他說,他是個罪人不是,我不知道．有一件事我知道．從前我是眼瞎的,如今能看見了.

9:26  他們就問他說,他向你作甚麼,是怎麼開了你的眼睛呢.

9:27  他回答說,我方纔告訴你們,你們不聽．為甚麼又要聽呢．莫非你們也要作他的門徒麼. (這瞎子很厲害呀!竟敢強嘴!)

9:28  他們就罵他說,你是他的門徒．我們是摩西的門徒.

9:29　神對摩西說話,是我們知道的．只是這個人,我們不知道他從那裡來.

9:30  那人回答說,他開了我的眼睛,你們竟不知道他從那裡來,這真是奇怪.

9:31  我們知道　神不聽罪人．惟有敬奉神遵行他旨意的,　神纔聽他.

9:32  從創世以來,未曾聽見有人把生來是瞎子的眼睛開了.

9:33  這人若不是從　神來的,甚麼也不能作. (這瞎子的神學觀念都很清楚,很厲害,他還能與這群的法利賽人舌戰!)

9:34  他們回答說,你全然生在罪孽中,還要教訓我們麼.於是把他趕出去了.

9:35  耶穌聽說他們把他趕出去．後來遇見他(新譯本翻作"找到他"),就說,你信　神的兒子麼.

9:36  他回答說,主阿,誰是　神的兒子,叫我信他呢.

9:37  耶穌說,你已經看見他,現在和你說話的就是他.

9:38  他說,主阿,我信．就拜耶穌.

9:39  耶穌說,我為審判到這世上來,叫不能看見的,可以看見．能看見的,反瞎了眼.

9:40  同他在那裡的法利賽人,聽見這話,就說,難道我們也瞎了眼麼.

9:41  耶穌對他們說,你們若瞎了眼,就沒有罪了．但如今你們說,我們能看見,所以你們的罪還在.

接著耶穌就說了第十章的一個比喻──羊圈的比喻,其中耶穌說:我就是羊的門,從我來的,就得生命,盜賊是不會從門而進來的...  耶穌又說:我是好牧人,我不是雇工,我顧念我的羊,為羊捨命.而雇工看見危險的時候,有狼進入羊群,就會跑掉了... 耶穌說:「

我就是門．凡從我進來的... 出入得草喫.」(約10:9)羊圈是羊的居所,牧人領羊出來是從門進來的,那裡有開門的人打開門讓牧羊人進來領羊出去吃草,晚上吃完草之後,帶羊回來,也是從門進去的.耶穌說,我就是羊圈的門,如果不藉著我,羊得不著草吃.但其他人──盜賊──並非從門進入羊圈(因為有看門的人),是由其他地方進去的,那是盜賊,不是主人.

這些說話是針對當時的法利賽人,猶太人而講的,這些法利賽人是猶太人的領袖──會堂的宗教領袖,他們的責任是引導百姓明白上帝的律法,認識上帝的救恩.當耶穌醫好生來就瞎眼的瞎子之後,他們竟然因為那天是安息日,就說耶穌破壞安息日是罪人,反而被醫好的瞎子說:「

我們知道　神不聽罪人．惟有敬奉神遵行他旨意的,　神纔聽他.」(約9:31)所以耶穌應該是來自上帝的.

瞎眼的人竟能看出耶穌是來自上帝,他是進過耶穌的那道門,眼睛被打開,得到青草吃,但這群宗教領袖,比耶穌早到的宗教領袖,卻硬著心否認那麼清楚的神蹟,他們並不關心這瞎子的福祉(我覺得他們應該知道這個人其實是很可憐的,來會堂的,應該不會有很多個生下來就瞎眼的),為何他們會如此硬心呢?他們理當是牧養這班群眾,神的選民,引導他們明白聖經,等候彌賽亞的,但他們卻是盜賊,狠心將這原本是瞎子的人趕出會堂.所以耶穌說,比我先來的是你們,是盜賊,羊是不會聽他的聲音的.這個打開眼睛被醫好的瞎子,他不聽從這班猶太人的說話,不會為了留在會堂而否認醫治他的耶穌基督是來自上帝,羊是不聽他們的聲音,因為他們是盜賊,他們並沒有好好的愛護這群百姓.

耶穌基督接著在第十章11-15節,由盜賊,羊的門的對比,轉到牧人和雇工的對比,雇工是為工價而工作──你能出價多少,我就為你幹事.雇傭兵就是這種情形,你付錢給我,我就為你去打仗；殺手也是如此的,只要你出得起價錢,不管你要殺誰,應不應該殺,我還是會去殺.最近我看到一個名詞,是以前未曾見過的,叫做「文化雇傭兵」──只要你能出得起價錢,你叫他寫甚麼都可以,且寫得還蠻好的,大家不要猜測是誰啦!

雇工是為金錢而工作的,但當雇工見到狼出現的時候,他很快地就跑走了,因為他犯不著要賠上自己的性命.但好牧人與羊有一個很特別的關係,耶穌說:「

我認識我的羊,我的羊也認識我. 正如父認識我,我也認識父一樣.」(約10:14－15)

認識是很重要的,認識是一種很嚴肅的事,乃關乎生死大事.前陣子政府就呼籲大家如果認識那兩位計程車司機的話,就趕快通知有關當局,因那兩位司機是曾經接載過一位確診染上甲型H1N1流感的墨西哥人,全香港的人都在尋找,認識這兩位司機,而且當時是以一個很嚴重的語調來發出呼籲.香港會否成為疫埠?會不會發生社區爆發?這就很在乎在這一,兩天裡面找不找到這兩位司機,結果找到了,是嗎?結果可能是一些認識他的行家,勸喻他們兩人趕快接觸有關當局.我的心很著急,如果找不到這兩位司機,怎麼辦呢?很危險的!所以認識很重要,關乎生死大事!認識上帝也是如此,是關乎永生,永死的!我們認識上主,上主認識我們,我們才有永生.我們感謝上主,這兩位司機最終有自我認識,接受隔離了.

「認識」這個字很重要,這個字是用來描述亞當認識夏娃而生了該穩,這是一個很親密的關係.我們認識上帝,上帝認識我們,是關永生,永死的.如果一個被主認識的人,對自我都有新的認識,就會顛倒雇工的觀念,會顛倒對工作的看法,他不再視自己為雇工,乃是好牧人,甚至為羊捨命的好牧人.當一些人為了自己的方便或利益,無視自己可能成為帶菌者,而逃避自我隔離的時候... 當我聽見墨西哥城發生醫護人員逃亡的時候,我特別懷念香港2003年沙士疫症時的醫護人員的專業操守,對他們肅然起敬... 我特別想起香港兒女,中大校友,92年醫科畢業的謝婉雯醫生.這本是謝婉雯醫生在這裡舉行追思禮拜的記念手冊,她是主動在當天與同事調換進入屯門醫院的疫症病房來搶救病人,原本她並不需要去的(當時候,許多在台灣的醫護人員都是逃亡,避開上班的),香港的醫護人員真的很漂亮,特別是第一位在沙士時捐軀生命的,她是自願請纓進去一間沙士病房裡,因此感染到沙士的病毒而離開世界,終年三十五歲.在2003年5月18日在這裡舉行的追思會裡面,有很多人出席,當時我是主持這個追思會的牧師,我很感動.中文大學的校友會發起萬人感謝的網頁,很多人上網留下他們致意的說話,其中有一段,我想唸出來:

「謝婉雯醫生,我真的衷心敬佩妳,看了妳一生的傳奇,神賜給妳的勇氣,在妳的生命裡是可等的美麗,完美,無悔,求神使你的靈魂得到安息,在世的工作已經完成了,妳教活人的生命之外,更徹底醫治了無數的心靈,使人因妳失去生命而得著生命,我淚水也流乾了.一次又一次聽到新聞有關死亡的消息,一次又一次心情沉了下來,而妳卻是點亮了生命,我每天都為著周圍的人的平安來禱告,亦為患病孤單在病房,獨力,無助地面對死神,希望他們知道天上的神一直垂聽他們的禱告.」

我不再唸了,有很多篇感人的話... 特別在今天唸出來的時候,使我們感到很敬佩,我們的主認識我們,我們認識我們的主,如果認識主的時候,我們對雇工的工作看法,是會不一樣的.我們不會視自己為雇工,只為了打一份工作,我們視之為牧人,我們牧養上帝交給我們的人,去愛護,幫助他們.

謝婉雯醫生當時在她的婚禮中獻上了一首詩歌【彩虹下的約定】,感動的是她先生患了血癌,但她仍遵守承諾與他結婚,結婚時所唱的一首歌就是【彩虹下的約定】,在那天的追思會當中,同一班醫學團契的人,也同樣唱著這歌:

【彩虹下的約定】我空虛的心靈,終於不再流淚,期待著雨後,繽紛的彩虹,訴說你我的約定.

我不安的腳步,終於可以停歇,主你已為我,擺設了生命的盛宴.

與你有約,是永恆的約,彩虹為證,千古不變.我要高歌,為生命喜悅,萬物歌頌你的慈愛,大地訴說你的恩典.

這首歌成為當時沙士時期少有的流行金曲,是榜上有名的金曲,就是因為這群基督徒的見證,不是視自己為雇工,逃避危險... 在沙士時間,很多基督徒的醫護人員,因著信仰的緣故去牧養那些患病的人,要拯救更多患病的人,他們得到很多香港人的尊敬,所以在那段時間,傳福音是比較容易的.我們感謝上主,我們可以因為認識上帝的緣故,我們對工作有新的認識.

這位瞎子被趕出去之後,耶穌怎樣說呢?約翰福音第九章35節:「

耶穌聽說他們把他趕出去．後來遇見他」, 「遇見他」應該翻成「找到他」(found him).耶穌知道他被趕之後,耶穌動了慈心,周圍去找他,找到他後,耶穌邀請他相信耶穌基督,作祂的門徒.這是很感動人的對話:「

誰是上帝的兒子,叫我信他呢. 耶穌說,你已經看見他,現在和你說話的就是他. 」(約9:36-37)這裡說到耶穌尋找羊群,歸回祂的羊圈,在第十章16-18節最後的比喻裡,耶穌說:「

我另外有羊,不是這圈裡的．我必須領他們來,他們也要聽我的聲音．並且要合成一群,歸一個牧人了. 我父愛我,因我將命捨去,好再取回來.沒有人奪我的命去,是我自己捨的.我有權柄捨了,也有權柄取回來．這是我從我父所受的命令.」

這裡說到愛和權柄的關係,我們很多時候找工作是為了「跳槽」,因為可以賺更多的錢,當錢賺多了的時候,我們就希望有更大的影響力,希望有更多的權柄,追求金錢,權柄是我們對「工作路徑」(Career Path)的一個很清楚的目標.但這裡耶穌說,捨命的愛與權柄是不能分開的.

愛──耶穌因為愛的緣故,順服上帝的差遣,拯救世人,父因為耶穌愛的緣故及順服,賜祂權柄從死裡復活,所以我們看見,我們可以因為愛的緣故,捨棄自己的生命甚至為了愛去牧養人的緣故,得著權柄被人尊重.

當想起愛與權柄的時候,使我想起我曾牧養的一間教會裡的一位姐妹,這位姐妹因著小兒麻痺症,行動很不方便,但是她很有愛心,很關心教會裡頭一些年老,年長的弟兄姊妹們,探望他們,陪他們下樓梯... 對新來的朋友也很關心.每一年我生日的時候,都收到一張賀卡,是她寄給我的(最近這幾年就沒收到了),我在1989年有半年的時間在那間教會當顧問的,但很多年都收到她寄來的生日賀卡,很有愛心.在他們50週年堂慶時,他們選舉最有愛心的會友,她就被選上了,她沒唸甚麼書,行動又不怎麼方便,但她在會友中是被受尊崇的一位,因為她有愛心──有愛心,有權柄,被人尊重.

耶穌說,有有很多羊不在這個羊圈內,我要將他們領回來.

這就是每一個認識耶穌的人,又是被主認識的人一個很重要的召命,我可以說我們的工作只有一個,就是領人歸回主的羊圈──This is our vocation我們的召命──Vocation只有一個,就是領人歸回上主的羊圈,而我們的方法就是我們透過不同的工作.在那裡我們不是把自己看作雇工,我們乃是看自己為牧羊人──牧養上帝給我們的羊,無論做甚麼工作,崗位,occupation,但記得要 pre-occupation,occupation 前的pre-occupation就是領人更接近那個羊圈.

我們如此行的時候,當你用愛心去牧養上帝給你的學生,同事,甚至乎是上司,要用愛心來對待他們,盡心做好這份工作,如同上帝把這些人託付給你,我們自然而然有一種權柄(吸引人的權柄),使人更接近上帝的權柄.

我們來禱告:天父上帝,當我們在疫症的期間,我們再一次地被提醒「人生命的軟弱與脆弱」,明天會如何,我們真的不知道,上主啊!我們感謝你,因為你認識我們,以至我們能不再看自己為一個雇工,我們是你的牧羊人,我們帶領很多人能夠進入這個羊圈內,求主賜我們有愛心,因為愛人的緣故,愛你的緣故,又甘心樂意地在我們的工作裡頭盡上自己的本份,顯揚你的香氣,奉耶穌的名禱告,阿們.



\newpage

\section{你不能缺少的一件事}
\label{sec:8}
\textbf{賀志勇博士 - 你不能缺少的一件事}
\newline
\newline
\hyperref[sec:7]{< < < PREV SERMON < < <}
~
\hyperref[sec:toc]{[返目錄]}
~
\hyperref[sec:9]{> > > NEXT SERMON > > >}


路加福音 10:38-42
\newline
\begin{longtable}{cl}
\hline
\hline
章節 & 經文 (和合本修訂版)\\
\hline
10:38 & \begin{tabularx}{0.7\textwidth}{X} 他們繼續前行，耶穌進了一個村莊。有一個女人，名叫馬大，接他到自己家裡。 \end{tabularx} \\ \\ \relax
10:39 & \begin{tabularx}{0.7\textwidth}{X} 她有一個妹妹，名叫馬利亞，在主的腳前坐著聽他的道。 \end{tabularx} \\ \\ \relax
10:40 & \begin{tabularx}{0.7\textwidth}{X} 馬大伺候的事多，心裡忙亂，進前來，說：「主啊，我的妹妹留下我一個人伺候，你不在意嗎？請吩咐她來幫助我。」 \end{tabularx} \\ \\ \relax
10:41 & \begin{tabularx}{0.7\textwidth}{X} 主回答說：「馬大，馬大，你為許多的事操心煩惱， \end{tabularx} \\ \\ \relax
10:42 & \begin{tabularx}{0.7\textwidth}{X} 但是不可少的只有一件。馬利亞已經選擇了那上好的福分，是沒有人能從她奪去的。」 \end{tabularx} \\ \\
[1ex]
\hline
\hline
\end{longtable}
日期:2009年8月9日

證道:賀志勇博士

題目:你不能缺少的一件事

讀經:路加福音10:38-42

10:38

他們走路的時候,耶穌進了一個村莊.有一個女人名叫馬大,接他到自己家裡.

10:39

她有一個妹子名叫馬利亞,在耶穌腳前坐著聽他的道.

10:40

馬大伺候的事多,心裡忙亂,就進前來說,主阿,我的妹子留下我一個人伺候,你不在意嗎?請吩咐她來幫助我.

10:41

耶穌回答說,馬大,馬大,你為許多的事,思慮煩擾.

10:42

但是不可少的只有一件.馬利亞已經選擇那上好的福分,是不能奪去的.

?resize=250\%2C183 and \#038;ssl=1" alt="" data-recalc-dims="1" />廣東話講:因愛成恨.這是我問了Joanne的,應該不會錯.意思就是說,愛得太多了,可能變得恨得也越多.當初愛雷曼的人,我估計現在恨得一定不少.我在生活中見過很多這樣的事.嚴正申明一下,不是講我自己拍拖這一類事情啊.

幾年前,我陪同幾個老師和一班學生一起玩,其中一個老師帶著師母.那個師母人很好,很Nice,大概50歲,同我也談得來.但是過了兩天,我發現了一個奇特現象,就是吃飯的時候,師母坐在哪里,其他人都爭先恐後的往她對面坐,包括她自己老公.算了,我看到這種情形,慈心大動,坐在她身邊.結果,那餐飯吃得我真是一世難忘.那個師母自己吃了兩口之後,就開始幫我夾菜,剛開始我還暗暗高興,不用我站起來夾好菜,多斯文啊.可是過了一陣,發現不妙,這個師母幫人夾菜不停手,看我吃完了,又夾了一碗.我又不好意思浪費,結果我剛吃完,她又夾了一碗.我終於明白,為什麼她老公都爭著往她對面坐.終於我都開始爭著往她對面坐,她很委屈.她同我抱怨,說,小勇,怎麼連你也不肯坐我身邊了?我心想,吃不消啊,我夠膽坐你身邊麼!你看看,好人做得過頭,做得太多,有時候反而變成了怨恨.

父母對子女也可能這樣.我就見過一個媽媽帶小朋友坐火車,她要那個小朋友坐好,但是小朋友就是不肯坐,到處跑.好不容易,媽媽把他捉過來,他又扭來扭去不肯坐,搞到媽媽都要發瘋了.最後媽媽大怒,提起他就打屁股.媽媽是好意,她愛她的小朋友,怕他出危險,可是因為她不知道怎麼去愛,擔心的事情太多,這份愛有時反而成為了怨恨,變成了暴力.

身邊其實經常能發現這種情況,我們愛得很辛苦,我們做了很多,但好像別人並不領情,並不體諒我們,這是為什麼?好,今天,我們借著聖經的亮光來看看,這到底是一種什麼情況,我們應該如何去面對這些問題.

馬大與瑪利亞的故事,其實有好幾個不合理的地方.耶穌剛剛講完了好撒瑪利亞人的故事,告訴人們怎麼樣去愛人.之後他來到了伯大尼的一個村莊,見到了他的老朋友馬大和瑪利亞.主耶穌上門,當然要好好招待,尤其馬大是一家之長,好好招呼主耶穌是理所當然的.你想想,如果胡錦濤來到你家做客,你會怎麼招待呢?肯定是一家大小忙個不停了,很正常啊.但是瑪利亞作為家裡的老二,卻什麼事都不做,就是坐在耶穌腳下聽耶穌講道,這是不是待客之道啊?是不是有點奇怪.這是其一.可能現在有的男人在家裡來客會這樣.其二,當做姐姐的要求妹妹也來幫手時,耶穌基督卻告訴她,你妹妹已經得了上好的福分,你卻沒得到.奇怪,做事那個沒有誇獎,不做事那個反倒得了上好的福分.算了,以後家裡來客,我什麼都不做,就陪住他聊天,直到他餓到要告辭就好了.是不是有點奇怪呢?

這段經文看起來有點奇怪.恩,是有點奇怪!耶穌為什麼要讚揚瑪利亞,而不是誇獎殷勤待客的馬大呢?我想,理解的焦點往往集中在這幾點上.

第一,耶穌並不是責備馬大伺候得太忙碌.10:41-42節:

10:41

耶穌回答說,馬大,馬大,你為許多的事,思慮煩擾.

10:42

但是不可少的只有一件.馬利亞已經選擇那上好的福分,是不能奪去的. 在聖經中,接待弟兄姊妹是一個很好的愛的表現,接待主當然是一種愛主的表現.主耶穌並不是不喜歡馬大的殷勤招待,主很享受馬大的招待,你看看,主並沒有說,馬大,你忙來做啥啊,我一頓不吃又不會餓死,麻煩你過來聽講道啊.主並沒有這麼說馬大.主只是說,馬大,馬大,你為許多的事思慮煩擾.什麼是煩擾?希臘文原文的意思,就是一盆平靜的水,你就攪動一下,水就不平靜了.心也是這樣,心裡邊有事,就不安寧了.主不是批評馬大太好客,招呼得太多,而是說她招呼得不專心.她招呼得不專心,有太多煩擾!  如果馬大雖然忙碌伺候,但是她懷著甘心樂意的心去招待主耶穌,主一點都不會批評她.問題是她的敬拜--或者說伺候,不專心.聽道是敬拜神,伺候主也是敬拜神.就好像今天在我們的禮拜堂裡,聽牧師講道是敬拜神,參加詩班也是敬拜神,做司事做主席都是敬拜神,哪怕我們走出教堂,我們的接人待物都可以是敬拜神.我們合在一起,有一顆合一的心,安寧的心去敬拜神,這才是神最悅納的敬拜.怕就怕的是我們有太多的思慮煩擾,我們的心不能安定,不能合一,不能專心.

有時候想想,馬大跟瑪利亞,我們更像哪一個? 我在詩班做禱告時,我常常強調,請神脫去我們一個星期以來的纏累.我們有太多的思慮煩惱.這裡還沒有坐穩,這裡想著哎呀,小孩子功課沒做好了,或者小孩子出去參加營會,不知道安全如何,哎呀,現在這個豬流感,真煩人,進門還要量體溫!或者剛剛踏進禮拜堂,突然之間想到,慘了,車窗沒關好,昨天同太太吵架了,今天早上都昏頭了,車窗都不記得關了；或者一邊聽講道,一邊想,哎呀,教堂對面那間餐廳關了,中午去哪里飲茶好呢?你看看,我們香港人真是忙啊.最有趣是,一邊聽一邊想,今天這個牧師講得不錯,回去我要同我老公分享一下,他怎麼這麼不開竅了,還不肯信主.--等自己一想完,牧師已經不知道講到哪里去了.思慮煩擾,思慮煩擾,身在教堂,心卻不專心敬拜.  生命與生活必須有焦點.大家玩過相機就知道,對焦對不准,就照不出好相片.弟兄姊妹,我們的生命,生活如果不專心,沒有好好對焦,我們的生命與生活也好不到哪里去.

第二,主耶穌是責備馬大把次序搞反了.主耶穌講,上好的福分只有一件!什麼是上好的福分?上好的福分不是六合彩中最大獎,上好的福分就是上帝的永生之道,聽上帝永生之道並行得出來就是上好的福分.馬大卻把次序搞反了,他以為幫主耶穌多做點事奉就是最愛神了.其實不是,最愛神就是聽神的話.耶穌並不是誇獎瑪利亞不做事,不是鼓勵懶惰,耶穌是看到他的道已經在瑪利亞心中生根,帶給她內心的平安.如果馬大有心聽神的話語,她就會安安心心去伺候.

但是馬大把次序搞反了,她把自己的事工看得比神的話還要大!雖然她愛神,但是次序不對.次序這個問題對我們來講很重要!我們往往把次序搞反了!比方說,我有個朋友,很愛她的女兒,希望她的女兒又聰明,又漂亮,又有修養,所以她帶個女兒,除了功課之外,星期一晚上她陪女兒學跳舞,星期二晚上陪女兒學鋼琴,星期四晚上陪女兒學畫畫,星期六晚上陪女兒學英語.為了鼓勵女兒更有自信,她還要幫她報演講班.有一天她很痛苦的跟我說,她的女兒罷工了,不幹了,要炒掉她這個媽媽.--這是她女兒的原話來的.為什麼,為什麼?她問我,她這麼辛苦,有時候陪女兒畫畫,陪女兒彈鋼琴,自己在一邊累得要睡覺,還要忍受女兒一遍一遍沒有練好的鋼琴聲.多辛苦!結果女兒還要炒她魷魚!

我就問,"你為什麼要做這麼多,你究竟是希望你的女兒幸福快樂,還是希望你女兒痛苦呢?你如果希望她幸福快樂,她這麼小,為什麼給她這麼大的壓力?你完全可以根據女兒的特性來栽培她,但是請不要把自己的想法,把自己未完成的夢想全部加在她身上,她還是個孩子." 現在很多孩子都不像孩子,像個提早成熟的小大人.

我們的愛有時候也是這樣,我們往往把愛的次序搞反了.我們愛孩子,往往愛那些表面的東西,卻忽視了孩子有他們自己的特點,有他們自己需要的快樂,忽略了神賜給孩子真正的恩典,乃是從神的話語而來的成長--堅定,勇敢,善良,有獨立的思考與能力,有責任感,生命最重要的不僅僅是眼前成功,更重要的是成長!正如神所說,你們不要為生命憂慮吃什麼,穿什麼... 你們要先求神的國,和他的義,這些東西都要加給你們了.但我們往往先是去愛飲食,愛衣服,愛權力,愛車子房子,愛這個世界,然後再去愛神的國和神的義.神通過聖經說:你們愛的次序錯了,上好的福分,乃是與神在一起的平安!如果你不懂得這個次序,你就是除了忙,就是煩.

第三,耶穌看到馬大內心的怨恨!愛居然會帶來怨恨?我告訴你,會的!大家可以再看看馬大說的話--10章40節:

10:40

馬大伺候的事多,心裡忙亂,就進前來說,主阿,我的妹子留下我一個人伺候,你不在意嗎?請吩咐她來幫助我. 其實聖經希臘文原文並沒有這麼客氣,並沒有這個請字,而且口氣還是帶著一點埋怨.主啊,你沒有看見我為了招待你正忙得要死麼?你怎麼留住我的妹子說話,卻不讓她來幫我呢?這個就是馬大的口氣.你聽聽,多麼不耐煩啊?而且這句話還有點指桑罵槐的意思,一部分是說給主耶穌聽,一部分是說給她妹妹瑪利亞聽. --妹妹,你有沒有一點良心啊?我在那裡忙碌招待主,你就在這裡偷懶. 你看看,馬大很愛主,所以為伺候主而忙碌.但因為愛的次序錯誤,讓她最後不僅埋怨妹妹,更大膽埋怨主來了.你想想看,如果胡錦濤坐在那裡跟你先生說話,你夠膽去埋怨麼?你敢跟他說:胡錦濤先生,我這麼忙,你就不吩咐一下我老公來幫我剁肉麼?你看看,我們不敢跟胡錦濤抱怨,可是馬大的怨讓她夠膽去埋怨主耶穌了.

主看到了馬大內心怨恨的種子.這種怨恨其實不是來自愛,而是來自我們的恐懼. 我們可以想像一下,當時馬大恐懼什麼呢?她可能恐懼貴客來了,自己招待得不好,客人不滿意；或者她怕自己招待不周,周圍鄰居會笑話她,或者她怕主耶穌更喜歡同瑪利亞說話,而忽視了她的存在.總之,這些恐懼催生了她心中怨恨的種子.

這些種子看起來小,但發出來就不小.以弗所書4章26節說,「

生氣卻不要犯罪.不可含怒到日落.」 主知道,恐懼,生氣,抱怨都不是罪,但是恐懼,生氣和抱怨很容易引發罪.我們常常會因為那一點點恐懼,抱怨而導致心裡的苦毒.尤其是這種因愛而來的恐懼,抱怨,更加是如此.有時候孩子太調皮,我們本是擔心他們出危險,但最後我們會恨不得打孩子一巴掌,多暴力!也許我們可以一時控制住脾氣,但是難保不會把這個脾氣轉移到配偶上,工作上,或者其他事情上.還有一些家庭的小事,夫妻之間怕傷害彼此感情,忍住不說,等到積怨一多,最後肯定小事變大,大事變爆.本來互相相愛的人,卻變成了互相傷害的人.還有一些大教會,本來大家一班基督徒,彼此相愛,但是可能因為彼此界線不清楚,又不能及時直接溝通,時間長了反而弄到肢體分散. 含怒到日落,就會出大事! 所以在這段經文中,主不是責怪馬大招待主而忙碌,主其實很享受一個愛主的人對他的事奉,主責怪馬大是因為她招待太多,被許多的事情煩惱,由恐懼反而生出怨恨來了.

主的智慧正在這裡顯明,她清楚看到了馬大怨恨的苗頭,他直接指出了馬大的問題,不讓馬大含怒到日落.他不是在批評,他是在傳道,是在愛馬大.他絕不做一個和事老,"哎呀,瑪利亞,你也別聽道了,你去幫幫你姐姐吧."他絕不這麼說.許多人以為做和事老是出於愛,但有的時候,做和事老不是出於愛,而是出於恐懼.許多人怕失去愛, 怕失去和人的關係,怕失去安寧,因此愛得戰戰兢兢. 但"愛裡是沒有懼怕的!" 主的愛不僅僅包含著體諒,也包含著管教.主同樣也希望馬大認識到她的問題,好叫她在愛裡沒有懼怕! 弟兄姊妹,如果你現在愛的很疲倦,愛得很辛苦,愛得很抱怨,那我提醒你,不是你的愛太多了,而是你愛的方式不對!正如我剛才說過的,第一,不要愛得不專心,不要把許多思慮煩擾,那些不屬於真愛的因素加在這樣一份愛中,愛得十分不甘心樂意.對子女,對他人,對神,我們付出了許多時間精力,可是事情一多,我們就開始計算划算不划算.我是到底要帶孩子去學英文呢,還是去崇拜呢?我們是出去度假呢,還是參加主日學呢?這樣的愛怎能不讓我們思慮煩擾呢!可能這樣的愛會讓我們一邊崇拜一邊想,哎呀,星期天怎麼這麼長,怎麼還沒過完呢!

第二,不要把愛的次序搞反了.愛子女是管教培養,讓他們發展自己的潛能,多一些呵護體貼,而不是整天說我很忙,我很忙,忙著賺錢 養家,忙著帶他們學這學那,卻忽視了我們的家人真正的需要.愛神是安心聽神的話,而不是老是說我很忙,我很忙,做得越多是為了將來奉獻的越多.  多一點親近神吧,那才是你愛的重心.

第三,愛裡並沒有懼怕.真正的愛是坦然無懼的愛,我們愛不是因為我們怕失去什麼才去愛.我們愛是因為神先愛我們,讓我們有能力有自信去傳遞這份愛.

當我們把愛真正搞懂時,我們才能愛得沒有疲倦,沒有懼怕,也就沒有抱怨,就如神愛我們一樣,因為我們的愛是清心的,是平安的!

不可少的只有一件,那就是神的愛.在世上,沒有什麼憂苦愁煩,能使我們與神的愛隔絕!我想在講道結束之前,請大家安靜聽一首歌:耶穌愛你!



\newpage

\section{曠野行旅嗎哪}
\label{sec:9}
\textbf{梁元生教授 - 曠野,行旅,嗎哪}
\newline
\newline
\hyperref[sec:8]{< < < PREV SERMON < < <}
~
\hyperref[sec:toc]{[返目錄]}
~
\hyperref[sec:10]{> > > NEXT SERMON > > >}


日期:2009年8月16日

證道:梁元生教授

題目:曠野,行旅,嗎哪

一,    引言今天講道的題目是「曠野,行旅,嗎哪」,題目中有三個重要的名詞,第一是曠野(Wilderness),第二是行旅(Pilgrims),第三是嗎哪(Manna).講道題目是伍牧師定的,可能也不是完全出自他個人主觀的主意,因為在基督教崇拜學中有一個傳統,讀經和講道都有一個普世的,宏大的設計,就是跟隨著基督教全年的時令和節日去編排題目與經課,好像今天是聖靈降臨期的第九個主日,全世界許多教會都會誦讀今天的經文和直接或間接地用今天的主題去進行崇拜.換言之,今天這個時候,有不少的教會和不少的基督徒和我們同時對今天的題目進行思考.

所以,我今天第一步要做的工作,是要把這個題目送給大家,讓大家一同進行反思.面對這樣一個題目──「曠野,行旅,嗎哪」,你會想到些甚麼?你會怎樣去處理這個題目?這個題目是牧師給我的題目是,也是給你的題目.甚麼是「曠野」?甚麼是「行旅」?甚麼是「嗎哪」?對一般人而言,「曠野」是不常見的,可能是一片沙漠,或是一片荒原,對香港人來說,這是個陌生的辭彙.甚麼是「嗎哪」,對於非基督徒而言,更加是個完全陌生的東西,但對熟悉聖經的基督徒來說,卻是個常常聽到的名字:聖經中讀到它,唱詩歌唱到它,牧師講道時提到它,但有很少人會說我們曾經親自見過它和嚐到它.三樣東西之中,我想大家都比較熟悉的,應該是「行旅」,因為香港人喜歡旅行,一有假期,很多香港人都會選擇旅行去消磨假日,遠者去歐洲,美國,近者去日本,東南亞,或者中國內地.但「行旅」又是否等同於「旅行」呢?如果我們仔細看伍牧師給我們的題目,除了中文之外,題目中尚包括了英文的三個字──Wilderness, Pilgrim, and Manna,「行旅」的英文是Piligrim, 而不是 Traveller,明顯的兩者的意義不盡相同,下面我等一下再詳細講述.

此時此刻我想大家一起來做的,是如何思考這樣一個

題目:「曠野,行旅和嗎哪」?

當我接到這樣一個題目,我最初的反應有二.首先,是把題目中的三種東西或三種元素,即是曠野,行旅,嗎哪連接起來講,或者說是把題目裡的三種東西合起來一起去講.這樣的處理方法,再我看來,是非常簡單而直接的,把它們連起來,信息就出來了,也不用我多加分析.我可以用幾句話就可以把信息說得清楚,甚至用一句話就把道理說完了.然而,我們也可以用另一種方法去處理這個題目,即是分拆式的方法.這個方法和連起來或合起來的方法不同,是要把題目裡的三種東西或三種元素分別開來,逐一處理和逐一分析,這樣的話,可能就要用較長的時間,才可以分析得清楚透徹.

二,    曠野,行旅,嗎哪:連成一起的信息先讓我們嘗試用第一個方法去思考今天的

題目:即是把題目中的三種東西──曠野,行旅,嗎哪──連成一起來思考.

如果大家接到這個題目,要你把這三種東西串連起來,即是把曠野＋行旅＋嗎哪,你會得到甚麼?意思是說請你去把三樣不同的事物聯繫起來,使之建立起一種內在關係,構成為一個主題,表達一個信息.你會怎樣去把它們連接起來呢?曠野和行旅有甚麼關係呢?旅人與嗎哪又有甚麼相干呢?可能你有你的說法,我有我的說法,但我相信,凡是對基督教道理或背景略有認識的人,無論大家從那一個途徑去思索,得出來的結論或信息都幾乎是一致的.

我相信大家面對今天的

題目:曠野,行旅,嗎哪,從三者相連的角度看,一定會得到相當一致的看法.根據聖經所提到的背景,題目中的第一個詞彙或第一樣東西是「曠野」,許多時候「曠野」在聖經中都代表著漂泊,而第二個名詞是「行旅」,「行旅」是「旅人」,是旅途當中的人,即是過客,也寓意著人生,而第三者「嗎哪」,是從天降下的白色的食物,即是天糧,意即上主所賜的恩典..對此三者,大部分的基督徒都似乎有著共識,爭議不大.換言之,若我們去思考今天主題中的三種事物,我們很容易得到相同的看法.我相信我的想法和各位的想法不會相差太遠.我們從這個題目得出來的啟示大致相同:那就是:在我們漂泊(曠野)的人生(行旅)之中,往往有上主的照顧,得蒙上帝的保守,按時賜予(嗎哪).如果用最簡單的一句話概括:那就是「漂泊人生中得蒙主恩典.」

這是第一種解讀的方法,把三者連成一起的解釋. 三,    把題目拆開來講:曠野,行旅,嗎哪若果大家認為我這樣去解析今天的題目已經足夠的話,接下來的講道便算是「畫蛇添足」,可以不理的了.但容許我仍然試一試用第二種方法去解題.至於第二種方法,就是把題目拆解,把三種元素逐一的分析.先說曠野,再說行旅,最後來談嗎哪.

今天題目中的三樣東西,或者說三種元素之中,我們最為熟悉的應該是「行旅」,至少我們直覺認為如此.至於「曠野」和「嗎哪」,對我們來說都有點陌生.試問有誰真的到過曠野?(可能只有李惠雄吧!) 又有誰食過真的嗎哪呢?但我們常常在基督教的圈子裡聽到人談「曠野」和「嗎哪」,便漸漸的覺得對這兩種東西也不陌生,甚至認為相當熟悉了.今天讓我們來反省一下,捫心自問,我們對這三樣東西究竟認識多少?

(1)    談「曠野」

先來談談「曠野」.

曠野是甚麼?是一片沙漠?抑或是一片荒原?是未曾開發的山林,抑或是了無人煙的大漠?曠野是沙是石?是有草沒草?是山崗,或是叢林?皆無定準.但無論如何,曠野總有一個特色,就是它是個「人跡罕至的地方」(a no-man"s land),是個「野獸出沒的地方」(a place of the wild).由於人跡罕至,野獸當道,推而廣之,我們可以說「曠野」象徵者「原始」和「獸性」,象徵者「未開發」和「野蠻」,是個「弱肉強食」,「適者生存」的世界,和「文明,進步」的城市人,需要互相調節,互相守護,互相忍讓和互相幫忙的城市文化,都市生活正好相反.在曠野,我們要看生命的原始,自然的變化,環境的限制,在城市,在現代社會,我們則講五德四美,文明規矩.

換言之,今天讓我們先從對立面看「曠野」和「城市」,「野蠻」和「文明」.站在這個互相對立的角度,我們自然地會對「曠野」存著戒心,產生恐懼,甚至會對之抗拒和排斥.另一方面,我們又會希望把「曠野」開發成為「文明」,用文明(特別是科技文明)去把「曠野」征服,把「曠野」改造成為「城市」,讓人聚居,讓人過文明生活.

在聖經之中,「曠野」的描寫,「曠野」的圖像,層出不窮,多次在舊約和新約中出現.我記得我在這裡也講過舊約以色列人出埃及的故事,摩西帶著他們離開埃及,出紅海,過西乃山,就進入曠野之地.他們在曠野漂流四十年,處處受到外族人的攻擊,受盡日曬雨淋,口渴飢餓之苦,最後才得進入神為他們預備的迦南美地.在新約中,我們也讀過四福音書中的記載,就是關於耶穌出來傳道之前的事蹟.聖經說:耶穌先到曠野去,經過四十天的挨飢抵餓生活以及三天的魔鬼的試探,才出來傳道.另外聖經也告訴我們,曠野裡有獅子遊行,四處尋找可以吞食的人.

換言之,就算從聖經去看,「曠野」也是一個充滿危險的地方,有魔鬼的試探,有遊行的獅子,有伺機攻擊的敵人,有自然環境的險阻,包括風霜雨露,雷電冰雹等.但是聖經也相當明顯地告訴我們,「曠野」是人生必經之路.若果我們要從「此岸」到「彼岸」,要出埃及進迦南,一定要經歷「曠野」之途.甚至要為主創造一番事業,就像主耶穌要出來傳道時一樣,也要經歷「曠野」的考驗.換言之,我們又不能完全把「曠野」與「城市」,「野蠻」與「文明」那樣放在對立面上去思考,因為「曠野」乃人生的一個階段,是一個從此岸到彼岸必須經過的旅程.站在對立面,我們只要選擇,但從發展的階段看,我們只能接受,無從選擇.循著這個途徑思考,我們對「曠野」就不能忽略,有必要要認識更多一些.因為每個人都需要經過人生旅途中的「曠野」,以及面對「曠野」的考驗.

除了是人生中的必經之途,「曠野」也是人心深處的幽暗之角.每個人,尤其是現代的文明人,心的深處總有一處留下對「曠野」的思慕和眷戀.曠野本來是原始和野蠻的地方,是野獸和魔鬼出沒的地方,但人對於其原始的本性,對自然的樸素與未鑿,總懷著一份眷戀之情,而這份潛藏心底的渴慕又往往成為突然迸發出來的動力.雖然在文明的壓抑之下,獸性之本能與對原始的渴望都受到克制,但並未根絕.時不時總要破網而出,表達及發洩一下.

在現代文明社會中,曠野與荒原並不多見,但其代替品卻不少.例如,在現代都市中的體育館與競技場,可以說是一種都市的「曠野」.一位魅力十足的美國歷史學家拿殊(Roderick Nash)曾經這樣分析美國的「曠野歷史」(history of wilderness): 他認為美國的體育運動和競技場是美國史上被人遺忘了的「曠野」,在早期歷史的「曠野」中,充滿了鬥爭,肉搏,競逐的精神和求生的意志.在「曠野」那裡,是野獸充斥和弱肉強食的地方.而在現代社會中,那種「開疆闢土」,「好勇鬥狠」的「邊陲精神」已經不再,現代人已經被文明教化和都市的群居生活弄得十分溫順,再也沒有昔日的曠野雄風.只有在競技場和體育館中,這種「邊陲精神」(Frontier spirit)才得以再次呈現.請看看美國足球和籃球:足球隊都多以勇猛的野獸為名,如「底特律獅子」(Detriot Lions),「芝加哥熊人」(Chicago Bears),「卡羅連納黑豹」(N. Carolina Panthers),「印第安納野馬」(Indianapolis Colts),「辛辛那提猛虎」(Cincinnati Bengals)等等,又或者用常常在邊陲與蠻荒出沒的「牛仔」(Cowboys),「紅番」(Redskins),「酋長」(Chiefs)作為名號；至於籃球場上,亦是由一群動物去鬥爭和廝殺,包括「木狼」,「灰熊」,「山貓」,「野馬」,「鷹隼」,「羚羊」等等,都在在表示競技場是一個「曠野」與「莽原」的翻版.在文明的時代和程式的生活中,「曠野」雖不能以蹴即就,但透過體育場館與競技搏鬥,至少在某個程度上可以使「邊陲再現,曠野重臨」.拿殊教授說:這是美國現代人保存邊陲及曠野時代的「開荒精神」的一個法子,一種文化.

換言之,用美國文化的角度去看曠野,那就是「邊陲精神」的延續了,從美國體育是角度看,曠野是美國文化(特別是西部文化)中的Frontier,是值得珍惜,值得保留的傳統.人的內心深處,往往流露著對曠野的嚮往,就是對原始的,本來的(或本能的)「我」的失落後的追思與懷想.

總的而言,我們今天談曠野,大概有三點:

一,「曠野」作為一個實體,在是聖經中常常出現,它是死海西邊的一塊沙漠山地,也是以色列人歷史地理的一個重要部分,時刻提醒著以色列人他們民族的苦難及奮鬥的經過,如何走來,如和逃脫埃及的統治,經過曠野四十年的漂泊,流浪,才有今天.

二,曠野做為人生旅途的一個中途站,是人生中必須經歷的一個階段,摩西河以色列人用了四十年在曠野漂流,而耶穌則再曠野過了四十天,你的人生曠野是長是短,是四十年還是四十天,我不知道,也許你也不知道,但我知道,這是你我必經之路.

三,曠野也在我們的心裡.向好的方面看,內心中的曠野的呼喚可以令我們想起我們的本性與童真,令我們重新振奮,恢復「邊陲」和「拓荒」精神,換另一方面去看,我們總部能完全擺脫我們原來的「老我」和原始的獸性,在文明洗禮之後,仍然時不時會反叛,會犯罪,會重尋舊夢,回到曠野去生活.

(2)談「行旅」

單單談題目中三種事物中的的一種──「曠野」,已經花了我們很多的時間,餘下我們只能簡單的談談「行旅」,至於「嗎哪」,看來今天再沒有時間處理了.我們大概都聽過「人生如客旅」這句話,但「行旅」和「旅人」又有沒有分別呢?大家以為會對題目中的「行旅」比較熟悉,因為我們往往會誤認「行旅」為「旅人」.上面已經說過,香港人喜歡旅行,故此對「旅人」經驗甚豐.然而,若你仔細看題目,留意英文的用詞,那麼,你可能會有另一番想法.

我們一般用「行旅」或「旅人」,都是指旅行的人(Traveller)但這裡用了Pilgrim一詞.一般上我們不會把Pilgrim翻譯成旅行者,而會把他翻作「朝聖者」,更準確的是「一個旅行的朝聖者」.朝聖者所不同於一般旅行的人在於其目的與心情皆很不一樣.一般的旅行者是看風景,采風問俗,增廣見聞,求的是眼目之怡悅和知識的增加,但朝聖者的旅行,則抱著一顆虔敬的心態,為的是要實踐一種人生的願望.

我們所知道的朝聖者,或者說在詞彙裡常用到「朝聖者」時,指的是伊斯蘭教的五個要求之一的「朝聖」,即去聖地麥加朝聖的行為.在伊斯蘭的教義中,信仰真主阿拉,每日五次祈禱,誦讀可蘭經,向窮乏人施賑,以及一生中至少一次需到聖地麥加朝聖,是五個基本的要求.故此,每年都會有數百萬乃至過千萬的回教徒會到麥加,完成其朝聖的工作.這種朝聖,需要多年的籌劃和準備,而且在生活上,衣著上,食物上,都要配合,才能手潔心清的去朝聖,一個伊斯蘭教徒往往為了朝聖的旅行,會花上一年甚至好幾年做準備.近年我們看到新聞報導,每年的麥加朝聖團都會因為人流擠擁或其他原因,令到意外頻頻發生,不少人因此而喪失生命.這樣麻煩,費時,而又危險的朝聖之旅,還是有無數的信徒參加.原因是朝聖是他們生命中的一個重要目標,是他們生命信仰中的一個重要的不可或缺的部份.換言之,朝聖之旅,並非一般的旅行,它與信仰有關,與生命攸關.

另外,我們常常聽到的Pilgrim,則和美國開國早期的「清教徒」有關,這些早期從歐洲大陸移民美洲的清教徒,不是一般的旅行者,他們是逃避傳統皇朝的壓制和天主教勢力的壓迫,要到一個遙遠的新世界去追尋他們的自由和理想,坐著「五月花號」漂洋過海,經過狂風巨浪,來到尚未開發的美洲新大陸,土地是一片荒蕪,野獸與「紅番」當道(在他們看來土著是野蠻和未開化的番人,膚色棕紅,故曰紅番),旅行者不會在此久居,甚至商旅也只是以此地為過度.但「清教徒」則不一樣,他們選擇了新大陸作為永恆的家園,他們的選擇是生命的抉擇.所以這批Pilgrim,雖不是伊斯蘭的朝聖者,但同樣具備朝聖之旅的特徵:就是對生命與信仰的承擔.

所以,我們先要知道,今天講題中的「行旅」是Pilgrim,是朝聖者,是不怕翻山越嶺,橫過大漠,經過荒原,與走過曠野的朝聖者.換言之,人生是朝聖之旅,而不是一般的旅行.朝聖者向著一個終極的目標進發,惶惶然帶著一顆顫兢,敬虔的心,不懼艱難險阻,也不怕途中敵人的攻擊和魔鬼的試誘.這種心態,與遊山玩水的旅行者有很大的分別.人生亦復如此.大家通行一條路,甚至大家一起經過同意各曠野,受到同樣環境的限制與影響,但內中應有分別.這就在乎其人生之目標──是達到一處終點,抑或漫無目的?前者符合了Pilgrim之義,後者則仍然是「旅人」!人生旅途上,漫漫長路上行人無數,但誰才是朝聖者?

(3)談「嗎哪」

最後,讓我們來談談「嗎哪」,一條「三及第」的題目內的第三種元素.如果「曠野」是人生旅途中的一種環境,一個階段,「行旅」是人生旅途上的一種心態,一種承擔,那麼,「嗎哪」在人生旅途中的又是代表著甚麼呢?

毫無疑問,大家若是對基督教稍有認識的話,都會知道「嗎哪」是一種食物,一種頗為神秘的食物.剛才我們度過的聖經《出埃及記》第六章中就對「嗎哪」有下列的描寫:「到了早晨,營的四周有一層露水,露水上升以後,就見野地上有細小麟狀的東西,像小白霜一樣細小的東西.以色列人看見了,就彼此對問說:這是甚麼?原來他們不知道那是甚麼.摩西對他們說:『這就是耶和華賜給你們吃的食物,耶和華吩咐的話是這樣的:你們要按著個人的食量收取,按著你們的人數為帳幕裡的人收取,一人兩公升.』 以色列人就這樣行了；有的多收,有的少收.他們用升斗衡量的時候,多收的沒有剩餘,少收的也不缺乏；各任按著自己的食量收取.摩西對他們說:誰也不可把一些留到早晨.他們卻不聽摩西的話,有人把一些留到早晨,就生蟲發臭了.摩西就向他們發怒.他們每天早晨都按著個人的食量去收取,太陽一發熱,就融化了.」(13-21節)

從上面的經文看,我們對這種天上降下來的食物還是不能有一個清楚的了解,只知道它是在早晨露水蓋地,太陽還未出來之前降下,它的體積細小,顏色雪白如霜.另外一處經文則多加上一點描述:「它像芫荽的種子,色白,味道像摻蜜的薄餅.」(31節) 又說:「以色列家給這食物起名叫嗎哪.」(31節)

我在這裡不會對嗎哪的成份和性質作出猜想,在此只擬提出關於「嗎哪」的兩點:一是「時」,一是「份」.是的,「嗎哪」是食物,可以抵肚腹之饑.但許多東西皆可以有相同的作用,「嗎哪」之異於何處?曰「時」.所謂「時」者,適時也,即使在最需要的時候出現,在緊要關頭出現,以解決我們的困難.英文中有Perfect timing 的說法,而「嗎哪」的出現,正是Perfect timing,是神宰我們最為需要的時候給予我們的幫助.我們從上述經文看到「嗎哪」的出現,是在以色列人出埃及,過紅海,用盡了帶來的家當(包括食物)之後,此外,它的出現是在以色列人經歷過「瑪拉」苦水之後,是上主提供給他們苦後的甘甜.何謂「份」?「份」者有二義,一是「份量」,每個人有每個人的份量,不能多,也不會少；二是「本份」,甚麼時候收取,甚麼時候多收,甚麼時候不收,都有一定的規矩,守規矩,即盡本分.要吃「嗎哪」,還真的要遵守這些規矩!

在我們的人生中,我們常常求上帝賜福,這種福氣當然包括了食物,溫飽,成功等等.但有否想一想:甚麼是你的「量」和你的「份」?

今天講「嗎哪」,就簡單的替這兩點!

四,    結語時間已經用的夠了,我不想在做詳細的總結.但大家會不會留意到今天講道的一個弔詭?那就是:一個常常的,複雜的,具備三個元素的題目,連成一起的去講,我只用了幾分鐘的時間,可以說簡單而直接.而把題目分開來講,我卻用了大半個重點的時間,好像也沒有把主題很清楚的講出來,留下更多的問題未解決,美其名是說留給大家思考.其實,人生往往也是這樣.我們以為很複雜的,很麻煩的,很多件同時要去處理的事情,可能是最簡單直接的；而往往我們認為要分開來,逐一了解,逐一分析之後才有答案的,又往往是作繭自綁,甚至是自尋煩惱的.



\newpage

\section{餓孩子飽孩子}
\label{sec:10}
\textbf{鄧瑞強博士 - 餓孩子,飽孩子(Hungry Child, Well-Fed Child)}
\newline
\newline
\hyperref[sec:9]{< < < PREV SERMON < < <}
~
\hyperref[sec:toc]{[返目錄]}
~
\hyperref[sec:11]{> > > NEXT SERMON > > >}


約翰福音 6:35,41-51
\newline
\begin{longtable}{cl}
\hline
\hline
章節 & 經文 (和合本修訂版)\\
\hline
6:35 & \begin{tabularx}{0.7\textwidth}{X} 耶穌對他們說：「我就是生命的糧。到我這裡來的，絕不飢餓；信我的，永不乾渴。 \end{tabularx} \\ \\ \relax
6:36 & \begin{tabularx}{0.7\textwidth}{X} 可是，我告訴過你們，你們已經看見我，還是不信。 \end{tabularx} \\ \\ \relax
6:37 & \begin{tabularx}{0.7\textwidth}{X} 凡父所賜給我的人，必到我這裡來；到我這裡來的，我總不丟棄他。 \end{tabularx} \\ \\ \relax
6:38 & \begin{tabularx}{0.7\textwidth}{X} 因為我從天上降下來，不是要按自己的意願行，而是要遵行差我來那位的旨意。 \end{tabularx} \\ \\ \relax
6:39 & \begin{tabularx}{0.7\textwidth}{X} 差我來那位的旨意就是：他所賜給我的，要我一個也不失落，並且在末日使他復活。 \end{tabularx} \\ \\ \relax
6:40 & \begin{tabularx}{0.7\textwidth}{X} 因為我父的旨意是要使每一個見了子而信的人得永生，並且在末日我要使他復活。」 \end{tabularx} \\ \\ \relax
6:41 & \begin{tabularx}{0.7\textwidth}{X} 猶太人因為耶穌說「我是從天上降下來的糧」，就私下議論他， \end{tabularx} \\ \\ \relax
6:42 & \begin{tabularx}{0.7\textwidth}{X} 說：「這不是約瑟的兒子耶穌嗎？我們豈不認得他的父母嗎？現在他怎麼說『我是從天上降下來的』呢？」 \end{tabularx} \\ \\ \relax
6:43 & \begin{tabularx}{0.7\textwidth}{X} 耶穌回答，對他們說：「你們不要彼此私下議論。 \end{tabularx} \\ \\ \relax
6:44 & \begin{tabularx}{0.7\textwidth}{X} 若不是差我來的父吸引人，就沒有人能到我這裡來；到我這裡來的，在末日我要使他復活。 \end{tabularx} \\ \\ \relax
6:45 & \begin{tabularx}{0.7\textwidth}{X} 在先知書上寫著：『他們都要蒙神教導。』凡聽了父的教導而學習的，都到我這裡來。 \end{tabularx} \\ \\ \relax
6:46 & \begin{tabularx}{0.7\textwidth}{X} 這不是說有人看見過父，惟獨從神來的，他才看見過父。 \end{tabularx} \\ \\ \relax
6:47 & \begin{tabularx}{0.7\textwidth}{X} 我實實在在地告訴你們，信的人有永生。 \end{tabularx} \\ \\ \relax
6:48 & \begin{tabularx}{0.7\textwidth}{X} 我就是生命的糧。 \end{tabularx} \\ \\ \relax
6:49 & \begin{tabularx}{0.7\textwidth}{X} 你們的祖宗在曠野吃過嗎哪，還是死了。 \end{tabularx} \\ \\ \relax
6:50 & \begin{tabularx}{0.7\textwidth}{X} 這是從天上降下來的糧，使人吃了就不死。 \end{tabularx} \\ \\ \relax
6:51 & \begin{tabularx}{0.7\textwidth}{X} 我就是從天上降下來生命的糧；人若吃這糧，必永遠活著。我為世人的生命所賜下的糧就是我的肉。」 \end{tabularx} \\ \\
[1ex]
\hline
\hline
\end{longtable}
日期:2009年8月23日

證道:鄧瑞強博士

題目:餓孩子,飽孩子(Hungry Child, Well-Fed Child)

經文:約翰福音 6:35,41-51

各位弟兄姊妹,主內平安.

上個禮拜梁元生院長講道,講「曠野,行旅,嗎哪」(The Wilderness, The Pilgrim and the Manna),他的講道像在密密麻麻的草叢中打開一條路,讓人窺見人生的另一面.若果你在效野旅行,走在草叢中,忽然見到草叢的一段被人動過手腳,後面開了一條小徑,你很難不受引誘走入去看看.我就是那些經不起引誘而走入去看的人.今日的講道,是走在梁院長開的路之上的.有人開了路,後來者走起來也容易一點.

梁院長上次講道結尾時,講到嗎哪,講到與嗎哪有關的「時」與「份」.今日的講道,題目是「餓孩子,飽孩子」,與食物有關,故我以嗎哪開始,我要講嗎哪的本質.當以色列人出了埃及,走在曠野時,他們沒有太多食物.神看顧他們,在早晨,野地地面出現白色的小圓物,可以吃的.以色列人靠此充飢.他們不知道這些東西是什麼,便說這是「嗎哪」.「嗎哪」在希伯來文的含意,就是:「這是什麼」?這是一種與慣常所見的東西完全不同的東西,一種不知名的東西,一種不能定義的東西.「嗎哪」不能保留,也不能多要,多要了會變壞,故「嗎哪」不單不被人理解也不由人控制,不由人儲存.面對這種東西,只能說:「這到底是什麼來的?」無以名之,這東西就叫「這是什麼」(「嗎哪」的含意).嗎哪的本質,就是「不知是什麼」?

上次梁院長講的「曠野」和「行旅」(朝聖)其實也有「嗎哪」這種「不知是什麼」的元素.我們住在城巿,相對而言,「曠野」就是我們所不熟悉的領域,這領域總在城巿之外,也即在我們城巿人所能理解,所能掌握的領域之外.城巿人在曠野,頗為尷尬,想飲水卻沒有水龍頭,想涼快卻沒有空調.對城巿人而言,曠野是空間的「嗎哪」,即:在空間上不知是什麼的地方.

「行旅」(Pilgrim)即朝聖,是離開日常熟悉的生活及環境,走在陌生的路上,追尋那陌生的神聖.朝聖的路是陌生的,神聖的臨在也是陌生的.故此,朝聖可說是生活上的「嗎哪」,離開知道的,去經驗不知道的.而最終極的不知道的東西,就是神.

問題是:我們很不願意身邊有「嗎哪」,很不願意身邊有些「不知是什麼」的東西.試想想:你身邊有隻不知名的小動物在活動,你會如何?再者,「不知是什麼」意味著不受你控制.人很怕有些不受控制的東西在身邊,就如你睡覺時,你會怕有一個未懂性的孩子在你身邊,因為你不知什麼時候,他會「剝了你棚牙」或「打扁你個鼻」.在心理學上,人有潛意識,潛意識的利害,在於你不知它的內容是什麼而它卻影響著你.故此,你可能會無緣無故地傷感,無緣無故地失落,無緣無故地恐懼.要擺脫潛意識的控制,要透過種種方法揭露潛意識的內容,當你知道它是什麼時,你才能控制它.不知道是什麼,你就無法控制它.

神賜「嗎哪」叫人生命飽足,若我們不喜歡「嗎哪」,則怎能飽足呢?

我們不喜歡「嗎哪」,我們喜歡可認知,可量度,可控制的東西.

我們不喜歡「曠野」,我們喜歡商場.曠野的野花,我們叫不出它們的名字.但商場裡的LV,我們老遠便認得.曠野的石頭,我們看不到它們的價值.但商場的商品,每樣東西的價值,我們都心中有數.我們喜歡商場,這是我們的世界.

我們也不喜歡「行旅」,我們喜歡「旅行」.「行旅」是會遭遇不測的.「不測」即「測不透」,即「不知是什麼」,即「嗎哪」.我們將「行旅」變成「旅行」.「旅行」,則一切在計劃之內.最好,連神聖的顯現,也包在行程內.其實,在現代什麼也可控制的科技下,要製造神聖的顯現的虛擬實在,又有多難?

若崇拜本身就像朝聖的「行旅」的話,則神可能會在不同環節中突然進入信徒的生命中.至於,神如何顯現,神如何工作,這有點像「嗎哪」,我們不知道,也不能掌握.但是,有人認為,將崇拜的「行旅」性質變成「旅行」,豈不更好?將神的感動變成可操控的環節,豈不更好?連神什麼時候出場也計算好,這豈不更好?於是,在不同的環節,刻意加入感動的元素,以音樂將情緒推向高峰.然後,保證每個人都有感覺,保證人人的情緒都得到滿足,並且將這感覺命名為神的同在.多麼完美!人喜歡這樣,正如我們喜歡一個吃喝玩樂皆滿分的旅行團一樣.我們不希望神的臨在像「嗎哪」一樣不可控制,我們希望連神的臨在也被控制.

約6:35記載,

耶穌說:「我就是生命的糧.到我這裡來的,必定不餓；信我的,永遠不渴.」

主耶穌就是生命的糧.到他那裡去的,就不飢餓；信他的,就不乾渴.

多麼簡單,就可以成為一個飽孩子.

只要領受生命的糧,就可以飽足,多麼簡單.

但問題是:耶穌是真正的「嗎哪」,而我們卻不喜歡「嗎哪」.耶穌不受控制,但我們渴望耶穌受我們控制.我們不要獨行奇事的耶穌,我們想要一個可利用他來感動我們的耶穌.我們不希望耶穌來挑戰我們,我們卻願意奉耶穌的名去挑戰他人.

耶穌是真正的神,但我們卻喜歡偶像,因為我們控制不到神,卻能控制偶像.若耶穌是「嗎哪」,則偶像就是我們製造的零食.我們可以調校零食的味道,你喜歡朱古力味就朱古力味,你喜歡芒果味就芒果味,零食最能滿足人的慾望,但對身體無益.耶穌講什麼變得不重要,最重要的,是耶穌聽我所講的.耶穌對我的要求變得不重要,最重要的,是耶穌實現我的要求.恩典太抽象,最後最好能變成可控制的感受.只要我流淚,就表示神的愛是真的,就表示神的恩典臨到了我.我的感受才是真的,神到底愛不愛我反成為其次.當然,人間的愛情也是這樣.

耶穌是「嗎哪」,是一種「不知是什麼」的東西,但是,我們想要的,是能被我們控制的東西.遠離「嗎哪」,人就很難接近耶穌.

6:41　　猶太人因為耶穌說「我是從天上降下來的糧」,就私下議論他,6:42　　說:「這不是約瑟的兒子耶穌嗎?他的父母我們豈不認得嗎?他如今怎麼說『我是從天上降下來的呢』?」

猶太人將耶穌看成完全可被掌握的對象,他們掌握耶穌的身世,掌握耶穌的社會背景,掌握耶穌的來龍去脈,他們以自己可掌控的東西來理解耶穌,當然看不到耶穌的神聖,故此,也不可能接受他.這些人很難得到生命的飽足.

每樣事物都有其「價值」(value),但「價值」像「嗎哪」,很難掌控,人為了方便掌控,便將「價值」變成「價格」(price),經過量化,一切便可比較,可被掌控了.一隻很美的狗,我們欣賞牠,牠有一種無法掌握的價值.但只要這隻狗有一個價格,如:值三百元,則我們便掌握了牠和其他狗隻在巿場上的價格差異,我們也知道我們是否有能力擁有牠,是否想擁有牠.

事物的價值原本像「嗎哪」,不能被我們掌控,但我們不喜歡「嗎哪」,我們不喜歡有一種「不知是什麼」的東西.我們試圖將「嗎哪」定一個價.一個女孩子的生命,有無限的價值.若有人向她問價,而她卻又還一個價的話,則一切價值失去了,只留下價格.生命的一切神秘也失去了,只留下可被操控的驅體.生命的豐裕失去了,只剩下赤裸裸的貧窮.朝聖的生命失去了,因為神聖隱藏了.曠野也失去了,因為只留下商場的一買一賣.「嗎哪」沒有了,那個「不知是什麼」的神秘沒有了,只留下沒有價值的價錢牌.失去了對事物價值的欣賞,生命將飢餓得無以復加.

有人反對在書展裡有美少女的身體展露,這些人想保持「嗎哪」的純潔.天父賜人不可保存的「嗎哪」,為叫人珍惜事物的永恆價值.有人認為書展有美少女,正常得不得了.這些人認為,「嗎哪」畢竟是有個價的,沒有價,怎能明白它是什麼.沒有了「嗎哪」,人怎能飽足?

哲學家西美爾(Simmel)寫了一本書,叫《金錢哲學》.他指出,在金錢世界裡,由於一切事物都以金錢標價,每樣事物的真正價值反而隱藏了.若果一罐可樂的價錢貴過一罐汽油,則我們很容易便認為,一罐可樂的價值貴過一罐汽油.若果一個人的巿場價值貴過另一個人,則我們很容易便認為,這個人的價值重要過那個人.我們若用價錢去理解所有事物的話,則我們將失去事物本身的價值.失去了事物的價值,我們的生活變得空洞,心靈亦變得飢餓.現代人總是覺得不知如何獲得豐足的生命,是因為我們失去了事物的價值,是因為我們將天父賜的「嗎哪」變成商品.

「嗎哪」是拒絕被訂價的,因為這是一「不知是什麼」的東西.對於一「不知是什麼」的東西,你如何能訂價呢?只有當我們放下價格的比較時,我們才能學會欣賞事物的價值.

主耶穌講過一個比喻,是這樣的:牧羊人有一百隻羊,其中一隻走失了,牧羊人會丟下九十九隻羊,找那迷失的一隻.我小時候聽到這比喻,總會想,那九十九隻羊怎麼辦?若有野獸來,怎麼辦?牧羊人這樣做,不化算,可能得不償失.我的這種想法,仍是一種定價錢,量高低的想法.但比喻中的牧羊人,他不是比較價錢,衡量得失,而是看到每隻羊的價值.一隻具有無限價值的羊失去了,其損失是無限的.找到這隻羊,獲得的價值也是無限的.

神天天賜人新鮮的「嗎哪」,人天天領受從天父而來的恩典.生命多麼飽足.但人不喜歡這些不被控制,不能收藏的「嗎哪」,人還是喜歡自己造餅,努力收藏,縱使這只是塑膠的餅,不能叫人飽足.

6:46　　耶穌說:這不是說,有人看見過父,惟獨從神來的,他看見過父.6:47　　我實實在在的告訴你們,信的人有永生.6:48　　我就是生命的糧.

主耶穌見過父,他的生命與賜「嗎哪」的天父完全合一,他的生命流露著天父的美意,他的生命也就是天父賜給世人的「嗎哪」,這就是天父無窮的大愛.愛,不能量化,不能掌控,卻肯定被愛者的無限價值.領受這真「嗎哪」的人,就是能愛的人,在施予中體會生命的飽足.

面對「嗎哪」,你如何領受?你是一個餓孩子,抑或是一個飽孩子?



\newpage

\section{沒有群羊的世界}
\label{sec:11}
\textbf{Prof. Dennis P. McCann - 沒有群羊的世界,牧者有何用呢?}
\newline
\newline
\hyperref[sec:10]{< < < PREV SERMON < < <}
~
\hyperref[sec:toc]{[返目錄]}
~
\hyperref[sec:12]{> > > NEXT SERMON > > >}


日期:2009年8月30日

證道:Prof. Dennis P. McCann

題目:沒有群羊的世界,牧者有何用呢?

羊與牧人二者的形象,構成了基督徒意想中很重要的一部分.其中最早的一個形象,就是逾越節羔羊所象徵的耶穌基督.耶穌在十字架上所受痛苦和死亡的意義,透過了猶太人逾越節禮儀中所獻上的無辜羔羊已經象徵出來了.耶穌復活之後,在委任彼得一個重大使命,要他關照基督徒會眾時,也親自提到羊:「西門彼得,你愛我比這些更深麼?你餧養我的小羊,牧養我的羊.」(約翰福音21:15-19)

照約翰福音前幾章所記,耶穌讓人想起祂曾宣稱自己是好牧人,而好牧人為他的羊捨命.(約翰福音10:1-18)在今天所唸的經文裡,我們看見牧人與羊的形象,已經深深的植根在古時猶太人的歷史裡.這可以從名為大衛王的詩篇23篇裡集中表現出來.「耶和華是我的牧者,我必不至缺乏 and \#8230;」,同時,我們也聽見先知耶利米譴責那些作為牧者的以色列君王的言詞.他們「趕散我的羊群,並沒有看顧他們」(耶利米23:2).

那些常來此處崇拜的人,如果覺得這種說法不過是老生常談的話,不妨回憶一下十個星期之前在"好牧人的主日"所聽到的.然而我們為什麼還舊調重彈,一再提到牧人與羊群呢?到底還能多說些甚麼呢?

也許我們需要更具體的思考一下有關羊與牧人的事,好使我們藉著它能夠更清晰的聽見神所要向我們說的話.世界上羊群的數量,並非分佈得那麼平均,你在香港曾看見過羊隻麼?山羊也許是有的,但卻不是綿羊.

根據聯合國食物與農業組織(FAO)2004年的記錄,中國目前是擁有最多綿羊的國家,總數幾乎達到1億5千8百萬隻.當然我們有興趣想知道,這大部分的羊群是否散佈在中國西部省分；在紐西蘭來說,和羊群接觸,簡直就是司空見慣的事,因為羊群的總數比起人口來,幾乎要多出四倍.同樣的,澳洲是繼中國之後,擁有第二大羊群數目的國家.所以你若是對羊群或牧人的事不熟悉的話,也毋須感到難為情.事實上,二者在香港是很罕見的.

在這一方面看來,我們今天的生活與經歷,與聖經所講論的世界是很不一樣的.由於我們身處於一個沒有羊的世界裡,我們對牧人可能沒有甚麼感覺.此外,雖然漢字中的"羊"字與"美"字有點關聯,如許慎的《說文解字》所謂「羊大衛美」,但在英文來說,"sheep(羊)"and"sheepish(像羊一般的)"二字的意義,卻是模稜兩可,有些矛盾.從某一方面看來,主在馬太福音25章31-33節裡曾設個比喻說:世界末日審判的時候,祂會把綿羊山羊分別出來.在英語的一般的慣用法上,綿羊都被視為良善的,因為牠們的天性馴良,易於駕馭,但山羊卻被視為邪惡的,因為牠們是野性難馴,不受羈絆的,因此,山羊的角有時甚至是被當作魔鬼的象徵.從另一方面說來,英語裡有一個字叫"sheepish"(像羊似的),含有阿諛奉承的意思.說一個人"像羊一般",是表示他像羊那樣溫馴,愚拙及懦弱,膽怯.如果你的舉止"像羊一般"(即所謂"羊手羊腳")的話,那就表示你做事尷尬異常,過分羞怯,甚至乎被人奚落,成為不受歡迎之人.其實像羊一般,舉止像羊,就表示幼稚,做事不經大腦,或是人云亦云,亦步亦趨,全無主見."像羊一般"一詞,在英語中既這種聯想,有誰還願意被視為"像羊一般"的人呢?

可是,牧人與羊二者,卻是並行不悖的.真正的羊是需要有人看管的,牠們需要有人來牧養,因為這是保謢牠們免被殺害(像被狼所吞噬)的唯一辦法.狼是相當精明的,牠知道只要將羊群驅散,就可以為所欲為了.當羊群被分散,孤立的時候,就容易成為牠的餐點了.羊固然很容易被率領,但牠們也容易因為沒有牧人把牠們聚攏在一起而迷失了,牠們也可能會驚慌而分散.在某些方面,羊可能有點笨拙,但牠們能夠取得真正牧人的愛謢與忠誠.一個牧人可以冒著性命的危險,把一隻羊從豺狼的虎口中搶救回來,這個說法,是不會太過離譜的.至於雇工,他們看守羊群,只不過是為了一天的工資,沒有人會企望他們盡心竭力的去牧養羊群.雇工究竟是雇工,他不能像真正牧人那樣去了解羊群.

那些研究這類事物的人告訴我們,羊是最早被人類所馴養了的動物中一種,距今至少在一萬年之前,始於中東或亞洲的西部,羊與牧人很容易就建立起一種彼此受惠的關係.雖然有一些的羊,尤其是一些幼小的羔羊,偶爾會被宰殺成為盤中餐,但只要牧人對羊照顧周全,並且不太貪心或不負責任的話,整個的畜群必會繁盛起來的,因為只要你關心羊,牧養羊,羊是多產的,牠的數目就會跟著加增,那時候照顧羊群的牧人之財富,也必然相應的加增.牧養羊群,從羊群身上剪取羊毛去製成各種有用商品以供人們享用,是屬於最古老的一種工業.畜牧的人進入由農民所建立的鄉村,鄉鎮以及城市中,就是為了貿易和謀生的緣故.由於從事畜牧與農耕的人所過的生活方式和需要不同,因而常常導致利益衝突,例如,早期的埃及人,不願意有羊在他們的田野間放牧,就把羊群及牧人局限在歌珊一地,歌珊位於尼羅河的東面,朝向西乃山,是一塊不適合農耕的地方.雙方雖有這些爭論,但也正因為有這些歧見,貿易就成為一種可行及渴求的事,最終人類的文明也隨之出現了.

當我們從以上的觀點來觀察聖經上許多有關羊與牧人的故事時,我們可能就領會到先知耶利米何以用很強烈的筆觸,將神對於背信棄義的以色列,猶大君王所發出的憤慨描述出來.因為他們沒有對神及百姓盡忠職守,神就會把他們的國交付給他們的敵人,藉以懲罰他們邪惡的行為,但在毀滅之餘,耶利米還預言:被擄剩餘的人最終還將歸回本地,且有新的牧人王者 and \#8230;一位新興的大衛王來治理他們,以取代過往的舊君.從聖經的觀點看來,我們得知羊與牧人的形象,確實是帶有政治性的,它們是依據聖經給人提供一種表達對腐敗和不義所發憤慨的方式.質言之,即對於那些矢言要為人群服務,卻墮入腐敗和不義深坑的人,要如何表達神的義怒.在此同時,它們也給人帶來一種表達希望的途徑,期望新的領袖們能遵循誓言,在他們的身上體現神所應許的公義.

聖經中含有政治性的意味,並非世俗觀念學家所發明,為了要分散我們敬拜神的注意力.在馬太福音書裡,作者記載耶穌餵飽五千人的故事,再一次使我們想起好牧人的意義.我們可以想像一下:當時的群眾,身穿粗衣,手提毛毯,成群結隊,簇擁著耶穌,從遠處望去,不正像一群被放牧在山坡上的羊嗎?可是,他們卻一直沒有食物可喫,直等到那位牧人前來餵養他們.出現如此大批的群眾,非同小可,絕不是一樁個人的私事.在這裡,作者刻意提出一個新的政治形像,在即將來臨之神的國度裡要出現的.耶穌餵飽了那些前來聽祂講道的人,正扮演了耶利米預言中所指的那位善牧君王.

新約以弗所書給我們提供了一幅有關這個新政治方案的發展程序和有承諾效果之圖像.神國度展開的時候,它會將分割猶太人和外邦人之間的歷史冤仇廢掉,並彌補二者之間的鴻溝:「因他使我們和睦,將兩下合而為一,拆毀了中間隔斷的牆.而且以自己的身體,廢掉冤仇,就是那記在律法上的規條.為要將兩下,藉著自己造成一個新人,如此便成就了和睦.既在十字架上滅了冤仇,便藉這十字架,使兩下歸為一體,與神和好了.並且來傳和平的福音給你們遠處的人,也給那近處的人.因為我們兩下藉著他被一個聖靈所感,得以進到父面前.這樣,你們不再作外人,和客旅,是與聖徒同國,是神家裡的人了.並且被建造在使徒和先知的根基上,有基督耶穌自己為房角石」(以弗所書2:14-20).在神的國度裡,將來只會有一群羊和一個牧人──那就是這個新政治的期望和目標.

或許我們認為自己並不是羊,也不需要一個好的牧羊人.不過,縱然羊與牧人的形像與我們個人的經驗尚有一段距離,但一旦這形像的意義被解開之後,我們便曉得他們的力量及智慧.我們也許不是羊,但我們是一群人,因而我們便需要有適當的領導人,要不然我們就會被分散,受干擾,受困惑,害怕那些與我們不一樣的人,甚至乎易於陷入驚慌與暴力之中.

我原是一個外來的客人,一個寄居者而已,但多年來我卻深愛著香港和中國,至於對某些事件的是非得失,我是不適宜加以評論的,例如樞機主教陳日君先生一年前在香港回歸十一年的記念日裡,引述聖經以西結書34章,告誡那些不忠心的牧人,質疑政府拖延在港全面推行民主制度一事.這個說法的對與錯,應留待你去作出判斷,同時,你的良心也會在你為正義奮鬥時給你作出正確的指引.再者,從以弗所書所描述有關猶太與外邦之間的冤仇看來,最引人注目的,莫過於今天新疆人民(漢族與維吾爾族)所受的.假使我們認真正視耶穌的門徒以族群合一為他們所期望之事的話,我們又該為新疆做點甚麼事呢?我們是否只為在暴亂中遇難者來禱告便了事,而不做任何可以做的事以確保事件不再重演呢?假若羊群當中有漢人及維吾爾族人,身為一位真正的王者牧人應該如何處理呢?你真的以為他只要把綿羊從山羊中分別出來,就大功告成了麼?假使新疆的事過於遼遠,難以考量,何不把焦點放在香港?為何最近一位菲律賓女傭因患了豬流感(H1N1)而被送到醫院去,就迅速引起社會上某些人強烈的抗議,把大家對菲傭的抱怨,牢騷(不管是真實的或是想像的)都一併的發洩出來?你是否真的相信我們近這幾星期以來目擊的種種反應,只單單表達一種對衛生本身的關注麼?

或許由於恐懼心理作祟,我們想不到自已根本就像羊一樣,可是,要成為一隻羊,而又不類似羊一般的膽怯,懦弱,真是談何容易,更進一步說,我們要是成為一個好牧人,誠心誠意的關心自己的羊,盡自己所能的去認識牠們,親近牠們,這便更富有挑戰性了.這就是神呼召我們和我們的領袖們所要去體現及實行的事.願神賜我們勇氣與智慧去聆聽祂的話語,並且靠著祂的恩典作出回應,使我們順著祂的帶領來走祂要我們走的路.阿們!



\newpage

\section{我的肉是可喫的}
\label{sec:12}
\textbf{盧龍光牧師 - 我的肉是可喫的,我的血是可喝的My Flesh Is True Food and My Blood Is True Drink}
\newline
\newline
\hyperref[sec:11]{< < < PREV SERMON < < <}
~
\hyperref[sec:toc]{[返目錄]}
~
\hyperref[sec:13]{> > > NEXT SERMON > > >}


日期:2009年9月13日

題目:我的肉是可喫的,我的血是可喝的My Flesh Is True Food and My Blood Is True Drink

證道:盧龍光牧師

經文:約翰福音6

請大家一同禱告:慈愛天父,求你的話語在我們當中,求你潔淨我們的心,又潔淨我們的思想,打開我們的心來迎接你的話,奉耶穌名求,阿們!

各位弟兄姊妹,如果大家回來在這裡崇拜的話,都知道我們在這段時間,特別是研讀約翰福音第六章的經文,剛才大家看的這段經文,不曉得大家有沒有甚麼特別的想法,因為這段經文有幾句話,好像很令人吃驚的,就是我今天所講的這個題目(也是這經文所出現的):我的肉是可喫的,血是可喝的.

不曉得大家唸到這經文時有甚麼想法?耶穌怎麼講得那麼狠呢?好像菇毛飲血的時代,未開化的人,為何會用「我的肉是可喫的,我的血是可喝的」呢?大家如果有聖經的話,請重新打開約翰福音第六章整段的經文,而我們今天要看的,特別是51-58節.

在這段經文當中,其實由第六章是「五餅二魚」,從餵飽五千人開始的,而整段經文都是強調耶穌不斷重覆一句說話:我是從天上降下來,生命的糧.這句話本身已經夠抽象,耶穌是天下降下的,我們如果信耶穌是主,是基督的人,當然我們會接受.因為約翰福音有一個主題,就是「道成為肉身,住在我們當中」,耶穌正正就是那個道,成為肉身的這位生命.但怎麼會講到祂是生命的糧?我們都會認為這只是一種比喻而已,但為何說得那麼確實呢?說到祂的肉是可喫的,祂的血是可喝?

如果大家留意的話,這種說法在其他的福音裡頭是沒有出現的,只有約翰福音特別用這個字眼來描寫耶穌的.如果大家唸聖經覺得枯燥呢,就是因為大家看不到聖經的多元化.每一次看聖經,這個故事就是這樣子來教訓我們的... 那你一定感到枯燥的,但你讀聖經看到這些經文時,你該警醒,為什麼祂的肉是可喫的,祂的血是可喝呢?為何只有約翰福音如此說,而其他福音不這麼來形容呢?正正因為耶穌只有一位,但在祂在世的時候,祂的門徒同樣接觸祂,但原來也有很多不一樣的描寫與經驗,他們自己也有他們自己關於耶穌的故事,而三卷的福音書,我們稱為「共觀福音」,是很相似的──馬太,馬可,路加.但你仔細地去看,在相同當中,亦有不少的差異.而約翰福音就與三卷的福音書是很不一樣的,大部份經文,並沒有在其他福音書中出現.其實不單止是第六章這段,「五餅二魚」──幾卷福音書都出現,但這五餅二魚,餵飽五千人的解釋,在約翰福音裡是截然不同.如果大家留意約翰福音的話,還有一樣也同樣令人震驚的,就是耶穌最後的晚餐.

大家知道馬太,馬可,路加所講有關「最後晚餐」,亦是我們今天聖餐的基礎.我們每一次吃聖餐的時候,都是記念著主最後的晚餐,我們擘餅,分杯,是代表耶穌釘十字架,在祂最後晚餐,最後的行動中,提醒我們,我們以後也會如此繼續的做.但是看看約翰福音,沒有特別地去記載這最後的晚餐.沒錯!都有提到祂釘十架之前,與門徒有一個晚餐,但它最重要的信息是說甚麼呢?是說到耶穌為門徒洗腳,而這個行動亦是其他三卷福音書所沒有記載的.所以有些聖經學者,就說或許是吃了兩次晚餐(每天都要吃的嘛),當然我們覺得這個說法,不是用來解釋那些差異的理由.

而很明顯在最後晚餐那裡,耶穌可能有很多的行動,祂可能亦有洗腳,祂當然有擘餅,不過在三卷福音書的傳統裡面,它們所強調的是耶穌的擘餅,分杯,是代表著逾越節的晚餐,這對猶太人來說是一個很重要的傳統,在那裡他們特別記載擘餅及分杯的經文,去代表耶穌基督的釘十字架.但在約翰福音的傳統裡,他強調耶穌這個洗腳的行動,可能就是最震撼的.所以我們發現,透過不同的經文當中,明白門徒們接觸耶穌的時候,有不同的重點和反應,但更加令人覺得奇怪的是,難到擘餅,分杯的事,在約翰的傳統裡顯得不重要嗎?其實我們今天唸的第六章,正正就是約翰福音當中,就是所謂擘餅和分杯的傳統.

這句話,很明顯的是指向耶穌,這個釘十字架,犧牲生命,亦是約翰福音的版本來解釋「最後晚餐」.但如果大家留意的話,那為什麼他不放在最後晚餐那裡說的呢?而要在第六章那麼早就提及呢?因為今天很多人唸聖經,以為福音書只是一種歷史,是耶穌的生平,如果你仔細地看的話,你當然知道,這當然不是耶穌的生平啦!因為耶穌的出生,只有在馬太,路加裡有記載,馬可及約翰福音是完全沒有記載馬糟的故事.就算是其他福音,也只是提到耶穌出來,當祂開始受試探,開始傳道,記載祂的說話,祂的行動,而打從祂出生直到三十歲左右的經文,只有在路加福音,提到祂在十二歲上聖殿的故事.所以我們知道福音書,不是記載耶穌的生平,而是記載耶穌的生命,祂透過說話,行動,去讓我們認識祂是誰.

今天我們對人的認識都是一樣,如果你是上學的,新同學去到新的學校,或許上班的去了新的工作崗位,甚至乎你第一次來這裡做禮拜的時候,你也不會認識太多人,你可能觀察周圍的人,聽一下他們在說些甚麼,看一下他們怎麼行動.我們中國說:聽其這,觀其行而去認識一個人.正正福音書是同一個生命,但它說的說話,行動,對不一樣的人就有不一樣的沖擊,對他的生命有不一樣的改變,對他的生命有不一樣的領受,雖然有許多共同點,但亦有許多差異點,就如同我們人一樣──我們當然有許多共同點,但我們的生命當中就有許多差異點.

而這個約翰福音,其實他將最後晚餐的擘餅,分杯,這生命的沖擊,分佈在他的福音書的內容裡,在第六章已將這個沖擊提出來了,這是第一樣我們該留意的,因為不單止只有約翰福音有這樣子的記載,只是提醒大家,約翰福音有許多不一樣,但仍有些相同的地方.譬如:五餅二魚餵飽五千人,約翰福音一樣有記載,同樣的,不單止記載五餅二魚餵飽五千人,亦說明耶穌如何透過五餅二魚餵飽五千人的記載裡,在六章十一節,你會發現有一句話是很共通的.

六章十一節這裡說:「耶穌拿起餅來,祝謝了,就分給那坐著的人．分魚也是這樣,都隨著他們所要的.」這裡說「祝謝了,就分開」,在五餅二魚的記載裡面,它這樣寫,在最後晚餐裡面,也是同樣的「祝謝了,就擘開分給他們」.而且在路加福音,還有另外一段也只有路加福音有記載(其他福音書皆沒有記載),講到耶穌復活以後,祂當時只向女的門徒們顯現,當時有兩個男的門徒,從耶路撒冷走向以馬忤斯的路上,這以馬忤斯路上的故事,是很令人震撼的,就是說他們在走出去的時候,他們討論耶穌復活的事,但他們沒有見過,他們只是聽說了,他們充滿了憂愁離開了耶路撒冷,往以馬忤斯去,在路上,耶穌向他們顯現,但特別的地方就是說,他們不認得耶穌,當耶穌問他們說,你們談得那麼熱烈又充滿憂愁的樣子,你們在談些甚麼呀?那兩個門徒就說,哎呀,你在耶路撒冷都不知道這重要的事,我們的老師釘了十字架,又聽說復活了,你怎麼連這重要的事都知道呀?聽起來很笑話!他們在問耶穌怎麼不知道耶穌的事,結果在路上耶穌跟他們談,從舊約開始談彌賽亞的事,直到已經天黑了,祂要繼續往前走,那兩個門徒就說,哎呀,不要啦!天已經黑了,我們一齊來吃晚餐罷!就是當耶穌拿起餅來祝謝了,擘開分給他們的時候,他們突然認出耶穌來.這個記載可以有不同的解釋,但是很明顯的就是,門徒見到耶穌的面卻不認識耶穌,祂只是像普通人一樣,只有當耶穌拿起餅來,擘開分給他們的時候,他們才認出耶穌來.

你發現這個拿起餅,祝謝,擘開...(可能孩子餓了,聽到有餅就很高興),所以你發現拿起餅,擘開,分給人,這個行動,不祗是在最後晚餐才記載,這個行動,在五餅二魚那裡有記載,在以馬忤斯的路上,在重要的那個點的時候,他們是憑這個行動,耶穌這個行動,認出耶穌來,在認出耶穌的那一刻的時候,他們的記載就是說,耶穌不見了.

原來我們今天所說的聖餐,不只是最後晚餐的記載,這個擘開的行動,不只是在聖餐裡面,也不是在最後晚餐裡面才是重要的,原來這是耶穌的記號,門徒認出耶穌這個行動很明顯,因為這個行動就像我們今天一些商品的商標一樣,你看到那個商標,你就想起那個產品.今天有許多人有不同的記號,他說的話,他的行動... 今天在聖經裡面,耶穌是誰,是最核心的問題,但這又可是知識性的,理論性的,但在聖經中這共同點,從甚麼行動能夠認出這是耶穌?就是祂拿起餅來,擘開,分給人的時候,所以這點不僅是個禮儀,不僅是最後的晚餐,也不是今天我們聖餐的一個行動,乃是問:這個行動代表了甚麼?而且這個行動是耶穌的記號,這個代表了,人在那在地方能夠認出耶穌.

所以第二點,在約翰福音裡面,他特別去解釋... 他這個解釋就是,耶穌是從天上降下來生命的糧.這個生命的糧是要讓我們得生命!今天甚麼叫「生命」?我們大家就認為,我們就是有生命,生物,其他的是死的,我們活著,我們「生命」想起來就是我們肉身的生命.所以在這裡就特別記載說,那些猶太人... 因為約翰福音很特別,他好像不是一個猶太人,所以說他們猶太人,那些猶太人,其實當時應該全部都是猶太人才對嘛!不過這是約翰福音的特點.他說,猶太人就彼此爭論,在第52節:

這個人怎能把他的肉,給我們喫呢.耶穌說,我實實在在的告訴你們,你們若不喫人子的肉,不喝人子的血,就沒有生命在你們裡面.

很明顯耶穌這個講法,不是講我們今天肉身活著的這個生命,他所提到的是我們裡面的生命,那甚麼叫「裡面的生命」?

今天我們人跟動物有很多不少的差別,但其中最重要的一樣就是,在所有的記載裡面,一隻貓,一隻狗在兩三千前的生活,形態,跟今天沒有甚麼兩樣,牠們沒有做衣服,蓋房子,牠就只是喵喵叫,有東西吃就很高興,牠就如此活著,幾千年都沒有變化,因為我們說,一隻貓,一隻狗,牠不會這種反思的能力,沒有想像的能力,所以牠的生活沒有甚麼變化.我在三十多年前在中文大學唸書的時候,我們神學院有很多門課可以上,也可以上別的系所開的課,當時我們其中有一位最重要的哲學家唐君毅先生要退休了,所以我去上他的課,到了今天,我最記得他說過的幾句話,其中一句就是:今天人跟其他動物最大的差異是甚麼?其他動物,像一隻貓,你不會跟牠說:貓你不乖,你要學做一隻貓；你要學做一隻狗,但我們人呢?從小孩子到現在,老了還在說:你都不會做人呢!你要學好做一個人!我們不是已經是一個人嗎?還要學做人呀?是啊!我們今天用哲學的講法就是,我們是一個being,原來人還有becoming,原來人會變化的,有不同的做法,想法,我們有不同的想像力,原來我們有一種超越的能力,我不只是今天的我,我會思想明天的我會是怎麼樣?而且我不只是今天人家認識了我,跟我自己認識的我有差異的.今天我們基督徒當然還有一個層次,在上帝面前的我是怎麼樣?甚麼是真我?我也不難想像我們會發展,所以今天,人一方面就是像貓,像狗一樣,我們要生活,吃飯,穿衣服,也怕天氣的改變,但貓天然就有牠的能力了,我們要用別的方法.我們每天的生活,每一個人都一樣,是基督徒,不是基督徒,都是一樣,因為我們就是一個人,我們為了生活,保存肉身的生命的這種需要.

但是人不同的地方就是,原來只有人會自殺的,為什麼會自殺呢?人不單止要問生存有甚麼意義.為什麼會自殺呢?因為覺得生存沒有意義.最近大概是會考放榜,現在很少人會為了會考而自殺了,我們以前唸書的時候,有些人因為考試考得太差了,沒有面目見老師,父母,覺得他自己沒有希望了,他就自殺了,那是很笨的.像我,很久已經沒有人問我會考的成績了!為一個當時很重要,之後不重要的東西而死,我們覺得,哎呀!真笨!有些年輕人為了愛情就自殺了,我覺得其實半年以後,你就忘記了他嘛,為什麼要為他而自殺呢?但人在當下覺得很重要,其實經過時間,可以變成不重要,因為人就是有這樣的特點.我們有今天,有現在,我們有肉身的生活以及我們的生命,但更重要的是問:我們活著有甚麼價值,有甚麼意義呀?

今天耶穌在記載中所說的「生命」,不只是這個生存的生命,乃是問你裡面的生命,因為裡面會問自己:我活著為什麼?有甚麼意義,有甚麼價值?我們生存就有兩個層次,所以今天在這裡,耶穌一方面在約翰福音裡面,祂特別用這比喻說:「我就是生命的糧」,因為你們不單需要食物,就像耶穌在馬太福音提到祂受試探的時候,四十天沒有吃東西,當撒旦來試探祂的時候:你假如是上帝的兒子,你可以將石頭變為食物呀!但耶穌好像回答得很奇怪!因為假如我們按照言語分析的話,耶穌既然被撒旦問:假如你是上帝的兒子... 假甚麼如呀!祂本來就是嘛!這怎麼算試探呀?祂說,我就是!不就行了嗎?繼續問:假如你是上帝的兒子,你有能力將石頭變為食物呀!那也算甚麼試探呀?假如換我來問這試探,就試試看能不能將石頭變成食物?當然是不能啦,但我仍想試試看!但耶穌,當然祂有這種能力呀!為什麼耶穌不直接回答呢?我們今天很多人對耶穌的認識就是... 耶穌說,是就說是,不是就說不是嘛!對不對?那你是神的兒子,就說我是就對了嘛!問你有沒有能力將石頭變成食物,我能!不就行了嗎?就變給牠看,這不就完了嗎?不是喔!耶穌怎麼回答?為什麼是一個試探?其實這個試探最主要的重點,就是耶穌整個生命的使命重點,因為當祂被釘十字架,祂死裡復活,在路上,門徒／人都問祂,你是誰?耶穌在該撒利亞,腓力比的地方,祂也問祂的門徒:人家說我是誰?他們之間有很多答案,有些人說你是以利亞,先知裡的一位,好像那個被殺的施洗約翰復活... 但耶穌就問:你們說我是誰?只有彼得敢回答:你是基督,永生上帝的兒子,馬太福音裡頭記載,耶穌馬上就稱讚他:哎呀!你答得真對呀!這不是你血肉之軀回答的,是天上的指示回答的.彼得當然很高興,馬太福音繼續講,天國的鑰匙等等... 但那個時候,耶穌開始說,我要被害,我要受很多的苦,我還要被釘在十字架上,會被害死,然後復活... 彼得立刻反應說:那是萬萬不可!怎麼可以這樣?那耶穌就說:撒旦,退到後邊去吧!搞甚麼?幾分鐘前才那麼厲害地稱讚他,幾分鐘後就變成撒旦了(雖然彼得他不在香港受教育,他是跟我們在香港受教育的情況差不多,因為我們的學生都很會答問題,但理解可能錯的).

彼得很明顯,他的心中:是啊!耶穌你是基督,永生上帝的兒子,但是怎麼樣能夠去承認這個耶穌,是基督,是上帝的兒子?耶穌就告訴他:我原來要受苦,受害,被釘在十字架上,這樣我才是基督,上帝的兒子.嘩!彼得說,萬萬不可!因為他心裡面所理解的基督跟永生上帝的兒子,不是這樣子的,應該是高頭大馬,應該坐在寶座上的,是很有權威的,怎麼可以釘死,被害?他抓住祂:萬萬不可!

今天這些經文跟所唸的都一致的,耶穌你是誰?撒旦就說:假如你是上帝的兒子,你就將石頭變成食物,那人家不就認識你是上帝的兒子了嗎?在整個記載耶穌的事蹟裡面,祂能夠有能力──「五餅二魚」餵飽五千人,而且不斷有神蹟出現,但耶穌從來沒有一次,無論是醫病,趕鬼的神蹟,或者是五餅二魚的神蹟,做完了以後,祂是靜悄悄地走了.假如我有這種能力的話,我當然應該說:你們都吃飽了嗎?誰願意跟從耶穌的舉手?大概4999個都舉手,如果你問為什麼有一個不舉手,通常總會有一些人是不舉手的嘛!耶穌沒有這樣子做啊!耶穌靜悄悄地走了!祂不為那些為餅而來的跟從祂.

其實我們今天對耶穌的認識,不只是從祂的教導──看聖經認識祂的.聖經裡面,很明顯地除了耶穌的講話以外,祂是透過行動來表達祂是誰.而祂最重要的記號,在該撒利亞,腓力比的時候,祂告訴祂的門徒說:我要被害,釘死,然後從死裡復活,這個才是重要的標誌,也是今天為什麼基督徒在不同的地方總有一個十字架作為記號.但今天這個記號越來越漂亮了,像我們這個禮拜堂的十字架也蠻漂亮的,但肯定耶穌所釘的十字架不是那麼漂亮的,是一個爛木頭來,充滿恐怖的感覺的,是代表死亡,痛苦,但今天基督徒,或是新教的十字架,沒有耶穌釘在上面.天主教與基督教的分別就是,有耶穌釘在十字架上的,就是天主教,所以今天有些基督徒在笑天主教或是反對天主教說:你看你們的耶穌還在釘十字架,我們的耶穌就不一樣了,我們的耶穌已經復活了,所以不再釘在十字架上了.但其實真正的耶穌,這兩樣東西都同樣重要!今天我們往往都忘了祂釘十字架,只看祂的復活,榮耀,就跟彼得很像,他所期望的基督,永生上帝的兒子,就是那樣的榮耀的基督.但忘記了基督必須需釘死,然後才從死裡復活,沒有見過十字架的苦,就沒有復活的榮耀.

當耶穌在約翰福音,祂特別提到,假如你不喫我的肉,喝我的血,你就不能得裡面的生命.裡面的生命是代表生命的意義,價值,為什麼而活?比只是活著更重要?所以耶穌回答撒旦的試探也是一樣:人活著不是單靠食物!你要靠食物才能生存,但不是單靠生存,乃是靠上帝口裡的話!

我們小時候,主日學老師就說,你們要多看聖經,上帝的話就是聖經,所以你對付撒旦就是用聖經來對付牠.但我想起,老師忘記第二個試探,因為第二個試探,原來撒旦也看聖經呀!撒旦用聖經來試探耶穌:聖經不是說,你從殿上跳下來,上帝的聖者會托住你,那你就在殿上跳下來給我們看看,不就立刻證明你是上帝的兒子了嗎?「假如你是上帝的兒子」,「經上記著說」... 其實那個試探很明顯,也是問耶穌,你用甚麼方法,途徑讓人認出你是上帝的兒子?而最好的辦法就是跑到聖殿的殿上,也是整個耶路撒冷的最高點,最引人注意的點,一跳下來竟然沒有死,哎呀!人家不都認識你是上帝的兒子了嗎?(前兩年,在香港有一個很大的佈道會──葛培理的兒子──「假如」他能夠跑到IFC──國際金融中心,全香港最高點,所有電視台都去拍這個鏡頭,全香港都注意他了,假如他跳下來而不死,立刻就會6百多萬人或是一半都舉手決志了,但他沒有這個能力.耶穌也沒有這樣做,因為耶穌不要用這所謂超自然的神蹟來吸引人,讓人認出祂是上帝的兒子.但今天我們很多基督徒,只是盼望一些超自然的東西來令我們信耶穌,今天在香港的教會裡面也一樣,很多強調神蹟,醫病,趕鬼... 這個沒有問題,但假如是用這個方法來吸引人跟從耶穌的話,就有問題了,因為耶穌自己就拒絕了!用這種神蹟最快的方法,一跳下來,就證明祂是上帝的兒子──耶穌沒有選擇這條路,耶穌所選擇的,是那條十字架的路!

今天我們看見耶穌說:我是生命的糧,你們要喫我的肉,喝我的血,甚麼意思呀?因為那個就是代表祂自己.今天我們基督徒的生命價值是用甚麼來衡量呢?不是用會考有沒有考取三十分來衡量,不是你唸中文大學來衡量,不是你做甚麼職位來衡量,不是看你的收入來衡量,乃是你的生命裡面有沒有十字架的記號,而十字架的記號是代表著耶穌生命中最重要的意義.我們要喫祂的肉,喝祂的血,所代表的乃是我們要完全認同祂,跟祂連結在一起,因為約翰福音第六章,十五章,都很強調:是我們的生命有沒有跟耶穌連在一起?我們有沒有在祂裡面,祂有沒有在我們裡面?完全在生命的認同.在其他幾卷福音書裡面很強調:假如你不背起十字架來跟從我,你不要跟我來!

今天很多人都喜歡跟從耶穌,是因為希望能夠從耶穌身上得到平安,有喜樂,有祝福,有很多好處,甚至有些禮拜堂雖然沒有講出來,在洗禮人當中有很多父母,因為他們想洗了禮以後,他們的兒女去唸基督教學校的機會就大了,香港政府也同意加5分給他,很多都是為了加5分來洗禮的；有些人也覺得基督教很多老人院,進老人院也方便一點,洗了禮可以進老人院；基督教有墳場也便宜一點.我們今天為了許多好處,福利,一些得著而來,但耶穌在福音書裡面記載的很清楚,祂行神蹟,的確吸引了很多人來；講道也很有能力,也吸引了很多人來,但耶穌特別在路加福音就說:轉過身來對他們說,你假若不愛我,超越你的父母,兄弟,姐妹... (很多人聽到這裡就無所謂,因為這些都馬馬虎虎愛的),但祂繼續講下去:妻子,兒女,甚至你的生命,你假如不愛我超越這些,你不要來跟從我.所以有些人就說,你算甚麼?你有甚麼權力如此要求人呢?是啊!假如祂只是一位老師,一個普通的人,祂沒有理由去要求這些.但今天我們假如承認祂是上帝的兒子；假如我們承認祂就是道成肉身,到我們當中,充充滿滿有恩典,有真理的話,我們要選擇:你要跟從耶穌嗎?你就要背起十字架來跟從祂!你要跟從耶穌嗎?你要喫祂的肉,喝祂的血,將祂的生命放在你的裡面,你要活出耶穌的生命,而這個活出,不是抽象的,是有十字架的記號.

所以在今天的經文裡面,不斷從五餅二魚的神蹟,一直到第六章結束的時候,同樣的主題,有些人看到耶穌,耶穌走了,他們想找耶穌,耶穌說,我不要為餅而來的人,你們要來的話,你們要認識我是誰,祂在這裡沒有提到上帝的兒子,沒有提到基督,你要認識我就是從天上降下來生命的糧,要將你們從肉體的生存,為生活奔波的生命,只追求物質,人的名譽,權力這種位置的生命改變過來!你要裡面得著生命,不是肉身活著；人只看重活著只為了食物,但你要上帝的話──上帝的話就是道,道就是代表意義,代表價值.今天人的價值觀怎麼來?其實都是人告訴我們,從小的時候,父母就告訴你:做功課,考第一很重要!所以你從小就知道這個很重要!你乖很重要,你不聽話,你就得受罰；你聽話,就有獎品.我們的價值觀都是人的話給我們的,小的時候就是父母的話影響我們一生,因為父母這樣子講的.甚至一些很聰明的人,考試很成功的人,他都脫離不了父母的話,往往有些充滿詛咒的話,一生就活在父母親詛咒的話裡面；長大了是老師的話很重要；再長大是社會,傳媒的話,人家的話.我們很介意人家怎麼講我,他怎麼說我,我們的價值往往建築在人的話,因為人的話就代表人付予我們的價值.上帝的話不僅是這本書而已,是因為聖經是上帝的話,它告訴我們甚麼是有價值,有意義的,甚麼是真正的生命,所以耶穌說:人活著不是單靠食物,乃是靠上帝口裡面的話,因為只有上帝給我們的價值,我們才有意義活下去.或者,正面的說,我們活著,就是為了上帝付予我們的價值.人看不起我,人壓你,這人的權力,不同的位置,對不同的人有不同的價值壓力,但人怎麼能超脫呢?我們都受別人影響,除非我們看到上帝告訴我們說:你是有價值的,不但是我創造你,為你釘十字架,你無論在甚麼境況底下,你的生命是活在上帝的價值裡面,那才是最有價值的,你可以怎樣維持這種生命呢?當然我們來吃聖餐,但只是聖餐是不夠的,所以耶穌用的話很具體:「你要喫我的肉,喝我的血」.是很具體的,當然今天很多神學家,因為這幾句話,就做了很多神學的文章,有些傳統就說:所以我們今天吃聖餐,餅不是餅,是肉... 有人就不懂,怎麼吃來吃去都是"威化餅",沒有變成肉嘛!不像牛肉乾,也不是豬肉乾的味道,當然沒有人肉的味道,但很多神學家說:耶穌這樣說,就是這樣!那個餅就已經是肉,那個杯就是血.有些人就很怕說不要喝,我不想吃人的肉,喝人的血.但有另外一個傳統就是說:我們只是記念,那個沒有變,那個還是餅,那個還是杯,只不過是用來記念的!而有一些是中間派的:那個又是,又不是,吃起來是餅,但實在是肉──這些都是那些神學家所寫的文章之內容(所以我們才有那麼多活要幹,但那些神學家做的,並非最重要,也已經做了一千多年了).

今天其實在經文裡面對我們信徒來講,最重要的就是:你跟耶穌基督的生命有甚麼關係?你的生活裡面還是跟別人一樣,活著就是為了食物,我每天就是找食物... 廣東話有時候有一些很精彩的地方就是,我們一些年青的時候就常問一些很忙的成年人:--    你那麼忙幹甚麼呀?--    「搵食嘛!」(如果翻譯成普通話就失去那個神韻──找吃的嘛!)--    「搵食啊!搵食艱難啊!」(為找吃的東西啊!生活艱苦啊!)所以我以前年青的時候就很奇怪地問:--    那你「不搵食」,你會怎樣呀?(假如你不找食物,你會怎麼樣?)--    不找食物,我就死了嘛!--    那你不死,你又幹甚麼?--    「唔死就繼續搵食囉!」(不死就繼續找食物囉!)今天這就是我們生存的形態,不斷地去找食物,因為怕死嘛,當然死了就不用找了嘛!不死就繼續找!繼續忙!今天很多人失去了,去找生命更高的價值,耶穌告訴我們說:我們做基督徒,不是單單靠食物,是靠上帝口裡的話!所以代表的是意義,價值,你是否活著一個豐盛的生命?不是用物質來衡量的,不是用人的眼光來衡量的,不是用你的銜頭,職位來確定,你就這樣活在上帝面前有甚麼價值?你有沒有活出上帝創造你的價值?上帝的價值就在於使別的人也活出上帝的價值來.你有沒有令到別人能夠認識上帝,聽見福音,使他生命改變?所以今天在這裡所提到的:我是生命的糧,我就是上帝的道,要活人的當中,充充滿滿的有恩典,有真理,你要喫我的肉,喝我的血.而且不止聖餐的時候吃,也不止聖餐的時候喝,所以第三方面這個重點是提醒我們自己的生命──第三個重點就是剛才所說的:別人從那裡得著耶穌基督的生命?別人怎麼認出耶穌來?就是當祂擘開祂的餅,分給別人的時候,人家就認出耶穌基督來,但以馬忤斯的記載,耶穌就不見了!甚麼意思呀?因為耶穌基督的生命,是靠這些跟從祂的人延續在世上.

在我們生命的記號裡面,有沒有十字架的記號?但更具體的就是,我們有沒有繼續擘開分給別人的記號?我們每一天的生活裡面,都希望找到利益的,每一個都為了自己的利益,權力去奔波,去苦腦,去爭取,但今天基督徒的申命記號,不單止代表著十字架面對著罪惡的鬥爭而不害怕,面對死亡而不害怕,面對痛苦不害怕,因為上帝叫我們從死裡復活,耶穌已經復活了,我們為什麼怕死呀?耶穌在十字架那麼苦,就是經過這個苦難才得到榮耀,這就是十字架最重要的意義.

今天經文所說的「擘開」,「分給」就不止是自己的經歷,是你的生命怎麼樣不斷地擘開呢?怎樣不斷地分給別人呢?這也就是今天跟從基督... 就如同那個有錢的年青人來到耶穌的面前說:我怎麼能夠承受永生呢?耶穌跟他講了話之後,就很愛他,就告訴他:你要變賣所有的,分別窮人,然後你就得永生了!他就憂憂愁愁的走了,因為他財業很多,他當時已經很多東西了,他覺得自己缺乏一樣,就是永生,他希望N+1就很滿足了.耶穌說,你也很笨,你也就是差一樣,就是減N得到1啦!不是N+1,原來他要不斷的變,不是拍賣啊!一變,不斷的分,就跟我們的聖餐,五餅二魚的神蹟,以馬忤斯路上都是一致的,就是你的生命的價值在那裡?上帝的話就是生命的價值,你讓生命怎麼樣不斷地擘開,分給別人,而且起點是你就先感謝上帝給你的生命,所以這句話:你要喫我的肉,喝我的血──因為我們的生命要學習完全認同耶穌基督的生命,讓耶穌基督的生命在我們裡面,使我們得著耶穌基督這種生命的價值與意義,有血,有肉的,不是抽象的觀念,話語,乃是用行動表達出耶穌的生命是道成肉身,住在人的當中,不斷地將自己的生命,包括你的時間,才幹,精神,上帝肯定你的價值,上帝有賜給你生命的恩賜,你怎麼學習不斷的擘開,分給那些還不認識上帝,還沒經歷福音,還不懂得生命的價值是從上帝而來,不是從人的話語而來的人.只有將我們的生命不斷的擘開,分給別人的時候,人家認出耶穌基督在我們的生命裡面,也讓人家能夠遇見耶穌基督,他也得著耶穌基督的生命,我們的生命豐盛,他們的生命也豐盛,這才是今天我們活著不是單靠食物,乃是靠上帝的話,靠耶穌基督的生命,繼續在世上不怕死亡,直到永生.

請大家同心禱告:慈愛天父,感謝讚美你!因為你是我們生命的主,今天我們常常這樣說的時候,我們只想起我們肉身的生命,因為我們的生命是你創造的,但更重要的,你將我們跟別的動物分別出來,原來我們還在整個生命的發展當中,我們不單止有今天,還有明天,我們明天不單止有肉身的活著,靠著這些物質的食物,我們今天活著不單止靠著人的話語,雖然那些也有一些重要性,而且影響著我們非常之大,但你提醒我們:我們活著不是單靠食物,不是單靠別人的話,乃是靠你的話,因為你的話代表你的價值與意義；你的話要叫很多人,詛咒的話,惡毒的話要消除,也讓我們活在你的話當中,我們才明白真正的價值與意義,求你今天讓我們常常記得勝餐不單止是一個禮儀,也不是禮拜中的程序,乃是在每天的生活裡面──無論在以馬忤斯的路上,在曠野五餅二魚的時候,你總是告訴我們說:生命最重要的是怎麼擘開,分給別人,求你就讓我們渴慕你的話,渴慕你的生命,我們喫你的肉,喝你的血,讓你的生命充滿我們的生命,讓我們有血,有肉的,在這個世上,彰顯你生命的大能,讓我們願意被擘開,不斷的分享,以致別人也能夠得著你的愛,你的價值,我們感謝你,因為你常常與我們同在,奉主耶穌名求,阿們!



\newpage

\section{我們還跟從誰呢}
\label{sec:13}
\textbf{賴品超教授 - 我們還跟從誰呢?To Whom Shall We Go?}
\newline
\newline
\hyperref[sec:12]{< < < PREV SERMON < < <}
~
\hyperref[sec:toc]{[返目錄]}
~
\hyperref[sec:14]{> > > NEXT SERMON > > >}


約翰福音 6:69-71
\newline
\begin{longtable}{cl}
\hline
\hline
章節 & 經文 (和合本修訂版)\\
\hline
6:69 & \begin{tabularx}{0.7\textwidth}{X} 我們已經信了，又知道你是神的聖者。」 \end{tabularx} \\ \\ \relax
6:70 & \begin{tabularx}{0.7\textwidth}{X} 耶穌回答他們：「我不是揀選了你們十二個嗎？但你們中間有一個是魔鬼。」 \end{tabularx} \\ \\ \relax
6:71 & \begin{tabularx}{0.7\textwidth}{X} 耶穌這話是指著要出賣他的加略人西門的兒子猶大說的；他本是十二使徒裡的一個。 \end{tabularx} \\ \\
[1ex]
\hline
\hline
\end{longtable}
日期:2009年9月20日

題目:我們還跟從誰呢?To Whom Shall We Go?

證道:賴品超教授

經文:約翰福音6:69-71

各位弟兄姊妹,主內平安!

非常高興再次有機會與大家一同來分享上帝的話語,好不好大家再唸最後兩節的經文,因為剛才我們讀經的時候,唸到約翰福音第六章69節,但希望講道的時候也講一講70-71兩節:

6:70 耶穌說,我不是揀選了你們十二個門徒麼．但你們中間有一個是魔鬼.6:71 耶穌這話是指著加略人西門的兒子猶大說的．他本是十二個門徒裡的一個,後來要賣耶穌的.

今天我們所講的題目是:我們還跟從誰呢?就是在69節裡頭,所出來的一句話.我們還跟從誰呢?聽出來好像一個問題,在程序表上是加上了一個問號,但這可能不是一個問題,其實是說:不跟從你還能跟從誰呢?所以基本上不是一個問題,反而是甚麼呢?乃是一個很嚴肅的認信.意思就是說,彼得在這裡認信:如果不跟從耶穌,我們還有何人可跟從呢?為什麼?因為「主阿,你有永生之道,我們還歸從誰呢.我們已經信了,又知道你是　神的聖者.」(6:68-69)

信耶穌得永生,我相信很多基督徒都知道,且是很簡單的道理,甚至乎不是基督徒的也知道,對吧?香港一些教會外面的招牌也會寫上:信耶穌得永生.似乎這是一個很基本的知識,在彼得這個認信當中說:「主阿,你有永生之道,我們還歸從誰呢.」好像是一個很簡單,普通的認信,好像是大家所知道的事,又有甚麼特別之處呢?但是如果當大家看一看這段經文的上,下文的時候,我們可能會發覺這不是一個簡單的認信.彼得作這個認信的時候,正正是許多門徒都退去的時候,在六章66節裡面提到,當時有許多門徒退去的,這裡所指到的門徒是十二門徒之外的那一批門徒,雖然他們不是最接近主耶穌的那十二位,但是這一批門徒他們也曾經聽過耶穌的講道,甚至跟隨了耶穌一段時間,也甚至經歷過耶穌曾行過許許多多的神蹟,特別是第六章提到的「五餅二魚」──用五餅二魚餵飽了五千人,他們有這個所謂的屬靈經歷,他們有屬靈的知識,聽過耶穌講道,但是最後他們都離開了,不再跟隨耶穌,為什麼?

曾經有一段時間,有許多人跟隨耶穌,特別是行了五餅二魚的神蹟之後,對吧?因為根據聖經所記載,當時有許多人... 嘩!原來跟隨耶穌有餅可吃,這相當不錯,如果談及投資回報,跟隨一段路就有餅吃,那也不錯嘛!可以解決溫飽的問題.甚至當時有些人跟隨耶穌是因為他們看見主耶穌所行的五餅二魚的神蹟,令他們想到當年摩西帶領以色列人出埃及的時候,同樣在曠野,沒有東西可吃的時候,摩西向上帝禱告,當時上帝照顧那些以色列人,給他們有嗎哪吃,因此當耶穌行五餅二魚的神蹟的時候,會使人聯想到... 會不會這位耶穌都可以像當年摩西一樣使我們有嗎哪吃,有溫飽,不單止如此,並且能像摩西般帶領以色人離開埃及建立自己一個獨立的國家?所以有人會期望,耶穌會不會能像摩西般帶領我們猶太人脫離羅馬人的統治,讓他們成為一個獨立的國家(因為當時以色列是羅馬的殖民地),他們希望耶穌成為那位政治的彌賽亞,可以解放猶太人.但是當耶穌指出,我們應該為永生的食物而不是那些必壞的食物來努力的時候,那些只想吃餅的人,他知道他找錯地方,他們就離開了.接著耶穌就和這些猶太人說,猶太人的祖宗在曠野喫過嗎哪後,還是會死的,唯有這個從天而降的糧,上帝的話語才能給人有永生,人要喫祂的肉,喝祂的血,祂與那些喫祂的肉,喝祂的血的人成為一體,祂在他裡面,他們也在祂裡面,這些人才能有永生.同樣地,有間接使一些原本想跟隨祂的猶太人,對祂很有期望的猶太人都會退去.特別當時有很多的猶太人非常尊重自己的祖宗,覺得自己的祖宗是很厲害的,他們是蒙神祝福的一群人,曾經在曠野吃過嗎哪,曾經使神用大能的手脫離埃及人的統治,難道這位耶穌比摩西更偉大嗎?祂是誰呀?我們都知道祂父母是誰?就是約瑟,又有甚麼特別呢?祂如此的宣稱,難道不是一種褻瀆嗎?彼得宣稱祂比摩西這偉大的先知還更偉大,到底祂是誰呢?當許多人聽見耶穌這番話的時候,他們都覺得是一件很難接受的一件事,而這段經文也曾經提講過的,「這些話是耶穌在迦百農會堂裡教訓人說的.他的門徒中有好些人聽見了,就說,這話甚難,誰能聽呢.」(6:59-60)

作為一個正統的猶太人,怎麼會接受耶穌這番的說話呢?耶穌好像說到自己比摩西更偉大,耶穌說那些祖宗吃了嗎哪之後還是會死,難到祂能給我們永生不死嗎?祂是誰呢?祂好像在侮辱我們的祖宗,傳統.當然更加難解釋的就是,為何要「喫祂的肉,喝祂的血」呢?相當血腥,非常不文明的一件事,但最難解釋的是究竟為何喫祂的肉,喝祂的血,祂就會在我們裡面,我們也會在祂裡面...「我又因父活著,照樣,喫我肉的人,也要因我活著」?說祂自己是從天降下來的糧,如何解釋呢?

當然我見到那些唸神學的學生,在唸這段經文的時候,會聯想到聖餐這東西是很難去理解,對吧?唸完的在那裡點點頭,的確是很難去理解,每逢講到聖餐的時候,都會想起基督教裡面有許多不同的神學理論,都在說那個餅和酒是變化了它本有的性質,成為耶穌基督的身體和血,還是說耶穌基督的身體和血加入了這個餅和酒這種性質呢?還是這餅和酒有某種象徵的意義,令人想起耶穌基督的身體呢?又有一些神學家說,其實耶穌基督的身體是很麻煩的,都不曉得在說些甚麼?有時講的是教會,有時講的是耶穌的身體,但有時又講的是餅和酒,明明耶穌是活生生的站在那裡,然後就說「這是我的身體」,祂是指著這個,還是前面的餅和酒呢?神學家就為了這個問題爭論不休,他們都不曉得那一個才是真實的答案.

但是對於這班門徒來說,最難理解的不是這些(當時候都還沒想過這種複雜的神學問題),他想到的就是說,究竟耶穌在說些甚麼?並且更重要的是他們發覺,耶穌所說的話,是很難使猶太人接受.就是說,門徒聽完耶穌所講的,他們都覺得是在「趕客」,如果有猶太人想跟隨你的,但聽了你這麼說之後,他們都會退去,甚至門徒們都覺得很難接受,就是連他們自己都難以接受,可能他們覺得耶穌不應該在迦百農的會堂,一個公開的場合裡面講這番話,在一個毫無必要,毫無原因的情況底下,失去這一大批的支持者.情況有如總統選舉的時候,理論上在總統選舉上應該講一些能吸引選民投他一票的話,好使他能得到最多的選票,以成為那一位的勝利者,總統,但偏偏在這個時候,耶穌剛行完一個神蹟,是一個最好的黃金時間來吸引人歸向祂的時候,祂就講了一番話,一番使人後退的話,祂告訴那些想得餅的人說,你不來跟隨我,因為我要給你的東西不是這些,而是永生的糧食,對著那些跟隨祂的猶太人說,最後要給你們的,不是你們所想像的那種,我不是要做政治的彌賽亞,以脫離羅馬人的統治,甚至講一些令猶太人覺得祂是侮辱祖宗,傳統的話語,所以如果你從總統選舉的角度來看的話,耶穌所說的話,基本上是自毀長城,祂因失言,使很多人失望或是憤恨,接著很正常地在總統選舉期間,如果兩位競選者在較力,其中一位如果民望很高,卻突然下滑,看似必輸無疑的時候,他原本的助選團會怎麼樣呢?挑選!不想再支持他,甚至轉至另一個陣營,希望另一個支持者最後成為總統的時候,自己也能謀求到一官半職.

這班的門徒因為聽到耶穌說了這番話的時候,他們覺得厭棄,在這段經文說他們厭棄,甚至跌倒了:「耶穌心裡知道門徒為這話議論,就對他們說,這話是叫你們厭棄麼.〔厭棄原文作跌倒〕」(6:61)正如主耶穌說的,他們中間有不信的(6:64),就是他們沒辦法接受主耶穌在迦百農會堂裡所說的那一段話,他們感受到失望或憤恨.但是這個時候,主耶穌並沒有用任何所謂的甜言蜜語來哄這群門徒說,剛才說的不算數,一時失言... 有時候有些政客也會這樣子說,其實我的意思不是這樣子... 目的是想挽回群眾的心.但相反地,主耶穌很直接的說,你們中間是有不信的,因這些話就跌倒,離棄我了嗎?緊接著又說:「倘或你們看見人子升到他原來所在之處... 」這裡可能就說,當你看到主耶穌復活升天的時候,就是說,主耶穌原本是從天上而來的那一位上帝的兒子,這不過是更加強調,強化祂在迦百農會堂裡所說的:祂不單單是一個人,祂是天上而來的樣式,可能比以色列人的祖宗所喫的嗎哪... 理論上很像,但不知道是從天上降下,或是從地而來的,沒有人知道.但是主耶穌只是強化在迦百農所說的話,祂並沒有說許多溫婉的話,以挽留那些的門徒,反而更加強祂的信息,使那群門徒最後離開了祂,所以在經文說:「從此他門徒中多有退去的」.

主耶穌當時所面對的是一個眾叛親離的局面,然後祂才問那十二個門徒:你們是否也要離開呢?這是一個非常強烈的語氣:你們是否也要像那些人一樣離開呀?是否連你們也要離開呀?(6:67)這是一個很嚴肅的問題.而彼得就在這個時候作出他的認信:「主阿,你有永生之道,我們還歸從誰呢.我們已經信了,又知道你是　神的聖者.」(6:68-69)假若我們從這個脈絡去看的話,彼得的認信是一個不簡單的認信,當許多許多人離開耶穌之後,他仍然堅持要跟隨耶穌的這種認信.根據我們中國古代的一個說法,這就是「舉世非之而不惑」──就是當人人都認為你錯的時候,人人皆跳船的時候,你依然堅持下去要跟隨耶穌.當然對耶穌來說,這句話你可以說是「雪中送炭」,對吧?這麼多人都離棄了你,不再跟隨你的時候,居然有一個人講了一些悅耳的話:你是永生上帝的兒子,你有永生之道,我們還歸從誰呢?

在最需要的時候,作出這種的認信,展示出彼得好像「眾人皆醉,唯我獨醒」,「力排眾議」... 代表著他那種堅定的信心,他不管其他人是如何地離棄耶穌基督,他只需在基督身上找到永生之道,他就跟隨這位耶穌.他記得他聽過甚麼道理──主耶穌給他有永生；他也經歷過神很多的恩惠──我們已經信了.他很清楚自己的信仰,很清楚──知道你是　神的聖者.耶穌本該覺得很安慰,彼得真是「汝子可教」,總算找著「知音人」,當所有人都離開我的時候,還有一個人能夠堅持信仰,要跟隨我,也不枉我教導你這麼長的時間.可能大家也會想起馬太福音裡頭的一段聖經,當彼得回答得很精彩的時候,說,你是... 你是基督... 人說你是先知裡的一位...你是永生上帝的兒子的時候,主耶穌就說,你是磐石上,我要將天國的鑰匙給你等等... (所以現在大家看到聖彼得的像,都拿著一串鑰匙,就是天國的鑰匙),但在這段經文裡面卻不是如此,主耶穌沒有立刻跟他說,你是磐石上,我要將天國的鑰匙給你,凡你在地上所捆綁的,在天上也要捆綁．凡你在地上所釋放的,在天上也要釋放.(太16:18-19)主耶穌沒有說這番的話,反而跟他說甚麼呢?祂說:剩下這十二個門徒,不過在這十二個裡面,還有一個是魔鬼會背叛祂(約6:70-71)換言之,即是堅持到最後的這批人,這十二個裡面,仍然有最後的失敗.在這段經文所突出的就是主耶穌知道萬人的心,哪些是信的,哪些是不信的,哪些會出賣祂,所以祂不在意到底要用甚麼話來挽留那些人,那些會出賣祂的人,因為祂知道,只有蒙父恩賜的人,才能到祂那裡去(約6:64-65).

當然我們唸神學的,毋須很快推論說,這是不是在講「預定論」呢?上帝一早預定的人就會揀選他得救,沒有上帝預定的,或是上帝預定下地獄的就注定不得救了!毋須要在這裡作出這種的推論.因為經文裡頭所說的,只不過是說:如果有人能夠信耶穌,跟隨耶穌到底,這是出於上帝的恩賜:「若不是蒙我父的恩賜,沒有人能到我這裡來.」(約6: 65)原來問題並不是那麼簡單說:我們還能跟從誰?而是另外一個問題:究竟我們為了甚麼跟隨耶穌?有些人為了得餅而跟隨耶穌,很快就會走掉；有些人為了得到政治彌賽亞,他們都會退去.但是不單單是「誰」的問題,或者「為何」的問題,還有一個就是「如何」(How)的問題:究竟跟耶穌到那裡?有些人問對了問題,跟從耶穌對了,但是最後是出賣耶穌的那一個,他沒有跟從耶穌到底.

我不曉得大家看到這經文會想到甚麼?我想到對於我們當代中國來說,也會有類似的一個情況... (我經常就要接觸一些大陸的學者,今天下午我就要坐飛機去武漢開會,天氣這麼熱... 四大火爐... 還要今天去,真的是選對了日子呀!)我們中國最近會有許多的所謂「基督教熱」,多了許多知識份子對基督教很有興趣,而去研究基督教.有些人非常仰慕基督教文化,覺得這是一個先進的文明,有機會帶給中國一個更加先進的文明,以建立民主,人權等等.甚至認為基督教有可能對人的心理健康有幫助,對經濟增長有幫助,對建立和諧社會有幫助,甚至對建立可持續發展的社會有幫助... 但是有多少是想在耶穌基督身上得到永生?我們有興趣的是在經濟上面「可持續發展」,我們希望經濟不斷增長,永續增長；我們希望我們這個和諧社會(如果它曾經存在的話)會繼續存在.但是有多少是想得到在基督耶穌裡面所找到的「永生」呢?

我們中國人很現實,我記得我年輕勸人信耶穌的時候,都很容易聽到這句話:信耶穌,耶穌拿飯給你吃呀?我們覺得如果信耶穌,耶穌拿飯給你吃的話,那你信耶穌是一件相當應該的事,與那些跟從耶穌,希望得到餅的那些人一樣沒分別.中國人真的很現實,如果你拿飯給我吃,你給我一些實際的利益,我跟從你沒問題,但是有多少是想在基督耶穌身上得到「永生」呢?是為永生跟隨耶穌,為永生之道跟隨耶穌呢?

我們香港的父母,很想我們的子女進入基督教,天主教的名校,但是有多少是有興趣進入上帝的國度,永生的國度呢?我們中國亦可能像一些猶太人一樣,很注重祖宗的傳統,有些人覺得信耶穌不合乎中國的傳統,基督教有好的東西,我們的祖宗怎麼會不知道呢?怎麼孔夫子都沒有提起呢?我們覺得這是不可以接受的,但問題是,我們是否為了迎合這些人的需要,口味,我們就不再講永生的信息呢?是否為了避免得失這些人,我們不敢對我們傳統的文化作出任何的批判,反思呢?原來我們同樣生活在一樣的年代,當你說一些很受歡迎的東西時,可能很多人跟隨你；或許你向人提供一些物質的需要的時候,很多人會跟隨你.但是當你告訴他們,你所要給的不是那些可以朽壞的食物,不是經濟的增長,不是和諧的社會,不是心理的健康,不是可持續的發展,而是真真正正永遠生命的時候,有多少人願意來接受?有多少人願意跟隨耶穌呢?

講到這裡,我有時候也覺得人的「信」是一個奇蹟,如果不是上帝的恩賜,沒有人可以信.就是因為信耶穌,沒有飯給你吃,信耶穌不一定令你很健康,信耶穌不一定令你很快樂,信耶穌不會令你富有,但仍然有許多人相信.信耶穌當中,包括許多難明的教義:怎麼會喫祂的肉,喝祂的血就得到永生呢?其實是真的很難去解釋,明白,但是竟然有那麼多人相信.但在我們身邊我們同樣看見一個情況,也有許多門徒是退去了.可能在我們身邊,人生的經歷裡頭,我們也見過許多不同的人──我們見過很多表面上信耶穌,實際上不是的人；我們也見過一些曾經很熱心,甚至做過教會領袖的人,他們最後都退去了.究竟我們對主耶穌的跟從是怎麼樣的呢?我們跟從祂是希望得到可以朽壞的食物呢?還是為了永生?我們跟隨主耶穌是跟隨到一半之後退去,還是跟隨到底?還是我們成為出賣主耶穌的那一個?我覺得這段經文給我們一個很深刻的反省,就是原來即使在十二個門徒裡面,他們很有屬靈的知識,經歷,但是最後可以有人出賣主耶穌.

原來我們要跟隨主,希望得到永遠的生命,但這不是憑著我們自己屬靈的知識,經歷,就說我們可以跟隨到底,最後我們仍然只能夠求主幫助我們,讓我們對準主耶穌基督,要祂裡頭得到永遠的生命,而不是那些可朽壞的食物,更加求主幫助我們,我們可以忠心地跟隨祂,並且為主的永生之道,來做見證.

在我們這個世代裡作基督徒,表面上作基督徒是件很容易的事,而我的學生也看見,以前作基督徒左顧右盼的,像吃飯禱告要在兩秒內搞定,或許你在桌子底下擦擦汗就沒事,就做了禱告.現在作基督徒就輕鬆自在很多,你要禱告人家也不覺得有甚麼問題,對吧?甚至有一些藝人也很公開說,我是信耶穌的,也甚至有一些也不曉得是否相信就自說是相信的,特別是被人揭發醜聞之後,就告訴人:我悔改了,是新造的人.信耶穌在我們的社會裡面,已經不再成為一個羞恥的事,要公開說自己信耶穌是一件很簡單的事,但是如果我們要真正跟隨耶穌到底,卻未必是一件很簡單的事,因為我們看見身邊有很多人退去,我們容易看見信仰裡面有很多我們未能明白,了解的事情,更加要緊的就是,我們很容易被一些可朽壞的東西吸引了,忘記主耶穌最後要給我們的是「永生之道」──在其他事物身上我們沒辦法找著的「永生之道」.

願神再次祝福我們,讓我們靠著主站立得穩,緊緊的跟隨祂,在祂裡面找著「永生之道」,不是靠著我們的屬靈知識和經歷,乃是靠著主的恩典,我們才可以跟隨到底,為主作美好的見證!



\newpage

\section{律法遺傳與良心}
\label{sec:14}
\textbf{戴國璋牧師 - 律法,遺傳與良心The Law, the Tradition and the Conscience}
\newline
\newline
\hyperref[sec:13]{< < < PREV SERMON < < <}
~
\hyperref[sec:toc]{[返目錄]}
~
\hyperref[sec:15]{> > > NEXT SERMON > > >}


馬可福音 12:28-34
\newline
\begin{longtable}{cl}
\hline
\hline
章節 & 經文 (和合本修訂版)\\
\hline
12:28 & \begin{tabularx}{0.7\textwidth}{X} 有一個文士來，聽見他們的辯論，知道耶穌回答得好，就問他說：「誡命中哪一條是第一呢？」 \end{tabularx} \\ \\ \relax
12:29 & \begin{tabularx}{0.7\textwidth}{X} 耶穌回答：「第一是：『以色列啊，你要聽，主—我們的神是獨一的主。 \end{tabularx} \\ \\ \relax
12:30 & \begin{tabularx}{0.7\textwidth}{X} 你要盡心、盡性、盡意、盡力愛主—你的神。』 \end{tabularx} \\ \\ \relax
12:31 & \begin{tabularx}{0.7\textwidth}{X} 第二是：『要愛鄰如己。』再沒有比這兩條誡命更大的了。」 \end{tabularx} \\ \\ \relax
12:32 & \begin{tabularx}{0.7\textwidth}{X} 那文士對耶穌說：「好，老師，你說得對，神是一位，除了他以外，再沒有別的了； \end{tabularx} \\ \\ \relax
12:33 & \begin{tabularx}{0.7\textwidth}{X} 並且盡心、盡智、盡力愛他，又愛鄰如己，要比一切燔祭和祭祀好得多。」 \end{tabularx} \\ \\ \relax
12:34 & \begin{tabularx}{0.7\textwidth}{X} 耶穌見他回答得有智慧，就對他說：「你離神的國不遠了。」從此以後，沒有人敢再問他甚麼。 \end{tabularx} \\ \\
[1ex]
\hline
\hline
\end{longtable}
日期:2009年9月27日

題目:律法,遺傳與良心The Law, the Tradition and the Conscience

證道:戴國璋牧師

經文:馬可福音12:28-34

很高興能夠在今天這一個機會,與大家一起在聖經的說話當中彼此訓勉,我們按著這一個主題,我們希望能夠從聖經一個我們熟悉的故事來思考,這個故事是記載馬可福音第12章28至34節,若果大家手上有聖經,可翻開聖經看看這個故事.

故事內容是這樣,他說有一個文士來,聽見他們辯論,曉得耶穌回答的好,就問他說,誡命中哪是第一要緊的呢?耶穌回答說,第一要緊的就是說,以色到啊,你要聽,主我們上帝是獨一的主,你要盡心,盡性,盡意,盡力愛主你的上帝,其次就是說要愛人如已,再沒有比這一,兩條誡命更大的了,那文士對耶穌說,夫子說,上帝是一位實在不錯,除了衪以外,再沒有別的上帝,並且盡心,盡意,盡力愛衪並愛人如已,就比一切的燔祭和各樣的祭祀好的多,耶穌見他回答的有智慧,就對他說,你離上帝的國不遠了,從此以後,沒有人敢再問他什麼了.

這個故事是我們很熟悉的福音故事,在幾本的福音書上都有提及到,雖然記述的內容有少許不同,其實我們在這個故事裡面,從耶穌與文士的一問一答,我們可以見到,神,人關係的兩種看法,怎樣從猶太傳統的看法,轉向基督教傳統的看法,其實我今天想以這個方式跟大家思想,「律法」,「傳統」與「良心」這一個課題,我們知道這一個情景,其實是一個為難的情景,當時耶穌其實被很多人圍攻,被法利賽人及文士圍攻,嘗試從不同方面去為難耶穌,希望讓耶穌懼怕,或者尋找耶穌話語的把柄,其中一個文士就以這方面去為難耶穌,對於這個文士來說,他心裡面可能沾沾自喜,心想:「這問題還不難到你?,你看!律法是很艱深,誡命是很多,多如牛毛,現在我要你以一句說話,來把誡命當中第一要緊的說出來,這樣還不難到你?」誰知我們發現到,耶穌很簡單便作出了回覆,而且按照馬可福音這個版本,令到為難他的文士,心裡也誠服,這顯示出我們主耶穌基督的智慧,同時顯示出,從主耶穌說話裡面,他點出了神,人關係的精要所在,怎樣從猶太傳統,轉回我們從耶穌基督身上,看到的基督教傳統.

其實耶穌基督這個回答,即第30節的回答,他說:「以色列啊!第一要緊的是,你要聽,主我們的上帝是獨一的主,你要盡心,盡意,盡性,盡力愛主你的上帝.」這一個回答,在當時的猶太人來說,這個答案已經存在,其實在耶穌基督的時代,那些研究聖經的教法師們,他們基本上已經成熟到一個程度,他們對誡命的看法已經成熟到一個程度,認為在眾多的誡命裡面,即出現在申命記,第6章第4至5節裡面的誡命,是最重要的誡命,他們已經有這樣的成熟程度,所以當耶穌一回答的時候,他們就覺得耶穌真的掌握到這個誡命的精義.

我地知道申命記第6章第4至5節這一句說話,在後來猶太人的傳統裡面,他們給了耶穌一個特別的稱呼,這個稱呼叫做「士馬」(SHEMA),「士馬」是什麼意思呢?「士馬」其實就是聽的意思,因為在這段經文一開始時,那段經文是說,以色列啊!你要聽,在希伯來語是你要聽,以色列啊!是你要「聽」,由「聽」開始的.「聽」就是「SHEMA」,其實就是將「聽」(SHEMA)這個動詞,轉化成名詞的形式,所以叫作「士馬」.在猶太人的傳統中,非常重視「士馬」,因為在他們傳統裡,不但講述在他們傳純裡,將誡命的重點刻劃了出來,甚至這些經文,是用一個經文盒被釘在他們房子的門框上,甚至一些虔誠的猶太人,會將它放在經文盒裡,然後戴在頭上,提醒自己時刻要記著這經文,這個就是「士馬」的傳統.

耶穌基督稱這個「士馬」為什麼呢?耶穌基督稱他為誡命中的第一,是什麼意思?這個意思其實解作是誡命中的誡命,猶太人很希望用???這樣的形容詞,來形容他們最重要亦最極致的表達.例如我們知道,聖經裡有一本書叫「雅歌」書,但其實原來的意思是「歌中之歌」,Song of songs,用這種方式表達,即是最美麗的詩歌.同樣,當主耶穌說這是誡命中的第一時,其實是用了這種重要的方式表達,這個是誡命中的最重要,除了我們剛才說有「歌中之歌」這種形式表達之外,聖經也有例如稱耶穌為「萬王之王」,「萬主之主」,其實都是指衪是最大的王,最大的主的意思,所以在和合本裡,便乾脆將它翻譯成「第一要緊」,這是「第一要緊」!我覺得和合本的翻譯,其實是取得原本的意思,於是我們就明白,當耶穌說誡命中的第一時,他的意思不是次序上的第一,第二,第三,第四...那麼簡單,其實是重要性上的第一,即是最大的誡命,或是最重要的誡命,這個在第31節裡,耶穌的說話便指示出來了,如是者我們便知道,盡心,盡性,盡力愛主你的上帝,這是誡命中的第一,盡心,盡性,盡力愛上帝,這是第一要緊的誡命,第一要緊的誡命,就是要盡心,盡性,愛主你的上帝.

大家可能很奇怪,我一句說話,為何要重複三次那麼多?其實如果大家有留意,我這三句說話的表達,有少許不同,這個少許不同,正正可以幫助我們,能夠看見有關這個誡命中的一個特點,所以接下來,我想跟大家按著這三句說話,去回看聖經提醒我們,與神,人關係的一上事情.第一,就是說,愛上帝就是我們生命中的第一,剛才我們已經說過,這個情景,其實是圍攻耶穌基督,質疑耶穌基督的情景,對於這個文士來說,他會覺得誡命很多條,多如牛毛,你怎麼可以在那麼多,又多如牛毛中的誡命中,你可以一下子掌握到最重要?所以對於他們來說,他們覺得這是一個不容易回答的問題.當他問:「什麼是誡命中的第一?」,即所有裡面的第一.

猶太人是一個守誡命的民族,守誡命對猶太人來說是很重要的一件事,虔守誡命是古代猶太人一生努力的目標,因為我們發現,在他們與上帝的相遇裡,是以聆聽上帝的誡命為開始,這一個成為了猶太人屬靈傳統起源的根據,所以當他們傳釋神和人的關係時,他們是以守上帝的誡命為開始,我們會在第二點加以詳細去講解,這個其實沒有問題.猶太人很重視守誡命,他們覺得這是虔誠猶太人應該做的,但到了耶穌基督的時代,我們發現他們將聖經一些誡命,作一些很仔細的分晰,將所有聖經,即摩西五經裡,摩西所給我們的律法傳統,他將之分作兩個大類別,這兩個大類別就是,有應作的,有不可作的,一共有613條,那不是小數目,是613條.

據一些經學家說,這613條誡命中,248條是應作的,即是你必定要履行的,他們將這248條比喻作,一個人的身體內外大小骨頭的總數.是否真的248條,我想真的很難說,要請教一些醫學上,特別有見地的弟兄姊妹才知道,但有人是這樣看的,他說我們人的身體上,有248條大大小小的骨頭,就暗合上帝給我們履行的誡命,我應作的,一共是248條.不應作的,有365條,即是說有365條是禁令來的,你不可以這樣,你不可以那樣...365條便暗合週天的數目,因為一年有365天.而且猶太人有一個很簡單的方法,幫助我們記著一共有613條誡命去履行,這613條誡命裡,248條,即我們人體的骨頭是應作的,365條,即一年的數目,一共是365條是不應作的,猶太人在古代是很有智慧的民族,他們懂得很多這樣的聯想,來幫助他們去熟悉有關他們的傳統.但某方面來說,或是岔開一點來說,我覺得這裡把248條比喻作人身體的骨頭,365條比作週天的數目,這個不會是偶然的巧合,或是一些聯想,來幫助我們去記憶一些東西,因我們知道在古代的世界,在很多民族裡面,包括猶太民族,希臘民族,甚至包括我們中國的民族,也不多不小有一些天人感應且相似的看法,就是人是一個小宇宙,天是一個大宇宙,天與人彼此之間,有一種感應在裡面,所以365條是賦予了週天的運行,其實是將他聯想起一個大宇宙,,而248條視作人骨頭的總數,其實就是將人看為一個小宇宙,天人感應,所以人便要去追尋上帝的律法,我相信是有這樣的內容,但這樣的發揮,或者容讓一些在神學或文學上,有研究的一些弟兄姊妹作探討.但無論如何,那時的猶太人已將整個律例誡命,細分作613條,但不是那麼簡單,根據猶太人的他以穆,即他們傳統的經卷,百課全書,他們把這613條再細分,細分後再細分,再細分成對神,對五經,對聖殿和祭司,對祭祀,對許願,對潔淨禮儀,對捐獻,安息年,食用的鳥獸,節期,社區生活的條例,偶像,戰爭,社會,家庭,執行刑罰的條例,奴婢,還有其他...一共18的條項,分了18的條項裡,又會教我們在這18的條項中,我們應該怎樣做,怎麼履行,才能把這些條項一一履行,所以猶太人這個百課全書的傳統,其實有很多很多這些誡命的仔細分析,其實我們知道,這些都已經離開了誡命的基本精神,已經變成了遺傳,什麼是遺傳?遺傳其實就是曆代以來的經學家,他如何去研究這些聖經,然後將這些聖經不停仔細分析,將他應用在每一個場合中,仔細地教導人如何做,所以你會察覺到,猶太人的生活指引很仔細,從某方面來說真的挺方便,因為有一本指南告訴我們怎樣做,但從另一方面來說其實非常綑綁,因為連你最細微的行動,也要告訴你要如何做,才能算是履行誡命,才能討主的喜悅,就是因為這原故,猶太人他們耗守窮經,分析很仔細,因為他們怕有遺漏,在他們傳統裡,他們以守誡命作為他們與神建立關係的起點,剛才我說過,是沒有問題的,但問題是以這作起點時,你心裡時常也會有一種懼怕,而這一種懼怕是,我守了很多很多誡命,但萬一我守漏了一條會如何?正如雅各會告訴我們,我們以為自己能守全律法,只是在一條上面我們有虧欠,我們便犯了信條,猶太人有這一種擔心,因此他們分析得非常仔細,人生的每一面,都被他們的遺傳及分析來覆蓋,但這種仔細是出於他們的敬虔,因為他們怕遺漏,遺漏便不蒙上主的喜悅.但這個正是他們的毛病所在,他們的毛病在哪?他們的毛病是犯了捨本從末,拆眉不知撮要.他們能夠作很仔細的分析,在很多的末節上,亦掌握得非常仔細深入,但最重要的是犯了這錯誤,遺漏了下來.

所以耶穌基督的回答,除了以他的智慧來回答文士的試問,令他們無話可說之外,耶穌其實還有一個心意,這個心意其實就是,藉藉這個回答,來提醒他們的傳統,由其是,由誡命以致到遺傳的傳統來看與上帝的關係,是行不通的,不單止人失去了自由,而且會把上帝所盼發的誡命精義,精要遺漏了,這個正是耶穌基督所說,法利賽人的毛病,他指那些法利賽人是如何的?他告訴他們,駱駝你們倒吞下去,蠓蟲你們就濾出來.你看看這個遺傳的網多麼仔細?仔細到一個地步...什麼叫蠓蟲?蠓蟲是細到幾乎看不到的蟲.以前古代人不知道有細菌或微生物這些更細小的東西,蠓蟲他們已覺得很細小,這個過濾網細小到一個地步,就連蠓蟲你們也可以濾出來,但是最重要的駱駝,你們反而遺漏了,吞了下去,這個就是本末倒置,捨本從末,拆眉不知撮要的毛病.因為我們看到,當猶太人的傳統,把613條分為18項的時候,不錯!他們是以神為首,剛才我們說過神是他們的重視首要,但大家又可以看看,他們對有關神的誡命,作了如何的細分.關於神的,他們細分為:相信神的同在,承認神的獨一,要愛神,要敬畏神,要事奉神,要戀慕神,單指神的名起誓,要效法上帝,尊神的名為聖等等...

其實在這上細分項裡,我們發現到只有一至兩項,是真正碰觸到耶穌基督首要著緊的精神,就是愛上帝,真的只有這裡碰觸到.所以耶穌是借這句說話來提醒他們,這一種細分律法的遺傳工夫,是不能把上帝表達的精義掌握到,除非他們明白,那精義是神與人的關係建立,就是愛上帝.

親愛的弟兄姊妹,或許我們也很熟悉猶太人這一種毛病,因為當我們上教會久了,特別是多讀福音書的時候,知道耶穌與法利賽人之間的毛病,耶穌對法利賽人評論的時候,我們或多或少也聽過.不過你們有否發覺到,我們作基督徒的時候,我們亦會犯下與這些猶太人相似的毛病,例如我們同樣可以,以讀經,禱告,事奉,見證,傳揚福音,社會關懷,職場服侍等等...的明目來代替一些重要事情,就是愛上帝.這樣說不是指我們屬靈的操練不重要,我相信你一回來教會的時候,自小回來教會,便告訴你讀經是如何的重要,對吧?連串教你如何讀經,靈修,禱告如何重要,事奉如何重要,這些不必多提及,甚至今天,很多時候在教會也宣揚,鼓勵大家在這上屬靈的操練上不要懈怠,我從來沒有說過這些操練不重要,我不是這樣說,我只是說,我們很容易會用這些屬靈操練,來代替那最簡單,最重要的,就是愛上帝.為何這樣說?因為我們可以愛讀聖經但不是愛上帝,真的,有一些人很愛讀聖經,因為他們覺得當他們研究聖經時,他們能看到一些獨特的亮光,當他們把這些獨特的亮光跟人們分享時,人們會覺得他們很聰明,原來這麼難的東西你也看到,於是便沾沾自喜,他們喜悅的,其實是他們自己解經的能力.有些人很愛研究聖經,聖經裡每一個細節他都研究得到,但這些研究是出於他對上帝的愛慕,還是只不過是出於他研經的興趣?這真的是值得疑問.同樣,禱告,事奉,見證都可以是這樣,我們可以把這一些屬靈的操練,代替了最簡單的,就是要愛上帝.

所以當耶穌基督回答這班法利賽人的時候,其實是把我們基督教裡很重要的重點,標誌了出來,這個標誌告訴我們,第一,是要提醒我們神與人的關係,才是最核心的關係,這個核心的關係是以愛去維繫,我們很熟悉保羅,在哥林多前書第13章第一節處說到,若沒有愛,我就成了鳴的鑼,響的鈸一般,若果沒有愛,我們剛才所說,我們很珍惜的屬靈操練,讀經,禱告,事奉,見證,傳揚福音,社會關懷等等,其實都變成沒有意思,我覺得在這個事情上面,可藉著耶穌基督有一些提醒.

第二方面,就是第二句說話,愛上帝是第一要緊的誡命.在第一方面,我們著重了什麼叫第一,現在我要說的是誡命,猶太人是一個守誡命的民族,遵守誡命曾經一度融入他們的血液中,這種遵守誡命,甚至遵守他們遺傳中最仔細的部份,其實可追溯到他們民族誕生的契機,即西乃山的屬靈經驗,因為他們在西乃山裡跟上帝相遇的時候,上帝給了他們十條的誡命,上帝要求持守這十條誡命,所以某方面來說,猶太人這種守誡命的傳統,是跟他們的祖先...特別是在西乃的這種傳統,是很有關係.

有這處我想點出,律法的精義及精要,律法是什麼?我們知道律法可分為三個大範疇,就是誡命,律例及典章,這個我們在出埃及記,申命記中會讀到,誡命,律例及典章其實是三個範疇,什麼是誡命呢?誡命通常是指道德的律例,這些道德律例是不受環境影響,不受情況影響的,它是普世通行,是恆久有效.例如在十誡裡,大部份要人遵守的,例如不可奸淫,不可偷盜,不可作假見證陷害人,這些有關貪婪,誠實這方面的事情,這些誡命是道德律例.律例其實是指什麼?律例其實是指Case Law,即與猶太人生活傳統有關的Case Law,因為在猶太人的社會變化裡,每一個時代有不同的社會環境,在不同的社會環境中,如何應用這些誡命,就成為這些律例,例如說,如果有人走失了一隻牛,你看到了該如何處理?因當時猶太人是一個農業的社會,有一個牛主,他不小心讓他的牛撞死了一個人,那撞死人以後,又要如何作出賠償呢?建一間屋,有天台,而天台要有圍籬,如果天台沒有圍籬的時候,有人有天台掉下來死了,那個業主便需要賠償了.這些都是一些Case Law,律例.而什麼是典章呢?典章就是指一些宗教上禮儀的條例,包括了各樣獻祭的條例,獻牛,羊等,這些各樣熟悉的條例,或者各樣潔淨的條例,例如耶穌跟法利賽人談論的洗手問題,這些不是誡命來的,這些不會牽涉律例,而是牽涉典章,即所謂宗教禮儀上的問題.猶太人把這些律例分成三個大分別,用這三個大分別去看上帝的事,我覺得沒有問題,也是可以的,雖然我們用誡命的方式去傳達有關愛的事情,在我們看來,我們會覺得好像很奇怪.當耶穌說愛上帝是第一要緊的誡命,感覺是怪怪的...為何呢?因為「愛」,看上去是一種自發的東西,「愛」不是一種規限,我不可以用一個規定,去規定你愛我,或規定妻子愛我,或者愛我們的妻子,「愛」是應該出於我們內心自發的,所以用誡命來詮釋上帝的愛,所以有時我們覺得怪怪的,但從另一角度來說,若果我們將愛看為要求的時候,我們亦會覺得很合理,為什麼呢?因為當我們愛一個人的時候,我們對一個人總會有一些要求,正如耶穌基督告訴我們,當我們愛上帝的時候,我們便會行衪所喜悅的事,做一些事是衪喜歡及高興的,所以我們會發現「愛」往往有給予,或者會有受落的,只有把愛給予,對方又受落,我們才能長久保持新鮮感,所以在基督教傳統裡,我們是以這一種「愛」,來解釋神與人的關係,而用誡命的傳統,即是說我們明白對方需要什麼東西,我們希望滿足他,這種方式是沒有問題的.

其實這樣正正提醒我們,在基督教傳統裡,來講述神與人關係,以愛去開展的時候,我們應該記得,我們要時常保持這種我的新鮮,來令到對方感受愛的施予.

在較早時,報章上讀到一段花邊新聞,暑假前有一天我駕車送兒女上學,通常我會聽聽電台消息,發覺電台節目主持人正在談論這段新聞,這段花邊新聞叫婚姻保鮮法,原來在英國有一個禮服服務網站,他向三千名已婚的男女進行了一個調查,就詢問他們有關婚姻幸福的秘訣,結果發現到,那上婚姻長久的夫妻,多數在婚前已經相戀三年半,而且多半有兩個共用興趣,丈夫又比他的妻子年長兩年多,而第一個小孩最好在婚後兩年零四個月出生,不知大家有否讀過這段新聞.如果有,那會否在擇偶時考慮這上條件?不單止這樣,他還有一些很有趣的,他說要令夫妻婚姻維持長久的,還有很多工夫要做,第一,就算工作如何忙碌,大家都要時刻彼此聯絡,每天最少要通電話三次,發短信或是電郵,每天要擁抱三次,親吻四次,最少要說一次我愛你,並且每週有三次性愛,每個月吃兩浪漫的燭光晚餐,或是一起外出看電影,每年到海外旅行兩次,我看了這些要求,感到自己實在沒有一件事合格,不過也好,讓我知道我不合格,可以做得更好,讓我妻子知道我對她的關心,愛護,這一段只是花邊新聞,大家聽了可能只會一笑置之,但他的確提醒了我們一件事,就是人與人之間,夫妻之間的感情時時刻刻要保鮮,真的!所以不要說,當天我追求她時已經表達了,不需要再表達.但其實不然,真的要多些表達,我發覺真的有效,因為你多些向你妻子表達關心及愛慕時,你會得到一個很好的回報.

其實這件究提醒我們,我們其實跟神的關係都是一樣的,我們應該常常想對上帝的事情,衪想我如何為衪呢?衪對我有什麼要求呢?誡命在某方面來說就是要求.我們不要將他了解成一個要求,我們不要這樣想,而是一個愛我們的主,向我們發出的要求,其實是因為衪愛我們,衪才不會對我們有要求,如果衪不愛我們,根本不需有這些要求,這個我們可從第二方面來看,所以由誡命的角度來看愛上帝,其實沒什麼問題,但我們現在要轉向看第三句說話,即誡命第一個要點,是愛上帝,我現在把重點放在要愛上帝這事情上.

為何基督教信仰的傳統,在維繫上帝與人的關係時,一定要以愛來定義?某方面來說,不是以守誡命來定義.以愛來定義神與人的關係,是不會取消守誡命的重要性,兩者不是相對的,耶穌基督教得很清楚,愛我的人就會守我的誡命,守我的誡命的人,上帝便會以愛回應,所以兩者不是相對的.但對基督教的傳統來說,是強調我們對神的愛,多於我們做這些,做那些來取悅上帝,大家有否發覺這是很奇怪的定義,我想問問大家,在所有的宗教裡,在他們的信眾與他們所信奉的神奇之間,從來也沒有用愛去定義,大家有否聽過佛教徒說,愛佛祖,愛菩薩?你有聽過嗎?大家又有否聽過信奉道教的人,我們要愛黃大仙,愛太上老君呢?我沒有聽過.印度教也不會愛濕婆神,來作為對信眾的要求,不會這樣的,即使與猶太人傳統接近的伊斯蘭教,這個可能有少許保留,但我仍然甚少聽到,他們的信眾要愛安拉.所有的信眾所膜拜的,與他們之間的關係,很多時候是用什麼來形容的呢?是用尊,敬,崇,拜,畏,怕,甚至最多的就是信,但我們很少聽到,是用愛來形容他們之間的關係,為何基督教要獨大,強調這種愛的傳統?有兩個可能的原因,其中一個是因為要藉此來表達神與人之間的關係,在於人這種哺乳類動物,進化最深的需要.

為何這樣說呢?我們時常聽到一句說話,舐犢情深,通常是形容父母對兒女的愛,其實這句說話,我們是借用一些對哺乳類動物的觀察,這個「犢」是指牛小孩,我看見那些牛媽媽,如何愛護那些牛小孩,就是會用他們的舌頭,為他們的牛小孩整理身體,那麼我們中國人便投射進去,形容為一種愛的關係,叫舐犢情深.

哺乳類動物是所有動物當中,在上一代跟下一代的關係中,是愛的關係裡最頂尖的動物,這個跟哺乳類動物的生理特質有關,因為哺乳類動物基本上都是胎生的,所以換句話說,媽媽是能夠在腹中感受兒子的成長,這種接觸能夠感受到生命的真實,這種愛的真實,所以當嬰孩出生,來到這世界時,哺乳類動物的上一代與下一代之間,可以維持很深的關係,我們知道很多哺乳類動物出生後,他們初期的生活都是跟隨父母成長的,一直長大至一個日子,就會像一些獅子,老虎離群相處,這是一種感情的關係.而人更特別,何解呢?因為很多的哺乳類動物出生時,便能夠四處走動,像牛,馬一出生就可以,但人卻不可以,人的嬰孩出生時,好像貓兒那樣細小,很軟弱的,他的腦袋需要一段時間才能成熟,且對四周的環境沒有抵抗能力,所以人類的嬰孩需要很長的時間,經由父母撫養他,保護他,有一些學習告訴我們,這就是人類能夠在哺乳類動物當中脫穎而出的主要原因,因為嬰孩小時候不需要花時間,精神找東西吃,父母會幫助他,所以他便可以專心地發展他的腦袋,如是者,人類便能夠成為萬物之靈.這些都是一些科學的說法,有道理的!於是讓我們看見,人更特別,我們跟父母的關係很微妙,很多時我們對兒女的照顧及愛,是無微不至的,要照顧兒女整個人生的方向,這正可看到感情上的表達,因為人有這方面的感情表達,以致用來說明我們跟上帝當中的愛,那就最方便了.所以當我們說,有一位愛我們的上帝,很像我們的父親愛我們,我們會很明白,因為在人世間我們也能感受相似的愛,我們只要將他轉化,便會明白上帝的愛是什麼意思.這是其中一個說法,來表達基督教裡,神與人的關係.但我覺得並不是這樣,因為我覺得這一解釋是倒轉了,我們將神與人的關係,來把人與人的關係來定義,其實在聖經裡告訴我們,我們人與人的關係是相反的,其實是用神與人的關係來定義及投射才對.因為上帝是愛的上帝,衪對我們的付出,基本上就是愛的上帝.

我們看到,耶穌基督從天上來到我們中間,衪為了什麼要來這這世間?就是要把父上帝的愛帶給我們,讓我們看見,原來父是一位愛的上帝,衪多麼親切地,希望我們能回歸到他們身邊,上帝不是要我們怕衪,上帝不是要我們思前想後,想盡辦法,來做一些討衪喜悅的事,本來衪是喜悅我們,上帝不是要我們日日夜夜活在恐懼當中,想辦法滿足衪,上帝不是叫我們的良心,每一天懼怕地生活,並擔心是否能滿足上帝所有要求,上帝要我們明白,我們要擁有一個無虧欠的良心,無虧欠的良心意思不是要滿足上帝所有要求,而是約翰所說的:「坦然無懼的信心」,我們可以坦然無懼地與上帝接觸.於是我們明白到,之所以用愛來定義神跟人的關係,其實是因為,上帝本身就是一位愛的上帝.因此這一位深深地愛著我們的上帝,衪要求你深深地回應衪,有什麼問題?因為,若果一個人,他不是真正愛你,他跟你毫無關係,他根本不會向你發出這個愛的要求,這個正是第34節,耶穌對有智慧的文士所說,耶穌見他回答的有智慧,就對他說,你離上帝的國不遠了,從此以後,沒有人敢再問他什麼了.耶穌稱讚他離上帝的國不遠了,但我們要留意,不遠的意思即還有一些距離,但已經不錯了,真的挺好.還有一個距離,這個距離你要踏進去才行.你已經掌握了誡命的重要性,誡命的重要其實是建立神與人的關係,當你掌握了那個重要性的時候,你不但只要知道,而是進入那關係當中,去享受上帝給你的慈愛,去以愛來回應,這個才是真真正正上帝的心意,而這一個心意,我們不是在西乃山上見到的,我們是在各各他的屬靈經驗看到的,我們同心禱告.

我們的救主耶穌,我們感謝你,因為你來到我們這個世間,將上帝那個真正的面目向我們顯露,你給我們看到上帝的慈悲,憐恤及樂意寬恕,及讓我們看到,上帝是多麼希望我們回歸到他身邊,去享受衪賜給我們無阻的交通,耶穌我求你給我們在你身上看見其中的重要,並學習如何回歸到上帝的身邊,去享受這份愛,並以愛去回應,感謝主,禱告奉耶穌基督的名字,誠心所願.



\newpage

\section{從桌子掉下來的餅碎}
\label{sec:15}
\textbf{伍渭文牧師 - 從桌子掉下來的餅碎The Scraps from the Table}
\newline
\newline
\hyperref[sec:14]{< < < PREV SERMON < < <}
~
\hyperref[sec:toc]{[返目錄]}
~
\hyperref[sec:16]{> > > NEXT SERMON > > >}


馬可福音 7:24-30
\newline
\begin{longtable}{cl}
\hline
\hline
章節 & 經文 (和合本修訂版)\\
\hline
7:24 & \begin{tabularx}{0.7\textwidth}{X} 耶穌從那裡起身，往推羅境內去，進了一家，他不願意人知道，卻隱藏不住。 \end{tabularx} \\ \\ \relax
7:25 & \begin{tabularx}{0.7\textwidth}{X} 立刻有一個婦人，她的小女兒被污靈附著，一聽見耶穌的事，就來俯伏在他腳前。 \end{tabularx} \\ \\ \relax
7:26 & \begin{tabularx}{0.7\textwidth}{X} 這婦人是希臘人，屬敘利亞的腓尼基族。她求耶穌從她女兒身上趕出那鬼。 \end{tabularx} \\ \\ \relax
7:27 & \begin{tabularx}{0.7\textwidth}{X} 耶穌對她說：「讓孩子們先吃飽，拿孩子的餅丟給小狗吃是不妥的。」 \end{tabularx} \\ \\ \relax
7:28 & \begin{tabularx}{0.7\textwidth}{X} 婦人回答：「主啊，桌子底下的小狗也吃小孩子的碎屑呀！」 \end{tabularx} \\ \\ \relax
7:29 & \begin{tabularx}{0.7\textwidth}{X} 耶穌對她說：「憑著這句話，你回去吧，鬼已經離開你的女兒了。」 \end{tabularx} \\ \\ \relax
7:30 & \begin{tabularx}{0.7\textwidth}{X} 她就回家去，見小孩子躺在床上，鬼已經出去了。 \end{tabularx} \\ \\
[1ex]
\hline
\hline
\end{longtable}
日期:2009年10月4日

題目:從桌子掉下來的餅碎The Scraps from the Table

經文:馬可福音7:24-30

證道:伍渭文牧師

猶太人有一個傳統:收割時不要拾起掉下來的麥穗,讓窮人和孤兒寡婦拿起來充飢.路得記第二章記載摩押女子路得在田間拾穗,被地主波亞斯看見,開始發展一個 浪漫愛情故事.今天的福音書也說到一位外邦女子,乞求從桌子上掉下來的食物,餵養她的女兒.她的女兒被污鬼附身,她俯伏在耶穌面前求耶穌趕走污鬼.

耶穌進入推羅西頓原本是避靜,不想人和道,推羅西頓是加利利北方沿海的外邦人港口.耶穌在世首先在猶太人中間工作,他們是上帝的選民,他就是選民所期待的 彌賽亞.所以耶穌對她說:『讓兒女先吃飽,不好拿兒女的餅丟給狗吃.』這話表面聽來很刺耳,耶穌好像很不近人情,而且語帶刻薄.但這裡的重點,不在兒女和 和狗的分別,乃在得到餅的先後次序.耶穌用當時社會用膳時熟悉的圖畫,說明救恩先臨到猶太人,後才來到外邦人,就像在桌上與父親共席的兒女先得餅食,在桌 下等待的狗才得飽足.重點並不是說外邦人是狗,猶太人是兒女.

上帝的國先在預備好了的土壤,才能生根發芽長大.我們都曉得,種稻米的人,先把種子下在秧田(seedbed) 培育 ,讓它發芽成為秧苗,才播在田里--插秧.猶太人是選民,有上帝的律法,但因為拜偶像,經歷諸般苦難.他們被亞述國人侵,強行雜婚,混亂種族宗教；又被擄 巴比倫,數百年顛沛流離,痛定思痛,才能澈底棄絕偶像,信奉獨一上主,虔守安息日和十誡.回歸後,多年遭羅馬人鐵蹄殖民統治,更使他們盼望彌賽亞的降臨. 福音種子很難在外邦人多神的異教土壤生長.

說外邦人是狗,是得餅次序先後有別,不是貶抑外邦人.再者,在當時社會,狗也是家中要員.狗不是寵物,給我們呵護.狗是守護者,忠誠的朋友,我們幼小兒女的玩伴,晚上睡覺時保護我們,防備盜賊和野獸入侵.所以,用膳時狗也圍繞在桌下,並期待分享應有的一份.

(一)    認識我們的界限和身份:『讓兒女先吃飽』 耶穌同意福音先臨到猶太人,以後才傳到外邦人,讓我們明白,在上帝的國度開展過程中,上主有他的計劃和我們的參與.我們各人有不同的本份,不同的界限,不 同的機會,行事有不同的優先次序.正如大法官李國能先生宣佈退休時引用傳道書三章一節所說:凡事都有定期,天下萬務都有定時.我們都有不同的界限:小時候 是兒女,將來是父母,再遲一些是祖父母.耶穌肯定猶太人和外邦人的界限,提醒我們應盡的責任.這些界限是暫時的,過了期限,便不能返回了.離開工作崗位, 便不能與我們的同事共事了；兒女長大後,到就不能常伴左右了；大學畢業後,便沒有在學時那麼多時間讀書了.機會不會永遠等待你,機會就像一個人從前面走 來,前額有頭髮,可以把握,後面是兀的,過後就捉不住了.當猶太人拒絕福音時,福音就傳到外邦.

今日是九月第一個主日,很多學校己開學了,我們再一次回到學生的身份,回到教師的身份.讓我們為上主守住自已的身份,看重上帝這段時間給我們的界限,曉得自已的優次,接受本份和責任.正如這開學禱文所說:主基督,禰小時候在聖殿專注聆聽智慧；叫我們學習的,不敢虛擲光陰.禰曾坐在摩西的座位訓示律法的奇妙；叫站在課堂的,看重教導者的尊榮和責任. (二)    認識周圍有需要的人:『狗也吃掉下來的碎渣兒』 上主賜我們不同的身份,不同的界限,不同的機會,目的使我們在不同的時間,接觸到不同的人,分享上主給我們的從天上有來的餅.我們要關心在我們周圍有需要 的人,他們就在我們桌子附近,桌子下邊.但很多時候,我們把他們移離視野之外,所以正生書院遷校困難重重.猶太人禁止把葡萄全部摘清,禁止拿走掉在地上的 穗子,使我想起小時候,茶樓飯館晚上快將關門時,每晚不少窮人,乞丐輪候殘餘飯菜.現今社會保障制度,比以前進步得多了,我們卻對周圍有需要的人冷漠起 來,把他們交給有關部門照顧,說我們己納了稅,間接已照顧了他們.

前一段日子,屋村兩位相依為命的精神病患姊妹出了事,姐姐如廁時跌倒失救,妹妹數日惘然守著屍體,後來才給人發現.我們會批評有關醫院,沒有好好跟進她 們.但我想,她們的親戚朋友在那裡呢?她們的鄰居呢?也許現在作鄰里的,期望福利署的家務助理,探訪隔壁的獨居長者,這是他們的責任.我們沒有讓桌子上的 餅碎掉落到地下.上帝給我們不同的界限,目的讓我們接觸到不同的人,在他們有需要時,給予援手.其實他們祇需要掉下來的碎渣兒,就是那麼多.撥一電話問候 一下,停下腳步聽聽他們,寫一個電郵關心,寄一張問候卡致意,花一個小時探訪.

(三)    認識我們內在的自由 .『因你這句話,你回去吧!鬼己經離開你的兒了.她就回家去,見小孩躺在床上,鬼己經出去了.』(29,30節)

誰是上主的兒女呢?信心使我們成為上帝的兒女.在前面我們說過,『讓兒女先吃飽,不好拿兒女的餅丟給狗吃.』這句話並不是指猶太人比外邦人優越,有特權, 乃福音傳播的先後次序.事實上,能夠與成為上帝的兒女與父親一起坐席,是這外邦女人,不是反對耶穌的猶太人.因著信心,這女人的女兒己被治癒了.『因你這 句話,你回去吧!鬼己經離開你的女兒了.』也許我們身份不同,界限不同,機會也不同,但我們有超越任何外在身份的內在自由(inner freedom),不受任何界限制約.對主的信心,超越一切.『並不分是猶太人,希利尼人,自主的,為奴的,或男或女,因為你們在基督裡都成為一 了.』(加3;28)

誰是不潔淨的呢? 這故事的背景,是猶太人責難耶穌門徒沒有遵守飯前清潔的禮儀.他們認為若不洗手,就不應吃飯；從市上回來也要洗手,表示不被外邦人所沾污.耶穌就指出,從 裡面的才能污穢人,從外面進入的,便落到茅廁裡.『從人心裡發出惡念,苟合,偷盜,兇殺,姦淫,貪婪,邪惡,.......』(21, 22節)

這些拒絕耶穌,反對耶穌的的猶太人,正被心裡的污穢所充滿,與被污鬼附的身的女孩沒有分別.不同的是女孩的污穢是可見的,明顯的；他們的污穢是隱藏的.最近政府宣傳吸毒的壞處,繪形繪聲；索k會生冰瘡,胡言亂語,小便失禁,樣子愈來愈醜陋.不可一,不可再.吸入的毒品的禍患,就寫在你的臉上.

但內裡出來的污穢更可怕,很難從臉相看出來,而且愈犯罪,臉相愈可能迷惑人.我曾經在講道時,提及王爾德(Oscar Wilde) 的The Picture of Dorian Gray(道林格雷的畫像).天生漂亮異常的道林格雷因見了畫家霍華德給他畫的真人一樣大的肖像,發現了自己驚人的美,但為自己韶華易逝,美貌難久感到痛苦,希望那幅肖像能代替自己承擔歲月和心靈的負擔,而讓自己永遠保持青春貌美.他愈犯罪愈年青,相貌比以前更俊朗,但在閣樓收起來的肖像開始變化,咀角漸漸有一絲陰陰的笑.但Dorian享受幽暗的生活,兩面人行為使他感到刺激興奮,他殺人二十四小時後,可以若無其事參加舞會.他的罪行和污穢不是寫在臉上,寫在收藏起來的肖像上.最後,他不能接受自已愈來愈醜陋的肖像,試圖用刀劃破肖像並瘋狂大叫,最後倒在地下死了.後來人發現的,是一具相貌醜陋的屍體,祇能從手上的戒指,才認出是道林格雷,而旁邊的肖像,這時候就回復原來的俊朗.

道林格雷的肖像並非虛擬的小說.當華爾街經紀人,納斯達克前董事會主席伯納德‧麥道夫(Bernard L. Modoff) 被揭發500億美元巨額欺詐時,我們發現不少富商巨賈,達官貴人,在案發之前,爭相求見親炙他.他們被麥道夫的氣派所迷倒,當然也是財迷心竅想發達.愈貪 婪愈使麥道夫看來愈有氣派,愈有自信.有人說,有婚外情的人,開始時有可能特別關心配偶,比以前更體貼(但有時可能是遲來的覺醒) .

但行過的路徑,總會留下痕跡.耶穌進到推羅,西頓本來『不願意人知道,卻隱藏不住.』同理,罪惡總有一天被揭發,位高權重的,也會下在獄中,衹是上帝有恩 典,等待我們趕快悔過.願上主在新的學年的開始,使我們的靈魂甦醒.主昔日願意逐出污鬼,今天也願意驅走我們的心魔,祇要我們渴想被潔淨.  結論:從 這故事,我們認識到,若我們承認自已的污穢,信靠耶穌尋求潔淨,必能得到生命的餅,可以與父一同坐席,被稱為上帝的兒女.這些人才是上帝真正的選民,才是 真以色列人.猶太人,外邦人的界限已因基督的來臨成為過去了.但我們仍然有自已人生的不同界限,不同的身份.這界限和身份,目的用來接觸周圍有需要的人, 我們己得到生命的餅,讓我們關心那些在我們桌子附近的人,在桌子下邊的人,他們是上帝交給我們照顧的人.



\newpage

\section{外人在摩押地}
\label{sec:16}
\textbf{陳南芬教授 - 外人在摩押地 Aliens in Land of Moah}
\newline
\newline
\hyperref[sec:15]{< < < PREV SERMON < < <}
~
\hyperref[sec:toc]{[返目錄]}
~
\hyperref[sec:17]{> > > NEXT SERMON > > >}


日期:2009年10月11日

題目:外人在摩押地 Aliens in Land of Moah

證道:陳南芬教授

大家好!不知道我們當中有沒有飄泊異地,客居他鄉的經驗呢?其實我估計在座當中,應該有很多都不是本地人,對吧?不知道大家在外地居住的經驗是正面的還是負面的呢?當時候有沒有住在那裡的本地人幫助你適應新環境呢?還是你要拚命地去找同鄉朋友幫助你呢?我想如果有語言不通的情況,可能是更加的困難吧!

在當今發達時代,要到一個陌生的地方去,當然沒有十年前那麼辛苦,我們可以在舒適的家裡或網吧上網,通過 google,百度等,就可以下載我們所需要的資料,譬如:天氣情況,當地歷史,經濟,貿易,購物,交通情況,以及飲食習慣和文化風俗等等；你甚至可以到google map輸入你要住的地址,就可以看到並熟悉一下周圍的環境如何；如果有言語差別的話,在網上也可以找到一些初級的語言課程來補習當地的語言,做好搬遷預備.

我想我們有很多大陸的學生來到神學院,其中有一位姐妹很厲害,沒到香港之前,就買了粵語的連續劇,學了一下廣東話,現在雖然到香港只有一年,但廣東話講得非常好,是我們神學院其中一個很好的榜樣.

這些都是我們在科技發達生活的好處,但是我們卻可以想像一下,在聖經時代,要到一個陌生的地方居住,是有多麼的困難呀!路得記除了告訴我們,這摩押女子怎樣當了大衛王的祖先,它也描述了當時的社會怎樣去接待外來的人,今天我們就從這個角度來讀路得記,讓聖靈來教導我們.

故事開始的時候,以色列國正遭遇飢荒,伯利恆有一家人就逃到摩押去,故事開頭表明當時正處於士師時代,那時以色列人民已經逃離埃及的掌控,並剛剛結束曠野時期,定居在應許之地,也就是說距離曠野飄泊的日子不是太久,要是這樣子的話,這一家四口選擇摩押去逃難,應該是一個非常艱難的決定,為什麼這麼說呢?

如果可以參考以色列的地圖,你會發現從伯利恆當時可供搬遷的地,只有亞捫和摩押,但能夠選擇的就只有摩押,因為士師記十一章告訴我們,亞捫是富有侵略性的,而摩押沒有攻擊以色列的意圖,然而這一家四口,遷居摩押卻並非沒有危險,因為通過民數記和申命記,我們知道以色列人和摩押人的關係曾經像仇敵一般,但是這一家四口到了摩押,卻把摩押當成了自己的家鄉,儘管以色列的飢荒過去有一段時期了,這一家卻選擇繼續留在摩押定居了來,原恩惠是甚麼呢?

我的家人陪我到英國深造的時候,遇見一連串不愉快的事件,就是跟我們租房子和跟BT(就是電話公司)接觸的一些負面的經驗,這些讓我們深深品嚐到被英國人欺壓的感覺.之前我們是宣教士,去過很多地方,住過至少有一年或是更長的經驗,都沒有那種負面的經歷,所以我的先生特別不喜歡英國,到今天還是在唸英國的不好,你現在叫他去英國,他才不肯呢!同樣的,如果你對某個地方有負面的經歷,而可以選擇到另一個地方去,你一定選擇使你不愉快的地方,但是如果你的經歷都是順利,愉快的,你就會對那個地方產生一種依戀不捨的感覺.

我們應該可以肯定這一家四口到了摩押之後,受到當地居民很好的招待,摩押人還把他們兩位樣貌出色的女子嫁給這一家的兩位公子,於是這一家四口一直沒有回國的念頭,並決定在那裡立地生根,摩押人也不介意這一家四口所敬拜的神,跟他們的不一樣,總而言之,摩押人好像對這家人很好,慷慨,又肯接納這家外來的人,使他們都不想回國了.不幸的是,這家的男丁都死了,家婆拿俄米在那痛苦的經歷之後,才決定回國,她兩位摩押媳婦卻是那麼的孝順,明知家婆要回國重新生活,她們卻因為忠誠而不顧自己生命危險,適應能力,生活情況,也沒考慮到以色列會排斥她們等等,一心一意堅持要跟著家婆.拿俄米是費了很多的淚水和勸解才能夠說服其中一個媳婦不要跟她,但是路得緊緊抱住她的家婆,拿俄米也沒辦法了,只好同意路得跟著她.

之前是有以色列人到摩押去,受到摩押人優越招待,現在我們來看一下,以色列人又怎樣招待路得這摩押來的寡婦?

在基督教傳統,一提起路得記就很快的斷定,路得能夠成為大衛王的祖先是她的榮耀,因為路得顯出她對以色列人的忠誠,所以上帝賜下這個福份給她,這無疑是以色列禮儀的角度所做出判斷,但是如果路得有機會發表她的心聲,她的故事又會是怎樣的呢?路得和她的家婆拿俄米一到達伯利恆,便遇到了幾個巧合.

第一個巧合:就是大麥收割的時候.路得向她家婆請安去收割──根據申命記二十四章19節:「你在田間收割莊稼,若忘下一捆,不可回去再取,要留給寄居的,與孤兒寡婦,這樣,耶和華你　神必在你手裡所辦的一切事上,賜福與你.」

第二個巧合:就是路得在拿俄米丈夫的親戚大財主波阿斯的田地裡收割.

第三個巧合:就是波阿斯又恰好從眾多的貧窮人,就是寄居的,孤兒寡婦中,偏偏注意到路得,並且恩待她.

故事就這樣的發展,收割一完成,拿俄米很快的就吩咐路得,要照她的計劃這樣,那樣做.第三章2節,拿俄米對路得說:「你與波阿斯的使女常在一處,波阿斯不是我們的親族麼．他今夜在場上簸大麥,你要沐浴抹膏,換上衣服,下到場上,卻不要使那人認出你來,你等他喫喝完了, 到他睡的時候,你看準他睡的地方,就進去掀開他腳上的被,躺臥在那裡,他必告訴你所當作的事.」就是要趁波阿斯成為拿俄米的親族,「掀開他腳上的被」是性交的微婉語,所以這個計劃是需要路得願意跟波阿斯過夜.在第三章1節中,拿俄米解釋她的動機,即是幫路得找到一個安身之處.

有幾樣東西可能必須解釋一下:第一,就是在第二章波阿斯與路得的對話當中,波阿斯不斷稱呼路得「女兒」,這裡包涵兩個事實── 一,波阿斯的年紀比路得應該大很多；二,波阿斯應該有妻子和自己的孩子了.第二,所謂救助人的意思又是甚麼呢?其實在申命記二十五章5-10節:「弟兄同居,若死了一個沒有兒子,死人的妻不可出嫁外人,他丈夫的兄弟當盡弟兄的本分娶他為妻,與他同房.婦人生的長子必歸死兄的名下,免得他的名在以色列中塗抹了.那人若不願意娶他哥哥的妻,他哥哥的妻就要到城門長老那裡說,我丈夫的兄弟不肯在以色列中興起他哥哥的名字,不給我盡弟兄的本分.本城的長老就要召那人來問他,他若執意說,我不願意娶他．他哥哥的妻就要當著長老到那人的跟前,脫了他的鞋,吐唾沫在他臉上,說,凡不為哥哥建立家室的,都要這樣待他.在以色列中,他的名必稱為脫鞋之家.」

(我第一次到大陸的時候,就發現大陸的人很喜歡吐沫,到處吐沫,那麼我每次讀到這裡的時候,不可以想像... 可能古代以色列人他們也很喜歡吐沫,只是現在找到機會可以吐沫在人的臉上,我自覺這非常有趣.)

我們回到這個特別的風俗(不是吐沫的風俗),乃是要幫兄弟傳宗接代,這個目的就是好讓死者可以有後代繼承他的產業,所以這個救助人他應該也有自己的孩子,繼承他自己的產業.

第三點要解釋的是,從某一方面我們可以說,拿俄米是為了她和路得兩個的將來著想,而編了這個計劃,但是在第三章裡,路得只有順服,沒有表達出她對這個計劃的意見；還有路得當晚在床上跟波阿斯過夜時,才發現其實還有一個更親的親戚有權救助拿俄米的,路得當晚可說是冒險犧牲了自己.

第四,不曉得你們有沒有想過,為什麼當天不是拿俄米跑到場上很認真跟波阿斯談判這件事呢?當時的社會,若一個女性有長輩的身份跟男性談話是允許的.但拿俄米沒有這樣做,只有一個原因,就是波阿斯肯定會拒絕,因為還有一個更加親的親戚有權救助拿俄米,所以說這個計劃可算是一個圈套,波阿斯若是一個君子,不能不為他喝醉時所做的一切不負責任,況且這個女子是路得,是他有責任照顧的遠親.路得,這外女,來到一個陌生的地方,她的真誠就是要服侍她的家婆,願意跟她相依為命過日子.恰巧她跟著收割的人都沒有欺負她,又受了波阿斯的恩典,可以跟著他的使女們收割,休息,吃飯,若拿俄米沒有想出這個妙計,可能路得繼續在波阿斯的田地裡幹活吧.

路得這女子,真的是一個外族的好榜樣,在她家鄉時,沒有虧待過外族的家婆；到陌生的地方時,也沒有違反外族的風俗習慣,她既樸實又順從,婦女們都稱讚路得的孝心比七個兒子更好,這稱讚也成為拿俄米的福份.路得對她外族的家婆忠誠,信任並孝順,從拿俄米對路得這外族女子的態度中,我們也找到同樣的尊重和招待嗎?

現在就讓我們反省一下,我們到底為了周圍外來的朋友們,做了些甚麼呢?我們怎樣招待他們呢?是像摩押人那樣,讓外來的,有難的人感受到歡迎和接納,還是使他們感到孤獨,冷漠和排斥呢?還以上分析可以看到路得記非常注重忠誠,信任和尊重等內在品格,但我們的品格又怎樣呢?是否在內心裡看不慣,看不起他們,甚至厭惡,冷漠他們呢?也許我們都知道應該接待外人,但是我們是出於甚麼的動機呢?是真的出於愛心和尊重嗎?或者是趁著他們對我們的制度,規則和風俗習慣還不太了解時,把他當作我們獲取自己利益的工具而利用他們呢?

結束之前,想提一句經文給各位弟兄姊妹,希伯來書十三章2節上:「不可忘記用愛心接待客旅．」

我們一齊禱告:上主,沒有你的恩典,我們今天沒有可能在這裡跟我們的弟兄姊妹一齊分享你的愛,慶祝你的救恩,今天你的聖靈所感動我們的,給我們力量去改變我們自己,幫助我們去體貼他人,使他人都能從我們的榜樣中榮耀你的名,我們這樣禱告,奉我們救主聖名,阿們.



\newpage

\section{領袖僕人和小孩}
\label{sec:17}
\textbf{伍渭文牧師 - 往耶路撒冷的路上--領袖,僕人,和小孩On the Road to Jerusalem-Leader, Servant, and the Child}
\newline
\newline
\hyperref[sec:16]{< < < PREV SERMON < < <}
~
\hyperref[sec:toc]{[返目錄]}
~
\hyperref[sec:18]{> > > NEXT SERMON > > >}


馬可福音 9:30-37
\newline
\begin{longtable}{cl}
\hline
\hline
章節 & 經文 (和合本修訂版)\\
\hline
9:30 & \begin{tabularx}{0.7\textwidth}{X} 他們離開那地方，經過加利利；耶穌不願意人知道， \end{tabularx} \\ \\ \relax
9:31 & \begin{tabularx}{0.7\textwidth}{X} 因為他正教導門徒說：「人子將要被交在人手裡，他們要殺害他；被殺以後，三天後他要復活。」 \end{tabularx} \\ \\ \relax
9:32 & \begin{tabularx}{0.7\textwidth}{X} 門徒卻不明白這話，又不敢問他。 \end{tabularx} \\ \\ \relax
9:33 & \begin{tabularx}{0.7\textwidth}{X} 他們來到迦百農。耶穌在屋裡問門徒說：「你們在路上議論的是甚麼？」 \end{tabularx} \\ \\ \relax
9:34 & \begin{tabularx}{0.7\textwidth}{X} 門徒不作聲，因為他們在路上彼此爭論誰最大。 \end{tabularx} \\ \\ \relax
9:35 & \begin{tabularx}{0.7\textwidth}{X} 耶穌坐下，叫十二個使徒來，說：「若有人願意為首，他要作眾人之後，作眾人的用人。」 \end{tabularx} \\ \\ \relax
9:36 & \begin{tabularx}{0.7\textwidth}{X} 於是耶穌領一個小孩過來，讓他站在門徒當中，又抱起他來，對他們說： \end{tabularx} \\ \\ \relax
9:37 & \begin{tabularx}{0.7\textwidth}{X} 「凡為我的名接納一個像這小孩子的，就是接納我；凡接納我的，不是接納我，而是接納那差我來的。」 \end{tabularx} \\ \\
[1ex]
\hline
\hline
\end{longtable}
日期:2009年10月18日

題目:往耶路撒冷的路上--領袖,僕人,和小孩On the Road to Jerusalem-Leader, Servant, and the Child

經文:馬可福音9:30-37

證道:伍渭文牧師

入直路了!入直路了!對很多初到香港的人,這話沒有甚麼特別意思,但對久居此地,特別是喜愛賽馬人士,『入直路了』是一場賽事的高潮.進入直路就不需再繞圈了,進入直路就見到終點；祇要加把勁,就可以奪標.在主耶穌的一生,甚麼時候是他的直路呢?進入耶路撒冷.四福音的高潮是主耶穌在地上最後的一週 and \#8212;騎著驢駒子進入耶路撒冷,與各方敵人公開衝突,直到被出賣釘死在十字架,復活,升天.

今天所讀的福音經課,門徒爭論誰為大,就發生在往耶路撒冷的路上.『你們在路上議論的是甚麼?門徒不作聲,因為他們在路上爭論彼此為大.』(可8:33,34)這語法很特別,其實可以略去在路上三個字.當青年財主尋求永生之道時,馬可也特別描述他在『耶穌出來行路的時候』(10:17),而同樣的事件馬太是『有一個人來求見耶穌』(太19:16)；路加是『有一個官問耶穌說』(18;18).馬可對往耶路撒冷的路,特別情有獨鍾.

為了預備這路的旅程,耶穌特別退隱到加利利以北的外邦人聚居的推羅,西頓,『進了一家,不願意人和道,卻隱藏不住.』(7:24)我上次講道:從桌子掉下來的餅碎,耶穌指出他的首要目標是猶太人,『讓兒女先吃飽,不好拿兒女的餅丟給狗吃』(7:27)他的目標就是往耶路撒冷的路.

這路的起點是該撒利亞腓立比,在那裡耶穌第一次公開承認自已是基督,也第一次預言基督要上耶路撒冷受苦受死.(8:27-31)跟著耶穌帶著雅各,約翰,彼得登上黑門山登山變像；下山時醫好被污鬼附身的孩子；回到靠近加利利的迦伯農,作第二次受苦預言,這是今天的講道經文九章30-32節.第三次預言受苦是接近耶利哥,耶路撒冷在望了!所以耶穌的神情愈凝重,令門徒害怕起來.『他們行路上耶路撒冷去.耶穌在前頭走,門徒就希奇,跟從的人也害怕.』(10:32)

耶穌三次預言受苦,但三次門徒都不明白.彼得雖然第一個人認信耶穌是基督,但不能接受他上耶路撒冷受苦,『彼得就拉住他,勸他』(8:32)耶穌就轉過身來責備他:撒旦,退我後邊去吧!因為你不體貼神的意思,祇體貼人的意思.第三次雅各和約翰爭取坐在耶穌左右；今天經文是第二次,也很糟糕,在路上門徒爭論誰為大,顯出他們的無知,畢竟人的路與上主的路是不同的.

(一)往耶路撒冷的路與往華爾街的路 and \#8212;領袖與僕人 往耶路撒冷的路是捨己的,甘心放下的.耶穌是領袖(Leader),但甘心為人捨去生命.應該生的,竟然死去,但應該死的,竟然活下來,而且活得比以前壯大.誰人?我是說華爾街的貪婪.前數天是金融海嘯一周年,我們發現貪婪並未死去,而且有捲土重來的苗頭.有些金融機構,大到不能沒有它們,像美國的AIG,經濟學人說了很諷刺的話:當今的問題是:不適者生存 (The survival of the unfittest),因為投鼠忌器,影響太大了,就逃過被砍掉的命運!

貪婪不是華爾街的專利,貪婪也在我們心中.我們不一定貪財,但會像門徒一樣,貪圖首位,並認為憑努力得到高位,是應得的報酬,沒有不妥.就如這些在華爾街的金融精英認為,在法律容許下,運用精算財技,為客戶賺取巨額金錢,也應分享高薪厚祿,何罪有之?

往華爾街的路,與往耶路撒冷的路是不同的.華爾街自有它的魅力,這裡代表財富,權力,地位,和才智.對,才智,因為最好的高考成績的一群,才能入讀全球金融糸,最好畢業成績加上運氣,才能進入國際投資銀行.所以門徒不明白往耶路撤冷的路,跟隨一個受死的基督是他們難以接受的.『門徒卻不明白這話,又不敢問他.』(9:32) (二)接待小孩與接待主

請注意,耶穌並沒有否定我們應有的雄心,力爭上游,出類拔萃.耶穌祇是顛覆了目標:『若有人願意作首先的,他必需作眾人末後的,作眾人的用人.』(9:35)作個服待人的領袖,這是劃時代的觀念.這觀念也啟發了現代管理學,學者認為,為僕的領袖servant leader,有十種特性:善於聆聽listening, 具同理心empathy, 醫治修復healing, 警覺性強awareness,具說服力persuasion, 觀念清楚conceptualization, 有遠見foresight, 有托付感stewardship, 建立他人commitment to the growth of others, 建立社群and building community.

具有這些素質的領袖,誰人不願跟隨呢?但耶穌在這裡所講的,是希望門徒成為令人心悅誠服的領袖嗎?為僕的領袖最終目標是有效的領導(leadersjip) ,而耶穌所說的是真正的僕人.

我們注意耶穌跟著的行動.耶穌坐下,叫十二個門徒來,說:『若有人願意作首先的,他必作眾人末後的,作眾人的用人.於是領一個孩子來,叫他站在門徒中間,又抱起他來,對他們說:凡為我的名接的待一個像這小孩子的,不是接待我,乃是接待那差我來的.』(9:36, 37)

小孩本身的需要:接待小孩純粹是為小孩的需要,耶穌時代的小孩沒有獨立性,是父母的附屬品,不像現代社會有保障小孩的法律.也不像現代社會,小孩往往是重要的消費對象,因為父母疼愛他們,一個自由行來香港的父母,一定會買奶粉給孩子.耶穌時代的孩子,是沒有甚麼地位的,接待他沒有回報的.耶穌講道時,有小孩子想親近耶穌,也被門徒隔開,他們沒有地位.然而耶穌說,讓隊孩子到我這裡來,抱起他們,祝福他們.

效法小孩的信心:請注意這裡用接待,不是接納,接待是積極的,接納是消極的.而且要像接待耶穌般接待他,而接待耶穌,就是接待差耶穌來的父.我們注意這幅圖畫:耶穌坐下,這是教導的姿態,小孩在他旁邊站,十二門徒在四周.這位小孩在十二使徒之列,而且最靠近耶穌,何等尊榮.

在這圖像中,誰是最大呢?最接近耶穌的小孩最大,因他最接近耶穌,貼近耶穌的胸懷,小孩子最能表明因信才能進入天國的道理.『凡為我的名接的待一個像這小孩子的,不是接待我,乃是接待那差我來的.』接待耶穌,就是接待差耶穌來的父.進入天國是信靠,沒有人可以靠自已進入天國,進天國的要像小孩子,因他們信賴的心最單純.我們都知道,不能隨便答應小孩子,你怎樣說他們就怎樣信.『我們得救是本乎恩,也因著信,這並不是出於自已,乃是神所賜的,也不是出於行為,免得有人自誇.』(Eph. 2: 8,9)

門徒在往耶路撒冷的路上彼此爭論誰是最大.但這幅圖畫告訴我們,那些接待小孩子,尊崇他像尊崇上主,學效法他們的信心,他們就是天國最大的,因為天國靠接受上主的恩典,不是靠人的努力.

結論:有一次,國際大炒家索羅斯(Soros) ,這位曾被讚譽為華爾街最出色的基金經理接受訪問,坦承他確曾1998年衝擊港元,但他說這是合法參與國際金融遊戲,沒有不妥,誰叫我們不好好堵塞漏洞.讚揚他的人說:他關心人權,並慷慨捐助各地區的人概權發展.他左手自金融投機市場上賺取大把鈔票,右手再親身參與將之用於慈善公益事業,以基金會組織培育反對人士,加速共產主義的崩潰,同時創造就業機會,來盡量紓解轉型所帶來的痛苦.

索羅斯用金錢暴力(也是一種暴力),用金融的領導地位,來成就他的理想,去塑造他的世界；耶穌承受人的暴力,以僕人身分,用自已的生命來救贖這世界.因他知道人類問題的徵結,不在任何政治制度,乃人的罪性,而若不流血,罪不得赦免.他要走上往耶路撒冷的路上,為人贖罪.

往耶路撒冷的路,是耶穌要走的路,也是每一個門徒要走的路.



\newpage

\section{拴在頸上的石磨}
\label{sec:18}
\textbf{林豪恩先生 - 拴在頸上的石磨(A millstone on the neck)}
\newline
\newline
\hyperref[sec:17]{< < < PREV SERMON < < <}
~
\hyperref[sec:toc]{[返目錄]}
~
\hyperref[sec:19]{> > > NEXT SERMON > > >}


馬可福音 9:38-50
\newline
\begin{longtable}{cl}
\hline
\hline
章節 & 經文 (和合本修訂版)\\
\hline
9:38 & \begin{tabularx}{0.7\textwidth}{X} 約翰對耶穌說：「老師，我們看見一個人奉你的名趕鬼，我們就阻止他，因為他不跟從我們。」 \end{tabularx} \\ \\ \relax
9:39 & \begin{tabularx}{0.7\textwidth}{X} 耶穌說：「不要阻止他，因為沒有人奉我的名行異能，反倒輕易毀謗我。 \end{tabularx} \\ \\ \relax
9:40 & \begin{tabularx}{0.7\textwidth}{X} 不抵擋我們的，就是幫助我們的。 \end{tabularx} \\ \\ \relax
9:41 & \begin{tabularx}{0.7\textwidth}{X} 凡因你們是屬基督，給你們一杯水喝的，我實在告訴你們，他一定會得到賞賜。」 \end{tabularx} \\ \\ \relax
9:42 & \begin{tabularx}{0.7\textwidth}{X} 「凡使這些信我的小子中的一個跌倒的，倒不如把大磨石拴在這人的頸項上，扔在海裡。 \end{tabularx} \\ \\ \relax
9:43 & \begin{tabularx}{0.7\textwidth}{X} 如果你一隻手使你跌倒，就把它砍下來； \end{tabularx} \\ \\ \relax
9:44 & \begin{tabularx}{0.7\textwidth}{X} 你缺一隻手進入永生，比有兩隻手落到地獄，入那不滅的火裡去還好。 \end{tabularx} \\ \\ \relax
9:45 & \begin{tabularx}{0.7\textwidth}{X} 如果你一隻腳使你跌倒，就把它砍下來； \end{tabularx} \\ \\ \relax
9:46 & \begin{tabularx}{0.7\textwidth}{X} 你瘸腿進入永生，比有兩隻腳被扔進地獄裡還好。 \end{tabularx} \\ \\ \relax
9:47 & \begin{tabularx}{0.7\textwidth}{X} 如果你一隻眼使你跌倒，就去掉它；你只有一隻眼進入神的國，比有兩隻眼被扔進地獄裡還好。 \end{tabularx} \\ \\ \relax
9:48 & \begin{tabularx}{0.7\textwidth}{X} 在那裡，蟲是不死的，火是不滅的。 \end{tabularx} \\ \\ \relax
9:49 & \begin{tabularx}{0.7\textwidth}{X} 因為每個人必被火像鹽一般醃起來。凡祭物必用鹽醃。 \end{tabularx} \\ \\ \relax
9:50 & \begin{tabularx}{0.7\textwidth}{X} 鹽本是好的，若失了鹹味，你們怎能用它調味呢？你們中間要有鹽，彼此和睦。」 \end{tabularx} \\ \\
[1ex]
\hline
\hline
\end{longtable}
日期:2009年10月25日

題目:拴在頸上的石磨(A millstone on the neck)

經文:馬可福音9:38-50

證道:林豪恩先生

一．    引言把「拴在頸上的石磨」形象化地想像一下,你會看見怎樣的圖畫呢?你看見的圖畫,可能與你對石磨的認識有關.石磨由兩塊石塊上下緊密貼合組成,下面的作為基座,上面的石塊能夠旋轉,當上面的石塊旋轉而與下面的基座磨擦時,就可以把穀物磨成細碎粉粒.現在常見的是家用石磨,有人稱之為「石磨仔」,這種「石磨仔」只是如一個生日蛋糕般大小,用人手攪動.若把這個石磨上半可旋轉的部份掛在頸項上,對於我來說當然感到辛苦,但對於如姚明那麼高大的人來說,或者可以視作飾物.

不過,聖經時代的石磨,是比我們現在的家用石磨要大很多倍.如此龐大的石磨,是不會用人來拉動的,而必須以動物拉動的.在中國,這樣的大石磨通常用牛來拉動；在聖經時代的巴勒斯坦,他們通常找驢來把石磨拉動.試想像,把這塊套在牛或驢的頸項上的石磨套在你幼小脆弱的頸項上,會是多麼可怕的狀況.「把大磨石拴在這人的頸項上」此句,是來自今天所講的福音經課──新約聖經馬可福音9章35-38節,當我們細閱經文時,會發現經文描述的狀況比我們想像的更極端及可怕.

二．    經文

38約翰對耶穌說:「夫子,我們看見一個人奉你的名趕鬼,我們就禁止他,因為他不跟從我們.」

39耶穌說:「不要禁止他；因為沒有人奉我名行異能,反倒輕易毀謗我.

40不敵擋我們的,就是幫助我們的.

41凡因你們是屬基督,給你們一杯水喝的,我實在告訴你們,他不能不得賞賜.」

42「凡使這信我的一個小子跌倒的,倒不如把大磨石拴在這人的頸項上,扔在海裡.

43倘若你一隻手叫你跌倒,就把它砍下來；

44你缺了肢體進入永生,強如有兩隻手落到地獄,入那不滅的火裡去.

45倘若你一隻腳叫你跌倒,就把它砍下來；

46你瘸腿進入永生,強如有兩隻腳被丟在地獄裡.

47倘若你一隻眼叫你跌倒,就去掉它；你只有一隻眼進入上帝的國,強如有兩隻眼被丟在地獄裡.

48在那裡,蟲是不死的,火是不滅的.

49因為必用火當鹽醃各人.

50鹽本是好的,若失了味,可用甚麼叫它再鹹呢?你們裡頭應當有鹽,彼此和睦.」

三．    釋經與應用

1．對比顯輕重　抉擇宜慎重

這段經文的第42節如此說:「凡使這信我的一個小子跌倒的,倒不如把大磨石拴在這人的頸項上,扔在海裡.」把大磨石套在一個人的頸項上,並不是要勞役這個人,而是要置他於死地；因為接著的行動是把這個人連人帶石磨拋下海中,如此,這個人必死無疑.這樣的處決行動,與我們古代的「浸豬籠」類似,但卻更迅速及徹底.耶穌描述「把大磨石拴在這人的頸項上,扔在海裡」這樣恐怖的狀況,為的是與「凡使這信我的一個小子跌倒」作強烈的對比.如此強烈的對比,在這個段落中共有四對之多.

對比一:「凡使這信我的一個小子跌倒的,倒不如把大磨石拴在這人的頸項上,扔在海裡.」(42)對比二:「倘若你一隻手叫你跌倒,就把它砍下來；你缺了肢體進入永生,強如有兩隻手落到地獄,入那不滅的火裡去.」(43-44)對比三:「倘若你一隻腳叫你跌倒,就把它砍下來；你瘸腿進入永生,強如有兩隻腳被丟在地獄裡.」(45-46)對比四:「倘若你一隻眼叫你跌倒,就去掉它；你只有一隻眼進入上帝的國,強如有兩隻眼被丟在地獄裡.」(47)

耶穌用「把大磨石拴在這人的頸項上,扔在海裡」,「砍下一隻手」,「砍下一隻腳」,「去掉一隻眼」這些激烈的行動來與「使一個小子跌倒」或「叫自己跌倒」作對比,明顯是要告訴聽眾「使人跌倒或叫自己跌倒」其實比這些激烈的行動更激烈,比那些可怕的後果更可怕.

在現實生活中,我們也運用對比,因為對比能夠逼使人們去思考事情的大小輕重.例如社會福利署宣傳反家庭暴力的口號是這樣的:「贏了一場架,輸了一個家.值得嗎?」輸一場架,故然不好受,但若輸一個家,則所有人也受到傷害,其破壞程度可以非常巨大.兩害取其輕,結論當然是不值得為一場架而輸一個家.另外一些處境,是兩項選擇都是有利益的,然而,就算純粹從利益的角度考慮,也必須清醒地比較一下兩者利益的大小,否則會跌入「因小失大」的陷阱之中.

對於一位警員來說,怎樣的交換條件能夠令他願意冒失去工作,失去宿舍,失去100萬元的長俸危險呢?合理估計,這一定是「大茶飯」,所得到的利益一定件一定要比現在的工作報酬及福利更優厚,事實卻令人意外.

貪200元拾遺不報警失百萬長俸累家人【明報專訊2007年7月7日】兩名熱心市民路不拾遺,把拾到的手提袋交給警員,警員卻監守自盜將手提包據為己有,偷去其中的200多元現金及禮券,結果被揭發.警員昨承認盜竊罪,求情時指事件已令他前途盡毀,失去工作,失去宿舍,及失去近100萬元長俸.

透過對比,我們發現有些行動是不值得的,有些事情是更重要的.

據福音書的記載,耶穌經常運用對比,逼使聽眾分辨事情的大小和輕重,進而要求聽眾列出優次,作出選擇.回到這段經文所提出的四對對比,相信門徒不難明白耶穌的意思,就是:如果你「使信我(耶穌)的一個小子跌倒」比失去自己的生命的後果更嚴重；如果使自己跌倒,比失去一隻手,或者失去一隻腳,去掉一隻眼睛的後果更嚴重.

「跌倒」是怎樣的情況呢?為甚麼「跌倒」的後果如此嚴重呢?我們可以參考馬可福音第四章「撒種的比喻」,耶穌解釋比喻時說:

「撒種之人所撒的就是道.那撒在路旁的,就是人聽了道,撒旦立刻來,把撒在他心裡的道奪去.那撒在石頭地上的,就是人聽了道,立刻喜歡領受,但他心裡沒有根,不過是暫時的,及至為道遭了患難,或是受了逼迫,立刻就跌倒了.」(4:14-17)

馬可福音第十四章耶穌預言彼得不認主的那一段也提及「跌倒」:

「耶穌對他們說:你們都要跌倒了,因為經上記著說:我要擊打牧人,羊就分散了.但我復活以後,要在你們以先往加利利去.彼得說:眾人雖然跌倒,我總不能.」(14:27-29)

由此可見,「跌倒」是一種象徵式的說法,所指的不是身體的,不是生物性及物理性的,而是靈性的,是與信仰上帝相關的,是與跟隨耶穌相關的.因此,「現代中文譯本」把「跌倒」翻譯為「離棄耶穌」.把失去性命以及失去肢體與「離棄耶穌」作對比,似乎是要求聽眾有需要時在「保存生命及身體」以及在「繼續跟隨耶穌」兩者之間排列優次及作出抉擇.

2．身體性命確可貴,信仰上帝更永久

把失去性命以及失去肢體這些令人恐懼和不幸的事情與「離棄耶穌」作對比,有沒有令作為聽眾的你感到不安,甚至感到震驚?「語不驚人誓不休」的效果,就是要刺激聽眾的反應,激盪聽眾的思想,挑戰聽眾固有的觀念,期望聽眾因此而作出改變.

為甚麼用失去性命和失去肢體做對比呢?因為關注自己的身體,關注自己的生命,是人的自然反應.保羅也說:「從來沒有人恨惡自己的身子,總是保養顧惜」(弗5:29).因此,自古至今人們都尋求各種和樣的方法去達至延年益壽,以及使身體更好看.有人認為,對身體的關注,現代人比古代人有過之而無不及,沒有一個時代的人比我們這個時代的人更重視身體.他們從下列這些現象得到如此的結論:沒有一個時代比現在的化妝品及護膚品的利潤更豐厚；沒有一個時代比現在的纖體或瘦身生意更蓬勃；沒有一個時代的保健產品比現在更有市場；在已發展國家,沒有一個年代的醫療開支比現在更龐大,以致無論在美國或香港,醫療改革都是燙手山芋.關注身體是人們的天性,是人們的慣性,對某些人來說,身體帶來的困擾甚至已經成為自己的枷鎖,成為自己的囚牢.

耶穌用「語不驚人誓不休」的對比,提醒門徒:第一,正面來說,既然你們那麼關注身體的狀況,是否應該更關注跟隨耶穌的狀況呢?因為身體有一天會朽壞,在世的生活有一天會總結,但與耶穌的關係卻延續至永恆,而且決定著永恆的歸宿.第二,反面來說,如果對身體的關注令人或令自己離棄耶穌,而離棄耶穌的後果卻比身體受損嚴重得多,在必要的時候,你們是否應該捨輕取重,棄車保帥呢?耶穌的聽眾可能會問,如此的抉擇,是否可能呢?在約伯記中,撒旦也質疑這樣的抉擇的可能性,因此向上帝下戰書,以人對身體的關注向上帝發出挑戰:

「人以皮代皮,情願捨去一切所有的,保全性命.你且伸手傷他的骨頭和他的肉,他必當面棄掉你．．．於是撒旦從耶和華面前退去,擊打約伯,使他從腳掌到頭頂長毒瘡．．．他(約伯)的妻子對他說:你仍然持守你的純正嗎?你棄掉上帝,死了吧!約伯卻對她說:你說話像愚頑的婦人一樣.難道我們從上帝手裡得福,也不受禍嗎?在這一切的事上約伯並不以口犯罪.」(約但記2:4-10)

約伯見證了耶穌的教導,身體的威脅或者身體本身都沒有令他離棄上帝.相反來說,持守對上帝的信任,幫助了約伯超越身體帶給他的困擾.若從這個角度來理解耶穌對門徒的說話,這並不是恐嚇,而是安慰及保證.馬太福音記載耶穌的一段相關的說話:「那殺身體,不能殺靈魂的,不要怕他們；能把身體和靈魂都滅在地裡的,正要怕他.兩個麻雀不是賣一分銀子嗎?若是你們的父不許,一個也不能掉在地上；就是你們的頭髮也都被數過了.所以,不要懼怕,你們比許多麻雀還貴重!」」(太10:28-31)那位連我們的頭髮也數過的天父,豈會不愛惜我們的身體呢?我們豈可以因為身體而離棄天父呢?

3．身體只是媒介　心態主宰起跌

話說回頭,如果我們那麼重視的身體都不容許來令人跌倒,還有甚麼能夠使我們跌倒,或者令別人跌倒呢?在某種意義上來說,身體只是一個媒介,一個執行意念的工具.且看箴言的精警教訓:

16耶和華所恨惡的有六樣,

連他心所憎惡的共有七樣:

17就是高傲的眼,

撒謊的舌,

流無辜人血的手,

18圖謀惡計的心,

飛跑行惡的腳,

19吐謊言的假見證,

並弟兄中布散紛爭的人.(箴言6:16-19)

難道去掉了眼睛就不可以高傲嗎?難道砍掉了手就不可以流無辜人的血?難道砍掉了腳就不可以作惡?有諸內而形諸外,令人跌倒,令自己跌倒的,很可能是隱而不見,無影無形的事物.

如果我們仔細閱讀經文,似乎找到證據支持上述的假設.耶穌在怎樣的處境之下向門徒提出這樣的警告呢?門徒怎樣的狀況,怎樣的表現今耶穌向他們提出「凡使這信我的一個小子跌倒的,倒不如把大磨石拴在這人的頸項上,扔在海裡」這警告呢?

讓我們回到經文,我們就會看見,耶穌說這些對比之前,門徒正自豪地向耶穌報告:「夫子,我們看見一個人奉你的名趕鬼,我們就禁止他,因為他不跟從我們.」(9:38)門徒禁止圈外的人奉耶穌的名趕鬼,會不會在處理信仰的事情上運用了我們今天的「專業主義」的自我保護心態──你不是我們這個專業的註冊會員,所以你不可以從事我們的相關工作?會不會在跟隨耶穌的事情上運用了自古以來「門弟主義」的精英心態──我們才是根正苗紅的入室弟子,你們這些旁末節怎能與我們相提並論呢?

讓我們再向後退一步看看,門徒向耶穌邀功之前,正為著在他們之中誰為大而爭論不休(9章33-34節),導致耶穌介入.爭論誰為大,即是爭論排名次序.為人排名,從來都是吃力不討好的事情,所以現今的刊物在列出人物名稱之後,總會加一句:排名不分先後.雖然寫了這一句,現實上還是無法避免些名字在前,有些名字在後,於是有人嘗試以圈圈的形式列出人名,避開了排名先後的問題.排名先後,有甚麼重要呢?反映甚麼呢?成就?權位?身份?名譽?若果耶穌的門徒當時正在為自己制定一個十二門徒排名榜,會不會暗示他們雖然跟隨這位一無所有的耶穌,但心底裡仍然放不下權位,成就,身份,名譽這些社會認可的價值,正在他們的小圈子之內為著這些眷戀而爭奪呢?

從耶穌提出「跌倒」這話題來討論的背景來看,甚麼使別人跌倒呢?甚麼使自己跌倒呢?會不會是門徒禁止那位「不跟從我們」的人奉耶穌的名趕鬼所反映出的「小圈子心態」,「排外心態」,及「屬靈精英主義心態」?會不會是門徒在排名榜上力爭上游而顯示出對權力,成就,身份,地位的眷戀及爭奪?

如果剛才所指出門徒所展現的心態及價值觀是引致「跌倒」的危險因素,我們如何能夠避開這些陷阱呢?讓我們從耶穌對門徒的回應中找尋指引.

耶穌回應門徒報告他們阻止門徒不認識的人奉耶穌的名趕鬼時,糾正他們的小圈子心態,說:「不要禁止他；因為沒有人奉我名行異能,反倒輕易毀謗我.不敵擋我們的,就是幫助我們的」(9:39-40).如果使人及自己跌倒的危險因素是小圈子心態及排斥別人的行為,避免跌倒的素質豈不是有容乃大的胸襟,海納百川的心懷?

耶穌繼續教導門徒說:「凡因你們是屬基督,給你們一杯水喝的,我實在告訴你們,他不能不得賞賜」(9:41)「給你們一杯水喝」,代表微不足道的事情,代表舉手之勞.對門徒做這樣微不足道的事情,耶穌也記念,並且要賞賜他們,這是何等的細緻,何等的慷慨.如果門徒斤斤計較是導致跌倒的危險因素,耶穌展示的細心及慷慨正是避免跌倒的因素.

回應門徒爭論誰為大的時候,耶穌教導他們說:「凡為我名接待一個像這小孩子的,就是接待我；凡接待我的．．．乃是接待那差我來的」(9:37).如果門徒對身份權位的爭奪是導致跌倒的危險因素,耶穌要求門徒把目光從自己的慾求轉向服務對象的需要,一視同仁地,胸無城府地按著別人的需要去服務人,就有助避免跌倒的危險了.

四．    總結

「倘若你一隻手叫你跌倒,就把它砍下來；你缺了肢體進入永生,強如有兩隻手落到地獄,入那不滅的火裡去.」(43-44)

「倘若你一隻腳叫你跌倒,就把它砍下來；你瘸腿進入永生,強如有兩隻腳被丟在地獄裡.」(45-46)

「倘若你一隻眼叫你跌倒,就去掉它；你只有一隻眼進入上帝的國,強如有兩隻眼被丟在地獄裡.」(47)

如果是我們的心態叫我們跌倒,做甚麼手術可以去掉心態呢?願耶穌與門徒的這段對話存留在我們的心中,常常成為我們的提醒.



\newpage

\section{人算什麼你竟眷顧他}
\label{sec:19}
\textbf{伍渭文牧師 - 人算甚麼,你竟眷顧他?What is Man That You Are Mindful of Him}
\newline
\newline
\hyperref[sec:18]{< < < PREV SERMON < < <}
~
\hyperref[sec:toc]{[返目錄]}
~
\hyperref[sec:20]{> > > NEXT SERMON > > >}


日期:2009年11月1日

主題:人算甚麼,你竟眷顧他?What is Man That You Are Mindful of Him

證道:伍渭文牧師

經文:希伯來書2

人算甚麼?在人類歷史上人曾三次發問這問題「人算甚麼?」因為人對自已的價值產生懷疑.一直以來,人被視為萬物之尺度,地球是宇宙的中心,直到哥伯尼(Copernicus)的天文學革命,認為太陽才是宇宙的中心,人深受打擊,覺得被邊緣化了.像不少大學生在自已的預科學校中,都是獨領風騷,但到了大學,群雄逐鹿,突然產生強烈的自卑感,因為自已不再是課堂的中心了.這是第一次人問:人算甚麼?到了加爾文(Darwin)提出物兢天擇的進化論,認為作為萬物之靈的人,與低級生物同是一源,人是原始生物千萬年的進化結果.相信進化論的人,自然失去人的尊貴感覺,人與猿猴原是近親.這是第二次人問:人算甚麼?人向四周觀看,覺得被邊緣化；人向過去觀看,又是出身卑賤,向裡面看呢?情況更糟.佛洛伊德(Freud)坦白告訴我們,人的內心是黝暗的,是被原始衝動力量所掌控的.人傷心極了,人並非如此良善美好的.人第三次人問:人算甚麼?

但講道的題目並沒有在此結束,『人算甚麼?』後還有一句:你竟顧念他!在新標點和合本中是用感嘆號,是一個讚美驚嘆,我不知道多少人會為著國慶的美好天氣而讚美驚嘆.前天國慶閱兵,我們都被盛大場面所感動,每方隊以同一步伐,同一時間操過司令台,參加的人都經過半年的艱苦練習,保證萬無一失,但有一樣是人無能為力的,就是天氣,萬事俱備了,但天氣如何呢?在國慶前夕,大家都非常關心天文台的天氣報告,幸好天公造美,藍天朗現.

國慶閱兵充份表現出人掌控的能力,但在國慶前夕相繼在菲律賓寶和印尼發生的天災,又提醒我們人的能力與大自然相比,實在有限.國慶前數次天台風"凱薩娜"(Ketsana)造成菲律賓40余年來最嚴重水災,數百人喪生,逾60萬人受到洪水影響.有人拍到片斷,看到汽車像木材一樣,被急湍水流漂走,有汽車被擱在屋頂,司機等待救援.在印尼,全球地殼活動最頻繁的太平洋「火環」,在國慶前夕兩度出現強烈8級地震,觸發高達7.5米的海嘯,上千人死亡,多人被埋在瓦礫下.

水災和地震提醒我們,人的生命在大自然的威力之下,顯得脆弱渺小,人算甚麼?但詩篇第八篇並沒有用自然的威力來震懾人,叫人感到渺小,詩人用天體來說明上帝的偉大,而且以天體按其軌跡運轉無間,說明上主的應許無誤不爽,他的眷顧無微不致.有人稱詩篇第八篇為天文學的詩篇,透過觀察天使體運行,人感到宇宙的浩瀚,人在其中何其渺小,更覺上主眷顧的恩典.我覺得這詩篇非常適合中秋佳節,遙望皎皎明月,高高懸掛,靜靜的銀光,輕輕覆蓋全地,就像上帝的恩典,覆蓋著我們.

人的渺小和上主的看顧,對我們中間一位弟兄來說,這星期最一個真實的經歷.本身是口腔專科醫生並大學醫學院兼任教授的余弟兄,本週發現腦部長了大瘤,令他寫字困難,平時十分鐘寫的報告竟然要一個多小時,記憶消退,開車回家也記不起方向,行路平衡也困難,作為口腔外科醫生,這樣子一定不能工作了,他也明白為何暑假在外地旅行開車時,有兩三次不能專心,差點發生意外,是上帝冥冥中保護了他一家.感謝上主,四個多小時手術成功.弟兄說:我行過死蔭的幽谷.

1.    基督成為肉身是上主的能力

『你因敵人的原故,從嬰孩和吃奶的口中,建立了能力,使仇敵和報仇的閉口無言.』(v.2)

上主的榮耀,不單顯明在諸天的星宿,更顯明於我們所處的大地 and \#8212;他的獨生兒子竟然降世為人,成就救恩.

『你因敵人的原故,從嬰孩和吃奶的口中,建立了能力,使仇敵和報仇的閉口無言.』(v.2) 這節聖經有兩重意思,一是應用在主耶穌身上,二是應用在信徒身上.上主的能力,透過有一嬰孩為我們而生,上主取了人的樣式,死在十字架上,為人贖罪.在十字架上,藉著死敗壞那掌死權的仇敵,就是撤但.女人的後裔要傷蛇的頭,蛇要傷他的腳跟.

我們可以換一個場景:明亮的星宿,嬰孩的降生,在那一個晚上,那應許的星,己劃破伯利恆的夜空.景星高懸,銀光瀉地,祥和像一張被褥,輕輕覆蓋全地,天使歌聲繚亮:在至高之處榮耀歸於上主,在地上平安歸於他所喜悅的人.

『你因敵人的原故,從嬰孩和吃奶的口中,建立了能力,使仇敵和報仇的閉口無言.』(v.2) 這句話己應驗在基督身上.

上主指頭陳設的星宿何其偉大,令人讚嘆,但同樣令人讚嘆的,是當人把行淫時捉拿的婦人,等待耶穌定她的罪,拿起石頭時,耶穌用指頭在地寫字不作聲,最後直起腰來說:誰沒有犯過罪,可以拿起第一磈石頭.人從老到幼一個一個離開,祇剩下那婦人,耶穌說,沒有人定你的罪嗎?我也不定你的罪,去吧,以後不要再犯罪了.陳設星宿的指頭,不是用來定我們的罪,乃用來指引我們歸回父的懷抱.

2. 嬰孩的軟弱顯出上主的能力另一方面,這話又應驗在門徒身上.當耶穌進入耶路撒冷時,小孩子在聖殿喊叫:和撤那歸於大衛的子孫!祭司長和文士就甚惱怒,對他說:這些人(小孩子)所說的你聽見了嗎?

『耶穌說:是的.經上說「你從嬰孩和吃奶的口中完全了讚美」的話,你們沒有念過嗎?』(太21:16) 『這些人(小孩子) 』 是輕藐的語氣,他們不過是小孩子,竟在聖殿大叫大喊,他們懂甚麼?耶穌卻指出這正應驗詩篇第八篇所預言的話.

過去數主日我們的講道的經課都和小孩有關.上主日提到要是我們絆倒小子中一個,令他離開信仰,不如把石磨拴在頸項,拋在海裡.再上一個主日,門徒在往耶路撒冷的路上爭論誰最大,耶穌叫小孩子站在十二使徒當中,並說為首的就要服待這小孩,要為首的就要服待最小的.

詩篇提醒我們,人看來沒有能力的小孩,被人輕藐,但因著他們的的信心,被上帝大大使用.保羅說:並不是我們憑自已可以承擔甚麼事,我們可以承擔的,乃是出於上主.(林後3:4) 我們有這寶貝放在瓦器裡,要表明這莫大的能力,是出於上主,不是出於我.(林後4:7)

在一些禮儀的教會,保留了古教會的一個傳統,在崇拜中安排一些小孩子作禮儀的襄助,我們也邀請一些小孩子,始禮時燃燭和禮終時消滅燭光,目的在提醒我們,來到父面前敬拜,要有小孩子的單純信心,因為要有孩子般單純的信心,才能進天國.自覺像小孩,需要依靠,才能得著上主的能力,被上主使用. 3. 透過受苦嬴取榮耀 vv.5-8詩篇又說到人極為尊榮,掌管一切:『

你叫他比天使微小一點,並賜他榮耀和尊貴為冠冕.你派他管理你手所造的,使萬物,就是一切的羊牛,田野的獸,空中的鳥,海裡的魚,凡經行海道的,-都伏在他的腳下.』(詩8: 5-8)希伯來書的作者,把這段經文應驗在基督身上；因著他的受苦,得了一切的榮耀,使萬物伏在他的腳下.

『惟獨那成為比天使小一點的耶穌(惟獨見耶穌暫時比天使小),因為受死的苦.就得了尊貴,榮耀為冠冕......因受苦難得以完全.』(來2;9, 10) 天使是按上主的形像所造,在上主面前侍立,等待差遣的靈體,甚為尊貴,在人之上.耶穌成為人,暫時比天使,但經過苦難,成為萬主之主.

跟著作者鼓勵在受苦中的信徒,祇管坦然無懼到基督面前,為要得憐恤,蒙恩惠,作隨時的幫助.我們不要在苦難中退縮,甚至放棄信仰,祇管仰望那創始成終的基督,他凡事受過試探,經歷諸般苦痛,明白我們的難處,他是我們的大祭司.苦難使生命得到磨練,行經苦難,留下尊貴的記號.

在著名的足印(footprint) 的故事中,作者回顧過去一生,發現在沙灘上有四個足印,是上主與自已同行的腳印,但有一段路祇有兩隻腳印,就是他經歷最艱難的日子,他有點埋怨上帝在這時候離他而去,讓他獨自上路.但主對他說:孩子,這不是你的腳印,是我的腳印,因你己經無力前行,是我抱起你行留下的腳印.

結語:1969年七月三十日,有史以來全球最多觀眾屏息對著電視機,觀賞兩位太空人Buzz Aldrin and Neil Armstrong首次踏足月球.68年聖誕前夕,太陽神八號環繞月球時,數位太空人因為在太空誦讀創世記,而被無神論者Madalyn Murray O"Hair控告太空總署(最後法院不受理) ,Aldrin受著很大的壓力,要他放棄在太空讀出有關創造的經文.他如何處理呢?他很有智慧說:『休斯頓,這是雄鷹(Eagle) 機長說話.我邀請大家有片刻靜默,Over.我邀請每位聽著,無論你身在何處,你是何人,請現在就剛才數小時眼睛所見到的一切,作出你認為合適的感恩.』跟著他對自己唸出約翰福音15:5 (我是萄葡樹:你們你們是枝子,離了我你就不能作甚麼......),領受聖餐.在回程進入大氣層之際,Aldrin 讀出詩篇第八篇第三,四節: 我觀看你指頭所造的天,並你所陳設的月亮星宿,便說,人算甚麼,你竟願念他!世人算甚麼,你竟眷顧他!



\newpage

\section{你總要付的代價}
\label{sec:20}
\textbf{賀志勇博士 - 你總要付的代價}
\newline
\newline
\hyperref[sec:19]{< < < PREV SERMON < < <}
~
\hyperref[sec:toc]{[返目錄]}
~
\hyperref[sec:21]{> > > NEXT SERMON > > >}


路加福音 14:25-33
\newline
\begin{longtable}{cl}
\hline
\hline
章節 & 經文 (和合本修訂版)\\
\hline
14:25 & \begin{tabularx}{0.7\textwidth}{X} 有一大群人和耶穌同行。他轉過來對他們說： \end{tabularx} \\ \\ \relax
14:26 & \begin{tabularx}{0.7\textwidth}{X} 「無論甚麼人到我這裡來，若不愛我勝過愛自己的父母、妻子、兒女、兄弟、姊妹，甚至自己的性命，就不能作我的門徒。 \end{tabularx} \\ \\ \relax
14:27 & \begin{tabularx}{0.7\textwidth}{X} 凡不背著自己的十字架來跟從我的，也不能作我的門徒。 \end{tabularx} \\ \\ \relax
14:28 & \begin{tabularx}{0.7\textwidth}{X} 你們哪一個要蓋一座樓，不先坐下來計算費用，看能不能蓋成？ \end{tabularx} \\ \\ \relax
14:29 & \begin{tabularx}{0.7\textwidth}{X} 免得安了地基，不能蓋成，看見的人都笑話他，說： \end{tabularx} \\ \\ \relax
14:30 & \begin{tabularx}{0.7\textwidth}{X} 『這個人開了工，卻不能完工。』 \end{tabularx} \\ \\ \relax
14:31 & \begin{tabularx}{0.7\textwidth}{X} 或是一個王出去和別的王打仗，豈不先坐下來酌量，他能不能用一萬兵去抵抗那領二萬兵來攻打他的嗎？ \end{tabularx} \\ \\ \relax
14:32 & \begin{tabularx}{0.7\textwidth}{X} 若是不能，他就趁敵人還遠的時候，派使者去談和平的條件。 \end{tabularx} \\ \\ \relax
14:33 & \begin{tabularx}{0.7\textwidth}{X} 這樣，你們無論甚麼人，若不撇下一切所有的，就不能作我的門徒。」 \end{tabularx} \\ \\
[1ex]
\hline
\hline
\end{longtable}
日期:2009年11月8日

主題:你總要付的代價

證道:賀志勇博士

經文:路加福音14:25-33

1980年的冬季奧運會是體育界的大事,那次奧運會,最讓大家喜歡看的就是冰上曲棍球賽.當時最弱的一個隊就是美國隊,隊員全部是來自美國中部、西部的籍籍無名的大學生,好像是臨時湊起來參加比賽的一樣.沒有一個人看好他們,沒有人認為他們進得了決賽,因為隊員全部都不是職業球員,都是些名不見經傳的大學生.但是你看看其他的隊伍--蘇聯,當時還是蘇聯隊,自1964年以來蘇聯隊就沒有輸過一場球,連續12屆世界錦標賽的冠軍；還有瑞典,芬蘭,捷克,全都是冰上曲棍球的勁旅.因此,當時沒有人看得起這支籍籍無名的美國隊.但不知怎麼搞的,它們一開始就同瑞典打和,然後連勝幾場力克勁旅捷克,在半決賽的時候他們既然打敗了從未曾被打敗過的蘇聯隊,而且在決賽的時候,在0-2落後的情況下,他們卻最終以4-2反敗為勝,戰勝了強大的芬蘭隊,奪得了金牌.整個體壇都震驚了!當時記者就去採訪球隊的教練,問他:你有什麼秘訣,你有什麼秘方,可以把這樣一班籍籍無名的後生仔,打造成世界冠軍?這個教練講了這樣一句話,讓世界難忘,他說"我們沒有任何突出的戰術,沒有任何秘方.但是我們能夠奪得世界冠軍,絕對不是冷門,絕對不是偶然.因為我們有一個成功很重要的因素,我們每一個球員,每一個管理,上至領隊,下至總務後勤撿球的--每一個球員,無論是正選的還是後備的,每一個人都是百分之百投身在球隊中,沒有半條心的.用英文來講就是"total-commitment".沒有半條心的!

其實你想一下,如果我們崇基禮拜堂要組織一個籃球隊,伍牧師做領隊兼教練,姚明做中鋒,Cuby Baile做後裔,喬丹做前鋒,全明星隊伍,都是很有天分的球員.但如果他們都懶懶散散的,喜歡就來練,不喜歡就不來練；天氣好就練多些,天氣不好就不來；遲到,早退,不用心,各有各打的,完全不投身在這個球隊裡.我相信,就算明星再多,恐怕我們這個隊連中大的校隊都打不過.為什麼啊?很重要的一件事,他們完全不投身啊!一班不投身的人,一個不投身的球隊,是永遠不能成為一個好球隊的.縱然你有多麼出色的球員,縱然你有多麼出色的教練,但隊員不投身的時候,這支隊伍就潰不成軍.但很可惜,在今天這個膚淺的社會中,在這個不講投身,不講責任的社會裡頭,投身,或者委身這個字似乎已經從很多人的字典中挪走了.

今天這段經文,就是講到委身的經文,也是令很多牧師都頭疼的經文,因為怕一講這段就嚇走人,就趕客.我以前在一家教會做義工服事.有一次我講道就講這段經文,聚會完了之後,有一個弟兄過來跟我講:"賀博士,我不喜歡你今天這篇講道."我心想,這個人也夠坦白的.別人不喜歡講道,最多在回應紙上不簽名的寫一點意見,像這位仁兄當面來指點的還真是少見.我說:"多謝弟兄啊,我欣賞你的坦白,請問有什麼指教啊?"他就講了三大理由,第一,他講,"你的廣東話不好,講得不夠正,我聽得不順耳"；我都承認我的廣東話馬馬虎虎,所以現在,廣東話講道之前我都要練個好幾次.第二,他講,"我不喜歡你講話太快",這個我都接受.第三,他講,"我不喜歡你講道太挑戰性,我星期一到星期六返工上班都這麼挑戰性,星期天到教堂來想舒服一下,聽牧師講點新聞笑話,讓我吐口氣,可以恢復一下元氣.可是一來到教堂,又被你這麼講,又被你挑戰,哇,我真是周身不舒服,所以我不喜歡你."我回答他說:"第一,弟兄,你講的好,我的廣東話的確要好好練習一下,不要讓弟兄姊妹一邊聽講道一邊猜,賀博士到底在講什麼啊；第二,我會好好控制一下語速的,本來廣東話不好,再講快了,豈不是帶大家坐過山車.第三,對不起,不是我講話太挑戰性,是耶穌基督的話太有挑戰性,我不能不直接講."是的,我們很少有人像耶穌基督這麼坦白,這麼坦白地講出要跟從他的代價.我們不講代價.我們怕講代價,怕講投身.就好像我有個朋友,她媽媽聽說她信主了,就勸她說:"女兒,你信信就好了,不過不要太認真啊!去去教堂就好了,要捐錢就不要了啊."

世上有一種人最可憐,那種人一輩子以為自己信耶穌,到耶穌基督再來的時候說:"耶穌啊耶穌,你記不記得我啊,我不是經常來崇基禮拜堂聚會麼?我不是做過好多事奉麼?你認不認得我啊?"耶穌說:"我從不認得你!"信耶穌,不簡單!信耶穌是一件絕不簡單的事.

究竟什麼是信耶穌?我邀請大家一起來看這一段經文.首先,請大家看25節,25節說"有極多的人和耶穌同行".什麼叫同行啊?同行這個字在英文解釋是"going along with",用香港話講就是"埋堆",用普通話講就是"走群眾路線."再說白一點,希臘文原文的意思,就是一起去看戲的意思.同行,就好像一起去看戲一樣,是完全不需要委身,不需要投身的.聽說耶穌行了好多神跡,名聲大噪了,去看一下又何妨了?去看一下,就跟去UA看戲一樣,有什麼所謂呢?如果表演得好,給多點掌聲給多點錢羅；表演得不好,去過另一場啦,反正UA有這麼多場,從早放到晚.再不行,拉隊走人啦.不需要任何責任,也不用付什麼大代價,天氣好就來,天氣不好就不來.講道講得好我就來,講得不好我就走人,閃!這就叫有極多的人與主同行.或者今天有人也是這樣,來教會就像看戲一樣,跟聖誕節去尖東看倒數一樣,娛樂放鬆嘛,完全不需要委身,完全不需要投身在其中.

但是對不起,耶穌不買這些假賬,耶穌轉過頭來,講出的是這大堆人極度不願意聽的說話,人家覺得極之難聽的說話.以至於我們會批評:"耶穌,有沒有搞錯啊,你這樣講是要嚇走人的,你有沒有學過教會增長術啊?"但是對不起,耶穌很誠實地告訴你,要你好好想想,想清想楚,再來做一個基督徒.在這一段經文中,耶穌一而再,再而三地提醒我們,如果你要做一個基督徒,你究竟要付上什麼樣的代價.信耶穌不是跟風,不是埋堆,信耶穌有你不得不付的代價,那就是委身.什麼叫做委身?在經文中,耶穌基督提到了委身的三層意思:

第一,26節,耶穌說:" 人到我這裡來,若不愛我勝過愛自己的父母,妻子,兒女,弟兄,姐妹,和自己的性命,就不能作我的門徒."這句話翻譯得已經算是很柔和.其實原文的意思是:人若不愛我恨自己的父母,妻子,兒女......."你說:"耶穌,你說什麼啊?"是啊,難道信耶穌就要六親不認,甚至是大義滅親?不是的,耶穌這句話根本不是這個意思,路加福音十章27節,耶穌吩咐我們要"愛鄰舍如同自己",他又怎麼會要我們去恨自己的父母妻子兒女呢?在這裡,耶穌並不是要我們真的去恨我們的父母妻子,耶穌只是想加強一下語氣,讓你的印象更深刻些.他是想強烈地突出兩種情感的對比,兩種不同的愛的對比.

我們為什麼愛自己的父母妻子兒女?因為那是我的父母,我的妻子,我的兒女,不是別人的父母妻子兒女,這種愛我叫做"自我中心的愛",因為一切都是我的我的我的,所以我才去愛.但是神要我們轉變,要我們把這種"自我中心的愛"轉變成"以神為中心的愛".也就是說,父母妻子兒女,我們不要把他們看成是自己的私有財產,想怎麼樣就怎麼樣.我們應該明白,他們是神賜給我們的產業,是神交給我們好好去管理,好好去對待的.對待老人,我們是敬重他們的智慧、閱歷,還有他們辛辛苦苦把我們養大呢,還是把愛老人當做一種無可奈何的責任,或者因為老人家還有油水可以榨取；有錢就有親情,沒錢就連回家吃頓飯都沒時間? 我聽說深水Po 有一個老人家,好不容易把幾個孩子都供到大學畢業,小孩做了律師商人之後,就再也沒去看過他.除了社工以外,十幾年,再沒有一個孩子去探望過他."自我中心的愛",就是這麼容易變質.夫妻之間呢, 是因為她順從我我就愛她,否則就要吵架,就要包二奶呢? 對待孩子也是啊,養兒防老,養兒防老--我們是抱著什麼樣的心態來培養孩子呢?我們要孩子學這個,學那個,給他買這個買那個,是為了孩子呢,還是為了我們自己呢?我們愛孩子,除了關心孩子的事業前途之外,我們還關心他靈性的成長麼,還關心他跟神的關係麼?我只看見父母為孩子選學校,我很少看見父母為孩子選教會!我不是說選名校不好,但當我們自我中心地去愛時,我們並不懂得真正去愛人--愛,不是因為自己喜歡才去愛,是因為愛神我們才去愛人.你說這話很難聽,是的,是很難聽,但耶穌基督很誠實,他誠實地告訴跟隨他的人:"你真的要跟隨我麼?如果你真的愛我,就以我的話語為中心,不要以你自己為中心來愛吧!"

第二,27節,耶穌說"凡不背著自己十字架跟從我的,也不能作我的門徒."做門徒要背個十字架,算了吧,這麼重,掛一個十字架不就好了!我們香港人這麼忙,這麼累,哪有時間來背十字架.什麼是背十字架?背十字架是父神要耶穌基督所做的工,耶穌基督要我們背十字架,就是要我們學習他,放下自己,去做主的工.為了做主的工,我們願意付上時間、能力,哪怕付上生命的代價都不怕--就像耶穌基督為父神所做的一樣.

大家有沒有看過得過奧斯卡獎的影片--《烈火戰車》,它講述英國短跑名將Eric Liddell(中文名叫李愛略)一生的故事.這是一個真實的故事.1924年,第八屆巴黎奧運會,Eric Liddell代表英國出戰,本來他很有實力奪得100米短跑冠軍,可是100米短跑是在星期天舉行.他拒絕出戰,因為他是基督徒,他要做禮拜.當時很多人笑他傻,勸告他,"做禮拜,其實任何一天都可以,或者你跑完了之後去補一下也好啊"--好像做禮拜跟吃補藥一樣.但是Eric Liddell還是放棄了比賽,因為在他心目中,在神的殿中敬拜神是最大的事.神為他開啟了另一扇門,幾天之後,他獲得了400米賽的奧運冠軍.1925年,就在Eric Liddell得到奧運冠軍第二年,他放棄一切榮譽到中國宣教,後來幫助中國人民抗日,最後死在日本人的集中營裡面--但哪怕是在集中營裡面,他仍舊幽默樂觀,像一盞燈一樣,鼓舞集中營裡的人對未來充滿希望.

今天,我們願意為我們的信仰付上什麼代價?今天許多人,不要說為主付上生命,就算就是讓他付上一點時間,他都捨不得.你問他,為什麼不來參加主日崇拜啊?理由一般都是以下幾個:我要做實驗,我要應付考試,我要接待朋友,我要送小孩,總之就是我很忙,忙得沒有時間來敬拜上帝.不好意思,這些理由,有的理由我自己也用過.

背十字架很難麼?弟兄姊妹,其實,當你向你身邊的弟兄姊妹微笑一下,打個招呼；當你每個星期打個電話給一個人,聽他傾訴一下他的憂傷,分享一下他的快樂；或者當你在團契中服事一下其他弟兄姊妹,還有樓下那些可愛的孩子,或者你在廣東話班做組長,你就在做主的工,你就在背十字架.哪怕是再小的事,甚至是你帶自己的孩子出去郊遊,或者給他講一個福音故事,約朋友出去吃個飯,你都可以在做主的工,背他背過的十字架.弟兄姊妹,我們不是星期天的基督徒,不是星期天才來背主的十字架,其實只要留心,我們每一天都可以放下自己,背他的十字架.

第三,33節,耶穌說"你們無論什麼人,若不撇下一切所有的,就不能作我的門徒."這句經文並不難理解,難的是在這句經文之前,耶穌講了兩個比喻:第一個是要人在蓋樓之前計算花費,恐怕蓋不成被人笑話；第二個是打仗的時候,看看打得過打不過,打不過就早點同敵人講和.這是啥意思?難道耶穌說的是,各位,你們好好計算一下,你們跟不跟得了我,如果吃不了這個苦,那就算了,趕快回家去,不要做基督徒啦!耶穌說的要是這個意思,我想跟的人肯定走掉一大半.誰願意吃這個苦!,又要恨自己父母妻子兒女,還要天天背十字架,我車子哪怕是賓士,也擺不下那個十字架啊!

但耶穌說的絕不是這個意思!難道我們跟撒旦打仗,打不過就去求和麼?不是,耶穌基督是想告訴我們,你好好計算一下,無論你怎麼計算,靠你自己,都是完不了工,都是打不贏的.若不是依靠上帝的大能,我們就絕不可能完工,也絕不可能得勝.正如詩篇127篇所說:"若不是耶和華建造房屋,建造的人就枉然勞力,若不是耶和華看守城池,看守的人就枉然警醒."弟兄姊妹,我們跟時間來鬥,我們贏得過時間麼,我們贏得過這個世間的得失成敗,生老病死麼?不能!所以耶穌基督要我們撇下一切跟從他.跟從他,就是不再倚靠自己的錢財,權力,學識去過日子,而是單單倚靠他去生活!

撇下一切,不是要我們賣掉車子賣掉房子,不是,那些都是外在的東西.撇下一切,就是要我們撇下我們自我中心的愛,撇下我們的野心,撇下我們的驕傲,撇下我們的內疚,撇下我們的勞苦愁煩.當我們撇下這一切,我們會發現,我們不但沒有失去什麼,我們還會獲得釋放,獲得從真理而來的自由.

我有個朋友,1949年跟父親逃到香港,因為他父親是國民黨.過了一年,他父親覺得香港日子很艱難,家裡有田有地,想回去又不敢,於是讓他母親帶著他的弟妹先回去看看.那年他5歲.他母親弟妹一去,就再沒有回來.後來才知道,他母親一回家就被公安逮捕了,然後公審批鬥,兩個星期後,他母親懸樑自盡了.自此以後,他父親終身未娶!他父親心裡很內疚,幾十年都沒怎麼笑過.一聽到他說母親,就會大怒,就會打他.所以他不再敢提他母親,他知道,那份內疚一直在他父親心裡.直到幾十年後,他信了主,做了牧師,帶著他父親也信了主,他父親臉上才重新有了笑容,有一天,他回家時,父親平靜地在椅子上睡著了,去世了.

這樣一份內疚,捆綁了他父親很多年,終於,在主耶穌那裡,他父親得到了平靜,得到了釋放,得到內心真正的自由.

弟兄姊妹,耶穌基督要我們做鹽,是要我們委身做他的門徒,是要我們與這個世界不同.不要以為耶穌整天在要求我們,整天給我們講律法--耶穌,好煩啊,我整天日理萬機,你還要求這麼多.然而,當我們放棄自我中心的愛,當我們背起他的十字架,當我們撇下一切時,我們會發現,我們賺的其實要多得多.表面看來,我們付出了很多,失去了很多,但我們最終得到的將是自由,釋放--用一句話說就是:我們失去的是枷鎖,但我們將得到整個世界!



\newpage

\section{你知道所求的是甚麼嗎}
\label{sec:21}
\textbf{林豪恩先生 - 你知道所求的是甚麼嗎?}
\newline
\newline
\hyperref[sec:20]{< < < PREV SERMON < < <}
~
\hyperref[sec:toc]{[返目錄]}
~
\hyperref[sec:22]{> > > NEXT SERMON > > >}


馬可福音 10:35-45
\newline
\begin{longtable}{cl}
\hline
\hline
章節 & 經文 (和合本修訂版)\\
\hline
10:35 & \begin{tabularx}{0.7\textwidth}{X} 西庇太的兒子雅各和約翰進前來，對耶穌說：「老師，我們無論求你甚麼，願你為我們做。」 \end{tabularx} \\ \\ \relax
10:36 & \begin{tabularx}{0.7\textwidth}{X} 耶穌對他們說：「要我為你們做甚麼？」 \end{tabularx} \\ \\ \relax
10:37 & \begin{tabularx}{0.7\textwidth}{X} 他們對他說：「在你的榮耀裡，請賜我們一個坐在你右邊，一個坐在你左邊。」 \end{tabularx} \\ \\ \relax
10:38 & \begin{tabularx}{0.7\textwidth}{X} 耶穌對他們說：「你們不知道所求的是甚麼。我所喝的杯，你們能喝嗎？我所受的洗，你們能受嗎？」 \end{tabularx} \\ \\ \relax
10:39 & \begin{tabularx}{0.7\textwidth}{X} 他們對他說：「我們能。」耶穌對他們說：「我所喝的杯，你們要喝；我所受的洗，你們也要受。 \end{tabularx} \\ \\ \relax
10:40 & \begin{tabularx}{0.7\textwidth}{X} 可是坐在我的左右，不是我可以賜的，而是為誰預備就賜給誰。」 \end{tabularx} \\ \\ \relax
10:41 & \begin{tabularx}{0.7\textwidth}{X} 其餘十個門徒聽見，就對雅各和約翰很生氣。 \end{tabularx} \\ \\ \relax
10:42 & \begin{tabularx}{0.7\textwidth}{X} 耶穌叫了他們來，對他們說：「你們知道，外邦人有君王作主治理他們，有大臣操權管轄他們。 \end{tabularx} \\ \\ \relax
10:43 & \begin{tabularx}{0.7\textwidth}{X} 但是在你們中間，不可這樣。你們中間誰願為大，就要作你們的用人； \end{tabularx} \\ \\ \relax
10:44 & \begin{tabularx}{0.7\textwidth}{X} 在你們中間誰願為首，就要作眾人的僕人。 \end{tabularx} \\ \\ \relax
10:45 & \begin{tabularx}{0.7\textwidth}{X} 因為人子來，並不是要受人的服事，乃是要服事人，並且要捨命作多人的贖價。」 \end{tabularx} \\ \\
[1ex]
\hline
\hline
\end{longtable}
日期:2009年11月15日

題目:你知道所求的是甚麼嗎?

經文:馬可福音10章35-45節

証道:林豪恩先生

一．    引言(網絡笑話)

有個男人非常厭煩每天都要早出晚歸工作,但羨慕老婆可以整天閒在家裡.於是,他向上帝禱告說:「親愛的主,我每天上班,但我的老婆只留在家裡,閒閒沒事做而已,我希望讓她知道我過的是怎麼樣的日子,所以請讓我們的身體交換!」

上帝以祂無比的智慧,成就了這個男人的願望.隔天,這男人一起床就成了個女人.他一起床,就忙著預備早餐,同時再三叫醒那些賴床的孩子們,疏洗,穿衣,執書包,送他們去學校.在回家路上,還要去洗衣店,去銀行,去街市買sung ,回家還要仔細記帳.

記完帳之後,他清了貓的砂盒,後來還替狗洗澡.他趕快餵完貓狗,準備午餐,然後已經到下午一點鐘了,就衝到學校去接小孩.回家之後,看管那些越來越沒大沒小的孩吃午餐,做功課,簡直聲嘶力竭.同時還要洗衣服,擦地板．．．

轉眼間,又是時候要預備晚飯,洗洗切切,煎炒煮炸,他在廚房內團團轉,孩子們在廚房外喧嘩不絕,令他不時要探頭出去喝罵一下.

晚餐後,那個死不要臉的老婆,剛下班一吃完飯,竟然坐在客廳看起足球賽來!

他則忙著整理廚房洗碗,把收好的衣服摺好,哄小孩睡覺,他已經累壞了,但是又擔心白天的雜務還很多沒作完,不管了,他終於衝上床想休息了,但是這時那個閒得發慌的老婆,竟然還興致勃勃提議要「嘿咻嘿咻」!

他心不甘情不願,但還是努力裝得很愉快,設法不要抱怨.隔天他一起床,沿著床邊跪下來,流淚跟上帝說:「主啊!我不知道我之前在想什麼,現在我知道我實在錯得離譜!竟錯到去妒嫉我老婆整天在家!求求您,讓我們換回來吧!」

主以祂無比的智慧,回答他:「孩子,我很高興知道你已學到了智慧,而且,我也很願意把你們兩人換回來,但．．．你還要等二百七十九天,因為昨天晚上你懷孕了.

二．    經文

35　西庇太的兒子雅各,約翰,進前來,對耶穌說:「夫子,我們無論求你甚麼,願你給我們做.」36　耶穌說:「要我給你們做甚麼?」37　他們說:「賜給我們在你的榮耀裡,一個坐在你右邊,一個坐在你左邊.」38　耶穌說:「你們不知道所求的是甚麼.我所喝的杯,你們能喝嗎?我所受的洗,你們能受嗎?」39　他們說:「我們能.」耶穌說:「我所喝的杯,你們也要喝；我所受的洗,你們也要受；40　只是坐在我的左右,不是我可以賜的,乃是為誰預備的,就賜給誰.」

41　那十個門徒聽見,就惱怒雅各,約翰.42　耶穌叫他們來,對他們說:「你們知道,外邦人有尊為君王的,治理他們,有大臣操權管束他們.43　只是在你們中間,不是這樣.你們中間,誰願為大,就必作你們的用人；44　　　　　　　　　　　　　在你們中間,誰願為首,就必作眾人的僕人.45　因為人子來,並不是要受人的服事,乃是要服事人,並且要捨命作多人的贖價.

三．    釋經與應用

1．    門徒不知其所求　耶穌接納並引導經文的上半段記述了耶穌和兩位門徒一段很特別的對話.十二門徒之中的一對兄弟雅各和約翰來到耶穌面前,向耶穌提出請求:「賜給我們在你的榮耀裡,一個坐在你右邊,一個坐在你左邊.」耶穌沒有肯定他們的祈求,也沒有指責他們的祈求.耶穌的回應是:「你們不知道所求的是甚麼.」

從雅各及約翰的角度來說,祈求的內容十分具體,就是當耶穌「坐正」的時候,他們兄弟作左右丞相,分坐第二及第三把「膠椅」,因此,他們當然知道自己所求的是甚麼.

奇怪的是,當事人明明清楚講出自己的要求,但對方卻指出「你不知道自己所求的是甚麼」.我所知的,竟然正是我所不知的,這是一種非常令人困惑的狀況.

在接受輔導訓練的過程中,有一次督導的經驗令我跌入類似的困惑之中.聆聽完一段我的輔導錄音後,老師問:「Walter,你知道自己在做甚麼嗎?」被如此一問,我感到驚愕,也有一點挫敗感.聆聽錄音之前,我不是已經說明了個案,也解釋了我對個案的理解,及所運用的策略了嗎?為甚麼老師仍提出這樣的問題呢?我除了再重覆一次之前的內容,還有甚麼可以澄清呢?老師指出:「你當然知道自己在做甚麼,你所做的,符合輔導理論.不過,問題也就在這裡.你太知道自己想要甚麼,以致你把理論套在受導者身上,錯過了受導者所說的重點,你沒有真正聽得明白受導者的說話.」

我的困惑,原來不是因為我不知道自己所要的是甚麼,反而是因為我太過知道自己所要的是甚麼,到頭來竟然產生「我不知道自己在求甚麼」的矛盾狀況.

帶著這樣的個人經驗,我以諒解的心情去看待雅各及約翰.或者,他們都是太過知道自己想要的是甚麼,阻礙了他們聆聽耶穌所說的是話,以致提出了耶穌認為「你們不知道自己所求的是甚麼」的請求.

在雅各及約翰提出「一個坐在你右邊,一個坐在你左邊」的請求之前,耶穌叫了十二個門徒過來,把自己將要遭遇的事告訴他們:

33　「看哪,我們上耶路撒冷去,人子將要被交給祭司長和文士,他們要定他的罪,交給外邦人.他們要戲弄他,鞭打他,殺害他.過了三天,他要復活.」

雅各和約翰聽到甚麼呢?從他們接著對耶穌的請求推測,他們很可能只聽到耶穌所說的第一句:「看哪,我們上耶路撒冷去」.人的自然傾向,都是聽到自己想聽的,聽不到自己不想聽的.雅各和約翰是三位有幸跟隨耶穌上山會見摩西和以利亞的其中二位,他們親自目睹耶穌在高山上改變形象的莊嚴和輝煌景象:

「(耶穌)衣服放光,極其潔白,地上漂布的,沒有一個能漂得那樣白.忽然,有以利亞同摩西向他們顯現,並且和耶穌說話．．．有一朵雲彩遮蓋他們；也有聲音從雲彩裡出來,說:這是我的愛子,你們要聽他.」(9:2-8)

雅各和約翰之前已經看見耶穌做出能人所不能的神蹟奇事,深知道耶穌是有能力的人；也看見耶穌挑戰宗教及社會權威,和被群眾蜂擁的情況,耶穌是他們心目中的英雄好漢；現在,竟然有幸在高山上看見耶穌改變形象會見舊約律法的代表摩西,及先知的代表以利亞,如此神聖超凡的景象,令雅各和約翰深深相信,他們一定沒有跟錯人,耶穌一定是那位天命所在者,是他們盼望以久復興以色列國的大君王.在這樣的理解及期盼之下,難怪雅各和約翰充滿著勝利主義的心態.有一次,他們路經撒馬利亞一條村莊,那裡的人不接待他們,這對兄弟向耶穌建議:「主啊,你要我們吩咐火從天上降下來燒滅他們嗎?」(路9:54)他們的氣焰比火更旺盛.

現在,耶穌快要進入京城耶路撒冷了.進入京城做甚麼?符合雅各和約翰期待的,進入京城當然是要作王了,怪不得這對兄弟捷足先登,向耶穌請求:「賜給我們在你的榮耀裡,一個坐在你右邊,一個坐在你左邊.」

我們向耶穌祈求甚麼呢?我們知道自己所求的是甚麼嗎?我們會不會好像雅各和約翰一樣,太過知道自己所求的是甚麼,反而令我們聽不到耶穌說的話,以致出現「其實你們不知道自己所求的是甚麼」這樣尷尬的情況.

話說回頭,「不知道自己所求的是甚麼」又如何呢?有甚麼大不了呢?如果我們把焦點集中在「知不知道自己所求的是甚麼」這個議題上,我們的目光仍是放在自己身上,而不是放在耶穌身上.縱使雅各和約翰「不知道自己所求的是甚麼」,耶穌沒有因此而嘲笑他們,沒有因此而輕看他們,也沒有因此而開除他們,甚至沒有因此而責備他們.耶穌知道自己的說話沒有被聆聽,也知道自己被誤解.耶穌對著這對「不聽話」的兄弟,對著這兩位誤解自己的門徒,仍是耐心地與他們對話,慢慢地引導他們,按著他們的狀況教導他們.當我看著耶穌如何對待雅各和約翰如此無知的祈求,我就更放心向耶穌祈求,更放心在耶穌面前顯露我的無知,更放心在耶穌面前面對自己的欲望,更放心承認自己對耶穌的誤解.各位朋友,耶穌也是如此地接納你,你放心來到祂的面前與祂坦誠交往嗎?

2．    門徒尋求獎項　耶穌肯定參與有人為雅各和約翰抱不平,認為他們不至於如此不濟,不同意他們跟隨了耶穌一段日子仍然如此地誤解耶穌.他們指出,雅各和約翰其實是明白耶穌的說話,並沒有期待耶穌進入耶路撒冷作王.他們祈求的是:將來你復活之後在天國作王的時候,賜我們兄弟一個坐在你的右邊,一個坐在你的左邊.然而,無論我們對雅各和約翰的請求視作現世的理解也好,視作來世的理解也好,他們的請求都隱含著對勝利的期望,對獎項的等待.從耶穌與他們的對話可見,如果門徒尋求的是獎項,耶穌肯定的是參與.讓我們再一次細聽他們使徒間的對話:

門徒:「夫子,我們無論求你甚麼,願你給我們做.」耶穌:「要我給你們做甚麼?」門徒:「賜給我們在你的榮耀裡,一個坐在你右邊,一個坐在你左邊.」耶穌:「你們不知道所求的是甚麼.我所喝的杯,你們能喝嗎?我所受的洗,你們能受嗎?」門徒:「我們能.」耶穌:「我所喝的杯,你們也要喝；我所受的洗,你們也要受；只是坐在我的左右,不是我可以賜的,乃是為誰預備的,就賜給誰.」

一般人都認為,坐在王帝的左右,才能顯示自己與王帝的關係,才能分享王帝的榮耀.原來,要與耶穌有關係,分享耶穌的榮耀,不一定要等待坐在耶穌左右的日子,反而是跟隨耶穌的腳步,走耶穌所走的路就是了.喝耶穌所喝的杯,受耶穌所受的洗,即是認同耶穌的使命,經驗耶穌所經過的挑戰.如此做,就能得到耶穌左右的位置嗎?可能得到,也可能得不到.然而,得到或得不到根本不是焦點所在,跟隨耶穌的腳步這個過程,本身就有足夠的意義.興奮,不需要等到領獎的時候；價值及意義,不需要由獎項來決定.

有報章以「遲來的肯定」來描述我們的前校長高錕獲得今屆諾具爾物理學獎.高錕接受訪問時,有這樣一個片段:(高錕太太)黃美芸問高錕,「你是不是光纖之父?」只見高錕一臉茫然,口中唸唸有詞︰「光纖......光纖之父......」.此時此刻,高錕獲頒諾具爾物理學獎,香港及中文大學,可能比高錕自己本人更有反應.高錕校長的興奮,不需要等到獲獎的時候；高錕校長的人生意義及貢獻,也不需要等到獲獎的時候才能肯定.

獲獎與否,不是當事人可以操控,這是眾多條件之下交織而成的結果,何必耿耿於懷呢?或者,到獎項降臨的時候,獎項已經不再重要,你正在享受人生另外一些美好的狀況.

縱然高錕或已忘卻大半生科研心血,但愛妻對自己的好,他記得相當清楚.記者跟高錕說,高太對他照顧得很細心,高錕就不停重複著說:「是,是,她很好的.」夕陽下,兩個瘦弱而纖細的身影,手拖手並肩而行,彷如他們相扶半世紀的寫照.

耶穌說:「坐在我的左右,不是我可以賜的,乃是為誰預備的,就賜給誰.」既然「為誰預備的,就賜給誰」,門徒就不需要為位置而操心,反過來說,操心也無補於事.門徒可以做的,應該做的,就正如耶穌所說:「我所喝的杯,你們也要喝；我所受的洗,你們也要受」,如此,就足夠了.

3．    門徒尋求權位帶來的掌控感　耶穌示範權位的服務性俗語說,講錢失感情；講權位可能更傷感情.耶穌與雅各,約翰兄弟的對話,本來到此結束,可惜,卻被其餘十位門徒聽見,事情就不能就此罷休.他們對於雅各和約翰的所作所為感到憤怒,甚至可能指責他們出賣兄弟.眼見這十二位門徒又起爭執,耶穌又要就地施教了.

42　耶穌叫他們來,對他們說:「你們知道,外邦人有尊為君王的,治理他們,有大臣操權管束他們.43　只是在你們中間,不是這樣.你們中間,誰願為大,就必作你們的用人；44　　　　　　　　　　　　　在你們中間,誰願為首,就必作眾人的僕人.45　因為人子來,並不是要受人的服事,乃是要服事人,並且要捨命作多人的贖價.從耶穌對門徒的教導,我們可以肯定門徒為著雅各和約翰向耶穌要求作左右丞相的事情而爭執,即是說,他們再一次為了「十二門徒排名榜」而有損和氣.「位置」,何以能夠令門徒妄顧兄弟之情,一次又一次地為此而動氣,為此而爭吵呢?因為門徒縱然跟隨了耶穌,仍然擺脫不了外邦人的權位觀念:位置越高就能夠統治越多人,權力越大,就能夠號令更多人.或者,門徒實在羨慕君王大臣擁有的掌控感,「大地在我腳下,國計掌於手中,哪個再敢多說話」的感覺,實在太「過癮」了.

有些人看出了權力架構是令人爭權奪位的罪魁禍首,是造成紛爭及破壞的事物,於是提出剷除權力架構的建議.既然沒有階梯,就沒有權位差異,沒有權位差異,就不會有追權逐位的爭執.英國作家喬治‧奧威爾(George Orwell)的《動物農莊》(Animal Farm)裡的動物,曾經進行過類似的實驗,牠們彼此不再稱為長官,改稱為同志.在教會歷史中,有些人也作過類似的嘗試,認為教階制度於聖經真理不合,必須消除.在教會內,不再有主教或牧師,所有人都是兄弟妹妹,大家都以兄姐互稱.

耶穌對於門徒為權位而爭執的回應,看來沒有那麼極端.他教導門徒說:「你們中間,誰願為大,就必作你們的用人；在你們中間,誰願為首,就必作眾人的僕人」.耶穌沒有因為有人爭權而剷除權力架構,耶穌所做的,是重新定義權位.門徒與耶穌對於「為大」,「為首」的定義有所不同,門徒以管轄或支配別人的廣度及深度來定義權位,耶穌以服務別人的廣度和深度來定義權位.

如果有人問耶穌:為甚麼要以服務,而不是以支配來衡量權位呢?耶穌的回答是:「因為人子來,並不是要受人的服事,乃是要服事人,並且要捨命作多人的贖價.」如果門徒尋求的是權位帶來的掌控感,耶穌卻示範權位的服務性.門徒選擇跟隨與否,不是依靠管理學的驗證,不是通過政治學的辦論,也不是等待社會科學研究的確認；門徒選擇跟隨與否,只能基於這是耶穌自己的榜樣.

各位「為大的」,「為首的」,你們辛苦了.你們如何勞心勞力,你們的服務對象不一定知道,甚至會誤解你們.然而,耶穌基督的眼睛是雪亮的,祂看見你們的任務有多繁重,祂能夠體會你們肩頭所負的有多沉重,因為祂也是以服務別人為己任,甚至捨命作多人的贖價.

四．    總結

「你知道自己所求的是甚麼嗎?」知不知道其實都不是太過重要.

如果有學生祈求上帝賜給他好成績,上帝會教他如何做一個好學生如果有人祈求上帝賜給他一位好女朋友,上帝會教他如何做一位好男朋友如果有人祈求上帝賜給他一位好太太,上帝會教他如何做一位好丈夫如果有人祈求上帝賜給他們好兒女,上帝會教他們如何做好父母如果有人祈求上帝賜他步步高昇,上帝會教他如何服務更多的人如果．．．放心祈求,上帝不會拒絕我們,祂會耐心地引導我們.



\newpage

\section{我要能看見}
\label{sec:22}
\textbf{鄧瑞強博士 - 「我要能看見」}
\newline
\newline
\hyperref[sec:21]{< < < PREV SERMON < < <}
~
\hyperref[sec:toc]{[返目錄]}
~
\hyperref[sec:23]{> > > NEXT SERMON > > >}


馬可福音 10:46-52
\newline
\begin{longtable}{cl}
\hline
\hline
章節 & 經文 (和合本修訂版)\\
\hline
10:46 & \begin{tabularx}{0.7\textwidth}{X} 他們到了耶利哥。耶穌同門徒並許多人離開耶利哥的時候，有一個討飯的盲人，是底買的兒子巴底買，坐在路旁。 \end{tabularx} \\ \\ \relax
10:47 & \begin{tabularx}{0.7\textwidth}{X} 他聽見是拿撒勒的耶穌，就喊了起來，說：「大衛之子耶穌啊，可憐我吧！」 \end{tabularx} \\ \\ \relax
10:48 & \begin{tabularx}{0.7\textwidth}{X} 有許多人責備他，不許他作聲，他卻越發喊著：「大衛之子啊，可憐我吧！」 \end{tabularx} \\ \\ \relax
10:49 & \begin{tabularx}{0.7\textwidth}{X} 耶穌就站住，說：「叫他過來。」他們就叫那盲人，對他說：「放心，起來！他在叫你啦。」 \end{tabularx} \\ \\ \relax
10:50 & \begin{tabularx}{0.7\textwidth}{X} 盲人就丟下衣服，跳起來，走到耶穌那裡。 \end{tabularx} \\ \\ \relax
10:51 & \begin{tabularx}{0.7\textwidth}{X} 耶穌回答他說：「你要我為你做甚麼？」盲人對他說：「拉波尼，我要能看見。」 \end{tabularx} \\ \\ \relax
10:52 & \begin{tabularx}{0.7\textwidth}{X} 耶穌對他說：「你去吧！你的信救了你。」盲人立刻看得見，就在路上跟隨耶穌。 \end{tabularx} \\ \\
[1ex]
\hline
\hline
\end{longtable}
日期:2009年11月22日

主題:「我要能看見」

經文:馬可福音10:46-52

證道:鄧瑞強博士

各位弟兄姊妹,主內平安.

近來香港一個最熱門的話題是:慳電膽.

這慳電膽有點怪,能使人慳了電卻又交多了電費.

當然,用少一點能量而又能看見,多好.

但有人認為,這個慳電膽政策是有人故意看不見某些應看見的東西.

又有人認為,世界這麼大,政策這麼多,什麼也不看而單看這慳電膽政策,是否也是一種盲目.

一個慳電膽,照出了人的盲目.

有些盲目的人,看不到自己的盲目,甚至還為自己沒有盲目而盲目辯護.

有些盲目的人,卻渴望光明.

一個人看見的世界,是否真是一真實的世界,這是希臘哲學早已探討的問題.

柏拉圖指出,很多人其實活在山洞裡,看見牆上的影,便以為是真實世界.他們完全不知自己看見的只是假象.他們一生就在虛假中度過.

這個主題在一齣戲中以高科技重現.在「The Matrix」(香港譯:二十二世紀殺人網絡)這齣戲中,電腦就是現代人的山洞,人活在電腦裡.電腦給的影像,人以為是真的.人完全活在虛擬世界裡.只要安於這虛擬世界,電腦就會供應你的生命養份.這齣戲的主角,稍稍懷疑這個虛擬世界,便被電腦世界排斥,當作廢物般被清除掉,卻被一些已走出虛擬世界的人救回.面對真實世界,人是需要適應的.畢竟,真實世界和虛假世界,相差太遠了.

或許,總有些人多多少少體會到自己看到的是假象.

或許,有些人想看到「真相」.

但是,這是真實世界的真相呢,還是虛假世界提供的真相呢?

我在搜索網站中鍵入「真相」兩字,發現這世界很多人提供很多真相給我們.有「九一一真相」,「六四真相」,「登陸月球真相」,「外星人真相」,「明星是否已婚的真相」,「議員辭退員工的真相」等等,其中一個網站的主題改得好,叫「何謂真相」.

人喜歡真實世界提供的真相呢,還是喜歡虛假世界提供的真相呢?

你想想:若你剛剛結婚,你問朋友你的終生伴侶如何?你渴望什麼答案?你的朋友又會給你什麼答案?你渴望什麼真相?你的朋友會給你什麼真相?

你的朋友不會太介意自己說什麼,反正,面對真相的是你.

美國杜魯門總統講過:「我從沒有給人壞透的東西,但只要我講真話,人們就認為這是壞透的了.」("I never gave anybody hell. I just told the truth and they thought it was hell.")

不願意面對真相,就是一種盲目,一種心靈的盲目.

問題是:心靈盲目的人未必知自己是盲目的,反而認為別人是盲目的.

在路上,偶然會有些橫衝直撞的人.他們會忽然衝過來,你冷不提防就會被撞到.他們可能認為你「不帶眼」行路,他們看不到,「不帶眼」的其實是他們.

知道自己是盲目的,並努力尋求能看見的人,多麼難得.

今日的講道經文,是馬可福音10:46-52.故事裡有一個盲眼的人,他呼求說:「我要能看見」.這是多麼難得的一件事.

要明白這段經文,不能單獨來看,必須將這經文放回馬可福音的整體脈絡裡看.故此,今日我們要多翻幾段馬可福音的經文.

先看馬可福音4:11-12.可 4:11　耶穌對他們說:「神國的奧祕只叫你們知道,若是對外人講,凡事就用比喻,可 4:12　叫他們看是看見,卻不曉得；聽是聽見,卻不明白；恐怕他們回轉過來,就得赦免.」

這段經文是馬可福音第一次嚴肅地處理「看見」的問題.經文的內容有違常理,故教會一般都不會讀這經文.

主耶穌剛剛講了一個撒種的比喻.門徒不太明白,私下問耶穌.　　耶穌就講了上述一番話.簡單而言,主耶穌說,有人些看見神國的奧秘,有些人看不見.看見的就看見,看不見的就讓他看不見.經文說:「恐怕他們(即那些看不見的人)回轉過來,就得赦免.」在我看來,這經文就是要強調:看不見的人,他們用任何方法,都是看不見的心靈的盲目,有一種頑固性,這些人不意會自己是盲目的,也走不出盲目的人生.

門徒理應是「看見」神國奧秘的人,但事實卻差強人意.

我們看第二段

經文:馬可福音6:47-52可 6:47　到了晚上,船在海中,耶穌獨自在岸上；可 6:48　看見門徒因風不順,搖櫓甚苦.夜裡約有四更天,就在海面上走,往他們那裡去,意思要走過他們去.可 6:49　但門徒看見他在海面上走,以為是鬼怪,就喊叫起來；可 6:50　因為他們都看見了他,且甚驚慌.耶穌連忙對他們說:「你們放心!是我,不要怕!」可 6:51　於是到他們那裡,上了船,風就住了；他們心裡十分驚奇.可 6:52　這是因為他們不明白那分餅的事,心裡還是愚頑.

這件事發生在五餅二魚神蹟之後.

門徒在加利利海,耶穌在岸上.因風不順,門徒搖櫓甚苦.耶穌要幫他們,就在海面上走過去.

他們看見耶穌,卻很驚慌.耶穌卻表明身分.耶穌上了船,他們心裡十分驚奇.

若你是門徒,是否會定睛望著耶穌,更深的看看他是誰?

但經文的記載很奇怪,門徒對耶穌行在水面上固然十分驚奇,但他們全部都想著另外一件事,這件事對比起耶穌行在水面上應是無關重要的事.門徒記掛的,竟是五餅二魚事件那些「餅」.

這就像你心儀很久的對象,有一晚和你共赴一個燭光晚餐.在燭光下,理應是你眼望我眼,我眼望你眼的時間,但你卻若有所失,心中想著的,竟是褲袋裡早上還在的一圓錢幣為何現在不見了.

我們應怎樣形容這個人?

聖經形容當時的門徒是:「心裡還是愚頑」.

他們看不到那些真正重要的東西.他們仍活在人生的山洞裡.

原來,不單外人,連這麼接近耶穌的門徒,也看不到耶穌是誰.

另一段

經文:馬可福音8:14-21可 8:14　門徒忘了帶餅；在船上除了一個餅,沒有別的食物.可 8:15　耶穌囑咐他們說:「你們要謹慎,防備法利賽人的酵和希律的酵.」可 8:16　他們彼此議論說:「這是因為我們沒有餅吧.」可 8:17　耶穌看出來,就說:「你們為甚麼因為沒有餅就議論呢?你們還不省悟,還不明白嗎?你們的心還是愚頑嗎?可 8:18　你們有眼睛,看不見嗎?有耳朵,聽不見嗎?也不記得嗎?可 8:19　我擘開那五個餅分給五千人,你們收拾的零碎裝滿了多少籃子呢?」他們說:「十二個.」可 8:20　「又擘開那七個餅分給四千人,你們收拾的零碎裝滿了多少筐子呢?」他們說:「七個.」可 8:21　耶穌說:「你們還不明白嗎?」

從這段經文,我們看到門徒都是相當「愚頑」的.

他們總是想著食物.

想著食物,想著地上人生的種種煩惱,困在人生的山洞裡,就看不見真實的世界了.

接著這段經文,就是一個醫治瞎子的故事:馬可福音8:22-25可 8:22　他們來到伯賽大,有人帶一個瞎子來,求耶穌摸他.可 8:23　耶穌拉著瞎子的手,領他到村外,就吐唾沫在他眼睛上,按手在他身上,問他說:「你看見甚麼了?」可 8:24　他就回頭一看,說:「我看見人了；他們好像樹木,並且行走.」可 8:25　隨後又按手在他眼睛上,他定睛一看,就復了原,樣樣都看得清楚了.

在古代英雄故事中,不能一次過便醫好人是少見的.正如武林超級高手,通常是一揮劍,敵人便倒下的.一揮劍,敵人不死,要走過去,多劈幾刀敵人才死.這算不了超級高手.

但耶穌醫治人的眼瞎,不是一下子的.

經過初步醫治後,這個人能看,但看不清.

要進一步的醫治,進一步的恩典,這個人才能真正看見.

門徒在這階段,是看到一點,但看不到的也多.

接下來的情節,更令人心酸.人真的盲目得緊要.

我們讀下去,就是上個禮拜的講道經文,馬可福音 10:35-37, 41-45

可 10:35　西庇太的兒子雅各,約翰進前來,對耶穌說:「夫子,我們無論求你甚麼,願你給我們做.」可 10:36　耶穌說:「要我給你們做甚麼?」可 10:37　他們說:「賜我們在你的榮耀裡,一個坐在你右邊,一個坐在你左邊.」

可 10:41　那十個門徒聽見,就惱怒雅各,約翰.可 10:42　耶穌叫他們來,對他們說:「你們知道,外邦人有尊為君王的,治理他們,有大臣操權管束他們.可 10:43　只是在你們中間,不是這樣.你們中間,誰願為大,就必作你們的用人；可 10:44　在你們中間,誰願為首,就必做眾人的僕人.可 10:45　因為人子來,並不是要受人的服事,乃是要服事人,並且要捨命作多人的贖價.」

門徒認識耶穌,又不認識耶穌.

他們對耶穌的認識,只是認為耶穌是一具有超能力的人,能幫他們成就地上人人追求的事情.他們看不到耶穌所看的,他們按自己的看法去看耶穌.

他們追求地上的權力和榮耀,他們停留在山洞裡,看不到真實的世界.

人性中有一種黑暗,將我們釘緊在世人眼前的目標裡.

原來,一般的責備,教育,規勸,鼓勵等,是不能叫門徒看見的.

在這裡,馬可福音插入另一個耶穌醫治瞎子的故事.　　馬可福音 10:46-52可 10:46　到了耶利哥；耶穌同門徒並許多人出耶利哥的時候,有一個討飯的瞎子,是底買的兒子巴底買,坐在路旁.可 10:47　他聽見是拿撒勒的耶穌,就喊著說:「大衛的子孫耶穌啊!可憐我吧!」可 10:48　有許多人責備他,不許他做聲.他卻越發大聲喊著說:「大衛的子孫哪,可憐我吧!」可 10:49　耶穌就站住,說:「叫過他來.」他們就叫那瞎子,對他說:「放心,起來!他叫你啦.」可 10:50　瞎子就丟下衣服,跳起來,走到耶穌那裡.可 10:51　耶穌說:「要我為你做甚麼?」瞎子說:「拉波尼(就是夫子),我要能看見.」可 10:52　耶穌說:「你去吧!你的信救了你了.」瞎子立刻看見了,就在路上跟隨耶穌.

這裡,有一個認識到自己盲目的人.他渴望得到醫治.

這個渴望是得不到鼓勵的.聖經說,他身邊的人,責備他,不許他出聲.

但他不願意停留在現在的光景,他不願意在虛假的真相中生活,他要能看見.

耶穌醫好他.聖經說,他成為一個在路上跟隨耶穌的人.

相對於盲目的門徒,他是真正的跟隨者.

在馬可福音裡,人不能憑自己的力量看見,縱使你多麼接近耶穌,仍無法看見真正的耶穌.

馬可福音強調耶穌是一個人不能掌握的謎.沒有他的醫治和恩典,沒有人能明白他,看見他,也不能看見他要我們看見的人生的真相.

這個故事提醒我們,我們多麼需要恩典.不要憑自己的經驗和知識去看,要憑主耶穌的恩典去看.你要突破世人對你人生的控制,他們不許你接近耶穌,你卻要專心,注目耶穌,要獲得他的醫治,要在他手中看到他的世界.

有一個故事是這樣的:　　在美國芝加哥機場裡,一群人趕著上機,而碰跌了一張小台,台上的蘋果散滿一地.一個女孩在哭.這台是她的,她在擺賣蘋果.很多人匆忙而過.有一個人停下來,看著這女孩.發現她是盲的.她在哭.這個人停下來,錯失了他那班機,為這小女孩執拾蘋果,將蘋果放回台上.他發現有些蘋果壞了,他掏錢給小女孩,買下這些蘋果.他要走了,女孩問:「先生,你是耶穌嗎?」

但願榮耀歸於聖父,聖子,聖靈.阿們.



\newpage

\section{你們把石頭挪開}
\label{sec:23}
\textbf{伍渭文牧師 - 你們把石頭挪開Take Away the Stone}
\newline
\newline
\hyperref[sec:22]{< < < PREV SERMON < < <}
~
\hyperref[sec:toc]{[返目錄]}
~
\hyperref[sec:24]{> > > NEXT SERMON > > >}


約翰福音 11:32-44
\newline
\begin{longtable}{cl}
\hline
\hline
章節 & 經文 (和合本修訂版)\\
\hline
11:32 & \begin{tabularx}{0.7\textwidth}{X} 馬利亞到了耶穌那裡，看見他，就俯伏在他腳前，對他說：「主啊，你若早在這裡，我弟弟就不會死了。」 \end{tabularx} \\ \\ \relax
11:33 & \begin{tabularx}{0.7\textwidth}{X} 耶穌看見她哭，並看見與她同來的猶太人也哭，就心裡悲嘆，又甚憂愁， \end{tabularx} \\ \\ \relax
11:34 & \begin{tabularx}{0.7\textwidth}{X} 就說：「你們把他安放在哪裡？」他們對他說：「主啊，請你來看。」 \end{tabularx} \\ \\ \relax
11:35 & \begin{tabularx}{0.7\textwidth}{X} 耶穌哭了。 \end{tabularx} \\ \\ \relax
11:36 & \begin{tabularx}{0.7\textwidth}{X} 猶太人就說：「你看，他多麼愛他！」 \end{tabularx} \\ \\ \relax
11:37 & \begin{tabularx}{0.7\textwidth}{X} 其中有人說：「他既然開了盲人的眼睛，難道不能叫這人不死嗎？」 \end{tabularx} \\ \\ \relax
11:38 & \begin{tabularx}{0.7\textwidth}{X} 耶穌又心裡悲嘆，來到墳墓前。那墳墓是個穴，有一塊石頭擋著。 \end{tabularx} \\ \\ \relax
11:39 & \begin{tabularx}{0.7\textwidth}{X} 耶穌說：「把石頭挪開！」那死者的姐姐馬大對他說：「主啊，他現在必定臭了，因為他已經死了四天了。」 \end{tabularx} \\ \\ \relax
11:40 & \begin{tabularx}{0.7\textwidth}{X} 耶穌對她說：「我不是對你說過，你若信就必看見神的榮耀嗎？」 \end{tabularx} \\ \\ \relax
11:41 & \begin{tabularx}{0.7\textwidth}{X} 於是他們把石頭挪開。耶穌舉目望天，說：「父啊，我感謝你，因為你已經聽了我。 \end{tabularx} \\ \\ \relax
11:42 & \begin{tabularx}{0.7\textwidth}{X} 我知道你常常聽我，但我說這話是為了周圍站著的眾人，要使他們信是你差了我來的。」 \end{tabularx} \\ \\ \relax
11:43 & \begin{tabularx}{0.7\textwidth}{X} 說了這些話，他大聲呼叫說：「拉撒路，出來！」 \end{tabularx} \\ \\ \relax
11:44 & \begin{tabularx}{0.7\textwidth}{X} 那死了的人就出來了，手腳都裹著布，臉上包著頭巾。耶穌對他們說：「解開他，讓他走！」 \end{tabularx} \\ \\
[1ex]
\hline
\hline
\end{longtable}
日期:2009年11月29日

主題:你們把石頭挪開Take Away the Stone

經文:約翰福音 11:32-44

證道:伍渭文牧師

今天十一月一日為教會的諸聖日(All Saints Day),昨晚是萬聖節,究竟諸聖日與萬聖節有何關係呢?是否萬聖節是世俗化了的諸聖日呢?剛相反,萬聖節原是異教的鬼節,鬼佬的盂蘭節,羅馬的教會在主後835 年,將原本於五月中旬舉行的歷代無名聖徒記念日,改往每年十一月一日舉行,並稱這日為「諸聖日」(All Saints" Day),把異教的鬼節用基督信仰聖化(Christianize),移風易俗.

古代西歐塞爾特民族(Celtic)在十一月一日過新年,因為收獲季節剛結束,準備迎接嚴寒冬天.此時冷風吹落葉,草木漸漸枯死,黑夜逐漸拉長,對死亡或鬼靈作怪,自然產生豐富的聯想.塞爾特人認為新年十一月一日的前一晚是一個奧祕的時刻,不屬於過去,也不屬於現今,是屬於另一個世界的時間.這個夜晚,所有鬼魂最容易出入人類世界的一晚.在這晚,塞爾特的年青人就會手持火炬,跑遍領土邊際以示驅逐鬼魂出境,或大聲喊叫,或將堅果丟入火中引起聲響來嚇走鬼靈.後來羅馬人佔領英國後,塞爾特的新年節慶就與羅馬收獲女神Pomona 的慶典結合,產生了以空心南瓜作成的「鬼火燈」(jack-o-lantern)來嚇走惡鬼的風俗.

十一月一日的前一晚就是萬聖日前夕(All Hallows Eve),稱「萬聖夜」(Hallowe"en).Hallowe"en在美國,又演變成扮鬼扮馬,縱情胡鬧的另類兒童節,在萬聖節大人陪同兒童逐戶索取糖果,否則扮鬼臉嚇人,treat or trick.由1965 年開始,聯合國兒童基金會更選定在這段兒童反斗遊玩的日子,進行週年募捐.近年, Hallowe"en 更演變為部分成年人的「化妝嘉年華會」.

異教的萬聖節,原來用來驅逐厲鬼或籠絡遊魂,帶來合境平安.教會的萬聖節,用來思念聖徒的嘉言懿行和信德.我們相信上主是掌管萬有的,無懼幽暗的世界,我們在主裡與離世聖徒連成一體,同屬一個大公教會,每主日我們認信:我信聖而公的教會,我信聖徒相通.在地的聖徒稱作戰的教會(church militant),離世的聖徒稱凱旋的教會(church triumphant).

今天諸聖日的福音書經課,選自約翰福音十一章耶穌使拉撒路復活,這段經文告訴我們:一,上主看聖民之死,極其寶貴；二,死亡不足懼:耶穌說,復活在我,生命也在我.離世的聖徒祇是睡了,將來會復活；三,對上主來說,沒有來得太遲的拯救.

一,上主看聖民之死,極其寶貴

耶穌哭了.今天的福音經課,有一節聖經最短的經節:耶穌哭了 (11:35),耶穌與哀哭的人一起哀哭.在這段經文中,我們看到耶穌滿有情感.伯大尼之家是耶穌常到的溪水旁,那裡有馬大的殷勤服侍,馬利亞細心傾聽,拉撒路的安靜伴隨,馬利亞在這裡用香膏抹耶穌.耶穌愛這家中每一個人,視他們為好朋友.在今天,伯大尼之家代表友善,接待,退隱避靜的好地方.所以當拉撒路病重時,馬大和馬利亞打發人兄見耶穌說:主啊,你所愛的人病了.當耶穌看到兩姊妹和前來安慰她們的人哭泣時,他心裡悲嘆,又甚憂愁,跟著哭了.(11:33) 旁邊的人看到他哭,說:你看他愛這人何等懇切.(11:35)

聖經說:上主看聖民之死,極其寶貴.詩人說:我流淚,你用皮袋裝.耶穌要叫拉撒路復活,不是去完成一項事工,乃去拯救他所愛的子民,抹去他們的眼淚.你要給他起名叫耶穌,因他要把他的百姓,從罪惡拯救出來；耶穌要進到他百姓因罪帶來的的痛苦,背負他們的重擔,耶穌哭了.

因愛而做,和為工作而工作是很不同的.前天晚上崇基學院在嶺南場舉行千人宴,學生會安排了多項表演.有些是邀請外邊專業表演者,有些是自已同學表演.我們發現前者非常專業,很準時出現,表演完畢就離開,可能趕往下一場,表演是一項工作.但作為同學的表演者就很不同,為了在其他同學,師長們面前粉墨登場,他們排演了一個半月.出場時,不少同學都為他們歡呼,未開始表演大家己感到興奮,他們與台下觀眾有一種關係,有一種連結,為了熟悉的同學,老師們,他們渾身解數,落力演出,他們之間有一種情.

當我預備今天的講章時,我想起過去牧會時不少信徒痛苦的臉容.我想起剛出來作傳道,探訪一中年喪妻,父兼母職照顧兩位幼兒,在小學作校役的中年弟兄.他說生活的壓力,令他透不過氣來.另一個場景,就是二十多年前,我趕到醫院一進病房時,一位剛信主新婚不久的年輕姊妹,擁抱著我哭不成聲,她年青丈夫剛剛離開.這些牧養經驗,每次都使我更想作傳道,代表基督表達關心他們的苦痛.

二. 死亡不足懼

a. 耶穌心裡激動.耶穌不單心裡難過,哭了,他還心裡悲嘆,中譯本作心裡激動(v.33: he was deeply moved in spirit and troubled) .耶穌心裡難過,憂傷引至哭泣,我們明白,但耶穌為何心裡激動哩!這字有發怒,憤慨的意思(anger, indignation) .他為何憤慨呢?我們看看三十八節.當耶穌來到拉撒路的墳前,這字再次出現,此時耶穌不再哭了,不再難過了,但面對墳墓,他再度憤慨,發怒.耶穌悲傷,因為看到死亡帶來人的苦痛,死亡奪去拉撒路,死亡使馬大,馬利亞失去兄弟拉撒路而悲痛；他激動憤慨,因為看到死亡對人的宰制.無論是才智聰穎的,愚昧無知的,富有的,貧窮的,都要在死亡面前俯首稱臣.有人嬰孩夭折,有人中年棄世,有人壽終正寢.耶穌來到墳前,就像一個摔角手,直視對手,要擊倒對手,要把對手--死亡--拿下來,所以在墳墓前,耶穌心裡激動憤慨.面對死亡,耶穌怒不可遏,不久他要進耶路撒冷,被釘十字架,藉著死敗壞那掌死權的,要釋放一生因怕死而為奴的人.

在整件事件中,耶穌都採取主動,他沒有按照馬大的要求,立刻趕往她的家醫治拉撒路的病,他仍舊在當地多停留兩天,他等待父的時間,就像約翰福音第一個神蹟--水變酒,耶穌母親馬利亞說酒用盡了,暗示他可以做點工作吧.耶穌回答說:我的時候還沒有到.他不會按照我們有限的知識的要求,有求必應,他等候天父的時間,把死亡消滅.

b. 顯出神兒子的榮耀.使拉撒路復活,是約翰福音中七個神蹟最後的一個,也是最重要的,因為此神蹟最能揭示耶穌進耶路撒冷的目的--被釘十字架--戰勝死亡.所以當馬大和馬利亞派人東告訴耶穌說拉撒路病了時,耶穌說了一些很特別的話:這病不至於死,而是為了上主的榮耀,使上主的兒子因此得到榮耀.(v.4) 為可說顯出神兒子的榮耀呢?在第九章耶穌醫治生來瞎眼的神蹟時祇說:要在他身上彰顯上主的作為(9:3),沒有說到顯出上主兒子的作為.為何這裡提到顯出上帝兒子的榮耀來呢?因為這第七個神蹟與耶穌的死和復活最有關係.

神蹟在約翰福音指記號,路標,神蹟不是使人注意神蹟本身,乃叫人認識基督的身份和工作:『但記這些事,叫你們信耶穌是基督,是上帝的兒子,並且叫你們信了他,就可以因他的名得生命.』(20:30) 約翰福音列舉了七個神蹟或標記,揭示耶穌的身份.第二章水變酒,耶穌是生命源頭.第四章醫好大臣的兒子,顯出耶穌是空間的主,大臣的兒子復元時,剛好是耶穌說:你回去吧,你兒子的病好了,耶穌沒有跟作父親回去就隔空醫好他兒子了.第三個神蹟,耶穌醫好病了三十八年,躺在池旁的癱子,他每次池水動時,都不能第一個下去得醫治,給人捷足先登,耶穌是時間的主(第五章) .第六章五餅二魚的神蹟,揭示耶穌是生命的糧.第六個神蹟,耶穌醫好生來瞎眼的,他是世界的光.最後第七個神蹟--拉撒路復活,說明耶穌勝過死亡,他是復活的主.耶穌對門徒說:拉撒路死了.我沒有在那裡就歡喜.這是為你們的原故,好叫你們相信.(11:14) 若耶穌在拉撒路未死前趕到伯大尼,門徒祇看到醫病的神蹟,他們看不到死人被召喚離開墳墓.耶穌到了伯大尼時,拉撒路己死了四天,猶太人相信人死了三天,相貌開始變化,靈魂認不出自已的的相貌,不能重返自已的身體,死亡就確定下來了.耶穌真的使死人復活,戰勝死亡.

三. 你們把石頭挪走

這段經文提醒我們,上主與我們一起經歷死亡的苦痛,他哭了；他看聖民之死,極其寶貴.但我們不用懼怕死亡,因基督勝過死亡.也許死亡作為生命的結束,對有些人來說,有點遙遠,但死亡的力量,我們每天都經歷.有時我們感到被困在墳墓裡,被墓石封閉,一片幽暗,生活的壓力,像不透氣的墓穴,使我們窒息.這段經文教導我們如何離開墳墓的困境.

a. 對上主有信心,上主的拯不會太遲.馬大說:主阿,己經四天了,他必定發臭了.(v.39) 何必揭開罐頭,讓蟲走出來,我怎可以面對發臭的屍體呢?事情己到了不能挽回的地步,不能逆轉了.兩位姊妹都希望耶穌在拉撒路未死前趕到,醫好他,但死後己四天,解決門題的時間己過去了.我們有沒有同樣的想法?關係已破裂到不可能挽救的地步了,怎能修補?我離開上帝太遠了,我的罪孽太深,深過上帝的恩典,真的嗎?對上帝來說,沒有太遲的拯救,祇有我們願不願意；這是意志問題,也是信心問題.If you would you could.

b. 努力推動石頭.耶穌叫猶太人推開石頭,人可以做的,上帝不會幫助我們,天助自助者.當然耶穌可以使大石移開,像復活節的清晨,但上帝不想我們成為觀賞者,他要我們成為參與者,落場打波.成聖需要功夫,我們抵擋罪還沒有到流血的地步,不運動,不調節生活飲食習慣,怎能減肥,保持身體健康呢?屬靈的生命也如是.推動封墳墓的石頭,是我們的意志,也是我們的責任.我們若願意,上帝會幫助我們.If we would we could.

c. 你們把石頭挪開.在整個故事中,拉撒路復活誠然是耶穌的能力,但上帝也使用一個班愛拉撤路的猶太人,他們從耶路撒冷來,為要安慰喪家 (v. 19), 與馬大馬利亞一起哭泣.他們協助把墓石挪開,他們協助把拉撒路身上的裹屍布解開.我們不曉得他們的名字,聖經沒有明言他們是耶穌的門徒,但他們有門徒的行為--順服耶穌,照耶穌的話去做.耶穌行了這個神蹟,聖經說有許多猶太人信了他(v.45).要挽回信徒,走出黑暗,遠離死亡的力量,使他得著新生,我們作為信徒群體也有責任互相守望.而因為拉撤路復活了,有人離開猶太人,信了耶穌(12:11) ,產生積極的連鎖反應.

在崇基教堂,雖然沒有甚麼緊密的組織--甚麼探訪部,關懷部,但當有弟兄姊妹有困難時,都看到電郵彼此知會代禱,通知教牧同工,互相探訪,我為著這樣自發的彼此關心感謝上主.有人需要找房子,找工作,或其他的需要,讓我們一起協助和關心,這樣我們就能看見上帝在我們中間的工作了.

把墳墓的大石挪開,是上主給我們是應許,上帝的拯救不會太遲,祇要你願意,祇要你努力,祇要我們同作聖徒的,一起彼此守望,石頭一定會挪開,我們可以走出黑暗,迎向光明.把石頭挪開,你願意嗎?耶穌願意,你願意嗎?You could if you would. 你若願意,你就能夠.你若願意,你就能夠.阿們.



\newpage

\section{我行過死蔭的幽谷}
\label{sec:24}
\textbf{胡美儀 - 我行過死蔭的幽谷}
\newline
\newline
\hyperref[sec:23]{< < < PREV SERMON < < <}
~
\hyperref[sec:toc]{[返目錄]}
~
\hyperref[sec:25]{> > > NEXT SERMON > > >}


日期:2009年12月6日

主題:我行過死蔭的幽谷

分享:胡美儀

各位早晨,希望神的靈能夠喚醒你們的內心深處.我感謝主,祂讓我知道自己很渺小,但我很寶貴.謝謝剛才伍牧師,他介紹,今年五月十二日,在紅勘體育館,主行了一個神蹟.在那兒舉行了一個五一二恩雨之聲的胡美儀音樂佈道會.在那兒,我深深感受到自己的渺小,但在那場音樂會中,神親自行了祂的神蹟,祂親自呼召了八千多人,坐滿了整個紅勘體育館.但我心知道,這是主耶穌給我的信心基石,因為當初售票的時候,賣得很少.第一天開票房,只賣得幾十張票,但距離佈道會的日子只有三個星期.我的經理人跟我說:「胡小姐,你以前開的音樂會,每場都爆滿,為什麼開佈道會,則只賣得幾十張票呢?難道你的歌迷全都是拜神的?」恩雨之聲的牧師知道之後,他們說:「如果真的是,就要感謝主.」但亦有人的憂慮在當中.我很感謝一群查經班的姊妹和教會的支持,但我可以告訴大家,我心裡是有憂慮的,但我很想跟大家分享,我們真的要細心聽聖靈在我們心中的教導和提醒.

三月初之時,我跟從恩雨之聲到加拿大幾個不同的地方,舉行音樂佈道會.當時,其實紅勘體育館這個場地還未獲批,我們又不可宣傳,當我在那巡迴佈道會,在加拿大的路上,我們用的是胡美儀生命分享,但是我可以跟大家說.當我站於台上,我辦了十多場音樂佈道會的時候,每一次都受到聖靈的責備,聖靈責備我:「你無須用其他的花巧去吸引人入場,你就是要告訴人,信耶穌不是羞恥的事.」經過很多次在台上,我分享見證,我唱歌,我唱的有詩歌,亦有職業的歌,但我不敢告訴牧師,我三番四次受到聖靈的責備.當紅勘體育館的檔期成功獲批,牧師跟我說:「我們應該用什麼名字呢?金曲分享,還是生命分享?」最後,我告訴牧師,是佈道會.大家很擔心,因為知道佈道會這三個字,很難售票.果然不出大家所料,並且有些朋友問:「為什麼佈道會要收錢?」所以有很多不同的聲音出現,但我感謝主,主耶穌就是要我在這些事件裡,看到祂的信實.

當真是用了音樂佈道會的時候,主耶穌是會負責的,衪會負起這個責任.五月十二日,真是座無虛席,而且看到聖靈運行在場會當中.衪使用台上每一個人,不單只我和毛俊輝.這個佈道會在我心靈上,打了一個信心的基石,我很相信,亦願主憐憫我,在往後我的生命裡,如果遇到任何困難,這個「五一二」,我回望都是信心的基石,亦是我信心的燈塔.

主為什麼會使用我?是因為我渺小,是因為祂告訴我,還需要更大的信心.收到伍牧師今天的講題,是走過死蔭的幽谷.我願主親自去跟大家心靈的深處講說話,我分享的,我心裡一直的祈禱,「聖靈,要分享的說話有很多,你告訴我,我應該說什麼?」但原來聖靈告訴我,如何走過死蔭的幽谷,不單只毛俊輝經歷過,衪說:「胡美儀,你忘記了嗎?你曾經走過一個很深的死蔭幽谷,你一定要跟大家分享.」這並不是我身世故事的各種片段,原來神說:「你要告訴大家,當天你學氣功的時候,如何走過死蔭的幽谷,然後,主耶穌如何拯救你.」

我記得,當我還是電視藝員的時候,每一天拍八組戲.我有一個花名,叫胡八組.一天睡兩小時,無論生理或心理都得不到很好的平衡.當時,有朋友告訴我,練氣功可以強身健體,更可以幫你開運,運氣會更好.於是,我便走到廣州拜師,正因為我做事的執著,所以學氣功也非常用心.起初學氣功,我以前的氣功師傅,打通我的任督二脈.這些大家看武俠小說都會看到,並且,我看到其他在學的徒弟,在氣功的狀態,可以做出很多自己不能做出來的動作.所以我看到他們練功,便心生嚮往,而且很快,我亦加入他們練功的行列.每一次練功,我們都會在觀音像前上香,然後就練功.當初,是師傅給我們氣,但慢慢他會帶我們去收氣.走到很多樹林,去收天然的氣.在練氣功的狀態當中,初期,當我合上眼,腦海中全都是誦經的和尚,非常祥和.我很舒服,並且,當我練功的時候,都有不同的靈體通過我的思想,告訴我他是觀音,他是華陀.師傅教我們的,就是如何去醫治人,並且跟我說:「你努力的學,你所學的東西是你在演藝中,其他人不懂的,你會比他們優勝.」我心很歡喜,但過了一段時間,我一直學氣功,當我合上眼,很多的情景就漸漸改變,很多沒有頭的鬼,男女性交污穢的場面一一湧現.並且,想教我練功的靈體越來越多,漸漸經我的口說出來.「謝謝我帶他到兆慶,謝謝我帶他到某個地方.」漸漸,我感謝天父,他突然叫醒我,我練了功幾年,很用心.但突然有一天,我猛然的醒悟,如果要做醫生,他們要讀書,要經過很長的時間,我怎能靠著氣功就可以去治病呢?我突然想到這些問題,並且開始想,如果我喜愛這些特異功能,即我把自己的靈魂作了賭注.因為當時,當人家睡著了,我看著他,就知道他在發什麼夢,夢中是什麼顏色,當人有特異功能的時候,他會心生歡喜,因為他覺得自己超脫了其他凡人.但那一天,不知為何突然感到心寒,我在想如果我還有幾十年的生命,是不是這些靈體,這些靈界就要跟著我幾十年,那是不是我把自己的靈魂作為賭注呢?我值不值得把自己的靈魂交出來給那些不知是什麼的靈體手上?並且,漸漸很多時會感受到另外有些負面的聲音,甚至乎為何我會醒悟?有一次在逛街,突然有聲音叫我走到馬路被車撞.要我了結我的生命,我覺得很可怕,這全都不是我胡美儀所想的.當我想到靈魂賭注的時候,我放下了,不再練氣功.我這個經歷,後來有長者跟我分享,靈界這東西我們不清楚,最好就不要碰.身體不適,我們去尋求醫生,靈界的醫生,我們最好不要給他們開門,因為當你一開門,你都不知道是什麼靈體來跟你打交道.我感謝主,那種失敗的感覺困擾了我很久,因為我很認真去尋找,我以為那些都是真的.怎知原來差點要了我的命.

之後,有些信了主的藝人跟我宣講福音,心裡想,今次我不會那麼容易中計.不會你說一說我便相信.但他們的愛心實在值得欽佩,他們鍥而不捨,三番四次的邀請.我敷衍的,也要應酬一下.於是在聖誕節,就去了他們的聖誕佈道會,因為我覺得勝誕節都不錯,而且去都不是因為耶穌,只是為了玩.但想不到,當我去的時候,一個演藝人,黃凱欣,她站在台上,分享一些她可以不說的軟弱的經驗.我很好奇,我覺得這個藝人,她無須在台上說這些事給我知道,為什麼她要這樣做呢?原來,她很希望她生命中所得的這份寶貴的禮物,可以跟大家一起分享.我被感動,但一些固執的人不會這麼容易就範.

我去藝人之家,每一次都是安靜的坐在最後我想看看他們怎樣.有一天,牧師講了一篇道,「當中你們那麼多演藝人,每天接觸電視,你有沒有想過在這一刻,好像在坐的各位,由你懂性到今,將你們的人生,製作成錄影帶,現在公開播放給大家看,怎樣呀?」我不敢,即時心裡就浮現了很多我根本不能告訴他人的罪惡.牧師繼續說:「其實每一個人的錄影帶,因為有很多事都不能讓人看,所以便播放,停,播放,停,但是,你認為當你在人前認了罪,你所犯的罪是否就沒有了?」我想想,不是,那些罪仍然在我心中.牧師繼續:「這個宇宙當中,有一位擁有赦罪權柄的真神,你們信主吧﹗你就是一個新造的人,舊時已過,都變成新的了,你可以重新做人.」我心想:「都幾十歲了,怎可能重新做人,無可能重新出生.我一邊聽,心裡不邊反駁.但牧師繼續說:「你認罪,你回家不用告訴人,就好像睡覺,但你心裡慢慢禱告,逐件事向主耶穌認罪.」

於是那天晚上回家,我真的試試這樣做,我當初認罪的時候,真的自覺有罪.但一直說,一真認罪,我發覺有些事件,我都有錯,於是所認的罪越來越多,眼淚不斷流,但那不是痛苦的眼淚,是我前所未曾試過,輕散的眼淚.固執的人當然不會立刻就範.認罪後覺得舒服,就跟主耶穌說:「睡覺醒來再算.」第二天睡醒,感到前所未有的輕鬆,想起昨晚所做的事,我想人與人之間都要守約,何況人與神之間呢?我接受了主耶穌,只在心中,並沒有告知其他人.很多人問我有沒有宗教信仰,我會回答沒有.但每一次我回答沒有的時候,心中感到很不安.好像兒子沒有認爸爸的.但長話短說,神很幽默的.有一次,講完道後,牧師沒有呼召,沒有叫人舉手,但我心裡就有一個很大很大,好像浪子要回頭的心態,我親自舉手,我站起來,說我要信耶穌.你說祂多幽默?

我很相信,神在每一個人身上,都有祂的計劃.我的死蔭幽谷,原來就是我接觸另外的一些靈界,我不知道在座各位,你們,你們有沒有這樣的經歷,或你現在還在經歷中,我可以跟你說,主耶穌救我胡美儀,祂同樣會救你,你快些放下,不要那些東西,又或許你們心靈的深處,有任何東西會防礙你和你家人的關係,你們呼叫主耶穌幫助你,祂一定會幫你.你信,你一定要信﹗你一定可以走過這個死蔭幽谷.

今天,其實我很高興可以跟丈夫一起作見證,他有他另一個死蔭幽谷,交給我的丈夫,毛俊輝.

分享:毛俊輝

大家好,很多人都知道,在二零零二年的時候,我患上了癌症,於胃部有一個惡性腫瘤,很突發的,要立刻切除.因為那個腫瘤的位置易散播,要立刻做手術切除,接著更要立即接受化療,因為是淋巴癌,很容易於全身散播.這個可能性很大,接著便接受了十個月的化療.這件事發生之後,神給我很大的恩典,保留了我的生命,這是很多人都知道的.但今天,在短短的時間內很想跟大家分享的是,神如何真的救贖我,容讓我成為一個基督徒.這是一個恩典,才是我最想跟大家分享的.

其實神很早之前已給我做基督徒的機會,在我的中學時代,已接觸到神的信息,因為我所就讀的是一間教會學校.但當時我拒絕了祂,因為在學校有聖經科,有崇拜,有主日學,可是年少時心高氣傲,校內牧師每次講道都說我們是罪人,都罵我們的,我便很抗拒,不服氣,你這樣說我是不會聽的.而師母每次講道都哭,很悲傷,為我們不停痛哭,我受不了.俗語說:「我;頂唔順!」我不能接受.所以很有趣的,當時有很多機會接近主,但我拒絕了.過了一段時間,便去到美國讀戲劇藝術.當我在美國讀戲劇,更有機會入行,加入戲劇這行業,參與美國戲劇的工作.之後的十多年,並實我遠離了神,我相信你們都明白,在我們戲劇這行業,戲劇工作這環境,是很容易遠離神的.為什麼呢?因為在我們心目中有自己的神,我們就是神.越多機會在戲劇,在藝術方面彰顯自己的榮耀,我們很容易以為自己是 神,所有東西的是自己的,是自己創造出來的.所以我都不知道自己原來這樣遠離神.直至回到香港,香港演藝學院邀請我回來教學,這是一件美好的事.我回來之 後,有一段日子,我很投入我的教學工作,神亦給矛我很多在藝術創造上得到滿足感的機會,但心內是空虛的.我不會告訴他人,但心裡卻是空虛,一種說不出的空虛.就算更多的榮耀,更多的成功,更多在藝術上表演自己,但你是知道欠缺了一些東西.

很難得,有機會認識我的太太,故美儀.她帶領我重新接近神,回到教會認識神的道理,一步一步的重新學習.我說我信主,在這個時刻有機會讓我接觸主,我是順服的.我會認罪的,我知道自己有很多地方需要認罪,去接受神,受洗信主.但話需如此,其實很多心底處的事尚未解決的,神就展現祂的神蹟.第一,在我工作上 最緊張,最重要,最關鍵的時候,我病倒,突然病倒.當時我擔當香港話劇團藝術總監,我正為話劇團預備一個大型的演出時,在演出前兩星期,我回到家中突然暈倒,先嘔血,後暈倒.差不多要完結的時候,幸好我的太太剛回家,立刻送我到醫院急救,就發現原來患上了癌症,那惡性腫瘤發作,是生死之間的.經過手術之後,割了那腫瘤,救回生命,繼續治療.以常情來看,其實就是神,祂已展示了第一個奇蹟,救了我生命,我肉身的生命保留了下來.當然,是非常非常感恩,但神真是了不起,祂要展示的神蹟並不只這些,這只是第一步.接著,祂讓我發覺,我之前所說的空虛,我怎樣,有什麼盼望呢?怎樣可以延續我生命?是不是有了肉身的生命就可以呢?什麼事情都解決呢?絕對不是,祂讓我明白,我的盼望其實是在祂手中,在祂帶領之中如何繼續我的生命.路程開始轉變,那個改變不是過去幾十 年來自以為是我,去安排的.

我記得當我康復,大約一年的時間,城市大學的校長張信光,邀請我出席他們的文化沙龍,希望我路大家分享一下藝術創作上的經驗.當時我要講一個主題,主題是 「放下自我,追求無我.」我其實說得不好,因為我覺得那個概念很矇糊,只是覺得有些事發生了,我只能形容是放下自我,這個我很清楚,幾十年來我很清楚,那 種自我是很強的,亦在創作上是很需要的,有一種我認為是的態度.但什麼是無我呢?其實我都不知道怎樣說,但我知道是有些東西在帶領我走.那個病,讓我重新 找到,引導我往某個方向走.但我不懂得表達,那東西一直放在我心上,是什麼呢?接著那幾年,由二零零二年出事,病了之後到今天,不短的時間,神是非常非常奇妙,一直帶領,一直教導我,不只是尋回生命,是不斷的教導我,打開了我的生命.原來我以前的生命是很狹窄的,是很有限的,我自以為有很寶貴的生命經驗, 在世俗眼中,自以為很了不起,讓我在藝術上有很多的接觸,學習的機會,但原來的很狹窄的.那些東西是並不足夠的,不足以支撐我的生命.我的生命是很不足, 但現在,在今天我站在這裡跟大家說,當我去年放下了香港話劇團藝術總監的工作,放下了,繼續走我生命的道路時,漸漸神就給予我很豐盛的前路.不是要做多 少,不是要做得多好,不是這些.而是,我有很豐富,很滿足的生命,帶著我走.如果我繼續走我的創作路,每一步都在祂的帶領下,是充滿喜悅的.有什麼困難, 都算不上困難,沒有問題的,因為祂帶領著我走,而我有機會接觸一個更大的層面,更大的空間.譬如說,在戲劇上,我現在可以做的戲劇,除了我個人的藝術上的 一些探討會繼續,其實是叫我在戲劇的創作上關心更多人的創作,更多人的前路.我們要做的事,其實可以跟很多人分享,但以前的分享很多時都集中於自己的想法 上,但現在的創作可讓更多人參與,更大的喜悅.如果有件事,是逐步去做的話,我希望是不停可以見證神的恩典,見證神的大能.那種感覺,不需形容為無我,或 什麼,並不重要,是我知道,感受到祂能力是更大的,影響是更大.成功失敗都不是問題,那個重點是,「可以分享」的喜悅是更大的.所以我非常感激神,祂不只 救贖了我的靈,給我的肉身更大的前路,如何行走,於在世時可以做的事,那種滿足感在我一生以前,從未經歷過,就算獲什麼獎,得什麼優惠,得到什麼可以彰顯 自己的機會,什麼得著,都是完全不同的.我不是一個很會說話的人,我只是從心底裡跟大家分享神給我的喜悅,的恩典.祂帶給我生命中一份很寶貴的喜悅,那份 喜悅亦包括我個人,家庭,在婚姻上,神給我的那份愛,妻子給我的那份愛,是我的家,是我意想不到的,但神給我.所以我只想我所做的事能夠彰顥祂的榮耀,能 夠見證祂的大能. 這就是我未來道路上的抱負.我說到這裡,把時間都交給我太太.

其實很感謝主,剛才我丈夫所講的,所分享他走過的死蔭幽谷,其中就是神賜給他的一場大病,這一場病,可能在眾人的眼中,似是一個禍,但原來是祂賜給我們夫婦的福.經過這一場的病,我們知道我們夫婦是相愛,我知道在婚姻生活上,夫婦相愛是很難求的,但只要我們真的信靠我們的神,我們的主,我們求,祂一定會給我們.願主祝福大家,希望大家絕對不要錯過或放過任何一個認識和接受主的機會.

非常感謝毛先生,毛太太,他們信主得救的經歷,給我們莫大的啟示,使我們能夠認識到神的道和旨意.我們再一次以掌聲來感謝他們.



\newpage

\section{石頭記}
\label{sec:25}
\textbf{陳龍斌教授 - 石頭記}
\newline
\newline
\hyperref[sec:24]{< < < PREV SERMON < < <}
~
\hyperref[sec:toc]{[返目錄]}
~
\hyperref[sec:26]{> > > NEXT SERMON > > >}


日期:2009年12月13日

題目:石頭記

證道:

大家早,回到了自已成長,讀書,求學的地方講道,是特別的有感情,聖殿是用石頭所建造的,可是在舊約的聖經裡面,聖殿裡的祭壇像是小鳥建位牠們的家一樣,一個很溫暖的地方,所以就像一個鳥巢一樣,這個使我想起去年我回來的時候,院長在開學禮的時候所說的一個叫"祟基"的地方,就像一個鳥巢的地方,是很多人都可以來棲身的地方,不知道大家每個禮拜天來到這個地方,有沒有感覺是回到神的家裡,像是回到一個鳥巢一樣,主在這個地方去牧養你,今天雖然我講道的題目是"石頭記",可是我們今天不是說[可福音]而是要跟大家說[彼得前書]二章4到8節,如果大家有聖經的話也可以一起翻開來看,我現在要念的是和合本的修定版聖經,可能會跟大家看的經文不太一樣,「親近主,衪是活石,雖然被人丟棄,卻是上帝所揀選,所珍貴的,你們作為活石,要被建造成為屬靈的殿,成為聖潔的祭司,藉著耶穌基督,獻上蒙上帝悅納的屬靈祭物,因為聖經上記著說,看哪,我把一塊石頭放在石安,一塊蒙揀選珍貴的「房角石」,信靠衪的人必不羞愧,所以這石頭在你們信的人,是珍貴的,在那不信的人卻有話說,匠人所棄的石頭,已作了房角的頭塊石頭,又說:『作了絆腳的石頭,使人絆跌的石頭.』」

今天的題目是"石頭記",當我準備這篇講章的時候,有時候也需要培養一下自己的感情,聽一下音樂,曾經聽了一首歌,是由"達明一派"所唱的,歌的名字是"石頭記",如果有人聽過這首歌,那你可能已經洩漏了你自己年齡的祕密,那你的年紀有可能是跟我差不多,相信比較年輕的,應該不懂這首歌的,這我在聽這首歌的時候,我3歲半的女兒問我,這首歌的名字,那我跟我女兒說這首歌的名字是"石頭記",我女兒竟然說,這就是紅樓夢裡面的歌嗎?組樓夢的作者是曹雪芹,有學者這樣子覺得,包括了胡適,魯迅,巴金等等...在他們的眼裡,石頭記在開始的80回,曹雪芹用了石頭去比喻自己的人生,他用了自己生的領悟,都寫在這個石頭記裡面,把整個石頭記都充滿了宗教的含意,可是今天主要不是說紅樓夢,可能你們會有點失望,可是我們文綜系的老師,下學期準備開一個科學,就是紅樓夢,可是原來紅樓夢跟宗教有關係的嗎?那如果你們有興趣的話,也可以去聽聽,會有一位教授講一下紅樓夢跟宗教的關係!

話說回來,我們今天要分享的是彼得前書,可能彼得前書裡的文學沒有像紅樓夢那麼好,可是早在一千七百年前,就是比紅樓夢成書的一千七百多年前,聖經彼得前書的作者己經在用琯頭這個比喻,去描寫我們基督徙的人生,所以今天想要透過紅樓夢去作一個引言,現在我們要說基督版本的石頭記,跟我們之間的關係,第一點跟大家分享的是,在彼得前書第二章四,六跟八節裡,也有提到我們的耶穌基督是我們可信靠的「房角石」,特別在第二章六節裡說,信靠衪的人必不羞愧,不知道弟兄姊妹們,你們的朋友怎麼去評價耶穌基督這一塊石頭呢?我們基督徒又怎麼去評價呢?

聖經裡面說,衪是棄石,衪是可棄絕的,在這邊開始呼應紅樓夢的主角,賈寶玉在紅樓夢的作者的筆下,所寫的就想當時社會的一塊石頭一樣,事實上耶穌基督在衪所身處的時代,也是一塊被放棄的石頭,是因為其他人針對耶穌基督嗎?那我們看要從那個角度去思考,如果耶穌基督當時是要參選健美先生的話,衪肯定會落選,是可以被棄絕的,因為耶穌基督是無佳形美容的,我們現在的男人也要謝謝衪,在猶太人的角度來說,耶穌同樣是可被棄絕的,因為那些猶太人在當時的羅馬社會,受到欺壓的時候,他們一心想要盼望的是什麼?驅除撻擄,把那些羅馬人驅出猶大的地方,剛愎他們的河山,耶穌如果帶領我們,去推翻這個政權,我們會跟隨衪,可是耶穌只有教導我們去愛仇敵,耶穌基督似乎有很多人去愛戴衪,可是衪沒有想要去推翻當時的羅馬政權,所以從一群愛國的份子來看,耶穌基督是可以被棄絕的,同樣,在這個羅馬的社會來看,耶穌是同樣的可以被棄絕的,羅馬的政權管理很多不同民族的人,對於這個羅馬的政權來說,他們最想要的是什麼呢?

社會和諧,民族之間能夠相處和諧,可是當耶穌基督出現以後,很多很多的人去追隨衪,說衪可以把幾千人的肚子填滿,不論你是怎麼去理解這個記載也好,聖經裡面也肯定給我們知道耶穌當時在地上的時候,衪每一之出來,就像現在的明星,歌星一樣,也有很多跟隨衪的人,可是當時的彼拉多,他並不是耶穌的追隨者,看到這個情況,彼拉多開始擔心耶穌會帶來暴亂,對政權來說,突然出現這麼一大群人,圍堵在一些地方的話,永遠都是危險的,所以在一些重要的日子,千萬不要集會,可惜當時沒有一些像是豬流感,要不然也可以要他們不可以集會,可是在聖經裡面告訴我們,耶穌在羅馬人的眼中,同樣他是可以被棄絕的,可是聖經也在另一方面提醒,鼓勵我們要詳細尋求考察,在一章十二節告訴我們,要我們詳細察看衪,我們才看到耶穌是好的,像紅樓夢裡面,如果在當時的社會,''賈''寶玉這塊並不是真的,這塊是假的,可是在這個學家來說,就是看到他與眾不同的個性,一種時代的特徵,可是我們看到耶穌是寶石嗎?是被揀選的,是一個活石嗎?當我們一路認識耶穌那麼多年,在你的生命裡面,耶穌真的是一塊寶石嗎?我本人並不是很懂得去欣賞這塊石頭,所以如果我去像是油麻地或是廟街這,所以我們也需要一些學問及時間去研究去纜賞,耶穌是好還是不好,可是在我們信徒裡面,我們花了多少時間去尋問一個問題,耶穌在我們的生命裡面,衪對我們來說有多寶貴呢?我曾經經歷過一件事情,相信大家也會覺得信仰是需要經過經歷的,在很多年前,我跟我的教會,一起去探訪一家老人院,我們跟那些老人家分享基督教的信仰還有福音,我們也跟那些長者分享我們信靠主耶穌的經歷,當中有一位長者,當我們正準備要跟他分享的時候,他沒有特別的反應,聰明的弟兄姊妹們,你們也應該知道發生什麼事吧,他遇到的問題就像我們以前的高昆校長一樣,原來那位長者患了老人痴呆症,當我們在問他一些問題的時候,他不懂得回答,可是他不停的重複在說''耶穌愛我'',在他的人生走到晚年的時候,當他很多事情都己經記不起來的時候,他記得耶穌愛他,在他的生命裡面,他體會到耶穌就是他的寶石,是他所寶貴的!

可能在我們的人生裡面,有很多值得我們驕傲的,可能在你還在念書的時候,你得了很多的獎狀,可能在你的事業裡,你得到了很大的成就,可能你的家庭也很快樂,有很多值你驕傲的事情,每當你回憶起來的時候都是開心的,可是如果有一天,人生有一個階段,我們的主讓我們把很多事情都給忙掉的時候,如果你的人生只能記得一句話我的時候,你盼望你會記得那一句話呢?在那個時候.你心裡面覺得什麼是最煩惱的呢?或是什麼最寶貴呢?在那個時候,你可能會不知不覺的說了出來,在我們的人生裡面耶穌對我們來說是真的最寶貴嗎?其實這個是很值得我們去尋問,不段認真的去思索的一個問題!

另一方面要跟大家一起去分享的是,在這一段的聖經裡面說,耶穌衪是一塊寶石,衪不單單只是一塊寶石,事實上衪是一塊活石,聖經裡面說,衪是一從死裡面活過來,帶給人有生命,力量的一塊石頭,衪提醒我們,人看衪只不過是一塊可以棄絕的石頭,可是衪卻叫人有生命,在彼得前書的第一章八節裡面也有提醒我們,衪說我們沒有看過衪,我們卻愛衪,我們愛這個耶穌基督,在彼得前書的第二章五節裡面說,我們也要成為活石,不單單耶穌是活石,原來衪也要我們成為活石,這塊活石並不像在祟基所看到的石頭一樣,這塊石頭很奇怪的,是有感情的,這塊石頭也有愛在裡面,是可以經歷愛,也可以經歷百般的試煉.

在我們去思考一下的時候,人生裡面我們可能會經歷很多不一樣的事情,就好像彼得前書的作者對一些讀者說,人生有很多不一樣的經歷,可是你是一塊真正的活石嗎?是可以經歷不同的試煉嗎?在紅樓夢裡面,賈寶玉經歷了整個家族的興亡,寶玉及黛玉的一生,當中也有很多坎坷的事情,但是裡面說,事實上就是因為這些事情,這本小說留給我們很多的空間,但今天,雖然我們經歷了金融海嘯,不知我們在座有多少朋友,弟兄姊妹,這個金融海嘯對你有多少影響,但似乎現在的影響對某些人來說不算太嚴重,但是這一個海嘯,最可悲的事情,不是它帶來金錢的損失,而是問我們,究竟人生的價值是什麼?我們活在一個資本主義的社會中,我們怎樣衡量自己的價值,我們很多時候,都是以市場來衡量自己,究竟我們可以有何成就,可以做些什麼事?我們可達到什麼,來衡量我們的價值?所以我們的長者,活在香港社會,所獲得的尊嚴,如果他不是自己在人生中累積很多財富...相比我在德國的生活,我看到很多長者,他們的生活可說完全是兩個世界,是一件很可惜的事,今天回看,我們如何能夠增值自己呢?我們可能不斷爭扎,如何令自己增值,但上帝容許社會有百般的試鍊,或容許一個人,在人生中有不同的考驗,其實是要問我們,人生中什麼才是重要的事情,像聖經所說,他們所過的,是像客旅,像寄居的生活,人生中原來是經過很多考驗,考驗我們看什麼是最有價值.

早陣子,我去學駕車,當我學習途中,我的師傅恰巧又是一位基督徒,我發現這位師傅他以前並不是教車師傅,想當年他自己有不錯的生意,他的生意不錯的時候,其實他對教會有很多奉獻,每個月奉獻的數目不少,如何反影是不呢?想當年他的奉獻是會讓牧師,執長,堂委也要時常找他聊天,但好景不常,當他的工作不順利時,他要結束整盤生意時,有些時間他很艱難,他甚至未必每一個月奉獻教會,人情冷暖,就在此時不知何故,牧師可能因為教會人數太多,不太有空跟他聊天,但弟兄姊妹,不要取笑這間教會,今天我們用什麼價值來看在座任何一位朋友,來到上帝的殿中,我們去看這是一個數字,還是這是與我無關的人?還是我們去看大家都是主所愛的人?我們不是憑資本主義來衡量,他究竟為我們帶來什麼價值,事實上,聖經提醒我們非常好,我們是活石來的,本來只是石頭...各位朋友,什麼是石頭?其實很簡單,現在我們神學院後面有工程,有很多石頭,我打算用那些石頭來換你的財富,你願意嗎?當然你不會那麼愚蠢...什麼是石頭?石頭是四處也可撿到的東西,不值錢的,在上帝面前,我們的價值在衪眼中,本來就是一無所有,我們不可以迷信,用資本社會主義來衡量自己價值,這並不是毫無意思,我們確是要在社會勤奮,對社會有貢獻,這個制度有它好的地方,但是當人與人之間相處時,我們必須反思,是否用這種社會的價值,來衡量我們的價值,甚至會否把這個價值觀,帶回教會當中,去衡量別人?這是我們需要非常小心的事情,當我們回想時,聖經教導我們,作為活石是需要有感情,是需要有愛,所以聖經彼得前書教導我們,這一個好像冷冰冰的比喻,卻是不斷提醒信徒,要彼此切實相愛,這一個主題貫穿整個彼得前書,在當中提醒我們,我們是否一班有感情的信徒.

正如紅樓夢一樣,其實有人看紅樓夢,只是述說男女間的感情,但在不少紅樓夢學家眼中來看,紅樓夢不只是愛情小說,特別以前八十回來說,主角與不同人物之間的感情,那種人與人之間的關係.今天,當我們的人生走完時,我們的評價如果同基督教愛的價值來看時,究竟我們會獲得什麼價值?在當中,我們是在一個公共領域的教堂,星期天中有很多不同的人,可能會來到我們教會中間,在這個地方聚聚散散,可能跟香港其他很多教會不同,雖然這樣,但建築物並不能反映信徒的生命,我們在這一個地方,是否能珍惜彼此的關係?令我們思想,在這個美麗的崇基校園中...有很多人希望在此結婚,連我在內亦然...而是在這個地方,人在這美麗校園裡,感受及經驗到另一種價值觀,一種生命的體會,這會否成為,我們都歡迎朋友,同事或香港整個社會,來到這個地方,表達我們愛的價值?當然,這視乎我們是否願意踏出多一步,作一個真正的活石,死的可頭是不會行動的,真正的活石是表示他能夠移動,即是說,當有生命來到你面前時,你願意作敞開心靈幫助別人的人?

今年我們的神學院,發覺一個現象,我們收內地學生人數越來越多,我相信在中文大學或是一些高等學府中,很多都會有此現象,但是我們責任是教書,跟其他學課沒有分別的,但我們跟老師之間有個體會,我覺得一個一個的生命,來到這地方讀神學,他們可能在學術思考上有一種挑戰,是新的理解,但當他們一個個生命來到時,我們如何與他有更多的交往?於是我們自發思想,倒不如在星期六的時間為他們安排一個團契,這個雖然不是我們份內事,但我們卻是自發地想去做,我們自己在外國留學過,知道在陌生地方有人支持他,在學術爭扎的過程中能有人與他同行,這是我們非常願意的事情.相同地,有很多人每一個星期都會到訪我們的教會,但我們的心是否準備好,來與這些人交心,分享我們的信仰?

最後一點,在彼得前書第二章第五節,提到基督成為房角石,衪盼望我們基督徒成為活石,衪有一個屬靈的計劃,是希望建造一座靈殿.事實上,如果你熟悉原文,在希伯來文來說,他們稱聖殿是用很大的屋苑等詞語形容,但這詞翻譯成希臘文時,就取了一個概念,是指像房屋或宮殿的詞語,但事實上,如果有一個希伯來背景人士,他會知道,那不只單單是屋苑或宮殿的意思,其實是指一個殿,這一個殿在當時對一般人民心中來說,殿是一個可見的東西,殿對四周的人來說,是一個見證來的,他們在耶路撒冷,會知道耶路撒冷的殿在哪,就算他們在散居的地方,同樣地會知道,什麼是神廟,因為在當時來說,全部均是可見的,所以聖經並不是排斥有一個可見性的殿,正如我們崇基的教會,建立在一個公共領域中,我們希望讓人知道,基督教在香港多元價值中是其中一份子,我們是有貢獻的一份子,當中提醒我們這件事情,但同樣地,聖經在後文提到,什麼是靈殿?這一種可見的殿最重要有一功能,就是一種屬靈的功能,這種功能是能夠帶人出黑暗,入奇妙的光明,是宣揚上主的美德.一個聖殿,在一個公共領域中,成為一個可見的殿,但可見的殿其實不是最吸引人的地方,可見的殿來來去去像德國教堂一樣,可能每個星期其實只是去參觀一下聖殿,真正可留住人的,是這個殿究竟有否這種超凡的特質,是否真的流露出這份美德,去吸引其他生命共同在這地方建設?這很需要我們共同努力,惜日紅樓夢中有他的大觀園,但大觀園可以是在封建社會的象徵,但今日的聖殿並不只表示我們的基督宗教,在一所大學裡面,我們也有我們的位置,一個公共空間,但重要也不是這樣,最重要是我們在這個空間中,可彰顯另類的價值觀,這才是整個殿的建設.

最後,當我們去看這段聖經,他提醒我們,耶穌基督在我們生命中,是否最寶貴的活石?我們必須要詳細考察,花上我們的心思,我們每一個人要很努力地活出耶穌的呼召,我們彼此相愛,在這個社會中一同建立另類的殿.

今天我再說最後的比喻,在很多年前,所謂廣義基督宗教,傳入我們中國時,相傳是我們景教時期...其實當年在唐朝時,很多的宗教都是外來宗教,但事隔多年,在今天中國人的心目中,佛教不再是一個外來宗教,但是基督教仍然是外來的宗教.惜日來說,在整個佛教當中,很多人願意付上很多代價,去為他們的信仰追尋,正如有一位很出名的人,叫唐玄奘,相傳說他在當時在中國薄有名聲,其實他沒必要去印度取經,但他覺得有需要,結果他花了人生很多光陰,在印度帶了經典回中國,去努力地作出翻譯,作了建設.他詳細尋求考察,他信仰的寶貴地方,所以今天,為何中國能寫出一本紅樓夢,其中因素,是因為佛教思想已經植根在中國文化的土壤當中,但今天,我們要問一個問題,有多少人願意為自己的信仲付出追尋?究竟耶穌是否一個寶石?現世代,我們在中國有不同宗教及體系去尋問時,基督教在未來的中國的建設中,會否成為一個文化成分?這要在乎我們在座的每一位弟兄姊妹,每一個對基督信仰有興趣的人,是否認真花上我們的心思?是否耶穌就是珍貴的寶石?但只有知性的追尋是不夠的,還有我們剛才所說,我們需要切實彼此相愛,在這個公共領域的殿中,活出一種另類的價值觀,願神祝福衪自己的說話.



\newpage

\section{十二月的星空}
\label{sec:26}
\textbf{梁元生教授 - 十二月的星空}
\newline
\newline
\hyperref[sec:25]{< < < PREV SERMON < < <}
~
\hyperref[sec:toc]{[返目錄]}
~
\hyperref[sec:27]{> > > NEXT SERMON > > >}


日期:2009年12月20日

題目:十二月的星空

證道:梁元生教授

經文:馬太福音2;路加福音2

各位弟兄姊妹早安!

從這個星期開始,我們就踏進了今年最後的一個月,白天的陽光會慢慢的短一點；黑夜的時間會長一點,睡得好一點；對於天氣,大家大概也感覺到了,會慢慢變得冷一點；外面的樹葉它青綠的顏色也慢慢淡了,有一些葉子可能會變黃；同時這個茂林也變得稀疏一點；我們掛在牆壁上的年曆,也是到了最後的一頁,這個時候,我們都應該知道,聖誕節就快要近了.

對於我們在這個在北半球居住的人來說,十二月是冬天的開始,天氣與樹葉的變化,通常都會帶來一個消逝,蒼涼的感覺,但是另外一方面十二月就知道聖誕節要來了,所以在這個瑟縮,寒冷之中,我們帶著歡樂,希望的期待!進入十二月,面對著快要來臨的聖誕節,我今天就選這個題目關於聖誕的.

我用兩段的經文帶大家一齊思想.一個是馬太福音第二章,另一個是路加福音第二章:

當希律王的時候,耶穌生在猶太的伯利恆．有幾個博士從東方來到耶路撒冷,說,那生下來作猶太人之王的在那裡.我們在東方看見他的星,特來拜他.希律王聽見了,就心裡不安．耶路撒冷合城的人,也都不安.他就召齊了祭司長和民間的文士,問他們說,基督當生在何處.他們回答說,在猶太的伯利恆．因為有先知記著說,『猶大地的伯利恆阿,你在猶大諸城中,並不是最小的．因為將來有一位君王,要從你那裡出來,牧養我以色列民.』當下希律暗暗的召了博士來,細問那星是甚麼時候出現的.就差他們往伯利恆去,說,你們去仔細尋訪那小孩子．尋到了,就來報信,我也好去拜他.他們聽見王的話,就去了．在東方所看見的那星,忽然在他們前頭行,直行到小孩子的地方,就在上頭停住了.他們看見那星,就大大的歡喜.(馬太福音2:1-10)

在伯利恆之野地裡有牧羊的人,夜間按著更次看守羊群.有主的使者站在他們旁邊,主的榮光四面照著他們．牧羊的人就甚懼怕.那天使對他們說,不要懼怕,我報給你們大喜的信息,是關乎萬民的．因今天在大衛的城裡,為你們生了救主,就是主基督.你們要看見一個嬰孩,包著布,臥在馬槽裡,那就是記號了.忽然有一大隊天兵,同那天使讚美　神說, 在至高之處榮耀歸與　神,在地上平安歸與他所喜悅的人.眾天使離開他們升天去了,牧羊的人彼此說,我們往伯利恆去,看看所成的事,就是主所指示我們的.他們急忙去了,就尋見馬利亞和約瑟,又有那嬰孩臥在馬槽裡.既然看見,就把天使論這孩子的話傳開了.(路加福音2:8-17)

從這兩段經文的記載來說,耶穌出生的時候,應該是一個相當混亂的時局,約瑟與馬利亞從外地回到他們的家鄉報名上冊,做戶口的登記,但是因為人很多,旅館沒有地方,所以最後來到一個馬棚,在那邊暫時的棲身,在這樣的一背景之下,我們來思考聖誕的問題.

在這裡出現了多過三組的人,我們說的聖誕三人組.第一組當然是約瑟和馬利亞,後來加上嬰孩耶穌；第二組是從遠方來尋找嬰孩的三博士,或幾個博士；第三組就是在野地裡放羊的牧羊人.同時我們又發現這幾段經文有一個不約而同的地方,就是這些人出現的時間,往往是在一個星空的夜晚.換言之他們是在同一個星空之下出現的幾組人,而星光把他們招聚在一起,至少是星光把博士引領到耶穌,約瑟,馬利亞的面前,又同時是這星光下面出現了天使把牧羊人帶到約瑟,馬利亞的身邊,所以這是星空下面發生的故事──聖誕節.

讓我們先來看博士的經歷:「博士」這字,我翻到新譯本的聖經,把他們翻成「占星人」──占卜星座的人.很多人並不完全同意,所以我還是用回原來的翻譯──「博士」.「博士」的意思是指一些尊貴的,有學問,有學識的人,這裡所說的三博士因為他們帶來了三樣的禮物(我們是從禮物看人數):一個帶來黃金,一個帶乳香,另一個沒藥──所以我們猜想是三個博士,三份禮物,但是否有三位,還是多過三位,我們都不知道.總而言之,這幾位的博士是從遙遠的東方國土而來,為的是要尋找一位新生的王.

他們怎麼會知道有一位新生的王會出生呢?為什麼會遠道而來呢?就是他們仰望天空,觀星的結果.這些人我們猜想到不是以色列地方的人,對於以色列民族的歷史也並不完全的知道.但是我們透過上帝星象的啟示,知道有一位君王要降生了.我們知道他們得到這個衝動是來自上面而來的啟示,有些人說這個可能就是「神喻」了.但另外一方面來說,是他們求學問的結果.一個求學問,學識豐富的博士,他在各方面追求,認識道理,所以他們看到天象的變化,於是產生了追求的心.所以無論是來自「神喻」還是一個自我尋求學問的衝動,我們並不太清楚,但同時我們知道這幾個博士雖然有學問,他們是甚麼背景,來自那一個國家,很多人有很多不同的說法,有些人是做研究的,希望真的把他們的身份研究出來,但是我們很難有一個定論,我們可從下面幾點作一個猜測:

第一,他們有充足的資源──他們從遙遠的東方出發,可能使用駱駝或其他運輸工具,把他們帶過沙漠,帶過森林,他們有充足的資源來到伯利恆的地方,那時候,也帶著貴重的禮物.其實他們來到耶路撒冷,找到了執政的,掌權的統治者希律王,我們知道一個在位,有權力的人不是容易被接見的,所以第二,他們大概也有一個尊貴身份,因此他們見王也很容易,故此猜測他們大概有一點來頭.當我們從他們貴重的禮物來看,他們大概也來自貴族的一個身份,因為有些人也計算過,這幾份禮物值多少錢呢?香港人最喜歡的啦!黃金現是多高?很貴重了,是吧?今天的值錢跟去年的值錢朿不一樣,所以很多人從這些地方來推算這三位來自一個尊貴的背景,大概有貴族的身份.今天我們不想在這裡把他們的禮物逐一地分析,那當然有些人喜歡問:沒藥在那個時候,價值多少?有甚麼用?在這裡只想強調一點,這些尊貴,有學問,來自遠方的人,他們在星光導引之下來到了伯利恆──新生王的地方,是來朝拜的,他們看到那顆星,就產生了一個追尋與叩問的衝動；看著那顆星,他們把學問化成行動.一般有學問的人,大概停留在看見,認識的層面,看見以後,接下來就永不休止的辯論,我們就一直去談,爭辯,最後我們的學問,學識會增加,但是我們不會化成行動.我們看見這幾位博士不同於一般哲人,知識份子的地方,就是他們並不停留在認知的層面,他們把這個認識,知識化成行動.看見了,知道了,就起而行──動身出發,翻山越嶺,過沙漠,越荒原,走上這個追尋與叩問之路.

這條路是信仰之路,他們不只看到天空上面的星,也看到星空之下,他們要實實在在看到星空之下的君王,朝見祂,拜祂,把禮物獻給祂.他們把叩問知識之路變成追尋信仰之路,換而言之,這博士能看見,也能追尋,既有了認知,也付出行動,這點是難能可貴的.至於他們尊貴的背景,貴族的身份,我們可以不必特別的強調,因為星光的導引並不限於貴族階級,星光的導引不限於有學問的階層,星光的導引是可以帶領不同的人,不同的階級,不同身份的人.好像我們從路加福音那裡的經文就可以看到,有另外一種人,也是在星光導引裡面來到朝拜君王的,他們就是牧羊的人.

牧羊人並不需要走很遙遠的路程,他們是本土的人,怹們不是來自尊貴的階層,他們來自這個基層的社會,他們也沒有資源,他們是靠勞動來糊口的,他們沒有貴族,王宮為伴,他們只有羊群為伍.他們在寒夜裡面也需要工作,當他們碰到一個特殊的現象的時候,好像突然有一個大的星光──天使,這些高級,非平凡的現象出現的時候,他們的習慣性是甚麼?懼怕!他們就生懼怕,這就是他們背景,也可能是一種階級的習慣性反應,因為他們是社會最下層的階級.

昨天在電視新聞的特集裡講述香港有一個牧羊人,最初人家送他三隻羊,他很喜歡,一養就是120多頭的羊,所以他就放棄他的工作,養起這批的羊,成為香港唯一的牧羊人.記者訪問他:為甚麼如此呢?他說羊在香港賺不到錢,只能每個月把幾頭羊賣給穆斯林人或飯店,數目不多,但是他很喜歡這工作,所以就帶著一根長杆,每天帶羊走在公路上,讓牠們在路邊吃草.

牧羊人是比較貧窮的,工作也比較辛苦的,在社會上面是一個邊緣的民族,也可以說這些牧羊人是勞動人民,居無定所的一批人,他們代表社會上面失去勢力的,沒有發言權的人,雖然不一定是失敗的,但一定是比較失落的一群.然而星光還是照他們,引導他們.星光照耀有財富,有學問的博士,同樣也照著這些目不識丁的,沒有資源,貧窮的牧羊人.

我們如果把馬太,路加福音的經文搬到舞台上面做為舞台的場景的話,那我們最主要看到黑夜,在黑夜之中,突然之間有兩個強烈的燈光照射在台上:一個照著幾個博士,另外一個燈光照著的是牧羊的人,同樣在一個舞台上面,他們都是平分秋色,不會因為他們沒有財富,就不受燈光的照射,這意思是表明後面的主宰對他們的看待是一樣的,星光照耀沒有分貴賤,貧富,階級,它引領有學問的人,也引領一些沒有受教育,在社會上面失落,失群的人一齊來走上信仰之路,這是叩問與追尋.

接下來,我想跟大家想想:星光下面也不是完全平等,也不一定引向聖誕的歡樂與平安,因為後面還有一些比較慘烈的故事.

十二月的星空下面,聖嬰誕生,博士來朝,牧人歡慶,天使唱歌... 唸這些當然是歡慶快樂的圖畫,很多人對聖誕節的期望也大概就是這樣,縱有不喜歡聽到天使報喜的歌聲,但誰人不喜歡過著快樂,和平的生活?當然這個也是上帝與天使宣告的重要信息.

十二月星光帶給我們歡喜的一個信息,因為一個不平凡的生命要來到這個世間,我們來慶祝祂的誕生,我們來歌頌祂的誕生,博士們喜悅地獻上他們的厚禮,牧羊人歡喜地到處傳揚這個新生王的誕生的消息,但是在相同的星空之下,也有一些人懷另外一種心情,另外一種籌算,他們仰望星空,但星光的引照對他們來說好像利刀,他們眼睛裡面看到的是一片的危險,在他們眼光裡面,透視出來的是一股殺氣,他們是誰?他們是以希律王為首的一夥人.

在同一個星空下面,他們不斷叩問與追尋,跟博士,牧羊人一樣,他們去追尋新生嬰孩的下落,希律王對博士說:「你們去仔細尋訪那小孩子．尋到了,就來報信,我也好去拜他.」我們當然知道同樣是一個叩問,追尋,但出自內心完全不同的動機,內心存著的是奸詐,藏著的是殺機,幸好博士在夢中得到啟示繞道回家,沒有去向希律王報告作奸人的工具.不過雖然他們繞道了,但是殺戮不會因為他們不遵守希律王的話就停止,殺戮是不能避免的.

十二月星光下,隨著這個嬰孩的誕生,牧羊人的喜悅... 一場腥風血雨正在展開,希律王不見博士的回報,就下令設查全地新生嬰孩的下落,如果找不到這個新生的嬰孩,就下令把伯利恆城一帶所有兩歲以下的嬰孩全部殺光,這個事情記載在馬太福音2:16-18節:

希律見自己被博士愚弄,就大大發怒,差人將伯利恆城裡,並四境所有的男孩,照著他向博士仔細查問的時候,凡兩歲以裡的,都殺盡了.這就應了先知耶利米的話,說, 『在拉瑪聽見號咷大哭的聲音,是拉結哭他兒女,不肯受安慰,因為他們都不在了.』 (太2:16-18)

黑夜來到,星斗滿空,星光照耀著那個出生,生命的來到,也同樣照耀著死亡,照耀著那個慘忍的殺戮.

聖誕到了,耶穌來了,為世界上帶來和平與歡樂,但耶穌的誕生不能棄掉死亡.一個嬰孩的出生,反而帶來成千上萬人(我不知道數目有多少,但至少有上千的數目)的死亡,這是人生一個很大的吊詭,也是神學上面常常讓我很難去理解的疑團.但是我們知道,信仰不是一條不用思索的道路,在信的人,當然有新的生命,但是不能免於死亡,以致死亡成為對自己,家人,朋友及其他人的威脅；一個不信的人,也要面對這個死亡的威脅.星空之下,有生,有死,有歡樂的歌聲,也有逃亡的痛苦,在那個夜闌人靜的晚上,約瑟帶著妻子馬利亞,嬰孩耶穌要逃離伯利恆,逃離希律兵丁的追捕,逃過刀劍,連夜奔波逃到南方的埃及,很有可能那個時候他們是著天空上面的星光,才找到方向,穿越軍防來到南方埃及的地方.

我們抬頭看到十二月的星空,感到一片的寧靜,安祥,但是應該會想到,在同一個星空底下,有許多人要流亡,要逃亡,有很多人被追捕,追殺.星光引照,有歌聲,有禮物,有出生的生命,也有逃難的經驗,有嬰孩的血流出與生命的死亡,同一星空之下,是生與死,善與惡,角力跟共存的同一世界.

上面所說的星光之下,有生有死,有樂有哀,有笑有淚,有恩慈也有慘酷,有尊貴也有貧賤,星光跟太陽一樣,要照好人,也照歹人,照智慧人,也照愚昧人,同樣聖誕是慶生,慶祝一個生命的來臨,但是卻引來殺戮,我們是否不應該過聖誕節呢?我們是否不應該慶祝這個聖誕呢?我們有信仰的人跟沒有信仰的人是否就沒有分別呢?在星光之下,走曠野的行旅,有許多時候會陷入迷惑,信仰跟人生實在上有太多的疑團,問號還沒有解決.

上兩個週會我們請來演講的嘉賓,這嘉賓是龍應台女士,她先我們的同學發了一篇文章,這文章裡面提到,人生就像一個迷宮.她用她自身的經驗──她在歐洲走過一個綠色的樹林迷宮,叫人在裡面團團轉,費盡了自己的力量也走不出來.人生,社會好龍應台教授所說的,是一個叫人迷惑,叫人徬徨的迷宮,既然是一個思想上面的迷宮,也可能是一個權力的迷宮,甚至說是一個生命的迷宮.我們被社會上面不黑不白,無是無非,價值翻來覆去的事情,弄到六神無主,很難把握,人生就是一個迷宮,有人在這個迷宮裡面,為了權力,享受或是人與一些東西,都不願意走出去,他情願在迷宮裡面.但有一些極力地尋求走出這個迷宮,但是往往找不到途徑.前面一種人是沉緬,後面一種人是無奈.龍應台大概屬於後面一種人,很希望能走出去!對她來說,她發現走出迷宮的方法──她認為哲學可能指點人家離開迷宮,走出迷宮的一條出路!甚麼是哲學?對她來說,就是仰觀天象,抬頭望天空.因為星斗可以指示你的方向,也認識星斗的位置,就可知道你的方向,這樣就有辦法走出人生迷宮.

我們或可以從龍教授的解釋得到人生一些的解釋／啟示.望觀星斗,走出迷宮,但是星象跟斗數不是從研習的哲學而來,但可能是從宗教信仰而來,好像博士的仰光星斗,牧羊人遙望夜星,他們都看到那個星,他們就追求這星,走出他們人生的迷宮.

最後我們看到博士,牧羊人,希律黨以外的人,在聖誕節就是這個場景出現的最後一批重要的人,就是約瑟,馬利亞,耶穌.在星光之下,他們也跟天使,博士,牧羊人相遇,他們都是有第一手聖誕經驗的人,我們可以在士身上學習到接受上帝的啟示:第一次追求君王的,第二次是離開伯利恆的,都是受到從神而來的啟示.從牧羊人身上我們可以學化懼怕成興奮,把人生改變,以前是看到甚麼事情都害怕的,現在可以很高興,坦誠,勇敢的跟其它人說:「現在我給你報這大喜的信息」.因為看到主把人生的勞變成出樂歌.

我們從馬利亞的身上又學到甚麼呢?馬利亞表現出來是無言無語,在整個場景裡面,最安安靜的一群人就是馬利亞,約瑟和耶穌.馬利亞表現出來的完全是一種內心的安寧與安祥.信仰不是一個激情的表現或是到處張揚,好像馬利亞一樣,聖經有一句話說:「她看到了,就存在心裡反覆思想」,也可能像在後面的一位老先生西面的身上一樣,看到了主耶穌,他說:「如今可以照你的話,釋放僕人安然去世．因為我的眼睛已經看見你的救恩.」(路2:29-30)一個長的人生,在作一個句號的時候,可以說我再沒有遺憾了.

在星空之下,在人生之中,我們應該向馬利亞跟西面一樣,追求是與主相遇,知道祂以後,得到內心的安祥與寧靜,那就是信仰的真實與信仰的真諦.



\newpage

\section{新生王的記號}
\label{sec:27}
\textbf{伍渭文牧師 - 新生王的記號The Marks of the New Born King}
\newline
\newline
\hyperref[sec:26]{< < < PREV SERMON < < <}
~
\hyperref[sec:toc]{[返目錄]}
~
\hyperref[sec:28]{> > > NEXT SERMON > > >}


路加福音 2:12
\newline
\begin{longtable}{cl}
\hline
\hline
章節 & 經文 (和合本修訂版)\\
\hline
2:12 & \begin{tabularx}{0.7\textwidth}{X} 你們要看見一個嬰孩，包著布，臥在馬槽裡，那就是給你們的記號。」 \end{tabularx} \\ \\
[1ex]
\hline
\hline
\end{longtable}
日期:2009年12月27日

主題:新生王的記號The Marks of the New Born King

經文:路加福音 2:12

證道:伍渭文牧師

``你們要看見一個嬰孩,包著布,臥在馬槽裡,那就是記號了."

剛過去的星期四是中文大學畢業禮,相信大家在傳媒的報導己認識到中大的畢業禮情況了,示威抗議差不多成為中大畢業禮的記號.我喜歡參加大學畢業禮,並不是因為喜歡看到示威場面,乃希望與畢業生和家長一起慶祝這大日子,學生等了三年,但家長等了二十二年了.除了希望分享那種歡樂氣氛外,我特別喜歡畢業禮的榮耀感.畢業禮當天,家長都穿得很隆重得體,畢業同學和師長們都穿上不同顏色的學士袍.醫學院紅色,文學院黃色,理學院紫色.......各老師穿起博士袍更是好看,因為不同大學有自己不同的顏色,所有中大哲學博士畢業生袍都是紫色,這是中大的顏色.而大學校長和數位副校長,更有特別的禮袍 and \#8212;中間紫色鑲金邊,坐在前列.進場時,主禮嘉賓,大學高層,講座教授,教授,副教授,助理教授依次而人；禮成離場時次序相反.加上會場的尊貴布置,莊嚴的奏樂,廁身其中,有一種榮耀感覺.

我們看到畢業禮中有很多記號,把一些重要的人物分別出來,他們穿的服飾,坐的位置,進場的次序都是記號,突顯他們的地位.那麼新生王來到世界上,上以甚麼作為記號呢?這些記號又表明些甚麼呢?讓我們看看聖經路加福音二章十二節:"你們要看見一個嬰孩,包著布,臥在馬槽裡,那就是記號了."

第一個記號是馬槽.包著布,臥在馬槽裡是一個甚麼的圖像呢?相對客店的房間,馬槽樸素簡約,沒有奢華花巧.我們現正在哥本哈根召開全球氣候峰會,關心地球暖化的危機,這圖像可以表示環保節約,反樸歸真,很有時代感,可以被視為前衛時尚,且有點浪漫.

但耶穌降生的馬槽本身並不浪漫,是晚上安放牛羊牲口的地方,其實是一個山洞,晚上牲口避過風霜雨露的地方.今天我們到聖地旅行,主誕堂就是在一個山洞內.約瑟,馬利亞最初也不是選擇馬槽的,他們來到客店時,已住滿人了,也沒有人讓出客房,祇好住在馬槽.

耶穌誕生的馬槽,以前有沒有人住過呢?耶穌藉童貞女懷孕,騎沒有人騎過的驢駒子進入耶路撒泠,葬在門徒財主亞利馬太人約瑟的新墳墓,聖經特別強調耶穌徵用從來沒有人用過,突出分別歸上主為聖.不過,這些被上主所用的都有共通點 and \#8212;甘心樂意,馬利說:我心尊主為大,我靈以我救主為樂,而客店的主人,並沒有為馬利亞的臨盆作出特別的安排.約瑟有否把聖靈感孕的事告訴店主,要求他特別安排地方?這點聖經沒有說.但無論如何,對將臨盆的母親,任何人也應給予特別照顧,他甚至可以讓出自己的房間,協助馬利亞平安誕下聖嬰,可惜他沒有這樣做,其中一個可能的原因,馬槽雖然不適合住人,但因為有人需要住而給人住過.在今日社會,我們有沒有人因為需要住,而住不適合人住的地方呢?我想大家都想起籠屋了.在中國大陸,也有人因為需要住,而住不適合人住的地方,很多離鄉別井在城市打工的農民工,就住在沒有窗戶的地下層,因為付不起昂貴的租金.馬槽這記號,提醒我們主認同被邊緣化的弱勢族群,我們也應關心他們的福祉.

他生的時候在最卑微的馬槽,死的時候降到最低的陰間,像我們每次誦唸使徒信經時所說的.他誠然為我們成了貧窮,使我們因他成為富足,就像婚姻中歡悅的對換.交換戒指時新郎先說:我以此戒指為徵記,與你結婚,我以我的身心尊敬你,我以我的財物,與你分享；女方也說同樣的話:我以我的財物,與你分享.基督把我們的罪,歸在他身上,他成為罪人,而把他的義歸在我們身上.基督為我們成為貧窮,使我們因他成為富足.

第二個記號是包著嬰孩的布.這是極其普通的一塊布,每一位作母親的都會用一塊布 and \#8212;襁褓包裹自己的嬰孩.相對於馬槽,這塊布亳不起眼,沒有令人驚喜的地方.嬰孩包著布,司空見慣,但嬰孩臥馬槽則令人驚訝,有點浪漫,但沒有這塊布,也不會影響整個聖誕故事.

這塊布使我想起耶穌平靜風浪的神蹟中那一個枕頭.(馬可四35-41) ``忽然起了狂風,波浪不斷的打進船來,艙裡積滿了水.耶穌卻在船尾靠著枕頭睡著了......." 整個故事沒有了那枕頭,也沒有損這神蹟的完整性,誰人會關心注意那枕頭呢?我們都會注意風浪大作,耶穌運用權柄止住風浪:" 不要作聲!安靜吧!" 但人睡覺都要枕著枕頭,孩子生下來就被一塊布所包裹；枕頭與襁褓雖不起眼,但上主的眼目沒有離開它們.

也許在主日回到教堂敬拜時,我們特別感到上主很親近,但每天上班擠地鐵和巴士,趕文件,趕功課,做家務的例常工作,就感到沉悶,覺得上主離我們很遠了.這塊布提醒我們,祇要我們為著愛人和榮耀上主而作,平凡的事亦能顯揚上主的香氣,也能成為王的記號.母親用襁褓包裹嬰兒,目的使他得到暖和,感到安全,顯示母親的愛；這塊布弄污了,就清潔它,破了就補它,直到嬰孩長大,這就是一樁了不起的工程,因為嬰孩就是這樣成長過來的.用愛心為人做一件平凡,例常的事,就成為王的記號,顯揚基督的香氣.

第三個記號是嬰孩.王怎麼會是嬰孩呢?嬰孩最多是王儲,直到嬰孩長大經過加冕登基典禮才能成為王.我十一月在英國倫敦開會時有機會到西敏寺大教堂參觀,看過英王加冕登基的王座(坐椅),是木造的,不過教堂一度沒落,這登基用的椅給不認識的學生用刀劃花了.參觀過程中,也看到很多下葬該教堂的王墓,其中不少死於宮廷鬥爭.為王不易,作王儲更加危險.在中國,有人隔代便指定成為王儲,成為王儲時年紀也不少了,但還要戰戰兢兢,循規蹈矩二十多年才能真正掌王權.若是嬰兒成為王儲,總有攝政的人,他才是有實權的攝政王.

嬰孩表示上主為我們成為軟弱的救主.我說成為軟弱(vulnerable),不是說他沒有能力,乃指他把能力收藏起來,用愛來親近我們.嬰孩是可親近的,很少人不喜歡嬰孩,我們很想抱抱他,亦可隨時放下他.同樣,我們可以拒絕耶穌,漠視他,輕忽他給我們的感動,這是他的軟弱(vulnerability).然而,他總不放棄我們,他像那位佇候在村口的慈父,等待浪子歸家,還預備了上好的牛犢,預備慶祝浪子回歸,將殘的燈火,他不消滅,壓傷的蘆葦,他不折斷.他為我們的原故,如羊被牽到宰殺之地,默默無聲.新生王是嬰孩,因為他來救贖,不是強勢統治.馬槽的嬰孩,就是十架上的基督.基督取代了人的樣式,為要藉著死,敗壞那掌成權的,要釋放一生因怕死而為奴的人.

不少人希望新生王的出現,革新制度,帶來盼望.今年捷克天鵝絨革命二十週年,前總統哈維爾最近接受亞洲週刊專訪,文章提到二十年過去了,人還是懷念哈維爾,因為現在捷克街頭的賭場充斥,且被黑社會用來洗黑錢,在經濟繁華的背後,捷克成為歐洲貪污重災區.另一篇外國報紙文章報導,伊朗最近大選過後,很多反對黨和異見份子被捕,令很多伊斯蘭國家失望,他們最初對推翻前親西方的巴列維王朝後成立的現政權懷有憧憬,伊斯蘭世界終於可以對美國說不了,但獲取權力後,一樣迫害異已,同前政權沒有分別,他們曾經對新生王懷有希望,但最後希望幻滅.耶穌來並沒有應許我們一個眼睛看得見的政權,他說他的政權就在我們心中,神的國就在我們心中,這國以嬰孩的軟臨到世界,我們可以拒絕這權柄,輕忽這權柄,新生王是嬰孩,上主為我們成為軟弱(vulnerable) 嬰孩,但被上主的愛感動的人,這軟弱就成為改變生命的大能.

最近在美國報刊看到一篇談論歐洲宗教動態的文章,提及瑞士公投禁止建造穆斯林宣禮塔(minaret)；歐洲人權法院 (European Court of Human Rights) 宣判,義大利在學校課室懸掛有基督在其上的十字架(crucifix)是違反1950年歐洲人權宣言的(European Convention on Human Rights),但義大利政府辯稱十字架是文化象徵,並非宗教.我想無論瑞士人民或和義大利政府,都不能忍受基督這位新生王是嬰孩,他們需要一位強者作王.瑞士人民懼怕伊斯蘭信仰在他們中間坐大,影響他們原有的基督教信仰,雖然他們他們並不認真對待基督教信仰.義大利的俗世政府,正如作者所觀察,也沒有真正尊重基督信仰,十架的基督,己淪為文化符號,安放在課堂,成為文化符號.事實上,教堂在歐洲是博物館多於敬拜上帝的殿堂,主日崇拜人數很少.有一次在瑞士主日參加教會崇拜,發現很多人,很興旺,但後來才發現該主日是堅信禮,親戚朋友都來祝賀,但下主席就不回教會了.很多人都想脫離教會,目的不想交教會稅.掛在課堂的十架基督,給義大利政府和人民有一種基督信仰臨在的感覺,儘管對真正的信仰已沒有興趣了.

嬰孩是新生王的記號,提醒我們,信仰並不是君臨天下,信仰不能由外力促成,我記起中學剛信時,很想我的同學也信主,於是邀請他出席佈道會,到講員呼召決志時,我見他還未舉手,我便把他的手舉起來.信仰最珍貴是內在的自由(inner freedom) ,出自內心.耶穌說:上帝的國在那裡?上帝的國就在你們心中.那裡多不多宣理塔,並不影響上帝的國；課室掛不掛十字架,並不影響上帝的國,上帝的國就在你們心中.讓那位新生王耶穌,今天生在我們心中,那麼我們無論做甚麼事,看來很平凡的事,像包裹嬰孩那塊布,都能成為王的記號,而不是掛在課堂成為記號,也不怕被人拿走.祇要基督在我們心中成形,就成為活的宣理塔,宣揚基督的真理.

馬槽作為記號,顯明上主認同卑微的人,站在被邊緣化,被拒絕的弱勢族群一邊,我們也應關心他們,行公義,好憐憫,與上主同行.

那塊布作為記號,表示上主眼目不會離開我們,在我們認為平凡的,例常的生計活動中,總有上主的記號,祇要我們留意,總看到上主的痕跡.

新生王是一位軟弱的嬰孩,但在信的人,這軟弱卻是上帝的大能.

``你們要看見一個嬰孩,包著布,臥在馬槽裡,那就是記號了."



\newpage

\section{我心尊主為大}
\label{sec:28}
\textbf{鄧瑞強博士 - 「我心尊主為大」}
\newline
\newline
\hyperref[sec:27]{< < < PREV SERMON < < <}
~
\hyperref[sec:toc]{[返目錄]}
~
\hyperref[sec:29]{> > > NEXT SERMON > > >}


路加福音 1:46-53
\newline
\begin{longtable}{cl}
\hline
\hline
章節 & 經文 (和合本修訂版)\\
\hline
1:46 & \begin{tabularx}{0.7\textwidth}{X} 馬利亞說：「我心尊主為大； \end{tabularx} \\ \\ \relax
1:47 & \begin{tabularx}{0.7\textwidth}{X} 我靈以神我的救主為樂； \end{tabularx} \\ \\ \relax
1:48 & \begin{tabularx}{0.7\textwidth}{X} 因為他顧念他使女的卑微；從今以後，萬代要稱我有福。 \end{tabularx} \\ \\ \relax
1:49 & \begin{tabularx}{0.7\textwidth}{X} 因為那有權能的為我做了大事；他的名是聖的。 \end{tabularx} \\ \\ \relax
1:50 & \begin{tabularx}{0.7\textwidth}{X} 他憐憫敬畏他的人，直到世世代代。 \end{tabularx} \\ \\ \relax
1:51 & \begin{tabularx}{0.7\textwidth}{X} 他用膀臂施展大能；他趕散心裡妄想的狂傲人。 \end{tabularx} \\ \\ \relax
1:52 & \begin{tabularx}{0.7\textwidth}{X} 他叫有權柄的失位，叫卑賤的升高。 \end{tabularx} \\ \\ \relax
1:53 & \begin{tabularx}{0.7\textwidth}{X} 他叫飢餓的飽餐美食，叫富足的空手回去。 \end{tabularx} \\ \\
[1ex]
\hline
\hline
\end{longtable}
日期:2010年1月3日

主題:「我心尊主為大」

經文:路加福音1:46-53

證道:鄧瑞強博士

各位弟兄姊妹,主內平安.

沒有耶穌的降生,這世界就像冰凍的荒原,永遠見不到太陽一樣.

耶穌的降生,叫所有活著的人,明白如何活,才能使這世界,不致淪落為一個悲慘世界.

今日的講道經文是:路加福音1:46－53

這裡有一個女人馬利亞.

在1:48,這女人自稱為「卑微的使女」,一個低下的奴婢.

這段經文記述一個低下女人在基督裡的故事.

或許將以上這句話改成這樣說,基督在一低下女人的生命裡而產生的故事,因為這時,基督的確在這女人所懷的胎裡.

基督改寫了她的人生.

沒有基督,貧窮人的故事是淒慘的.

看過一則這樣的新聞報導:

這是一個女孩子的故事.她叫楊元元,原籍湖北宜昌.六歲時喪父,她的母親含辛茹苦地養育她.不負母親期望,她以優異成績考入武漢大學,畢業後,到上海讀碩士.這時,她的母親的工廠破產,她的母親失去了住的地方.這女孩子唯有帶母親到上海,因無錢在校外租屋,故只好將母親偷偷安放在學生宿舍裡.她多次向學校申請,讓母親暫住在學生宿舍裡,因無錢在外面租房子.但校方的回應是:「沒錢不應來讀書.」校方又罵其母親為「鄉下人」,禁止她再踏入學生宿舍.這位母親就坐在學校禮堂前在寒風中過夜.這女孩子深感愧對母親,多日無睡,最後精神崩潰,以吊頸方式結束了生命.很多人指,這是上海樓巿發瘋所引起的又一宗悲劇.有人概嘆:「窮人沒資格活在這偉大的時代.」(參考:\url{http://www.chinapress.com.my/content_new.asp?dt=2009-12-12 and amp};sec=world and amp;art=1212wf31.txt)

這是一個什麼偉大的時代,偉大得叫窮人不能好好活下去.

在網上看到另一個故事:

「有錢的老闆到「洗腳店」讓人替他洗腳.而所謂洗腳竟然是:一個盲眼的女孩子用舌頭去舔他的腳,並用牙齒輕咬後跟作為一種按摩.」

這則故事的網上評論是:

「這樣的羞辱已經到了變態的程度,但更讓人深思的是背後冷酷的邏輯:沒有這樣的服務形式,這些雙目失明的女孩就可能連生活都維持不了.」

(參考:\url{http://big5.ce.cn/gate/big5/blog.ce.cn/html/24/108524-204665.html})

這是什麼邏輯?

貧窮的人,在大都巿裡,在大時代裡,的確很難生活下去.這裡的生活需要錢.而這裡的錢,很多時候是需要用先前的錢去賺後來的錢的.從來沒有錢的人,很難賺錢,很難生活.能用錢去賺錢的工作,才是值錢的.勞力的工作,賺的錢也只能維持最基本的生活.

飽食足,然後知榮辱.很多吃不飽的人,只能放下尊嚴去討生活.

問題是:已經吃飽的人,為什麼竟要別人沒有尊嚴,才給予施捨?

前幾天,在旺角街頭,見有一乞丐伏在地上,不斷將頭撞向地面,以叩響頭的方式向人行乞.他以接近自殘的方式,叩問人的良心.但,為什麼他要這樣做?他自己必然是服膺這種想法,也相信別人也有這種想法,這想法是:我放下了尊嚴,活得不像人,那麼,那些有尊嚴的人便會施捨給我了.他這樣做,這樣想,是因為這的確反映了很多人的想法.

前幾天,在尖沙咀的隧道裡,見有一位女士,唱著中國民歌.你可以向她點唱,付一點錢,她便唱出那首你喜歡的歌.她有尊嚴地唱,開心的唱.你開心的聽,不用以施捨的方式將錢給她.她相信,人不需失去尊嚴,也會得到別人的幫助.她相信,總有人能以互相尊重的方式去彼此欣賞和幫助.這位女士,其實相信的,正是「聖誕節」的信息.

我們回到路加福音書的經文.

福音書裡的這位卑微的奴婢馬利亞,在基督裡,她看見人與人相處的新邏輯.

路 1:51　他用膀臂施展大能；那狂傲的人正心裡妄想就被他趕散了.路 1:52　他叫有權柄的失位,叫卑賤的升高；路 1:53　叫飢餓的得飽美食,叫富足的空手回去.

基督的偉大,在於他沒有因為有錢,叫貧窮的盲眼女孩子去舔他的腳.他沒有因為自己有權,叫別人像狗一樣在他面前討食.他沒有因為自己有勢,叫別人毫無尊嚴地阿諛奉承.

他叫卑賤的人抬起頭,有尊嚴地唱歌.他叫貧窮的人,能有尊嚴地吃一餐飽飯.在基督裡,人不需以尊嚴為代價去乞取生活.在基督裡,沒有人需要放下尊嚴去博取人的同情.在基督裡,那種有意無意奪取人的尊嚴而抬舉自己的威勢的人和社會,要被神審判.

曾經讀過一篇文章,說到內地某城巿的公安公告,提示巿民,切勿同情露宿街頭者.這公告提示:「在車站碼頭,高架橋下或偏僻路段,以及城郊結合部地區,露宿著一些拖兒帶女的外來者.他們衣衫襤褸,髒亂不堪,好逸惡勞,靠拾荒,乞討混日子,嚴重影響了大城市的文明形象,給這個城市治安帶來了很多問題.」

這個城巿,以城巿之力,拒絕貧窮人.

這種公告標示著:「窮人沒資格活在這偉大的城巿裡.」

在基督裡,卑微的小女子馬利亞,體會神尊重她的尊嚴.

我們被訓練,去求人,就要低下頭,認自己無用.

基督提醒我們,就算是最無用的人,神都要保護他的尊嚴.

我們要努力,讓貧窮人不會為了討食,而必須出賣尊嚴.

這極之需要我們學習尊重貧窮人.

2006年諾貝爾和平獎的得主是孟加拉的經濟學者尤努斯(Muhammad Yunus),他的信念很值得我們重溫.他得諾貝爾獎是實至名歸的,不像今年的得獎者那樣引起多方討論.

一篇訪問尤努斯的報導的文章標題是: ``To protect the dignity of the poor, give them credit not charity"「保護窮人的尊嚴,最好給他們信貸而不是施捨」.

尤努斯發覺,窮人很脆弱,一次天災,便會死很多人.但窮人也和其他人一樣,充滿創意,只是他們缺乏資金去謀生.尤努力開設「鄉村銀行」(Grameen Bank, ``Village Bank"),為貧窮人士提供「微額信貸」.農村婦女,其實只需要很少資金,五十元至一百美元,便可開始一些小生意,如養一些家畜或編織草蓆.就這樣簡單,這些窮人的生活便開展一良好的循環.他貸款的原則是:「你越窮,就享有更優先的待遇.」越窮的人,越能享受Priority Banking.一般的銀行信念是:你越窮,越沒有財富也就越沒有風險,越不講信用,也越不能借到錢.反過來,你越有錢,越要承受損失的風險,越講信用,也就越能借到錢.結果是:銀行是典型嫌貧愛富的機構.越有錢的人,越能借更多錢.越無錢,越不能指望銀行幫助你.窮人因缺乏那怕是很微少的資金,便無法改善生活.

尤努斯的偉大,是他堅信:「一個人無論再窮,他都有權利被人信任」.這些無人信任的窮人,是值得信任的.於是,乞丐也能借到錢.別的銀行不屑一顧的人,在尤努斯那裡,有尊嚴地享受Priority Banking.

天災來臨時,別的銀行只顧叫客人快些還錢,這使客人們只會為失去的東西哭泣.但尤努斯卻給他的客人新的借貸,叫窮人們不要為失去的哭泣,而能盡快重新開始.救濟並不能解決問題.救濟叫人失去前進的動力.救濟叫人失去尊嚴.我們要保護人的尊嚴.要如此做,必須給人機會,只要有機會,窮人是會為自己的尊嚴而奮鬥的.

最後,尤努斯向世人證明,窮人是值得信任的,他們的貸款償還率超過百分之九十九.

人生在世,有人生於富裕國家,有人生於貧窮國家；有人生於富人家庭,有人生於貧窮家庭；有人有手有腳,有人無手無腳；有人有高智慧,有人智商較低；有人生活平順,有人諸多不幸.成為貧窮,有點無奈.但只是有人願意有尊嚴地對待他們,給他們機會,他們便有可能離開貧窮.貧窮的存在,有點無可奈何.但若叫貧窮的人沒有人性尊嚴地活,則是我們「沒有人性」的殘酷造成的.

我們又回到路加福音書的經文.

小女子馬利亞體會到基督降生的意義,她明白,基督要來保護窮人的尊嚴.作為一個卑微的人,她體會到,這正是神偉大之處.這對人而言,是一福音.她要宣告這福音,為要我們知道,基督降世的重要性.這是一個重拾人性的故事,這是一個富人和窮人互相尊重的神聖要求.我們看看馬利亞對基督降生的體會:路加福音 1:46-53

路 1:46  馬利亞說:我心尊主為大；路 1:47  我靈以神我的救主為樂；路 1:48  因為他顧念他使女的卑微；從今以後,萬代要稱我有福.路 1:49  那有權能的,為我成就了大事；他的名為聖.路 1:50  他憐憫敬畏他的人,直到世世代代.路 1:51  他用膀臂施展大能；那狂傲的人正心裡妄想就被他趕散了.路 1:52  他叫有權柄的失位,叫卑賤的升高；路 1:53  叫飢餓的得飽美食,叫富足的空手回去.

一個小故事:

在大陸某個城巿,一個民工上了一公車.他衣衫襤褸,身上發出一陣味道.售票員叫他買票.他拿出一張舊得不能再舊的二十元來,遞給售票員.售票員說,沒零錢找贖,請給零錢.那民工說,我沒零錢.售票員說:沒零錢,則下車.那民工求情.那售票員堅持.在售票員快趕那民工下車之際,有人說:我替他付.

「我替他付」,這就是聖誕節的精意.

但願,你能學懂:「我替他付.」

但願榮耀歸於聖父,聖子,聖靈.阿們.



\newpage

\section{誰是老闆}
\label{sec:29}
\textbf{鄧瑞強博士 - 「誰是老闆?」}
\newline
\newline
\hyperref[sec:28]{< < < PREV SERMON < < <}
~
\hyperref[sec:toc]{[返目錄]}
~
\hyperref[sec:30]{> > > NEXT SERMON > > >}


約翰福音 18:33-38
\newline
\begin{longtable}{cl}
\hline
\hline
章節 & 經文 (和合本修訂版)\\
\hline
18:33 & \begin{tabularx}{0.7\textwidth}{X} 於是彼拉多又進了總督府，叫耶穌來，對他說：「你是猶太人的王嗎？」 \end{tabularx} \\ \\ \relax
18:34 & \begin{tabularx}{0.7\textwidth}{X} 耶穌回答：「這話是你說的，還是別人論到我時對你說的呢？」 \end{tabularx} \\ \\ \relax
18:35 & \begin{tabularx}{0.7\textwidth}{X} 彼拉多回答：「難道我是猶太人嗎？你的同胞和祭司長把你交給我。你做了甚麼事呢？」 \end{tabularx} \\ \\ \relax
18:36 & \begin{tabularx}{0.7\textwidth}{X} 耶穌回答：「我的國不屬於這世界；我的國若屬於這世界，我的部下就會為我戰鬥，使我不至於被交給猶太人。只是我的國不屬於這世界。」 \end{tabularx} \\ \\ \relax
18:37 & \begin{tabularx}{0.7\textwidth}{X} 於是彼拉多對他說：「那麼，你是王了？」耶穌回答：「是你說我是王。我為此而生，也為此來到世界，為了給真理作見證。凡屬真理的人都聽我的話。」 \end{tabularx} \\ \\ \relax
18:38 & \begin{tabularx}{0.7\textwidth}{X} 彼拉多對他說：「真理是甚麼呢？」 \end{tabularx} \\ \\
[1ex]
\hline
\hline
\end{longtable}
啟示錄 1:4-8
\newline
\begin{longtable}{cl}
\hline
\hline
章節 & 經文 (和合本修訂版)\\
\hline
1:4 & \begin{tabularx}{0.7\textwidth}{X} 約翰寫信給亞細亞的七個教會。願那位今在、昔在、以後永在的神，與他寶座前的七靈， \end{tabularx} \\ \\ \relax
1:5 & \begin{tabularx}{0.7\textwidth}{X} 和那忠信的見證者、從死人中復活的首生者、世上君王的元首耶穌基督，賜恩惠和平安給你們。他愛我們，用自己的血使我們從罪中得釋放 ， \end{tabularx} \\ \\ \relax
1:6 & \begin{tabularx}{0.7\textwidth}{X} 又使我們成為國度，作他父神的祭司。願榮耀、權能歸給他，直到永永遠遠。阿們！ \end{tabularx} \\ \\ \relax
1:7 & \begin{tabularx}{0.7\textwidth}{X} 「看哪，他駕雲降臨；眾目都要看見他，連刺他的人也要看見他；地上的萬族要因他哀哭。」這是真實的。阿們！ \end{tabularx} \\ \\ \relax
1:8 & \begin{tabularx}{0.7\textwidth}{X} 主神說：「我是阿拉法，我是俄梅戛 ，是今在、昔在、以後永在的全能者。」 \end{tabularx} \\ \\
[1ex]
\hline
\hline
\end{longtable}
日期:2010年1月10日

主題:「誰是老闆?」

經文:約翰福音18:33-38;啟示錄1:4-8

證道:鄧瑞強博士各位弟兄姊妹,主內平安.

近來,香港政制改革的問題,引起不同政治勢力的角力.這涉及「誰是老闆」的問題.這也是今日的講題.「有得揀,你至係老闆」,這口號,有人曾經用來競選特首.這口號意圖標示出,誰才是社會的老闆,而老闆的特質就是有選擇自由.關於自由,有一個這樣的故事:

「美國人和前蘇聯人,爭論誰的國家更自由.美國人說:我能在白宮前,喊叫打倒美國總統,之後不會有什麼後果.前蘇聯人說:這有什麼出奇,我也可以在克里姆林宮前,喊叫打倒美國總統,之後也不會有什麼後果.」這笑話隱晦地指出:哪個才是真自由的社會.

對政治來說,揀一種政制,選一個政治人物,採納一個政治藍圖,這關乎巿民的福祉,這都是巿民應該盡可能參與抉擇的範圍.但有另一種層面的範圍的抉擇,對人的生命意義,影響更深遠.這就是人生的王權屬誰的問題.今日是基督教聖曆裡的「基督君王日」(Christ the King Sunday),高舉基督為王.基督的王權,當然涉及人的生命主權,但這主權,和世上的政治權力,不在同一層面.單單看到世上的權力,視野有限,也容易犯錯.能看到天上有一比人間更大的權力,就能讓人警惕,小心犯錯.

今日的講道

經文:約翰福音 18:33-38

18:33　彼拉多又進了衙門,叫耶穌來,對他說:「你是猶太人的王嗎?」18:34　耶穌回答說:「這話是你自己說的,還是別人論我對你說的呢?」18:35　彼拉多說:「我豈是猶太人呢?你本國的人和祭司長把你交給我.你作了什麼事呢?」18:36　耶穌回答說:「我的國不屬這世界；我的國若屬這世界,我的臣僕必要爭戰,使我不至於被交給猶太人.只是我的國不屬這世界.」18:37　彼拉多就對他說:「這樣,你是王嗎?」耶穌回答說:「你說我是王.我為此而生,也為此來到世間,特為給真理作見證.凡屬真理的人就聽我的話.」18:38　彼拉多說:「真理是甚麼呢?」

這是彼拉多和耶穌的對峙.他們的話題聽起來頗具政治性,涉及王權的問題.

在這經文裡,彼拉多是在內庭審問耶穌的.這可以看做是彼拉多內心與耶穌相遇.彼拉多的內心不是平安的,因為在外面,有群眾的要求,這班群眾要求彼拉多處死耶穌.看起來,耶穌與當時有權的群眾並不合拍.耶穌潔淨聖殿,干擾了有權者的經濟秩序,這經濟方式欺壓人.任何令弱者被欺壓的經濟制度,可以推想,都是耶穌要潔淨的制度.耶穌接納罪人,保護了卑微者的人性尊嚴,可以推想,任何叫人性尊嚴被取消的行動,耶穌都反對.耶穌在比喻中,欣賞富有愛心的撒瑪利亞人,藉此反對對外人的歧視,可以推想,任何製造歧視的想法,耶穌都不同意.我的經驗發現,在巴士上,外籍傭工很樂於讓座給有需要的人.若耶穌再講故事,他可能會說:「在一天,有個老人家上巴士,他真的需要一個座位.一位身上掛著十字架的人假裝睡覺,一位大學生見到這情況仍自顧自的打遊戲機,只有一位外籍傭工讓座給這有需要的人.你們說,誰是那位有需要的人的鄰舍呢?誰能承受永生呢?」這些故事,總有人不喜歡聽.人不喜歡被人揭露自己潛在地欺壓人,看低人,歧視人的真相,人反對揭露真相的力量.人反對耶穌.彼拉多承受著外面群眾的壓力,他在自己的內在世界裡,那裡有耶穌,也有羅馬政權給他的要求.羅馬政權要求他好好管理好這片土地.對彼拉多而言,最真實的權力就是羅馬的權力,一種現世最大的權力,除此,他不認識別的權力.然而,面對耶穌,他看來要面對一種他不熟悉的權力,其實,這權力我們又何嘗真的熟悉?彼拉多問耶穌:「你是猶太人的王嗎?」耶穌是否世上的政治勢力的一種?耶穌是否一種擾亂公眾安寧的勢力?

越是獨裁政體,越怕人越出正常軌道.政治笑話,卻是在軌道上的越軌行為.有一個故事是這樣的:「在某一個獨裁政體裡,人民不准講笑話取笑領袖.有一次,一個愛講笑話的人被捕了.他被帶到領袖面前.那領袖對他說:『你可知道,取笑我是不會有好結果的.你知嗎,我的政權會持續千秋萬代.』那人聽後,說:『嘩,你這笑話我還是第一次聽到.』」

耶穌的存在,的確是一種笑話,讓人看到沉悶不堪的生活背後,有一值得歡笑的力量.

耶穌回答彼拉多說:「我的國不屬這世界；我的國若屬這世界,我的臣僕必要爭戰,使我不至於被交給猶太人.只是我的國不屬這世界.」耶穌是王,但他的國不屬這世界.關於這國度,彼拉多一無所知.的確,生活可以屬於另一個國度,這國度的秘密很多人不知.若果說,一個人放棄了一份高薪的職業而選擇了一份薪金低微卻能幫助他人的職業,他是屬於那神秘的國度的.若果說,一個人沒有權選擇而只能在一份薪金低微的工作上工作,卻仍能甘心樂意藉這工作服事人的,他是屬於那神秘的國度的.若果說,一個人沒有工作,只能靠救濟生活,但他仍主動以僅有之物去救濟同樣有需要的人的,他是屬於那神秘的國度的.

德蘭修女見證過這樣的一件事,她拿米糧去救助一個貧困家庭.這貧困家庭卻拿得到的米糧去救助貧困的鄰舍.若彼拉多只想著如何回應群眾,只注意現世的國度,只關心自己的權力和利益,則他便無法看到耶穌的國度了.同樣,若你想像不到,你的時間除了花來讓自己增值外,還有什麼用?若你想像不到,錢除了令自己有更好的享受外,還有什麼用?若你想像不到,一生除了為自己打算外,還能怎麼樣?若你這樣沒有想像力,則很難看到耶穌的國度了.

早幾天,我和學生行山,行八仙嶺,走當日八仙嶺大火師生們走過的路,走到仙姑峰下面的「馬騮崖」,一個頗難越過的懸崖,在這裡,我試圖想像一下那兩位老師的內心世界.大火到了,自己帶領著十多位同學.若各自走,則老師應能脫火海.但是,這刻,可能沒有想到其他東西,沒有想過自己是否會死,沒有想過死後家人的傷心,沒有想過一生的努力就此完結,可能什麼也沒有想,只想救那些跟著自己的學生.但也可能什麼也有想過,就算火燒到自己,仍選擇留下.當他們感到火燒的痛苦時,他們卻想著無助的同學們.當他們的求生本能驅使他們離開時,他們抑制了這本能而協助同學們逃生.當他們生命即將完結時,他們的眼睛看到生命的永恆.他們活在耶穌所講的那國度中.

18:37　彼拉多就對他說:「這樣,你是王嗎?」耶穌回答說:「你說我是王.我為此而生,也為此來到世間,特為給真理作見證.凡屬真理的人就聽我的話.」18:38　彼拉多說:「真理是甚麼呢?」

耶穌是王,但這個人的王權,是建立在說真理這事上.我有一位朋友,他說,他從電影裡黑社會講數的橋段中,學到的功課是,連黑社會講數也要講道理.當然,道理不是真理,真理是揭露真相的.有一個故事是這樣的:

「在前蘇聯,發生了一次暗殺領導人的事件.但這次是一次失敗的暗殺事件,子彈打不中該領導人.於是,獨裁政府組成調查小組,去調查事件的真相.最後,知道子彈打不中領導人的原因.真相是這樣的:當那位行刺者拿槍出來對準領導人時,他身旁的人問他,你做什麼?他告訴這人,要行刺領導人.這消息傳開後,太多人走過來爭著開槍.由於太多人爭的緣故,子彈便打不中了.」

在美國,發生另一個故事:

「話說,有一個庸碌而胡亂開戰的美國總統,有次坐空軍一號,飛翔於美國上空.他忽發奇想,說:在這戰爭期間,人心不安,若我將一張一百元鈔票丟下去,豈不會至少有一個人開心.他的夫人說:若丟十張十元鈔票下去,豈不有至少十個人開心.副總統說:若丟一百張一元鈔票下去,豈不有至少一百個人開心.隨行記者說:總統先生,不如你跳下去,這樣,全國都開心.」

這些故事,揭露著某種真相.揭露了前蘇聯的獨裁,也揭露了美國的霸權.

耶穌來,要為真理作見證,揭露出人生的真相.耶穌要說,世界是有神的,神的作為在人間.神審問誤用權力的人,卻叫受欺壓的人得安慰.耶穌為揭露真理而生,也為之而死.在世間的王面前,高舉另一種王權,真理的王權,是危險的.

緬甸的昂山素姬(Aung San Suu Kyi)也是一個為真理而生,為真理而死的人.她說過:「(英文原文在下面)在緬甸決意爭取民主的人,是確信:在壓迫人的社會裡,為了堅持人權而承受危險,好過在為奴的安逸生命裡享受安全.我們的非暴力運動,堅持一信念,就是相信人喜愛公平和憐憫.有人說,人基本上只是經濟動物,只顧自己的物質快樂.這看法實在太看扁人了.有無數男女,為了堅守普遍公認的信念和原則,準備承受無情的迫害.在我的國家裡,今天仍有這些男女,這是我的榮幸,也是激勵我的力量.」(``Those of us who decided to work for democracy in Burma made our choice in the conviction that the danger of standing up for basic human rights in a repressive society was preferable to the safety of a quiescent life in servitude. Ours is a nonviolent movement that depends on faith in the human predilection for fair play and compassion. Some would insist that man is primarily an economic animal interested only in his material well-being. This is too narrow a view of a species which has produced numberless brave men and women who are prepared to undergo relentless persecution to uphold deeply held beliefs and principles. It is my pride and inspiration that such men and women exist in my country today."── ``Please Use Your Liberty to Promote Ours" in the International Herald Tribune)

這裡揭露的,是一種力量,一種抵抗現世不義力量的力量.這不是戰爭的力量,而是一種和平力量.這力量指向一條路,這條路引領人活得更豐盛.這不是為自己的物質享受和個人安舒而活了一生,而是為了他人,為了超越現世的另一個國度.耶穌是這國度的王.

18:38　彼拉多說:「真理是甚麼呢?」

這不是一種求真理的態度,這是一種放棄真理的態度.面對耶穌,卻失之交臂.多麼可惜.

關於耶穌所講的王權,今日的書信經課啟示錄1:4-8有多一點的說明.請大家看這段

經文:啟1:4-8

1:4　約翰寫信給亞西亞的七個教會.但願從那昔在,今在,以後永在的神,和他寶座前的.

神是「昔在,今在,以後永在的神」,在歷史的流逝裡,唯有神在歷史之上,引導歷史的前進.他坐在王者的寶座上,審問世人.「七靈」指聖靈,聖靈和所有在神國度裡生活的人同在.他給人力量,在謊言處處的地方講真話,在物質的安逸裡指出屬靈價值的可貴.

1:5　並那誠實作見證的,從死裡首先復活,為世上君王元首的耶穌基督,有恩惠,平安歸與你們!他愛我們,用自己的血使我們脫離罪惡,1:6　又使我們成為國民,作他父神的祭司.但願榮耀,權能歸給他,直到永永遠遠.阿們!

這裡有較大篇幅論述耶穌.耶穌是那位「誠實作見證的」,他堅持真理.他是那「從死裡首先復活」,人面對死亡,會變得懦弱,會甘願成為強權下的奴隸.但這裡,有一位勝過死亡的王者.人講真話,可能會死,但這裡有一位復活主.耶穌是「世上君王的元首」,他是王者之中的王者.這位萬王之王,高高在上者,你以為他遠離人群.你錯了.「他愛我們,用自己的血使我們脫離罪惡」,他走到人群中,不是像「政治家」為了選舉而巡區做秀,他在人群中,承擔了人間疾苦.他「使我們成為國民,作他父神的祭司」,他呼召我們實現神的國度的能力,將這國度顯現出來.他呼召我們成為祭司一般的人物,在人間顯示神聖.

1:7　看哪,他駕雲降臨!眾目要看見他,連刺他的人也要看見他；地上的萬族都要因他哀哭.這話是真實的.阿們!1:8　主神說:「我是阿拉法,我是俄梅戛,〔阿拉法,俄梅戛:是希臘字母首末二字〕,是昔在,今在,以後永在的全能者.」

「看哪,他駕雲降臨!」他要審判一切行惡者.他是「我是阿拉法(Alpha),我是俄梅戛(Omega),是昔在,今在,以後永在的全能者.」他是一切美善的根源,他是創始成終者.

一個故事:「在美國,在一個制定政策的會議上,兩名議員互罵起來.當時的美國總統卡特在場,大吼一聲,說:安靜,我想你們可能興奮得完全忘了我還在這裡.」

誰是老闆?在擾擾攘攘的人生中,耶穌大喝一聲:安靜,我才是老闆.

但願榮耀歸於聖父,聖子,聖靈.阿們.



\newpage

\section{一件美事}
\label{sec:30}
\textbf{戴國璋牧師 - 一件美事}
\newline
\newline
\hyperref[sec:29]{< < < PREV SERMON < < <}
~
\hyperref[sec:toc]{[返目錄]}
~
\hyperref[sec:31]{> > > NEXT SERMON > > >}


日期:2010年1月17日

主題:一件美事

證道:戴國璋牧師首先很感謝伍牧師安排我今天在這裡跟大家見面,很高興能夠在教會裡面跟那麼多屬主的會眾一齊敬拜,但願上帝賜恩給大家,不但大家在生活上,蒙上帝奇妙的恩,在我們教會發展上,服侍上也可以看到上帝的恩典,常常使用這個地方.

今天早上會讀一個大家也很熟悉的故事,這個故事其實是要講述耶穌基督在世的時候,當他還在肉身裡,他最後一次去享受,一些愛他的人去為他所準備的服侍,一個筵席,因為根據聖經告訴我們,過了那天晚上以後,沒多久以後耶穌已被釘在十架上面了.這個是一個很熱鬧的筵席,因為按照 馬可福音跟約翰福音的描述告訴我們,當天的筵席其實是一個很重要的筵席,有許多的人也參加了這個筵席,那當然這個筵席的主角是我們的救主耶穌,當筵席到達 了高潮的時候,大家也在那個酒筵當中,大家都彼此歡欣的時候,忽然發生了一件事情,原來有一位女子走到耶穌的身邊,這件事情本來並不值得去注意的,因為對 很多人來說,這個只不過是一位侍女,在這個筵席裡面,這位女子要為耶穌基督去搬弄一些杯子,盤子跟碗碟,可是當這位女子走到耶穌的身邊的時候,突然從她的懷裡拿 了一樣東西出來,然後「澎」一聲的,所有筵席上的人都被吸引而看過去,之後大家也都聞到一陣香味,一種很強烈的香氣,原來這位女子把她懷中的一瓶香膏打破 了,然後把這個香膏澆灌在耶穌基督的頭上面,不單這樣子,當時的婦女為了方便她們自己做事,她們會把她們的頭髮用頭巾包裹好的,可是這位婦女突然把自己的 頭巾扯下來了,把她的頭髮散亂了,然後她用她的頭髮塗上香膏,替耶穌擦腳,整個行動吸引了大家的注意力,同時也引發了我們現在要讀的故事,馬可福音裡面並沒 有提供這位婦女的資料,我們不知道她是誰,可是我們知道對比約翰福音的記述的時候,我們就知道這位女子就是馬利亞,就是那位一家三口,包括她的姐姐馬大, 弟弟拉撒路,住在伯特利村莊的一位婦女,也是耶穌所關心,耶穌所關愛的一個家庭,這件事情招惹了很多的非議,很多人的批評,可是在我們的耶穌基督的身上卻得 到了好好的評價,主耶穌說,她為耶穌所做的,是一件的美事.

我們都是上帝的人,都是屬於耶穌的人,我們都希望我們所做的事,特別是我們為耶穌基督所做的事,能夠得到耶穌基督的稱讚,成為一件美事,可是到底有什麼素質,能夠讓耶穌基督去稱讚馬利亞所做的事,是一件美事呢?

這個是值得我們去細想的,或許能夠緊更加提醒我們今天為耶穌所做的事物,如果能夠包括這些素質的話,就能夠讓耶穌基督稱讚 為一件美事,我們感謝上帝,因為祂並沒有在這裡隱藏了一個原因,沒有告訴我們這個事情上面包含了什麼素質,以致成為了一件美事,祂是這樣告訴我們的,所以我們可以從後面耶穌基督所講的一些話裡看到這些素質,首先,在第一方面,耶穌在十四章六節裡面很清楚的告訴我們,她(馬利亞)在祂身上作了一件美事,在這句話裡 面,已經有提到為什麼這是一件美事,因為她是作在耶穌基督身上,凡作在耶穌基督身上的都是一件美事.

我們再看看馬利亞的行動的時候,我們可以發現馬利亞所作的,獨獨地作在耶穌基督身上,這個獨獨地為耶穌而作,把她的事奉突顯了出來,我們已經講過,那個筵席,其實是為耶穌所設的一個筵席,所以在整個筵席裡面耶穌是主角,很多人為耶穌去準備這個筵席.

我們知道這個筵席裡面的東道主,是在第三節所提到的一位長大痳瘋的西門,雖然有關於西門的資料不多,可是這個人都是尊重耶穌,因為他愛護耶穌,所以為耶穌設下這個筵席,他很慷慨,因為他不單招待了耶穌,他還招待了很多人,有門徒,有很多從其他地方慕名來見耶穌的人,甚至約翰福音告訴我們,在耶路撒冷一些人,因聽見耶穌行許多神蹟,使拉撒路從死裡復活,於是都來看看祂. 所以有很多人在這裡,西門很慷慨地用他自己的地方去設這個筵席,來招待耶穌及許多的人,這個真的是值得稱讚.

我們也知道馬大也在這個筵席裡面,而且他在這個筵席裡面也很殷勤的服侍,帶領很多婦女進進出出的去準備各樣的東西,讓每位在筵席裡面的人,都能夠享用最好的東西,所以馬大就在那兒服侍.

約翰福音也告訴我們,其實拉撒路也在裡面,拉撒路作為一個曾經經歷從耶穌基督裡面死裡復活的人,可以見證耶穌基督是充滿恩典及能力的人,他就坐在那邊陪伴著耶穌,讓耶穌可以好好享用一個很開心的宴會,這些人都是服侍耶穌的,這些人所作的一切都是為了耶穌,都是值得我們欣賞.

可是如果我們要把馬利亞為耶穌所作的,跟剛才那些為耶穌服侍的人而作出比較的時候,我們仍然可以發現馬利亞為耶穌服侍特別 的地方,這個特別的地方是其他人的服侍作在耶穌基督的身上,並不是為耶穌獨獨的享用,因為招待耶穌的也有招待其他的人,陪伴著耶穌的人,也有陪伴其他的人,唯獨馬利亞的這個香膏,她打破了香膏的瓶子,完全的倒到耶穌身上,這裡我們看到馬利亞一心一意,把她所有的一切東西都放在耶穌基督的身上,耶穌成為了 她唯一重視的.

這個給我們看到服侍真正的義意,我們也可以試試看,用一個比較神學的角度去表達這個義意,這個特別的講法是以神為中心,馬利亞的服侍是一個以上帝為中心的一個服侍,她要把一切的服侍都放重在上帝的身上,可是今天也有人把人為中心,因為現在的人都覺得人是很重要的,人是很值得 我們去關懷的,今天我們也覺得應該服侍那些屬上帝的人.

作在人的身上便是作在神的身上,這個觀念是沒有錯的,事實上耶穌基督也是這樣子教導我們,耶穌基督曾經教導我們,祂說:「作在一個小子身上,也是作在我的身上」,所以即使是給一杯涼水,就像以前的人從水井裡面拿出來的水給一位小子,耶穌說也是作在祂身上,我們也了解到我們 用什麼去表達我們對那位看不到的上帝呢?就是去愛一些看得到的弟兄,所以今天我們覺得這個是很重要的,我們是透過服侍人去服侍我們的上帝,在教會裡面我們 所做的一切都是為了神,但是當我們認識這個道理的時候,會還需要像馬利亞的行動一樣,去提醒我們,當我們去服侍弟兄姊妹而去服侍上帝的時候,我們真的需要 仍然以上帝為中心,我們是為上帝服侍,我們是因為對上帝的愛,對上帝的關心,所以我樂意參與其中,不介意沒有任何的回報,這個是我們可以從馬利亞身上所學 到的事情.

我講到這裡,我想起基督界一位很難得姊妹,一位已經過世的姊妹,在基督教的大家庭裡面,這一位就是德蘭修女,我們知道德蘭 修女在印度做了很多美麗的服侍,她所做的工作是很多人都會覺得很厭惡的,可是她為了上帝的原因,幫助了印度很多很貧窮的百姓,她做了很多美麗的事情,所以 當她離開這個世界的時候,很多人都捨不得,天主教也有人問,可不可以為德蘭修女封聖?可是,封聖在天主教會裡面有很多規矩的,所以要做並不是一件容易的事 情,但這個可以証明德蘭修女受到很多人的尊重以及敬愛.

我們知道自從德蘭修女離開我們以後,開始有很多人研究德蘭修女的生平,有些人可能拿到一些特別的資料,了解到德蘭修女的內 心世界,有些人就說,德蘭修女留下了一些文獻,在她的內心裡面也有恐懼不安的時候,身邊的人看起來可能會覺得不好,因為可能是一些心理衛生的不健康,時刻 都很擔心自己所做的一切,是否上帝所喜悅的,她所做的一切,神真的是喜愛的嗎?常常落在這種擔心,恐懼當中,有些人會作出批評,覺得德蘭修女的屬靈生命或心理 衛生有問題,大家如果曾經看過或是聽過這些講法,請你不要相信,為什麼呢?

因為其實當我們了解或是研究一些在基督教出現過的一些屬靈巨人的生路歷程的時候,很多這些偉大的屬靈偉人,都有一些類似的經驗,他們生活在上帝面前的時候,或是在服侍上帝的過程的時候,也會感到自己的不足,會有強烈的恐懼,覺得做什麼都不能夠滿足上帝的時候,為什麼呢,因為上帝是多麼的神聖,是多麼的美好,很多前輩也有這種經驗,這些前輩的經驗實在是能夠鼓勵一些後輩.

就像馬丁路德也有這些經驗,他才能展現出因信稱義的思想,所以其實這個是很自然的,為什麼呢? 因為當我們要服侍主的時候,當服侍的重點都放在上帝的當中,在這個時候,我們很多時候都會感覺到自己的不足,越是感覺到自己的無能,其實這是好的,因為會 提醒我們,在我們每次服侍裡面,我們都會注意到以上帝為中心,以討祂喜悅為中心.這個從耶穌基督裡面所看到的,作在我身上面,就是馬利亞之所以可以蒙上帝 的稱讚為一件美事的一個行為原因.

我們再看看第二方面,第二個素質在哪呢? 其實耶穌基督在十四章八節裡面說:「他所作的,是盡他所能的」,他是已經盡力了,因為他盡了他所能的,所以他所作的是一件美事,這是耶穌基督的講法,可是 很可惜,耶穌基督沒有再解釋清楚,這個盡他所能的''能''是什麼意思,是那一方面的事呢,所以我們只能夠去作一些的猜測,這個會是指經濟上的能嗎,因為他已 經在他經濟上盡他所能,用了最貴的香膏在我身上,這會是一件美事嗎,也有可能的,我們知道這是最貴的香膏,這是一個很珍貴的香膏.

據說這個哪噠香膏是從印度喜馬拉雅山上,一種植物的根部所提煉出來的,在古時代,特別是西亞,一個在亞洲最西面的地方巴勒 斯坦,對他來說印度是一個很遠的地方,對他們來說就像東部的地極一樣的遠,事實上當時的人覺得,有點知識的人會知道,只會知道印度這個地方,再遠一點的地 方他們跟本不曉得,從一個那麼遙遠的地方的植物所提煉而成的香膏,再從那邊運送過來,那你們覺得他的價錢會是多少呢?

所以聖經告訴我們,就是因為哪噠香膏,是那麼貴重的原故,所以她使用了一個玉瓶,放那些香膏,我們都知道玉瓶都是一些比較珍貴 的石頭,所以這些香膏的價值得很名貴,賣三十兩銀子,當時是相當於當時那些工人一年的工價,可是馬利亞她只不過是一個鄉下的女子,她竟然為了耶穌而用了一個這 麼貴重的禮物,所以我們說,如果她是盡所能的,經濟上面的說法是可以的,可是會不會除了經濟上的可能以外,也是感情上的可能呢,是因為馬利亞為耶穌傾出一 切的感情上的可能,其實我們也知道當時一些比較富裕的家庭,特別是一些富裕的女子,在哪噠香膏的預備是有特別的目的,很多時候她們會用盡家財去買這些香 膏,作為她們人生大事上唯一一次的用途,這個人生大事就是說當她出嫁的時候,當她要出嫁的時候,她會好好的沐浴,洗乾淨自己的身體,然後把這些那麼珍貴的 香膏塗抹在自己身上,然後當她跟她的丈夫完房的時候來取悅她的丈夫,當時有錢的婦女就是這樣子用的,我們不知道馬利亞買這個香膏時的用意是不是這樣,可是 在這一刻裡面,她用盡了香膏在耶穌身上,表明她覺得在這個時刻裡面.她最愛慕,最尊重,最值得去使用這個香膏的,不是她那個將來可能會出現的愛郎,而是她 現在所愛的主,那位在不久以前使她所愛的弟弟從死裡復活的主,她願意全部都放在祂身上,這個不是全然感情上的投身嗎?

所以這個''能''可以是指經濟上的投身,也有可能是指感情上的投身,可是我想還有一個可能,這個可能才是馬利亞所''能''真正 所指的所在,這個所能,其實是指意志上,就是她要愛耶穌,要在意志上顯出她對主耶穌的愛,因為意志是管理行動的,如果我們沒有強烈的意志,就難以去實踐一個艱難的行動,馬利亞的行動招徠了很多非議,在第四節說:''何用這樣枉費呢'',似乎這個是當時共有的一個非議.

我們知道其實猶太人的傳統是很重視對窮人的施捨,所以猶太人每天會有一些救濟的行動,會去救濟一些窮乏的,流浪的,沒有食 物的人,去為他們預備一些食的給他們,每星期也會有一些救助的行動,為他們預備一些衣物給一些的窮人,特別是我們知道當時是節期的時候,馬可福音在第一 節告訴我們,過兩天便是逾越節,逾越節就在前面了,猶太人相信在節期上面行善,是會特別得到上帝對他們的祝福及重視,特別是當行善在耶路撒冷上面的時候, 因為耶路撒冷上面有很多很特別窮乏的人,如果在節期的時候在耶路撒冷行善事,是特別得到上帝的祝福的.

我們知道當時有這麼一個背景,所以在聖經裡面告訴我們,當耶穌跟祂的門徒,在享用聖餐的當天晚上,耶穌跟猶大談了幾句話,猶大便走出去,準備這個賣主的行動,當時的門徒並不知道發生什麼事,門徒以為猶大外出是去施捨,因為當時是節期,在節期的時候是特別的令人提醒,原來施捨是 這麼的重要,所以猶太人很重視這個施捨的行動.

今天基督教也很重視施捨,所以很多的慈善基督教也不敢太落後,坦白說,基督教這種樂意行善,樂意施捨,其實也是從猶太人學 回來的,我們接受了這個傳統,我們也發揚了這個傳統,聖經告訴我們,何用這樣枉費呢,可是當我們細想一下,其實這個後面,是包含了很多深層的嫉妒,我們可以從這些 人,隱藏著在他們面容後面心靈所聽到的聲音,聽到的評語是另一類的.

比如說,這個香膏價值三十銀錢,是很貴的,馬利亞只不過是一個鄉下的女子,她去那裡賺錢買呢,她何來這種能力呢,有些人就開 始說,她可能從一些見不得光的地方去賺錢來買那麼貴的香膏,也有可能我們會聽到一些聲音說,這個女人憑什麼,?耶穌那麼多的門徒在,可是他們也沒有向耶穌 表示那麼強烈的愛,馬利亞竟然搶先佔有了這件事,可能我們又會聽到另一類的評語,這個是什麼女人呀,為什麼她可以那麼不知廉恥去靠近一位陌生男子,這個都 不是一位健康女子適合的舉動,到低她想要怎樣呢,可能我們會知道或是聽到,感受到許多這些惡毒的批評.

其實你覺得馬利亞有沒有想過她自己的行動,可能會招徠這樣的批評呢,我猜馬利亞是知道的,馬利亞早就想到,可是她並不害怕,沒事情可以阻撓她對耶穌的那份尊重,所以即使她知道會有這種非議,她還是會盡她最大意志上的力量,把香膏全都用在耶穌基督的身上,這個是馬利亞對耶穌愛 的表現,馬利亞真的做到愛上帝的原則,要盡心,盡性,盡力,盡意,把你的意志都放上.

我們並不是做每一樣的事情都會得到別人的讚賞,有時候我們耗盡我們的心力去設計一些東西,或是預備了一些東西,想說這些東 西會讓上帝的工作更順利,能讓上帝的子民更喜愛,可是很多時候,當我們把這些東西做出來以後,之後就會聽到很多不一樣的批評及意見,我們可能心裡面會受傷 害,有可能我們聽到以後就會退下來,不願意再去做了,我們真的很容易會有這個毛病,當我們有這個毛病的時候,讓我們從馬利亞的身上去學習,只要這是一件對 上帝有益的事情,或是這是一件對上帝子民有益的事,不管有多麼的批評,可是只要我們專心,用我們的意志去服侍,蒙上帝的喜悅.

這個時候讓我想起,德蘭修女的一些事情,德蘭修女在世的時候,在慈善家來說,她是一位人道主義者,因此她得到很多人的尊 重,所以很多時候,有一些很重要或是一些國際性的人道會議也會邀請德蘭修女去參加,她的來臨真的增光不少,我們在電視機上面看到的德蘭修女都是很謙卑的,她什麼人也不會拒絕,都會跟他們握手,只要有人想要跟她握手,她也會跟他們握,所以又出現了很多的批評,那些人會覺得為什麼德蘭修女什麼人都會接見呢?那些人覺得應該要有分別的,應該只有跟一些好人握手,可是當中有一些南美洲的獨裁者,在國家裡面壓迫人民的,可是當人們看見德蘭修女的時候,他們又會拖著德蘭修女,要求拍照,因為他知道這樣子就可以讓他們有機會曝光,因為德蘭修女所以他們就會有機會出現在報紙上面,有很多有野心的人就會藉以這個機會來為自己買 廣告,去推廣他們自己,很多人會去批評,為什麼這些人的手,德蘭修女也會去握,為什麼還要跟他拍照呢?

可是我覺得這樣子的批評是不公平的,因為我們要知道德蘭修女並不是一位政治家,她不會也不應該從政治這一面去考慮與人接 觸,她只不過是一位幫助人的人,她只是上帝的侍女,她只是為上帝服侍,如果有任何人想要她的幫助,德蘭修女也不會拒絕的,所以德蘭修女想用她自己的手,去表達上帝的愛,所以各位弟兄姊妹,當我們在聚會或是要離開的時候,我們都可以跟認識或是不認識的弟兄姊妹握個手,打個招呼,笑一個,為彼此祝福,所以我們不 要害怕任何的批評,要把最好的獻給上帝.

我們再看看第三方面,為什麼這個行動會成為一件美事呢,第八節告訴我們說:「他是為我安葬的事,把香膏預先澆在我身上」, 這一句是更加難明白,馬利亞怎麼可能比耶穌的門徒更早知道在耶穌面前的遭遇,為什麼馬利亞不可能像耶穌基督所講的而捕捉到這位主快要離開我們了,所以`她 要把最好的都獻給耶穌.

過了不久,耶穌基督被釘死在十字架上,釘死在十字架以後,人們很快的把祂的屍體埋葬,因為當日已經是安息日了,所以他們很快的把耶穌安葬了,有人用一百斤的香膏放在耶穌身上,可是耶穌已經聞不到了,所以馬利亞這個行動能夠被耶穌基督稱讚,因為她早就把香膏塗在耶穌基督的身上.

我們知道馬利亞有一個特點,在路加福音有說,當耶穌進到馬大及馬利亞的家庭的時候,馬大很忙亂的服侍耶穌,因為他很愛耶穌,當他看到他的妹妹什麼都不做的時候,他對馬利亞有點埋怨,就跟耶穌說,夫子呀,請你吩咐我的妹妹跟我一齊做吧,因為馬大實在太多事情要去做,太忙了,當時耶穌回答他說,馬大,馬大,你為很多的服侍去忙亂,可是最大福份就是你妹妹已經坐在我面前聽我講的道,馬利亞給我們最深刻的印象就是她很喜歡聽耶穌講 道,一個喜歡聽道的人,一個留心聽道的人,他能夠從耶穌基督的話裡面,能夠知道一些事情快要發生在耶穌身上,為什麼不可能呢?

事實上這個也解釋了在往後一些的迷,比如說,往後耶穌基督所受許多的苦害,描述的裡面,在任何的福音書裡面,提到當日耶穌也 有這些顧慮,可是祂並沒有提到馬利亞這位女子,到了復活的清晨,有幾位婦女一同跑到耶穌的墓,因為在早一兩天,當耶穌的屍體被送到墓裡面去的時候,因為太匆忙了,所以跟本沒有時間去為耶穌準備香膏,所以她們買了香膏以後,想要用香膏,膏抹在耶穌身上,當時大家正在擔心的時候,有大石頭阻擋,誰可以幫助她們拿開 石頭,讓她們可以進去墓穴裡面呢?

我想一個唯一可以解釋,一位那麼愛耶穌的婦女會在那裡出現的原因,就是因為她一早就知道這件事情,而為耶穌基督作準備,這個是馬利亞被耶穌基督所稱讚的第三件美事的原因,而她能夠做到這樣子,是因為她對耶穌基督的真心,真理,留心以及傾聽.

親愛我弟兄姊妹,我們很樂意為主耶穌去做很多事情,我看到大家,當中有很多的恩賜,當我來到大家的聚會當中,看到大家的讀 經,講道,主席,唱詩,也都發覺大家有許多的恩賜,真的感謝上帝賜給這些恩賜大家,也希望大家能夠好好的使用這些恩賜去服侍上帝以及弟兄姊妹,讓教會能夠 更加榮耀上帝,不過在我們許多作在上帝以及作在人身上以前,原來留心聆聽也很重要的,留心聆聽在每次講道裡面,牧師為我們準備的信息,在當中向上帝捕捉上帝要跟我們說的話,去改變我們的生命,去改變我們對事奉的態度,原來這也是上帝喜悅的東西,在這裡我要謝謝大家,在這一個小時的忍耐,聽我的道,相信上 帝一定會祝福大家的.

請大家一起同心禱告:我們的主,我們希望我們能夠學習馬利亞,能夠從這位女子身上,去學習她愛你的心,求你給我們機靈,機敏,看見我們應該怎麼去投身你的服侍,求你常在我們心靈當中受感動,就像當天你感動馬利亞一樣,願意我們把最好的東西,獻上給你,願你喜悅,聽我們的祈 禱,奉耶穌基督的名字,誠心所願!



\newpage

\section{問道的小孩}
\label{sec:31}
\textbf{林豪恩先生 - 問道的小孩(The inquisitive Lad)}
\newline
\newline
\hyperref[sec:30]{< < < PREV SERMON < < <}
~
\hyperref[sec:toc]{[返目錄]}
~
\hyperref[sec:32]{> > > NEXT SERMON > > >}


路加福音 2:41-52
\newline
\begin{longtable}{cl}
\hline
\hline
章節 & 經文 (和合本修訂版)\\
\hline
2:41 & \begin{tabularx}{0.7\textwidth}{X} 每年逾越節，他父母都上耶路撒冷去。 \end{tabularx} \\ \\ \relax
2:42 & \begin{tabularx}{0.7\textwidth}{X} 當他十二歲的時候，他們按著過節的規矩上去。 \end{tabularx} \\ \\ \relax
2:43 & \begin{tabularx}{0.7\textwidth}{X} 守滿了節期，他們回去，孩童耶穌仍舊在耶路撒冷。他的父母並不知道， \end{tabularx} \\ \\ \relax
2:44 & \begin{tabularx}{0.7\textwidth}{X} 以為他在同行的人中間，走了一天的路程才在親屬和熟悉的人中找他， \end{tabularx} \\ \\ \relax
2:45 & \begin{tabularx}{0.7\textwidth}{X} 既找不著，就回耶路撒冷去找他。 \end{tabularx} \\ \\ \relax
2:46 & \begin{tabularx}{0.7\textwidth}{X} 過了三天，他們發現他在聖殿裡，坐在教師中間，一面聽，一面問。 \end{tabularx} \\ \\ \relax
2:47 & \begin{tabularx}{0.7\textwidth}{X} 凡聽見他的人都對他的聰明和應對感到驚奇。 \end{tabularx} \\ \\ \relax
2:48 & \begin{tabularx}{0.7\textwidth}{X} 他父母看見就很驚奇。他母親對他說：「我兒啊，為甚麼對我們這樣做呢？看哪，你父親和我很焦急，到處找你！」 \end{tabularx} \\ \\ \relax
2:49 & \begin{tabularx}{0.7\textwidth}{X} 耶穌對他們說：「為甚麼找我呢？難道你們不知道我應當在我父的家裡嗎？」 \end{tabularx} \\ \\ \relax
2:50 & \begin{tabularx}{0.7\textwidth}{X} 他所說的這話，他們不明白。 \end{tabularx} \\ \\ \relax
2:51 & \begin{tabularx}{0.7\textwidth}{X} 他就同他們下去，回到拿撒勒，並且順從他們。他母親把這一切的事都存在心裡。 \end{tabularx} \\ \\ \relax
2:52 & \begin{tabularx}{0.7\textwidth}{X} 耶穌的智慧和身量，並神和人喜愛他的心，都一齊增長。 \end{tabularx} \\ \\
[1ex]
\hline
\hline
\end{longtable}
日期:2010年1月24日

題目:問道的小孩(The inquisitive Lad)

證道:林豪恩先生

經文:路加福音2章41-52節

一．    引言

耶穌出場的時候是人生的甚麼階段呢?馬可福音及約翰福音的耶穌出場時已經是三十歲的成年人.馬太福音及路加福音的耶穌出場時卻是一位初生嬰兒,或者可以說只是腹中的胎兒.因此,如果沒有馬太福音及路加福音,我們便不會有聖誕節的情景及想像了.不過,縱使嬰孩耶穌出現於馬太福音及路加福音的記述中,篇幅卻只像過場一般的短促,鏡頭一轉,三十年就過去,這位嬰孩耶穌再出場的時候,就已經是成年人了.福音書的確沒有提供多少耶穌從嬰孩至成人之間的資料,因此,引來不少猜測,成為有趣的野史.

有一個傳說如此說:耶穌和一群兒童玩耍,有個孩子從二樓墜落摔死,其他孩子都嚇跑了,唯耶穌獨自留下.孩子的父母前來探視,只見死去的孩子和耶穌,便一口咬定是耶穌把孩子推下摔死的.耶穌雖然一再否認,他們堅決肯定是耶穌把那孩子推下摔死的.耶穌便從房頂上跳下,走到屍前高叫道:則農(孩子的名字),起來,告訴他們是不是我把你推下來的.那孩子立時起來說:你沒有把我推下來,你把我叫來了.大家都目瞪口呆.(「多默福音擇要」).

另一個傳說:有一次,耶穌叫一棵棕櫚樹彎下來輕撫媽媽.又有一次,耶穌用泥土製作了十二隻燕子,然後拍拍手,十二隻燕子就飛走了.

還有其他的傳說,都是神奇有趣,娛樂性豐富.在新約聖經的四卷福音書中,只有路加福音有一個小段落記述了耶穌從嬰孩至成人的其中一個片段.這個唯一的片段所展示的耶穌是一位怎樣的耶穌呢?與這些傳說相比,又有甚麼不同呢?路加福音所展示的這位問道的小孩的形象,對我們有甚麼意義呢?

二．    經文

41每年到逾越節,他父母就上耶路撒冷去.

42當他十二歲的時候,他們按著節期的規矩上去.

43守滿了節期,他們回去,孩童耶穌仍舊在耶路撒冷.他的父母並不知道,

44以為他在同行的人中間,走了一天的路程,就在親族和熟識的人中找他,

45既找不著,就回耶路撒冷去找他.

46過了三天,就遇見他在殿裡,坐在教師中間,一面聽,一面問.

47凡聽見他的,都希奇他的聰明和他的應對.

48他父母看見就很希奇.他母親對他說:「我兒!為甚麼向我們這樣行呢?看哪,你父親和我傷心來找你!」

49耶穌說:「為甚麼找我呢?豈不知我應當以我父的事為念嗎註?」

50他所說的這話,他們不明白.

51他就同他們下去,回到拿撒勒,並且順從他們.他母親把這一切的事都存在心裡.

52耶穌的智慧和身量註, 並上帝和人喜愛他的心,都一齊增長.

三．    釋經與應用

1．    問道的小孩令人:驚喜有時,擔憂有時

46他在殿裡,坐在教師中間,一面聽,一面問.

47凡聽見他的,都希奇他的聰明和他的應對.

一位十二歲的小孩,不是在聖殿外玩耍,不是在路上追逐,不是旁觀著成人的討論,而是在聖殿中與聖殿教師切磋交流,聽得明白之餘,還能問得合適,如此的應對能力,令飽讀經書的聖殿教師都驚喜萬分,寵愛有加,甚至讓這位問道的小孩坐在他們中間,有與他們平起平坐的意味.這位問道的小孩教聖殿教師驚喜的同時,卻也教他的父母擔憂.讓我們從另一個鏡頭看看.

41每年到逾越節,他父母就上耶路撒冷去.

42當他十二歲的時候,他們按著節期的規矩上去.

43守滿了節期,他們回去,孩童耶穌仍舊在耶路撒冷.他的父母並不知道,

44以為他在同行的人中間,走了一天的路程,就在親族和熟識的人中找他,

45既找不著,就回耶路撒冷去找他.

他母親對他說:「我兒!為甚麼向我們這樣行呢?看哪,你父親和我傷心來找你!」

節期過後,拿撒勒的鄉民一同回鄉.叔伯兄弟,路上互相照應,十二歲的耶穌沒有跟在父母身旁,不是甚麼奇怪的事情.然而,整天都不見兒子的蹤影,耶穌的父母要尋找兒子了.第一天找不到,第二天找不到,稍為代入,都感受到耶穌的父母的擔心和焦慮,也難免胡思亂想,往往越想越傷心,越想越擔憂.

這位孩子耶穌,有令人驚喜的時候,竟然也有令人擔憂的時候.他像誰呢?像我們的孩子,像我們每一個人的孩子階段.

2．    問道的小孩:聰穎有時,輕率有時

46坐在教師中間,一面聽,一面問.

47凡聽見他的,都希奇他的聰明和他的應對.

49耶穌說:「為甚麼找我呢?豈不知我應當以我父的事為念嗎註?」

這位問道的小孩與聖殿教師的談話中,顯出他的聰明及應對能力,令聖殿教師也大開眼界；這位孩子耶穌回答父母的詰問時,也似乎大有道理.對於敬畏上帝的人來說,豈能隨便否定「以天父的事為念」這崇高的心意呢?然而,這位令人激賞其聰穎的孩子耶穌,卻同時又令人苦惱其輕率.

43守滿了節期,他們回去,孩童耶穌仍舊在耶路撒冷.他的父母並不知道,

無論刻意不告訴父母也好,忘記告訴父母也好,樂而忘返也好,在父母不知情之下,獨自留在距離家鄉幾天路程的耶路撒冷,怎樣說也不是令父母提心吊膽的行為.孩子耶穌難道想不到父母會為他擔心的嗎?體諒不到父母的關注及感受的嗎?

這位一方面聰穎,說起話來頭頭是道,同時又輕率,沒有兼顧父母的緊張及感受的孩子耶穌像誰呢?像我們的孩子,像我們每一個人的孩子時期.

3．    問道的孩子令人:激心有時,窩心有時

43守滿了節期,他們回去,孩童耶穌仍舊在耶路撒冷.他的父母並不知道,

他母親對他說:「我兒!為甚麼向我們這樣行呢?看哪,你父親和我傷心來找你!」

49耶穌說:「為甚麼找我呢?豈不知我應當以我父的事為念嗎註?」

兒子失蹤了幾天,相信約瑟及瑪利亞已經擔心及焦急得不得了.終於找到耶穌了,母親於是對他說:「我兒!為甚麼向我們這樣行呢?看哪,你父親和我傷心來找你!」.瑪利馬在此時如此說話,相信天下父母都認為是合理及適當的.然而,孩子耶穌竟然反問:「為甚麼來找我呢?」如果約瑟和瑪利亞如我們一般的情商(EQ),恐怕已經被氣得兩眼冒火,鼻孔噴煙.約瑟和瑪利亞果然有過人之處,他們不明白孩子耶穌所言所行的意思,卻把這些事藏在心裡,沒有大動肝火.

51他就同他們下去,回到拿撒勒,並且順從他們.

雖然這位孩子令父母激心,但卻馬上跟隨父母回拿撒勒的家鄉.路加福音描述孩子耶穌整體來說是順從父母的,孩子對父母順從,最能夠令父母感到窩心.

這位有時令父母感到激心,有時又令父母感到窩心的孩子耶穌像誰呢?像我們的孩子,像我們每一個人的孩子時期.

4．    問道的孩子:身體成長有時,心靈成熟有時

52耶穌的智慧和身量註, 並上帝和人喜愛他的心,都一齊增長.

路加福音指出,孩子耶穌不是一夜長大,他的身體按著自然的速度發育,他的智慧也從生活經驗中累積起來,他很可能也需要透過待人接物的成熟態度去獲得人們的喜愛,透過敬畏和順服得到上帝的稱許.

這位身體成長需時,心靈成熟需時的孩子耶穌像誰呢?像我們的孩子,像我們每一個人的成長過程.

5．問道的孩子令人感到驚訝:像人多於像神到底是偽經及傳說所描述的孩子耶穌像神的兒子,還是路加福音所記載的孩子耶穌像神的兒子呢?面表看來,當然是偽經及傳說所描述的,因為偽經及傳說描繪出一位擁有超凡能力的神奇小子,路加福音反而記載著孩子耶穌如鄰家的小孩一般地平凡.上帝來到人間,人們左看右看,前看後看,他都不像神,反而與一般人無異,這或許可以解釋為甚麼耶穌那個時代的人對他感到迷惑,甚至拒絕他.耶穌的同鄉,看著耶穌成長的拿撒勒人對此最有共鳴.耶穌長大出外傳道,有一次回到拿撒勒的家鄉,遭到自己同鄉輕薄地對待:

1耶穌離開那裡,來到自己的家鄉；門徒也跟從他.2到了安息日,他在會堂裡教導人.眾人聽見,就很驚奇,說:「這人哪來這本事呢?所賜給他的是甚麼智慧?他手所做的是何等的異能呢?3這不是那木匠嗎?不是馬利亞的兒子雅各,約西,猶大,西門的長兄嗎?他姊妹們不也是在我們這裡嗎?」他們就厭棄他.(馬可福音6:1-3)

耶穌在外面傳道時顯出的智慧和異能沒有受到同鄉熱烈歡迎,反而令他們質疑及感到厭煩.拿撒勒人可能會互相議論說:「耶穌不可能是神.如果他是神,那次他怎會如我的兒子一樣患病呢?如果他是神,那次他怎會在切割木材時弄傷自己呢?如果他是神,他為甚麼不變多些錢給父母,改善家庭生活呢?如果他是神,為甚麼他不是一夜長大,反而如我們的孩子一樣逐步成長?如果他是神,為甚麼不能叫所有人都喜歡他?如果他是神．．．總之,耶穌不可能是神,神絕對有能力令自己不受身體及時空的限制,又有能力免除自己於各種疾病及人際的困擾,而這位耶穌卻如我們一樣承受著這些限制及困擾.」

到臨死的一刻,耶穌也令人質疑他作為上帝的身份.百姓和官長都嘲笑他說:「他救了別人；他若是基督,上帝所揀選的,可以救自己吧!」(路23:35)與他同針十字架的其中一個犯人也說:「你不是基督嗎?可以救自己和我們吧!」(23:39)這些人的意思是:「耶穌,你雖然曾經行神蹟幫助過人,但你的神能不單止沒有為自己帶來甚麼好處,甚至不能為自己解除困境.我們綑綁你,你不能掙破；我們打你,你一樣痛；我們刺你,你一樣流血；我們把你打十字架,你不能自己跳下來.你只不過是如我們這些血肉之軀一樣的人而已.耶穌,其實我們曾經對你心存盼望,但你令我們失望,我們以為你來人間做神,但你竟然來人間做人.人間已經太多人了,不需要多你一個.人間需要的是一位神,一位可以幫助我們突破作為人的限制,免除人的軟弱,解除自然環境,人間疾苦的神.」

5．    神成為人,鼓勵人成為人神竟然成為人,令所有人都感到難以致信.神成為人令人難以致信,因為有些時候,連人自己都認為做人實在太多限制,太多煩擾,太辛苦,連人自己都不喜歡做人,不想做人.

「唉 今日返學..心情好down 唔想做人... 今日返學派卷 跟住卷1唔合格...佢要我drop科之後重派埋drop科紙比我但係我想讀落去 佢重話遲下會見我家長...打黎屋企...唉..我咁廢...唉..」

「我失業到唔想做人我今年24歲 失業了3個月　我覺得年青人搵工好難　讀展翅同毅進計劃又係垃圾　又冇用　讀完冇野學 間公司又唔請你」

「我同人吵架了,,,　仲要在GAME係咁SD信比我....　大家一定要相信我...完來是我最好朋友都心淡了....　睇來我要改死主意...　我SD左封SORRY信佢都唔肯完X我...　心淡了...　想自殺了...　唔想做人了...」

「唔想做人,感覺做人好辛苦,好沈重,生活好大壓力,做人冇咩快樂,情緒好down」

(節錄自網上討論區留言)

人不想做人,因為作為人,存在著太多限制,也有太多事情不能如願,不能控制.有時考試成績不如願,有時人際關係不能控制；有時工作不如願,有時身體狀況不受控制；有時計劃不如願,有時情緒不受控制．．．還有太多的不如願及不能控制的狀況.有些人不禁發出慨嘆:「做人太麻煩了,太疲累了.」

對於人們不想做人的慨嘆,對於人們對自己的限制及處境諸多投訴,上帝沒有義正詞嚴地指責,沒有咄咄逼人地說服人做人的意義,甚至沒有拿出波板糖來引誘人繼續做人.上帝的回應,是道成肉身,自己成為人.作為人,祂沒有運用特權免除自己遇到的阻撓,沒有運用能力取消自己血肉之軀的自然感覺及反應,沒有運用能力去控制別人對他的態度.他如眾生一樣,需要懷胎十月才出生,需要照顧及時間才長大,心靈受傷會哭泣,身體受傷會流血,甚至免卻不了有人稱之為「大限」的死亡.

腓立比書2:6-8如此描述耶穌所經的路:

他本有上帝的形象

不以自己與上帝同等為強奪的

反倒虛己

取了奴僕的形像

成為人的樣式

既有人的樣子,就自己卑微

存心順服,以至於死

且死在十字架上

當人有時不想做人的時候,上帝不是來教人如何成為神,而是來教人如何做人；當人有時不想接受作為人的各種限制及無能為力的時候,上帝不是來教人如何變超人,而是來與人一齊經歷各種的限制及無能為力的感覺；當人處身死亡的恐懼之時,上帝不是來派發長生不老藥,而是來親自死在十字架上.道成內身的上帝,被稱為「以馬內利」,意思就是「上帝與我們同在」.上帝與我們同在,不再是概念,而是事實.上帝作為人的時候,也接納了人的各種限制及不能自控的際遇,並示範給人們看如何在這眾多的限制及不確定中按著上帝的旨意生存,我們又豈需要感到不忿及害怕呢?

希伯來書4:14-16節說:「

我們既然有一位已經升入高天尊榮的大祭司,就是上帝的兒子耶穌,便當持定所承認的道.因我們的大祭司並非不能體恤我們的軟弱.他也曾凡事受過試探,與我們一樣,只是他沒有犯罪.所以,我們只管坦然無懼地來到施恩的寶座前,為要得憐恤,蒙恩惠,作隨時的幫助.」(希伯來書4:14-16)

上帝成為人,他對我們的處境的明白,不再是想像,而是親身經歷,所以,希伯來書告訴我們,他能體恤我們的軟弱,鼓勵我們坦然無懼地來到施恩的寶座前,為要得憐恤,蒙恩惠,作隨時的幫助.

四．    總結

剛過去的星期日,學生團契參與「路向四肢傷殘人士協會」的冬至探訪活動,代表協會走訪了二十多位四肢傷殘人士,聆聽了二十多個學習生存及學習生活的感人故事.

其中一位朋友在游泳池邊做熱身運動時,失去平衡,飛身跌入池水中.從此,頭部以下的身體就失去知覺.經過一段時間留院醫治,差不多要出院的日子,他猶豫著以後怎樣呢?最後,他給自己答案是「世界,社會就是現實,既然死不了,那就生存下去吧!」於是,他接受各樣訓練,重投學習,經過會考,進入大專.

常言道,家家有本難念的經,各人有各人的難處.我們也不用比較誰的處境艱苦,誰的處境不困難.無論在甚麼處境,正如那位四肢傷殘的朋友所說,既然死不了,那就生存下去吧!既然生存下去,就盡力生存得豐富及精彩吧!

耶穌說:「我來了,是要叫人得生命,並且得的更豐盛.」(約翰福音10:10)願人類社會也彼此扶助,使大家都能夠生存,並且生存得珍惜生存.



\newpage

\section{仰望星空}
\label{sec:32}
\textbf{陳龍斌教授 - 仰望星空}
\newline
\newline
\hyperref[sec:31]{< < < PREV SERMON < < <}
~
\hyperref[sec:toc]{[返目錄]}
~
\hyperref[sec:33]{> > > NEXT SERMON > > >}


以弗所書 3:14
\newline
\begin{longtable}{cl}
\hline
\hline
章節 & 經文 (和合本修訂版)\\
\hline
3:14 & \begin{tabularx}{0.7\textwidth}{X} 因此，我在父面前屈膝— \end{tabularx} \\ \\
[1ex]
\hline
\hline
\end{longtable}
日期:2010年1月31日

題目:仰望星空

經文:以弗所書3:14

證道:陳龍斌教授

弟兄姊妹們早安,今天要講的經文是以弗所書一章3到14節,這段聖經可以說是在整個新約聖經裡面其中一段,思想非常偉大的,所涉及的層 面是很廣泛的一段聖經,在一章三到14節作為一個書信的開始,它也據有保羅的書信典型讚美詩的特色,而今天誦經的時間,除了誦讀以弗所書,其實弟兄姊妹也一起讀過詩篇第一百四十七篇,詩篇第一百四十七篇也是一首讚美的詩,在這首讚美的詩裡面,它是讚美耶 和華為能力,美善,慈愛,在兩方面去頌讚,一方面是頌讚耶和華是怎樣賜福給耶路撒冷,另一方面是頌讚上帝怎樣去保守,創造這個世界,其 中一首詩篇裡面,有一節的聖經,我們剛才是沒有讀的,就是第四節,詩篇一百四十七篇第4節,經文是這樣說的,它說:「耶和華他數點星宿的數目,一一稱他的名」,基於這一節的聖經,我構思了今天的講題,那今天要講的講題是仰望星空.

最近當我跟崇基神學院的同學去探訪一家教會的時候,在公禱的時間,我聆聽到一位基督徒的朋友,他在公禱的過程裡面,他四次在禱文裡面 四次的用到我仰望星空,因為我沒有在禱文裡面聽過仰望星空,以為他是自己想到這樣子祈禱,在中間那位朋友並沒有提到主的名 子,那我覺得仰望星空聽起來的感覺很不錯,也覺得蠻有意思的.

其實在2009年10月二號,慶祝我們建國六十年的晚宴上,壓軸的一首獻唱,是由張學友跟一位小妹妹一起唱的一首名叫仰望星空的歌,其 實這首歌的歌詞,是源自於我們國家總理所說的一首詩叫《仰望星空》,這首詩是這樣子說的,有一些民族有一些關注天空的人,他們才有希 望,一個民族只是關心腳下的事情,他是沒有未來的,我們的民族是一個大有希望的民族,我希望同學們經常仰望天空,學會做 人,學會思考,學會知識和技能,做一個關心世界及國家命運的人.

而這個《仰望星空》的內容是這樣的:我仰望星空,它是那樣寥廓而深邃；那無窮的真理,讓我苦苦地求索,追隨.我仰望星空,它是那樣莊嚴而聖潔；那凜然的正義,讓我充滿熱愛,感到敬畏.我仰望星空,它是那樣自由而寧靜；那博大的胸懷,讓我的心靈棲息,依偎.我仰望星空,它是那樣壯麗而光輝；那永恆的熾熱,讓我心中燃起希望的烈焰,響起春雷.弟兄姊妹們,當你聽到這首詩詞的時候,不知道你們在心裡面有什麼思想,最近我有機會到四川參加一個聖經的學術會議,可能是怕人數不足,主 辦單位給當地一些的信徒出席這個聖經的學術會議,其中有一位還蠻年長的太太,那天當我在做一個關於耶穌在生的報告的時候,耶穌所牽起的一場運動,是一場屬靈的運動,對於當時的社會來說,是一場有價值的運動.

當我讀完我的報告以後,那位太太,她多次的跟我聊天,可是當時其實我並不是很明白,那位太太所講的話,因為她的普通話夾了一些四川的口 音,當時她不段的告訴我,我每天都為溫家寶總理跟我們的國家去祈禱,當時我沒有聽得很明白,可是當我要走的時候,她給了我一 紙張,這紙張天天陪她入睡,這張紙到低是什麼呢?

這張紙就是我剛才讀給大家聽的,印在人民報上面的這篇《我仰望星空》,當她一路在聽的時候,她聽得明白我報告裡面的內容,對於總理的這篇《我仰望星空》,她也很有感觸,在面對我們國家屈起的時候,我們需要什麼呢?我們可以把太空人帶到太空上,我們也掌握真理,為 國家富強的能力,我們不在需要去面對別的國家的臉色,在這個時候,我們的國家需要什麼呢?

不久之前,一則新聞,有關一位姓劉的朋友,劉曉波先生被逮捕,有很多不一樣的時事評論人員去評論這件事情,在最近一次聚會的時候,有 一位同事這樣子說,其實這個新聞炒作了一個星期,香港人不會再去關心了,當我聽到他這樣子提及,我並沒有回應他,第一,如 果他所所描述的是事實的話,我對香港人也會覺得很可惜,可是.在另一方面來說,有不少的時事評論員也去說一些關於這件事的事情,他 們說國家在這個時間,跟昔日當那些外國的元首來探訪我們中國的時候不一樣,以前我們是要看他們的臉色,所以在那些時候我們 總會釋放一些異見人士,可是今天我們不需要了,我們國家已經掌握財富,已經掌握真理,我們國家已經有能力飛上太空,在這個時候我們的 國需要什麼呢?

在2008年,有另一位歌手名叫張杰,他也唱了一首歌曲,名叫《仰望星空》,這一天,我開始仰望星空發現,星並不遠,夢並不遠,只要你踮起腳 尖,張開雙臂,和你一起,飛的更高,看的更遠,尋找生命中最璀璨的亮點,在面對這新的一年,兄弟姐妹們,在我們的生命中,有沒有勇氣,我們要飛得更高,看得更遠,我們要尋找生命中最璀璨的亮點.

今天的經文,是會說以弗所書一章三到14節,在這個讚美的詩篇來說,其實整個以弗所書有不少屬靈的前輩也說,以弗所書是講神奧秘 的計劃,這個計劃,用一個很濃縮的字眼去表達神在創世以來直到永恆,上帝的計劃究竟是怎麼樣的,而這個計劃,對這個掌握奧 秘的作者來說,他不禁在生命裡面發出一個頌讚,他讚美主的計劃是何等的偉大.

弟兄姊妹們,在你們的生命裡,有多少次翻開或是聆聽過以弗所書的經文,特別是第一章,在回顧過去的時間,對於這一段聖經,我也很感 動,有很多個早上,當我跪下向上帝祈禱的時候,我心裡面不禁的禱告,也很懺悔,為什麼在我生命裡面,沒有像這個作者一樣,對 上帝的計劃有同樣程度的讚美呢,有時候會覺得,這些都只不過是說說而已,像在詩歌頌讚的時候,有一些教會很興盛,不停的唱哈利路亞, 其實有沒有內容的呢?可以沒有內容,可是這個作者,在生命裡面發出了一個頌讚,特別是在整個第一章,他三次的發出頌讚,在第六節裡面 說,使祂的榮耀得著頌讚,第12節,要祂的榮耀可以得著稱讚,跟第14節,使祂的榮耀得著頌讚,在整個的以弗所書有三重的 頌讚,當然每個頌讚也有它的重點所在,這三個頌讚也都呼應著整個以弗所書這三重結構,無論是緊接著整個計劃的展開之後,他有一個三重 的禱告,跟著從三個的向度去講,耶穌基督,使徒跟教會怎樣去成就這個計劃,談到新人這個角色的時候,從三個特別屬靈的角 色去演繹,最後在這個全副軍裝的時候,也有三個方面的禱詞,整個以弗所書充滿了三重頌讚的呼應.

第一方面要跟大家一起看的是,這個更高的福氣,父上帝在第三節說,在基督裡面,曾賜給我們在天上各樣屬靈的福氣,弟兄姊妹不知道你們有沒有覺得自己是一個很有福氣的人,最近在看新聞的時候,有一些人真的很有福氣,這個新聞是關於迪士尼樂園的,過往的兩年有不少的人在迪士尼樂園裡面成為了主角,因為有人向這些主角在迪士尼樂園裡面求婚,新聞報導說一百四十九對情侶,百發百中,全部在迪士尼樂園裡面求婚成功,可能是群眾的壓力吧,可是每一位女主角在迪士尼樂園裡面都成為一個夢幻的主角,因為有她最心愛的卡通人物,跟那些男士一起向那些女主角去求婚.

當時的主角都會感到自己是有福氣的人,可是其實這段聖經同樣告訴我們,其實我們是很蒙福的,父上帝為我們去準備這個福音,是經過長年 累月的安排,從創世的時候,已經揀選你和我,一定要去蒙福,是上帝的一種心意,祂要把這個福氣帶給我們,可是這種天上各樣 屬靈的福氣,是不是我們作為一個人所羨慕的呢,在這個人世間裡面,有很多人覺得自己其實是沒有福氣的,覺得自己人生也蠻坎坷的,因為 很多人也只是重視腳下的事情.

正如總理所說的,他說如果一個民族只關心腳下的事情,他是沒未來的,一般人當他們面對這個講法的時候,面對這個腳下的事情,他擔心, 徬徨,也有很多的憂慮,其實是可以理解的,因為事實上,我們活在世界上,我們知道,不是每個人都會有地上的福氣,特別是面對香港的前景,很多人最擔心的,最盼望的是在未來一年可以怎樣去舒解香港的貧富懸殊呢?

有些人在物質上很豐富,有些人不段的掙扎,他們可以享有的其實是很少的,對於一段人來說,我們需要關心腳下的事情,是非常可以理解 的,這並不是一件幼稚的事情,是一件很真實生命裡面的掙扎,正如國家一樣,先讓一些人富有起來,以後其他人也有機會富有,這 句話,我們的總理講了很久了,可是還是有很多人依然落在非常貧窮的境況,其實香港人也都是一樣的.

一些有錢的商人,常說我們要共渡時艱,可是當我們真正要去面對的時候,其實還是貧者越貧,富者越富,這一段的聖經說,在基督的裡面, 我們是體會天上各樣屬靈的福氣,基督徒可不可以在這個社會裡面去體會到我們其實是一個很有福氣的人,我們的福氣不只是單單腳下的事 情,我們的福氣是因為我們有這種屬靈的福氣,是一種另類的福氣,不只是地上的事情,當總理說,有一些民族關心天空的人,他們才有希望,總理的這一句話,我相信他心目中所想的,不希望只有一班基督徒有這種想法,可是作為基督徒的,我們是不是真的有一種仰望星空的氣魄,一種的渴望.

最近我去買了我很少機會會買的唱片,因為我自己不是很在乎明星的,可是最近因為一位歌手,我很想買來聽聽,支持她,這位歌手收錄了一首歌,歌名是《有一種快樂》,它裡面是說這一位明星,曾經去過一個很貧窮的國家,名叫老撾,去探那裡的小朋友,的歌詞是這樣的,從我投入你我不相 似的世間,便不段的慨嘆,還曾經弄濕了雙眼,這個明星很難去明白那裡的小朋友,存活為什麼會這樣艱苦的,艱苦的地步就像是 攀一些高山一樣的辛苦,然後,回望我這非必要的襯衫的時候,我換到什麼呢?換到的是讚嘆,才明白你縱使苦難不斷,微笑卻比起一加一更簡單,人生原來有這樣的一種人,她體會到能夠扶持你,能夠扶著你,是最高的高興.

弟兄姊妹,當你在思考這件事的時候,什麼是你生命中最高的福氣呢,你是希望在你的人生事業中不斷的攀一個又一個的高峰,當然這個是值得驕傲的事情,是一個作為有上進心,在這個社會不應該放棄的事情,可是除此以外,人生裡面有沒有更高的福氣呢?

另一方面跟大家一起去看的,是一個更遠的價值,在這一段聖經裡面第十四節,上帝給我們聖靈住在我當中,成為一個保障,就像一個存款一 樣,保證我們可以得到一個,從上帝為我們安排的一個基業,它是一種更遠的價值,人生裡面可能我們所追求的是很多在地上,在 眼前的基業,可是聖靈住在我們生命裡面,祂是要應許我們,給我們知道,原來生命不只是地上,眼前的一些價值,還有一些更遠的價值,值 得我們去追尋.

在以弗所書一章13節裡面說這個使人「得救的福音」,特別是基督徒所聽的福音,因為這個原因,所以聖靈住在我們裡面,可是請大 家注意第12節,在我們這些「在基督裡有盼望的人」,原來去領受同一種福音的信徒來說,不是每個人都必然有生命裡面的這個盼 望,我們都是得著聖靈住在我們生命裡面,我們基督徒是這樣子相信的,可是不是每個基督徒都必然有同樣的盼望,相信有這個價值,這 個更遠的價值,人生有很多人都是追求眼前的價值.

今天有很多人面對工作,特別是一些大學畢業生,是非常的艱難的,其實當找到一份工作,最近有一位朋友的女兒,從外國回來找暑期的工 作,在僱主來看,這個女生真的很不錯,英文很好,又很聰明,所以當僱主遇到一個那麼好的員工,當然什麼都要她幫忙,所以每天晚上11 點才可以下班,有很多人為生命中一些價值去勞碌,很難找工作,這個也是國內很多年輕所要面對的問題.

最近在網上的一段報導,有一些學生,因為要面對畢業,有些年輕人這樣子說,願意用他兩年的青春來換一份工作,我心想,這樣值得嗎?可以換一份怎麼樣的工作呢?今天在80年代後的年輕人,在香港要置業也是很困難的,可是聖經告訴我們,我們並沒有辦法可以立刻去解決這個人生基業的問題,在聖經裡面提醒我們一件事情,聖靈住在我們心裡面,給我們去渴望一種永恆的基業,如果真的像基督徒所 說,上帝是三位一體的,這個聖靈住在我們生活裡面,上帝愛我們,所以把聖靈放到我們生命當,祂常常的感動我們,告訴我們,聖 靈提醒我們,要有另外一種價值.

香港在97年回歸以後,社會有很多大大小小的風浪,其實香港人是蠻苦幹的,我們經歷過SARS,我們經歷過金融風暴,我們也經歷 過金融海嘯,可是很可惜,在面對社會危機的時候,有多少香港人知道應該去追求人生一些另類的價值呢?這種另類的價值,是一 種可持續的價值,是一種有意義的社會價值,不知道有沒有弟兄姊妹覺得要面對未來一年裡面,我們除了要面對累積財富之外,我們有沒有想幫助我們自己或是幫助這個社會去追尋一些可持續的價值呢?

昔日耶穌的運動,是一種有價值的社會運動,可是今天有多少人在這個社會裡面,累積財富的時候,也想到這些可以持續的價值呢?其中有不 少香港的基督徒,盼望透過社會企業的運動,為香港人帶來一些消費的反思,我們活在一個消費主意的社會裡面,其實有很多的 消費都是換來不必要的襯衫,可是有沒有想過,用我們的消費去累積一個更長遠,一個永恆,值得上帝去加許的基業,換來一種對社會更有意 義的價值呢?

最後一點,跟大家一起去思考的,是一種更璀璨的愛,有沒有想過,有什麼值得我們去犧牲的,聖經裡面在第5節說,父上帝「愛我們」,第7節說「藉這愛子的血,得蒙救贖」,真的很難理解,有時愛到匪夷所思,聖經裡面說,祂愛我們,愛到把祂所愛的兒子,一個真正 寶貴的兒子,來跟我們這種生命去交換,父上帝是瘋了嗎?我相信是的,換我們值得嗎?如果上帝問我們,像是約翰福音二十一章裡面說,你 愛我比這些更甚麼?如果有一天我們要去面對上帝這樣問我們的時候,我們生命裡的這些是什麼呢?

面對這個世界,有很多人,是很願意去作出犧牲的,特別是伊斯蘭教徒,當我去德國留學的時候,跟我最要好的朋友是一位巴基斯坦藉的男 生,我對他們的宗教信仰也心全一種敬畏,可是今天有一些比較極端的伊斯蘭教的教徒,他們願意去追尋生命裡面他們應為的價 值,他可以流血,他可以犧牲自己的生命去換取他們自己想要追尋的事情,今天我們稱這些恐怖主義,可是他們是願意去犧牲他們自己,在 這個社會裡面,當他們爭取不到的時候,他們沒辦法去發表他們的意見,所以有些人會用這種暴力去表達他們對強權社會的不滿,用 戰爭去入侵別的國家.

昔日以弗所書所面對的處竟,羅馬的帝國就是用軍事去建立整個國家,不少猶太人或是其他民族要去抗衡這些意識形態的時候,也 都走上戰場去抗衡這種文化,可是以弗所書裡面所提到的一個屬靈軍裝,作為以弗所書的結束,這個結束,今天有不少的教會,特別 是一些比較靈恩的教會很著重屬靈的爭戰,可是當我們把這一段聖經放在當時的處境去理解的話,面對羅馬帝國真的走上戰場的一種意識形 態,面對要去抗衡這種世界的意識形態,走上戰場的一種革命,這個是在宣揚另外一種價值觀,宣揚為世人持續的一種福音,一個 基督為我們作出的犧牲,我們要穿上全副軍裝,目的是要宣揚這個寶貴的福音.

最近有機會參加社會企業的會議,裡面有一對夫婦,他們每年的年假在星期三放假,因為他們想要配合一個企業的活動,這裡的弟兄姊妹出錢又出 力的去陪伴一些有問題的青少年,那些更新的人士,希望不只是用嘴巴跟他們傳福音,也是表達對他們的關懷,有時候這樣子真的會換來好的 結局嗎?其實是不一定的,當年青人做得不對的時候,當這對夫婦想要教他們騎腳踏車的時候,當中有一位年青人學了三個小時也 學不會,如果怎樣呢,他連續三個小時去做同一事情,就是不斷的用髒話去罵她,想想看,如果你是那一位姐妹,你會怎麼樣呢?可是我把信這對夫婦做到了我們剛剛所讀的經文,什麼是「道成肉身」,耶穌基督來到這個世間,充充滿滿的有恩典,有真理!

聖經詩篇裡面的第八篇最後是這樣結束的,裡面也有一首仰望星空的詩,可是它裡面問了一個很重要的問題:「耶和華我們的主阿,你的名在全地何其美...人算甚麼,你竟顧念他．世人算甚麼,你竟眷顧他...」耶和華所造的全地何其美,我們 算什麼,你竟顧念我們,人本來很微小,耶穌基督卻來到這個世間去拯救我們.

弟兄姊妹,面對新的一年,我們的盼望,我們的人生是什麼呢?願主賜福祂的話!



\newpage

\section{上好的福份}
\label{sec:33}
\textbf{戴國璋牧師 - 上好的福份}
\newline
\newline
\hyperref[sec:32]{< < < PREV SERMON < < <}
~
\hyperref[sec:toc]{[返目錄]}
~
\hyperref[sec:34]{> > > NEXT SERMON > > >}


日期:2010年2月7日

題目:上好的福份

證道:戴國璋牧師

實在是很高興今天能夠在大家當中,與大家一起共度這早上的優美時間.回想起上次和大家見面的時候,似乎還是不久之前,因為也還是零九年年底的時候.只是在短短的一段日子裡,我們又有機會在主內再見.

在我上一次來到大家的中間的時候,我選用了一個故事,不知道大家是否還記得,就是有關《聖經》裡所傳頌的馬利亞用香膏抹耶穌的一個故事.我們看到愛耶穌的這個馬利亞,她怎樣來將自己所有的一切,全然的奉獻在耶穌基督的身上,服事在耶穌基督的身上.她所作的,被耶穌稱讚成一件美事.我們從當中學習我們可以怎樣透過服事耶穌,以至得到耶穌基督這個成為美事的稱讚.對於今天再和大家在這裡見面的時候,我感到有點意猶未盡,我覺得還可以再讀一些《聖經》裡談到的這位馬利亞的故事,所以我們便選擇了記載在路加福音內的這個簡單的記述,講述馬利亞怎樣坐在耶穌的腳前聽道,被耶穌稱讚她選擇了「上好的福份」.在服事上她所作的是美事,在她選擇的是上好的福份,實在值得我們來思想這位姊妹留下給我們的見證.

或者值得我們在這裡簡單介紹一下馬利亞和她的兄弟姊妹這個家庭,因為這個家庭在《聖經》內和耶穌基督有著一個非常親近的關係,尤其是在約翰福音內的記述告訴我們,他們不只是和耶穌基督十分熟絡,而且他們每一位都是耶穌基督所愛的,甚至他們每一位也都是深深地愛著耶穌的.這個家庭其實是十分簡單的,《聖經》內告訴我們,馬利亞這個家庭只有三位成員,就是她和她姐姐馬大(亦即是馬大,馬利亞兩姊妹),另外還有一個兄弟(雖然我們不知道這位究竟是哥哥還是弟弟,在《聖經》中是很難去分辨是哥哥還是弟弟的),我們知道他的名字叫拉撒路.有關拉撒路的事,我們最熟悉的,便是在後來他因為某些原因離開了世界,後來耶穌基督以他的大能大力讓他從死裡復活,施行了這個神蹟.他們這個家庭,其實是住在耶路撒冷城以東的橄欖山的東坡上(即是不是向著聖殿山的那邊,而是向著死海那邊)的一條叫伯大尼的小村莊.這條小村莊距離耶路撒冷城並不遠,大概只有兩公里,一點也不遠(我自己居住在旺角,兩公里的路程,其實連尖沙嘴也還未到,所以是一個很近的地方).耶穌和這一家人有著很親密的關係.所以有很多時候,耶穌就好像一般的猶太人一樣(我們知道一般的猶太人要上耶路撒冷去過節,猶太的男性全都有這一責任,一年三次),耶穌也會經常履行這個責任.當耶穌有機會北上往耶路撒冷過節的時候,我們發現到,他很多時候都會選擇寄住在馬大,馬利亞,拉撒路的家中,接受他們的款待.甚至我們知道,當《聖經》講述到耶穌基督在耶路撒冷度過最後一星期的時候(我們知道這個星期叫Passover Week,受苦的一星期),其實每一天耶穌在耶路撒冷城中很忙碌應付各種事情,他安歇的時候,他總是會選擇離開耶路撒冷,在耶路撒冷東面,橄欖山的東坡上的這條小村莊-伯大尼村,馬大,馬利亞,拉撒路的家中去歇息.從此我們可以知道,這個家庭是怎麼樣和耶穌基督有著一個親近關係,這個家庭是蒙福的,因為他們全家都和耶穌基督建立著一個良好的關係,而且成為耶穌基督所愛的家,是一個接待耶穌基督的家庭.

我們剛才讀到的故事,路加福音第十章三十八至四十二節,就是記述在某一情況之下,耶穌又來到了這家庭；耶穌又接受了這個家庭的款待.在這個款待的過程當中,妹妹馬利亞十分殷勤的服侍耶穌.我們知道,在猶太的社會裡面,服事的工夫(尤其是接待客人)很多的時候都會由女性負責的；男士(譬如說,該家庭的男士或主人)很多時候的主要工作就是陪伴客人,陪伴著他,與他聊天,令他有一種賓至如歸的感覺.而那些雜務,例如鍋碗瓢勺,預備食物,這些事情則通常由婦女負責.馬大覺得她在這個家庭是一個大姐姐,又是一個婦女,所以她好像當時那些接待耶穌或接待朋友的人一樣很殷勤的接待耶穌.《聖經》告訴我們,這個接待是多麼的忙亂,以致於有煩擾的出現.其實,接待的工夫真的一點兒也不容易做,因為有許許多多瑣碎的事情要處理,所以這過程中發生忙亂和煩擾,其實不足為奇.在這過程中,馬大的心裡其實有一點點的埋怨:她埋怨的地方就是,為何這裡只剩下我一個去為服侍耶穌基督的事情忙亂呢?我不是還有一個妹妹嗎?妹妹不是也應當參與在這服侍中嗎?所以,她有一點點的埋怨.我們可以從第四十節裡馬大所說的話中,感到她有些微的怨言.她向耶穌說:「主阿,我的妹子留下我一個人伺候,你不在意麼?」她沒有將這怨氣直接地發洩在她妹妹的身上,反而,她用繞圈子的方式,來到耶穌的跟前,向他投訴.她覺得:「耶穌,我們的夫子,我的妹妹不應該只坐在你的腳前,甚麼也不用做.你是不是應該叫我的妹妹進來,去減輕我的忙亂呢?」從馬大的這話中,我們感受到她對這光景,對這事頗有微言.不是只是頗有微言,我們從馬大的話中,感受到她所表達的一種情感,就是她覺得她是有道理的,覺得她所說的是對的:接待耶穌是一個如此重要的工夫,不應該只是由她獨力完成,她的妹妹也應該參與在其中.我們可以感受到這個表達,但是為何會有這個表達呢?其實我們知道,在猶太人的傳統裡(這一家人和耶穌都是猶太人),猶太人的傳統是將接待客人視為他們天職的一部份.換言之,他們會覺得接待客人是一件非常重要的事情.所以,每當有客人來到他們當中的時候,他們都會竭盡所能去接待這位客人,令客人有一種賓至如歸的,像回到家的感覺.這是他們的習慣.這一個傳統,其實已經有許多年的歷史了,所以這個習慣,在某程度上其實已經成為了猶太文化的一部份.希伯來書十三章二節這段《聖經》提醒早期和現在的基督徒們,其實我們是有些我們應該跟隨的,優美的傳統.這段經文說:「不可忘記用愛心接待客旅；因為曾有接待客旅的,不知不覺就接待了天使.」希伯來書第十三章是如此說的.

希伯來書的作者的想法,令我們想起了《聖經》內的兩個故事:第一個故事是在創世記十四章裡,記述亞伯拉罕是如何接待天使的(但亞伯拉罕所接待的,其實是耶和華).如果我們還記得那故事的話,我們知道當時正是亞伯拉罕在帳棚外乘涼的時候,他忽然看見數人,於是亞伯拉罕便向他們走去,向他們下拜,並且邀請這些人到自己的帳棚內,還給他們預備了豐富的筵席.不過我們知道,他所接待的其中一人是耶和華.後來,耶和華還向亞伯拉罕作出應許,應許他在新的一年裡,他將會有一個兒子出生,而以撒就是這孩子的名字.我們記得這個故事.可能這個作者所擁有的,其實就是和這個故事有關的傳統.我們又記得,在創世記第十六章的時候,就記述到有關所多瑪和俄摩拉被毀滅的事情.在那裡,故事告訴我們,羅得在城外接待了兩個客旅,後來才發現原來這兩位是天使.他們要將這「滅命」消息帶往所多瑪和俄摩拉,和通知這個義人羅得,要將羅得和他的家族救離這個毀滅的地方.

所以,希伯來書作者的這個想法,其實是利用了這個典故,然後提醒我們也應該接待客旅,因為在接待客旅的過程中,你可能會不知不覺間接待了天使,而說不定就是因為這緣故,你會有機會蒙福.所以,在基督教的傳統中,我們會發現我們其實是有一個接待客人的傳統,而在《聖經》中亦有不少有關這方面的吩咐和教導.在早期,有一本叫做《十二使徒遺訓》的書籍,內裡就有很大的篇幅提醒信徒們,有關接待客旅的各方面的提醒和規矩.但是,我們知道,這個在基督教早期裡面十分重視的一個(接待客旅的)傳統,其實是源自猶太人的,因為猶太人覺得在接待客旅的過程中,可能會接待了天使,這已經成為了他們文化的一個部份.

我記得以前我在神學院的時候,我曾聽過一個老師的分享他的一些很有意思的,印象深刻的經驗.我這個老師在攻讀博士學位的時候,就是研修《舊約》的博士學位的時候,他在美國選擇就讀一所猶太人的大學,並在那裡完成研究的工作.因為猶太人在研究舊約聖經有著很高的造詣,所以在那裡很多教授都是猶太人.他自己身為一個香港人,來自遠方的香港,在那大學裡可能也有著為數不多的香港人進修,他就是在這麼一個環境下讀書.當有節慶的時候,客在異鄉是會帶給人一種十分孤單的感覺(因為身邊的人都不在),而且節日的時候,總是會有著一些假期,而這些會令人覺得更孤單.當他在那裡就讀的時候,有一位教授十分關照這些從異鄉來的同學們,因為這位教授知道這些同學們當節慶的時候,他們會躲在宿舍中,不知道做甚麼才好.所以,每當節日來臨的時候,這教授便會邀請這些外來的,外地的學生到他家中,與他的家人一起度過這個節日.所以,我的這個老師經常都很感激這位教授所給予他的家庭溫暖(這倒是蠻好的).如果我們教會中有些單身的弟兄姊妹,「每逢佳節倍思親」,在節日裡他們也會感到十分孤單的.如果我們已經擁有自己的家庭,而我們又不介意的話,其實我們是可以邀請這些單身的弟兄姊妹與我們一同度過這些佳節,讓他們有一種溫馨的感覺,其實這是十分優美的.在過程中,我覺得我們是真的會有所學習的.

所以這位老師就對這些留下了很深的印象,因為這個大學教授時時刻刻都像父親一樣照顧他們,讓他感到有家庭溫暖.後來這位老師要畢業了,畢業之後他便要返回香港,他覺得這位教授讓他吃了這麼多頓免費的飯,他在離開之際是很應該作東邀請這位教授吃一頓豐富的.於是,他便邀請這位教授,希望作東請客,和這位教授吃一頓飯.由於他是一位單身人士,還沒有建立家庭,所以他選擇了那裡一間十分著名的中國餐館,預備請這位教授和他的家人吃一頓飯.這位教授和他的家人一同到來,所有的人都很愉快地用餐.當菜餚上桌的時候,大家便很快樂地一邊談天,一邊用餐.其中一道菜餚是蒸鵝(這個蒸鵝對中國人是再也熟悉不過了,但是對外國人來說,這道蒸鵝是一道頗新鮮的中國菜),這位教授便挾了一塊鵝紅來吃.吃在口中的時候,他覺得它的味道頗特別,於是他便問這位老師這些呈塊狀的果凍狀的東西是甚麼來的.我的這位老師便直接和那教授,就直接翻譯說那是鵝紅,goose red.甚麼叫鵝紅呢?若是你和我的話,當然一下子便會知道是甚麼回事,但是那猶太的老師,卻是不知道甚麼是鵝紅,甚麼是goose red.我的這位老師便說這其實是鵝的血.但正當這句說話脫口而出之際,我的這位老師冷不防突然記起,糟糕!為甚麼?因為,原來猶太人是不吃動物的血的,對不對?他們對血是有所禁忌的.他知道這下子麻煩了,因為他一時忘了這個事實,而只顧介紹一些比較有特色的中國菜,並沒有留意到中菜中有些食物是猶太人所忌諱的.不過事到如今已是沒有辦法了,這位老師只好硬著頭皮向這教授解釋這是鵝的血.他說他這位教授十分大方,並沒有因為這位老師令他犯了禁忌而立刻大發雷霆,反而十分得體地「噢!原來如此.」不過,這位教授由始至終都沒有再碰這道菜一下.

這個故事也是蠻有趣的,所以在我的腦海中留下了很深刻的印象.這故事一方面讓我們看到其實文化差異是確實存在的,不過,更重要的是在這故事的背景中,是有所表達的.所表達的,就是猶太人其實真的很好客,所以當有從遠方來的中國的同學的時候,他們會覺得與其這些同學孤單一人,不如讓他們到我們的家,一同享受家庭的溫暖.

因為猶太人有這個傳統的關係,所以馬大覺得接待客人,接待耶穌是一件很重要的事情,怎麼可以只有我為這些事情忙亂呢?我的妹妹動也不動,其實是不應該的.所以馬大會覺得是理直氣壯,要求耶穌吩咐她的妹妹幫忙做這些事情,做這一個服事的工夫.不過我們知道,路加透過耶穌的話把這故事記述下來,其實正是要將她們作一比較,將馬大和馬利亞作一比較.他所希望突顯的,其實是接待固然是一個很重要的,美好的服事,接待耶穌是一件更加重要的東西,而耶穌從來都不覺得接待祂是沒有用的.不過,在這一刻,耶穌覺得選擇在祂的面前聆聽上帝的道理,和在體力上為接待而勞動,煩擾,前者其實遠比後者來得重要.

事實上,我相信馬利亞並不是甚麼也沒有做的,馬大的話也留下了一些線索,馬大說:「主阿,我的妹子留下」.既然是「留下」,可能是在開始之前,馬利亞也有幫忙一同接待耶穌的事務,不過當耶穌甫開始講道的時候,馬利亞便立即將這些事務放下,坐在耶穌的面前,細心聆聽耶穌的講道.耶穌說她選擇了一個上好的福份.其實,能夠細心聆聽,是一件很舒服的事情.細想大家回到教會,大家坐在這裡,享受美麗的音樂和詩歌,享受能讓我們感動的見證,又能享受講道的人所準備的道的筵席,一同的分享.我們就是坐在那裡靜靜的聆聽,靜靜的享受,你想在我們的心靈裡,那是多麼大的福份.再且,我們在上帝的真道裡,在不知不覺間我們可能會聽到一些能夠警惕我們生命,影響我們和上帝關係的東西,就好像我上次和各位見面的時候提出了的問題:為何後來在耶穌基督受苦的這幾天中,我們在《聖經》中完全讀不到任何有關於這位馬利亞的事呢?

因為馬利亞在很早的時候,就已經知道耶穌基督要離開他們,所以她事先已經用香膏抹耶穌,在上次的故事已提及了.這表明了馬利亞在平常細心聆聽耶穌的話語裡,她聽到了一些特別的事情,而這些事情是令她有很大的得益的.

我們亦很重視在教會中同心的服事,因為若不是有這麼多人將他們的好東西拿出來奉獻給耶穌的話,教會是不可能成長的；而我們的崇拜亦當然不可能有這麼精彩美好的內容.不過,除了我們眾多的奉獻外,聆聽其實也是很重要的.

另一方面,讓我們亦留意路加福音所記述的這個故事所在的位置.在路加福音中,這故事是放在第十章裡,剛才開始的時候我不是這樣說嗎?我說馬大,馬利亞和拉撒路這個家庭在哪裡?他們是在耶路撒冷城以東的橄欖山的東坡有一條叫伯大尼的小村莊.換言之,在講述這故事的時候,耶穌已經前往耶路撒冷了.不過,親愛各位,如果我們細看路加福音第十章附近的時候,若按照耶穌當時的行程,其實他當時是應該在往耶路撒冷的途中,而並未抵達耶路撒冷.路加福音九章五十二節告訴我們,當時耶穌還在撒瑪利亞的地方,還有許多的路程才到耶路撒冷.在之後所講述的故事裡,仍然是集中述說在途中的故事.所以,當路加記述有關馬利亞在這裡聆聽上帝真道的故事的時候,他其實並不是按照事件發生的時序來編排.如果是按時序來編排的話,那麼這故事應該出現在更後期的地方.如果耶穌當時還在路途中,或是還在撒瑪利亞的話,那麼為何路加在此處加插一個在耶路撒冷所發生的故事呢?其實,他是想透過這個故事來突顯某些東西的.在這故事前所講述的,究竟是甚麼故事呢?在這之前記述的,其實是另一個著名的《聖經》故事,好撒瑪利亞人的故事,對不對?因為我們幼時已經在教會聽了許多,所以我在此也不重覆這故事了.這其實是耶穌在某一場合裡說有關好撒瑪利亞人的故事.將這個故事和好撒瑪利亞人的故事放在一起,其實路加可能是別有用心的.好撒瑪利亞人的故事的焦點是在服侍之上,當然,在那故事中的的焦點,是上帝對我們的服侍,對不對?上帝對我們的服侍,就好像撒瑪利亞人一樣這麼徹底.不過另一方面,這好撒瑪利亞人的故事,是在教導我們對我們的鄰舍也應有這種愛心.不過,當我們讀這故事的時候,可能我們會很容易產生誤會,可能會誤解救恩(或神的福音)原來是十分著重於自身的良好行為,因為好撒瑪利亞人的故事所強調的,其實就是我們要有良好的行為.可能是因為這原因,路加覺得需要另一個故事,於是他就在這之後放進了這個故事作為比對,再提醒我們,我們在上帝面前最大的恩福,其實並不是靠著行為而得著的,而是靠著我們在上帝面前坦然領受祂給我們的福祿,福份,就像馬利亞一樣,她就是坐在那裡聆聽,領受上帝給她的福份.從此,我們更能看見聆聽,坦然接受上帝給我們的福份,其實是多麼重要和珍貴.

在一次偶然的機會下,我在一間教堂看到一張單張,單張的標題是「多上教堂人長壽」.當我看到這單張的時候,覺得蠻有趣的:多上教堂居然可以令人長壽?於是我便取了一張這樣的單張,打算細心看看怎樣多上教會會令人更長壽.在這裡,我想和大家分享一些這單張上所記的資料.據這單張記載,美國有一研究指出上教堂的和熱心參與教會事務的人的壽命是較長的,而他們的身體亦會較為健康.原來,這個研究調查了四千四百名北卡羅萊納州的居民,發現這在過去六年內,有一千七百七十七人離開世界,而在這些人當中,熱心參與教會事務的人佔了22.9\%,在比例上這些信徒的死亡率是較低的.

據說美國加州大學柏克萊分校也進行過相似的研究,這些研究人員用了二十八年的時間,訪問了五千名二十一至六十五歲以上的人.他們發現,在這些人當中,熱心參與教會事務的人的死亡率比非信徒低23\%.這些都印證了單張的說法,但是到底是不是真的呢?有教友問我:「牧師,多上教會是不是真的會更長壽?」我呆立一會,不懂回答.反覆思量下,我只能回答的是,在《聖經》中好像是沒有關於多上教會可令人長壽的教導.《聖經》內的確是有提及許多敬虔的,愛上帝的人,他們亦是相當長壽的,我所事奉的教會亦是一間有許多人瑞的教會,那麼可能說不定這是有道理的,因為有很多我們送別的長者,他們是以一個較高的年齡離開這個世界,我們就是有一個這樣的情況.從這個這樣的觀察,我想這也是有些道理的.

其實這單張所說的,也是有它的道理的.為甚麼呢?其實它所說的,恰巧印證了一位教授的說法:參與教會活動的人有長壽的可能,為甚麼呢?他說,信徒(參與教會活動的,大多是信徒)一般都有健康的生活,而且通過聚會,人際間的社會網絡關係亦會有比較好的發展,而這些社會網絡是會令人有更多支持的.再者,壓力在很多時候也是令人不能長壽的很重要的原因,因為人在社會上許多時候要面對佷多不同種類的壓力,不過回到教會其實是一個舒緩壓力的時候,因為(我們時常這麼教導信徒們)回到教會時將一切的思慮,憂愁,擔心全都放下,託付給那位為我們背負重擔的主,所以你可以安詳你的心靈來到神的跟前.再加上教會是個寧靜的地方,我們有很好的詩歌來帶領我們來唱,來聆聽,如此種種都可以幫助我們舒解壓力.所以,我相信如果你是患有高血壓的話,回到教會後說不定你的血壓會下降.當然,我們作為信徒,很多時候《聖經》教導我們注重身體,不沾煙酒.這樣對我們的身體是好的；有好的身體,比較長壽也是一個很合理的結果.所以我說,可能有一些間接的證據.不過,我這樣回覆那位詢問我的弟兄:「多上教會人長壽,我沒有辦法在聖經中找到確切的證據.不過,『多上教會人蒙福』,我覺得倒是真實的.」在《聖經》讓我們看到這是很真實的.我們知道,當我們回到教會(尤其是我們參加教會聚會)的時候,聚會其實是一個我們與上帝會面的時間.上帝應許我們,在這會面中,祂要給我們各樣的恩典及恩福.有些人認為,我已經信了耶穌,這就已經足夠了；我信了耶穌,我已經得救了,我信了耶穌後,我可以自己讀經,自己祈禱,這是基督教經常強調的；所以我並不是一定要上教會的,我不在教會,我一樣可以做一個基督徒.

真的,不上教會,其實也可以成為一個基督徒:基督徒是我們時刻的身份,並不只限於教會內.不過,我覺得這樣將基督徒日常生活中怎樣去蒙恩典,和我們在教會中蒙恩典混為一談,其實並不很正確.我們知道,我們今天在這地上生活,在這世上生活的時候,我們在享受上帝恩典的事情上,並不是每時,每處都均勻.我的意思是說,並不是處處一樣,時時一樣,因為在有些時候,有些地方,上帝會給我們特殊的恩典,恩惠,而在哪些時候,上帝給予我們的恩典,是每時每處都一樣呢?我相信是在將來,我們整個宇宙都歸於基督內,上帝充滿在我們的宇宙中的時候,我相信就會是這樣子.不過,這是在將來,並不是在今天.

今天,我們有時會受上帝的呼召,來到教會中參加聚會.在這聚會的地方,我們奉主的名敬拜,這地方彷如人間與天堂的接合點.在這地方,上帝有特別的恩典要賜給我們,尤其當我們在這裡宣講上帝的話,而我們又用我們的信心來接受的時候,這是特別真實.

所以,再次鼓勵大家不可停止聚會,當我們被上帝的靈召喚的時候,我們就來到教會中,打開我們的心扉來聆聽,我相信,上帝必然透過他的話,對我們有更多的賜恩.但願我們能像馬利亞一樣存著一顆受教的心,每當上帝將道放在我們面前的時候,我們都靜靜的聆聽,並從中得福.

我們一同禱告:我們的主耶穌,我們感激你,因為我們知道在生命裡,你有許多的恩典,恩福.你從來都是那位慷慨的上帝,你未嘗有一樣好處不分給那些行動正直的人.我們都是羨慕你的好處,渴望得著你的好處.我們知道,只要我們羨慕,只要我們渴望,只要我們來到你跟前,你重來也不拒絕,你從來都會將最好的賜給我們.我們要感激你,因為在每星期,當你準備豐盛的筵席,呼召我們到來的時候,我們都是飽飽的離開.我相信今天一樣,將來都是一樣.我們在此再次感恩.但願我們來到你面前,就被你的道喚醒,我們的生命得著你的提醒,激勵,甚至得著改變,以至我們能結果子,將來我們見面的時候,我們能有所呈獻.再次多謝你,禱告奉耶穌基督的名字,誠心所願.



\newpage

\section{隱藏有時顯明有時}
\label{sec:34}
\textbf{賀志勇博士 - 隱藏有時  顯明有時}
\newline
\newline
\hyperref[sec:33]{< < < PREV SERMON < < <}
~
\hyperref[sec:toc]{[返目錄]}
~
\hyperref[sec:35]{> > > NEXT SERMON > > >}


路加福音 8:40-56
\newline
\begin{longtable}{cl}
\hline
\hline
章節 & 經文 (和合本修訂版)\\
\hline
8:40 & \begin{tabularx}{0.7\textwidth}{X} 耶穌回來的時候，眾人迎接他，因為他們都等候著他。 \end{tabularx} \\ \\ \relax
8:41 & \begin{tabularx}{0.7\textwidth}{X} 有一個會堂主管，名叫葉魯，來俯伏在耶穌腳前，求耶穌到他家裡去， \end{tabularx} \\ \\ \relax
8:42 & \begin{tabularx}{0.7\textwidth}{X} 因為他有一個獨生女，約十二歲，快要死了。耶穌去的時候，眾人簇擁著他。 \end{tabularx} \\ \\ \relax
8:43 & \begin{tabularx}{0.7\textwidth}{X} 有一個女人，患了經血不止的病有十二年，在醫生手裡花盡了一生所有的，但沒有人能治好她。 \end{tabularx} \\ \\ \relax
8:44 & \begin{tabularx}{0.7\textwidth}{X} 她來到耶穌背後，摸他的衣裳繸子，經血立刻止住了。 \end{tabularx} \\ \\ \relax
8:45 & \begin{tabularx}{0.7\textwidth}{X} 耶穌說：「摸我的是誰？」眾人都不承認。彼得說：「老師，眾人擁擁擠擠緊靠著你。」 \end{tabularx} \\ \\ \relax
8:46 & \begin{tabularx}{0.7\textwidth}{X} 耶穌說：「有人摸了我，因為我覺得有能力從我身上出去。」 \end{tabularx} \\ \\ \relax
8:47 & \begin{tabularx}{0.7\textwidth}{X} 那女人知道瞞不住了，就戰戰兢兢地俯伏在耶穌跟前，把摸他的緣故和怎樣立刻痊癒的事，當著眾人都說出來。 \end{tabularx} \\ \\ \relax
8:48 & \begin{tabularx}{0.7\textwidth}{X} 耶穌對她說：「女兒，你的信救了你。平安地回去吧！」 \end{tabularx} \\ \\ \relax
8:49 & \begin{tabularx}{0.7\textwidth}{X} 耶穌還在說話的時候，有人從會堂主管的家裡來，說：「你的女兒死了，不要勞駕老師了。」 \end{tabularx} \\ \\ \relax
8:50 & \begin{tabularx}{0.7\textwidth}{X} 耶穌聽見就對他說：「不要怕，只要信！她必得痊癒。」 \end{tabularx} \\ \\ \relax
8:51 & \begin{tabularx}{0.7\textwidth}{X} 耶穌到了他的家，除了彼得、約翰、雅各，和女兒的父母，不許別人同他進去。 \end{tabularx} \\ \\ \relax
8:52 & \begin{tabularx}{0.7\textwidth}{X} 眾人都在為這女孩哀哭捶胸。耶穌說：「不要哭，她不是死了，是睡著了。」 \end{tabularx} \\ \\ \relax
8:53 & \begin{tabularx}{0.7\textwidth}{X} 他們知道她已經死了，就嘲笑耶穌。 \end{tabularx} \\ \\ \relax
8:54 & \begin{tabularx}{0.7\textwidth}{X} 耶穌拉著她的手，呼叫著：「孩子，起來吧！」 \end{tabularx} \\ \\ \relax
8:55 & \begin{tabularx}{0.7\textwidth}{X} 她的靈魂就回來了，她立刻起來。耶穌吩咐給她東西吃。 \end{tabularx} \\ \\ \relax
8:56 & \begin{tabularx}{0.7\textwidth}{X} 她的父母非常驚奇；耶穌吩咐他們不要把所發生的事告訴任何人。 \end{tabularx} \\ \\
[1ex]
\hline
\hline
\end{longtable}
日期:2010年2月14日

題目:隱藏有時  顯明有時

證道:賀志勇博士

經文:路加福音8:40-56

我們今天要講的是路加福音八章四十節到五十六節,我們一起來看一下這段

經文:

40 耶穌回來的時候,眾人迎接他,因為他們都等候他.

41 有一個管會堂的,名叫睚魯,來俯伏在耶穌腳前,求耶穌到他家裡去．

42 因他有一個獨生女兒,約有十二歲,快要死了.耶穌去的時候,眾人擁擠他. 43  有一個女人,患了十二年的血漏,在醫生手裡花盡了他一切養生的,並沒有一人能醫好他．

44 他來到耶穌背後,摸他的衣裳繸子,血漏立刻就止住了. 45  耶穌說,摸我的是誰．眾人都不承認,彼得和同行的人都說,夫子,眾人擁擁擠擠緊靠著你.〔有古卷在此有你還問摸我的是誰麼〕

46 耶穌說,總有人摸我．因我覺得有能力從我身上出去. 47  那女人知道不能隱藏,就戰戰兢兢的來俯伏在耶穌腳前,把摸他的緣故,和怎樣立刻得好了,當著眾人都說出來.

48 耶穌對他說,女兒,你的信救了你,平平安安的去罷. 49  還說話的時候,有人從管會堂的家裡來說,你的女兒死了,不要勞動夫子.

50 耶穌聽見就對他說,不要怕,只要信,你的女兒就必得救. 51  耶穌到了他的家,除了彼得,約翰,雅各,和女兒的父母,不許別人同他進去.

52 眾人都為這女兒哀哭捶胸.耶穌說,不要哭．他不是死了,是睡著了. 53  他們曉得女兒已經死了,就嗤笑耶穌.

54 耶穌拉著他的手,呼叫說,女兒,起來罷. 55  他的靈魂便回來,他就立刻起來了．耶穌吩咐給他東西喫.

56 他的父母驚奇得很．耶穌囑咐他們,不要把所作的事告訴人.

最近香港鬧的高鐵事件,讓香港的市民又開了一次眼界,之所以說讓大家開了一次眼界,是因為之前曾蔭權「慳電膽」事件,「慳電膽」在普通話裡頭叫做「節能燈」,以前曾特首喜歡打領結,我的朋友都喜歡叫他「煲呔曾」,現在都喜歡叫他「電膽曾」,其實呢,公眾並不是懷疑曾特首從中拿回扣,拿回佣,公眾只是希望,曾特首做事要公開,透明!

我想高鐵事件,也有這一方面的元素,為什麼會鬧成這個樣子呢,其實也是跟不公開,不透明有關,公平,公正,公開,這是市民的要求,可是呢,是不是什麼事情都要公開了,我想呢,要不竟然如此,我聽說,有一個教會有一個傳道人,一家三口本來很幸福的,不幸的,這個傳道人愛上了教會另外一個姐妹,發展了一段婚外的戀情!

後來傳道人悔改了,向牧師,執事會承認錯誤,但執事會認為,傳道人的做法不只褻瀆了神,也得罪了整個教會,應該向整個教會公開認錯,這個傳道人為了表明真心認錯,沒有辦法,只好答應了公開認錯,可是他沒有跟妻子,兒女商量好,最後呢,在某個主日崇拜之前,他突然出現在台上認錯,妻子,兒女在台下面,無地自容,然後不但家庭分裂,妻子女兒也從此不肯回教會,傷害太大了!

中國人講,天理不離人情,人情是什麼?人情就是人之常情,事事都要公開,事事都要講透,不但不可能,還可能弄得天下大亂,可是有些事情你不講透,不公開,又會讓人胡亂捉摸,眾多猜疑,也容易弄得天下大亂,所以做事情也好,講話也好,該顯明的顯明,該隱藏的隱藏,這是個大學問!

關於這方面的學問了,中國有一位大學問家,叫做「孔子」,外國也有一位大學問家,就是我們的「耶穌基督」,我們再看看剛才那段經文,基督的智慧,路加福音八章四十到五十六節,講了兩個神蹟,兩個小故事,一個是血漏十二年的婦女,花光了錢都沒辦法把病治療好,但是她摸了一下耶穌的衣裳繸子之後,病就好了!另一個呢,管會堂的睚魯,他的女兒已經去世了,可是耶穌一句話,就讓她活過來了!

這兩個小故事看起來,意思簡單明瞭,連在一起看,就會出現幾個奇怪的地方,第一,在一個人擠人的地方,很多人的地方,耶穌大叫一聲,誰摸了我的衣服,就好像早上八點我去上班,地鐵裡面人多得不得了,連腳都沒有地方放了,突然有人大叫一聲,誰摸我?你覺得有人會理他嗎?所以彼得說,夫子你有沒有弄錯呢?這麼多人擠來擠去,你還問是誰摸你,這還要問嗎?

第二,耶穌為什麼一定要迫這個女人出來,公開承認是她摸了耶穌的衣服,公開這個神蹟,為什麼呢?這好像並不像耶穌慈悲愛人的常理,我們知道血漏這種病,一聽到名字就很嚇人,當然這是一種什麼病,我也不清楚,可是在利未記第十五章,就說在當時的環境,血漏的女人會被認為不潔淨的,看到別人她就應該要躲起來,否則就會把她的不潔淨傳染給別人,律法要制裁她,這樣一個女人,她混在人群當中,已經是犯了律法,她還摸了一個男人,耶穌的衣服,那更加是犯了律法,也就是說,萬一她比人發現了,她會被羞辱,會被懲罰,耶穌肯定也知道這個女人的處境,但卻迫她出來,不但在人家面前說她得了這種病,還要她承認是她偷偷摸了耶穌的衣服,作為一個女人,是不是特別難看呢!耶穌為什麼一定要迫她這麼做?

第三,這個女人當時不是立刻痊癒呢嗎?耶穌基督為什麼要她公開承認之後,然後才說呢,女兒你的信救了你,平平安安的去吧,她不是病好了嗎?

第四,耶穌要這個血漏的女人公開她是如何得救的,但是當

祂

救這個小女孩的時候,就不讓別人進去,並且救活了這個小女孩之後呢,耶穌又阻止他們,不要把

祂

所做的事情告訴別人,在這裡,耶穌為什麼又不肯公開呢?還要他們保守祕密呢?聖經中很多經文起初看了都沒什麼問題,可是緻細一想,問題還真的不少,我想,天理也不離人情,許多問題的答案,其實都在事態人情當中,為什麼我這樣子說呢,大家可以想一下,這個女人雖然治好身體的病,但如果她就這樣偷偷的走了,她的心病治好了沒有?

人家說,偷來的鑼,敲不得,她雖然身體的病治好了,可是她不能跟別人說她治好,為什麼,第一,別人可能不相信她,既然不相信她,就不敢接近她,因為怕她沒痊癒好,怕被傳染了；第二,就算有人相信她,敢去接近她,也會問她是怎樣治好的,她可以怎麼說呢,是她偷耶穌的能力治好的嗎?如果她不想實話實說,她就撒謊,她已經犯了律法了,偷偷接近耶穌,難道她繼續犯律法去撒謊嗎?

我想她不會,因為她是一個誠實的女人,你看她,她再難看,最終還是走出來,跪下來,這就說明她是一位非常誠實的女人,所以大家想想看,如果她不能被別人知道,她的病已經好了,就算她身體的病已經好了,她還是沒有辦法被別人接受,她還是沒辦法融入社會,就算她身體的病好了,她還是要避開人群,孤孤單單沒人關心,沒人理她,這種孤獨的痛苦呢,會影響到她後半輩子,不緊如此,她也不知道這個病會不會復發,大家知道得了病的人最怕是什麼,就是好了又復發,比方說,癌症,心臟病,皮膚病,管它是治好了,還是會提心吊膽,它到低會不會復發呢?

我的手就有一點「主婦手」,很容易脫皮,開裂,好幾年了,好一點又差一點,痛苦得很,這個女人也是如此,她雖然被治好了,可是她經歷了十二年血漏的苦楚,等她回到家裡頭,你覺得她會不會提心吊膽這個會不會復發呢?提心吊膽的滋味好不好受呢?有時候提心吊膽的痛苦比身上的痛苦更加折磨人,耶穌就是知道這種種的人情世理,

祂

是真正的人情大師,要這個婦人站出來去承認的這個行動,不是為了羞辱她,而是要切切底底的治好她的心病,因為當這個女人當眾承認,是她摸了耶穌,得到耶穌能力的時候,也就是當眾宣佈,她的病已經完全好了!

而且耶穌也接受了她,告訴她,她的心已經得救了,於是當眾宣佈,這個能力不是這個婦人偷走的,而是耶穌自願賞給她的,並且這個婦人也不用擔心,因為救世主已經大聲宣佈,她已經得救了,她的信救了她,從此這個女人的內心,有了真正的平安,她就可以堂堂正正的站在人的面前,被大家接納,以後就跟別人一樣,生活在陽光底下,可以大大方方去接受愛,也可以大大方方去愛別人,不用擔心別人的眼光,她的生命是完全不一樣了,因為救世主已經宣佈醫治了她,給她建立了信心,你們覺得耶穌基督是不是一位治病的大師呢?

耶穌也治了睚魯的心病,其實耶穌要這個女人公開自己的病,不但是要堅定這個女人的信心,也是要堅定管會堂睚魯的信心,為什麼這樣說呢?因為當耶穌跟這個女人說話的時候,有人跟睚魯說,不要勞動夫子了,你的女兒已經死了,可是耶穌聽見了,就跟睚魯說,不要怕,只要信,你的女兒就必得救,你們看,我們我主耶穌是不是洞悉一切呢?

如果祂不讓這個女人出來做見證,睚魯就不會看到神的大能,如果睚魯沒有看到神的大能,他聽到女兒死了,他就只會在那邊傷心痛苦了,他不會有任何心情,也不會有任何信心再邀請主耶穌去治療他的女兒,所以你想想,神的作為是不是洞察先機呢?有的時候,當我們為我們的未來,買下所有的保險,為我們的未來絞盡腦汁,存心積累,我跟你講,不如把這份焦慮交給神吧!

祂

早為你洞察先機了,至於在治療小女孩的時候,耶穌不讓人公開,這其實是耶穌一貫的做法,

祂

救治一些瞎子,痳瘋病的人,都不讓他們出去說,之所以這樣做,是因為

祂

不想大家關注

祂

治病的能力,把

祂

當成一位有魔力的魔法師,

祂

希望人所關注的是福音,而不是治病或神蹟,耶穌基督知道,人往往關心身體的需要,就忘記了心靈的需要,往往關心當前的需要,卻忘記對生命最重要的需要,耶穌基督在這裡隱藏了真正體貼人的軟弱,

祂

知道人是軟弱的,所以

祂

的作法是希望人能夠注目於福音,你看看我們的主耶穌是不是多麼的細心,多麼的有人情味,

祂

不但治病,

祂

也治心,看醫生呢,就是要看這樣子的醫生.

我有一個朋友三十年前在香港大學讀書,他說那個時候,香港大學有一個醫務所,有一位很有趣的醫生,是一位洋人醫生,名字好像是Mr. Smith,Mr. Smith看病很少會幫人家看好的,有時候呢,不看他還好,看了反而更慘,但是學生都喜歡找他看病,為什麼呢?因為他態度最好,最有人情味,中國人很喜歡講人情味,就看冬天喝了一碗好湯,就像下雪天食熱騰騰的涮羊肉,會讓人心裡頭很溫暖的,主耶穌做事就是這樣子,他做事,隱藏有時

,顯明有時,不論是顯明或是隱藏,都包含著濃濃的人情味,真正的人情味呢,是一種體諒的愛,愛不是一件簡單的事情,有時候我們會很想愛,但往往我們不懂得怎麼去愛.

體諒的愛是什麼樣的愛?就是多從對方的角度去想的愛,這種愛不居高臨下,不氣勢凌人,不拖拉懶惰,這種愛是從你的眼睛而不是從自己的眼睛來看代事情,這種愛讓你樂於接受,我認識一位學生,他叔叔給他好多的幫助,他來香港讀書的錢,全都是他叔叔給的,可是回到家呢,沒有必要他絕對不會去叔叔家裡頭,為什麼?他說因為他看不慣,他叔叔把他當成下屬訓話般的口氣,人就是這樣子,人不單需要幫助,人更需要一份帶著體諒的幫助,人不僅需要愛,我們更喜歡愛是帶著體諒的.

我以前在北京讀書的時候,最怕長輩親戚請我食飯,因為我要陪笑臉,還要聽教訓,要神情專注聽長輩教訓,看到了好食的菜,不停的流口水,但是筷子都不敢動,所以後來長輩找我去食飯呢,我都推這推那,不想去.我想這種經歷不僅我有,蠻多人都有的,比方父母跟小孩子,小孩子考試未考好,最希望聽見呢是父母的鼓勵,但是聽到最多的是責罵,這種愛雖然深卻少了一種體諒,孩子很難接受的,體諒是什麼呢?在火車上呢,看到老人,孕婦,給她們讓座是一種體諒,不是因為你想要上TVB所以讓座,看到身邊的人不舒服,給一種問候也是一種體諒,對孩子多一點溝通少一點責備也是體諒,對妻子多一點讚美少一點挑剔也是一種體諒.

我們不是缺乏愛心,但往往像保羅所說,立志行善由得我,行不行得出來就由不得我,體諒是一種很細心的愛,當德蘭修女去幫助那些貧窮的人的時候,她的臉不是居高臨下的,是一種溫柔的體諒,為什麼?因為我們的神也是如此.主耶穌基督不僅僅體諒這個血漏的女人,不僅僅體諒失去女兒的一家,

祂

更體諒我們人的軟弱,

祂

為我們死在十字架上面的時候,

祂

還說,父阿,赦免他們的罪吧,因為他們所做的,他們不曉得.顯明有時,隱藏有時,可是不管是顯明也好,還是隱藏也好,但都是因為這份從神而來愛的體諒,一個體諒的人是一個充滿平安喜樂的人,是神所愛的,一個體諒的教會是一個充滿喜樂平安的教會,是神所愛的教會,弟兄姊妹,十二年長不長呢?對於血漏的女人,十二年的血漏,十二年的病,實在是太長了.那對於睚魯的女兒來看,十二年的青春也實在是太短了.但無論是長還是短,她們都因為得到了主耶穌基督的愛,而有了新的生命.雖然血漏婦人是膽顫心驚的信心,可是她還是得到了治療,雖然睚魯的信心是半信半疑的,可是他的女兒得到了復活.

有時候,我們的信心有可能是膽顫心驚的,也可能是半信半疑,但不要怕,只要信,因為基督體諒我們的軟弱,因為經文說,這福音本來就是神的大能,叫一切信

祂

的,不至滅亡,反得永生.我們一起禱告,親愛的天父,感謝你賜給我們教會,讓我們彼此相愛,我們還要愛教會以外的人,像是海地的災民,求你都給我們一個體諒的心,不是一個律法的心,看到你讓我們看到的世界,這一切都在你的手中,奉耶穌的名而求.



\newpage

\section{綑縛與釋放}
\label{sec:35}
\textbf{林豪恩先生 - 綑縛與釋放}
\newline
\newline
\hyperref[sec:34]{< < < PREV SERMON < < <}
~
\hyperref[sec:toc]{[返目錄]}
~
\hyperref[sec:36]{> > > NEXT SERMON > > >}


路加福音 4:14-21
\newline
\begin{longtable}{cl}
\hline
\hline
章節 & 經文 (和合本修訂版)\\
\hline
4:14 & \begin{tabularx}{0.7\textwidth}{X} 耶穌帶著聖靈的能力回到加利利，他的名聲傳遍了四方。 \end{tabularx} \\ \\ \relax
4:15 & \begin{tabularx}{0.7\textwidth}{X} 他在各會堂裡教導人，眾人都稱讚他。 \end{tabularx} \\ \\ \relax
4:16 & \begin{tabularx}{0.7\textwidth}{X} 耶穌來到拿撒勒，就是他長大的地方。在安息日，照他素常的規矩進了會堂，站起來要念聖經。 \end{tabularx} \\ \\ \relax
4:17 & \begin{tabularx}{0.7\textwidth}{X} 有人把以賽亞先知的書交給他，他就打開，找到一處寫著： \end{tabularx} \\ \\ \relax
4:18 & \begin{tabularx}{0.7\textwidth}{X} 「主的靈在我身上，因為他用膏膏我，叫我傳福音給貧窮的人；差遣我宣告：被擄的得釋放，失明的得看見，受壓迫的得自由， \end{tabularx} \\ \\ \relax
4:19 & \begin{tabularx}{0.7\textwidth}{X} 宣告神悅納人的禧年。」 \end{tabularx} \\ \\ \relax
4:20 & \begin{tabularx}{0.7\textwidth}{X} 於是他把書捲起來，交還給管理人，就坐下。會堂裡的人都定睛看他。 \end{tabularx} \\ \\ \relax
4:21 & \begin{tabularx}{0.7\textwidth}{X} 耶穌對他們說：「你們聽見的這段經文，今天已經應驗了。」 \end{tabularx} \\ \\
[1ex]
\hline
\hline
\end{longtable}
日期:2010年2月21日

題目:綑縛與釋放

經文:路加福音4:14-21

證道:林豪恩先生

弟兄姊妹平安,今天講道所講的經文是剛才所讀的福音書,路加福音第四章14-21節,請大家翻開這段經文,我會從三個段落來分享這段經文.

第一個段落,我給它一個名字,就叫「被綑綁的人歡迎釋放的信息」.經文內提及耶穌回到拿撒勒,就是他長大的地方,在安息日祂照常進入會堂內,祂站起來讀經.有人將以賽亞書交給祂,祂讀出了以賽亞書的其中一段經文.但是當祂讀完這段經文後,請大家看看四章二十二節:「眾人都稱讚他,並希奇他口中所出的恩言.」聽祂讀經的聽眾定睛看著祂,牢牢的望著祂,然後稱讚祂.究竟耶穌所讀的是甚麼經文,可以讓祂吸引到眾人的目光?究竟耶穌讀了甚麼經文,可以令他們這麼讚歎,這麼驚喜?耶穌讀的經文,主要出處是在以賽亞書六十一章一至二節.耶穌讀了這個內容.主的靈在我身上；因祂用膏膏我,叫我傳福音給貧 窮的人,差遣我報告,被擄的得釋放,瞎眼的得看見,叫那受壓制的得自由,報告上帝悅納人的禧年.耶穌的這一段宣告,宣告了一個新時代的來臨,一個由以賽亞先知所開始盼望的新時代的來臨.宣告了一個已經等待了數百年的新時代的來臨,在這個新時代裡,貧窮的人得福音,被擄的得釋放,在黑暗的會得到光明,受壓制的會得到自由.這些狀況為何會引起拿撒勒人這麼大的反應?為何會引起拿撒勒人這麼積極的迴響?這些狀況和拿撒勒人有何關係?拿撒勒人全都是貧窮的人嗎?拿撒勒人全都是被擄嗎?拿撒勒人全都是瞎眼的嗎?拿撒勒人全都是受欺壓的嗎?誰才是真正貧窮的人?為何拿撒勒人會對號入座?誰是貧窮人這個問題,就已經可以引起足夠的爭論了.是經濟上的貧窮,還是心靈上的貧窮?如果是經濟上的貧窮,那麼收入多少,擁有多少資產才叫貧窮?如果是心靈上的貧窮,是不是窮得只剩下錢也叫貧窮?貧窮已是一個很難定義的詞語.我們還要定義誰是被擄,誰是瞎眼,誰是受欺壓.不同的定義在影響我們對這經文的理解亦影響著我們對信仰的實踐.經文中所說的狀況,讓我們從一個群體的角度,而不是個人的角度；用一個生存的狀況的處境,而不是一些條件來衡量它,我們會發現另外一個圖畫.拿撒勒人是整體上,他們從心底裡認同他們就是經文所描繪的那夥人.他們是在經文中所描繪的,一群正等候盼望的人.他們不需要爭論,到底他們是心靈貧窮,還是經濟貧窮；他們亦不需爭論,他們到底是精神受壓,還是武力受壓.對於一個長期被外族統治的民族,一個生活在列強的鐵蹄之下的人民,一群要將他們的稅項上繳至外族的統治者的人民來說,哪些是貧窮人?哪些是被擄的?哪些是生活在黑暗內?哪些是被欺壓的?答案是不言而喻的.我們並不是說每一個拿撒勒人都貧窮,拿撒勒人中可能有人比他人富裕；我們也不是說每一個拿撒勒人都被人囚禁,不過,和羅馬所派遣的統治階層相比,哪一邊是較貧窮的?和羅馬所派遣的軍隊相比,哪邊是被囚禁的?我們不是說每一個拿撒勒人都被欺壓,他們當中亦可能有互相欺壓的情況.不過,和外族的統治者相比,誰是被欺壓的?相信生活在早期殖民地香港的人都會身同感受.有錢,有權力,有自由,有體面的,一定是洋人,洋行,洋服,和操洋語.沒錢,沒權力,沒自由,市井的,通常都是那些小漁港內的居民,和從北方湧進來的難民.這種劃分,基本上是一種共同的經驗,是一種共同的意識.對於這群在政治,經濟,宗教,文化,社會生活上各層面都感到自己被綑綁的人,他們聽到貧窮人得福音,被擄的得釋放,黑暗的得自由,受欺壓的得自由,這種訊息是何等的珍貴,何等的中聽,是何等的激勵,及何等的欣然.對於拿撒勒人,甚至是整個以色列人來說,釋放的宣告是非常適合他們的現實處境,亦是呼應著他們內心的盼望.

第二,第一我說被綑綁的人歡迎釋放的信息,第二是宣告釋放的人反而被綑綁.世事是往往出人意外的,亦是充滿著弔詭的.宣告好訊息的人,是不一定有好報的.耶穌在拿撒勒會堂宣告了一個釋放的訊息,幾分鐘後,這位宣告釋放訊息的人,就被綑綁了.我們一同看看四章二十八節.四章二十八節說:「會堂裡的人聽見這話,都怒氣滿胸,就起來攆他出城,他們的城造在山上；他們帶他到山崖,要把他推下去.他卻從他們中間直行,過去了.」是不是很特別?這位宣告釋放訊息的人,轉眼間他就被綑綁.憤怒的群眾圍著耶穌,甚至想將祂置諸死地.那一次,耶穌很奇妙的從群眾的熙攘中逃脫了.不過所謂能避一時,不能避一世,最終耶穌都不能逃離被綑綁至行刑台的遭遇.路加福音二十三章十三節這樣說:「彼拉多傳齊了祭司長和官府並百姓,就對他們說:你們解這人到我這裡,說他 是誘惑百姓的.看哪,我也曾將你們告他的事,在你們面前審問他,並沒有查出 他甚麼罪來；就是希律也是如此,所以把他送回來.可見他沒有做甚麼該死的事.故此,我要責打他,把他釋放了.眾人卻一齊喊著說:除掉這個人!釋放巴拉巴給我們!彼拉多願意釋放耶穌,就又勸解他們.無奈他們喊著說:釘他十字架!釘他十字架!彼拉多第三次對他們說:為甚麼呢?這人作了甚麼惡事呢?我並沒有查出他甚麼該死的罪來.所以,我要責打他,把他釋放了.他們大聲催逼彼拉多,求他把耶穌釘在十字架上.他們的聲音就得了勝.」就這樣,一位宣告自由的人就失去了他的自由了.根據美國一位著名的基督徒作家楊腓力的憶述,他說他兒童年代的美國是這樣的:「我們曾經生活在種族隔離政策年代.雖然亞特蘭大的黑人數目差不多跟白人一樣多,我們卻是在不同的餐廳進食,在不同的公園玩耍,上不同的學校和教堂.有時我會看見一些招牌寫著『黑人與狗不得進內』.法例規定黑人不得成為陪審員,不得送孩子到白人公立學校,不得用白人專用的廁所,不得在黑人旅館度宿,不得坐在電影院的大堂部份,不得在白人泳池游泳.因為阿拉巴馬州的度假地方並不招待黑人,馬丁•路德•金結婚當晚住在一個他可以使用而又最接近公眾旅館的地方,就是他家族朋友所擁有的一所殯儀館.」楊腓力這篇文章中提到的這一位在結婚當晚住在殯儀館的馬丁•路德•金,就是在一九六三年八月八日發表著名演說《我有一個夢》的馬丁•路德•金牧師.他在演說中這樣宣告,他說:「我夢想有一天這個國家會站立起來,真正實現其信條的真諦,信條就是我們認為這些真理是不言而喻的,人人生而平等.我夢想有一天,在喬治亞的紅山上,昔日奴隸的兒子將能夠和昔日奴隸主的兒子坐在一起,共聚兄弟情誼.我夢想有一天,甚至連密西西比州這個正義匿跡,壓迫成風,如同沙漠般的地方,也將變成自由和正義的綠洲.我夢想有一天,我的四個孩子將在一個不是以他們的膚色,而是以他們的品格優劣來評價他們的國度裡生活.我夢想有一天,阿拉巴馬州能夠有所轉變,那裡的黑人男孩和女孩將能與白人男孩和女孩情同骨肉,攜手並進.當自由之聲響起時,讓自由之聲從每一個大小村莊,每一個州和每一個城市響起來時,上帝的所有兒女,黑人和白人,猶太教徒和非猶太教徒,基督徒和天主教徒,都將手攜手,合唱一首古老的黑人靈歌:『終於自由啦!終於自由啦!感謝全能的上帝,我們終於自由啦!』」.這一位宣告自由,平等,友愛的的黑人牧師,曾經三次被捕,三次被行刺,要綑綁他的人,不單是那些和他膚色不同的白人,有些甚至是和他同樣膚色的黑人.有些黑人不能忍受馬丁•路德•金牧師所堅持的「打不還手」的非暴力原則,於是轉向武裝反抗,迫他下台.這位宣告自由,平等,友愛的牧師,在一九六八年四月四日,在田納西州參與領導罷工運動的時候,在旅館的陽台被刺客開槍擊中致死,年僅三十九歲.一個宣告平等,自由和友愛的牧師,他自己就失去了自由,失去了平等的對待.

我接下來要說的是劉曉波先生.在二零零九年十二月二十五日,在聖誕日,原北京師範大學老師劉曉波先生被判處煽動顛覆國家政權罪名成立,判監十一年.這是劉曉波先生第四次成為國家的囚犯.有一位民運人士寫了這樣的一首獄中詩,他說:「我怕孤獨,但連自己影子都難得一見.我怕黑暗,卻只能在鐵窗後面仰望藍天.我只靠夢生活,但夢中卻永遠只是飄著染血的辮子；我的罪名,卻只是對自由的渴望.」有評論者認為,劉曉波先生有份草擬的《零八憲章》所提的十九項主張,包括修改憲法,分權制衡,立法民主,司法獨立,公器公用,人權保障,公職選舉,城鄉平等,結社自由,集會自由,言論自由,宗教自由,公民教育等等,都是十分溫和,是公民最基本的,正當的權利,是在提倡用一個溫和改良的方式推動社會進步.這位評論者說看不到有任何顛覆國家的地方,亦看不到何罪之有.劉曉波先生所有的罪名,只不過是對一個民主自由的中國的渴望而已.評論者最後的一句「劉曉波先生所有的罪名,只不過是對一個民主自由的中國的渴望而已」說出了不少人的心聲.既弔詭又令人心痛的現實是,宣告釋放的人自己往往是被綑綁的.

第三,人可以被綑綁,但是真理是不能被綑綁的,真理是必然被釋放的.我們再回頭看看耶穌.第四章二十九節:「會堂裡的人聽見這話,都怒氣滿胸,就起來攆他出城,帶他到山崖,要把他推下去.」到底會堂裡的聽眾聽了些甚麼,以致他們要憤怒得綑綁這位宣告釋放訊息的耶穌呢?我們看看二十五節,看看耶穌究竟對他們說了些甚麼:「我對你們說實話,當以利亞的時候,天閉塞了三年零六個月,遍地有大饑荒,那時,以色列中有許多寡婦.以利亞並沒有奉差往他們一個人那裡去,只奉差往西頓的撒勒法一個寡婦那裡去.先知以利沙的時候,以色列中有許多長大痲瘋的,但內中除了敘利亞國的乃縵,沒有一個得潔淨的.」到底是甚麼意思呢?拿撒勒人聽到了甚麼呢?耶穌宣告了釋放的訊息後,說今天這已經應驗在你們耳中了.

沒錯,這個宣告的訊息,是說給拿撒勒人聽的.但是,祂並不是只說給拿撒勒人聽的.耶穌說先知以利亞的事蹟,說先知以利沙的事蹟,正在顯示著上帝的眼目並不是只在看以色列人.最難接受的是,上帝的眼目竟然也都看以色列人的敵人.上帝的恩典並不單單惠及以色列人,最難受的是,上帝的恩典竟然都惠及以色列的敵人.拿撒勒人要綑綁耶穌,因為耶穌是他們的「猶奸」,就正如我們中國人的漢奸一樣.拿撒勒人憤怒,因為他們接受不了上帝不單是他們的上帝,同時也是他們敵人的上帝；他們接受不了先知不單是以色列的先知,他們其實是上帝的先知；他們接受不了彌賽亞不單是猶太人的彌賽亞,竟然是全人類的彌賽亞.沒有人不喜歡釋放的訊息,沒有人不喜歡被釋放；不過,他們未必喜歡其他人像自己一樣得到釋放.沒有人不喜歡平等,不過未必喜歡其他人得到和自己一樣的平等.沒有人不喜歡自由,不過未必喜歡其他人享受和自己一樣的自由.這就是我們.這就是我們的遊戲規則.所以,粵語長片中的媳婦,永遠都等待著自己變成奶奶的那天；唱著《半斤八兩》的工作階層,永遠都等待著自己變成老闆的那天；被統治的人,永遠都等待著自己變成統治者的那天.到那天,大家都翻身了；到那天,位置都逆轉了；到了那天.......

彼拉多曾經私下審問耶穌,你是猶太人的王嗎?你是不是猶太人的王?耶穌回答 說:「這話是你自己說的,還是別人論我對你說的呢?」彼拉多說:「我豈是猶太人呢?你本國的人和祭司長把你交給我.你做了甚麼事呢?」耶穌回答說:「我的國不屬這世界；我的國若屬這世界,我的臣僕必要爭戰,使我不至於交給猶太人.只是我的國不屬這世界」.彼拉多就對他說:「這樣,你是王麼?」耶穌回答說:「你說我是王.我為此而生,也為此來到世間,特為給真理作見證.凡屬真理的人就聽我的話.」耶穌是王,不過祂的國不屬於這世界；所以,祂不跟隨這世界的遊戲規則.耶穌最終因此而被綑綁,但是祂所啟示的真理卻是不能夠被綑綁的.耶穌最終被釘十字架而死,但是祂所帶來的真理卻是綿延萬世,惠及萬民；跟隨真理的人繼續發出掙脫綑綁,繼續發出掙脫欺壓,繼續發出掙脫纏擾的這種呼聲,繼續發出平等,自由,友愛的呼聲.這種呼聲在人類歷史上從未消失過,呼喚著每一個人的良知的回應.

馬丁•路德•金牧師從黑人民權運動那裡被推下台之後,他眼見暴亂不停在洛杉磯,芝加哥和克林這些地區爆發；他穿洲過省提醒民眾,道德的改變不能透過不道德的手段達致.他深信只有一個基於愛的運動,才可以使壓迫者不至於成為壓迫他的人的翻版.他想改變白人的心,但是他不願意在這過程中讓他所帶領的黑人變了心硬.他相信非暴力是會讓黑人不至於以一個暴君來取代另一個暴君.一九六四年,馬丁•路德•金牧師獲得諾貝爾和平獎.一九八六年,美國規定每年一月的第三個星期一為馬丁•路德•金紀念日,定為法定假日.評論劉曉波事件的評論者,他這樣說,他說:「獨裁者將會慢慢發現他們無法杜絕任何批評,因為即使你審查,甚至禁制了書刊,一切有正義感的作家將有辦法找到自己的表達方式.劉曉波如今是被判入獄了,但我堅信歷史將宣判他無罪.」

平等,自由,友愛,是每一個人都歡迎的訊息,是每一個人都追求的狀況.不過,當有一些人拒絕讓另一些人得到平等,自由,友愛的時候,又或者是一些人堅持比另一些人享受更多的平等,自由和友愛的時候,綑綁就出現了.不單是綑綁了那些缺乏平等,自由,友愛的人,其實也是綑綁了自己.不讓自己經歷平等所帶來的平安和安全,亦妨礙了自己享受自由的那種舒暢,和令自己缺乏了和人與人之間的那種友愛,那種手足之情.被綑綁的要釋放,綑綁人的其實也要釋放.耶穌為此而來,亦為此而死.但是耶穌死而復活,鼓舞了那些跟隨著祂腳步的人,努力地邁向平等,自由,友愛的新世界.



\newpage

\section{獨食難肥}
\label{sec:36}
\textbf{鄧瑞強博士 - 「獨食難肥」}
\newline
\newline
\hyperref[sec:35]{< < < PREV SERMON < < <}
~
\hyperref[sec:toc]{[返目錄]}
~
\hyperref[sec:37]{> > > NEXT SERMON > > >}


路加福音 4:16-30
\newline
\begin{longtable}{cl}
\hline
\hline
章節 & 經文 (和合本修訂版)\\
\hline
4:16 & \begin{tabularx}{0.7\textwidth}{X} 耶穌來到拿撒勒，就是他長大的地方。在安息日，照他素常的規矩進了會堂，站起來要念聖經。 \end{tabularx} \\ \\ \relax
4:17 & \begin{tabularx}{0.7\textwidth}{X} 有人把以賽亞先知的書交給他，他就打開，找到一處寫著： \end{tabularx} \\ \\ \relax
4:18 & \begin{tabularx}{0.7\textwidth}{X} 「主的靈在我身上，因為他用膏膏我，叫我傳福音給貧窮的人；差遣我宣告：被擄的得釋放，失明的得看見，受壓迫的得自由， \end{tabularx} \\ \\ \relax
4:19 & \begin{tabularx}{0.7\textwidth}{X} 宣告神悅納人的禧年。」 \end{tabularx} \\ \\ \relax
4:20 & \begin{tabularx}{0.7\textwidth}{X} 於是他把書捲起來，交還給管理人，就坐下。會堂裡的人都定睛看他。 \end{tabularx} \\ \\ \relax
4:21 & \begin{tabularx}{0.7\textwidth}{X} 耶穌對他們說：「你們聽見的這段經文，今天已經應驗了。」 \end{tabularx} \\ \\ \relax
4:22 & \begin{tabularx}{0.7\textwidth}{X} 眾人都稱讚他，並對他口中所出的恩言感到驚訝；他們說：「這不是約瑟的兒子嗎？」 \end{tabularx} \\ \\ \relax
4:23 & \begin{tabularx}{0.7\textwidth}{X} 耶穌對他們說：「你們一定會用這俗語向我說：『醫生，你醫治自己吧！我們聽見你在迦百農所做的事，也該在你自己的家鄉做吧。』」 \end{tabularx} \\ \\ \relax
4:24 & \begin{tabularx}{0.7\textwidth}{X} 他又說：「我實在告訴你們，沒有先知在自己家鄉被人接納的。 \end{tabularx} \\ \\ \relax
4:25 & \begin{tabularx}{0.7\textwidth}{X} 我對你們說實話，在以利亞的時候，天閉塞了三年六個月，遍地有大饑荒，那時，以色列中有許多寡婦， \end{tabularx} \\ \\ \relax
4:26 & \begin{tabularx}{0.7\textwidth}{X} 以利亞並沒有奉差往她們中任何一個人那裡去，只奉差往西頓的撒勒法一個寡婦那裡去。 \end{tabularx} \\ \\ \relax
4:27 & \begin{tabularx}{0.7\textwidth}{X} 在以利沙先知的時候，以色列中有許多痲瘋病人，但除了敘利亞的乃縵，沒有一個得潔淨的。」 \end{tabularx} \\ \\ \relax
4:28 & \begin{tabularx}{0.7\textwidth}{X} 會堂裡的人聽見這些話，都怒氣填胸， \end{tabularx} \\ \\ \relax
4:29 & \begin{tabularx}{0.7\textwidth}{X} 就起來趕他出城。他們的城造在山上；他們帶他到山崖，要把他推下去。 \end{tabularx} \\ \\ \relax
4:30 & \begin{tabularx}{0.7\textwidth}{X} 他卻從他們中間穿過去，走了。 \end{tabularx} \\ \\
[1ex]
\hline
\hline
\end{longtable}
日期:2010年2月28日

主題:「獨食難肥」

經文:路加福音4:16-30

證道:鄧瑞強博士

各位弟兄姊妹,主內平安.

今日講道的經文,是今日的福音書經課:路加福音4: 21-30.

不知大家有否留意,今日的福音書經課,和上個禮拜的福音書經課,其實是同一個故事的上下段.這個故事的上段和下段連得很緊,故講故事的上段時,不期然會講到下段.講這故事的下段時,又不期然地講到這故事的上段.故此,上個禮拜林豪恩先生講道時,已講完我要講的經文.並且他講得很好,他講的內容也正是我要講的內容.我已無話可說.我想過,最好的解決方案是:我病了,然後失蹤.然後,找人再拿林先生的講章,再唸一遍.我想,這和我講一篇道的效果也差不多.但既然病不成,故我只能上來講道.我的主旨和林先生的,沒有大分別.若我講得好,除了歸榮耀給神之外,也應將 credit 歸給林先生.幸好,按我過去多年聽道的經驗,聽過的道,通常在下午已忘掉一半,到晚上則全忘了.故此,我再講林先生講過的道,但用另一番話來講,等你以為是新的.若你覺得是新的,則我的奸計便得逞.話說回來,同一番道理,聽完再聽,有一個好處,就是讓各人到了今晚,仍不會完全忘記,到明天還記得一點.

讓我再讀一遍

經文:路加福音4:16－30

4:16　耶穌來到拿撒勒,就是他長大的地方.在安息日,照他平常的規矩進了會堂,站起來要念聖經.

4:17　有人把先知以賽亞的書交給他,他就打開,找到一處寫著說:

4:18　主的靈在我身上,因為他用膏膏我,叫我傳福音給貧窮的人；差遣我報告:被擄的得釋放,瞎眼的得看見,叫那受壓制的得自由,

4:19　報告神悅納人的禧年.

4:20　於是把書捲起來,交還執事,就坐下.會堂裡的人都定睛看他.

4:21　耶穌對他們說:「今天這經應驗在你們耳中了.」

4:22　眾人都稱讚他,並希奇他口中所出的恩言；又說:「這不是約瑟的兒子嗎?」

4:23　耶穌對他們說:「你們必引這俗語向我說:『醫生,你醫治自己吧!我們聽見你在迦百農所行的事,也當行在你自己家鄉裡』」；

4:24　又說:「我實在告訴你們,沒有先知在自己家鄉被人悅納的.

4:25　我對你們說實話,當以利亞的時候,天閉塞了三年零六個月,遍地有大饑荒,那時,以色列中有許多寡婦,

4:26　以利亞並沒有奉差往她們一個人那裡去,只奉差往西頓的撒勒法一個寡婦那裡去.

4:27　先知以利沙的時候,以色列中有許多長大痲瘋的,但內中除了敘利亞國的乃縵,沒有一個得潔淨的.」

4:28　會堂裡的人聽見這話,都怒氣滿胸,

4:29　就起來攆他出城(他們的城造在山上)；他們帶他到山崖,要把他推下去.

4:30　他卻從他們中間直行,過去了.

這段經文給我的第一個教訓是:做傳道人很危險.除非你深信有神「照著」,否則不要隨便做.傳道人講一番話,可能即刻引來眾人衷心的讚賞.但當你再多講一句話,可能又即刻引來眾人「衷心」的殺機.經文記載,那些人聽了耶穌多講一點話之後,即刻想推耶穌落山崖.故此,我很怕走到山崖.就算在地鐵站月台,作為傳道人,我也怕遭人暗算,推我落月台.

故此,傳道的人,務要小心,這是一件危險的工作.因為,這個世界,很渴望聽道,也很渴望反抗道.

這段經文給我的第二個教訓是:人在種種困局中,多麼渴望解救,而神的說話能帶來解救.

主耶穌打開以賽亞書,讀出經文.路加福音4:18 那裡說:「

主的靈在我身上,因為他用膏膏我,叫我傳福音給貧窮的人」.「福音」就是好消息.窮人能得到什麼好消息呢?在古代社會,在現代社會,窮人得到的消息通常都是壞消息.你要種地,無地你種.你要找工作,沒有工作給你做.你沒有人工加,但通脹又來了.巴士加價,對富人沒有影響,但對窮人影響甚大.在香港,窮人能指望政府什麼呢?最多只是財政司派幾粒糖.若果有一個世界,貧窮人也能有好消息,這真是天國.

經文接著說:「

差遣我報告:被擄的得釋放」.古代的戰俘,是典型被擄的人.戰俘能生存,但沒有講話的自由,沒有行動的自由,沒有尊嚴.現代有沒有戰俘呢?當一個世界,容許你生存,或者說,你的生存權受到重視.但不容許你自由地講話和行動,不容許你追求人性的尊嚴時,這也就是一個人性被擄的世界.我時時搭乘公車.公車上有電視,播放著廣告.若那電視聲浪太大的話,我便變成廣告的俘擄.我坐在那裡,無處可逃,任廣告魚肉.

經文接著說:「

瞎眼的得看見,叫那受壓制的得自由」.上兩個禮拜,我的神學院開學前去大嶼山一個營舍舉作退修會.舉目向天,天是無限的,可以向四面開展.晚上,星會自然走出來照綴黑夜.偶然,會有流星劃過長空,提示你,你是可以有願望的.對比之下,忽然覺得,平時看的天是有限的.若你四周都是高樓大廈,你會發現,你像在井中,看的天是有限的.天不能向四面開展.這時,你會發現自己多麼受壓制,多麼渴望自由.這時,你會發現,我們受城巿的空間設計壓制著,我們受城巿空間設計背後的經濟規律壓制著,我們受城巿空間設計背後的經濟規律背後的生存模式壓制著.這種生存模式,像一條高速鐵路,穿過並摧毀你安靜地耕種的大地.在這種生存模式裡,人是瞎眼的,很難看到星星,很難看到流星,很難再有願望.

若果有一個世界,「

被擄的得釋放」,再沒有人被擄,再沒有人因為自由地說話而被禁制,則這真是天國.若果有一個世界,「

瞎眼的得看見,受壓制的得自由」,人在其中,能看到生存的意義,看到星,能再有願望,則這真是天國.

主耶穌在會堂中,宣告說:「

今天這經應驗在你們耳中了.」

主耶穌乃是來到人間的神.他的到來,正是要將每個人所需要的恩典帶來.因著耶穌,生命會不再一樣.

因著耶穌這番好聽的話,眾人都稱讚他.

這段經文給我的第三個教訓是:要帶來人性的解放,耶穌要做的,是越過某種界線.這種界線是一種很複雜的界線.人用這界線去界定自己特殊的身分,去保障自己的特權,去認識自己的命運,去令自己有一種存在的穩妥.但耶穌卻是動搖這種界線.不動搖這種界線,人間是不能有真正的解救的.

耶穌要做什麼?

在路加福音4: 25-27,耶穌說:「

我對你們說實話,當以利亞的時候,天閉塞了三年零六個月,遍地有大饑荒,那時,以色列中有許多寡婦,以利亞並沒有奉差往她們一個人那裡去,只奉差往西頓的撒勒法一個寡婦那裡去.先知以利沙的時候,以色列中有許多長大痲瘋的,但內中除了敘利亞國的乃縵,沒有一個得潔淨的.」

話說得明白,耶穌要打破「神的選民」和「非神的選民」兩者之間的界線,要打破「我是神聖的」和「你是不潔的」兩者之間的界線,要打破「我是有權利獲得上天的眷顧的」和「你是無權獲得上天的眷顧的」兩者之間的界線.有人說,這是一條神聖的界線,不容打破；耶穌說,根本無這界線.

有一個故事是這樣的:

在一個考試裡,哲學教授將一張椅子放在台上,說:「請各同學以各種學過的哲學理論,說明這張椅子是不存在的.」各同學於是拼命的寫.有一個同學,只寫了幾秒,簽上大名,交了卷,便施施然走出課室.過了一段日子,教授公佈考試結果.出乎意外,那答題不超過十秒的同學拿了最高分.他的答卷只有兩個字,寫著:「What chair?」

人間很多界線,但耶穌的回答是:「什麼界線?」單是這一表態,已為耶穌招來殺身之禍.但是,不移走這些界線,人間又怎能獲得解救?

大作家茨威格(Stefan Zweig)寫了一本書,叫《昨日的世界》(The World of Yesterday),他寫到以前的歐洲,人可以隨意在國與國之間漫遊.沒有證件,不需護照.就算去到美國,也可隨意入境,隨意找份工作,隨意留下來.但是,經過第一次世界大戰,經過第二次世界大戰,在國與國之間入境和出境,就需要花大量時間搞不同的證件.界線豎立起來了.

越是不信任,越是透過戰爭解決問題,越是「自己人」和「非自己人」清楚分開時,越多的界線便設立起來了.

在《阿凡達》(Avatar)這齣電影裡,外星人的世界和地球人的世界的其中一個分別,是外星人的世界是一個處處可感通的世界.外星人可以和野獸感通,可以和樹感通,可以和大地感通.外星人有條辮,只要這條辮和其他事物連結,便能和它們感通.這條辮裡面看起來有點似光纖,這應該多謝中大的前校長高錕發明了光纖,外星人也拿來用.在這個感通的世界裡,是無需以征服來解決問題的.

在《切古華拉少年日記》(The Motorcycle Diaries)這齣電影裡,講到南美的革命人物切古華拉(Che Guevara)的少年事蹟.有一次,他到一個痲瘋病院裡去做一點義工.痲瘋病院規定工作人員要戴手套工作,但在第一日,切古華拉已放棄手套,親手和病人握手.病人的第一個反應是:這真是一個人.那個痲瘋病院分佈在河的兩岸.工作人員在一方,病人在河的另一方.要工作,就先坐船到對岸.有一個晚上,是切古華拉的生日,員工們為他舉行了一個生日會.但他覺得,這生日會不能沒有對面病人的參與.晚了,沒有船,他就游水到對岸,這是一個異常危險的舉動,因他有哮喘,晚間的河水也潛伏著各種危險,但他不顧一切地游過去了.這舉動,打破了以河為界的界限.他走到了被認為不宜參與他的生日會的人當中,與他們分享自己的生日.其實,這是生命的真正分享.不打破這界線,生命是不能分享的.

另一齣電影,《甘地傳》(Gandhi).甘地,印度的一個解放英雄.要解放印度,就要打破種種人為的界線.其中一幕,他的太太和他爭論.他太太不甘願去洗廁所,因為在「種姓」制度中,這是下等賤民(The Untouchable)才作的事,而這位太太是屬於高級「種姓」的.甘地對太太說,你不作這些事,便離開這裡.單單這件事,甘地便越過了印度教裡一條不可越過的界線.

當甘地由印度教的世界要踏進回教的世界時,一個槍手開槍奪了他的命.當耶穌越過義人和罪人的界線時,自以為義的人也奪了他的命.

作家Jack London寫了一本小說,叫《白牙》(White Fang).說到一隻狼「白牙」的故事.有一個人馴養了牠,真心關心牠.牠也愛這個主人.「白牙」最喜歡吃雞,而偏偏這個主人就是開養雞場的.見到雞,「白牙」的本能反應便是撲向雞.這時,牠的主人開口制止,牠便放下雞.如是者,經過一段日子.主人的爸爸說:「白牙是無法改變吃雞的天性的.」但「白牙」的主人卻對「白牙」有信心.基於「白牙」對主人的愛和順從,這隻狼終於克服了與生俱來的天性和慾望.狼與雞,本來有一條不可克服的界線,但基於主人對狼的愛,也基於狼對主人的愛,這界線被克服了.

人間仇恨的界線,幾時才能因愛而取消呢?

但願榮耀歸於聖父,聖子,聖靈.阿們.



\newpage

\section{我可以差遣誰呢}
\label{sec:37}
\textbf{伍渭文牧師 - 我可以差遣誰呢?Whom Should I Send?}
\newline
\newline
\hyperref[sec:36]{< < < PREV SERMON < < <}
~
\hyperref[sec:toc]{[返目錄]}
~
\hyperref[sec:38]{> > > NEXT SERMON > > >}


以賽亞書 6:1-8
\newline
\begin{longtable}{cl}
\hline
\hline
章節 & 經文 (和合本修訂版)\\
\hline
6:1 & \begin{tabularx}{0.7\textwidth}{X} 當烏西雅王崩的那年，我看見主坐在高高的寶座上。他的衣裳下襬遮滿聖殿。 \end{tabularx} \\ \\ \relax
6:2 & \begin{tabularx}{0.7\textwidth}{X} 上有撒拉弗侍立，各有六個翅膀：兩個翅膀遮臉，兩個翅膀遮腳，兩個翅膀飛翔， \end{tabularx} \\ \\ \relax
6:3 & \begin{tabularx}{0.7\textwidth}{X} 彼此呼喊說：「聖哉！聖哉！聖哉！萬軍之耶和華；他的榮光遍滿全地！」 \end{tabularx} \\ \\ \relax
6:4 & \begin{tabularx}{0.7\textwidth}{X} 因呼喊者的聲音，門檻的根基震動，殿裡充滿了煙雲。 \end{tabularx} \\ \\ \relax
6:5 & \begin{tabularx}{0.7\textwidth}{X} 那時我說：「禍哉！我滅亡了！因為我是嘴唇不潔的人，住在嘴唇不潔的民中，又因我親眼看見大君王－萬軍之耶和華。」 \end{tabularx} \\ \\ \relax
6:6 & \begin{tabularx}{0.7\textwidth}{X} 有一撒拉弗向我飛來，手裡拿著燒紅的炭，是用火鉗從壇上取下來的， \end{tabularx} \\ \\ \relax
6:7 & \begin{tabularx}{0.7\textwidth}{X} 用炭沾我的口，說：「看哪，這炭沾了你的嘴唇，你的罪孽便除掉，你的罪惡就赦免了。」 \end{tabularx} \\ \\ \relax
6:8 & \begin{tabularx}{0.7\textwidth}{X} 我聽見主的聲音說：「我可以差遣誰呢？誰肯為我們去呢？」我說：「我在這裡，請差遣我！」 \end{tabularx} \\ \\
[1ex]
\hline
\hline
\end{longtable}
日期:2010年3月7日

主題:我可以差遣誰呢?Whom Should I Send?

經文:以賽亞書6:1-8

證道: 伍渭文牧師

賽六章1-8節這段經文是舊約聖經最為人知曉及引用的經文.很多宣教士和作獻身作傳道的,都被這段經文所感動:我可以差遣誰呢?誰肯為我們而去呢?自動對號入座:我在這裡,請差遣我.這段經文對了解基督教的崇拜神學的核心經文,因為這段經文的場景就是聖殿的敬拜.不少崇拜的禮文和詩歌都從這段靈感,就像我們今天聖餐時唱的凡有血氣歌,是正教會Orthodox Church聖餐時所唱的禮儀詩歌:凡有血氣,皆當靜默,敬畏謙恭同肅立,莫想絲毫世俗事物,因主基督臨世界,親手分賜天福神恩,接受我眾恭敬拜,詩歌左上角寫上Liturgy of St. James, 聖雅各禮儀,第四世紀.採用這禮儀的正教會從第四世紀到現在,每主日聖餐都唱這首詩歌.詩歌第三節引伸經文說:萬千天軍,排列成行,前隊先鋒一路長,周圍環繞光中真光,永晝之中天堂降；地獄威權終止消除,驅走黑暗滅死亡,第四節更是節重複

經文:伏地腳跟,恆常警醒,撒拉弗與基路伯-兩翅遮面,高聲讚美,頌歌騰響門庭震；阿利路亞,,阿利路亞,阿利路亞主至聖.

1. 崇拜是與上主對話

A 崇拜的抒發作用

.我們今早也從崇拜的角度來思考這段經文.首先我們看到的是一個對話的體裁,崇拜是與上主對話.英文崇拜是worship, 從worthship一字而來,崇拜重點在崇敬上主,把他的榮美,威嚴,超越,權能,良善,恩典藉詩歌,祈禱,奉獻顯揚,是上主本身堪當接受我們給他的敬拜,是從人給上主的服事.一詞似更能表達賽六章的情景.有些教會甚至把崇拜開始時那十五分鐘的唱詩時間稱為敬拜,,我收到很多在崇基教堂結婚的程序表都是把唱詩那段時間為敬拜,當然我會要求改為唱詩或頌詩,頌讚,因為整個婚禮都是一個敬拜(wedding service) ,整個崇拜都是敬拜,像我現在的講道,以前的讀經,跟著的奉獻,甚至我們彼此握手表現和好問安都是敬拜的行動.英文敬拜worship著重崇拜中的抒發作用expression,我們藉詩歌,祈禱,奉獻抒發對上主的情感.但崇拜不單有抒發作用,更有引發作用和誘發作用.詩歌,祈禱,奉獻不單可以抒發對上主的情感,也可以誘發對上主應有的情感.

B 崇拜的誘發作用

.多年前有一次我在一間教會專題講崇拜學,到答問環節時有一位青年人問我:每逢唱到一些跟隨主時,我心裡有難處,我覺得自己沒有資格唱,所以我不唱,我看到他時心裡感恩,因為對自己有要求,對自己的靈性景況有自覺.我讚賞他後對他說,祇要我們有願意作的心,上主就會悅納.誰人可以在誦唸主禱文的時候說自己真的做到願上主的國降臨呢?我們每天都記起上主的國度,並願上主的旨意行在地上嗎?恐怕我們平常就忙這忙那,不多時間想起上主的國度.我對他說,詩歌不單用來抒發我們對上主的情感,也是上主用來誘發我們對他應有的情感,你應唱,藉唱詩歌堅定自己的信心.

基督徒常常生活在應然和實然的張力中,我們知道應該如何活出應有的生活,但此刻我們做不到,認罪文就提醒我們,我們沒有做應該做的,也做了不應該做的(sins of omission and sins of commission) .上主就藉崇拜提醒我們,激勵我們:從實然走向應然,所以我喜歡德文崇拜一字,Gottesdienst (service of God) 上主的服事有兩層意思,可繙作人給上主的服事,像英文worship,但亦可作從上主而來給人的服事,這字 更能 比較表達以賽亞六章1-8的意思.我們看看第一節:當烏西雅王崩的那一年.為何提及這一年呢?這一年與先知看到主高高坐大寶座上有何關系呢?烏西雅王崩那一年是令先知非常沮喪,舉國上下感到悲傷.上主知道先知的需要,藉這異像服事他,叫他明白上主仍然坐著為王,他不需擔心.歷代志下二十六章16-19節對烏西雅王的突然長大麻瘋有很戲劇性的描述.烏西雅王本是好王,原是英明的君主,帶領猶大中興,可惜後來心高氣傲,擅自進入聖殿燒香,大祭司亞撒利雅帶同八十個勇敢的祭司試徒擋住他對他說:烏西雅阿,給耶和華燒香不是你的事,乃是亞倫子孫接聖職的事.你出聖殿吧!.....烏西雅就發怒,手拿香爐要燒香.他向祭司發怒的時候,額上忽然發出大麻瘋.

C 每項禮序都重要

.亦應是上主透過每一項禮序來服事我們.我記起出來在教會作傳道不久,崇拜後有一位姊妹對我說:伍先生,多謝你今天的祈禱,幫助了我,但對我的講道集字不提.所以我們來到教堂崇拜,可以祈禱求上主按你真正的需要幫助你.崇拜是給上主的服事,也是上主給我們的服事,崇拜有抒發作用,也有誘發作用.其實崇拜每一項都有獨特的意義,有多層次的宗教情感,整全的屬靈操練.唱第一首詩是湧溢讚美adoration,認罪是反思生命,學習謙卑求諸己；代求擴張屬靈的視野,操練對世界的悲憫；奉獻誘發感恩的心；握手問安表達和好的祝願；聽道操練對上主話語的愛慕；聽道後的回應詩一般是立志向高處行；領受祝福叫我們學習信賴.

2. 崇拜是信仰的表白 以賽亞六章1-8節的對話體裁,其中有很長的篇幅說到先知認罪,被潔淨.他在聖潔的上主面前感到不配,他是嘴唇不潔的人,天使用炭火潔淨他的跆嘴唇,所以認罪在崇拜的禮儀非常重要.我們的程序表也有認罪這一項:召請,認罪,宣赦,回應詩歌.因為崇拜是信仰的表白,保羅在林前十一章26節論聖餐時說:你們每逢吃這餅,喝這杯,是表明主的死,直等到他來.

認罪,被潔淨,得救贖就是福音的信息

.墮落後,人被逐出樂園,上帝在伊旬園東面安設基路伯(天使)用轉動發火焰的劍,把守生命樹的道路.火代表上主的聖潔,人不能接近聖潔的上主,人不能靠自己重返樂園,人不能靠自己得到生命樹.但在這裡,天使撒拉弗(seraphine火的意思)用火炭潔淨人的罪.上帝主動親近人,走向人,潔淨人的罪.

很多學者提到究竟以賽亞在聖殿見到的是誰呢?是上帝以他的本體對先知顯現嗎?但根本人沒有可能面對面見到上主而能存活.約翰在福音書說:

從來沒有人見過神,祇有父懷裡的獨生子將他表明出來.為此很多解經家認為在舊約中以天使形象對人顯現的神,就如在博雅渡口雅各與天使摔交,這天使都是成為肉身前的聖子基督.凡有血氣皆靜默那首東正教聖餐時的詩歌就是這樣說的:萬千天軍,排列成行,前隊先鋒一路長,周圍環繞光中真光,永晝之中天堂降；地獄威權終止消除,驅走黑暗滅死亡.詩歌說到基督從高天走向人間,他就是先知所見到在高高坐在寶座的那一位,被天使天軍簇擁的那一位就是基督,他降落凡塵,消滅地獄.

我們看跟著的經文1-13節,就明白這裡是提到上主的拯救.第九節說你將被差遣,百姓卻不接受你的信息.他們硬心到何時呢?直到城邑荒涼無人居住,他們將被遷到遠方,像樹被砍剩下樹根,但仍然存留聖潔的種類.第七章跟著提到將有童女懷孕生子,這是拯救的兆頭.當時北國以色列和亞蘭(敘利亞) 大軍壓境,迫使猶大國結盟對抗更強大的亞述國(伊拉克) ,以賽亞叫當時的亞哈斯(烏西雅的孫子) 不要依靠亞蘭和北國聯軍或亞述,不要歸邊,要依靠上帝的拯救,但亞哈斯不聽,並停止先知的職務.先知這時極失望,他對百姓失望,也看不到上帝的作為而失望.上主用以前他在崇拜中見過的異象來提醒他:上帝有恩典,上帝有拯救,雖然人不配,但上主有憐憫,必有那童女懷孕,他名稱為以馬內利--上主和我們同在,這就是福音.

所以崇拜若是信仰的表白,認罪,赦罪就是不能缺了的禮儀

.我們每主日透過認罪,被宣赦,愈來愈感到福音的甘甜,悔改不是信耶穌那一剎那的決定,乃每天的經驗,與基督同死同復活是我們每天的經驗,透過崇拜的禮儀重新演繹,深化我們的信仰.

3. 崇拜與事奉 and \#8212;我在這裡,請差遣我這段對話的末後說到先知聽到上主的呼喚:我可以差遣誰呢?誰可以為我而去呢?在崇拜結束時牧師也會說:你們離開這裡,進入世間,但記得你們不屬於世界.崇拜天主教稱彌撒Mass,宗教改教時期改教者也稱崇拜為彌撒,這是歷史悠久的名稱.彌撒一字的意思是

你們離開吧

!拉丁文Ite Missa  est (Thus you are dismissed) ,因為早期教會崇拜分兩時段,先是慕道班,通常要三年慕道後才被接受洗禮,我們現在大學畢業後讀三年神學就作傳道.後半段時候舉行聖餐,像一般我們現在的安排.但當時聖餐祇容許領洗的人出席,所以聖餐開始前襄禮的說:請你們離開吧!Ite  missa  est.離開教會就是崇拜的開始,這話甚有啟發.我們在崇拜結束時有接受祝福和被差遣.

祝福是美善的話語,英文benediction.進入世界前我們面對世界很多可能性,明天如何,我們不知道,我們像無限信賴的小孩,信賴掌管明天的上主,信賴他的恩慈,始祖因懷疑上主的美善話語而墮落,上帝說:你不可吃分別善惡樹的果子分別善惡就是代表成長獨立,你若獨立於上帝,與生命的源頭割裂,你就死亡,被逐出樂園.但重返樂園的人,被上帝救贖的人,對上帝的話語有無限的信賴,視他的話語為美善benediction.

接受祝福後,我們與人分享福氣,被差遣進入世界,目的是去服事有需要的人.十誡由兩塊法板組成,與上帝有關的宗教和與人有關的倫理.愛上帝與愛人不能分開.我們這教堂的設計別具匠心,這前面的玻璃幕牆使我們看到外面的世界,教會不能對世界的現況視若無睹.耶穌在曠野被魔鬼試探,魔鬼要耶穌拜他,就把整個世界給耶穌,不需要上十字架.耶穌對撒旦說:當拜主你的神,單要事奉他.敬拜和事奉不能分離,虔誠和善行不能分開.雅各書說:真正的虔誠,就是看顧在患難中的孤兒寡婦,並且保守自己不沾染世俗.(雅1:26)

誰能被主差遣呢?誰有資格呢?在這段經文中,我們看到那些感到自己的不配,深覺自己不潔的人,才適合被主差遣.那些體會上主對世界的心腸的人,才被主差遣.那些感到被赦免的人,心被恩感,才會被上主所差遣.差遣的路不是容易走的,但因為看到上主的榮耀,經歷被赦免的釋放和喜樂,他們願意說:我在這裡,請差遣我.阿們.



\newpage

\section{光明在人間}
\label{sec:38}
\textbf{鄧瑞強博士 - 「光明在人間」}
\newline
\newline
\hyperref[sec:37]{< < < PREV SERMON < < <}
~
\hyperref[sec:toc]{[返目錄]}
~
\hyperref[sec:39]{> > > NEXT SERMON > > >}


路加福音 9:21-36
\newline
\begin{longtable}{cl}
\hline
\hline
章節 & 經文 (和合本修訂版)\\
\hline
9:21 & \begin{tabularx}{0.7\textwidth}{X} 耶穌切切吩咐他們，命令他們不可把這事告訴任何人； \end{tabularx} \\ \\ \relax
9:22 & \begin{tabularx}{0.7\textwidth}{X} 又說：「人子必須受許多的苦，被長老、祭司長和文士棄絕，並且被殺，第三天復活。」 \end{tabularx} \\ \\ \relax
9:23 & \begin{tabularx}{0.7\textwidth}{X} 耶穌又對眾人說：「若有人要跟從我，就當捨己，天天背起自己的十字架來跟從我。 \end{tabularx} \\ \\ \relax
9:24 & \begin{tabularx}{0.7\textwidth}{X} 因為凡要救自己生命的，必喪失生命；凡為我喪失生命的，他必救自己的生命。 \end{tabularx} \\ \\ \relax
9:25 & \begin{tabularx}{0.7\textwidth}{X} 人就是賺得全世界，卻喪失了自己，或賠上自己，有甚麼益處呢？ \end{tabularx} \\ \\ \relax
9:26 & \begin{tabularx}{0.7\textwidth}{X} 凡把我和我的道當作可恥的，人子在自己的榮耀裡，和天父與聖天使的榮耀裡來臨的時候，也要把那人當作可恥的。 \end{tabularx} \\ \\ \relax
9:27 & \begin{tabularx}{0.7\textwidth}{X} 我實在告訴你們，站在這裡的，有人在沒經歷死亡以前，必定看見神的國。」 \end{tabularx} \\ \\ \relax
9:28 & \begin{tabularx}{0.7\textwidth}{X} 說了這些話以後約有八天，耶穌帶著彼得、約翰、雅各上山去禱告。 \end{tabularx} \\ \\ \relax
9:29 & \begin{tabularx}{0.7\textwidth}{X} 正禱告的時候，他的面貌改變了，衣服潔白放光。 \end{tabularx} \\ \\ \relax
9:30 & \begin{tabularx}{0.7\textwidth}{X} 忽然有摩西和以利亞兩個人同耶穌說話； \end{tabularx} \\ \\ \relax
9:31 & \begin{tabularx}{0.7\textwidth}{X} 他們在榮光裡顯現，談論耶穌去世的事，就是他在耶路撒冷將要完成的事。 \end{tabularx} \\ \\ \relax
9:32 & \begin{tabularx}{0.7\textwidth}{X} 彼得和他的同伴都打盹，但一清醒，就看見耶穌的榮光和與他一起站著的那兩個人。 \end{tabularx} \\ \\ \relax
9:33 & \begin{tabularx}{0.7\textwidth}{X} 二人正要和耶穌分離的時候，彼得對耶穌說：「老師，我們在這裡真好！我們來搭三座棚，一座為你，一座為摩西，一座為以利亞。」他卻不知道自己在說些甚麼。 \end{tabularx} \\ \\ \relax
9:34 & \begin{tabularx}{0.7\textwidth}{X} 說這些話的時候，有一朵雲彩來遮蓋他們；他們一進入雲彩就很懼怕。 \end{tabularx} \\ \\ \relax
9:35 & \begin{tabularx}{0.7\textwidth}{X} 有聲音從雲彩裡出來，說：「這是我的兒子，我所揀選的。你們要聽從他！」 \end{tabularx} \\ \\ \relax
9:36 & \begin{tabularx}{0.7\textwidth}{X} 聲音停止後，只見耶穌獨自一人。當那些日子，門徒保持沉默，不把所看見的事告訴任何人。 \end{tabularx} \\ \\
[1ex]
\hline
\hline
\end{longtable}
日期:2010年3月14日

主題:「光明在人間」

經文:路加福音9:21-36

証道:鄧瑞強博士

各位弟兄姊妹,主內平安.祝大家新年快樂,生活有目標,生命有意義.虎年像老虎一樣生猛,但不要咬人.在這情人節裡,祝願大家找到你愛的人,也有人愛你.

今日講道的經文,是今日的福音書經課,再加其前面的一段

經文:路加福音9:21-36.

讓我讀一次

經文:路加福音 9:21-36

9:21　耶穌切切的囑咐他們,不可將這事告訴人,

9:22　又說:「人子必須受許多的苦,被長老,祭司長,和文士棄絕,並且被殺,第三日復活.」

9:23　耶穌又對眾人說:「若有人要跟從我,就當捨己,天天背起他的十字架來跟從我.

9:24　因為,凡要救自己生命的,必喪掉生命；凡為我喪掉生命的,必救了生命.

9:25　人若賺得全世界,卻喪了自己,賠上自己,有什麼益處呢?

9:26　凡把我和我的道當作可恥的,人子在自己的榮耀裡,並天父與聖天使的榮耀裡降臨的時候,也要把那人當作可恥的.

9:27　我實在告訴你們,站在這裡的,有人在沒嘗死味以前,必看見神的國.」

9:28　說了這話以後約有八天,耶穌帶著彼得,約翰,雅各上山去禱告.

9:29　正禱告的時候,他的面貌就改變了,衣服潔白放光.

9:30　忽然有摩西,以利亞兩個人同耶穌說話；

9:31　他們在榮光裡顯現,談論耶穌去世的事,就是他在耶路撒冷將要成的事.

9:32　彼得和他的同伴都打盹,既清醒了,就看見耶穌的榮光,並同他站著的那兩個人.

9:33　二人正要和耶穌分離的時候,彼得對耶穌說:「夫子,我們在這裡真好!可以搭三座棚,一座為你,一座為摩西,一座為以利亞.」他卻不知道所說的是什麼.

9:34　說這話的時候,有一朵雲彩來遮蓋他們；他們進入雲彩裡就懼怕.

9:35　有聲音從雲彩裡出來,說:「這是我的兒子,我所揀選的,你們要聽他.」

9:36　聲音住了,只見耶穌一人在那裡.當那些日子,門徒不提所看見的事,一樣也不告訴人.

在9:22,主耶穌說:「

人子必須受許多的苦,被長老,祭司長,和文士棄絕,並且被殺,第三日復活.」說話簡單,道理卻深入.人間有種種權力,敵對他們不接受的東西.生命有意義,但難免艱難.死亡的力量真實無比,但又不是終極的,有復活的力量.

9:23　耶穌又對眾人說:「若有人要跟從我,就當捨己,天天背起他的十字架來跟從我.

9:24　因為,凡要救自己生命的,必喪掉生命；凡為我喪掉生命的,必救了生命.

9:25　人若賺得全世界,卻喪了自己,賠上自己,有什麼益處呢?

主耶穌指出,人人都想生命豐盛,但走錯了路,則生命不但不能豐盛,反令生命萎縮.這就像人人都想「美」,但錯誤地整容,就反為不美.看Michael Jackson的面目,總覺得有點假.這就像人人都想「健康」,但錯誤地注射藥物,則損害健康.主耶穌指出,有些人令生命豐盛的方式,其實是害了自己生命的.很多人努力去「賺得全世界」,努力加添一些生命外在的東西,但忘記了自己,遺忘了意義,生命仍是空虛萬分的.有些人月入十多萬,卻因買不到心目中的「豪宅」而不開心.有些人做一些簡單的義工,覺得生命豐盛得不得了.「人若賺得全世界,卻喪了自己,賠上自己,有什麼益處呢?」 在9:26,主耶穌說:「

凡把我和我的道當作可恥的,人子在自己的榮耀裡,並天父與聖天使的榮耀裡降臨的時候,也要把那人當作可恥的.」

這裡談到神聖的榮耀.當人輕視神聖的榮耀時,他自己生命應有的榮耀也消失了.現代家庭出現很多問題.當人結婚時,總許下結婚的誓言,這誓言是神聖的.當夫婦努力愛對方的時候,神聖的榮耀便在這家庭彰顯了.當一方倒瀉了一杯水,對方不會說:「你永遠都是這樣的.」只要對方不將「一次」看成「永遠」時,神聖的榮耀便出現了.當雙方起了很大衝突時,吵架吵得很大聲時,你卻用更大的聲音說:「我是不會和你離婚的,因為我愛你.」這時,神聖的榮耀便出現了.但當我們說,或心裡說:「我想你死」,則一切光明都消失了,神聖的約被遺忘了,這時,我們生命應有的榮耀也消失了.

在這世上,我們能看見這神聖的榮耀嗎? 在9:27,耶穌說:「

我實在告訴你們,站在這裡的,有人在沒嘗死味以前,必看見神的國.」 不必死後,才可見「神的國」.「神的國」是在人間的.神的榮耀是在人間的.人看不見,只是人用種種方式消滅這神的榮耀而已.光明實在是在人間.

接下來的經文,便講到神的光明在人間,但人會不自覺地將之隱藏起來.

9:28-31記載,神的光明在人間顯現.

9:28　說了這話以後約有八天,耶穌帶著彼得,約翰,雅各上山去禱告.

9:29　正禱告的時候,他的面貌就改變了,衣服潔白放光.

9:30　忽然有摩西,以利亞兩個人同耶穌說話；

9:31　他們在榮光裡顯現,談論耶穌去世的事,就是他在耶路撒冷將要成的事.

在山上,在接近神的地方,耶穌祈禱親近神.他的面貌改變,衣服放光.舊約記載,摩西因接近神的緣故,面貌放光.這時,人反映著神的光明.人的生命,應如鏡一樣,反映神的光明.但這生命的鏡若蒙塵,或這鏡背向神,便反映不到神的光明了.耶穌這時流露著神的光明.他和舊約最偉大的先知摩西和以利亞在一起,他們談論耶穌受死的事.耶穌的受死,乃是不執著自己的生,而甘願為人而死的事.人生在世,不求獲取,只求分享.不望自己成功,只望讓神的愛流露出來.這是耶穌的生命,這是神的光明在人間.

各位弟兄姊妹,不要灰心失望,因為神的光明在人間.

然而,門徒的反應,卻是令人失望的.

9:32　彼得和他的同伴都打盹,既清醒了,就看見耶穌的榮光,並同他站著的那兩個人.

9:33　二人正要和耶穌分離的時候,彼得對耶穌說:「夫子,我們在這裡真好!可以搭三座棚,一座為你,一座為摩西,一座為以利亞.」他卻不知道所說的是什麼.

聖經的記載很諷刺,當神聖的光明展露時,他們是在「沉睡」中.彼得說了一番話,但聖經的評語是,「他卻不知道所說的是什麼」.

以彼得為首的門徒,他做錯了什麼.若彼得的反應也代表我們的反應,則我們做錯了什麼. 彼得的反應是在這山上搭三座棚,如此,這永恆的時刻便被保留了,這神聖的光明便被固定了,這神聖的臨在便被永遠擁有了.問題正正出在這裡. 當你試圖固定一有價值的東西,將之放在手上被你掌握時,這東西的價值便消失了,這東西便變成一商品,而你卻變成消費這商品的消費者.

你喜歡見到人在你面前「笑」.這很好.笑真是值得欣賞的.笑是無價的.但是,若你想用一個棚,去捉住「笑」,供你隨時把玩時,這時,便變成「買笑」.「笑」的價值消失了,神聖消失了,而你則不再是笑的欣賞者,而是「買笑」的消費者.

我們喜歡美少女.這很好.美少女是值得欣賞的.但若我們想捉住她,將之變成「美少女攬枕」,則美少女的美消失了,她變成消費品,供消費者消費.

「神聖的事物」與「模仿神聖事物的消費品」之間的分別,在於「神聖的事物」吸引我們走出自私的自我,走出自我的狹窄空間,去活在無限的空間裡,去接近神聖,去體會生命的光明.但「模仿神聖事物的消費品」則是叫我們更自我,更自私,將外在有價值的東西供我們去奴役,視外在有價值的東西為為我們而存在的東西,自我更膨脹,外在有價值的事物的價值更被漠視.越消費,越看不到價值,生命越萎縮.

神太偉大,我們控制不到,最好將之縮少,使我們能操控得到.神的光明時隱時顯,我們控制不到,最好加以掌握,使神的光明能留在我們能操作的範圍內.神的要求太複雜,我們辦不到,最好將之變成能增加我們健康和財富的金石良言.這樣,我們便會甘心樂意去實行.簡言之,神呀,你最好成為我的消費品.我要你來的時候,你就來.不要你的時候,你就走.我要消費你的時候,是希望你帶給我們愉快的享受,給我們滿足,千萬不能對我們有任何捨己的要求.若你不能滿足我的情感要求,我的投訴是合理的,正如我投訴不合格的消費品一樣.故此,若我祈禱,你要我等待,我是要投訴的.因我買一個消費品,要得到即時的滿足,怎能等待?難道我點一個菜,要我等十年嗎?我是消費者,你沒有聽說嗎:「顧客是上帝」.我不開心時,你要叫我開心.請不要和我講什麼意義,我只求開心而已.我買一張高價按摩椅回來,難道我是要求什麼意義嗎?我只求開心而已.

當你用一個棚將神聖的光明留住時,神聖的光明便消失了.消費文化的一個特點,就是人藉消費「虛擬地」將他想要的事物留住,供他享用.青春留逝了,不怕,用這神奇液體啦,它會將你的皮膚回復十年前的光景.你的頭髮消逝了,不怕,用這神奇產品啦,它會召回你失落的頭髮.你失去了童年,不怕,去迪士尼啦,你失落的童年留在那裡,待你認領.只要你消費,一切都會回來.永恆就在這裡.這樣的消費,令人遠離了真實的生活.遠離了真實的生活,遠離山下,就不能再體會神聖的光明了.

彼得希望留在山上,以一個棚盛載神聖的光明.但是,耶穌所展示的神聖光明,其實是在山下的十架路上.在看似艱難的人生中,與窮人同在,與受苦者同行,與釘在十架上的犯人分享永生的盼望.遠離這一切,怎會見到神聖的光明呢?光明在人間,在真實的人間,不是在山上,在可供消費的消費品裡.

消費文化產生的虛擬世界,讓人在消費中得到虛擬的滿足,卻遠離了生命應有的豐盛.情人節,是消費的高峰時刻.有人買一百枝玫瑰,送給女朋友.買的人,認為這是愛的表達.接受的人,認為這是愛的送贈.雙方在這一百枝玫瑰裡,獲得幸福和滿足.然後,雙方等下一次的一百枝玻瑰.沒有下一次的一百枝玻瑰,雙方的空白日子怎麼過? 但是,愛根本不在這裡.愛,是在沒有玫瑰的日子,是在對方不再好看的日子,是在相處的感覺不再良好的日子,仍努力以捨己的行動,去澆灌對方的心田,耐心等待開花的日子.愛,在沒有玫瑰的日子裡.

當然,我們能掌握的,能控制的,是在我們高興時,買一百枝玫瑰.我們能把握的,是消費給我的力量.至於,我能否有捨己的愛的品德,這是控制不了的事情,在消費文化裡,這也是人容易遺忘的事情.只求表面,追求可掌控,拒絕真情付出,是消費文化的特點.但是,擁抱真實的人生,關懷真實的人,生命有確實的投入,這個人才會與神聖的光明有真實的相遇.能像消費品一樣被控制的神聖,不再是神聖.唯有不再由你控制,卻需要你付出,這時神聖才會出現.聖經指彼得「不知道所說的是什麼」,是他的想法不單不能留住神聖,反而遠離神聖.

為了留住想要的東西,消費文化鼓勵不斷消費.你要美貌,請不斷買我們的神奇液體.你要留住時尚,請不斷消費我們的時尚產品.你要留住那種成功的感覺,請繼續買我們的奢侈品,讓人羨慕你視錢財如糞土.消費品日新月異,買之不盡,以消費品去滿足自己,總令人花九牛二虎之力卻一無所得.到頭來,買了全世界,卻失去了真正的自己,這有什麼益處呢?

這個故事的結尾:

9:34　說這話的時候,有一朵雲彩來遮蓋他們；他們進入雲彩裡就懼怕.

9:35　有聲音從雲彩裡出來,說:「這是我的兒子,我所揀選的,你們要聽他.」

雲彩代表神的臨在,這雲彩遮蓋人,它比人更大. 這是人無法掌控的神秘.天上的聲音說,耶穌的生命是神揀選的,我們要聽他.我們不能將在耶穌身上顯示的光明放在我們的棚裡,乃是我們要放下我們的框框,聽從他. 主耶穌在人間中行走,光明在他那裡.尋找光明的人,請跟從他.

但願榮耀歸於聖父,聖子,聖靈.阿們.



\newpage

\section{再細看試探的臉龐}
\label{sec:39}
\textbf{林豪恩先生 - 「再細看試探的臉龐」}
\newline
\newline
\hyperref[sec:38]{< < < PREV SERMON < < <}
~
\hyperref[sec:toc]{[返目錄]}
~
\hyperref[sec:40]{> > > NEXT SERMON > > >}


路加福音 4:1-13
\newline
\begin{longtable}{cl}
\hline
\hline
章節 & 經文 (和合本修訂版)\\
\hline
4:1 & \begin{tabularx}{0.7\textwidth}{X} 耶穌滿有聖靈，從約旦河回來，聖靈把他引到曠野， \end{tabularx} \\ \\ \relax
4:2 & \begin{tabularx}{0.7\textwidth}{X} 四十天受魔鬼的試探。在那些日子，他沒有吃甚麼，日子滿了，他餓了。 \end{tabularx} \\ \\ \relax
4:3 & \begin{tabularx}{0.7\textwidth}{X} 魔鬼對他說：「你若是神的兒子，叫這塊石頭變成食物吧。」 \end{tabularx} \\ \\ \relax
4:4 & \begin{tabularx}{0.7\textwidth}{X} 耶穌回答說：「經上記著：『人活著，不是單靠食物。』」 \end{tabularx} \\ \\ \relax
4:5 & \begin{tabularx}{0.7\textwidth}{X} 魔鬼又領他上了高山，霎時間把天下萬國都指給他看， \end{tabularx} \\ \\ \relax
4:6 & \begin{tabularx}{0.7\textwidth}{X} 對他說：「這一切權柄和榮華我都要給你，因為這原是交給我的，我願意給誰就給誰。 \end{tabularx} \\ \\ \relax
4:7 & \begin{tabularx}{0.7\textwidth}{X} 你若在我面前下拜，這一切都歸你。」 \end{tabularx} \\ \\ \relax
4:8 & \begin{tabularx}{0.7\textwidth}{X} 耶穌回答他說：「經上記著：『要拜主—你的神，惟獨事奉他。』」 \end{tabularx} \\ \\ \relax
4:9 & \begin{tabularx}{0.7\textwidth}{X} 魔鬼又領他到耶路撒冷去，叫他站在聖殿頂上，對他說：「你若是神的兒子，從這裡跳下去！ \end{tabularx} \\ \\ \relax
4:10 & \begin{tabularx}{0.7\textwidth}{X} 因為經上記著：『主要為你命令他的使者保護你； \end{tabularx} \\ \\ \relax
4:11 & \begin{tabularx}{0.7\textwidth}{X} 他們要用手托住你，免得你的腳碰在石頭上。』」 \end{tabularx} \\ \\ \relax
4:12 & \begin{tabularx}{0.7\textwidth}{X} 耶穌回答他說：「經上說：『不可試探主—你的神。』」 \end{tabularx} \\ \\ \relax
4:13 & \begin{tabularx}{0.7\textwidth}{X} 魔鬼用完了各樣的試探，就離開耶穌，再等時機。 \end{tabularx} \\ \\
[1ex]
\hline
\hline
\end{longtable}
日期:2010年3月21日

主題:「再細看試探的臉龐」

經文:路加福音4:1-13

証道:林豪恩先生

經文

1耶穌被聖靈充滿,從約旦河回來,聖靈將他引到曠野,2四十天受魔鬼的試探.那些日子沒有吃甚麼；日子滿了,他就餓了.3魔鬼對他說:「你若是上帝的兒子,可以吩咐這塊石頭變成食物.」4耶穌回答說:「經上記著說:『人活著不是單靠食物,乃是靠上帝口裡所出的一切話.』」5魔鬼又領他上了高山,霎時間把天下的萬國都指給他看,6對他說:「這一切權柄,榮華,我都要給你,因為這原是交付我的,我願意給誰就給誰.7你若在我面前下拜,這都要歸你.」8耶穌說:「經上記著說:當拜主－你的上帝,單要事奉他.」 9魔鬼又領他到耶路撒冷去,叫他站在殿頂上,對他說:「你若是上帝的兒子,可以從這裡跳下去；10因為經上記著說:主要為你吩咐他的使者保護你；11他們要用手托著你,免得你的腳碰在石頭上.」 12耶穌對他說:「經上說:『不可試探主－你的上帝.』」13魔鬼用完了各樣的試探,就暫時離開耶穌. 引言教授生涯的三個試探:一位教授的讀經領受1．石頭變食物新入職教授的挑戰是「生存」(survival).如何能夠把文章變為食物,維持自己的生命,保留自己的教席,是年輕教授的關注.學術界的生存規則十分嚴厲:Publish or Perish (不出版,就出局).面對這項挑戰,有人辛辛苦苦,認認真真地做研究,把研究結果出版,以文章換取生存所需；有人隨隨便便,甚至虛作數據,含混過關.這位教授認為,教授是要靠文章維生的,為了生存而作虛弄假,就是試探的所在.

2．萬國的權柄和榮華教授生存的問題解決了,教席坐穩了,往往也累積了相當的資歷.在這個時候,主動或被動地,教授都開始增添了行政的任務,承擔了某層次的管理角色.這些行政任務及管理角色賦予了教授權力,審核其他同事的表現,甚至決定其他同事的升遷及調動.權力是管治的必須,但同時又帶來引誘.萬人之上的感覺,可以是很大的試探.該教授憶述一次應邀往一個國際會議演講的經驗:演講廳萬多座位座無虛席,他的照片被投射到多個超大的螢幕,加上一萬人的鼓掌聲,突然間令他有大地在我腳下的感覺.這次感覺,令他產生危機感,反省自己,到底仍然在事奉上帝,還是已經轉變為敬拜自己.教授認為,這是老授生涯的第二個試探.

3．從殿頂跳下來生存早已不成問題,而且地位穩固,又有知名度的時候,教授指出,另一個階段的試探又來了.在這個時候,有人可能想到要做一些石破天驚的事情,一些能人所不能的事情,令舉世觸目.有些人就在這個試探中失足,好像幾年前,南韓某位科學家早於歐美地區的同行發表某項複製技術的研究結果,為南韓及亞洲爭光,成為英雄.可惜,這位科學家後來卻要為作假而付上沉重的代價.超越常人,能人所不能,教授認為這些慾望為某些尖端教授製造試探.

釋經與應用耶穌的三個試探1．1．石頭變食物

3魔鬼對他說:「你若是上帝的兒子,可以吩咐這塊石頭變成食物.」 如果魔鬼用食物來試探人,對某些人來說的確可以構成試探,然而,有些人卻能夠堅持「士可殺不可辱」的精神,寧死不屈.如果魔鬼也用食物來試探耶穌,未免太輕視耶穌了吧!按福音書的記載,當人們對耶穌亳不認識的時候,魔鬼及污靈已經深知耶穌是誰,所以,魔鬼應該不會無知得以為耶穌會臣服於他的波板糖之下.退一萬步來說,就算耶穌真的把石頭變成了食物,又有甚麼問題呢?這根本是他能力範圍之內的事情,完全不需要求助於魔鬼.那麼,石頭變食物在甚麼意義上是耶穌的試探呢?到底在試探甚麼呢?

1．2．天下的萬國

5魔鬼又領他上了高山,霎時間把天下的萬國都指給他看,6對他說:「這一切權柄,榮華,我都要給你,因為這原是交付我的,我願意給誰就給誰.7你若在我面前下拜,這都要歸你.」 天下的萬國,對不少野心家來說,的確很吸引,以致願意為之而生,為之而死.如果魔鬼以此來試探人類當中的野心家,很可能有效.如果魔鬼以此來試探你,或者沒有作用,因為天下的萬國對你來說可能不是獎品,而是重擔.然而,魔鬼以此來試探耶穌,會否又輕看了耶穌呢?難道魔鬼以為耶穌是野心家嗎?退一萬步來說,如果耶穌得到萬國的權柄及榮華,又有甚麼問題呢?如果耶穌是更好的管治者,萬國的權柄及榮華由他擁有,豈不是比由那些暴君來擁有更好嗎?那麼,天下的萬國在甚麼意義上是耶穌的試探呢?到底在試探甚麼呢?

1．3．從殿頂跳下

9魔鬼又領他到耶路撒冷去,叫他站在殿頂上,對他說:「你若是上帝的兒子,可以從這裡跳下去；10因為經上記著說:主要為你吩咐他的使者保護你；11他們要用手托著你,免得你的腳碰在石頭上.」  從殿頂跳下其實是充滿刺激性及娛樂性的項目,這些行動,通常都是魔術師(如大衛高拍飛),或者是蜘蛛俠的表演.對一般人來說,寧願作觀眾,而不願接受這樣的挑戰.如果魔鬼以此來試探你,我相信魔鬼不會得逞,一來你很可能不敢冒險,二來你會認為不值得冒這個險,因為這行動對你來說沒有甚麼意義.所以,我們實在不明白為甚麼魔鬼以為從聖殿的高處跳下來可以成為耶穌的試探.難道魔鬼以為耶穌要做特技人?難道魔鬼以為耶穌要入娛樂圈?又再退一萬步來說,如果耶穌真的從殿頂跳下而絲毫不損,又有甚麼問題呢?正如把石頭變食物一樣,這根本就是耶穌能力所及的事情.然則,從殿頂跳下在甚麼意義上能夠成為耶穌的試探呢?到底在試探甚麼呢?

三個試探對耶穌的意義馬太,馬可,路加三卷福音書共同見證,耶穌受試探之前,剛剛接受過施洗約翰的洗禮.三卷福音書都清楚記載耶穌受洗時所出現的特殊景象:

「天就開了,聖靈降臨在他身上,形狀彷彿鴿子；又有聲音從天上來,說:你是我的愛子,我喜悅你.」(路加福音3:21-22)按福音書的記載,這是耶穌第一次從上帝那裡公開地聽到自己的身份──上帝的兒子.上帝的兒子這個身份意味著甚麼呢?上帝的兒子的使命是甚麼呢?上帝的兒子在人類歷史中出現所為何事呢?這才是魔鬼要試探耶穌的地方.

2．1．把石頭變成食物──試探上帝的兒子放棄道成肉身

3魔鬼對他說:「你若是上帝的兒子,可以吩咐這塊石頭變成食物.」

魔鬼的弦外之音可能是:「如果你不能把這塊石頭變成食物來填飽自己「打鼓」的肚腹,你還是上帝的兒子嗎?」耶穌剛剛聽到天上來的聲音說自己是上帝的兒子,上帝的兒子吩咐石頭變成食物,只是舉手之勞.問題是,要不要運用神能來證明自己是神的兒子呢?要不要運用神能來解決自己的困難呢?如果要運用神能,耶穌一定有這個能力,而且人生會變得一帆風順,心想事成.不過,同時也顯示出耶穌重新做回神,而放棄了做人.如果耶穌取回神能,放棄做人,那麼,道並沒有變成肉身,上帝的拯救計劃不攻自破.魔鬼關心的,當然不是耶穌肚餓還是肚飽,魔鬼害怕的,是上帝的拯救計劃能否實現.魔鬼建議耶穌在肚餓時運用神能把石頭變成食物,其實是在引誘耶穌要做神,而不要做人.如果魔鬼這個試探成功,就足以令耶穌放棄使命,放棄他來到世上的目的.

2．2．求取天下的萬國──試探上帝的兒子重掌王權

5魔鬼又領他上了高山,霎時間把天下的萬國都指給他看,6對他說:「這一切權柄,榮華,我都要給你,因為這原是交付我的,我願意給誰就給誰.7你若在我面前下拜,這都要歸你.」 魔鬼以「天下的萬國」作為條件來引誘耶穌,但是,誰說「天下的萬國」是屬於魔鬼的,以致魔鬼可以視之為私有財產來出讓呢?「天下的萬國」屬於魔鬼,只是魔鬼自己說的.魔鬼如此的宣稱是否值得懷疑呢?事實又是怎樣的呢?新約聖經歌羅西書21章15-16節如此宣告:「愛子是那不能看見之上帝的像,是首生的,在一切被造的以先.因為萬有都是靠他造的,無論是天上的,地上的,能看見的,不能看見的；或是有位的,主治的,執政的,掌權的；一概都是藉著他造的,又是為他造的.」即是說,「天下的萬國」根本就不屬於魔鬼的,魔鬼的行為等同於銷售一件不屬於他的物件,諷刺的是,這物件其實是他的銷售對象本身所擁有的.

我們或者以為天下的萬國是屬於魔鬼的,所以經常有以與魔鬼交易為主題的文學及影視作品面世.耶穌應該不會有這樣的誤會吧!耶穌應該知道天下的萬國不是屬於魔鬼的,而是屬於他自己的.魔鬼在提醒耶穌:「耶穌,天下的萬國本來就是你的,你要不要確認你的權柄呢?耶穌,其實你不需要做乞丐王子,也不需要扮奴僕,你大可以宣認你的權柄,奪回你的王位.耶穌,我知道你有拯救萬民的崇高使命,然而,你無權無勢,孤掌難鳴,憑甚麼拯救萬民呢?如果你手握絕對的政治權力,指揮萬國的軍事力量,調配所有的經濟資源,豈不是更直接及更有效地實踐你的使命嗎?何必受僕人之苦呢?何必透過受刑罰來使人們得平安呢?何必透過受鞭傷來使人們得醫治呢?何必透過死亡來赦免人們的罪孽呢?(以賽亞書53)」如果耶穌接受魔鬼的建議,立即接管天下萬國,就等於放棄「奴僕的形象」,不願「謙卑自己」,拒絕「存心順服」,逃避「死在十字架上」的安排(腓立比書2:7-8),即是等於引誘耶穌放棄使命,放棄神的兒子來到人間的目的,對耶穌來說,這才是致命的.

2．3．從殿頂跳下──試探上帝的兒子對上帝的信任

9魔鬼又領他到耶路撒冷去,叫他站在殿頂上,對他說:「你若是上帝的兒子,可以從這裡跳下去；10因為經上記著說:主要為你吩咐他的使者保護你；11他們要用手托著你,免得你的腳碰在石頭上.」  相信大家都同意,耶穌不是特技人,不需要練習高空跳躍；耶穌不是娛樂圈中人,不需要做特殊動作吸引人的注意；耶穌又不是某產品的代言人,不需要做市場推廣.如果耶穌要從殿頂跳下來而絲毫無損,唯一的目的是證明自己仍得到上帝的保護.為甚麼要如此證明自己仍然得到上帝的保護呢?因為在第一個試探中,耶穌堅持道成肉身,「成為人的樣式」(腓立比書2:7),因此,如果從殿頂跳下,必定會如一般人一樣不死則傷.然而,從殿頂跳下而不受傷,證明自己仍然得到上帝的保護,證明上帝的應許仍然信得過.

耶穌貴為上帝的兒子,本來擁有無限的神能,也擁有無限的權力及資源,來到人間,成為人的樣式,令人聯想到電影中的警察臥底,交出配槍,放下委任證,放棄所有保護自己的裝備,變成無權無械.他們的人身安全,完全依賴警方的保護.警方的保護信得過嗎?警方會遵守他們的承諾嗎?三年又三年,要不要測試一下警方是否仍然知道自己的存在呢?要不要測試一下警方是否仍然重視自己呢?要不要肯定一下自己仍然是警察的身份呢?如果要測試,反映臥底與警方怎樣的關係呢?如果要測試,對行動又會有甚麼影響呢?

要不要測試上帝對自己的感情呢?要不要測試上帝對自己的重視呢?要不要測試上帝對自己的保護呢?要不要測試上帝的應許呢?如果上帝的兒子也不信任上帝,這是徹底的失敗,所有的使命及意義都不用多談了.這樣的試探,對耶穌來說是致命的,足以摧毀耶穌道成肉身的使命,使耶穌降臨人間成為一個笑話.

3．耶穌的試探對人類的意義如果我們從上述的分析來看魔鬼對耶穌的試探,試探雖然有不同的層次,透過不同的行動,但試探的目的就是要更改耶穌作為上帝的兒子這個身份的意義及使命,離間上帝的兒子與上帝的關係,引誘上帝的兒子懷疑上帝,挑戰上帝,最後叛逆上帝.如此,魔鬼就能夠摧毀上帝交給耶穌的使命,取消耶穌道成肉身的意義,阻止上帝對人類的拯救行動.

路加福音告訴我們,魔鬼上述的計謀曾經成功,魔鬼似乎相信「橋唔怕舊,最緊要受」的道理,照版煮碗地運用在耶穌身上.為甚麼如此說呢?讓我們看看路加福音記述耶穌受試探的前文:第3章21-22節記述耶穌受洗,結束的句子是:「你是我的愛子,我喜悅你.」接著是讀者通常會跳過的段落,這個段落記述耶穌的家譜.若我們錯過這個段落,就錯過了重要的信息.請留意家譜的最後一句:「亞當是上帝的兒子.」下文記述魔鬼對耶穌的試探,就是環繞著「上帝的兒子」這個身份.

究竟上帝的兒子會否信任上帝呢?上帝的兒子會否聽從上帝呢?上帝的兒子會否執行上帝的使命呢?這是自從有人類以來魔鬼要試探人的地方.且看上帝的兒子,人類的代表亞當和夏娃如何在這個試探失足:

1耶和華上帝所造的,惟有蛇比田野一切的活物更狡猾.蛇對女人說:「上帝豈是真說不許你們吃園中所有樹上的果子嗎?」2女人對蛇說:「園中樹上的果子,我們可以吃,3惟有園當中那棵樹上的果子,上帝曾說:『你們不可吃,也不可摸,免得你們死.』」4蛇對女人說:「你們不一定死；5因為上帝知道,你們吃的日子眼睛就明亮了,你們便如上帝能知道善惡.」6於是女人見那棵樹的果子好作食物,也悅人的眼目,且是可喜愛的,能使人有智慧,就摘下果子來吃了,又給她丈夫,她丈夫也吃了.(創世記3章1－6)創世記告訴我們,上帝的兒子亞當失足於魔鬼的試探,福音書告訴我們,上帝的兒子耶穌成功擊退了魔鬼類似的試探.亞當作為代表,展示人類如何破壞與上帝之間的信任；耶穌也作為代表,展示人類可以如何信任上帝.保羅在羅馬書如此總結:

「因一次的過犯,眾人都被定罪；照樣,因一次的義行,眾人也就被稱義得生命了.因一人的悖逆,眾人成為罪人；照樣,因一人的順從,眾人也成為義了.」(5:18-19)

4．總結耶穌被魔鬼試探的事件,讓我們得到甚麼啟發呢?

魔鬼對耶穌的試探項目是為耶穌而設的,魔鬼應該不會照樣地試探我們.難道魔鬼會當我們肚餓時建議我們把石頭變成食物嗎?難道魔鬼會把萬國的權柄和榮華交給我們嗎?難道魔鬼會帶你到殿頂叫我們表演高空跳躍嗎?魔鬼會如何試探我們?甚麼能夠成為我們的試探?可能沒有人比我們自己更清楚.魔鬼試探耶穌的項目本身不一定是邪惡,不過,試探的目的是要耶穌放棄上帝交給他的使命,放棄他的人生目標,取消他的人生意義.上帝讓我們來到世上,生活於此時此地,我們的人生目標是甚麼?我們的人生使命是甚麼?如果能夠令我們亳無意義地虛度一生,對自己充滿怨憤及滿懷遺憾地離開,可能反映魔鬼對我們的試探成功了.人生使命,人生目標,生存意義這些詞語不一定是抽象的概念,不一定是難以捉摸的事情.耶穌的人生使命,人生目標,及生存意義,是與他的身份──上帝的兒子有關於,也與他的稱號──基督(即是彌賽亞)有關的,當他實踐他的身份的時候,他就在經歷他的人生使命,人生目標及生存意義.或者,我們也可以從自己的身份開始思考,從實踐我們的身份開始行動,擊退魔鬼的試探.耶穌的身份不是抽空的,也不是隨意的；耶穌的身份是與上帝相關的,是上帝所賜予的.能否接受這身份,願意否實踐這身份,反映著耶穌對上帝的信任.我們能否從這個角度來看待自己的身份?我們能否發現自己的身份與上帝的關係?我們能否信任上帝賜給我們這些身份,並且無論順境逆境都能安心立命,坦然承擔?



\newpage

\section{認信背棄和重建}
\label{sec:40}
\textbf{賀志勇博士 - 「認信,背棄和重建」}
\newline
\newline
\hyperref[sec:39]{< < < PREV SERMON < < <}
~
\hyperref[sec:toc]{[返目錄]}
~
\hyperref[sec:41]{> > > NEXT SERMON > > >}


日期:2010年3月28日

主題:「認信,背棄和重建」

經文:路加福音22:31-62

証道:賀志勇博士

經文

22:31 主又說,西門,西門,撒旦想要得著你們,好篩你們,像篩麥子一樣．22:32 但我已經為你祈求,叫你不至於失了信心．你回頭以後,要堅固你的弟兄.22:33 彼得說,主阿．我就是同你下監,同你受死,也是甘心.22:34 耶穌說,彼得,我告訴你,今日雞還沒有叫,你要三次說不認得我.22:49 左右的人見光景不好,就說,主阿,我們拿刀砍可以不可以.22:50 內中有一個人,把大祭司的僕人砍了一刀,削掉了他的右耳.22:54 他們拿住耶穌,把他帶到大祭司的宅裡．彼得遠遠的跟著.22:55 他們在院子裡生了火,一同坐著．彼得也坐在他們中間.22:56 有一個使女,看見彼得坐在火光裡,就定睛看他,說,這個人素來也是同那人一夥的.22:57 彼得卻不承認,說,女子,我不認得他.22:58 過了不多的時候,又有一個人看見他,說,你也是他們一黨的．彼得說,你這個人,我不是.22:59 約過了一小時,又有一個人極力的說,他實在是同那人一夥的．因為他也是加利利人.22:60 彼得說,你這個人,我不曉得你說的是什麼．正說話之間,雞就叫了.22:61 主轉過身來,看彼得.彼得便想起主對他所說的話,今日雞叫以先,你要三次不認我.22:62 他就出去痛哭.

今天要講的經文是路加福音22章31-34節,以及54-62節,請大家翻到路加福音22章31-34節,以及54-62節.

``滾滾長江東逝水,浪花淘盡英雄.是非成敗轉頭空,青山依舊在,幾度夕陽紅." 這個世界上,做英雄容易不容易?當然不容易!當英雄,有時候是要拿命來換的.比方我們的候任校長沈祖堯先生,被人稱為"SＡＲＳ英雄",那可是人家在香港最危難的時候,冒著生命危險去救人所得的榮譽.再比方抗日戰爭時期,我們有很多搞情報工作的地下英雄,被日本人抓住之後,無論是用電擊,還是用針刺,甚至是殺頭,都不出賣自己人,真是了不起的英雄!老實說,小時候看抗日的電影,我一方面看得熱血沸騰,為英雄叫好；另一方面,我很慶倖那個被打的不是我,很慶倖我是生在現在這個幸福的時代,不需要擔心自己被日本人抓起來.否則不要說用針刺我、用電打我啊,可能也就是拿鞭子抽個幾十下,我就什麼都招了.做英雄,不容易啊!

或許你跟我都不太想做英雄,可是世上偏有些人想做英雄.比方說項羽,比方說彼得.　彼得就是個想做英雄的人,或者說,彼得是個行動派,是個積極的行動派.他行動很迅速.當耶穌第一個呼召彼得時,他本來是在打魚的,一聽到呼召,立刻放下網就跟從了耶穌,你看看,決心很大吧!要是現在牧師呼召你做傳道人,你肯馬上辭掉工就跟從麼?某個夜晚,當耶穌在加利利海面行走時,是彼得第一個"從船上下去,在水面上走,要到耶穌那裡去."　在凱撒利亞腓利比,當耶穌問眾人,"你們說我是誰?"又是西門彼得第一個回答說:"你是基督,永生神的兒子." 連耶穌都誇獎他是有福的.彼得行動很迅速,反應很快,做事總是沖在前面,這已經具備了做英雄彼得的素質.

你說,這些都只能說明彼得有幹勁,可以做個好領導,但說他想做英雄未免誇張.好,我們再來看看我讀的後兩段,在耶穌基督被賣的那夜,當耶穌基督為彼得的信心擔心時,

22:33  ``彼得說,主阿．我就是同你下監,同你受死,也是甘心."你看看,多麼豪言壯語.而且,我們的彼得不僅僅是如此說,他也的確是行動了.當大祭司派的兵丁來捉拿耶穌時,

22:49    左右的人見光景不好,就說,主阿,我們拿刀砍可以不可以.     22:50   內中有一個人,把大祭司的僕人砍了一刀,削掉了他的右耳.這個拔刀削掉人家耳朵的那位,約翰就明確指出,是彼得!面對強敵,第一個拔刀出來保護耶穌的還是彼得.還有,在門徒驚惶四散的時候,彼得也還是捨不得耶穌,遠遠地跟去看.你說,在當時當地,你有沒有勇氣第一個拔刀出來保護耶穌?!你說,在當時當地,彼得算不算英雄?

可是,非常奇怪的是,就在短短一兩個小時之後,面對這麼多兵丁都敢拔刀出來保護耶穌的彼得,卻三次不認耶穌.我說大概一兩個小時之後,是因為我走過那段路,從客西馬尼園走到大祭司住的地方,不超過一個小時.這讓我們很奇怪,這個彼得,還是不是那個信誓旦旦說同耶穌坐牢,同耶穌一起受死都甘心的彼得?還是不是那個面對強敵都要拔刀保護耶穌的英雄彼得?

更為諷刺的是,當彼得面對的威脅越來越小的時候,他的否認卻一次比一次強.第一次,是大祭司的一個使女說彼得跟耶穌是一夥的,危險性當然是最大；第二次,有個人又指認彼得跟耶穌是一夥的,其他幾卷福音書指明,那個人還是一個使女,她不過是跟院子裡其他人說彼得跟耶穌是一夥的,小道消息,危險性降低了；第三次,有一個人是憑藉彼得的加利利口音,來判斷彼得跟耶穌是一夥的,這就好像人家憑藉我會說四川話,判斷我是個四川人一樣,做不得准,威險性也是最小的.可是,面對的危險性是越來越小,彼得不認主的口氣卻是越來越激烈.第一次,彼得說:"女子,我不認得他."第二次彼得聲調就提高了,"你這個人!我不是."第三次,彼得更是氣急敗壞了,彼得說,"你這個人!我不曉得你說的是什麼!"你聽聽,當危險性逐步下降的時候,我們的彼得不認主的聲調卻越來越高.

以前,每次讀到這裡,我都非常奇怪!我常常想,這是彼得麼?這是那個放下漁網跟從耶穌的彼得麼,這是那個大聲說"你是永生神的兒子"的彼得麼?這是那個拔刀保護主耶穌的英雄彼得麼?這怎麼可能!這怎麼可能!這麼愛主的一個人,就這麼一兩個小時,就開始說他不認識主了!面對人家的質問,他完全可以不吭聲,他完全可以不回答,可是他回答了,而且是一次比一次大聲的不認主!怎麼回事,怎麼回事,英雄彼得怎麼突然間變得這麼軟弱.這跟我印象中的英雄太不一樣了.以前看抗日英雄的電影,人家抗日英雄可是怎麼折磨都不屈服的啊!就算是像我這樣的,也起碼要挨了幾鞭子打才有可能屈服嘛.怎麼一個這麼愛主,又這麼勇敢的彼得,面對這樣一般性的危險,鞭子都沒挨,就不認主了!

只有信主這些年後,我才漸漸明白了生活與電影的差別.我才明白了聖經上的那個彼得不僅僅是彼得,那也是我們,是活在這世上的每一個人!我常常看到決志的時候興奮得不得了的人,決志之後因為一點挫折,又迅速離主而去；我常常看到團契裡的學生互相發的Email,很多都是說,因為我很忙,因為我要去參加別的什麼什麼活動,所以很遺憾很抱歉,團契活動就不能參加了--是啊,也許前一個小時我們還想著要愛主,要參加團契,這一個小時我就在忙著實驗,或者忙著去參加各種學生活動了.有時候,這個世界的誘惑比主的吸引力更大,這個世界常常讓我們前一個小時還熱血沸騰的說要愛主,下一個小時就把主的名拋在一邊.我們常常是剛剛讀完"愛是恆久忍耐,"10分鐘後就因為一點小事跟自己的太太吵架或者生氣；也常常背誦"生氣但不要動怒,不可含怒到日落"的金句,但轉頭看見孩子不聽話,就怒從心中起,恨不得給他們兩巴掌；每天我們把許多時間花在電腦上,網絡上,請客吃飯聊天上,卻擠不出二十分鐘來閱讀聖經.

這就是彼得,這就是我們!我們都是愛主的人,可是我們常常忙著自己的事情,卻常常把主的名,把主給我們的使命忘記了.某種意義上講,我們也常常不認主,而且我們不認主不需要一個小時,也許就是一轉念的事情.不信你試試看,如果你過年回家,同親戚朋友一起喝酒吃飯,有人批評基督教這也不好那也不好,你是埋頭不說話呢,還是站出來說"我就是個基督徒"呢?認不認主,有時候就在這一轉念之間!

但主對我們不離不棄,他沒有對彼得說:"你這個沒良心的傢夥,你不認我,你就下地獄吧!"他仍舊是愛彼得,路加在這裡強調了主耶穌的一個動作--主耶穌只是轉過身來,靜靜地看著彼得.彼得想起主說過的話,就出去痛哭!--這個痛哭的彼得,才是真正的英雄!

什麼是真正的英雄?第一,真正的英雄乃是能夠看清自己的人.彼得許多事情都要搶在前面,而且,耶穌誇他是磐石,說要建立教會在磐石上,彼得就心情大好了；所以他跟眾門徒一樣,也會爭論排名高低.彼得認為自己是個勇敢堅定的人,他完全忘記了自己是血肉之軀,也是一個會軟弱的人.所以當耶穌說他"雞叫之前,你會三次不認我,"彼得卻信誓旦旦說"眾人雖然為你的 緣故跌倒,我卻永不跌倒".(馬太26:33)他看不到自己的盲點,看不到自己的軟弱之處.我們何嘗不如此,我們忙這個忙那個,忙得沒有時間參加崇拜團契,我們安慰自己說不是我們不愛主,實在是很忙,其實真實的原因是因為我們缺乏安全感,我們對自己的身份地位錢財缺乏安全感,我們怕一旦不忙,這些的東西就得不到,就失去了；如果我們看不到真正的安全感在哪里,我們就很容易跌倒.我曾經看過一篇見證,是一位美國名校的量子力學博士寫的,他很年輕就從中國到美國波士頓求學,在見證中他說:"我本來一直覺得我愛我父母,也愛我的妻子,但其實我發現我最愛的是自己,根本不愛別人.以前我就像瞎了眼一樣,不能認識自己的本相.認識上帝后,才看清自己原來是這麼自私的一個人."當他看清楚了自己的罪,他才開始尋求上帝.

彼得也是如此.彼得之所以出去痛哭,是因為在主回頭看他的那一刻,他才真正看清了自己:那是一個軟弱的自己,一個膽怯的自己,一個自私的自己；那是一個喜歡誇口,喜歡爭名次,卻在信仰上還遠遠不夠成熟的自己!弟兄姊妹,我們有看清自己麼?

第二,真正的英雄,乃是一個願意悔改的人.悔改是一種行動.但我們往往是只有後悔,卻沒有行動--沉迷於打遊戲機的人知道沉迷不好,可是一到遊戲機邊上就控制不住自己的手；沉迷於煙酒、毒品的人也很後悔自己的行為,可是後悔歸後悔,卻沒有決心去戒掉；沉迷於工作的人也是如此,他們何嘗不知道過分的工作會傷害身體健康,會讓他們少了很多跟朋友家人的正常交流,可是只要有機會,任何地點都是他的辦公學習場所.

這個世上有許多人是悔而不改.知道做錯了,心裡很後悔,但就是不改.後悔藥吃多了會傷心肝脾肺腎,悔而不改,後悔多了,就會變成心裡的包袱.有時候,這個包袱甚至會要人的命.猶大就是這樣,他出賣了耶穌,他後悔了,但他寧願上吊,也不悔改.彼得則不同,雖然他三次不認主,雖然他錯了很多,但當他知道自己錯了的時候,他就會出去痛哭.他再見到主面的時候,他也沒有躲進辦公室.雖然萬般羞愧,他還是來到耶穌面前,真正聽從了耶穌的話,去牧養主的羊,直到為主獻出生命.

彼得肯悔改,神也沒有因為彼得犯錯而拒絕他.同樣,神也不會因為我們犯錯而拒絕我們.我經常問一些朋友,怎麼來了教會這麼久還沒有信主啊,他們一般都是回答我:你們基督徒要求太高了,做人太高尚,我還沒有達標,不敢做基督徒.其實不然,我們之所以是基督徒,不是因為我們個個都很聖潔,個個頭上都有個小光環.我們到這裡來,是因為我們知道我們都是罪人,但我們雖然是罪人,卻是蒙恩的罪人,因為我們肯認罪,我們肯悔改.弟兄姊妹,當我們來到主面前,認罪悔改的時候,主一定不會拒絕我們,一定會擦去我們的眼淚.我們是悔改的英雄,因為神是這樣愛我們!

第三,一個真正的英雄,是個可以聽得進意見的人.或者說,一個真正的英雄,是聽神的話聽得進去的人.彼得以前聽不進意見,耶穌說他的話,他總喜歡反駁,總喜歡表達一下自己的獨立見解,可是這一次,神的話語重重的擊中了他.如果說彼得的悔改重建了一個新彼得,那重建的根基就在於神的話語--彼得之所以會悔改,是因為他想起了主的預言,主預言他雞叫之前三次不認主.這一次,主的話語讓他不僅僅認清了自己,也真正知道了主是"至高神的兒子". 我們基督徒沒有什麼可以值得誇口的,也沒有什麼值得倚靠的.唯一能夠倚靠的,那就是神的話語,就是這本聖經上所見證的話語.只有當我們聽神的話聽得進去的時候,你才真正體會神的大能有多大.改變天地算得了什麼,他還能改變你的心,你的生命!

弟兄姊妹,真正的英雄,不是那些從來不曾失敗,從來不曾跌倒的人.真正的英雄,是那些失敗跌倒之後,認清自己,改正自己,重建自己的人.可怕的不是跌倒,可怕的是跌倒之後卻賴在那裡不爬起來.弟兄姊妹,讓我們按著神的話語重新建立自己,向著標杆前行,哪怕經常跌倒,但卻能站起來,一直走到神給我們的目的地.



\newpage

\section{使人和睦的人有福了}
\label{sec:41}
\textbf{鄧瑞強博士 - 使人和睦的人有福了}
\newline
\newline
\hyperref[sec:40]{< < < PREV SERMON < < <}
~
\hyperref[sec:toc]{[返目錄]}
~
\hyperref[sec:42]{> > > NEXT SERMON > > >}


馬太福音 5:9
\newline
\begin{longtable}{cl}
\hline
\hline
章節 & 經文 (和合本修訂版)\\
\hline
5:9 & \begin{tabularx}{0.7\textwidth}{X} 締造和平的人有福了！因為他們必稱為神的兒子。 \end{tabularx} \\ \\
[1ex]
\hline
\hline
\end{longtable}
日期:2010年4月4日

主題:使人和睦的人有福了

經文:馬太福音5:9

証道:鄧瑞強博士

各位弟兄姊妹,主內平安.今日講道的題目是:「使人和睦的人有福了」.主要的經文是馬太福音5:9,經文說:

「使人和睦的人有福了!因為他們必稱為神的兒子.」

小朋友時常打架,因為他們只是小朋友,其實,這通常只是「玩」,因為他們沒有「隔夜仇」,打完,很快便再玩過.奇怪的是大人們也時常打架,這通常「不是講玩」,並且,其中的仇隔幾代也隔不清.關於大人們的這種現象,最簡單的解釋是,大人們比小朋友們更未長大,更不成熟.

「和睦」這題目,對未長大的大人們,是一個永恆的主題.「和睦」這個字是聖經中最重要的一個字.舊約希伯來文的用字是 ``Shalom",新約希臘文的用字是 ``eirēne",即英文名 ``Irene" 這個字.中文聖經翻譯這兩個字,有時譯「和睦」,有時譯「和平」,有時譯「平安」,泛指人和人之間的和洽,國與國之間的和平,人和大地的和諧,人和神之間的和好.種種光景,是指一種生命的完滿,這是神創造萬物及人類的秩序的至終完成.中國哲學家牟宗三寫了一本書叫《圓善論》(「圓」是「圓滿」的「圓」,「善」是「美善」的「善」),講到中國的聖人,能憑一己之力,達到美善的圓滿境界.聖經講的「和睦」或「平安」,其實也是一種「圓善論」,一種「基督教的圓善論」,但這種生命的美善的完滿,不能憑一己之力而得.這種生命的完滿,需要神的愛去接觸我們,需要我們以愛去回應,需要我們以神的愛去愛他人,需要我們與神,與人,與大地建立正義的關係.這不是一個人關了房門,戴著耳筒,聽著美妙的音樂,享受著一種「平安祥和」的感受而獲得的；這需要我們走出大門,以堅毅的生命,在人間建立和平而得的.

生命的「和睦」,「和平」,立足於何處?

曾經,在蘇聯和美國冷戰的時代,「和平」立足於彼此的核武加給對方的威懾力,並彼此對對方的恐懼.這種「和平」建基於威嚇和恐懼.這就像武林上的兩大高手,拿著置對方於死地的劍,互相對峙,大家心中想,敵不動我不動.雙方在僵持.在不知底蘊的人看來,他們看似是在「和睦」的狀態中.事實上,他們是在最「不和睦」的光景裡.當然,很多看似「和諧」的家庭都在這「和諧」中.這些武林高手,在威嚇他人和被他人威嚇的恐懼中彼此對峙著.這能是「和睦」,「和平」嗎?事實上,恐懼下的和平,只是壓抑著的戰爭.

聖經中有一個故事,講到兩兄弟.弟弟欺騙了哥哥,引起哥哥的殺機.弟弟逃亡了,然後,經過一段時日,弟弟也強大了,要返回自己的老家.兄弟相殘,可以是極其殘酷的.兩人即將相遇,他們的結局將如何?舊恨會否帶來新仇?這是以掃和雅各兩兄弟的故事.兩人即將相遇.弟弟雅各派人去查探哥哥那邊的情況.我們看一看這故事的一些經文.

創世記32: 6-12:32:6　所打發的人回到雅各那裡,說:「我們到了你哥哥以掃那裡,他帶著四百人,正迎著你來.」32:7　雅各就甚懼怕,而且愁煩,便把那與他同在的人口和羊群,牛群,駱駝分做兩隊,32:8　說:「以掃若來擊殺這一隊,剩下的那一隊還可以逃避.」32:9　雅各說:「耶和華──我祖亞伯拉罕的神,我父親以撒的神啊,你曾對我說:『回你本地本族去,我要厚待你.』32:10　你向僕人所施的一切慈愛和誠實,我一點也不配得；我先前只拿著我的杖過這約旦河,如今我卻成了兩隊了.32:11　求你救我脫離我哥哥以掃的手；因為我怕他來殺我,連妻子帶兒女一同殺了.32:12　你曾說:『我必定厚待你,使你的後裔如同海邊的沙,多得不可勝數.』」

以掃帶著四百人前來,來意不善,凶多吉少.

聖經怎樣形容雅各?創32:7形容雅各「甚懼怕」,而且「愁煩」.在強大武力威懼下,雅各懼怕了,他害怕別人殺他和全家.他也「愁煩」,這個字指他入了「死胡同」,苦無出路.但是若人「趕狗入窮巷」,狗會反咬人的.有學者說,雅各的懼怕,是他怕被殺；但他被迫入窮巷,他便有理由動殺機.他在被殺的懼怕和有理由去殺人的內心衝突中.論力,他可能不及以掃；論智,則應綽綽有餘.

這可以是進入冷戰的關鍵時刻.他可以運用智力,用自己的武力震懾對方.但在這裡,雅各卻訴之神的慈愛和誠實.他呼求神以慈愛救他脫離凶惡,他也仰賴神誠實地落實神對他的應許.他在神裡面,以神的愛和神給予的希望,來面對這看來必死的局面.

之後,他往哥哥那邊送去大批的牛羊牲畜作和解的禮物,並向哥哥「認低威」.

最後,

創33:1-4說:33:1　雅各舉目觀看,見以掃來了,後頭跟著四百人,他就把孩子們分開交給利亞,拉結,和兩個使女,33:2　並且叫兩個使女和她們的孩子在前頭,利亞和她的孩子在後頭,拉結和約瑟在儘後頭.33:3　他自己在他們前頭過去,一連七次俯伏在地才就近他哥哥.33:4　以掃跑來迎接他,將他抱住,又摟著他的頸項,與他親嘴,兩個人就哭了.

本來,這可能是一場血鬥,但因著雅各倚靠神的慈愛,他以禮物和自甘卑下去示愛；也因他相信神的應許,他在窮巷中卻不反咬,他在「愁煩」中相信神給予的出路,而放棄了因恐懼而生的攻擊性,如此,化解了一段人間的冤仇.生命的「和睦」,「和平」,立足於何處?

立足於神的愛,立足於因這份愛而來的希望.面對人間的仇恨,新約有一段經文,表達了神如何成就人間的和睦.

以弗所書2:14-15:2:14　因他使我們和睦("eirēne")〔原文是:因他是我們的和睦〕,將兩下合而為一,拆毀了中間隔斷的牆；2:15　而且以自己的身體廢掉冤仇,就是那記在律法上的規條,為要將兩下藉著自己造成一個新人,如此便成就了和睦("eirēne").2:16　既在十字架上滅了冤仇,便藉這十字架使兩下歸為一體,與神和好了.

耶穌的生命,就是神向人顯現的生命.耶穌十字架的作為,就是神在人間顯明的作為.

弗2:14說:「耶穌就是我們的和睦,平安」.他的生命就是「圓善」(圓滿的善)的生命.他藉著為他人犧牲的愛,去化解人間的冤仇.沒有以暴力去威攝人,只有平等地和人講道理.沒有在恐懼中合理化自己的反擊,只有以正義去堅定自己的承擔.中國的維權人士劉曉波被判入罪,他的期望是:「希望自己是中國最後一個以言入罪者.」多麼輕省,多麼沉重.十字架是輕省的,十字架是沉重的.

耶穌是我們的和睦,平安.他「將兩下合而為一,拆毀了中間隔斷的牆」.所謂「兩下」,即人間敵對雙方.敵對雙方之間,當然有一將人分隔的牆.「柏林圍牆」是這種牆的典型.「種族隔離」是這種牆的另一典型.不公平的階級分野,是這種牆的另一典型.以暴力威懾人,叫人沉默,是這種牆的再另一典型.主以十字架的承擔,去拆毀這一「中間隔斷的牆」.

這些「圍牆」的拆毀,意味著更正義的社會的出現.使人「和睦」,或者更精確的說,使人間有真正的生命的圓滿,不單是追求個人的無災無痛,更是追求社會的正義得以落實.故此,在聖經裡,「和睦」或「和平」,常和「正義」連在一起.最出名的經文是

詩篇85:10:「公義和平安彼此相親.」英文NIV譯本的翻譯是: ``righteousness and peace kiss each other." 正義和平安彼此親嘴.兩者本是相連的,追求「和睦」,就是追求「正義」.

弗2:15指出,敵對雙方的冤仇,是人間種種「律法上的規條」造成的.所謂「律法上的規條」,用現代的講法,可以說,就是令人分隔開的種種「意識形態」.當人以「意識形態」認為自己是站在真理一方,別人是站在錯誤一方,而視打擊對方是正當時,人間的冤仇便無日無之了.基督當然視自己是站在真理一方,但這沒有使他視打擊對方為理所當然,反之,既然自己是站在真理一方,就當為這真理付上十字架的代價.讓不明白真理的人,因這十架的愛,而打破「意識形態」的魔障.

這樣,原先因仇恨分隔的群體,才能結成一新群體,成為新人類.弗2:15說,這就是真正的「和睦」("eirēne").

弗2:16說,這就是人和神的真正「和好」(reconciliation).

關於化解敵對雙方的冤仇,在歷史上,發生了這樣一件事.

在第一次大戰後,戰敗的德國政壇上,有一個重要人物,一個猶太裔的德國人,他叫拉特諾(Walther Rathenau),1922年成為德國的外交部長.他力圖重建德國的經濟.為了打破戰後德國外交上的孤立,他簽署了一些看來是向外國讓步的條約.這引致他最後被政治極右主義者及反猶主義者的暗殺.三個暗殺者,兩個自殺了,一個被捕,他的名字叫 Ernst Werner Techow .三日後,被殺者的母親,寫信給暗殺者的母親,信中說:

「在無可言喻的悲傷中,我向你伸出我的手,你是女人中最值得同情的.請告訴你的兒子,假若你的兒子在地上的法官面前,完全誠實地承認他的罪行,......並在天上的法官面前懺悔,我會寬恕他,我以他所殺害的人的名字及精神,寬恕他,正如神會寬恕一樣.」("In grief unspeakable, I give you my hand and \#8212;you of all women the most pitiable. Say to your son that, in the name and spirit of him he has murdered, I forgive, even as God may forgive, if before an earthly judge your son makes a full and frank confession of his guilt...and before a heavenly judge repents.")

這個暗殺者在1927年被釋放,其後,在第二次大戰期間,他協助超過700猶太人逃離納粹的魔掌.他這樣說:「正如Rathenau太太在寫那封寬恕的信時,征服了自己；我同樣嘗試克服自己.我只希望我有機會,去糾正我做過的錯事.」

殺人者與被殺者,德國極右主義者與被認為是出賣國家利益者,反猶主義者與猶太人,應是敵對雙方,但卻因一個承擔了十字架的女人的寬恕,打破了兩者中間隔斷的牆,成就了和睦,帶來了和平.

神以基督的生命,向所有人展示這種十架的愛的力量.在這裡,仇恨的罪被指明了,被定罪了；但也表明愛比恨更大,寬恕比定罪帶來更大的出路.這種愛,不是浪漫的愛,而是十架的愛.這種寬恕,不是無奈的接受,而是看到真正的人性前景.沒有這一切,人就會在仇恨的循環中互相煎熬下去了.

馬太福音5:9說:「使人和睦的人有福了!因為他們必稱為神的兒子.」

「神的兒女」是屬於神的,反映神的性情的,有神的尊榮的.這些人明白生命的真意義,也活出這種意義.但若要成為神的兒女,你必須「使人和睦」,「使人得平安」,「使人生命完滿」.你願意成為一個使人和睦的人嗎?為了寫一個關於「平安」的故事作結尾(「和睦」即「平安」),我在電腦的搜尋網站中問了一個問題:「如何能得平安?」回應是:買一份「平安」保險.

另一回應是:安裝一個「平安鐘」.另一回應是:要在平安夜有真平安,要做好防範措施.

最後一個回應是:平安是被裁的工人再次被聘用.平安是殘障的人可隨意無障礙地活動.平安是農民可保存他們的土地.如此這般,便可以使人得到和平,使人和睦.多麼簡單,多麼艱難.你願意成為一個使人和睦的人嗎?若願意,你是神寶貴的兒女.

「使人和睦的人有福了!因為他們必稱為神的兒子.」

但願榮耀歸於聖父,聖子,聖靈.阿們.



\newpage

\section{在橄欖山與聖殿山之間}
\label{sec:42}
\textbf{林豪恩先生 - 橄欖山和聖殿山}
\newline
\newline
\hyperref[sec:41]{< < < PREV SERMON < < <}
~
\hyperref[sec:toc]{[返目錄]}
~
\hyperref[sec:43]{> > > NEXT SERMON > > >}


路加福音 19:28-41
\newline
\begin{longtable}{cl}
\hline
\hline
章節 & 經文 (和合本修訂版)\\
\hline
19:28 & \begin{tabularx}{0.7\textwidth}{X} 耶穌說完了這些話，就走在前面，上耶路撒冷去。 \end{tabularx} \\ \\ \relax
19:29 & \begin{tabularx}{0.7\textwidth}{X} 快到伯法其和伯大尼，在名叫橄欖山的地方，他打發兩個門徒， \end{tabularx} \\ \\ \relax
19:30 & \begin{tabularx}{0.7\textwidth}{X} 說：「你們往對面村子裡去，進去的時候會看見一匹驢駒拴在那裡，是從來沒有人騎過的，把牠解開，牽來。 \end{tabularx} \\ \\ \relax
19:31 & \begin{tabularx}{0.7\textwidth}{X} 若有人問為甚麼解開牠，你們就這樣說：『主要用牠。』」 \end{tabularx} \\ \\ \relax
19:32 & \begin{tabularx}{0.7\textwidth}{X} 被打發的人去了，所遇見的正如耶穌對他們所說的。 \end{tabularx} \\ \\ \relax
19:33 & \begin{tabularx}{0.7\textwidth}{X} 他們解開驢駒的時候，主人問他們：「為甚麼解開驢駒？」 \end{tabularx} \\ \\ \relax
19:34 & \begin{tabularx}{0.7\textwidth}{X} 他們說：「主要用牠。」 \end{tabularx} \\ \\ \relax
19:35 & \begin{tabularx}{0.7\textwidth}{X} 他們把驢駒牽到耶穌那裡，把自己的衣服搭在上面，扶耶穌騎上。 \end{tabularx} \\ \\ \relax
19:36 & \begin{tabularx}{0.7\textwidth}{X} 他前進的時候，眾人把衣服鋪在路上。 \end{tabularx} \\ \\ \relax
19:37 & \begin{tabularx}{0.7\textwidth}{X} 他將近耶路撒冷，正下橄欖山的時候，一大群門徒因所見過的一切異能，都歡呼起來，大聲讚美神， \end{tabularx} \\ \\ \relax
19:38 & \begin{tabularx}{0.7\textwidth}{X} 說：「奉主名來的王是應當稱頌的！在天上有和平；在至高之處有榮光。」 \end{tabularx} \\ \\ \relax
19:39 & \begin{tabularx}{0.7\textwidth}{X} 人群中有幾個法利賽人對耶穌說：「老師，責備你的門徒吧！」 \end{tabularx} \\ \\ \relax
19:40 & \begin{tabularx}{0.7\textwidth}{X} 耶穌回答：「我告訴你們，若是這些人閉口不說，石頭也要呼叫起來。」 \end{tabularx} \\ \\ \relax
19:41 & \begin{tabularx}{0.7\textwidth}{X} 耶穌快到耶路撒冷，看見那城，就為它哀哭， \end{tabularx} \\ \\
[1ex]
\hline
\hline
\end{longtable}
日期:2010年4月11日

主題:橄欖山和聖殿山

經文:路加福音19章28-41節

証道:林豪恩先生

引言

橄欖山詩篇125篇有詩句如此說:

「眾山怎樣圍繞耶路撒冷,耶和華也照樣圍繞他的百姓」.耶路撒冷位於中央山脈最高處的平台上,四面被群山環繞,該詩篇詩意地描寫了耶路撒冷的地理位置.

圍繞在耶路撒冷東面的山稱為橄欖山,山上因有茂密的橄欖樹而得名.不過,橄欖山的歷史事跡乏善可陳.舊約第一次提及橄欖山,是大衛王因兒子押沙龍背叛而逃亡的時候.大衛逃離耶路撒冷,從東面登橄欖山,再往前走下裂谷去.

新約時代,橄欖山也不是甚麼繁華的地方,聖經提及的,只是兩個村落:伯大尼和伯法其.

聖殿山與橄欖山距離不遠,相對而望的聖殿山就不同凡響了,從古至今,它都在政治及宗教上扮演著舉足輕重的角色.聖殿山位於耶路撒冷城內的東部,相傳就是亞伯拉罕獻以撒為祭的地方.正如前文所說,耶路撒冷被群山環繞,外敵不易接近,但交通卻便利,是一個天然的統治中心.大衛王高瞻遠矚,攻取當時稱為耶布斯的耶路撒冷,建造大衛城,從此以後,就成為了王國的京城.

大衛的兒子所羅門接續王位,在王朝盛世的年代,他在聖殿山的山巔,展開建造聖殿的工程,歷時七年竣工.從大衛王朝開始,耶路撒冷成為以色列人的政治中心；所羅門王建成聖殿之後,聖殿山確立了其宗教中心的絕對地位.

公元前587年,巴比倫王攻陷耶路撒冷,所羅門王建造的聖殿被徹底摧毀.公元前536年,波斯王國取代巴比倫王國,被擄的猶太人開始歸回耶路撒冷,在百廢待興的艱困情況下,他們堅持重建聖殿.公元前516年,第二聖殿正式啟用.幾百年間,聖殿因戰爭而有所破壞.公元前20年,羅馬派來的統治者希律開始擴建聖殿,至公元後64年才完工.

馬可福音記載

:

耶穌從殿裡出來的時候,有一個門徒對他說:「夫子,請看,這是何等的石頭,何等的殿宇!」(13:1)門徒所讚嘆的,就是這個擴建中的聖殿.

今天福音經課所記載的耶穌進入耶路撒冷的事件,即是從橄欖山往聖殿山的過程.橄欖山與聖殿山對耶穌的態度,以及耶穌在這兩座山的遭遇,對我們有甚麼啟發呢?

經文

28耶穌說完了這話,就在前面走,上耶路撒冷去.29將近伯法其和伯大尼,在一座山名叫橄欖山那裡,就打發兩個門徒,說:30「你們往對面村子裡去,進去的時候,必看見一匹驢駒拴在那裡,是從來沒有人騎過的,可以解開牽來.31若有人問為甚麼解牠,你們就說:『主要用牠.』」32打發的人去了,所遇見的正如耶穌所說的.33他們解驢駒的時候,主人問他們說:「解驢駒做甚麼?」34他們說:「主要用牠.」35他們牽到耶穌那裡,把自己的衣服搭在上面,扶著耶穌騎上.36走的時候,眾人把衣服鋪在路上.37將近耶路撒冷,正下橄欖山的時候,眾門徒因所見過的一切異能,都歡樂起來,大聲讚美上帝,38說: 奉主名來的王是應當稱頌的! 在天上有和平； 在至高之處有榮光. 39眾人中有幾個法利賽人對耶穌說:「夫子,責備你的門徒吧!」40耶穌說:「我告訴你們,若是他們閉口不說,這些石頭必要呼叫起來.」 41耶穌快到耶路撒冷,看見城,就為它哀哭,42說:「巴不得你在這日子知道關係你平安的事；無奈這事現在是隱藏的,叫你的眼看不出來.43因為日子將到,你的仇敵必築起土壘,周圍環繞你,四面困住你,44並要掃滅你和你裡頭的兒女,連一塊石頭也不留在石頭上,因你不知道眷顧你的時候.」

釋經與應用聖殿山是中心,眾人以為神聖所在；橄欖山是邊緣,耶穌選為留宿之處 聖殿山是猶太人的宗教中心,是每一位猶太人都想去的地方.一般民眾都會好像路加福音第2章所記述耶穌的父母一樣:

「每年逾越節,他父母就上耶路撒冷去」.無論是過節,或是因事獻祭,居住於不同地區的猶太人都四方八面朝聖殿山而來.另一些奔向聖殿山的,是宗教從業人士,包括祭司,長老及文士等.經文所記述的事件發生在逾越節之前,耶穌與他的門徒很可能也在往聖殿過節的行列中.福音書記載,他們從東面而來,經過耶利哥,來到與聖殿山東鄰的橄欖山.

對於人流而言,橄欖山不是目的地,只是目的地的外圍；目的地是聖殿山,那裡才是招聚從四方八面而來的人流的中心.因此,若把這兩座彼鄰相望的山對比,毫無疑問聖殿山是中心,橄欖山只是相對於中心的邊緣.中心,最少有兩項特別的地方,第一,是人群所流向的地方,就如香港的中環一帶,每天早上人群從港九新界以及大小離島湧往那裡上班.中心除了有「吸入」的能力之外,第二項特別的對方,可能也是不為人所留意但卻影響更深遠的,就是它的「輸出」,這「輸出」,就是它的價值觀.在香港,有人稱中環的「輸出」為「中環價值」.當「中環價值」向外輸出的時候,它能夠渡過維多利亞港,主導西九龍文化區的發展,也可以跨越獅子山,影響新界鄉村的命運.「中環價值」有強大的伸展能力,它可以進入及重整社福界,也可以影響及指導教育界,這可能就是近年令社福界及教育界同工感到不舒服的其中一個原因.

聖殿山作為宗教的中心,的確凝聚著信徒的心靈,牽引著信徒的腳步.信徒對聖殿山這個宗教中心的迷戀及執著,引發了不少戰爭與血淚.新約時代的聖殿如耶穌所言,不久之後就被徹底毀滅.隨後,猶太人也被驅逐出故土,浪跡天涯.伊斯蘭信徒在這個當年的猶太聖殿遺址上,建造了宏偉的阿克薩清真寺和金頂清真寺.伊斯蘭信徒視聖殿山為聖地,猶太教徒更視聖殿山是他們的聖地.猶太教徒及伊斯蘭信徒對聖殿山的爭奪,轉化為宗教的紛爭和歷史的積怨.這個小山頭,承載著太多的仇殺和血淚.

現在,耶穌來到橄欖山,與聖殿山越來越近了.面對聖殿山,面對以敬拜包括祂在內的三位一體上帝之名而建造的聖殿,耶穌的態度是怎樣的呢?從耶穌的行跡來看,他對這個「宗教中心」並不是趨之若鶩,反而似乎刻意與之保持一點距離.馬可福音如此描述耶穌的行程:

「耶穌進了耶路撒冷,入了聖殿,周圍看了各樣物件.天色已晚,就和十二個門徒出城,往伯大尼去了.」(11:11)作者繼續刻意說明:

「每天晚上,耶穌出城去.」(11:19)在地理上,耶穌選擇棲身於橄欖山,而不留在聖殿山；在意識形態上,耶穌似乎也刻意與「聖殿價值」保持距離.他提醒他的門徒說:

「你們要防備法利賽人的酵」(12:1).

神聖在哪裡呢?當人們以為神聖必然在聖殿山的時候,神聖卻在橄欖山；當人們以為神聖必然住在聖殿的時候,神聖卻住在村莊民居；當人們以為神聖必然在宗教中心的時候,神聖卻在邊緣.如果純粹為接近神聖的緣故,何必一定要爭奪聖地呢?何必要爭得不是你死就是我亡呢?

橄欖山的信仰簡樸而熱情,聖殿山的信仰專業而權威路加福音記述,耶穌下橄欖山的時候,騎在一匹驢駒之上.耶穌能夠騎在一匹驢駒之上,反映著橄欖山簡樸及互信的氣氛.耶穌吩咐他的門徒說:

「你們往對面的村子裡去,進去的時候,必看見一匹驢駒拴在那裡,是從來沒有人騎過的,可以解開牽來.若有人問為甚麼解牠,你們就說:主要用牠.」驢駒的主人查問的時候,聽到門徒說

「主要用牠」,就借驢駒給耶穌使用了.(30-34)

耶穌騎著驢駒下山,路加福音描述:

「門徒把衣服搭在上面」(35),

「眾人把衣服鋪在路上」(36),

「眾門徒都歡樂起來,大聲讚美神」(37),並且高呼:

「奉主名來的王是應當稱頌的!在天上有和平；在至高之處有榮光.」(38)這熱鬧及喜氣揚揚的氣氛充滿著橄欖山,熱情及歡樂的人們懷著盼望伴隨著耶穌步向聖殿山.

在這興高采烈的時候,人們的熱情突然被「撥冷水」,有嚴肅而刺耳的聲音劃破愉快的氣氛,說:「老師,命令你的門徒安靜吧!」這是從何而來的掃興?這是何人發出的指責?把人們的熱情叫停的,正是人群中的幾個法利賽人.福音書把法利賽人與祭司,長老和文士視為相同性質的群體,他們在聖殿山學藝,承繼聖殿山的傳統,他們是聖殿山的代表.為甚麼是他們出聲指責歡迎耶穌的群眾呢?因為他們擁有分辯是非對錯的權威.為甚麼他們多次派人去審查耶穌的言論及行為呢?因為他們是宗教的專業人士,只有他們才有能力去判斷正邪.為甚麼他們多次派人去挑戰耶穌呢?因為他們是聖殿山的守護者,他們有權力去清除一切不容於聖殿山的事物.

當耶穌走下橄欖山向聖殿山而去的時候,橄欖山的群眾投以信任及幫助,熱情高興地隨行.然而,聖殿山的反應則完全是另外一個模樣,聖殿山已經準備好發揮它的功能,對耶穌作出最嚴謹及專業的審查.耶穌對於聖殿山的審查不以為然,他回應法利賽人說:

「我告訴你們,若是他們閉口不說,這些石頭必要呼叫起來.」(40)石頭都要呼叫,橄欖山對耶穌的歡呼是合情合理的,為甚麼聖殿山無動於衷呢?

橄欖山讓耶穌彰顯身份,聖殿山認不出它的主人為甚麼耶穌下橄欖山走向聖殿山的時候,突然間要騎在一隻驢駒之上呢?眾人都認為,不是因為耶穌疲倦而要以驢駒代步.那麼,耶穌騎著驢騎從橄欖山進入聖城,這景象有甚麼啟示呢?舊約時代的先知撒迦利亞曾經如此預告:

「鍚安的民哪,應當大大喜樂；耶路撒冷的民哪,應當歡呼.看哪,你的王來到你這裡!他是公義的,並且施行拯救,謙謙和和地騎著驢,就是騎著驢的駒子.」(撒迦利亞書9:9)又有預言說:

「那日,他的腳必站在耶路撒冷前面朝東的橄欖山上.」(15:4)耶穌從橄欖山騎著驢駒進入聖城的景象,呼應著先知對彌賽亞來臨的預言.

眾門徒大聲呼喊的內容,出自詩篇118篇的詩句:

「奉耶和華的名來的是應當稱頌的!」(118:26)118篇是被虜前的以色列王每年的重新嘉勉禮儀所使用的詩篇,當王帝進入嘉勉場所的時候,人們就頌唱起這首詩歌來.所以,當耶穌騎著驢駒,從橄欖山步向耶路撒冷,眾人又唱起詩篇118篇的名句,這個情景的含意呼之欲出了.由此可見,對耶穌而言,橄欖山是一怎樣的地方呢?橄欖山是一個讓耶穌彰顯彌賽亞身份的地方,是一個承載著王者來臨的地方.

相對而言,對這位奉主名來的王來說,聖殿山又是一個怎樣的地方呢?路加福音記述耶穌望著聖殿山,悲嘆地說:

「但願你今日知道那有關你和平的事,可是你現在那樣的盲目．．．因為你沒有認出上帝拯救的時機」(41-44現中)聖殿為敬拜上帝而建,是宗教中心；聖殿山是聖山,是代表著上帝同在的地方,非常諷刺的是,聖殿山竟然認不出上帝的來臨(did not recognize the time of your visitation from God).這真是一個非常值得探討的問題,也是一個十分值得信徒反省的現象.

耶穌視橄欖山為家,在此留情；耶穌更視聖殿山為家,為此赴義路加福音清楚指出:

「耶穌每日在殿裡教訓人,每夜出城在一座山,名叫橄欖山住宿.」(21:37)耶穌似乎視橄欖山為家.家,是一個可以休息的地方,一個可以讓人放鬆的地方,橄欖山,就有著這些舒服而隨和的氣氛.家,更重要的,是有愛你的人,以及你所愛的人.橄欖山上,有一位接待耶穌的馬大.路加福音10章如此記載:

「他們走路的時候,耶穌進了一個村莊.有一個女人,名叫馬大,接他到自己的家裡.她有一個妹子,名叫馬利亞,在耶穌腳前坐著聽他的道.」(10:38-39)橄欖山上,有一位馬利亞,約翰福音指出,

「這馬利亞就是那用香膏抹主,又用頭髮擦他腳的.」(約11:2)橄欖山上,也有一位耶穌所愛的拉撒路.有一次,拉撒路病重,當耶穌趕到的時候,拉撒路已經去世了.耶穌看見眾人哀哭,聖經說:

「耶穌也哭了」,以致眾人說:

「你看他愛這人是何等懇切.」(約11:36)橄欖山是耶穌的家,耶穌在此留情.

耶穌對橄欖山有情,耶穌對聖殿山其實也是有情有義.路加福音19章41節如此描述:

「耶穌快到耶路撒冷,看見城,就為它哀哭」.耶穌為橄欖山的拉撒路而哭,耶穌也為聖殿山的遭遇而哭.愛得深,痛得也深,而責之也切,耶穌對聖殿山的感情不比對橄欖山輕.

橄欖山歡迎耶穌,耶穌喜愛橄欖山,耶穌視橄欖山為家是合理的.聖殿山呢?聖殿山更應該是耶穌的家啊!還記得耶穌十二歲時曾經離隊獨留自己在聖殿的事件嗎?當耶穌的父母找到他時,還記得他對父母說甚麼嗎?

「為甚麼找我呢?豈不知我應當在我父的家裡嗎?」(2:49)可惜,現在聖殿山認不出耶穌,聖殿山不歡迎耶穌,聖殿山甚至要除去耶穌.話說回到,天大地大,耶穌並不是無處容身,為甚麼仍要去一個不歡迎自己的地方呢?

有些決定,不能純粹基於個人的喜好,也不能只考慮個人的感受.有些時候,因為責任,不能因為不被歡迎而不去；有些時候,因為使命,不能因為感覺不好而不去；有些時候,因為其他人的需要,不能因為個人安危而不去.耶穌回去聖殿山,暴露出這個宗教中心的異化,拆除架設在創造主與被造者之間的障礙,讓人可以與上帝回復關係.在人世間,也常常有「你愛他,但他不愛你」的狀況.有人的回應是「既然你也不愛我,我也不愛你」,有人的回應是「縱使你不愛我,我也不放棄愛你」,對家庭如是,對國家如是.

1988年6月,有一位青年剛在北京師範大學取得博士學位,同年到美國進修.1989年春季,由高校學生發起的民主運動在北京展開.這位青年馬上回國,投身在學運的熊熊烈火之中.當時,劉賓雁如此評論這位青年的行動:「他這次回去,明明是飛蛾撲火!」果然他在6月13日被捕.接著的二十多年,當局希望他離開國家,自動消失.而去,他卻戀戀不休,為了使命,為了理想,為了人民,寧願坐牢,也要留在國內而不願出走.

總結總括而言,橄欖山與聖殿山的故事對我們有甚麼啟發呢?

讓我們問一間自己,基督徒的聖殿山在哪裡呢?基督徒的聖殿在哪裡呢?若說,基督徒沒有聖殿,因此也沒有聖殿山,這是對的.保羅卻從另一個角度指出,基督徒的聖殿不是一座建造物,而是基督徒這個人.保羅在哥林多前書6章19節說:

「豈不知你們的身子就是聖靈的殿嗎?這聖靈是從上帝而來,住在你們裡頭的」.

從這個角度來看,橄欖山和聖殿山的對比對我們就有意義了.

雖然我們是聖靈的殿,但基督徒不應該像聖殿山一樣視自己為中心,為主角,甚至以作為上帝的代表自居；反而應該像橄欖山一樣,作為一個隨時接待上帝的地方.

基督徒可以學習聖殿山一樣的嚴緊而專業,但不是用來審查上帝,而是用來檢討自己,檢視教會,反省神學；對上帝的態度,應該學習像橄欖山一樣,簡樸而熱情,信任而投入.

基督徒應該像橄欖山一樣,作為上帝彰顯自己的媒介；而提防如聖殿山一樣的盲目,把上帝的探訪及眷顧拒之千里之外.

最後,讓我們緊記:當我們能夠討上帝歡心的時候,上帝對我們有情；但當我們靈頑不靈的時候,上帝對我們有義.或達或窮,上帝對我們的情義都是那麼濃厚.



\newpage

\section{為義受逼迫的人有福了}
\label{sec:43}
\textbf{伍渭文牧師 - 為義受逼迫的人有福了}
\newline
\newline
\hyperref[sec:42]{< < < PREV SERMON < < <}
~
\hyperref[sec:toc]{[返目錄]}
~
\hyperref[sec:44]{> > > NEXT SERMON > > >}


題目:為義受逼迫的人有福了

日期:2010年4月18日

香港中文大學(前任)校長沈祖堯教授曾說過,希望中大學生能出一位諾貝爾和平獎得主.為甚麼是和平獎,不是其他科研的得主呢?我想這與沈教授一路來強調大學不單追求科研的卓越,更要培育學生的品格和使命感有關.諾貝爾和平獎的得主,不能單單埋首實驗室,圖書館,他要進入人群,作個使人和睦的人.

今天是第四屆崇基基督教文化節最後一天,我們講道的題目『為義受逼迫的人有福了』,這題目和『使人和睦的人有福了』有何關係呢?是否使人和睦就會被人逼迫呢?我翻看過去二十年諾貝爾和平獎的得主,很多人沒有受過逼迫,當然被提名的人,不少是被打壓或被拘禁的異見份子,提名作為和平獎的候選人己成為對他們的支持了.不少和平獎的得主,其實都身世顯赫,握有權力,他們運用自己的影響力,去降低及試圖消弭地球上的危機和衝突.因為他們的努力,世界變得更安全,族群變得更融和,他們是使人和睦的人.

對伊斯蘭世界表示友善的奧巴馬,促成以色列和埃及大衛營協議(Camp David Accords)前美國總統卡特(Jimmy Carter),聯合國和前秘書長安南(Annan),國際原子能委員會(International Atomic Energy Agency),國際政府氣候變化委員會(Intergovernmental Panel on Climate Change)和前美國副總統戈爾(Al Gore),無國界醫生(Medecins Sans Frontiers),窮人銀行(Grameen Bank) 家Muhammud Yunus；他們一開始工作就受到嘉許,沒有人逼迫他們.無國界醫生的前線志願者,都得到敵對相方的保護.

但九一,九二,九三年三位得主亞洲的昂山素姬,中美洲危地馬拉的Rigoberta Menchu,和南非的孟特拉就受到極大的迫害.Menchu女士的父親,母親和兄長都被軍政府拘禁,施酷刑和殺害.昂山素姬至今仍被軟禁,孟特拉坐了二十七年牢.為何有這麼不同的待遇呢?因為前一班人沒有觸及公義的問題,而昂山素姬(Aung San Suu Kyi),Menchu和孟特拉(Mandela)都是民權鬥士,他們為了公義而抗爭,所以受到逼迫.

1每個基督徒都是和平的使者

基督教文化節的主題『使人和睦的人有福了』取自登山寶訓,是登山寶訓的綱領,列舉進入上帝國的生命素質,所以首尾兩福都有『因為天國是他們的』,這是生命盛宴的邀請,因為種子己撒在我們的心田,會結出在上主面前的敬畏,自省己罪的哀慟,與人交往的溫柔謙遜,進而饑渴慕義,得著屬靈的力量,開出對社會的關懷,憐恤有需要的人,使人和睦,締造和平,而在過程中保持良心的清潔,甚至為此受到逼迫,義無反顧.八福層層相扣,而其中使人和睦一福極為重要,因為和睦在原文指和平shalom,和平貫徹大地就是shalom,上帝的旨意行在地上如同行在天上就是shalom,基督就是我們的和平.

政治現實主義者(political realists) 所說的和平是威脅的平衡(balance of terror)下,彼此相安無事,因懼怕敵人核武的報復而有的和平,而八福所說的是復和,轉仇恨為友愛.我們說和平是建立在公義之上,而公義是愛的均分,因為愛每一個人,我們保護他們應有的權利,這就是公義.沒有公義就沒有真正的和平,長久的和平,而沒有愛和憐恤,就不是八福中所說的公義.以弗所書第二章十四,十五節:因他使我們和睦(他是我們的和睦) ……以自己的身體廢掉冤仇……為要將兩下藉著自己造成一個新人,如此便成就了和睦.和睦是從愛的犧牲和救贖而來的,是十字架的愛,帶來和睦.因此經歷上帝的愛的人,被召進入神國的人,最適合作使人和睦的人.

憐恤人因為看到人的苦難,但使人和睦使我們正視苦難的原因.梁燕城博士在上主日給我們看到大陸很多城市令人神傷的圖像.每個晚上華燈初上的時候,在城市邊緣的貧民區,很多少女打扮得性感亮麗,走過馬路的大廈的夜總會,卡拉OK上班.這些從農民工子女因為沒有戶籍,不能享受城市義務中學教育,她們的父兄參與建築華廈,但他們在城市付不起昂貴的租金,被迫住在貧民區,或沒有窗戶的地下室.他們沒有分享到經濟發展的成果,但承受著經濟發展帶來的高物價和租金. 針對這不公義的戶籍制度,前一段時間十三間報紙發表了共同社論:

戶籍制度分割了城市和鄉村.”農民工”是對那些戶籍在農村而身在城市打工的人群的特定稱謂,最早的一代農民工,為城市的發展付出了自己的勞動,可是,他們 的下一代仍然沒有辦法解決身份認同,他們的子女仍然背負著上一代的困惑,他們生活的城市仍然無法接納他們.即便在城市中,戶籍制度也分割了城市的居民.在同一座城市中,儘管我們與其他人一樣為這座城市的建設奮鬥多年,我們與其他人一樣納稅,但沒有戶口讓我們無 法與其他人一樣享受平等的就業機會,享受同等的醫療,教育,養老等等社會保障.因此,夫妻被迫兩地分居,年老的父母無法與子女團聚,孩子無法獲得良好教 育.我們要問,這樣的隔離究竟還要持續幾代人?

八福是每個信徒生命的素質,也是每個信徒應有的責任,我們每個人都有責任去成為和平之子,參與社會,維繫公義.但甚麼是公義呢?為義受逼迫又是甚麼意思呢?是否每個基督徒都會受到逼迫呢?

2 為基督的名受逼迫因為敢於與眾不同

為義受逼迫一詞常常被人引用,合理化自己的抗爭行為.說自己為義受逼迫,就可否不聽反對我們的人的意見呢?我們遭遇的難處,可能因為自己沒有智慧惹來的,亦可能由於冒進狂熱,引起人的反感和排斥.我以前認識一位姊妹,上班時常常看聖經引起老闆不滿,多次轉工,這不是為義受逼迫.有一次耶穌在拿撒勒的會堂講道完畢,猶太人怒氣滿胸,把他趕出城外,並把他帶到山崖,想推他下山,『他卻從他們中間直行,過去了』(路4:29, 30) 耶穌不會招惹不必要的苦難,在成長過程中聖經形容神和人都喜歡他:『耶穌的智慧和身量.並神和人喜愛他的心,都一齊增長.』(2:52) 保羅又教導我們,選立教會的長老和執事,候選人要在教外有好名聲,與人和睦,得眾人的喜愛敬重.

聖經說:『人若因我辱罵你們.逼迫你們,捏造各樣壞話毀謗你們,你們就有福了!』(太5:11) 我們受到逼迫因為我們敢於和世人不同.世人說自我感覺良好,我們說在上帝面前虛心謙卑,對自己的罪敏感哀慟；我們說溫柔的人有福了,但世人說溫柔是弱者,強者才可出人頭地；我們說要飢渴慕義才得飽足,但世人要飽嘗一切山珍海錯,才不枉此生；世人看政治的現實,祇有威脅的平衡,才帶來和平,我們說以愛去化解仇恨,贏取敵人,勝過贏取戰爭.

張曉風老師在『從閱讀到寫作』的公開講座,分享她與眾不同的價值觀,作為她寫作的原則.「我曾經寫一些跟愛情有關的文章,大概三百字,寫了兩篇.有一個人覺得很好,說一定會大賣,要幫我出書.我很感謝他,他給我很好的條件.但過了一天,我拒絕了他,因我不再想寫這一類受歡迎的愛情題材,這些都有商品的意味,但我想做的是更有意義的事情.我們要有原則,既然掌握了文字的力量,便要小心用,要想 想對別人有甚麼意義.」因為這與眾不同,可能引起譏笑,也可能使走大眾路線的作家們不開心.

在今天若我們堅持一夫一妻,至死不渝,可能被當作不能與時俱進了,這些話祇有牧師在主持婚禮時才說吧.教導祇有結婚後才可以有性關係,已被視為落伍的思想.結婚真的要一男一女嗎?為何不能有同性婚姻呢?若說同性戀是可以矯正的,也有人矯正了,更加被認為是政治不正確.沒有絕對的,一切都是相對的；其實說這話的人的潛台詞是,沒有絕對的,除了這話. 我們受到排斥,攻擊,甚至逼迫,因為我們跟隨與眾不同的基督.

3 公義要與憐恤調和

因基督的名受到逼迫提醒我們,追求公義的過程中,我們不能和其他七福的生命素質割裂,我們曉得公義的基礎乃是憐恤,公平分配憐憫就是公義,我們的敵人也應得到我們的憐恤.所以對掌權者陳述我們的意見為民請命時,可以多一些尊重和體諒,溫柔而堅定,不需使用文字暴力.

梁燕城博士在復和與救贖的公開講座以個人的見證,分享他與中國的復和.以前他就批評,嘲諷共產政權,但九三年以海外學人身份開始進入大陸參與辦義學扶貧,十七年來協助了近五千農民子弟完成中學,而且不少考進大學.同時不斷和相關官員陳情進言,倡議改革,其中不少意見漸漸被採納.他說,要救贖中國,要用點點滴滴的愛,聚涓成流.

當孟特拉成為總統後,他成立由大主教圖(Tutu) 主持的『真相與和解委員會』 (Truth and Reconciliation Commission),讓受種族隔的受害者傾訴苦情,又讓施行罪行的人公開尋求饒恕,使對立的群族,得以和好.這委員會讓公義昭彰,同時又有憐恤,對受害者和施行罪行的人同樣有憐恤.憐恤與公義並舉,公義是憐恤的均分,沒有憐恤,尋求公義可能成為仇恨和報復,就如一七八九年的法國大革命；革命黨人為了推行自由,平等,博愛的新理想,徹底剷除舊社會,把大量貴族送上斷頭台.一位革命黨人羅蘭夫人(Madame Roland) 因為黨爭也被送上斷頭台,在行刑前她直視斷頭台大聲說:自由,自由,多少的罪惡假汝之名而行!(O liberty, liberty, how many crimes are committed in thy name!)

虛心的人有福了……為義受逼迫的人有福了,因為天國或神的的國是他們的.阿們.



\newpage

\section{讓復活與我們走過每一天}
\label{sec:44}
\textbf{Dr. Kuruvilla Abraham - 讓復活與我們走過每一天}
\newline
\newline
\hyperref[sec:43]{< < < PREV SERMON < < <}
~
\hyperref[sec:toc]{[返目錄]}
~
\hyperref[sec:45]{> > > NEXT SERMON > > >}


日期:2010年4月25日

題目:讓復活與我們走過每一天

“主復活了”.”祂確實復活了.”早期教會的會員每禮拜天聚集做禮拜時,彼此間都會這樣打招呼.對他們來說,復活節的紀念不是一年只有一次.他們在每個禮拜天都是慶祝復活的主與他們同在.為了這樣一個每週一次的紀念,他們特地把崇拜之日從禮拜六改為禮拜日.其實對他們來說,主的復活是每天都在發生的現實.

對許多人而言,復活是我們死後發生的事,或是一年一度人們在教堂里慶祝的節日.它很少影響我們的日常生活.且看門徒們的經驗:在主耶穌被釘死以後他們都害怕,失望.有的復歸他們原先的職業,而因為跟他們的同行們已有相當一段時間脫離聯系,所以就無法前進.他們失去對自己的希望,並且對他們曾跟隨的師傅的希望.

不過主的復活就完全改變了他們的人生.可怕的命運已經讓整個世界上下顛倒.他們並沒有等在死後的世界才體驗復活的現實,而是在宣揚當下的生命中復活的主與我們同在.主的復活改變了他們眼前的生活.

那麼主的復活跟我們日常生活有甚麼關係呢?

第一,一套嶄新的由關係和義務承諾而組成的世界會進入我們的生活.上主復活,最先向一群曾跟隨祂到底的婦女們顯現:這絕不是偶然的事情.她們無保留的信心得到了證實.這幾位在福音書中被提及的婦女,都能清楚地回憶起上主與他們在一起時的親切和慈悲.復活的主就顯現于我們的友愛關系當中,復活的主就透過人們彼此間相互負責的關係向我們顯現.

我們今日的生活完全為科技牽制,我們很少意識到自己也是信仰世界的一部份.但主的復活提醒我們,我們的人生必須從信仰的價值觀中吸取營養.

讓我舉個日常生活中的例子.我妻子去看病時,她必須受臨床身體檢查.醫生想要用客觀測驗查明她血液中的礦物質,鹽,和葡萄糖的比例.他也查她的心跳.

許多科學或生物學的原理支配著我們的身體,而我們的健康有賴於它們的相互協調.不過當我跟妻子相處的時候,那些原理以及相關的臨床結論都不重要.

我們之間的關系基於互相的愛和信任.人類的許多活動都需要科技原理的指導,但是若想把人生的一切都還原於科技,就會產生很大的歪曲.為了活得更有意義,我們需要的是愛,同情,互相尊敬,互相關懷等從信仰而生的價值觀.

今日我們面臨很多全球性危機:金融風暴,生態危機,家庭和社群的瓦解.後者可以稱之為人際關系的危機.

且聽一位科學家的言辭:

“我從前以為世界上最嚴重的環境問題就是全球變暖,環境惡化以及生態系統的崩潰.可是我錯了.真正的問題並非此三項,而是貪心,自私和冷漠,而為了解決這些問題需要精神和文化的轉變.我們科學家並不知道怎樣達成這種轉變.我們需要你們.”

他演講的對象是一群教會的領袖.堅信並宣揚主的復活是請人們重視這個人際關系的世界.我們所有的人際關系都能因神的大能而轉化.當我們把他人當作要操縱的客體而對待,以自己的利益或舒適為重時,我們就會毀壞人性,毀壞神的創作.

復活的主帶來的轉化,會讓我們看重人際關系以及做出的承諾,而這就是復興,更生之所由來.透過對他人的承諾,復活就成為日常生活中的現實.信仰的社群,教會,虔誠地致力于培養這種價值觀,使人們願意承擔責任去照看他人和自然

第二,復活的主創造一個愛的社群.門徒們害怕時就失去了他們人生的目標.可是復活的主造就新的團結,創造了一個充滿責任感的社群,一個尤其關注弱者和邊緣化人的社群.哪裡有愛的社群,哪裡就有復活的主.我想把這個作為我個人的證詞.

一九七七年的時候,我在班加羅爾一個信徒頗多的教會服侍.當時我妻子生了一個女嬰.出生後不久這女嬰生病,雖然有很好的醫療護理,但腦炎對她的大腦仍造成了很嚴重的損害.這對我們來說是莫大的精神創傷.我們問神,為甚麼讓這個事情發生?我們感激我們的女兒能夠活下去,但是照顧一個大腦受了損傷的女孩兒,並不是很容易的事.我們努力去理解這件事情背後的意義.

這個經歷讓我們有機會接近許多擔負同樣責任的父母.像任何處於這種情況的人一樣,我們向神禱告,祈求醫治.不過神賜給我們的,並非奇跡性的治愈,而是更大的奇跡.神帶領我們,為大腦受過損傷的孩子提供療養之所,名叫麗薩的家.現在我們大約照看有二十個兒童.有一部分和我們的女兒一樣不能說話,有的需要幫助才能夠走動.他們來自不同的宗教背景.不過他們一起生活,而且生活在一個愛的社群裡.

當人們來訪問我們療養所的時候,我總會告訴他們,我們社群的成員們沒有甚麼特別才能,而且不善於講任何人類的語言,但是他們卻精通一種語言:神的語言,就是愛.神不講希伯來語,不講希臘語,不講漢語,不講印度話.神只知道一種語言,而我們療養所的孩子們好像都擅長這種語言.在他們的懷抱中,他們的親吻中,他們的舞蹈中,他們活出神的樣子.今天是復活節,所以他們肯定早已聚集在一起贊美神.我們和他們一起擘餅,分享這餅的時候總會說,”在破碎了的人們當中,破碎了的神”.對我們來說這就是復活的意義.一個破碎了的人的社群在神的愛當中結合為一.復活若非此義,請問有何義哉?



\newpage

\section{願你們平安}
\label{sec:45}
\textbf{陳曉東博士 - 願你們平安}
\newline
\newline
\hyperref[sec:44]{< < < PREV SERMON < < <}
~
\hyperref[sec:toc]{[返目錄]}
~
\hyperref[sec:46]{> > > NEXT SERMON > > >}


約翰福音 20:19-31
\newline
\begin{longtable}{cl}
\hline
\hline
章節 & 經文 (和合本修訂版)\\
\hline
20:19 & \begin{tabularx}{0.7\textwidth}{X} 那日（就是七日的第一日）晚上，門徒因怕猶太人，所在的地方門都關了。耶穌來，站在當中，對他們說：「願你們平安！」 \end{tabularx} \\ \\ \relax
20:20 & \begin{tabularx}{0.7\textwidth}{X} 說了這話，他把手和肋旁給他們看。門徒一看見主就喜樂了。 \end{tabularx} \\ \\ \relax
20:21 & \begin{tabularx}{0.7\textwidth}{X} 於是耶穌又對他們說：「願你們平安！父怎樣差遣了我，我也照樣差遣你們。」 \end{tabularx} \\ \\ \relax
20:22 & \begin{tabularx}{0.7\textwidth}{X} 說了這話，他向他們吹一口氣，說：「領受聖靈吧！ \end{tabularx} \\ \\ \relax
20:23 & \begin{tabularx}{0.7\textwidth}{X} 你們赦免誰的罪，誰的罪就得赦免；你們不赦免誰的罪，誰的罪就不得赦免。」 \end{tabularx} \\ \\ \relax
20:24 & \begin{tabularx}{0.7\textwidth}{X} 那十二使徒中，有個叫低土馬的多馬，耶穌來的時候，他沒有和他們在一起。 \end{tabularx} \\ \\ \relax
20:25 & \begin{tabularx}{0.7\textwidth}{X} 其他的門徒就對他說：「我們已經看見主了。」多馬卻對他們說：「除非我看見他手上的釘痕，用我的指頭探入那釘痕，用我的手探入他的肋旁，我絕不信。」 \end{tabularx} \\ \\ \relax
20:26 & \begin{tabularx}{0.7\textwidth}{X} 過了八日，門徒又在屋裡，多馬也和他們在一起。門都關了，耶穌來，站在當中，說：「願你們平安！」 \end{tabularx} \\ \\ \relax
20:27 & \begin{tabularx}{0.7\textwidth}{X} 然後他對多馬說：「把你的指頭伸到這裡來，看看我的手；把你的手伸過來，探入我的肋旁。不要疑惑，總要信！」 \end{tabularx} \\ \\ \relax
20:28 & \begin{tabularx}{0.7\textwidth}{X} 多馬回答，對他說：「我的主！我的神！」 \end{tabularx} \\ \\ \relax
20:29 & \begin{tabularx}{0.7\textwidth}{X} 耶穌對他說：「你因為看見了我才信嗎？那沒有看見卻信的有福了。」 \end{tabularx} \\ \\ \relax
20:30 & \begin{tabularx}{0.7\textwidth}{X} 耶穌在他門徒面前另外行了許多神蹟，沒有記錄在這書上。 \end{tabularx} \\ \\ \relax
20:31 & \begin{tabularx}{0.7\textwidth}{X} 但記載這些事是要使你們信耶穌是基督，是神的兒子，並且使你們信他，好因著他的名得生命。 \end{tabularx} \\ \\
[1ex]
\hline
\hline
\end{longtable}
日期:2010年5月2日

題目:願你們平安

經文:約翰福音20:19-31

証道:陳曉東博士

弟兄姊妹在主裡面向大家問安.

在我今天準備這個講道的時候,我再一次重新複習了一下,一位美國的神學家尼布爾(Karl Paul Reinhold Niebuhr),他曾經寫過一篇非常出名的禱告,是一個平安的禱告(Serenity Prayer),他的禱告是這樣說:上帝啊,求您賜我寧靜的心,接受我所不能改變的事,賜我勇氣,改變我所能改變；賜我智慧,分辨兩者的差別.神啊,求你教導我,珍惜每一天,享受其中的每一刻,以苦難作為通往平安之路,如耶穌看這世界如其所是,而非如我所願.深信衪使萬物妥當,只要我降服於上主的旨意.而我在此世將有合宜的幸福,他日與祂共享無宇倫比的幸福,阿們.God, give us grace to accept with serenitythe things that cannot be changed,Courage to change the thingswhich should be changed,and the Wisdom to distinguishthe one from the other.Living one day at a time,Enjoying one moment at a time,Accepting hardship as a pathway to peace,Taking, as Jesus did,This sinful world as it is,Not as I would have it,Trusting that You will make all things right,If I surrender to Your will,So that I may be reasonably happy in this life,And supremely happy with You forever in the next.Amen.我們今天所看的約翰福音,有許許多多不同的解釋方式,很多的弟兄姊妹或者是聖經學者都詫異這卷書卷,和許多希臘的一些的古典哲學的概念之間有相同之處,所以在很長的時間以來,約翰福音被解讀為是深受希臘哲學思想影響的一卷聖經的書卷.也有人把它稱之為一卷是屬靈的福音書,好像是說這一卷書卷是在一種非常獨特的信仰經驗裡面被啟示出來,它背後沒有任何的歷史或者是社會的脈絡,但是我們今天更多的聖經學者,聖經的研究都告訴我們這卷福音書在它背後其實有非常深刻的社會和政治方面的意義,我想每一位熱愛聖經的弟兄姊妹當您打開聖經的時候,一定也會常常有一些很基本,很簡單的問題,從我們腦袋裡面崩出來,這一卷書,到底這個作者的身份是誰啊?到底聖經裡面一而再,再而三所說「主所愛的門徒」,這是一個什麼的稱呼呢?到底約翰福音再加上約翰書信再加上啟示錄有多少個約翰呢?到底主所愛的這一位門徒和約翰書信裡稱呼的那位長老,他們之間是怎麼樣的關係?每一個讀者當我們打開聖經的時候,我們都會很自然的去問寫作的年代,寫作的場景,寫作的對象,資料的來源,社會的環境.我想許多的聖經學者他們所以被稱為聖經學者也是跟我們一樣,就是因為熱愛這本的書卷,所以很希望能夠看到這卷書背後,神到底是透過怎麼樣的方式把祂自己的心意啟示給初期的教會,也啟示給我們今天的世界.

1940年代死海古卷的發現,為我們這卷約翰福音提供了更多的一些的可以參照的資料,很多的學者他們更願意今天用一種文學的方式,文學的技巧,形式批判(form criticism) 等等的方式,來試圖去理解這個約翰福音,理解它背後整個的社會脈絡是甚麼樣子,在今天愈來愈多的研究顯示,在這個文本後面,它不是一個孤絕的人在拔摩島上面把這樣的一卷書來寫下來而已,而是說在整個的約翰這個稱謂後面,有一大群的一個約翰的群體在裡面.在初期的教會當中見證這位又真又活的主跟這個世界之間的關係,其實不是僅僅是講個人很神秘的,很宗教的一種體會,約翰福音告訴我們一件事情,耶穌的生平,和他相信耶穌的這樣一個生命群體,他所面對的這個生活世界之間有一個非常緊密的關係.

在今天我們所看到的約翰福音第20章19-31節,有三個對我來說非常有價值有意義的三個要點,在這樣的一段經文裡面,我們看到耶穌重複的,重複的來去向他的門徒們來去問候,發出一個願你們平安的問候.也許有問候的這樣的一個字眼還有一點的誤解,也許更好的一個方法是說,耶穌在這樣的一段經文裡面用了三次的宣告,向信仰耶穌的這一個生命群體發出三次的宣告,每一次的宣告都有它一個特定的脈絡把這樣的一個宣告來宣講出來.第一次的宣告是在一個危機和恐懼中來帶出一個平安的宣告；第二次的宣告是在一個宣教使命的承擔當中帶出這樣的一個平安的宣告；第三次的宣告是在一個疑惑和一種宣認,一種的confession或者是認罪這樣的一個脈絡當中帶出一個平安的宣告.這是因為耶穌他作為彌賽亞這樣一個很獨特的宣信,對初期教會當它接受了這樣的一個獨特的宣認的時候,與他們所身處的世界,與他們所身處的猶太群體產生了一個非常大的一種的不同的意見,而在他與外邦和好這樣的信息當中,是面對著喪失他自己一個主流的身份和傳統的一個危機.

我們每一次讀到約翰福音的時候,每一次讀到新約的書卷的時候,也許我們在座今天的弟兄姊妹能夠再一次的受到提醒,就是提醒我們:教會在它自己創立之初曾經也是一個受到主流的排擠和受到權貴壓迫的一個生命群體,這樣的記憶也許是要提醒我們不要忘記自己的草根性,不要忘記自己的生命裡面有這些相通的聖徒們,也就是曾經飽受逼迫的一群的信徒,這樣的意義對我們今天不再受到逼迫的信仰群體來說是一個非常重要的學習,也許讓我們自己在安逸當中的時候,我們不要忘記有許許多多的人還在窮苦當中,還在生命掙扎的邊緣,還在被主流的意見所完全忽視的世界當中.而聖經,福音提醒著我們,我們是跟這一群人是在一起的.

當我們看到今天所講的這三點的時候,第一點說到:在恐懼中的一個平安的宣告的時候,我們看到在這個十九節,第二十章的十九節及二十節,很有意思的來表達,在這樣的一個環境裡面,書卷告訴我們這是一個黃昏,告訴我們這一群的信徒在房子裡面把門緊緊的關著.有一個字深深的打動我的心,就是說這個經文裡面他們說他們很害怕,害怕猶太人fear of the Jews,在這樣一個環境裡面,耶穌來了,耶穌站在他們當中向他們宣告:願你們平安或者是平安與你們同在,然後祂把祂的手,祂的肋旁指給他們看.也許黃昏給我們一個很重要的提醒,因為他是一個非常重要的象徵,告訴我們這是一個白晝將盡的一個環境裡面,也許這樣的一個門戶的緊閉也是非常的重要,他告訴這樣的一個群體是被外界所排斥的,他自己也是沒辦法把自己跟別人溝通的,而隨之而起的就是這樣一種的恐懼的心,這樣一個心靈緊緊的閉鎖的光境,在這樣的一個驚恐裡面,我們看到主在生命裡面的顯現,和平安的問候.手和肋旁的傷口,這是第一次,這是在復活之後第一次向他的門徒來去指示的,傷口告訴我們什麼?傷口告訴我們十字架,它所帶來這樣一種的創傷,如果我們自己想像一下,我們可以看到,恐懼的確是我們人性裡面很大的敵人.

約翰福音裡面的恐懼是一樣被迫害,被邊緣化的恐懼,但是我想恐懼這樣一種的心靈的狀態或者是恐懼這樣一種存在的狀態,並不僅僅是被迫害的時候被邊緣化的時候才存在的.在我們今天,我們所生存的世界當中,我們自己的社會裡面也是充滿了恐懼的氣氛,有一位哲學家他說:我們所生處的世代是一個技術化的時代Technological Age,人活在一個競爭和失敗的一個恐懼當中,所以我們常常在新聞報導裡面,或者是一些學術的刊物裡面,我們也是看到很多這樣的恐懼情緒的散播,城市和城市之間的競爭,今年香港在這樣的城市競爭裡面排名第幾位?背後是一個失敗的恐懼；學校和學校之間的競爭,當整個的適齡兒童大幅度下降的時候,學校怎麼樣能保存他們的競爭力,怎麼樣保持他們的存活率,怎麼樣保持被免於殺校這樣的一個的情況... 每一個的校長,每一位的老師,每一位的家長都活在這樣的一種恐懼當中.宗教和宗教之間也在競爭,當有人在說為什麼有這麼多的這個復活節的假期,拿一天在做這個孔誕不是很好嗎?我們這個信徒之間心裡面的恐懼就會出現,宗教和宗教之間開始彼此競爭.人和人之間不斷的在競爭,我們怕失敗,我們怕被邊緣化,我們怕被剝奪既得的權益,我們怕自己所做的一切努力付諸東流.

聖經在這裡宣告說,耶穌在門徒中間出現,把平安帶給了我們,和合本在這裡翻譯說:願你們平安,原文裡面沒有這個「願」字,這不是一個意願,這不是說耶穌「願意」我們平安,這是耶穌基督向我們的「宣告」,因為這是一件成就了的世界,當耶穌在我們當中向我們來宣告平安與我們同在的時候,這是一個與創世記當中宣告安息的日子,以這樣的一個高度的宣告,因為他事實具有上主他自己施工的成就,這樣的一個高度和廣度在當中,讓我們能夠明白整個的世界包括了死亡,失敗,恐懼,一切在我們心裡面勝不過的軟弱... 都包含在上主的肯定和平安當中.

弟兄姊妹這是一種何等的宣告.當我們願意自己去凝神默想注視這樣一個奧秘的宣告的時侯,我們會與保羅一齊的來去凱旋:死啊!你的毒勾在那裡,死啊!你的權勢在那裡,有甚麼樣的可怕,有甚麼樣的恐懼能夠勝過死亡對我們生命的威脅.但是在今天耶穌宣告了這樣的一個宇宙性的事實:平安與我們同在.

初期的教會領受了這樣的一個奧秘,領受了這樣的一個信心的道理.所以耶穌第二次的來問候他們的時候,平安與我們眾人同在的時侯,他下面的一句話就是一個差遣的話語.他說父怎樣的差遣了我,我也怎樣去差遣了你們.所以我們看到第二次的問侯,第二次的宣告是把耶穌基督的生命和教會的群體,教會的整個的宣教的工作連繫在一起,我們看到主耶穌當祂這樣的一個宣告的時候,祂不僅僅是宣告完了之後就離開了,而是祂宣告完了之後,祂給我們力量,給我們能力,所以祂向他們吹祂的生命氣息,祂向他們來去發出一個確認:你們來接受,你們來領受聖靈.

耶穌的氣息,耶穌的話語,耶穌讓我們在整個宣教的生命裡面是帶出這樣的一個平安的信息.這樣的,這樣的一個命令,任務,把我們和這個世界當中一個赦免的權柄帶在了一起.祂提醒我們一件事情,福音是一個平安的福音,福音是一個的醫治的福音,福音是一個包紮傷口的福音,福音是一個安慰的福音,福音不是帶著一種的恐嚇或者是操縱的信息,也許這是我們每一位的基督徒,每一個的宣教者常常面對的試探和引誘.我們怎麼樣用我們自己的技巧,用我們自己的方式或是用我們自己的口舌,能夠把人勸到教會裡來.

但是今天我們從耶穌基督的生命裡面看到的是,耶穌要我們把平安和生命帶進這一個世界當中,耶穌的平安與教會受差遣的召命是互相關係的.特別是在我們今天所生活的世界裡面,當我們生活在二十一世紀的時候,我們看到了一種唯物主義,一種資本主義和一個緊緊接合的二十一世紀的時候,我們更是要在這樣的一個世界當中,大聲的呼喊耶穌基督平安的福音,當這個世界如果相信一切的事物是唯物的,世界所有的事情所有的物資都是我們的資源,是供我們盡情去消費去享用的自然資源的時候,是我們可以隨便去的棄置的東西的時候,耶穌的顯現,耶穌的平安是告訴這個世界,這是這個世界是上主的傑作,它是上主去彰顯祂神聖的載體,它不可以被剝削,不可以被胡亂破壞的.教會有耶穌的氣息,有耶穌的話語,教會作為一個生命的群體,是要為造物界帶來平安和祝福的.

也許我們會有一點點的困難去看到,為什麼這樣一個平安的福音是可以帶著一個赦免的權柄,它讓我們在回顧到在教會的歷史裡面,我們曾犯下的許許多多的過犯和錯誤；它讓我們看到在歐洲的歷史裡面,這樣一個宗教裁判所的可怕；它讓我們看到在這個在教會的歷史裡面,人的這種私慾被這個站在一個領袖的位置上面對世界所帶來的傷害.為什麼主的話語裡面依然付給教會裡面有這樣一個的赦免的權柄.我相信是因為主信任我們,信任我們就算在罪過當中,在過犯當中,我們依然可以在主面前回轉,我們依然是帶著前人所犯下的錯誤來到主的面前求赦免,求幫助,求聖靈的充滿,求平安的同在.乃至於我們可以繼續成為上主的同工,在這個世界裡面,把和好,平安與世界分享.

所以當我們在平安中差遣的時候,我們用一位神學家艾魯所說的話,他說當這樣的一個平安帶出來的時候,就是要解除當代我們所身處的這個價值的神聖性,他告訴我們耶穌所差遣的平安,向信徒顯現的上帝對安息,對平安,對整個的世界,對整個人類的一個平安的一種肯定,我們所做一切事情,如果沒有為這個世界,沒有為人帶來平安的話,我們所做的事情是枉然的,是錯誤的.

第三方面,我們從今天的信息裡面,從今天的經文裡面看見了一件事實,就是平安勝過了懷疑.因為在二十四節裡面告訴我們,有一位在教會裡面非常聰明的人物,多馬Thomas-使徒多馬,他的聰明不只是他曾經在傳說裡面去到印度宣教,把福音帶到了印度,他的聰明是他能夠在信仰當中發出一個很大膽的懷疑聲音,所以我們在聖經裡面看到一件事情就是說,當耶穌顯現的時候,多馬並不在他們當中,而別的門徒向他說,耶穌向他們顯現的情形,耶穌所傳講的信息的時候,告訴他我們看到了主了.他說我除非我親眼看見了,除非我的感官,五官看見了,我的手摸過了(這個多馬可以說是教會歷史上面最大的經驗主義者),除非我看見了我不會相信.所以當耶穌在他生命當中顯現的時候向他說:願平安與你們同在.然後祂向多馬發出了一個邀請,祂說:請把你的指頭探進我手上的釘痕,請把你的指頭探進我肋旁的傷口,不要懷疑,總要信.

我們可能在這個話語裡面,我們的腦袋接不上來,為什麼多馬看見了這樣的事實之候,他整個人發出一個宣告:我的主我的神.耶穌這樣的一個傷口,他的懷疑,他所需要的經驗和他所需要感官的經驗,有什麼樣的關連在當中?我們接不上來!也許我們花一點點的時間來思想一下,這樣一個懷疑對我們的生命的價值和意義在那裡?也許把這樣的一個懷疑的精神,推到現代文化的高度,我們可以從笛卡兒的哲學思想裡面找到一種的共鳴感和找到現代多馬所要問的問題,我們可能看到一件事情,當我們去懷疑的時侯,我們是把我們整個人的思考的基礎放在「我」-- 放在「我」這樣的一個概念裡面,把「我」放在一個絕對判斷秩位置上面,「我」是一個思考的主體,有了「我」我才能夠思考,有了「我」我才能夠確定我的存在.所以笛卡兒曾經發出了一個非常出名的一個格言  and \#8212; 我思故我在(拉丁語:Cogito, ergo sum)笛卡兒的「我思」開啟了我們整個現代人的這種的懷疑時代.但是我們這種的懷疑不是一種毀滅性的懷疑,這是一種建設性的懷疑,在這樣的一個懷疑的過程裡面,人類的技術是突飛猛進,所以現代的思想家告訴我們,技術不僅僅是我們所想像的applied science ,技術和科學是這一種雙生子的關係.

我們有了這樣的一個懷疑的精神,我們有了這樣的一個大膽的把上帝擺在一個括弧裡面的一個這樣子的勇氣,我們可以在自己的生命當中去確定自己的位置的時候,整個人類,近代的文明裡面是突飛猛進,我們看見我們能夠在我自己的生命裡面,在我的雙手裡面,在我的理性裡面,我可以發展一系列技術的能力,所以我們人類可以征服自然,我們不怕和自然去發生巨大的衝突,我們可以想愚公一樣可以把這個山都移掉,把海都填平,把人類的DNA可以分解出來,把人類的這個最深奧的這個宇宙大爆炸,這個最初的億萬分之一秒的瞬間用現代的技術能夠把它掌握出來,我們看到人在確立自己的族體性之後,這一種的信心,爆炸,能力,人的慾望,人對權力,人對控制這樣的一個渴望,人的計算,一切一切... 到了一個極點  and \#8212; 我是王,我可以判斷一切,我可以做一切的事情.

所以我曾經跟一位的教授在聊天的時候,他告訴我,他說:在我們今天的大學裡面,我們不是用學科的價值來去規範一個大學,我們是用英文的字母來去規範一個大學,所以大學裡面最高級的學科叫做administrationn,最後的一個學科可能做theology,因為我們現在的人,我們是用技術性的頭腦來去思考問題,我們是用字母來去規範我們的世界.但是耶穌在今天面對這多馬,宣告的平安告訴我們,世界比我們的技術能力要大,要廣,要寬.我們當然可以與多馬一起來懷疑,我們當然可以看得到多馬的懷疑的價值,但是,在今天的經文入面我們看到了一件事情:耶穌邀請多馬去摸主的傷口.在摸祂的傷口裡面來去認識祂來,宣告祂是主祂是我們的神,這是怎麼樣的意思?

也許在結束的時候,我用一位我非常尊敬的作家  and \#8212; 李家同教授,他是台灣科學界非常受人尊重的教授.他提到他自己的一個經驗,他曾經放下自己的教職去到印度加爾各答在仁愛傳道會那裡,跟這個德雷莎修女一起去服侍了一段時間,當然他知道自己的身份依然是在一個學院裡面工作的人,但是他有一個文章裡面有一段的話,對我去理解耶穌邀請多馬去摸祂肋旁的傷口有一個非常大的啟發.他的故事是這樣:我(李家同教授)離開加爾各答的那天早上,旅館門口又多了一個老人乞丐,他衣不蔽體的躺在地下已經病得奄奄一息了,有一位日本的義工,單腿跪下,握住了這位老人的手等待計程車來,他為什麼單腿跪下呢?無非是想表示對這老乞丐的一點尊敬罷!這位日本義工大概二十歲左右的大學生,大個子,一定超過了一百八十公分的高度,他擔白承認他愛玩,我也兩次在街上看到他喝啤酒,可是他說,當他看到了乞丐馬上就變成了很嚴肅的人,他說他屬於那種叛逆型的人,很不願意表示對別人的尊敬,但是當他看到垂死的窮人的時候,他會很自然的跪下.

弟兄姊妹我想在這樣的一個記錄裡面沒有那一個的研究比這樣一段的記述,這樣的一個文章是更好的來解釋:為什麼耶穌邀請懷疑的多馬,要用感官的世界去認識上主的多馬來去摸祂肋旁的傷口,來去摸祂手上的傷口.

也許在這樣一個的復活節的日子裡面,當我們願意在這個世界裡面繼續地去摸耶穌的傷口,繼續摸祂肋旁的傷口的時候,我們的生命能發出多馬同樣的聲音:我的神,我的主!

也許我們可以發出如保羅一樣的呼喚,主啊!當我摸這個世界這些毀爛的傷口的時候,我渴望與你同行,我渴望... 雖然我帶走了我的懷疑,但是我渴望能夠在這樣的一個傷口的碰觸當中,知道你是那一位釘十字架的上帝,你是那位用生命來去醫治世界創傷的上帝,就是因為這樣,我們在今天經文才說:你因看見了我才信,那沒有看見就信的是有福了.

願意我們每一個人沒有看見就信了,而成為有福的人.



\newpage

\section{轉眼的怒氣一生的恩典}
\label{sec:46}
\textbf{伍渭文牧師 - 轉眼的怒 氣,一生的恩典}
\newline
\newline
\hyperref[sec:45]{< < < PREV SERMON < < <}
~
\hyperref[sec:toc]{[返目錄]}
~
\hyperref[sec:47]{> > > NEXT SERMON > > >}


詩篇 30:1-12
\newline
\begin{longtable}{cl}
\hline
\hline
章節 & 經文 (和合本修訂版)\\
\hline
30:1 & \begin{tabularx}{0.7\textwidth}{X} 耶和華啊，我要尊崇你，因為你救了我，不讓仇敵向我誇耀。 \end{tabularx} \\ \\ \relax
30:2 & \begin{tabularx}{0.7\textwidth}{X} 耶和華－我的神啊，我呼求你，你醫治了我。 \end{tabularx} \\ \\ \relax
30:3 & \begin{tabularx}{0.7\textwidth}{X} 耶和華啊，你救我的性命脫離陰間，使我存活，不至於下入地府。 \end{tabularx} \\ \\ \relax
30:4 & \begin{tabularx}{0.7\textwidth}{X} 耶和華的聖民哪，你們要歌頌他，要頌揚他神聖的名字。 \end{tabularx} \\ \\ \relax
30:5 & \begin{tabularx}{0.7\textwidth}{X} 因為，他的怒氣不過是轉眼之間，他的恩典乃是一生之久；一宿雖然有哭泣，早晨便必歡呼。 \end{tabularx} \\ \\ \relax
30:6 & \begin{tabularx}{0.7\textwidth}{X} 至於我，我凡事順利，就說：「我永不動搖。」 \end{tabularx} \\ \\ \relax
30:7 & \begin{tabularx}{0.7\textwidth}{X} 耶和華啊，你曾施恩，使我穩固如山；你轉臉不顧，我就驚惶。 \end{tabularx} \\ \\ \relax
30:8 & \begin{tabularx}{0.7\textwidth}{X} 耶和華啊，我曾求告你；我向耶和華懇求： \end{tabularx} \\ \\ \relax
30:9 & \begin{tabularx}{0.7\textwidth}{X} 「我被害流血，下到地府，有何益處呢？塵土豈能稱謝你、傳揚你的信實嗎？ \end{tabularx} \\ \\ \relax
30:10 & \begin{tabularx}{0.7\textwidth}{X} 耶和華啊，求你應允我，憐憫我！耶和華啊，求你幫助我！」 \end{tabularx} \\ \\ \relax
30:11 & \begin{tabularx}{0.7\textwidth}{X} 你將我的哀哭變為跳舞，脫去我的麻衣，為我披上喜樂， \end{tabularx} \\ \\ \relax
30:12 & \begin{tabularx}{0.7\textwidth}{X} 使我的靈歌頌你，不致緘默。耶和華－我的神啊，我要稱謝你，直到永遠！ \end{tabularx} \\ \\
[1ex]
\hline
\hline
\end{longtable}
日期:2010年5月9日

題目:轉眼的怒 氣,一生的恩典

經文:詩篇30:1-12

証道:伍渭文牧師

為何在復活節期間安排詩三十篇作為經課,這詩篇與基督的復活有何關係呢?詩篇 三十篇是一篇典型的感謝詩篇,詩人身罹重病,命懸一線,最後獲上主醫治重獲生命而頌詩感謝他.第一節是題目的宣告:

耶和華阿,我要尊崇你,因你曾提拔我,不叫仇敵向我誇耀.第二,三節進一步擴充宣告:

你醫治了我,你曾把我的靈 魂從陰間救上來,使我存活,不至於下坑.第四至十二節再詳細說明整件事,八,九節再說下到坑中.

Life redeemed蒙拯救生命 詩篇三十篇2-3

對前一段時間山西那班被困深坑的礦工,這段經文來得非常真實,他們都是從深坑中被救起來的,一百多部救護車在場侍候, 每部車都有醫生護士,往醫院的五十分公里路,都戒嚴佈防,通行無阻,醫院都預備好一切了.從深坑升上來的人,他們的靈魂也從陰間奪回來了.為了對比深坑的 下墮,詩人用高高揚起讚美來尊崇你(exalt)  .復活節就是基督降到陰間,拯救我們脫離深坑,然後被父神升起,高過諸天,成為萬王之王,凡信靠他為救主的,被稱為義,與他一同坐在父神的右邊.詩篇三十 篇是復活節的詩篇,因為這詩篇論及蒙拯救的生命(life redeemed) .同時這詩篇也是靈魂得著甦醒的詩篇.

Life awakened覺醒的生命 詩篇三十篇6-7

「至於我,我凡事平順,便說:我永不動搖耶和華阿.你曾施恩,叫我的江山穩固,你掩 了面我就驚惶.」(詩篇三十篇6-7)

凡事平順,永不動搖:詩人以為平順亨通(prosperity) 就可永不動搖:「

至於我,我凡事平順,便說:我永不動搖」(

詩篇三十篇

6),他以此為 傲,直到上帝藉著重病使他看到自己的限制和朽壞(mortality)  ,他才發覺以前所有的都是恩典,在重病中得到屬靈的覺醒,看到自己的有限,生命極其脆弱,他不過是血肉之軀.我記起我神學院一位老師說:他經過極大的手術 後,極其口渴,就是想喝水,你給他一顆鑽石都沒有用,不能解渴,祇想喝一口水,他願意以寶石來交換一杯水來解渴.一顆寶石的價值,不及一杯水,祇有曾經游 走生死邊緣的人才能體會到.

應作而沒有作的罪:

「耶和華阿,你曾施恩,叫我的江山穩固,你掩了面我 就驚惶.」(

詩篇三十篇7

) 這裡沒有說明作者犯了何罪,不像大衛明顯的犯了殺人丈夫,奪人妻子的大罪:

因為我知道我的過犯,我的罪常在 我面前,我向你犯罪,惟獨得罪了你,在你眼前行了這惡.(詩51:3-4) 這是明知故犯的罪行(sins of commission)  ,但詩篇三十篇的作者,以乎是犯了應作而沒有作的罪(sins of omission)  .在舊約時期,有兩種祭祀,一是贖罪祭,針對犯了具體的罪來求赦免,另一種是燔祭,祭牲全部燒掉,成為香氣上騰歸上帝,燔祭是尊崇上帝的祭,全歸上帝的 祭.詩人在繁華昌盛時沒有感恩,以為甚麼都是靠自己,沒有想到上帝,到上帝把手放開,他便傾覆了.他命懸一線時看到了自己的朽壞(mortality)謙 卑下來,是上帝的恩典他才興盛,他的罪就是自以為是:我凡事平順,我永不動搖.他沒有當上主為上主來敬拜他.他沒有愛上帝,他祇愛自己的財富平順.

無知的財主:這事使我想起路加福音12:16-19.那位無知的財主.他說:

「我的出產沒有地方收藏,怎麼辦呢?......要把我的倉房拆了,另蓋更大的,在那裡好收藏我的糧食和財物,然後對我的靈魂 說:靈魂哪,你有許多財物積存可作多年的費用,祇管安安逸逸地吃喝快樂吧!但神卻對他說:無知的人哪,今晚我要你的靈魂,你所預備的要歸誰呢?」

從北方漂來的微塵:詩人對財富的誇耀,以為有了經濟的成就,就可以永不動搖,也使我想起最近從北方漂來的微塵.當我們 誇耀擁有世界最長的高鐵,貫穿沿海和中部城市,惦記著半天的大都會生活圈,我們西南的土地,卻乾涸龜裂,寸草不生,己經窮困的農民,無地可耕.當上海博覽 會籌備得如火如荼,當我們正想向世人顯示我們的世博是歷年最大最好的,突然從北方刮起前所未有的沙塵暴,把上海網羅住,使我們看不清前面要走的路.沙塵顯 出我們國家的mortality,限制,和脆弱.因為人希望凡事平順,永不動搖,肆意如取如攜,使草原沙漠化,地下水大量流失,才造成如此乾旱.

熱帶雨林消失:據傳媒報導,為了更大的經濟效益,很多雲南的熱帶雨林被砍伐後改種造紙的高大桉樹,但桉樹對當地鄉土 的,原產,原生的物種有極大的抑制性.它生長了,其他物種就不能生長,而且會慢慢地退縮,最後造成桉樹林都是地表光禿禿的,不能儲存地下水.但原始雨林有 多層的林木,下面有灌木,中有喬木,有二三十米的,面有六七十米的摩天樹,還有籐科植物,地下水蒸發到一米,就被一米的灌木吸收了,上升到五米,就被五米 高的樹冠壓住了,整個生態是濕潤的.原始森林的根系大量積蓄地下水,打水井很淺,不用深挖就有水冒出來.由於大量種植桉樹,熱帶雨林愈來愈少,土地沙漠化,地下水嚴重流失.

一天從晚上開始:江山穩固,永不動搖,不是因為人的平順(prosperity)  ,努力發展,乃是因為上帝施恩.人祇看到自己的努力和生產力,忘記上帝的恩典,可以從對一天的理解看出來.人計算一天是白天 and \#8212;可以工作的白天,沒有計 算晚上,因晚上是睡覺的,不能工作.一天過後,便到晚上,好像一天並不包括不能工作的晚上.但聖經計算一天是從晚上開始:

起初神創造天地,地是空虛混沌,淵面黑暗,神的靈運行在水面,神說:要有光!就有了光,神看光是好的,就把光暗分開了. 神稱光為晝,稱暗為夜,有晚上,有早晨,這是頭一天 (創 1-3) .上帝在黑暗中工作,造出光和晝,這是二十四小時的一天.當然現代社會的晚上也是不停工作,不夜天已成為發達城市的記號,但對人的身體來說,晚上十一時就 睡覺最合適人的生物時鐘.

有晚上,有早晨,這是一天的開始,一天是從接受上帝恩典開始,不是從麥當勞的早餐開始.太陽下山後,氣溫下降,空氣中 的水氣漸漸形成露水,滋潤大地.當我們睡覺不能工作的時候,大自然乃然在工作,上帝在晚上讓大地生生不息,也讓我們的身體好好竭息,上帝在我們休息時修補 我們的細胞,修復我們的器官,使我們經過休息,白天可以有力量繼續工作.上帝在我們休息的時候默默的做了大量工作,雖然我們不察覺,也不明白,這是恩典, 上帝用他的權能托住萬有,我們才永不動搖.是上帝的恩典,我們才永不動搖.

詩人經歷重病,來到死亡的門檻,他看清楚是上帝的恩典,使他永不動搖,不是凡事平順使他永不動搖,他有了屬靈的覺醒. 詩篇三十篇不單說到得拯救的生命,覺醒的生命,還提及以上帝為樂的生命.

以上帝為樂的生命(Joyful Life) 詩篇三十篇9-12

「我被害流血,下到坑中,有甚麼益處呢?塵土豈能稱讚你,傳說你的誠實 嗎?」(

詩篇三十篇

9)  「你己將我的哀哭變為跳舞,將我的麻衣脫去,給我披上喜樂,好叫我的靈歌頌你,並不住聲.耶和華我的神阿,我要稱謝你到永遠.」(

詩篇三十篇

11-12)

我們留意到詩人以下到坑內來形容死亡,而聖經說到撒旦和跟隨撒旦的人將會被扔在無底坑,這是甚麼意思呢?下到坑內是遠 離地上,更是遠離高天 and \#8212;上帝的寶座,死亡就是遠離生命的主；與上帝緊密相連,連死亡也不能分開,這就是永生.

「我被害(新譯本沒有被害)流血,下到坑中,有甚麼益處呢?塵土豈能稱讚你,傳說你的誠實嗎?」(

詩篇三十篇

9)New International version繙作

What gain is there if I am silenced in my going down  into the pit? Can the dust praise you? Will it proclaim your  faithfulness? 希西家的祈禱也說:

原來陰間不能稱謝你,死亡不能頌揚你,下坑 的人不能盼望你的誠實,祇有活人,活人必稱謝你,像我稱謝你一樣,為父的,必使兒女知道你的誠實.(賽38: 18-19) 而瀕臨死亡的人,沒有讚美稱頌,祇有像燕子呢喃,白鶴鳴叫,鴿子哀鳴.生命之可貴,因生命能顯揚上帝的榮美,生命的意義就在乎透過我們,上帝得著頌揚.

他的恩典乃是一生之久:這數節經文也告訴我們:生命是一首詩歌,是頌歌,喜樂的詩歌,不是哀歌.生命的基調是禮讚稱頌,不是悲傷絕望,在人生中也許有死亡般的經歷,但最後是凱歌,一宿雖然有哭泣,早晨便必歡呼,他的怒氣不過是轉眼之間,他的恩典乃是一生之久.這種人生 和仇敵的人生剛相反,仇敵的人生是嘲諷(詩篇三十篇1  誇耀gloat over me) ,幸災樂禍,犬儒的,懷疑的,不信的.

有人指出,現代人其中一種病態,就是沉溺在沒有上帝的不幸和苦痛中.當我們說:上帝已死,甚麼都可能了,根本沒有甚麼 對與錯.同時,上帝真的死了,甚麼都不可能了,我們不可能有甚麼作為.兩次大戰後,悲觀的存在主義盛行,很多人沉溺在沒有上帝的不幸中和苦痛中,像毒蠍子 把自己有毒的拑來刺痛自己,自得其樂.很多現代的小說,像羅得的妻子,回顧被火焚燒的所多瑪,俄摩拉兩個罪惡的城市.她回眸扭曲的人性,屬靈的空虛,畸形 的社會,還有那污穢的床,雖然她臭到上天上降下來的硫磺,但羅得的妻子,未能夠從失敗和懼怕中回過頭來,結果她成了鹽柱.

但詩篇三十篇告訴我們,上帝的怒氣是轉眼之間,他的恩典是一生之久,一宿雖然有哭泣,早晨便必歡呼.去年復活節,我在 講道中提及潘霍華的見証,劉曉波先生讀到他在納粹集中營所寫的書柬,沒有苦毒,充滿喜樂,令他感到驚訝,也得到啟發和鼓舞.他說:「學會敬畏,學會謙卑, 學會在危險和恐懼中坦然面對厄運,並對未來保持樂觀的信心,學會以發自內心的坦誠與人相處,這就是潘霍華聖徒人格的啟示.......潘霍華始終平靜地對自己微 笑,彷彿他是一個從一出生就只會笑的怪物.......他之所以能夠保持絕境中的希望,發出開朗而令人驚奇的笑,就在於他始終信仰著.他的文字不只是用來表達悲哀 罷了,更是表達歡樂的,其份量沉得比任何絕望都豐富.正如他在臨刑前向獄友告別時所說:「這,就是終點.對我來說,是生命的開端.」」

在今天主復活的日子,讓我們從新出發,更新心志,不再回頭,因為他已把我們從坑中拯救出來,他的怒氣是轉眼之間,他的 恩典是一生之久.阿們.



\newpage

\section{恩惠慈愛隨著我}
\label{sec:47}
\textbf{鄧瑞強博士 - 「恩惠慈愛隨著我」}
\newline
\newline
\hyperref[sec:46]{< < < PREV SERMON < < <}
~
\hyperref[sec:toc]{[返目錄]}
~
\hyperref[sec:48]{> > > NEXT SERMON > > >}


詩篇 23:1-6
\newline
\begin{longtable}{cl}
\hline
\hline
章節 & 經文 (和合本修訂版)\\
\hline
23:1 & \begin{tabularx}{0.7\textwidth}{X} 耶和華是我的牧者，我必不致缺乏。 \end{tabularx} \\ \\ \relax
23:2 & \begin{tabularx}{0.7\textwidth}{X} 他使我躺臥在青草地上，領我在可安歇的水邊。 \end{tabularx} \\ \\ \relax
23:3 & \begin{tabularx}{0.7\textwidth}{X} 他使我的靈魂甦醒，為自己的名引導我走義路。 \end{tabularx} \\ \\ \relax
23:4 & \begin{tabularx}{0.7\textwidth}{X} 我雖然行過死蔭的幽谷，也不怕遭害，因為你與我同在；你的杖、你的竿，都安慰我。 \end{tabularx} \\ \\ \relax
23:5 & \begin{tabularx}{0.7\textwidth}{X} 在我敵人面前，你為我擺設筵席；你用油膏了我的頭，使我的福杯滿溢。 \end{tabularx} \\ \\ \relax
23:6 & \begin{tabularx}{0.7\textwidth}{X} 我一生一世必有恩惠慈愛隨著我；我且要住在耶和華的殿中，直到永遠。 \end{tabularx} \\ \\
[1ex]
\hline
\hline
\end{longtable}
日期:2010年5月16日

主題:「恩惠慈愛隨著我」

經文:詩篇23:1-6

證道:鄧瑞強博士

各位弟兄姊妹,主內平安.今日講道的經文是:詩篇23篇.這是一篇很詳和,對生命很有洞見的詩篇.我將之分成三個段落來講.

(1)一個牧者就足夠了

23:1　耶和華是我的牧者,我必不致缺乏.

這裡有一個簡明的比喻.若我們是無助的羊,神就是牧人.

若將這比喻推而廣之,則:若我們是入不敷支的工人,神就是出價遠高於時薪二十元的好老闆；若我們是坐輪椅而要上巴士的人,神就是那位毫無不悅之色而幫助我們上巴士的司機；若我們是無助地臥床的病人,神就是那鼓勵及安慰我們的護士或醫生.

「耶和華是我的牧者,我必不致缺乏.」

這節經文最直接的解釋就是:人生有種種需要,神看到了,他幫助我們.

生命無缺,生活無憂,是人生一大樂事.

由於「耶和華是我的牧者,我必不致缺乏」這節經文裡,「我必不致缺乏」到底「不缺乏」什麼,經文沒有交代,故此,對這經文,我們可以有另外一種理解.我們可以將之理解為:有了耶和華這位牧者,我就不缺乏牧者了.

這裡,重點不在於:有了神,我們就不缺乏人生之種種需要了.

重點在於:有了這位神,我們就不需要別的神了.

三角戀之痛苦,在於「你愛這個,又愛那個」,拿不定主意,愛情不知落足何處.

做人的痛苦,在於「你想有錢,又想有時間」,但到你有錢的時候便通常無時間,到你有時間的時候便通常無錢.兩種好處不能兼得.

人的痛苦,很多時就是你想要這個神,又想要那個神.有了耶和華這個牧者還不夠,還想要另一個牧者.

禮拜日返教會時,你崇拜基督教的神.

到禮拜一返工時,你崇拜森林定律的神,不惜犧牲他人的成功而令自己上位.你又崇拜金錢這位神,為保住飯碗,不惜出賣某些道德原則.你又崇拜成功的神,為此奉獻了家庭的幸福,遺忘了別人的尊嚴.

我們想拜的神太多了,我們想要的好處太多了,我們的心太多了.

有否想過:要一個便夠呢?

有一個故事是這樣的:有一個財主,病得快死.他唯一的兒子遠在外地,身邊只有一個奴僕.奴僕問:「主人,若你死了,你的財產怎麼辦?」主人知道這是一個不忠的奴僕,於是便立了一份遺囑給這奴僕.遺囑說:「我的所有財物都歸我的奴僕,但我的兒子卻可從中取其中一樣.」財主死了.奴僕很高興地拿這遺囑去尋找主人的兒子.找到了,奴僕對財主的兒子說:「你父親的一切現在都歸我了,你只能從中取其中一樣東西.你要什麼呢?」兒子說:「你是我父親的財產的一部分,我要的一樣東西就是你.」這兒子很聰明,因為有了這奴僕,就有了一切.當然,他的父親也很聰明,若不立這樣的一份遺囑,這奴僕已靜靜地獨吞這財富了.

有了一個,便有了全部.有一位真神,其他的可以不要了.

詩篇23:1:「耶和華是我的牧者,我必不致缺乏」,向我們發出一個詢問.詩篇的作者找到生命的立足點了,他找到一個牧者,一位神,他已不再需要別的神了.你呢?

(2)生活的三個境況

我將詩篇23:2-4,看成是生活的三個境況:

一,平順的境況

23:2　他使我躺臥在青草地上,領我在可安歇的水邊.

人生有時會遇到順境,這是值得感恩的.這節詩句的用字,全是令人感到平靜的,休閒的.

「躺臥」:不用急,不用趕.不必追趕死線,也沒有必須要完成的工作.停下來,享受片刻.

躺臥在哪裡?

「在青草地上」.在新長出來的,幼嫩的草上.單單想這樣的一塊草地,已感到很舒暢了.這樣的一片草地,對羊而言,是美味大餐.

「可安歇的水邊」:這不是死水,而是緩流的水.緩流的水易於飲用,令人舒暢.羊怕驚嚇,太急的流水,水聲太大,他們會驚.但在這裡,一切都是這般緩慢,一切都是這般平安.一定不會令人患上「驚恐症」.

這樣平靜的日子,詩人想到,這是神的帶領.

二,生命低沉時

23:3　他使我的靈魂甦醒,為自己的名引導我走義路.

人生有失意時,情緒會低落,意志會消沉.這時,最容易想歪了,最容易走歪路,這時候也最容易自殺.這時,我們需要重新得力,需要有人引導我們不走歪路.哪個神能叫我們靈魂甦醒呢?哪個神能叫我們走人生的正路呢?「森林定律的神」?「金錢的神」?「成功的神」?

詩人想,唯有耶和華神,唯有這位愛人類但又對人類有正義要求的神,能叫我們心靈堅強起來,能叫我們在失意時想起人間的正路.

毛澤東說:「天若有情天亦老,人間正道是滄桑.」詩篇的詩人在此卻說:「天若有情人不老,世間正道亦康莊.」

三,死蔭幽谷中

23:4　我雖然行過死蔭的幽谷,也不怕遭害,因為你與我同在；你的杖,你的竿,都安慰我.

行過「死蔭的幽谷」,這是人生最艱難的時刻.人生有順境,有低迷,也有極艱難的時刻.「死蔭的幽谷」是危險之地,潛伏著猛獸,路也不好走.

在這些危險關頭,神和我們的關係有深刻的變化.在這節經文,詩人忽然不再用「他」去指稱耶和華,而轉用「你」這個字.字眼的改變,意謂著關係的改變,意謂著神學思想的改變.

人在順境時,在不太危急時,神只是一個「他」.一個我們認識但關係遙遠的神.甚至,我們不必與他有直接的關係.

你在彌敦道上走,無數人走過,是無數的「他」,直到有人叫「你」.原來是你的朋友,他走過來,問你:「你近來如何?」你就需要回應.你可以不管任何一個「他」,但你必須真實面對一個「你」.

神是一個「你」,則我們可以向他傾訴,他也必須回應.在生活的「死蔭幽谷」裡,神是一個具體的「你」.「你」這個字,代表著親密的關係.

生命的危機,在詩人看來,可幫我們真實地面對神,和神對話.

再者,「你的杖,你的竿,都安慰我.」

神手中有杖有竿.這裡「杖」是棍.神手中的棍和「神棍」無關.神會用手中的棍,責打不聽話的羊.你可能會說,他放下手上的棍就好了.但他不會放下的,因為這棍同時是用來打野獸,趕走野獸的.這棍也是用來保護羊的.

至於「竿」,是有彎位的長棍.當羊在幽谷裡失腳時,打撈牠回來的.

生命的「死蔭幽谷」,是艱難時刻,人會犯錯,會失腳,也會遇到外在的凶險.但是,神手中的杖和竿會帶來安慰.野獸來,有神的杖；自己滑腳失足,有神的竿.

這樣說來,人生最艱難的時刻,可能正是與神最親近的時刻.我們可能跌倒,但總有神扶持.我們可面對凶險,但神總會幫助.

(3)回望一生:無仇無恨,只是福份

23:5　在我敵人面前,你為我擺設筵席；你用油膏了我的頭,使我的福杯滿溢.

23:6　我一生一世必有恩惠慈愛隨著我；我且要住在耶和華的殿中,直到永遠.

詩篇描繪的場景改了.這是一場宴會.神是請客的主人家.神請你赴宴,多麼高興.但坐在你對面的,竟是你幾十年來的敵人(譬如說:情敵),多麼掃興.這場宴會是否是一「和頭酒」?當然,有釋經學者認為,這是神為你打勝仗後的盛宴,面前的敵人是戰俘,擺在這裡是被你玩弄的.若是這樣理解,這就是你生命的最大「慶功宴」.

在這盛宴中,神「用油膏了我的頭,使我的福杯滿溢」.用中國人的講法,這主人夾隻雞脾給你,還為你斟酒,甚至滿瀉.若果這是一「和頭酒」,主人的盛情能否使你下氣,和敵人化解怨仇?若果這是一場「慶功宴」,則主人的盛情真令你一生的努力得到最完滿的成果.

我認為,這首詩的結尾部分,是一個飽經世事的老者,回望一生,詳和地講出他的信仰結論.他認為:我們一生充滿仇怨,種種惡人惡事纏在我們心中,但神作為主人,他會主持公道的.在神手中,緊纏著我們的,不應是痛苦仇恨.那又是什麼呢?經文說:「我一生一世必有恩惠慈愛隨著我」.

這裡的「隨著我」的用字,意指「緊纏著我不放,追趕著我」.我們日日逃避,日日掙扎,以為過去的痛苦抓住我們不放,但其實,纏著我們生命不放的,不是什麼惡人惡事,而是神的恩惠慈愛.追著我們的,只是神的種種祝福.

這是詩人的一種善意的諷刺.我們逃避的東西,會否其實只是神的祝福?帶著生命的無憾和滿足,詩人的心就以神為居所,進入永恆.「我且要住在耶和華的殿中,直到永遠.」

但願你被神的恩惠慈愛捉住.但願你明白,緊隨你一生的,是神的恩惠慈愛.

但願榮耀歸於聖父,聖子,聖靈.阿們.



\newpage

\section{尋找一棵香柏樹}
\label{sec:48}
\textbf{梁元生教授 - 尋找一棵香柏樹}
\newline
\newline
\hyperref[sec:47]{< < < PREV SERMON < < <}
~
\hyperref[sec:toc]{[返目錄]}
~
\hyperref[sec:49]{> > > NEXT SERMON > > >}


以西結書 17:1-3,22-24
\newline
\begin{longtable}{cl}
\hline
\hline
章節 & 經文 (和合本修訂版)\\
\hline
17:1 & \begin{tabularx}{0.7\textwidth}{X} 耶和華的話臨到我，說： \end{tabularx} \\ \\ \relax
17:2 & \begin{tabularx}{0.7\textwidth}{X} 「人子啊，你要向以色列家出謎語，設比喻， \end{tabularx} \\ \\ \relax
17:3 & \begin{tabularx}{0.7\textwidth}{X} 說，主耶和華如此說：有一隻大鷹，翅膀大，翎毛長，羽毛豐滿，色彩繽紛；牠飛到黎巴嫩，啄去香柏樹梢， \end{tabularx} \\ \\ \relax
17:4 & \begin{tabularx}{0.7\textwidth}{X} 啄斷它頂端的嫩枝，叼到貿易之地，放在商業城中。 \end{tabularx} \\ \\ \relax
17:5 & \begin{tabularx}{0.7\textwidth}{X} 牠又從這地取了一些種子，種在肥沃的田裡，栽於豐沛的水源旁，如種植柳樹。 \end{tabularx} \\ \\ \relax
17:6 & \begin{tabularx}{0.7\textwidth}{X} 它漸漸生長，成為低矮蔓生的葡萄樹；樹枝伸向那鷹，根部在牠下面。這樣，它就長成了一棵葡萄樹，生出枝子，長出枝幹。 \end{tabularx} \\ \\ \relax
17:7 & \begin{tabularx}{0.7\textwidth}{X} 「有一隻大鷹，翅膀大，羽毛多。看哪，葡萄樹從栽種它的苗圃向這鷹伸出根來，長出枝子，期盼從牠得到澆灌。 \end{tabularx} \\ \\ \relax
17:8 & \begin{tabularx}{0.7\textwidth}{X} 這棵樹栽於肥田豐沛的水源旁，原是為了生枝、結果，成為佳美的葡萄樹。 \end{tabularx} \\ \\ \relax
17:9 & \begin{tabularx}{0.7\textwidth}{X} 你要說，主耶和華如此說：這棵葡萄樹豈能發旺呢？鷹豈不拔出它的根來，摘光它的果子，使它枯乾，連長出的嫩葉都枯萎了嗎？要把它連根拔除，並不需要費大力或動用許多人。 \end{tabularx} \\ \\ \relax
17:10 & \begin{tabularx}{0.7\textwidth}{X} 看哪，葡萄樹雖然栽種了，豈能發旺呢？一經東風擊打，豈不全然枯乾了嗎？它必在生長的苗圃中枯乾了。」 \end{tabularx} \\ \\ \relax
17:11 & \begin{tabularx}{0.7\textwidth}{X} 耶和華的話臨到我，說： \end{tabularx} \\ \\ \relax
17:12 & \begin{tabularx}{0.7\textwidth}{X} 「你要對那悖逆之家說：你們不知道這些事是甚麼意思嗎？你要這樣說，看哪，巴比倫王曾到耶路撒冷，把其中的君王和官長帶到巴比倫去， \end{tabularx} \\ \\ \relax
17:13 & \begin{tabularx}{0.7\textwidth}{X} 又從以色列王室後裔中選取一人，與他立約，令他發誓，又擄走國中有勢力的人， \end{tabularx} \\ \\ \relax
17:14 & \begin{tabularx}{0.7\textwidth}{X} 使王國衰弱，不再強盛，只能靠守盟約方得生存。 \end{tabularx} \\ \\ \relax
17:15 & \begin{tabularx}{0.7\textwidth}{X} 他卻背叛巴比倫王，差派使者前往埃及，要求埃及人給他馬匹和許多人。他豈能亨通呢？這樣做的人豈能逃脫呢？他背了約豈能逃脫呢？ \end{tabularx} \\ \\ \relax
17:16 & \begin{tabularx}{0.7\textwidth}{X} 主耶和華說：我指著我的永生起誓，他定要死在巴比倫，就是巴比倫王所在之處；因為巴比倫王立他為王，他竟輕看向王所起的誓，背棄王與他所立的約。 \end{tabularx} \\ \\ \relax
17:17 & \begin{tabularx}{0.7\textwidth}{X} 當敵人建土堆，築堡壘，要殲滅許多人時，法老雖有強大軍隊和大批人馬，在戰場上還是不能幫助他。 \end{tabularx} \\ \\ \relax
17:18 & \begin{tabularx}{0.7\textwidth}{X} 他輕看誓言，背棄盟約，看哪，雖已投降，卻又做這一切的事，他必不能逃脫。 \end{tabularx} \\ \\ \relax
17:19 & \begin{tabularx}{0.7\textwidth}{X} 所以主耶和華如此說：我指著我的永生起誓，他既輕看我的誓言，背棄我的約，我必使這罪歸到他頭上。 \end{tabularx} \\ \\ \relax
17:20 & \begin{tabularx}{0.7\textwidth}{X} 我要把我的網撒在他身上，他就被我的羅網纏住。我要帶他到巴比倫，在那裡因他背叛我的罪懲罰他。 \end{tabularx} \\ \\ \relax
17:21 & \begin{tabularx}{0.7\textwidth}{X} 所有逃跑的軍隊必倒在刀下；剩餘的也必分散四方。你們就知道說這話的是我---耶和華。」 \end{tabularx} \\ \\ \relax
17:22 & \begin{tabularx}{0.7\textwidth}{X} 主耶和華如此說：「我要從香柏樹高高的樹梢摘取並栽上，從頂端的嫩枝中折下一嫩枝，栽於極高的山上， \end{tabularx} \\ \\ \relax
17:23 & \begin{tabularx}{0.7\textwidth}{X} 栽在以色列高處的山上。它就生枝、結果，成為高大的香柏樹，各類飛禽中的鳥都來宿在其下，宿在枝子的蔭下。 \end{tabularx} \\ \\ \relax
17:24 & \begin{tabularx}{0.7\textwidth}{X} 田野的樹木因此就知道是我---耶和華使高樹矮小，使矮樹高大，使綠樹枯乾，使枯樹發旺。我---耶和華說了這話，就必成就。」 \end{tabularx} \\ \\
[1ex]
\hline
\hline
\end{longtable}
日期:2010年5月30日

題目:尋找一棵香柏樹

經文:以西結書17:1-3,22-24

証道:梁元生教授

我記得大概在兩個月前,有一個人叫奚淞從台灣來到香港,在香港大學藝術館舉行了一個畫展,主題叫做「尋找一棵菩提樹」,我看了心裡有一點觸動.　後來又看到他在城市大學做了一場的演講,題目又是「尋找一棵菩提樹」.　奚淞這人是誰呢?他的畫展跟他的演講是怎麼一回事呢?

奚淞是台灣的一位畫家,一位藝術家並作家.他是以畫觀音像出名的,是一個著名小說家白先勇的朋友.白先勇籌備,改編他的第二部青春版的昆劇《玉簪記》上演的時候,這舞台的佈景,需要一些佛教的畫像,於是白先勇就想起以畫觀音著名的奚淞,奚淞用毛筆白描的方法,就替這《玉簪記》的舞台的背景畫了很多佛手和觀音,為《玉簪記》的演出帶來非常好的效果.配合這《玉簪記》在香港的上演,白仙勇就請了奚淞來香港,以尋找一棵菩提樹為主題,在香港大學的藝術館展出他的佛手觀音的畫像和書法.他來到香港以後,遇見他小學時候的一個同學並好朋友,就是城市大學中國文化中心的主任鄭培凱教授,他就請奚淞做演講,他就講「尋找一棵菩提樹」.這跟我們有什麼關係呢?奚淞的話跟奚淞的生命的追求,都跟佛教修行有關係的.奚淞有一件事對我很感動的,就是奚淞他的媽媽的死.他的爸爸很早就去世了,所以他一直以來跟媽媽相依為伴的,感情非常好.年青的時候奚淞因為要學藝術,出國到了歐洲,去了法國,離開家庭,離開母親.後來他的媽媽因為患癌症,他就停了他的學業回到台灣侍候媽媽.每天為他媽媽誦讀佛經,和畫觀音像,就發現他母親稍微好了一些,但最後他媽媽還是死了.他媽媽死的時候,他心裡非常悲慟,他每天連續不斷的三十三天,為她造觀音像衣服.　這一次的展覽裡面,就是展出三十三幅的觀音圖,主題就是以慈悲的眼睛看世界.奚淞說他作畫三十三天,只是在體會這種憐憫與慈悲的一個修行,不是為藝術的創作.換而言之,他的「尋找一棵菩提樹」,是在痛苦之中去體會愛,慈悲,憐憫,是人生修行裡面的一種磨鍊.雖然他的淵源是彿教,我覺得可以把這個故事帶到來我們中間,反醒,借鏡,這是一顆宗教的心.

我今天想借著奚淞「尋找一棵菩提樹」的話題,談談人生的歷練和修行,但是是用一個基督教的背景來談,所以這個菩提樹不是我要尋找的,我要跟大家一齊去尋找的,是一棵香柏樹.

香柏樹是我們今天要講的,在尋找,栽種... 基督教常常談到這香柏樹的時候,聖經也有很多給我們思考之處,當然聖經裡也提到很多種不同的樹,最常見的是我們今天本來要唱的尋找一棵葡萄樹.耶穌說:我是葡萄樹,你們是枝子,我父是栽培的人.聖經也常常提到無花果樹,你看祂所有的門徒,最後看到的一個門徒拿但業是在無花果樹的.聖經也有講到橄欖樹(下次如果談這橄欖樹的話,我要請台灣的歌星來唱這橄欖樹).我們也有大麻樹和桑樹.

今天我們特別講到的是香柏樹.聖經也常常提到黎巴嫩的香柏樹,如果從這個本土的神學的方面去想,我們也可能去伸延這題目到尋找一棵白楊樹,尋找一棵山西的老槐樹,尋找一棵老榕樹,這在我們華南的校園最多,或是櫻桃樹,越南的神學來說可能尋找一棵木瓜樹.在我們想像當中,這些樹如人生,身邊任何的樹都可以帶給基督徒一種反思,一種啟示,都可以讓我們人生的修行是有幫助的.但今天不可能都把這些樹,生命樹的系列都講完了,將來有時間,作這有系統思考以後,跟大家分享討論.今天我要順著奚淞在香港牽動的一點點的思潮,改動了一下題目,以這尋找一棵香柏樹為題跟大家談,本來要後面加一句的,尋找一棵香柏樹,復活的樹,因為我們還在復活節的期間,我也希望從這方面多思考一下.

現在我先介紹一下這香柏樹,在中東的地區,特別在以色列,黎巴嫩一帶,香柏樹是一種常見的灌木,樹幹高聳如山,枝葉繁茂覆蓋,木材堅實耐久,非常良好的建材.那個地區很多的房屋都用這香柏樹作為材料.我們知道聖經所說的,連古時候所羅門建的聖殿,都以佳美的香柏樹,良好的香柏樹為主要的棟樑.所以長久以來,香柏樹在黎巴嫩的民族,歷史裡面,是一個文化的符號.今天黎巴嫩的國旗也是用香柏樹作為她的標誌.今天思考的香柏樹,是從什麼地方開始?

我請大家來看一段聖經,是舊約的以西結書十七章:

17:22　主耶和華如此說,我要將香柏樹梢擰去栽上,就是從儘尖的嫩枝中折一嫩枝,栽於極高的山上．

17:23　在以色列高處的山栽上．他就生枝子,結果子,成為佳美的香柏樹,各類飛鳥都必宿在其下,就是宿在枝子的蔭下.

17:24　田野的樹木都必知道我耶和華使高樹矮小,矮樹高大,青樹枯乾,枯樹發旺．我耶和華如此說,也如此行了.

聖經裡面還有其他地方提到香柏樹的,剛才也讀過:詩篇104:16佳美的樹木,就是利巴嫩的香柏樹,是耶和華所栽種的,都滿了汁漿. 撒迦利亞書11:1利巴嫩哪,開開你的門,任火燒滅你的香柏樹. 耶利米書22:23 你這住利巴嫩在香柏樹上搭窩的,有痛苦臨到你,好像疼痛臨到產難的婦人,那時你何等可憐.今天特別要大家注意的香柏樹,就是剛才以西結書十七章22-24節的香柏樹.除了剛才給大家特別讀的22節以外,實際上以西結書的十七章,最前面也說了很大段是關於香柏樹的.

17:1　耶和華的話臨到我說,

17:2　人子阿,你要向以色列家出謎語,設比喻,

17:3　說,主耶和華如此說,有一大鷹,翅膀大,翎毛長,羽毛豐滿,彩色俱備,來到利巴嫩,將香柏樹梢擰去．

17:4　就是折去香柏樹儘尖的嫩枝,叼到貿易之地,放在買賣城中．

跟剛才我們讀的都有相同的故事.一個是大鷹來到,把香柏樹樹林的嫩枝折斷,然後拿走.後段就是說到耶和華來到,把這香柏樹的嫩枝擰下折斷,都是同一個動作.對於香柏樹來說,這是一個最重要的動作.香柏樹本來有一片的林,我叫它香柏樹林.這香柏樹林是比喻今天我們每一個人的過去.每一個人都有一片香柏樹林,每一個基督徒在回顧的時候,他的過去,就是一片香柏樹林.

請大家注意這經文裡面,兩段的經文裡面,都有三個動作:1.    擰去2.    折去3.    栽上

今天跟大家講的這個故事,擰下,折斷的嫩枝再植,再栽的一個故事.一段嫩枝被擰下來,被折斷,然後這斷枝再獲得一個復活重生的生命,這就是基督徒的人生,像一棵香柏樹的枝子.

每個人包括基督徒都有一個過去,都有一個童年,都有一段從小成長的日子,以前我在我的母校講過的,那一片香柏樹林.我跟大家一同回想,一同回憶.當然剛才彈琴的小朋友,回憶的時候他不容易看到那一片的綠.這童年跟過去,如果當我們集中回憶的時候,我想大家腦袋裡面都會有一片的綠,就是那一片香柏樹林.聖經說到大鷹飛臨之前的那片香柏樹林.我們信主以前的人生,基督徒的過去,都在一層香柏樹之中,有一些在那邊變成是綠葉,襯托整個樹林,讓它發出一片的青綠,也有一些是嫩枝,希望能夠在那裡發芽,繼續的生長.但是我們知道在這個香柏樹林裡面,有一些是開枝散葉,有一些最後隨風飄散,也有一些就留在原地,也有一些葉子慢慢的枯乾了.那一片香柏樹林,我們的過去,我們的童年,這回憶當中的過去,我們就會看到有一些好像枝子一樣,有一些好像綠葉,大家共同襯托青蔥的綠色的樹林,我們擠在一齊,散發一團的綠.我們的過去就叫做青蔥歲月,然後我們各散東西,有一些好像斷枝零落,留下像殘葉飄散.有一些枝子,在聖經所說的,就是被耶和華選上了,如果不是基督徒,當然是說命運安排他離開這個樹林了.他要折枝叼去,有一些就留在原地.

我用這個比喻說,有一些是像葉子,葉子經過歲月的洗禮,就變得枯乾,逐漸的衰殘,最終在秋天來的時候,葉子落下枯死.我們也一樣的,經過了一段漫長的歲月,童年要過去,少年要過去,青年也要過去,甚至我們的盛年也會過去,就如那一片香柏樹林,也都會成為記憶,但是我們在回憶之中,還是可以留一片的綠.彷彿還是有一片香柏樹林在那邊.但現實上,我們現在都長大了,至少對我來說,年華老去了,走上遲暮,好像殘了枯乾的葉子,葉子殘了枯黃以後,掉下「化作春泥更護花」.殘葉是一個盡頭,沒有下一個重生的故事,所以大家對這個香柏樹林一部份的綠,也應該了解就變成葉子死了,就是他的盡頭.但是這枝子卻是另外一個故事,是一個復活的故事.今天所以就不再談葉子,單單談枝子,香柏樹的枝子.

經文是這樣的記載的.耶和華說:祂會將一些枝子折斷,然後把這些枝子栽在高山上.一個信主的人,是一段再續的斷枝.耶和華來到香柏樹林,還是大鷹來到這個香柏樹林,是要把嫩枝折斷.可能是耶和華折斷的一段枝子,但耶和華選上了,要把你選擇上了,就是信主的人,就變是上帝的選民,不一定表示以後的生命就順風順水,或是平步青雲,這一切的順利,實際上被選上,對我自己個人的經驗來說,那個才是受苦的開始.這一點,我要常常提醒初信的基督徒,以為接受了這個基督教,接受了這個信仰,接受了這個救恩,以後就一切都平坦,一切就安穩,一切就順利.上主把你呼召,把你選上,不一定給你安排一條平穩安樂的一條吉路,更加有可能的話,選上就是痛苦的開始.聖經中所說的三個動作,其中至少一看這兩個都不容易承受的: 第一個是擰下,英文是Twist,我們常常用擰,用扭這個動作,來對付我們很頑皮的小孩,擰他一下,扭他一下,我太太說咬他一下,這都是要讓他痛一下的.這擰下的動作,是製造一點痛苦的動作.擰你的耳朵,擰你的臉龐,擰你的手臂,另外跟按摩的動作是不一樣的.一個是帶來舒服的感覺,一個是帶來痛苦的行為. 另外一個是折斷,英文是breaking off,是更加嚴重的動作,要把你折斷了,身體某部分的折斷,還是你的生命受到被折斷的壓力.折斷與分離是生命中不容易承受的重.這比擰下更加的痛,更加的壓力大.有的時候耶和華的選上,就是帶來這兩下的動作,首先把你生命擰下,扭動,折斷,分離.對於那些期望信主以後帶給他一切平安,一切順利的基督徒來說,大概不喜歡我今天的分享.尤其是對很多家庭安逸,家庭溫暖的人.選上了被耶和華的大鷹來到,折斷嫩枝,擰下樹梢的這種動作.我要在這裡向大家做一個提醒,本來我有很大段自己個人的見證,什麼是擰下的經驗?什麼是痛苦的經驗?什麼是折斷的經驗?今天因時間的關係,有機會再跟大家分享.我回顧我的童年,經過不少的搬遷,不少的考驗,然後不少的擰下,折斷的經驗,之後選上生命帶來的同時候的一種經歷. 把枝子擰下折斷以後,還有栽種.栽種當然是一件好事情,也沒有看到有什麼特別痛苦的牽連.他代表了上帝選擇你,然後培育你的心意.耶和華把香柏樹的枝子栽在什麼地方呢?聖經告訴我們,耶和華把香柏樹的枝子栽在高山之上,栽在寒冷之中,栽在孤獨之處.只有高山才會寒冬,只有高山才會人跡罕至,只有高山才會是一種孤立獨力自存,「高山巍巍我獨寒,白雪凱海漫我身」,這是我自己隨便寫的打油詩.天地之間無依傍,斷枝栽栽說奈何.為什麼要栽在高山?如果是上帝給你選上作為一個基督徒的生命,為什麼先讓你經歷擰下的痛苦,折斷的沉重,還要栽在寒冷高山,栽在孤獨無人的地方?高山的經歷也是每一個基督徒需要經過的經歷.上帝選擇帶領枝子折斷以後,是要離開他的香柏樹林,他要移植到高山上不與同群,受寒風受雨雪,冰雪侵凌之後,我們看到香柏樹長出來新芽,泥土下面長了根,枝幹慢慢的粗壯,高山上面長出了一棵香柏樹,是不怕風雪飄搖,不怕孤山孤寂,不怕獨立無伴在朝陽之下,在春風之中,在這露水的滋潤和大地的孕育之下,香柏樹變成一棵活潑的生命樹.他在不斷的生長茁壯,奮鬥中成為一棵非常高大的,聖經說:佳美的香柏樹.他的努力和奮鬥,在這兩方面,明顯的表現出來.聖經的話來說:一個是生枝子,一個是結果子.在一方面他生出其他的嫩枝,另一方面嫩枝會長大結出果子,自己明明是有大樹折斷下來的一根嫩枝,如今在他枝子的生命上,又在生出很多的枝椏.這在我們念文化,念歷史的人說:叫做傳承.枝椏上面結出果實,讓生命得以延續,讓生命變成不同的形式,不同的形象出現,讓生命變得更加的多元,更加的豐富,這個叫做轉化(也是我們文化歷史學者的術語).創造的轉化,在這枝子栽植之後,他身上活出了傳承,轉化的香柏樹,是一棵我們要尋找的,有復活生命的香柏樹.這就是我們希望見到基督徒,斷枝,栽植以後的表現.他的身上起了變化,他的身上既能傳承一種上帝付予的生命,也可以變化,發展出一個多元的,不同的生命的模式.我自己當然也經歷過神的選擇,神的考驗,然後帶來生命的變化.當然我不能說我已經成為大樹,只是今天跟大家分享.

香柏樹經過這個階段以後,就變成一棵佳美的香柏樹,基督徒進一步的追求,就是把自己轉化,建立成為一個大樹的生命.經過風霜雨雪,受盡了孤獨,嚴寒.香柏樹從斷枝到栽植.我記得唐君毅教授也說過,香港有很多很多的文化人,就像飄零在海外的果子,變成連根栽植的一段故事.香柏樹從斷枝到栽植的新苖,再變成一棵長成的大樹,日出日落,月明月缺,經過了多番的歲月,多少的光陰,香柏樹現在成為一棵葉茂幹粗的大樹.這種的描述也在聖經上很清楚的看到.聖經這樣說:香柏樹就生枝子,結果子,成為佳美的樹,各類飛鳥都必宿在其下,就是宿在枝子的陰下.長大以後的香柏樹有以下的特徵:香柏樹,枝條榮美,影密如林,極其高大,樹尖插入雲中.眾水使他生長,他高大超過田野諸樹,發旺的時候,枝子繁多．因得大水之力枝條長長.空中的飛鳥,都在枝子上搭窩．田野的走獸,都在枝條下生子．所有大國的人民,都在他蔭下居住.樹大條長,成為榮美.(以西結31章3-7節)這都是聖經裡面的話,我把他讀出來讓大家可以感受一下,可以想像一下,這香柏樹長大以後的樣子是怎麼樣.

我特別是從事教育的人,跟作為基督徒生命都可以感受到這栽種再長的香柏樹.我們都希望變成枝條繁茂,綠葉成蔭的樣子,佳美的香柏樹的樣子.因這是活潑的,豐富的,堅持的生命.但是最重要的,眾鳥棲息的地方.我們辦教育的,辦學校的,如果辦得不好,學生們不在這裡停留,棲息,就算在這裡停留,棲息,也沒辦法在這裡築巢作窩,可以有一個安樂的地方,我們不單是要為他們提供學位,提供一個活動的空間而已,我們希望樹上長出的葉子,是可以遮蔭的.我們希望樹上長出的枝子,是可以替他們作巢的,讓每一個人來到香柏樹下,都有一個家的感覺.每一天在這樹下面,天空之下,可以自由的飛翔,但是在飛翔之後,讓他們可以回到這棵大樹,可以在這棵大樹的枝子上面停留,可以在這樹上所造的巢安居.同樣的,一個基督徒,生命好像佳美的香柏樹,生枝子,結果子,是讓別人來遮蔭,讓鳥來築巢,讓貧乏的人可以得到保護,讓需要的人可以得到安慰.

大家身處在中文大學,或者可以想想,你也可以是一棵大樹.這禮拜堂是屬於崇基學院的.在某程度上是一棵香柏樹,大樹成蔭,眾鳥來棲.我們二千多的學生,每年有新同學七百多人.當然我們希望把這個地方,像一棵香柏樹一樣開放,讓他們都可以來,好成績的人可以,成績差的人可以,家庭背景富有的人可以,貧困的人可以,眾鳥來棲.我們長成的香柏樹,是一個開放的香柏樹.當然這個教堂也是我們的香柏樹.不少人有心靈的問題,不少人有精神的問題,不少人有家庭的問題,希望在這裡可以找到造窩的地方,可以找到安息之地,可以找到作巢之處,給你一個家的感覺.希望這個是香柏樹的故事.

今天跟大家一齊回到過去,想一想無論你們有多大年紀,總會有一個過去.你一閉眼睛,你就看不到一片綠色的香柏樹林,一片的綠,那是你們童年,青春的歲月.想一想有沒有經過擰下的考驗,想一想有沒有經過折斷的考驗,想一想有沒有經過栽在高山的經歷,如果有的話,你應該現在看到新芽,看到新枝,應該長出了一個華蓋一樣大的樹葉,你的葉子可以給很多需要的飛鳥來在下面遮蔭.你的枝子可以承托不同的巢,不同的窩,給人有一個家的感覺,我們尋找這棵香柏樹.



\newpage

\section{離開祇留下愛}
\label{sec:49}
\textbf{林豪恩先生 -}
\newline
\newline
\hyperref[sec:48]{< < < PREV SERMON < < <}
~
\hyperref[sec:toc]{[返目錄]}
~
\hyperref[sec:50]{> > > NEXT SERMON > > >}


約翰福音 14:15-31
\newline
\begin{longtable}{cl}
\hline
\hline
章節 & 經文 (和合本修訂版)\\
\hline
14:15 & \begin{tabularx}{0.7\textwidth}{X} 「你們若愛我，就會遵守我的命令。 \end{tabularx} \\ \\ \relax
14:16 & \begin{tabularx}{0.7\textwidth}{X} 我要求父，父就賜給你們另外一位保惠師，使他永遠與你們同在。 \end{tabularx} \\ \\ \relax
14:17 & \begin{tabularx}{0.7\textwidth}{X} 他就是真理的靈，是世人不能接受的。因為他們既看不見他，也不認識他；你們卻認識他，因他常與你們同在，也要在你們裡面。 \end{tabularx} \\ \\ \relax
14:18 & \begin{tabularx}{0.7\textwidth}{X} 我不會撇下你們為孤兒，我必到你們這裡來。 \end{tabularx} \\ \\ \relax
14:19 & \begin{tabularx}{0.7\textwidth}{X} 再過不久，世人不再看見我，你們卻會看見我，因為我活著，你們也要活著。 \end{tabularx} \\ \\ \relax
14:20 & \begin{tabularx}{0.7\textwidth}{X} 到那日，你們就會知道我在父裡面，你們在我裡面，我也在你們裡面。 \end{tabularx} \\ \\ \relax
14:21 & \begin{tabularx}{0.7\textwidth}{X} 有了我的命令而又遵守的人，就是愛我的；愛我的人，我父要愛他，我也要愛他，並且要親自向他顯現。」 \end{tabularx} \\ \\ \relax
14:22 & \begin{tabularx}{0.7\textwidth}{X} 猶大（不是加略人猶大）問耶穌：「主啊，為甚麼親自向我們顯現，而不向世人顯現呢？」 \end{tabularx} \\ \\ \relax
14:23 & \begin{tabularx}{0.7\textwidth}{X} 耶穌回答他說：「凡愛我的人就會遵守我的道，我父也會愛他，並且我們要到他那裡去，與他同住。 \end{tabularx} \\ \\ \relax
14:24 & \begin{tabularx}{0.7\textwidth}{X} 不愛我的人就不遵守我的道。你們所聽見的道不是我的，而是差我來之父的。 \end{tabularx} \\ \\ \relax
14:25 & \begin{tabularx}{0.7\textwidth}{X} 「我還與你們在一起的時候，已對你們說了這些事。 \end{tabularx} \\ \\ \relax
14:26 & \begin{tabularx}{0.7\textwidth}{X} 但保惠師，就是父因我的名所要差來的聖靈，他要把一切的事教導你們，並且要使你們想起我對你們所說的一切話。 \end{tabularx} \\ \\ \relax
14:27 & \begin{tabularx}{0.7\textwidth}{X} 我留下平安給你們，我把我的平安賜給你們。我所賜給你們的，不像世人所賜的。你們心裡不要憂愁，也不要膽怯。 \end{tabularx} \\ \\ \relax
14:28 & \begin{tabularx}{0.7\textwidth}{X} 你們聽見我對你們說過，我去了還要回到你們這裡來。你們若愛我，就會因我到父那裡去而喜樂，因為父比我大。 \end{tabularx} \\ \\ \relax
14:29 & \begin{tabularx}{0.7\textwidth}{X} 現在事情還沒有發生，我預先告訴你們，使你們在事情發生的時候會信。 \end{tabularx} \\ \\ \relax
14:30 & \begin{tabularx}{0.7\textwidth}{X} 我不再和你們多說了，因為這世界的統治者將到，他在我身上一無所能。 \end{tabularx} \\ \\ \relax
14:31 & \begin{tabularx}{0.7\textwidth}{X} 我這麼做是照著父命令我的，為了讓世人知道我愛父。起來，我們走吧！」 \end{tabularx} \\ \\
[1ex]
\hline
\hline
\end{longtable}
日期:2010年6月6日

經文:約翰福音14:15-31節

証道:林豪恩先生

離開,難以理解的事實有人稱約翰福音14-17章為「離別的談話」,記載了耶穌向門徒所說「最後」的話.讓我們把今天的福音經課擴闊一點,提前至第15節開始.當我們讀這經文的時候,嘗試感受當中的離別意味.

15「你們若愛我,就必遵守我的命令.

16 我要求父,父就另外賜給你們一位保惠師,叫他永遠與你們同在,

17就是真理的聖靈,乃世人不能接受的；因為不見他,也不認識他.你們卻認識他,因他常與你們同在,也要在你們裡面.

18我不撇下你們為孤兒,我必到你們這裡來.

19還有不多的時候,世人不再看見我,你們卻看見我；因為我活著,你們也要活著.

20到那日,你們就知道我在父裡面,你們在我裡面,我也在你們裡面.

21有了我的命令又遵守的,這人就是愛我的；愛我的必蒙我父愛他,我也要愛他,並且要向他顯現.」

22猶大(不是加略人猶大)問耶穌說:「主啊,為甚麼要向我們顯現,不向世人顯現呢?」

23耶穌回答說:「人若愛我,就必遵守我的道；我父也必愛他,並且我們要到他那裡去,與他同住.

24不愛我的人就不遵守我的道.你們所聽見的道不是我的,乃是差我來之父的道.

25「我還與你們同住的時候,已將這些話對你們說了.

26但保惠師,就是父因我的名所要差來的聖靈,他要將一切的事指教你們,並且要叫你們想起我對你們所說的一切話.

27我留下平安給你們；我將我的平安賜給你們.我所賜的,不像世人所賜的.你們心裡不要憂愁,也不要膽怯.

28你們聽見我對你們說了,我去還要到你們這裡來.你們若愛我,因我到父那裡去,就必喜樂,因為父是比我大的.

29現在事情還沒有成就,我預先告訴你們,叫你們到事情成就的時候就可以信.

30以後我不再和你們多說話,因為這世界的王將到.他在我裡面是毫無所有；

31但要叫世人知道我愛父,並且父怎樣吩咐我,我就怎樣行.起來,我們走吧!」

在此之前,耶穌已經多次向門徒透露自己不久之後就會離開,只是,門徒對耶穌的預告似乎不以為然.當耶穌對彼得說:

「我所去的地方,你現在不能去．．．」彼得回應說:

「主啊,我為甚麼現在不能跟你去?我願意為你捨命!」(約13:36-37)可見以彼得為代表的門徒對於耶穌預告自己的離開似懂非懂.耶穌要離開的預告及臨別的囑咐未能令門徒領悟他們的夫子即將離開的事實,反而令他們感到不解及迷糊.

「還有不多的時候,世人不再看見我,你們卻看見我」,「我去還要到你們這裡來」,這些說話是甚麼意思呢?耶穌是離開還是不離開呢?耶穌離開是真實還是幻覺呢?從人們的經驗來說,離開,從來都是一種難以捉摸,也不容易確實的狀況.門徒難以理解及接受耶穌離開的事實,因為按著自然及常理,才三十而立的耶穌方華正盛,漸露頭角,與百年歸老之期相距甚遠,也很難聯想在一起.不按自然發展的離開,有異於常態的離開,違反人們期望的離開,一定是不能理解及難以接受的.戀人提出分手,同學來哭訴:「他為甚麼要走?他為甚麼要離開?不會是真的,他心情好些就會再來找我.他一定會來找我．．．」按著常理及自然的期望,談戀愛,當然是預期長相廝守.違反期望的離開,怎可能輕易接受呢?耶穌離開,不切合門徒的想像及期望,然而,耶穌語重心長地對門徒吩咐後事,讓人看到,縱使門徒不能理解,也未必願意接受,但耶穌即將離開卻是真實.

離開,不等於遺棄現在,耶穌尚未離開,門徒尚未能感受到他們的夫子的離愁別緒.當耶穌真正離開的時候,門徒可能會出現多種負面的感受及混亂的思緒.在門徒可能有的眾多的感受及思緒中,我們可以從耶穌對門徒說的話中得知耶穌特別關注的層面.耶穌對門徒說:

「我不撇下你們為孤兒」(v18),「你們心裡不要憂愁,也不要膽怯」(v27)離開,留下的人即時最直接及最強烈的感覺是被撇下,被遺棄.失戀的人哭叫:「為何拋棄我?不要拋棄我!」在生離死別的地方,也常常聽到類似的號啕大哭.在此刻,滾滾不絕的淚水,怎也衝不走心裡的混亂,驚惶,和憂愁.

耶穌接受門徒會有被遺棄的感受,也接受門徒有憂愁及驚懼的反應.這是人之常情,也是合理的反應.不過,縱使接納門徒有被遺棄的感受,耶穌仍然對門徒說明,也希望門徒明白,他的離開絕對不是要撇下門徒.即是說,留下的人有被遺棄的感覺是真實,但離開的人沒有遺棄留下的人的動機也是真實.反過來說,離開的人其實期望著留下的人不要沈溺在被遺棄的感受中.離開,不等於撇下.

我們有一位學生在兩歲的時候,媽媽因病去世.他長大後,姨媽把媽媽的錄音帶交給他.媽媽透過錄音對他說:「仔仔,媽媽和你說話就是這一次了.你從未聽過媽媽的聲音,因為年紀還小.你兩歲多媽媽便不在你身邊了.你要生性做人啊,因媽媽不能負起母親的責任,但婆婆,姨媽,和其他親戚都會很疼痛你的.雖然這份愛是有點缺憾,但你一樣會很幸福的．．．」(《我們從這裡說起》,頁58)

媽媽豈會為了撇棄兒子而離開呢?縱然要離開,都寄望兒子不要有被遺棄的感覺,努力追求他的幸福.

二十一年前的差不多這個日子,即是一九八九年五月中,一群大學生在北京天安門廣場絕食.想到自己即將離開,他們在絕食書如此向他們的父母解釋:「當我們挨餓時,爸爸媽媽們,請不要悲哀；當我們告別生命時,叔叔阿姨們,請不要傷心,我們只有一個願望,那就是讓你們能更好地活著」.

兒女們豈會為了撇棄父母而離開呢?縱然要離開,都寄望父母能夠活得更好.

耶穌即將離開門徒,耶穌也不希望門徒以為他要撇棄他們,因此,他清楚地對門徒說明:「我不撇下你們為孤兒」,「你們心裡不要憂愁,也不要膽怯」.

離開,不是被擊敗,而是另有安排對於門徒來說,為甚麼憂愁呢?為甚麼膽怯呢?他們將會看到宣揚真理的師父被宗教領袖拘捕,他們將會看見本來滿有能力的師父被士兵鞭打,他們將會看見叫死人復活的師父被釘死在十字架上.師父被擊敗了,師父也輸了,怎可能不憂愁呢?怎可能不膽怯呢?是甚麼力量能夠連如此滿有能力的師父也打敗呢?這種力量一定很神秘,一定很可怕,一定很凶悍,一定很邪惡；這種力量一定令人感到恐懼,令人感到威脅,令人感到被操控,令人感到任其魚肉.

耶穌離開,門徒以為耶穌被神秘力量打敗,這實在也是人們常見的想法.我們的語言,也反映著這些觀念,例如某某人不敵病魔,某某人被死神打敗,某某人被神秘力量苦害．．．因此,有些人為了保住性命,保住戀人,保住工作．．．要討好某些神秘力量,拜祭某些鬼神.

耶穌知道門徒會這樣想像,會這樣感覺,因此向門徒澄清:

「以後我不再和你們多說話,因為這世界的王將到.他在我裡面是毫無所有」.耶穌希望門徒不要單單憑眼睛所看見的來理解自己的離開,而要以耶穌自己的解釋來看待耶穌的離開.門徒親眼看見的事情不容否認,然而,門徒需要學習相信耶穌所啟示的更是真實.是的,門徒將會看見世界的統治者把他們的師父處決,但其實世界的統治者對於耶穌是無能為力的.即是說,耶穌若要反抗,是有能力反抗的；耶穌若要逃走,是有能力逃走的.耶穌要門徒弄清楚,他離開,不是因為他輸了,不是因為他被打敗,更不是被甚麼神秘可怕的力量作弄.因此,耶穌安慰門徒無需憂愁,也不用膽怯.

既然「世界的王對我是無能為力的」,耶穌為甚麼不反抗呢?耶穌如此回答:

「父怎樣吩咐我,我就怎樣行」(v31)耶穌離開,是尊重天父的安排,聽從天父的吩咐,完成天父交托給他拯救的使命.離開,並不是被甚麼神秘力量操控,也不是被死神打敗的結果.

很多時候,離開,都是人們測不透,想不通的.有些時候,離開,甚至令人大惑不解,令人苦思不得其法.為甚麼離開呢?為甚麼在這個時間離開呢?信仰的先賢先哲也苦思這些現象,在傳道書中留下這段文字.

「凡事都有定期,天下萬務都有定時．生有時,死有時．栽種有時,拔出所栽種的,也有時．殺戮有時,醫治有時．拆毀有時,建造有時．哭有時,笑有時．哀慟有時,跳舞有時．拋擲石頭有時,堆聚石頭有時．懷抱有時,不懷抱有時．尋找有時,失落有時．保守有時,捨棄有時．撕裂有時,縫補有時．靜默有時,言語有時．喜愛有時,恨惡有時．爭戰有時,和好有時.」(傳道書3:1-8節)

最後,作者有如此的領悟:

「上帝造萬物,各按其時成為美好」(傳道書3:11節)

黃福榮先生結束自己的生意,在身患長期疾病的時候離開香港,但他並不是被疾病打敗,而是獨自從香港步行到北京為中華骨髓庫籌款.黃福榮先生在青海玉樹地震中離開玉樹,並他並不是被死神擊倒,而是因為自願折返現場救出幾名仍被困的師生.離開,不是被擊敗,而是另有任務,另有意義,另有安排.

離開,是另一種形式的同在耶穌的對門徒所說的臨別贈言,一方面充滿著離別的意味,但同時卻又充滿著同在的應許.

「我要求父,父就另外賜給你們一位保惠師,叫他永遠與你們同在」(v16)

「我不撇下你們為孤兒,我必到你們這裡來」(v18)

「到那日,你們就知道我在父裡面,你們在我裡面,我也在你們裡面.」(v20)

「但保惠師,就是父因我的名所要差來的聖靈,他要將一切的事指教你們,並且要叫你們想起我對你們所說的一切話.」(v26)

耶穌的確要離開門徒,但耶穌指出,他的離開,並不是永遠的離開,而是另外一種形式的同在.耶穌離開之後,聖靈降臨,突破時間和空間的限制,永遠與門徒同在.聖靈的同在,不但沒有把門徒與耶穌分隔,反而令門徒與耶穌更親密,因為聖靈要叫門徒想起耶穌對他們所說的一切話.

是甚麼能夠令門徒與離開的耶穌以另一種形式繼續同在呢?這正是猶太在22節提出的問題:

「主啊,為甚麼要向我們顯現,不向世人顯現呢?」耶穌從另外一個角度回答這個問題,很值得細味:

「人若愛我,就必遵守我的道；我父也必愛他,並且我們要到他那裡去,與他同住.」耶穌離開,但誰能夠永遠與他同在呢?這不是耶穌要向誰顯然而不向誰顯現的問題,而是人們愛不愛他的問題.我們可否如此推想:沒有愛,離開就是永別；有愛,離開是另一種形式的同在.

年近八十的教育界及民主運動前輩使徒華(華叔),他對香港及中國民主運動的投入廣為人知,但原來他的愛情故事也非常感人.三月份的一集《星期二檔案》介紹,華叔年輕時曾經拍過拖,但女朋友婚前因病離世.華叔遵守與女朋友的諾言,受洗成為基督徒,以及一生都沒有結婚.在訪問中,華叔如此表白:「假如你心裡有一位你深深愛的人,也有一個人曾經深深愛過你,即使現在大家天人相隔．．．但這份感情在自己的心裡,都使自己感覺到幸福.」

有愛,離開是另一種形式的同在.

離開,祇留下愛耶穌向門徒解釋自己的離開時,澄清自己並不是被世界打敗,指出

「世界的王在我裡面亳無所有」(v30).那麼,何以耶穌甘願束手就擒呢?耶穌接著如此說:

「但要叫世人知道我愛父」(v31).耶穌離開,留下甚麼榜樣給門徒呢?留下甚麼形象給世人呢?英雄?義士?哲人?先知?革命家?宗教家?在這裡,耶穌希望門徒及世人認識的,是他愛天父.耶穌對門徒的最後囑咐,也就是對他這種愛的承傳.

在玉樹孤兒院做義工的黃福榮離開了,他留下甚麼給世界呢?

黃福榮家人給記者的一封信中這樣寫道:「請大家不要把福榮變成一個烈士或一個完美的人」.或者,這段網友的留言與黃福榮家人的心意比較接近:「福哥你真是人間的天使,你已完成任務離開這裡,你的慈悲教識每個香港人甚麼是愛,你的精神會繼續下去,祝你在另一世界會更快樂.」另有一段內地網友的留言如此說:「黃哥,你用生命愛汶川,愛玉樹,愛大陸.雖然你走了,但你卻活著,你永遠活在人民的心中.」

人有悲歡離合,月有陰晴圓缺.離合散聚,是人生的常態.與幼稚園同學相聚又離別,與小學同學相聚又離別,與中學同學相聚又離別,與大學同學相聚又離別,進入某個工作場所又離開某個工作場所,進入某些關係又離開某些關係,來到世界又離開世界,我們留下了甚麼給別人呢?別人又留下了甚麼給我們呢?耶穌離開,留下愛和愛的榜樣給門徒及世界,我們能否以此來疏理自己的經歷,以及調教自己的目標呢?若果有一天,當我們回想離別的經歷,情緒起伏過後,複雜的心情沈澱後,心裡只留下愛；如果有一天,別人想到我們的離別,只想起愛,那會是多麼美好,多麼完滿的人生呢!



\newpage

\section{我沒有醉}
\label{sec:50}
\textbf{鄧瑞強博士 - 「我沒有醉」}
\newline
\newline
\hyperref[sec:49]{< < < PREV SERMON < < <}
~
\hyperref[sec:toc]{[返目錄]}
~
\hyperref[sec:51]{> > > NEXT SERMON > > >}


使徒行傳 2:1-21
\newline
\begin{longtable}{cl}
\hline
\hline
章節 & 經文 (和合本修訂版)\\
\hline
2:1 & \begin{tabularx}{0.7\textwidth}{X} 五旬節那日到了，他們全都聚集在一起。 \end{tabularx} \\ \\ \relax
2:2 & \begin{tabularx}{0.7\textwidth}{X} 忽然，有響聲從天上下來，好像一陣大風吹過，充滿了他們所坐的整座屋子； \end{tabularx} \\ \\ \relax
2:3 & \begin{tabularx}{0.7\textwidth}{X} 又有舌頭如火焰向他們顯現，分開落在他們每個人身上。 \end{tabularx} \\ \\ \relax
2:4 & \begin{tabularx}{0.7\textwidth}{X} 他們都被聖靈充滿，就按著聖靈所賜的口才說起別國的話來。 \end{tabularx} \\ \\ \relax
2:5 & \begin{tabularx}{0.7\textwidth}{X} 那時，有從天下各國來的虔誠的猶太人，住在耶路撒冷。 \end{tabularx} \\ \\ \relax
2:6 & \begin{tabularx}{0.7\textwidth}{X} 這聲音一響，許多人都來聚集，各人因為聽見門徒用他們各自的鄉談說話，就甚納悶， \end{tabularx} \\ \\ \relax
2:7 & \begin{tabularx}{0.7\textwidth}{X} 都詫異驚奇說：「看哪，這些說話的不都是加利利人嗎？ \end{tabularx} \\ \\ \relax
2:8 & \begin{tabularx}{0.7\textwidth}{X} 我們每個人怎麼聽見他們說我們生來所用的鄉談呢？ \end{tabularx} \\ \\ \relax
2:9 & \begin{tabularx}{0.7\textwidth}{X} 我們帕提亞人、瑪代人、以攔人，和住在美索不達米亞、猶太、加帕多家、本都、亞細亞、 \end{tabularx} \\ \\ \relax
2:10 & \begin{tabularx}{0.7\textwidth}{X} 弗呂家、旁非利亞、埃及的人，並靠近古利奈的利比亞一帶地方的人，僑居的羅馬人， \end{tabularx} \\ \\ \relax
2:11 & \begin{tabularx}{0.7\textwidth}{X} 包括猶太人和皈依猶太教的人， 克里特人和阿拉伯人，都聽見他們用我們的鄉談講論神的大作為。」 \end{tabularx} \\ \\ \relax
2:12 & \begin{tabularx}{0.7\textwidth}{X} 眾人就都驚奇困惑，彼此說：「這是甚麼意思呢？」 \end{tabularx} \\ \\ \relax
2:13 & \begin{tabularx}{0.7\textwidth}{X} 還有人譏誚，說：「他們是灌滿了新酒吧！」 \end{tabularx} \\ \\ \relax
2:14 & \begin{tabularx}{0.7\textwidth}{X} 彼得和十一個使徒站起來，他就高聲向眾人說：「猶太人和所有住在耶路撒冷的人哪，這件事你們要知道，要側耳聽我的話。 \end{tabularx} \\ \\ \relax
2:15 & \begin{tabularx}{0.7\textwidth}{X} 這些人並不像你們所想的喝醉了，因為現在才早晨九點鐘。 \end{tabularx} \\ \\ \relax
2:16 & \begin{tabularx}{0.7\textwidth}{X} 這正是藉著先知約珥所說的： \end{tabularx} \\ \\ \relax
2:17 & \begin{tabularx}{0.7\textwidth}{X} 『神說：在末後的日子，我要將我的靈澆灌凡血肉之軀的。你們的兒女要說預言；你們的少年要見異象；你們的老人要做異夢。 \end{tabularx} \\ \\ \relax
2:18 & \begin{tabularx}{0.7\textwidth}{X} 在那些日子，我要把我的靈澆灌，甚至給我的僕人和婢女，他們要說預言。 \end{tabularx} \\ \\ \relax
2:19 & \begin{tabularx}{0.7\textwidth}{X} 在天上，我要顯出奇事，在地下，我要顯出神蹟，有血，有火，有煙霧。 \end{tabularx} \\ \\ \relax
2:20 & \begin{tabularx}{0.7\textwidth}{X} 太陽要變為黑暗，月亮要變為血，這都在主大而光榮的日子未到以前。 \end{tabularx} \\ \\ \relax
2:21 & \begin{tabularx}{0.7\textwidth}{X} 那時，凡求告主名的都必得救。』 \end{tabularx} \\ \\
[1ex]
\hline
\hline
\end{longtable}
日期:2010年6月13日

主題:「我沒有醉」

經文:使徒行傳2:1-21

証道:鄧瑞強博士

各位弟兄姊妹,主內平安.

在教會的聖曆裡,今日是「五旬節」,亦即「聖靈降臨節」,在舊約的世界,這是慶祝收割的日子,也是向神獻上初熟的五穀的日子(參出23:16；民28:26).辛勤的勞碌,終於得到回報.這是開心的日子,這是感恩的日子,這是分享的日子.大家應留意到,講台的禮儀用色改為紅色,顯出今日的喜慶色調.

今日講道的經文是:使徒行傳2:1-21.

這是一段關於五旬節的故事.這是主耶穌復活後第一個五旬節,主耶穌救贖的辛勞,在這日慶祝其初熟的成果.在使徒行傳的記載裡,這日發生了一件大事,主耶穌的救贖帶來的果效,在這日被世人看見了.

在喜悅中飽嘗基督救贖果子的信徒,在這日有異乎尋常的表現.世人以為他們醉了.一個過份興奮的人,有時看起來真的有點醉.這種「醉」,是一種「陶醉」,他們「陶醉」在一個凡人看不到的世界裡.「醉」的人總常說自己「無醉」,或許,他們真的無醉,醉的是凡人,凡人的「醉」,是一種「沉醉」,凡人「沉醉」在世俗的生活裡.單單講錢的人,當然看不見一塊菜田的寶貴.在他們看來,那些守著這塊菜田不肯拿去發展的人,是傻的,是醉的.但對那些「陶醉」在田園生活的人而言,他們認為自己是清醒的,其他人才是醉了.

我想過,在講這篇道之前,飲一杯酒,帶點醉意,然後走上講台第一句話便說:「我沒有醉」.事實上,相信理想的人,相信價值的人,講的話,總帶點醉意,總帶點不正常.信耶穌的人,應相信「施比受更有福」,應相信「生命的內容比成功更重要」,應相信「這可見的世界之外有另一個世界」,若如此,則說話應帶點醉意.這種醉,只因為你看到不一樣的世界.

「我沒有醉」,就是今日講道的主題.

(1)從聖靈而來的新生命

2:1　五旬節到了,門徒都聚集在一處.2:2　忽然,從天上有響聲下來,好像一陣大風吹過,充滿了他們所坐的屋子,2:3　又有舌頭如火焰顯現出來,分開落在他們各人頭上.

那日是五旬節,主耶穌的門徒聚在一起.

「忽然」,這是意料之外的事.有非常的事發生.

「天上有響聲下來,好像一陣大風吹過」.「風」是一種能力,一種生機.這是一種偉大的生命力,從天而來,充滿信徒身處的地方.

真正的大風是會吹動人的,這裡的風,是能推動人的大風.

又有火舌顯現,落在各人頭上.火,代表神的審判和能力.火和風,在舊約很多時都是代表神的出現.這裡提到火的形狀像舌頭,看來,神的出現要表現在人的說話上.

這個在主耶穌復活後的第一個五旬節,要顯出耶穌救贖事工結出的果實.這果實就是神同在所帶來的生機.

人在罪中死了.人在自私的慾望中失去了喜樂,人在自我的迷戀中失去了方向.世俗的麻醉,像香港的污濁的空氣一樣,使人看不到前景.我們需要一些清新的大風,吹走這些迷霧一樣的空氣.我們需要從天而來的風.

文化學者Neil Postman 寫了一本書,叫 ``Amusing Ourselves to Death"(《娛樂至死》),講到電視為我們塑造了一個講求娛樂而不善思考的文化.電視的一代,將什麼也娛樂化.賬災需要娛樂,沒有娛樂則籌不到款.批判時事需要娛樂,沒有娛樂化的情節,則根本不去理會時事.我們喜歡綽頭,不喜歡道理.我們需要一些審判的火,讓人眼目明亮,明辨是非,驅走過份的娛樂.我們需要從天而來的火.

耶穌的救贖,將我們從沉醉中救出來.從天而來的風和火,叫我們清醒地生活,豐盛地生活,自由地生活.

(2)聖靈下的感通

2:4　他們就都被聖靈充滿,按著聖靈所賜的口才說起別國的話來.2:5　那時,有虔誠的猶太人從天下各國來,住在耶路撒冷.2:6　這聲音一響,眾人都來聚集,各人聽見門徒用眾人的鄉談說話,就甚納悶；2:7　都驚訝希奇說:「看哪,這說話的不都是加利利人嗎?2:8　我們各人,怎麼聽見他們說我們生來所用的鄉談呢?2:9　我們帕提亞人,米底亞人,以攔人,和住在米索不達米亞,猶太,加帕多家,本都,亞細亞,2:10　弗呂家,旁非利亞,埃及的人,並靠近古利奈的利比亞一帶地方的人,從羅馬來的客旅中,或是猶太人,或是進猶太教的人,2:11　克里特和阿拉伯人,都聽見他們用我們的鄉談,講說神的大作為.」2:12　眾人就都驚訝猜疑,彼此說:「這是甚麼意思呢?」2:13　還有人譏誚說:他們無非是新酒灌滿了.

基督救贖裡的新生命,帶來一特別現象,就是在聖靈的充滿裡,門徒講起別國的話來.

五旬節是猶太人朝聖的日子,很多人在耶路撒冷聚集.這些人來自五湖四海.

來自東面的有「帕提亞人,米底亞人,以攔人」,這些人住在現在伊朗地區.有「住在米索不達米亞」的人,這些人住在現在伊拉克地區.有「猶太」人,這些人住在現今巴勒斯坦.

來自北面的有「加帕多家,本都,亞細亞,弗呂家,旁非利亞」的人,這些人住在現今的土耳其.

來自南部的有「埃及的人,並靠近古利奈的利比亞一帶地方的人」,這些是北非的人.

來自西部的有「從羅馬來的客旅中及克里特人」.還有南部的阿拉伯人.

這些人來自五湖四海,但都聽到門徒用他們的家鄉話,講說神的大作為.這到底是什麼一回事呢?

在舊約,一件重大的關乎人類語言的事件是「巴別塔」事件.

所謂「巴別塔」事件,就是當時的人以為自己了不起,想蓋一座可以登天的神殿.人要走上聖壇,扮演神的角色.人要宣揚自己的偉大.神看見人這種狂妄,於是變亂了人類的語言.人與人不能溝通,蓋通天塔的事便作罷了.過去,一般認為,神變亂他們的語言,是要懲罰他們,好叫他們不能合起來圖謀狂妄的事.現在,神學家認為,神變亂他們的語言,重點可能不是懲罰,而是保護人類,好讓人類語言和文化的多元性得以保存.

人世間多少次造神運動,多少狂人意圖一統天下,這都是近代的「巴別塔」事件.在強力的,政治的巴別塔之下,人不能發出自己的語言,而必須將自己要說的話統一在神化的獨裁者的說話之下.你要說自己的話,則必須要為此付上代價,一種封殺你的說話權而要你承受的代價.巴別塔是獨裁的,滅絕人性的,蒙蔽真相的.神看見了,神要保護人性尊嚴,神要保護人說話的權利,神變亂他們的語言.神保護了人講不同說話的權利.

如今,在這五旬節,聖靈降臨.神沒有統一人類的語言.五湖四海的人,有五湖四海的語言.但是,聖靈降臨,卻讓每個人都能聽懂別人的說話.我太太說,懂得聆聽便是最大的溝通.在這裡,人的語言縱使不同,但無損溝通和明白.我們不必統一語言才能溝通,只需要在神的大能裡,人縱使講不同語言,仍是能溝通的.

巴別塔事件,是人間的專制事件,表面語言統一,實則不能溝通.

五旬節事件,是人間嘗試彼此了解的事件,表面語言不通,但在聖靈裡,人和人能溝通,能明白神的作為.

這些人說:他們醉了.

(3)我沒有醉

2:14　彼得和十一個使徒站起,高聲說:「猶太人和一切住在耶路撒冷的人哪,這件事你們當知道,也當側耳聽我的話.2:15　你們想這些人是醉了；其實不是醉了,因為時候剛到巳初.2:16　這正是先知約珥所說的:2:17　神說:在末後的日子,我要將我的靈澆灌凡有血氣的.你們的兒女要說預言；你們的少年人要見異象；老年人要做異夢.2:18　在那些日子,我要將我的靈澆灌我的僕人和使女,他們就要說預言.2:19　在天上,我要顯出奇事；在地下,我要顯出神蹟；有血,有火,有煙霧.2:20　日頭要變為黑暗,月亮要變為血；這都在主大而明顯的日子未到以前.2:21　到那時候,凡求告主名的,就必得救.彼得說:我沒有醉.

在這新時代裡,兒女要說預言,少年人要見異象,老年人要做異夢,僕人和使女也有權講出他們心底的渴望.在基督的救贖下,在聖靈的充滿裡,每個人都可以講出自己的預言,講出自己的意願,講出自己的理想,講出自己的心聲.

在基督的救贖下,在聖靈的充滿裡,每個人都能明白對方的渴望,都能努力實現對方正直的夢想.

這是喝醉了嗎?我沒有醉.

我們年少的兒女們有什麼渴望呢?香港小童群益會(The Boys"  and amp; Girls" Clubs Association of Hong Kong)有一個「香港小特首」計劃.他們從小學生中選出一「小特首」,讓他從小學生的角度去發表一「小小施政報告」.我特別喜歡其2008年的「小小施政報告」,這裡講出了小學生們的心聲.這份施政報告提到:

「父母不應以錢代替關心,不應溺愛子女,給他們太多錢,這會令子女們建立錯誤的金錢觀.學校要進行課程改革,應多一點體育課,以增強體；多一點校外活動,以增廣見識.大人要學懂尊重小孩子的想法,不要只要求小孩子滿足大人的期望.要建設香港成為一「兒童友好」城巿,要設立地區兒童議會,讓小朋友發聲.......」這些小學生醉了嗎?他們沒有醉.

我們的「八十後青年」有什麼渴望呢?他們的渴望很懸殊.但有些青年,他們大聲呼喚社會公義,呼喚我們不要以賺錢能力去衡量他們,呼喚我們尊重不同的人的生活方式,呼喚我們不要只想著發展而要想想保育.他們欣賞街頭的賣藝者,這些賣藝者雖然可能在社會的底層掙扎,卻活得有色彩,他們保持著一種傳統文化,維持著本土經濟,並保存著人和人之間的基本互動.我們的「八十後青年」,呼喚著打破中環價值的生活形式,呼喚著自由.這些青年人醉了嗎?他們沒有醉.

老年人有什麼渴望呢?他們渴望社會不視「老年」為一種累贅而是一種寶貴的「資產」.渴望電視廣告不是只強調將老年人經化妝後變成美少女,而是學曉欣賞老人家的本身的美.渴望上公車時,人人主動讓座.渴望當在商場緩慢地前行時,其他人能慢下來而不會碰到他們.渴望社會不是時時趕著去淘汰他們,而是重視他們.渴望當政府提出增加老人「生果金」時,全民一致贊成.這些老年人醉了嗎?他們沒有醉.

我們沒有醉.我們只是渴望每個人都能被人性地對待.我們只是渴望這個社會更好地分配資源,讓人能有尊嚴地活下去.雖然語言不同,但我們渴望每個人都能聽懂別人的心聲,而不只是自己獨白.這是酒後的醉語嗎?沒有,我們沒有醉.這是基督救贖的世界,在聖靈的充滿下,這是我們每個人應有的渴望.

但願榮耀歸於聖父,聖子,聖靈.阿們.



\newpage

\section{施行拯救的三一上帝}
\label{sec:51}
\textbf{賴品超教授 - 施行拯救的三一上帝}
\newline
\newline
\hyperref[sec:50]{< < < PREV SERMON < < <}
~
\hyperref[sec:toc]{[返目錄]}
~
\hyperref[sec:52]{> > > NEXT SERMON > > >}


日期:2010年6月20日

主題:施行拯救的三一上帝

經文:羅馬書5

證道:賴品超教授

非常高興再次有機會跟大家來思想神上帝的話語,真巧合今天安排我在三一主日講道,其實有一些尷尬,因為三一主日講道不是那麼容易做的一件事.為什麼呢?對於很多的人來說,三一的教義在三一主日,通常都必定會講這位三一的上帝.但這是不容易講的,因為在一方面固然是重要,因為人們說基督教信的上帝,和其他宗教信的上帝有什麼不同?通常最容易講的,我們信的上帝是三位一體的上帝,與一般的一神論的宗教,無論是猶太教,伊斯蘭教是不同的.事實上,在三一信仰,在我們崇拜當中,我們非常清楚地體驗,剛才我們所諗誦的使徒信經裡,很清楚的講到聖父,聖子,聖靈.甚至在我們很多的禮文裡面,我們頌讚的都是聖父,聖子,聖靈,三一的上帝.正如有一位天主教的神學家Karl Rahner,他說他也有一點尷尬地承認,事實上在我們信徒的生活當中,基本上與這位三一的上帝,是毫無關係的,反而我們只信的是一位上帝,獨一的上帝(註:教父說三位一體是 Three Hypostases in One Ousia -- 三個位格於一個神性本質之內).

在三一主日中,通常我們一定要講三一的信仰,但是很多時候會出現很多理性上的問題,有些人覺得,不知道基督教的人在說什麼?因為有時在講他們的上帝,好像有三個,有時有好像在講一個.到底在講一個,還是在講三個呢?好像很難用邏輯去解釋當中有什麼關係?因為我自己是教神學的,曾經我也教過學生,在早期教會裡的三一教義裡講的是什麼?首先我必定會對他們講一些很簡單,又不是很容易理解的那些拉丁文,或是希臘文,什麼是substantia(本質)?什麼是Persona(拉丁文:位格)?什麼是Prosopon(希臘文 : 位格),什麼是Hypostasis(拉丁文:指神的三個位格或屬性)?在這些觀念之間有什麼關係?什麼叫做實體?原來實體裡面,在希臘文有不同的字,在拉丁文有不同的字,不是完全對等的.以至於早期的神學家,在爭論到底是那一個字,如何的去解釋.什麼是三位?這個位是什麼意思呢?拉丁文和希臘文各有一個不同的用字,裡面有不同的意思,其實是很難解釋的.我記得在我學生時期,在College Assembly崇基學生週會時,有一位物理學教授,對那些中三之後,從來沒有修讀物理的同學,大談什麼是相對論.他第一件事說:以下所講的,如果聽得懂的,是大家聰明,聽不懂的,是小弟丟臉,跟各位的智力無關,所以大家不用感到自卑,我大概明白他為何這樣說,因為我是其中一個不明白的.所以我今天也不會嘗試跟大家講以上那些拉丁文,希臘文是在講什麼?因為那些只會讓各位覺得,我是在拋書包,也在丟自己的臉,我最後也可能在講大家不明白,那幾個觀念是什麼?

歷代的教會都嘗試用很多不同的方法,用很多不同的比喻,有一點像主耶穌講天國的比喻,用很多不同的比喻,來講這三位一體的信仰是什麼?有些人說:三個人共同分享相同的人性.有些人會講,好像在人的裡面,在心靈裡有知,情,義,甚至我聽有人說,就如港式奶茶裡的三合一,有茶,有糖,有奶,或是有水,諸如似類的來解釋什麼是三位一體的上帝.但我只能告訴大家,任何這類的比喻,引喻,最多只能夠講明三位一體裡面,我們所認識的上帝的某一方面.而在另一方面,這些比喻,總是會有它的限制,不能盡說清楚.反而我自己覺得重要的是,其實這些比喻,不同的哲學觀念去講什麼叫做三位?什麼叫做一體?往往有時候會掉進一個陷阱當中,跟我們的日常生活,好像沒有什麼關係,跟我們的信仰實踐越來越疏遠,所以有時候在三一主日講道是有一些困難,特別每年都要講一次的時候,都不是很多東西可以說的時候.所以我相信伍牧師刻意叫我來講,不用他自己來講.

以下的時間,我不想跟大家再講什麼哲學觀念或甚麼,也不想提一些新的比喻讓大家思考,反而我想讓大家來去想清楚,為何在基督教裡,我們要宣認這三一的上帝?正如我剛才引用那位神學家Karl Rahner所講的,三一上帝這一個信仰,其實是上帝給我們的恩典,給我們的拯救,是直接有關的.我們為何要宣認一位好像很難明白的三一上帝?為何我們要在講道裡講?我們的講道,不錯是要讓會眾得造就,但也同要重要的是,要將榮耀歸於上帝,而二者是不能分割的.因為只有我們將榮耀歸給上帝的時候,我們的會眾才能夠得著造就.三一的上帝,是我們認信的對象.但是最重要的,是關乎我們所領受的恩典,也是我們信徒生活的依據.為何這麼說呢?

如果我們考據一下聖靈,我們的聖靈是很有智慧的,三一主日通常一定是在復活節後,聖靈降臨節之後,才到三一主日.為何呢?因為是我們一定經歷了聖子的受死,受難,復活,升天等,也經歷聖靈降臨之後,我們才有一個更加清楚的確據,來講這個三一的信仰,為何這樣的呢?我們看看今天幾段的經文,詩篇第八篇裡面所講的,上帝的創造和對人的顧念,怎樣將人放在伊甸園裡,使他管理天空的飛鳥,走獸.讓人想起上帝在創世記第一章裡所說的,上帝是照自己的形象樣式造人──造男造女,這是講上帝的創造,和對人的顧念,並且願意給人有一種獨特的地位,獨特的榮耀,祂在創造當中,是非常獨特的一個部分.甚至將上帝創造其他事物,叫或說就已經造成了.唯有當造人的時候,是照著祂的形象樣式造人.在箴言書八章講到上帝的智慧,講到祂向人說話,講到祂從太初就有,講到祂不是上帝受造之物,而是上帝藉著祂創造其他的生物和事物.這裡更加的提到,上帝對人的照顧,給人住在伊甸園,讓智慧是可以住在人間裡,人群的當中.同樣的說到一件事,上帝樂意將祂的智慧分享給人,讓智慧來呼籲人,來聽祂智慧的言語,同樣表達一個信息,上帝是願意把祂一些獨特的東西,特別是祂的智慧,來跟人分享,希望人能得著上帝這樣的智慧.

當我們進入新約聖經,我們今天要看的約翰福音十六章12-15節,講到當時主耶穌要告別門徒,提到一件事,聖靈要引導門徒進入一切的真理,或是明白一切的真理.這真理也正是主耶穌所講的,只不過人們當時還沒有能力去接受主耶穌的話語,而主耶穌也快要離開了,去受死,復活,升天等.因此,聖靈就有一個責任使命,來接續主耶穌的工作,但祂跟主耶穌的工作有所不同,主耶穌是有形有體,上帝的聖道成為人,有形有體捉得到,摸得到,甚至聞得到的形式向人啟示上帝,讓那不可見的上帝成為可以見,是一個上帝的形象.但是聖靈來,祂作的方法不一樣,祂是進入.主耶穌是在人群當中來生活,among us成為我們當中的一群一員.但聖靈是進入不同的信徒心裡in us,來引導我們的生活.好像在以弗所書所講的,上帝是超乎在眾人之上,貫乎在眾人之中,也住在眾人之內的上帝.因此,聖靈引導我們進入真理,並不是主耶穌所要講的真理之外的事,只不過當時主耶穌沒有時間講完的事,而讓聖靈來接續.祂們所作的工,不是分開的,不是間斷的,不是主耶穌作一件,聖靈作一件,因為聖靈是要將父所受於我的,告訴你們.也就是聖靈啟示給你們的,同樣也是主耶穌要啟示給我們的,也同樣是聖父要啟示給我們的.三者是環環相扣的.所以在聖經裡,有時在我們教會禮文裡,我們的聖父,聖子,聖靈是同尊同榮,而同樣的聖靈是要榮耀這位聖子.

我們今天要看的羅馬書第五章的經文,更加清楚的提醒我們,我們對拯救要有一個全面的理解.在羅馬書第五章1節-5節:我們既然因信稱義,就藉著我們的主耶穌基督,得與神相和.我們又藉著祂,因信進入現在所站的這恩典中,並且歡歡喜喜盼望神的榮耀.不但如此,就是在患難中,也是歡歡喜喜的,因為知道患難生忍耐,忍耐生老練,老練生盼望,盼望不至於羞恥.因為所賜給我們的聖靈,將神的愛澆灌在我們心裡.不知道大家在看這段經文的時候,有什麼感覺?很多時候我們在教會裡,特別在佈道會,很喜歡講因信稱義,你信耶穌就得拯救.好像我們的拯救,就是在我們信耶穌的那一刻.我們的罪得著赦免,我們可以在天堂占有一席位.將來死後,可以進入天堂.我們覺得那拯救就是這樣子.但我們看這段經文,原來不只是這樣.為什麼呢?因為經文告訴我們,已經因信稱義之後,不錯我們已經與上帝復合,藉著耶穌基督的工作.但是我們仍然在恩典中,歡歡喜喜的盼望上帝的榮耀.也就是說,上帝給我們恩典,不只是這些,上帝的恩典也包括了所賜給我們的聖靈,將上帝的愛澆灌在我們心裡,讓我們有盼望,有忍耐,有老練.換言之,上帝給我們的拯救,給我們的救恩,不單是免除了上帝的忿怒,消極的來說,消除了罪的刑罰.更加積極的是,上帝藉著聖靈將祂的愛澆灌在我們心裡,讓我們分享到上帝的愛和上帝的生命.叫我們在患難中,我們可以有老練,有忍耐,有盼望,可以歡歡喜喜的盼望上帝的榮耀.甚至我們再讀第五章後面的經文,我們更可以藉著祂以神為樂,也就是我們可以在地上與上帝同行.

我們試想一想,如果上帝的拯救,只不過讓人的罪可到赦免,而不需要受到懲罰,那這是不是一個完備的救恩呢?我不知道大家有沒有遇過不信的朋友,跟大家有這樣的質疑.他覺得基督教有時所講的事情,他覺得不適合他聽,為何呢?因為他覺得不是很公道,有很多人做好人一輩子,不信耶穌就要下地獄；反而有人臨死之前,不知道為何,突然想到我需要悔改.好像黃飛鴻的電影裡,石堅臨死之前悔改一樣.悔改之後,他就可以升天堂了,就便宜了他,好像上帝很不公道.最重要的,你會問,是不是信了耶穌之後,你不用在生活上有任何的改變呢?

最近我看到一個奶粉的廣告,我認為這廣告有不恰當的行為,需要家長指引.為何這樣說呢?因為那奶粉廣告說,有些小朋友有偏食的習慣,有偏食的習慣,使他營養不均衡,但我們要買這牌子的奶粉,可以補足他不夠的營養,表面上這牌子的奶粉很好,因為平常有偏食的習慣,不足夠的營養,你食這牌子的奶粉,就可以補充欠缺的營養.但是他沒解決他偏食的問題,鼓勵他繼續的偏食,是否要他吃這牌子的奶粉,直到他八,九十歲,百年歸老的時候呢?不錯現在也有老人的奶粉,給年紀大的人飲用的.是不是鼓勵他這樣子,還是你應該留意一下,他這個偏食習慣本身是需要被改變的呢?

最近有一間大學的中醫學院,研究出一個方法,使那些吸毒者解決某些問題,例如尿頻,政府的廣告說,十五分鐘就要上洗手間,所以吸毒就會使你身體不健康的,使你的生活受到很大的影響,每十五分鐘就要上洗手間,連搭長途車去遠一點的地方也不方便.最近這間大學的中醫學院,找到一個中醫的療法,使之可以改善這方面的問題.同時我在報紙看到,千萬不要亂說這樣的事情,為什麼?因為吃了這中藥後,就算吸毒後也不會對你的身體有問題,不會再有尿頻.這很容易給他們一個錯誤的信息,以為可以繼續的吸毒,沒什麼問題的,只要吸毒後,吃我們的中藥,就不會有尿頻的問題了.

上帝的拯救,如果是這樣,是否一個完備的拯救?如果上帝的拯救,上帝的恩典,只不過免除罪的刑罰,並沒有改變那個人,使那人有新的生命──是一個愛上帝,厭惡罪的生命,這不是一個完備的救恩.相反的羅馬書告訴我們,上帝給我們的,不只是免除上帝的忿怒,也賜給我們新的生命,上帝的靈,將上帝的愛澆灌在我們心裡,使我們有愛上帝的心,愛公義,愛和平,愛人群這樣的心,這才是一個完備的救恩.

我有一位認識了二十年以上,近這十幾年,他從事戒毒的工作.他看到一個戒毒的工作,如果你只是幫那人處理毒品對身體的傷害,在短期內戒掉毒癮,不再吸了,可以離開戒毒所.這是否幫到他呢?他發覺不是,因為那個吸毒的問題,他發現年青人,當你解決了他戒毒的問題,其實那只不過是一個表徵,他吸毒背後有許多問題,可能是家庭的破碎,可能他沒有一群好的朋友來支援他,他認識的都是帶壞他的朋友,也可能他對前途毫無希望,對人生毫無盼望的時候,他就用這吸毒,成為他的一個表徵,而不是真正的原因.如果你不能解決他真正的原因,可能是他家庭的問題,對前途的迷茫,對人生失去了意義,或是受到朋輩的影響,如果他的問題沒有得到解決的時候,只是短期內使他不吸毒,是沒有用的,他很快就會再次吸毒.之後,我朋有就嘗試,用一個新的方法: 原本在求學的,就給他機會讀書,等他有機會去考上好的學校,使他對前途再次有信心,因為當一個人對他自己的前途沒信心的時候,他就很容易尋求吸毒這東西來麻醉自己. 幫助他建立朋輩的社交圈子,帶他去教會,讓他認識正當的朋友. 幫助他處理他家庭的問題,使他和家人如何的復和. 他發覺原來這問題是很深層次的,如果你不去解決那些問題,只是解決表面吸毒的問題,很快他便會再去吸毒,不是K,可能是R或M.今天這段經文提醒我們,上帝為我們所預備的完備救恩,不單只講我們的罪得赦免,不需要下地獄,接受永遠的刑罰,不單只是因信稱義,這只不過是我們與上帝之間的關係,入門的初階,並非終極的目的.終極的目的是什麼呢?是領受上帝給我們的生命,仁愛,喜樂,和平,聖潔的生命,是一種厭惡罪惡的生命,這才是上帝的救恩,那終極的目的.有時候我們說,上帝給我們永生,但有時候我們只看到,永生是死後,人進入永恆,永遠的生命.

有一次我到大陸去講學,跟這題目沒有關係的,講的是大乘佛學的觀念,去解釋基督教的基督論.恰巧有一位學生,他問我一個問題,通常學生問的問題,跟你講的題目是沒什麼關係的,總之,他有什麼問題他就發問.他的問題是:為什麼我們要那麼關注「拯救」這個問題呢?這是死後的事.那次我所講的內容,是完全跟主題沒有任何關係的.但是我還是要回答:按照我的理解,基督教所講的拯救,不是在死後才得拯救,不是死後得永生,而是那永遠的生命,現在已經開始了,在你接受主耶穌的一刻,你已經開始跟上帝建立一個新的關係,領受一個新的生命,所以我告訴他,基督教的拯救,不是在你死後才開始,是你現在就可以開始得到的.換言之,其實基督教的拯救,最棒的部份,不是你死後的那一部份,而是你死前的那部份.原來我們今天有機會,可以在今生,很快樂的與上帝同行,以上帝為樂,飽嚐上帝的慈愛,正如在詩篇九十篇:願我們早早飽得上帝的慈愛.原來上帝為我們所預備的救恩,不單只是除掉罪的刑罰,而是給我們一個新的生命,這生命是現在開始可以經驗得到的,以至我們在地上可以預嚐天國的福樂,我們可以感受到上帝要給我們一個怎樣的生命,這才是創天造地,偉大無比的上帝,但也顧念卑微的人,願意將祂自己的形象樣式來給人的上帝,為人所預備的奇妙又完備的救恩.我們原本是塵土所造的,但是上帝卻是願意,按著祂的形象樣式造我們,並且在基督的拯救中,來恢復上帝的形象.

在二世紀,有一名很出名的神學家,叫做Irenaeus,他用一個觀念去解釋什麼是拯救?他用聖經裡,講到yiu 匠手裡的器皿弄壞了再重做,他說:上帝利用祂的左右手,什麼是左右手?就是聖子,聖靈,或是聖道和聖靈,來塑造人性,使這原本敗壞的,破壞的或是扭曲的,當然他是分得很清楚,什麼是形象?什麼是樣式?上帝的形象受到破壞之後,在基督裡面來重新恢復,在基督裡,從無形的上帝,有形的展現出來.將這形象展現出來是怎麼樣,而聖靈幫助我們,來重新獲得這形象,藉著祂在我們心裡面工作,引導我們來到聖子面前,由聖子引導我們來到上帝面前.簡單的說,上帝給我們拯救,要將那失落,扭曲的形象重新恢復.

在羅馬書這段經文所講的,患難生老練,老練生忍耐,忍耐生盼望等等.也讓我想到另一處經文,是在彼得後書一章,我們有了知識,敬虔,忍耐,愛心... 他說:原來上帝為我們所預備的,是關乎生命和敬虔的事,最終的目的,是要我們脫離世上的情慾,並且得與上帝的性情有份.如果我們理解這一點的時候,我們大概也可以理解到,為什麼我們會承認這位三一的上帝?因為上帝所預備給我們的拯救,是一個完備的拯救.祂不單藉著聖子,在十字架上的犧牲流血,也藉著祂在人間活出人的樣式,是作為一個人該如何來事奉上帝的樣式.也藉著祂的復活升天,差遣聖靈,由聖靈來引導信徒,進入一切的真理,最後目的是要我們完完全全的來分享上帝的聖潔,愛,和平的生命.這才是完備的救恩,不只是告訴我們罪的刑罰得以赦免,而是給我們新的生命,並且要長成基督長成的身量,這才是上帝的目的.

原來我們所承認這位三一的上帝,因為我們所經歷拯救的恩典,正是上帝有計劃的策劃,上帝從起初到終末,是同一個的計劃,在這計劃中,是需要由聖父,聖子,聖靈,扮演著不同的角色,來配合,來成就,目的是讓原本是由塵土所造,很卑微的人,也領受上帝形象的人,最後得到上帝完美的形象,可以分享到上帝的生命,這才是上帝給我們的拯救.是不可分割,整全的計劃,一個終極的目的,這目的不是停留於我們所講的因信稱義,停留於我們的罪得赦免,而是積極的分享上帝的生命.

有時我們在教會宣講中,在佈道會中,急著要人信耶穌,就用永刑,永火,地獄來嚇人,使他們快快的來信耶穌,但這只不過是上帝要拯救的初階,是一個入門.上帝更希望在我們的生命中,在我們現世的生命中,怎樣學效基督,在我們在世為人的時候,我們能夠活出基督形象的榮美,榮耀上帝,能夠分享到,聖靈給我們上帝的愛,聖靈所給我們的智慧,然而我們再一次來思想,藉著這三一主日,再次來思想上帝給我們的各樣恩典,和那完備的救恩,我們為此來感恩,更加求主使我們在世為人,順從聖靈的引導,活出基督的樣式,以至榮耀歸於天父.阿們.



\newpage

\section{漫漫朝聖路}
\label{sec:52}
\textbf{伍渭文牧師 - 漫漫朝聖路}
\newline
\newline
\hyperref[sec:51]{< < < PREV SERMON < < <}
~
\hyperref[sec:toc]{[返目錄]}
~
\hyperref[sec:53]{> > > NEXT SERMON > > >}


日期:2010年6月27日

主題:漫漫朝聖路

証道:伍渭文牧師

大家有沒有參加過巡行(parade)?大家有沒有看過巡行?雖然我們不一定參加過巡行,但我想在座每一過人都會看過巡行.香港旅遊發展局(前稱香港旅遊協會)於二零零一年開始,每年農曆年初一舉辦"繽紛巡遊賀新歲"花車巡遊.因為首次舉辦,故大事宣傳謂:巡行共有13部花車,12隊國際和26隊本地表演團體沿途獻技,上海電視台經由衛星頻道向內地超過1300個廣播站現場直播巡遊盛況.

同是行列,同是人頭湧湧,朝聖行列與巡行行列就很不同.(1) 巡行沒有目的地:巡行沒有一個要到達的目的地,巡行的過程就是目的,巡行的人從自我表演中找到滿足,所以全程要歡樂開心,充滿嘉年華氣氛.但朝聖者有一種強烈的召喚要前往目的地,朝聖者不在乎此刻的歡樂,也許過程是艱辛的,但想到最終追尋的目的地就心得鼓舞,繼續堅持下去.

(2) 朝聖者有很強的團結:巡行來得熱哄哄,參加巡行的人是一時結集,巡行完畢就各自散去.但朝聖有相同的召喚感,而且路途漫長,險阻重重,過程中彼此幫助,所以大家有較深的團結.西遊記記載玄裝與門徒到西域取經,面對諸般試探和危險,遇到流沙,火焰山,和諸般妖怪迷惑和誘惑,但為了取經,互相幫助,排除萬難.同樣,基督徒的生命也是一個朝聖旅程,荊棘滿途,容易氣餒和迷失,需要信心和毅力,也需要同行者的鼓勵和扶持,直抵天家.教會誠然是朝聖的群體,而朝聖的起始就是受洗歸入基督的名下的一刻. 一  洗禮是新旅程  羅6:3, 4對十位接受水禮弟兄姊妹來說,今天是一個朝聖旅程的開始.今天不是你們來到屬靈的高峰 (高峰是跟著下坡的),乃來到屬靈生命的高原,來到屬靈的新台階.羅馬書說:豈不知我們這悉受洗歸入基督耶穌的人,是受洗歸入他的死麼.所以叫我們一舉一動有新生的樣式.(羅6:3,4)  所以洗禮不是一次的決定,乃是一生的態度,每天的經歷.我們不會否認,洗禮的確是心志的的決定,早期教會要經過三年慕道才能領洗,除了用三年學習義理,也用三年時間觀察考驗他們的心志.在聖餐開始時他們便被請離開,Ite Missa est! 祇有受洗的人才能領受聖餐.慕道的人,若等不到洗禮便殉道,他們流出的血,就被視為給他們洗禮了.在文化大革命時,願意洗禮的,就是表示願意為主殉道.印度聖徒孫大信為自己洗禮寫的一首詩歌The world behind me(世界在背後),表明接受洗禮者的決心,歌詞說:The world behind me, the cross before me, no turning back, no turning back(世界在背後,十架在前頭,永不回首,永不回首).聖經也以過紅海,離開為奴的埃及地,永不回頭,作為水禮的比喻.

哥林多前書說:弟兄們,我不願意你們不曉得,我們的祖宗從前都在雲下,都從海中經過,都從雲裡,海裡受洗歸了摩西.(林前10: 1, 2)  但我們若留心細看,這段經文所強調的,不單是一次離開埃及過紅海的決心,乃強調在西乃山從摩西領受律法,一生不斷學習作上帝的子民.我們要留意這裡說:都從雲裡,海裡受洗歸了摩西.這是甚麼意思呢?是指離開埃及過紅海,到西乃山接受十誡,與上帝立約,作上帝的子民.

不錯,從人的角度看,洗禮是表明我們的信仰,我們經過在此一年聚會,領受教訓,並參與數月的慕道班,下了決心接受水禮,在人面前為主作見證.因為是決心,是信心一個臨界點,我們有時猶豫自己未夠資格洗禮,亦懼怕洗禮之後再犯罪,歷史上君士坦丁(Constantine) 就在臨終前在病榻領洗的,認為在樣就保證到領死後不再犯罪了.然而聖經說,我們的決心和信心也是上帝賜的.我們得救是本乎恩,也因著信,這並不是出於自己,乃是上帝所賜的,免得有人自誇.(弗2: 8, 9)

從上帝的角度看,水禮是上帝給我們的恩典.我們強調聖禮 and \#8212;看得見的道,和聽得到的道 and \#8212;講道,都是外在於我們,賜給我們的,不是我們內心想出來的.我們聽見或看見神的話語,誘發出信心,回應上帝的恩典,像兩手相握.耶穌在約但河洗禮時,天上有聲音說,這是我的愛子,我所喜悅的.同樣今天上主對領洗的人說:某某,你是我所揀選的,我曾使用某人在某時引導你,直到今天.

所以,洗禮不是一次的決定,是一生的態度,每天的經歷與主同埋葬同復活,一舉一動都有新生的樣式.所以今天對領洗的人來說,你們不是來到屬靈的高峰,乃來到屬靈的高原,洗禮是漫漫朝聖路的開始.

二  洗禮是進入上帝的家 彼前3: 20, 21巡行和朝聖不同,朝聖隊伍彼此有很深的團結.我們想像一下,一班從遠處前往耶路撒泠朝聖的駱駝隊香客,坐駱駝應是老弱,幼小和病患的,年壯的步行引路.日薄西山時,他們紮起營幕,圍在一起用膳,用餐完畢點起火來,唱歌跳舞,彼此鼓勵晚上好好休息,明天上路.我們受洗歸入基督名下的,也是歸入基督的身體,彼此互為肢體.洗禮是進入方舟,加入上帝的家.

當時進入方舟,藉著水得救的人不多,祇有八個人.這水所表明的洗禮-現在藉耶穌基督的復活,也拯救你們.這洗禮本不在乎除掉肉體的污穢,祇求在上帝制面前有無虧的良心.(彼前3: 20, 21)

很多古老的教會在進入教堂時都放置領洗盆,象徵藉洗禮進入上帝的家,而稱教會的大堂為Nave船,教會就是方舟,進入其中得到拯救,而得拯救的人成了一家人.又如歌羅西書2: 11, 12所說:洗禮是新的割禮,烙在我們身上的印記,同屬上帝的子民標記.你們在他裡面,也受了不是人手所行的割禮,乃是基督使你們脫去肉體情欲的割禮.(西2: 11, 12) 因此,我們憂戚與共,患難相扶,一個肢體受苦,全體都受苦,一個肢體得榮耀,全體都得榮耀,我們要關心上帝的國.

我盼望大家此後更多關心教會,支持教會,發揮恩賜,彼此服待,在愛中建立自己.教會是不是巡行的行列,每主日回到教會湊熱鬧,認識朋友,一起茶聚,一起午膳,一起活動,我們分開後還彼此記掛,互相幫助,一起在信仰上互勉,彼此代求.



\newpage

\section{你可以行多一步嗎}
\label{sec:53}
\textbf{賀志勇博士 - 你可以行多一步嗎?}
\newline
\newline
\hyperref[sec:52]{< < < PREV SERMON < < <}
~
\hyperref[sec:toc]{[返目錄]}
~
\hyperref[sec:54]{> > > NEXT SERMON > > >}


路加福音 7:36-50
\newline
\begin{longtable}{cl}
\hline
\hline
章節 & 經文 (和合本修訂版)\\
\hline
7:36 & \begin{tabularx}{0.7\textwidth}{X} 有一個法利賽人請耶穌和他吃飯，耶穌就到那法利賽人家裡去坐席。 \end{tabularx} \\ \\ \relax
7:37 & \begin{tabularx}{0.7\textwidth}{X} 那城裡有一個女人，是個罪人，知道耶穌在法利賽人家裡坐席，就拿著盛滿香膏的玉瓶， \end{tabularx} \\ \\ \relax
7:38 & \begin{tabularx}{0.7\textwidth}{X} 站在耶穌背後，挨著他的腳哭，眼淚滴濕了耶穌的腳，就用自己的頭髮擦乾，又用嘴連連親他的腳，把香膏抹上。 \end{tabularx} \\ \\ \relax
7:39 & \begin{tabularx}{0.7\textwidth}{X} 請耶穌的法利賽人看見這事，心裡說：「這人若是先知，一定知道摸他的是誰，是個怎樣的女人；她是個罪人哪！」 \end{tabularx} \\ \\ \relax
7:40 & \begin{tabularx}{0.7\textwidth}{X} 耶穌回應他說：「西門，我有話要對你說。」西門說：「老師，請說。」 \end{tabularx} \\ \\ \relax
7:41 & \begin{tabularx}{0.7\textwidth}{X} 耶穌說：「有兩個人欠了某一個債主的錢，一個欠五百個銀幣，一個欠五十個銀幣。 \end{tabularx} \\ \\ \relax
7:42 & \begin{tabularx}{0.7\textwidth}{X} 因為他們無力償還，債主就開恩赦免了他們兩個人的債。那麼，這兩個人哪一個更愛他呢？」 \end{tabularx} \\ \\ \relax
7:43 & \begin{tabularx}{0.7\textwidth}{X} 西門回答：「我想是那多得赦免的人。」耶穌對他說：「你的判斷不錯。」 \end{tabularx} \\ \\ \relax
7:44 & \begin{tabularx}{0.7\textwidth}{X} 於是他轉過來向著那女人，對西門說：「你看見這女人嗎？我進了你的家，你沒有給我水洗腳，但這女人用眼淚滴濕了我的腳，又用頭髮擦乾。 \end{tabularx} \\ \\ \relax
7:45 & \begin{tabularx}{0.7\textwidth}{X} 你沒有親我，但這女人從我進來就不住地親我的腳。 \end{tabularx} \\ \\ \relax
7:46 & \begin{tabularx}{0.7\textwidth}{X} 你沒有用油抹我的頭，但這女人用香膏抹我的腳。 \end{tabularx} \\ \\ \relax
7:47 & \begin{tabularx}{0.7\textwidth}{X} 所以我告訴你，她許多的罪都赦免了，因為她愛的多；而那少得赦免的，愛的就少。」 \end{tabularx} \\ \\ \relax
7:48 & \begin{tabularx}{0.7\textwidth}{X} 於是耶穌對那女人說：「你的罪都赦免了。」 \end{tabularx} \\ \\ \relax
7:49 & \begin{tabularx}{0.7\textwidth}{X} 同席的人心裡說：「這是甚麼人，竟赦免人的罪呢？」 \end{tabularx} \\ \\ \relax
7:50 & \begin{tabularx}{0.7\textwidth}{X} 耶穌對那女人說：「你的信救了你，平安地回去吧！」 \end{tabularx} \\ \\
[1ex]
\hline
\hline
\end{longtable}
日期:2010年7月4日

主題:你可以行多一步嗎?

經文:路加福音7:36-50

証道:賀志勇博士

有一個法利賽人,請耶穌和他吃飯．耶穌就到法利賽人家裡去坐席.那城裡有一個女人,是個罪人．知道耶穌在法利賽人家裡坐席,就拿著盛香膏的玉瓶,站在耶穌背後,挨著他的腳哭,眼淚濕了耶穌的腳,就用自己的頭髮擦乾,又用嘴連連親他的腳,把香膏抹上.請耶穌的法利賽人看見這事,心裡說,這人若是先知,必知道摸他的是誰,是個怎樣的女人,乃是個罪人.耶穌對他說,西門．我有句話要對你說．西門說,夫子,請說.耶穌說,一個債主,有兩個人欠他的債．一個欠五十兩銀子,一個欠五兩銀子．因為他們無力償還,債主就開恩免了他們兩個人的債.這兩個人那一個更愛他呢. 西門回答說,我想是那多得恩免的人.耶穌說,你斷的不錯.於是轉過來向著那女人,便對西門說,你看見這女人麼．我進了你的家,你沒有給我水洗腳．但這女人用眼淚濕了我的腳,用頭髮擦乾.你沒有與我親嘴,但這女人從我進來的時候,就不住的用嘴親我的腳.你沒有用油抹我的頭,但這女人用香膏抹我的腳.所以我告訴你,他許多的罪都赦免了．因為他的愛多．但那赦免少的,他的愛就少.於是對那女人說,你的罪赦免了.同席的人心裡說,這是什麼人,竟赦免人的罪呢.耶穌對那女人說,你的信救了你,平平安安的回去罷.(路加福音7:36-50)

前些日子聽說唐崇榮牧師來香港講道,中大有一些朋友就過去聽.回來之後跟我說唐牧師講得多好多好啊,我就問他們說唐牧師講了什麼啊."哎呀,我忘了."他拍拍腦袋說,不過他又說:"我記得他講了幾個故事." 呵呵!有時候的確是這樣,長篇大段的講道理人家很難記得住,可是幾個小故事,倒是讓人很難忘.所以耶穌就經常給我們說比喻,講故事.而不是拼命給我們背金句.下面這個故事,也是一個大家耳熟能詳的故事了--一個有罪的女人用玉瓶裝香膏膏抹耶穌的故事.當時,法利賽人西門邀請耶穌到家裡吃酒席,一個有罪的女人用眼淚洗耶穌的腳,用頭髮擦乾.還用香膏膏耶穌的腳.法利賽人西門不高興耶穌接近罪人,但耶穌對他說,赦免多的人愛也多,因為這女人愛耶穌多,所以這女人的罪就被赦免了.

相似的故事還出現在馬可福音14:3-9節,經文說在伯大尼,在西門家的宴會上,一個女人打破貴重的玉瓶,把裡面價值三十兩銀子的真哪達香膏,澆在耶穌頭上.三十兩銀子是什麼概念?有經文解釋,三十兩銀子大概是當時平均一個工人一年的工資.哇,你說用一年的工資買真哪達香膏,就好像我們現在用一年的工資來買一瓶法國香水,是不是太奢侈啦?這還沒算上那個可能更加貴重的玉瓶的價錢了.

那這段經文到底要講什麼呢?如果對照馬可福音,我們發現這段經文有幾個很獨特的地方:在路加福音中,耶穌是被法利賽人西門邀請到家裡來赴宴的,這就有點奇怪了.我們發現耶穌跟法利賽人關係並不好,耶穌經常在安息日治病救人,又掀翻聖殿買賣東西的錢櫃,還經常批評法利賽人假冒為善,是偽君子,(參馬太,馬可相關章節)所以法利賽人並不喜歡耶穌.但很明顯,這段經文裡的法利賽人西門還是相當尊重耶穌的,在古時候,邀請人家來家裡吃飯表示關係還很不錯.法利賽人相當有階級觀念,不會隨便邀請人到家裡去.其實不但法利賽人這樣子,今天北京的上海也都是這樣子.西門邀請耶穌去家裡坐席,說明他還算尊重耶穌；另外,他稱呼耶穌夫子,說明他也認為耶穌是一位拉比.在當時,拉比不但受人尊重,他還是教導人的,也就是說,某種意義上,西門還是接受耶穌的教導的,40節耶穌說:"西門,我有句話要對你說."西門說:"夫子,請說." 在這裡,夫子原文可以翻譯成老師,或者說是拉比,這說明西門在一定程度上是尊重耶穌的.但這一點就很奇怪,因為當很多法利賽人恨耶穌的時候,西門卻請耶穌到家裡來吃飯,他就不怕接近耶穌會惹禍,會遭人排擠麼?

一個有罪的女人,挑選一個吃飯的時候跑到西門家裡去,給耶穌又洗腳又膏抹,時機是不是有點不對呢?她總不至於希望西門也順帶請她吃飯吧?當時的這個場景很戲劇化,就好像某天有人請伍牧師吃飯,還請了其他幾個有身份的朋友一起作陪,結果跑來一個不正經的女人,又是給伍牧師洗腳,又是給伍牧師膏抹,這肯定會嚇到身邊的人.因為這樣的場景太匪夷所思了,一邊是牧師以及一些有身份的人,另一邊是個不正經的女人,一個罪人,這個罪人還做出這麼讓人吃驚的舉動,這個場面如果我在場,我都會嚇到.但耶穌卻沒有被嚇到,他還坦然接受這個女人的舉動,並且為這個女人的舉動辯護,甚至宣稱他赦免了這個女人的罪.耶穌這樣一個舉動,怎麼不會讓人議論紛紛了--畢竟,這個女人不太正經,是一個罪人.中國人說"瓜田李下",要注意影響嘛,可是耶穌一點都不在意,還坦然接受一個不正經的女人給他洗腳,給他膏抹.你說奇怪不奇怪?

這個膏抹的舉動很戲劇性,不僅僅是因為這個女人的身份,還有一點就是女人所用的東西很不尋常.首先,她不是用水,她是用她的淚水來給耶穌洗腳--西方人說,女人的淚水是珍珠,我們暫且不說這珍珠有多麼寶貴；可是你再想想,淚水能有多少水啊,這女人用眼淚來洗腳,她得要流多少眼淚啊?所以可想而知,當時當地,這個女人又是多麼動情!是什麼讓她這樣動情呢?其次,她居然打破珍貴的玉瓶,用價值一年工錢的貴重香膏來膏耶穌的腳,這樣做值得麼?大家請想像一下,你會用一年的薪水來買一大瓶法國香水,然後送給太太來沖涼麼?這一次還不是用一點點香水,這一次是倒光了整瓶法國香水--一年薪水買的法國香水,還不是用來洗澡,是用來洗腳!說到這裡,我們就可以體會,聖經描繪的這一幕是多麼富有戲劇性,多麼讓人匪夷所思,多麼讓人震驚!

更富戲劇性的是耶穌的舉動,如果這是個戲劇,耶穌的舉動就是整幕劇的最高潮.耶穌並沒有因為膏抹他的是個女人而不好意思,因為一般猶太男人是不直接面對自己親人之外的女人的；我去耶路撒冷的時候,女導遊跟我說,她去鄰居家裡借東西,鄰居家的猶太男人都面朝牆壁,背對著她,因為不想看到她的臉.

而且,耶穌更沒有因為這個女人是罪人就說,"不好意思啊,Miss,我現在在吃飯,你能不能跟我約個時間,我們再私下談談,傾一下解啦.或者你Email給我也行.現在實在太不方便啦."耶穌沒有!他不但坦然接受這個女人的舉動,而且誇這個女人愛她比西門更多,因此她的罪得赦免了.這不是很奇特麼?

愛是什麼?愛難道是在街上買菜,還可以稱斤論兩,還可以計算一下,誰的愛更多,誰的愛更少麼?難道天國也是賣門票的,你給的愛越多,你拿到入場券的希望就越大麼?耶穌明明說的是"被赦免的越多,愛也就越多,"但47節卻說"她許多的罪都赦免了,因為她的愛多."很明顯,耶穌說的前後不一致.他前面說的是先赦免,後有愛.現在卻又說因為這個女人愛多,所以她被赦免了.這又是為什麼?

在杭州的時候,我總在思索這個問題:這樣一個故事,實在是太富戲劇性了.

耶穌經常教育人要低調,可是在面對這個有罪的女人時,行事卻為什麼要這麼高調? 他完全可以換個時間去接見這個女人,但是他卻在這樣一個公眾場合,不顧身份去接受一個女人這樣方式的愛；而且他還說,愛他多的人得到的赦免也多,難道耶穌也像現在的戀人一樣,還要計較誰付出的多,誰付出的少麼?這些問題困擾著我,直到有一天,我想去西溪濕地玩,朋友說門票很貴,但是本地人進去就可以買半價票,他可以幫我借張別人的杭州市民卡,這樣我可以省點錢,冒充本地居民混進西溪濕地.

我當時很欣然接受了他的提議,因為以前在國內讀書的時候都是這樣做的嘛.晚上我正在準備東西的時候,我突然想:我在做什麼,賀志勇,難道為了省這麼一點小錢,我就把自己的原則出賣了麼?我口口聲聲稱那個女人是罪人,我還不是一樣,因為一點點小錢就做了罪人.那個時候我突然明白了,我之所以不太理解那段經文,是因為我帶著有色眼鏡在看它.

當我想耶穌為什麼要這麼高調地接受一個有罪的女人的時候,我其實跟法利賽人一樣,在帶著有色眼鏡在看待這個可憐的女人.我沒有看到她的眼淚,沒有看到她的傷痛,沒有看到她的悔改.一個女人,在這樣的環境下來到耶穌面前,在眾目睽睽下用眼淚來洗耶穌的腳,還用這麼昂貴的香膏來膏耶穌,是需要多麼大的勇氣!她的悔改又是何等的勇敢,何等的真切了!可是在我的有色眼鏡裡頭,在法利賽人的有色眼鏡裡頭,我們看不到這女人的傷痛,這女人的勇氣,這女人的悔改--這女人在悔改中對主莫大的愛.我們看不到!我們只看到一個曾經犯錯的罪人."我們看得到別人眼中的刺,卻看不到自己眼中的梁木."西門就是這樣子,他覺得自己很虔誠,很愛神,卻看不到自己其實跟這個女人一樣,欠神也是50兩.有時候,當我們自己覺得自己已經很不錯了,很愛神的時候,往往也是離神還很遠的時候.

這個世界還有很多這樣的有色眼鏡,比方說,種族的有色眼鏡--雖然馬丁路德金取得了勝利,但美國到現在還流行種族有色眼鏡；還有,家鄉地域的有色眼鏡--中國人常說,縣裡人看不起鄉下人,城裡人看不起縣裡人,大城市人看不起小城市人.我聽說,上海人喜歡把非上海人都看做是鄉下人,所以有個上海同學來香港讀書的時候,香港朋友問他 ``你鄉下是哪里啊?"他就傻了.

這世上這種有色眼鏡很多,外貌的,地位的,錢財的有色眼鏡--看人先看穿得怎麼樣,名片上印的是什麼職位,開的是什麼車子,住的是什麼地段,卻往往忘記了從神的眼鏡來看人,所以看不到人真正的價值.

但神從來不這樣看,神看我們每個人只比天使微小一點,都是他的兒女,都是他所愛的.他甚至為了愛我們,把他獨生愛子的生命都犧牲了.在他的眼裡,我們每個人都是這樣寶貴.所以,在耶穌當時的眼裡,他看到的不是一個罪人,他看到的是一隻曾經流離飄蕩的小羊,一個憂傷痛悔的婦人,一個為了自己的罪而深深悔過的婦人."憂傷痛悔的心,神必不輕看."透過婦人撕心裂肺的哭泣,他看到的是那婦人雖然處處被人拒絕,卻頑強悔改的心.--神是我們這些罪人的朋友,所以他對罪人的愛也特別的深!

只有透過這點我才明白,不是因為這個女人愛神多,神才赦免她；不是因為她的眼淚她的昂貴香膏,神才赦免了她.早在她踏步進入西門家之前,早在這女人鼓足勇氣要到耶穌面前來認罪悔改的時候,耶穌基督就已經赦免了她.不是因為她愛耶穌多,所以她蒙赦免,愛不是賺取赦免的條件,愛是已經被赦免的證明,愛是對赦免的回應.她的愛恰恰證明瞭,神已經賜給她的恩典有多麼的大!當這個女人認識了神,決心悔改的時候,她就蒙了赦免.神的赦免帶給了這個女人全新的生命,從此舊事已過,在基督裡她就成了新造的人.所以她才會不怕丟面子來到耶穌面前,才會哭的這麼動情,才會愛得這麼無所畏懼!這份愛,恰恰是這個女人得蒙赦免的強大證明.所以主耶穌說的沒錯:得恩免多的人,她的愛也多.

想到這裡,我心裡就特別感動.今天,我能夠站在這個講臺上講道,不是因為我有多優秀,不是因為我只欠神5兩銀子,而底下聽道的在座諸位就欠了神50兩,不是的!是神要使用我這個罪人,是神的愛讓我站在這裡.

同樣,今天我們聚集在這裡唱詩,聽道,讚美,奉獻；今天我們在這裡做主席,做司事事奉上主,今天我們在這殿裡受洗,一起成為愛的團契,不是因為我們想用我們的行動,我們的愛來博得主的歡心,來賺取主的赦免--赦免是賺不來的,恰恰是因為我們已經悔改,已經得蒙赦免,我們才能在這裡愛我們的神."我們愛,因為神先愛我們!"

從前幾節經文來看,耶穌行了好多神跡,他醫治羅馬的軍官,使拿因城寡婦的獨子復活,早就彰顯了他的身份--他就是那位拯救世人的救主,他的到來,乃是要全世界因他而得生命.這個婦人來膏抹耶穌,是因為她已經認清了耶穌是誰,所以無論多麼艱難險阻,她都要邁出這一步,不怕人嘲笑,不計算代價,只為了向主表達她最誠摯的悔改,所以在馬可福音中,耶穌誇獎這一類的美事,說"她是為我安葬的事把香膏預先澆在我身上."

但很諷刺的是,一個罪人能夠明白的事情,素來博學虔誠的法利賽人卻搞不懂.尤其是西門,他聽說過耶穌的神跡大能,所以他冒著風險請耶穌回家吃飯,他尊稱耶穌為老師,他尊重耶穌的教導,可是到最後,他的有色眼鏡還是讓他沒有認出耶穌是誰,當這有罪的女人邁出她生命中最重要的一步的時候,博學的西門還在那裡疑慮重重.

今天的教會裡也如此啊,有的人在頭腦裡、知識上知道耶穌是主,但在心裡頭,在行事為人上卻還不知道耶穌是主.有的人來了教會很久了,卻遲遲不肯決志,帶著疑惑來到教會,又帶著疑惑離開教會；還有的人雖然決志了,卻總是長不大,生命被許多思慮煩擾給擠住了,不能像那個有罪的女人一樣把自己委身給主.所以,我想問一問,是什麼樣的有色眼鏡讓我們疑慮重重呢?是什麼樣的困擾讓我們不能邁前一步呢?我們每次來到這裡,卻每次與主耶穌擦肩而過,是多麼可惜的事啊!

信心,Faith這個字,在希臘文原文裡頭的意思就是"跳躍",Jump.也就是說,你只有跳下水去,你才能學會游泳.我們只有把心一橫,才能真正接受主的恩典、主的赦免.就像那個有罪的女人,只有當她向前邁出一步,去向耶穌認罪悔改時,她的生命才真正得到了釋放,才真正得到了自由.這種自由讓她不再從人的眼色來衡量她自己,而是從神的眼鏡來衡量她自己,讓她活得充滿勇氣,充滿力量,充滿愛.弟兄姊妹,今天,你願意向前邁出一步麼?



\newpage

\section{一個人值多少}
\label{sec:54}
\textbf{林豪恩先生 - 一個人值多少?}
\newline
\newline
\hyperref[sec:53]{< < < PREV SERMON < < <}
~
\hyperref[sec:toc]{[返目錄]}
~
\hyperref[sec:55]{> > > NEXT SERMON > > >}


路加福音 8:26-39
\newline
\begin{longtable}{cl}
\hline
\hline
章節 & 經文 (和合本修訂版)\\
\hline
8:26 & \begin{tabularx}{0.7\textwidth}{X} 他們到了格拉森人的地區，就在加利利的對面。 \end{tabularx} \\ \\ \relax
8:27 & \begin{tabularx}{0.7\textwidth}{X} 耶穌上了岸，就有城裡一個被鬼附的人迎著他走來。這個人好久不穿衣服，不住在屋子裡，而住在墳墓裡。 \end{tabularx} \\ \\ \relax
8:28 & \begin{tabularx}{0.7\textwidth}{X} 他看見耶穌，就喊叫著俯伏在他面前，大聲說：「至高神的兒子耶穌，你為甚麼干擾我？我求你，不要叫我受苦！」 \end{tabularx} \\ \\ \relax
8:29 & \begin{tabularx}{0.7\textwidth}{X} 這是因耶穌曾吩咐污靈從這人身上出來。原來這污靈屢次抓住他；他常被人看守，又被鐵鏈和腳鐐捆鎖，他竟把鎖鏈掙斷，被鬼趕到曠野去。 \end{tabularx} \\ \\ \relax
8:30 & \begin{tabularx}{0.7\textwidth}{X} 耶穌問他：「你的名字叫甚麼？」他說：「群」；這是因為附著他的鬼多。 \end{tabularx} \\ \\ \relax
8:31 & \begin{tabularx}{0.7\textwidth}{X} 鬼就央求耶穌不要命令他們到無底坑裡去。 \end{tabularx} \\ \\ \relax
8:32 & \begin{tabularx}{0.7\textwidth}{X} 那裡有一大群豬正在山坡上吃食，鬼央求耶穌准他們進入豬裡；耶穌准了他們。 \end{tabularx} \\ \\ \relax
8:33 & \begin{tabularx}{0.7\textwidth}{X} 於是鬼從那人出來，進入豬裡，那群豬就闖下山崖，投進湖裡，淹死了。 \end{tabularx} \\ \\ \relax
8:34 & \begin{tabularx}{0.7\textwidth}{X} 放豬的看見這事就逃跑了，去告訴城裡和鄉下的人。 \end{tabularx} \\ \\ \relax
8:35 & \begin{tabularx}{0.7\textwidth}{X} 眾人出來，要看發生了甚麼事；到了耶穌那裡，發現那人坐在耶穌腳前，鬼已離開了他，穿著衣服，神智清醒，他們就害怕。 \end{tabularx} \\ \\ \relax
8:36 & \begin{tabularx}{0.7\textwidth}{X} 看見這事的人把被鬼附的人怎麼得醫治的事告訴他們。 \end{tabularx} \\ \\ \relax
8:37 & \begin{tabularx}{0.7\textwidth}{X} 格拉森周圍地區的人，因為害怕得很，都求耶穌離開他們；耶穌就上船回去了。 \end{tabularx} \\ \\ \relax
8:38 & \begin{tabularx}{0.7\textwidth}{X} 鬼已從身上出去的那人懇求要和耶穌在一起，耶穌卻打發他回去，說： \end{tabularx} \\ \\ \relax
8:39 & \begin{tabularx}{0.7\textwidth}{X} 「你回家去，傳講神為你做了多麼大的事。」他就走遍全城，傳揚耶穌為他做了多麼大的事。 \end{tabularx} \\ \\
[1ex]
\hline
\hline
\end{longtable}
日期:2010年7月11日

題目:一個人值多少?

証道:林豪恩先生

經文:路加福音8章26-39節

引言:當我們用錢來衡量人

你有沒有問過自己值多少錢呢?你會用甚麼來計算自己的價格呢?人籌保險的賠償金額?工作年期的收入總和?聽說在第二次世界大戰時,希特勒曾拿猶太人做實驗,分析一個人所含的物質元素,到底值多少錢,結論是人體內所含的油脂,能做一塊肥皂,所含的磷質,能做一盒火柴,所含的石灰質也不過值幾塊錢.不過,後來卻有人為人的價值平反,指出人體身上的血管,合計長達十萬公里,可以繞地球二圈半,而一公分長的人造血管,值新台幣270元,光是人體內的血管,就值27億新台幣.這樣的平反,是要說明人是值錢的東西.

或者如希特勒的傳說所言,人只值幾塊錢；又或者如那位平反者所言,單單是血管就值27億台幣；又或者用人籌保險的賠償金額以及工作收入來計算,總括來說就是用錢來衡量人的價值.把人貼上一個價格標示(in terms of money),就能夠製造一種一致的比較標準,有助於計算盈虧,便利於作決定.

豆腐渣建築工程在四川8級大地震中埋葬數以千計學生.在驚愕悲痛的時候,舉國上下都誓言要重建合規格的樓房,保障人民性命.為甚麼地震後兩年,記者卻發現竟然有大規模的豆腐渣民房重現於災區呢?災民慨嘆:「好不容易逃出來,又要我住進豆腐渣去!」為甚麼會有如此荒謬的事情呢?很可能是因為把人命與工程比較之後,人命仍然不太值錢.建築商對記者透過,建築商自己要賺錢,又要給政府官員塞錢,「一個工程,幾十萬一百萬,至少要給5萬塊政府.」因此,唯有盡力壓低成本,偷工減料是在所難免的了.【明報專訊2010年5月12日】

2008年「三聚氰胺」問題奶粉毒害30萬名嬰兒.毒奶粉對人的傷害瀝瀝在目,為甚麼兩年之後,問題奶粉再度現身市場呢?衛生部通報,已有3種「毒奶粉」流入廣州市場.業內人士估計,在三鹿奶粉事件之後,全國市場上仍有約10萬噸毒奶粉未被銷毀.【明報專訊2010年2月5日】為甚麼有如此荒謬的事情呢?很可能是經過格價之後,相對來說,嬰兒並不太值錢.根據報章報導,『結石寶寶』賠償方案標准是:『結石寶寶』死亡病例補償20萬元,重癥病例補償3萬元,普通癥狀補償2000元.如果企業清盤,就可能一分錢也不用賠.這是非常客觀的計量,十分精明的盤算.

一個人值多少錢呢?給予任何一個數目,都是危險的.以錢來衡量人,意味著把人非人化,把人變為一件商品,或者一件工具.各位生活於香港的朋友,我們沒有受到豆腐渣工程以及毒奶粉的威脅,不過,你有沒有感到自己也有被以錢來評值的壓力呢?有沒有感到自己被視為一件物品或者一件工具的危機呢?

今天福音經課記載了一個拯救的故事,一位被非人化的人,失去人的身份的人,因著與耶穌基督的相遇而重捨人性.

經文(路加福音8章26-39節)

他們到了格拉森〔有古卷作加大拉〕人的地方,就是加利利的對面.

耶穌上了岸,就有城裡一個被鬼附著的人,迎面而來．這個人許久不穿衣服,不住房子,只住在墳塋裡.

他見了耶穌,就俯伏在他面前,大聲喊叫,說,至高　神的兒子耶穌,我與你有甚麼相干．求你不要叫我受苦.

是因耶穌曾吩咐污鬼從那人身上出來．原來這鬼屢次抓住他,他常被人看守,又被鐵鍊和腳鐐捆鎖,他竟把鎖鍊掙斷,被鬼趕到曠野去.

耶穌問他說,你名叫甚麼．他說,我名叫群．這是因為附著他的鬼多.

鬼就央求耶穌,不要吩咐他們到無底坑裡去.

那裡有一大群豬,在山上喫食．鬼央求耶穌,准他們進入豬裡去．耶穌准了他們.

鬼就從那人出來,進入豬裡去．於是那群豬闖下山崖,投在湖裡淹死了.

放豬的看見這事就逃跑了,去告訴城裡和鄉下的人.

眾人出來要看是甚麼事．到了耶穌那裡,看見鬼所離開的那人,坐在耶穌腳前,穿著衣服,心裡明白過來,他們就害怕.

看見這事的,便將被鬼附著的人怎麼得救,告訴他們.

格拉森四圍的人,因為害怕得很,都求耶穌離開他們．耶穌就上船回去了.

鬼所離開的那人,懇求和耶穌同在．耶穌卻打發他回去,

說,你回家去,傳說　神為你作了何等大的事．他就去滿城裡傳揚耶穌為他作了何等大的事.

釋經與應用－一個人被非人化的過程

經文描述這個人已經「許久不穿衣服」,又指出這個人「不住房子,只住在墳塋裡」,並且描述他「常被人看守,又被鐵鍊和腳鐐捆鎖,他竟把鎖鍊掙脫」.這是一個怎樣的人呢?不穿衣服,意味著失去了人的尊嚴；不住房屋而住在墳場,等於與人群隔絕；被鐵鍊捆鎖,即是失去自由；竟然連鎖鍊也掙脫,等同怪物,非常嚇人.總括而言,這不是一般的人,甚至不再是人.有一位如此不像人的人留在社群中,對該社群來說,是很「頭痛」的事情,如何處理呢?故事記載說:「這人被鬼趕到曠野去」.有意無意地,鬼這個行動幫上了一個忙,清除了社群的威脅,化解了人們的困擾.

這個人本來是怎樣的呢?我們可以合理假設,他本來是穿衣服的,因此經文描述他「許久不穿衣服」；我們可以想像,他本來不是住在墳地,而是住在人群中的；我們又可以假定,他本來是自由而沒有被人用鐵鍊綑鎖著的.即是說,他本來是一個人:一個有尊嚴,一個與其他人有關係,一個自由自主的人.是甚麼剝奪了這個人的人性,把他改變為不像人呢?這個故事中的罪魁禍首是污鬼.經文指出,污鬼附在這個人身上,控制著這個人,令他失去作為人的尊嚴,把他與人群隔離,剝奪了他的自由及自主.污鬼,控制了這個格拉森人,把他非人化,把這個人折磨得不似人形.

今天,在我們的社會中,是甚麼力量企圖控制人呢?是甚麼力量企國把人非人化呢?是甚麼力量令人活得不似人呢?有先賢說:「忽視魔鬼的存在是天真,但把所有事物看為魔鬼的傑作則未免抬舉了牠」.有時我們會聽到污鬼侵犯個別人士的個案,但能夠廣泛地把人集體地非人性化的力量,恐怕連污鬼自愧不如.

最近,媒體上流傳著一個新名詞,稱為「窮忙族」.窮忙族概念來自日本知名經濟學者門蒼貴史,他把窮忙族界定為「即使揮汗拚命工作,卻仍然無法擺脫最低水準生活的一群人」.窮忙族會愈忙愈窮,愈窮愈忙,愈來愈沒有競爭力.作者估計,日本有四份一勞動人口屬於窮忙族.作者分析,主因是日本企業減少了正式員工人數,改為聘請合約或兼職.非正式員工的薪資及保障明顯較差,終生所得只及正式員工的五份之一左右.窮忙族其實不但在日本出現,二岸三地的華人社會都有很大迴響.

窮忙族人工低,工時長,生活就是「返工－放工－睡覺－再上班」的循環,上班的目的就是維持餓不死以繼續上班.一位大學畢業已經工作了四年的青年任職廣告行業,每天起床時父母已經返工,每天放工回家父母已經睡覺.一位三十而立的青年,工作了十多年,一個合約接一個合約,結果現在的薪酬比十年前還低.一位二十多歲的少女,在髮型屋工作多年,底薪三千,洗一個頭八元,連勤工獎,月入約四,五千.這位少女為自己只能給媽媽一千元家用而感到羞愧.

就兩岸三地而言,數字顯示資本正在不斷擴大,特別是內地,GDP每年都在增長,外匯儲備更令全球矚目,突破兩萬億美元大關,已達到第二大儲備國日本的兩倍.在這樣的情況下,人民應該更有條件活得更像人,生活應該可以更具人性,為甚麼人們的實際經歷卻不是如此呢?

水能載舟,亦能覆舟,或許「乾水(錢)」也如此.前香港金融管理局總裁任志剛指出,資金除了有優點(Virtue)外,卻同時存在數量大(Volume),變化多端(Variety),速度快(Velocity),反覆不定(Volatility)和不懷好意(Viciousness)等特性.資本只有一個目標,就是追逐更大的利潤；資本只有一個方向,就是流向有利潤的地方.因此,無論是勞動工人,管理人員,以及各大小老闆,即是說所有人,都承受著資本的壓力,被資金的流竄所威嚇著.資金流竄,弱肉強食,大規模的資本會越來越大,細規模的資本會被吞噬.資本的肚永遠填不飽,永遠都要求更高的回報,壓迫著各大小商業機構以及公私營部門止境地開源節流,製造著永無休止的惡性競爭.於是,在盡量減低人力成本的原則下,工人的實質工資不升反跌；那些要向上交數,而不得不向下擠壓的管理階層就有苦自己知；世界經濟一體化,市場二十四小時運作,連富甲天下的大老闆也睡不安寢.強橫的資本,漸漸消滅其他價值,成為衡量事物及作決定的唯一標準.資本把人變成它的僕役,在不同程度上削弱人的尊嚴,離間人的關係,剝奪人的自由.資本把人視為工具,因此工廠以「機械人管理」模式來管理工人,它看不見作為人的情感需要,家庭生活需要,安全及穩定的需要,發展及前途的需要等.資本去解決事情的方法仍然不離資本的本質:給你多一點錢,你就要留住你的性命,繼續為我服役,千萬不要再提作為一個人的種種需要.

耶穌把一個被非人化的人恢復人性

污鬼操控著那位格拉森人,卻沒有取其性命,只是令其生活得不似人.耶穌遇見這位可憐的格拉森人,立即吩咐污鬼從那人身上出來.按經文的記述,污鬼完全不是耶穌的對手.污鬼在耶穌面前,只能夠掙扎著懇求:

「至高神的兒子耶穌,我與你有甚麼相干?求你不要叫我受苦!」(v28),只能夠「央求耶穌准他們進入豬群裡去」(v32).得到耶穌的批准,污鬼立即離開那人,進入豬群裡去,豬群闖下山崖,跌入湖裡淹死了.

讓我們看看現在經文如何描述那位格拉森人:

「坐在耶穌腳前,穿著衣服,心裡明白過來」(v35).這位思想一度被操控的人現在神志清醒過來了；這位一度失去尊嚴的人現在穿上得體的衣服了；這位一度失控的人現在能夠安靜地受教於耶穌了.這位被非人化的格拉森人現在回復人性了.

社會拒絕把非人化的人恢復人性的行動

耶穌與污鬼交戰,高下立見,耶穌於靈界的能力是無容置疑的,也是我們可以信任及信靠的.值得置疑的,不是耶穌的能力,而是人們對於耶穌的行動的反應.耶穌幫助了那位被非人性化的格拉森人回復人性,幫助那位活得不似人的格拉森人重新做人,你會為他高興嗎?合理的假設是格拉森社群為這位康復的兄弟而感到安慰,並且感激拯救他們的兄弟的耶穌.然而,故事的發展卻非常出人意外:

「眾人出來要看是甚麼事．到了耶穌那裡,看見鬼所離開的那人,坐在耶穌腳前,穿著衣服,心裡明白過來,他們就害怕.」(v35)那位格拉森人的鄉親父老,叔伯兄弟看見他們的同伴回復正常,他們的反應竟然是害怕!這反應實在非常奇怪.格拉森社群了解過事情的來龍去脈後,沒有改變觀感,反而

「害怕得很」(v37).他們不但不歡迎這位拯救了他們的兄弟的耶穌,竟然好像見了鬼一樣

「求耶穌離開他們」(v37).

或者,格拉森的鄉親父老,叔父兄弟出來要看的,並不是那個被鬼附的人,而是要看那群衝入湖淹死了的豬；或者,他們要了解的並不是他們的兄弟如何康復,而是那群豬如何失去.一個人的康復與否事小,整個社群的經濟損失事大.經過比較,格拉森社群清楚意識到,那位耶穌的救人方法成本實在太高,完全不合乎經濟效益.他們一致評定:這個人不值那麼多錢,甚至沒有人值那麼多錢.若果要用那麼多錢才幫助一個人重新做人,再多幾個這樣的個案,豈不是把我們的經濟拖垮?

格拉森人寧願繼續活在非人性化的狀態中,都不願意接受耶穌助人回復人性的方去,他們把耶穌當作過街老鼠,希望耶穌「有咁快走咁快」,盡快在他們的社群消失.

原來,不想人活得像人的,除了污鬼之外,還有人類自己的社會及其社會價值觀.

幾十萬人的工廠,如果容許工人發表意見,會如何混亂呢?如果容許工人發洩情緒,會如何失控呢?如果容許工人休閒娛樂,會浪費多少工作時數呢?如果容許工人自由發展人際關係,談婚論嫁,會如何影響他們的工作狀態呢?如果容許工人組織工會,會增加多少管理風險呢?其實,工廠只需要運作金屬機器的人肉機器,過去,我用九百元把你們變做機器,如果你們認為不足夠,我可以提高你們的價格,不過,我需要的始終是機器,而不是人.可能有人會問,當地政府為甚麼任由她的人民被當作機器而不加以保護呢?或者,連政府也被資本的流竄所震懾,如果資本撤去越南或柬埔寨,對當地的GDP有甚麼影響呢?地方政府能夠承擔得起責任嗎?資本撤走,或者資金流走,對現代社會簡直是最大的威脅.

人性是甚麼呢?人性可以當飯吃嗎?人性值多少錢呢?小小的資本,引誘耶穌時代的格拉森社群選擇了犧牲人性而保護資本；大大的資本,逼令現代人變成工具,任其魚肉,甚至互相撕殺.

在洪流中發出吶喊

按經文的脈絡,故事至此可以結束.為何如此說呢?因為故事開始如此描述:「他們到了格拉森〔有古卷作加大拉〕人的地方,就是加利利的對面.耶穌上了岸,就有城裡一個被鬼附著的人,迎面而來．」故事的發展告訴讀者,耶穌已經把污鬼從那人身上趕走,解決了要處理的問題,第37節總結呼應文首:「耶穌就上船回去了.」

那麼,第38及39節存在的目的是甚麼呢?

「鬼所離開的那人,懇求和耶穌同在．耶穌卻打發他回去,說,你回家去,傳說　神為你作了何等大的事．他就去滿城裡傳揚耶穌為他作了何等大的事.」那位被耶穌拯救的人,懇求耶穌收他為徒,好像其他門徒一樣跟隨耶穌出入,這是合情合理的選擇.誰會願意留在一個不重視自己的地方呢?誰會享受處身於與自己的價值觀衝突的群體呢?不過,耶穌對這個人的懇求的反應又是出人意外,他竟然「打發他回去」,而不讓他跟隨.我們不用猜想耶穌看不起這個人,反而很可能是耶穌非常重視這個人,因此把重任交給他,要他在不重視人性的社會中傳講自己如何被耶穌恢復人性的故事.這實在是重任,孤身隻影,在洪流中發出另類聲音,宣揚與主流價值相左的信息,連舊約時代的先知也「灑手拎頭」.經文記述這位被耶穌拯救的格拉森人

「就去滿城裡傳揚耶穌為他作了何等大的事」.他能夠承擔重任,獨排眾議,擇善固執,這位曾經被輕視,被趕逐,被綑鎖的格拉森人現在活出人的美善,值得尊敬及欣賞.

總結

或許我們都或輕或重地曾經被非人性化地對待,或許我們都或多或少地曾經感到自己變成了工具,或許我們曾經因為人性尊嚴被否定而感到憤怒,或許我們曾經被人排擠而感到不安及抑鬱,或許我們曾經被評為沒有價值而感到沮喪,耶穌與我們相遇,告訴我們創造我們的上主並不是如此看我們,在創造主的眼中,每一個人都有本然的價值,每一個人都是無價的.耶穌與我們相遇,除去那掛在我們身上的價錢牌,恢復我們豐盛及美善的人性.

讓我們從對待身邊的人開始,可能是我們的配偶,子女,父母,兄弟姊妹,學習以尊重的態度對待他們,嘗試承認及滿足他們對親密,安全,信任等人性的需要,互相引發創造主已經賜予給我們的美善.我們是為此而被創造,也因此而生存.



\newpage

\section{重價的恩典}
\label{sec:55}
\textbf{鄧瑞強博士 - 重價的恩典}
\newline
\newline
\hyperref[sec:54]{< < < PREV SERMON < < <}
~
\hyperref[sec:toc]{[返目錄]}
~
\hyperref[sec:56]{> > > NEXT SERMON > > >}


路加福音 9:57-62
\newline
\begin{longtable}{cl}
\hline
\hline
章節 & 經文 (和合本修訂版)\\
\hline
9:57 & \begin{tabularx}{0.7\textwidth}{X} 他們在路上走的時候，有一個人對耶穌說：「你無論往哪裡去，我都要跟從你。」 \end{tabularx} \\ \\ \relax
9:58 & \begin{tabularx}{0.7\textwidth}{X} 耶穌對他說：「狐狸有洞，天空的飛鳥有窩，人子卻沒有枕頭的地方。」 \end{tabularx} \\ \\ \relax
9:59 & \begin{tabularx}{0.7\textwidth}{X} 他又對另一個人說：「來跟從我！」那人說：「主啊，容許我先回去埋葬我的父親。」 \end{tabularx} \\ \\ \relax
9:60 & \begin{tabularx}{0.7\textwidth}{X} 耶穌對他說：「讓死人埋葬他們的死人，你只管去傳講神的國。」 \end{tabularx} \\ \\ \relax
9:61 & \begin{tabularx}{0.7\textwidth}{X} 又有一人說：「主啊，我要跟從你，但容許我先去辭別我家裡的人。」 \end{tabularx} \\ \\ \relax
9:62 & \begin{tabularx}{0.7\textwidth}{X} 耶穌對他說：「手扶著犁向後看的人，不配進神的國。」 \end{tabularx} \\ \\
[1ex]
\hline
\hline
\end{longtable}
日期:2010年7月18日

題目:重價的恩典

証道:鄧瑞強博士

經文:路加福音9章57-62節

各位弟兄姊妹,主內平安.

現代的產品,大多數都盡量做到「用家友善」(user-friendly).以前的電視,要人手調校頻道；現代的電視,卻能自動追台,只要按掣,便能將一切頻道記錄下來.多麼方便!以前的汽車,是「棍波」車,「轉波」很麻煩；現代的汽車,是「自動波」,「轉波」自動化,多麼方便!

當然,若你學駕車的時候,是學「自動波」的車,則你面對一輛「棍波」車,可能是束手無策的.若然全世界的車都是user-friendly的「自動波」車,則當然無問題；然而,世事不能盡如人意,偶遇一輛「棍波」車,一輛不friendly 的車,則那些只能操控user-friendly的「自動波」車的人,只能乾著急.

由於我們已習慣了user-friendly的產品,不user-friendly的,只能被淘汱.想想:若你做一個推銷員,推銷一部電視.顧客問你,這電視好不好用?你說:「這電視要手動選台,還有,這電視沒有記憶程式,故此,每次你開機時,總要重新再調一次,像拉小提琴前,每次都要重新調一次一樣.」聽到推銷員這樣講,你會買這user-enemy的產品嗎?

我們喜歡user-friendly的產品.那麼,信仰又如何?

我們會否怕沒有人相信基督,而將基督變得像「靚模」一樣去吸引人?我們會否為了推銷信仰,將信仰變得過份user-friendly?我們會否為了推銷耶穌,將耶穌變成個人的寵物?我們會否為了推銷十字架,將十字架鍍金鑲鑽石,忘記了十字架意味著苦難與死亡?

德國神學家潘霍華(Bonhoeffer)寫了一本書,叫《追隨基督》(Discipleship),第一章的題目是:「重價的恩典」,這也是今日講道的題目.這神學家指出,現今的教會,宣揚廉價的恩典,cheap grace.將信仰弄cheap了.什麼是廉價的恩典?廉價的恩典就是人不需真誠悔改便可獲得寬恕；加入教會卻不接受教會的管教；領受聖餐卻不用真正與基督結連；手拿十架卻只有好處而沒有負起信仰的代價.

美國神學家Richard Niebuhr(尼布爾)一語講出基督教的迷失狀況,他說,有些人傳講的福音,是 ``A God without wrath brought men without sin into a kingdom without judgment through the ministrations of a Christ without a cross."「一個沒有震怒的神,藉著一個沒有十架的基督,帶領那些沒有罪的人,進入一個沒有審判的天國.」(The Kingdom of God in America)

如此這般,我們是否將信仰變得過份user-friendly呢?

若你只懂駕user-friendly的「自動波」車,怎能駕較複雜的「棍波」車呢?若你的信仰這般user-friendly,怎能應付複雜的人生呢?

今日講道的經文,是今日的福音書經課:路加福音9: 57-62.在主耶穌口中,跟從他,毫不user-friendly.

路加福音 9:57-62

9:57　他們走路的時候,有一個人對耶穌說:「你無論往哪裡去,我要跟從你.」

9:58　耶穌說:「狐狸有洞,天空的飛鳥有窩,只是人子沒有枕頭的地方.」

9:59　又對一個人說:「跟從我來!」那人說:「主,容我先回去埋葬我的父親.」

9:60　耶穌說:「任憑死人埋葬他們的死人,你只管去傳揚神國的道.」

9:61　又有一人說:「主,我要跟從你,但容我先去辭別我家裡的人.」

9:62　耶穌說:「手扶著犁向後看的,不配進神的國.」

這裡有三個人,三個都想跟從耶穌,他們分別代表三種意向.

這段經文的背景脈絡,是9:51:

「耶穌被接上升的日子將到,他就定意向耶路撒冷去」.耶穌將要完成天父的旨意,他定意向耶路撒冷去.耶路撒冷意味著十架和死亡.他在步向十架的路上.就在這路上,這三個人走過來,要跟從耶穌.

(1)你知道什麼是「跟從主」嗎?第一個人走過來,對耶穌說:「你無論往哪裡去,我要跟從你.」多麼堅決,很好.但這是否也顯示出,這人根本不明白跟從基督是什麼意思.這就像有人對你說:「你無論要求什麼,我都應承你.」這只能顯明,這個人根本不明白什麼叫做「應承」.他太「口輕輕」了.

你若對主耶穌說:「你無論往哪裡去,我要跟從你.」你期望主耶穌怎樣回答你?我想,我們大多數人,都期望主耶穌說:「好,很好.」主耶穌對這人的回應,有點令人詫異.主耶穌說:「狐狸有洞,天空的飛鳥有窩,只是人子沒有枕頭的地方.」這段對話用英文來表達,更突顯主耶穌的用意.那人說,我要跟從你,wherever you go.耶穌卻說,nowhere I lay my head.看來,那個人想跟從基督,還帶有某些自己隱藏的想法,there is somewhere to go.但對耶穌而言,所謂跟從基督,就是走上十架路.在這裡,你要放下一切自己的想法.跟從基督,你就沒有你自己立足之地了.There is nowhere to lay yourself.

一個故事:一個男人想討一個女人的歡心.那男人說:「無論你要什麼,我都會買給你.若你要大屋,我會買最好的大屋給你.若你要鑽石,我會買最閃的鑽石給你.若你要玫瑰花,我會每日送你一百枝玫瑰花.你到底想我買什麼送給你?」那女士說:「我想的,就是你不要時時想著藉買些什麼來討好我.我就是不要你以這種想法來對我.」

這男士的想法背後,有一個隱藏的想法,總能買些什麼,令她動心.There should be something to make her happy.但那女士的回應是:重點不在這些東西.It is nothing of this kind.

這個男士的問題,是完全放錯重點.他以自己的想法,去想別人的想法.那女士要的,卻是希望他放下自己的想法.

在福音書裡,那個想跟從基督的人,是否也完全放錯重點?他是否以為,仍可以有自己的觀點?

主耶穌的回答,向他顯明,他還未完全明白跟從基督是什麼意思.

跟從基督,就是完全放下自己,只有基督.這條路是一條十架路,在人看來,這是死路；在基督看來,這是完成天父心意的路,是一條真正的生路.

你願意放下自我嗎?

(2)什麼是更大的責任?

9:59　又對一個人說:「跟從我來!」那人說:「主,容我先回去埋葬我的父親.」

9:60　耶穌說:「任憑死人埋葬他們的死人,你只管去傳揚神國的道.」

這第二個人不是不想跟從耶穌,而是有其他人間的責任在身.人生有很多責任.為人父牧者,有父母的責任.為僱員者,有僱員的責任.為公民者,有公民的責任.責任無窮.

耶穌的回答,是讓這人明白,在人間一切責任之上,有更重要的事要做,要傳揚神的國,或者說,讓神顯明在人間.人生不能不承擔責任,但責任不能使人見不到神.

對主耶穌來說,讓人見到神,比單純的承擔責任更重要.每個人的責任範圍不同,一個負責任的人,可以負責任地拒絕責任範圍外的要求.

一個在醫院裡工作的人,可以負責任地拒絕醫院外的醫療事故的要求.一個照顧自己孩子的人,可以負責任地拒絕照顧他人的孩子.一個住在香港的人,可以負責任地拒絕照顧住在非洲的有需要的人.

就算比鄰而居,我們也可以負責任地拒絕鄰舍的要求.

但耶穌的講法,是指出要跟從他,就要超出這一切自我設定的責任範圍的限制.我們要大膽地將生命交出,不為自己設定的責任而活,乃為神國的責任而活.我們不能再以已經盡上自己的人間責任來為自己辯護.若要跟從神,則要承擔神給你的責任.

你是否願意承擔神給你的責任呢?

(3)跟從基督,不能講條件

9:61　又有一人說:「主,我要跟從你,但容我先去辭別我家裡的人.」

9:62　耶穌說:「手扶著犁向後看的,不配進神的國.」

這裡有一個人,想跟從基督,卻要講條件.但主耶穌說,真正的跟從,就是不可講條件,專心的跟從.

一個故事:(改編自:葛愛麗編:《心靈故事》,頁85-86)

有兩間教會,常在競爭狀態中.

有一次,甲教會的牧師教導耶穌為門徒洗腳的道理.為了讓年青信徒體會這道理,他吩咐他們到社區裡去做點服務人群的工作.

年青信徒很熱心,他們分組到社區做不同的服務.事後,他們回來報告.有一組去了醫院探望病人,其他組員為他們所作的感謝神.有一組買了些禮物送給窮人,其他組員為他們所作的感謝神.最後一組報告時,其他組員報以噓聲.這一組原來走去找與他們競爭的乙教會的牧師,問他可有什麼事需要幫忙.乙教會的牧師著他們去幫助一位老人家.於是,他們找到這位老婦人,幫助她整理花園.辛勞的工作完成了.那老婦人說:「你們這群乙教會的年輕人,真是我的救星.」這時,甲教會的牧師插嘴說:「你們至少應該告訴她,你們是甲教會的會友.」這個小組的成員說:「我們認為這並不重要.」

我們有時會為我們的善行,設定一些條件.

如:若要我在慈善晚會中捐些錢,則我要拿張大支票站在台上拍張照才行.但是,這是重要的嗎?或者,我們會想:我要跟從主,但我必須健康,財富,快樂,一無所缺.若是這樣,這還算是跟從主嗎?對跟從主而言,這一切重要嗎?主耶穌對這個人的回答,指出:真正的跟從,是無條件的.

我們在跟從主的過程中,是否設下條件?

(4)結語有一齣電影,叫《十月圍城》,片中講到孫中山要來香港,商討革命事業.有人要暗殺他,故此需要有人假扮孫中山,以引開暗殺者的注意.有一個富商的獨生子,擁護革命.他願意假扮孫中山.有人反對,因為這是危險之事,他可能會死,那富商會失去獨生子.然而,這位年輕人卻說,在革命面前,人人平等.誰也無權為誰去死誰不去死設下條件.既然信奉革命,就只能放下一切人間的責任,放下自己的計算,從容就義,完成生命的最大意義.最後,這獨生子死了,他完成了他的信仰帶來的最大責任.

我們是否明白跟從主的代價?是否明白這需要無條件地放下自我的一切?面對主的呼召,你又如何回應呢?

但願榮耀歸於聖父,聖子,聖靈.阿們.



\newpage

\section{沒有甚麼能傷害你們了}
\label{sec:56}
\textbf{伍渭文牧師 - 沒有甚麼能傷害你們了}
\newline
\newline
\hyperref[sec:55]{< < < PREV SERMON < < <}
~
\hyperref[sec:toc]{[返目錄]}
~
\hyperref[sec:57]{> > > NEXT SERMON > > >}


路加福音 11:17-24
\newline
\begin{longtable}{cl}
\hline
\hline
章節 & 經文 (和合本修訂版)\\
\hline
11:17 & \begin{tabularx}{0.7\textwidth}{X} 他知道他們的意念，就對他們說：「一國自相紛爭，必定荒蕪；一家自相紛爭，就必敗落。 \end{tabularx} \\ \\ \relax
11:18 & \begin{tabularx}{0.7\textwidth}{X} 撒但若自相紛爭，他的國怎能立得住呢？因為你們說我是靠著別西卜趕鬼。 \end{tabularx} \\ \\ \relax
11:19 & \begin{tabularx}{0.7\textwidth}{X} 我若靠著別西卜趕鬼，你們的子弟趕鬼又靠著誰呢？這樣，他們要作你們的判官。 \end{tabularx} \\ \\ \relax
11:20 & \begin{tabularx}{0.7\textwidth}{X} 我若靠著神的能力趕鬼，那麼，神的國就已臨到你們了。 \end{tabularx} \\ \\ \relax
11:21 & \begin{tabularx}{0.7\textwidth}{X} 壯士全副武裝，看守自己的住宅，他所有的都很安全； \end{tabularx} \\ \\ \relax
11:22 & \begin{tabularx}{0.7\textwidth}{X} 但有一個比他更強的來攻擊他，並且戰勝了他，就奪去他所倚靠的盔甲兵器，又分了他的掠物。 \end{tabularx} \\ \\ \relax
11:23 & \begin{tabularx}{0.7\textwidth}{X} 不跟我一起的，就是反對我；不與我一起收聚的，就是在拆散。」 \end{tabularx} \\ \\ \relax
11:24 & \begin{tabularx}{0.7\textwidth}{X} 「污靈離了人身，走遍無水之地尋找安歇之處，卻找不到。就說：『我要回到我原來的屋裡去。』 \end{tabularx} \\ \\
[1ex]
\hline
\hline
\end{longtable}
日期:2010年7月25日

題目:沒有甚麼能傷害你們了

証道:伍渭文牧師

經文:路加福音11章17-24節

傳說希臘著名數學家,物理學家阿基米德(Archimedes)有一次從浴室跑到街上,欣喜若狂大叫:我找到了!我找到了!(Eureka!Eureka!) 當時他真的得意忘形 and \#8212;他忘記穿衣服.他發現甚麼呢?他發現浮力定律──阿基米德定律,簡單地說,就是物體浸在水中所失去的重量,等於它所排開的液體重量.

因為國王吩咐珠寶工匠為他製造一頂純金王冠,但又恐防那頂王冠摻雜銀,所以國王吩咐阿基米德找一個不會損毀王冠,而又可作鑑別的方法.阿基米德為此苦思多日,直到一天在浴盆洗澡時發覺,當身體浸入水中時,水溢出盆外,並感覺到身體輕了,阿基米德突然想通了箇中意義,高興得連衣服也不穿,光著身子發狂的跑到宮中,並大喊:"我發現了!我發現了!"

阿基米德的鑑別的方法是把相同重量的黃金,放入裝滿水的容器中,如溢出的水比王冠小,表示王冠有摻雜密度較低的銀.

今天的福音書經課路加十一章17-24節說到耶穌的門徒同樣的興奮莫名,七十個門徒到各城傳道回來,歡歡喜喜向耶穌匯報工作,同樣欣喜若狂的說:"我們發現了!我們發現了!" .主啊,因你的名,連鬼也服了我們.阿基米德因發明了物理的定律興奮莫名,門徒發現一個更大的秘密 and \#8212;從沒有人發現屬靈的定律 and \#8212;因耶穌的名,連鬼也服了他們.

阿基米德定律可以使國王頭頂的皇冠保持純金,但這定律不能保證他腳下的臣僕永遠對他忠誠,國王仍然生活在被欺騙,被出賣的懼怕中,恐怕被出賣,被欺騙的焦慮終日綑縛著他.不竟,戴在頭頂的皇冠是假東西A貨,實在很沒有顏面.沙士比亞名劇凱撒大帝中有一幕,描述凱撒在元老院階梯上被刺殺時,說出發人深省的名句:tu brutu! (you too, Brutus!布魯圖,你也參與其中!) .凱撒大帝致命的一刀,竟來自極其信任的好朋友布魯圖,他心碎了."布魯圖,你也參與其中!" 永遠成為掌有權力者的惡夢,終日綑鎖著他的焦慮,愈是專制的掌權者,焦慮愈大.

基督的定律卻使人成為自由的君王.主啊,因你的名,連鬼也服了我們.連鬼也服了我們,我們還怕誰呢?馬丁路德說:基督徒連罪惡也勝過,我們就是最自由的君王了.

門徒奉耶穌的名趕鬼醫病,連鬼也服了他們.

``要醫治那城裡的病人"(10: 9) 福音所到之處,都有神蹟奇事伴隨,上帝的國度開展,黑暗的權勢節節敗退.耶穌說:我已經給你們權柄可以踐踏蛇和蠍子,勝過仇敵一切的能力.蛇在聖經代表撒旦和罪:女人的後裔要傷蛇的頭 and \#8212;是致命的,蛇要傷他的腳跟 and \#8212;是輕微的；以色列人因發發怨言犯罪被火蛇所咬(民21: 6, 7) ,被蛇咬的人可仰望摩西舉起的銅蛇得到醫治,被罪綑鎖的人仰望被十架舉起的基督得到醫治,如果大家留意,世界所有醫學會的徽號都是繞在杆上的銅蛇.

一    福音使我們得著脫離內心怨憤和仇恨的自由

願這一家平安.(10:5)用蛇和蠍子代表罪是蠻有意思的.毒蛇和蠍子是毒物,是至命的東西.毒物不同猛獸,遇到猛獸你會警覺,會對抗,但我們可能不經意地中了毒,毒氣慢慢漫延全身而不自覺,終於毒發身亡.毒藥可能披上糖衣,入口甘如飴；毒藥可能非常誘惑,蛇蠍美人,使人不能自己.但毒物仍然是毒物,傷害我們,致我們於死地,罪也如是.

耶穌的福音第一句是:

願這一家平安.(10:5) 平安正是世人需要的福音,福音能消解我們內心的怨憤和仇恨的毒害.故事的開始九章51節至56節,就指出撒馬利亞人因耶穌往耶路撒冷去,就不接待耶穌.由於宗教信仰不同,本是同更新的撒馬利亞人,拒絕耶穌和門徒,而門徒反應強烈,求耶穌降火燒滅他們,藉此顯明他們才是真正的信仰.耶穌說:

人子來,不是要滅人的性命,是要救人的性命.(9: 56) 因為被拒絕,要求降火滅絕對手,這想法最近是否耳熟能詳呢?本是同更新的政黨,因為路線不同,產生敵視,甚至仇恨.但福音把敵視和仇恨除掉,不讓這些有害的情緒傷害我們.這個故事的開始是撒馬利亞人拒絕耶穌和門徒,但故事的結束十章23-33節,路加卻以好撒馬利亞人的故事作為逆轉.

一個猶太人被強盜奪去一切,被打到重傷臥地,祭司經過不理,利未人經過不理,因怕接觸將要死的人,根據律法,承擔祭祀的人,不能接觸不聖潔的死人.但被猶太人岐視為雜種,不潔的撒馬利亞人,雖然趕路,卻停留下來,上前用油為他裹傷,帶他到客店療養,並付擔所有的費用.

因著福音的能力,我們可以勝過怨憤和仇恨.有需要的人,就是我們的鄰舍,我們要愛每一位有需要的人,不分種族,不分信仰.

耶穌說:我已經給你們權柄可以踐踏蛇和蠍子,勝過仇敵一切的能力,斷沒有一樣能傷害你們了.(10: 19)今天甚麼東西正傷害著你呢?是否多年在心中積存的怨憤,不能寬恕的心,像蛇纏著我們,像蠍子牢牢的拑著我們?這怨憤正毒害著我們的心靈.今天耶穌說:願你們平安；願你平安.主要釋放你,給你饒恕的心.從今天起,怨憤,仇恨不能再傷害你.

二    福音使我們勝過外在一切黝暗的力量

耶穌說:我已經給你們權柄可以踐踏蛇和蠍子,勝過仇敵一切的能力,斷沒有一樣能傷害你們了.(10: 19)

相信耶穌的的人,勝過黝暗的力量,罪不能傷害,撒旦也不能傷害我們.我在大學工作剛過十年,發現不少碩學鴻儒,與白衣布丁,一樣相信運程,認為冥冥中有一種力量支配著人生.剛來大學不久,參加全港高校教職聯會到四川旅行,有一位友校的理科教授拿著一個水晶球,到廟宇求主持和尚替水晶球開光,希望帶來祝福.一位多年在大學工作的資深教授,竟然在大年初三到沙田車公廟求簽,轉轉風車,換換運程,因他一年來身體不好,先後兩位親人離世.他跟我說:其實他的老師很早就跟他說,當人長大後,你就會發現謀事在人,成事在天,不由你不相信命數.

對人生沒有把握的人,其實有一種不安全的感覺.很多人相信運程命數,但運程命數是陌生的,不可知的,但又掌控著我們的生命,這才使我們懼怕.做生意的人,因為面對很多不明朗的因素,或虧或贏,不能掌按,他們覺得風水很重要,風水就像一種屬靈的力量,可能對自己不利的,所以我們試圖藉著一些擺設,把這些屬靈的力量收為己用,消災解難,帶來福氣.

耶穌說:我已經給你們權柄可以踐踏蛇和蠍子,勝過仇敵一切的能力,斷沒有一樣能傷害你們了.(10: 19) 門徒進入人家傳福音第一句就是:願這一家平安.

我信全能的父,創造天地的主.整個天地被上帝所造,而上帝是父,滿有慈愛,沒有一件事是失按的,祇是我們不明白,

萬事都互相效力,叫愛上帝的人得益處.(羅8:28)隱密的事屬於上帝,明顯的事屬於我們和我們的子孫,沒有一件事上主不知道,祇是我們不明白.(申29:29) 當聖經宣告:起初神創造天地,就把物質世界凡俗化.物質世界 and \#8212;太陽,月亮,星宿,山川不是我們膜拜的神靈,是我們研究的對象,科學因而興起.沒有一種力量不伏在上帝之下,"這是天父世界,我心滿有安寧......".

各位,今天你是否對前途擔心呢?是否總覺得有些東西對自己不利呢?耶穌說:我已經給你們權柄可以踐踏蛇和蠍子,勝過仇敵一切的能力,斷沒有一樣能傷害你們了.甚至在威脅我們生命的頑疾面前,我們也不懼怕.

剛才領洗的樂第兄在見証中分享說,他首次踏進教堂就有一種平安保護著他,因福音的安慰,勝過對頑疾和死亡的恐懼.耶穌說:"那殺身體不能殺靈魂的不要怕它,惟獨能夠殺身體又殺靈魂的上帝才要懼怕." 但上帝的兒子因著愛我們,已為我們捨命,我們與上帝和好了,並稱上帝為阿爸父,我們無所懼怕.

另一位領洗的余老太見証中,分享失去珍愛的珠寶時的失落和傷心,那時雖然對耶穌不太認識,但懇求耶穌為她找回寶貝,當找到時開心極了.我有機會聽她分享這事,也被她開心的緒所感染." 我找到了,我找到了,我尋回失而復得的寶貝,我心愛的寶貝."

門徒興奮,因為連鬼也服了他們.耶穌興奮,因為門徒認識上帝.我想到耶穌的歡樂,因門徒認識自己是上帝差來的救主而歡樂,因我們被上帝尋回而歡樂.

``正 當 那 時 , 耶 穌 被 聖 靈 感 動 就 歡 樂 , 說 , 父 阿 , 天 地 的 主 , 我 感 謝 你 , 因 為 你 將 這 些 事 , 向 聰 明 通 達 人 就 藏 起 來 , 向 嬰 孩 就 顯 出 來 ． 父 阿 , 是 的 , 因 為 你 的 美 意 本 是 如 此 ." (10: 21) 在宗教改革時期,致力調和信義宗與改革宗信仰的《海德堡問答》第一條即問: 你或生或死惟一的安慰在哪裡?

答:我或生或死,身體靈魂皆非己有,而是屬於信實的救主耶穌基督；祂用寶血完全補償了我的罪債,並且救了我脫離魔鬼一切的權勢；祂保守我,若天父不許,我的頭髮一根也不能掉下；祂又叫萬事互相效力,使我得救,是以祂藉聖靈也使我確知有永生,並且使我從此以後誠心為祂而活.

問:你或生或死惟一的安慰在哪裡?

答:我或生或死,身體靈魂皆非己有,而是屬於信實的救主耶穌基督；祂用寶血完全補償了我的罪債,並且救了我脫離魔鬼一切的權勢；祂保守我,若天父不許,我的頭髮一根也不能掉下；祂又叫萬事互相效力,使我得救,是以祂藉聖靈也使我確知有永生,並且使我從此以後誠心為祂而活.



\newpage

\section{使生命完滿的愛}
\label{sec:57}
\textbf{曹偉彤博士 - 使生命完滿的愛}
\newline
\newline
\hyperref[sec:56]{< < < PREV SERMON < < <}
~
\hyperref[sec:toc]{[返目錄]}
~
\hyperref[sec:58]{> > > NEXT SERMON > > >}


日期:2010年8月1日

主題:使生命完滿的愛

証道:曹偉彤博士

在我們所讀的經文中:

路加福音的十一章三-四節:我們日用的飲食,天天賜給我們.赦免我們的罪,因為我們也赦免凡虧欠我們的人.

馬太福音五章四十三-四十八節:你們聽見有話說,『當愛你的鄰舍,恨你的仇敵.』 只是我告訴你們,要愛你們的仇敵．為那逼迫你們的禱告.這樣,就可以作你們天父的兒子．因為他叫日頭照好人,也照歹人,降雨給義人,也給不義的人. 你們若單愛那愛你們的人．有甚麼賞賜呢．就是稅吏不也是這樣行麼. 你們若單請你弟兄的安,比人有甚麼長處呢．就是外邦人不也是這樣行麼.所以你們要完全,像你們的天父完全一樣.

這段經文是路加福音十章二十七節:他回答說,『你要盡心,盡性,盡力,盡意,愛主你的　神．又要愛鄰舍如同自己.』 我們同心禱告:親愛的主,我們懇求你將你的話語賜給我們,讓我們每個人都有能聽的耳,能聽到,能明白,能記得,能相信,常常在你的恩典中,我們這樣的禱告是奉主耶穌基督的名字,阿們.

我們都是生活在一個尋覓的年代,尋尋覓覓的年代,尋尋覓覓的經驗,不單是這個年代的特性,也可以說是人的特性.在我們很熟識的家裡,很多時候我們都會呼喊,我借了一本書,今天到期了,我要去還,那本書現在在那裡呢?我的車匙在那裡?當我們在找某一些東西的時候,在平時我們是可以看見的,但當我們要找這些東西的時候,偏偏我們找不著,人生就是這樣的尋尋覓覓.

馬太福音的智者,都是有一種尋尋覓覓的經驗,他們從熟識的地方去到不熟識的地方,他們跟隨星星的引導,去到伯利恆,尋找耶穌基督.這應驗以賽亞先知的預言,就是有一嬰孩由童貞女所生,祂的名字是以馬內利,就是神同在的意思.神同在:就是神道成肉身的愛.因為愛的原故,神要求祂的兒子,降世成人.耶穌基督因為愛的原故,順服神的心意,為罪人死.我們也曾經按著星星的指引,來到崇基學院的禮拜堂中,我們經歷到神的同在,能夠找到真正的愛.愛應該是我們每一個人很熟識的東西,但當我們需要愛的時候,你們我們能否找到愛呢?要回答這問題,今天早上讓我們一同去思想,在崇基學院的禮拜堂中,我們有否盡心的去愛神,盡心的去愛他人,如同愛自己一樣.在這裡是否充滿了神的愛.一方面,我們每一個人都是需要愛的,愛是我們每個人的基本需要.另一方面,當我們信耶穌到一段時候,我們在教會生活了越長的時候,我們發覺愛是一個美麗的理想,愛是一個難以實現的理想.為何我們會有這樣的經歷呢?從人間的經驗來分析,亞里士多得告訴我們,我們不能夠經歷到愛,我們不能夠順心的愛,因為我們很多時候,愛是公利式的,也是寬裕式的.公利式的愛,引申出來的是公利式的友情.公利式的愛,公利式的友情,就是當我們相處的時候,我們能夠得到一些好處的時候,我們就去交往.公利式的愛,公利式的友情,就是告訴人,你知道我背後有那些人是我認識的,我認識曾特首的,這是公利的友情.亞里士多得告訴我們,公利的友情是怎樣的呢?公利式的愛,公利式的友情是很多的,是好像蒼蠅那樣的多,當你擺設筵席的時候,你發覺很多的蒼蠅飛到你的廚房,來享用那筵席,但他們享用完了那筵席後,或是你再沒有什麼給他享用的時候,蒼蠅就會飛走.當你的友情,寬裕的友情同樣是一樣的,寬裕的友情是ｆｒｉｅｎｄｓｈｉｐ　ｏｆ　ｐｌｅａｓｕｒｅ,或是寬裕的愛,意思是當我和他們在一齊相處的時候,我覺得你和我很談得來,或是你的思想,你的背景和我很相近的時候,我和你相處是非常開心的.但亞里士多得告訴我們,寬裕的愛,公利的愛,這會是眾多的,但是不能夠長遠的.那些友情是眾多的,但那些友情是膚淺的.因為這樣的原故,因為公利的愛,寬裕的愛,我們會發覺難以去愛,難以去經歷去愛.但在另一方面,從一個教會的經驗去看的時候,我們發覺為何我們不能去愛呢?發覺愛只是好像一個,我們常常呼喊的理想呢?發覺愛是我們所熟識的東西,但當我們要尋索的時候,為何偏偏尋索不到愛呢?從信仰的經驗當中告訴我們,我們是不能愛的,因為我們都有很多生命的傷口,我們為何會有生命的傷口呢?因為在我們生命裡面,仍然有很多很多的情慾,有很多很多的慾念.在基督教的傳統中去講這些慾念,這些成為我們生命傷口的慾念,這叫做七宗罪.七宗罪是指什麼呢?七宗罪是指驕傲,妒忌,忿怒,貪心,貪吃,怠惰,好色.這些慾念是很平常的,人之常情,是你和我都有的,但我們的先賢告訴我們,當這些慾念在我們生命裡面滋生的時候,是會成為了我們生命的傷口.比喻講到驕傲,驕傲是一種迷失了的自愛,就是自己愛到一個地步,是迷失了方向,這種迷失了的自愛,在創世記中看到,當亞當夏娃吃了樹上的果子的時候,他們希望是能夠與神同等,全能,全知,全是,有神的能力.但他們吃了樹上的果子的時候,他們是赤身露體,當神向亞當講話的時候,就如爸爸媽媽跟兒女們講話的時候,當兒女犯了錯,爸爸媽媽都會微聲細語的說:孩子阿,你在那裡阿?你做了什麼事?其實父神的心意,是希望亞當夏娃站出來,我在這裡,我錯了.但我們看到人的反應是什麼呢?亞當的反應是什麼呢?亞當的反應就是說:是你所造的女人,使我這樣的.這種就是迷失了的自愛,意思是愛自己愛到一個地步,愛自己多於其他人,愛自己多於神.

妒忌是什麼?妒忌是聖徒的罪,是基督徒的罪,甚至有時是牧者的罪.當修道者討論這七宗罪的時候,其實他們有一種是追求神的心態,為此修道院的建立,是有一班朋友,他們希望去親近神,於是他們就離開他們的工作,離開他們的家人,到沙漠曠野中,一段短的時間,或是一段長的時間中,能夠多多的經歷神.有些弟兄有好的經歷,就有其他的朋友知道,他們就一同的去到曠野沙漠,去一同的經歷神.因為有這個原故,修道院就產生.其實在修道院中,是好得無比的,你說是不是,因為大家在修道院中,大家在神面前是平等的,每個人穿的都是同一類的衣服,吃的是同一類的食物,睡同一樣式的床,其實大家在神的面前是平平等等的,快快樂樂的,但人的慾念,偏偏就是我必需要給其他人知道,我們跟你們是不一樣的,我們是與眾不同的,我們是配受尊崇的.在教會中,很多時候我們看到有很多的爭競.我們一同來到神的殿中,我們一同敬拜,我們是平等的,我們是一式一樣的,我們有不同的性格,但是我們穿不同的衣服,但我們在神面前是平等的,但是我們在神的面前,偏偏我們不甘願大家是同一樣的.因為這樣的原故,我們在當中就有爭競.因為爭競的原故,我們產生很多很多的傷口,這講出我們心中的慾念,這慾念雖然是零亂的,但是它們慢慢的滋生,成為至死的罪,成為我們屬靈生命死亡的罪.因為這樣的原故,我們是難以去愛,因為這樣的原故,我們發覺在教會裡面,愛是一個難以實現的理想.親愛的主耶穌基督,祂是明白這個道理,明白這實況,因為祂其中的門徒猶大,因為貪婪的原故,買主買友,因為這種情慾的原故,離開神的,親愛主耶穌基督是明白的.親愛主耶穌基督仍然給我們有啟示,給我們生命的藥方.祂的啟示的藥方,要去除這肉體的慾念的藥方是什麼呢?耶穌基督的藥方就是祂的命令.正如我們剛才所讀的馬太福音五章48節所講:你們要完全,好像你們的天父完全一樣.

約翰福音十五章12節:你們要彼此相愛,像我愛你們一樣,這就是我的命令.耶穌基督說:你們要盡心,盡性,盡意,盡力愛主你的神,並且要愛鄰舍如同愛自己.弟兄姊妹們,這是耶穌基督給我們生命的藥方.我們要如何的去理解這藥方呢?我們要如何的去理解,在馬太福音五章47-48節呢?你們若單請你弟兄的安,比人有甚麼長處呢．就是外邦人不也是這樣行麼.所以你們要完全,像你們的天父完全一樣.

在48節有一個很重要的字眼:完全.什麼是完全?完全不是指道德上的完全,是指我們完全沒有罪,完全是指成長的過程,是一個完滿的過程.當我們再詳細的看經文的前文後理的時候,我們再深一度發覺,完全是指什麼呢?完全是指一種完全的愛.這種完全的愛,就是父神所彰顯完全的愛.這種完全的愛,就是父神的完全,這種父神的完全,就是父神完全的愛.父神完全的愛是怎樣呢?父神完全的愛是說:將祂的獨生子差到世上,為我們罪人的原故,釘死在十字架上拯救我們.父神完全的愛,是連我們這群敵擋他的罪人,惡人,他都愛.所以弟兄姊妹們,完全的愛,是指一個生命成長的過程,是能夠一點一點的成長,就能夠完滿.當我們有完全的生命的時候,我們就有完全的愛,當我們有完全的愛的時候,我們就有完滿的生命.

何謂完全的愛呢?完全的愛就是連敵人都愛.意思就是當我們繼續成長,繼續去愛的時候,我們是不斷的成長,我們可以愛很多的人,甚至愛到一個程度,連我們的敵人都去愛,這種是完全的愛.要完全的愛,要生命完滿,我們要聽從耶穌基督的命令:你們要盡心,盡性,盡意,盡力愛主你的神,並且要愛鄰舍如同愛自己.我們要有完全的愛,弟兄姊妹們,我們要盡心,盡性,盡意,盡力的愛,其實是一種意志的愛.很多時候我們因為因信稱義這重要的教導,而忽略了意志.但如果我們只強調意志,而忽略了一切的恩典,這是一個很多的問題,一個很大的危險.另一方面,當我們在恩典裡的時候,我們必需知道,原來愛的經歷,愛的實踐是需要意志的,盡心,盡性,盡意,盡力愛主你的神,並且要愛鄰舍如同愛自己.

什麼是盡心的愛?什麼是意志的愛?意志的愛是你和我都可以有的愛.問題是我們有沒有操練,有沒有學習.意志的愛在我們的父母身上可以看到的.好的父母是怎樣的呢?好的父母是當他每天工作後,非常的勞碌,非常的辛苦,或是要面對很多人際關係的問題.但好的父母,當他們工作完畢,一回到家的時候,看到兒女們有需要的時候,幫他們預備足夠的飲食,幫助他們解決功課的問題,每當兒女不聽話,不受教,反叛的時候,好的父母會繼續的用意志去愛.同樣美好的婚姻,都是需要有意志的愛.好的婚姻就是我們二人成為一體,當我們行完紅地毯後,我們經歷這很有情感的愛之後,原來發覺我們要接受對方缺點的時候,我們需要繼續有意志的去愛.所以弟兄姊妹們,意志的愛,是讓我們生命能夠慢慢的經歷到愛.所以今天在教會中,我們尋尋覓覓愛的時候,找不著愛的時候,我就說:所以弟兄姊妹們,我們有沒有盡心,盡性,盡意,盡力的去愛.這是一個生命成長的過程,是成長到一個地步,是連你的敵人都愛,這就是一種意志的愛.這種意志的愛,是要一點一點的學習,培養而得的.我們要鍛鍊這種意志的愛的時候,我們要生命完全,要有意志的愛,今天早上我要提醒你們,你們需要有好的朋友,你們需要在崇基學院禮拜堂中,你們要有很好的屬靈朋友.屬靈朋友是神給我們的禮物,是叫我們操練意志的愛,屬靈的好朋友是我們生命成長完滿的引擎.所以弟兄姊妹們,反過來,如果你在教會中,在你生活中,如果你是沒有朋友,我會說:你屬靈的生命可能是有問題的.相反的,如果你在教會裡,在崇基學院禮拜堂中,你會有越來越多的屬靈好朋友,我為你感恩,你的生命是在完滿的過程中.屬靈的朋友,好的朋友是叫我們生命能夠完滿,能夠完全的.屬靈的好朋友是怎樣呢?他是會給我們看到,我們屬靈生命的實況.好的朋友,是當我們有一些屬靈的恩賜的時候,他們會為我們歡呼,為我們快樂的,在當中肯定我們的恩賜.當我們有開心的事情的時候,好的朋友是會與我們同的歡喜快樂,一同的慶祝.好的朋友也會指出我們生命的缺點,指出我們生命的傷口.好的朋友他不會給我們失落在罪惡中.好的朋友是會拉著我們,不會讓我們去犯罪.不知不覺我己經回來香港有十一年,在十一年期間,在事奉中,我感謝神給我有一些屬靈的好朋友.

我記得在兩年多前,我遇到一位朋友,他是在基督教機構當主管.他在商界裡,工作了一段很長的時間,有很多行政的經驗.他希望來神學院跟我有多一點的認識,希望我能成為他的屬靈遮蓋(ｍｅｎｔｏｒ),但當我細心的去想的時候,我根本不能做他的屬靈遮蓋(ｍｅｎｔｏｒ),但做朋友也不錯,互相的相交,所以我們就每個月我們都見面,一同的禱告.每次見面禱告的時候,他也會跟我分享行政管理的心得,我也會從信仰和神學角度去分享我一些的看法,然後我們一同禱告.每次禱告,我都是很歡喜,因為他的禱告會用上很多的經文,每次禱告的時候,我能夠在當中和他一同經歷神的心意.

我記得在二ＯＯ八年五月份,當我要作出一些生命的決定時,我跟他分享,因為他是我的好朋友.我跟他分享我生命的方向.他對我說:曹博士,在過往的日子中,你有否看到神在你生命中的軌跡.當他這樣問的時候,我不知道他是想問我什麼?於是我就想辦法回答,在這半年我怎樣的有生命的軌跡,看到神在我生命中的軌跡.我回答完,我就問他:你是否想知道的就是這些.於是繼續的說:在先前的時間,你講出你生命的決定方向,是否有些血氣?我聽到的時候,一方面我很慚愧,但在另一方面我是很開心的.因為他是我一個非常好的朋友.弟兄姊妹,在教會中,你有沒有這樣的朋友?你有沒有這樣的屬靈朋友?可以讓你盡心,盡性,盡意,盡力的去愛.除了我們可以在教會中認識一些屬靈的好朋友,在中間是用我們的心,盡心的去操練,學習意志的愛去愛.我們還有一些人,是我們要繼續的操練.比喻,有一些人是你不喜歡的,可能他的背景跟你不同,甚至他會背後說你閒話,或是不喜歡你,或是漠視你,或是不知不覺成為你的仇敵,你有沒有這樣子的肢體呢?弟兄姊妹,意志的愛是成長的愛,是需要一點一點的操練學習.

我們除了要有屬靈的好朋友之外,我們也需要有意志的愛,操練的愛,甚至操練到一個地步連敵人都可以愛.當我們可以連敵人都可以愛的時候,當然不是今天的事,也不是明天的事,而是一生一世持續操練的過程.總有一天,你會操練到有意志的愛去愛你的敵人,當然這也是神的恩典在其中,去愛你的敵人時,使你的生命得以完全.

一位愛爾蘭的婦人,她是很喜歡幫助人的婦人,她丈夫是在監獄裡做主管的.有一天早,當她丈夫出門上班,預備開車的時候,有兩名愛爾蘭共和的軍人,來到他身旁,在他的臉上開了五槍.他太太聽到這槍聲之後,跑出門去看,見到她的丈夫躺在血中.這位婦人在血中抱著她的丈夫,在背誦主禱文.她懇求神赦免這些兇徒的罪,因為主已經赦免她的罪,她懇求主能夠幫助她,去寬恕這些殘害她丈夫的兇徒.以後她每天都向神獻上這寬恕的禱告,求主幫助她不要被忿怒去催殘她屬靈的生命.每天親愛的神給她這寬恕的恩典,叫她能夠繼續的前進.雖然這兩個兇徒還是找不著逃之夭夭.但這婦人每天學習寬恕的功課,學習意志的愛,她的生命是沒有被催殘的,她的生命反而得著釋放.

我有一位很好的朋友,她也有類似生命的經歷.在起初的時候,她認識了一個朋友,這朋友是很有恩賜的,樣貌也不錯,又能幹,又能寫,又能講.兩人惺惺相惜,但過了一段日子,她發現原來他是個偽君子,他其實不是那麼有學問的,他只是一個很貪心的人,很喜歡跟有錢的女士交往,希望能夠得到些好處.我的朋友是很不喜歡這個人的.所以你看到,當你對一個人不喜歡的時候,對方也會知道你是不喜歡他的.於是到一個地步,大家可以在同一部電梯時,大家都不看對方的.其實在她的生命裡,她是很痛苦的,因為她知道愛是重要的,她知道是要寬恕七十個七次的,但在另一方面,據實對方是一個錯的人,對方是一個不好的人,對方是一個惡的人,對方是一個偽君子.所以我這位朋友一方面想去學習去寬恕,另一方面想去証明那人是一個偽君子,是一個偽善的人.我朋友本該她的傷口是痊癒的,可以從新開始的,但幾天後,她又把傷口打開看再蓋上,不斷的打開又蓋上,讓其他人看看,傷得多利害,於是又蓋上.我朋友的生命是極其的痛苦.這好比一個人,他心靈裡有很多很多的雜聲,沒有真正的平安,她看到樓台,鐘台有很多的鐘聲在響的時候,她就跑上去鐘台,去制止這鐘聲,於是用手拉著鐘的繩子,要制止這鐘聲,但心靈的鐘聲繼續的在響起.有一天,當她發覺,她願意在這鐘台上將她手於下的時候,不拿著鐘的繩子時,ｌｅｔ　ｇｏ放手的時候,她發覺心靈的雜聲越來越小,越來越沒有.經過十年的時間,我這位朋友學習了寬恕的功課,生命能夠得著釋放,能夠邁向完全,邁向完滿.

弟兄姊妹,完滿的生命,完全的愛,是需要操練的,需要學習的.一個最好的指標是在你生命中,去學習愛一些你不喜歡的人,甚至是你的敵人.除了要有屬靈的朋友,除了在教會當中,去愛一些我們不喜歡的人,一些頂撞我們的人,甚至是敵擋我們的人,因為這樣的原故,我們可以經歷到神的愛去跟大家分享.我們要操練意志的愛,也包括愛那些陌生的罪人.

希伯來書講到:不可忘記用愛心接待客旅,因為曾有接待客旅的,不知不覺就接待了天使.另一個翻譯是: 不可忘記用愛心去接待陌生人,因為曾有接待陌生人的,不知不覺就接待了天使.在愛的實踐中,如果我們是接待那些陌生人的罪人的時候,我們發覺神是與我們同在,並在當中為我們加油的.

在美國一間浸會大學Ｅａｓｔｅｒｎ　Ｕｎｉｖｅｒｓｉｔｙ,有一位教授,叫做Ｐｒｏｆｅｓｓｏｒ　Ｔｏｎｎｙ　Ｋａｎｐｏｌｏ,有一次他出差到海地工作,在他離開的最後一天晚上,他在所住之假日酒店的大堂,遇見三個少女,最年長的是十五歲,那少女她看到Ｔｏｎｎｙ　Ｋａｎｐｏｌｏ,就走過去跟他說:十塊錢一個晚上,你明白我的意思,於是Ｔｏｎｎｙ　Ｋａｎｐｏｌｏ看著另外一位少女,你是不是也是十塊阿,再問第三個少女,但她用很輕蔑的態度,表示跟她們都是同一個價錢.於是Ｔｏｎｎｙ　Ｋａｎｐｏｌｏ告訴她們房間的號碼,叫她們十分鐘後到他房間.當Ｔｏｎｎｙ　Ｋａｎｐｏｌｏ到他房間的時候,他就打電話給酒店的管理員,請他們預備最好的雪糕,和找一些dik 士尼的影片.十分鐘之後,雪糕和影片送到,女孩子也都來了,他們四個人在房間裡面,一同吃雪糕,一同的看影片,是非常的開心,到一點鐘的時候,這三位少女都很累了,就很開心的在床上睡著了.

Ｔｏｎｎｙ　Ｋａｎｐｏｌｏ看著這些少女,沒有什麼會改變的,因為這樣的少女孩,明天起來的時候,她們同樣的被那些壞男人,污穢不堪的男人去糟塌,不會有改變的,真是沒什麼改變的,但讓我們有一天的晚上,能夠做回純真的少女,都是好的.對Ｔｏｎｎｙ　Ｋａｎｐｏｌｏ來說,愛那些陌生人,愛那些陌生的妓女,是非常重要的.

弟兄姊妹,這就是完全的愛,這就是成長的愛,這些就是能夠使生命完滿的愛,這樣的愛是主給我們的恩典.當我們有這樣愛的時候,我們就能不斷的成長.

我用另一個圖畫跟大家分享這屬靈的道理,思想崇基學院禮拜堂的弟兄姊妹們,你們圍成一個圓圈,而每一個人有一個距離,以這中心是耶穌基督,當你們每一位願意向旁邊的人走近一步的時候,這圓圈會怎樣?會縮小了,每一個人是更加的靠近這個圓心,更加靠近耶穌基督.當每個人願意向圓心走近一步時,靠近耶穌基督,那圓圈就會縮小,而每個人的距離又會靠近一點,這是一個屬靈的道理.我希望在今早的崇拜中,大家都有好的耳朵,能夠聽得進去,能夠明白,能夠記得,能夠相信,就讓我們在這裡,一同的操練,學習去學盡心的愛,以至我們的生命能夠完全.

另一個角度來說,當我們每一個人,願意這麼的單純聽耶穌基督的教導,聽耶穌基督的命令,聽耶穌基督的藥方,我們是盡心的愛,努力的去愛,有意志去愛的時候.弟兄姊妹,你們留意一下,你發覺你生命的傷口,越來越小.你試想一個人,如果他是肯去愛,在主的恩典裡,用意志去愛的時候,這些人還會不會驕傲?這些人還會不會嫉妒?這些人還會不會貪心?這些人還會不會忿怒?你會發覺當人願意,聽耶穌基督的命令去愛的時候,很奇妙的,那些生命的傷口,他是不知不覺的會挪開.深盼我們每一個人,繼續的去聽主耶穌基督的命令,我們一同去愛,能夠一同實踐生命完滿的愛.



\newpage

\section{如果愛是真的}
\label{sec:58}
\textbf{林豪恩先生 - 如果愛是真的}
\newline
\newline
\hyperref[sec:57]{< < < PREV SERMON < < <}
~
\hyperref[sec:toc]{[返目錄]}
~
\hyperref[sec:59]{> > > NEXT SERMON > > >}


路加福音 10:25-37
\newline
\begin{longtable}{cl}
\hline
\hline
章節 & 經文 (和合本修訂版)\\
\hline
10:25 & \begin{tabularx}{0.7\textwidth}{X} 有一個律法師起來試探耶穌，說：「老師！我該做甚麼才可以承受永生？」 \end{tabularx} \\ \\ \relax
10:26 & \begin{tabularx}{0.7\textwidth}{X} 耶穌對他說：「律法上寫的是甚麼？你是怎樣念的呢？」 \end{tabularx} \\ \\ \relax
10:27 & \begin{tabularx}{0.7\textwidth}{X} 他回答說：「你要盡心、盡性、盡力、盡意愛主—你的神，又要愛鄰如己。」 \end{tabularx} \\ \\ \relax
10:28 & \begin{tabularx}{0.7\textwidth}{X} 耶穌對他說：「你回答得正確，你這樣做就會得永生。」 \end{tabularx} \\ \\ \relax
10:29 & \begin{tabularx}{0.7\textwidth}{X} 那人要證明自己有理，就對耶穌說：「誰是我的鄰舍呢？」 \end{tabularx} \\ \\ \relax
10:30 & \begin{tabularx}{0.7\textwidth}{X} 耶穌回答：「有一個人從耶路撒冷下耶利哥去，落在強盜手中。他們剝去他的衣裳，把他打個半死，丟下他走了。 \end{tabularx} \\ \\ \relax
10:31 & \begin{tabularx}{0.7\textwidth}{X} 偶然有一個祭司從那條路下來，看見他就從另一邊過去了。 \end{tabularx} \\ \\ \relax
10:32 & \begin{tabularx}{0.7\textwidth}{X} 又有一個利未人來到那裡，看見他，也照樣從另一邊過去了。 \end{tabularx} \\ \\ \relax
10:33 & \begin{tabularx}{0.7\textwidth}{X} 可是，有一個撒瑪利亞人路過那裡，看見他就動了慈心， \end{tabularx} \\ \\ \relax
10:34 & \begin{tabularx}{0.7\textwidth}{X} 上前用油和酒倒在他的傷處，包裹好了，扶他騎上自己的牲口，帶他到旅店裡去，照應他。 \end{tabularx} \\ \\ \relax
10:35 & \begin{tabularx}{0.7\textwidth}{X} 第二天，他拿出兩個銀幣來，交給店主，說：『請你照應他，額外的費用，我回來時會還你。』 \end{tabularx} \\ \\ \relax
10:36 & \begin{tabularx}{0.7\textwidth}{X} 你想，這三個人哪一個是落在強盜手中那人的鄰舍呢？」 \end{tabularx} \\ \\ \relax
10:37 & \begin{tabularx}{0.7\textwidth}{X} 他說：「是憐憫他的。」耶穌對他說：「你去，照樣做吧！」 \end{tabularx} \\ \\
[1ex]
\hline
\hline
\end{longtable}
日期:2010年8月8日

題目:如果愛是真的

經文:路加福音10:25-37

証道:林豪恩先生

經文:路加福音10:25-37

25有一個律法師起來試探耶穌,說:「夫子!我該做甚麼才可以承受永生?」

26耶穌說:「律法上寫的是甚麼?你念的是甚麼呢?」

27他回答說:「你要盡心,盡性,盡力,盡意愛主──你的神；又要愛鄰捨如同自己.」

28耶穌說:「你回答的是；你這樣行,就必得永生.」

29那人要顯明自己有理,就對耶穌說:「誰是我的鄰捨呢?」

30耶穌回答說:「有一個人從耶路撒冷下到耶利哥去,落在強盜手中.他們剝去他的衣裳,把他打個半死,就丟下他走了.

31偶然有一個祭司從這條路下來,看見他就從那邊過去了.

32又有一個利未人來到這地方,看見他,也照樣從那邊過去了.

33惟有一個撒瑪利亞人行路來到那裡,看見他就動了慈心,

34上前用油和酒倒在他的傷處,包裹好了,扶他騎上自己的牲口,帶到店裡去照應他.

35第二天拿出二錢銀子來,交給店主,說:「你且照應他；此外所費用的,我回來必還你.」

36你想,這三個人哪一個是落在強盜手中的鄰捨呢?

37他說:「是憐憫他的.」耶穌說:「你去照樣行吧.」

釋經與應用

1．愛──永恆的主題

耶穌所講的這個「好撒馬利亞人」故事大家都耳熟能詳,「撒馬利亞人」也成為了通用的詞彙,香港一個防止自殺的機構也採用了「撒馬利亞」為其名字.

根據經文記載,耶穌講出「好撒馬利亞人」這個故事,是因為有一位律法師前來以試探的心態問了一個問題:「夫子!我該做甚麼才可以承受永生?」耶穌以問代 答,律法師回答律法的要求是:「你要盡心,盡性,盡力,盡意愛主──你的神；又要愛鄰捨如同自己.」耶穌肯定律法師的答案.不過,律法師卻糾纏下去,提 出:「誰是我的鄰捨?」這個問題,於是耶穌就講出了「好撒馬利亞人」這個故事.

從律法師與耶穌的對答中,雙方都承認「愛上帝及愛鄰捨如同自己」與「得著永生」有著密切的關係.基督徒普遍認同愛是信仰的核心,社會上絕大部份人也同意愛 是美好東西.愛,實在是永恆的追求.從前人們比較含蓄,愛,往往只存於心而不宣之於口.有一首我父親年代的情歌,取自元朝才女管道升的一首詩,名為「我儂 你儂」,這首詩完全不提「愛」字,可是卻明白地顯示出她的深切之愛.

「你濃我濃,忒煞情多,情多處深如火,

把一塊泥,撚一個你,塑一個我,

將咱倆個一齊打破,用水調和,再撚一個你,再塑一個我,

我泥中有你,你泥中有我,

我與你同生一個衾,我與你同在一個槨.」

現在人們鼓勵「我口說我心」,傾向把愛字掛在口邊.我聽歌的年代,有一首流行曲名為「愛到你發狂」,「愛的宣言」說到極致.

我愛你愛到發狂發曬狂 我愛你愛你愛你

我愛到我快發狂實發狂 我愛你愛你愛你 我愛到發狂要發狂 我愛我要愛到你發了狂 發曬狂要你愛你要你

教會也與時並進,不後於人.有些教會聚會時會作出這樣的邀請:「弟兄姊妹,向你前後左右四個人說:上帝愛你,我也愛你!」

耶穌講的這個「好撒馬利亞人」故事,也是一個愛的故事,因為這個故事是回應「愛鄰捨如同自己」說的.在這個段落裡,除了耶穌之外,有幾個重要角色,包括:律法師,祭司,利未人及撒馬利亞人.作者對這些角色的描述,提供了重要的資料給我們進一步的探索「愛為何物」.

如果愛是真的,就不會只停留於「口水戰」

故事中第一位出場的人物是律法師.律法師負責教導及解釋律例典章,因此,他們是猶太社會的律法專家.路加福音對這位律法專家的描述並不正面.對於律法師向 耶穌提出的問題:「夫子!我該做甚麼才可以承受永生?」作者在25節直接指出他的動機是「試探耶穌」,而不是真誠地來向耶穌請教永生之道；問題解答後,耶 穌吩咐他付諸實行的時候,律法師反問耶穌:「誰是我的鄰捨呢?」作者在29節又指出他的動機是「想證明自己有理」.

律法師其實知道「愛上帝及愛鄰捨如同自己」是承受永生的不二之法,他的知識與耶穌的宣講吻合.不過,律法師的來意並不是「請教」耶穌,而是「試探」耶穌, 因此繼續與耶穌爭辨下去.其實,律法師是「靠把口搵食」的,好爭辨一點也不意外,反而是耶穌吩咐把「你要盡心,盡性,盡力,盡意愛主──你的神；又要愛鄰 捨如同自己.」這所知的實行出來,就令律法師感到難為.耶穌向律法師的吩咐暗示,愛不單單是知識,也不單單是辨論,愛,是要行出來的.真理可能越辨越明, 但愛卻不會越辨越真.對於一眾「靠把口搵食」的人,這可是一項挑戰.現代社會越來越來「靠把口搵食」的行業,例如推銷員,政治人物,外交官員等.

德蘭修女曾經分享一次經歷,她說:「我在孟買有過一次不尋常的經歷.那裡正舉行一個有關飢餓問題的龐大會議.我本該出席該會,但我迷了路.不知怎樣,我後 來到達那會場.而就在那會場的門前,我發現了一個垂死的人,與會的人士就在那裡高談闊論有關食物與飢饉的問題.我把那垂死的人帶回善終會.他在那裡死去, 他死於飢餓.而在那會堂裡面的人則仍在高談闊論,講及如何在十五年內生產足夠的糧食,足夠的 and \#8230; and \#8230;,足夠的 and \#8230; and \#8230;,而那飢餓的人卻死去.」

有些事情的確需要從詳計議,不過,討論的目的是甚麼呢?為了救人於水火之中,還是為了推卸責任呢?就好像世界各國的官員開會回應全球暖化問題的時候,往往 變成誰要負多少責任的爭論.當各國代表繼續在冷氣大廳爭論不休的時候,地球上有不少人及生物正受著氣候變化所傷害.如果把會議場地移師在乾旱得龜裂的農地 上進行,或者在山洪暴發的險境中進行,與會者的爭論時間可能會大大縮短.

當耶穌吩咐律法師實踐「愛上帝及愛鄰捨如同自己」的時候,律法師反問:「誰是我的鄰捨呢?」這反問其實有一定的道理.「全世界那麼多人?我能夠去愛哪一個 呢?」「如果要選擇?為甚麼是這個,而不是那個呢?」資源有限,界定服務對象在所難免；各有分工,釐清責任,可以免除混亂.然而,界定,可以是漫長的「口 水戰」,也可以隱含自保的動機.離港不超過多少天的人才可以申請綜緩呢?離港不超過多少天的長者才合資格領取生果金呢?這些定義,或者可以等待法院澄清. 然而,若有人在醫院門外昏迷,到底是否醫院的責任呢?界定服務對象及責任,似乎在所難免,但有時卻可以變成埋沒愛心的陷阱.

聖經作者對這位律法師的負面評價,或者正在告訴我們,如果愛是真的,就不會只停留在口舌之爭,也不會被界定及責任來掩蓋愛心.

如果愛是真的,就不會只停留於「做好呢份工」

這個「愛的故事」中另外一對負面的角色,是祭司和利未人,他們都是在聖殿裡主持獻祭的神職人員.在耶穌的時代,祭司仍然是上帝與人之間的橋樑.路加福音記 載,耶穌有一位祭司親戚,他是施洗約翰的父親──撒迦利亞.路加福音第一章如此記錄祭師撒迦利亞的工作:「撒迦利亞按班次在上帝面前供祭司的職分,照祭司 的規矩掣籤,得進主殿燒香.燒香的時候,眾百姓在外面禱告.有主的使者站在香壇的右邊,向他顯現.」(路1:8-11)此外,第五章有一個段落讓我們看見 祭司重要的角色:「有一回,耶穌在一個城裡,有人滿身長了大痲瘋,看見他,就俯伏在地,求他說:「主若肯,必能叫我潔淨了.」耶穌伸手摸他,說:「我肯, 你潔淨了吧!」大痲瘋立刻離了他的身.耶穌囑咐他:「你切不可告訴人,只要去把身體給祭司察看,又要為你得了潔淨,照摩西所吩咐的獻上禮拜,對眾人所證 據.」(路5:12-14)

在耶穌目前所講的故事中,祭司和利未人都先後行走在這條從耶路撒冷下耶利哥的山路,他們都先後遇上一個被強盜打劫,連衣服也被剝去,並被打得半死躺在路上的受害者.他們共同的反應都是看了一眼就繼續上路.

在猶太人群體中,祭司及利未人被認為是接近上帝的人,他們代表人向上帝獻祭,也代表上帝宣告人得潔淨.然而,在耶穌的故事中,祭司及利未人的表現卻出乎意 料之外.為甚麼眾人以為最有可能活出「愛上帝及愛鄰捨如同自己」的宗教領袖卻置這位受害者於不顧呢?人們嘗試為他們的行動提供解釋:

有人認為,祭司及利未人可能擔心這位受害者不是律法所界定的「鄰捨」,而現場狀況又不容許他們確認對方是否受割禮的人,因此還是留給其他人協助為妙.

又有人指出,律法禁止祭司觸摸屍體,如果摸了,就是不潔淨.若沾染不潔,就不能主持獻祭,不能吃祭物,而且要行潔淨禮.祭司和利未人或者是為了潔淨的緣故而對路旁的不幸者避之則吉.

還有人認為,可能是祭司和利未人擔心劫匪還埋伏在附近.若停下來幫助受害者,自己也可能成為劫匪的目標.

上述的猜測都有可能,但也都不能被肯定,因為故事並沒有為祭司及利未人沒有救助受害者提供解釋.故事只是報告祭司及利未人看見這位受害者卻不救助的事實.故事中沒有直接批評或指責祭司和利未人的文字,但其負面的角色則不言而喻.

祭司和利未人按著他們的身份,完成了在耶路撒冷聖殿的職責,踏在回家的路上.或者他們也按著禮儀的要求,保持自己潔淨以執行聖職.不過,他們以負面的角色 出現在耶穌所講的這個故事中,彷彿告訴著人們,如果愛是真的,就不會只停留於「做好呢份工」.只是附合身份及禮儀的要求,也有可能成為埋沒愛心的陷阱.

如果愛是真的,就讓憐憫激發行動

有一個人在路上遇劫,被打得半死躺在路旁,祭司經過,沒有出手相助,利未人經過,也沒有出手相助.聽故事的律法師心裡可能在想,最有可能出手相助的人都沒 有行動,還有還會出手相助呢?會不會是一位律法師呢?出乎意料之外,這出手相助的竟然是一位撒馬利亞人.對於我們來說,這位出手相助的人是猶太人還是撒馬 利亞人是亳無分別的,但對於這位來挑戰耶穌的猶太律法師來說,可能卻感到自己的民族感情受到傷害.

撒瑪利亞人是何許人也.撒瑪利亞人認為自己的祖籍是猶太人中的嗎哪西和以法蓮支派,但在耶穌時代的猶太人的眼中,他們是「非純種」的猶太人.公元前 722 年猶太人亡國後,亞述帝國把巴比倫等外族遷進撒瑪利亞城,因此撒馬利亞人摻了外族血統.在宗教上,猶太人指責撒瑪利亞人在基利心山敬拜,捏造聖殿的發源 地,是分裂者.更因為撒馬利亞人只有摩西五經,外加自己的撒瑪利亞經,更被指為異端.在公元 6至9 年,在猶太人與撒馬利亞人之間,發生了一件火上加油的事件,撒瑪利亞人在一個逾越節的晚上,把死人骨頭丟進聖殿的院子裡.從此,猶太人和撒瑪利亞人更是勢 不兩立了!

在耶穌的故事中,這個「非純種」猶太人,這個分裂者,這個異端,竟然以最正面的角色出現.大家比較看看耶穌講這個故事時,花了多了篇幅在祭司身上,在利未 人身上,以及在這位撒馬利亞人身上.耶穌只用了一句:<font color="\#0000ff">「31偶然有一個祭司從這條路下來,看見他就從那邊過去了.」來提及祭司,又用了相似的一句來提及利 未人:「32又有一個利未人來到這地方,看見他,也照樣從那邊過去了.」但耶穌卻仔細地描述這位撒馬利亞人救助受害者的行動,並為他加上對話:「33惟有 一個撒瑪利亞人行路來到那裡,看見他就動了慈心,34上前用油和酒倒在他的傷處,包裹好了,扶他騎上自己的牲口,帶到店裡去照應他.35第二天拿出二錢銀 子來,交給店主,說:「你且照應他；此外所費用的,我回來必還你.」」

這位撒馬利亞人「把油和酒倒在他的傷處」,這行動好像一個獻祭的動作.油和酒是在聖殿敬拜時用的獻祭材料,「倒」也是一個敬拜或獻祭的用詞.這應該是祭司 的行動,然而在有需要的人面前,行出來的不是祭司,而是撒馬利亞人.獻祭本來就是象徵人對上帝的慈愛的正確回應,撒瑪利亞人在這位不幸者的身上獻上了聖潔 及神所喜悅的祭.

為甚麼是撒馬利亞人呢?撒馬利亞人從哪裡學識「愛上帝及愛鄰捨如同自己」的呢?更重要的是,撒馬利亞人何以能活出「愛上帝及愛鄰捨如同自己」的吩咐呢?耶 穌指出,這位撒馬利亞人對受害者動了「慈心」(新譯本:「憐憫之心」,呂振中譯本:「憐憫的心」).這位撒馬利亞人的憐憫之心,連律法師也感受得到.當耶 穌問他:「這三個人哪一個是落在強盜手中的鄰捨呢?」,律法師回答:「是憐憫他的.」因為動了慈心,這位撒馬利亞人停了下來看顧受害者；因為動了慈心,這 位撒馬利亞人動手替受害者包扎傷口；因為動了慈心,這位撒馬利亞人用自己的牲口來運送傷者；因為動了慈心,這位撒馬利亞人為傷者尋找棲身之所並徹夜照顧 他；因為動了慈心,這位撒馬利亞人自掏腰包為客店照顧受害者的費用「包底」.

孟子曰:「惻隱之心, 人皆有之」.任何人看見一個小孩爬近井邊,在這生死關頭,都不會考慮那小孩是不是仇人的兒子,也不會考慮名譽地位等問題,而會第一時間去拯救那個小孩離開 危險.孟子進一步指出,有「不忍人之心」,就會有「不忍人之政」.就算在現代社會,為政者以不忍心令人民受苦的態度施政,仍然是人民的期望.世事紛亂,人 心險詐,憐憫的心很容易被擠壓掩蓋.其實一般人民盼望的,只不過是安居樂業,值得為了追求某種政治理想而令人民受苦,甚至受苦到一個程度連做手術時到用不 到麻醉藥嗎?一般人民盼望的,只不過是父慈子孝,平安度日,值得為了經濟增長,推高GDP,而令嬰孩因奶粉而受苦,父親在礦場受難呢?一般人民出外旅遊, 無非都是尋找快樂,導遊因目標未達到而心情欠佳,但何必為了發洩一己的情緒而謾罵及恐嚇遊客,掃遊客的興呢?如果主事者都好像故事中這位撒馬利亞人一樣, 動了慈心,就不會輕易做出令人受苦的事情.

總結

故事中這位被不幸者動了慈心的撒馬利亞人,有沒有令我們想到耶穌自己呢?福音書記載耶穌經常被人動了他的慈心,因而採取行動.剛才我們提及耶穌醫治長大痲 瘋的人的事件,馬<font>可<font>福音(1:40-41)是這樣記載的:<font color="\#0000ff">「有一個長大痲瘋的來求耶穌,向他跪下,說:「你若肯,必能叫我潔淨了.」耶穌動了慈心,就伸手 摸他,說:「我肯,你潔淨了吧!」」馬太福音綜合地說(9:35-36):「耶穌走遍各城各鄉,在會堂裡教訓人,宣講天國的福音,又醫治各樣的病症.他看 見許多的人,就憐憫他們；因為他們困苦流離,如同羊沒有牧人一般.」

或者,我們有時都好像那位被打得半死的受害者一樣,躺在路旁等待憐憫.七月初受洗的一位弟兄的水禮見證按「好撒馬利亞人」的故事結構可以如此演繹:

「有一個人正值精壯之年,他的運動神經系統突然出現退化,自此右手完全不能發力活動.有一個醫生替他檢查,告訴他無法醫治；又有一個醫生替他檢查,告訴他 無藥可救.這人的意志一直消沈,感到極度驚恐,頓覺人生絕望.他的太太的二姐和姐夫帶她參加主日崇拜,在那裡他重新認識主耶穌,感到異常平靜,身心靈都感 受到平和而實在,使他漸漸走出消極的陰霾.」(參崇基禮拜堂水禮見證集──樂俊偉先生)

讓我們回到律法師的第一個問題:「我該做甚麼才可以承受永生呢?」

「憐憫人的人有福了!因為他們必蒙憐憫.」



\newpage

\section{信心與權柄}
\label{sec:59}
\textbf{賀志勇博士 - 信心與權柄}
\newline
\newline
\hyperref[sec:58]{< < < PREV SERMON < < <}
~
\hyperref[sec:toc]{[返目錄]}
~
\hyperref[sec:60]{> > > NEXT SERMON > > >}


路加福音 7:1-10
\newline
\begin{longtable}{cl}
\hline
\hline
章節 & 經文 (和合本修訂版)\\
\hline
7:1 & \begin{tabularx}{0.7\textwidth}{X} 耶穌對百姓講完了這一切的話，就進了迦百農。 \end{tabularx} \\ \\ \relax
7:2 & \begin{tabularx}{0.7\textwidth}{X} 有一個百夫長所器重的僕人害病，快要死了。 \end{tabularx} \\ \\ \relax
7:3 & \begin{tabularx}{0.7\textwidth}{X} 百夫長風聞耶穌的事，就託猶太人的幾個長老去求耶穌來救他的僕人。 \end{tabularx} \\ \\ \relax
7:4 & \begin{tabularx}{0.7\textwidth}{X} 他們到了耶穌那裡，切切地求他說：「你為他做這事是他配得的； \end{tabularx} \\ \\ \relax
7:5 & \begin{tabularx}{0.7\textwidth}{X} 因為他愛我們的民族，為我們建造會堂。」 \end{tabularx} \\ \\ \relax
7:6 & \begin{tabularx}{0.7\textwidth}{X} 耶穌就和他們同去。離那家不遠，百夫長託幾個朋友去見耶穌，對他說：「主啊，不必勞駕，因你到舍下來，我不敢當。 \end{tabularx} \\ \\ \relax
7:7 & \begin{tabularx}{0.7\textwidth}{X} 我也自以為不配去見你，只要你說一句話，就會讓我的僮僕得痊癒。 \end{tabularx} \\ \\ \relax
7:8 & \begin{tabularx}{0.7\textwidth}{X} 因為我被派在人的權下，也有兵在我之下。我對這個說：『去！』他就去；對那個說：『來！』他就來；對我的僕人說：『做這事！』他就去做。」 \end{tabularx} \\ \\ \relax
7:9 & \begin{tabularx}{0.7\textwidth}{X} 耶穌聽到這些話，就很驚訝，轉身對跟隨的眾人說：「我告訴你們，這麼大的信心，就是在以色列，我也沒有見過。」 \end{tabularx} \\ \\ \relax
7:10 & \begin{tabularx}{0.7\textwidth}{X} 那差來的人回到百夫長家裡，發現僕人已經好了。 \end{tabularx} \\ \\
[1ex]
\hline
\hline
\end{longtable}
日期:2010年8月15日

主題:信心與權柄

証道:賀志勇博士

經文:路加福音7章1-10節

耶穌對百姓講完了這一切的話,就進了迦百農. 有一個百夫長所寶貴的僕人,害病快要死了. 百夫長風聞耶穌的事,就托猶太人的幾個長老,去求耶穌來救他的僕人.他們到了耶穌那裡,就切切的求他說,你給他行這事,是他所配得的．因為他愛我們的百姓,給我們建造會堂.耶穌就和他們同去.離那家不遠,百夫長托幾個朋友去見耶穌,對他說,主阿,不要勞動．因你到我捨下,我不敢當．我也自以為不配去見你,只要你說一句話,我的僕人就必好了.因為我在人的權下,也有兵在我以下,對這個說去,他就去．對那個說來,他就來．對我的僕人說,你作這事,他就去作.耶穌聽見這話,就希奇他,轉身對跟隨的眾人說,我告訴你們,這麼大的信心,就是在以色列中我也沒有遇見過. 那托來的人回到百夫長家裡,看見僕人已經好了.(路加福音7章1-10節)

引言--權柄濫用的可怕最近看了電影《葉問》,打得很好看；聽說還出了一部《葉問前傳》,打得也很好看.我想,中國人喜歡看武俠片,香港拍得最好的電影之一,就是《黃飛鴻》系列.我小學就喜歡看金庸、古龍的小說,那時候,最羡慕大俠了,倒不是羡慕大俠身邊有很多美女,那時我才8、9歲,美女跟我同桌的女生有什麼不同,我也不大分得清.那時候我最羡慕大俠是:第一,大俠到哪里吃飯好像都不要給錢,吃得很好,身上還從來不缺銀子；第二,大俠輕功很好,飛來飛去,過海關不用回鄉證通行證.第三,大俠想救誰就救誰,想殺誰就誰.生殺予奪,在我一人之手.那種感覺多好啊,中學時候被化學老師罵,我當時想,如果他再罵我,我就去找九陰真經,然後給他嘗一下我的九陰白骨爪.當然他再也沒有罵我,沒有給他惹來殺身之禍.

我想很多人小時候都做過大俠夢吧.尤其是最後一點,別人的生殺大權都在大俠手中,這種感覺非常之好.現在回過頭來,我發現大俠夢,其實跟皇帝夢是同一種類的夢.做大俠其實跟做皇帝一樣,都擁有不受約束的絕對權力.想做皇帝的人都以為自己會做一個好皇帝,但天下有幾個好皇帝?就算是李世民,殺起對手來也絕不手軟.做大俠也一樣啊,無論他是什麼大俠,他可以殺得了別人,而別人卻殺不了他.很多人的大俠夢,骨子裡其實就是一種絕對的、不受約束的權力夢.

但現實情況是,當有人擁有絕對權力時,是一件好恐怖好恐怖、好讓人驚恐的事情.英文有句話說:"absolute power corrupt absolutely!"絕對的權力絕對會將人腐化.在不同的戰爭裡頭,這些權柄的悲劇一幕一幕重現,幾千年了,絲毫不變.我記得德國人拉貝在他的日記裡頭,記錄了日本人南京大屠殺.日本兵以殺人來取樂,比賽誰殺中國人殺得多；他們為了打賭女人肚子裡懷的是男孩還是女孩,就把孕婦的肚子刺開來看看,究竟是男的還是女的.權柄,絕對的權柄,很容易就讓人變成了野獸.當一個人做什麼事都不需要負責任時,是一件很恐怖很恐怖的事情.

不單止是德國人屠殺猶太人,不單止是日本人屠殺中國人.你看看今日的世界,911事件,美軍攻打伊拉克事件,歷史總是在重演.而且,不單止在戰爭中間,你看看今日的世界,在商場,在家庭,在公司何嘗不是一樣了?今天,被虐待的小孩與老人不在少數；今天,被壓逼的雇員也不在少數.我看到最近的一個社會調查說,香港虐待兒童和老人的數量比前兩年多了一倍--權柄濫用,會造成多麼多麼可怕的後果.

2,百夫長的權柄但今天我們再看一個故事,一個很美麗的故事,講的是一個很有權柄的人,但是因為他敬畏上帝,就完全不同了.首先我們來看看路加7章第一節至第二節.

(255, 255, 255)耶穌對百姓講完了這一切的話,就進了迦百農.有一個百夫長所寶貴的僕人,害病快要死了.百夫長是什麼人呢?第八節說,百夫長讓人去就去,來就來；讓人做事,就要去做事.很明顯,百夫長是一個很有權柄的人.他命令手下,去,就要去；來,就要來；坐在那裡,就要坐在那裡啦；端杯水給我,就端杯水來；斟杯茶,就斟杯茶來；拿個拖鞋來,就拿著拖鞋過來.好像叫婢女(妹仔)一樣.他是有權柄的人,不得違命.他命令一發出,手下就要去做.在羅馬的兵制裡邊,百夫長是軍隊的骨幹.他通常是管理一百個兵,但到耶穌時期,他大概要管理1000個左右.歷史學家潑利比爾斯這樣描述一個羅馬的百夫長:"百夫長是一個危險的人,因為他很可能貪功,危及手下的安全；很可能貪生怕死,視手下的生命如芥草...."所以一個兵跟了個不好的百夫長,分分鐘都可能丟掉性命.因為百夫長要他去這裡,他就要去這裡；要他去那裡,他就要去那裡,萬一百夫長指揮錯了,他就會去送死.羅馬的百夫長是整個軍隊裡頭的支柱.他擁有絕對的權柄,用今天的軍制來講,他相當於今天的少校.放到公司裡頭,他就是個基金經理.

很有趣的是,聖經裡面幾乎每次提到羅馬的百夫長都是好人.路加福音23章47節,使徒行傳20章2節,22章26節,23章27節,在新約聖經,所有關於羅馬百夫長的經文都說百夫長好,好奇怪哦.我不是說所有羅馬百夫長都是好人,但是在聖經裡邊記載的羅馬百夫長都是好人.今天我們看的這個,是在迦百農作百夫長.迦百農是羅馬人佔領的地方,就好像今天的殖民地一樣,是直接歸羅馬人統治的地方.大部分不是羅馬人,而是佔領區的亡國奴.對於一個羅馬的百夫長來講,這種地方更加是可以為所欲為的地方.如果百夫長對著的是羅馬公民,他還會有所顧忌；但是他現在對著的是一班亡國奴,就好像德國人對著佔領區的猶太人一樣,他其實想做什麼就可以做什麼.在迦百農,他有絕對的權柄,是個危險人物,隨時可以搞出人命的,而不需要付任何責任.

但很奇怪,我們看第二節,說百夫長所寶貴的僕人害病就要死了.僕人這個字,希臘文有兩個詞都可以表示僕人:第一個意思了就是請回來的,你的菲傭,你的印傭來的,是被請回來的.他隨時可以走人,隨時可以說,我賠錢給你,我不幹啦,走了啊.他是一個自由人,free man來的.她不是你的奴隸,不過是同你打工罷了；但另外一個詞的意思,就是,她是你買回來的,她沒有自己的主權,沒有生命的權柄,主人要她死她就要去死,要她嫁就嫁,主人對她了有絕對的權柄.當年香港1920年代,有英國人告香港紳士家裡養女僕(妹仔),就是養奴隸；香港人不服氣,說妹仔不是奴隸.於是在法院裡頭吵了一大輪.結果香港人輸了,不得不取消了妹仔這個制度.聖經裡頭的僕人就是第二個意思,就是奴隸的意思.在當時的羅馬社會,這些奴隸是沒有任何權力的.

你想想,如果在今天的社會裡頭,我們擁有這麼大的權柄,好像百夫長一樣,我們會怎樣待我們的手下呢?我們會怎樣看他們呢? 我記得前兩年我在一個朋友家裡小住,他們家請了個菲傭打理一下家務.朋友說,住在這裡不要客氣啊!我剛開始還有點客氣,過了兩天,我就開始放鬆了.一邊在那裡看書,一邊喊菲傭"julia,給我倒杯水來；"看電視的時候,我又懶嘛,就喊"Julia,遙控在哪里啊?給我找一下遙控過來."剛一說完,我心裡也有些不好意思,心想,人家請的傭人,又不是我自己請的,使人不使本麼?多不好意思啊.我再轉頭一想,我不過這點權柄,就這樣使喚人,這個德性了,要是我有絕對的權柄,叫人做什麼就做什麼,還怎麼得了?權柄是個好危險的東西.

3,百夫長的敬畏神但是我們再看看這個百夫長,我們看看三到五節.

百夫長風聞耶穌的事,就托猶太人的幾個長老,去求耶穌來救他的僕人.他們到了耶穌那裡,就切切的求他說,你給他行這事,是他所配得的．因為他愛我們的百姓,給我們建造會堂.

聖經這樣形容這個百夫長,他雖然掌握著迦百農一帶的猶太人的生殺大權,但卻寶貴他的僕人,而且很愛猶太的百姓.在希臘文當中,這個愛Agapaw是一種犧牲的愛的意思,就是非常非常的愛.他怎麼愛的呢?他在迦百農起了個猶太會堂啊.耶穌剛剛出來傳道,就在迦百農會堂那邊,那個會堂就是這個百夫長起的.那個會堂的遺址現在還在,在加利利海邊,雖然破敗了,可是還是可以看得出來,相當雄偉.為什麼百夫長要建會堂呢?我相信,他可能受了猶太教的一神教影響.羅馬人拜很多神的,滿天神佛.很可能這個百夫長是個很有道德的人,很可能他覺得羅馬的宗教太過淫亂,他很嚮往猶太人那種道德觀念,很嚮往猶太人一神教的宗教思想,甚至有可能他已經歸化猶太教.這些我不能確定,但最低限度,他是一個敬畏上帝的人.聖經裡頭經常用一個字"怕神"來形容這樣的人,縱然他不是一個猶太教徒,他起碼是個怕上帝的人,是一個敬畏上帝的人.他不敢亂來的,因為他怕上帝.

這是一件很要緊很要緊的事.當你做一件事,你知道"哇,我不可以亂來的",這是一件很重要的事.Charles Windoor講過一個故事:當他18歲那年,他考到車牌.他沒考到車牌之前,學車的時候,他老爸就坐在他身邊,他一開快車,他老爸就說:"莫開快車,莫開快車!"在身邊說三道四,攔頭阻尾的.如今他考到牌了,阿爸又不在身邊.現在自由啦,哇,多爽啊!於是把車開到高速公路,嘩嘩聲開得不知有多響.身邊沒人吵他了,自在了.當他正要不停的踩油門的時候,有個聲音在他耳邊響起:"你以為上帝是瞎子啊?上帝看不到你啊?上帝不知道啊?"他這麼一想,就不敢去踩油門了.他把這個事就講給老爸聽,老爸說:"好嘛,感謝上帝!我不可能天天24小時監督你,把你交給上帝,上帝24小時看住你,多好啊."我想我們做爸爸媽媽的都是這樣,你天天看住小孩子做功課,做這個做那個,你能看得住多少呢?他一出去你就看不住了.將他交給上帝,由上帝來看住,多好啊.

4,敬畏神與尊重人怕神--一個人最重要是敬畏神,一個敬畏神的人是不會亂來的.這個羅馬的百夫長就是這樣啊,他怕神,他是個敬畏神的人.因為他知道我們有一天都要向上帝交賬,所以他敬畏上帝,所以他循規蹈矩,寶貴他的僕人,又愛他的百姓.這是件很有趣的事,一個敬畏神的人,他自然就會去尊重別人.一個不敬畏神的人,他就完全不管人家的感受.我最大嘛,我想怎麼用你就怎麼用你啦."啊,Julia,端杯水來!啊,Julia,拿拖鞋過來!"用人可以用到盡的.但是當你敬畏上帝的時候,你會覺得:這好像不太好吧.人家是傭人,又不是女僕(妹仔).自己的事情還是要自己做吧.當我們怕神的時候,我們就會尊重人.尊重人,我們就不會在大浪西灣鏟平樹林,不管人家抗議不抗議.百夫長也是啊,因為他怕神,所以他就愛、就尊重那些毫無地位的亡國奴猶太人.

事實上,敬畏神十分十分重要.所羅門王就說:"敬畏耶和華是智慧的開端,認識至聖者便是聰明!"對我們基督徒也是一樣,什麼是基督徒呢?基督徒就是敬畏神的人.敬畏神就是敬畏神的權柄,當神就在這裡.(當神有到)你莫當神是死的,你莫當神是聾的,你莫當神是盲的.你每一句話每一件事,人家聽不到,上帝聽到!人家看不到,上帝看到!一個敬畏上帝的人,他知道他哪怕把自己都騙了,都騙不到上帝.無論是善是惡,上帝都必審問.

有一個探險隊7個人去北極探險,迷了路.走了幾夜,最後只剩下一個麵包.大家商量先睡覺,第二天大家分著吃,看看能不能走出去.到了晚上,有個人餓得實在睡不著,就爬到隊長那邊,偷到了麵包.他正要吃,不幸他是個基督徒,基督徒要做謝飯禱告嘛,他就禱告說"感謝上帝讓我偷到麵包".這一感謝,心想不對哦,這樣子不行啊.於是又放下麵包.但是實在肚子餓,他又爬過去,剛偷出麵包,又想起了上帝,只好把麵包又放回去.這樣子半個小時都沒吃到麵包.第二天,他跟大家講這件事,大家都說"幸虧你是個基督徒,你要是個共產主義者,我們就完蛋了."靠著那一點麵包,這7個人終於等到了救兵.

敬畏神是什麼?敬畏神就是知道神的權柄是什麼!我們以為神不知鬼不覺,當我們報稅的時候；當我們恨別人的時候；當我們罵小孩的時候；當我們有很多私心雜念的時候,我們以為沒人知道,我們當上帝是死的,是盲的,是聾的.但上帝其實都知道.如果這個世界都敬畏上帝的權柄,就天下太平了.

5,百夫長的信心這個百夫長除了敬畏神之外,還有另一個驚人的表現,就是他驚人的信心.我們看看6-10節就可以看到了.這個百夫長不但不敢輕易勞動耶穌,甚至覺得他都不配見耶穌.人的政治就好像猶太長老那個樣子,這個長老說"百夫長配得上耶穌出手幫忙."人也是這樣子啊,你要找人寫推薦信進拔萃中學,推薦人肯定要誇你智商200以上,琴棋書畫樣樣都好,還熱愛公益事業,是配得上拔萃中學的.你不會說這個孩子雖然樣樣都好,不過這些都算不得什麼,不過是神的恩典罷啦.我們找工作也是啊,儘量要誇自己多麼優秀,多麼能幹,有多麼多麼豐富的經驗,意思就是,我是配得上你給的這份工的.但這個百夫長卻恰恰相反,耶穌來到的時候,他卻說,"我不配","我不敢當","雖然我有權,叫這個叫那個,做這個做那個,但是你的權更加大."奇怪哦!

不過我轉念一想,我就想到什麼是祈禱,什麼是禱告.向耶穌求就是禱告了,我不知我們中間有幾個人禱告說,我是配得上你幫我的,耶穌啊!我昨天晚上預備了8個鐘頭講道,耶穌啊,我配得上你讓我把它講好啊.有沒有人這樣祈禱啊?不會啊,人會怎麼祈禱?就好像百夫長這樣祈禱.Helesibe說:"祈禱最重要有兩樣東西:第一,你有信心；第二,你知道自己搞不定.什麼叫信耶穌?信耶穌就是知道自己不行,所以就求耶穌了.你相信耶穌有權柄,搞得定,所以你信他.什麼叫祈禱?祈禱就是發現自己不配,來到神面說,主你幫下我,我是不配的,你是配的.

6,信心與權柄我們是不配的,我們是軟弱的,神的大能就在人的軟弱上顯得完全.這個就是信心.對神有信心,就看到自己的軟弱.你看一下這個百夫長,雖然他知道羅馬凱撒大帝的權柄,他也知道羅馬帝國賦給他的權柄,但他卻認為耶穌基督的權柄比這些權柄都更大,大得多,大得讓他覺得請耶穌到家裡來都不配."沒錯,耶穌基督,我有權柄可以叫人做這個做那個,但是我戰勝不了死亡,而耶穌基督,你才能戰勝死亡."他深深明白自己的卑微,雖然在人的面前他似乎滿有權柄,但他知道在神面前,他是完全不配.

基督教就是這麼奇怪,你越是認清楚神的權柄,越是怕神,越是覺得自己不配,就越是有信心.基督教的權柄就是受苦的十字架,就是耶穌幫門徒洗腳的大毛巾,就是耶穌基督對天父的順服；而基督教的信心不是來自刀,不是來自槍,不是來自銀行裡的存款；卻反而是來自這十字架,這洗腳的大毛巾,來自自己覺得不配,來自對神的順服.

百夫長認清了神的權柄,所以怕神,覺得自己不配,所以他得了神的大恩典.他是個羅馬人,更是一個神的子民.今天,我們是一個羅馬人,還是一個神的子民呢?我們有沒有怕錯對象?--怕不能出人頭地,怕強積金縮水,怕找不到工作,怕健康不好,怕子女不聽話,怕蓋房子別人來抗議.有這些擔心,你不如好好去怕神,去敬畏神.到那時,我們就會發現自己完全不同了.因為信心的含義在於:人的盡頭,就是神的開始.



\newpage

\section{窮得祇剩下錢}
\label{sec:60}
\textbf{伍渭文牧師 - 窮得祇剩下錢}
\newline
\newline
\hyperref[sec:59]{< < < PREV SERMON < < <}
~
\hyperref[sec:toc]{[返目錄]}
~
\hyperref[sec:61]{> > > NEXT SERMON > > >}


日期:2010年8月22日

主題:窮得祇剩下錢

証道:伍渭文牧師

經文:路12:13-21

2009年六月一期《紐約時報》刊出一幅漫畫,主人公是一位繼承了巨額遺產的富家公子駕駛著遊艇,遊艇上還停著私人直升機.他洋洋得意地對身邊的秘書說:"我每天晚上都在想,自己 是多麼幸運啊--因為我的爸爸不是沃倫-巴菲特(Warren Biffett)!"

2007年蓋茲夫婦(Bill  and amp; Melinda Gates)與巴菲特開始推行捐獻承諾運動,敦促億萬富豪把至少一半財產捐出來做慈善事業.他們請富豪撥出時間想清楚他們的後代需要多少錢,並決定如何處理剩下的財產.巴菲特以身作則,保證捐出99\%財產的承諾.他認為每個有錢人都會想到如何處理自己的財富,但要把握時機.「如果拖到90多歲才留遺囑,他們的腦力和意志力不可能比現在更強.......我的一切都來自美國人民,我沒有理由不把這一切交還給美國人民.」

來自中國大陸一位作者以此比較中美富豪,認為後者是真正的富足.他說:

``自從來到紐約之後,我一直目睹著美國歷代富豪慷慨捐贈的遺跡；聯合國總部的地皮來自石油巨頭小約翰-洛克菲勒(Rockefeller, John, 1839-1937) 的無償捐贈；現代藝術博物館和古根海姆博物館 (Guggenheim Museum)分別來自洛克菲勒家族和古根海姆家族的慷慨解囊；鋼鐵大王卡內基(Carnegie, Andrew, 1835-1919)為這裡的每一個公立學校捐獻了圖書館；比爾-蓋茨則給所有的公立學校圖書館免費安裝了網路連接.

這一切在中國大陸是很難想像的.我們的文化藝術機構長期缺乏運營資金, 我們的珍貴文物缺乏最起碼的保護,我們的公立圖書館系統基本不存在,更不要提網路連接的奢望.

與此同時,《華爾街日報》報導,中國是世界奢侈品消費的第三大市場,價值100萬美元的賓利(Benz)豪華轎車在中國長期供不應求.中國的富人以消費10萬美元一瓶的葡萄酒,5000萬美元一匹的賽馬或者1億美元一套的超豪華住宅"為榮."

他問:為什麼中美富豪對捐獻有這麼大的分別呢?為何中國根本沒有產生洛克菲勒,卡內基的土壤呢?

不少評論者都認為美國富豪的慈善捐獻,與猶太人和基督信仰的捐獻傳統有密切的關系:我們要行公義,好憐憫,存謙卑的心與上主同行.一般教會都教導要什一奉獻,有些教會並設立慈惠捐獻,特別幫助有需要的人.這種捐獻文化深入民心.慈善家卡內基就說過:"在巨富中死去是一種恥辱" ,與富人的榮譽不相乎.

耶穌在今天的福音經課,談論到一位財主,一生祇儲藏財富而不為他人利用財富,不單是恥辱,更是無知."無知的人哪,今晚我要你的靈魂,你所預備的要歸誰呢?(路12: 20) " 這人無知甚麼呢?為何耶穌說這個無知財主的比喻呢?我們看看這段經文的背景.

1.    Don"t Be a Fool and \#8212; Taking the Urgent as Significant 不要無知 and \#8212;把緊急的當為重要的 (12:13, 14)

``

眾人中有一個人對耶穌說:夫子,請你吩咐我的兄長和我分開家業.耶穌說:你這個人,誰立我作你們斷事的官呢?" ( 路12: 13, 14)

福音經課的開始,耶穌正對成千上萬人(新譯本by the thousands) 聽道,甚至彼此踐踏,耶穌提醒聽眾小心防備法利賽人,他們假冒為善,虛有宗教形式,模糊了信仰的焦點,以致使人遠離上主,他們的偽善將會被上帝審判.12: 2, 3節:掩蓋的事,沒有不露出來的；隱藏的事,沒有不被人知道的.因此,你們在暗中所說的,將要在明處被人聽見；在內室附耳所說的,將要在房上被人宣揚.

審判是頂嚴肅的信息,但聽到這資訊的人竟然靈機一觸,突然要求耶穌為他和兄長作家產的宙判官,而不反思在上帝的宙判面前,應如何自處.為甚麼?

我想分家產一事,一定終日煩擾著他,不竟家庭糾紛,不單是利益,還有潛藏的情緒牽引,與兄長家產之爭,肯定太太和嫂嫂也扯進去了.對他來說,家庭紛爭是一件緊急的事,耶穌作他家產的宙判官,可以解開他此刻的困擾,而這困擾使他看不清楚更重要的宙判--最後的宙判.

今天我們面對耶穌時,是否犯了同樣的錯誤,把緊急的關懷,取代了至為重要的關懷呢?你向耶穌禱告,為你解決難處,但有沒有向耶穌祈禱,赦免你的罪呢?你願意接受耶穌作你的救主,但有沒有接受他為你生命的主呢?

2.    Don"t Be a Fool and \#8212;The Soul Does Not Feed on Excessive Wealth 不要無知 and \#8212;靈魂的富足不在乎家道豐厚vv. 12: 15-20

求耶穌為他和兄長分家產的人,最關心是自己得到合理的財富,錢財使他有安全感.但在此之前(12: 4-12),耶穌豈不是說過上帝是看顧人的,不要懼怕迫害,那殺身體不能殺靈魂的不要怕,當怕那殺了身體以後又有權柄把我們丟在地獄的上帝,而這位上帝愛我們的；我們的頭髮都被數過,不要懼怕,你們比麻雀貴重得多.就算被押到官府被宙,不要思慮怎樣申辯,聖靈要指教我們該說些甚麼.(12:11, 12) 但他過還是很擔心,擔心沒有足夠的財富,怎能好好過活.其實耶穌知道我們的憂慮,在這比喻之後,耶穌繼續講論不要憂慮的資訊(12: 22--34) .

耶穌即時拒絕他的要求並對眾人說:你們要謹慎自守:免去一切的貪心,因為人的生命不在乎家道豐厚.就用比喻對他們說(v. 15).......故事中的財主雖然田產很豐富,原有的倉房己儲得滿滿,但還不滿足,拆掉倉房,蓋更大的,收藏更多的田產.他沒有以上主的應許(assurance) 作為保証,他要掌控自己的生命,蓋更大的倉房收藏更多的田產作保証(insurance) .他對靈魂說:靈魂哪,你有許多財物積存,可作多年的費用,祇管安安逸逸地吃喝快樂吧!

但人算不如天算,倉房還未蓋好,生命就在想不到的時候就結束了.他為儲蓄更多田產,蓋更大的倉房謀算預備,但卻沒有為著迎見上帝作任何準備.

但何以見得這財主是貪心呢?增加財富,發展業務是雄心是每個從牽的人的應盡之義,何貪有之?貪心和家道豐富有何相關呢?

人的生命不在乎家道豐富,豐富一字,其實有多出來的意涵.路加在9:17節:五餅二魚吃飽五千人,把剩下的零碎收拾起來,裝滿了十二藍子.15: 17:節浪子故事中,浪子打算回歸時說:我父親有多少的雇工,口糧有餘,我倒在這裡餓死嗎?21: 4節談到聖殿的奉獻.眾人奉獻都是自己有餘,唯有窮寡婦把養生的都奉上.有餘就是過份了(excessive) ,逾越界限,侵犯了其他人的神聖靈域.

我們採用八達通,付出五十元為押金,心甘情願,商戶接納八達通,付出傭金也是心甘情人願.但八達通為了每年多賺八千八百多萬,把八達通的持有人的私人資料當商品賣給客戶 and \#8212;應說與他們合作,賺取利潤,這是過份(excessive),是貪心.

結了婚的人,用情不專,到處留情,也是貪心,傷害了配偶.抄襲他人的文章,把人的學術成果據為已有,也是貪心.

十誡最後一誡提到:不可貪戀人的房屋,也不可貪戀人的妻子,僕婢,牛驢,並他一切所有的,似乎重複了以前的不可姦淫,不可偷盜,其實是總結禁令和誡命的核心 and \#8212;不可貪心.貪心是一種強烈的熾欲promiscuity.該隱殺了亞伯,上帝說:罪就伏在門前.它必戀慕你,你卻要制伏它.(創4: 7) 罪有牽引戀慕,這是罪的誘惑,但祇要我們停留在安全的境域 and \#8212;不出門口,我們可以制伏它.

無知的財主的問題不在發展其事業,增加其財富,乃在要掌控自己的生命.但他不是生命的主,怎能掌控自己的生命呢?他對靈魂說:靈魂哪,你有許多財物積存,可作多年的費用,祇管安安逸逸地吃喝快樂吧!

人怎可以說這些僭越的話呢?祇有上帝才知道人的靈魂將歸何處,將歸何時.說這些話的人,就是僭越上帝的權柄,是貪心.若財主蓋更大倉房的目的為要幫助有需要的人,不是為自己,就不是貪心了.所以問題不在財富多少,財富多可以感到富足,財富少的也可以感到富足.失樂園的作者約翰.米爾頓在失樂園中說得好:"There is nothing that makes men rich and strong but that which they carry inside of them. Wealth is of the heart, not of the hand.( 除了心之所繫,沒有一樣東西令人感到富足和壯大,富足是心靈的,不是握在手中的.)"

3.    Don"t Be a Fool and \#8212;True Wealth is Soul Wealthy 不要無知 and \#8212;真正的財富是心靈的滿足 vv 20, 21

神卻對他說,無知的人哪,今夜必要你的靈魂,你所預備的要歸誰呢?凡為自己積財,在神面前卻不富足的,也是這樣.(12:21, 22)

耶穌在這裡沒有稱呼無知的財主,稱他為無知的人哪.在一位國王的喪禮上,講道的牧師開始就說:Before death all bow(死亡面前,我們都要屈膝下跪.) 是的,無論是富有的,貧窮的,聰穎的,愚拙的,都要在死亡面前,俯首稱臣.

今天的詩篇經課四十二篇就提到:無論上流下流,；富足貧窮-都當留心聽.(49: 2) 留心聽甚麼呢?就是那些倚仗財貨自誇錢財的人,一個也無法贖自己的弟兄,也不能替他將贖價給神,叫他長遠活著,不見朽壞.(vv. 7, 8, 9) 死亡不單使人平等,也使人和畜生平等.他必見智慧人死,又見愚頑人和畜類一同死亡.(v. 10) ；人居尊貴中不能長久,如同死亡的畜類一樣(v. 12) ；人在尊貴中而不醒悟,就如死亡的畜類一樣.(v. 20) 人與其他生物不同,因為人有靈魂,沒有靈魂的人,祇是行屍走肉.

但怎樣才有心靈的滿足呢?心靈的滿足來自與上帝和好的關系.

耶穌對腰纏萬貫的財主說:

今夜必要你的靈魂,你所預備的要歸誰呢? (路 12: 20)

耶穌對身懸十架,身無分文的強盜說:

我實在告訴你,今日你要同我在樂園裡了.(路23: 43)

一句是嚴厲宙判的話；一句是盼望安慰的話.一個是與上帝隔絕,一個是與上帝結連.一個是殺了身體,又殺了靈魂；-個是殺了身體,保存了靈魂.

那一個人比較富有?

那一個人比較貧窮?

無知的人哪,今夜我要你的靈魂,你所預備的要歸誰呢?無知的富人哪,你其實窮得祇剩下錢.無知的富人哪,你其實窮得祇剩下錢.阿們.



\newpage

\section{看得見的與看不見的}
\label{sec:61}
\textbf{張美珍博士 - 看得見的與看不見的}
\newline
\newline
\hyperref[sec:60]{< < < PREV SERMON < < <}
~
\hyperref[sec:toc]{[返目錄]}
~
\hyperref[sec:62]{> > > NEXT SERMON > > >}


日期:2010年8月29日

主題:看得見的與看不見的

經文:路加福音12

証道:張美珍博士

我和外子於兩年前的暑假從加拿大回流香港,回流後十天,我便開始在中大上班.很多朋友知道我回流,過去兩年陸陸續續約我見面；但由於工作太忙,都被我一一推卻.兩個多月前,又收到一些老朋友的電話,相約見面；我說仍然很忙而再度推卻.老友劈頭一句:學生都放假了,不用教書還那麼忙?我回答說我在忙一些別人看不見的事──批改試卷及作業,準備往國內教學備課,為找工作的學生寫推薦信等.在中大教書,看不見的工作比看得見的工作更加重要,亦更令人疲累及忙碌.

我在十多年前曾在香港循道衛會牧會八年,記得有段時間,哥哥每逢周末將剛升讀中學的女兒交托予我,讓她參與我教會的團契及崇拜.那時我的教會在沙田一所中學內,主要的活動都在周末或主日舉行；我姪女第一次參加完崇拜後,聽完教會的活動報告,對我說:姑姐,我這份工作豈不是很清閒,只在星期六和星期日才需要上班?我那時回答她:牧師平日做很多事情,都並非人可以看得見的(例如對會友及社區的關懷,婚姻及個人輔導,準備會議及講章,計劃教會的發展等等).同樣,看不見的東西比看更見的東西更忙碌,亦更加重要.

很可惜,在這世上,大部份人都只著眼或重視看得見的東西,忽略了看不見的東西,而這些看不見的比看得見的更加重要.當我在溫哥華居住的時候,有次接待了卡加利一位朋友(是單親媽媽)的兒子(剛進大學不久),他來溫哥華看醫生,早機來,晚機走.當他看完醫生後,由於尚有數小時才上機,他提出想買件禮物送給母親,因為還有個多星期便是母親節.我記得陪著他在購物商場逛了兩個多小時,看了很多一般女性都喜愛的東西,例如香水,手袋,銀包,廚房用品,小擺設等等,但他說他母親這些統統都已經有了；結果他帶著很惆悵的心情上了飛機.

數月後,我有機會去探訪這位友人,在他們家裡吃飯；我想起這件小事,便問她兒子最後買了什麼送給母親,誰知他面有難色,不想再提.她母親插咀說,她兒子送給她的禮物,是理髮代用券一張,花了一百多元加幣,簡直是浪費金錢,因為普通的理髮,只需二十多元.她兒子告訴我,他返回卡加利後,適逢大學舉行義賣籌款,他既然為母親節禮物躊躇,不如拍賣一些特別的東西,送禮物之餘又可做善事,何樂而不為!結果他所拍賣得來的這張理髮券反令他被母親責怪了一頓.問題就是出於他的母親太著重那「看得見」的禮物,而忽略了那「看不見」的兒子的苦心.

我知在座中有很多是中大的學生,大家讀大學為了什麼?一紙文憑?入職資格?書本的知識及技能?是為了看得見的還是看不見的東西?我在中大教了兩年,前後教了十六科；印象中,有相當多學生很重視分數與成績,尤其是選修科及通識科.在開學時講解完課程大綱之後,學生發問的大部份問題,都與功課或考試有關.記得有一科有學期考試,他們問了很多關於考些什麼內容的問題後,有個學生衝口而出:「老師,聽你這樣說,這科很難tip題目je !」

我常對學生說,隨著互聯網面世後,人類對知識的倚賴將越來越低──只要能夠接上互聯網,世界各地,各學科的知識都唾手可得；唸大學,最重要的是學習掌控知識及技術的能力,亦即是訓練自己的思維方式及判斷能力,還有就是待人接物的處世態度,建立正確的人生觀和價值觀,這些「看不到」的東西,才能令學生一世受用.

我在美加居住了十五年,回來香港兩年,與在加拿大的生活相比,最強烈的感覺,就是香港的物質主義非常盛行,物質及金錢主宰著很多人的生命──彷彿擁有越多的物質,生活就越快樂.結果很多人就是為了追求這些看得見的物質而出賣時間與勞力(甚至人格),拼命賺錢.在加拿大的生活完全不一樣.大部份加拿大人過的是非常簡樸及平庸的生活,他們情願賺少點錢,多點時間來享受家庭生活.記得當我在加拿大生活時,有段時間幫家人打理一所西式快餐店.由於位處購物商場,除了聖誕及元旦外,所有公眾假期必須開門營業.記得在招聘員工時,有位白人女士來見工,她有很多年快餐店經驗,但卻明言在公眾假期不會開工；我正在猶豫是否應聘用她時,她叫我不要擔心,因為很多華人及亞裔人士會搶著在假期開工,原因是假期開工的工資是平時的一倍半.結果被她言中,我發現每逢假期,大部份在商場的快餐店上班的,都是華人.是我們華人特別勤力,捱得,還是我們特別願意做金錢的奴隸?

在香港,賺錢的法可謂層出不窮.為了錢,有大財團出賣你我私隱,發展商賣樓想盡辦法「呃秤」,地產中介可閃閃縮縮隱瞞資料；近日名噪一時的導遊阿珍,亦是為了要賺取更多的回佣而破壞了香港的聲譽及形象!其實也很難全怪阿珍,因為自從香港開放了自由行之後,很多內地同胞來到香港一擲千金,尖沙咀的國際名店經常見到咱們同胞排隊購物,近年聽說你與我都難以負擔的香港豪宅,都成為祖國同胞的獵物.可能這令阿珍以為同胞們來香港旅遊,都是為了購物,為了這些「看得見」的東西,所以沒有想過旅遊可以是為了「看不見」的東西,例如單純地休息,或者體驗當地的風土人情及文化.

今日讀經的福音書(路十二章),主耶穌告訴我們,上帝賜給我們這麼寶貴的生命,不是為了追求那看得見的金錢或物質,而是為了實現那看不見的上帝國.

路十二:29-32如此說:「你們不要求喫其麼,喝甚麼,也不要罣心.這都是外邦人所求的,你們必須用這些東西,你們的父是知道的.你們只要求他的國,這些東西就必加給你們了.你們這小群,不要懼怕,因為你們的父樂意把國賜給你們.」求上帝國是什麼意思?每星期我們唸主禱文時說,「願你的國降臨」,到底什麼是上帝國?

早幾天,如果大家有留意國際新聞的話,相信都會知道以色列與黎巴嫩再起衝突,雙方軍隊都有開火.新聞說是因為一棵樹──據說以色列軍想鋸走一棵樹,而黎巴嫩人則說那棵樹是屬於他們的,阻止以軍鋸樹,因而開火.以,黎兩國出現衝突,真是為了「看得見」的這棵樹嗎?其實在中東,以色列與巴勒斯坦經常都為樹而起衝突.紐約時報專欄作家Thomas Friedman,亦即是"The World is Flat"《地球是平的》一書作者,他寫了很多書批評全球一體化所帶來的破壞及影響；去年最新出版的書為"Hot, Flat, and Crowded" 《地球又熱,又平,又擠》,可能大家都曾聽聞.他在二千年曾出版一本書,名為"The Lexus and the Olive Tree"《凌志汽車與橄欖樹》,當中提到過去數十年,以,巴兩國人民都經常為樹而起爭執,搞到這麼大件事,當然不是真的為了這棵「看得見」的橄欖樹,而是為了「這橄欖樹屬誰」──即是看不見的這樹的所在領土主權而爭執.

當主耶穌教我們追求上帝國,求上帝的國降臨,意思即是實現上帝的主權.而上帝國的領士在那裡?就在我們每一個人的心內.早在二千年前,主耶穌已經警告過我們人類:

「一個人不能事奉兩個主,不是惡這個愛那個,就是重這個輕那個.你們不能事奉上帝,又事奉瑪門(瑪門是錢財的意思).」(太六:24)看得見的錢財,或許能夠解決我們生活上的需要,甚至帶給我們豐富的物質享受,但卻不值得奉它為主,出賣我們的生命,因為這些都是短暫的,沒有永恆價值的.上帝國雖然是肉眼看不見的,但當我們接受上帝的主權,讓上帝來管理我們的人生,在你身邊的人是可以感受得到的,因為上帝就是愛.一個充滿愛的人,不論何時何地都能夠令人感到溫暖；一個充滿愛的人,是善良的,是不會傷害人的,更不會做出損人利己的事.唯有愛,才有永恆的價值,亦即是主耶穌說積財寶在天的意思.

想知道什麼有永恆價值,大家如果有去過安息禮拜便會曉得.我覺得在安息禮拜中,是最容易體會什麼是永恆價值.每當有親友悼念死者或講述死者生平,所說的一定不會是死者擁有過多少「看得見」的物質──賺過多少錢,住過那些豪宅,開過什麼名車,嚐過多少美酒佳餚；反而會講述死者為家庭,別人或社會作出過什麼貢獻,付出過多少愛心.

主耶穌說,「你的財寶在那裡,心也在那裡」.今天,我們的心是放在那些看得見的,會朽壞的物質上,還是那看不見的但可感受到的上帝國那裡?誰在主宰我們的人生?是瑪門,還是上帝?願望我們時刻儆醒,不要在不知不覺間出賣了自己的靈魂予那些看得見,沒有永恆價值的東西.



\newpage

\section{戰爭與和平}
\label{sec:62}
\textbf{賀志勇博士 - 戰爭與和平}
\newline
\newline
\hyperref[sec:61]{< < < PREV SERMON < < <}
~
\hyperref[sec:toc]{[返目錄]}
~
\hyperref[sec:63]{> > > NEXT SERMON > > >}


路加福音 12:49-53
\newline
\begin{longtable}{cl}
\hline
\hline
章節 & 經文 (和合本修訂版)\\
\hline
12:49 & \begin{tabularx}{0.7\textwidth}{X} 「我來是要把火丟在地上，假如已經燒起來，不也是我所希望的嗎？ \end{tabularx} \\ \\ \relax
12:50 & \begin{tabularx}{0.7\textwidth}{X} 我有當受的洗還沒有受，在這事完成之前，我是多麼地焦急！ \end{tabularx} \\ \\ \relax
12:51 & \begin{tabularx}{0.7\textwidth}{X} 你們以為我來是要使地上太平嗎？不！我告訴你們，是使人紛爭。 \end{tabularx} \\ \\ \relax
12:52 & \begin{tabularx}{0.7\textwidth}{X} 從今以後，一家五個人將要紛爭，三個和兩個相爭，兩個和三個相爭： \end{tabularx} \\ \\ \relax
12:53 & \begin{tabularx}{0.7\textwidth}{X} 父親和兒子相爭，兒子和父親相爭；母親和女兒相爭，女兒和母親相爭；婆婆和媳婦相爭，媳婦和婆婆相爭。」 \end{tabularx} \\ \\
[1ex]
\hline
\hline
\end{longtable}
日期:2010年9月5日

主題:戰爭與和平

証道:賀志勇博士

經文:路加福音12:49-53

我來要把火丟在地上．倘若已經著起來,不也是我所願意的麼.我有當受的洗．還沒有成就,我是何等的迫切呢.你們以為我來,是叫地上太平麼．我告訴你們,不是,乃是叫人紛爭． 從今以後,一家五個人將要分爭,三個人和兩個人相爭,兩個人和三個人相爭．父親和兒子相爭,兒子和父親相爭．母親和女兒相爭,女兒和母親相爭．婆婆和媳婦相爭,媳婦和婆婆相爭.

一,舊耶路撒冷:戰爭的無處不在耶路撒冷,希伯來語的意思就是平安之城.公元70年,羅馬提多將軍(Titus)圍困耶路撒冷長達幾個月.城裡斷糧斷水,饑餓、疾病蔓延,有婦人餓得把自己的孩子都宰來吃.當羅馬軍隊攻陷城牆時,發現全城幾乎沒有活人,遍佈屍體,連悲痛的哀號都聽不見.這是一場慘絕人寰的戰爭,共死110萬人,九萬多人被捉拿.提多將軍幾乎將全城鏟平,整個聖殿變成廢墟,"沒有一塊石頭留在另一塊石頭之上."--在人類歷史中,哪一次戰爭不慘烈呢?!

這世界有兩樣東西無處不在,一個是交稅,另一個是戰爭.只要有政府,就要交稅；只要有人,就會有戰爭.二戰離香港不到70年.而中國大陸,建國60年以來,光大的戰役就有九次之多--抗美援朝,抗美援越、中越邊境自衛反擊戰,等等等等.小的衝突就更不用說,只是報紙一般不賣罷了.這個世界並不太平.我們和平的、歌舞昇平的歷史表面底下,其實是暗流洶湧.

今天,各種形式的戰爭無處不在.家人、同事、朋友間的戰爭,政治、經濟、文化的戰爭,甚至自己與自己的戰爭,只要你留心,通街都是,你到處都看得到.《創世記》裡面朋友間的戰爭,《宮心計》裡面女人之間的戰爭,巴士阿叔跟小青年的戰爭,女導遊跟遊客的戰爭,龔如心弟妹跟風水佬陳振聰的戰爭,韓國跟中國爭孔子是不是韓國人的戰爭.我不用一一舉例,太多了!現實生活永遠比蘋果日報、明報的新聞要複雜一千倍一萬倍.可是戰爭的結果是什麼?不要說二戰死傷將近1點9億人.哪怕是小戰爭,你看看《創世記》裡頭的葉榮添與許文彪,他們的戰爭讓他們失去了友情、生命.再看龔如心,雖然贏了官司抓到了錢,最後還是因為那些錢,死後還被人公開爆料,真真是入土難安啊.所以古人講:甯為太平犬,不做亂世人.人人都希望這個社會是個太平盛世,家是個太平之家,公司是個太平公司,當然,醫院太平間要除外.

今天,我們就要看一段有關耶穌到世間引起戰爭的經文,一段很特別很特別的經文.請大家跟我一起來看路加福音12章49-53,57-59節.

二,耶穌是什麼樣的?--一個點火受洗的人.(12:49-50)我想提一個問題,在你心目中,你是怎麼看待耶穌的?我想,你可能很少會去想這個問題.但如果真要你去想,你可能會搬出小時候主日學學過的答案:耶穌是神的兒子；耶穌是一個慈愛的牧羊人,溫柔地抱著小羊.有的人或者會說:"我覺得他是一個好人,說了許多對人有益的話,做了很多好事."我以前問一個小朋友,"耶穌是什麼樣子啊?"他告訴我耶穌有白鬍子,喜歡帶著紅帽子,背著個紅袋子.我猜想他可能把耶穌跟聖誕老人搞混了.我的同事很特別,她說,她覺得耶穌是個很瘦、很慘的男人,很慈祥,但是釘在十字架上受苦.這些回答概括起來,耶穌基本上是一個老好人,即使不是笑咪咪的,但肯定是很慈祥、很有愛心的樣子.

很少有人記得耶穌的另一種形象,在啟示錄中,耶穌眼目如同火焰,面貌如同烈日放光.他除了帶給人慈愛之外,也帶給人末日審判的宣告.馬太福音7章21節就說:

``凡稱呼我"主啊!主啊!"的人,不能都進天國；惟獨遵行我天父旨意的人,才能進去."路加福音13章5節,耶穌斬釘截鐵地說:

``我告訴你們,不是的.你們若不悔改,都要如此滅亡."非常非常嚴厲!

我們的世界已經習慣了慈愛的耶穌、恩典的耶穌,有許許多多神跡加給我們、讓我們覺得靈驗的耶穌,但還沒有習慣那位審判的耶穌、收取的耶穌、對我們掩面不顧的耶穌.約伯說:"賞賜的是耶和華,收取的也是耶和華."可是又有幾個人願意面對收取呢?所以當耶穌說"我來是要把火丟在地上,...我有當受的洗還沒有成就..."時,這種語氣讓人何等的懼怕.耶穌所說的火,就是上帝審判的火；當受的洗,就是耶穌被釘死在十字架上的苦難以及救恩的完成.當耶穌說這話時,心情是何等複雜!他一面嚴厲地提醒世人,末世已經來到,不悔改,審判必要降臨；另一面,他知道他即將釘死在十字架上,但為了完成他的救恩,他又迫切希望那個時刻快點到來.他的憐憫又是何等的深!我想做父母的也是這樣.一方面惱孩子不爭氣、不上進,要嚴厲管教；一方面心裡又愛又痛得不得了.我記得我小時候調皮搗蛋,我爸爸狠狠打了我一頓,過後他自己都流眼淚,又給我西瓜吃,又給我說他是為了我好才打我的.天父的心也是這樣啊.

三,真理要我們改變--主耶穌所帶來的紛爭(12:51-53)如果這段經文前兩句話已經讓人不舒服,那麼後面幾句就讓人更難懂了.耶穌居然說:

``12:51, 你們以為我來,是叫地上太平麼．我告訴你們,不是,乃是叫人紛爭."這好像不是耶穌的行事為人吧,他不是說"使人和睦的人有福了"麼?現在他又說他是來製造麻煩的,搞搞震.這些話讓我們聽不懂.但真正的耶穌,帶著權柄又充滿憐憫的耶穌,卻真真實實會帶來紛爭.因為他所宣告的真理是絕不讓步、是不打折扣的.他之所以會帶來紛爭,是因為他的真理不會改變,而我們卻需要改變.

改變容易不容易?不容易!你看我家裡,收埋好多舊東西.搬家時,小弟跟我說,"你好多東西要扔掉了",我總捨不得.每次要扔的時候,總有點不安全感,怕什麼時候還會用到吧,又收回來.我是個怕變的人.到現在還是無辣不歡,每次出去,都要找找哪里有辣椒菜.我發現人都是這個樣子啊!我跟香港的朋友出去旅行,才過兩三天,他們就到處找哪里有茶餐廳,哪里可以飲茶了.

很多人都怕新,怕變,因為新跟變都是一種危機.一切都照舊,當然給人安全感.人家說要改變一個人,先要改變他屁股底下坐的椅子--他坐得太久太舒服了,換個新的不容易啊.椅子是什麼?椅子就是你的習慣.所以你要一個人在跟人說話時,忘掉他自己是博士、碩士,是教授、董事,是名校畢業,很不容易,因為那是他的驕傲；聖經上說"快快地聽,慢慢的說,慢慢地動怒",但我們往往是"慢慢地聽,快快地說,快快地動怒",因為那是我們的脾氣.要改變習慣是很不容易的.

但耶穌來了,就完全不同了!我們不是改變一點小東西就可以過關了,我們要完全改變,改變我們的習慣、改變我們的生活去跟從他,這樣肯定會帶來紛爭.正如保羅所說"...因 你 們 已 經 脫 去 舊 人 , 和 舊 人 的 行 為 , 穿 上 了 新 人 ,... " (西3:9-10)

我們要脫去舊人,所以我們會與過去的我爭戰；我們穿上新人,我們又會與這個屬過去的世界爭戰.在這個末世中,政治、經濟、文化的戰爭人人看得到,但有的戰爭並不是人人都看得到.比方說,家庭的戰爭.中國現在最大的問題是家庭問題.在中國大城市,25-40歲之間的夫妻,離婚率高達百分之四十.香港的離婚率也高達百分之37.這還不包括沒離婚卻家庭不和的夫妻.你想想,現在家庭問題有多嚴重!我聽內地朋友說,他們同學畢業十周年聚會,大家的問候都是,"你離婚了嗎?"我還聽說某個香港富豪的別墅裡,每個兒子都有自己的飯廳,一家人從不一起吃飯--``從今以後,一家五個人將要紛爭,三個人和兩個人相爭,兩個人和三個人相爭．"是耶穌帶來家庭紛爭麼?不是!是我們的罪,是我們的驕傲、貪心讓我們與我們最親愛的人,最親密的關係紛爭.所以,不是基督把紛爭帶入世界,而是基督清楚知道:當他來到世界,這世界的罪惡就必然會反對他,逼害他.真理與邪惡的鬥爭不可避免.主對蒙召的門徒這樣說:

``我差你們去,如同羊進入狼群."(太十 16)真正使家庭紛爭、親人不和的原因乃是這世上的罪,乃是如狼群一樣的罪人.

這種紛爭在基督徒與世界之間更明顯.聽我的同事阿芝說,她的牧師講道時候跟會眾說:"你們老說自己窮,可是你們個個都有幾部手機.你看看你衣櫃裡有多少件衣服,你看看你鞋櫃裡有多少對鞋?如果你們能改變生活方式,就不會哭窮了."但我想,改變生活方式不容易.你如果很節省,可能你過去的朋友就不帶你玩了,就排斥你了；在大陸做生意,你如果不跟客戶喝酒,就接不到生意了.世界有世界的遊戲規則,我們一旦改變,紛爭就不可避免.以前我在浸會大學,有個重慶的姐姐老是諷刺我,後來很偶然的機會她來了查經班,很奇妙地信了主.她跟我說:"小勇,以前沒信主,老覺得你怪怪的,覺得基督徒怪怪的；現在信了主,才發現不是基督徒怪怪的,是基督徒太正常了."是啊,我們太正常了,我們與這個瘋狂的、賣毒奶粉的世界差別太大了.世界之所以與我們紛爭,是因為我們跟這個世界不同.我聽說在國內,如果你不是貪官,你就沒法做官；因為別的貪官都懼你、怕你、恨你,因為你與他們不同.我們基督徒也是這樣啊,世界懼我們、怕我們、恨我們,是因為我們不同,是因為我們不肯同流合污.我們是做光做鹽的,但正因如此,從基督降臨的那個晚上,世界與我們的戰爭就從來沒有停過.耶路撒冷就是一個縮影,兩千年來,這個平安之城從來就沒有平安過.

四,與對頭和好--主耶穌將帶來和平(12:57-59)平安--Peace--這個詞,,在聖經裡頭還可以翻譯成和平,或者和睦,是同一個詞來的.可是這樣一種和平或者和睦,絕不是不講原則、縱容作惡,或者忍氣吞聲的和平.比方說在家庭當中,如果沒有了原則,就容易縱容子女,表面的和睦終有一天會害了孩子.神所帶來的和平絕不是表面的一團和氣,他的和平是有原則的,他不僅是慈愛的神,也是公義的神、審判的神.他所帶來的和平,一定是與罪惡經過鬥爭的和平.每一次真正長期的和平,都是經歷過鬥爭獲得的.就好像中國古代的盛世,漢朝的"文景之治",唐朝的"貞觀開元盛世",明朝的"永宣之治",都是中國人經歷過長期鬥爭才獲得的結果.

但有時候,我們的敵人不是在外部,很多時候,我們的敵人往往是我們自己.所以在57-59節中,耶穌教導我們要與我們的敵人和好,表面上看起來,耶穌是要我們對別人好,但奇怪的是,最終受益者卻往往是我們自己.簡單來講,和敵人和好能夠減少我們的敵人；但更深一層講,如果我們內心藏著怨恨,用抗爭和仇視來對待與我們有矛盾的人,付出巨大代價的則是我們的健康、我們的安寧和我們的生活質量.如果我們對別人做過的事始終耿耿於懷,怨恨帶來的是越積越深的沮喪.我們將不得安寧!當我們和別人敵對的時候,我們往往不是把自己,而是把別人當成了我們的生活中心.我們這樣做是和對手一起來殘害自己,不同的是,對手的攻擊來自外部,而自身的殘害卻來自內部.

Eric是華爾街很出名的經紀人,賺了很多錢.有一天他疲倦了,他厭倦了華爾街中間人與人的欺騙,人與人的爭鬥.他放棄了工作,去了山區的一家修道院靜修.每天早上他5點半起床去靈修,教堂裡只有他一個人.哈!多清淨,多安寧,他這樣開心地過了三天.到第四天,他5點半去的時候,發現已經有個人比他還早就跪在那裡靈修了.他心想這樣不行；第五天,他提早15分鐘,5點15分就到了教堂.不錯啊,教堂就他一個人,他又覺得安寧滿足了；可是第六天,他發現那個人也提早了,比他早到.他心裡很不服氣,第七天,他又提早15分鐘,5點鐘就到了教堂.終於到了第九天,他為了贏得這場鬥爭,不得不4點半起床.那天天還沒亮,黎明一點點浮現在天空中,他跪在黑濛濛的教堂裡,他突然發現,由於他的驕傲,這個修道院變得跟華爾街已經沒有什麼分別了.經過這番戰爭,他終於看清了他自己,也終於改變了他自己.--有時候,不是外部環境的問題,是心的問題,是我們內心的問題.

所以,當我們與敵人和好的時候,實際上也是與自己和好.當我們放下了對敵人的怨恨的時候,我們也放下了自己心中的苦毒.怨恨,乃是借人家的錯誤來懲罰我們自己.主耶穌要我們放下怨恨、放下驕傲,與敵人和好,其實,也就是要我們放下自己內心的問題,與自己和好,與主耶穌和好,與他所創造的這個世界和好.但這一切容易不容易?罪總是引誘我們去戰爭,只有靠著神的話語、神的恩典與公義,才能讓我們認清自己,重獲和平.

五,結論--新耶路撒冷:末世的戰爭與和平啟示錄說,將來,舊耶路撒冷將成為過去,新的耶路撒冷將從天而降.在新耶路撒冷,神將擦去我們一切的眼淚,我們將獲得和平.但在這和平之前我們會先經歷戰爭,末世的戰爭.不要以為末世還沒有到,大家吃好喝好玩好睡好.其實當主耶穌降生的那個夜晚,末世就開啟了.末世戰爭每天都在進行,每天我們都在跟自己作戰,跟怨恨、驕傲、苦毒、不公義作戰,但我們基督徒滿有希望,神必定會掌權.靠神改變之後,我們失去的是戰爭,獲得的將是真正的和平.



\newpage

\section{你們給他們吃吧}
\label{sec:63}
\textbf{王群女士 - 你們給他們吃吧}
\newline
\newline
\hyperref[sec:62]{< < < PREV SERMON < < <}
~
\hyperref[sec:toc]{[返目錄]}
~
\hyperref[sec:64]{> > > NEXT SERMON > > >}


馬可福音 6:30-44
\newline
\begin{longtable}{cl}
\hline
\hline
章節 & 經文 (和合本修訂版)\\
\hline
6:30 & \begin{tabularx}{0.7\textwidth}{X} 使徒們聚集到耶穌那裡，把一切所做的事、所傳的道全告訴他。 \end{tabularx} \\ \\ \relax
6:31 & \begin{tabularx}{0.7\textwidth}{X} 他就說：「你們來，同我私下到荒野的地方去歇一歇。」這是因為來往的人多，他們連吃飯的時間也沒有。 \end{tabularx} \\ \\ \relax
6:32 & \begin{tabularx}{0.7\textwidth}{X} 他們就坐船，私下往荒野的地方去。 \end{tabularx} \\ \\ \relax
6:33 & \begin{tabularx}{0.7\textwidth}{X} 眾人看見他們走了，有許多認識他們的，就從各城步行，一同跑到那裡，比他們先趕到了。 \end{tabularx} \\ \\ \relax
6:34 & \begin{tabularx}{0.7\textwidth}{X} 耶穌出來，見有一大群的人，就憐憫他們，因為他們如同羊沒有牧人一般，於是開始教導他們許多事。 \end{tabularx} \\ \\ \relax
6:35 & \begin{tabularx}{0.7\textwidth}{X} 天已經很晚，門徒進前來，說：「這地方偏僻，而且天已經很晚了， \end{tabularx} \\ \\ \relax
6:36 & \begin{tabularx}{0.7\textwidth}{X} 請叫眾人散去，他們好往四面的鄉鎮村莊去，自己買些東西吃。」 \end{tabularx} \\ \\ \relax
6:37 & \begin{tabularx}{0.7\textwidth}{X} 耶穌回答他們說：「你們給他們吃吧！」門徒對他說：「我們要拿兩百個銀幣去買餅給他們吃嗎？」 \end{tabularx} \\ \\ \relax
6:38 & \begin{tabularx}{0.7\textwidth}{X} 耶穌說：「你們有多少餅？去看看。」他們知道後就說：「有五個，還有兩條魚。」 \end{tabularx} \\ \\ \relax
6:39 & \begin{tabularx}{0.7\textwidth}{X} 耶穌吩咐他們，叫眾人一組一組地坐在青草地上。 \end{tabularx} \\ \\ \relax
6:40 & \begin{tabularx}{0.7\textwidth}{X} 眾人就一群一群地坐下，有一百的，有五十的。 \end{tabularx} \\ \\ \relax
6:41 & \begin{tabularx}{0.7\textwidth}{X} 耶穌拿著這五個餅和兩條魚，望著天祝福，擘開餅，遞給門徒，擺在眾人面前，也把那兩條魚分給眾人。 \end{tabularx} \\ \\ \relax
6:42 & \begin{tabularx}{0.7\textwidth}{X} 他們都吃，並且吃飽了。 \end{tabularx} \\ \\ \relax
6:43 & \begin{tabularx}{0.7\textwidth}{X} 門徒把餅和魚的碎屑收拾起來，裝滿了十二個籃子。 \end{tabularx} \\ \\ \relax
6:44 & \begin{tabularx}{0.7\textwidth}{X} 吃餅的男人共有五千。 \end{tabularx} \\ \\
[1ex]
\hline
\hline
\end{longtable}
日期:2010年9月12日

主題:你們給他們吃吧

証道:王群女士

經文:馬可福音6:30-44

引言:剛才誦讀的福音書經課是大家都耳熟能詳的故事,我們一般籠統地稱它為「五餅二魚」.我想如果要舉辦一個聖經故事選舉的話,論知名度,論代表性,論吸引性五餅二魚一定榜上有名,甚至名列前茅.因為無論你年紀多大,信主年日有多長,我相信這個神蹟故事一定是你最早接觸的聖經故事之一,甚之連未信主的朋友都聽過這個故事.透過它你我認識到耶穌基督是大有能力的,祂可以把原本僅有的變成足夠的,變成有餘的,在一首現代詩歌裡其中幾句歌詞就是這樣說的:「無論是多少,全是不緊要.在我主手中,一切也奇妙.」你又會發現不少有基督教背景的機構,活動,網站都愛用「五餅二魚」作為名稱,有跟餐飲業有關的,有跟慈善扶貧有關的,亦有跟搵工招聘,資源共享等等有關的.「五餅二魚」就如我們從小就聽到的寓言故事一樣,一提起它,我們自然意會到特定的內容,同時亦會聯想到特定的意義,例如「狼來了」是教訓我們不要說謊,又例如「龜兔賽跑」是教訓我們驕兵必敗的道理.這些寓言故事的其中一個特點是,我們在口耳相傳的過程中,已經不太需要知道這個故事原本的情節是怎樣發展的,我們就能一跳跳到結論教訓的部份.如果你有機會認真聆聽別人講述這兩個故事的話,你會發覺一個有趣的情況,就是不同人所講的內容都有所不同,每個人都在細節上有自己編輯,改動過的痕跡,彷彿故事的細節變得不重要,重要的只是教訓的道理是什麼,教訓必須一致,否則就不是那個故事了.因此一旦我們掌握了那個特定的教訓,這些故事就沒有重看的必要.「五餅二魚」是否已經成為其中一段我們不需要經過細心閱讀,就能直接跳到它要帶出意義的經文呢?

最近有一部荷里活電影《潛行凶間》(Inception)在香港上映,單憑它的中文戲名並不能吸引我走入電影院觀看,但後來實在聽到太多讚譽,很多朋友都表示這是一套值得一睇再睇的電影,不少專欄,影評更談論到這套電影掀起了本港近年少見的「翻睇」熱潮,就是觀眾願意花真金白銀,一而再,再而三地走入電影院.而我之所以會進入電影院看第一次亦是因要陪我妹妹看第二次.電影看完之後,我發現故事並不如想像中的艱深難明,心裡正暗喜自己頗有慧根,不用看第二次,第三次,就能掌握到故事的關鍵,講的不就是現實與夢境的關係,這不是二千多年前莊周夢蝶故事要說的嗎?然而後來,跟妹妹談論到劇情,又閱讀到一些評論的時候,我才發現原來電影中有一些重要的細節是我沒有留意到的,我亦發現原來吸引大部份人重看這套電影的並不是故事的橋段,而是故事的細節,而這些細節卻往往是提示出導演想說的根本訊息.我希望今天大家都乘著這股「翻睇」熱潮,懷著可以睇真D ,睇多D 的期盼一起閱讀這段熟悉的經文,拿走「這個故事係要教訓我地......」的框框,讓上主親自跟我們重述這個故事,帶領我們看今天早上祂想我們看到的重點.

經文背景:讓我們一齊來看看這件事蹟是在怎樣的情況下發生的?第30節是緊接著前文第7-13節而發展的,門徒奉耶穌的差遣兩個兩個出去,他們靠著主賜給他們的權柄去傳道,趕鬼,醫病,從第12-13節的記載,我們知道他們的工作成果很理想,他們趕出了許多鬼,又治好了許多病人,在座每一位應該都可以理解到這些繁重工作為他們帶來了身心疲累.因此當他們向主耶穌匯報完一切後,耶穌就建議他們要離開人群,退到曠野的地方歇一歇,在第31節的下半節,福音書的作者補充了一句:「這是因為來往的人多,他們連吃飯也沒有工夫」,由此顯示出這個退下來的計劃對他們來說是何其必須及重要的.然而他們的「如意算盤打不響」,這個退修休息的計劃很快就被從各城湧至的人群打亂了.這大批人群比耶穌及門徒更早到達他們要去的地方,34節記載耶穌上了岸,「見有許多的人」,就「憐憫」他們,然後就「開口教訓他們許多道理」.耶穌因為看見群眾靈裡的需要,他們如同沒有牧人的羊,就產生了憐憫的情感,繼而帶來具體的行動,就是甘願犧牲祂自己都好需要的休息時間,改變原本計劃好的日程,把道理教訓給他們.接下來我們發現門徒同樣有因住看見,而產生情緒,接著帶出實際行動,不過他們的反應卻與耶穌的反應是南轅北轍的.那門徒看到什麼?我們從35-36節門徒的說話中可知,他們見到天已晚了,他們又身在野地,更嚴重的是面前結集了大批的民眾,門徒察覺他們正面對一大能題,就是要怎樣解決這班群眾的膳食問題,又或者準確些說是如何不讓這班群眾的膳食問題發展至成為他們的問題.雖然經文沒有直接交代,但我們可以想像門徒面對著既失天時,又無地利,還欠人和的局面,會產生怎樣的憂慮不安的情緒.這種情緒驅使門徒要求耶穌叫眾人散開,因為他們希望群眾可以自己解決食物的問題.如果我們從門徒的角度出發設身處地想一想的話,我們應該也不會認為他們的憂慮是多餘的,從客觀條件來看,天時,地利,人和三方面都完全不配合,無辦法中的辦法就是化整為零,各自為政,則所謂的「自己顧自己」.他們這種思維模式跟我們現代人頗為接近,所以在一定程度上我們對他們的想法有不少的共鳴.大家習慣的做法是:看環境,個人的條件,面前的數字及背後牽涉到的成本等等,經過一輪計算,門徒清楚看見自己的限制,他們對於解決群眾膳食的問題上是無能為力的,所以最好的做法就是及早驅散人群.我們在生命中是否都看見甚至經驗到很多個人的限制呢?當我們的年紀漸長,身體機能卻漸退；當我們累積到一定人生經驗後,肩頭上的責任亦越來越重；加上我們在座的都有一定的宗教生活,我們部份人更好像門徒一樣參與不少的事奉.家庭,學業,事業,信仰佔據了我時間表上的絕大部份,我已經疲於奔命,甚至連個人休息時間也貢獻出來了,主啊!你還想我怎樣呢?

概念:儘管我們是軟弱的,不足的,有限制的,但只要我們願意順服主的吩咐,我們可以成為主向人施恩的器皿.

基於前面提及過的種種限制,第36節門徒向耶穌提出一個合理的要求:「請叫眾人散開,他們好往四面鄉村裡去,自己買什麼吃.」但他們得到的卻是一個看似不合理的回覆,耶穌回答說:「你們給他們吃吧!」耶穌在此不單沒有順應門徒的要求,反過來用了一個命令的語氣,直接要求門徒給他們吃,耶穌把原來由門徒拋過來的「波」交回給門徒.耶穌不把讓群眾吃飽的責任加諸自己身上,也不拋回給群眾,視作為他們個人的問題,不干我事,而是將這個餵飽眾人的重大責任轉交給門徒,要求門徒給他們吃.我們細心想的話,就會知道耶穌這個要求有多奇怪的,在六章30-31節已經交代了耶穌清楚知道門徒忙得「連吃飯也沒有工夫」,他們原本的計劃是暗暗地退到曠野去歇一歇,他們又怎會隨身帶備足夠大批群眾的食物呢?難怪門徒會有如此的反映:「我們可以去買二十兩銀子的餅給他們吃嗎?」(v.37下)如果我們明白「二十兩銀子」到底是什麼意思的話,我們便能更準確掌握門徒這句反問所附帶的情緒.新漢語譯本將「二十兩銀子」翻譯成「二百個得拿利」(denarii),「得拿利」是當時的貨幣單位,一個得拿利相當於一個普通工人一日的工資,二百個得拿利就等於一個工人二百日的工資.如果按照我們今天的生活水平,200日的工資又等如多少錢呢?雖然社會及政府對最低工資的定額尚未達成共識,商界希望將上限設定在24-26元,而勞工界則進取達到時薪33元,我們即管取其中間位,以30元1小時作計算,每日工作8小時,每日得240元,而200日的工資就是四萬八千元,我不知道這個數額對你的意義有太多,但對於草根階層來說,特別是身上只剩下五個餅二條魚的門徒來說,它絕對是一個巨額數字,那是超過普通人半年的人工總額,他們怎有可能隨時拿出來呢?難怪他們的語氣有點「倔」(rude).另外,我們從門徒的回應可以知道,他們的焦點仍然沒有離開眼前的人數,及由人數所兌換成的金錢數額,這些都是他們所缺乏的,是他們真實地面對著的限制.然而難道耶穌不明白這一切限制嗎?

我們的主當然明白這一切,事實是祂比任何人都更了解我們真正的限制呢!門徒真正的限制是他們對主耶穌心意的認知不足,他們對所跟隨的這位主的信心不足,他們亦對群眾真正的關愛不足.今天讓我們的生命裹足不前的真正限制又是什麼呢?我們的視野是否被局限在自己和自己所能看到的事情上面呢?主耶穌不單沒有否定我們的限制,反而繼續邀請我們一同參與祂的事工.經文的第38節,耶穌問門徒有多少個餅；之後39節祂吩咐門徒叫眾人有秩序地坐好；41節記載:「耶穌拿著這五個餅,兩條魚,望著天祝福,擘開餅,遞給門徒擺在眾人面前,也把那兩條魚分給眾人.」耶穌這一連串的動作是一般猶太家庭的一家之主會在餐前做的舉動,他們向上帝獻上頌讚,表明萬物都是從主而來的,為到上帝賜下人們生活所需的而獻上由衷的感恩；之後,擘開分給其他人.但在這裡經文清楚交代耶穌是「遞給門徒」,然後由門徒分發給群眾,門徒在此分享了耶穌在宴席中主人的身份,對待群眾如客人一樣,要盡力地款待他們.而第42-43節就交代眾人都吃,並且吃飽了,最後門徒盡義務把碎餅碎魚收拾起來,收拾的足有十二個籃子之多.我們在這個過程中發現一個頗有趣的地方,同時亦是關鍵所在,就是雖然門徒在37節對耶穌的回應有些撞衝的意味,但是對於耶穌的吩咐,他們卻仍然選擇聽從及順服,他們按照耶穌的指示,一步一步地完成,他們在做的過程中亦不見得對整件事有多深切的理解,只是他們依然聽從耶穌的吩咐去做.最後,難題被解決了,福音書作者在第44節才點出「吃餅的男人,共有五千.」五千這個確實的數目在這裡才首次被提及,而放在這個位置確實有種畫龍點睛的效果,原來門徒眼前的實際困難是比他們想像中的更加困難,但是因為他們聽話,願意聽從主的吩咐,參與主的事工,主便使用他們成為施恩的管道,亦讓因他們而得福的人數超出他們能想像的.

這件事如果能在這裡劃上句號,似乎會更加完滿,門徒的形象亦會更加討好.但是事實卻是我們意料之外的,在接著的記載v.45-52,耶穌與門徒暫時分開,祂獨自上山禱告,之後晚上海面起了波浪,門徒整夜辛苦地撐船,在天將快破曉,開始有些微光透露之際,耶穌從海面上走向他們的位置,門徒看不清楚,第一個反應就是以為是鬼怪,大呼小叫起來.這班表現得軟弱無能的不是別人,乃是曾經奉耶穌的差遣去傳道,趕鬼,醫治,有過所謂「成功」的事奉經驗,又親身見證著主耶穌餵飽五千人的神蹟,而且參與其中的門徒,我們以為他們必會因此上述種種的經歷就對主耶穌有進一步的認識和理解.這時候,在52節作者忍不住要從幕後走到台前,向我們說明門徒這些不合理的表現是「因為他們不明白那分餅的事,心裡還是愚頑.」在此我們進一步確認門徒的真正最大的限制是他們對主的認識不足,這個又何嘗不是你我的限制呢?而這個限制更是我們一日在世,一日都要面對的,保羅很透徹地明白我們當前的局限,同時亦為我們預告了出路,他說:「我們現在所知道的有限,先知所講的也有限；等那完全的來到,這有限的必歸於無有了.」(林前十三9-10)接著他繼續說:「到那時,就要面對面；我如今所知道的有限；到那時就全知道,如同主知道我一樣.」(林前十三12下)雖然我們的限制是何其的巨大,何其的實在,現時我們尚未能完全明白主的奧秘,但從門徒的故事看到,這些限制卻不會局限了主的工作,亦不會局限了主使用我們來完成祂的工作.門徒在還未完全弄清楚到底是一件什麼事情前,卻因著他們願意聽從主耶穌的話,主便使用了他們,叫很多人得到祝福.今天我們又是否願意謙卑地回到聖經裡,聆聽主對我們的說話,又按著祂親自所啟示的,實踐我們的信仰,叫更多的人得著從主而來的祝福呢?

結語:最後,我想跟大家分享一個見證,那是我在七月上大陸短宣時聽到的,這個故事的特別之處是它既是一個真人真事,卻又有很留白的地方.在這次短宣中,我們去到廣東省英德市其中一個鎮──含洸鎮探訪當地的教會,教會領袖跟我們分享他們教會的歷史,他特別提到一名叫祈姑娘(Miss Qi) 的西方宣教士,這位祈姑娘就是在上世紀20-30年代,一手建立含洸教會的人,但是由於教會經歷了過去的政治運動,有關這段歷史的文字紀錄已經散佚,知情者亦不在人世,他們現在所掌握的資料又零碎,又少,簡單如這位祈姑娘真正的姓名都不能肯定,亦不知道她到底是來自歐洲大陸,還是來自美國,更不知道她是屬於哪個差會,只知道這位無名的宣教士單身一人,坐船而來,上岸後就在岸邊買了一個細小的地方,開始聚會,後來她傳福音給一位姓林的先生,他之後成為了當地的傳道人,福音亦逐漸在含洸這片不出名的土地上zhak 根.雖然我們再沒有其他有關祈姑娘的資料,但我們可以肯定的是祈姑娘與同時期來華的西方宣教士一樣,所面對的困難和限制都是超出我們能想像到的,一名女子,飄洋過海,孤身一人,來到一個言語文化完全不通的落後地方,一切從頭學起,所有東西都超出她的掌握以中,但她仍然作出這樣選擇,唯一的原因就是她體會主耶穌的心腸,她又願意聽從主的吩咐,來尋找及餵養屬主的羊；同時亦因為她這份對主話語的順服,主耶穌就透過她叫當地的人在屬靈的需要上得著飽足.大半個世紀過去了,我相信祈姑娘自己也估計不到在含洸這個小鎮上,仍然有人在講說她的故事,而且會繼續傳流下去.

弟兄姊妹,今天主耶穌仍然向我們發出挑戰,說:「你們給他們吃吧!」而我們的回應會是怎樣呢?求聖靈在我們生命裡動工,讓我們明白儘管自己是軟弱的,不足的,有限制的,但只要我們願意順服主的吩咐,我們可以成為上主向人施恩的器皿.又求聖靈賜我們信心,加我們勇氣,讓我們可以堅定地回答:「好!我們給他們吃吧!願你的心意得到滿足,阿們.」



\newpage

\section{戰雲飯局}
\label{sec:64}
\textbf{林豪恩先生 - 戰雲飯局}
\newline
\newline
\hyperref[sec:63]{< < < PREV SERMON < < <}
~
\hyperref[sec:toc]{[返目錄]}
~
\hyperref[sec:65]{> > > NEXT SERMON > > >}


路加福音 14:1-24
\newline
\begin{longtable}{cl}
\hline
\hline
章節 & 經文 (和合本修訂版)\\
\hline
14:1 & \begin{tabularx}{0.7\textwidth}{X} 安息日，耶穌到一個法利賽人的領袖家裡去吃飯，他們就窺探他。 \end{tabularx} \\ \\ \relax
14:2 & \begin{tabularx}{0.7\textwidth}{X} 這時在他面前有一個患水腫病的人。 \end{tabularx} \\ \\ \relax
14:3 & \begin{tabularx}{0.7\textwidth}{X} 耶穌回答律法師和法利賽人，說：「安息日治病合不合法？」 \end{tabularx} \\ \\ \relax
14:4 & \begin{tabularx}{0.7\textwidth}{X} 他們卻不說話。耶穌扶著那人，治好了他，叫他走了。 \end{tabularx} \\ \\ \relax
14:5 & \begin{tabularx}{0.7\textwidth}{X} 耶穌對他們說：「你們中間誰有兒子或有牛在安息日掉在井裡，不立刻拉他上來呢？」 \end{tabularx} \\ \\ \relax
14:6 & \begin{tabularx}{0.7\textwidth}{X} 他們對這些事不能反駁。 \end{tabularx} \\ \\ \relax
14:7 & \begin{tabularx}{0.7\textwidth}{X} 耶穌見所請的客人選擇首位，就用比喻對他們說： \end{tabularx} \\ \\ \relax
14:8 & \begin{tabularx}{0.7\textwidth}{X} 「你被人請去赴婚宴，不要坐在首位上，恐怕主人請了比你尊貴的客人， \end{tabularx} \\ \\ \relax
14:9 & \begin{tabularx}{0.7\textwidth}{X} 請了你和他的那人前來，對你說：『請讓座給這一位吧。』你就羞羞慚慚地退到末位去了。 \end{tabularx} \\ \\ \relax
14:10 & \begin{tabularx}{0.7\textwidth}{X} 你被請的時候，去坐在末位上，好讓主人來對你說：『朋友，請上座。』那時，你在同席的人面前就有光彩了。 \end{tabularx} \\ \\ \relax
14:11 & \begin{tabularx}{0.7\textwidth}{X} 因為凡自高的，必降為卑；自甘卑微的，必升為高。」 \end{tabularx} \\ \\ \relax
14:12 & \begin{tabularx}{0.7\textwidth}{X} 耶穌又對請他的人說：「你準備午飯或晚餐，不要請你的朋友、弟兄、親屬和富足的鄰舍，免得他們回請你，你就得了報答。 \end{tabularx} \\ \\ \relax
14:13 & \begin{tabularx}{0.7\textwidth}{X} 你擺設宴席，倒要請那貧窮的、殘疾的、瘸腿的、失明的， \end{tabularx} \\ \\ \relax
14:14 & \begin{tabularx}{0.7\textwidth}{X} 你就有福了！因為他們沒有甚麼可報答你。到義人復活的時候，你要得到報答。」 \end{tabularx} \\ \\ \relax
14:15 & \begin{tabularx}{0.7\textwidth}{X} 同席的有一人聽見這些話，就對耶穌說：「在神國裡吃飯的有福了！」 \end{tabularx} \\ \\ \relax
14:16 & \begin{tabularx}{0.7\textwidth}{X} 耶穌對他說：「有人擺設大宴席，請了許多客人。 \end{tabularx} \\ \\ \relax
14:17 & \begin{tabularx}{0.7\textwidth}{X} 到了坐席的時候，他打發僕人去對所請的人說：『請來吧！樣樣都已齊備了。』 \end{tabularx} \\ \\ \relax
14:18 & \begin{tabularx}{0.7\textwidth}{X} 眾人異口同聲地推辭。頭一個對他說：『我買了一塊地，必須去看看。請你准我辭了。』 \end{tabularx} \\ \\ \relax
14:19 & \begin{tabularx}{0.7\textwidth}{X} 另一個說：『我買了五對牛，要去試一試。請你准我辭了。』 \end{tabularx} \\ \\ \relax
14:20 & \begin{tabularx}{0.7\textwidth}{X} 又有一個說：『我才娶了妻子，所以不能去。』 \end{tabularx} \\ \\ \relax
14:21 & \begin{tabularx}{0.7\textwidth}{X} 那僕人回來，把這些事都告訴了主人。這家的主人就發怒，對僕人說：『快出去，到城裡大街小巷，領那貧窮的、殘疾的、失明的、瘸腿的來。』 \end{tabularx} \\ \\ \relax
14:22 & \begin{tabularx}{0.7\textwidth}{X} 僕人說：『主啊，你所吩咐的已經辦了，還有空位。』 \end{tabularx} \\ \\ \relax
14:23 & \begin{tabularx}{0.7\textwidth}{X} 主人對僕人說：『你出去，到大街小巷強拉人進來，坐滿我的屋子。 \end{tabularx} \\ \\ \relax
14:24 & \begin{tabularx}{0.7\textwidth}{X} 我告訴你們，先前所請的人沒有一個可以嘗到我的宴席。』」 \end{tabularx} \\ \\
[1ex]
\hline
\hline
\end{longtable}
日期:2010年9月19日

主題:戰雲飯局

証道:林豪恩先生

經文:路加福音14:1-24

引言有些飯局是溫馨情切的,例如一家人吃團年飯；有些飯局是歡愉暢快的,例如良朋知己把酒談心；有些飯局是尷尬勉強的,例如某電視台高層親自主持飯局節目,被邀者不出席好像落對方面子,出席就等於被公開地窺探私隱；有些飯局是暗藏殺機的,例如公元前206年楚漢相爭時期項羽為暗殺劉邦而擺設著名的鴻門夜宴.今天的福音經課,提及耶穌參與一個飯局.這個飯局看來不像是溫馨情切的家人聚餐,也不似是把酒言歡的良朋共聚；這個飯局雖然不像鴻門放宴一樣有即時的生命危險,卻也是戰雲密佈,火藥味甚濃,我們暫且稱之為「戰雲飯局」.

經文

1安息日,耶穌到一個法利賽人的首領家裡去吃飯,他們就窺探他.2在他面前有一個患水臌的人.3耶穌對律法師和法利賽人說:「安息日治病可以不可以?」4他們卻不言語.耶穌就治好那人,叫他走了；5便對他們說:「你們中間誰有驢或有牛,在安息日掉在井裡,不立時拉牠上來呢?」6他們不能對答這話. 7耶穌見所請的客揀擇首位,就用比喻對他們說:8「你被人請去赴婚姻的筵席,不要坐在首位上,恐怕有比你尊貴的客被他請來；9那請你們的人前來對你說:『讓座給這一位吧!』你就羞羞慚慚地退到末位上去了.10你被請的時候,就去坐在末位上,好叫那請你的人來對你說:『朋友,請上座.』那時,你在同席的人面前就有光彩了.11因為,凡自高的,必降為卑；自卑的,必升為高.」12耶穌又對請他的人說:「你擺設午飯或晚飯,不要請你的朋友,弟兄,親屬,和富足的鄰舍,恐怕他們也請你,你就得了報答.13你擺設筵席,倒要請那貧窮的,殘廢的,瘸腿的,瞎眼的,你就有福了!14因為他們沒有甚麼可報答你.到義人復活的時候,你要得著報答.」 15同席的有一人聽見這話,就對耶穌說:「在上帝國裡吃飯的有福了!」16耶穌對他說:「有一人擺設大筵席,請了許多客.17到了坐席的時候,打發僕人去對所請的人說:『請來吧!樣樣都齊備了.』18眾人一口同音地推辭.頭一個說:『我買了一塊地,必須去看看.請你准我辭了.』19又有一個說:『我買了五對牛,要去試一試.請你准我辭了.』20又有一個說:『我才娶了妻,所以不能去.』21那僕人回來,把這事都告訴了主人.家主就動怒,對僕人說:『快出去,到城裡大街小巷,領那貧窮的,殘廢的,瞎眼的,瘸腿的來.』22僕人說:『主啊,你所吩咐的已經辦了,還有空座.』23主人對僕人說:『你出去到路上和籬笆那裡,勉強人進來,坐滿我的屋子.24我告訴你們,先前所請的人沒有一個得嘗我的筵席.』」

釋經與應用這場「戰雲飯局」共分四個回合,個別回合好像是獨立的,原來又是互相關連的.爭論的話題以梅花間竹的編排出現:第一個回合關於主人,第二個回合關於客人,第三個回合又關於主人,第四個回合又關於客人.

第一個回合(關於主人):哪裡有待客如己,待己如客的主人呢?

1安息日,耶穌到一個法利賽人的首領家裡去吃飯,他們就窺探他.2在他面前有一個患水臌的人.3耶穌對律法師和法利賽人說:「安息日治病可以不可以?」4他們卻不言語.耶穌就治好那人,叫他走了；5便對他們說:「你們中間誰有驢或有牛,在安息日掉在井裡,不立時拉牠上來呢?」6他們不能對答這話.

這次宗教領導的飯局,有兩位「稀客」.其一是經常批評宗教領導的異見人士──耶穌.其二是沒有資格也不合適地出現在這尊貴午宴──一個患水臌的病人.我們可以想像,對於作為客人的耶穌來說,這個飯局很可能是一個「局」,對那位患水臌的病人來說,他很可是這個「裝耶穌的局」的「餌」.這就是律法師和法利賽人擺出的戰陣.

耶穌會如何回應呢?如果耶穌對眼前的病人視若無睹,就是見死不救；如果耶穌出手醫病,就在一眾宗教領導面前公然違反安息日條例.行動之前,耶穌向法利賽人提問:「安息日治病可以不可以?」對於法利賽人來說,這是多餘的問題,所以他們拒絕回答.可能他們心想:「你做了之後就知道可以還是不可以!」

耶穌果然出手,醫好了那位患水臌的病人.那位病人的戲份完成,就退出舞台了.在法利賽人提出指控之前,耶穌再次連消帶打地提問:「你們中間誰有驢或有牛,在安息日掉在井裡,不立時拉牠上來呢?」法利賽人和律法師無言以對.為甚麼無語以對呢?因為大家心裡都有矛盾及尷尬的答案.

耶穌的第一個問題:「安息日治病可以不可以?」是原則性的問題,法利賽人和律法師按著條文規例來回答,當然是不可以.耶穌的第二個問題:「你們中間誰有驢或有牛,在安息日掉在井裡,不立時拉牠上來呢?」耶穌這一問,把一個原則性的問題變成一個處境性的問題,並且具體地問到對手身上,答案就會不同.香港旅行團在菲律賓遭槍手劫持的時候,香港特首曾經致電菲律賓總統,很可能要求總統特別重視及處理事件.菲律賓總統沒有接聽電話,也沒有回電.菲律賓警方的處理方法真的非常特別──特別令人失望.或者,菲律賓總統,甚至相關官員都認為這些事情應該按既定的政策及程序回應,無需特別關注.不過,如果總統或官員的家人也在旅遊車內呢?還會堅持按既定程序處理嗎?還是會特事特辦呢?法利賽人和律法師面對耶穌的兩個問題,答案為甚麼不一致呢?因為患病的是「那個人」,而掉在井裡的,是「我的」驢和「我的」牛.處理別人的事情,就運地用原則和程序,處理自己的事情,則視作個別處境和例外情況.如此看來,律法師和法利賽人所制定及執行的律例典章,條文和規則,原來只用來規範別人,並不應用於自己.

律法師及法利賽人本來的策略是以「安息日治病」為武器攻擊耶穌,但耶穌卻把焦點轉移到律法師及法利賽人的雙重標準及誠信問題.「安息日可否治病」可以是一個議題,不過,耶穌卻指出,那些制定及執行規則的人實行「安息日不可以治你的病,而只可以治我的病」這種雙重標準是一個更嚴重的問題.這是公平的問題,是誠信的問題,是信仰的問題,他們違反了愛鄰舍如同自己的律法核心精神.難怪耶穌在馬太福音23章4節指責律法師和法利賽人說:「他們把難擔的重擔捆起來,擱在人的肩上,但自己一個指頭也不肯動.」

人們親疏有的傾向,是可以理解的.有一個笑話這樣說:「有一位太太對另一位太太說,我的女兒很幸運,她嫁到一位好男人,只要她一開聲,丈夫就做妥所有家務.但我的兒子就不幸了,他娶了一個懶女人,甚麼家務也要我的兒子做.」不過,若果掌管公共權力及公共資源的人處事雙重標準,親疏有別,甚至徇私枉法,這就不是笑話,而是罪惡.律法師和法利賽人或者都不一是面目可憎的人,但他們把自己和群眾分開,自己和群眾彷如是兩種不同地位及不同等級的人:「我們負責編制律法,你們負責學習律法；我們負責執行律法,你們負責遵守律法.我們不畏艱辛,承擔這吃力不討好的任務,並不是為我們自己,而是為了你們.因為,我們和你們是不同的.」

如此的心態或言論,是否似曾相識?在哪裡聽聞過?可能在官場,可能在商界,可能在學校,可能在教會,甚至可能在我們自己的心裡.

有一位立法會議員評論香港現行的教育制度時這樣說:「我認為要看香港的教育制度是否成功,可以看政府高官的兒女是否在香港的教育制度下成長,成才．．．如果連政府高官也都寧願送子女出國,而不願他們留下接受他們自己所建立的教育制度,那就說不過去了──難道對自己所訂立及推廣的制度也沒有信心?」(涂謹申 HYPERLINK ''\url{http://www.tokunsun.org.hk/atoi/page8.htm and quot}; \url{http://www.tokunsun.org.hk/atoi/page8.htm})議員的質疑雖然沒有提出證據,但足以引發市民的不公平感覺及搖動對制度的信任.

耶穌指責律法師及法利賽人作為宗教的領袖,「規別人的矩,行自己的方便」.這些雙重標準,假冒為善的行為,沒有人會認為正確.不過,在我們伸出指責的手指的時候,我們會否發現,如果有足夠的條件,我們心底裡會否也有「以政策程序迎接別人,留彈性例外給自己」的雙重標準呢?到底哪裡有待客如己,待己如客,無分彼此的主人呢?

2．第二場回合(關於客人):哪裡有自甘卑微的客人呢?

7耶穌見所請的客揀擇首位,就用比喻對他們說:8「你被人請去赴婚姻的筵席,不要坐在首位上,恐怕有比你尊貴的客被他請來；9那請你們的人前來對你說:『讓座給這一位吧!』你就羞羞慚慚地退到末位上去了.10你被請的時候,就去坐在末位上,好叫那請你的人來對你說:『朋友,請上座.』那時,你在同席的人面前就有光彩了.11因為,凡自高的,必降為卑；自卑的,必升為高.」

首度交鋒,法利賽人及律法師被技術擊倒,尷尬之際,幸好正是入席之時,大家乘機把焦點轉移離開耶穌,投入「爭凳仔,霸頭位」的遊戲.法利賽人和律法師知難而退,想不到耶穌竟然在這個時候出招,用比喻批評他們在筵席上爭頭位的失禮行為.

既然我與你的不同如此重要,我們與你們的分別如此敏感,哪麼,都底你重要還是我重要呢?是你的民族優越還是我的民族優越呢?在筵席上,到底誰是VIP呢?如果你是VIP,他是VVIP,那麼,我一定是VVVIP.每個人都渴望自己被重視,希望自己是最重要的那一位.每一個民族,都宣稱自己優秀,值得在地球上佔一席位.反過來說,用點心做比喻,如果我不是頂點,也希望是特點,如果不是特點,也希望是大點,如果不是大點,也希望是中點,怎會甘心做小點呢?不被重視的感覺實在太難受了,多餘及無用的感覺實在太傷害了.人們怎能不為證明自己重要而競爭呢?怎能不為證明自己民族的優越而大動肝火,甚至以其他民族為代價呢?

耶穌用比喻告訴筵席上「爭凳仔,霸頭位」的客人,「爭得到,也未必好」,甚至會適得其反,因為筵席的座位,是主人安排的.作為客人,如果漠視主人的安排,自己去爭去搶,到頭來,「爭到凳仔」,也不一定「霸到頭位」,「霸到頭位」,也不一定坐得舒服安穩.

耶穌用一個聽起來充滿矛盾的格言作為第二回合的總結:「凡自高的,必降為卑,自卑的,必被升高.」正在玩「爭凳仔,霸頭位」遊戲的律法師及法利賽人被「串」個正著,再次啞口無言.不過,他們心裡可能會感到不服氣.現實上,的確有人在前,有人在後,有人名成利就,有人三餐不飽.生存在這個世界,若然不用盡每一分力「爭凳仔,霸頭位」,落在人後的時候,有人同伴嗎?三餐不飽的日子,有人照顧嗎?來到這個世界,誰能做自甘卑微的客人呢?

3．第三回合(關於主人):哪裡有純粹施恩的主人呢?耶穌對筵席上「爭凳仔,霸頭位」的客人「開火」後,轉而挑戰筵席的主人.

12耶穌又對請他的人說:「你擺設午飯或晚飯,不要請你的朋友,弟兄,親屬,和富足的鄰舍,恐怕他們也請你,你就得了報答.13你擺設筵席,倒要請那貧窮的,殘廢的,瘸腿的,瞎眼的,你就有福了!14因為他們沒有甚麼可報答你.到義人復活的時候,你要得著報答.」

在「爭凳仔,霸頭位」的遊戲中,總有人「爭到凳仔」,也有人「霸到頭位」.鄧小平1985年10月23日會見美國客人時講了影響深遠的話:「我們的政策是讓一部分人,一部分地區先富起來」.對於國家當時的局勢來說,這是劃時代的宣言,突破均貧的困局.接著的二十多年,一部分人富起來的速度及幅度,可能大大超出鄧小平的所想所求.《福布斯》雜誌公布2010年全球富豪排行榜,中國有64人上榜,人數僅次於美國.先富起來的一部分人以及先富起來的一部分地區,應該如何運用他們的財富及資源呢?鄧小平當時有這樣期望:「一部分地區,一部分人可以先富起來,帶動和幫助其他地區,其他的人,逐步達到共同富裕.」到目前為止,鄧小平這美好的盼望連政府也承認未能如願,也不能保證甚麼時候才能實現.

無論是耶穌時代的人,還是現代的人,無論是耶穌時代的宗教領袖,還是現代中國的富豪,都信奉同一個道理,玩同一套遊戲,就是財富及資源是用來賺取更多財富及資源的.原來,在未有股票市場的時候,人們已經懂得投資,所以,擺設筵席的時候,當然是請朋友,弟兄,親屬,和富足的鄰舍為客,而不會請貧窮的,殘廢的,瘸腿的和瞎眼的.

期望先富起來的一部份人和一部份地區自動自覺地帶動和幫助其他地區和其他的人,達到共同富裕,鄧小平的期望會否太過理想化呢?要求擺設筵席的主人不要請朋友,弟兄,親屬,和富足的鄰舍這些有機會回報的人為客,反而要請貧窮的,殘廢的,瘸腿的和瞎眼的這些沒有能力回報的人為座上賓,耶穌的指示會否也是太過理想化呢?還是人類太過分得清你與我,你們與我們,而把人類本身作為手足,作為同伴的現實忘記了呢?還是那些先富起來的,把他們的手足忘記了,把他們的同伴忘記了,把他們那些落在後面的兄弟當作與自己截然不同的另一類人看待呢?

當人們都懂得為財富和資源尋找有回報的投資的時候,哪裡有無條件付出的主人呢?哪裡有純粹施恩的主人呢?

4．第四回合(關於客人):誰是天國筵席的客人呢?耶穌擊退了法利賽人和律法師的攻擊,接著批評了同場的客人,又挑戰了筵席的主人,場面應該非常尷尬.同席有一個人聽到耶穌提及「義人復活的時候」,聯想到將來彌賽亞來臨的時候,就乘機把話題由談論在場客人及主人的尷尬局面轉移到遙遠的將來.

15同席的有一人聽見這話,就對耶穌說:「在上帝國裡吃飯的有福了!」

怎知耶穌沒有就此罷休,更打蛇隨棍上,承接著這位客人的話題而講另外一個比喻.耶穌對著律法師和法利賽人講這個比喻,差不多等於「反台」.

16耶穌對他說:「有一人擺設大筵席,請了許多客.17到了坐席的時候,打發僕人去對所請的人說:『請來吧!樣樣都齊備了.』18眾人一口同音地推辭.頭一個說:『我買了一塊地,必須去看看.請你准我辭了.』19又有一個說:『我買了五對牛,要去試一試.請你准我辭了.』20又有一個說:『我才娶了妻,所以不能去.』21那僕人回來,把這事都告訴了主人.家主就動怒,對僕人說:『快出去,到城裡大街小巷,領那貧窮的,殘廢的,瞎眼的,瘸腿的來.』22僕人說:『主啊,你所吩咐的已經辦了,還有空座.』23主人對僕人說:『你出去到路上和籬笆那裡,勉強人進來,坐滿我的屋子.24我告訴你們,先前所請的人沒有一個得嘗我的筵席.』」

相傳耶穌時代的巴勒斯坦人請客,分兩次邀請.首先發出請帖,客人回覆後確實名單.到筵席準備就緒的時候,再派人去邀請答應了出席的客人.因此,在這個比喻中推辭的客人,其實是已經應邀的客人.這些客人「臨時甩底」,是情有可原,還是罪無可恕呢?頭一個的理由是::『我買了一塊地,必須去看看.』人們買棵菜,都看過才買,何況買塊地,怎會買了才去看的呢?另一個人的理由是:『我買了五對牛,要去試一試.』同樣道理,人們買隻小蒼鼠作為寵物,都會先撫摸一下,何況買五對牛,怎會買了才去試呢?另一個人的理由是:『我才娶了妻,所以不能去.』舊約的確有條文寫著娶妻頭一年不用當兵,但不致於吃頓飯也不能去吧!總括而言,不是不能赴宴,而是不想赴宴.如果不想赴宴,又何必答應邀請呢?難怪筵席的主人大動肚火.既然筵席已經預備好,而本來應約的客人無福消受,就給其他人意外驚喜吧!

聰明的律法師和法利賽人應該聽得明白耶穌這個比喻的意思.「在上帝國裡吃飯的有福了!你們說得對,不過,你們卻無福消受.」耶穌在法利賽人首領的筵席上向法利賽人和律法師說這樣的話,無疑等於「反台」.

總結最後,讓我們回看耶穌與律法師及法利賽人在這「戰雲飯局」所帶出的幾個問題:

誰是待客如己,待己如客的主人呢?人們雖然推崇「愛人如己」的理想,但難免會有「親疏有別」傾向.只有天父是那位待客如己,待己如客,亳不徇私的主人.何以見得呢?我們被安排生活於這個世界裡面,我們的兒女生活於這個世界裡面,天父竟然也安排祂的獨子耶穌基督生活於我們及我們的兒女所生活的同一個世界裡面.人類生活於這個世界所遭遇到的,祂的獨子耶穌基督也沒有幸免,無論這些遭遇是愉快的,還是痛苦的,美好的,還是邪惡的.因此,我們可以相信,天父有多愛惜祂的獨子耶穌基督,就有多愛惜我們；有多重視祂的獨子耶穌基督,就有多重視我們.

哪裡有自甘卑微的客人呢?天父的獨子耶穌基督來到我們所生活的世界作客,沒有證明自己重要,也沒有爭奪首位,反而甘心作一位卑微的客人,甘心作一位被視作多餘的客人.腓立比書2章如描述這位客人:「6他本有上帝的形像,不以自己與上帝同等為強奪的；7反倒虛己,取了奴僕的形像,成為人的樣式；8既有人的樣子,就自己卑微,存心順服,以至於死,且死在十字架上.」耶穌作為客人知道最終的評價並不在於其他在世作客的人,而是這個世界的主.如何擺脫位置及名望的壓力呢?這是信仰的問題,耶穌的示範是信任「上帝要把自高的降為卑微,又高舉自甘卑微的人」.(現中14:11)

哪裡有不求回報,純粹施恩的主人呢?「奇異恩典」電影中有一個令人動容的片段:年老的牛頓牧師雙眼失明,在偏遠的小教堂裡打掃洗抹的時候,在國會裡推動廢除販賣奴隸條約的威伯福斯議員來探訪他.牛頓說:「比起我眼明時盲目的罪行,我的內心從未如此澄明!」牛頓提出將自己的罪惡寫成呈堂證據,以支持廢止奴隸制度.他說:「只有上帝的恩典,才能接受像我這樣無理,褻瀆,又做過奴隸交易的人,並使我轉變成為神的孩子.」《奇異恩典》歌詞中「前我失喪,今被尋回,瞎眼今得看見」正是這首詩歌的作者的生命寫照.如牛頓一樣,不少人都親身體驗到上帝是那位充滿恩典的主人,樂意施恩給所有人.如何能夠突破回報的吸引和壓力,學習做慷慨施予的主人呢?耶穌的建議仍然是信仰性的:「在義人復活的時候,上帝要新自報答你.」(現中14:14)

誰是天國筵席的客人呢?最後,誰是天國筵席的客人呢?剛買了一塊地的可以是,剛買了五頭牛的可以是,剛成家立室的可以是,貧窮的,殘廢的,瞎眼的,瘸腿的可以是,路過的可以是,籬笆外面的都可以是.筵席的主人似乎比我們更熱心,祂彷彿對我們說:Come, be my guest!



\newpage

\section{個人理財與社會公義}
\label{sec:65}
\textbf{賴品超教授 - 個人理財與社會公義}
\newline
\newline
\hyperref[sec:64]{< < < PREV SERMON < < <}
~
\hyperref[sec:toc]{[返目錄]}
~
\hyperref[sec:66]{> > > NEXT SERMON > > >}


阿摩司書 8:4-7
\newline
\begin{longtable}{cl}
\hline
\hline
章節 & 經文 (和合本修訂版)\\
\hline
8:4 & \begin{tabularx}{0.7\textwidth}{X} 你們這些踐踏貧窮人、使這地困苦人衰敗的，當聽這話！ \end{tabularx} \\ \\ \relax
8:5 & \begin{tabularx}{0.7\textwidth}{X} 你們說：「初一幾時過去，我們好賣糧；安息日幾時過去，我們好擺開穀物；我們要把伊法變小，把舍客勒變大，以詭詐的天平欺哄人， \end{tabularx} \\ \\ \relax
8:6 & \begin{tabularx}{0.7\textwidth}{X} 用銀子買貧寒人，以一雙鞋換貧窮人，把壞的穀物賣給人。」 \end{tabularx} \\ \\ \relax
8:7 & \begin{tabularx}{0.7\textwidth}{X} 耶和華指著雅各的驕傲起誓說：「他們這一切的行為，我必永遠不忘。 \end{tabularx} \\ \\
[1ex]
\hline
\hline
\end{longtable}
提摩太前書 2:1-7
\newline
\begin{longtable}{cl}
\hline
\hline
章節 & 經文 (和合本修訂版)\\
\hline
2:1 & \begin{tabularx}{0.7\textwidth}{X} 所以，我勸你，首先要為人人祈求、禱告、代求、感謝； \end{tabularx} \\ \\ \relax
2:2 & \begin{tabularx}{0.7\textwidth}{X} 為君王和一切在位的，也要如此，使我們能夠敬虔端正地過平穩寧靜的生活。 \end{tabularx} \\ \\ \relax
2:3 & \begin{tabularx}{0.7\textwidth}{X} 這是好的，在我們的救主神面前可蒙悅納。 \end{tabularx} \\ \\ \relax
2:4 & \begin{tabularx}{0.7\textwidth}{X} 他願意人人得救，並得以認識真理。 \end{tabularx} \\ \\ \relax
2:5 & \begin{tabularx}{0.7\textwidth}{X} 因為只有一位神，在神和人之間也只有一位中保，是成為人的基督耶穌。 \end{tabularx} \\ \\ \relax
2:6 & \begin{tabularx}{0.7\textwidth}{X} 他獻上自己作人人的贖價；在適當的時候這事已經證實了。 \end{tabularx} \\ \\ \relax
2:7 & \begin{tabularx}{0.7\textwidth}{X} 我為此奉派作傳道，作使徒，在信仰和真理上作外邦人的教師。我說的是真話，不是說謊。 \end{tabularx} \\ \\
[1ex]
\hline
\hline
\end{longtable}
路加福音 16:1-13
\newline
\begin{longtable}{cl}
\hline
\hline
章節 & 經文 (和合本修訂版)\\
\hline
16:1 & \begin{tabularx}{0.7\textwidth}{X} 耶穌又對門徒說：「某財主有一個管家，有人向主人告管家浪費他的財物。 \end{tabularx} \\ \\ \relax
16:2 & \begin{tabularx}{0.7\textwidth}{X} 主人叫他來，對他說：『我聽到了，你做的是甚麼事？把你所經管的交代清楚，你不能再作我的管家了。』 \end{tabularx} \\ \\ \relax
16:3 & \begin{tabularx}{0.7\textwidth}{X} 那管家心裡說：『主人辭我，不用我再作管家，我將來做甚麼呢？鋤地嘛，沒有力氣；討飯嘛，怕羞。 \end{tabularx} \\ \\ \relax
16:4 & \begin{tabularx}{0.7\textwidth}{X} 我知道怎麼做，好叫人們在我不作管家之後，接我到他們家裡去。』 \end{tabularx} \\ \\ \relax
16:5 & \begin{tabularx}{0.7\textwidth}{X} 於是他把欠他主人債的，一個一個地叫了來，問頭一個說：『你欠我主人多少？』 \end{tabularx} \\ \\ \relax
16:6 & \begin{tabularx}{0.7\textwidth}{X} 他說：『一百簍油。』管家對他說：『拿你的賬，快坐下，寫五十。』 \end{tabularx} \\ \\ \relax
16:7 & \begin{tabularx}{0.7\textwidth}{X} 他問另一個說：『你欠多少？』他說：『一百石麥子。』管家對他說：『拿你的賬，寫八十。』 \end{tabularx} \\ \\ \relax
16:8 & \begin{tabularx}{0.7\textwidth}{X} 主人就誇獎這不義的管家做事精明，因為今世之子應付自己的世代比光明之子更加精明。 \end{tabularx} \\ \\ \relax
16:9 & \begin{tabularx}{0.7\textwidth}{X} 我又告訴你們，要藉著那不義的錢財結交朋友，到了錢財無用的時候，他們可以接你們到永遠的住處去。 \end{tabularx} \\ \\ \relax
16:10 & \begin{tabularx}{0.7\textwidth}{X} 人在最小的事上忠心，在大事上也忠心；在最小的事上不義，在大事上也不義。 \end{tabularx} \\ \\ \relax
16:11 & \begin{tabularx}{0.7\textwidth}{X} 若是你們在不義的錢財上不忠心，誰還把那真實的錢財託付你們呢？ \end{tabularx} \\ \\ \relax
16:12 & \begin{tabularx}{0.7\textwidth}{X} 如果你們在別人的東西上不忠心，誰還把你們自己的東西給你們呢？ \end{tabularx} \\ \\ \relax
16:13 & \begin{tabularx}{0.7\textwidth}{X} 一個僕人不能服侍兩個主；他不是恨這個愛那個，就是重這個輕那個。你們不能又服侍神，又服侍瑪門。」 \end{tabularx} \\ \\
[1ex]
\hline
\hline
\end{longtable}
日期:2010年9月26日

主題:個人理財與社會公義

証道:賴品超教授

經文:阿摩司書8:4-7;詩篇113;提摩太前書2:1-7;路加福音16:1-13

我不是財務方面的教授,不會教人如何投資,保值,賺更多的錢；今天只就著這幾段的經文,分享一下我們應如何運用我們的資產,用得合乎上帝的心意.

路加福音16:1-13不義的管家的比喻──表面上似在鼓勵不恰當的行為,此人平時很可能已是"財費主人的財物"(故為不義的管家),被吞發也不反擊,大概是自知其身不正,無法掩飾或辯駁,只能謀定後路,而更不義的是,在剩餘的日子,仍以權謀私,做假數,出賣主人的利益,討好久債的,希望日後有個人的回報.聖經不是鼓勵人欺騙,它肯定此不義的管家是不義的,是今世之子(俗稱"世界仔")而不是光明之子,只是用此比喻來指出,這不義的管家的"聰明",就是怎樣用錢,就是運用短暫,可朽壞的金錢,得到長遠,永恆的回報(16: 9).在16:10-11更明確地點出,我們要在最小的事上忠心,包括怎樣使用錢.

不時被一些經紀接觸,游說買基金,物業,要保值,增值,單靠公積金,退休金,不夠.最近,我對那些經紀說,如果計過條數,如無意外,若能介做到60歲退休,退休金及現有的投資,已足夠退休,那又如何?有個經紀講,要預鬆一點云云；有另一個經紀,我覺答得比較好的.他回答是,投資不一定為退休,多些投資回報,可以捐多些錢給教會或做其他善事,他的一些顧客也是這樣,聖經教人要為做好管家.(我不知那句做好管家是否他的僱客告訴他的,也我不知他是否教徒,更不知他是否知道我是教徒.)

曾想過在講道題目用"私人財產" (Private Property),最後改為"個人理財"(Personal Finance),主要有兩方面的考慮.(1)對於我們的財產,正如路加福音所講,我們是受託付負責"管理"這些財產的管家,但非真正"擁有"這些財富的主人."理財"強調的是"管理",若用"財產"則更強調的是"擁有"；而"property"易聯想起房地產,這是我們不少人所沒有的,所以還是用"理財"更適合.(2)至於如何管理／理財,我想提出的是,it is a personal but not private matter.我們有時將"personal／個人"等同"private／私人".若是"私人"的事情,其他人無權過問或規管."私人(擁有的)物業,非請勿進"."私人"的東西可以隨意運用,不必問任何人；一講是私人的事,就是不必講給你聽,別人也無權過問.但英文的personal(而不是中文的"個人")與private不同.該字源於拉丁文的 Persona,是神學上的一個核心概念,"三位一體"中的"位格",是行動的中心.Person 可以是natural person or artificial person(法人,亦稱fictitious person; juristic person; legal person; 中文"財團法人"; moral person),但總是能行動及負責任的.正如路加福音所講,我們是受託付負責管理這些財產的管家,但卻並非真正擁有這些財富的主人；理財是 personal 而非private的事情,所指的是,不是"私人"可以隨意運用,不必問任何人；而是個人要為此而負上責任.

那麼我們要怎樣使用金錢才是忠心,才是義的管家?對此,很多教會也許會馬上提供十分具有的建議,包括多些奉獻給教會.金句是:努力賺錢,努力儲錢,努力奉獻.但如果我們將路加福音這段經文,與今天所讀的其他三段經文連在一起看,可能會有另一方面的提醒.

舊約聖經,如Psalm 113:7,表明上帝要從抬舉貧寒／窮乏的人.這種對貧窮人的關注,並非局限在給予即時的幫助,"餓了給他吃",而是轉到對社會公義的訴求,因為發覺貧窮的出現,也許有一些是因為個人的懶惰(在簽言中有不少勸喻),但更多的情況是因受到欺壓.阿摩司書 8: 4-7,更堅決批判欺詐,不公平的交易,批評,囤積居奇(炒ｘ炒ｙ),"呃稱"(發水樓?),賣過期食物等,是"吞吃窮乏人","使困苦人衰敗".除了反面的批評外,聖經也表達了去對一個安定的社會的嚮往,甚至認為信徒也應為此而努力 (提前 2:1-7, 尤2).

再回到路加福音的經文,一個值得問的問題是,是否及如何可以,將"不義的錢財"用於公義事情上,使社會更加公義,包括讓貧窮的人可以免於飢餓等的痛苦,甚至進一步可以脫貧.(這正是今天的講題的由來.)也就是說,我們如何理財,不單要向上帝交代(answerable to God),也要對社會負責(responsible to the society),回應社會的需要,尤其社會公義.值得問,努力奉獻給教會後,教會是否也有做足夠的貢獻,促進社會公義尤其扶助貧窮人的事上呢?個人方面呢?似乎沒有既定方式.

以往教會過去做得最多的可能是慈惠救災的工作,有人饑便給魚,後來發現這不足夠,並長貧難顧,改為不是給魚,而是教捉魚(例如:教育,發展社企),但最近又發現,若不保護生態,保持有魚,跟本無魚可捉；有不少人有技術但無用武之地,因此要有一個良好的就業環境,合理的工資等.(也許更應談的不是最低工資,而是家庭工資／維生工資.)也有人懷疑,直接給了人財物,那人不懂理財,用不得其所,未必是最好的方法,反而保持一部份繼續投資,可能更好,可更持續.因此教會也要有投資.有些教會在投資上,堅持不參與與煙草及賭業,寧選社企或環保事業.甚至,政治捐獻也可以是一種,支持某政黨或某些致力改變社會制度的機構,使社會有更好的制度,也可以是促進社會公義.

在這些事上,聖經並沒有既定的具體指引,只有各自考量,人人不同.

最近看一本小書,是教父的著作,覺得很有啟發.St. Basil the Great, On Social Justice (2009).　當中有不少神學上的思想,我只講一小部份.聖巴西略(St. Basil the Great, ca 329-378, 另譯大巴西流)的一些講章,是由後人編輯而是,基本上是環繞社會公義,貧窮的問題.聖巴西略除了是著名的神學家,稍為讀過神學的也應該聽過他是Cappadocian fathers之一,也是一位主教(Bishop of Caesarea今之土耳其).他曾變賣所有財物分給窮人,之後到曠野隱修.後來,他改變那種脫離座世的,個人的隱修,並發展出一種群居,並且著重對社會服務的修院制度,除自給自足外,更照顧飢餓,無家可歸,患病等的人.這除了成為日後東方教會修院制度的藍本外,更可能是醫院的開端.對耶穌基督呼召少年財主散盡家財跟從主,當時有人認為應"屬靈"地解釋(Clement of Alexandria),也有人認為照字面做(St. Anthony),更有人認為是只應用於少數追求完全的人,一般信徒不適用,但他認為所有信徒應不斷學習以此為目標.他有一句名句:「你未曾食用的麵包,就是饑渴者的麵包；你懸掛在衣櫃裡的長袍,就是赤身露體者衣服；你不穿的鞋子,就是赤足者的鞋履；你保持深鎖的錢財,就是窮人的錢財；你不願踐行的愛德,就是諸多你所干犯的罪過.」(中譯引自網上資料) 聖巴西略在他的講道中,不少地方是批評當時的富人,尤其是放債的人(也勸人不要借錢,以免憂慮,更免成為奴隸或被羞辱,因自由與尊嚴更可貴),未有好好幫助貧窮人,他們的財富是來自對貧窮人的欺壓.在St. Basil the Great, On Social Justice (2009)一書中有一篇是聖巴西略的追隨者寫的,當中提及"Acts of charity made from unjust gains are not acceptable to God; nor are those who refrain from injustice praiseworthy if they do not share what they have." 簡言之就是 ``possess with justice and dispense with mercy" (p.103).取之有義,用之有憐,可能比中國人講的"取之有道,用之有道"更為精彩. 另有一段很有趣的文字(pp.55-56),講有些人不想在生前捐出財產,以為立遺囑便可死後捐出,他提醒,只要稍一更改,落入別人手中,(我想大家都可能會想到一些相關的例子),最好還是生前做,可以肯定財物是用於自己所指定的用途,直接享受施予的快樂. 我想起的,不僅是千億女富豪的爭產案,更是前兩天新聞.美國兩大富豪想邀請某些中國的富豪出宴會,有些不想應邀,美國兩位富豪,呼籲富豪捐最少一半,一種"半裸捐"(相對於捐盡所有的"裸捐").某些中國富豪,既不想承諾"半裸",也不想公開說不想"半裸",為免尷尬不如不去.

理財／用錢只是(路加福音講的)"小事",但我們應從這小事上學好.今天我們當如何賺錢,如何理財,有一天都要向上帝交代,究竟我們已經開始學習是否按上帝的心意來用錢?



\newpage

\section{普世的福音}
\label{sec:66}
\textbf{姚西依博士 - 普世的福音}
\newline
\newline
\hyperref[sec:65]{< < < PREV SERMON < < <}
~
\hyperref[sec:toc]{[返目錄]}
~
\hyperref[sec:67]{> > > NEXT SERMON > > >}


日期:2010年10月3日

主題:普世的福音

經文:使徒行傳10

証道:姚西依博士

請大家再一次把聖經翻到使徒行傳第10章,就是在新約的第140頁,今天我們一起來思想,使徒行傳第10章所講彼得和哥尼流的故事.

剛才我們讀了使徒行傳第10章1至16節,這個故事我相信很多朋友都很熟習,內容是說哥尼流是位沒有受過割禮,他不是位猶太人,也沒有遵守摩西律法當中對食物聖俗的規則,他是位外邦人,作為外邦人,可不可以接受耶穌基督,這是彼得在這個故事當中所面對的最大掙扎,也是我們今天要思想的核心故事.

彼得在掙扎時思考要不要領這位外邦人信主,這個信仰是不是也對外邦人開放的,在第16節我們看到彼得獲得一個來自神的異象,指猶太人對於食物清潔的問題上應該要有什麼的態度,請大家看經文第10章的34和35節「彼得就開口說,我真看出神是不偏待人,原來各國中,那敬畏主行義的人,都為主所悅納.」再從44節一直讀到48節「彼得還說這話的時候,聖靈降在一切聽道的人身上.那些奉割禮和彼得同來的信徒,見聖靈的恩賜也澆在外邦人身上,就都希奇,因聽見他們說方言,稱讚　神為大.於是彼得說,這些人既受了聖靈,與我們一樣,誰能禁止用水給他們施洗呢.就吩咐奉耶穌基督的名給他們施洗.他們又請彼得住了幾天.」

有的解經家說,這些話語是彼得整個宣教工作中,其中一個最重要的事件,其實這也是早期教會所面對最重大問題的挑戰之一,就是怎樣去處理耶穌基督的福音和猶太傳統的關系?怎樣去對待教會上越來越多的外邦信徒,作為耶穌基督的門徒,是不是要遵守一個特定的文化傳統,一個特定的生活方式,還是可以保留自己的生活和文化方式,這個福音和文化傳統之間的關系,是不是也要成為猶太人?這是五旬節之後,早期教會所經歷過的決定性教會未來命運和走向的關鍵最大事件.耶穌基督被釘十字架復活升天之後,最早期的這批門徒都聚集在耶路撒冷,以聖殿為中心來崇拜和宣道.

這些早期的教會的門徒,包括耶穌基督本人都是什麼人?在身份上全都是猶太人,那時也有小數歸依了猶太教的外邦人,可是這早期的群體他們猶太人的身份是非常明確的,他們行割禮,遵守猶太的律法,和種種的規定,包括食物清潔的規定,期間我們看不出早期的教會在這個問題上做了什麼主動的事,他們沒有主動的去思考怎樣接納這些外邦人,問題就出來了,隨著福音的廣傳,外邦信主的人越來越多,那麼這福音要不要傳給他們?問題就非常現實地放在早期的群體面前,所以在這個問題上有很多的思考,很多的辯論,我們看到這個的當中如果沒有神的帶領,沒有聖靈的帶領,沒有聖靈的光照,這個問題是找不出突破,走不出這個猶太傳統的圈子,感謝上帝,因著祂的帶領,看到彼得身上所發生的事,同時在保羅身上所發生的事,在耶路撒冷的教會當中,雖然有過很激烈的辯論,甚至是鬥爭和磨擦.在這使待行傳第11章接著往下看,就描述彼得給哥尼流施洗後,回到耶路撒冷,遭到了質問說:你和那些沒有行割禮的人一起吃飯嗎?你為什麼到他家裡去呢?彼得很有智慧的回答:這是來自於神的帶領,我們又怎樣去阻擋呢?

在早期教會他們在聖靈的帶領下,開始意識到原來這福音不單是給猶太人的,也給普世萬民的,給外邦人的,因此,彼得說我真的看出神是不偏待人,原來各國中那敬畏神的,行義的,都為主所悅納.這福音是開放給世上各國各族不同文化的人,不管是猶太人,不管是外邦人,你也不需要放棄自己的文化傳統,文化的身份,自己的生活方式,也不需要行割禮,不需要遵守食物聖俗的規定,這福音是給你的,當然你的生活要有更新,你在文化上要有改造,所以做為基督徒還可以是位中國人,希臘人,羅馬人,不需要變成猶太人,有些學者把這個過程總結了,我自己觀察,把這個總結歸納的最佩服,最好,最到位的就是耶魯大學神學院的拉明‧賽奈(Lamin Sanneh)教授,他說:這整個的過程可以說是外邦人的突破,經過這個突破,福音就終於從早期耶路撒冷教會,只是猶太的傳統的範疇之中,走進希臘,羅馬...的文化傳統,就真正的變成一個普世性,國際性的宣教運動.

這位學者接著說:其實在這個當中,突顯出一個新生的信仰,對普世各國民族文化的一個傳統基本態度,當中有兩點,第一,在上帝面前世上各地文化傳統都有其價值,都是值得尊重,福音的本質是超越世上的,階級上,國家上,種族上和文化上的界限,所以它是跨文化的,在上帝裡面沒有一個文化可以是神聖的地位,或是比其它的文化較優秀,也沒有絕對的控制權和主宰權；另外,上帝尊重普世文化,同時要進入這個文化當中,在這個文化當中道成肉身,在這個文化當中生根發芽成長,也要更新這個文化,改造這個文化.

從早期耶路撒冷教會的領袖彼得他們心中所發生的掙扎,猶太人的身份和外邦人信主問題的掙扎,他們意識到其實上帝的福音是給普世萬民的,我們知道這個過程完成最後應該是在使徒保羅開始,公元49年,在耶路撒冷會議以後,教會領袖就確立保羅所提倡的對外邦基督徒的基本立場,拒絕不再把猶太的傳統強加於外邦人身上,並拒絕耶路撒次教會作為普世教會的中心,那跨越文化界限,跨越國家民族傳播的趨勢,便一發不可收拾,很快便進入希臘,羅馬,埃及等地的文化當中,並且出現了不同地域特色教會文化中心,所以,我們通過今天的經文,回到早期教會的掙扎,便能再次提醒自己,不管你是那國家人,不管你的文化傳統的差異,這福音是尊重你自己文化的身份,尊重你自己國家的身份.教會在這二千年的歷史當中,也可以找到許多的例子,就是基督教的宣教士,或是教會的宣教士,或是教會的領袖,有些地方,有些時候,他從自己母國的文化,到另一個文化當中去宣教的時候,他帶著一份很強烈文化的優越感,他自己充當了一個把自己本國家文化的輸入者,或是成為把當地的文化纖滅的一份子,那又怎樣去解釋?我們絕對不會迴避這事實,在這二千年的教會歷史當中,教會在這信仰,福音和文化的問題上,其實犯過很多的錯誤,走過許多的偏路.

我想特別在中世紀的歐洲羅馬帝國裡,確立了一個至高無尚的地位,這維持長達一千年之久,與西方文化和文明緊密的結合為一體,教會在這時候所犯的錯誤和偏路是得明顯的,甚至是很嚴重的,那時候分不清福音的核心信息,和本族文化特有的東西,有時候會把文化傳統的一部份東西,變成基督教信仰很重要的一部份,所以每當傳福音的時候,也把這傳統連福音核心信息一起帶出去,而且感覺自己這個本國本族的文化是蒙神悅納的,這是高等優越的,我們現在還常常聽到基督教文化和基督教文明,我不能對這詞很籠統說它對與錯,可是我們要小心的使用,使用不對便產生文化的優越感,無論如何,儘管有這些錯誤,福音最終被傳開了,福音的核心信息還是傳開了.

在二千多年的歷史教會上,特別是在十八-十九世紀的時候,是福音運動大爆炸的時間,是最興旺的時期,到了二十世紀初時,信仰已從地中海走出至歐美這核心地帶,遍及全世界七大洲,特別是透過第二次世界大戰之後,老生常談,到現在我們看到有一個現象,這信仰在南半球在亞非拉地區過去是宣教地域,非常興旺,傳播快速,教會成長非常迅速,信主人數遠超原來基督教最強的地區和國家,也是我們傳統說的西方國家,則北平球,因此,今天引起許多專家對這現象的討論,有人說這是基督教的全球化,世界化,有人說基督教的中心從北半球向南半球的轉移,等不同的說法,也有各種的統計數字指出,大概在二十世紀初的時間有百份之七十的基督徒都在這北平球生活著,特別是在歐洲,到了二十世紀未的時候,歐洲的基督徒在全世界基督徒人口中的比率不到百份之廿八,但是拉丁美洲和非洲的基督徒佔全世界的百份之四十三,而且還在增長,不管統計有多少,數字有多少,有個事實很清楚,就是大部份的基督徒都生活在這個傳統基督教世界文化之外,所為西方世界之外,在內亞非拉這地的基督徒是最多的,這個事實正在改變我們的信仰傳統面貌,很多大學尤其是西方大學於宗教性質方面看,在定位的時候,還把基督教定位為西方宗教,我個人認為這見解已脫離了現實,會不會在未來的日子上將會變成歷史,基督教的面貌越來越不西方化,所以今天世界上各種主要信仰當中我們可以非常肯定地說,基督教的信仰它的國際性最強,它的跨越民族界限最多的,這個事實本身也告訴我們其本質是什麼,它對待普世文化的基本態度也是一個印證,今天我們為什麼有必要回到早期教會,乃至於廿十世紀的教會今天的現象來思考這問題,來提醒我們教會普世的性質,教會的信息,跨文化的性質,因為經過二百年福音在中華大地的傳播,很遺撼中國人對這信仰的看法,仍然有著十分嚴重的洋教意識感覺,西方的形象.廿一世紀,中國教會面對一個重要的挑戰,依然還是怎樣讓我們的信仰普世的特質突顯出來,在中國文化社會當中,讓它呈現出一個真正中國的面貌,廿一世紀,其實今天我們都很清楚,如果說中國教會面對最重要的挑戰之一,其實已經不是無神論,不是科學主義,不是理性主義,而是來自於傳統中國文化,傳統中國宗教種種挑戰和質疑.

如果說在1980年代的時候,中國大陸剛剛改革開放的時候,整個時期,中國人的心態對世界比較開放,對福音比較開放,很急於向西方其他國家文化上去學習,可是到了今天,好像馬克斯主義沒有多少人信了,但是傳統的儒,佛,道全面的回巢,從政府到民間的推動,從高層的學術討論,以至小學生的背誦背經 -- 背論語,背孔子,背大學,背中庸,今天前往內地的小學觀察活動廣告,往往都是招收國學班,所以傳統文化興起得非常旺盛的,傳統文化的表達人物都帶著一個動機,他們動機是什麼?他們為什麼復興這個傳統文化?就是要針對改革開放以後,在中國傳播十分迅速的洋教,因為在他們的眼中,基督教是西方的文化洋勢力,所以越多的中國人接受福音,基督教,對未來中國文化是一個大災難,所以他們有著一個使命感或危機感,要盡快的復興的傳統文化來對抗,這就告訴我們,雖然經過二百年的努力傳福音,今天雖然教會與福音都非常興旺,仍然我們在中國人心目中,我們的形象和二百年前相差無幾,至少在相當的部份中國人心目中,所以我們在如何實踐讓福音真正在中國社會文化當中道成肉身,讓它這個普世性能夠突顯出來,我們還有漫長的路,這就是為什麼我們要藉早期教會的歷史和經驗,並通過聖經的教導來不斷提醒自己,這福音從耶穌宣講帶給世人,祂就是給普世萬民的,不論你是平信徒,或是教會領袖,我們必須要在事奉和見證當中,思想我們所言所行有什麼可以幫助改變國內同胞對我們信仰根深蒂固的成見,如果你是一位對信仰有興趣的朋友,也許你今天來這裡抱著一種好奇的心態,或許是透過基督教認識西方文化,我希望你有天能夠走出對這信仰有更深的認識,我希望你有天信主的時候,你還是位中國人,還可以愛自己的文化,愛國家,因為上帝就是愛,以色列人,猶太人也愛中華民族,雖然這個當中對每個民族都有祂的管教,有祂的審判在其中,我們今天通過回顧使徒行傳這故事,讓我們再次記著這信仰什麼對待世上的文化基本態度,讓我們更好的在這個文化多元,社會價值多元的世代,我們即把握真理,同時能夠做神很有智慧和能力的見證人,感謝神,有機會與大家分享,謝謝各位.



\newpage

\section{一片燒魚和一塊蜜房}
\label{sec:67}
\textbf{戴國璋牧師 - 一片燒魚和一塊蜜房}
\newline
\newline
\hyperref[sec:66]{< < < PREV SERMON < < <}
~
\hyperref[sec:toc]{[返目錄]}
~
\hyperref[sec:68]{> > > NEXT SERMON > > >}


日期:2010年10月10日

主題:一片燒魚和一塊蜜房

証道:戴國璋牧師

我相信我們大家都知道,耶穌復活是耶穌傳奇故事當中最高潮的一個部份；我們也知道耶穌復活成為一切活著及睡著之人的主,這個也是基督教信仰裡頭最核心的部份.

因為耶穌已經從死裡復活,所以我們知道也相信耶穌已經成為基督,正如彼得提醒我們的,就是因為這位的主從死亡當中復活,所以上帝已經將祂高舉,將祂放在宇宙萬有最高的位置上面.祂已經成為基督,在沒有任何事情的發生,並不在祂的管理與掌握中間.我們又知道,因為耶穌從死裡面復活,所以我們認識耶穌基督是神的兒子,只有那位擁有不朽生命的人,祂才可以衝破死亡,向我們顯明生命勝過死亡這個重要的真理.由此,我們知道耶穌不是一個普通的凡人,祂擁有著上帝永恆的生命,祂用衝破死亡向我們見證了這個真理.

我們又知道,因為耶穌基督復活的緣故,耶穌已經成為救主,祂是世上人類最終極的希望,因為祂可以將祂這個不朽的生命與所有仰望祂,信賴祂的人來分享.因此所有的福音書都不像一般英雄傳奇故事那樣,以一位英雄的離世作為結束,它乃以耶穌的復活作為最後的終局.路加福音的記述也是一樣的,它將耶穌的復活作為它記述耶穌故事的最後終局.

剛才我們唸到有關路加的記述,是講述當耶穌復活之後,初次向門徒顯現.耶穌在禮拜天的早上從死裡復活,所以那次的記述是在晚上的時間向門徒顯現的故事.我們相信這個講述也是約翰福音20章19-23節所講述的那一次的顯現,故此我們可以將它們彼此來比較,對當晚所發生的時間,所發生的事,有多方面的了解.根據福音書的資料告訴我們,在那晚十個門徒在耶路撒冷某一地方裡面聚集,為何是十個門徒呢?因為當時猶大已經離隊而去了,多馬不曉得甚麼原因在那天晚上並沒有出席那個約會,所以根據約翰告訴我們當耶穌顯現的時候,多馬沒有看見,多馬心中存疑,表示說我必須見到這位主,摸祂的釘痕手,摸祂肋旁的傷痕,我才願意相信,結果八天之後,耶穌再向他們顯現,向多馬證實祂自己復活的真實.

那時候不單單有十個門徒在那裡,因為經文告訴我們,原來還有兩個早上往以馬忤斯路上的兩個門徒,他們在路途上遇見復活的主,當黃昏時他們與主在客店裡一同坐席擘餅的時候,他們認出這位是復活的主,於是他們急忙地趕回去耶路撒冷,向門徒們報信.當他們在談論著這幾天所發生許多的事,有許多使他們疑惑的事情,因為他們的主被殘酷的殺害,軀體被埋葬在墳墓裡面,這些的事,是他們所知道的,甚至他們當中有些人親眼看見這些事情的發生,他們因為這個緣故,曾經感到非常的沮喪與哀傷,因為畢竟這位是他們所愛的主,是三年多以來與他們一同生活的主,當然除了這個沮喪與哀傷之後,他們也感到有絕望,因為他們已經將他們許許多多的希望押注在祂的身上,他們相信祂是基督,祂是要帶來以色列國的復興,所以他們已經將一切都押注在祂身上,甚至他們本來已經在那裡籌算著,當有那麼的一天,以色列國真的復興的時候,我可以從我這位主的身上得到甚麼呢?但這一切都已經成為絕境,因為這位他們視之為彌賽亞的主,已經離開這個世界,不過在早上的時候,有些婦女曾跑過來報信,他們宣聲她們看見這位復活的主,又有甚麼可能呢?人已經死了,怎可能再活起來呢?我們知道,從來婦女的見證都不能成為見證的,法庭上也不會被接納的,相信這些婦女一定是眼睛昏花,才作出一些幻想吧!後來當中的彼得與約翰,也因為這個緣故曾經跑到埋葬耶穌的墳墓裡去查察,他們看見墓門打開,墓穴是空置的,屍體不見了,但是他們不敢肯定,究竟這個原因是如同這些婦女所說的,夫子已經從死裡面復活,還是因著某一些的原因屍體被人用了某些方式挪動了呢?他們沒有辦法肯定.在這個時候,忽然有兩個耶穌的其他的跟隨者,他們早上去一個離耶路撒冷不遠的村莊,叫做以馬忤斯,就在路上,他們竟然說他們遇見了復活的主,又有甚麼可能呢?相信他們只是眼睛昏花而已!又或許他們思憶成狂,所以幻象猶然而生?抑或主真的是復活了呢?

我們發現當時候門徒的心情,是非常複雜的,他們的心情既有疑惑,又有相信,疑惑與相信交織著；既有沮喪又有興奮,沮喪與興奮彼此互換著,他們真的不知道應該抓住那一個信息,證據,就在這個時候,耶穌突然之間出現在他們當中,給了他們一些安撫,給他們心靈裡頭最重要的答案.

上個月(8月)23號,我相信對於許多香港人來說,是一個很難忘的日子,因為我們真的沒法估計到在一個普普通通的旅行事件,竟然會發展成為有許多人傷亡,有許多親人離世,有許多家庭破落的一個事件,我所指的就是在馬尼拉黎剎記念公園裡面的挾持事件,後來變成殺人事件.我還記得早上的時候,我當時在看NowTV的新聞,NowTV報告說,在馬尼拉的地方有一個旅行團被挾持,我很奇怪,為何會有挾持事件的發生呢?一路觀察動態的發展時,心裡頭總覺得應該沒甚麼事的,因為似乎是那麼的平靜,那殺手是那麼的友善,沒多久又聽見他釋放了兩位老人家,沒有久又釋放了一位太太與她的兒女和另一位小朋友離開了,情況就如同普普通通的一個挾持事件,應該是可以很容易來解決的,但是沒想到到了下午的時候,整件事會峰迴路轉,到了傍晚七點多的時候,我們聽見開槍的事情,也在電視上看到一連串的看起來啼笑皆非的拯救行動,最後導致的是一連串悲慘的事情,以及許多香港人傷心的回憶!

我相信在這個過程裡面,每一個注意的香港人,心情都是很激動的,可能在心靈裡面都泛起許許多多不同的問題,以及許多的不明白.例如有一次,特首在某一個場合裡頭,他表達了他對這件事的一個問題,他說,他不明白,不明白為何一個人可以變得那麼邪惡,一個人怎麼可以將自己的不滿,不高興發洩在別人的身上,然後令許多人受傷,令許多家庭破碎,他說他不懂,不明白.對於我來說,我也有一些不明白,我不明白為何幸福是那麼的脆弱的呢?其實我們經常都會去旅行,坐旅行團為我們安排的旅遊巴士都不曉得有多少次,都是這樣子的上上下下... 那一天,我在家中,那時候我正值休假,因為我們慣常在七月或八月底休假,而且許多時候,我們也會與家人一同外出旅行,好幾年都如是,不過今年就比較特殊,今年因為有一些特別的事要去做,我兩個女兒也有一些特別的活動,所以我們沒有選擇外遊,那段期間我就在家休息,以及整理一下房子,因為許久沒去整理房子了,留在家中,沒想到就看見這件事!

我們經常都高高興興地去旅行,我們隨時都可以回憶在最後一天的行程會是如何,早上吃完早飯之後,行李放上車,接著就會隨團出去遊玩最後一個景點,走完後,可能就會在市區某一些百貨公司裡逛一下,接著去吃飯,又或許連吃飯也沒有就直接去機場,上飛機,然後在飛機上面吃午餐... 我們經常如此,就在這架旅遊巴士上面上上下下,通常在那些日子裡,心情是有點複雜的,一來玩了幾天,要離去真有點不捨,而另一方面又覺得,要走了,要回香港,回到家就可以舒舒服服的啦... 我們都有這兩種心境──又不捨,又想快點回家... 但是又有誰會猜得到,對於這些家庭來說,當他們抱著這種心情走上旅遊巴士的時候,他們將要經歷十多個小時的難受經歷呢?對我們來說(或對我來說),我的問題是,為何幸福是那麼的脆弱的呢?我不期然想起以前唸過白居易的一首詩:「大都好物不堅牢,彩雲易散琉璃脆」,好的東西,為何那麼的不堅牢呢?彩雲真的很漂亮,但沒多久你會發現它很容易就消散了,琉璃也漂亮,但也是很脆弱,容易破碎.

人世間好的事物,為何是那麼的脆弱呢?當然也有一些與我一樣有信仰的人,心裡面也是很疑惑,我們知道我們一心倚靠上帝,我們也指望上帝在我們生命當中賜下平安與祝福,我們也相信上帝真的也能夠如此做.但在這次的行動裡面,我們很疑惑,因為我們看見一些倚靠上帝的人,他們蒙上帝的拯救,但也有一些倚靠上帝的人,上帝似乎沒有做任何事情,為何呢?譬如,我們都知道,最早被釋放的一對李姓老夫婦,他們是基督徒,據這位李老太回憶當時的情況,她說,當搶手上車挾持他們的時候,她很害怕得很厲害,不斷地禱告說:神啊!求你救救我們,救救我們... 在這個時候,她突然間覺得肚子痛,痛得很厲害,很難受,於是她只好跟導遊說,我很難受,快受不了了,肚子很痛,我要上廁所!這導遊幫她翻譯給這位搶手聽,這搶手竟然說,既然如此,你帶她離開吧!他就允許當地的導遊帶這位李老太太離開,就是如此,這位老人家就脫險了.當她脫險之後,她遇見紅十字會的工作人員,她跟他們說,我先生他有糖尿病,藥也已經吃完了,再不吃就會有危險... 很奇怪這些紅十字會的工作人員,將這個消息轉達給這位搶手知道,搶手就立刻釋放了她的先生.他們都直說這是一個奇妙的恩典!很奇怪,很奇怪,原來肚子痛都是上帝的祝福,她是這麼的說.我們很為他們感恩!一對倚靠上帝的家庭,老人家能提早被釋放.我們也知道,當中被釋放的傅太太,帶著兩個兒女離開,後來她也很機智地將另一個汪姓的小弟弟帶在她身邊一起離開的那位,她的家庭也是一個很虔誠的基督徒,我們也為他們感恩.這兩個家庭蒙上帝的恩典,我們也知道,姓梁的一家,雖然他們是天主教徒,但對於我們中華基督教徒來說,天主教會他們對上帝的信心,與我們都是一樣的人,但他們全家裡頭有三個人離開世界,現在還有一個躺在醫院裡,我們默默的為他們禱告,希望他們能夠康復.

為何上帝對一些祂的百姓那麼寬大,充滿各種的恩,對某一些的人,祂又不是這樣子去處理呢?我們心裡頭有許多不明白,當然也會有一些不相信的朋友,他們會對基督教的信仰有一些的攻擊,他們會再一次地說:你們的上帝,你們的上帝是一個無能的上帝,如果祂是一位有能力的上帝的話,為何祂看到這些事情的發生,卻甚麼也不做呢?當然一些的誣蔑,控告都發生在一些大事情之後,而基本上我們都會從一些未信的人的口中聽得到,我們心裡頭很多時候都很不服氣的.因為我覺得如果你從來都不相信有上帝,上帝只是子虛烏有,那為何你還要去懷疑一個子虛烏有的上帝的能力,工作呢?因為祂既然不存在,我們還有甚麼可以譏笑祂的地方呢?但人很多時候就是如此,一方面否定祂的存在,另一方面又責怪祂的無能!

許許多多的問題在我們心靈裡面交織,就如同當日的門徒一樣,疑惑與相信交織；沮喪與興奮交換,我們應該怎麼樣找到一個定位呢?當日耶穌出現在門徒中間,耶穌的顯現平伏了他們這種複雜的心情.我在想,如果今日我們一同來唸復活這故事的時候,會不會就在這故事當中找到一些的答案呢?我相信,或許我們可以的,我們會把那個焦點放在路加福音這個故事裡頭的一個地方,這個地方就是在所有復活故事當中,路加有一個很獨特的記載,甚至我們都會奇怪,一些不起眼的事,又有甚麼理由值得花筆墨來把它寫下來呢?但路加卻為我們記載了,我們就以這個作為我們今天的焦點.

這個就是,耶穌竟然向門徒取食物.為何耶穌要向他們取食物呢?是不是因為耶穌餓了?但是當我們仔細地去想的時候,我們都會記得,耶穌餓了這事,是一點也不希奇的,因為上一頓飯已是星期四晚上的事情──星期四晚上與門徒吃完這個「最後晚餐」之後,在整個星期五凌晨耶穌受審,在星期五早上被交付釘十字架,在星期五傍晚祂被釘十字架,然後埋葬在墳墓裡,星期六的一整天裡頭,祂都在墳墓當中,而星期天出來,到現在,祂都還沒有吃過任何東西.啊-當然啦,會不會祂在以馬忤斯的路上,那時候不是說和那兩個門徒一同坐席嗎?不過故事告訴我們,當擘餅的時候(開始吃飯的時候),他們認出耶穌來,但耶穌突然之間就不見了!如果真的如此的話,其實耶穌基督一路以來都沒有吃過東西,超過24個小時當中,耶穌沒有吃東西,而現在要求食物,是一件不希奇的事.於是,他們就給了耶穌一片燒魚以及一塊的蜜房,耶穌就吃了.

不過,當我們細心地去想的時候,我們總覺得有點不對勁,為什麼?因為這裡明明地告訴我們,耶穌其實是以復活之身來向門徒顯現,既然是以復活之身向門徒顯現的時候,為何一個復活了的身體,難道仍然繼續需要一些食物來維持祂的生命嗎?難道還需要物質的食物來給予能量的供應嗎?所以當我們細心地去想,耶穌要食物的原因,並不是因為祂餓了,事實上,如果真是如此,這20多個小時或更長的時間的饑餓,一片燒魚以及一塊的蜜房又怎麼能夠滿足祂的饑餓呢?耶穌這個行動,是一個故意的行動,而且明顯地是一個象徵性的舉動.讓我們看看這一片燒魚以及一塊的蜜房,可以提醒我們甚麼事?

首先我們說,這一片燒魚以及一塊的蜜房其實是代表著一種事物,它代表我們世俗生活裡的平凡一面,而且是最最最平凡的一面.在人的世俗生活當中,吃東西是我們每一天都要做的事,對不?有些人一天吃三餐,有些人一天可能吃超過三餐.「吃」是我們生命裡頭每一天都必須做的事,因為我們不吃,生命就不能得以維持,我們的生命就缺乏能量讓我們活下去.所以「吃」是生活裡頭一件很平凡的事.而在食物的裡頭,燒魚及蜜房是食物當中最最最平凡的一個食物.

我們大家都知道,在古代的世界,吃肉其實是一件非常困難的事,也是一件奢侈品,普通的一個家庭,一年到頭來是很難得可以吃肉的,只有在一些大節目裡頭,才有些機會能夠吃點肉,否則許多時候,都只不過吃一些菜蔬而已,只有一些有錢的人,他們才可以有多些機會吃肉.所以為了補充所缺乏的肉類,蛋白質,我們發現魚是其中一個,既廉價又容易得到的肉類蛋白質的來源,因為你只要到海邊去打撈一下就可以了.因此,我們有時候聽到鄉下有些人,當他們吃飯的時候,因為沒有肉,有甚麼可以加菜呢?就是去河邊釣一條魚回來,那天就有肉可以吃了.所以吃魚是一個方便又便宜的肉食蛋白質的來源.

在耶路撒冷的東北角,有一個城門叫「魚門」(fish gate),為何叫「魚門」呢?因為那裡有一個魚市場的,每天一大早的時候,就有魚販挑一些魚,從魚門進入耶路撒冷,在魚市場內擺賣,城裡的人就可以在魚市場當中買到他們所需要的各種魚類,可能門徒當日也在這個魚市場買了這些魚,然後製作成燒魚.當然,每一樣的食物,多簡單的食物也好,也可以變得很複雜.你可以用一個很複雜的烹調方法,加上一些很複雜的佐料,配料,煮一餐很複雜的一餐.不過魚有一個很簡單的方法來處理,這個最簡單的方法就是「燒」.任何一個地方,只要你能夠起一個溝火,譬如找一些柴枝,樹枝,起一個溝火,你就可燒一些魚來吃.猶太人習慣用一支尖的木棒,將魚串成一串,然後放在溝火上面烤,門徒所說的烤魚,有可能就是這一種煮食的方法.或許當他起了一個炭火,將魚平放在炭火之上燒,聖經裡頭所說的燒魚,應該就是指著這些炭火了.

我們有理由相信這一片的燒魚,是當日門徒們的晚餐來的.其實當天門徒的心情是相當複雜,既有悲傷亦有疑惑,又恐懼那些猶太人會來傷害他們,所以他們把門窗都關起來了,他們那有心情在當天晚上大吃大喝呢?所以他們的食物是越簡單越好.這一片燒魚可能是他們當晚所吃剩的,當耶穌要的時候,他們就隨手把這片燒魚遞給耶穌,祂就吃了.

蜜房也是如此,也是一個簡單的食物.很多時候我們以為我們所吃的蜂蜜,是一些新鮮的花蜜,原來事實並非如此.原來所有的蜂蜜,都是經過蜜蜂作了一些基本的處理,然後才有的一個成品.當蜜蜂採花蜜的時候,將花蜜帶回去蜂巢裡頭,好餵飼那些小蜜蜂(幼蟲),這一些小蜜蜂除了吃工蜂帶回來的花蜜外,牠們會將這些的蜜在牠們的體內作一些初步的處理,例如,把牠們嘴裡頭的分泌物混和在這些蜂蜜裡面,再將蜂蜜給吐出來,把它黏在所住的蜂巢內的巢洞壁裡,用以保暖.而我們所吃的蜂蜜就是黏在巢洞壁內的,被吐出來結成晶體的東西.人類就很本事,能夠把這些蜂蜜拿出來,溶在水裡,而成為今天的蜂蜜了.

一個蜜房的營養價值並非很高,不過很多時候卻能提供我們對糖的需求,是一種瞬間補充能量而很有效的東西.如同我們中間有一些喜愛遠足的人,總要帶一些巧克力,因為當肚子餓沒力量的時候,吃一些巧克力就可以瞬間補充能量.糖分是可以幫助我們的.我心裡頭在想,會不會這些蜜房是當日那些門徒們準備隨時逃亡用的呢?因為你從不知道環境的變化會如何?何能隨時都要離開,有一些簡單的食物,如同蜜房一般能夠提供能量是最好的了.

耶穌伸手接過蜜房,隨後就吃了,這個行動有甚麼象徵呢?耶穌很明顯是想藉著這個行動向門徒證明,祂是百分之百的一個有血有肉的耶穌,祂是百分之百真真實實的耶穌,祂百分之百從前與他們一齊生活,都不曉得一共吃了多少次,喝過多少次的那一位耶穌.我們知道,路加如此寫這個故事,有這種的記述,是有特別原因的.因為當教會開始傳播的時候,隨著教會的有一個叫「洛斯底主義」的異端,這異端滲入教會裡頭,利用教會許多的教義來宣傳他們許多的說法,其中的一個說法就是:物質是邪惡的,身體因為是物質,所以也是邪惡的,如果耶穌是神的兒子的話,祂不可能有一個物質的身體,因為身體是邪惡的.所以一直以來活著的耶穌,其實只是一個幻象來的,這就是「洛斯底主義」裡頭的「幻影說」(Docetism).當時候的教會與這個幻影說來爭辯,門徒們要清楚地見證當日與他們一齊生活的耶穌,以及後來從死裡頭復活的耶穌,都是有血有肉的耶穌,路加很明顯的是要強調這一點.所以你會發現,當耶穌向他們顯現的時候,對他們說我不是鬼魂:「摸我看看,魂無骨無肉,你們看我是有的」(路24:39)耶穌繼而取了食物,藉著食物(不是祂需要食物)的行動來證明,祂真是那位有血有肉的耶穌.當然今天我們不會再說這個有關「洛斯底主義」的事,當我們研究聖經背景的時候,已經成為一種學習的課題.

但今天這個復活的真理,耶穌是身體復活的真理,是我們每一個都必須要相信,這是我們信仰裡頭的內容.大家記得當我們背誦「使徒信經」的時候,最後的一句話:「我信身體復活,我信永生」.基督教相信的不單單是靈魂的不朽,基督教相信的更加是身體的復活,這個故事就是告訴我們,當門徒切切盼盼這位耶穌,以為祂已經離他們而去(死了),耶穌卻向他們顯現,然後上帝藉著這個行動把絕對原本的耶穌還給他們,祂不是鬼魂,耶穌很清楚的說:「魂無骨無肉,你們看我是有的」(路24:39)我並不鬼魂回歸.

當時候在猶太人的民間信仰裡,也有所謂「守衛天使」的說法,他們認為每一個人,每一位神的百姓,都會有一個「守衛天使」代表他常在天上見天父的面,他們有一個很特別的看法,這個「守衛天使」的樣貌與那人的樣貌是一樣樣的,像我戴國璋就有一個「守衛天使」,他的樣子跟我是一樣的,所以他在天父前面常見天父的面.那這個會不會就是耶穌的「守衛天使」呢?如同我們在使徒行傳唸過一個故事,當彼得被收在監牢裡頭的時候,希律王想傷害他,教會的人以為彼得會被害,於是切切為他禱告,後來彼得離開監牢,去敲門徒的門時,當時有一個使女來開門,聽見是彼得就很害怕說,彼得站在門外,他們卻說,怎麼可能會是彼得,「必是他的天使」(使12:15).就是有這麼一個的背景.他們以為你所見到,所聽到的,是像彼得的天使.

會不會上帝差了耶穌的「守衛天使」來裝扮成耶穌,為的是安慰門徒呢?不是!那是耶穌本人.今天我們知道那些複製技術是很厲害的,我們已經有複製羊,犬... 等等.再過沒多久,就會有複製人.如果只要拿到人的DNA資料,將來可能就可複製一個人出來,而上帝就更加有可能如此作,祂可否複製一個耶穌,再把耶穌還給你呢?不過那只不過是「複製」的耶穌,不是「原本」的那一個.聖經藉著這個行動要向門徒顯示,上帝把完完全全,百分之百原來的耶穌展現在他們眼前.所以門徒喜得不敢信.

今天這個也是我們信仰的真理,我相信我們每一個人都有機會離開這個世界,不過不打緊,身體雖然毀壞,不過有一天我們在主裡面,會重新得回這個身體,上帝會把原來的那個「我」給我們.會不會在馬尼拉這件事情中,這個復活是其中一個答案?上帝沒有解釋為何祂這樣子做,為何沒這樣子做?上帝只是告訴我們,任何的邪惡對我們的毀滅,縱然毀滅我們的身體,但最終上帝能夠把那個百分之一百的親人交回給我們.我們相信,昨日出殯的梁先生及他兩個女兒,與他太太暫時的離別,但他們在基督裡有真實信仰的話,上帝會把原來的先生,女兒交回給她,是原原本本的那一個.所有對上帝有復活盼望的人,我們都相信我們真的能夠看見.這是我們看見的第一方面,讓我們接下來看看第二個方面.

燒魚與蜜房代表著甚麼呢?燒魚與蜜房代表著世事的回味,提醒我們在生命當中許許多多寶貴的回憶,這些回憶都是真實又寶貴的.又代表著我們對未來的憧憬.

耶穌並非第一次吃燒魚.這燒魚讓我們想起許多許多次耶穌在提比哩亞海邊與門徒在一齊的生活片段.耶穌的門徒中有幾個都是加利利的漁夫,他們這群人不曉得有多少次在加利利海那裡泛舟湖上,經過平靜的湖面,也經過波浪翻騰的湖面,耶穌也不曉得有多少次看到門徒在那裡捕魚,也不曉得有多次停船在提比哩亞海邊,門徒撿了些柴枝起火,把捉回來的魚燒來吃,然後在那閒話家常或講天國的道理,許許多多這類的經歷.

各位你們記得約翰福音21章,約翰特別提到耶穌復活之後向他們顯現,是他們不能忘懷的.在提比哩亞海邊門徒照往常一樣去打魚,但他們有點心灰意冷,因為他們還未能從他們的灰心當中被振奮起來,他們一夜都沒打著甚麼,然後他們看見在遠遠的岸上有人呼喊他們,問他們一個問題:「你們有喫的沒有.」(約21:5)這個問題,真的很奇怪!(在我們剛才唸的路加福音裡,也是如此,聽見耶穌這個問題:「你們這裡有甚麼喫的沒有.」(路24:41))門徒就回答耶穌,我們一夜勞力,並沒有打著甚麼.

親愛的各位,這個故事很明顯想讓門徒去回想他們以前很重要的經歷,經驗,這是記載在路加福音4章那一次打魚的故事裡.同樣,他們在那裡打魚,打不著甚麼魚,然後耶穌跟他們說:「把船開到水深之處,下網打魚」(路5:4)門徒聽耶穌的話,姑且試一試,就就圈住許多魚.耶穌在約翰福音21章那裡說:「你們把網撒在船的右邊」(21:6),左邊,右邊又有甚麼分別?左邊不會有,只有右邊有呢?我們也不明白,但聽你所說的,結果也打了許許多多的魚.就是這種雷同的經驗,挑動約翰,約翰說:「是主」(21:7),彼得一聽見,衣服也不穿,就跳在海裡,游上岸,要見他的主耶穌,那真的是主耶穌!當門徒上了岸,發現主主耶穌已經在岸上為他們預備了燒魚及餅為早飯,後來就聽見耶穌問彼得:「你愛我比這些更深麼」這些故事.對基督徒來說,這故事我們應該不會不記得,或是從未唸過這個故事.

這些故事給我們看見,耶穌很重視回憶,祂藉著這個餅和魚讓門徒回憶以前許許多多舊有的片段,在一起的日子是多麼的甜蜜,溫馨,利用這些回憶重新振作:當日我如何召喚你們,叫你們放下這些魚船,魚簍,成為一個得人漁夫,今天,你們要去餵養我的群羊.

甜蜜的回憶很多時候是我們生命裡頭的動力,所以很多人有時候都很喜愛把舊有的照片拿出來,從照片當中重尋舊有的回憶.而事實上學者告訴我們,記憶是人類生命裡頭最寶貴,最寶貴的恩賜,我相信這是神給我們的.當然,動物可能也會有回憶的,但不會有人類那麼清楚的回憶.在我們生命裡頭許許多多的片段,看起來或許不是甚麼重要的片段,但也可以從我們的回憶當中尋索而得.這些回憶讓我們找著生命的寶貴,那一種的甜蜜.

吃魚能夠帶給他們回憶,吃蜜房就更加... 因為蜜房是甜的,復活的耶穌,祂的舌頭是可以嘗到甜味,多麼的溫馨,奇妙.我們不單單可以從回憶裡頭找到我們珍惜,寶貴的東西,原來回憶也是我們計劃未來,憧憬未來的根本的所在.對未來還未發生的事,我們可以想像,是因為我們從回憶裡頭,可以有許多的元素以幫助我們構建未來.

在25那晚,電視台播放了一套記錄片,這記錄片就是研究有關人類的回憶,這記錄片當中有一個研究是蠻特別的.它找了一些志願者來做實驗,而為了觀察他們腦波的變化,志願者頭上都放置了一些電極,然後就把他們送進去一個類似磁力共振的儀器當中,觀察員就在旁邊的電腦屏幕上觀察他們的變化,他們的實驗就是回憶與想像的關係.觀察員要求實驗者先去回憶一下從前的生活片段,接著在電腦屏幕上觀察他們腦波的變化.這部份之後,就請他們計劃一下未來,想一下將來未發生的事,想想將來打算如何,之後也在電腦屏幕上觀察他們腦波的變化,最後將這些圖片拿來比較的時候,發現一個很奇怪的事,他們覺得這是一個很震驚的事實.原來當我們回憶一些事,以及當我們去想像一些的事... 回憶是以過的事實,想像或計劃未來是屬於還沒發生的事,當我們做這些行動的時候,原來都在我們腦部的同一個區域裡頭活動著(這是可以從那些剖析圖中看到的),那就得出一個結論,我們腦部同一個區域是管理著我們的記憶以及計劃,這是不希奇的,因為當我們能夠憧憬未來的時候,一定需要有一些東西好讓我們能夠構築將來,用甚麼東西去構築呀?就是利用我們在回憶裡得到許許多多的片段.

這件事情讓我們看見,回憶寶貴,憧憬也寶貴,我們憧憬有一天與上帝一起,在神國裡頭喝新的,在神國的裡頭,再沒有眼淚,再沒有死亡,我們可以與那些離開我們,暫時離別的家人再次在主裡見面,到那一天,是多麼的興奮與快樂呢!

我相信路加福音的這個故事,很可能有這方面的提醒,再次告訴我們,上帝會再次將我們的親人交回給我們,以及我們在基督裡有復活的寶貴應許,因著這個原因,所以今天也要好好珍惜我們每一個生活裡的片段,因為我們生活上每一個經歷都是寶貴的,且帶我們進入永恆裡面.

在25日的早晨,我留下一個剪報的內容,蠻有意思的.有一個寫諷刺漫畫家尊子先生,他畫了一個四格漫畫,這漫畫以當日馬尼拉事件為題,蠻能顯示出我們香港人的特點.他第一格漫畫說有幾個典型的香港人去搶一個「\$」符號的東西,香港人是很努力地去賺錢,有的說,這是我的；有的又說,是我的,我首先發現的... 大家都在爭奪....第二格漫畫有一個人說,不要搶,安靜,有新聞特別報導.於是在第三格漫畫裡,大家聚集地在看到底是甚麼新聞,原來就是馬尼拉脅持事件,新聞說,有關脅持事件流血告終,死者有.... .第四格漫畫就看見這群人四散而開,各人從各人的口袋中拿出手機打電話,有人說,媽,我今晚回來跟你吃飯,老婆你沒事吧,沒事就好,兒子呀,我今晚陪你到公園玩....他藉著這個事件來提醒我們,平常我們覺得一些很平凡的東西,不值得注意的東西,但其實這些都是幸福,是我們各人都很需要去珍惜的.

一位中大神經科學的教授接受訪問時,他對這件事的評語,其中一句話說:或許這是我們香港人重新去思考,究竟我們生命裡頭那一些東西是最重要?是我們的事業,生意,學業,還是家庭?求主讓我們看見,原來我們今生與親人間的寶貴經歷,都是上帝所珍惜,是能夠帶進去未來的.

最後我們看燒魚.我們知道,魚在往後的基督教裡頭成為一個很重要的符號,因為他代表了基督徒,我們知道後來的教會,由於環境的因素,受了許多迫害,而許多基督徒都無法用公開的身分來生活,所以當他們需要暗通身分的時候,他們就會畫一條魚,如果對方是基督徒的話,他就會明白,也畫上一條魚,大家就可以走在一起分享神的真理.當時候的教會沒有固定的活動地方,都經常在一些秘密的地方,例如,在墳墓裡敬拜.而每當需要聚會的時候,如何呼喚信徒呢?原來就是在市集的某一地方掛上一條魚,就表示教會有活動要聚集.魚成為基督徒很重要的一個標誌.因為在希臘文「魚」這個字,如果把字給拆開,每一個字也成為另一個新的字的話,就會成為一句話:「耶穌基督,神的兒子,救主」.耶穌真的是救主.

這條魚成為基督徒的標誌符號的原因,與我們今天路加福音24章耶穌吃的這條燒魚是毫無關聯.更多的是與耶穌所呼召的門徒,他們大部份都是漁夫,後來耶穌應許他們要他們成為「得人的漁夫」的典故有關.不過這條魚,象徵耶穌是救主的故事,是再次提醒我們所有問題的答案,都在耶穌基督祂帶給我們不朽生命的分享那裡.

今天我們人世間有許多的悲傷,有許多不幸,這乃提醒我們人世間許多的幸福都是很脆弱的,所以我們要靠上帝來守著,你和我都要和上帝一同來守著我們的幸福,幸福並不是必然的.不過即使這個邪惡在我們人生中將我們的幸福粉碎,上帝已經有辦法,在將來把這些幸福再一次地賜給我們,因為祂應許我們在基督裡有永遠復活的生命,我們謝謝上帝,請大家同心的禱告.

我們的主我們感謝你,因為我們已經知道你已從死裡復活,將不朽的生命向我們顯露,也讓我們看見死亡已經被徹底的打破,今天生命已經可以很完全地向死亡誇勝,就讓我們在你裡頭緊緊地抓住這個盼望.人世間有許多悲歡離合,求主讓我們靠著你的恩典,好好地守著所有能夠與我們共幸福的人在一齊,來享受幸福,即使有一天幸福遭到毀壞,我們仍然相信上帝必能補救.感謝讚美你!聽我們的禱告,祈禱奉我主耶穌的名字,誠心所欲!



\newpage

\section{乞丐與財主}
\label{sec:68}
\textbf{林豪恩先生 - 乞丐與財主(The Beggar and the Rich Man)}
\newline
\newline
\hyperref[sec:67]{< < < PREV SERMON < < <}
~
\hyperref[sec:toc]{[返目錄]}
~
\hyperref[sec:69]{> > > NEXT SERMON > > >}


路加福音 16:19-31
\newline
\begin{longtable}{cl}
\hline
\hline
章節 & 經文 (和合本修訂版)\\
\hline
16:19 & \begin{tabularx}{0.7\textwidth}{X} 「有一個財主穿著紫色袍和細麻布衣服，天天奢華宴樂。 \end{tabularx} \\ \\ \relax
16:20 & \begin{tabularx}{0.7\textwidth}{X} 又有一個討飯的，名叫拉撒路，渾身長瘡，被人放在財主門口， \end{tabularx} \\ \\ \relax
16:21 & \begin{tabularx}{0.7\textwidth}{X} 想得財主桌子上掉下來的碎食充飢，甚至還有狗來舔他的瘡。 \end{tabularx} \\ \\ \relax
16:22 & \begin{tabularx}{0.7\textwidth}{X} 後來那討飯的死了，被天使帶去放在亞伯拉罕的懷裡。財主也死了，並且埋葬了。 \end{tabularx} \\ \\ \relax
16:23 & \begin{tabularx}{0.7\textwidth}{X} 他在陰間受苦，舉目遠遠地望見亞伯拉罕，又望見拉撒路在他懷裡， \end{tabularx} \\ \\ \relax
16:24 & \begin{tabularx}{0.7\textwidth}{X} 他就喊著說：『我祖亞伯拉罕哪，可憐我吧！請打發拉撒路來，用指頭尖蘸點水，涼涼我的舌頭，因為我在這火焰裡，極其痛苦。』 \end{tabularx} \\ \\ \relax
16:25 & \begin{tabularx}{0.7\textwidth}{X} 亞伯拉罕說：『孩子啊，你該回想你生前享過福，拉撒路也同樣受過苦，如今他在這裡得安慰，你卻受痛苦。 \end{tabularx} \\ \\ \relax
16:26 & \begin{tabularx}{0.7\textwidth}{X} 除此之外，在你們和我們之間，有深淵隔開，以致人要從這邊過到你們那邊是不可能的；要從那邊過到這邊也是不可能的。』 \end{tabularx} \\ \\ \relax
16:27 & \begin{tabularx}{0.7\textwidth}{X} 財主說：『我祖啊，既然這樣，求你打發拉撒路到我父家去， \end{tabularx} \\ \\ \relax
16:28 & \begin{tabularx}{0.7\textwidth}{X} 因為我還有五個兄弟，他可以警告他們，免得他們也來到這痛苦的地方。』 \end{tabularx} \\ \\ \relax
16:29 & \begin{tabularx}{0.7\textwidth}{X} 亞伯拉罕說：『他們有摩西和先知的話可以聽從。』 \end{tabularx} \\ \\ \relax
16:30 & \begin{tabularx}{0.7\textwidth}{X} 他說：『不！我祖亞伯拉罕哪，假如有一個人從死人中到他們那裡去，他們一定會悔改。』 \end{tabularx} \\ \\ \relax
16:31 & \begin{tabularx}{0.7\textwidth}{X} 亞伯拉罕對他說：『如果他們不聽從摩西和先知的話，就是有人從死人中復活，他們也不會信服的。』」 \end{tabularx} \\ \\
[1ex]
\hline
\hline
\end{longtable}
日期:2010年10月17日

主題:乞丐與財主(The Beggar and the Rich Man)

証道:林豪恩先生

經文:路加福音16:19-31節

乞丐與財主共存,似乎是由來已久,非常普遍的社會現象.古代有乞丐,也有財主,現代有乞丐,也有財主；已發展國家有乞丐,也有財主,發展中國家有乞丐,也有財主.富裕地區有乞丐,也有財主,貧窮地區當然有乞丐,但都有財主.由於「八卦雜誌」每年都編制富豪排名榜,財主的名字,人人都知曉,財主的關係,人人想結交,連國家主席都公開說:「不管時間長短,我總是要見一見李先生」.至於乞丐,最出名的一定是金庸武俠小說筆下的丐幫.丐幫領袖,武功高強,行俠仗義.誰人不懷念降龍十八掌的威風?誰人不回味打狗捧法的痛快?不過,小說歸小說,現實世界中乞丐通常都是無名無姓,無個人身份的,如果有,他們的名字仍然叫做乞丐.乞丐與財主共存的社會現象,擴大來說是貧富共存的社會現象.貧富共存的社會現象被關注的時候,就被稱為貧富懸殊問題.近日,貧富懸殊問題再度被關注,成為香港社會的熱門話題.今天的福音經課,耶穌講了一個乞丐與財主的故事,即是一個貧富懸殊的故事,讓我們思考這個故事對我們的處境有甚麼啟發.

經文

19「有一個財主穿著紫色袍和細麻布衣服,天天奢華宴樂.

20又有一個討飯的,名叫拉撒路,渾身生瘡,被人放在財主門口,

21要得財主桌子上掉下來的零碎充飢,並且狗來舔他的瘡.

22後來那討飯的死了,被天使帶去放在亞伯拉罕的懷裡.財主也死了,並且埋葬了.

23他在陰間受痛苦,舉目遠遠地望見亞伯拉罕,又望見拉撒路在他懷裡,

24就喊著說:『我祖亞伯拉罕哪,可憐我吧!打發拉撒路來,用指頭尖蘸點水,涼涼我的舌頭；因為我在這火焰裡,極其痛苦.』

25亞伯拉罕說:『兒啊,你該回想你生前享過福,拉撒路也受過苦；如今他在這裡得安慰,你倒受痛苦.

26不但這樣,並且在你我之間,有深淵限定,以致人要從這邊過到你們那邊是不能的；要從那邊過到我們這邊也是不能的.』

27財主說:『我祖啊!既是這樣,求你打發拉撒路到我父家去；

28因為我還有五個弟兄,他可以對他們作見證,免得他們也來到這痛苦的地方.』

29亞伯拉罕說:『他們有摩西和先知的話可以聽從.』

30他說:『我祖亞伯拉罕哪,不是的,若有一個從死裡復活的,到他們那裡去的,他們必要悔改.』

31亞伯拉罕說:『若不聽從摩西和先知的話,就是有一個從死裡復活的,他們也是不聽勸.』」

釋經與應用一扇可以跨越但沒有跨越的門耶穌所講的這個故事可以分為三個部份,第一部份由第十九節至第二十一節.

19「有一個財主穿著紫色袍和細麻布衣服,天天奢華宴樂.

20又有一個討飯的,名叫拉撒路,渾身生瘡,被人放在財主門口,

21要得財主桌子上掉下來的零碎充飢,並且狗來舔他的瘡.

故事中有兩個對比鮮明的人物,一個是財主,一個是乞丐.耶穌對財主的描述符合典型的財主形象:「穿著紫色袍和細麻布衣服(即是華麗的衣服)」,「天天奢華宴樂(即是過著窮奢極侈的生活)」,對乞丐的描述也符合典型的乞丐形象:「渾身生瘡(患病)」,「被人放在財主門口(被遺棄)」,「要得財主桌子上掉下來的零碎充飢(乞食)」,「並且狗來舔他的瘡(沒有自衛能力)」.縱使耶穌對財主和乞丐的描述有上述符合典型的地方,有一點卻是「非典型」的,他沒有提及財主是誰,卻給予乞丐名字──拉撒路.耶穌如此說這個故事,值得我們留意.

到底財主是如何變成富有的呢?拉撒路是如何變成乞丐的呢?耶穌的故事沒有探討這個問題.耶穌的故事的第一個部份只是描繪了一幅圖畫,一種情景,一個現象.耶穌描繪的這幅圖畫告訴我們,財主和乞丐在空間上十分接近,觸手可及.既然如此接近,財主的豐足應該很容易向拉撒路的欠缺流動,是甚麼阻礙著這種流動呢?讓我們仔細地看看圖畫,有沒有發現在財主與乞丐之間有一扇門,就是這扇門把他們分隔開.拉撒路不能跨越這扇門,接收財主多餘的東西；財主也不能跨越這扇門,把自己多於所需要的與拉撒路分享.不知道耶穌時代巴勒斯坦的建築技術有多高超,但我們可以想像,分隔財主和拉撒路的這扇門,一定是非常堅固的門,因為財主和拉撒路在有生之年,都跨不過這扇門.

富裕與貧窮,原來只是一扇門之隔.在我們的世界中,你認為這扇門是敞開的,還是關閉的呢?是開得越來越闊,還是關得越來越窄呢?是越來越容易跨越的,還是越來越難跨越的呢?在座當中不少人都在這個城市生活好了一段日子,相信大家都有獨到的觀察,以及親身的體會.

有人擔心這個城市的貧富懸殊狀況,已經燃點起人們仇富及仇商的情緒.有人急於滅火,澄清香港人既不仇富,也不仇商.何以見得呢?他們指出,第一,富豪從來都是香港人的偶像；第二,各大專院校與商業有關的科目仍然大受歡迎.不過,火熱似乎越燒越旺,連商界也感到「燒到埋身」.最近,有商界人士發表文章,指出這扇分隔貧富的門是政府關緊的,他們自己也成為受害者.

「商界克盡本分做生意,納稅,創造就業,推動經濟發展,政府利用稅收去扶貧,改善弱勢社群市民的生活,本是應有之義.可惜政府只顧累積儲備,取之於商,卻未有用之於民,市民無法享受經濟成果.社會貧富懸殊日益嚴重,市民怨氣無處宣泄,責怪政府之餘,也將不滿指向商界,漸形成一股反商情緒,商界被標籤為「官商勾結,輸送利益,無良僱主」等等,我感覺商界有冤無路訴「被屈食死貓」.(田北俊:《2.2萬億儲備再多無用 解決深層次矛盾更重要》明報2010年9月3日)

到底誰把貧富之間的門越關越緊呢?貧富之間的門是如何越來越堅固的呢?對於這些問題,不同人有不同的見解.然而,貧富之間的門越關得緊,越造得堅固,對所有人都沒有好處:政府被埋怨,富有的一方被質疑,貧窮的一方繼續受苦.

貧富之間的門越來越堅固的時候,我們的社會在貧富之門以外設立了一些緩衝的門.這些門姿態上是打開的,但在某些情況下有需要的人仍然碰壁.有一位朋友告訴我這樣的經驗:她租房獨居,又因病失業,暫時領取綜援維生.在綜援發放日子前三天,驚覺金錢用盡,得知仍有食物銀行可以提供短期支援,於是厚著面皮致電相關部門求助,解決燃眉之急.這位朋友得到的回應是:「你已經是綜援的受助者,不可以再往食物銀行領取食物.對不起,我幫不到你.」結果,這位朋友忍受了三天肚餓.

在耶穌的故事中,乞丐與財主之間的門,本質上是可以跨越的,不過,卻沒有被跨越.貧窮人士與富裕人士之間的門,豈不也是可以跨越的嗎?一扇可以跨越的門,但沒有被跨越,可以有甚麼後果呢?讓我們繼續閱讀故事的第二部份.

一個不能跨越但渴想跨越的深淵

22後來那討飯的死了,被天使帶去放在亞伯拉罕的懷裡.財主也死了,並且埋葬了.

23他在陰間受痛苦,舉目遠遠地望見亞伯拉罕,又望見拉撒路在他懷裡,

24就喊著說:『我祖亞伯拉罕哪,可憐我吧!打發拉撒路來,用指頭尖蘸點水,涼涼我的舌頭；因為我在這火焰裡,極其痛苦.』

25亞伯拉罕說:『兒啊,你該回想你生前享過福,拉撒路也受過苦；如今他在這裡得安慰,你倒受痛苦.

26不但這樣,並且在你我之間,有深淵限定,以致人要從這邊過到你們那邊是不能的；要從那邊過到我們這邊也是不能的.』

那位乞丐拉撒路死了,他很可能在某個寒冷天氣警告的日子凍死的,或者在某個酷熱天氣警告的日子熱死的,又或者在財主的桌子沒有零碎掉下來的日子餓死的.我們可以作這樣的猜測,乞丐拉撒路很可能是因為那扇分隔貧富的門關得過緊而死去的.世界上某些人的經驗就是如此,他們沒有能力扭轉命運,沒有能力改善生活,連買個中六合彩的盼望也沒有.他們在生的時候,沒有人關注,死去的時候,也沒有人記念.如果人生只是如此,真是何等的悲哀!無盡的淒涼!耶穌的故事指出,這不是乞丐拉撒路的最後結局.乞丐拉撒拉路無人埋葬,天使親自來把他接去；乞丐拉撒露無人記念,亞伯拉罕把他放在懷裡,親自安慰.乞丐拉撒路所受的痛苦,上天是知道的,有反應的,並且有補償的.

有一位崇基校友寫了一首廣受歡迎的現代詩歌,名為《尋著你的那天》,詩歌首兩句如此說:「傳說人生只不過是個夢:無方向,無意義,沒內容；貧窮或富足,結局都一樣相同,如炊煙飄散埋在糞土中!」貧窮或富足,結局真的都是一樣嗎?都是相同嗎?耶穌的故事卻指出貧窮和富足的結局可以是非常不同的.

財主也死了,他很可能被風光大殮.當他再有意識的時候,卻發現自己落在陰間受著痛苦.感受到極其痛苦的財主,舉目望見亞伯拉罕,又望見那位從前在他門外乞食的拉撒路,似乎燃起他的盼望.說好說歹,拉撒路都是靠吃我「桌子上掉下來的零碎充飢」的,要求亞伯拉罕打發他來為我效勞也不為過吧!就在這個時候,財主才驚覺他與拉撒路之間竟然隔著一個深淵.財主非常渴望能夠跨越這個深淵,但這卻是一個「人要從這邊過到你們那邊是不能的；要從那邊過到我們這邊也是不能的」(v.26)的深淵.在耶穌所講的這個乞丐與財主的故事中,死後這個分隔著財主與乞丐的深淵,彷彿回應著生前那扇分隔著財主與乞丐的門.不同的是,生前的那扇門,是可以跨越的,而死後那個深淵,卻是不能跨越的.令人感慨及遺憾的是,一扇可以跨越但沒有跨越的門,換來的原來是一個渴求跨越但不能跨越的深淵.

貧富共存之道去者已矣,來者可追.聰明的財主知道既然死後的深淵是不能跨越的,要改變結局,就唯有跨越生前的那扇門.於是,財主與亞伯拉罕關於這個話題的討論,構成了故事的第三部份.

27財主說:『我祖啊!既是這樣,求你打發拉撒路到我父家去；

28因為我還有五個弟兄,他可以對他們作見證,免得他們也來到這痛苦的地方.』

29亞伯拉罕說:『他們有摩西和先知的話可以聽從.』

30他說:『我祖亞伯拉罕哪,不是的,若有一個從死裡復活的,到他們那裡去的,他們必要悔改.』

31亞伯拉罕說:『若不聽從摩西和先知的話,就是有一個從死裡復活的,他們也是不聽勸.』」

對於如何能夠令人踏出腳步去跨越貧富之間的門,或者如何能夠令人打開這扇門,財主和亞伯拉罕有不同的意見.財主要求亞伯拉罕派遣死去的拉撒路重回人間,警告他的五個兄弟,他們就會悔改.亞伯拉罕則認為他們老早就知道應該怎樣做,因為摩西和先知早已告訴了他們,因此,若他們不去做已經知道應該做的事情,從死裡復活的人也不能對他們產生甚麼作用.

財主和亞伯拉罕所提出的方法背後的觀念和價值,十分值得我們深思.第一,財主要求一個死了的人去執行任務,是不尋常的,極端的,特殊的方法；亞伯拉罕提出聽從他們手上已經有的聖經教導,是尋常的,可理解的,人人都有的,公平的方法.第二,財主要求一個死了的人去執行任務,帶有警告及恐嚇的意味；亞伯拉罕提出聽從聖經的教導,是學習及教育的取向；第三,財主的方法著眼於後果,因為後果而改變行為,亞伯拉罕的方法著眼於行為的本質,擇善固執而不因利忘義.

財主所提出的建議,會否反應著他平時處理事情的態度呢?近日有人公開勸勸商界要「取之有道」,會否反映出某些公司已經「為求盈利,不擇手段」呢?近日又有聲音要求商界不要「賺到盡」,會否反映某些公司正在採取「渣乾渣剩」的營商策略呢?間中有些聲音說把投資撒出香港,把總部遷離香港,是否對香港人的警告,甚至恐嚇呢?有商人公開說:「商人治港是應該的,最重要是有能力.」(曾憲梓回應「胡李會」),會否反映出某些階級的特殊勢力揮之不去呢?如此種種的處事態度及處事方法,正在打開貧富之門,還是關緊貧富之門呢?正在拆除貧富之門,還是鞏固貧富之門呢?

財主與他的兄弟,手頭上都有摩西和先知留下的說話.

申命記明言:「地上的窮人永不斷絕」(15:11),所以上帝對以色列人有如此的吩咐:

1「每逢七年末一年,你要施行豁免.2豁免的定例乃是這樣:凡債主要把所借給鄰舍的豁免了；不可向鄰舍和弟兄追討,因為耶和華的豁免年已經宣告了．．．7「在耶和華－你上帝所賜你的地上,無論哪一座城裡,你弟兄中若有一個窮人,你不可忍著心,揝著手不幫補你窮乏的弟兄.8總要向他鬆開手,照他所缺乏的借給他,補他的不足.9你要謹慎,不可心裡起惡念,說:『第七年的豁免年快到了』,你便惡眼看你窮乏的弟兄,甚麼都不給他,以致他因你求告耶和華,罪便歸於你了.10你總要給他,給他的時候心裡不可愁煩；因耶和華－你的上帝必在你這一切所行的,並你手裡所辦的事上,賜福與你.11原來那地上的窮人永不斷絕；所以我吩咐你說:『總要向你地上困苦窮乏的弟兄鬆開手.』」

總結貧窮人與富有人共存,是自古以來常見的社會狀況,今時今日的香港也不例外.富有的人是香港社會的一部份,他們也可以視香港為家；貧窮人是香港社會的一部份,他們也有權在此生活.我們的社會不能,也不應該忽視任何一方.教會的本質是越超貧富的界線,超越階級的差距,教會可以思想如何在貧富差距越來越大的香港社會發揮作用,協助打開分隔貧富之間的門,創造及鼓勵彼此之間的流動及交往.反過來說,教會也可以自我反省,避免自己加入鞏固分隔貧富之門的隊伍,否則,貧窮者今生受苦,富有者來世煎熬的責任,教會也難以抽身而退.

教會呼籲社會開放分隔貧富的門的時候,若教會能夠打開自己的門,就是一個美好的示範,也可以盡僅有的能力為有需要的人解即時之困.

剛才提及的那位朋友的故事尚有續集:那一個月她捱了三天肚餓,尚未被天使接去.禍不單行,接下來的一個月,她弄壞了所租住的房間的一些簡單設備,維修的支出令她出了綜援半個月就身無分文.上次被拒絕的經驗令她猶豫再去尋找社工.我得到她的同意之後,替她聯絡她居住地方附近的教會,然後安排她往該教會申請幫助.教會當天就把她接下來半個月所需要的食物費用交給她.

如果教會的使命是幫助人們不用落入那不能跨越的深淵,耶穌所講的這個故事提醒我們,用各種方法,在不同層面,協助人們去跨越這扇本來能夠被跨越的門,以免人們墮入不能回頭的深淵.



\newpage

\section{卑微者要擁有土地}
\label{sec:69}
\textbf{鄧瑞強博士 - 卑微者要擁有土地(The Poor Shall Inherit the Land)}
\newline
\newline
\hyperref[sec:68]{< < < PREV SERMON < < <}
~
\hyperref[sec:toc]{[返目錄]}
~
\hyperref[sec:70]{> > > NEXT SERMON > > >}


詩篇 37:1-40
\newline
\begin{longtable}{cl}
\hline
\hline
章節 & 經文 (和合本修訂版)\\
\hline
37:1 & \begin{tabularx}{0.7\textwidth}{X} 不要為作惡的心懷不平，也不要嫉妒那行不義的人。 \end{tabularx} \\ \\ \relax
37:2 & \begin{tabularx}{0.7\textwidth}{X} 因為他們如草快被割下，又如綠色的嫩草快要枯乾。 \end{tabularx} \\ \\ \relax
37:3 & \begin{tabularx}{0.7\textwidth}{X} 你當倚靠耶和華而行善，安居地上，以他的信實為糧； \end{tabularx} \\ \\ \relax
37:4 & \begin{tabularx}{0.7\textwidth}{X} 又當以耶和華為樂，他就將你心裡所求的賜給你。 \end{tabularx} \\ \\ \relax
37:5 & \begin{tabularx}{0.7\textwidth}{X} 當將你的道路交託耶和華，並倚靠他，他就必成全。 \end{tabularx} \\ \\ \relax
37:6 & \begin{tabularx}{0.7\textwidth}{X} 他要使你的公義如光發出，使你的公平明如正午。 \end{tabularx} \\ \\ \relax
37:7 & \begin{tabularx}{0.7\textwidth}{X} 你當安心倚靠耶和華，耐性等候他，不要因那道路通達的和那惡謀成就的心懷不平。 \end{tabularx} \\ \\ \relax
37:8 & \begin{tabularx}{0.7\textwidth}{X} 當止住怒氣，離棄憤怒；不要心懷不平，以致作惡。 \end{tabularx} \\ \\ \relax
37:9 & \begin{tabularx}{0.7\textwidth}{X} 因為作惡的必被剪除；惟有等候耶和華的必承受土地。 \end{tabularx} \\ \\ \relax
37:10 & \begin{tabularx}{0.7\textwidth}{X} 還有片時，惡人要歸於無有；你就是細察他的住處，也不存在。 \end{tabularx} \\ \\ \relax
37:11 & \begin{tabularx}{0.7\textwidth}{X} 但謙卑的人必承受土地，以豐盛的平安為樂。 \end{tabularx} \\ \\ \relax
37:12 & \begin{tabularx}{0.7\textwidth}{X} 惡人設謀要害義人，向他咬牙。 \end{tabularx} \\ \\ \relax
37:13 & \begin{tabularx}{0.7\textwidth}{X} 但主必笑他，因見他受罰的日子將要來到。 \end{tabularx} \\ \\ \relax
37:14 & \begin{tabularx}{0.7\textwidth}{X} 惡人刀已出鞘，弓已上弦，要砍倒困苦貧窮的人，要殺害行為正直的人。 \end{tabularx} \\ \\ \relax
37:15 & \begin{tabularx}{0.7\textwidth}{X} 他們的刀必刺入自己的心，他們的弓必折斷。 \end{tabularx} \\ \\ \relax
37:16 & \begin{tabularx}{0.7\textwidth}{X} 一個義人所有的雖少，強過許多惡人的富餘。 \end{tabularx} \\ \\ \relax
37:17 & \begin{tabularx}{0.7\textwidth}{X} 因為惡人的膀臂必折斷；但耶和華扶持義人。 \end{tabularx} \\ \\ \relax
37:18 & \begin{tabularx}{0.7\textwidth}{X} 耶和華知道完全人的日子，他們的產業要存到永遠。 \end{tabularx} \\ \\ \relax
37:19 & \begin{tabularx}{0.7\textwidth}{X} 他們在患難的時候必不致羞愧，在饑荒的日子必得飽足。 \end{tabularx} \\ \\ \relax
37:20 & \begin{tabularx}{0.7\textwidth}{X} 惡人卻要滅亡。耶和華的仇敵要像草地的華美；他們要毀滅，在煙中消失。 \end{tabularx} \\ \\ \relax
37:21 & \begin{tabularx}{0.7\textwidth}{X} 惡人借貸卻不償還；義人恩待人，並且施捨。 \end{tabularx} \\ \\ \relax
37:22 & \begin{tabularx}{0.7\textwidth}{X} 蒙耶和華賜福的必承受土地；他所詛咒的必被剪除。 \end{tabularx} \\ \\ \relax
37:23 & \begin{tabularx}{0.7\textwidth}{X} 義人的腳步為耶和華所穩定；他的道路，耶和華也喜愛。 \end{tabularx} \\ \\ \relax
37:24 & \begin{tabularx}{0.7\textwidth}{X} 他雖失腳也不致全身仆倒，因為耶和華攙扶他的手。 \end{tabularx} \\ \\ \relax
37:25 & \begin{tabularx}{0.7\textwidth}{X} 我從前年幼，現在年老，卻未見過義人被棄，也未見過他的後裔求乞。 \end{tabularx} \\ \\ \relax
37:26 & \begin{tabularx}{0.7\textwidth}{X} 他常常恩待人，借貸給人，他的後裔也必蒙福。 \end{tabularx} \\ \\ \relax
37:27 & \begin{tabularx}{0.7\textwidth}{X} 你當離惡行善，就可永遠安居。 \end{tabularx} \\ \\ \relax
37:28 & \begin{tabularx}{0.7\textwidth}{X} 因為耶和華喜愛公平，不撇棄他的聖民，他們永蒙保佑；但惡人的後裔必被剪除。 \end{tabularx} \\ \\ \relax
37:29 & \begin{tabularx}{0.7\textwidth}{X} 義人必承受土地，永居其上。 \end{tabularx} \\ \\ \relax
37:30 & \begin{tabularx}{0.7\textwidth}{X} 義人的口發出智慧，他的舌頭講說公平。 \end{tabularx} \\ \\ \relax
37:31 & \begin{tabularx}{0.7\textwidth}{X} 神的律法在他心裡，他的步伐總不搖動。 \end{tabularx} \\ \\ \relax
37:32 & \begin{tabularx}{0.7\textwidth}{X} 惡人窺探義人，想要殺他。 \end{tabularx} \\ \\ \relax
37:33 & \begin{tabularx}{0.7\textwidth}{X} 耶和華必不把他交在惡人手中，當審判的時候，也不定他的罪。 \end{tabularx} \\ \\ \relax
37:34 & \begin{tabularx}{0.7\textwidth}{X} 你當等候耶和華，遵守他的道，他就抬舉你，使你承受土地；你必看到惡人被剪除。 \end{tabularx} \\ \\ \relax
37:35 & \begin{tabularx}{0.7\textwidth}{X} 我見過惡人大有勢力，高聳如本地青翠的樹木。 \end{tabularx} \\ \\ \relax
37:36 & \begin{tabularx}{0.7\textwidth}{X} 有人從那裡經過，看哪，他已不存在，我尋找他，卻尋不著了。 \end{tabularx} \\ \\ \relax
37:37 & \begin{tabularx}{0.7\textwidth}{X} 你要細察那完全人，觀看那正直人，因為和平的人有好結局。 \end{tabularx} \\ \\ \relax
37:38 & \begin{tabularx}{0.7\textwidth}{X} 至於罪人，必一同滅絕，惡人的結局必被剪除。 \end{tabularx} \\ \\ \relax
37:39 & \begin{tabularx}{0.7\textwidth}{X} 義人得救是出於耶和華，在患難時耶和華作他們的避難所。 \end{tabularx} \\ \\ \relax
37:40 & \begin{tabularx}{0.7\textwidth}{X} 耶和華幫助他們，解救他們；他解救他們脫離惡人，把他們救出來，因為他們投靠他。 \end{tabularx} \\ \\
[1ex]
\hline
\hline
\end{longtable}
日期:2010年10月24日

主題:卑微者要擁有土地(The Poor Shall Inherit the Land)

証道:鄧瑞強博士

經文:詩篇37:1-40

各位弟兄姊妹,主內平安.今日講道的經文是詩篇37篇,今日的詩篇經課.詩篇37篇的主題,和詩篇第1篇的主題,有很多相似的地方.

詩篇第1篇是全部詩篇的引言和總綱,那裡提到惡人和義人不同的道路；提到惡人思想種種害人的計謀,義人卻思想神的律法；提到惡人像糠秕被風吹散,義人卻像栽在溪水旁的果樹；提到惡人的路至終引向滅亡,義人卻蒙神賜福.這些主題,都在詩篇37篇再次出現.

信仰,就是在不義的世界,對抗不義；面對不敬虔的人,堅守敬虔.這是何等的難.在一個以享受物質為快樂基礎的世界,講求心靈的滿足,何等艱難.在一個以賺錢多少來衡量人的價值的世界,要真心尊重貧窮人,何等艱難.世俗的不義像洪流,我們在洪流上坐著船.順著洪流而下,易；逆水行舟,難.故此,詩篇每隔不久便會重提詩篇第1篇的主題.今天要講的詩篇37篇,可說是詩篇第1篇的回響.

(1)面對惡人的成功,不要嫉妒詩篇37篇的開頭,開宗明義是一個勸告.詩 37:1

不要為作惡的心懷不平,也不要向那行不義的生出嫉妒.

面對惡人,特別是面對成功的惡人,詩人勸喻信徒,不要「心懷不平」.但看來「惡人」是這世界的常態,他們看來很成功,人很易跟從他們的腳步行,故這勸告在這詩篇中反覆出現.詩 37:7

你當默然倚靠耶和華,耐性等候他；不要因那道路通達的和那惡謀成就的心懷不平.詩 37:8

當止住怒氣,離棄忿怒；不要心懷不平,以致作惡.

詩人這樣再三叮嚀,像母親不厭其煩地勸導兒子.詩人越囉唆,越發顯示信徒很易對成功的惡人「心懷不平」.「心懷不平」,有「怒火中燒」的含意,是一種有「熱度」的「憤怒」.但這種「憤怒」同時又夾雜著「嫉妒」.就如:你見到人作弊,你很憤怒,但同時心中又羨慕一個人可以這樣輕而易舉地成功.你見到一個人犧牲了一點色相便得到一個高位,你很憤怒,但同時你又妒忌那一點可以出賣的色相.你見到一個有錢人在樓宇的炒賣中獲得暴利,你很憤怒,但同時心中又羨慕他可以這樣輕而易舉地賺這麼多錢.面對這個走歪了的世界,詩人說:詩 37:1

不要為作惡的心懷不平,也不要向那行不義的生出嫉妒.

惡人做了什麼惡事呢?這詩篇有四節經文提到惡人做的惡事.詩 37:12

惡人設謀害義人,又向他咬牙.詩 37:14

惡人已經弓上弦,刀出鞘,要打倒困苦窮乏的人,要殺害行動正直的人.

詩 37:21

惡人借貸而不償還.詩 37:32

惡人窺探義人,想要殺他.惡人的惡,是他們的思想,針對義人,咬牙切齒地向義人發怒.犯錯的政府,一定不想別人揭露他的錯,揭露的人可能犯上顛覆國家罪.犯錯的有錢人,一定不想別人揭露他違規買賣,揭露的人是阻他發財.面對「困苦窮乏的人」,惡人不會幫助他們,因為弱者正是被他們魚肉的最好對象.在學校裡,同學間的欺凌事件,被凌虐的都是勢孤力弱的弱者.學校如是,社會也如是.「惡人借貸而不償還」,更突顯他們只會取而不會給.近來美國巴菲特及蓋茨在中國擺豪門宴,有些中國富豪不敢赴宴,是因為怕被勸捐身家.其實,我們也會同樣怕被人勸捐.我們不覺得需要向社會作出償還.這是否也是惡人的一種根深柢固的思想方式呢?

邪惡作為一種力量,是有針對性的,它針對正義.惡人不會幫助弱者,反魚肉他們.惡人不會施與,只會拿取.惡人不會容忍正義,只望正義從此消失.邪惡想毀滅正義.惡人的世界,不講神的正義,不講憐憫,不講彼此幫助.當一個社會,只講強者的賺錢機遇而忽略草根階層的經濟機會,只講經濟效益而忘記人間道義,只講物質的享樂而無視心靈的需要,這是否正是一個惡人社會的寫照?在惡人的世界,一個人很容易「心懷不平」與「嫉妒」.我們明白人生的價值不在擁有,但我們卻羨慕擁有的人.我們明知屬靈品格的長進是好的,但我們花更多精力去爭奪世間的權力.我們明知不勞而獲是不好的,卻時常注目六合彩那過億元的彩金.這是我們的墮落.

(2)要倚靠神世界是充滿惡人惡事的,但我們要信靠神.什麼是信靠神?詩人用多角度來描述.詩 37:1

不要為作惡的心懷不平,也不要向那行不義的生出嫉妒.詩 37:2

因為他們如草快被割下,又如青菜快要枯乾.

詩 37:35

我見過惡人大有勢力,好像一根青翠樹在本土生發.詩 37:36

有人從那裡經過,不料,他沒有了；我也尋找他,卻尋不著.

惡人很有勢力,但從信仰講,這是短暫的.信仰的人要有一信仰的視界,看到什麼是短暫的,什麼才是永恆的.將眼前的成敗,當作是最重要的,這是惡人的看法.問題是:信徒也沾上這種做人態度.詩人提醒我們,惡人的繁榮是空洞的.我們應追求永恆的事.有什麼東西具有永恆價值呢?詩 37:3

你當倚靠耶和華而行善,住在地上,以他的信實為糧；詩 37:4

又要以耶和華為樂,他就將你心裡所求的賜給你.倚靠神的人,要行善.「以神的信實為糧」:即相信神應許的,不要相信世俗所應許的.世俗應許我們,葯物可帶來健康,美食可帶來快樂,外表可帶來成功.我們的確需要葯物,美食,外表,但我相信神的應許更真實,倚靠神而心中平安的人,更健康；願意分享施與的人,更快樂；遵行神法度的人,更成功.「以神為樂」:這就多於倚靠神了,這是積極地以神為福樂的根源.沒有這正面追求,很難克服心中的「不平」與「嫉妒」的.

詩 37:7

你當默然倚靠耶和華,耐性等候他；不要因那道路通達的和那惡謀成就的心懷不平.詩 37:8

當止住怒氣,離棄忿怒；不要心懷不平,以致作惡.倚靠神的人,要忍耐,要等候,要放下怒氣,要放下忿怒.若耿耿於懷,則最終你便可能會作惡,以暴易暴.你對「作威作福」的人心懷不平,你可能就會期望爬得比他高,然後向他「作威作福」.你變成你原先反對的惡人.

詩 37:21

惡人借貸而不償還；義人卻恩待人,並且施捨.相對於惡人只懂取而不懂給,倚靠神的人卻懂得分享和施與.這涉及對人的基本態度.其他人在惡人眼中,是競爭對手,是他們欺負的對象.簡單而言,他人的存在,是為了自己的好處的.他人只是惡人的工具,只此而已.倚靠神的人,詩篇稱他們作義人,其他人是他們的鄰舍,是朋友,是人生的同伴.生命是分享的.

(3)倚靠神的人的盼望:承受地土這詩篇有一個獨特的主題,就是信靠神的人會擁有土地.我們不知道這詩篇的確實背景,可能是某些人強佔了弱者的家園,又可能是強國奪了他們的土地.惡人扮演了強者的角色,在他們面前,信靠神的人只是弱者.弱者在世上仍有盼望嗎?詩人清楚說:有,弱者要承受地土.詩 37:9

因為作惡的必被剪除；惟有等候耶和華的必承受地土.詩 37:11

但謙卑人必承受地土,以豐盛的平安為樂.詩 37:22

蒙耶和華賜福的必承受地土；被他咒詛的必被剪除.詩 37:29

義人必承受地土,永居其上.詩 37:34

你當等候耶和華,遵守他的道,他就抬舉你,使你承受地土；惡人被剪除的時候,你必看見.

地土,是人賴以為生的最基本元素.有了土地,才能耕作,才能生產,才能有食物,才能安居.工作,食物,居所,這是生活最基本的元素.在惡人的世界,人最卑微的生活所需都被剝奪.上一期＜時代論壇＞(2010年9月26日)的封面故事,就是三個「露宿者」的故事.這些就是無工作,無食物,無居所者的故事.其中一個故事是這樣寫的:「曾經在麥當勞,露宿在沙田城門河附近一年多的Grace,多少日子飢腸轆轆,在擔心惶恐間,原來還有人關心她.【Grace說】:『有以前中學同學經過,知道我露宿,就每日返工時來到城門河,給我一個包和一包飲品,即使她趕不及買,都會來到,於低二十元給我.她還抱歉未能請我到她家裡睡.』一杯涼水,今天回想,Grace無言感激.」故事的主人翁變成露宿者,固然有她的前因,摧毀我們生命的力量在我們之外,也在我們之內.但叫她重新振作的力量,源自無條件的愛.所謂信靠神,就是相信愛而不相信恨,相信神應許的明天而不相信可悲的過去,相信天地有情,而我們可以有工作,食物,居所,有一片生活的地土.神會作工,神也需要信靠他的人與他同工.神需要更多人改變人生態度,神呼喚更多人去改變邪惡的結構.不再視人為工具,而是視人如朋友；不再榨取別人,反向他們施與；不犧牲他人而建立自己,而是互相甘願犧牲而共度難關.後來主耶穌進一步說:「溫柔的人有福了!因為他們必承受地土.」(太5:5)溫柔的人,是那些扶持他人而不高抬自己的人.地土在他們手中,就不是用來抬高地價,而用來作為所有人賴以為生的憑藉.地土在他們手中,才能造福所有人.

(4)結語我以兩節經文作結.詩 37:12

惡人設謀害義人,又向他咬牙.詩 37:13

主要笑他,因見他受罰的日子將要來到.這詩句使我想起蘇東坡＜赤壁懷古＞裡的兩句詞句:「談笑間,強虜灰飛煙滅」.赤壁一役,面對曹操大軍,周瑜卻在談笑間從容面對.在這個貧弱者越來越難找到生活空間的世界,邪惡的力量看似當道,詩人卻看到那顧念世人的神,他在高天,俯視眾生,他不是無情地觀看,而是審判惡人.在談笑間,惡人便要灰飛煙滅.神是真實的,他在掌權.故此,「不要為作惡的心懷不平」,不要灰心失望；要樂善好施,要行走正路.讓我們相信:在神的恩典裡,卑微者要擁有土地.

但願榮耀歸於聖父,聖子,聖靈.阿們.



\newpage

\section{基督的特使}
\label{sec:70}
\textbf{王群女士 - 基督的特使}
\newline
\newline
\hyperref[sec:69]{< < < PREV SERMON < < <}
~
\hyperref[sec:toc]{[返目錄]}
~
\hyperref[sec:71]{> > > NEXT SERMON > > >}


哥林多後書 5:11-21
\newline
\begin{longtable}{cl}
\hline
\hline
章節 & 經文 (和合本修訂版)\\
\hline
5:11 & \begin{tabularx}{0.7\textwidth}{X} 既然我們知道主是可畏的，就勸導人；但是神是認識我們的，我盼望你們的良心也認識我們。 \end{tabularx} \\ \\ \relax
5:12 & \begin{tabularx}{0.7\textwidth}{X} 我們不是向你們再推薦自己，而是要讓你們有誇耀我們的機會，使你們好面對那憑外貌、不憑內心誇耀的人。 \end{tabularx} \\ \\ \relax
5:13 & \begin{tabularx}{0.7\textwidth}{X} 如果我們癲狂，是為神；如果我們清醒，是為你們。 \end{tabularx} \\ \\ \relax
5:14 & \begin{tabularx}{0.7\textwidth}{X} 原來基督的愛激勵我們；因我們這樣斷定，一人既替眾人死，眾人就都死了； \end{tabularx} \\ \\ \relax
5:15 & \begin{tabularx}{0.7\textwidth}{X} 並且他替眾人死，是叫那些活著的人不再為自己活，乃為替他們死而復活的主活。 \end{tabularx} \\ \\ \relax
5:16 & \begin{tabularx}{0.7\textwidth}{X} 所以，從今以後，我們不再按照人的看法來認識人，縱使我們曾經按照人的看法認識基督，如今卻不再這樣認識他了。 \end{tabularx} \\ \\ \relax
5:17 & \begin{tabularx}{0.7\textwidth}{X} 所以，若有人在基督裡，他就是新造的人，舊事已過，都變成新的了。 \end{tabularx} \\ \\ \relax
5:18 & \begin{tabularx}{0.7\textwidth}{X} 一切都是出於神；他藉著基督使我們與他和好，又將勸人與他和好的使命賜給我們。 \end{tabularx} \\ \\ \relax
5:19 & \begin{tabularx}{0.7\textwidth}{X} 這就是：神在基督裡使世人與自己和好，不將他們的過犯歸到他們身上，並且將這和好的信息託付了我們。 \end{tabularx} \\ \\ \relax
5:20 & \begin{tabularx}{0.7\textwidth}{X} 所以，我們作基督的特使，就好像神藉我們勸你們一般。我們替基督求你們，與神和好吧！ \end{tabularx} \\ \\ \relax
5:21 & \begin{tabularx}{0.7\textwidth}{X} 神使那無罪的，替我們成為罪，好使我們在他裡面成為神的義。 \end{tabularx} \\ \\
[1ex]
\hline
\hline
\end{longtable}
日期:2010年10月31日

主題:基督的特使

証道:王群女士

經文:哥林多後書5:11-21

主內弟姊妹,平安!

引言最 近讀到丹麥哲學家祈克果(Soren   Kierkegaard)一句話:「生命只有在回望中才能了解,只有往前看才活得下去.」(中文翻譯似乎未能把其中的精粹表達出來)(Life can   only be understood backwards; but it must be lived   forwards.)覺得很有意思,這句短小精警的說話,卻涵蓋了時光長河的三個向度,除了文意上顯然而見的過去及將來外,亦包括了此時此刻,因為它所談   論的根本就是當下的生活態度.人漸入中年之境,開始逐步體會到很多人生中的境遇是要經過了,回頭望,再思想,才能稍稍參透箇中的意義.然而人又不能老是回   頭望,只管著想從前,否則活在當下的我們就會原地踏步,失卻走下去的動力,因此我們必須同時又要被對將來的憧憬和盼望所拖動向前.今天我們選讀的書信經課  林後五章11-21節,同樣向我們展示這三個時間的向度,保羅提醒我們的將來和過去都是影響和塑造住今天的你我.

一,面向將來──我們在主面前,在人面前都是赤露敞開的在   11-13節,我們可以從字裡行間感受到保羅正面對極大的挑戰,在哥林多教會裡,他的使徒身份受到質疑,保羅當然要盡力捍衛,辯護其個人的品格,聲譽及身   份,保羅的辯才了得是不爭的事實,但更重要的是,他不單是要在言辭上說服哥林多教會,他更需要擺上他的整個生命作見證,才能說服別人接受他是使徒的身份,   他整個生命,包括:他的行事,為人,思想,言語,態度等等,生命是一體而連結的,並不是割裂分離的,正如我們不可以一方面立志要盡忠事奉,但另一方面就偷   懶,卸膊,走後門；又或者這邊廂我們立志要愛人如己,那邊廂就講是非,玩杯葛,搞小圈子.我們的生命是一個整體,第11節下半部更提醒我們,我們整個的生  命「

在主面前是顯明的」,新漢語譯本把「顯明的」翻譯作「赤露敞開的」,要 突  顯出所指的意思不單是出現或顯露,程度上是更為強烈的,是指整個人被公開,被揭露,是完全無可遮掩的.在11節上半部保羅清楚知道「主是可畏的」,要明白  他何出此言,我們一定要看上文,因為在第10及11節之間原文是有一個連接詞,新漢語譯本是如此翻譯第11節的,請聽我讀出:「

所以,我們既然知道主是可畏的,就勸勉人.」在此「所以」這個連接詞是要突顯出接著下來的論述是承接上文第10節(或者更準確地說是1-10節這一段落)而歸結出來的.讓我們一齊看第10節,保羅明確指出:「

因為我們眾人必要在基督臺前顯露出來,叫各人按著本身所行的,或善或惡受報.」這裡所描述的是在終末審判的日子,我們都要在主面前呈上自己的生命,向祂交賬,並且要按我們在世時的所做的,得著從上主而來的相應獎賞或者懲罰.在審判的主的面前,我們生命的裡裡外外都被攤開,無一樣東西可以被收藏,掩飾.

保羅在11節的下半句不單說「

我們在神面前是顯明的」,他繼續講:「

盼望在你們的良心裡也是顯明的」,   我們除了將來要向審判的主交賬外,雖然現在那審判日子尚未來到,但是我們的生命也是無法收藏得妥妥貼貼的.我們身邊周圍的人,他們跟我們一起生活,一起學   習,一起工作,一起事奉,我們講的那一套和實際做的那一套,他們都聽在耳中,看在眼內,他們都是驗證我們基督徒的身份和我們所作所為是否一致的見證人.記   得有一次,我和朋友去一間大教會參與聚會,由於出席的人數眾多,門口一早站滿衣著整齊得體,笑意盈盈的司事人員,列隊歡迎所有來賓.此時,我的朋友在我耳   邊輕說地說:前面最左邊的那個司事,我曾跟他一起上過一個課程.我從朋友的語氣中已感愛到一份不屑的味道；接著他繼續嘀咕:想不到他原來也是基督徒,他上   課遲到早退,做小組功課時又推卸膊,又吞泡,「做D 唔做D 」.對於他的投訴,我沒有什麼可以回話,只在心裡牢牢地警惕自己,我要謹慎自己的言行,原來我的   生命不單要向上主交賬,而我的一舉手一投足同時亦影響身邊的人如何看我這個基督徒呢?因此,無論我們願意與否,是有心栽花,還是無心插柳,我們在主面前,  在人面前都是赤露敞開的,正如保羅在林前四章9節用人生如戲來打方,他說:「

我們成了一臺戲,給世人和天使觀看.」而我們又有沒有忠心地演好上主交託予我們角色呢?而這個角色是終生制的,並且是屬於全方位的角色,我們在生命中的任何一個場景都要落力地演.還有這不是一場獨腳戲,而是一場一個也不能少的群戲,是整個信仰群體一齊參演的戲.

各位弟兄姊妹,當我們被預告知道我們要走向一個如何的終局時,我們這個認知有沒有影響或者指導我們當下如何生活呢?今天我們活出的又呈現一個怎樣的生命狀態呢?

二,全因過去──我們今天乃為替我們死而復活的主活保羅立定志向一生單為主而活,為上帝對他的召命而奔走,除了因為他清楚知道將來要向審判的主交賬之外,亦是因為主耶穌基督在過去所成就的一切,那份大愛強力地催逼著他要以生命來回應.我們一齊看14-15節:「

原來基督的愛激勵我們.因我們想一人既替眾人死,眾人就都死了；並且他替眾人死,是叫那些活著的人不再為自己活,乃為替他們死而復活的主活.」我們現在享有的生命是耶穌基督過去用祂的性命來買贖回來的,叫我們不致於淪亡,而耶穌基督這個舉動就把祂對我們的愛具體顯明出來,從此我們的生命不再屬於我們的,而是屬於那位救贖我們,又賜下永生盼望給我們的救主耶穌的.第17節進一步描述我們所擁有的生命之的特點:「

若有人在基督裡,他就是新造的人,舊事已過,都變成新的了.」   保羅強調我們「在基督裡」的話,我們的生命就不再一樣,我們成為新造的人,我們的身份改變了,我們跟上帝與人的關係都改變了,我們的人生軌跡改變了,我們  換上了一個全新的鏡片,一個與世界價值觀不同的鏡片來透視這個世界,但我想繼續問的是,這個改變只是停留在你頭腦上的認知,還是真正帶來你生命的徹底改   變?我認識一位朋友,他是來自內地的,他本身在北京最高學府學士,碩士畢業,他品學兼優,加上身家清白,所以順利入黨,在人眼中可謂是踏上青雲路,基本上   已經半隻腳踏入官途,前景一片光明.後來他有機會來香港進修攻讀博士課程,在這個期間,他與耶穌基督相遇,他雖然自小就是一位無神論者,但是他最終被基督   完全的愛所折服,他的生命被徹底改變,他做了一個世人眼裡看為「憨居」的選擇,他決定要獻上自己年輕的生命,為主而活,一生事奉上主,所以他留在香港,放   棄了原本的專業,放棄了已經鋪好了康莊大道,他入讀神學院,重新學習,學做一個傳道者.我相信這個不是一個很特別的例子,只不過它是個真實的例子,告訴我  們基督的愛是何其的實在,是何其的強而有力!

我想問在座的弟兄姊妹:耶穌基督在過去為我們所作的犧牲,對我們有什麼意義呢?是否只停留在決定信主那一刻的感動,感動過後,我仍然是我,我仍然是我的主人,我的生活仍然圍著我自己來旋轉呢?第15節:「

並且他替眾人死,是叫那些活著的人不再為自己活,乃為替他們死而復活的主活.」   保羅提醒我們,今天我們的生命是因為耶穌基督過去為我們所做的一切而變的不一樣,而這種變化不是我們在頸項上由原本的水晶,鑽石項鍊變成一條十字架項鍊；   也不是吃飯前由不假思索起筷下箸變為會先閉目謝飯祈禱；更不是每個星期天由原本去茶樓飲茶,或者繼續懶在床上,而變成按時來到禮拜堂敬拜上帝.這種改變是   生命焦點的改變,我們由原本為自己而活,變成為耶穌基督──那位為我們死而復活的主──而活.這會觸及我們的價值觀,這會影響到我們人生的優先次序,這亦  會改變我們行事為人的模式,這種變化是由信仰引發的,那是既外顯又內蘊的,徹底而又持久的.

三,忠於現在──我們要作基督的使者傳和好的信息保羅在第18節說:「

一切都是出於神,他藉著基督使我們與他和好,又將勸人與他和好的職分賜給我們.」這裡清楚指出上帝在基督裡賜予我們獨特而尊貴的身份,我們的身份就是20節所言的「

基督的使者」,   新漢語譯本把「使者」翻譯成「特使」,其中的意思可能會更加突顯一點,「特使」所擁有的權力並非來自身的,而是來自委派他的那一位,情況就如我們在電視熒  幕中常見的,在古裝片集裡,一個地位卑微的小官或者太監帶住皇上的聖旨去到一品大官的府第,當聖旨打開,讀出:「奉天承運   皇帝詔曰」,這位身居高職的官員便要馬上向面前的特使,屈身下拜,洗耳恭聽.作為基督的特使我們的身份是尊貴的,當然有威風八面的時刻,但保羅亦沒有掩飾  我們作為特使要冒上的風險,要付上的代價,請聽我讀出以弗所書六章20節上半句:「

(我為這福音的奧秘作了帶鎖鏈的使者)」在此保羅用上同一個字,但卻描述自己是帶鎖鏈的使者,表示他正是身陷監牢的.而特使的任務就是無論他遇上任何的情況,他代表委派他的那位,只傳講委派他那位所要他傳講的信息.那麼上主要我們傳講怎樣的信息呢?保羅在林後五章18節及19節告訴我們所要傳揚的是「

和好的道理」,這是我們親身經歷過在主基督裡得以與神與人和好的寶貴經驗,我們一方面是受惠者,另一方面被委託去傳講這信息,透過我們的生命見證,叫其他人也藉著基督與神和好.

換   一個角度來思想,要驗證我們有沒有在此時此刻盡忠職守,按上主的吩咐成為基督特使的方法,就是檢視一下我們身邊的人,到底他們因為我們當下的生命見證,而   與神的距離拉近,達至和好?還是他們因為我們而與上帝的距離越來越遠呢?所謂「成也蕭何,敗也蕭何」,我相信大家也接觸過一些朋友,他們始終不願意接受福   音的原因,就是曾經遇上一些很不濟的基督徒,這叫他們在福音的門外卻步.原本應當傳揚和好信息的基督特使,反而成為別人與神和好的絆腳石,這是值得我們深   切反省的,到底我們有沒有忠於所託呢?然而,我們亦聽過更多弟兄姊妹的得救見證,其中引領他們一步一步走近主耶穌的是身旁的基督徒,可能是他們的家人,朋   友,同學,同事,無論是哪種身份,但是他們都有一共通點,就是他們的美好見證都能吸引人思想:到底影響這人能夠活出如此生命的,是一位怎樣的上主?這讓我   想起馬尼拉人質事件中的倖存者──梁太,她在一瞬間失去了至愛的丈夫和兩個乖巧的女兒,而唯一的兒子則頭部重創,大家都為她落下同情,傷心的眼淚,大家更   為她一家的遭遇感到憤慨,在所有人的眼中梁太是最有資格憤怒,發怨言,提出控訴的人,但是我們卻在這個瘦小,看來弱不禁風的女子身上看見強而巨大的力量,   她不單沒有口出一句惡言,反而她一而再地反過來安慰港人,叫我們不要再為他們一家傷心流淚,又鼓勵香港人要一起堅強生活下去,這力量是來自她內裡的寬容,   來自她對上主的信靠,以及她對將來的盼望,她深信兒子在醫生及她的悉心照料下,會逐漸康復過來,會有一日能夠親自答謝香港人的關心；她亦肯定在終末耶穌基   督再臨掌權的時候,她能夠和丈夫及兩個女兒在天家重聚.她雖然沒有用言語去講保羅所言「與神和好」的道理,我們卻見到她用生命來演繹出這個道理,她對上主  的篤定,信靠,順服,確實拉近了不少人(包括了基督徒在內)與上主的距離.

結語今年的諾貝爾和平獎得主已經塵埃落定,根據評選 準  則,和平獎是頒給「為促進民族團結友好,取消或裁減常備軍隊以及為和平會議的組織和宣傳盡到最大努力或作出最大貢獻的人」.我並非打算在講壇上評論劉曉波   的獲獎是否實至名歸,但我想指出的是:有和平獎的設立可能正正反映出真正的和平尚未臨到世上,同時大家對和平獎的重視亦表達了世人對和平的渴望.作為基督   徒我們的身份就是基督的特使,我們今天在世上的任務是奉上主之命,向世界宣告和好的信息,並以生命忠心地見證出,真正的和平是始於我們在基督裡先與上主和  好.

願上主因為我們今天所作祂的心意得著滿足,又願真正的和平早日臨到世上,阿們.



\newpage

\section{路加福音的最後財主}
\label{sec:71}
\textbf{林豪恩先生 - 路加福音的最後財主}
\newline
\newline
\hyperref[sec:70]{< < < PREV SERMON < < <}
~
\hyperref[sec:toc]{[返目錄]}
~
\hyperref[sec:72]{> > > NEXT SERMON > > >}


日期:2010年11月7日

主題:路加福音的最後財主

証道:林豪恩先生

經文:路加福音12:16-21,16:19-31,18:18-30

一．    引言路加福音記述了四位財主的故事,其中兩位是耶穌講述的故事中的角色,另兩位是耶穌遇上的人物.耶穌曾經說:「

有錢財的人進神的國是何等的難哪!駱駝穿過針的眼比財主進神的國還容易呢!」(18:24-25)有錢財的人進神的國真的是如此艱難嗎?這些艱難是如何構成的呢?把這四位財主的故事拼湊一起,又有另一番觀感.讀者無論自覺是財主,或不是財主,都可以在他們的故事中得到一點啟發.

二．    第一位財主12:16

(耶穌說)「

有一個財主田產豐盛；12:17

自己心裡思想說:『

我的出產沒有地方收藏,怎麼辦呢?』12:18

又說:『我要這麼辦:要把我的倉房拆了,另蓋更大的,在那裡好收藏我一切的糧食和財物,12:19

然後要對我的靈魂說:靈魂哪,你有許多財物積存,可作多年的費用,只管安安逸逸地吃喝快樂吧!』12:20

神卻對他說:『無知的人哪,今夜必要你的靈魂；你所預備的要歸誰呢?』12:21

凡為自己積財,在　神面前卻不富足的,也是這樣.」

耶穌口中的第一位財主,被評為「無知」,因為他用「財物」去款待「靈魂」.如果以為人生的所有需要就是財物,因此以物質去滿足非物質的需要,就好象緣木求魚,徒勞無功,甚至弄巧反拙,得不償失.若我們細心觀察,財主這種「以財物去餵養心靈飢渴」的無知,日常生活中隨處可見.所以,當我們準備嘲笑財主的時候,不妨自我反省一下.

例如:進食的基本功能是滿足身體的需要,然而,有些時候,肚腹填飽了,還是要吃；「儲備」過多了,還繼續吃.大家有沒有留意,人們不能自制地吃的時候,通常處於甚麼狀態?有一位朋友有這樣的經驗:

「一年半前,愛情,家庭,事業,朋友...... 竟同時令我徹底失望和迷網,真的再承受不下．．．以為狂食可以逃避問題,但越食越喊,越喊越有不想生存的想法．．．接受專業幫助後,發覺原來從小到大,有些不愉快的事其實一直纏繞著我.可能基於過份自卑,不敢動身改變；又或者過份自信,從來都看不見自己的問題.所以情緒稍一波動,便選擇暴食去逃避最根本的問題.(進食失調患者心聲\url{http://www.heda-hk.org/chi/08case_02patient.php})

環顧四週,我們會發現更多以物質去解決非物質需要的例子:例如有些父母因工作或個人原因未能經常陪伴孩子,在情感虧欠下,往往不自覺地以金錢物質作為缺席的補償.又例如有些政府避免處理不公平的「遊戲規則」,往往以「派糖」去安撫市民.又例如有些政權以GDP數字,城市化的速度,建築物的高度來取代人權,自由,民主的發展.只懂得或者只能夠以物質去解決非物費需要的人,往往感到大惑不解:為甚麼供給了孩子那麼多金錢和物質,仍然感到與孩子疏離呢?為甚麼屢次「派糖」,民怨仍沸騰呢?為甚麼成為了經濟大國,仍得不到人們衷心的尊重呢?有些人說:「錢可以解決的問題,就不是問題!」原來,人生還有很多問題,不是錢可以解決的.

耶穌總結這個無知財主的故事時說:「凡為自己積財,在神面前卻不富足的,也是這樣.」(v21)奧古斯丁在《懺悔錄》中如此說:『神啊!你創造我們的心靈,留下了空位,除了你之外,世間沒有任何東西可以滿足!』

某位城中富豪曾有這樣的經歷:「1994年我和太太剛旅行回港,有點jetlag,睡得不好,又有點坐立不安.太太盡量安慰和幫助我舒緩心情,但我仍總感到有股空虛、惘然,甚至恐懼.我知道在許多人眼中,我甚麼都有了:父親打好根基的一番事業由我們兄弟繼續營運,一個美好的家庭,賢慧聰明的太太,精乖健康的子女,自己身體也頗健康．．．但總是經歷著那股無法,也無物可填補的空虛.我便向太太說:「我們不如去旅行吧.」太太說:「我們剛剛回來呢!」(郭炳江\url{http://www.alpha.org.hk/tchi/peoplec/newcreationc/thomaskwokc.htm})

意識自己有物質以外的需要,承認這些非物質的需要不是物質所能滿足,體會到財物不足恃,是打破無知的開始.八福的第一福宣告:「心靈貧窮的人有福了!因為天國是他們的.」(馬太福音5章3節)

三．    第二位財主16:19 「

有一個財主穿著紫色袍和細麻布衣服,天天奢華宴樂.16:20

又有一個討飯的,名叫拉撒路,渾身生瘡,被人放在財主門口,16:21

要得財主桌子上掉下來的零碎充飢；並且狗來舔他的瘡.16:22

後來那討飯的死了,被天使帶去放在亞伯拉罕的懷裡.財主也死了,並且埋葬了.16:23

他在陰間受痛苦,舉目遠遠地望見亞伯拉罕,又望見拉撒路在他懷裡,16:24

就喊著說:『我祖亞伯拉罕哪,可憐我吧!打發拉撒路來,用指頭尖蘸點水,涼涼我的舌頭；因為我在這火焰裡,極其痛苦.』16:25

亞伯拉罕說:『兒啊,你該回想你生前享過福,拉撒路也受過苦；如今他在這裡得安慰,你倒受痛苦.16:26

不但這樣,並且在你我之間,有深淵限定,以致人要從這邊過到你們那邊是不能的；要從那邊過到我們這邊也是不能的.』16:27

財主說:『我祖啊!既是這樣,求你打發拉撒路到我父家去；16:28

因為我還有五個弟兄,他可以對他們作見證,免得他們也來到這痛苦的地方.』16:29

亞伯拉罕說:『他們有摩西和先知的話可以聽從.』16:30

他說:『我祖亞伯拉罕哪,不是的,若有一個從死裡復活的,到他們那裡去的,他們必要悔改.』16:31

亞伯拉罕說:『若不聽從摩西和先知的話,就是有一個從死裡復活的,他們也是不聽勸.』」

面對真相,如果第一位財主的反應可能是:「原來係咁既咩?我無唸到先!」,第二位財主的反應就是:「真係咁既咩?我以為講下架炸!」所以,如果第一位財主是「無知」的話,第二位財主就不能夠再稱為「無知」,他只是「扮無知」.第一,他不知道陰間和他的先祖亞伯拉罕所在之處是不相通的；第二,他不知道如何避免落在陰間受苦.回應財主提出打發拉撒路從重回人間去向他的兄弟發警告的要求,亞伯拉罕指出:「他們有摩西和先知的話可以聽從.」既然有摩西和先知的話,財主沒有理由不知道上帝對他的期望和要求,也不會不知道通往上帝之路與通往陰間之路的分別.所以,他們不能算作「無知」,最多只可以稱為「扮無知」.財主可能被眼前「紫色袍和細麻布衣服」矇閉了眼睛,對摩西的律法視若無睹；同時被「奢華宴樂」充塞了耳朵,對先知的警告聽而不聞.所以,亞伯拉罕認為,就算拉撒路重回人間去警告財主的兄弟,他們的回應可能是:「唔好嚇我啦,摩西同先知都係講下既即,如果真係咁認真,我地而家就唔可以繼續榮華富貴,大魚大肉啦.」

創世之初,狡猾的蛇也是以類似的語氣試探夏娃:「上帝果真不許你們吃那棵樹上的果子嗎?吃了一定會死嗎?或者不一定死呢?就算一般人吃了會死,但你可能不同呢!」夏娃面對美食當前,色香味俱全,兼有提精補腦的作用,「扮無知」,吃了才算.

以第二位財主為代表的這種「知道而當作不知道」的態度,在今天的社會也是隨處可見的.

約三年前,一輛貨車失事打橫翻側,壓著一輛新界的士,並拖行約50米.的士司機和5名乘客當場死亡.受輕傷的肇事貨車司機酒精測試超標逾3倍.(星島2009年1月23日)約二年前,一名老師早上回學校參加聖誕聯歡會,途經觀塘利安道時,一輛由醉酒超標兩倍的年輕司機駕駛的私家車失控剷上行人路,老師被撞飛拋出七十呎跌入水渠內,當場慘死,另一名女途人則被撞傷.(星島2009年12月25日)

酒後駕駛,甚至吸食毒品後駕駛,並沒有因為上述的意外而減少.難道這些司機不知道酒後不應該駕駛的嗎?難道沒有足夠的資訊讓他們認識酒精及毒品對駕駛能力的損害嗎?或者,只要車禍一日未發生在自己身上,他們都把這些警告視作耳邊風；又或者僥倖心理作祟,以為意外總不會發生在自己身上.財主和他的兄弟們可能心想:摩西和先知真嘮叨,比媽媽更「長氣」.現在我們一切都很好,何必那麼認真呢?何必談論這些掃興的事情呢?看來上帝都是說說而已,不一定執行呢!

時候未到,財主和他的兄弟們「馬照跑,舞照跳,樓股一齊炒,那理乞丐門外跑」.但當他們要面對真相的時候,已經太遲了.第二位財主的故事提醒我們不要「扮無知」,很多事情,其實我們知道甚麼是對,甚麼是錯；甚麼重要,甚麼沒有那麼重要；甚麼是自己真正需要的,甚麼是可有可無的.當我們還未走到「太遲」的地步,還未走到追悔不及的日子前,就當回頭,或者調校自己的方向盤,走回正確的路上.

四．    第三位財主18:18

有一個官問耶穌說:「良善的夫子,我該做甚麼事才可以承受永生?」18:19

耶穌對他說:「你為甚麼稱我是良善的?除了　神一位之外,再沒有良善的.18:20

誡命你是曉得的:『不可姦淫；不可殺人；不可偷盜；不可作假見證；當孝敬父母.』」18:21

那人說:「這一切我從小都遵守了.」18:22

耶穌聽見了,就說:「你還缺少一件:要變賣你一切所有的,分給窮人,就必有財寶在天上；你還要來跟從我.」18:23

他聽見這話,就甚憂愁,因為他很富足.18:24

耶穌看見他,就說:「有錢財的人進　神的國是何等的難哪!18:25

駱駝穿過針的眼比財主進　神的國還容易呢!」18:26

聽見的人說:「這樣,誰能得救呢?」18:27

耶穌說:「在人所不能的事,在　神卻能.」18:28

彼得說:「看哪,我們已經撇下自己所有的跟從你了.」18:29

耶穌說:「我實在告訴你們,人為　神的國撇下房屋,或是妻子,弟兄,父母,兒女,18:30

沒有在今世不得百倍,在來世不得永生的.」

第三位出場的財主並不是故事中人,而是前來向耶穌問道的真實人物.若這位富有的官與前兩位耶穌口中的財主相比,他顯得完美無瑕.從他來向耶穌請教「我該做甚麼才可以承受永生」,可以看出他並不像第一位財主般無知,以為物質能夠滿足靈魂的需要；從他向耶穌表示從小就遵守了耶穌提出的誡命來看,他並不像第二位財主一般把摩西和先知的教導視若無睹.

第三位財主的困難,可能就是因為他太過完美.『不可姦淫；不可殺人；不可偷盜；不可作假見證；當孝敬父母』這一切的誡命他都遵守了,可見他克盡己身,奉公守法,而且宗教上虔誠；同時,官運亨通,仕途順利,很可能反映他既有能力又勤力；此外,他很富足,可以看出他應該善於理財,投資有道.對於這位行為,人品,能力,態度都優秀的人,有甚麼是他不能透過自己的努力,運用自己的才能而不成達到的呢?所以,他問耶穌:「請你告訴我,在甚麼事情上我再努力一些,我就能夠得到永生?」

怎知,耶穌的回應卻是:「你不需要再努力再做多一些甚麼事,只需要放下一些東西,與人分享.」耶穌這出乎意料之外的建議,真是難為了這位富有的官.「能力是我幾經艱辛磨練得來的,品格是我忍受試探操練得來的,學位是我捱更抵夜讀回來的,成就是我努力拼搏爭取得到的,錢財是我一點一滴累積下來的．．．一切成果都有著我的血和汗.為甚麼要勤勞的人津貼懶惰的人?為甚麼要奉公守法的人承受胡作非為的人的爛攤子?為甚麼不獎勵好人,反而要好人付出更多?」

第三屆崇基基督教文化節邀請了崇基早期的師兄,香港大學中文系榮休教授陳耀南博士作為主題講員.陳教授憶述,他出生不足一歲,就被交給別人,由養父養母照顧.憑著自己的努力及才華,讀完小學,中學,進入大學,完成博士學位,繼續鑽研中國文化,成為才子教授.教而優則演,曾經在電視,電台主持文教節目,盡顯其努力發揮才華的成果.人到中年,陳教授某次閱讀哥林多前書4章7節,深受觸動:「使你與人不同的是誰呢?你有甚麼不是領受的呢?若是領受的,為何自誇,彷彿不是領受的呢?」陳教授指出聖經的說話如暮鼓晨鐘,發人心省,令他一生受用.

對於富有的官來說,錢財是我的,房屋是我的,妻子是我的,弟兄是我的,父母是我的,兒女是我的.當他堅持所有的一切,都是「我的」的時候,他就把自己與他人分隔開來,在自己與他人之間築起高牆,把自己困在「我的」之內.最後,他只能夠獨自在私家泳池游泳,而不能與大伙兒一齊在沙灘嬉戲；只能夠獨自玩自己的名貴玩具,而不能夠與其他小朋友一齊在遊樂場追逐；只能夠回到自己的豪宅,而沒有穿家過戶的樂趣,沒有守望相助的社群,沒有四海之內皆兄弟的安全.如果富有的官根本不想要這一切,也不享受這一切,永生對他來說可能是痛苦而多於享受.富有的官垂頭喪氣地走,可能因為他終於遇到挫敗,有些東西竟然不能透過自己的本錢而得到.

好人與永生都失之交臂,多麼可惜啊!

五．    最後一位財主19:1

耶穌進了耶利哥,正經過的時候,19:2

有一個人名叫撒該,作稅吏長,是個財主.19:3

他要看看耶穌是怎樣的人；只因人多,他的身量又矮,所以不得看見,19:4

就跑到前頭,爬上桑樹,要看耶穌,因為耶穌必從那裡經過.19:5

耶穌到了那裡,抬頭一看,對他說:「撒該,快下來!今天我必住在你家裡.」19:6

他就急忙下來,歡歡喜喜地接待耶穌.19:7

眾人看見,都私下議論說:「他竟到罪人家裡去住宿.」19:8

撒該站著對主說:「主啊,我把所有的一半給窮人；我若訛詐了誰,就還他四倍.」19:9

耶穌說:「今天救恩到了這家,因為他也是亞伯拉罕的子孫.19:10

人子來,為要尋找,拯救失喪的人.」路加福音記述與耶穌相遇的最後一個財主,也是一個官,他是一位稅吏長.無論猶太人基於民族感情如何看不起這些為殖民地政府收稅的人,稱他們為「稅棍」,但也不能拒絕他們為官的權力,也不能否認他們高強的「搵錢」能力.最後的財主似乎對耶穌非常有興趣.路加描述「他的身量又矮,所以不得看見,就跑到前頭,爬上桑樹,要看耶穌」 (19:4-5) .他如此大動作,竟然純粹是「要看看耶穌是怎樣的人」(19:3).他與耶穌似乎惺惺相惜,耶穌主動提出往他的家去作客.在他們的相遇中,耶穌沒有要求撒該做甚麼,撒該竟然自動提出:「

主啊,我把所有的一半給窮人；我若訛詐了誰,就還他四倍.」(19:8)撒該也沒有向耶穌要求甚麼,耶穌就主動宣告:「

今天救恩到了這家,因為他也是亞伯拉罕的子孫.」(19:9)

如果被第三位同樣為官的財主知道這件事,他一定大呼不公平.為甚麼撒該把所有的一半給窮人,就可以得到救恩,而卻要求我變賣一切所有的分給窮人?耶穌怎可以對撒該和對我提出不同的要求?不公平!絕對是不公平!

如果第三位財主認為耶穌不公平,完全是因為他以「做生意」的思維來理解耶穌與人相遇的性質.他與耶穌的相遇,其實只是一場交易.財主自覺本錢豐厚,於是來向耶穌買永生,要求耶穌開個價.既然財主豪語自己實力雄厚,耶穌開天殺價,財主連落地還錢也不嘗試,挫敗地放棄交易.

最後一位財主爭取與耶穌相遇的機會,並不是來耶穌面前交換甚麼,也不是來多拿一枚勳章,錦上添花,他純粹想認識耶穌,渴望一瞻聖容.超乎他的所想所求,耶穌竟然停下來與他交個朋友.耶穌熱情相待,令他感動不已,受寵若驚.我是誰呢?我何得何能呢?耶穌竟然不嫌棄握我的手,還屈身光臨寒舍.或者是這份恩寵,觸動了最後一位財主的良知,令他醒覺到他所擁有的,其實不只是他個人努力的成果,而是很多人自願或不自願地付出過的,甚至被他犧牲過的.他所擁有的,有他自己的血和汗,其實也有別人的血和汗,更可能有別人被他犧牲時所流下的眼淚和傷痕.與耶穌的相遇,最後一位財主被耶穌的接納與恩慈溶化了第三位財主的執著和自困,開放自己,以致「

今天救恩到了這家」(19:9).

六．    總結聽完路加福音記述的四位財主的故事,我們有甚麼感想呢?第一,讓我們提醒自己不要如第一位財主般「無知」,以物質來解決心靈的需要；第二,小心自己學第二位財主一樣「扮無知」,以致追悔莫及；第三,不要學第三位財主一樣以自己知道的市場運作來看待永生,以自己的本錢來要與耶穌交易；最後,不妨好像最後一位財主一樣,保持對耶穌的濃厚好奇,找他,不為甚麼,只為交個朋友.雖自知不配,但耶穌總不嫌棄,向他開放自己,就有精彩的事情發生.



\newpage

\section{感恩的痲瘋病人}
\label{sec:72}
\textbf{陳曉東博士(普通話) - 感恩的痳瘋病人}
\newline
\newline
\hyperref[sec:71]{< < < PREV SERMON < < <}
~
\hyperref[sec:toc]{[返目錄]}
~
\hyperref[sec:73]{> > > NEXT SERMON > > >}


哥林多後書 5:11-21
\newline
\begin{longtable}{cl}
\hline
\hline
章節 & 經文 (和合本修訂版)\\
\hline
5:11 & \begin{tabularx}{0.7\textwidth}{X} 既然我們知道主是可畏的，就勸導人；但是神是認識我們的，我盼望你們的良心也認識我們。 \end{tabularx} \\ \\ \relax
5:12 & \begin{tabularx}{0.7\textwidth}{X} 我們不是向你們再推薦自己，而是要讓你們有誇耀我們的機會，使你們好面對那憑外貌、不憑內心誇耀的人。 \end{tabularx} \\ \\ \relax
5:13 & \begin{tabularx}{0.7\textwidth}{X} 如果我們癲狂，是為神；如果我們清醒，是為你們。 \end{tabularx} \\ \\ \relax
5:14 & \begin{tabularx}{0.7\textwidth}{X} 原來基督的愛激勵我們；因我們這樣斷定，一人既替眾人死，眾人就都死了； \end{tabularx} \\ \\ \relax
5:15 & \begin{tabularx}{0.7\textwidth}{X} 並且他替眾人死，是叫那些活著的人不再為自己活，乃為替他們死而復活的主活。 \end{tabularx} \\ \\ \relax
5:16 & \begin{tabularx}{0.7\textwidth}{X} 所以，從今以後，我們不再按照人的看法來認識人，縱使我們曾經按照人的看法認識基督，如今卻不再這樣認識他了。 \end{tabularx} \\ \\ \relax
5:17 & \begin{tabularx}{0.7\textwidth}{X} 所以，若有人在基督裡，他就是新造的人，舊事已過，都變成新的了。 \end{tabularx} \\ \\ \relax
5:18 & \begin{tabularx}{0.7\textwidth}{X} 一切都是出於神；他藉著基督使我們與他和好，又將勸人與他和好的使命賜給我們。 \end{tabularx} \\ \\ \relax
5:19 & \begin{tabularx}{0.7\textwidth}{X} 這就是：神在基督裡使世人與自己和好，不將他們的過犯歸到他們身上，並且將這和好的信息託付了我們。 \end{tabularx} \\ \\ \relax
5:20 & \begin{tabularx}{0.7\textwidth}{X} 所以，我們作基督的特使，就好像神藉我們勸你們一般。我們替基督求你們，與神和好吧！ \end{tabularx} \\ \\ \relax
5:21 & \begin{tabularx}{0.7\textwidth}{X} 神使那無罪的，替我們成為罪，好使我們在他裡面成為神的義。 \end{tabularx} \\ \\
[1ex]
\hline
\hline
\end{longtable}
日期:2010年11月14日

主題:感恩的痳瘋病人

証道:陳曉東博士(普通話)

經文:哥林多後書5:11-21

弟兄姊妹平安!

我想我們每個人剛才在聽這樣美麗音符的時候,許多的心靈都會被感動,因為這樣的音符在敲擊著我們心靈中很深很深的地方.在這樣一個很深很深的地方,來呼喚著一個的聲音,我們喜歡這樣的活水泉源嗎?我們喜歡回到這樣一個生命的泉源之地,能夠在那裡可以把我們乾渴的心靈得到滋潤,能夠把我們疲倦的生命得到復甦,能夠讓我們覺得沒有力氣再往前走的時候,因為有這樣一個活水的源泉,我們可以站立起來,能夠往前走嗎?

在今天我們所讀的聖經福音書裡面,耶穌問了一個問題,讓我們心裡面有一種抽痛的感覺.在路加福音17:17:「

耶穌說,潔淨了的不是十個人麼．那九個在哪裡呢」.

今天如果我們再一次思想,默想所讀的福音書,我們所看見的是幾個問題?病痛,醫治,回頭,感謝,而在這一切所發生的事情上面,耶穌卻問一個問題:「那九個在哪裡呢?」那些喝了生命源頭的水,重新得到力量,那些生命因為遇見了耶穌,生命得到一個完全改變以後,這些的人他們喝了水,但是又不回到生命的源頭當中,那九個在哪裡呢?

我最近有機會讀到中國傳統民間裡面的一些籤文,其中一首是媽祖崇拜裡頭的一首籤詩,在媽祖籤詩裡頭共有60首,其中第54首是這樣子說:「病中若得苦心勞,到底完全總未遭,去後不須回頭問,心中事務盡消磨.」這些籤文其實也不是甚麼詩歌,只是一些文字的短句,它的意思是說:一個病了的人還有操勞,實在太辛苦了,人間的事情那有十全十美的?過去的事,你就不要再回頭去問了,你心裡面還有甚麼打算呢?這都取消吧!

我想一個在病中的人,若他抽到這首籤文的話,我想他實在是一種哀傷,無奈,一種非常複雜的心靈就會湧流出來.我想當任何的一個人,他有生病的時候,特別是當我們有機會病得特別嚴重的時候,大概都有過類似這樣的一種心靈,非常複雜,一圈環著一圈的一種的經驗.

疾病就好像我們今天在福音書裡記載的大痳瘋病,或是在人類歷史裡面,許許多多名醫大國手也束手無策的病,不論是以前發生過的黑死病,還是在十九世紀,所發生的肺病,或者是我們在二十世紀早期所發生的癌症,還有我們今天所面對的愛滋病等等的病,都是一種讓我們生命致死的一些疾病,所有一些發生在我們身上的疾病,我們大概第一個反應就是說:為什麼是我?為什麼這件事會這麼不幸遭遇在我的身上?每一種的疾病帶給我們身體上面的痛苦,身體上面的一些捆鎖,說不出來的難過.但是在這身體所承受的苦難以外,其實還有更多的苦難來伴隨著我們.包括我們身,心,靈所受的苦難以外,還有一種罪疚的感覺,一種羞恥的感覺,一種自卑的感覺,一種自責的感覺,渴望自己能夠得到醫治,但是得到的答案,又是讓自己失望的感覺,更加讓我們難過的是,我們常常所患的疾病,不單帶給個人很大的痛苦,也帶給我們自己的家庭很多的痛苦,在這些痛苦裡面,整個的社會如果對這些的事情沒有了解的時候,我們更要承受一種雙重的打擊,包括我們社會對這些病痛都是一種的排斥,不諒解,好像每個人都是污穢的人,每一個這樣的病人,都應該在社會裡面消滅.這是我們所看見的世界裡,常常遭遇的問題所在.

在我們今天所讀的經文裡面,我們看見這樣的一個圖畫:耶穌往耶路撒冷去,但祂經過撒瑪利亞和加利利交界的地帶的時候,聖經說祂進了一個村子,然後有十個長大痲瘋的迎面而來,但只是遠遠地站著,也許這一群大痲瘋病人就像我們所描述那樣,在社會裡面已經被邊緣化了,在社會裡面已經被擠壓到主流之外,這些的人在我們的世界裡面,是不應該存在的,這些人在我們當中會引起我們恐慌,不愉快的感覺,也會引起我們很多很多負面的聯想,恐懼說這樣的病患也會侵擾到我們,所以最好的方式,就是請他們遠遠的不要在我們世界裡存在,我們眼看不見為乾淨.這一群的人也許就是生活在這樣一種的困難裡面,他們看見耶穌,他們迎面而來,遠遠的站著,但是他們發出一個高聲的禱告,大聲的呼喊:耶穌!耶穌!耶穌!夫子!夫子啊!可憐我們吧!

他們所呼喊的夫子,是一個甚麼樣的人呢?他們所面對的夫子,是一個甚麼樣的人呢?在我們所看到路加的經文裡面,他寫得很平實,但是在這樣平實的語言後邊,卻是一種波濤洶湧的生活世界.在我們所看見的耶穌的生命當中,路加很清楚的記載了一件事情:耶穌往耶路撒冷去.這不單止是描述著一個人的行動,在路加的寫作當中,他表達了對耶穌生平,生命的了解,耶穌祂整個的處境,生命,祂就是在前往耶路撒冷的路上,這耶路撒冷等著祂,是怎麼樣一個的等著呢?不錯,是有一些人在夾道歡呼,當祂走過的時候,有驢駒載著祂,有路邊在大聲地呼喊著:和撒那!和撒那!但是我們更知道的一件事情,耶穌走向耶路撒冷的路上,是有可怕的各各他在那裡等著祂,有可怕的十字架在前頭等著祂,有著眾叛親離的結局等著祂... 當祂在客西馬尼園大聲呼喊,祈求生命的苦杯可以挪開的時候,祂所面對的門徒是倒頭而睡,沒有任何感覺.耶穌往耶路撒冷去,祂經過撒瑪利亞和加利利的交界,這是一個生命當中,旅程當中,從一個處境轉到另外一個處境的交界地帶,在那裡祂需要進入一個村子.一個旅行的佈道者,一個波瀾壯闊耶穌運動的創始者,這樣帶著門徒風塵僕僕的,在加利利的地方來宣傳天國福音的先知,祂這個時候需要進到一個村子裡面,為甚麼要進到一個村子裡面?這是人之常情的事情,告訴我們一件事:祂渴了!耶穌渴了!耶穌餓了!祂疲倦了!祂需要休息!祂需要有人來去接待祂,需要有人打開他的房門,家庭說:耶穌請進來!你累了!喝一碗水,吃一塊餅吧.所以耶穌進到一個村子裡面,但是祂沒有得到祂所需要的,祂所面對的卻是一個人的需要──十個長大痲瘋的,十個在生命當中飽受憂患的人,十個在社會主流裡已經被排擠到邊緣的人,十個沒有人願意去接觸的──賤民.他們遠遠的站著,但卻是迎面而來,向一個滿身疲倦的,飢渴的,餓了的耶穌來呼喊:耶穌,可憐我們吧!

我們看到這樣的一段經文,其實好像是一個喜劇的結局,因為路加很快就告訴我們說:耶穌看見了.耶穌雖然有祂自己的疲倦,有祂自己的口渴,有祂自己的飢餓,但是祂看見這樣更多飢餓的人,口渴的人,被排斥的人的時候,祂看見了.祂看見了就說:你們去罷.你們把身體給祭司去看.他們去的時候就潔淨了.

這是一個故事,如果停在這裡的時候,我想路加不必把它記載下來.因為這樣的事情,在耶穌自己的工作裡面,發生太多太多次了.路加繼續的記述這件事,這看似大團圓喜劇的結局,好事呀!耶穌實行了一個大神蹟,又為了祂的名聲加添了一個信用狀,又讓祂的名聲可以到處宣揚了.我想這不是路加要記載的目的,也不是因為這樣的一件事情,耶穌的名聲揚開了就值得把它大大地記載聖經當中.我們看到一件事情就是說,當這樣一件事發生完了,15節開始進入一些最困難的地方:「內中有一個見自己已經好了,就回來大聲歸榮耀與　神．又俯伏在耶穌腳前感謝他」,讓人意外的是「這人是撒瑪利亞人」.耶穌喜歡跟撒瑪利亞人交往,因為撒瑪利亞人是雜種,因為撒瑪利亞人從猶太的傳統來看,是一群已經混雜了的人,是一群已經離開信仰的人,是一群跟摩西純正律法沒有相干的人,這一群人是一群燒不著的餅,一邊有神的律法,另外一邊,他的生命裡充滿了雜質,這樣的人不配來跟猶太人交往的.但是我們看到,耶穌無論在撒瑪利亞的井邊和婦人談道,還是在這裡頭所記載的撒瑪利亞人回頭的故事,所講的是一個撒瑪利亞人回到了耶穌的面前.

15節告訴我們幾件的事情.第一件,這樣的人已經好了,他得到了一個醫治,就像其他九個得到醫治的人一樣,他的生命從一種的被歧視,排斥,從一種自卑,傷痛,低人一等的生命光景裡面,能夠走出來了.哈利路亞!謝謝主!感謝主!但是在這裡,不單單他看見自己已經好了,更重要的是,15節用了一個字眼:他回來.我想我們多少人的經驗裡,當我們在醫院裡面,我們身體能夠得到醫治的時候,有多少人是迫不及待的離開醫院呢?又有多少人會回頭跟照顧我們的醫生,護士來說一聲:謝謝?但是更重要的是甚麼呢?更重要的是有甚麼人願意說:我來向醫生,護士謝謝了,之後還會說:醫生啊!我決定了,我要去讀醫,我要做護士,我要跟你們在一起.我們在我們的經驗裡面,這樣的人實在是很少,很少.但卻在這裡發生了.這裡有一個撒瑪利亞人,他自己好了,聖經說,他就回來了,然後大聲的歸榮耀與神.這裡記載了幾件事情,第一,他得了醫治；第二,他回到耶穌那裡.他不單單回到耶穌那裡,第三,他大聲地讚美神.他把自己生命裡面所受的恩惠,感謝的事,大聲地來去表達出來.第四,他俯伏在耶穌腳前.他能夠在耶穌面前,有一種完全的感謝.俯伏乃表達身體的一種動作,然後他感謝耶穌的醫治.路力在這裡加了一句:「這人是撒瑪利亞人」.這個人與我們不一樣的,與我們所理解的人不一樣,他不是一個猶太人,他不是一個能持守摩西律法到今天還是一個聖潔的人,這個人,他的生命裡充滿了問題.

我們都知道以色列本來就有十二個支派,所以當所羅門的兒子羅波安登位以後,十二個支派就分裂成南北兩個國,北方叫以色列,南方就是猶太,所以北國的國民,在主前721年的時候,被擄往亞述,亞述就把巴比倫地方的百姓遷徙到北國的以色列,代替了以色列人,所以你想像,這七百多年的歷史裡,以色列地區變成了一個很混雜的地方,所以遷來以色列北方的巴比倫人都變成混血,因此撒瑪利亞是他們的首都.列王記下17:33「

他們又懼怕耶和華,又事奉自己的神」.難怪在耶穌的時代,他們被稱為撒瑪利亞人.所以猶太人真是很難願意跟他們在一起,但是在我們今天所看的經文裡,這裡記載他是一個「撒瑪利亞人」.對!就是我們看不起的人,就是我們不想碰的,就是我們聽見了頭都搖一搖.耶穌問的問題是:「

潔淨了的不是十個人麼．那九個在哪裡呢.除了這外族人,再沒有別人回來歸榮耀與　神麼.」所以「

就對那人說,起來走罷．你的信救了你了.」

在我們今天所讀的這樣一個故事裡,弟兄姊妹,我們要問的問題是:這是一個怎麼樣的故事呢?這是一個講醫治的故事嗎?如果講醫治的故事,為甚麼還要再重覆一次?難道在聖經,福音書中,這類醫治的故事不夠多嗎?還需要再加添一個嗎?再讓耶穌的信用再增加一點嗎?再讓祂榮耀的光環再大一點嗎?還是這是榮耀神的故事呢?這麼樣的故事是怎麼樣的榮耀?

當我自己來去默想這段經文的時候,我看到了一點,也許我們可以一起來去默想思考的.在我看來,這裡是一個邀請,這是為每一個念這故事的讀者,無論他是猶太人,還是撒瑪利亞人,是外邦人,還是聖經記載屬以色列的子民... 是一個邀請,這是一個甚麼故事呢?這是一個外族人,如同耶穌所說的,他是個外族人,他回頭,陪伴耶穌的故事.我們記得耶穌走到村子裡面,祂走到那兒,有祂的需要,祂渴了,祂餓了,祂疲倦了,祂希望有人接待祂,祂把祂自己需要工作,需要的服侍做完以後,有一個外族的人,回頭陪伴著耶穌.我們想想看,誰是路加的讀者?也許在這樣的一個讀者群中,有許許多多這樣的一個撒瑪利亞人,就是這樣一群被稱為外邦的人.路加在這裡,鼓勵每一個在生命當中曾經遭遇到困苦,飢餓,疾病,乾旱的一群讀者,一同來回頭,來陪陪耶穌.

經文在這裡有一些沒有講完的話,但是當我們仔細的聆聽,我們聽到了.我們這一群人,會回來陪伴耶穌嗎?這樣的問題帶來進一步的問題... 也許我們心裡面不服氣說:我常常回來呀!我常常回到主的殿中,大聲的唱詩歌；我常常回到主日學裡打開聖經,來閱讀主的話語；我常常回來與我的兄弟姐妹們聚集在一齊,但是在今天的經文內,這樣回來陪伴耶穌的外族人,來去問一個更深的問題:我們在那裡可以遇見耶穌呢?我們在那裡可以去陪伴耶穌呢?這也許引發我們再一次去思考,聖經當中有一些的話語,對我們有很深的提醒.

如果我們不善忘的話,我們大概記得在馬太福音25章,在最後審判日子裡,那位坐在寶座上的王,問這些受審判者,我飢餓的時候,你們在那裡?我乾渴的時候,你們在那裡?我赤身露體的時候,你們在那裡?我在監獄裡的時候,你們在那裡?「這些事你們既作在我這弟兄中一個最小的身上,就是作在我身上了.」「這些事」──當你看見一位弟兄,他嘴巴乾裂,沒有一口水喝的時候,而遞上一杯涼水時,你作在我身上了.這些的事,當你看見,我赤身露體的時候,被人羞辱的時候,被人看不起的時候,你拿起來一塊的毛毯,遮蓋在我身上的時候,你作在我身上.兄弟姐妹們呀!這是耶穌!

我想這樣一個福音的呼召,在歷世歷代中,給許許多多人聽見了.所以我們才在教會歷史裡,有如此多我們稱之為聖人的人,這些的聖人並非他們的品格有怎麼樣的高超,而是說這些的聖人,他們的聖是因為他們聽見了耶穌的呼召.所以我們知道這些閃光的名字,當有人病痛的時候,我們聽見了史懷哲醫生在非洲,與這些貧賤的非洲的病人們住在一起的事情；我們聽見德蘭修女,在街邊那些窮苦的,垂死的乞丐那裡,用毛巾把他們抹得乾乾淨淨,讓他們能夠有尊嚴的離開世界；我們聽見許許多多無名信徒們,他們在身無分文的情況下,收留了多少的孤兒,把他們帶到自己的生活世界與他們一齊生活.但是不夠,不只是這樣子,耶穌在那裡呢?我們去那裡陪伴耶穌呢?在馬太福音那裡告訴我們,那裡有飢餓,那裡乾旱,那裡赤身露體的時候,那裡有政治破壞的監獄的時候──這些事你們既作在我這弟兄中一個最小的身上,就是作在我身上了.

我最近有機會重新溫習,從我書架上看到2006年諾貝爾和平獎得主Yunus(尤努斯)是一位孟加拉人,他最出名的事情是他在孟加拉最赤貧的地方,發展小額貸款,讓許許多多在貧困當中的婦女以及在貧困當中站立不起來的窮人,透過這樣子的小額貸款,重新找到他們生活的方式,重新透過他的工作,能夠獲取他生命的尊嚴.今天2010年,我們知道在遼寧的錦州監獄裡,還有一個人在那裡得到了諾貝爾的和平獎.

我想我們這一群基督徒,在這個世界上常常尋找耶穌,我們去那裡看見耶穌呢?我們在黑暗的世界裡,當看見一根火柴點起來的時候,我們知道耶穌在那裡.我們知道在赤窮的地方,沒有一顆米,一杯水的時候,我們一起的過去,送上了我們這一杯的水,我們一起分享糧食的時候,耶穌在那裡.

主在這個世界的醫治工作,從來沒有停止過,但是,「潔淨了的不是十個人麼．那九個在哪裡呢.」感謝主,「內中有一個見自己已經好了,就回來大聲歸榮耀與　神．又俯伏在耶穌腳前感謝他」.

弟兄姊妹,讓我們一起來做這一個撒瑪利亞人.我們禱告:耶穌!我們今天來到你面前陪你,求你讓我們的心能夠向著你.我們知道,主你也有你疲倦的時候,乾渴的時候,赤身露體被掛在十字架的時候,你也有被放在監獄裡的時候.主啊!在今天我們這一群生活在安逸中的信徒,祈求你的話語能夠震撼我們的心靈,讓我們能夠在一個舒適的小世界當中,知道這個世界有許許多多的呼聲,需要我們來去靠近,因為我們知道我們靠近這些地方的時候,就是靠近你.求主幫助我們,在我們恐懼害怕的時候,加添力量給我們,奉靠耶穌基督聖名祈求,阿們!



\newpage

\section{你的應許之地}
\label{sec:73}
\textbf{沈祖堯教授 - 你的應許之地}
\newline
\newline
\hyperref[sec:72]{< < < PREV SERMON < < <}
~
\hyperref[sec:toc]{[返目錄]}
~
\hyperref[sec:74]{> > > NEXT SERMON > > >}


日期:2010年11月21日

主題:你的應許之地

証道:沈祖堯教授

經文:約書亞記1

中文大學是一個好特別的大學,由於在這所大學中,竟然有一個禮拜堂在內,同時有這麼多的人在每週都前來敬拜上帝,感謝牧師給我這次機會與大家分享,事實上對我本說是很大壓力的,因為我從來沒有學過講道,也不曾受神學訓練,如果你邀請我講腸胃出血,我可以不停的講三小時,至於傳染病也只至少可講二小時之久,可是要我講道實感有些壓力.

今天我不是以校長的身份,而是以一位弟兄的身份與大家分享,過去一年是我人生中,最重要的轉接點,2009年是我50歲生日的日子,今天我透露了我的秘密給大家聽,那時我正想再過多十年就可以退休了,我與太太商量退休後的生活,我們曾經去過澳洲,那裡很美,計劃就在那裡買房屋住下來,對我們來說住下來需要符合三個條件,第一,天氣要好,不要經常掃雪,凍又凍死,熱又熱暈,第二,要有好醫生在附近,由於環境如何美麗,可是遇到心臟病發,要坐直升機才能到達醫院,那就完了,第三,能與好友方便相聚的,要符合以上三點,那麼退休後住那兒呢?我們經常考慮這問題,我亦與上帝商討,我以醫生身份事奉你已廿五載,再做多十年,應該足夠了,而且在”沙士”時期與你一起並肩作戰,期間不單使我增添了許多白髮,還讓我從此心跳不規則,然而最珍貴的給我有機會與祂緊密地走了這段路程,現在應該允許我退休了,讓我舒服地過渡我的餘生.

突然,就在這段時間,學校方面宣布上任校長退休,自覺退休找人代替便成,與我毫無關係,可是身邊開始有人叫我考慮這份職位,當刻我側著頭對他說:「請不要這樣害我.」真的,我不是開玩笑,他們重覆二次對我說,我心內憤怒地想罵他們,為何這麼黑心的對待我?我正準備退休,請不要讓我陷入這情況中,事實我從末考慮過,安排過和計劃過.

回想於2009年8月,有位從加拿大來的牧師在我教會中講道,當時我身在外國開學和教學,他講的就是剛才你們讀的約書亞記,第一章說摩西死了,上帝要找一個人去帶領以色列做祂要做的事,我太太聽了講道後,就把那次的分享藏起來,沒有告訴我,她深知如果讓我知道了的後果是什麼,因為由於裡面的部份內容與我有關,這章講道內容有個特別的地方,當時約書亞聽到神向他呼召,要他帶領祂的僕人過約旦河,並前往神要他們去的地方,「凡你們腳掌所踏之地,我都照著我所應許摩西的話賜給你們了.從曠野,和這利巴嫩,直到伯拉大河,赫人的全地,又到大海日落之處,都要作你們的境界.

你平生的日子,必無一人能在你面前站立得住…」書1:3-4下,用Eugene Peterson的英文譯本寫著”cross to the country I am giving to the people to the Israel,I am giving you every square inches of the lands you set your foot on, strength, courage, you are going to led …., giving it everything you have heart and soul make sure you carry out … that Moses commend you every bit for it.“Eugene Peterson採用了現代的術語講出,你要剛強壯膽”strength  and  courage”,你要帶那些人去承受地士,無人可以在你面前站立得住”every inches of land that they have been promised“,當上帝與約書亞講述以上的說話時,約書亞年已經80歲,他大約40歲跟隨摩西出埃及前往曠野,在沙漠中使以色列人轉來轉去地走了40年,故計算下來他便約80歲,約書亞不是青少年,然而原本猶太人的年曆80歲,即西方計算的50歲,他是位有人生閱歷的人,不過我不知約書亞有沒有想過退休,如果是我一定希望能找另一位人繼承這位置,讓自己可安享晚年,然而神說不是這樣,而是要繼續地帶領祂的子民承受地士；因此,我太太便把這講道分享藏起來,深怕我會與約書亞一樣繼續地事奉,直至後選校長提名名單列出後,我正猶疑當中,這絕對不是我原初的計劃,我的計劃是休息(set back and release),校長這名銜對我並不太大吸引,因為我十分滿意現在的象況,那時我是逸夫書院的院長,不單每天為病人診證,還可以教授課程,當時我在各方面的情況下,都樂在其中,與病人的關係非常融洽,學生對我的評價亦不錯,那麼為何要轉工呢?我非常討厭做行政工作,較喜愛施手術,這帶給我許多滿足感,我有個享福家庭,工作又能幫助別人,夫復何求,只期待著退休的休閒生活,行政工作並不是我所追求的(is not my cup of tea),不是我的喜好,亦不是我的專長,最重要是當我成為校長後,我將會面對我一切不熟習的環境,遠離我的安全地帶 (comfort zone),每個人都有個人的安全地帶,各人不同,有此在家庭,辦公室,職業等,校長的工作並不是我喜悅的,譬如需要和一些不太熟的人扮作很熱情,還要籌款,做醫生從來無需找人取錢,完全與我的日常生活規律有隔離,故我再重申,這不是我的強處,亦不是我的計劃.

本來太太收藏得很好,但奇怪地教會牧師?出奇地把當天講道的硬盤副本寄給我,上帝要你聽的話語,你是逃不過的,聽後,我問自己如果要做這職份時,我是否準備好什麼?我要達成的是什麼?神原來已經為我作出安排了,第一個準備,回想我從2003年”沙士”(Sara)開始,那時我只專注於科研方面,然而讓我開始改變認為教學的重要性,還記得當年有位學生成為”沙士”病人,於病房中住了六星期,他是位五年級醫學院生,本將畢業,因此,他心裡又擔心又憤怒,經常問為何在畢業前讓我患上疾病,使他不能在圖書館裡讀書,他是由於實習時在醫院看個案中受感染,他每天看著自己身體溫度表,就十分生氣,經常抱怨,最終,他寫了一封信和一首詩給我,信中寫著他在這六星期中,反倒學會許多事,在這五年的醫學課程中,從來沒有人教他如何做醫生,但在這六星期裡對醫護人員的觀察,他學懂許多,看到醫護人員所作的事,正是他這五年的醫學院學習中無法學到的一課.病房中許多人也因此信了基督,我們並沒有拿著聖經去傳福音,那時如果有醫生拿著聖經走向你時,我想你一定很害怕,換上我亦同樣懼怕,所以我們沒有這樣做,他的說話讓我不停反思,發現教育並不是單靠課程,考試,功課等,而是以生命影響生命,讓他們明白真正的價值觀是什麼?講述真信仰,真行為是如何實踐出來的,事實上從那時起,我便特別留心學生們的行為,故我曾邀請十位學生參加「逸夫書院」的千人宴,同時亦發受到我都很受「逸夫書院」的學生愛載,慢慢讓我感覺如果我留在中文大學,是應該將我的價值觀,信仰,將我的人生經歷與他們分享,這是神給我的第一個預備,慢慢地將我在醫院的工作,改變去服待學校工作.

第二次準備,在今年3月我和幾位朋友一起前往以色列旅行,遊完加利利海,再去耶路撒冷,當中還有位朋友在約旦河施洗,約旦河很細,它比不上我們的城門河,後經過約旦河前往死海,約旦再往士耳其,整個旅程中,有兩個經歷帶給我很深的印象,於死海的鹹度使生物不能生存,以前從電視曾見到泳客浮在海中,現在自己親歷其境,便立刻跑去海中,這時同樣地有一群小孩子跑到我面前,我並不認識他們,將死海的泥塗在我身上與我玩耍,可是我又不能與他們溝通,他們不懂英語,後發現他們的媽媽全身穿著黑衣,連穿著黑衣一起下水,原來他們是巴勒斯坦人,那些巴勒斯坦的小孩子好像膏抹我一般,並把我扶到海中,我印象中的巴勒斯坦的小孩,是拋酒樽,丟石頭,放火等去危害他人,原來巴勒斯坦的小孩都可以那麼和平可愛；接著我們到達以弗所,位於士耳其中部,在以弗所的舊城區遇上另一群小孩子,但不知為何我與他們好像認識好久,非常熱烈地拍照留念,這兩個經歷讓我感覺到,我與由幾歲至十幾歲的兒童是沒有隔膜的,並帶有親切感,回頭看,我應該選擇兒科最佳,不知為何選擇了老人科,這又是一個預備,讓我明白上帝對我有另一安排.

此外,我也諮詢一些屬靈的長者商量,究竟是否應該考慮這份職位,蔡元雲醫生給了一個很好的提醒,他將八福讀給我聽「使人和睦的人有福了,因為他們必稱為上帝的兒子.」,你是否能作一個使人和睦的人呢?這是很需要的,因為學校很需要這原素,我們經常有許多不和睦的事情發生,聖經說「清心的人有福了,因為他們必得見上帝.」經文在馬太福音5:8,他問我知不知正確的詮釋是什麼?而Eugene Peterson講”You are blessed when you get your inside world, your mind  and  heart put right then you can see God in the outside world.“這裡的意思是如果你處理到你自己內心思想和靈魂,清楚知道自己的動機,清楚知道自己內心的環境,那麼你便能清楚地看見外面世界所發生的事情.這讓我心思,如果我的動機不在於榮譽,出外有司機,薪金和名利等,當動機清晰,便能看出世界的環境,我正處於這情況下,明白上帝給我的呼召,然而為何神再三強調”strength  and  courage”呢?當我答允做這職位時,才深深體會這剛強壯膽的重要性,由於現今校長的工作已不是設計課程,要求建設校舍,或籌款,卻包含許多政治問題,人事問題,這可構成工會與校方,校友與校方,學生與校方的爭持,是需要許多的力量,勇氣去解決的,除此以外,每天都必須跟隨摩西的律法,不可偏離左右,事實上從七月至今,我的屬靈情況比過往30年間的階段還好,因為每天早上我都很安靜小心地查看聖經,很懇切地禱告,因為經上指出虛心的人有福了,虛心是指自己無得救的人,”you are blessed when you are end of the road“,當你覺得自己處理不到的時候,便是你虛心的時候,這些人便能稱為神的兒子,過去的三個月挑戰,不比”沙士”時期弱,然而再次讓我能與上帝走得很親密,同時,我感受到身邊有許多人的支持和幫助,特別是一群弟兄們,當我接受這任命後,律政師黃仁龍弟兄請我到他家飲茶,事實我與他不熟,他說政府部門有一群基督徒的高官,他們定期每週舉辦祈禱會,他忠告我記緊要開設同樣的祈禱會作為支持,否則很容易跌倒的,故我接受了這建議,現時每週四與幾位教授一起祈禱,我亦知教會和學校中有許多人對我的支持.

直到今天我仍然不知流奶與蜜之地是什麼?究竟是在約旦河旁,還是深圳那裡建設校院,或於城門河旁建造一所醫院,如何處理?我完全不知,但上帝說你要帶他們經過這河,前往應許之地,各位我今天與你們分享以上的話語,對你都是很有意義的,想在結束前向大家分享,我的經歷與你們是同樣地真實的,你是否想著退休,或許你有同感已事奉多年,應該休息了,沒有這麼容易,上帝說退休與否,是祂作主,不是你自己,應該想上帝還有什麼差派我們去做的.另外,當我們去做的時候,就像這經文所說”giving everything you have heart and soul”,將你的心,靈魂,感情全然奉上,有時會受傷害,如果你只為了兩餐而做,這便不會有傷害,不過上帝要你全心付出,如你順從上帝就與你同在,不要膽怯,不要惶恐,不要沮喪,雖然短短三個月,我都有沮喪的時候,如何處理這樣的環境?那些問題遠超我能力控制範圍之內,是否自己領受錯了這聲音?神其實沒有吩咐你去?你自己做錯了決定,是否應該請辭呢?立即聖言提醒,以弗所書6:11「要穿戴　上帝所賜的全副軍裝」,英文是”be prepared you up against far more then you can handle on your own“你會面對你無能力處理的事,”take up all you helped you can get”,找所有能協助你的,”every weapon God his issued, so that when it is all over …. you still be on your feet.” 例如古時代的鬥獸場中許多人倒下,只餘下你站著,那些武器是什麼?是真理,正義,平安,信心,和拯救,這些都勝過一些文字,就是這麼簡單,學習如何運用它們,上帝的話語如兩刃的利劍,十分真實的臨在我身,取回過往所認知的聖言聖語,不再是文字,仍是要抓緊它才能生存,每天靠著神的話語,才能生活,這不單是我,你亦如此,所以不要與神說退休,不知祂的心意如何?然而當神使用你時,要剛強壯膽,你今天的生活,是為了將來美好的生活,因此,要常存忠心地事奉上帝,盼望你為我禱告,為中文大學禱告,為香港青年人禱告,為香港禱告,我未明白那應許之地是什麼?不過我開始看到是否中國大地,我開始感覺到不單只是在沙田,或我們的教育,信念,或是醫療技術,或醫療制度,有此事是可以貢獻社會,貢獻國家,在這個過程中我們要忠心服侍上帝,謝謝各位.



\newpage

\section{給你喝一杯生命之水}
\label{sec:74}
\textbf{周小安牧師 - 給你喝一杯生命之水}
\newline
\newline
\hyperref[sec:73]{< < < PREV SERMON < < <}
~
\hyperref[sec:toc]{[返目錄]}
~
\hyperref[sec:75]{> > > NEXT SERMON > > >}


約翰福音 4:13-14
\newline
\begin{longtable}{cl}
\hline
\hline
章節 & 經文 (和合本修訂版)\\
\hline
4:13 & \begin{tabularx}{0.7\textwidth}{X} 耶穌回答，對她說：「凡喝這水的，還要再渴； \end{tabularx} \\ \\ \relax
4:14 & \begin{tabularx}{0.7\textwidth}{X} 誰喝我所賜的水，就永遠不渴。我所賜的水要在他裡面成為泉源，直湧到永生。」 \end{tabularx} \\ \\
[1ex]
\hline
\hline
\end{longtable}
日期:2010年11月28日

主題:給你喝一杯生命之水

証道:周小安牧師

經文:約翰福音4:13-14

周牧師現於溫哥華的Amazing Grace Christian Fellowship 佳恩基督教會當傳道人,1986年取得粒子物理學碩士,兩年後再取得粒子物理學博士學位.同年信主,蒙恩得救.之後進入神學院學習,1998年作傳道人,直到現在.   and \#8212; and \#8212; and \#8212; and \#8212; and \#8212; and \#8212; and \#8212; and \#8212; and \#8212; and \#8212; and \#8212; and \#8212; and \#8212; and \#8212; and \#8212; and \#8212; and \#8212; and \#8212; and \#8212; and \#8212; and \#8212;-我很榮幸, 今天能夠跟崇基學院禮拜堂的會眾一同敬拜神,一同來領受神的話.教會今天給我的題目是《給你唱一杯生命之水》.我想在這裡結合個人的見證,跟約翰福音四章十三至十四節的經文,來分享這個題目.在我們讀這段經文之前,我們先來禱告:「天父上帝,我們滿心感謝讚美你.因為你是唯一配得我們尊崇,敬拜的神.我們今天能夠一同在這裡自由地敬拜你,乃是你的恩典.求你的聖靈,今天在這裡自由地運行.解開你的話語,賜給我們生命的活水.讓我們的心在你裡面得以滿足.奉主耶穌基督的名禱告,阿們.」 讓我們一同來讀約翰福音四章十三和十四節:「

耶穌回答說,凡喝這水的,還要再渴．人若喝我所賜的水就永遠不渴．我所賜的水,要在他裡頭成為泉源,直湧到永生.」讀過約翰福音的人都知道,約翰福音四章十三,十四節這段經文是主耶穌對一個到井邊打水的撒瑪利亞婦人所說的一段話.在這一段話裡,主耶穌談到了兩種水:一種是喝了要再渴的水,另一種是喝了就永不渴的水.並且,祂應許要把那個喝了就永遠不渴的水,賜給凡渴慕尋求祂的人.喝了就永遠不再渴的水,就是生命之水.只有主耶穌能夠,並且願意把這活水賜給人. 讓我們按照次序來看耶穌所說的三句話.第一句話,耶穌說:「

凡喝這水的,還要再渴.」這一句話真是道盡了人間真情.美國作家海明威,因他寫的一篇中篇小說《老人與海》,而獲得諾貝爾文學獎.讓我們用耶穌的這一句話,來解讀一下老人與海.這一篇小說的情節非常簡單,主要描寫一位老漁民一次出海捕魚的經歷.這一個經歷,正是我們人生的縮影.

小說大體分為三個段落:第一段描寫一個老漁民,由於他的運氣不佳,常常捕不到魚,因此窮愁潦倒.第二段描寫這個老漁民,一次獨自的出海捕魚,終於釣到一條大魚.由於魚太大,船太小；這一條大魚帶著魚鉤將老人的船,一直往深海游去.經過兩天兩夜的拼搏,老人經過了千辛萬苦；終於捕到了這一條大魚.第三段描寫這個老漁民,拖著所捕獲的大魚,往岸邊搖船.但一路上,不斷受到鯊魚的攻擊,搶奪,最後,終於精疲力竭.所捕獲的大魚,被鯊魚吞食得只剩下一副骨架子.這老漁民出海捕魚的經歷,正是我們人生的寫照. 在開始的時候,人們往往苦於沒有賺錢的機會,沒有成功的機會.等到賺了錢,成功了,又被工作和事業,緊緊地拖著走.最後創業成功了,又要為守業來拼搏,直到用盡最後的一點力氣.然而,人生到頭來卻是衰老,疾病和死亡,最終不過是水中撈月一場空.另一位得過諾貝爾文學獎的作家,說過一句名言,他說:「人生注定有兩種痛苦:一種是目標達不到的痛苦；另一種是達到了目標,卻發現人仍然不滿足的痛苦.」 請問各位哪一種痛苦更加嚴重呢?我想是後一種痛苦更嚴重.因為第一種痛苦還有希望；第二種痛苦卻使人絕望.耶穌所說的正是這個意思.這喝了還會再渴的水,就是我們所嚮往和追求的世上的目標.它們雖然能夠使我們得到一時的滿足,但是卻不能夠使我們得到真正長久的滿足.你認為這是人間的實情真相嗎?在信主之前,我的人生哲學,就是追求成功.我相信人生的義意就是──成功.成功就有意義,不成功就沒有意義.我在1969年下農村插隊落戶,幾年以後就進到工場,77年恢復高考,我成為文革十年後,第一批的大學生.畢業後留校任教,又考上了研究生,研究生剛畢業,我就出國留學,供讀粒子物理博士學位.我雖然取得了一個又一個的成功,但是總是得不到長久的滿足.我反而覺得生活的憂慮越來越多,生活的擔子越來越重,人生的道路越走越窄.我在認識主耶穌之前,整整奮鬥了十八年,這十八年的經歷,可以成為主耶穌這一句話的註腳:「凡喝這水的,還要再渴」.六四事件改變了一切,我原打算學成回國,回到原校(母校)當教授,但是因為六四不得不在國外留下來.當時,我雖然獲得博士學位,卻找不到適合的工作.所以心中沒有一點成功的快樂,凡而充滿了失落感和對前途的憂慮.我覺得雖然自己經歷了多年的奮鬥,獲得了物理博士；但這個世界並不需要我,我在這個世界是一個多餘的人.一點沒價值.我深深感到人生的虛空,這正是主耶穌所說的:「

凡喝這水的,還要再渴」. 但是感謝神!主耶穌祂不只是停留在此,祂繼續說了第二句話:「

人若喝我所賜的水,就永遠不渴」.這一句話,是人類最偉大,最重要的真理!物質主義者告訴我們,如果我們擁有了財富,房屋,汽車,高薪的工作等等,就一定有幸福.但許多的人生經歷告訴我們,這不過是一句謊言.共產主義告訴我們,只有我們推翻了私有的制度,理想的社會就一定會到來.但二十世紀的歷史告訴我們,這不過是一種烏托邦.但是,耶穌在兩千年前說了這一句話:「

人若喝我所賜的水,就永遠不渴」.兩千年來,無數的人聽見了這句話,接受了這句話,並且親身經歷了這句話.他們真正的喝到了永遠不渴的生命之水. 我自己的生命經歷就是一個見證,過去我不斷追求世上的成功,卻總是得不到長久的滿足.自從我信了耶穌以後,祂使我的心得到了真正的滿足.無數基督徒的生命經歷,也證明了耶穌的話是千真萬確的.在1990年的夏天,受朋友的邀請,我參加了一個家庭查經班.第一次查考聖經,經文是路加福音十二章,耶穌談到不要憂慮的那一段；這一段經文和那天晚上的查經,改變了我對聖經和基督教的看法.我以前雖然從來沒有讀過聖經,卻頑固地認為聖經是一本過時的古書,不值得花時間一讀.我記得我出國留學三年多,每週末,都有人邀請我去參加福音聚會.但每次我都婉言拒絕,我認為那是浪費時間.我覺得時間寶貴,要抓緊學習,做研究.只有一次,有一個假期,我的基督徒朋友邀請我星期六到他家吃飯,然後在他們家裡過了一夜；第二天,他就邀請我們早上去禮拜堂,那是我信主前唯一的一次到禮拜堂. 在禮拜堂裡,優美的聖歌(敬拜)當中,我就在那裡呼呼大睡.當我張開眼睛時,我看見周圍的人都用一種的眼光看著我,那是我唯一一次在信主前,進教堂的經歷.但那天晚上,我覺得主耶穌的話,句句都進到我的心裡,也是我從來沒有想過的.耶穌深刻的指出,人憂慮的根源,是不認識天地的主宰,耶穌稱祂為天父.但人把人生的根基建立在自己的身上,於是便有了無盡的憂慮.耶穌一方面提醒我們,觀察天上的飛鳥,百合花和野地的小草,它們雖沒有人那麼大的能耐,卻比人活得更加逍遙自在.由此,我們相信必定有一位大能的護理者,耶穌稱祂為天父.另一方面,耶穌又向我們指出人的有限,人連自己的身量,就是身高,或者壽命也不能增加一點.怎麼能夠主宰自己的命運?所以,人克服憂慮的關鍵,就是承認自己的有限,並且跟宇宙的主宰建立良好的關係,並且全心的信靠衪.耶穌的話,句句針對人內心深處的憂慮和不安全感.耶穌的獨特之處,在於祂處處以神的眼光看待事物,展現出造物者的性情.

後來我才知道整本聖經,都是以神的眼光看待事物,以神的口吻在說話.但是,我有一點後悔,為什麼這樣遲才接觸到聖經的教會.從此,每星期五的晚上,在牧師家的查經聚會,就像磁貼般吸引了我.當時我居住的地方離牧師的家很遠,但每一週我都期待著查經的夜晚上到來.是什麼東西對我具有這麼大的吸引力呢?是聖經的話語,弟兄姊妹彼此的分享和關懷,是一種雖然不能看見,卻能夠感受到神的同在.在查經班上,我聽到了主耶穌在十字架上拯救罪人的福音,神的愛觸摸了我的心,我就在聖靈的感動下決志信主；並且不久就受洗了.感謝神,自從我成了基督徒以後,耶穌改變了我的人生.回想起來,我信主前後最大的改變,就是對人生的看法,發生了根本的改變. 信主之前,我看人生是一個競技場,有時候是戰場,自我奮鬥,是我的人生哲學；追求成功,是我人生唯一的目標,也是我人生的義意所在.信主後,我回到了天父的家,認識了天父,我的人生觀完全改變了.我發現原來人生是一個恩典!我不是一個多餘的人,沒有內在的價值──必需要靠外在的成功和財富肯定自己的價值... 恰恰相反,因為我是神的孩子,天父愛我,甚至把祂的獨生子耶穌基督賜給我,為我捨身流血.這就是我的價值所在.我所生活的世界,是天父為我們精心設計的家園.我們能夠生活在這個世界,完全是天父的恩典.原來世界的一切,並不是偶然的,也不是人的勞動和努力所創造的,而是天父精心的安排.不僅我們居住的環境,包括空氣,陽光,水份,雨水,生態環境,都是由天父所提供的.就是我們每天生活的大事小事,都在天父的掌管之中.我們生活在天父的恩典之中,不需要為明天憂慮,也不需要追求個人的成功；我們只需要一心來追求天父的旨意,並成全.這樣的生活是何等的平安,喜樂和充實.一方面,世界對我失去了吸引力；另一方面,耶穌對我的吸引力越來越大.我飢渴的心在耶穌那裡得到飽足,浮燥的情緒得到了平安,空虛的人生變得越來越充實.我越來越體會到,主耶穌所說的那句話:「

人若喝我所賜的水,就永遠不渴.」 最後,我們來看主耶穌所說的第三句話:「

我所賜的水,要在他裡頭成為泉源,直湧到永生.」耶穌應許每一個信祂的人說:「

我所賜的水,要在他裡頭成為泉源,直湧到永生.」為什麼喝了主耶穌所賜的水,就永遠不渴呢?原來耶穌所賜的水,不是世上的事物,也不是外面的事物；而是心靈裡面永不會枯竭的泉源.這裡頭的泉源,又是什麼呢?就是主耶穌衪自己的生命,衪透過聖靈,住我們裡面,成為我們裡面活水泉源,直湧到永生.我在教會信主有二十年,我牧會也有十二年,我看到身邊有許許多多的人,信主以後,生命確沒有得到很大的改變.看到世人所遇到種種的問題,他們也同樣遇到；世人有怎樣的回應,他們也大體至上有怎樣的回應.例如他們對生活,工作充滿了憂慮,一旦面臨失業就憂心忡忡,坐立不安.例如他們夫妻的關係,仍然充滿了衝突；和世人一樣,有很多矛盾.這樣的問題,究竟出在什麼地方呢? 原來,我們心中要湧出活水,有一個前題條件,主耶穌不僅要在我們裡面,而且,祂要在我們生命的中心,祂要坐在我們生命的寶座上.所以,只有讓耶穌成為我們內裡活水的泉源；才能夠使這個活水在我們裡面湧流.主耶穌祂必要成為我們生命的主人,佔據我們生命當中的核心,活水才能夠從我們的裡面湧流出來. 我在信主後的一段時間,我並沒有讓主耶穌坐在我的生命的寶座上；我還是自己坐在生命的寶座上,只有在需要的時候才請主耶穌來幫我一下忙.在這樣的時候,我會發現我生活的模式沒有根本的改變.後來,在神的光照下,讓我看到罪的本質──就是自己要作主,自己要作神,自己要坐在生命的寶座上.所以我們要真正的悔改,我們必需要從生命的寶座上下來,讓主耶穌坐在我們生命的寶座上.這樣衪不僅在我的生命裡,而且在我們生命的中心；這樣,生命的活水,才可以在我們的裡面湧流.許多基督徒,雖然決志信主受了洗,但是他跟我初信主時一樣,沒有讓主耶穌來作生命的主宰,仍然穩坐在自己生命的寶座上.這樣的基督徒,心中肯定是湧不出生命的活水來.所以當我在神的光照下,我在神的面前降服自己,我自己退下生命的寶座,讓主耶穌作我生命的主.從此,我把生命的主權交給祂,祂就成為我生命裡頭的活水的泉源,我就看到內裡有活水湧流出來,能夠幫助到周圍的人.自從我把自己生命的主權交給了主耶穌,衪就成為我裡面的活水泉源,甚麼時候我讓祂在我的生命中行使祂的主權；就有活水在我的生命當中湧出來.甚麼時候我的老我冒出來要作主,那麼我看到活水的湧流就慢慢地止息.我意識到,這個老我並不會甘心地,一次性地退下生命的寶座；它常常企圖重新回到它原來的位置.所以,我常常要跟自我來爭戰,我要再一次把自己降服在主耶穌的面前,讓主耶穌重新來坐在我生命的寶座上.我發現我的生命就是這樣一次又一次讓老我下來,讓主耶穌來登上這個生命的寶座,生命就是在一個螺旋式的戰爭的過程.在這裡,我衷心的祝願在這裡每一位都從主耶穌那裡領受生命的活水,願我們每一個人都把自己生命的主權交給主,好讓活水的湧流永不止息,阿們.



\newpage

\section{腓利的旅程}
\label{sec:75}
\textbf{梁元生教授 - 腓利的旅程}
\newline
\newline
\hyperref[sec:74]{< < < PREV SERMON < < <}
~
\hyperref[sec:toc]{[返目錄]}
~
\hyperref[sec:76]{> > > NEXT SERMON > > >}


使徒行傳 8:26-40
\newline
\begin{longtable}{cl}
\hline
\hline
章節 & 經文 (和合本修訂版)\\
\hline
8:26 & \begin{tabularx}{0.7\textwidth}{X} 有主的一個使者對腓利說：「起來！向南走，往那從耶路撒冷下迦薩的路上去。」那路是曠野。 \end{tabularx} \\ \\ \relax
8:27 & \begin{tabularx}{0.7\textwidth}{X} 腓利就起身去了。不料，有一個埃塞俄比亞人，是個有大權的太監，在埃塞俄比亞女王甘大基的手下總管銀庫，他上耶路撒冷去禮拜。 \end{tabularx} \\ \\ \relax
8:28 & \begin{tabularx}{0.7\textwidth}{X} 回程中，他坐在車上，正念著以賽亞先知的書， \end{tabularx} \\ \\ \relax
8:29 & \begin{tabularx}{0.7\textwidth}{X} 聖靈對腓利說：「你去！靠近那車走。」 \end{tabularx} \\ \\ \relax
8:30 & \begin{tabularx}{0.7\textwidth}{X} 腓利就跑到太監那裡，聽見他正在念以賽亞先知的書，就說：「你明白你所念的嗎？」 \end{tabularx} \\ \\ \relax
8:31 & \begin{tabularx}{0.7\textwidth}{X} 他說：「沒有人指教我，怎能明白呢？」於是他請腓利上車，與他同坐。 \end{tabularx} \\ \\ \relax
8:32 & \begin{tabularx}{0.7\textwidth}{X} 他所念的那段經文是這樣：「他像羊被牽去宰殺，又像羔羊在剪毛的人手下無聲，他也是這樣不開口。 \end{tabularx} \\ \\ \relax
8:33 & \begin{tabularx}{0.7\textwidth}{X} 他卑微的時候，得不到公義的審判，誰能述說他的身世？因為他的生命從地上被奪去。」 \end{tabularx} \\ \\ \relax
8:34 & \begin{tabularx}{0.7\textwidth}{X} 太監回答腓利說：「請問，先知說這話是指誰，是指自己，還是指別人呢？」 \end{tabularx} \\ \\ \relax
8:35 & \begin{tabularx}{0.7\textwidth}{X} 腓利就開口，從這段經文開始，對他傳講耶穌的福音。 \end{tabularx} \\ \\ \relax
8:36 & \begin{tabularx}{0.7\textwidth}{X} 二人正沿路往前走，到了有水的地方，太監說：「看哪！這裡有水，有甚麼能阻止我受洗呢？」 \end{tabularx} \\ \\ \relax
8:37 & \begin{tabularx}{0.7\textwidth}{X} 腓利說：『你若是一心相信，就可以。』他回答：『我信耶穌基督是神的兒子。』 \end{tabularx} \\ \\ \relax
8:38 & \begin{tabularx}{0.7\textwidth}{X} 於是他吩咐把車停下來，腓利和太監二人一同下到水裡，腓利就給他施洗。 \end{tabularx} \\ \\ \relax
8:39 & \begin{tabularx}{0.7\textwidth}{X} 他們從水裡上來，主的靈把腓利提了去，太監再也看不見他了，就歡歡喜喜地上路。 \end{tabularx} \\ \\ \relax
8:40 & \begin{tabularx}{0.7\textwidth}{X} 後來有人在亞鎖都遇見腓利；他走遍那地方，在各城宣揚福音，一直到凱撒利亞。 \end{tabularx} \\ \\
[1ex]
\hline
\hline
\end{longtable}
日期:2010年12月5日

主題:腓利的旅程

証道:梁元生教授

經文:使徒行傳8:26-40

今天我想與大家分享一段經文,剛才已經讀過了,這是一篇耳熟能詳的經文,內客提及初期教會的故事,主角也是一位大家相當熟悉的人物──腓利.經文記載在新約中的使徒行傳第八章26-40節,大概內容是說初期教會的使徒腓利,去了南方傳道給一位太監──埃提阿伯的經過,我希望大家一同來思考這段

經文:「

有主的一個使者對腓利說,起來,向南走,往那從耶路撒冷下迦薩的路上去.那路是曠野.腓利就起身去了．不料,有一個埃提阿伯〔即古實見以賽亞十八章一節〕人,是個有大權的太監,在埃提阿伯女王干大基的手下總管銀庫,他上耶路撒冷禮拜去了.現在回來,在車上坐著,念先知以賽亞的書.聖靈對腓利說,你去貼近那車走.腓利就跑到太監那裡,聽見他念先知以賽亞的書,便問他說,你所念的,你明白麼.他說,沒有人指教我,怎能明白呢.於是請腓利上車,與他同坐.他所念的那段經,說,『他像羊被牽到宰殺之地,又像羊羔在剪毛的人手下無聲,他也是這樣不開口.他卑微的時候,人不按公義審判他．〔原文作他的審判被奪去〕誰能述說他的世代,因為他的生命從地上奪去.』太監對腓利說,請問,先知說這話,是指著誰,是指著自己呢,是指著別人呢.腓利就開口從這經上起,對他傳講耶穌.二人正往前走,到了有水的地方．太監說,看哪,這裡有水,我受洗有甚麼妨礙呢.〔有古卷在此有腓利說,你若是一心相信,就可以．他回答說我信耶穌基督是　神的兒子．〕於是吩咐車站住,腓利和太監二人同下水裡去,腓利就給他施洗.從水裡上來,主的靈把腓利提了去,太監也不再見他了,就歡歡喜喜的走路.後來有人在亞鎖都遇見腓利,他走遍那地方,在各城宣傳福音,直到該撒利亞.」使8:26-40

先讓大家思想這住主角腓利,他的英文名是Philip,我相信大部份基督徒對這段的故事都非常熟習,腓利是初期教會其中的一位領袖,掃羅是當中一個逼迫初期教會給的人,腓利和司提反等7個使徒被選舉出來管理教會,這表示他們有管理的能力和屬靈的智慧,這段經文往往被採用作傳道的信息,然而對我而言,它有另一種深刻的意義,你們上兩個星期也曾聽到,約書亞記對沈祖堯校長有深層的意義,我相信每個人讀經的時候,會有幾節經文對你有特別的有意義,很吸引,印象特別深,也常留在腦海中無法忘懷,這段經文對我來說,就是有這種的感受,效果,每次當我翻閱到使徒行傳的時候,它就有跳出來的效果,要讓我注意它,因為腓利的英文名字是Philip,與我的英文名相同,有許多人不知道我叫Philip(因為我在香港很少用),事實上這名字跟著我已有40多年了,這名字是在我去英國的時候改的,這內中也有一個小插曲,當我離開香港之際,還記得滕近輝牧師(前建道神學院院長)曾給了我一個英文名叫Horace,但當我下飛機到了南加州的聖巴巴拉,初期我住在艾里德牧師(建道神學院副院長)家中,他剛巧亦在加州大學進修博士學位,當他聽了我的英文名字,他就說不太適合,因為太英國化了,他便提議採用滕牧師的英文名,這就是Philip了,當你習慣了,便會產生認同感和代入感,這便是我對使徒行傳第八章有特別感受的由來,聖經上的Philip有許多,然而這個腓利就很少出現了.

印象最深刻的一次,是在我完成博士學位的時候,當時我正忙著找尋工作,而在我面前就有兩份的工作,一是留在南加州,另一份是前往新加坡,那時我便需要在當中作出一個選擇,這是我生命中其中一個重要決擇,這節經文便影響著我,好像對我說Philip起來往南走,新加坡在南邊,我在北美遇到一位在馬來西亞來的詩人朋友,我問他從那裡來,他用了兩句詩來回應,寫著:「問君家何處?南斗更南邊.」而新加坡還比馬來西亞位置更偏向南方,當我讀到這裡,便覺得神的使者要往南方走的吩咐,最後在1980年的時候作了前往新加坡的決定,不經不覺已經三十年,這段經文與我實際的生活連接上了,那麼經文的下一節,就使我十分震撼了:「那路是曠野」.

曠野是孤單,炎熱,遇到許多挑戰的地方,生活不是容易過的,我在新加坡起初的幾年,無論在生活環境上都適應不來,環境陌生,人面生疏,當然現在的新加坡與以前已完全的不一樣了,我還記得朋友提醒我,要帶著前往非洲打獵時配戴的帽子,由此,你們可想像那時的狀況,處處都讓我不適應,所以用曠野的生涯是很貼心的.如果你陪著我往下讀這段經文,我相信你也會與我一樣,有這份擔憂,甚至產生恐懼,那就是最後的部份,經文記載腓利完成與太監施洗的任務後,從水裡上來,主的靈把腓利提了去,太監就再也看不到他,基督徒讀到被提的事,感覺好像被接到天家,好像啟示錄裡寫著以諾與神同行三百年後被提,和以利亞同樣地沒有經歷死亡,被提升天,所以當聖經描述到腓利完成工作後被提,我便想是否南行時被主接回天家,當然在我們讀到最後那節,與以上的不同,柳暗花明又一村,腓利與舊約的以諾和以利亞都不同,他被提後再次出現,第八章最後的經文,提到「後來有人在亞鎖都遇見腓利,他走遍那地方,在各城宣傳福音,直到該撒利亞.」腓利沒有因為被提,一去不返,還在不同的地方為主工作,讀到這裡心內較為安心,當你有那代入感時,你便能產生共鳴,這是我個人與這段經文的關係與經歷.

今天我想取這經文與大家分享,當重新閱讀時,希望大家集中注意腓利在這次南下的行動上(因為近來校方經常要求我們寫"行動計劃",所以從行動上就產生不同的觀點角度):

1    起身上路(get up and go South)經文前4節,描述腓利在耶路撒冷事奉,他被推舉出來治理教會,後受到掃羅的迫害,便去了撒瑪利亞傳道,腓利所作的都蒙上帝的祝福,非常順利,他施行神蹟,宣揚福音,獲得許多人擁戴,「

在那城裡,大有歡喜」使8:8,經文還提及有人見到大大驚奇,無論是歡喜,或驚奇,都是描寫腓利在撒瑪利亞的工作,這工作十分成功,然而在這成功的時候,有主的使者對他說:「

起來,向南走,往那從耶路撒冷下迦薩的路上去.那路是曠野.腓利就起身去了」這正是腓利南行旅途的開始,要離開一個熟悉的地方,一個安全的環境,與及一個成功的事業,需要前往一個不知明的地方,一個炎熱的環境,這是一個很勇敢的行動,get up and go South to the road, to the desert road,盧博士可以就這幅圖畫譜寫一首歌曲,主的使者講述十分清晰,那路是曠野,而腓利毫不猶疑地起程,起行是腓利的第一個行動,他不因撒瑪利亞的成功,工作的順利而躊躇滿志,或貪圖安舒,反而,順服於神的吩咐,回應上帝的呼召,請大家留意這行動「起身」──「起身」前的狀態,有的坐著,有的躺臥,甚至不動,處於靜止的狀態,為何不動?或許是由於工作上處理得十分妥善,又或十分疲倦,一點力氣也沒有了,無論是好是壞都是不動的,exhausted很疲憊,complacent自滿都呈現無法向前的情況,「起身」正是針對這情況而行動的,行動者與空言者不同,行動者是一位實事求事的人,空言者普遍是文弱書生或理論家,從這方面看腓利是位坐而起行的人,作為一個行動者都有其分別,有的謀定而後動,有的魯莽衝動的,有的和腓利一樣的,他的行動來自於背後其他的因素,腓利的原動力是來自於聖靈,是主的使者向他囑咐.那他與以上兩種人的分別,在於腓利聽從主的聲音,並願意前往那曠野,他是位有勇氣且帶著承擔感的人,無論道路如何艱辛,神的天使時刻相隨如「

聖靈對腓利說,你去貼近那車走.」從中我們可學習到,基督徒的行動是出於屬靈的帶領(Sprit guide the action),並非一般盲目與衝動,同時也有方向的指引,往南走,雖然他不知其終點在那裡,事實上聖經也沒有清楚交代終點在那裡,只是向著南方前進,讓大家思考自己的方向,不要一定往南走,或長或短,最重要是尋回方向,我們不能要求神給我們有預知的能力,如果我們已知未來將發生的事,那麼人生便沒有驚喜,可是人生需要擁有目標,這才有意義,不是枉然.

2    貼近和聆聽(get close and listen)第二個行動,在腓利南行的時間,我們不太清楚他用了多長的日程,日出日落,然而終於起了變化,如同刻板的人生,毫無變化地生活,或許有一天好像腓利一般,遇到可以改變人生,改變自己,改變別人的人或事,對於腓利來說,向南行這條路程上,看見一台車和在車上坐著的人,這就改變了這人的一生,同時也改變了自己的一生,腓利遇到一位乘車的貴冑,一位權高位重的太監,這是埃提阿伯女王干大基手下的權臣,他掌管國家的銀庫,相等於現代的國家財政首長,我們相信他是位黑色皮膚的猶太人,因為他曾往耶路撒次裡參與崇拜,正在回家的路途上,他坐在車上讀舊約聖經的《以賽亞書》,他們兩位就在曠野中相遇,這次的相遇就改變了雙方的人生,經文記載:「

聖靈對腓利說,你去貼近那車走.」當我們唸了這段經文,是否產生疑問,為何神的靈不吩咐那車停下來呢?或是為何不吩咐腓利攔住那車,讓自己上車呢?為何那麼轉折?貼近車子走,這不是容易做的動作,在崇基讀書的舊生們就比較有同感,那時每一小時才有一班火車,他們會貼著火車跑,需要在10碼內跳上火車,否則就要等一小時才有下一班車,這不單不容易,而且相當危險,雖然腓利追的車不是火車,可是依然相對比人走的快,唸到這裡我不禁地問:「神的用意在什麼?」當然人生中往往出現許多問題,神不一定要回應你的,讓我們繼續唸下去,就會發現,原來要腓利聽到那位太監讀聖經的聲音,如果腓利沒有貼近車子走,就聽不到太監讀經文的聲音了,「車轔轔,馬蕭蕭」是很難聽到車上的聲音,這段經文是《以賽亞書》53章,「貼近」這行動是一種追求明瞭,追求了解的動作,首先腓利要把自己的身體盡量的靠近,和遷就正在移動的車子,這就表示出要努力地付出,目的就是要尋找更深的了解,一切的付出和努力都是出於自動和自願的,聖靈雖然發聲引導,但從「貼近」這裡看到人都需要努力付出來配合,人才能聽到埃提阿伯這位太監讀經的聲音,而這聲音是一個要求幫忙的呼喚,如果不是用心地貼近,很難聽到 the cry of need──這需要幫助的聲音,由於車輛的聲音很容易蓋過心靈呼喚的聲音,需要的是細心耹聽背後心靈的訴求,「貼近」是放下自己去遷就別人,甚至願意付出犧牲的行動,最終獲得很大的回報和果效.

3    上車與分享(get on and get off)第三個行動,腓利往南行遇到一輪車子,並聽到上帝的吩咐,上前貼近車子,且用心的聆聽太監讀經的聲音,我們估計這太監在車上誦讀經文,故當腓利貼近車旁才能聽到他的聲音,腓利第一次開聲說:「

你所念的,你明白麼.他說,沒有人指教我,怎能明白呢.於是請腓利上車...」太監喜歡讀書,雖然不明白當中的意思,可是仍然的繼續閱讀,是位不懈的追求者,而腓利也是位滿腹經綸的讀書人,當他聽了太監誦讀的經文,便能識別它的出處,背景,含意與及寓意,由此可見,腓利是熟讀經書,他對經曲的熟習,和真意的了解也是被上帝選上的因素之一 ── 一個長久的「準備」是經年的「累積」,「上車」前的準備事工──答案很簡單,便是熟讀經書,明白上帝的話語,才能達到隨時的幫助.我們要留心,雖然腓利是順服聖靈的引導,但到了聖經翻開,太監誦讀時,聖靈便沒有發聲,而是由腓利自己去解答太監的提問,這情況我稱之為「自發性的腓利」,是自動的,而且已經有足夠的能力去應付他當前的問題,(今天我只集中探討腓利的行動,而不是他向太監解說經文那部份),他的「上車」帶來了太監的得救,這構成了行動與救恩的關係,這讓我聯想到佛教中「大乘」與「小乘」的爭論,所謂「大乘」是用一輪大車去接載皈依的信眾,普渡眾生,來者不拒,只要他們口中說:「阿彌陀佛」,便能同登極樂,「小乘」的車比較少,容量有限,所以要求上車的人做許多修行,誦經,守誡,做各方面個人的努力,才能收成正果,「大乘」與「小乘」在國內引起許多的爭論,對信眾的要求相差甚大,一是普及性的,另一則是個人的救恩,基督教論救恩也有集體與個人的分別,從腓利「上車」這故事,基督教的救恩,不是在於車的大小,或是上車和下車的分別,腓利上車是為了傳福音,太監坐在車上不是為了得救,而是因為讀經產生的啟發,加上腓利的解說,太監便打開心門,一心相信,故基督教的救恩,不注重「上車」,「下車」,也不在乎於「大車」,「小車」,眾人還是個人,救恩的重點在於每顆心靈,是否願意打開心窗接受耶穌基督,相信救主,本來「下車」以後,還有許多故事可以說的,get off 做什麼呢?是為了施洗,我覺得生命完滿的時候,是一個很重要的禮儀.

現在作一個簡單的總結,腓利的南行之旅,在某方面上可以反映我們現實人生的寫照,他本來是位成功人士,也取得神的悅納,在人生階段上神會希望我們為祂做一些事,無論是大事,或是奇事,祂只期望我們將事情完成,譬如腓利於南行前成功帶領許多人信主,而聖靈叫他去讓一人信主,這證明了成功的定義不在於人數的多少,乃是有否出於主的吩咐,起身向南走,開始走上孤獨難行之路,但腓利終於找到上帝的召命,就是傳福音給膚色不同,文化有異,階級有別的外邦人,且是位太監.腓利聽從聖靈的吩咐,貼近車旁,就聽到讀經的聲音,上帝的經典,讓倆個不認識的人相遇相知地聚在一塊,尋找生命的真諦,兩個生命的結連,從一起翻開聖經來閱讀,最終得到上帝的救恩,並以施洗的行動完成了救贖的過程.當然施洗後的腓利被聖靈提到另一地方,叫亞鎖都這地方,繼續為主作工,腓利是位跟隨聖靈,熟習聖經,工作有成效的人,我希望藉著腓利南行的幾個行動的表現,學習到應有的功課,想一想,「起身」,「上路」,「貼近」,「聆聽」,「上車分享」,「下車施洗」,為主而活,為主作工.



\newpage

\section{耶和華是我的牧者}
\label{sec:76}
\textbf{王建宣先生 - 耶和華是我的牧者}
\newline
\newline
\hyperref[sec:75]{< < < PREV SERMON < < <}
~
\hyperref[sec:toc]{[返目錄]}
~
\hyperref[sec:77]{> > > NEXT SERMON > > >}


詩篇 23:1-6
\newline
\begin{longtable}{cl}
\hline
\hline
章節 & 經文 (和合本修訂版)\\
\hline
23:1 & \begin{tabularx}{0.7\textwidth}{X} 耶和華是我的牧者，我必不致缺乏。 \end{tabularx} \\ \\ \relax
23:2 & \begin{tabularx}{0.7\textwidth}{X} 他使我躺臥在青草地上，領我在可安歇的水邊。 \end{tabularx} \\ \\ \relax
23:3 & \begin{tabularx}{0.7\textwidth}{X} 他使我的靈魂甦醒，為自己的名引導我走義路。 \end{tabularx} \\ \\ \relax
23:4 & \begin{tabularx}{0.7\textwidth}{X} 我雖然行過死蔭的幽谷，也不怕遭害，因為你與我同在；你的杖、你的竿，都安慰我。 \end{tabularx} \\ \\ \relax
23:5 & \begin{tabularx}{0.7\textwidth}{X} 在我敵人面前，你為我擺設筵席；你用油膏了我的頭，使我的福杯滿溢。 \end{tabularx} \\ \\ \relax
23:6 & \begin{tabularx}{0.7\textwidth}{X} 我一生一世必有恩惠慈愛隨著我；我且要住在耶和華的殿中，直到永遠。 \end{tabularx} \\ \\
[1ex]
\hline
\hline
\end{longtable}
日期:2010年12月12日

主題:耶和華是我的牧者

証道:王建宣先生

經文:詩篇23:1-6

大家平安.

我們每個人都跟旁邊的人擊掌一下,說:平安,喜樂!

我們基督徒最大的安慰就是我們享有平安喜樂,我們享有屬地的平安,我們享有屬天的盼望,天地都有了,所以我們有平安和喜樂.今天很高興有機會來跟大家分享上帝的信息.在前不久,我們有一個聚會幾個朋友聊到牧師就像牧羊人的話題.

約翰福音 10:11「

我是好牧人.好牧人為羊捨命. 」

在英文中牧師「pastor」就跟拉丁文「牧人」是同一個詞,所以牧師就是牧人.所以我們常常稱:牧長,在聖經裡面也有多處稱耶穌為牧長,所以牧師其實就是牧養羊的人.所以耶穌說:「

我是好牧人.好牧人為羊捨命.」

聖經上說,人能為人捨命,愛心沒有比這個更大的.和捨命相比,捨錢還算容易,我的錢可以統統都給你,但如果要讓你為我捨命,這就很難了,不得不讓人考慮再三.

很多年前,在台灣有一位國小的教師,帶學生去郊遊遠足.進到山區就遇到一群虎頭蜂,小朋友嚇得不得了,四散哭跑,老師也嚇了一跳,趕緊指揮大家逃離,在萬分危急的時刻他就趕緊脫去自己的上衣,袒胸露背來吸引狂蜂的注意,好讓學生們有時間逃離.最後學生們都安全了,很可惜,老師卻因為中毒太深,最後搶救無效死掉了.這個老師就是孩子的好牧人,因為他愛自己的學生,天天相處在一起,就好像自己的孩子,面對威脅他能捨掉自己的命來救回孩子,這是發生在台灣一件很感人的事.這件事也影響了我們做教育,當老師的人,要愛自己的學生,好像愛自己的孩子,好像耶穌說的:我是好牧人,我為羊捨命.

在耶穌時代的巴勒斯坦,有很多人以牧羊為業,牧羊人身上有三件必備的東西.第一是皮囊用來裝食物和淡水；另外一個東西叫機弦,機弦也就是投石帶,由一個小兜,兩端繫繩組成,使用時,把石子放在兜裡,一根繩纏繞在手腕上,另一根繩子用手捏住,然後快速旋轉兜子,最後鬆開繩子,石子飛出去擊打目標.聖經裡大衛戰勝歌利亞就是用這種機弦打中非利士人頭額的.牧羊人身上最後一樣東西:杖,可以趕羊和驅趕野獸.在牧羊的時候遇到狼,牧羊人就用機弦和杖來保護羊,但牧羊人也時常因為保護羊而喪命狼口.所以耶穌就用他身邊的事給當時的人打比喻:好牧人為羊捨命.

我們通常講道理都會採用身邊的事來做比喻,例如保羅寫給哥林多教會的信中比喻說:地上的帳棚若拆毀了,必得神所造,不是人手所造,在天上永存的房屋.這裡,他就用帳棚來舉例,因為保羅從前是做帳棚的.

我認識一個媽媽,有一次她帶著自己的小寶寶在路上散步,孩子在前面邊玩兒邊跑,突然迎面奔來了一隻大狼狗,小女孩猛得嚇得哭起來,狼狗也愣了一下,媽媽看到狗這麼大又沒拴鏈子,也怕得要死,但是大家想想媽媽會怎麼樣呢?媽媽會拔腿就跑嗎?不會!因為媽媽對孩子來說就像耶穌對我們一樣,好牧人為羊捨命,好媽媽為孩子捨命.所以媽媽這個時候都是把孩子抱起來摟在懷裡,她自己可能會被狗咬到,但是嘴裡還是不斷安慰孩子說「咬不到,寶寶不要怕......

我是好牧人；好牧人為羊捨命.若是雇工,不是牧人,羊也不是他自己的,他看見狼來,就撇下羊逃走；狼抓住羊,趕散了羊群.(約翰福音 10:11-12)

如果,我們是僱人看孩子,遭遇狼狗雇工可能叫,「小寶快跑!」,然後自己也跑--有可能這樣哦,因為雇工是拿工錢的.聖經上說,這種情況他看見狼來,就撇下羊逃走,他是雇工就不顧念羊,我是好牧人,我顧念我的羊,我為羊捨命.

幾十年前政府派我到美國哈佛大學留學,當時學生沒有錢大家都租住在一起.有個住隔壁的同學運氣好找了個兼職--晚上幫人看屋做警衛.我就問他,要是真遇到有壞人來你怎麼辦啊,是跟他打呢,還是報警阿.那個同學半開玩笑的說,「我當然跑啊!我一個學生,打工又沒有拿他多少錢,來的壞人當然是很凶悍的；我打呢,八九成會受傷甚至會死掉,這個家又不是我的......」 為什麼會這樣呢?這就是雇工啊.而我們的主耶穌不是雇工,他是我們的好牧人,我們都是他的羊.

雇工逃走,因他是雇工,並不顧念羊.我是好牧人；我認識我的羊,我的羊也認識我.(約翰福音 10:13-14)

我們現在養的羊主要是用來吃的,但是在耶穌時代的巴勒斯坦牧羊不是用來吃,主要是用來剪毛的,羊今年剪一次毛,明年剪一次毛,後年再剪一次,一直剪,所以牧羊人對羊有好幾年的感情,他就會認得這個羊,雖然一個牧人會牧羊很多羊,但每一隻羊他都可以叫得出名字,就像我們按照自己寵物的特點給他們取名字一樣.這就是好牧人,他認識自己的羊,叫得出名字.好的老師也這樣,他叫得出每個學生的名字；在教會,我們要給牧師一個重擔,就是他要做一個好牧人,就要叫得出所有人的名字--你可以把它當玩笑話,也可以認真去對待--這是一件很難的事,特別像我這樣年齡大,記憶力不好的人.

不過在巴勒斯坦的牧羊人,他們能夠叫得出自己羊的名字；通常他們在山上牧羊,自己走在羊群的前面帶領羊走該走的路,而羊也認得他.牧羊人是各自牧羊自己的一群羊,有時候會混在一起,比如,晚上或者下雨的時候把羊領進一個山洞裡躲避,到了第二天要分開的時候,就會有一個牧羊人走開然後喊自己的羊,他的羊認得他的聲音,就都跟他走,剩下的就是另一群羊.後來有人做實驗,模仿一個牧羊人的聲音去分散兩群羊,結果不但沒有分別出來,反而把所以羊都嚇走了,這就說明了羊也認得自己牧人的聲音.

我是好牧人；好牧人為羊捨命.我認識我的羊,我的羊也認識我--這就是我們跟耶穌的一個非常親密的關係,他為我們釘在十字架上,他真的為我們捨命了.他真的認識我們,我們向他禱告,他都清清楚楚.我們60億人,還包括已經過世的人加起來人數多如星辰,我們向他禱告,這麼多人他都分得清楚,沒有說搞錯了的；他是神,他都認得我們每一個人.

有一節聖經是大家最喜歡的詩篇 23:1:(大衛的詩.)

耶和華是我的牧者,我必不至缺乏.在座的各位弟兄姊妹,大家有感覺自己缺乏什麼嗎?是的,也許你會覺得自己有很多缺乏,比如房價這麼貴我們都買不起.

這次我到香港來,遇到港澳台灣慈善基金會正在頒發今年的愛心獎.一共要頒發10個,台灣和香港各有5個,他們大都在生活裡遇到很多非常的磨難.比如其中有一位獲獎者,她有一個9歲的兒子,孩子先天性殘障--眼睛完全失明,聽力很好但是不會講話,也不能自己行動,媽媽從小照顧他長大.我們本以為這些獲獎遭遇這些磨難可能生活會過得比較淒慘,會向上帝抱怨生活的不公,甚至會自怨自艾地放棄自己的生命.但是讓我們讚歎的是他們讓生活充滿了愛,活得非常積極,因此基金會頒獎給他們.

這位母親用樸實的話說:「我只看到我有的,看不到我沒有.」她看到自己有一個孩子,她看到孩子聽力還不錯,她看到孩子的笑容,她看到孩子慢慢的成長......她看到了很多已經擁有的東西,感覺自己很富足,沒有缺乏什麼.我們身邊有很多人,專門看見自己沒有的東西.自己明明有一所可以安身立命的房屋,但他看不見；能看到的可能是別人的房間比自己的大,像這樣的人將永遠是缺乏的.

耶和華是我的牧者,我必不至缺乏.我們基督徒在世界上有一個優點就是非常滿足,非常感恩.以前我在台灣做財政部長的時候,有一次剛上一輛計程車,司機就很熱情的和我打招呼,我們立刻就攀談起來.他說自己現在每天都很喜樂,各位弟兄姊妹知道他為什麼喜樂嗎?司機接著告訴我說:「現在每天早晨一起來,睜開眼就發現自己還活著,真是太好了!阿們!感謝主!接下來呢,是有早飯可以吃,感謝主!」現在還有好多的人因為窮或者身體的疾病,沒有飯吃或自己吃不進去,這位計程車司機看到的是自己能吃,也有飯吃,而且吃得很香,他就很喜樂的感謝主!阿們!再接下來呢,他還有工作可以做,感謝主!阿們!

各位弟兄姊妹可能說,我還是有很多缺乏,這個沒有,那個沒有.但是如果你像這位計程車司機一樣,眼睛一睜開就阿們,一吃飯就阿們,一工作就阿們,也會變得很喜樂的.再說一件我自己的事:現在我年紀大了,之前每天早晨起來上衛生間要搞很久--揉揉肚子,慢慢的才能解決問題,最近因為改善了飲食和作息習慣居然早晨再上衛生間,一分鐘就解決了.所以我每次沖馬桶的時候,都要:阿們!感謝主!大家見笑了,其實這個新陳代謝對人體很重要,能吃沒有什麼了不起的,如果人不能通就完蛋了,所以大家如果看到這點,以後再沖馬桶的時候也可以像我這樣事事感恩的說:阿們!感謝主!

詩篇23:1

(大衛的詩.)耶和華是我的牧者,我必不至缺乏.詩篇23:2

他使我躺臥在青草地上,領我在可安歇的水邊.

羊最重要的就是吃草,喝水,有這兩樣,羊就很滿足,可以健康長大.這句詩的意思就是我們人所需要的,神都會預備.但是你不能一直看自己沒有的,你這樣看你就沒有滿足的感覺,你就常常住在缺乏當中.越南戰爭時,有個越方士兵被打斷了一條腿,戰爭結束後他就拄著枴杖在一家餐廳當服務員,每天都很喜樂的,在餐廳裡也會唱歌讚美主.而富有的餐廳老闆卻每天悶悶不樂.老闆感覺很奇怪,因為自己這麼富有開了很多餐廳,這個夥計斷一條腿,工資也不高,他為什麼這麼快樂呢?就把他找來問原因.這名服務員說,我很開心是因為我還有一條腿啊,你不知道在戰場上送了命的戰友數不勝數,斷了兩條腿的也很多,我還好啦還有一條腿呢.很多斷了兩條腿的人還很羨慕我啊,看我可以駐這枴杖到處走,可以自己上衛生間,可以自食其力.

所以不樂觀的人就會說:我只有一條腿.而樂觀的人會說:我還有一條腿.那位富有的餐廳老闆因為沒有這種樂觀的生活態度,雖然有很多錢,但是總拿自己和更富有的人相比,就會天天生氣,如果他不改掉這種不樂觀的態度就會一輩子痛苦.

詩篇23:3

他使我的靈魂甦醒,為自己的名引導我走義路.

「走義路」對我們來說就是走正當的路,走好的路,比如夫妻之間,沒有發生外遇就是「走義路」；如果去搞外遇,離婚這些就是走了不好的路.在台灣有個名人出去招妓被人偷拍,搞得沸沸揚揚,妻離子散,這就是走偏了路,沒有走義路.

耶和華是我們的牧者,如果有羊走偏了,牧羊人就叫那隻羊的名字或者用牧羊的手杖把羊拉回來；羊要是偏得更遠牧羊人就用投石帶投擲石塊打到羊的旁邊把他趕回來.有節聖經說:丈夫也當照樣愛妻子,如同愛自己的身子.所以如果愛讀聖經的丈夫失手打了妻子,他就知道自己不對,責備自己並且以後不會再打了.

我們基督徒也會走偏路,即便沒有人告訴你,但是神的話會把你拉回來的.我們的一生就在這條軌道上,走偏然後被拉回來,再走偏再被拉回來,但總要行在正路上.這樣我們就不會做出各樣傷害自己,傷害別人,傷害國家,丟人現眼的事.

詩篇 23:4

我雖然行過死蔭的幽谷,也不怕遭害,因為你與我同在；你的杖,你的竿,都安慰我.

就像我們開始講到的,當小孩子遇到大狼狗會很害怕.如果媽媽跑掉,不和小孩子同在,那孩子一定更加害怕.但是媽媽會抱住孩子說:不要怕,媽媽與你同在.我們也是這樣,當遇到困難就向神禱告,神會說:不要怕,我是你的神,我一定會用公義的右手來扶持你,我必與你同在.聖經講:「我必與你同在」講了很多次,我們有這樣一位創造宇宙的神同在當然很滿足,喜樂.

詩篇23:5

在我敵人面前,你為我擺設筵席；你用油膏了我的頭,使我的福杯滿溢.

「福杯滿溢」就是我們的福氣多到像水盛滿杯子甚至溢出來一樣.很多弟兄姊妹可能說我哪有這麼多福氣啊?我缺乏得很!但是我們看看前面講的那位得獎的媽媽,她這麼困難但是看到的是這樣有,那樣有,於是她福氣就多得裝不下了.

詩篇23:6

我一生一世必有恩惠慈愛隨著我；我且要住在耶和華的殿中,直到永遠.

「

有恩惠慈愛隨著我」在聖經原文的意思是有恩惠慈愛追著我,「追著」和「隨著」有很大的區別,如果有東西追著你,你是跑不掉的.為什麼我們有飯吃阿?因為有飯吃的恩惠慈愛追著.我們也看到現在非洲很多國家的人民沒飯吃,所以追著你的恩惠慈愛你就應該看到.

詩篇23:6

我一生一世必有恩惠慈愛隨著我；我且要住在耶和華的殿中,直到永遠.

最後一句很重要.「耶和華的殿」就是天家,我們有一天要回天家去.所以以弗所書 2:19說:「這樣,你們不再作外人和客旅,是與聖徒同國,是神家裡的人了」.我們信上帝的人相信靈魂就是生命,是看不見的,是一口氣:「生命之氣」,肉體只是裝載靈魂的.有一天靈魂會從肉體裡出來,耶穌在十字架上講:父啊!我將我的靈魂交在你手裡.然後就斷氣了.這就是生命之氣從肉體裡出來,靈魂與身體斷絕關係.我們中國人說「死」,也說「斷氣」,也是一個意思.靈魂是不會死的,我們人嚥氣只是靈魂脫離肉體到天家去得到永生,我們脫離世間的勞苦到耶和華的殿中是多美好啊.

我這次來香港住在旅館,在那裡住幾天--很短,旅館不是我的家,我的家在台北.其實台北的家也不是真正的家,我不過在那裡住幾十年.旅館幾天,台北的家幾十年對我來說都是暫住,只不過是時間長短而已.我們地上只是做客旅,我們信上帝就不在作客旅和外人,是與聖徒同國,是神家裡的人了.神就是天,我們是人,成為神家裡的人,就是「天人合一」--達到中國人所說的最高境界.

我以前做信徒分享時,常開玩笑問大家:有沒有結婚?有沒有找對象?我們台灣馬英九先生還有兩個女兒沒有出嫁.如果你有興趣可以去追,如果追到以後就不得了!你就是台灣地區領導人的家人了!多麼榮耀的事情啊,大家也都將你的事傳為佳話.不過做領導人的家人真的很厲害嗎?我們也看到前領導人陳先生現在一家人都弄得很慘.我們是神的家人,神是創造宇宙萬物的主宰,我們變成祂家裡的人,這才是真正厲害.(詩篇23:6　我且要住在耶和華的殿中,直到永遠.)我們作為基督徒心裡多麼的平安,世界算得了什麼,你有再多的財富,到離開世界的時候都一樣.如果醫生告訴你:還有三個月的壽命.你會因為只剩三個月就可以回天家了而阿們,感謝主嗎?如果大家聽到這個消息不平安的話,說明信仰還不夠強.住三個月和住三十年都是要回家的,想通了就會很快樂.

如果我們週末都來禮拜,每天都讀聖經,每晚都禱告,多讀屬靈的書,我們屬靈的生命會提高,就能夠有這樣平安喜樂的感覺.當我們的靈命足夠高,再聽醫生那樣說,你就會謝謝他告訴你這個好消息,因為還有三個月就能息掉世上勞苦,到天家享受天堂永生之福.你還可以調侃他說:不起醫生,你還要繼續在地上受苦受難,替人看病.我馬上就可以住在耶和華的殿中了,多麼的美好啊.

最後講最重要的一點.我常常在追思禮拜上做15分鐘的「短講」,我都會講這個「天人合一」和「天家」的事.比如有次一個媽媽過世了,我就講:你們不要替這個老姊妹難過,她現在回天家,在那裡非常快樂.你們要替自己難過,你們信上帝沒有啊?如果沒有信將來的結果就會是滅亡,而你媽媽在天家,你們就永遠不能相會了.

我常看到大殮蓋棺之即,兒女痛哭不已,甚至都會衝上去不准蓋上棺蓋,因為沒有蓋好像人還在那裡,一蓋上人就要被埋葬掉,永遠沒有了.我都會勸他們不要難過,如果他們也信上帝,這就只是暫時的別離,早晚大家都會去再相會了,就像追思禮拜常唱的《再相會》.

說到《再相會》,我想起一件逸事.一次聖誕節,我們的主辦人沒有太注意,居然安排唱《再相會》.後來大家都說議論:他到底懂不懂聖經啊,怎麼安排這個歌!聖誕節這麼快樂的日子怎麼唱追思禮拜的歌呢,估計是他聖經讀得少吧,感覺像個外邦人在搞聖誕節聚會似的.後來我轉念一想,也對!聖誕節也可以唱《再相會》啊,因為耶穌來為我們釘十字架,我們就沒有罪了,我們才可以和家人再相會--所以聖誕節唱這首歌也很適合.

還有另一件逸事.一家花店要為銀行喬遷和葬禮分別派送花籃,居然給弄混了,把「喬遷之喜」送到葬禮上,本以為會鬧出亂子結果反而歪打正著:送對了.因為葬禮確實是喬遷之喜--是從地上遷到天上,確實應該道喜.所以我們有一天可以在天家和家人相會,對基督徒來說是極大的恩典,最大的盼望.親人離開我們,真的很難過,但信上帝可以天家再相會,想到這點整個人就會變得非常的愉快,世界也變得那麼的可愛.

最後我們一起禱告:我們在天上的父,謝謝你.今天在你的主日,在你的大日子裡面,我們聚集在這裡來思想你的話語.我們謝謝你,你是我們的好牧人,你為我們捨命,你認識我們每一隻羊,我們也都認識你；我們聽得懂你的聲音,你帶領我們躺臥在青草地上,安歇的水邊,有一天,我們還要住在你的殿中.我們感謝主,你給我們這麼好的盼望.滿心感謝!我們也求主帶領,我們今天還有很多弟兄姊妹他們的父母,孩子還沒相信的,求主你保守他們,求主你開啟他們的心門,讓主耶穌你進入他們的心中,有一天我們不管先離開這個世界還是後離開這個世界,我們都會到你的殿中再相會.求主你保守我們大家,我們謝謝主.這樣的禱告奉主耶穌基督的聖名.阿們.



\newpage

\section{如此我信}
\label{sec:77}
\textbf{羅勝強教授 - 如此我信}
\newline
\newline
\hyperref[sec:76]{< < < PREV SERMON < < <}
~
\hyperref[sec:toc]{[返目錄]}
~
\hyperref[sec:78]{> > > NEXT SERMON > > >}


日期:2010年12月19日

主題:如此我信

証道:羅勝強教授

I. 導言各位弟兄姊妹好!伍牧師請我來講的時候,要我講一講我信主的經過和見證.但是我想一想,自己信主的經過其實是蠻簡單的.我沒有見過什麼極大的神跡,也沒有 什麼很感動人的人生經歷.我的前半生是很暢順和平淡的.我想最值得講的,是我信主的過程中的一點掙扎.它們也許是知識份子信主的一些障礙.希望透過這個分 享,可以為慕道的朋友提供一些新的看法.

追尋人生的意義 -- 我感謝神,不知道為什麼,自小開始我就有一份無名的追求「終極真理」的心.所以,自從我第一個女朋友(今天是我的太太)向我傳福音開始,我一直徘徊在基督教的門外.因為我是一個非常理性的人,我一直覺得基督教的信仰非常不理性.從我第一次聽到福音,我掙扎了八年的時間才最後信主.最後是一位美國的神學博士 生,在我在美國讀博士班的時候引領我信主的.

II. 「理性」與「信心」的衝突我記得當時我遇到最大的困難就是所謂「理性」與「信心」的衝突.其實,我是一個非常理性的人.我相信一個理性的,或者是對自己負責的人,應該是明白,瞭解 某些道理後,才會去接受它的.不然的話,就變成了「盲信」或是「盲從」了.那就和無知的人拜一些毫無根據的「泥菩薩」一樣了.但是偏偏聖經和基督教裡面講 的卻是對一個「看不見,摸不到的上帝」,和一些「將來要發生的事情」的信心.所以,我一直相信一個理性的人是沒有可能會信耶穌的.因為「信心」和「理性」 是衝突.但是,和這個美國的神學博士生每星期五晚上在他家裡的討論,經過大約一年的時間,我慢慢的有了一點亮光.

我慢慢開始想通了一個問題,就是什麼叫做「理性」呢?「理性」其實就是建基於一些「基本假設」的「邏輯分析」.這是什麼意思呢?讓我用一個例子來解釋.如 果你告訴我昨天晚上你見到一個人,突然在你面前消失得無影無蹤.我不但不會相信,我更會說你這個人整天疑神疑鬼,癡人說夢一樣.我之所以說你不理性,主要 是我對「人不會突然出現或消失」這個假設的一份信心.基於這個假設,我邏輯的結論就是這一定是你的幻覺.

III. 為什麼可以對上帝有「信心」?當我明白到這一點以後,我就看到「理性」和「信心」是沒有衝突.相反,理性是基於一些假設的信心而出發的.一個人可以非常的冷靜和理性,但同時對一個摸不 著,看不見的上帝的存在有極大的信心.我們如何能做到這一點呢?很簡單,只要我們謙卑下來就可以了.我覺得要我們認識上帝的無限是很難的,但是要認識我們 自己的有限卻比較簡單.讓我來舉兩個例子解釋一下.

第一,如果我們是生存在兩維的空間裡,當在三維空間中一個球跌下來的時候,我們首先是什麼都看不見.突然,我們看見一個圓圈,然後轉瞬間圓圈不見了.現在 我們生存在三度空間裡,對於第四度空間已經很難掌握了.如果到了第五,第六空間的事情,我們幾乎是無從知曉的.所以有些東西在我們的三維空間突然出現或消 失是可能的.我們的經驗中很少見,是因為我們是活在有限的三維空間裡.所以,只要我們承認自己的無知,我們是可以理性的接受上面的所謂奇跡的.

第二,我們的經驗是時間是單向的往前的,所以事情才有過去和未來.有沒有可能有些事情是沒有過去的呢?今天的宇宙學運用了對相對論和量子力學,科學家告訴 我們只要「質量」夠大的話,時間是有可能停止,甚至是雙向的.意思就是「過去」可以和「將來」同時存在的.或者時間是可以不存在的.所以,只要我們承認自 己的無知,我們是可以理性地接受預知未來,或是上帝不是由其他東西創造出來的種種可能性了.

以上的例子說明了一件事,我們覺得很多事情不合理,本來就是一個極度「自大」的表現.因為這樣一個心態,就好像把我自己看成是一個無所不知,無所不明的 「全知者」一樣.因為凡是我覺得不合理的,或是我經驗中沒有見過的事情,都是錯的,假的.但是對於一個有限的人來說,怎麼可能把所有我們不能理解的事情都 看成不是真的呢?所以,我們不需要明白上帝是如何的無限,我們只要謙卑下來,承認我們是很有限的,我們就可以看得見上帝.所以主耶穌在約翰福音3:12節 說: 「

我對你們說地上的事,你們尚且不信；若說天上的事,如何能信呢?」箴言9:10講:「

敬畏耶和華是智慧的開端,認識至聖者就是聰明.」確實是一句發人心醒 的話.今天我會如何定義一個有智慧的人呢?一個有智慧的人是知道自己是「有限」的,而且是如何有限的人.

小結:所以,建基於一些不同的「基本假設」,一個理性的人可以「相信」凡是我們沒有經驗過的事情都不是真的,因而不接受基督的信仰.但是另外一個理性的人可以謙卑地「相信」我沒有經驗過的事情仍然有可能是真的,因而接受聖經的內容的.

IV. 如何可以對上帝有「信心」?(邏輯分析能力)但是,如果所有自己不能明白的事情都「訴諸自己的無知」的話,那跟迷信有什麼分別呢?如果我可以用「所謂」信心接受聖經的內容,我同樣也可以用同樣的信心 接受佛經,可蘭經,甚至是邪教的經典.到底「信心」與「理性的分析」有沒有一個平衡點的?如果有的話,應該如何的平衡呢?在這裡就帶出了我認為「信心」與 「理性」是沒有衝突的第二個原因.

其實,我覺得某程度上,科學與信仰都是一樣,都是由「邏輯 + 信心」所構成的.以宇宙學中現在幾乎一面倒的勝利的「大爆炸」理論為例,一般的科學家都  「相信」我們現在這個宇宙是經過一場驚天動地的大爆炸而產生的.但是,我們是否「證明」過「大爆炸」呢?自然沒有.因為它只發生過一次.除非是時間倒退回 一百三十七億年以前,不然我們是沒有機會經驗到「大爆炸」理論的.但是為什麼這麼多科學家會「相信」「大爆炸理論」呢?因為他們看見有一定的觀察證據支持 這個理論(例如「恆星光譜的紅移」和宇宙中的「背景微波輻射」等).基於這些「經驗上的觀察」和「邏輯上推論」,他們就有「信心」這個理論是真的.同樣 地,聖經裡面的真理是有很強的歷史學,考古學,預言的驗證,哲學上的合理性,道德倫理上的高超性,和初期殉道者們用他們的血來寫成的見證等等的支持的.但 是你問我有了這些東西,是不是我們已經「證明」了上帝在聖經中講的救贖的工作是真的呢?跟大爆炸理論一樣,我們沒有證明.因為這不是一個可以重複的實驗. 但是,這些不同的「證據」總合起來,是可以給我們的信仰作為一個很強的理性的根基的.

其實,聖經的作者都是非常審慎,認真的人.路加福音的作者路加醫生被譽為極為嚴謹的歷史學家.他在路加福音1章1-4節說:「

提阿非羅大人哪,有好些人提 筆作書,述說在我們中間所成就的事,是照傳道的人從起初親眼看見又傳給我們的. 這些事我既從起頭都詳細考察了,就定意要按著次序寫給你,  使你知道所學之道都是確實的.」保羅在寫給哥林多教會的信內說:「

我當日所領受又傳給你們的:第一,就是基督照聖經所說,為我們的罪死了,而且埋葬了；又 照聖經所說,第三天復活了,並且顯給磯法看,然後顯給十二使徒看；後來一時顯給五百多弟兄看,其中一大半到如今還在,卻也有已經睡了的.」(林前 15:3-6)

其實,科學研究本身很大程度就是一個信心的行為.譬如在十八世紀,所有的人都相信牛頓的經典物理學,而把它看成是舉世皆准的一個範式和金科玉律.這是根據 一些觀察的資料而產生的一份「信心」.但是,到了十九世紀愛恩斯坦的相對論的出現,我們才知道經典物理學是有限制的.相對論為我們帶來一個新的範式 (paradigm).我們首先看見的是有一定的證據支持了牛頓的經典物理學,但是我們還是要有一定的「錯誤的」信心去相信它是放之於四海而皆准的.同樣 的,我們對上帝的信心也可以基於聖經的合理性和種種的客觀證據開始,而慢慢經過自己和上帝的靈的交往的經驗,因而信心日漸成長.

小結:所以,我相信只要我們抱著一個向上帝謙卑的心,人是可以用比較理性的角度來認識上帝的.我們的主耶穌基督沒有反對從理性的角度來認識祂.約翰福音20章 25節記載,當門徒多馬說:「

我非看見他手上的釘痕,用指頭探入那釘痕,又用手探入他的肋旁,我總不信.」主耶穌向他顯現,並沒有責罵他的多疑,並且對他 說:「

伸過你的指頭來,摸我的手；伸出你的手來,探入我的肋旁.不要疑惑,總要信.」在歷史歷代的信徒中,自然有不少追求透過禱告和默想去直接認識上帝的 偉人,但是亦不乏如奧古斯丁(Augustine),阿圭那(Thomas  Aquinas)等透過自己的理性思維,同時藉著上帝的啟示,去認識上帝的人.我深信在這個兩種過程中聖靈都可以帶領我們走向上帝.

V. 結語在過去二十多年傳福音的經驗中,我看到自己在信仰上理性的反思和學習,不單不是浪費,而且更帶領了很多慕道者到主的面前.在過去十多年內,我深深的感到上 帝把我放在香港科技大學和中文大學中,給我一個教授的位置,讓我接觸一群從國內來的研究生,把福音傳給他們,是有特別的心意的.在我的體驗中,從國內來的 中國同學在信仰的追求上,都有一些共通點.比如第一,他們都自小就受著無神論的熏陶.無神的進化論對他們來說是有強而有力的科學作為基礎的真理.這是無可 置疑的.對於一個受了二十幾年無神論教育的人,傳講這個世界是有一個創造主,是極其困難的.第二,他們很多都是來香港讀研究院的.其實無論你是哪一個專 業,做研究的一個共同特點就是要理性,客觀,講求事實,證據.如果要向他們傳講像主耶穌說的「像小孩子一樣的信心」也是極其困難的.神一步一步的把我放在 今天這個崗位,讓我有機會向這些學生傳講福音,可能也正因為我在還沒信主的時候,就正是這樣的一個人.我感謝主給我這個機會服侍這個群體,更感謝主給我今 天有機會站在這裡做一個祂的見證人.今天因為時間關係,我沒有時間講上帝如何的一步一步預備我,直到把我放在今天的崗位上.並且我過去是如何向這些同學朋 友傳講福音的.或者將來有機會再跟大家分享這方面的經歷吧.謝謝.



\newpage

\section{聖誕節第一個問安}
\label{sec:78}
\textbf{伍渭文牧師 - 聖誕節第一個問安}
\newline
\newline
\hyperref[sec:77]{< < < PREV SERMON < < <}
~
\hyperref[sec:toc]{[返目錄]}
~
\hyperref[sec:79]{> > > NEXT SERMON > > >}


路加福音 1:26-45
\newline
\begin{longtable}{cl}
\hline
\hline
章節 & 經文 (和合本修訂版)\\
\hline
1:26 & \begin{tabularx}{0.7\textwidth}{X} 到了第六個月，天使加百列奉神的差遣往加利利的一座城去，這城名叫拿撒勒， \end{tabularx} \\ \\ \relax
1:27 & \begin{tabularx}{0.7\textwidth}{X} 到一個童女那裡，她已經許配大衛家的一個人，名叫約瑟；童女的名字叫馬利亞。 \end{tabularx} \\ \\ \relax
1:28 & \begin{tabularx}{0.7\textwidth}{X} 天使進去，對她說：「蒙大恩的女子，你好，主和你同在！」 \end{tabularx} \\ \\ \relax
1:29 & \begin{tabularx}{0.7\textwidth}{X} 馬利亞因這話就很驚慌，又反覆思考這樣問候是甚麼意思。 \end{tabularx} \\ \\ \relax
1:30 & \begin{tabularx}{0.7\textwidth}{X} 天使對她說：「馬利亞，不要怕，你在神面前已經蒙恩了。 \end{tabularx} \\ \\ \relax
1:31 & \begin{tabularx}{0.7\textwidth}{X} 你要懷孕生子，要給他起名叫耶穌。 \end{tabularx} \\ \\ \relax
1:32 & \begin{tabularx}{0.7\textwidth}{X} 他將要為大，稱為至高者的兒子；主神要把他祖先大衛的王位給他。 \end{tabularx} \\ \\ \relax
1:33 & \begin{tabularx}{0.7\textwidth}{X} 他要作雅各家的王，直到永遠；他的國沒有窮盡。」 \end{tabularx} \\ \\ \relax
1:34 & \begin{tabularx}{0.7\textwidth}{X} 馬利亞對天使說：「我沒有出嫁，怎麼會有這事呢？」 \end{tabularx} \\ \\ \relax
1:35 & \begin{tabularx}{0.7\textwidth}{X} 天使回答她說：「聖靈要臨到你身上；至高者的能力要庇蔭你，因此，那要出生的聖者要稱為神的兒子。 \end{tabularx} \\ \\ \relax
1:36 & \begin{tabularx}{0.7\textwidth}{X} 況且，你的親戚伊利莎白，就是那素來稱為不生育的，在年老的時候也懷了男胎，現在懷孕六個月了。 \end{tabularx} \\ \\ \relax
1:37 & \begin{tabularx}{0.7\textwidth}{X} 因為，出於神的話，沒有一句不帶能力的。」 \end{tabularx} \\ \\ \relax
1:38 & \begin{tabularx}{0.7\textwidth}{X} 馬利亞說：「我是主的使女，願意照你的話實現在我身上。」於是天使離開她去了。 \end{tabularx} \\ \\ \relax
1:39 & \begin{tabularx}{0.7\textwidth}{X} 在那些日子，馬利亞起身，急忙前往山區，來到猶大的一座城， \end{tabularx} \\ \\ \relax
1:40 & \begin{tabularx}{0.7\textwidth}{X} 進了撒迦利亞的家，向伊利莎白問安。 \end{tabularx} \\ \\ \relax
1:41 & \begin{tabularx}{0.7\textwidth}{X} 伊利莎白一聽到馬利亞問安，所懷的胎就在腹裡跳動。伊利莎白被聖靈充滿， \end{tabularx} \\ \\ \relax
1:42 & \begin{tabularx}{0.7\textwidth}{X} 高聲喊著說：「你在婦女中是有福的！你所懷的胎也是有福的！ \end{tabularx} \\ \\ \relax
1:43 & \begin{tabularx}{0.7\textwidth}{X} 我主的母親到我這裡來，為何這事臨到我呢？ \end{tabularx} \\ \\ \relax
1:44 & \begin{tabularx}{0.7\textwidth}{X} 因為你問安的聲音一入我耳，我腹裡的胎就歡喜跳動。 \end{tabularx} \\ \\ \relax
1:45 & \begin{tabularx}{0.7\textwidth}{X} 這相信的女子是有福的！因為主對她所說的話都要應驗。」 \end{tabularx} \\ \\
[1ex]
\hline
\hline
\end{longtable}
日期:2010年12月26日

題目:聖誕節第一個問安 - 你要懷孕生子 (The First Greetings of Christmas-You Should Conceive in Your Womb and Bear a Child)

証道:伍渭文牧師

經文:路加福音1:26-45

有人問愛因斯坦:你最想提問的一個問題是甚麼?(你最想知道的是甚麼?)他說:我很想知道這宇宙是否對我們友善.作為比任何人更了解宇宙的神秘,浩瀚無際,愛因斯坦明白,相對這宇宙,地球太渺小了,人太脆弱了,若這宇宙對人類充滿敵視,我們就有禍了.很多災難電影的主題,如隕石衝擊地球,星球位置移動做成引力改變,引至洪水氾濫或冰河時期重臨,外星人襲地球等等,就表露人類對這問題的不安和焦慮.

今天的福音經課路加1:26-45提到來自外域的天使做訪地球,不是宣告審判,乃宣告福音,是聖誕節第一個的問安.天使對馬利亞問安宣告:「蒙大恩的女子,我問你的安,主和你同在了!……你在上帝面前己經蒙恩了.你要懷孕生子,可以給他起名叫耶穌…….」這是第一次聖誕節的問安.(路1:28-31)

為何是第一個聖誕節的的問安呢?在伯利恆的野地那晚上天使向牧羊人不是第一次問安嗎?何是說這次對馬利亞問安,才是第一次呢?伯利恆聖嬰誕下來是結果,聖靈感孕是原因.台灣總統馬英九不是說過:我是香港人,是香港製造,台灣付交.他是在香港成孕,台灣誕生.在教會的節期中,天使對馬利顯現問安,說她要聖靈感孕,懷有救主,被稱為宣告節(Feast of Annunciation),出生時是聖誕節,宣告節就說明教會認為這是聖誕節第一個問安.不過我們要留意,英文題目問安是眾數,因為懷孕與生子是一件事的兩個階段,就如將臨期和聖誕期,都是預備上帝兒子的來臨.我們經過將臨期的自省等候,就會更能體會聖誕的歡樂.

孕懷期是預備時期,父母在子女成孕在胎中成長的日子,滿心期待著孩子的誕生.誕生時的喜樂,是多個月期待的沉澱和結晶所造成的.好像剛過去的畢業禮,同學領文憑時的喜樂,是多年孕育的成果,若是不用讀書,用金錢買來文憑,就沒有這種喜悅了,有的可能是羞愧和怕人揭發的恐懼.教會的將臨期,就是要我們好好預備迎接耶穌的降生,思考他為何降生,愛慕他的誕生帶來的救恩.

1. 將臨期使我們操練自省

教會用四個星期的將臨期來迎接聖誕,目的使我們好好預備自己接受這信息.從將臨節福音經課安排,我們明白將臨期重點是審判的信息.首主日提及末日的審判,第二,三主日提及施洗約翰,重點都是預備自己,儆醒等候基督降臨.施洗約翰在曠野呼喊:預備主的道,修直他的路!一切山窪都要填滿大大小山崗都要削平!…..凡有血氣的,都要見上帝的救恩!(路二:4-6)

將臨節重點不在紀念耶穌的降生,乃在說明基督的為何降.生,及我們如何預備自己等候他的降生,將臨節是一個等候自省的時間,認識自己的罪,認識上帝的聖潔,律法的嚴厲,接受上帝的拯救,得著福音的安慰.

十二月二十五日的聖誕節本是中古羅馬帝國異教用來紀念太陽的回歸,古羅馬曆法該天為冬至,太陽直照南回歸線,北半球白天最短,以後太陽好像得著新生命,向北移動,北半球白天愈來愈長.教會以公義的太陽-基督的誕生來聖化異教風俗.但現在另一個異教-消費主義-卻把教會的聖誕節異化了.聖誕強調個人消費享樂:聖誕禮物,聖誕老人,聖誕大餐.這一切都是使我們個人感覺良好,我們希望擁有聖誕佳節,不需理會佳音,因這佳音說出律法的嚴厲,要求我們悔改,接受上帝的救贖.我們要一個沒有佳音的佳節,沒有律法的福音.

然而,律法是輕漫不得的,觸犯律法就要承擔嚴重的後果.任意碌卡購物誠然感覺良好,但最終還要埋單付卡數.不理會自己患有糖尿病的警告,飲食不戒口,隨意大快朵頤,最終可能失去雙腳,繼而失去生命,前一段時間一位電視藝人就是這樣離開我們的.在將臨期愈操練等候和期待,我們的聖誕就會愈歡悅.好像結婚,若相方婚前克制自律,保守自己的貞潔,直到結婚那一刻,他們的歡悅也最大.

時代論壇今期引述南華早報David Evans一篇太太的禱文,諷刺名牌的拜物教:我們在置地廣場的Armani,願人都尊你的名為聖.願你的Parda降臨,願你的店在市中心,如同在巴黎.請賜給我們日用的老公金卡,求你讓老公寬恕我們的帳戶餘額,如同我們寬恕他人收取的利息.不叫我們進入三越百貨,救我們脫離永安百貨,因為香奈兒,Gaultier,凡賽斯,Dolce and Gabbana, 全是你的.願美國運通卡與我同在,真到永遠.

這些名牌能否真的滿足我們呢?馬利亞說:但我心尊主為大,我靈以我救主為樂.」巴哈把這尊主頌的第一句話作成千古名曲,作品147:「耶穌,我靈所樂」(Jesus, Joy of Man”s Desiring) .真的,祇有耶穌才是我靈所樂.奧古斯丁數百年後用了有個禱告,重複馬利亞話說:上主阿,在我心靈深處,你為我造了一個空間,而且祇有你才能填滿它.

2. 將臨節使我們學習等候

將臨期就像懷孕時期,要求我們忍耐等候,但等候是現代人失落的藝術.消費心態強調立時的滿足,要排隊等位嗎?我們就即時到另一間餐館消費吧.2004年香港旅遊協會在中環遮打花園,以極速時間,搭建營造白色聖誕小鎮風情,一夜之間,遮打花園「白雪」皚皚.但節期一過,整個立刻出現的聖誕小鎮也突然消失,來去了無痕.

但真正的白色聖誕需要等候.我們等候一年中某個季節,溫度慢慢下降,相對濕度飽和,空氣中水氣凝結為極細粉狀的雪花,翩翩飄在地面.一層又一層,愈來愈厚.有的飄在屋頂,有的掛在樹梢.這是真正的白色聖誕.天氣回暖,雪開始融解,是靜悄悄的.樹身由白轉綠,屋簷滴下冰寒的水,起初點點滴滴,跟著細涓長流.白色聖誕慢慢的臨近,邂逅期待它的人,像情人的相約.分離時漸漸遠去,背影愈來愈小,但總不離開目光的牽引.還有,消失後的回憶.

教會少談將臨期的等待,祇注目聖誕快快來到,也是受到要求立時歡樂的消費主義思想所影響.兩個星期前曾給大家一篇有關將臨期的文章,在其中我提及等候的操練:

等候是接受距離的事實,尊重區隔的空間.等候的人順著日月的推移,季節的運轉,氣候的變化,把距離拉近.兩個渴想見面而分隔的人,因等候而產生期待,因為期待,產生愛慕.愛慕是欲得而未得,欲及而未及的牽引.猶如新郎站在地毯的另一端,用視線牽引著新娘,挽著父親的臂彎,踏著白色的婚紗,向自已迎來.這牽引使雙方的心靈彼此靠近.所以愛不在朝朝暮暮,因為朝朝暮暮容不下距離,也沒有能耐等候；容不下彼此的神聖空間,所以也難產生心靈移向對方的愛慕.

崇拜工具理性,認為科技萬能的社會,很難學會等候.科技的核心精神是掌控.漁農作物和禽畜的自然成長時間太長了,來不及給我們享用消費,人便大量使用化學肥料,飼料,壓縮成長期,管它對人體有害.人把禽畜密密擠在籠內,每天按時往牠們口裡填進飼料,管它患病時會否交叉感染,衍生病毒變種.眾多的魚養在困迫的空間,為了防止皮膚損傷彼此感染,便放進孔雀石綠來殺菌.我們相信利用科技,可以拂逆自然規律,掌控生物的生長節奏,結果造成嚴重的食物中毒問題.

我們作為消費者也有責任.若我們仍然整天想著享受消費,而且追求以最低廉的價錢,謀求穿得,吃得最好,也是一種貪婪,間接造成發展中貧窮國家眾多農民和工人被僱主剝削,也造成眾多漁農禽畜產品,附有對人體有害的殘留農?和催生激素.我們希望藉將臨期有關審判經課的反思,心靈得到覺醒和淨滌.

3等候需要信心

作為耶穌的母親,馬利亞備受尊崇.因為面對否認基督的神性的異端,教會稱馬利亞為上帝之母.她是第一位聽聞福音的信徒,教會以她對上帝愛慕和順服,作為教會貞潔的榜樣.但作為一個信徒,她經歷所有跟隨耶穌的人的考驗.她要忍受被人譏笑,未婚成孕,暗結珠胎.約瑟也感到莫大的壓力,想暗暗的休了她.三十年來她等待機會,想耶穌公開他是上帝獨生子的身份,証明她真的是聖靈感孕而生的.在迦拿婚宴那一個晚上,沒有酒了,她催促耶穌顯出本領,以證真身.”「酒用盡了,耶穌的母親對他說,他們沒有酒了.耶穌說:母親,我與你有東麼相干?我的時候還沒有到.」 (約二3, 4)到耶穌真的顯明他救主的的身份,高懸十架為人贖罪時,又令她肝腸寸斷,正如路加稍後引用先知西面所說:你自己的心也要被刀剌透.(路二35)

施洗約翰也有類似的考驗,作為耶穌的先鋒,他在約但河為耶穌施洗,也看見聖靈彷彿像鴿子降在耶穌身上,也聽到天上傳來的聲音:「你是我的愛子,我喜悅你」(可1:11).但上帝的公義遲遲未有彰顯,因批評希律王和王后希羅底約翰被關在牢中,他打發門徒問耶穌:「那將要來的是你嗎?還是我們等候別人呢?」(路七20)

今天你是否正等候上帝顯出作為,而上帝又沒有按著你的時間表答應你呢,你是否因此感到困惑?馬利亞不愧為信徒的表率,對於上帝的旨意,雖然不明白,也感到為難:「我沒有出嫁,怎麼有這事呢?」(路1:34).她對報喜訊的天使說:「我是主的使女,情願照你的話成就在我的身上.」(路1:38)

在信心的試驗中,上帝並沒有讓馬利亞獨自上路,孤身作戰,他早就安排了親戚撒迦利亞和以利莎白這對夫婦,他們的經歷-過了生育期仍能生育,給馬利亞很強的印證.天使終會離開,誰人會相信她遇見天使呢?更不用說相信天使給她令人難以置信的問安了.馬利亞拜訪以利莎白時,她聽到馬利亞問安,六個月的胎兒就雀躍跳動,又是另一個明證.馬利亞與以利莎白同住了三個月,一起談論聖經,思考信仰,信心得著堅固.在這裡我們看到信徒之間的互相勉勵,何其重要.在這裡,我們也看到上帝給我們一些蒙恩的經歷,不單為了自己,也要為著造就建立其他人.

結論:聖誕第一個問安:你要懷孕生子,是生命的誕生和生命得救贖.這應許不單給馬利亞這第一位信徒,也給每一位跟隨耶穌的人,參與上帝的救贖計劃.想大家都知道我的眼睛最近做了手術,這事令我體會到,隨著歲月的流逝,我們的身體就會自然衰敗,但人的生命是否一生下來,就慢慢死亡消失呢?正如哲學家所說:生命就是死亡的過程.今天聖誕節第一個問安:你要懷孕生子,說明生命可以是多個生命誕生,和多個生命成長的過程.上星期有機會參加一個教會六十週年鑽禧感恩晚宴,五位教會現任,前任主席都在會中,他們相繼發言時,都提及前任對他們靈命和事奉心志的陶造和影響.前人的腳步,孕育了他們,使他們見賢思齊.

今天在你生命中,正孕育著甚麼東西呢?雅各書說:私慾既懷了胎,就生出罪來；罪既長成,就生出死來.(雅15) 孕育期間可以隱藏,但時候到了就會顯出來.私慾尚未生出罪來,乃有機會改弦易轍,罪尚未成長致死,仍可懸崖勒馬.但願我們效法馬利亞:「我心尊主為大,我靈以我救主為樂.」「我是主的使女,情願照你的話成就在我的身上.」願我們為天國多結出生命的果子.「我心尊主為大,我靈以我救主為樂.」阿們.



\newpage

\section{聖言哪裡去了}
\label{sec:79}
\textbf{林豪恩先生 - 聖言哪裡去了?(Where has the word of God gone?)}
\newline
\newline
\hyperref[sec:78]{< < < PREV SERMON < < <}
~
\hyperref[sec:toc]{[返目錄]}
~
\hyperref[sec:80]{> > > NEXT SERMON > > >}


路加福音 3:1-6
\newline
\begin{longtable}{cl}
\hline
\hline
章節 & 經文 (和合本修訂版)\\
\hline
3:1 & \begin{tabularx}{0.7\textwidth}{X} 凱撒提庇留在位第十五年，本丟‧彼拉多作猶太總督，希律作加利利分封的王，他兄弟腓力作以土利亞和特拉可尼地區分封的王，呂撒聶作亞比利尼分封的王， \end{tabularx} \\ \\ \relax
3:2 & \begin{tabularx}{0.7\textwidth}{X} 亞那和該亞法作大祭司。那時，撒迦利亞的兒子約翰在曠野裡，神的話臨到他。 \end{tabularx} \\ \\ \relax
3:3 & \begin{tabularx}{0.7\textwidth}{X} 他就走遍約旦河一帶地方，宣講悔改的洗禮，使罪得赦。 \end{tabularx} \\ \\ \relax
3:4 & \begin{tabularx}{0.7\textwidth}{X} 正如以賽亞先知書上所記的話：「在曠野有聲音呼喊著：預備主的道，修直他的路！ \end{tabularx} \\ \\ \relax
3:5 & \begin{tabularx}{0.7\textwidth}{X} 一切山窪都要填滿；大小山岡都要削平！彎彎曲曲的地方要改為筆直；高高低低的道路要改為平坦！ \end{tabularx} \\ \\ \relax
3:6 & \begin{tabularx}{0.7\textwidth}{X} 凡有血肉之軀的，都要看見神的救恩！」 \end{tabularx} \\ \\
[1ex]
\hline
\hline
\end{longtable}
日期:2011年1月2日

題目:聖言哪裡去了?(Where has the word of God gone?)

証道:林豪恩先生

經文:路加福音3:1-6

一．    經文「

凱撒提庇留在位第十五年,本丟‧彼拉多作猶太巡撫,希律作加利利分封的王,他兄弟腓力作以土利亞和特拉可尼地方分封的王,呂撒聶作亞比利尼分封的王,亞那和該亞法作大祭司.那時,撒迦利亞的兒子約翰在曠野裡,神的話臨到他.他就來到約旦河一帶地方,宣講悔改的洗禮,使罪得赦.正如先知以賽亞書上所記的話,說:在曠野有人聲喊著說:預備主的道,修直他的路!一切山窪都要填滿；大小山岡都要削平!彎彎曲曲的地方要改為正直；高高低低的道路要改為平坦!凡有血氣的,都要見　神的救恩!」路3:1-6

四卷福音書都有記述施洗約翰在曠野傳道的行動,其中以路加福音的版本最特別.路加福音的記述特別之處,在於把施洗約翰放置於具體的歷史場景之中,而其他三卷都沒有這樣的做法.

「

凱撒提庇留在位第十五年,本丟‧彼拉多作猶太巡撫,希律作加利利分封的王,他兄弟腓力作以土利亞和特拉可尼地方分封的王,呂撒聶作亞比利尼分封的王」(3:1).這些當時在位的羅馬王帝及巴勒斯坦一帶的統治者的名字,標示著施洗約翰處身的政治時空.「

亞那和該亞法作大祭司」(3:2).當時在位的大祭司的名字,標示著施洗約翰處身的宗教時空.在大時代中,在政要權貴星光熠熠之中,施洗約翰在哪裡呢?「

那時,撒迦利亞的兒子約翰在曠野裡」(3:2).

路加福音這短短的幾句,帶出無限的意境及想像.第一重想像是地理上的位置:凱撒,巡撫,分封王,這些統治者的官謂,令人聯想到王宮；大祭司,令人聯想到聖殿；與金碧輝煌的王宮,以及莊嚴宏偉的聖殿相比,施洗約翰在哪裡呢?在貧乏荒蕪的曠野.第二重想像是權力及資源上的分別:凱撒,巡撫,分封王,各擁有在其範圖內絕對的權力,支配著管轄地區的人民及財產；大祭司擁有宗教及社會地位,主導著猶太社群的信仰及倫理生活.與這些位高權重的人物對比,施洗約翰是誰呢?施洗約翰只不過是一位普通祭司的兒子,在大漠中吃蝗蟲野物為生的邊緣人物.

最令人感到意外的,是接著的那一句:「

神的話臨到他」(3:2).鏡頭的焦點竟然不是在王宮,不是在聖殿,而在曠野；鎂光燈照射著的竟然不是政治領袖,不是宗教領袖,而是在曠野中無權無勢的小人物.上帝的話並不是臨到王宮的政要貴冑,不是臨到聖殿的神職人員,竟然臨到在曠野的小人物.這真是耐人尋味,又意味深長.

二．    釋經與應用1．    聖言不在王宮雖然世上不少帝王都以「天子」自稱,甚至把自己的喻令定性為「奉天承運」,但上帝的話沒有臨到王宮.為甚麼上帝的話不臨到王宮呢?大衛王臨終前對繼位的兒子所羅門說:「我現在要走世人必走的路.所以,你當剛強,作大丈夫,遵守耶和華－你上帝所吩咐的,照著摩西律法上所寫的行主的道,謹守他的律例,誡命,典章,法度.這樣,你無論做甚麼事,不拘往何處去,盡都亨通.耶和華必成就向我所應許的說話:「

你的子孫若謹慎自己的行為,盡心盡意誠誠實實地行在我面前,就不斷人坐以色列的國位.」」(列王記上2:2-4)要當權者盡心盡意誠誠實實地遵守上帝的吩咐,實在太不可思議了.正所謂「創業容易守業難」,坐上了王位的人,最大的威脅是失去王位,因此,精明的帝王必然盡力運用自己的權力去保住江山,延續王位；昏庸的帝王更不會遵守上帝的吩咐,他們只會濫用權力去滿足一己慾望.當保住王位成為最重要的考慮的時候,「上帝的說話」很有可能被變為保住王位的其中一件工具.

2005年3月19日,美國及其盟軍以伊拉克擁有大殺傷力武器為理由,向伊拉克開戰.當時的美國總統說:「

我們為解放伊拉克人民自由而戰,是一場正義之戰.」當時的伊拉克總統說:「

發動聖戰是每個伊拉克人的責任.攻擊他們,戰勝他們,因為他們是被真主詛咒的邪惡侵略者.我們將會取得勝利,他們將會被消滅.」

若上帝的說話落入了王宮,若上帝的說話落在凱撒,巡撫,分封王手中,在保持王位的優先考慮下,在延續權力的推動之下,由王宮再宣告出來的,會是原裝版本,不折不扣的上帝的說話嗎?

上帝的話不臨到王宮,但王宮卻極力令人相信它擁有上帝的說話,用一切方法塑造自己成為真理及正義的化身.蘇聯共產黨統治時期,他們的黨報稱為「真理報」,要壟斷真理的意圖「畫公仔畫到出腸」.最近,內地一名三聚氰胺毒奶粉受害兒童的父親,聯同其他三聚氰胺毒奶粉受害者的家長,尋求公道,爭取公平,卻被控告「尋釁滋事罪」,由原告被改變為被告,判監兩年半.尋求公道,爭取公平的人被不公道及不公平對待,令不少人都感到錯愕.這很可能反映著公道和公平被當權者壟斷了的後果.

提庇留就任羅馬帝國王位十五年了,還委派了本丟‧彼拉多作猶太巡撫,希律作加利利分封的王,腓力作以土利亞和特拉可尼地方分封的王,呂撒聶作亞比利尼分封的王,不過,上帝的話沒有臨都他們任何一人.他們雖然擁有權力,擁有軍隊,但他們不代表真理,也不代表正義,更不代表上帝.因此,尋求真理,尋找正義,仰望上帝的人,不需要因他們而感到害怕,也不需要因他們的行動而放棄盼望.

2．    聖言不在聖殿上帝的話不臨到王宮,尚算可以理解.上帝的話不臨到聖殿,就令人感到大惑不解.因為按著傳統,摩西在西乃山從上帝領受的兩塊寫上十誡的法版,就是存放在聖殿中,上帝的話不就在聖殿嗎?路加福音指出上帝的話不臨到聖殿,令人驚訝之餘,也令人摸不著頭腦.

何西亞書第4章有一段論及祭司:「

祭司越發增多,就越發得罪我；我必使他們的榮耀變為羞辱.他們吃我民的贖罪祭,滿心願意我民犯罪.將來民如何,祭司也必如何；我必因他們所行的懲罰他們,照他們所做的報應他們.他們吃,卻不得飽；行淫,而不得立後；因為他們離棄耶和華,不遵他的命.」(4:7-10)

放在聖殿裡上帝的話去了哪裡呢?是甚麼把上帝的話語擠掉了呢?

列王記下18章記述一個這樣的片段:「

希西家王十四年,亞述王西拿基立上來攻擊猶大的一切堅固城,將城攻取.猶大王希西家差人往拉吉去見亞述王,說:「我有罪了,求你離開我；凡你罰我的,我必承當.」於是亞述王罰猶大王希西家銀子三百他連得,金子三十他連得.希西家就把耶和華殿裡和王宮府庫裡所有的銀子都給了他. 那時,猶大王希西家將耶和華殿門上的金子和他自己包在柱上的金子都刮下來,給了亞述王.」(18:13-15)由此看來,聖殿存放著的,除了摩西的法版之外,還有很多銀子和金子.聖殿成為了王室的夾萬,祭司看管著國家的金庫.在這樣的政教關係下,祭司關心的是甚麼呢?如果上帝的話臨到祭司,祭司豈能不考慮眾多政治及經濟因素才決定如何宣告,以及宣告甚麼嗎?

聖殿的確存放著上帝的說話,而且以存放著上帝的說話自居.上帝的說話存放在哪裡呢?上帝的說話被收藏在無人能觸摸的至聖所之內,上帝的說話在聖殿很可能已經變成文物,成為歷史遺產.現在,聖殿身兼多職,肩負著眾多的政治及社會任務,因此判斷事情難免有很來政治性及社會性的考慮.道成肉身來到聖殿的門口,聖殿如何處置耶穌呢?「

祭司長和法利賽人聚集公會,說:「這人行好些神蹟,我們怎麼辦呢?若這樣由著他,人人都要信他,羅馬人也要來奪我們的地土和我們的百姓.內中有一個人,名叫該亞法,本年作大祭司,對他們說:「你們不知道甚麼.獨不想一個人替百姓死,免得通國滅亡,就是你們的益處.」」(約翰福音11:47-50)由此可見,聖殿判斷事情的根據,而經不是上帝的說話,而是其他眾多的考慮.

「亞那和該亞法作大祭司」,他們坐鎮在聖殿中,自居為人民與上帝的中間人,但路加福音指出,上帝的話不臨都他們.因此,若果我們看見聖殿腐敗跌倒,看見聖殿同流合污,請不要對上帝失望,也不要以為上帝已經放棄了這個世界,事實上,上帝仍然關心這個世界,上帝仍然對這個世界說話.若果上帝仍有話要說,人們在甚麼地方可以聽得到呢?

3．    聖言在曠野「

那時,撒迦利亞的兒子約翰在曠野裡,神的話臨到他」(3:2)

第一,上帝的說話是施洗約翰為唯一考慮.路加福音接著說:「

他就來到約旦河一帶地方,宣講悔改的洗禮,使罪得赦.」(3:3)上帝的說話臨到這位在曠野的約翰,他立即行動,有話直講.他呼籲悔改,一視同仁.對一般百姓如是,對宗教領袖如是,對達官貴人也如是.他甚至指責希律王說:「

你娶你兄弟的妻子是不合理的.」(馬可福音6:18)就算因此被收在監裡,也沒有違背從上而來的吩咐.

第二,成全耶穌基督是施洗約翰唯一目的.「

預備主的道,修直他的路!一切山窪都要填滿；大小山岡都要削平!彎彎曲曲的地方要改為正直；高高低低的道路要改為平坦!」(3:4-5)施洗約翰就是為此而來到世上,也為此存在於世上.世上有鬥志及有能力的人,都會如施洗約一樣,努力削平人生路上的大小障礙,跨越所有低谷,披荊斬棘,克服困難.人們努力奮鬥,為的是「

預備自己的道,修直自己的路」,富豪們為的是「預備富二代的道,修直富二代的路」,目的都離不開展示自己的能力,製造自己的成就.施洗約翰努力奮鬥,為的卻是「

預備主的道,修直他的路」.施洗約翰把自己的聲望及成就置之度外,有人問施洗約翰:「

拉比,從前同你在約旦河外,你所見證的那位,現在施洗,眾人都往他那裡去了.」(約翰福音3:26)這個把耶穌與約翰比較的問題,並沒有成為約翰的「酸葡萄」,因為約翰清楚自己的身份及使命,他回答說:「

娶新婦的就是新郎；新郎的朋友站著,聽見新郎的聲音就甚喜樂.故此,我這喜樂滿足了.他必興旺,我必衰微.」(約翰福音329-30)「

他必興旺,我必衰微」,這就是施洗約翰看待自己與耶穌基督的關係的態度.

第三,人人得見救恩是洗約翰唯一的關注.「

凡有血氣的,都要見　神的救恩!」(3:6)王宮關注王權的穩定及延續,實屬合理,所以上帝的話不臨到王宮,很可能是一件好事.如果聖殿關注的,也都是自身的地位及利益,神的話不臨到聖殿,實屬無奈,也教人失望.幸好世上仍有心繫萬民的人,仍有視全人類的福祉為首要的人,神的話臨都這樣的人,被忠心及如實地傳講出去的機會應該比只記掛著自己的權力及利益的人高得多.

第四,信仰的傳統是施洗約翰唯一的根源.「

正如先知以賽亞書上所記的話,說:在曠野有人聲喊著說」(3:4)上帝的話沒有臨到王宮,沒有臨到聖殿,而臨到曠野裡的約翰,並不代表約翰或約翰傳講的信息反傳統,也不暗示約翰或其信息離經叛道.不獨是路加福音,其他三卷福音書都同時指出施洗約翰承繼著舊約的傳統,是以賽亞書所預告的人物.其實,聖殿更應該是信仰傳統的承繼者,只可惜聖殿已經把信仰傳統封存起來,如歷史文物一般妥善收藏.在曠野的施洗約翰卻站穩在信仰傳統的基礎上,把信仰傳統活化起來.

三．    總結按照教會年曆,現在已經進入了將臨期,是教會迎接耶穌基督降生的日子.教會當如何迎接成為肉身的「道」呢?教會當以甚麼態度去對待聖言呢?路加福音記述上帝的話臨到施洗約翰的這個段落,給教會一個警惕,一個反省,四個學習.

第一,警惕自己與政權的關係.王宮是會極力說服人們聖言在他們那裡的,但其實聖言不臨在王宮.若教會以為聖言在王宮,教會很容易失足犯錯.好像二次大戰前夕的德國教會,支持希特拉的國家社會主義,甚至表明國家社會主義是正面的基督教作風,在希特拉身上彰顯的是聖靈工作的方式,也是神的旨意.

第二,反省行動的動機.縱使人們合理地假設聖言在聖殿,但原來聖言可以不在聖殿,因此,教會的行動不必然都是正確的,尤其是當涉及自身利益的時候,有機會迷失自己.例如半個世紀之前,當馬丁路德金在在美國爭取廢除種族隔離及種族岐視政策時,南部的教會力證白人比有色人種優越,努力維護種族隔離及岐視政策.就在此時此地,也有聲音提醒者教會界,教會承辦了眾多的:「社福機構與學校,這些全辦得不錯,但我們的整體質素卻不見得提升甚多.我們有很多事工與活動,會議與宴會,忙亂一番,仍是積習難改,文化與生命仍是不變,我們就得反省 :「事工是為了組織的延續,或是事工是為了國度 ?」」

(胡志偉:「事工「機構化」的危機」2010/11/17)

教會向學習施洗約翰學習的四點:第一,當以聖言為唯一考慮；第二,當以耶穌基督為唯一的目的；第三,當以萬民的福祉為唯一的關注；第四,承繼並活化深厚豐富的信仰傳統.



\newpage

\section{及時行善}
\label{sec:80}
\textbf{鄧瑞強博士 - 「及時行善」}
\newline
\newline
\hyperref[sec:79]{< < < PREV SERMON < < <}
~
\hyperref[sec:toc]{[返目錄]}
~
\hyperref[sec:81]{> > > NEXT SERMON > > >}


路加福音 3:7-18
\newline
\begin{longtable}{cl}
\hline
\hline
章節 & 經文 (和合本修訂版)\\
\hline
3:7 & \begin{tabularx}{0.7\textwidth}{X} 約翰對那出來要受他洗的眾人說：「毒蛇的孽種啊，誰指示你們逃避那將要來的憤怒呢？ \end{tabularx} \\ \\ \relax
3:8 & \begin{tabularx}{0.7\textwidth}{X} 你們要結出果子來，和悔改的心相稱。不要自己心裡說：『我們有亞伯拉罕為祖宗。』我告訴你們，神能從這些石頭中給亞伯拉罕興起子孫來。 \end{tabularx} \\ \\ \relax
3:9 & \begin{tabularx}{0.7\textwidth}{X} 現在斧子已經放在樹根上，凡不結好果子的樹就砍下來，丟在火裡。」 \end{tabularx} \\ \\ \relax
3:10 & \begin{tabularx}{0.7\textwidth}{X} 眾人問他：「這樣，我們該做甚麼呢？」 \end{tabularx} \\ \\ \relax
3:11 & \begin{tabularx}{0.7\textwidth}{X} 約翰回答：「有兩件衣裳的，就分給那沒有的；有食物的，也該這樣做。」 \end{tabularx} \\ \\ \relax
3:12 & \begin{tabularx}{0.7\textwidth}{X} 也有稅吏來要受洗，對他說：「老師，我們該做甚麼呢？」 \end{tabularx} \\ \\ \relax
3:13 & \begin{tabularx}{0.7\textwidth}{X} 約翰對他們說：「除了規定的數目，不要多收。」 \end{tabularx} \\ \\ \relax
3:14 & \begin{tabularx}{0.7\textwidth}{X} 也有士兵問他說：「我們該做甚麼呢？」約翰說：「不要勒索任何人，也不要敲詐人；自己有糧餉就該知足。」 \end{tabularx} \\ \\ \relax
3:15 & \begin{tabularx}{0.7\textwidth}{X} 百姓期待基督的來臨；他們心裡猜測，或許約翰是基督。 \end{tabularx} \\ \\ \relax
3:16 & \begin{tabularx}{0.7\textwidth}{X} 約翰對眾人說：「我是用水給你們施洗，但有一位能力比我更大的要來，我就是給他解鞋帶也不配。他要用聖靈與火給你們施洗。 \end{tabularx} \\ \\ \relax
3:17 & \begin{tabularx}{0.7\textwidth}{X} 他手裡拿著簸箕，要揚淨他的穀物，把麥子收在倉裡，把糠用不滅的火燒盡。」 \end{tabularx} \\ \\ \relax
3:18 & \begin{tabularx}{0.7\textwidth}{X} 約翰又用許多別的話勸百姓，向他們傳福音。 \end{tabularx} \\ \\
[1ex]
\hline
\hline
\end{longtable}
日期:2011年1月9日

主題:「及時行善」

証道:鄧瑞強博士

經文:路加福音3:7-18

各位弟兄姊妹,主內平安.按教會聖曆,現在是「將臨期」,是等待救主來臨的日子.這是一個等待的時刻.這是一個對將來有期待的時刻.

哲學家海德格(Heidegger)指出,人有別於萬物,在於人有一種「將來」的意識.人會未雨綢繆,為了將來作出種種計劃.只有人才會想到將來有一天要退休,一隻狗不會想到退休.只有人才會想到要好好退休便要供「強積金」,一隻狗沒有想過要將部分骨頭交給主人.亦只有人才會意識到供了「強積金」也不能好好退休,並為這個尚未來到的將來而煩惱.一隻狗最難明白的,是人為何會為明天而煩惱.

聽說,一個正在自殺的人,若想起明天有些重要事情要做；或者,有人告訴他,明天他約了好朋友,則會減輕他自殺的念頭.只要他意識到明天,意識到將來還有一點東西等著他,則他的生命力量便會回來.一個在森林裡迷路的人,若覺得根本無可能找到出路,便會放棄.願意走下去的人,是那些仍抱著希望的人.最失敗的人,是那些沒有想過明天會成功的人.

在但丁(Dante)的《神曲》(Divine Comedy)裡說,地獄門口寫著:「進入此處的人,要放下所有希望.」沒有了明天,生命才能進入地獄.或者說,人若沒有了明天,便如在地獄中生活.有明天的人,他無論如何還未能進入地獄.所以當我爸爸臨去世時,我跟他說:「天堂再見了」.他點了點頭.他是懷著希望離開的.懷著明天的人,進不了地獄.或者,今日我們可以在這裡互相約定:「將來,我們在天堂相會」.懷著希望的人,進不了地獄.

每個努力活下去的人,都對明天有種種希望,有種種夢想,有種種期待.少時候,每個人都有夢想.我的夢想是能造一台收音機.這個夢想到今日還未實現.今日的青年人,夢想有一天能擁有一個住宅單位,但社會的經濟政策令他們這個夢想很難實現了.我希望這些年青人,不要放棄希望.也渴望那個時常嚷著「建造美好明天」的政府,能有實際政策,將一個應許過的明天帶來.

活著的人,總有希望,總有期待.問題是:你希望什麼?你期待什麼?

有人期待明天遇著愛自己的公主和王子.有人期待明天多一些頭髮,少一些皺紋.有人期待早些賺夠錢以退休.據網上資料顯示,人的十大期望是:不用再與人競爭,富有,自己做老闆,有自由時間,能在家中工作,與愛人在一起,幫助他人成功,擁有一全球企業,環遊世界,活出夢想.然而,卻沒有多少人,期待生命的救贖,期待救主的來臨?

路加福音第三章講到,有一個人在曠野中呼喊.我們一般稱這個人為「施洗約翰」,以別於使徒約翰.「

他就來到約旦河一帶的地方,宣講悔改的洗禮,使罪得赦.正如先知以賽亞書上所記的話,說:在曠野有人聲喊著說:預備主的道,修直他的路!一切山窪都要填滿；大小山岡都要削平!彎彎曲曲的地方要改為正直；高高低低的道路要改為平坦!凡有血氣的,都要見神的救恩!」路3:3-6

施洗約翰大聲疾呼,叫人悔改.在曠野裡,他要預備救主的路.在人世間,他要修築人接近神的大路.他講說神救贖世人的事.

這把聲音,注定是一把曠野的聲音.

他要人關注的,是遙不可及的明天的事,是人的終極命運的事.人縱使總會為明天籌劃,但會想到那麼遙遠嗎?我們現代有最利害的望遠鏡,可探視宇宙的邊緣,但是現代人同時卻是短視得只看到當下自己.我們還重視生命的救贖嗎?我們還期待與主相遇嗎?我們還期待生命進入永恆嗎?

有一次,乘地鐵,閘門一開,一個小朋友直衝空位那裡去,幾乎碰跌旁人.他坐在空位上,招呼他的媽媽來.他媽媽讚他:「叻仔」.若我們也如此教我們的孩子,我們怎能說重視生命的救贖.

當一個人在金融巿場短炒一番,賺到超乎想像的財富時,整個社會讚他:「叻仔」.若我們也如此界定這就是「叻」的時候,我們怎能說期待與主相遇.

現代社會是一個娛樂化的社會,我們將資訊娛樂化,將新聞評論娛樂化,將信仰娛樂化.我們可以想像如下情節,一個大學生跟另一大學生談話.甲問:「你近來讀什麼書?」乙答:「我讀莊子.」甲問:「有什麼心得?」乙便徐徐在背包中取出一本《漫話莊子》,然後煞有介事地說,說:「你看這主角的形象多Cute.」當我們視充滿娛樂的生命為豐盛的生命時,我們怎能說期待生命進入永恆.

施洗約翰在曠野,宣講另一種期待.不是期待比別人「醒目」,不是期待耀目的財富,不是期待感性的歡娛.他呼喚人期待救主,期待生命的救贖.

面對那些對救主毫無期待的人,施洗約翰頗為「勞氣」.路加福音3:7-9說:「

約翰對那出來要受他洗的眾人說:「毒蛇的種類!誰指示你們逃避將來的忿怒呢?你們要結出果子來,與悔改的心相稱.不要自己心裡說:『有亞伯拉罕為我們的祖宗.』我告訴你們,神能從這些石頭中,給亞伯拉罕興起子孫來.現在斧子已經放在樹根上,凡不結好果子的樹就砍下來,丟在火裡.」

施洗約翰為什麼這般「勞氣」呢?就著剛才讀的經文,我想講四點.

因為他面對一群自命期待神的人,但事實上這群人毫無期待神的心.施洗約翰稱他們為「毒蛇的種類」,即是一群離棄神,不忠於神的人.這群人對自己生命已遠離神的實況毫不自覺. 這群人怎會是離棄神,不忠於神的人呢?他們不是主動去到施洗約翰那裡,準備受施洗約翰的洗禮嗎?施洗約翰指出,他們的生命毫不結果,與悔改的心毫不相稱.說是要接近神,但實質是遠離神.這有點似Michael Jackson的MoonWalk,人是向前行,實質是向後退.這也有點似某些政府,說是要越來越民主,實質卻是將爭取民主的人收監.生命的行為與悔改的心不相稱. 為何人會講一套,做一套呢?因為他們有恃無恐.施洗約翰對那些言行不一的人說:「不要自己心裡說:『有亞伯拉罕為我們的祖宗.』我告訴你們,神能從這些石頭中,給亞伯拉罕興起子孫來.」當時那些以色列人,以為自己是亞伯拉罕的子孫,便自動成為「神的子民」,便自動能迎接救主,便自動能獲得救贖的生命.施洗約翰說,這些人的想法錯了.一個人的身分,不等於一個人的生命實質.正如:一個大學生的身分,並不等於他有學問.一個已婚人士的身分,並不等於他重視他的婚姻.一個人有教會的會員身分,不等於他就是屬基督的人.在期待基督這件事上,生命的實質比身分重要. 施洗約翰說:「現在斧子已經放在樹根上,凡不結好果子的樹就砍下來,丟在火裡.」每個人都期待明天,但明天對某些人來說,是美好的明天；對某些人來說,是黑暗的明天.考過會考的人很明白這事,會考放榜之日,對某些人來說,是美好的一天；對另一些人來說,這是審判的日子.面對基督的來臨,情況也一樣.對某些人來說,生命將得到更完滿的救贖；對另一些人來說,這是審判的日子.   我們面對將來,有很多計劃,但有否計劃過將來要面對神?我們當然想過,要面對神,但是否如施洗約翰所講的,我們的期待是假的,我們沒有適切的,配合期待的行動?

我們想像以下情節: 你問一個男人:「你期望什麼?」這個男人將手中的香煙放到嘴邊,然後深深一吸,然後徐徐吐出一圈的煙,然後跟你說:「我期望健康.」我們會認為,他的行為與他的期望不相稱.

你問一個年輕人,你在等什麼?他埋頭埋腦地打手中的遊戲機,手與眼半刻也沒有離開遊戲機,然後跟你說:「我在等女朋友」.然後,她的女朋友到了.他仍如痴如醉地打機.他的女朋友叫他.他有氣沒氣地應一聲:「不要吵,緊張關頭.」他的女朋友走了,他仍懵然不知.我們會認為,他的行為與他的等待不相稱.

又或者,某一日,你怒氣沖沖,只為了一個同事頂撞了你,而你卻整天苦思如何進行報復大計.又或者,某一日,你悶悶不樂,是因為你家中的工人病了,你不是為她的病而憂心,而是覺得便宜了她,她又少做一天的工,而自己又浪費了一天給她的工錢.又或者,某一日,你以權謀私,卻沒有人發現,而整天沾沾自喜.那位背負十架的主耶穌,走到你身邊,問你:「你期待什麼?」你答:「我期待與主相遇,我期待救贖.」這時,你會發現,主耶穌已悄然離去.因為,你的行為與你的期待不相稱.

我們當作什麼,以致我們對救主的期待,是一真的期待.

不同的人走到施洗約翰那裡,問了同一個問題:「我們當作什麼呢?」

路3:10-14:「

眾人問他說:「這樣,我們當作什麼呢?」約翰回答說:「有兩件衣裳的,就分給那沒有的；有食物的,也當這樣行.」又有稅吏來要受洗,問他說:「夫子,我們當做什麼呢?」約翰說:「除了例定的數目,不要多取.」又有兵丁問他說:「我們當做什麼呢?」約翰說:「不要以強暴待人,也不要訛詐人,自己有錢糧就當知足.」

施洗約翰沒有強求人作一些非凡的事,他的要求,都是平凡生活中人人能作的事.眾人問:「

我們當作什麼呢?」

施洗約翰答:「

有兩件衣裳的,就分給那沒有的；有食物的,也當這樣行.」很簡單的要求:「分享」.「擁有」會增加我們的社會身分,但只有「分享」才能改變我們的生命實質.香港有不同機構設有「食物銀行」,讓窮人領取食物以維生.捐助「食物銀行」,可實踐「分享」.

稅吏問:「我們當作什麼呢?」施洗約翰答:「除了例定的數目,不要多取.」

在某意義上,稅吏是做生意的人.他承諾上繳多少稅款,以取得這個職位.然後,他盡辦法回本多賺.施洗約翰要求他們,不要見錢開眼,不要見利忘義,不要為了私利而損害他人.

兵丁問:「我們當做什麼呢?」施洗約翰答:「不要以強暴待人,也不要訛詐人,自己有錢糧就當知足.」對於有權力的人,施洗約翰要求他們,不要以權謀私,不勒索他人.總的而言,施洗約翰要求每個人,善待他人,視他人為人.不懂分享的人,擁有將成為他的累贅.漠視他人的人,將失去自己的人性.看不到別人尊嚴的人,也將看不到基督.

若我們等待基督,等待救贖,我們的確需要有合宜的行動,與我們等待的心相稱.

關於期待基督,施洗約翰說:約翰說:「

我是用水給你們施洗,但有一位能力比我更大的要來,我就是給他解鞋帶也不配.他要用聖靈與火給你們施洗.他手裡拿著簸箕,要揚淨他的場,把麥子收在倉裡,把糠用不滅的火燒盡了.」路3:16-17

基督的來臨,像火的來臨.他要驗出我們的生命,是否真實的等待他.近日香港出現假金,在烈火下,假金便現形了.但真金不怕紅爐火,在烈火下,真金越發顯出其是真金.基督的到來,要顯出誰活出了永恆的生命.

面對絕症,大部分的人都願意花盡一切去醫治,因為他明白,生命才是最重要的.他同時會明白,什麼才是最寶貴的,他會重視家庭多於工作,重視生命多於金錢,重視友誼多於成功.同樣,基督來臨的日子,也像是人生的大限,他會讓我們明白,什麼才是最寶貴的.原來,分享強於擁有,尊重人強於征服人,服事人強於奴役人.若我們明白了,又「及時行善」,則我們便是真正等待基督來臨的人.

但願榮耀歸於聖父,聖子,聖靈.阿們.



\newpage

\section{為上主的事著迷}
\label{sec:81}
\textbf{盧龍光牧師 - 為上主的事著迷}
\newline
\newline
\hyperref[sec:80]{< < < PREV SERMON < < <}
~
\hyperref[sec:toc]{[返目錄]}
~
\hyperref[sec:82]{> > > NEXT SERMON > > >}


路加福音 2:41-52
\newline
\begin{longtable}{cl}
\hline
\hline
章節 & 經文 (和合本修訂版)\\
\hline
2:41 & \begin{tabularx}{0.7\textwidth}{X} 每年逾越節，他父母都上耶路撒冷去。 \end{tabularx} \\ \\ \relax
2:42 & \begin{tabularx}{0.7\textwidth}{X} 當他十二歲的時候，他們按著過節的規矩上去。 \end{tabularx} \\ \\ \relax
2:43 & \begin{tabularx}{0.7\textwidth}{X} 守滿了節期，他們回去，孩童耶穌仍舊在耶路撒冷。他的父母並不知道， \end{tabularx} \\ \\ \relax
2:44 & \begin{tabularx}{0.7\textwidth}{X} 以為他在同行的人中間，走了一天的路程才在親屬和熟悉的人中找他， \end{tabularx} \\ \\ \relax
2:45 & \begin{tabularx}{0.7\textwidth}{X} 既找不著，就回耶路撒冷去找他。 \end{tabularx} \\ \\ \relax
2:46 & \begin{tabularx}{0.7\textwidth}{X} 過了三天，他們發現他在聖殿裡，坐在教師中間，一面聽，一面問。 \end{tabularx} \\ \\ \relax
2:47 & \begin{tabularx}{0.7\textwidth}{X} 凡聽見他的人都對他的聰明和應對感到驚奇。 \end{tabularx} \\ \\ \relax
2:48 & \begin{tabularx}{0.7\textwidth}{X} 他父母看見就很驚奇。他母親對他說：「我兒啊，為甚麼對我們這樣做呢？看哪，你父親和我很焦急，到處找你！」 \end{tabularx} \\ \\ \relax
2:49 & \begin{tabularx}{0.7\textwidth}{X} 耶穌對他們說：「為甚麼找我呢？難道你們不知道我應當在我父的家裡嗎？」 \end{tabularx} \\ \\ \relax
2:50 & \begin{tabularx}{0.7\textwidth}{X} 他所說的這話，他們不明白。 \end{tabularx} \\ \\ \relax
2:51 & \begin{tabularx}{0.7\textwidth}{X} 他就同他們下去，回到拿撒勒，並且順從他們。他母親把這一切的事都存在心裡。 \end{tabularx} \\ \\ \relax
2:52 & \begin{tabularx}{0.7\textwidth}{X} 耶穌的智慧和身量，並神和人喜愛他的心，都一齊增長。 \end{tabularx} \\ \\
[1ex]
\hline
\hline
\end{longtable}
日期:2011年1月16日

主題:為上主的事著迷

証道:盧龍光牧師

經文:路加福音2:41-52

各位弟兄姊妹早晨:

在寒冷的主日,弟兄姊妹來到整年最後一個主日,我們能夠一同敬拜,聽上帝話語,實在是上帝的恩典.

今天伍牧師給我的題目是《為上主的事著迷》,我問他,為何給我這個題目?他說:「因為我時常都很著迷」.其實為一些事著迷,相信大家都有自己的經驗.我不知道大家想想有沒有?你這一生當中或到今天為止,有沒有一些事或東西令到你著迷呢?

今年裡,今天是最後一個主日,你有沒有想過在這一年裡,有沒有甚麼事令你著迷呢?我一日過一日沒有甚麼事令我著迷,當然今天很多青少年最著迷的,就是電子遊戲機.你會認為「老套」,現在不是稱呼電子遊戲機,最近著迷是ipad.因為ipad有很多好玩的遊戲,當然我們看到我們的同學,希望他們是為了讀書而著迷,為了學問而著迷,為了交功課而著迷(內容是否著迷就不曉得了),至少到了交功課的最後限期時就會著迷,因為希望能儘快完成.

但人生當中的時間,應該在不同時候為不同的事情而著迷.我記得我年少時,為了郵票而著迷.當時集郵不用花錢,只要周圍問人拿信封便可以了.當時父母親要求我們集郵,那時候很開心,很著迷.當然那個年代,相信也有不少人集郵,現在就沒有甚麼郵票可以收集了,人們很少寄信,也沒有信封.當時候我們小學六年級到中學,同學不是賭錢,而是集郵；那時同學之間鬥翻書,翻開多少頁,合共的頁數最大的便能贏多少張郵票,真是為集郵之事非常著迷.

當長大後,大家有的為工作著迷,為情人著迷...當然也有的為了所謂的「四子」政策而著迷──我們希望中大畢業後,我們可以賺到錢,要找一架車子；然後要買間房子,才可以娶到妻子；最後要生一個孩子──這就是我們所講的「四子」政策,現在可能仍然有效的.

不過,現在買房子可能會比較困難,要買了房子才結婚,所以很多人會晚婚.人生當中所謂的「著迷」,是因為有些東西是你認為,覺得很重要,很有興趣的,而你願意為那些東西付出不同的代價,因為那些東西真的可以令你整個人沉迷下去.當然今天很多父母正為了自己的兒女而擔心,因為今天有很多各式各樣的東西令人著迷,像打電動,上網等等.但究竟作為一個基督徒,或在你的人生裡頭,有沒有真正思考過,究竟有哪些東西是令你著迷,那些東西又是甚麼呢?最頭痛的是你連自己著迷了還不曉得,以致於你不知道,你竟然為了那些東西付出不同程度的代價,而你也沒去問,這些東西是否值得你付出這麼多的代價?

今天我們所唸的經文,特別是福音書,有一段描述耶穌著迷的事,可使我們再次去思考.大家打開福音書 -- 路加福音第二章41 -- 52節,這段經文,可以說是整本聖經當中,記載耶穌少年時期的唯一一段.也可以說是耶穌出生之後(聖誕節過後),除了講耶穌在聖殿當中,被奉獻之後唯一一處,是在祂成長之前所記述的.

這段經文記載有東西令耶穌著迷,著迷到了一個甚麼樣的地步呢?很明顯的是著迷到連祂的父母去了那裡也都不知道；就像今天很多青少年打遊戲機一樣,可以廢寢忘餐,以致父母親去了那裡?不知道!過了多少時間?不知道!原來要回家了,也不知道!這些就是所謂著迷的表現.所以你想想,有甚麼事情使你著迷了呢?你尚且用這些指標來問問自己,有些甚麼事情可以使你廢寢忘餐呢?時間過了,你不記得了,有甚麼人在你旁邊,你也是不太記得了...

當然我們看見耶穌所著迷的,是在聖殿當中.耶穌十二歲時到了耶路撒冷,祂父母都不曉得祂去了那裡,過了三天才發現,原來祂在聖殿裡.經文四十六節講祂在做些甚麼?祂乃坐在教師中間,一面聽一面發問.凡聽見祂的人,都對祂的聰明和應對感到驚訝!祂父母看見,就很驚奇.祂母親對他說:「我兒,為甚麼對我們這樣做呢?看哪,你父親和我很焦急,到處找你.」耶穌說:「為甚麼找我呢?難道你不知道我應當在我父的家裡嗎?」祂所說的話,他們不明白,於是就同他們下去回到拿撒勒,並且順從他們,祂母親把這一切的事都存在心裡.耶穌的智慧和身量,並上帝和人喜愛祂的心,都一齊增長.

這段經文當然是耶穌復活之後,才寫成文章.所以這段經文當中所記述的,表達了兩個衝突性的感覺.第一,他父母很焦急地找自己的兒子,因為他父母走了一整天卻不見兒子,(當時候的年代,不像現在的香港,會有陌生人拐兒子).第二,他們與村子,家族的人約好,一同到耶路撒冷朝聖,但回程需要幾天的路程,到了晚上,父母才在人群中找孩子.

但在這種的情況之下,孩子不見了,當然會很焦急.對於愛兒女的父母來說,多為自己的兒女著迷,認為兒女是他們最重要的東西.最近,我們知道英國有大風雪,以致香港學生沒有機位回港,要在機場留宿；或許航空公司已安排他們到飯店內留宿,但父母仍然非常緊張,一有飛機回港,很多父母早早到達機場,眼泛淚光,接自己的兒女.這些事每年都有,不過這是三十多年來最寒冷的一天,但這些環境卻給了青少年人們一個學習的機會.可是很多父母卻不以為然,產生許多擔心,驚恐,直到見到兒女時就眼泛淚光,像是心肝寶貝被冷壞了.

甚至今天的教育,最重要的不是大學放榜的入學試,乃是小學一年級的放榜.因為今天認為,只要小學一年級的時候入了名校,便能得到一條龍式的教育──由小學,中學,大學,所以最緊的就是小學一年級的放榜,為此,有些人就要去更改地址等等... 因此,今天兒女成為父母最著迷的,最願意付代價的；甚至,會擔心他們會產生甚麼樣的事而睡不著覺.

耶穌的父母原本是很安心的,直到找不著耶穌時,他們心裡非常焦急,在親戚朋友中也尋不見祂,他們心裡真的急壞了.他們只好返回耶路撒冷去尋找,找了祂三天(耶路撒冷不是很大,半天已經走完了,可是他們萬萬也沒有想過十二歲的耶穌,會在聖殿中出現,並跟拉比和教師在一起).他們這麼著急,找了三天才找到祂,所以母親所講的話,真的是人之常情:兒子,你為何要這樣對我們呢?你看看,父親和我很焦急,到處找你(大部份時候,母親總是說話的那個,父親只在旁邊靜默不語,約瑟在新約裡,只跟天使說話,不過也沒有開口說過任何的話).

但耶穌竟然沒有認錯,還要還嘴說:「為何找我呢?」很正常的,今天的兒女也是這樣對待父母.「你為何找我?」我不應當在我父的家裡嗎?這句話說得太深奧了,今天兒女只會講,為何找我?我去了玩遊戲機；為何找我?去了找同學玩.不過,耶穌的回答是祂的母親所不明白的,但祂母親既沒有打祂也沒有責罵祂.反而,祂的母親把這些話放在心裡去思想.這裡,當然提到耶穌是很乖,很順從父母,並跟他們回家.不過我想這裡有一個格格不入的矛盾,耶穌不高興,因為祂原本高高興興地在聖殿裡──為何你們要找我呢?母親也是不太高興,她焦急地前來找耶穌,找著了卻被還嘴.那時侯應該會發生口角,但聖經卻沒有記載,總而言之,就是要回家了.

這裡描寫耶穌不是一個故意要去頂撞父母的人,但很明顯路加福音記載這事就是要讓我們知道,原來耶穌少年時曾經發生過這樣的事,這些事也是使我們今天作為父母傷心的原因:父母為自己的兒女著迷,但兒女並沒有為我們著迷,他們另有使他們著迷的事.

原來在人生當中,這樣的事比比皆是.作為夫妻也好,同事也好,兄弟姊妹也好,每個人都為著不同的事而著迷.當這些的著迷起了衝突的時侯,該怎樣去處理?怎樣去化解呢?很顯然,馬利亞的化解方法是她所不明白的,我為你著迷,你是我的心肝寶貝,你是最重要的,我這樣辛苦地找你,但原來你並沒有找我,你為另外的東西而著迷,她並沒有在那刻責罵耶穌,反之她把一切放在心裡,所以你會發現今天的天主教教會,對馬利亞有一定的尊重.

雖然有些人以為天主教教會,拜馬利亞；但在正統的天主教教會裡,並不拜馬利亞,他們很尊重馬利亞,甚至在他們的神學上認為,耶穌是無罪的,她是從聖靈懷了孕,馬利亞應該也沒有罪,以所謂「無原罪」的觀念來理解馬利亞.

事實上,今天更正教對馬利亞也太忽略,甚至我們沒有去理解,欣賞馬利亞(所以許多教會也不是太著重婦女).馬利亞這位母親,不像在聖誕節時,我們說普世歡騰那般值得慶祝,其實第一個聖誕是很慘的,因為這位母親未婚懷孕,有誰知道這個嬰孩是從聖靈來的?除了聖靈,沒有人知道,約瑟也不知道,聖經只記載約瑟聽到天使對他說:「你娶了她回來啦!」

約瑟本想把她給休了,一個私生子要降生,爸爸是誰也不知道,只有馬利亞和聖靈知道,在不知道該怎辦的情況底下,這個想休了她,那個未婚懷孕,甚至根據路加福音所記載的,耶穌要在馬槽上誕生,因為找不著地方,無人容下他,是這樣開始的一個故事.今天聖誕節的偉大,但當年在這樣的危機下,在這樣的羞辱當中開始了耶穌基督的生命.耶穌的生命將一切都改變,馬利亞的起點就是一個未婚媽媽,她一開始可能受到丈夫的羞辱... 一生當中,除了這次的頂撞外,耶穌其實也常頂撞祂的母親.如果你看看聖經你便知道,昔日耶穌傳道,傳道到了一個很著迷的地步,連飯也顧不著吃,然後母親很關心地來找祂:哎地!兒子,食飯啦!耶穌是怎樣回應祂的母親?

門徒對耶穌說:「你母親來找你」,耶穌回答說:「誰是我的母親?」門徒說:「誰是你的母親也不知道?她這麼辛苦,滿有羞辱地撫養你,誰是你的母親也不知道?」耶穌說:「凡遵行父旨意的人,就是我的父親母親,兄弟姊妹了.」──中國文化對這段經文也很不接受.

中國人認為,豈有此理,誰是你的父母也不知道?誰是你的兄弟姊妹也不知道?耶穌在釘十字架時,祂母親在十字架下面,傷在兒身痛在娘心──馬利亞為她的兒子著迷,以致無論如何她也要誕下這個兒子,耶穌十二歲時頂撞她,長大了也頂撞她──釘十字架時,最痛的是她,但她卻一直在耶穌生命當中與祂同行.

很多事情她不明白,但她心尊主為大,她願意放在心裡頭去思考.我們會為著一些東西著迷,但究竟為著甚麼而著迷呢?當我們發現身旁的人所著迷的跟我們不一樣的時候,你會頂撞他,否定他,還是將別人所著迷的事放在心裡頭去思考呢?

另一方面,我們看見耶穌所著迷的事是很正常的,原來祂為上帝的事而著迷,就是我們今天的講題「為上主的事著迷」.那時所描寫的,只有兩個字,就是「坐在教師中間,一邊聽,一邊問」.

今天我們的確不明白為何祂這樣著迷,如果你未玩過遊戲機,你真的很難明白.近來我也學習玩遊戲機,這才明白為何會著迷,真的很容易著迷.除了畫面很吸引人之外,你有如置身現場一樣,看見,聽見,手便動了,輸了想贏,贏了很高興,又會繼續玩下去,在遊戲裡又可以升級,每個人也為著升級而著迷.

原來電子遊戲機是人生的寫照,今天我們為何會著迷?因為我們看見許多使人著迷的東西,用佛家的思想來說:在這個紅塵的世界裡,有許多五顏六色的東西,許多看得見的東西,的確是很吸引人的.

小朋友看見郵票很美,被郵票吸引,被不同國家的郵票吸引；長大後,當然是被異性吸引──哪個女孩很漂亮,哪個男孩很英俊,所以青少年可以為了失戀而自殺.通常我是不會責怪那些人,因為自殺是人類最高的一個尊嚴,動物不會自殺的,只有人類才會.很可惜,我們會為很愚蠢的事跑去自殺.

今天我們會為了一些我們認為很重要的東西而生存,也願意為了那些東西而放棄生命.然而,當我們喜歡一個人的時候,縱然她的聲音是吵啞的,你也感到很磁性,可是當我們一旦不喜歡的時候,就會發現那聲音很難聽.

因此當你玩遊戲機的時候,你整個人會被聲音,視覺所吸引,因為你與那東西結上了關係,尤其是它有升級的功能(在這個社會當中,我們很著重升級；讀書要升班,教授也要升級),這就使你更多的投入,更多的著迷.

耶穌不一樣的地方是,祂一邊聽,一邊問.其實今天我們很少認真懂得如何發問,在香港的教育制度下,很多時候我們只是聽,看,做和答題目,並沒有教導我們如何發問問題.當我的女兒還小時候,她所回答的跟老師不一樣,老師就說她錯了,可是我看見她所作答的都蠻好,「沒用的啦!老師說我錯了」.我認為不是的,女兒卻說不可以,要按老師的要求回答.我便要求她回答兩次,一次回答給老師看,一次回答給爸爸看,但她確不願意回答.沒辦法,因為要按老師的要求才能拿分數.

今天的教育制度,是要按對方的要求而作答.人生當中,不會思考,不會發問,不會尋求理解和不同的可能性,所以人著迷一樣東西時,只會為了那樣東西,為了看見,聽見,然而卻忽視了其它的東西.

耶穌很簡單,只是坐在老師面前數天,但祂一邊聽一邊問,簡簡單單的兩句話,正正表明了祂所著迷的,那我們要問:為甚麼要著迷?那東西是否值得我們去著迷?是否只有那樣東西才令我們著迷?相反的,祂母親與兒子所著迷的東西不一樣,但這位母親會想想兒子,把她不明白的事放在心裡.

馬利亞和耶穌為不同的事而著迷,母親為兒子,耶穌為著上帝的事情.我們知道馬利亞之所以偉大,正是因為她聽到天使告訴她要從聖靈懷孕,即使從人來看,她要受羞辱,但她卻能唱出《尊主頌》:我心尊主為大.在她的一生中,難道馬利亞不正是為上主的事而著迷嗎?她不是停留在兒子身上嗎?

為何她竟然能夠容忍兒子多次的頂撞她,直到耶穌釘死在十字架上,她都默然地支持自己的兒子?她是為著上主的事才能如此.耶穌所著迷的跟馬利亞所著迷的本是沒有衝突,但在外人看來,那是不一樣甚至是頂撞的；耶穌跟馬利亞都有一個共通著迷的地方,馬利亞也會問,為甚麼?耶穌為為何會如此回應她?是否有一些更好的理由呢?

今天我們不是迷信於基督教信仰,我們會批評民間的信仰,認為他們所追求的,是有求必應.很多人迷信的原因,是因為「靈驗」.在古裝電影／電視劇中,常出現許多破廟,為何廟會破嗎?不靈驗!所以很慘的,不靈驗就沒有人去拜祭!今天基督教信仰,要是靈驗就不得了,因為猶太人從舊約開始,他們的上帝都不是很「靈驗」的,除了我們看見出埃及,在曠野兜了四十年才進入迦南地,在迦南地又遇到各式各樣的困難等等... 甚至在二千多年前,看見以色列被亡國,被俘擄,做了二千年亡國奴.

但你會發現這個民族,認為上帝很偏心,揀選這個民族為選民,這個民族很聰明,很有智慧,很有本事,今天不只是紐約的金融被猶太人控制,就連奧斯卡一半的獎項也是給猶太人所奪取的,因此你會發現這個民族有很多特點,很多天賦.但在他們的生命當中,尤其是跟上帝的關係裡頭,他們往往選擇背叛上帝,所以上帝並非靈驗與否,上帝會拯救,會審判,為的是要他們行在正路中.今天我們看見這個迷信,因為上帝給予我們聰明智慧,我們不會思考,我們有這種心靈,我們不會擴闊,收窄才看到我們著迷的東西,看不見為何別人會這樣做.

耶穌的「一邊聽,一邊問」也是我們基督教信仰的特色,別人也常說我們迷信,因為我們可以信得很迷,但為何我們值得如此「迷」呢?一個人只有保持冷靜,保持距離,當然就不會玩電子遊戲機,但卻也不能享受別人打電子遊戲機時的樂趣.

人生也一樣,我每天都是吃飯,睡覺,起床,工作... 有些人就是平平淡淡地過,他可能也很著迷這些平淡的日子.著迷是因為在人的一生當中,你覺得有些東西是重要的,有些東西使你投入,有些東西使你覺得要付出代價,使你覺得「這是值得的」.

今天有人甚麼東西都不著迷,因著某些原因,他對很多東西都很疏離,很沒興趣.人生當中,總是有些東西是要著迷的,問題是你是否知道為甚麼著迷?有甚麼攔阻你去著迷?那些究竟是甚麼?

今天我們知道基督教信仰從耶穌出生開始計算,至今已二千多年,當中因為有無數人被著迷,而著迷的竟然是誕生於馬槽,在別人眼中為私生子,且被釘死在十架上的一個罪犯.為何有這麼多人為祂著迷?為祂犧牲生命?很明顯,因為祂叫著迷者的生命被改變,不是改得一塌糊塗,乃是藉由改變的生命,繼續改變他人的生命.耶穌改變很多人的生命,將一些被人看不起的人,邊緣化的人,無知識的人,在困苦中無希望無朋友的人或痲瘋病人,妓女,社會上被遺棄的人,這些人通通得到重新的生命和希望.

或者像保羅這樣有學問的人,他為猶太的信仰著迷,但當他遇見耶穌的時候,他才知道以前以為上帝所恨惡的外邦人(猶太人看不起的外邦人),原來耶穌基督在十架上已化解了所有的冤愁,化解了猶太人在律法所咒詛的,因為外邦人也是上帝所愛的.

各位弟兄姊妹,你知不知道有些東西使你著迷,使你沉溺當中而不自覺的呢?聖經叫我們要著迷,但要弄清楚「著迷」甚麼?做父母的,做兒女的或今天你是一個領袖,甚至在事業當中是一個很成功,或是失敗的人,你要有所著迷的東西,你不要放棄自己,但你要弄清楚這樣東西是否有價值,是否有生命的價值?耶穌基督「一邊聽,一邊問」的時候,祂的生命被整個的「道」所吸引,然後祂才體會到父神的心意.今天很多人以為,耶穌一出生就知道祂是基督,上帝的兒子.

你看看聖經就會知道,耶穌跟我們一樣是人,祂在尋找上帝在祂身上的心意,特別是在曠野四十天受試探時,魔鬼問祂「假如你是上帝的兒子」.當祂一路走的時候,祂需要聽,需要學,需要問,然後祂尋找,發現上帝在祂身上的心意,到了客西馬尼園時,祂沒有為自己的生存而著迷,祂願意順服上帝的心意,願意接受十字架的苦難,所以聖經上說,上帝就將這個耶穌升為至高,祂本與上帝同等,但卻不與上帝同等為強奪,反取了奴隸的形像,至終上帝叫祂從死裡復活,證明祂是上帝的兒子.

在祂的生命當中,有很僕實的一面,就像你我一樣,但在祂的生命當中,就是「一邊聽,一邊問」,為上主的事著迷,然後祂體會到祂就是上帝的兒子,能夠改變其他人的生命,我們是否都找到上主的事是甚麼?一些令我們著迷的事?以致我們一生當中,都以這事為我們生命的目標!

新一年又開始,希望大家都有一些令你著迷的東西.可能你我所著迷的和不相同,但都是為了上主的事而著迷.讓我們同心禱告:

慈愛天父,我們感謝讚美你,因為你是我們生命的主宰,我們誕生在這個世界,可以是茫然不知,我們不知自己的方向,自己的喜好,不知道有甚麼令我們著迷的事.但今天我們從耶穌基督的生命,從馬利亞的生命,我們看見他們各有不同的東西,但原來他們都是為了你的事而著迷.

主!就求你今天引導我們,在我們人生當中,有許多東西不知不覺中常使我們著迷,在不知不覺中有些東西成為我們最終極的關懷,有些東西就成為我們的上帝.

主!你讓我們有安靜,省察的心,不要只被我們今天所看到的各式各樣的東西,或屬這個世界的東西,吸引著我們而不自知.又或是對所有事物也不感興趣.

主!求你的靈感動我們,知道我們的生命是你所賜予的,在人一生的旅程中,求你使我們不要迷失方向,好叫我們可以為你的事著迷.因為我們知道我們的生命可以被你使用,以致其他人的生命也可被你的愛所激勵.

主!今天我們在這年的最後的一個主日,讓我們反思,有沒有甚麼東西讓我們不知不覺地被著迷.或者當我們決心為你的事,奉獻生命的時候,卻往往失落,跌倒,主!求你調整我們的方向,使我們一生當中,在不同的境遇,不同的崗位上,都願意為你而著迷.將我們的生命獻上給你,求你使用,我們為過去一年感謝,更為新一年盼望,求你與我們同在,祈禱奉耶穌基督名求.阿們!



\newpage

\section{這個冬天不算冷}
\label{sec:82}
\textbf{伍渭文牧師 - 「這個冬天不算冷」}
\newline
\newline
\hyperref[sec:81]{< < < PREV SERMON < < <}
~
\hyperref[sec:toc]{[返目錄]}
~
\hyperref[sec:83]{> > > NEXT SERMON > > >}


提摩太後書 4:
\newline
\begin{longtable}{cl}
\hline
\hline
章節 & 經文 (和合本修訂版)\\
\hline
4:1 & \begin{tabularx}{0.7\textwidth}{X} 我在神面前，並在將來審判活人死人的基督耶穌面前，憑著他的顯現和他的國度鄭重地勸戒你： \end{tabularx} \\ \\
[1ex]
\hline
\hline
\end{longtable}
日期:2011年1月23日

主題:「這個冬天不算冷」

証道:伍渭文牧師

經文:提摩太後書4章

我們以為一個老人家,當他們晚年時,應該是享清福的時候.不錯,假如他是兒孫滿堂,身體健康,晚年確是享清福的時間,但如果一個人獨居,身體開始漸漸虛弱,那個滋味就不同了.今天我們所讀的經文,記述一位年過六十的老人,不單是獨居,更是上了枷鎖,被囚禁在地下只有四立方公尺的斗室中,微弱的陽光,祇能透過室頂窄小的窗戶射進來.然而,這位老人沒有自怨自艾,這位老人就是使徒保羅,他是因傳揚福音而下獄.面對同工的變節和離棄,他心中也沒有任何憤懟.冬天很快就來到,因為地中海氣候,冬雨夏乾,屆時囚室四壁潮濕,陣陣冷風,寒氣迫人,而且冬天過後,保羅自己預感將面對身首異處的處決:「

我現在被澆奠,離世的日子到了.」(提後四6)隨著年日的流逝,他身體開始衰弱,他需要麼加衣禦寒:「

我在特羅亞留於加布的外衣,你來的時候可以帶來.」(四13)在人看來,這是一個寒冷的冬天,但對保羅來說,這個冬天不算冷.

一    這個冬天不算冷,因為他看見許多人被上主使用窩心暖意.保羅沒有覺得這是冬天太嚴寒,因為很多事使他感到窩心.不錯,他看到人的軟弱,離棄信仰,但同時有更多的人,持守真理,勇往向前.

當他初次被提訊時,沒有人前來幫助他；他期望站在他旁邊的人,竟然離開他,特別是底馬:「

因為底馬貪愛現今的世界,就離開我往帖撒羅尼迦去了.」(四6)底馬是誰?為何保羅感嘆底馬離他而去呢?底馬就是保羅在腓利門書所提及曾經和他出生入死的同工.「

與我同工的馬可,亞里達古,底馬,路加也都問你安」(門24)你有沒有好友,多年和你在團契一起事奉,為信仰大發熱心,甚至是領你信主,但現在己沒有返教會,他沒有放棄信仰,但金劍已沉埋,壯志蒿萊.

世界的魅力.現今的世界,永遠閃爍著迷惑,叫天路客放軟了腳步,模糊了遠象,偏離了方向.貪愛一字,顯示了世界永恆的魅力 and \#8212;眼目的情欲,肉體的情欲,和今生的驕傲.有一首歌這樣描述現今的世界:我們活在這時代,除了追求金錢,成功,名譽,炫耀外,追求其他的價值會被嘲弄,甚至遭消滅；唯有金錢,成功,名譽,炫耀,因為我們活在唯有物資的世界.(We are living...in the age, in which the pursuit of all values other than money, success, fame, glamour have either been discredited or destroyed. Money, success, fame, glamour, for we are living in the age of the thing.)

貪愛一字,也表明信徒不是經過誠實的思考而放棄信仰,乃是被吸引隨流飄離信仰.魯益師(C. S. Lewis)說:「我們要常常被提醒所信為何,沒有一樣信仰會自動在我們的頭腦中茁壯成長的,它需要培育滋養.事實上,若你好好查驗一百位離棄基督信仰的人,我懷疑多少是經過誠實思考而離開信仰,很多都是隨流飄走,不是嗎?」(We have to be continually reminded of what we believe. Neither this belief nor any other will automatically remain alive in the mind. It must be fed. And as a matter of fact, if you examined a hundred people who had lost their faith in Christianity, I wonder how many of them would turn out to have reasoned out of it by honest argument? Do not most people simply drift away?- C.S. Lewis, Mere Christianity) .

闊大視野.但保羅不會因底馬貪愛世界離開他而意志消沉,因他看到更多的人被主的愛所吸引,愛慕主的顯現多於眼目的情欲,羨慕將來公義的冠冕多於現今的名和利.他看到革勒士往加拉太,是為事工的需要,提多往撻馬太也是為了事工,而且路加一路都陪伴左右,還有馬可,他就是馬可福音所記載,耶穌被賣時,有一少年人被兵丁追問:你可否認識耶穌?他急忙把身上的布拋開逃跑,他是一位十分怛怯的人.:「

有一個少年人,赤身披著一塊麻布,跟隨耶穌,眾人就捉拿他.他卻丟了麻布,赤身逃走了.」(可十四51)第一次保羅與巴拿巴到歐洲傳福音時,他中途在旁非利亞怕太艱辛而離隊,所以第二次保羅到歐洲傳道時,堅持不攜帶馬可.(徒十三13；十五36~38)但勸慰子巴拿巴慧眼識英雄,有培訓的恩賜,曾經栽培保羅,在他悉心提攜下,馬可得到接納和磨練,最終成為幹練的宣教士:「在傳道的事上於我有益.」

保羅的眼目,不單看到令他失望的底馬,他看到更多被主改變生命的同工.在教會,祇要我們細看,我們也會看到上帝在不少人的生命中工作,生命得著激勵.

寬廣的心.保羅也有寬廣的心,面對人的攻擊和離棄.雖然猶太人領袖銅匠亞歷山大多多苦害自己,保羅沒有報復反擊,祇說:「

主必照他所行的報應他.你也要防備他,因為他極力敵擋了我們的話.」(四14, 15)面對被同工離棄,也沒有怨憤:「

我次申訴,沒有人來幫助,竟都離棄忘我；但願這罪不歸與他們」(四16)面對逆境,總看到上帝的恩典.因為被囚,獄卒得聞福音信主,福音被傳開；雖然被囚,但上主保守了自己的生命:我也從獅子口裡被救出來.(四17)所以,這個冬天不算冷.

二    這個冬天不算冷,因為有人為他帶來外衣保羅第一次被囚是軟禁,可見客人,但這次上了枷鎖,而且身體已不及當年勇了.然而,冬天寒凍,有衣物禦寒；禁見訪客,可以閱讀,亦可寫作.「

我在特羅亞留於迦布的那件外衣可以帶來,那些書也要帶來,更要緊的是那些皮卷.」(四13)我們若接受自己的限制,在制約的環境中,也可以從容自若,盡量做好自己.保羅曾經年輕過,那時他拿著大祭司的文書,進入大馬士革,搜捕基督徒,那時候保羅何其強壯.史提反被擲石塊,行刑的人,都把衣服放在保羅的腳前,那時他何其威武；壯年時,保羅曾風塵僕僕,遠赴歐洲傳道,未嘗言倦.他曾靠主行神蹟,使死人復生,病人得愈,驅魔趕鬼,如今天氣寒冷,難道不可以求上主賜他金剛不壞之身?

保羅明白,人的限制和自然規律,是上主為我們立下的界限,不能逾越,他甘心順服,要求提摩太攜來外衣,抵禦風寒.其實很早的時候,他曾三次求主醫治他身體的軟弱 and \#8212;拿走身上的的刺,這刺很可能指眼疾,加六11:「

請看我親手寫給你們的字是何等的大呢?」但上主不允,他就順服,毫無怨言,藉此認識到自己的限制,學到謙遜.在不同的人生季節,都有不同的限制,萬物各按其時,在諸般的限制中,盡量做好自己.

這星期有機會在午餐與一位退了休的長輩同席,談起退休後的生活.他說退休後生活多姿多彩,他義務協助我們神學院緬甸同學書寫英文,每天做一個小時健身,以前是沒有做健身習慣的.每星期一次和一班朋友學吹色士風(saxophone) ,被邀請表演時,大家都穿上緊身的黑色衣服,我們中文大學前醫學院院長鍾尚志教授也在其中,他學色士風目的是呈獻給在太太作為生日禮物.

年老退休,可學色士風娛樂老伴；坐牢禁見,亦能看書寫作.在諸般人生限制中,學曉接受現實,與限制共生共舞,雖然冬天寒凍,也不算冷.

三    這個冬天不算冷,因為春天就快出現了冬天己經來到了,春天還會遙遠嗎?保羅很早就知道,人生季節最後的一站是冬天,所以他很早就栽培提摩太為接班人,他對提摩太說:「

但你所學習的,所確信的,要存在心裡；為你知道是跟誰學的.」(三14)保羅知道他快要澆奠離世了,巴拿巴怎樣訓練自己,培育自己,他也要如何訓練提摩太,接續他未完成的工作.人要過去,但上主的工作會持續發展.保羅不是為了事工而做事工,他是透過不同人的生命,去推動事工,所以保羅在人生的晚年,特別珍惜同工,除了我們提過的提摩太,馬可,路加,還有許多人名,提後四19,21:百基拉亞居拉夫婦,亞尼色弗一家人,以拉都,友布羅,布田,利奴,革老底亞等等.

家學淵源提摩太.在不同的人物類型中,我們看到一路順暢,預備接班的提摩太,他年少時己受教於敬虔的袓母羅以和母親友妮基.「

想到你心裡無偽之信,這信是先在你外祖母羅以和你母親友妮基心裡的,我深信也在你的心裡.」(一5)我們看到家庭教育對人材的培育何其重要.羅以 and \#8212;友妮基--提摩太,提摩太家學淵源,得天獨厚,保羅來料加工,便成大器.盼望我們在新的一年,也重視家庭宗教教育,食飯前謝恩,臨睡前講一簡短聖經故事,然後合手謝謝天父睡覺.

大器晚成馬可.另一位是起步緩慢的馬可,馬可也有敬虔母親馬利亞,她開放家庭作為禱告的地方 and \#8212;馬可樓(徒十二12).可能長於富裕家庭,馬可難耐艱辛,他也生性怯懦,但經過生命導師巴拿巴悉心訓雕琢,鼓勵誘導,給他第二次機會實踐磿練,最終也能大器晚成,在傳道的事上於保羅有益.

摸不透的底馬.最後是摸不透的底馬.底馬隨流飄去,究竟是外在的壓力令他退縮,或是靈魂深處的潛伏的牽引使他迷失呢?作為馬其頓省省會的帖撒羅尼迦的大都會生活,何時開始對他有特別的吸引力呢?能與保羅,路加等使徒同工並受他們肯定,底馬表現令人激賞,但在靈魂深處,卻孕育了靈性的危機,底馬已被現今世界的閃爍炫耀所吸引.就如我上一次講道所說:你要懷孕生子這句話,可以是生命的應許,也可以是嚴厲的警告.「

私欲既懷了胎,便生出罪來,罪既長成,就生出死來.」(雅一15)

究竟底馬何時讓私欲在心中成孕呢?保羅不知道,但底馬自己應該知道,我們的靈魂渴慕甚麼,其他人可能不知道,但我們自己應該知道.保羅開始不用馬可,但後來愈來愈欣賞馬可,我想馬可有一顆羨慕-聖工的心,他的靈魂深處,渴想事奉上帝,被主的愛所吸引,愛慕主的顯現多於眼目的情欲,羨慕將來公義的冠冕多於現今的名和利.

今天是2 0 1 1年第一個主日,是一年的起跑點,我們不知道在新的一年中,會否遇到寒冷的冬天,刮起冷風,上帝好像隱藏起來,不理我們.這時候,請你仔細看看,在諸多聖徒的生命中,上帝的工作何其真實；在我們諸般的限制中,上帝的恩典也何其真實,我們的日子如何,力量也如何.

願在新的一年,我們的心被主的愛所吸引,愛慕主的顯現,多於眼目的情欲,羨慕將來公義的冠冕,多於現今的名和利.阿們.



\newpage

\section{盡諸般的義}
\label{sec:83}
\textbf{賀志勇博士 - 盡諸般的義}
\newline
\newline
\hyperref[sec:82]{< < < PREV SERMON < < <}
~
\hyperref[sec:toc]{[返目錄]}
~
\hyperref[sec:84]{> > > NEXT SERMON > > >}


馬太福音 3:13-17
\newline
\begin{longtable}{cl}
\hline
\hline
章節 & 經文 (和合本修訂版)\\
\hline
3:13 & \begin{tabularx}{0.7\textwidth}{X} 當時，耶穌從加利利來到約旦河，到了約翰那裡，請約翰為他施洗。 \end{tabularx} \\ \\ \relax
3:14 & \begin{tabularx}{0.7\textwidth}{X} 約翰想要阻止他，說：「我應該受你的洗，你怎麼到我這裡來呢？」 \end{tabularx} \\ \\ \relax
3:15 & \begin{tabularx}{0.7\textwidth}{X} 耶穌回答他：「暫且這樣做吧，因為我們理當這樣履行全部的義。」於是約翰就依了他。 \end{tabularx} \\ \\ \relax
3:16 & \begin{tabularx}{0.7\textwidth}{X} 耶穌受了洗，隨即從水裡上來。天忽然為他開了，他看見神的靈降下，彷彿鴿子落在他身上。 \end{tabularx} \\ \\ \relax
3:17 & \begin{tabularx}{0.7\textwidth}{X} 這時，天上有聲音說：「這是我的愛子，我所喜愛的。」 \end{tabularx} \\ \\
[1ex]
\hline
\hline
\end{longtable}
日期:2011年1月30日

主題:盡諸般的義

証道:賀志勇博士

經文:馬太福音 3:13-17

「

當下,耶穌從加利利來到約但河,見了約翰,要受他的洗. 約翰想要攔住他,說,我當受你的洗,你反倒上我這裡來麼.耶穌回答說,你暫且許我．因為我們理當這樣盡諸般的義．〔或作禮〕於是約翰許了他. 耶穌受了洗,隨即從水裡上來．天忽然為他開了,他就看見　神的靈,彷佛鴿子降下,落在他身上. 從天上有聲音說,這是我的愛子,我所喜悅的.」太3:13-17

一,    引言(人生的意義)做人容易嗎?做人到底是為什麼? 前些日子看電影《歲月神偷》,讓人唏噓不已.電影裡面,爸爸是個鞋匠,他看著鞋字說"鞋字半邊難啊",因為鞋字左邊是個革字,而難字的繁體字,左半邊也是個革字；太太搖搖頭說"鞋字還有半邊佳了",原來鞋字右半邊跟佳美的佳字是一樣的.做人,是不是跟做鞋一樣,半邊難,也半邊佳呢?《歲月神偷》裡的爸爸,辛辛苦苦做鞋子,賺不了幾個錢,還要給員警交保護費；好不容易供大兒子讀名校皇仁書院,大兒子樣樣拿第一,想不到又突然得了血癌.爸爸辛苦給大兒子湊藥費,又突然間來了場大颱風,把整個家都吹了個稀爛,鞋子飛了一天一地.最後,大兒子還是死在了醫院.沒幾年,爸爸也積勞成疾過了身.嗨,做人真不簡單啊.

我們基督徒呢?做基督徒是不是處處有主恩典,事事有主保守呢?很明顯,不是如此.我以前去的那家教會人比較多,要報告的事項也多.報告單張上左邊印著誰誰誰結婚啦,誰誰誰喜得貴子啦,誰誰誰出國留學啦,感謝主的恩典,主的眷顧.阿們阿們!  右邊又印著誰誰誰得了癌症,誰誰誰要做手術,誰誰誰去了菲律賓沒回來,求主醫治,求主垂憐.阿們阿們! 每次聽到這些,我就覺得我整個人都要分裂了.我想問神,神啊,你到底要我們做什麼?對於人,你到底安排了什麼呢?

人在這世上行走幾十年,回過頭一看,做得最多的事情就是睡覺、吃飯、喝水,當然,還有工作.不忙的時候就煩；不煩的時候就忙.前幾天查經班,一個姊妹分享說,她忙得都忘記給孩子交午餐錢了,開學第一天,學校打電話問她,你兒子這個學期還要不要吃午飯啊?她跟我說:"我才40歲,記性就這麼差,胃又不好,腰椎間盤又突出,真是命苦啊.我常常問神,神啊,你讓我降生到這個世上,到底是為了啥啊?"

人生的問題,千百年都有人問.答案如何,千姿百態.佛陀悟到了空,老子找到了道,馬克思想到了唯物主義,我們基督徒又找到了什麼?對於人生,要知道基督徒找到了什麼,就要先看看我們的基督找到了什麼.好,讓我們一起來看看這段非常離奇的經文.馬太福音3:13-17節.

二,    解經(神為什麼喜悅耶穌)這段經文之所以離奇,是因為它講的是耶穌還沒有傳道之前,去找施洗約翰受洗的故事.這並沒有什麼離奇.最離奇的是,當他受洗之後.天開了,神的靈如鴿子落在他身上,天父說:"這是我的愛子,我所喜悅的."三位一體的上帝同時出現在一個畫面,這在新約中是絕無僅有的."

這是我的愛子,我所喜悅的.``同樣的話還出現在馬太十七章5節,那時耶穌登山變相,從雲中有聲音說:"

這是我的愛子,我所喜悅的."

如果說耶穌登山變相時,已經傳播了福音,大行了神跡,廣收了門徒,所以得父神的喜悅.那這一次呢,神說喜悅耶穌,究竟是喜悅耶穌什麼?神說這話的時候,耶穌還沒有開始傳道,他是在受洗,勝過試探之後,才開始傳道的.所以這個"喜悅",一定不是喜悅耶穌傳道三年半的工作表現.那個時候耶穌還沒有行神跡、醫病、講道,也沒有什麼大隊群眾跟隨他,還沒有什麼工作表現了,為什麼神會說"這是我的愛子,我所喜悅的."?

我想,大概有如下幾個原因:

第一,神喜悅他到約翰那裡受洗,並說"我們理當這樣盡諸般的義".耶穌基督在這些禮儀小節上也願盡人的本分,這證明他降世為人是甘心情願的.他以一個人的身份來到約翰面前,受這個悔改的洗禮.這並不是說,祂承認自己有罪,須要領受悔改的浸禮,而是表明,祂要站在人的身份上,履行神對所有人的規定.

約翰是什麼人?約翰就是那個在曠野大聲呼喊"預備主的道,修直他的路" 的人,他只是神的開路先鋒.開路先鋒是做什麼的?開路先鋒是為後面的元帥掃清障礙的人.約翰脾氣不小.當他看見法利賽人跟撒都該人也來受洗,他就罵他們:"毒蛇的種類!誰指示你們來逃避將來的憤怒?"他脾氣這麼大,如果今天想按立做牧師,恐怕不容易.你想想看,如果今天有個牧師站在臺上,對著大家說"毒蛇的種類!誰指示你們來逃避將來的憤怒呢?"台下的聽眾肯定嚇跑了一半.約翰為什麼這麼嚴厲?因為他所行的乃是悔改的洗,乃是要人痛下決心悔改,"要結出果子來,與悔改的心相稱."所以他會很嚴厲.

猶太人一般是地位高的人給地位低的人施洗.但耶穌基督並沒有罪,並不需要悔改,他為什麼到地位比他低的約翰那裡受洗?我想,他來約翰這裡受洗.是因為他要順服天父；他謙卑自己,完成天父定給這世界的禮儀."盡諸般的義"的義字,原文的意思可以做禮儀來解.禮者,義也.上帝的禮儀--律法--就實實在在體現了上帝的義.這點跟中國人很相像.我們中國人很強調"守禮",一個人做事不講規矩我們會罵他"無禮".猶太人也是這樣,他們很注重禮,很注重律法書上說的那些"禮".耶穌來到約翰面前受洗,表明自己是一個完完全全的猶太人,所以也按律法守禮.

我想耶穌來到約翰面前,把他高貴的頭顱低下來,這一幕非常感人.因為這樣,他不但守了律法,這還表明,他願意守神給人定的禮儀,他願意認同我們罪人,做人的一份子,而不是做一個高高在上的神.我們的神是個有血有肉的人,這一點非常感人.

基督教非常感人的地方就是,我們的神有血有肉,我們的神同時也是個愛我們這些罪人的人.大家記不記得拉撒路復活的故事.耶穌知道拉撒路死了,他看見瑪利亞哭,就"心裡悲歎,又甚憂愁".(路11:33) 他並沒有非常得意地把眾人分開,說"讓開讓開,讓我來,讓我來,哭個啥啊,我會讓他復活的."他並沒有,他看其他人哭了,他跟其他人一樣哭了.看到的人就說"你看他愛這人是何等懇切."

這就是基督教的愛,基督教的神並沒有站在道德高地,指責我們說"你們自己要犯罪,死是你們自找的."沒有,他完全沒有!他自己降世為人,謙卑地體驗人的苦難,人的感情,人的掙扎.並且他受洗--置自己於死地.這種義不是那種居高臨下的義,這種義是一種滿帶恩慈,滿懷順服的義.這就是耶穌蒙天父喜悅的地方.

第二,我想,神不僅喜悅耶穌順服律法,以人的身份到約翰面前受洗.神更喜悅耶穌出來傳道之前的三十年,都"盡了諸般的義"--也就是說,耶穌在前三十年,守全了神給猶太人的禮儀.他受過割禮；他敬重神的僕人,在十二歲那年,就坐在教師中間,一面聽一面問--在人的地位上受教.更重要的是,他道成肉身,在拿撒勒這個小地方,一個連加利利人都看不起的小鄉村,做一個平凡的人.他三十年在家盡孝,無聲無息,卻處處都盡義,這才是最不容易的事.

人是很麻煩的.耶穌是個木匠,他當時肯定要辛苦工作；一方面,他要事奉母親瑪利亞(當時約瑟死了沒有,研究資料還不清楚),不過很清楚的是,耶穌曾經把母親託付給約翰,所以約瑟在世的日子不會太長；另一方面,耶穌是長子,還要照顧5個弟弟妹妹.一家的長子不好做啊.以前很多人家裡都有弟兄姊妹好幾個,長子長女總是最吃虧的那個.父母生第一胎,沒有經驗,不太會照顧小孩,所以長子長女吃的苦最多.等到是小兒子,父母有經驗,就不同了.所以聖經裡的小兒子都是有點嬌慣的.

耶穌是長子,他的成長肯定吃了不少苦.我記得看電影《耶穌受難曲》,裡面就演到耶穌在做木匠活.做木匠活很辛苦.小時候我就看見有木匠挑著沉重的擔子,到處去找要做傢俱的人.木匠的擔子很重,但更重的是一家的生活擔子.我想,耶穌的擔子肯定不輕.但這麼重的擔子,在電影裡頭,他還為母親精心打造了一張桌子.雖然瑪利亞說耶穌打的款式不流行,但從瑪利亞的眼神裡頭,我們看到了他對耶穌孝心的喜悅.耶穌就是這樣盡他做長子的本分--盡諸般的義.

除此之外,耶穌生長在小鄉村,那種地方自然環境簡單,可人的環境卻不簡單.我小時候在鄉村過暑假,發現鄉村裡頭是非多；大家不做事又不賭錢的時候,就開始講是非了.三姑六婆,前村後村,誰家偷了菜地的菜；誰家偷了鄰居的雞,誰家的孩子不學好、沒出息,誰家的女人偷懶不做事,誰家欠錢不還錢.聽得多了,你會覺得這世上沒一個好人.但你想想看,我們的神耶穌在小鄉村長了三十年,聽到了多少人世間的是是非非,看到了多少人性的醜惡呢?有時候,大的打擊容易忍耐,長痛不如短痛；但是小折磨不斷,卻實在讓人煩心不已.

整整三十年,耶穌看著他身邊人的種種小善微惡,看著這些鐵石心腸,將來要把他推下山谷的人,我不知道耶穌的心情會是怎樣? 我記得幾年前,我第一次去杭州做福音事工,心裡面滿是事奉內地同胞的熱情.我去深圳坐飛機,安檢完之後到候機廳登機,正在排隊時,居然有人插隊.我跟他說"登機又不用趕時間,你為什麼要插隊呢?"他看著我,對我說:"因為我素質差啊."我當時真是又好氣又好笑.在飛機上,我心想,這就是我要去事奉的人群,這就是我要去傳播福音的人群麼?這樣的人,我能愛得起來麼?我當時心裡一陣彷徨.

現在想,耶穌要面對人三十年,他的心情只會比我更差.但他還是要去愛人,還是要按神的話,無聲無息把自己隱藏了三十年；又按神的話,出來傳道三年多.直到釘死在十字架.之所以如此,是因為他順服神,努力盡諸般的義.這也正是神喜悅他的地方.

第三,神喜悅耶穌謙卑自己,來做萬民的榜樣.耶穌說"因為我們理當盡諸般的義".他說的是我們,不是我.主耶穌這句話不是單指著他自己說的,也是指著我們說的.這句經文告訴我們,我們也需要信而受洗,同耶穌基督一樣,同死同復活,靠著他的榜樣向前走.

基督教與猶太教很不同的地方,是耶穌帶給我們的不是律法,是榜樣.榜樣不是規定我們一定要做這個,一定不能做那個；榜樣不是規定,榜樣是路標,是讓我們去學習的.中國人說"有樣學樣".就好像父母教育孩子,你規定他不能做這個,一定要做那個,他未必聽；但你自己是不是以身作則,你自己是不是抽煙喝酒講大話,孩子是一定會學習的.你就是你孩子的榜樣!孩子的性格行事有父母的痕跡,不僅僅是遺傳,也是他們跟你朝夕相處中學回來的.

耶穌基督就給了我們這樣的榜樣,他三十年無聲無息地學習,孝順父母,忍耐罪人,並愛他們.他雖然沒犯過任何罪,卻願意順服天父的律法,來到約翰面前受"悔改的洗禮"；他雖然擁有神的能力,卻處處謙卑自己,無能地死在十字架上.這就是他的榜樣.聖經就是一本有關榜樣的書,正如希伯來書說"亞伯拉罕因著信,在還不知道前路如何的情況下就離開了本家本族；因著信,在試驗的時候就把兒子以撒獻上；因著信,挪亞在未見洪水之前建了方舟；因著信,摩西的父母冒著生命危險收藏了摩西",更不用說因著信,瑪利亞就童女生子了.無論是先知,還是妓女喇合,無論是大衛王,還是摩押女子路得,聖經裡到處是這樣信心的榜樣.而耶穌基督,他盡諸般的義,就是要給我們最順服,最美好的榜樣.這也正是神喜悅他的地方.

三,    結論(生命的意義)講到這裡,我想再回到我們開始的問題:做人到底是為什麼?或者說,我們基督徒做人到底是為什麼?現在,基督給了我們答案:跟隨他的榜樣,盡諸般的義.   在世為人真不簡單.每天我們在辦公室忙碌,我們擠地鐵,我們做研究功課,我們給孩子做飯,我們看病,我們送別親友.我們探訪、關顧.這些瑣碎小事,天天在這個小小星球上上演著.日復一日,主耶穌還沒有再來.夜深人靜的時候,我們或許常常會想,這就是我們的一生麼?但借著聖經的話語,我們明白,我們在世為人辛苦所做的一切,婚喪嫁娶,生離死別,都是為了盡諸般的義；都是要像主耶穌一樣,謙卑順服,守人間的禮,盡諸般的義.總有一天,我們會如保羅所說:「

那美好的仗我已經打過了．當跑的路我已經跑盡了．所信的道我已經守住了．

」(提後4:7)那時天上必有聲音說:"你是所喜悅的."



\newpage

\section{你住在哪裡}
\label{sec:84}
\textbf{林豪恩先生 - 你住在哪裡?(Where are you staying?)}
\newline
\newline
\hyperref[sec:83]{< < < PREV SERMON < < <}
~
\hyperref[sec:toc]{[返目錄]}
~
\hyperref[sec:85]{> > > NEXT SERMON > > >}


約翰福音 1:29-42
\newline
\begin{longtable}{cl}
\hline
\hline
章節 & 經文 (和合本修訂版)\\
\hline
1:29 & \begin{tabularx}{0.7\textwidth}{X} 第二天，約翰看見耶穌來到他那裡，就說：「看哪，神的羔羊，除去世人的罪的！ \end{tabularx} \\ \\ \relax
1:30 & \begin{tabularx}{0.7\textwidth}{X} 這就是我曾說『那在我以後來的先於我，因為在我以前，他已經存在』的那一位。 \end{tabularx} \\ \\ \relax
1:31 & \begin{tabularx}{0.7\textwidth}{X} 我先前不認識他，如今我來用水施洗，為要使他顯明給以色列人。」 \end{tabularx} \\ \\ \relax
1:32 & \begin{tabularx}{0.7\textwidth}{X} 約翰又作見證說：「我曾看見聖靈彷彿鴿子從天降下，停留在他的身上。 \end{tabularx} \\ \\ \relax
1:33 & \begin{tabularx}{0.7\textwidth}{X} 我先前不認識他，可是那差我來用水施洗的對我說：『你看見聖靈降下來，停留在誰的身上，誰就是用聖靈施洗的。』 \end{tabularx} \\ \\ \relax
1:34 & \begin{tabularx}{0.7\textwidth}{X} 我看見了，所以作證：這一位是神的兒子。」 \end{tabularx} \\ \\ \relax
1:35 & \begin{tabularx}{0.7\textwidth}{X} 又過了一天，約翰同兩個門徒站在那裡。 \end{tabularx} \\ \\ \relax
1:36 & \begin{tabularx}{0.7\textwidth}{X} 他見耶穌走過，就說：「看哪，神的羔羊！」 \end{tabularx} \\ \\ \relax
1:37 & \begin{tabularx}{0.7\textwidth}{X} 兩個門徒聽見他的話，就跟從了耶穌。 \end{tabularx} \\ \\ \relax
1:38 & \begin{tabularx}{0.7\textwidth}{X} 耶穌轉過身來，看見他們跟著，就問他們說：「你們要甚麼？」他們對他說：「拉比，你在哪裡住？」（「拉比」翻出來就是老師。） \end{tabularx} \\ \\ \relax
1:39 & \begin{tabularx}{0.7\textwidth}{X} 耶穌說：「你們來看。」他們就去看他在哪裡住。這一天他們就跟他同住；那時大約是下午四點鐘。 \end{tabularx} \\ \\ \relax
1:40 & \begin{tabularx}{0.7\textwidth}{X} 聽了約翰的話而跟從耶穌的那兩個人，其中一個是西門‧彼得的弟弟安得烈。 \end{tabularx} \\ \\ \relax
1:41 & \begin{tabularx}{0.7\textwidth}{X} 他先找到自己的哥哥西門，對他說：「我們遇見彌賽亞了。」（「彌賽亞」翻出來就是基督。） \end{tabularx} \\ \\ \relax
1:42 & \begin{tabularx}{0.7\textwidth}{X} 於是安得烈領西門去見耶穌。耶穌看著他，說：「你是約翰的兒子西門，你要稱為磯法。」（「磯法」翻出來就是彼得。） \end{tabularx} \\ \\
[1ex]
\hline
\hline
\end{longtable}
日期:2011年2月6日

主題:你住在哪裡?(Where are you staying?)

証道:林豪恩先生

經文:約翰福音1:29-42

一．    引言

每年一度由香港電台主辦的「聲音檔案選舉」,挑選該年二十件社會及國際大事,讓市民投票選出其中十段.入選二零一零年二十大的其中一段聲音檔案顯示「你住在哪裡」這個問題的重要性.事緣青年協會邀請運輸及房屋局局長鄭汝樺與青年對話,中文大學的一位學生在對話中他向局長訴苦說:「因為我女朋友前排同我講,『你買唔到樓我唔嫁你,我唔住公屋』.」這段說話立即被傳媒廣泛報導,紅遍全港.「你住在哪裡」,不單單是「住」的問題,更與終生幸福掛鉤.

或者局長認為發言的同學是個別例子,但有調查證實「住在哪裡」其實是普遍年輕人的關注.智經研究中心在去年8至10月訪問了近1000名16至35歲的年輕人,了解他們的個人抱負和工作期望.調查發現,最多年輕人視買樓為未來10年的奮鬥目標.「住在哪裡」這個問題,形象化地想像一下,彷彿是打在馬兒背上的馬鞭,又是吊在馬兒面前的紅蘿蔔.

時下年輕人關注「住在哪裡」,剛巧我們今天誦讀的福音經課告訴我們,二千多年前的年輕人也關注「住在哪裡」,他們見到耶穌,第一個問題就是:「

拉比,在哪裡住?」(1:38)幸好他們沒有補充說:「我地唔住公屋架,如果你住公屋,我地唔跟你架.」門徒對於「住在哪裡」的關注,以及耶穌對他們的關注的回應,能夠帶給我們甚麼啟發及反省呢?讓我們細閱今天的福音經課.

二．    經文 29     次日,約翰看見耶穌來到他那裡,就說,看哪,　神的羔羊,除去〔或作背負〕世人罪孽的.     30     這就是我曾說,有一位在我以後來,反成了在我以前的．因他本來在我以前.     31     我先前不認識他．如今我來用水施洗,為要叫他顯明給以色列人.     32     約翰又作見證說,我曾看見聖靈彷彿鴿子,從天降下,住在他的身上.     33     我先前不認識他．只是那差我來用水施洗的,對我說,你看見聖靈降下來,住在誰的身上,誰就是用聖靈施洗的.     34     我看見了,就證明這是　神的兒子.     35     再次日,約翰同兩個門徒站在那裡.     36     他見耶穌行走,就說,看哪,這是　神的羔羊.     37     兩個門徒聽見他的話,就跟從了耶穌.     38     耶穌轉過身來,看見他們跟著,就問他們說,你們要甚麼.他們說,拉比,在那裡住.(拉比翻出來,就是夫子.)     39     耶穌說,你們來看.他們就去看他在那裡住,這一天便與他同住,那時約有申正了.     40     聽見約翰的話,跟從耶穌的那兩個人,一個是西門彼得的兄弟安得烈.     41     他先找著自己的哥哥西門,對他說,我們遇見彌賽亞了,(彌賽亞翻出來,就是基督)     42     於是領他去見耶穌.耶穌看著他說,你是約翰的兒子西門,〔約翰馬太十六章十七節稱約拿〕你要稱為磯法．(磯法翻出來,就是彼得.)

三．    釋經與應用1．「住」的空間向度那兩位本來跟隨施洗約翰的門徒聽到約翰介紹耶穌說:「

看哪,這是　神的羔羊!」他們就聽從師傅的介紹轉而跟隨耶穌.耶穌回頭看見有人跟著自己,很自然就問:「你們想要甚麼?」如果耶穌這樣問你,你會如何回答呢?根據剛才我們提及的那個「青年10年目標」調查的結果,最多人回答的是:「我想要一層樓!」；第二多人回答的是:「我想賺到第一桶金」；接著是:「我想要得到更高學歷」,「我想升職」,「我想有自己的生意」.這兩位門徒似乎不懂得珍惜這個難得的機會,提出自己的要求.不過,不立即接出要求,也可能是聰明的表現.他們以反問取代回答耶穌的問題,說:「你住在哪裡?」如果你是耶穌,你會如何回答呢?有人建議江湖味比較濃厚的劇情發展可以是這樣的,耶穌回答:「我住在哪裡與你何干?英雄莫問出處,英雄以四海為家.」門徒回應:「師傅,請不要介意我們解構你那瀟灑的說話,四海為家其實即是無處為家.我們要先知道你住在哪裡,才可以決定向你提出甚麼要求.我們對於住在半山區豪宅的人的期望,當然不同於住在深水Po 板間房的人.如果你住在深水Po 板間房,你能夠給予我們甚麼呢?如果你住在半山區豪宅,你從手指隙漏些出來,也能夠實現我要想一個小單位的期望.」

聖經學者們認為,當時的那兩位門徒沒有我們那麼深謀遠慮,也沒有我們那麼勢利,他們比較單純,他們問耶穌「你住在哪裡?」,一方面可能因為他們未想到如何簡單回答耶穌「你們想要甚麼」的問題,另一方面反映他們當下的需要.經文指出,「那時約有申正了」,即是一天將要結束,漫漫長夜快要來臨,「今夜宿誰家」,的確是一個他們當下適切的話題.

耶穌說,「

狐狸有洞,天空的飛鳥有窩」(路9:58)以此類比,人也應該有棲身的地方.棲身之所,可以是性命尤關的.每年寒流襲港期間,都會看到有露宿者凍死街頭的新聞.每次發生地震等災難,救災的行動之一,就是為災民提供臨時棲身的帳篷；災後重建的最主要項目,就是重建房屋.由此可見,居所作為人的基本需要是合情合理,被廣泛認同的.不過,我們的城市卻把房屋視為如黃金一樣的奢侈商品和投資產品,使房屋偏離了作為人類基本需要的本質及功用,令生活於這個城市的人為了一個棲身之所這基本的需要而感到極大的焦慮和壓力.

棲身之所本身是一個物理空間,這個物理空間如何變成奢侈商品及炒賣品呢?如此的改變,樓盤廣告應記一「功」.小弟從前就讀的中學位於何文田村,經常步行下土瓜灣馬頭圍道一帶吃飯,購物及乘車.在小弟的印象中,那裡是一個舊區.最近的某年某月某日,在媒體上看到一個座落於於馬頭圍道的樓盤廣告,令人眼前一亮.在廣告裡面,一男一女外國人依偎在一個會所泳池旁邊.女模特兒身材玲瓏浮凸,穿著性感泳衣,面帶笑容,作小鳥依人狀；而男模特兒則身材健碩,優哉悠哉地讓美人依傍.泳池鋪滿紅玫瑰花瓣,既浪漫,又瑰麗.在一片晴朗蔚藍的天空下,泳池旁邊有白色的太陽傘,桌子,椅子和樹木.泳池只有他們二人,彷如他們的私人空間,非常奢華.而廣告的標題「花得起,矜貴地」說明只要有錢購買這個樓盤的單位,便表示你的地位也矜貴起來.廣告更一再強調這個樓盤「精緻」,「時尚」,「優越」,「豪華」.樓盤廣告不斷地告訴人們,有樓就有伴侶,有樓就有幸福,有樓就有舒適生活,有樓就有尊貴地位...所以,為甚麼有人說:「你買唔到樓我唔嫁你」.在一次婚宴上,我問坐在旁邊的校友,何時結婚?他回答說:「尚未買樓,何談結婚?」可想而之地產商的宣傳策略何等有效.地產商把人們棲身的基本空間包裝成奢侈品,同時對這件奢侈品賦予「精緻」,「時尚」,「優越」,「豪華」等眾多的質數.當人人都被吸引,或者被動地投入這個遊戲的時候,誰人必定是贏家呢?福布斯香港富家排行榜已經把答案告訴大家了.

2．「住」的人際向度「住在哪裡」有其物理空間的向度,地產商以物理空間等同於眾多人們渴望的事物,然後暗示大家只需要追求這物理空間就得到一切.然而,「住在哪裡」除了物理空間的向度,還有人際關係的層面.門徒問:「拉比,你住在哪裡?」耶穌沒有回答問題,卻邀請他們去看.看過後,門徒就與耶穌同住了.到底耶穌住的地方有甚麼吸引門徒呢?會否因為坐落於地理位置極佳的「香港正中心」呢?會否因為是高居臨下的「八十八樓」呢?會否因為擁有六十一萬平方尺屋苑會所呢?只可惜作者沒有就此留下線索給我們去猜測.或者作者刻意不提耶穌住處的物理空間,似乎期望讀者留意更重要的地方.

到底門徒見到甚麼令他們決定與耶穌同住呢?到底門徒見到甚麼令他們不單只自己與耶穌同住,還介紹自己的兄弟去與耶穌同住呢?安得烈的說話透露了重點:「

我們已經遇見彌賽亞了」(1:41).作者暗示,門徒最後決定與耶穌同住,不是考慮「住所」的物理空間條件,而是考慮「同住的是誰」.俗語有云:「相見好,同住難」,安居之所的含意,不單單是一個物理空間,而且包括同一屋簷下的相處狀況.有網友留言:「老豆老母又嘈,左怨憤右粗口,乜都講晒出黎.老實講,真係極度影響我既心理,唔想結婚,唔想長大.」有網友回應,安慰說:「我屋企都係甘啦,唔止家嘈屋閉,仲打埋交添,扮下間歇性失聰啦.」好像這些網友的居住環境,經常「家嘈屋閉」,關係惡劣,縱使有一處棲身的空間,但與安居之所相距仍遠.

「住」這個字在約翰福音及約翰書信中出現次數相當多,除了在門徒的問題

「夫子,在哪裡住?」出現之外,以另一個層面的意義在32-33節出現.「約翰又作見證說:「

我曾看見聖靈,彷彿鴿子從天降下,住在他的身上.我先前不認識他,只是那差我來用水施洗的,對我說:『你看見聖靈降下來,住在誰的身上,誰就是用聖靈施洗的.』在這裡,我們看見「住」的人際關係層次.門徒發現他們的師傅是有聖靈同在的彌賽亞,因此,不論這人住在甚麼地方,都是「跟得過」的.約翰福音15章用了「住」這個字很多次,不過中文聖經對這個字有另一種翻譯.

4     你們要常在我裡面,我也常在你們裡面.枝子若不常在葡萄樹上,自己就不能結果子．你們若不常在我裡面,也是這樣.     5     我是葡萄樹,你們是枝子．常在我裡面的,我也常在他裡面,這人就多結果子．因為離了我,你們就不能作甚麼.     6     人若不常在我裡面,就像枝子丟在外面枯乾,人拾起來,扔在火裡燒了.     7     你們若常在我裡面,我的話也常在你們裡面,凡你們所願意的,祈求就給你們成就.     8     你們多結果子,我父就因此得榮耀,你們也就是我的門徒了.     9     我愛你們,正如父愛我一樣．你們要常在我的愛裡.(約15:4-9)「你住在哪裡?」,耶穌呼籲我們不一定要想著九龍西還是港島東,某某豪庭還是某某山莊,可以想想這樣回答:「我住在耶穌裡面」,或者:「我住在神的愛裡面」.若以這個角度去問自己:「我住在哪裡呢?」可以有更豐富的發現:有人住在「家嘈屋閉」中,有人住在「不抽不睬」中,有人住在「冷言冷語」中,有人住在「爭家產」中；也有人住在父慈子孝中,有人住在家庭溫暖中,有人住在手足互助中．．．「你住在哪裡呢?」這個問題不是能完全由住的空間層面去回答,如果誤信樓盤廣告,以為房屋代表了一切,於是耗盡一生的時間及精神去追逐那幾塊磚石,有可能得不償失呢!

3．「住」的意義向度「住」除了有物理空間的向度,以及人際關係的向度外,還反映著個人對人生意義的追求.有些人突破樓盤廣告宣傳的價值,不追求「精緻」,「時尚」,「優越」,「豪華」,對於「富甲半山」的引誘也不以為然,因為他們的人生另有追求.例如二千多年前的耶穌,他說:「狐狸有洞,天空的飛鳥有窩,只是人子沒有枕頭的地方.」(路九58)又例如劉曉波教授,為了國家的進步和同胞的福祉住進監獄之中.又例如華叔使徒華弟兄,中文大學政治及行政學系的蔡子強導師於華叔安息主懷後如此描述他的居所:「華叔清心寡欲,比較我們大多數人,生活簡單,因為,他信無欲則剛的道理.過去一年,我多次到過他旺角家中探望他,發現他家裡裝修簡單,大部份空間,由客廳到睡房,由床頭到飯台,大都用來放書,沒有多餘的擺設．．．其實,清苦與堅持,就是華叔的一體兩面.」(頭條日報2011-01-03)

「

狐狸有洞,天空的飛鳥有窩」,耶穌住在哪裡呢?耶穌說:「

我在父裡面,父在我裡面」(約14:10)耶穌住在天父裡面.住在天父裡面的耶穌無畏無懼實踐他的人生使命:「

傳福音給貧窮的人」,報告「

被擄的得釋放,瞎眼的得看見,叫那受壓制的得自由,宣告上帝悅納人的禧年」(路:4:18-19節).歷世歷代無數矢志實踐上帝使命的人都追隨耶穌,與耶穌同住.劉曉波住在哪裡呢?他住在諾貝爾和平獎的肯定裡面,住在世界各地追求公平,公義,人權的人們心中,住在期望國家進步及人民幸福的同胞心中.華叔住在哪裡呢?住在一群長期戰友的心中,住在眾多香港人－無論是政治上的左,中,右,為官的還是平民百姓心中.報章報導說:「在離開人世前的一段日子,華叔身體雖然痛楚,但有一班班的學生,戰友絡繹不絕探望,有黨友更親自到華叔故鄉,為他拍攝短片解鄉愁,一段段甜密的回憶,陪伴華叔走完最後一段路.連華叔的主診醫生莫樹錦教授也說:「我真的很慶幸能作為他的醫生!今日他的離開,我們相當傷痛,但我們知道他的精神會繼續與我們在一起.」(明報2011年1月3日)」

四．    總結人生在世數十載,我們的社會何必逼人民為了一個棲身之所而變成「房奴」呢?我們的社會能否把空間還給人民,讓人們不需要為了「居住」的物理空間而犧性「居住」的人際關係,也不需要埋沒人生的理想及意義呢?你住在哪裡呢?人生在世數十載之後,又住到哪裡去呢?華叔在離留之際,陪伴的家人朋友唱起了《安居主懷》這首詩歌,彷彿是「你住在哪裡」這個問題的最後,也是「至安心」的答案.

安穩在耶穌手中,安穩在主懷內；因主慈愛常覆翼,我靈在此得慰.無論遇何等危險,惟在主懷藏躲；有主聖臂常護庇,患難焉能害我!

安穩在耶穌手中,心平安無掛礙,在此不怕世誘惑,在此不被罪害.能脫離疑惑迷霧,脫離各般憂戚；世上惟餘幾患難,眼淚惟餘幾滴.

安穩在耶穌手中,耶穌曾為我死,主乃永生大磐石,我必永遠靠此.現今當忍耐等候,等過今世黑暗,等到見清晨曙光,照在榮華金岸.

副歌:安穩在耶穌手中,安穩在主懷內,因主慈愛常覆翼,我靈在此得慰.



\newpage

\section{無法抗拒的呼喚}
\label{sec:85}
\textbf{鄧瑞強博士 - 無法抗拒的呼喚(Irresistible Calling)}
\newline
\newline
\hyperref[sec:84]{< < < PREV SERMON < < <}
~
\hyperref[sec:toc]{[返目錄]}
~
\hyperref[sec:86]{> > > NEXT SERMON > > >}


馬太福音 4:12-23
\newline
\begin{longtable}{cl}
\hline
\hline
章節 & 經文 (和合本修訂版)\\
\hline
4:12 & \begin{tabularx}{0.7\textwidth}{X} 耶穌聽見約翰下了監，就退到加利利去； \end{tabularx} \\ \\ \relax
4:13 & \begin{tabularx}{0.7\textwidth}{X} 後來離開拿撒勒，往迦百農去，住在那裡。那地方靠海，在西布倫和拿弗他利地區。 \end{tabularx} \\ \\ \relax
4:14 & \begin{tabularx}{0.7\textwidth}{X} 這是要應驗以賽亞先知所說的話： \end{tabularx} \\ \\ \relax
4:15 & \begin{tabularx}{0.7\textwidth}{X} 「西布倫，拿弗他利，沿海的路，約旦河的東邊，外邦人的加利利— \end{tabularx} \\ \\ \relax
4:16 & \begin{tabularx}{0.7\textwidth}{X} 那坐在黑暗裡的百姓看見了大光；坐在死蔭之地的人有光照耀他們。」 \end{tabularx} \\ \\ \relax
4:17 & \begin{tabularx}{0.7\textwidth}{X} 從那時候，耶穌開始宣講，說：「你們要悔改！因為天國近了。」 \end{tabularx} \\ \\ \relax
4:18 & \begin{tabularx}{0.7\textwidth}{X} 耶穌沿著加利利海邊行走，看見兩兄弟，就是那叫彼得的西門和他弟弟安得烈，正往海裡撒網；他們本是打魚的。 \end{tabularx} \\ \\ \relax
4:19 & \begin{tabularx}{0.7\textwidth}{X} 耶穌對他們說：「來跟從我，我要叫你們得人如得魚一樣。」 \end{tabularx} \\ \\ \relax
4:20 & \begin{tabularx}{0.7\textwidth}{X} 他們立刻捨了網，跟從他。 \end{tabularx} \\ \\ \relax
4:21 & \begin{tabularx}{0.7\textwidth}{X} 耶穌從那裡往前走，看見另外兩兄弟，就是西庇太的兒子雅各和他弟弟約翰，同他們的父親西庇太在船上補網，耶穌就呼召他們。 \end{tabularx} \\ \\ \relax
4:22 & \begin{tabularx}{0.7\textwidth}{X} 他們立刻捨了船，辭別父親，跟從了耶穌。 \end{tabularx} \\ \\ \relax
4:23 & \begin{tabularx}{0.7\textwidth}{X} 耶穌走遍加利利，在各會堂裡教導人，宣講天國的福音，醫治百姓各樣的疾病。 \end{tabularx} \\ \\
[1ex]
\hline
\hline
\end{longtable}
日期:2011年2月13日

主題:無法抗拒的呼喚(Irresistible Calling)

証道:鄧瑞強博士

經文:馬太福音4:12-23

各位弟兄姊妹,主內平安.

一個故事:「一個漆黑的晚上,一個盲人提著一個燈籠在街上走.旁人好奇的問他:你是盲的,用得著燈籠嗎?盲人說:我當然用不著燈籠.這個燈籠是讓那些開眼的人看見我的.」誰是開眼的?誰是盲的?誰能看見人性的盲點?誰看不見人性的盲點?是否能看見人生的真相,如何能看見人生的真相,這常常是讓我們思索的問題.

今日講道的經文,是馬太福音4:12-23,這是主耶穌剛剛開始傳道時的一個福音記載. 12     耶穌聽見約翰下了監,就退到加利利去．     13     後又離開拿撒勒,往迦百農去,就住在那裡．那地方靠海,在西布倫和拿弗他利的邊界上.     14     這是要應驗先知以賽亞的話,     15     說,『西布倫地,拿弗他利地,就是沿海的路,約但河外,外邦人的加利利地.     16     那坐在黑暗裡的百姓,看見了大光,坐在死蔭之地的人,有光發現照著他們.』     17     從那時候耶穌就傳起道來,說,天國近了,你們應當悔改.     18     耶穌在加利利海邊行走,看見弟兄二人,就是那稱呼彼得的西門,和他兄弟安得烈,在海裡撒網．他們本是打魚的.     19     耶穌對他們說,來跟從我,我要叫你們得人如得魚一樣.     20     他們就立刻捨了網,跟從了他.     21     從那裡往前走,又看見弟兄二人,就是西庇太的兒子雅各,和他兄弟約翰,同他們的父親西庇太在船上補網．耶穌就招呼他們.     22     他們立刻捨了船,別了父親,跟從了耶穌.     23     耶穌走遍加利利,在各會堂裡教訓人,傳天國的福音,醫治百姓各樣的病症.主耶穌在加利利作一個周遊四方的傳道者,福音書的作者說,這是應驗了先知以賽亞的話:「

那坐在黑暗裡的百姓看見了大光；坐在死蔭之地的人有光發現照著他們.」這經文形容當地的人是「在黑暗裡的人」,是「在死亡的影子下的人」.光明對在黑暗中的人是很重要的.生命的意義對在死亡邊緣的人是很重要的.問題只是:人意識到自己在黑暗中嗎?人意識到自己在死亡邊緣嗎?每個時代都有每個時代的黑暗,每個時代都有令人失去生命意義的方式.我們活在一個消費主義的世界裡.有人說,「我消費,故我存在.」消費主義將一切都變成買賣.

我打一份工,是因為有人出了一個價,我覺得他的價出得好,我願意付出時間和勞力去換這筆錢.我們很難將這份工看成是神對我的呼喚(calling),我們很難想像這份工是神給我機會去榮耀他,去開發我的生命,去服事其他人.我有一個學生,曾在異地作過宣教服務.她說,宣教工作要適應文化差異,又沒有假期,天天起床後,便是工作了.在消費主義的思考習慣裡,沒有多少人願意作這種付出多而回報少的「買賣」.若你有兒有女,他要到偏遠地區作宣教士,你會贊成嗎?在消費文化裡,愛情也是一種買賣.女孩子說:「你要娶我,你要先買樓,我不住公屋.」這種講法真的很現實,一種現實的買賣.這句話的意思是說:我的價值高於公屋,你若只能付出公屋的價錢,那這買賣便告吹了.現代愛情關係衝突多多,是因為雙方都很計較.付出了半斤,便期望取回八兩.

甚至,連信仰神也變成一種買賣.我為神付出了這麼多代價,神也不祝福我.我每個禮拜都返崇拜,神也不祝福我個仔派到名校.我們很少想,為了神,為了生命的永恆價值,我不問代價,不計收成,完成那些有意義的事情.我們每個人都欣賞德蘭修女,而我們每個人都說只有她做得到她做的事情.

是否有一種無形的黑暗在我們生命裡頭呢?我們的生命意義是否在無聲無息地消逝呢?我們是否不意識自己在黑暗中呢?

主耶穌的呼聲,像是黑暗中的光明.黑暗中的光明是易見的.如:在家裡,吃完飯,通常的問題是:「誰去洗碗?」若有一天,吃完飯,有人主動說:「我去洗碗.」我們會發現在黑暗裡,有光照耀.主耶穌周遊傳道,說:「

天國近了,你們應當悔改!」這是神的聲音.這是光明所在.

你們要悔改了.悔改,就是整個人的全然改變.為何要全然改變呢?因為天國近了.天國是什麼?天國是人的心思明白神的心思,人願意降服在神的真理和引領下,以致神的計劃能在人的生命裡顯明出來.天國是神的價值迫近人而人必須抉擇的時刻.我是依然故我,抑或放下自己而降服神.

主耶穌走在人間.他顯明回應天父心意的生命是什麼樣的.主耶穌在日後,講了很多「天國」的比喻.其中一個是這樣的:一個葡萄園的主人清早請人做工,講好一天多少工資.早上九點的時候,又請了另一批人.十二點時及下午三點時,又請了一批.最後,下午五點時,他見到巿場上仍有人閒來無事地站著,於是又請了他們.最後,發工資時,先發最後來的,發給他們一天的工錢.最先來的工人以為會多領一點,怎知仍是同樣的一天的工錢.那些最先來的工人有些不滿.主人對他們說:難道我無權使用自己的錢嗎?為了我待人慷慨,你就嫉妒嗎?(參考馬太福音20: 1-16)

在這個故事裡,主人的做法最不合消費文化,也搞亂了整個公司秩序.看來不合理,但這只是不合我們的理,而合神的理.這個故事講的,是對弱者的憐憫.在人力巿場裡,到最後還沒有人請的,肯定是力有不逮的人.這個主人請他們,是讓這些無力的人被尊重.這個主人請他們,是讓這些人自力更生.這個主人請他們,不是給他們最低工資,而是給他們和精壯的人同樣的工資.再者,這主人教導那些早來的工人,不是從「我付出多便拿得多」來思考人生,而是從「每個人在世上都當彼此相助」的角度來思考人生.被恩典憐憫的人,不是你的敵人,而是你的鄰舍.在這世上,不單是「我沒有敵人」,而更是「每個人都是我的朋友」.

主耶穌見到的,是顛倒現實世界的另類秩序.這是天國,一個人必需改變整個生命才能看到的天國.你能否看到天國的寶貴?你能否在黑暗中看到光明?天國需要願意降服神的人去顯明,猶如燈需要由人放在燈台上.主耶穌呼召人作他的同伴,神需要人作他的同工.

主耶穌向人發出呼喚.馬太4:18-2318     耶穌在加利利海邊行走,看見弟兄二人,就是那稱呼彼得的西門,和他兄弟安得烈,在海裡撒網．他們本是打魚的.     19     耶穌對他們說,來跟從我,我要叫你們得人如得魚一樣.     20     他們就立刻捨了網,跟從了他.     21     從那裡往前走,又看見弟兄二人,就是西庇太的兒子雅各,和他兄弟約翰,同他們的父親西庇太在船上補網．耶穌就招呼他們.     22     他們立刻捨了船,別了父親,跟從了耶穌.     23     耶穌走遍加利利,在各會堂裡教訓人,傳天國的福音,醫治百姓各樣的病症.打魚,在當時,算是一份好職業.主耶穌呼喚他們,叫他們放下這些魚網,因為得到人的生命,彼得到魚更重要.馬太福音說,這幾個漁夫,「立刻」捨下一切,跟從主耶穌.他們意識到這是黑暗中的光明,他們意識到這光明的重要性,他們意識到追隨這光明的重大意義.他們意識到,除非是甘願留在黑暗中,否則應有所行動,去跟隨這光明行走.對他們而言,這是一個無法抗拒的呼喚.這猶如昔日的亞伯拉罕,放下一切,追隨神的聲音.一個近代的人物托爾斯泰(Leo Tolstoy),在82歲高齡時,在寒冷的冬夜,靜悄悄的離家出走,最後病死在一個火車站附近.火車站是一個富有意義的象徵,一個求道的人死在求道的路上.

很多人去思考為什麼這個老人家要放下一切離家出走.很明顯,這個老人家再不能忍受日常生活中的一切.他不能忍受富裕的生活方式,他不能忍受眾人世俗的期望.他內心聽到一種呼喚,呼喚他過一種另類的生活,這呼喚重要到一個地步,他忘記了自己的年紀,他忘記自己生命的脆弱,忘記了自己的一切.對他而言,這是一個無法抗拒的呼喚.他知道,若不聽從這呼喚,他將無意義地度過餘生.

海明威(Ernest Hemingway)有一篇短篇小說,叫＜屈力馬扎羅山的雪＞(The Snows of Kilimanjaro).小說的開場白是這樣的:「屈力馬扎羅山是一座海拔19,710呎的雪山,據說是非洲最高的山.它的西峰,馬賽語(the Masai)稱之為『神的居所』(the House of God).在接近西峰處,有一具枯乾而凍僵的豹子屍體.沒有人知道,這隻豹子在那高處究竟尋找什麼.」

劉再復寫了一篇散文,叫＜屈力馬扎羅山的豹子＞.他如此寫:「我相信這隻擁抱雪峰的豹子一定有一種人間智力還察覺不到的靈魂.牠是在尋找食物嗎?庸俗的眼睛大約會這樣看.牠是尋求丟失的同伴與兄弟嗎?如果是,牠是一種多麼有情的生命!但是,在山頂上怎麼會有像牠一樣勇敢的生命也走得那麼遠,值得牠如此獻身如此尋找呢?那麼牠尋覓無上的光榮與無上的地位嗎?牠也像人類那樣知道佔據頂峰是一種榮譽並且由此可以讓萬千同類抬頭仰望和俯首膜拜嗎?豹子恐怕沒有人類那麼複雜,牠的強大生命一定是單純的.」(引自劉再復:《我對命運這樣說》,頁51.)

之後,劉再復講到自己生命走向高處的飄泊,他說:「真的,我生命的一切現象都源自愛:我的沉思,我的寫作,我的苦痛,我的歡樂,我的告別,我的漂流,全都源於愛,源於我酷愛陽光下美的生命,酷愛洋溢著歌聲與故事的土地,山巒,河流和白雪.屈力馬扎羅山上的豹子,一定也酷愛這一切,一定酷愛雄奇的山巒與閃著銀輝的白雪.」(引自《我對命運這樣說》,頁52-53.)

劉再復沒有留意,這隻豹很大可能要走到「神的居所」那裡,要朝見神.在人生中,劉再復體會過世間的黑暗,也提示我們要努力尋求光明.

再引用他的文字,他在＜尋找的悲歌＞一文裡說:「我害怕黑洞.害怕文化世界中那些離奇莫測的黑洞.遙深的黑洞,佈滿著不會閃光的物質.那是死亡的本體.那是具有無窮吸引力的巨大深淵.」講完黑暗的力量,他講光明:「我曾見到一隻頂著夜的狂風與狂雨還在草間與樹間飛翔的螢火蟲.風雨打擊著她的小燈,企圖剝奪她的光明.但她捍衛著自己微弱的光明.依然一閃一閃地低飛著,依然堅韌地高擎著自己的生命之火,繼續著尋找的事業....我喜歡螢火蟲,儘管她身上的光明是弱小的,但她穿越著龐大的黑暗,並在黑暗中實現自己的發光發熱的靈魂.」(引自《我對命運這樣說》,頁127, 134-135.)

我在神學院教神學,有時問同學們為什麼走來神學院讀書.有同學說:「要來學習聖經.」很好.我們的確花很多時間學了很多別的東西,但很少真的研究聖經.有同學說:「在人生中,抽點時間出來,去認清自己的本相,去克服自己的罪性.」很好.我們的確花很多時間學了很多令我們在世界成功的東西,但很少花時間去克服自己的罪性.有同學說:「來學習服事神.」很好.我們的確花很多時間追逐這個世界的價值,但很少花心思去服事神.

這些同學都暫時放下事業,去回應神的呼召.你可能是一隻聆聽到上天呼喚的豹,要走到雪山之顛,在「神的居所」那裡朝見神.你可能是一隻聆聽到上天呼喚的螢火蟲,在風雨飄搖的日子,堅持發光.你可能是一個聆聽到主呼喚的人,呼喚你重新思索人生.呼喚你放下不重要的一切,學習過一種不一樣的生活.世界的黑暗常在,只怕我們意識不到.從天而來的光明常在,只怕我們不再願意去看.那呼喚我們放下黑暗走向光明的呼喚常在,你是否願意回應?「

天國近了,你們應當悔改.」

但願榮耀歸於聖父,聖子,聖靈.阿們.



\newpage

\section{與神同行奔天路}
\label{sec:86}
\textbf{高志棻博士 - 與神同行奔天路}
\newline
\newline
\hyperref[sec:85]{< < < PREV SERMON < < <}
~
\hyperref[sec:toc]{[返目錄]}
~
\hyperref[sec:87]{> > > NEXT SERMON > > >}


彌迦書 6:8
\newline
\begin{longtable}{cl}
\hline
\hline
章節 & 經文 (和合本修訂版)\\
\hline
6:8 & \begin{tabularx}{0.7\textwidth}{X} 世人哪，耶和華已指示你何為善。他向你所要的是甚麼呢？只要你行公義，好憐憫，存謙卑的心與你的神同行。 \end{tabularx} \\ \\
[1ex]
\hline
\hline
\end{longtable}
日期:2011年2月20日

題目:與神同行奔天路

經文:彌迦書6:8

證道:高志棻博士

我不曉得你們有沒有聽過這個小故事?許多年前英國倫敦時報出了一個有獎問答遊戲,題目是:從倫敦到羅馬,哪一條路線最短?

為了得到這個獎品,許多人於是就查看地圖來找答案,結果通通都落選了,但最後有一個人獲獎,他的答案你們可能猜不到,他說:有一個好朋友陪伴著你,一路說說笑笑,那從倫敦到羅馬哪裡會遙遠呢?豈不是一條最短的路線嗎?

所以弟兄姊妹,在你的人生道路上,你有沒有因為耶穌基督作你的朋友,陪伴你,一路同行,因此就充滿光明,喜樂和力量?剛才我們所讀的彌迦書六章八節:「世人哪,耶和華已指示你何為善.他向你所要的是甚麼呢.只要你行公義,好憐憫,存謙卑的心,與你的　神同行.」

這裡的說的「世人哪」,原文是「人哪」,是單數,不是多數──表示對我們每一個「個人」所說的話:張弟兄,王同學哪,你有沒有行公義,好憐憫,存謙卑的心,與你的　神同行?──所以「與　神同行」是一個命令,那同時也是一個邀請!

在這裡我要先澄清一個觀念:甚麼叫做「與　神同行」?

「行走」,在聖經裡是指我們的「行事為人,生活方式」.你我如何過生活呢?「與　神同行」就是與神採取同一個步調往前走,與神一同處理生活中每一件事情,這是一種生活方式.你信主得救是一回事,那你生活中有沒有與神同行,是另外一回事.若是我們的生活與神的步調一致,那你走起來就會充滿能力與平安.

先知彌迦的使命就是要教導神的百姓──「與　神同行」是非常重要；同時他也要警告百姓:若是你的心很剛硬,你選擇離棄神,不與神同行,那麼有一天,你可能會活出那敗壞的生命,導致上帝對你的管教.

為了要明白彌迦書六章一到十六節這段經文的教訓,我們需要理解當時的時代背景.先知彌迦的世代,大約在西元七百年前左右,當時以色列分為北國以色列及南國猶大,先知彌迦同時向北國,南國來傳遞上帝要對百姓說的話(當時的社會很像今天的香港非常繁榮,物質的條件特別的好,也很像今天的中國內地,百姓的眼目向前看).當時許多人很有錢,可是人就貪婪,貪心,不滿足,欺詐窮人,賄賂長官,好讓自己多賺一筆.感謝神!今天香港很努力維持一個廉潔的政府,但當時候的社會,真的是沒有公平,誠信可言,老百姓把上帝當作成是祝福他們的「工具」,他們並不敬畏神,所以先知彌迦住在耶路撒冷,他觀察整個社會,祭司,百姓的生活,他就非常的痛心大聲吶喊的說:以色列啊!你們已經偏離神,你們要回轉,要與神同行.

弟兄姊妹,你知道當你一信主以後,就走在一條通往天上的道路──天路.我們的一生就是要學習如何與神用同一個步調奔跑這一條天路.彌迦書的經文教導我們三個秘訣／原則,怎樣「與神同行」.若我們堅守這三個原則,我們就能成功的跑完這條天路.

彌迦書六章一節:「以色列人哪,當聽耶和華的話.要起來向山嶺爭辯,使岡陵聽你的話.」這裡怎麼會提到高山,丘陵呢?讓我們來想像一下:現在耶和華神,祂要跟祂的百姓爭辯,祂要責備以色列的不是,可是祂需要一些證人.所以耶和華就邀請一些人來做見證,來判斷祂跟以色列人的恩恩怨怨,而先知彌迦就如同一位律師.律師宣告證人入席,那誰是證人呢?在這裡就是高山,丘陵,大地,所以我們看到第二節接著說:「山嶺和地永久的根基阿,要聽耶和華爭辯的話.」高山,丘陵和大地都是證人,現在就站出來做見證,它們要見證些甚麼呢?它們要見證歷史.證人曾經看見過,上帝過去如何以大能拯救以色列,祂可以說是「救命恩人」,但沒有想到以色列竟然拋棄上帝,所以在第三節,神很激動地質問:「我的百姓阿,我向你作了甚麼呢.我在甚麼事上使你厭煩.」我看這句話會很感動,你看見嗎?我們的上帝非常激動,我們的神,祂不是道貌岸然,沒有感情.我們的神,祂的感情非常豐富,祂很傷心,而以色列是祂的心肝寶貝,祂很疼惜他們,拯救他們,但現在他們居然覺得厭煩,所以上帝就說:為什麼呢?我做錯了甚麼呢?你們有甚麼理由厭煩我呢?上帝遭受遺棄,祂心很痛,所以祂要求證人,快快地拿出歷史的證據,來證明上帝過去是如何恩待以色列.

所以弟兄姊妹,我們與神同行的第一個原則,就是要數算上帝的恩典.因此經文第五節就說:「我的百姓阿,你們當追念...」.「追念」甚麼?就是追念上帝在歷史中的救贖行動,數算上帝的恩典對以色列很重要,對我們基督徒更重要,當你遇見困難的時候,如果數算過去的恩典,就可以產生堅韌的力量,堅持與神同行,所以第四,五節就是記載上帝的恩典,過去以色列在埃及作奴隸,上帝如同老鷹一般,用大能的膀臂把以色列背在翅膀上,帶領他們飛到自由的土地上,毋須再做奴隸,還把領導人賜給他們,像摩西,亞倫,米利暗帶領他們走曠野；當以色列遇到仇敵要詛咒他們,上帝就出手干預化詛咒為祝福；當他們來到約旦河,上帝又差遣約書亞來引導他們化險為夷.所以一切的恩典,以色列不難數算,高山,丘陵,大地都能作證.你知道嗎?不僅當年如此,今天也要拯救,祝福你和我.你可能會遭遇到許多的困難,但是我們慈愛的天父要帶領你一一地去克服.常常數算神的恩典,培養一個感恩的心,會讓你更有信心,平安!

有一次,我參加了一個姐妹的聚會,有十位姐妹出席,而帶領我們的師母就發了九張紙條給每一個人,她就要求大家在每一張紙條上寫下一位姐妹的名字,然後如果你是她,你有甚麼值得感恩?最後師母就把所有的紙條和答案都收齊了之後,就高聲的朗讀.這很有趣,每一位姐妹在別人心目中總有一些情形令人羨慕,有的人會彈鋼琴,會管教孩子,有的特別會讀聖經,會燒菜,有的人的氣質特別迷人.我們所得到的祝福,總是覺得理所當然,可是在別人心目中,是那麼的值得羨慕,感恩,所以上帝對你特別恩寵的地方,你要常常記念.

二十多年前,我在美國加洲讀書,第一次參加查經班,我很清楚記得當初唱了一首詩歌叫《數主恩惠》,歌詞說:要數算主的恩惠,一一的數.那個時候,我不認識神,不明白上帝給了我甚麼恩惠呢?但二十多年下來,我就明白了.在我們成長環境,是很不容易相信耶穌基督的,我能夠得救,都是上帝的恩典,我信主以後,我的丈夫有六年半之久參加了大大小小的查經班,佈道會,培靈會達三百多次,但是因為他的科學知識和理性,他就拒絕相信耶穌基督.後來神的憐憫臨到他,神特別興起一個環境,讓他遇見神,接受耶穌基督作救主,那是非常奇妙的!今天我的丈夫在中文大學教書,研究的工作也蒙神的祝福,我們真的很感恩.我兒子小的時候,皮膚非常的過敏,非常嚴重,他每天晚上都發癢不能睡覺,連我也不能睡覺.但神安慰我,給我力量來陪伴孩子走過來,這是恩典.等到我兒子十三歲的時候,他開始叛逆,不想唸書,他整天無所事事,我常常為他流淚禱告.很奇妙,他十八歲的時候,他突然間又變了一個人,又開始願意讀書,喜歡聚會,學習事奉,這個也是恩惠,也是恩典.

常常數算過去的祝福,我們就不會懼怕未來,因為我們會明白生活中的點點滴滴,都是經歷恩典的機會.所以以色列人過除酵節,逾越節,住棚節都是為了記念神的拯救,都是在數算上帝的恩典.美國人過感恩節,基督徒過聖誕節,復和節,也是有著同樣的意義.

我們家庭有一個傳統,每年年底時都會舉行感恩聚會,我們一家四口,每人發一張聖誕卡,在卡上寫下當年感恩的地方,然後就交換分享.我也鼓勵你也找一個日子,特別的日子,全家來感恩.如果你的單身的話,你也可以找一個日子,譬如說你的生日,專門用那一天在你的日記上或跟你的朋友數算上帝的恩典,這樣你就會覺得,上帝就好像你多年的好朋友,有這麼深厚的友誼,怎麼可以輕易的拋棄呢?理當要珍惜這份友誼,繼續與神同行.

因此,數算神的恩典是與神同行的第一個秘訣.

當以色列聽了先知的勸告,他們怎麼樣回應呢?在彌迦書六章六節那裡,以色列人就覺得:「我朝見耶和華,在至高　神面前跪拜,當獻上甚麼呢.」意思就是說:先知你責備我,我不該拋棄上帝,上帝不開心了,那我給上帝甚麼禮物呢?該獻上甚麼樣的祭物才能搞定呢?如同一個人,他做錯了事情,他不想認錯,不想改正,只想:如果老闆對我不滿,那麼就送個禮吧!然後繼續他的錯誤.以色列人就是如此,他們不想悔改,他們只是想用祭物來賄賂神以遮掩他們的過錯.他們沒有留心上帝藉著先知要教導他們甚麼樣的功課,這是一個很大的諷刺,上帝不渴也不餓,祂不需要祭物,人若不聽上帝的話,不留心祂的教訓,怎麼能知道上帝的心意,來與神同行呢?

所以我們講到與神同行的第二個原則,就是要留心神的教訓,就是剛才唸的六章八節裡頭的教導:「世人哪,耶和華已指示你何為善.他向你所要的是甚麼呢.」這裡說神已經指示了祂的要求是甚麼,就是「行公義,好憐憫,存謙卑的心,與你的　神同行.」以色列人只要留心先知的話,就知道甚麼事情該做,甚麼不該做.舊約的時代,神的教訓已經透過先知來傳達；新約的時代,神的教訓很清楚,很明白的寫在聖經上,並透過聖靈向我們說話.上帝冀望我們基督徒要好好的研讀聖經.讀聖經可以明白神做事的法則,明白神的教訓,如此,你就會理解到有一些事情,上帝不要你做；有一些場合,上帝不要你去；有一些人,上帝不要你跟他們混在一起；有一些東西,上帝不要你有非份之想.所以常常讀聖經,留心上帝的教訓,你就知道如何與神用同一個步調來走天路.

在教會歷史裡馬丁路德很有名,他很重視上帝的話.十六世紀,1517年的時候,馬丁路德把一張告白貼在教堂的大門上(當時的教會,用金錢來賣贖罪券),他問:這合乎聖經,合乎真理嗎?他就列出95條來反對當時教會的作法,就引起軒然大波.四年以後,德國的皇帝就把馬丁路德叫來審訊他說:你要不要收回你的看法呢?馬丁路德就回答說:你讓我回去想一想.皇帝就答應了.

可能我們會想,馬丁路德是不是被皇帝的威嚴給嚇壞了?我看了一本書,他就評論這事,他說:其實不然!馬丁路德心裡是很明白的,坐在高高的寶座上是上帝,是天上的皇帝,有一天,地上的皇帝也要面對上帝,上帝說的話比地上皇帝說的更值得敬畏.所以經過整夜禱告以後,馬丁路德就站在皇帝的代表面前莊嚴地回答說:我的良心受聖經約束,被神的話掌管,所以我不能撤回我的意見.所以他講了一句非常嚴肅而有名的話:這是我的立場(Here I stand),阿們!因此馬丁路德的立場,就是聖經的立場,他這麼一喊「這是我的立場」,就震動了世界,改變了歷史.

其實馬丁路德也是一個平凡人,跟你我一樣,有許多的掙扎,他實在不願意反抗當時候的天主教教會,但是他要按照聖經的教訓與神同行.這一點點小小的堅持,就使他成為一個不平凡的人.所以我們與神同行,是要留心上帝的話語,注重神的教訓是非常重要的.

弟兄姊妹,你是否寶貴神的話語呢?我鼓勵你每一天至少要有一個固定的時段來讀聖經,如果你沒有這個習慣,那你可以從十五分鐘開始,慢慢加多至半小時,一小時,都很好,你也要學習順從聖經中的命令.我信主二十多年,經歷到很多次上帝藉著聖經跟我說話.

神的話真的是很活潑,真的是很有力量,那我們如果留心上帝的話,就能夠與神同行；相反的,如果我們不肯聽上帝的話,偏離上帝的道路,神因著愛我們,就會管教我們.所以與神同行來到第三個原則,就是要敬畏上帝的管教.

彌迦書六章九節:「耶和華向這城呼叫,智慧人必敬畏他的名．你們當聽是誰派定刑杖的懲罰.」你們要知道是神要懲罰,要審判,所以智慧人你們要敬畏神.所以六章九至十六節是宣告了上帝的審判.

我們的神的特別憐憫貧窮的人,祂要祂的子民也來憐憫貧窮的人.神說:「憐憫貧窮的,就是借給耶和華．他的善行,耶和華必償還.」(箴19:17)但彌迦書就記載,這些以色列的財主充滿強暴,說謊,犯罪,行事詭詐,欺壓窮人.有的時候,我們會讀到一些特別的報導,也會看見社會上窮人被欺壓,很不公平令人很憤怒也很痛心.以色列不數算上帝的恩典,也不留心神的教訓而欺壓窮人,那審判就臨到他們.六章十三節:「因此,我擊打你,使你的傷痕甚重」,十四節:「你要喫,卻喫不飽．」,十五節(新譯本的翻譯):「你要撒種,卻不得收割．搾橄欖油,卻不得油抹身．你要搾葡萄汁,卻不得酒喝．」因為上帝的審判,即使以色列努力工作,也是白費力氣,一無所得.最後十六節說,以色列啊,你還要受羞辱.所以犯罪不悔改,以色列就遭到上帝的審判與棄絕.

今天,我們身為基督徒,主耶穌已經在十字架上代替了我們受審判,因為上帝的公義,就不再定我們的罪,然而基督徒如果繼續叛逆神,不肯與神同行,那聖潔的神因為愛我們,就會管教我們,如同父母親管教孩子一樣,管教是出自神的愛,要引導我們回頭.以色列受審判,基督徒受管教.

看見以色列受審判,讓我們明白上帝是聖潔的,使我們產生敬畏的心,就不敢做惡事了,若是跌倒了,蒙神管教,我們就要快快回頭.基督徒與神同行時,第三個原則就是要敬畏神的管教.

現在大家很留心一本書,叫做《標竿人生》,它的作者是華理克牧師 (Rick Warren),有一次我到他在加洲的教會裡去聚會,我聽到他說:過去,我曾經在男女的關係上受到誘惑,但是我害怕神,我害怕祂的管教,我敬畏上帝,我不敢放肆,我只愛一個女人,就是我的妻子,我要跟她白頭偕老.因為敬畏上帝的管教,就讓華理克牧師敵擋了誘惑.

有時候想想很奇妙:這位創造宇宙的神,祂竟然會掛念你和我這種小人物是不是聽祂的話?是不是誠實?做生意的時候有沒有欺騙客戶?如此想一想,我會覺得很溫暖,你看六章十到十一節:「惡人家中不仍有非義之財,和可惡的小升斗麼.」一方面這是出於神的聖潔公義,要商人誠實的做買賣,另一方面,這也是出於神的愛.父母親的愛,會讓父母注意孩子的一舉一動,好管教他,塑造他的品格.所以我們的神也是如此,祂很仔細地觀察我們,如果我們基督徒偷偷摸摸,用一些不正當的方法做生意,我們偷偷摸摸對我們的配偶不忠實,那你等著看,出於上帝的聖潔與慈愛,祂一定會管教的.

有一次,我在美國的華人教會聚會,在教會用餐的地方裡,擺了一個魚缸,裡面的熱帶魚很漂亮,我覺得好稀奇,因為我從來沒有看見教會放魚缸,原來它背後有一個故事:魚缸的主人,年青的時候就相信耶穌基督,可是等他到了美國,他的生意越做越大,他就漸漸遠離神,金錢成為他的神,事業成為他的偶像,他很貪婪,很揮霍,很不可一世,很驕傲,他把上帝置之腦後,當時他擁有五幢公寓,他自己有一幢很大的房子,但因為他很喜歡熱帶魚,家裡就擺著各式各樣,大大小小的魚缸,擺在教堂裡面的,只是中小型的魚缸.但是因著神愛他,就管教他,就興起一個環境,結果他的生意完全失敗,他那五幢公寓在一夜之間要全部賣掉,好拿來還債,而他的大房子也換成了小房子,所有的魚缸都賣掉,只剩下兩個,一個擺在家裡,一個擺在教堂.這位弟兄受到管教的時候,他心裡明白,這是出自神的手,他就產生一個敬畏的心回轉歸向神.而且從那以後,他就很熱心傳揚福音.所以魚缸就是一個見證,見證著上帝的信實,慈愛；魚缸也是一個教訓,提醒大家要敬畏神,否則會受到管教.

有一句話對於與神同行的罪人很重要,我要請大家仔細聽:上帝祂絕對不會放過你的,上帝也絕對不會放棄你.

上帝不會不會放過與祂同行的人:你我犯罪,上帝一定會管教.但上帝的管教是為了愛,要帶祂的兒女走回正途,所以祂也絕對不會放棄你和我.希伯來書十二章六節:「主所愛的他必管教.」就是這個意思.

因此以色列行惡,受到審判,可是到了彌迦書第七章八至九節:「我雖跌倒,卻要起來．我雖坐在黑暗裡,耶和華卻作我的光.我要忍受耶和華的惱怒．因我得罪了他．直等他為我辨屈,為我伸冤．他必領我到光明中,我必得見他的公義.」

你看到嗎?與神同行是有盼望的,即使我們受到管教,只要敬畏祂,回轉歸向祂,繼續相信上帝的信實與慈愛,就會得到拯救.

所以親愛的弟兄姊妹,神在古時候呼召以色列與神同行,今天祂也同樣地呼召你和我要來與神同行,先知彌迦教導百姓要記念上帝在歷史中的救贖作為──那我們基督徒與神同行的時候要數算神的恩典；先知彌迦提醒百姓要留心神已經有的指示,活出神的要求──基督徒的領受,就是要留心聖經的教訓；彌迦又宣告因為百姓不悔改,要面臨神的審判──基督徒的領受,就是要敬畏神的管教.

因此數算神的恩典,留心神的教訓,敬畏神的管教──這是可以每一天都來操練的.你晚上睡覺前可以向神禱告,說「謝謝」,「請」,「對不起」:

1)   「謝謝」──謝謝神是數算神一天的恩典,為今天發生的事情,沒有產生的事情來感恩；2)   「請」──留心上帝的教訓,請祂幫助我來順服聖經的教導；3)   「對不起」──順服與敬畏上帝當天的管教,為我們一天的過犯向神悔改.

「謝謝」,「請」,「對不起」就是用來數算上帝的恩典,留心上帝的教訓,敬畏上帝的管教,如此,我們就可以天天與神同行.

願上帝祝福各位!



\newpage

\section{和睦是福}
\label{sec:87}
\textbf{陳榮超牧師 - 和睦是福講員}
\newline
\newline
\hyperref[sec:86]{< < < PREV SERMON < < <}
~
\hyperref[sec:toc]{[返目錄]}
~
\hyperref[sec:88]{> > > NEXT SERMON > > >}


馬太福音 5:9
\newline
\begin{longtable}{cl}
\hline
\hline
章節 & 經文 (和合本修訂版)\\
\hline
5:9 & \begin{tabularx}{0.7\textwidth}{X} 締造和平的人有福了！因為他們必稱為神的兒子。 \end{tabularx} \\ \\
[1ex]
\hline
\hline
\end{longtable}
日期:2011年2月27日

主題:和睦是福

講員:陳榮超牧師

經文:馬太福音5:9

耶穌基督登山寶訓裡的其中一「福」 -- 「使人和睦的人有福了．因為他們必稱為　神的兒子」.和合本聖經所寫的「兒子」,我相信無論男女都包括在內.當我看見人們不斷祈福時,背後總帶著一種不安或對前途的信心,希望能求到一個好彩頭,那之後的三百六十四日都不用去擔心,所以今年劉皇發先生為香港求了一支中簽,但中簽就有不同的人從不同角度解簽.無論怎樣解也好,相信人心裡面祈福,是因為對心裡面的福不太安心,希望能多求一個好福,這樣我們就可以心安理得去做事了,就有如拿了一顆安心丸後,我們不需再為這個福思念太多.

於暴亂當中,我們見到人的心對於這個福,安定,安穩,平靜的祈願,每天有百萬香港人希望這個國家能有好的前途,我們就知道人心裡面實在需要有這一樣東西的.但是我們作為基督徒的,從聖經上學習,就該知道我們是有從神而來的「福」,所以我們要去祈福,求福.

我們不但需要祈福,求福,我們還需要做「福」──這是我們與其他信徒在行動和思念中有少許不一樣的地方.我們不是驕傲或高人一等地說,而是在我們的信仰,實際的行為上去做,我們應該將已經得到的福,在人群當中去造「福」人群.

我喜歡聖法蘭西斯的一個禱文,我曾聽王祖藍唱這首歌.我相信把歌詞讀出時,大家都會哼這首歌的調子和意思.『使我作你和平之子』 -- 聖法蘭西斯的禱文中文的翻譯是,「神啊!讓我作你和平之子,在怨恨中,使用你的愛；在憂傷之中,傳送你的寬恕；在懷疑之中,顯出你的信實.使我作你和平之子,失望之處,帶出盼望；在罪惡黑暗,發出你的光；在難過心境,放下你的喜樂.使我作你和平之子!」重覆又重覆,我們要造福於人群的信徒,使我們能夠為神當一個和平的使者,願寬恕別人好像你寬恕,願施予別人好像你施予,願燃燒生命得永遠生命.副歌是:幫助我,不要單單求安慰,更願意努力安慰幫助人.

我們不需要求,我們是一群有福的人群,我們的責任是要把這個福幫助人,要努力做安慰,做平安,造福人群,「不求被了解,但求了解人」,接受愛,更加願意付出愛──法蘭西斯的禱文:使我們作和平之子.

各位弟兄姊妹,朋友,是我們在基督裡的一個本份,一個使命,不只是我一個人的,乃是整個教會的,要將這個和睦,和平帶給眾人,而教會在我們這個社會中是一個使人蒙福的地方.所以耶穌在八福當中,有很多是針對我們的態度,但這福是要我們有行動的.清心的人有福,是關於我們對神的態度；但使人和睦,不是我們的態度,而是行動上的一個要求.使人和睦的人是很幸運,「福」在耶穌年代,是代表很幸運的人,這是有福的原因.

為何我們這樣幸運呢?為何我們今天這樣有福?因為我們懂得使人和睦；為何使人和睦是多麼幸運呢?因為我們可以被稱為神的兒女.我知道有些人常常都很幸運,笑面迎人,猶如一棟大廈在他旁邊塌下來,也與他無關,可以很開心.

有些人天生是開朗的,但並不是每個人生出來都是開朗的.耶穌這句說話是不理會你是開朗或低潮的人,不理會你是個怎樣情緒的人,你都是很幸運,很有福的.因為你有一份工作要做,你要完成這份工作──使人和睦.將和睦,平安帶給別人,跟我們的情緒無關,你是一個很幸運,很開心,很開懷的人,因為神在這裡祝福你.

我們看看這個要緊的『福』,八福當中,這個福是充滿活力的,能夠使我們作為信徒.在教會當中的弟兄姊妹,甚至前來慕道的朋友,我們可以做一個這麼幸運的人.這個幸運並不是我們中了彩票,原來神已給予我們開心的本質,我們要使用出來.

我們來摸索一下耶穌基督,思想一下這段經文對我們的挑戰:照祂所說,我們怎樣可以成為有福的人呢?我們怎樣成為使人和睦的使者?因為造平安,造福人群的人,必被稱為神的兒女.因為我們被稱為為神的兒女,所以使人和睦是關於人際關係的問題.

因為有人的地方,就可能有磨擦.人際關係是講關係,但今天一講人際關係,就出現許多問題.人際關係就是人與人之間的相處,但很多時我們會很自然地或在潛意識裡,填寫了下半句:人際關係就是人際問題,在我們的潛意識中人與人之間似乎好像帶來一點爭吵.

人與動物沒有這麼多的爭吵,回家不開心,把怨氣發洩於狗的身上,牠也不會吠你,縱然會,但與你之間也沒甚麼好爭執的.人與動物是沒有爭執的,這是一種特殊的關係.人與人之間有很多誤會,不了解,甚至帶出很多恩恩怨怨,問題的關鍵出於在「關係」.

第一,不知自己是誰.當我們處理不好的話,往往在和睦相處中,就會產生一些人際上的爭執,因為我們不認識自己是誰?我有姓名,知道自己住在那裡；為何不知自己是誰呢?其實我們沒有一個人知道自己是誰.你知道自己的樣子是怎樣子的嗎?你沒有沒看過自己的樣子,你說看過了,其實你只看過鏡中的你.我在美國教會有個姊妹是失明的,她很幽默,每次我跟她打招呼的時候,她就會說:「靚仔牧師來找我了.」我回應,你是失明,怎能看見我是靚仔呢?

你的樣子是別人告訴你,你自己是看不到的；你看不到自己的樣子,雙眼是怎樣,你是看不到的.你看見自己的樣子,是因為有朋友為你拍照,然後靠著你的記憶,大腦神經幫你拼圖,把樣子拼回出來；神經錯亂的人,失去拼圖能力,所以可以拼出很多樣子.我們是不知自己是誰,如果我們沒有一個好的人際關係,對自己的真正形象就會失去,我們憑著自己是不知自己是誰的.好像一個很簡單的問題,但仔細想想卻感到很複雜.

人與人之間建立關係時,首先要問問自己是誰,我的身份,境況如何?誰最了解我呢?了解我的人能告訴我,我是誰.了解我的人,不但能說得出我的樣子,還有心理,情緒,性情等等.從我出生到十二歲,我父母是最了解我的,但今天我父母已經不了解我了.今天我父母想要了解我,就要問我太太,我的兒子是個怎樣的人呢?因為我父母了解的我,是從出生到十二歲；十二歲到十八歲就沒有人了解我；十六,七歲時常發脾氣和生氣,是一個全世界沒人了解的年青小伙子,特別是失戀的人,所以我們需要一個很好的人際關係.

在青少年,讀大學時,我的同學了解我；拍拖時,我的另一半了解我；當另一半變成配偶時,這十年來我的太太了解我最深.你只不過是認識我的表面,聽到我的聲音,看到我的笑容... 除此之外,我的弱點,生活上的細節,你不了解,而了解我的人,是我的太太.

今天,在你的生命裡,你對創造主的不認識,你就不會認識自己的真我.因為若我們不認識自己,在人際關係中便會成為絆腳石.當你不明白自己的時候,你就會在社團,社區,家庭中成為爭執的絆腳石.但我們要成為傳達和睦訊息的基督徒,你要在人群,社團中製造和平,如果你成為人群中的絆腳石,就要很小心地去處理.要回到神那裡,讓祂幫助我們了解「我是誰?」.如果神不告訴我們是誰,我們對自己的認識仍然不足.今天縱使我的太太認識百分之九十的我,但餘下的百分之十,神若不告訴我,我還對自己的認識就打了一個很大的折扣.

這樣,我怎樣可以使人和睦?如果我與神沒有和好,我怎能與人有和睦呢?如果我得不到神完全的接受,我怎可以接受別人呢?今天我們要成為一個使人和睦的使者的話,首先我們要回到神的面前,求神接納,讓祂告訴你「我是誰?」,所以,和睦的使者是從神開始.

耶穌說:「使人和睦的人有福了」,因為祂是神的兒子,你是神的兒女,你是神所愛的兒女.昔日你可能是浪子,但今天你回到天父的家裡,在祂家中是一位兒女.

做和睦的工作是要主動地去做.在八福中這福是充滿活力的,因為是要求你去做一些事.在八福中,你要這福成為你的福,你就要來採取行動,要在眾人中成為一個和睦的人.使人和睦是一個字,英文叫「peacemaker」,原文是一個組合性的字,叫「平安」.使人和睦是製造和睦的人,在我們當中做和平者.

南非一位退休總統,是第一位黑人總統 -曼德拉,在一九九三年,他獲得諾貝爾和平獎.獲獎後,他於一九九四年至九九年成為南非第一位黑人總統.但當他未成為總統之前,被認為是一個「反動份子」,且被南非的白人囚在牢裡幾十年.去年退休時九十二歲,那時南非正值舉辦世界盃足球賽,他健康不好但仍掙扎出席世界盃賽並擔任頒獎嘉賓.前陣子,還把他其中一段的生平製作成電影《不敗雄心》,就是介紹他剛成為總統時的故事:那年南非舉辦英式橄欖球大賽,南非球隊差不多是最差的一隊,但曼德拉在比賽中,親自鼓勵全白人的橄欖球球員.當時候的內閣成員跟他說:「你剛成為總統,黑白歧視正在爭吵中.」曼德拉回應:「我現在成為總統,看不見黑與白的分別.」因為曼德拉堅持的態度,心態,思想,透過橄欖大賽慢慢一步一步的推前,親自來到白人管轄的地方.這套電影主要將他的心境,對黑白種族的歧視,態度和做法表達出來.南非的橄欖球隊在大賽中,竟然得到冠軍成為電影的高潮,曼德拉在大賽中走到頒獎台前,親自把獎狀從黑人手中頒發給白人的隊長,這一幕帶來使人感動的畫面.

曼德拉的一生中,因為種族歧視,仇恨,爭取黑白平等,被囚禁幾十年... 他最好的青春,最有力的日子,都在牢獄中渡過.一般人遇到相同的境況,總是會充滿著仇恨,盤算著出獄後的報復,革命,煽動等等,為的就是要報復那些曾經傷害過自己的人.

曼德拉在牢獄的幾十年裡,他跟自己說:「如果我是一個不能饒恕別人的人,我會白費了坐牢的日子」.他離開牢獄的那天,他步行離開,當牢獄還在他背後時,他心裡頭就有一個聲音對他說:「如果你不能將悲痛和怨恨留在背後,這樣你仍是坐在牢獄當中.」後來他寫傳記時,這句說話就成了一句名言.

曼德拉是循道公會派的一個信徒,雖然他平日在公開場合沒有提及信仰,但在自己的傳記中說得很清楚,『基督在我心中,從來沒有離開過.』

因為他的態度,南非現今成為和平共處的種子.所以和睦的工作,是要有人在主裡面找到平靜,然後在人群當中,作為一個撒種的人.我們要將和平的種子種於不同的地方,種在人群中,那並不是一朝一夕就會看見果效；但如果是種子的話,和平的種子是會發芽生長的.

另一方面,我們使人和睦的工作,是要將神的智慧帶出來.雅各書三章17 -- 18節:「惟獨從上頭來的智慧,先是清潔,後是和平,溫良,柔順,滿有憐憫,多結善果,沒有偏見,沒有假冒；並且使人和平的,是用和平所栽種的義果.」並且要人際關係中,幫助爭執的人,我們要把每件事的正面的訊息帶給他們.

為何我們會起爭論呢?因為我們所看見的,所解釋的皆是負面,每個事實都是透過你的解釋,如果不從正面去解釋的話,你便成為使人和睦的絆腳石.當你能在每件事上,從正面的方向去解釋,你便能成為和睦的使者,所以我們要不斷地來到神面前承認自己的錯失,失誤.

在埃及事件中,其中一位埃及牧師- 沙瑪說:兩堂教堂旁的建築物已燒得平平滑滑,且倒塌了,但教堂仍完好無缺；教會裡的信徒們在那十天當中,分為不同的小組,到混亂的市區裡去清理垃圾,治理衛生,把瓦礫推開使道路暢通.另外一些小組,到傷者家中,為他們包紮傷口；到貧窮人的地方,把食物帶給他們,幫助他們；在受傷人當中,服侍他們.在一片混亂,前途未卜的境況底下,信徒們在當中成為和平的使者,在人心茫茫的時刻,種下和平的種子.

這就是耶穌基督在八福中,給予我們很大的勸勉和挑戰,我們要出去,在人際的關係中扮演一個角色,做一個實際的工作,就是成為一個「使人和睦的信徒」:使人和睦的人有福了,因為你必被稱為神的兒女.願上主祝福我們眾人.

讓我們一同低頭禱告:天父上帝願我們發出的話語,不會空空盪盪的回到你當中,讓我們成為有福的,使人和睦的人,這福是已經賜予給我們了,就讓我們實實在在的將你的福帶到人群當中,作一個使人和平,和睦的使者,因為我們是你的兒女,這是我們的本分,奉主名求,阿們.



\newpage

\section{西乃山下與獅子山下}
\label{sec:88}
\textbf{林豪恩先生 - 西乃山下與獅子山下}
\newline
\newline
\hyperref[sec:87]{< < < PREV SERMON < < <}
~
\hyperref[sec:toc]{[返目錄]}
~
\hyperref[sec:89]{> > > NEXT SERMON > > >}


利未記 19:9-18
\newline
\begin{longtable}{cl}
\hline
\hline
章節 & 經文 (和合本修訂版)\\
\hline
19:9 & \begin{tabularx}{0.7\textwidth}{X} 「你們在自己的地收割莊稼時，不可割盡田的角落，也不可拾取莊稼所掉落的。 \end{tabularx} \\ \\ \relax
19:10 & \begin{tabularx}{0.7\textwidth}{X} 不可摘盡葡萄園的葡萄，也不可拾取葡萄園中掉落的葡萄，要把它們留給窮人和寄居的。我是耶和華－你們的神。 \end{tabularx} \\ \\ \relax
19:11 & \begin{tabularx}{0.7\textwidth}{X} 「你們不可偷盜，不可欺騙，也不可彼此說謊。 \end{tabularx} \\ \\ \relax
19:12 & \begin{tabularx}{0.7\textwidth}{X} 不可指著我的名起假誓，褻瀆你神的名。我是耶和華。 \end{tabularx} \\ \\ \relax
19:13 & \begin{tabularx}{0.7\textwidth}{X} 「不可欺壓你的鄰舍，也不可偷盜。雇工的工錢不可在你那裡過夜，留到早晨。 \end{tabularx} \\ \\ \relax
19:14 & \begin{tabularx}{0.7\textwidth}{X} 不可咒罵聾子，也不可將絆腳石放在盲人面前。你要敬畏你的神。我是耶和華。 \end{tabularx} \\ \\ \relax
19:15 & \begin{tabularx}{0.7\textwidth}{X} 「你們審判的時候，不可不公正；不可偏護貧窮人，也不可看重有權勢人的臉，總要公平審判你的鄰舍。 \end{tabularx} \\ \\ \relax
19:16 & \begin{tabularx}{0.7\textwidth}{X} 不可在百姓中到處搬弄是非，不可陷害鄰舍的性命。我是耶和華。 \end{tabularx} \\ \\ \relax
19:17 & \begin{tabularx}{0.7\textwidth}{X} 「不可心裡恨你的弟兄；要指摘你的鄰舍，免得因他承擔罪過。 \end{tabularx} \\ \\ \relax
19:18 & \begin{tabularx}{0.7\textwidth}{X} 不可報仇，也不可埋怨你本國的子民。你要愛鄰如己。我是耶和華。 \end{tabularx} \\ \\
[1ex]
\hline
\hline
\end{longtable}
日期:2011年3月6日

題目:西乃山下與獅子山下

證道:林豪恩先生

經文:利未記19章9-18節

一．引言三千多年前,有一大群以色列人浩浩蕩蕩地離開埃及.他們進入曠野,來到西乃山下,暫時安營.作為一個數目龐大的群體,應該如何一起生活呢?人與人之間應該怎樣彼此相待呢?西乃曠野只是中途站,到達目的地之後,這群人要建立一個怎樣的社會呢?當以色列人暫時停留於西乃山下的時候,他們尋索自己的身份,確立自己的信仰,以及探討將來要建立怎樣的人際社會.在尋索,探討,確立的過程中,這個群體經歷挑戰及掙扎,有時是天災,有時是人禍,有時是內憂,有時是外患,有時是由吃飯喝水等日常事件演變而成的問題.

以色列人在西乃山下的故事,令人聯想到我們在獅子山下的故事.黃霑先生透過《獅子山下》一曲,如此道出獅子山下的故事.

人生中有歡喜,難免亦常有淚,我地大家在獅子山下相遇上,總算是歡笑多於唏噓.人生不免崎嶇,難以絕無掛慮,既是同舟在獅子山下且共濟,拋棄區分求共對.放開彼此心中矛盾,理想一起去追,同舟人,誓相隨,無畏更無懼.同處海角天邊,攜手踏平崎嶇,我地大家用艱辛努力寫下那,不朽香江名句.

獅子山下的人們也發出類似的問題:作為一個群體,我們應該如何一起生活呢?我們要建立一個怎樣的社會呢?《獅子山下》這首歌唱了那麼多年,人群間的「唏噓」似乎不比「歡笑」少,社會上「矛盾」的聲音似乎比「共濟」還要響亮.《獅子山下》一次又一次地重唱,我們的社會一年又一年地在爭論,到底我們應該如何彼此相待呢?到底我們盼望一個怎樣的社群呢?

我們今天誦讀的舊約經課,記載了上帝就著這些問題給西乃山下的以色列人一些指引.三千多年後,獅子山下的人們重讀,會否也是寶貴的資源,合時的提醒呢?二．釋經與應用利未記19章9-18節論及以色列人作為一個重新建立的群體,在人際社會的層面上,應該持守怎樣的價值,以及如何彼此相待.這段經文可以分為五個段落,記錄了五組指引,每個段落以「我是耶和華」作為劃分.

1．第一組指引:讓大家都有謀生的機會19:9 「在你們的地收割莊稼,不可割盡田角,也不可拾取所遺落的.19:10 不可摘盡葡萄園的果子,也不可拾取葡萄園所掉的果子；要留給窮人和寄居的.我是耶和華你─們的　神.」

在古代的農業社會,土地是人們賴以為生的資源.中國唐代詩人李紳如此描述農民的生活狀況:

「鋤禾日當午,汗滴禾下土.誰知盤中飧,粒粒皆辛苦.」(《憫農二》)

對於古代農業社會的人們來說,雖然「粒粒皆辛苦」,但有田地可作耕種仍是幸運的,因為這是「盤中飧」的唯一來源.

在一個以士地為生產資源的農業社會,存在著一些沒有士地的「窮人和寄居者」,或者可以稱為「無產者」,這個社會應該如何對待他們呢?上帝給予西乃山下的人民的指引有值得我們留意的地方:祂不是要求莊園園主收割莊稼後賙濟窮人,也不是要求果園的園主摘盡果子後施捨給寄居者,而是要求他們在收割時「不用那麼乾淨」,讓窮人和寄居者去執拾收割後的「漏網之魚」來維生.捕捉漏網之魚,也可以稱為工作.看來,上帝認為窮人和寄居者這些沒有土地的人,透過工作自力更生,總比失業而接受救濟更為理想,因為靠自己雙手養活自己,能夠讓人體會到作為人的尊嚴及價值.

今日獅子山下的經濟模式與當日西乃山下的農業社會相距甚遠,不過,相同的是,兩個社會都存在著「有產者」及「無產者」.上帝指示西乃山下的社會,給予所有人謀生的機會,製造工作機會給「無產者」,讓他們也能體會自力更生的尊嚴,這樣的呼籲對今日獅子山下的社會仍然適用嗎?

近年來,有越來越來聲音指出,獅子山下的「莊園園主」收割得太「乾淨」,甚至連「所遺落的也要拾取」,「果園園主」摘葡萄摘得太盡,甚至連掉下來的也不放過.近期某集《星期日檔案》介紹一位「向地產商說不」的市民,透過他的嘗試,證實生活在獅子山下,若要向地產商說不,可以說是難以生存的.地產商可能偷取了教育界「全人教育」,或者醫護社福界「全人關顧」的概念,也對獅子山下的人們實行「全人照顧」,提供「一條龍服務」:除了「照顧」人們「住」的需要,還「照顧」人們「衣,食,行」,甚至醫療的需要.在這種資本快速擴張,而且高度集中的過程中,獅子山下的生存環境改變了.我讀小學時,有同學的父母推著小販車養活一家；我讀中學時,有同學的爸爸在當看更的工廠大廈樓梯底賣汽水和小食供養兩個孩子完成學業；又有一位學師做五金的長輩,曾經自立門戶在政府的工廠大廈靠自己一雙手謀生,可惜這些個體戶都被淘汰得七七八八.人們被排擠離開市場後,又發現沒有機會去車衫和剪線頭,因為製衣廠沒有了；沒有機會去做裝配工作,因為電子廠消失了；連電話接線生及後勤文書的工作也遷移離境了.

過去,我們的社會為資本高速的擴張及集中而沾沾自喜；我們社會對於弱肉強食視為正常競爭；與此同時,我們的社會寄望食飽「弱肉」的「強者」發財立品,施捨捐贈.現在,越來越多人發現這個遊戲的惡果,越來越多人感到如此發展下去不甚妥當.我們社會中的「窮人和寄居者」－那些缺乏金錢資本,也缺乏知識資本的「無產者」又如何生活呢?上帝對西乃山下的社群的指示,讓「窮人和寄居者」也有自力更生的空間和機會,而不是對他們施捨和救濟,值得今日獅子山下的社群思考.

2．第二組指引:謀生要取之有道19:11 「你們不可偷盜,不可欺騙,也不可彼此說謊.19:12 不可指著我的名起假誓,褻瀆你　神的名.我是耶和華.

上帝沒有阻止西乃山下的地主以田地來謀生,只是肯定沒有田地的人也有謀生的需要.接者,下帝提醒這個社群的所有人,就算謀生是必須的,但卻不能不擇手段.從上帝對西乃山下的人們的指示可見,謀生最大的試探是誠信和廉潔,因為祂特別提到「偷盜,欺騙,說謊,起假誓」等有違誠信廉潔的行為.

謀生,獅子山下的人稱之為「搵食」.這裡的名言是:「搵食即,使唔使咁呀?」即是說,「搵食」是天經地義的首要任務,不應加以太多限制.其實,我們的社會對於如何「搵食」的確是最少限制.美國傳統基金會及華爾街日報聯合公布的「經濟自由度指數」顯示,香港以2.5分之微略高於新加坡,蟬聯全球最自由的經濟體系.報章報導此項消息時,多把香港與新加坡作比較.與新加坡作比較,有一個項目值得留意,就是廉潔程度(Freedom from corruption),新加坡取得92分,而香港只取得82分,足足落後10分之多.(\url{http://www.heritage.org/index})應該如何解讀這份報告,意見紛紜；應否採取更放任的經濟政策以保持最自由經濟體系的地位,也存在爭議.不過,透過改善廉潔程度來增加分數,應該不會有人反對的.事實上,人們在獅子山下「搵食」的誠信廉潔狀況,正引起關注.國際調查機構「政經風險評估」於2009-2010年訪問了亞太區及美國共十六個國家及地區共二千一百七十四名商界及行政人員,了解他們去年對各國的廉潔觀感.香港排名比前一年下跌.調查機構指,香港廉潔度評級下跌,相信與地產發展商影響力強大,在樓宇銷售上推出不少市場策略,但透明度卻不足有關.報告指出,這些引起爭議的樓宇買賣運作方式,其他市場也不容許.(\url{http://orientaldaily.on.cc/cnt/news/20100310/00176_041.html})

上帝禁止西乃山下的人民「搵食」的時候「偷盜,欺騙,說謊,起假誓」,因為這些破壞誠信廉潔的行為對相關的個人,甚至整個社會的傷害實在太大,太深遠了.這些行為傷害人與人之間的信任,也侵蝕著人民對制度的信任.餵哺嬰兒的時候要擔心奶粉中有三聚氰胺或者破舊皮鞋物料；揭發這些罪惡的時候又要擔心會否遭到法律懲處,生活於如此缺乏誠信廉潔的社會,人民怎能不充滿猜疑和驚恐呢?因此,獅子山下的人們應該思想,是的,我們都要「搵食」,我們都不會「阻人發達」,然而,如果我們要安居樂業於此,如果我們希望生活於一個「宜居」的地方,我們就要盡力維護這個地方的誠信和廉潔.

3．第三組指引:不可恃強凌弱19:13 「不可欺壓你的鄰舍,也不可搶奪他的物.雇工人的工價,不可在你那裡過夜,留到早晨.19:14 不可咒罵聾子,也不可將絆腳石放在瞎子面前,只要敬畏你的　神.我是耶和華.

任何社群都有著強弱懸殊的狀況,西乃山下的群眾也不例外.有些強弱差異出現於身體上,例如有些人五官齊整,有些人卻失聰或失明；有些強弱差異存在於勢力,例如早段日子網上流行的「我爸是李綱!」；有些強弱差異存在於資本,所以有人是顧主,有人是顧員.如何看待這些差異別?上帝對西乃山下人們的要求看來並不很高,只是吩咐他們不可恃強凌弱.鄰舍本來已經比你弱,你搶奪他的東西,這樣不是令他更弱嗎?工人依賴顧主發放的工資,顧主延遲發放人工,工人等著「開飯」的一家大小怎麼辦呢?你咒罵聾子,你將絆腳石放在瞎子面前,聾子和瞎子何來有能力保護自己呢?所有這些以強凌弱的行為,都在製造強者更強,弱者更弱的結果.如果是個人以強凌弱,我們稱之為「大蝦細」,而對付「大蝦細」的民間方法是:「俾屎餵」.然而,如果做成這種強者越強,弱者越弱,富者越富,貧者越貧的是經濟制度和社會政策,又可以如何對付呢?

農曆新年期間,其中一項引起討論的社會話題是:在天水圍買菜買過在灣仔買菜.有團體搜集了天水圍二十個食品及日用品項目,結果整體售價較灣仔區貴一成三.灣仔是全港最富有的地區,平均家庭入息比天水圍高出一倍.窮人要以更貴的價格購買食物和日用品,後果是貧者更貧.有學者指出,做成如此的現象,是因為天水圍的街市被領匯壟斷,沒有如灣仔,深水Po 等地區仍有街邊市場.獅子山下另一個製造富者越富,貧者越貧的現象出現於居住範疇.去年,當山頂豪宅租金每呎平均約70元,南區貝沙灣呎租要50多元的時候,深水Po 的板間房及籠屋呎租最高達93.3元.(明報2010-2-5)

西乃山下的人們存在著強弱的差異,獅子山下的人們這方面的差異可能更顯著.我們應該如何看待這些差異呢?相信沒有人會認同瑞銀投資研究執行董事王震宇的意見,認為「低增值人士」應該離開香港,以騰出空間讓高增值人士來港的建議.(\url{http://paper.wenweipo.com/2011/01/25/YO1101250002.htm})既然我們接受強者是獅子山下的一份子,弱者也是獅子山下的一份子,我們如何回應製造強者越強,弱者越弱的社會及經濟制度呢?上帝並沒有要求西乃山下的人們犧牲自己去幫助弱者,只是禁止他們「恃強凌弱」,「落井下石」,相信獅子山下沒有人會反對如此基本的要求吧!

4．第四組指引:處理糾紛時要公平19:15 「你們施行審判,不可行不義；不可偏護窮人,也不可重看有勢力的人,只要按著公義審判你的鄰舍.19:16 不可在民中往來搬弄是非,也不可與鄰舍為敵,置之於死(原文是流他的血).我是耶和華.

有人的地方就有糾紛,西乃山下的社群也在不能倖免.如何處理紛爭呢?第15節與審判有關,要求西乃山下的人們施行審判時,「不可偏護窮人,也不可重看有勢力的人」,即是要「不偏不倚」；16節則與作證有關,要求證人「不可搬弄是非,也不可與鄰舍為敵,置對方於死地.」

獅子山下「不偏不倚」的司法精神,形象化地呈現於香港現在的立法會大樓,即是從前的最高法院.大樓樓頂樹立了一個女神雕像,一手執劍,一手持天秤,雙眼蒙了起來.雕像手中的劍是判斷是非的權力,手中的天秤是平等,公正的表徵；矇上雙眼,表示大公無私,不問何人何黨派.在一個人際社會中,糾紛難免,若施法不公,會製造更多糾紛；若施法公正,則有助把糾紛劃上句號,制止糾紛漫延.在公平公正的司法制度之下,人們生活有所依循,遇到衝突時也知道如何處理,生活於此的人就會感到安全及安心.

普遍認為,獅子山下的司法制度比鄰近地區較為公平和公正.不過,由於法律費用高昂,社會中每一個有需要的人都能夠享用法律服務嗎?都有條件爭取公平的審訊嗎?有些留意這個範疇的團體及議員,常常希望我們不要忽略在目前的司法制度及法律程序之下,仍有吃虧的人,提醒我們法治精神和司法獨立需要我們維護,法律服務也需要繼續改善.

5．第五組指引:既同舟,應共濟19:17 「不可心裡恨你的弟兄；總要指摘你的鄰舍,免得因他擔罪.19:18 不可報仇,也不可埋怨你本國的子民,卻要愛人如己.我是耶和華.

最公平公正的司法制度,都只能夠把犯法的人處以罪有應得的刑罰,不過卻未必能夠化解相關人士心裡的恨意.有些犯人聞判後指著另一些人說:「我恨你一世!」；有些受害者聞判後被訪問時說:「我這一生也不會原諒他!」法庭處理了法庭能夠處理的部份,法庭不能處理的又如何呢?17-18說,衝突的可能是弟兄,可能是鄰舍,可能是本國的子民,即是同一群體中的一份子,大家仍要生活在同一個社會裡,難道真的能夠「河水不犯井水」嗎?上帝給西乃山下人們的吩咐是用愛去取代恨,用指正去取代報仇和埋怨.箴言27章5節說:「當面的責備,勝過隱藏的愛」.恨,總有原因；埋怨,總有理由；仇；總有要報的基礎,怎可能要求人們以德報怨呢?上帝沒有給予理由,只是再一次強調,那些被恨的,被埋怨的,被報仇的,畢竟都是弟兄,都是鄰舍,都是本國的子民.上帝勸募西乃山下的人們,要放棄的是報仇的衝動,而不是犯錯的弟兄；要清除的是心裡的恨意,而不是「黑人憎」的鄰舍.當獅子山下的人們面對群體內的矛盾,紛爭,埋怨,仇恨時,上帝對西乃山下人們的第五組指引有沒有可以參考的地方呢?

三．總結:為甚麼我們要如此相待?聽過了上帝在利未記19章9-18節的吩咐,人們可能會問,為甚麼我們要如此相待?有資本的,為甚麼不能賺到盡,而要「留給窮人和寄居的」?「搵食」的時候,為甚麼要以廉潔誠信來為規範自己?有能力的,為甚麼要顧及能力較弱的?對付異己,為甚麼要以公平公正來限制自己?對待「黑人憎」的鄰舍,為甚麼要以德報怨?

19章後段,記述了上帝對西乃山下的人們的提醒,34節說:「因為你們在埃及地也作為寄居的」(19:34)；36節又說:「我是耶和華－你們的上帝,曾把你們從埃及地領出來的.」(19:36)上帝希望他們仍然記得,作為一個群體,他們是從哪裡來的,他們曾經行過怎樣的路.曾經在埃及作過寄居的,難道不明白寄居者的需要嗎?曾經在埃及作過奴隸的,難道不知道被欺壓的苦況嗎?曾經在埃及為二等公民的,難道體會不到被不公平對待的痛苦嗎?在上帝的大能介入之下,他們才能夠離開壓迫者手；扶老攜幼,長途跋涉,千山萬水,這些集體回憶瀝瀝在目,難道脫離被欺壓之後,自己要成為另一位欺壓者嗎?脫離不公平之後,大家反而要互相欺負嗎?脫離了敵人之後,大家反而互相成為敵人嗎?

獅子山下的人們,難道不也是一同走過一條出埃及的道路嗎?我們這個社群從何而來的呢?難道我們忘記了獅子山下本來是一個荒涼的地方嗎?難道我們忘記了這片荒涼之地本來是走難者的聚居處嗎?難道我們忘記了當年鄉親們如何互相接濟,互相收留嗎?獅子山下的故事,曾不是一個如何「發達」的故事,而是一個群體如何相互幫助令大家都能夠生存下去的故事；是一個群體逃離壓逼,爭取自由的故事；是一個群體共同奮鬥建設家園的故事.難道我們眼白白地任由這個獅子山下的故事被改寫為一個弱肉強食的故事嗎?容讓它變成一個恃強凌弱的故事嗎?忍受它變成一個爭鬥仇恨的故事嗎?相信這不是獅子山下的人們樂見的,更不是上主的願望.願上主賜給此時此地獅子山下人們勇氣和毅力,使獅子山下那反抗壓逼,爭取自由,同舟共濟,同建家園的故事延續下去.



\newpage

\section{失去味的鹽味道如何呢}
\label{sec:89}
\textbf{陳曉東博士 - 「天國」.比如}
\newline
\newline
\hyperref[sec:88]{< < < PREV SERMON < < <}
~
\hyperref[sec:toc]{[返目錄]}
~
\hyperref[sec:90]{> > > NEXT SERMON > > >}


日期:2011年3月13日

講員:陳曉東博士

今天我們有個很有趣的題目《失去味道的鹽》,那這樣的鹽是什麼味道呢?

我們在開始之前,請先讀3章經文,馬太福音第5章,第6章和第7章,這是耶穌基督一段非常豐富,非常複雜的教訓,這3章經文是一個整體,首先是天國「8福」的教訓,這段教訓在馬太福音第5到7章裡有極其重要地位和價值.

馬太福音5-7章,教導我們以下幾件事情.

第1,是邀請和提醒我們神在救恩計劃裡「神成全律法,所表達的意義」第2,是邀請和提醒我們去追求卓越的人格.第3,耶穌透過一系列的講道,來喚醒我們去選擇追求道德的生活.

在這幾章經文裡重複著一個

主題:「天國」.比如:

馬太 5:3  虛心的人有福了!因為天國是他們的.

馬太 5:10  為義受逼迫的人有福了!因為天國是他們的.

馬太 5:20  我告訴你們,你們的義若不勝於文士和法利賽人的義,斷不能進天國.

馬太 6:10  願你的國降臨；願你的旨意行在地上,如同行在天上.

馬太 6:13  不叫我們遇見試探；救我們脫離兇惡(或作:脫離惡者).因為國度,權柄,榮耀,全是你的

馬太 6:20  只要積　財寶在天上；

馬太 6:33  你們要先求他的國和他的義

在這裡我們可以看到「天國」兩個字重複的出現.耶穌在簡述了「8福」之後,後面的經文又告訴了我們一個重要的事情:耶穌基督是律法的成全者.

馬太 5:17  「莫想我來要廢掉律法和先知.我來不是要廢掉,乃是要成全.」

再接下來耶穌用一系列短講來道出人生相關的主題,例如:「怒氣」,「和好」,「情慾」,「社會關係」等等的問題.在5 章13節到28節的地方出現了「你們要完全」的教導.這裡福音書作者馬太用了「施捨」,「禱告態度」等一些教導來呼應「我們如何完全」.一直到6章13的 主禱文,我們可以看到主禱文是和「8福」,「如何完全」,「生命的成長」等經文息息相關的.主禱文之後,耶穌教導我們一個饒恕的理由,以及和天父之間互動 的關係,他教導我們宗教的飲食,宗教禁食的意義等等.

在馬太 5:13-16節,「8福」之後的一段評語.如果我們成為「8福」中描述的人(虛心的人,哀慟的人,溫柔的人,飢渴慕義的人,憐恤人的人,清心的人,使人和睦的人,為義受逼迫的人)那麼我們就是世上的鹽,世上的光.

關鍵的問題是如果我們有鹽和光這麼美好的目標,但不去行,會怎麼樣呢?這裡耶穌用了一連串的反問

馬太 5:13  「你們是世上的鹽.鹽若失了味,怎能叫他再鹹呢?以後無用,不過丟在外面,被人踐踏了.」

我們如果是一群願意跟隨天國「8福」教訓的人,就好像:

馬太 5:15  人點燈,不放在斗底下,是放在燈台上,就照亮一家的人.

耶穌沒有要我們去照亮整個世界這麼偉大,但是他說我們可以照亮一家的人.

馬太 5:16  你們的光也當這樣照在人前,叫他們看見你們的好行為,便將榮耀歸給你們在天上的父.

之後,耶穌又說

馬太 5:17  莫想我來要廢掉律法和先知.我來不是要廢掉,乃是要成全.

這樣的教導成為教會處理文化,社會關係的基礎:我們來這個世界不是要廢掉它,而是要成全它,因為這是天父所造,所愛的世界,所以耶穌來不是要廢掉它,而是要轉化,帶它進入一個新天新地的行動.

弟兄姊妹當大家聽到這些信息的時,也許會很激動,或是很感動,或是心生嚮往.但當我們默許這段經文時,其挑戰又是何等 巨大呢?每當我默想這段經文時,就有這樣一些的感受:要麼這樣的經文對我們生命是一個巨大的衝擊,挑戰,它將重新塑造我們的生命,讓我們有一個新的起點； 要麼這樣的挑戰誠意太高了,所以我們讀完之後只是笑而束之高閣吧.

試想一下自己的人生,有多少時間真的是虛心?有多少的時間,見到別人哀慟時,我們也跟著哀慟?又有多少時間我們的心是 充滿了溫柔,充滿了對其他人的接納呢?更不要說這個世界本就充滿了暴力,充滿了兇殺,充滿了不公義,我們自己怎麼堅持著不參與,怎麼杜絕成為幫兇而不自知 呢.

昨天聚會有一個討論,一群香港的弟兄姊妹在為印尼籍,菲律賓籍家庭助理爭取權利而奔走,從這些弟兄姊妹的分享裡,我得知在香港有幾十萬外籍傭工,他們面對了很多不公義,剝削和虐待,甚至這些情形還出現在基督徒的家庭裡.

我們在主日裡讀這些經文,思考神的話語.但是回到家裡,面對自己所用的工人,卻又充滿了凶暴和虐待.所以大家要想到這 些教導:憐恤人的人有福了,清心的人有福了,使人和睦的人有福了,為義受逼迫的人有福了.我們想到這些美好道理的時候,又怎麼能只把它們放在道理的層面, 而不應用在生活的層面裡呢?  我們聽到這些教訓,可能會感覺很痛心,因為自己還有很多的虧欠和不足；更有甚者為了自己好過,寧可關閉心門,讓這些教導左耳進,右耳出,不存在心裡.

所以耶穌問我們:

馬太 5:13  「你們是世上的鹽.鹽若失了味,怎能叫他再鹹呢?

馬太 5:14  你們是世上的光.城造在山上是不能隱藏的.

馬太 5:15  人點燈,不放在斗底下,是放在燈台上,就照亮一家的人.」

在這些教導的背後,也是對我們生命的挑戰和指責.請問大家,願意接受這樣的改變嗎?

現在有很多知識淵博的學者在研究天國「8福」,並將它與猶太教文化傳統對比,發現當中有很多的教導是相互印證的.將耶 穌的教導和希臘,羅馬的文化傳統相比,也發現許多意思相近的地方.換而言之,學者們認為耶穌的教導可能來源於這些文化傳統,因為耶穌其實並未說過,這些話 是他原創的.

恰恰

馬太寫到 5:17  「莫想我來要廢掉律法和先知.我來不是要廢掉,乃是要成全.」

這就告訴我們,神在他創造的次序裡面,已經把這美好的教導安放其中.在耶穌來之前,人和耶穌相處,耶穌去之後,這虛心 的教導,有哀慟心的教導,溫柔的教導,追求公義心的教導,憐恤人的教導,清心的教導,和平的教導等等,不是基督教獨家擁有的,不是耶穌獨家擁有的.但關鍵 是:

馬太 5:17 我來不是要廢掉,乃是要成全.

馬太 5:18  我實在告訴你們,就是到天地都廢去了,律法的一點一畫也不能廢去,都要成全.

當我們聽到這些教導後,請大家思考我們該怎麼去回應呢?

從5章第21節開始到第7章,耶穌通過一系列的短講來教導我們應該怎麼去回應,怎麼去面對生命中遇到的挑戰和問題.我們來讀其中一段:

馬太 5:21  「你們聽見有吩咐古人的話,說:『不可殺人』；又說:『凡殺人的難免受審判.』

馬太 5:22  只是我告訴你們:凡(有古卷在凡字下加:無緣無故地)向弟兄動怒的,難免受審判；凡罵弟兄是拉加的,難免公會的審斷；凡罵弟兄是魔利的,難免地獄的火.

馬太 5:23  所以,你在祭壇上獻禮物的時候,若想起弟兄向你懷怨,

馬太 5:24  就把禮物留在壇前,先去同弟兄和好,然後來獻禮物.

馬太 5:25  你同告你的對頭還在路上,就趕緊與他和息,恐怕他把你送給審判官,審判官交付衙役,你就下在監裡了.

馬太 5:26  我實在告訴你,若有一文錢沒有還清,你斷不能從那裡出來.」

馬太 5:27  「你們聽見有話說:『不可姦淫.』

馬太 5:28  只是我告訴你們,凡看見婦女就動淫念的,這人心裡已經與他犯姦淫了.

馬太 5:29  若是你的右眼叫你跌倒,就剜出來丟掉,寧可失去百體中的一體,不叫全身丟在地獄裡.

馬太 5:30  若是右手叫你跌倒,就砍下來丟掉,寧可失去百體中的一體,不叫全身下入地獄.」

馬太 5:31  「又有話說:『人若休妻,就當給他休書.』

馬太 5:32  只是我告訴你們,凡休妻的,若不是為淫亂的緣故,就是叫他作淫婦了；

如果我們仔細讀上面的經文,會發現一個非常有趣的排比句式.目前我們不知道這樣的排比句式是耶穌的原話,還是作者馬太在收集,記錄耶穌話語時寫成了排比句式.這是一個非常有趣句式,每句話都先點題,然後用排比來說三點,並做判詞.例如:

馬太 5:21  「你們聽見有吩咐古人的話,說:『不可殺人』；又說:『凡殺人的難免受審判.』」 and \#8212;-這句話是在點題.

馬太 5:22  只是我告訴你們:凡(有古卷在凡字下加:無緣無故地)向弟兄動怒的,難免受審判；凡罵弟兄是拉加的,難免公會的審斷；凡罵弟兄是魔利的,難免地獄的火. and \#8212;-排比說三點.

從這樣的句式結構裡,我們可以看到耶穌講話的一種風格,耶穌就是用這樣風格的話語來對我們進行教導的.但是這些教導到 底是律法,還是一種神學的修辭呢?如果這是律法,我們跟隨主的人就要竭盡所能的完成律法的要求.如此說,福音是可以成為律法的.但也有人不同意這樣的講 法,說這是一種神學的修辭,不但要看字面意思,還要看其背後的含義.

馬太 11:12  天國是努力進入的

我們要這樣理解神對我們生命的要求.例如耶穌說不可殺人,不單是行為上,外表的要求,還包含了內心裡,當我們和人衝突,當我們怒氣發生的時候,不就已經干犯了神對人的要求了嗎?

所以我們要思考:那怎麼辦?--這才是關鍵.我們每個人都有這樣的情況,難道都要死光,都要在神的地獄之火裡被懲罰嗎?這就我們今天的

題目:如果鹽失去了味道意味著什麼?

如果我們這些基督徒,或是號稱是基督徒的人,活不出神的樣式,那意味著什麼?如果這個鹽的味道真的失去了,這就意味著耶穌渴望我們成就的人格不能被達到.耶穌教導的天國「8福」的各樣生命品格,我們就沒有動力去達成.

如果我們這些鹽失去味道,即是在耶穌短講裡各個場景中失去了道德抉擇的自由.我們也就沒有能力去實現「8福」的品格.請大家不要忘記,耶穌的語重心長,他的邀請是:我們是世上的鹽.我們是世上的光.

我們只能捶胸頓足,羞愧難當的說,「我罪!我罪!我的重罪!主啊-我知道,我是世上的鹽,我是世上的光,但是我達不到 主你要我做的事情.我的生命裡面做鹽,但是沒有鹹味；我的生命裡面作光,但是我把光放在斗底下.我心裡失去了勇氣,能力,還有渴望,我達不到你的標準,活 不出你的形象,活不出你的義.主啊-我怎麼辦.」

當我們默想這些經文的時候,請大家注意,不要將自己陷入到律法主義當中.耶穌在這裡告訴我們要達到的目標,以及他的要 求；但他沒有說我們現在就要達成,他發出的是一個邀請,我們現在可以選擇一個改變,可以再一次將我們發酸的腿,下垂的手靠著神的恩典舉起來.當我們願意再 一次定睛仰望耶穌基督的十字架時,十字架會發出救贖的光,來進入我們生命深處,賦予我們能力,愛心,盼望.希望大家能夠通過今天的經文,能夠看到這個重要 的提醒.

在我們生命的成長裡面,會看到很多的師傅,很多的榜樣.有的師傅叫阿波羅,有的師傅叫保羅,有的師傅叫彼得.我們很容 易陷入一種盲目:用這些師傅替代了神,把這些師傅放得比神還大.但是保羅提醒我們:每個生命,讓它成長的是神.所有的師傅是幫助我們靠近神的,與神建立更 深的連接.當我們明白這個道理後,從天國「8福」裡就會有更深的體會:是的!主啊-你告訴我:虛心的人有福了,但是我做不到.主啊-你告訴我:哀慟的人有 福了,我也做不到.溫柔的人,飢渴慕義的人,憐恤人的人,清心的人,使人和睦的人,為義受逼迫的人這些我都做不了....

各位弟兄姊妹,也許正因為如此的坦誠,也許正因為我們看到了自己生命中的不足,我們才能夠再一次的回到神面前,求能 力,求寬恕,求出路,求成長,失去鹹味的鹽可不可以在再一次有味道呢?　我相信是可以的,當我們再一次回到神面前,把我們的手掌打開,求神按著他的真光來 照亮我們的時候,求神用他無條件的愛再次充滿我們的時候,我們可以行公義,好憐憫,能夠謙卑的與主同行.阿們-



\newpage

\section{飛鳥野花神的國}
\label{sec:90}
\textbf{鄧瑞強博士 - 飛鳥,野花,神的國}
\newline
\newline
\hyperref[sec:89]{< < < PREV SERMON < < <}
~
\hyperref[sec:toc]{[返目錄]}
~
\hyperref[sec:91]{> > > NEXT SERMON > > >}


馬太福音 6:24-34
\newline
\begin{longtable}{cl}
\hline
\hline
章節 & 經文 (和合本修訂版)\\
\hline
6:24 & \begin{tabularx}{0.7\textwidth}{X} 「一個人不能服侍兩個主；他不是恨這個愛那個，就是重這個輕那個。你們不能又服侍神，又服侍瑪門。」 \end{tabularx} \\ \\ \relax
6:25 & \begin{tabularx}{0.7\textwidth}{X} 「所以，我告訴你們，不要為你們的生命憂慮吃甚麼喝甚麼，或為你們的身體憂慮穿甚麼。生命不勝於飲食嗎？身體不勝於衣裳嗎？ \end{tabularx} \\ \\ \relax
6:26 & \begin{tabularx}{0.7\textwidth}{X} 你們看一看那天上的飛鳥，也不種也不收，也不在倉裡存糧，你們的天父尚且養活牠們。你們不比飛鳥貴重得多嗎？ \end{tabularx} \\ \\ \relax
6:27 & \begin{tabularx}{0.7\textwidth}{X} 你們哪一個能藉著憂慮使壽數多加一刻呢？ \end{tabularx} \\ \\ \relax
6:28 & \begin{tabularx}{0.7\textwidth}{X} 何必為衣裳憂慮呢？你們想一想野地裡的百合花是怎麼長起來的：它也不勞動也不紡線。 \end{tabularx} \\ \\ \relax
6:29 & \begin{tabularx}{0.7\textwidth}{X} 然而我告訴你們，就是所羅門極榮華的時候，他所穿戴的還不如這些花的一朵呢！ \end{tabularx} \\ \\ \relax
6:30 & \begin{tabularx}{0.7\textwidth}{X} 你們這小信的人哪！野地裡的草今天還在，明天就丟在爐裡，神還給它這樣的妝飾，何況你們呢？ \end{tabularx} \\ \\ \relax
6:31 & \begin{tabularx}{0.7\textwidth}{X} 所以，不要憂慮，說：『我們吃甚麼？喝甚麼？穿甚麼？』 \end{tabularx} \\ \\ \relax
6:32 & \begin{tabularx}{0.7\textwidth}{X} 這都是外邦人所求的。你們需要這一切東西，你們的天父都知道。 \end{tabularx} \\ \\ \relax
6:33 & \begin{tabularx}{0.7\textwidth}{X} 你們要先求神的國和他的義，這些東西都要加給你們了。 \end{tabularx} \\ \\ \relax
6:34 & \begin{tabularx}{0.7\textwidth}{X} 所以，不要為明天憂慮，因為明天自有明天的憂慮；一天的難處一天當就夠了。」 \end{tabularx} \\ \\
[1ex]
\hline
\hline
\end{longtable}
日期:2011年3月20日

主題:飛鳥,野花,神的國

證道:鄧瑞強博士

經文:馬太福音6:24-34

各位弟兄姊妹,主內平安.

做人時時面對抉擇.有些抉擇很重要,影響日後整個生命,如:是否和這個人結婚?(之後是禍是福就看這決定了.)是否將一個可能有問題的小孩子生下來?(這可能決定了日後整個家庭的責任承擔.)是否移民?(這影響日後整個生活的方式.)「抉擇」舖設了日後要走的路.劉曉波不逃避坐牢,因為他抉擇了為民主而奮鬥,而他明白坐牢是這抉擇的後果.他的「抉擇」舖設了他日後要走的路.

在網上讀到劉曉波的一篇文章,叫＜獄中重讀《獄中書簡》＞.他坐牢時讀另一個坐牢的人寫的書.那本書叫《獄中書簡》(Letters and Papers from Prison),作者是一個德國神學家,叫潘霍華(Dietrich Bonhoeffer),大陸翻譯叫朋霍費爾.這個神學家在納粹時代在德國參與反納粹的活動而坐牢.劉曉波寫:「朋霍費爾的榜樣正在逼視和召喚,坐牢正是參與塵世苦難的一種方式,無論如何,我們都不會放棄的,縱使我們改變不了什麼,但我們的行為起碼可以證明耶穌精神仍然活在人間；在無上帝的現代世界,耶穌精神是唯一能夠抗衡人類墮落的信仰力量.正如朋霍費爾所說:『產生行動的並不是思想,而是願意承擔責任的準備.』...在苦難中發現希望和意義,人的存在的品質視其接近上帝(神聖)的程度而定.沒有神聖價值參照的生命只是一種深淵似的貧乏.因為人的生命能夠從神的能力和信念中獲得存在的豐富性和高貴性.」

神學家潘霍華為何要走向牢獄,因為他作了一個反抗暴政的抉擇.劉曉波為何要走向牢獄,因為他作了一個抗衡人類墮落的抉擇.他們的「抉擇」顯出了他們的生命內涵,同時也舖設了他們日後要走的路.

今日講道的經文是:馬太福音6: 24-34

這經文的主題,涉及人的抉擇,人的終極抉擇,一種改變我們整個生命的抉擇,一種叫生命獲得高貴性和永恆意義的抉擇,一種叫我們過新生活的抉擇.經過這抉擇,一切都會變成新的,新的世界,新的人生.

這段講道經文的第一句,即馬太福音6:24說:「一個人不能事奉兩個主；不是惡這個,愛那個,就是重這個,輕那個.你們不能又事奉神,又事奉瑪門〔瑪門:財利的意思〕.」

經文要我們面對一個生命的真相,這真相是:「一個人不能事奉兩個主」.因為主人的要求不同,你不能同時滿足兩個不同的要求.兩個主人指出兩條不同的路,你不能同時走兩條路.一個主人要求你肥一點,一個主人要求你瘦一點.你能如何做?總不能變成「半肥瘦」.

主耶穌在這裡將人生的兩個主人分別稱為「神」和「瑪門」(Mammon).「瑪門」一般翻譯為「錢財」,然而,我們較少留意到「瑪門」(Mammon)這個字其實和「阿們」(Amen)這個字在希伯來文字根上是相關的.「阿們」指「確實的,穩妥的」.我們祈禱時講「阿們」,意指我們的祈禱是「確實如此的」,我們是「誠心所願」如此的.而「瑪門」(Mammon)這個字的字根意義也是「確實的,穩妥的」.

「瑪門」一般解作「錢財」,「錢財」也的確讓人感到「確實,穩妥」.但若直接將「瑪門」理解為「錢財」,就很難和下文相接.下文馬太福音6:25開頭講「所以...」,即是說,下文是作為一個例子講解「神」與「瑪門」的分別的,但下文並不講「錢」,而講「生命的終極所賴」,故此,我們可以理解,這裡講「瑪門」(Mammon),重點在這字的字根意義,即生命的「穩當性」的問題.生命的兩個主人,就是人生的兩條大路,就是兩條不同的尋求生命的「穩當性」的路.一條路是通過信賴神,將「穩當性」寄托於神；另一條路是通過信賴自己,將「穩當性」寄托於我們用自己的手去握住叫我們覺得穩當之物.

你將生命的「穩當性」寄托於何處,是一個重要的問題.這是人生的抉擇所在,這是你要作的終極抉擇所在.

馬太福音 6:25-34對這生命的終極抉擇作出一些提示:6:25    「所以我告訴你們,不要為生命憂慮吃什麼,喝什麼；為身體憂慮穿什麼.生命不勝於飲食嗎?身體不勝於衣裳嗎?6:26    你們看那天上的飛鳥,也不種,也不收,也不積蓄在倉裡,你們的天父尚且養活牠.你們不比飛鳥貴重得多嗎?6:27    你們哪一個能用思慮使壽數多加一刻呢〔或譯:使身量多加一肘呢〕?6:28    何必為衣裳憂慮呢?你想野地裡的百合花怎麼長起來；它也不勞苦,也不紡線.6:29    然而我告訴你們,就是所羅門極榮華的時候,他所穿戴的,還不如這花一朵呢!6:30    你們這小信的人哪!野地裡的草今天還在,明天就丟在爐裡,神還給它這樣的妝飾,何況你們呢!6:31    所以,不要憂慮說:吃什麼?喝什麼?穿什麼?6:32    這都是外邦人所求的,你們需用的這一切東西,你們的天父是知道的.6:33    你們要先求他的國和他的義,這些東西都要加給你們了.6:34    所以,不要為明天憂慮,因為明天自有明天的憂慮；一天的難處一天當就夠了.」

這段經文的開始是「所以...」,顯示這段經文是承接上文的.這段經文講的是上文討論的生命的終極主宰的問題.經文結尾部分提到「神的國」,同樣顯示這段經文涉及我們的生命應服在什麼權柄之下的問題.故此,這段經文講的是人生的終極抉擇問題,而不是處理如何不再憂慮的心理學問題.換句話說,這段經文討論的是神學,而不是心理學,縱使神學的討論會有心理學的意涵.

這段經文可分析出三個小論證.第一個論證:6:25    所以我告訴你們,不要為生命憂慮吃什麼,喝什麼；為身體憂慮穿什麼.生命不勝於飲食嗎?身體不勝於衣裳嗎?6:26    你們看那天上的飛鳥,也不種,也不收,也不積蓄在倉裡,你們的天父尚且養活牠.你們不比飛鳥貴重得多嗎?

主耶穌叫我們看天空的飛鳥.在平凡的眼光裡,天空的飛鳥只是平凡之物.但主耶穌叫我們「看」,看看這平凡之物背後所隱藏的神聖世界.在這神聖世界裡,有一主人,這主人在這裡被稱為「天父」,這「天父」看顧著這平凡的飛鳥,這「天父」珍惜著世上被看為無關痛癢的生命.而這小鳥,沒有試圖抓住什麼,沒有試圖將生命的「穩當性」放在自己手中,牠們倚靠的,只是天父的仁慈；牠們倚靠的,是上主的恩典.牠們的生命,顯出了上主無限的愛.

一隻小鳥也獲得上主的眷顧,何況是比飛鳥貴重得多的人呢?

這是一種「由弱到強」的推論.即:弱者如飛鳥也做得到,強者如人是否應做得更好呢?

面對生命的「穩當性」問題,人抉擇倚靠神,抑或倚靠自己製造出來的「瑪門」(Mammon)?我們的心仰望創造天地的神,抑或仰望「瑪門」這種虛假的「穩當性」?

一隻小鳥尚且明白生命全依賴神的仁慈,人是否應該明白得更多呢?

第二個論證:6:27    你們哪一個能用思慮使壽數多加一刻呢〔或譯:使身量多加一肘呢〕?6:28    何必為衣裳憂慮呢?

這是以人自己的生命體驗來論證.

我們誰能用自己的方法使壽數多加一刻呢?或者更準確的翻譯是:誰能用自己的方法使身量多加一肘呢?或用現代的講法,誰能用自己的方法使皮膚皺少一點?誰能用自己的方法使基因完美一點?

生命的長短是我們控制不了的,身體的特徵是我們控制不了的,既然這些如此重要的東西我們也管不了,我們為何為「衣裳」這些外在的東西而擔憂呢?

我們想將生命的「穩當性」牢牢地握在手中,但這可能嗎?生命裡最重要的東西,我們都握不住,那麼,握住那些無關痛癢的東西,又有何益呢?

這是一種「由強到弱」的推論.即:生命的長短如斯重要的東西,我們都管不了,只能仰賴上天的眷顧；那麼,弱如穿什麼衫這類事情,是否更應仰賴上天的眷顧呢?

第三個論證:6:28    你想野地裡的百合花怎麼長起來；它也不勞苦,也不紡線.6:29    然而我告訴你們,就是所羅門極榮華的時候,他所穿戴的,還不如這花一朵呢!6:30    你們這小信的人哪!野地裡的草今天還在,明天就丟在爐裡,神還給它這樣的妝飾,何況你們呢!

主耶穌叫我們仔細研究「野地裡的百合花」,思想「野地裡的草」.

同樣,在生命的「穩當性」的議題裡,「野地裡的百合花」和「野地裡的草」,是最無「穩當性」的,最不值得時時想著「穩當性」的人去理會的.但主耶穌卻偏偏叫我們去留心這些最無「穩當性」的生命.

「野花」和「野草」,不紡線,不穿戴,也不妝飾,但在天父的看顧裡,它們的生命比所羅門王更美麗.

這也是一個「由弱到強」的推論.即:弱者如花草的生命,在天父的看顧裡,也有其「莊嚴」和「穩當性」,強者如人,不是更應體會到天父的眷顧,而從中找到生命的「穩當性」嗎?

「飛鳥」也不種,也不收,也不積蓄在倉裡.可以說,「飛鳥」對努力工作和積蓄的「男士們」作出提醒.

「野花」和「野草」,不紡線,不穿戴,也不妝飾.可以說,「野花」和「野草」對努力買衫和努力裝扮的「女士們」作出提醒.

講完以上的論證,主耶穌作出結論:6:31    所以,不要憂慮說:吃什麼?喝什麼?穿什麼?6:32    這都是外邦人所求的,你們需用的這一切東西,你們的天父是知道的.6:33    你們要先求他的國和他的義,這些東西都要加給你們了.6:34    所以,不要為明天憂慮,因為明天自有明天的憂慮；一天的難處一天當就夠了.

「外邦人」就是不明白生命主權屬於神的人,他們膜拜「瑪門」,他們想用自己的方法去尋求生命的「穩當性」.就是這種想法,將他們的人生牢牢地鎖在「憂慮」的深淵裡.他們看不見飛鳥背後的神聖世界,他們想不透野地的花和草訴說著的神的看顧.他們痛苦地,拼命去抓,想抓住一些東西讓自己覺得穩當一些.他們抓住的東西,就像是一塊又一塊的石頭,這些石頭裝在生命的行囊裡,只會越背越重,最終將人壓得透不過氣來.

面對未來,不明白「神國」道理的人,會用種種方法去控制未來,儲蓄多一點,裝扮靚一點,去克服面對未來而有的恐懼.殊不知,未來不是我們掌控的領域,面對不可知的領域而試圖控制之,只會令人越發「憂慮」.

「明天自有明天的憂慮」,明天不是屬於我們的,明天是屬於神的.「神掌管未來」.我們試圖用自己的手去掌管明天,去操縱生命的「穩當性」,自然徒勞無功.只有信賴掌管明天的主,我們的生命才能安頓下來,找到真正的「穩當性」.

「先求他的國」,就是作出終極的明智抉擇,以神為神,放棄「瑪門」,放棄自求「穩當性」,而相信神的應許,信賴天父的仁慈.你會倚靠神的恩典,而不是人的謀算.你會倚靠神的大愛,而不是在人間戰場中謀求勝利.這生命會引導你活出生命的「義」.這種「義」源自看到生命終歸之所屬,看到「神」而放下「瑪門」.明白財富的獲得,並不能叫人獲得生命的「穩當性」,而在分享財富時,享受生命的豐盛.明白人生之所得,並不能叫人獲得生命的「穩當性」,而當生命為他人的福祉而活時,生命進入神聖.

在這時候,一隻小鳥是值得珍惜的,一個社會上的失敗者是值得幫助的,一個仰賴他人支援的窮人是值得我們去愛的.在這時候,你會發現一朵野花是美的,一個拾紙皮維生的老人家是美的,一個殘障的人是美的.在這時候,你就不會介意你擁有多少,你就不再介意皺紋多了幾多條.因為,這時候,你的生命自由了.這時候,你在天國中.這時候,你不再憂慮.這時候,你明白人活著是靠賴天父的恩典.這時候,你與主耶穌面對面地活著.

生命的路是需要抉擇的.「一個人不能事奉兩個主；不是惡這個,愛那個,就是重這個,輕那個.你們不能又事奉神,又事奉瑪門.」

主耶穌很渴望你選擇那回歸天國的路.

但願榮耀歸於聖父,聖子,聖靈.阿們.



\newpage

\section{登山變像命途}
\label{sec:91}
\textbf{伍渭文牧師 - 登山,變像,命途}
\newline
\newline
\hyperref[sec:90]{< < < PREV SERMON < < <}
~
\hyperref[sec:toc]{[返目錄]}
~
\hyperref[sec:92]{> > > NEXT SERMON > > >}


馬太福音 17:1-9
\newline
\begin{longtable}{cl}
\hline
\hline
章節 & 經文 (和合本修訂版)\\
\hline
17:1 & \begin{tabularx}{0.7\textwidth}{X} 過了六天，耶穌帶著彼得、雅各和雅各的弟弟約翰，領他們悄悄地上了高山。 \end{tabularx} \\ \\ \relax
17:2 & \begin{tabularx}{0.7\textwidth}{X} 他在他們面前變了形像，他的臉明亮如太陽，衣裳潔白如光。 \end{tabularx} \\ \\ \relax
17:3 & \begin{tabularx}{0.7\textwidth}{X} 忽然，有摩西和以利亞向他們顯現，與耶穌說話。 \end{tabularx} \\ \\ \relax
17:4 & \begin{tabularx}{0.7\textwidth}{X} 彼得回應，對耶穌說：「主啊，我們在這裡真好！你若願意，我就在這裡搭三座棚，一座為你，一座為摩西 ，一座為以利亞。」 \end{tabularx} \\ \\ \relax
17:5 & \begin{tabularx}{0.7\textwidth}{X} 說話之間，忽然有一朵明亮的雲彩遮蓋他們，又有聲音從雲彩裡出來，說：「這是我的愛子，我所喜愛的。你們要聽從他！」 \end{tabularx} \\ \\ \relax
17:6 & \begin{tabularx}{0.7\textwidth}{X} 門徒聽見，就俯伏在地，極其害怕。 \end{tabularx} \\ \\ \relax
17:7 & \begin{tabularx}{0.7\textwidth}{X} 耶穌進前來，拍拍他們，說：「起來，不要害怕！」 \end{tabularx} \\ \\ \relax
17:8 & \begin{tabularx}{0.7\textwidth}{X} 他們舉目，不見一人，只見耶穌獨自一人。 \end{tabularx} \\ \\ \relax
17:9 & \begin{tabularx}{0.7\textwidth}{X} 下山的時候，耶穌囑咐他們說：「人子還沒有從死人中復活，你們不要把所看到的告訴人。」 \end{tabularx} \\ \\
[1ex]
\hline
\hline
\end{longtable}
日期:2011年3月27日

題目:登山,變像,命途

證道:伍渭文牧師

讀經:馬太福音17:1-9

假如今天的題目改成一個字:「登機,變像,命途」,你會想起甚麼呢?我們想起今年一月的易容偷渡駭人新聞:一內地青年戴矽膠頭套易容變像,扮成白人老翁,偷渡到加拿大.這人以真正的身份入境及出境本港,抵達機場禁區後,變像易容,使用他人的登機證及假護照登機.希望到了加拿大,改變自己的命途.

今天的題目是「登山,變像,命途」,這題目也使我聯想另一新聞.報載近年瘦身纖體熱潮席捲全港,不少女士不惜一擲數萬至十多萬元,以償苗條心願.有專責協助客戶重整債務的律師指出,近期有多名客戶因參與纖體療程,被纖體公司一次過過數而弄致「碌爆信用卡」,無力還款,陷入破產邊緣,該律師擔心若這股瘦身風氣持續,因瘦身而導致破產將成為新趨勢.

這則新聞指出不少人願意花錢纖體,弄至債台高築,像被一坐大山壓在自己身上.但為何要這樣犯險登上債山呢?我想這些人希望登債山以致能變像 and \#8212;使自己看來更美麗動人,改變命途；能夠找到如意郎君,如意妻子；從事演藝行業的,變像後接更多廣告,拍更多電影,開更多演唱會.

無論登機變像或登債山變像,他們就想改變自己的命途,而今天福音經課登山變像,剛剛相反,是耶穌顯明不會偏離自己的命途.

1.    命途耶穌變像,談論離世的事.太十六21記載耶穌在彼得公開認信他是彌賽亞後,立刻表明他要上耶路撒冷受死.彼得就拉著他:「

主啊,萬不可如此!這事必不臨到你身上.」(太十六22)耶穌對彼得說:「

撒旦,退我後邊吧!你是絆我腳的,因你不體貼神的意思,祇體貼人的意思.」(太十六23)這段經文緊隨上一章耶穌對自己命途的宣告,是上帝藉耶穌變像,明確表示他真是彌賽亞,而彌賽亞的命途就是上十字架受死.太十七3:忽然有摩西,以利亞向他們顯現,同耶穌說話.說些甚麼呢?馬太沒有說明,但路加則說得很清楚是耶穌要上耶路撒冷受死的事.路九30「

忽然有摩西,以利亞兩個人同耶穌說話:他們在榮光裡顯現,談論耶穌去世的事.」

各人命途不同,但終點相同.每人自有不同的人生路徑,但始終也要來到相同的終點.白先勇前一段時間在崇基周會演講,有同學問他的小說「台北人」的主題意涵.他回答:寫的是「美人遲暮,英雄垂老.」台北人寫的是大陸四九年隨國民政府播遷到台灣的上海人,書中「金大班最後的一夜」最能表達這意思.上海豪華百樂門的夜夜笙歌,已成絕響,明天就離開台北夜巴黎,從良嫁人的金大班,此刻已徐娘半老；鉛華盡洗,明燈轉暗,我們都被時間拋在後面,最後歸回塵土.登機易容變像,成功移民,在彼邦努力打拼,成為英雄；塑身變像,成為絕代美人,但最終也會英雄垂老,美人遲暮,歸回塵土.

聖灰日的揭示.今天是顯現期最後的主日,下星期開始就是預苦期；除卻歡樂主日,從星期三聖灰日開始計,直到受苦節一共四十天,預備自己,紀念基督的受難.教會在聖灰日的崇拜中有塗灰禮,主禮把去年棕枝主日(受苦節前一主日)  and \#8212;耶穌進耶路撒冷時 and \#8212;人們歡迎揮動的棕枝,燒成灰燼,以十架記號塗在領灰者的額上或掌心說:「父上帝,你用地上的塵土創造我們,求你施恩,使這些灰燼向我們揭示生命的有限,並提醒我們必須悔罪,因為只有仰賴你的恩典,我們才能領受在我們的救主耶穌基督裡頭的永生.」眾答:阿們.在塗灰時,主禮向領灰者說:「你要記住,你本是塵土,仍要歸於塵土.當遠離罪,效忠基督.」

命途沒有路線,只有路標.耶穌變像時,摩西和以利亞的出現,談論他去世的事,有甚應意思呢?摩西和以利亞二人代表著律法和先知, 耶穌的死是聖經的預言和應驗.耶穌在以馬忤斯路上,「

從摩西和眾先知起,凡經上所指萬自己的話都給他們講解明白了.」(路二十四27)因為耶穌的解釋,本來隱密難明的,終於真相大白,說明命途是有意義的,而這意義來自對聖經的認識和了解.「去世」原文為出埃及(exodus),耶穌受死是一件救贖事件,像昔日摩西帶領以色列人離開為奴之地,進入迦南應許之地.  聯合國第二任秘書長韓馬紹 Dag Hammarskjold 說:我們不需要尋求死亡,死亡會找上門的,但我們要找到往死亡的路,使死亡顯得有意義 (Do not seek death. Death will find you. But seek the road which makes death a fulfillment) .他又說:「命運本身不能逃避,也不需響往,它不是沒有原委的隱密,因為它說明這世界和人類歷史的軌跡,都有意義.」(Destiny is something not be to desired and not to be avoided, a mystery not contrary to reason, for it implies that the world, and the course of human history, have meaning. )

我們要找到往死亡的路,使死亡顯得有意義,而意義來自聖經的話語,聖經是我們人生的路徑的路標.登山和進入醫院的指引是不同的,醫院利用地上不同顏色的路線引領我們,但登山遠足只有路標.登山者不想踏著地面劃上的路線前行,他們認為透過觀察,推理,在路標指示之下,作出自己的抉擇才有樂趣,雖然有時會為著尋找正確方向而迷惘,作為天路客的基督徒有時也會為到前路感到焦慮.

2. 登山登山臨絕頂,一覽群山小.太十七1:「

過了六天,耶穌帶著彼得,雅各和雅各的弟兄約翰暗暗的上了高山.」「

耶穌帶著彼得,約翰,雅各上山去禱告.正禱告的時候,他的面貌就改變了,衣服潔白放光.」(路九28, 29)

耶穌在那一個山變像呢?聖經沒有說明,只說是高山,還強調暗暗的帶同三位親密的門徒為著禱告而上高山.高山本身是有吸引力的.數星期前鄧瑞強博士講道時,當時的題目為「無法抗拒的呼喚」,提及海明威(Ernest Hemingway)有一篇短篇小說,叫《屈力馬扎羅山的雪》(The Snows of Kilimanjaro).小說的開場白是這樣的:「屈力馬扎羅山是一座海拔19,710呎的雪山,據說是非洲最高的山.它的西峰,馬賽語(the Masai)稱之為『神的居所』(the House of God).在接近西峰處,有一具枯乾而凍僵的豹子屍體.沒有人知道,這隻豹子在那高處究竟尋找什麼.

但高山也有它的試探,登山臨絕頂,一覽群山小,有大地就在我腳下,有君臨天下的感覺.耶穌受試探時,「

魔鬼又帶他上了一座最高的山,將世上的萬國與萬國的榮華都指給他看,對他說:你若俯拜我,我就把這一切都賜給你.」(太四8)

耶穌登山為的是避靜,追求獨處,過濾思想,沉澱情緒,聆聽父神的聲音,使生命重新聚焦,調較人生方向.二十世紀中期,隱修士多瑪斯‧牟敦 (Thomas Merton) ,寫成的經典靈修自傳,已有二十多種譯文,就以山命名:《七重山》(The Seven Storey Mountain) .書中的七重山是但丁神曲中淨界 (purgatory) 七宗罪的種種試探.在人生步向死亡的路途,我們常常面對驕傲 (pride),嫉妒 (envy),暴怒 (anger),懶惰 (sloth),貪財 (greed),貪食 (gluttony),貪色 (lust) .勝過這些試探,就能登上七重山的頂峰樂園.其實,七宗罪是愛的錯配,罪是沒有愛上帝,愛其他多於上帝.

在《七重山》將結束時,牟敦總結他登山目的,就是回到上帝那裡:「我希望獨處的唯一原因 and \#8212;忘掉一切被造之物,它們對我來說,已經死了,毫不認識,因為它們使我與你遠離」(That is the only reason why I desire solitude and \#8212;to be lost to all created things, to die to them and to the knowledge of them, for they remind me of my distance from You) .3. 變像 耶穌和三個門徒登山變像,「

臉面明亮如日頭,衣裳潔白如光」(太十七2),山頭頓成天上人間,難怪彼得想把時間停住.「

主啊,我們在這裡真好!你若願意,我就在這裡搭三座棚:一座為你,一座為摩西,一座為以利亞.」(太十七4)但就在此刻；「

忽然有一朵光明的雲彩遮蓋他們,且有聲音說:這是我的愛子,我所喜悅的,你們要聽他.」(太十七5)耶穌沒有叫彼得停留在山上的美好時光,連在山上發生甚麼也不要告訴人.

接著他們下山了,迎面來的是一位父親,他兒子被鬼所附,並害癲癇病.(14-18節)我們知道,癲癇病發作時,臉容極其扭曲難看,與山上變像耶穌的臉面成強烈的對比.耶穌斥責污鬼和醫好孩子,使他恢複原來的臉相.耶穌沒有解釋為何是這孩子被鬼所附,為何他有癲癇病,他只憐憫他,治好他.

在人生的命途中,我們也許像這個家庭,遇到飛來橫逆,造成變像,沒有原因,看來極其荒謬.馬尼拉人質被殺事件正進行死因聆訊,其中一位幸存者作証時一直帶上口罩.這位女士本來相貌娟好,但被子彈擊碎下顎,事發後三個月無法說話,至今不敢不戴口罩外出,怕被說是「醜八怪」.我們真你不明白為何發生馬尼拉人質被殺事件,這是一件看來荒謬的事情,一個人的迷失,層層的錯誤,弄至數個家庭破碎,最不可能發生的竟然發生了,我們說真是荒謬.

說到人生的荒謬,使我想起最近一套電影《127小時》,內容提及喜愛登山探險男主角攀山途中意外墮入深坑,被石塊壓住右手 127 小時的求生經過.在這期間,主角對過去生命的作出深刻的回顧和反思,最後以中國製造的劣質童軍刀壯士斷臀.對於這塊壓住他的大石,他有一段藉得我們反思的說話:「這大石一生都在等著我,我一生也等著它......我的一生就向著它走,我的 DNA 帶我到這裡,從我一出生,每一口氣,每一個動作都領我至此,這地表的縫隙.」但在看似荒謬的痛苦經歷,卻把平時逗人發笑,對生命毫不在乎的主角,引發出生存意志,成為一個死過翻生的「英雄」.

但說到荒謬,甚麼比上帝因著愛讓他獨生愛子成為肉身,死在十字架上更為荒謬呢?十字架的道理,對自以為聰明的人是愚拙,但對相信的人,卻是上帝的大能.因著愛,上帝願意被看成荒謬,以至在我們可以在荒謬的事情中,看到上帝的愛.是的,上主的作為,確實隱密難明,但上主在基督裡的愛,卻是顯而易見.基督愛我們,沒有放棄他的命途,甘心上耶路撒冷受死.

意外可以改變我們的形像,疾病可以改變我們的形做,歲月也自然改變我們的形像.就在這時候,讓我們轉目凝視「臉面明亮如日頭,衣裳潔白如光」的基督,因為相信基督的人,也有一日,像他擁有榮耀的身體,榮耀的臉相.

在結束時,讓我們再一次聽聽韓馬紹對命運或命途的睿見:我們不需要尋求死亡,死亡會找上門的,但我們要找到往死亡的路,使死亡顯得有意義；命運本身不能逃避,也不需響往,它不是沒有原委的隱密,因為它說明這世界和人類歷史的軌跡,都有意義.阿們.



\newpage

\section{人活著不是單靠食物}
\label{sec:92}
\textbf{朱建忠先生 - 人活著不是單靠食物}
\newline
\newline
\hyperref[sec:91]{< < < PREV SERMON < < <}
~
\hyperref[sec:toc]{[返目錄]}
~
\hyperref[sec:93]{> > > NEXT SERMON > > >}


馬太福音 4:1-11
\newline
\begin{longtable}{cl}
\hline
\hline
章節 & 經文 (和合本修訂版)\\
\hline
4:1 & \begin{tabularx}{0.7\textwidth}{X} 當時，耶穌被聖靈引到曠野，受魔鬼的試探。 \end{tabularx} \\ \\ \relax
4:2 & \begin{tabularx}{0.7\textwidth}{X} 他禁食四十晝夜，後來就餓了。 \end{tabularx} \\ \\ \relax
4:3 & \begin{tabularx}{0.7\textwidth}{X} 那試探者進前來對他說：「你若是神的兒子，叫這些石頭變成食物吧。」 \end{tabularx} \\ \\ \relax
4:4 & \begin{tabularx}{0.7\textwidth}{X} 耶穌卻回答說：「經上記著：『人活著，不是單靠食物，乃是靠神口裡所出的一切話。』」 \end{tabularx} \\ \\ \relax
4:5 & \begin{tabularx}{0.7\textwidth}{X} 魔鬼就帶他進了聖城，叫他站在聖殿頂上， \end{tabularx} \\ \\ \relax
4:6 & \begin{tabularx}{0.7\textwidth}{X} 對他說：「你若是神的兒子，就跳下去！因為經上記著：『主要為你命令他的使者，用手托住你，免得你的腳碰在石頭上。』」 \end{tabularx} \\ \\ \relax
4:7 & \begin{tabularx}{0.7\textwidth}{X} 耶穌對他說：「經上又記著：『不可試探主—你的神。』」 \end{tabularx} \\ \\ \relax
4:8 & \begin{tabularx}{0.7\textwidth}{X} 魔鬼又帶他上了一座很高的山，將世上的萬國和萬國的榮華都指給他看， \end{tabularx} \\ \\ \relax
4:9 & \begin{tabularx}{0.7\textwidth}{X} 對他說：「你若俯伏拜我，我就把這一切賜給你。」 \end{tabularx} \\ \\ \relax
4:10 & \begin{tabularx}{0.7\textwidth}{X} 耶穌說：「撒但，退去！因為經上記著：『要拜主—你的神，惟獨事奉他。』」 \end{tabularx} \\ \\ \relax
4:11 & \begin{tabularx}{0.7\textwidth}{X} 於是，魔鬼離開了耶穌，立刻有天使來伺候他。 \end{tabularx} \\ \\
[1ex]
\hline
\hline
\end{longtable}
日期:2011年4月3日

題目:人活著不是單靠食物

經文:馬太福音4:1-11

證道:朱建忠先生

引言:這是一段很有意思的經文,講到「耶穌被聖靈引到曠野,受魔鬼的試探」.在這段故事中,耶穌三次回答撒旦時都是引用舊約 申命記第六到第八章裡面話(4節＝申八3；7節＝申六16；10節＝申六13).當時,以色列人正面臨爭戰應許地的使命,但是他們多次忘了神的吩咐,神借 著摩西提醒他們「要謹慎,免得忘記耶和華你的神,不守他的誡命,典章,律例」.(Deu  8:11)這一次,神的獨生子耶穌基督,在即將開始他的使命之時,在曠野也遇到了試探.因為時間的關係,我們今天只講魔鬼對耶穌的第一個試探.

一,耶穌是站在人的處境中來面對試探Mat 4:3   那試探人的進前來,對他說:「你若是神的兒子,可以吩咐這些石頭變成食物.」在這裡,「你若是」意思不是「你可能是」,「或者是」,「也許是」,而是「既 然是」.它表達的是「你既然是神的兒子,那就將石頭變成食物以充飢吧.」魔鬼並沒有懷疑耶穌是神兒子的身份,它所要做也並不是要從耶穌那裡得到證明,而是 想借此破壞神的兒子與神的關係及阻撓耶穌進行彌賽亞的使命.耶穌作為神的兒子,三位一體上帝的第二位格,當然有能力將石頭變成食物.在迦南的婚宴中,耶穌 曾將水變為酒；在伯賽大城邊的曠野地,耶穌曾將五餅二魚餵飽五千人.更何況,耶穌作為創始成終的神,「衪說有,就有；命立就立」.(詩33:9)衪甚至可 以使無變為有.但,既然是神的兒子又有何必要將石頭變成食物?既然是神的兒子又何必聽從撒旦的吩咐?不過,耶穌的回答是卻站在人的處境當中,他沒有抽離出 這個受苦的世界:「人活著,不是單靠食物.」聖經在此之前,已經講到當時的背景:耶穌經過四十天禁食後,他餓了.(Mat 4:2)作為人性的耶穌,他是真實的餓了.如果我們不能先明白這一點,我們就無法知道這試探的沉重.

例:一次探望的經歷.(略)

二,試探的深層次是針對人性的慾望這位中年弟兄的兒子,是以金錢為他的信仰；他的妻子,一直是以她的兒子作為也的信仰；他本人,在掙扎了一番之後,也是 無奈的站到了他兒子的一邊.我相信,主仍有他自己的憐憫.在座的各位可能很難遇到我剛才講的情況,我們也不希望面對如此艱難的處境.但是在我們的日常生活 中,試探仍然存在.食物的試探,在這裡僅僅是一個外在的形式,而更深深層次的卻是對各種人性慾望的試探.「把石頭變成食物」不就是任意行使權能的試探! 「從殿上跳下來而有天使托住」不就是對於宣揚名聲的試探!「只要拜一拜魔鬼就可以得到榮華富貴」,不就是出賣靈魂的試探!人,可以為著爭奪巨額的財產而拼 個你死我活；可以為了保持權力的控制而不顧億萬民眾的生命；可以以國家的口號而展開侵略殖民；也可以為著信仰的名義而進行聖戰行動.當你沒有權位,又沒有 機會晉陞的時候,或許權位對你構不成試探；當你金錢上沒有缺乏,又沒有機會面對巨額財富的時候,可能金錢對你也構不成試探；當你婚姻幸福,從來都不感到孤 單,婚外情可能也不會成為你的試探.但只要這些慾望存在,一旦時機來臨,試探也就隨之而來.魔鬼很會瞅機會,而且他的試探也往往真假難辨.亞當和夏娃原來 在伊甸園,我想他們生活的一定很幸福.不料一天,蛇對夏娃說:「神豈是真說不許你們吃園中所有樹上的果子嗎?」夏娃就開始有點暈頭轉向了.蛇接著對夏娃 說:「你們不一定死；  因為神知道,你們吃的日子眼睛就明亮了,你們便如神能知道善惡.」夏娃就開始動心了,她又看到「那棵樹的果子好作食物,也悅人的眼目,且是可喜愛的,能使 人有智慧,就摘下果子來吃了,又給他丈夫也吃了.」人的慾望往往就是這樣一步一步的被引誘,最後去犯罪.

德國偉大作家哥德作品中的浮士德,「他感覺有兩種意念在內心中搏鬥:一個要執拗地守著塵世, 沉溺在迷離的愛慾之中；   另一個要猛烈地離開凡塵,向一個崇高的境界飛馳.」惡魔見到浮士德正處在這種思想鬥爭當中後,便馬上找上門來,表示願意替浮士德解除煩惱,尋找歡樂,但條 件是一旦浮士德表示滿足的時候,就要成為他的僕人.在故事中,我們看到浮士德經過的各個的追求,但凡是從肉體慾望出發的,到來都得不到真正的滿足.主耶穌 曾在敘加城邊的一個井旁遇上一個中午來打水的撒馬利亞婦人.當時,主耶穌覺得渴了,就向這個婦人要水喝.這個婦人覺得奇怪,一個猶太人怎麼向撒瑪利亞人要 水喝.耶穌便對她說:「你若知道神的恩賜,和對你說『給我水喝』的是誰,你必早求他,他也必早給了你活水.」於是這個婦人便轉而向他要活水喝.但耶穌的回 答卻是:「把你的丈夫叫過來.」這話大概只有這位撒瑪利亞的婦人最能明白,也最感到震驚.但她先是裝糊塗,她回答說:「我沒有丈夫.」耶穌卻對她說:「你 說沒有丈夫是不錯的.但  你已經有五個丈夫,你現在有的並不是你的丈夫.」我想或許,在我們心中都有那一個稱不上丈夫的丈夫.但,主耶穌卻叫我們要帶他到主的面前來.並且他也告訴 我們,「凡喝了那井裡的水,還要再渴；人喝了我所賜的水就永遠不渴.並成為泉源,直流到永生.」 (Joh  4:13－4)人有肉體的需要,但人同樣有心靈的需要.當人以擴張肉體的需要來代替心靈的需要時,結果卻只能是喝了還要再渴.奧古斯丁說「人心中有一真 空,無法填滿.人若不安息在主的裡面,就永遠得不到安息.」

三,    得勝試探是因為我們的「所是」對面魔鬼的試探,主的回答並不只是消極的否定,「人活著不是單靠食物」,他還指出了積極的出路:「人活著是靠神口裡所 出的一切話.」這一處聖經原文的出處,是在申命記的第八章第三節.當時神藉著衪的僕人摩西提醒以色列民:「他苦煉你,任你飢餓,將你和你列祖所不認識的嗎 哪賜給你吃,使你知道,人活著不是單靠食物,乃是靠耶和華口裡所出的一切話.」但這句話,到底怎麼來理解呢?　　申命記第八章的第11－14節給我們解 答:「「你要謹慎,免得忘記耶和華你的神,不守他的誡命,典章,律例,就是我今日所吩咐你的；恐怕你吃得飽足,建造美好的房屋居住,你的牛羊加多,你的金 銀增添,並你所有的全都加增,你就心高氣傲,忘記耶和華你的神,就是將你從埃及地為奴之家領出來的.」其實,對人而言,當我們有的時候,能的時候,恐怕 「心裡說:『這貨財是我力量,我能力得來的.』」(Deu8:17)而把那位「厚賜百物的神」給忘了.

其實,魔鬼的試探,說穿了就是要達到使我們不聽從神的教訓,疏遠與神關係的目的.感謝神,我們之所以今日可以在這裡聚 會,不是因為我們的「所能」,而是因為我們的「所是」.神在造人的時候,吹了一口氣在他裡面,人就成了有靈的活人.從這一層面上來講,人是因與神的關係而 成為人,人是信仰的存在.魔鬼的試探無非就是在「你若是...就...」裡面耍花招；無非就是想讓耶穌停滯於,甚至是濫用神性超越性的一面,而忘記衪作為受苦彌賽 亞的使命.但耶穌的回答卻並非著眼於「上帝兒子」的權能而是在「上帝兒子」的順服與信靠上.同樣,基督徒的「所是」並非是因著我們的名聲,能力,或地位, 而是因神恩典的揀選.以弗所書一章四到五節告訴我們:「神從創立世界以前,在基督裡揀選了我們,使我們在他面前成為聖潔,無有瑕疵；又因愛我們,就按著自 己意旨所喜悅的,預定我們藉著耶穌基督得兒子的名分」.我們的社會總是過多的看人的所能,你在學校工作是這樣,在公司上班也是如此.但,也有一些地方卻是 看人的「所是」,最明顯的是家裡.有哪個遠方的遊子會嫌棄慈母的叮嚀?又有哪個真正的丈夫會離棄已經花容不再的髮妻?當浪子變賣家產,在外面吃喝玩樂而耗 盡了產業,最後失魂落魄的向父家走回時,他的慈父給他的不是斥責與嫌棄,乃是遠遠看見,就跑過去連連的與他親嘴,並叫僕人拿來上好的袍子,鞋子和戒指給他 穿戴.詩篇103篇中講:「父親怎樣憐恤他的兒女,耶和華也怎樣憐恤敬畏他的人!因為他知道我們的本體,思念我們不過是塵土.」 (Psa  103:13－14)

結語:　　最後,我想以之前的一次經歷作結束.記得有一次,在國內的時候,我與溫州的幾位教會同工一起在餐館裡面吃飯.當飯 菜預備好了,我們正準備開始謝飯禱告的時候,忽然看到這個餐館的筷子上都寫著一句中國的古語:「民以食為天」.我當時就想,若我們基督徒開一個餐館,那麼 上面寫著的就應該是:「人活著不是單靠食物」了.其實,在我們每次做謝飯禱告的時候,不正是在提醒自己:人活著不能以食為天,以物為天,或是以己為天,而 是要靠耶和華口裡所出的一切話.」

禱告:天父,我們感謝你,藉著今天經文提醒我們,「人活著,不是單靠食物,乃是靠你口裡所出的一切 話」,因為你是賜日用飲食的主,我們的動作,存留,氣息都在乎你,你也是領我們出埃及為奴之地的神.求你憐憫我們,幫助我們,不叫我們陷在試探當中,救我 們脫離兇惡,因為國度,權柄,榮耀全是你的,直到永遠.阿們.



\newpage

\section{起跑線上定輸贏}
\label{sec:93}
\textbf{林豪恩先生 - 起跑線上定輸贏?}
\newline
\newline
\hyperref[sec:92]{< < < PREV SERMON < < <}
~
\hyperref[sec:toc]{[返目錄]}
~
\hyperref[sec:94]{> > > NEXT SERMON > > >}


約翰福音 3:1-17
\newline
\begin{longtable}{cl}
\hline
\hline
章節 & 經文 (和合本修訂版)\\
\hline
3:1 & \begin{tabularx}{0.7\textwidth}{X} 有一個法利賽人，名叫尼哥德慕，是猶太人的官。 \end{tabularx} \\ \\ \relax
3:2 & \begin{tabularx}{0.7\textwidth}{X} 這人夜裡來見耶穌，對他說：「拉比，我們知道你是由神那裡來作老師的；因為你所行的神蹟，若沒有神同在，無人能行。」 \end{tabularx} \\ \\ \relax
3:3 & \begin{tabularx}{0.7\textwidth}{X} 耶穌回答他說：「我實實在在地告訴你，人若不重生，就不能見神的國。」 \end{tabularx} \\ \\ \relax
3:4 & \begin{tabularx}{0.7\textwidth}{X} 尼哥德慕對他說：「人已經老了，如何能重生呢？豈能再進母腹生出來嗎？」 \end{tabularx} \\ \\ \relax
3:5 & \begin{tabularx}{0.7\textwidth}{X} 耶穌回答：「我實實在在地告訴你，人若不是從水和聖靈生的，就不能進神的國。 \end{tabularx} \\ \\ \relax
3:6 & \begin{tabularx}{0.7\textwidth}{X} 從肉身生的就是肉身；從靈生的就是靈。 \end{tabularx} \\ \\ \relax
3:7 & \begin{tabularx}{0.7\textwidth}{X} 我說『你們必須重生』，你不要驚訝。 \end{tabularx} \\ \\ \relax
3:8 & \begin{tabularx}{0.7\textwidth}{X} 風隨著意思吹，你聽見風的聲音，卻不知道是從哪裡來，往哪裡去；凡從聖靈生的也是如此。」 \end{tabularx} \\ \\ \relax
3:9 & \begin{tabularx}{0.7\textwidth}{X} 尼哥德慕問他：「怎麼能有這些事呢？」 \end{tabularx} \\ \\ \relax
3:10 & \begin{tabularx}{0.7\textwidth}{X} 耶穌回答，對他說：「你是以色列人的老師，還不明白這些事嗎？ \end{tabularx} \\ \\ \relax
3:11 & \begin{tabularx}{0.7\textwidth}{X} 我實實在在地告訴你，我們所說的是我們知道的，我們所見證的是我們見過的，你們卻不領受我們的見證。 \end{tabularx} \\ \\ \relax
3:12 & \begin{tabularx}{0.7\textwidth}{X} 我對你們說地上的事，你們尚且不信，若對你們說天上的事，如何能信呢？ \end{tabularx} \\ \\ \relax
3:13 & \begin{tabularx}{0.7\textwidth}{X} 除了從天降下的人子，沒有人升過天。 \end{tabularx} \\ \\ \relax
3:14 & \begin{tabularx}{0.7\textwidth}{X} 摩西在曠野怎樣舉蛇，人子也必須照樣被舉起來， \end{tabularx} \\ \\ \relax
3:15 & \begin{tabularx}{0.7\textwidth}{X} 要使一切信他的人都得永生。 \end{tabularx} \\ \\ \relax
3:16 & \begin{tabularx}{0.7\textwidth}{X} 「神愛世人，甚至將他獨一的兒子賜給他們，叫一切信他的人不致滅亡，反得永生。 \end{tabularx} \\ \\ \relax
3:17 & \begin{tabularx}{0.7\textwidth}{X} 因為神差他的兒子到世上來，不是要定世人的罪，而是要使世人因他得救。 \end{tabularx} \\ \\
[1ex]
\hline
\hline
\end{longtable}
日期:2011年4月10日

題目:起跑線上定輸贏?

經文:約翰福音3:1-17節

證道:林豪恩先生

一．    引言現今社會上有幾句說話深深地影響著父母,或者阻嚇著人們不敢做父母.出自香港的其中一句是:「原來,養大一個BB要成四百萬」,出自內地的那句更要命:「不要讓您的孩子輸在起跑線上」.這些影響深遠的說話來自何處呢?原來,它們都是來自商業廣告.為甚麼「不要讓孩子輸在起跑線上」這句話對父母來說更要命呢?因為這句說話警告著父母,就算擁有四百萬,也不能保證你的孩子不會輸在起跑線上,因此鞭策著家長們前赴後繼地為孩子四處奔波.「起跑線」從中學推前至小學,從小學推前到幼稚園,從幼稚園推前到胎教.在媒體的傳播下,在學習產品生產商及教育服務供應商的推波助瀾下,不但令家長們躁動不安,也引起內地教育界的討論,甚至教育官員的關注.

今天的福音經課記述耶穌和尼哥德慕的對話,他們似乎也在討論「起跑線上定輸贏」這項話題.一個人的終極命運,是否決定於起跑線上呢?堅持這種信念與否又意味著甚麼呢?

二．    經文

3 : 1    有一個法利賽人,名叫尼哥德慕,是猶太人的官.

3 : 2    這人夜裡來見耶穌,對他說:「拉比,我們知道你是由神那裡來作老師的；因為你所行的神蹟,若沒有神同在,無人能行.」

3 : 3    耶穌回答他說:「我實實在在地告訴你,人若不重生,就不能見神的國.」

3 : 4    尼哥德慕對他說:「人已經老了,如何能重生呢?豈能再進母腹生出來嗎?」

3 : 5    耶穌回答:「我實實在在地告訴你,人若不是從水和聖靈生的,就不能進神的國.

3 : 6    從肉身生的就是肉身；從靈生的就是靈.

3 : 7    我說『你們必須重生』,你不要驚訝.

3 : 8    風隨著意思吹,你聽見風的聲音,卻不知道是從哪裡來,往哪裡去；凡從聖靈生的也是如此.」

3 : 9    尼哥德慕問他:「怎麼能有這些事呢?」

3 : 10耶穌回答,對他說:「你是以色列人的老師,還不明白這些事嗎?

3 : 11我實實在在地告訴你,我們所說的是我們知道的,我們所見證的是我們見過的,你們卻不領受我們的見證.

3 : 12我對你們說地上的事,你們尚且不信,若對你們說天上的事,如何能信呢?

3 : 13除了從天降下的人子,沒有人升過天.

3 : 14摩西在曠野怎樣舉蛇,人子也必須照樣被舉起來,

3 : 15要使一切信他的人都得永生.

3 : 16「神愛世人,甚至將他獨一的兒子賜給他們,叫一切信他的人不致滅亡,反得永生.

3 : 17因為神差他的兒子到世上來,不是要定世人的罪,而是要使世人因他得救.

三．    釋經與應用經文記述了耶穌與尼哥德慕一次「夜話」的內容.他們這次「夜話」到底談論甚麼呢?我們聽到耶穌說:「

就不能見神的國」(v3)；過了一會又聽到耶穌說類似的話:「

就不能進神的國」(v5).看來,他們在談論信仰上一個非常重要的主題,關於「見神的國」或者「進神的國」的話題.

讓我們再聽清楚一點,原來耶穌的整句話是這樣說的:「

人若不重生,就不能見神的國」(v3).然後,耶穌換句話說進一步解釋:「

人若不是從水和聖靈生的,就不能進神的國」(v5).如此看來,耶穌認為「見神的國／進神的國」與「重生／從水和聖靈生」有著密切的關係.

尼哥德慕明白耶穌所說的話嗎?從尼哥德慕的反應,可見他滿腹疑惑,摸不著頭腦:「

人已經老了,如何能重生呢?豈能再進母腹生出來嗎?」(v4)就算耶穌進一步解釋,他只有驚訝得張大嘴巴說:「

怎麼能有這些事呢?」(v9)

對尼哥德慕來說,要透過「重生／從水和聖靈生」來「見神的國／進神的國」,是不可思議的事情.經文介紹尼哥德慕是法利賽人,但卻不是普通的法利賽人,他是猶太人的官,即是猶太議會的成員之一,也是律法教師.他的資歷及背景足以反映著當時法利賽群體,以及猶太領導層的信念.他們認為「見神的國／進神的國」的條件是甚麼呢?我們可以從施洗約翰對他們的指責看出端倪:

「

約翰看見許多法利賽人和撒都該人也來受洗,就對他們說:「毒蛇的孽種啊,誰指示你們逃避那將要來的憤怒呢?你們要結出果子來,和悔改的心相稱.不要自己心裡說:『我們有亞伯拉罕為祖宗.』我告訴你們,神能從這些石頭中給亞伯拉罕興起子孫來.」(馬太福音3:7-9)

法利賽人和撒都該人宣告:「我們有亞伯拉罕為祖宗」,其意味與那位駕車撞死人的「官二代」,對執法者說:「我爸是李剛」有點相似.他們都以為,輸贏在起跑線上已經決定了.一個人出生在那裡,就決定了這個人的終極命運.帶著這樣的信念,尼哥德慕怎能明白耶穌的說話呢?怎能不為耶穌的說話而感到驚訝呢?

尼哥德慕不明白,尼哥德慕驚訝.然而,耶穌卻認為他不應該驚訝,也不應該不明白.他挑戰尼哥德慕說:「

你是以色列人的老師,還不明白這些事嗎?」(v10)為甚麼尼哥德慕應該明白呢?耶穌解釋:「

摩西在曠野怎樣舉蛇,人子也必須照樣被舉起來,要使一切信他的人都得永生」(v14-15).耶穌告訴尼哥德慕,現在說的並不是甚麼新鮮的事情,其實你們的先祖早已經歷過,你們熟讀的律法書也有記載.摩西在曠野舉蛇的事件,在民數記21章4-9節找到:

「

他們從何珥山起行,繞過以東地往紅海那條路走.在路上,百姓心中煩躁.百姓向神和摩西發怨言,說:「你們為甚麼把我們從埃及領上來,使我們死在曠野呢?這裡沒有糧食,沒有水,我們厭惡這淡而無味的食物.」耶和華派火蛇進入百姓當中去咬他們,於是以色列中死了許多百姓.百姓到摩西那裡,說:「我們有罪了,因為我們向耶和華和你發怨言.求你向耶和華禱告,叫蛇離開我們.」於是摩西為百姓禱告.耶和華對摩西說:「你要造一條火蛇,掛在杆子上.凡被咬的,一望這蛇就必存活.」摩西就造了一條銅蛇,掛在杆子上.凡被蛇咬的,一望這銅蛇就活了.」

如果尼哥德慕接受「

凡被蛇咬的,一望這銅蛇就活了」這事實,為何不能接受「

人子也必須照樣被舉起來,要使一切信他的人都得永生」(v14-15)呢?對尼哥德慕這位律法老師來說,這圖象是熟識的,這事件在歷史中是具體的,他要接受耶穌的啟示,還有甚麼難處呢?

如果我們能夠設身處地為尼哥德慕想一想,就會發現難處實在非常之大.要尼哥德慕接受「人子也必須照樣被舉起來,要使一切信他的人都得永生」,等於要求他放棄「我們有亞伯拉罕為祖宗」的優勢,等於要求他放棄「輸贏決定在起跑線上」的觀念.因為耶穌的說話為當時的猶太人帶來了兩項革命性的衝擊:第一,「見神的國／進神的國／得永生」原來並不取決於生在那裡,祖先是誰,而是「信他」；第二,「合資格人士」並不是局限於誰人的後裔,而是「一切人」.神的國竟然人人可進,永生的路竟然人人可走,上天的福氣竟然人人有份,對於以尼哥德慕為代表的猶太領袖來說是何等的震驚,是何等大的威脅,怎能不感到憤憤不平呢?今天,我們作為非猶太人基督徒,可能會認為當時猶太領袖的反應難以理解.然而,類似的觀念其實並不陌生.試想想今時今日的香港社會,如果來自任何家庭的小朋友也擁有均等的機會走進傳統名校,會否有人因此而感到威脅呢?如果人人也吃到那粒六千元的「糖」,會否有人因此而感到苦澀呢?

尼哥德慕被耶穌「衝擊」過後,有沒有因此而改變呢?約翰福音的作者在這裡賣個關子,卻留下伏筆.在第7章,尼哥德慕再次出場,彷彿是延續著「輸贏決定在起跑線上」這項討論.

7 : 32    法利賽人聽見群眾對耶穌這樣議論紛紛,祭司長和法利賽人就打發聖殿警衛去捉拿他.

7 : 45    警衛們回到祭司長和法利賽人那裡.他們對警衛說:「你們為甚麼沒有帶他來呢?」

7 : 46    警衛回答:「從來沒有像他這樣說話的!」

7 : 47    於是法利賽人說:「你們也受了迷惑嗎?

7 : 48    難道官長或法利賽人中有信他的嗎?

7 : 49    但這些不明白律法的眾人是被詛咒的!」

7 : 50    其中有尼哥德慕,就是從前去見過耶穌的,對他們說:

7 : 51    「不先聽本人的口供,查明他所做的事,難道我們的律法還定他的罪嗎?」

7 : 52    他們回答他說:「你也是出自加利利嗎?你去查考就知道,加利利是不出先知的.」

從尼哥德慕和他所屬的領導層關於如何處理耶穌的討論來看,他的態度似乎變得較有彈性了.他質疑祭司長和法利人對耶穌那種「未審先判」的立場:「不先聽本人的口供,查明他所做的事,難道我們的律法還定他的罪嗎?」面對尼哥德慕的挑戰,祭司長和長法利賽人更強硬地堅持「輸贏決定於起跑線上」的立場,他們反駁尼哥德慕說:「你也是出自加利利嗎?你去查考就知道,加利利是不出先知的.」他們以權威的聲音再次重申,一個人生在那裡,決定了這個人的終極命運.對於視「我們有亞伯拉罕為祖宗」為優勢的人,要放棄「輸贏決定於起跑線上」的觀念,真是談何容易.如果尼哥德慕個人要改變,除了更新自己的信念之外,還要面對他身處的群體加給他的巨大壓力.這種「人在江湖,身不由己」的壓力,或者生活在此時此地的我們也能體會.

到目前為此,關於「輸贏決定於起跑線上」的討論,約翰福音的作者讓我們看見祭司長和法利賽人所代表的體制性的強硬,也讓我們感受到尼哥德慕所代表的個體性的無奈.當難題彷如沒有出路的時候,作者立即記述一件看似沒有關係的事情,然而,就是這件看似沒有關係的事情,卻透露了化解難題的玄機.

8 : 2清早,他又回到聖殿裡.眾百姓都到他那裡去,他就坐下,教導他們.

8 : 3文士和法利賽人帶著一個犯姦淫時被捉的女人來,叫她站在當中,

8 : 4然後對耶穌說:「老師,這女人是正在犯姦淫的時候被捉到的.

8 : 5摩西在律法書上命令我們把這樣的女人用石頭打死.那麼,你怎麼說呢?」

8 : 6他們說這話是要試探耶穌,要抓到控告他的把柄.耶穌卻彎下腰,用指頭在地上寫字.

8 : 7他們還是不住地問他,耶穌就直起腰來,對他們說:「你們中間誰沒有罪,誰就先拿石頭打她!」

8 : 8於是他又彎著腰,用指頭在地上寫字.

8 : 9他們聽見這話,從老的開始,一個一個都走開了,只剩下耶穌一人和那仍然站在中間的女人.

8 : 10耶穌就直起腰來,對她說:「婦人,那些人在哪裡呢?沒有任何人定你的罪嗎?」

8 : 11    她說:「主啊,沒有.」耶穌說:「我也不定你的罪.去吧!從今以後不要再犯罪了.」

祭司長,文土和法利賽人所帶的眼鏡,把看到的人分為兩類,一類是亞伯拉罕的後裔,另一類就是不潔的外邦人；就算是同為亞伯拉罕後裔,他們也是只看到兩類人,一類是該用石頭打死的罪人,另一類就是如他們一樣自以為無罪的審判者.在這事件中,耶穌逼使他們換另一副眼鏡看人:「

你們中間誰沒有罪,誰就先拿石頭打她!」(8:7)放下只看見差異的眼鏡,換上看見共同之處的眼鏡,效果非常震撼:「

他們聽見這話,從老的開始,一個一個都走開了」(8:9).原來,我們都是一樣的人.被審判的是罪人,審判人的也不能否定自己是罪人.如果人人都是罪人,則人人都有著被拯救的需要.既然人人都有被拯救的需要,我們就明白第三章16-17節所宣告的了:「

神愛世人,甚至將他獨一的兒子賜給他們,叫一切信他的人不致滅亡,反得永生.因為神差他的兒子到世上來,不是要定世人的罪,而是要使世人因他得救.」人所需要的,不是定罪,而是拯救,這宣告正好呼應著第8章2-11節耶穌與祭司長和文士關於如何處理一位犯姦淫時被捉到的人的討論.

四．    總結看過第三章耶穌與尼哥德慕的對話,以及在第七章及第八章延伸的討論,我們能夠發現耶穌和當時的猶太領袖有甚麼分別嗎?猶太領袖心目中的起跑線,其實是用作「資格」的判別:誰有資格進入神的國,誰沒有資格進入神的國.祭司長,文士和法利賽人看重資格,耶穌則看重「需要」.拯救不是因為有些人有資格而出現,拯救是因為回應所有人的需要而出現.

耶穌與尼哥德慕的討論帶給我們不輕的挑戰:對於個人來說,我們戴著「資格」的眼鏡去看待別人,還是戴著「需要」的眼鏡去看待別人呢?例如作為顧主,我們可能認為自己有資格被家庭傭工尊重,然而,我們會否也願意接受對方也有被尊重的需要呢?對於社會來說,我們以資格為考慮去分配公共資源,還是以需要為考慮去分配公共資源呢?我們以資格為前題去制定醫療,教育,經濟等政策,還是以需要為前題去制定這些政策呢?例如,「那粒六千元的糖」,應該按資格還是需要分派呢?誰人有資格,誰人沒資格的爭論,若然產生撕裂社會,令鄰舍互相為敵的後果,又豈是社群之福呢?

我們不致於天真得以為人間就是天國,我們不得不認識我們身處的地方越來越重視資格,因此起跑線的影響力也越來越顯著.讓我們互相提醒,不要只看資格而忽視需要,盡力締造一個平等及友愛的環境.

一位出生時只有三隻手指的女孩,如果用起跑線來看,她被拋離多遠呢?她曾經為不能玩「包剪鎚」遊戲而沮喪,曾經擔心沒手指戴結婚戒指,也以為自己必然與鋼琴絕緣.今日,她是如何成為鋼琴演奏者的呢?她是如何成為中文大學其中一位音樂老師的呢?因為她身處的社會有醫療技術及服務,透過多次矯型外科手術,切割了另外五雙小手指出來；因為有鋼琴老師花心思和時間去教導這位「獨特」的學生.自幼患有青光眼及眼球萎縮症導致雙目完全失明的小男孩,在起跑線上落後多遠呢?他有機會發展天賦嗎?他後來是如何成為香港樂施會總幹事為社會服務的呢?別人都把他的成就歸功於他的個人毅力,但他卻念念不忘一位修女.在還未有充足學習資源給失明人士的年代,這位修女學習點字,為的是協助這位失明的青少年繼續學業.



\newpage

\section{井旁的邂逅}
\label{sec:94}
\textbf{鄧瑞強博士 - 「井旁的邂逅」}
\newline
\newline
\hyperref[sec:93]{< < < PREV SERMON < < <}
~
\hyperref[sec:toc]{[返目錄]}
~
\hyperref[sec:95]{> > > NEXT SERMON > > >}


日期:2011年4月17日

主題:「井旁的邂逅」

經文:約翰福音4:3-26

證道:鄧瑞強博士

各位弟兄姊妹,主內平安.

今日講道的經文,是約翰福音第四章主耶穌和撒馬利亞婦人在雅各井旁相遇的故事.

為了多一點了解撒馬利亞人,我好奇地在網上查看現代撒馬利亞人的生活.現代的撒馬利亞人人口很少,約有1000人.住在以色列的基利心山(Gerizim),Holon及納布盧斯(Nablus).而納布盧斯被視為巴勒斯坦的一個恐怖活動中心,經常被封鎖.撒馬利亞人認為自己是以色列十二支派裡的嗎哪西(Manasseh),以法蓮(Ephraim)和利未支(Levi)派的後人.他們至今仍持守著由出埃及時代傳下來的古老信仰.　　他們的信仰是一種古老的希伯來信仰,有五條基本信念: 只有一神,就是亞伯拉罕,以撒,雅各的神； 摩西是最後一位先知； 唯有摩西寫的「五經」(penteteuch)是聖經； 「基利心山」是神聖的地方,是聖殿的所在； 在末日,彌賽亞會來,惡人遭報應,義人被賞賜.歷來,他們都與猶太人不和.

按猶太人的講法,當以色列被亞述征服時,有些人留在以色列.這些人與外族通婚,失去了種族的純正性.後來,當以色列人被擄歸回時,他們在耶路撒冷重建聖殿,撒馬利亞人想加入,但被拒絕,以後,便形成兩民族的分裂.

但撒馬利亞人自己的講法是:亞伯拉罕獻以撒的地方是基利心山,這才是建聖殿的所在.撒馬利亞人的祭司是亞倫的後人,擁有祭司的永遠職分.他們在基利心山上的獻祭,才是純正的信仰.但猶太人的祭司,卻選了耶路撒冷作為聖殿的所在,在撒馬利亞人眼中,這是不正統的.聖殿的選址,涉及信仰的純正性的問題.現代的撒馬利亞人,左右做人難.他們夾在以色列和巴勒斯坦之間.在他們心中,他們是偏向「以色列」的.但在以色列眼中,他們是屬於巴勒斯坦的.故此,他們住的社區,屬於被隔離和封鎖的社區.他們的行動不完全自由,因為有軍事哨站看守著.有時,以色列政府會截斷他們的水源,他們會有飲水的難題.飲水的難題困擾著約翰福音裡的撒馬利亞婦人,也同樣困擾著今日的撒馬利亞人.今日,若一個正統的猶太人越過種種關卡,去接觸一個撒馬利亞人,是令人側目的,引起困惑的,令人不安的.而這情況,今日如是,過去也如是.當主耶穌接觸一個撒馬利亞婦人時,同樣是令人側目的,引起困惑的,令人不安的.那女人會困惑,這個猶太人為何會和我說話.甚至耶穌的門徒也會困惑,老師為何會與這撒馬利亞婦人傾談.然而,當人與人之間有不當的界線時,當民族與民族有不當的隔膜時,主耶穌卻主動走過這界線,走過軍事的關卡,走過建造起來的圍牆,走過種種審判的眼睛,走過社會既定的法則,他走向另一個人,明白她生命的悲苦,明白她心底的渴望,明白她的懇求.如此,井旁的邂逅故事便開始了.我們看看約翰福音的記載:

約 4:3    他就離了猶太,又往加利利去.

約 4:4    必須經過撒馬利亞,

約 4:5    於是到了撒馬利亞的一座城,名叫敘加,靠近雅各給他兒子約瑟的那塊地.

約 4:6    在那裡有雅各井；耶穌因走路困乏,就坐在井旁.那時約有午正.主耶穌從猶太地北上加利利,其實是「不必須」經過撒馬利亞的.很多猶太人便不會經過撒馬利亞.這裡寫「必須經過撒馬利亞」,是神學的寫法,是神學的「必須」,表達主耶穌的生命,神學上必然地走過種種圍牆,去接觸社會邊緣上的生命.主耶穌對人的愛,使他作出這必然的行動.主耶穌的救贖使命,使他作出這必然的行動.他「必須經過撒馬利亞」.

約 4:7    有一個撒馬利亞的婦人來打水.耶穌對她說:「請你給我水喝.」

約 4:8    (那時門徒進城買食物去了.)

約 4:9    撒馬利亞的婦人對他說:「你既是猶太人,怎麼向我一個撒馬利亞婦人要水喝呢?」原來猶太人和撒馬利亞人沒有來往.有一個女人出來打水.一般而言,打水不應在「午正」時分,因為太熱；也不應一個人出來打水,因為這太危險.這個女人如此打水,顯示她有難言之隱.主耶穌主動問她要水喝.種族的隔膜,社會的隔膜,宗教的隔膜,社會的隔膜,男女的隔膜,一下子被打破了.這舉動引起這女人的反感.女人稱耶穌為「猶太人」,是有一種「貶義」的,就正如有人稱呼中國人為「chinaman」一樣.這女人說他們撒馬利亞人和猶太人是誓死不相往來的.

約 4:10    耶穌回答說:「你若知道神的恩賜,和對你說『給我水喝』的是誰,你必早求他,他也必早給了你活水.」

約 4:11    婦人說:「先生,沒有打水的器具,井又深,你從哪裡得活水呢?

約 4:12    我們的祖宗雅各將這井留給我們,他自己和兒子並牲畜也都喝這井裡的水,難道你比他還大嗎?」

主耶穌用兩件事引起這女人的好奇.一,是「神的恩賜」,即「神給人的禮物」.二,是「耶穌的真正身分」.若她知道這兩件事,就會反求主耶穌給她活水.「活水」,在希伯來信仰裡,是一豐富的象徵.它象徵使人走正路的律法,象徵神話語的智慧,象徵甦醒人生命的聖靈,象徵生命的潔淨,象徵神醫治的大能,象徵救贖的終極福樂.主耶穌用信仰裡的終極福樂吸引著這女人的注意.這女人的反應,有點似一個香港人初聽福音時的反應.一個香港人可能說:「救恩是什麼來的?我不太需要你的福音了.我有我做人的方式.我有我的賺錢哲學.賺來的錢帶給我很大的快樂.我的生活很成功.我看不到你的福音如何能比我的生活哲學更好.」這女人的反應,與這反應相差不遠.這個女人如今稱呼主耶穌為「先生」,禮貌得多了.她說:「你沒有打水的器具,井又深,你從哪裡得活水呢?」這固然是當時實際的情況,主耶穌沒有器具,是打不到水的.但是,這回答也可以理解為「你那一套行不行」?撒馬利亞人以這個井自豪.他們認為這是先祖雅各留下的神聖信物.這個井,銘刻著先祖的偉大事蹟.時間消逝了,但這個井仍未乾涸.無數世代的人來了又去,這個井仍流著活水,滋潤著世人.這個井很「深」,這個井的故事很深,這個井承傳的信仰很深,這個井的神聖傳統很深.撒馬利亞人的生存,全賴這個井.這個深不可測的井如今仍滋潤著每一個撒馬利亞人,這位「先生」,你什麼也沒有,你看來與這傳統毫無關係,你能有益於我們嗎.女人問了一個很震撼的問題:「難道你比我們的祖宗雅各還大嗎?」我們的生活有很多神明,耶穌,你比這一切還大嗎?在投資的世界,我們有股神巴菲特(Warren Buffett),難道你比他還大嗎?在企業的運作裡,我們有台灣的「經營之神」王永慶,難道你比他還大嗎?「主耶穌,你比他們都大嗎?」這是很多人心底的問題.

約 4:13    耶穌回答說:「凡喝這水的還要再渴；

約 4:14    人若喝我所賜的水就永遠不渴.我所賜的水要在他裡頭成為泉源,直湧到永生.」

約 4:15    婦人說:「先生,請把這水賜給我,叫我不渴,也不用來這麼遠打水.」

主耶穌進一步和這女人討論「如何才能不再乾渴」的問題.當這問題一提出,這女人即刻明白,她的生命的確是常常乾渴,常常要打水,常常要用種種方式去止住乾渴,但常常是只有辛勞卻沒有長遠成果.她渴望一些永恆的東西,她渴望生命不再乾渴,而是有活水從生命中湧出來.

約 4:16    耶穌說:「你去叫你丈夫也到這裡來.」

約 4:17    婦人說:「我沒有丈夫.」耶穌說:「你說沒有丈夫是不錯的.

約 4:18    你已經有五個丈夫,你現在有的並不是你的丈夫.你這話是真的.」

要讓生命不再乾渴,便要正視生命,縱使這生命是殘缺的.主耶穌要這女人正視她的婚姻狀況,要這女人正視她生命的混亂和不忠,要這女人明白她生命的無所寄托.唯有這樣,這女人才能意識到新生命的可能.如今,若我們要得到神的禮物,神的救恩,神的活水,我們同樣要正視我們生命的混亂,正視我們生命的無依無靠.我們的乾渴在哪裡呢?

約 4:19    婦人說:「先生,我看出你是先知.

約 4:20    我們的祖宗在這山上禮拜,你們倒說,應當禮拜的地方是在耶路撒冷.」

約 4:21    耶穌說:「婦人,你當信我.時候將到,你們拜父,也不在這山上,也不在耶路撒冷.

約 4:22    你們所拜的,你們不知道；我們所拜的,我們知道,因為救恩是從猶太人出來的.

約 4:23    時候將到,如今就是了,那真正拜父的,要用心靈和誠實拜他,因為父要這樣的人拜他.」

約 4:24    　神是個靈,所以拜他的必須用心靈和誠實拜他.

若果撒馬利亞人的信仰傳統只承認摩西為先知,那麼,這女人稱耶穌為「先知」是一種突破.按她的信仰傳統,真正敬拜的地方是基利心山(Gerizim).她本來以自己的信仰傳統自豪,但當生命的乾渴被揭露時,她明白她需要「活水」,而這「活水」是她的傳統無法給予的.她追問真正的敬拜應在何處的問題.她追問在哪裡可找到活水.主耶穌的答案是:「那真正拜父的,要用心靈和誠實拜他」,或者譯:「在心靈裡和在真理裡敬拜他」.在約翰福音裡,耶穌就是「真理」,故此,在真理裡敬拜,即在基督裡敬拜.如今,真正的敬拜,重點不在外在的地方,而在你的心靈是否在基督裡.一個人的生命是否連於基督,這才是敬拜的重點.

約 4:25    婦人說:「我知道彌賽亞(就是那稱為基督的)要來；他來了,必將一切的事都告訴我們.」

約 4:26    耶穌說:「這和你說話的就是他!」撒馬利亞人的信仰傳統,等待救主的來臨.當救主來時,人才能得到終極的救贖.主耶穌對她說:我就是你等待的救主.

這是雅各井旁發生的故事.在創世記,雅各的愛情故事,正是發生在井旁.在井旁,他看見拉結(Rachel),並且一見鍾情.如今,在這個井旁,同樣發生了動人的愛情故事,是神和人的愛情故事,是耶穌接納罪人的愛的故事.這女人有無盡的渴望,她的感情生活混亂,她渴望愛,她渴望被接納.主耶穌也有渴望,他渴望讓這女人得到神的愛,他渴望這女人的渴望得到永恆的滿足,他渴望這女人的生命能活在甘霖裡.主耶穌這種愛的渴望是很大的,這渴望迫使他走上一條「必然的路」,必然的愛的路,必然的十架的路.在約翰福音裡,在十架上,耶穌說「我渴了」.驅使耶穌走上十架的,耶穌在十架上念念不忘的,就是這種渴望,愛的渴望.為了這渴望,縱使「撒馬利亞」是不潔的,他仍要走過去；縱使阻隔他與「撒馬利亞女人」接觸的圍牆是高大的,他仍要越過去；縱使兩者之間的宗教鴻溝是闊的,他仍要跨過去.因著主耶穌的渴望,因著主耶穌一句「請你給我水喝」,那女人被觸動了.一個久被違忘的渴望在她心裡重生了,她渴望「生命的活水」.她正視她的生命,她重新思考信仰,她在耶穌裡找到一種新的接近神的途徑.在耶穌身上,她體會愛,體會神的愛,體會到神在人間深深地愛她.在這份愛中,她找到活水,找到不再乾渴的秘密.今日,我們重溫這個井旁邂逅的故事.你會在這井旁,與耶穌相遇嗎?

但願榮耀歸於聖父,聖子,聖靈.阿們.



\newpage

\section{棕樹枝革命}
\label{sec:95}
\textbf{鄧瑞強博士 - 「棕樹枝革命」}
\newline
\newline
\hyperref[sec:94]{< < < PREV SERMON < < <}
~
\hyperref[sec:toc]{[返目錄]}
~
\hyperref[sec:96]{> > > NEXT SERMON > > >}


馬太福音 21:1-11
\newline
\begin{longtable}{cl}
\hline
\hline
章節 & 經文 (和合本修訂版)\\
\hline
21:1 & \begin{tabularx}{0.7\textwidth}{X} 耶穌和門徒快到耶路撒冷，進了橄欖山的伯法其時，打發兩個門徒， \end{tabularx} \\ \\ \relax
21:2 & \begin{tabularx}{0.7\textwidth}{X} 對他們說：「你們往對面村子裡去，會立刻看見一匹驢拴在那裡，還有驢駒同在一處，解開牠們，牽到我這裡來。 \end{tabularx} \\ \\ \relax
21:3 & \begin{tabularx}{0.7\textwidth}{X} 若有人對你們說甚麼，你們就說：『主要用牠們。』那人會立刻讓你們牽來。」 \end{tabularx} \\ \\ \relax
21:4 & \begin{tabularx}{0.7\textwidth}{X} 這事發生是要應驗先知所說的話： \end{tabularx} \\ \\ \relax
21:5 & \begin{tabularx}{0.7\textwidth}{X} 「要對錫安的兒女說：看哪，你的王來到你這裡，謙和地騎著驢，騎著小驢—驢的駒子。」 \end{tabularx} \\ \\ \relax
21:6 & \begin{tabularx}{0.7\textwidth}{X} 門徒就照耶穌所吩咐的去做， \end{tabularx} \\ \\ \relax
21:7 & \begin{tabularx}{0.7\textwidth}{X} 牽了驢和驢駒來，把他們的衣服搭在上面，耶穌就騎上。 \end{tabularx} \\ \\ \relax
21:8 & \begin{tabularx}{0.7\textwidth}{X} 許許多多的人把自己的衣服鋪在路上，還有人砍下樹枝來鋪在路上。 \end{tabularx} \\ \\ \relax
21:9 & \begin{tabularx}{0.7\textwidth}{X} 前呼後擁的人群喊著說：「和散那歸於大衛之子！奉主名來的是應當稱頌的！至高無上的，和散那！」 \end{tabularx} \\ \\ \relax
21:10 & \begin{tabularx}{0.7\textwidth}{X} 耶穌進了耶路撒冷，全城都驚動了，說：「這是誰？」 \end{tabularx} \\ \\ \relax
21:11 & \begin{tabularx}{0.7\textwidth}{X} 眾人說：「這是從加利利的拿撒勒來的先知耶穌。」 \end{tabularx} \\ \\
[1ex]
\hline
\hline
\end{longtable}
日期:2011年4月24日

主題:「棕樹枝革命」

經文:馬太福音21:1-11

證道:鄧瑞強博士

各位弟兄姊妹,主內平安.

按照基督教的聖曆,復活節前一個星期天是「棕枝主日」(Palm Sunday),是記念人群拿著棕樹枝迎接主耶穌進入耶路撒冷的日子.在某些教會,參與今日崇拜的人,會拿著棕樹枝進入教堂.表面看來,這是一個令人高興的日子.但在高興的背後,隱伏著艱難和苦難,因為今日同時是主耶穌的受難週的開始.進入耶路撒冷後,主耶穌和當時的各種勢力正面交鋒,神的世界和塵俗的世界激烈交盪,上帝的心意和世人的心意彼此角力,而最終,將主耶穌釘在十字架上.

我們看一看這段

經文:馬太福音21: 1-11

21:1　耶穌和門徒將近耶路撒冷,到了伯法其,在橄欖山那裡.

21:2　耶穌就打發兩個門徒,對他們說:「你們往對面村子裡去,必看見一匹驢拴在那裡,還有驢駒同在一處；你們解開,牽到我這裡來.

21:3　若有人對你們說甚麼,你們就說:『主要用牠.』那人必立時讓你們牽來.

21:4　這事成就是要應驗先知的話,說:

21:5　要對錫安的居民說:看哪,你的王來到你這裡,是溫柔的,又騎著驢,就是騎著驢駒子.」

21:6　門徒就照耶穌所吩咐的去行,

21:7　牽了驢和驢駒來,把自己的衣服搭在上面,耶穌就騎上.

21:8　眾人多半把衣服鋪在路上；還有人砍下樹枝來鋪在路上.

21:9　前行後隨的眾人喊著說:和散那歸於大衛的子孫!奉主名來的是應當稱頌的!高高在上和散那!」

21:10　耶穌既進了耶路撒冷,合城都驚動了,說:「這是誰?」

21:11　眾人說:「這是加利利拿撒勒的先知耶穌.」

這段記述散發著節日的氣氛,卻隱不住那種神和人相遇時會有的張力.包圍著這個節日一樣的情節的,是文字沒有表述出來的種種緊張.主耶穌進入耶路撒冷的日子,按聖經記載,是逾越節的前夕.逾越節是猶太人最重要的日子,最重要的宗教節日.逾越節記念的,是整個民族被神救贖的大事.這件救贖大事,不是一件心靈事件,不是一個人坐在深山裡的覺悟,而是一個民族在受苦,在哀求,而神聽到他們的哀求,解救他們的大事.昔日,這個民族在埃及受苦,他們受政治的奴役,受苦工的奴役,他們的哀求,代表著千千萬萬在極權下受苦者的哀求,代表著千千萬萬在工作中受苦的人的哀求.神聽到他們的哀求,神解救他們.這個救贖的日子,是一個宗教節日,也是民族解放的解放日.慶祝這個宗教節日,同時也紀念著民族的解放,或者可以說,紀念著人性的解放,紀念著人性尊嚴的重獲,紀念著能再次體會生命的意義和福樂.

就在這多種意義的逾越節裡,這意味著解放的節日裡,主耶穌進入耶路撒冷.這個日子,總是敏感的,總是充滿張力的,總是叫權力中心的人緊張的.當時的宗教領袖和當時的政治領袖,在這緊張的日子,總是千方百計地維持和諧的局面的.當時的宗教領袖,和政治力量聯手,希望一切和諧,他們繼續做他們的宗教生意,繼續享受既得利益.可以想像,當主耶穌進入耶路撒冷後,便進入聖殿,搗亂他們的生意,他們是多麼忿怒.當時的政治領袖,是極權的彼拉多,他以絕對優勢的羅馬軍事力量,粗暴鎮壓一切可能的反對勢力.可以想像,在這個容易出亂子的逾越節前夕,耶路撒冷是佈滿重兵的.宗教的力量在監視著一切,彼拉多的線眼亦注視著一切.主耶穌要進入耶路撒冷朝聖了.主耶穌在橄欖山上一條小村莊裡.沿著山路,慢慢下山,山腳就是客西馬尼園了.在這裡過了汲淪谷(Kidron Valley),就是耶路撒冷了.主耶穌離開耶路撒冷只有約2至3公里路程.按朝聖者的慣常做法,應走路入耶路撒冷.按主耶穌慣常的做法,也是走路.按這裡的距離,平常人也是走路.本來,若主耶穌走路入耶路撒冷,則只會湮沒在朝聖人群中,平靜地進去,平靜地出來.但是,主耶穌今次卻故意選擇坐驢入耶路撒冷.如此,一件平常事就變得不平常了.舊約的先知撒迦利亞曾經講過:「

錫安的民哪,應當大大喜樂；耶路撒冷的民哪,應當歡呼.看哪,你的王來到你這裡!他是公義的,並且施行拯救,謙謙和和地騎著驢,就是騎著驢的駒子.」(亞 9:9)先知的信息,留在人民的記憶裡.當人民看見有人騎著一隻小驢進入耶路撒冷時,他們的記憶就會被喚起,他們的希望就會被喚起,他們的力量就會被喚起.主耶穌刻意喚起他們的記憶,他們的希望,他們的力量.人是容易失憶的.當世界不斷叫我們快點賺夠錢買一個住宅單位時,我們忘記了,除此以外,我們還可以追求什麼.人是容易失去希望的.當生活的壓力日復一日地,重重地壓在我們身上時,我們不再希望生命有別的可能.人是容易失去力量的.當我們想不到有別的可能,當我們認為眼前的一切已是生命的全部時,我們便失去改變的力量.主耶穌放棄步行,刻意將先知講的,變成實在的圖畫,去喚起人民的記憶,希望,力量.除了先知撒迦利亞在預言中談過驢,先知以賽亞也曾在預言中談過驢.先知以賽亞說:「

天哪,要聽!地啊,側耳而聽!因為耶和華說:我養育兒女,將他們養大,他們竟悖逆我.牛認識主人,驢認識主人的槽,以色列卻不認識；我的民卻不留意.」(賽1:2-3)這是人最大的悲哀.連牲畜也能認識主人,但人卻不認識生命的真正主人.我們的罪,是因我們忘了天地的主.社會的不義,是因人忘了上帝要求的正義.世界的不安,是因世人忘了神的誡命.主耶穌要用那隻驢時,派他的門徒說:「主要用牠」.聽起來,主耶穌才是這隻驢的「主人」.這隻驢認識牠的真正主人,但世人呢?他們認識主嗎?主耶穌騎著驢,進入耶路撒冷.他要提示世人,先知撒迦利亞講的預言,應驗了.神來到人間,進行救贖的日子到了.同時,他要提示世人,先知以賽亞的警告,又再發出了.一隻驢也能認識牠的主人,但我們認識神嗎?主耶穌的舉動是引起注意的.可以想像,有人向大祭司匯報,主耶穌的舉動.可以想像,彼拉多派出更多的軍隊,以防萬一.群眾的熱情,是可以想像的.這是逾越節,是最大的宗教節日,民眾的情緒是高漲的.他們見到周圍的軍隊,也見到騎著驢的主耶穌.見到軍隊,他們意識到他們活在鎮壓之下.見到主耶穌,他們記起先知的預言,他們燃起希望,他們渴望被解放.他們懷著高漲的情緒,迎接主耶穌進入耶路撒冷.聖經說:

21:8　眾人多半把衣服鋪在路上；還有人砍下樹枝來鋪在路上.

21:9　前行後隨的眾人喊著說:和散那歸於大衛的子孫!奉主名來的是應當稱頌的!高高在上和散那!「和散那」這個字,按字面解,即「請你救助我們」.這個字現在變成一個讚美的字眼.民眾認為,這就是解放我們的救主.但是,主耶穌就是他們心目中的那種救主嗎?在當時的群眾心中,他們渴求的,是一個戰爭的彌賽亞.要解放,就是打一場解放戰爭.他們以為,主耶穌就是這場戰爭的領袖.軍事的壓迫力量越大,他們這種渴求則越熱切.他們大聲歡呼,「和散那歸於大衛的子孫!奉主名來的是應當稱頌的!高高在上和散那!」在羅馬兵的耳中,這些歡呼聲可能意味著挑釁.在群眾的心中,這些歡呼聲可能同樣意味著挑釁.暴力對抗,一觸即發.在歡呼聲中,主耶穌騎著驢.他聽到群眾的歡呼,他也抗拒群眾的歡呼.多少人,在這些群眾的歡呼聲中失去了理性,也失去了靈性.一個人獨自不會做的事,在群眾的歡呼聲中便可能會做出來.一個學生不會打另一個學生,但若有一群同學在旁吶喊:「打他,打他」,則這個學生很可能打出他的一拳.一個人可能不會覺得「造反有理」,但當群眾吶喊「造反有理」時,則「造反」便變得「有理」了.這群眾向主耶穌歡呼,但是,也是群眾的力量,呼喊著釘耶穌在十架上.主耶穌在想,如何才能叫這班熱情的人,不以暴力,而以和平之力,去對抗羅馬的壓迫.對抗強權,不是以暴力,而是以信仰,那種相信上主會主持正義的信仰,那種相信真理來自上主而不是來自強權的信仰,那種相信正直的犧牲會結出仁愛的果子的信仰.縱使身邊的羅馬軍隊可能騎在戰馬上,主耶穌依然騎在驢子上.這是一大對比,也是一種對抗.想像一個畫面,背景是當日波斯尼亞和塞爾維亞軍隊在薩拉熱窩的激戰,前景是住在薩拉熱窩的一個波斯尼亞女子和一個塞爾維亞男子的擁抱.這是一大對比,也是一種對抗.這種對抗,是一種非軍事的對抗,一種精神意義的對抗,一種強弱懸殊的對抗.當然,所謂強弱懸殊,是指騎驢的是強者,騎馬的是弱者.可以想像,緬甸的昂山素姬站在荷槍實彈的軍隊前,強弱立見.強者是昂山素姬,她一無所懼.弱者是軍隊,他們之所以荷槍實彈,正因為他們恐懼.群眾揚起棕樹枝,他們渴望發起棕樹枝革命.主耶穌卻騎著驢,他想著驢的革命.主耶穌明白他進入耶路撒冷的目的.他要對抗假宗教對人的壓迫,他要對抗不正義的制度對人的壓迫,但他不是用暴力的方式,而是藉著向人展示真愛的方式.這種方式不會一下子改變世界,卻讓所有人看到可以改變世界的力量何在.近日,有一則小新聞,一個70多歲的老人家,自力更新,在街頭賣「雞蛋仔」,卻屢被檢控,甚至無錢即時繳交罰款.我想,沒有人想犯法的.香港的舖租這麼貴,這樣的一個弱勢的人的弱勢經營,怎能捱得下去?在這高壓社會,要自力更生,這個老人家可以怎樣做?聽說,有一間酒樓的負責人,願意在酒樓的範圍內,劃出位置,免費讓這位老人家繼續經營下去.這位酒樓的負責人,就像騎著驢子,在充滿軍隊的耶路撒冷街頭走過一樣.或者,有一日,地產商會留下地皮的一角,讓窮人在那裡謀生.或者,有一日,政府的城巿規劃能留下空間,讓窮人可以呼吸多一點空氣.這有可能嗎?至少,我們見到,有人騎著驢子經過.不要少看騎驢子的人.聖經說:

21:10　耶穌既進了耶路撒冷,合城都驚動了,說:「這是誰?」這裡的「驚動」,是一強烈用字,是用來表達地震的.一個和平的舉動,竟引發全城的震動.這到底是誰?

21:11　眾人說:「這是加利利拿撒勒的先知耶穌.」

啊,又是拿撒勒,這是窮鄉僻壤,沒有出過什麼好東西.但是,就是這樣的一個出身低微的人,他要讓人明白什麼是偉大.

但願榮耀歸於聖父,聖子,聖靈.阿們.



\newpage

\section{抹大拉的馬利亞也復活}
\label{sec:96}
\textbf{林豪恩先生 - 抹大拉的馬利亞也復活}
\newline
\newline
\hyperref[sec:95]{< < < PREV SERMON < < <}
~
\hyperref[sec:toc]{[返目錄]}
~
\hyperref[sec:97]{> > > NEXT SERMON > > >}


約翰福音 20:11-18
\newline
\begin{longtable}{cl}
\hline
\hline
章節 & 經文 (和合本修訂版)\\
\hline
20:11 & \begin{tabularx}{0.7\textwidth}{X} 馬利亞卻站在墳墓外面哭。她哭的時候，低頭往墳墓裡看， \end{tabularx} \\ \\ \relax
20:12 & \begin{tabularx}{0.7\textwidth}{X} 看見兩個天使穿著白衣，在安放耶穌身體的地方坐著，一個在頭，一個在腳。 \end{tabularx} \\ \\ \relax
20:13 & \begin{tabularx}{0.7\textwidth}{X} 天使對她說：「婦人，你為甚麼哭？」她對他們說：「因為有人把我主移走了，我不知道他們把他放在哪裡。」 \end{tabularx} \\ \\ \relax
20:14 & \begin{tabularx}{0.7\textwidth}{X} 說了這些話，她轉過身來，看見耶穌站在那裡，卻不知道他是耶穌。 \end{tabularx} \\ \\ \relax
20:15 & \begin{tabularx}{0.7\textwidth}{X} 耶穌問她：「婦人，你為甚麼哭？你找誰？」馬利亞以為他是看園子的，就對他說：「先生，若是你把他移了去，請告訴我，你把他放在哪裡，我去把他移回來。」 \end{tabularx} \\ \\ \relax
20:16 & \begin{tabularx}{0.7\textwidth}{X} 耶穌對她說：「馬利亞。」馬利亞轉過身來，用希伯來話對他說：「拉波尼！」（「拉波尼」就是老師的意思。） \end{tabularx} \\ \\ \relax
20:17 & \begin{tabularx}{0.7\textwidth}{X} 耶穌對她說：「不要拉住我，因為我還沒有升上去見我的父。你到我弟兄那裡去告訴他們，我要升上去見我的父，也是你們的父，見我的神，也是你們的神。」 \end{tabularx} \\ \\ \relax
20:18 & \begin{tabularx}{0.7\textwidth}{X} 抹大拉的馬利亞就向門徒報信：「我已經看見了主。」她又把主對她說的話告訴他們。 \end{tabularx} \\ \\
[1ex]
\hline
\hline
\end{longtable}
日期:2011年5月1日

題目:抹大拉的馬利亞也復活

經文:約翰福音20章11-18節

證道:林豪恩先生

一．引言

福音書記載了復活的耶穌多次向門徒顯現的片段,描述得較詳細的分別有與多馬的相遇(約翰福音20:24-29),在以馬忤斯路上與兩位門徒的相遇(路加福音24:13-35),在提比哩亞海邊與彼得的相遇(約翰福音21:1-23).耶穌向多馬顯現,讓他檢查釘十字架遺留的傷痕,以滿足他對證據的需要.耶穌向以馬忤斯路上的兩位門徒顯現,「從摩西和眾先知起,凡經上所指著自己的話都給他們講解明白」(路24:27),證明「基督這樣受害,又進入他的榮耀,豈不是應當的嗎?」耶穌在提比哩亞海邊向彼得顯現,呼喚彼得回來跟從他,委以牧養群羊的重任.今天的福音經課敘述復活的耶穌向抹大拉的馬利亞顯現,他們的相遇,與上述的幾個片段並不相同,不是提供證據,不是解釋事件,也不是委以重任.如果我們留心經文,就會發現耶穌與抹大拉的馬利亞的相遇是充滿情感的,可以說,他們在情感世界中相遇.

二．經文

20 : 11馬利亞卻站在墳墓外面哭.她哭的時候,低頭往墳墓裡看,

20 : 12看見兩個天使穿著白衣,在安放耶穌身體的地方坐著,一個在頭,一個在腳.

20 : 13天使對她說:「婦人,你為甚麼哭?」她對他們說:「因為有人把我主移走了,我不知道他們把他放在哪裡.」

20 : 14說了這些話,她轉過身來,看見耶穌站在那裡,卻不知道他是耶穌.

20 : 15耶穌問她:「婦人,你為甚麼哭?你找誰?」馬利亞以為他是看園子的,就對他說:「先生,若是你把他移了去,請告訴我,你把他放在哪裡,我去把他移回來.」

20 : 16耶穌對她說:「馬利亞.」馬利亞轉過身來,用希伯來話對他說:「拉波尼!」(「拉波尼」就是老師的意思.)

20 : 17耶穌對她說:「不要拉住我,因為我還沒有升上去見我的父.你到我弟兄那裡去告訴他們,我要升上去見我的父,也是你們的父,見我的神,也是你們的神.」

20 : 18抹大拉的馬利亞就向門徒報信:「我已經看見了主.」她又把主對她說的話告訴他們. 三．釋經與應用

20 : 11馬利亞卻站在墳墓外面哭.她哭的時候,低頭往墳墓裡看,

20 : 12看見兩個天使穿著白衣,在安放耶穌身體的地方坐著,一個在頭,一個在腳.

20 : 13天使對她說:「婦人,你為甚麼哭?」她對他們說:「因為有人把我主移走了,我不知道他們把他放在哪裡.」

20 : 14說了這些話,她轉過身來,看見耶穌站在那裡,卻不知道他是耶穌.

20 : 15耶穌問她:「婦人,你為甚麼哭?你找誰?」馬利亞以為他是看園子的,就對他說:「先生,若是你把他移了去,請告訴我,你把他放在哪裡,我去把他移回來.」

復活的耶穌與抹大拉的馬利亞相遇的故事,充滿著眼淚.經文說:「馬利亞站在墳墓外面哭」.一位女子孤伶伶地站在墳墓外面哭,此情此景,不用多問,都想像到是生離死別場面了.

馬利亞的眼淚,可能已經流了好幾天.兩天前,馬利亞看著耶穌被釘在十字架上的時候,她已經流下了心痛的眼淚；聽到耶穌在十字架上掙扎呼吸的痛苦呻吟,也流下了不忍的眼淚；耶穌在十字架上斷氣的時候,馬利亞哀慟地痛哭.這一切痛苦而且荒謬的事情發生在耶穌身上,馬利亞眼白白地看著,但毫無能力阻止,她的無助及傷痛非言語能夠形容.

安息日近了,大家趕忙把耶穌的屍體埋葬,馬利亞只能夠待安息日過去才再拿香膏去膏抹耶穌的身體.那一天的安息日特別漫長,每一分每一秒都是煎熬.馬利亞好不容易才等到安息日過去.天還未亮,馬利亞已經急不及待,帶著那哭得浮腫的雙眼,摸黑走向耶穌的墳墓,找機會好好地哀悼耶穌,讓自己為耶穌的遺體做點最後的事情.就在這個時候,本來已經傷心欲絕的馬利亞碰著最匪夷所思的事情,耶穌的身體竟然不見了.馬利亞飛奔回去把這消息告訴門徒,那兩位男性門徒竟然看過空墳墓就回去,剩下馬利亞一個人孤伶伶地對著空墳墓,既傷痛,又無助,除了流淚,有甚麼話可說呢?除了哭泣,還有甚麼事可做呢?

有聲音問馬利亞:「婦人,你為甚麼哭?」「因為有人把我主移走了,我不知道他們把他放在哪裡.」馬利亞哭訴說.又再有聲音問馬利亞:「婦人,你為甚麼哭?你找誰?」馬利亞哀求說:「先生,若是你把他移了去,請告訴我,你把他放在哪裡,我去把他移回來.」

出於善意的人可能會如勸導馬利亞:「人都已經死了,何必為一具遺體而糾纏呢?而且這是一具囚犯的遺體,難道你不怕被牽連嗎?其他門徒都躲藏起來了,你快回去吧!」

如此勸導馬利亞的人,可能不明白馬利亞與耶穌的關係,也體會不到耶穌對馬利亞的重要.馬利亞可能回答說:「耶穌不是壞人,他是我的恩人.我曾經被七隻惡鬼附身,它們控制著我,令我失去尊嚴和自主；它們把我和人群分隔,令我無人為伍,孤立無援.自從耶穌把我身上的惡鬼趕走後,我重獲自由和尊嚴,恢復感情和關係；自從我跟隨耶穌之後,她找到了同伴,找到了接受愛和付出愛的對象；自從我照顧耶穌及門徒後,找到了生存的意義.」對於馬利亞來說,耶穌不單單是上帝的兒子,也不單單是是她的老師,更重要的,耶穌是她的恩人,甚至是比家人更親密的親人.看著自己的恩人受委屈被嘲弄,看著自己的親人受苦受害,自己無能為力；現在,連遺體也失蹤,對馬利亞來說這是非常巨大的打擊.這樣的打擊足以摧毀馬利亞整個人.

第一,她跑去墳墓,是要哀悼耶穌,以香膏膏抹耶穌的遺體.現在,連為耶穌的遺體做一點事情的機會也沒有,連哀悼的機會也沒有,這是極大的哀傷和無助；第二,耶穌被殺,馬利亞失去了情感相依的對象,也失去了群體；第三,這同時意味著她的信仰被挑戰,她所跟隨的,那位一直教導愛與寬恕,曾經展現非凡能力的神子,竟然被判刑處死；第四,也意味著她人生方向的消失,她所跟隨及照顧的恩人已經不再存在.

耶穌的死,耶穌遺體的失去,對馬利亞的打擊,誰能明白呢?誰能體會呢?或者,天安門母親會明白多一些馬利亞的悲傷,他們的兒女死於非命的時候,他們也是連為遺體做一點事情的機會也沒有,連正式地哀悼自己兒女的機會也沒有.直至現在,他們的兒女仍背負著反革命暴亂的污名.有一位音樂人創作了一首歌表達她們的心聲,名為「無處告別」:

已經多少個清明

他們不讓我哭泣

不讓我去你墳前

那怕燒一柱香

還要多少個清明

我才能夠去看你

在你墳前的青草地

開始埋鍋造飯

我真想掘開這墓穴

看看那真實的眉睫

你永遠不要瞑上眼

除非已剜掉那膿血

我無處告別

六月飛雪漫天

我無處告別

大地靜默無言

我無處告別

呼天不應,叫地不聞.

馬利亞只是一位弱小女子,何來有辦法改變猶太議會對耶穌的判決呢?何來有能力阻止羅馬政權處死耶穌呢?現在,連哀悼的機會也沒有,馬利亞還能夠做甚麼呢?對馬利亞來說,天地從沒有如此昏暗,一切都彷彿完結了.在傷痛絕望之中,馬利亞憑著自己的眼睛,沒法看見那兩個坐在耶穌墳墓裡的白衣天使；憑著自己的耳朵,也沒法分辨得出耶穌的聲音.此刻,孤立無援的馬利亞,完全感覺不到自己有任何面對當下處境的能力.

除了馬利亞和天安門母親,還有多少人處身於無助和絕望之中呢?還有多少人在死蔭幽谷之中期望復活呢?

最近,在小女幼稚園派發的一本刊物上讀到這樣一個故事:「某晚,丈夫徹夜未歸,我苦等了整夜.隔天早上他一進門,劈頭第一句話竟說:「我要離婚.」我如晴天霹靂,完全摸不著頭緒,只能不斷問:「為甚麼?發生甚麼事?」丈夫甚麼也不說,態度卻很堅決,並且收拾行李準備離去.不知所措的我,像瘋了般跑到窗旁,打開窗子說:「你要離婚,我就從六樓跳下去!」丈夫竟冷冷的回答:「請便!」．．．我的生活陷入狂亂絕望,每天幾乎只能躺在床上流淚,沒辦法爬起來做任何事情．．．只有在孩子大哭時,才勉強擠出一絲力氣起來看看她.」(「幽谷之光」《蒲公英》145期)

上月,在某個電台節目聽到一個感人的訪問,受訪者名叫鄧英蘭.她本來是一位活躍好動的幼稚園老師,一九八八年發生了嚴重車禍.鄧英蘭清醒後,發現自己失去了左手和左腳,陷入極度絕望之中.躺在病床上,鄧英蘭用她剩下的一隻手,不止一次把插在身體上維持生命的喉管拔走.護士來勸她:「你不要再拔走喉管了.你再拔喉,我們就連你剩下的一隻手也捆綁起來.」生命是維持了下來,不過,失去了的一隻手和一隻腳,如何活下去呢?鄧英蘭心裡充滿悲慘,徬徨和絕望.出院回家後,她把自己鎖在房間,過著以淚洗臉的生活.

甚麼能夠為悲傷的人帶來安慰呢?甚麼能夠為無助的人帶來力量呢?甚麼能夠為絕望的人帶來希望呢?甚麼能夠為死蔭幽谷的人帶來曙光呢?

20 : 16耶穌對她說:「馬利亞.」馬利亞轉過身來,用希伯來話對他說:「拉波尼!」(「拉波尼」就是老師的意思.)

當馬利亞失去了看見天使的洞察力,也失去了認出耶穌的觸覺的時候,耶穌用她熟識的聲音呼喚她:「馬利亞.」這親切的呼喚超越了馬利亞的眼睛,越過了馬利亞的耳朵,直接進入馬利亞的心窩,馬利亞從心坎裡回應:「拉波尼!」以深情對待耶穌的馬利亞,耶穌同樣以深情回應.耶穌不再追問再馬利亞在哭甚麼,而是親切地呼喚她的名字.這親切的呼喚,安定了馬利亞驚惶失措的心,安慰了馬利亞哀痛絕望的靈,為馬利亞昏暗的天空劃出一道光線.

復活的耶穌為馬利亞昏暗的天空所劃出的那道光線,出現在「幽谷之光」故事的主角的身上:「某晚,我走進常去的食店,老闆娘突然端了盤小菜到我面前,我搖頭表示送錯了,老闆娘卻說:「你太瘦了,要多吃點；今天老闆娘請客,儘管吃!」她接著逗弄我懷裡的孩子說:「真可愛,多大啦?」孩子被逗得露出了笑臉,就這樣,我和老闆娘聊了起來．．．並不習慣和陌生人交談的我,不知為何,卻被她的親切和溫柔所吸引,和她談起話來也沒有壓力.談話間,我提及前夫的種種行徑,忍不住泛出淚水．．．她繼續默默地聽我說,直到最後,我以為她會像大部份人一樣只勸我想開一點,沒想到她卻說:「抱歉,我實在不知道該說話來安慰你,但我可以為你禱告嗎?」我一時不知該怎麼反應,便點了點頭,她牽著我的手,溫柔低語著.我其實不太了解她說了甚麼,只覺得一股暖流流過心中,眼淚就這樣止不住的落下．．．即使現在我仍像在黑暗幽谷,不知多久才能走出來,但我深知,在黑暗中有一束亮光引領我前進,終有走出幽谷的一天.」(「幽谷之光」《蒲公英》145期)

復活的耶穌呼喚馬利亞的親切聲音,也臨到鄧英蘭的心中,為她的生存燃起了希望.

在最無助,最絕望的時候,鄧英蘭的妹妹每天對她說:「媽媽著緊你,爸爸著緊你,家裡每一個人都著緊你.」她的男朋友每天來對她說:「上帝愛我,我會用上帝的愛來繼續愛你.」

20 : 17耶穌對她說:「不要拉住我,因為我還沒有升上去見我的父.你到我弟兄那裡去告訴他們,我要升上去見我的父,也是你們的父,見我的神,也是你們的神.」

20 : 18抹大拉的馬利亞就向門徒報信:「我已經看見了主.」她又把主對她說的話告訴他們.

按著這個段落的描述,我們可以想像兩個截然不同的景象:第一個景象是馬利亞一個人孤伶伶地在耶穌的墳墓外哭泣,第二個景象是馬利亞驚喜萬分地飛跑回去把喜訊告訴門徒.是甚麼令馬利亞產生如此強烈的改變呢?可能有人認為,因為耶穌復活,馬利亞失而復得,這安慰了她的哀傷.不過,這又未必是真相的全部.馬利亞的確想耶穌回來,因此她伸手要拉住耶穌,期望回復以往的一切:跟隨耶穌走遍大城小鄉,照顧耶穌和門徒的起居飲食.不過,復活的耶穌卻不讓馬利亞拉住,而且告訴她,自己復活不是要回到馬利亞身邊,而是要離開她,回到父神那裡.耶穌復活,對馬利亞來說,是耶穌回來了,但同時也是耶穌要離開了.那麼,耶穌復活,對馬利亞有甚麼意義呢?

馬利亞拿清晨著香膏跑去耶穌的墳墓,顯示出她只是去哀悼耶穌,完全沒有估計耶穌會復活.事實上,沒有人期待耶穌復活,把耶穌釘十字架的人沒有這樣的期待,就是跟隨耶穌出入多年的門徒也沒有這樣的期待.就算馬利亞找到耶穌的遺體,膏抹了耶穌的遺體,完成了她的哀悼,她很可能仍然擺脫不了耶穌的死帶來的沉痛打擊,很可能在失落及幻滅中渡過餘生.如果說耶穌的復活使馬利亞也復活,這很可能是因為耶穌的復活還了馬利亞一個公道,肯定了馬利亞的信仰,確認了馬利亞的人生方向.耶穌雖然被猶太議會判罪,被羅馬政權處死,但馬利亞從自己與耶穌的相遇及生活經驗中,無法否認的事實是,耶穌是趕走她身上七隻惡鬼的恩人,耶穌是愛與真理的老師,耶穌是道成肉身上帝的兒子.否定這些事實,即是否定她的親身經歷,即是宣佈強權戰勝了真理,即是宣佈世上沒有公道可言.耶穌復活,對馬利亞的意義,不在於耶穌回到她的身邊,而在於肯定還耶穌一個公道,還馬利亞一個公道,讓她肯定自己對耶穌的經驗及認識,對耶穌所啟示的真理的執著與堅持是正確的.這樣的肯定,足以讓馬利亞從哀傷和絕望中復活,確認自己過去以及餘下的日子為耶穌而生存是值得的.

四．總結

耶穌復活,馬利亞也復活,我們可以從馬利亞對耶穌復活的經歷中,得到甚麼啟發幫助處身於死蔭幽谷的人呢?有些人需要有人來到他們那裡,安慰他們,支持他們,好像耶穌親切地呼喚馬利亞的名字,又好像那位食店老闆娘向那位一個人帶著嬰兒的母親送上的關懷和代禱,又好像鄧英蘭的家人和男朋友以上帝的愛不離不棄的同行,使這些人能夠走出死蔭幽谷.另外一些人則需要還給他們公道,好像耶穌復活還給耶穌的公道,還給馬利亞的公道,還給真理和愛的公道,使他們在真理和愛中帶著盼望繼續生存下去.



\newpage

\section{看見真相的代價}
\label{sec:97}
\textbf{伍渭文牧師 - 看見真相的代價}
\newline
\newline
\hyperref[sec:96]{< < < PREV SERMON < < <}
~
\hyperref[sec:toc]{[返目錄]}
~
\hyperref[sec:98]{> > > NEXT SERMON > > >}


日期:2011年5月8日

題目:看見真相的代價

讀經:約翰福音9

證道:伍渭文牧師

看見真相的代價,首要是小心保護眼睛.眼睛出現異常的圖像或飛蚊時,立刻看眼科醫生.視網膜有撕裂的現象時,立時進醫院動手術,車往醫院途中要慢駛,不能太快.等候手術的晚上要全身躺直,鼻尖向上.手術後一星期內要俯著睡,坐著時不可讓鼻孔向上,此期間不能乘撘飛機,攀登高山,否則眼睛會因眼壓變大弄至目盲.我半年前右眼進行了視網膜手術,幾天後我把紗布眼拿開時,發現一片漆黑,我感到害怕(traumatic),經過一個多星期,黑暗漸漸轉成光明,心情愈來愈輕鬆和感恩.這段時間,我常常想起約翰福音九章那位生下來就失明的瞎子.

今天的福音經課,我們看到一位生來瞎眼討飯的人,給耶穌醫好後,受到當時的宗教領袖所盤問,最後他堅持自己看到的真理,被趕出所屬的社群會堂.他被驅逐,祇因他堅持看到的真相:我以前我是瞎眼的,現在被他醫好了,可以看見東西了,而「

從創世以來,未曾聽見有人把生來是瞎子的眼睛開了.這人若不是從神來的,甚麼也不能做.」(約九32)這人麼該從上帝而來.

1. 他看到一個屬靈的世界 (vv.1-12)我們要注意,聖經沒有描述這瞎子得到光明後,被這彩色繽紛的世界吸引而興奮,也沒有因為自己獲得光明而欣喜若狂,忙著告訴人:「我眼睛開了!我叫眼睛開了!」上星期我們談到那位中午在雅各井旁打水的撒馬利亞女人,當她發現耶穌是彌賽亞時,她立刻放下水罐子,急不及待跑進城中奔走相告:「

你們來看,有一個人將我素來所行的一切事都給我說出來了,莫非這就是基督嗎?」(約四29)

但這瞎子不同,第八,九節說鄰舍和認識他的人彼此對問:「

這不是那從前坐著討飯的人嗎?有人說:是他.又有人說:不是,卻是像他.他自己說:是我.」在他們再三追問下:你的眼睛是怎麼開的呢?他才如實把過程說出來.當他們追問:那人在那裡?他也是簡單的說:我不知道.

在整本舊約聖經,真的沒有記載瞎眼的被醫治,醫治瞎眼得見光明是上帝親自的工作,是彌賽降臨的記號.上帝在西乃山對摩西顯現,叫他帶領百姓離開埃及,摩西說他是拙口笨舌不能勝任,上帝回答:「

誰人使人口啞,耳聾,目明,眼瞎呢?豈不是我 and \#8212;耶和華嗎?」(出四11) 以賽亞描述當彌賽亞來臨時:「

那時,聾子必聽見這書上的話；瞎子的眼必從迷矇中得以看見.」(賽二十九18) 所以福音書記載耶穌施行的眾多神蹟,以瞎子被醫治得看見最多,因為這是彌賽亞的記號.「

從創世以來,未曾聽見有人把生來是瞎子的眼睛開了.這人若不是從神來的,甚麼也不能做.」

叫這瞎子震驚欣喜,不是他親眼見到了花草樹木,飛鳥白雲的世界,乃是他看見彌賽亞已經來到,萬有將被更新的世界,他的心被這認知攝住了,他的感動超越了情緒的興奮.對生出來是瞎眼的他,他每天討飯不是就等候這一刻嗎?因為彌賽亞來的時候,瞎子的眼會被打開,他盼望又盼望這天的來臨,然而當他眼睛明亮了,他又有一種反興奮的沉思.

作為主持婚禮的牧師,我偶然也遇見有人有這種反興奮的沉思.當然,結婚令多數的人興奮不已,因為贏得美人歸,找到如意郎君；他們憧憬如何弄一個畢生難忘浪漫的婚禮.新娘想著如何佈置會場,如何開始減肥穿美麗的婚紗.新郎想著開始練歌,在地毯的那一邊唱歌迎接新娘進場,或婚禮完畢抱起新娘回來致謝(這項危險動作,我建議不宜嘗試) .但有一些人沒有這麼興奮,他們一想到這是不能逆轉的抉擇,直到死亡才結束的盟誓,結婚後還有養兒育女的責任,整個世界剛剛開始,他們窺探到就快要進入到一個不能估計的世界,不是不歡喜,乃是嚴肅認真起來了.

劍橋大學中古和文藝復興時期英國文學教授,著名作家護教學祭酒魯益師(C.S. Lewis)形容自己信主那一刻是全英國最低沉和最不興奮的歸依者.(the most dejected and reluctant convert in all England) .1929年經過多年的上下求索,多方閱讀,師友砌磋,他來到要作抉擇的一刻；其實他不想面對他看到的真理,他曾把這些關在門外,因他知道若把門打開,後果無可估計,但他終於順服,面對不能逃避的真理 and \#8212;他相信上帝.「低沉和最不興奮的歸依者」的形容詞,不是不認真,乃非常認真,變得冷峻,因他知道門一打開,後果無可估計.

2. 他要個人面對這屬靈的世界 (vv.13-34)故事中這位得見光明的瞎子,個性鮮明.耶穌吐唾沫混合地上的泥土,抹在他眼上,並叫他往西羅亞池子去洗.聖經詳細描述瞎子如何得醫治,就像上帝如何帶來新的創造.「

耶和華神用地上的塵土造人,將生氣吹在他的鼻孔裡他就成了有靈的活人.」(創二7)這池子名叫西羅亞,奉差遣的意思,顯示耶穌就是奉差遣的彌賽亞.

瞎子得醫治,令反對耶穌的猶太人非常尷尬,也引起內部不同意見,甚至紛爭.一派說:這個人不是從神而來的,因為他破壞安息日.又有人說:一個罪人怎能行這樣的神蹟呢?反對耶穌的藉口似乎佔了上風,他們希望沒有發生過神蹟,瞎子根本不是真正的瞎子,他們問瞎子的父母:「

這是你們的兒子嗎?你們說他生來是瞎眼,如今怎麼能看見了呢?」(約九19)父母如實說,孩子生來是瞎眼的,但他們怕猶太人:「

因猶太人已經商議定了,若有識認耶穌是基督的,要把他趕出會堂.」(約九22)會堂是猶太人的身份認同,趕離會堂就是失去族群的身份,是非常嚴厲的懲罰,他們叫猶太人直接問瞎子,因他已經成人了.

瞎子要自己面對信仰,畢竟,信仰是個人與上主之間的關係,我這個人面對上帝,我自己要為自己的良心負責任,我的父母不能為我作任何決定,沒有一個人可以為我作決定,這是更正教一個非常重要的觀念 and \#8212;在上帝面前(coram deo):教會的傳統,其他聖徒的功德,不能幫助你,我們在上帝面前要孤身面對他,正如我們要孤身面對死亡.

馬丁路德在1521年首次為他所傳講所寫的福音信仰接受公開審訊,被要求收回所有與當時教會不能接受的言論.他說:「我的良心已被上帝的話語所捊擄,我不能也不會收回任何東西,因為違背良心是不安全和不對的.」(

My conscience is captive to the Word of God. I cannot and I will not retract anything, because it is neither safe nor right to go against conscience.

) 同樣,這位復見光明的瞎子沒有在巨大的壓力下退縮,他為著自己的良心,堅持真理；他和猶太人辯論,最後總結陳詞:「

他開了我的眼睛,你們竟不知道他從哪裡來,這真是奇怪!我們知道神不聽罪人,惟有敬奉神,遵行他旨意的,神才聽他.從創世以來,未曾聽見有人把生來是瞎子的眼睛開了,這人若不是從神來的,甚麼也不能做.」(約九31)

最後瞎子被趕離會堂,他願意為信仰付上代價,因為他看見了的一個真實的屬靈世界,他要追尋這比眼睛看得見的世界來得更真實的屬靈世界.

3. 他要成為引導人認識基督的標記 (vv.35-41) 在故事開始的時候,耶穌看見這瞎子,走到他跟前醫好他,使他得著光明,故事結束時,耶穌聽見他被逐出會堂,後來遇見他(found him) and \#8212;尋找他,中間一段時間耶穌似乎不理會他,讓他獨自面對挑戰,熬煉他的信心.每個人的信心歷程不同,我們都有自己要走的路,但上帝的耳目總不離開我們.經過與猶太人的辯論,這人的信心預備好了,耶穌發出邀請.「

你信神的兒子嗎?他回答說:誰是神的兒子,叫我信他呢?耶穌回答說:你已經看見他,現在和你說話的就是他.他說:主啊,我信!就拜耶穌.」(約九35-38)

要我們從聖經選出一段感人至深的對話,相信其中一段非這段莫屬了.我們之前說過:瞎子眼睛被醫好後,他沒有忙著興奮狂喜,到處告訴人,反而變得深峻,沉思嚴肅起來,因為瞎子看到了一個更重要的屬靈的世界.

從來沒有人把生來就瞎眼的醫好,整本舊約沒有記載這樣的神蹟,但有預言過:彌賽亞來臨的時候神的國降臨的時候,瞎子會看見,啞巴會說話,我的眼被醫好了,豈不是彌賽亞已來臨了,但他那在那裡呢?我可以親眼見到他嗎?誰是神的兒子,叫我信他呢?你已經看見他,現在和你說話的就是他.主啊,我信!就拜他.在故事開始時,瞎子的不幸成為門徒的困擾,為何這樣呢?是他自己犯了罪,還是父母犯了罪呢?門徒所問的問題,也許瞎子本人也多次問過,耶穌沒有提供這苦難的原因,他沒有解釋苦難,他要解決苦難,這苦難要成為標記 (sign),顯出上帝的作為來,瞎子被醫好的神蹟就像約翰福音其他的六個神蹟或標記,指向一個目標 and \#8212;耶穌是彌賽亞,上帝的兒子.醫好生來是瞎子的神蹟,是標記,叫我們認識耶穌是世界的光,帶給人光明,使我們不再走在黑暗中.

每一個信耶穌的人,都是上帝的標記,叫人認識上帝,引人歸向上帝.路標可能有不同大小,我們認為小的路標,也可以為上帝所用,影響整個歷史,我們千萬不要自我菲薄.現代管理學之父 彼得.杜拉克 (Peter Drucker,1909-2005) 接受一份雜誌訪問,談及甚麼是有效的服待時說到馬丁路德一件逸事:

馬丁路德人生中有一次感到非常沮喪,他來到奧古斯丁修院的院長面前尋求幫助.院長對他說:馬丁弟兄,陷溺在沮喪中是罪.你們試想想,要是沒有這次見面,這位年青的修士將會如何呢?相對(路德)來說,院長的智力,甚至靈性都較遜色,但他說對了話.他花時間接見了路德,且利用了那一刻的機會.他沒有解答路德的靈性和神學問題,卻改變了他整個人,這是一樁有效服侍的範例.

(There was an episode in the life of Martin Luther when he was in deep despair. He went to his Augustinian prior, who said, ``Brother Martin, it is a sin to be in despair." For Martin, that was the important thing to say. What would have become of that young monk without that moment? The prior intellectually was surely his inferior, and probably spiritually, too, but he said the right word. He was available and used the moment of opportunity. It did not answer any of Martin"s spiritual and theological question, but it totally changed him. It was an example of effective ministry.)

總結:今天的福音經課提醒我們,有一個屬靈的真實世界,我們每一個人都要面對這世界,進入這世界,也要成為別人的路標,引導他們探索這世界.在結束的時候,我想分享去年崇基基督教文化節特邀講員北京人民大學何光滬教授,在崇基神學院周四晚會中作的見證.何教授是中國大陸著名研究基督教學者,繙譯了幾本重要的著作.他說學習英語不單是學習一個語言,對當時文革時期下放到農村的生活的他,英語是另外一個世界的召喚.在文革時期的農村,生活非常艱苦,不單是物質匱乏,更是精神極度苦悶.整個中文的世界就是虛偽的口號,仇恨的口號,世界好像停止了.他看到英文,好像看到另一個世界,他要尋求認識這世界.然而,農村沒有老師,他知道有一位英文老師在城中,為了學好英語,他行上幾天的山路,走出城市找老師.幸運的時候,剛有牛車到城中請求搭過便車,但牛車不好坐,山路崎嶇,幾次翻車,有一次掉落車下,傷得很重.但為了學好英語,認識另一個世界,他排除一切艱難,付上極大的代價.因為認識英語,接觸到有關基督信仰的文字,後來信了基督.他不甘心當時眼見到的悲慘世界,他響往一個更美的世界,他尋找又尋找,終於找到基督.

這兩個多星期,日本地震,海嘯和核電災使我們看到甚麼真相哩?我們可否藉這些災難,看到人生如此無常,繁榮如此脆弱,而得著智慧,追求那比肉眼看得見的世界更真實,更宏大的屬靈世界嗎?誰是神的兒子,叫我信他呢?你已經看見他,現在和你說話的就是他.主啊,我信!就拜他.阿們.



\newpage

\section{人生的馬拉松}
\label{sec:98}
\textbf{賀志勇博士 - 人生的馬拉松}
\newline
\newline
\hyperref[sec:97]{< < < PREV SERMON < < <}
~
\hyperref[sec:toc]{[返目錄]}
~
\hyperref[sec:99]{> > > NEXT SERMON > > >}


希伯來書 12:1-2
\newline
\begin{longtable}{cl}
\hline
\hline
章節 & 經文 (和合本修訂版)\\
\hline
12:1 & \begin{tabularx}{0.7\textwidth}{X} 所以，既然我們有這許多見證人如同雲彩圍繞著我們，就該卸下各樣重擔和緊緊纏累的罪，以堅忍的心奔那擺在我們前頭的路程， \end{tabularx} \\ \\ \relax
12:2 & \begin{tabularx}{0.7\textwidth}{X} 仰望我們信心的創始成終者耶穌，他因那擺在前面的喜樂，輕看羞辱，忍受了十字架的苦難，如今已坐在神寶座的右邊。 \end{tabularx} \\ \\
[1ex]
\hline
\hline
\end{longtable}
日期:2011年5月15日

題目:人生的馬拉松

經文:希伯來書12:1-2節

證道:賀志勇博士

1    我們既有這許多的見證人,如同雲彩圍著我們,就當放下各樣的重擔,脫去容易纏累我們的罪,存心忍耐,奔那擺在我們前頭的路程,

2    仰望為我們信心創始成終的耶穌(註:或作"仰望那將真道創始成終的耶穌").他因那擺在前面的喜樂,就輕看羞辱,忍受了十字架的苦難,便坐在神寶座的右邊.

一,引言 : 訓練有素

這次日本大地震讓全世界都震驚.非常有秩序!地震來了,沒有騷亂、沒有人聲鼎沸、沒有人奪路而逃.地震過了,地鐵停下來時,大家都整齊上車,沒有人鼓噪,沒有人亂擠.雖然路上塞車許多個鐘,有的司機甚至累得睡在駕駛座位上,卻沒有人按一下喇叭來催促.看到這些報導,我非常非常感動.再回頭看看我們自己,搶購日本明治奶粉；搶購碘鹽.我聽說武漢有個人搶購了6點5噸鹽,可以夠他們一家從唐朝開始吃,吃到現在,還吃不完.

面對災難,日本人為什麼如此鎮定?答案很簡單,平時訓練有素.日本經常發生地震,所以日本國民平時就已經訓練有素,鎮定,守秩序,平靜地面對危險.訓練有素,就是一種素質.

日本人面對災難訓練有素,那我們呢?我們面對生活是不是訓練有素呢?對於生活,我們是不是總是做一個業餘演員呢?每天上班下班,每天上學放學,每天餵奶洗尿片,你會不會覺得心累,你還有心思去訓練嗎?香港人把這種狀態叫做""心疲",就是心裡面累了.有的人踏入中年了,就有了中年危機.不像年輕時有那種傻勁,去闖去拼；總覺得生活失去了新鮮感,開始有點苦悶了,開始有婚外情,有七年之癢了.有的人做論文做到一半,越寫論文就越覺得煩,有論文危機了.找工作如是,做人如是,做基督徒也如是.我們基督徒也會經歷信仰低潮.返教會返慣了,聽道聽膩了,牧師一開口,哎,我都知道你要講什麼啦.靈修麻木了,不再每日讀經了,每況愈下,一年不如一年.漸漸屬靈生命就乾枯了,"burn out"了.最近有個傳道人跟我分享,她說每次小組快要結束,她就請大家講一下代禱事項,然後請大家回去禱告,她說沒幾個人回家會禱告.我說你怎麼知道啊,她說"因為我自己回家也經常忘了代禱".所以傳道人最怕兩種狀態,一種就是等退休；還有一種就是等上天堂,天天望著天等主耶穌再來.如果我們真的在這樣的境況當中,我們就要小心了,因為我們現在該跑的路還沒有跑完了,弟兄姊妹,我們還在生命的長跑當中了.

二,解經及應用1,    長跑的準備:放下重擔前段時間香港跑馬拉松賽,何巧蟬姊妹還特意為她學校的殘障運動員做了分享,讓人很感動.人生是什麼?人生不是跑一百米,一百公尺,槍一響,砰一聲,10秒,9秒9就跑完了.人生不是100米,人生是馬拉松.

經常跑長跑的人知道,跑步,一定要放下重擔,輕裝上陣.希臘人就很懂得這個道理,希臘人跑步比賽,什麼衣服褲子都不穿,裸體的.所以也不讓女人進去看比賽.除去一切重擔,除去一切不必要的東西,希臘人就跑向終點.

你有沒有看到過一個人跑馬拉松,穿著四寸高的高跟鞋,穿著旗袍呢?你會說,當然沒見過.但是我去外面旅行的時候,見過很多人穿著高跟鞋,穿著好漂亮的西裝去爬山.辛苦不辛苦啊?聖經裡頭說,你要跑步的時候,輕輕鬆松去跑.事實就是如此,有不少信徒,背了很多重擔,擔了好多包袱去跑人生路.這些包袱可能是你以前的創傷,可能是你的自卑,可能是你很想討人的喜悅,可能是你自己的出身--可能你出身低微,總想用什麼成就來證明自己,可能是你家人之間的關係,也可能是你跟導師的關係.背了這麼多重擔,我們心裡有很多的掙紮,很多的憤怒,很多"不能寬恕"藏在了心裡頭,你說這條人生路跑得會怎樣呢? 有沒有人會綁住沙包沙袋去跑馬拉松?沒有.但有沒有人綁住沙包沙袋跑這條人生路呢?我發現,有不少.難怪我們這麼累.香港人經常說濕氣重,人很累；這些重擔就像濕氣一樣,讓我們很累.有多少過往的創傷、失敗、失望、灰心,有多少不能寬恕的人際關係,壓住我們,讓我們不開心,讓我們不輕鬆.值不值呢? 聖經說,當放下各樣的重擔!

但聖經說這樣還不夠,還要"脫去纏累我們的罪".纏累這個字,希臘文的原文就是"野藤"的意思,哇,纏住了,纏死了.我鄉下的老樹,被藤纏住了,怎麼拉都拉不完.有的樹被纏久了,就慢慢枯死了.  那些藤是什麼?那些藤就是纏累我們的罪,它綁住我們,困住我們,叫我們不得自由,叫我們不能成長,不能踏出去.   弟兄姊妹,現在,有什麼東西纏住了我們?有什麼東西讓我們不能輕輕鬆松跑這條人生路?大家可以回家想想,你不一定現在有答案,回家了,自己安靜,可以好好想想這個問題.

2,長跑訓練:存心忍耐,奔那前面的路程卸下重擔,輕裝上陣,我們開始跑步了.經常長跑的人知道,跑長跑不是百米衝刺,你一開頭就使勁沖,肯定不行.我中學的時候,學校運動會比賽,我們班都是些書生,跑不動,只好讓我這個班長上陣.槍一響,其中一個運動員就百米衝刺,使勁沖,其他人一看不對,就跟著他使勁跑.我心想我又不拿名次,慢慢來吧,結果跑到第四圈,其他人都被那個運動員拖得累死了,一個個都下場了.但是我咬牙忍住,忍住,一個一個超上去,我每超過一個運動健將,我就說,加油,再超一個,結果我這個倒數第二名拿了全校第二名.

我班主任很高興,讓我在班上做報告,總結一下成功經驗.我的總結就一個字:忍."百忍成鋼".希臘文忍耐這個字,有兩個意思,第一個是 ``remain  here", 就是雖然上面有千斤重擔,你都保持站在當地,站穩了,不放棄；上面再重,你都要頂住.這是對環境的忍耐.還有另外一個意思,就是"keep patience",就是對人的忍耐.聖經要我們存心忍耐,不但要忍耐環境,還要忍耐人.我記得香港馬拉松論壇有個人這樣說:"今年第三次跑馬拉松啦,最難唔係體力,係耐性.跑到尾段,就覺得胸口悶,個胸好似要著火,覺得個人不是累死,是燒死.一路跑一路就想訓dik ,真係唷不到啦!"我可以用普通話翻譯一下.....

人生最怕是什麼?最怕是跑著跑著就胸悶了,就想躺下來睡覺了,就灰心了.有的時候,體力不是問題,是心力.心疲,是人生長路的最大敵人.做論文做到心疲,看護病人看護到心疲,找工作找到心疲,教育子女教育到心疲,家庭關係相處到心疲.這個時候,就很想很想放棄了.信仰也是這樣子啊,信了很多年,慢慢地心疲了；慢慢地沒有靈修禱告了,慢慢地離主越來越遠了,這種情況不僅僅別人有,有時候我自己也會有.但聖經告訴我們,我們要存心忍耐,不僅對自己的際遇,同時對別人,都要操練這份忍耐,去奔擺在前面的路程.

3,長跑目標:仰望為我們信心創始成終的耶穌為什麼要忍?很簡單,因為人生這個長跑,沒有中場休息,你也沒法中途退場.人生不得不跑那擺在前面的路程.但長跑不等於逃亡,逃亡是沒有目標的.有一次大地震,有一個人駕著車拼命開,有路就走有路就走,走走走,開了幾個小時,停下來時,發現都不曉得自己開到哪里了.他很害怕,因為他迷路了.一個人沒有了方向感,一個人不知道要去哪里的時候,就會非常驚慌.你不知道明天會睡在哪里,你不知道明天到底要做什麼,你怕不怕呢?當你沒有人生目標時,大海茫茫,你想一下,是不是很慘呢?當你每天返工放工睡覺,返工放工睡覺,每次來禮拜堂,一次兩次三次,都不知道為了什麼的時候,是不是很慘呢,是不是很容易累呢?香港人說有兩種人生觀:一種是追只狗,一種是被狗追.我想我寧願還是追只狗好點.

當我們有目標的時候,我們的忍耐就有了意義,我們的忍耐就是積極的.就好像孕婦一樣,孕婦要忍耐啊.沒有哪個孕婦說,我受不了啦,我要他六個月就出生.不可能嘛!懷孩子很辛苦啊,肚子大了,睡覺壓著腰骨痛；經常容易餓醒,因為媽媽一個人吃,兩個人消耗啊；睡覺翻身都不容易,要抱著肚子,慢慢移身啊.生孩子也很辛苦.我有個同學,生孩子生了八個鐘才生出來,她說她累到都想跟醫生說"醫生,能不能不生了啊?" 但即使是這樣痛苦,很少有母親說我不忍了,我放棄啦.至少我敢肯定,在座朋友的母親都沒有放棄.為什麼?因為有目標有盼望,這種忍耐就不是消極的"死忍".永遠有盼望,這就是懷孕母親偉大的地方.

我們基督徒的目標是什麼?是耶穌基督.希伯來書要我們"仰望為我們信心創始成終的耶穌".創始"Arkeigos"這個字,希臘文原文是領袖、元帥的意思.比方說蔡元雲醫生,蘇恩佩姊妹,他們是突破機構的開創者,就是突破的"Arkeigos".但這個字在希伯來書裡頭還有一個很特別的意思,它是一個領航的術語.當條船撞到了礁石,船就要沉的時候,唯一的生機就是有個人,背住條繩,遊到岸上,綁住繩子,然後所有其他船員就順著繩子游到安全的地方.第一個背著繩子,不怕風雨,不怕礁石,為了大家的生命遊到岸上的人,就叫"Arkeigos",創始.耶穌基督就是那個不懼風雨,不怕危險,為我們遊到對岸的人.

耶穌基督不但為我們創始,還為我們成終.成終就是完成.耶穌基督為我們輕看羞辱,忍受了十字架的苦難,不顧生命遊到對岸去,就是為我們作成一條橋樑,一條繩子,叫我們可以借著他遊到神那裡去.就像腓利比書4:13節所說:"靠著那加給我力量的,凡事都能行."當我們仰望基督,有了生命的目標,靠著他加給我們的力量,我們怎麼不能夠忍耐到底呢?4,長跑支援:見證人不過,忍耐到底,不是靠我們自己一個人死忍.自己一個人,很容易就忍到崩潰了.跑長跑的人知道,一個人跑步最無聊,也最容易放棄,要有一幫人跟你一起跑.你累了,有人鼓勵你不要放棄；你跑慢了,有人也跑慢等著你；你腳扭傷了,有人過來看護你,給你搽黃道益活絡油.總之,陪伴你,不放棄你!就像聖經說的那樣,像雲彩一樣圍繞你.圍繞這個字的意思,就是伴隨的意思.伴隨你的人不是觀眾,不是來看你表演的；他們不會說你跑得這麼差啊,怎麼這麼慢啊,你跑步姿勢不好看啊.不會!他們陪伴你,一起跑.他們就是何巧嬋姊妹說的那些跟殘障學生綁著手一起跑的人,不畏困難,不怕落後,一直跑到終點的那群人.

長跑需要這樣一群人,我們可以稱他們為弟兄姊妹,也可以稱他們朋友.朋友就是希伯來書所說的像雲彩一樣的見證人.就像亞伯拉罕,像約瑟,像摩西,像我們的主耶穌,在我們前面給我們領路.

人這一生,朋友很重要.我小時候我爸媽說,你要跟開心上進的人交朋友,不要跟脾氣不好的人交朋友.為什麼?因為如果你的朋友不是像雲彩一樣,是一樣烏雲一樣圍繞著你,那就慘了.所以西方有句話說,要知道你是什麼樣的人,看你的朋友是些什麼樣的人就好了.

有一位美國教友,在班康街教會長大,多蘿絲(Dorothy,上帝的贈禮)是他(或她)教會主日學中班的永久成員.教會中每個小朋友都知道:只要你升到班康街教會主日學中班,你就會遇到多蘿絲.當他們有些人的父母還在那個中班時,多蘿絲已經在那個班了.多蘿絲負責分發鉛筆,按名冊點名,然後在主日學結束時把鉛筆收回去.大家都把多蘿絲看成是助教.大家都知道,多蘿絲患有唐氏綜合症,智力永遠長不大.可當年在教會時,大家都當多蘿絲是一名助教.當多蘿絲五十多歲死時--一般患有唐氏綜合症的人命都不長,整個教會都為她的葬禮出動了.沒有人提到多蘿絲智力遲鈍或者飽受折磨.許多人做見證說:認識多蘿絲,他們感到多麼幸運! 我要說,我們的多蘿絲也很幸運,她有這樣一幫雲彩一樣的弟兄姊妹,不丟下她,陪她一起跑.教會是什麼?這就是教會,教會就是誰也不丟下誰,向著標竿直跑,直跑,跑下去的一群見證人.

三,結論:人生的馬拉松弟兄姊妹,你同我都一同在跑這條人生路,你是不是已經累了?你是不是心疲啦?是不是想放棄? 讓我們檢查一下,我們是不是背著好多的重擔在跑這條人生路,我們是不是有很多的罪纏住我們,叫我們放不下呢?我們回去可以細細想想.我們有沒有朋友--支援系統呢?如果沒有,我告訴你,那這條路很危險.還有,最重要的是,我們人生的目標、方向是什麼?是討神的喜悅,還是討世界跟別人的喜悅?還是我們都不知道自己跑到哪里,都迷路了?"下垂的手,發酸的腿,直起來吧! ``弟兄姊妹,仰望為我們信心創始成終的領航員耶穌,靠著神的話語,我們訓練有素,跑下去,一直跑下去,無論你現在是人生的上半場,還是下半場!



\newpage

\section{在途上}
\label{sec:99}
\textbf{鄧瑞強博士 - 「在途上」(On the Way)}
\newline
\newline
\hyperref[sec:98]{< < < PREV SERMON < < <}
~
\hyperref[sec:toc]{[返目錄]}
~
\hyperref[sec:100]{> > > NEXT SERMON > > >}


路加福音 24:13-35
\newline
\begin{longtable}{cl}
\hline
\hline
章節 & 經文 (和合本修訂版)\\
\hline
24:13 & \begin{tabularx}{0.7\textwidth}{X} 同一天，門徒中有兩個人往一個村子去；這村子名叫以馬忤斯，離耶路撒冷約有二十五里。 \end{tabularx} \\ \\ \relax
24:14 & \begin{tabularx}{0.7\textwidth}{X} 他們彼此談論所發生的這一切事。 \end{tabularx} \\ \\ \relax
24:15 & \begin{tabularx}{0.7\textwidth}{X} 正交談議論的時候，耶穌親自走近他們，和他們同行， \end{tabularx} \\ \\ \relax
24:16 & \begin{tabularx}{0.7\textwidth}{X} 可是他們的眼睛模糊了，沒認出他。 \end{tabularx} \\ \\ \relax
24:17 & \begin{tabularx}{0.7\textwidth}{X} 耶穌對他們說：「你們一邊走一邊談，彼此談論的是甚麼事呢？」他們就站住，臉上帶著愁容。 \end{tabularx} \\ \\ \relax
24:18 & \begin{tabularx}{0.7\textwidth}{X} 兩人中有一個名叫革流巴的回答：「你是在耶路撒冷的旅客中，惟一還不知道這幾天在那裡發生了甚麼事的人嗎？」 \end{tabularx} \\ \\ \relax
24:19 & \begin{tabularx}{0.7\textwidth}{X} 耶穌對他們說：「甚麼事呢？」他們對他說：「就是拿撒勒人耶穌的事。他是個先知，在神和眾百姓面前，說話行事都大有能力。 \end{tabularx} \\ \\ \relax
24:20 & \begin{tabularx}{0.7\textwidth}{X} 祭司長們和我們的官長竟把他解去，定了死罪，釘在十字架上。 \end{tabularx} \\ \\ \relax
24:21 & \begin{tabularx}{0.7\textwidth}{X} 但我們素來所盼望要救贖以色列民的就是他。不但如此，這些事發生到現在已經三天了。 \end{tabularx} \\ \\ \relax
24:22 & \begin{tabularx}{0.7\textwidth}{X} 還有，我們中間的幾個婦女使我們驚奇：她們清早去了墳墓， \end{tabularx} \\ \\ \relax
24:23 & \begin{tabularx}{0.7\textwidth}{X} 不見他的身體，就回來告訴我們，說她們看見了天使顯現，說他活了。 \end{tabularx} \\ \\ \relax
24:24 & \begin{tabularx}{0.7\textwidth}{X} 又有我們的幾個人往墳墓那裡去，所發現的正如婦女們所說的，只是沒有看見他。」 \end{tabularx} \\ \\ \relax
24:25 & \begin{tabularx}{0.7\textwidth}{X} 耶穌對他們說：「無知的人哪，先知所說的一切話，你們的心信得太遲鈍了。 \end{tabularx} \\ \\ \relax
24:26 & \begin{tabularx}{0.7\textwidth}{X} 基督不是必須受這些苦難，然後進入他的榮耀嗎？」 \end{tabularx} \\ \\ \relax
24:27 & \begin{tabularx}{0.7\textwidth}{X} 於是，他從摩西和眾先知起，凡經上所指著自己的話都給他們作了解釋。 \end{tabularx} \\ \\ \relax
24:28 & \begin{tabularx}{0.7\textwidth}{X} 他們走近所要去的村子，耶穌好像還要往前走， \end{tabularx} \\ \\ \relax
24:29 & \begin{tabularx}{0.7\textwidth}{X} 他們卻強留他說：「時候晚了，天快黑了，請你同我們住下吧。」耶穌就進去，要同他們住下。 \end{tabularx} \\ \\ \relax
24:30 & \begin{tabularx}{0.7\textwidth}{X} 坐下來和他們用餐的時候，耶穌拿起餅來，祝福了，擘開，遞給他們。 \end{tabularx} \\ \\ \relax
24:31 & \begin{tabularx}{0.7\textwidth}{X} 他們的眼睛開了，這才認出他來。耶穌卻從他們眼前消失了。 \end{tabularx} \\ \\ \relax
24:32 & \begin{tabularx}{0.7\textwidth}{X} 他們彼此說：「在路上他和我們說話，給我們講解聖經的時候，我們的心在我們裡面豈不是火熱的嗎？」 \end{tabularx} \\ \\ \relax
24:33 & \begin{tabularx}{0.7\textwidth}{X} 於是他們立刻起身，回耶路撒冷去，看見十一個使徒和與他們正在一起的人聚集在一處， \end{tabularx} \\ \\ \relax
24:34 & \begin{tabularx}{0.7\textwidth}{X} 說：「主果然復活了，已經顯現給西門看了。」 \end{tabularx} \\ \\ \relax
24:35 & \begin{tabularx}{0.7\textwidth}{X} 於是，兩個人把路上所遇到，和耶穌擘餅的時候怎麼被他們認出來的事，都述說了一遍。 \end{tabularx} \\ \\
[1ex]
\hline
\hline
\end{longtable}
日期:2011年5月19日

主題:「在途上」(On the Way)

經文:路加福音 24:13-35

證道:鄧瑞強博士

各位弟兄姊妹,主內平安.

郁達夫寫過一篇文章,叫《一個人在途上》.講到他的兒子「龍兒」剛剛過身.講到他和妻子的心情.他這樣寫:

上床就寢,把電燈一滅,兩人只有緊抱的痛哭,痛哭,痛哭,只是痛哭,氣也換不過來,更那裡有說一句話的餘裕?受苦的時間,的確脫煞過去的太悠徐,今年的夏季,只是悲嘆的連續.晚上上床,兩口兒,那敢提一句話?可憐這兩個迷散的心靈,在電燈滅黑的黝暗裡,所摸走的荒路,每會湊集在一條線上,這路的交叉點裡,只有一塊小小的墓碑,墓碑上只有「龍兒之墓」的四個紅字. 在這文章裡,郁達夫講到他和兒子過去生活的樂事,講到他無力滿足兒子簡單的渴望,講到妻子對兒子的思慕.他這樣寫:

有一次和女人在那裡睡午覺,她驟然從床上坐了起來,鞋也不拖,光著襪子,跑上了上房起坐室裡,並且更掀廉跑上外面院子裡去.我也莫名其妙跟著她跑到外面的時候,只見她在那裡四面找尋什麼.找尋不著,呆立了一會,她忽然放聲哭了起來,並且抱住了我急急的追問說:「你聽不聽見?你聽不聽見?」哭完之後,她才告訴我說,在半醒半睡的中間,她聽見「娘!娘!」的叫了幾聲,的確是龍的聲音,他很堅硬的說:「的確是龍回來了.」這篇文章的結尾是這樣的:

「現在去北京遠了,去龍兒更遠了,自家只一個人,只是孤零丁的一個人.在這裡繼續此生中大約是完不了的飄泊.」

郁達夫經歷人生的悲苦,飽嘗愛兒離世的淒慘,體會人的無能為力,明白人間情感的無法把持.在生命的路上,只感到無盡的孤寂,腳下虛若浮萍,前景渺茫如霧.

這種感受,或許,也是人生的常態.

今日講道的經文是:路加福音24: 13-35.

這段經文講到兩個人在途上,他們心中苦悶,有無盡的失落.

24:13　正當那日,門徒中有兩個人往一個村子去；這村子名叫以馬忤斯,離耶路撒冷約有二十五里.

24:14　他們彼此談論所遇見的這一切事.

24:15　正談論相問的時候,耶穌親自就近他們,和他們同行；

24:16　只是他們的眼睛迷糊了,不認識他.

24:17　耶穌對他們說:「你們走路彼此談論的是什麼事呢?」他們就站住,臉上帶著愁容.

這兩個人在途上,他們心中的愁苦流露在面上.

他們對遭遇的事情,有很多的不明不白.

經文說,他們「談論相問」,他們不單「談論」,也「相問」.「相問」就是他們試圖探究事情內中的意義.我們遇到悲苦的事,會試圖明白事情是否能構成另一層次的意義.若不能找到另一層次的意義,則我們會苦上加苦.如:「塞翁失馬,焉知非福」,這是說,「塞翁」失去一匹馬,這是不好的,但從另一角度看,這可能是好事呢?焉知這不是一種我們現在未能明白的福氣?

若能找到事情的意義,我們較能接受事情的發生.但是,有些事情,的確是我們無法參透的.有些事情,違反了我們認定的常理.有些事情,根本不合理,其意義在道理之外.

南非的曼德拉(Nelson Mandela)坐牢坐了很久,受到無理的對待.放監之後,人們問他,是否對監禁他的人懷恨於心.他只說,唯有寬恕,才能釋放自己,也釋放他人.

曼德拉受的苦,沒有道理可言.按道理而論,正義的要求是「以牙還牙,以眼還眼」.而他卻選擇寬恕.「寬恕」是沒有道理可言的.他放下正義的「道理」的要求.他在道理之外,他轉化自己的生命,他活出道理以外的意義.

在這往以馬忤斯的路上,這兩個門徒想不通,猜不透.

然而,縱使他們想不通,猜不透,耶穌卻成為他們路上的同行者,對話伙伴.他們認不出耶穌,耶穌卻「親自就近他們」.他們未注意到耶穌,耶穌卻注意到他們.他們用自己的方式參不透他們的遭遇,耶穌卻主動讓他們明白.

耶穌問:「你們走路彼此談論的是什麼事呢?」

24:18　二人中有一個名叫革流巴的回答說:「你在耶路撒冷作客,還不知道這幾天在那裡所出的事嗎?」24:19　耶穌說:「什麼事呢?」他們說:「就是拿撒勒人耶穌的事.他是個先知,在神和眾百姓面前,說話行事都有大能.24:20　祭司長和我們的官府竟把他解去,定了死罪,釘在十字架上.24:21　但我們素來所盼望,要贖以色列民的就是他!不但如此,而且這事成就,現在已經三天了.24:22　再者,我們中間有幾個婦女使我們驚奇；她們清早到了墳墓那裡,24:23　不見他的身體,就回來告訴我們,說看見了天使顯現,說他活了.24:24　又有我們的幾個人往墳墓那裡去,所遇見的正如婦女們所說的,只是沒有看見他.」

就著這長長的回答,我想講兩點:

一,這兩人曾經將人生的盼望寄托在耶穌身上.他們說,「我們素來所盼望,要贖以色列民的就是他!」曾經有希望,如今失去了這希望,這比從沒有過希望,會感到更大的失望.「失戀」的痛苦比從未戀愛過,感到更大的痛苦.我們聽過失戀的人自殺,但很少聽一個人因為未戀愛過而自殺.

我們的希望,可以是多方面的,如:政治,經濟,社會,民生,愛情,事業,人類的救贖,天國的實現.若我們曾有過希望,當希望幻滅時,我們會感到無比的失落.

這兩個人的希望,曾落在耶穌身上,但也是耶穌的際遇,令他們失望.

二,這兩人對耶穌的認識,和我們的「信經」對耶穌的描述相差不遠.他們知道耶穌是「先知」,是神的使者.知道耶穌被釘在十架上,受死,埋葬.知道耶穌第三天復活了.他們甚至認識那些見過復活主顯現的婦女.他們對這一切,瞭如指掌.但是,這一切知識,對他們的人生,完全沒有作用.

一切的認信,原來可以對一個人毫無意義的.對耶穌的這一切認識,原來可以是完全對生命毫無承托力的.他們看來已知曉耶穌的一切,但看來也對耶穌一無所知.

在人生的路上,在信仰的路上,或許,我們就是這兩個門徒.

24:25    耶穌對他們說:「無知的人哪,先知所說的一切話,你們的心信得太遲鈍了.24:26    基督這樣受害,又進入他的榮耀,豈不是應當的嗎?」24:27    於是從摩西和眾先知起,凡經上所指著自己的話都給他們講解明白了.

主耶穌的說話顯示,我們對某些事情,總是很難接受的.這兩個門徒不能明白,耶穌在死亡中進入榮耀,是應當的.我們不能接受耶穌死去,我們希望耶穌以君臨天下的形式剷除世上的邪惡.我們不能接受神的無言無語,我們不能接受神不顯神蹟,我們不能接受神採取弱者的行徑.我們不能對這位神寄予厚望.我們有種種渴望,但耶穌為何不滿足這些渴望,而要死在十架上?

耶穌說,這一切是應當的.「於是從摩西和眾先知起,凡經上所指著自己的話都給他們講解明白了.」十字架,是神行事的方式.

這兩個人明白不明白呢?我們明白不明白呢?

或者,他們需要一些機遇,讓他們不單在思想上明白耶穌的話,也能在生命上進入耶穌的生命中.

24:28　將近他們所去的村子,耶穌好像還要往前行.

為什麼耶穌好像還要往前行呢?

或許,耶穌是製造一個機遇,讓他們的生命進入耶穌的生命中.

天色將晚,一個陌生人,若走在陌生的路上,是危險的.

這兩個門徒,若不關心這個陌生人,則會失去耶穌了.

24:29　他們卻強留他,說:「時候晚了,日頭已經平西了,請你同我們住下吧!」耶穌就進去,要同他們住下.

明顯,這兩個門徒關心這個陌生人.他們「強留」這個陌生人.他們為了保護這個陌生人,用盡辦法留住這個陌生人.

這是一個美麗的接待陌生人的故事.

在路加福音裡,有一個獨有的接待陌生人的故事.在路加福音10: 25-37裡,記載了一個我們稱之為「好撒馬利亞人」的故事.在途上,有人被強盜打得半死.在這途上,有祭司經過,有利未人經過,他們熟讀聖經,但毫無反應.有一撒馬利亞人,走在這途上,他接待了這陌生人,救助了他.耶穌說,這撒馬利亞人做的,能承受永生.

這個撒馬利亞人,額外付出了金錢和勞力,打擾了生活的節奏,只為了幫助一個陌生人.在途上,他背起了一個十字架,他認為:這是應當的.

如今,耶穌在途上,天色已晚,若無人示意接待他,他便要繼續往前行了.

這個陌生人會肚餓,露宿荒野時會凍,會遭遇盜賊的危險.

而這兩個門徒,他們看到這個陌生者的需要.世上有無數的陌生者,需要無數的門徒去接待.

主耶穌講過:「

凡為我名接待這小孩子的,就是接待我；凡接待我的,就是接待那差我來的.」(路9: 48)

這兩個門徒的確在接待陌生人時接待了耶穌.

24:30　到了坐席的時候,耶穌拿起餅來,祝謝了,擘開,遞給他們.24:31　他們的眼睛明亮了,這才認出他來.忽然耶穌不見了.

到坐席時,被接待的陌生人,看來不再是客人,而成為分享的主人.在一個懂得接待有需要者的人的生命裡,耶穌是他生命的主人.耶穌做的,是經典的聖餐的場面.這個餅被擘開,是表明生命的破碎.破碎自己,是為了他人的需要.活著,是為了他人.

「強留」一個陌生人,刻意幫助一個有需要的人,這本身已是一種聖餐的行動,而這行動,正就是耶穌的十字架行動.如今,這兩個門徒不單是知道耶穌上十架的事蹟,他們是進入了這十字架的行動中.他們明白了,他們認出了耶穌.忽然耶穌不見了,但他們明白耶穌活在他們的生命中.他們活著,就是基督在他們生命中活著.

在人生的路上,我們仍將有無數的困惑,無數不可解的事情,無數不想面對的痛苦,無數不能滿足的渴望.這一切將仍是不可解的.只有將眼睛望向他人,望向街上露宿的,望向無倚無靠的,並愛他們,將餅分給他們.或許,只有當我們放下自己,走向他人時,我們才能放下愁苦.或許,在十架下,我們會明白生命的意義.

24:33　他們就立時起身,回耶路撒冷去.

「耶路撒冷」代表神的使命的中心.這兩個門徒,重新回到神的使命中去了.

我們的「耶路撒冷」在哪裡呢?

但願榮耀歸於聖父,聖子,聖靈.阿們.



\newpage

\section{虛假與代替}
\label{sec:100}
\textbf{伍渭文牧師 - 虛假與代替}
\newline
\newline
\hyperref[sec:99]{< < < PREV SERMON < < <}
~
\hyperref[sec:toc]{[返目錄]}
~
\hyperref[sec:101]{> > > NEXT SERMON > > >}


日期:2011年5月22日

主題:虛假與代替

經文:詩篇16

證道:伍渭文牧師

網上滿流行一個黑色笑話,諷刺內地造假的風氣.有一位農夫吃老鼠藥輕生,因為老鼠藥是假的,死不了；既然老天不想我死就好好活下去吧,買一瓶酒回來慶祝新開始吧,豈料酒有甲醇,便一命嗚呼了.最近一項造假是把豬肉浸在磞砂和加上色素,當作牛肉來買.造假風氣之盛行,禁之不絕,有人戲謔:甚麼都是假的,祇有騙子是真的.

虛假是冒認真實,但代替是換掉真實.前者仍然覺得真實的重要,因為重要才去冒認它；但代替則認為真實並不是那些重要了,可以有其他的取代它了.母親自己乳養嬰兒最好,因為可以不單把養料給孩子,而且還可以把抗體傳給嬰兒.但奶粉商為了商業利益,宣傳奶粉可以代替母乳.事實上很多第三世界貧窮地區的婦女,因為這些宣傳而放棄了乳養嬰兒,失了機會把自己的抗體傳給嬰兒,還因為調較奶粉的水不夠潔淨,引致疾病.世界衛生組織(WHO) 在1981年通過國際母乳代用品銷售守則(International Code of Marketing of Breast milk Substitutes),聲明奶粉不能取代母乳, 1988年危地馬拉依守則立例,兒童夭折率明顯下降.WHO更在2001年發表宣言,鼓勵出生六個月應以純母乳餵嬰兒,之後才添加副食品,直到兩週歲.

然而財力雄厚的跨國奶粉商,不單宣傳奶粉可以代替母乳,並謂奶粉加入了神奇的添加劑PHD令嬰孩聰明.PHD是phospholipid,形聖所有細胞膜的重要元素,包括腦細胞,父母就想當然吃了會令孩子聰明,因為PHD使人聯想到博士.所以代替也有虛假成分,因它聲稱可以代替不能代替的.

今天詩篇十六篇也提到代替,以別神代替耶和華.「

以別神代替耶和華的,他們愁苦必加增.」(v.4) 其實代替在這裡不是否定相信上帝,乃相信上帝同時又相信偶像.以色列人的試探不是放棄上帝,乃在上帝之外,加上對偶像的敬拜,混合的信仰.經文的結構,第一節可和二,三節平行的,第一節以第一身說:神啊,求你保佑我,因為我投靠你.第二節「

我的心哪

」下有虛線表示原文沒有的,可以繙作以第二身說:「

你曾對耶和華說:你是我的神,我的好處不在你以外.論到世上的偶像(聖民holy ones) :你又美又善,是我最喜悅的.」

[著名版本耶路撒泠聖經(Jerusalem Bible) 就把這對比繙出來:Look after me, God, I take shelter in you. (v.1)To Yahweh you say, ``My Lord you are my fortune, nothing else but you."(v.2)Yet to those pagan deities in the land, ``My princes, all my pleasure is in you."(v.3)]

一 真實信仰是滲透生活每一個層面

vv. 5-8「

耶和華是我的產業,是我杯中的份；我所得的,你為我持守,用繩量給我的地界,坐落在佳美之處；我的產業實在美好.我必稱頌那指教我的耶和華；我的心腸夜間也警戒我.我將耶和華常擺在我面前,因他在我右邊,我便不至動搖.」(vv.5-8)

上帝信仰是排他的,無添加的.十誡首誡「

除了我以外,你不可有別的神.」(出二十3) 他們以為添加了偶像的敬拜,可以左右逢源,但結果是愁苦加增.我們看看偶像敬拜和上帝信仰有何不同.偶像需要獻祭和頌禱等宗教儀禮作為供奉:「

他們所澆的血我不獻上,我嘴唇也不提別神的名號.」(4下)但信仰上帝是遵守他的誡命.「

我必稱頌那指教我的耶和華；我的心腸夜間也警戒我.」(v.7)

奧圖(Rudolf Otto) 在《神聖的概念, 1917；English translation, 1923》(

The Holy -- On the Irrational in the Idea of the Divine and its Relation to the Rational

) 一書強調宗教的神聖感覺；祭祀的人透過禮儀,進入和凡俗割裂的陌生領域,產生神聖感,感到敬畏,渺小；這種神聖感是所有宗教所共有的,這種神聖的感覺就是除去悟性和德性的上帝；但除去悟性和德性就不是聖經中的上帝了.

上帝以外加上偶像的敬拜,本來是希望信仰增值,多一種助力,但慢慢添加上去的偶像崇拜蠶食對上帝的信仰,我們也以為上帝單單喜悅我們萬萬的油河,千千的公羊,忘記上帝要的是行公義,好憐憫,存謙卑的心與上帝同行.(彌六6-8)

真實的信仰是滲透到生活每一個層面,不單嘴唇頌讚,宗教儀禮.在上一篇詩篇(十五篇)說得非常清楚:「

耶和華啊,誰能寄居你的帳幕?誰能住在你的聖山?就是行為正直,作事公義,心裡說實話的人,他不以舌頭讒謗鄰里.他眼中藐視匪類,卻尊重敬畏耶和華的人,他發了誓,雖然自己吃虧,也不更改.他不放債取利,不受賄賂以害無辜.行這事的人必不動搖.」這是職場與人相處之道.

在以馬忤斯路上,甚麼時候兩位門徒認出耶穌呢?是他們願意接待一位陌生人到家中,並分享食物時,才認出這是復活的基督.「

時候晚了,日頭已經平西了:請你同我們住下吧!耶穌就進去,要和他們住下.到了坐席的時候,耶穌拿起餅來,祝謝了,擘開,遞給他們.他們才認出他來.」(路二十四29-31)

詩人得到上帝祝福時,沒有忘記上帝:「

耶和華是我的產業,是我杯中的份；我所得的,你為我持守,用繩量給我的地界,坐落在佳美之處；我的產業實在美好.」(v.5)信仰滲透到他的經濟活動,參與他的財務管理,他不被所擁有的所吸引,因為上帝才是他的產業,他的祝福.

二 真實信仰能直視生命的苦痛和死亡 vv.9, 10「

我的肉身也要安然居住.因為你必不將我的靈魂撇在陰間,也不見你的聖者見朽壞.」(vv. 9下, 10)

詩人描述上帝如何祝福他,讓他昌盛,但他要面對身體的危疾,甚至曾經來到死亡的邊緣.「

我的肉身也要安然居住.因為你必不將我的靈魂撇在陰間,也不見你的聖者(faithful ones信靠你的人)見朽壞.」(vv. 9下, 10) 這是重病得醫治的感恩讚美,也說明我們作為人,要面對隨生命而來的諸般苦難和死亡.但人對死亡有一種不安和顧忌,不肯真實地面對死亡,避諱死亡的名字,稱遺囑為平安紙,買棺木的地方為長生店,反對在鄰近地區安置骨灰龕,也是忌諱死亡,多於衛生或噪音問題.

教會復活節後的主日選讀這詩篇,說明復活不是取消死亡,乃幫助我們直視死亡,超越死亡,看到死亡是通到更榮耀生命的路徑.「

基督特要藉著死,敗壞那掌死權的,就是魔鬼,並要釋放一生因怕死而為奴僕的人.」(來二14,15)

每個人都要經過生,老,病,死,祇有電影中的吸血疆屍,才能永遠年輕.上星期和一位弟兄談話時,他題到以前很少運動,但兩年前聽到一位年長牧師一席話,開始每星期有幾天運動了.這位牧師說:「我要好好運動,以致日後好好照顧太太,或不用太太照顧自己,因為我們年紀老了.」我跟他說,這是愛的表現.他說,也是責任.之後,有一晚睡在床上,師母熟睡了,我左手握著她的右手,兩手都是暖暖的,我想如果時間現在慢慢流逝,誰的手將會首先變得冷冰呢?想到這裡,我緊緊握著她的手,她仍然熟睡.

因著基督己從死裡復活,我們就無懼死亡,故我們不需避諱或顧忌死亡,不用代替,也不用虛假.盧雲(Henri J. M. Nouwen)說:「我們可否願意不逃避痛苦,無能為力時不故作忙碌,站在哀傷的人旁邊一起面對死亡.」我們可以如此,因我們相信:「

我的肉身也要安然居住.因為你必不將我的靈魂撇在陰間,也不見你的聖者(faithful ones信靠你的人)見朽壞.」(

Are we willing to not run away from the pain, to not get busy when there is nothing to do and instead stand rather in the face of death together with those who grieve?

)

三 真實的信仰帶來真實的滿足 v.11「

你必將生命的道路指示我；在你面前有滿足的喜樂,在你右手有永遠的福樂.」

推崇代替品的人,可能是出於欺騙,但接受代替品的人,可能出於真心.以別神代替耶和華的人,有時真的以為崇拜的偶像,就是可以信賴的真神.有人付出高價買了以為是真的假藥,有人把終身託付給一位後來才發現是薄幸的郎君.當人買了虛假的東西,損失的是錢財,吃了虛假的食物,傷害的是健康,但當人跟隨一個人,後來發現他是虛假,背棄了起初的誠信,他就有被出賣的感覺.我在網上看到一篇提及前菲律賓靈總統馬可斯的文章:

馬可斯(Marcos)於 1965 年當選菲律賓總統,執政初期,確實也是致力於農業,工業,教育發展,民眾支持率高達七成.之後馬可斯變本加厲地將軍政大權一手獨攬,任用親信,大肆收刮財產,揮霍無度,政治貪污腐敗,擔任"總統"近 20 年期間,家族累積的財產近一百億美元!馬可斯夫人伊美黛(Imelda),超愛名牌鞋子,單單收藏的鞋子就有三千雙!

文章題目是〈

陳水扁與馬可斯,驚人的巧合!下場相同?

〉,文章內容不用說了.1980年白樺的電影劇本《苦戀》,描述主角一家在文化大革命中遭受迫害,但作為主角的父親反對女兒出國,女兒反問父親,"你愛國家,國家愛你嗎?"

不少人追求理想,跟隨擁戴他們所崇拜的偶像,認為走對了路,最後發現路是虛假的,後悔太遲了.但詩人在這裡說,信靠上主是生命的路:「

你必將生命的道路指示我；在你面前有滿足的喜樂,在你右手有永遠的福樂.」 (v.11)

使徒行傳記載耶穌復活後,使徒們走上街頭,公開傳講福音,不怕下在獄中,甚至賠上生命,因他們相信,上帝必不將他們的靈魂撇在陰間.保羅在林前十五章談到基督的復活,作了一結論:「

所以,我親愛的弟兄們,你們務要堅固,不可搖動,常常竭力多作主工；因為知道,你們這勞苦在主裡面不是徒然的.」(林前十五58)

約翰的門徒士每那(Smyrna)主教八十六歲時被捉押到刑場被火燒,主審官試圖游說坡旅甲否認基督,說:我們尊重你年事已高了,祇要你以凱撒的名字發誓,悔改,我會釋放你,拒絕基督,就立刻得到自由.坡旅甲說:我服事他八十六年,他從沒有傷害我,我怎能褻瀆我的王和救主呢?

結論:上帝說:「

除了我以外,你不可有別的神.」不錯,這是絕對主權的宣告,同時也是慈愛的召喚,免得當我們以虛假為真實,以別神代替耶和華時,生命充滿愁苦.當我們說:「我的心哪,你曾對耶和華說,你是我的主,我的好處不在你以外.」我們就不能對其他的偶像說:「你又美又善,是我最喜悅的.」這是代替,而這代替是虛假的,不真實的.

我們祇能說:「我的心哪,你曾對耶和華說,你是我的主,我的好處不在你以外.」阿們.



\newpage

\section{黑門的甘露}
\label{sec:101}
\textbf{王教佺先生 - 黑門的甘露}
\newline
\newline
\hyperref[sec:100]{< < < PREV SERMON < < <}
~
\hyperref[sec:toc]{[返目錄]}
~
\hyperref[sec:102]{> > > NEXT SERMON > > >}


日期:2011年6月5日

題目:黑門的甘露

証道:王教佺先生

經文:詩133:1-3

看哪,弟兄和睦同居,是何等地善,何等地美!

這好比那貴重的油澆在亞倫的頭上,流到鬍鬚,又流到他的衣襟.

又好比黑門的甘露降在錫安山,因為在那裡有耶和華所命定的福,就是永遠的生命.

引言:

今天的分享,是關於見證的.

我們基督徒的生活就是見證的生活,基督徒的一生也是作美好見證的一生.那麼,這見證的含義如何呢?聖經上這方面的教導很多,今天我們要分享的主題經文是詩篇133:1-3.

理解這段經文,不需要所謂高深的神學.經文以一種明白樸素的話語,告訴我們一個明白樸素的道理:基督徒要作的見證,就是在上帝的愛中,彼此相愛.

詩中提到一地方--黑門山.雖然這座山海拔並不算太高(最高峰2,814米),然而它對以色列民族來說,有著特別的意義.

我們知道,黑門山在以色列的北部,是約旦河的重要源頭.約旦河水就是從這裡一路南下,流向加利利海,流向以色列國,迦 南的腹地,流向死海,流向紅海......在漫長的以色列歷史中,他們體會到,上帝的愛就這樣如同黑門山的水,黑門山興起的雲,降下的風霜雨雪,滋潤著他們.(備 用:對以色列駐華農業參贊薩瑞格博士的採訪)

那麼,我們又當如何回饋上帝的愛呢?詩中又說,弟兄和睦同居是何等的善,何等的美!並且,「好比那貴重的油澆在亞倫的 頭上,流到鬍鬚,又流到他的衣襟.」就是說,這彼此相愛,不單單是我們的回應,乃是上帝對我們的使用,就如同用貴重的油膏抹亞倫一樣膏抹我們；這彼此相 愛,不只是付出,而是一種莫大的祝福.

那麼,這所謂的彼此相愛,意味著什麼呢?

一,見證,是鄰舍般的,用全部生活去付出.

路得記上記著一個人,一個不會講道,也沒有多少知識的婦人,她名叫拿俄米.當拿俄米正當年華的時候,她的國家遇到一場 空前浩劫,尼布加尼撒王的大軍讓以色列人從此成了亡國奴.拿俄米這樣一個平凡的婦人,被迫帶著兩個兒子被擄到別人的國家,在那裡勉強住下來.其生活的景況 可想而知.

後來,拿俄米總算為兩個兒找了媳婦,日子才有個開頭,可是這兩個兒子卻一個接一個相繼去世.就當拿俄米決定回到自己的國家,過一個人的孤苦生活的時候,一個讓她意想不到的事情發生了:

她的兒媳婦路德要跟她走,路德說:「

你往哪裡去,我也往那裡去；你在哪裡住宿,我也在那裡住宿；你的國就是我的國,你的　神就是我的　神.」(德1:16)

拿俄米一家客居外邦,雖然地位低下,無錢無勢,但外邦的兒媳婦路得,卻從她們身上看出背後的真神.基督徒的見證,就是這樣,用她的一生,用她生命中每一個細節「作」出來的.

在我老家的農村教會裡,就有一位姊妹,她的丈夫一度不許她禮拜,聚會,在家裡也處處為難她.但是幾年後,她的丈夫居然也信主了.怎麼回事呢?

自從她信了主,她就決定,在家裡要把一切的家務做好,但是禮拜一定要參加.可是既便是這樣,丈夫還是逼迫她.記得有一 天,禮拜結束,她很高興來找長老說,每個禮天早晨,她都是把家裡打裡好,豬餵好,飯做好,還把丈夫的飯乘好,然後再偷偷溜出門,可是今天,她正要出門,卻 被丈夫正好碰上.她跟長老說:感謝主啊!要是平時,丈夫一定要打我,可是今天,丈夫沒打我,今天丈夫說了一句話:「去吧!」

我們信主的人往往有一個習慣性的認識,就是認為,信主就要旅行耶穌基督的大使命,大使命就是要傳福音,傳福音就是派單 張,請人聽布道會.結果讓人感覺基督徒好像是一群誇誇其談的人,以為自己擁有真理.然而許多的見證告訴我們,福音是靠行出來的.是一種鄰舍般的,用全部生 活所參與的付出.

我們還可以看到,基督徒的見證,不僅是要靠行出來,而且,它是一種滋潤,一種感化,它是讓人自己情願.

二,見證,是「讓他自己情願」,而不是強加與人.

上帝很愛中國,這中間有一個很大的標誌,就是中國是一個以無神論為主流意識的國家.這種環境其實也是一種恩典.上帝是要讓我們這些中華兒女,有機會在兩個截然不同的心靈世界裡,體會,品味他的崇高,他的能力,他的之真之善之美.

上上個主日,我們內地生的團契一起在下面教工餐廳吃午飯,其中有一位新來的弟兄,他也信主不常,他的父親,母親也在 坐,當時,這位同學的父親告訴我,他很高興自己的兒子入教(雖然他還不信),他說,孩子受了近二十年的無神論教育,可是來到香港,才幾個月就信上帝了.這 是內地教育的失敗呢,還是因著一種不可抗拒的大能呢?.

在北京的海澱堂,我記得,常常會看到一個剃著光頭男孩坐在那裡,讚美,禱告,聽道很認真,全然不像他那幅看上去玩皮的樣子.他告訴我,沒有人勸他信主,是他自己來到這裡以後,不想走了.他說學校裡聽到的都是假的,這裡是真的.

我們團契的另一位弟兄也有同樣的經歷,他告訴我,當時他是一個人在大街上閉走,看到一座教堂,就毫無目的的走進去.當他坐下來,不知道怎麼了,眼淚一下子就流下來.

雅歌2:7:耶路撒冷的眾女子啊,我指著羚羊或田野的母鹿囑咐你們,不要驚動,不要叫醒我所親愛的,等他自己情願.

--上帝揀選我們,不是用金錢誘惑,不是用政治的高壓,他是要我們自己情願.--我們在內地經常看到一個商品的廣告詞:「只要是美的,人人都喜歡」

那麼,我們的神是不是那種總是討好人的神呢?我們的見證,是不是就是要投人所好呢?

三,見證,是堅持,守望,而不是中庸,妥協.

上帝的愛給人以自由的意志,但不意味著他是處處總是隨人的心意.

當所多瑪和俄摩拉這兩個地方惡貫滿盈的時候,既便亞伯拉罕怎樣向上帝求告饒恕他們,但是滅亡之火還是臨到了他們.

最近幾年來,我們越來越多地看到中國內地不斷發生一些令人痛心的事情,「毒奶粉」,「黑芝麻」,「計生幹部賣小孩」,現在又有「全國70\%的轉基因食用油」等等,愈演愈烈.科技創新的消息越來越少,可是邪惡的勾當卻是層出不窮.許多的事情就發生在我們的身邊.

作為一個普通的基督徒,我們又能作什麼呢?我們當然只能祈求上帝的寬恕?但是,我們這樣禱告,並不意味著在實際生活中無所作為,更不是妥協,中庸,同流合污.

我所在的海澱堂就在被稱為中國硅谷的中關村,教會裡許多弟兄姊妹在那裡上班,就有一些人是作老闆的.他們告訴我,他們 最大的困惑就是目前商界的「潛規則」,許多的大買賣都是政府採購,這種業務很容易碰到所謂「極個別」的腐敗分子,本來同行之間競爭就很激烈,小企業已很難 生存,如果你要是拒絕使用灰色手段的話,可以說,根本無法做生意,無法在那裡生存.所以就出現兩種情況:一種是遠離.就有一對夫婦,十年前賺了第一筒金之 後,不幹了!於是,在中關村買了四套房子(當然,他們沒想到今天房價漲得這麼高),目的就是一個--遠離世俗.不看電視,不上網,專心靈修禱告......這是少 數例子.

而更多的是第二種情況,那就是繼續奮鬥,打拼,但無比的痛苦,掙扎,可想而知.因為他們既要守上帝的律法,又要面對商業場上不許多不規範的競爭.

這條道路能不能走的通呢?現在有很多好的情況開始出現.就在我們崇基內地生的團契,有一位律師,他工作在深圳.他告訴 我,他們的律師事務所有好幾位基督徒,他們決定,不用灰色的不正當的手段打官司,不打灰色的官司.但是很感恩,他說,他們最近的工作很順利,並且,很多人 慕名請他們辦案.

這條道路的前景如何,我們不知道,但是有一點應該很清楚,它一定是一條十字架的道路.這條道路,沒有掌聲,沒有歡呼,需要的是忍耐和等候.我相信,對大部分中國基督徒來說,這條道路對他們來說最為現實.

其實,看看中國的歷史,中國不僅僅需要有良心有使命的政治家,最最需要的是一群有責任,使命,有尊嚴的公民.而中國基督徒的使命,就應該去努力的造就.雖然這條道路看上去有平淡,但我們生命的永恆就在其中.如賽42:1所說:

「

看哪,我的僕人,我所扶持,所揀選,心裡所喜悅的,我已將我的靈賜給他,他必將公理傳給外邦.他不喧嚷,不揚聲,也不使街上聽見他的聲音.壓傷的蘆葦,他不折斷；將殘的燈火,他不吹滅.他憑真實將公理傳開.」(如賽42:1-3)

作為一個公民,作為一個基督徒,我們無須擔憂我們在歷史中為什麼這樣有限,這樣的無足輕重,我們不要以為我們作的是一件懦弱的事.因為,當時候到了的時候,上帝會自己站出來的.以賽亞書這裡接著說:

「

耶和華必像勇士出去,必像戰士激動熱心,要喊叫,大聲吶喊,要用大力攻擊仇敵. 「我許久閉口不言,靜默不語,現在我要喊叫像產難的婦人,我要急氣而喘哮.」(賽42:13-14)

結語:

在新約聖經使徒行傳一章那裡有這樣的記述:當時,耶穌已經復活了.於是,他的門徒們一下子來了精神,就在聚集的時候,問耶穌說:「

主啊,你復興以色列國就在這時候嗎?」 (徒1:6)主耶穌看出他們的心思,知道他們想搞政治.這個心願雖然是美好的,但主耶穌卻告訴他們說:「

父憑著自己的權柄所定的時候,日期,不是你們可以知道的.但聖靈降臨在你們身上,你們就必得著能力；並要在耶路撒冷,猶太全地和撒瑪利亞,直到地極,作我的見證.」(徒1:7-8)

主耶穌是在說:「我是你們良心的主,生命的主,不是一個小小的政治家」,主耶穌也在告訴我們:在這個悖謬的世代,當人 與人的關係變得遙遠和淡漠的時候,我們基督徒使命,就是要站出來,走出去,作愛的見證.叫人看到基督,讓人心得滋潤.如同黑門的干露,滋潤出一片流奶與蜜 之地.

(備用:我想起小時,老一輩的基督徒......文革的時候,大家悄悄地在一起聚會.姥姥從棺材裡拿出她唯一一件新衣服說:這件衣服我現在不穿,等到教會復興的時候才穿.現在想來,她一個不識字的老太太,怎麼會有這樣大的信心呢?)

相信中國人最能理解「黑門甘露」的念義:春夜喜雨(杜甫) 好雨知時節, 當春乃發生.隨風潛入夜, 潤物細無聲.野徑雲俱黑, 江船火獨明.曉看紅濕處, 花重錦官城.



\newpage

\section{愁眉得展}
\label{sec:102}
\textbf{林豪恩先生 - 愁眉得展}
\newline
\newline
\hyperref[sec:101]{< < < PREV SERMON < < <}
~
\hyperref[sec:toc]{[返目錄]}
~
\hyperref[sec:103]{> > > NEXT SERMON > > >}


約翰福音 14:1-14
\newline
\begin{longtable}{cl}
\hline
\hline
章節 & 經文 (和合本修訂版)\\
\hline
14:1 & \begin{tabularx}{0.7\textwidth}{X} 「你們心裡不要憂愁；你們信神，也當信我。 \end{tabularx} \\ \\ \relax
14:2 & \begin{tabularx}{0.7\textwidth}{X} 在我父的家裡有許多住處；若是沒有，我就早已告訴你們了。我去原是為你們預備地方去。 \end{tabularx} \\ \\ \relax
14:3 & \begin{tabularx}{0.7\textwidth}{X} 我若去為你們預備了地方，就必再來接你們到我那裡去，我在哪裡，叫你們也在哪裡。 \end{tabularx} \\ \\ \relax
14:4 & \begin{tabularx}{0.7\textwidth}{X} 我往哪裡去，你們知道那條路。」 \end{tabularx} \\ \\ \relax
14:5 & \begin{tabularx}{0.7\textwidth}{X} 多馬對他說：「主啊，我們不知道你去哪裡，怎麼能知道那條路呢？」 \end{tabularx} \\ \\ \relax
14:6 & \begin{tabularx}{0.7\textwidth}{X} 耶穌對他說：「我就是道路、真理、生命；若不藉著我，沒有人能到父那裡去。 \end{tabularx} \\ \\ \relax
14:7 & \begin{tabularx}{0.7\textwidth}{X} 既然你們認識了我，也會認識我的父。從今以後，你們就認識他，並且已經看見他了。」 \end{tabularx} \\ \\ \relax
14:8 & \begin{tabularx}{0.7\textwidth}{X} 腓力對他說：「主啊，將父顯給我們看，我們就知足了。」 \end{tabularx} \\ \\ \relax
14:9 & \begin{tabularx}{0.7\textwidth}{X} 耶穌對他說：「腓力，我與你們在一起這麼久了，你還不認識我嗎？看見我的就是看見了父，你怎麼還說『將父顯給我們看』呢？ \end{tabularx} \\ \\ \relax
14:10 & \begin{tabularx}{0.7\textwidth}{X} 我在父裡面，父在我裡面，你不信嗎？我對你們所說的話不是憑著自己說的，而是住在我裡面的父在做他的工作。 \end{tabularx} \\ \\ \relax
14:11 & \begin{tabularx}{0.7\textwidth}{X} 你們要信我，我在父裡面，父在我裡面；即使不信，也要因我所做的工作信我。 \end{tabularx} \\ \\ \relax
14:12 & \begin{tabularx}{0.7\textwidth}{X} 我實實在在地告訴你們，我所做的工作，信我的人也要做，並且要做得比這些更大，因為我到父那裡去。 \end{tabularx} \\ \\ \relax
14:13 & \begin{tabularx}{0.7\textwidth}{X} 你們奉我的名無論求甚麼，我必成全，為了使父因兒子得榮耀。 \end{tabularx} \\ \\ \relax
14:14 & \begin{tabularx}{0.7\textwidth}{X} 你們若奉我的名向我求甚麼，我必成全。」 \end{tabularx} \\ \\
[1ex]
\hline
\hline
\end{longtable}
日期:2011年6月12日

題目:愁眉得展

證道:林豪恩先生

經文:約翰福音14章1-14節

一．經文

14 : 1「你們心裡不要憂愁；你們信神,也當信我. 14 : 2在我父的家裡有許多住處；若是沒有,我就早已告訴你們了.我去原是為你們預備地方去.14 : 3我若去為你們預備了地方,就必再來接你們到我那裡去,我在哪裡,叫你們也在哪裡. 14 : 4我往哪裡去,你們知道那條路.」 14 : 5多馬對他說:「主啊,我們不知道你去哪裡,怎麼能知道那條路呢?」 14 : 6耶穌對他說:「我就是道路,真理,生命；若不藉著我,沒有人能到父那裡去. 14 : 7既然你們認識了我,也會認識我的父.從今以後,你們就認識他,並且已經看見他了.」 14 : 8腓力對他說:「主啊,將父顯給我們看,我們就知足了.」 14 : 9耶穌對他說:「腓力,我與你們在一起這麼久了,你還不認識我嗎?看見我的就是看見了父,你怎麼還說『將父顯給我們看』呢? 14 : 10我在父裡面,父在我裡面,你不信嗎?我對你們所說的話不是憑著自己說的,而是住在我裡面的父在做他的工作. 14 : 11你們要信我,我在父裡面,父在我裡面；即使不信,也要因我所做的工作信我. 14 : 12我實實在在地告訴你們,我所做的工作,信我的人也要做,並且要做得比這些更大,因為我到父那裡去. 14 : 13你們奉我的名無論求甚麼,我必成全,為了使父因兒子得榮耀. 14 : 14你們若奉我的名向我求甚麼,我必成全.」

三．釋經與應用1．有人先行到目的地打點14章開始,耶穌就對門徒說:「你們心裡不要憂愁」,即是說門徒正處於憂愁的狀況.對憂愁的人來說,憂愁總有原因,讓我們閱讀上文,了解令門徒憂愁的事情.13章33節記載了耶穌向門徒宣佈自己即將離他們而去,而且明言門徒不能跟著去:「孩子們!我與你們同在的時候不多了；你們會找我,但我所去的地方,你們不能去.」此乃第一項令門徒困惑的背景原因.耶穌說要離他們而去,門徒自然地問:「主啊,你去哪裡?」即是說,門徒連耶穌要去哪裡也不知道,此乃第二項令門徒困惑的背景.想到耶穌之前說過有人出賣他,加上現在要離門徒而去的宣佈,門徒似乎預感耶穌將預不測,因此彼得代表眾門徒挺身而出,保護耶穌,說:「主啊,為甚麼我現在不能跟你去?我願意為你捨命.」可惜,卻遭到耶穌「撥冷水」:「

你願意為我捨命嗎?我實實在在告訴你,雞叫以先,你要三次不認我.」(13:37-38)因此,門徒怎能不憂愁呢?

耶穌要離開,門徒問:「你去哪裡?」,「我為甚麼現在不能跟你去?」這些問題,其實是很自然的問題,而且是熟識的問題,不過,卻不是容易回答的問題.

例如父母離家上班時,幼小的兒女很可能會追著問:「你去哪裡?去做甚麼事情?為甚麼不帶我去?」

曾幾何時,為數不少的香港人舉家移民外國,其中一些父親隻身回港工作,變成「太空人」,留下妻兒子女在外地.雖然妻兒子女未必會問,但如果他們問:「爸爸,你去哪裡?為甚麼我們不能跟你去?」父親可以怎樣回答呢?

讓我們想像,成千上萬的菲律賓媽媽及印尼媽媽離鄉別井去照顧別人的孩子,如果他們的孩子問:「媽媽你去哪裡?為甚麼我們不能跟你去?」這些媽媽可以怎樣回答呢?

又讓我們想像,當趙連海為毒奶粉受害兒童申訴而被拘捕時,他年幼的兒子問:「爸爸,你去哪裡?為甚麼不帶我去?」趙連海可以怎樣回答呢?

耶穌又如何回答他的門徒這些問題呢?

14 : 1「你們信神,也當信我. 14 : 2 在我父的家裡有許多住處；若是沒有,我就早已告訴你們了.我去原是為你們預備地方去.14 : 3 我若去為你們預備了地方,就必再來接你們到我那裡去,我在哪裡,叫你們也在哪裡.

難道耶穌厭棄他的門徒嗎?他要離門徒而去,豈不是為了門徒的長遠利益嗎?難道父親捨得變成「太空人」離開自己的妻兒子女嗎?要與自己的妻兒子女別離,豈不也是為了改善他們的生活條件嗎?難道菲律賓和印尼媽媽不想親自照顧自己的兒女嗎?要離開自己的兒女去照顧別人的兒女,豈不也是為了兒女們更好將來嗎?

「你去哪裡?為甚麼我們現在不能跟你去呢?」耶穌回答門徒的話可以這樣重述:「我去為你們做一些事情,事情辦好後,就會回來接你們.」的確,耶穌是去為門徒預備地方,然而,這只是真想的一部分,坦白一點來說,耶穌是去為門徒受苦；再坦多一點來說,耶穌其實是去為門徒受死.門徒真的要跟著去嗎?耶穌忍心讓門徒跟著去嗎?

難道爸爸們忍心讓孩子跟著去建築地盤門工嗎?就算能帶孩子來,難道當家庭傭工的外地媽媽忍心讓自己的兒女看著自己的媽媽照顧其他孩子而不能照顧自己嗎?難道趙連海忍心讓小兒子跟著自己去坐牢嗎?他們可以對兒女說甚麼呢?會否都類似耶穌安慰門徒的話:「你們不要擔心,我去做一些事情,都是為了你們的好處,辦好事情就回來了.」

先知以賽亞描述上帝如何看待我們時說:「

婦人焉能忘記她吃奶的嬰孩,不憐恤她所生的兒子?即或有忘記的,我卻不忘記你.」(以賽亞書49:15)詩人也有類似的回響:「

父親怎樣憐恤他的兒女,耶和華也怎樣憐恤敬畏他的人!」(詩 103:13)

耶穌請門徒不要憂愁,雖然有些事情暫時不明白,但耶穌請他們相信,他的行動,是為他們的益處而做的.

2．有人接載去目的地

14 : 4 我往哪裡去,你們知道那條路.」 14 : 5 多馬對他說:「主啊,我們不知道你去哪裡,怎麼能知道那條路呢?」 14 : 6 耶穌對他說:「我就是道路,真理,生命；若不藉著我,沒有人能到父那裡去.

彼得問耶穌:「主啊,你去哪裡?」耶穌回答:「

我所去的地方,你現在不能跟我去,以後卻要跟我去.」(3:36)耶穌接著說:「

我在哪裡,叫你們也在哪裡.我往哪裡去,你們知道那條路.」(14 :3- 4)耶穌這樣的回答,令多馬大惑不解:「主啊,我們不知道你去哪裡,怎麼能知道那條路呢?」

多馬的疑惑是可理解的,門徒不知道耶穌到哪裡去了,又如何可能找到路去會合耶穌呢?我們使用巴士公司的網上查詢系統,都必須要提供目的地,或者巴士路線,系統也能夠運作.

最令多馬費解的是,多馬認為耶穌沒有提供目的地和路線圖,但耶穌卻指出門徒其實是知道目的地和路線圖的.門徒與耶穌之間,為甚麼存在著如此大的認知上的差異呢?

對於門徒來說,目的地是一個具體的地方,因此,路也是一條具體的路.若果耶穌告訴他們一個具體的地方,他們就可以尋找路線,尋到路線後,就可以想辦法,例如:步行,坐船,乘車,到耶穌所在的那個地方去.

從耶穌回答多馬的說話來看,多馬誤解了耶穌所說的目的地,也誤解了耶穌所說的道路.耶穌說:「

我就是道路,真理,生命；若不藉著我,沒有人能到父那裡去.」原來,耶穌所說的目的地,並不是一個地方,而是一位對象,就是天父；耶穌所說的通往目的地的道路,也不是某一組路線,而是一個人,這就是耶穌自己.

以多馬的想法,他要往目的地,就要靠自己.有資源的,可以使用交通工具,沒有資源的,就憑著毅力一步一步走往目的地.多馬的想法,會否也是我們的想法呢?我們的目的地在那裡呢?天堂?通往天堂的路是如何走的呢?我們走了多少步呢?還有多少步呢?有否擔心自己半途而廢呢?有否憂慮自己最後走不到目的地呢?

耶穌澄清,根本就沒有一條讓我們憑自己的資源或能力可以走到目的地的道路；耶穌安慰我們,如果要用「道路」作比喻,他自己就是那通往目的地的「道路」；如果要乘車,耶穌就是那車；如果要搭船,耶穌就是那艘船.往天父那裡去的人,都是坐在耶穌的臂膀上,或者靠在耶穌的背上,不需要自行前往的.

3．目的地竟然就在這裡

14 : 7 既然你們認識了我,也會認識我的父.從今以後,你們就認識他,並且已經看見他了.」 14 : 8 腓力對他說:「主啊,將父顯給我們看,我們就知足了.」 14 : 9 耶穌對他說:「腓力,我與你們在一起這麼久了,你還不認識我嗎?看見我的就是看見了父,你怎麼還說『將父顯給我們看』呢? 14 : 10 我在父裡面,父在我裡面,你不信嗎?我對你們所說的話不是憑著自己說的,而是住在我裡面的父在做他的工作. 14 : 11 你們要信我,我在父裡面,父在我裡面；即使不信,也要因我所做的工作信我.

面對耶穌宣佈離去,門徒各有憂愁.彼得要求現在就跟著耶穌去,不但得不到批准,更被耶穌挫敗,斷言他只是匹夫之勇,意氣用事.耶穌保證了有重聚的時候,多馬的憂愁是不知道耶穌的目的地,也不知道圖線圖,根本沒有辦法去找耶穌重聚.耶穌進一步澄清了目的地就是天父的時候,就引起了腓力的關注,他對耶穌說:「主啊,將父顯給我們看,我們就知足了.」對於腓力來說,如果目的地是天父,可否現在就預先看一看呢?

腓力的期望,其實是可以理解的.在香港開售的樓宇,通常都只是一個地盤,因此,發展商會設計示範單位,讓客人先看一看他們購買的大概是甚麼東西.推銷食品的策略也類似,售貨員鼓勵客人試食,先淺嘗一口,知道自己捧回家的大包小包是甚麼味道.買樓之前沒有示範單位看,只憑尺寸資料想像,怎能安心呢?買零食之前能先嘗一口,就放心了.耶穌的門徒可能也有類似的期望,到天父哪裡之前,能夠預先看一眼,就滿足了.

雖然腓力的期望對我們來說是可以理解的,但卻令耶穌感到驚訝,因此他反問腓力:「腓力,我與你們在一起這麼久了,你還不認識我嗎?看見我的就是看見了父,你怎麼還說『將父顯給我們看』呢?」約翰福音沒有記載腓力的反應,不過,我們可以猜想,聽到耶穌說,看見他就等於看見天父,門徒可能感到非常驚訝.耶穌進一步向門徒解釋,門徒跟隨耶穌那麼多年,應該從他的言行之中,知道耶穌所說的話,所做的工作都是天父的,也應該知道耶穌在天父裡面,天父在耶穌裡面.

縱使耶穌繼續解釋,門徒很可能繼續感到驚訝.事實上,怎可能不驚訝呢?驀然回首,那人正在燈火闌珊處.包括門徒在內敬拜上帝的猶太人,一直渴望看見天父,想不到看著這位每天與自己朝夕共對的耶穌,就等於看見了天父.有些時候,我們四處尋找一些東西,以為在將來某個時候才會遇上,竟然就在身旁.就好像有一個人尋找眼睛,在睡房找不到,在客廳找不到,在洗手間也找不到.原來,這個眼鏡已經戴在眼睛前.又好像某些人四出尋找真愛,其實真愛就是一直守候身旁的配偶.

腓力說:「主啊,將父顯給我們看,我們就知足了.」到底腓力會否感到知足呢?腓力會否感到知足,並不決定於耶穌會否答應他把天父顯給他看,而是決定於他是否願意接受看見眼前的耶穌,就看見天父.基督徒又會否感到知足呢?還是充滿憂愁呢?答案很可能也決定於我們自己.曾幾何時,宣道會的創辦人宣信博士雖然事工果效不錯,但卻感到自己不夠聖潔,也不夠屬靈能力.有時,他似乎覺得已經得著了聖潔及能力,他就奮力抓住它,惟恐失掉,以致誠惶誠恐.為了追求靈力,他特別趕往千哩外的芝加哥,參加慕迪主講的培靈會.在慕迪講道前,有一位弟兄見證說:「我來此地,盼望慕迪弟兄幫助我,但昨夜我看見主耶穌,並一直注目瞻仰祂,我就覺得今後一生除主以外我也無求.」當宣信聽到這個見證時,觸及他的深處,「你所需要的就是主耶穌,到祂那裡就好了!」那晚,他未等慕迪講道結束,就趕著回家去.他回到自己房間,關上門,懇切禱告,寫下了「今要主自己(Himself)」這首詩歌.詩歌的第一節如此表達宣信當時的體會:

前所要是祝福,今所要是主；前所憑是感覺,今靠主言語；前所要是恩賜,今要賜恩主；前我尋求醫治,今要主自己.

4．承繼耶穌的工作

14 : 12 我實實在在地告訴你們,我所做的工作,信我的人也要做,並且要做得比這些更大,因為我到父那裡去. 14 : 13 你們奉我的名無論求甚麼,我必成全,為了使父因兒子得榮耀. 14 : 14 你們若奉我的名向我求甚麼,我必成全.

耶穌請門徒不要憂愁,因為門徒根本沒有時間憂愁.門徒的時間用來做甚麼呢?耶穌吩咐他們接續他所做的工作,而且要做得比耶穌在世上時更大.如何做好耶穌的工作就佔據了門徒的所有心思意念,那再有時間去憂愁呢?當門徒承繼耶穌的工作時,耶穌答應按他們的祈求成全,使這些工作能夠延續,透過這些工作讓人們得知上帝的美善.

耶穌的工作,有些是只有他做才有意義,其他人做是沒有效用的,例如釘十字架,我們當中任何一個人死在十字架上,都沒有救贖的效用.然而,耶穌復活後對門徒的吩咐,無疑是門徒應該繼承的重要工作之一:「

所以,你們要去,使萬民作我的門徒,奉父,子,聖靈的名給他們施洗,凡我所吩咐你們的,都教導他們遵守.看哪,我天天與你們同在,直到世代的終結.」(28 : 19-20)

耶穌最後的吩咐指出,當門徒繼承耶穌的工作時,耶穌就天天與他們同在,即是說,門徒到達天父那裡的時候,不是遙遠的將來,而是現在,就是門徒做耶穌所做的工作的當下.此時此地信耶穌的人們,我們是否在等候遙遠的將來上天堂去見天父呢?原來,天父不受天堂的限制,他在我們做耶穌所做的工作時就與我們同在了.

四．總結憂愁總有原因,人生不如意事十常八九,怎能不憂愁呢?耶穌沒有輕視令人憂愁的事情,也沒有否認憂愁的時候.那麼,他請門徒不要憂愁的基礎在哪裡呢?「憂愁」這個字,在約翰福音出現了五次,其中兩次用在門徒身上,分別在14章1節:「

你們心裡不要憂愁」,以及重覆在14章27節:「

你們心裡不要憂愁」.另外三次竟然用在那位叫人不要憂愁的耶穌身上,第一次在11章33節,形容耶穌看見剛死去的拉撒路的家人為拉撒路哀哭的心情:「耶穌看見她哭,並看見與她同來的猶太人也哭,就心裡悲嘆,又甚憂愁」；第二次在12章27節,耶穌用來表達自己快要走上十字架時候的心情:「

我現在心裡憂愁,我說甚麼才好呢?說『父啊,救我脫離這時候』嗎?但我正是為這時候來的.」最後一次出現在13章21節,形容耶穌快樂被出賣的心情:「

耶穌說了這些話,心裡憂愁,於是明確地說:『我實實在在地告訴你們,你們中間有一個人要出賣我.』」

門徒怎能不憂愁呢?除了上述我們探討過的耶穌在與他們對話中的解釋和安慰之外,最重要的,是耶穌讓自己憂愁,換取可以對門徒說:「你們心裡不要憂愁」.

若我們也落在眾多憂愁之中,可能會聽到耶穌說:「你們心裡不要憂愁」,我們心裡會問:「怎可能呢?」耶穌可能會回答我們:「因為我憂愁了,以致你們可以不用憂愁.你們繼續做我做的工作就足夠了.」

有時,父母拖著疲倦的身軀及勞累的心靈下班回家,兒女們問:「爸爸媽媽,工作不開心嗎?」父母可能會回答:「不用擔心,你們好好讀書就是了.」



\newpage

\section{參悟死生}
\label{sec:103}
\textbf{鄧瑞強博士 - 參悟死生}
\newline
\newline
\hyperref[sec:102]{< < < PREV SERMON < < <}
~
\hyperref[sec:toc]{[返目錄]}
~
\hyperref[sec:104]{> > > NEXT SERMON > > >}


約翰福音 17:1-11
\newline
\begin{longtable}{cl}
\hline
\hline
章節 & 經文 (和合本修訂版)\\
\hline
17:1 & \begin{tabularx}{0.7\textwidth}{X} 耶穌說了這些話，就舉目望天，說：「父啊，時候到了，願你榮耀你的兒子，使兒子也榮耀你； \end{tabularx} \\ \\ \relax
17:2 & \begin{tabularx}{0.7\textwidth}{X} 因為你曾賜給他權柄掌管凡血肉之軀的，使他把永生賜給你所賜給他的人。 \end{tabularx} \\ \\ \relax
17:3 & \begin{tabularx}{0.7\textwidth}{X} 認識你—獨一的真神，並且認識你所差來的耶穌基督，這就是永生。 \end{tabularx} \\ \\ \relax
17:4 & \begin{tabularx}{0.7\textwidth}{X} 我在地上已經榮耀你，你交給我做的工作，我已完成了。 \end{tabularx} \\ \\ \relax
17:5 & \begin{tabularx}{0.7\textwidth}{X} 父啊，現在求你使我在你面前得榮耀，就是在未有世界以前，我同你享有的榮耀。 \end{tabularx} \\ \\ \relax
17:6 & \begin{tabularx}{0.7\textwidth}{X} 「你從世上賜給我的人，我已把你的名顯明給他們。他們本是你的，你把他們賜給我，他們也遵守了你的道。 \end{tabularx} \\ \\ \relax
17:7 & \begin{tabularx}{0.7\textwidth}{X} 現在他們知道，你所賜給我的一切都是從你那裡來的； \end{tabularx} \\ \\ \relax
17:8 & \begin{tabularx}{0.7\textwidth}{X} 因為你所賜給我的話，我已經賜給他們，他們也領受了，又確實知道，我是從你出來的，並且信你差了我來。 \end{tabularx} \\ \\ \relax
17:9 & \begin{tabularx}{0.7\textwidth}{X} 我為他們祈求，不為世人祈求，卻為你所賜給我的人祈求，因他們本是你的。 \end{tabularx} \\ \\ \relax
17:10 & \begin{tabularx}{0.7\textwidth}{X} 凡是我的都是你的，你的也是我的，並且我因他們得了榮耀。 \end{tabularx} \\ \\ \relax
17:11 & \begin{tabularx}{0.7\textwidth}{X} 我到你那裡去；我不再留在世上，他們卻在世上。聖父啊，求你因你的名，就是你所賜給我的名，保守他們，使他們像我們一樣合而為一。 \end{tabularx} \\ \\
[1ex]
\hline
\hline
\end{longtable}
日期:2011年6月19日

主題:參悟死生

經文:約翰福音17:1-11

證道:鄧瑞強博士

各位弟兄姊妹,主內平安.

我今日的講道,是講論主耶穌對死亡和生命的看法.但講主耶穌的看法前,我要講一講其他死亡和生命的事件.

「六四」,一個在歷史上將會永遠被記起的日子.無數人紀念這日子,無數人紀念在這日子死去的人.這些人的死,不是生老病死這自然規律裡的自然的死.他們的死,飽含著對生命應該如何活下去的叩問.他們的死,充滿著對國家應該如何走下去的期盼.他們的死,叫那些不去思考生命的意義為何的人不能再不思考生命的意義為何,或者說,他們的死,將生命的重量放入其他活著的人的生命裡,叫活著的人不再活得輕省.

昨天已逝,已進入記憶中的永恆,它將永遠訴說著昨日的故事.死亡,叫人不再忽略生命的重量.

昨天已逝,明天又如何呢?

端午節故事的中心,是一個詩人的自殺事件,一個被稱為愛國詩人屈原的自殺事件.詩人不是普通人,詩人是觸覺比人敏銳,思想比人深遠,視域比人廣闊的人.詩人不是看到事情的表面,他們看到事物的實質.詩人不是人云亦云,他們指出生命的前景.詩人不是訴說現況,他們訴說真理.所以詩人的自殺不是一件普通事件,它意味著一種「前景」的崩塌,它意味著一種「真理」的死亡.這種「真理」的死亡,叫那些對這「真理」認真的人,不得不自殺.這是面對生命無所寄托,無所依歸,也不知向誰傾訴之下的痛苦的死亡.屈原面對這個絕境,他在這絕境中走投無路,他只能自殺.

中國學者劉小楓寫了一本書,叫《拯救與逍遙》(上海人民出版社),開首的篇章討論屈原的自殺事件,他稱屈原「是歷史上第一位自殺的大詩人.他的偉大首先於他敢於自殺.」(頁60)以下我對屈原的講論,見解源自這書.

屈原的自殺,固然有其歷史的前因後果,今日要講的,是這事件的精神意義.屈原是詩人,但更根本地說,他是一個政治家.他早年深得楚懷王的信任,是楚懷王的政治顧問,在那個合縱連橫的世界,他使楚國能在秦國面前有一時的風光.

作為政治家,屈原骨子裡持守的,是儒家精神.生命的意義,在於肯定自己,實現一完美的人格.在儒家「內聖外王」的格局下,要成就個人完美的人格,就要成為一出色的政治家.若果政治的路途順利,他會很滿意,覺得生命的意義得到充分發揮了.問題是:若果政治的路途不順,則生命的意義能在何處安頓?

若果遇到昏君,你要徹底質問昏君的昏庸,你要在昏君以外有一善的標準去質問?若果遇到暴君,你要徹底質問暴君的殘暴,你要在暴君以外有一善的標準去質問?但是,屈原唯一可以持守的,是自己對完美人格的追求,而這種追求,能落實的地方,就是服務於現實的君主政體.違反這種政體,質問這些君主,就變成不愛國.作為一個「愛國」詩人,屈原想像不到在這君主體制以外能有什麼堅實的基礎去質問這種君主體制,縱使這君主體制已腐敗不堪.

屈原的痛苦是因為他「愛國」,但他只有這個世上的現實的國,他想像不到這世界的現實的國以外的國.他的生命意義,只能寄托於服務這個現實的國.他的「內聖外王」理想不能想像另一種「愛」,另一種「國」.他只能擁抱這個現實的國.但這現實的國已腐敗不堪了,這個現實的國已放逐他了,這個現實的國已發展出一種力量毀滅他了,但是,他只能「忠」於這個國,他只能日夜思念這個國,否則,他無處發揮他的「忠」,無處發揮他的「愛」,無處實現他的生命意義.縱使楚國已放逐他,縱使他不能腳踏自己國家堅實的土地,他仍不能走向另一塊堅實的土地上.他仍眷戀現實的楚國,縱使腳下已空無一物.他生命意義寄托之所在,現在已被懸空,但他只能維持在這懸空中,他找不到另一塊堅實的土地.

畢竟,屈原除了是政治家,他也是詩人.他試圖在現實的痛苦中找到出路.他寫出《天問》這詩篇.《天問》這詩篇全是問題,一共問了170多個問題.第一個問題是:「遂古之初,誰傳道之?」即:在太初之始,誰引導萬物的開展?屈原問了很多萬物何所立足的問題.之後,他問了很多自然現象及人事的問題,叩問個人生死的意義問題.最後的一個問題是:「何試上自予,忠名彌彰?」即:為何自己欣賞自己告誡君主,這是因為要顯揚自己忠義之名嗎?最後,屈原落腳於現實政治.一系列的天問,一連串的意義追尋,他仍走不出要在現實政治裡留下忠義之名.

一個詩人,一個對生命意義抱嚴肅態度的人,當自己的生命意義遇到極大危機時,向天發問,是正常不過的.向天發問,可以是一個契機,讓人從自己的自我滿足走出來,讓人對現實政治的腐敗找到一個屬天的標準去審問,讓人走出純粹現實的維度而思索永恆的維度.簡言之,一個嚴肅的天問,可以打開塵世的人和天對話的大門.天問,可以變化為對塵世腐敗的質問,可以打開另一個世界,讓生命懸空的人找到另一個落腳點.

但是,屈原這個詩人畢竟太「愛國」了,他眷戀塵世的現實的國.他這個「天問」是白問.孔子說過:「天何言哉?」上天是無語的.屈原的「天問」終歸只是絕望的人的哀鳴,他的「問」從不期望「答」,因為在他看來,天是無言對應人間的痛苦的.

劉小楓這樣說:「整個《天問》都在詢問諸現象(宇宙,人生,社會歷史)的終極根據,彷彿是在進行一場超驗的追問,而實際上卻一再徘徊於諸現象,並沒有真正在進行超驗的追問.這實際上等於,《天問》表面上在作超驗的追問,而事實上卻是什麼也沒有問.」(頁149-150)

一個痛苦的詩人,他只想到現實的自我和政治,他想透過政治去肯定自我的意義,但政治否定了他.他發出天問,但他沒有想過「天」會回答,他也不試圖走向「天」的超越境地,他心中的「天」,來來去去只是人間的「地」.他渴望超越人間的腐敗,卻又只能在腐敗中找超越腐敗的力量.他注定失望.

本來,中國的政治人物在政治失意時,可走向道家.陶淵明如是,蘇東坡也如是,可寄情於山水,忘記現實的社會與國家.但是,屈原堅持「愛國」,他堅持儒家的理想,他不甘願走向道家.他堅持要在君主體制中完成他的人生意義,縱使這體制已腐敗.他堅持要在人間體制中完成人生理想,雖然他根本找不到超越體制罪惡的基礎.他發出「天問」,縱使他自己準備作答.

在汩羅江邊,他不能像蘇東坡那樣說:「小舟從此逝,江海寄餘生」.他堅持愛國的理想,他無力超越他愛的國對他的拒絕.他只能自殺,以表達他的忠心.這就是他的《天問》的最後問題,如何能保證自己忠義之名.

我們回到今日的講道

經文:約翰福音17: 1-11.

這段經文表達了主耶穌在受死前夕的內心思念.

這不是「天問」,而是「祈禱」,是和「天上的父」的內心談話.

主耶穌活了相當日子,他明明白白地知道他的生命是與塵世的罪惡相對抗的,他明明白白地預計到這種生命的下場是死亡.他正在面對這種死亡的來臨.

我們看看主耶穌內心和天父的對話.

17:1　耶穌說了這話,就舉目望天,說:「父啊,時候到了,願你榮耀你的兒子,使兒子也榮耀你；

17:2　正如你曾賜給他權柄管理凡有血氣的,叫他將永生賜給你所賜給他的人.

17:3　認識你─獨一的真神,並且認識你所差來的耶穌基督,這就是永生.

17:4　我在地上已經榮耀你,你所託付我的事,我已成全了.

17:5　父啊,現在求你使我同你享榮耀,就是未有世界以先,我同你所有的榮耀.

有否留意,這段經文最常出現的字眼是「榮耀」?

主耶穌是否像屈原那樣,在人生意義幻滅下等待死亡嗎?

不是,他完全體會到生命的意義,體會到生命的意義要經歷一種為他人而死的死亡去呈現出來.

他是腳下空無一物,要以自己的力量去抵抗人生的橫逆嗎?

不是,他與天父保持緊密連繫.他的生命的意義緊緊地安放在天父手中.他到世上來,不是要完成自己,而是要完成天父的心意.他面對塵世的種種體制,不是要依附體制去完成自己的生命,而是要將永生賜給「凡有血氣的」.永恆的生命,生命的完滿意義,不是從自我實然而來,不是從體制而來,而是從天而來.這個「天」,不是「天何言哉」的天,不是無言的天.天上有父,這是天父世界.這天父俯視人間的痛苦,他也與人共苦.他聆聽人的生命之問,他作出回應.他差來主耶穌,他進入人間苦難的深淵,他以愛的行動去解救苦難中的罪人.

人的終極問題源自罪.這不是自我實現得不夠,這不是遇不到明君以致不能實現生命意義的問題,這是人的視野狹窄得只有塵世的問題,這是人不願意與天父交往的問題,這是人只走向自己而不願走向天父的問題.約翰福音說:「認識你─獨一的真神,並且認識你所差來的耶穌基督,這就是永生.」這是一種生命的轉向,認識自身以外的神,不以自己為神,這才是永生.

主耶穌這段禱文,常常提及「榮耀」.主耶穌並不像屈原那樣,在絕望中自殺.主耶穌將要受死,這是對抗塵世的罪惡的死,這是向世人表明神愛我們的,因愛而有的死,這份愛是對所有真誠地作出「天問」的人的真誠的回答.主耶穌實現天父的愛,並視實現這份愛而有的死亡為生命的榮耀.

主耶穌繼續祈禱說:

17:6　「你從世上賜給我的人,我已將你的名顯明與他們.他們本是你的,你將他們賜給我,他們也遵守了你的道.

17:7　如今他們知道,凡你所賜給我的,都是從你那裡來的；

17:8　因為你所賜給我的道,我已經賜給他們,他們也領受了,又確實知道,我是從你出來的,並且信你差了我來.

這段禱文最常出現的字眼是「賜給」.

一切都是天父賜的.主耶穌是天父賜給世人的,世人的永生是天父賜予的.一切都是天父賜的禮物.

主耶穌在死亡的前夕,想到的是,世間充滿天父的禮物.當然,他自己也是天父賜給世人的禮物.他想到他的死,也是天父賜給世人的愛的禮物.死,不是無奈,不是在死胡同中用命拼出來的出路.在主耶穌眼中,死是一份禮物.他的死,是對在生者的禮物,讓他們看到天父的慈愛.有些人的死,的確是一份禮物,他們將重量加入我們的生命中,讓我們不得不思索生命到底為何.主耶穌作為天父給我們的禮物,他的死最少讓我們明白,神對我們的愛是真的,他是願意回應我們生命的「天問」的.

之後,主耶穌為他的門徒祈求.

17:9　我為他們祈求,不為世人祈求,卻為你所賜給我的人祈求,因他們本是你的.

17:10　凡是我的,都是你的；你的也是我的,並且我因他們得了榮耀.

17:11　從今以後,我不在世上,他們卻在世上；我往你那裡去.聖父啊,求你因你所賜給我的名保守他們,叫他們合而為一像我們一樣.

主耶穌即將離去,他特別記掛那些跟從他的人.這些人應該在主耶穌裡體會過天父大愛的,但他們仍活在世上.活在世上,對門徒而言,是一種使命,一種挑戰,但同時是一危機.

作為一種使命,門徒體會到他們被差派到世上,像天父差派主耶穌到世上一樣,他們需要與天父對話,將天父的愛實現在人間.

但這也是一種挑戰,因為門徒的生命並未能與天父緊密結合,像天父和主耶穌合而為一一樣.若未能與天父合一,則我們便不能將生命,將死亡,視為對他人的禮物一樣向他人奉獻出來.

活在世上,對門徒而言,也是一大危機.我們會被世界同化,我們會忘記天父而只顧塵世,我們會忘記天國而沉溺於地上的君主體制,我們會忘記天父的愛而追逐世上的成敗.如此,我們可能再走屈原的舊路,在理想幻滅之際,不能回到天父的懷抱.

主耶穌離世前夕,很擔心我們.

今日,我們明白主耶穌的死,領受他的恩典和大愛,我們又將如何活下去呢?我們會讓主耶穌擔心嗎?但願榮耀歸於聖父,聖子,聖靈.阿們.



\newpage

\section{都市牧羊}
\label{sec:104}
\textbf{梁元生教授 - 都市牧羊}
\newline
\newline
\hyperref[sec:103]{< < < PREV SERMON < < <}
~
\hyperref[sec:toc]{[返目錄]}
~
\hyperref[sec:105]{> > > NEXT SERMON > > >}


日期:2011年6月26日

題目:都市牧羊

證道:梁元生教授

各位弟兄姊妹,早上好:

今天是耶穌復活四十日後升天的日子,是聖靈降臨節.在耶穌升天前,曾經對門徒交待說:「你們要去傳福音給萬民聽,由耶路撒冷,到撒瑪利亞,直到地極.」這個傳福音的吩咐,就是一般基督徒所說的「大使命」.我們常常會覺得由耶路撒冷,到撒瑪利亞,再到世界的地極,是一個不斷往外延伸的過程,是一個由核心到邊沿,由繁華到荒漠的過程,越走越遠,越走越辛苦,但為了天國,為了完成「大使命」而必須吃苦,開荒的過程.今天我想在這裡和大家反思耶穌「大使命」中所說的「地極」.我想提出的是:「地極」 and \#8212;-基督徒從耶路撒冷走來,從撒瑪利亞走來,可能並非邊遠之地,並非荒漠之地；這地極,就二十一世紀的情況來說,可能就在我們身邊不願之處,更有可能就在我們身處之地 and \#8212;-城市.今天我要講的題目是「都市牧羊」,目的想要探討一下在現代城市化的大潮下,基督徒如何在城市的氛圍裡過基督徒的生活,及如何實踐傳福音的大使命.

一,    史上規模最大的「城市化」我是學歷史的,而研究的興趣主要是城市史；因此,對於人口遷移,房屋發展,工業商業等等與城市有關的題目,都略有涉獵.据我所知,在過往的兩個世紀中,有過不同形式和不同規模的「城市化」運動；而最受人注意並且規模最大的則要算十九世紀末到二十世紀初數十年間北美洲(特別是美國)的「工業城市」興起的現象(Rise of the Industrial Cities).由於機械發明和工業興起(多得該時代的許多發明家如發明電燈的愛迪生,發明電話的貝爾等),以及全國鐵路網絡的建成,造成大規模的人口由農村向新興工業城市流動.一些曠源豐富和位于鐵道網絡樞紐的工業城市,如芝加哥,匹玆堡,奇里夫蘭,辛辛那提等迅速膨脹,人口劇增,力逼傳統的商業城市波士頓和紐約.這個「城市化」現象,帶出許多美國史上有關社會改變和文化改變的問題,包括倫理的,宗教的,生活上的適應,以及城鄉矛盾,以致在過渡和轉折期間出現的種種問題.今天我對此不會作詳細的分析,只想借此帶出近代歷史上另一次的「城市化」運動,而這次的「城市化」規模更大,牽連更廣,對人口流動,社會變化,經濟結構,和生活方式,都有更為廣大而深遠的影響.這次「城市化」運動,要比美國十九世紀末期的「城市化」,來得更重要,卻是在我們身邊發生的.我想要說的,就是過去這二,三十年來中國本土上的「造城運動」.只要大家回頭北望,就知道我們的「邊城」深圳,三十年前只有羅湖口岸的幾條街道和幾十萬人口,規模不大,但二十年間深圳已經變成一座大城市,人口超過千萬,高樓聳立,百貨紛陳,熱鬧非凡.連年產經濟總值,也已經逼近香港了.聽說今年或明年就可能超過香港!深圳只是一個例子,其餘國內的許多地區,成千上百的鄉鎮小城,在過去二十年間也都迅速發展,大興土木,迅速地變成人口數百萬,甚至上千萬的大城市.大量農村人口流向城市,把城市本來較為擠迫的空間都填得滿滿的,甚至高架橋下,火車站的廣場和空地,都是容納農村流民的城市空間；當然,也有廠商為勞工而建的簡陋房舍,和農民自己臨時架搭的棚屋和帳房,形成了現代中國城市的一幅奇異的景觀.

根據最近的資料,現時中國內地由於城市化而出現的人口數百萬或以上的都市,有超過六百個之多,而城市人口也佔了全國人口四成左右.已往中國全國人口大概有九成是住在鄉村的,現在那所謂六成鄉居的人,也多數是住在鄉鎮裡的,真的住在農村裡的農民,已經不多了.這是中國歷史中史無前例的,規模最大的「城市化」現象,對中國社會,經濟,文化,都會帶來重要的變化.

以下我擬從基督教的角度去探討一下中國城市化的問題,以及從聖經中學習在城市中牧養和傳福音的經驗.

二．都市作為大牧場都市牧羊,是城市教會需要面對的問題.我們所講的羊,當然是指著人來說的.他有兩種:一種是「無牧的羊」,即在都市裡迷失而沒有信仰的人；另一種是「主的羊」,是都市教會的人,是信徒,是跟隨耶穌的人.我們思考「都市牧羊」這個問題時,即是說在現代城市環境中信徒如何生活,以及如何傳福音的問題.

首先,我們來思考一下這群無牧的羊與主的羊所共同生活的地方 and \#8212;-都市.在昔日,甚至在聖經的描寫中.牧羊的地方和今日有很大的差異.最佳的牧場,當然是草原,是青草地,溪水旁.牧人帶領羊群覓食,到有水草之處,一邊放牧,一邊高歌,又或手上一枝竿,口邊一支笛,在群山之中,綠草之上,儼然一幅美麗安詳的圖畫.這個情景可能發生,但另一種情況也可能出現,就是山路崎嶇,草叢稀少,有羊在走路不小心時滑下山徑,咩咩哭啼；也有可能走失一只,要牧羊人四處尋找.這些情況在聖經中常常提到,大概是市場發生的事.再有一些情形,就是經過曠野山林,難免遇到豺狼猛虎,或者其它怪獸,對於羊群,隨時有生命的威脅,而作為牧羊人,也就戰戰兢兢,注視環境,希望所有羊群能夠安全渡過險境.

現代的人,對荒漠草原的牧羊生活,都不容易體會和琢磨了；但現代都市中的教會,現代都市中生活的信徒,其實險阻和困難比起昔日山林野地之間的群羊並不少,而都市中的困難和危險,和野地牧羊,情況還有許多相似之處呢!

在石屎森林中,在幽暗夜燈下,豺狼四伏,毒蛇圍繞,要尋找可吞吃的群羊；在大自然我們看見的猛獸,張牙舞爪,大聲咆哮的時候,我們是較為容易嗅到或見到身邊的危險的；換言之,在自然的曠野和牧場裡,我們較為容易辨識野獸和敵人,也較為容易了解甚麼是平安,甚麼是兇險；但是在大都市中,危險和狼毒,誘惑和陷阱,往往會躲在濃厚的化妝及閃耀的霓虹燈光背後,不容易為人察覺.都市牧場,機關處處,不單止羊群容易涉險,連牧羊人也常常會受到欺騙,或自顧不暇.在過去研究都市文化的書籍中,對於城市化最大的指責是使到年輕人道德墮落,失去一些農村或傳統賦予他們的基本價值.城市中充滿各樣的聲色犬馬的引誘,容易令人迷失.這個「罪惡之城」(the wicked city)的說法,一直非常普遍.今天許許多多的中國人進入大城市,尋找工作和生活,也碰到十九世紀後期美國農村青年遷入工業大城市時相同的問題,但情況更為嚴重,因為農村裡所給予他們的基本價值也是支離破碎和殘缺不全的,到了城市後,他們就更為空虛迷失.這種無牧的羊群,在中國城市中,遍地皆是.是今日中國基督教,特別是中國城市教會面對的最大的挑戰.我今天早上沒有太多時閒去分析中國城市教會的種種問題,只想強調一下:城市,是大使命中所說的地極,城市化,是當前中國基督教要嚴峻面對的問題,我們應多加留意,和多加思考.今天在八問題提出之後,我想回到聖經,用保羅的城市經驗來反思我剛才所提出的都市牧羊和相關的一些的問題.

三．都市中的牧羊人 and \#8212;-保羅使徒時代的保羅位複音的緣故走過不少地方,他的足跡大概踏遍了現在中亞西亞和歐洲南部地區,其中所到之處當然少不了各地的大小城市.根據使徒行傳記載,他到過的城市由大馬士革開始,有路司德,安提阿,耶路撒冷,雅典,尼亞坡里,哥林多,以弗所,特羅亞,羅馬等等.以下我們想透過保羅的城市經驗,反思一下在現代城市中作為信徒及作為牧人的種種問題和挑戰.

由於保羅到過的城市不少,今天我們只選擇其中以一個城市中保羅的際遇來思考一下.這個城市就是雅典.

雅典是座大城.在當時可以說是百方輻輖,商旅雲集,各種文化和族裔皆匯聚於此,可以說得上是個國際大都會.雅典的商業和文化,和現代許多城市很接近,所以保羅的經驗對現代城市教會和城市中生活的基督徒也很有參考意義.雅典的興起,在主前第四,五個世紀,它不但是地中海貿易的一個含蘊中心,也是文化和政治的發源地.是西方哲人蘇格拉底和柏拉圖的故鄉,這裡有文士,學者,哲人,貴冑,藝術家,商旅,工匠,以及各式各樣和各個階層的外來人.城市就是容納不同人口,不同文化和不同階級的地方.雅典就是這樣一個城市.香港也是這樣一個城市.

保羅來到雅典,會見了住在那兒的基督徒,也和當地的虔誠的猶太人談論.他在雅典的事跡,記載在使徒行傳第十七章.

在這個繁華熱鬧的都市中,充滿了各樣的挑戰,尤其是要為主傳福音的保羅.他來到這個城市,令他印象最為深刻的是甚麼?

首先是他踫到一批學者.18節那裡說,那裡有伊壁鳩魯和斯多亞兩門的學士,與他爭論.城市中往往會有許多知識分子,會有很多不同的學派和意見, 對人生,對宗教,對終極關懷都有其獨特的看法.這樣的一個依賴知識和理性的對基督教的挑戰,不但在雅典如是,我相信在香港也如是.即使在我們的校園裡,理性的和知識對宗教的挑戰也是十分普遍的.所以,城市教會和城市基督徒要面對的第一個挑戰,是知識的挑戰,也可以說是理性的挑戰和科學的挑戰,the challenge of intellectualism. 相信你身邊就會有不少這樣的人,像基督教做出不同形式的學術問難.

其次,保羅在雅典所見所聞的,令他印象深刻的,就是城中的人,「

都不顧別的事,只將新聞說說聽聽.」(21節)這是另一種我們熟悉的城市文化:把新聞說說聽聽.政府的新聞,社會的新聞,名人的新聞...大家都沒有用嚴肅的態度看待新聞,就是身邊發生的大小事情,甚至關乎國家社會的事情!城市人只喜歡把新聞拿來作閒聊的話題,說說聽聽,就是gossip的意思.城市中的教會也有許多猜疑和糾紛,而許多麻煩大概也是從說說聽聽的gossip而產生的.因為城市中可以收聽到不同消息和新聞的渠道很多,造成眾說紛紜,莫衷一是的情形,很難形成共識,對於眾人一心的信仰歸屬感也不容易形成.這是城市挑戰的第二种:閒話及疏離the challenge of alienation and gossip.

但最令保羅心焦的,當然是雅典城裡滿城的偶像.16節:「

保羅在雅典等候他們的時候,看見滿城都是偶像,就心裡焦急.」

雅典城中有很多偶像.因為雅典的文化傳統中,也有許多天將鬼神,並不為怪.如天神宙斯,太陽神阿波羅,女神雅典娜,亞底米等等,都有自己的神廟.都在希臘傳統中有重要的地位和在民眾間有很大的影響力.

保羅在雅典和他們日夜相處,深深感到偶像對人的影響,所以心裡焦急.他為此除了心焦,還常與人辯論.就是因為面對各種偶像的挑戰,使到保羅發表了他那篇著名的衛道神學講章:24-32節:「

創造宇宙和其中萬物的神,既是天地的主,就不住人手所造的殿；也不用忍受服侍,好像缺少甚麼；自己倒將生命,氣息,萬物,賜給世人.他從一本造出萬族的人,住在全地上,並且預先定準他們的年限,何所住的疆界；要叫他們尋求神,或者可以喘摩而得,其實祂離我們各人不遠；我們生活,動作,存留都在乎祂....我們既是神所生的,就不當以為神的神性像人用手藝,心思,所雕刻的金,銀,石.世人蒙昧無知的時候,神並不鋻察,如今卻吩咐各處的人都要悔改.因為祂已經定了日子,要藉著祂所設立的人,按公義審判天下；並且叫祂從死裡復活,給萬人作可信的憑據.」這篇道影響甚大,鄰近的地區如路司得,以弗所和馬其頓的人都深受影響,例如有人藉著造神盦和廟宇發財的就因此對保羅恨之入骨,認為他斷人財路.

這篇充滿神學內涵的演說,我今天不能在此作仔細的分析,只能指出一點,那就是這是向城市人說的,因為其理據和一般神學篇章不同,一般都會從神的創造談起,但這篇講章則強調人的結局,亦即是從神的審判談起.明顯的,這是針對城市文化和人的罪行而展開的神學反思.

不過,我認為保羅在雅典的最大挑戰,不是學者們的問難,保羅自己有著豐富的學識,及良好的辯論口才,舌戰群儒之下,表現絕不遜色；也不是滿城的偶像,和眾多的閑說(gossip),清者自清,保羅及其跟隨者,包括青年人提摩太,在雅典城這個大染缸裡,就如一度濁市清流.那麼,對保羅來說,雅典城內最大的挑戰又是甚麼?

是物欲麼?是權力麼?是學問麼?是虛無麼?

不,都不是.城市中最大的試誘不是這些,而是自己.

使徒行傳十七章提及保羅在雅典以及鄰近城市中繼續為主工作,宣講復活的信息,並且治病趕鬼,得到一些人的跟隨.但也有一些人把他看成希臘傳統中的神靈再世,要跟隨他,把他看成宙斯或亞底米,換言之,他面對一群奉他為神的宗教盲目跟隨者.這是一個被偶像化的過程.我認為這也是城市文化中最能令人迷惑和使人陷落的事情.也是保羅最大的挑戰.

四,結語 and \#8212;-看,那城市和牧場上無牧的羊總括而言,保羅在雅典城中遇到的幾個挑戰,仍然值得近日城市教會參考:它們是:知識化或學派化的挑戰閒話化或gossip化的挑戰偶像化及被偶像化的挑戰最後,讓我們回到耶穌升天前的大使命作為今日的總結:你們要往普天下去,由耶路撒冷,撒瑪利亞,到猶太全地,直到地極.這地極,以今日看來,不在曠野,不在荒原,不在邊疆,不在深山,不在無人之地,乃在城市之中.都市牧羊,應該是二十一世紀的特色,也是末世的挑戰.



\newpage

\section{真葡萄樹與紅酒的故事}
\label{sec:105}
\textbf{伍渭文牧師 - 真葡萄樹與紅酒的故事}
\newline
\newline
\hyperref[sec:104]{< < < PREV SERMON < < <}
~
\hyperref[sec:toc]{[返目錄]}
~
\hyperref[sec:106]{> > > NEXT SERMON > > >}


約翰福音 15:1-5
\newline
\begin{longtable}{cl}
\hline
\hline
章節 & 經文 (和合本修訂版)\\
\hline
15:1 & \begin{tabularx}{0.7\textwidth}{X} 「我就是真葡萄樹，我父是栽培的人。 \end{tabularx} \\ \\ \relax
15:2 & \begin{tabularx}{0.7\textwidth}{X} 凡屬我不結果子的枝子，他就剪掉；凡結果子的，他就修剪乾淨，使枝子結果子更多。 \end{tabularx} \\ \\ \relax
15:3 & \begin{tabularx}{0.7\textwidth}{X} 現在你們因我講給你們的道已經潔淨了。 \end{tabularx} \\ \\ \relax
15:4 & \begin{tabularx}{0.7\textwidth}{X} 你們要常在我裡面，我也常在你們裡面。枝子若不常在葡萄樹上，自己就不能結果子；你們若不常在我裡面，也是這樣。 \end{tabularx} \\ \\ \relax
15:5 & \begin{tabularx}{0.7\textwidth}{X} 我就是葡萄樹，你們是枝子。常在我裡面的，我也常在他裡面，這人就多結果子，因為離了我，你們就不能做甚麼。 \end{tabularx} \\ \\
[1ex]
\hline
\hline
\end{longtable}
日期:2011年7月3日

題目:真葡萄樹與紅酒的故事

經文:約翰福音15:1-5

證道:伍渭文牧師

去過雲南的人大都曉得雲南紅這品牌的紅酒,和有關這紅酒的傳說.「雲南紅是個絕色美女:出生在法國,有高貴的血統,又在雲南紅土地上孕育出了具有雲南民族文化特徵的美貌.」 聽說用來釀造「雲南紅」的葡萄,是一種來自法國的古老品種,名為「玫瑰蜜」(ROSEHONEY),這個品種相傳是一百多年前由法國傳教士引入雲南,目前在法國本土因蟲禍已絕跡.為了配合這歐洲風情概念,有關酒莊旅遊服務有限公司,更安排基督教堂舉行西式婚禮,請牧師親自主持,見證愛情的神聖.他們又說:亞當和夏娃因偷吃禁果,被趕出伊甸園赤身來到荒漠大地上時,他們唯一的所有便是亞當初遇夏娃時,纏繞在她髮間的青色葡萄籐.

夏娃間髮有青色葡萄籐,明顯是附會,但傳教士到雲南傳道時帶來葡萄樹的種子,應該是事實.但傳教士帶來葡萄樹的良種,不是因為他們喜歡飲葡萄酒,乃因為葡萄酒是用作聖餐的聖血,傳教士帶同葡萄樹良種,因為他們憧憬有一天可以和這些山區少數民族一起領聖餐,他們想到基督的救贖寶血.當耶穌說:我是真葡萄樹時,也是在一個聖餐的場景.

耶穌被賣那一夜,和門徒同進最後的晚餐,席中拿起餅來,祝謝,就擘開,說這是我的身體,為你們捨的；飯後,也照樣拿起盛載葡萄酒的杯來,祝謝,就遞給門徒,說:這是我的寶血,為你們流的.跟著耶穌和門徒有很長的講論.約翰福音十三章談及最大的誡命 and \#8212;彼此相愛,叫人認識我們是主的門徒；十四章應許聖靈帶領我們認識真理；十五章論到主是葡萄樹,我們是枝子,與主結連,就能為主結果子；十六章論到聖靈的伴隨,不要憂愁,聖靈將會把真理教導門徒；十七章為大祭司的禱告,我怎樣被父差遣,門徒也要被耶穌差遣.

貫穿幾章離世前講論有一

主題:聖靈與門徒同在,被主差遣,為主結果子.這些離世前的講論,沒有半點離情別緒的宣洩,而是天國藍圖的宣示.今天是聖靈降臨第三個主日,我們用一些時間,思念教會一些宣教先賢,讓我們見賢思齊.

普世的福音首先,為何說耶穌是真的真葡萄樹呢?難道葡萄樹有真和假的嗎?一路以來,聖經都以上帝的選民為葡萄樹:「萬軍之耶和華的葡萄園,就是以色列家.」(賽五7)但這葡萄園結出卻是野葡萄,令主人失望.耶穌說:「我是真葡萄樹,你們是枝子」,其實是新時代的宣告:所有與耶穌結連的人,不分任何種族,都是上帝的選民；每一位分享主聖血的,喝這葡萄酒的,都成為真葡萄樹的枝子,被差遣結果子.

使人得到尊嚴榮美的福音葡萄樹本身不需要很多種植面積也可以生長的植物,因它向上生長；繞著一枝柱子,一面牆,或一道籬笆就可以枝葉茂密.所以在雲南的山區,葡萄樹隨處可見,因為氣候適合.很多不起眼的籬笆和欄杆,都以葡萄枝葉來美化,不好看的場景,也利用豐茂的葡萄樹濃蔭遮擋.推介紅酒的企業,說夏娃髮間有青色葡萄籐,而葡萄酒就來自這葡萄籐,誠然是浪漫的神話,但葡萄籐給不起眼得到冠冕,使醜陋的得到美麗,就如福音使被擄的得釋放,被囚的得到自由.

怒江僳僳族傳道的富懷仁(James Fraser, 1886-1938)中國沿海都會林立,繁榮富裕；內河的交匯,民豐物阜.但深入雲南貴州山區邊陲荒野,就要長途跋涉,遠離文明,營養缺欠,體力消耗,加上蠻巒瘴氣,蛇蟲蚊蠅滋咬,容易生病.到怒江僳僳族傳道的富懷仁(James Fraser, 1886-1938) 就曾患瘧疾,傷寒,腸胃病,最後死於腦性瘧疾,終年五十八歲.他為了福音的原故,深入山區,尋找這些原本文化落後,生活困苦,被人遣忘的少數民族,為他們創製文字,繙譯聖經,穿上衣服,學曉種植葡萄和各種農作物,使他們生活得到改善,文化得到提升.雲南少數民族多信奉基督信仰,七種少數民族傳教創製的文字, 除東僳僳文和載瓦文未能推廣外,其餘基本沿用,或者在原有基礎上修改後繼續使用,教堂就是學校.今天我們若到昆明少數民族文化村,我們就看到這些信仰基督的少數民族,穿著美麗民族服飾來崇拜,載歌載舞,每天幾場崇拜聚會,成為吸引遊客的項目.

貴州烏蒙山區傳道的柏格理 (Samuel Pollard, 1864-1915)這些宣教士的感人事蹟,不單使信仰基督的少數民族畢生感戴,也感動不少藝術工作者.2009年底,曾有藝術家「以靈命愛中華--當代藝術巡迴展」,紀念宣教士富能仁和柏格理,分別於北京,上海,深圳與英國等地八個月依次展出.有商人信徒出資,讓藝術家先到這些傳教工作過的地方訪問參觀,得到感動後為展覽創作作品.

柏格理創辦了近百所學校,培養了上萬名苗,彝族學生,為烏蒙山區的苗族人創製了苗文,並譯出苗文《聖經》,創建了足球場,游泳池,醫院,教堂,孤兒院,麻瘋病院,十年後,柏格理為救治患傷寒的石門苗,彝,漢族學生,被感染病逝在那片荒涼,貧寒的高原上,年僅五十一歲.

胡錦濤在貴州任省委書記時,1985年曾向貴州的幹部倡導學習柏格理的奉獻精神.他說:「公元1904年,一個名叫柏格理的英國人,來到了貴州畢節地區威寧縣的一個名叫石門坎的小村,那是一個非常貧窮荒涼的地方.他帶來了投資,就在這塊土地上蓋起了學校,修起了足球場,還建起了男女分泳的游泳池.他用英文字母仿拼當地的老苗文,自編了「我是中國人,我愛中國」這樣的教材,免費招收貧困的學生.後來一場瘟疫,當地的百姓都逃走了,他卻留下了,呵護他所愛的中國學生.最後,瘟疫奪走了他的生命.柏格理去了,在中國一個荒涼的小村裡,留下了他的一個墳墓,留下了他培育出來的一代中華精英.有人統計這裡出過三個博士,培養出中共廳級幹部二十名.他傳播了知識和西方文化,留下了奉獻和敬業精神.至今這個小村老人們,儘管不識字,居然能說上幾句英語.柏格理用實踐告訴人們:進步的科學文化和艱苦創業,可以在貧困的落後地區,實現教育的超常規科學發展.」

葡萄酒的香氣吸引著福音的使者今天我們看到在少數民族落聚的山區,一座一座的教堂屹立,伴隨著這些教堂,也有一個一的山墳 and \#8212;傳教士和他們的子女的山墳.這些傳教士好像五旬節當天被新酒灌醉,遠離文明,舒適,富裕的都會,深入不毛.也許他們被葡萄酒的香氣所吸引著,他們攜帶著葡萄樹的種子,憧憬著有一天與這些被人迴忘,生活艱苦的山區少數民族,一起共祝聖餐,分享基督的身體和寶血.他們也知道,祇要這些不為人所知,為人輕看的山野之民,祇要連結基督,就是真葡樹的枝子,也可以為主結果子.他們明白,作為基督的葡萄樹枝子,要給貧乏的人,被人輕看的人,看來並不起眼的人,加上冠冕,使他們看來美麗.

在今天,我們的召喚也許再不是雲貴山區的少數民族.兩個星期前的聖靈降臨節主日,院長提到城市牧人,當今的地極就是我們居住的城市,在城市中,有很多人被遺忘,很多人看來並不起眼,但作為真葡萄樹的枝子,我們要給他們冠冕知美麗,讓他們因信仰帶來盼望和尊嚴.

亞洲孤兒成為德國副總理在最近一期亞州週刊,記載了一個感人的報導,主角是現任德國副總理Philipp Roseler.令人驚訝的不是Philipp以三十八歲之齡,掌管重要的經濟事務,令我們驚訝的是Philipp本是一位越戰時被棄的孤兒,沒有名字,被一對當時年青德國夫婦收養.也許這對年青夫婦,像其他德國知識份子,也走上街頭示反越戰,但他們想為戰爭的受害者多做一些具體的工作,收容一位被遺棄的孤兒,給他冠冕,令他美麗,雖然他們已有兩位可愛的女兒.Philipp的義父,是一位直升機機師訓練員,有一次他有機會升職,但為了更多時間陪同兒子而放棄升職的機會.年幼時,義父都讓Philipp都回到自己的辦公室做功課,盡量爭取親子相聚的時間,因他後來與妻子離異,成為單親家長.因為Mr. Roesler的愛護和悉心裁培,一個亞洲的孤兒,竟成為歐盟最強大國家第二位人物.

我們給予冠冕的人,使他更美麗的人,也許不能成為名人,要人,但一定能夠成為基督的枝子,為上主結果子的枝子,基督的身體,聖靈的殿.真葡萄樹與紅酒的故事,不單發生在富懷仁和柏格理這些宣教偉人的身上,也要發生所有願意成為城市牧人的身上.我們,也可以有自己真葡萄樹與紅酒的故事.阿們.



\newpage

\section{在愛與真理中}
\label{sec:106}
\textbf{戴國璋牧師 - 在愛與真理中}
\newline
\newline
\hyperref[sec:105]{< < < PREV SERMON < < <}
~
\hyperref[sec:toc]{[返目錄]}
~
\hyperref[sec:107]{> > > NEXT SERMON > > >}


約翰福音 14:21-24
\newline
\begin{longtable}{cl}
\hline
\hline
章節 & 經文 (和合本修訂版)\\
\hline
14:21 & \begin{tabularx}{0.7\textwidth}{X} 有了我的命令而又遵守的人，就是愛我的；愛我的人，我父要愛他，我也要愛他，並且要親自向他顯現。」 \end{tabularx} \\ \\ \relax
14:22 & \begin{tabularx}{0.7\textwidth}{X} 猶大（不是加略人猶大）問耶穌：「主啊，為甚麼親自向我們顯現，而不向世人顯現呢？」 \end{tabularx} \\ \\ \relax
14:23 & \begin{tabularx}{0.7\textwidth}{X} 耶穌回答他說：「凡愛我的人就會遵守我的道，我父也會愛他，並且我們要到他那裡去，與他同住。 \end{tabularx} \\ \\ \relax
14:24 & \begin{tabularx}{0.7\textwidth}{X} 不愛我的人就不遵守我的道。你們所聽見的道不是我的，而是差我來之父的。 \end{tabularx} \\ \\
[1ex]
\hline
\hline
\end{longtable}
日期:2011年7月10日

題目:在愛與真理中

證道:戴國璋牧師

經文:約翰福音14:21-24

謝謝上帝能夠讓我再一次在這個地方跟大家見面.

今天我們在真理上的學習,是以福音書裡頭耶穌的說話作為我們思考的對象.經文記載在約翰福音第14章,我想大家跟我一起打開這段

經文:

21 有了我的命令又遵守的,這人就是愛我的．愛我的必蒙我父愛他,我也要愛他,並且要向他顯現. 22 猶大問耶穌說,主阿,為甚麼要向我們顯現,不向世人顯現呢. 23 耶穌回答說,人若愛我,就必遵守我的道．我父也必愛他,並且我們要到他那裡去,與他同住. 24 不愛我的人就不遵守我的道．你們所聽見的道不是我的,乃是差我來之父的道.(約翰福音14:21-24)

在耶穌短短的話語中,發現到有三件事情:第一,耶穌談到愛祂就是愛上帝；第二,談到守道,即是要遵守祂的命令；第三,談到如果我們能夠愛祂和遵守祂的道的時候,祂要給予我們賞賜.

我想從這段經文裡跟大家來看看愛主,守道,與得賞賜彼此之間的關係.

首先,我們看耶穌基督談到賞賜的事──在這段談及我們要在真理裡彼此持守,在愛中學習的經文中,耶穌以賞賜作為鼓勵,讓我們在這門功課裡做得更加美好.從耶穌基督的話裡我們看到一些事情,這些事情也是我們基督徒生命裡常常經歷到的事,就是我們發現我們的上帝是一位樂意賞賜的上帝.這是祂的本質,所謂賞賜就是要賜予禮物,上帝本來就是一件樂意賞賜,樂意賜予禮物的一位,我們知道這些的賞賜有不同的形式,譬如可以是物質的,但也可以是指著各種屬靈的福氣.

從前,五穀豐收,生產豐盛是猶太人所羨慕的,這個是屬於物質的賞賜.到現今,我們在基督裡的乃渴求賜予永生和承受天國,這個我們知道是屬靈福氣上的賞賜.不過,我們發現賞賜也可以是非條件性,也就是我們所說的「白白的」.例如上帝賜給我們有救恩,這是白白的,非條件性.但賞賜也可以具有條件性,取決於我們對上帝召命的回應.當我們對上帝特殊的感召,吩咐給予正面的回應,順從的時候,我們就可以得著這些賞賜,不過無論這些賞賜的性質是屬於物質的,還是屬靈的福氣,它是屬於非條件性的,或是有條件性的,這些賞賜其實也是表達了上帝豐富和慷慨.聖經上有許多的應許是告訴我們這乃是上帝的本質.例如詩篇84:11:「因為耶和華　神是日頭,是盾牌,要賜下恩惠和榮耀．他未嘗留下一樣好處,不給那些行動正直的人」.先知賽亞也提及過:「看哪,耶和華曾宣告到地極,對錫安的居民說,你的拯救者來到．他的賞賜在他那裡,他的報應在他面前」(以賽亞書 62:11).這裡所指的報應其實就是指著,當我們回應,追求上帝獎賞的時候,祂會以美麗的事物向我們回報,所指的就是美麗的事.

新約聖經也有類似的經文鼓勵我們,希伯來書11: 6:「人非有信,就不能得上帝的喜悅．因為到　神面前來的人,必須信有神,且信他賞賜那尋求他的人.」我們對上帝的信心其中最強烈的表現回應之一,就是我們相信上帝是一個賞賜的上帝.以色列的始祖亞伯拉罕所體驗到的上帝是「我是你的盾牌,必大大的賞賜你」(創世記15:1).許許多多的舊約,新約的經文都告訴我們上帝是一位樂意賞賜的上帝,這也是我們對上帝的第一個經驗.故此,我們也渴求上帝給我們的賞賜以作為我們對上帝樂意賞賜的第一種回應.其實在我們生命不同的歷程裡,上帝也會以不同的感召來呼喚我們,這些感召可能是召喚我們去預備得著祂為我們所預備的終極救恩,可能是召喚我們在我們生命裡某些階段裡的靈命成長和更生.假若我們聽從上帝這個召喚,這樣我們就會進入一個順從的反應當中.那麼我們就會進入一連串的彼此的互動裡頭──action and reaction(行動與反應).在這些關係裡,我們往往可以享受到上帝的福氣是這麼多,這麼深.正如我們剛才讀到耶穌基督的話語,假若我們能夠回應更愛祂,愛上帝的話,上帝應許告訴我們,祂必然給予我們更多和更深的回報.耶穌提醒我們說我們將會得到父和子的愛,我們也會得到祂向我們顯現和與祂同住的應許.為什麼耶穌基督會在這裡給我們這些賞賜的應許呢?就是因為上帝了解到一件事,賞賜其實是我們成功學習屬靈功課裡的一個相當重要的因素或手段.

賞賜可以是教學過程裡一個很重要的手段.我們知道在所謂教育心理學(就是從心理的角度來研究怎樣才可以幫助學生在學習的過程中更成功),我們發現在教育心理學上,教導我們該怎樣去了解才能夠促進我們在學習上的果效,很多時候都可以借用一些的手段,而賞賜就是其中一個最好的方法.在早期的心理學的研究上,就曾經有過一些的發現,譬如在行為主義學的心理學上,從他們的研究裡就發現說,如果一個行為伴隨著一個令人滿意的情境時──譬如當一個動物,牠做了一個恰當的行為後,我們就給予牠一些食物作為獎賞──當再次遇到類似的情境的時候,重複這種可能性的行為就會增加.相反,如果在這個學習的過程中所伴隨著的是不滿足的行為,例如做完一件事後,被試驗者是受到懲罰(譬如是經歷一些疼痛的話),那再次重複這種行為的機會就會相對的減少.在心理學上則稱之為「效果率」──賞賜或懲罰對某方面的學習過程是有幫助的.

另一個比較著名的例子,就是所謂「巴普洛夫的狗」的故事.巴普洛夫其實本來是俄國的一位研究消化系統的著名生理學家,但偶然地在他的研究的過程當中,發現到一些特別的東西,於是他就將他的研究轉向改為研究心理學的條件反射,如果聽過這個故事的,就會更加清楚知道我在說些什麼.不過我在這裡不打算重覆地說有關巴普洛夫的條件反射的著名故事,反而想借著另一個事例來說明有關學習裡怎樣給予回報的重要性.

這個例子據說是有關研究兒童對攻擊性行為模仿的一個經典試驗.這個試驗是在二十世紀六十年代的時候,由一個著名的學習心理學家所設計出來的,那個方法就是用三群小朋友,或者是將一群小朋友們分為三組,然後安排他們觀看同一段的影片,影片所描述的是一個人在毆打一個不倒翁的玩偶,當然在毆打的過程當中,這個人還會咆哮大叫地責罵那個不倒翁玩偶.三組兒童所觀看的影片基本上是一模一樣的.唯一不同的就是他們所看到的結局.譬如第一組兒童所看到的結局是這個毆打不倒翁玩偶的人是會得到人的讚賞和賞賜；第二組兒童所看到那個結局就是那個毆打不倒翁玩偶的人會受到責罵和懲罰；而第三組兒童所見到的結局就是既沒有責罵,也沒有賞賜.三組兒童在觀看影帶之後,就進入另一個執行的階段,就是讓這三組兒童一個一個的與這不倒翁玩偶獨處,然後觀察他們會怎樣對待這個不倒翁玩偶.他們觀察的時候發現到一些令人頗感意外的表現:他們發現第一組小朋友們(就是影帶主角在毆打玩偶後會得到人的獎賞)會傾向對那個不倒翁玩偶做出同樣地的拳打腳踢和囂喝；至於第二組小朋友(主角在毆打玩偶後會受到懲罰),他們就會傾向於比較靜止,反應沒有那麼強烈；而第三個組別的小朋友和第二組別的小朋友的表現一樣,比較不去模仿那種行為.在第三個階段,就是在所有的小朋友面前讚賞第一組的小朋友(就是那些會對不倒翁公仔拳打腳踢的小朋友),他們發現這些小朋友竟然全部受到影響,無論是第一,第二,還是第三組的小朋友都會起來對這個不倒翁公仔進行拳打腳踢和囂罵.這個結果得出一個什麼樣的結論呢?這個結論就是,如果一種行為得到人的獎賞,那些行為是會在人群當中得以強化,即使那種行為是惡的,但假使得不到任何的懲罰或者反面糾正的話,小朋友也會傾向去做這些事或者這些行為.如果將這個結論帶到現實的生活裡,譬如電視節目的內容裡往往有一些壞蛋的出現,但如果這些壞蛋的結局是受到懲罰的話,小朋友往往就會明白到這些壞蛋的行為是不應該學習的,但是假使這壞蛋的行為沒有得到任何的懲罰,他們就會很擔心小朋友其實是會學了,然後將這些行為宣洩在其他人的身上.這個例子其實是一個很著名的例子讓我們看到,原來在一個學習的過程當中,該怎樣去鼓勵一個人,或者是怎樣防止一個人的學習.這就使我們明白到為什麼我們的救主耶穌在約翰福音第14章裡講到「愛主」與「守道」的時候,提到有關賞賜的問題,其實祂是鼓勵我們,提醒我們,希望我們能夠在愛上帝,守真道的事情上有更多的盡心,更多的成長.

繼而我想跟大家看看這段經文的「愛主」和「守道」的關係.究竟愛上帝和守真理有什麼關係呢?我們發現到愛上帝與守真理或守誡命都連接在我們要作順命的兒女的吩咐上面.我們知道耶穌基督曾經用過一個猶太人的「示瑪」(Shema)──「盡心,盡性,盡力愛耶和華你的上帝」──來提醒我們要將愛上帝視為誡命當中的第一條,如此,就將愛上帝和守誡命聯接在一塊.我們知道藉著守誡命來表達愛上帝的這份心意是猶太人的傳統,我們也覺得沒有問題,只要我們記得愛上帝是根本,守誡命是實踐的手段.我們知道後來猶太人的遺傳出了毛病,就是因為他們把這兩個的次序倒轉了,他們視守誡命為根本,愛上帝就被消化了,消失了,後來慢慢地變成了一個僵化,死守遺傳的一個民族.其實我們基督徒也應該要小心,因為基督徒同樣地也會犯這個毛病.很多時候我們基督徒會將耶穌基督的福音,神人和好的消息變得「誡命化」,使每一個基督徒都要怎樣屬靈,虔敬,熱心,而忽略了最重要的乃是要以單純的心來接受「上帝的愛」這個事實.

另一方面,我們基督徒又可以犯另一種完全相反的毛病,就是我們明白並且接納上帝在耶穌基督裡無條件的愛,但我們卻以為我們不需要再遵守誡命,因為這些都是律法.但耶穌的話清楚地讓我們看見並不是的,兩者之間其實有一定的關係.在21節裡,耶穌說:「有了我的命令又遵守的人,這人就是愛我的」.換句話來說,只有那些肯遵守上帝命令的人,才有資格去講我們對上帝的愛.其實愛上帝與守上帝的命令彼此之間並不是排拒的.若果我們能夠明白到將愛或誡命視為愛的要求,我們就不會覺得誡命是一些行為,誡命只是「愛的要求」.

在愛裡有要求是一件很自然的事.昨晚我們教會的一班夫婦團契的弟兄姊妹,到一位執事家中做家訪,其中,我們做了一個這樣的活動:將他們分為兩組,弟兄一組,姊妹一組,當中有許多夫婦,大家要將對配偶的期望列出來,像是希望配偶能做些什麼來互相配搭,好達到自己的一些期望.列出後就在這十多項的期望裡頭,選出五個你覺得最重要的來與大家分享.弟兄跟弟兄分享,姊妹跟姊妹分享,分享完畢後,大家再坐在一起,弟兄們就分享哪五個期望是弟兄最希望配偶能配搭的,姊妹也選五個.然後我們發現到一件奇怪的事,原來弟兄所列舉的五個跟姊妹所列舉的五個都頗吻合,當然個別的人在次序裡的重要性有少許的不同,但總體上表明了一件事,原來我們作弟兄的希望配偶給我們的和姊妹希望弟兄給她們的其實都很一致.這同時印證了耶穌基督的話:「

你們願意人怎樣待你們,你們也要怎樣待人」(路6:31)──我們希望我們的配偶怎樣對待我們,原來她們也同樣地期望我們能夠怎樣去對待她們.這些期望是因為我們覺得真正的愛,真正的關心是帶著期望的.如果我們明白其實誡命也是一些的期望,當上帝要我們守這些誡命的時候,其實是因為祂愛我們.

其實愛主,守誡命兩者之間並無矛盾,兩者之間可聯繫在一起的.上帝愛我們,祂盼望我們能夠順命；我們愛上帝,故此我們樂意去作,行祂所喜悅的事.正如耶穌基督自己在約翰福音8:29說:「

那差我來的,是與我同在．他沒有撇下我獨自在這裡,因為我常作他所喜悅的事.」祂能夠回應「我常作父所喜悅的事」,就是叫使徒明白這個道理.所以使徒給了我們同樣的吩咐,約翰說:「

因為我們遵守他的命令,行他所喜悅的事.」(約翰壹書3:22)如果我們愛上帝,我們就樂意做上帝所喜悅的事.

我想起一次經歷,有一天,我下班回家後坐在沙發上很優閒地看電視,太太因為約了丈母娘逛街,晚了點回來,當她回來的時候,我覺得她神情有點古怪,為什麼呢?因為她在我面前走來走去,又好像在偷笑,我心想應該發生了些甚麼事,就問她到底發生了什麼事?怎麼你看想來好像很開心?她就展示她手上的一枚戒指,我問她怎麼會多了一枚戒指呢?她說是她媽媽買給她的.我說為什麼丈母娘突然之間會買一枚小戒指送給妳呢?她說因為她愛媽媽,媽媽知道而且也愛她,所以就送她一枚戒指作為賞賜的回應.若我們愛惜上帝的時候,我們同樣樂意地去做上帝所喜悅的事,耶穌在十字架上的犧牲對神,對人,也以此原則來行事.

當耶穌說到愛祂的人,祂就給予賞賜,那賞賜是什麼呢?在21節,耶穌這樣說:「

有了我的命令又遵守的,這人就是愛我的．愛我的必蒙我父愛他,我也要愛他」.當我們切實地去愛上帝的時候,我們必然得到愛的回應和回報,絕不石沉大海般沒有消息,我們不單得到必然的回報,耶穌的應許是說:「我父要愛他,我也要愛他」,耶穌講完後沒有停在那裡,他還說:「並且要向他顯現」.

我們相信這個顯現的意思,當然不是指耶穌基督以祂榮耀的本體來向我們顯現,乃是指我們的生命裡常得到上帝恩典的顯現,給我們許許多多的恩惠.但再進一步地,不單是顯現,耶穌更談到「

並且我們要到他那裡去,與他同住.」(14:23)我們相信顯現是一個非常美好的回應,但同住是一個更深,更美好的回應.如果顯現是代表著耶穌基督間歇性的到訪,這樣同住就是一個長久的在一起.耶穌用這個方法來鼓勵我們,提醒我們當我們愛祂,願意守祂誡命的時候,我們會得著祂各種奇妙的恩典,這是何等的美好呢?

我們細細地想一想,如果耶穌基督真的給我們同住的應許,常與我們在一起的時候,這是一個何等嚴肅的事情呢?我們曾經記得當上帝降臨在西乃山的時候,當時的以色列人是何等的恐懼.當上帝要臨到的時候,他們心裡充滿各樣的恐懼.同樣的,當上帝與我們同住的時候,會否喚醒我們要在生命當中儆醒,如果我們能夠有「上帝同在」的這個的福氣,我們就要在我們生命上用聖潔,儆醒來迎接上帝的來到.

申命記23:12-14,記載了一段我覺得頗特別的一個故事(其實也是一個吩咐).當時的以色列人在曠野漂流,所以他們是住帳棚安營的.在申命記的23章,有一個有趣的吩咐,當以色列人出外征戰打仗的時候,祂要求每一個以色列人在隨身的器械裡要預備一把鍬,也就是一把小小的鏟,用以鏟土,為何每人都要預備一把鍬來鏟土呢?後面就有解釋,當晚上安營的時候,有時候突然之間肚子痛,要起來便溺,那當然不能在營中,營是人居住的地方,所以要在營外有一個便所,但是當你便溺之後,要用這個鍬鏟土,轉身掩蓋.我們知道這個純粹是為了衛生的原因,應該是這樣的,下面就有一個理由,為什麼以色列人要這樣做呢?「

因為耶和華你的　神常在你營中行走,要救護你,將仇敵交給你,所以你的營理當聖潔,免得他見你那裡有污穢,就離開你.」這我們就明白了,因為如果上帝在我們中間陪伴著我們,我們的生活是否應當要聖潔呢?令我感到特別的地方就是這個「聖潔的要求」不單是「屬靈的聖潔」,竟然涉及到有關「衛生方面」的條件,這個是要以色列人在他們平日的生活當中也當重視他們的衛生,如果物質的污穢尚且如此的時候,那麼我們屬靈的生命更加要有何等的聖潔呢?所以當我們念到耶穌基督的這番話的時候,我們真的要了解上帝的苦心,祂提醒我們真正愛上帝的人是不會將上帝的誡命丟在一邊,而不加以理睬,我們反而是要很樂意地去遵行,因為我們願意做祂所喜悅的事.我們相信當我們有這種正面回應的時候,上帝就會有一個更正面的回應,就是給予我們生命裡各樣美麗的賞賜,以致於我們能夠有更多的成長.

讓我們同心來祈禱:我們的主,我們謝謝你!因為你在這段經文當中向我們闡述了有關「愛主」,「守道」和「得賞賜」的事情,讓我們看見原來真理在愛上帝的人的生命當中是這麼的重要.主啊!我求你讓我們明白,也求你激勵我們有一個正面回應的心.我們很軟弱,我們對真理的渴慕常常不夠,當我們立志要去行的時候,我們又受著我們肉體許多的困擾,以致我們常常不能成功,我們唯有祈求上帝繼續在我們生命當中不斷地提醒,不斷地鼓勵,好叫我們每天有進步,謝謝你,奉耶穌基督的名字,誠心所願.



\newpage

\section{在愛中建立自己}
\label{sec:107}
\textbf{陳以諾牧師 - 在愛中建立自己}
\newline
\newline
\hyperref[sec:106]{< < < PREV SERMON < < <}
~
\hyperref[sec:toc]{[返目錄]}
~
\hyperref[sec:108]{> > > NEXT SERMON > > >}


以弗所書 4:15-16
\newline
\begin{longtable}{cl}
\hline
\hline
章節 & 經文 (和合本修訂版)\\
\hline
4:15 & \begin{tabularx}{0.7\textwidth}{X} 我們反而要用愛心說誠實話，各方面向著基督長進，連於元首基督， \end{tabularx} \\ \\ \relax
4:16 & \begin{tabularx}{0.7\textwidth}{X} 靠著他全身都連接得緊湊，百節各按各職，照著各體的功用彼此相助，使身體漸漸增長，在愛中建立自己。 \end{tabularx} \\ \\
[1ex]
\hline
\hline
\end{longtable}
日期:2011年7月17日

題目:在愛中建立自己

證道: 陳以諾牧師

讀經:以弗所書4:15-16

15 惟用愛心說誠實話,凡事長進,連於元首基督． 16 全身都靠他聯絡得合式,百節各按各職,照著各體的功用,彼此相助,便叫身體漸漸增長,在愛中建立自己.

各位親愛的弟兄姊妹平安!非常謝謝伍渭文牧師的邀請,讓我有機會在這裡與大家一同分享神寶貴的話語.我們感謝上帝的恩典和帶領,讓我們有這一次的機會來到祂自己的面前敬拜和事奉祂,祂也保守我們廣州教會的二十位教牧同工,能有機會來到崇基學院神學院分享和學習,這是一個非常難得的機會,雖然我們兩地的教會在文化和處境上都各有同,但我們藉著上帝的愛能夠走在一起,因為神的愛是不分彼此,沒有隔閡,能夠彼此認識的,原本我們是素不認識的,但藉著神的愛讓我們能夠走在一起,能夠彼此認識,也能互相交流,學習,並在愛裡彼此建立.

聖經告訴我們「在愛中的建立」可以有三方面的層次:

第一──就是我們每一位弟兄姊妹,每一位神的兒女都可以與上帝同工合力.因為上帝是我們每一位屬祂兒女的學習榜樣,我們可以透過我們與神的同工合力,從中學習到神要我們所學習的功課.因為我們的耶穌基督是教會的元首,是教會的根基,而教會是基督的身體,所以我們應當與主聯絡得合式,這種的關係是什麼樣的關係呢?就是我們這屬神兒女的與阿爸父彼此同工合力的那種親密關係,在救贖人類的事情上,我們都很清楚耶穌一直都是這樣的與神同工.

耶穌在世上傳道的時候,祂揀選門徒,與門徒一同傳福音,這個就是人與上帝彼此同工合力的一個最好的表現.所以我們看到在救贖人類的計劃裡面,我們的上帝是希望我們每位屬神的兒女,是不分彼此的,都能參與這一場救贖靈魂的計劃,能與我們的神彼此同工,成就上帝在我們身上偉大的聖工.所以在福音的事工上,上帝不是沒有能力去擔負起這個救贖的工作,才需要我們與祂同工,而是上帝給我們有選擇的自由意志,為我們樹立了一個榜樣,使我們懂得怎樣與祂同工,與祂同行,去完成上帝交付給我們的使命.

在這裡我要講一幅的圖畫,一幅叫做【腳印】的名畫.這幅【腳印】有著非常豐富的屬靈意思,說明了我們在許多順境的時候,覺得主耶穌是和我們一起行走的,因為地上是有兩對平行的腳印,但我們很快的就會發現只剩下一對的腳印,這是什麼意思呢?在我們人生的道路中,有時候是從晴天轉為陰天,這時候我們就會發怨言,以為主耶穌離開了我們,但是主耶穌並沒有離開,祂愛我們就愛我們到底,雖然看上去好像是只有一對的腳印,因為神的愛,祂在我們最困難的時候,背著我們繼續向前走,所以地上只有一對的腳印.

弟兄姊妹,我們是要感謝神每天的同在,感謝上帝每天保守我們走天路,感謝上帝對我們的愛,使我們能夠與祂同工,亦能夠與祂同行,並且學習怎樣去愛我們的上帝.當我們與上帝同行,同工的時候,也學習到怎樣與我們的神,去經歷各種的苦與樂,學習怎樣去愛神,愛人.使我們確信無論遇到怎樣困難的事,我們都懂得去倚靠主,倚靠上帝,過一個信心的生活.

在福音的事工上,人的作用是不容忽視的,因為上帝要透過我們人來完成某些事工.但是弟兄姊妹我們應該要保持非常清醒的頭腦,不能夠從此就認為我們是比上帝聰明,比上帝厲害及有本事,因為我知道在福音的事工上,我們離開了上帝我們就不能夠作什麼.

聖經都教導我們──「

不是倚靠勢力,不是倚靠才能,乃是倚靠耶和華的靈,方能成事.」(撒迦利亞書4:6)

對於福音的事工,我們作為神的兒女有我們自己的工作,我們的上帝亦有祂的工作,我們與上帝同工,同行的時候,要怎樣學習呢?要學習順服,倚靠,禱告,交托和等候等等的信心功課.我們不能夠在與上帝同工同行的時候,還去搶奪屬於上帝的工作,把屬於上帝能力範圍內的工作都做了,如果是這樣子的話,我們不單做不好,反而還會將問題複雜化,像廣州話的口頭譂所說:「好心做壞事」,所以我們不要在這個聖工上越界,把上帝的工作都做了,這樣當然辛苦,而且我們都沒有這樣子的能力,不是嗎?

所以弟兄姊妹,我們與上帝同工同行的時候,我們應該怎樣做呢?我們要明確自己的使命,又是怎樣去定位的呢?工作和使命又是怎麼一回事呢?

我們的工作和使命就是撒種,就是傳福音.

上帝的工作又是什麼?就是讓我們所撒的種籽可以生長,開花,結果.

所以我是不同意那些作為父母的,只懂得強迫小朋友要信耶穌,你是強迫不了的,反而會爭吵使氣,離家出走,斷絕父母關係,在教會的聖工上,我們很多時候是需要等待的,當然我們的責任是禱告,但上帝有祂自己的時間,在上帝的時間和預定當中,祂必定成就在我們每一個人身上那種美好的旨意.所以弟兄姊妹,與神同工合力的時候,要知道上帝永遠是我們的上帝,是我們的元首,我們是祂的僕人,乃是上帝所使用的器皿,因此我們要倚靠一位無所不能的主,凡事高舉我們的主耶穌基督,並要謙卑在主的面前,過一個順服主的生活,這樣我們才能夠活出基督耶穌的樣式,能夠為主的名作出美好的見證,這是第一個層面.

第二──我們可以透過人與人之間的那種彼此同工合力達到什麼目的?達到我們在肢體之間的相顧,互助,互愛.因為教會是基督的身體,並且彼此互相為肢體,而各肢體有不同的功用,我們藉著這些肢體之間的同工合力,彼此相助,彼此倚存,如此,我們可以使這個身體,就是我們所說的教會,能夠得以建立起來.

主耶穌起初出來傳道時候,祂首先設立了一個屬靈團契,當十二個使徒出去傳福音的時候,明顯地我們看到他們被分成了十二個小組,他們是乃是兩個,兩個的出去傳福音.

有什麼小組呢?聖經馬太福音講得非常清楚,是有六個小組,門徒有十二個,但是有六個小組,即是兩個,兩個一隊,首先就是彼得和安得烈,雅各和約翰,腓力和巴多羅買,多馬和馬太,亞勒腓的兒子雅各和達太,奮銳黨的西門和加略人猶大.這些都是一對一對的,一共分成了六個小組,當主耶穌基督揀選了他們,把權柄分派給他們,好叫他們有能力去救人的靈魂,也能救人的身體,門徒果然有這樣子的能力,工作也如此的有果效,這並不單單是靠著一個人的能力,乃是他們彼此相助,彼此相愛,彼此同工的結果.

這次非常難得有機會來到崇基神學院進修學習,在課程當中也聽到崇基神學院劉宗鳴教授在上課的時候談到有關「自我牧養」這方面,他是這樣說的:「自我牧養是可以分成兩個層面,一是屬於自己個人層面,而另一個層面是屬於小組群體的層面.他說處理自己的時間,事工,家庭和情緒是一個最基本的問題,也是容易讓人迷失方向,所以我們要懂得怎樣來自我設定界限,懂得知道什麼事可以做?什麼事不應該做?給自己一些界限,一些限定,一些底線,目的就是維持一個在事奉上的支持系統.」我聽後很有感動,也讓我們看見在我們牧養的過程當中,我們其實是需要彼此牧養,當然耶穌基督是我們的牧者,我們每一個人都要接受主的牧養,但我們信徒之間亦離不開一個群體的「支持系統」──彼此的相交,彼此的扶持──離不開小組群體的牧養,因為群體牧養得好的時候,我們在事奉上就能夠更加得力,有效.

在這裡想和大家分享一些廣州的情況,去年在廣州舉辦了第十六屆亞洲的運動會,簡稱【亞運會】.在香港的弟兄姊妹可能也很清楚,舉辦亞運會跟我們有什麼關係呢?因為亞運會有來自世界各地的來賓,而運動員內有些是信耶穌的,當然除了信耶穌亦有來自五大宗教的信徒,所以作為教會或是其他宗教,都要自願性地去提供這個的接待服務.廣州市的基督教堂會不是很多,只有十四間的教堂,但其中有九間的教堂是作為重點的接待堂會,那要找什麼人去作接待呢?當然是找信徒去接待,政府是不會派人來的,但這些接待的人是要經過培訓,英文較好的就是正式的自願者,不懂英文的是後備自願者,但不論是正式或是後備的自願者,他們在接待的工作上都非常配合,他們忘我似的勞動,完全體驗出在主裡面的同心合力.這次亞運會我們基督教總共接待的來賓共有六千四百多人次,接待這些外賓的來訪,自願者發揮了他們積極的作用,當一個自願者忙不過來的時候,其他的自願者就自動的補上了,他們雖然是後備,但他們不但不去計較名與利,甚至也不計較報酬,只積極地參與接待的服務.亞運會有一句口號,其中一句就是:「人人都是自願者」.「人人都是自願者」──不是說你經過培訓才是自願者,總之有需要,你就有責任提供一個優質的服務讓來訪的運動員,外賓有一種賓至如歸的溫暖感覺.我們是教會,讓他們能夠感受到教會弟兄姊妹是何等的美,何等的善.所以亞運會我們做得非常好,受到政府有關部門和來訪客人的高度讚揚,我們亦透過亞運會這個服務為上帝作了美好的見證,因為我們不像香港可以走到街上傳福音,我們有我們的國情,但我們藉著這個服務在社會上產生了巨大的影響,為上帝作了美好的見證,實際上也是間接地去傳福音為基督教增添了光彩.

所以弟兄姊妹一個團隊的精神是很要緊的,教會裡面不是說沒有這個,那個就不行,而是沒有大家都不行.教會復興離不開教會內弟兄姊妹的共同參與,互相的支持,幫助,團結友愛以及發揮每個弟兄姊妹的恩賜,這樣才能成全信徒以建立基督的身體.

第三──我們要透過普世的教會之間的相互合作及交流來實現愛的建立.

經文內談到的「全新」是指什麼呢?乃是指全教會和普世的教會.所以不論什麼地方的教會,在基督裡都是祂的身體,並且有互為肢體的關係,我們是普世教會裡的一份子──共同建立一個普世教會,拓展上帝的國度──也是彼此互為同工的關係,所以我們應該在愛中建立,互相聯絡,互相扶持,共同來推動整體教會的發展.

在聖經裡我們看見昔日耶路撒冷的教會也得到各地教會的幫助和支持,因為昔日的耶路撒冷教會的信徒是面臨著貧窮和飢餓的威脅,所以使徒保羅為他們籌集捐款,他建議信徒們在所得的款項內撥出部份款項交給教會,以便集中在一起等候保羅再來的時候帶給耶路撒冷教會的信徒.後來保羅到了哥林多教會就帶走了信徒籌集的捐款,並且帶到耶路撒冷去發給那裡有需要的信徒.使徒保羅與哥林多教會,腓立比教會,以及以弗所教會都有友好的來往.而這些教會之間亦在主裡面有一種非常友好的交往和彼此分擔,保羅為他們在主裡面有這種同工相助,彼此興旺福音而向神獻上感恩和讚美.所以弟兄姊妹今天我們兩地教會,特別是我們廣州的教會和香港的教會,也應該在主裡面多多交流和合作.

崇基學院神學院葉清華教授和我們上課時談到神學院是教會的後勤部.他說教會就是戰場,所以剛才伍牧師說我們的教牧同工是站在戰鬥場上的第一線,是在戰場上為神工作的.我們這次來的同工都是非常基層的教牧同工,他們真的是默默地殷勤的為主作工.葉教授談到神學院是教會的後勤部──為教會提供彈藥資源的地方；教會是戰場──是衝鋒陷陣的地方.我聽後就比較清楚,原來神學院的功用是這樣子的,它為我們提供彈藥,資源,訊息,靈糧的地方,我們進來進修就是為了這個目的,因此我們要去充實,裝備,好讓我們能更好地牧養信徒,辦好教會.所以在福音事工上,我們都彼此互為同工,不分你我,因為我們都是一家人,同一個戰場上的戰友,所以我們要攜手並進,共同努力去成就上帝偉大的聖工.

感謝上帝的恩典,在祂的保守與帶領下,多年來我們廣州的教會與崇基學院神學院都有很多的交流和合作的機會,這次能夠來到中文大學崇基神學院學習,也讓我今天有機會在這裡講道,這是一個莫大的突破,也足以看到上帝在廣州,香港所作的奇妙恩典和見證.我記得在九七年或以前,崇基神學院每年都會派神學生到廣州教會來實習或交流,在01年起每年崇基神學院更派牧師或教授來廣州教會給教牧同工授課,而且課程內容都是非常豐富的.

感謝神的恩典,今年我們又有一個機會親自來到學院這裡進修學習,所以在這裡非常謝謝神學院為我們預備了這麼好的進修學習的機會,也謝謝各位的牧長,教授給我們講了那麼好的課程,使我們能夠透過這麼短的學習時間,開闊了眼界,提高了認識,並且也增強了我們事奉的信心,因此我們非常盼望日後能夠有更多類似的學習機會.所以各位弟兄姊妹,今後我們兩地有更多交流學習的機會,好使我們在主裡面能夠有更多的聯絡,在資源上能夠互補,發揮各自的優勢,不斷地拓展上帝國度,為基督身體的建立,盡上我們一份的力量,願上帝美好的旨意成就在你我身上.

請我們同心的禱告:慈愛天父,我們滿心歡喜快樂地感謝你!因為你愛我們,揀選了我們,你在我們每個人身上都有你的大使命,就是我們們要往普天下去傳福音,為萬民作見證.親愛的主啊!我們都是神你所揀選,你所喜愛的兒女,今天我們在這個神學院裡接受造就,我們完全將我們每位弟兄姊妹都擺上,主,你既然揀選了我們,就求主你親自來造就我們,使用我們,使我們更能夠得著裝備,更能夠有靈,有力地來事奉你,好讓我們的事奉能夠更討神你自己的喜悅.我們求主你的聖靈來加添我們的能力,充滿我們,使我們都能夠同心合意地去復興神你的教會,傳講神你的福音.主啊!我們感謝讚美主!求主你繼續帶領我們,帶領我們面前的道路,使我們能為主的名打美好的仗,作出更多美好的見證,願上帝你美好的旨意成就在我們各人身上,我們這樣子同心祈禱交托,是奉主耶穌基督的聖名,阿們!



\newpage

\section{無債一身輕}
\label{sec:108}
\textbf{林豪恩先生 - 無債一身輕(Release from Debt)}
\newline
\newline
\hyperref[sec:107]{< < < PREV SERMON < < <}
~
\hyperref[sec:toc]{[返目錄]}
~
\hyperref[sec:109]{> > > NEXT SERMON > > >}


羅馬書 8:12-25
\newline
\begin{longtable}{cl}
\hline
\hline
章節 & 經文 (和合本修訂版)\\
\hline
8:12 & \begin{tabularx}{0.7\textwidth}{X} 弟兄們，這樣看來，我們不是欠肉體的債去順從肉體而活。 \end{tabularx} \\ \\ \relax
8:13 & \begin{tabularx}{0.7\textwidth}{X} 你們若順從肉體活著，必定會死；若靠著聖靈把身體的惡行處死，就必存活。 \end{tabularx} \\ \\ \relax
8:14 & \begin{tabularx}{0.7\textwidth}{X} 因為凡被神的靈引導的都是神的兒子。 \end{tabularx} \\ \\ \relax
8:15 & \begin{tabularx}{0.7\textwidth}{X} 你們所領受的不是奴僕的靈，仍舊害怕；所領受的是兒子名分的靈，因此我們呼叫：「阿爸，父！」 \end{tabularx} \\ \\ \relax
8:16 & \begin{tabularx}{0.7\textwidth}{X} 聖靈自己與我們的靈一同見證我們是神的兒女。 \end{tabularx} \\ \\ \relax
8:17 & \begin{tabularx}{0.7\textwidth}{X} 若是兒女，就是後嗣，是神的後嗣，和基督同作後嗣。如果我們和他一同受苦，是要我們和他一同得榮耀。 \end{tabularx} \\ \\ \relax
8:18 & \begin{tabularx}{0.7\textwidth}{X} 我認為，現在的苦楚，若比起將來要顯示給我們的榮耀，是不足介意的。 \end{tabularx} \\ \\ \relax
8:19 & \begin{tabularx}{0.7\textwidth}{X} 受造之物切望等候神的眾子顯出來。 \end{tabularx} \\ \\ \relax
8:20 & \begin{tabularx}{0.7\textwidth}{X} 因為受造之物屈服在虛空之下，不是自己願意，而是因那使它屈服的叫他如此。 \end{tabularx} \\ \\ \relax
8:21 & \begin{tabularx}{0.7\textwidth}{X} 但受造之物仍然指望從敗壞的轄制下得釋放，得享神兒女榮耀的自由。 \end{tabularx} \\ \\ \relax
8:22 & \begin{tabularx}{0.7\textwidth}{X} 我們知道，一切受造之物一同呻吟，一同忍受陣痛，直到如今。 \end{tabularx} \\ \\ \relax
8:23 & \begin{tabularx}{0.7\textwidth}{X} 不但如此，就是我們這有聖靈作初熟果子的，也是自己內心呻吟，等候得著兒子的名分，就是我們的身體得救贖。 \end{tabularx} \\ \\ \relax
8:24 & \begin{tabularx}{0.7\textwidth}{X} 我們得救是在於盼望；可是看得見的盼望就不是盼望。誰還去盼望他所看得見的呢？ \end{tabularx} \\ \\ \relax
8:25 & \begin{tabularx}{0.7\textwidth}{X} 但我們若盼望那看不見的，我們就耐心等候。 \end{tabularx} \\ \\
[1ex]
\hline
\hline
\end{longtable}
日期:2011年7月24日

題目:無債一身輕(Release from Debt)

經文:羅馬書8:12-25

證道:林豪恩先生

一．引言有一對情侶拍拖多年,決定「拉埋天窗」.他們兩人的收入比上不足,比下有餘,稍稍高於本年度的家庭入息中位數19100元.他們幾經辛苦,找到一個售價三百萬元以下的住宅單位,終於能夠共築愛巢.蜜月過後,雙眼明亮過來,驚覺自己剛結婚就欠債二百多萬.他們做了為期二十五年的按揭貸款,除了維持生命的基本開支,例如:一日三餐,水電煤,交通費等,所有的收入都要拿去還債.欠債還錢,本是天公地道的事情,但想都未來二十五年,即是差不多整個能夠工作的有生之年,都要為還債而工作,為還債而生活,他們立即感到自己變了奴隸,就是內地流行的用語所說的「房奴」－房屋的奴隸.

這對夫妻欠下的債務,操控著他們的生活.因為要還債,他們不能失業；因為要還債,他們不能去旅遊渡假；因為要還債,他們要延後生養兒女；因為要還債,他們甚至要逃避參加朋友的婚宴.他們感到自己的人生已經被債務支配,活著,就是為了還債.

今天的書信經課,保羅以「欠債」的比喻去描述人們的生存狀況.保羅觀察及體會到,人們生存得好像被債務操控著一般.人們都底欠了誰的債呢?對於陷入這種債務困擾的人,有甚麼出路呢?讓我們細看保羅在羅馬書8:12-25節分享的體會及啟發.

二．經文

8 : 12弟兄們,這樣看來,我們不是欠肉體的債去順從肉體而活.

8 : 13你們若順從肉體活著,必定會死；若靠著聖靈把身體的惡行處死,就必存活.

8 : 14因為凡被神的靈引導的都是神的兒子.

8 : 15你們所領受的不是奴僕的靈,仍舊害怕；所領受的是兒子名分的靈,因此我們呼叫:「阿爸,父!」

8 : 16聖靈自己與我們的靈一同見證我們是神的兒女.

8 : 17若是兒女,就是後嗣,是神的後嗣,和基督同作後嗣.如果我們和他一同受苦,是要我們和他一同得榮耀.

8 : 18我認為,現在的苦楚,若比起將來要顯示給我們的榮耀,是不足介意的.

8 : 19受造之物切望等候神的眾子顯出來.

8 : 20因為受造之物屈服在虛空之下,不是自己願意,而是因那使它屈服的叫他如此.但受造之物仍然指望從敗壞的轄制下得釋放,得享神兒女榮耀的自由.

8 : 21併於上節

8 : 22我們知道,一切受造之物一同呻吟,一同忍受陣痛,直到如今.

8 : 23不但如此,就是我們這有聖靈作初熟果子的,也是自己內心呻吟,等候得著兒子的名分,就是我們的身體得救贖.

8 : 24我們得救是在於盼望；可是看得見的盼望就不是盼望.誰還去盼望他所看得見的呢?

8 : 25但我們若盼望那看不見的,我們就耐心等候.

三．釋經與應用

1．人人都債務纏身

8 : 12弟兄們,這樣看來,我們不是欠肉體的債去順從肉體而活.

8 : 13你們若順從肉體活著,必定會死；若靠著聖靈把身體的惡行處死,就必存活.

保羅指出,無論有沒有借錢,人人都好像債務纏身.債主是誰呢?債主不是銀行,不是「大耳窿」,而是自己的「肉體」.「肉體」這個字在保羅的筆下有多重意義:可以指人類或動物的血肉之軀,也可以泛指人,在這種意義下,「肉體」這個字是中性的.這中性的「肉體」,在某些文意之下,其含意是抽象而負面的.例如在羅馬書第八章,「肉體」所指涉的是與靈性相對的人性.保羅看見,不少人都生活得好像欠了肉體的債一樣.保羅同時發現,若生存只是滿足肉體,人們遲早都會成為肉體的奴隸,失去自己.令保羅感慨的是,肉體對人的操控力量,連他自己也不能倖免.他深刻地體會到肉體好像他的債主般向他追債,令他非常痛苦.保羅在前文如此分享他的掙扎:

「

我所願意的,我並不做；我所恨惡的,我反而去做．．．立志為善由得我,只是行出來由不得我．．．因為,按著我裡面的人,我喜歡上帝的律,但我看出肢體中另有個律和我內心的律交戰,把我擄去,使我附從那肢體中罪的律.我真是苦啊!誰能救我脫離這必死的身體呢?」(羅馬書7:15-24)

保羅的困擾不在於他不能夠選擇,而是他行不出自己的選擇,甚至做出自己也討厭的事情.為甚麼會這樣呢?他感受到有一股力量藉著他的身體把他「俘擄」,因此,他慨嘆:「誰能救我脫離這必死的身體呢?」

學生都知道上課不應該打蓋睡.然而,往往坐進舒服清涼的演講室一會,「睡魔」就「急Call」.為了對抗「睡魔」,有些人捻自己的大腿,有些人拍打自己的臉,有些人甚至咬自己的舌頭,刺激自己清醒.老師們未必需要因為這些學生上課睡覺而自責,大家都知道,在課堂上睡著了的大學生大多都在還債──因為晚上沒睡得夠,他們的身體向他們追討睡眠的債.學生難以抗拒「睡魔」的「急Call」,會否因為這根本就是他們身體的需要呢?保羅體會肉體的力量如此巨大,會否因為肉體牽引出我們的需要,同時也在滿足我們的需要呢?

上課時,我們好像欠了睡眠的債一般,不能自控地違背自己的意願在課堂上睡覺；在人生的路途上,我們有否好像欠了誰的債一般,過著不由自主的還債式生活呢?若果債務纏身,逃避並不是辦法,以債蓋債更會令債台越築越高.保羅接著的講論,帶引我們去思考人生的其中三項債務,同時啟發我們清除這些債務的正確方向.

2．從尋找歸屬的債務中得釋放

8 : 14因為凡被神的靈引導的都是神的兒子.

8 : 15你們所領受的不是奴僕的靈,仍舊害怕；所領受的是兒子名分的靈,因此我們呼叫:「阿爸,父!」

8 : 16聖靈自己與我們的靈一同見證我們是神的兒女.

保羅在上文勸募信徒不要活得好像欠了肉體的債一樣,被肉體支配著自己的生命,接著力證信徒作為神的兒女的身份,這兩件事情好像風馬牛不相及,放在一起有甚麼啟發呢?

我們呼叫神為「阿爸,父!」,顯示出我們與神親切的父母與子女的相屬關係,這種相屬關係,其實對人非常重要.不少人終其一生,就是在尋找這種歸屬關係.有些欠缺這種歸屬關係的人,好像欠下巨債一樣,不斷地找各樣事物來填補.有時,填補不得其法,令自己泥足深陷.

有一位報紙的專欄作家,得聞好友的一位弟弟因為吸食K仔三年而失去腎功能,發出如此的感概:「為甚麼?年輕人,人生有好多事情,你們還未經歷過,開心的,不開心的,精彩的,沉悶的,統統本來都應該屬於你的,就是因為吸毒,你將本來一個人應該擁有的所有經歷,變成了這樣的一個結局,一個令人痛心的結局.」(\url{http://orientaldaily.on.cc/cnt/lifestyle/20090628/00329_001.html})下面這位朋友的經驗,或者能夠提供線索給上述那位專欄作家.有一位吸毒者走頭無路,到戒毒機構去接受福音戒毒.工作人員問及這位戒毒者的家庭狀況,他怒氣沖沖地叫嚷:「如果他們關心我,就根本不會有今日.」(\url{http://www.gospelherald.com.hk/news/min_2279.htm})眾多的經驗告訴我們,不少沈溺行為都帶有以沈溺來填補歸屬關係缺乏的根源.

我被關心嗎?我被重視嗎?我被接納嗎?這些都是人的基本需要.嬰兒的大部份行為,或哭或笑,就是要引起父母注意,期望父母重視.當父母在自己的嬰兒面前抱著另一位嬰兒,自己的嬰兒立即有反應,因為他的歸屬關係受到威脅──父母是屬於自己的,自己是屬於父母的.

有一位年輕朋友「界手」.我問她「界手」那麼痛,為甚麼仍「界手」呢?她說,「界手」可以令她忘記男朋友移情別戀的心痛.每個人心裡都有這些問題:到底我屬於誰?誰又屬於我?誰與我彼此相屬呢?為甚麼有些人在親人或戀人面前自殺呢?尚未死心的,希望親人或戀人緊張他,重視他；已經老羞成怒的,要令親人或戀人心痛──心痛,證明你們心裡面仍有我.

有心理學者指出,父母與子女的關係是最基礎的歸屬關係,如果子女不能夠從與父母的關係中得到歸屬的滿足,他們就好像欠債一樣,耗用大量的資源去在人生各種場境中尋找這位歸屬關係.他們一方面尋求歸屬,同時間又害怕被拒絕；他們一方面討好別人以求被接納和肯定,同時卻因為得不到別人的接納及肯定而感到委屈,甚至憤怒.誰可以與我們發展這份親切,穩定及充滿接納的歸屬關係,以滿足我們的渴求呢?

保羅說:「你們所領受的不是奴僕的靈,仍舊害怕；所領受的是兒子名分的靈,因此我們呼叫:「阿爸,父!」」稱呼神為「阿爸,父!」,好像是老生常談,若我們細心想一想,其實這是劃時代的改變.在古代的世界,誰有資格稱神為「父」呢?稱神為父,自古都被皇帝所壟斷.保羅指出,事實不是這樣.神並非只屬於皇帝,也不只是皇帝才屬於神,神屬於祂所創造的每一個人,所有祂所創造的人也屬於神；神與人那份彼此相屬的關係,能夠修補我們在人生中所經驗到歸屬關係的欠缺,清除我們在歸屬關係上的債務.

那位指責家人不關心自己的戒毒人士建立了與天父的關係後,能夠以另一種態度看待家人.他說:「咁大個仔,都未好好做完一件事俾屋企人睇,我希望今次可以做好呢件事.」由此可知,他仍然非常重視家人,現在,他能夠在與天父的關係的安全基礎之上,以正面及健康的方向去回應家人.

3．從尋找保障的債務中得釋放

8 : 17若是兒女,就是後嗣,是神的後嗣,和基督同作後嗣.如果我們和他一同受苦,是要我們和他一同得榮耀.

神與人那種父母與子女的親切而穩定的關係,回應著人對歸屬關係的需要.人作為神的後嗣這宣告,又回應著我們的甚麼渴求,以及因這渴求而產生的債務呢?

不少人終其一生都在為自己及家人的生存及生存的條件尋求保障.有人在網上討論區發出這樣的問題:「而家香港地,要養大一個細路分分鐘要四百萬架,如果無呢一筆錢,就唔會俾到保障d小朋友.你地會唔會生呀?」能否為子女提供保障這個問題,嚇怕了不少人.有一位網友回應說:「我不會生孩子,小孩很本事,很難教,又怕他學懷,自己有空閒時會計算一下母親投資了幾多在我身上:從嬰孩開始,幼稚園,小學,中學,我讀書不算棒,要去澳洲念4年大學,還未計算一些醫葯費,托兒,補習,課外活動,興趣班,從小到大在家中的吃住費用．．．我都不曉得這筆賬該怎麼去算,但母親能得到的回報就．．．我自問要進修,要交租,要開支,哪有多餘的錢給家裡?當我算完這筆賬時,就不想生孩子,自己就是一個模板,我與男朋友都不想生,養了3隻貓,打算讓貓咪陪我過這輩子.」(\url{http://hk.knowledge.yahoo.com/question/question?qid=7007041900714})

有一對夫婦很勇敢,差不多「餐搵餐食餐餐清」(找一餐清一餐),也生了一個女兒.財務顧問「恐嚇」他們:「如果你不想愛女承受香港教育之苦,就要購買教育基金,為她提供較理想的教育保障.」結果,這對本來已經拮据的夫婦,某些月份要向家人借錢以支付女兒教育基金的供款.「一身兒女債」都是為了將來尋找一點保障.

人人都需要保障,但保障又好像是無底深潭,永遠填不滿.你的退休保障足夠嗎?有人問:要儲多少錢才能夠退休?這人退休時將有155萬元強積金供款,以及433萬元其他儲蓄,總計近600萬元.對不少人來說,這已經是天文數字,但理財顧問的意見是:「如果你會活到香港人的平均壽命,你仍欠約123萬才夠退休.」我們本已有著尋找保障的自然傾向,四周的聲音又不停地「恐嚇」著我們,令我們感到保障永遠都不足夠,鞭策著我們去追逐那永遠都不會足夠的保障.

在怎樣的情況下才得到一世無憂的保障呢?不少人的回應是,再投胎出生在富豪家庭.有一位時裝界朋友如此抱怨:「你有沒有曾經抱怨自己投錯胎．．．含著金鎖匙出世的人,天生註定住洋樓養番狗,不是ABC就是BBC,畢業是名牌大學,用的是Birkin手袋與Chanel,所以東西得來全不費功夫.若然父母是甚麼大明星或知名人士更不得了,尤其當你在跟紅頂白的時裝界與娛樂圈打滾,甚麼都跟你講名氣,你企定定已經有人會主動向你招手.」(\url{http://hk.lifestyle.yahoo.com/fashion/women/article.html?id=art_4ab9d538})

人們都認為,有足夠的資源,包括財富資源及社會資源,才能有足夠的保障.保羅指出,其實我們也承續著非常豐富的資源,因為我們就是神的後嗣.神所擁有的一切,我們都有份.作為神的後嗣,可以得到甚麼保障呢?這話題可以引發各種的想像,保羅提示我們參考基督的例子,他指出我們與基督同作後嗣,將會與基督一同得到榮耀.

「保障」這個債主,永遠也不滿足,企圖永遠勞役我們.保羅提醒我們,與其身心疲累地去填飽這個永不飽足的債主,不如轉眼望清楚我們在神裡面的身份,數算神已經為我們預備了的保障.

4．從害怕苦楚的債務中得釋放

8 : 18我認為,現在的苦楚,若比起將來要顯示給我們的榮耀,是不足介意的.

8 : 19受造之物切望等候神的眾子顯出來.

8 : 20因為受造之物屈服在虛空之下,不是自己願意,而是因那使它屈服的叫他如此.但受造之物仍然指望從敗壞的轄制下得釋放,得享神兒女榮耀的自由.

8 : 21併於上節

8 : 22我們知道,一切受造之物一同呻吟,一同忍受陣痛,直到如今.

8 : 23不但如此,就是我們這有聖靈作初熟果子的,也是自己內心呻吟,等候得著兒子的名分,就是我們的身體得救贖.

8 : 24我們得救是在於盼望；可是看得見的盼望就不是盼望.誰還去盼望他所看得見的呢?

8 : 25但我們若盼望那看不見的,我們就耐心等候.

在人們的經驗中,如債主一樣追著我們不放的,而我們又極力想擺脫它的,苦楚必定是其中之一.有些人一生的心思及精力,就是尋求擺脫苦楚.有人問:「如何擺脫失戀的苦楚?」；又有人問:「如何擺脫人際衝突的苦楚?」；可能最多人問的是:「如何擺脫身體疾病的苦楚?」疾病可以對人產生非常大的威嚇,它侵襲我們的身體令我們痛苦,它侵襲我們的家人的身體更令我們痛苦.它好像隨時到來,避無可避；有時又好像糾纏不休,驅之不散.疾病的威嚇,可以令人生活得提心吊膽,千金散盡,仍買不到安心.

面對著人類逃避痛苦的自然傾向,但又避無可避的現實,保羅的回應不是派止痛藥,而是鼓勵我們心存盼望.保羅勸導我們一方面不需要顧影自憐,因為,「一切受造之物一同呻吟,一同忍受陣痛」,我們不是唯一經驗苦楚的一個,在敗壞的世界得到完全的救贖之前,全人類,以至所有受造之物都活在相同的狀況之中.另一方面,保羅建議我們把眼光放長一點,我們就會看見苦楚總有過去的一天,並且美好的結局已經預備好了.若有這樣盼望,「現在的苦楚,若比起將來要顯示給我們的榮耀,是不足介意的」.

保羅自己也經歷各種的苦楚,他說:「

坐過更多次監牢,受過無數次的鞭打,常常冒死．．．我勞碌困苦,常常失眠,又飢又渴,忍飢耐寒,赤身露體.除了這些外表的事以外,我還有為眾教會操心的事天天壓在我身上.」(林後11:23-28)保羅還說有「一根刺」加在他的身上,這可能是疾病.他求主使這根刺離開他,但沒如願.對此,他有這樣的領受:「他對我說:「

我的恩典是夠你用的,因為我的能力是在人的軟弱上顯得完全.」(林後12:7-9)苦楚沒有把保羅擊敗,因為保羅心中的盼望沒有被苦楚奪走.

四．總結神與我們那種那種親切又穩定的父母子女關係,填補了我們在歸屬關係上的欠缺；我們作為神的後嗣的身份,給予我們資源回應現世保障的不足；那已經成就的盼望賦予我們韌力去經驗眾生同遇的苦楚.道理就是這樣,但改變如何發生呢?

著名基督徒作家Max Lucado寫了一本給兒童故事書,名叫You Are Special.故事的主角胖哥生活在一個每天互相貼貼紙的社會,得人喜歡的就被貼金星貼紙,不被欣賞的就被貼灰點貼紙.胖哥努力要得人喜歡,爭取更多金星貼紙,但事與願為,得回來的都是灰點貼紙.有一天,胖哥遇到路西亞,她身上竟然沒有金星貼紙,也沒有灰點貼紙,甚麼貼紙都不能貼在路西亞身上.胖哥感到非常驚訝,向路西亞請教.路西亞告訴他:「很簡單啊,我每天都去找伊萊,就是那位創造者.」克服了猶豫,經過了掙扎,胖哥終於鼓起勇氣去找伊萊,在回家的路上,有幾張貼紙從胖哥的身上掉下來了.

改變如何發生呢?改變在接觸中發生,改變在關係中發生.如果有一天,我們能夠以「阿爸,父!」這親切的稱呼呼叫神,我們對於歸屬關係的欠缺,對於保障的欠缺,對於苦楚的害怕會逐漸減少,生活可以越來越輕省.



\newpage

\section{知與情的交戰}
\label{sec:109}
\textbf{賴品超教授 - 知與情的交戰}
\newline
\newline
\hyperref[sec:108]{< < < PREV SERMON < < <}
~
\hyperref[sec:toc]{[返目錄]}
~
\hyperref[sec:110]{> > > NEXT SERMON > > >}


約拿書 4:1-11
\newline
\begin{longtable}{cl}
\hline
\hline
章節 & 經文 (和合本修訂版)\\
\hline
4:1 & \begin{tabularx}{0.7\textwidth}{X} 這事令約拿大大不悅，甚至發怒。 \end{tabularx} \\ \\ \relax
4:2 & \begin{tabularx}{0.7\textwidth}{X} 他就向耶和華禱告，說：「耶和華啊，這不就是我仍在本國的時候所說的嗎？我知道你是有恩惠，有憐憫的神，不輕易發怒，有豐盛的慈愛，並且會改變心意，不降那災難。我就是因為這樣，才急速逃往他施去的呀！ \end{tabularx} \\ \\ \relax
4:3 & \begin{tabularx}{0.7\textwidth}{X} 耶和華啊，現在求你取走我的性命吧！因為我死了比活著更好。」 \end{tabularx} \\ \\ \relax
4:4 & \begin{tabularx}{0.7\textwidth}{X} 耶和華說：「你這樣發怒，對嗎？」 \end{tabularx} \\ \\ \relax
4:5 & \begin{tabularx}{0.7\textwidth}{X} 約拿出城，坐在城的東邊，在那裡為自己搭了一座棚。他坐在棚子的蔭下，要看看城裡會發生甚麼事。 \end{tabularx} \\ \\ \relax
4:6 & \begin{tabularx}{0.7\textwidth}{X} 耶和華神安排了一棵蓖麻，使它生長高過約拿，影子遮蓋他的頭，使他免受苦難；約拿因這棵蓖麻大大歡喜。 \end{tabularx} \\ \\ \relax
4:7 & \begin{tabularx}{0.7\textwidth}{X} 次日黎明，神卻安排一條蟲來咬這蓖麻，以致枯乾。 \end{tabularx} \\ \\ \relax
4:8 & \begin{tabularx}{0.7\textwidth}{X} 太陽出來的時候，神安排炎熱的東風，太陽曝曬約拿的頭，使他發昏，他就為自己求死，說：「我死了比活著更好！」 \end{tabularx} \\ \\ \relax
4:9 & \begin{tabularx}{0.7\textwidth}{X} 神對約拿說：「你因這棵蓖麻這樣發怒，對嗎？」他說：「我發怒以至於死，都是對的！」 \end{tabularx} \\ \\ \relax
4:10 & \begin{tabularx}{0.7\textwidth}{X} 耶和華說：「這棵蓖麻你沒有為它操勞，也不是你使它長大的；它一夜生長，一夜枯死，你尚且愛惜； \end{tabularx} \\ \\ \relax
4:11 & \begin{tabularx}{0.7\textwidth}{X} 何況這尼尼微大城，其中不能分辨左右手的就有十二萬多人，還有許多牲畜，我豈能不愛惜呢？」 \end{tabularx} \\ \\
[1ex]
\hline
\hline
\end{longtable}
日期:2011年7月31日

題目:知與情的交戰

證道:賴品超教授

讀經:約拿書4:1-11

今日主要和大家思想約拿書裡面的經文,在第四章開始講述一件事,相信大家在主日學裡都聽過,這是關於上帝呼召約拿去尼尼微城宣講福音訊息,而約拿卻逃跑離開了,所以有很多主日學的故事書都形容他為「逃跑的先知」──故事講述他是如何掉進海裡,再被大魚吃進肚裡,然後回到岸上去.

這段經文在教會中亦很重視,因為聖經提到約拿曾在魚腹中渡過三日,而後來馬太褔音十二章亦提及這件事是這世代最重要的神蹟,所以這真是一件很重要的神蹟.不過縱然這件事很有趣,但當你仔細地閱讀整段經文的時候,你會發現這件事並非經文的重點,整卷書的高潮非但不是在於約拿曾在魚腹中渡過三日,亦不是在尼尼微城人的信主悔改,因為如果這些事是經文的重點,那麼在第三章的時候就便可完結整卷約拿書,因為尼尼微城的人在第三章的時候悔改──合城人以及牛羊都披麻蒙灰──如果重點是在於先知怎樣完成他的工作,而當他的職責已經完成了,那麼還有什麼話是要說的呢?為什麼會有第四章呢?

第四章的存在並非附錄,乃是整卷書的高潮所在.前面三章所描述的,只不過是為第四章而鋪排,預備,情況就如同脫口秀(talk show)的妙語(punchline)般,在第四章最後,最終的一節,上帝對約拿說了一句話,就結束整卷書,當中沒有約拿的回應.

「

何況這尼尼微大城,其中不能分辨左手右手的有十二萬多人,並有許多牲畜,我豈能不愛惜呢.」這是第四章的最終的一節話,似乎整卷書的總結,主題,最後的底牌和訊息,就包含在這句說話當中,而之前所講及的事,似乎只是一個佈局.又有人說,書中所提到的事蹟,可以歸類為約拿綜合症:好像時下年青人容易逃避責任,上帝托負責任給他,他卻逃跑去了,後又被上帝捉回來,並領他完成責任.

你會發現整卷書不是單單敘述「順服」與「反叛」的問題,而是他能否明白上帝普愛世人甚至牲畜的心意,這才是事件的重點.如果要討論約拿的順服,最終約拿是真的順服了,縱然心裡頭是不順服,但行為上倒是順服了,因為他最終去了尼尼微城宣講訊息,只是他心裡未能過自己的關口,這可能是民族感情的關口,為什麼這樣說呢?先說說約拿這個名字,約拿字義是解作「鴿子」,而鴿子使人聯想起和平及傳遞訊息,這實在與先知身份絕對配合.在列王記下十四章廿三至廿七節亦有提及到一個叫作約拿的人物,但我們不能證實他們是否是同一個人,有可能只是同名同姓而已,而事實與否亦不重要.我們要注意的重點是:為什麼約拿要逃避上帝的呼召,不願前往尼尼微城去傳講福音呢?如果你有留心的話,其實在第四章上已有交代.當約拿看到尼尼微城的人悔改,上帝就沒有將所預言的災禍降到尼尼微城的人身上,約拿就因此大大的不悅.如果我們不去考慮其他的因素,我們就會覺得這個很荒謬,因為作為一個先知,傳講了上帝的福音,聽的人都願意悔改,他理應高興才對,因為聖經也記載,天使天軍都為那些悔改的人高興和讚美,但相反的,約拿卻不因此而高興,這會否是一件很奇怪的事呢?同一道理,作牧師的,因傳講福音而得到多人悔改信主,都會為此而高興,反之,如果聽道的人沒有為福音所動,這才會令牧師不高興.

要了解約拿的反常反應,我們要先了解那個時代的背景.尼尼微城當時是化表亞述帝國,對猶太人／以色列人來說是具有威脅性,即是說他們是敵國的人民,從「情」這個點上,就大概明白約拿不願到尼尼微城的因由.在「知識」上,約拿認識神的性情,明白到他是應該去那裡的.經文記載,約拿因為發怒,於是就禱告神:「

．．．耶和華阿,我在本國的時候,豈不是這樣說麼．我知道你是有恩典,有憐憫的　神,不輕易發怒,有豐盛的慈愛,並且後悔不降所說的災.所以我急速逃往他施去.」這就說明了他要逃往他施的原因──為了逃避上帝的呼召.明顯地,他是清楚認識上帝的──上帝是有豐盛的慈愛,又後悔不降所說的災──然而他仍然決定要逃走,是因為他不希望這事情真的發生.知識上,他是明白上帝的；但情感上,他不希望這事發生在尼尼微城人的身上,因為尼尼微城人正正是他自己國家,民族的敵人.他認為上帝不應該寬恕他們,因為尼尼微城的人罪大惡極,上帝是應該對他們施行審判的.

我曾經到過外國的博物館參觀有關亞述帝國的展覽,從當時某些城市發掘出來的文物圖像中,我了解到,作為一個猶太人如此惱怒他們是正常的,何解?原因是亞述那裡到處都是偶像,且罪大惡極,而他們也正在欺壓著自己的民族.所以縱然約拿明白上帝是一位有慈愛,有恩典的上帝,願意赦免他們的罪,但約拿卻不願意傳講福音訊息給尼尼微城的人,他希望尼尼微城的人繼續不悔改,好讓上帝來懲罰他們,使尼尼微城全然被毀滅.

在此我們明白雖然約拿知道上帝的性情,卻不願意按照上帝的心意而行,因為有他自己國家民族的感情.有時候,我也會明白他的感受,特別是當自己本來就不同意那事時,但卻因某些要去作那事的原因時,這實在是很難受的.如果我是約拿,我會跟神說:『你要原諒他們是有沒關係的,但你可否差派其他的人去,而不是我呢?』作為猶太人,神卻要差他去宣講訊息,好達成祂的心意,在感情上的確是很難順從的.

我們文學院很快地將要搬家了,所謂的搬家就是成立一個稱為人民樞紐,把文學院裡不同的學科放置在兩橦樓房當中.這是我上一任院長所建議的,而當時我是第一個舉手反對的,原因是我認為:第一點:『上屋搬下屋,不見一籮穀』(廣東俚俗),這樣搬家一定會花掉一筆錢；第二點:搬家必會製造相當的建築廢料,並不環保；第三點:同事搬近了,但並不代表會有更好的溝通,因為我與隔壁的同事一年也只會碰上五,六次,亦即是每年學系規定的系務會議和指定的會議上碰碰頭而已,其餘時間其實是不會遇見的,所以我很懷疑就算將不同學系置於同一個地方上面,會否因此就能增加學術上的合作；再者我從一些唱反調的舊書籍中看到,在近代的大學裡,不同的老師和不同的學系之間的主要聯系,就是中央冷暖氣系統,而更現代一點的說法就再加上互聯網,所以我認為是否把我們放在同一處並不重要.但最諷刺的是,當新上任院長要實行這個建議時,卻委派了我去執行這事(當然他並不知道我是持反對意見的),在感情上,要我去執行我所反對的事,你可以想像這當中的矛盾,我是第一個反對建議的,但卻被安排要去執行這件事,當中更要和不同學系的同事們討論地方分配的事宜,真的浪費了不少時間.理性上,政策既然已經決定了,我就要照著執行；但感情上,還是覺得非常難受,總認為我明明是反對此事的,理應找其他人來執行.所以,當完成聯絡事務工作後,我就提議應由專業人士來負責具體事項,終於找到同系一名建築師的同事.

約拿就是有著這樣一個的想法.尼尼微城對他來說本來就是一個罪大惡極的城市,對之他非但痛恨,且是自己國家民族最大的敵人.如果上帝要毀滅這城,這不就是對自己國家民族最大的解放嗎?這簡直是一個莫大的拯救.但現在卻要去為他們傳講福音,同時,也明白上帝會寬恕他們,讓他們悔改,所以,在感情上,這是很難接受的.事實上,還有其他先知可以被差派去的,為什麼要找自己呢?再者,除了差派先知之外,上帝也可採用其他不同的方法,例如聖經曾記載神用指頭在牆上寫字的神蹟,這也可使人悔改呀,為什麼一定要自己去呢?因此,約拿非常不悅,且要逃跑離開,當那城的人最終悔改了,他更大大地不悅.一開始的時候,約拿已表明自己的不情願,為什麼還要強迫他呢?於是他乘船往相反的方向逃離,但卻被同船的人丟進海裡,再被大魚吃進肚子裡,最後到了尼尼微城,他心裡頭真的非常不悅,所以最後他發怒了,認為死了比活著還好,寧死不看倒好.但上帝對他說:「

你這樣發怒合乎理麼」?接著約拿沒有回答上帝,因為這不是道理的爭論,而是情感上無法接受這件事.

然後經上記載約拿在東面搭了一座棚遠觀尼尼微城,上帝在這時候藉蓖麻樹安慰約拿.當你仔細地閱讀時,你會發現這段經文的情節猶如卡通片劇情,開始的時候,沒有這棵樹,接著一下子長起來,甚至高過約拿為他遮蔭,約拿就「

因這棵蓖麻大大喜樂」.但翌日黎明,上帝卻安排一條蟲子咬這蓖麻樹,以致樹下的約拿被烈日曝曬,使他發昏,於是約拿又向上帝求死說:「

我死了比活著還好」.上帝就對約拿說:「

你因這棵蓖麻發怒合乎理麼」?約拿回答說:「

我發怒以至於死,都合乎理」.

其實我大概能明暸他這種的心情.在佛學中,我們常說「無常」──就是要去避免太多執著.試想想,如果蓖麻樹從來都不存在過,而約拿也一直地被烈日曝曬著,那就不會有甚麼不妥,因為這會被日頭曝曬的位置是他自己所選取的.正如一個人在窮困中突然發大財,但又在一夜間突然地失去,這實在是很難接受的；又好像一個人在窮困中突然中了六合彩,在高興之際,卻發現中的不是該期的六合彩,那樣是很受傷害的；又像政府在今年落實了的優惠政策,但次年就取消,市民會因期望落空,而深感難受.約拿的情況,就是這個樣子,蓖麻樹當時根本就不存在,他原本坐在那裡等待,就是要看尼尼微城的傾覆,而無懼烈日的曝曬.但蓖麻樹的出現與枯萎,好像上帝在開約拿玩笑,一開始約拿為蓖麻樹而感恩,但第二天卻枯萎了,約拿感到被上帝玩弄著.上帝好像帶盼望給約拿,卻馬上又拿走,所以約拿自覺就算「

發怒以至於死,都合乎理」.於是「

耶和華說,這蓖麻不是你栽種的,也不是你培養的.一夜發生,一夜乾死你尚且愛惜」──上帝問約拿,他對這蓖麻樹有多少感情?他只不過與這蓖麻樹相處了一天的時間──上帝藉著這蓖麻樹作實例教材,讓約拿明白,蓖麻樹一夜長成,又一夜枯死,他尚且愛惜它,那麼他是否又能體會上帝的心情?那些在尼尼微城中不能分辨左手右手的有十二萬多人,而且還有其他牲畜,上帝豈能不愛惜呢?

上帝兜了一個很大的圈子來讓約拿經歷實際的經驗,使他明白上帝同樣愛惜尼尼微城的人以及牲畜.你會發現約拿書的主題不是尼尼微城的人如何悔改,或是在魚腹三天三夜的神蹟,而是上帝如何管教祂的先知──約拿.祂讓約拿能更深地認識自己,而不是單單停留在頭腦上的知識,對上帝性情的認知──「

．．．我知道你是有恩典,有憐憫的　神,不輕易發怒,有豐盛的慈愛．．．」──但卻無法體會上帝的心情,他內心滿有對自己國家民族的情感,認為尼尼微城的人充滿罪惡,又是自己國家的敵人,這一些的想法迫使他在情感上不能明白上帝的心腸.

有些時候我們也會這樣子,就是情與知識未能配合上.我有一個同事是研究香港電影的,據他說,香港某電影公司有一段時間所出品的電影,後期製作參與的工作人員當中,有不少是日本人,可是在字幕上,觀眾是不會知道的,因為他們的名字是用中文顯示,看來像是中國人的名字,當然那只是他們的別名.在六,七十年代的香港,人們普遍不能接受日本人,所以他們要把自己的名字漢化,以避免觀眾發現有日本人參與,而選擇不看該電影,由此,你就可以想像,那時候人們心裡頭對日本人有著什麼樣的心態和感情.另外,有一次和那位同事的師父閒談,我問他:香港的電影,特別是在六,七十年代,電影中出現的日本人角色,絕大多數都是扮演壞蛋,奸角的角色,很少是做好人的.他很肯定地回答說:「什麼絕大多數?當中根本沒有「好」日本人的角色.」所以,你可以想像,電影反映出一個什麼樣的情緒.最近日本大地震事件發生後,網上有些留言是這樣寫:「該死的．．．他們在南京的大屠殺．．．他們又篡改教科書．．．．」.由於時代已經不同了,我亦看到其他不同的回應:「南京大屠殺的共犯都已經死了,現代的人與上一代的人是不應該拉在一起討論,今天這件事是關乎人命傷亡,我們不應該幸災樂禍地認定這是因為祖宗犯了罪,所以招致天遣.」時代真的已經不一樣了,反日的情緒不再那麼的高漲了；時代真的轉變了,不但是這樣,當翻查一些歷史資料時,發現早在很早以前,已有一位牧師──趙世光牧師,曾到日本做了很多宣教的工作(似乎教會在這方面較當時候的人容易跨越民族感情的障礙,對日本人產生同情,憐憫和關懷的感情,不會因民族感情的障礙,而不去愛他們).當然,時至今日,我們已經對日本的仇恨已漸漸淡忘,我最近甚至看到有一間中學門外拉起了布條,上面雖然寫上日文,不容明白內容,但可意會到的是為日本這次地震的災打打氣氣,鼓厲他們可以很快地渡過難關,重建家園等等.我相信在教會裡,我們可以慢慢克服民族情感障礙,並體會上帝的愛.然而我們還有其他的難關是需要跨越的,就是我們對「萬物」的愛惜.

我們多年來談論環保的問題,我自己也寫過一些有關環保神學,但是最後我發現,原來不是我們理智上知道的有多少,特別是在大學裡的神學教育,通常都只在乎知識的傳授,甚至一般神學院都有這方面的傾向──注重知識的傳遞,教導我們上帝是怎樣,怎樣的,而忽略了情感上的培養,甚至教會也如是.試想想,我們在教會多做些什麼?就是坐在那裡,聽聽道.聽道也是一種知識的傳遞,而情感上是沒有改變的,所以我們需要想想如何在教會中把情感給培養出來.情感不但需要跨越民族感情的障礙,且是要把上帝的愛伸延至經文所提及的牲畜,甚至那棟蓖麻樹.究竟我們對那些牲畜,身邊的植物和環境的生態有多少的情感呢?會否體會那些生物都是上帝所鍾愛的呢?還是依舊停留在知識層面上的明白?

從前和合本翻譯本裡寫的『上帝愛世人』,其實就應該翻譯成『上帝愛世界』.上帝不單單是愛人,而是愛整個世界,我們在情感上,能否真的體會到呢?願意今天透過經文,讓我們思想,我們到底有多體會上帝的心腸呢?上帝可以吩咐不同的僕人去完成祂的工作,但祂卻選用了一個迂迴的方法,利用到大魚,蓖麻樹,蟲子和船上的水手,硬要約拿去接受這次教育的機會,好讓他明白上帝的心腸.

上帝不是需要人去完成祂的工作,因這只是僕人的職份.上帝需要朋友,這種朋友是要明白上帝的心腸──不僅要完聖所要完成的工作,且要仔細明白工作背後的因由,更加要明白背後是一個什麼樣的心腸.

上帝不止要僕人,更要朋友.我們是否願意成為上帝的朋友,不但替祂完成祂的工作,也體會上帝普愛世人和普愛眾生的心腸.

願主祝福我們,讓我們的愛在祂的愛裡日漸成長,成為祂的朋友,明白祂的心腸.



\newpage

\section{種子土壤與撒種的人}
\label{sec:110}
\textbf{戴國璋牧師 - 種子,土壤與撒種的人講員}
\newline
\newline
\hyperref[sec:109]{< < < PREV SERMON < < <}
~
\hyperref[sec:toc]{[返目錄]}
~
\hyperref[sec:111]{> > > NEXT SERMON > > >}


馬太福音 13:10-16
\newline
\begin{longtable}{cl}
\hline
\hline
章節 & 經文 (和合本修訂版)\\
\hline
13:10 & \begin{tabularx}{0.7\textwidth}{X} 門徒進前來問耶穌：「對眾人講話，為甚麼用比喻呢？」 \end{tabularx} \\ \\ \relax
13:11 & \begin{tabularx}{0.7\textwidth}{X} 耶穌回答他們說：「因為天國的奧祕只讓你們知道，不讓他們知道。 \end{tabularx} \\ \\ \relax
13:12 & \begin{tabularx}{0.7\textwidth}{X} 凡有的，還要給他，讓他有餘；凡沒有的，連他所有的也要奪去。 \end{tabularx} \\ \\ \relax
13:13 & \begin{tabularx}{0.7\textwidth}{X} 我之所以用比喻對他們講，是因為他們看卻看不清，聽卻聽不見，也不明白。 \end{tabularx} \\ \\ \relax
13:14 & \begin{tabularx}{0.7\textwidth}{X} 在他們身上，正應驗了以賽亞的預言：『你們聽了又聽，卻不明白，看了又看，卻看不清。 \end{tabularx} \\ \\ \relax
13:15 & \begin{tabularx}{0.7\textwidth}{X} 因為這百姓的心麻木，耳朵發沉，眼睛閉著，免得眼睛看見，耳朵聽見，心裡明白，回轉過來，我會醫治他們。』 \end{tabularx} \\ \\ \relax
13:16 & \begin{tabularx}{0.7\textwidth}{X} 但你們的眼睛是有福的，因為看得見；你們的耳朵也是有福的，因為聽得見。 \end{tabularx} \\ \\
[1ex]
\hline
\hline
\end{longtable}
日期:2011年8月7日

主題:種子,土壤與撒種的人

經文:馬太福音13:10-16

講員:戴國璋牧師

今天的講題是「種子,土壤與撒種的人」,出於耶穌一個很著名的比喻,就是有關撒種的比喻.在今天的經課裡,福音書的部分我們已經讀過有關這個撒種比喻的內容,但是在這裡我想跟大家補充兩段聖經,其實也是跟耶穌基督撒種的比喻是有關的,因為今天我們的重點不是放在了解撒種的比喻,反而是想從耶穌基督撒種的比喻裡了解上帝很想我們能夠在聆聽上帝的道上能有更進深的功課,故此有兩段經文很想跟大家一起再讀.第一段經文出於馬太福音13:10-16,也是我們今天早上福音書經課剛好中間的一段經文,如果大家有聖經的話很歡迎大家一起翻開那頁,經文是這樣說的:「10 門徒進前來,問耶穌說,對眾人講話,為甚麼用比喻呢. 11 耶穌回答說,因為天國的奧秘,只叫你們知道,不叫他們知道. 12 凡有的,還要加給他,叫他有餘．凡沒有的,連他所有的,也要奪去. 13 所以我用比喻對他們講,是因他們看也看不見,聽也聽不見,也不明白. 14 在他們身上,正應了以賽亞的預言,說,『你們聽是要聽見,卻不明白．看是要看見,卻不曉得．』 15 因為這百姓油蒙了心,耳朵發沉,眼睛閉著．恐怕眼睛看見,耳朵聽見,心裡明白,回轉過來,我就醫治他們. 16 但你們的眼睛是有福的,因為看見了你們的耳朵也是有福的,因為聽見了.」

另一段補充的經文是路加福音8:16-21,原因是在路加福音也有與馬太福音相同的比喻,就是撒種的比喻,不過當中的說話是馬太沒有記述下來的,所以很想跟大家看看這段經文裡耶穌的補充.路加福音8:16-21經文如此說:「 16 沒有人點燈用器皿蓋上,或放在床底下,乃是放在燈臺上,叫進來的人看見亮光. 17 因為掩藏的事,沒有不顯出來的．隱瞞的事,沒有不露出來被人知道的. 18 所以你們應當小心怎樣聽．因為凡有的,還要加給他．凡沒有的,連他自以為有的,也要奪去. 19 耶穌的母親和他弟兄來了,因為人多,不得到他跟前. 20 有人告訴他說,你母親,和你弟兄,站在外邊,要見你. 21 耶穌回答說,聽了　神之道而遵行的人,就是我的母親,我的弟兄了.」如果我們讀過馬太福音的故事,這些經文也是我們所熟悉的,尤其是在第13章裡天國的比喻,或者是路加福音裡的故事,我們也發覺這些經文也是大家很熟悉的,但是如果我們將這些經文放在一起的時候,我們就發現他提醒我們有關這個撒種的比喻一些很奇怪的事實.我們知道撒種的比喻是馬太福音第13章天國比喻裡的第一個,馬太將很多很多耶穌基督有關天國的比喻組合起來放在一章裡去解決,這些比喻在其他福音書裡很多時候就散見在耶穌的講論中間,馬太福音13章裡第一個比喻就是撒種的比喻,可見這個撒種的比喻其實是很重要的,因為它提醒我們要怎樣去聽耶穌的教訓,怎樣去聽上帝的話,以此作為一個前導讓我們更重視後面的接著要說的許許多多跟天國有關的比喻.有關這個撒種的比喻,我們不需要太過花時間去解釋,原因是馬太福音在後面的地方,譬如是在我們的經課裡我們也特別將耶穌基督對這個比喻解釋放在裡面,所以我們可以很清楚地了解到有關撒種比喻裡的稻種落在各種土地上,土壤的意義.

今日我很想跟大家特別留意反而是當耶穌說畢這個比喻後,到他向門徒解明之間,中間是有少少的環節,經文告訴我們當耶穌在群眾的面前說畢後就沉默了,然後他就離開群眾走入屋內,這個走入屋內既可以是字面上的理解,表示他已經離開了群眾,進入到一間屋的裡面；但是也可以理解為耶穌其實是進入了跟門徒私下的空間中,將群眾排除了,將一些特別重要的訊息讓門徒知道.事實上,當耶穌進到屋內,門徒就問耶穌:「對眾人講話為何要用比喻呢?」門徒這個問題問得很好,耶穌往後的解釋就讓我們更加深層次地了解到上帝傳道的心意,尤其是有關撒種比喻裡進一步的消息.門徒查問為什麼會用比喻,其實就引發出耶穌在這個剛才我們所讀聖經原文裡的講論及後面有關這個比喻的正確解釋.或許門徒的這個問題我們會覺得有點多餘,為什麼呢?因為我們自己都覺得自己應該明白為什麼耶穌基督會用比喻,道理很簡單,就是因為對群眾講話總是要用一些較淺白的道理；淺白的道理才能讓人聽得明白,所以目的是要讓群眾更明白天國的奧秘,這是我們的理解.事實上,我們知道用比喻或者用寓言也是常見的文學表達方式,在古今中外都是一件常事.其實比喻和寓言也是修辭手法,耶穌作為一個成功的真理教師,所以他也時常使用這種手法,而且他往往能夠靈活又生動地帶出一個教訓來,這是當我們常常讀到福音書裡關於耶穌的教訓的時候會得到的印象,所以比喻和寓言只不過是以故事的手法闡揚人生的真理,它既有吸引力,因為可以引發我們的興趣,有很多時候,它又有說明的能力,因為它能夠幫助我們去明瞭.譬如我們知道在公元前六世紀的時候,西方世界古希臘裡有個很著名的作者叫伊索,伊索本來是一個奴隸,不過他很聰明,也能善於講故事,也留下很多很多的故事,後來的人輯成為「伊索寓言」,我們知道「伊索寓言」裡的故事都不是真實的,它只不過是藉著動物裡的對話來講論人生,講論人性,所以當我們讀「伊索寓言」的時候,尤其是在我們小時候讀這些故事的時候,我們覺得這些故事很親切,因為它牽涉很多很多小動物,同時我們又覺得道理很清晰,因為往往讀完這些故事後,我們也能領會其中提醒我們做人的事.其實不單是公元前六世紀有伊索,在公元前四世紀在東方世界也有一個人叫莊周,也就是莊子,他也是一個善於使用寓言的人,他也用了很多著名的寓言來講述一些很深奧的道理.這些都是例子,無論在東方世界還是西方世界,在古代也會有善用寓言的教師,耶穌也是一個很成功的教師,因此他能夠用寓言來說明一些奧秘的說話.當耶穌要將天國的奧秘跟普羅大眾說明時,他唯有借助我們日常生活的經驗去比附超出我們經驗的真理,耶穌這種心意我們是可以了解的.但當我們細讀耶穌怎樣回答門徒這個我們看來比較多餘的問題時,我們就發現很奇怪,因為耶穌的回答絕對不是我們剛才所講的,我們所理解的方面.耶穌的回答是怎樣的呢?

在第十三到第十五節的回答裡,耶穌很清楚地表達:「其實我說比喻的目的是為使群眾不能夠明白天國的奧秘.」這個是令我們覺得奇怪的地方.剛才我們已經說過了,我們覺得用比喻其實是讓群眾更加明白天國的奧秘,為何耶穌反而說,他說比喻的目的其實是故意地使群眾不能明白天國的奧秘呢?耶穌是這樣說的,耶穌自言他用比喻的因由其實是想將這個真理變得模糊,變得難以明白,故意叫人聽不見,看不見.剛才我們所讀的很清楚是這樣說.耶穌的話使我們產生兩個問題,第一個問題就是「如何能」的問題,意思就是為何用比喻會反而是叫人聽得模糊呢?第二個問題才是「為何要」,為什麼要這樣做呢?讓我們嘗試去處理這兩個問題.首先是了解下「如何能」,因為剛才我們已經說其實聽比喻應該是更容易明白,為何是這樣呢?我們知道中國有很多成語故事,其中有一個成語故事發生在春秋戰國時代,叫「買櫝還珠」,不知道大家有沒有聽過這個故事.故事是出自「韓非子」,很多很多年前的春秋時代,有一個鄭國的商人去到某地方做生意,在市集裡將他帶來的商品一一陳列在地上,歡迎那些客人參觀選擇,其中一個客人見到他一件商品,就是一個盒裝著一顆很美很美的珠,他看中這件商品的時候,他就跟這個鄭國商人討價還價,就問這顆珠子和這個盒子賣多少錢,這個商人看見有客人看中,當然稱讚他有眼光,強烈推銷這粒珠子漂亮,便宜,最後大家商量了一個價錢後,這個客人就決定要買了它,於是就按照約定的價錢,將錢給了這個鄭國商人,拿了商品.當他拿到商品後,誰知道他很奇怪地將顆珠子拿出,然後還給商人,原來他看中的不是這顆珠子,而是裝這顆珠子的盒子.想到這件事的時候我也想到一個笑話,只是昨天的事情,發生在我家.我的女兒找到我太太一個盒,我的太太在學校教書教了很多年了,所以最近學校頒了一個獎給她,我的女兒的年紀還小,還有少少天真味,她拿了那個盒出來就打開見到有一件東西,然後就問媽媽那件東西是什麼,媽媽就說這件東西是學校送給她的一份紀念品,她就說我喜歡這個盒,給這個盒我吧,這個獎品我還給你.她喜歡那個盒勝過那個獎品.「買櫝還珠」故事裡的客人也一樣,原來花了這麼多時間,他看中的只是那個盒而不是那顆珠,我們中國就有一個這樣的比喻,叫「買櫝還珠」,這些比喻通常也有一些特別的道理在裡面,這個比喻其實是描述一個人缺乏眼光,見識短淺,所以當他處理事情的時候不懂得分輕重,往往棄貴而重賤,只是著重包裝而不著重內涵的珍貴,甚至輕視外在的平凡,賤看內在的包涵,其實這個「買櫝還珠」有點像耶穌基督所講的這個比喻,當耶穌用比喻來去講的時候,很多群眾也聽了,不過他們聽完就作罷,並沒有因為這樣而引發他們更深層次的興趣,他們也不會再去追問,很多時候,他們對上帝真理的明白就永遠只能得到這個比喻所給他們表現的層面.

我們回答「如何能」的時候,其實是想講,有很多時候用比喻的方式令那些人厭倦,覺得這個方式太顯淺了.如果我們要回答「為何要」,就是為什麼耶穌基督要講一些說話是令到真理變得模糊呢?耶穌就引用了14-15節,在這裡是引旨以賽亞書6:9-10的話,我們發現耶穌基督在這裡所講的這番話是將他自己比作當日的以賽亞先知一樣,耶穌自比為先知傳道向那些心硬不順從的百姓去宣講真理,他自知他所說的一切都將會是徒勞無功的；他又將自己比作先知所預言的那位救主,這位救主來到自己的地方卻會被自己的百姓拒絕,所以這句說話是表達了耶穌基督那份無奈的心情,他要無奈地面對著群眾的膚淺,以顯示出這種膚淺其實上帝早已知道.如是者,我們將我們第一點裡所理解的(我們覺得用比喻可以將天國的奧秘講得更加明白)和第二點耶穌的回答(他故意將它弄得更模糊)結合起來的時候,我們就會了解上帝的心意,就是上帝希望群眾當中那些願意相信的,他們能夠更進一步,更進深去追尋；至於那些不信的,他們就會失去他們的機會.其實門徒的查問是一個更好的說明,門徒跟群眾一樣聽完這個比喻後,他們心裡總是覺得有一些事我還沒有弄清楚,於是當耶穌離開群眾時,他們就走前來詢問耶穌,為何你要說比喻呢?如是者,耶穌基督就講得更多,所以門徒的查問就說明了,若果我們不是停留在表面的功夫,而是進深一步來了解,來知道的時候,其實我們就可以知道得更加多；但是如果我們只是停留在表面的功夫,聽了就算了,那只不過好像用一個漏水的花瓶儲水,儲多少也沒有用,因為一邊儲水,花瓶裡的水就一邊漏走.其實12節是很關鍵的一節,耶穌在第12節裡是這樣講的:「凡有的,還要加給他,叫他有餘．凡沒有的,連他所有的,也要奪去.」這句話耶穌不是只在這裡講過,我們在其他的福音書的經文中也能找到耶穌曾經這樣講過.他的意思要看放在什麼場合上應用,我們才能真正能掌握到.其實這就正正道出恩典的原則,這個恩典的原則就是:凡是對於那些渴慕的,祈求得著更多的,越有渴慕就越能得著更多,因為上帝還要加給他；但是如果是不渴慕的,就此了結的話,他所得的就只有這麼多的了,而且他所得的還要被奪去,他即使有聽到那個故事,有聽過那個道理,但轉頭就會忘記這個道理,於他一點幫助也沒有.所以這個恩典的原則其實正好說明,當我們聆聽上帝的道的時候,我們需要求上帝給我們一個熱切渴慕的心,不是只停留在聆聽的層面,而是更深地求知的層面.

接著我想跟大家去看路加福音的補充,在路加福音的故事裡是多了一些經文的.當我們從這些經文看的時候,就會更看見上帝渴慕讓我們明白真道的心意.在馬太福音,耶穌講解時,我們總是覺得有點「那個」,為什麼耶穌故意用比喻要讓他們耳朵發沉,讓他們看是看不見,聽是聽不見呢?難道這是上帝對子民的心意?若果我們加上路加福音第8章裡的故事,我們就會發現並不是這樣.其實路加福音的這個記述,加多了一些比喻.這些比喻是從另一個角度來令我們能夠了解到耶穌講比喻或者是撒種比喻的特點.路加保留的這段獨特的講論其實是要闡明上帝的兒子向人宣道那份熱切的心境.譬如我們先看路加福音8:16裡著名的點燈之喻,是一個很短的比喻,耶穌說「沒有人點燈用器皿蓋上,或放在床底下,乃是放在燈臺上,叫進來的人看見亮光」.這個比喻其實是要向我們說明上帝向人類宣示祂自己心意的那股衝動.其實這個比喻我們也很熟悉,不過我們熟悉它的地方不是在路加福音,而是在馬太福音；我們記得在馬太福音5:15,耶穌曾經用過這個比喻,用在那個「作鹽作光」的比喻當中.我們知道這個比喻在馬太福音裡的主旨其實是要教導門徒要將好行為展示在人的面前,正如人點燈,不是將它放在斗底下或者床底下,要展露在人的面前一樣,所以我們也應該將我們的好行為展示在人的面前.不過這句話可能是一個很常用的諺語,這個諺語有一個基本的意思,就是人點著燈為要照明事物,我們很明白這個道理,這是點燈基本的意思.這個基本的用法可用在不同的場合中,路加福音就將這個基本的用法用在耶穌基督講比喻的那件事裡,由此衍生出一個新的看法.路加的用法是要指明就像人點燈是要使人得見亮光,所以上帝講道的目的其實是想人明瞭,而不是想人糊塗,從此我們就可以了解到耶穌基督心裡的矛盾:他既看見人的不信,以至厭煩他們,因為他知道無論他做多少事,這些群眾還是會拒絕他的；但是另一方面,他內心又有一種極度的衝動向他們將上帝的道講得明白,將天國的真理來表達讓他們聽,他很希望他們明白,這種矛盾的心境在耶穌基督的身上顯露出來,路加的話就印證了耶穌用比喻的目的是要使普羅大眾更明白其中的意義.這句話從側面地了解就是上帝要向人宣示祂心意的熱情,而這種熱情在整個救贖的歷史裡也是非常清楚的.我們從亞當開始,我們記得耶和華呼喚那人,也就是亞當,問他到底在哪裡(創世記3:9)；然後亞伯,耶和華對他的兄弟該隱說,亞伯在哪裡呢?(創世記4:9)；向挪亞,亞拍拉罕等等的列祖,在荊棘裡祂呼叫說「摩西,摩西」(出埃及記3:4)；會幕建成後,上帝住在百姓中間,但是利未記1:1說耶和華仍然從會幕中呼喚摩西,因為祂有事是藉著摩西傳達給以色列人.孩童撒母耳,我們記得他在睡夢中三次聽見有人呼喚他,最後他得到他施父以利教他怎樣回應,他後來成為了一代的先知.這個睡夢中的呼喚,「撒母耳,撒母耳,撒母耳」,也就是表達上帝對人的呼喚；上帝在殿中呼喚以賽亞(以賽亞書第6章),以賽亞看見的異象和聽見聲音:「我可以差遣誰呢?誰肯為我們去呢?」那段經文也就是剛才我們所讀的耶穌引用以賽亞那番話的出處所在,然後我們發現在末世的時候,上帝藉著祂的兒子曉諭我們(希伯來書1:2),清楚地告訴我們上帝就在這末世,藉著祂兒子曉諭我們.我們從這些許許多多的話語和聖經的故事裡,我們就彷彿看見上帝內心的孤寂,祂很想向人傾訴祂的話,向人道出真理的話,但是祂找不到溝通的途徑,因為沒有人肯願意去聽.

這個思考令我想起荷里活的一套電影,叫「Cast Away」,是Tom Hanks(湯漢斯)很多年前的一套電影,在電視也重播了很多遍,也是因為在電視播放我才有機會看到.我記得這套電影很長,由九點半播到十二點多,我記得有一幕,Tom Hanks很淒慘地流落到一個孤島上,就是沒有人跟他談話,他又想保持著會說話,因為如果他不再說話,一直不說話,他害怕他會忘記說話,不再懂得說話,他無時無刻不展望自己會離開荒島,回到人間,然後可以跟人溝通.他用什麼方法呢?在隨他飄流到島上的飛機上的物資中--因為他本來是一個飛機速遞員,因飛機失事掉在海上飄到荒島上--在那些貨物中,他找到一個類似足球的東西,他就將這個足球畫上人的樣子,然後又幫這個足球起了一個名字,就當是他的partner,他由朝到晚就跟那個足球談話,無論吃飯,睡覺前還是睡醒後,都跟那個足球談話,彷彿這個足球就是他的知音人,他要盡量將他的心事向它傾訴.在孤島上的幾年間,沒有任何人跟他溝通.我看到這個故事的時候就想到上帝也有這種孤寂,祂也很想跟我們說話.既然是這樣,上帝不會將祂的心意隱藏,祂總是會用各樣的方法向我們說明,引導我們來到祂面前願意聆聽,這是我們從這個點燈之喻能夠看到的.特別的除了點燈之喻外,還有另一個的,就是接著的17節,路加福音8:17:「因為掩藏的事,沒有不顯出來的．隱瞞的事,沒有不露出來被人知道的.」這句話我們在馬太福音也有讀到的,這句是記載在馬太福音10:26,不過是兩個完全不同的場合,但是用了一句相同的說話或者是相同的比喻,當時耶穌基督是這樣說的:「所以不要怕他們．因為掩蓋的事,沒有不露出來的．隱藏的事,沒有不被人知道的.」在馬太福音的背景裡,耶穌當時是被人誹謗,當他完成了趕鬼的行動後,有人說他是靠鬼王趕鬼,所以耶穌說「人既罵家主是別西卜」--他常常被人這樣說--「人既罵家主是別西卜,何況他的家人呢」.接著耶穌就說,「所以,不要怕他們」.耶穌講這句話的意思其實是想鼓勵門徒不要認為被迫害的事是奇怪,因為正如我們剛才所說,耶穌早已預料他將會好像以賽亞先知一樣,所有的工作也是徒然,會是被人拒絕的,所以不要為被迫害的事而奇怪,反而要接受一個真理,這個真理就是我們剛才所說的第17節,那句可能是一個諺語,那句諺語的基本意思本來是「真相最終也會水落石出」,所以今日我們被人誤會不緊要,因為有一天真相總會水落石出的.如果我們了解到這句說話的基本義,然後將它放在路加福音第8章這裡耶穌用的這個地方的時候,我們就了解到耶穌在這裡是要表明他來傳道其實是要將藏在上帝裡最大的奧秘傳給我們知道,所以耶穌是這樣說的:「我在暗中告訴你們的,你們要在明處說出來.你們耳中所聽的,要在房上宣揚出來.」這是馬太福音10:27也有用到這個講論.上帝要尋找傾訴對象,要將祂的奧秘向我們揭示.其實不用,有很多天使,為什麼上帝不找天使傾訴,而要找我們呢?所以在我們心中或者在我們身上,我們可以見到我們在上帝心目中的地位是何等的重要.憋在心中的感覺是有事要說又不知道怎樣說或者是怎樣去表達,那種感覺其實真是很難受,有時候我們也難免會有這樣的心情.今天,我們的主也鼓勵我們要找多些人來成為我們朋輩間作互相傾訴的對象；在傾訴中,我們能夠得到溝通,得到安慰.我們鼓勵夫婦間,親人間,朋輩間,我們也有傾訴與聆聽,因為我們知道這種傾訴與聆聽本來就是重要的.如果是這樣,為什麼我們不去領會一下在神人之間也是同樣重要呢?

最後,我想跟大家看看路加福音,剛才我們所讀的經文裡,我們故意讀多了一些,就是有關「聽而行者,乃為至親」的這個講法.故事告訴我們當耶穌講完這些比喻後,就發生耶穌的母親和他的兄弟來找他而不得來見的情境,有人將此事告訴他說:「你母親和你弟兄站在外邊,要見你」.當時「耶穌回答說,聽了上帝之道而遵行的人,就是我的母親,我的弟兄了」.這個情境也是我們熟悉的,如果你讀過馬太福音也一定讀過一個類似的故事,兩者似乎也是同一個情境,不過奇怪的地方是大家放的位置有點不同.馬太其實是將這段故事放在第12章的結尾,就是撒種比喻之前,在耶穌講天國的比喻之前；但是路加就將這個故事放在耶穌講撒種比喻之後,我相信有特別的心意,而且結語也稍有不同:馬太福音說「凡遵行我天父旨意的人,就是我的弟兄姊妹和母親了」,路加的講法是故意地特別強調,「聽了上帝之道而遵行的人,就是我的母親,我的弟兄了」.這個特別的強調其實是間接地想突顯前面主耶穌所講撒種比喻的重要性.綜合這麼多的講解,我們要好好珍惜上帝給我們聽道的機會,不要聽了就算了,而是聽了之後,我們還要進一步的問「為什麼」,上帝我求你給我們更多.而在我們渴求更多的過程中,我們必然能得到更多,「因為凡有的,還要加給他,叫他有餘．凡沒有的,連他所有的,也要奪去」.

請我們同心祈禱:「我們的主,我們知道你的心意,我們知道你渴求將你自己的心意向願意的人揭露,我們也盼望我們每一個人能夠成為那些願意的人在你的面前來聆聽,聆聽你給我們的真道,滿足我們的內心,也能滿足你自己的心意,在這種相交中,我們既成為有智慧的人,又能成為上帝揭露心志的對象,願主的話發出且並不徒然返回,求你垂聽我們的祈禱,奉耶穌基督的名字,誠心所願.」



\newpage

\section{禱告的能力}
\label{sec:111}
\textbf{周志豪牧師 - 禱告的能力}
\newline
\newline
\hyperref[sec:110]{< < < PREV SERMON < < <}
~
\hyperref[sec:toc]{[返目錄]}
~
\hyperref[sec:112]{> > > NEXT SERMON > > >}


羅馬書 8:26-27
\newline
\begin{longtable}{cl}
\hline
\hline
章節 & 經文 (和合本修訂版)\\
\hline
8:26 & \begin{tabularx}{0.7\textwidth}{X} 同樣，我們的軟弱有聖靈幫助。我們本不知道當怎樣禱告，但是聖靈親自用無可言喻的嘆息替我們祈求。 \end{tabularx} \\ \\ \relax
8:27 & \begin{tabularx}{0.7\textwidth}{X} 那鑒察人心的知道聖靈所體貼的，因為聖靈照著神的旨意替聖徒祈求。 \end{tabularx} \\ \\
[1ex]
\hline
\hline
\end{longtable}
日期:2011年8月14日

題目:禱告的能力

證道:周志豪牧師

讀經:羅馬書8:26-27

首先代表遠東廣播公司良友電台--我事奉的機構--在主裡問候各位弟兄姊妹,也代表我們的教會--在深水Po 的中華聖潔會,它有八十多年歷史,是一家學校的教會--我們的弟兄姊妹問候今天在崇基學院的禮拜堂每位弟兄姊妹主內平安.願今天我們所看的聖經跟神的話親自給我們的內心帶來力量,帶來安慰.

剛才我們的弟兄姊妹讀聖經的經文有羅馬書第八章,我們看裡面第二十六節和二十七節兩節經文.我再為大家讀一遍,羅馬書第八章二十六,二十七兩節,使徒保羅這樣說:「況且我們的軟弱有聖靈幫助,我們本不曉得當怎樣禱告,只是聖靈親自用說不出來的歎息,替我們禱告.鑒察人心的,曉得勝靈的意思．因為聖靈照著　神的旨意替聖徒祈求.」我們今天在這個非常寶貴的時間裡面,思想一個很簡單的題目,也是很重要的題目,就是禱告的能力.受著這裡一個重要的字「軟弱」的影響,我本來的題目是「弱者的禱告」,但是很謝謝伍牧師,他很坦白,很不客氣地給我一個回應,他說這個題目好像不太恰當,因為在大學裡會有廣告宣傳,叫我再想一想,那我就把它反過來,變成是「禱告的能力」.

我們都有禱告的經驗.禱告雖然有很多理論,很多神學,但是更重要的是它是在生活裡一件實際的事情.農村裡的老太太可能不認識字,但是她可以虔誠地禱告.她所求的,無論是為家人,兒女,親戚,有一天都會有改變.我們聽過很多教會偉人的見證,母親的禱告,老人的禱告,改變他們的一生.在香港,每年夏天有很多很多人學習禱告,觀摩禱告,有上千的人組織旅行團到韓國去觀摩,體驗韓國教會的禱告.如果包括台灣,新加坡,東南亞,每年夏天真的是有很多人到韓國去.香港某一個靈恩的宗派舉辦了二十多年,每一年都有很多人參加.最近這幾年有大的宗派也有定期地到韓國去觀摩禱告.

因為我曾經長時間在電台做宣教士,我們不但在菲律賓有花時間,在韓國也有電台,所以我跟我內人在韓國長時間事奉,有二十多年的時間,跟韓國的教會和傳道人接觸.我們在韓國電台的總幹事是一個很有恩賜的教會領袖,他常常向人傳福音.我也有機會在他們的教會事奉過,早上真的是四點半就有二百人,三百人一起禱告.有一次,我被安排早上帶領禱告會.禱告的崇拜是每天都有的,四點半要到教會,所以我就把鬧鐘弄到四點半.因為是在另一個城市,我心裡面又興奮又緊張.早上四點半去韓國教會帶領這個禱告會,只是講十分鐘,剩下的半個鐘都是大家禱告.太興奮了!我早上一點半醒了一次,兩點半醒了一次,然後三點半左右又好像醒了一次,就乾脆不要睡了,開車去到那裡.看到天還沒有亮,大約二,三百人在那裡一起禱告.他們有一個共同的關心,大部分都是自己的家人,兒女的前途,還有就是有疾病的,有各種疾病的家人和親戚,他們在那裡禱告.

在那邊事奉的時候,我們那個電台的台長,一位韓國牧師,他有一年去荷蘭參加一個國際巡迴佈道者的會議,是Billy Graham(葛培理)召集的,那是1986年.這位金牧師是講員之一,那一次他的題目是「我們所需要的復興」(The Evangelist and the Revival We Need),因為是巡迴佈道者的會議,他在那裡用韓國會議的特色來講了一篇信息.五個特色包括:第一,他說他們的教會是禱告的教會；第二,他們的教會是悔改的教會；第三,他們的教會是宣教的教會；第四,他們的教會是讚美唱詩的教會；第五,他們的教會是受苦的教會.當時他大約是這樣說的:「Korean church is marked as praying church, propagating church, purified church, praising church, and the persecuted church.」最後一點講到韓國教會的逼迫跟苦難.過去兩百年,我們中國的歷史裡面是苦難的歷史.韓國教會的歷史也同樣是苦難的歷史,但是兩者有一個很大的不同,那就是我們雖然被列強侵略,我們仍然是在有租界跟沒有租界部分的情況.但根據我在事奉時觀察,韓國教會在過去一百年,國家有半個世紀完全被另一個政權統治,就是被日本合併,那個用另一種字眼講是完全的滅亡,在世界的地圖裡面沒有了那個國家,但中國不是.這種民族的自尊完全地喪失,跟我們所經歷的,兩個比較以下,就是這個不同.這個我相信是每一年去韓國觀摩教會,觀摩他們禱告的,我們的,中國的,香港的,或是海外華人的基督徒,所不能帶回來的一個特色.

禱告,甚麼樣的人禱告呢?突然間發現你有癌症,本來你是很理性的,你有不能夠解決的疾病的時候,你會呼天.前幾天我看完一本書,在四川的成都買的,台灣作家柏楊的回憶錄,原來他是基督徒,其中他在最苦難的時候,在監獄的時候,在台灣的綠島(監獄)裡面,不能夠承受的那一種痛苦.當中他寫了一句話:「基督啊,求你救我!」甚麼樣的人禱告?就是找不到前途的人,在沒有盼望的時候.可能你能夠有的最後一句話就是「神啊,求你救我」.

看剛才那句經文「況且我們的軟弱」.第一,我們要留意軟弱這現實,在羅馬書裡面好幾次提到軟弱.當我們軟弱的時候,耶穌基督為罪人死,這是羅馬書第五章.在第八章開始的時候,他說,我們肉體的軟弱有聖靈幫助,這個軟弱的現實提出了救恩.福音的真理跟罪人的真理,是因為我們的軟弱.在羅馬書第五章,你的軟弱是跟罪人和罪連在一起.人是罪人.我們讀歷史的時候發現我們所認識的人,包括領袖,原來都有軟弱,虧欠,失誤,跌倒.中文大學最近出了一本書,「現代中國民間教會」,是美國一位歷史系的華人教授寫的,關於由太平天國到最近幾十年中國民間教會情況,其中有一段講到宋尚節博士.他是上一個世紀中國教會的領袖佈道家,完全為主燒盡,我們那麼的尊敬他.但是因為這本是歷史學家所寫的書,所以他寫得很客觀,優點跟缺點都寫.有一次宋尚節博士遇到交通意外,船沉了,要淹死的時候,他拼命地把自己的聖經跟自己寫的日記弄好,游泳逃了出來,卻忘了救他的太太.我相信如果他太太是香港人的話,回到岸裡面,可能會罵他一頓.人有軟弱,所以我們需要耶穌基督的救恩.就在羅馬書,我們看到了第一方面.

第二方面,就是保羅自己的經驗,他是在哥林多寫羅馬書.哥林多教會是保羅人生裡面最痛苦的教會,最不容易有成績的教會.我們讀新約聖經的時候知道,保羅大約寫了四封信給哥林多的教會.他在哥林多前書裡面說:「我在你們那裡,又軟弱,又懼怕,又甚戰兢.」(林前2:3)第二種軟弱是情緒上的,心靈上的壓力.我不知道今天這裡有沒有大學的學生,你要交功課,你要考試,一定是有壓力的.現在還好,現在是放暑假了,壓力已經過了,不過下學期還是會有壓力.如果你生病,你的家裡面有別的事情,通常這些事情是一起來的,不會給你足夠的時間,精神狀態最好的時候預備考試,很少有這樣的情況,將來做事情也是這樣.很奇怪,你很喜歡做的一件事情,偏偏是在你做那件事情之前發生別的事情,加上自己身體又不好,但是要熬夜,要這樣子把事情趕出來.精神上有壓力,這是人軟弱事實的第二方面.

第三方面,人人都會經歷的,那就是疾病,衰老,我們都要面對.為甚麼在醫院的臨終的那一科,很多人願意聽一聽教會的朋友們給他講福音,願意邀請教會的傳道人牧師來給他禱告?因為身體的軟弱讓他知道,自己不得不謙卑下來,接受軟弱的現實.回到我們本段經文的下文.你看,如果是看這裡第八章上文的一段,十八節「我想,現在的苦楚」.其實第八章是羅馬書裡面的巔峰,很好的一章,得勝的一章.剛才我們讀的時候到後面「得勝有餘」三十一節到三十九節,這是羅馬書的凱歌,是我們每一位讀羅馬書,讀聖經的弟兄姊妹都應該好好掌握的.但是在之前他講甚麼呢?我們與耶穌基督一同受苦,將來一同得榮耀,這是八章十七節.然後,接下來他就講幾次的苦難.我們看一下:「現在的苦楚」(八章十八節)；八章二十節另外一個用詞「服在虛空之下」；下一節八章二十一節「脫離敗壞的轄制」；八章二十二節一個重要的詞彙,關鍵詞,就是「歎息,勞苦」,然後八章的二十三節「自己心裡歎息」.你有留意十八節到二十五節的話,他的歎息包括了「世界萬物的歎息」,「伏在空虛之下」.第二,教會的歎息,那就是說弟兄姊妹基督徒神的兒女的嘆息.八章二十二節那裡說「歎息,勞苦」,我們的和合本有一個字沒有翻譯出來,那就是「生產之苦」.和合本的修訂本,原文聖經,英文聖經都有翻譯出來的.和合本的修訂本是這樣翻譯的:「一切受造之物一同呻吟,一同忍受陣痛.」這一段經文給我們聽到歎息,世間萬物的歎息,就連教會的弟兄姊妹也有歎息.然後,保羅講到我們的軟弱.

第二方面,軟弱的程度.軟弱到一個甚麼程度呢?「我們本不曉得當怎樣禱告」.很多解經家也有解釋,讓我為大家介紹幾位比較流行的解經家他們的解釋.第一位是美國教系統神學的約翰 and bull;莫維莊(John Murray).他說,這不是關於對禱告的正確態度有所不知,反而是對禱告最恰當的內容的無知,是對於我們的處境有甚麼迫切需要的無知.基本上呢,解經家在「怎樣」跟「甚麼」(how跟what)之間做一個選擇的時候,不是我們不知道怎樣去禱告,而是不知道我們要禱告甚麼.另外一位解經家鄧雅各(James D. G. Dunn),他講的解釋也是類似,他說並不只是他們要禱告祈求,但不知道如何表達；他們不知道神對他們的心意是甚麼,也不知道所在的社會的處境.回到宗教改革時期的神學家約翰 and bull;加爾文,他講得最簡單,最直接.四五百年了,他的解經書還被很多人參考.約翰 and bull;加爾文這樣說,保羅在此不過是想指出,我們向神祈禱的說話是盲目的,因為我們雖然經驗各種的苦難,我們的心卻是極其的困惑和迷惘,以至於我們不知道甚麼是當求的,甚麼事是對我們有益的.

弟兄姊妹,我猜想我們當中大部分是基督徒,或者有一些是慕道的朋友,可能有少部分是還沒有信主的,我不太知道,因為我只是幾年才有一次的時間來跟大家一起敬拜.如果我叫你問自己,「今天我最大的禱告是甚麼」,你有嗎?你知道你應該禱告的是甚麼嗎?你會不會說,「上帝啊求你讓我活到120歲,但是我的樣子仍好像是20歲的樣子.」這個是沒有意思的,卻有很多廣告公司,化妝品的公司可能用這樣的廣告,但這不是我們的禱告.可能你也會關心今天早上的新聞,在福建溫州開車的火車由於交通意外,六輛的車廂掉出了高架橋以外,三十多個人死亡.如果有你的親戚,你的朋友,你心靈會傷痛.我們會關心,但可能這不是你今天心裡最迫切的禱告.我們最迫切的禱告是甚麼呢?保羅說,人啊,你是很有限的,你是很軟弱的,你是需要救恩的,你是有精神壓力的,你是會衰老,得疾病的.「我們的軟弱有聖靈幫助,我們不知道怎樣禱告.」那個真正應該禱告的內容究竟是甚麼,我們不知道.很感謝主,我今天是從崇基學院來到大學,每次來都感到心裡很舒服.這麼美麗的校園環境,我相信正正是這樣的地方,幫助我們去認識一個社會,認識一個國家的將來,認識社會的將來,知道我們所求的是甚麼.在你小小的領域裡面,去讀的去研究的,你知道就是這一個,它有所貢獻,它的方向應該是這樣走,那個大概是你要禱告的內容.

過去三十多年,我們從事廣播.時代變了很多,單單講我們在錄音室所用的機器,剛開始的時候我們用的是盤帶,一條一條的磁帶,我們做剪接的時候是用刀片.接著有卡帶,然後有數碼卡帶(Digital tape),再有Mini Disk,還有Cartridge,一直輾轉到今天甚麼都不要的,就是一個電腦然後convert,送到發射站,甚至直接做直播就可以了.在香港直接做廣播,透過菲律賓或其他的發射站也好,可以同時地跟中國大陸裡的聽眾朋友交通交流.有很多的改變,很多的進步.下一步我們要做甚麼呢?良友電台在網上廣播,網上直播之外呢,要做新平台手機廣播,現在已經不算是很早了.我們要學手機廣播,要有相關的部門,都是很年青的人.手機廣播之後下一代又是甚麼呢?那麼多東西隨時地變,要隨時地觀察.可是有樣不變的,就是我們對所知道的,的確有限.

這裡保羅說,我們軟弱的程度是不知道自己要為甚麼禱告.我們看看馬丁 and bull;路德他怎麼解釋呢?他的解釋是最簡單,最形象化,平民百姓能聽懂的.他說我們禱告得太少,太低,神所給的呢,是太高,太好.好像一個孩子寫信跟爸爸要一個銀幣,但是他父親收到這封信,把信扔掉了.哎呀!這個孩子就為這件事暴躁起來,傷心難過,但是他不知道他爸爸可以給他許多的金銀.馬丁 and bull;路德就運用了一個例子,西庇太的兩個兒子跟耶穌基督祈求的時候,主耶穌跟他們說「你們不知道你們所求的是甚麼」.「況且我們的軟弱有聖靈幫助,我們本不曉得當怎樣禱告」這就是軟弱的程度.

最後今天的福音--軟弱時的安慰:聖靈幫助.聖靈怎樣禱告呢?「說不出來的歎息,替我們禱告」.「我們的軟弱有聖靈幫助」,「替我們禱告」.作者是保羅,他在哥林多寫羅馬書,羅馬書之前是哥林多後書,哥林多後書之前是哥林多前書.研究教會領袖的教授有一個人叫Bobby Clinton(博比 and bull;克林頓),他說保羅如果要在他一生裡面要辭職不幹了,要放棄了,最痛苦的時候就是在寫哥林多後書以後,也就是寫羅馬書之前,那個是保羅人生最困難的時候,大約是五十歲中間,五十五歲左右.因為哥林多的情況,令保羅非常痛苦,常常流淚.但是他怎麼能夠走過去呢?他發現他的軟弱本身就是倚靠神的動力.禱告的能力從哪裡出現?發現自己最軟弱的那一刻.發現自己最沒有能力的那一刻,發現自己解決不了的時候,「神啊,幫助我!聖靈啊,求你進入我的生命,進入我的世界,你幫助我.」聖靈幫助,這是軟弱時的福音,軟弱時的安慰,聖靈本身就是能力.所以保羅發現,我甚麼時候軟弱,甚麼時候就剛強.因此我誇我的軟弱,讓基督的能力覆庇我,這是基督教的信息.

「聖靈的幫助」中的「幫助」這兩個字是甚麼意思?我們看解經家告訴我們--很形象的--A. T. Robertson(A. T. 羅拔遜)說,伸出一隻手,拉住另一隻手,每當我們軟弱的時候,聖靈就伸出手拉我們一把,而且在時間上都是恰到好處,不會太晚.聖靈無所不在,祂願意幫助你和我.為甚麼有些人得到幫助,另一些人又沒有得到幫助呢?差別就是有的人有把手伸出來,他知道他需要幫助；有的人不需要伸出手,根本不相信這個世界有聖靈,根本不相信聖經是神的話,他用他自己的方法,甚至呼吸到最後一口氣,這就是差別.另外一位我們華人的聖經學者鮑會園牧師說,聖靈和我們一起,從旁邊幫助我們,不是代替我們.聖靈不代替我們做我們要做的事,祂在旁邊幫助我們,和我們一起做,我們仍然在這裡做,聖靈推你一把,祂幫助你,「況且我們的軟弱有聖靈幫助」.

弟兄姊妹,各位朋友,你今天願意讓聖靈幫助你嗎?你願意的時候,禱告的能力就彰顯出來了.今天早上,我們聽到讀詩篇119篇136節.我想了一下136節那裡,作者說我流眼淚成河,這是大衛的詩.因為周圍的人不遵守神的律法,那種流眼淚成河,詩人誇張的作法,使詩篇的作者寫了最長最長的一篇詩:119篇.但是你知道在詩篇裡面,最早的第3篇裡面,大衛王有一句話,他說「有許多人議論我說」,議論甚麼呢?「他得不著神的幫助.」這是一個最痛苦的批評.可是大衛在下一句說,「但你耶和華是我四圍的盾牌．是我的榮耀,又是叫我抬起頭來的」.軟弱者的安慰是聖靈的幫助.打開你的雙手,聖靈願意幫助你.

我們做一個總結.八章二十六節,二十七節,聖靈用嘆息幫助我們禱告.整段的上下文給我們看到,在苦難的人生當中,萬物的嘆息當中,教會的弟兄姊妹的嘆息當中,有一個盼望.盼望的字眼從下面這幾句話,這幾個用詞可以看到.八章十八節「將來的榮耀」；十九節「切望等候」；二十一節「指望脫離」；二十一節「享受神兒女自由的榮耀」這意思；將來二十三節「等候得著身體得贖」；二十四節「得救在乎盼望」八章二十四至二十五節.你看有多少盼望?我讀給你聽:「我們得救是在乎盼望．只是所見的盼望不是盼望．誰還盼望他所見的呢.但我們若盼望那所不見的,就必忍耐等候.況且我們的軟弱有聖靈幫助.」所以這一段經文是兩條線交叉的,今天軟弱痛苦的現實,跟今天等候榮耀的盼望.兩條線同時進行,聖靈就在這個交叉點,給你,給我盼望.基督徒跟非基督徒最大的差別是甚麼?不是我們沒有難處,壓力,疾病,而是我們有永遠不朽壞的盼望.

結束的時候,我想介紹上個月看到的一份雜誌的文章,不是基督教的,是Time(時代週刊)的.在六月初,報導了一個很有意思的主題文章.因為我們不是幹這一行,所以我看了特別有興趣.他說人有樂觀的傾向(optimism bias),看了文章大概就是說,心理學家做很多實驗,它對樂觀的結果,是在大腦的結構裡面,有樂觀的元素使人樂觀.都說離婚率愈來愈高,但訪問很多剛結婚的人,覺得自己會否離婚,大部分剛結婚的人都不相信他會離婚.心理學家在最近的一項研究中發現,當人吸取新的資訊的時候,大腦的神經元能夠一絲不苟地,把人樂觀的正面信息完全吸收進來到你的腦袋裡面.但對於那些突如其來而又不合我們心意,反面的,不好的資訊,大腦的神經元卻沒有辦法容納它.整篇文章從稍為科學的一個角度講,人有一種本能,那就是樂觀.各位弟兄姊妹,你跟我是神所造的,希望你跟我都是正常,像這篇文章所講的,我們都是樂觀的.為甚麼呢?25節說我們有盼望,所以忍耐等候,況且我們的軟弱有聖靈幫助.打開我們的心,聖靈啊,來幫助我,我跟你一起走.禱告的能力.我們一起低頭做一個禱告.

「懇求耶穌基督的靈藉著羅馬書第八章二十六節二十七節的經文,再次教導我們,開導我們的心,讓我們裡面的生命,讓我們的信仰能夠健康長大.如果我們在座的弟兄姊姊朋友們,心中有切願的,求聖靈幫助,以你的嘆息說不出來的言詞替我們禱告,讓我們在生活當中剛強,讓我們往前面走的時候有美好的見證,求主幫助我們.謝謝你話語的應許,給我們力量.這樣的禱告是奉主耶穌基督聖名,阿們.」



\newpage

\section{禱告初階}
\label{sec:112}
\textbf{賀志勇博士 - 禱告初階}
\newline
\newline
\hyperref[sec:111]{< < < PREV SERMON < < <}
~
\hyperref[sec:toc]{[返目錄]}
~
\hyperref[sec:113]{> > > NEXT SERMON > > >}


日期:2011年8月21日

主題:禱告初階

讀經:創世記28:10-22節

證道:賀志勇博士

一,    引言:人的盡頭,禱告的開始很多年前,美國華盛頓州的聖海倫火山爆發過一次,當時有一個記者去拍攝.突然間火山開始噴發,爆出來很多的濃煙跟岩漿,嚇得記者使勁跑.但他跑的時候忘記關掉錄影機了,所以他跑的過程全部錄下來了.每一句話都在錄影機裡聽得清清楚楚.起初的時候聽見他喘氣,沒多久就聽見他喊道"上帝啊上帝,求你拯救我,求你拯救我,周圍好熱啊,好黑啊.......我快要死啦.你幫我啊,上帝,你幫我啊." 這個故事讓人看到一個事實:當人在危險的時候,一切保障一切安全都失去的時候,人的盡頭,就開始懷念上帝了；就開始禱告了.換句話說,為什麼很多人平時不會想到宗教信仰呢?為什麼人沒有想到要禱告呢?很簡單,因為他們覺得沒這個需要.當保障仍然在我們手裡的時候,當安全仍然在我們手裡的時候,我們為什麼要求一個看不見,摸不著,都不知道是真是假的上帝?但是當他們發覺,這些保障不管用,這些安全沒有了的時候,他們就會想,冥冥中或者有個上帝可以幫我--似乎這就是禱告的開始.

但是,如果我們發現有人真的是因為這樣去禱告時,我們很可能心底裡說,我們會輕視它,會譏笑它是"臨急抱佛腳."這種做法,不是一個虔誠的人應該做的,有事就求,沒事就放一邊,這絕對不是一個敬虔人的作為.況且我們會覺得,這樣的祈禱太自私,這樣的祈禱太自我中心、太功利、太利用上帝啦.好像去黃大仙燒香一樣,這絕對不是禱告的好榜樣.所以現在我們發現有兩種做法,一種是有需要才求,才禱告；另一種看法是認為,這樣的禱告未免太自私太功利,向神求東西才去禱告是不虔誠的表現.

其實看深一層,在骨子裡面,這兩種看起來完全不同的看法,裡面包涵著同一種態度,就兩個字:驕傲.第一,當一個人覺得自己夠好,夠厲害,夠安全,他就不需要禱告時,這是自我中心,是驕傲,因為他以為自己可以搞得定；第二,以為自己這樣禱告是不夠好,不夠純正,動機不夠虔誠,沒資格祈禱,所以不去祈禱時,其實裡面同樣包含著一個驕傲的問題.今天我就不分析第一種態度了,因為對於基督徒來講,第二種態度更隱蔽,更具有危害性.我今天想講講的就是第二個問題,什麼人才有資格祈禱?什麼叫合乎聖經的禱告(或者說禱告初階)呢?讓我們一起來看看剛才讀那段聖經,我們看看雅各是怎麼禱告的.

二,解經:什麼是禱告初階?(1)雅各的禱告 經文裡頭雅各這樣禱告說:上帝啊,如果你給飯給我吃,給衣服給我穿,給房子我住,這些東西都給我了,我肯定答應,給回十分之一給你.這就好像我同鄧博士講,鄧博士,你給我100塊,我肯定答應給回10塊給你.這就是雅各的禱告!這樣的禱告是不是很離譜啊?有沒有搞錯,這樣子還不是利用神?討價還價,自私自利.這樣的禱告,絕對不是我們信徒禱告的好典範、好榜樣.這簡直是反基督教教育的最佳題材,我都搞不懂聖經為什麼要記錄雅各的禱告.不過後來想一想,我發現我以前禱告不是差不多!我以前什麼時候禱告最勤快啊?就是我坐飛機出去旅行的時候.那個時候我最自覺,一上飛機坐好了,我就禱告,求神保守飛機師睡得好不要打瞌睡；沒有跟太太吵架,情緒穩定；求神保守天氣好,不要被雷劈,信號系統失靈；禱告完了,就開始看雜誌.然後等到飛機快要降落時,又開始禱告,求神保守降落平安,不要被追尾.大家看看,其實我以前禱告跟雅各是差不多的.

(2) 摩西的禱告但是我們再看看聖經,我們會發現,那些聖經裡的信心偉人,你同我在主日學裡學過的、非常非常尊敬的信心偉人、信心英雄,全部是這樣祈禱的.不信你看看!摩西是這樣祈禱的,大衛也是這樣祈禱的,以利沙、約書亞、路得,還有以利亞,甚至是亞伯拉罕,都是這樣祈禱的.你不信,我給點例子給你看看.摩西在民數記11章11節到12節,你看他怎樣禱告的--``11摩西對耶和華說:"你為何苦待僕人?我為何不在你眼前蒙恩,竟把這管理百姓的重任加在我身上呢?12, 這百姓豈是我懷的胎,豈是我生下來的呢?" 你看看,就好像你吩咐你孩子去洗碗,他卻強嘴說:"阿爸,為什麼要我去洗碗啊?今天又沒有輪到我洗碗,怎麼要我去洗碗!飯又不是我一個人吃的,碗又不是我一個人用的,為什麼你不叫阿姐去洗呢?"非常非常幼稚.摩西不就是這麼幼稚? 或者有人晚上打電話給我,問:"傳道人,可不可以探訪一下我啊?" 於是我就禱告說:"神啊神啊,為什麼你要我去探訪這個人呢?這個人又不是我生的,又不是我養的,為什麼你這個時候要我去呢?"你說我要這樣禱告,幼稚不幼稚呢?當然,不但幼稚,還很自私.(3)以利沙和大衛的禱告再看看以利沙,列王記下2:24節,以利沙更加離譜了.以利沙就因為那些小孩子笑他是光頭佬,啊光頭佬啊；於是他就奉耶和華的名祈禱,咒詛他們,於是就有兩隻母熊從林中出來,撕裂了四十二個童子.我們讀到這裡,我們想,有沒有搞錯啊,有這樣去禱告的麼?這樣的禱告還不離譜?

還有好多例子,你們再來看看大衛.詩篇第140篇10節,大衛怎麼祈禱呢? 他對神說:"願炭火落在他們身上,願他們被丟在火中,拋在深坑裡,不能再起來."這是祈禱哦!我不知道你們中間有沒有人為老闆這樣祈禱的:願炭火落在他們身上,願他們被丟在火中,拋在深坑裡,不能再起來. 我想,大衛祈禱的時候肯定握緊拳頭,摩拳擦掌去禱告啊.其實了,這樣的例子多不勝數.你打開聖經看看,詩篇啊,舊約啊,甚至是新約,你都可以看到好多這樣的祈禱.雅各已經算好的啦,給回十分之一給上帝.   你發現,原來在聖經裡面,有好多好多這樣的禱告,簡直成為聖經裡頭禱告的某種典範,如果有人問我說,這算什麼典範啊?我想我會把它命名為禱告初階.

對於我們來講,這樣的祈禱好難接受--太血氣,太不屬靈,太自私,太幼稚.我們有一種錯覺.什麼人才有資格祈禱啊?不自私、沒功利,非常虔誠的人,才有資格祈禱.我們的禱告內容,一定要是無私的,一定要是敬虔的,屬靈的；求屬世的事,就是不好的禱告,不像樣.求上帝處理你的老闆,讓炭火落在他的頭上,這就不對,就不蒙悅納.所以你要祈禱說:上帝啊,求你幫助我的老闆良心發現.  你的祈禱一定要非常漂亮,非常屬靈.一定要沒有私心.明明很生氣,卻要對上帝說,上帝啊,我不生氣啦,求你讓我祝福他吧,為他禱告吧.我們以為這樣就算敬虔,這樣才算一個合格的禱告.如果我們這樣看祈禱的話,就可能大錯特錯了.禱告初階,不單止是禱告的開始,也是禱告的重心.耶穌基督在主禱文中教導我們,我們可以為每天的飲食禱告.你有沒有為你每天的麵包禱告了?你有沒有禱告說,賀博士,你不要講得太長啊,我胃不好,想早點去飲茶了.--你有沒有為這個祈禱過呢?你可以這樣禱告的.

三, 應用:禱告初階的三種特色或者我們要界定一下,什麼叫禱告初階?剛才我們讀的經文,恰恰是禱告初階的典型例子.我們看看,雅各的禱告有三個特色:

第一個特色:體驗、承認自己無助無援,非常弱小雅各為什麼會流落到這步田地?是因為他用計騙了哥哥以掃長子的名分,還有父親的祝福,以掃要殺他,他只好一個人出來逃命.他自小嬌生慣養,被老媽呵護慣了,從來沒有離開過家.現在出來逃命,逃到荒山野嶺的伯特利,四周黑漆漆,曠野還有野獸,雅各這個時候是最無援最無助的時候.就在這樣的曠野裡,雅各做了這樣的禱告.他現在孤單單一個人在危險的曠野,連個落腳的地方都沒有,非常淒涼,這種心情,可能出去海外、或者來香港讀書工作的人,會有很深的體會.如果你現在一個人,半夜三更還在西貢的深山裡,是個什麼滋味呢?沒安全,沒保障,生命受到威脅,無援無助,你會怎麼樣?--人的盡頭,神的開始.只有當我們體驗並承認自己無援無助的時候,承認自己貧乏的時候,我們就開始真正的禱告了.

耶穌基督在登山寶訓裡頭說:靈裡貧窮的人有福了,哀慟的人有福了.一個人之所以不禱告,是以為自己可以搞的定,不需要去求神.我們很容易就"有事鐘無鹽,沒事夏迎春."想想看,我們每一個人,同雅各相差不是好遠.我們都是靈裡貧窮的人,好像保羅所說,"立志行善由得我,作不作得出來就由不得我."我不想做的壞事我偏偏去做,我想做的好事卻偏偏做不好.保羅說"我真苦啊."但是我們沒有保羅這麼大膽,我們不敢告訴別人聽,我們要死頂,硬撐,明明心裡苦,還要在人前裝笑臉.什麼人下水容易淹死呢?不是那種一落水三尺就喊救命的.到膝頭水嘛.那種人一定不會死.最怕是那種腳抽筋還要說"不怕不怕",死撐的,怕丟臉的,怕人家小看的人,這種人才是最危險.

一個人如果知道自己很軟弱,然後肯去禱告,這就有福了.但人很奇怪,有些朋友跟我說,我現在不洗禮,因為我還不夠虔誠,沒資格洗禮.等我變好一點,我就有資格洗禮了.不然我做了基督徒,就會把上帝的名聲帶壞了.我就問他,一個病人會不會說"我暫時不去看醫生,等我病好一點再來看吧,費事把醫生的名聲搞壞了".不會嘛!世界上沒有人會說等我好點再去看醫生吧.一個人為什麼要看醫生?他不會說,醫生,我沒事,我是來探訪你的.為什麼我們要到上帝面前禱告啊?耶穌基督說他來不是來召義人的.你以為夠敬虔才來禱告,或者你要做一個最敬虔的禱告,可是神並不要你來表現你的敬虔.神希望通過你的禱告,讓你看到自己的軟弱、無助,就好像雅各看到自己的軟弱無助一樣.雅各的性格就是什麼都要抓住,雅各這個名字的意思就是抓住,什麼都要抓在手裡頭.但是當雅各發現他抓住的所有東西都是肥皂泡的時候,他就在上帝面前開始祈禱.禱告初階的第一個特色:你看沒看到自己的軟弱? ``我沒事啊,我自己可以啊.我搞得定啊."這種人最容易出事.禱告初階的第一步,認清自己的弱小,向上帝求助.第二個特色:體驗上帝的活現什麼是上帝活現?雅各做了一個夢,夢見上帝向他說話.當雅各醒的時候,就怕了,大衛震驚,就說: ``上帝既然在這裡,我居然不知道."

禱告初階的第二個特色,就是當神是活生生的神,就在這裡,當神有到.神不僅僅在教堂裡才現身.神是在任何時候任何地方都掌權的.在我們等車等得不耐煩的時候,在我們開車找不到停車位的時候,在我們管教小孩的時候,在我們回大陸被人插隊打尖就要罵人的時候,或者都不是這些,在我們開心得不得了的時候,我們都可以禱告,因為神就在我們身邊陪伴著我們.很多時候我們不禱告,是因為我們覺得神離我們很遠.哎,神不會管我這些小事的!那你錯了.神是活生生的神,"上帝既然在這裡,我居然不知道."

第三個特色:全盤向神傾述我們看看雅各的禱告,他要耶和華給他飯吃,給他衣穿,保護他,讓他平平安安回到父親的家裡,他就以耶和華為神.雅各就是這樣:全盤向神傾述.祈禱,希臘文的意思就是向神說話.祈禱不是演講.啊,我今天要跟神講一篇很漂亮很漂亮的禱文--你以為神是你演說的聽眾麼?祈禱就是你把你心裡所想的,毫無保留向神傾述.

耶穌教導我們,要學像小孩子的樣式.我們聽完這句話,很快就會神話了小孩子--小孩子是不自私的,是很傻很天真的.我想問一下,我們中間有多少人認為小孩子是不自私的?是這麼認為的人可以舉手.我們中間做過父母的人就知道,小孩子,絕對不是我們認為的那麼純真,有時候還很自私,你讓他同別人分享一點東西,比登天還難.看到人家手裡的玩具,一把搶過來,"這個是我的,"玩完之後,情願丟在一邊也不給別人玩.

小孩子是這個樣子,那為什麼耶穌還要我們學習小孩子的樣式?我想,不是因為他們敬虔,他們無私；而是因為,他們不虛偽,他們一點都不虛偽.他們想說什麼就說什麼,想搶什麼就搶什麼,不分好歹,不問是非.不會裝扮成另外一個不是他的人.他們就是這樣啊!他們只會把心裡面所求所想講出來,沒有深思熟慮,沒有經過裝扮,沒有半點掩飾,是就是,不是就不是.我記得我認識的一個人,有一次他到教堂這邊來玩,看到有伴娘跟新娘在拍婚紗照,他就大聲說,"左邊那個姐姐好漂亮啊!"左邊那個伴娘就很開心,然後他又說:"右邊那個姐姐好難看啊."結果右邊那個新娘就氣死了.小孩子就這樣,他不管大人的人情世故,虛偽圓滑,什麼都講得出來的.我們大人就不是了,我們要留有餘地,我們要扮扮樣子,明明不感興趣對方講什麼,還要睜大眼睛說我會為你祈禱的.

禱告初階,就是要我們不要扮樣子,不要裝成自己更天真,扮成自己更敬虔,而是向神坦誠地講出自己內心所想.禱告初階,就是我們不要以為這些是瑣碎的、自私的要求,就不說給上帝聽.其實,任何事情都可以祈禱.等巴士可以祈禱,趕面試卻堵車了可以祈禱,不見了鑰匙可以祈禱,同太太吵架了可以祈禱.可能你一聽我講道就要釣魚了、要睡覺了 ,你也可以為此祈禱.芝麻綠豆的事,都可以祈禱.哪怕是你對別人憤怒了、傷心了、疲倦了,你都可以祈禱.你以為這些都是芝麻綠豆,上帝不會理這些事,上帝只會理大事,那你就錯了.可能你想買件漂亮衣服,或者買個Iphone5,你又不捨得,你都可以去祈禱.不要以為這些祈禱太露骨,太不屬靈.這些小事,你都可以禱告詢問上帝.上帝連你每根頭髮都數過,你不用擔心他會煩躁.

有時候,我們的禱告可能會帶著不好的動機,可能我們的禱告中有溫柔但也有怨恨,但是上帝會大度去聽我們的.凡求告祂的,祂都會垂聽.我們不要以為上帝這麼小氣,他不會接納我們的瑕疵.他會的.

四,結論:人的盡頭,神的開始弟兄姊妹,不管我們以前是怎樣禱告的,不管我們喜不喜歡禱告,這都不要緊.今天,神的話臨到了我們,他要我們去學習雅各,第一,看清楚自己軟弱無助,第二,體驗到神的活現；第三,坦誠地、毫無保留地向神全盤傾述.哪怕你很不想禱告,或者你想不到怎麼禱告好,你都可以跟神說:神啊,我今天很不想禱告,或者不知道怎麼禱告好,求你原諒我吧. 這樣都可以.因為只有當我們在神面前全然坦誠時,我們就真正明白了:人的盡頭,神的開始.



\newpage

\section{當你走到無力}
\label{sec:113}
\textbf{林豪恩先生 - 當你走到無力(When You Down and Troubled)}
\newline
\newline
\hyperref[sec:112]{< < < PREV SERMON < < <}
~
\hyperref[sec:toc]{[返目錄]}
~
\hyperref[sec:114]{> > > NEXT SERMON > > >}


列王記上 19:1-18
\newline
\begin{longtable}{cl}
\hline
\hline
章節 & 經文 (和合本修訂版)\\
\hline
19:1 & \begin{tabularx}{0.7\textwidth}{X} 亞哈把以利亞一切所做的和他用刀殺眾先知的事都告訴耶洗別。 \end{tabularx} \\ \\ \relax
19:2 & \begin{tabularx}{0.7\textwidth}{X} 耶洗別就派使者到以利亞那裡，說：「明日約這時候，我若不使你的性命像那些人的性命一樣，願神明重重懲罰我。」 \end{tabularx} \\ \\ \relax
19:3 & \begin{tabularx}{0.7\textwidth}{X} 以利亞害怕，就起來逃命，到了猶大的別是巴，把僕人留在那裡。 \end{tabularx} \\ \\ \relax
19:4 & \begin{tabularx}{0.7\textwidth}{X} 他自己在曠野走了一日的路程，來到一棵羅騰樹下，就坐在那裡求死，說：「耶和華啊，現在夠了！求你取我的性命吧，因為我不比我的祖先好。」 \end{tabularx} \\ \\ \relax
19:5 & \begin{tabularx}{0.7\textwidth}{X} 他躺在羅騰樹下睡著了。看哪，有一個天使拍他，對他說：「起來吃吧！」 \end{tabularx} \\ \\ \relax
19:6 & \begin{tabularx}{0.7\textwidth}{X} 他觀看，看哪，頭旁有燒熱的石頭烤的餅和一壺水，他就吃了喝了，又再躺下。 \end{tabularx} \\ \\ \relax
19:7 & \begin{tabularx}{0.7\textwidth}{X} 耶和華的使者回來，第二次拍他，說：「起來吃吧！因為你要走的路很遠。」 \end{tabularx} \\ \\ \relax
19:8 & \begin{tabularx}{0.7\textwidth}{X} 他就起來吃了喝了，仗著這飲食的力走了四十晝夜，到了神的山，就是何烈山。 \end{tabularx} \\ \\ \relax
19:9 & \begin{tabularx}{0.7\textwidth}{X} 他在那裡進了一個洞，在洞中過夜。看哪，耶和華的話臨到他，說：「以利亞，你在這裡做甚麼？」 \end{tabularx} \\ \\ \relax
19:10 & \begin{tabularx}{0.7\textwidth}{X} 他說：「我為耶和華－萬軍之神大發熱心，因為以色列人背棄了你的約，毀壞了你的壇，用刀殺了你的先知，只剩下我一人，他們還要追殺我。」 \end{tabularx} \\ \\ \relax
19:11 & \begin{tabularx}{0.7\textwidth}{X} 耶和華說：「你出來站在山上，在耶和華面前。」看哪，耶和華從那裡經過。在耶和華面前有烈風大作，山崩石裂，耶和華卻不在風中；風後有地震，耶和華也不在其中； \end{tabularx} \\ \\ \relax
19:12 & \begin{tabularx}{0.7\textwidth}{X} 地震後有火，耶和華也不在火中；火以後，有輕微細小的聲音。 \end{tabularx} \\ \\ \relax
19:13 & \begin{tabularx}{0.7\textwidth}{X} 以利亞聽見，就用外衣蒙臉，出來站在洞口。聽啊，有聲音向他說：「以利亞，你在這裡做甚麼？」 \end{tabularx} \\ \\ \relax
19:14 & \begin{tabularx}{0.7\textwidth}{X} 他說：「我為耶和華－萬軍之神大發熱心，因為以色列人背棄了你的約，毀壞了你的壇，用刀殺了你的先知，只剩下我一人，他們還要追殺我。」 \end{tabularx} \\ \\ \relax
19:15 & \begin{tabularx}{0.7\textwidth}{X} 耶和華對他說：「去吧，從原路回去，往大馬士革的曠野去。到了那裡，你要膏哈薛作亞蘭王， \end{tabularx} \\ \\ \relax
19:16 & \begin{tabularx}{0.7\textwidth}{X} 又膏寧示的孫子耶戶作以色列王，並膏亞伯‧米何拉人沙法的兒子以利沙作先知接續你。 \end{tabularx} \\ \\ \relax
19:17 & \begin{tabularx}{0.7\textwidth}{X} 將來逃過哈薛之刀的，必被耶戶所殺；逃過耶戶之刀的，必被以利沙所殺。 \end{tabularx} \\ \\ \relax
19:18 & \begin{tabularx}{0.7\textwidth}{X} 但我在以色列中留下七千人，是未曾向巴力屈膝，未曾親吻巴力的。」 \end{tabularx} \\ \\
[1ex]
\hline
\hline
\end{longtable}
日期:2011年8月28日

題目:當你走到無力(When You Down and Troubled)

經文:列王記上19:1-18

證道:林豪恩先生

一．引言公元前202年,西楚霸王項羽敗走烏江邊,遇到烏江亭長.亭長對項羽說:「江東雖然小,但土地縱橫各有一千里,民眾有幾十萬,也足夠稱王啦.希望大王快快渡江.現在只有我這兒有船,漢軍到了,沒法渡過去.」項羽笑了笑說:「天之亡我,我何渡為!且籍與江東子弟八千人渡江而西,今無一人還,縱江東父兄憐而王我,我何面目見之?縱彼不言,籍獨不愧于心乎?」項羽想到自己曾經「力可撥山,氣可蓋世」,現在落難江邊,大勢已去,無可奈何,於是舉劍自刎於烏江邊.

今天舊約經課提及的以利亞先知,也可以用「力拔山兮氣蓋世」來形容.他曾經在迦密山上,一個人打敗四百五十個巴力先知,舉國震撼.轉過頭來,他向耶和華呼喊:「求你取我的性命吧!」偉大先知何以突然如此無可奈何?充滿能力的先知何以突然生無可戀?在古代中國的項羽走到無力時在江邊遇上烏江亭長,在古代近東的以利亞則在山上遇到耶和華,耶和華如何回應他呢?耶和華對以利亞的回應又給予我們甚麼啟發呢?

二．經文

19 : 1亞哈把以利亞一切所做的和他用刀殺眾先知的事都告訴耶洗別.

19 : 2耶洗別就派使者到以利亞那裡,說:「明日約這時候,我若不使你的性命像那些人的性命一樣,願神明重重懲罰我.」

19 : 3以利亞害怕,就起來逃命,到了猶大的別是巴,把僕人留在那裡.

19 : 4他自己在曠野走了一日的路程,來到一棵羅騰樹下,就坐在那裡求死,說:「耶和華啊,現在夠了!求你取我的性命吧,因為我不比我的祖先好.」

19 : 5他躺在羅騰樹下睡著了.看哪,有一個天使拍他,對他說:「起來吃吧!」

19 : 6他觀看,看哪,頭旁有燒熱的石頭烤的餅和一壺水,他就吃了喝了,又再躺下.

19 : 7耶和華的使者回來,第二次拍他,說:「起來吃吧!因為你要走的路很遠.」

19 : 8他就起來吃了喝了,仗著這飲食的力走了四十晝夜,到了神的山,就是何烈山.

19 : 9他在那裡進了一個洞,在洞中過夜.看哪,耶和華的話臨到他,說:「以利亞,你在這裡做甚麼?」

19 : 10他說:「我為耶和華－萬軍之神大發熱心,因為以色列人背棄了你的約,毀壞了你的壇,用刀殺了你的先知,只剩下我一人,他們還要追殺我.」

19 : 11耶和華說:「你出來站在山上,在耶和華面前.」看哪,耶和華從那裡經過.在耶和華面前有烈風大作,山崩石裂,耶和華卻不在風中；風後有地震,耶和華也不在其中；

19 : 12地震後有火,耶和華也不在火中；火以後,有輕微細小的聲音.

19 : 13以利亞聽見,就用披風蒙臉,出來站在洞口.聽啊,有聲音向他說:「以利亞,你在這裡做甚麼?」

19 : 14他說:「我為耶和華－萬軍之神大發熱心,因為以色列人背棄了你的約,毀壞了你的壇,用刀殺了你的先知,只剩下我一人,他們還要追殺我.」

19 : 15耶和華對他說:「去吧,從原路回去,往大馬士革的曠野去.到了那裡,你要膏哈薛作亞蘭王,

19 : 16又膏寧示的孫子耶戶作以色列王,並膏亞伯‧米何拉人沙法的兒子以利沙作先知接續你.

19 : 17將來逃過哈薛之刀的,必被耶戶所殺；逃過耶戶之刀的,必被以利沙所殺.

19 : 18但我在以色列中留下七千人,是未曾向巴力屈膝,未曾親吻巴力的.」

三．釋經與應用1．當以利亞走到無力

19 : 1亞哈把以利亞一切所做的和他用刀殺眾先知的事都告訴耶洗別.

19 : 2耶洗別就派使者到以利亞那裡,說:「明日約這時候,我若不使你的性命像那些人的性命一樣,願神明重重懲罰我.」

19 : 3以利亞害怕,就起來逃命,到了猶大的別是巴,把僕人留在那裡.

19 : 4他自己在曠野走了一日的路程,來到一棵羅騰樹下,就坐在那裡求死,說:「耶和華啊,現在夠了!求你取我的性命吧,因為我不比我的祖先好.」

很多人都不明白,為甚麼有能力獨自戰勝四百五十個巴力先知的以利亞,轉眼間竟然落荒而逃.迦密山上的英雄何以害怕?呼火喚雨的先知何以逃命?

有人批評以利亞沒有信心,有人指責這位先知失職逃跑.不過,持平一點來看,世上豈真有長勝將軍嗎?股神神機妙算,但也有「損手」的時候；球王百發百中,但卻有「上帝之手」的時候；「天王天后」如何努力塑造自己清純的形象,但也有被手提電腦出賣的時候；何況今次以利亞遇到的挑戰與上次不同了,今次的對手與上次的對手不可同日而語.上次,以利亞以耶和華上帝先知的身份與四百五十個巴力先知對決,雖然人數對比懸殊,但性質上仍然是先知與先知之間的「比武」,屬於宗教界的範疇.現在,以利亞的對手是王后耶洗別,等於政權及其軍隊.這不再是宗教界的事情,而是政治及軍事範疇的力量.要先知一個人獨自迎戰整個政權及軍隊,難道我們以為以利亞刀槍不入嗎?王后耶洗別動用整個國家的政治及軍事力量來對付以利亞,簡直就是「大石壓死蟹」,以利亞感到害怕而逃命,也是人之常情吧!

不久之前,有鐵路意外在某城市發生.引起人們注意的,不單單是意外本身,媒體的激烈反應更令人們感到驚訝.媒體展示出對生命的珍惜,對公義的追求,以及對責任的追究,燃起人們對媒體及對社會進步的盼望.不過,當媒體的壓力不是來自鐵道部門,而是來自最高層權力階層的時候,在「大石壓死蟹」的狀況下,無可奈何,也要收起文稿,重新排版.

新聞工作者甘心嗎?媒體甘心嗎?在龐大的壓力及威嚇之下,不甘心又可以怎樣呢?以利亞甘心落荒而逃嗎?以利亞甘心放棄先知的召命嗎?就算以利亞不甘心,又可以怎樣呢?

人生走到無力,不一定因為走路太長而雙腳疲勞,可能是因為對手太強,壓力太大；也可能是因為遊戲規則太過不公平,無論你如何努力,都不可能在這場遊戲中得到甚麼.還記得《少林足球》電影中其中一場足球比賽,對陣球員都帶著武器上陣,從褲袋跌出土巴拿的球員若無其事地說:「作為一個汽車維修員,隨身攜帶士巴拿,是很正常的事.」期望公平公正地踢足球的球員,以為可以透過技術取勝的球隊,怎能不感到無力呢?怎能不氣餒呢?這種「大石壓死蟹」的無力感,這種不公平狀態下無可奈何的感覺,相信處身於類似狀況中的人們是能夠體會的.

以利亞甘心嗎?從以利亞對耶和華上帝所說的話,他似乎不甘心.他向耶和華上帝說:「耶和華啊,現在夠了!求你取我的性命吧」.他的意思是,如果遊戲是要這樣玩的話,我受夠了,我不玩了.我寧死都不繼續玩這樣的遊戲!

這場遊戲出了甚麼問題,以致先知寧死也不願意繼續玩下去呢?在道出詳細原因之前,以利亞指出,這場遊戲已經令他的自我形象及尊嚴都大受打擊.以利亞是有使命和理想的人,現在被逼成為逃兵,內心的愧疚和掙扎非言語所能形容,只能向耶和華呼喊:「求你取我的性命吧,因為我不比我的祖先好.」

王后耶洗別動用國家機器追殺以利亞,究竟對以利亞產生怎樣的打擊呢?且聽以利亞自己怎樣說.

19 : 9他在那裡進了一個洞,在洞中過夜.看哪,耶和華的話臨到他,說:「以利亞,你在這裡做甚麼?」

19 : 10他說:「我為耶和華－萬軍之神大發熱心,因為以色列人背棄了你的約,毀壞了你的壇,用刀殺了你的先知,只剩下我一人,他們還要追殺我.」

19 : 13．．．聽啊,有聲音向他說:「以利亞,你在這裡做甚麼?」

19 : 14他說:「我為耶和華－萬軍之神大發熱心,因為以色列人背棄了你的約,毀壞了你的壇,用刀殺了你的先知,只剩下我一人,他們還要追殺我.」 耶和華兩次問他:「以利亞,你在這裡做甚麼?」他兩次回答:「我為耶和華－萬軍之神大發熱心,因為以色列人背棄了你的約,毀壞了你的壇,用刀殺了你的先知,只剩下我一人,他們還要追殺我.」我們有否從他的控訴感受到剛發生的事情令以利亞內心充滿鬱結呢?

第一,以利亞戰勝四百五十個巴力先知,為國家斬妖除魔,為以色列帶來雨水,舒緩了三年的旱災,但國王沒有肯定和感激他,反而追殺他.白樺在《苦戀》中的對白或者可以道出以利亞心中的鬱結:「你愛國家,國家愛你嗎?」

第二,「我為耶和華－萬軍之神大發熱心」,竟然落得「只剩下我一人,他們還要追殺我」的結局.是否不應該如此相信耶和華呢?還是不應該如此熱心呢?如果不是,為何耶和華不保護祂的先知呢?

第三,「以色列人背棄了你的約,毀壞了你的壇,用刀殺了你的先知」.耶和華到底有沒有看見人民背信棄義呢?為甚麼容許這樣的事情發生而不出手扭轉局勢呢?耶和華果真是上帝嗎?正義果真有穩固的基礎嗎?

類似以利亞的鬱結,歷世歷代不少人可能都有不同程度的同感.有些人為國家及社會的進步而努力,卻被指控為國家和社會的敵人；有些人與公司共渡時艱,但公司卻不與他們分享成果；有些人堅持有教無類的教育理想,但關注及資源總是不投向他們；有些社工苦口婆心勸導邊緣者臨崖勒馬,但往往下次再見他們就在報紙新聞上；有些人關心家人,卻被誤會為管束及干預．．．這些狀況,都容易令人感到既憤怒又無奈,既不甘又無能為力.

耶和華問藏在山洞的以利亞:「你在這裡做甚麼?」我們想像,以利亞除了重覆之前的那段控訴之外,可能心裡還有聲音說:「世界如此黑暗,先知勢孤力弱,除了藏在山洞,還有甚麼可以做呢?」以利亞像受傷的野獸一樣躲到山洞裡,獨自默默舔拭傷口.

2．走到無力時遇上耶和華杜牧的《題烏江亭》一詩道:「勝敗兵家事不期,包羞忍恥是男兒.江東子弟多才俊,捲土重來未可知.」這是認為,項羽應該渡江,以圖東山再起.

不過,有論者認為,項羽跑到江邊,的確打算渡江.可惜來迎接他的,只是一位亭長,而不是一支大軍.項羽見到只有一舟一槳一亭長,東山再起的信心立即消逝.

以利亞先知雖然落荒而逃,甚至有人認為他慌不擇路,不過,他最終卻走到耶和華的山.在那裡,他遇到耶和華.當人走到無力的時候,遇到不同的人物,可以有很不同的效果.

2．1．以利亞一個人在「逃」上,耶和華上帝「陪著你走」

19 : 3以利亞害怕,就起來逃命,到了猶大的別是巴,把僕人留在那裡.

19 : 4他自己在曠野走了一日的路程,來到一棵羅騰樹下,就坐在那裡求死,說:「耶和華啊,現在夠了!求你取我的性命吧,因為我不比我的祖先好.」

19 : 5他躺在羅騰樹下睡著了.看哪,有一個天使拍他,對他說:「起來吃吧!」

19 : 6他觀看,看哪,頭旁有燒熱的石頭烤的餅和一壺水,他就吃了喝了,又再躺下.

19 : 7耶和華的使者回來,第二次拍他,說:「起來吃吧!因為你要走的路很遠.」

19 : 8他就起來吃了喝了,仗著這飲食的力走了四十晝夜,到了神的山,就是何烈山.

我們可以用郁達夫的文章「一個人在途上」來形容以利亞,不過這應該是逃走的「逃」,而不是路途的「途」；另一方面,我們可以用盧冠廷的歌「陪著你走」來描述耶和華對他的回應.在成功路上,眾人擁戴不足為奇,但在落難逃命之時不捨不棄則是難能可貴.耶和華的使者好像「毅行者」比賽中的補給員,為以利亞設立補給站,使他不致於餓死在荒山野嶺.在人生的旅途上,我們也可能有一個人在「逃」上的時刻,並且有意無意地把自己置身於危險境況.我們有否發現上帝陪著我們走呢?我們有否發現上帝的使者在我們逃走的過程中的補給呢?在人生走到無力時,在巨大困難的壓迫之下,我們不會想到將來,我們的反應如以利亞一樣:「夠了!」但耶和華的使者卻對他說:「起來吃吧,因為你要走的路很遠.」「你要走的路很遠」這句說話一語雙關,以利亞腳上接著走的路的確很遠,同時,以利亞的人生也任重而道遠.

2．2．以利亞懷疑真理的根基,耶和華展示基礎穩固

19 : 11耶和華說:「你出來站在山上,在耶和華面前.」看哪,耶和華從那裡經過.在耶和華面前有烈風大作,山崩石裂,耶和華卻不在風中；風後有地震,耶和華也不在其中；

19 : 12地震後有火,耶和華也不在火中；火以後,有輕微細小的聲音.

19 : 13以利亞聽見,就用披風蒙臉,出來站在洞口.聽啊,有聲音向他說:「以利亞,你在這裡做甚麼?」

以利亞意圖做「隱蔽中年」,淹沒於鬱結之中,就在這時候,耶和華呼喚他出來.出來做甚麼?出來解決他與耶和華之間的矛盾.局勢的發展似乎搖動了以利亞的信念及使命的基礎.耶和華也以搖動來回應以利亞,不過,耶和華搖動的是氣流,使烈風大作；搖動的是大地,使山崩石裂,火焰飛揚.以利亞是耶和華的先知,當然知道這是耶和華新自顯現的現象,因此蒙著臉出來見耶和華.以利亞看見局勢的混亂而產生信仰危機,耶和華向他展示他的信仰基礎絲毫無損,仍然穩固如昔.當人淹沒於鬱結的時候,當人藏進山洞獨自舔拭傷口的時候,很容易以為自己是「最後戰士」,以為世界已經完全黑暗,「出來站在山上,在耶和華面前」睜開雙眼,黑暗並沒有掌管世界.

2．3．以利亞妄自菲薄,耶和華委以重任

19 : 15耶和華對他說:「去吧,從原路回去,往大馬士革的曠野去.到了那裡,你要膏哈薛作亞蘭王,

19 : 16又膏寧示的孫子耶戶作以色列王,並膏亞伯‧米何拉人沙法的兒子以利沙作先知接續你.

以利亞的痛苦是:國家不需要他,人民不需要他,他成為了亡命天涯的先知,彷彿是世上最多餘的一個人.耶和華對以利亞重新委以重任,彷彿對他說:「雖然你感到人民不需要你,感到國家不需要你,但我耶和華上帝需要你.雖然你看見黑暗遮蔽大地,就是在黑暗蓋地的時候,大地更需要光明.雖然你看見人民背信棄義,就是在人民背信棄義的時候,信實和公義更顯得可貴.耶和華是全地的主,不是廟宇裡的神靈,沒有任何政治及軍事力量比耶和華更長久及真實.你回去,繼續實踐你的使命,並且把這使命傳遞下去.」耶和華對以利亞的回應,也是對歷世歷代所有致力實踐上主使命的人的鼓勵.

5．以利亞以為自己被放棄,耶和華視他為戰友

19 : 17將來逃過哈薛之刀的,必被耶戶所殺；逃過耶戶之刀的,必被以利沙所殺.

19 : 18但我在以色列中留下七千人,是未曾向巴力屈膝,未曾親吻巴力的.」

以利亞自身難保,難免想像先知團隊已經全軍覆沒,懷疑自己已經被遺棄.耶和華不單重新確立以利亞的信仰和使命的根基,再次委以重任,更把以利亞當作伙伴,當作戰友,向他透露「軍情」,把局勢的發展與他分享.耶和華也把尚有七千未變節的隊友這秘密向以利亞揭露.孤軍作戰很自然感到無力及被遺棄,不過,在我們不知道的地方,原來存在著我們不認識的戰友.更重要的是,賦予我們使命的上主,不是把我們當作一個傳聲筒,而把我們當作戰友,與我們共商「軍情」.

四．總結漫漫人生長路,難免有走到無力的一刻.無力前行的時候,不妨參考以利亞的經驗,暫時離開戰場,走向信仰的根基,使命的源頭,真理的基礎去重整自己.變幻不定的環境沒有能力回應我們內心深層次的矛盾和鬱結,只有耶和華上帝,那個給予我們信仰和使命的源頭才能重建我們繼續上路基礎和動力.



\newpage

\section{夏生命盛放的季節}
\label{sec:114}
\textbf{賀志勇博士 - 夏──生命盛放的季節}
\newline
\newline
\hyperref[sec:113]{< < < PREV SERMON < < <}
~
\hyperref[sec:toc]{[返目錄]}
~
\hyperref[sec:115]{> > > NEXT SERMON > > >}


羅馬書 12:1-2
\newline
\begin{longtable}{cl}
\hline
\hline
章節 & 經文 (和合本修訂版)\\
\hline
12:1 & \begin{tabularx}{0.7\textwidth}{X} 所以，弟兄們，我以神的慈悲勸你們，將身體獻上當作活祭，是聖潔的，是神所喜悅的，你們如此事奉乃是理所當然的。 \end{tabularx} \\ \\ \relax
12:2 & \begin{tabularx}{0.7\textwidth}{X} 不要效法這個世界，只要心意更新而變化，叫你們察驗何為神的善良、純全、可喜悅的旨意。 \end{tabularx} \\ \\
[1ex]
\hline
\hline
\end{longtable}
日期:2011年9月4日

題目:夏──生命盛放的季節

證道:賀志勇博士

讀經:羅馬書12:1-2

1, 所以弟兄們,我以神的慈悲勸你們,將身體獻上,當作活祭,是聖潔的,是神所喜悅的；你們如此事奉,乃是理所當然的.2, 不要效法這個世界,只要心意更新而變化,叫你們察驗何為神的善良,純全,可喜悅的旨意.

一,    引言:改變之難

有一位香港牧師,以前住在美國.有一次他帶女兒去看電影,就去買票.買票的地方有條高度線,如果小朋友在8歲以下,又不超過那個高度線的話,就可以買半價票.牧師的女兒10歲了,但是華人看起來年紀小,好像還不到8歲的樣子,可惜就是高了一點點.牧師買了自己的票,就把女兒壓下去一點,說"再買張半價兒童票."他女兒聽見了,很不高興地望著他說:"daddy!" 他一下子驚醒了,很羞愧地買了張全價票.後來他跟我們講,人的本性要完全改變真不容易,一不小心,就像孫猴子,露出猴子尾巴來了.

二,解經

但保羅在這短短兩節經文中,就極為強調,做基督徒是要轉變的,要更新,要變化的.保羅通過這短短兩節,告訴了我們做基督徒的三個訣竅,是哪三個訣竅,我們可以從經文中細細來看.

(1)    第一個訣竅:將身體獻上,當做活祭

保羅一開頭,就以神的慈悲來勸我們.神的慈悲,就像上次賴品超教授講的,就是神的救恩、神的憐憫.保羅這個勸誡,是因為神已經通過耶穌基督的寶血救贖了我們,所以保羅才能夠來勸我們.但這個勸字,是個很嚴厲的字.就好像有個小孩飆車出事了,他爸爸不會這樣溫柔地說:"兒子啊,你以後不要飆車了,很危險的哦!"他爸爸一定會這樣說:"你想死啊!以後不許飆車了!聽到了沒有!" 這是父親帶著愛的勸誡,這也是保羅勸誡的口氣.非常嚴厲,因為他知道,在這個世界做基督徒很不容易.我們在面對一個很複雜、很危險的社會,一不小心,我們基督徒就會被社會給吞沒了.

所以保羅要這樣嚴厲地勸誡我們,要我們將身體獻上,做聖潔的活祭.什麼叫活祭?活祭就是不用殺死了燒給神的.我有個貴州的好朋友,他是苗族人,他們那裡也會吃蜘蛛啊,蜈蚣啊.還說他們苗族過春節要祭天,要殺牛向上天謝罪,討來年的豐收.他不停邀請我去他們家過苗族春節,我就很懷疑地對他說:"你不是騙了我去,把我燒了祭天吧." 過去猶太人要燒牛跟羊去去祭上帝,去贖罪；但是耶穌基督來了之後,我們不用再被燒了去贖罪,我們要活生生的服侍神,事奉神.所以獻作活祭,跟後面"理所當然的事奉"是緊緊聯繫在一起的.

事奉這個字,在希臘文裡,其實跟敬拜這個字是同一個字.而理所當然這個字,原文是"合乎理智"的意思.也就是說,對於基督徒來講,理所當然的事奉,就是合乎理智的敬拜.事奉就是敬拜,而且必須是合乎理智的敬拜.

敬拜到底是什麼?有的人以為敬拜就是星期天到教堂裡唱唱歌,聽聽管風琴,聽聽牧師講道.講得好就多聽兩句,講得不好就回家去睡覺.還有的人追求一種所謂的敬拜的感覺,追求唱詩唱得很high,要把手高高舉起來；追求聽道一定要感動得流眼淚,沒有流淚,就算是沒聽到神的道.但保羅告訴我們,敬拜乃是理性的、是合乎理智的.如果你對敬拜沒有經過好好理解,深刻理解,那你的敬拜就不完全.就好像腓利對那個太監說的:"你所念的,你明白嗎?"神的道不僅僅要我們在感情上感動,還要我們用心去理解,因為它是合乎理智的.

其次,敬拜就是事奉,事奉就是敬拜.敬拜事奉不僅僅是星期天我們聚在教堂裡的事,其實也是我們每天生活裡的事.每天你的生活、工作、學習,都是事奉的好機會,也是敬拜上帝的好機會.有個老姊妹,以前是幼兒教師,退休了,想在自己的社區傳福音,但是傳來傳去人家都不聽.後來她就想了個辦法,在社區裡面義務幫人家帶小孩,開辦兒童主日學,再給家長傳福音,結果真的在社區裡開了家教會.她把她的事奉叫做"湊仔佈道法(帶小孩佈道法)".這就是事奉的生活.

有人曾經跟我講過:看一個基督徒成熟不成熟,就看他肯不肯主動去事奉.這很有道理!有的弟兄姊妹說:我跟神關係很好,我經常讀經、禱告,靈修；我說,這很好,但還不夠,那還不是一個完全合乎聖經的敬拜.如果你想真的跟神很好,那就主動去事奉神,也事奉人.有的人說,事奉神我明白,但是事奉人是一定需要的嗎?我說,是的.當你從心裡真正去事奉人時,你才能真正去事奉神.基督誇獎好撒瑪利亞人,並沒有誇他有多麼多麼愛神；他只誇他"憐憫別人",這種憐憫,乃是從神而來的事奉,也就是從神而來的敬拜.所以耶穌只有一句話:"你去照樣行吧!"

事奉並不吃虧!真正的服事並不吃虧!因為當你去服事別人的時候,你自己就會被神改變了.只有當你能對別人微笑的時候,你才會感覺到善意的美好.只有你肯為別人付出的時候,你才會感受到人際關係間的那種美好.不自信的人才不敢付出!如果我是李嘉誠,我怕捐一百萬麼?所以當我們去為別人付出的時候,神就在中間動工,他就在中間建立了我們的自信,還有平安.吃虧就是賺便宜,神的國跟這世界還真有點不一樣.

(2)第二個訣竅:不要效法這個世界,要心意更新而變化

在這一節中,效法這個字,希臘文原文也有"變化"的意思,效法這個世界,就是跟著這個世界的潮流去變化,是隨大流,是跟風.我們生在一個潮流的世界,香港是一個非常潮的都市,所以我們很容易去隨潮流變化.有朋友開玩笑說:香港這個都市,讓女人少買一件時尚衣服,或者讓男人不去買Iphone4,就好像讓他們背十字架一樣.因為不時尚,毋寧死!所以人家買樓,我也要買樓；人家小孩學鋼琴,我家小孩也要學鋼琴；人家罵董建華,我也要罵董建華.買東西不買好的,只買貴的,因為那是taste.我們怕離開世界的潮流,怕不合群,怕受世界排斥.世界就像塊磁鐵,總要把我們吸回去.

我以前在教會裡,聽一個姊妹分享過一件事,她說她不小心跟另外一個姊妹有點誤會.她本來想解釋,可是那個姊妹看見她,卻當她不存在一樣,一點都不搭理她.她聽著那個姊妹對別人分享,說自己怎麼怎麼被神拯救,怎麼怎麼內心變得柔和謙卑,怎麼怎麼充滿愛,那一刻,她覺得好假啊!她不想再解釋了,她很無力,她覺得基督徒跟這世界一樣虛偽.

是啊,教會是什麼?教會不是一個建築物,弟兄姊妹,你們就是教會.沒有什麼比教會的虛偽更可怕,更傷害人了.所以我跟她說,不管那個姊妹如何對你,你都應該去向她解釋清楚.也許你遇到了冷暴力,但是你引起了誤會,就有責任向人家解釋清楚,去跟人家和好.我不在中間做第三者,你直接去面對.她能不能接受你是她的事情,但你作為基督的門徒,你有你解釋的責任.

是的,我們這個世界迷信暴力.武器的暴力,經濟的暴力,政治的暴力,冷暴力.就好像這次倫敦的暴動,你會很驚奇,很多受過良好教育的小孩都可以這麼暴力,這麼殘忍.再看看我們的網絡上,發表評論大部分都是罵人,赤裸裸的語言暴力.也許在真實世界,人們還要裝扮一下,可是一到了網絡世界,許多人就根本沒有顧忌,就變得非常暴力了.

保羅要我們不要效法世界的道路.世界以為暴力意味著真正的權力,但上帝真正的權力,是他的忍耐,是他的受苦,受死,最後復活.世界的道路是一條引向一條沒有平安、沒有喜樂,只有死亡的道路,所以保羅告訴我們,不要效法世界,要心意更新而變化.   在原文裡頭,更新和變化兩個字意義差不多,都有新的意思.這個新字的意思,在希臘文裡不是買了件新衣服的新,新衣服只是外面的新.而聖經裡面的更新,是要我們從裡面更新.就好像一個蠶蛹,有一天破繭而出,變成飛蛾一樣.這種變化,叫做蛻變,這才是聖經裡頭所說的更新變化.更新變化,是個持續的動作,就是說我們要不住地更新,不住地變化.就是說,我們要藉著不住學習、不住事奉、不住禱告,去天天經歷神的新恩惠,學習神的新教導.只有這樣不住更新,我們才不會被這世界的潮流牽著鼻子走.

弟兄姊妹,我們問問自己,我們有多少天沒有讀聖經靈修了?我們有多少天沒有跟人和好了?我們有多少天沒有打過電話去關心一下別人了?這些小小的動作,都是每日的更新變化,只有這樣,我們才會不效法這個世界,真正成為新造的人.

(3)第三個訣竅:察驗神的善良,純全,可喜悅的旨意

保羅要我們常常察驗神的旨意,察驗,就是考察之後,然後心裡接納的意思.保羅要我們不要亂想,亂下決定,看到什麼,就立即下決定.明明是自己的意思,卻還要說是神的意思.這就不是察驗!

有一個團契,其中的弟兄姊妹喜歡用"異象"這個字,許多事開口閉口就講異象,非常屬靈.有一次有一個姊妹快畢業了,然後分享說神給了她"異象",讓她去內地山區服事貧窮的孩子,當時大家聽了非常受鼓舞,非常感動.然後那個姊妹就去了內地山區服事,去了不久就開始經常打電話回來哭訴,說蚊子特別多,沒有自來水,用水不方便,廁所不乾淨.終於有一天,她說神給了她"異象",讓她不要在山區服事了,讓她快點回香港來.這個戲劇性的變化,讓弟兄姊妹都要昏倒了.

你聽了這件事,或許你會問,我怎麼知道神的旨意是什麼呢?保羅在提摩太后書說,"聖經都是神的默示",那說明神的心意會在聖經中顯明；另外,保羅在歌羅西書中,要我們"恆切禱告,警醒感恩",因此我們要借著經常的禱告,去明白神的心意.

說到這裡,我想起了戴德生牧師的故事.戴牧師早年在英國帶了一個窮人的小組,向他們傳福音.有一次,他的雇主忘了給他發工資(出糧),他身上只剩下一塊值兩先令六便士的硬幣,也快要沒飯吃了.這個星期天,戴牧師帶完小組,有一個窮人請戴牧師去為他的妻子做臨終禱告,因為他的妻子快要死了,雖然他是個天主教徒,但神父要十八個便士才肯替他妻子禱告,他沒錢,而他妻子就快要死了.戴牧師心情沉重地來到那家人中,他發現那女人其實是快餓死了.他心裡就想:如果我身上有兩塊錢幣,而不是只有一塊錢幣,我就一定施捨一塊給這個女人.但是不幸他身上只有一塊,而他自己也在挨餓,他就捨不得把錢拿出來.然後他就跪下幫那個女人禱告,當他開口說:"我們在天上的父......"他的良心就對他說:"你敢跪下來稱祂為父?你口袋裡不是還藏著那塊錢嗎?"他幾乎禱告不下去了.他本來很喜歡禱告,但那一次,他非常艱難才禱告完,等他起身時,那個窮人對他說:"先生,你看到了我們的處境,為著上帝的緣故,請你幫幫我們把!"就在此刻,主的話閃進戴牧師的腦海中:"有祈求的,就給他們!"王者的話是帶有能力的!戴牧師慢慢從口袋裡抽出那塊錢,交給了那個窮人.他告訴那窮人,那其實也是他唯一的一塊錢了.當他說那些話時,他感到無比的真實,喜樂如同浪潮一樣卷過來.真實就是一種力量,王者的話是帶有能力的!戴牧師感覺到,捆綁他的捆鎖已經消除了,永遠消除了.因為他已經按照耶穌所吩咐的"你照樣去做吧!"去做了.

正是在那一刻,戴牧師全然真誠地面對神,察驗了神的善良,純全,可喜悅的旨意!他的生命也真正做到了改變.所以,要察驗神的旨意,那就多禱告,並借著聖靈的工作、聖經的亮光去提醒我們.

結論:盛放的生命

弟兄姊妹,保羅告訴了我們做基督徒的三個訣竅.第一個訣竅就是當做活祭,努力去事奉上主,服事別人；第二個訣竅是通過事奉上主,服事別人,讓我們的生命全然改變；第三個訣竅是通過生命改變,察驗上主善良,純全,可喜悅的旨意.每個人的生命都充滿著神的恩典.在這個教堂裡,有多少碩士博士,有多少音樂藝術才華橫溢的藝術家,有多少醫生教師等專業人士呢?恩典是豐富的.

主耶穌曾經大聲說過:"我來了,是要叫人得生命,並且得的更豐盛."弟兄姊妹們,我希望靠著神的話語,在走出教堂的時候,你們的生命就與過往不同,因為你們已經不住更新而變化了.泰戈爾有一句詩說:"生如夏花般燦爛."我也期望,弟兄姊妹的生命,一如這夏天的花一樣,開得燦爛!因為,我們是君尊的祭司,是聖潔的國度,是屬神的子民.



\newpage

\section{救恩的奧秘}
\label{sec:115}
\textbf{賴品超教授 - 救恩的奧秘}
\newline
\newline
\hyperref[sec:114]{< < < PREV SERMON < < <}
~
\hyperref[sec:toc]{[返目錄]}
~
\hyperref[sec:116]{> > > NEXT SERMON > > >}


日期:2011年9月11日

題目:救恩的奧秘

證道:賴品超教授

非常感謝主讓我再次有機會跟大家去思想上帝的話語,今天伍牧師給了我幾段的經文,那些經文也是跟這個主題有關──救恩(我發覺不知道是什麼原因,我很少會這麼近期,在一,兩個月內再來講第二次的道,這麼巧這次的主題與上一次的又有一定的關係),今天幾段經文(似在延續上次的講道-約拿書)都指向救恩問題──不單以色列人,外邦人也在上帝的救恩之內:

詩67:肯定萬民得知上帝救恩. 賽56:耶和華要彰顯他的公義和救恩,外邦人也要事奉耶和華. 太15:雖然耶穌來本是尋找以色列家的迷羊,但外邦婦人也因信而蒙拯救. 羅11:外邦人的大量歸主,是否意味上帝放棄以色列人?

外邦人,也就是以色列之外的民族,他們也會得知上帝的救恩,上帝要彰顯祂的公義,彰顯祂的救恩,臨到外邦人當中,而外邦人也要去事奉耶和華,馬太福音15章的經文裡,就有再進一步很具體的事例給我們看見即使主耶穌在他的生平的時候似乎主要是尋找以色列人家的迷羊,他來到就是為了這樣的目的,但是外邦的婦人(迦南的婦人),因為她有信心,同樣地得蒙上帝的拯救.理論上是一個很好的圖畫,似乎上帝的名得到高舉,在舊約裡所期盼的就是外邦的人也以認識上帝的救恩,在主耶穌基督身上開始展現出來,但是當我們看羅馬書第11章就會發覺保羅問了一個很有趣的問題,當這麼多外邦人來歸主的時候,是否意味著上帝放棄了以色列人呢?他直接問了這個問題,上帝現在是否已經放棄了以色列人呢?棄絕了他們呢?這就是我們所讀的經文裡第11章所提出的,他要很清楚地說上帝沒有棄絕祂的百姓,那為何保羅會有這樣的問題呢?為何保羅要問上帝是否棄絕了以色列人呢?我們需要知道當時的歷史背景,就是當門徒在五旬節之後,他們四散了,當然部分是因為過節要到耶路撒冷,過節後就要走；但是也有一部分的因為門徒受到迫害,要四周逃命的時候,就將這個福音由耶路撒冷帶到猶太全地,撒馬利亞去到很多外邦人的當中,接著歷史的發展就是越來越多的外邦人歸向耶穌基督.

透過新約聖經,我們知道很多書信所提及的地方,有許多都不在猶太的境地裡,那你就會有這樣的一個印象就是當時的基督徒大多數是以外邦的基督徒為主,猶太的基督徒反而成為了少數,直到今日也是這樣.因此,基督教給人的印象是「西方人」的宗教,你會覺得一個典型的基督徒是西方人,而不會是猶太人.猶太人的基督徒的是很少的,我曾經在外國教會教過一位猶太人基督徒,這是很少出現的情況,當時保羅也看到這個情況,為什麼會這樣的呢?本來救恩是從猶太人而出,他在羅馬書裡也說這個救恩臨到萬邦,先是猶太人,後是希利尼人,為何會有一個吊詭的情況出現?當許多外邦人信主的時候,反而那麼少猶太人信主,是否就代表:上帝棄絕了以色列人呢?

這跟我上次跟大家分享的約那的故事剛好不同,約拿是逃避上帝的呼召到利未未城宣講訊息,他想逃避,逃跑；保羅他沒有,他接受上帝的呼召去做外邦人的使徒,專向外邦人傳福音,但是這也帶給他一個身分上的掙扎,所以保羅所問的問題不單只是一個神學上的問題:究竟上帝是否已經棄絕了以色列人?上帝已經將一個以外邦人為主的教會來取代了原本的那批選民猶太人呢?

這不單只是神學上的問題,上帝是否後悔,不再跟猶太人立約呢?是否已經棄絕了他們呢?這是一個神學的問題,但對於保羅來說,則是一個身分認同的問題.如果上帝離棄了以色列人,那麼我作為一個以色列人來說,又有什麼意義呢?同樣地,他另一個身分是向外邦人傳道的一個使徒,他同樣也會問:我是猶太人的身分和我是外邦人使徒,這兩個身分是否有衝突的呢?神學上來說,你可能會問「公不公平?」以色列人一直等候上帝的拯救者,彌賽亞,基督,等候了這麼多年了,為什麼耶穌基督來的時候反而猶太人不相信呢?在猶太人的眼中,這些外邦人從來沒有追尋上帝,反而等著了耶穌基督,公不公平呢?就好像有人打完井水後,自己卻沒有飲用,而被別人喝了一樣.上帝是否不太公平呢?救恩既然是從猶太人而出,但將這個救恩帶去外邦的時候,為什麼猶太人反而好像被棄絕了?那麼少猶太人相信耶穌?公不公平呢?努力的反而得不到,不努力的,不追求的外邦人反而得到上帝的救恩,究竟公不公平呢?上帝會否後悔祂的揀選而離棄以色列人呢?

對保羅來說,這個當然是作為一個猶太人,以色列人身分的民族感情認同的一個問題,在第9章開始的時候,他開始就問這個問題,他是很有感情來去講,表達他對民族的關注,民族的感情:

1 我在基督裡說真話,並不謊言,有我良心被聖靈感動,給我作見證． 2 我是大有憂愁,心裡時常傷痛.3 為我弟兄,我骨肉之親,就是自己被咒詛,與基督分離,我也願意.  他關心以色列人的前途問題,他甚至覺得為此我被咒詛與基督分離我也願意,他關心他自己的民族是否還在上帝的救恩之內,所以他接著開始就說上帝是揀選以色列人的,一直都在歷史中揀選以色列人,但是他不能不面對這個事實,今天相信上帝,相信基督的猶太人,以色列人是很少,如果主耶穌快要再來的時候,以色列人的命運會是怎樣呢?他們是否已經被上帝拋棄了,不再是上帝的救恩裡的一部分呢?作為一個外邦使徒的保羅同樣問一個這樣的問題,如果我是一個外邦人的使徒,我又是一個猶太人,這兩個身分應該怎樣去處理呢?如果我是這麼關心我猶太的同胞是否能得到這個福音的時候,為什麼我還要去做一個外邦人的使徒,為什麼我不回頭去做猶太人的使徒,向猶太人傳福音呢?這兩個身分是否互相衝突呢?

有時候我們中國人也會問類似的問題:為什麼那麼少中國人相信耶穌呢?如果大家翻查歷史,最早的記載在唐朝的時候,有大秦景教開始在中國佈道,如果看西方的基督教歷史,那段時期就是日耳曼民族開始接觸福音的時候(日耳曼民族就是現在德國人的祖宗).在同一個時期,為何一個民族對福音有這麼好的接受,有那麼多人歸向它；但是我們中國民族似乎對這個福音很抗拒呢?還要抗拒了這麼多年,一千多年.在我們的民族中很少人信耶穌,是否上帝對我們不公平呢?是否上帝對我們沒有揀選呢?

不是,在我們的古代典籍也有提到上帝這兩個字,古代時候的中華民族對上帝是有某種的觀念,為什麼我們會這麼少人相信呢?在教會裡,我們有時也會問:現在中國有十多億的同胞,裡面有很多人也未相信耶穌,為什麼我們還要派宣教士去非洲傳道呢?為什麼我們不將我們的資源集中在中國的市場?現在很多的資金也是這樣,將資金集中在中國的市場裡,希望有更高的回報,我們不應很功利地去看這個問題,但是我們作為中華民族的一員,為什麼我們不將福音的資源,教會的資源放在中國的佈道事工上?我們反而差派宣教士到如日本這些與我們為敵的人去做福音的工作?為何我們差派人到一些與我們沒有什麼親戚關係的非洲去做福音的工作呢?(其實可能不是沒有親戚關係,如果那些生物學家說的是真的時候,我們的祖宗可能是由非洲來的),但我們仍會問同樣的問題,究竟我們中華民族在上帝救恩的計劃裡頭,是扮演著什麼樣的角色呢?上帝是否離棄我們?如果不是的話,為什麼我們福音的道路是這麼艱難呢?保羅在羅馬書第9-11章也一直在掙扎這個問題,就是究竟猶太人或者以色列人在上帝的救恩中扮演什麼角色呢?

如果我們讀完整段經文,就會發現很有趣的,到結束的時候,他由最初很苦惱,很激動；但到結尾的時候,他同樣地激動,不過他所激動的是他讚嘆上帝豐盛的智慧和知識.在11章33-36節──他讚嘆上帝的足跡是何其難尋,祂的判斷是何其難測,誰能知道主的心,誰做過祂的謀士；誰是先給了他,使他後來償還呢.因為萬有都是本於他,倚靠他,歸向他──中間的這段,究竟保羅由第9章開始到第11章結束的時候,他所掙扎的問題,他想到什麼呢?看這段經文的時候,保羅在第9章裡率先地肯定了上帝對以色列人的揀選是不會動搖的,祂是曾經揀選過,在歷史當中不斷地揀選,並且祂也很清楚地說沒有所謂公平與不公平的問題,你選了A,也就是沒有選B；這是很正常的一個情況,他也很肯定以色列人對上帝是有熱心,只是沒有按著真正確實的知識,特別是信的態度.到第11章,保羅首先所說的11:1-2的上半部:

1 我且說,　神棄絕了他的百姓麼．斷乎沒有．因為我也是以色列人.

如果上帝真是棄絕以色列人,應該一個以色列人也沒有是信上帝的,但是我信上帝,並且上帝差遣我作外邦人的使徒.這證明上帝對以色列人的揀選是從來沒有動搖過,在後面的經文裡更加講述上帝是不會對祂自己的揀選而後悔,他繼續說這一切也是出於上帝恩典的揀選.如果是這樣,11:11-24,就是25節之前的上文,他就會提出一件事,原來他看到以色列人的不信反而帶出一個很好的結果.

那個很好的結果就是將救恩帶到外邦人當中.所以保羅才會在25節,就是我們剛讀的那段經文說道這個救恩的奧秘是什麼.這個救恩的奧秘就是──以色列人有幾分硬心──什麼是硬心呢?在舊約聖經的出埃及記就是說法老怎樣硬心,就是同樣的硬心.無論他看到多少個災劫,到了第十災的時候才肯讓以色列人離開,在以色列人離開之後,他又後悔想追回以色列人,追殺他們,這就是硬心.你用很多方法想嘗試去說服他,使他改變心意,你發覺很困難(保羅就以此為同樣的表達方法).

以色列心硬直到何時呢?直到外邦人的數目添滿,於是以色列全家也要得救,背後的想法是什麼呢?舊約聖經說上帝要外邦人也來敬拜祂.當時保羅所看的,有誰可使外邦人也歸向耶和華的呢?不是摩西,不是所羅門,而是耶穌基督.我們發現今天研究聖經最多的人是什麼人呢?不是猶太人,而是那些所謂的外邦人,並且是從基督教的角度去研究,今天最多人看的舊約聖經不是猶太教裡的人去看的,他們人數是有,但是遠比基督徒更少.如果在一千多,二千年之後再看,哪個猶太人是最能吸引人去歸向舊約聖經裡所講的耶和華,是誰呢?可能就是耶穌基督.保羅看見當越來越多外邦人信主的時候,這就是一個證明,證明耶穌就是那個被揀選的彌賽亞,就是上帝的拯救者,就是那個基督.保羅看見一個很有趣的情況,就是如果越多的外邦人歸向主,就越能夠證明耶穌就是以色列人素來所盼望的基督,原來他自己的外邦使徒身分和他是以色列人的身分不一定是矛盾的,也就是當他越努力作一個外邦人的使徒去傳福音,又越來越多外邦人信主的時候,就越能夠證明耶穌就是那個基督.他的盼望就是當以色列人看到耶穌就是那個基督的時候,以色列全家也要得救,原來他看到的是這件事.

後來,他為何會有一個讚嘆呢?他看到兩件事.我們可以從兩方面去說,第一方面是救恩的本質是什麼?我們習慣了在教會裡講,我站在一個宗派的立場的時候,會有一個標準答案,就是因信稱義,本來這是一個一定對的答案,神學院也一定是這樣子教導的──「因信稱義」──但是當我們講到因信稱義的時候,我們就將重點從上帝那裡轉到我們自己本身,好像是因為我有這樣的信心,所以上帝一定要稱我為義,但是保羅去思考以色列人的命運的時候,他發現原來這一切最終,最重要的部分不是人的回應是什麼,而是那個定義的上帝.所以在第9-11章有很多的經文也是引用舊約聖經的,就是在乎定義的那位.簡單地說就是得救與否不是以人的主觀意志為轉移(這是我經常在內地聽到的說法),而是在乎那位定義的上帝.

給大家一個簡單的邏輯問題思考一下,這是小學中學所學的簡單邏輯問題,如果我作為一個老師跟一個學生說如果你今次考試得A的話,我就送你一份禮物(通常不是老師做這樣的事,是父母才會做這些利誘的工作去引誘子女考試考得好一點).邏輯上有四個可能性會出現:

學生考試拿了A,老師送他一份禮物(他沒有食言,按照他所應許的去做,非常公道).

學生考試拿不到A,老師沒有送禮物給他.(他沒食言,也很公道,沒欺騙).

學生考試拿A,老師食言不送他禮物.(不公道,這個老師一定被學生投訴).

學生考試拿不到A,老師仍然送一份禮物給他.(老師有沒有食言?沒有).

我告訴大家如果你拿到A我送禮物給你,但是我沒有說你拿不到我不送給你,大家不要多想.即使你考不到A,我也可以照樣送禮物給你,我沒有欺騙你,也沒有食言,那有沒有違反公道的原則?可能也沒有.但是這是出於什麼呢?是老師的憐憫,也就是即使猶太人不信不代表上帝已經取消了祂的揀選,即使猶太人不信,他們仍然在上帝的救恩裡,這是完全可能的,因為最終的救恩是在乎上帝的應許,在乎上帝的憐憫,而不是那個人值不值得上帝去救.不是因為他信,所以他值得上帝去救,因為他的信也是上帝所給的,他硬心與否也是上帝在背後所謂的「作弄」.人能否得救最後是在上帝的手中,猶太人的得救也與外邦人的得救一樣,他也是因為上帝的憐憫,而不是猶太人有什麼特別的東西,因此在保羅看來,即使現在我看到的以色列人,他們很多也是不信的,但不代表上帝不再憐憫他們,不等如上帝已經收回祂的恩典.羅11:29:上帝的恩賜和選召,是沒有後悔的.既然上帝在歷史中已經選召他們,上帝會繼續堅持下去,不會放棄以色列人.

接著保羅不只從救恩的本質上,也從第二個角度去看上帝的恩典這個問題,在同一個救恩上來看,原本以色列人的不信既然帶來一個好的後果,就是外邦人也可以信耶穌,這個救恩臨到外邦人的時候,保羅就說以色列人的硬心,他甚至說作福音的仇敵,如果是讀原來的文字,他所用的表達方式是被動式,簡單地說就是:不是我想硬心,是你使我硬心,你會覺得很不合理?在舊約聖經裡,為什麼法老要到第十災的時候才肯釋放以色列人呢?不是因為法老本身硬心,而是背後有一個上帝使他硬心.簡單地說,保羅就是看見以色列人的硬心是出於上帝的計劃,是上帝工作的一部分,是上帝使他們硬心,以色列人的硬心就好像法老的硬心,是上帝特別的護理,為要成就祂自己某個特別的旨意.如果是這樣,保羅可以很合理地盼望當目的達成後──就是福音臨到外邦,許多人信主的時候──當外邦人得救的人數添滿的時候,上帝就會回頭去憐憫以色列人,讓以色列人全家得救,他覺得這是合理的人,如果不是,上帝就是「打完齋不要和尚」的上帝.既然上帝選用以色列人作為工具,使福音臨到外邦人,成就上帝在舊約聖經裡所講的遠象,以色列人的不信是出於上帝的計劃,當這計劃的目的完成的時候,上帝自然就會憐憫以色列人.在保羅眼中,以色列人的不信是有一個目的,就是使外邦人可以蒙憐恤.11:30「

你們從前不順服　神,如今因他們的不順服,你們倒蒙了憐恤」──以色列人是為著福音的原故,是為了外邦人的原故,成為了好像福音的仇敵,好像敵擋著福音,但是這只是要成就上帝給他們的一個特別的歷史任務,當這個歷史任務完成後,上帝從來沒有後悔對他們的揀選,從沒有後悔對他們的憐恤,時候到了,以色列人全家自然會得救.

保羅沒有具體地說是怎樣得救,但是他相信這是出於上帝恩典的計劃,是一個很奇妙的計劃.你會覺得這計劃不是很合理,大家想想是嗎?怎麼會使原本是選民的不信,反而使福音臨到外邦人?也就是透過一些人的不順服,好使恩典顯多,透過一些人的不信,使更加多人相信.大家推銷的時候應該也不會用過這樣的方法,是嗎?如果你去到一間酒樓,伙計跟你說我不會來這間,我去旁邊的那間,你還會在這間酒樓吃飯嗎?你會很有保留,伙計自己也不在這間酒樓吃飯,你還會光顧?救恩由猶太人而出,耶穌是猶太人,猶太人也不信,為什麼你們這群外邦人會相信耶穌呢?是否很不合理呢?上帝就是要人相信,為何故意地使某些人硬心,使到他們不信,反而使另一班人相信呢?在我們的邏輯底下,這好像很不合理,不合乎我們所認為的常理.如果上帝要使外邦人信耶穌,怎會先使那班猶太人不信呢?連猶太民族也不相信耶穌,你怎樣使外邦人相信呢?如同我們推廣中國文化,連中國人也不再接受我們所推廣的中國文化的時候,我們怎能使其他外國人接受中國文化的精粹?你會覺得是一個很荒謬的方法,但是上帝就用了一個我們認為像是荒謬的方法──出自猶太人的主耶穌,卻是被猶太人拒絕的猶太人──使更加多的外邦人歸向祂.兜兜轉轉,很吊詭的一個方法,當許多外邦人歸向主後,上帝會再使那些猶太人,以色列人再次歸向祂.

保羅想來想去,最後他不能夠不讚嘆上帝的工作是多麼奇妙的一個計劃,就是用一個我們認為好像不是很合理的方法去達成祂的計劃.當保羅想到這點的時候,他發覺我猶太人的身分,以色列人的身分原來跟我作一個外邦人的使徒的身分是完全沒有衝突,我為著我同胞的得救,我應該更努力向外邦人傳福音,因他們越多人信,以色列人更早得救.所以接著的第12章,他就很清楚地說不需要問太多這類問題,我們更需要做的是將自己的身體獻上當作活祭,奉獻給上帝.在他讚嘆後他的反應就是這樣,因為他看見原來所做的一切也都在上帝恩典的計劃中,我們只是他恩典計劃中的一顆棋子.

作為基督徒,我們有時很喜歡自誇,誇我們是信主的人,那些是不信主的人下地獄是應該的,該死的；覺得我們有信心,應該上天堂... 但是當我們讀保羅所寫的這段經文的時候,我們就會發覺原來我們信或者不信有時原來不是在自己的掌握當中,有一些人不是想不信的,我曾經接觸過一個醫學院的教授,他說我真的很想信主,但是總是信不下去,你有沒有辦法幫助我呢?原來有一些人真的是這樣,他很想信,但總是覺得自己信不了.有一些人是不想信的但又信了,情不自禁,好像墮入愛河一樣.這個說法真好,情不自禁,不自覺地踩入陷阱裡,所以愛情是一個陷阱,不自覺地,不能自主地墮下去,信也是一樣.有時是你不自主地信了,你沒有辦法自主自己信與不信,信與不信也是出於上帝的恩典.

神學家保羅田麥克就關於什麼是因信稱義有這樣的說法,就是「唯獨信心」.不是我們有了這個信心,這信心就是我們自己擁有一件東西.這是一個恩,上帝見我們有這個信的時候,祂就把恩典給我們.不是!他說我們稱義是一個恩惠恩典藉著信心來稱義,因為所有稱義的,是出於上帝的恩典,我們只是憑信心來去接受而已.所以,與其他人相比,我們更沒有任何優越的地方,我們不可能透過這件事就說我們比其他人更加敬虔,更加值得上帝的恩寵.如果我們自以為有信心就應該上天堂,其他人應該下地獄,可能跟當時候的猶太人沒有分別(只不過是五十步笑百步的想法)──因為我們守律法,所以上帝應該把恩典給我們,那些外邦人沒有守上帝的律法,他們應該死──保羅在這段經文裡帶出一個訊息:原來我們今日所擁有的一切,甚至包括我們的信心,原來也是上帝給我們的恩賜,我們在這件事上,一無所誇.所以我們不可能跟上帝講:我們覺得我們有遠古的文化,應該值得上帝去揀選我們,我們覺得我們有這麼深厚的文化,幾千年前已經開始追問上帝,為什麼上帝不揀選我們呢?

原來沒有所謂公平不公平的問題,我們所謂的成就在上帝眼中可能不值一提,但是我相信我們的民族仍然是蒙上帝的揀選,至於我們今日為何只有這麼少基督徒,坦白說,我也不知道.正如保羅可能也不知道究竟外邦人的數目何時才填滿,我們可以看到的是什麼呢?原來當上帝施行祂恩典的時候,往往是出乎我們的意料,出乎我們的期待,甚至是我們沒有期待,沒有想過上帝會這樣作工的時候,祂就是這樣作工.上帝可以透過人的軟弱去顯出祂的剛強,祂透過以色列人的不信去將救恩帶到外邦人,去成就自己的計劃.

我覺得我們不需把得失成敗看得這麼重要,沒錯,我們要努力去傳福音,但有時候,上帝是透過我們的失敗來成就祂的救恩,因為祂可以用一些我們意想不到的方法來成就祂的計劃.我們想想,父上帝也是藉著基督的受苦,看似軟弱,一個被猶太人離棄的所謂猶太人的王,是將上帝救恩帶給外邦人.有時當我們失敗,我們受苦,我們遇到挫折的時候,這不等如上帝的失敗,因為上帝是可以藉著我們的軟弱,我們的痛苦,我們的失敗去成就祂的救恩.真正信靠上帝的,不是相信我們可以為主作什麼?可以做許多很偉大的事跡,而是我們相信上帝是可以藉著我們的失敗或者我們的成功去成就祂的美意.

我個人一個小分享:當我96年從英國回香港的時候,最初我也覺得自己要去教導那些神學生(就是那些想著讀完書後到教會事奉的人),但是事隔十五年,大概十五年,我再回頭看的時候,原來上帝也可以使用我去做另一件事,我發覺我越來越少接觸的就是這類神學生,但我更多接觸到的是對信仰有興趣,不知道自己是信還是不信,介乎在信與不信之間,有時信有時不信,對信仰有懷疑的人.在過去的一段時間裡,我發現原來我到中國大陸裡參與講座的時間,可能比在香港教書的時間還要長.這是我最初回來的時候完全沒有想過的事情,我覺得上帝是可以在我們身上做出很多很奇妙的事,是我們沒有想過的.願意我們再次思考以色列人命運的時候,我們同樣地去想想,我們不知道上帝在我們身上最後會成就什麼事,但是我們願意好像保羅在羅馬書第12章裡這樣說:我們將我們的身體獻上當作活祭.

當我們願意將我們的身體獻上當作活祭的時候,神總是在祂救恩計劃中用得著我們,無論是這件事,還是那件事；或是大事,或是小事；或是藉著我們的成功,或是藉著我們的失敗,可能是藉著我們的強處,也可能藉著我們的弱點去成就祂的救恩,最重要的是我們願意擺上我們自己為主所用,願主祝福大家,阿們.



\newpage

\section{伏在門前的野獸}
\label{sec:116}
\textbf{伍渭文牧師 - 伏在門前的野獸}
\newline
\newline
\hyperref[sec:115]{< < < PREV SERMON < < <}
~
\hyperref[sec:toc]{[返目錄]}
~
\hyperref[sec:117]{> > > NEXT SERMON > > >}


日期:2011年9月18日

題目:伏在門前的野獸

證道:伍渭文牧師

讀經:創4:1-16

美作家霍桑(Nathaniel Hawthorne, 1804-64) 是美洲新大陸清教徒移民的後代,在他的名著

緋紅字母 (The Scarlet Letter)

開始時有一段發人深省的話:

建立新殖民的尋夢者,縱使他們原本對人類的德行和快樂,有烏托邦的無限憧憬,但總認識到,面對初期眾多現實需要中,也得在其處女土地上,分出一幅作為墓地,另一幅作為監獄.(The founders of a new colony, whatever Utopia of human virtue and happiness they might originally project, have invariably recognized it among their earliest practical necessities to allot a portion of the virgin soil as a cemetery, and another portion as the site of a prison.)

.

霍桑用文學的手法,道出人生的現實 and \#8212;面對罪惡和死亡.新殖民來到美洲新大陸,雖然獲得宗教信仰的自由,但他們仍然不能擺脫罪惡和死亡,需要監獄和墓地.今天舊約創四1-16記載該隱殺弟弟亞伯的故事,是聖經首次把罪惡和死亡同時呈現在我們眼前,並說明罪的本性和勝過罪惡的方法.

1. 罪必戀慕你 and \#8212;罪惡的誘惑和牽引

耶和華對該隱說:你為甚麼發怒呢?你為甚麼變了臉色呢?你若行得好,豈不蒙悅納?你若行得不好,罪就伏在例門前.它必戀慕你(its desire is for you , NRSV；新譯本:它要纏住你),你卻要制伏它.」(創四6, 7)

該隱殺亞伯的故事中,對罪有一個非常生動的形容:罪像伏在我們門前的野獸,靜靜等候我們,祇要我們走出門口,它便有機可乘,把我們吞噬.當然,若我們知道,這野獸若是面目可怕,看來使人毛骨聳然,我們就不會開門,走出街外,送羊入虎口.但聖經說到罪這野獸時,提到有一種牽引,我們要很費力才能擺脫它.「

它必戀慕你(its desire is for you , NRSV；新譯本:它要纏住你),你卻要制伏它.」(創四7)這野獸開始時看來可能是可愛的寵物,抱在懷中,令我們傾心,愛慕不已,甚至神魂顛倒,不能自己,但一搭上它,就被它纏住,不放過我們,我們也不想擺脫它,因為罪可以帶來歡娛--罪中之樂,成為心魔,揮之不去,但成為這野獸最終露出憎獰本相,要把我們吞噬.

這段經文使我想起上星期一則法庭新聞.報載:「已婚中年富商5年前與酒吧女公關經理邂逅,迅即「撻著」發展孽緣,惟因錢銀轇轕終令兩人分手,之後女公關疑「獅子開大口」勒索富商1.39億元,更涉請私家偵探跟蹤對方,再以「辣招」恐嚇他及其妻女,揚言要他們「一家三口家破人亡」.富商數度「投降」,終令她成功取得千萬分手費.」(明報八月二十五日) 對富商來說,這野獸是好色,但對女被告來說,這野獸是貪心,被告己得到一千萬,應可收手,但貪心纏上了被告,她想1.39億元,最終被刑事檢控.

對該隱來說,伏在門前的野獸是甚麼呢?為甚麼他要殺亞伯呢?有甚麼力量牽引著他呢?「

耶和華對該隱說:你為甚麼發怒呢?你為甚麼變了臉色呢?你若行得好豈不蒙悅納?你若行得不好,罪就伏在門前.」(vv. 6, 7)伏在門前的野獸是心中的嫉妒.該隱不能忍受弟弟竟然蒙上主喜悅,我和他一起在祭壇獻祭,但上主接納他而不接納我的祭品,亞伯祭壇上騰的煙火是濃厚,而我的煙火稀薄.一路以來,該隱不單是長子,而且是獨子,他名叫該隱,這名字意思:「耶和華使我得了一個男子」,是夏娃的驕傲,也是她信心的憑證,他是將來家族的唯一盼望,擁有父母全部的愛,直到弟弟的出現.

一路以來,其實該隱並不覺得弟弟是一個威脅,亞伯名字意涵為脆弱,無有,是弱者的名字；他沒有說過甚麼,給比他強很多的該隱引到田野,手起斧落,生命就完結了.一路以來,該隱可能祇視亞伯為工具,協助他完成長子的功業.所以,當他知道上主接納亞伯而不接納自己時,「就大大發怒,變了臉色」(新譯本:非常忿怒,垂頭喪氣) ；這兩種極端情緒同時出現,表示他的自尊受到傷害,甚至破碎.其實,他應該躬身反思,為何如此?是我的崇拜的態度出了問題?是弟弟有些好東西我過去沒有察覺嗎?上主既然是上主,我要順服他,雖然我不明白.

他沒有躬身自問,嫉妒充塞他的心,他因妒成恨,下了殺機,要把亞伯除掉.嫉妒,仇恨好像野獸伏在門前,靜靜等候該隱.上主知道他身處險境,告誡他:小心那野獸!不要踏出上主為你設定的界限,亞伯是你的弟兄,你看守的對象,回頭吧!不要踏出門外!他聽不進上主的警告,因他不能忍受比他脆弱的弟弟勝過他.該隱纏上了嫉妒,他他戀上了嫉妒,因為嫉妒能給他動力和目標,他計劃除掉亞伯.

我認識一位朋友,三十多年前初到美國德州,在一商場買東西時,可能英文不太流利和是黃色膚色,感到被奚落和岐視,以後用幾年時間,努力賺錢,把那商場買下來,還以顏色,吐氣揚眉.當被報復的心纏上時,我們的心也戀慕報復,而報復給我們動力和目標,報復成了我們揮之不去的心魔.

2. 守望自己的心

伏在門前的野獸,提醒我小心罪的潛伏性 and \#8212;就在我們心中.我們要注要該隱的殺機起於獻祭的地方,獻祭的時候,就是在進行崇拜的時候該隱竟然犯罪,因為該隱的心遠離上主.上主沒有說:該隱,你的祭品不及格,乃說:

你若行得好豈不蒙悅納?

詩篇是舊約公共崇拜的禱文和詩章,今天詩篇經課一一九:33-40就談到上主的律法,這是最長的詩篇,共176節分二十二段,每段依次序用同一希伯來文字母起頭,非常優美,內容全是談及我們的心要愛慕上主的誡命 and \#8212;我們要遵守的界限.經文說:「

求你使我的心趨向你的法度,不趨向非義之財(selfish gains).求你叫我轉眼不看虛假(worthless things),又叫我生活在你的道中.」(詩119: 36, 37)

罪是有誘惑和牽引的,勝過這牽引需要更大的牽引,就是上帝的律法.要勝過病毒和細菌祇有一個方法 and \#8212;增強免疫能力.試探和誘惑隨時隨處都在,就如我們沒法離開地心吸力的牽引.我們會趨向非義之財,也容易被虛假的閃爍和炫耀所吸引,但在上主話語的亮光中,我們得著屬靈的透視,明白到非義之財 (selfish gains) 和虛假 (worthless things) 都是過雲煙,轉眼成空.在親近上主時,我們也被上主的話語所吸引,這屬靈的吸引力,幫助我們擺脫罪的誘惑,就如飛機若不受地心吸力牽引下墮,就要被更大的氣流動力所牽引,停留在高空.

3. 彼此守望

罪是伏在門前的野獸,指出罪的嚴重後果,罪就像要吞噬我們的野獸,但這比喻也指出保護自己免被罪攻擊的方法 and \#8212;不逾越界限,這野獸就不能傷害我們,主禱文中赦免我們的債(trespasses),就有逾越界限的意思；觸犯上主誡命,就是冒犯上主,虧欠了上主的榮耀,所以用債一字表示罪.罪這野獸祇能伏在門外,不能進入家中,祇要我們不逾越界限,這野獸就不能傷害我們,「

這家就是神的教會,真理的柱石和根基」 (提前3:15)

今天福音書經課馬太十八15-20提到信徒之間要彼此守望,互相提醒,一起祈禱；書信經課雅各書五19, 20說要作守望者,使人從迷路中回轉:「

我的弟兄們,你們中間若有人迷失真道的,有人使他回轉,這人該知道,叫一個罪人從迷路中轉回,便是救一個人靈魂不死,並且遮蓋很多的罪.」(雅五19, 20)

在上主的群體中,我們學習彼此認罪,一起追求成聖,得到屬靈的保護.當該隱知道自己罪孽深重時,人人看到他時就要追討他的罪:「

我的刑罰太重,過於我所所能當的.」(創四13)上主對悔罪的人,永遠給予接納和保護,這是信仰群體的挽回事工.

在屬靈的追求上,我們要勉力而行,不可疏懶,不想聚會時,可能就是我們最需要親近上帝的時候.我小時候,還沒有即食麵時和微波爐時,父母常對我們說:吃飯時不要趕著玩而不進食,要強飯加衣,保守身體健康.聖經教導我們不可停止聚會:「

又要彼此相顧,激發愛心,勉勵行善.你們不可以停止聚會,好像那些停止慣了的人,倒要彼此勸勉.既知道那日子臨近,就更當如此.」(來十24, 25)

數年前有一對夫婦在教會聚會,完了就離開,很少到下邊茶聚,更少一起午膳,後來兩夫婦先後完成神學訓練,在教會工作了.有一次他們告訴我:牧師多謝你讓我們在崇基教堂安靜聚會,沒有要他們參與這事奉或那事奉,也不用崇拜散會後開這會開那會,可以靜靜療養心靈的傷痛,我們在這裡得到非常需要的空間.

我要讚賞他們,在心力交瘁時,他們沒有停止聚會,繼續親近上主,聆聽上主的誡命.船停泊在港灣,是維修整頓,儲備糧水,但內心仍渴想那無際的汪洋,等待風起雲湧時,揚帆遠去.

這種在基督裡友愛守望的關係,與該隱的世界成了強烈的對比.該隱說:「

我必流離飄蕩在地上,凡遇見我的必來殺我」(創四14)這句話使我訴想起利比亞內戰政府軍人現時的遭遇,殺人者,現今懼怕被人追殺；決定開始戰爭的人,不能決定何時結束戰爭.

因為罪使人與人失去信任,彼此懼怕對方,到創世記十一章,他們來到偌大的示拿平原,他們懼怕分散全地,但上帝豈不是告訴亞當夏娃:要生養眾多,遍滿全地嗎?因為罪使人與人關係割裂,人要採用一種言語,集中居住同一地方,以便統治；他們認為不同的語言,居住不同的地域,會引起不穩定,不和諧,因罪使人與人彼此懼怕,失去信任；罪使人邁向專權,要建造巴別塔,塔頂通天,取代上帝.他們要統一發新聞稿,發一樣的官方圖片,但上帝使人發明微博,tweeter,把語言混亂,上帝不是作弄我們,乃保護我們.

但在上主的家,我們看看一個逆轉,在這裡我們彼此相愛,為對方的靈魂儆醒,在這基督為元首,在這裡我們彼此謙卑服待.不要停止聚會,強飯加衣.一顆碳放在一堆碳中就燒得長久,但離開這堆碳就很快熄滅.我非常高興看到幾個月前開始了周四晚上的祈禱會,讓在香港工作的內地生,彼此有更多生命的分享和建立.

4. 為社會守望

教會不單在內彼此守望,教會作為社會的良心,也要為社會守望.該隱殺亞伯不單是家庭的仇殺,亦是人類社會的寫照,罪破壞人與人的信任,罪是自我中心,罪是放棄看守弟兄的責任,該隱說:「

我不知道!我豈是看守我兄弟的嗎?」教會作為被救贖的群體,天國的子民,就要履行為弟兄守望的責任.

上主日我們認識到,是教會透過基督教工業委員會的守望,爭取到工人有薪公眾假期和女性分娩假期,以前工人一年才在農曆初一到初三放假.在香港的勞工子弟學校,是教會的一位主教創辦,不是左派工會創辦.在審判的日子,耶穌清楚告訴我們,能進天國的不是你信些甚麼,乃你那沒有實行所信的,進天國的人是守望弟兄的人:「

我餓了,你們給我吃；渴了,你們給我喝；我作客旅,你們留住我；我赤身露體,你們給我穿；我病了,你們看顧我；我在監裡,你們來看我.」(太二十五34-36)他們的名字叫亞伯--弱者,無有,這是上帝最關心的人.今天,上主也問我們:「你弟兄亞伯在那裡?」

當我們為社會守望時,也許不能改變社會,但起碼不容易被社會改變我們.

結論:罪像伏在門前的野獸,伺機把我們吞噬,一戀上了罪,它就纏上了你,但祇要我們活在上帝的誡命中,不逾越界限(trespasses) ,我們就不會被罪傷害我們,我們要保守自己常在上主的話語中,被上主的話吸引,也藉著信仰群體的彼此守望,不至迷失,在迷失時也得到挽回；在為社會守望時,我們也不斷儆醒,不隨波逐流.我們也許會偶然被過犯所勝,但祇要回轉,就得著醫治,挽回,重新建立,再次揚帆出海.阿們.



\newpage

\section{若有人要跟隨我}
\label{sec:117}
\textbf{馮煒文先生 - 若有人要跟隨我}
\newline
\newline
\hyperref[sec:116]{< < < PREV SERMON < < <}
~
\hyperref[sec:toc]{[返目錄]}
~
\hyperref[sec:118]{> > > NEXT SERMON > > >}


馬太福音 16:21-28
\newline
\begin{longtable}{cl}
\hline
\hline
章節 & 經文 (和合本修訂版)\\
\hline
16:21 & \begin{tabularx}{0.7\textwidth}{X} 從那時起，耶穌才向門徒明說，他必須上耶路撒冷去，受長老、祭司長和文士許多的苦，並且被殺，第三天復活。 \end{tabularx} \\ \\ \relax
16:22 & \begin{tabularx}{0.7\textwidth}{X} 彼得就拉著他，責備他說：「主啊，千萬不可如此！這事絕不可臨到你身上。」 \end{tabularx} \\ \\ \relax
16:23 & \begin{tabularx}{0.7\textwidth}{X} 耶穌轉過來，對彼得說：「撒但，退到我後邊去！你是我的絆腳石，因為你不體會神的心意，而是體會人的意思。」 \end{tabularx} \\ \\ \relax
16:24 & \begin{tabularx}{0.7\textwidth}{X} 於是耶穌對門徒說：「若有人要跟從我，就當捨己，背起自己的十字架來跟從我。 \end{tabularx} \\ \\ \relax
16:25 & \begin{tabularx}{0.7\textwidth}{X} 因為凡要救自己生命的，要喪失生命；凡為我喪失生命的，要得著生命。 \end{tabularx} \\ \\ \relax
16:26 & \begin{tabularx}{0.7\textwidth}{X} 人若賺得全世界，賠上自己的生命，有甚麼益處呢？人還能拿甚麼換生命呢？ \end{tabularx} \\ \\ \relax
16:27 & \begin{tabularx}{0.7\textwidth}{X} 人子要在他父的榮耀裡與他的眾使者一起來臨，那時候，他要照各人的行為報應各人。 \end{tabularx} \\ \\ \relax
16:28 & \begin{tabularx}{0.7\textwidth}{X} 我實在告訴你們，站在這裡的，有人在沒經歷死亡以前，必定看見人子來到他的國裡。」 \end{tabularx} \\ \\
[1ex]
\hline
\hline
\end{longtable}
日期:2011年9月25日

題目:若有人要跟隨我

經文:馬太福音16:21-28

證道:馮煒文先生

各位弟兄姊妹今天很高興有機會和大家分享,特別是分享馬太福音十六章的一段經文,當我預備今天的講章時,使我記起不久前有一次的經歷,香港有一群很愛主的基督徒商人,他們組織了不同的商人團契,當中一位弟兄邀請我參加他們一次的福音聚會,在中環一所很高級的酒店貴賓廳以午餐形式舉行,在那一小時的聚會中我很受感動.

聚會的形式與一般商人的交流會差不多,開始大家喝著雞美酒,互相交換名片,時候到了就坐下用餐.聚會的人數約有廿多人,負責的弟兄有三,四位,當中有一位分享見證,講述他怎樣做生意,我相信他是在金融業工作的,他分享開始時很順利,他怎樣用盡各樣的方法,才幹和人脈來建立他的事業,然而在當他忽略了他的家庭,忽略了他的太太,兒女,他陷入了一個很大的困境,很大的心靈痛苦當中,後來他接觸了信仰,開始經歷神,結果他悔改,他改變了自己的人生和生命.

很簡單十至十五分鐘的見證,很真誠,我為他高興亦受到感動,我相信在座其他的朋友大部份都是未信主的,同樣受到感動,完結前另一位負責的弟兄上台帶領說:「我們平日聊天喝酒開開心心便過去了,今天我們來一個例外的好嗎?一起為我們的家人祈禱!」結果聚會中廿多位菁英就低頭祈禱.

我很感動,在聚會中我隱約領略到人對基督的那份的渴望,我亦領略到信仰是實際的,是和我們的生命有關的,信仰是有效的,結束會議後我心裡面有點陰影存在.邀請我參加的那兩位弟兄詢問我對聚會的意見,我們在酒店的咖啡廳談了半小時,他們希望知道我的反應,我就將剛才分享的說了出來.

在這一小時多的見證中我看見我所認識的信仰影子,我認識到的信仰有一點成份在裡面,但我認識的耶穌是在馬糟出生,活在獻祭之外,行十字架的道路,又叫跟隨祂的人要背起十字架來跟隨祂,當我準備今天的講章時就想起這個經歷.並不是互相矛盾,互相對敵的,兩種對信仰不同的看法,但是一個很不同的信仰描述和敘述的方法,亦都帶來了怎樣作基督徒的不同.我依然很欣賞當天我所經歷的,但我的信仰和我所認識的主是有些不同的地方.

馬太福音十六章21-28節的背景就是耶穌要上耶路撒冷受苦犧牲,所以祂要門徒跟隨祂,跟隨祂是什麼意思呢?受苦,犧牲,拾己,背十字架來跟隨主,這就是當時的情境所得到的詮釋,可以說馬太十六章這段經文是反映基督教那種悲情的本質,如果我們認為耶穌的死和十字架是基督教信仰的核心,這樣基督教的本質就有這個悲情的成份.但如果我們看耶穌一生(祂三至四年的生涯教導),我們就會發現祂這句話「來跟隨我」,很多時候都沒有這個悲情的成份存在的.

最明顯就是馬可(符類福音的第一本),亦可以說是四福音書的原本,馬可第一章記載耶穌開始出來工作就沿加利利海邊傳道,祂看見西門的弟兄安得烈在打魚,耶穌就說:「來跟隨我」,西門,安得烈就放下魚網跟隨主.在這段記述裡面和其它相同的記述裡,耶穌所發出的只有來跟隨我,並沒有提到要拾己,要拿起你的十字架.所以我相信來跟隨我,就是耶穌呼召門徒福音的內容,今天經過保羅和猶太文化的洗禮,我們很少用來跟隨我作為傳福音的信息內容,我們喜歡用信耶穌(保羅的影響).這兩者並沒有矛盾之處,但是因為語言的不同,所帶來的形象,要求和所涉及的哲學都很不同,或許日後有機會‥‥今大不能處理這個問題,「跟隨耶穌」和「信耶穌」這兩個呼召的不同之處.

今天我只處理「跟隨耶穌」,似乎這句說話是耶穌對祂碰到每個人所講的──「來跟隨我」.在祂生命最後的階段,祂加入了悲情的成份,你要捨己,跟隨我的要捨己,要背十字架,然而在其它場合,背景祂沒有作出這個悲情的要求,祂只是對西門彼得,其他很多的人說:「你來跟隨我」.當我讀大學的時候,當時六,七十年代在香港大學有個很大的復興,基督徒團契很興旺,每年都有夏令會,我記得夏令會最後一天就是呼召,總是用這段經文──「要捨己,背十字架跟隨主」,當時的詮釋就是我們要去做宣教,我們要進入神學院,要去東南亞,要去非洲參與那些宣教團體去宣教.

當時在我們的團契的裡面亦有很多人回應這個呼召,這段經文很有意思,但這是一個呼喚基督徒去受苦,放棄自己.在當時六十年代來說,可以說是菁英的事業,有時我想想這當然是真的,也是對的,所以要我受苦怎可說是喜訊呢?做基督徒有時不能不受苦,可是你要我一直受苦而這還是喜訊呢?我就很難說出口,日後我在香港的新蒲崗,荃灣區的工廠區工作時我就開始明白.原來「來跟隨我」的確有可能是涉及受苦,但它本身來說真的是一個很大的喜訊,這是一個邀請,邀請你進入一個愛的群體；這是一個邀請,邀請你每天的生活是有方向的；這是一個邀請,邀請你生命是有意義的；是一個邀請,是進入一個群體裡面.因為耶穌從來都沒有單單叫一個人去跟從祂,祂叫一群人去跟從祂,跟從耶穌就成為弟兄姊妹,就和耶穌一起,和神一起在世界上做天國的事情.

其實這跟我們的自我理解有關,當時我們在香港大學讀書,別人說我們是菁英,自己或多或少都有點這樣的感覺,還沒畢業就已經有許多公司請你去工作,執著這個身份,這個自我理解,聽到耶穌說:「來跟從我」,這真是一個受苦的呼召.不可以按照自己的心意去賺很多錢,不可以按自己的心意進入政府當公務員,甚至升上部長.確實對於一位自認了不起的人來說,這個呼召真的是一個悲情的呼召.一個很大的挑戰,要你捨己,但是,如果你不在這種情況,如果你是一年三百六十五日中要工作三百六十二日,只有三天是過年時放假的；如果你是一位懷孕的女工,法律給予你有分娩假期,可是這個假期是沒有工資的.你聽到這句說話,「來跟從我」,這就是一個很大的喜訊,是一個邀請,邀請你和其他人跟隨耶穌成為一群愛和公義的群體.拿起你的十字架再不是一個重擔,責任不是一個道德上的問題,責任不是一個法律問題,責任是一個恩典,責任成為生命豐富的原因.

所以為什麼「來跟從我」是一個喜訊,亦是福音的內容?今天我們面對這個世界,我們有什麼訊息?你可以選擇保羅的方法──「信耶穌」,這是好的.但我會選擇──「跟隨耶穌」,「跟隨耶穌」是一個邀請,邀請一個朋友進入這個群體,這句說話是當時耶穌在祂三,四年中事奉的理解,今天我們讀這段經文,用今天的背景,今天的解說,我們怎樣理解呢?我相信我們每一個人都需要自己去問神,問耶穌,你要我跟隨你是什麼意思?我是不能代你回答.

但有一件事我很清楚知道的,就是當你跟隨耶穌的時候,你同時亦是不跟隨權勢,你不用作聲的,當你走向耶穌十字架的時候,你就不是走向權勢的,當你走向跟隨耶穌的時候,走向十字架的方向的時候,你就不是走向金錢瑪門那個方向,你不用選擇,不用作出一個有意識的選擇說:「我拒絕權勢,我拒絕金錢」,當你有意識地,認真地選擇跟隨耶穌,你已經是很認真地拒絕了權勢,拒絕了瑪門.因為你不能同時走向兩個方向,我相信這就是學術界(從事文字工作的人)最大的危險,我們很喜歡說正面的說話,我們很怕說反面的說話,我們很喜歡說祝福,我們很怕說有禍,在福音書裡耶穌基督講八福,隨後祂就說七禍；祂不像今天許多教會的領袖專講祝福,在他生命上好像完全沒有咒詛的,專講祝福,這個祝福裡面是沒有任何拒絕的成份.同樣「和階」,「穩定」都是很美麗的文字,每個人都想要的,我都想要,我亦舉手贊成的,但當你說我相信「和諧」,我相信「穩定」時,你同時亦在說我不相信什麼什麼的.

就如我上次在這個禮拜堂一個婚禮上講道,那個婚禮中的誓言,我xxx以你xxx為我終身的伴侶,這是一個肯定,一個確定,亦都是一個認信,而這個肯定當中實則包含了很多的否定,即是說這兩位戀人在這個認信當中同時向全世界其他的俊男,美女否定你不是我生命中所愛的,你不是我的妻子,你不是我的丈夫.

任何的肯定必然存在這個拒絕的成份,若然不是那個肯定就是假的,我們的肯定,我們的認信要由我們的否定來証實.我信全能主,亦等於我不信權勢,我不信瑪門,如果你那個認信我信全能主不包括這方面的拒絕,這個認信──「我信全能主」是沒有意思的,今天我就是用這個來作出我對馬太十六章這段經文對於今天的意義.

今天是一個大時代,雖然是一個很細小的香港,818事件作為香港大學的校友我深受感動,我被徐立之感動,我被李同學感動,我被當天參加這個聚會的人感動.我感覺到這是一個大時代,所以我們要拜一個大時代的神,讓我們不要將我們的神固限於我們的事業,家庭.事實上我們的家庭很重要,我愛我的妻子四十六年,如果今天她病了我就不會來這裡,這個先後次序在我生命中很清楚的,但讓我們不要將神固限於事業；個人的事業；個人的家庭幸福.這些全都是神所寶貴的,是我們所寶貴的.

在大時代我們要拜一個大時代的神,當我們六十年代,七十年代時我們看什麼書呢?我們看英國【基督徒學聯團契(ＩＶＰ)】出版的《鐵證待判》一書,早前　史陶德(John R.W.Stott)逝世,他就是那個時代的人,當時有一本很出名的書,大家都看的,就是　李德爾(Paul Little)的《你的神太小了》(Your God is Too Small),他當時推動的方向就是宣教,離開自己舒適的環境,進入未聽聞福音的地方.今天我們可以有不同的詮釋,但讓我們所信的神是一位大的神,大時代的神,讓我們不要使我們的神變得渺小.

所以今天以這經文大家互相勸勉,我們跟隨耶穌是有悲情的成份,同時我亦看到一個大的喜訊,這是一個邀請,邀請我們進入一個群體,同基督一起在這個時代,這個世代與神一起作神的大事.



\newpage

\section{七十個七次}
\label{sec:118}
\textbf{林豪恩先生 - 七十個七次}
\newline
\newline
\hyperref[sec:117]{< < < PREV SERMON < < <}
~
\hyperref[sec:toc]{[返目錄]}
~
\hyperref[sec:119]{> > > NEXT SERMON > > >}


馬太福音 18:21-35
\newline
\begin{longtable}{cl}
\hline
\hline
章節 & 經文 (和合本修訂版)\\
\hline
18:21 & \begin{tabularx}{0.7\textwidth}{X} 那時，彼得進前來，對耶穌說：「主啊，我弟兄得罪我，我當饒恕他幾次呢？到七次夠嗎？」 \end{tabularx} \\ \\ \relax
18:22 & \begin{tabularx}{0.7\textwidth}{X} 耶穌說：「我告訴你，不是到七次，而是到七十個七次。 \end{tabularx} \\ \\ \relax
18:23 & \begin{tabularx}{0.7\textwidth}{X} 因為天國好像一個王要和他僕人算賬。 \end{tabularx} \\ \\ \relax
18:24 & \begin{tabularx}{0.7\textwidth}{X} 他開始算的時候，有人帶了一個欠一萬他連得的僕人來。 \end{tabularx} \\ \\ \relax
18:25 & \begin{tabularx}{0.7\textwidth}{X} 因為他沒有甚麼償還之物，主人下令把他和他妻子兒女，以及一切所有的都賣了來償還。 \end{tabularx} \\ \\ \relax
18:26 & \begin{tabularx}{0.7\textwidth}{X} 那僕人就俯伏向他叩頭，說：『寬容我吧，我都會還你的。』 \end{tabularx} \\ \\ \relax
18:27 & \begin{tabularx}{0.7\textwidth}{X} 那僕人的主人就動了慈心，把他釋放了，並且免了他的債。 \end{tabularx} \\ \\ \relax
18:28 & \begin{tabularx}{0.7\textwidth}{X} 那僕人出來，遇見一個欠他一百個銀幣的同伴，就揪著他，扼住他的喉嚨，說：『把你所欠的還我！』 \end{tabularx} \\ \\ \relax
18:29 & \begin{tabularx}{0.7\textwidth}{X} 他的同伴就俯伏央求他，說：『寬容我吧，我會還你的。』 \end{tabularx} \\ \\ \relax
18:30 & \begin{tabularx}{0.7\textwidth}{X} 他不肯，卻把他下在監裡，直到他還了所欠的債。 \end{tabularx} \\ \\ \relax
18:31 & \begin{tabularx}{0.7\textwidth}{X} 同伴們看見他所做的事就很悲憤，把這一切的事都告訴了主人。 \end{tabularx} \\ \\ \relax
18:32 & \begin{tabularx}{0.7\textwidth}{X} 於是主人叫了他來，對他說：『你這惡奴才！你央求我，我就把你所欠的都免了； \end{tabularx} \\ \\ \relax
18:33 & \begin{tabularx}{0.7\textwidth}{X} 你不應該憐憫你的同伴，像我憐憫你嗎？』 \end{tabularx} \\ \\ \relax
18:34 & \begin{tabularx}{0.7\textwidth}{X} 主人就大怒，把他交給司刑的，直到他還清了所欠的債。 \end{tabularx} \\ \\ \relax
18:35 & \begin{tabularx}{0.7\textwidth}{X} 你們各人若不從心裡饒恕你的弟兄，我天父也要這樣待你們。」 \end{tabularx} \\ \\
[1ex]
\hline
\hline
\end{longtable}
日期:2011年10月2日

題目:七十個七次

經文:馬太福音18章21-35節

證道:林豪恩先生

18 : 21那時,彼得進前來,對耶穌說:「主啊,我弟兄得罪我,我當饒恕他幾次呢?到七次夠嗎?」 18 : 22耶穌說:「我告訴你,不是到七次,而是到七十個七次. 18 : 23因為天國好像一個王要和他僕人算賬. 18 : 24他開始算的時候,有人帶了一個欠一萬他連得的僕人來. 18 : 25因為他沒有甚麼償還之物,主人下令把他和他妻子兒女,以及一切所有的都賣了來償還. 18 : 26那僕人就俯伏向他叩頭,說:『寬容我吧,我都會還你的.』 18 : 27那僕人的主人就動了慈心,把他釋放了,並且免了他的債. 18 : 28那僕人出來,遇見一個欠他一百個銀幣的同伴,就揪著他,扼住他的喉嚨,說:『把你所欠的還我!』 18 : 29他的同伴就俯伏央求他,說:『寬容我吧,我會還你的.』 18 : 30他不肯,卻把他下在監裡,直到他還了所欠的債. 18 : 31同伴們看見他所做的事就很悲憤,把這一切的事都告訴了主人. 18 : 32於是主人叫了他來,對他說:『你這惡奴才!你央求我,我就把你所欠的都免了； 18 : 33你不應該憐憫你的同伴,像我憐憫你嗎?』 18 : 34主人就大怒,把他交給司刑的,直到他還清了所欠的債. 18 : 35你們各人若不從心裡饒恕你的弟兄,我天父也要這樣待你們.」

一．無限量的饒恕

18 : 21那時,彼得進前來,對耶穌說:「主啊,我弟兄得罪我,我當饒恕他幾次呢?到七次夠嗎?」 18 : 22耶穌說:「我告訴你,不是到七次,而是到七十個七次.

饒恕雖然重要,但總要有個限度,彼得希望耶穌為饒恕定下具體次數,讓他們「有法可依」.為甚麼彼得建議七次呢?他知道耶穌的要求比文士高,若文士認為饒恕三次已經仁至義盡,彼得提出七次這個代表完滿的數目,心想應該會得到耶穌的讚賞.出乎意料之外,耶穌的要求竟然是彼得提出的七十倍.七十個七是簡單的數學題,但門徒明白耶穌的意思並非要他們計算饒恕的次數,而是打破門徒心裡面「饒恕總有個限量」的假設,告訴他們真正的饒恕是沒有限量的.為甚麼無限量地饒恕呢?我們可以從醫學的角度說明饒恕對身體的好處,也可以從心理學的角度展示饒恕對心理的好處.現在,讓我們透過耶穌與門徒在上下文的對話,探討耶穌要求門徒饒恕七十個七次的原因.

二．因為耶穌不願意失去信仰群體中的任何一位如果我們留意上文,彼得提出這個問題的時候,耶穌正與門徒談論教會紀律.

18 : 15「若是你的弟兄得罪你,你要去,趁著只有他和你在一起的時候,指出他的錯來.他若聽你,你就贏得了你的弟兄； 18 : 16他若不聽,你就另外帶一個或兩個人同去,因為『任何指控都要憑兩個或三個證人的口述才能成立』. 18 : 17他若是不聽他們,就去告訴教會；若是不聽教會,就把他看作外邦人和稅吏.

任何群體都難免有人與人之間的衝突,耶穌預計這群跟隨他的人也不會例外.動物面對衝突的自然反應通常是「打或走Fight or Flight」.人類可能會用比較文明的方法,試圖令對方消失,如果不能夠令對方消失,那麼就自己消失.只要衝突的任何一方消失,就不需要面對衝突了.然而,若果要其中一方消失,即是說要群體失去成員.對於某些性質的群體,成員離開沒甚麼大不了,例如某些商業機構,以減低成本的方法去增加利潤,就算你盡忠職守,公司也要趕你走.問題不是你對公司有沒有貢獻,而是在這一刻,公司是否需要你,當公司需要你的時候,就說你是公司的寶貴資產,當公司不需要你的時候,你就被說成公司的負累.耶穌教導他的門徒,跟隨他的信徒群體並不是這樣看待群體中的成員,耶穌不願意失去群體中的任何一位,用保羅的說話,就是說身體不能失去任何肢體.如果當時耶穌與門徒正在討論教會紀律,再看看上文耶穌的言論,我們可以合理假設,耶穌是在信徒群體不能失去任何一位成員的意願上去討論教會紀律.

18 : 6「凡使這些信我的小子中的一個跌倒的,倒不如把大磨石拴在這人的頸項上,沉在深海裡. 18 : 7這世界有禍了,因為它使人跌倒；絆倒人的事是免不了的,但那絆倒人的有禍了! 18 : 8如果你一隻手或是一隻腳使你跌倒,就把它砍下來扔掉.你缺一隻手或是一隻腳進入永生,比有兩手兩腳被扔進永火裡還好. 18 : 9如果你一隻眼使你跌倒,就把它挖出來扔掉.你只有一隻眼進入永生,比有兩隻眼被扔進地獄的火裡還好.」 18 : 10「你們要小心,不可輕看這些小子中的一個；我告訴你們,他們的天使在天上,常見我天父的面.

在一個群體中,總有些是「多你一個唔多,少你一個唔少」的小子,這些小子的「小」,未必因為他們身材短小,可能是他們人微言輕,因此,群體可能未必留意到他們,或者不把他們放在眼內.然而,耶穌卻指出,天父留意他們,也把他們放在眼內,因此,縱使他們不起眼,群體不能「當他們無到」,不能因為他們「小」而輕看他們,不能因為他們「軟弱」(保羅用語)而絆倒他們.

18 : 12「一個人若有一百隻羊,其中一隻走迷了路,你們的意見如何?他豈不留下這九十九隻在山上,去找那隻迷路的羊嗎? 18 : 13若是找到了,我實在告訴你們,他為這一隻羊歡喜,比為那沒有迷路的九十九隻歡喜還大呢! 18 : 14你們在天上的父也是這樣,不願意失去這些小子中的一個.」 耶穌教導門徒,不能絆倒信仰群體中的一個小子,否則天父會追究.然而,就算該小子不是被絆倒,而是自己走失,也不等於群體可以就手旁觀.耶穌並沒有與門徒討論迷羊走失的原因,也不是討論走失的責任屬誰,而是強調要去找那隻走失的迷羊回來.為甚麼為了一隻迷羊如此費神呢?他對群體有何貢獻呢?耶穌把原因說得清楚不過了:「你們在天上的父不願意失去這些小子中的一個.」

如果耶穌不願意信徒群體失去任何一位成員,我們當如何理解耶穌接著對門徒如何處理信徒群體間的衝突的教導呢?

耶穌說,第一,「

若是你的弟兄得罪你,你要去,趁著只有他和你在一起的時候,指出他的錯來.他若聽你,你就贏得了你的弟兄.」(v.15)通常人們對耶穌的教導都是聽一半,不聽一半.聽的那一半是指出得罪我們的那個人的錯,不聽的那一半是我們對所有人指出那人的錯,唯獨不接觸當事人.我們這樣做,很可能是要證明自己對,對方錯；但耶穌教導門徒要直接溝通,而不是背後指責,目的不是證明誰對誰錯,而是「贏得你的弟兄」,避免一拍兩散.

第二,「

他若不聽,你就另外帶一個或兩個人同去,因為『任何指控都要憑兩個或三個證人的口述才能成立』.」(v.16)耶穌承接舊約的教導,不容許我們私下判斷,也不可以執行私刑.為了好好解決紛爭,應該要找別人協助.在第三者的參與之下,衝突可能以兩個當事人沒有想像到的方式解決.

第三,「

他若是不聽他們,就去告訴教會；若是不聽教會,就把他看作外邦人和稅吏.」(v. 17)如果得罪你的弟兄連教會也不理睬,是否就表示我們可以把這人逐出教會呢?這可能是我們的意思多於耶穌的意思,因為我們都不想「麻煩人」留在群體中,若這些人消失,大家都舒服得多.如果耶穌在這裡指示教會驅逐這些成員,豈不是與他之前的教導──他不願意失去任何一位有矛盾?而且我們不要忘記,耶穌開始這個話題時,就定下了目標:「贏得你的弟兄」.那麼,在這樣的前文後理之下,我們如何理解耶穌「把他看作外邦人和稅吏」這指引呢?「把他看作外邦人和稅吏」是否等於清理門戶呢?是否等於驅逐出教呢?回答這個問題,或者我們需要回想一下耶穌是如何對待外邦人和稅吏的.當我們以這個角度回望,我們會發現,耶穌經常接觸他們,並且因此而被當時的宗教領袖指責他與稅吏等人混在一起.當我們向前望,我們會看見那位拒絕接觸外邦人的彼得,上帝要用特別異象令他願意進入外邦人的家；而保羅就想盡辦法令外邦人可以成為耶穌的門徒,甚至因此而被他的猶太同胞誤會.如果說耶穌的意思是吩咐教會與那些連教會也不聽的信徒割蓆,不如說耶穌吩咐教會去尋找他們,接觸他們,把他們帶回群體中.

18 : 18「我實在告訴你們,凡你們在地上所捆綁的,在天上也要捆綁；凡你們在地上所釋放的,在天上也要釋放. 18 : 19我又實在告訴你們,若是你們中間有兩個人在地上同心合意地求甚麼事,我在天上的父必為他們成全. 18 : 20因為,哪裡有兩三個人奉我的名聚會,哪裡就有我在他們中間.」

耶穌賦予給教會的權責實在非同小可,教會自己不可輕看,因此,教會必須正確地按著耶穌基督的教導去運用被賦予的權力和承擔責任.在這樣的背景之下,彼得問:「主啊,我弟兄得罪我,我當饒恕他幾次呢?到七次夠嗎?」我們就明白為甚麼耶穌回答:「我告訴你,不是到七次,而是到七十個七次.」因為耶穌不願意跟隨他的群體失去任何一人.

三．因為耶穌不願意我們失去自己無限量的饒恕,除了因為耶穌不願意失去信徒群體中任何一人之外,也因為耶穌不願意我們失去自己.

18 : 23因為天國好像一個王要和他僕人算賬. 18 : 24他開始算的時候,有人帶了一個欠一萬他連得的僕人來. 18 : 25因為他沒有甚麼償還之物,主人下令把他和他妻子兒女,以及一切所有的都賣了來償還. 18 : 26那僕人就俯伏向他叩頭,說:『寬容我吧,我都會還你的.』 18 : 27那僕人的主人就動了慈心,把他釋放了,並且免了他的債. 18 : 28那僕人出來,遇見一個欠他一百個銀幣的同伴,就揪著他,扼住他的喉嚨,說:『把你所欠的還我!』 18 : 29他的同伴就俯伏央求他,說:『寬容我吧,我會還你的.』 18 : 30他不肯,卻把他下在監裡,直到他還了所欠的債. 18 : 31同伴們看見他所做的事就很悲憤,把這一切的事都告訴了主人. 18 : 32於是主人叫了他來,對他說:『你這惡奴才!你央求我,我就把你所欠的都免了； 18 : 33你不應該憐憫你的同伴,像我憐憫你嗎?』 18 : 34主人就大怒,把他交給司刑的,直到他還清了所欠的債. 18 : 35你們各人若不從心裡饒恕你的弟兄,我天父也要這樣待你們.」

這個比喻中的僕人欠下主人巨債,本來要全家賣為奴隸,但蒙主人赦免所有債項,恢復自由.然而,這位被主人赦免巨債的僕人,欲不肯寬容他的同伴相對來說的一點小債.主人知道後,收回當初給予他的赦免,把應當的刑罰還給他.耶穌總結這個比喻說:「你們各人若不從心裡饒恕你的弟兄,我天父也要這樣待你們.」

我們當初以為,不饒恕只是會失去對方,耶穌藉著這個比喻揭露給我們知道,不饒恕原來更會令我們失去自己.這不是你贏我輸,或者你輸我贏的遊戲,而是雙方都輸的遊戲.耶穌不願意失去信仰群體中的任何一人,耶穌也不願意我們失去自己.耶穌珍惜一個小子,珍惜一隻迷羊,珍惜稅吏和外邦人,當然也珍惜你,他怎會忍心看到你失去自己呢?因此,他回答彼得說:「我告訴你,不是到七次,而是到七十個七次」,不只是表達對彼得的要求,同時間也表達對彼得的關懷.我們聽到嗎?我們聽到耶穌要求我們饒恕的同時,正正是表達對我們的珍惜嗎?

四．饒恕的鑰匙耶穌教導門徒,若果我們不能失去對方,也不能失去自己,我們就需要饒恕.然而,如何打開饒恕之門呢?開啟饒恕之門的鑰匙是甚麼呢?

18 : 24有人帶了一個欠一萬他連得的僕人來. 18 : 25因為他沒有甚麼償還之物,主人下令把他和他妻子兒女,以及一切所有的都賣了來償還. 18 : 26那僕人就俯伏向他叩頭,說:『寬容我吧,我都會還你的.』 18 : 27那僕人的主人就動了慈心,把他釋放了,並且免了他的債.

是僕人「寬容我吧,我都會還你的」這償還的承諾開啟了饒恕之門嗎?不可能吧!僕人欠下主人的巨債等於當時一個工人16萬年的人工,主人當然知道這僕人任何清還債務的承諾都不可能對兌.那麼,主人為甚麼釋放他,並且免了他的債項呢?經文說:「那僕人的主人就動了慈心」.開啟饒恕的鑰匙,並不是償還的承諾,而是動了慈心,不忍對方受苦.那麼,不肯饒恕的原因又是甚麼呢?

18 : 28那僕人出來,遇見一個欠他一百個銀幣的同伴,就揪著他,扼住他的喉嚨,說:『把你所欠的還我!』 18 : 29他的同伴就俯伏央求他,說:『寬容我吧,我會還你的.』 18 : 30他不肯,卻把他下在監裡,直到他還了所欠的債. 18 : 31同伴們看見他所做的事就很悲憤,把這一切的事都告訴了主人. 18 : 32於是主人叫了他來,對他說:『你這惡奴才!你央求我,我就把你所欠的都免了； 18 : 33你不應該憐憫你的同伴,像我憐憫你嗎?』

一般相信被饒恕的人會饒恕別人,這個比喻中的僕人告訴我們,不一定的.如果說主人饒恕僕人是因為動了慈心,僕人對同伴的不饒恕就是因為沒有憐憫.比喻刻意把僕人自己的巨債獲得赦免,而不寬待同伴的小債作對比,使聽眾認同比喻中旁觀者的反應:「同伴們看見他所做的事就很悲憤」,而主人的指責:「你不應該憐憫你的同伴,像我憐憫你嗎?」說出了饒恕的鑰匙.

彼得應該聽得出耶穌說這個比喻的意思,如果把他對其他人的饒恕與上帝對他的饒恕比較,實在微不足道,七十個七次又算得甚麼呢?如果他深刻地體會到上帝對他的憐憫,他怎能對弟兄收起憐憫之情呢?我們又聽得明白這個比喻的意思嗎?但願這個比喻把我們的注視從得罪我們的人的錯轉向上帝對我們饒恕,從計算別人的惡轉向體會上帝對我們的憐憫.

五．總結馬太福音18章在信徒群體或教會的處境之下討論饒恕.今天的舊約經課(創世記50:15-21)則把這個課題放在一個家庭的處境.饒恕放於家庭的處境中,可能令我們有更多的感受及聯想.

50 : 15約瑟的哥哥們見父親死了,就說:「也許約瑟仍然懷恨我們,會照我們從前待他一切的惡,重重報復我們.」 50 : 16他們就傳口信給約瑟說:「你父親未死之前曾吩咐說: 50 : 17『你們要對約瑟這樣說:從前你哥哥們惡待你,你要饒恕他們的過犯和罪惡.』現在求你饒恕你父親的神之僕人們的過犯.」他們對約瑟說了這話,約瑟就哭了. 50 : 18他的哥哥們又來俯伏在他面前,說:「看哪,我們是你的奴隸.」 50 : 19約瑟對他們說:「不要怕,我豈能代替神呢? 50 : 20從前你們的意思是要害我,但神的意思原是好的,要使許多百姓得以存活,成就今日的光景. 50 : 21現在你們不要害怕,我必養活你們和你們的孩子.」於是約瑟安慰他們,講了使他們安心的話.

無論過去有多少恩怨情仇,上帝讓約瑟清楚知道,如果他不饒恕他的兄弟,他就會失去他的兄弟,同時也會把自己埋葬於怨恨之中；唯有饒恕他的弟兄,他的父家才能夠存活下去.福音經課和舊約經課都暗示,饒恕不是錦上添花的事情,而是生存的必須.如果有一天我們在掙扎著饒恕與否,讓我們思想一下約瑟的領悟:「神的意思原是好的,要使許多百姓得以存活」,學習約瑟說:「我豈能代替上帝呢?」



\newpage

\section{關於權威的討論}
\label{sec:119}
\textbf{鄧瑞強博士 - 關於權威的討論}
\newline
\newline
\hyperref[sec:118]{< < < PREV SERMON < < <}
~
\hyperref[sec:toc]{[返目錄]}
~
\hyperref[sec:120]{> > > NEXT SERMON > > >}


馬太福音 21:23-27
\newline
\begin{longtable}{cl}
\hline
\hline
章節 & 經文 (和合本修訂版)\\
\hline
21:23 & \begin{tabularx}{0.7\textwidth}{X} 耶穌進了聖殿，正教導人的時候，祭司長和百姓的長老來問他：「你仗著甚麼權柄做這些事？給你這權柄的是誰呢？」 \end{tabularx} \\ \\ \relax
21:24 & \begin{tabularx}{0.7\textwidth}{X} 耶穌回答他們說：「我也要問你們一句話，你們若告訴我，我就告訴你們我仗著甚麼權柄做這些事。 \end{tabularx} \\ \\ \relax
21:25 & \begin{tabularx}{0.7\textwidth}{X} 約翰的洗禮是從哪裡來的？是從天上來的，還是從人間來的呢？」他們彼此商議說：「我們若說『從天上來的』，他會對我們說：『這樣，你們為甚麼不信他呢？』 \end{tabularx} \\ \\ \relax
21:26 & \begin{tabularx}{0.7\textwidth}{X} 若說『從人間來的』，我們又怕眾人，因為大家都認為約翰是先知。」 \end{tabularx} \\ \\ \relax
21:27 & \begin{tabularx}{0.7\textwidth}{X} 於是他們回答耶穌：「我們不知道。」耶穌也對他們說：「我也不告訴你們，我仗著甚麼權柄做這些事。」 \end{tabularx} \\ \\
[1ex]
\hline
\hline
\end{longtable}
日期:2011年10月9日

題目:關於權威的討論

經文:馬太福音 21:23-27

證道:鄧瑞強博士

各位弟兄姊妹,主內平安.

今日講道的經文,是今日的福音書經課:馬太福音 21:23-27.在這裡,祭司長和民間的長老與主耶穌展開一番討論,是一場關於權威的討論.祭司長和民間的長老不同法利賽人.法利賽人基本上關心信仰和靈性.祭司長和民間的長老卻是手握權力的人,他們不是民間宗教慈善團體的人,他們是當地行政會議的人.他們手握權力,也關心權力.他們與主耶穌展開一場關於權力來源的討論.他們問了一個問題:「你仗著什麼權柄做這些事?給你這權柄的是誰呢?」(太21: 23)他們問:你有什麼權這樣做?

近日,在網上流行一段短片.短片中,一個小孩子在一個商場裡,大聲責罵他的父親.引來很多遊人圍觀.這小孩子和他父親展開一段有關權力的爭論.這小孩子大致說:「我不喜歡來這裡,你無權硬要我來?...我有人權....我有權報警....」簡言之,這小朋友對他爸爸說:你仗著什麼權柄做這些事?給你權柄的是誰?

這小朋友強調的是個體的人權,聽他說來,這是絕對的權威.的確,我們要尊重每個個別的人.但是,每個個別的人也生活在一個合宜的群體裡.「人權」的重點,不是叫每個人最終變成互不相關的孤島,而是在互相尊重裡面結成互愛的群體.這個小朋友需要個體人權,但其實他也需要他的家庭支持他的存在.這個小朋友講的人權,看來是指自己想如何便如何的自由,他能以這自由之名,公開羞辱他的爸爸嗎?他能視他的家庭關係與他完全無關嗎?他所依賴的權威,最終建立他的生命,抑或摧毀他的生命?

近日,某國家高層訪港,引發不少事件.有居民因為穿著印有敏感政治字眼的Ｔseuk ,被黑衣人抬走.這件事留下不少「黑影」.不知什麼力量,將「真實的人」變成「黑影」.也不知什麼力量,將「黑衣人」留下的另類「黑影」抹走.光天化日之下,黑影處處,殊不尋常.黑衣人可能會問這居民:「你憑什麼權柄穿這Ｔseuk ,你有什麼權這樣做?」當然,這居民也可反問:「你仗著什麼權柄做這些事?給你權柄的是誰?」

難道,這個居民是闖入了某種「聖殿」,某種被封鎖線封鎖的「權力祭壇」,做了些衝擊權威的動作,因而被黑衣人抬走?這麼一個微少的居民,穿了這麼一件卑微的Ｔseuk ,竟然觸動了權威人士的神經.一件普通的Ｔseuk ,看來殊不普通.它不是普通的衣著,而是一隻碰向牆壁的雞蛋.的確,我們可能也會問:這位小巿民,他究竟憑著什麼力量,有足夠勇氣做這事?

主耶穌也做了某些事,引得祭司長和民間的長老問他:「你仗著什麼權柄做這些事?給你這權柄的是誰呢?」

主耶穌做了什麼?

我們要做一點回顧.

我們翻開聖經:馬太福音21: 12-15.

太21:12　耶穌進了神的殿,趕出殿裡一切做買賣的人,推倒兌換銀錢之人的桌子,和賣鴿子之人的凳子；太21:13　對他們說:「經上記著說:我的殿必稱為禱告的殿,你們倒使它成為賊窩了.」太21:14　在殿裡有瞎子,瘸子到耶穌跟前,他就治好了他們.太21:15　祭司長和文士看見耶穌所行的奇事,又見小孩子在殿裡喊著說:「和散那歸於大衛的子孫!」就甚惱怒...

原來,主耶穌走入聖殿,干擾了人家做生意.

主耶穌的行動不能算是暴力,因為在這裡沒有人受傷.其實,在神聖之處做買賣,很多時才是暴力的根源.主耶穌只是趕出聖殿裡做買賣的人,推倒做買賣的台凳.主耶穌關心的是,人在聖殿中,能向神禱告,能體會天地的主的眷顧,能體會人在世上能實踐天父的心意.但是,當這一切美好的東西滲透著買賣時,一切便變質了.這就像美好的愛情滲透著買賣時,愛情便變質了.這看來是愛情,但實質已是沒有愛情的貨物交易了.

主耶穌稱這買賣場所為「賊窩」,在這裡,沒有公平交易.真誠來這裡祈禱的人,他需要獻祭的處境被別人利用為公然剝削的機會.他在這裡,應該體會到天父的眷顧,卻更體會到被欺壓的真實.這是一種買賣的暴力,阻隔人接觸神.

不單如此,在聖殿裡,主耶穌接觸瞎子和瘸子.很奇怪,聖殿裡應是沒有瞎子和瘸子的,這些人是禁止進入聖殿的.可能,這些人跟從主耶穌進入了聖殿.主耶穌將這些不能進入聖殿的人帶入了聖殿,還醫好他們.在聖殿規矩的把關人眼中,這是破壞規矩的.

還有,在聖殿中,有小孩子的歌聲,他們稱頌大衛的子孫.小孩子無名無分,怎能大聲唱頌.在祭司長眼中,社會上被遺棄的人,或社會上無用的人,都不能進入聖殿,都不能唱歌.這是禁止孩童唱歌的地方.

經濟掛帥,排斥殘障者,禁止弱者唱歌,當然不止是昔日的聖殿,也是我們的社會.香港是否「地產霸權」,我不知道,但我知道的是,租金的昂貴,將民間的生機扼殺.我們再難聽到小孩子的歌聲.

祭司長就是聖殿的守護者,但與其說他是守衛著聖殿,不如說他是守衛著聖殿中的買賣.他能守衛這裡的買賣,是因為他同時是政治權力的代理人,守衛著政治權力.這是一個「宗教,經濟,政治」同盟.政治力量強化經濟的剝削,經濟剝削的利益強化著政治,而宗教力量卻神聖化這種合謀.

現在,主耶穌趕走了做買賣的人,解除這同盟的宗教神聖光環,這就表示他面對政治的力量.這不是政治行動,他只是渴望人在聖殿裡能向神祈禱,能尋求神,能實現神的心意.這個單純的渴望,在一扭曲的世界,變成對權威的正面挑戰.

我們回到今日的講道

經文:馬太福音21: 23-27.

太21:23　耶穌進了殿,正教訓人的時候,祭司長和民間的長老來問他說:「你仗著什麼權柄做這些事?給你這權柄的是誰呢?」太21:24　耶穌回答說:「我也要問你們一句話,你們若告訴我,我就告訴你們我仗著什麼權柄做這些事.太21:25　約翰的洗禮是從哪裡來的?是從天上來的?是從人間來的呢?」他們彼此商議說:「我們若說『從天上來,』他必對我們說:『這樣,你們為什麼不信他呢?』太21:26　若說『從人間來』,我們又怕百姓,因為他們都以約翰為先知.」太21:27　於是回答耶穌說:「我們不知道.」耶穌說:「我也不告訴你們我仗著什麼權柄做這些事.」

對於祭司長和民間的長老來說,除了他們那「宗教,經濟,政治」同盟以外,還怎能有別的權威?

有效的管治,是需要「合法性」(legitimacy)的.祭司長和民間的長老認為他們的權力是唯一合法的權力.在這個世界裡,難道「權柄」有另一「合法」的來源?若有,則這「權柄」大得能批判他們現實權力的「合法性」嗎?

他們問主耶穌:「你仗著什麼權柄做這些事?給你這權柄的是誰呢?」

主耶穌卻反問一問題,讓他們思考權柄的終極基礎.

主耶穌問:「約翰的洗禮是從哪裡來的?是從天上來的?是從人間來的呢?」

「約翰的洗禮」基本上是一個叫人悔改的洗禮.約翰說:「天國近了,你們應當悔改.」(太3: 2)他呼喚人,「你們要結出果子來,與悔改的心相稱.」(太3: 8)這種悔改,要求人整個生命的改變.人與人的關係,若已墮落成欺壓的買賣,則需要改變成生命的互相餽贈.應放下現世任何形式的欺壓,要相信有一明察世情的神,他的國度要透過主耶穌的大愛來實現.世上不應有謀取暴利而不理他人死活的買賣.不應有與謀取暴利同謀的政權.不應有宗教將光環加給欺壓人的政治和經濟制度.「約翰的洗禮」要求人的悔改.

主耶穌問祭司長等人,「約翰的洗禮」從哪裡來?

這是迫使他們思考,在他們的權力以外,是否有另一種權力?當然,這不是一般的政治和經濟權力,而是一種要求人轉化生命的呼喚.主耶穌讓祭司長等人想清楚,這種轉化生命的力量,是否來自天上?

聽說,在前蘇聯,有一個人批評政府無力把一個叫人受苦的制度改變成一個叫人活得幸福的制度.這批評使他即時被送到法院.法官問:你知否為何國家現今的制度不能變成你講的那種社會制度?這人答:「這是因為領導人是白痴,沒有想像力.」這人即時被判監二十年.罪名是:洩露國家機密.

祭司長等人,是否有想像力,去設想施洗約翰透過洗禮開展的那個世界?

祭司長等人不敢承認施洗約翰是從天上來的,若他們承認的話,則他們怕主耶穌的反問:「你們為什麼不信他呢?」

的確,祭司長等人不可能相信施洗約翰,他們不可能按約翰的信仰原則處理他們權力範圍內的事.

但是,他們能說施洗約翰的權威來自人間嗎?

他們是這樣認為的,但不夠膽這樣說,因為民意認為施洗約翰是從天上來的.

民意有時是很奇怪的.民意有時不滿壓迫人的現況,但又支持維持現況的背後的權力結構.民意有時不滿壓迫人的買賣制度,但同意卻支持令這買賣制度成功的結構.民意會反對地產霸權,但同時民意卻會以買到樓,賺到錢為榮.民意支持施洗約翰,因施洗約翰批判當時腐敗了的道德生活；但民意也支持釘死主耶穌,因主耶穌看來動搖了使腐敗的道德生活賴以維持的制度.

民意常變,但民意支持施洗約翰,認為他是從天上來的.

主耶穌問祭司長等人,「約翰的洗禮」是從哪裡來的?

祭司長等人的思考,全無面對「約翰的洗禮」從哪裡來的事實,他們只是盤算如何回答得到最大利益,最能維持他們的權力.

他們無言以對.

太21:27　於是回答耶穌說:「我們不知道.」耶穌說:「我也不告訴你們我仗著什麼權柄做這些事.」

祭司長等人的權力,毫無反省力,毫無想像力,他們只想維持聖殿裡的賊窩般的買賣,他們不想見到社會上的弱者,他們不想聽到孩子頌讚救主的歌聲.這樣的權力,怎能明白主耶穌.

我們又如何呢?

我們能否理解到主耶穌的另一種權威?他像施洗約翰一樣,呼喚我們悔改,並過一種與悔改相稱的生活.

但願榮耀歸於聖父,聖子,聖靈.阿們.



\newpage

\section{基督徒的品牌}
\label{sec:120}
\textbf{梁元生教授 - 基督徒的品牌}
\newline
\newline
\hyperref[sec:119]{< < < PREV SERMON < < <}
~
\hyperref[sec:toc]{[返目錄]}
~
\hyperref[sec:121]{> > > NEXT SERMON > > >}


日期:2011年10月16日

題目:基督徒的品牌

證道:梁元生教授

各位弟兄姊妹,早安.

今天與大家思考一個非常功利,庸俗的題目,換句話說也是一個非常「港氣」的題目,究竟基督徒這個身份,招牌值多少錢?假使我們不用金錢來估量,可以換另一個角度來問:究竟基督徒是否是一面的品牌?它的價值有多少?──It"s a question of brand name. And how much is the ``Christian brand" worth? ──香港才會問這些問題.

提到這個品牌,先讓我考考大家對這個品牌的認識.我手上有兩個清單,一張清單上面有十個名字,大家應該都會很熟悉的；另外一張清單,可能有些人很懂,但怎麼讀都讀不清楚,我試一試讀出來給大家認識一下.

第一個清單:Apple,google,amazon,BMW,Mircosoft,coca-cola,Walmart,Mercedes-Benz,Toyota,facebook.這些是2011年世界性十大品牌 (Top World Brand),是我從Business Week 裡面取得的.這些品牌本身就很值錢了,這個價值也不是靠公司出產的產品,貨品來衡量的,就是這個標誌 (logo),品牌的本身,在國際上是非常值錢,但這個會活動的.我看六年前(即2005年),十大品牌的公司有點不同,有些相同,例如:coca-cola.沒有的像是IBM,GE,已不再2011年的清單裡面了,現在有許多被跨網公司所取代了.

另外一張清單:Louis Vuitton,Herm and egrave;s,Gucci,Chanel,Hennessy,Rolex,Mo and euml;t  and amp; Chandon,Cartier,Fendi,Tiffany  and amp; Co.,Christian Dior.這張清單大家都很熟悉,是高價的品牌 (Luxury product──有一些還被稱呼為 most powerful brand).當然我們一般來說,都稱這些為時尚品牌 (fashion brand),我相信在座當中,許多人對這些品牌的認知,比我還要多.有一些人認為我羡慕名牌,因為有時候,他們看到人家送我的衣服裡面,有三個字母"YSL".告訴大家一個秘密,這個YSL不是Yves Saint Laurent,那個名牌啦,我慣用這個YSL是因為那個就是我的名字,一直以來,我書寫的名字就是YSL.

今天為什麼從名牌講起呢?就要回到幾個星期之前,崇基有個年度的開學禮,我與新進崇基大學的新同學說:你們能夠進來中文大學,又在許多書院裡頭選擇了崇基,是非常難得及不容易的,因為崇基在中大九個書院裡頭,成為了一面的品牌,同學們對它的品質有一定的信心.我們要建立這個品牌,打造這個品牌,是需要經過許多年的努力,和許多人一起合作的成果.今天我不談崇基這個品牌,但從打造崇基這個品牌這一個點上,我想到一個相關的問題,就是今天我們基督教或是基督徒,會否成為一面的品牌活在人的面前?一面良好的品牌?在社會人士的眼光當中,我們是否是一面的品牌?在基督徒自己的眼光之中,是否也是如此?直接地問:我們怎麼看基督徒這個身份?是作為一種身份的象徵,是具有良好品質的代號,使您引以為傲,為榮的記號,還是一個不敢見人的東西呢?基督徒會否是一個品牌呢?

首先,我們認識到的基督徒是指那些信奉基督教的人,那些跟隨基督耶穌的人.早期基督徒是指一小部份受到耶穌基督呼喚的人,包括十二個門徒,後來有成千上萬跟隨祂的人,其中有漁夫,文士,稅史,婦女,後來當然也有保羅,提摩太,以及從耶路撒冷,猶大全地,撒瑪利亞直到地中海一帶,有許多跟隨耶穌的人,他們就被稱為「基督徒」.到今天,「基督徒」泛指很多,很多跟隨信仰耶穌基督的人,也包括當中的你和我.但「基督徒」就是那些口中說「跟隨耶穌基督的人」嗎?耶穌自己曾經說過:那些口稱耶穌為主的,到了最後(審判日子)的時候,耶穌就對一部份的人說「

我不認識你們」(太7:22-23).就是說一個被稱為「基督徒」的人,並非是他嘴巴上說他是「基督徒」,或者他是猶太人,希利尼人(希臘人),也不因為他是富有的,貧窮的,不是他做了好事,善事,又或者是他愛國,不愛國,如何分辨「真誠的基督徒」與「假冒的基督徒」?怎麼叫他相信基督徒背後是一個良好的品牌呢?

聖經說:一個人被稱為基督徒,是因為他身上戴著基督徒的印記,這個印記,對人而言,就是品牌 (brand).印記這個字,在原文上是 ``brand",就是現在的「品牌」.以弗所書一章說:上帝用印印了我們,並賜聖靈在我們心裡面作憑據(註一).這個印記是有憑據的,也就是說你們留下了作基督徒的印記.保羅的意思是說:每一個基督徒(大概信主以後),身上就應該戴著基督透過聖靈打下來的印記,如果我們身上沒有找到這個印記 (brand) 的話,可能就是一個膺品.因此對於基督徒來說,我們今天要談的問題,是一個重要的問題:基督徒身上戴著的印記 (The branding of Christiasn)──基督的印號──這是一個非常重要的事情.

剛才我們提到的名牌／品牌,那個背後有許多膺品.今天我們面對一個嚴肅的問題,基督徒是否有真的,假的,有原裝的,有偽冒的,有打印的,沒打印的?這樣子的話,我們看看自己吧!看看自己頭上,身上有甚麼樣子的「印」?

聖經告訴我們,每一個基督徒的身上,都應該有印的──有上帝的印,有聖靈的印,有基督的印記.保羅說:「

你們既聽見真理的道,就是那叫你們得救的福音,也信了基督,既然信他,就受了所應許的聖靈為印記」(弗1:13).

既然我們認為真正的基督徒身上有耶穌基督的印記,那印記是怎麼樣的一個印記呢?是可以看見的「水印」,還是肉眼所不能見的印呢?我希望大家看看自己,看看對方,旁邊的人,看他們有沒有「印」?我們從兩方面看基督徒打印記這回事.

第一,打印對基督徒來說,並不表示「名牌」.反過來看,「打印」在最早的時候,是一個羞辱的記號──它不是一個被公認為好東西的品牌,不是優質的代表,不是名牌的象徵.在西方羅馬的時代,我們知道被打印者是奴隸.奴隸多半在額上,臉上打印,好表示這個人的身份,以讓其他人認出來.古代的中國,也是如此.在春秋戰國時代,墨家的墨翟臉上有印,所有弟子臉上都有印,表示他們是奴隸出身；漢朝的時候也打印,那叫做「刺青」,在犯人的額上印上一個記號,現在叫做「紋身」,同樣表示在社會上,這個人不是好人,是罪犯的意思.大家都很熟悉宋朝的故事──108個梁山伯的好漢──他們大部份的臉上都有刺青,這些人大概都是犯過罪,或許誤判也好,臉上都打上刺青.所以我們說,做個「打印的基督徒」,實際上沒有要把他們的身份提高,做一種名牌的看待,這乃是一個受辱的記號──做基督徒也是甘於作為一個奴隸,受人支配,吩咐,服侍他人的人──最初的意思,非說跟隨主就要離開群眾,低賤,而進入名牌的行列,做一個社會尊敬的人物,反之「打印」就是基督徒要與低微的人,為奴的人同負一軛,要同等,同行,這是基督徒接受印記的一個時代背景,也是我們今天談到基督徒做為一種品牌的初衷(原來的本意).

第二,另外一個意思,當然就是說,一個受辱的記號,經過時代,歷史後,現在逐漸的成為一個光榮的記號了.基督徒既受了印記,就不以與基督同行為苦,為辱,為恥,至終基督徒靠著基督,聖靈,把我們的身份改變,作為一種榮耀的品牌.

接下來,我們再看一看,基督徒打印的一個過程 (The branding of a Christian is a process) ──是一個慢慢的,漸漸的過程,不是一觸即發,一下子改變我們的身份,我們就從受辱變成榮耀──實際上是一個漸進,漫長的過程(一個慢慢學會像主,像基督的過程).保羅說,基督徒身上要有基督的印記.他的意思就是說,我們怎麼樣把這個印記慢慢地表明出來?根據聖經裡有關保羅的書信,至少有三個方面.

一,有印記的基督徒,是一個要像基督的人──Christ like or Be like Christ──有一個「學」或「效法」的意思.這裡也有兩層的意思,一個是「外貌」的相似.我們慢慢地像基督的外表,但更重要,「像基督」是在性格上,氣度上,品質上像祂,因此不完全在外表──頭髮,面容,衣服上面.聖經很清楚的告訴我們,耶穌基督的外貌不是太多人喜歡去效法的:「

他無佳形美容,我們看見他的時候,也無美貌使我們羨慕他.」(賽53:2)即然如此,我們如果要學像基督,當然是指著基督的人格,個性.也就是說,「效法基督」就是從初信到像基督的一個過程,這不是一天就能完成的,因此是一個漸進,漫長的過程.從信主的那一刻開始,基督徒就應該開始學基督與效法基督.希望經過不斷的努力,改進,最終的目標雖不能完全像祂,至少能夠活得像祂.叫其他人看出他有一個基督的模樣,越來越像基督的過程,是一個不斷追求的態度 (increasingly or progressingly Christ-like).

二,有印的基督徒,是有心的基督徒.要像基督,要有基督的心,保羅說:「

你們當以基督耶穌的心為心.」(腓2:5)除了有外在的行為表現,還要表裡如一──效法,學像 (the act of imitation).有時候,我們只能做到貌像,神似,但不一定源自內心,像許多唱歌的人,希望效法／學貓王 (Elvis Presley),怎樣學他呢?頭髮梳貓王的髮型,穿他的衣服,甚至學他唱歌的聲音,但終究我們學他這個人,但不是貓王.聖經告訴我們要「真」像基督.「真像」不是「貌似」,也不單單「神似」而已,乃是從心裡頭出來的學習,是內外如一的仿傚.要成為一個有印記的基督徒,要在我們心裡頭時刻帶著一顆基督的心.甚麼才是「基督的心」?我認為主要是心裡面常常有心熱,渴慕與追求屬靈的事情,要熱切的服事主,這是一個方面.另一個方面是有推己及人的愛心,渴望與人分享福音,在英文,前一個方面是 passion(熱情),後面是 compassion(愛心).基督的心是兩者皆有.

三,要像基督,有基督的心,還要誇基督.保羅說:「

但我斷不以別的誇口,只誇我們主耶穌基督的十字架．」(加6:14)所謂誇基督,不是要我們自誇,也不是要我們吹噓,吹水.以基督誇口,實際上是以基督為榮的意思 (to take pride in Christ, be proud of Christ).我們知道在古時候,打印 (branding) 是恥辱的記號,例如十字架同樣是一個羞恥的記號,但是因為有了基督,我們常常向人宣示這個記號,它成為了一個榮耀的標誌.作為基督徒,不應該常有難以啟齒,引以為羞的感覺.誇基督一方面不是常在口頭上說:主啊,主啊!稱呼主名的人,另一方面也不是存心難以口宣揚的沉默的基督徒.我個人認為,誇基督是一種態度,是以作為基督徒為光榮的自豪感,而這種作為基督徒的自豪感,是十分重要的,也是許多基督徒現在常常不願意表示出來的.如果你身上穿的是一件名牌的衣服,你駕駛的是一台名車,你會常常向人展示一下.作為一個基督徒,你會不會常常感到自豪?有多少的時候,你會以基督徒的身份感到驕傲?這種是誇基督的自豪感,並不相同於自誇,自驕,乃是一種對自己基督徒身份的自尊,自重與自覺.有這種感覺,你就不會隨風搖擺,飄盪,因這種自豪感所帶來身份的認同與自我的肯定,是基督徒在這個人海浮沉的世界上的一種定位.

如果從以上三方面來看基督徒印記的表現,我可以總結一下:基督徒的生命應該包括三個「P」:Purpose:基督徒的生命是以基督為對象及目的的.Progress:基督徒的生命是不斷變化及進步,越來越像基督的.Pride:基督徒的生命是充滿著自豪感,以主為榮,以基督徒的身份為傲的.今天,我們生活在一個追逐名牌的世界,你,我是否察覺到基督徒的身份是一個尊貴的名牌?因為這個是天父兒子的身份,是一個進入永恆國度的依據.

我希望各位身上帶著基督的印記,又效法基督,以祂的心為心,並以基督為榮.

願大家珍惜這個耶穌基督用寶血買回來的身份,經過歷世歷代聖徒努力的打造,才成為光榮的品牌.

註一:

你們既聽見真理的道,就是那叫你們得救的福音,也信了基督,既然信他,就受了所應許的聖靈為印記．這聖靈,是我們得基業的憑據,〔原文作質〕直等到　神之民〔民原文作產業〕被贖,使他的榮耀得著稱讚.(弗1:13-14)



\newpage

\section{在後的必要在先}
\label{sec:121}
\textbf{陳曉東博士 - 在後的必要在先--再思上主的邊緣思維}
\newline
\newline
\hyperref[sec:120]{< < < PREV SERMON < < <}
~
\hyperref[sec:toc]{[返目錄]}
~
\hyperref[sec:122]{> > > NEXT SERMON > > >}


日期:2011年10月23日

講題:在後的必要在先--再思上主的邊緣思維

證道:陳曉東博士

主內弟兄姊妹,大家好!伍牧師還沒有介紹一件事情,就是我結婚的時候是伍牧師幫我証婚的,哈哈.

今天跟大家分享的題目是:在後的必要在先(註一).

我們一起思考一下,一起默想一下「上主的邊緣思維」.這個字是比較奇怪,我待會兒解釋.在講這個題目的時候,這幾年我有一個新的學習經驗.我讀了很久的神學,從1994年一直到讀完了Boston college(波士頓大學)之後,十幾年的時間,當然也是在牧會,也是在讀神學,可都一直在這個領域裡.但是這幾年,我更想打開自己去聆聽,去學習在不同的領域當中,許許多多很精彩,很傑出的人物對我們生命的一些教導和學習.

幾年前,我看一個音樂錄影帶,看到歌星在唱歌的時候,他的背景特別漂亮；那是一座教堂,但這座教堂是在水上的.我當時就很奇怪,說:這個教堂好漂亮!到底在哪裡可以找到呢?後來我發現這個教堂是在日本,而且是在一座酒店裡.那我就更有興趣想知道,是誰設計這個漂亮的教堂?所以在這個因緣底下,我認識了一位日本很傑出的建築大師,他的名字叫安藤忠雄.

認識了他以後--當然我這個認識是從書本上認識他,我很有興趣想知道這樣的建築師,他是怎樣去思考建築的?他背後的思想,理念是什麼?他用什麼材料去建築?他整個建築想表達什麼意境?對自己的生命是怎樣的信念?這樣的一個建築師帶給我很多的趣味,為什麼?因為當我認識他越多一點,就發現他是一個很有趣的人,因為他在他整個建築的生涯裡,最喜歡使用的建築材料只有一樣,就是清水和水泥.就是這兩樣東西--水和水泥,是他整個許許多多出名的建築物標誌性的特徵.

有一次我看他專輯的時候,他有一個建築物特別吸引了我的注意力.故事是這樣的:有一次,他有一個在出版社工作的朋友去找他,說:安藤忠雄先生,我是一個基督教徒,我們的教堂最近要改建,好希望能請來你這樣的建築大師幫我們設計教堂.安藤忠雄之前己經設計了兩座教堂,一座叫「風之教堂」,另外的就是我之前介紹的「水之教堂」.他聽到計劃之後也蠻有興趣的,可是再聽下去,發覺這個案子是一個非常困難的案子,為什麼?因為這個教會是一個非常小的教會,人數也不過100人,他們的教堂要重建的地方也非常小,在一個住宅林立的地方裡,大概只有100平方米的空間.更要命是這個教堂其實是沒有錢的,告訴安藤的時候說,他們的預算其實很緊絀,教堂能不能蓋起來也很有問題,但是因為安藤先生是一個建築大師,所以把這麼困難的案子找上門來跟他談談,請他想辦法.

所以今天,如果你到日本大阪的話,你一定要去看看這個教堂.這個教堂是安藤忠雄非常出名的教堂,是幾乎在絕境之中建立出來的,它拿了20世紀一個教堂建築大獎,叫「光之教堂」.教堂在1980年左右完成,成為安藤先生非常著名的代表作,為什麼?因為這座教堂,你進去之後,看到的幾乎是空無一物,可是在這個空無一物的水泥盒子裡邊,安藤先生在前邊開了一個十字架的洞,當你整個進去的時候,在這個黑暗的教堂裡,你看到的就是一個光的十字架向你迎面而來.所以,每一個去參觀這教堂的人發現,這座教堂的魅力不是在外面所看到的,而是你要進到裡邊去,當你去摸那非常堅實硬厚的清水混凝土,就是水泥的圍合,當你去感受一個被創造出來的黑暗空間,你一進去的時候,會感受到整個人跟外間的隔絕；但是就在這個時候,當你一看這個牆體的時候,整個陽光就從水平交差的開口裡邊湧流進來,成為出名的「光之十字架」.聖潔,簡單,純淨,震撼,這是安藤忠雄給教堂所帶來的一份禮物.

當然,當我看到這個案子的時候,自己在思考:這個安藤不是一個神學家,但是他用一個建築物去表達一個神學思想的時候,用一個建築語言去模仿,模擬上帝在想什麼的時候,我就想到耶穌所講的這句話:在後的必要在先,在前的必要在後.有時候,我們這些讀了一輩子神學的人跟這樣一位智者比起來的話,可能更是有很多學習的方式,因為在這座簡單的教堂裡,我們看到了一件事,就是我們怎樣去學習上帝思維的方式.

其實從今天看到的福音經文裡面,我們也在學習上帝的思考和行動的方式.我們這個時代,其實是非常迫切需要思考,重新學習上帝思維方式.我們都會承認一點,我們的思考方式決定了我們的生活方式.例如,我們今天的社會非常強調一種觀念,就是個人的權利.它所衍生出來的生活方式,和我們所賦被父輩的生活方式非常不一樣.如果我們仔細去探討這個道理,可能需要很長時間,但是,有一次,我的朋友跟我分享,我們今天這麼強調個人權利,它所衍生出來的問題真叫我們啼笑皆非.事情是這樣的:有一次,我的朋友住在美國的Maryland(馬里蘭),他的女兒在馬里蘭大學(Maryland University)讀書.有一次生急病,因為她自己住在Boston(波士頓),所以這個女兒就簡單地打電話告訴我朋友說:爹,我進了醫院了.她也沒有說快來看我,可是你知道為人父母,心裡急,就馬上趕到醫院去,可是去到醫院竟不能進去.為什麼?因為美國在幾年前通過一條法案,叫個人私隱法,就是說如果你要到醫院去探望病人的話,哪怕是你的妻子,你的孩子,只要你的妻子,你的孩子沒有簽一封文件,同意你去看她的話,你是進不了這個大門的.

我們今天聽起來覺得很奇怪,很荒唐,但是它背後的那個精神,那個理念是這樣的:每個人都有他的私隱權,當這個私隱權願意打開的時候,它才能打開；如果它不願意打開的話,沒人能打開.

當我們想到這樣的問題的時候,我們發現,個人的權利真的影響了我們今天的生活方式.但是從今天的福音經文裡我們看到上帝的一種思考方式,我把它簡單叫作上帝的邊緣性思考方式.為什麼是邊緣?因為這是被人所忽視的.為什麼是邊緣?因為這是另類的.但是為什麼是邊緣?因為它具有一種典範性的轉移力量,我們看到的福音故事其實很簡單,耶穌講故事,簡簡單單,沒有什麼懸疑之處.有一個故事談及葡萄園裡工人得薪水的故事(註二).地點在那裡?葡萄園.什麼事情發生?就是請工人工作.工作的薪水是多少?一天一錢銀子.這個故事有趣的地方是那個園主作了4次的招聘工作.第一次在什麼時間?經文說是巳初就去招請人了,巳初是什麼時間了?就是我們現在曆法的上午9點.所以從他幾次的招聘行動看到園主在上午的9點招了一批工人來,在中午12點再招一批,下午3點又再招一批工人,然後下午5點再一批.園主給薪水是在一個小時後,就是6點.所以在這個故事裡,一共發生了4次招聘行動,結賬的時間是在6點.

結果發生了什麼事?我們看到的圖畫是:先來的和後來的,他們都得到了同樣的薪水.很簡單的故事,沒什麼懸疑的地方,可是有趣的是這個故事最後幾句話.第9節說:約在酉初雇的人來了(最後一批雇的人來了),各人得了一錢銀子；接著先雇的人來了,他們以為必要多得,應該有4錢銀子吧,誰知他們也是各得一錢的銀子--這是結果.有趣的是這個反應:他們得了,就怒氣沖沖,帶著非常大的抱怨對園主說:我們整天勞苦受熱,那後來只做了一個小時,你竟叫他們和我們一樣嗎?!這是工人們對個人權利的伸張,對公義的呼喊,對不對?但是我們看到家主回答說:朋友,我沒有虧負你,你跟我講定的不是一錢的銀子嗎?我們有合約在手呀!拿著你的工錢走吧!我給後來的人和你的一樣,這是我的意願,我的東西難道不可以隨我的意思用嗎?因為我做好人,你就紅了眼嗎?所以耶穌在這個故事後就說了一句話:在後的將要在前,那在前的將要在後.我相信在這個故事裡的應用會惹來我們這些神學家的大麻煩,因為如果我們把這樣的故事放在今天這個社會的倫理裡,公平工資法的環境中,我相信我們這些神學家,牧師,肯定若不是被人家丟得滿頭包,就是被人說成是社會裡的極右保皇份子.

但是我們看到經文中是有一個脈絡在當中的,這個脈絡是在前邊.它前邊講到有一位富有的少年人去找耶穌問:我該做什麼善事才可得永生(註三).所以整個故事的脈絡是放在怎樣得永生這件事情上,然後我們當然知道這故事中耶穌說了:我實在告訴你們(因為這個少年人是很有錢的),財主進天國是難的,我告訴你,駱駝穿過針眼比財主進神的國還容易.門徒聽了就非常絕望,說:那誰還能得救呢?耶穌就說了一句很重要的話:在人是不能的,在神凡事都能.這是個脈絡,之後我們看到在馬太福音第20章17節裡,耶穌預言自己的死.所以整個經文的脈絡在說什麼呢?我相信如果我們把它引申到今天社會的公義環境裡或是倫理環境裡,耶穌在葡萄園的比喻就不好解釋了.當然,我也看過一些的著作,有些學者,特別是受到解放神學影響的學者,他們恰恰是在這裡伸張一個非常重要的想法,說:主是偏好窮人的.所以今天在天主教的倫理教導裡有一句很嚮亮的口號:神向窮人靠近--就是社會上的弱者,社會上被剝削的人,社會上被忽略的人,教會是要跟他們站在一起的.這是一個衍生,但我們在經文處境中更清楚看到它的前文和後文的脈絡是在說神的救恩方式.

神怎樣救人呢?神思考祂救恩的時候,跟我們有什麼不同呢?今日我們讀的約拿書是一個有趣的例子,因為在所有舊約的經文裡,很少有一卷書像約拿書那樣清楚地指出神是顧及整個世界的人.我們讀舊約的時候,很多時候聽到以色列民族在他們自己的苦難和流浪歷史中,對亞伯拉罕的神,以撒的神,雅各的神的那種呼喊,以至他們把自己這種貧窮,被剝削的生命和這位亞伯拉罕,以撒,雅各的神緊緊地連結在一起.但是這種渴望,這種信心,有一個沒辦法克服的盲點,就是有一個問題--那其他人怎麼辦?埃及人怎麼辦?亞述人怎麼辦?波斯人怎麼辦?其他整個世界的人怎麼辦?所以我們看到約拿書是藏著回答這問題的一卷很重要的書卷.但是它沒有答案,它只是問一個問題,問所有聆聽這故事的人--你說神該怎樣思考?你說我們這一群跟隨主的人該怎樣辦?該怎樣思考?

在今天我們所看到的一個故事中,我們發現這個園主讓我們看見一個很重要的圖畫,就是當園主願意把這個禮物,這個恩惠給人的時候,在他有絕對的自由,在他絕對的主權,在他有絕對的裁判權.每一個幫祂工作的人,我們能得到自己的那一份,但當我們看別人能不能得到同樣的一份的時候,我們常常在人生當中產生許許多多這樣的妒忌,掙扎和吵鬧.但是今天我們看到葡萄園的比喻裡很重要的問題,很重要的答案是:拿你的走吧!我願意給後來的和你的一樣,這是我願意的,我的東西難道不可以隨我的意思而用嗎?因為我做好人,你就紅了眼嗎?這句經文,其實不太容易回答,也不太容易解釋.我們聽故事容易,我們聽比喻也容易,但把它變成一個信仰思考的時候就不太容易,教會對這樣的事情也不太容易能完全明白透徹.但從今天這樣的一個故事,我們看到上帝的思考方式.

上帝告訴我們每一個人:我們每個人自己的權利,自己的自由,不是絕對的.個體的自由,個體的權利,有一個依從性在當中.它和我們的自然,我們的社群和我們的上帝,有一種彼此交依的關係.我們在生活中的遺忘,最可惜的是遺忘了上帝的思考方式.所以,當我們忘了上帝的思考方式的時候,我們常常帶來生命中許許多多的不安,愁苦和焦慮.

如果我們看聖經,我們看主的話,我們會看到最另類,最邊緣的思考方式--其實就是神愛世人,甚至是把祂的獨生子賜給了世界,叫一切信祂的不至滅亡,反得永生(註四).這是人類完全沒辦法思考的.在今天我們讀的保羅的經文裡,我們看到了上帝另類的思考方式.「祂本有神的形像,不以自己與神同等為強奪的,反倒虛己,取了奴僕的形像,成為人的樣式.既有人的樣子,就自己卑微,存心順服,以至於死,且死在十字架上」(註五).這樣的上帝,夠另類了吧?這樣的上帝,夠邊緣了吧?這樣的上帝的思考方式在社會的主流裡邊,哪裡去找?!

基督徒信仰的一個核心教導是去「完全」,乃至於我們像上帝的「完全」一樣,馬太福音第5章裡就是這樣說的.在這麼崇高的一個目標裡,保羅進一步解釋這樣的一個「完全」,落實在我們生命當中,去效法耶穌.但是效法基督的話,我們要追問一個問題:我們要效法什麼?其實保羅許許多多的時候藉著自己的經驗,藉著自己的生命,告訴我們怎樣效法耶穌的生命.今天我們所讀的經文裡,除了有一個這樣的葡萄園的故事之外,還有保羅自己對基督信仰的表達和體會.從保羅的經驗裡,從保羅寫的經文裡,我想與弟兄姊妹分享的有3點:

如果我們說去學習上帝的思考方式,我們從保羅他自己的生命和他自己的見證裡去學習幾點.第1點:他說我們要去學習一種無我的生命,沒有自己的生命,這聽起來很「東方」,但是我們從保羅自己的生命當中看到一個真實的圖畫.為什麼?因為保羅他自己說:我活著就是基督,我在我的肉身活著,為你們更是要緊的(腓立比書1章24節).保羅在他自己的生命裡是尋求一個為基督和為他人而活的生命.當我們是為基督而活,為他人而活的時候,一種無我的生命就在我們裡面產生了.當我們擁有這種無我的生命時,我們就會看輕自己內在要擁有的權利,要擁有的利益.

我們從今天保羅在腓立比書所講的經文裡,第2點,看到一個怎樣的圖畫呢?這是保羅他自己用一生的生命來見證「沒有我保羅的生命,只有基督的生命」這樣一個話語.「沒有我保羅的生命,只有為了你們這些弟兄姊妹而活的生命」的時候,他願意我們所有聽見的人,看到第2個道理--生命是一個彼此依存的生命.所以,在腓立比書第1章27節,當保羅講完「我這樣的一個生命,願意向你們敞開,為耶穌來敞開,因為我知道我活著就是基督,我死了就有益處」的時候,他講一句話,第27節:只要你們行事為人與基督的福音相稱,所以當我來見你們的時候,知道你們同有一個心志,站立得穩,為所信的福音齊心努力.這是彼此連結,彼此扶持,彼此依存的生命,同得一個心志,站立得穩,為所信的福音齊心努力.保羅在別的地方也說過像這樣一個身體的生命,一個互相為肢體的生命.我想這樣一個異象告訴我們:我們每一個人不是一個孤島,每一個人都要學習向別人敞開,每一個人要邀請別人進入自己的生命中,也願意把自己的生命傾倒在別人的生命當中.這是上帝想我們學習的思考方式,這是保羅的生命見證,這是在教會當中歷時歷代告訴我們什麼叫無我的生命,什麼叫一個彼此依存的生命.

當然,在這裡,我們看到保羅給了我們一個很深的提醒--上帝的思考方式是一個怎樣的思考方式?上帝這個邊緣性的思考方式是怎樣的思考方式?祂是要我們「活在當下」.「活在當下」這幾個字在最近已經都聽到爛掉了,很多人也說「活在當下」,但什麼叫「活在當下」,怎樣去「活在當下」,在什麼樣的時間狀態下「活在當下」?我們帶領很多的進修活動,帶領很多靈修,時常教導入門的弟兄姊妹去體驗我們當下的時候說:好,把你的注意力集中在你的呼吸,一呼,一吸,一呼,一吸......不只是基督徒教這樣的道理,一行法師也教這樣的道理,各大衛大的宗教真理也教這呼吸的法門.但是,基督教信仰的「活在當下」是活在什麼樣的當下呢?

在聖經研究裡,常常有一句話,叫做 「Already but not yet」--已然未然.已經發生的和還沒發生的事.對基督徙來說,什麼是「活在當下」?什麼是神要我們學習的「活在當下」的思考方式?我想聖經裡還是那三句話:活在信心當中,活在愛心當中,活在盼望當中.活在信心當中的時候,我們知道,從創世到今天,神從古時曉諭列祖列宗,神從古時曉諭歷時歷代的聖賢,把這個信心之道傳給了我們；活在盼望當中讓我們知道,雖然我們整個生命沒有在這個世間完全被展現,整個生命沒有在這個世間完全得完美,但是主應許,祂必然有再一次的來臨.但是在這一種信心和盼望的張力當中,「活在當下」是什麼當下?活在愛的當下,活在分享的當下,活在把我們的生命傾倒出來,好像保羅這樣一個無我的生命活出來,容不容易做?我相信非常難做.這是為什麼我們被造作門徒的原因,如果這麼容易做,幹嘛要造你來學,幹嘛要造我們來學?所以當主把這樣一個生命之道,把祂的思考之道來教導我們之後,我們就在一生的功夫裡來學習,去實踐這樣的生命.

第3點:上帝要我們怎樣去學習祂的思考?上帝是那一位創造宇宙萬物的主,上帝是那一位,祂的主權,祂的自由,祂的揀選,祂的判斷是完全自由的主.在今天,我們這一群跟隨耶穌的人,也許只是一位慕道者,也許只是一位過客,也許是非常認真地思考這個得救法門的門徒.葡萄園的故事帶給我們什麼啟發?葡萄園的方式帶給我們一個對上帝祂自己的心胸廣大,那種長闊高深的廣大一個更深的默想.祂的長,是從創世到整個世界的末了的長；祂的闊,是無限的,在整個新天新地裡,祂不會以人的先後次序把人招聚在祂的寶座前；祂的高,讓我們知道神行事奧妙難測,祂使萬事互相效力；祂的深,讓我們默想,耶穌甚至把祂的生命埋入死亡當中,在死亡當中祂做一些什麼事,我們不知道,但是我們知道祂復活的日子必是召聚列國列邦進入祂自己的國中.

這個道理很難懂,但是當我們願意去默想上帝思考的方式的時候,我們起碼能夠學會一點點功課,就是說我們的心胸變得非比尋常,我們能夠透過上帝的思考方式來改變我們自己的思考方式,我們能夠用上帝的廣大修改自己的偏窄,我們能夠用上帝的恩慈去勝過我們自己對自己權利的伸張.葡萄園主的話非常難明白,但是在難明白裡又不是完全沒辦法了解,起碼我們看到一點:在神的國度當中,我們進來的時間無論是在9點,12點,3點還是5點,神以祂絕對的主權,以祂絕對的自由,以祂無比的恩慈來款待眾人,來歡迎眾人.

更重要的是,我們這一群在葡萄園中的工人,不要紅眼,不要因為別人比你來得晚,承受了跟你一樣的報酬,你心裡就埋怨.因為當我們心裡埋怨指責的時候,我們還沒摸著這位葡萄園主的心意.但願我們每個人能夠學習祂的心意,跟隨祂的心意.

我們禱告:主,我們把自己交託在你面前.讓我們這一群跟隨主的人,心胸寬闊,讓我們能夠看見萬事互相效力,而我們在你園子裡所享用的一切救恩和享用的一切成聖的果子,都是出於你的禮物.讓我們帶著一個感恩的心來使用,來生活；也讓我們自己內在那種喋喋不休的妒忌,心胸的狹窄,在你自己恩典慈光裡能夠煙消雲散,乃至於這樣的生命的群體,在這個世界裡活出美好的見證.奉耶穌聖名求,阿們.

(註一)然而有許多在前的將要在後,在後的將要在前.(太19:30)

(註二)因為天國好像家主,清早去雇人,進他的葡萄園作工.和工人講定,一天一錢銀子,就打發他們進葡萄園去.約在巳初出去,看見市上還有閒站的人.就對他們說,你們也進葡萄園去,所當給的,我必給你們．他們也進去了.約在午正和申初又出去,也是這樣行.約在酉初出去,看見還有人站在那裡．就問他們說,你們為甚麼整天在這裡閒站呢.他們說,因為沒有人雇我們．他說,你們也進葡萄園去.到了晚上,園主對管事的說,叫工人都來,給他們工錢,從後來的起,到先來的為止.約在酉初雇的人來了,各人得了一錢銀子.及至那先雇的來了,他們以為必要多得．誰知也是各得一錢.他們得了,就埋怨家主說,我們整天勞苦受熱,那後來的只做了一小時,你竟叫他們和我們一樣麼.家主回答其中的一人說,朋友,我不虧負你．你與我講定的,不是一錢銀子麼.拿你的走吧．我給那後來的和給你一樣,這是我願意的.我的東西難道不可隨我的意思用麼．因為我作好人,你就紅了眼麼.這樣,那在後的將要在前,在前的將要在後了.〔有古卷在此有因為被召的人多選上的人少〕(太20:1-16)

(註三)有一個人來見耶穌說,夫子,〔有古卷作良善的夫子〕我該作甚麼善事,才能得永生.(太19:16)

(註四)神愛世人,甚至將他的獨生子賜給他們,叫一切信他的,不致滅亡,反得永生.(約3:16)

(註五)他本有神的形像,不以自己與神同等為強奪的,反倒虛己,取了奴僕的形像,成為人的樣式.既有人的樣子,就自己卑微,存心順服,以至於死,且死在十字架上.(腓2:6-8)



\newpage

\section{因我們只有一個老闆}
\label{sec:122}
\textbf{鄧瑞強博士 - 無得揀,因我們只有一個老闆}
\newline
\newline
\hyperref[sec:121]{< < < PREV SERMON < < <}
~
\hyperref[sec:toc]{[返目錄]}
~
\hyperref[sec:123]{> > > NEXT SERMON > > >}


馬太福音 22:15-22
\newline
\begin{longtable}{cl}
\hline
\hline
章節 & 經文 (和合本修訂版)\\
\hline
22:15 & \begin{tabularx}{0.7\textwidth}{X} 於是，法利賽人出去商議，怎樣找話柄來陷害耶穌， \end{tabularx} \\ \\ \relax
22:16 & \begin{tabularx}{0.7\textwidth}{X} 就打發他們的門徒同希律黨人去見耶穌，說：「老師，我們知道你是誠實的，並且誠誠實實傳神的道，無論誰你都一視同仁，因為你不看人的面子。 \end{tabularx} \\ \\ \relax
22:17 & \begin{tabularx}{0.7\textwidth}{X} 請告訴我們，你的意見如何？納稅給凱撒合不合法？」 \end{tabularx} \\ \\ \relax
22:18 & \begin{tabularx}{0.7\textwidth}{X} 耶穌看出他們的惡意，就說：「假冒為善的人哪，為甚麼試探我？ \end{tabularx} \\ \\ \relax
22:19 & \begin{tabularx}{0.7\textwidth}{X} 拿一個納稅的錢給我看！」他們就拿一個銀幣來給他。 \end{tabularx} \\ \\ \relax
22:20 & \begin{tabularx}{0.7\textwidth}{X} 耶穌問他們：「這像和這名號是誰的？」 \end{tabularx} \\ \\ \relax
22:21 & \begin{tabularx}{0.7\textwidth}{X} 他們說：「是凱撒的。」於是耶穌說：「這樣，凱撒的歸凱撒；神的歸神。」 \end{tabularx} \\ \\ \relax
22:22 & \begin{tabularx}{0.7\textwidth}{X} 他們聽了十分驚訝，就離開他走了。 \end{tabularx} \\ \\
[1ex]
\hline
\hline
\end{longtable}
日期:2011年10月30日

主題:無得揀,因我們只有一個老闆

證道:鄧瑞強博士

經文:馬太福音22:15-22

各位弟兄姊妹,主內平安.

在以前蘇聯的時代,一個蘇聯人和一個美國人,討論誰的國家更自由.

美國人說:我的國家很自由.蘇聯人說:何以見得?美國人說:我們美國人可以在白宮前,大罵美國的總統.蘇聯人說:這有什麼特別,這種事我們蘇聯人也做得到.美國人說:真有此事.蘇聯人說:的確如是,我們蘇聯人同樣可以在克里姆林宮前,大罵美國的總統.身在極權下,自由會被扭曲的.

政治議題很多時都是敏感的.若身在極權社會內,要對政治議題表態,則很困難.困難在難於暢所欲言.暢所欲言者,很難避免身陷麻煩或險境.

主耶穌身在極權力量下,被迫要對政治議題表態.他如何回應,以顯示他內心的自由呢?

今日講道的經文,是今日的福音書經課:馬太福音 22:15-22.

在這裡,主耶穌面對一個含有政治意味的兩難問題.

我們再讀一次這段經文.

22:15　當時,法利賽人出去商議,怎樣就著耶穌的話陷害他,22:16　就打發他們的門徒同希律黨的人去見耶穌,說:「夫子,我們知道你是誠實人,並且誠誠實實傳神的道,什麼人你都不徇情面,因為你不看人的外貌.22:17　請告訴我們,你的意見如何?納稅給凱撒可以不可以?」22:18　耶穌看出他們的惡意,就說:「假冒為善的人哪,為甚麼試探我?22:19　拿一個上稅的錢給我看!」他們就拿一個銀錢來給他.22:20　耶穌說:「這像和這號是誰的?」22:21　他們說:「是凱撒的.」耶穌說:「這樣,凱撒的物當歸給凱撒；神的物當歸給神.」22:22　他們聽見就希奇,離開他走了.

這裡講到,法利賽人設局陷害耶穌.

法利賽人基本上是宗教人士,只關心宗教.他們以敬虔者自居,也教人以敬虔.他們處處審視人是否敬虔,好保護社會的敬虔,也凸顯自己的敬虔.

在一般人看來,他們不是什麼壞人,只是有時他們為了顯出自己的敬虔,對別人過份挑剔而已.

他們為何要害耶穌?

因為耶穌針對了他們以敬虔自居的死穴.耶穌偏偏指出了他們的敬虔正正是不敬虔,而他們認為極不敬虔的事,在耶穌看來,卻是無傷大雅的.

法利賽人認為人在安息日不應醫病,但耶穌偏偏在安息日醫病,並指出在安息日不救助病人,正正是不敬虔.法利賽人要在吃飯前洗手,但耶穌卻不洗手,並指出洗不洗手並不重要,重要的是人內心的意念.

耶穌刺痛了法利賽人的自尊心.他們賴以為生的敬虔,被主耶穌揭露為一文不值.當你最珍惜的東西,被人當場揭露為毫無價值,你一定極不開心.

試想想,當你一生賴以為榮的就是自己的慈悲心,有一天卻被人當面指為只是沽名釣譽的虛偽,你一定極不開心.

這故事提到希律黨人.希律黨人是依附希律的人,希律是羅馬的傀儡王,在當時統治加利利地區.希律黨人是支持羅馬政權的人.法利賽人一般都不喜歡希律黨人.但為了剷除耶穌,他們和希律黨人合作.宗教的狂熱加上政治的陰險,就像失去自制力的人拿著機關槍,常常是最危險的.

在希律黨人的監視下,在這些暴政的執政黨人的監視下,他們問了耶穌一個政治問題:「請告訴我們,你的意見如何?納稅給凱撒可以不可以?」

「可以不可以」在英文KJV或RSV等譯本裡譯為 ``Is it lawful?"這樣做,是否「合法」?問題是:合什麼法呢?是指合羅馬的法呢,還是指合他們祖宗的宗教律法呢?法利賽人守衛著他們祖宗的宗教律法,而希律黨人守衛著羅馬的律法.一個政治問題,背後牽涉神的要求,難答之極.

「納稅給凱撒,是否合乎律法?」訴諸民意,民意一定是說這不合乎神的律法.但這樣回答,便是犯了羅馬的律法.法利賽人要希律黨人在場,就是希望他們見證耶穌這樣回答.耶穌沒有隨民意回答說「不」,但耶穌也沒有回答說「是」.耶穌沒有回答說:納稅給凱撒,是合乎律法的.不這樣答,不是因為耶穌怕背逆民意,乃是因為這的確違背了神的律法.

耶穌的回答乃是:「拿一個上稅的錢給我看!」這個稅錢將會說明誰是誰非.羅馬人不單要求他的屬地納稅,還要求他的屬地用羅馬人的錢納稅.這正是問題所在.猶太人未必不願意納稅,但他們不願意用這羅馬錢幣納稅.

因為,這錢幣的一面刻著當時羅馬皇帝Tiberius的像,另一面刻著「凱撒提庇留,神聖的奧古士督的顯赫的兒子,大祭司」(Tiberius Caesar, august son of the divine Augustus, high priest).〔參考:Daniel J. Harrington, The Gospel of Matthew (Collegeville: Liturgical Press, 2007), p.310〕用這錢幣,就多少承認凱撒是神.猶太人不想用這錢幣,就正如我們不想用「溪錢」來奉獻一樣.

主耶穌沒有這錢幣,這幫人身上有這種錢幣.我們不知道這錢幣來自希律黨人,抑來自法利賽人.若來自法利賽人,則是相當耐人尋味的.

這錢幣一亮出來,足以叫那些自命信仰純潔的法利賽人慚愧.他們口中常常掛著摩西的律法,而摩西的十誡不是明明禁止為人造像的嗎?他們若用這錢幣納稅,不是已經說明自己的所謂聖潔是多麼虛偽!一個自問信仰純潔的人拿著刻有偶像的錢,就像你對著愛人拿著前度愛人的相片一樣,罪行不證自明.這些法利賽人,還敢用納稅的問題去陷害耶穌.

對於希律黨人而言,這些錢幣猶如他們的黨徽.凱撒,正是他們效忠的對象.他們在凱撒那裡獲得權力,身分和利益.凱撒的法則,是順我者生,逆我者亡的法則.希律黨人守護著這種權力.

主耶穌問:「這像和這號是誰的?」他們說:「是凱撒的.」耶穌說:「這樣,凱撒的物當歸給凱撒；神的物當歸給神.」

在這裡,耶穌是說政治的事按政治的法則處理,神的事按神的法則處理嗎?耶穌是說政治應有一獨立性,能獨立於神嗎?

對這節經文,過去一般人的理解是政治有政治的領域,按政治的法則處理.凱撒的領域,當歸給凱撒.推而廣之,經濟有經濟的領域,只應按經濟的法則處理.心理有心理的領域,只應按心理的法則處理.當然,信仰也有其領域,但這只是一非常私人的領域.神是一個私人領域的神.這位神,不應過問政治的事,那是屬於凱撒的.這位神,不應過問經濟的事,經濟的考慮只在乎獲取最高利潤.這位神,最好不要過問世事,因為世事並不屬於他.如此,這位神還有什麼領域?他只有人類心靈的一個幽暗角落,用來提供心靈慰藉.

若是這種理解,一個基督徒會將生活分裂成兩個領域,一個是屬於凱撒的,一個是屬於神的.這個基督徒回到一個人人落井下石以爭上位的公司時,他便會理直氣壯地採取落井下石的方式爭上位.他會說,這裡是屬於凱撒的,要用凱撒的法則.這個基督徒上到人人爭座位的公車時,他也會毫不猶疑地爭座位,爭到時在心裡會毫不猶疑地舉起勝利的手勢.他會說,這裡是屬於凱撒的,要用凱撒的法則.這個基督徒在一個「專走精面」而不會勤奮工作的社會裡,他也同樣會「走精面」而不會腳踏實地.他會說,這裡是屬於凱撒的,要用凱撒的法則.

耶穌的說話,真是這樣理解嗎?

若耶穌的說話,真是這樣理解,他就不會在凱撒權力的執行者彼拉多手下受難了.

若耶穌的說話,真是這樣理解,他就只會埋首在聖殿裡祈禱,而不會隨處指斥人的墮落了.

在耶穌眼中,難道有什麼領域不屬於神嗎?

這是天父世界,所有的領域都是屬於神的.

我們不能在天父的世界裡揀一些東西屬於神,又揀一些東西屬於凱撒.無得揀,因我們只有一個老闆在天上.

「凱撒的物當歸給凱撒；神的物當歸給神」,主耶穌這句話的意思,在我看來,應是說:讓凱撒的暴政歸還給凱撒,讓凱撒的個人崇拜歸還給凱撒,讓凱撒那種倚仗軍事力量去對付異己的態度歸還給凱撒,讓凱撒那種徹底屬世的力量歸還給凱撒.面對世界,面對政治,面對經濟,面對異己,面對人生,這一切都是屬於神的,不能用凱撒的方式,只能用神的方式.這一切屬神之物全歸給神,不能容讓半點凱撒滲透其中.〔參考:Stanley Hauerwas, Matthew (Grand Rapids: Brazos Press, 2006), pp.189-191〕

面對一個政治問題,主耶穌將之轉化為一信仰問題,他將問題的重心放在人是否以神的法則去處理人生的種種問題之上.

如果上述的理解是正確的話,則主耶穌是要我們用神的法則去克服凱撒的法則.由於我們生活的方方面面都滲透著凱撒的法則,則我們便要在生活的方方面面裡進行信仰的對抗.「將屬神的一切事物,全歸給神」.

我見過有車在公路上爆了輪胎,司機將車泊在路旁,他招手向路過的車求助,有人停了車,幫助他.一般人說,時間就是金錢.用盡一切時間去賺金錢,這是凱撒的法則.用盡一個家庭傭工的每一刻時間,同樣也是凱撒的法則.但有人停下來,花時間去幫人,這是神的法則.在不同地方有很多人花時間做義工,這是神的法則.

有一次,我在超級巿場買東西.前面的那位顧客和收銀員吵起來.原來收銀員入錯了一條數,後來發覺了,已多番賠罪,但那顧客仍罵過不停.當他罵得那收銀員尊嚴盡失後,而自己也罵得「無氣」了,便帶著勝利的微笑離開.有人說,顧客就是上帝.這位顧客真的以為自己是上帝.他和凱撒一樣,以為自己是上帝.他的行為,盡是凱撒的法則.真正的上帝,會來到人間,體諒人的錯失.這是神的法則.

有一次,神學家莫特曼(Moltmann)問克羅地亞裔的神學家沃弗(Miroslav Volf)一個問題:「你能寬恕那些屠殺你們族人的塞爾維亞人嗎?」沃弗的回答大意是:「若主耶穌要求我們寬恕,我們就只能寬恕.」這回答真好.唯有寬恕,才能瓦解凱撒的法則.唯有神的真理,才能拒絕世俗的潮流.唯有背起十字架跟從主,我們才不會拿著刀劍,帶著算盤跟從凱撒.

「凱撒的物當歸給凱撒；神的物當歸給神.」願凱撒的法則隨凱撒而去,而我們以神的法則,活在這個天父世界裡.

但願榮耀歸於聖父,聖子,聖靈.阿們.



\newpage

\section{這福音是上帝的大能}
\label{sec:123}
\textbf{張伯笠牧師 - 這福音是上帝的大能}
\newline
\newline
\hyperref[sec:122]{< < < PREV SERMON < < <}
~
\hyperref[sec:toc]{[返目錄]}
~
\hyperref[sec:124]{> > > NEXT SERMON > > >}


羅馬書 1:16-17
\newline
\begin{longtable}{cl}
\hline
\hline
章節 & 經文 (和合本修訂版)\\
\hline
1:16 & \begin{tabularx}{0.7\textwidth}{X} 我不以福音為恥；這福音本是神的大能，要救一切相信的，先是猶太人，後是希臘人。 \end{tabularx} \\ \\ \relax
1:17 & \begin{tabularx}{0.7\textwidth}{X} 因為神的義正在這福音上顯明出來；這義是本於信，以至於信。如經上所記：「義人必因信得生。」 \end{tabularx} \\ \\
[1ex]
\hline
\hline
\end{longtable}
日期:2011年11月6日

主題:這福音是上帝的大能

經文:羅馬書1:16-17

證道:張伯笠牧師

弟兄姊妹,平安!

的確很激動,在中國的領土內,第一次參加崇拜.對我來說,這是一個很特別的時刻,能夠與香港的弟兄姊妹,朋友們一齊來敬拜我們的神.

今天我用的經文在羅馬書1:16-17:我不以福音為恥．這福音本是　神的大能,要救一切相信的,先是猶太人,後是希利尼人.因為　神的義,正在這福音上顯明出來．這義是本於信,以致於信．如經上所記,『義人必因信得生.』

保羅說:「這福音本是神的大能,要救一切相信的」.當保羅被上帝在大馬士革的路上耶穌向他顯現,呼召出來以後,保羅將一生把他所有智慧,才華用於傳福音的路上,所以保羅過去是一個希利尼人,保羅是一個非常有才華的希伯來人,那麼保羅熟讀舊約的律法,但是他接受福音以後,他做了一個宣告:「這福音本是神的大能,要救一切相信的」.當我們從一個舊的文化中,進入到福音中來的時候,我們發現我們有很多的事情,跟我們過去的格格不入.保羅說:在他的同胞中也好,在福音的路上也好,他是不以福音為恥的.

他說:「這福音本是神的大能,要救一切相信的」.因為在當時情況下,在猶太人,希臘人的心目中,他們覺得「福音」是一個非常不切實際的東西,他們說福音是一個不可能的東西,他們覺得福音是沒有什麼文化的,他們也不能夠尊重福音.所以保羅說他不以福音為恥,這個很了不起的宣告.其實我們中國基督徒也是如此,我們中國有五千年的文化,我們覺得我們中華文化博大精深.所以當我們一進入福音的時候,講到十架福音的時候,我們就對我們的信心產生了懷疑.福音就是這樣簡單嗎?信就能得救這樣簡單嗎?所以保羅告訴我們說「我不以福音為恥．這福音本是神的大能,要救一切相信的」.

這個在保羅的神學中,他常常提到這個十字架的道理是一個絆腳石.在希臘文是(skandalon)的意思,這個(skandalon)的意思就是說,是個絆腳石的意思.很多人當他認識耶穌時候,當他看到十字架就知道,這是一個不能認識到福音的救恩中很重要的東西,保羅說:因為十字架的道理,對那些愚拙／滅亡的人是一個愚拙的道理,在我們得救的人則是神的大能.為什麼呢?因為「十字架」在當時的文化中是一個咒詛.十字架把人分為兩類:一類是得救的,一類是滅亡的；一類是愚拙的,一類是聰明的.可是在當時的文化中就不是如此,當時侯的文化,是分羅馬人,野蠻人；猶太人,外邦人；奴隸和自由人,等級深嚴.可是當耶穌釘在十字架上,祂的寶血遮蓋世人的罪的時候,世人就因著祂可以成為上帝的兒女,他們在文化中發現了新的突破,這就是基督的文化,聖經的文化,聖經的世界觀.所以保羅告訴我們,在傳統的文化中,因為恥辱的東西成為神等大能.這個對於那些比較聰明的人,是愚拙的,但是保羅把這些聰明的定為滅亡的人,對我們這些得救的人卻是神的大能,我們感謝如此.

因為我們中國人也是如此,我們這一代的中國人常常覺得自己是文化人,我們有博大精深的中國文化,其實想來想去,我們這一代人蠻可憐的,我們對中國文化並不了解太多.我常常在美國跟一些中國的博士生,碩士生,這些大學的老師談中國文化,他們其實對中國文化只有了解一點點而已.我們只懂得一些皮毛(但是香港中文大學可能比海外的條件好一點).我們這一代人對中國文化…

孔子我們不知道多少,我們沒有讀過孔子的書,我們沒有讀過孔子的《四書五經》,基本沒有讀過,我們小時候可能讀過《百家姓》,但可能沒有讀完,基本上大概屬於這種狀態.我們對老子也了解不了多少,我們覺得,我們對佛教了解也不多.中國的主文化主要來自儒釋道.儒家我們叫「人本主義」；釋迦摩尼我們叫「虛無主義」；老子,老莊哲學叫「消遙主義」,其實這些我們都不太了解,我們只是了解一個皮毛而已.

我們想起孔子,就知道這是以人本為主的東西,孔子說:「敬鬼神而遠之」──不要碰這個東西,我們強調人的智慧,有一天,人可以達到一個聖賢的程度.但老子更有意思啦!老子講這一切都是虛無的,沒有意義的.所以我們中國人常常把道家跟道教放在一起,其實這是不對的,道教跟道家很大的區別,這跟狗跟熱狗的區別.當我們對釋迦摩尼和佛教也了解不了多少.我們常常看到中國人到佛堂裡去尋求心裡的平安,其實真正對佛教文化我們也了解不多,其實我們現代人了解的是馬克思主義的哲學,多少懂一點.懂馬克思主義也不是全部,我們說我們是馬克思主義的哲學者,又或是馬克思主義的信徒.但是我們卻沒有讀過馬克思多少書,《資本論》是沒有人讀過的,其實我們所知道的是列寧哲學其中的一部份──暴力革命──一個接一個,另一個推翻另一個的暴力的行動,我們小時候的文化影響,應該是列寧的不是中國傳統的,所以當我們聽了福音的時候,我們就發現我們與十字架的道理,有一個很大的距離.

我們常常用我們自己的觀點來衡量聖經,當我們看聖經,哎!這段不對,這一塊寫得沒有道理,耶穌祂說:祂就是「道路,真理,生命」,若不藉著祂,就沒有人得救.

這個為什麼這麼狂妄?對不對?!我們會對聖經提出很多質疑,因此在美國大學的查經班裡,會看到特別的景象──我們大陸弟兄姊妹問的問題,台灣跟香港的牧師們常常不知道該怎樣回答.我們也曾經問過牧師:在人類犯罪之前,上帝沒有咀咒蛇的時侯,蛇是用什麼走路的?我們問過很多問題,牧師都說不知道,有一天你們問上帝去吧?這個問題我們也回答不了.所以我們常常用我們的文化來對比十字架的道理.

保羅說:我們這些道理都不重要,唯有十字架的道理.因為這福音本是神的大能,他認為福音是一種能力.這個能力對我們這些滅亡的人,是屬於愚拙的.因為我們想到這個福音,是想到它是個絆腳石,因為十字架在當時的文化中,這個十字架是一個詛咒的記號,犯了最大惡極的人,他才可以釘在十字架上,而且還不是羅馬人,因為羅馬人就算犯了最大惡極的罪──殺人掠貨──也不會釘在十字架上,因為他有一個公民的身份,這個是給羅馬之外的外國人.所以一個釘在十字架上的人,那麼軟弱的人,祂怎麼能夠成為人類的得救主呢?祂怎麼是福音呢?十字架怎麼是個福音呢?如果從人類的眼光看,十架應該是個禍延.當我們看到十字架,我們想到那是一個死亡,不僅僅是死亡,而且是一種詛咒.

我們看到耶穌受難記那部電影,你會看到耶穌釘在十字架上的時候,他們鞭打祂,羞辱祂,用槍去捅祂,不僅是這個羞辱.十字架釘的人不是馬上死的,他們用三根釘子,把犯人的身體釘在十字架上,所以他的血是一點點流,一點點流,他非常痛苦的展現在別人,親人面前,他的父親,母親在下面看著他在受難,他的妻子,兒女可能在那看著他在受難,而且把他拔得清光,身上沒有一絲的布可以摭擋他身體任何一個部份.我們今天知道說,當今耶穌的藝術品,不管畫家也好還是電影家也好,給祂圍了一塊布而已,其實沒有的,沒有的.十字架在猶太人的文化中是一種詛咒的記號,犯了罪的人才可以掛在木頭上.那麼耶穌釘在十字架上,卻成了人類的福音,這是一個很奇怪的想法,所以我們用中國的文化說,這種叫作一種「反合性」,有的時候說叫「吊詭性」,很特別的.

所以當我們想到這裡的時候,我們知道當時的社會對於這個福音的認識是錯誤的,所以保羅提出說,這個福音,這個十字架的道理,是神的大能要救一切相信的.他也引用了以賽亞書二十九章十四節說:「我要滅絕智慧人的智慧,我要棄絕聰明人的聰明.」換句話說:十字架的道理就是神要按照神的方法,不是按照人的方法,是用神的方法,去做神要做的事情,藉著十字架,一個詛咒的記號,使人類的罪滅絕,使人類的罪在那一刻都釘在耶穌的身上,耶穌從死裡復活,勝過死亡,帶來人類的福音.什麼叫作「福音」呢?「福音」就是好消息.

當你得了癌症,醫生說我沒辦法啦!化療吧!你說化療能活多久?他想了想,三年.如果不化療呢?他想了想,也差不多三年.這是我在院裡醫生對我說的.大概就是這個樣子,那你這三年做什麼呢?你不會再讀大學了,你就覺得沒勁,沒意思了.那你幹甚麼呢?你天天在死亡的咒詛和恐懼之中,你的人生就一片黑暗,你不會有什麼喜樂,你不會有什麼光明,除非你認識那永能的神,你知道你未來到那裡去,可能比現在還好,你才能夠有喜樂,所以我們叫作福音.如果有一個人對你說,好了我有一個藥方,你這個醫生沒有用的我有,我有這個藥方你吃了你就會好了,這個就叫作「福音」.人類在死亡的咒詛中沒有出路,摩西在詩篇九十篇中說:人類的生命很短暫的,我們像早晨的草,早上生長發芽,晚上割下枯乾.他說我們生命的長度就像一聲嘆氣.唉一聲就過去了,就那麼短!我今天說因為你年青,你還沒有覺得歎氣,我做牧師常常會為了一個病人,一個教會的會友臨走的最後一刻伴隨他.那個時候我覺得摩西說得很真實,他看了我一眼,他就走了,那就是一種歎息.摩西說不僅是這個短暫,而且還苦難.他說我們這個人生無非是勞苦愁煩,轉眼成空,我們便如飛而去了.我們人生有什麼盼望呢?人們的盼望在於福音,人們的盼望在於十字架上的耶穌基督,祂給我們帶來了一個全新的觀念,我們信靠祂的人,可以出死入生,進入永恆,讓有限的生命跟上帝的永恆連在一起,所以保羅在羅馬書第六章這樣說,我們受洗歸入耶穌基督的人,是受洗歸入祂的死.我們說我們既然在死的形狀上與他聯合了,那我們就在復活的形狀上與他聯合.所以我們就有了永恆的生命,這個對人類來說是福音,這個佛教沒有,道教沒有,儒教沒有,在中國的傳統裡沒有福音.

所以我記得中國有一位將軍叫作劉亞洲(大家我不知道你們知不知道這個人?我年青寫作的時候,他也在寫作,那個時候,我們開作家會也見過面,我也很喜歡讀他的東西,後來他當了官,我坐了牢,他當了官,人走的路還不同,當然他有才華,他是李先念國家主席的女婿),有一天他有一篇講演,叫《信念與道德》──他談到了中國的文化,他說中國的文化主要儒釋道,這三個的文化不能給我們中國的復興提供精神的資源,基本上他認為是沒有意義的,套用他的話,是無法振興中華.但是他對基督教的文化非常尊崇,他說你看我們中國的神跟西方的神很不同,西方的神是受苦受難的形像,西方是上帝在受苦,人在享福；東方的神在享樂,人民在受苦.比如說我們廟宇裡的神,每一個都大腹便便,無憂無慮,嘻皮笑臉,食盡人間煙火,個個吃得腦滿腸肥.他講得有一定的道理,他說:我在美國教堂門口坐了一整天,我發現去教堂的人都是愁眉苦臉的進去,但神情輕鬆喜樂地出來,如果你久而久之是這樣子,就會變成健康的人.共產黨員都有這樣的看見,而且是共產黨的一個將軍,國防大學的政委,有這樣的一個看見,這樣是很奇特的.所以人們在尋求真理的時候,我們一定要來到十字架的面前,所以我今天用三點與大家弟兄姊妹來分享:

第一,福音有赦罪的大能；

第二,福音有醫治的大能；

第三,福音有勝過死亡的大能.

你們就是福音等大能!所以保羅說:這福音本是神的大能,要救一切相信的.

我第一次接觸福音是一九八九年,對我來說,那年是我人生最黑暗的一年,也是我最幸運的一年.因為那一年我什麼都沒有了,但是我有了耶穌.我記得一九八九年我被政府通緝,我們開始逃亡,不知道往那裡逃,聰明的人往南方逃,因為會廣東話就比較方便,我就往北方逃,因為我不會廣東話(只會東北話),所以當我逃到東北方的時候,六月十三號中國政府在中央電視台公開通緝王丹等二十一名學生,我是其中的一個.那時候我開始覺得人生沒有盼望了,人生就是死,我知道我奮鬥的一切都沒有了,不僅沒有了,你還要面對一個你不知道的苦難在前面.感謝主!那時我逃到中國最北方的一個農村,我藏在一個農民的家裡,這個農民是信耶穌的,我很感謝主,神很愛我,祂把我放到那個環境,祂若不把我放到那個環境,今天沒有張伯笠牧師,而且今天我可能也不會活著站在這裡,我不知道我會怎樣?有怎麼樣的選擇?那個時候這個姐妹向我傳福音,因為藏在她家裡我不敢說我是誰,我就欺騙她.沒有信主之前,中國人都會騙人的,有的是善意的欺騙,我給自己起了個名字叫王老四,如果說我是張伯笠她不敢幫助我的.當然有一天聖靈還是感動我,叫我應該說實話.因為我有危險,如果在她家,我被抓了她也會坐牢的,這是有危險性的,所以有一天我問她:姐姐,你知道我是誰嗎?我想告訴她,人家其實是知道的,人家不好意思揭發你而已,她就說,你不是張伯笠嗎?你怎樣知道的?你不是上電視了嗎?她看到通緝了.我就是,我就是…

那不就是通緝令嗎?不是不好的事,她說管你什麼通緝今,能上電視就不太容易.

她說如果你不是上電視,你不被通緝怎麼逃我家來了?我相信上帝把你帶來的,我第一次從一個農民的口中說「上帝」,我覺得這個農民還很迷信.因為我們不信上帝,我們信自己,我們唱《國際歌》不是唱得很清楚嗎?我們不要什麼神仙皇帝,我們只倚靠我們自己,這是我們的文化.可是後來她就開始跟我們傳福音,你沒有出路,你也不會有喜樂,你唯有信靠耶穌,祂是你唯一的道路,真理和生命.所以一個北京大學中文系的一個學生,聽一個不識字的人來傳揚天國的福音,是很有意思的一件事情.如果我不被通緝,如果不是在她家有求於她,我可以天天聽一個不識字的人跟我講人生的哲理?我想挺難,挺難的!人,當你面對一個環境,你才知道謙卑下來,所以她向我傳福音的方法很特別,她有一個筆記本抄的聖經是《約翰福音》,唯一的一卷書,我再也沒有讀過其他書,我就信耶穌了.所以打開《約翰福音》看,它很奇妙的,《約翰福音》開篇就說:「太初有道,道與神同在,道就是神,這道太初與神同在.」我想農民讀懂這書嗎?講的是:道!道!道!道!對不對?

我唸北京大學都不懂,我在北京大學上了一門哲學的課,叫《道法如墨》,我的老師叫作陳鼓應從台灣大學來的,是個老莊的專家,講了一個學期的「道」,越講我們越糊塗.要考試的時候,我就找他談,問他什麼是「道」?他說我也不知道!你也不知道你講什麼「道」?(今天我講道我知道是什麼道,我講的道是什麼?耶穌基督.祂是道,道成肉身住在我們中間,充充滿滿的有恩典,有真理.)所以我就為難他,你不知「道」,那你講什麼課?北京大學的學生挺會難為人的,那個時候我們很狂妄,什麼都不知道,就知道跟老師唱對台戲,所以我就說:那你告訴我什麼是「道」?他說我不知道.但是他說,你是學生對不對?我說是,你有「為徒之道」,你要懂得「為徒之道」,不要跟老師亂較勁.我是老師我有「為師之道」,經商有「經商之道」,從政有「從政之道」.連強盜也有道──盜亦有盜,如果你不懂道,你別活了,這個很有意思.可是我還是不知道.但我一讀《約翰福音》:太初有道──LOGOS, WORD──沒有宇宙之前,這個道就存在了,祂是自有永有的,道與神同在,道就是神,道就是耶穌.祂來到我們中間充充滿滿的有恩典,有真理,祂為我們釘死在十字架上,祂的寶血為我們流出,這就是福音.當我們信靠祂的,我們與祂連合,我們不僅在十字架上死的形狀與祂連合,我們在復活的形狀上也與祂連合,我們有永恆的生命,死亡不再成為我們的咒詛,我們可以出死入生.

在中國文化裡,我們的革命文化講的是「出生入死」.我第一次聽到出死入生,覺得這個比出生入死好多了,我們可以「出死入生」,這就是「福音」.AMEN!哈利路亞! 我是在中蘇邊境,冰天雪地上決定信主的.其實雖然是我一個決定,但是上帝已經揀選了我.逃走那一天是一九八九年聖誕節,我越境去蘇聯,也是我有生以來第一次出國,沒有護照,沒有VISA,什麼都沒有,而且還冒著生命的危險去偷渡.

大雪很大,零下40度,解放軍的巡邏隊在巡邏,蘇聯也可能有巡邏隊,不知道.反正我就從雪地爬到蘇聯,我在走之前,這個基督徒,她跟我講一句話:你再往前走姐姐不能夠幫你了,但耶穌可以幫助你,你遇到任何的難處,你要禱告,祂是聽禱告的神.對我來說,我還未能完全相信祂是神,所以我當然也不會禱告,但是我看到她眼淚一行,一行的…

她為我幹麵條,眼淚就往麵裡頭掉,我就特別感動,人世間是有真情的.我就說,姐姐我記住你的話,你放心別難過.雖然當時我是敷衍她,沒有認真.但當我逃到蘇聯一個無人區的地方,方圓幾百里跟本沒人的地方,看地圖以為離城市很近,走起來根本… 後來大雪就把我埋起來,我想回國也找不到回國的腳印,大雪把我腳印完全覆蓋,我躲在草堆裡快要凍死的時候,我想到姐姐的話:你要禱告,祂是你得救主.那一天我就禱告說:主啊!如果你存在,你讓我活過這一天,我就把自己獻給你當作活祭.那個時候不知道活祭是什麼意思,說聽他們基督徒農民禱告時愛說一句話:主啊!我要把自己獻給你當作活祭,是聖潔的是,是你所喜悅的.我當時以為這個活祭…

因為我知道祭物總是生畜嗎,對不對?我以為這是吃的,上帝喜歡吃活祭.所以祂不想吃死的,所以我說:主啊!我要把自己獻作活祭,讓我活過這一天,沒想到那個禱告,成為我一生中上帝給我的一個應許.我把自己獻給祂作活祭,二十四小時以後我被挖出,一天一夜過去,挖出來,我沒有凍死,這是個神績.是我人生中第一個神績,不可能沒有凍死的,但是我就沒有凍死.

挖草的一個月來拉一次草,當時的雪很大,就是那一天他來拉草,他晚一天也拉不到我,那個草堆巨大,比那個教堂還大,他挖偏一點也挖不到我,他就把把我救了.那天我相信有神績,AMEN!猶太人喜歡看神績,希臘人求智慧,中國人看神績也看智慧,中國人把猶太人和希臘人都包在一起,這是我們的智慧.那一天我看到了神績,我認識了神.他們把我關到蘇聯KGB的監獄裡面,我一進監獄,給我的感覺是極大的我喜樂,極大的平安,為什麼?頭一天晚上差點沒被凍死,這裡面很溫暖,頭一天晚上差點被狼給吃了,這裡很安全.監獄門一鎖上,不要說是狼,誰也進不來,因為外邊有衛兵在把守,你就感恩了,感謝主!那天我懂什麼叫自由.自由不是環境,自由是心境,自由是信心,當你有了這個信心的自由,在什麼監獄裡面,你都會有喜樂,所以我就在裡面唱歌,我就唱一首基督教的歌,我就只會唱那首歌,其他我不會,四句話:你的頭髮已被神數算,你的重擔主已替你擔,你不要為將來的事情去盤算,主內有真的平安.那天我懂得信耶穌的好處,信耶穌以後我就跟主禱告,認罪悔改,把自己所有得罪都向主訴說.當我把所有的罪都訴說完了,就像劉亞州將軍所說的,人就輕鬆了,我們不要背負我們的重擔,除掉一切的苦毒,除掉一切的仇恨,除掉一切的傷害,在基督裡面有一個新的生命,感謝主!這是我第一次體會到十字架的大能.所以保羅說,在十字架面前,智慧人是愚拙的,那些愚拙人反而得了智慧,祂說祂要廢去聰明人的聰明,滅絕智慧人的智慧,所以當我們看到十字架的時候,想到我們自己是何等的需要十字架那個思典臨到我們.

一九八九年我認識耶穌,祂伴隨我走了二十多年的路,這二十多年我走了很多曲折的路,但我始終記得耶穌對的應許:信靠祂的就不至於羞愧,叩門的就開門,尋找的就尋見.所以當我來到美國的時候,我更深地來經歷到上帝的大能.我記得我是九一年來到美國,我在中國藏了兩年,在蘇聯我要求去美國,但他們不同意:你覺得我們國家那麼隨便就可以過去,你也沒有簽證,對不對?我說美國不行,我到法國也可以,法國不行,敘利亞也行,敘利亞也不行,那我留在蘇聯可以嗎?留蘇聯也不行,後來他們做了決定把我送到中國,那時候我也不太害怕了,送就送吧!認識了主,槍斃了我還有永生,這個很重要啊!今天講你可能覺得很虛無,在那個時候我真的覺得非常重要,如果我死了復活了,你想我在北京大學走一圈,老師也嚇壞了,那些人尖叫:嘩!回來了!這個很好,這個信仰太好了,這個信仰可以從死裡復活.對我來說,今年你覺得很虛無,但對我來說,就在那一刻非常實際,後來我就回國了,他們把我送回中國,但他們給我留了一條路,他沒有送給中國政府,讓我自己走回去,他們把我驅逐出境.我在中國的山裡生活了兩年,那兩年,我所以為的驕傲都沒有了,上帝很奇妙!以為你驕傲嗎?你沒有了!北京大學的人都很驕傲,開除了!作為共產黨員也很驕傲,我們是無產階級先鋒隊,開除了!作為自由人也很驕傲,通緝命你沒有自由了!還有一個最有意思的,我最後還剩一項驕傲的,我還健康呢!這可以驕傲的.你不要以為人就是驕傲的,人的本領就是驕傲,對不對!健康都會令你驕傲,你甚麼知識都沒有,我有健康可以打工,也能驕傲.有一天健康會沒有的.我九一年來到美國,那時候逃到香港,那是我第二次出國逃亡,因為香港還是歸英國管轄,我不告訴你從哪上岸,反正我上來了,我上來以後去美國大使館,後來就把我弄到美國.到了美國後就去到普林斯頓大學,我的老師是香港中文大學余英時先生,在他那裡我們有一個課程 (program),所以我覺得不錯,師生也挺好,能和余先生在一起生活幾年,我覺得一生挺幸福的.因為那時候我和柴玲年齡最小,其他的年齡比較大,見到先生都必恭必敬,我倆也小不太懂得禮貌,推門就進去,進去就跟他聊天,也不管他在幹甚麼.我們聊夠了,我們就走,先生從來不說我們倆.其他人不可以的,因我們比較小…

學了很多東西.後來有一天,余先生說你們想念書嗎?想讀什麼?我說當然選文學啊!跟你學歷史也可以啊!這個不錯你可以研究近代歷史.那個時候我對台海關係比較有興趣,所以我就研究台海關係,寫了很多東西,看了很多書,參加很多坐談會,認識很多台海關係的專家,我對這個很感興趣.結果有一天,我生病住進了醫院,我進醫院,一進就四年,同學們都畢業了,我在醫院裡還沒出院,那天我才發現健康也沒有可以驕傲的,健康也可以忽然失去,所以人沒有出路.

從一九八九年到一九九五年我是在苦難中,在沒有盼望,在死亡的咒詛中,在孤獨中渡過的.我現在常想起那六年.人生六年不長也不短,可是那六年是我人生最美年華的六年,我是這樣過去.一進醫院,醫生跟我化療──腎癌,癌細胞第三期.醫生說:”你要化療”,”能活多久?”,”三年,頂多三年”,”不化療呢?”,”也差不多三年”,”為甚麼要化療?”,”因為化療還有一絲希望,不化療一點希望都沒有”.

那個時候,你可以在幹嘛呢?所以我就在醫院裡寫了一本書,就是你們看到那個《逃亡者》──人之將死,其言也善──寫了很多自己的軟弱,對女兒的思念,對父母的愧疚.化療做得還不錯,癌細胞不見了,但是腎臟衰竭了,我必需換腎,所以我就去了台灣,當我到了台灣,就住進了榮民總醫院,換腎等了一年都沒有腎源,我就禱告:主啊!你給我一個腎.每天就這樣禱告,後來有一個老將軍基督徒提醒我說:你這樣禱告,不合神的心意,你的意思是,主啊!死個人,死個人…

你知道嗎?主啊!死個年青一點的,因為我不要老人的腎.那我問應該怎樣去禱告?那你應該合主的心意去禱告,你要主醫治你.

這十字架是醫治大能的十字架,當耶穌的寶血流到你身上的時候,你的疾病會得醫治,你的傷害會被除去.我們不僅要醫治身體的疾病,也要醫治心靈的疾病.我們這一代人,心靈都被破碎了,那個人心靈沒有病?我有一次到美國國會裡有一個講演,出列者中有一個很有名的參議員,他前幾年去世了.他就問我,你要不要找個心理醫生?我認識一位朋友,可以免費的為我服務,我當時還挺不高興的,我覺得你什麼意思啊?你覺得我心理有病啊?後來我就想,一個人與世隔絕兩年,在醫院折騰六年,我估計心理也不會太正常.你覺得會正常嗎?每當我在醫院裡打開窗子,我看到柴玲開著小車去普林斯頓上課,背著書包,我就很難過.我就說:上帝為什麼你愛她,你不愛我?對不對?每當李錄從紐約來看我,告訴我他一天可以掙四十塊美金之後,我又想,上帝你愛他為什麼不愛我啊?我還信你,他們都不信你!這是很奇怪的一種現像,對不對?神不愛我了!可是我又不能離開祂,因為除了祂以外,根本沒有人愛我了.在醫院待了六年,根本沒有朋友來看你,久病無孝子,久病無朋友,看你一兩次也差不多了,人家也該去學習,讀書,工作.孤獨,苦難伴隨你,我非常痛苦,但神有醫治的大能,後來我的禱告變成了一個緊急的禱告:主耶穌我不再要腎了,請你醫治我的腎.因為你是出死入生的主,你是給我們生命氣息的主,你要醫治就能醫治.所以禱告變成緊急了.感謝主!所以神醫治了我,今天你看我站在你們面前,身體很好,身體所有的器官也正常了,上個月是我從一九八九年以來,腎臟最正常的一個月,已進入正常指標之內,在過去是正常指標之外,但是現在越來越正常.我來到香港做培訓,一天就有兩三個小時的休息,感謝主!給我們力量,上帝就醫治我們.AMEN!哈利路亞!所以這福音本是神等大能,祂不僅有赦罪的大能,還有醫治的大能.聖經說我們若認自己的罪,神是信實的,公義的,祂必赦免我們的罪,洗淨我們一切的不義.當我們遇到難處來到祂的面前,祂也必醫治我們.聖經說:因他受的鞭傷,我們得醫治；因為祂受的邢罰,我們得平安.(註一)

弟兄姊妹,相信福音的大能,在一個人的一生中,使你變為完全不一樣的人.

所以很感謝主!這福音也是個奧秘.在哥林多前書一章二十一節:「世人憑自己的智慧,既不認識　神,　神就樂意用人所當作愚拙的道理,拯救那些信的人．這就是　神的智慧了.」神的智慧跟我們的智慧不一樣,我們有人本主義的觀念去分析神,可是上帝用我們覺得愚拙的道理來拯救你,這就是福音的道理,是奇特的.我們中國人喜歡高深,其實真理是最簡單的,一加一等於二,不高深,這是真理.所以當我們來到上帝面前,發現上帝的大能就是要拯救我們.這福音標誌著神與苦難,罪惡深中的人類的認同.

人類都會經過苦難,你可能什麼苦難都沒有,但有一天你會死,死就是苦難.死不是苦難嗎?死是隔絕,死就是離開你的親人,死就是離開你熱愛的生活.我做牧師,常常上午證婚,下午就去主持追思禮拜.你看人多有意思,上午笑著進入一個新的人生,下午就離開了這個世界.你有沒有想到,人就是這樣,人類有什麼希望,盼望?我們經歷那麼多苦難,我經歷了苦難,感得苦難是與人有義的.聖經說,祂讓我們經歷苦難,是讓我們認識祂的律例；我們經過苦難,就可以安慰在苦難中的人,因為耶穌在十字架上所經歷的苦難,就成為我們人類的盼望,成為我們的救主.如果沒有耶穌在十字架上的苦難,我們今天毫無盼望,一點盼望都沒有,所以苦難並不都是壞事.苦難也是一個化了裝得祝福.當我們不求苦難,我們可以像雅比斯一樣不求苦難,雅比斯說:讓我不要遇見患難,你就祝福我就好了(註二).可是上帝還是給我們苦難,是為了讓我們經過苦難,被祂磨練得更純潔,更愛祂.有些人經歷苦難就不愛主了,可是當我們經歷苦難後,能更愛主的人,這就是上帝要揀選的.

第三點,我要講的是這福音有永生的大能,因為這福音是恩典.人類在死亡咒詛的時候,我們有了這樣一個恩典領導我們.所以以弗所書說:我們得救本乎恩啊!也因著信,不是出於自己,也不是我們行為怎麼好(註三),因為這樣我們會自誇,會驕傲.是因為這是神的大能,哈利路亞!所以經過苦難,後來九五年我身體好了回到美國,我記得我到台灣治病的時候,那時普林斯頓裡沒有一個朋友認為我會回來,我的醫生告訴他們,我去了就不會再回來,會回到中國了!畢竟台灣也是中國的一部份嘛.回到中國,我也感到很安慰,埋在台灣也總比美國強,當時候想的,都是這些.我記得余先生到台灣去講演,還到醫院看過我,他看了我說:誒!你看來沒有那麼差,樣子還不像馬上要走的樣子.其實他們那邊已經準備好我要走了!很奇妙的,我記得有一天,我覺得我真得很不想活了.在台灣那段時間,我就打電過去美國給我的牧師說,我覺得活不過天亮,這個苦難對我來說,有點越不過去,你趕緊過來給我受洗吧!這個世界上我什麼都沒有了,什麼都沒有撈到,死了我不能把天堂給錯失,如果這個都沒有了,那我這一生一定淒淒慘慘…(我問牧師,你看現在幾點?已經半夜三點了.如果你在醫院一,兩年就不會知道時間,醫生叫你,你就起來,醫生檢查,就起來,他不來你就睡覺,你跟本不知道是白天還是晚上.)那天牧師沒有來,但是來了一對馬來西亞夫妻,他給我看一段經文在路加福音二十三章(註四)──他說,如果你真的相信耶穌基督,你就得救了!你不要懼怕死亡──耶穌到一個地方叫髑髏地,就在那裡把耶穌釘在十字架上,然後兩邊有兩個強盜,一個左邊一個右邊,耶穌說,父啊!赦免他們,他們所做的他們自己不饒得,那兩個強盜呢?

那位聰明的強盜,智慧的強盜,有名的強盜就嘰笑耶穌說,:”哎!你不是上帝嗎?對不對?你可以救自己啊!你可以馬上下來,既然我們釘在一起,你也可以救我們啊!”

這個比較聰明他有理論,他知道祂是上帝,他希望祂能救自己也能救他們,他接受光也被得救了.可是另外一個沒有這樣說:”我們都一樣受刑的,你還不怕神嗎?祂沒有犯過一樣的罪,不像我們,我們是罪人.”在上帝面前不管你是老人,小孩,外國人,中國人,愚拙人,智慧人,有文化,沒有文化；在世界分有錢的,富足的,貧窮的,各式各樣的人…

但是在十字架的面前,人類就是兩種人:滅亡的,的救的；聰明的,愚拙的；認識耶穌的,拒絕耶穌的,就是兩種人.聰明的強盜,他拒絕耶穌,那個比較愚拙的強盜,他卻認識了耶穌,他說祂沒有犯過一樣的罪,證明他知道神,當猶太人說他沒有犯過罪,那只有神,對不對?在猶太人的文化觀念,沒有義,一個義人都沒有,唯有上帝,祂是義的本身,所以你要做那一個強盜?你選一個?我當然要做那個得救的強盜,我怎可以做那個聰明的強盜呢?我要做那個認識耶穌的強盜,因為祂在樂園裡,耶穌說:”你現在就跟我在樂園裡”,當然我希望在樂園裡啊!所以感謝主,九五年我進入神學院,我的病好了一些,就回到了美國進入神學院.那時神學院比醫院好多了.

把自己獻給主當活祭,也終於認識了十字架的大能,能夠讓自己走一個全新的道路,十字架在前頭,十字架在後頭.進神學院很有意思,神學院的同學跟世界大學同學完全不一樣.一分享的時候,他不是牧師的兒子,就是長老的女兒,因為我是共產黨員的兒子,小時候沒讀過聖經,讀的都是馬列之類這種的東西,到神學院就非常的辛苦,因為你的背景不如他們好.但是非常的奇妙,他們也分享說,當他們奉獻給主的時候,是放下了世界,他們放下世界什麼呢?他們有的是大公司的總裁,他們讀了世界最好的專業博士,他們放下來服侍主,他們認為這是最好的選擇.他們問我,”伯笠你放下了什麼呢?”我想一想我放下了什麼呢?我放下了苦毒,傷害,疾病…

我放下的都是負面的,我什麼都沒有,所以我沒有什麼重擔,神呼召我跟祂走,我就走了,因為我什麼都沒有──無產者無所畏懼──真奇妙感謝主,而且我放下了…

我以為世上所有的人都不會找你了,你就往前走好了,可是你還有問題,一些老朋友經常跟你打電話,最後連電話我也不接,因為我跟他講屬靈的話,他們也聽不懂.最後醫院還要找你,因為欠錢.

我在台灣欠了一千多萬台幣,我在美國欠了一百多萬美金,來到海外沒掙幾個月工資,就欠了那麼多的錢.但是我沒太在意,因為醫生說我只有三年命,我三年就死了,誰還還不還,對不對?死就是沒反應,死就是隔絕與帳單隔絕,結果,沒有死,慘了!問題來了,活著就要面對苦難,活著要面對挑戰,活著就要靠耶穌,否則你沒有喜樂.我接到寄來的帳單的時候,打開帳單真的嚇了我一跳,那麼多的項目啊!我還拿一個計算機算算,拍!拍!拍!一個帳單三十多萬美金,一個十多萬美金,一個二十多萬美金,最少的八萬美金,我再一算,已經一百多萬啦,我想我在美國真的掙不了一百萬,我能掙一百萬嗎?我在餐館打工,才十幾塊錢一個小時,我能掙一百萬嗎?我的人生有什麼前途?沒有前途了,沒有盼望了,很痛苦的.跟老師說,老師也解決不了,跟牧師說,牧師說:”你看天上的飛鳥,也不種,也不收,也不積蓄在倉裡.我們天父尚且養活它,你不是比飛鳥更貴重得多嗎?”我說,我是比它貴重啊!可是我欠錢啊!飛鳥不欠錢呀!對不對?另外,它欠了也沒關係,它飛走,我飛不走呀,我往裡飛?我有帳單,對不對?禱告!每天禱告看著十字架的耶穌基督,我總想起一句話,祂得恩典會夠我用的.主啊!你的恩典夠我用,那你可以替我還帳嗎?我們很實現的,對不對?醫院打電話找你,原本上教會敬拜很喜樂,一回到家中,一接電話是醫院的,醫院很有意思說,你要列一個還款計劃,從你現在開始,你的百份比多少,就是一個月還多少錢,還到你死為止.工資高,多還；工資少,少還.對不對?你要負責,因為你欠的,你該欠的.我那時候也後悔,但人沒有後悔藥…

我記得住醫院的時候,我住兩個人的病房.後來,美國總統布殊,給我寫了一封信(因為沒有電郵只有傳真),傳真到醫院裡,院長一看,嘩,總統的信,還說了那麼多祝福的話,他說我是最勇敢的人,美國人不願意把勇敢給外國人的,你知不知道?院長幫我換了全醫院最貴,最好的病房,那個時候一個晚上一千六百美金.那我說多少錢?他問我:你有錢沒錢?我說我沒錢!那你問他幹嗎?你問多少錢也沒有意義吧!你沒錢你就住吧!反正你住不上多長,沒想到我住一年,然後到台灣,回來了,要錢.哎ye !心裡面沒有喜樂,覺得自己一生全是苦難,一九八九年學運,我去錯了,我去幹嘛?好好的,對不對?都是同學們,老師們鼓勵我們去,都是老師的錯,最後去了,出了這麼多的問題,欠了那麼多的債.結果還是要禱告的,神總是有恩典的.我一天禱告很喜樂地回家,結果接到電話,我一下子心又涼下來,我就很不高興,我說,以後不要跟我打電話 “I don’t speaking English”,後來第二天,他找了一個香港人跟我要錢:張先生你要還錢,我是醫院的.我說這個醫院很厲害,問我要帳,我說我不懂英文,他就找個香港人,你不可能說你不懂中文吧!對不對?後來一天他給我打電話,我說你以後不要找我啦,那我找誰?你找我爸去,我阿爸有錢,你阿爸在國裡做官了?發財了?他有多少錢啊?我阿爸在天上,哦!你爸死了!我說,你爸就死了.我爸叫”GOD”──上帝祂是我阿爸父,哈利路亞!結果你說他說什麼?他也是我阿爸,他也是基督徒.那就簡單了,咱們是同一個爸,那要什麼錢,對不對!他說,”那不知道啊!””阿爸給不給錢我怎知道啊?牧師沒教導啊!你那個教會啊?你要想了解阿爸的心意你要禱告的”,他說”是啊!是禱告啊!那為什麼你欠錢我禱告啊?”我給他一段經文馬太福音十八章十九節,我說:”兩個人同心合意,無論求什麼,上帝都成全”(註五),他問,”哪兩個人?””你和我啊!”所以他再不打電話,他一打電話,我就跟他開禱告會.我在神學院畢業那年是兩千年,二零零零年六月十五號,我畢業,上帝給我兩個禮物,第一個上帝免了我一切的罪債,也免了我地上一切的帳單,醫院打來電話,這個姐妹說:”張弟兄阿爸給錢了.”我說,”給多少?”她說,”全部.”我說,”誰給的?””阿爸給的,我們都不知道誰給的,反正你不再欠債了,給你傳真一個表格,填一下,簽上你的名字,你就不再欠債了.”那天我流淚了,我想到兩個字叫「恩典」.恩典!什麼是恩典?我們欠了要還,不要了就是恩典,我們犯了那麼多的罪,在上帝面前,污穢不堪,耶穌的寶血,洗淨了,潔淨了,我們成為新造的人,這就是恩典.上帝不再追討我們的罪債,在基督裡面成為上帝兒子,成為神國家裡的人,這就是恩典!不再欠債,這就是福音的大能.Amen!

使我們脫離咒詛,進入永生的大能；使我們脫離苦難,進入平安的大能；使我們脫離憂愁,進入喜樂的大能.這就是福音的大能,祂使變為一個全新的人,一個新造的人.那天我知道我們的神,是多麼的真實,祂知道我任何一點的憂愁.祂來幫助我解決,祂不要我背著這個債,走到福音工場去,祂要親手的幫我走進福音工場,為祂建立教會,收割莊稼,牧養羊群,我非常的感恩.

第二個神給我的禮物,是我的女兒.我的女兒八九年被通緝,她十八個月,我好多年沒見到她,一直想來到美國,一直來不到,拿不到護照.但就是在我畢業前的一個禮拜,她拿到了護照,飛到了美國,她就睡了一個晚上,參加爸的神學院畢業典禮,我知道我們的神是現實的.Amen!

神知道我們一切細微的需要的,祂就是那樣的真實,不需要你去做什麼,祂都替你想好了.就在我的畢業典禮,祂把我一切的重擔,最大的包袱,全部都解決了.我女兒,在洛杉機下飛機,我去接她,一看到一個大姑娘站在我的面前,我眼淚就流出來了,人家一個眼淚都也沒有掉,就跟我講:爸!我要回家.感謝主!哈利路亞!

親愛的弟兄姊妹,你能認識上帝福音的大能嗎?祂在你的生命中,祂可以改變你,聖經說,若有人在基督裡,他就是新造的人,舊事已過,都變成新的了.(哥林多後書 5:17)

講一個小故事結束來結束今天的信息.信主後,還沒有讀神學院之前,也常常的問:我所信的主,祂真的就有這個大能嗎?

這個福音真的有這個大能嗎?除了我們基督徒尊敬祂,世界的人會不會尊敬祂呢?我們基督徒,有的時侯,朋友看到我都會封住我說,”你幹嗎?你去做個傳道人,做個基督徒,你還有沒有出息了?”有時候我跟這些在美國做大學老師,這些天文的朋友在一起的時候,他們都會嘲笑我,他們都覺得,”你是最沒出息的”,當然也有人會很理解,”他也就這樣,他還能幹嘛?他這個身體,信耶穌!求個心裡平安,也就得了,”對不對?這是一種很正常的看見,我每次跟柴玲傳福音,她都會說:”Are you sure? (你確定嗎?)你講的東西是真的嗎?你真信嗎?你是不是逃避到一個地方,找個心裡平安,你也就這樣,行了!你去吧,去吧!你看我們在世上做很多的事,你什麼都做不了呢!你還能幹嘛呢?”前幾天她的自傳出版了,我帶了一本,她讓我送給中文大學,我為她寫了序,因為她現在是基督徒了,她認識了這個福音本是神的大能,使她謙卑在十字架的面前,她在書裡面認了很多很多我都不知道的罪.最後她說,認識耶穌基督是她一生中最好的禮物!超過她的丈夫,超過她的兒女,所以我們認識福音,是要我們謙卑放下自己,認識祂的大能.

九四年我去美國國會,參加總統的早餐會,那時侯身體還不是特別好,我從醫院裡出來,國會兩院請我去參加早餐祈禱會,那一天我真是認識到了,十字架的大能,認識我的主,就在那一刻我生命翻轉.我進入會場的時候,好多人,他不是請我一個人吃飯,請了好多人,大概有三千人.全世界各國都有代表團到來(總統,參眾兩院的議員,還有世界各國的政要代表團),我坐到第七桌,他就介紹了幾個總統和國王,王子,其他都沒有介紹,當他介紹到加納國王的時候…(因為介紹的人很多,當介紹那個人的時候,那個人的圖像就印到屏幕上面,大家就鼓掌就好了,因為時間的關係),可能是攝影機的錯誤,一下子把我的樣子放到鏡頭上去了,因為我離得很近,第七桌在前面,我還不知道是我,因為有時候自己從來不照鏡子的,我兩臉浮腫,腫得很嚇人,不照鏡子的,真是我嗎?大家在鼓掌…

結果我旁邊這個人站起來了,他一站起來,鏡頭轉向著他,他才是真正的國王(他坐在我旁邊,我都沒認出他是國王),他站起來了,大家當然笑得一塌糊塗,笑了攪錯了,我自己也笑,我就拍拍他肩膀,坐下來的時候,總統跟全體說,讓我們全體起立,讓我們在萬王之王,萬主之主,耶穌的面前起立,讓我們為世界的和平來祈禱,我們就拉起手來起立,他拉著我的手,我拉著他得手,我不知道他是不是基督徒,他可能不是,加納好像不是基督教的國家,可是當我們禱告的時候,我就發現,他低下頭了!他是國王,他低下頭了!克林頓也低下頭了!他是總統,也低下頭了,我發現所有人都低下頭了,在耶穌基督的面前,在釘在十字架耶穌基督面前低下頭,那一刻,我的眼淚一下子流出來,我才知道什麼叫做基督徒,我承認我的主是多麼的尊貴,榮耀,祂真的是萬王之王,萬主之主,所有的王都要在祂的面前低下頭來,因為祂是和平之君.那一天,我才懂得做基督徒的身份是何等榮耀,上帝兒女的身份是何等榮耀,像彼得在彼得前書二章九節說:

「惟有你們是被揀選的族類,是有君尊的祭司,是聖潔的國度,是屬　神的子民,要叫你們宣揚那召你們出黑暗入奇妙光明者的美德.」我們是天國的使者,天國的大使,帶著耶穌基督的福音向世界宣告,世界沒有盼望,人類沒有盼望,唯有祂,釘在十字架上的耶穌基督,唯有這福音是神得大能.靠著祂,我們可以得救,可以成聖,可以過全新的生活,有永遠的盼望.哈利路亞!

我們來禱告,在禱告之前,我想說一句話,結果你今天願意,讓福音的大能進入你心中,願意十架的寶血來遮蓋你,除掉的重擔,願意接受,耶穌做你生命的主,你舉手我為你禱告,有沒有這樣的人?感謝主!還有沒有?感謝主!還有沒有?…

感謝主!我看到兩位從中國大陸來的,那天我在餐聽吃飯碰到他的.感謝主!還有沒有?…

今天你願意接受耶穌的,你站起來,我為你做個禱告!

感謝主!感謝主!當你認識這個福音大能的時候,祂的寶血會遮蓋你,祂會洗淨你的罪,祂赦免你一切的罪,祂也可以醫治你身體的軟弱,心靈的傷害,祂可以給你永生的盼望,因為這福音就是上帝的大能,祂要救一切相信的,祂的愛今天要臨到你,把祂的盼望給你,你不再一個人孤獨的行走人生的道路,是有耶穌基督伴隨你,引導你走一生的,永恆的路.還有沒有?感謝主!哈利路亞!還有沒有?

我們一齊來來禱告:「天父我感謝你,在這個時刻,我把每個站立在你面前的,打開心接受你作為生命救主的弟兄姊妹,交到你的手上.主啊!請你聖靈的大能進入他們裡面,讓他們一生跟隨你,讓你的寶血潔淨他們,讓他們一生更像你,事奉你,跟隨你,榮耀你,哈利路亞!我們也為香港中文大學來禱告,我們也為香港來祝福,每當我們來到這塊土地,我們都心存感恩,每當我來到這裡,我會感覺到,其實這裡是我的家!祝福香港!祝福中國!奉主耶穌的名禱告,Amen!」

註一:那知他為我們的過犯受害,為我們的罪孽壓傷．因他受的刑罰我們得平安．因他受的鞭傷我們得醫治.(以賽亞書 53:5)

註二:雅比斯求告以色列的　神說,甚願你賜福與我,擴張我的境界,常與我同在,保佑我不遭患難,不受艱苦．　神就應允他所求的.(歷代志上 4:10)

註三:你們得救是本乎恩,也因著信,這並不是出於自己,乃是　神所賜的．(弗2:8)

註四:到了一個地方,名叫髑髏地,就在那裡把耶穌釘在十字架上,又釘了兩個犯人,一個在左邊,一個在右邊.當下耶穌說,父阿,赦免他們．因為他們所作的,他們不曉得.兵丁就拈鬮分他的衣服.百姓站在那裡觀看.官府也嗤笑他說,他救了別人．他若是基督,　神所揀選的,可以救自己罷.兵丁也戲弄他,上前拿醋送給他喝,說,你若是猶太人的王,可以救自己罷.在耶穌以上有一個牌子,〔有古卷在此有用希利尼羅馬希伯來的文字〕寫著,這是猶太人的王.那同釘的兩個犯人,有一個譏誚他說,你不是基督麼．可以救自己和我們罷.那一個就應聲責備他說,你既是一樣受刑的,還不怕　神麼.我們是應該的．因我們所受的,與我們所作的相稱．但這個人沒有作過一件不好的事.就說,耶穌阿,你得國降臨的時候,求你記念我.耶穌對他說,我實在告訴你,今日你要同我在樂園裡了.(路加福音 23:33－43)

註五:我又告訴你們,若是你們中間有兩個人在地上,同心合意的求甚麼事,我在天上的父,必為他們成全. (馬太福音 18:19)



\newpage

\section{上主仍是我們911後的}
\label{sec:124}
\textbf{沈祖堯教授 - 上主仍是我們911後的}
\newline
\newline
\hyperref[sec:123]{< < < PREV SERMON < < <}
~
\hyperref[sec:toc]{[返目錄]}
~
\hyperref[sec:125]{> > > NEXT SERMON > > >}


日期:2011年11月13日

題目:上主仍是我們911後的避難所

經文:詩篇46

證道:沈祖堯教授

剛才那首詩歌非常優美,非常動聽,它亦與我們今日所讀的聖經很有關係.

在9月11號那日,我剛剛與我們的大學(中文大學)副校長,校友會辦事處的負責人一起到北美洲探望中文大學的校友.沒想到我們會選擇了這麼特別的一個日子──9月11日起飛.後來我到達了加拿大的時候還在想,是否真的要在9月11日 (註1) 坐飛機去美國呢?但結果我們還是去了美國.9月11日的上午,我住在三藩市一間酒店裡,因為 jet lag(時差)的緣故,我很早便醒來,當時是當地時間早晨5時,我打開電視,當時也是紐約時間的早晨8時,我從電視上看到Obama(美國總統奧巴馬)在讀一段聖經,原來當時他們正在紀念911事件發生的十週年,Obama總統與布殊(美國前總統)及他們兩位夫人到紐約911的現場主持紀念活動.在會場,如果你今日有機會到紐約 Ground Zero 的地方,即當時世界貿易中心雙子塔的所在地,你會看到原本那兩座大廈的位置已變成了兩個洞,兩個水池,水池是四方形的,大小正正等於之前原本兩座大廈的面積,所以有兩個很大的四方形水池,水池四面也有水流入,流到底部的時候有一個洞,水便向洞的方向流去,你不會看見水流去那裡,這個特別的設計是喻意著當日911,有很多的生命傾倒,墮進深淵裡,今日,原本雙子塔的位置不會再建立任何樓宇,他們重建了另一幢很高的樓宇在外面,奧巴馬總統及布殊總統就在那裡主持911十週年的紀念活動.

我現在重讀紀念典禮當日總統所讀的聖經(詩篇46)給大家聽,因為他以英文讀,我先以英文重讀 (American Standard Version): 1God is our refuge and strength, A very present help in trouble. 2 Therefore will we not fear, though the earth do change, And though the mountains be shaken into the heart of the seas; 3 Though the waters thereof roar and be troubled, Though the mountains tremble with the swelling thereof. Selah 4 There is a river, the streams whereof make glad the city of God, The holy place of the tabernacles of the Most High. 5 God is in the midst of her; she shall not be moved: God will help her, and that right early. 6 The nations raged, the kingdoms were moved: He uttered his voice, the earth melted. 7 Jehovah of hosts is with us; The God of Jacob is our refuge. Selah 8 Come, behold the works of Jehovah, What desolations he hath made in the earth. 9 He maketh wars to cease unto the end of the earth; He breaketh the bow, and cutteth the spear in sunder; He burneth the chariots in the fire. 10 Be still, and know that I am God: I will be exalted among the nations, I will be exalted in the earth. 11 Jehovah of hosts is with us; The God of Jacob is our refuge.

或者我們用中文讀一次當日奧巴馬總統為紀念911而讀的聖經──詩篇第46:1-11:

神是我們的避難所,是我們的力量,是我們在患難中隨時的幫助.

所以地雖改變,山雖搖動到海心,其中的水雖匉訇翻騰,山雖因海漲而戰抖,我們也不害怕.

有一道河,這河的分汊使神的城歡喜；這城就是至高者居住的聖所.

神在其中,城必不動搖；到天一亮,神必幫助這城.

外邦喧嚷,列國動搖；神發聲,地便鎔化.萬軍之耶和華與我們同在；雅各的神是我們的避難所.

你們來看耶和華的作為,看他使地怎樣荒涼.他止息刀兵,直到地極；他折弓,斷槍,把戰車焚燒在火中.

你們要休息,要知道我是神.我必在外邦中被尊崇,在遍地上也被尊崇.

萬軍之耶和華與我們同在；雅各的神是我們的避難所.

你知不知道為什麼我要你們一起讀,因為這段聖經至少我有三個字不會念的,哈哈.不過,這段聖經的意思給了我很深的感受,特別就911事件而讀.當日當總統讀完後,接著的項目是每位遇難者的家庭出來讀出他們遇難家人的名字,這部份每年也進行,但今年有些不同,因為在911大型的monument(紀念碑)上,即我剛才提及的兩個水池旁,刻上了3000多名死難者名字的紀念碑,所以他們一邊讀一邊會摸著碑上刻著的名字,有些死難者家屬更用紙筆印上碑上遇難者的名字,有些人在紀念碑送上玫瑰花,場面很令人感動,如是者每位遇難者家屬輪流讀出他們死去親人的名字,用了數小時的時間將3000多位死難者的名字讀出.然後我看到電視鏡頭突然一轉,播出另一個畫面,原來在911當日,事發時有很多人拍攝到當日發生的情景,有用流動電話拍的,有用相機拍的,但這些鏡頭在電視上從沒有播出過,當日是第一次播出,所以感覺很新鮮,很震撼,完全是第一手資料.有些人住在Downtown Manhattan(曼哈頓市中心)的附近,其中有一位女士,她從家裡的窗可以望見雙子塔,在當日8時15分,當第一架飛機撞入世貿(世界貿易中心)第一幢樓的時候,她拍攝到當時情景,鏡頭裡聽到她震驚的大叫:Oh, my God, my God, this is not real! 不會吧,是否我眼花!這不是真的!發生什麼事?

那架飛機怎可能這樣撞向世貿!她以為那是意外,從沒想到那是有預謀的,因為根本沒可能想像到會發生恐怖襲擊；但再過大約半小時,當她看到第二架飛機撞向大廈,她便驚叫起來:oh, This is insane! 發神經!無可能!怎會有兩架飛機撞入大廈裡!?

若你當年911有收看電視的話,你還會記得當時那震撼的情景:在兩幢樓還沒倒塌前,上面沖天大火,電視上看到有人從高樓上跳下來,你還記得嗎?因為當時被火燒的大廈太熱了,有些地方溫度更可能超過1000度,所以中間的鐵柱才會溶化,所以在大廈裡的人才會受不了,一個個的跳下來,當時拍攝到這些畫面的人,心裡很難受,很懼怕…為何大廈變成鍊獄一樣,為何人間鍊獄會在紐約最繁榮的地方出現!?而當我再看到這些畫面也目定口呆,因為很久也沒看到這些情景,我不太記起當時的細節,但我和他們的感受是一樣的,到底為什麼?全世界人完全沒有心理準備,整件事根本不像是真的；記得事發當日,紐約天氣很好,天色很藍,但突然間卻像拍電影一樣,簡直不可想像,不可理喻,但這就是911.

另外,在今年3月(2011年),我還記得當時我身在辦公室,正在準備當日下午需要教的通識教育課程,我桌上有兩部電腦,一部是普通用的,另一部作播出有線新聞用,由於不停的轉播新聞,突然間,我看到日本地震,之後更有海嘯.我看到大水不斷沖上來,當時橋上有人還在交叉著手,看著水沖上去的情景,但水越沖越近,剎那間,船被沖上橋,沖向汽車,汽車就像玩具一樣被沖走,還記得當時有架巴士看見沖來的大水,便想向後退,但我們在鏡頭看到,它背後也有大水沖向它,結果,那架巴士被淹沒了.當我埋頭準備教材的時候,突然看到此情此景,我簡直難以相信,怎會發生這樣的事?這段聖經說: The Ocean roars, 海好像在大叫,海好像在大怒,日本海嘯,地震,跟著發生核電洩漏.還有,今年發生很多的事情,都令我們非常驚訝,由年初一月開始,我們看到一個一個的國家發生叛亂,一個一個的國家發生所謂的革命,一個一個的政府倒台,在英國最繁榮的倫敦街上竟然發生暴亂,大家如果去倫敦,可能也會去Oxford Street購物,因為那裡是全倫敦最多名店的地方,但誰會想像到Oxford Street那裡也會有人打破玻璃搶東西,有人燒汽車等等,誰會想像到這些不可能發生的事也會發生?在這個月及上個月,我們看到歐洲整個金融體系的collapse(瓦解),我們怎能相信這些事情會發生?一個又一個的國家面臨破產,經濟評級被下降,更有人說,歐元會貶值得被取消,歐盟可能因此而瓦解.這怎可能會發生?整個歐洲怎可能在短短的幾個月,幾個星期內便破滅?還有很多很多的事情在我們身邊不斷發生,就像這段聖經所說:山也在動搖,海也在翻騰,又因為海的漲,整個地方也被淹沒.This is happening everyday, this is happening just around us(這些每日都在發生,這些不斷在我們身邊發生).

或者你會說,這些不關我們的事.香港沒有這些事情發生,沒有地震,沒有海嘯,沒有恐怖襲擊,暫時也沒有人上街打破玻璃,燒車等等,雖然每個星期日也有人抗議示威──所以現在星期日最好避免到中環,因為會遇到塞車,但這些也不算嚴重.不過,在你的生活裡,你的家庭裡,是否也可能發生過,甚至最近正在發生像海嘯的事情?最近有沒有人告訴你患了嚴重的疾病?可能你自己也不能相信原來一條血管已塞,冠狀動脈只餘下百分之十的通行程度,要做手術放一條鋼絲進身體內?這些可能就是你的海嘯；又或者,你的男朋友或女朋友對你不好,拋棄你,你的丈夫或太太對你不忠,有了第三者?這些可能對你來說,等於一個原子彈的殺傷力!又或者你最近發現你的子女做了一些事情是你難以想像的──他們在「索k」(吸毒),誤交一些你難以接受的損友!作為一個父母,如果發生這些事的話,對你來說可能就像海嘯沖來一樣.你又或者正遇上一些官非,或因為你正在選舉區議員而遭人妒忌,因此有人因你的陽台過大對你提出訴訟,又因為你沒有繳稅而去揭發你.當然相對911,這些沒那麼嚴重的死亡,但對你來說,這些事對你的impact(衝擊)可能不下於911的嚴重性.

在過去12個月,作為中文大學的校長,我的911是什麼呢?不單只有911,我們還有818,及很多其他的號碼曾經出現…(但818其實也不是關於我們的).但我們有女神像事件,中山像等等難題,我們有數宗法庭訴訟要面對.例如可能大家也知道,最近要到終審庭才完結的,關於是否要採用中文,英文或普通話作為教學的事件,這訴訟不單只花了大學數年的時間,更耗費了龐大的人力,物力.

還有,我記得一天上午,我的司機突然告訴我,有一個同學跳樓,後來知道原來這個同學還沒入讀的,但不知為何,參加了迎新營後在家裡的陽台跳樓自殺.這些突發事件令我難以想像,難以預計.每一件都可以成為911事件,或者小型的911事件,看似不嚴重,但對我或對學校是嚴重的事. 我還記得大約在兩星期前….每日我的秘書也會給我行程表確定當日要做的事,但因為自己也有些偷懶,已沒有預先看明天及後天的行程,只於早上看當天行程確定當日需與誰會面.但那天大約上午7時30分回到辦公室,一看行程表才發現當日要到香港電台接受訪問,訪問我的人更是出名言詞尖銳,針對政府,針對大學的主持人,於是心裡有些遲疑,考慮缺席,但時間又不能容許失場.但再細看資料之下,發現那牙尖咀利的主持人還邀請了其他幾位嘉賓,但這些嘉賓卻不是接受訪問而是到來一起訪問我的!這幾位人士比該主持更激進,其中兩位是在其他大學憤然辭職,然後在專欄裡常常罵他們的大學.我開始雙腿發軟,心想要不要借病藉詞推卻,因總比失言好,但想清楚不能如此不負責任,因為節目要播放,而且這不單是在電臺(只聽聲音),而且也是香港電台電視部的節目之一(要上鏡頭).我未曾看過該節目,但當初他們邀請時,只說該訪問是關於《大學教育之道》,所以沒有仔細弄清楚下我便一口答應了.我覺得他們有點兒陷害我,因為他們沒有提及另外還有三名嘉賓是來訪問我,而非一起接受訪問的.但也沒辦法,因為9時節目便要開始,我惟有硬著頭皮去.於是8時30分出車,但到達電台的時候,更把我多嚇一跳,原來除了提及的主持,三位激進的專欄作家外,還看見香港大學姓李的同學!原來她說是特意申請參加該節目的.我心想,怎麼辦?誰知再進到錄影廠看見那裡有一大群的大學生,也是來參加錄影的,原來這節目是可以大家一起加入訪問而非單單主持人提問,對於我來說,這簡直就是一個快要覆蓋過來的海嘯,怎麼辦?怎去招架一連串問題?結果我現在不告訴你,請你於11號自己收看該節目你便會知道我的慘況.但我只想 illustrate (說明)一點, 我不知道你會不會這樣,但我自己每日都戰戰兢兢,因為每日也要面對突發事件；若然你沒有這種感覺,那麼你的危機感可能不足夠.其實每日我們都有可能會發生一些意想不到的事,可能是一個大災難,可能你會無法相信會發生在你身上,試想想,誰會料到,到北京旅行也會無緣無故地遇上車子後溜,而使三人喪命的事件!很多事情,真的不是我們可以預測得到的.

回頭看這段聖經,詩篇第46篇,這裡說:當我們遇到這些事的時候有什麼想法呢?

第一,它說:上帝是我們的避難所.我不知道你是怎樣?但當我遇到困難的時候,我會問:上帝,你去了那裡?為何你不在這裡?為何正當我遇到這些困難,這些人的時候,要面對這些鏡頭的時候我找不到你,你在那裡?

在2003年沙士事件,我曾經就有這個非常強烈的感覺.當時立法局追查何以醫院開門又關閉事件,立法局傳令我們在一班議員面前被質問,我心想:怎麼辦?事隔一段時間,很多詳細情況也忘掉了,我們該說什麼?我心裡就問上帝:你到底有沒有在幫助我們?在危難中你有幫助我嗎?當那件事發展下去的時侯,所得到的訊息,正如我之前也曾提及,報告裡會針對4人接受批評:包括當時的衛生局局長,衛生處處長,副處長及我,在政府的報告裡會被點名批評,我可以想像緊隨下來會惹來很多官非,很多投訴.於是有一晚,我在電郵中向我的同學呻嘆了一句:God, why do you forsake me…上帝,你為什麼你離棄我?你記得這句話是誰說的嗎?是耶穌說的,在祂被釘十字架的時候,祂說:我的神,我的神,為什麼離棄我?Why did you forsake me?為何你把我留下去面對這些?為何要離棄我?

當我們遇到困難的時候,我們很容易也會有這樣的感覺.但這裡說,上帝是一個安全的避難所,雖然在前面的山會倒下,海也不斷沖來,但我們也不懼怕.為何可以這樣呢?怎可以做得到呢?這裡還說,上帝可以製造奇蹟.它說:雖然山也倒下,海不斷沖來,雖然周圍的國家也有叛亂,有動盪,但有一條河在城裡,上帝是住在聖所,住在這城裡,這城不動搖.到了第二天,他們便得救了,跟著它說:有很多列國在爭吵,時局動盪,但上帝發聲,一切就溶化,刀槍被折斷,弓也被消除,最後回復平靜,最後也取得和平,然後上帝的名得以稱頌. It is miracle(這是神蹟)!突然間你的麻煩,突然間你的disaster(災難),突然間你的疾病,你不能解決的事,神都會神蹟地幫你解決.上帝是可以有我們超乎想像的能力!我不知道你讀此聖經的時候有否留意,經文裡兩次重覆了一個人物的事蹟,如果以Eugene Peterson的聖經,其實是重覆了三次:這個人物是誰?你留意到嗎?他就是 Jacob(雅各),聖經說了三次雅各與神摔跤的事,在創世記32章:提及有一晚,雅各發夢,當時他在獨睡,為何他一個人獨睡呢?因為他快要面對他的敵人,所以他先把他的妻子兒女送走,以免遭到敵人的迫害,然後當晚他獨自在曠野睡覺,作了一個夢,夢裡看到自己與上帝或上帝的使者摔跤.其實當時雅各先佔上風,所以聖經說:當時上帝的使者見自己不能勝雅各,便在雅各的大腿窩摸了一下,而我們念醫科的人,就很有興趣地去研究到底是摸了大腿那個部位?是那塊muscle(肌肉)出事? 還是hip joint(髖關節)脫臼?總之摸了一下後,雅各便跛了不能動,躺在那裡.但當雅各輸了後,他不讓使者離去,說:我不會讓你走,直到你祝福我.雅各問使者是誰,但上帝使者沒有回答他,只是祝福了他.最後,當雅各醒來的時候,便給那地方起名叫毗努伊勒(神的意思),意思是:我面對面見了神,我的性命仍得保全 (I still survive, survive quite well actually).雖然我在這裡跛了,但上帝祝福了我.這個事蹟在我剛才所讀的詩篇裡提及了三次,是要告訴我們:在你最危險的時候,最困難的時侯,你是否也像雅各一樣,用盡自己的能力,方法,聰明想去 overcome(克服)困難呢?但上帝卻要摸你一下,然後廢你武功,於是乎你發覺,無論你讀了多少書,儲存了多少錢,找多少人幫忙,想了多少方法,都不能解決,於是乎你會感到在一個非常無助的環境裡,然後你會捉緊上帝,對上帝求助說:怎麼辦?我已經完全沒有辦法,求你救我.在那個時刻,上帝就會祝福我們,他提了三次這事,詩人在他極大的困難中不能應付的時候,他提了三次,他提到困難大如列國的戰爭,大如金融海嘯,大如身患絕症,但令我深深受益的是在第10節:你們要休息,要知道我是神,英文譯過來是be still──不是叫你睡覺,而是叫你不要動,要鎮定──”be still and know that I am God”,原來當我們越是恐慌,越是慌亂的時侯,我們沒辦法知道上帝在那裡,於是我們便會說:Why are you forsaking me…為何你離棄我?但be still and you know that I am God,你就會知道我是上帝,be still and you will see that my name be exalted,我的名就會被高舉.

這一段聖經,給予我很大的鼓勵.雖然我沒有經歷911事件,但經歷過沙士；雖然我沒經過金融海嘯(因為我沒有投資,所以沒有虧損),但這些日子有很多人非常不開心,像世界末日般；雖然我沒身患絕症,但沙士一役之後我也患了一個病,正如我之前提及,我的心跳很不正常,間中會多跳一下或少跳一下,這病常常在remind(提醒)我,我有心跳是上帝賦予的,不是我所能控制的.但這裡告訴我們:be still and know that I am God.由去年7月1 日開始,我每天都在學習這話:be still,要安靜,要平靜,然後你才知道我是上帝.當我搬到校長官邸,即宿舍的時候,發現那裡是一處非常特別的地方,因為原來當我站在草地上的時候,我可以望到整個中文大學的校園.哈哈,所以假若你上班或上學遲到的話我是看到的.我望到整個校園,然後每天早上,在上班前我盡量都會在那裡先安靜一下,然後為大學祈禱,為我們正遇到的問題或將會遇到的問題,或今天會出現的問題,未有答案的問題去祈禱.因為在這世界上實在有太多東西不是我們的能力所能控制,所能預測,所能應付.但上帝就是要我們學習這樣的功課:讓你遇到困難的時候,你才知道我們的能力是這麼有限的,遇到困難你才知道你是要依靠神,你才能渡過的；當你安然渡過後,你才會發現,雲上的太陽根本沒有挪移過,只是我們在雲下看到狂風大雨,但你有沒有想過,只要我們能夠transcend(升高)自己在雲上面或者乘坐飛機飛到雲上面,你便會看到原來雲上面的東西從來沒有改變過──Be still and know that I am God.未來的日子裡面,還有很多富挑戰性的事情,有很多東西是你從來未準備過的,caught you by surprise, unprepared(出奇不已,難以準備的).我相信我一定會遇到很多這些事情,但希望這段聖經能夠繼續給予你,我很大的支持,無論你是在讀書或在大學裡工作的,又或是為人父母的,我希望這段聖經都可以給予你很大的祝福.

希望我們可以再唱一次雲上太陽…

我們再唱兩次這首歌來提醒我們:無論山怎樣動搖,海怎樣泛濫,列國怎樣的動盪,學校怎樣的作反,我們也不用懼怕,因為在雲上的太陽,我們的上帝告訴我們:Be still and know that I am God.

註1:911事件,又叫911恐怖襲擊,由阿爾蓋達組織發動連串針對美國ge 自殺式襲擊.

2001年9月11號在美國本土,數人騎劫了美國航空同聯合航空四架客機,兩架撞入紐約市世界貿易中心,一架撞入五角大樓.另外一架聯合航空93號班機也被騎劫,但機內乘客與劫機者博鬥,阻止客機撞入人多地方,雖然客機最後成功避免地面傷亡但但客機內40人全部犧牲.整個襲擊一共有 2,986 人死,美國紐約世界貿易中心雙塔在內的六幢建築物全部被毀,美國經濟大受打擊.



\newpage

\section{為福音而活的人生}
\label{sec:125}
\textbf{伍渭文牧師 - 為福音而活的人生}
\newline
\newline
\hyperref[sec:124]{< < < PREV SERMON < < <}
~
\hyperref[sec:toc]{[返目錄]}
~
\hyperref[sec:126]{> > > NEXT SERMON > > >}


日期:2011年11月20日

題目:為福音而活的人生

證道:伍渭文牧師

當喬布斯離世後,很多蘋果迷悼念他,人們不斷談論他,為他英年早逝而婉惜.當然有人認為他對待員工過分嚴苛,脾氣暴躁,尖酸刻薄的態度不敢恭維,對他少作慈善捐獻也有微言,更有人批評他為了經濟利益,無視蘋果產品生產商對工人的壓榨；但無可否認,他的產品改變了世界.

喬布斯認為人生是一個旅程,能夠踏上這旅程本身就是一個報酬；在這旅途中,人尋覓他在世界的角色:為何我來到這世界,在這世界我要辦演甚麼角色哩?他認為他的角色,就是藉通訊科技改變世界,藉蘋果產品改變世界.

前美國總統克林頓接受時代雜誌視頻訪問時說,從他身上我們學到領袖的特質:當然他不可能作一個政治領袖,因他沒有制衡,在公司沒有人對他說不.但他認識自己的恩賜gift,培育它,專心一志,別無旁鶩,好好使用它.

喬布斯令我想起保羅,他也有非常強烈的召命感,也非常專心去完成這召命.但保羅與喬布斯不同的是,保羅教我們如何改變自己,因為改變世界先從自己開始,他本身也經歷了生命的改變,這改變的力量來自福音.如果說喬布斯為蘋果產品而活,保羅是為福音而活.One  lives for apple the other for gospel.從今天書信的經課帖前二章,我們看到甚麼是一個為福音而活的人生.

歷史背景:

我們首先看看帖撒羅尼迦教會和保羅的關係.這教會是保羅在第二次往歐洲傳道時所建立,使徒行傳十七4, 5:「他們中間聽了勸,就附和了保羅和西拉,並有許多虔敬的希臘人,尊貴婦女也不少.但那些不信的猶太人心裡嫉妒,招聚了些市井匪類,搭伙成群,聳動合城的人,闖進耶孫的家,要將保羅,西拉帶到百姓那裡.」

保羅迫於形勢,逃到庇哩亞繼續傳神的道,但猶太人也追到庇哩亞,聳動攪擾眾人 (v.13) ,並散播謠言,中傷保羅,說保羅像當時流行的江湖術士,憑著三寸不爛之舌,存心詭詐,諂媚人心,所說全是錯歪理.(二3) 這些初信的信徒根基不深,認為耶穌很快再回來,就放下工作,在橄欖山等候主再來.他們又認為若主再來之前死了,就沒有盼望,所以保羅說:「

若有人不肯工作,就不可吃飯.」(帖後二10) 而且他們正受到迫害(三3),保羅恐怕他們誤信謠言,動搖信仰,除了差派提摩太勸慰他們外,他更寫了這信,剖白自己作為傳道者的內心世界.這信真摯感人,使我們窺探到一個被福音改變的人,如何獻身專心一致為福音而活.在這段經文,我們曉得 保羅深信福音有改變生命的大能 保羅深信福音有永恆的價值 保羅有強烈的受托使命感

1. 1.      福音有改變生命的大能

「

我們也不住地感謝神,因你們聽見我們所傳神的道,就領受了；不以為是人的道,乃以  為是神的道.這道實在是神的,並且運行在你們心中.」(帖前二13)

在帖前二章這段經文中,最常出現的字是福音,神的道.「

在大爭戰中把神的福音傳給你們」(二2) ；「

我們既是這樣愛你們,不但願意將神的福音給你們,連自己的生命也願意給你們,因你們是我們所疼愛的.」(二8) 「

你們紀念我們的辛苦勞勞碌,晝夜作工,傳神的福音給你們,免得叫你們一人受累.」(二9) 第一章五至七節:「

我們的福音傳到你們那裡,不獨在乎言語,也在乎權能和聖靈,並充足的信心......你們在大能中蒙了聖靈所賜的喜樂,領受真道......因為主的道從你們已經那裡出來.」

保羅在開始一,二章重複又重複福音,神的道是有原因的.保羅信主前迫害基督徒,司提反殉道時,猶太人把衣服放在保羅(掃羅) 跟前,用石頭擲向司提反:「

作見證的人把衣裳放在一個少年人名叫掃羅的腳前.」(使七58) 他是當中的領袖,迫害信徒並不手軟.他不單在耶路撒冷如此,還拿著大祭司簽署的文書,直奔大馬色(大馬士革) 搜捕基督徒,信徒聽到掃羅的名字都懼怕.保羅享有領導地位,因為他是著名拉比加瑪列的大第子,被猶太人視為明日之星,復興猶太教就是他.他熱心信仰,以維持猶太教正統信仰為己任,他迫害基督徒因為他認為基督信仰是異端.他自詡「

是希伯來人的希伯來人.就律法說,我是法利賽人,就熱心說,我是逼迫教會的,就律法說,我是無可指的.」(腓三5, 6)

但信主後,整個猶太教的建制,從支持他反過來迫害他,他被視為數典忘祖,成為被猶太人迫害的對象.現在保羅一無所靠,祇靠所傳播的福音.他看見福音有能力,使沉溺邪淫拜偶像的的帖撒羅尼迦人,擺脫多年的綑縛,迎向基督.帖撒羅尼迦滿城廟宇,其中充斥男女祭司,與敬拜偶像的人,藉行淫進行宗教禮儀.色欲的歡愉,或偶像的迷信本身已是很難放棄,當邪淫被賦與宗教意義時,沉溺其中就更難擺脫.但福音有能力使他們離棄偶像,過聖潔的生活.保羅知道他們信了主後,在這方面仍有很大的試探,就勸勉他們:「

成為聖潔,遠避淫行,要你們各人曉得怎樣用聖潔,尊貴守著自己的身體,不放縱私欲的邪情,像那不認識神的外邦人.」(帖前四3, 4)

福音不單改變了他們,更改變了保羅自己.以前的保羅搜捕信徒為己任,直視司提反被石頭擲死,面不改容；拿著大祭司的文書,直奔大馬士革,理直氣壯.但認識基督後,變得像母親乳養孩子般溫柔.「

我們作基督使徒的,雖然可以叫人尊重,卻沒有向你們或別人求榮耀,祇在你們中間存心溫柔,如同母親乳養自己的孩子,我們這樣愛你們,不但願意將神的福音給你們,連自己的生命也願意給你們,因你們是我們所疼愛的.」(二6-8) 另一方面,保羅也為著他們的好處,勸勉他們,囑咐他們,像父親對兒女的管教,毫不避諱,「

要叫你們行事對得起那召你們進他國,得榮耀的神.」(二12)

保羅親自經歷福音的能力,福音改變了他的生命,也改變了他和信徒的關係.「

弟兄們,我們暫時離別,是面目離別,心裡卻不離別我們極力地想法子見你們的面.」(二17) 福音改變了我們彼此的關係,我們稱為弟兄姊妹.一人得榮耀,全體歡喜,一人有難處,彼此關心.耶穌說,你們要彼此相愛,就使人認出你們是我的門徒.

Apple 真的改變了世界,當香港年青人曉得將獲政府派發六千元時,很多人說他們要買iphone, ipad.但gospel改變人心,最怕被遺返中國的賴昌星,因信了福音,平靜接受被遺返中國.前天新聞報導,十二年斬殺報販霞姐的兩位兇手,去年六月落網,其中一位接受福音後,感到難過,要供出幕後主謀.福音使坐牢的感到不安,又使懼怕坐牢的感到平安.

1. 2.      福音有永恆的價值

「

我們的盼望和喜樂並所誇的冠冕是甚麼呢?豈不是我們主耶穌再來的時候,你們在他面前站立得住嗎?因為你們是我們的榮耀,我們的喜樂.」(二19)

傳記作者Walter Isaacson與喬布斯作了超過四十次訪談,最後一次在他家中的後花園進行,是陽光普照的下午,當時喬布斯身體很不舒服,喬布斯把話題轉到宗教和死亡.當問及他是否相信死後生命afterlife,是否相信有上帝時,他說:「五十五十.這是大奧秘,人永遠不知道,但我希望死不是生命的終結,我但願累積的智慧不會隨死亡歸於無有,可以永續.但他跟著說:但可能它真的像開關制,click一下,你就離開了；停頓一秒後他笑說:所以我們蘋果產品沒有開關制的.」

這是一個為apple而活的人生,如何看待生命和死亡.

現在讓我們看一個為gospel 而活的人生,如何看待生命和死亡.

「

我們的盼望和喜樂並所誇的冠冕是甚麼呢?豈不是我們主耶穌再來的時候,你們在他面前站立得住嗎?因為你們是我們的榮耀,我們的喜樂.」(二19)

保羅相信死亡不是終末,基督再來才是終末.他是充滿盼望的.「

論到睡了的人,我們不願意弟兄們不和道,恐怕你們憂傷,像那些沒有指望的人一樣,我們若信耶穌死而復活了,那已經在耶穌裡睡了的人,神也必將他們與耶穌一同帶來.」(四13)

基督再來parousia一字,就像當時希臘的君主駕臨自己的王國,當地的人民為著迎接君王,把最好的珠寶黃金,造成一個冠冕,呈獻給君王.保羅說,你們就是耶穌基督回來時,我呈獻給他的冠冕.我以你們為是我的榮耀和喜樂.保羅對自己累積的智慧已看作糞土,管它永續不永續,但信徒的生命是他最關心的,因為他是受上帝所托傳揚福音,這福音的果子,將成為基督頭上的冠冕.

美國總統尼克遜,因竊聽民主党總部水門事件被國會彈劾,重要幕僚多人下獄,著名的首推被稱為總統刀斧手高信(Charles Colson),他在獄中信了主,出獄後成立並獻身監獄事工(Prison Ministry),其回憶錄「重生」(Born Again) 有一段描述他重返熟悉的總統辦公室(Oval Office) 時,感觸萬千.他說:政權可以改變 (尼克遜已下台了),朝代有更替,文明有興衰,祇有一個人的靈魂,才有永恆的價值.

1. 3.      福音是上帝給我們的托付

「

但神既然驗中了我們,把福音托付我們,我們就照樣講,不是要討人喜歡,乃是要討那察驗我們心的神喜歡.」(二4)

保羅對福音的熱誠,是因為感戴上帝的愛而有的回應.掃羅在大馬士革被天上的光擊倒,有聲音對他說:掃羅,掃羅,你為甚麼逼迫我?他說:主啊,你是誰?主說:我就是你所逼迫的耶穌.究竟保羅看見甚麼呢?為何他說,主啊你是誰呢?難道他看到上主榮耀的臉相,還不知道是榮耀的主嗎?

有聖經學者解釋掃羅之困惑,是因他見的是到一個在十架上,臉龐因苦難而扭曲的基督.掃羅迫害基督徒,因為他不能接受基督列在罪犯之中,死在十字架上,上帝的兒子應該帶有榮耀的從天而降,以權能更新萬有.所以掃羅見到受苦的基督的臉相時,他不明白；他問,主啊,你是誰?他稱他為主,但不明白為何主要受苦.

主說:我就是你所逼迫的耶穌.他立時就沉默了,三天之久,掃羅不能看見,也不吃,也不喝.太震撼了,基督竟然為我而死,時間好像停止了,他被上帝的愛深深感動,生命改變了,以前看為至寶了,現在視為糞土,為的是餘生成就上帝交付他完成的工作.

有人改變世界,藉此證明他是與別不同,他是特別的.喬布斯傳記述六歲的時候,他的同年鄰居曉得他是收養回來的便告訴他:「你父母遺棄你,不想要你.」立時小喬布斯跑回家問過究竟,養父母對他說:「你是特別的,我們選上你.」這事令喬布斯感到自己與眾不同.我在想,會不會喬布斯一生為了證明自己與別不同,他追求完美,追求卓越,甚至諸多挑剔,辱罵同事平庸不達標,使同事敬而遠之呢?

保羅因為基督的愛,使他感到自己的價值,以前的掃羅是為一個傳統-猶太人信仰傳統而活,但現在為福音而活.福音是甚麼呢?福音就是耶穌基督,上帝的兒子成為人,為我們釘在十字架上,為我們贖罪.保羅說:「

我已經與基督同釘十字架,現在活著的不再是我,乃是基督在我裡面活著；並且我如今在肉身活著,是因信神的兒子而活,他是愛我,為我捨己.」(加二20)

結論:各位,今天你為甚麼而活呢?你為何為這些而活呢?讓我們為福音而活.因為「

他是愛我,為我捨己.」(加二20) .因為「

他是愛我,為我捨己.」(加二20)



\newpage

\section{看不見耶穌的時候}
\label{sec:126}
\textbf{林豪恩先生 - 看不見耶穌的時候}
\newline
\newline
\hyperref[sec:125]{< < < PREV SERMON < < <}
~
\hyperref[sec:toc]{[返目錄]}
~
\hyperref[sec:127]{> > > NEXT SERMON > > >}


馬太福音 25:31-46
\newline
\begin{longtable}{cl}
\hline
\hline
章節 & 經文 (和合本修訂版)\\
\hline
25:31 & \begin{tabularx}{0.7\textwidth}{X} 「當人子在他榮耀裡，同著眾天使來臨的時候，要坐在他榮耀的寶座上。 \end{tabularx} \\ \\ \relax
25:32 & \begin{tabularx}{0.7\textwidth}{X} 萬民都要聚集在他面前。他要把他們分別出來，好像牧人分別綿羊、山羊一般， \end{tabularx} \\ \\ \relax
25:33 & \begin{tabularx}{0.7\textwidth}{X} 把綿羊安置在右邊，山羊在左邊。 \end{tabularx} \\ \\ \relax
25:34 & \begin{tabularx}{0.7\textwidth}{X} 於是王要向他右邊的說：『你們這蒙我父賜福的，可來承受那創世以來為你們所預備的國。 \end{tabularx} \\ \\ \relax
25:35 & \begin{tabularx}{0.7\textwidth}{X} 因為我餓了，你們給我吃；渴了，你們給我喝；我流浪在外，你們留我住； \end{tabularx} \\ \\ \relax
25:36 & \begin{tabularx}{0.7\textwidth}{X} 我赤身露體，你們給我穿；我病了，你們看顧我；我在監獄裡，你們來看我。』 \end{tabularx} \\ \\ \relax
25:37 & \begin{tabularx}{0.7\textwidth}{X} 義人就回答：『主啊，我們甚麼時候見你餓了，給你吃；渴了，給你喝？ \end{tabularx} \\ \\ \relax
25:38 & \begin{tabularx}{0.7\textwidth}{X} 甚麼時候見你流浪在外，留你住；或是赤身露體，給你穿？ \end{tabularx} \\ \\ \relax
25:39 & \begin{tabularx}{0.7\textwidth}{X} 又甚麼時候見你病了，或是在監獄裡，來看你呢？』 \end{tabularx} \\ \\ \relax
25:40 & \begin{tabularx}{0.7\textwidth}{X} 王回答他們說：『我實在告訴你們，這些事你們做在我弟兄中一個最小的身上，就是做在我身上了。』 \end{tabularx} \\ \\ \relax
25:41 & \begin{tabularx}{0.7\textwidth}{X} 「王又要向那左邊的說：『你們這被詛咒的人，離開我！進入那為魔鬼和他的使者所預備的永火裡去！ \end{tabularx} \\ \\ \relax
25:42 & \begin{tabularx}{0.7\textwidth}{X} 因為我餓了，你們沒有給我吃；渴了，你們沒有給我喝； \end{tabularx} \\ \\ \relax
25:43 & \begin{tabularx}{0.7\textwidth}{X} 我流浪在外，你們沒有留我住；我赤身露體，你們沒有給我穿；我病了，我在監獄裡，你們沒有來看顧我。』 \end{tabularx} \\ \\ \relax
25:44 & \begin{tabularx}{0.7\textwidth}{X} 他們也要回答：『主啊，我們甚麼時候見你餓了，或渴了，或流浪在外，或赤身露體，或病了，或在監獄裡，沒有伺候你呢？』 \end{tabularx} \\ \\ \relax
25:45 & \begin{tabularx}{0.7\textwidth}{X} 王要回答：『我實在告訴你們，這些事你們沒有做在任何一個最小的弟兄身上，就是沒有做在我身上了。』 \end{tabularx} \\ \\ \relax
25:46 & \begin{tabularx}{0.7\textwidth}{X} 這些人要往永刑裡去；那些義人要往永生裡去。」 \end{tabularx} \\ \\
[1ex]
\hline
\hline
\end{longtable}
日期:2011年11月27日

題目:看不見耶穌的時候

經文:馬太福音25:31-46

證道:林豪恩先生

一．    經文及背景

25 : 31當人子在他榮耀裡,同著眾天使降臨的時候,要坐在他榮耀的寶座上. 25 : 32萬民都要聚集在他面前.他要把他們分別出來,好像牧羊的分別綿羊山羊一般, 25 : 33把綿羊安置在右邊,山羊在左邊. 25 : 34於是王要向那右邊的說:『你們這蒙我父賜福的,可來承受那創世以來為你們所預備的國； 25 : 35因為我餓了,你們給我吃,渴了,你們給我喝；我作客旅,你們留我住； 25 : 36我赤身露體,你們給我穿；我病了,你們看顧我；我在監裡,你們來看我.』 25 : 37義人就回答說:『主啊,我們甚麼時候見你餓了,給你吃,渴了,給你喝? 25 : 38甚麼時候見你作客旅,留你住,或是赤身露體,給你穿? 25 : 39又甚麼時候見你病了,或是在監裡,來看你呢?』 25 : 40王要回答說:『我實在告訴你們,這些事你們既做在我這弟兄中一個最小的身上,就是做在我身上了.』 25 : 41王又要向那左邊的說:『你們這被咒詛的人,離開我!進入那為魔鬼和他的使者所預備的永火裡去! 25 : 42因為我餓了,你們不給我吃,渴了,你們不給我喝； 25 : 43我作客旅,你們不留我住；我赤身露體,你們不給我穿；我病了,我在監裡,你們不來看顧我.』 25 : 44他們也要回答說:『主啊,我們甚麼時候見你餓了,或渴了,或作客旅,或赤身露體,或病了,或在監裡,不伺候你呢?』 25 : 45王要回答說:『我實在告訴你們,這些事你們既不做在我這弟兄中一個最小的身上,就是不做在我身上了.』 25 : 46這些人要往永刑裡去；那些義人要往永生裡去.

耶穌這段說話,是回應門徒在24章3節發出的提問.「

耶穌在橄欖山上坐著,門徒暗暗地來說:「請告訴我們,甚麼時候有這些事?你降臨和世界的末了有甚麼預兆呢?」」

初期的信徒對於耶穌再來很可能比我們切身和敏感得多.使徒行傳1章如此記錄:「他們(門徒)正看的時候,他(耶穌)就被取上升,有一朵雲彩把他接去,便看不見他了.當他往上去,他們定晴望天的時候,忽然有兩個人身穿白衣,站在旁邊,說:「

加利利人哪,你們為甚麼站著望天呢?這離開你們被接升天的耶穌,你們看他怎樣往天上去,他還要怎樣來.」(1:9-11)門徒知道及經歷了耶穌的離開,也得到保證,以及深信不疑耶穌必定回來.耶穌再次降臨對初期信徒的現實生活的影響,現代的信徒是難以想像得到的.耶穌會否再來,以及何時再來,對我們來說是離身又遙遠的問題,我們都是「學照上,工照返」,或者「樓照供,股照炒」.然而,對於初期的信徒,這卻是一個非常逼切而又重要的關注.

使徒行傳2:44-46節如此描述初期信徒的生活狀況:「

信的人都在一處,凡物公用,並且賣了田產,照各人所需用的分給各人.他們天天同心合意恆切地在殿裡,且在家中擘餅,存著歡喜,誠實的心用飯」.當時的信徒相信耶穌必定會回來,而且很快回來,這種盼望改變了他們的生活方式.

同時,耶穌回來與否,以及甚麼時候回來,甚至與當時信徒的性命息息相關.由門徒在耶路撒冷傳道開始,這個初期信徒群體從沒停止被兩面夾攻,包括從猶太宗教及社會而來的攻擊,以及從政治力量而來的壓逼.從新約聖經的最後一卷書──啟示錄對當時的信徒要「至死忠心」(2:10)的勸勉來看,可見他們所受的逼害達到危及生命的程度.啟示錄的結束語表達了這群被逼迫的信徒對耶穌再來的渴望:「

主耶穌啊,我願你來!」(22:20)

耶穌的確保證會回來,但卻沒有說明甚麼時候回來.當時受壓逼的信徒,熱切盼望耶穌再次降臨.然而,日子一天一天過去,耶穌仍未回來.這些等候耶穌回來的日子應該如何過呢?等候耶穌回來的日子應該如何生活?這些問題對初期信徒來說是非常切身的,甚至也令他們感到焦慮和疑惑.馬太福音24-25章記錄了耶穌關於他的再臨的講論,並且教導信徒如何過這些等候他再來的日子.今天的福音經課──25章31-46節,就是關於這主題的教導的最後一段.

二．    釋經與應用

1．    看不見耶穌的時候不應該做甚麼耶穌往天上去了,耶穌還未回來,這些看不見耶穌的日子怎樣過呢?看不見耶穌的時候,應該做甚麼呢?耶穌在馬太福音24-25章的講論暗示,看不見耶穌的時便,人們面對的試探往往是估計耶穌甚麼時候回來.然而,在看不見耶穌的日子,花時間和精力去估計耶穌回來的日子,並不是耶穌所期望跟隨他的人所做的事情.

耶穌在24章45-51節講述了一個估計耶穌遲回來的故事.這個人估計:「

我的主必來得遲」(48),於是「

動手打他的同伴,又和酒醉的人一同吃喝」(49),結果在這個人想不到的時候,主人回來了,這人被主人重重地處罰.按這個故事的暗示,估計耶穌遲回來的人可能會放縱自己,為所欲為,以為自己還有很多時間及機會收捨殘局,在耶穌回來之前「化好個妝」.故事中這位估計耶穌遲來的朋友的心態,會否也反映著現實中某些人的生活態度呢?

媽媽對孩子們說:「我去街市買菜,你們留在家做功課,不可以看電視.」媽媽走後,一位孩子說:「媽媽走了,開電視看了媽媽也不會知道.」另一位孩子擔心地說:「如果媽媽回來看見我們看電視就會責罵我們了.」又有一位孩子為他解憂:「不用怕,媽媽沒有那麼快回來.如果擔心的話,我們可以派一個人貼者大門站崗,聽到媽媽回來的腳步聲就立即關電視.」於是,他們就既喜又驚地圍著電視來看,連那位負責站崗的孩子,也只聽到電視的聽音,而聽不到媽媽回來的腳步聲.當媽媽推門進來時,站崗的孩子才驚慌地大叫:「媽媽回來了!」其他孩子立即關電視,裝作做功課的樣子.然而,無論孩子們的動作如何敏捷,都瞞不過聰明的媽媽.當大家都鴉雀無聲的時候,媽媽走返電視機,伸手一摸,熱得還燙手的.於是媽媽對孩子們說:「看來我真的在街市花了不了時間,讓你們看了這麼久電視!」

估早了又怎樣呢?美國有一位傳教士名叫哈洛德康平(Harold Camping),喜歡預測最後審判(judgment day)的日子,最近的預測是今年五月二十一日.有些信徒非常嚴肅的看待他的預測.有些人辭掉了工作,專心等候.有一個人名叫羅伯特菲茨帕特里克(Robert Fitzpatrick ),為了給「最後審判」的大量廣告活動提供資金,花掉了畢生的積蓄.2011年5月21日下午6時,這是康平預言將會開始「最後審判」的時間,菲茨帕特里克站在紐約市的時代廣場,滿懷期待地等待著.當「最後審判」沒有出現的時候,菲茨帕特里克驚呆了.他說:「我不明白,為什麼什麼也沒有發生.這不是一個錯誤.我做了我應該做的.我做了聖經上說的.」康平並未就此罷休,而是根據聖經創世記所做的一個新分析,堅稱真正的天地終結日,事實上應為十月二十一日.然而,今天已經是十一月二十日了.

最近,某八卦雜誌報導,本港某教會牧師預測有顆彗星會撞落太平洋,繼而引發海嘯撲向香港.數以百計信徒收執包袱帶齊糧水,拖男帶女兼領狗逃到大帽山,飛鵝山甚至廣州避難.該雜誌記者直擊一批信徒避走大帽山的過程,並寫道:「眾人呆等幾小時後,不但彗星沒出現,連一隻烏鴉亦無飛過.」

估早了耶穌回來的人,以為自己虔誠,卻把自己變成一個笑話,花費了上主所賜寶貴而有限的資源在沒有意義的事情上,甚至也因此而誤導了其他人對基督教信仰的觀感.

2．    看不見耶穌的時候應該做甚麼看不見耶穌的時候,不要花時間精力去估計他甚麼日子回來,那麼,應該做甚麼事情呢?耶穌在25章31-46節其實說得相當清楚.

「

你們這蒙我父賜福的,可來承受那創世以來為你們預備的國.」(25:34)這些是人做了甚麼事情呢?是估中了耶穌回來的日子嗎?是最快走到上大帽山上嗎?答案在耶穌接著說的話中找到:「

因為我餓了,你們給我吃,渴了,你們給我喝；我作客旅,你們留我住；我赤身露體,你們給我穿；我病了,你們看顧我；我在監裡,你們來看我.」(25:35-36)

另外一些人為甚麼接到以下的判決:「

你們這被咒詛的人,離開我!進入那為魔鬼和他的使者所預備的永火裡去!」(25:41)那又為甚麼呢?耶穌指責他們說:「

因為我餓了,你們不給我吃,渴了,你們不給我喝；我作客旅,你們不留我住；我赤身露體,你們不給我穿；我病了,我在監裡,你們不來看顧我.」(25:42-43)

耶穌舉例列出了六項人生的狀況,前四項代表著人類生存的基本需要,就是我們常說的「衣食住行」的「衣食住」；另外兩項──患病及下監,代表著人生艱難無力的時刻.有些人回應了耶穌這些需要,有些人拒絕了耶穌這些需要.奇怪的地方是,無論是回應了耶穌的需要的人,還是拒絕了耶穌的需要的人,他們都不知道自己所做的已經回應了耶穌的需要,也不知道自己所做的其實拒絕了耶穌的需要,因為他們都異口同聲地反問:「耶穌,我們甚麼時候見過你呢?」他們根本沒有見到耶穌,他們怎可能在耶穌有需要的時候回應了或者拒絕了他的需要呢?耶穌的回答肯定了他們的認知,他們的確沒有見到耶穌肚餓,口渴,流浪,赤身露體,患病,坐監,甚至他們根本沒有見過耶穌,關鍵在於「

這些事你們既做在我這弟兄中一個最小的身上,就是做在我身上了」(25:40),「

這些事你們既不做在我這弟兄中一個最小的身上,就是不做在我身上了」(25:45).

如果你看見耶穌肚餓,你當然會給他食物.然而,除了當年跟隨耶穌出入的門徒,其他人根本沒有機會直接接觸耶穌.如故事中的兩批人一樣,我們都沒有看見耶穌,看不見耶穌,是我們絕大多數人的共同經驗,但是,耶穌卻非常重視我們在看不見他的日子所做的事情.如何能夠避免Out of sight變成out of mind呢?或者,我們可以透過德蘭修女的祈禱得到啟發.

耶穌,我受苦的主,願你賜我恩典,讓我今天,每一天都透過這些久病的人看見你,在我服事他們的時候,讓我能服事到你.就算遇到那些要求苛刻,不合理的人,求你也讓我認出你,然後能夠照樣說:我受苦的耶穌阿,能服事你,是多麼甘美呀.

德蘭修女在服務貧病人士的時候看見耶穌,並視之為對耶穌的服務.在我們身處的世界中,其實充滿著服務耶穌的機會.我們身處的世界,現在有多少人飢餓呢?最近,聯合國糧農組織報喜,世界飢餓人群的數量今年將減少9.6\%,即是有多少人飢餓呢?是9億2千5百萬人.按著耶穌提出的項目檢視:我們的世界現在還有多少人連清潔衛生的飲用水也欠缺呢?有多少人衣不蔽體呢?有多少人流離失所呢?有多少人受到病患的困擾呢?有多少人因種種原因而被隔離於家人和朋友呢?

這不是說人人都要去做修女,也不等於所有人都去做義工.反過來說,如果要做修女,或者做義工,才能夠使餓者得飽,渴者得水,飄泊者得安定,寒者得衣,病者得照顧,被隔離者得連繫,即是說我們的日常工作及生活與此完全沒有關係.如果人類的工作,經濟活動,國際交往,人際關係都不能使餓者得飽,不能使渴者得水,不能使飄泊者得安定,不能使寒者得衣,不能使病者得照顧,不能使被隔離者得連繫,有一天,耶穌問我們:你們興建那麼多工廠所為所事呢?你們發明那麼先進的現代科技所為所事呢?你們成立那麼多銀行所為所事呢?你們花十幾二十年時間去讀書所為何事呢?你們的領袖飛來飛去開這個會議開那個會議又所為何事呢?你們成立各種國際及地區組織又所為何事呢?

有一位孟加拉人名叫穆罕默德尤努斯,他的確用了很長時間去讀書,他又花了很多精力去成立一間銀行,他成為了2006年諾貝爾和平獎的得主,因為他於1976年創辦鄉村銀行組織,提供小額信貸,在30年之間至少貸出了50億美金,幫助500萬的孟加拉人脫離貧困.諾貝爾評審委員會出:「只有當大批人口找到消除貧困的途徑,才能取得永久的和平.」

逢星期日黃昏,都從收音機聽到一個點唱節目.最初聽的時候,只把它當作是一個普通的點唱節目,但漸漸卻發覺點唱的內容有點奇怪:「老婆,我今日第一次搭馬鐵轉車出尖沙咀,好快呀,等你出來我再帶你搭!」「我好掛住你呀,唔經唔覺仲有半年你就可以出來,到時我們就可以結婚喇,放心啦,我一定會等你!」「爸爸,你幾時返黎呀?我肥jor 呀,你快D 返黎抱我啦!I love you!」原來這是一個讀出在囚人士及其家人朋友的來信,並接聽囚犯家人的電話的電台節目.我越聽下去,越覺得這個節目很有意義,也很令人感動.這個節目把在囚人士和他們的家人朋友連繫起來,也讓社會人士聽到,在囚人士雖然因為犯錯而被囚,但他們也是一個需要關係,需要親情的人.這個節目已經做了十年,有人這樣形容這個節目:「每個星期日黃昏,全港的監房都會從喧鬧之中忽爾靜謐下來,能令一班各有自己故事的在囚人士全神貫注的,正是葉韻怡每一百六十六小時一次的電台節目.近十年來,從空氣傳播真心,發展到身體力行地協助更生人士,她發放了數之不盡的祝福,安慰了無數孤單困苦的心靈,讓在囚人士及其家屬的創傷得以修補.」(劉孝偉《守護牢房的天使--葉韻怡》)

有人問德蘭修女:「世界上飢餓人口如此龐大,我怎能餵飽他們呢?」德蘭修女回答:「如果你能夠餵飽一個人,去就去餵飽一個人吧!」

三．    總結

看不見耶穌的時候,我們的信仰成長追求甚麼呢?

一個無可反對的建議是:追求與耶穌契合,甚至到一個看見耶穌的境界.德蘭修女的經歷值得我們參考.

隨著修女會逐漸取得豐碩成果,德蘭修女內心的痛苦日益加深.以前與天主,耶穌的共融消失了,現在只有苦澀的枯乾.她不覺得天主臨在,不覺得與耶穌結合.她極渴望與天主合一,但卻感到很遙遠.後來,德蘭修女的生命進入一個新境界.她不再探究或質疑內心黑暗的奧秘,她接受黑暗,接受天主的一切安排,並以微笑接受一切.她學會與痛苦為友,甚至愛上痛苦.德蘭修女在一篇禱文中說:「耶穌,我接受任何你給我的,我獻上任何你所要的.」

在看不見耶穌的時候,耶穌看重的,似乎不是我們對他的追捧,而是我們對於別人的基本需要的反應.



\newpage

\section{果樹與糠秕}
\label{sec:127}
\textbf{林豪恩先生 - 果樹與糠秕}
\newline
\newline
\hyperref[sec:126]{< < < PREV SERMON < < <}
~
\hyperref[sec:toc]{[返目錄]}
~
\hyperref[sec:128]{> > > NEXT SERMON > > >}


日期:2011年12月4日

題目:果樹與糠秕

經文:詩篇1

證道:林豪恩先生

一．    經文

1 : 1    不從惡人的計謀, 不站罪人的道路, 不坐褻慢人的座位, 1 : 2    惟喜愛耶和華的律法,晝夜思想,這人便為有福!1 : 3    他要像一棵樹栽在溪水旁, 按時候結果子, 葉子也不枯乾. 凡他所做的盡都順利.  1 : 4　惡人並不是這樣, 乃像糠秕被風吹散. 1 : 5    因此,當審判的時候,惡人必站立不住； 罪人在義人的會中也是如此. 1 : 6    因為耶和華知道義人的道路； 惡人的道路卻必滅亡.

詩篇第一篇是一首「美麗與智慧並重」的詩.它描繪出美麗的景象:「

一棵樹栽在溪水旁,按時候結果子,葉子也不枯乾」；它也道出了人生的智慧:「

不從惡人的計謀,不站罪人的道路,不坐褻慢人的座位,惟喜愛耶和華的律法,晝夜思想,這人便為有福」.今天,讓我們在欣賞這首「美麗與智慧並重」的詩的時候,同時反思我們的生活態度.

二．    釋經與應用

1．實在與虛浮讓我們想像兩種景象,第一種是「

一棵樹栽在溪水旁,按時候結果子,葉子也不枯乾」,第二種是「

糠秕被風吹散」.「

一棵樹栽在溪水旁,按時候結果子,葉子也不枯乾」的景象,給人一種實在的感覺,我們可以環抱樹幹,可以嗅嗅果子的味道,可以摸摸樹葉的紋理；相比之下,「

糠秕被風吹散」給人一種虛浮的感覺,沒有可環抱的,沒有可品嚐的,也沒有可觸摸的,當下就隨風飄散.詩篇第一篇以這兩種景象來比喻兩種人生,人們很容易就意識到,「

一棵樹栽在溪水旁,按時候結果子,葉子也不枯乾」象徵實在的人生,「

糠秕被風吹散」代表虛浮的人生.在比喻的世界,實在與虛浮對比強烈,人們輕易地作出選擇,取「栽在溪水旁的樹」而捨「被風吹散的糠秕」.

回到現實的世界,虛與實卻不再是那麼分明.有些人把實與虛玩弄得似是而非,實則虛之,虛則實之,令人無從判斷.有退休人士,拿著退休金去銀行請教:如何能夠穩穩陣陣地既保本又有利息收呢?金融衍生產品過於虛浮,退休人士不宜沾手,債券就是穩健實在之選,既保證本金安全,又有利息收.後來,他們痛苦地認識到,在金融世界,以為實在的,原來也可以是很虛浮的.以他們的遭遇讀詩篇第一篇,可能有這樣一個版本.

不從 I Band 的計謀不站 accumulator 的道路不坐「孖展」的座位惟喜愛實在多於虛浮晝夜挑選這人便為安全他要把積蓄放在雷曼債券上按時候收利息本金也不枯乾凡他所要的基本開支盡都供給現實並不是這樣乃像肥皂泡被風吹散因此,當爆破的時候,本金也站立不住保本與不保本也是如此因為財技魔術師知道虛實互換的道路冒險與穩健也都滅亡.

人們除了把虛實弄得似是而非,令人難以分辨之外,更進一步指虛為實,貶實為虛.詩篇指出,好像「栽在溪水旁的樹」一樣實在的人,他們「喜愛耶和華的律法,晝夜思想」,而好像「糠秕被風吹散」般虛浮的人,則「從惡人的計謀,站罪人的道路,坐褻慢人的座位」.矛盾就在這裡出現.人們同意「栽在溪水旁的樹」為實在,但在現實中卻往往視「耶和華的律法」為不實在；人們同意「被風吹散的糠秕」為虛浮,但在現實中卻認為「惡人的計謀」更實在,因為「惡人的計謀」比「耶和華的律法」往往令人更有效地達到目的.例如耶和華的律法要求人們要忠誠和正直地生活,人們的回應則是「忠忠直直,終需乞食」.有些人認為忠誠與正直是捉不實,抓不住的,填飽肚腹的感覺更是實在.按這樣的邏輯,誠信實在,還是入讀名校實在呢?誠信是看不見的,摸不到的,但名校是可以看見的,可以坐幾年時間的.因此,怎可能為了虛無飄渺的誠信而犧牲入讀名校的機會呢?相同的邏輯,發展是硬道理,價值是點綴品；經濟增長必須保證,品格培育談談就算；搵錢最實際,仗義相助請往武俠小說裡尋.當一個社會視價值,品格,仗義相助為虛無縹緲得可有可無,同時視發展,經濟增長,搵錢為最實在的時候,審判就臨到這個社會.當人們在媒體看著那位二歲的小女童被車輛輾過兩次,十八位路人視而不見,匆忙走過而沒有停下來施以援手的影像,審判不單臨到那剛好被攝入鏡頭的那十八位路人,同時也臨到整個社會.

但願這些使人不寒而慄的不幸影像,喚醒人們不要總是把能掌握在手中的才視為實在,而把其他的都貶為虛幻.詩篇第一篇指出,耶和華的律法雖然手不能握,但卻是實在的,因為這是創造主為被創造者設下的生存之道,是人類社會能夠按創造的美意而存在的基礎.上主說:「

天地要廢去,我的話卻不能廢去」(馬可福音13:31).甚麼是實在,甚麼是虛幻,我們是時候要反思一下.

2．時序與即時讓我們再想像「栽在溪水旁的樹」這幅圖畫.詩篇描述「栽在溪水旁的樹」的其中一個特徵是「按時候結果子」.「按時候結果子」對古代人來說,是美好及合理的期待.美好,因為果樹克服了風雨及蟲害的挑戰,終於結成果子；合理,因為果樹必須按著自然定律,在到某個生長週期才會結果子.然而,現代人未必認同「按時候結果子」是美好的和合理的,反而感到這是令人難以忍耐的狀況.

從前,人們依循著四季運轉的時序而生活.例如,我們在夏天吃西瓜,在冬天吃火鍋.現在,無論甚麼季節,想吃西瓜就有西瓜吃,無論甚麼天氣,想吃火鍋就有火鍋吃.「按時候結果子」表達一種有時序的生存狀況,人類配合著四季運轉,配合農作物生長期,配合家禽家畜成長日子而調節自己的生活起居和飲食習慣.現在,人們要求即是滿足欲望,不想等候,不想愛到大自然和時序的限制.那麼,人們會怎樣做呢?例如,我們不想等到夏天才有西瓜吃,可以怎樣做呢?

人民網的記者有這樣的觀察和報導:「綠油油的西瓜,黃澄澄的鴨梨,紫珍珠似的葡萄,紅艷艷的荔枝．．．冬季里各色反季節水果飄散著陣陣或濃烈或淡淡的果香,源源不斷地走進千家萬戶.可就在人們品嘗美味的同時,也不斷傳出一些不和諧音:濟南,西安相繼有三四歲男孩女孩因吃早熟的"激素水果"過多而出現男孩長胡子,女孩乳房鼓的性早熟症狀．．．目前植物生長調節劑在蔬菜,水果種植業已較多應用,一株果樹從幼苗至成熟可以使用十几種激素．．．使用激素后可增產20％左右.(陳海:「注射催熟劑膨大劑防腐劑,激素水果危害大」,《人民網》)

人們忍耐不了「按時候結果子」的等候,想要吃雞就吃雞,想吃牛肉就吃牛肉,然而,雞如何能長得那麼快,去滿足人們一日三餐雞排呢?牛如何能長得那快多,去滿足人們一日五餐牛排呢?要突破時序的限制,要消除等候的難耐,人類總有辦法,把生長激素打進家禽家畜體內,人們便能夠以最快的速度得到即時的滿足.

「按時候結果子」是一種生活態度,這種生活態度要求人們與大自然配合,要求人們與時序配合,因此要求人們等待；以最快的速度達到即時滿足,也是一種生活態度,這種生活態度好像讓人擺脫自然與時序的限制,消除了等待的難耐,不過,也不是必然美好,有時可能帶來災難的.

日本和德國花了多少年日,才能夠把火車提升至現在的行走速度呢?我們只花了三,四年時間,就迎頭趕上了.2010年12月3日,京滬高鐵「和諧號CRH-380A」在山東南部試車成功,飆出486.1公里時速,再度刷新全球鐵路營運試驗的最高速.在人們都追求更高速度,追求即時滿足的時候,有人忠言逆耳地建議我們學習等待,體會等待的情操.

中國,請停下飛奔的腳步.等一等你的人民,等一等你的靈魂,等一等你的道德,等一等你的良知!不要讓列車脫軌,不要讓橋樑坍塌,不要讓道路成陷阱,不要讓房屋成危樓.慢點走,讓每一個生命都有自由和尊嚴,每一個人都不被『時代』拋下,每一個人都能順利平安地抵達終點!

詩篇第一篇提示我們,「喜愛耶和華的律法,晝夜思想」的人,未必會得到即是滿足,速度也未必比別人快,他們得到的,只是「按時候結果子」的應許.然而,詩人同時指出告訴我們,其實,能夠「按時候結果子」,「這人便為有福」.要不要尊重和接受四季運轉的時序,動植物成長的時序,我們的兒女成長的步伐,並且以此來調整我們的生活態度?這些將會是我們越來越不能逃避的問題.

3．可持續發展與一殺那光輝讓我們又回到「栽在溪水旁的樹」,以及「糠秕被風吹散」這兩幅圖畫.栽在溪水旁的樹每日一點一滴地生長,變化緩慢,要看它的花朵,可能要等一年,要見它的果子,可能也要等一年.有時,植根於泥土的果樹,看見隨風飄揚的糠秕,在陽光照耀之下散射出的光芒,有時可能會羨慕它的瀟灑和不羈.

隨風飄揚的糠秕的瀟灑和不羈,畢竟只是殺那間的光輝,它本身沒有生命,也不可能延續生命.栽在溪水旁的樹,雖然變化緩慢,但它本身卻是實在的生命,而且透過所結的果子延續生命.就是這棵有生命,並且延續生命的樹,為不少人帶來了希望,其中一位包括《安妮的日記》的作者──安妮．法蘭克.

二戰期間,安妮和家人為了躲避納粹軍隊的搜捕,只好藏匿在狹小閣樓裡,與外間斷絕一切聯絡.每天,一家人都得忍受無止竟的黑暗恐懼,以及從收音機傳來死亡人數的噩耗.而能支撐安妮求生意志的,就是看見窗外一棵隨著四季成長變化的板栗樹.「我們望著窗外的藍天和那棵板栗樹,光禿的樹枝上還滴著水珠．．．樹木彷彿是銀色的,我們都被這幅圖畫感動得說不出話來.」兩年多的密室生活日記,透露了安妮渴望自由的心情,她始終堅信殘酷的屠殺終將結束.「板栗樹是那麼茂盛,甚至比去年更美麗．．．」(《浦公英》2011年10月)

一棵有生命的樹能夠為別人的生命帶來盼望,一棵有生命的樹即使倒下,仍然能夠延續生命.數十年後,這棵170歲的板栗樹在一次風雨中倒下了.然而,樹上所結出的150顆種子,卻被分送出去,如今在各地發芽茁壯,繼續為世界帶來希望,如同當時為安妮帶來希望一樣.

詩篇第一篇再一次邀請人們思考及選擇生活態度.「惡人的計謀」,「罪人的道路」,「傲慢人的座位」有時會令人得到隨風飄揚的快感,讓人感到從心所欲的瀟灑和不羈,有時會令人享受到一殺那的光輝.「喜愛耶和華的律法,晝夜思想」的人,可能好像植根於溪水旁的果樹一樣,沒有那麼瀟灑和不羈,而且變化緩慢,不肯仔細觀察,或者沒有耐性待候它開花結果的人,甚至以為它千年如一日.不過,這卻是能夠延續生命的生命.

曾經何時,為了滿足欲望,人類遇樹砍樹,遇山開山,遇海填海；那裡有燃油就去抽取,那裡有礦物就去發掘,那裡有市場就大量製造.人類為自己能夠戰勝大自然而沾沾自喜的時候,自豪地炫耀自己的成就的時候,突然發現,森林和草原變成了沙漠,清潔的水源變得罕有,清潔的空氣越來越差稀少,礦物能源快將用盡,氣候越來越極端．．．原來,我們的成就只是一殺那的光輝,我們的生活方式取消了很多事物的生存機會.當大自然以災難的方式告回饋我們的時候,我們應該醒覺這是對我們的生活方式的審判.

詩人說:「惟喜愛耶和華的律法,晝夜思想,這人便為有福」.耶和華的律法不單為人類共處定下原則,同時也為人類與大自然共處建立規範.創世記告訴人們,人類出現之前,天地日月的運轉已經設定,海洋陸地已經分佈,動物植物已經滿遍.我們可以用這個角度解讀詩人的意思:若上帝所創造的人類能夠與上帝所創造的大自然和平共處,這人便為有福.當我們維持大自然的生命,讓大自然生生不息的時候,我們就能夠維持人類的生命,並且延續人類的生命.

三．    總結

詩人指出,如「栽在溪水旁的樹」一樣的人,是「有福的人」,因為他是實在而有根有基的,有生命而且延續生命的,更會按時候結果子的.「栽在溪水旁的樹」結出的果子是怎樣的呢?那就視乎它是甚麼樹,橙樹長出橙,蘋果樹長出蘋果.「喜愛耶和華的律法」的人結出的果子是怎樣的呢?保羅在加拉太書有這樣的描述:「聖靈所結的果子,就是仁愛,喜樂,和平,忍耐,恩慈,信實,溫柔,節制」.願我們的心靈被詩篇第一篇的美麗所滋潤,願我們的生活被詩篇第一篇的智慧所指引.



\newpage

\section{神啊復興我們}
\label{sec:128}
\textbf{鄧瑞強博士 - 神啊,復興我們}
\newline
\newline
\hyperref[sec:127]{< < < PREV SERMON < < <}
~
\hyperref[sec:toc]{[返目錄]}
~
\hyperref[sec:129]{> > > NEXT SERMON > > >}


日期:2011年12月11日

主題:神啊,復興我們

經文:詩篇80

證道:鄧瑞強博士

各位弟兄姊妹,主內平安.

我們每週讀的經課,都有詩篇,但我們較少以詩篇為講道經文.今日算是例外,我打算以今日的詩篇經課(詩篇80篇)為講道經文.

我們再讀一遍這首詩篇:詩篇80篇

80:1　〔亞薩的詩,交與伶長.調用為證的百合花.〕領約瑟如領羊群之以色列的牧者啊,求你留心聽!坐在二基路伯上的啊,求你發出光來!80:2　在以法蓮,便雅憫,嗎哪西前面施展你的大能,來救我們.80:3　神啊,求你使我們回轉〔或譯:復興〕,使你的臉發光,我們便要得救!80:4　耶和華──萬軍之神啊,你向你百姓的禱告發怒,要到幾時呢?80:5　你以眼淚當食物給他們吃,又多量出眼淚給他們喝.80:6　你使鄰邦因我們紛爭；我們的仇敵彼此戲笑.80:7　萬軍之神啊,求你使我們回轉〔或譯:復興〕,使你的臉發光,我們便要得救!80:8　你從埃及挪出一棵葡萄樹,趕出外邦人,把這樹栽上.80:9　你在這樹根前預備了地方,它就深深扎根,爬滿了地.80:10　它的影子遮滿了山,枝子好像佳美的香柏樹.80:11　它發出枝子,長到大海,發出蔓子,延到大河.80:12　你為何拆毀這樹的籬笆,任憑一切過路的人摘取?80:13　林中出來的野豬把它糟踏；野地的走獸拿它當食物.80:14　萬軍之神啊,求你回轉!從天上垂看,眷顧這葡萄樹,80:15　保護你右手所栽的和你為自己所堅固的枝子.80:16　這樹已經被火焚燒,被刀砍伐；他們因你臉上的怒容就滅亡了.80:17　願你的手扶持你右邊的人,就是你為自己所堅固的人子.80:18　這樣,我們便不退後離開你；求你救活我們,我們就要求告你的名.80:19　耶和華─萬軍之神啊,求你使我們回轉〔或:復興〕,使你的臉發光,我們便要得救!

今日是「將臨期」的第一主日,「將臨期」的經課內容,多少帶有一種張力,即身處此刻卻盼望將來.此刻,可能是孤寂,可能是失敗,可能是艱難；所盼望的將來,是團圓,是勝利,是解救.「將臨期」代表的時間,又點似:你受傷了,卻已打了電話召喚救護車.救護車正來,但未到.而你,在受苦中,但知道再多等一會,便有人救你,你是有盼望的.但這救援卻仍未到,此刻,你仍在苦中.這夾縫時間,就是「將臨期」.

剛才讀的詩篇,是一集體哀求的禱文.他們身在危難之中,我們不能確定這危難的性質,但看來似是國破之苦.這是一首在苦中仍不失盼望的詩歌,一首很合宜地放在「將臨期」裡思想的詩歌.

我將這首詩分成四段.第一段:1-3節.這是一段緊急的哀求.

80:1　〔亞薩的詩,交與伶長.調用為證的百合花.〕領約瑟如領羊群之以色列的牧者啊,求你留心聽!坐在二基路伯上的啊,求你發出光來!80:2　在以法蓮,便雅憫,嗎哪西前面施展你的大能,來救我們.80:3　神啊,求你使我們回轉〔或譯:復興〕,使你的臉發光,我們便要得救!

「亞薩」應是詩班長.「為證的百合花」是詞牌名.

這裡用了大量意象.神是誰?他是「牧人」,他是帶領羊群的牧人.他在哪裡?他坐在二「基路伯」之上.「基路伯」是有翼的神獸,後來理解為天使.這二「基路伯」是約櫃頂上的神獸.藉此,詩人提醒自己,這位神是與自己有關的,神與人是有約的關係的.人在艱難中,不是孤身面對,而是有一位與人有關的神幫助人去面對的.神以約櫃為寶座.這約櫃放在耶路撒冷的聖殿裡,神的寶座不在高天之上,卻在人間.對身在苦難的人來說,神在人間是重要的.距離太遠,便遠水不能救近火了.你若需要急救,但救護車需要從廣州急調過來,則你便會絕望了.在生命的黑暗裡,最需要的是光明,所以詩人求神發出光來.

第2節,詩人求神「施展大能」,所謂「施展」,其實意指神像睡了覺一樣,現在叫醒他,叫他快快回復體力.Eugene Peterson 的 ``The Message" 的翻譯很傳神,他這樣翻譯這節

經文: ``Get out of bed  and \#8212; you"ve slept long enough! Come on the run before it"s too late." 這裡有點似福音書裡的一個故事情節.耶穌和門徒在船上.忽然起了暴風,船入了水快沉,但耶穌卻睡著了.門徒叫醒他,看來帶點埋怨,你看,危難當前,你還睡覺嗎?詩人在這裡表達同樣的情緒－－神啊,你是否睡了覺?你看不到我們的苦況嗎?你快快「施展大能」,快快「恢復體力」吧!

神是親近我們的,他是牧人,坐在百姓中間.但神看來也是很遠的,他像看不見我們苦難似的,他像睡了覺似的.

若神真是牧人,若神是為羊捨命的好牧人,而羊現在受苦,就可能不是牧人睡了覺,而是羊故意走入迷途了.

第1及2節,用了四個imperative的動詞.第1節:「求你留心聽」可直解為「聽啊」；「求你發出光來」可直解為「發光啊」；「施展你的大能」可直解為「起身啦」；「來救我們」的「來」可直解為「來啊」.我們感受到這些受苦的人的迫切.他們呼求神,聽啊,發光啊,起身啦,來啊.似乎在說:「救命呀」,「救命呀」,「救命呀」,「救命呀」.

現代人有問題,會第一時間想起 google,因為相信 google 會提供一切問題的答案.詩人在生命的苦難中,卻尋求一位有情,有義,有約的牧人,一位在人間的神.

第3節是「副歌」的第一次出現.

80:3　神啊,求你使我們回轉〔或譯:復興〕,使你的臉發光,我們便要得救!

「回轉」與「復興」是同一個字.走迷的人,需要神帶領他回轉；陷在生命幽谷中的人,需要神給他復興的力量.

詩人尋求神的光明,勝於尋求世上的智慧和力量.

第二段:4-7節.首段的哀求以後,便是苦況的自述.

80:4　耶和華──萬軍之神啊,你向你百姓的禱告發怒,要到幾時呢?80:5　你以眼淚當食物給他們吃,又多量出眼淚給他們喝.80:6　你使鄰邦因我們紛爭；我們的仇敵彼此戲笑.80:7　萬軍之神啊,求你使我們回轉〔或譯:復興〕,使你的臉發光,我們便要得救!

第5節,他們的食物是淚水；飲的,也是淚水.

我想起集中營裡的人,在絕望裡,在痛苦裡,吃的,飲的,也是淚水.

第6節,苦況以外,再加上旁人的嘲笑與敵意.

你病了,這是苦.但若旁人說,「看你幾時死」,則更苦.你破產了,這是苦.但若旁人說,「還不等到你有今日」,則更苦.你失敗了,這是苦.但旁人說,「還不輪到你跌一回」,這更苦.

這些苦,是鄰舍加的,是敵人加的.

主耶穌希望人能愛鄰舍,愛敵人.但這裡的鄰舍和敵人,只懂雪上加霜.

第4節,詩人發出哀聲,這一切,要到幾時呢?

最近探訪過一個住在板間房的家庭,他們也問這個問題:要到幾時呢?

他們申請了公屋,等了若干日子,但,還要等多久呢?

問這問題的人,在一苦境中.但他們仍在等,因他們仍抱一絲希望.

這真是一個「將臨期」的天問:「要到幾時呢?」

人生的苦難,幾時才會完呢?

第4節說:「耶和華──萬軍之神啊,你向你百姓的禱告發怒,要到幾時呢?」這裡的「發怒」,原來的用字是指「煙」.詩人用這意象,去表達神的怒氣如煙.神被這煙遮住了眼睛,故看不到他們的苦況.這要到幾時呢?

但是,我們也可按這意象作一點推想,煙熏目則會流淚.在人間,受苦的人認為自己日日以淚為飲食,但是,這難道不也是神的眼淚?在Mel Gibson 導演的＜受難曲＞(The Passion of the Christ)裡,當人間的基督死在十架上時,一大滴眼淚從天而降,代表著神的無盡悲哀.人間的眼淚,難道不也是神的眼淚嗎?

若人問:「要到幾時呢?」神也會問:「要到幾時,人才懂得愛鄰舍和愛敵人呢?」

第8節,副歌第二次出現.但這次,稱神為「萬軍之神」,更強調神的力量了.

80:7　萬軍之神啊,求你使我們回轉〔或譯:復興〕,使你的臉發光,我們便要得救!

第三段:8-14節.受苦的人,想當年.

80:8　你從埃及挪出一棵葡萄樹,趕出外邦人,把這樹栽上.80:9　你在這樹根前預備了地方,它就深深扎根,爬滿了地.80:10　它的影子遮滿了山,枝子好像佳美的香柏樹.80:11　它發出枝子,長到大海,發出蔓子,延到大河.80:12　你為何拆毀這樹的籬笆,任憑一切過路的人摘取?80:13　林中出來的野豬把它糟踏；野地的走獸拿它當食物.

想當年,神帶領以色列人出埃及,在迦南地茁壯成長.曾有輝煌的日子.西面,去到大海,即地中海.東面,去到大河,即幼發拉底河.即當時從東到西之處.但是,俱往矣.這個國家被人毀了.這棵好樹的籬笆被人拆毀了,任人予取予攜.一片江山,被以色列人最討厭的豬摧毀了.豬,當然就是最粗暴的外邦人了.

中國南唐的李後主,在亡國後,也寫了很多想當年的詞句.我們熟識的＜虞美人＞說:「春花秋月何時了,往事知多少」.大自然的春花秋月,是常有的,這對比著人事的無常.「小樓昨夜又東風,故國不堪回首月明中」.東風是常有的,這對比著故國生活的不常有.回首這不常有的過去,是不堪回首的.「雕欄玉砌應猶在,只是朱顏改」.過去的宮殿,或許不變吧,這對比著人的歲月的不能不變.「問君能有幾多愁,恰似一江春水向東流」.春水向東流,是永恆的,在人生中只有愁,也如此永恆.

李後主的故國回憶,只有愁,因李後主不懂「將臨期」.詩篇80篇的故國回憶,卻是喚起希望的回憶.在詩篇的詩人眼中,現在,的確是一片荒涼,故國不再,但是,他不停在這裡,他回憶,他進行救贖的回憶.過去,荒涼之前,是美好的故園.這故園,這過去的生命和生活,是神栽種的,是神培植的.這一切背後,神是牧人,他和人有親密的,約的關係.故此,詩人越回憶,他越接近神,他的心越堅定,他越明白救贖的大路.他明白,一切不是愁,一切是有希望的.神是希望之源.

第四段:14-19節.這是希望的宣告.

80:14　萬軍之神啊,求你回轉!從天上垂看,眷顧這葡萄樹,80:15　保護你右手所栽的和你為自己所堅固的枝子.80:16　這樹已經被火焚燒,被刀砍伐；他們因你臉上的怒容就滅亡了.80:17　願你的手扶持你右邊的人,就是你為自己所堅固的人子.80:18　這樣,我們便不退後離開你；求你救活我們,我們就要求告你的名.80:19　耶和華─萬軍之神啊,求你使我們回轉〔或:復興〕,使你的臉發光,我們便要得救!

第14節是副歌的一個變奏.這次,不是求神使我們回轉或復興,而是求神回轉.「神啊,你回來吧!」只要你回來,一切便會復興的了.我們曾是你手提攜的人,我們的生活是你建立的,我們的生命是你指引的.現在,我們誤入歧途,我們走入死角,我們苦不堪言.神啊,求你回來吧,繼續在我們生命裡,完成你未完成的工程.

第17節:「願你的手扶持你右邊的人,就是你為自己所堅固的人子」.「人子」即人,不必想到新約的基督.

第18節:「這樣,我們便不退後離開你；求你救活我們,我們就要求告你的名.」這是一種悔改的決志.若神回頭再看我們,我們便不能再離開神了.

在詩人看來,一切苦難,來自人背叛神.希望來自神,只要人回到神那裡,他就能懷著希望地活下去.

第19節,最後一次副歌.這次,稱神為「耶和華─萬軍之神」.「耶和華」是神的名字,是讓以色列人叫神的名字.這名字不是泛泛的一個神的名字,是尊指那位與他們有關係,與他們有約,並至終看顧他們的神的名字.但願,你在艱難中,能對神有信心,看到神和你的關係,你也能叫得出神的名字.不是泛泛的一個神的名字,而是和你生命有愛的連繫的神的名字.

結語:

我們的生命是由神栽種的,請回到神那裡找生命的意義.

你若在逆境中,請檢視一下你的生活態度,請看看你離開了神多遠,請悔改,請回轉,並期待復興.

你的生命可能面對逆境,但逆境不是你生命的主,神才是.請不要失望,要仰望神,要有「將臨期」的期待精神.我們不知要等多久,但請保持等待.

但願榮耀歸於聖父,聖子,聖靈.阿們.



\newpage

\section{有的加給沒有奪去}
\label{sec:129}
\textbf{李炳光牧師 - 有的加給,沒有奪去}
\newline
\newline
\hyperref[sec:128]{< < < PREV SERMON < < <}
~
\hyperref[sec:toc]{[返目錄]}
~
\hyperref[sec:130]{> > > NEXT SERMON > > >}


日期:2011年12月18日

題目:有的加給,沒有奪去

證道:李炳光牧師

各位早晨,非常高興,今天早上有機會和大家一同分享聖經的話語.我相信大家都覺得這個題目很特別,「有的加給,沒有的奪去.」如果在今天香港這麼喜歡抗議的氣氛下,可能又出來遊行了,可能有很多抗議,但是這些話是耶穌基督說的.講到這裡的時候,我立刻想到,我們今天很多時候都希望有的更有,還要更多.因為我們很多時候的祈禱,都是求上帝讓我們恩上加恩,力上加力.恩上又要加恩,力上又要加力,你會覺得人好像貪得無厭.

我相信大家對於城中發生的事,都非常留意.最近一位大富豪的兒子鬧離婚,原來他的生活奢華得這麼過份,一年有四分之一時間住在遊艇上,而且那位準備和他離婚的太太,每個月的花費要一千五百萬.你問為何他們享受這麼多?為何他們能夠一個人佔有這麼多?或者你更留意的,是一位女性大富豪,因為死了而引發爭家產.當然,這個爭家產的主角,大家都知道是輸了官司.人們說這是一個道德的審判,或者說是輿論的審判.但無論如何,大家對他已經有了定案,問題是太貪心了,為何有了還要更多.已經拿了二十多億,為什麼還要更多?

今天在這個世界上,很多人都希望能夠錦上添花,沒有人願意雪中送炭.錦上還要添花,是否很過份?但聖經有一句說話,是耶穌基督親自講的,不止一次提過,就是說:「凡有的,還要加給他；沒有的,連他所有的也要奪過來.」凡有的加給,沒有的奪去.如果這不是錦上添花,又是什麼?今天我們所讀的一段經文,大家讀的時候,便會覺得這段經文的主題,是說神按著人不同的恩賜,給他一些工作.有些人真的運用恩賜去賺回來,有些人則把它收藏起來.結果,神責備那個收起來的,稱讚那個有的,甚至從沒有的拿去給有的.

耶穌說:「凡有的,還要加給他；沒有的,連他所有的也要奪過來.」其實我們讀這些經文的時候,可能很困擾,神是否偏心?聖經不是說神是不偏待人的嗎?祂降雨給好人,也給不好的人,日頭曬好人,也曬壞人.「神愛世人,甚至將祂的獨生子,賜給他們.」神說給所有人,為何祂竟然說凡有的加給?沒有的就奪去,是否很不公平?「恩上加恩」是我們基督徒祝福,禱告的口吻,但是不是過分的要求?是否貪得無厭呢?

其實恩上加恩,美上加美,福上加福,甚至錦上添花,確有聖經根據,而且更加配合一個我們更深了解神心意的主題,包括了很實際的真理.恩上可以加恩,美上可以加美,福中可以加福,錦上可以添花.不是別人加上去,是神對忠心事奉祂的人所施的恩惠,亦是基督徒盡心,盡性,盡意,盡力愛主愛人的報酬.這是神的恩典,視乎我們是否努力實踐.

在福音書裡,有句話耶穌基督說了四次,今天和大家思考一下.首先看路加福音19章11節,我不詳細讀經文內容,大家可以慢慢看.這裡說:「耶穌因為將近耶路撒冷,又因他們以為　神的國快要顯出來,就另設一個比喻,說:「有一個貴冑往遠方去,要得國回來,便叫了他的十個僕人來,交給他們十錠銀子,說:『你們去做生意.』」」這個比喻和剛才那個很不同,十錠銀子,個個平均分,然後主人打發他們出去.結果這十個僕人去做生意,頭一個一錠銀子賺了十錠.

「主人說:『好!良善的僕人,你既在最小的事上有忠心,可以有權柄管十座城.』第二個來,說:『主啊,你的一錠銀子已經賺了五錠.』主人說:『你也可以管五座城.』又有一個來說:『主啊,看哪,你的一錠銀子在這裡,我把它包在手巾裡存著.我原是怕你,因為你是嚴厲的人；沒有放下的,還要去拿,沒有種下的,還要去收.』主人對他說:『你這惡僕,我要憑你的口定你的罪.』」結果這段經文下面這樣說:「就對旁邊站著的人說:『奪過他這一錠來,給那有十錠的.』他們說:『主啊,他已經有十錠了.』主人說:『我告訴你們,凡有的,還要加給他；沒有的,連他所有的也要奪過來.』」

這是耶穌第一次講這句說話,和剛才所講按才幹受責任有些不同.為什麼耶穌在不同場合同樣講這個真理?如果我們詳細看這段經文,我覺得十錠銀子分給十個人,每人一錠,基數相同.但按靈意來看,每個基督徒都有一錠銀子,這錠銀子預表什麼?其實我們每一個稱為神的兒女,神都給我們一位聖靈,住在我們心內,証明我們是神的兒女.有了這錠銀子,我們便重生得救.聖靈只有一位,普天下的基督徒都住在聖靈裡面,聖靈也住在我們褢面.這錠銀子在我們褢面生長,開花,結果,所以有稱為聖靈的果子.

加拉太書第5章說,聖靈的果子有九種,大家都知道,就是仁愛,喜樂,和平,忍耐,恩慈,良善,信實,溫柔,節制.但當你讀聖經,詳細研究原文時,聖靈的果子,很奇怪是用單數.九個果子,但是用單數,其實意思是說聖靈只有一位.一錠的銀子,會結出九個果子,或者結出五個果子,或者三個果子,或者十個果子.這個有一錠銀子的人,把它包起來,沒有賺回來,等於一個基督徒,受了聖靈,卻沒有見證一樣.

聖經裡談及結果子的原則是什麼?在約翰福音15章5節,葡萄園與枝子的比喻,耶穌基督說:「常在我裡面的,我也常在他裡面,這人就多結果子；因為離了我,你們就不能做甚麼.」耶穌又說:「我實實在在地告訴你們,一粒麥子不落在地裡死了,仍舊是一粒,若是死了,就結出許多子粒來.」我們在主裡面,主在我們裡面,到底我們有沒有為祂的緣故結出果子呢?其實有愛便有喜樂,便會結出和平,忍耐,恩慈,良善,信實,溫柔,節制等等九種果子.

「凡有的,還要加給他；沒有的,連他所有的也要奪過來.」今天我們從這段經文看到什麼?我們看到究竟我們付出了多少.神給我們有一定的恩賜,平均的,沒多也沒少,究竟我們付出多少?有一位作家描寫了一個很感人的故事:有一次他去到英國一個很落後,或是說很荒涼的鄉村,一直行走迷了路.他在欣賞大自然美景時,突然發現自己迷失了方向,不知到了哪裡.在最徬徨的時候,一路找,突然見到一條小路,路口是一條小徑,上面有一個木牌寫著:「請進來,喝一杯清涼的水 (Come in, have a cold drink)!」真好,一路行走,又口渴又迷路,被請進去喝一杯涼水.

他便沿著小徑走,一直走,見到一間茅屋,屋前有一個很小但很清澈的池塘.池塘旁有一個木牌,掛著一個水殼,寫著:「請隨便喝 (Help yourself)!」這位遊客作家覺得很奇怪,設法找茅屋的主人,結果找到兩位老人家,一對年老的夫婦.便問他們:「為什麼這樣做?很有意思,請人進來喝杯涼水.然後見到泉水的時候,請人隨便喝.這麼清涼的泉水,飲了確是很舒服.」這對老人家說:「先生,因為我們很貧窮,不能幫助什麼人.上帝給我們的恩賜,我們要盡量和人分享.」我覺得這個故事,實在非常感人.

有人說,我們好像一條船,這條船有兩個艙口,這邊入貨,那邊出貨.同樣今天神賜給我的,一定要拿出來,你愈付出愈有.另有一個相似的故事:在一個花園門口,有一個籃,放滿鮮花,寫著:「不用買,隨便拿 (Help yourself)!」旁人問花園園主:「為什麼,你不會賠本嗎?種花出來任人取.」那個女人說:「當我愈剪便愈有 (The more I cut,the more I have).」我覺得這樣很有意思,付出才有收穫 (Giving is receiving).今天我們很多時忽略了這個真理,神給我們每個人有同樣的恩賜,但是我們有沒有付出?如果收藏起來,可能會失去,連你所有的也拿走,這是非常重要的一個道理.

今天第二段經文,馬太福音第25章所說的,聖靈的恩賜是愈用愈多,聖靈的果子是愈結愈多.剛才我講的比喻,五千,二千,一千,耶穌說的是不同的恩賜.我把五千給你,把二千給你,把一千給你,我們不要妒忌.其實上帝給我們的恩賜是不同的,但每一個恩賜,如果你使用的話,一樣發揮作用.聖經說,收起來不用的,便將恩賜拿走.今天有很多人在教會裡有很多抱怨:「千萬不要答應,愈做愈多,不做最好,來到聽聽講講就好了.」有些人回到教會裡,愈做,牧師便給他更多的工作去做,大家很多都學會了「耍太極」,推掉算了.

俗語話,能者多勞.不要抱怨,可能是你任重道遠.神給你愈多的恩賜,祂希望我們用得愈多,否則有一天來算帳的時候,我們怎樣面對主呢?我知道在我自己的教會裡,有很多弟兄姊妹,在星期日忙得不得了,由朝忙到晚.有人說:「星期日你應該休息,為何用上這麼多時間?」但他說:「你不知道,這是我最開心的一天,因為我能夠把我所做的事貢獻出來.」

今天有詩班,有讀經,又有主席.我們每個人有機會參與的時候,是愈做愈多.我們不要抱怨,因為上帝給你多少恩賜,上帝會給你更多回報.當你愈有便愈多,收藏起來便失去.大家都知道我已經退休,但我還主辦一個佈道團,四處去.有人說:「你退下,休息一下豈不是更好.」但我發現,原來我愈做愈開心,愈做愈多事做.現在我講壇的事奉,比我未退休前還要多.差不多每星期我都要講,不但在香港,甚至到世界各地,都要負責講道.

有一天我的兒子跟我說:「父親 (Daddy),我和你談心事好不好?」很難得我的兒子跟我談心事.他說:「我很欣賞你.」我很開心,為什麼欣賞我呢?突然間他說:「父親,我很欣賞你,為什麼你這麼喜歡你的工作?」我不知道怎樣回答他,便說:「你不喜歡你的工作嗎?」「是的,做哪行厭哪行,真是煩.」但是我自己回想,我覺得我很開心去工作.我不覺得疲倦,也不覺得辛苦,為什麼?因為我覺得愈付出,愈擁有,愈做愈有力量.所以在座中,快要退休的朋友,請你不要說,退休後便躲起來.要愈做愈有能力,生命力愈來愈強,這是聖經所講的真理.

在第三個比喻裡,耶穌同樣講過這些話.大家翻到馬可福音第4章,這一次是另一個場合.剛才十錠銀子,每人一錠,又給五千,二千,一千.這是另外一個比喻,在馬可福音4章21節:「耶穌又對他們說:「人拿燈來,豈是要放在斗底下,床底下,不放在燈臺上嗎?因為掩藏的事,沒有不顯出來的；隱瞞的事,沒有不露出來的.有耳可聽的,就應當聽!」又說:「你們所聽的要留心.你們用甚麼量器量給人,也必用甚麼量器量給你們,並且要多給你們.因為有的,還要給他；沒有的,連他所有的也要奪去.」」

奇怪,耶穌又在這場合,突然講這句話.可見耶穌很重視這個道理,連續在四個不同場合中--以後還有一個--都是講這個道理.原來這個道理讓我們知道,神給我們的光是愈照愈明,愈照愈亮.主耶穌說:「我是世上的光.」我們要做光,為祂的緣故而發光.你們亦是世上的光,光照在世界上面,光來到黑暗便離去,光是愈照愈亮,有的還要加給他.如果你不發光的時候,就會慢慢暗起來,就沒有光.你愈作見證,就愈有見證,愈不為主作見證,就愈沒有見證.

真的,各位,你有沒有想過,神在你身上有沒有給你特別的恩典?如果你想一下,是愈想愈多.很多時我教人祈禱,牧師為什麼這麼多事情祈求?重複祈求上帝賜給我們一些東西便算了吧!但我說:「不是的,祈禱的內容很豐富.首先你要學習認罪,跟著要數算神的恩典.你試著數一下,愈數愈多,然後你告訴別人時,便愈講愈多.」今天我們的生命要發出光輝,我們愈為主的緣故去作見證,愈見證主的愛和主恩典的時候,你的機會愈多,你可以為主作工更多.今天我希望大家思考一下,能夠為主的緣故發出什麼光輝?其實,愈照愈明,愈發出光輝的時候,愈有力量.這是耶穌基督的吩咐,燈不是放在斗底下,是放在燈台上,就可以照亮一家人,今天我們有沒有為主的緣故而發光?

最後一次耶穌所講同樣的話,是在馬太福音13章,剛才是15章,現在是13章12節.「門徒進前來,問耶穌說:「對眾人講話,為甚麼用比喻呢?」耶穌回答說:「因為天國的奧祕只叫你們知道,不叫他們知道.凡有的,還要加給他,叫他有餘；凡沒有的,連他所有的,也要奪去.」」同樣講這句話,剛才是講聖靈的果子,愈結愈多.聖靈的恩賜,愈用愈多.接著是講神的光輝,愈照愈亮.現在這個是什麼?這個是天國的奧秘,愈知愈多,天國的奧秘是愈尋找愈有.

今天我們對於聖經的道理,有沒有覺得「又是那一套」,又是沉悶呢?我講道講了差不多五十年,有人說牧師是否還有話說?我說:「還有得講,而且愈講愈多.」我現在逢星期四晚有查經班,這個查經班差不多有一百人參加,幾十年沒停過.嘩!這麼多東西講嗎?真的,不是說謊.詩篇23篇,只有六節,但我講了23堂.真的嗎?牧師有沒有說謊?真的,錄了音的.原來這麼多東西講,聖經的道理,愈發掘愈多.如果我們停留在表面上,覺得這不過是老生常談,就永遠沒有辦法再在聖經上,在真理上,在天國的奧秘上發掘得更多.

有一次,王明道先生,一位中國很著名的傳道人,到一間禮拜堂講道.有一位長老跟他說:「王先生,你不要希望用聖經來感動我們,因為這本聖經我讀了幾十年,我什麼都知道.」也就是說,他不想再去發掘更多這個奧秘.其實天國的奧秘是愈知愈多.今天讓我們不要忽略神的說話,你愈有便愈加給你,愈沒有,連所有的也拿走.耶穌提及這四點,是萬古常新,放之四海而皆準的道理.

神賜給人肉體和靈魂,祂賜給我們聖靈.這是最平等的,但我們不加運用便會退化失靈,或者枯乾腐朽.神賜給我們不同的才幹,如果我們不加以利用,不經常操練的話,便因此退步.一個音樂家,他不多些練習,彈奏或獻唱就退化,是非常可惜的.神給我們的見證也一樣,愈有見證,愈多見證,愈不見證,愈無見證.神的話,天國知識的奧秘,是愈尋找愈多,是無止境的.如果大家都同意這真理的話,我們就必須注意以下幾點:

第一,神是一切恩賜的賜予者,祂有權賜予,也有權收回.這很重要,這是第一個道理.第二,我們有沒有加以利用?盡力勞其筋骨,苦其心志.這是恩上加恩,力上加力的道理,不是貪婪,而是當你愈利用時便愈多.第三,上帝給我們的恩賜,是用之於己,還是施之於人?其實你留心一下,耶穌基督所講的八福,沒有一樣是為自己,都是為他人.當你付出給別人時,福氣便加在你身上,關心他人,為主為人所得的酬報,是愈來愈多的.

今天在教會和社會裡,需要更多付出的人,有服務貢獻的人.台灣有一位著名的女作家,她寫了一個比喻:這世界上有三種人,一種人覺得世界欠他很多,所以他盡量巧取豪奪,盡量去爭取.這種人很不開心,因為覺得這個世界欠他太多.這世界也有第二種人,這種人是人不欠我,我不欠人.我是中規中矩,我不得罪人,人也不要得罪我.我過著一種自我滿足的生活,但這種人很平凡,只為自己而活,但他很幸運.這世界並不是只有這兩種人,另外一種人就是覺得自己欠這世界很多,他盡量付出,而且愈付出愈有.

各位,感恩的人,是恩上加恩.各位想不想恩上加恩?大家都想,是不是?多利用,多付出,可以恩上加恩.你想不想錦上添花?當然想,多關心別人,多雪中送炭,你就會錦上添花,神一定在我們的人生裡添上美麗的花朵.有一班學生,跟隨老師--這位老師年紀已相當大--去到一個荒山空谷旅遊.突然有一個學生大聲叫,發現空谷有回音.他再大聲,回音(feedback)愈大,再大聲,好像立即有人回應.那班學生愈玩愈開心,突然,老師見他們在忘形開心的時候,說道:「各位同學,這就是人生,你愈大聲,回應愈大,你愈付出便愈有.」

我們同心祈禱:「感謝主給我們有機會,思考聖經的話.求主讓我們知道,神給我們每個人有同樣的靈,亦有同樣的恩典.有時大,有時小,但無論如何,我們要多運用.主啊,求你叫我們明白這句說話,你重複多次講,有的便加給,沒有的便奪去.讓我們清楚知道,我們愈付出愈有,收起來的便失去.主啊,幫助我們,用你的話語激勵我們,使我們願意付出更大的愛,而我們所得的愛更多.當我們愛你更深,我們的愛更大,因為你答應過,用你永遠的愛愛我們.感謝主,願你的話語成為我們的力量,我們願意將你的話語藏在心裡.奉耶穌基督聖名祈禱,阿們.」



\newpage

\section{光明與黑暗}
\label{sec:130}
\textbf{鄧瑞強博士 - 光明與黑暗}
\newline
\newline
\hyperref[sec:129]{< < < PREV SERMON < < <}
~
\hyperref[sec:toc]{[返目錄]}
~
\hyperref[sec:131]{> > > NEXT SERMON > > >}


約翰福音 1:1-14
\newline
\begin{longtable}{cl}
\hline
\hline
章節 & 經文 (和合本修訂版)\\
\hline
1:1 & \begin{tabularx}{0.7\textwidth}{X} 太初有道，道與神同在，道就是神。 \end{tabularx} \\ \\ \relax
1:2 & \begin{tabularx}{0.7\textwidth}{X} 這道太初與神同在。 \end{tabularx} \\ \\ \relax
1:3 & \begin{tabularx}{0.7\textwidth}{X} 萬物都是藉著他造的，沒有一樣不是藉著他造的。凡被造的， \end{tabularx} \\ \\ \relax
1:4 & \begin{tabularx}{0.7\textwidth}{X} 在他裡面有生命，這生命就是人的光。 \end{tabularx} \\ \\ \relax
1:5 & \begin{tabularx}{0.7\textwidth}{X} 光照在黑暗裡，黑暗卻沒有勝過光。 \end{tabularx} \\ \\ \relax
1:6 & \begin{tabularx}{0.7\textwidth}{X} 有一個人，是從神那裡差來的，名叫約翰。 \end{tabularx} \\ \\ \relax
1:7 & \begin{tabularx}{0.7\textwidth}{X} 這人來是為了作見證，是為那光作見證，要使眾人藉著他而信。 \end{tabularx} \\ \\ \relax
1:8 & \begin{tabularx}{0.7\textwidth}{X} 他不是那光，而是要為那光作見證。 \end{tabularx} \\ \\ \relax
1:9 & \begin{tabularx}{0.7\textwidth}{X} 那光是真光，來到世上，照亮所有的人。 \end{tabularx} \\ \\ \relax
1:10 & \begin{tabularx}{0.7\textwidth}{X} 他在世界，世界是藉著他造的，世界卻不認識他。 \end{tabularx} \\ \\ \relax
1:11 & \begin{tabularx}{0.7\textwidth}{X} 他來到自己的地方，自己的人並不接納他。 \end{tabularx} \\ \\ \relax
1:12 & \begin{tabularx}{0.7\textwidth}{X} 凡接納他的，就是信他名的人，他就賜他們權柄作神的兒女。 \end{tabularx} \\ \\ \relax
1:13 & \begin{tabularx}{0.7\textwidth}{X} 這些人不是從血生的，不是從情慾生的，也不是從人的意願生的，而是從神生的。 \end{tabularx} \\ \\ \relax
1:14 & \begin{tabularx}{0.7\textwidth}{X} 道成了肉身，住在我們中間，充充滿滿地有恩典有真理，我們也見過他的榮光，正是父獨一兒子的榮光。 \end{tabularx} \\ \\
[1ex]
\hline
\hline
\end{longtable}
日期:2011年12月25日

主題:光明與黑暗

經文:約翰福音1:1-14

證道:鄧瑞強博士

各位弟兄姊妹,主內平安.

「從前,有一個國王,他有三個兒子.他正在考慮,要將王位傳給哪一個兒子.他希望將王位傳給一個最有智慧的兒子.於是,他召齊三個兒子.對他們說:『我給你們每人一兩銀.我們王宮外,有一個貨倉,你們用這些金錢,去買些東西回來,填滿那貨倉.誰花最少金錢而能填滿那貨倉,我就將王位傳給誰?』於是,三個王子帶著金錢,出去各自買東西.過了不久,三個王子都回來了.國王問大王子:『你買了什麼?』大王子說:『我用盡你給我的錢,買了很多沙泥.』沙泥雖然平,但一兩銀買回來的沙泥,仍只佔貨倉的一個很少位置.國王問二王子:『你買了什麼?』二王子說:『我用盡你給我的錢,買了很多禾稈草.』二王子買回來的禾稈草,比剛才大王子買回來的沙泥,真的多很多.雖然如此,仍不足以填滿一個貨倉.最後,國王問小王子:『你買了什麼?』小王子說:『父王,請多等一會,到了晚上,我才能告訴你我買了什麼.』等到夜幕低垂的時候,國王問小王子:「你到底買了什麼回來呢?」小王子這時帶著各人,進入漆黑一片的貨倉,然後點燃一枝燭光.頓時,全貨倉充滿了光明.小王子對父王說:『我用幾文錢,買了這枝蠟燭.』國王很開心,因為小王子不是花盡金錢去買些無謂的東西,也因為小王子明白一個道理,要填滿這麼大的一個空洞空間,物質是無用的.國王於是將王位傳給小王子.」

我們每個人的內在生命,多少像一個偌大的空間.裡面很深,很多層次,也很幽暗.這個空間,是物質無法滿足的,是物質無法填滿的.只有光明,才能填滿這空間.

今日講道的主題,是「光明」來到我們的黑暗世界的故事.

講道的

經文:約翰福音1: 1-14

1:1　太初有道,道與神同在,道就是神.1:2　這道太初與神同在.1:3　萬物是藉著他造的；凡被造的,沒有一樣不是藉著他造的.1:4　生命在他裡頭,這生命就是人的光.1:5　光照在黑暗裡,黑暗卻不接受光.1:6　有一個人,是從神那裡差來的,名叫約翰.1:7　這人來,為要作見證,就是為光作見證,叫眾人因他可以信.1:8　他不是那光,乃是要為光作見證.1:9　那光是真光,照亮一切生在世上的人.1:10　他在世界,世界也是藉著他造的,世界卻不認識他.1:11　他到自己的地方來,自己的人倒不接待他.1:12　凡接待他的,就是信他名的人,他就賜他們權柄,作神的兒女.1:13　這等人不是從血氣生的,不是從情慾生的,也不是從人意生的,乃是從神生的.1:14　道成了肉身,住在我們中間,充充滿滿地有恩典有真理.我們也見過他的榮光,正是父獨生子的榮光.

這段經文的開首,回應著舊約創世記的開首,講到萬古之先的事情.

1:1　太初有道,道與神同在,道就是神.1:2　這道太初與神同在.

在原初,在太古,在萬有之先,有「道」.

「道」就是「說話」,神有話要說,神有所啟示.神的說話及啟示,充滿神自己的本性.「道就是神」.這猶如一個絕對誠實的人,他說的話就是他本性的流露.他的本性,他的心思,他的感情,都在他的話裡流露出來.「話如其人」.

1:3　萬物是藉著他造的；凡被造的,沒有一樣不是藉著他造的.1:4　生命在他裡頭,這生命就是人的光.

萬物的存在都依賴這「永恆的道」,萬物的意義都由這「永恆的道」去支撐.至於萬物之靈的人類,更需要「永恆的道」的指引.人是尋索意義的動物,人渴望生命有內容,人希望生命有方向.經文說:「生命在他裡頭,這生命就是人的光.」人要尋索的生命內容,就在這「永恆的道」裡面.人渴望的光明,就在這「永恆的道」裡面.

在萬古之先,一切都隱沒有黑暗裡.然後,「永恆的道」像探射燈一樣發出光明.想像一下,在漆黑的舞台裡,探射燈照到那裡,那裡的物件便顯出來了.這燈光照著你,你便出現在人生的舞台裡.你在這舞台裡如何行,需要這光的照耀.這光照在你前面,你便知道踏向前是安全的.這光一熄滅,你便會不知所措.「永恆的道」讓你在人生的舞台上有方向.

在神的創造裡,一切都應該如此美好.人在「永恆的道」那裡找到真正的生命的模樣,人按照道的指引而行,活出真正的人性.

但是,問題出現了.

1:5　光照在黑暗裡,黑暗卻不接受光.(《新譯本》的翻譯是:「光照在黑暗中,黑暗不能勝過光.」)

人生命裡頭有黑暗,這黑暗不喜歡光明,這黑暗不接受光明,這黑暗與光明對抗.縱使黑暗不能勝過光明,仍是對抗光明.人性的黑暗有一種魔力,叫人明知光明是好,也拒絕光明.這就像我們都想身體健康,我們都知道不要吃那些高脂肪,高鹽的食物,但是,只要面前有炸薯片,炸薯片仍會自然地走入口中.我們就像電影《魔戒》的主角,我們很想毀滅那魔戒,但魔戒的魔力仍無處不在,就是在最後關頭,也差點陷在魔戒的黑暗力量中.

早前,內地佛山巿有一個小女孩被車撞倒.這原先應該只是一意外.司機知道意外發生,應停車處理,召喚救護車.但若這司機知道發生了意外,卻思前想後,想到自己可能要接受的刑罰,想到自己要作出的賠償,想到這裡沒有人看見自己,想到逃之夭夭是最化算的處理方法的話,則他內心的黑暗便與人性的光明在鬥爭了.他最後走了,他的黑暗充滿了他.但是,「永恆的道」的光明仍在,顯出他的黑暗.這女孩在地上,有人走過,這個人或許看不見,或許他看見.若他看見,他可能思前想後,想到這是否是一個騙局,想到過去很多人因想幫人而被人屈,惹得一身蟻,想到多一事不如少一事.這途人過去了.如是者,十多名途人都過去了.

這原先是人性光輝流露的事件,只要有一人為這女孩操心,打一個救援電話,光明便照在這漆黑的晚上了.光明在人心呼喚,但黑暗卻拒絕光明.一個途人又一個途人,一次又一次,光明在呼喚,黑暗卻拒絕.有人說:「若我是那些途人,我也會悄悄離開.在這個充滿黑暗的世界,在這個充滿騙局的社會,在這個道理講不清的國家,離開是上策.」黑暗不單在人的心靈深處,也在社會的制度深處,亦在整體的價值取向的深處.這一切加起來,是深不見底的黑暗.

很多國家財政出現困難,解決的方式,我想,只能是大家節約一點,過一點簡樸的生活,並且各人貢獻多一點,納稅多一點,讓國家財政慢慢復完.但是,這些國家很多人民反對這些緊縮措施.他們不想自己的利益受損,他們不願意只吃半個蘋果而讓半個蘋果給別人吃,他們只想自己有蘋果吃而不想將蘋果留給下一代吃.這樣一個只顧自己,沒有他人,不顧將來的生活方式,看不見光明,只剩人性的黑暗.

面對這無邊的黑暗,神並不袖手旁觀,神有話說,神有「道」要言傳,神有所行動.

1:9　那光是真光,照亮一切生在世上的人.

《現代中文譯本》將這節經文譯作:「那光是真光,來到世上照亮全人類.」

《現代中文譯本》的翻譯,重點在:真光要來到世上.

神有「道」要說,他的「說」是他的行動,神要說的「道」成為具體的生命,來到世上.「道」要成為生命的光,照亮人性及世界的黑暗.

不知什麼時候開始,香港人說「講耶穌」.我們的「道」只是用來「講」的.但神要言說他的「道」時,他不是「講」,而是將「道」差派到人間,成為人間具體的生命.這不是「講耶穌」,而是「活出耶穌的生命」.這不是「講道」,而是「活出道」.

真正光明,要透過具體的生命活出來.

德蘭修女(Mother Teresa)講過:「Love until it hurts」.仁慈的施予,這是很好的.但人不單需要你的施予,更需要你的寬恕,接納,擁抱,肯定.這是需要你犧牲的,這是會令你受傷的,這是令你難受的.真愛,是需要為愛而受苦的.你施予金錢給一個乞丐,你不會有多大痛苦；但你擁抱他,則需要很大的勇氣.為一個痳瘋病人祈禱,你不會損失太多；但你若握著他的手,你便需要衝破很多心理障礙.去愛一個不太可愛的人,你可能需要付出一點心力；但去寬恕你的仇人,則會令你受苦了.「Love until it hurts」,這種愛,是來到人間的永恆的道要呈獻的愛,這是主耶穌的愛,是有救贖作用的.我們如此去愛,能讓我們與基督的生命結連,活出生命的光明.

在越黑暗的地方,越容易發現光明.天越黑,星越明.塵世最黑暗的地方,莫如集中營.有一個住過納粹集中營的精神病學家Viktor Frankl,寫了一本書 ``Man"s Search for Meaning",記述在黑暗中所看見的光明的事.在集中營,每人得到的食物是很少的,人常在飢餓的狀態中.但是,在生理的,心理的匱乏與壓力下,人是否就只能是環境制約的產物?人在此是否就失去自由抉擇的力量?Viktor Frankl看見,就算在這極端的生存條件下,仍然有人願意將最後一片麵包給另一個人吃.他發現,就算一個人在集中營裡被貶抑為非人,被視為純粹會動的生產機器,被剝奪一切,人仍然保有其內心的自由,仍可憑這自由,發揮生命的光.一個看來已被完全剝奪人性尊嚴的人,因著他生命的光,他的人性尊嚴顯得偉大無比.

「Love until it hurts」,越是艱難,越是黑暗,光明顯得越耀眼.

讀過一個故事是這樣的:

「1933年,當美國出現經濟危機的時候,Harrison 紡織公司(Harrison Textile Companies)因一場大火化為灰燼.3000名員工悲觀地回到家裡,等待公司的董事長宣佈公司破產,也等待失業的來臨.在無望而漫長的等待中,他們終於接到了董事會的信:全體員工能繼續支薪一個月.

在全國經濟蕭條時,能有這種消息傳來,員工們深感意外.

一個月後,正當他們為未來的生活發愁時,他們又接到公司的第二封信,董事長宣佈,再支付全體員工一個月薪酬.3000名員工接到信後,不單是驚喜,更是熱淚盈眶.在人人失業的年代,能得到如此照顧,誰不會感激萬分呢?第二天,全體員工紛紛返回公司,自發地清理災場,擦洗機器,有人主動去聯絡被中斷的貨源.三個月後,公司重新運作.有報章這樣評論:員工們上下一心,日夜不懈地賣力工作,恨不得一天工作25小時.曾勸公司董事長領取保險金,然後一走了之,或批評董事長感情用事,缺乏商業精神的人,開始服輸.」【參看張健鵬,胡足青編:《虛掩的門》(香港:三聯,2001),頁185-186.】

在黑暗裡,就是因為有些人,就算自己有所損失,也堅持讓人看到光明,這些人成為了世上的光明.在黑暗的世界裡,這些光明顯得很耀眼.

1:10　他在世界,世界也是藉著他造的,世界卻不認識他.1:11　他到自己的地方來,自己的人倒不接待他.1:12　凡接待他的,就是信他名的人,他就賜他們權柄,作神的兒女.1:13　這等人不是從血氣生的,不是從情慾生的,也不是從人意生的,乃是從神生的.

神看見人間的黑暗,他沒有置之不理.他來到人間,他給人光明,但這光明也顯明了人的黑暗.有些人寧願在黑暗裡,拒絕光明.有些人卻願意回應,願意與基督結合,願意成為光明的一分子,縱使成為光明之子,需要承受十架的艱難.這些願意回應的人,他們是神的兒女.

你願意成為光明之子,成為神的兒女嗎?

但願榮耀歸於聖父,聖子,聖靈.阿們.



\newpage

\section{恩典與真理}
\label{sec:131}
\textbf{伍渭文牧師 - 恩典與真理}
\newline
\newline
\hyperref[sec:130]{< < < PREV SERMON < < <}
~
\hyperref[sec:toc]{[返目錄]}
~
\hyperref[sec:132]{> > > NEXT SERMON > > >}


約翰福音 1:14-18
\newline
\begin{longtable}{cl}
\hline
\hline
章節 & 經文 (和合本修訂版)\\
\hline
1:14 & \begin{tabularx}{0.7\textwidth}{X} 道成了肉身，住在我們中間，充充滿滿地有恩典有真理，我們也見過他的榮光，正是父獨一兒子的榮光。 \end{tabularx} \\ \\ \relax
1:15 & \begin{tabularx}{0.7\textwidth}{X} 約翰為他作見證，喊著說：「這就是我曾說：『那在我以後來的先於我，因為在我以前，他已經存在。』」 \end{tabularx} \\ \\ \relax
1:16 & \begin{tabularx}{0.7\textwidth}{X} 從他的豐富裡，我們都領受了恩典，而且恩上加恩。 \end{tabularx} \\ \\ \relax
1:17 & \begin{tabularx}{0.7\textwidth}{X} 律法是藉著摩西頒佈的；恩典和真理卻是由耶穌基督來的。 \end{tabularx} \\ \\ \relax
1:18 & \begin{tabularx}{0.7\textwidth}{X} 從來沒有人見過神，只有在父懷裡獨一的兒子將他表明出來。 \end{tabularx} \\ \\
[1ex]
\hline
\hline
\end{longtable}
日期:2012年1月1日

題目:恩典與真理

經文:約翰福音1:14-18

證道:伍渭文牧師

太初有道

約翰一1-18是一首頌揚道的詩章,首句「太初有道,道與神同在,道就是神,這道太初與神同在.」就把人攝住.道來自在那裡?道來自何方?望向時間,道不生在時間之內；望向蒼茫宇宙,也看不到道的開端.未有時間,未有空間,在一切之始,就有了道.詩歌的開始,令人感到敬畏.到了十四節:「道成了肉身,住在我們中間,充充滿滿的有恩典有真理.我們也見過他的榮光,正是父獨生子的榮光.」聽的人更覺驚懼.

在此之前,約翰的讀者容易明白道在太初之先,道創造萬物.他們明白,因為希臘文化都這樣相信.眼看到的世界是從眼看不到的道而來,世界和這道是區分的,這看得到的世界,是道的拷貝,因為是翻版,所以會改變,會朽壞.我結婚前腰圍二十八吋,但現在不用提了.但到了十四節「道成了肉身,住在我們中間」,就不容易明白了.

道怎可能成為肉身呢?肉身意味著限制,肉身在這教堂,就不能在家中了.我留在香港工作,就不能常常見到移民北美的家人,為了工作成了太空人.肉身也會漸漸衰殘朽壞,美人遲暮,明日黃花,所以魯益斯(C. S .Lewis) 說,最大的神蹟不是復活,是道成肉身.道若能成為肉身,為甚麼不能水變酒呢?為甚麼不能使瞎眼看見,跛腿行路,死人復活呢?相比之下,這些神蹟輕而易舉哩.

基督成為肉身,就像一個人,深深呼吸,閉著氣,潛入深海至深之處,尋找失去的寶貝.愈潛入海底,就遠離陽光,愈來愈黑暗,愈來愈寒冷,最後到達了那幽黑冰凍的海底,把淤泥撥開,找到埋在下面的一顆明珠,然後緊緊握著這明珠,慢慢升上海面,高舉這顆明珠,使它原來的光澤重現世間.我們就是被埋在深海至深之處的明珠,罪沾污掩埋了我們原來的光澤.基督成為肉身,把我們尋回.

住在我們中間

道不單成為肉身,並住在我們中間.「住」和「帳幕」是同一個字,指道在人間安設帳幕.舊約時候,上帝的子民進入迦南應許之地前,曾在曠野飄流四十年,他們住在流動的帳蓬,以會幕為聖殿,他們移動,會幕也跟著他們移動,上帝同他們一起飄流,他們前進時,白天有雲柱遮蓋太陽,晚上有火柱照引前路.「住在我們中間」表示上帝和我們一起走人生的路.因為道成為肉身,馬利亞童女聖靈感孕,經歷母親生產之痛.道成為肉身,約瑟,馬利亞帶著耶穌,逃難埃及,躲避希律王的刀劍.道成為肉身,耶穌凡事受過試探,與我們一樣,祇是他沒有犯罪,所以耶穌能體恤我們,他是我們的大祭司,我們祇管到他面前,求憐恤得幫助.

充充滿滿恩典有真理

經文又說:「充充滿滿的有恩典有真理.」恩典和真理,就是舊約時代上帝的慈愛和信實.信實指上帝的應許不會落空,他怎樣拯救子民離開埃及為奴之地,也會引導他們行走曠野,帶領他們進入應許之地,他的應許確實無誤；他頒下的律法是真理,引導我們走在義路,不致迷失.在我們干犯律法時,他有赦免的恩典,把我們挽回,他是慈愛的上帝.在人生路上,我們需要真理指引,跌倒時需要恩典挽回.

一位行淫時被捉拿的婦人,被帶到耶穌面前,試探耶穌的人問:可否按摩西的律法用石頭打死她.耶穌說:「你們中間沒有罪的,誰就可以先拿石頭打他.」(約八7) .「他們聽見這話,就從老到少,一個一個的都出去了,祇剩下耶穌一人.還有那婦人仍然站在當中.」這故事說明每人都有罪,沒有行真理.留下的婦人有罪,離開的人也有罪,祇有耶穌沒有罪,祇有他可以施行審判.但耶穌說:「我也不定你的罪.去吧!從此不要再犯罪了.」(約八11)

耶穌赦免這婦人,我們看到恩典；但耶穌為了真理,他走上十字架,為這婦人贖罪.這婦人因為赦罪的恩典,開始過新生.恩典不單遮蓋真理的過失,恩典更能挽回人,建立人；恩典使人有勇氣面對真理,愛慕真理.

不離不棄的愛

最近一套非常叫好叫座的電影-作死不離三兄弟(3 Idiots) 就說明,朋友不離不棄的愛,可以改變了一個人的生命.三兄弟之一的Raju因為醉酒犯校規被校長勒令退學,無法面對家人(老父臥病；姐姐貌不揚,沒有嫁妝,嫁不出去；媽媽退休,沒有工作),從校長辦公室的窗戶躍下,在醫院兩個月,兩位好朋友用盡方法使他甦醒過來.因為朋友的愛,他覺得自己有價值,有自信.

Raju申請跨國公司職位面試時,被問及為何成績那麼差呢?他坦白說出真相.「我小時候成績很好,但父母希望我能夠進有名的大學脫貧,令我很懼怕,因為這裡人人都像瘋了一般爭取好成績,如果你不是第一,你就不算甚麼,我愈來愈懼怕.......斷了十六節骨給我兩個月反省,終於我醒悟過來.今天我不祈求上蒼給我工作,能夠活下去我就感謝上蒼了.如果我被拒絕,我不懷怨,我想這一生仍可以做一些有價值的事.」

面試他的說:「你這樣坦白,對公司來說並不是好事,我們需要圓通 (diplomatic)的人與客戶聯系,你太直率了,但如你向我們證明可以控制一下這態度,也許會聘請你.」Raju說,「我不能改變,你保留你的工作,我保留我的態度.」面試他的說:「過去二十五年我面試無數的人,他們為了工作,面試時都變成唯唯諾諾,你從那裡來的?」兩位朋友不離不棄的伴隨和關愛,改變了Raju的人生.

這個聖誕,我要燃點燭光

在新的一年,我們需要真理引路,更需要恩典支持...

Howard Thurman (1899-1981)I will light Candles this Christmas;Candles of joy despite all sadness,Candles of hope where despair keeps watch,Candles of courage for fears ever present,Candles of peace for tempest-tossed days,Candles of grace to ease heavy burdens,Candles of love to inspire all my living,Candles that will burn all the year long.

這個聖誕,我要燃點燭光；喜樂的燭光,在悲傷中搖曳,盼望的燭光,照亮四周的沮喪,勇敢的燭光,直視懼怕亦步亦趨,平安的燭光,在狂飆中柔暉穩妥,恩典的燭光,使重擔輕盈一點,仁愛的燭光,激勵我的生活,從歲首到年終,燭光常在.

但從歲首到年終,燭光真的常在嗎?過去一年,大家都知道我做過右眼視網膜的手術,這手術令我有很深的體會.手術過後三天解開紗布,眼睛完全看不到東西,令我吃驚,這是我從未有過的經驗,但數天後視力慢慢回來了,每一天都比前一天好.我領會到一個真理:現在壞不重要,最重要是每天好過一天.不過我知道,我眼睛恢復視力後,將來也會因為自然規律愈來愈差,最後老眼昏花；我的身體也會隨時日漸漸衰敗,最終成為一坯塵土.但我知道,外在身體雖然朽壞,但內在的新人,一天新似一天.我等候基督再來,身體復活.因著信,從歲首到年終,燭光常在；因著信,從現在到永遠,燭光常在.因為

「太初有道,道與神同在,道就是神,這道太初與神同在......道成了肉身,住在我們中間,充充滿滿的有恩典有真理.」阿們.



\newpage

\section{義薄坤輿}
\label{sec:132}
\textbf{伍偉亨教授 - 義薄坤輿}
\newline
\newline
\hyperref[sec:131]{< < < PREV SERMON < < <}
~
\hyperref[sec:toc]{[返目錄]}
~
\hyperref[sec:133]{> > > NEXT SERMON > > >}


日期:2012年1月8日

主題:義薄坤輿

證道:伍偉亨教授

基督教「將臨期」的一個特點,就如同伍渭文牧師最近與大家分享的一篇文章《將臨的已經來到了》──是一個等候的操練.伍牧師很能一針見血地講出,我們活在現代科技的社會裡頭,等候是很不容易.在「將臨期」時,這值得我們深刻地來思考:如何等候?

我今天的分享,主要想作一點的補充,因為普遍地當我們想到的「等侯」,是一種完全的靜態,無所事事的等候.假使你細心地聽今天所宣讀的經文時,你會發現聖經裡頭所提到的等候,是一種很深,很主動的預備.這一種為主降臨的預備,在聖經的舊約,新約都重複地出現:「預備主的道,修直他的路」.

我們如何準備自己呢?這也是我想與大家一同來思考,尤其是來思考有關「義」這個課題.

彼得後書三章13節:「但我們照他的應許,盼望新天新地,有義居在其中.」使徒強調「義」是作為天國社群的特質.這種公義不單單是在彼岸的事情,也是基督徒在此生的召命所在──「誠實從地而生,公義從天而現」(詩85:11)──這個是屬天的公義,也期待展現在人間；這公義是上帝對祂的兒女們在塵世中的要求──以公義預備上主的道路──「公義要行在他面前,叫他的腳蹤,成為可走的路.」(詩篇85:13)

一般來說,我早上很早就上班,而每一次離開地鐵站的出口處時,總是會碰到一些七,八十歲的老婆婆,老伯伯,他們慣常性地用一個很卑微以及乞求的態度,期待著能夠獲取乘客手中那份免費的報紙,而他們之間常是互相競逐地爭奪著.他們為的是甚麼?只不過為了能多賺一塊,半毛的錢,每一次當我看到這個畫面的場景,我想起我們中國傳統的《中庸》,《禮運大同篇》裡,描述到一個大同的社會,應該是一個老有所終的社會──這彷彿對我們生活在香港這個社會來說,是一個很深的譴責.怎樣才可以有一個「老有所終」的社會呢?香港相對上是一個比較富裕的社會,但不少老人家終老的地方就在地鐵站閘口的角落裡,在那裡競逐著一些不值錢的,破爛的報紙.一個資本主義市場最血淋淋的圖畫,就在我每天上班的時候,很清楚地展現在我眼前.

我相信這個很遠離基督教信仰對一個公義社會的要求.從佔領美國華爾街以至本土社會上的抗爭,很明顯地指向「公義」這個主題. 或許我們當中各有不同的立場,或許你不完全同意一些左翼社群對公義這課題所提出的答案.但我們很難否認,他們對公義的渴求,是相當深層的,且也貼近上主或聖經對義的要求.

在六十年代的基督教傳統裡,南美洲的解放神學是很強調社會的公義,它以馬克斯主義的社會分析,作為基督教神學的建構綱領.作為一個唸神學的人來說,我對這種做神學的方法,在我個人的專業裡頭,我對這個存在著許多的疑惑,批判和問題.解放神學很強調的一個詮釋導向,是以貧窮人為優先.我個人認為是相當有聖經的基礎,特別是新約聖經.在新約裡,耶穌特別向貧窮的人,社會上卑微的人有憐恤.對比香港這個社會,主要是傾向一些富裕的,既得利益的人,那我們就很難看到那種的差異.學習傾向卑微的人,是基督信仰裡對「義」的實踐,也是香港一些「中產教會」所要學習的一個功課.

就我個人對信仰的理解而言,一個忠於基督信仰,聖經啟示的信徒,你很難去擁抱香港這個資本主義的社會.如果你和我同樣相信香港現在所走的方向,不全然是上帝對我們的期待時,我們作為基督徒,是有責任去改變──「預備主的道,修直他的路」.

我出生於香港,長於香港,粗俗的一句是:喝著資本主義的奶水長大的一代.要我們放下數十年生活的形態,我們習慣用市場,效率,計算等的運作來決定哪些事要做,不要做,當我們習以為常的時候,是很難去作改變的,很難以慈愛,公義為判斷的依歸.但正如先知在聖經裡頭說:「在曠野預備耶和華的路,在沙漠地修平我們　神的道.」(賽40:3 )困難從來就不應該讓我們放棄要承擔責任的藉口.香港教會需要更多的人,弟兄姊妹的努力,在這裡頭明辨心思,甚至不是個別信徒可以做得到的,我們需要思考,如何在一個活生生,有血有肉的世界裡,活出我們的信仰.我相信有許多香港的弟兄姊妹也開始看見這種的景像,覺得社會需要改變.

我最近參加了一些弟兄姊妹的群體,都是學有專精的人──有的是從事法律,經濟,心理學,神學──我們大家走在一起,重新的來思考:在香港這個社會,身為一個基督徒,在公眾的一個領域裡,我們如何活出我們的信仰?但在這以外,未必每一位弟兄姊妹都適合做這些時,但我們總有許多的空間可以有一些的改變──重新學習如何建立一個合乎公義的社群.

記得當我在德國時,是一位窮學生,每一次去超級市場,是不需要多思考,只要是價錢便宜的就可以了.當住久了,生活開始安頓下來,與一些德國朋友交往之後才明白,原來在德國的超級市場,雞蛋價錢的差異可以遠超一倍之多,原因是因為有四種不一樣的雞蛋,這種的分類是以被飼養的雞的活動空間大小為正比,越便宜的雞蛋,就是被困在一個很狹小的籠子裡頭,最貴的雞蛋,就是「走地雞」(牠們能夠在農場裡,隨意的跑動)的雞蛋.我看見不少的德國學生,他們買雞蛋從來不會買最便宜那種,而會買那種給予許多空間讓雞活動的雞蛋,這樣才合乎他們對道德的判斷.這使我想到,我們在香港習慣以市場,效率,計算等的運作來決定一些思考的模式,我們能否在當中有一些的更改呢?當我們買一樣物品的時候,不單單以價錢作為考慮的準則... 這幾年提倡公平貿易,它背後的理念,與這個很相似.資本主義的魔咒,很深地籠罩在香港(我想國內來的弟兄姊妹,更加明白這種文化的差異).

這改變不單是追求聖經上的教導,也是基督徒靈性更新的重要一步.改變一些是無意識中會做的事,在意識上去判斷,甚麼時候需要節儉?甚麼時候我們的付出是應該憐恤貧窮人的需要?如此的話,當討論最低工資時,就能夠有一個完全不一樣的空間去考慮.

讓信徒走出移風易俗的第一步,改變一個像香港這樣那麼根深蒂固的資本主義社會,又缺乏靈性的社會,使這社會,因著我們的改變,可以得以重新看見上帝的義.

義──作為基督徒去追求社會公義時,我們該清楚地意識到,基督信仰是我們一切努力的基礎及準繩.這信仰使我們在追求社會公義的同時,應該與其他人顯得有點不一樣.公元五世紀時,教父奧古斯汀在他的著作《上帝之城》(the City of God)裡,思考永恆與塵世的關係,他很清楚地表明國家政權的設立,是上主給墮落之人的恩眷,所以人是得享平安.因此就算一個政權有多極權,奧古斯汀起碼是努力想達致讓老百姓生活在一個平安的環境裡.這平安對奧古斯汀來說,是眾善之首.奧古斯汀也指出,平安的出現,很大的程度是依賴國家政府去踐行公義.如同我們今天所唸的詩篇85:10:「公義和平安,彼此相親」.這是奧古斯汀按照聖經的訊息,對今天在中國社會所倡議和諧社會的提醒.和諧不是至義而有的,只有在公義的社會裡,平安才會在裡頭顯現.奧古斯汀同時很清楚無誤地指出,塵世的公義,平安的追求都是相對的.一切的努力,都有其限制.

上帝的救恩是臨近敬畏祂的人,詩篇85:9:「他的救恩,誠然與敬畏他的人相近,叫榮耀住在我們的地上.」信徒的責任是讓公義走在上帝的面前,預備祂的道,修直祂的路.我們很清楚知道,我們一切的努力,本身並非是上帝的義,為此,我們不相信地上任何一切的制度,是可以代替上帝的恩典成為人最終的救贖.作為基督徒,我們不應該為塵世任何的制度而背書.塵世一切的制度,亦須受上帝的義所批判,是上帝的義臨近大地,決不是我們的義替代祂的義,這也是我今天的題目《義薄坤輿》的「薄」所要呈現的意思.

從基督教神學而言,在香港的這幾年來,我們強調著社會的關懷,社會上人仕所追求的義,並非能夠等同於基督信仰對於公義的理解.因此,當我們與其他社群同時追求公義的時候,教會當以辯正的態度來作出一些批判,是以一個不卑不亢的態度來表達彼此間的差異.特別在這個現代社會,一個左翼的傳統裡面,他們對現代人的理解,坦白一點的來說,與基督教信仰對人的理解,有很大的差異.他們背後有一個從啟蒙運動帶出來的「人觀」,這就是基督教信仰於十八世紀時與當時候的社會思潮各自分道揚鑣的原因.

當我們爭取社會公義的同時,我們也知道我們站在哪一邊.無論我們對解放神學有多少的欣賞與同情,作為一個基督徒,我們最終所效忠的應該是基督的啟示,而不是馬克斯主義的意識形態.我相信,特別對於參與社運的信徒,我們一分面表達欣賞他們,同時也特別地去提醒他們,要學習分辨的功課──痴言,良語──哪些是出於信仰,哪些是與我們炯然有別.

最後,我以吳經熊先生(蔣介石年代的老人家),他把詩篇翻譯成典雅的文字(我很喜歡他的翻譯),我用詩篇85:13的翻譯作為結束,這譯文與原本的意思有點出入,但對今天的講道而言,它卻是一個很好的總結,也讓這幾句的話放在我們心裡,好讓聖靈來提醒我們,是一個四句的四言詩: 以仁為居,以義為路,康莊大道,眾庶同步.



\newpage

\section{救恩就在平凡處}
\label{sec:133}
\textbf{鄧瑞強博士 - 救恩就在平凡處}
\newline
\newline
\hyperref[sec:132]{< < < PREV SERMON < < <}
~
\hyperref[sec:toc]{[返目錄]}
~
\hyperref[sec:134]{> > > NEXT SERMON > > >}


路加福音 2:22-35
\newline
\begin{longtable}{cl}
\hline
\hline
章節 & 經文 (和合本修訂版)\\
\hline
2:22 & \begin{tabularx}{0.7\textwidth}{X} 按摩西律法滿了潔淨的日子，他們就帶著孩子上耶路撒冷去，要把他獻給主。 \end{tabularx} \\ \\ \relax
2:23 & \begin{tabularx}{0.7\textwidth}{X} 正如主的律法上所記：「凡頭生的男子必歸主為聖」； \end{tabularx} \\ \\ \relax
2:24 & \begin{tabularx}{0.7\textwidth}{X} 又要照主的律法上所說，用一對斑鳩，或用兩隻雛鴿獻祭。 \end{tabularx} \\ \\ \relax
2:25 & \begin{tabularx}{0.7\textwidth}{X} 那時，在耶路撒冷有一個人，名叫西面；這人又公義又虔誠，素常盼望以色列的安慰者來到，又有聖靈在他身上。 \end{tabularx} \\ \\ \relax
2:26 & \begin{tabularx}{0.7\textwidth}{X} 他得了聖靈的啟示，知道自己未死以前必看見主所立的基督。 \end{tabularx} \\ \\ \relax
2:27 & \begin{tabularx}{0.7\textwidth}{X} 他受了聖靈的感動，進入聖殿，正遇見耶穌的父母抱著孩子進來，要照律法的規矩而行。 \end{tabularx} \\ \\ \relax
2:28 & \begin{tabularx}{0.7\textwidth}{X} 西面就把他抱過來，稱頌神說： \end{tabularx} \\ \\ \relax
2:29 & \begin{tabularx}{0.7\textwidth}{X} 「主啊，如今可以照你的話，容你的僕人安然去世； \end{tabularx} \\ \\ \relax
2:30 & \begin{tabularx}{0.7\textwidth}{X} 因為我的眼睛已經看見你的救恩， \end{tabularx} \\ \\ \relax
2:31 & \begin{tabularx}{0.7\textwidth}{X} 就是你在萬民面前所預備的： \end{tabularx} \\ \\ \relax
2:32 & \begin{tabularx}{0.7\textwidth}{X} 是啟示外邦人的光，是你民以色列的榮耀。」 \end{tabularx} \\ \\ \relax
2:33 & \begin{tabularx}{0.7\textwidth}{X} 孩子的父母因論耶穌的這些話就驚訝。 \end{tabularx} \\ \\ \relax
2:34 & \begin{tabularx}{0.7\textwidth}{X} 西面給他們祝福，又對孩子的母親馬利亞說：「這孩子被立，是要叫以色列中許多人跌倒，許多人興起；又要成為毀謗的對象， \end{tabularx} \\ \\ \relax
2:35 & \begin{tabularx}{0.7\textwidth}{X} 叫許多人心裡的意念顯露出來；你自己的心也要被劍刺透。」 \end{tabularx} \\ \\
[1ex]
\hline
\hline
\end{longtable}
日期:2012年1月15日

主題:救恩就在平凡處

經文:路加福音2:22-35

證道:鄧瑞強博士

各位弟兄姊妹,主內平安.

今日,是新的一年的開始,祝願各人在尋索生命意義的路上,又踏前一步.

一個故事:

一個鄉下人,來到城巿.晚上,住在城巿的朋友帶他到繁華的鬧巿逛逛.忽然間,他停下來,帶點好奇地說:「你聽到嗎?」他的朋友說:「聽到什麼?」這位鄉下人說:「你真的聽不到嗎?我聽到蟋蟀的叫聲.真奇怪,在這鬧巿裡,竟然有蟋蟀.」他的朋友說:「你是說笑吧!我只聽到車聲,哪裡有什麼蟋蟀的叫聲.」他們繼續走.忽然間,這次,這位城巿人說:「你聽到嗎?」他的鄉下朋友說:「聽到什麼?」這時,這位城巿人說:「你真的聽不到嗎?我聽到有人丟了一個錢幣.」

人心注意什麼,就會留意什麼,也會聽到什麼.人心像收音機,不同的人調校不同的頻道,就接收到不同的信息.走過一條街,注視美男美女的人,會留意到那裡有美男美女.喜歡美食的人,會留意到那裡有美食.喜歡讀書的人,會留意到那裡有書局.喜歡賽馬的人,隨便你講出哪匹馬,他都能講出這匹馬過去跑了多少次第一.對外人而言,這是匪夷所思的.鑽研股票的人,隨便你講出哪個股票號碼,他都能即時講出這股票的業績.對外人而言,這是匪夷所思的.注目於神的人,卻能在平凡的事物中,看到神的作為.對外人而言,這也是匪夷所思的.但話說回來,這又有什麼匪夷所思呢?專注某事的人,便能看到那事的底蘊.問題只是:你專注什麼?

今日,我們看看一位在平凡事物中看見神作為的人.

講道

經文:路加福音2:22－35

2:22　按摩西律法滿了潔淨的日子,他們帶著孩子上耶路撒冷去,要把他獻與主. 2:23　(正如主的律法上所記:「凡頭生的男子必稱聖歸主；」) 2:24　又要照主的律法上所說,或用一對斑鳩,或用兩隻雛鴿獻祭. 2:25　在耶路撒冷有一個人,名叫西面；這人又公義又虔誠,素常盼望以色列的安慰者來到,又有聖靈在他身上. 2:26　他得了聖靈的啟示,知道自己未死以前,必看見主所立的基督. 2:27　他受了聖靈的感動,進入聖殿,正遇見耶穌的父母抱著孩子進來,要照律法的規矩辦理. 2:28　西面就用手接過他來,稱頌神說: 2:29　主啊!如今可以照你的話,釋放僕人安然去世； 2:30　因為我的眼睛已經看見你的救恩── 2:31-32　就是你在萬民面前所預備的:是照亮外邦人的光,又是你民以色列的榮耀. 2:33　孩子的父母因這論耶穌的話就希奇. 2:34-35　西面給他們祝福,又對孩子的母親馬利亞說:「這孩子被立,是要叫以色列中許多人跌倒,許多人興起；又要作毀謗的話柄,叫許多人心裡的意念顯露出來；你自己的心也要被刀刺透.」

約瑟和馬利亞,抱著嬰孩耶穌,到耶路撒冷的聖殿獻祭.

在這裡,有一個人物出現,他名叫「西面」.經文對他有很正面的評價.他「又公義又虔誠」,並且素常盼望以色列得到安慰,得到救贖.他活在盼望中.他看到人間的逆境與難處,但他知道世間不應是這樣,他對神的作為充滿期待.

路2: 25-27連續三節,都提及聖靈.聖靈在他身上,聖靈啟示他,聖靈感動他(或者說,聖靈引導他).西面活在聖靈的指引下.一個人沉浸在什麼東西裡,久而久之,生命便有該等氣質.在武俠小說裡,高手身上散發著一種「劍氣」.在這裡,西面散發著「靈氣」,他的生命內容是聖靈,他看的東西是聖靈啟示給他看的,他的行動是聖靈引導他行的.

我想,我們的社會很需要這種人.我們的社會,什麼人才也有.我們有很多知識,但很少人向我們啟迪生命的意義.我們有很多技術去操控現實,但很少人教我們生活的智慧.我們有很多資訊,但很少人教我們在資訊中看到神走過的痕跡.

西面來到聖殿,見到嬰孩耶穌,用手接過他來.

我們並不知道西面是否一老人家,但設想他是老人家,也合乎常理.這個老人家抱著一個初生的嬰孩.聽說,老人家抱著自己的孫兒,是最感滿足的.他雖然知道自己生命之將盡,卻也體會到生命的某種延續,他能看到將來,這是令人滿足的.或許,西面抱著耶穌時,也有這種感想.

他稱頌神,講出一段被後世稱為「西面頌」(Nunc Dimittis)的話:

2:29　主啊!如今可以照你的話,釋放僕人安然去世； 2:30　因為我的眼睛已經看見你的救恩── 2:31　就是你在萬民面前所預備的: 2:32　是照亮外邦人的光,又是你民以色列的榮耀.

孔子說過:「朝聞道,夕死可矣.」＜論語‧里仁＞

一個人得聞真理,死可無憾了.

有這種追尋真理的決心,才能至終尋到真理.

西面在這裡表達同樣的心境.能夠看見神的救恩,死可無憾了.

西面,作為一個「又公義又虔誠」的人,仍有不能安然之處.在這個又不公義又不虔誠的世界裡,若不能看見神的救贖,就「死不瞑目」了.這有點像一般的父母,若看不見自己的子女長大成人,獨立自處,便總覺遺憾的.西面,他關心的是社會的光景,民族的前途.他站在城巿的高處,為城巿守望.他等待他的民族得到安慰,得到救贖.唯有這種人,著緊神的事情.唯有這種人,能在平凡事物中,看到神的作為.唯有這種人,能為神對人的救贖,懷有無限感謝.當他看見嬰孩耶穌時,他看到神的救恩,他滿懷安慰.

在聖靈的幫助下,他在這嬰孩生命裡,看見神的救恩,看見神給萬民的光明.一般人,要在神很轟烈的作為裡,才能見到神的救恩,但西面卻只在一嬰孩裡,已看到神的救恩.一般人,要等事情成功了,才會感恩,但西面卻只看到神的作為的起初的一點小火種,已深深感恩.西面,在小事物裡,看到其偉大.

若有一天,你要加班,很晚才放工.回到家裡,有一杯熱的花旗參茶放在你台面.若你的心夠敏感,你會在這微小的事物裡,發現偉大的愛.愛情,不必在你和他母親都墮海,而他卻拼命的救回你們兩人這種大事上顯明的.很多人在你的生命裡,留下很多微小的事物,例如:一句安慰的話,一次誠懇的握手,一個微笑,一個擁抱...,都可以包含神偉大的救贖.西面,在小小的嬰孩耶穌那裡,看到救贖.

一個老人家抱著一個嬰孩,看到救贖,希望,安慰.這真是一幅美麗的圖畫.我在神學院教書,面對很多學生.我雖然不是老人家,我的學生也不是嬰孩,但若學生們不怕辛勞地學習,不怕吃虧地服事,不隨波逐流而能注目於神,不被罪惡誘惑而精進向前,我也會覺得,他們能將救贖,希望,安慰帶給他們將來牧養的教會.而我,也死而無憾.

抱著耶穌,西面看得很遠.他看到這個帶來救贖的耶穌,會引起一場人性的革命.

2:34-35　西面給他們祝福,又對孩子的母親馬利亞說:「這孩子被立,是要叫以色列中許多人跌倒,許多人興起；又要作毀謗的話柄,叫許多人心裡的意念顯露出來；你自己的心也要被刀刺透.」

只要一個人活出了真正的人性,他就確立一人性的高度.其他人在這人性的標誌竿下,其人性的高低便顯明出來.很多醫生參與「無國界醫生」(Medecins Sans Frontieres)這個組織的工作,他們的專業應該可以為他們在很多社會裡賺取地位與金錢,但他們選取了義務的救援工作.他們選擇活出意義,多於考慮物質的回報.他們的生命,叫我們反省,我們的存活方式,會讓我們失去什麼.

耶穌將來的生命,將會成為人生路上的路標.他指出的方向,引人走向天父的懷裡.走相反的方向,令人遠離天父.面對這個路標,我們不得不選擇.這個路標指示的方向,是叫人肩負生命的責任,而不是不負責任；叫人愛人如己,而不是將身邊的人視為潛在的敵人；叫人不要將自己視為神,而是將眼目注目永恆的神.對很多人而言,這是艱難的.對很多人而言,這是十字架.對很多人而言,十字架只意味著失敗.如經文所說,耶穌成為他們「毀謗的話柄」,他們反對耶穌這個生命的路標.

若果生命只有一條路,我們只需要直走下去便行了.但生命有分叉路,耶穌在分叉口作為一個路標,指向良善的路.相反,便是惡路了.在這裡,我們不得不選擇.如何選擇,便顯明我們心中真正的意念了.我們是向善的呢?抑或向惡的呢?未到分叉口,我們不用抉擇.但在分叉口,我們的意念便顯明出來了.有時,你去買東西,售貨員找多了錢.你完全知道,他找多了錢.這時,我們便到達生命的分叉口了.耶穌在這裡作為一路標,指向良善的路.但是,我們是否必然走向良善的路呢?在這路口,我們的心念便顯明了.

西面最後對馬利亞說:「你自己的心也要被刀刺透.」一般的理解,是認為這是指耶穌將來的命運會叫他母親心如刀割.但是,現代的釋經學者認為,這是指馬利亞的生命,也要被審判的劍剖示.這即是說,雖然馬利亞和耶穌有特別的關係,但是,她是否走向善路,是否走向天父,是否走向永生,仍是懸而未決的.她是否合乎人性的標準,是否走向善路,仍不得不被耶穌這標誌竿及路標所衡量.沒有人,可以單憑外在的條件,能走回天父的懷裡.我們不能單單因為返教會而屬於神,我們不能單單因為奉獻了一定數額的金錢而屬於神,我們屬於神,是因為我們活出耶穌的生命.

西面,抱著耶穌,他看到這是神救贖計劃的開始,也看到這救贖計劃對人性帶來的挑戰.他感到滿足,也為人的反應擔憂.聖誕節過了,很有節日氣氛,我們都很快樂.然而,耶穌來了,他對我們生命作出的挑戰,同樣叫我們快樂嗎?

一個故事:(以下故事的英文原文:\url{http://bible.org/illustration/final-bid})

在英國,有個很富有的男爵,叫費茲哲羅(Baron Fitzgerald).他太太很早過身,留下一個兒子.後來,兒子也病逝.他將大量的財富,用來購買「名家」的藝術品.

很來,他死了.他立下遺囑,將一切藝術品拍賣.由於他藏品豐富,品質也高,吸引了很多人到拍賣會,準備競投.

拍賣前,所有藝術品都展示出來.其中一幅畫,幾乎沒有人表示興趣.那是一個不出名的畫家的作品,品質粗劣,畫的是費茲哲羅的獨生子的肖像.

拍賣開始了,拍賣官請律師宣讀費茲哲羅的遺囑,遺囑說:第一幅被拍賣的畫是「我最鍾愛的兒子」.

沒有人出價買這幅粗劣的畫.沒有人出價.除了一個人,一個老人家.他是費茲哲羅家中的老僕人,他認識畫中的人,他服事過他,他愛他.基於感情因素,他出價購買.但因他不富裕,他只能出價一英磅.但由於沒有其他人願意出價,最後,這老僕人買下這幅畫.

接著,拍賣官請律師繼續讀費茲哲羅的遺囑.遺囑說:「誰買下我兒子的畫像,誰就獲得我全部的收藏品.拍賣完畢.」

誰願意擁有神的兒子,將擁有天父給予人的一切福樂.「拍賣完畢」.

但願榮耀歸於聖父,聖子,聖靈.阿們.



\newpage

\section{禱告進階聆聽上主}
\label{sec:134}
\textbf{賀志勇博士 - 禱告進階}
\newline
\newline
\hyperref[sec:133]{< < < PREV SERMON < < <}
~
\hyperref[sec:toc]{[返目錄]}
~
\hyperref[sec:135]{> > > NEXT SERMON > > >}


詩篇 29:1-11
\newline
\begin{longtable}{cl}
\hline
\hline
章節 & 經文 (和合本修訂版)\\
\hline
29:1 & \begin{tabularx}{0.7\textwidth}{X} 神的子民哪，你們要將榮耀、能力歸給耶和華，都歸給耶和華！ \end{tabularx} \\ \\ \relax
29:2 & \begin{tabularx}{0.7\textwidth}{X} 要將耶和華的名的榮耀歸給他，要敬拜神聖榮耀的耶和華。 \end{tabularx} \\ \\ \relax
29:3 & \begin{tabularx}{0.7\textwidth}{X} 耶和華的聲音在眾水上，榮耀的神打雷；耶和華打雷在大水之上。 \end{tabularx} \\ \\ \relax
29:4 & \begin{tabularx}{0.7\textwidth}{X} 耶和華的聲音大有能力，耶和華的聲音滿有威嚴。 \end{tabularx} \\ \\ \relax
29:5 & \begin{tabularx}{0.7\textwidth}{X} 耶和華的聲音震碎香柏樹，耶和華震碎黎巴嫩的香柏樹。 \end{tabularx} \\ \\ \relax
29:6 & \begin{tabularx}{0.7\textwidth}{X} 他使黎巴嫩跳躍如牛犢，使西連跳躍如野牛犢。 \end{tabularx} \\ \\ \relax
29:7 & \begin{tabularx}{0.7\textwidth}{X} 耶和華的聲音使火焰分岔。 \end{tabularx} \\ \\ \relax
29:8 & \begin{tabularx}{0.7\textwidth}{X} 耶和華的聲音震動曠野，耶和華震動加低斯的曠野。 \end{tabularx} \\ \\ \relax
29:9 & \begin{tabularx}{0.7\textwidth}{X} 耶和華的聲音驚動母鹿落胎，樹林也脫落淨光。凡在他殿中的，都述說他的榮耀。 \end{tabularx} \\ \\ \relax
29:10 & \begin{tabularx}{0.7\textwidth}{X} 耶和華坐在洪水之上為王；耶和華坐著為王，直到永遠。 \end{tabularx} \\ \\ \relax
29:11 & \begin{tabularx}{0.7\textwidth}{X} 耶和華必賜力量給他的百姓，耶和華必賜平安的福給他的百姓。 \end{tabularx} \\ \\
[1ex]
\hline
\hline
\end{longtable}
日期:2012年1月22日

題目:禱告進階:聆聽上主

證道:賀志勇博士

讀經:詩篇29:1-11

1 神的眾子啊,你們要將榮耀能力歸給耶和華,歸給耶和華. 2 要將耶和華的名所當得的榮耀歸給他,以聖潔的妝飾(註:"的"或作"為")敬拜耶和華. 3 耶和華的聲音發在水上,榮耀的神打雷,耶和華打雷在大水之上. 4 耶和華的聲音大有能力；耶和華的聲音滿有威嚴. 5 耶和華的聲音震破香柏樹；耶和華震碎黎巴嫩的香柏樹. 6 他也使之跳躍如牛犢,使黎巴嫩和西連跳躍如野牛犢. 7 耶和華的聲音使火燄分岔. 8 耶和華的聲音震動曠野；耶和華震動加低斯的曠野. 9 耶和華的聲音驚動母鹿落胎,樹木也脫落淨光；凡在他殿中的,都稱說他的榮耀. 10 洪水泛濫之時,耶和華坐著為王；耶和華坐著為王,直到永遠. 11 耶和華必賜力量給他的百姓,耶和華必賜平安的福給他的百姓.

一,引言 (跟父親的溝通)

我發現同我老爸溝通很不容易,我說東,他就說西；我跟他講道理,他又一聲不吭,我根本不知道他到底在想什麼.打個比方說,有一次我跟我老爸打電話說:「老爸,你血脂高,我同你買些螺旋藻吃吧.」我老爸就說:「太貴了,不要.」我就說:「我都沒告訴你價錢,你就說太貴了；其實不貴的.」他還是說:「太貴了,不要.」我就著急了,跟他說:「我都已經買好了.」他還是說:「太貴了,不要.」 說得我一點辦法都沒有了.我發現,我太理性,不懂得跟老人家溝通.

跟地上的爸爸溝通不容易,我發現,跟天上的爸爸溝通也不容易.有時候真不知道跟天父說什麼好,翻來覆去都是那幾句話,總是求祂醫治,求祂帶領,求祂改變一下某個剛硬的人的心思意念,好像祂就是個大能的Helper,隨時聽孩子們的祈求,隨時準備出手救火.

還有一個問題,那就是向神禱告了很多次,我又覺得好多事神好像基本沒有什麼直接回應.祂從來不給我打電話,也從來不給我發電郵；祂好像從來沒有托一個夢給我,也不在我腦海裡說話.當年天使會給瑪利亞報信,現在信差換成了我們這些凡人,年年去報佳音的都是我們這幫凡人.我有時候一邊禱告一邊想,天父啊,你有沒有在聽我禱告呢?尤其是去年好朋友得了癌症,我們一幫同學不住替她禱告,可最後她還是去世了.我想做牧師做傳道人最慘了,因為有一件事情很痛苦,就是發現自己其實很多事都無能、無力、軟弱無助.很多次去探望一些重病,或者垂死的病人,我都不知道怎麼去安慰他.我告訴他說,上帝會看顧你和你的家人,上帝的旨意是美好的--我這樣講,自己都覺得自己不切實際.

好多事情,我都希望神親自來指導,就算不派天使來報信,起碼也托個夢、給個明確的消息吧.可是沒有!每一次次禱告之後,要麼是心想事不成,要麼是結果跟自己想要的貨不對版.無論什麼結果,我都要告訴自己:神給我的是最好的,最適合我的.嘴裡是這麼說,心裡卻半信半疑.我常常會想起年老的約伯,他真的好無奈啊.其實約伯這個名字原文的意思就是「爸爸,你在哪裡啊」,有時候我也會這麼問:天父,你在哪裡啊? 到底要怎麼禱告,到底要怎樣跟神溝通,才能聽到神的回應呢?或者讓我們一起來看看詩篇第二十九篇.

二,解經

1, 1-2節 將榮耀歸神

詩篇篇篇都是人向神禱告的記錄,所謂禱告,就是人向神說話,並聆聽神的說話.而詩篇 29 篇,就描繪了聽神說話的幾個重要原則.

在1-2節中,詩人宣召眾天使都來讚美神,將榮耀歸給神.這讓我想起了主禱文,主禱文第一句就說:「

我們在天上的父,願人都尊你的名為聖.」原來禱告的第一步就是要把榮耀尊貴都歸給神.

榮耀是什麼意思?是不是頭上有個光環,或者穿金戴鑽石,搞得全身發光那樣子. 不是,榮耀,原文的意思就是「彰顯」,「顯明」,歸榮耀給神,就是我們的所作所為,要把上帝顯明出來.但我們的禱告,一開始往往不是彰顯上帝,一開始往往是彰顯我們的需要.求神讓孩子進一家好中學；求神讓老闆不要太嚴厲苛刻,求神讓我早點拿到博士學位,求神讓我身體健康,跑得、跳得、食得、玩得.這些都不是壞事,這些都可以求,但看我們是以什麼心態來求.撒母耳的母親哈拿沒有兒女,被丈夫的另一個有兒女的妻子常常氣她.她向神求一個兒子,不再被人氣,但她還加了一句:「

我必使孩子終生歸於耶和華.」(撒1:11)神就應許了她.

向神禱告,態度很重要.我們可以求很多東西,但耶穌告訴我們,「

你們要先求祂的國和祂的義.」(太6:33)什麼叫先求祂的國和祂的義?是不是要我們每次禱告都要先求世界和平?不是,國與義是神給我們立的規則法度,是每個基督徒的終極理想,它就記載在聖經裡.也就是說,當我們求那些日常的東西時,我們要明白,我們所求的東西是跟「國和義」的目標是一致的.我們之所以求那些東西,是因為我們要通過那些日常的小事,經歷神的同在,與神同行.如果求那些東西,僅僅為了滿足我們自我中心的需要,神就會告訴我們:「

你們求也得不著,因為你們妄求」.(雅各書4:3)

你或者會抗議說:小勇哥,我求的可是好東西,我經常求神讓我自律,不要貪玩打遊戲、不要貪食、求神讓我有愛心、有忍耐之心、寬容之心,不要對我的太太孩子發脾氣,但神從來不幫我,我從前怎樣,現在還是怎樣.我說「

你得不著,是因為你不求」(雅各書4:3),因為你雖然也禱告,但其實你並不太相信這個禱告,你並不信這個禱告真的管用!你禱告完,然後你的生活行事還像從前一樣.這就叫「星期天基督徒」.他們禮拜天來做禮拜,然後其他六天活得像無神論者一個樣.他們也禱告,但是禱告完了之後,不是配合神去行動,而是繼續按自己的方式去行動.禱告只是一種儀式,一種情緒的發洩,他們的理性告訴自己:神不會管這些小事的,還是讓我們自己動手吧!

當我們的禱告是以這些自我需要、自我性格習慣為中心時,當我們的禱告首先和首要的東西不是彰顯神,榮耀神時,我們又怎麼能聽到神微小的聲音呢?

2,3-5節 神在大自然中說話

問題在於我們對於人的聲音過分敏感,而對於神的聲音過分不敏感了.我記得有一次我跟一些學生去海邊玩,海天一色,海鷗飛翔,我說:「上帝的造物多美妙啊!」我的學生繼續玩他的 iPhone,然後說:「小勇哥,讓一下,我想拍那只鳥!」 同樣是美景,我想到讚美神,他想到是用 iPhone 拍只大鳥,發在 Facebook 裡展示一下,告訴大家他到此一遊.我們聽到的聲音好不同啊.

但在詩篇 3-5 節,詩人告訴我們,耶和華的聲音在大水之上,在雷聲之中,在香柏木震裂時,在整個自然裡面.哪怕是野地的花,空中的飛鳥,都會傳揚祂的聲音.我們離真正的自然太遠了!以前我帶了一群香港的中學生回我鄉下,他們看到河邊上活的雞鴨,活的牛羊,興奮得不得了,都沖過去追,追得雞飛狗跳.很多學生說,他們見過炸子雞,白切雞,從來沒見過活雞.其實我們自己也是,在這個城市中,我們有多少時候聽見鳥鳴,狗吠,牛麼麼叫的聲音呢,有多少時候能聽見風吹過樹林的聲音呢,還有溪流的聲音呢?我們聽見的是汽車聲,火車聲,電視手機聲,冷氣機轟隆隆的聲音,當然還有,各種各樣好聽不好聽的人的聲音.我在火車裡經常看見人帶著耳機,一個人搖頭擺尾的,我還以為他在自言自語了,後來才發現是在打電話.我們也是這樣,如果我們被自己的心思意念遮蔽,就像戴了耳機那樣,除了我們電話裡的聲音之外,我們還能聽見什麼呢?"這日到那日發出言語",但是我們聽不見.

我有一位朋友,他說他從小在大埔長大,每次他心情煩惱,或者要做一些重大決定時,他發現他在家裡完全沒法安靜.他就回大埔,在海邊慢慢走,什麼話都不說,什麼都不想,只是聆聽大自然的聲音,慢慢走,慢慢走,直到內心安靜,一切都豁然開朗為止.

我們有種錯覺,覺得禱告就是要說話,要不停的說,要引用很多聖經的話語,說得越多越屬靈,其實未必.有時候我們說得太快太多,前面究竟說了什麼,禱告完都忘記了.禱告不是聖經知識問答,也不是上帝的隨堂測驗,看看我們到底能引用多少金句.禱告是我們全人全心的跟神交流,它不一定要說話；禱告也可以是倒空自己,靜靜聆聽.靈修神學家 Joyce Huggett 提醒我們:我們把天主教"默觀"的功夫全部丟掉了.她說她本人是福音派信徒,比較理性,但有一次,她到一個天主教修道院靜修,那個地方非常自然,非常幽靜.她跟一群修女一起用餐,大家默默無聲地吃完一餐飯；然後一群人一起禱告,卻沒有人開口說一句,她說她從來沒經歷過吃飯、禱告不說一句話的.一切都非常安靜.Huggett 說:「安靜的力量太讓人震驚了!」在安靜中,她感覺到自己的心開始鬆軟起來,越來越松,不再那麼緊了.在這個大自然當中,她感覺到有新的種子播種在她的心上.在那一刻,她聽清了神的話語.弟兄姊妹,這是不是我們可以學習的一種好方法呢?到自然中去,安靜地等候神的聲音.

3,6-9節 神在以色列歷史中說話

當然,神說話不是單聲道的,它是立體聲的；他不僅透過自然來說話,他還有很多講話途徑.本篇6-9節講到神的聲音使火焰分開,在加低斯曠野震動,這些表面講的是景,但實際暗含著以色列的歷史.黎巴嫩、西連、火焰分開、加低斯的曠野,都是以色列人深深熟悉的景象.

讀過舊約的人就知道,加低斯的曠野,就是當年摩西領著以色列人漂流了四十年的曠野.正是在這個曠野,以色列人向神求食物,神每天降下瑪哪,但不允許收藏過夜,但偏偏有人就要收藏過夜,結果第二天瑪哪就壞了.神的話語明明講給以色列民聽了,他們卻還是要建造金牛犢來敬拜.以色列民就是這樣,明明看見了神的大能,聽見了神的話語,可是一有不滿足就發怨言.神向他們說話,但是他們聽到了,卻聽不懂.但摩西並不這樣,他經常禱告,跟神溝通.當神因為金牛犢事件要滅絕以色列民的時候,摩西「

像從前一樣俯伏在耶和華面前四十晝夜,沒有吃飯,也沒有喝水.」(申命記9:18)然後再跟神講道理求情.摩西懂得對神敬畏、謙卑、降服.他的行動就是他的禱告,他懂得安靜地向神勸說,俯伏在神面前四十晝夜,而不是滔滔不絕地說個不停.更不像以色列民,一有不滿就向神發怨言.

加低斯的曠野,也是耶穌經常禱告的地方.馬可就記載耶穌清晨去曠野禱告；馬太跟路加都記載了耶穌被聖靈引領在曠野禁食禱告四十晝夜--大家發現沒有,又是一個四十晝夜,就好像摩西謙卑俯伏在神面前四十晝夜.當耶穌禱告四十晝夜之後,魔鬼的聲音就來了,都是些名聲、權聲、利聲.耶穌勝過這些聲音後,天使就來伺候他了.無論是摩西在加低斯曠野的禱告,還是耶穌在加低斯曠野的禱告,我們都忽視了一位隱藏的主角--聖靈!Llewellyn 指出,只有聖靈能教導我們禱告,也保守我們聽到神的話語.因為祂是中保,是仲介.

像我這樣的現代科學理性人,談聖靈的時候總會有點尷尬.人家一說神跡奇事、 聖靈醫治,我就皺眉頭,覺得他是不是靈恩派的人啊?有一次,我的一個好朋友悄悄跟我說,他會說方言,我看了他半天,確定他神智還清醒,心裡才有了點小變化.最近這幾年,經常能聽人分享神怎麼醫治了她,我都是笑笑說:「你真有福氣!」我總是半信半疑.但最近張伯笠牧師來訪,讓我震動了.他親身講述了他的經歷:他得了癌症,化療後腎臟衰竭,沒得救了,但是他憑著禱告,終於感動神出手醫治了他.以前看他的書,只是覺得「他真有福氣」,現在跟他真人接觸多了,才覺得書上寫的是活生生的東西.我發現,當我在學習放下理性,敞開自己的心的時候,我才真的感覺到緊繃的心開始鬆動了,變軟了,我才真正感覺到了神的話語的力量.

加低斯的曠野,是讓我們放下自己的纏累,敞開自己,靠神的靈引導去禱告的地方.每個人都需要一個這樣的曠野,去靠著聖靈聽到神的話,聽懂神的話.神的話有時候是隱藏的,恐怕你們有耳朵就聽見了!  結論:10-11節 倚靠聖靈等候神的話語

本篇10-11 節很像主禱文最後兩句.詩篇29篇最後兩句不僅是讚美神,同時也體現了神的子民對神的信心.今年是2012 年了,這種信心,在這個經濟動盪的年代尤其重要.我好多朋友都跟我說,2012了啊,你的方舟船票買好了沒有?

有人說這是蘋果的年代,是 Facebook 的年代,這標誌著人們更喜歡表現自己、更喜歡講話,卻不再注意去聆聽--我們很可能逐漸在喪失聆聽的能力.但詩篇 29 篇卻提醒我們,我們要放下我們的自我中心,先歸榮耀給神.敞開自己,讓我們被纏累的心先松下來,軟下來,靜靜去等候神的聲音.很有可能我們會像以色列民一樣,漂流四十年才聽到答復；也很有可能神通過自然向我們說話,或者通過聖經中的歷史向我們說話,但請記住,當我們的心安靜下來等候時,「

聖靈會親自用說不出的歎息替我們禱告」(羅8:26),到那時,我們只有一項選擇:順服聖靈的感動,那時候,神的聲音會格外清晰.



\newpage

\section{知遇之恩}
\label{sec:135}
\textbf{林豪恩先生 - 知遇之恩}
\newline
\newline
\hyperref[sec:134]{< < < PREV SERMON < < <}
~
\hyperref[sec:toc]{[返目錄]}
~
\hyperref[sec:136]{> > > NEXT SERMON > > >}


約翰福音 1:43-51
\newline
\begin{longtable}{cl}
\hline
\hline
章節 & 經文 (和合本修訂版)\\
\hline
1:43 & \begin{tabularx}{0.7\textwidth}{X} 又過了一天，耶穌想要往加利利去。他找到腓力，就對他說：「來跟從我！」 \end{tabularx} \\ \\ \relax
1:44 & \begin{tabularx}{0.7\textwidth}{X} 這腓力是伯賽大人，是安得烈和彼得的同鄉。 \end{tabularx} \\ \\ \relax
1:45 & \begin{tabularx}{0.7\textwidth}{X} 腓力找到拿但業，對他說：「摩西在律法書上所寫的，和眾先知所記的那一位，我們遇見了，就是約瑟的兒子拿撒勒人耶穌。」 \end{tabularx} \\ \\ \relax
1:46 & \begin{tabularx}{0.7\textwidth}{X} 拿但業對他說：「拿撒勒還能出甚麼好的嗎？」腓力說：「你來看。」 \end{tabularx} \\ \\ \relax
1:47 & \begin{tabularx}{0.7\textwidth}{X} 耶穌看見拿但業向他走來，就論到他說：「看哪，這真是個以色列人！他心裡是沒有詭詐的。」 \end{tabularx} \\ \\ \relax
1:48 & \begin{tabularx}{0.7\textwidth}{X} 拿但業對耶穌說：「你從哪裡認識我的？」耶穌回答他說：「腓力還沒有呼喚你，你在無花果樹底下，我就看見你了。」 \end{tabularx} \\ \\ \relax
1:49 & \begin{tabularx}{0.7\textwidth}{X} 拿但業回答他：「拉比！你是神的兒子，你是以色列的王。」 \end{tabularx} \\ \\ \relax
1:50 & \begin{tabularx}{0.7\textwidth}{X} 耶穌回答他說：「因為我說在無花果樹底下看見你，你就信嗎？你將看見比這些更大的事呢！」 \end{tabularx} \\ \\ \relax
1:51 & \begin{tabularx}{0.7\textwidth}{X} 他又說：「我實實在在地告訴你們，你們將要看見天開了，神的使者在人子身上，上去下來。」 \end{tabularx} \\ \\
[1ex]
\hline
\hline
\end{longtable}
日期:2012年1月29日

主題:知遇之恩

經文:約翰福音1:43-51

證道:林豪恩先生

經文

43　又次日,耶穌想要往加利利去,遇見腓力,就對他說:「來跟從我吧.」

44　這腓力是伯賽大人,和安得烈,彼得同城.

45　腓力找著拿但業,對他說:「摩西在律法上所寫的和眾先知所記的那一位,我們遇見了,就是約瑟的兒子拿撒勒人耶穌.」

46　拿但業對他說:「拿撒勒還能出甚麼好的嗎?」腓力說:「你來看!」

47　耶穌看見拿但業來,就指著他說:「看哪,這是個真以色列人,他心裡沒有詭詐的.」

48　拿但業對耶穌說:「你從哪裡知道我呢?」耶穌回答說:「腓力還沒有招呼你,你在無花果樹底下,我就看見你了.」

49　拿但業說:「拉比,你是上帝的兒子,你是以色列的王!」

50　耶穌對他說:「因為我說「在無花果樹底下看見你」,你就信嗎?你將要看見比這更大的事」；

51　又說:「我實實在在地告訴你們,你們將要看見天開了,上帝的使者上去下來在人子身上.」

一．聞名不如見面 剛過去的聖誕節,有十多位弟兄姊妹接受洗禮.其中一位在洗禮見證集如此分享:「當自己得到這麼美好的東西,實在是忍不住要向自己所愛的人介紹.於是我開始火熱地向家人和朋友傳福音,而且是毫不客氣.可是讓我難過和悲痛的是,這個世界真的是拒絕福音,為什麼會是這樣?就連我的家人也是.我的父母和朋友以為我遇到什麼難過的坎,所以尋求寄託,安慰我說不要太認真.當我執意想說服他們時,他們就一下子爆怒了,爸媽大聲訓斥我腦子壞了,也不自己努力,天天靠什麼神,命令我不許去教會了.」

類似的遭遇,耶穌的門徒腓力也經驗過.

43　又次日,耶穌想要往加利利去,遇見腓力,就對他說:「來跟從我吧.」

44　這腓力是伯賽大人,和安得烈,彼得同城.

45　腓力找著拿但業,對他說:「摩西在律法上所寫的和眾先知所記的那一位,我們遇見了,就是約瑟的兒子拿撒勒人耶穌.」

46　拿但業對他說:「拿撒勒還能出甚麼好的嗎?」

對於心裡火熱,忍不住要向好朋友介紹耶穌的腓力,拿但業的反應真是「一盤冷水照頭淋」.腓力所發現的「至寶」,拿但業竟然當作「禾稈草」,用諷刺的口吻回應:「

拿撒勒還能出甚麼好的嗎?」

你也遭遇過類似的回應嗎?

「信耶穌?耶穌俾飯你食呀?」

「我們是中國人,為甚麼要信一個西方的宗教?」

「信耶穌?不如信自己啦!信基督教?不如信訓覺啦!」

這些回應隱含的質疑是:「耶穌根本不認識我,耶穌與我沒有關係,也不知道我的需要.」

曾幾何時,我們也以類似的態度回應向你介紹耶穌的人,是嗎?

拿但業對耶穌的態度何時改變呢?是甚麼令拿但業改變對耶穌的態度呢?是因為他的朋友腓力能言善辯,說服了拿但業嗎?看來不是,腓力對於拿但業把他的「至寶」貶抑為「禾稈草」,並沒有動氣,也沒有反駁,只對他說:「你來看!」.

二． 知我者莫若你

47　耶穌看見拿但業來,就指著他說:「看哪,這是個真以色列人,他心裡沒有詭詐的.」 48　拿但業對耶穌說:「你從哪裡知道我呢?」耶穌回答說:「腓力還沒有招呼你,你在無花果樹底下,我就看見你了.」

腓力邀請拿但業去看耶穌,其實不是拿但業看見耶穌,而是耶穌看見拿但業.拿但業對耶穌態度的改變,不是因為他看透了耶穌,而是他被耶穌看透了,以至他驚訝地問耶穌:「你從哪裡知道我呢?」

震撼拿但業心靈的,是耶穌竟然看出他是一個怎樣的人.常言道,相識滿天下,知己無一人.人們見面交換名片,所認識的,無非是名字和職銜,有多少人願意了解你的內心和品格呢?有多少人有興趣了解你的職銜所承擔的任務及因此而所承受的壓力呢?

不少人如拿但業一樣疑質耶穌,因為他們以為耶穌與他們無關.同樣,不少人好像拿但業一樣改變對耶穌的態度,因為他們驚訝地發現,原來耶穌認識他們,而且耶穌對他們的認識,超乎想像的深入與私下.在剛才提及的洗禮見證集裡有另一位受洗者如此分享他對耶穌的態度轉變的經驗:

「向來好勝的我．．．忽然感到自己已經奔走得很疲累了.打從初執教職開始,一直擔任行政職務,還不斷進修,三十三歲當了學校主管,抱理想,承重責,期間勞累與壓力,不言而喻.事業雖非朝夕可期,但我自恃堅毅奮勇,並不介意一力承擔.當晚佈道會的那位高級行政人員,分享了他在職場的成長歷程:崗位上的孤單,責任的重擔,業績的逼迫,人事的糾纏......他問自己:我是何苦由來?信主後,他才認識到,所謂事業,不過是主人交付給僕人的工作.若那是管家的工作,自然便要按時分糧給其他僕人.「

因為多給誰,就向誰多取；多託誰,就向誰多要.」(路加福音12章48節)從幼及長,就是因為要做自己的主人,叫我勞累孤單.其實,沒有上帝,再努力也是徒勞；因為「

栽種的,算不得甚麼,澆灌的,也算不得甚麼；只在那叫他生長的神.栽種的和澆灌的,都是一樣,但將來各人要照自己的工夫得自己的賞賜.」(哥林多前書3章7-8節)一下子,我感到輕鬆了很多,並且聽到很溫暖的一句呼喚:「

凡勞苦擔重擔的人可以到我這裡來,我就使你們得安息.」(馬太福音11章28節)那天晚上,我決志了.」

這位弟兄的心靈被震撼,因為他感受到耶穌認識他──不單認識他的職位,不單認識他的成就,更了解他高處不勝寒的孤單,任重而道遠的疲累,業積的壓力及人事的困擾.

中國古人說:「士為知己者死,女為悅己者容」.拿但業這位質疑者,震驚於耶穌竟然看得見他「心裡沒有詭詐」的品格,從此跟隨耶穌.

現代人不缺乏朋友,連足不出戶的,都有幾百位朋友,甚至要多少就可以有多少,只要登入Facebook add朋友的朋友就可以了.雖然有那麼多朋友,甚至不乏傾慕者,但人們渴求的,仍然是一個懂自己的人.電影《那些年,我們一起追的女孩》的女主角沈佳宜,雖然享受被男主角柯景騰追求,然而,令觀眾感到遺憾的,是他最終沒有選擇與青梅竹馬的柯景騰一起.沈佳宜知道自己終生交托的,不單單是一位喜歡自己的人,而且是一位懂自己的人.柯景騰雖然對沈佳宜一往情深,但他並不懂沈佳宜的心.

徐志摩這位浪漫詩人一句「悄悄的我走了,正如我悄悄的來,我揮一揮衣袖 不帶走一片雲彩」擄掠了不少少女的心,其中一位是十六歲的林微音.林微音雖然一度被徐志摩的性格,熱忱和他對自己的狂戀所吸引,但對於徐志摩的熱烈追求,並沒有引起對等反應.為甚麼連浪漫才子的追求也不接受呢?林微音後來說:「徐志摩當時愛的並不是真正的我,而是他用詩人的浪漫情緒想像出來的林徽音,可我其實並不是他心目中所想的那樣一個人.」

人們其實都不希望別人以他們的想像來塑造自己,也不想別人單單以自己的成敗得失來衡量自己,人們心底裡都渴望著真正的自己能夠被了解,知道自己是是一個怎樣的人.

古人管仲說:「生我者父母,知我者鮑子也」.何以管仲對鮑子牙有這樣的評價呢?他說:「我貧困時,曾與鮑叔合夥經商,分錢時自己分得多,但鮑叔不認為我貪婪,因為他知道我貧窮.我曾為鮑叔辦事,結果越辦越窘迫,鮑叔不認為我愚昧,因為他知道時機有利與不利之別.我曾經數度當官,但都被國君斥逐,鮑叔不認為我不肖,因為他知道我遭遇不利的時運.我曾經參與戰爭,結果三戰三逃,鮑叔不認為我膽怯,因為他知道我家有老母.公子糾失敗之後,召忽自殺,我則幽囚受辱,鮑叔不認為我無恥,因為他知道我不羞小節,而以功名不能顯揚於天下為恥辱.」鮑子牙如此知道管仲,真的可稱為管仲的「知己」.

其實,很多人都期望自己如管仲一樣,有人如鮑子牙般看見真正的自己.

有一位丈夫做錯了一個投資決定,他的妻子幾十年來仍然指責他「無腦」,他等待著有人了解他真正的動機是為了改善家人生活,等待著有人體諒他偶然的「失手」.有一位父親在外地工作,而且與太太的關係欠佳,因此被妻兒子女視為「不顧家」,他等待著子女有一天知道他離家原是為了家,等待著子女認識他並非拋妻棄子之輩,等待著子女體諒自己有家歸不得之苦.有一位職員較為健忘,有時因為忘記上司的吩咐而被指責「無心工作」,他等待著有人知道他對自己的工作別無異心,等待著有人體諒他只是一時忘記而己.有一位學生在課堂上回應比較慢,被老師和同學以為不知道答案,或者不投入課堂學習,他等待著有人了解他需要仔細思考,確定答案後才發表意見的性格,等待著有人體諒他被反應快的同學奪去了參與機會的無奈.

耶穌的入室弟子西門彼得曾經跪在耶穌面前說:「

主阿,請你離開我吧,我是個罪人.」(路5:8)耶穌當然知道彼得是個罪人,彼得是一個怎樣的罪人,犯過甚麼罪,耶穌都清楚知道.然而,耶穌給彼得的回應是:「

不要怕,從今以後,你要得人了.」在耶穌被審判的那一夜,彼得為自己三次不認主而痛哭.耶穌並沒有因此而誤會他為「叛徒」.「你是彼得,是磐石；在這磐石上,我要建立我的教會,甚至連死亡的權勢也不能勝過它.」(馬太16:18)彼得可能不知道耶穌怎能對自己如此不離不棄,可能不明白耶穌為何如此看得起自己.事實上,耶穌認識彼得甚至比彼得認識自己更透澈,耶穌知道彼得的心思甚至比彼得自己更深入.

「

腓力還沒有招呼你,你在無花果樹底下,我就看見你了.」是的,其他人可能看不見你,就算看得見你,都難免夾雜了他們的想像和期望,甚至,有時連你自己也看不見自己,然而,耶穌看得見你,並且看得見真真實實的你.被看見,是一種心靈的震撼和欣慰,這種心靈的震撼和欣慰,改變了拿但業,也改變了世上眾多經驗到自己被耶穌看見的人.

三．難得知己,願意知彼嗎?

49　拿但業說:「拉比,你是上帝的兒子,你是以色列的王!」 50　耶穌對他說:「因為我說「在無花果樹底下看見你」,你就信嗎?你將要看見比這更大的事」； 51　又說:「我實實在在地告訴你們,你們將要看見天開了,上帝的使者上去下來在人子身上.」

耶穌竟然看得出自己是一個怎樣的人,拿但業對耶穌的看法也立即改觀.他對耶穌的評價由「

拿撒勒還能出甚麼好的嗎?」變為「

你是上帝的兒子,你是以色列的王」.耶穌沒有否認拿但業對他的認識:「

上帝的兒子,以色列的王」,不過,耶穌告訴他這樣的認識只是一個開始,精彩的好戲尚在後頭.

大家還記得在洗禮儀節中自己所作的認信嗎?

問:你是否相信永生之子耶穌基督? 答:我信. 問:你是否承認耶穌為主,為拯救者,並決定跟從耶穌? 答:我承認.

拿但業正確地宣告了耶穌是「上帝的兒子,以色列的王」這個身份,然而,這個身份的實質內涵是甚麼呢?我們認信耶穌為「主」的地位,認信耶穌是「拯救者」的身份,同樣地,作為我們的「主」,作為我們的「拯救者」的實質內涵是甚麼呢?耶穌並沒有在當下向拿但業定義「上帝的兒子,以色列的王」這個身份,事實上,任何定義都不能夠涵蓋這個身份,而且任何定義都有限制這個身份的意味,唯一能夠體會這個身份的精彩之處,耶穌的建議是來「看」,唯有拿但業願意繼續來看,他才有機會「看見比這更大的事」；唯有拿但業願意來看,他才會有機會「看見天開了,上帝的使者上去下來在人子身上」；唯有拿但業願意繼續看,他才會認識到耶穌不單單是一個身份,而是天地之間,神人之間的溝通.天地之間的溝通展開了就沒有停止,神人之間的溝通開始了仍在進行,因此,唯有繼續張開眼睛看,才能體會「耶穌是主」在我們身上的精彩,才能體會耶穌作為「拯救者」在我們生命中的豐富.

有網友寫道:「要找到真愛,便要找一個懂我的人.這個人也許並不是十全十美,但因為他懂我；所以,我就認為他是十全十美的,就是這麼一個「懂」字.」有一個懂自己的人,是非常窩心的事情,難怪古人說:「人生得一知己,死而無憾」.當我們期待著一位懂自己的人的時候,我們會否同時看見,身邊的人也同樣地在期待著一位懂自己的人呢?

又有一位網友寫道:「今生只為尋找一個懂我的人」

我微笑的時候 如果你懂,只要你握緊我的雙手對我微笑,就夠了 我哭泣的時候 如果你懂,只要借我一個肩膀,靜靜地陪著我,就夠了 我委屈時候 如果你懂,只要給我你的懷抱,讓我在你的面前脆弱,就夠了 我任性的時候 如果你懂,就會包容,因為是你,所以我才會任性 我沈默的時候 如果你懂,就會逗我,因為我喜歡,喜歡你逗我的樣子 鬧彆扭的時候 如果你懂,就會哄我,因為我只是想做回小孩,被你疼愛 吵架的時候 如果你懂,就請先喊停,因為你要知道是生氣重要,還是感情重要 (\url{http://blog.udn.com/Elisa37/5774955}\#ixzz1j2e6OrML)

當我們享受著有人懂我們的時候,我們懂得身邊的人微笑的時候需要甚麼回應嗎?哭泣的時候需要甚麼回應嗎?委屆的時候需要甚麼回應嗎?任性的時候,沈默的時候,發脾氣的時候,吵架的時候．．．需要甚麼回應嗎?當我們享受著被看見的時候,我們願意張開眼晴去看看身邊的人嗎?

或多或少,我們都有過拿但業被耶穌看見的經驗,心靈被震撼,原來耶穌認識我們,知道我們.今天,耶穌呼喚我們,正如當天呼喚拿但業一樣,不要停留於被耶穌看見,耶穌期望我們也看見他；不要停留於那一次被耶穌看見的經驗,耶穌期望我們願意繼續被他看見,也期望我們繼續看他,因為「上帝的使者上去下來在人子身上」是動態的,是繼續進行的.



\newpage

\section{如何看待富裕和貧窮}
\label{sec:136}
\textbf{伍渭文牧師 - 如何看待富裕和貧窮?}
\newline
\newline
\hyperref[sec:135]{< < < PREV SERMON < < <}
~
\hyperref[sec:toc]{[返目錄]}
~
\hyperref[sec:137]{> > > NEXT SERMON > > >}


日期:2012年2月5日

題目:如何看待富裕和貧窮?

經文:創世記47

證道:伍渭文牧師

法國傳媒大亨巴拉昂 1998 年去世時在科西嘉人報刊登了他的遺囑:我曾是一個窮人,去世時卻是以富人的身分走進天堂的.在跨入天堂門檻之前,我不想把我成為富人的秘訣帶走,秘訣就鎖在銀行一保險箱內,誰能回答:窮人最缺少的是甚麼?猜中將拿走一百萬法郎,那是我給他的掌聲.共有 48561 人來信競猜,有人說缺少技能,有人說缺少幫助...,但一位小姑娘答中了:窮了缺了野心或雄心 (ambition).窮人最缺少的是創造財富的慾望.

要是渴想致富不單是慾望,是一個信仰:致富是上主對人的祝福,甚至是對人的期望,就成為吸引人的成功神學──上主的本意是叫人健康,得著財富,我們要有積極的思想,活潑的信心.成功神學在強調醫治,勇於開創的巨型靈恩教會和福音機構頗為盛行,很多美國基督教電台節目,都有這種思想.事實上,這些巨型教會和教會機構,本身就是財富,成功的標記.其滿有魅力的創會領袖,就憑其積極的思想,頑強的意志,建立起巨型教會和教會機構.成功神學 (prosperity theology) 顯然認為神聖的牛,可以有好的漢堡包──Sacred Cows Make Good Hamburger.

從某種角度來說,宗教信仰的確可以有助我們致富,因為信仰使我們更積極面對人生,對上主有信心,不輕言放棄.而且上主給我們智慧；有了智慧,投資創業,處人處事自然優勝過人.不是嗎?有一本書名叫《腦袋決定錢袋》,開宗明義說:思路決定出路,腦袋決定錢袋,在當今全球經濟一體化的今天,智慧經濟已被人們普遍接受... .「創業的智慧」讓他們調動沉睡在心中的那些優異,獨特的品質與淵博的知識；「競爭的智慧」讓他們輕鬆躲避了看不見的硝煙和刀光劍影；「創新的智慧」為他們搭建了一座通往財富的橋樑.但神聖的牛(宗教信仰),真的可以有好的漢堡包?──Sacred Cows Make Good Hamburger? 主主給我們的智慧,是用來賺取更多的財富的智慧,或是處理財富的智慧?信靠上主的人是否一樣會經歷貧窮呢?還是上主給我們的智慧,使我們能夠合宜看待富裕和貧窮的智慧.有了上主的智慧,我們可以處富裕,也可以處貧窮.

自2007年四月開始,我們主日崇拜開始講解創世記.先是創世記前十一章講萬物的源起,人的墮落,罪如何破壞世界....十二章開始講列祖人物:亞伯拉罕,以撒,雅各.跟著講約瑟生平.在上一次講道時,我提到約瑟這人物與之前的亞伯蘭,以撒,雅各,有所不同.上主都是透過異象 (vision),直接對列祖他們顯現,對他們說話.但約瑟則透過夢境,再給他圓夢的智慧 (wisdom) 去尋求意義.得著意義可能經過細緻的觀察,縝密的思考,多方的推敲,頓然有所覺悟,這是聖經的智慧傳統.約瑟是聖經中第一位滿有智慧的人.他貴為埃及宰相時,適時埃及經歷最富裕的七年,最窮困的七年.連續七年豐收,連續七年失收.而約瑟利用積聚起來的財富,在飢荒的七年,推行土地改革,振興經濟.但上主賜約瑟智慧,不是為自己賺取財富,他的前半生,更經歷被賣為奴,被誣告下獄,被人過橋抽板,刻意遺忘,是窮困的前半生.上主賜約瑟智慧,乃藉著他拯救他人,脫離饑饉.事實上,上主的選民也一樣經歷饑災,貧窮.

1. 遍地發生飢荒,上主選民也在當中經歷困苦(47:13)

饑荒甚大,全地都絕了糧,甚至埃及地和迦南地的人,因那饑荒的緣故,都餓昏了.(47:13)

上主呼召亞伯拉罕離開吾珥,引領他到應許之地.雅各在伯特利夢見天梯,在雅博渡口與上主的使者摔交,他們是上主的選民,有深刻的屬靈經歷,但當全地都絕了糧,他們都遭遇同樣的窮困,信仰並沒有使他們免去飢荒.當年沙士期間,不少信徒與其他市民一樣,感染沙士病毒,甚而死亡.大家都會想起謝婉雯醫生,一位自動請纓進入沙士病房,因而感染病毒而殉職的好姐妹,她的追思禮拜就在這教堂舉行.上主給我們的智慧,不是來穩賺財富,否則甘泉航空公司就不會破產,因為甘泉老板在主日是教會的義務牧師.然而,信仰和上主給我們的智慧,可以在破產時,極度貧窮時,有力量面對困局.信仰不是減少我們遭遇貧窮,疾病的機會,信仰是增加我們消滅貧窮,疾病的承擔.

2. 雅各給法老祝福(47: 7, 10)

雅各與眾子因為迦南地飢荒到埃及投靠法老.但雅各見了法老時,就給法老祝福──「

約瑟領他父親雅各進到法老面前,雅各就給法老祝福」.作者特別提到約瑟領他父親,把約瑟──雖然改了埃及名字──放回上主選民的身份來看.約瑟是雅各的兒子,祖父是以撒,曾祖父是亞伯拉罕,是上主的選民,是埃及及鄰近地區的拯救者.上主賜約瑟圓夢的智慧,辦事的能力,在七隻肥羊的豐年儲蓄糧食,在飢荒的日子拯救眾人.選民領受上主來的祝福,目的使別人得福,這是亞伯蘭蒙召的目的.(創 12:1-3)

12:1　耶和華對亞伯蘭說,你要離開本地,本族,父家,往我所要指示你的地去. 12:2　我必叫你成為大國．我必賜福給你,叫你的名為大,你也要叫別人得福． 12:3　為你祝福的,我必賜福與他,那咒詛你的,我必咒詛他,地上的萬族都要因你得福.

十年一次的全球聖公宗蘭柏 (Lambeth) 會議,出席六百七十位來自世各地的主教,在會議結束當天舉行反貧窮大遊行,向英國首相遞交信函,要求遵守 2000 年世界領袖禧年發展目標 (Millennium Development Goals) 承諾,協力推動 2015 年,世界貧窮減半 (Keep the Promise Half Poverty by 2015).信函提及:在人類大家庭中,(2015 年貧窮減半)失敗會繼續弄至上百萬婦女,兒童,因生產而和原來可預防的疾病而死亡.上千萬的人希望活下去,但因貧窮引發的衝突,沒有受教育的機會,性別歧視,使希望渺然；新一代更難擺脫受滋病的綑鎖.氣候變化,全球貿易和經濟都與貧窮問題息息相關.(In our human family the continuing consequence of failure is the death of millions of women and children from childbirth and preventable childhood diseases... Millions more have hopes for sustainable life diminished because of poverty-fuelled conflict, lack of education and gender inequality, all of which undermine a new generation and \#8217;s liberation from HIV... Climate change, global trade and economics were also haviung an impact on poverty.)

Keep the Promise Half Poverty by 2015(2015 年貧窮減半).六百七十位主教的呼喊,也是信仰群體應有的呼喊.上主選召亞伯拉罕,豈不是說:我必賜福給你,叫你的名為大；你也要叫別人得福嗎?

3. 推行土地改革,發展經濟(創 47:14-27)

請注意,約瑟的拯救工程,不是施捨救濟,因為坐食也會山崩.當飢民用盡賣糧食的金銀,約瑟要求他們變賣牲畜,換取糧食.最後讓他們曉得,他們還有生產資源──土地和勞動力.他們可以賣地,並為法老農耕.

47:18　那一年過去,第二年他們又來見約瑟說,我們不瞞我主,我們的銀子都花盡了,牲畜也都歸了我主,我們在我主眼前除了我們的身體和田地之外,一無所剩. 47:19　你何忍見我們人死地荒呢．求你用糧食買我們的我們的地,我們和我們的地就要給法老效力．又求你給我們種子,使我們得以存活,不至死亡,地土也不至荒涼. 47:20　於是約瑟為法老買了埃及所有的地,埃及人因被饑荒所迫,各都賣了自己的田地．那地就都歸了法老. 47:21　至於百姓,約瑟叫他們從埃及這邊,直到埃及那邊,都各歸各城.

有人以現時的觀念,批評約瑟趁機會欺壓埃及人,令他們賣地為奴,一味鞏固法老王的勢力.但細看經文,是埃及百姓主動要求約瑟,並感謝約瑟使他們存活.當時很多人為了糧食離開家鄉,不少農田荒癈.約瑟收購他們的土地,使他們成為法老王的佃農,目的使他們返回家園,投入生產,並供應種子,應許祇要納五分一作為稅收,五分四歸他們,多勞多得.這是一場土地大改革,願意投入生產力的,就有收獲.這是具有創意的救援工程,建立被援者的自信.Muhammad Yunus 在 Bangladesh 創立的窮人銀行 (Grameen Bank),已低息小額貸出了七億美金 (US\$ 7 billion),使窮人自力更生.蓋茨 (Bill Gates) 在 08年八月一期時代週刊撰文 ``How to Fix Captitalism"(如何做好資本主義)提及面對全球貧窮問題,不能單靠政府和志願團體.企業家需要有創意的資本主義 (creative capitalism) 思維,在增加財富之餘,亦能使更多人分享到財富,這是企業的社會責任.企業的社會責任,其實能夠提高企業形象,增加企業的競爭力.全球抗愛滋基金正游說企業,加入紅色行動 (Red),把產品部分收益捐給基金.我們也要運用智慧,管理我們的財富,幫助有需要的人.

4. 以色列人聚居歌珊地,為出埃及事件埋下種子 (47:4; 36:28-34)

從約瑟的一生,我們看到上主冥冥中的掌管歷史.他少年被兄長們所賣,後被波提乏主母誣告下獄,在獄中為酒政解夢,被復位的酒政推荐為法老解夢,立為宰相,處理將要來的七年饑荒,說明上主在掌管歷史.雅各舉家因饑荒下到埃及,適時遇到外來的法老當政,一反過去厭惡畜牧的人,善待放牧的雅各一家,淮其在歌珊地集中發展,這是全埃及最好的地方:「

埃及地都在你面前,只管叫你父親和你弟兄住在國中最好的地,他們可以住在歌珊地．你若知道他們中間有甚麼能人,就派他們看管我的牲畜 」(創 47:6).後來埃及人用牲畜與約瑟交換糧食,約瑟把得來的牲畜,帶到尼羅河三角洲蘭賽 (Rameses) ／歌珊 (Goshen) 地.以色列人在歌珊沃土大大發展起來,成為埃及人的威脅.出埃及記第一章記載:

1:5　凡從雅各而生的,共有七十人．約瑟已經在埃及. 1:6　約瑟和他的弟兄,並那一代的人都死了. 1:7　以色列人生養眾多並且繁茂,極其強盛,滿了那地. 1:8　有不認識約瑟的新王起來,治理埃及. 1:9　對他的百姓說,看哪,這以色列民比我們還多,又比我們強盛． 1:10　來罷我們不如用巧計待他們,恐怕他們多起來,日後若遇甚麼爭戰的事,就連合我們的仇敵攻擊我們,離開這地去了. 1:11　於是埃及人派督工的轄制他們,加重擔苦害他們．他們為法老建造兩座積貨城,就是比東,和蘭塞.

創世記和出埃及記都分別記載了兩件拯救事件:約瑟拯救埃及人,上主拯救祂的選民.在出埃及記,上主降下十災,裂開紅海,施行異能奇事作為拯救.但約瑟的拯救工程,則運用人的智慧,未雨綢繆,土地改革,誘發生產力,推動經濟發展,增加財富.我們不能像上卜土行神蹟展大能,但可效法約瑟,運用上主給我們的智慧,消滅飢饉,貧窮.而兩件事情,對約瑟來說,沒有任何關聯.但在歷史長河的回願中,我們卻發現:當人回應上主,盡上應盡的本份,就是參與上主隱密的救贖計劃.

5. 甚麼是真正的富裕?

失樂園的作者米彌頓說得好:There is nothing that makes men rich and strong but that which they carry inside of them. Wealth is of the heart, not of the hand. (除了心之所繫,沒有一樣東西令人感到富足和壯大,富足是心靈的,不是握在手中的).前一段時間,我買了《腦袋決定錢袋》這書的同時,也買了 Helen Keller 海倫‧凱勒的自傳《假如給我三天光明》.這位被譽為廿世紀美國十大英雄偶像的盲人作家,她希望如果她有三天光明,她會去經歷尋找她的富裕生活.

第一天:我要看人.我不知道甚麼是透過靈魂之窗──從眼睛看到朋友的內心.我只能用手指尖來「看」一個臉的輪廓.我能發覺歡笑,悲哀前其他明顯的情感.首先,我希望長凝視我親愛的老師,當我還是個孩子時,她為我打開了外面的世界.

第二天:我要在黎明起身,去看黑夜變為白晝的動人奇蹟.我將懷敬畏之心,仰望壯麗的曙光,看太陽如何喚醒沉睡了的大地.我將向世界,向過去和現在的世界匆忙瞥一眼.我想看看人進步的奇觀,那變化無窮的萬古千年.這麼多的年代,怎能被壓縮成一天呢?當然是通過博物館.

第三天:第一天我奉獻給我有生命和無生命的朋友；我二天向我顯示了人與自然的歷史,今天我將在當前的日常世界中度過,我將從我的家出發,到曼哈頓的第五街環城遊覽,到貧民窟,到工廠,到孩子玩耍的公園,我還要到外國人居住的地區,作一次不出門的海外遊行.

這是海倫‧凱勒所期望富裕的三天.富裕是心靈的,不是握在手裡的.正如她說:我的身體雖然不自由,但我的心是自由的.

在創世記人物裡面,不少信心的健兒在人生的跑道上前進.他們在賽程中有跌倒,失敗過:亞伯拉罕,以撒曾懼怕死亡,說謊話把妻子說成妹子,將她們陷在好色的君王手中；雅各更被稱為奸狡之父,雅各 (Jacob),抓著每一個機會,得著最大的利益.約瑟秀雅俊美,滿有智慧,但要一個過程,才能原諒出賣自己的兄長.但他們對信仰有執著,鎖定目標,奮力前行.在滿有憐恤恩典的教練一呼召他們主──悉心培育塑造下,都成為我們的榜樣.他們被上主呼召,蒙上主賜福,目的是成為別人的祝福.約瑟的智慧,不是為自己賺取財富,乃為有缺乏的人,分配財富.

七個豐年,終會過去.就如十八天的北京奧運,也會結束.除了心之所繫,沒有一樣東西令人感到富足和壯大,富足是心靈的,不是握在手中的.

富足是心靈的,不是握在手中的.



\newpage

\section{人鬼神}
\label{sec:137}
\textbf{伍渭文牧師 - 人,鬼,神}
\newline
\newline
\hyperref[sec:136]{< < < PREV SERMON < < <}
~
\hyperref[sec:toc]{[返目錄]}
~
\hyperref[sec:138]{> > > NEXT SERMON > > >}


馬可福音 1:21-28
\newline
\begin{longtable}{cl}
\hline
\hline
章節 & 經文 (和合本修訂版)\\
\hline
1:21 & \begin{tabularx}{0.7\textwidth}{X} 他們到了迦百農，耶穌就在安息日進了會堂教導人。 \end{tabularx} \\ \\ \relax
1:22 & \begin{tabularx}{0.7\textwidth}{X} 他們對他的教導感到很驚奇，因為他教導他們正像有權柄的人，不像文士。 \end{tabularx} \\ \\ \relax
1:23 & \begin{tabularx}{0.7\textwidth}{X} 當時，會堂裡有一個污靈附身的人，他在喊叫， \end{tabularx} \\ \\ \relax
1:24 & \begin{tabularx}{0.7\textwidth}{X} 說：「拿撒勒人耶穌，你為甚麼干擾我們？你來消滅我們嗎？我知道你是誰，你是神的聖者。」 \end{tabularx} \\ \\ \relax
1:25 & \begin{tabularx}{0.7\textwidth}{X} 耶穌斥責他說：「不要作聲，從這人身上出來吧！」 \end{tabularx} \\ \\ \relax
1:26 & \begin{tabularx}{0.7\textwidth}{X} 污靈使那人抽了一陣風，大聲喊叫，就出來了。 \end{tabularx} \\ \\ \relax
1:27 & \begin{tabularx}{0.7\textwidth}{X} 眾人都驚訝，以致彼此對問：「這是甚麼事？是個新的教導啊！他用權柄命令污靈，連污靈也聽從了他。」 \end{tabularx} \\ \\ \relax
1:28 & \begin{tabularx}{0.7\textwidth}{X} 於是耶穌的名聲立刻傳遍了全加利利周圍地區。 \end{tabularx} \\ \\
[1ex]
\hline
\hline
\end{longtable}
日期:2012年2月12日

題目:人,鬼,神

讀經:馬可福音1:21-28

證道:伍渭文牧師

1:21 到了迦百農,耶穌就在安息日進了會堂教訓人. 1:22 眾人很希奇他的教訓．因為他教訓他們,正像有權柄的人,不像文士. 1:23 在會堂裡有一個人,被污鬼附著．他喊叫說, 1:24 拿撒勒人耶穌,我們與你有甚麼相干,你來滅我們麼．我知道你是誰,乃是　神的聖者. 1:25 耶穌責備他說,不要作聲,從這人身上出來罷. 1:26 污鬼叫那人抽了一陣瘋,大聲喊叫,就出來了. 1:27 眾人都驚訝,以致彼此對問說,這是甚麼事,是個新道理阿．他用權柄吩咐污鬼,連污鬼也聽從了他. 1:28 耶穌的名聲,就傳遍了加利利的四方.

今天講道的

題目:人,鬼,神,有點聊齋小說的味道.中國人想到鬼神,就想到靈異的世界,一個我們陌生但有能力影響我們的世界,所以要敬鬼神而遠之,又要敬畏,又要遠離,不好惹.作為人重要的是做好自己,平生不作歪心事,半夜敲門也不驚.富蘭克林 (Benjamin Franklin) 說了一句合符聖經又精闢關於人,鬼,神的話:To err is human, to repent divine, to persist devilish, 犯錯皆凡人,悔改靠神恩,沉溺是魔障.人會犯錯,上帝使我們醒悟悔改,魔鬼使我們沉溺錯誤.人,鬼,神的互動不是在靈異的世界,就在現實的道德世界.魔鬼利用人性的軟弱,試探誘惑人；上帝使人覺醒悔改,離開罪惡,就如今天召請認罪的經文前一節所說:

「務要抵擋魔鬼,魔鬼就離開你們逃跑了,你們親近神,神就必親近你們.」(雅四7, 8 )

(一)沉溺是魔障 每個人都會犯錯.有些錯誤出於無知,有些錯誤出於疏忽,有些錯誤是為了利益明知故犯,聖經說信徒也會偶然被過犯所勝.(加六1) 最近義大利郵船觸礁沉沒,引至多人死亡失蹤,因為船長錯誤偏離航道,繼而棄船,沒有履行指揮救援旅客,正面臨刑事起訴.我們說這船長犯了不可饒恕的錯誤,甚至是觸犯了嚴重的罪行,但不會說他被污鬼附身.

沉溺罪行是魔障.經文中沒有交待為何這人會被污鬼附身,是否他沉溺在某種靈界迷信遊戲,染上邪惡的癖好,走火入魔而被污鬼所控制?像吸毒的人,酗酒的人,嗜睹成聖的人,財迷心竅的人,被色欲所迷的人,久而久之,已失去常性,讓污鬼有機可乘.驅魔是教會的傳統,是頂嚴肅的事,天主教會設有驅魔的神父,「該人須生命虔敬,有知識,審慎,和德行.法典1172條」(a presbyter endowed with piety, knowledge and integrity of life, Canon law, no. 1172) 有古老的驅魔禮儀,但轉介個案給驅魔神父前,需徵詢心理學家,精神科或一般醫生的專業意見,排除是身體或精神疾病的可能,神父在施行驅魔禮前要得到主教批准.真正鬼附的人,在為他祈禱時,會露出猙獰面相.我們要關心的不在驅魔的過程,雖然很多電影以此為招徠吸引觀眾,我們要警覺的是因沉溺罪惡,而引至靈性被鬼魔所壓制 (suppressed),甚至控制 (possessed).

迷信.我認識有一位大學三年級的女生,很想一畢業就找到一份高薪的工作,但多份申請都石沉大海,內心愈來愈焦慮,去拜黃大仙,廟裡的人教她賣一隻玉手鐲轉運,穿上了這手鐲後,果然找到工作,她深信這手鐲給她帶來好運.這是一個開始,我擔心她為了更多的財利尋求靈異的力量,迷信愈來愈深.大門已打開了.

色欲.我又認識一對夫婦,丈夫最近一年多背著妻子北上發展婚外情,太太最近四個月才發現,被揭發後丈夫沒有悔意,反而繼續沉溺下去,變成另一個人.結婚三十多年了,太太突然發現身邊有一個熟悉但又感到陌生的人,天天嚷著把房子賣掉,分一半財產,繼續現在的生活,因他知道沒有金錢,深圳年青的情人就不當他是「老公」了.很明顯,這丈夫的靈魂已被色欲的心魔所盤踞.不單大門打開,而且魔鬼就坐在大廳.

仇恨.三十年前我在美國深造神學,有一位在越南退役的同學,他跟我說:「在戰場最初殺人時感到很困難,因為自己是基督徒,但當殺人愈來愈多時,就變成另一個人,不像自己.殺人變成發洩,甚至成為嗜好,有一種魔力附身,非常可怕.」魔鬼已從大廳睡在床上.

制度也可能被鬼附.我們看看這被污鬼所附的人如何出現:

「在會堂裡有一個人被污鬼附著.他喊叫說:拿撒勒人耶穌,我們與你有甚麼相干?你來滅我們嗎?我和道你是誰,乃是神的聖者.」(可一21-23)耶穌在會堂教訓人的時候,這人突然出現,好像從外邊闖進會堂,因一般被污鬼附身的人,不會到會堂崇拜聽道,他們早就被人隔離,像格拉森 (Gerasene) 被鬼附的人(可五).還是這人當時就在席中?當天如是,過去如是,鬼附的事實秘而不宣,表面與常人沒有分別.是耶穌的臨在,是他教訓人時帶出的權柄,令污鬼感到不安,露出本來的面目,是屬靈的權柄,挑釁邪惡的力量.耶穌潔淨聖殿,趕走在那裡做買賣的人,他指出法利賽人是粉飾的墳墓,耶穌揭發黑暗,挑戰罪惡,使隱藏的鬼附現身,他是神聖的尋釁滋事者.

這事使我想起前一段時間,某神職人員嚴厲批評某地產商的商業霸權行為無孔不入,包攬所有衣,食,住,行,而且分毫必賺,點滴不漏,與聖經規定收割麥穗,不可拾遺,留給窮困者充饑的教訓相去甚遠,忘記耶穌先祖路得就因為拾麥穗,存活下來 (路得記二) , 儼然是魔鬼附身的行為.所以,污鬼也可附在為利是圖的商業手法上,附在壓制人的政治制度上.作為跟隨耶穌的人,我們也要活出上帝兒女的權柄(約一12),挑戰不義的事,挑戰壓制人性的制度,為上帝驅魔.當然,作為驅魔者,「該人須生命虔敬,有知識,審慎,和德行.法典1172條」(a presbyter endowed with piety, knowledge and integrity of life, Canon law, no. 1172) .驅除制度的鬼附需要更多的知識,因問題綜錯複雜.審慎幫助我們不冒進弄壞事情,慎審是把握時機,擊中要害.虔敬使我們動機純良；德行使我們不懷苦毒怨憤.制度的驅魔是愛慕公義光明,不是嘲弄攻擊妖魔化他人.  (二)在限制中顯出能力的權柄

「到了迦百農,耶穌就在安息日進了會堂教訓人.」(可一21)

上帝不是活在靈異的世界,讓我們敬而遠之,上帝就在我們中間行走.馬可福音的開始,耶穌到約但河接受施洗約翰的洗禮；行到海邊,呼召打魚為生的西門,安得烈,雅各,約翰為門徒；進入迦百農的會堂,在安息日教訓人.他總是去到我們中間,接觸我們,服侍我們.迦百農的會堂,人數應該不會很多,參加的人,背景也不會顯赫,但耶穌的偉大事工,就以卑微的迦百農會堂開始.

「耶穌就在安息日」.「就在」有立刻的意思,指出耶穌把握時間 ,上帝的兒子活在時間內,受時期約制,不能虛擲光陰.「安息日」指耶穌順著當時的習慣,規章和限制來傳道,彰顯福音的能力.今天是除夕,提醒我們時間的寶貴,這隻兔跑得真快,明天就是龍年了.雖然對新年我們都很多願望,很多迷信的人會到沙田公車廟轉轉風車,換換運程,但可能仍然事與願違,我們還要面對種種限制,甚至變本加厲,我們要效法耶穌,把握機會,在限制中做好自己,不要逃避困難,諉過環境而裹足不前,磋跎歲月.

事實上,沒有任何外在的環境可以攔阻福音的能力,因福音是關乎權柄,不關乎權力.「眾人很希奇他的教訓；因為他教訓他們,正像有權柄的人,不像文士.」(可一22)權力和權柄不同,權力是影響力,來自權位和金錢,人可能迫於無奈屈從權力,或因利誘趨附權力.權柄可譯作權能,是改變生命的能力,使人甘心歸附的吸引力.沒有權位和金錢,也可以因為屬靈的品格,和認識聖經帶來的智慧和洞察力,使人感到有一種權柄,樂於聽從.

(三)使人希奇驚訝的權柄

「眾人都驚訝,以致彼此對問說,這是甚麼事,是個新道理阿．他用權柄吩咐污鬼,連污鬼也聽從了他.」(可一7)

耶穌教訓人,帶有權柄,使污鬼感到惶恐,懼怕被逐；但在人的心中,卻引起希奇(astonished),驚訝(amazed) ,這是從來沒有的發生過的事,新時代來臨了,密封的世界開了通道,我們窺探到一個新境界,帶來新盼望,我們不用絕望,也毋需犬儒嘲弄.今年三月第五屆崇基基督教文化節主題為「盼望」,但甚麼是盼望?我很喜歡耶穌會會士,澳洲著名難民,人權,社會公義維權人士馬可.雷珀(Mark Raper)對盼望的描述:「盼望並不是樂觀.樂觀是希望事情會好轉；盼望是忍受苦難的美德,衍生力量的恩典.盼望是心靈深處開出的應許,引領我們在茫茫前路中穩步邁進.」(Hope is not optimism. Optimism expects that things will get better. Hope is a virtue grounded in suffering. It is a grace that gives strength. Hope is a promise that takes root in the heart and is a guide to an unknown future-Mark Raper S.J.) 作為社會公義維權人士,雷珀神父明白到不公義的制度非常頑強,我們不會樂觀,但心存盼望,因耶穌已經打開密封的世界,我們已看到福音帶來的希奇和驚訝.

我曾在一次講道中提到福音使人生命得奇妙的改變:最怕被遺返中國的賴昌星,因信了福音,平靜接受被遺返中國,令人希奇.十二年斬殺報販霞姐的兩位兇手,去年六月落網,其中一位接受福音後,感到難過,要供出幕後主謀,使人感到驚訝.福音使坐牢的感到不安,又使懼怕坐牢的感到平安.

也許在新的一年,當人人為著經濟低迷擔心時,我們敬虔知足的心和平靜安穩的態度,使認識我們的人,都感到希奇,因為在我們心靈深處有一個盼望,引領我們在茫茫前路中穩步邁進.

「到了迦百農,耶穌就在安息日進了會堂教訓人.」(可一21)我們迎接新的一年,從初一到初三有三天公眾假期,可以和親人,友好相聚,也可以去旅行休憩,到了初四,上班了.阿們.



\newpage

\section{保障期永遠有效}
\label{sec:138}
\textbf{李慶宏先生 - 保障期永遠有效}
\newline
\newline
\hyperref[sec:137]{< < < PREV SERMON < < <}
~
\hyperref[sec:toc]{[返目錄]}
~
\hyperref[sec:139]{> > > NEXT SERMON > > >}


詩篇 121:1-8
\newline
\begin{longtable}{cl}
\hline
\hline
章節 & 經文 (和合本修訂版)\\
\hline
121:1 & \begin{tabularx}{0.7\textwidth}{X} 我要向山舉目，我的幫助從何而來？ \end{tabularx} \\ \\ \relax
121:2 & \begin{tabularx}{0.7\textwidth}{X} 我的幫助從造天地的耶和華而來。 \end{tabularx} \\ \\ \relax
121:3 & \begin{tabularx}{0.7\textwidth}{X} 他不叫你的腳搖動，保護你的必不打盹！ \end{tabularx} \\ \\ \relax
121:4 & \begin{tabularx}{0.7\textwidth}{X} 保護以色列的必不打盹，也不睡覺。 \end{tabularx} \\ \\ \relax
121:5 & \begin{tabularx}{0.7\textwidth}{X} 保護你的是耶和華，耶和華在你右邊蔭庇你。 \end{tabularx} \\ \\ \relax
121:6 & \begin{tabularx}{0.7\textwidth}{X} 白日，太陽必不傷你；夜間，月亮也不害你。 \end{tabularx} \\ \\ \relax
121:7 & \begin{tabularx}{0.7\textwidth}{X} 耶和華要保護你，免受一切的災害，他要保護你的性命。 \end{tabularx} \\ \\ \relax
121:8 & \begin{tabularx}{0.7\textwidth}{X} 你出你入，耶和華要保護你，從今時直到永遠。 \end{tabularx} \\ \\
[1ex]
\hline
\hline
\end{longtable}
日期:2012年2月19日

題目:保障期永遠有效

經文:詩篇121:1-8

證道:李慶宏先生

祝賀今天一同聚會崇拜的主內肢體,朋友-春節蒙福,主恩滿溢!2012年第一個月過去了,覺得一切順利嗎?在踏入新一年之前,我們就聽到很多鄭重的警告,預測世界經濟將面對艱難險峻的一年,日子不會好過,小心你的錢包和資產；至於政治方面,新一年也是改朝換代的一年,全球有超過五十個國家地區要改選或進行領袖更替,其中包括中,美,俄,韓,台灣......國際政局充滿變數.當泰坦尼克號船難事件在2012年適逢百周年紀念之際,意大利豪華郵輪協和號在一月中就先作了一場預告演出.年初三下午,我和家人到東涌和機場走了一趟,慶幸沒有踏上360吊車,總算是避過小小的一劫.

無論是2012年或者是壬辰龍年,你可能有不同的挑戰和疑慮,心中總會產生一個問題:「我怎樣可以得到足夠的保障?」這問題影響著我們的生活,學習,工作,家庭,與人的關係,也影響著人生的方向與目標:「我怎樣可以得到足夠的保障?」,你有答案嗎?要尋找答案的話,讓我們現在來看看詩篇121篇.(誦讀經文,詩篇121篇1-8節)

詩篇121,是上行之詩,是以色列人前往耶路撒冷朝聖旅程上所用的詩歌,今日我們也在天路歷程中,不斷親近敬拜神,走天路,目的也是更高處行.在這上行的路上,你行走時,是越走越有信心,還是越來越缺少安全感?讓我們從聖經找到安慰鼓勵!請記住,這安慰是給予那些真正願意為主奔走天路的人的.

(一)　造物主宰的保障:  我要向山舉目；我的幫助從何而來? 我的幫助從造天地的耶和華而來. 朝聖者在旅途中,登上耶路撒冷這山城,需要經過崎嶇山路,過程中有軟弱,有疲乏,也帶著生活中各種尋問,來到聖殿尋求幫助.朝聖者的眼目看對了方向,就是抬頭仰望.我們在生活的沖擊壓力下,有多少時間是向上仰望呢?我們常常會向後回望,向兩邊張望,懂得向前遠望已經了不起,但最重要的是向上仰望!朝聖者向群山舉目,有什麼感覺呢?詩人自問自答,我的幫助到底從何而來?向山舉目,這動作只不過是過程,不是目的!山可以代表蔭庇之所,是崇高穩定的象徵,但山也可能是敵對者藏身之處,此行是否平安,前路有沒有兇險?在外邦人信仰中,群山經常是偶像的居所,而我是否要投靠這些山以得幫助呢?

抬頭仰望,自然見到前面的高山,但山不是目標,山有何能力?山外的天又有何能力?我們從望山開始,進一步望見山外的天地,從遼闊的天地再仰望那位高天之上的創造者,耶和華神才是我們的幫助和保障!「我的幫助從造天地的耶和華而來」.我們尋找幫助時,向何處舉目呢?是否認定某些人,某些事,某些力量可以投靠,以為這些是不可搖動的靠山?我們投靠的不是被造之物,乃是造天地的耶和華,祂有無限能力,是超越一切的終極保障!我們相信的是造物之主,祂不是山,卻是真正的「靠山」!

(二)　全年無休的保障: 他必不叫你的腳搖動；保護你的必不打盹! 保護以色列的,也不打盹,也不睡覺. 這一組的兩節經文,從「不」的觀念來介紹我們的保護者,祂「不作」什麼!打盹睡覺都不是我們的神的特性.這裡是故意針對外邦偶像而說的,因為那些偶像被視為會打盹會睡覺的,善信要稟告這些偶像神明,否則這些「神」不會明白我們的境況.我家住在大衛車公廟附近,春節假期都會見到大排長龍的善男信女,來進行一年一度的報告,他們與車公的平日接觸大概不多,因此期望也不會過高.耶和華全年無休,祂可以保護我們的每一步,每分每秒在路上都不致跌倒失足.我們深知打盹和睡覺是什麼一回事,聽道好容易會打盹睡覺,講道者無論自認為多麼精彩,都發現有人在打盹睡覺,只希望與「我」無關吧.

原來我們的神不會睡覺,連打盹也不會.不打盹不睡覺的保護,是隨時隨地的最好關懷!一方面給你平安,給你安慰,不過另一方面,你要認真想想,你一言一行,這位不睡覺不打盹的神都一清一楚,我們是否要多一點敬畏?多一點謹慎?多一點責任感?在遭遇變故或者磨練的時候,不要以為神睡著了,疏於照顧,不會的,祂一直在我們身邊,連你睡著祂也不會打盹!放心交託!請相信,神在任何時地都是「保護你的」.祂是個人的保護者,祂是整個以色列的守護者.神保護我們每一個基督徒,祂也不睡覺不打盹的保守教會群體向前行!

(三)　同在爭戰的保障:   保護你的是耶和華；耶和華在你右邊蔭庇你,  白日太陽必不傷你；夜間月亮必不害你. 3-4節可以視為日常生活的一般保守,而5-6節就是兵凶戰危下的非典型保障.在右邊蔭庇,這裡指的是一種爭戰狀況,神不是在平安無事時,才能顯出保護,就算在生活爭戰中,神也有保護,而且是貼身的保護.什麼是「右邊蔭庇你」?古代作戰時,戰士通常右手持刀,左手持盾牌,於是右邊就露出破綻.為了保護一些大將,會有專人拿著盾牌在他的右邊,以保安全.神也同樣悉心看顧祂所愛的兒女,就算我們自己不能兼顧,經歷不自覺,不自知的危險時,神也在身邊蔭庇你.

這種蔭庇,目的是令我們不受傷害.白日和夜間,這種兩極的用語,表達兩極中所包含的一切範圍,意思指無論日夜,我們都平安!太陽,是指日間上路的灼熱,月亮是指晚上的陰暗,另一方面,太陽和月亮都是古代一些偶像的形像.在神的蔭庇下,與我們為敵的邪惡勢力對我們絲毫無損,於是我們放膽上路.弟兄姊妹,最近有沒有爭戰,無論來自生活哪一方面,如果你行走向上親近神的道路,神都會在身邊幫你打仗,讓你不受傷害!

(四)　永遠有效的保障:   耶和華要保護你,免受一切的災害；他要保護你的性命.  你出你入,耶和華要保護你,從今時直到永遠. 除了日常的保護,爭戰的保護之外,性命安危-「生死」,應該是我們所有人最大的關心和擔心.生命總有限期,縱使人類千方百計把大限延長,但誰也不敢斷言壽數的長短,當中的災病意外,更在人力控制之外.這裡說,神的保護可以免去一切災害,我們的性命不會被禍患所奪去.這並不是說我們出入平安到一個地步,一定絕處逢生,衝出馬路也沒有危險,對什麼疫症都可以不作預防.這裡說的是,我們不用為死亡擔心和懼怕,有時怕死比真死還要恐怖.不要怕,因為沒有別的力量和災害可以取去我的性命,請聽清楚,如果信徒行走天路期間結束世上的生命,只有一個解釋,這是神的工作,而不是災害的力量,地上的日子結束,不過是永恆人生另一階段的開始,你的性命永遠在神的保守之中.這是我們可以安然面對人生的最大安慰!

「出入平安」,是春節常見的賀詞祝福,惟有超越時空地域的神能夠給你這樣的應許!你出你入,神要保守你,從今時直到永遠!這裡說的,比5-6節的日夜更為廣遠,所指的是過去,現在與未來,超越了時間限制,這種永遠的保護只可能從永恆的神而來!到商場超市購物,消費者都會留意有效日期,精挑細選,哪怕只是多一天兩天的保存期,若上面打印著「永遠有效」,大概你會認定這是假劣貨品,騙人的,因為人間沒有可能提供這樣的保障,保險公司只能給你出事之後的賠償,絕對不敢保証你的平安.但我們有這樣一位真神,只要你以信心行走向上的天路,每一天,我們都可以經歷:「我的幫助從造天地的耶和華而來」!神要我們知道,祂定意要給我們幫助,祂的幫助不會改變,不會縮減,這是造物主宰的保障,是全年無休的保障,是同在爭戰的保障,這保障的有效日期是:「從今時真到永遠」!



\newpage

\section{跨越界限}
\label{sec:139}
\textbf{林豪恩先生 - 跨越界限}
\newline
\newline
\hyperref[sec:138]{< < < PREV SERMON < < <}
~
\hyperref[sec:toc]{[返目錄]}
~
\hyperref[sec:140]{> > > NEXT SERMON > > >}


馬可福音 1:40-45
\newline
\begin{longtable}{cl}
\hline
\hline
章節 & 經文 (和合本修訂版)\\
\hline
1:40 & \begin{tabularx}{0.7\textwidth}{X} 有一個痲瘋病人來求耶穌，向他跪下，說：「你若肯，你能使我潔淨。」 \end{tabularx} \\ \\ \relax
1:41 & \begin{tabularx}{0.7\textwidth}{X} 耶穌動了慈心，就伸手摸他，說：「我肯，你潔淨了吧！」 \end{tabularx} \\ \\ \relax
1:42 & \begin{tabularx}{0.7\textwidth}{X} 痲瘋病立刻離開他，他就潔淨了。 \end{tabularx} \\ \\ \relax
1:43 & \begin{tabularx}{0.7\textwidth}{X} 耶穌嚴嚴地叮囑他，立刻打發他走， \end{tabularx} \\ \\ \relax
1:44 & \begin{tabularx}{0.7\textwidth}{X} 對他說：「你要注意，千萬不可告訴任何人，只要去，讓祭司為你檢查，又因為你已經潔淨，獻上摩西所吩咐的祭物，作為證據給眾人看。」 \end{tabularx} \\ \\ \relax
1:45 & \begin{tabularx}{0.7\textwidth}{X} 那人出去，倒說許多的話，把這件事傳揚開了，使耶穌不能再公開進城，只好留在外邊曠野地方，人從各處都到他跟前來。 \end{tabularx} \\ \\
[1ex]
\hline
\hline
\end{longtable}
列王記下 5:1-14
\newline
\begin{longtable}{cl}
\hline
\hline
章節 & 經文 (和合本修訂版)\\
\hline
5:1 & \begin{tabularx}{0.7\textwidth}{X} 亞蘭王的元帥乃縵在他主人面前是一個偉大的人，得王的喜悅，因為耶和華曾藉他使亞蘭人得勝。他雖然是大能的勇士，卻染上了痲瘋。 \end{tabularx} \\ \\ \relax
5:2 & \begin{tabularx}{0.7\textwidth}{X} 亞蘭人成群出征的時候，從以色列地擄了一個小女孩，她就服事乃縵的妻子。 \end{tabularx} \\ \\ \relax
5:3 & \begin{tabularx}{0.7\textwidth}{X} 她對女主人說：「我希望主人去見撒瑪利亞的先知，他必能治好主人的痲瘋。」 \end{tabularx} \\ \\ \relax
5:4 & \begin{tabularx}{0.7\textwidth}{X} 乃縵去告訴他主人說，從以色列地來的女孩如此如此說。 \end{tabularx} \\ \\ \relax
5:5 & \begin{tabularx}{0.7\textwidth}{X} 亞蘭王說：「你可以去，我也會送信給以色列王。」於是乃縵手裡帶十他連得銀子、六千舍客勒金子和十套衣裳去了。 \end{tabularx} \\ \\ \relax
5:6 & \begin{tabularx}{0.7\textwidth}{X} 他帶著這信給以色列王，說：「現在你接到這信，看哪，我派臣僕乃縵到你這裡來，你要治好他的痲瘋。」 \end{tabularx} \\ \\ \relax
5:7 & \begin{tabularx}{0.7\textwidth}{X} 以色列王讀了信就撕裂衣服，說：「我豈是神，能使人死使人活呢？這人竟派人來，叫我治好一個人的痲瘋。你們要知道，看，這人是找機會來跟我吵架的。」 \end{tabularx} \\ \\ \relax
5:8 & \begin{tabularx}{0.7\textwidth}{X} 神人以利沙聽見以色列王撕裂衣服，就派人到王那裡，說：「你為甚麼撕裂衣服呢？讓那人到我這裡來，他會知道以色列中有先知。」 \end{tabularx} \\ \\ \relax
5:9 & \begin{tabularx}{0.7\textwidth}{X} 於是乃縵帶著車馬到了以利沙的家，站在門前。 \end{tabularx} \\ \\ \relax
5:10 & \begin{tabularx}{0.7\textwidth}{X} 以利沙派一個使者，對乃縵說：「去，在約旦河中沐浴七次，你的肉就必復原，你會得潔淨。」 \end{tabularx} \\ \\ \relax
5:11 & \begin{tabularx}{0.7\textwidth}{X} 乃縵卻發怒走了。他說：「看哪，我以為他必定會出來，到我這裡，站著求告耶和華－他神的名，在患處上搖手，治好這痲瘋。 \end{tabularx} \\ \\ \relax
5:12 & \begin{tabularx}{0.7\textwidth}{X} 大馬士革的亞瑪拿河和法珥法河豈不比以色列的一切水更好嗎？我難道不可以在那裡沐浴而得潔淨嗎？」於是他生氣，轉身走了。 \end{tabularx} \\ \\ \relax
5:13 & \begin{tabularx}{0.7\textwidth}{X} 他的僕人近前來，對他說：「我父啊，先知若吩咐你做一件大事，你豈不做嗎？何況是吩咐你去沐浴，得潔淨呢？」 \end{tabularx} \\ \\ \relax
5:14 & \begin{tabularx}{0.7\textwidth}{X} 於是乃縵下去，照著神人的話，在約旦河裡浸了七次。他的肉復原，好像小孩的肉，他就潔淨了。 \end{tabularx} \\ \\
[1ex]
\hline
\hline
\end{longtable}
日期:2012年2月26日

題目:跨越界限

証道:林豪恩先生

經文:馬可福音1:40-45;列王記下5:1-14

今天的福音經課和舊約經課所記載的,都是醫治大痳瘋的故事.閱讀這兩個醫治大痳瘋故事的時候,我們需要留意這不是醫學文獻,因此,不需要以現代醫學的角度去檢視故事記載的醫治方法；由於這是宗教文獻,我們值得花時間思考故事所啟發的人生及生活意義.

第一個故事:跨越潔淨與不潔淨的界限

第一個故事記載在馬可福音1章40-45節:

1 : 40 有一個長大痲瘋的來求耶穌,向他跪下,說:「你若肯,必能叫我潔淨了.」 1 : 41 耶穌動了慈心,就伸手摸他,說:「我肯,你潔淨了吧!」 1 : 42 大痲瘋即時離開他,他就潔淨了. 1 : 43 耶穌嚴嚴地囑咐他,就打發他走, 1 : 44 對他說:「你要謹慎,甚麼話都不可告訴人,只要去把身體給祭司察看,又因為你潔淨了,獻上摩西所吩咐的禮物,對眾人作證據.」 1 : 45 那人出去,倒說許多的話,把這件事傳揚開了,叫耶穌以後不得再明明地進城,只好在外邊曠野地方.人從各處都就了他來.

今次讀這個故事的時候,讓我們留意一下大痳瘋患者和耶穌對話的用語.大痳瘋患者跪下求耶穌時並不是說:「你若肯,必能醫治我.」而是說:「你若肯,必能叫我潔淨了.」對付大痳瘋,患者用的語言是「潔淨」,而不是「醫治」.耶穌回應說:「我肯,你潔淨了吧!」,他的用語,也是「潔淨」,而不是「醫治」.

相對於身體疾病,我們一般用「醫治」這個詞彙,因此,用「醫治」來談論大痳瘋本是合適的.「潔淨」所相對的是「不潔淨」,這個詞語用來描述身體健康狀況,它的含義超出了身體健康狀況,而包括了宗教及社會意義.

舊約聖經利未記第十三及十四章詳細記載了大痳瘋的種類和症狀.不過,Christian Medical Fellowship的醫學專家仔細考察利未記所描述的大痳瘋病,發現這其實不是現代醫學所指的痳瘋病,而是一系列展現於皮膚的疾病,例如象皮病,和各種癬病.無論如何,這些皮膚及外觀與眾不同的人,都被歸入「不潔淨」的類別.被歸類為「不潔淨」的人,有甚麼社會含意呢?利未記13章45-46節記載了加諸對他們身上的的規範:「身上有長大痲瘋災病的,他的衣服要撕裂,也要蓬頭散髮,蒙著上唇,喊叫說:『不潔淨了!不潔淨了!』災病在他身上的日子,他便是不潔淨；他既是不潔淨,就要獨居營外.」由此可見,若有人一旦被歸類為「不潔淨」,他們便失去正常人的權利,也得不到正常人應有的對待.因為「潔淨」與「不潔淨」是不能接觸的,若接觸,「潔淨的」會被沾污變成「不潔淨」.起初,人們認為沾污是經由身體接觸.後來,人們越來越謹慎,漸漸相信「空氣傳染」也是沾污途徑.例如,一個不潔淨的人站在樹下,而一個潔淨的人經過,這個潔淨的人就會變成不潔淨.除了「人傳人」之外,人與物品也可以互相沾污,例如,大痳瘋患者進入屋內,屋內的物品及人都會被變為不潔淨.因此,社會不能讓「不潔淨」的人出現於「潔淨」空間,免得「潔淨」的人和物都被沾污.

把這個故事放回當時的宗教及社會背景去閱讀,我們會發現耶穌醫治大痳瘋的故事是一個跨越界限的故事.

第一個跨越界限的人是那位痳瘋病患者.為甚麼說他跨越界限呢?如前所述,他被歸類為「不潔淨」,就需要存在於社會為他劃定的特定範圍,不應該在「潔淨」的空間出現.現在,他在加利利的會堂出現,是一個跨越界限的行動.若他要進入「潔淨」的空間,他就要撕裂衣服,蓬頭散髮,一邊走一邊喊叫說:『不潔淨了!不潔淨了!』,提醒「潔淨」的人退避三舍.現在,他走到耶穌的跟前提出要求,明顯地跨越了「不潔淨」的人與「潔淨」的人之間的界限.這位痳瘋病患者雖然被定義為「不潔淨」,被貶為次等人,被排拒於正常生活之外,但他以乎並不屈服於當時宗教及社會為他所設的界限,為了自己,也許也為了其他如他一樣的人,他要跨越界限.這是一位勇敢地跨越界限的痳瘋病人.

或者,勇敢地跨越「潔淨」與「不潔淨」界限的人並不只他一個,不過,他是幸運的,因為他遇到的,是一位同樣願意跨越界限的耶穌.就算在二千多年後的現代社會,若有被標籤為「不潔淨」的人勇敢地跨越界限,也不一定遇到同樣願意跨越界限的「潔淨」的人.

幾年前,有一間服務戒毒康復青少年的學校申請遷往一間空置學校,但遭遇學校所在地區的居民反對.校長及老師於是帶同三十位學生往該地區與居民見面,以為讓居民直接與學生接觸,認識這些曾經誤入歧途,現在改過自身的青少年其實並沒有甚麼異樣,因而接納他們遷到較為安全的校舍繼續讀書.

當天,約有 300名當地居民早有準備去「迎接」這30位學生.大部份居民手持反對遷校的標語及紙牌,甚至有人內高呼:「唔要吸毒仔,吸毒女!」那30位勇敢地走出來的學生聽到聲勢浩大的噓聲,一個一個垂下頭來,有些女學生因此而忍不住眼淚,哭了出來. 其中一位16歲的女學生說:「我一落船已經聽到有人叫我地做吸毒妹,我覺得好難受.從碼頭到社區會堂被反對的聲音包圍,我真係好驚,成個人震晒.我地好想軟化佢地個心,希望佢地可以接納我地.」(蘋果日報2009年6月15日)可惜,這群勇敢地跨越界限的戒毒康復青少年的行動,卻遇到保護界限的人的強硬回應,拒絕他們使用該空置校舍,以免該地的「潔淨」被沾污變成「不潔淨」.

生活於二千多年前古代社會的這位大痳瘋患者比生活於二千多年後文明社會的這群學生幸運,因為耶穌回應這位痳瘋病人的越界,並不是定義他的行為自私,沒有指責他為了一己的利益走出來沾染大眾；同時,耶穌的反應也不是明哲保身,退後三丈,然後「報警」把他捉回「營外」.耶穌沒有保護那條界線,反而如那位痳瘋病人一樣勇敢地跨越界線,他接納對方的要求,承諾說:「我肯,你潔淨了吧!」道成肉身的上帝伸出手來回應,成就了這個跨越界限的故事.耶穌基督以「潔淨」之身去接觸「不潔淨」,打破了當時宗教及社會「潔淨」與「不潔淨」的界限,並且粉碎了「潔淨的」與「不潔淨的」接觸會被沾污為「不潔淨」的恐懼.

面對這位被歸類為「不潔淨」的痳瘋病人的請求,聖經說「耶穌動了慈心,就伸手摸他」.是甚麼動了耶穌的慈心呢?是痳瘋病人被痳瘋病折磨的慘況嗎?或者是.不過,正如Christian Medical Fellowship的醫學專家判斷,聖經說的大痳瘋不等於現代醫學診斷的痳瘋病,那麼,觸動耶穌的慈心,可能是這些被歸類為「不潔淨」的人所遭遇的次等人的對待及生存狀況.

「潔淨」或「不潔淨」這些詞彙,現代社會已經不再使用,不過,這並不代表人們心裡沒有把人作類似的區分.人類仍然面對試探,特別是在非常時期,往往把另一些人歸類為次於正常人的次等人,以便合理化不以正常人的待遇去對待他們.

二十多年前有一套電影名叫「黑太陽七三一」.電影中的日軍將領要求年輕軍兵以中國及俄羅斯平民進行各種化學武器實驗.有些年少的軍兵望者實驗對象也是如自己一樣的人,不忍下手.於是,他們的長官訓斥他們,要求年少軍兵不把視實驗對象為如自己一樣的人,要稱實驗對象為「馬魯達」,即是木材.如此,這些年少軍兵就能夠容許自己殘酷地對待這些被捉回來的中國和俄羅斯平民.這是戰爭時期人類犯下最殘酷的罪惡之一.

與戰爭時期相此,此時此地,相對和平盛世.在和平的日子,我們假設有更大的空間去跨越社會上不必要的界限,重視檢視及拿走被判定為「不潔淨」的群體的標籤,拿走被定義為「不正常」的人們的標籤,讓他們回到社會中過一般人的生活；我們假設有較為和平及對重的方法去處人群之中的糾紛,冷靜客觀地去分析及解決問題,而不是情緒化地貶斥別人是「狗」,或者是「蝗蟲」.

昔日,耶穌藉著醫治大痳瘋,與大痳瘋患者一起創造了一個跨越「潔淨」與「不潔淨」的故事.今天,當現代文明社會面對困難時,但願這個故事提醒我們,不要墮入隨便標籤別人為「不潔淨」的陷阱,也不要以為貶低別人作為解決問題的引誘.

第二個故事:跨越敵我的界限

馬可福音的大痳瘋患者得醫治,因為有人跨越了潔淨與不潔淨的界限,列王記下那位亞蘭國的元帥的大痳瘋得醫治,又是一個跨越甚麼界限的故事呢?

這個故事發生於以利沙作先知的時代.當時,以色列和亞蘭是鄰國,兩國經常互相攻打,戰事連年.因為戰爭,有人被殺,有人喪失家人,有人家園被毀,打敗仗的更被擄到敵國,變為奴隸.戰爭往往帶給交戰兩國人民深刻的傷害,極大的憤恨,以及強烈的敵意.

5 : 1 亞蘭王的元帥乃縵在他主人面前為尊為大,因耶和華曾藉他使亞蘭人得勝；他又是大能的勇士,只是長了大痲瘋. 5 : 2 先前亞蘭人成群地出去,從以色列國擄了一個小女子,這女子就服事乃縵的妻.

故事首先出場的,是亞蘭國戰績輝煌的大能勇士,貴為元帥的乃縵.第二位出場的,卻是一位無名無姓的「小女子」,這位「小女子」正正是這位元帥攻打以色列時擄掠回來當下人的.這兩位人物的關係是甚麼呢?他們是交戰兩國的人,即是敵人.

現在,敵方的元帥竟然患了大痳瘋,這位以色列的小女子可以如何反應呢?常言道,對敵人仁慈,就是對自己殘忍.得知這位國家頭號敵人患大痳瘋,就算因為身為奴隸處身敵國而不能慶祝,心裡也禁不住幸災樂禍:「你都有今日啦!」就算這位以色列小女子對亞蘭元帥沒有個人的仇恨,也不應該同情或幫助對方,敵我分明,對方是國家敵人,同情或幫助國家敵人,即是背叛自己的國家.

5 : 3 她對主母說:「巴不得我主人去見撒馬利亞的先知,必能治好他的大痲瘋.」

故事的發展有違常理:這個以色列小女子竟然向元帥的妻子介紹醫治大痳瘋的秘訣.可能當時的以色列沒有「國民教育」,她竟然不知道在戰爭狀況之下向敵人提供幫助是叛國的行為,是要被軍法處置的.她跨越了敵我的界限,身為一位被擄者,向敵人伸出援手.

被自己擄掠回來的以色列小女子跨越了敵我界限,向自己伸出援手,亞蘭的元帥能夠同樣跨越敵我界限回應嗎?

5 : 4 乃縵進去,告訴他主人說,以色列國的女子如此如此說. 5 : 5 亞蘭王說:「你可以去,我也達信於以色列王.」於是乃縵帶銀子十他連得,金子六千舍客勒,衣裳十套,就去了； 5 : 6 且帶信給以色列王,信上說:「我打發臣僕乃縵去見你,你接到這信,就要治好他的大痲瘋.」 5 : 7 以色列王看了信就撕裂衣服,說:「我豈是神,能使人死使人活呢?這人竟打發人來,叫我治好他的大痲瘋.你們看一看,這人何以尋隙攻擊我呢?」

那位以色列小女子明明介紹亞蘭元帥去撒馬利亞找先知,為甚麼元帥卻去了找以色列王呢?軍政人物總是以國家的概念去理解事情,因此亞蘭的元帥及國王自動地把醫治大痳瘋的事情提升到國與國的層次.把事情提升到國與國的層次,對於交戰的以色列和亞蘭來說,即是敵我關係,難怪被打得喘不過氣來的以色列王會如此驚慌失措.以敵我關係去處理目前的情況,似乎只是進一步強化敵我界限,而無助於醫治亞蘭元帥的大痳瘋病.這位亞蘭元帥能夠跨越敵我的界限回應以色列小女子的善意嗎?

習慣了以敵我關係看事物的人,真的不容易跨越敵我的界限.這位亞蘭元帥最終到到了先知以利沙的門前,不過,一聽到先知的指示:「你去在約旦河中沐浴七回,你的肉就必復原,而得潔淨」,他就大發雷霆,以為對方藉此侮辱自己.這位元帥一面走一面罵:「我想他必定出來見我,站著求告耶和華－他神的名,在患處以上搖手,治好這大痲瘋.大馬士革的河亞罷拿和法珥法豈不比以色列的一切水更好嗎?我在那裡沐浴不得潔淨嗎?」(5:11-12)

元帥為甚麼抗拒先知醫治他的大痳瘋的方法呢?因為元帥仍然以敵我的意識去理解先知的指示,「為甚麼要在你們的河裡面沐浴?我們的河水怎會比不上你們的河水?簡直豈有此理!」幸好,他有一位對他忠心及能夠跨越敵我界限的下屬,否則,他會錯失了醫治大痳瘋的良機.這位為他著想的下屬冒險勸他:「我父啊,先知若吩咐你做一件大事,你豈不做嗎?何況說你去沐浴而得潔淨呢?」這位下屬知道重點不是在那條河,因此不需要比較,也不需要提升到國家意識,重點是先知,以及他所說的話.他能夠分辨得出先知所代表的是上帝,而不是某個政權,或者某個國家,因此,可以放下亞蘭和以色列之間的敵我狀態去回應先知的指示.

亞蘭元帥的大痳瘋得到醫治的故事,是一個跨越敵我界限的故事.是誰跨越了敵我界限呢?不是元帥,不是亞蘭國王,不是以色列國王,而是無名無姓的小人物,一位被擄的以色列小女子,以及一位元帥的僕人.

國與國之間的衝突,人與人之間的糾紛,製造及遺留下很多敵我界限.我們要繼續維持,甚至強化這些敵我界限嗎?如果不是,又如何面對及化解這些敵我界限呢?靠甚麼力量可以去化解這些敵我界限呢?

過去的聖誕節及新年期間,我在泰國參加短宣時,遇到一位從台灣來的宣教士.原來,這位宣教士被差派往泰國之前,已經在日本作了十年的宣教工作.她說,有些台灣人想起日治時期,會想到「海角七號」般的愛情故事,但她想到日治時期,想到的是自己的家人被日軍殺害,自己的親戚被捉去當慰安婦,在她心中湧出的情緒是憤恨及悲痛.因此,當聽到上帝呼喚她去日本宣教時,他回應上帝說,去甚麼地方也可以,一定不可以去日本.經過多年的掙扎,這位宣教士被上帝的慈愛化解了心中的悲憤,踏出跨越敵我界限的腳步,與她的丈夫走上往日本宣教的旅程.在日本的十年時間,這對宣教士夫婦親身接觸日本民眾,與日本教會建立互相信任的關係.記念第二次世界大戰結束五十年的時候,日本教會草擬了向中國人的道歉信,請這對宣教士夫婦翻譯為中文,期望中國人能夠閱讀到及接受他們的道歉.

總結 今天,我們閱讀了兩個醫治大痳瘋的故事.相信大家都會發現,這兩個故事珍貴的地方,不在於教導我們如何醫治痳瘋病,而在於啟發我們跨越阻礙人性發展美善的界限,這些界限可能是「潔淨」與「不潔淨」,可能是敵我的態度,也可能是某些集體潛意識,我們需要一個比這些界限更高的角度,比這些意識更廣闊的視野,以及比個人利益及群體利益更超越的上帝,才能跨越這些界限帶給人們的限制及傷害.



\newpage

\section{吃還是不吃}
\label{sec:140}
\textbf{賴品超教授 - 吃還是不吃}
\newline
\newline
\hyperref[sec:139]{< < < PREV SERMON < < <}
~
\hyperref[sec:toc]{[返目錄]}
~
\hyperref[sec:141]{> > > NEXT SERMON > > >}


哥林多前書 8:1-13
\newline
\begin{longtable}{cl}
\hline
\hline
章節 & 經文 (和合本修訂版)\\
\hline
8:1 & \begin{tabularx}{0.7\textwidth}{X} 關於祭過偶像的食物，我們曉得「我們都有知識」，但知識使人自高自大，惟有愛心能造就人。 \end{tabularx} \\ \\ \relax
8:2 & \begin{tabularx}{0.7\textwidth}{X} 若有人自以為知道甚麼，他其實仍不知道他所應當知道的。 \end{tabularx} \\ \\ \relax
8:3 & \begin{tabularx}{0.7\textwidth}{X} 若有人愛神，他就是神所認識的人了。 \end{tabularx} \\ \\ \relax
8:4 & \begin{tabularx}{0.7\textwidth}{X} 關於吃祭過偶像的食物，我們知道「偶像在世上算不得甚麼」；也知道「神只有一位，沒有別的」。 \end{tabularx} \\ \\ \relax
8:5 & \begin{tabularx}{0.7\textwidth}{X} 雖然在天上或地上有許多所謂的神明，就如他們中間有許多的神明，許多的主， \end{tabularx} \\ \\ \relax
8:6 & \begin{tabularx}{0.7\textwidth}{X} 但是我們只有一位神，就是父，萬物都出於他，我們也歸於他；並只有一位主，就是耶穌基督，萬物都是藉著他而有，我們也是藉著他而有。 \end{tabularx} \\ \\ \relax
8:7 & \begin{tabularx}{0.7\textwidth}{X} 可是，不是人人都有這知識。有人到現在因拜慣了偶像，仍以為所吃的是祭過偶像的食物；既然他們的良心軟弱，也就污穢了。 \end{tabularx} \\ \\ \relax
8:8 & \begin{tabularx}{0.7\textwidth}{X} 其實，食物不能使我們更接近神，因為我們不吃也無損，吃也無益。 \end{tabularx} \\ \\ \relax
8:9 & \begin{tabularx}{0.7\textwidth}{X} 可是，你們要謹慎，免得你們這自由竟成了軟弱人的絆腳石。 \end{tabularx} \\ \\ \relax
8:10 & \begin{tabularx}{0.7\textwidth}{X} 若有人見你這有知識的在偶像的廟裡坐席，而這人的良心是軟弱的，他豈不放膽去吃那祭過偶像的食物嗎？ \end{tabularx} \\ \\ \relax
8:11 & \begin{tabularx}{0.7\textwidth}{X} 因此，基督為他死的那軟弱弟兄，也就因你的知識沉淪了。 \end{tabularx} \\ \\ \relax
8:12 & \begin{tabularx}{0.7\textwidth}{X} 你們這樣得罪弟兄，傷了他們軟弱的良心，就是得罪基督。 \end{tabularx} \\ \\ \relax
8:13 & \begin{tabularx}{0.7\textwidth}{X} 所以，食物若使我的弟兄跌倒，我就永遠不吃肉，免得使我的弟兄跌倒了。 \end{tabularx} \\ \\
[1ex]
\hline
\hline
\end{longtable}
日期:2012年3月4日

題目:吃還是不吃

經文:哥林多前書8:1-13

證道:賴品超教授

因為工作的需要,經常與不同宗教(包括佛教,印度教,伊斯蘭教)的人吃飯,甚煩,諸多規條.表面上看,基督教很自由,講恩典而不是律法,甚麼都可吃,因此 變成不太留意,但其實很重要.因為"吃還是不吃",不是一個抽像的問題,而是一個十分實存 (existential),就有如 ``to be or not be?" 一樣.就如人類始祖面對的試探,也是一個"吃還是不吃"的問題,而對於這個問題的抉擇,影響了全人類.基督受試探也與吃有關.對人類學的研究來說,飲水文 化很重要:"we become what we ate/eat." 吃還是不吃?吃甚麼?怎樣吃?跟誰吃?這些都會與基督徒的身份有關；例如:耶穌與施洗約翰的一個重要區分也在吃,而基督教最重要的禮儀(聖餐)也是與吃有 關,七宗罪中之一是貪吃(Gluttony),神學史上不少爭議(包括宗教改革)也跟吃有關.

今天的經文也是處理這樣的一個十分實際的問題:對於曾經獻給偶像(異教的神)的祭肉,作為一個基督徒,是"吃還是不吃"?

為何有此問題?背景:在哥林多這一類的城市,異教崇拜盛行,  不少(但不一定全部)在市場上買到的肉都曾獻祭,基督徒可否買來吃? 參加(未信)朋友的宴會,也有機會吃祭肉,不吃好像很不友善,不禮貌,吃又會使人誤會(基督徒為了美食,也吃祭過異教的神的肉,信仰可以很隨便),吃還是不吃? 在節期(神誕)可在廟堂裡吃免費的祭肉,"太公分豬肉"人人有份,市民基本福利／權利,不吃白不吃,這甚至是窮人少有的吃肉良機,但坐在廟裡吃會使人以為自己也拜偶像,作為基督徒"吃還是不吃"?對於第一個問題,在林前10 : 25-26的回答是:只管照買照吃,反正不知來源(是否祭過)就不要問(當時沒有產地來源認證／標籤,眼不見為淨),更重要的是"因為地和其中所充滿的都屬乎主".

對於第二個問題,林前10:27-30的答案相似:不必問,只管吃；但若有告訴你是祭過偶像的,那就不要吃,而這是為了別人良心的緣故.

對於第三個問題,在林前10: 14-22,嚴斥與異教的神(鬼)有相交,吃鬼的筵席,答案似乎就是簡單直接的"不要吃".

然而在今天讀的經文(林前8:1-13)似乎有點不一樣,基本原則是與林前10:31-33一致,但不是如此直截了當的講"吃還是不吃",因為要先處理哥 林多教會中某些自以為很有"知識"的人的立場,就是8:4-6所講的神學立場.(有點像:酒肉穿腸過,佛祖心中留.)這神學立場其實與林前 10:26所講的基本上是一致的.保羅並不反對此一立場或知識,但他要指出的是,我們的行動／抉擇取捨,不一定是直接取決於我們的這種特殊的,屬靈的"知 識".

林前8:1率先聲明,在這樣的問題上,優先考慮的是"愛心"而不是行動者個人的"知識"(10:1b-3).只有"知識"而沒有愛心只會自高自大,自以為 ``有知"的其實"無知",這種"知識"是虛幻而不實在,對自己,對別人都無益；　對比出"愛心"能造就人,而愛上帝的人更是上帝所知道的.

在上基本原則上,再要考慮其他人不一定有此"知識",他們,尤其拜慣偶像的,會認為這些肉是祭過偶像,是污穢的(8:7).如果他們看見你們這些自以為有 知識的人,也在廟裡坐席吃祭偶像的肉,就會放膽去吃,以為基督徒拜偶像也沒有問題,最後是變相鼓勵其他信徒拜偶像,使他們沉淪.你的自由成為絆倒他們的絆 腳石,你雖不殺伯仁,但伯仁卻是為你而死.你得到甚麼好處?就是那些"不能叫上帝看中我們"的食物,就是"不吃也無損,吃也無益"的肉而已!這值得嗎?! 因此,為了別人的好處,還是不吃!

保羅並不否定這種知識的本身,但換言之,保羅是容許信徒有不同的神學理解,但在實踐上,則比較嚴謹,在此,保羅是說,即使是按照你這自以為有知識 (gnosis, 某種特別的屬靈的知識,慧根)的人的理解／慧解／洞見(但這不一定等於保羅自己的立場),有點像老子／道德經講的"知其雄,守其雌",為免爭執,我們不一 定要堅持完全按照我們所知的來做.

搶包山,本來是祭鬼,給鬼吃的,現在是人特意買來吃基本原則,按著愛,而不是知,accommodation ethics, 路德:基督徒．不在律法下,是自由的,但也是眾人的僕人.同樣原則應用在"工人得工價",以愛去運用自由,權利.

更多的是考慮的,不是個人的喜好／自由／權力(不是老子花得起錢去吃,就是對,might =X+ right),而是也要考慮個人的行為對弟兄姊妹(信仰群體)他人的影響.

吃不吃曾獻祭的肉,只是其中之一,去街市獻祭的,可能更平(付出代價),如果所有都曾獻祭,是否改吃素,what"s the problem! 原則如此,應用不易.

今天德國不少天主教地區仍保留此傳統──周五吃魚.最初不明白,後來在一書找到,早期教會有人寧吃較平的(例如魚──不是為鮮味),甚至不吃,以賙濟窮人 (Carter Lindberg, Love: A Brief History through Western Christianity, p.43).

到穆斯林群體中的宣教士,很可能要放棄吃豬肉.同樣道理,重點在有節制.

Why? 貪食者,傾向於以自我為中心,不理別人,失去對其他事物的興趣,失去自控的能力,容易在試探中失敗跌倒.因此,作為靈性操練,克服肉體的情慾,學習自我卑微,攻克己身叫身服我,這是重要的操練.

人的欲望有很多,為何特別提出食?食色性也.但性方面,沒有也不危害健康,也很少人會天天面對這方面的試探.但貪食,界線難定,人人需要不同,食量多少因 人而異,難客觀去定,更不一定有即是效果,一餐吃多一點不易察覺(尤其中國人食法),人又不可能不吃,因特別難自控,這是每天都面對的試探.

或許有人覺過時,教廷最近提出了新的七宗罪(more modern / social)沒有了貪食: Environmental pollution, Accumilating excessive wealth, Inflicting poverty, Drug trafficking (with a k) and consumption, Morally debatable experiments, Violation of fundemental rights of human nature, Genetic modification.但我認為,貪食的誘惑,並不過時,不單與己相關──身體與靈性,並且與社會十分相關,也與新的七宗罪有關.

箴言:先想想,在誰面前吃,有何陰謀(吃政治飯,背脊骨落).美食當前,其實是陷阱,特別小心.其實我們身處的社會文化也是充滿陷阱.

首先,中國人 -- 凡是有腳的背向天者,要不是桌, 便是可以吃; 凡會飛的, 只要不是飛機, 便是可吃；煎 炒 炸 燉 炆 / 燜 蒸 灼 泡 焗 爆 燴 煙 煲 煨 涮 醃(還有很多其他)；中國文化之全球化＝飲食文化.「民以食為天」,如何使十多億人吃得飽,也用來講人權,人有吃飯的基本人權.

其次,現代生活,想一想每日接收多少廣告,電視.我們的病,有多少與飲食習慣有關,又要花多少錢和時間在食,和與食有關的事(預備及善後,如:減肥,處理垃圾).我們的問題,不是為衣食憂慮,而是缺乏節制.

我們是身處一個極度鼓勵人狂吃,又多又挑剔的社會文化.

貪吃,引來對身體的傷害(慢性疾病)是眾所周知,也會使人花費不少的時間去為了美食,浪費生命,(一個德國朋友:以中國人之聰明,如果少花點時間在吃,多 花點時間在科學上,成就肯定不止如此.)也帶心理上的問題:自我形象低.身心兩方面,互相影響.(心理學家指出,eating disorder暴食與厭食皆暴露人的心理問題, 不安, 缺乏安全感,往往難自控；要先解決心理問題, 才能解暴食與厭食問題, 太得閑不好, 無聊便吃, divert attention, understand the impact.)

貪食的問題並非僅為一個"個人修養"／健康的問題,也是構成嚴重的社會政治問題.古代羅馬, 為多吃, 故意嘔, 甚至另闢一室去嘔,反映其奢侈,對殖民地的剝削.在當今世界上仍有那麼多人生活在飢餓之中,如果我們的飲食(尤其中國人的飲宴)是帶來不必要的浪費,實在 缺德.

upiscentia,要將萬有都據為己有,不單是對具體事物,更可以是對知識,性愛以至權力.the unlimited desire to draw the whole of reality into one"s self. [Tillich, ST2:52-3] 這種情慾concupiscence最清楚的表現在食欲.究竟我們的飲食習慣是在強化這種文化,還是在抗衡它?

更重要的是,而是對生態的破壞.這不單是指食量太大所引起的糧食短缺,也是太咀尖,專吃一些奇珍,令某些生物瀕臨絕種.曾見海報,上面是"鴻門夜宴"下面 是"滅門慘案,大型珊瑚魚,鯊,犀牛,虎,熊.基因改造,也多少為了吃!過多包裝的食物,用了太多能源去運送的食物,也會增加對環境的傷害.這些都可以是 在我們不知不覺中發生!



\newpage

\section{天上曠野人間}
\label{sec:141}
\textbf{鄧瑞強博士 - 天上,曠野,人間}
\newline
\newline
\hyperref[sec:140]{< < < PREV SERMON < < <}
~
\hyperref[sec:toc]{[返目錄]}
~
\hyperref[sec:142]{> > > NEXT SERMON > > >}


馬可福音 1:9-13
\newline
\begin{longtable}{cl}
\hline
\hline
章節 & 經文 (和合本修訂版)\\
\hline
1:9 & \begin{tabularx}{0.7\textwidth}{X} 那時，耶穌從加利利的拿撒勒來，在約旦河裡受了約翰的洗。 \end{tabularx} \\ \\ \relax
1:10 & \begin{tabularx}{0.7\textwidth}{X} 他從水裡一上來，就看見天裂開了，聖靈彷彿鴿子降在他身上。 \end{tabularx} \\ \\ \relax
1:11 & \begin{tabularx}{0.7\textwidth}{X} 又有聲音從天上來，說：「你是我的愛子，我喜愛你。」 \end{tabularx} \\ \\ \relax
1:12 & \begin{tabularx}{0.7\textwidth}{X} 聖靈立刻把耶穌催促到曠野裡去。 \end{tabularx} \\ \\ \relax
1:13 & \begin{tabularx}{0.7\textwidth}{X} 他在曠野四十天，受撒但的試探，並與野獸同在一起，且有天使來伺候他。 \end{tabularx} \\ \\
[1ex]
\hline
\hline
\end{longtable}
日期:2012年3月11日

主題:天上,曠野,人間

證道:鄧瑞強博士

經文:馬可福音1:9-13

各位弟兄姊妹,主內平安.

今日是基督教聖曆裡的大齋期(Lent)的開始,大齋期是紀念主耶穌為了承擔人類的罪,背起十架,成就救贖的故事.這是一個偉大的故事.一個卑微的人,面對令人沮喪的政治形勢,遭遇叫人心寒的人性醜陋,以一種順從天命的勇氣,背起人類罪惡的十架.沒有這個人,塵世仍暗淡無光.沒有這種生命的承擔,苦難的世界只會叫人絕望.沒有人背負這十架,罪惡的力量將橫行無度.

面對偉人的人物,我們很想知道他的過去是怎樣的.有人喜歡喬布斯,便很想了解他如何在一間車房裡開展他的「蘋果」事業.耶穌,這個偉大的生命,他的生命又是如何開始的?今日,我們看看馬可福音的記載,看看耶穌的事業,是如何開展的.

經文:馬可福音1: 9-13

1:9 　那時,耶穌從加利利的拿撒勒來,在約旦河裡受了約翰的洗. 1:10　他從水裡一上來,就看見天裂開了,聖靈彷彿鴿子,降在他身上. 1:11　又有聲音從天上來,說:「你是我的愛子,我喜悅你.」 1:12　聖靈就把耶穌催到曠野裡去. 1:13　他在曠野四十天,受撒旦的試探,並與野獸同在一處,且有天使來伺候他.

經文很短,但意義深長.

(1)「天上」 耶穌的使命,始於他的受洗.他來到約旦河那裡,受施洗約翰的洗.他從水裡一上來,就看見天裂開了,聖靈彷彿鴿子,降在他身上.

受洗,是一件關鍵事件.這是人歸屬神的一個外在記號.耶穌以洗禮表明他屬於神.這時,天裂開了.「天裂開」,是神的行動,意味著神和人直接的溝通.「天裂開」,也可以是標誌著神行事的關鍵時刻的到來.先知以賽亞曾說,在歷史的關鍵時刻,神會裂天而降(參賽64: 1).耶穌的受洗,是歷史的關鍵時刻,神在這刻有所作為.「天裂開」了,主耶穌的生命,此刻也豁然開朗.他腳踏大地,心繫上天,他立於天人之間.一個心靈連於上天的人,他活在塵世,卻不會容讓自己被塵世的力量控制.他活在世上,卻會活出屬天的氣質.面對一個偉大的人,我們總覺得,他頭上的天是打開了的.他的生命,散發出上天的色彩.

「聖靈彷彿鴿子,降在他身上」.鴿子,象徵著神的降臨.在日本,神廟有「鳥居」(Torii),「鳥」是天神的使者,故神廟有「鳥居」,讓「鳥」降臨.耶穌受洗時,聖靈彷彿從天而來的使者,降在他身上.這解釋了什麼力量驅動耶穌的生命.從天而來的聖靈,驅策著耶穌的生命.

偉大的生命,很多時都跟隨著一種神秘的力量.

藝術家創作時,有時是被「靈感」驅動的.作家茨威格(Stefan Zweig)講述他閱讀莫札特的手稿時的感受.他說,莫札特的手稿幾乎毫無改動的痕跡.看起來,莫札特心裡已有整篇樂章,他不假思索,一揮而就,便將整篇樂章寫下來.看來,有一種偉大的力量駕馭著莫札特,他只是順著這力量,便寫出天籟一樣的樂章.(參茨威格:《昨日的世界》)

魯迅有篇文章叫＜過客＞.文章講到,有一個過客,走了很長的路.在路上,他遇到一個老翁.過客問老翁,「前面是什麼?」老翁答:「前面是墳.」老翁看這過客已很疲累,便對他說:「回頭吧.」過客答:「那不行,我只得走.回到那裡去,就沒一處沒有名目,沒一處沒有地主,沒一處沒有驅逐和牢籠,沒一處沒有皮面的笑容,沒一處沒有眶外的眼淚.我憎惡他們,我不回轉去.」老翁說:「你真的要繼續走.」過客答:「是的,我只得走了.有聲音常在前面催促我,叫喚我,使我停不下來.」老翁說:「太陽下山了,你還是休息一會吧.」過客說:「那前面的聲音叫我走.」有些人的人生,跟從著一種屬天的呼喚.

魯迅寫透了那些偉大的生命,那些生命總是因為回應屬天的呼喚而離開腐爛不堪的光景,而不畏死亡地走向光明的境地.

耶穌,他將聽從屬天的呼喚,走向他人生的十架.馬可福音在這裡明確點出,驅動耶穌生命的內裡的聲音,就是聖靈.問題只是,我們中間,又有多少人能夠不聽世界的聲音,而聽聖靈的聲音而行.

在耶穌的受洗裡,最後,有聲音從天上來.馬太福音記述同樣的情節時,講到從天上有聲音說:「這是我的愛子,我所喜悅的」.馬太福音的那種講法,多少將這裡發生的事件描述為一客觀事件.但馬可福音這裡的講法,講到天上的聲音說:「你是我的愛子,我喜悅你」,這是神對耶穌直接講的.馬可福音將這描述為發生在主耶穌生命裡的一件主觀事件.他面對神,他聽到天上的神對他說話,他確定他的身分,他明白他的召命.

在洗禮裡,耶穌完全領悟他的身分.他日後的人生,就是這身分的展現.

身分認同,的確是現代社會的一大問題.在香港,「香港人的身分」和「中國人的身分」,有時也會引起某種緊張.為什麼某種身分的人可以在某名店前拍照,其他身分的人便不可以?為什麼某種身分的人可以在這裡出生及定居,其他身分的人便不可以?身分認同所引起的身分政治問題,將繼續成為社會的話題.

在耶穌的受洗裡,他確定他的終極身分.他是「神的愛子」.這是一個超越種族,膚色,地域的身分.他將以這一身分活著.他活在地上,有地上的身分,會服從地上的政府,但當政府不仁不義時,他會持守「當順從神而不順從人」的原則,他會活出那「神的愛子」的身分.他活在世上,有地上的家人,會愛他的家人,但這不會妨礙他視其他人為他的家人,不會妨礙他對其他人的愛.他那「神的愛子」的身分,凌駕一切其他身分之上.他以這身分,活出他的一生.

問題是:我們以什麼身分活在世上?我們是學生,是老師,是僱員,是老闆,然而,我們有多少人能以「神的兒女」的身分,凌駕這一切身分之上?我們有多少人,明白自己真正的身分是「神的兒女」,並活出這身分?

一個人能體會自己是「神的兒女」,並活出這身分,多麼難得.

今日的講道題目是:天上,曠野,人間.「天上」指的,是我們都像耶穌一樣,面對天上的神,倚靠屬天的力量,活出屬天的身分.我們的身分,是「屬天」的.這是我要講的第一點.

(2)「曠野」

1:12　聖靈就把耶穌催到曠野裡去.

在耶穌的救世事業開始之先,他內裡的聖靈催迫他到曠野.這是不得不走的一步.經過曠野的洗禮,耶穌才能開始他的事業.

「曠野」是聖經常用的意象.當以色列人出埃及時,他們經過曠野.這是離開「為奴之地」進入「自由之地」要經過的地方.這是生命轉化必需經過的階段.在曠野裡,以色列人面對種種艱難和考驗,有人回頭,有人埋怨,有人反抗神.住在沙漠裡的靈修大師提示人,要進入沙漠或曠野,必須千萬小心,因為在這種一無所有的地方,人會面對真正的自己.人性裡面的驕傲會跑出來吞噬人,人性深處的孤獨會使人瘋狂,人性的各種慾念會浮現出來叫人不能自拔.

在這一無所有的「曠野」裡,人不單面對真實的自己,同時,人也會面對真實的神.沒有城巿的光害,在晚上,我們便能看見天上的星星.同樣,在曠野,一切人為的阻隔皆沒有了,我們也能看見真實的神.當然,人性的內在弱點會蒙蔽人的眼睛,叫人看不見神.但是,只要人倚靠神的恩典,將生命全交給神,「曠野」就會變成人和神相遇的最親近的地方.

在這一無所有的「曠野」裡,人一無所有,也一無所靠,只能倚靠神.「曠野」是學習倚靠神的地方.在這裡,人不能靠自己的財富,人不能靠自己的學識,人不能靠自己的社會地位,人只能謙卑下來,回復生命的一無所有,只能仰賴神,只能在神那裡領受生命的無限價值.

聖靈催迫耶穌進入「曠野」,看來是要讓耶穌經歷「曠野」的考驗,讓自己成為一個一無所靠而只靠神的人.唯有這樣,他才能進入塵世的生活而不被塵世的生活牽引.很快,他便會進入人群,受群眾擁戴.唯有經歷「曠野」的考驗,他才能不為群眾的歡呼聲掩蓋,而能發出屬於「曠野」的真理的呼聲.很快,他便會與僵化了的宗教力量周旋,受他們挑戰.唯有經歷「曠野」的考驗,他才能擺脫制度的纏累,而向人指出真實的神.很快,他便會面對政治的不義,受其審判.唯有經歷「曠野」的考驗,他才能輕視強權的威嚇,而將生命單單交在神的手中.

在「曠野」,人變得一無所有,一無所靠,卻明白到神的支撐才是我們生命的終極支撐.在「曠野」,人像命懸一線,而神就是那不會斷的線.

現代人不喜歡真實的「曠野」.我們會將「曠野」改建聖城巿,我們會將「曠野」變成旅游的新奇經驗,我們會將「曠野」變成消費項目.我們不喜歡變成一無所有.我們的經濟活動就是要將我們從一無所有變成樣樣都有,我們的科技將多樣化的產品變成現實,我們的消費活動將多樣化的享受變成必需.在這個沒有「曠野」的世界,我們要抓住很多東西來肯定自己.但這些東西,卻成為心靈的「光害」,使我們看不見天上的星星.

或許,若我們要活出「神的兒女」的身分,我們要學習一種「活在曠野」的心態.在這裡,沒有炫耀,沒有驕傲,沒有將生命寄托在無意義的事情上.生命的雜質除掉了,留下最純淨的生命.這個生命坦然地活在神面前,心無掛礙.唯有這種生命,能一心一意完成生命中的天命.

今日的講道題目是:天上,曠野,人間.「曠野」指的,就是這種存活態度.能在花花綠綠的世界,持守「曠野」的心境,才能活出我們「屬天」的身分.這是我要講的第二點.

(3)人間

1:13　他在曠野四十天,受撒旦的試探,並與野獸同在一處,且有天使來伺候他.

耶穌在曠野四十天.在舊約,「四十」似乎是一個特別數字.以色列人出埃及時,在曠野四十年,受盡考驗,但也經驗神的供應.摩西在西乃山上四十天,與神面對面相遇.先知以利亞,走了四十晝夜,到了神的山,尋求神的幫助.看來,「四十」與考驗有關,也與面對神有關.

耶穌在曠野四十天,受撒旦的試探.敵對神的力量,在我們實踐真理的時刻,伴隨著我們.

奇怪的是,馬可福音說,有野獸和耶穌在一起.一般而言,野獸都不是善類.但在這裡,這些野獸看來都是馴服的.人和野獸和平共處,這是天堂的呈現.在伊甸園,人和野獸同在.先知以賽亞講到末世的景象時,也提到人和野獸同處.看來,耶穌在人間,是要帶來天堂.再者,這裡有天使伺候,這是天人合一的景象.

這是耶穌使命的方向.縱使撒旦常在左右,但這無礙耶穌在人間實現天堂一樣的境況.在這裡,人和獸,或者說,人和最凶殘的異己,能在天使的幫助下,和諧共處.這是愛的呈現,這是天堂的美景.

我們活在世上,這不也是我們生命的方向嗎?

今日的講道題目是:天上,曠野,人間.「人間」指的,就是我們的人間使命,就是帶來寬恕,包容,愛.這是我要講的第三點.

馬可福音的耶穌受洗故事,點出了耶穌「屬天」的身分,「曠野」的生活方式,「人間」的使命.但願,這也是我們生命的故事,堅定我們的「屬天」身分,活出「曠野」的心境,實現「人間」的天國.

但願榮耀歸於聖父,聖子,聖靈.阿們.



\newpage

\section{心靈誠實與祭祀}
\label{sec:142}
\textbf{陳曉東博士 - 心靈,誠實與祭祀}
\newline
\newline
\hyperref[sec:141]{< < < PREV SERMON < < <}
~
\hyperref[sec:toc]{[返目錄]}
~
\hyperref[sec:143]{> > > NEXT SERMON > > >}


日期:2012年3月18日

主題:心靈,誠實與祭祀

證道:陳曉東博士

上世紀二十年代,有一位挪威的宣教士,名叫艾香德,在湖南傳道.他在傳道的過程中,有一個很深刻的悔改經歷.就是當他作為一位西方的宣教士,進入了中國的土地文化裡,被中國文化的豐富和豐厚,深深震撼.他當時有一個經歷,就是在不同的名山大川,特別到佛教的寺院去探訪.在他自己所寫的日記裡記載著,他被一班中國基層僧人純樸的笑容,和那份慈悲仁愛的心,深深震撼.他自己坐下來反省,在基督的信仰裡,有沒有這種慈悲,接待遠人,愛鄰舍,將上好的生命與人分享的道理在其中?

他發現有很多很多,所以他立了一個心願,就是我們作為一個基督徒,怎樣能夠有一個款待遠人的地方 ( house of friendship ),使生命旅途中的雲遊者──不是發白日夢的魂遊,而是佛教詞語中「四處去尋找真理的人」──有一個歇腳的地方,我們可以坐下來談天論道.乃至於每個人的生命,可以在彼此的豐富和啟迪中,尋找到上主那份真理,美善,和祂的美麗.因此他開始了這個事工,最初在南京.不過南京的工作結果很慘痛,因為軍閥混戰而被摧毀.當時他改了一個名字,叫「景風山」,但事工失敗了.後來他輾轉去到杭州,杭州的工作剛開始,建了一個「天風山」,但又失敗告終.

在三零年代他和同工來到沙田──當時一個不毛之地,在標號五字的山丘上,用了三千墨西哥銀洋買下這個山丘.當時這三千元也很多,這樣便成就了今天「道風山」這事工.當然這是很久以前的歷史,中間經歷了很多變遷.但在今天,山上有很多事工.大家去到道風山的時候,可能見到信義宗神學院在那裡,可能見到漢語文化研究所在那裡,基督教文化研究所在那裡,也見到我們很樸實的同工在打理場地,那是道風山服務處.但歷史最悠久的機構,就是道風山基督教叢林.

我們依然在那裡服侍,依然秉持先賢他們所持定的異像,繼續在信徒和教牧心靈的發展,作靈修指導.有幾件事工和大家相關的,包括藝術的繼續發展,山上有藝術的工作,正在繼續開展.並且我們有一個遠像,這遠像和香港眾教會關注的問題相關,就是促進香港作為人民社會,有人性的社會發展.我們一同協手努力,令耶穌基督的名字,上主的慈愛,在這城市裡,令人看見,並觸摸到.令這位道成肉身的主,在我們這世代裡,透過眾教會的努力,能使眾人見到以為美,對這社會有好的貢獻.

當我自己分享這歷史的時候,再看今天我們所講的福音書,我看到歷世歷代,不同時代的信徒,對今天所講的,記載在約翰福音耶穌的這事件,都有那種和應.今天我們所讀的福音書,其實是一個生命得到釋放的故事.經學家告訴我們,這事情發生在一個正午時份.為什麼說是正午時份呢?因為在第四章第八節講到一件事情,是門徒進城買食物去,應該不是去買晚飯,也不是買早飯,而是一個買午飯的時間.同時正是烈陽高照的時候,在一個雅各井旁,有一個婦人來這裡打水.聖經特別透過這婦人的口,講明一件事,就是這婦人其實並不是在她自己的村莊打水.她故意走了很遠,因為當她和耶穌對話時,她說:「如果你早些告訴我,我便不用走那麼遠,來這裡打水.」所以我們知道,她走了很遠來打水.

第三方面,我們知道這婦人生命中充滿了憂傷,是生命充滿了痛苦的一個故事,為什麼?雖然聖經輕輕題過這事情,但當我們仔細聽耶穌指出她生命的問題時,我們知道這婦人在生命所遭遇的故事,並不是普通人可以承受,因為耶穌和她說:「你去叫你丈夫也到這裡來.」婦人就說:「我沒有丈夫.」我們知道她是一個寡婦,耶穌跟著和她說:「你說沒有丈夫是不錯的.你已經有五個丈夫,你現在有的並不是你的丈夫.你這話是真的.」這對話非常震撼,乃至那個婦人整個人,因為耶穌這句說話反轉過來.

想像一個撒馬利亞的少女,在那年代西亞的文化裡,大約十四五歲,當一個少女身體成熟,便可以出嫁.在一個農牧的社會裡,女性能夠繼續生活的方式,只能是找一個好人家,嫁了依附在一個男權的家庭裡.十四五歲的少女,花樣年華,對生命充滿憧憬.現在父親母親將她出嫁,為了她的未來,為了她的前途,嫁入一個家庭裡.可能依循傳統風俗,相夫教子,能夠安安樂樂過生活,這是每一個正常的人,心裡的渴望.但是嫁了一個丈夫,我們不知道這個丈夫為什麼離開她.最合理的推斷,是患了疾病,離開了世界.如果不是,她就是犯了大錯,給人休妻回家.

但是故事無論是丈夫離開世界,或是她犯了什麼大錯被休回家,生命須要繼續前行.所以在這個家庭裡,將一個十四五歲的少婦再嫁出去,可能第一次是如花似玉,第二次已經不值錢.對不起我這樣說,那時的女人的確是和錢掛鉤的.一次的出嫁,被送回家.兩次的出嫁,被送回家.三次的出嫁,可能連父母也不想要這個女兒.四次的出嫁,父母可能年紀老邁,無力養她,也可能已離開世界.五次的出嫁,再返回她自己的家鄉,還有誰想要她?無怪乎耶穌在這裡說:「你現在有的並不是你的丈夫.」是的,生命裡勉勉強強,隨隨便便支撐著,可以生活下去便算了.這是生命,這是生涯.

我們很多人的生命,其實並不如我們所期望的,朝著一個美麗的方向,一個非常成功的方向──我們現在常聽到的一個字眼──去發展.生命裡原來有這麼多慘痛,而且這並不是一個小數目.「你現在有的並不是你的丈夫.你這話是真的.」乃至於婦人遇到一個真正明白她生命的人,她從心底裡就說:「先生,我看出你是先知.我們的祖宗在這山上崇拜,你們倒說,應當崇拜的地方是在耶路撒冷.」是的,你們的宗教,你們的信仰,門庭高大,聖殿輝煌.你們去參加崇拜的人,個個衣冠楚楚.我們這些是被生活,生命逼迫到牆角,頭也抬不起做人的一班人.你是先知,你信的宗教很偉大.不過這與我無關,多謝你.我繼續要打水,我繼續要返回我自己寄居的家裡,繼續去生活.

但是耶穌在今天的經文裡,講了一番說話,是令到整個生命,或是令到整個人類歷史,心靈得釋放的一件事情.祂說:「時候將到,如今就是了,那真正拜父的,要用心靈和誠實──或是要用心靈和真理拜他,因為父要這樣的人拜他.」在教會很長的歷史裡,當我們讀到這段經文的時候,我們便有一種錯覺,說耶穌正在講祭事被廢除了,從此大家自由了.所以在任何一個地方,只要有二三個人,大家用真理和誠實拜祂,神便會在我們中間.這說法不錯,不過是有偏頗的地方.

我們看到當耶穌替這婦人做心靈醫治的時候,耶穌帶領這婦人,從那種生命幽暗和痛苦中走出來,依然肯定一件事情,就是崇拜是很重要的.來到神的殿中,和眾弟兄姊妹,和眾人一同敬拜這位創天造地的主,是非常重要的.但這個一同敬拜,一同崇拜的生命,有一個要求,是用我們的心靈,在真理裡一起崇拜祂.今天我們來到神的殿中,我們一同默想神的話語,有舊約的經課,有新約的經課,有詩篇幫助我們思想,有福音書幫助我們去明白.在剛才我們仔細去聆聽今天的經文時,有一個圖像,不斷在這幾節經文中出現,這圖像是「山」.

耶穌要帶領這婦人,重新進入屬靈的聖山上.將她的視野重新打開,將她生命的內涵重新去挖掘出來,將她的心態重新去調教,將她自己對世界的判斷重新修正,將她自己人生的行動重新導航,將她自己已經心如死灰的生命,重新賦予生命的盼望.我相信這是主耶穌──當祂和這婦人談道的時候,同樣對我們歷世歷代,願意去耹聽這位大先知,耹聽這位大祭司話語的人,使我們能夠從祂那裡學到的功課.祂呼喚著我們每一個世代的人,重新在祂裡面找回生命的視野,找回生命的內涵.重新幫助我們調教我們自己的心態,去作出一個好的判斷,去作出一種好的行動.在遇到困難時,生命不會因此而被打沉,而是有一種充滿盼望的生命在內,這真是一件不容易做的事情.

今天早上,我閱讀台灣聯合報的社論時,很特別,是關於一位剛去世的藝人,名叫鳳飛飛.或者我們在座的人都不認識鳳飛飛,她的名字在我們的生活世界很陌生.但如果曾經在七十年代,八十年代生活過,認識台灣流行曲的年青人,或者是中年人,曾經被一把聲音,有一些土氣,有一些沙啞,甜美起來但不似鄧麗君般甜到漏油,被這麼一把聲音吸引.同時會記起一位鄉下姑娘,時常穿得很鮮艷,最喜歡頭上載著一頂百變的帽子.她這個形像我覺得很有趣,她叫「鳳飛飛」.今早台灣聯合報的社論,說鳳飛飛選擇了用最沉默的方式離開這世界,告別了她自己的歌迷.

但這個被人遺忘的名字,觸發了整個台灣,一個很深的反省.一方面引起不同世代,對她很深的懷念,但另方面對我們身處的世界,帶來很深的反省.因為作為一個鄉下少女,她唯一的天賦,是她自己的歌喉,成為一個成功的藝人,用她的聲音歌唱她那一個時代.但她退隱之後,作為一個藝人,她相夫教子.她尊重曾經愛護過她的歌迷,到一個地步,當她面對人生最後一刻──她離世時六十歲,癌症離開,她選擇不驚擾她至愛的親人和歌迷,她跳脫一切世俗的禮儀,她用一種安靜的完美作一個告別.所以整個世界在她死後四十多天,才知道她離開了.現在她的骨灰被送回鄉下靜靜地安葬,她用她的細微,體貼和周到去報答曾經愛護過她的歌迷.

當進行這個反省的時候,其實是帶出一個我們生活世界裡所面對的囂張,喧囂,唯恐不用其極地告訴人:「我在這裡.」我要你注意到我是一個重要的人,我要你告訴全世界的人:「沒有我在這個位置上,你們全都要死.」這種價值,這種心態,瀰漫在我們現代社會裡,瀰漫在我們自己所身處的社會當中.每天打開報紙的時候,你所見到的都是不同的名人,不同的娛樂界中人,不同的名星.無論是政壇名星,或娛樂界的名星,或商界的名星,他們告訴我們所有人:「你要集中注意在我身上,你要看著我,我要告訴你,我多麼成功,你要來膜拜我.」我們在一個安靜離世的藝人身上,有很多深思的地方,對我們今天基督教的信仰,同樣有一個很深的反省和思考.

在今天我們所讀的詩篇裡,我們看到詩人呼喚神的子民,作出一種最真實的敬拜.在詩篇第五十篇第一至第六節,當我自己仔細去默想這段經文的時候,有幾個字出現在我心靈裡.是的,在我們今天生活這個城市當中,是的,我們作為一班耶穌的跟從者,學習者,信徒,我們應該可以從詩篇五十篇,幫助我們看到一種生命的視野,我們可以從那裡學習到一種生命的內涵,我們可以調教到作為耶穌跟從者的一種心態,我們能夠作出一種好的判斷,行動和盼望.有人說:「高度決定深度.」我以前不明白這句說話,直至當我自己開始學懂爬山的時候.去登山的時候,我對這句說話,體會愈來愈深.其實我覺得香港人很有福氣,因為我們住在很多高山上,只要走不多遠,便找到很高的山.嘗試去攀爬的時候,我們便明白何謂「高度決定深度」.

今天其實雖然是主顯節最後一個主日,但在教會的年歷,也是耶穌登山變像的節氣.所以在教會的年歷裡,我們同時呼籲弟兄姊妹,一齊去登山.登上高山,去學習,去體會耶穌基督在高山上變像的經驗.不過今天大家選擇了坐在禮拜堂裡,其實我們依然可以用心靈去到高山上.當我們擁有一個視野的時候──你怎樣擁有一個視野?你須要站得高,站得高的時候,我們看問題就不會把自己陷於枝節裡.這種經驗,無論作為家庭主婦,或是作為要日理萬機的重要人物也一樣.我亦明白當我們在一個高度裡,看問題的時候,我們不容易有偏差,為什麼?當我們在一個高度的時候,我們能夠不被細節將我們引入歧路.

所以詩篇第五十篇第一節,告訴我們,我們的眼目,常常望著神的時候,我們生命的高度,就能夠有保障.大能者神──耶和華,是我們生命裡的焦點.因為我們所做的事情,是跟從著上主所賞賜給我們的智慧,能夠和祂同行.所以祂發言呼召天下,從日出到日落之處.在這個從日出到日落之處,生命本身要帶來什麼內涵呢?第二節告訴我們一個很重要的字眼,就是「美」,生命美不美,我們所做的事情美不美.你知道聖經裡講到美的時候,它同時顯露出它的真和善.所以第二節說:「從全美的錫安中,神已經發光了.」

上帝希望我們每一個人能看到,作為跟從祂的人,我們所追求的人生,無論是什麼形態的人生,是一個鳳飛飛藝術的人生也好.或是在撒馬利亞雅各井旁,落泊得在中午逃離自己的村莊,在井旁遇到耶穌的一個生命.無論是什麼生命,當我們和耶穌相遇的時候,我們所期望的,是在我們的相遇中,帶來生命的美麗,帶來生命的善良,帶來生命的真實.每一個跟從主的人,不是只說:「我拜你,我拜耶穌你.」不是的,而是在耶穌的生命裡,領受著一種從神而來的美,從神而來的真,從神而來的善.我們就是上主在這個世間的見證,我們就是上主在這個世界的美麗.當我們的生命不向著這高度而去的時候,我們便落在一個枝節的生命裡.

同樣聖經告訴我們,我們在神面前明白這道理的時候,第三節提醒我們,這位神是值得我們去敬畏的神.「我們的神要來,決不閉口.有烈火在他面前吞滅；有暴風在他四圍大颳.」乃至於當我們見到這位創天造地的主宰,在我們生命裡臨在的時候,我們曉得以敬畏來度過我們的人生.在我們今天所生活的世界裡,敬畏或者已經被邊緣化至一地步,我們認為我們所敬拜的那一位神,只是和我們家裡所養的土地公,或者和其他偶像差不多地位.在太多環境裡,我見到當我們進入神的聖殿時,我們好像返回客廳一樣.在很多的場合,我們見到信徒生命中失去對神的敬畏.弟兄姊妹,「我們的神要來,決不閉口.有烈火在他面前吞滅；有暴風在他四圍大颳.」

當我們生命裡,常存一種對上主的敬畏時,它能夠令我們儆醒,能夠令我們謹慎,能夠令我們不張狂,能夠令我們在祂面前存著謙卑的心,去愛神和做人.這個心態很重要,乃至於我們能夠知道,在人生裡,上帝有一天要審判我們.審判我們,我相信不是只看我們有什麼惡行沒有.如果是這樣,我們所跟隨的這位神,和一個訓導主任差不多.不是的,當我們講到上主呼召祂的子民,在這個世界裡做祂榮美的見證時,我相信一件很重要的事情就是:到底在我們生命裡,有沒有活出上主所期望的那份真理?有沒有活出上帝所期望的那份美麗?有沒有活出上帝所期望的那份善良?乃至於我們在神的台前站立的時候,就好像第四節:「他招呼上天下地,為要審判他的民.」這個審判祂民的時候,是看在我們的生命程現當中,有沒有這一份美,有沒有這一份真,有沒有這一份善.

所以無怪乎祂說:「招聚我的聖民到我這裡來,就是那些用祭物與我立約的人.」我們只看這裡,一定不明白,唯有當我們去思想,一個生命和耶穌基督相遇的時候,思想在撒馬利亞井邊,這個婦人和耶穌生命相遇的時候,所帶來的改變.拜祂的當用你美麗的心靈,當用你生命的真理來拜祂,我們便明白,當上主把祂聖民,召聚到祂面前,用生命的祭獻,好像保羅所講,將生命獻上那祭獻,與耶穌基督立約的時候,看到聖經為何把跟從耶穌的人稱為「基督徒」.因為這個生命是被邀請,披戴起耶穌基督的生命,進入了耶穌的生命.

今天我們一同來到神的聖殿,去記念,去思想,耶穌基督豋山變像的日子.在這山上,耶穌將祂的聖容,向歷世歷代跟從主的人顯現.但是這顯現,是須要用心靈和誠實去看見.在哥林多後書第四章第三節,保羅說:「如果我們的福音蒙蔽,就是蒙蔽在滅亡的人身上.此等不信之人被這世界的神弄瞎了心眼,不叫基督榮耀福音的光照著他們.」這句說話,說得雖然沉重,但提醒我們每一位願意在屬靈生命成長裡,耶穌的跟隨者,是的,這個世界有一個叫做世界之神.他最大的渴望,就是令每一位在這個世界,無論是信徒或非信徒,眼睛都黑暗了,乃至於基督榮耀的光照射的時候,我們看不到.

但是在今天主顯節的日子,我們一同在神面前,再一次定睛注目在那個叫人出死入生的十字架上.被十字架的沉重,十字架的希望,再一次塗抹我們瞎了的眼睛,乃至於上主的真光,可以照射在我們生命當中.可以給我們再一次得到生命的視野,再一次得到上帝所期望我們那種真,那種善,那種美的內涵.再一次幫助我們在神面前,有一種敬畏,乃至於我們能夠在上帝面前站立得住.要我們每一天生命的祭事,生命的行動,和永不止息的盼望,成為祂在這世間最美好的見證.

最後請容許我用保羅所講的這番說話,彼此勸勉,在哥林多後書第三章,保羅這樣說:「我們既有這樣的盼望,就大膽講說,不像摩西將帕子蒙在臉上,叫以色列人不能定睛看到那將廢者的結局.這帕子在基督裡已經廢去了.主就是那靈；主的靈在哪裡,那裡就得以自由.我們眾人既然敞著臉得以看見主的榮光,好像從鏡子裡返照,就變成主的形狀,榮上加榮,如同從主的靈變成的.」讓我們彼此互勉,在生命的旅程裡,榮上加榮.



\newpage

\section{婚禮的禮服}
\label{sec:143}
\textbf{伍渭文牧師 - 婚禮的禮服}
\newline
\newline
\hyperref[sec:142]{< < < PREV SERMON < < <}
~
\hyperref[sec:toc]{[返目錄]}
~
\hyperref[sec:144]{> > > NEXT SERMON > > >}


馬太福音 22:1-14
\newline
\begin{longtable}{cl}
\hline
\hline
章節 & 經文 (和合本修訂版)\\
\hline
22:1 & \begin{tabularx}{0.7\textwidth}{X} 耶穌又用比喻對他們說： \end{tabularx} \\ \\ \relax
22:2 & \begin{tabularx}{0.7\textwidth}{X} 「天國好比一個王為他兒子擺設娶親的宴席。 \end{tabularx} \\ \\ \relax
22:3 & \begin{tabularx}{0.7\textwidth}{X} 他打發僕人去，請那些被邀的人來赴宴，他們卻不肯來。 \end{tabularx} \\ \\ \relax
22:4 & \begin{tabularx}{0.7\textwidth}{X} 王又打發別的僕人，說：『你們去告訴那被邀的人，我的宴席已經預備好了，牛和肥畜已經宰了，各樣都齊備，請你們來赴宴。』 \end{tabularx} \\ \\ \relax
22:5 & \begin{tabularx}{0.7\textwidth}{X} 那些人不理就走了，一個到自己田裡去，一個做買賣去。 \end{tabularx} \\ \\ \relax
22:6 & \begin{tabularx}{0.7\textwidth}{X} 其餘的抓住僕人，凌辱他們，把他們殺了。 \end{tabularx} \\ \\ \relax
22:7 & \begin{tabularx}{0.7\textwidth}{X} 王就大怒，發兵除滅那些兇手，燒燬他們的城。 \end{tabularx} \\ \\ \relax
22:8 & \begin{tabularx}{0.7\textwidth}{X} 於是王對僕人說：『喜宴已經齊備，只是所邀的人不配。 \end{tabularx} \\ \\ \relax
22:9 & \begin{tabularx}{0.7\textwidth}{X} 所以你們要往岔路口上去，凡遇見的，都邀來赴宴。』 \end{tabularx} \\ \\ \relax
22:10 & \begin{tabularx}{0.7\textwidth}{X} 那些僕人就出去，到大路上，凡遇見的，不論善惡都招聚了來，宴席上就坐滿了客人。 \end{tabularx} \\ \\ \relax
22:11 & \begin{tabularx}{0.7\textwidth}{X} 王進來見賓客，看到那裡有一個沒有穿禮服的， \end{tabularx} \\ \\ \relax
22:12 & \begin{tabularx}{0.7\textwidth}{X} 就對他說：『朋友，你到這裡來怎麼不穿禮服呢？』那人無言可答。 \end{tabularx} \\ \\ \relax
22:13 & \begin{tabularx}{0.7\textwidth}{X} 於是王對侍從說：『捆起他的手腳，把他扔在外邊的黑暗裡；在那裡他要哀哭切齒了。』 \end{tabularx} \\ \\ \relax
22:14 & \begin{tabularx}{0.7\textwidth}{X} 因為被召的人多，選上的人少。」 \end{tabularx} \\ \\
[1ex]
\hline
\hline
\end{longtable}
日期:2012年3月25日

主題:婚禮的禮服

證道:伍渭文牧師

讀經:馬太福音 22:1-14

在 崇基堂主持婚禮,有一些難忘的趣事.有一次花仔領著伴娘和新郎行在通道時,突然躺在地上哭起來,引起交通阻塞,大衛緊張,後用波板糖利誘走畢全程,以後我 就定下規例:花仔,花女要四歲以上.中途阻塞,可以疏通,但另一宗趣事就較為嚴重了,是所有主禮團差不多到齊,還有十五分鐘便開始,新人才發現在匆忙中忘 記拿戒指.幸好新郎的家就在附近,結果婚禮遲了十五分鐘開始.沒有戒指,婚禮很難開始,因婚禮的核心是新人彼此的盟誓,而戒指就是盟誓的信物.藉看得見的 信物,把盟誓加上信印印在心中(ring is the sign of the vow, seal on our heart),像聖餐藉看見的杯和餅把恩典印在心中.婚禮忘記帶戒指,令到眾人都很焦急.我真的想過校牧室應買一對戒指備用,以防不測.

今天的經文,也談到婚禮的忘記.有一位赴王子的婚宴的賓客,沒有穿上禮服,被國王捉拿,綁起來,趕出去.這比喻明顯不是趣事,是一個悲劇.這位被逐的賓客,是否真的忘記穿禮服嗎?還是他以為穿上了禮服,可以赴王子的婚宴呢?

古 代國王的婚宴,國王會準備禮服給受邀的賓客,沒有王的禮服,根本不能進入宴會廳,禮服就是入場券.這人進了會場,坐在其他賓客當中,應該沒有人懷疑他沒有 禮服,甚至他可能都以為穿上了禮服.但當王進來「觀看」賓客,就發現那人沒有穿上禮服,從上文下理,我們可以推論,這「觀看」是一種審判:

一,赴王子的婚宴是迎見上主──是審判的時刻

(基督再來是羔羊婚宴,每人到審判台前被審) 請看馬太福音 22:12－13

「就對他說,朋友,你到這裡來,怎麼不穿禮服呢.那人無言可答. 於是王對使喚的人說,捆起他的手腳來,把他丟在外邊的黑暗裡．在那裡必要哀哭切齒了.」

(猶太人棄絕基督,在主後70年,提多將軍攻耶路撒冷,猶太人被殺及滅城) 再看上文:馬太福音 22:4－7

「王又打發別的僕人說,你們告訴那被召的人,我的筵席已經預備好了,牛和肥畜已經宰了,各樣都齊備．請你們來赴席.那些人不理就走了．一個到自己田里去．一個作買賣去．其餘的拿住僕人,凌辱他們,把他們殺了.王就大怒,發兵除滅那些兇手,燒燬他們的城.」

其 他人看不出這人沒有穿禮服,因為人是看外猊,上主看內心,這人是內外不一致的.這比喻使我想起奧斯卡．王爾德 (Oscar Wilde,1854-1900),著名的愛爾蘭作家,詩人,戲劇家,英國唯美主義運動的倡導者的小說《格雷的畫像》(The Picture of Dorian Gray).這小說諷刺英國維多利亞時代的道德虛偽,因為經多年清教徒的共和統治,社會很多清規戒律,到君主制度復僻維多利亞女皇當政時,就有一種反動, 特別縱情追求肉體的愉悅,但表面仍然是尊重道德清規.

小說內容提及的主角格雷(Dorian) 是一位俊朗秀美的少年人,畫家裴素(Basil) 被這位滿有迷人魅力的王子格雷所迷倒,為他畫了一幅肖像,因為這肖像,自已藝術生命重新激發起來.但格雷被裴素一位朋友亨利 (Henry) 勳爵的享樂主義影響,認為人生的目的是對酒當歌,行樂及時,最大的滿足是官感的暢快和欣愉.格雷喜愛自己的肖像,留住最絢爛的美態,但又認識到他也會美人 遲暮,人會變老,我可否樣子不變呢?他自己許了願,希望給人看到的,都是美好迷人的樣子,但他知道相由心生,歪念惡行也會影響俊朗的外表,但可否這些轉移 到肖像身上,由肖像承擔呢?他說:「希望我的外形能保持青春常駐,若有衰老,邪惡由你替我繪畫的肖像承受.每逢做了一件惡事,我仍可保持無邪,迷人的笑 容.」

有一次他拋棄了曾火熱追求的女演員西比爾.雲妮( Sibyl Vane).這位西比爾因為真正格雷的愛,竟然演不好莎士比亞的羅蜜歐與朱麗葉,而被格雷拋棄:「我愛你,只因為你演戲多麼好,你不能演出了,我為何仍愛 你呢?」結果西比爾因為被拋棄而自殺了.格雷後來回到上了鎖的閣樓,看看曾令他心醉的畫像,竟然發現畫像中的他,嘴角有陰陰的奸笑,赫然驚怕,真的每逢干 了一宗壞事,不單良心不安,而且鎖在閣樓沒有人知道的畫像,都會顯露出來.不錯,我們行經的路,都會留下痕跡.但另一方面,他有恃無恐,因為靈魂的傷害不 會形於色,何懼有之?跟著十多年,格雷生活變得放蕩,荒誕不羈,惡事做盡,以求一己肉慾的愉悅,而自己的容猊永遠秀美青春迷人.然而,他的靈魂,像關在閣 樓的俊美少年,每天愈來愈衰老,邪惡和難看.有一天他對畫他肖像的格雷說:「為了這像--我漂亮的像,我可以重新來過嗎?我愛的的肖像,我要好好做人.」 不過,肖像還是愈來愈差.為什麼?因為格雷並不是猛然醒悟悔改,乃是為了保持肖像的完美,他仍是追求虛榮和自戀,不是悔改.最後,他竟然用刀剌死為他繪像 的畫家裴素,並用同一張刀,刺破他不能接受--又老,又醜,又邪惡的肖像.在這時候,可怕的結局終於來到了.格雷突然變得瘋狂,在閣樓傳來恐布的叫聲,令 人不寒而憟,僕人急急召來警察,門開了,發現地上躺了一位又老又醜的死屍,若不是手上的戒指,根本認不出是格雷,而畫架上,出現一位秀美迷人的少年肖像.

我 們每一個人都要赴王兒子的筳席,我們已被邀請了,我們走到生命的終結,要面對我們的上帝.祂不像人,祇見到外猊,不少人被格雷俊朗的外猊騙了,我們可騙過 守門的人,甚至混進賓客當中,像他們一樣.但上主是看格雷那個鎖在閣樓,一般人不能看到的靈魂深處,他看出我們並沒有穿上他為我們預備了的禮服.

二,婚禮服──信心及衍生出來的行為

禮服是什麼呢?

歷代以來,對禮服有不同的看法,總的來說:一種認為是義的行為.馬太所強調的,在末日審判時,我們以具體的善行受審.馬太 25 章是一貫的看法－我餓了,你給我吃；我坐牢,你探訪我－這是可以進天國的.

在前面所說兩個比喻,兩種兒子,有行動的稅吏和妓女,他們到約但河領受施洗約翰悔改的洗禮,比口頭同意,但沒有行動的法利賽人,更早進入天國.

第二個比喻提及兩種園戶,結果子和不結果子,祇有結果子的才配得佔用上帝的葡萄園.加爾文就指出,禮服是聖潔的飾,詩 96:9:當以聖潔的妝飾,敬拜耶和華.

另一方面,馬丁路德說禮服是信心－人非有信,就不能到上帝面前.保羅說,我們因信而稱義,在基督裡的身份(in Christ)我們被上帝接納了.

但 信心和由信心衍生出來的行為其實是一個銅錢的兩面.加爾文強調在「我裡面的基督」(Christ in me) , 馬丁路德強調「為我的基督」(Christ for me),兩者有分別但不可分離.稱義是地位上,成聖是稱義的延伸和期望,信心沒有行為是死的,信心在愛心中才能保持不墮.

我想用婚禮的戒指誓言說明,作為例子.

信 心與善行,稱義和成聖就像盟誓中的一句話的前半段和後半段.誓言後,牧師問:有沒有信物表徵你們的盟誓呢?新人拿出戒指和對方說:「我以此戒指為徽記,與 你結婚,我以我身心尊敬你,我以我的財物,與你分享.」每逢念到下半句,有些新人,特別是女方都微笑,彷彿說:記住這樣,這是公開承認.有時台下的賓客也 笑起來.有一次新郎的工作是印鈔票的,女的笑容更燦爛.因著進入結合,夫妻有一種歡愉的對換,甘心樂意的,就像基督與新婦教會的聯合:我們的罪歸在祂身 上,祂的義歸在我們的身上；在地位上,我們立刻擁有對方的一切.

但前半句:我以我的身心尊敬你.不是立刻就能夠實現,是身份改變後的承 諾 和實踐.我們要逃避試探,不站危樓之下,才能以身心尊敬我們的配偶.我們心中惦記配偶,才能做到這承諾:我以我的身心尊敬你是漫長的路,經歷很多的回轉, 認罪,饒恕,立志,是一生的追求,是加爾文強調的成聖過程.但這漫長的旅程的起始,持續不斷,和完成,是從聯合的盟誓就開始了,所以聖經提到夫婦關係時 說:「這是極大的奧秘,但我是指著基督和教會所說的.」(弗五32 )戒指的誓言,正好說明稱義和成聖的關係.我的財物與你分享,是立刻生效的稱義,擁有基督的一切；但我以我的心尊敬你,是一生的態度和承諾.

再說,基督教婚姻的一夫一妻的「一」,是優先次序的一(ordinal one),不單是數字的一(numerical one) ；我們信仰獨一的上主,就要盡心,盡意,盡性愛上主.優先的一, 是不斷的追求.

三,邀請──被召和選上,信仰的考驗

我 們看過婚宴是審判的時刻,而禮服是與信心相應的行為,我們會否感到沮喪呢?因為我們恐怕不能赴上主的筵席,被上主拒諸門外.但感謝上主,比喻的基調,不是 灰色的,這種對得救的焦慮也是一種確據:我們對信仰非常認真,正如君士坦丁皇帝一路遲遲不肯洗禮,直到在臨終,因他對信仰對領洗非常認真；認為洗禮後生活 要聖潔,但他沒有絕對把握,所以臨終時才洗禮,俾靈魂清潔進入天堂.

但我們再看聖經,比喻的基調是歡樂的,婚禮的筵席是喜樂的,不是恐懼戰競,誠惶誠恐,患得患失,像到法庭接受審判的感覺:是被定罪?或是當庭釋放呢?我們要注意,主人的邀請都是恩典:

「所以你們要往岔路口上去,凡遇見的,都召來赴席. 」(馬太 22:9)他們都答應了,但到了關鍵時刻要行動時:

「牛和肥畜已經宰了,各樣都齊備．請你們來赴席」, 就諸多理由推搪,爽約.我們在大學都有這種經驗,邀請函來了,很想聽一場演講,參加一個開幕禮,回覆了,會出席.但到了那時刻,發現很多事情放不下,都不 去了.其實,想深一層,那些事可以暫時放下的,但就是不看重答應了的約會.現在,變得謹慎,不隨便答應,答應了,就要盡力出席,是一個承諾,也是紀律.

信仰需要有考驗的時刻,撒種比喻中有些種子落在淺土和石頭上,沒有根,一到困難,就放棄信仰的,但起初是歡歡喜喜接受的.

我 們不知道何時是我們的考驗時刻,我們要基督或要瑪門,我們要選擇誰是我們效忠的對象.也許考驗時刻就是我們患了頑疾,或感情受到重挫,或失去工作,或投資 失敗的時候.沒有穿禮服的原因,我想大部分不會像格雷一樣:為了放縱肉體的官能享受而與魔鬼交易,讓人看不見自已的美麗純潔肖像－我們的靈魂,受罪行帶來 的傷害變得醜陋.但我們卻很容易像比喻中所說:我沒有空,我要到田里工作,我要作買賣.當工作或職業成為營營役役(occupation 成為 pre-occupation),我們就進入信仰的危機了.當我們主日為了工作而沒有時間或心情就近上帝,我們的禮服就漸漸褪色了,我們閣樓的畫像開始產 生微妙的變化了,因為 occupation 之前應該有 pre-occupation,兩者不能混而為一；我們記掛的(pre-occupation) ,先於我們的工作(occupation) .這在比喻中,我們看到輕視上帝的筵席,就不配赴上帝的筵席.

如果我們相信上帝是獨一的真神,我們就要盡心,盡性,盡意愛主我們的上帝,我們不要讓 occupation 成為 pre-occupation,我們不棄絕信仰,但會拖延,推卻信仰的行動.

然而,這婚宴的比喻是喜樂的,因著與基督的聯合:我們有歡愉的對調,讓我們對基督說:我以我的身心尊敬禰.



\newpage

\section{驢以載道}
\label{sec:144}
\textbf{林豪恩先生 - 驢以載道}
\newline
\newline
\hyperref[sec:143]{< < < PREV SERMON < < <}
~
\hyperref[sec:toc]{[返目錄]}
~
\hyperref[sec:145]{> > > NEXT SERMON > > >}


馬可福音 11:1-11
\newline
\begin{longtable}{cl}
\hline
\hline
章節 & 經文 (和合本修訂版)\\
\hline
11:1 & \begin{tabularx}{0.7\textwidth}{X} 耶穌和門徒快到耶路撒冷，來到伯法其和伯大尼，在橄欖山那裡。耶穌打發兩個門徒， \end{tabularx} \\ \\ \relax
11:2 & \begin{tabularx}{0.7\textwidth}{X} 對他們說：「你們往對面村子裡去，一進去的時候會看見一匹驢駒拴在那裡，是從來沒有人騎過的，把牠解開，牽來。 \end{tabularx} \\ \\ \relax
11:3 & \begin{tabularx}{0.7\textwidth}{X} 若有人對你們說：『為甚麼做這事？』你們就說：『主要用牠，但會立刻把牠牽回到這裡來。』」 \end{tabularx} \\ \\ \relax
11:4 & \begin{tabularx}{0.7\textwidth}{X} 他們去了，看見一匹驢駒拴在門外街道上，就把牠解開。 \end{tabularx} \\ \\ \relax
11:5 & \begin{tabularx}{0.7\textwidth}{X} 在那裡站著的人，有幾個說：「你們解開驢駒做甚麼？」 \end{tabularx} \\ \\ \relax
11:6 & \begin{tabularx}{0.7\textwidth}{X} 門徒照著耶穌的話說，那些人就任憑他們牽去了。 \end{tabularx} \\ \\ \relax
11:7 & \begin{tabularx}{0.7\textwidth}{X} 他們把驢駒牽到耶穌那裡，把自己的衣服搭在上面，耶穌就騎上。 \end{tabularx} \\ \\ \relax
11:8 & \begin{tabularx}{0.7\textwidth}{X} 有許多人把衣服鋪在路上，還有人把田間的樹枝砍下來鋪上。 \end{tabularx} \\ \\ \relax
11:9 & \begin{tabularx}{0.7\textwidth}{X} 前呼後擁的人都喊著說：「和散那！奉主名來的是應當稱頌的！ \end{tabularx} \\ \\ \relax
11:10 & \begin{tabularx}{0.7\textwidth}{X} 那將要來的我祖大衛之國是應當稱頌的！至高無上的，和散那！」 \end{tabularx} \\ \\ \relax
11:11 & \begin{tabularx}{0.7\textwidth}{X} 耶穌到了耶路撒冷，進入聖殿，看了周圍的一切。天色已晚，他就和十二使徒出城，往伯大尼去。 \end{tabularx} \\ \\
[1ex]
\hline
\hline
\end{longtable}
日期:2012年4月29日

題目:驢以載道

經文:馬可福音11章1-11節

證道:林豪恩先生

一．引言

一位擁有10多萬下屬,管理3000多億財政預算的上司,出外旅遊的時候應該使用甚麼交通工具呢?應該住在甚麼地方呢?如果這位上司那10多萬下屬中最基層的一位都是乘搭民航機去日本,或者乘搭客輪往澳門,作為10多萬員工的上司,不乘搭私人飛機和私人遊艇,怎能顯出他身份的尊貴和地位的崇高呢?如果只入住普通客房,怎能顯出他和一般人的分別呢?

讓我們把鏡頭轉到那位將近耶路撒冷的耶穌,他被前行後隨的人呼喊:「奉主名來的以色列王是應當稱頌的!」(約12:13)；又在腓立比書被宣稱「本有上帝的形象」,「與上帝同等」,以他的身份和地位,進入本屬於他的聖城耶路撒冷的時候,又應該乘搭甚麼交通工具呢?他選擇甚麼交通工具,又有甚麼意義呢?

二．經文

11 : 1耶穌和門徒將近耶路撒冷,到了伯法其和伯大尼,在橄欖山那裡；耶穌就打發兩個門徒, 11 : 2對他們說:「你們往對面村子裡去,一進去的時候,必看見一匹驢駒拴在那裡,是從來沒有人騎過的,可以解開,牽來. 11 : 3若有人對你們說:『為甚麼做這事?』你們就說:『主要用牠. 』那人必立時讓你們牽來.」 11 : 4他們去了,便看見一匹驢駒拴在門外街道上,就把牠解開. 11 : 5在那裡站著的人,有幾個說:「你們解驢駒做甚麼?」 11 : 6門徒照著耶穌所說的回答,那些人就任憑他們牽去了. 11 : 7他們把驢駒牽到耶穌那裡,把自己的衣服搭在上面,耶穌就騎上. 11 : 8有許多人把衣服鋪在路上,也有人把田間的樹枝砍下來,鋪在路上. 11 : 9前行後隨的人都喊著說:和散那!奉主名來的是應當稱頌的! 11 : 10那將要來的我祖大衛之國是應當稱頌的!高高在上和散那! 11 : 11耶穌進了耶路撒冷,入了聖殿,周圍看了各樣物件.天色已晚,就和十二個門徒出城,往伯大尼去了.

三．釋經與應用

今天是棕枝主日,記念耶穌進入耶路撒冷.不過,馬可福音記述耶穌進入耶路撒冷事件,主要篇幅用在驢駒身上,而不在棕樹枝身上,那隻驢駒的形象比棕樹枝的形象更要突出.棕樹枝,是人們用來歡迎耶穌的,驢駒,是耶穌自己選擇騎上的.耶穌為甚麼要在這個時候找一隻驢駒去騎呢?會不會是耶穌體力不支而需要坐騎呢?還是耶穌趕時間而需要坐騎呢?

當時除了耶穌和門徒之外,其實往耶路撒冷的路上行人絡繹不絕,因為猶太人的重要節期快到,人們都要去耶路撒冷過節.一般平民百姓大多數都如耶穌和門徒一樣,步行前往耶路撒冷.現在,耶穌和門徒已經來到位於橄欖山的伯法其,這伯法其在耶路撒冷的東面,距離耶路撒冷只有1.6公里.其實只要多行一會就到達耶路撒冷了.路加福音記載,耶穌十二歲時,已經跟著父母上耶路撒冷過節,而且自從公開傳道開始,踏遍大城小鄉,耶穌從不言倦,現在怎會連走這1.6公里的路也力有不逮呢?另一方面,從上下文來看,耶穌並沒有要趕時間入耶路撒冷的狀況,如果真的要趕時間,就應該找走得快一點的交通工具,如果要找驢,也應該找一隻成年的驢,而不是一隻驢駒.

如果不是現實上的需要,耶穌為甚麼要找一隻驢駒來騎著走這段進入耶路撒冷的路呢?耶穌這行動令人們聯想到古時的先知,先知有時會做出一些特別的行動,行動的目的是要表達信息.例如,以西結曾經用刀剃光自己的頭髮和鬍鬚,然後用天平把鬚髮平分為三份,把一份在城中用火焚燒,把一份在城的四圍用刀砍碎,把一份任風吹散(以西結書5:1-4).以西結削髮剃鬚,並非他個人的需要,而是要藉此表達以色列被審判的信息.耶穌騎驢進入耶路撒冷,這行動本身也是一項信息.他要藉著這項行動表達甚麼信息呢?馬太福音指出,「這事成就是要應驗先知的話」(12:4).這位先知是誰呢?所說的是甚麼話呢?這位先知就是撒迦利亞,在撒迦利亞書中如此預告說:「錫安的民哪,應當大大喜樂；耶路撒冷的民哪,應當歡呼.看哪,你的王來到你這裡!他是公義的,並且施行拯救,謙謙和和地騎著驢,就是騎著驢的駒子.」(9:9)

耶穌騎驢進入耶路撒冷的信息是:「我就是先知撒迦利亞所預言要來的王.」即是說,耶穌的行動是要展示身份.令人困惑的地方是,如果耶穌要展示王者的身份,就不應該選擇一隻驢來騎上,因為對於當時的人們來說,驢是步行以外一種最平民化和普通的交通工具,就正如現在我們去澳門乘搭飛翔船,去日本乘搭民船機一樣.如果要展示王者尊貴的身份,最低限度都要騎馬,這樣,才能把自己從平民百姓中分別出來.

有趣的地方是,耶穌的門徒和有些去耶路撒冷過節的路人,竟然配合耶穌的行動,參與在這場獨幕劇中,把耶穌當作王一樣來歡迎.

11 : 7他們把驢駒牽到耶穌那裡,把自己的衣服搭在上面,耶穌就騎上. 11 : 8有許多人把衣服鋪在路上,也有人把田間的樹枝砍下來,鋪在路上. 11 : 9前行後隨的人都喊著說:和散那!奉主名來的是應當稱頌的! 11 : 10那將要來的我祖大衛之國是應當稱頌的!高高在上和散那!

不過,約翰福音指出:「這些事門徒起先不明白,等到耶穌得了榮耀以後才想起這些話是指著他寫的,並且眾人果然向他這樣行了.」(約12:16)路加福音也藉著記述同行中的幾個法利賽人的反應,指出這行動的吊詭性.當耶穌騎在驢駒上,門徒在大聲呼喊的時候,途人中的幾個法利賽人對耶穌說:「夫子,責備你的門徒吧!」(路加福音19:39)即是說,有些途人把耶穌和門徒騎驢進城的行動當作鬧劇來看.

如果耶穌有公關團隊,就不會容許他騎驢進城,而很可能會為他預備駿馬,或者戰車,因為眾人都認為,驢是平民百姓使用的,不能突顯王者的身份.能夠乘載君王的,應該是駿馬,或者戰車.耶穌為甚麼偏偏不選用駿馬或戰車呢?

馬太福音指出,耶穌騎驢進入耶路撒冷是應驗先知撒迦利亞的預言.為甚麼不是駿馬或戰車,原來也在撒迦利亞書中找到提示.

9 : 10 我必除滅以法蓮的戰車 和耶路撒冷的戰馬； 爭戰的弓也必除滅. 他必向列國講和平； 他的權柄必從這海管到那海, 從大河管到地極.

這位應許要來的王不騎戰馬進入耶路撒冷,因為他來是要除滅戰馬；這位拯救者不站在戰車上進入耶路撒冷,因為他來正正是要除滅戰車.他的確是王,途人沒有說錯,他是應當稱頌的,因為他有權柄「必從這海管到那海,從大河管到地極」,但先知指出,他的素質是「謙謙和和」,他的原則是「公義」,他的行動是「拯救」,他的方法是「和平」.正因如此,平民百姓日常使用的交通工具──驢,比侯王將相,達官貴人使用的駿馬和戰車,更適合用來乘載這成為肉身的道.

11 : 11耶穌進了耶路撒冷,入了聖殿,周圍看了各樣物件.天色已晚,就和十二個門徒出城,往伯大尼去了.

在伯法其往耶路撒冷路上這獨幕劇演完了,一切都回復平靜.奇怪的地方是,耶穌不是要進入耶路撒冷稱王的嗎?為甚麼在夜幕低垂的時候帶同門徒離開耶路撒冷往伯大尼去呢?伯大尼只是位於橄欖山東側的的一條小村莊,而且距離耶路撒冷有3公里,每天來回耶路撒冷也不太方便呢!耶穌竟然選擇在伯大尼的民居留宿,而不入住耶路撒冷的「尊貴套房」.如果他入住耶路撒冷的「尊貴套房」,又選擇與耶路撒冷的利益系統交朋友,他明天如何有正義能量和道德位置進入聖殿挑戰祭司系統和商業系統的勾結呢?

耶穌認為,用來乘載道,一般百姓可以用到的驢比侯王將相用的駿馬和戰車更適合；用來接待道,小村莊的民居比耶路撒冷的府第更適合.

在地上代表耶穌的教會,又用甚麼來乘載道呢?教會為道去選擇「交通工具」的時候,考慮的是甚麼呢?藉此向世界宣告甚麼信息呢?這是值得教會思考的事情.

在1930-1945年任北京大學校長的蔣夢麟對於基督教當時選擇用來載道的媒介有這樣的評論:「中國人與基督教或任何其他宗教一向沒有甚麼糾紛,不過到了十九世紀中葉,基督教與以兵艦做靠山的商業行為結了夥,因而在中國人的心目中,這個宣傳愛人如己的宗教也就成為侵略者的工具了.人們發現一種宗教與武力形影不離時,對這種宗教的印象自然就不同了.而且中國人也實在無法不把基督教和武力脅逼相提並論.慢慢地人們產生了一種印象,認為如來佛是騎著白象到中國的,耶穌基督卻是騎在炮彈上飛過來的.」《西潮》

蔣夢麟校長所描述的圖畫──耶穌基督騎在炮彈上飛過來,希望已經成為過去.今天,教會用甚麼來盛載道呢?

有一位信徒在2007年去美國參觀水晶宮大教堂,參觀完後,她有這樣的感想:「由於是晚上到達這所座堂,未能一睹這所教堂的宏偉,然而這所座堂的洗手間卻擁有六星級的裝潢及設備,相信這是對神的殿予以應有及最崇高的敬意,亦是此行感覺最特別之處.回程時適逢遇到Disney Land放煙花,更因此可以享受免費的餘興節目.」有六星級洗手間的六星級大教堂,向這位慕名而來的信徒發放了關於信仰的甚麼信息呢?顯赫,輝煌,豪華,尊貴,享受?如果是這樣,和香港的樓盤廣告所宣傳的信息非常接近.與此同時,這六星級的大教堂又向世界宣告了甚麼關於基督教的信息呢?有報章報導這大教堂時這樣說:「馳名世界的水晶宮大教堂自2009年以來就遭遇財務危機.據悉,水晶宮大教堂拖欠小企業及員工750萬美元,包括駱駝和其他動物租用費,幾家電視廣播公司和設備租賃公司也起訴教堂欠款不還,估計總債務高達5千萬美元.」

或者,教會在表達自己的時候,應該多一些參考耶穌表達自己的方式,而減少追逐市場價值和市場標準.

教會面對選擇,信徒也面對選擇,其實,耶穌也面對選擇.騎驢?騎馬?還是坐車?住耶路撒冷府第,還是伯大尼民居?選擇的背後,反映著選擇者的自我定位及存在的追求.耶穌選擇騎驢,即是選擇平民,耶穌選擇伯大尼的民居,即是選擇百姓.耶穌把自己視為眾生之一,分享著一般人的命運與掙扎.我們的社會卻總是鼓動我們不要做一般人,引誘我們擺脫一般人的命運與掙扎.不久之前,有一位高官勸籲青年自我反省:「為何我做不到下一個李嘉誠?」如果所有人都那麼有「大志」,要做下一個李嘉誠,那麼,所有人的目光和努力都只是關注如何能夠做下一個李嘉誠,而現實上云云眾生中絕大部份人做不到李嘉誠,那麼,還有多少目光和努力會關注平民百姓呢?作為平民百姓還有甚麼意義呢?平民百姓的議程有沒有被認真對待呢?

最近有報導說,外國人不願意來香港工作的其中一個原因是不能夠為他們的子女在香港找到國際學校的學位.為甚麼國際學校的學位如此缺乏呢?原來,香國際學校有一半學生都是土生土長的本地孩子,其中不少是政府官員的子女.當大家的努力及行動都是把自己的孩子從本地普及教育系統中「拯救出來」,在這樣的情況下,有人說,「你的孩子在那裡,你的心就在那裡」,又有人說,「屁股指揮腦袋」,主事者改善這個本地普及教育系統的著緊程度和切身感會否不同呢?

教會參與在這場遊戲之中,如何為自己定位呢?如何選擇呢?這些選擇反映教會在追逐甚麼呢?耶穌選擇了平民用的驢,選擇了百姓的民居,選擇一般人的命運,耶穌基督的教會選擇甚麼呢?跟隨耶穌基督的門徒選擇甚麼呢?



\newpage

\section{撻著生命的動力}
\label{sec:145}
\textbf{陳榮超牧師 - 撻著生命的動力}
\newline
\newline
\hyperref[sec:144]{< < < PREV SERMON < < <}
~
\hyperref[sec:toc]{[返目錄]}
~
\hyperref[sec:146]{> > > NEXT SERMON > > >}


路加福音 24:13-32
\newline
\begin{longtable}{cl}
\hline
\hline
章節 & 經文 (和合本修訂版)\\
\hline
24:13 & \begin{tabularx}{0.7\textwidth}{X} 同一天，門徒中有兩個人往一個村子去；這村子名叫以馬忤斯，離耶路撒冷約有二十五里。 \end{tabularx} \\ \\ \relax
24:14 & \begin{tabularx}{0.7\textwidth}{X} 他們彼此談論所發生的這一切事。 \end{tabularx} \\ \\ \relax
24:15 & \begin{tabularx}{0.7\textwidth}{X} 正交談議論的時候，耶穌親自走近他們，和他們同行， \end{tabularx} \\ \\ \relax
24:16 & \begin{tabularx}{0.7\textwidth}{X} 可是他們的眼睛模糊了，沒認出他。 \end{tabularx} \\ \\ \relax
24:17 & \begin{tabularx}{0.7\textwidth}{X} 耶穌對他們說：「你們一邊走一邊談，彼此談論的是甚麼事呢？」他們就站住，臉上帶著愁容。 \end{tabularx} \\ \\ \relax
24:18 & \begin{tabularx}{0.7\textwidth}{X} 兩人中有一個名叫革流巴的回答：「你是在耶路撒冷的旅客中，惟一還不知道這幾天在那裡發生了甚麼事的人嗎？」 \end{tabularx} \\ \\ \relax
24:19 & \begin{tabularx}{0.7\textwidth}{X} 耶穌對他們說：「甚麼事呢？」他們對他說：「就是拿撒勒人耶穌的事。他是個先知，在神和眾百姓面前，說話行事都大有能力。 \end{tabularx} \\ \\ \relax
24:20 & \begin{tabularx}{0.7\textwidth}{X} 祭司長們和我們的官長竟把他解去，定了死罪，釘在十字架上。 \end{tabularx} \\ \\ \relax
24:21 & \begin{tabularx}{0.7\textwidth}{X} 但我們素來所盼望要救贖以色列民的就是他。不但如此，這些事發生到現在已經三天了。 \end{tabularx} \\ \\ \relax
24:22 & \begin{tabularx}{0.7\textwidth}{X} 還有，我們中間的幾個婦女使我們驚奇：她們清早去了墳墓， \end{tabularx} \\ \\ \relax
24:23 & \begin{tabularx}{0.7\textwidth}{X} 不見他的身體，就回來告訴我們，說她們看見了天使顯現，說他活了。 \end{tabularx} \\ \\ \relax
24:24 & \begin{tabularx}{0.7\textwidth}{X} 又有我們的幾個人往墳墓那裡去，所發現的正如婦女們所說的，只是沒有看見他。」 \end{tabularx} \\ \\ \relax
24:25 & \begin{tabularx}{0.7\textwidth}{X} 耶穌對他們說：「無知的人哪，先知所說的一切話，你們的心信得太遲鈍了。 \end{tabularx} \\ \\ \relax
24:26 & \begin{tabularx}{0.7\textwidth}{X} 基督不是必須受這些苦難，然後進入他的榮耀嗎？」 \end{tabularx} \\ \\ \relax
24:27 & \begin{tabularx}{0.7\textwidth}{X} 於是，他從摩西和眾先知起，凡經上所指著自己的話都給他們作了解釋。 \end{tabularx} \\ \\ \relax
24:28 & \begin{tabularx}{0.7\textwidth}{X} 他們走近所要去的村子，耶穌好像還要往前走， \end{tabularx} \\ \\ \relax
24:29 & \begin{tabularx}{0.7\textwidth}{X} 他們卻強留他說：「時候晚了，天快黑了，請你同我們住下吧。」耶穌就進去，要同他們住下。 \end{tabularx} \\ \\ \relax
24:30 & \begin{tabularx}{0.7\textwidth}{X} 坐下來和他們用餐的時候，耶穌拿起餅來，祝福了，擘開，遞給他們。 \end{tabularx} \\ \\ \relax
24:31 & \begin{tabularx}{0.7\textwidth}{X} 他們的眼睛開了，這才認出他來。耶穌卻從他們眼前消失了。 \end{tabularx} \\ \\ \relax
24:32 & \begin{tabularx}{0.7\textwidth}{X} 他們彼此說：「在路上他和我們說話，給我們講解聖經的時候，我們的心在我們裡面豈不是火熱的嗎？」 \end{tabularx} \\ \\
[1ex]
\hline
\hline
\end{longtable}
日期:2012年5月6日

題目:撻著生命的動力

經文:路加福音24章13-32

證道:陳榮超牧師

I. 引言 感謝主.在我繼續之前,讓我們做一個活動.把拇指中指snap .因為今天的資訊是一個動的資訊,我希望能夠把撻著大家的心,撻著基督復活的能力.我們這個手指活動就是象徵著一個撻著的復活生命.

好像我們要開動一部車,我們先要把車撻著.但是有些時候,就撻不著.最普遍的原因有兩種.沒有電油或者電池幹了.沒電油,需要加油.電池幹了,就可以跳電,來撻著部車.

所以今天早上,我想問問各位信徒,各位朋友,你的生命是否已經撻著了?或跟進一步問,有否被基督復活大能撻著?因為聖經哥林多前書,使徒保羅這麼講:

15:13 若沒有死人復活的事,基督也就沒有復活了.14若基督沒有復活,我們所傳的便是枉然,你們所信的也是枉然.所以,基督的復活就是撻著我們生活動力的能力.

我選了路加福音24章兩位門徒的經歷,因為他們有過撻著了基督復活能力的見證.

II. 迷糊-盲點:固執,偏見,傲慢

A. 以馬忤斯路上的迷糊 今天的經文一開始講到兩位門徒離開了耶路撒冷﹐向一條村子叫以馬忤斯的方向走.這村子離開耶路撒冷大約7英里路(~10公里).他們的情緒不大好﹐臉上帶著愁容.好像一部撻不著的車.

這兩位門徒的心境是很失望﹐他們的情緒是萎縮﹐他們的精神是很憂抑.這兩位是曾經跟隨耶穌的門徒.他們一直跟耶穌到了耶路撒冷.但是耶路撒冷帶給他們的是一片失望.

24:18二人中有一個名叫革流巴的,說...這幾天在[耶路撒冷出了事].就是拿撒勒人耶穌的事．他是個先知,在 神和眾百姓面前,說話行事都有大能.24:20祭司長和我們的官府,竟把他解去定了死罪,釘在十字架上.24:21但我們素來所盼望要贖以色列民的就是他.

B. 看而看不見的迷糊 經文記載了更有趣的事情.

15正談論相問的時候,耶穌親自就近他們,和他們同行；16只是他們的眼睛迷糊了,不認識他.

當我念到這句經文的時候,我第一個感覺就是"失之交臂"啊.如果這次他們與耶穌擦肩而過,不是人生一種悲劇嗎?什麼問題令到他們和耶穌差點擦肩而過?

16節:只是他們的眼睛迷糊了.

這位相處了三年的老師,竟然出現在眼前都認不得.迷糊了?不是因為眼膜發炎,青光眼,白內障.也不是沙塵風暴,大霧遮蓋.迷糊的意思是看卻看不見.不過你不要以為奇怪.其實這類的現象很普遍.也常常發生在我們日常生活上.

關鍵就是這兩位門徒自己所講的話:我們素來所盼望要贖以色列的就是他.而關鍵的兩個字就是:素來.因為"素來"這個意念往往令我們在人生路上與我們有意義的事情失之交臂.

C. ``素來"盲點 聽聽他們對耶穌"素來"的看法.他是個先知,在 神和眾百姓面前,說話行事都有大能.我們素來所盼望要贖以色列民的就是他.不過,現在一切都成為過去了.

24:20祭司長和我們的官府,竟把他解去定了死罪,釘在十字架上.素來認為的救贖主,竟然死在十字架上.他們是"素來"成為了他們對眼前耶穌的盲點.

D. 盲點的特徵 我們沒人都有盲點.形成盲點的特徵大約有三種,固執,偏見和傲慢.有人說,男人開車都是盲點的在開.特別迷路的時候,盲點特別大.坐在旁邊的太太發覺老公好像在兜圈,車好像在迷宮走不出來的一樣.提議不如停下來問路.大男人多半不停的.今天在座的男人,你會聽老婆話馬上停去問路的請舉手?不會的,因為大男人就是固執.

偏見會做成盲點.我跟朋友談福音.他說基督教信仰上,最不能接受的就是耶穌復活.我說為什麼不能接受?他說,那除非你能夠叫耶穌馬上再復活一次.這位朋友一定是科學腦袋的人.如果復活是真的話,一定可以重複測試.因為這是他的偏見.但也是成為了他的盲點.

傲慢更是做成盲點的因素.一位朋友和我說,我不信就不信.就算耶穌現在在我面前,我都不要.因為我的朋友都在地獄.我幹嘛去天堂.我也沒話可以說了.

E. 小結 這兩位門徒的眼目迷糊了.是因為心中的盲點.可能不會傲慢,但是固執和偏見妨礙了他們對這位和他們同行,已經復活了的耶穌看得清楚.

你們呢?你對認識耶穌基督有盲點嗎?

III. 遲鈍- A. 如鈍刀 耶穌一邊陪他們走,聽他們講在耶路撒冷的灰心.同時,耶穌慢慢把他們所講的事情,用聖經串聯起來.耶穌說:

26基督這樣受害,又進入他的榮耀,豈不是應當的麼?27於是從摩西和眾先知起,凡經上所指著自己的話都給他們講解明白了.

耶穌其實沒有講新的東西.都是講他們都認識的聖經.只是他們心遲鈍.

25耶穌對他們說:無知的人哪,先知所說的一切話,你們的心信得太遲鈍了.

遲鈍就如沒有鋒的刀.工若善其用,必先利其器.一把鈍的刀,只是一塊爛鐵.連洋蔥都切不進去.

B. 悟性鈍- 失去推理的本能 耶穌說他們遲鈍,是他們的腦袋鈍了.腦袋生鏽了.聽聽這兩位自己和自己的解釋.

24:22再者,我們中間有幾個婦女使我們驚奇,他們清早到了墳墓那裡．24:23不見他的身體,就回來告訴我們說,看見了天使顯現,說他活了.24:24又有我們的幾個人,往墳墓那裡去,所遇見的,正如婦女們所說的,只是沒有看見他.

他們講完都好像冇講一樣.事實都擺在眼前,就是沒有悟性.

C. 靈性鈍- 失去聖靈的默示 或者是靈性鈍.我記得一個見證.

范博士的見證.當我讀本科的時候,一位博士後的學者來參加我們校園的團契.他是物理化學的博士研究生.他參加我們星期五晚上舉行的基督徒團契.不過他表面來的態度是希望勸我們這群基督徒不要誤信宗教.因為是沒有神創造這回事.他說他已經可以在他的實驗室做了試管細胞.

他是非常聰明的學者.博士後研究生.我們這群本科生那有能力來反駁他.大家都很納悶.不過來了一次,都沒來了.我們也松一口氣.

學期終,他出現了.我們都緊張起來.他站起來,表面他今天是基督徒了.我們都很驚訝.又不知道怎樣回應.他作見證. 透過我們這團契一位老導師,和他禱告.祈求聖靈的能力來啟導他的思想.

這位博士不是悟性鈍,乃是靈性鈍.失去了聖靈的默示.

我盼望你今天不要作一個靈性遲鈍的人.免至你和耶穌失之交臂.

IV. 開門- 邀請耶穌坐席 A. 門徒的邀請

24:28將近他們所去的村子,耶穌好像還要往前行.24:29他們卻強留他說,時候晚了,日頭已經平西了,請你同我們住下罷.耶穌就進去,要同他們住下.

這是猶太人的好傳統.聖經教導他們要款待客人.雖然他們認為耶穌是外地人,陌生的人,禮貌上,他們應該邀請他過個晚上.給客人晚餐.

B. 耶穌的邀請

24:30到了坐席的時候,耶穌拿起餅來,祝謝了,擘開,遞給他們.24:31他們的眼睛明亮了,這纔認出他來．忽然耶穌不見了.24:32他們彼此說,在路上,他和我們說話,給我們講解聖經的時候,我們的心豈不是火熱的麼.

突然間,晚餐的時候,客人反客為主.在門徒眼前重映一個頗為熟悉的錄影.最後﹐他們就覺得眼睛明亮﹐而且﹐回想到在路上的時候﹐心不是已經覺得火熱嗎.

C. 要撻著生命,打開心門 我們的心不是已經覺得火熱嗎?這句話非常有意思.就是最遲鈍的心也曾經感受過"心中的火熱".只是忽略了靈性悟性的反應.我希望今年的復活節,我們眾人能夠把心中的火熱撻著.撻著生命的能力.

像這兩位門徒,當他們打開門,邀請耶穌入住的時候,也就是他們生命撻著的時候.所以要撻著生命,就要打開心門.

V. 結論: 兩位門徒的經歷,似曾相識,是嗎?他們對耶穌的遲鈍悟性遲鈍,靈性遲鈍,令他們 和耶穌差點失之交臂.我們過去也不是存著固執,偏見或者傲慢來看耶穌的復活.

今天,讓我邀請你把固執,偏見和傲慢放下來.讓聖靈光照你,讓你有一顆敏銳的心靈來接受耶穌是復活的主,復活的基督.這位復活的主就在我們當中.在你心門外敲門.

啟示錄:3:20節.20看哪,我站在門外叩門,若有聽見我聲音就開門的,我要進到他那裡去,我與他他與我一同坐席.

讓我們一同來撻著復活的生命.



\newpage

\section{你們為何愁煩}
\label{sec:146}
\textbf{伍渭文牧師 - 轉夕陽為清晨 -- 你們為何愁煩?}
\newline
\newline
\hyperref[sec:145]{< < < PREV SERMON < < <}
~
\hyperref[sec:toc]{[返目錄]}
~
\hyperref[sec:147]{> > > NEXT SERMON > > >}


路加福音 24:13-35
\newline
\begin{longtable}{cl}
\hline
\hline
章節 & 經文 (和合本修訂版)\\
\hline
24:13 & \begin{tabularx}{0.7\textwidth}{X} 同一天，門徒中有兩個人往一個村子去；這村子名叫以馬忤斯，離耶路撒冷約有二十五里。 \end{tabularx} \\ \\ \relax
24:14 & \begin{tabularx}{0.7\textwidth}{X} 他們彼此談論所發生的這一切事。 \end{tabularx} \\ \\ \relax
24:15 & \begin{tabularx}{0.7\textwidth}{X} 正交談議論的時候，耶穌親自走近他們，和他們同行， \end{tabularx} \\ \\ \relax
24:16 & \begin{tabularx}{0.7\textwidth}{X} 可是他們的眼睛模糊了，沒認出他。 \end{tabularx} \\ \\ \relax
24:17 & \begin{tabularx}{0.7\textwidth}{X} 耶穌對他們說：「你們一邊走一邊談，彼此談論的是甚麼事呢？」他們就站住，臉上帶著愁容。 \end{tabularx} \\ \\ \relax
24:18 & \begin{tabularx}{0.7\textwidth}{X} 兩人中有一個名叫革流巴的回答：「你是在耶路撒冷的旅客中，惟一還不知道這幾天在那裡發生了甚麼事的人嗎？」 \end{tabularx} \\ \\ \relax
24:19 & \begin{tabularx}{0.7\textwidth}{X} 耶穌對他們說：「甚麼事呢？」他們對他說：「就是拿撒勒人耶穌的事。他是個先知，在神和眾百姓面前，說話行事都大有能力。 \end{tabularx} \\ \\ \relax
24:20 & \begin{tabularx}{0.7\textwidth}{X} 祭司長們和我們的官長竟把他解去，定了死罪，釘在十字架上。 \end{tabularx} \\ \\ \relax
24:21 & \begin{tabularx}{0.7\textwidth}{X} 但我們素來所盼望要救贖以色列民的就是他。不但如此，這些事發生到現在已經三天了。 \end{tabularx} \\ \\ \relax
24:22 & \begin{tabularx}{0.7\textwidth}{X} 還有，我們中間的幾個婦女使我們驚奇：她們清早去了墳墓， \end{tabularx} \\ \\ \relax
24:23 & \begin{tabularx}{0.7\textwidth}{X} 不見他的身體，就回來告訴我們，說她們看見了天使顯現，說他活了。 \end{tabularx} \\ \\ \relax
24:24 & \begin{tabularx}{0.7\textwidth}{X} 又有我們的幾個人往墳墓那裡去，所發現的正如婦女們所說的，只是沒有看見他。」 \end{tabularx} \\ \\ \relax
24:25 & \begin{tabularx}{0.7\textwidth}{X} 耶穌對他們說：「無知的人哪，先知所說的一切話，你們的心信得太遲鈍了。 \end{tabularx} \\ \\ \relax
24:26 & \begin{tabularx}{0.7\textwidth}{X} 基督不是必須受這些苦難，然後進入他的榮耀嗎？」 \end{tabularx} \\ \\ \relax
24:27 & \begin{tabularx}{0.7\textwidth}{X} 於是，他從摩西和眾先知起，凡經上所指著自己的話都給他們作了解釋。 \end{tabularx} \\ \\ \relax
24:28 & \begin{tabularx}{0.7\textwidth}{X} 他們走近所要去的村子，耶穌好像還要往前走， \end{tabularx} \\ \\ \relax
24:29 & \begin{tabularx}{0.7\textwidth}{X} 他們卻強留他說：「時候晚了，天快黑了，請你同我們住下吧。」耶穌就進去，要同他們住下。 \end{tabularx} \\ \\ \relax
24:30 & \begin{tabularx}{0.7\textwidth}{X} 坐下來和他們用餐的時候，耶穌拿起餅來，祝福了，擘開，遞給他們。 \end{tabularx} \\ \\ \relax
24:31 & \begin{tabularx}{0.7\textwidth}{X} 他們的眼睛開了，這才認出他來。耶穌卻從他們眼前消失了。 \end{tabularx} \\ \\ \relax
24:32 & \begin{tabularx}{0.7\textwidth}{X} 他們彼此說：「在路上他和我們說話，給我們講解聖經的時候，我們的心在我們裡面豈不是火熱的嗎？」 \end{tabularx} \\ \\ \relax
24:33 & \begin{tabularx}{0.7\textwidth}{X} 於是他們立刻起身，回耶路撒冷去，看見十一個使徒和與他們正在一起的人聚集在一處， \end{tabularx} \\ \\ \relax
24:34 & \begin{tabularx}{0.7\textwidth}{X} 說：「主果然復活了，已經顯現給西門看了。」 \end{tabularx} \\ \\ \relax
24:35 & \begin{tabularx}{0.7\textwidth}{X} 於是，兩個人把路上所遇到，和耶穌擘餅的時候怎麼被他們認出來的事，都述說了一遍。 \end{tabularx} \\ \\
[1ex]
\hline
\hline
\end{longtable}
日期:2012年5月13日

題目:轉夕陽為清晨 -- 你們為何愁煩?

經文:路加福音 24:13-35

證道:伍渭文牧師

夕陽無限好,祇是近黃昏.不錯,夕陽是無限好的,如果你有時間,有心情欣賞的話.但如果你要趕十一公哩的路回家,而四周是曠野,隨時有野獸出現,心情憂傷,失望,夕陽便顯得催迫,山色暮掩,樹向晚風搖,四周顯得蒼茫,而你正慢慢步入黑暗中.這正是以馬忤斯路上兩個門徒心情的寫照,也是當人感到絕望,憂傷的寫照.就在這時候,耶穌在路上向兩位門徒顯現,並對他們講解聖經,給他們帶來盼望,後來他們曉得,這位伴隨他們的同路人,就是復活的主,他們非常興奮,急不及待,驅走疲倦,披星帶月,再走十一公哩,趕回耶路撒冷,對十一位門徒說:「主果然復活,已經現給西門看了.」(路二十四34)

1. 隱藏的基督陪同我們行走人生的窄路 (13-24節)

「他們彼此談論所遇見的這一切事.正談論的時候,耶穌親自就近他們,和他們同行；祇是他們的眼睛迷糊了,不認識他.」(15, 16節)

經文給我們看到,兩個門徒談論耶穌去世的事入了神,他們需要藉著談論彼此支持.耶穌親自就近他們,不動聲色,和他們同行,伴隨他們,聽他們談論一段時間後才問他們:「你們走路彼此談論的是甚麼事呢?」他們因為「眼睛迷糊」不認得是耶穌,新譯本作「眼睛卻給蒙蔽了」,這是上帝的工作,不是因過度傷心看不清是耶穌,因為十一公哩是不短的路程,也許最初因情緒憂傷看不清,但當耶穌開始解聖經時,他們應該愈來愈清醒:「於是從摩西和眾先知起,凡經上所指著自己的話都給他們講解明白了.」(27節) 他們聽的時候,心裡火熱,他們也明白,因明白而興奮起來,所以他們的眼睛不是因憂傷而看不清楚,是上帝的工作,上帝是隱藏的上帝,這位陌生人,根本不像他們見過的耶穌.

為何耶穌隱藏起來?不讓他們看見復活的真身呢?

因為耶穌成為同路人,可聆聽他們心中的苦痛.我們看聖經,耶穌不是突然出現,攔住去路來問他們,乃和他們行了一段路後才問他們.「他們就站住,臉上就帶著愁容.......一個叫革流巴的回答說:「你在耶路撒冷作客(守逾越節的香客),還不知道這幾天在那裡所出的事嗎?」(路二十四章十七,十八節)

耶穌與他們同行,贏取聆聽他們的權利

他們停下腳步,面對耶穌,並認為耶穌是從耶路撒冷作客,應該知道發生甚麼事.兩位門徒因為耶穌陪著他們走了一段路程,聆聽他們,認定耶穌是從耶路撒冷而來的同路人,願意分享內心的憂傷:「我們素來所盼望,要贖以色列人的就是他...... 而且這事成就已經三天了.再者,我們中間有幾個婦女使我們驚奇；她們清早到了墳墓那裡,不見他的身體.......」(路二十四21) 兩人開始很長的描述.

耶穌忍耐細心聆聽,沒有打斷他們,進入他們的悲傷情緒中:愁煩,失望,無奈.「已經三天了」是絕望的嘆息,猶太人認為死後人的靈魂離開身體,死後三天身體開始腐爛,改變原來的形狀,靈魂認不出是自己的身體,不能返回身體,死亡就成為一個不能逆轉的事實了.再者,他的屍首不見了,還有幾個婦女看到天使說,耶穌復活了.他們希望耶穌不會死,但三天過去了,他真的死了；屍首不見了,連回憶的連繫也消失了,何等無奈.

耶穌細心聆聽,贏取對他們說話的權利

跟著耶穌對他們說:「無知的人哪,先知所說的話,你們的心信得太遲鈍了.......」(路二十四25)於是從摩西和眾先知起,凡經上所指著自己的話都給他們講解明白了.(路二十四27)耶穌在這裡給我們作了福音使者的榜樣:他伴隨,他聆聽,他分享；一步一步贏取兩人的信任,得到分享信息的權利.

穌隱藏自己,以一位朝聖同路人的身份,詮釋生命的意義

在這段經文,耶穌以朝聖者同路人的身份,對兩位失望的門徒,詮釋生命和死亡的意義:基督這樣受害,又進入他的榮耀,豈不是應當的嗎?當我們低沉沮喪時,感到上主離我們不顧時,基督其實就在我們左右,他常隱藏在同路人中,伴隨著我們,鼓舞我們.天路客並不孤單,因有聖徒相伴隨.

蘇恩佩:將生命主權交出來

時代論壇復活節四月八日一期,有蘇恩佩姊妹逝世三十週年專輯；四月十二日有團體舉行「把火種撒在地上──蘇恩佩生命反思會」,七百多人出席,多位教會領袖分享蘇恩佩如何影響他們的生命.在北美,香港,台灣,星加坡,甚至大陸,蘇恩佩透過親身接觸,和所寫的文字,影響無數生命.蘇恩佩姊妹關心香港青少年工作,創辦「突破雜誌」,後來「突破雜誌」成為「突破機構」,開始多媒體事工,之後發展輔導事工,及青少年的訓練事工,成為極具影力的教會機構,香港青年事務委員會主席,突破機構榮譽總幹事蔡元雲分享她一九七八年的日記:

我就早已作最重要的決定,將生命的主權交出來......十七歲開始經歷失敗,挫折,病患,禍劫......神將我一切最重要的都奪去,除去我 自己的驕傲,虛榮,自我中心......後來有一位來探望我的宣教士說:『若神拿走你一些東西,只是因為祂想給你更好的』......是祂的愛鼓勵我,叫我背起十字架.......

我剛出來做傳道時,有幸參加「突破雜誌編委會」,有機會和擔任主席的她開會,一起晚膳.那時我看到她夏天時還頸項圍上絲巾,就猜想可能是遮蓋手術的傷痕,她做過甲狀腺癌的手術.她看來誠然是一古典美人,很有氣質,雖然身體纖瘦,但說話清脆,溫柔有信念,充滿理想.蘇恩佩在1981年出版的「死亡,你別狂傲」的序言說到她的病:

從青春期開始,我步上死亡的旅程.帶著癌活了超過二十年到底是一種怎樣的生活,是我在這裡要講的故事.科技進步帶來的治療到底是禍是福,治療會不會比疾病本身更可怕,也是我引出的問題.不過我在這裡最重要的是做一個見証人,見証在步向死亡的旅程中,在人類共同受咒詛的厄運中,上帝如何插手,將祝福注入咒詛,將一個人軟弱的受害者變成有用的器皿.......透過死,我不斷在學習生之智慧.

「死亡,你別狂傲」一書,在八十年代癌症治療方法沒有現今那麼進步時,幫助了無數癌病患者,得到鼓舞和安慰.作為天路客同路人,她為我們詮釋死亡和生命的意義.研究蘇恩佩思想的著名上海民間劇作家兼傳道人姜原來,認為蘇恩佩的作品有現代意義.她的真誠,大愛,和勇敢,正好幫助我們面對當今日益複雜破碎的世界.唯有真誠才能穿透現實,唯有大愛才能傳遞救贖；遠離至真的求索,愛便不能至大至聖,信仰充其量只能成為「小信」,與現實無關的花言巧語.

林書豪:我打球是為了榮耀上帝

上帝在我們朝聖的旅途中,隱藏在不少的朝聖同路人中,為我們詮釋生命的意義.我們不祇有離開我們三十年,體弱多病的蘇恩佩,還有生龍活虎,年謹二十三歲NBA球星,哈佛大學高材生林書豪.最新一期時代周刊選出當今最有影響力的一百人,榜首就是Jeremy Lin.一位網上基督教報刊編輯告訴我,他們因為有一篇林書豪的訪問報導,在高峰期全球一天收到近三十萬訪客.

一夕之間,「林旋風」橫掃全球,他上了全美20多家最有名的報章雜誌,他的球衣銷量連續高居榜首,他的新聞發佈會幾乎變成了他的信仰見證平臺,全世界不斷從他口中聽到「上帝是我的幫助和引導」,「我打球是為了榮耀上帝」,「基督信仰是我最大的財富和力量源泉.」這些在媒體不常見到的言論,引發各行各業對價值,成就和人生作出和以前不同的反思,「林書豪傳奇」 (Linsanity) 成為當今年輕人所喜愛的勵志故事.在時代周刊為林書豪寫介紹的是美國教育部長Arne Duncan,他說:「他以傳統的方法得到成功:他贏取成功,勤奮工作,保持謙虛,正確生活,正確比賽.......林書豪的故事告訴我們,如果你顯示出決心,紀律,誠實,你就能夠戰勝困難.」(In fact, he achieved the old-fashioned way: he earned it. He worked hard and stayed humble. He lives the right way; he played the right way...Jeremy"s story tells us that if you show grit, discipline and integrity, you too can get an opportunity, to overcome the odds.)

我們感謝上帝在朝聖的路途上,以眾多同路人,為我們詮釋甚麼是死亡,甚麼是生命,甚麼是有意義的生活,叫我們在悲傷,失望,無奈中看到盼望,他們的生命見證,使我們看到復活的基督,生命得到鼓舞.

2. 上帝以他的話語,引導我們走人生路 (vv.25-27)

「無知的人哪,先知所說的話,你們信得太遲鈍了.基督這樣受害,又進入他的榮耀,豈不是應當的嗎?」於是從摩西和眾先知起,凡經上所指著自己的話都給他們講解明白了.(路二十四25-27)

兩位門徒眼睛給蒙蔽了,因耶穌想與他們同行,聆聽感受他們的苦痛,他是我們的大祭司,擔負我們的重擔.另一個更重要的原因,耶穌要他們明白聖經.要是耶穌以復活的真身顯現,他們一定進入渾然忘我的興奮境界,不可能細聽耶穌解釋聖經.但無論屬靈的經驗如何寶貴,隨著日子的流逝,這些經驗會愈來愈模糊.而且,經驗是個人的,聽你分享經驗的人,可以羨慕你,為你感恩,但不能擁有你相同的經驗.但明白聖經就不同了,這兩位門徒立刻返回耶路撒冷,把耶穌的對聖經解釋,重複一次講給使徒聽.

在這裡我們看到一個非常重要的屬靈原則:明白聖經而來的火熱,是持久的火熱,可以傳遞的火熱.以馬忤斯的故事,叫我們重視聖經,尋求明白聖經.我想蘇恩佩和林書豪,他們能夠有這樣大的生命力,因他們勤讀聖經,明白聖經,實行聖經.

林書豪在哈佛大學時,與好友成立查經班,邀請同學參與,與一班基督徒同學常常彼此提醒,謹慎自己的生命,他說:「大學生活總是充斥著酒和各式各樣不同的女孩,我們常常討論怎麼克制自己,怎麼保持對上帝的虔誠,避免落入惡者圈套.」在教會,他參與教主日學.

他熟悉聖經,去年他分享說:「過去這一年對我來說是很艱難的一年,但也是我在神面前領受最多的一年.我越是困難的時候,神真的每一步都與我同行,他帶領我走過想放棄的每一刻.」

「聖經裡面有很多我喜歡的章節,例如羅馬書5章3-5節,講到苦難如何磨練品格,產生忍耐......每次當我被擊倒,我會嘗試再站起來重新開始,我不喜歡放棄.」「每次比賽後當我坐下來,我就禱告求問神,他的目的是什麼?我為什麼會在這裡?所以我一直和神在一起,他一直教導我要謙虛.」

聖經給我們智慧和屬靈的能力,面對人生的挫折和逆境.

3. 基督等候我們邀請他坐席,與他相交, vv. 28-31

「將近他們所去的村子,耶穌好像還要往前行,他們卻強留他說:" 時候晚了,日頭已經平西了,請你同我們住下吧!" 耶穌就進去,要同他們住下.到了坐席的時候,耶穌拿起餅來,祝謝了,擘開,遞給他們.他們的眼睛明亮了,這才認出他來.忽然耶穌不見了.」(路二十四 28-31)

在前面(15節) 耶穌親自就近他們,與他們同行,聆聽,垂詢,引導,啟發,但在這裡卻沒有要求兩位門徒留住他,他等待我們自願邀請他,進一步與他有更深入的相交.復活的基督,要門徒甘心樂意跟隨他:「當你掌權的日子,你的民要以聖潔為妝飾,甘心犧牲自己,你的民多如清晨的甘露.」(詩110: 3)

兩位門徒為何邀請耶穌留下來,跟耶穌一同坐席呢?是律法要求嗎?不錯,接待客旅是律法的要求,日已平西,前面愈來愈黑了,請留下來晚膳和過夜吧；但耶穌好像還要往前行,夤夜趕路,若是這樣,他們也不須強留耶穌.我想他們很想繼續認識聖經,不是律法要求,是渴想更深認識聖經,渴慕上帝的話語.在他們力邀下,耶穌留下和他們坐席.

結論:以馬忤斯的故事開始時是愁煩,失望,無奈的路程,兩位門徒走在路上,夕陽漸漸西沉,天色愈來愈暗,他們正步向黑暗.故事在結束時,是一個筵席.

各位朋友,各位弟兄姊妹,今天你是否好像以馬忤斯路上的兩個門徒,感到絕望,無奈,覺得上帝離開你很遠?其實上帝沒有離開你,上帝隱藏在很多人的身上,給你支持,他伴隨著你,給你力量.

或者,不是你覺得上帝離你很遠,而是你覺得你離開上帝很遠,你感到你的靈性正走向黑暗,但上帝並有放棄你.「他好像還要往前行」,其實祇要你願意,耶穌會留下,和你坐席:「看哪,我站在門外叩門,若有人聽見我聲音就開門的,我要進到他那裡,我與他他與我一同坐席.」(啟三20)

我們可以轉夕陽為清晨-你們為甚麼愁煩呢?

我們可以轉夕陽為清晨-你們為甚麼愁煩呢?阿們.



\newpage

\section{狼來了怎麼辦}
\label{sec:147}
\textbf{鄧瑞強博士 - 狼來了,怎麼辦?}
\newline
\newline
\hyperref[sec:146]{< < < PREV SERMON < < <}
~
\hyperref[sec:toc]{[返目錄]}
~
\hyperref[sec:148]{> > > NEXT SERMON > > >}


約翰福音 10:11-18
\newline
\begin{longtable}{cl}
\hline
\hline
章節 & 經文 (和合本修訂版)\\
\hline
10:11 & \begin{tabularx}{0.7\textwidth}{X} 「我是好牧人，好牧人為羊捨命。 \end{tabularx} \\ \\ \relax
10:12 & \begin{tabularx}{0.7\textwidth}{X} 雇工不是牧人，羊不是他自己的，他一看見狼來，就撇下羊群逃跑；狼抓住羊，把牠們趕散。 \end{tabularx} \\ \\ \relax
10:13 & \begin{tabularx}{0.7\textwidth}{X} 雇工逃走，因為他是雇工，對羊毫不關心。 \end{tabularx} \\ \\ \relax
10:14 & \begin{tabularx}{0.7\textwidth}{X} 我是好牧人；我認識我的羊，我的羊也認識我， \end{tabularx} \\ \\ \relax
10:15 & \begin{tabularx}{0.7\textwidth}{X} 正如父認識我，我也認識父一樣；並且我為羊捨命。 \end{tabularx} \\ \\ \relax
10:16 & \begin{tabularx}{0.7\textwidth}{X} 我另外有羊，不屬這圈裡的，我必須領牠們來，牠們也要聽我的聲音，並且要合成一群，歸一個牧人。 \end{tabularx} \\ \\ \relax
10:17 & \begin{tabularx}{0.7\textwidth}{X} 為此，我父愛我，因為我把命捨去，好再取回來。 \end{tabularx} \\ \\ \relax
10:18 & \begin{tabularx}{0.7\textwidth}{X} 沒有人奪去我的命，是我自己捨的；我有權捨棄，也有權再取回。這是我從我父所受的命令。」 \end{tabularx} \\ \\
[1ex]
\hline
\hline
\end{longtable}
日期:2012年5月20日

主題:狼來了,怎麼辦?

經文:約翰福音10:11-18

證道:鄧瑞強博士

各位弟兄姊妹,主內平安.

無論是遠古的蠻荒,抑或是現代的大都巿,都有狼出現.

早幾日,有一個返幼稚園的小朋友,在放學時,險些被四個男人強行擄走.這是發生在光天化日之下的,這是發生在鬧巿裡的,這是發生在現代文明的社會裡的.你們說,這些人「狼不狼」?所以說,我們這個社會真是有狼的.

狼來了,怎麼辦?這要看「羊」身邊有沒有牧人,也要看這牧人到底是一個怎麼樣的牧人.有報章說,那日,那小朋友不被擄走,是因為接送她的那位司機,冒死保護她.也有報章說,有人及時呼叫,嚇退了歹徒.想想,若沒有人守望,這小朋友的下場,便很難預計了.這隻小羊身邊,有守望羊的好牧人.我們多麼希望,當狼要吃我們的時候,有人願意冒死救我們.

當我們軟弱時,需要別人保護.當別人軟弱時,我們需要保護他們.生活,應是這樣子的.當我們軟弱時,卻沒有人願意提供保護,這是悲慘的.當別人軟弱時,我們不願意保護他們,這是無情的.有誰願意活在悲慘世界裡?有誰願意活在無情的人間?

只要世界仍有無知的,手無寸鐵的,容易被欺負的,像「羊」一樣的人,就會有狼.無論古代,無論現代,都出現狼食羊的事.所以,無論何時,都會有「狼來了」的故事.

今日,我要講三個「狼來了」的故事.那些羊都是無助的,那些狼都是「很狼」的,故事的分別,在於守望那些羊的牧羊人的素質.

第一個「狼來了」的故事,是「伊索寓言」的故事.

從前,有個牧羊的小孩,他在鄉村附近,牧放一群羊.他日日看著這群羊自由自在地吃草,不覺得滿足,反覺得無聊.無聊的人,為了搞事而搞事.他的心,不是記掛羊,而是尋找自己的快樂.於是,他大叫:「狼來了!狼來了!」村民是很齊心的,為了保護這小孩及那些羊,個個都出來,氣喘喘的跑上山.

有人為了快些跑上山救這小孩及那些羊,褲還沒有穿好,甩甩漏漏.有人為了快些跑上山救這小孩及那些羊,仆在地上,還仆崩個鼻.人齊了,發現沒有狼.那小孩子卻指著每個人的狼狽相逐一取笑.村民很不開心地下山回家.之後,久不久,那牧羊小孩便大叫:「狼來了!狼來了!」村民還是很齊心的,為了保護這小孩及那些羊,個個都出來,氣喘喘的跑上山.但還是發現沒有狼.

有一次,狼真的來了.牧羊小孩於是便大叫:「狼來了!狼來了!」儘管他今次叫得特別大聲,但村民想,這個牧羊小孩子又為了自己的快樂,搞些花樣來取樂.今次,沒有村民上山了.狼於是將羊逐隻逐隻吃掉了.

每次講這故事,大人都會說,這故事教訓我們,不要講大話.

這個牧羊小孩子固然是講了大話,但講大話應該只是他的表面問題,他還有深層次問題.他的深層次問題是他竟然覺得悶.「悶」是一個深奧的神學問題.作為一個守護羊的牧羊人,牧羊是他的天職,羊的命運在他手中,守護羊群是件多麼有意義的事,但他竟看不到這事的意義重大,竟感到悶.當一個人守護著眾多無助的人的生命,承擔著這件意義重大的事情時,竟感到悶,並看重自己的快樂感受多過羊的需要,這真是一種沉淪.他為了自己的快樂而犧牲別人,失去了別人對他的信任,最後所有羊都被犧牲掉.

他是一個牧羊人,卻只顧在牧羊的過程中尋歡作樂.他沒有珍惜自己牧羊人的角色,他沒有珍惜別人對他的信任,他沒有珍惜羊群對他的倚賴.他失去一切.

第二個「狼來了」的故事,是「麥兜」(McDull)的版本.(參《麥兜故事:豬來了》,\url{http://hi.baidu.com/starforever99/blog/item/4bc685442bf3a886b2b7dc8c.html})

我在網上讀這故事,這故事的版權屬「麥兜」的創作人.

從前有隻小豬,名叫麥嘜(McMug),他是一隻牧羊豬,負責牧放村裡的羊.若有狼出現,牠便要大叫「狼來了!狼來了!」藉此召喚山下的村民來幫手.

有一天,麥嘜大叫「狼來了」.山下的村民聽到麥嘜呼叫「狼來了」,都連忙拿起武器,急忙跑上山,去幫助麥嘜趕狼.

怎知辛辛苦苦跑到上山崗,卻看不見狼,只看見麥嘜在笑.麥嘜指著這些狼狽的村民說:「你們真是大傻瓜.我那裡有叫狼來了呢,我只是自己在唱歌而已.」接著,麥嘜便大聲唱起歌來,「啦啦啦,啦啦啦」.

村民很不高興,但也不能對麥嘜怎麼樣.

日子過了不久,麥嘜又重操故技,大聲叫:「狼來了!狼來了!」

聽到呼叫,村民便又急忙的跑上山.今次,像上次一樣,又被麥嘜嘲笑一番.

「狼來了」的惡作劇次數多了,受騙上山的村民變得越來越少.村民開始不相信麥嘜了.到最後,當麥嘜叫「狼來了」的時候,再沒有人,再沒有一個村民跑上山.但是,還有一個,那就是麥嘜的朋友,是全村最肥,最大汗,跑得最吃力,最慢的麥兜(McDull),牠帶著武器,氣喘喘的跑到山上來.

麥嘜嘲笑他:「你真是全村最傻的人!」

全條村現在就只有這個全村最肥,最大汗,跑得最吃力,最慢,並且是全村最傻的麥兜,每次聽到麥嘜大叫「狼來了!狼來了!」的時候,都會很慢,很辛苦,很大汗,很蠢的跑上山來,然後被麥嘜嘲笑,然後知道被騙.

故事不斷循環.「狼來了!」跑上山,被騙.「狼來了!」再跑上山,再被騙.

最後,麥嘜望著面前這位氣喘喘的朋友,這位朋友的每根豬毛都夾著幾粒大汗珠.麥嘜問牠:「你為什麼這麼傻?你為什麼到現在還相信我?」為什麼呢?事實上,麥兜再傻,也不會相信「狼來了」這謊話了.但是,為什麼牠依舊每次都跑上山來呢?

以下是麥兜當時對麥嘜講的一番話.

「我現在已不再相信你的謊話了.開始時,當然相信；現在,不相信了.但麥嘜,我不能不跑上來啊!萬一,萬一真的狼來了,只要有一次,真的狼來了,要是村民不跑上來幫助你,要是我也不跑上來幫助你,你很可能被狼吃掉了!我受騙,當然憤怒,但這不過是白跑了一場.可是,只要有一次,只要有一次我不相信你,你可能會沒命啊!」

聽過麥兜的說話,麥嘜拉著他的手,心裡又高興,又慚愧.

麥嘜覺得,這夜的山坡很美麗!

之後,麥嘜再沒有叫狼來了,麥兜也不用跑上山.由於這樣,麥兜便缺乏運動,變得越來越肥.

為了麥兜的健康,於是麥嘜改口叫:「豬來了!豬來了!」當麥兜聽到「豬來了」的時候,便跑上山,和朋友麥嘜一起,一起摘果子,一起欣賞那美麗的山坡.

作為牧童,麥嘜是完全不稱職的.

他沒有珍惜自己牧羊人的角色,他沒有珍惜別人對他的信任,他沒有珍惜羊群對他的倚賴.看來,他將失去一切.

故事的轉捩點,是因為他有一個豬嘜朋友:麥兜.

縱使被欺騙,縱使跑到身水身汗,縱使很憤怒.但是,他愛惜朋友,他為朋友著想,他寧願自己受虧損,也要叫別人得益處.他拯救了麥嘜,也拯救了整個故事.否則,這個狼來了的故事,便會再重複第一個「狼來了」的故事了.

麥兜,是一個好朋友.他是一個牧人的好牧人.他牧養了一個牧人.

友誼的最高境界,是為朋友捨命.

主耶穌說:「人為朋友捨命,人的愛心沒有比這個大的.」(約15:13)

第三個「狼來了」的故事,是「主耶穌」的版本.

經文:約翰福音10:11－18

10:11　我是好牧人；好牧人為羊捨命. 10:12　若是雇工,不是牧人,羊也不是他自己的,他看見狼來,就撇下羊逃走；狼抓住羊,趕散了羊群. 10:13　雇工逃走,因他是雇工,並不顧念羊. 10:14　我是好牧人；我認識我的羊,我的羊也認識我, 10:15　正如父認識我,我也認識父一樣；並且我為羊捨命. 10:16　我另外有羊,不是這圈裡的；我必須領他們來,他們也要聽我的聲音,並且要合成一群,歸一個牧人了. 10:17　我父愛我；因我將命捨去,好再取回來. 10:18　沒有人奪我的命去,是我自己捨的.我有權柄捨了,也有權柄取回來.這是我從我父所受的命令.

在村外,耶穌牧放一群羊.另有別的牧羊人,牧放另外的羊群.

有一天,狼來了.別的牧羊人,全是雇工,全跑了.

只有耶穌留下,守護他的羊.

狼來了一大群,牠們只想食羊,牠們發出低鳴,威脅著耶穌這個牧羊人,只想嚇走他.但是,這個牧羊人站著,站著,死命的守護著他的羊.

以下是一段模擬對話.

在天堂,一隻小羊與牧羊人耶穌閒話家常.

小羊:「我的牧人,你還記得那天嗎?」

牧羊人:「哪一天?」

小羊:「狼來了的那一天.」

牧羊人:「噢,你還記得那天.我記得,那一天,狼來了.」

小羊:「那天,狼來了,我們看見別的羊群的牧羊人,全走了.我們多麼擔心你也會走.」

牧羊人:「我怎會走.狼來了,我更加不會走.」

小羊:「那天,狼來了,你為什麼不走?別的牧羊人,全走了.」

牧羊人:「我為什麼不走?應該問,我為什麼要走?我不是雇工,我是你們的牧羊人.我走了,你們會變成怎樣?我不走,因為你是我的羊,我是你的牧人,我有責任照顧你們.天父將你們交給我,我怎能叫天父失望.」

小羊:「在世上,已沒有什麼責任不責任了.大能臨頭各自飛.」

牧羊人:「牧養你們是我的責任,這責任是我的天職.大能臨頭,並不是各自飛的時候；患難,更加讓人明白,生命不是為自己而活.患難,更讓人的生命緊密結連.」

小羊:「人人都是顧惜自己的,你不顧惜自己嗎?」

牧羊人:「我是牧羊人,理所當然顧惜我的羊.」

小羊:「你的心裡只有我們?」

牧羊人:「我心裡有你們.但我也顧念那些雇工牧放的羊,牠們很可憐,我多麼希望牠們能成為你們中的一分子.我多麼渴望能同時保護牠們.」

小羊:「狼來了,你不意識到危險嗎?」

牧羊人:「我意識到危險,但我更是意識到你們的危險.我不能留下你們,讓你們單獨面對危險.」

小羊:「你不怕狼吞噬你嗎?」

牧羊人:「怕,怕得汗像血點一樣滴在地上.但我更怕狼吞噬你們.」

小羊:「那天,我見到狼群咬你.當狼咬你的時候,你想著什麼?」

牧羊人:「我只想著,希望狼咬我便夠了,不要咬你們,不要傷害你們.」

小羊:「你在痛苦時,難道沒有想過逃跑嗎?」

牧羊人:「我怎能逃跑,我只知道,我和你的生命是連成一體的.我要保護你,正如我保護自己一樣.」

小羊:「那天,想不到狼真的來了.你的犧牲,是一件無可奈何的事.」

牧羊人:「不要說無可奈何.我是自願的.若再面對同一光景,我也會作同樣的事.能有機會牧養你們,是我人生的福氣.能有機會愛你們,是我三生有幸.能為你們受一點苦,是我生命的榮耀.我像一個戰士,為天父的大愛而作戰.若有一天,我倒下了,但若這倒下像種子一樣倒在地上,則我會懷著滿足而去,因為我知道,種子落在地上,便會結出子粒來.作為一個牧羊人,最有意義的事,莫過於保護你們免受傷害.這不是什麼無可奈何的悲劇.這是我生命必然要面對的事情.作為牧羊人,我為此而生.我能對其他所有牧羊人說,我完成了牧羊人的天職,我死而無憾.而你的生命就是我的生命.我將在你的生命裡活著.我受一點苦,而你卻體會愛,這是一件好得無比的事.我留下了愛.這種愛,將會幫助你克服生命種種的憂慮,最終,會幫助你克服死亡的恐懼.畢竟,我是好牧人,好牧人為羊捨命.」

三個「狼來了」的故事講完了.

現代社會推行「生命導師」(mentor)的做法.在不同行業,在不同學校,都有人為後輩找一個指導他的前輩.這前輩牧養著這後輩,是他的牧者.能為後輩留下深刻印象的,是那些無私的,燃燒了自己,照亮後來者的人.

主耶穌是我們每個人生命的牧人,他為我們捨命,讓我們看見生命的大道.他呼喚我們,作為基督的門徒,在人間,成為別人的牧者.狼無處不在,但只要有生命堅定的牧者在他們身邊,他們的生命便能得到保護,也看到生命的大道.

但願:榮耀歸於聖父,聖子,聖靈.阿們.



\newpage

\section{願你們平安}
\label{sec:148}
\textbf{陳曉東博士 -}
\newline
\newline
\hyperref[sec:147]{< < < PREV SERMON < < <}
~
\hyperref[sec:toc]{[返目錄]}
~
\hyperref[sec:149]{> > > NEXT SERMON > > >}


約翰福音 20:19-31
\newline
\begin{longtable}{cl}
\hline
\hline
章節 & 經文 (和合本修訂版)\\
\hline
20:19 & \begin{tabularx}{0.7\textwidth}{X} 那日（就是七日的第一日）晚上，門徒因怕猶太人，所在的地方門都關了。耶穌來，站在當中，對他們說：「願你們平安！」 \end{tabularx} \\ \\ \relax
20:20 & \begin{tabularx}{0.7\textwidth}{X} 說了這話，他把手和肋旁給他們看。門徒一看見主就喜樂了。 \end{tabularx} \\ \\ \relax
20:21 & \begin{tabularx}{0.7\textwidth}{X} 於是耶穌又對他們說：「願你們平安！父怎樣差遣了我，我也照樣差遣你們。」 \end{tabularx} \\ \\ \relax
20:22 & \begin{tabularx}{0.7\textwidth}{X} 說了這話，他向他們吹一口氣，說：「領受聖靈吧！ \end{tabularx} \\ \\ \relax
20:23 & \begin{tabularx}{0.7\textwidth}{X} 你們赦免誰的罪，誰的罪就得赦免；你們不赦免誰的罪，誰的罪就不得赦免。」 \end{tabularx} \\ \\ \relax
20:24 & \begin{tabularx}{0.7\textwidth}{X} 那十二使徒中，有個叫低土馬的多馬，耶穌來的時候，他沒有和他們在一起。 \end{tabularx} \\ \\ \relax
20:25 & \begin{tabularx}{0.7\textwidth}{X} 其他的門徒就對他說：「我們已經看見主了。」多馬卻對他們說：「除非我看見他手上的釘痕，用我的指頭探入那釘痕，用我的手探入他的肋旁，我絕不信。」 \end{tabularx} \\ \\ \relax
20:26 & \begin{tabularx}{0.7\textwidth}{X} 過了八日，門徒又在屋裡，多馬也和他們在一起。門都關了，耶穌來，站在當中，說：「願你們平安！」 \end{tabularx} \\ \\ \relax
20:27 & \begin{tabularx}{0.7\textwidth}{X} 然後他對多馬說：「把你的指頭伸到這裡來，看看我的手；把你的手伸過來，探入我的肋旁。不要疑惑，總要信！」 \end{tabularx} \\ \\ \relax
20:28 & \begin{tabularx}{0.7\textwidth}{X} 多馬回答，對他說：「我的主！我的神！」 \end{tabularx} \\ \\ \relax
20:29 & \begin{tabularx}{0.7\textwidth}{X} 耶穌對他說：「你因為看見了我才信嗎？那沒有看見卻信的有福了。」 \end{tabularx} \\ \\ \relax
20:30 & \begin{tabularx}{0.7\textwidth}{X} 耶穌在他門徒面前另外行了許多神蹟，沒有記錄在這書上。 \end{tabularx} \\ \\ \relax
20:31 & \begin{tabularx}{0.7\textwidth}{X} 但記載這些事是要使你們信耶穌是基督，是神的兒子，並且使你們信他，好因著他的名得生命。 \end{tabularx} \\ \\
[1ex]
\hline
\hline
\end{longtable}
日期:2012年5月27日

經題:願你們平安

經文:約翰福音20:19-31

證道:陳曉東博士

我們一起來思想,一起來默想,作為一個教會的群體,生命的群體,我們怎麼樣實踐復活的生命?我最近在看一本書,是一位作者Eugene Peterson寫的一本新書叫"Practice Resurrection"「實踐復活的生命」,他在這一本書中有一段落這樣說:我們的禱告是一種非常有個性的語言,或者甚麼也不是.

神是一位有性情的神,這是毫無疑問的,祂真的是一位有性格的神,而且是三而一的性情.當我們自己沒有用一種有性格的語言描述,與這位有性格的上帝溝通的時候,我們的語言是蒼白無力的.但是只要我們想一想,在詩篇中那些有性格的語言,動情之處或者是讓我們覺得非常的不安,或者是讓我們覺得震驚,但是絕對不會讓我們感覺到蒼白無力的.這不就是禱告的要點嗎?沒有性格的,沒有性情的的禱告,當中的遣詞造句,只是說給我們自己聽而已.有性格的祈禱是打開心門,向上帝述說.

我們今天讀的福音書,是一段有性格的福音書,我們聽見多馬的禱告,是一個有性格的禱告,同樣在我們所看到這樣一個生命的群體當中,與耶穌基督的相遇,也是有血有肉,是真實的,是有性格的,因為復活的耶穌,祂是有血有肉的臨格在信仰的群體中,所以我們今天所讀的經文,我們看到了一句耳熟能詳的經文,耶穌在這經文裡向我們連續講了三次的「願你們平安!」這樣的一個話語,不是一個蒼白無力的,不是一個空洞的問候,而是一個有性情的,有性格的問候,是帶著一個力量的問候,是帶著一個改變我們生命的問候.如果我們觀察一下,耶穌基督在這樣一段經文裡所說的三個問候,它的場景,它的背景,你會發現,祂每一次的平安的問候, 是在不同的處境當中一個的問候.

第一次當祂說「願你們平安」的時候, 我們看到這一群的門徒是在恐懼當中, 耶穌向他們發出了一個問候──「願你們平安」；第二次向他們發出一個「願你們平安」的這樣的一個問候的時候,我們看到是一個是在使命當中的一個問候,耶穌基督的問候是帶給了這些門徒,信徒有一個生命的使命和生命的方向；第三次的問候,也是我們所熟悉的多馬和耶穌基督所相遇的一個經驗,這樣的一個問候我們看見了耶穌基督的一個邀請,邀請多馬在這樣一個平安當中,用他的手來探入聖子耶穌基督的傷口.

約翰福音──它是出自主後大概四十年到一百或一百二十年左右,以使徒約翰為首一個信仰的群體,我們知道在新約中,它是一個統一的經卷,但同時它也是一個多元的聲音,我們在新約中聽見了不同的信仰群體與耶穌基督的相遇的記載.約翰福音,約翰一二三書,啟示錄等等這歸類在約翰名下的著作,跟以約翰為首的信仰群體是非常密切的.經學家他們推論這個群體是怎麼樣成型,大概的共識是在耶路撒冷的會堂中成型的,也是在耶路撒冷所發生的大迫害中四散.最近的一些研究指出,這樣的一個約翰的群體它最後在以弗所一帶發展教會,建立信徒的團體,約翰福音的成書與四散在各地的約翰的群體是大有關係.這歷史背景對我們有什麼亮光?有什麼意義?因為從這背景裡我們可從經文中折射出他們生活在充滿衝突,危險和孤絕的狀態中.這樣生命的經驗和體會,其實距離我們今天距離衝突,危險和孤絕的生命狀態並不遠.這讓我們再思考,再體會耶穌基督向我們說「願你們平安」的時候,在我們今天以及在經文世界當中有更深的連結, 更深的共鳴與感觸,我們看到這事情是發生在七日的頭一日,是在晚上,是新一週的開始 (20:19).

「七日的頭一日」是安息日之後的「第一日」, 本來這是猶太人按著聖經舊約的立法每一週所過的宗教的日子,但是「七日的頭一日」在耶穌基督復活的日子, 卻成為主的日子.如果我們看見門徒所經驗的世界,這裡告訴我們門徒所在的地方,因怕猶太人門都關了.這樣的一個群體並不是我們今天所理解的基督的群體,這樣的一個群體也不是我們以為的教會在今天的狀況,更大的是反映著信仰上的群體在一種自閉的,門戶封閉,自我鎖閉的狀態中,這種經驗跟我們許多現代人的經驗其實是相當相近的.

有一位挪威的哲學家 Lars Svendsen幾年前寫了一本書叫「恐懼的哲學」,二零零八年出版的.在這本小小的,薄薄的書裡他提到我們現代人生活在恐懼的經驗,就像我們今天看的經文裡「怕」──怕這個字,恐懼這個經驗,在我們的世界裡彌漫著,成為一種病毒.如果大家的經驗和記憶還新鮮的話,你可能還記得在一九九九年的時候,整個世界陷在一種非常大的焦慮以及恐懼當中,因為當時我們已經進入了資訊的世界,但據說千年蟲會把整個資訊世界擾亂崩潰,你可能記得許許多多的專家學者在探討著這樣的大災難發生的時候,當整個銀行系統崩潰的時候,當整個股票系統崩潰的時候,當整個通訊系統崩潰的時候,我們該怎麼樣來應變?我們怎麼樣在家裡儲備足夠的食水和食糧,來面對千年蟲的危險? 我們知道那件事情沒有發生,但是我們對當中這恐懼的經驗記憶猶新,不只是千年蟲,事實上在我們今天所生活的世界裡,這樣一種恐懼的經驗,這樣一種害怕的經驗,深深的根植在我們的生活當中,許多人在今日很習慣把我們的資料儲在電腦的硬碟中,但是許多專家學者提醒我們可能一個新的黑暗世界會來臨,當一個新的黑暗世界來臨的時候,這些儲蓄在硬碟中資料可以在一瞬之間全部消失. 整個我們人類的努力中所建立的知識,智慧,記憶可以在一瞬間消失,這是一種非常大的恐懼經驗,乃致於許許多多辛勤的電腦工程師們想盡辦法,千方百計確保這樣的事情不會發生,但是在這樣的努力背後,我們深深的怕核災難,我們經驗到整個的社會發展,科技發展帶給我們的害怕,生態危機在我們身邊每天發生著,當我們看見整個氣候的變化,我們心裡很深很深的害怕；金融災難──我們很深很深的害怕；恐怖主義──我們很深很深的害怕；政治的失勢,茉莉花的革命──各地政權的不穩定,讓我們心裡也擔憂著,害怕著；流行的病毒又在夏天裡開始出現,我們心裡害怕,因為二零零三年的記憶距離我們並不遠,更別說我們所說的食物的安全,癌症的流行,愛滋病的氾濫,青少年的濫藥的問題,甚至在這樣一個情慾失控的時代裡,攣童狂的出末,讓我們傳統上以為是價值最後的堡壘──學校,甚至教會都失守.

英國的統計是這樣講,在英國今天這個國家裡,每十四個人就擁有一部監控器,整個國家是架設在監控器下,來防止,來確保我們的安全,背後是很深很深的害怕,這位哲學家說,諷刺的世界是這樣的事,二十一世紀,其實是人類歷史上比任何的世代更安全,更健康的一個世代.但是為什麼我們生活在相比任何的一個世代都安全,都健康的時代裡,我們心裡面這麼害怕呢?哲學家說,那是因為我們的恐懼源自於當代我們的消費型的社會生活:Our fear is a by-product of luxury.傳媒帶給我們很多的害怕,壓力團體告訴我們這個世界,這個社會制度,這個環境是多麼的可惡,多少的陰謀家試圖去毀滅我們美好的生活,各種流行在世界上的恐懼現象,讓我們在一個放大的狀態裡,每一個人感受到那力量的強大,事實上害怕是我們的本能,害怕它最正面的力量,他可以拯救生命的.所以我們教孩子小的時候,常常教他,站在馬路邊,看著車子來來去去的時候,我們千叮囑,萬叮囑跟孩子說:小心啊,小心,你不小心的話,車子就會把你撞倒的.我們用他這恐懼的經驗來保護他,來讓他自己在生活裡去小心,關注自己的安全,這本來是好的,但恐懼失去了控制的時候,它能奪去我們內心中,我們世界中的自由感,它可以扼殺我們人的創造力,它可以使我們行善的動機都遭到毀滅,我們從每天所發生的經驗和新聞中,我們並不是沒有這樣深的體會.就像當我們看到新聞裡,有孩子在路上被車子撞倒了,一個行人走過,看見了不理；兩個行人走過,看見了不理；三個行人走過,看見了不理；所有的行人走過,看見了不理.不是說他們沒有善良的心,他們因為在傳媒的報導下,在法律法規嚴控下,他們心裡有很深的害怕,如果我去幫這個孩子的時候,他的父母會不會過來反告我傷害了這個孩子呢?我幫助人以後會不會讓被幫助者過來反控我,欺詐我,剝削我?我們看到現代人生命中這種的怕,甚至到了一個地步,扼殺了我們自己內在的這種善良的性情.

害怕,恐懼帶來了暴力,這更是我們耳熟能詳的事情.一個不受控制的恐懼,一個沒有節制的害怕,讓這個社會成為黑森林.讓我們自己人與人之間,鄰裡的關係是崩潰狀況,從我們自己的經驗,你走過所有的大廈,你看到每一棟房子,是一層一層的伽鎖.我們在今天生活在管理非常完善的屋村裡,很好的管理和保安公司在管理地方,但是我們的心裡還是很怕,怕什麼?所以買了一個厚厚的木門裝了第一個門,然後再買了一個厚厚的鋼閘再裝上,有時候我們走過我們生活的世界,也不禁問自己到底是住在監獄中,或是在生活在鄰裡彼此關懷的世界中?怕──沒有節制的怕,使我們的生命世界處在非常混亂的狀態裡.

那天的晚上,門徒所在的地方,因為怕猶太人,門都關了,耶穌來了,站在當中對他們說:「願你們平安!」復活之主臨在我們生命當中,向我們發出一個很深切的問候,但是這一個問候不是沒有實質的問候,我們看到經文二十節──「說了這話就把手和祂的肋旁他指給他們看」.耶穌基督祂受傷的手和祂的肋旁,成為門徒喜樂的原因,成為耶穌這些門徒從害怕中得釋放的原因,為什麼是這樣的事情?我們看到了下面的經文,耶穌繼續對他們說:「願你們平安!」父怎樣差遣了我,我也照樣差遣你們,說了這話就向他們吹一口氣說「你們受聖靈」.所以在今天我們看到的經文裡,第二個給我們大的啟示,給我們大的體會,以及什麼是生命方向?這樣的平安,在恐懼害怕裡能釋放出來的喜悅,是因為當我們在使命中實踐耶穌基督的生命的時候,我們就能在生命中得到大的平安.所以經文第二點告訴我們,在使命中我們能得平安,從耶穌基督的話語裡,我們看到了祂所說的幾個簡約但對教會的經驗來說,是非常非常重要的話語,祂說:「願你們平安!父,天父上帝怎樣差遣了我,我也照樣差遣你們」.

從這裡我們看到了古教會所謂三一信仰中的差遣,上帝是在祂的三一的性情中, 在祂三而一的位格中差遣祂的聖子,而聖子在這三一共融的關係中差遣門徒,所以真正的平安是怎麼樣的平安?因為二十二節──祂所說這話,就向他們吹了口氣.這個氣是什麼?我們知道這氣是聖靈的氣,是耶穌基督的氣息,是耶穌基督的聖靈.

在約翰福音裡,我們看到整個生命的平安,能力,,喜悅,力量,能使我們從恐懼中,害怕中走出來的,是因為我們選擇了一種更高生命的價值,更大的生命力量,這生命力量是來自父在基督裡所賞賜給我們的氣息.這樣的氣息,就是我們所說的聖靈.歷世歷代許許多多的牧者,傳道,神學家試圖去明白這樣一個「聖靈的氣息」究竟是什麼的一回事.

有一位的古教父奧骨斯丁,他說:「你看,在這裡的氣息其實是緊緊的扣緊了約翰福音,約翰一二三書裡所講的一個非常的關鍵」.這關鍵詞是什麼?就是愛. 就是我們詩班剛唱詩歌所唱出來的關鍵詞──因為愛.因為愛,所以我們的生命能克服恐懼；因為愛,所以我們的生命可以勝過那些黑暗的力量；因為愛,我們走在黑暗中不必吹口哨,但走得心平氣和,心安理得,因為耶穌基督向祂的子民吹下了一個愛的氣息,讓我們在這世界上成為祂的顯現,成為祂的見證,成為祂的光,因為這光就是愛的光,因為這光就是帶給我們力量,希望,信心的力量.

基督教的使命是一個什麼樣的使命?基督教的愛是一個什麼樣的愛?基督教的福音是一個什麼樣的福音?歸根究底就是耶穌基督所吹下的這一個氣息,就是使我們的生命不論在多麼的破敗當中,不論是在多麼的有限當中,不論是在多麼的軟弱當中,我們向上帝呼求這樣的氣息,這樣的能力.

在今天我們所看的經文裡第三方面,我們看到了另外一個不一樣的圖畫,這圖畫是告訴我們,如果我們在這愛的當中,如果我們在耶穌基督的氣息當中,我們就更能體會聖子耶穌的傷口是什麼的意思,也給了我們勇氣和能力去面對釘痕的手,這在肋旁的傷口.多馬,教會神學傳統喜歡叫他「多疑的多馬」──一位懷疑者的多馬,一位就算聽見了門徒的見證也不相信的多馬.常常我們對多馬的評價是覺得他是比較負面的人物,一個沒有信心的人物,小信的人物,但我們用不同的角度來看多馬,我們會為多馬對信仰中真實的實踐有更深的體會.二十四節聖經告訴我們,多馬稱為低土馬,因為他有個雙生子.耶穌來的時候,他沒有與他們同在,門徒就對他說,我們已經看見主了.多馬卻說,我非看見祂手上的釘痕, 用手指探入那個釘痕,又用手探入祂的肋旁,我總不信.我們常常說,唉!你看這多馬,多狂妄,多不敬,多離譜!我們聽見主的名,我們已經肅然起敬,他竟然說我要把手探到祂的傷口中?

我有一個朋友,他在大學裡做教授,他也在教會裡參加了很多年的聚會,但是他常常聽到我們牧者傳道分享的時候,他總是覺得言不及意,總是覺得隔靴捎癢,覺得這只是一些空話,飄浮在空氣中嗡嗡作響,沒有意義的聲音.有一年的學校暑假的時候,他去了印度的加爾各答 Calcutta 那邊,他去參加一個義工的活動, 去參加特蕾莎修女的仁愛修道會.你知道仁愛修道會的工作非常簡單,只是在印度的街道上,把那些垂死的,全身非常骯髒的,非常臭的,全身都是流膿流血的,或是餓到已經是皮包骨頭的那些將死的人撿回家,幫他們清潔乾淨,讓他們離開這世界的時候,走的體面,走的有尊嚴,走的像一個人一樣,就是這麼簡單的工作.我這位朋友他到了加爾各答,進入到仁愛修道會,他本來想去靈修,他要去跟特蕾莎修女學習靈修的法門,然後到了仁愛修道會的時候,大家忙得不可開交,然後他說我可以幫些什麼的?負責接待他的修女說,到街上幫忙好嗎?他是第一次嘛,什麼事情也不知道,那就去了,他便跟著另一個從美國來的朋友到街上撿人去.第一次碰到了將臨死亡的,看到了蒼蠅在身上到處爬,皮膚在流膿,一接近他,你知道加爾各答很熱,身體發臭再加上嚴熱的天氣,你可以想像發出來的氣息比我們香港四月的梅雨天你那衣櫃所發出來的衣服的臭味,相比之下是千倍以上.bang !一下子臭味,熱氣,流膿的傷口,腐爛的肌膚,將臨的死亡,所有這些的元素加在一起,在他的生命中,展現出巨大衝擊的力量,bang 一下撞上來... 他說我在那裡,就在那人類的流膿,創傷,死亡那邊──我遇見了神,我遇見了神.

我想這是多馬當他把手探進耶穌基督的傷口的時候,探進耶穌基督的肋旁的時候,這樣的經驗.他只是在生命中最深的傷口裡,他發出了這樣的呼聲──我的主,我的神!

弟兄姊妹,復活節帶給了我們許多的喜悅和盼望,但是這樣的喜悅,這樣的盼望,不是一個沒有深度的喜悅和盼望,不是一種膚淺的喜悅和盼望,這樣的喜悅,這樣的盼望,這樣的信心,是深深的根植在耶穌基督祂釘痕的手以及肋旁的傷口裡.整個基督的信仰,是讓我們緊緊的依附在耶穌基督釘痕的手,創傷的肋旁裡.祂讓我們這些已經生活在中產世界的弟兄姊妹,讓我們的生命不會再不知不覺中,跟著這世界,陷入了沒有邊際,沒有節制的恐懼中.而是呼喚著我們再一次的在耶穌基督釘痕的手,肋旁的傷口那邊思想我們的人生,信仰,我們在耶穌基督裡,是否常常伴隨著這位有釘痕的主,在這世界上活出一份的愛,活出一份的生命的相信,活出一份生命的盼望.我們這個世界不斷的向我們散佈恐懼的訊息,但復活節告訴我們每一個弟兄姊妹,讓我們行在耶穌基督的恩典裡,與耶穌基督的傷口同行,乃致於這世界卑微的聲音,受苦的聲音,能讓我們聽見.因為在那裡,我們可跟多馬一起宣告──我的主,我的神──願主成就祂的話語.



\newpage

\section{常在我心間}
\label{sec:149}
\textbf{林豪恩先生 - 常在我心間}
\newline
\newline
\hyperref[sec:148]{< < < PREV SERMON < < <}
~
\hyperref[sec:toc]{[返目錄]}
~
\hyperref[sec:150]{> > > NEXT SERMON > > >}


約翰福音 15:1-8
\newline
\begin{longtable}{cl}
\hline
\hline
章節 & 經文 (和合本修訂版)\\
\hline
15:1 & \begin{tabularx}{0.7\textwidth}{X} 「我就是真葡萄樹，我父是栽培的人。 \end{tabularx} \\ \\ \relax
15:2 & \begin{tabularx}{0.7\textwidth}{X} 凡屬我不結果子的枝子，他就剪掉；凡結果子的，他就修剪乾淨，使枝子結果子更多。 \end{tabularx} \\ \\ \relax
15:3 & \begin{tabularx}{0.7\textwidth}{X} 現在你們因我講給你們的道已經潔淨了。 \end{tabularx} \\ \\ \relax
15:4 & \begin{tabularx}{0.7\textwidth}{X} 你們要常在我裡面，我也常在你們裡面。枝子若不常在葡萄樹上，自己就不能結果子；你們若不常在我裡面，也是這樣。 \end{tabularx} \\ \\ \relax
15:5 & \begin{tabularx}{0.7\textwidth}{X} 我就是葡萄樹，你們是枝子。常在我裡面的，我也常在他裡面，這人就多結果子，因為離了我，你們就不能做甚麼。 \end{tabularx} \\ \\ \relax
15:6 & \begin{tabularx}{0.7\textwidth}{X} 人若不常在我裡面，就像枝子被丟在外面，枯乾了，人撿起來，扔進火裡燒了。 \end{tabularx} \\ \\ \relax
15:7 & \begin{tabularx}{0.7\textwidth}{X} 你們若常在我裡面，我的話也常在你們裡面，凡你們想要的，祈求，就給你們成全。 \end{tabularx} \\ \\ \relax
15:8 & \begin{tabularx}{0.7\textwidth}{X} 你們多結果子，我父就因此得榮耀，你們也就是我的門徒了。 \end{tabularx} \\ \\
[1ex]
\hline
\hline
\end{longtable}
日期:2012年6月10日

題目:常在我心間

經文:約翰福音15章1-8節

證道:林豪恩先生

一．經文

15 : 1「我就是真葡萄樹,我父是栽培的人. 15 : 2凡屬我不結果子的枝子,他就剪掉；凡結果子的,他就修剪乾淨,使枝子結果子更多. 15 : 3現在你們因我講給你們的道已經潔淨了. 15 : 4你們要常在我裡面,我也常在你們裡面.枝子若不常在葡萄樹上,自己就不能結果子；你們若不常在我裡面,也是這樣. 15 : 5我就是葡萄樹,你們是枝子.常在我裡面的,我也常在他裡面,這人就多結果子,因為離了我,你們就不能做甚麼. 15 : 6人若不常在我裡面,就像枝子被丟在外面,枯乾了,人撿起來,扔進火裡燒了. 15 : 7你們若常在我裡面,我的話也常在你們裡面,凡你們想要的,祈求,就給你們成全. 15 : 8你們多結果子,我父就因此得榮耀,你們也就是我的門徒了.

二．釋經與應用 1．枝子連接葡萄樹,人連接神

耶穌把自己與門徒的關係以葡萄樹與枝子作比喻,雖然是耳熟能詳,但若我們多想一想,意義相當深遠,也是令人感動的.

神學家巴特說過一句聽來很簡單,但對當代神學界影響深遠的說︰「上帝就是上帝」.這句說話要強調的是上帝與人類的差異.人與上帝之間的差異,不是量的差距,不是上帝的能力大一些,人的能力少一些；人與上帝之間的差異,是本質上的不同.神是神,人是人,這本質上的距離令上帝和人類之間並無任何的連續,人不可能透過修練把自己變成神,也不可能透過探索自己的內心而發現神.人與上帝本來是沒有可能連接在一起的,然而,耶穌卻在這裡對門徒說:「我就是葡萄樹,你們是枝子」,即是把人和上帝連接在一起,把本來不可能連接的連接起來.如果人類有一點自知之明,都不會否認自身的醜陋與罪污,如何能夠想像自己與聖潔的上主連接在一起呢?完美無瑕的耶穌基督與滿身罪污的人類連接,他會有怎樣的感覺呢?他會感到委屈嗎?他會感到羞恥嗎?

學友社在2009年的一項調查中,訪問了600多個年輕學生,被訪者竟有30\%的人願意做英國人,而選擇願意做中國人的竟然與選擇做日本人的比例一樣有四分一.為甚麼在香港居住的中國人不願意與中國拉上關係呢?有報章評論如此認為:「答案其實呼之欲出,中國形象確係太差.中國形象之差,罪魁禍首主要有兩大類,第一類係黑心食品,第二類係白色恐怖.黑心食品五花八門無處不在,毒奶粉,坑渠油,瘦肉精,罄竹難書.北京有一位小妹妹,只不過食jor 街邊炸雞,居然劇毒攻心小命不保.至於白色恐怖,又係令人悲憤莫名.維權人士只不過發表幾篇文章,居然構成一大座以言入罪ge 文字獄,罪名係極之誇張ge 煽動顛覆國家政權.你係毒奶粉受害家長想討回公道咩?對唔住,你犯jor 尋釁滋事罪.所以咯喎,真係難怪市民唔想認自己係「中國人」,正如姓秦ge 人從來唔會認秦檜係祖先.」(東方日報2011年06月29日)

香港學生感覺到自己與中國現況的差距,因此,羞於承認自己與中國的關係.耶穌又如何看待自己與我們連接在一起這狀況呢?希伯來說如此說:「他稱他們為兄弟也不以為恥」(2:11).雖然耶穌不以為恥,但神與人始終存在著極大的差異,耶穌與人連接,是貶低了神,還是抬高了人呢?其實兩者皆是.腓立比書說:「他本有上帝的形象,不以自己與上帝同等為強奪的；反倒虛己,取了奴僕的形像,成為人的樣式；既有人的樣子,就自己卑微,存心順服,以至於死,且死在十字架上.(2:5-11)為了成全神與人的連接,耶穌首先降卑,透過降卑,然後把人抬高.腓立比書繼續說:「他要按著能叫萬有歸服自己的大能,將我們這卑賤的身體改變形狀,和他自己榮耀的身體形狀相似.」(3:21)

如果我們誠實現對自己,縱使沒有保羅如此強烈的感覺,「在罪人中我是個罪魁」,也不得不承認自己的軟弱及過錯.我是誰呢?配得與耶穌基督連接在一起嗎?細心思想,天父竟然把我們與他的愛子連在一起,怎能不叫我們感動呢?

2．連接的效果:結果子

耶穌繼續說:「枝子若不常在葡萄樹上,自己就不能結果子；你們若不常在我裡面,也是這樣．．．常在我裡面的,我也常在他裡面,這人就多結果子」耶穌用枝子與葡萄樹連接長出葡萄的自然現象,指出門徒與主耶穌連接,也應該會結出果子.葡萄樹技子所結的果子,當然是葡萄；門徒所結的果子又是甚麼呢?我們會否想到保羅在加拉太書的說話呢?「聖靈所結的果子,就是仁愛,喜樂,和平,忍耐,恩慈,良善,信實,溫柔,節制.」(5:22)

門徒與耶穌連結結出果子,是自然的事情,同時,也是非常有意義的事情.我們祈禱說「願人都尊你的名為聖」,也在四處掛上寫著「榮歸上帝」的字句,我們如何能夠「榮歸上帝」呢?原來是透過我們生命的成長.耶穌說,「你們多結果子,我父就因此得榮耀」.耶穌把自己比喻為真葡萄樹,把天父比喻為栽培的人,我們結出屬靈的果子,天父看見自己的努力和心思沒有白費,就感到心滿意足,老懷安慰.另一方向,結出果子對門徒自己來說也是一種肯定.耶穌說,「你們多結果子．．．就是我的門徒了」.我們可以透過在自己身上的屬靈果子體認自己與耶穌連結的狀態.

葡萄不是一天生長,一晚成熟的,門徒的屬靈品格也不是突然出現的.屬靈品格的培育是生命成長的旅程,是人生的成熟過程.摩西年少氣盛的時候,看見他的希伯來人同胞受到埃及人欺壓,他為兄弟出頭,執行私刑,暗殺了那個埃及人,並且毀屍滅跡.萬料不到,抓為他痛腳的竟然是他拔刀相助同胞,他害怕得落荒而逃.逃亡期間遇上耶和華上帝.幾十年之後,民數記的作者如此描述摩西:「摩西為人極其謙和,勝過世上的眾人.」(12:3)這是摩西所結出的果子,是他的屬靈品格的成長.

我們可以檢視一下自己,與耶穌連結後,有沒有少了一點自我中心,多了一點仁愛呢?有沒有少了一點憂愁,多了一點喜樂呢?有沒有少一點製造紛爭,帶來多一點和平呢?有沒有少一點暴躁,多了一點忍耐呢?有沒有少一點殘酷,多了一點恩慈呢?有沒有少一點惡毒,多了一點良善呢?有沒有少一點詭詐,多了一點信實呢?有沒有少了一點操控,多了一點溫柔呢?有沒有少了一點放縱,多了一點節制呢?保羅在以弗所書說:「我們原是他的工作(藝術品／傑作),在基督耶穌裡造成的,為要叫我們行善,就是上帝所預備叫我們行的.」(2:10)我們願意讓上帝這位藝術家在我們的生命中動工,把我們變得美善一點嗎?

一個晴朗的日子,有兩張白紙在做日光浴.這兩張白紙都潔白如雪,沒有瑕疵,他們以自己的純淨為傲.

地平線上出現一個人影,引起了白紙的注意.

白紙甲:「那是誰呢?」 白紙乙:「他手上拿的是甚麼呢?」 那個人手上拿的是一個調色盤和一隻畫筆.他流露出好奇,夢幻的眼神,他的心裡抱著一個夢想. 白紙甲:「你覺得他想做甚麼?」 白紙乙:「他好像要拿我們來畫畫,你不覺得嗎?」 白紙甲:「我絕對不容許他在我身上畫圖.沒有一個畫家可以破壞我的純淨.」

藝術家來到這兩張紙前面,並且徵求他們准許他把夢想畫在他們純白的身上.

白紙甲說:「絕不!我要保持潔白,直到我死的那一天.」於是她繼續保持著潔白,可惜最後仍不敵長期的日晒雨淋而變得稀爛.

白紙乙說:「你盡情地畫吧!我信任你.我把自己交託給你.」結果,藝術家使他變成了一幅名畫──那是一幅獨特,絕美的傑作,是藝術家心中夢想的呈現.

許多人來欣賞藝術家的這幅畫,在它的深邃與優美之下,人們重新發現自己的夢想.(「藝術家」《一生必讀的100個睿智故事》)

我們呢?我們絕對不是純淨的白紙,只不過是一個充滿自義,論斷,驕傲,怨恨,苦毒,虛榮,貪婪．．．的人,為甚麼還要保衛自己的醜陋,讓自己及自己身邊的人受到傷害呢?天父把我們連結於他的愛子耶穌基督,並且親自栽培我們,我們何不把自己交給那位創造我們的藝術家,讓他來把我們變得少一點醜陋,多一點美善呢?

3．如何連接:常在我裡面

枝子連接在葡萄樹上,我們可以想像得到是怎樣的狀況.我們與連接耶穌,到底是怎樣一回事呢?耶穌說,「常在我裡面的,我也常在他裡面,這人就多結果子」.常在主裡面,就是我們與耶穌的連結,然而,常在耶穌裡面,到底又是怎樣的一回事呢?

有人以為在主裡面會有特殊的精神狀態,例如出神狂喜.有人認為在主裡面就是與上帝神秘契合,會有甜蜜和親暱的感覺.有人以為在主裡面會有特殊的能力,例如醫和趕鬼,能預知未來.有人認為在主裡面會特別蒙福,例如得到健康和財富.

耶穌如此說:「你們若常在我裡面,我的話也常在你們裡面」.原來,常在主裡面就是主話充滿我們的狀態.當主的話常常在我們裡面,我們的生命就會結出屬靈的果子.主的說對於結果子有兩方面的功能,耶穌說:「現在你們因我講給你們的道已經潔淨了」.主的說話好像把葡萄樹枝子修剪乾淨的剪刀一樣,除去阻礙我們結出屬靈果子的騷擾.同時,主的說話指導著我們應該做一個怎樣的人.各位弟兄姊妹,主的話有常在你們裡面嗎?我們如何對待主的話呢?

從前有一個會造火的人來到一條鄉村,把造火的技術教導村民,他們可以使用這項技能在寒冷的冬天取暖,在黑夜照明,用來煮食物,甚至用來燒陶,製作盛裝食物的器皿.如果在荒野迷了路,他們可以用這項技能發出求救信號.

很快地,造火的消息從一個村莊傳到另一個村莊,大家很需要造火人,所以他到每個村子去,把這項技能傳授給想要學習的人.

這個造火人似乎給予村民他們非常需要的東西,以及更充實,更快樂的生活.因為造火人廣受村民崇拜,令領導人感到不安.於是,他們決定採取強硬的手段,在某個黑夜把造火人逮捕,並且將他殺死.

造火人不在,村民漸漸忘記造火技能,也不再造火,不過,對造火人仍念念不忘.他們建造紀念館來,用鮮花,繪畫和雕像來裝飾紀念館.他們按時聚集追憶造火人,把造火人教導的技能寫成小冊子,在聚會中誦讀.然而,由於他們不再造火,人們又再次習慣了寒冷和黑夜,吃生冷食物,也忘記了如何在火光下唱歌和跳舞.他們世世代代都對造火人懷著敬意,可是沒有火.(「造火人」《一生必讀的100個睿智故事》)

上主的說話不是放在祭壇上被恭奉的,而是存在我們心裡被反覆思想,並且除去阻礙我們成長的騷擾,指導我們活出上帝創造我們的美意.

奧古斯丁(354-430)是一位重要的神學家.他出生於北非的阿爾及利亞.雖然喜好求學,卻生活放蕩.公元386年夏天,已擔任修辭學教授兩年的奧古斯丁三十二歲.在看到基督徒的行為見證以後,他心中出現了厲害的掙扎.有一天,他坐在朋友亞立普的花園中,正為自己內心的矛盾痛苦萬分的時候,聽到隔壁一個孩子哼唱著:「拿起來讀,拿起來讀.」奧古斯丁就隨手拿起保羅書信,視線落在這幾節上:「不可荒宴醉酒,不可好色邪蕩,不可爭競嫉妒,總要披戴主耶穌基督,不要為肉體安排,去放縱私慾」(羅13:13-14).他在自傳《懺悔錄》中說:「我沒有讀下去,也沒有必要再讀下去,剛讀完這句話,一道明亮的光充滿了我的心,一切黑暗疑雲剎時消失了.」奧古斯丁被聖經改變了他的個人操守及生活態度.

程翔先生在大陸被囚禁時,開始認真地閱讀聖經.重獲自由後,他分享在讀聖經的過程中,影響他至深的一些章節.

一,「如今常存的有信,有望,有愛這三樣,其中最大的是愛.」(林前十三13) 首先,整本聖經給我最大的觸動,就是「信,望,愛」,這三個字不但幫助我解脫這個困境,而且我曾嘗試用這三個字去鼓勵一些在獄中曾經試圖要自殺的朋友.「信,望,愛」對一切有困難的人,都是指路明燈.「信,望,愛」實在鼓勵了很多人去戰勝逆境,使人重拾自信和尊嚴.

二,「你們要靠主常常喜樂.我再說,你們要喜樂.」(腓四4) 在讀聖經的過程中,我得到很大的啟發,使我知道怎樣去面對逆境和排除恐懼,又使我尋回喜悅的心.我覺得喜悅的心,對在順境和逆境中人都同樣重要.我每天祈禱,都感謝上帝賜了一顆喜悅的心給我.我就是憑藉這顆喜悅的心,令我可以以微笑去面對痛苦；以禮貌面對冷酷；以謙卑去面對狂傲；以激勵面對踐踏；以寬恕去面對迫害.所以,時刻保持一顆喜悅的心,可以幫助我們渡過痛苦的日子.

三,「你們饒恕人的過犯,你們的天父也必饒恕你們的過犯.」(太六14) 「愛,恕,謙卑」也是我從讀聖經中領受到的.我在讀聖經的過程中,學懂了以愛去包容所有的人,包容曾經傷害我的人,包括今次令我身陷囹圄的人；我也學懂用恕來化解一切的仇恨,這些都是上帝給我的恩典.很多朋友都擔心我,出獄後的我會否大肆鞭撻北京呢?據我所知,這也是政府提出釋放我與否的一個考慮.但我對這些朋友說,因為我相信了基督,基督教我用恕來化解仇恨,因此,我會說我的事情,但不會抱著報仇的心態,我只會抱著促進國家民主,自由,公正,法治,尊重人權的角度去述說我的故事.這件事不會有絲毫仇恨的味道.我很感恩,我相信了主,祂教我放下仇恨.所以,「恕」是上帝給我最大的恩典之一.(時代論壇1090期)

各位弟兄姊妹,當我們思考如何對待身邊的人的時候,如何回應我們的處境的時候,指導著我們的態度及原則的,是誰的說話呢?如果上主的說話在這些時刻對我們的反應有一些影響力,屬靈的果子就漸漸在我們生命中成長了.

三．總結 最後,想起一首簡短的詩歌,名為「陶造我生命」,歌詞說:

神願你來陶造我,使我在靈裡果子更多. 奇妙聖靈,請將心窩潔淨,剔出污穢願能成聖. 恩主我求能像你,心內柔和與謙卑. 憑熱愛熱誠感染萬千生命,彰顯美善靈性.

詩歌提到在我們生命成長的過程中聖靈的幫助.聖靈如何幫助呢?約翰福音記載耶穌如此說:「但保惠師,就是父因我的名所要差來的聖靈,他要將一切的事指教你們,並且要叫你們想起我對你們所說的一切話.」(約14:26)願我們也常常想起上主對我們說的話,讓上主的說話常在我們心中.



\newpage

\section{崇拜失樂園的逆轉}
\label{sec:150}
\textbf{伍渭文牧師 - 崇拜-失樂園的逆轉}
\newline
\newline
\hyperref[sec:149]{< < < PREV SERMON < < <}
~
\hyperref[sec:toc]{[返目錄]}
~
\hyperref[sec:151]{> > > NEXT SERMON > > >}


約翰福音 4:19-26
\newline
\begin{longtable}{cl}
\hline
\hline
章節 & 經文 (和合本修訂版)\\
\hline
4:19 & \begin{tabularx}{0.7\textwidth}{X} 婦人對他說：「先生，我看你是一位先知。 \end{tabularx} \\ \\ \relax
4:20 & \begin{tabularx}{0.7\textwidth}{X} 我們的祖宗在這山上敬拜神，你們倒說，應當敬拜的地方是在耶路撒冷。」 \end{tabularx} \\ \\ \relax
4:21 & \begin{tabularx}{0.7\textwidth}{X} 耶穌對她說：「婦人，你要信我。時候將到，你們敬拜父，既不在這山上，也不在耶路撒冷。 \end{tabularx} \\ \\ \relax
4:22 & \begin{tabularx}{0.7\textwidth}{X} 你們所敬拜的，你們不知道；我們所敬拜的，我們知道，因為救恩是從猶太人出來的。 \end{tabularx} \\ \\ \relax
4:23 & \begin{tabularx}{0.7\textwidth}{X} 時候將到，現在就是了，那真正敬拜父的，要用心靈和誠實敬拜他，因為父要這樣的人敬拜他。 \end{tabularx} \\ \\ \relax
4:24 & \begin{tabularx}{0.7\textwidth}{X} 神是靈，所以敬拜他的必須用心靈和誠實敬拜他。」 \end{tabularx} \\ \\ \relax
4:25 & \begin{tabularx}{0.7\textwidth}{X} 婦人對他說：「我知道彌賽亞—就是那稱為基督的—要來；他來了，會把一切的事都告訴我們。」 \end{tabularx} \\ \\ \relax
4:26 & \begin{tabularx}{0.7\textwidth}{X} 耶穌對她說：「我就是，正在跟你說話呢！」 \end{tabularx} \\ \\
[1ex]
\hline
\hline
\end{longtable}
日期:2012年6月17日

講題:崇拜-失樂園的逆轉

經文:約翰福音 4:19-26

證道:伍渭文牧師

你想一想,有那一次和一個人的談話,改變了你的一生?會否就是與你現在的配偶第一次見面,談得很投契,被他的態度和內涵所吸引,開始了友誼,情愛,結果成為夫婦.

今天福音書一位婦人,一位很怕面對群眾的撒瑪利亞婦人,在烈日當空的時候,在沒有人打水的敘加井旁,和耶穌的談話,扭轉她的生命.她確實改變了,她掉下水罐子,跑到城中,對眾人說:「你們來看!有一個人將我素來所行的事都給我說出來了,莫非這就是基督嗎?」(約四29) 和上帝相遇,與上帝對話,生命得到改變,正是崇拜的意思,從這段經文我們是認識到關於崇拜的真理.

1.  上主尋找人敬拜

「然而時候將到,現在就是了,那用心靈按真理敬拜父的,才是真正敬拜的人,因為父在尋找這樣敬拜他的人.」(約四23, 24, 新譯本)

整個故事的背景是門徒都離開到城裡買食物去了,耶穌獨自在敘加的井旁等候這婦人.這段是經典論及崇拜的

經文:崇拜是上帝主動親近人,賜恩典給人,崇拜是上帝與人的對話.我們提到崇拜,就想到對上帝的榮頌,讚美,給上帝本身配當的敬拜,英文崇拜worship (worth ship)不是就有這意思嗎?有些教會就把崇拜開始的唱詩一段時間稱為敬拜,由敬拜隊帶領會眾唱詩,唱完詩說「敬拜完畢」.

但德文敬拜一字「Gottesdienst」service of God-上帝的服事,就指出崇拜是相向的:崇拜既是人給上帝的服事,也是上帝給人的服事.「因為父在尋找這樣敬拜他的人」,指出是上帝主動親近人,像耶穌主動親近這撒瑪利亞婦人.整個崇拜禮儀,是我們給上帝的服事,也是上帝透過禮儀給我們的服事.而我們對上帝的服事,是回應上帝先給我們的服事,這也是舊約時代列祖的崇拜經驗.

挪亞築壇獻祭,因為看到上帝的救恩,上帝曉喻他洪水將至,要預備方舟,進入方舟才能得救,洪水過後,挪亞帶領一家築壇獻祭.上帝呼喚亞伯蘭離開吾珥,使他成為萬國之父,他離開吾珥,跟隨上主,築壇獻祭.上帝呼喚摩西,帶領在埃及為奴的百姓,來到西乃山,頒下律法,作上帝的子民,君尊的祭司,敬拜上帝.

「你將百姓領出來之後,你們必在這山上事奉我；這是我打發你去的證據.」(出三12)所以敬拜是回應上帝的召喚,親近上帝,我們首項禮儀宣召(call to worship) 就說明這點.

上帝是施恩的主.上帝主動召喚我們親近他,不是因為他需要我們的敬拜,乃因為他願意施恩給我們.他創造人是祝福,分享他的愛.在婚禮中,我常常提醒新人,婚姻帶來父母的天職,你們要為著下一代,好好彼此相愛,使兒女在愛的氛圍中成長,使愛外溢overflow給兒女,生兒育女是分享愛.但有人生兒育女不是因為愛,是因為不滿足；自己很想學琴,但事與願違,就計劃生一個孩子,訓練他成為鋼琴家.配偶對自己不好,生一個孩子,為自己吐氣揚眉.一個不成熟的人,當然可以藉著婚姻生活,自己更趨成熟,生命得到建立,但配偶就要承受磨難了；一個成熟的人,進入婚姻時,能夠讓其配偶得到祝福和建立,配偶的磨難較少.

我初中時剛信耶穌,問我教會的傳道人:為何上帝創造人?他說:上帝造人,與人分享他的愛,好像一個小女孩,把愛分享給布娃娃,為這布娃娃梳頭穿衣服,餵娃娃奶,藉著關心布娃娃,這女孩子得到滿足,也得到成長.但這比喻神學有問題,女孩藉布娃娃得到生命的成長,但上帝是否需要創造才得到滿足?沒有世界,沒有上帝本體之外的世界,從永遠到永遠,上帝如何經驗愛呢,因為愛是對他者的關懷；沒有他者,如何經驗愛呢?希臘神話中納西斯(Narcissus)在水池看到自己俊朗的臉龐的倒影,顧影自戀,不願意離開水池,終於餓死,但沒有他者的自戀不是愛.

感謝上主,上帝本身就是愛,因我們的上帝是三位一體,在創造世界之先,父,子,靈就有「你」和「他」,就有愛的團契.耶穌在約翰福音十七章大祭司的祈禱中說:

「我在他們裡面,你在我裡面,使他們完完全全地合而為一,叫世人知道你差了我來,也知道你愛他們如同愛我一樣.」(約十七23)

所有的豐盛已在上帝裡面,他創造世界,創造人類,目的是祝福和分享他的愛.上帝尋找人敬拜他,不是因為上帝需要我們,乃我們需要上帝.上帝不需要藉著我們對他的敬拜供奉得到圓滿,敬拜是獲得神恩的機會,敬拜是上帝給我們的服事.

讓我們看看以賽亞的崇拜經驗.

「當烏西雅王崩的那一年,我看見主坐在高高的寶座上.」(賽六1)聖經特別提及當烏西雅王崩的那一年,那年舉國悲傷,因為烏西雅曾經是一位好王,使國家中興,但後來心高氣傲,擅自進入聖殿燒香,國王雖為一國之尊,但祇有祭司才能進入聖殿燒香供奉,烏西雅王僭越了祭司的職份,被上帝以大麻瘋擊殺,舉國哀慟,沒有王,國家何去何從?上帝就針對以賽亞本身的需要,讓他知道,君王有崩逝,但上主還是坐著為王.藉著崇拜,上帝服事了以賽亞.

跟著上帝引導他明白更深層的需要,上帝對以賽亞展示上帝的權能和威嚴:

「聖哉!聖哉!聖哉!萬軍之耶和華!他的榮光充滿全地.」(賽六3) 在聖潔的上帝面前他是不潔的,他需要面對自己的罪:

「我有禍了,我滅亡了,因為我是嘴唇不潔的人,又住在嘴唇不潔的民中.」回應以賽亞認罪,上帝赦免他:

「看哪!這炭沾了你的嘴唇,你的罪孽就除掉了.」(賽六7)

各位,崇拜是上主給我們的服事,上帝藉崇拜滿足我們屬靈的需要；也藉崇拜對我們說話,讓我們明白更深層的需要,我們要好好預備自己,虛心的人有福了!因為天國是他們的!饑渴慕義的人有福了!因為他們必得飽足!清心的人有福了!因為他們必得見上帝.

2. 要明白為何如此敬拜-悟性的敬拜

耶穌對撒瑪利亞婦人說:

「你們所拜的,你們不知道；我們所拜的我們和道,因為救恩是從猶太人出來的.」(約四22)

我們著重崇拜的抒發作用,較少著重崇拜的誘發作用；強調崇拜是我們給上帝的服事,較少注重崇拜是上帝給人的服事,因為我們強調要用「心靈和誠實」來敬拜神.耶穌豈不是說:

「時候將到,如今就是了,那真正拜父的,要用心靈和誠實拜他.」(約四23)

我們認為心靈和誠實就是敬拜者主觀的虔誠.其實剛相反,誠實原文指真理,心靈指聖靈,兩者都是外在於我們的恩賜,上帝賜真理,上帝賜聖靈,真理和聖靈是客觀的恩典,與主觀虔敬沒關.這婦人是虔敬的,她與耶穌討論崇拜,在撒瑪利亞城的基利心山上的聖殿敬拜對呢?還是在耶路撒冷錫安山上的聖殿敬拜對呢?當然她認為前者,耶穌的意見呢?耶穌認為在耶路撒冷的聖殿才對,因為這聖殿預表基督,他是神人相遇的聖殿,藉著耶穌人才能親近上帝,所以耶穌說:救恩是從猶太人出來的.若在任何地方都可以建立聖殿就失卻唯一的意義了.

婦人不是沒有虔誠,乃沒有真理依托的虔誠,信仰異端,異教的人也是非常虔誠,在我的鎖匙包十條鎖匙祇有一條可以開我的辦公室,拿了其它鎖匙,無論你如何誠心誠意,也不能打開我的辦公室那道門,找對了才能開門,這是客觀對錯,與主觀意願無關.真理來自上帝的啟示,人根本不能明白.你可以明白罪的字義是矢不中的,但沒有聖靈的感動,你不可能像以賽亞一樣承認:

「我有禍了,我滅亡了,因為我是嘴唇不潔的人,又住在嘴唇不潔的民中.」真理需要真理的靈光照才能明白,約翰福音中的靈指聖靈,不是人的心靈.聖靈在我們心中光照,引導我們明白真理,順服真理,聖靈就是真理的靈,

「他要將一切的事指教你們.」(約十四26)耶穌要我們有一個悟性的敬拜.

首先,我們要明白崇拜是集體的敬拜.

崇拜有模塑靈命的功能,因集體敬拜能克制自我中心.近代禮儀學者Virgil Michel就指出,崇拜的禮儀是對付個人主義毒素的解藥,在崇拜中我們經歷上主的愛,也因為上主的愛而愛其他人,而愛是靈命成長不可或缺的.他說:「在(崇拜中) 我們再一次與基督的身體連結在一起,淋浴在上主在基督裡和對世人的大愛中......因而我們重新得力,去抗衡異教的個人主義,這些自私自利,充斥在每日的思想,言語和行為.」[1]

哥林多前書十四章論到敬拜時,明言敬拜就是集體的聚會:「你們聚會的時候」(25節) ,

「婦女在會中」(34節) ,聚會時要關心參與者的福祉,保羅列舉崇拜的數項規範,就充份說明崇拜不能自我中心: 言語要人人明白,甚至圈外人也了解(十四16, 23)  要造就人(十四16)  規規矩矩按次序行(十四40, 33)  彰顯與深化基督徒的合一與相交(參十17)  每位參與者積極的投入(十四26)''詩歌,教訓,啟示,方言''崇拜是聚會,集體的聚會需要共同的禮儀,否則祇是不同的個體,在同一地點,各自修行,禮儀把不同的人造成為一個群體.我們來到教會,在崇拜開始之前都會省察認罪,這是個人認罪,但在崇拜的禮儀,應有集體認罪的時間.我們可以在進入教會前或離開教會時把獻金放進奉獻箱,這是個人的奉獻,但在禮儀中奉獻一項,是集體的奉獻,禮序中的認罪和奉獻,就有模塑我們靈命的作用,是屬靈操練.

其次,我們要明白崇拜是信仰的表白和宣示.

認罪和宣赦:我們崇拜禮儀中有認罪和宣赦.認罪幫助我們釐清人和上帝的距離,讓上帝作上帝,沒有距離,就沒有心靈移向上帝,尋求上帝的赦免；不認識自己的罪何其深厚,就不覺得恩典何其浩大,認罪是重要的操練.受洗是每天的屬靈狀況-與基督同死同復活(羅六4) ,崇拜藉認罪的禮儀,提醒自己唯獨恩典,唯獨信心,唯獨基督,才能得救.每主日的崇拜認罪禮儀,不單祈求上主赦免我們該周所犯的罪,也把我們領受過的水禮,重演一次.認罪完畢,作主席讀出約翰一書一章九節作為宣赦:

「我們若認自己的罪,神是信實的,是公義的,必要赦免我們的罪,洗淨我們一切的不義.」是上帝的應許使我們得蒙赦罪.

經課的誦讀和認信.聖經一路來是公開誦讀,給人聆聽的,直到德國人古騰堡發明活版印刷,1545年出版第一本給人閱讀的聖經.宗教改革前一千五百多年來,教會都依遵猶太會堂傳統,在崇拜中誦讀該天按節期編成的經課.耶穌在迦伯農會堂也按此傳襲

讀經:

「在安息日,照他平日的規矩進了會堂,站起來要念聖經.有人把先知以賽亞的書交給他,他就打開,找到一處寫著說:主的靈在我身上,因為他用膏我,叫我傳福音給貧窮的人.....」(路四16-18).

崇基教堂承繼了這經課傳統,以基督生平作年曆,崇拜中讀出舊約,詩篇,書信,和福音書的經課,沿用至今的有三代經課,以三年一個循環來編訂每主日的四段經文.我們現在採用的一九八三年《修訂三代經課》,為天主教,聖公宗,信義宗,長老宗,改革宗和循道衛理教會等共同採用.讀經完畢,我們唱出榮耀頌.因為四段經課分別顯出聖父(舊約) ,聖子(福音書) ,和聖靈(書信) .公元325年尼西亞會議,判定亞流主義為異端,教會通令崇拜唱榮耀頌 (Gloria Patri) 時加上去的:但願榮耀歸與父,子,聖靈, 加上一句「 起初這樣,現在這樣,以後也這樣,永無窮盡.」英文:As  it [trinity] was in the beginning,, is now, and ever shall be, world without end；「it」就是三位一體,每主日崇拜的榮耀頌,便成為正統信仰的衛道干城.

3.  崇拜是失樂園的逆轉

最後我想提的是,崇拜為何以宣召啟其端,以祝福總其成?這樣的禮儀設計,背後有何原因?何以說崇拜是失樂園的逆轉呢?

墮落是撕裂長夠信任與愛的關係,聖經不是早上創造,晚上墮落,上帝不是墮落時才來到樂園和人算帳,上帝常常到樂園和始祖有相交.分別善果子的樹表示人拒絕上帝的主權,分別善惡在希伯來文化中表示獨立自主,我們知道除了分別善惡的樹,還有生命樹,當人學習順服上帝,經過試驗,學懂對上帝的順服,就可以獲得生命樹,上帝不是要我們最終墮落,他要我們分享生命樹,正如老師給我們考試,目的要我們過關及格,不是要我們失敗.

但始祖不信任上帝的說話:

「你們不可吃,也不可摸,免得你們死.」(創三3)蛇對女人說:

「你們不一定死；因為神知道,你們吃的日子眼睛就明亮了,你們便如神知道善惡.」(創三4)墮落後上帝來到樂園,

「那人聽見神的聲音,就藏在園裡的樹木中,躲避耶和華的面.耶和華呼喚那人,對他說:你在那裡?他說:我在園中聽見你的聲音,我就害怕；因為我赤身露體,我便藏了.」(創三8-10)但在崇拜中,當聽到宣召經文,我們歡喜快樂以詩歌回應上帝的召喚,結束崇拜時,接受差遣並上帝美善的話語benediction,我們像進入天國的小孩,擱置懷疑,以童真你心,信賴上帝的慈愛.

領受祝福時,我們看到分派聖餐的聖桌十字架兩旁的的洋燭,靜靜散發出盼望的柔輝；再望過去,蒼蒼鞍山,吐露洋洋,訴說上帝信實如山堅定,上帝慈愛如海豐盈；綠樹交蔭,小鳥鳴叫,上帝的恩典環繞著我們.

接受差遣時,我們看到落地玻璃幕牆的大十字架是在外邊的,向外的十字架,不是向內的十字架；上帝關心這世界,擋住我們半邊風景的新教學大樓,提醒我們這教堂為學生而設立,馬鞍山前的民區,是我們的鄰舍,也是我們的牧區.

像以賽亞一樣,崇拜結束時我們說:我在這裡,請差遣我.阿們.

and \#8212; and \#8212; and \#8212; and \#8212; and \#8212; and \#8212; and \#8212; and \#8212;[1] ``We are united anew with all the members of the Mystical Body of Christ; we are imbued with the all embracing love of Christ for God, and of God for man....and we are freshly equipped to withstand the pagan way of selfish individualism in our daily thoughts, words and deeds", Virgil Michel, ``The Social Nature of Communion," in The Social Problem, Book Four: The Mystical Body and Social Justice (Collegeville, Minn.: Liturgical Press, 1939), 77, as quoted in Christian Scharen, Public Worship and Public Work: Character  and amp; Commitment in Local Congregational Life (Collegeville, Minn.: Liturgical Press, 2004), 23.



\newpage

\section{我知錯了}
\label{sec:151}
\textbf{鄧瑞強博士 - 我知錯了}
\newline
\newline
\hyperref[sec:150]{< < < PREV SERMON < < <}
~
\hyperref[sec:toc]{[返目錄]}
~
\hyperref[sec:152]{> > > NEXT SERMON > > >}


約翰福音 16:8-11
\newline
\begin{longtable}{cl}
\hline
\hline
章節 & 經文 (和合本修訂版)\\
\hline
16:8 & \begin{tabularx}{0.7\textwidth}{X} 他來的時候，要為罪、為義，為審判，指證世人； \end{tabularx} \\ \\ \relax
16:9 & \begin{tabularx}{0.7\textwidth}{X} 為罪，是因他們不信我； \end{tabularx} \\ \\ \relax
16:10 & \begin{tabularx}{0.7\textwidth}{X} 為義，是因我到父那裡去，你們將不再見到我； \end{tabularx} \\ \\ \relax
16:11 & \begin{tabularx}{0.7\textwidth}{X} 為審判，是因這世界的統治者已受了審判。 \end{tabularx} \\ \\
[1ex]
\hline
\hline
\end{longtable}
日期:2012年6月24日

主題:我知錯了

經文:約翰福音 16:8-11

證道:鄧瑞強博士

各位弟兄姊妹,主內平安.

有個故事是這樣的:

一個少女回到教會,找著牧師,向牧師說:「牧師,我犯了罪,請問天父會否原諒我?」牧師問:「你做了什麼事呢?」少女說:「牧師,我忍不住要時常照鏡,然後,我又忍不住對自己說,我真漂亮.我覺得,這是極大的虛榮心.我犯了虛榮的罪.」牧師仔細地看看這少女,然後說:「啊,親愛的姊妹,在我看來,這不算是罪,很明顯,這只是對事實的錯誤理解.」

不是「罪」而認「罪」,可能只是一個笑話.

是「罪」而不認「罪」,卻是一種墮落.敢於認罪的人,是勇敢的.

很多分析說,中國人是較為不願意認錯,不願意認罪的.

易中天有一篇文章,叫＜中國人為什麼不認錯＞(參:\url{http://news.163.com/10/0717/16/6BQBEJ4S00012Q9L.html}).裡面講到,中國人「越是地位高,名氣大,粉絲多,就越難〔認錯〕.要麼矢口否認,要麼一聲不吭,要麼倒打一耙,要麼把水攪混,甚至把質疑他的人統統說成是 ``文化殺手" .痛痛快快說聲 ``對不起,我錯了" 的,幾乎沒有.」

為何如此?

易中天的分析是,中國社會很早就只有終極的兩個等級:皇帝和臣民.大家看電視劇集都見到,臣子面對皇帝,誠惶誠恐.只要自己的意見與皇帝的意見不同,臣子便說:「請皇上恕罪.」晚上看劇集＜三國＞,有一集講到主子覺得某大臣的說話不中聽,便說:「你這樣跟主子說話的嗎?」無論那大臣的官有幾大,學問有幾高,道理有多好,那大臣的即時反應是:「奴才知罪,罪該萬死.」這也可能是我們有些人小時候的經驗,父親的意見我們不同意,父親不是跟我們討論,而是一巴掌的打過來.

其實,意見不同,應是講道理的時刻,透過討論,將真理辯明.但是,在權力畸形的社會裡,意見不同,不是透過討論來處理,而是透過權力使其屈服.於是,認罪的重點,其實不在乎真實的「有罪」或「無罪」,而在乎在權力結構裡應如何作「認罪」表演.易中天說:「無罪而稱死罪,哪有真實可言?不過是一種 ``姿態" .真正的錯誤,也就不會有人去認.結果,認罪也好,認錯也好,便都變成了 ``表演".」

在文化大革命的時候,全民都在作這種「認罪表演」.易中天說:「每個人也都要檢討自己,批判自己,甚至痛罵自己.那些把自己罵得最凶的,往往能得到領導表揚,群眾肯定,在掌聲中體面地下台.」

這種認罪表演,讓人在畸形的權力結構中苟存性命,並不讓人真誠悔改,也不會有生命重生.

寫《醜陋的中國人》的柏楊說:「假如你聽到中國人說: ``這件事我錯了."你就應該為我們國家民族額手稱慶.我女兒小的時候,有一次我打了她,結果是我錯怪了她,她哭得很厲害,我心裡很難過....我抱起她來,我說: ``對不起,爸爸錯了,爸爸錯了,我保證以後不再犯,好女兒,原諒爸爸."她很久很久以後才不哭.這件事情過去之後,我心裡一直很痛苦,但是我又感到無限驕傲,因為我向我的女兒承認自己錯誤.」

寫了一大段文字,只是寫一個爸爸向他的女兒認錯,這麼簡單的一件事,讀起來卻覺得這件事無限艱難.一件這麼簡單的事也做得這麼艱難,而做完後卻竟然感到無限驕傲,我們便看到,要一個中國人認錯,多麼艱難.

五旬節是記念聖靈降臨的日子,講道的主題理所當然是「聖靈」.今日,要講聖靈的其中一項工作,祂要叫人知罪及認罪.

罪在何處?

我們看一看約翰福音16: 2-3.主耶穌說:

16:2　人要把你們趕出會堂,並且時候將到,凡殺你們的就以為是事奉神.16:3　他們這樣行,是因未曾認識父,也未曾認識我.

主耶穌在這裡很簡單的指出了世人的罪的表現,也指出罪的表現背後的罪的本質.

罪的表現是什麼?

「人要把你們趕出會堂」,「會堂」對猶太人而言,類似我們的「祠堂」.這是宗教,宗族,文化,傳統等的集合點.被人趕出「會堂」,意味著從宗教上排斥你,視你為判教者；意味著從親朋戚友的關係網中離棄你,視你為不孝子；意味著從民族的認同中否定你,視你為外人.被人趕出會堂,意味著失去自己的身分,變得無家可歸.這有點似將一個未成年的小孩子趕出街.

不單如此,「時候將到,凡殺你們的就以為是事奉神」.人以為殺害你們,是事奉神的一種表現.世人犯罪,不會以犯罪的名義去犯罪的.世人犯罪,總會找一高尚理由.當某些國家要開槍殺害手無寸鐵的平民時,他們會提出神聖的理由,他們說:「這是為了國家未來的發展.」當某些民族要屠殺另一民族時,他們會提出神聖的理由,他們說:「為了人類的純潔,我們要將邪惡的種族消滅.」某些偉大的人物清除異己時,他們會提出神聖的理由,他們說:「為了要完成上天給我的召命.」多少人,以「真理」,「道德」,「神聖」,「天命」的名義,使他人離開他們賴以維生的「會堂」.一個人高舉著「事奉神」的名義,懷著「事奉神」的喜悅與熱情,去殺另一個人.這是否很悲哀?

這些都是罪惡的表現.耶穌接著說,「他們這樣行,是因未曾認識父,也未曾認識我.」人竟然以愛神,事奉神的名義行惡,是因為他們根本不認識那位創造天地的父.這位父,以聖子耶穌的生命,流露出他的生命內容.耶穌,一個自甘降卑的人,一個以愛去還報仇恨的人,一個看到人類的前景在乎人類能放下種種成見去擁抱他人的人.耶穌,他將天父的心流露出來.人類的罪,正是遺忘了這顆心.

聖靈來了,祂要指出這種人性的墮落.祂要叫人知罪.

16:8　祂〔即:聖靈〕既來了,就要叫世人為罪,為義,為審判,自己責備自己.16:9　為罪,是因他們不信我；16:10　為義,是因我往父那裡去,你們就不再見我；16:11　為審判,是因這世界的王受了審判.

聖靈要來,處理人性生命三個最重要的靈性光景:罪,義,審判.

首先,聖靈來,處理我們的「罪」的問題.

罪,是生命最基本的問題.人的生命,內注「神的形象」(Image of God),這是人類的神聖素質.這種神聖素質,嚴格講,不是一個人內在完備的.這素質,必需一個人合宜地愛另一個人的時候,才會出現.一個人與另一人相愛,愛情才會出現.單獨一個人,是無所謂「愛情」的.同樣,一個人,也無所謂「神聖素質」的.所謂「罪」,就是面對他人時,遺忘愛,也因此消滅了生命的神聖.

其次,聖靈來,處理我們的義的問題.

所謂「義」,就是行為的正當性.我們認為怎樣行,才算是正當?

第三,聖靈來,處理我們的「審判」的問題.

所謂「審判」,即 ``judgment",我們對人生與世事的判斷,是否能經得起神的判斷?我們說,我的行為對得起「天地良心」.這「天地良心」會否只是私心?這「天地良心」能否經得起真正的「天良」和「天心」的判斷?

16:8　他〔即:聖靈〕既來了,就要叫世人為罪,為義,為審判,自己責備自己.

這經文講的「責備」,指的就是像先知一樣,揭發罪,指斥罪,不容讓對罪姑息,也不容讓將罪掩埋.

聖靈來,就是揭露和指斥我們在「罪」,「義」,「判斷」這三方面的錯.

16:9　為罪,是因他們不信我.

我們的生命活不出神聖,問題出在哪裡?

問題出在我們根本「不相信」耶穌.

想想,當有人無端端打你右邊面時,你想著什麼?你很想打他一身.你根本不相信將左邊面也給他打能有什麼生命的光輝.當你看到公司裡的「衰人」做錯文件時,你想著什麼?你很想他跌一跤.你根本不相信愛你的仇敵是在現實裡可行的.當你的舊愛變得越來越不可愛,而又有新歡出現時,你想著什麼?你想棄舊迎新,撇掉舊愛.你根本不相信愛是甘願將自己釘在十架上.你忘記哥林多前書第十三章裡所講的:「愛是恆久忍耐,又有恩慈,...愛是不求自己的益處」.

聖靈指出我們根本不相信耶穌的生命就是回歸天父的大道.這是我們的罪.

我們的罪,是我們用種種「神聖」的理由,放棄生命的「神聖」.

16:10　為義,是因我往父那裡去,你們就不再見我.

有人偷懶不工作,以為這是「著數」.有學生寫論文時抄襲,以為這是「醒目」.有傳道人將福音稀釋,為了吸引更多的人,以為這是「榮耀神」.這一切,都正如約翰福音所講,有人殺其他人,還以為是事奉神.

耶穌被人殺害,這不表示他被神遺棄,反而,唯有他能往天父那裡去.

他的復活升天,是顯明他的生命是蒙天父悅納的.敵對他生命的,是與天父為敵的.

聖靈讓我們看見,我們的生命,到底是站在耶穌那一方,抑或是站在殺害耶穌那一方?

不過,人類心靈的盲目,有時會盲到極點,在殺害耶穌時,還以為是事奉耶穌.

16:11　為審判,是因這世界的王受了審判.

耶穌死了,然後,復活,升天.能復活,能升天,顯明他的判斷是對的.

「耶穌升天」這事實顯明,他對「生命應如何活」的「判斷」是對的,他要背負十字架這「判斷」是對的,他看到人的迷失,而以自己捨己的愛,作為指路的燈,這判斷是對的.

相反,那些宗教人士,為保持自己的既得利益和自鳴清高,而殺害耶穌,他們的判斷是錯的.那些政治人士,為討好群眾,為拿選票,而殺害耶穌,而犧牲無辜,他們的判斷是錯的.那些群眾,毫無正直,左搖右擺,被宣傳煽動,而殺害耶穌,他們的判斷是錯的.

聖靈來,叫我們以耶穌的生命去檢視我們自己的生命,叫我們在耶穌面前,在天父面前,真誠的說:「我知錯了.」

一個簡單的故事(參:\url{http://jameslau88.com/please_come_home.html}):

巴西近郊,一間簡陋的房舍,住著兩母女.母親叫Maria,女兒叫Christina.這個家簡單,卻溫暖.Christina的爸爸,很早便過世了.媽媽Maria當女傭,收入微薄.Christina不甘心留在這裡,她夢想著大城巿的生活.有一天,她離家走了.她媽媽心碎.這位媽媽知道,這女兒在大城巿根本無法維生,但這女兒個性倔強,不輕言放棄.然而,當驕傲遇上飢餓,她就會做愚蠢的事.Maria明白這一點,她帶著自己的照片,去到里約熱內盧這個大城巿,在酒吧,酒店,夜總會,和妓女會出沒的地方,在每處留下自己的照片.在照片背後,寫上幾句說話.很快,這母親的照片用光了,錢用完了,她只好回家.

幾個禮拜後,她女兒走下一間酒店的樓梯.她形容憔悴,面上再沒有歡笑.在樓梯盡頭,她看見一張相,一張熟悉的臉.她看見她母親的臉.她把照片拿下來,看到背後寫著:「無論你做了什麼事,無論你成了什麼樣的人,都沒有關係,請回家吧.」就這樣,她回家去了.

這位少女,以種種「神聖」的理由離家.然而,她母親的一張照片,喚醒了她,讓她看到整個人生都錯了.她遺忘了生命的神聖,她的行為根本沒有正當性,她的生命作出錯誤的判斷.她知錯了.她回家去.

我們是否願意回家?

但願:榮耀歸於聖父,聖子,聖靈.阿們.



\newpage

\section{柏拉奇圖士}
\label{sec:152}
\textbf{戴國璋牧師 - 柏拉奇圖士}
\newline
\newline
\hyperref[sec:151]{< < < PREV SERMON < < <}
~
\hyperref[sec:toc]{[返目錄]}
~
\hyperref[sec:153]{> > > NEXT SERMON > > >}


日期:2012年7月1日

主題:柏拉奇圖士

證道:戴國璋牧師

今天主題「柏拉奇圖士」是跟大家思考有關聖靈的事情.

「柏拉奇圖士」並不是什麼特別的密碼,只是聖靈其中的一個身份──保惠師(「柏拉奇圖士」是從希臘文翻譯而成,是保惠師的意思).

這名字是我們的天父以及主所授與的聖靈,以陪伴在我們身邊的身份和職份.「柏拉奇圖士」第一個意思是中保(在約翰一書第二章中,約翰曾以中保來稱呼主耶穌,這是中文的翻譯法),即是指中間人.在聖經中提及兩位保惠師,兩位中保,把我們與上帝連結起來.一位是救主耶穌基督,另一位就是我們的聖靈.耶穌基督在上帝面前,聖靈在我們的旁邊,成全了「柏拉奇圖士」的職份.

「柏拉奇圖士」的意思是甚麼?職份又包涵了甚麼內容?由於這字眼在聖經出現的次數甚少.據學者的研究,「柏拉奇圖士」(保惠師的希臘語)在新約中只出現過五次,有四次在約翰福音十四章至十六章,在耶穌介紹有關聖靈身份時用的.有一次在約翰一書,用在耶穌基督的身上,我們就明白「柏拉奇圖士」是中保的意思.因為這名稱出現的次數很少,所以要思想這字所表達有關聖靈的職份的意思並不容易.感恩的是這名稱還有以其他的形式出現,最多是以「動詞」的形式出現.我們可以從它的動詞中了解,聖經告訴我們聖靈是我們的保惠師,祂在我們身上的職份有甚麼內容?

這個動詞是有召喚的意思,即「叫喚來」.換一句話說,它的名詞是「被召喚來者」的意思,當我們說祂是「被召喚來」的時候,我們卻未明祂被召喚來作些什麼?從聖經裡面,上帝安排聖靈作我們的保惠師,祂是隨時在我們生命中,可以向祂召喚,請祂來幫助我們的人.故此,我們思考「柏拉奇圖士」第一個意思是「幫助者」.

人總會遭遇到一些處境,渴望得到一些的幫助,無論多成功的人也是如此.大約在三年前的星期一(是我的假期),因有特別的事需要處理,我需要到西貢北潭涌處理一些事務.我家住旺角,路程也不短,為了省時,我決定自行駕車去,這不是第一次去北潭涌.平常會經獅子山隧道,途經馬鞍山及西沙公路,再由大網仔方向去北潭涌,這是我認為較順暢的方法.完成事情後,沿著同一條路程回去,途中,車內只有我一人,我看看車子溫度的顯示器,心一驚,因為引擎顯示溫度的數字偏高,即引擎出現問題,我留心指針一直顯示越來越高,我知道引擎有過熱情況.大概到達烏溪沙的時候,路上車很少,只好把車停靠在一邊,嘗試看看問題所在.發現車內的一個引擎風扇出現問題,但怎麼辦?心裡很煩惱,哪裡可找到幫助?我只好默默禱告,求天父教導我該怎樣去處理.當時想到手機中有一個修理車房的聯絡電話,便與他們聯絡,我依照師傅的指示方法處理,待車降溫後,用水灌滿水箱,待車降溫後才可行,我按師傅的方法處理了,後來也要把車送到車房維修.這經驗告訴我,人有一個時刻總希望能得到幫助,因為適時的幫助就好像久旱甘露,令人感到舒暢,但不是每個人隨時隨地作我們適時的幫助,就好像修車師傅不能即時來到我身邊幫助我.聖經告訴我們,上帝願意成為我們適時的幫助,在希伯來書第四章十六節有神的應許,「所以我們只管坦然無懼的,來到施恩的寶座前,為要得憐恤,蒙恩惠作隨時的幫助」.

這句話雖然指出耶穌基督是我們的大祭司,但三位一體的上帝願意隨時隨地成就祂們各自的職事來成為我們的幫助,聖靈就是上帝召喚來作我們隨時的幫助的那一位.我們的救主耶穌知道我們生在世上,總有需要幫助的時刻,所以當祂要離開的時候,祂告訴我們,祂去但會差派聖靈來,聖靈在我們當中就成為我們隨時的幫助者,這就是聖靈──召喚者的意思,我們可以隨時隨地召喚祂,請求祂的幫助.故此,當我們有需要的時候,不要自己焦急,也不要怕麻煩上帝.我們只管向上帝禱告求聖靈的幫助,適時的幫助就會臨到.

第二個「柏拉奇圖士」的意思是「安慰者」.

「柏拉奇圖士」作為動詞的時候亦有安慰的意思,特別在舊約七十二譯本中常出現這意思.有一本很著名的希臘文舊約聖經是七十二譯本,相傳是由七十二位著名的猶太學者把希伯來話翻譯成希臘文,所以可作參考,在譯本中這個名稱可作為安慰.例如詩篇71:21:「求你使我越發昌大,又轉來安慰我」.這「安慰」在七十二譯本是用「柏拉奇圖士」的動詞形式；約伯記16:2:「這樣的話我聽了許多；你們安慰人,反叫人愁煩」.這是約伯埋怨朋友的話,安慰也是用上「柏拉奇圖士」；以賽亞書40:1:「你們的　神說:你們要安慰,安慰我的百姓」.所有都用上「柏拉奇圖士」這詞的動詞形式,所以可以說「柏拉奇圖士」是安慰者的意思.

在新約聖經──約翰福音十四章──耶穌基督介紹有關保惠師時,英文聖經用「comfortor」,中文譯作「訓慰師」,這指聖靈是有一個安撫和安慰的作用.我們說,人在某些時刻特別需要適時的幫助,同樣在某些時刻需要適時的安慰,聖靈就是隨時隨在,在我們身邊,危難中作出適時的幫助,以及在傷心的時候作適時的安慰的那位.

有一個關於美國民權領袖馬丁路德金牧師的故事,他在一次聚會中被一個白人主義者槍殺了,由於他很被愛戴,所以在安息禮拜中,有很多名人向他的遺孀慰問,他的遺孀傷心欲絕,但所有的慰問卻不及一個很沉默的人發出的慰問,她就是積奇蓮甘迺迪,是前總統的遺孀.積奇蓮甘迺迪只與丁路德金牧師的遺孀緊緊擁抱,哀傷痛哭,因為她們有同一個經歷,能了解失去丈夫的哀痛,她能了解和安慰,所以緊緊的擁抱成為了深深的安慰,比言詞更適切的安慰.

難怪聖經保羅教導我們,上帝使我們遭遇患難是叫我們能夠用祂所賜給我們的安慰,去安慰遭遇各種患難的人(林後1:4),保羅讓我們知道當我們感同身受的時候,更能用相同的同情來支持和安慰那些有相同遭遇的人.多受苦楚的人,更懂得安慰,也更能夠安慰人.聖靈就是這位訓慰師.除了馬丁路德金牧師的故事外,還有關於林肯總統的事情,近日有一本相關的翻譯書籍面世,作者講述林肯的寬大,包括邀請政敵管理國家,他只要相信這人可以協助治理國家便作出邀請.林肯總統是一位很有名的基督徒總統,他的演說都被記下,其中包括一封安慰的信,是寫給一個普通的婦人.那時正是南北戰爭,林肯總統知道這婦人的五個兒子都在參與北軍與南軍的戰事中犧牲,便立刻寫信給這婦人,一方面安慰她,同時亦讚揚她的兒子.林肯總統知道人在哀傷中需要適時安慰.後來知道婦人的兒子中,只有兩位是陣亡,其中一位成了俘虜,另外兩位成了逃兵,但仍可見林肯總統看到安慰的重要.當我們傷心時需要別人安慰,聖靈了解,祂就是我們的訓慰師,在我們身邊給予我們適時的安慰.

「柏拉奇圖士」第三個可能的意思是「訓慰師」.

在約翰福音14:16 提及保惠師中,還有一個小字──或作訓慰師──之前提及英文是「comfortor」,為什麼不譯作安慰師而是訓慰師?其實訓慰師比英文的譯來得更有意思,因為相信神在安慰的過程中有祂特別的目的.人給人的安慰是順應對方的情況給予令人舒服的慰問,上帝的安慰不單平伏憂傷的情緒,更重要的是使我們有力量從新站起來,因此上帝的安慰包括了訓誨.約伯記4:4「你的言語曾扶助那將要跌倒的人；你又使軟弱的膝穩固.」這是約伯的朋友鼓勵他,因他從前仍是這樣安慰朋友.這裡可見安慰的真正目的不只是平伏情況,更重要是使人從困苦的環境中再得力量,重新站起來.人的安慰有時只是順應人的情緒而發出來,但很多時候卻是不足夠.

聖靈的安慰包括有鼓勵和催促,好像官長鼓舞士兵在戰場上不氣餒,又好像運動場上啦啦隊,為自己的隊員打氣,鼓勵隊員在球場上發揮果效,來一場精彩的運動表演.聖靈的安慰包括訓慰,是要鼓勵我們.所以耶穌說到聖靈的工作──祂既來了,為要叫世人為罪為義為審判自己責備自己──是我們反省的內容,為叫我們在迷路上回轉,聖靈有時候會用責備的話語.

我想起一套勵志的日本劇集──熱血警校──故事關於一群青年因著不同的原因,理由要加入警隊受訓三個月,當中與教官發生的故事,他們加入的原因很多,有些是生活艱難,有的為顯威風,各有不同原因,教官以嚴厲方式訓練他們,什至乎言詞之間亦十分苛刻,他們覺得教官是故意針對,認為他是沒有愛心的地獄教官,但其實教官是努力幫助他們尋找人生目標及成為警察的目的,在嚴厲的言語中顯出關心,只是要透過這樣,好讓他們明白現實的殘酷,讓他們明白作警察時常會遇到受辱的環境,他們要有足夠的能耐去面對受辱才能成為好警察,這正如近日我們的警察局長打算動議設立辱警罪,因警察經常被辱罵.

這讓我們明白成功背後必需要有嚴格的訓練.聖靈也一樣,所以聖靈的安慰不單是讓我們感到舒暢,而是包括鼓勵,提醒,警剔和責備,叫我們在軟弱的環境中再次站起來,重新得力.聖靈以什麼話語,方法作出安慰,提醒,鼓舞或警剔呢?就是在聖經.

耶穌說:「因為他不是憑自己說的,乃是把他所聽見的都說出來,並要把將來的事告訴你們 .」聖靈是謙卑的,上帝也是謙卑的,耶穌基督也是謙卑來到世間就知道,聖靈是謙卑的祂不憑自己說話,只說從上帝所聽見的,祂用聖經的真理提醒,鼓勵我們,所以我們應該認真讀聖經,把神的話藏在心裡,因為這是聖靈用來安慰,鼓勵,提醒,警剔和責備的來源,唯有這樣聖靈便成為我們的「柏拉奇圖士」──保慰者.

我們一同禱告:我們的主,謝謝你要離開我們返回天上時,你說必不撇下我們為孤兒,你將聖靈差來成為我們身邊的保惠師,而且當時候到了,你會再來與我們相會,這是我們所得到的信息和深深相信,今天我們仍在等候你的再來,在等待的過程中,我們有不同的遭遇,在這些的遭遇中,我們需要特別的幫助者,安慰者和鼓舞者,不單帶我們渡過遭遇中的危難,也成為我們生命的幫助,使我們生命得以成長,我們向上帝呼求,向聖靈呼求.聖靈呀,請你來,作我們的幫助,成為我們的保惠師,求你垂聽我們的禱告,奉耶穌基督聖名求,阿們!



\newpage

\section{他癲狂了}
\label{sec:153}
\textbf{鄧瑞強博士 - 他癲狂了}
\newline
\newline
\hyperref[sec:152]{< < < PREV SERMON < < <}
~
\hyperref[sec:toc]{[返目錄]}
~
\hyperref[sec:154]{> > > NEXT SERMON > > >}


馬可福音 3:20-35
\newline
\begin{longtable}{cl}
\hline
\hline
章節 & 經文 (和合本修訂版)\\
\hline
3:20 & \begin{tabularx}{0.7\textwidth}{X} 耶穌進了屋子，眾人又聚集，甚至他連飯也顧不得吃。 \end{tabularx} \\ \\ \relax
3:21 & \begin{tabularx}{0.7\textwidth}{X} 耶穌的家人聽見，就出來要拉住他，因為他們說他癲狂了。 \end{tabularx} \\ \\ \relax
3:22 & \begin{tabularx}{0.7\textwidth}{X} 從耶路撒冷下來的文士說：「他是被別西卜附身的」，又說：「他是靠著鬼王趕鬼的。」 \end{tabularx} \\ \\ \relax
3:23 & \begin{tabularx}{0.7\textwidth}{X} 耶穌叫他們來，用比喻對他們說：「撒但怎能趕出撒但呢？ \end{tabularx} \\ \\ \relax
3:24 & \begin{tabularx}{0.7\textwidth}{X} 一國若自相紛爭，那國就立不住； \end{tabularx} \\ \\ \relax
3:25 & \begin{tabularx}{0.7\textwidth}{X} 一家若自相紛爭，那家就立不住。 \end{tabularx} \\ \\ \relax
3:26 & \begin{tabularx}{0.7\textwidth}{X} 撒但若自相攻打紛爭，他就立不住，必定滅亡。 \end{tabularx} \\ \\ \relax
3:27 & \begin{tabularx}{0.7\textwidth}{X} 沒有人能進壯士家裡，搶奪他的東西；除非先綁住那壯士，否則無法搶奪他的家。 \end{tabularx} \\ \\ \relax
3:28 & \begin{tabularx}{0.7\textwidth}{X} 我實在告訴你們，世人一切的罪和一切褻瀆的話都可以得到赦免； \end{tabularx} \\ \\ \relax
3:29 & \begin{tabularx}{0.7\textwidth}{X} 凡褻瀆聖靈的，卻永不得赦免，而要擔當永遠的罪。」 \end{tabularx} \\ \\ \relax
3:30 & \begin{tabularx}{0.7\textwidth}{X} 因為他們說：「他是被污靈附身的。」 \end{tabularx} \\ \\ \relax
3:31 & \begin{tabularx}{0.7\textwidth}{X} 那時，耶穌的母親和他兄弟來，站在外邊，打發人去叫他。 \end{tabularx} \\ \\ \relax
3:32 & \begin{tabularx}{0.7\textwidth}{X} 有許多人在耶穌周圍坐著，他們就告訴他說：「看哪！你母親、你兄弟和你姊妹在外邊找你。」 \end{tabularx} \\ \\ \relax
3:33 & \begin{tabularx}{0.7\textwidth}{X} 耶穌回答他們：「誰是我的母親？誰是我的兄弟？」 \end{tabularx} \\ \\ \relax
3:34 & \begin{tabularx}{0.7\textwidth}{X} 就環視那周圍坐著的人，說：「看哪，我的母親，我的兄弟！ \end{tabularx} \\ \\ \relax
3:35 & \begin{tabularx}{0.7\textwidth}{X} 凡遵行神旨意的人就是我的兄弟姊妹和母親。」 \end{tabularx} \\ \\
[1ex]
\hline
\hline
\end{longtable}
日期:2012年7月8日

主題:他癲狂了

證道:鄧瑞強博士

經文:馬可福音 3:20-35

各位弟兄姊妹,主內平安.

若果你身邊的人個個都「呃人」,而你卻指著他們說:「你呃人」,他們會認為你「黐線」.

若果你身邊的人個個都「先敬羅衣後敬人」,而你卻禮貌地對待一個乞丐,他們會認為你「神經病」.

若果你身邊的人個個都重視名牌,而你竟連那些這麼貴重的名牌都不感興趣,他們會認為你「戇居」.

在一個倒轉了的世界裡,那些直立的人,是被認為是倒轉的.

在一個顛倒是非的世界裡,那些正直的人,是被認為是頭腦不正常的.

在一個沒有理想的世界裡,堅持理想的人,是被認為是過份天真的.

如此,當主耶穌來的時候,很難避免人家說他是癲狂的.

今日講道的經文,是今日的福音書經課:馬可福音3:20－35.

3:20　耶穌進了一個屋子,眾人又聚集,甚至他連飯也顧不得吃.

3:21　耶穌的親屬聽見,就出來要拉住他,因為他們說他癲狂了.

3:22　從耶路撒冷下來的文士說:「他是被別西卜附著」；又說:「他是靠著鬼王趕鬼.」

3:23　耶穌叫他們來,用比喻對他們說:「撒旦怎能趕出撒旦呢?

3:24　若一國自相分爭,那國就站立不住；

3:25　若一家自相分爭,那家就站立不住.

3:26　若撒旦自相攻打紛爭,他就站立不住,必要滅亡.

3:27　沒有人能進壯士家裡,搶奪他的家具；必先捆住那壯士,才可以搶奪他的家.

3:28　我實在告訴你們,世人一切的罪和一切褻瀆的話都可得赦免；

3:29　凡褻瀆聖靈的,卻永不得赦免,乃要擔當永遠的罪.」

3:30　這話是因為他們說:「他是被污鬼附著的.」

3:31　當下,耶穌的母親和弟兄來,站在外邊,打發人去叫他.

3:32　有許多人在耶穌周圍坐著,他們就告訴他說:「看哪,你的母親和你弟兄在外邊找你.」

3:33　耶穌回答說:「誰是我的母親?誰是我的弟兄?」

3:34　就四面觀看那周圍坐著的人,說:「看哪,我的母親,我的弟兄.

3:35　凡遵行神旨意的人就是我的弟兄姊妹和母親了.」

這段經文有一明顯的「三明治」結構.

這段經文的開首3:20－21提到耶穌的親屬,他們出來要拉住耶穌,因為「他們說他癲狂了」.

這段經文的結尾3:31－35提到耶穌的家人,他們站在外邊,「打發人去叫他」.

經文的頭尾,點出了經文的主題.明顯,這段經文是討論「誰是耶穌真正的親屬和家人」.在這段經文最後一節裡,主耶穌說:「凡遵行神旨意的人就是我的弟兄姊妹和母親了.」那些親屬,那些家人,以為自己與耶穌有關的,卻在耶穌眼裡,是無關的人.真正與耶穌有關的,是那些遵行神旨意的人.

「三明治」中間的「餡」,令整塊三明治有特色,有味道.這段經文的中心部分,點出了「誰是那些遵行神旨意的人」.

有人認為耶穌癲狂了.

他與社會邊緣人士一起吃飯,接納妓女,接觸痲瘋病人.他是癲狂的.

若在今日,他會與露宿者一起露宿,請執紙皮維生的老婆婆吃一餐好飯,接觸愛滋病患者.他是癲狂的.

耶路撒冷的正派人士說,耶穌是被鬼附的,他著了魔.用我們的話說,「他撞邪」.若不是「撞邪」,他怎會這樣「黐線」.

在一個倒轉了的世界裡,那些憑神的靈行事的人,是會被認為是「撞邪」的.

那些人說,耶穌是被鬼王充滿的,他能趕鬼,是因為鬼王的能力.

耶穌的回應很長,我將之分成三段.

首先,耶穌說:「撒旦怎能趕出撒旦呢?若一國自相分爭,那國就站立不住；若一家自相分爭,那家就站立不住.若撒旦自相攻打紛爭,他就站立不住,必要滅亡.」

耶穌指出,若撒旦的領域「鬼打鬼」,則必會瓦解.現在,撒旦的領域沒有瓦解,也就說明,耶穌不是被鬼王附體,也就沒有「鬼打鬼」了.

若撒旦的領域沒有瓦解,耶穌用來辯論的論證,是言之成理的.但問題是,當福音書記載耶穌趕鬼時,正正是意味著撒旦的國度正在瓦解.

事實上,耶穌回應的第二段落,正正指出撒旦的國度正在瓦解.

耶穌說:「沒有人能進壯士家裡,搶奪他的家具；必先捆住那壯士,才可以搶奪他的家.」

很明顯,耶穌是指出,他是比撒旦更強的強者.他能硬闖撒旦的領域,捆住撒旦,瓦解他的國度.

若果撒旦的國度正在瓦解,問題又回來了,這是撒旦領域內裡的「鬼打鬼」嗎?耶穌不是剛剛說過,若「鬼打鬼」,則這國度便瓦解.而這國度現在真的在瓦解中,那麼,這不就是「鬼打鬼」嗎?

究竟,耶穌是以什麼力量趕鬼?

當耶穌受洗時,聖靈降在他的身上.

在旁人看來,這聖靈要耶穌做的,有點反常.

這聖靈引導耶穌去到曠野.在馬可福音的用字裡,聖靈其實是迫令,甚至趕逐耶穌到曠野(參可1:12: ``ekballō" 「催」的原文).在那裡,耶穌要受撒旦的試探.將人趕到撒旦那裡去的,怎會是「神聖」的靈呢?你想想,若有一種力量不是將你帶離試探而是將你陷入試探裡,你會認為這是聖靈嗎?

在這聖靈的力量下,耶穌接觸的人,都是社會上的邊緣人,政治上的受害者,宗教上的不潔者.在一般人的眼中,這怎會是聖靈的成為呢?聖靈不是引導人遠離不潔,在社會上步步高升,在政治上獲得權力的嗎?

聖靈一步一步引領耶穌走向死地.在十架上,馬可福音記載耶穌呼喊著:「我的神,我的神,為什麼離棄我?」一個靈若引領人走向絕望,使人感到神已離棄了自己的,這個靈能算是「神聖的靈」嗎?

文士的疑問是:究竟,耶穌是以什麼力量趕鬼?耶穌身上的靈,不是像邪靈一樣,迫使耶穌變得癲狂嗎?這個靈,令耶穌癲狂了.

耶穌回應文士的第三段話是這樣的:「我實在告訴你們,世人一切的罪和一切褻瀆的話都可得赦免；凡褻瀆聖靈的,卻永不得赦免,乃要擔當永遠的罪.」這話是因為他們說:「他是被污鬼附著的.」

所謂「褻瀆聖靈」,在這裡指認不出耶穌生命裡的靈是「聖靈」,反認之為「邪靈」.這個人不單不認識聖靈,也不認識耶穌,也不認識差耶穌來的天父.他拒絕神,怎能得到赦免呢?

誰是耶穌真正的親屬和家人呢?

是那些遵行神旨意的人.

誰是遵行神旨意的人呢?

就是那些承認推動耶穌生命的那個靈,是聖靈的人.

面對耶穌,有人說:「他癲狂了」.文士說:「他著了魔,撞邪」.

究竟,是耶穌癲狂了,抑或是世人癲狂了?

電影＜十月圍城＞講到,1906年10月,孫中山來香港,與同盟會的13省代表會面,籌劃各地的革命起義.

支持革命黨的陳少白對商人李玉堂說:「革命快開展了,整個中國會地動山搖.這麼偉大的行動,發源地就在我們這裡,而你就是其中一分子.」

後來,李玉堂發現自己的獨生子李重光上街示威,便去找陳少白理論,希望陳少白不要將李重光帶進革命事業裡.李玉堂對陳少白說:「重光是我兒子,我就這麼一個兒子.清廷派了殺手.我是商人,我兒子將來也是商人.」

陳少白說:「你是革命黨,從你為革命捐出第一分錢的那一天起,你已是革命黨.」

李玉堂回答說:「你瘋了.」

從李玉堂的角度看,陳少白那種不要命而要去改變國家的心靈,是「撞邪」,是癲狂了.

後來,警司史密斯要查封李玉堂的報社.史密斯對李玉堂說:「清政府不惜任何代價,要除掉孫文,我的英國老板已經下令,不許插手,後天的事誰去誰死,你可別忘了你是生意人.」

李玉堂回答說:「你也別忘了,你是中國人.」

史密斯對李玉堂說:「李玉堂,你瘋了.」

一個說人「瘋了」的人,如今,被人說是「瘋了」.

從政治的角度,一個商人願意放棄利益而捲入一場革命,是「撞邪」,是癲狂了.

再後來,李玉堂的獨生子李重光要做孫中山的替身,坐在轎車上,轉移清廷殺手的視線.車伕阿四不斷地拒絕,但李重光不住的堅持.在阿四眼中,這位少爺瘋了.他甘願犧牲自己,成全他人.

這齣電影裡,很多人看來都是著魔了,瘋了.很多學生瘋了.李玉堂這有錢商人瘋了.他的獨生子瘋了.那車伕瘋了.很多仗義的升斗巿民瘋了.唯有有人願意被人看為癲狂,才能改變世界.

耶穌來了.人也認為他是癲狂的,是著了魔的.但是,就是這癲狂的人,改變了世界.誰認出這人生命中的靈,是神聖的靈,這人有福了,他是屬於神的.

你是一個為主癲狂的人嗎?

但願:榮耀歸於聖父,聖子,聖靈.阿們.



\newpage

\section{神的永在人的永生}
\label{sec:154}
\textbf{鄭漢文博士 - 神的永在,人的永生}
\newline
\newline
\hyperref[sec:153]{< < < PREV SERMON < < <}
~
\hyperref[sec:toc]{[返目錄]}
~
\hyperref[sec:155]{> > > NEXT SERMON > > >}


約翰福音 3:1-17
\newline
\begin{longtable}{cl}
\hline
\hline
章節 & 經文 (和合本修訂版)\\
\hline
3:1 & \begin{tabularx}{0.7\textwidth}{X} 有一個法利賽人，名叫尼哥德慕，是猶太人的官。 \end{tabularx} \\ \\ \relax
3:2 & \begin{tabularx}{0.7\textwidth}{X} 這人夜裡來見耶穌，對他說：「拉比，我們知道你是由神那裡來作老師的；因為你所行的神蹟，若沒有神同在，無人能行。」 \end{tabularx} \\ \\ \relax
3:3 & \begin{tabularx}{0.7\textwidth}{X} 耶穌回答他說：「我實實在在地告訴你，人若不重生，就不能見神的國。」 \end{tabularx} \\ \\ \relax
3:4 & \begin{tabularx}{0.7\textwidth}{X} 尼哥德慕對他說：「人已經老了，如何能重生呢？豈能再進母腹生出來嗎？」 \end{tabularx} \\ \\ \relax
3:5 & \begin{tabularx}{0.7\textwidth}{X} 耶穌回答：「我實實在在地告訴你，人若不是從水和聖靈生的，就不能進神的國。 \end{tabularx} \\ \\ \relax
3:6 & \begin{tabularx}{0.7\textwidth}{X} 從肉身生的就是肉身；從靈生的就是靈。 \end{tabularx} \\ \\ \relax
3:7 & \begin{tabularx}{0.7\textwidth}{X} 我說『你們必須重生』，你不要驚訝。 \end{tabularx} \\ \\ \relax
3:8 & \begin{tabularx}{0.7\textwidth}{X} 風隨著意思吹，你聽見風的聲音，卻不知道是從哪裡來，往哪裡去；凡從聖靈生的也是如此。」 \end{tabularx} \\ \\ \relax
3:9 & \begin{tabularx}{0.7\textwidth}{X} 尼哥德慕問他：「怎麼能有這些事呢？」 \end{tabularx} \\ \\ \relax
3:10 & \begin{tabularx}{0.7\textwidth}{X} 耶穌回答，對他說：「你是以色列人的老師，還不明白這些事嗎？ \end{tabularx} \\ \\ \relax
3:11 & \begin{tabularx}{0.7\textwidth}{X} 我實實在在地告訴你，我們所說的是我們知道的，我們所見證的是我們見過的，你們卻不領受我們的見證。 \end{tabularx} \\ \\ \relax
3:12 & \begin{tabularx}{0.7\textwidth}{X} 我對你們說地上的事，你們尚且不信，若對你們說天上的事，如何能信呢？ \end{tabularx} \\ \\ \relax
3:13 & \begin{tabularx}{0.7\textwidth}{X} 除了從天降下的人子，沒有人升過天。 \end{tabularx} \\ \\ \relax
3:14 & \begin{tabularx}{0.7\textwidth}{X} 摩西在曠野怎樣舉蛇，人子也必須照樣被舉起來， \end{tabularx} \\ \\ \relax
3:15 & \begin{tabularx}{0.7\textwidth}{X} 要使一切信他的人都得永生。 \end{tabularx} \\ \\ \relax
3:16 & \begin{tabularx}{0.7\textwidth}{X} 「神愛世人，甚至將他獨一的兒子賜給他們，叫一切信他的人不致滅亡，反得永生。 \end{tabularx} \\ \\ \relax
3:17 & \begin{tabularx}{0.7\textwidth}{X} 因為神差他的兒子到世上來，不是要定世人的罪，而是要使世人因他得救。 \end{tabularx} \\ \\
[1ex]
\hline
\hline
\end{longtable}
日期:2012年7月15日

主題:神的永在,人的永生

經文:約翰福音3:1-17

證道:鄭漢文博士

如果你打從火車站進入香港中文大學這個校園,當眼所見的是一尊民主女神像,紀念的是一九八九年六月四日在中國北京天安門廣場,學生死傷無數卻未能自由哀悼的民主運動.其實這校園裡師生稱之為「民女」的民主女神像當天進入校園是有一段故事的,現在很多時候我們在這裡集合──師生集合,外邊的集合這地點,比方你收到我發給你的一個短訊──「Roger call 0,六月四日1835 民女等,去維園,穿黑衣.」

回復的人會寫「民女見」.民主女神像被中文大學的學生會迎接入這校園這地方,當時引起校方一些不批准的爭議,那晚上我跟我的大兒子都在現場,目睹民女由貨車放下,放置在現在她樹立的地方,見證了一個時代的精神,從此中大多了一個言論自由的集會點.因為我住在校園裡,每一個晚上都得見街燈殘照之下的身軀在倒映,倒映著那一年的天安門前天不安.六四,二十有三年,真是「多年生死兩茫茫,不思量自難忘」!

二十三年不能止息的是一種精神,有人認為是香港人良心好過點的,每一年的哭墳,有人認為香港堅持平反六四的活動正是中國難得的自由燭光,不能熄滅.

究竟永永遠遠,生生不息是什麼?

正如聖經作者常常會這樣說「當天我所見所聞的事情,今天說給大家知道,正是要叫你們得永生.」

何謂永生?誰不求死後能夠永生,常存不朽?因為我們的身軀無法永生,人就會追求精神的長存,渴望流芳百世,最怕唯臭萬年.人究竟渴望怎樣的精神能夠永在?二十三年前的六四,歷史百年功過誰屬?誰得千古萬名,誰的精神永存.正入三國演義所描繪「江山依舊在,幾度夕陽紅」.昂山素姬因為父親昂山將軍的愛國精神得到了自身時代一息尚存的空間,在這空間活出了可敬的人生,感動世人,回應著更大的精神.昂山素姬不單被視為緬甸的國民領袖,更成為當今世代,當今世界的民主女神.我在情人節當天跟內子一塊看”The Lady昂山素姬傳奇”

這一套電影,動人的不單單是愛國精神,不單單是與民同苦的情操或是堅持民主的女神,對父母親的親愛,對不離不棄但要天國一方的夫婦感情,對子女不能近身的養育,但遠程給了以身作則的教育,只有情永在.楊子瓊女士當天拍攝這套電影的時候,昂山素姬還在家中被軟禁,到了我們去看首輪電影的時候,昂山素姬已經能走出來競選國會議員,今天昂山素姬作為國會議員已經能夠出國到其他國家訪問,她能否回國?世事難料,人面對有限便會思索無限,盼望美好能無限的永在下去.人面對什麼有限性?人的生命面對時間的有限和認識上的有限,有生有死,所以有些人求長生不老,其實長生都很難保,生命有長有短,可能有些活得長一點兒,有些活得短一點兒,各有不同,但相同的是你我的生命都有時有候,大限臨頭,人不能超越時限,此生之後會否有來世?故此有些人追求靈魂不滅,這不滅不朽所存的地方,有些人叫為「天國」,好像難以捉摸,莫明底蘊,完全超乎我們認識的界限.永生像是比長生更長,但長生只是在量上將永生演繹,永生是在質上超過量上意義的長生,永生不單單是存在死後的來生,即便說永生是將來的事情,而是存在當下,究竟你是否感覺到現在你就活在永生當中?

今日經課中約翰福音3章1-17節,記載了一個感受不到當下活在永生的生命,這個當事人叫尼哥底母.讓我們看一看聖經如何記載這一段追求超越有限進入永生的對話:

3:1 有一個法利賽人,名叫尼哥底母,是猶太人的官.

3:2 這人夜裡來見耶穌,說,拉比,我們知道你是由神那裡來作師傅的,因為你所行的神跡,若沒有神同在,無人能行.

3:3耶穌回答說,我實實在在的告訴你,人若不重生,就不能見神的國.

這故事的開始談到尼哥底母,為何這個猶太族群中非常顯赫的人要夜欄人靜時來見耶穌?會不會是因為他懼怕其他人的眼光?就算是怕,尼哥底母依然到訪耶穌,因為他心中有一個謎,需要向這位他相信是有神同在的耶穌求個明白.耶穌看明白他心裡是想追問神的國的事情,就是追求生命的最後終極,追求如何超越此生而得進永恆的神的國.

耶穌的回答給他一個新觀念,就是「重生」.但這個猶太人的視野無法站入耶穌的視野裡,兩人正是在兩個不同的層次上展開關於生命的對話:

3:4 尼哥底母說,人已經老了,如何能重生呢．豈能再進母腹生出來麼.

3:5 耶穌說,我實實在在的告訴你,人若不是從水和聖靈生的,就不能進神的國.

3:6 從肉身生的,就是肉身．從靈生的,就是靈.

3:7 我說,你們必須重生,你不要以為希奇.

3:8 風隨著意思吹,你聽見風的響聲,卻不曉得從那裡來,往那裡去．凡從聖靈生的,也是如此.

尼哥底母是以字面義去看待重生:生命的開始是在母腹中,重生豈不是要從母腹中再生出來?耶穌指出的是生命的本體性,這樣從母腹中再生出來的生命只能是肉身的生命,能進神的國的生命是要從靈裡面生的而成為靈的生命,就是本體上要改變生命的性質.這個靈性義上的重生,對人來說,是難以明白的.從聖靈生的,究竟是如何可能?

3:9 尼哥底母問祂說,怎能有這事呢?

3:10 耶穌回答說,你是以色列人的先生,還不明白這事麼.

耶穌的對話,許多時候在現代人的看法都是頗能揶揄他人的,好像在說你怎麼會不明白?然後所給的答案往往是不會跟同一線索,同一水準開展.

祂說:

3:11 我實實在在的告訴你,我們所說的,是我們知道的,我們所見證的,是我們見過的．你們卻不領受我們的見證.

3:12 我對你們說地上的事,你們尚且不信,若說天上的事,如何能信呢.

3:13 除了從天降下仍舊在天的人子,沒有人升過天.

3:14 摩西在曠野怎樣舉蛇,人子也必照樣被舉起來．

3:15 叫一切信祂的都得永生.〔或作叫一切信的人在祂裡面得永生〕

3:16 神愛世人,甚至將祂的獨生子賜給他們,叫一切信祂的,不至滅亡,反得永生.

3:17 因為神差祂的兒子降世,不是要定世人的罪,〔或作審判世人下同〕乃是要叫世人因祂得救.

以上一段生命的對話啟示出一個救恩的真理.這個真理是關乎耶穌的身分,是耶穌借尼哥底母究問神的國而透露自己與聖父和聖靈的關係.耶穌開悟尼哥底母的故事並沒有就此完結,約翰福音是一個的很大敘事體,當中包含了不少小敘事體,其中的兩段,三段加起來可以構作出尼哥底母信仰成長故事的起承轉合,借此讓我們反思我們的生命:

一,起首於那兒?約翰福音沒有記載尼哥底母如何成為法利賽人和猶太人的官的過去.值得思考的是究竟他如何成長?你又如何成長到今天?

二,承接于約翰福音3:1-17:尼哥底母夜訪耶穌,究竟一生何求?剛才我們思考過,所求的是如何超越有限而得進入神的國.你又一生何求?

三,轉至約翰福音7:45-52第三個段落:尼哥底母為耶穌公開抗辯,大家看看經文便知道,耶穌開始被人指控的時候,尼哥底母願意在眾人前,在自己群體面前公開自己的信仰,公開自己為耶穌抗辯.當下為何他有這般勇氣為自己的信仰作出如此的事情?你信耶穌後有否在某些時刻你要鼓起勇氣做一些事情?

四,最後于約翰福音19:39-42:耶穌死了,尼哥底母為耶穌死後裹屍,他以身犯險,可會因耶穌的死而受死.若真的要死,你又會死而無憾呢?有甚麼在生命終結前你要做的以致你死而無悔呢?

耶穌開悟尼哥底母的故事同時啟示了一個奧秘:「從聖靈生的重生,就是相信這個從聖父所差來的聖子,因此而得永生.」

各位,這一句說話表白了我們所信仰的是一位三一的神.現在讓我們一同思考這句話的意思,以及對我們生命的意義.耶穌開悟尼哥底母的同時,祂有意地啟示了一個關於神的奧秘,後世神學家稱之為三一神觀(Trinity / Trinitarian God).我們基督徒所信仰的神是一個三位一體的神.

是一個神,但有三個位格.例如我們的榮耀頌是這樣唱頌的:「但願榮耀歸於聖父,聖子,聖靈,起初這樣,現在這樣,以後也這樣,永無窮盡,阿們.」「三位一體」這個奧秘我們難於明白,但可以理解.今天我們不會進入系統神學的討論,因為可能只會越搞越難明白.這三一神的信仰其實已經發展成使徒信經,將來有機會再跟大家研討這信經,但是,我們可以這樣理解這個奧秘:聖父,聖子,聖靈是一種關係上的契合.我們所信的神是講究關係.這個信仰是讓我們有三種方式與神相連,重新發生關係.讓我們從這種理解重新思考永生和重生的意義.

永生的相反是甚麼?最極致的就是永死,永死是任何時間都與神隔絕；重生就是重新建立關係.

首先,是與創造主聖父如何再聯繫上?以致能夠復原持續創造的生命.這是上帝按著祂的形象創造我們的生命.其次,是與救贖主聖子如何再聯繫上?以致因著十架救恩而恩上加恩.剛才的經文,耶穌揭露祂的身份,就是祂要被舉起以致我們不至滅亡,反得永生最後,是與保惠師聖靈如何再聯繫上?以致聖靈充滿而得以更新成聖.聖靈降臨了,其實聖靈跟我們關係非常密切,以致耶穌回答尼哥底母的時候提醒他,看不見的聖靈,其實是他生命中重生必須的開始.在有生之年,會死的人與永在的神重新建立關係,不單不致滅亡,而且進入永生.我們所說的進入,像是空間觀念,這空間像是在未來──與神的國有分,但永生還有一個時間的意義,是在當下時間與永生有關──與神的靈有分,通過神的靈與神發生關係,因為永在的神就在當下與你同在.神的永在注入人的今生,以致創造生命的聖父,讓我們今天就成為上帝子民在此崇拜天父.救贖生命的聖子,讓我們今天就成為耶穌基督的門徒,一瞬間一同領受聖餐.更新生命的聖靈,讓我們今天就成為靈裡重生的心靈,時時在靈修.以致剛才可能聖靈在你的生命裡已感動你,所以單單剛才十點半到十二點崇拜完結,神的永在不斷灌注入你當下的今生.

神是三位一體的神,是創造,救贖和更新的永在,參與其中而有分的就是活在永生中.人的永生有賴神的永在,永在的神以聖父,聖子,聖靈的存在方式讓人參透永生.與神的永在發生生命的關係,就是人得永生的核心,就是重生的開始.重生就是生命不再如此活,而是活出另一種生命.是否一頭霧水?

總結:我和大家分享的訊息是以下這二十多個字:人在靈裡重生,然後進入永生(這是耶穌回答尼哥底母那段說話的意思)永在的神與你我同在,活出基督精神,當下就是活出永生.讓我們再讀出3章16節:「神愛世人,甚至將他的獨生子賜給他們,叫一切信祂的,不至滅亡,反得永生.」

3章16節是我們常常用作佈道的金句,有被視為耶穌福音的摘要.這段摘要所出現的上下文,剛才我們曾經思考,正是耶穌在透露自己的身分,祂是愛世人的聖父所差來的聖子並留下聖靈讓我們得以重生；回應的是盼望得進神的國的生命關懷. 你關懷你自己的生命嗎?正如希臘哲學家亞裡斯多德(Aristotle)說,沒有人不追求美好人生的,每個生命都渴望活得更好.分別只在於:甚麼是生命的最好?

為了鼓勵我今天第一次講大學講章,其實1980年我進入了中文大學在崇基念大學後,從來沒有機會在這裡講話,謝謝牧師,謝謝他給我機會跟大家分享,但是非常之緊張,所以昨晚兩個兒子,內子和我一家四口在家共聚,一起齊吃刺身壽司──預祝一下,刺身壽司不多但我們依然快樂.曾經嘗過美好,真想再渡好時光.但是再多的刺身也留不過夜,共渡好時光卻成了人生的幸福回憶和美好盼望.

一對熱戀情人卻不在乎吃甚麼,所謂有情飲水飽,情到濃時最不願的是分離.關於分離,柳永這詩人曾寫詩說:「兩情若是久長時,又豈在朝朝暮暮.」其實詩人是一種自我安慰.因為心底裡依然渴慕久長的愛情.時時能相見,因為求之不得,唯有珍惜當下.

同在有多好.與相愛的人同在有多好,何況是與愛你的神同在,愛你的神,這永在的神今天與你同在,這不是好得無比,而是好得無法相比?

今天大家一同到臨崇拜,但願從今開始,這永在的神與你同在,直到永遠.

我們重生後,神的同在在我們生命並不止息,永永遠遠流動的是基督精神的永在我們生命間.

基督因是神的兒子,活出神的大能與神的大恩,受苦,復活,讓世人得識基督精神背後的神.

六四天安門之子,受苦受死,點燃愛國民主的精神,喚發二十三年不息的民主精神.因著他們肉身的消磨,我們認識到他們背後代表的精神.同樣基督因為祂道成肉身,肉身受苦受死,但復活以致我們見到我們看不見的神.

基督徒也愛國,首先愛天國,然後天國降臨,活出永恆.基督徒也追求自由民主,因為相信每個人都是按神的形像所造的,在合理的制度下讓每個生命自由發聲,就是聆聽上主的聲音.民中有主,主在民中.作為聖子基督的門徒,如何可能活出聖父的天國精神?因為有聖靈,能夠教我們適時感應神的同在,當下活出基督精神.

甚麼是基督精神?比如:在失信之時,堅守信德；在絕望之境,點燃盼望.在乏愛之處,流轉仁愛；當我們在六四晚會,為一位身傍將殘就滅的燭光,因著基督精神,重燃生命的信望愛時,我們都在活在永生,活出永生.

但願榮耀都歸給聖父,聖子,聖靈,永在的三一神.阿們.



\newpage

\section{哭泣與歡呼流淚的禱告}
\label{sec:155}
\textbf{賀志勇博士 - 哭泣與歡呼}
\newline
\newline
\hyperref[sec:154]{< < < PREV SERMON < < <}
~
\hyperref[sec:toc]{[返目錄]}
~
\hyperref[sec:156]{> > > NEXT SERMON > > >}


詩篇 30:1-12
\newline
\begin{longtable}{cl}
\hline
\hline
章節 & 經文 (和合本修訂版)\\
\hline
30:1 & \begin{tabularx}{0.7\textwidth}{X} 耶和華啊，我要尊崇你，因為你救了我，不讓仇敵向我誇耀。 \end{tabularx} \\ \\ \relax
30:2 & \begin{tabularx}{0.7\textwidth}{X} 耶和華－我的神啊，我呼求你，你醫治了我。 \end{tabularx} \\ \\ \relax
30:3 & \begin{tabularx}{0.7\textwidth}{X} 耶和華啊，你救我的性命脫離陰間，使我存活，不至於下入地府。 \end{tabularx} \\ \\ \relax
30:4 & \begin{tabularx}{0.7\textwidth}{X} 耶和華的聖民哪，你們要歌頌他，要頌揚他神聖的名字。 \end{tabularx} \\ \\ \relax
30:5 & \begin{tabularx}{0.7\textwidth}{X} 因為，他的怒氣不過是轉眼之間，他的恩典乃是一生之久；一宿雖然有哭泣，早晨便必歡呼。 \end{tabularx} \\ \\ \relax
30:6 & \begin{tabularx}{0.7\textwidth}{X} 至於我，我凡事順利，就說：「我永不動搖。」 \end{tabularx} \\ \\ \relax
30:7 & \begin{tabularx}{0.7\textwidth}{X} 耶和華啊，你曾施恩，使我穩固如山；你轉臉不顧，我就驚惶。 \end{tabularx} \\ \\ \relax
30:8 & \begin{tabularx}{0.7\textwidth}{X} 耶和華啊，我曾求告你；我向耶和華懇求： \end{tabularx} \\ \\ \relax
30:9 & \begin{tabularx}{0.7\textwidth}{X} 「我被害流血，下到地府，有何益處呢？塵土豈能稱謝你、傳揚你的信實嗎？ \end{tabularx} \\ \\ \relax
30:10 & \begin{tabularx}{0.7\textwidth}{X} 耶和華啊，求你應允我，憐憫我！耶和華啊，求你幫助我！」 \end{tabularx} \\ \\ \relax
30:11 & \begin{tabularx}{0.7\textwidth}{X} 你將我的哀哭變為跳舞，脫去我的麻衣，為我披上喜樂， \end{tabularx} \\ \\ \relax
30:12 & \begin{tabularx}{0.7\textwidth}{X} 使我的靈歌頌你，不致緘默。耶和華－我的神啊，我要稱謝你，直到永遠！ \end{tabularx} \\ \\
[1ex]
\hline
\hline
\end{longtable}
日期:2012年7月22日

主題:哭泣與歡呼:流淚的禱告

經文:詩篇30:1-12

證道:賀志勇博士

前言

我在杭州認識一個朋友,他是溫州人,他家好幾代都是基督徒.很不幸的是,前年他母親得了胰腺癌.有一次我在杭州講課,他來找我,他說他心裡有種擔憂,擔心他母親上不了天堂.我大吃一驚,問他為什麼?他說他母親是名虔誠的基督徒,剛開始知道自己得了癌症,他們那邊認為癌症是撒旦在作怪,所以請傳道人天天禱告趕鬼.他母親跟著傳道人天天禱告,天天懇求天父趕鬼治病；可是禱告好像無效,他母親身體一天差過一天,終於有一天,他母親就開始埋怨上帝了.他母親覺得上帝很不公平；她這樣愛主,她才四十多歲,主為什麼不恩待她,為什麼要讓她吃這些苦,受這些罪.他母親很生氣,連話都不願多說,也不禱告了.我的朋友問我:"是不是埋怨上帝就沒有信心了,就上不了天堂了呢?" 我說:"不一定啊,大衛在詩篇裡不是有很多篇都是哀呼,抱怨的禱告麼,詩篇三十五篇,詩篇八十六篇,大衛抱怨和懇求的語氣都非常強硬啊,神不是照樣拯救了他麼!" 我正想跟他讀讀詩篇三十五篇安慰他,安慰一下他母親,但腦海裡,卻突然想到這一篇:詩篇三十篇.詩篇三十篇這樣寫到:

神的恩典乃一生之久

1, 耶和華啊,我要尊崇你,因為你曾提拔我,不叫仇敵向我誇耀.

2, 耶和華我的神啊,我曾呼求你,你醫治了我.

3, 耶和華啊,你曾把我的靈魂從陰間救上來,使我存活,不至於下坑.

4, 耶和華的聖民哪,你們要歌頌他,稱讚他可記念的聖名.

5, 因為他的怒氣不過是轉眼之間；他的恩典乃是一生之久.一宿雖然有哭泣,早晨便必歡呼.

6, 至於我,我凡事平順,便說:"我永不動搖."

7, 耶和華啊,你曾施恩,叫我的江山穩固；你掩了面,我就驚惶.

8, 耶和華啊,我曾求告你.我向耶和華懇求說:

9, ``我被害流血,下到坑中,有甚麼益處呢?塵土豈能稱讚你,傳說你的誠實嗎?

10, 耶和華啊,求你應允我,憐恤我!耶和華啊,求你幫助我!"

11, 你已將我的哀哭變為跳舞,將我的麻衣脫去,給我披上喜樂.

12, 好叫我的靈(註:原文作"榮耀")歌頌你,並不住聲.耶和華我的神啊,我要稱謝你,直到永遠!

I will exalt you, Lord, for you lifted me out of the depths and did not let my enemies gloat over me.2 Lord my God, I called to you for help, and you healed me.3 You, Lord, brought me up from the realm of the dead; you spared me from going down to the pit.

4 Sing the praises of the Lord, you his faithful people; praise his holy name.5 For his anger lasts only a moment, but his favor lasts a lifetime;weeping may stay for the night, but rejoicing comes in the morning.

6 When I felt secure, I said, ``I will never be shaken."7 Lord, when you favored me, you made my royal mountain[c] stand firm;but when you hid your face, I was dismayed.

8 To you, Lord, I called; to the Lord I cried for mercy:9 ``What is gained if I am silenced, if I go down to the pit?Will the dust praise you? Will it proclaim your faithfulness?10 Hear, Lord, and be merciful to me; Lord, be my help. ``

11 You turned my wailing into dancing; you removed my sackcloth and clothed me with joy,12 that my heart may sing your praises and not be silent. Lord my God, I will praise you forever.

某些時候,用詩篇三十篇這樣的感恩詩來安慰病重的人,好像並不是很切實際.因為擺明瞭他現在病得很重,並沒有得到醫治,又哪里會有心情感謝神拯救了他呢?但是很奇怪,在我朋友說到他母親的時候,我突然想到了這首詩；之所以會想到它,是因為這首詩不僅僅是一首讚美詩,不僅僅是大衛感謝神醫治、拯救的讚美詩,它同時也非常真實地表達了大衛經歷大災大能時的的心境,還有大衛信心的掙紮過程.這是一首充滿感情的詩篇.我想透過分析這首詩,來看看大衛的信心的心路歷程.

解經一, 蒙拯救和"流淚的禱告"(三十:1-5)

在1到5節,詩人不僅僅自己感謝神,還帶領大家一起感謝.在第一節,大衛感謝神提拔了他.提拔原文是一個非常生動的字,就好像打水,把水桶從深井裡頭拉起來.大衛感謝神,因為神就像打水一樣,把他從深井里拉起來.這個深井可不是深井燒鵝的那個深井,這個深井裡沒有燒鵝,只有仇敵--大衛命運中的仇敵.這裡提到的「仇敵」是什麼意思呢?原文這個字是一個眾數(enemies)字,用普通話講就是一個複數字,包括後面提到的「疾病」(v.2)陰間(v.3)被害流血(v.9)下坑(v.3,9),或者像唐佑之博士講的那樣,仇敵甚至還包括當你遭難時,那些看你笑話、不同情你的人.

疾病是我們的仇敵!香港人經常講:"得閒死毋得閒病."(有空去死沒空去病)對於我們這些天天忙碌的人來講,生病很難受.火車上的廣告,有將近三分之一是跟健康有關的.但我發現,人一過三十歲,病就一個接一個找上門來了.今天脖子酸,剛剛痊癒,背又痛了；背剛剛好,腸胃又不好,腸胃剛剛好,皮膚又不好；皮膚剛剛好了,又血壓血糖偏高.我以為生病是我這種三十歲以上的人的專利,結果發現不是,現在很多本科生也經常肩膀酸痛,背部酸痛了.我們本來很美好的生活,就被疾病這些小麻煩困擾著；就好像你要睡覺,蚊子不停在你耳邊飛一樣,麻煩得很.這還是小病,最近這幾年癌症增多,幾乎隔幾天就會聽到癌症這個詞.

在面對疾病時,在面對死亡的威脅時,大衛懂得呼求神,我們會不會這樣做呢?我發現我自己生病時多數是求告醫生,而不是求告神.我一有病就去求告醫生,然後吃藥；當然看的多數是中醫,吃的多數是中藥.因為中藥據說更健康.弟兄姊妹,你吃藥之前會不會禱告呢?我很少會,我想也沒幾個人吃藥之前會先禱告一下.

我們中間有沒有人吃藥之前會禱告一下?你說吃藥之前禱告做什麼?!難道跟神說,"神啊,求你讓藥一定有效果把!"我們通常不會這樣禱告,因為我們相信,醫生總可以幫到忙的,藥總會有效果的,起碼會有一定效果.我們是單單懷著對醫生、對藥的信心去看病的；當我們覺得人的手可以搞定這些事情,我們覺得靠自己可以搞定事情的時候,我們就不會向神求告.小病,我們覺得靠醫生、靠自己就能治好,所以不禱告；大病,我們覺得醫生都不管用,神又能怎樣呢?所以我們也不禱告.但神其實不喜歡這種態度.列王記下講,亞哈謝王生病了,他去求問所謂的"神醫"巴力西蔔,而不是去求告耶和華神,神就讓以利亞去宣佈"你必不下你所上的床,你必定死."

大衛王不像我們這樣,你看詩篇30篇,他就因為疾病向耶和華禱告；其實你看看詩篇,幾乎篇篇都有向神懇求的地方.聖經並不是認為治病可以不需要用藥,在好撒瑪利亞人的故事中,就是用橄欖油治病.馬可福音6:13節,雅各書5:14節也提到用橄欖油治病.所以我猜想,大衛並不是不看醫生不吃藥；大衛也看病也吃藥,但是大衛明白,無論醫生也好,無論醫藥也好,如果神不讓那些東西起作用,就一定起不了作用.

所以從大衛的詩篇裡頭,我們首先看到的是一種態度:一種懇求的態度,一種敬畏神的態度.這種態度,就是禱告的態度."凡事謝恩,凡事禱告",並不是要你為了今天晚上吃白菜還是吃蘿蔔禱告,也不是要你出門時為了穿這雙鞋還是那雙鞋禱告.凡事禱告乃是讓你學習一種態度,一種凡事都要敬畏神的態度.哪怕我們還在困境中,哪怕神還沒有救我們脫離苦海,我們仍舊要記得敬畏神、向神呼求,這就是禱告的態度.就好像詩篇35:16-17節,大衛大聲向神說:"我在患難中,他們卻歡喜,......主啊,你看著不理要到幾時呢?"

哪怕神不理大衛,大衛仍舊懂得要向神禱告.就好像今天讀的詩篇30篇,大衛經歷了非常多、非常大的苦難,大衛也會害怕、也會流淚,也會大聲抱怨,因為神並不是事事都提拔他,不是處處都醫治他,但無論什麼樣的光景,他都會向神呼求,向神感謝,因為他敬畏神.這樣的禱告,我稱它為"流淚的禱告".

二,愚人的自誇和"流淚的禱告"(三十:6-10)

弟兄姊妹,很多時候我們為什麼不去禱告呢?是不是因為我們還不懂得敬畏神,還不知道自己的危險呢?但詩篇30篇--也就是"流淚的禱告"--提醒我們,不要做個愚蠢人.什麼是愚蠢人?6到10節講,當遭受到仇敵和疾病攻擊之前,大衛的生活非常平順.平順這個字可以翻譯成安逸.那個時候,詩人有錢、有享樂、有才智,又有健康,日子過得非常安逸.當時他就在心裡說:這都是出於我自己的本領.我可以靠著自己的力量,永遠不動搖,江山也穩固--現在很多事業成功、身體也還算健康的人都是這種心態.誰知道,正在這個自滿自大的時候,戰爭、魔鬼和疾病,接踵而來,都來襲擊他.詩人這才發現,自己的生命原來十分危險.他十分驚惶害怕.當他發現自己是個愚蠢人時,他就開始流淚禱告,流淚求告耶和華了.

前些日子,Hardy跟我談到胡雪岩.(Hu Xueyan.) 胡雪岩是晚清江南第一首富,在杭州有所大宅子,但只住了八年就不得不搬出去了,破產了.胡雪岩窮苦人起家,有錢之後花錢無度,終於弄到抄家沒產.沒有敬畏的人生實在太多跌宕起伏了.(這個可以不用)

不過在9-10節,詩人的禱告並沒有一下子全然認錯,這時候他還有點情緒,9到10節,他向耶和華懇求說:9, ``我被害流血,下到坑中,有甚麼益處呢?塵土豈能稱讚你,傳說你的誠實嗎?10, 耶和華啊,求你應允我,憐恤我!耶和華啊,求你幫助我!

你說這怎麼可以,怎麼可以同神講數,怎麼可以同神討價還價呢?怎麼可以一邊懇求神的憐憫,又一邊抱怨神呢?這不是一個虔誠的禱告,這不是一個合格的禱告!

我們認為禱告要非常虔誠,要說非常非常順服神、非常非常愛神的話,要非常非常屬靈.有的人自己一個人禱告的時候,可能還會跟神說那麼一點心裡話,但是跟弟兄姊妹一起禱告,就變成了另外一個人.如果你跟他說,禱告可以跟神抱怨,跟神討價還價,跟神開點小玩笑,他眼睛就會瞪得大大的說:"這怎麼可以!這太不虔誠了!impeity" 這種信心有點像猶太人的信心,猶太人相信神只是威嚴遙遠、殺氣騰騰的神,我們表現好了,他就獎賞我們；表現不好,就懲罰我們.所以禱告的時候,我們只能報喜不報憂.神不是慈愛的父親、友善的朋友、溫柔的安慰者,神只是高高在上的王.那就是了,對人還不能多說一句話,說得越多錯得越多,更何況是對神說話了.還有的人信心剛硬,覺得輸贏都該自己負責.成功是因為個人的能力,失敗是因為個人的運氣.人要靠自己,別奢望神來幫你.

但大衛不是這樣,他面對神的時候就像個孩子,想到什麼就說什麼,想求什麼就求什麼.雖然有時候他因為凡事平順,讓他暫時遠離了神,但很快他就會知道自己錯了,就會悔改,重新歸向神.因為他知道他跟神的關係,正如他在詩篇二十三篇所說:"耶和華是我的牧者,我必不至缺乏.......在我的敵人面前,你為我擺設筵席."大衛知道,神是看顧牧養他的牧者,是他天上的父親,所以雖然他會抱怨神、會討價還價,也會生氣發怒,但他在神面前不會隱藏自己.誰會跟自己的父親說客氣話呢?大衛雖然軟弱過,愚蠢過,但他浪子回頭時,他懂得他要求告神.這是出於什麼?這是出於信心.我們之所以有"流淚的禱告",是出於我們對神的信心.

三,歡慶和"流淚的禱告"(三十:11,12)

弟兄姊妹,我們怎樣才能不做愚蠢人呢?詩篇30篇告訴我們:敬畏神,用真正的信心去禱告.那什麼是真正的信心呢?有一次,有一個弟兄跟我說:"我總是不確定神是不是在看顧我、拯救我.有時候我甚至在想祂一定存在麼!"我問他:"那你還查經,還靈修,還會跟神說話麼?"他說"會".我就告訴他:"信心就是雖然你懷疑、你害怕、你抱怨,但是你還是不會背棄神,你還是願意沿著主耶穌的道路繼續走下去."有的人以為信心就是毫無疑惑,毫無情緒,確確實實抓在手裡毫不動搖.但一件事情,你已經確實抓在手裡了,那就不需要靠信心了,信心就是"所望之事的實底, 是未見之事的確據." Now faith is being sure of what we hope for and certain of what we do not see當我們有許多懷疑、害怕、抱怨,我們還願意求告神,還願意向神禱告、向神感謝的時候,我們就經歷了信心.信心乃是一種經歷."流淚的禱告"就是活生生地經歷了信心.

大衛就是這樣,他大病剛好,剛剛解決掉各種仇敵,甚至是那些在他困苦時還無情譏笑他的人,他經歷大災大能,他的信心其實還在非常矛盾掙紮之中.人都是這樣,剛剛經歷了跟親人的生離死別、經歷了大災大能,你的心情不可能像剛剛從海洋公園或者麥當勞回來一樣,你肯定有很多的情緒掙紮在裡頭.大衛的心情也是這樣!但就在這樣矛盾掙紮的時刻,大衛開始向神承認錯誤,向神獻上感謝,當然,也沒有忘記向神抱怨兩句.他不掩飾他做人的軟弱情緒,他也不矯情地顯示他的虔誠,因為他知道:"耶和華啊.我舌頭上的話,你沒有一句不知道."當大衛暫時放下這一切,他就得了自由,就可以放聲讚美耶和華了.

``流淚的禱告"就是這種懷著種種情緒、種種掙紮的禱告.有信心的禱告,就是哪怕我們還很疑惑、很軟弱,我們還能繼續往前走的禱告.弟兄姊妹,在我們的生活中,無論是疾病,或者是教育孩子的問題,或者是跟他人相處的問題等等,我們能不能學習大衛,放下自己的掙紮,努力讚美神、努力向神禱告呢?哪怕你覺得神並沒有聽你的禱告,也請繼續向神說話,因為"耶和華啊.我舌頭上的話,你沒有一句不知道."

結論:一宿雖然有哭泣,早晨便必歡呼

當我再次見到我那個朋友時,已是他母親的葬禮.他告訴我說,在他母親臨終之前的某天,母親對他講:她願意接受這一切,她相信這都是神的安排."賞賜的是耶和華,收取的也是耶和華." 然後她母親靜靜為自己的死亡預備.聽他說完,我心裡感慨萬千.一方面為他母親過世而難過,另一方面也很感動.因為他母親終於學會了不再向神問"為什麼?為什麼?",她開始學會向神問"神啊,你究竟要我怎麼做?"這份信心來之不易!在面對死亡的威脅時,他母親學到了這樣一個功課:"他的怒氣不過是轉眼之間；他的恩典乃是一生之久."

弟兄姊妹,人生多風雨,但是"晚上雖然有哭泣,早晨便必歡呼."還記不記得在某個逾越節的晚上,主耶穌被毒打、侮辱,受盡苦難；但在某個主日的清晨,他宣佈他戰勝了死亡,他從死裡復活了.這還不夠有恩典,這還不夠有盼望麼?讓我們學會在苦難掙紮中繼續禱告,哪怕這禱告有疑惑、抱怨,都請我們繼續去禱告、懇求那位帶我們出黑暗入光明的奇妙者,因為"一宿雖然有哭泣,早晨便必歡呼."感謝神,感謝神賜下的盼望.



\newpage

\section{一根刺的恩典}
\label{sec:156}
\textbf{林豪恩先生 - 一根刺的恩典}
\newline
\newline
\hyperref[sec:155]{< < < PREV SERMON < < <}
~
\hyperref[sec:toc]{[返目錄]}
~
\hyperref[sec:157]{> > > NEXT SERMON > > >}


哥林多後書 12:1-10
\newline
\begin{longtable}{cl}
\hline
\hline
章節 & 經文 (和合本修訂版)\\
\hline
12:1 & \begin{tabularx}{0.7\textwidth}{X} 雖然自誇無益，我還是不得不誇。我現在要提到主的異象和啟示。 \end{tabularx} \\ \\ \relax
12:2 & \begin{tabularx}{0.7\textwidth}{X} 我認識一個在基督裡的人，他在十四年前被提到第三層天上去；或在身內，我不知道，或在身外，我也不知道，只有神知道。 \end{tabularx} \\ \\ \relax
12:3 & \begin{tabularx}{0.7\textwidth}{X} 我認識的這樣的一個人—或在身內，或在身外，我都不知道，只有神知道— \end{tabularx} \\ \\ \relax
12:4 & \begin{tabularx}{0.7\textwidth}{X} 他被提到樂園裡，聽見隱祕的言語，是人不可說的。 \end{tabularx} \\ \\ \relax
12:5 & \begin{tabularx}{0.7\textwidth}{X} 為這人，我要誇口；但是為我自己，除了我的軟弱以外，我並不誇口。 \end{tabularx} \\ \\ \relax
12:6 & \begin{tabularx}{0.7\textwidth}{X} 就是我願意誇口也不算狂，因為我會說實話；只是我絕口不談，恐怕有人把我看得太高了，過於他在我身上所看見所聽見的； \end{tabularx} \\ \\ \relax
12:7 & \begin{tabularx}{0.7\textwidth}{X} 又恐怕我因所得的啟示太高深，就過於高抬自己，所以有一根刺加在我身上，就是撒但的差役來折磨我，免得我過於高抬自己。 \end{tabularx} \\ \\ \relax
12:8 & \begin{tabularx}{0.7\textwidth}{X} 為了這事，我曾三次求主使這根刺離開我。 \end{tabularx} \\ \\ \relax
12:9 & \begin{tabularx}{0.7\textwidth}{X} 他對我說：「我的恩典是夠你用的，因為我的能力是在人的軟弱上顯得完全。」所以，我更喜歡誇耀自己的軟弱，好使基督的能力覆庇我。 \end{tabularx} \\ \\ \relax
12:10 & \begin{tabularx}{0.7\textwidth}{X} 為基督的緣故，我以軟弱、凌辱、艱難、迫害、困苦為可喜樂的事；因為我甚麼時候軟弱，甚麼時候就剛強了。 \end{tabularx} \\ \\
[1ex]
\hline
\hline
\end{longtable}
日期:2012年7月29日

題目:一根刺的恩典

經文:哥林多後書12章1-10節

證道:林豪恩先生

一．神人身上也有一根困擾的刺

前言

12 : 1雖然自誇無益,我還是不得不誇.我現在要提到主的異象和啟示.12 : 2我認識一個在基督裡的人,他在十四年前被提到第三層天上去；或在身內,我不知道,或在身外,我也不知道,只有神知道.12 : 3我認識的這樣的一個人-或在身內,或在身外,我都不知道,只有神知道-12 : 4他被提到樂園裡,聽見隱祕的言語,是人不可說的.12 : 5為這人,我要誇口；但是為我自己,除了我的軟弱以外,我並不誇口.12 : 6就是我願意誇口也不算狂,因為我會說實話；只是我絕口不談,恐怕有人把我看得太高了,過於他在我身上所看見所聽見的；

12 I must go on boasting. Although there is nothing to be gained, I will go on to visions and revelations from the Lord. 2 I know a man in Christ who fourteen years ago was caught up to the third heaven. Whether it was in the body or out of the body I do not know-God knows. 3 And I know that this man-whether in the body or apart from the body I do not know, but God knows- 4 was caught up to paradise and heard inexpressible things, things that no one is permitted to tell. 5 I will boast about a man like that, but I will not boast about myself, except about my weaknesses. 6 Even if I should choose to boast, I would not be a fool, because I would be speaking the truth. But I refrain, so no one will think more of me than is warranted by what I do or say, 7 or because of these surpassingly great revelations. Therefore, in order to keep me from becoming conceited, I was given a thorn in my flesh, a messenger of Satan, to torment me. 8 Three times I pleaded with the Lord to take it away from me. 9 But he said to me, ``My grace is sufficient for you, for my power is made perfect in weakness. " Therefore I will boast all the more gladly about my weaknesses, so that Christ"s power may rest on me. 10 That is why, for Christ"s sake, I delight in weaknesses, in insults, in hardships, in persecutions, in difficulties. For when I am weak, then I am strong.

你曾經得到上主的異象和啟示嗎?得到上主的異象和啟示的是甚麼人呢?你曾經被上主提到第三層天上去嗎?你曾經被上主提到樂園去嗎?上過天堂又能夠回到人間的是甚麼人呢?你曾經聽過隱秘的言語,天使的聲音嗎?擁有這超凡宗教經驗的這個人,蒙上主賦予神聖能力的這個人,我們或者可以稱之為「神人」.吊詭的地方是,這位有通天堂之能,起死回生之術的神人,竟然身上也有一根令他困擾的刺.我們這些凡夫俗子呢?有沒有感覺到有根令我們不舒服的刺總是跟著我們,揮之不去呢?這根令我們感到不舒服的刺又是甚麼呢?是身體的病痛?還是對疾病的恐懼呢?是工作間的某位同事?是過去發生過的某件事情?無論是甚麼,彷彿人人都有一根刺,或在內,或在外,或在前,或在後,令我們感到美中不足,感到掃興,感到不舒服,甚至感到困擾.我們如何對待這根刺呢?保羅又如何對待這根刺呢?

二．人人都想把刺除之而後快

12 : 7又恐怕我因所得的啟示太高深,就過於高抬自己,所以有一根刺加在我身上,就是撒旦的差役來折磨我,免得我過於高抬自己.12 : 8為了這事,我曾三次求主使這根刺離開我.

7 or because of these surpassingly great revelations. Therefore, in order to keep me from becoming conceited, I was given a thorn in my flesh, a messenger of Satan, to torment me. 8 Three times I pleaded with the Lord to take it away from me.

人生往往就是如此不完美和不協調.保羅樣樣皆好:正統希伯來人又擁有羅馬國籍,才智過人又受教於名師之下,滿懷理想而且執行力強,被上主揀選又賜予能力,然而,卻有一根刺在他身上跟著他,令他感到困擾.

令保羅困擾的刺是甚麼呢?保羅自己沒有說清楚,因此,留下了空間給後人猜想.約翰加爾文認為這刺是靈性的試探,引誘保羅疑惑和躲避使徒的責任,令保羅的良心受譴責而刺痛.馬丁路德認為這根刺是指他面對的反對與逼迫,經常與破壞他工作的人鬥爭.中古時代的修士們認為這根刺是指肉身的引誘.當僧侶或隱士關在修道院裡,在斗室中最難馴服的,是性慾.他們想驅除它,可是往往去了又來.近代的人通常視保羅身上的刺為身體上的疾病,有位精神病學家名叫安德森湯普森(Anderson Thompson)估計保羅患有癲癇症,即是香港稱為腦癇症的疾病,傳聞蘇格拉底,柏拉圖,凱撒大帝,畫家梵谷,音樂家海頓,都是癲癇症的患者.這位精神病學家認為,保羅在大馬士革路上突然失明的狀況可能是癲癇發作.另外,也有不少人認為保羅的刺是眼睛的疾病.保羅在加拉太書的結尾寫著:『請看我親手寫給你們的字,是何等的大呢!』(加六11).要寫那麼大的字,似乎是一個眼睛看不大清楚的人.

到底保羅身上的刺是甚麼,我們無從得知,但可以肯定的,是這根刺令保羅有除之而後快的渴求,因此,他三次求主除去他身上的這根刺,免得他繼續被困擾.

要把那根困擾我們的刺除之而後快,這是人類的共同反應.有沒有試過被木刺或鐵絲刺入手指呢?我們會感到不舒服,非要拔它出來不可.這些刺雖然小,沒有威脅生命,但留著它卻令人感到十分困擾,影響正常行動,因此,我們都想盡辦法把刺除去,才感到安心.

如果令人感到不舒服的刺是身體的痛楚,現代人比較有辦法,因為我們有有止痛藥.美國有關當部門報告,2000年至2010年這十年間,止痛藥的銷量大幅上升,有些城市上升了十六倍.(中央社紐約2012年4月5日綜合外電報導)不過,身體某處感到痛楚就馬上服用止痛藥止痛,這反應是否對身體最好的呢?

如果,令我們感到不舒服的刺,是我們身邊的人,又可以如何除去呢?

有一位男士被妻子以離婚威脅挾持去見婚姻輔導員.妻子不滿甚麼呢?是丈夫有婚外情嗎?不是,是丈夫沒有給家用嗎?又不是.原來,是妻子不滿丈夫沒有情趣,不能良好溝通,說得白一點,是妻子嫌棄丈夫是老粗一名.輔導員觀察這對夫妻,外來看來的確有點不相襯,丈夫是典型基層怒漢,似是地盤工人；妻子有點兒中產味,說話溫文爾雅,好像個文化人.然而,他們是如何走在一起的呢?原來,他們已經結婚二十年,大家都是中學同學,學生時期已經開始談戀愛.男的尚未畢業就出來工作,女的畢業後不久就嫁給他.婚後即為人父母,為了照顧孩子,女的成為全職主婦,一家四口的生活重擔就落在男的肩頭上.

時光飛逝,孩子漸漸長大,邁向獨立.這位太太在朋友的帶領下,開始參加教會.對她來說,真是大開眼界.她上了很多成長課程,學習溝通技巧,氣質也慢慢改變過來.她很喜歡教會生活,因為她覺得那裡的都是斯文人.而更重要的,是她明白了甚麼是美滿的婚姻,就是夫妻能夠彼此分享,溝通心事,在靈命上一同成長.

眼界開了,自然不甘現狀.她學了一套如何向未信主的丈夫傳福音的方法,除了每天祈禱之外,她十分積極地帶丈夫上教會,上成長課程,可惜丈夫只是偶然應酬.難得假日,他寧可在家中休息,也沒有興趣去聽耶穌,也沒有興趣去學習甚麼溝通技巧和夫婦之道.妻子失去耐性,尤其是當她自覺成長了,而丈夫卻是原地踏步,她看丈夫越來越不順眼,慢慢變為看不起他.這根刺,要除去嗎?

有些刺,不單刺激個人,更刺激著整個社群,社群的自然反應也是要他們消失於眼前.

六月底,我與十多位崇基同學往台灣參加華人校牧協會主辦的交流營.某天下午,營會安排了我們往收容成人嚴重智障人士的療養院進行學習服務.兩輛旅遊車載著八十位師生,從輔仁大學出發,穿越大城小鎮,到了八里山下.龐大的旅遊車開始上山,在窄小彎曲的山路上,顯得份外笨重.在上山的路程中,有同學看著窗外的墓地,好奇地問,為何台灣的山墳每個都建造得如此寬敞,比香港的大得多.我們笑說因為台灣人前住的空間比香港人寬敞,所以墳墓也比香港的大.旅遊車穿過墓地,終於到達療養院.療養院裡能夠出來參加活動的幾十位院友,都坐在輪椅上,欣賞同學們的表演,以及接受同學們協助他們吃晚飯.這次服務學習讓同學們有很多體會,其中一位同學問:「為甚麼智障人士宿舍要建在交通如此不方便的山上?而且建在墓地之中呢?」這會否反映出當地人如何看待智障人士呢?社群不想看見的刺,是否要把它遷到視線之外,死人之地呢?

有些刺,刺激著當權者,當權者自然的反應,也是要把刺除掉.

湖南邵陽市玻璃廠有一位工人名叫李旺陽,六四事件之後因政治罪名共被囚二十二年之久.2011年5月5日,李旺陽刑滿出獄,他的病情惡化,雙耳已經完全失聰,只能抬回家中,其後轉往大祥醫院接受治療.2012年5月22日,在朋友幫助下,李旺陽脫離了監視,接受香港有線新聞台訪問.訪問在2012年6月2日播出,引起海外關注,當地政府亦加強了對李旺陽的監控.6月4日,李旺陽獲全美中國學生學者自治聯合會頒發自由精神獎.6月6日,李旺陽的妹妹李旺玲收到大祥區人民醫院通知李旺陽在醫院病房自殺身亡.令當權者不舒服的刺,是否一定要除去呢?

這些可能令人的身體感到不舒服的刺,令人感到失望的刺,令社群感到不光彩的刺,令當權者感到尷尬的刺,是否一定要除去呢?如果不去掉,還有其他態度去對待這些刺嗎?

三．恩典把軟弱變為剛強

12 : 9他對我說:「我的恩典是夠你用的,因為我的能力是在人的軟弱上顯得完全.」所以,我更喜歡誇耀自己的軟弱,好使基督的能力覆庇我.12 : 10為基督的緣故,我以軟弱,凌辱,艱難,迫害,困苦為可喜樂的事；因為我甚麼時候軟弱,甚麼時候就剛強了.

9 But he said to me, ``My grace is sufficient for you, for my power is made perfect in weakness. " Therefore I will boast all the more gladly about my weaknesses, so that Christ"s power may rest on me.保羅三次祈求上帝除去令他困擾的刺,但上帝沒有應允,反而帶領他以另一種態度去對待這根刺.刺令人不舒服,令人困擾,令人尷尬,甚至阻礙人前進,是負面的東西,是討厭的事情,但上帝卻建議保羅以恩典的態度去看待這根刺,發掘這根刺出現在他身上的正面意義,以及這根刺的存在帶給他對信仰和人生的體會.當保羅能夠接納這根刺的時候,也就是他經歷上帝的恩典的時候；當保羅以恩典的態度去看待這根刺的時候,也就是保羅讓世人在他身上看見上帝的恩典的時候.

那根令我們困擾的刺,那根令我們除之而後快的刺,原來有時不一定需要拔掉,因為這根刺令我們看見一些重要但遺漏的事情,如保羅的體會,這根刺一方面令我們感到困擾,但同時間在保護我們,免得我們跌進自毀的陷阱.

粗魯的丈夫「死不悔改」,仍然像蠻牛一隻,令妻子失望和氣憤得發出離婚威脅.有一天,輔導員對那位妻子說,她今天擁有的斯文談吐和舉止,是建立在丈夫的粗魯之上,妻子聽後感到莫名奇妙.於是,輔導員向她說了下面這個故事.

昔日有一對十分要好的朋友,他們都想成為雕塑家,但他們的家境一樣清貧,何來錢完成心願呢?他們想出了一個辦法,就是其中一人工作,供給另一人去接受訓練.待訓練完成後,他們交換位置,這樣,他們兩個人都可以完成心願,只是時間長了一點.後來,首個接受訓練的朋友真的成為了雕塑家,並決定履行承諾.但當雕塑家告訴朋友輪到他去接受訓練時,朋友伸出粗糙的雙手問:「我還可以成為雕塑家嗎?」朋友默默地犧牲,成就了一個雕塑家,究竟那個已經成為雕塑家的人,要反過來嫌棄那個幹粗活的朋友,還是終身感激朋友無私的付出呢?

那位妻子聽完這個故事,無言以對.沒有丈夫辛勞的工作,把賺得的每一分錢花在家庭上,而不是用於修飾自己之上,妻子可以有這樣的空間和機會去變得斯文和優雅嗎?(黃麗彰:「修養」,《宣訊》第105期)

如果妻子能夠視在那根刺裡找到恩典,並以恩典回應,這可以成為多麼美麗的故事呢?

在回程的飛機上,坐在我旁邊的同學似乎仍念念不忘那次往智障院舍服務的經驗,他問:「社會應該如何看待這些人呢?他們沒有生產生力,也沒有市場價值.」他疑惑除了生產力和市場價格之外,還可以有甚麼參考去衡量人的價值呢?我們談到基督教的信念,每一個人都是創造主所創造的,因此,每一個人都有同樣的價值.一個真正強大的群體,是可以容納軟弱者於其中,與軟弱者共同生活的.我請他試想想,如果那所智障人士療養院不在山上而在山下的城鎮中,這個城鎮向世界發出了甚麼不同的信息呢?對於院友及其親戚朋友又有甚麼不同呢?這個城鎮告訴世界,他們是一個有恩典的社群,是一個方便院友和親戚朋友接觸的社區,是一個能夠把困難變為恩典,把軟弱變為剛強的群體.這根刺,可以保護一個社群避免沉淪在弱肉強食的殘酷之中,讓一個社群建立互相扶持,守望相助的人倫社會.

掘起的大國,應該如何看待那些刺激著當權者的刺呢?他們能夠看得出這些刺其實對國家和人民有正面作用嗎?他們能夠意識到這些刺其實正正是在保護國家和人民免受更大的苦難嗎?

從前,有一個很愛漂亮的國王,一天到晚命令王宮裡的裁縫師,替他做不同式樣的新衣裳,然而,沒有一件新衣令國王感到滿意.於是,國王決定到街市上張貼布告,上面寫著:「如果有人能做出讓國王滿意的新衣,可以得到一大筆獎金.」有一天,王宮裡來了兩個裁縫師,一高一矮,他們信心十足的說:「我們會一種魔法,可以做出世界上最神奇的衣服.」國王聽了,非常高興,急忙問:「那是什麼樣的款式,什麼花色?為什麼是世界上最神奇的衣服?有什麼特別的呢?」矮個兒裁縫師說:「它的顏色比花朵鮮艷,質料比雲彩輕柔,而且款式獨一無二,最特別的是,只有聰明人才看得見它.」「啊!」國王驚訝的瞪圓眼睛說:「那麼,愚笨的人就看不見這衣服囉?」「是啊,所以說它是一件神奇的衣服.」「太好了!」國王高興的說:「只要我穿上那件衣服,就可以分辨我的臣子,哪個是聰明人,哪個是不可信賴的笨蛋.」

過了幾天,兩位裁縫師捧著新衣去拜見國王.他們得意的說:「陛下,您的新衣做好了,您瞧,這是一件多漂亮的衣服呀!相信穿在您的身上,一定會顯得更年輕英俊,是世界上穿著最高貴的國王.」裁縫師一邊說,一邊將衣服翻來翻去,展示給國王看.國王雖然什麼也看不見,仍然裝出很滿意的表情,誇讚兩位裁縫師.於是兩位裁縫為國王穿上新衣.國王只覺得空空的,什麼也摸不到.不過,他仍然站在穿衣鏡前面,左看右瞧,一直誇讚說:「嗯,像羽毛一樣輕盈,穿起來真涼爽,你們的手藝的確高超,我要好好賞賜你們.」國王召喚眾臣來,問大家說:「怎樣,我這件衣服是不是真的很好看?」臣子們愣了一下,就同聲讚美說:「哇,好高貴的衣服呀!陛下您穿起來真是合適.」

國王穿著新衣,得意洋洋的走出王宮,到街上遊行.國王一邊走,一邊望向兩旁的人群,心想:「我這件衣服是給聰明人看的,不聰明的人就看不見,不知道這群人當中,有多少人看得見呢?這件衣服真是神奇.」其實每個人所看見的,是一個脫得光光,只剩下一條內褲的國王,只是沒有人敢說出來.就在這個時候,一個小男孩從人群中跑出來,指著國王哈哈大笑說:「哈哈哈!國王沒穿衣服,差差臉!」接著他指著國王的肚子說:「哇,你的肚子好大哦,你不穿衣服,不怕感冒嗎?」

這位小男孩如一根刺刺在國王身上,刺在群臣之中,國王應該如何回應呢?國王拉著小男孩的手匆匆跑回王宮.路邊的群眾都擔心小男孩會不會被國王處死.

沒想到,國王不但沒有處罰小男孩,還賞賜他許多禮物,並且誇讚他說:「你是個誠實的孩子,全國只有你對我說實話,讓我明白自己是個被人愚弄的傻瓜國王.」

在大國中揭露真相,說誠實話的人雖然被視為一根刺,但願有一天,這些刺能夠得到故事中小男孩的待遇.以恩典回應令人不舒服的刺,這才顯出真正的剛強.



\newpage

\section{信仰路上的四句話}
\label{sec:157}
\textbf{林豪恩先生 - 信仰路上的四句話}
\newline
\newline
\hyperref[sec:156]{< < < PREV SERMON < < <}
~
\hyperref[sec:toc]{[返目錄]}
~
\hyperref[sec:158]{> > > NEXT SERMON > > >}


馬可福音 4:35-41
\newline
\begin{longtable}{cl}
\hline
\hline
章節 & 經文 (和合本修訂版)\\
\hline
4:35 & \begin{tabularx}{0.7\textwidth}{X} 那天晚上，耶穌對門徒說：「我們渡到對岸去吧。」 \end{tabularx} \\ \\ \relax
4:36 & \begin{tabularx}{0.7\textwidth}{X} 門徒離開眾人，耶穌已在船上，他們就請他一同去；也有別的船和他同行。 \end{tabularx} \\ \\ \relax
4:37 & \begin{tabularx}{0.7\textwidth}{X} 忽然狂風大作，波浪打入船內，以致船灌滿了水。 \end{tabularx} \\ \\ \relax
4:38 & \begin{tabularx}{0.7\textwidth}{X} 耶穌在船尾上，枕著枕頭睡覺。門徒叫醒他，說：「老師！我們快沒命了，你不管嗎？」 \end{tabularx} \\ \\ \relax
4:39 & \begin{tabularx}{0.7\textwidth}{X} 耶穌醒了，斥責那風，向海說：「住了吧！靜了吧！」風就止住，大大平靜了。 \end{tabularx} \\ \\ \relax
4:40 & \begin{tabularx}{0.7\textwidth}{X} 耶穌對他們說：「為甚麼膽怯？你們還沒有信心嗎？」 \end{tabularx} \\ \\ \relax
4:41 & \begin{tabularx}{0.7\textwidth}{X} 他們就非常懼怕，彼此說：「這到底是誰？連風和海都聽從他。」 \end{tabularx} \\ \\
[1ex]
\hline
\hline
\end{longtable}
日期:2012年8月5日

題目:信仰路上的四句話

經文:馬可福音4章35-41節

證道:林豪恩先生

經文

4 : 35當那天晚上,耶穌對門徒說:「我們渡到那邊去吧.」4 : 36門徒離開眾人,耶穌仍在船上,他們就把他一同帶去；也有別的船和他同行.4 : 37忽然起了暴風,波浪打入船內,甚至船要滿了水.4 : 38耶穌在船尾上,枕著枕頭睡覺.門徒叫醒了他,說:「夫子!我們喪命,你不顧嗎?」4 : 39耶穌醒了,斥責風,向海說:「住了吧!靜了吧!」風就止住,大大地平靜了.4 : 40耶穌對他們說:「為甚麼膽怯?你們還沒有信心嗎?」4 : 41他們就大大地懼怕,彼此說:「這到底是誰,連風和海也聽從他了.」

在耶穌平靜風浪這事件中,馬可福音記錄了耶穌和門徒之間的四句說話,這四句說話刻劃了門徒這次驚險的渡海旅程.若我們把門徒這次渡海經歷擴闊來看,其實也描繪著耶穌跟隨者信仰人生的路程.讓我們仔細閱讀這四句說話,可能道出了我們曾經經過的狀況,或者我們現在正經歷的狀況.

一．「我們渡到那邊去吧.」門徒的驚險旅程,由耶穌的邀請開始.

4 : 35當那天晚上,耶穌對門徒說:「我們渡到那邊去吧.」4 : 36門徒離開眾人,耶穌仍在船上,他們就把他一同帶去；也有別的船和他同行.

按馬可福音的記載,耶穌曾經帶著門徒在加利利海的西岸四處醫病趕鬼,傳道教導.門徒看見他們的夫子滿有能力,能言善辯,現在,他們的夫子向他們發出邀請「我們渡到那邊去吧」,他們與耶穌一起,帶著期望起行,向著目的地進發.

曾幾何時,我們當中不少人也因為聽到耶穌的邀請而起行.有身心勞累的人聽到「凡勞苦擔重擔的人可以到我這裡來,我就使你們得安息.」(太十一28)而轉向跟隨耶穌；有盼望活得精彩的人聽到「 我來了,是要叫人得生命,並且得的更豐盛. 」(約十章 10節)而昂首闊步；有尋找人生終極意義的人聽到「我就是道路,真理,生命,若不藉著我,沒有人能到父(神)那裡去.」(約十四6)而踏上通往彼岸的道路；有徘徊於良心責備的人聽到「惟有基督在 我們還作罪人的時候為我們死,神的愛就在此向我們顯明了.」 (羅:5:8)而放下重擔,輕省上路.無論在甚麼地方,甚至時候,聽到甚麼邊請而起行,我們都好像門徒一樣,帶著盼望,懷著憧憬,向著應許的彼岸前進.

本地聖樂界前輩楊伯倫寫了一首詩歌,名為「思念信主初」,藉此鼓勵信主日久的信徒.詩歌描述了人們初信主時非常熱心,滿懷盼望,祈禱讀經,生活喜樂的狀況,這會否也是我們的集體回憶呢?

為何常問為什麼,強把身心來苦磨.皆因貪愛次等福,妄將目標來認錯.應思信主初,常算主恩多,求主眷顧我,常為我操舵,長沐主愛常謳歌.為何常問為什麼,須知歲月易蹉跎.仰望主顏思主愛,自能順服悔己過.再念信主初,毋忘主恩多,與神同行常謳歌.不應常問為什麼,誤將歲月來蹉跎.應思主訓無懶惰,處處莊稼處處禾.學習信主初,常念主恩多,殷勤事奉常謳歌.

「我們渡到那邊去吧」,起行的時候,我們往往都懷著有盼望,有方向,有應許的狀況開步.用甚麼可以比喻這種心情呢?讓我們回憶或想像有位我們傾慕已久的對象來對我們說:「我們拍拖吧」；又或者小學派位放榜日,心儀的學校校長來對我們說:「讓你的孩子來我們學校讀書吧」,又或者尋求升學的學生得到知名的教授說:「來跟我做研究吧」,又或者兩情相悅的戀人終於勇敢地說:「我們結婚吧」.帶著盼望,我們起行了.

門徒也起行了.讓我們繼續看他們在旅途中的經歷.

二．「我們喪命,你不顧嗎?」

4: 37忽然起了暴風,波浪打入船內,甚至船要滿了水.4 : 38耶穌在船尾上,枕著枕頭睡覺.

俗語說,「行船跑馬三分險」,耶穌的門徒當中,不少人是漁夫,怎會不知道海上航行有遇著風浪的危險呢?不過,今次的渡海旅程不可以與以往的出海打魚相比,因為今次的旅程非常特別,特別得不應以常理來估計.因為,第一,這次航程,是耶穌的吩咐.即是說,行動是耶穌發起的,門徒只是應邀而行.第二,這次航程有耶穌同在,從頭到尾,耶穌都在船上.由耶穌所發起的行程,並且有耶穌同行的旅程,應該與從前海打魚的旅程不同的,也應該與其他不是由耶穌發起的,也沒有耶穌同在的行程不同的,這樣的想像應該是有基礎的.

不過,有基礎的想像不一定是事實.縱使行程由耶穌發起,縱使耶穌也在船上,暴風仍可以忽然出現,波浪仍可以打入船內,船隻仍然可以陷入沉沒的邊緣.如果門徒曾經以為跟隨耶穌,有耶穌同行的這次渡海之旅,與他們從前出海打魚的行程完全不同的話,他們很可能會感到迷惑.為甚麼耶穌所邀請的行程,竟然也會遇到危難的呢?為甚麼有耶穌同在的旅程,竟然也有「搭沉船」的危險呢?

「我們渡到那邊去吧」,門徒應邀起行了,現在,門徒「在水中央」,大風大浪,後無退路,前行吉凶未卜,如何是好呢?起行時的期望,應許,雀躍,自信彷彿仍然飄盪於空氣中,然而,眼前發生卻是另外一種狀況,非常不協調,也非常令人疑惑及失望.類似的境況,我們是否曾經相識呢?

有一對男女兩性相悅,他們確認上主的帶領,懷著興奮和期待的心情共築愛巢.踏過紅地毯,站在聖壇上,堅定地,深情地在上主面前,在眾親友面前互許承諾.晚上婚宴結束回到新居,他們依偎在客廳柔軟的疏化上,望見掛在當眼處寫著「基督是我家之主」的掛飾,既感恩,又溫馨.

時間一天一天過去,一年一年過去,丈夫心裡暗叫不妙.他越來越發覺太太「貨不對版」,睡在旁邊的這位女士,與自己結婚前的想像相距太遠.不滿和失望難免流露於言語和態度,批評引發吵架,輕視激起憤怒,積壓下來的傷心和失望令兩個人都沉默起來,同床異夢,有時各自偷泣.這對夫妻焦慮地感受到他們這隻船正在下沉,然而,其他人似乎看不見,教會的弟兄姊妹仍視我們為模範夫妻.這隻船已經入滿了水,災難隨時發生.要「跳船」嗎?要讓它自然沉沒嗎?還有更多的可能選擇,他們都一一想過.他們沒精打彩地坐在客廳的那張疏化上,望見那個寫著「基督是我家之主」的掛飾仍默默地高懸在客廳的當眼處.他們疑惑:如果「基督是我家之主」,為甚麼我們的婚姻會如此艱難呢?

「我們渡到那邊去吧!」這邀請令人感到有方向,有盼望,有應許.然而,彼岸未達,現在卻有客死途中的凶險.門徒如何回應呢?

4 : 38門徒叫醒了他,說:「夫子!我們喪命,你不顧嗎?」

「跳船」,可以是一個選擇.不過,卻不是唯一的選擇,也未必是可行的選擇.

當那位「完全正直,敬畏上帝,遠離惡事」(伯1:1)的約伯在人生旅程的某一刻遇到橫禍時,他的太太鼓勵他「跳船」:「你仍然持守你的純正嗎?你棄掉上帝,死了吧!」(伯2:9)約伯拒絕「跳船」,他堅持要與上帝爭辯,要上帝回答他.

先知耶利米也曾感到自己跌入了這種陷阱之中,他沒有選擇「跳船」,而回去向呼召他的耶和華投訴:「耶和華啊,你曾勸導我,我也聽了你的勸導.你比我有力量,且勝了我.我終日成為笑話,人人都戲笑我.我每逢講論的時候,就發出哀聲,我喊叫說:有強暴和毀滅!因為耶和華的話終日成了我的凌辱,譏刺.」(耶20:7-8)

耶穌的門徒選擇了去叫醒耶穌,亳不客氣地責問耶穌.他們對耶穌的責問是有合理的,說到底,是耶穌發起這次行程的.現在危急存亡的時候,門徒都在與風浪搏鬥的時候,這位邀請他們啟船的耶穌竟然還在睡覺,難怪門徒質疑他到底是否關心他們,是否置他們的性命於不顧.

在危難的焦躁中,上帝竟然沉默不語,「跳船」,可能是一項「好過等死」的無奈選擇.耶利米沒有選擇「跳船」,約伯沒有選擇「跳船」,耶穌的門徒都沒有「跳船」.門徒去弄醒那位在船尾睡覺的耶穌,約伯去質問上帝,耶利認為自己只是遵照上主的吩咐而行,因此認為上主應該要為目前的狀況負責.當他們去叫醒耶穌,質問上帝,對上帝死纏不放的時候,他們對身處的危難,對自己及對上帝都有了不同的體驗.

三．「為甚麼膽怯?你們還沒有信心嗎?」

耶穌平靜了風浪,卻不讓門徒剛才遇到的危難經驗流走,否則,門徒就白白地辛苦了.耶穌藉著反問,引導門徒探索自己的心靈狀態.難道門徒不知道自己「膽怯」嗎?難道門徒不知道自己「沒有信心」嗎?這並不奇怪,當門徒在水上遇到凶險的時候,當船入滿了水的時候,一切的行動都是即是反應:趕著穩定船隻,忙著把傾倒入了船的水,慌著去叫醒耶穌,憤怒地問責.其實,大多數人在危難中的行為都只是反應,有些時候,甚至連日常行為也是反應,甚少體會到反應背後的心靈狀況.

「直升機父母」這一個詞彙源於美國,同來形容經常偵測子女行蹤的父母,如直升機般在子女頭上盤旋.若果問父母,為甚麼要每日偵測子女的行蹤呢?他們會回答,當然是要保護子女,難道我空閒得沒事作嗎?其實,父母的內心充滿懼怕,擔心子女遭綁架,終日惶恐不安.

有本地社工指出,不少家長對子女的行為也是一種反應.有家長知道孩子的同學找補習老師,於是立即為自己的孩子安排補習；有家長看見別的家長為孩子報某項課外活動,又立即為自己孩子搜羅課外活動；有家長得知孩子的朋友參加交流團,馬上張羅為孩子安排交流團.如果要解釋,每一項行動都可以有合理的原因,例如補習可以幫助孩子取得更好成績,課外活動如何幫助孩子增加升學機會,交流團可以擴闊令孩子增廣見聞.然而,若果父母進入自己的心靈世界,感受自己內心的聲音,可能會聽到自己對於孩子吃虧的擔憂,感受到對孩子的競爭能力的焦慮,以及對孩子的前途的不安.然而,父母卻不容易察覺可能是自己內心的擔憂,焦慮和不安驅使自己行動,而未必是孩子此時此地的需要.

「為甚麼膽怯?你們還沒有信心嗎?」耶穌這提問令門徒把目光從外在的風浪轉向內心的風浪.啊,原來我們正在膽怯,原來我們還沒有信心.門徒可以自我合理地解釋為甚麼膽怯,為甚麼還沒有信心.不過,若果問徒合理化自己的反應,把膽怯及沒有信心的心靈狀況解釋去掉,門徒又失去了一個進入自己的內心世界去體會自己的機會了.

甚麼狀況才能逼使我們進入自己的心靈世界呢?甚麼事情才能使我們願意面對內心的風浪呢?

前文提及的那位在婚姻路上進退兩難的夫妻,從來遇到一項不能夠不面對自己的事情.某年某月某日,太太被診斷患上癌症.對著這個消息,丈夫的心情十分複雜.某天,太太對丈夫說:「這樣也好,我們不需要再糾纏下去了,你可以名正言順去找另一位你滿意的人了.」聽到身患絕症的太太如此說,丈夫心裡揪著痛,他走進房間,獨自向上帝祈禱.他發現自己對太太諸多不滿,其實自己也好不了多少,甚至比太太還要差得多.他向上帝承認自己寡情薄義,向上帝承認自己不懂得愛,雖然在婚禮上讀過「愛是恆久忍耐,又有恩慈．．．」這愛的詩篇,但自己對太太沒有忍耐,更沒有恩慈.他跪在地上,請求上帝赦免及憐憫.這個晚上,他們夫婦兩人握著手,跪在上帝面前,就像婚禮時握著手跪在聖壇前一樣.丈夫向太太承認自己的帶給對方的傷害,太太也坦承自己因為欠缺安全感,經常懷疑上帝是否愛自己,也經常懷疑丈夫是否真心愛自己,而為婚姻添上困難.這個晚上,他們終於能夠面對自己,彼此諒解,並請求上帝賜給他們恩典及力量一齊去面對太太的癌症.

四．「這到底是誰,連風和海也聽從他了.」這位丈夫憶述說:「在婚姻遇到困難的時候,我們曾經要想迴避上帝,做自己想做的事,但祂卻耐心地等待我們回歸.回想過去,除了感恩之外,沒有可表達的,因為神對我們不離不棄,並用重價買贖我們回來.」

馬可福音以門徒經歷過這海上危難後所說的話作結:「這到底是誰,連風和海也聽從他了.」看來,這次驚心動魄的事件衝擊了門徒對耶穌的認識.

那位承受著巨大打擊也不肯「跳船」的約伯,後來上帝出來為他和他的朋友的爭辯作結,並且回應約伯,約伯因此而說出傳頌後世的說:「我從前風聞有你,現在親眼看見你.」(伯42:5)

其實,這對夫妻聖經中讀過關於上帝的恩典,上帝的拯救,也在教會中聽過上帝對人不離不棄,等待人回歸的心懷,然而,走過婚姻路上的波折起伏,上帝的恩典和拯救成為了他們的親身體會.

耶穌的門徒也看見過耶穌的神蹟,又聽過耶穌的教導,不過,他們經歷這次風浪,所有這些看見的和聽見的,都變得切身了.

有一位同學分享說,某天,她的男朋友來跟她說分手,她如晴天霹靂,質問上帝為甚麼祈過禱才開始的戀愛也要分手,為甚麼上帝不保守他們的關係.她恨天父到一個地步,不肯讀經,不想祈禱,好像女兒跟父親鬥氣.每當拿起聖經,拉開拉鍊只及一半,便氣憤的再拉上,把它擱在一旁.甚麼停止到教會聚會,想逃避所有人.經過一段日子,有一次她接觸到溫偉耀博士《跨越苦難》的錄音帶,聽到溫博士在失去了深愛的妻子,照顧一個弱智女兒,面對種種痛苦,仍相信天父和信靠的分享.溫博士的故事,感動了這位同學.她開始再翻閱聖經,在何西阿書讀到天父對她的不離不棄,她感受到莫大的安慰與關切.她深深體會到,縱然她向天父發脾氣,但天父忍耐,等候她,默默以愛,接納她回到身邊.

「這到底是誰,連風和海也聽從他了.」這到底是誰?上帝遠超過我們對他的認識,也往往令我們有意想不到的驚訝,以至我們在跟隨他的人生路上,可以不斷地發出:「這到底是誰」的驚嘆.

五．總結:第五句不是對我們說的話

「我們渡到那邊去吧!」我們帶著期待,朝著目的地起程.

「我們喪命,你不顧嗎?」旅程未必一帆風順,在途中,我們可能會遇到危難.有危難,不一定代表走錯了路,不一定要「轉馱」,也不一定要「跳船」,重要的,是耶穌仍在船上.

「為甚麼膽怯?你們還沒有信心嗎?」提示我們不要單單回應外在的風浪,更要進入自己的心靈,平息內心的風浪.這樣,才能夠認清方向,培養耐力,繼續向著目的地前行.

「這到底是誰,連風和海也聽從他了.」在旅程中,上主讓我們透過親身的經驗,更新對他的認識,也把把對他的認識變得切身.最後,讓我們留意那句在事件中耶穌不是對門徒說,而是對風和海說的話:「住了吧!靜了吧!」無論我們經歷的事情有多凶險,耶穌邀請了我們起行,就總有辦法把我們帶到彼岸.



\newpage

\section{別讓聖靈擔憂你}
\label{sec:158}
\textbf{賀志勇博士 - 別讓聖靈擔憂你!}
\newline
\newline
\hyperref[sec:157]{< < < PREV SERMON < < <}
~
\hyperref[sec:toc]{[返目錄]}
~
\hyperref[sec:159]{> > > NEXT SERMON > > >}


以弗所書 4:25-30
\newline
\begin{longtable}{cl}
\hline
\hline
章節 & 經文 (和合本修訂版)\\
\hline
4:25 & \begin{tabularx}{0.7\textwidth}{X} 所以，你們要棄絕謊言，每個人要與鄰舍說誠實話，因為我們是互為肢體。 \end{tabularx} \\ \\ \relax
4:26 & \begin{tabularx}{0.7\textwidth}{X} 即使生氣也不要犯罪；不可含怒到日落， \end{tabularx} \\ \\ \relax
4:27 & \begin{tabularx}{0.7\textwidth}{X} 不可給魔鬼留地步。 \end{tabularx} \\ \\ \relax
4:28 & \begin{tabularx}{0.7\textwidth}{X} 偷竊的，不要再偷；總要勤勞，親手做正當的事，這樣才可以把自己有的，分給有缺乏的人。 \end{tabularx} \\ \\ \relax
4:29 & \begin{tabularx}{0.7\textwidth}{X} 一句壞話也不可出口，只要隨著需要說造就人的好話，讓聽見的人得益處。 \end{tabularx} \\ \\ \relax
4:30 & \begin{tabularx}{0.7\textwidth}{X} 不要使神的聖靈擔憂，你們原是受了他的印記，等候得救贖的日子來到。 \end{tabularx} \\ \\
[1ex]
\hline
\hline
\end{longtable}
日期:2012年8月12日

題目:別讓聖靈擔憂你!

證道:賀志勇博士

經文:以弗所書4:25-30

4:25 所以你們要棄絕謊言,各人與鄰舍說實話.因為我們是互相為肢體.4:26生氣卻不要犯罪.不可含怒到日落.4:27 也不可給魔鬼留地步.4:28 從前偷竊的,不要再偷.總要勞力,親手作正經事,就可有餘分給那缺少的人.4:29 污穢的言語,一句不可出口,只要隨事說造就人的好話,叫聽見的人得益處.4:30 不要叫神的聖靈擔憂.你們原是受了他的印記,等候得贖的日子來到.

一,引言:真貨與真言

最近我有很多朋友來香港買東西,買的不是化妝品、LV,買什麼呢?奶粉,洗頭、沖涼水,甚至還有醬油、油,有一個甚至特意買了一大箱益力多回去.我問他們買這麼多做什麼?他們都回答說,香港的東西真嘛!以前就聽過大陸一個朋友講過一個笑話,講一個農夫同太太吵架,一氣之下就喝了農藥,結果沒有死成,因為那個農藥是假的；他非常高興,送了一塊匾給農藥公司,然後出去慶祝喝酒,結果喝了一瓶酒就死了,因為酒也是假的,有毒.

假貨,大家都知道不好,但是假話--謊言呢?你說:我從來不講假話!這句話就是一句大話言.英國心理學家講,人講話經常有灰色地帶,每人每天平均要講25句不真實的話,其中有五句一定是假話.我舉個例,流行的假話有這麼幾句:"我今天來不了,我有點事."--其實沒什麼事；"得閒找你飲茶啊,遲點同你聯絡!"遲點多數不會聯絡."我給你打過電話,但是打不通."--多數是沒有打過電話的."你看起來瘦了哦!"--其實瘦沒瘦我們還真沒看出來.

或者你說講假話的人一定沒信仰.那可不一定了.聖經裡面有個主要人物,叫做亞伯拉罕,他就不止一次講大話.大家都知道,他老婆太漂亮,他怕人家搶他老婆還要殺他,就兩次對別人說他老婆是他妹妹.之後的事就不用我八卦了.

二,經文分析

(1) 25節,棄絕謊言

為什麼我們要說謊話了?說假話不是一個人的事,其實是兩個人之間的事,有的人講假話是為了討別人歡心,讓人家感覺好一點；有的人是為了誇耀自己和裝派頭；第三類是為了自我保護. 說穿了,說假話意味著兩個人的關係還不夠信任.我有個朋友的兒子上中學,每次問他:"你今天在學校考試怎麼樣啊?"他兒子就答他:"Fine!"有一次他就問:"為什麼你次次都是Fine?"他兒子說:"你天天都要問我.我說不Fine,你肯定要罵死我,不如講Fine啦!"

所以說,要彼此不講謊言,首先要建立彼此信任的關係.25節要我們棄絕謊言,棄絕這個字,原文就是"脫掉"的意思.舉個例子,你上街踩到狗屎了,回家第一件事就是即刻脫掉鞋子,趕快洗乾淨.謊言就是那對有狗屎的鞋子.你想想,一個人天天穿著有狗屎的鞋子出去逛街,是什麼感覺?神提醒我們,我們不是敵人,不是陌生人；我們是鄰舍,是肢體.屈梭多模就講:"眼睛看到蛇,會不會欺騙腳?舌頭嘗到苦味,會不會欺騙胃?"所以如果我們視身邊的人是肢體,我們就要彼此講實話.

你或者會說,賀博士,你不食人間煙火.事事都要講老實話,豈不是自尋死路.老闆問你:"我這個方案好不好?"你夠膽講不好?太太問你:"我穿這件衣服靚不靚?"你夠膽講不靚.你說不靚,她即刻出去買多幾件給你看看.是的,講實話不容易；但是了,憑著愛心講實話更加不容易,這個不用急,這個是要慢慢鍛煉的.聖經要我們脫去舊人,這個要慢慢來.當你發現自己假話越講越少的時候,你就在一點一點變成新人了.

(2) 26-27生氣但不要犯罪

脫去舊人,還要一點一點脫掉我們的臭脾氣. 有的人以為脫掉臭脾氣,就是不可以生氣.回家晚了,太太講多了兩句,明明不高興,嘴裡還要扮樣子.太太問他:"你生氣了啊?""沒有啊!""不高興你就說嘛!""沒有啊!"

聖經從來沒有說不可以生氣.耶和華都有怒氣.沒可能人家欺負你小孩子,你都不生氣的.但是聖經說,生氣但不要犯罪.人家罵了你的兒子,你不要衝過去打一頓他的兒子.太太講多你幾句,你不要用汙言穢語彈回去.聖經教導我們"要快快地聽,慢慢地說,慢慢地動怒."有個女士說:"我脾氣很直的,我發了脾氣之後就沒事了."她的朋友回答她:"是哦!就好像放了一炮,不過你知不知道放了一炮後果有多大?"

聖經不是說生氣了不要表達出來.聖經講,不可含怒到日落,意思就是說如果你有氣就要表達出來.既然已經生氣了,你不可能當沒事一樣.但是關鍵是如何表達你的怒氣,這個是需要技巧的.有的人喜歡以"你"字開頭來表達:"你怎麼這點小事都做不到啊?""你怎麼會這麼差啊?"你字一開頭,對方就覺得你是在指責他,那他聽完了就不是含怒,那就一定是發怒了.要表達,我提議多以"我"字開頭,"我不高興了,因為...","我覺得受到了傷害,因為..."你用我字開頭,表達自己的感受,並且解釋了原因,好好溝通就可以了,指責的話,一句不可以出口.

生氣了就要表達,人家說"夫妻兩個是當天帳當天算",因為不表達,氣鼓鼓埋在心裡,容易給魔鬼留地步.魔鬼,原文是撒旦,也就是"誹謗者"的意思.神發現,一個人如果含怒,不同當事人當面表達,就會私下跟其他的人表達,就容易生出流言蜚語,是是非非來了.是非一多,就容易變成誹謗了.所以詩篇說:"在床上的時候,要心裡思想,並要肅靜."我們彼此相愛,就不要給"誹謗"留地步.

(3) 28節不要再偷,做正經事

脫掉舊人,28節提醒我們,還要不再偷竊.十幾年前我去到老廣州火車站,印象就是人多、小偷多.所以廣州的朋友都特別提醒我,到了火車站要小心啊.偷東西,就是說我是打工的,小偷變成老闆,我替小偷打工了.沈保羅牧師在東京開學生佈道會,講的就是這段經文.講完之後,有個青年當時被聖靈感動,站起來當眾承認自己在東京念博士7年,坐地鐵從來沒有買過票.他站起來認罪,散會之後他就去地鐵公司一次性付清了這筆賬.

有的朋友說,一個博士怎麼會做這樣的事呢?的確,小偷這個詞,好像離我們挺遠的.不過,或許有的人不會去盜取人家的錢,但是會不會盜取人家的時間,偷人家的名氣呢?小說《圍城》裡面有個讀哲學的,就經常跟人家說大哲學家羅素等等是他的朋友.你說這算不算偷呢?其實,想多一層,我們有沒有盜取神的時間、神的榮耀呢?很多教會都有崇拜遲到的現象,牧師們一再提醒,效果也不是很明顯.這算不算盜取屬神的時間呢?

所以以弗所書告誡我們,不要再偷竊了,要我們親手做正經事,還要分給缺乏的人.正經這個字,原文就是"有用、有益、、美好、慈善"的意思,同中文正經不正經這個字關係其實不大.我們崇基禮拜堂有不少會友在做各種慈善事業,這值得大大鼓勵.但如果你說,我錢不多,分不了什麼給窮人,那你就錯了.我知道,崇基有的弟兄姊妹會收集衣物去派給露宿者,這就是非常好的分享.其實,平時去探訪一下弟兄姊妹,或者打個電話同朋友分享溝通一下,你就把你的時間分給他了.我認得一位牧師,他經常讚賞別人的優點,懂得欣賞別人,還經常推薦人,他就把他的眼光、威望都分給別人了.分給人的不一定是物質,可以是心情,可以是眼光,可以是觀點,可以是時間.總之所有有益、美好的東西都可以分享.

(4) 29節 不說汙穢言語,說造就人的話

當然,美好的東西可以分給人,汙言穢語就不要分給人了.汙穢這個字,原文指腐爛、壞的東西,聖經裡一般用來形容變壞的魚、腐爛的水果.變質的海鮮、爛了的水果還能不能吃?不可以吃!那麼變質、腐爛語言可不可以講給別人聽,呢?汙言穢語不一定是下流的語言,北京人就有這種本事,就是罵人不帶一個髒字,但是聽完之後你一定想打人.我小時候聽老師罵一個學生,老師講:"你哥英文這麼好,你英文怎麼這麼差,你同你哥是不是親生的啊?你是你爸媽親生的麼?"這種話千萬不能說,因為一說出口,有的關係就難以挽回了.

不但污言穢語不能說,聖經進一步要求我們,要凡事說造就人的話.說造就人的話,不一定是說好聽的話這麼簡單.造就人的話,原文的意思是"與場合相配的好話."好聽的話,如果與場合不相稱,就沒辦法造就人；不好聽的話,如果與場合相稱,也一樣可以造就人.所以箴言15:23節說:"話合其時,何等美好!"其實,我們的主耶穌就是非常有講話藝術的人,當面對法利賽人問他:"納稅給凱撒可以不可以?"他就回答"凱撒的物當歸凱撒,神的物當歸給神."(太22:17-21)還有一次,抓住一個行淫的婦人,法利賽人問耶穌是不是要用石頭打死這個婦人,耶穌就說了一句話:"你們中間誰沒有罪的,誰就第一個拿石頭打她."那些人很羞愧,沒有人敢做第一個,就一個一個溜掉了,然後耶穌對那個婦人說:"我也不定你的罪,去吧!從此不要再犯罪了."你看看,說話多有智慧、多有藝術.他一句話就讓眾人知道羞恥,一句話就救了這個女人的命,但是他也一句話就責備了這個女人,讓她反省、知錯.所以箴言25:11節講:一 句 話 說 得 合 宜 , 就 如 金 蘋 果 在 銀 網 子 裡 .

當Steve Joe還沒有大大成功之前,他曾經邀請百事可樂的總裁約翰$\,\cdot\,$斯卡利(John$\,\cdot\,$Scully)做蘋果的銷售總裁,當時Scully年薪是兩千萬美金,所以Scully不理睬Steve Joe.於是Steve Joe對Scully說:"你留在百事可樂,五年後只可以多買些糖水給小朋友.你是想一輩子賣糖水,還是跟我一起改變世界?"聽了這句話,Scully就去了蘋果公司做總裁.

所以說,造就的話不一定是好聽的話,它可以是讚揚,也可以是批評；可以是分享,也可以是啟發.只要場合適宜,叫聽見你講說話的人得到益處、得到神的恩典就好了.

三,結論:30節:別讓聖靈擔憂你

我這個人有時喜歡對號入座.比如沙斯那年,2003年,我那幾個月幾乎每天都量體溫,一看到體溫表顯示體溫升高一點,我就非常有去看醫生的衝動.不過想一下,醫院病人多,去了可能反而更加不安全,所以只好控制好自己的衝動.

今天我這個對號入座的脾氣又犯了.30節講"別讓聖靈擔憂你",我就發現讓聖靈擔憂的有這樣幾種現象:第一種現像是經常把"彼此相愛,互為肢體"掛在口邊,但是其實從來沒把身邊的弟兄姊妹當肢體,所以也談不上有多"彼此相愛"；第二種現像是經常有情緒、有想法,但是都埋在心裡,人家根本不知道他心裡想什麼.不單單含怒到日落,可能含怒一個星期都有可能.第三種現像是大小活動都喜歡遲到,而且答應的事不負責任.偶爾遲到我們可以理解,但是大小活動都喜歡遲到,那就擺明是要偷人家、偷神的時間了；第四種現像是講話喜歡擺著個面孔,好像生怕人家不知道他脾氣不好,或者生怕人家不知道他有料那樣.不講話都不造就人,一講話就更不造就人了.

弟兄姊妹,千萬要記住剛才的四種現象,因為那是讓聖靈擔憂的四種現象.我們都是屬神的人,額頭上都有屬神的印記.但願世人都看到"你們的好行為,便將榮耀歸給你們在天上的父"(太五:16).願我們脫去舊人,活出聖經教導我們的新生命、活潑的生命.如果你真的能活出新的生命,你就會發現那不單只是責任,更帶著神的祝福!



\newpage

\section{單純與委身}
\label{sec:159}
\textbf{梁燕城博士 - 單純與委身}
\newline
\newline
\hyperref[sec:158]{< < < PREV SERMON < < <}
~
\hyperref[sec:toc]{[返目錄]}
~
\hyperref[sec:160]{> > > NEXT SERMON > > >}


約翰福音 16:8-11
\newline
\begin{longtable}{cl}
\hline
\hline
章節 & 經文 (和合本修訂版)\\
\hline
16:8 & \begin{tabularx}{0.7\textwidth}{X} 他來的時候，要為罪、為義，為審判，指證世人； \end{tabularx} \\ \\ \relax
16:9 & \begin{tabularx}{0.7\textwidth}{X} 為罪，是因他們不信我； \end{tabularx} \\ \\ \relax
16:10 & \begin{tabularx}{0.7\textwidth}{X} 為義，是因我到父那裡去，你們將不再見到我； \end{tabularx} \\ \\ \relax
16:11 & \begin{tabularx}{0.7\textwidth}{X} 為審判，是因這世界的統治者已受了審判。 \end{tabularx} \\ \\
[1ex]
\hline
\hline
\end{longtable}
日期:2012年8月19日

主題:單純與委身

經文:約翰福音 16:8-11

證道:梁燕城博士

最近大概是四,五月的時候,中國很多暴雨,我到了廣西很遠很遠的山上──廣西百色這地方,就是鄧小平革命的地方,百色是一個很遠很遠的地方,我們從南寧坐車要六個小時才到,到了百色還要找一個叫隆林縣的地方.在那裡我們有一些工作,隆林縣是一個更遠山上非常貧困的一個縣,我們要探望一些我們資助的同學.從百色到隆林也要一個多小時,從隆林再行車一個多小時才到山裡面.那個時候暴雨之後,路很爛但也是非常乾旱,農田已經乾旱了三年,暴雨忽然又來,把路都破壞,後來我們走另外一段路,到資源這個地方,我們要轉二百多個彎,那些塌山塌石的有五十多處,中國政府很快就把路打通了,但路也是不容易走.當我到了隆林那個鄉──沙梨鄉,到那村裡也試探望我們自資助的一些貧困學生.

我們見到一位女孩子,她的名字叫蓉,我們和她談話,問她情況如何?她就痛哭出來,我們同行的一些太太立時把她抱著,她說因為家很窮,農田亦種不出東西,爸爸跑到外面工作,但是跑出去已四年都沒有消息,爸爸是生是死沒有人知道,爸爸出去後是病?受傷?還是跟其他女人結婚都不知道?四年了,爸爸沒有消息.我們對著這個女孩我們真的無能為力,惟一我們只能幫她一點錢讓她讀書,惟一我們只有一天去看她,來對她表示一種鼓勵和安慰.

在中國貧困的大眾面前,我們發現我們真的是無能,另外我們也看其他的農民,一位我們供他讀高中快要考大學的孩子,他也是農田種不出東西,一家都沒有收入,他們的家真的是破爛得不得了.另外想起蓉,她並不是一個漂亮的女孩,且成績只有三百多分,可能考不上大學,也不知道前途如何?後來我們回到隆林縣見另外一批學生,每個學生都有血淚的故事在後面.其中有一位他說,因為窮,爸爸要到外面打工,可是老板不給工資,他們去理論的時候,爸爸給老板派來的人給打傷,然後受了傷又病,身上只有三十塊錢,媽媽要跑到城裡面把爸爸帶回來,她是一個女孩子叫雪巧,她看到爸爸受了傷又病回來,她心中非常疼,但是作為一個孩子也是無能為力,她又說看到我們很遠很遠的地方來的人,愛他們,關懷他們,她本來以為自己沒有機會,沒有錢讀高中,但是我們支助她讀書,她很激動的說她明白這個叫「愛」.

這愛是無條件的,我們從很遠很遠的地方來,我從加拿大飛十二個小時到香港,香港再飛兩個小時到南寧,南寧要開車六個小時到百色,一個多小時到達隆林縣,再一個小時才能到山裡面見他們,花了差不多二十個小時來見他們.我告訴他們我們那麼遠來只是為了愛你們──愛是沒有條件的,不認識的人也要愛──在這個時候,我想起作為一個孩子看到家庭的困難,是如此無能為力,面對著家庭的困難,他只能夠哭；有些是爸爸媽媽死了,有一些是爸爸媽媽離婚了,他們都是如此的無能為力,我們能夠給予的幫助很有限,我們只是給他們某一種愛的關懷而已,他們已經哭得不得了.

我想起我從前也不是這樣嗎?當我還年青的時候,還是孩子的時候,我爸爸在我小的時候很好,他很有學問,常常跟我們講歷史,他牽著我手去看電影,回家的時候經過一個玩具店,他總是牽我過去買一些小的車子,我記得我很高興,每次爸爸買個小車子我就很高興.但是當我約十歲的時候,中國大躍進失敗,三年的飢荒來臨；很多人餓死在中國,很多飢餓的人從邊界跑到香港,那是在中國最艱難的日子.

我爸爸就患了精神病,他們都是知識份子,我爸爸媽媽都懂英語的,他們用英語吵架的,在那個年代他們是高知識份子,他們受了很多訓練,但沒辦法回報祖國,他們跑到香港也是鬱鬱不得志.爸爸患上精神病,常常認為有人害他,常常跟媽媽吵架,每逢吵架爸爸就很兇的把東西丟到地上,小的時候最怕就是爸爸媽媽吵架,爸爸開始變得很兇,也不工作,天天在家抽煙,自言自語.作為孩子我們好像真的是無能為力,爸爸變成一個很可怕的人,我是如此的無能為力,然後媽媽去工作,她懂英語找到工作,還能維持家庭,爸爸常常欺負媽媽,我也無能為力.

我知道爸爸的痛苦難受,我們國家的悲痛,後來文化大革命,在很多的悲痛時間裡面,對中國我們也是無能為力的.這個大概是我年青還單純的時候所經過的經驗,同時我更明白,記得孩子的時候,看到哥哥有一本近代史課本,裡面有些圖畫,我問這個是什麼?哥哥說:「這個叫八國聯軍」,我看到軍隊打進北京；後面是日本「九一八事變」,從那些圖畫我知道,我們這個民族過去經過很多很多悲痛的事情,從小到青年 and \#8230; 我們都是如此的無能為力.

不過當我有機會回中國,我一定盡我能力,讓痛苦的人減少他們的痛苦,只要你願意用真誠單純的愛去關懷他們,中國是願意讓你回去幫助他們的,而且中國願意聽你的意見,因為多年來,我都把民間的苦難,民間的疾苦,我看到的,我有十年在農村,也到了城市的農民工當中,我知道他們的苦,我就把教育,醫療,體制等等的問題,還有官僚,濫權等等都寫成建議書給中央,我希望中國可以改好一些體制的問題,讓人民活得更好,而且我也寫了有關人權,有關中國自由 and \#8230;反正我在海外 and \#8230;最初我交給領事館,領事館說你講的都很好,你是為了中國好而寫的,但領事館說告訴他們也沒用,他們不是很高的官員,只不過是外交部的,起不了什麼作用,不如介紹我到國務院去反映意見.從那時候我就到了國務院,後來就到了黨中央,我總是想辦法寫我知道的情況,希望政府可以改好,所以最近李旺陽事件,我也準備寫.希望中國政府能夠對下面的官員,教導他們尊重每個人民的生命,每個人的尊嚴是比一切更重要,人民獲得尊重,國家才真正有穩定,多年來我就是憑著我的真誠和單純寫下這些.

耶穌曾經說過在祂遇到很多批評,很多困難的時候,在馬太福音第十一章裡面,祂特別講到「耶穌說:父阿,天地的主,我感謝你,因為你將這些事,向聰明通達人,就藏起來,向嬰孩,就顯出來,父阿,是的,因為你的美意本是如此.」這一段的聖經,馬太福音十一章廿五節,耶穌說在祂最困難的時候,給人看不起,給人批評的時候,祂回到上帝的親情裡面,」父阿」,就是一個親情的稱呼,祂不是說:「最高的上帝阿,我求你什麼」.而是上帝你就是我的父親,我是如此的愛你,你如此的愛我,所以在最艱難的時候,祂就回到父的面前,而在父的面前祂有一種親情相愛的關係,「父阿天地的主」,第二句才說,是你掌管的,宇宙萬物你是看到的,世界上有些發生的事情好像很不公義,很不公平,很多困難的事情,人沒有愛,但是祂說:父阿!天地的主,你還在掌管,只要我知道你在,這個世界一定會有前途,一定會有盼望,一定會變得更好.所以祂說:天地的主,我感謝你,雖然我遭遇困難,但我還是感謝你!因為只要有一個感謝的心,你才能面對所有的批評,所有的困難.然後祂說:你把這一切向聰明通達的人就藏起來,向嬰孩就顯出來.這一句最重要,嬰孩 and \#8230; 維持單純的心你才能理解上帝；維持單純的心你才能夠用真的愛去愛上帝,而且以真的愛去愛其他的人.

這麼多年我要學習的就是:你還能維持你的單純嗎?當你年紀越大,你的地位越高,錢財跟權力的機會越來越多,你能夠維持你的單純嗎?這麼多年來,我都要不斷的反省,你看到一些不重要的人你能夠愛他嗎?這世界上有很多不重要,看上去沒有什麼影響力的人,你會愛他嗎?你會關懷他的需要嗎?當我們看到蓉的時候,你會否因為她不重要而不理她嗎?還是因為她不重要,你更要愛她,更要告訴她你是非常重要的,你是上帝心中的寶貝,你是我們的寶貝.每次我到少數民族那裡,我都告訴他們你是我們的寶貝,我們最愛的就是你們.

維持某一種單純才能夠以真誠來對人,所以一位著名的加拿大卑詩省前省督林思齊,他說他要學習 and \#8230; 當他當了總督後地位很高,他是大有錢的人,但是他學習以真誠,單純來對人,他說「我每見一個人,我先以真誠對他,如果他回報以真誠,我就給他多一分真誠,他若又回以真誠,我就回更高的真誠,他越真誠我就更真誠,如果我這樣對他他不理會,他沒有真誠,你就沒辦法」.但是如果你單純真誠對人,你會得到很多朋友,雖然這個世界很險惡,雖然人很醜惡,但是你還是能夠有單純的心去跟人交朋友,當我每逢遇到困難的時候,我發現很多朋友就出來,當別人在後面無故加害你的時候,無故在你後面攻擊你的時候,總是很多朋友出來為你說話,那就是過去我以真誠看待他們(縱使是不重要的人).所以當我有機會進入到中國的領導階層,在會見他們的前一天,我先見那些最給人不重視的──貧困人.我要學習不要因為能接觸領導人而覺得有多麼的了不起!而是重視那些不重要的人,尊重他們,愛他們,這才是我最大的學習.

當我單純的時候,我才能夠看那個天空 and \#8230; 我小的時候在陽台放風爭,看到太陽下山的時候,夕陽西下的那種雄偉的景象,我心非常震撼,然後天慢慢黑,那時候媽媽把飯菜帶到陽台一同吃,天黑的時候,星星慢慢出來,媽媽說:你看到第一顆星出來你許的願一定成功.我和哥哥馬上說:希望年年考第一,結果沒有成功.可是當看到天上星星出來那種偉大,哥哥問了一個問題:如果我們坐飛機坐火箭飛到太空,什麼時候才到盡頭?媽媽說沒有盡頭的,這個時候我忽然想:什麼叫沒有盡頭?什麼叫宇宙的宏大?什麼叫無限?我更想到這麼宏大的宇宙,如果沒有的時候,原本是什麼?當一切不存在的時候,什麼叫沒有,什麼叫無?然後它後來變成有,無怎麼變成有?我小的時候覺得這個非常奧妙,那是一個單純的心的一個想法.

這個宇宙太偉大了,當我仰觀天上的星星時,我總覺得在深層的宇宙深處一定還有一些完美的東西等待我去追求的,所以我一生都在研究這個課題.然後當我爸爸變了,人生遭遇痛苦的時候,但是我知道宇宙還是很美麗的,舉頭看天,每逢黃昏還是有夕陽,晚上還是有星星,這個宇宙何等的美,我知道一定有完美的真理存在.當我十三歲的時候,爸爸開始有問題帶我到書店,看到一本莊子,我就把莊子買回家,哥哥曾經跟我講過莊子的秋水篇,裡面講到一條河流,越流越大就以為很了不起,看到大海才知道自己是如此的有限,然後它跟大海的神說,俗語說「懂一點點,就以為自己了不起!就是我這種人」.那個河伯一看到海就說海偉大,但海說你不要以為我很偉大,我在宇宙也是很渺小.我覺得這個說法很精彩,我自少就看到宇宙多麼大,人多麼小.當我買了莊子回去第一章打開就是《逍遙游》,不太明白.從前有駒電影叫「青春消遙遊」,就是謝霆鋒這位明星他的爸爸是大明星的時候,很多電影就是他英俊的爸爸和一些很美麗的女明星演這齣「青春消遙遊」.

難道莊子是講這些嗎?一開始就說「北冥有魚,其名為鯤.鯤之大,不知其幾千里也.化而為鳥,其名為鵬.鵬之背,不知其幾千里也；怒而飛,其翼若垂天之雲」聽得懂嗎?莊子一開始就說:北海有條很大的大魚,後來轉化成一隻大鵬鳥,一飛飛到太虛之上,牠的翅膀有雲那麼大 and \#8230;.我當然不完全明白,後來我想到──「有」──什麼叫「有」?心靈能游於宇宙,心靈自由地飛到天上,心靈的飛翔是什麼?那個時候我突然之間想到,雖然人間很多的痛苦,雖然我爸爸有問題,國家民族有很多難題,香港亦很不公平,但是我的心可以飛到天上,當那個大鳥飛到天上,從太虛回看人間是什麼樣子?心靈的飛翔,那個時候我心靈忽然之間,追求一個超過現實最後的真理.後來我讀柏拉圖,柏拉圖講到:為什麼有數學的存在?數學是一個完美的系統,它不是經驗看得到的,但是我們可以用數學建房子,證證明數學是真的,但是它看不見的；它不是經驗的,那麼它的存在在那兒呢?柏拉圖指出,數學存在在一個完美的系統裡面,那麼那個完美的系統是在另外一個完美的宇宙裡面,宇宙有完美的,他說通過數學就知道,宇宙有完美的存在,這是一個很大的訊息.

後來我就追求佛教,十五歲我就打坐,十五歲的心很單純,很清,所以一打坐就進入狀態,一個人到五十歲打坐什麼都弄不出來,五十歲的人都很複雜的,十五歲的心比較清,單純,所以很快開天眼,跟著不單看見其它空間的存在,而且最後有個領悟的經驗,發現宇宙真的是變幻無常,一切都因緣起滅,一切皆空,得到這個領悟的時候,真的很高興,我看透了,一切都變幻無常.所以你看到一個美女當你很喜歡她,但是你想想五十年後,變幻無常,她變成其她樣子,你就不會再有慾望,你這樣不是想通了嗎?所以當我娶我太太的時候我也在想,可能廿年後,她會變成一個又胖又囉嗦的女人.我認識太太在中文大學火車站,有一次等她下火車,那時候火車站後面就是大海,那個年代很浪漫的,太太下車看到我在馬路的對面等她,很高興地跑過來,風吹起她的頭髮非常美麗,不過結婚之後,她頭髮吹不起來,因為結婚後她不會跑過來,她會說:「你過來,洗盤子」.所以我想二十年後,肯定變成不好看的女人,我都看透了,一切皆空,所以有一天她變得難看,我說我早就知道,我已經看透了一切皆空,那麼我覺得我很有智慧.可是奇怪的是過了二十年,太太還是很漂亮,結婚到了三十年,她還是很漂亮,唉!我說這不是佛教的原理,這是基督教的原理,上帝有神蹟出現,忽然之間上帝用神蹟把太太變成三十年都漂亮給我看,我看透一切的變幻是受佛教的領悟.

接著我再有個領悟,就是道家思想.莊子常講「上與造物者遊」──你要飛到天上,跟宇宙造物最後的真理一同飛,一同遊,自由地在宇宙四周的心靈一同飛翔,道家也講到宇宙每處都有規律的存在,忽然之間我又明白.大概十八歲的時候有一個領悟,就是看那個花,樹,草,蟲,動物,當你觀察牠們,會發覺每個東西都不同,因為牠每一個都有牠的「道」在那裡,宇宙是有一種安排的.道家也說,你只要回到像孩子一樣的心,才能夠看透宇宙後面的安排,維持你的單純,要虛心,沒有你自以為是的理論,才能看透宇宙後面的安排,自那個時候忽然之間,我打通了.老子說:「有物混成,先天地生」──天地之前,有一個偉大的安排已經在那裡.但是天地的母親,我也不知道她的名字──「吾不知其名,字之曰道」,就稱為「道」吧!宇宙有一個你不知道的「什麼開始」安排在那裡,那宇宙的安排是什麼?宇宙的安排 and \#8230;當我太太懷孕的時候,我想了一個問題,她肚子裡面的,你怎麼知道是人呢?裡面只是一顆細胞,細胞分裂不一定變成人,他可以變成其它東西,為甚麼它一定變成人呀?因為裡面有個「道」,現在科學說有個DNA,一個後面的基因,有個密碼安排了它是人,但是我想如果那個密碼排錯了,變成豬的DNA,它在裡面分裂變成一隻豬,你怎麼知道它一定變成一個人呢?如果生了一隻豬出來,你怎麼處理牠?吃牠還是不吃牠,成為一個問題呢!最近我女兒也生孩子還是雙胞胎,我想如果它分裂弄錯了密碼,生了兩隻豬,那你吃還是不吃?很奇怪結果生出來都是人,為甚麼?宇宙有他的安排,所以從那裡,我從佛教講到無常變幻講空,我就領悟到宇宙有個偉大的智慧在安排,從單純的心──你回到你的虛跟靜的心,沒有你的理論,只是單純看宇宙萬物的變幻,你就知道宇宙後面有個偉大的東西在那裡.老子說我不知道他叫做什麼名字?叫他做「道」吧!

後來更進一步,大概十八,九歲開始讀儒學,就是儒家的思想,孔孟之道,裡面我更有一個領悟,就是:人怎麼知道宇宙?因為人的心靈就是宇宙的部份,我們不單關切這個花,草,樹木,宇宙自然萬物來了解宇宙,還要反省自己的內心,才能了解宇宙,我的內心就是宇宙部份,而且可能是非常高的一部份,當我反省內心,我發現儒家講:當看到別人的痛苦,為什麼會不安,不忍?孟子講到惻隱之心,看見別人痛苦時,為什麼心中不忍?因為你的內心是有一個良知,一個仁愛的心.我發現儒家講得很對,人是需要有仁愛的心,我心中有仁愛的心,為什麼我看到他們痛苦,我心中不安?不單是我的家庭有痛苦 and \#8230; 我媽媽小的時候帶我看望貧困的人,媽媽是做社會工作,我看到香港有很多貧困的人很痛苦,後來知道世界到處都有打仗也很多痛苦,看見這些痛苦我心中總是不安,他們痛苦我也不安,動物受苦我也不安,那是因為人有仁愛的心.然後儒家講到這個仁愛心,從中庸裡面說:「天命之謂性」──人性這種善是從天而來.那麼那個時候我就知道,天應該是仁愛的,如果宇宙有一個仁愛的天,他是不是會知道我們的苦楚?來安慰我們的憂愁,他會不會跟我們說話呢?孔子說:「天何言哉」──上帝沒有說話,上帝不需要講話,你看萬物這麼偉大,知道有個很偉大在那裡,這個偉大,可能是老子講的那個「道」的那種智慧.也可能是孟子講的仁愛──仁愛的心來自天,天是仁愛.如果天是仁愛,他應該知道我們的苦難,他會不會跟我們說話?在那時候我就開始追求,如果真的有仁愛的天,可能就是中國古人所謂的上帝,因為中國古書 and \#8230; 我在高中畢業那一年,我看詩篇,發現裡面常常提到上帝(但我讀書那個時候中國還未有開放),後來中國開放,我們可以到北京的天壇,那個天壇裡面沒有偶像的,在祈年殿裡只有一個牌子,寫著「皇天上帝」.「皇」是最高的,最高的天有個主宰──皇天上帝.所以從那個時候,我就開始想到:這個宇宙應該有一種最後的,我們祖先也相信的,一個主宰在那裡.祂應該是有情感的,明白人間的悲歡離合,所以我向祂禱告:我是如何能夠知道你?如果真的有上帝,你大概知道人間那麼多苦難,那麼多難題,你有甚麼方法看人類?你能夠跟我溝通嗎?我向祂禱告,我又為人類的苦難向祂禱告.

結果很奇特的我就遇見基督徒,問我星期六下午有沒有空?基督徒一般問這個問題都有陰謀,要帶你去教會,不過我想基督教加天主教是全世界最多人相信的,她一定有道理吧,不會只是笨笨的一批人嘛,那麼林肯是基督徒,華聖頓是基督徒,牛頓是基督徒,孫中山也是基督徒,他們不是笨的,他們大概一定有些偉大東西在那裡.我就去聽聽,他們查聖經,然後是新來賓介紹,介紹我的時候一同站起來唱「我們歡迎你,真歡迎你 and \#8230; 」和我握手,基督徒都很熱情,在我來說熱情還是熱情,我有理性的問題,我立時問:其她宗教 and \#8230;我是研究其她宗教,我修了很多年,基礎差不多,都是殊途同歸,那些基督徒是香港大學的,但是他們沒有讀中國文化,因為他們讀英語學校的,讀法文的,所以他們不懂中國文化,香港有一批這樣的人,然後他們對著我,他們是很好人,但是我一出中國文化他們沒辦法,臉都青了,只能夠回去為我禱告,那個時候我很喜歡虐待基督徒,折磨他們,只要我問一下難題,他們就很痛苦,不過後來慢慢發現,基督徒也挺奇怪,面對這些人你跟他們辯論,他們不一定懂,但是他有一種經驗,他們體會過上帝,因為他們體會過上帝,他們很快樂,他們充滿愛心,也很單純的對人.

我覺得這批人倒是挺好的人,跟我現實上見到的都不一樣,我也感到奇怪,因為我也是通靈界的,我也通各種宗教的境界,只是不通上帝.他們什麼都不通,但是通到頂找到上帝,我奇怪我知道上帝,但不知道該如何上去?因為我知道我內心還有一種醜惡在內面,讓我不能接近上帝,因為上帝是完全的完美,我修了多年表面上是很好,我一點壞事都沒有做,都是好人,但是我的心總會想一點壞的東西,我沒有做,但我會想.我沒有殺人但看到那些混蛋,也想砍他兩刀,我沒有姦淫,但是我可以幻想得很精彩,我知道心裡面的醜惡,所以我沒辦法接近上帝,除非上帝從那邊過來.那麼基督徒他們怎能遇到上帝?難道他們的罪已經洗乾淨嗎?我就決定開始研究,便開始看約翰福音,我最初看「太初有道,道與神同在 and \#8230;」,我說這是道家思想,東西文化都差不多,但再看下面「 and \#8230;道成了肉身,住在我們中間,充充滿滿的有恩典有真理」,這句話怎麼解?「道」是很偉大,宇宙萬物的規律,如果祂有情的,是不是祂會變成一個人來到人間?上帝是全能的,當然祂可以來到人間.所以我又想,如果上帝真的來過人間,這種「道」是非同小可的事情,然後聖經說:耶穌基督是上帝來到人間,宇宙有上帝我覺得是合理的,因為那個時候,科學最新的發展就是證證明宇宙有起源.宇宙如果有起源,一個大科學家Georges,他也曾經說過:我好像看到了上帝.當他最新的人造衛星計劃,研究到宇宙的起源可以找到證證據,後來人問他你為什麼說好像找到上帝?他說:我找到上帝創造宇宙的「指紋」(Finger Prints).因為宇宙一開始的時候,都有非常複雜的設計在那裡,如果沒有這個安排我們今天不可能存在,生命能夠存在,人能夠存在；因為第一秒鐘,所有這些規律都安排好了,看來是有上帝,而且還有人推論出宇宙如果從「無」變成「有」,已經證證明了.「無」──如果真的「無」不可能變成「有」的.我們說一切都自然而來,但是宇宙如果根本連自然都不存在,又怎麼可能自然產生一切呢?所以從「無」變成「有」,且都很有規律,讓生命可以存在,這些生命存在第一秒已安排好,而我們已經存在顯明第一秒已經安排,看來真的是有一個上帝.

從科學來說,大概這個時候根據相信有一位上帝是符合理性的,我們不是說科學證明有上帝,我是說當你信有上帝,你會發現你的信仰跟科學是一致的,是符合理性的.然後從儒家來說,也應該有上帝,因為人的仁愛是來自天的話,宇宙應該是有一個仁愛的真理,但是我就想問:這個真理真的來過人間嗎?耶穌基督就是這個真理?但祂來到人間經歷各種痛苦,上帝為什麼要受苦?祂不是全能的嗎?祂把什麼東西變好就可以了:「上帝呀!你向我顯現就最好,你告訴我明天買什麼股票?歐元前途如何?李旺陽的問題也有沒有解答?(我們見到上帝都想問這些的).但是上帝不是來滿足我們的慾望,如果說你要什麼就給你,我說:「上帝我想吃蝦餃」,上帝就說:「不用到對面；現在把咪高風變成一個蝦餃給你吃下去,是熱的 and \#8230;」,你看上帝是全能的,祂能夠滿足你吧?然後我說:「上帝我想吃鮑魚」,上帝說:「好,你剛才上過洗手間,你到洗手間把蓋子打開,裡面全都變成鮑魚,要用大勺子呢,上面還有鮑魚湯」.上帝當然可以變出任何我們想要的東西,但是這樣子的話,上帝就很膚淺,祂只是滿足我們的需要,上帝不是來以很膚淺的方式來滿足我們需要.

上帝來到人間跟我們共同渡過苦難,當我研究耶穌基督生命才明白:耶穌基督受很多的苦,給人侮辱,鞭打,他又很有上帝能力,用五餅二魚可以使五千人吃飽,但是他又受到人的侮辱鞭打還釘在十字架,為什麼?祂選擇一個痛苦的道路來告訴人類:你受的痛苦我與你共同擔當,我來不是用全能變好這個世界,我來是跟你共同渡過苦難,你受苦你不孤單,我跟你一起走,你的受苦我是與你同行,我跟你一同上路.原來祂受很大的痛苦只為承擔我們的苦難,而且更進一步人類很多的罪惡,我們千百年來人類的罪惡多得不可勝數,我活了數十年發現,醜惡的人比好的人多很多,好人也有很醜惡的,醜惡的好人卻不多,我發現人有很多醜惡,這麼多的罪如果上帝要刑罰人沒有人能逃脫,因為上帝要刑罰人包括你的心想過的,你沒有做壞事,可是你想過也要受到上帝的審判,但是上帝不願意刑罰我們,所以祂通過耶穌基督承受我們的罪,通過祂的痛苦來接受一個刑罰,這個刑罰代替我們受了罪應該受的刑罰,以致我們得到寬恕,得到自由,原來基督教有一個很深的道理,就是人的罪可以得到寬恕.而且我們也要寬恕其他人,這個世界就沒有那麼多衝突了.這樣基督教有個威力,就是通過寬恕來處理問題,不是通過跟你鬥,一直鬥下去,永遠仇恨就下去,現在很多種族衝突都是這樣,但是上帝是以寬恕,與你共同受苦來處理人類兩大問題:「苦難」,「罪惡」問題.

然後祂最後死了,但是祂死而復活,三天之後無數人看到他的復活,顯明第三個人類的大能題:人死後到哪裡?我們都是從某些理論來猜,印度人告訴我們有輪迴.佛教從印度來的,講的是輪迴,是真的有輪迴嗎?我們不知道!如果真的有輪迴就麻煩,這樣人類便得永遠受苦了,因為成佛,成菩薩,成羅漢的很少,大部份人就是輪迴,到了天還是要重新輪迴,因為人總有個「孽」那裡,這個「孽」沒有清理,還是要繼續輪迴,成佛的很少數.但是基督死裡復活,忽然之間把這個問題處理了,祂把人的罪清掉──用佛教的名稱就是:那個孽除掉後,你就再沒有輪迴的問題,你清掉了所有的罪,那麼你所有的罪除掉之後,你就直接進到永生當中.當耶穌復活顯明,人死後有一個永恆的生命,死亡只是一個門,這個門一打開,原來是永恆,而不是不斷的輪迴這種說法.

當人親眼看到復活就知道這個是真的,不只是說的,不只是理論,而且祂復活也顯明祂就是上帝,那個時候我也有思考這是真的嗎?耶穌復活是真的嗎?是不是神話?但是當我看到聖經,保羅在哥林多前書講到(保羅寫此書的時候是公元五十到五十五年),耶穌復活的時候向門徒顯現,跟著也向五百多人顯現,五百多人少部份是睡了死掉,大部份還在,看過耶穌復活的有五百多人目擊的一個事實,不單是人造出來的神話,所以一開始傳說祂復活這個說法,是eyewitness tradition(目擊者傳說),不是oral tradition(口述傳說),不僅僅是口述,是有人看見的,五百人現在還在,寫哥林多前書時候在公元五十到五十五年,在歷史上是可靠的,保羅寫這封信的時候離耶穌復活才二十多年,二十多年的事我們一定記得很清楚.好像香港最近紀念六四,我們實際回憶還是很清楚,二十多年記憶猶新.所以當保羅寫這封信很多人都親眼看見耶穌復活的,而且他說,最後耶穌基督還看他顯現,他也親眼看見耶穌基督復活,否則我不用浪費時間做一個笨蛋到處傳福音,保羅是一位很有地位,學問的人,他也是羅馬人,他根本地位很高,他不需要自己受苦去傳福音,就是他親眼看見復活的耶穌向他出現,他不能不一生傳這個事實.

然後當我繼續研究這方面,聖經記載那個年代事情時──耶路撒冷好像有兩個水池,一個畢士大池,一個西羅亞池,這兩個水池的功能基本都在公元七十年被羅馬人給毀掉了,整個耶路撒冷沒有了,水池也找不到.但考古學在1888年找到畢士大池,2005年找到西羅亞池,看來寫聖經的人都是當年那個年代看過這個環境的人,而不是後來的人只聽人傳而隨便寫的,所以當我研究聖經的時候,我感到這不是簡單的,只是後人不相信就說這都是傳說罷了.當年基督教能夠出現,一定是有人看到一些很重要的東西,所以耶穌基督復活.如果耶穌基督祂沒有復活,祂出來傳道只有三年半,大概也不會有什麼影響吧?這個老總給抓啦,甚麼人都跑掉.而且那個時候傳福音還是要被殺害的,如是假的誰願意?就是當年的人親眼看到,所以他們堅持被人殺害,迫害也堅持下來,教會才一直存到今天.

再者,看看那些基督徒他們都體會過基督,所以他們也發神經向人傳福音,這個時候我發現基督教很合理,只差一點就是我沒有相信,但是多年來我都堅持惟持單純一點像孩子一樣的心,我希望我真的能夠遇到上帝,所以我非常真誠的去探索基督教,最後我發現原來我願意相信,但是我也沒有很深的經驗,我是理性上覺得祂可信,我還沒有很深的經驗,那個時候基督徒都覺得我怪怪的,過去我是個怪人,我讀大學是穿長袍的,還拿個扇子,然後扇子上寫了很多詩句,哲學古人的句,下午也打著把油紙傘,常常在草地打坐,我常常穿長袍,晚上的時候有人走過,常把他們嚇了一跳,這樣一個打坐的人,怎麼忽然之間信耶穌?基督徒都感到很奇怪,後來邀約我參加一些查經,結果我在查經裡講佛學的解經法,他們也嚇了一跳.

後來就邀我參加一個夏令會,夏令會我才明白不單有單純還有一點,就是畏有「委身」.原來接受上帝不只是普普通通我信了些東西,然後我得到永生,不是為了利益,而是要付出你生命代價的.原來從前因信稱義也不是這麼簡單,從前相信的時候,是要付代價的,因為人要殺死你的,基本都會受到許多的攻擊.當我參加了夏令會我明白了,聖經說:「將你身體獻上當作活祭,是聖潔的事,是神所喜悅的事」,你要把身體獻上去,不只靈魂,你的身體,你具體的人也要獻出來成為一個活的祭,像一個祭獻給上帝,活活的我活著就是一個祭,把上帝的火燃燒,那個時候在夏令會查經,帶領查經的那個人叫溫偉耀,大概都有人知道他是著名的講員,也是一位著名的神學家,他那時候還是一位年青人,不過不是特意派他來對付我的,因為我問題特別多.他說:一個少年官來見耶穌,他說我怎可以得到永生?耶穌說你把你的財產買掉捐給窮人來跟從我,少年做不到很憂愁的離開.我聽了我覺得這個容易,我做得到,把我的財產奉獻,可以,因為我沒有財產,這個很容易呢!但溫偉耀說:不是這個意思!耶穌的意思是說,你心裡抓著最重要的東西,你能夠把它放下來嗎?你心裡面抓著最重要的是什麼?對那個少年官是他的財產,對你呢?那個時候我忽然間明白了,在我心中上帝告訴我:我把我的學問抓得很緊,我自以為了不起,我從十三歲讀莊子,一直讀中國哲學,而且打坐了四年多,我以為自己很了不起,也是我唯一可以跨耀的,上帝說:你願意為我放下這些嗎?這個很難其它都可以,甚麼叫「委身」?就是把你最不願意放下來的,放下來.那個時候上帝在我心中作工,終於到最後一個晚上,我還在思考如果真的有上帝,你當然應該放下一切的,想到最後,我遇到溫偉耀,他也在想他的前途,我就問他:「我都說信了上帝,信了耶穌,為甚麼我還未有經驗你們那種喜樂?你的體會是很真的嗎?」「對!我體會上帝是很真的!」「那為什麼我沒有你們那個體會?」「你越想得自己多,就越想得神小,你越想自己,越自我中心,就越覺得上帝小.」講完那句話,他就學了我那種沉重的方法──很智慧的轉身而去,沒有再跟我講話──你越想得自己多,就越想得上帝小,因為我是自我中心,我信上帝也是為了滿足自己,現在我終於明白想通了.

如果真有上帝,祂對一切有祂的使命,並交付給我們,有祂的安排,而且祂的安排一定有祂的祝福,我忽然間明白我就願意.第二天早晨,我就對著初升的太陽,對著天地,對著上帝說:「我願意放下一切跟從你,我們偉大的天父,我的學問,我的知識,智慧,自大全都放下來,我就一生跟從你」,就在那個時刻,忽然間,我本來心中有一道圍牆圍著我自己,我自以為很了不起看不起其他人,瞬間那個圍牆溶解了,上帝就如風的進來,上帝進來的時候那種喜樂,我明白了,那種喜樂那種宇宙的親情,聖靈一進來那種美,那種善,我明白了,我心什麼都倒空了!這不就是佛教追求那個空的境界嗎?人追求那個空我現在可以得到,信耶穌可以得到.然後道家講宇宙那個大道,我信基督也馬上體會,天地萬物是有祂的安排.然後那孔孟之道講的,那種善,那種仁愛,我信基督時候那個仁愛真的出來,我的罪給拿去洗乾淨,那一刻仁愛重新出來,我就明白儒道佛最後的追求,在信基督時都得到了.但是更重要的上帝的親情跟慈愛一進來的那一刻,我就明白:為什麼基督徒總是發神經整天叫人信基督,我也發神經了,我明白了!我發神經直到今天,今天不發神經為什麼來這裡?本來今天難得有一天不用講道,遇到伍牧師他把我拉來,因為講道有錢?我告訴你沒什麼錢的,你靠講道賺錢肯定完蛋,肯定不會富足.

是因為發神經,讓我們覺得我們中華民族受苦太多,彼此傷害太多,我們鬥爭太多,還不夠嗎?我們需要寬恕,需要愛,需要仁愛,讓更多中國人明白這種仁愛,體會過上帝的仁愛,你才會做一個有道德的人,所以在那個時刻我就把生命獻出去.我二十多歲就獻出,二十多歲什麼都沒有,所以獻出很容易的；到三十多歲難一點,三十多歲自己充滿前途,很多人很重視,你以為你很了不起,上帝要從新把我拉到祂的國度裡面,那個時候我終於再放棄一切跟從祂,那個時候我三十歲,我覺得我應該可以當一個虛偽的人,我一生都不看色情電影,現在年紀大應該可以看看吧,為什麼當基督徒不可以做這些事,那個時候在美國沒有人認得我,在香港還怕有人認得我,到美國讀書的時候我就不理了,就去做壞事,表面是基督徒就可以了,後來上帝又從新把我重新帶回來,晚上讓我醒來問我:你在做什麼?你到那裡去?我說我在為名為利,上帝說你願意放下再跟從我嗎?最後我說好,以後我得名或利,要榮耀你我才要,不能榮耀你的我就不要.那一次是三十多歲的奉獻.四十多歲更麻煩,四十多歲的時候,你地位越高,害你的人越厲害,你遇到一些高手害你的時候你心中很痛苦,想用自我能力來報仇,我完全有能力把他打下來,但上帝教導我要寬恕,結果通過這個過程,我痛過這個寬恕之後,心中的忿怒去掉我才看到前面,上帝把中國前途放在我前面.

九十年代1993年中國很艱難,全球在罵她,你跟普通人一樣的罵她嗎?你能夠愛他們嗎?包括你最恨的共產黨能愛他們嗎?上帝是愛他們,這個我真的很大學習,真的很難,從小聽父母講的都是共產黨不好 and \#8230;怎麼去愛?實在很難!但這個不是愛國的愛,是基督的愛,無盡的愛,這個愛能夠醫治他們過去很多很多的苦難,能夠帶來前途,而且愛能夠讓中華民族更美好,過去我們受苦太多,民族受傷害太深,先從愛開始,公義才能夠出來.

上帝呼召我,我就終於順服開始文化更新研究中心,當我開始的時候還要放棄剛剛開始的一份工作,一百二十萬年薪聘請我的,九十年代的一百二十萬就等同現在三,四百萬的年薪,這個是最難的!一個是錢,一個是你要愛你本來不愛的,不過那個時候是一個最大的突破,終於我選擇了為十三億同胞能夠重建他們的尊嚴,重建他們一個彼此相愛的關心,我放棄了那一百多萬的薪酬.愛不能單單憑口說,乃是要行出來.我就回去中國十多年,現在已十九年了,最初是大學講課,後來十年我就跑到農村到山上理解貧困人的需要,縱使你改不了整體的中國,起碼你幫助了一批人讀書,讓他們有前途,我們前後幫過六千多人讀書,最基本我愛這些人,我不能把中國制度立刻改變,我也沒有這個能力,但最少我幫過六千多人讀書,我在大學推動過仁愛的思想,然後跟官員很多接觸,能夠把民間的苦難告訴他們,希望他們也能夠有仁愛的心,而且有公義,公平尊重人的心來帶領人民.就這樣子我過了幾十年到現在,從單純的心,從那個時候開始,我也盡量維持自己要真誠一點,不講假話,不騙人,不拿好處 and \#8230; 學習不拿好處,我回中國十多年沒有做過一單生意,沒有賺過一毛錢.我們只有在大學講課有點錢,不在中國營利,維持這個基本原則,你不拿好處,你講的話才真,然後維持委身,你向上帝完全獻出你自己,上帝的使命來臨,你就會進入自己國家獻出自己,希望國家最終能夠變得公義,公平,更民主,更法治,這些是中國自己要求的,我們希望讓她變得更走上這條路,我們以完全的愛來做這件事情.結果我覺得中國還是在改變當中,雖然每次改變都有些不好的事情被發現,但整體不好事情還是會改好的,所以我們維持多年對中國堅持的關懷和愛,對他們貧困的人關懷,理解他們的情況,也向高層的人反映,我們能力就到這裡,我們沒辦法把中國整個改變,但是你能夠做一點就一點,這是我們努力的方向.

所以獻出了生命,我沒有高的薪酬,但是我活得很快樂,而且每年我們需要的金錢上帝都足夠地供應我們機構,也不會有盈餘,每年剛好,我就明白神真的與我們同在,祂的祝福也與我們同在,也願意我們學習像孩子一樣單純在上帝面前,也學習獻出自己的生命,來把上帝所給的使命給活出來.無盡的喜樂就在裡面,別人的批評就變成我的祝福,為義受迫逼的人有福了,越有批評,越罵我的,我就越高興,因為我得的福氣是大的,因為我知道我不是靠別人來評價,乃是靠上帝來評價,當我走上歷史上這條路,我就非常清楚我怎麼走,也願意我們每個人都能夠獻出他的生命,為了使命活下去的心,讓我們活著,使人能夠得到更多的愛和快樂,讓人得到上帝的喜樂,這個世界就變得更美好,我們一同低頭禱告:

親愛的天父!在這個主日的清晨,天上又下雨,我們紀念在中國南方有三年多的乾旱,有雨對他們是好的,當大地被雨水所滋潤的時候,我們紀念遠方有很多很多悲慘的事情,在敘利亞很多人死亡,在埃及有很多動亂,我們紀念中國還有很多貧困的人是需要人關懷,尊嚴,我們能力很有限,主求你讓我們以我們有限的能力,行出無限的愛,懇求聖靈在這個時刻臨到我們當中,復興教會,感動我們的心,讓我們以無盡的愛,以有限的人行無盡的愛來對其他的人,讓這個世界上,社會裡面,在遠方山上很多貧困的人,在香港很多給人看不起的人,讓我們心中有關懷,賜我們有能力把這個世界帶來更大的慈愛和公義,你賜我能力減少人間的鬥爭,殘忍,暴力和霸權,讓我們以一個單純仁愛的心,像你一樣的柔和謙卑,讓我們背負的重擔,勞苦擔重擔來到你面前,把重擔放在你面前,也求你讓我們維持一點單純,特別當我們年紀越大,我們越複雜的時候,叫我們以單純的心,真誠的心來對人,也讓我們能夠獻出自己,為你而活甚至於為你而死,讓我們的存在能帶給別人祝福,而不是帶給別人痛苦,咒詛,讓我們活得更美麗,我們感謝你,這樣禱告奉主耶穌之名而求,阿們!



\newpage

\section{豁免年}
\label{sec:160}
\textbf{陳志強牧師 - 豁免年}
\newline
\newline
\hyperref[sec:159]{< < < PREV SERMON < < <}
~
\hyperref[sec:toc]{[返目錄]}
~
\hyperref[sec:161]{> > > NEXT SERMON > > >}


日期:2012年8月26日

證道:陳志強牧師

上年和今年被受注意的問題,就是債務問題.對很多人來說,這是很重要的問題.現代的經濟體系中很難避免借貸問題,不論個人,公司,以至國家,都牽涉這些問題.借貸的情況不單是關乎窮人,在現今社會,有錢人便可能需要有更大的借貸.因此,不單窮的國家有問題,連富有的國家也有這些情況.借貸不能避免,它跟資源運用很有關係,有資源才有借貸活動的出現.聖經時代的借貸活動跟今天的情況不同,舊約時以農業為主,最重的資源是土地和牲畜,所以借貸是窮人才有需要,特別在遇到天災人禍時影響了他們的生產,他們需要借貸,而能夠提供借貸的是有錢人,包括土地和種子等.

聖經申命記提及一個重要日子,就是安息年的設立.安息年常令我們想及不用工作,現代人都希望不用工作,因為工作太辛苦,所以希望不用工作,退休,有收入但不用工作,所以安息年使我們想到不用工作.有什麼工作是工作六年便能放假一年?只有牧師,但亦不是所有牧師都有安息年,像我服侍的教會卻是沒有安息年的制度,但有些教會有,那真好,可以工作六年便能放假一年,但我們卻忽略了聖經中提及的安息年並不單指不工作和耕田,最重要是指讓土地休息,不是不做工,並不是那一年去遊玩,去環遊世界,那一年有很多事情要做,包括維修家居等很多事情可處理,但不可以耕田,最重要是讓土地休息,不是人休息.在經文中把安息年也稱為「豁免年」,是處理債務問題.豁免年是讓我們知道世界上有些需要的人,聖經提到這世上總有「痛苦窮乏的弟兄」,總有人有困難,每個人的生活和遭遇不同,在沒辦法計算下我們可能會出現天災人禍,例如疾病,有些人一出生便有不能治癒的疾病,影響他們的謀生和生活,當然有些人的窮乏與他們的抉擇有關(有人欠債因為亂花錢──今天報章報導前政務司司長許仕人欠債七千多萬,這肯定是他的抉擇,並不是什麼天災人禍所造成).

缺乏是一個常有的現象.耶穌說在地上常有窮人,缺乏是歷代歷世都會出現的現象,物質上缺乏的人在所有時代都會出現,這個時代我們是富有的,在廿一世紀中,普遍的人是富有了,古代的社會就只有貴族才不缺乏,大部份的人卻活在缺乏的狀態.在廿一世紀的這裡,我們又比其他人富有,研究經濟的人可以告訴我們香港人的購買力在全界的排位,雖然我們比較富有,但在當中是否沒有窮人?我們能否自己界定自己不是窮人?這世界中人是缺乏,是由人的罪性產生,不單因為個人選擇和社會問題,還有在我們生命中的很多缺乏.不一定是物質,在我們的心靈裡也有很多缺乏,很多人生活中沒有平安,常常活在恐懼中,他們的生活中沒有愛,沒辦法跟他人建立正常關係,他們覺得自己沒有用,覺得活在世上沒有價值和意義,人總是有很多的缺乏.從前人因為貧窮就比較留意物質的缺乏,但當我們生活進步時,卻又發現我們有更多的缺乏,這些缺乏的人就在我們身邊圍繞我們,只要看看身邊的,了解一下,物質的缺乏可從他所吃的或是物質指標便可知道他的營養標準,身體狀況,而心靈的缺乏是需要我們睜開眼,了解他的內心思想,他所渴望,恐懼,擔心的便知道他們所缺乏的是什麼.但很多時候我們以為別人都生活得可以的,因為他們三餐溫飽,收入足夠為生,便忽略了他們的需要,我們因為自己有工作有收入而忘記了自己的需要,看著銀行存摺的保障忘記,也忘了自己也是心靈貧乏的人,只是我們沒有去留心自己的缺乏而已.豁免年讓我們留意需要的問題,這世界是有需要的；豁免年教導我們去幫助有需要的人,豁免年的制度是七年的最後一年,要將別人的債務豁免,借錢只借六年,到了安息年便不用歸還可獲豁免.

聖經還有一個禧年的制度,土地需要歸還(香港從前被賣九十九年,很久之前的事吧!)聖經中也有提及賣地是四十九年,禧年一到,所有土地便要物歸原主.豁免債務的意義是幫助某些人,雖然有些人有能力償還,但有些面對天災,或者在人生中有控制不了的災難,他便要一生負債.現代因債務屬於自己,可以宣佈破產,但古時的債務卻有可能使到後裔不能翻身之日,若債務太重,最容易的出路就是被賣作奴隸,這樣便延續到下一代.禧年又包括了釋放奴隸的意義.

今天所提的豁免年,對當時大部份以色列人面對借貸是很有幫助的.許多工種和經濟活動需要借貸來運作,聖經中提及的豁免年是借貸背後的一種幫助,借貸是幫助那些沒能力生活的人,所以到了第七年仍然沒能力償還的,聖經提出要給予豁免,不要再拖延下去了,因為他根本就沒月能力來償還.對於放債的人,若以幫助別人的心,卻又認別人對自己虧欠,也不是一件舒服的事.就如請別人吃飯,若常覺得別人欠了自己一餐飯,心裡就很不快樂.在我們幫忙別人時,心裡總想別人回報我們,特別中國人以恩報恩,好像欠了別人人情就要回報,或者對人好的時候,也期望別人回報,因而積存了很多的「人情債」.豁免讓我們知道要放開,有些事情不用去記念和爭取,對人的幫助,到了一個時限就要放開.為什麼一開始不就送給他?卻要在豁免年處理?我相信不要「濫用」這份心意或慈善,聖經提到的是到安息年時要「鬆手」,即是放手.抓緊對雙方也沒幫助,人都是自私的,我們都顧著自己的益處,我們把自己的事情都捉得很緊,覺得這些都是屬我們的,不管是東西或對人的好處,我們總希望有回報,特別是現代社會我們都講求投資回報,我們做很多事甚至結交朋友也期望回報,像是作為一種的準備.

聖經提醒我們不要抓得太緊,因為我們永遠都不會滿足.我們抓得緊是因為沒有安全感,就好像孩子緊緊的抓緊大人一樣,無論我們對人,對事,甚至善事</font><font>都是抓得太緊,因為我們感到不安全,心裡沒有滿足和平安.豁免年提醒我們要放手,把我們有的與人分享,分享後不計較回報.同時,別人也可以與我們分享,他們分享便欣然接受,並不需要覺得欠了人情,這是人與人之間的關係.上帝愛我們,差祂的兒子來救贖我們,為我們而死,上帝並不希望我們覺得欠了祂人情,聖經中說我們不是欠債而活,保羅在羅馬書提及上帝給我們的是真自由,所以我們也應給別人自由,沒有人欠了我們什麼,我們有的是從上帝而來,所以我們願意與人分享,不,一定是金錢.我們可以與人分享我們的心和愛,去欣賞和支持身邊的人,讓他們得到滿足.滿足不是由金錢而來,耶穌說:你若要人怎樣待你,你也要怎樣待人.我們知道自己有缺乏,所以也希望得到滿足,我們也知道身邊的人也有所缺乏,我們是可以給予滿足,問題是其他人在我們心中的地位,抑或我們目中無人,只看見自己,我們看到自己的缺乏和需要,並很努力去滿足自己的需要,我們的傳統教導我們不要去求別人,所以有些很窮困的人也不取綜綬,我們要自食其力,不能與他人扯上關係,我們不願意幫忙和被幫助.問題在我們心中沒有他人,聖經說:神愛世人.我們也會說愛人,但很空泛,如果你能說出一個你所愛的人的名字就具體得多了.我們會說愛人,但卻不愛認識的人,因為有些人愛不下.不認識的人我們可能會愛,就如捐獻,覺得已幫上了,心感滿足,但卻不知道愛的是誰,他的需要是否只是金錢?我們心中有無看重一些人,是能使我們鬆開手把我們所有的分享,除非我們認識他們,把他們視如自己般看待,與他分享.

我們會擔心幫了他人之後,自己又如何?我們會先看自己的需要,看自己重要.聖經中提到很多祝福:在耶和華你神所賜你為業的地上,耶和華必大大賜福與你；因為耶和華你的神必照他所應許你的賜福與你.我們相信上帝是賜福我們的.在舊約中能借貸的人已被賜福,因為他能夠借錢和幫助別人已是很有福氣,能夠愛人已是很有福氣的,因為普通人做不來的.我們有福氣,因為神的愛充滿我們,我們學習認識身邊的人.若我們能鬆開手放下時,上帝有更多的福氣臨到我們,我們相信神在暗中察看,喜悅和記念我們.其他人或許並不知道,但在幫助他們,藉祈禱把他們的需要帶到神面前時,這種心意就只有神知道並喜悅,當我們願意鬆開手與人分享,上帝必定賜福,雖然我們並不知道祂使我們心裡得到滿足的方法.耶穌說,為自己賺取生命,會失去生命；為天國的緣故放棄生命的,可得著生命.神跟我們的關係是如此「詭秘」,我們發現我們越想抓緊,越抓不到.就如用手抓水,把手打開就能抓到,反而用手去抓就一滴也抓不到,當我們願意付出,神會記念,我們對神有沒有信心呢?沒有的話,向神祈求,求神賜予信心給我們.

當我們相信,願意付出和放手時,神的祝福便會臨到,上帝會幫助我們,我們的心靈會更滿足,生活更快樂,因為我們知道上帝充滿我們,不靠我們所擁有的和自己所努力的,乃是靠神自己來充滿我們.

讓我們一同禱告:主耶穌謝謝你,你來到世上,為我們捨棄一切,求你幫助我們能體會你的愛,體會你的付出,讓我們在這世界裡,面對許多人的時候,讓我們學習認識他們,讓我們學習了解他們的真正需要,以致我們願意放手,在我們的豐富裡與他們分享,幫助他們,而這幫助是不要求回報,或得到某些特別的價值,這分享是因為體會耶穌的心腸,以致我們相信你要幫助我們,滿足我們,給我們力量,幫助我們在這眾多需要的世界裡,讓我們成為你流通的管子,把你的愛向更多人分享,謝謝你聽我們祈禱,奉耶穌基督的名求,阿們.



\newpage

\section{躑躅邊城}
\label{sec:161}
\textbf{梁元生教授 - 躑躅邊城}
\newline
\newline
\hyperref[sec:160]{< < < PREV SERMON < < <}
~
\hyperref[sec:toc]{[返目錄]}
~
\hyperref[sec:162]{> > > NEXT SERMON > > >}


使徒行傳 16:1-10
\newline
\begin{longtable}{cl}
\hline
\hline
章節 & 經文 (和合本修訂版)\\
\hline
16:1 & \begin{tabularx}{0.7\textwidth}{X} 後來，保羅來到特庇，又到路司得。在那裡有一個門徒，名叫提摩太，是信主的猶太婦人的兒子，他父親卻是希臘人。 \end{tabularx} \\ \\ \relax
16:2 & \begin{tabularx}{0.7\textwidth}{X} 路司得和以哥念的弟兄都稱讚他。 \end{tabularx} \\ \\ \relax
16:3 & \begin{tabularx}{0.7\textwidth}{X} 保羅要帶他同去，只因那些地方的猶太人都知道他父親是希臘人，就給他行了割禮。 \end{tabularx} \\ \\ \relax
16:4 & \begin{tabularx}{0.7\textwidth}{X} 他們經過各城，把耶路撒冷使徒和長老所決定的規條交給門徒遵守。 \end{tabularx} \\ \\ \relax
16:5 & \begin{tabularx}{0.7\textwidth}{X} 於是眾教會信心越發堅固，人數天天增加。 \end{tabularx} \\ \\ \relax
16:6 & \begin{tabularx}{0.7\textwidth}{X} 因為聖靈禁止他們在亞細亞講道，他們就經過弗呂家、加拉太一帶地方。 \end{tabularx} \\ \\ \relax
16:7 & \begin{tabularx}{0.7\textwidth}{X} 到了每西亞的邊界，他們想要往庇推尼去，耶穌的靈卻不許。 \end{tabularx} \\ \\ \relax
16:8 & \begin{tabularx}{0.7\textwidth}{X} 他們就越過每西亞，下特羅亞去。 \end{tabularx} \\ \\ \relax
16:9 & \begin{tabularx}{0.7\textwidth}{X} 夜間，有異象向保羅顯現。有一個馬其頓人站著求他說：「請你過來，到馬其頓來幫助我們！」 \end{tabularx} \\ \\ \relax
16:10 & \begin{tabularx}{0.7\textwidth}{X} 保羅既看見這異象，我們就立即設法往馬其頓去，認為神呼召我們傳福音給那裡的人。 \end{tabularx} \\ \\
[1ex]
\hline
\hline
\end{longtable}
日期:2012年9月2日

主題:躑躅邊城

證道:梁元生教授

經文:使徒行傳16:1-10

``保羅來到特庇,又到路司得；在那裡有一個門徒,名叫提摩太,失信主之猶太婦人的兒子,他父親卻是希利尼人.路司得和以哥念的弟兄,都稱讚他.保羅要帶他同去,只因那些地方的猶太人,都知道他父親是希利尼人,就給他行了割禮.他們經過各城,把耶路撒冷使徒和長老所定的槼條,交給門徒遵守.於是眾教會信心愈發堅固,人數天天加增.

聖靈既然禁止他們在亞西亞講道,他們就經過弗呂家,加拉太一帶地方,到了每西亞的邊界,他們想要往庇推尼去,耶穌的靈卻不許；他們就越過每西亞,下到特羅亞去.在在夜間有異象現與保羅,有一個馬其頓人,站著求他說:請你過到馬其頓來幫助我們.保羅既看見這異象,我們隨即想要往馬其頓去,以為神召我們傳福音給那裡的人聽..."

一,引言:人生的驛旅

人生如客旅.Life is a journey!  我們的一生就好像幾十年的旅程.有些人的旅途很長,要走八,九十年,甚至過百年；也有些人的人生旅程較短,只有三數十年,甚至更短.然而,人生的旅途也並不是長的就好,短的就差；基督和佛祖,在人間世不過短短三十多年,但所成就的是千古的功業；而有些人的人生旅途漫長而顛簸,舉步維艱,一步一腳印,一步一艱難；走得辛苦,並不值得羡慕.人生的旅途,除了有長短之分,也有動靜之別.有些人一生中到處移動,流徙四方,不願長居一處；這些人喜歡流浪,喜歡旅行,到不同的地方,穿州過境,迎接新事物,接受新環境的挑戰,人在旅途,才覺得人生可貴.但也有人不太喜歡人在旅途的生活,所以不願到處走動,而只願長期停留於一個地方,生於斯土,長於斯土,並歸於斯土.他的一生,不願意事常移動,而樂意長期停留於一個地方:最好是出生之地也是成長之地,甚至也是終老之地,越少移動越好,因為移動帶來轉變,也帶來適應；他的comfort zone,就是他最熟悉的環境.

在初期教會的信徒中,有些人是長期定居一城的,離尤其是住在耶路撒冷的許多信徒,他們長期定居聖城,雖然常常接待每年來自不同地方到耶路撒冷來朝聖的人,但他們自己卻不太願意到處走動；但也有喜歡到處去的人,例如保羅.有些人常在驛途,可能是從商,從軍,或是出於喜愛流浪的個性,但保羅卻是為了傳揚福音的緣故.

起初他的腳跡只在中亞西亞的猶太人世界,往返于耶路撒冷,安提阿,以哥念,加拉太和以弗所之間,很少踏足馬其頓和希臘等外邦人的世界.

在今天所讀經文當中,我們看到一個突破:跨越邊界(border crossing)的突破.就是保羅決定踏出他所熟悉的地域,來到一個從未經驗過的新的地方,那就是馬其頓.馬其頓的呼聲,或馬其頓的異像,這是我們常常聽到的宣教的信息.今天我不准備以此為主題,今天我要講的只能是保羅往馬其頓宣教的"前傳",主要重點集中在保羅往馬其頓之前在每西亞邊界停留的那段短暫的日子.二,來到了邊城 (We"ve come to a border town)

在使徒行傳十五六章的開端,我們看到保羅,西拉和提摩太在以弗所,加拉太等地方遊行佈道之後,來到了每西亞.聖經告訴我們,那裡是邊界.過了這個地方,就是另外一個羅馬帝國的管轄區.大致來說,在羅馬大帝國版圖覆蓋之下,各個地區,除了管轄的人有所不同,連地方上的居民,種族,文化和生活習慣也有所不同.換言之,馬其頓代表著的是一個外邦人合異文化的世界.保羅和提摩太,來到每西亞的邊界上,他們在邊境的一個城市城市稍作停留.讓我們來思想一下:他們到了邊城.他們在這裡做些甚麼?有甚麼樣的經歷?這是今天第一個我想和大家思考的問題.

首先,我想大家注意的是保羅的自覺性(self-awareness).保羅清楚知道自己到了邊城.他明白自己的處境,也知道自己需要面對一個重要的抉擇,那就是:往前走,跨過邊境,抑或掉轉頭,找回頭的路.邊城,就是一個需要抉擇的時候:是跨越邊界,或退後不前?到了邊城,就需要抉擇.人生的旅途上,有規範,也有邊界.我們到了邊城,那時,我們需要作出選擇和決定.Its time for decision,或者說It"s the deciding moment.問題是:很多時候,我們走人生的路,一路漂泊,也不知道人在那裡,身在何處,不知道自己的處境；換言之,我們到了邊城,卻不知道到了邊界,不知道應該要做的抉擇.所以我想借這個故事先問問大家:在你的人生旅途上,你知否當前身在何處?在你的人生旅途上,你是否來到了邊城?你有沒有需要做前進或後退的決定?

我來到了邊城!I"ve come to the border town.這是一個對自身處境有所認識,有所了解之後的的覺悟.It"s a sense of awareness, self-awareness.是一種自覺性. 如果你對自身處境一無所知,一無所感,那些發生在保羅身上的事情,當然也不會發生在你的身上.

保羅,西拉和提摩太來到了每西亞的邊界,他們都清楚身在邊城.聖經告訴我們,他們有這樣的經歷:首先,16章7節說:"他們想往庇推尼去,但神的靈卻不許."然後,第8,9節說:"他們就越過每西亞,下到特羅亞去.在夜間有異象現與保羅,有一個馬其頓人,站著求他說:請你過到馬其頓來幫助我們."

以下我想和大家仔細思想一下保羅在邊城的那段日子的經驗,特別是在特羅亞城裡的那一夜.

我們從聖經看到,保羅來到邊城時,第一個反應是不要越境,而要轉往他處.他們想要去的地方是庇推尼.我查了一下地圖,發現庇推尼是在每西亞以北的地方,同屬於小亞細亞,但馬其頓和腓立比,則在海的西邊,現在是屬於歐洲的範圍.換言之,保羅來到歐亞的邊境.他初時不要渡海到歐洲,而願北行到庇推尼,即黑海的南岸.就在這時,來自上帝的干擾,讓他改變了主意.聖經沒有把細節說得清楚,只說"(他們)到了每西亞的邊界,他們想要往庇推尼去,耶穌的靈卻不許."(7節)

你想這樣做,但耶穌卻不許.這是保羅在邊城的第一種經驗--禁制的,拂逆的,負面的經驗.

來到人生邊城,你不能順著自己的意思走,你會聽到上帝的靈向你說不.Do you hear God saying ``No"to you? 上帝的干擾(God"s intervention)不是很舒服的經驗.其實,不要說耶穌說不,任何人向我們說不,也都是難受的經驗.

躑躅邊城,第一種經驗,是聽到上帝說不!

躑躅邊城,第二種經驗,是看到異象.對保羅來說,是馬其頓的異象.

聖經說:保羅他們越過每西亞,到了特羅亞.特羅亞更是頻海的邊城了,再往前行,就要越境,到達羅馬人在馬其頓建立的管治區了.第9節這樣說:"在夜間有異象現與保羅:有一個馬其頓人,站著求他說:請你過到馬其頓來幫助我們.保羅既看見這異象,我們隨即向要往馬其頓去,以為神召我們傳福音給那裡的人聽."

一個邊城的長夜,那被上帝說不之後的難過仍然鬱結在胸,保羅是肯定難以安寢入眠的了.就在充滿鬱悶,彷徨和不安的時候,躑躅邊城,保羅看到了異象:跨越邊界,到馬其頓去!

聽到神說不,他們就不去；看到神使者說去,他們就"隨即往馬其頓去".保羅在每西亞及特羅亞的屬靈經驗,很值得我們細心體會和學習.當我們細心聆聽之時,我們會聽到耶穌的聲音,引導我們走前路；當我們睜開眼睛之時,我們會看到異象,幫助我們鼓起勇氣,跨過邊界,接受新事物和新挑戰.

三,跨過了邊界(after border crossing, what"s next?)

一直以來,解經的學者和講道的牧師講道"馬其頓的異象"或"馬其頓的呼聲"這段經文時,所注目和強調的是向外邦人傳道的這一點.今天我想從另外的一些角度來看這個故事.它們共同的背景是"跨越邊界".一般而言,保羅王馬其頓去,象徵著跨越國界和跨越文化的樊籬,把福音帶出猶太人之外的世界.但在我看來.可能還有更深一層和更廣闊的意義.

讓我們來看看保羅一眾離開特羅亞,進入馬其頓境之後的幾件事:

第一,    他們在腓立比遇見一個賣紫色布的婦人.(跨越性別的樊籬)

由十六章十一節起,就記載著保羅坐船離開特羅亞,進入馬其頓境,第一個站是撒摩特喇,但沒有停留,就經過尼亞波利,來到腓立比.12節說:腓立比是馬其頓的頭一個大城,也是羅馬的駐防城.他們在這裡住了好幾天.就在這裡,他們遇到了幾件事,也是我要拿來和大家討論的事.首先是賣紫色布的婦人.十四節那裡說:有一個賣紫色布疋的婦人,素來敬拜神；她聽見了,主就開導她的心,叫她留心聽保羅所講的話.她和她一家因而都領了洗.

這是保羅在馬其頓的第一顆果子.婦人呂底亞,商人呂底亞.保羅跨越邊界,在馬其頓接觸到的異邦人,是一個做生意的女人.這一定是在保羅意料之外的.怎麼這樣說呢?因為在異象中保羅所看見的馬其頓人,他站著求他說:請你來馬其頓幫助我們.看著段經文的前文後理,我們推斷他很可能是個男人.見到呂底亞,是一個沒有料到的事.同時,這也反映出保羅到馬其頓,不單使福音跨過了政治上的樊籬,也跨過了性別的樊籬.

第二,    他們在同一個地方遇到一個被鬼附身的婦女.(跨越階級的樊籬)

十六節說:他們往禱告的地方去,有一個使女迎面而來,她被巫鬼所附,用法術叫她主人大得財利.保羅把鬼從她身上趕出來,令到她的主人十分氣憤,教唆人把保羅和友人都抓到監牢裡.這也是一個女人,一個無名無姓的女人,一個卑下的女人.保羅把鬼趕出來,讓她重拾新生,過正常人的生活.這個故事告訴我們,跨過國界的保羅,不單跨過性別的樊籬,也跨過階級的樊籬.

第三,    他們被關在監牢裡,遇見禁卒,並帶領他全家得救.(二十至四十節)這裡,我們還看到另外一層樊籬:監獄的樊籬,或鐵柱的樊籬.這個故事大家都耳熟能詳,我的時間也不多,所以現在暫時就不展開了.

四,結語

最後,讓我總結一下今天的講道如下:

人生就像旅途,我們常常都在旅途上.途中我們有很多不同的際遇,也到過不同的地方.今天所要講的題目叫「躑躅邊城」,講的是在邊城發生的事.那裡會有上帝的提醒和督責,甚至祂會向你說不；那裡會有神的使者帶來異象,要你跨越邊界:文化的邊界,國族的邊界,政治的邊界.但這還不夠,當你跨越之後,還要進一步跨越其他的鴻溝:性別的邊和階級的邊界,最終,我們可能要面對的是肉身的囚禁:被監牢的鐵枝困住也罷,或者就被自己這個肉身的臭皮囊困住也罷,最終我們要跨過這個監禁的邊界,才會得到真正的釋放.

躑躅邊城.請你反思一下你所走的人生路,是否到了邊城?如果你自覺是來到邊緣地,來到border town,那麼,請你傾耳靜聽主示,請你睜眼等待異象.

聆聽,睜眼,反思.躑躅邊城.



\newpage

\section{做個有心人}
\label{sec:162}
\textbf{張美珍博士 - 做個有「心」人}
\newline
\newline
\hyperref[sec:161]{< < < PREV SERMON < < <}
~
\hyperref[sec:toc]{[返目錄]}
~
\hyperref[sec:163]{> > > NEXT SERMON > > >}


日期:2012年9月9日

主題:做個有「心」人

證道:張美珍博士

大家有沒有「心」?有多少個?在哪裡?

你的「心」是什麼顏色?最好不是「黑」色,也不是「灰」色；

什麼顏色最好?紅色──壯志「雄」心；紫色──紫心勳章:美國軍隊的榮譽獎章,頒與對戰事有貢獻,或參戰時負傷的軍人.

說到這裡,相信大家都明白我所說的心,並不是指肉身上或醫學上的心──即中間偏左,負責輸送血液的器官.我們在日常生活上,除了形容身體上的病痛──例如心臟病,心交痛等會提到「心」這個字外,我們絕大部份時間所提到的「心」；都並非這個肉身上的「心」.

「心」這個中文字,變化萬千；是英文字的heart無法相比.心,可以形容不同的狀態──痛心,痴心,細心,粗心,忍心,小心,苦心,死心,恆心,清心,硬心,童心,空心.......

當某某與某人暗地裡不和,我們會形容他們有「心病」；對於熱戀中的男女來說,最希望對方「真心」而又「不花心」；到了談婚論嫁階段,最重要就是「一條心」；若你已為人父母,而又有超過一個子女的話,千祈不要「偏心」,遇上快樂或難過事情,我們會覺得「開心」或者「傷心」；有時看到一些倫常慘案或謀殺案的新聞,我們又會慨嘆怎麼那樣「狠心」.......

今天讀的四段經文,都有提及「心」字；這個「心」又是是指什麼「心」?

申四:9「你只要謹慎,殷勤保守你的心靈,免得忘記你親眼所看見的事,又免得你一生這事離開你的心；總要傳給你的子子孫孫.」

詩十五:1－2「耶和華啊,誰能寄居你的帳幕?誰能住在你的聖山?就是行為正直,做事公義,心裡說實話的人. 」

雅一:21；　2621. 所以你們要脫去一切的污穢和盈餘的邪惡,存溫柔的心領受那所栽種的道,就是能救你們靈魂的道.26. 若有人自以為虔誠,卻不勒住他的舌頭,反欺哄自己的心,這人的虔誠是虛的.

可七:6；　21－236.耶穌說:以賽亞指著你們假冒為善之人所說的預言是不錯的.如經上說:這百姓用嘴唇尊敬我,心卻遠離我.21. 因為從裡面,就是從人心裡,發出惡念......23. 這一切的惡都是從裡面出來,且能污穢人.

很明顯,這四段經文所說的「心」,當然都不是肉身的「心」,而是以指一個人的心思意念(mind).心思意念,我們看不見,觸不到,在莎劇《麥克白》有一句對白:「世上沒有一個技術,能從一個人的臉上看出他的內心.」

心思意念,主宰我們的一切外在的行為及活動.很可惜,社會的法律,大學的校規,可以規範一個人的行為,卻管束不了一個人的心.熱戀中的情人,若然一個變了心,縱使結了婚有一紙婚書,留到他的人,卻留不到他的心.我有個朋友結婚多年,懷疑丈夫變心,每天查看他的手機,又限制他放工後一小時內回到家,結果摩擦更多,最終離婚收場.心已變,行為也會相應改變.

在立法會選舉裡,作為基督徒,我們也應該盡我們的公民責任,投票選賢能.近日城中到處掛滿候選人的宣傳品,拉票時與我們講心,說他們如何熱心公益,為民請命.然而當選之後,很多時都「貨不對辦」,是他們變了心,還是根本無心?

我曾經聽過一個關於選舉的笑話:

一位候選人到市集拉票,問其中一位途人:「我可以得到你的支持嗎?」途人答道:「對不起,我已經答應支持另一位候選人了.」政客嘻皮笑臉,死纏不休地道:「不要緊,『答應』與『實行』是兩回事,你仍然可以支持我.」途人勉為其難答道；「那麼,我『答應』支持你好了.」這些與政治選舉一關的笑話,大多諷刺政客心口不一,講一套,做一套.

今天讀的福音書經文,就記載了主耶穌當年以極為嚴厲的口吻,斥責法利賽人「假冒為善」；他們訂下很多宗教上的律例規條,以為守這些便能得到上帝的喜悅,結果這些規條反而成為欺壓弱小,阻擋別人親近主耶穌的武器!所以主耶穌說:「這百姓用嘴唇尊敬我,心卻遠離我.」(太六:6)

我在中大社工系所任教的,主要是大學通識教育的課程,我的學生來自全中大不同的學系.在香港這個競爭激烈的教育制度之下,大部份學生都被訓練成競賽者,爭名次,爭成績,爭表現,爭資源；因此緣故,造成學生之間互相競爭,有時甚至出現不誠實的行為.

曾經有一位在外國讀中學經非聯招而進入中大的商科學生,在修讀我任教的科目時,學期前後判若兩人,成績每下愈況.到了學期快結束時,我請她來我辦公室,她甫坐下便哭不成聲,原因是她其中一科主修目需要同學分組作課堂報告；她那一組得分很低,同學們怪罪於她,認為她所負責的部份做得不好,說她是外國回來,追不上香港的水平.自此她便對自己失去信心,所讀的科都讀得不好.

如果我們今天所教出來的學生,能夠為分數而傷害同學,甚至欺騙老師(最常見的是呃簽到,抄功課等)；難保他們將來出來社會做事,會否金錢而做出一些傷天害理的事情?若我們今天的大學生與當年的法利賽人一般,只是追求外表風光,欺壓弱小,失去憐憫關懷之心,我們的社會將會變成怎樣?怪不得現今社會的金融騙案,醫療失誤的情況愈趨嚴重.

前匯控集團主席葛霖(也是聖公會的義務牧師),在零八年的金融海嘯後出版了一本書,名為"Good Value: Reflections on Money, Morality and an Uncertain World", 中文版:「美好價值:金錢,道德與不確定世界的省思」；他在導言如此說:

「日益全球化的世界已墜入一場可能綿延多年的危機之中,情況之嚴峻是我們多數人此生未見的.至少四分之一個世紀以來,我們視為理所當然的一切,驟然間似有分崩離析之虞.信任大規模崩壞:對金融體系的信任,對銀行家的信任,對企業的信任,對企業領袖的信任,對從政者的信任,對媒體的信任,以至對整個全球化過的信任,全部嚴重損壞,富國窮國皆然.」(p. 11)

......「要面對這些問題,我們首先得承認,有一些教訓必須記取,集體與個人皆然.大家必須願意學習,社會才有可能恢復進步；必須有增長智慧的決心,世界才有希望；而大家將會發現,必須誠實尋找美善之道,智慧才有可能增長.」(p. 12)

這位資深行家認為,只有回歸基本,尋找美善之道,智慧才可增長,這個世界才有希望.什麼是美善之道?很多人說,宗教都是導人向善,多行善事,多種善因,多得善果；信什麼教都無所謂；但基督教與其他宗教最不同的地方,就是她所著重的,並非外在行為上的所謂善事(to do good),而是要求追隨基督的人本身成為善(to be good).我常告訴學生,一個人成功與否,不是在乎他擁有多少學位或多高的分數,而是在乎他是一個怎樣的人.一個人若缺乏了珍惜生命,尊重異己,關懷別人的心,有可能會為這個世界製造更多的災難.

當年我剛進入崇基讀神學時,最震撼我心靈的,就是她的校訓──「止於至善」,描繪人生所追求的一種完美境界.「至善」,亦正是基督教精神的核心所在.主耶穌當年在登山寶訓(太五:48)中亦向祂的追隨者說道:「所以你們要完全,像你們的天父完全一樣.」只有天父才是至善者,我們每個人天生都有人性的軟弱(神學上稱此為原罪),主耶穌當年在地上所傳的道,就是要改造人的內心,因為主宰我們行為的,正是我們的心.箴四:23:「你要保守你心,勝過保守一切,因為一生的果效,是由心發出.」聖經常常提醒我們,上帝所看重的,是我們的內心,所以要保守自己的心──因為如果我們的心與上帝連在一起,謹記上帝的教訓,愛主愛人,行為自然會有所改變.

說到「善」,我最喜愛聖經在彌迦書六:8,亦是伍牧師多次在祝福時提到的:「世人哪,耶和華已指示你何為善.他向你所要的是什麼呢?只要你行公義,好憐憫,存謙卑的心,與你的神同行.」

今日的

講題:做個有「心」人；你的心是那一個心呢?是雜念纏身的心,還是追求上帝公義的心?是自私自利的心,還是有憐憫心腸的心?求主幫助我們,讓我們有的一顆心,是基督的心.阿們.



\newpage

\section{我是生命的糧}
\label{sec:163}
\textbf{伍渭文牧師 - 我是生命的糧-我的食糧,我的生命}
\newline
\newline
\hyperref[sec:162]{< < < PREV SERMON < < <}
~
\hyperref[sec:toc]{[返目錄]}
~
\hyperref[sec:164]{> > > NEXT SERMON > > >}


約翰福音 6:51-58
\newline
\begin{longtable}{cl}
\hline
\hline
章節 & 經文 (和合本修訂版)\\
\hline
6:51 & \begin{tabularx}{0.7\textwidth}{X} 我就是從天上降下來生命的糧；人若吃這糧，必永遠活著。我為世人的生命所賜下的糧就是我的肉。」 \end{tabularx} \\ \\ \relax
6:52 & \begin{tabularx}{0.7\textwidth}{X} 因此，猶太人彼此爭論說：「這個人怎能把他的肉給我們吃呢？」 \end{tabularx} \\ \\ \relax
6:53 & \begin{tabularx}{0.7\textwidth}{X} 耶穌對他們說：「我實實在在地告訴你們，你們若不吃人子的肉，不喝人子的血，在你們裡面就沒有生命。 \end{tabularx} \\ \\ \relax
6:54 & \begin{tabularx}{0.7\textwidth}{X} 吃我肉、喝我血的人就有永生，並且在末日我要使他復活。 \end{tabularx} \\ \\ \relax
6:55 & \begin{tabularx}{0.7\textwidth}{X} 我的肉是真正可吃的；我的血是真正可喝的。 \end{tabularx} \\ \\ \relax
6:56 & \begin{tabularx}{0.7\textwidth}{X} 吃我肉、喝我血的人常在我裡面，我也常在他裡面。 \end{tabularx} \\ \\ \relax
6:57 & \begin{tabularx}{0.7\textwidth}{X} 永生的父怎樣差我來，我又怎樣因父活著，照樣，吃我肉的人也要因我活著。 \end{tabularx} \\ \\ \relax
6:58 & \begin{tabularx}{0.7\textwidth}{X} 這是從天上降下來的糧，不像你們的祖宗吃過嗎哪還是死了；吃這糧的人將永遠活著。」 \end{tabularx} \\ \\
[1ex]
\hline
\hline
\end{longtable}
日期:2012年9月16日

主題:我是生命的糧-我的食糧,我的生命

證道:伍渭文牧師

讀經:約翰福音 6:51-58

51     我是從天上降下來生命的糧．人若喫這糧,就必永遠活著．我所要賜的糧,就是我的肉,為世人之生命所賜的.     52     因此,猶太人彼此爭論說,這個人怎能把他的肉,給我們喫呢.     53     耶穌說,我實實在在的告訴你們,你們若不喫人子的肉,不喝人子的血,就沒有生命在你們裡面.     54     喫我肉喝我血的人就有永生．在末日我要叫他復活.     55     我的肉真是可喫的,我的血真是可喝的.     56     喫我肉喝我血的人,常在我裡面,我也常在他裡面.     57     永活的父怎樣差我來,我又因父活著,照樣,喫我肉的人,也要因我活著.     58     這就是從天上降下來的糧．喫這糧的人,就永遠活著,不像你們的祖宗喫過嗎哪,還是死了.

開學了,社會熱燙燙的新聞仍然是開學時應否撤回新設立的國民教育科,其實六零年代的中學有一科叫公民科(civic), 但內容不是教我們如何做一個好公民,乃講論香港的社會和政府組織,是資訊多於道基督教育,但教我的老師是在香港大學唸哲學的,第一課就給我們上了非常有啟發性的道基督教育,他一進課堂,就問我們中一的學生:究竟我們吃喝,目的為了生存,還是為了活下去,我們吃喝呢?(Do we eat to live or do we live to eat?) 吃喝是手段,還是吃喝是目的呢?

當然有些人活著-很辛苦的幹活,在烈日下汗流夾背,就是盼望收工後,一夥兒喝杯冰凍啤酒,大魚大肉,這就是他們的享受,整天辛苦就是為了恢復體力,大吃一頓.他們可能很少買衣服,更少會買古龍水,活著就是能吃得痛快.我以前住德福花園,鄰近是牛頭角,旁晚在大牌檔見到很多這些體力勞動工人.也有人活著是要吃得精緻,他們花錢旅行,重點不在尋訪名勝,乃要嚐盡各地名菜,並說可以從名菜中體會不同文化.

從提比哩亞來的人

耶穌說這段話,因為有一班猶太人,看見耶穌施行五餅二魚神蹟,吃飽五千人,他們要找耶穌.「然而有幾只小船從提比哩亞來,靠近主祝謝分餅給人吃的地方.眾人見耶穌和門徒都不在那裡,就上了船,往迦百農去找耶穌.」(約六24) 他們尋找耶穌呢是否為了美食?嚐過五餅二魚,味道不錯,食過翻尋味.耶穌不在那裡了,神蹟結束,曲終人散,他們就應該回家,然而他們鍥而不捨,要到迦伯農找耶穌.似乎他們不單是為了食物來找耶穌,他們找耶穌可能另有目的.

我們留意經文特別指出他們一夥人來自同一個城市-提比哩亞,而第六章開始時也特別提及提比哩亞海:「這事以後,耶穌渡過加利利海,就是提比哩亞海」(六1) 為何作者提及提比哩亞呢?提比哩亞是有濃烈的政治意涵的,加利利海已改名為提比哩亞海了.提比哩亞是加利利的首府,政治商業行政中心,提比哩亞是大希律王的兒子希律安提帕(Herod Antipat) 管治加利利時建立的首府,用以紀念繼承該撒亞古士督(Caesar Augustus)的羅馬皇帝提庇留(Tiberius).耶穌降生時,代理羅馬人統治巴勒斯坦是大希律,他死後國土一分為四,管治加利利和約但河東比利亞的就是分封王希律安提帕.希律安提帕取悅羅馬主子,鐵蹄管治猶太人,並殺害施洗約翰.施洗約翰因指責希律安提帕不應奪去兄弟腓力的妻子希羅底(Herodias)而坐獄,最後希羅底懷恨在心游說希律安提帕要施洗約翰的人頭,是他把施洗約翰斬首,也是他把耶穌釘在十字架.

從提比哩亞來的幾條船的猶太人來找耶穌,顯然並不祇希望再嚐五餅二魚那麼簡單,他們把希律安提帕的劣行看在眼裡,從耶穌身上他們看到希望；耶穌既然能夠解決糧食問題,也可以成為他們的政治領袖,帶來民族的解放,脫離羅馬的鐵蹄統治.

你們找我,並不是因見了神蹟

耶穌對他們說:「你們找我,並不是因見了神蹟,乃是因餅得飽.不要為那必壞的食物勞力,要為那存到永生的食物勞力,就是人子要賜給你們的,因為人子是父神所印証的.」(約六27) 神蹟這字,原文指記號或路標,路標是指向目的地,路標不是叫我們留意本身,五餅二魚神蹟顯示耶穌是基督,他是從天上而來的天糧.其實,他們己看見兩個神蹟了,另一個神蹟就是耶穌履海.請留意他們問耶穌甚麼時候到了迦伯農:「既在海那邊找著了,就對耶穌說:拉比,是幾時到這裡來的?」昨晚加利利海不是狂風大作嗎?那麼惡劣的天氣,根本不可能回航,我們的船也是今早待風平浪靜後才趕來,你就不見了,你甚麼時候回到迦伯農呢?我想他們一夥人,應該從門徒那裡聽到昨晚另一個神蹟-耶穌履海.

耶穌對他們說:你們應該藉著這些神蹟,明白我是從上帝而來的,你們要關心是那不能朽壞的國度,希臘帝國會過去,羅馬帝國也會過去,但上帝的國度永遠長存.

事實上,研究糧食的專家告訴我們,糧食問題不是生產問題,乃分配問題,浪費問題；個人利益,自我中心問題.世界上仍有許多人民處於赤貧狀態,但每年竟有三分之一糧食白白浪費.聯合國糧農組織(FAO, Food and Agriculture Organization)2011年五月出版的《全球糧食耗損及糧食浪費報告》指出,每年損失或浪費的食物高達13億公噸,約佔生產量的三分之一,其中以蔬菜和水果的耗損比例最高.在歐洲和北美,每個消費者每年所浪費的食物高達95─115公斤；在香港,每天產生廚餘達3,200公噸.

同理,解決貧窮不能單靠是增加財富,使經濟增長,乃要處理財富分配不均,因大部分資源被少數人所壟斷.在香港中小企業生意好一些,就要面臨加租.在中國大陸,要是沒有三公的揮霍,根本不需要希望工程.三公消費,指政府部門人員在因公出國(境)經費,公務車購置及運行費,公務招待費產生的消費, 2011年中央行政單位,事業單位和其他單位的"三公經費"支出,合計93.64億元人民幣.

耶穌要來自提比哩亞的猶太明白,惟有吃了永生的食物,改變生命,我們才有真正的盼望.他來並不是建立地上的國度,乃在我們心中,建立屬靈的國度.這是友愛的國度,生命對人敞開,為他者而活的國度.

愛與開放的生命

耶穌說我是從天上來的糧食,吃他肉,喝他血的人有屬天的生命,不朽的生命.猶太人不明白如何吃耶穌的肉,喝耶穌的血:「這個人怎能把他的肉給我們吃呢?」耶穌解釋吃他的肉喝他的血就是與他的生命聯合,進入愛與契通的生命.「吃我血的人常在我裡面,我也常在他裡面.永活的父怎樣差我來,我也因父活著；照樣,吃我肉的人也要因我活著.」(約六56, 57)

這聯合可以用婚禮中以戒指為盟誓來解釋:「我以此戒指為徵記,與你結婚,我以我的身心尊敬你,我以我的財物與你分享.」這是歡愉的對換.基督把我們的罪擁攬在身上,我們得著他的生命,脫離罪的綑縛；更因著這生命,走出自我的世界.基督被父神所差遺,為順服父,為拯救世人,為他者,基督死在十字架.我們吃了從天而來的糧食,也得著這屬天的生命,走出自我的世界,進入他人的生命,為他者而活,也彼此相愛,這是上帝國度的記號.羅馬帝國的鬥爭和虛妄

相對於耶穌要建立的友愛,真理的國度,羅馬帝國就充滿仇恨,虛妄.加利利海改稱提比哩亞海是紀念羅馬皇帝提庇留,提庇留晚年暴虐無道,把關在獄中的政敵殘殺並禁止親人為他們舉哀,暴屍荒野,不得殮葬或焚化.到提庇留臨終時,舉國百姓歡騰不已,後來說他又活過來,百姓沉默不語(非常失望,感到沒趣),最後聽到提庇留被政敵焗死,再度興奮莫名.

繼位的卡利古拉(Caligula) 行為更為狂妄荒誕,淫亂兇殘.自詡是神靈顯世,在羅馬為自己立廟,並計劃在耶路撒冷聖殿立像,叫人膜拜.為了一時之快,因為不夠罪犯放進鬥獸場,竟把部分觀眾推入場內鬥獸.一位凡人把自己看成是上帝,但上帝的兒子卻甘願成為凡人,為人贖罪.一個國度充滿恐懼,權謀,另一個國度要棄絕謊言,以愛相繫.

生命的糧不單給我們相愛的生命,也帶給我們感恩祝謝的人生.

感恩祝謝的生命

耶穌是生命的糧,我們要吃他的肉,喝他的血,不單是一件已過去的救贖歷史,也是崇拜禮儀的核心,教會稱這禮儀為感恩祭(Eucharist) ,主餐,聖餐.路加福音記載耶穌復活後,在以馬忤斯路上對兩位門徒講解聖經及和他們一起用餐,正是復活的主每主日臨到信徒當中要做的事.透過傳道者解釋聖道,藉著領受聖餐,耶穌分享他的生命:吃他的肉,喝他的血.

古教會的崇拜分兩部分:聖道與聖餐,甚至視聖餐部分才算真正的崇拜,慕道者要上慕道班三年,才可以領洗,聽完聖道,襄禮的執事便高聲說:你們離開吧!Ite Missa est(Thus you are dismissed) 未領洗的慕道者離開後,便開始崇拜了,領受聖餐了,彌撒開始了,彌撒mass就是離開的意思.所以到現在每主日在天主教,正教(Orthodox church),聖公會(Anglican Church),信義會(Lutheran Church)都保存每主日都領聖餐的傳統.宗教改革時期,加爾文(Calvin)直言,應該每主日都有聖餐,一月一次是撒旦的意見,遷就人的軟弱.天主教,正教和聖公會稱崇拜為感恩祭Eucharist.我們稱聖餐為感恩祭,因為耶穌拿起餅來,祝謝了,就擘開,祝謝原文是Eucharist.

崇拜是感恩祭,我們進行聖餐的禮儀,就貫穿感謝的觀念.在開始主禮和會眾的對答:願主和你們同在(應:也和你同在) ；你們的心當仰望主(應:我們的心仰望主) ；我們應當感謝我們的主上帝(應:感謝我主上帝是應該的) .跟著我們讀出耶穌設立聖餐的話,然後說:「在基督的祭司職份裡我們聯合在一起,向你獻上這祭禮.」基督以感謝父的心,順服至死,拯救我們,我們效法基督,以感謝的心,把身體獻上,當作活祭.約翰福音六51-58的經文已在歷史成就了,又每主日藉教會聖餐的禮儀,重新演繹.

西方教會強調基督的人性,在聖餐或感恩祭中,基督被描述為上帝的羔羊(Lame of God),在中古的羅馬天主教,聖餐桌是祭壇(altar) ,崇拜就是重獻基督,我們更正教(Protestantism) 雖然強調聖餐桌是聖桌(Table),不是獻祭的祭壇,但紀念主餐時還是有一種哀慟,想念基督在十字架上為我們受死.東正教崇拜的感恩祭(Eucharist) 強調基督的神性,基督被描繪成掌管萬有者(Pantocrator), 是基督從天降臨到我們當中,分享生命.崇拜是一個旅程(Journey) ,我們進入神的國度,而基督從天堂臨到我們中間,分派永生食糧.東正教在圖像幛與信徒分隔,圖像幛象徵天堂,讀福音書時,主禮從圖像幛(Iconostasis)的門出來,宣讀福音書,稱小進門(Small Entrance) ；把聖餐的餅酒端出來時稱大進門(Great Entrance) ,聖餐是分派基督神聖身體,我們祝謝,然後分領主的聖體寶血.

祝謝是基督信仰的中心思想,聖餐是喜樂的筵席,基督在聖桌旁分享他的聖體寶血,基督不是在聖桌上,供我們默念懷想.He is at the table , not on the table.正教會(Orthodox Church)聖餐的詩歌「凡有血氣歌, 首頌441首」的歌詞這樣描述:

凡有血氣,皆當靜默,敬畏謙恭同肅立,莫想絲毫世俗事物,因主基督臨世界,手分賜天福神恩,接受我眾恭敬拜.

王之王為聖母所生,早已臨世悠久長,主之主竟取人樣式,血肉為體掩神光,自願將其聖身寶血,賜與信者作天糧.

千萬天軍,排列成行,前隊先鋒一路長,周圍環繞光中之光,親從永晝天堂降,光之光從無極領域,親臨世間從天降.

俯伏主足下常警醒,撒拉弗與兩翅遮面,大聲讚美,頌歌騰響門庭震,哈利路亞,哈利路亞,哈利路亞,主至聖.阿們.

作為天糧的基督,藉聖餐來到我們中間,分享他的生命,聖餐重點不在懷想過去基督的死,憂傷哀慟,聖餐是上帝的筵席,是歡樂,得勝的.「這是從天上來的糧.吃這糧的人,就永遠活著.」(約六58)

分享的生命

祝謝的生命也是分享的生命:「祝謝了,就擘開,遞給他們」；對上帝感恩的生命,也是分享的生命；感恩祝謝,因明白生命是上帝所賜,我們活著存留都在乎他,我們不過是管家,上帝是是終極的審計長,和我們算賬.在這裡我想起胡秀英教授,我鼓勵大家散會後買「冬青之母胡秀英」,這書有胡教授寫的一篇文章「芥菜種子一粒」,簡單勾勒她的一生,文章充滿感恩,她一生就是想到他人.其實文章沒有提到的是她中年喪偶,她沒有自己的兒女,但她為上帝關心了無數子子孫孫.

祝謝的生命就是赤子的生命,對上帝信賴,願意與人分享的生命,是這樣的生命,啟動了五餅二魚的神蹟.當耶穌對門徒說:「我們從那裡買餅叫這麼多人吃呢?」門徒正憂愁時,一個小孩就把自己僅有的五餅二魚給耶穌.院長在這書的序中第二首詩「赤子心」這樣描述胡秀英教授:

崇基的校園,常見一個個子細小的老婦人,滿頭銀發,逐漸背有點兒傴僂,身體雖然逐漸地老去,慢慢地衰殘,但她的內心,從來就祇有童真.

逾百歲的老人,竟然是一個赤子.每一日她懷著那顆赤子心,工作從來不計算報酬,待人從來不存著機心,她的信仰,是一面白璧.......

就像小孩子對父母的相信.因為她聆聽並且記住耶穌對她說的話:讓小孩子到我這裡來.進天國的要像小孩子.現今她歡然回到天家,接受耶穌的擁抱和祝福.我們為她高興,但我們也會時常懷念她.我的食糧,我的生命

為何胡教授有這生命呢?因她重視生命的糧食,她勤讀聖經,勤於聚會.我們食甚麼食物,就有甚麼的生命.

真的,我們食甚麼食物,就有甚麼的生命.美國科羅拉州電影院槍擊慘案中的失意博士生James Holmes智商很高,受電影虛擬暴力影響,在蝙蝠俠:夜神起義(The Dark Knight Rises) 首影當晚,穿上該電影上集歹角小丑(Joker)的服飾行凶,小丑智商極高但失意,是蝙蝠俠的死敵,常常做一些瘋狂毀滅性的事,折磨蝙蝠俠.在電影中小丑曾說過:「瘋狂是逃生門,走出外邊,把門關上,把可怕的事情掉在後面,永遠鎖起來.」(``Madness is the emergency exit. You can just step outside, and close the door on all those dreadful things that happened. You can lock them away... forever.")

很惋惜一個智商那麼高的博士生,竟然作出這樣瘋狂毀滅性的行徑.很明顯,他受了電影中小丑的影響,以為瘋狂就是逃生門,可以把失敗關掉在背後,上鎖.我們吃甚麼糧食,就有甚麼的生命.報載凶手有返教會的,但似乎他受電影中虛擬暴力影響,多於受聖經真理影響；這件事也提醒我們,人在教堂,但我們的心是否被上主的道所吸引呢?怎樣的糧食,怎樣的生命；你看些甚麼書,整天想些甚麼,就成為甚麼的人.

胡秀英教授提到她母親雖然不識字,但信了耶穌,接受了生命的糧,眼光開擴了,怎樣辛苦也讓胡教授受好的教育,胡教授是整條村第一個完成中學的女生,因而進入金陵女子大學,被金陵校訓厚生(Abundant Life, 豐盛的生命) 深深所影響:「我來是要叫羊得生命,並且得的更豐盛.」(約十10) 豐盛的生命就是祝謝感恩的生命,也是對上帝,對人開放的生命,對上帝感謝,對人分享的生命.

不同的食糧,不同的生命；一個是瘋狂,毀滅性的瘋狂；一個是厚生,澤及千萬人的厚生.

結論:

究竟我們吃喝,目的為了生 存,還是為了活下去?(Do we eat to live or do we live to eat?)

有些人真的覺得吃喝玩樂就是人生的目的,保羅在以弗所傳道時,在以弗所的人就這樣說:「我們就吃吃喝喝吧!因為明天死了」(林前十五32) ,但耶穌耶穌說:我就是生命的糧,吃這糧的人就永遠活著.

讓我們領受那生命的糧,得著愛與契通的生命,面對人生的橫逆,也可以感恩；物質匱乏時,亦可以分享,因為基督在我裡面,我們在基督裡面.阿們.



\newpage

\section{你們也要離開嗎}
\label{sec:164}
\textbf{林豪恩先生 - 你們也要離開嗎?}
\newline
\newline
\hyperref[sec:163]{< < < PREV SERMON < < <}
~
\hyperref[sec:toc]{[返目錄]}
~
\hyperref[sec:165]{> > > NEXT SERMON > > >}


約翰福音 6:56-69
\newline
\begin{longtable}{cl}
\hline
\hline
章節 & 經文 (和合本修訂版)\\
\hline
6:56 & \begin{tabularx}{0.7\textwidth}{X} 吃我肉、喝我血的人常在我裡面，我也常在他裡面。 \end{tabularx} \\ \\ \relax
6:57 & \begin{tabularx}{0.7\textwidth}{X} 永生的父怎樣差我來，我又怎樣因父活著，照樣，吃我肉的人也要因我活著。 \end{tabularx} \\ \\ \relax
6:58 & \begin{tabularx}{0.7\textwidth}{X} 這是從天上降下來的糧，不像你們的祖宗吃過嗎哪還是死了；吃這糧的人將永遠活著。」 \end{tabularx} \\ \\ \relax
6:59 & \begin{tabularx}{0.7\textwidth}{X} 這些話是耶穌在迦百農會堂裡教導人的時候說的。 \end{tabularx} \\ \\ \relax
6:60 & \begin{tabularx}{0.7\textwidth}{X} 他的門徒中有好些人聽見了，就說：「這話很難，誰聽得進呢？」 \end{tabularx} \\ \\ \relax
6:61 & \begin{tabularx}{0.7\textwidth}{X} 耶穌心裡知道門徒為這話私下議論，就對他們說：「這話成了你們的絆腳石嗎？ \end{tabularx} \\ \\ \relax
6:62 & \begin{tabularx}{0.7\textwidth}{X} 如果你們看見人子升到他原來所在之處，會怎麼樣呢？ \end{tabularx} \\ \\ \relax
6:63 & \begin{tabularx}{0.7\textwidth}{X} 聖靈賜人生命，肉體毫無用處。我對你們所說的話就是靈，就是生命。 \end{tabularx} \\ \\ \relax
6:64 & \begin{tabularx}{0.7\textwidth}{X} 可是你們中間有些人不信。」耶穌起初就知道哪些人不信他，哪一個要出賣他。 \end{tabularx} \\ \\ \relax
6:65 & \begin{tabularx}{0.7\textwidth}{X} 於是耶穌說：「所以，我對你們說過，若不是蒙我父的恩賜，沒有人能到我這裡來。」 \end{tabularx} \\ \\ \relax
6:66 & \begin{tabularx}{0.7\textwidth}{X} 從此，他門徒中有很多退卻了，不再和他同行。 \end{tabularx} \\ \\ \relax
6:67 & \begin{tabularx}{0.7\textwidth}{X} 耶穌就對那十二使徒說：「你們也要離開嗎？」 \end{tabularx} \\ \\ \relax
6:68 & \begin{tabularx}{0.7\textwidth}{X} 西門‧彼得回答他：「主啊，你有永生之道，我們還跟從誰呢？ \end{tabularx} \\ \\ \relax
6:69 & \begin{tabularx}{0.7\textwidth}{X} 我們已經信了，又知道你是神的聖者。」 \end{tabularx} \\ \\
[1ex]
\hline
\hline
\end{longtable}
日期:2012年9月23日

題目:你們也要離開嗎?

證道:林豪恩先生

經文:約翰福音6章56-69節

一．引言

多年前看過一本書,名叫《為甚麼我不敢告訴你我是誰》.作者對這個問題的回答道出了很多人的心聲:「我不敢告訴你我是誰,因為,如果我告訴了你我是誰,你可能會不喜歡這個我.」這樣的擔心是實在的,今天的福音經課讓我們看到,當耶穌告訴人們多一點真正的自己的時候,人們開始不喜歡他,並且離他而去.

二．經文

56     喫我肉喝我血的人,常在我裡面,我也常在他裡面.     57     永活的父怎樣差我來,我又因父活著,照樣,喫我肉的人,也要因我活著.     58     這就是從天上降下來的糧．喫這糧的人,就永遠活著,不像你們的祖宗喫過嗎哪,還是死了.     59     這些話是耶穌在迦百農會堂裡教訓人說的.     60     他的門徒中有好些人聽見了,就說,這話甚難,誰能聽呢.     61     耶穌心裡知道門徒為這話議論,就對他們說,這話是叫你們厭棄麼.〔厭棄原文作跌倒〕     62     倘或你們看見人子升到他原來所在之處,怎麼樣呢.     63     叫人活著的乃是靈,肉體是無益的．我對你們所說的話,就是靈,就是生命.     64     只是你們中間有不信的人.耶穌從起頭就知道,誰不信他,誰要賣他.     65     耶穌又說,所以我對你們說過,若不是蒙我父的恩賜,沒有人能到我這裡來.     66     從此他門徒中多有退去的,不再和他同行.     67     耶穌就對那十二個門徒說,你們也要去麼.     68     西門彼得回答說,主阿,你有永生之道,我們還歸從誰呢.     69     我們已經信了,又知道你是　神的聖者.

三．釋經與應用

1．因誤會而結合

究竟耶穌表露了關於自己的甚麼,以致人們不喜歡他呢?

56     喫我肉喝我血的人,常在我裡面,我也常在他裡面.     57     永活的父怎樣差我來,我又因父活著,照樣,喫我肉的人,也要因我活著.     58     這就是從天上降下來的糧．喫這糧的人,就永遠活著,不像你們的祖宗喫過嗎哪,還是死了.     59     這些話是耶穌在迦百農會堂裡教訓人說的.     60     他的門徒中有好些人聽見了,就說,這話甚難,誰能聽呢.     61     耶穌心裡知道門徒為這話議論,就對他們說,這話是叫你們厭棄麼.〔厭棄原文作跌倒〕

為甚麼耶穌向人們說出「吃我肉,喝我血」這麼難以理解的意象去啟示自己呢?原來,吃喝這個主題,是從前文「五餅二魚」的事件引申下來的.約翰福音第6章1-2節說:「耶穌渡過加利利海,就是提比哩亞海.有許多人因為看見他在病人身上所行的神蹟,就跟隨他.」「眾人就坐下,數目約有五千.耶穌拿起餅來,祝謝了,就分給那坐著的人；分魚也是這樣,都隨著他們所要的.他們吃飽了．．．第二日．．．眾人見耶穌和門徒都不在那裡,就上了船,往迦百農去找耶穌.」(6:10-24)這批追隨耶穌的群眾,是剛剛在五餅二魚神蹟中吃飽的人,耶穌順理成章地繼續以食物和吃喝作為主題和意象教導他們.

一般來說,人人都喜歡被歡迎,享受被肯定,而且,跟隨者越多,代表著越受歡迎,越得到認同,也代表著越成功.然而,尷尬的地方是,如果當事人深知跟隨者只是出於誤會以及錯誤的期望而追隨,也意識到如果追隨者知道真正的自己,就會失望而鳥散的話,這個人應該如何自處呢?他應該徇眾要求,修改自己?或者利用群眾的誤會,乘勢而上?還是冒險面對群眾的期望,澄清誤會呢?

耶穌選擇了直接面對群眾的期望,刺破誤會的泡沫.他把人們的期望反映給他們知道:「我實實在在地告訴你們,你們找我,並不是因見了神蹟,乃是因吃餅得飽.」(6:26)

如果你是人群中的一員,你會如何回應耶穌呢?可能有人會這樣想:「我們找你,當然是因為見了神蹟.如果不是見了神蹟,我們怎知道你能夠醫好我們的病,餵飽我們的肚呢?」

誤會,有時就是這樣產生.耶穌或者沒有意圖去令群眾對他有醫病和「醫肚」的期望,他行神蹟,是希望人們認識「道成了肉身,住在我們中間,充充滿滿地有恩典有真理」(1:14).然而,當耶穌有所行動的時候,他卻不能控制群眾以自身的需要去理解他的行動,以及因此而產生與他們的需要相關的期望.隨著跟隨的人越來越多,投放在耶穌身上的期望越來越大.難道群眾不可以按著自己的需要而對耶穌投以期望嗎?當然可以.就正如我們按著自己的需要而對其他人投以期望一樣,或者其他人按著他們的理解對我們投以期望一樣.這只是自然而合理的生活現象.既然對他有期望是自然而合理的,耶穌應該如何回應呢?他應否被群眾的期望去改變他道成肉身的使命呢?他應否被群體的視線去改寫他的人生呢?然而,如果耶穌選擇不去滿足群眾的期望,他會否感到愧疚呢?會否感到有負於這些對他有期望的人呢?會否擔心人們會對他失望而離棄他呢?令人失望,被人厭棄,對很多人來說,都是重大的心理負擔及威脅.

從幼兒時期開始,我們已經直接或間接地體會到,若我們不去滿足別人的期望,或者令別人失望,後果可以是非常不好受的.有些孩子做不到成人的要求,或者達不到成人的期望時,他們會感受到成人的不滿和憤怒,嚴重的甚至會被趕出門外,令這些小孩子感到做不到要求或者達不到期望是會被遺棄的.到學校讀書的時候,有些同學發現了若要老師喜歡自己,就要按著老師的期望表現自己.交了朋友,成長中的青少年感受到朋輩的力量原來可以非常龐大的,若不同聲同氣,不幸地成為了議論和杯葛的對象,一定比考試不合格更難受.進入職場,人人都深知上司的期望有生殺大權；以為有朝一日坐了上司的位置,就可以自由自在地做回自己,怎知坐上了,才體會到位置越高,看著自己的眼睛也越多,加在自己身上的期望也更大.最令人難受的,也是最令人感到威脅的,是別人對自己失望的眼神和態度.這些負面的反應,令人感到自己一無是處,是沒有價值的人,是應該被放棄的人.陷入這種困境的人,會感到思想混亂和情緒不穩.

有一位同學本來開開心心地在某所排名較後的大學修讀喜歡的學科.她的父親是在排名第一的大學的畢業,母親是在排名第二的大學畢業,他們要求女兒第二年以非聯招入學辦法再報讀這兩所大學,並且策略性地填報她的成績有可能入讀的學系.結果,不知道是幸運還是不幸,她於八月底開學前收到取錄通知.她猶豫,因為這不是她想讀的學科,而是也不是她擅長的學科.她的父母勸告她,只要捱過第一年,就可以轉系往她有興趣的學科.由於她錯過了所有的迎新活動,開學後面對人生路不熟的孤單,同時發現自己對於現在修讀的學科的確力有不逮,她越來越感到惆悵.最令她困擾的是,原來第二年要轉去她心儀的學系,她必須要在第一年取得非常突出的成績,這要求對於力有不逮的她,簡直是不可能的任務.她向本來就讀的學校查詢,學校容許她在指定日期之內恢復學籍.然而,她清楚知道如果她從目前這所排名較前的大學退學,回到那所排名較後的大學,她的父母會非常失望.她一方面不想父母失望,也不能面對父母的失望,另一方面對著自己沒有興趣又力有不逮的科目,而且轉系無望,好像度日如年,充滿挫敗感地一日一日捱下去.

群眾期待的,是飽肚的餅和魚,而耶穌提供的,卻是太初的道化成的肉和血；或者說,群眾從耶穌身上要得到的,是他提供的物質,而耶穌提供的,卻是他自己.耶穌必須面對群眾的期望和他自己的本質和使命之間的分岐.他不能夠取消別人對他寄予各種期望的自由,甚至不能夠禁止別人對他失望,但他也有責任保持自己按著天父的美意而存在,有責任幫助人們認識真正的自己不是變食物出來的魔術師,而是那永恆的道,是「道路,真理,生命」.

2．因了解而分開

60     他的門徒中有好些人聽見了,就說,這話甚難,誰能聽呢.     61     耶穌心裡知道門徒為這話議論,就對他們說,這話是叫你們厭棄麼.〔厭棄原文作跌倒〕     62     倘或你們看見人子升到他原來所在之處,怎麼樣呢.     63     叫人活著的乃是靈,肉體是無益的．我對你們所說的話,就是靈,就是生命.     64     只是你們中間有不信的人.耶穌從起頭就知道,誰不信他,誰要賣他.     65     耶穌又說,所以我對你們說過,若不是蒙我父的恩賜,沒有人能到我這裡來.     66     從此他門徒中多有退去的,不再和他同行.     67     耶穌就對那十二個門徒說,你們也要去麼.

果然,當耶穌啟示自己越清晰的時候,人們就越清楚知道耶穌與他們的期望之間的分歧,人群就失望而散去了.可惜及傷痛的是,不單止吃過五餅二魚的群眾離開,連他的門徒中也有不少人因此而離他而去.為甚麼說這是可惜和傷痛呢?能夠從群眾分別出來稱為門徒的人,應該與耶穌有某程度上的關係,也有某程度上的承諾.耶穌願意向他們躺開自己,亳不掩飾地啟示自己,本來是他們使徒之間增進了解,建立更親密的關係的機會.然而,某些門徒來卻因此而厭棄他們的老師,離他們的老師而去.若果你向一個朋友躺開自己,讓關係進深,他卻因此而厭棄你,離你而去,怎能不叫人感到可惜和傷感呢?這種傷害可以嚴重地打擊一個人的自信心,令人意志消沉,也可以令人自憐,埋怨自己為甚麼那麼愚蠢讓自己受傷害,更會令人失去對別人的信任.如果這種傷害表現為憤怒,把厭棄自己的人理解為背叛,就會怒不可遏.倘若這些事件發生在武俠小說裡面,接著的劇情很可能就是清理門戶.

能夠向帶著誤會的群眾澄清自己並不是為了滿足他們的期望而來,這已經不容易的事情；當群眾,甚至門徒失望而離棄自己,仍然能夠不亢不卑,有節有理地繼續表達自己,繼續說明自己的身份和使命,這更是難得.能夠如此的人,必定是對自己有清楚的自我了解,以及充份的自我肯定.

61     耶穌心裡知道門徒為這話議論,就對他們說,這話是叫你們厭棄麼.〔厭棄原文作跌倒〕     62     倘或你們看見人子升到他原來所在之處,怎麼樣呢.

群眾沒有興趣理會耶穌從那裡來,往那裡去,因為他們只想在耶穌身上得到自己想要的東西,但耶穌自己不能忘記自己從那裡來,往那裡去,因為從那裡來,往那裡去,就定義了他自己,與他的身份息息相關,與他的使命息息相關.他必須提醒自己,不能被群眾的歡呼聲充昏頭腦,也不能被群眾離棄的威嚇而改變方向和目的地.人子終有一天要升到他原來之處.這種永恆的視野與肯定,令耶穌在被讚揚與被批評之間,被需要與被輕視之間,被擁戴和被離棄之間,仍然頂天立地,站立得穩.

我們有沒有想過,是甚麼令我們在讚賞與批評之間被拋來拋去呢?是甚麼令我們在被需要與被輕視之間折騰呢?是甚麼令我們在被擁戴與離棄之間要生要死呢?有一天,我們都要回到原來所在之處.我們原來所在之處在哪裡呢?我們從哪裡來呢?最終往哪裡去呢?如果我們仍記得,我們是從天父那裡而來,最終會回到天父哪裡去,或者,我們就比較有力量在起伏不定的人生遭遇中,在變化不定的價值與評價中,不亢不卑地活出天父創造我們,活在天父創造我們的美意.

58     這就是從天上降下來的糧．喫這糧的人,就永遠活著,不像你們的祖宗喫過嗎哪,還是死了.

63     叫人活著的乃是靈,肉體是無益的．我對你們所說的話,就是靈,就是生命.

知道自己從那裡來,往那裡去,能夠幫助人安身立命；知道自己為甚麼而出現於世上,能夠幫助人生活得有方向.

群眾肚餓要吃,這是自然和合理的需要,不過,這需要卻不一定要由道成肉身的耶穌來解決,因為上帝創造天地的時候,已經提供了生產食物的環境,也交托了「治理全地」的任務給人類.耶穌知道自己成為肉身來到人間,不是要取代土地提供食物,而是要喚起人關注生命,關注永恆的生命,關注人類與創造主的關係.

昔日,耶穌對從他身上尋求食物的群眾澄清並呼籲他們除了關注食物外,也要關注生命,除了關注肚腹的需要,也要關注心靈的需要.因為人是很容易和自然地以物質去滿足心靈需要的.耶穌曾經說一個關於財主的故事,最能代表這種傾向:「有一個財主田產豐盛；自己心裡思想說:『我的出產沒有地方收藏,怎麼辦呢?』又說:『我要這麼辦:要把我的倉房拆了,另蓋更大的,在那裡好收藏我一切的糧食和財物,然後要對我的靈魂說:靈魂哪,你有許多財物積存,可作多年的費用,只管安安逸逸地吃喝快樂吧!』(12:16-19)財主以為有充足的物質就能夠使自己的靈魂得安穩,能夠滿足心靈的需要.今天,我們的社會也不停地引誘我們用樣的方法去回應心靈的需要.如果你想擁有一個美滿的家庭,廣告建議你去買一楝更豪華的房子；如果你想擁有甜蜜的睡眠,廣告建議你去買一張更高貴的睡床；如果你想受歡迎,廣告建議你穿某名牌衣服．．．現代消費社會深深知道人們需要的其實是心靈的滿足,因此,對症下藥地把商品描述成解決心靈需要的出路,引導人們以物質去回應心靈的需要.

耶穌希望他的門徒意識到,單單以物質去回應生命的需要是沒效果的.能夠回應生命需要的,就是這成了肉身的生命之道.對於耶穌這樣的呼籲,除了厭棄而離開之外,還可以有甚麼回應呢?

3．「自反而縮,雖千萬人吾往矣」

67     耶穌就對那十二個門徒說,你們也要去麼.     68     西門彼得回答說,主阿,你有永生之道,我們還歸從誰呢.     69     我們已經信了,又知道你是　神的聖者.

透過五餅二魚神蹟吃飽的群眾當然不能接受真正的耶穌,他們鳥散了.就連跟隨耶穌的門徒,也因耶穌的自我表白而清醒過來,覺令是而昨非,一個一個離開.現在,剩下十二個門徒,耶穌問:「你們也要離開嗎?」

「你們也要離開嗎?」或者,這也是對我們的問題.還記得少年時代與你一齊在學校的佈道會中站起來決志要跟隨耶穌的那大半班同學嗎?一路上退去的有多少呢?還記得在大球場的佈道會回應呼召,和你一齊從坐位走下去站滿草地,那些決志要跟隨耶穌的人群嗎?還有多少人仍然與你一齊走這條路呢?還記得當年主日學坐在你旁邊的同學,青年團契與你一齊唱詩歌的團友,還有多少仍然與你一齊唱到今天呢?看著人們離開的時候,我們心裡會否也問自己:「我也要離開嗎?」會否也聽到耶穌問我們:「你們也要離開嗎?」

眾人離去,你仍要留下,就算你能夠獨排眾議,也總要給自己一個留下的理由.你如何回答自己這個問題呢?你如何回答耶穌這個問題呢?彼得代表當時仍沒有離耶穌而去的十二位門徒回答:「你有永生之道,我們還跟從誰呢?」彼得繼續跟隨耶穌,不是因為耶穌有醫病之道,不是因為耶穌有「醫肚」之道,也不是因為耶穌有復國之道,而是認出耶穌是成為肉身的道,是永生之道.

四．總結

或者,我們也是好像群眾一樣,好像門徒一樣,帶著各眾期望去跟隨耶穌,要在他身上得到一些我們想要的東西.那些只想在耶穌身上獲得醫治之道的群眾,始終會失望而去；那些只想在耶穌身上獲得「醫肚」之道的群眾,始終都會失望而去；那些只想從耶穌身上獲得復國之道的群眾,最終也會失望而去.然而,耶穌似乎不怕別人對他失望,因為在失望之時,耶穌期望人們認清甚麼才是永恆的盼望,期望人們找到永恆的盼望.因此,若有一天當我們感到耶穌令我們失望的時候,不一定要馬上變臉,也不一定要立即離他而去.這可能是一個機會讓我們進深認識自己的需要,進深認識我們跟隨的對象的機會.



\newpage

\section{神的豐盛充滿我們}
\label{sec:165}
\textbf{陳天祥博士 - 神的豐盛充滿我們的三步曲}
\newline
\newline
\hyperref[sec:164]{< < < PREV SERMON < < <}
~
\hyperref[sec:toc]{[返目錄]}
~
\hyperref[sec:166]{> > > NEXT SERMON > > >}


日期:2012年9月30日

題目:神的豐盛充滿我們的三步曲

講員:陳天祥博士

弟兄姊妹平安! 今天要跟大家分享的經文就是以弗所書第三章.

我這幾年都很少被邀請講道,所以時間未必一定掌握得那麼準確,可能講長了點也不一定,預先說明.可能我也不常列在教會的邀請類別之中,所以很少被邀請.在未曾講道之前,我首先說說我講道的模式,使大家可以掌握.我每一次講道的時候,我都想想會眾有什麼需要?我是從會眾的需要來出發.我在序內介紹了這需要,然後講解聖經,以聖經為本,希望能按正意,分解真理的道,同時我也給大家我的大綱,使大家明白整體的結構,我覺得這樣會有助於明白那些細節.由於是以一個靈性問題作為出發點,所以有些時候會眾會覺得我在針對著他們,其實不是,這只不過是以聖經為本,純粹是因為聖經的道理,聖靈的工作,有人覺得扎心,因為「聖經都是神的默示,於教訓督責,使人歸正,教導人學義,都是有益的」.我平常在聽道的時候我很認真,我一定會想到幾條問題,結束之後問問講員,不過可惜沒多少個講員喜歡這些聽道的人,不過我講道的時候也很認真,目的就是希望大家結束以後能夠發問題,如果大家在問題中提出你內心的掙扎,思考一些理據,以及你的疑問的話,我想這對我也有很大的幫助.

我們現在一同祈禱:天父,當我們要讀你的話語的時候,我們想起你聖經的應許,你說你的話語一被解開,就會發出亮光,使我們能夠通達,聖經又說聖靈能夠引導我們明白,進入聖經的真理,求你在這時間讓你的應許實現,我們這樣祈禱是奉主耶穌基督的名,阿們.

有一間茶樓來了一個茶樓「怪客」,經常坐了一會便離開,對於在茶樓原來是吃東西的地方,他完全沒有這種的認識,因此就沒有吃東西的動機,行為.學校,球場,家庭,任何這些地方,對這些地方的認識都影響了人的動機和行為表現.神在基督耶穌裡設立了教會,是要將人帶進榮耀裡,倘若我們沒有這個認識,可能成為教會的「怪客」的.其實,要找一個茶樓「怪客」也相當困難,但教會的「怪客」說不定就有許多,甚至已成為了「常客」. 因此我們要藉著耶穌基督所設立的教會有深切的認識,產生正確的「動機」與「行為」,這樣我們才能進入神的豐富之內,表現出我們跟其他非基督徒有著不一樣的地方.今天我們要分 享兩大重點:一,神的豐盛充滿的教會──「充滿」是神設立教會的目的,同時「充滿」也有一個程序,步驟.我們要談的就是:「充滿」是「目的」和「步驟」 在目的之下,我們可從三個角度來看:從保羅的禱告,保羅的人生(第三章)及整卷以弗所書是「充滿」的「目的」. 在「步驟」之下我們要看看三個詞:第一個詞便是「剛強」──第十六節 「求祂按祂豐盛之路,藉著祂的靈,叫你們心裡的力量剛強起來」；第二個是「內住」(就是基督住在我們心裡)──第十七節「使基督因你們的信,住在你們心裡,叫你們的愛心有根有基」； 第三個是「一同明白」──第十八節「能以和眾聖徒一同明白基督的愛是何等長闊高深」.

我們首先來看「目的」──這目的就是使外邦的教會(如以弗所教會等),是充滿神一切的豐盛.以弗所書有一個很清楚的主題,貫徹在整卷的以弗所書之上,這個主題是什麼?就是「充滿萬有者的榮耀」──將神的豐盛充滿整個教會.

我們看看以弗所書　第一章第廿二至廿三節「又將萬有服在他的腳下,使他為教會作萬有之首. 教會是他的身體,是那充滿萬有者所充滿的.」然後第三章十九節「並知道這愛是過於人所能測度的,便叫神一切所充滿的,充滿了你們.」第四章沒有提到「充滿」這個詞,但在第一章我們看見基督是教會的頭,教會是基督的身體,所以如果教會是滿有基督長成的身量的時候,這就是「充滿」.我們再看第四章十三節「直等到我們眾人在真道上同歸於一,認識神的兒子,得以長大成人,滿有基督長成的身量」──所以這個主題「充滿」是以弗所書內的主題.我們不單可從以弗所書內看見這主題,目的,我們還可以在保羅的禱告內看見這目的.在第三章第十四節至第廿一節裡,第十六節「求他按著他豐盛的榮耀」──提到榮耀.最後的時候又提到榮耀,就是第廿一節「但願他在教會中,並在基督耶穌裡,得著榮耀,直到世世代代,永永遠遠」.當中包括了第十九節「便叫神一切所充滿的,充滿了你們.」所以保羅在禱告裡,神的「榮耀」和「充滿」是息息相關的.因神榮耀的緣故,所以神的豐盛能夠充滿教會；又因為神的豐盛充滿教會,神就得榮耀──所以這「榮耀」與「充滿」息息相關,「充滿」是神設立教會的目的.

在第三章中保羅講到他的使命,我們看見他將第二章所說的,神將外邦人和以色列人合而一個教會,在第三章中保羅稱之為「奧祕」──他的人生就是要使這奧祕能夠達成──保羅將外邦人得蒙基督測不透的豐富,稱之為「奧祕」,他深知這「奧祕」,並傳揚這「奧祕」為神的恩賜,而且他以被囚為榮,他說就連天使也要細察外邦人怎麼樣蒙神的恩典,教會──人得恩典是宇宙的中心.保羅忠於他所託,他深深掛念這奧祕的達成,這個達成就是神的「豐盛」能夠「充滿」教會,這是他禱告的中心,是他人生使命的中心,也是以弗所書整卷書的中心.

在他的禱告裡(第廿一節可以說是一個大結局,一個榮耀頌),廿一節說「但願他在教會中,並在基督耶穌裡,得著榮耀,直到世世代代,永永遠遠 and \#8230;」當神的「豐盛」能夠充滿教會的時候,就是大結局,就是「目的」了.

我們從以弗所書整卷書來看,從保羅的祈禱來看,或從保羅的人生使命來看:神的豐盛充滿教會是神設立教會的目的(是保羅的盼望).但要達成這個目的就需要有些步驟,就是剛才所講到的三個詞,三個不同的步驟.第一就是剛強──剛強有力(英文就是 ``empowered" 或是 ``strengthened with power" ).──神藉著祂的靈,使內心的力量剛強起來.在中文的翻譯裡,怎樣解釋「內心剛強」呢?是外面不剛強嗎?剛強不是描述一個人的嗎?讓我們看看聖經第六章「...作剛強的人.」是那個「人」作剛強. 約書亞記1:9 「 and \#8230;.你當剛強壯膽!」這剛強像似外面的表現,但也是由內心而發出的.因此提摩太後書1:7 說「因為神賜給我們,不是膽怯的心,乃是剛強,仁愛,謹守的心.」在原文裡有一個寫法說:靈是在我們的內人裡(無論你是看KJV 好,是NAS也好,是NIV好,全都是這樣說的),都是譯作「在內人裡的靈」使我們剛強有力:第一就是使「人」,第二就是有一個叫「內人」(inner  man), 當然「內人」是很難去解釋的,即便說在你內裡又一個人,何解?中文翻譯成「內人」就避開了這個問題.「內人」是否可解作「內心」?「內心」或「內人」在原文上,不是``heart",這是人性或人   Anthropos (Ανθροπος) ,用於聖經上也有很多不同的用法,在羅馬書6:6譯作「舊人」這個字,許多時候這個字只是表現一個範圍；在馬可福音7:15 「從外面進去的不能污穢人,惟有從裡面出來的乃能污穢人.」外面, 面都有``man"「人」這個字,所以這個「人」有時候是說一個範疇,很明顯這裡的內住不是說我們因信得救的時候,聖靈在我們身體裡,我覺得不是這個意思,因為這裡沒說到一個空間來安置聖靈,反而是反應聖靈工作果效的之所在.如果聖靈改變了我們的性格,思想,我們就說聖靈是我們的性格,聖靈在我們的思想裡.我們都知道聖經告訴我們聖靈有那些工作?聖靈有很多工作:包括提醒我們的記憶,在人的思想上引導人明白真理,或作引導,教訓,或安慰,又幫助人確認自己是神兒女的身份,又在人的良知,良心內動工,這全都是心理的範疇,所以我是以心理學的角度來研究信仰,如此,就能多一點明白聖靈的工作.因此我的解釋就是,藉著這些內心的渠道,影響一個人怎樣持守信仰,在他困難,危險,損失的時候的取向,在信仰中能表現出剛強,這外面的表現是從心而發的,假使按照這個角度來看,這「內人」也可以解釋為「內心」──都是從心而發的.這裡告訴我們是按著神的榮耀,怎樣按著神的榮耀使我們剛強呢?有何解釋?

羅馬書 8:16-18 給我們一個很完美的答案:「聖靈與我們的心同證我們是神的兒女；既是兒女,便是後嗣,就是　神的後嗣,和基督同作後嗣.如果我們和他一同受苦,也必和他一同得榮耀.」下面的一句很重要:「我想,現在的苦楚若比起將來要顯於我們的榮耀就不足介意了.」──如我們有人因信仰而受苦待,因親人朋友的反對受藐視,事業上受了損失,或是我們克制肉體覺得很辛苦,很付代價,但便如此,我們要堅定不移的信靠神,因為比起將來要顯於我們的榮耀是微不足道,不足介意.所以可以按著祂的榮耀,藉著祂的靈,叫我們剛強起來.以弗所書的寫法,好像這個「剛強起來」是被動的,完全是聖靈作主導的,沒人的責任,對不對?聖經真是這樣子寫的,像是完全由來作聖靈主導,不過我們可以借用其他的經文來看看,是否完全由聖靈來主導,完全沒人的責任?如果聖經說「人要」或「人當」的時候,這是很明顯是在說人的責任,對不?你看看聖經有沒有?有啊,在那兒出現?「要剛強」──舊約出現了六次,新約出現了一次；「當剛強」──舊約出現了十六次.除了「當剛強」,「要剛強」出現過之外,還有其他的說法,提摩太後書1:6說,我們有關的工作要持續和加倍,要「如火挑旺起來.」 哥林多前書 16:13,以弗所書6:10「 and \#8230;.要剛強,在真道裡站立得穩」；約翰一書 2:14 說「 and \#8230; and \#8230;.神的道常存在你們心裡；你們也勝了那惡者.」；希伯來書 12:3-5說「我們持守不灰心」；彼得後書說「我們更加殷勤,使恩照,揀選堅定不移」；歌羅西書說「堅定不移,不失福音的盼望」；以弗所書4:14「 and \#8230;..不被異教之風搖動,飄來飄去」 and \#8230; 這都說明人的責任,人是有責任的!沒錯,聖靈作工,人要配合.如果你說神是供給我們避難所,你也要去避才成,你不避便不成.如果說神給你兵器,你也要拿起兵器才成,兵器不會飛進你的手中.所以我們要配合:神有恩典,我們有責任配合.當我們配合的時候,我們真能做到聖經所說的:按照神的榮耀,藉著祂的靈使我們剛起來.

當我們在信仰中持守的時候,有什麼事情發生?我們便會有愛心,原來是信心產生愛心的.這兒是如此說的:「是基督因你們的信,住在你們心裡,叫你們的愛心有根有基」,這亦是聖經的道理.有一個次序的問題,先有信心,後產生愛心.加拉太書補充說「使人增發仁愛的信心」: 提摩太前書 1:5「但命令的總歸就是愛．這愛是從清潔的心,和無虧的良心,無偽的信心,生出來的.」所以我們有了這無偽的信心,持守這信心的時候,我們便慢慢地產生愛.但很奇怪,當聖經提到這「愛」的時候,它不是說聖靈,乃是基督因你們的信住在你們的心裡. 換了一個主角──基督,何解?難道基督跟聖靈是不同的?我們有兩個不同的存在我們心裡面?一個是靈,一個是基督?其實當聖經說聖靈是住在我們裡面的時候,亦等同基督也住在我們裡.是基督藉著靈,與我們同在,但重點就很不一樣. 首先我們談到前面第十六節所說的「心」是人性,而這裡所說的「心」是另一個字,是我們的思想和感情.基督住在我們的思想與感情裡,重點是不同的.當我們相信耶穌的時候,就是接納了神的愛,這位主所表現的,祂的降卑,祂的容忍,祂的犧牲,當我們想到主的時候,你想到什麼?你很容易想到這一位是愛的主,祂在我們的思想和我們的感情上,這自然產生了愛的反應,因為耶穌基督是信的對象和愛的樣式,所以當祂提到愛心的時候,就是說基督的內住.很多時候,當我們看聖經的時候,提到愛的時候,便提到耶穌基督為我們死,像羅馬書5:8 說「惟有基督在我們還作罪人的時候為我們死,神的愛就在此向我們顯明了.」聖經又說「我們能夠相愛,弟兄姊妹相愛,是神先愛我們」.所以這裡用基督住在我們心中,在我們的思想和感情裡產生了愛的果效.但奇怪的是,要持續我們的信仰,基督才內住,明顯的「使基督因你們的信,住在你們心裡」──如果我們的心不剛強持續的話,基督就不內住.所以這裡提到的「內住」,不是你一信了耶穌就有的內住,乃是說靈的工作和果效.因著靈的工作和果效,我們在思想和感情裡能夠使我們活出愛,這就是基督內住的果效了.

而這「住」是一個很特別的字.這「住」是長久永久作為家,耶穌基督在我們的心中是很舒服,能夠自主,作王的意思.如果我們在我們的思想感情中,這位愛的主是能夠作王,能夠作主導的話,當然我們活出來的,就一定有愛.大家試想想,人的愛是天生的,還是培養出來?不過從心理學的角度來說,人的感情有很多是「原感情」──即便一生下來就會產生的感情.有許多心理學家談及不一樣的「原感情」,但很多都沒提及「愛」,似乎「愛」是後來結合而產生的.一個人為何懂得愛?就是因為他被愛,他被愛以致能懂得愛.一個小孩子出生後,如果他不被愛,被擁抱,愛惜,他生存不了,會死掉的.大家可知道?那些有少許愛就能夠生存的動物(先說動物,不說人),牠完全不懂愛,牠是不懂得社交的,牠很孤伶伶的,躲在角落的；假使是女性的話,她不懂得性交,並非一出生便曉得,偶然是成功性交的話,她不會懂得育兒.所以愛小孩子也不是天生的,都是需要培養,是有後天的成份在內(當然有先天的成份).

當我們信耶穌的時候,這位愛的主,在我們的思想和感情裡的時候,我們就有愛的本能,能力,所以聖經說我們能在愛心裡有根有基.當我們真的有根有基的時候,我們就可以跟其他的信徒一同去活出愛,一同明白.一同明白並不是指著我們開聖經班(bible study) ,不是頭腦上的明白,因為下面第十九節都說「並知道這愛是過於人所測度的」──是我們思想測度以外所能夠明白的.所以「一同明白」並不是說我們一起去開研討會,一同查經來明白這愛是什麼?而是我們既然在愛裡有根有基,我們便一齊來彼此相愛,在從相愛裡明白愛是什麼?「一同明白」是有一個社交的因素的.宗教是個人的,還是社交性的呢?當然有個人的,亦有社交性的,沒社交性,人們怎把福音傳給你,你怎去相信?一定要有社交性,但當你實踐這宗教信仰的時候,亦有社交性,這愛便很清楚了.一同明白基督的愛是何等的長闊高深的.所以躲藏起來不去教會,不能作一個真正的基督徒,我們要去教會,過一個群體的生活,才能作一個基督徒,不過這社交似乎是一種方法,不是一個最終的目的,如果我們最終的目的是社交的話,這是錯誤的.

我們崇基裡有很多的配套可以給大家弟兄姊妹在相交,社交下互相扶持,互相認識,彼此相愛.基督的愛是要實踐才能明白的,這說法是很有根據的,第一,我們可從上下文中看,上面已說了「有根有基」,然後說我們可以與眾聖徒「一同明白」基督的愛.而第二是從原文裡看,在新約聖經裡應用這個字是有四個完全不同的譯法,第一個譯法就是「捉弄」(跟這個意思好像完全風馬牛不相及)；第二個是「捉拿」；第三個是「臨到」；第四個是「發現底蘊」(看出,查明等等的意思). 所以綜合這些說法這個詞語「一同明白」應該是:捕捉,得著,發現愛的底蘊.我們跟弟兄姊妹一同過彼此相交,愛的生活的時候,我們便能捕捉,得著和發現愛的底蘊.同時其他的聖經學者也都有這樣的解釋,John Stock說有能力去愛,領悟基督的愛,兩者是不能分開的,要明白祂的愛,唯一的方法,是愛他人.從一般的經歷裡,亦可以支持這說法:有很多人不是說,他做人家的兒女,他不明白父母的愛,唯有當他成為父母的時候,才發現父母是如何地愛他的.所以要去愛,一同去愛的時候,才能明白基督的愛是何等的長闊高深.是這麼的長──我們無論有多少次軟弱失敗,耶穌的愛,基督的愛都跟著我們──是愛到底的愛；是多闊──我們的愛可能只可容納少數的人,我們的親屬,或是仁人,或是義人,但基督的愛是包括那些敵對祂的人,如果我們能夠包容敵對我們的人,我們會有些兒明白基督的愛,是何等的闊,是何等的高,沒窮盡的.你所有的,接受的愛,還可有更大的愛,是沒盡頭的；有多深──無論你墜落的地步,基督的愛都追隨著你,不會使你失落.所以我們可以一同明白基督的愛是何等的長闊高深,當我們一同明白的時候,其他人的愛就可使我們明白基督的愛.比方說,有一個人離開了教會很久了,但有弟兄姊妹帶他回教會,用了很多時間照顧他,關心他,為他禱告,這樣的愛心使他感受到神的愛,是從人的愛之中明白神的愛這點保羅是最能夠明白的,因為他從前是逼迫教會的,是罪人當中的罪魁,他感受到比眾聖徒當中的最小還小,但神都賜恩典給他,他感到這是神的大愛,甚至你給別人傷害的時候,你都可以明白基督的愛,如果教會裡有一個人,你很難忍受他的,你要容忍他,這時候你就明白基督的愛.

在很多年前,我在大專學校裡做團契的事奉工作,我花了許多心血,投放了許多時間,而且栽培了一個低班的同學接我的班,做了一個新的模式出來,用上了很多時間,但這位被我栽培出來的弟兄,後來反對我,杯葛我,而且聯合了一些導師一齊要復舊,即是丟棄新的制度,我就好像是一個小孩子在沙灘上用泥沙砌了一個堡壘,另一個小孩子就跑過來一腳踢倒了,這小孩子的失望情緒是何等大!當我與導師開完會回教室時,旁邊的同學問道「你病了?」我真是心裡有很多壓力,病了.當我在禱告的時候,在神面前我很憤怒,在很憤怒中,聖靈便提醒我,聖經的說話,耶穌基督說「同我吃飯的用腳踢我」當耶穌基督這樣說,是說猶大出賣祂,當時我頓然多多少少明白多一點耶穌基督的愛.儘管在教會裡生活你受氣,你如果因為受氣而離開教會是愚蠢的,你如果受氣的時候,你會因為這種的受氣而更加明白基督的愛.所以無論如何,我們是一同明白基督的愛是何等的長闊高深.接著,祂還說,「知道這愛是過於人所能測度的」,是什麼意思?有兩個可能性,一個是我們剛才說過的,這份愛不是你本身會明白的,是要經歷才能明白的.而第二個可能性就是說,根本上是很大的,是無窮無盡,你揣摩不到,你明白不了的,究竟是什麼意思呢?其實亦可兩者都有.從一個例子看見兩者都有:有沒有聽過有人說,牛頓說人怎樣發現知識?他說他就像一個小孩子在沙灘裡找一些圓滑的石頭,而他未曾探討的知識,就像汪洋大海一樣,是未曾接觸的.當科學家去發現出一些自然的真理的時候,他發現了一些真理,道理,但同時他亦覺悟到那個未曾發現汪洋大海的道理.

當我們因為經歷的緣故,去明白基督的愛時,我們可以揣摩到基督的愛,但同時我們會感受到基督的愛好像汪洋大海一樣 ``out of  reach" -- 還不能明白,是無窮無盡的.當我們在此境況的時候,聖經說:「便叫神一切所充滿的,充滿了我們」.當我們這樣彼此互相相愛的時候,就能揣摩到耶穌基督的愛的時候,神的豐盛就充滿了我們教會,保羅很多時候他愛主,事奉主的動機,都是想到主的愛,哥林多後書 5:14-15「原來基督的愛激勵我們；因我們想,一人既替眾人死,眾人就都死了,並且他替眾人死,是叫那些活著的人不再為自己活,乃為替他們死而復活的主活.」所以當我們有一同明白基督的愛,而我們能夠揣摩,又能看見原來基督的愛是無窮無盡的時候,我們應該有很大的動機愛主,起來事奉,在這時候神的豐盛便充滿教會.你看看,如果每一個人都明白基督的愛,每個人都感受到基督的愛是無窮無盡,因為這緣故有很大的動機,起來愛主和事奉,教會怎會不興旺?對不?所以這裡說「便叫神一切所充滿的,充滿了你們」.當然我們再看看第四章的時候,看見怎樣因愛的緣故,使神的豐盛充滿整個教會.第四章經文中提到要「竭力保守聖靈所賜合而為一的心」,「唯用愛心說誠實話,凡事長進,連於元首基督 and \#8230; 便叫身體漸漸增長,在愛中建立自己.」所以我們看見聖經上一個很清楚的道理:我們要神的榮耀彰顯在教會當中,我們就要彼此相愛,從愛裡看見基督的愛,當我們看見基督的愛是何等的長闊高深,當我們看見是無窮無盡的,是遠超過我們所能明白的,因此而激勵出每一個信徒積極愛主,事奉,教會就有不同的境況.

剛才我們所說的,我們看見要使神的豐盛充滿教會,是神設立教會的目的,我們從以弗所書內我們看見,從保羅的人生使命看見,我們從保羅的祈禱我看見這目的,要達成這目的,我們有三個步驟,大家說這三步驟是什麼?第一是兩個字「剛強」；第二個是「基督內住」；第三個「一同明白」,讓神的真理在我們打開的時候,我們便得到亮光,使我們能夠有智慧.

我們祈禱:天父,在這時候,我們好像保羅這樣祈禱說,求你的靈在我們內心工作,使我們剛強起來,以致基督能夠在我們的感情和思想裡作主,使我們一同明白基督的愛是何等的長闊高深,當我們一同活出這愛的團體的生活的時候,又明白基督的愛是何等的遠超過我們能明白的時候,神的豐盛就充滿了整個教會,讓這應許成就在香港的教會裡,我們祈禱奉主耶穌基督的名,阿們.



\newpage

\section{位置決定我們的腦袋}
\label{sec:166}
\textbf{伍渭文牧師 - 位置決定我們的腦袋}
\newline
\newline
\hyperref[sec:165]{< < < PREV SERMON < < <}
~
\hyperref[sec:toc]{[返目錄]}
~
\hyperref[sec:167]{> > > NEXT SERMON > > >}


雅各書 2:1-17
\newline
\begin{longtable}{cl}
\hline
\hline
章節 & 經文 (和合本修訂版)\\
\hline
2:1 & \begin{tabularx}{0.7\textwidth}{X} 我的弟兄們，你們信奉我們榮耀的主耶穌基督，就不可按著外貌待人。 \end{tabularx} \\ \\ \relax
2:2 & \begin{tabularx}{0.7\textwidth}{X} 若有一個人戴著金戒指，穿著華麗的衣服，進入你們的會堂，又有一個窮人穿著骯髒的衣服也進去， \end{tabularx} \\ \\ \relax
2:3 & \begin{tabularx}{0.7\textwidth}{X} 而你們只看重那穿華麗衣服的人，說：「請坐在這裡」，又對那窮人說：「你站在那裡」，或「坐在我腳凳旁」； \end{tabularx} \\ \\ \relax
2:4 & \begin{tabularx}{0.7\textwidth}{X} 這豈不是你們偏心待人，用惡意評斷人嗎？ \end{tabularx} \\ \\ \relax
2:5 & \begin{tabularx}{0.7\textwidth}{X} 我親愛的弟兄們，請聽，神豈不是揀選了世上的貧窮人，使他們在信心上富足，並承受他所應許給那些愛他之人的國嗎？ \end{tabularx} \\ \\ \relax
2:6 & \begin{tabularx}{0.7\textwidth}{X} 你們卻羞辱貧窮的人。欺壓你們，拉你們到公堂去的，不就是這些富有的人嗎？ \end{tabularx} \\ \\ \relax
2:7 & \begin{tabularx}{0.7\textwidth}{X} 毀謗為你們求告時所奉的尊名的，不就是他們嗎？ \end{tabularx} \\ \\ \relax
2:8 & \begin{tabularx}{0.7\textwidth}{X} 經上記著：「要愛鄰 如己」，你們若切實守這至尊的律法，你們就做得很好。 \end{tabularx} \\ \\ \relax
2:9 & \begin{tabularx}{0.7\textwidth}{X} 但你們若按外貌待人就是犯罪，是被律法定為犯法的。 \end{tabularx} \\ \\ \relax
2:10 & \begin{tabularx}{0.7\textwidth}{X} 因為凡遵守全部律法的，只違背了一條就是違犯了所有的律法。 \end{tabularx} \\ \\ \relax
2:11 & \begin{tabularx}{0.7\textwidth}{X} 原來那說「不可姦淫」的，也說「不可殺人」。你就是不姦淫，卻殺人，也是成為違犯律法的。 \end{tabularx} \\ \\ \relax
2:12 & \begin{tabularx}{0.7\textwidth}{X} 既然你們要按使人自由的律法受審判，就要照這律法說話行事。 \end{tabularx} \\ \\ \relax
2:13 & \begin{tabularx}{0.7\textwidth}{X} 因為對那不憐憫人的，他們要受沒有憐憫的審判；憐憫勝過審判。 \end{tabularx} \\ \\ \relax
2:14 & \begin{tabularx}{0.7\textwidth}{X} 我的弟兄們，若有人說自己有信心，卻沒有行為，有甚麼益處呢？這信心能救他嗎？ \end{tabularx} \\ \\ \relax
2:15 & \begin{tabularx}{0.7\textwidth}{X} 若是弟兄或是姊妹沒有衣服穿，又缺少日用的飲食； \end{tabularx} \\ \\ \relax
2:16 & \begin{tabularx}{0.7\textwidth}{X} 你們中間有人對他們說：「平平安安地去吧！願你們穿得暖，吃得飽」，卻不給他們身體所需要的，這有甚麼益處呢？ \end{tabularx} \\ \\ \relax
2:17 & \begin{tabularx}{0.7\textwidth}{X} 信心也是這樣，若沒有行為是死的。 \end{tabularx} \\ \\
[1ex]
\hline
\hline
\end{longtable}
日期:2012年10月7日

題目:位置決定我們的腦袋

證道:伍渭文牧師

經文:雅各書2:1-17

今年6月發射升空的神九號宇航員景海鵬,劉旺以及中國首位女航天員劉洋訪問香港時,其中一位說:從太空望看香港,景色美極了!翌日有報紙即說,這是奉承香港人逗我們開心的假話,從太空那麼高,根本不能看到香港.有人說,在外太空的白天,我們唯一認得的地上建築是萬理長城；在晚上,我們可以見到美國拉斯維加斯,因為這賭城是從大遍沙漠中打造出來,晚上燈火通明,而四周沙漠漆黑,可以識別.

其實,從外太空望向地球,地球不同的高山,祇是不同圖案,不同顏色的平面.從至高上帝的眼中看每一個人,都是同樣的尊貴,我們不可按外貌偏待人.我們已與基督一同坐在天上:「他又叫我們與基督耶穌一同復活,一同坐在天上.」(弗二6)不單不要輕視窮困的人,更要關心他們的需要,我們的位置,決定我們的腦袋,我們要有屬靈的眼光看事物.

雅各書的主題談及邁向成熟的靈命,成熟的屬命有透析事物的能力,不被事物表象所迷惑.雅各書一開始就提及人生中我們有很多試煉,在人生的舞台我們會面對很多不同的爭戰:「我的弟兄們,你們落在百般試煉中,都要以為大喜樂；因為知道你們的信心經過試驗,就生忍耐.」(雅一2)在人生種種爭戰中,魔鬼誘惑我們離棄上帝,這是試探；但上帝要藉此熬煉我們的信心,這就成為試煉,試探和試煉是同一個字.

作者要我們要有屬靈的識見:在人生舞台的戰爭中,要留意人生舞台的後台,我們內心的靈慾之爭,因為勝負在此決定.我們要用上帝的道,來治死內心盈餘(prevalent)的邪惡.「惟有詳細察看那全備,使人自己之律法,並且時常如此,這人既不是聽了就忘記,乃是實行出來,就在他所行的事上必然得福.」(雅一25)要人生舞台得勝,秘訣是實行聖經的教訓,克制內心的邪惡,使靈命成長.

一個成熟優美的靈性有三個記號:勒住舌頭(第三章) ,看顧窮乏人(第二章),不隨從世俗(第四章).「若有人自以為虔敬,卻不勒住他的舌頭,反欺哄自己的心,這人的虔誠是虛的.在上帝我們的父面前,那清潔沒有玷污的虔誠,就是看顧在患難中的孤兒寡婦,並且保守自己不沾染世俗.」(雅一26, 27) 這三個記號都和屬靈的識見有關. 勒住舌頭是不隨意作論斷；不要按外貌待人,就是不要按表象作判斷；隨從世俗就是祇用今生來作判斷,忘記來世的審判,也忘記上帝可以介入今生,干預我們的規劃.若我們已與基督一同坐在父上帝的右邊,就要有屬靈的洞見；位置決定我們的腦袋.

今天的經文提醒我們,要有屬靈的洞見,不要按外貌待人.

一   認識真正的榮耀

「你們信奉我們榮耀的主耶穌,便不可按外貌待人.若有一個人帶著金戒指,穿著華美的衣服,進你們的會堂去；又有一個窮人穿著骯髒衣服也進去；你們就看重那穿華美衣服的人,說『請坐在這好位上』；又對那窮人說『你站在那裡』,或『坐在我腳凳下邊.』這豈不是你們偏心待人,用惡意斷定人嗎?」(雅二1-4)

有人說:教會對社會最大的貢獻,是活出與世界不同的價值觀,令世界不能自圓其說,教會以基督信仰,打開了閉封的世界.當世界說:「我們吃吃喝喝吧,因為明天要死了.」(林前十五32) 保羅就說我們都要復活,「所以,我親愛的弟兄們,你們務要堅固不可搖動,常常竭力多作主工；因為知道,你們的勞苦在主裡面不是徒然的.」(林前十五58)

我們雖然成為天國公民,但沒有好好接受天國公民教育,常常受內心盈餘的邪惡影響我們的價值觀,世界先敬蘿衣後敬人,教會也如是.會堂是敬拜上帝,聆聽聖道的地方,但我們用惡意以外貌待人,對衣飾華美的奉承,對衣著寒酸的冷落.我們以為華美悅目是榮耀,忘記基督給我們的榜樣:服事他人和憐恤有需要的人才是真正的榮耀.

榮耀來自服事和犧牲

我們要注意聖經特別提及榮耀的主耶穌,來對比華美的衣服和金戒指.耶穌的榮耀就是他的紆尊降貴,他無表象的佳形美容,叫我們羨慕他,耶穌親近被人輕看的稅吏和妓女,和他們一起吃飯.祇要他們願意尋求上帝的國,耶穌對他們滿有憐憫,和他們做朋友,不避嫌疑,不怕批評.耶穌的榮耀就是他的憐憫,耶穌說:人子得榮耀的時候到了,得榮耀就是被高舉釘十字架的時候；在十字架,上帝顯出憐憫,也顯出他的榮耀.

憐憫和審判是不同的,審判是大公無私,按著我們客觀的價值作判斷,法官判案,不看人面,在法院門前代表公正的人像,手裡拿著天秤,眼睛蒙布,審判祇看我們客觀的價值(Merit),但憐憫則看我們主觀的需要.看客觀分數,我們並未達標,拿不到這工作,但我們因為經濟原因,極需要得到這份工,雇主出於憐憫,雇用了我們.

上帝的恩慈和憐憫,也是上帝的榮耀

當上帝囑咐摩西要引領以色列人離開埃及,摩西求上帝顯出他的榮耀,好叫百姓看到憑據,曉得他是上帝差來的.「摩西說:求你顯出你的榮耀給我看.耶和華說:我要恩待誰就恩待誰；要憐憫誰就憐憫誰.」為何求顯榮耀而上帝說恩慈和憐憫呢?因為上帝以他的憐憫為他的榮耀.「你不能看見我的面,因為人見我的面不能存活.耶和華說:看哪,在我這裡有地方,你要站在磐石穴中.我的榮耀經過的時候,我必將你放在磐石穴中,用我的手遮掩你,等我過去,然後我要將我的手收回,你就得見我的背,卻不得見我的面.」甚麼是上帝的背呢?為何我們祇能見到上帝的背呢?很多解經家都認為上帝的背,就是道成肉身的基督,藉著耶穌的犧牲,我們才能來到父神面前.基督是上帝的恩慈和憐憫,也是上帝的榮耀.

二   認識憐憫隱含公義

遵守愛人如己的律法

在上帝眼中,榮耀不在華美的衣服和金戒指,乃效法基督,憐恤有需要的人.所以按外貌待人,就是不遵守愛人如己的律法.「經上記著說:要愛人如己.你們若遵守這至尊的律法才是好的.但你們若按外貌待人,便是犯罪,被律法定為犯法的.」(雅二8, 9) 愛人如己是律法,律法是責任,律法不是感情,不由我們喜歡不喜歡,是律法就要遵守.我相信沒有人收到超速駕車的告票,會欣然接受,都會看成今天倒霉了,無論你開不開心,也要交罰款.愛人如己是律法,指出這誡命超越我們個人的情感喜好,是道德的責任,所以愛人如己也涉及公義,公義是均分的愛.

事實上,不少志願者從憐恤救援開始,轉到追求爭取公義結束.作為醫療救援志願者,起初照顧在醫院的病友,跟著關心公共醫療政策,進而連繫到環境保護,追溯到資源分配不公,思考公義問題.

「經上記著說:要愛人如己.你們若全守這至尊的律法才是好的,但你們若按外貌待人,便是犯罪法的.因為凡遵守全律法的,祇在一條上跌倒,他就是犯了眾條.原來那說不可姦淫的,也說不可殺人；你就是不不姦淫,卻殺人,仍是成了犯律法的.」(雅二8-11) 在列出的不可淫姦和不可殺人兩條律法,我們信徒大多不會觸犯(sin of commission),因有政府的刑法對付我們,這兩條律法屬於個人倫理範疇,但愛人如己則屬於社會倫理,不愛人如已也不觸犯刑法,我們比較忽略(sin of omission),但聖經說,這是至尊的律法.

解經家提出經文的處境,不是進行崇拜的會堂,可能是進行訴訟的會堂.「那富足人豈不是欺壓你們,拉你們到公堂去嗎?」(雅二6) 進行訴訟時,我們更不可按外貌待人,對權貴奉承,欺負沒有權勢的貧窮人.

華美衣服與權力

為何人喜歡華美的衣飾?因為華美的衣飾代表財富,而財富代表擁有權力.財富和權力是孖生兄弟,難怪國家主席胡先生在深圳時,公開表示想見見香港最富有的李先生.但教會作為信仰群體,小心靠攏權貴,教會對社會最大的貢獻,是活出與世界不同的價值觀.貧窮不一定指財物的貧窮,也可以指權力的貧窮.當世界要管制人的思想,要人順從專制的國家機器時,教會應仗義執言,為異見份子提供保護；教會要為執政掌權者祈禱,也為受壓迫的人民祈禱；教會要尋求真相,也要推動修和解怨,棄絕清算報復.

兩年多前第四屆崇基基督教文化節的主題為「使人和睦」,我最初邀請程翔先生講演,但他推卻說「我的靈命尚淺,未能分享」,但當我待他訪問前東德地區教會後,特別是萊比錫(Leizip)的聖尼古拉斯教堂(St. Nicholas Church)回港後,再度邀請他,他立刻答應.圍牆倒下前,聖尼古拉斯教堂一連十三個星期一舉行和平禱告會,每次幾千名群眾聚集.東德是共產鎮壓得非常厲害的國家,唯獨教會讓人有自由的思想和行為.該教會牧師以登山寶訓八福開始帶領禱告會,以和平,理性,非暴力的形式爭取自由,群眾亦以八福克服恐懼.在圍牆倒下前一個月,群眾層層包圍著教會,坦克則層層包圍著群眾,當時國防部說:「我們已準備一切,但沒有想到燭光和禱告」(We are ready for everything, except candles and prayer) .統治者面對如此的群眾,根本無法流血鎮壓.這件事傳遍德國,一個月後柏林圍牆便倒下了.程翔指東西德統一背有很多複雜的因素,但也確實是一件以祈禱開始的奇蹟.

作為教會,我們要追求公義和散播憐恤.我們在基督裡的位置,決定我們的信念.

三   認識真正的富足

「我親愛的弟兄們,請聽,上帝豈不是揀選了世上的貧窮人,叫他們在信上富足嗎,並承受他所應許給那些愛他的人的國嗎?」(雅二5) .「卑微的弟兄升高,就該喜樂；富足的降卑,也該如此.」(雅一9)

真正的富足不是穿華美衣服,戴金戒指,乃是信仰上的富足.因為信仰基督,貧窮的弟兄已升到天上,和基督一同坐在父神的右邊；富有的人,因著悔改歸信基督,也得謙卑下來.信上富足不單給貧窮人,也給富有的人,要是沒有信仰,富人也窮得祇剩下錢.

追上的富足不受限制

信上富足的人,雖然身體殘缺,但活得精彩.生下來沒有手腳的力克.胡哲(Nick Vujicic)憑藉著自己的努力,學會了游泳,打高爾夫球,衝浪,踢足球,樂觀幽默,成為甚受歡迎的演講家.他說:「我們常常覺得人生真不公平,然而《聖經》說:「你們落在百般試煉中,都要以為大喜樂.」那是我掙扎了許多年想學會的功課.我在YouTube上最受歡迎的影音,是一段拍攝了我在溜滑板,衝浪,聽音樂,打高爾夫球,跌倒,爬起來,跟聽眾說話,跟許多很棒的人擁抱等等的影片.總的來說,這些不過就是一般人也做得到的平常之事,不是嗎?那你覺得為什麼這樣的影片卻可以吸引幾百萬人點閱呢?我認為,大家之所以想看,是因為儘管生理上有重重限制,我卻活得似乎完全不受限.」他寫的一本書〈神采飛揚〉: 沒有肢體的生命-從無肢體到無限制(Verve: Life Without Limbs- From NO Limbs to No Limits) 成為多種譯本的暢銷書.今年二月力克結婚,娶了一位漂亮的太太.

信上的富足永恆不變

「若是弟兄或是姊妹,赤身露體,又缺了日用的飲食；你們中間有人對他們說:平平安安地去吧!願你穿得暖吃得飽,卻不給他們所需用的,這有甚麼益處呢?這樣,信心沒有行為就是死的.」(雅二15)

以外貌待人,見到穿華美衣服的給他們上痤,衣服寒酸的坐在自己的腳下,是重視財富.本來財富是用來分享給有需要的人,但這些人見到有需要的窮人竟然口惠而實不至,這些人是沾染世俗,他們忘記有一天要站在上帝的審判台前接受審判.成熟的靈命其中一個記號是不沾染世俗.世俗主義是甚麼?世俗主義就是今生主義,世俗(secular)一字來自拉丁文今世( saeculum) .祇有今生,沒有來世,沒有想到來世上帝的審判,也把上帝排除在今生的規劃,就是沾染世俗.世俗主義因為祇有今生,所以注重消費享樂:「你們求也求不著,是因為你們妄求,要浪費在你們的宴樂中,你們這些淫亂的人哪,豈不知與世俗為友的,就是與神為敵了.」(雅四3)以外貌待人或看事物,看產品,也是典型的消費主義,很多廣告就以外貌吸引客戶,鼓動即時消費.所以看到弟兄的需要而沒有憐憫的心,口惠而實不至,要面對上帝的審判.「因為那不憐憫人的,也要受無憐憫的審判.」(雅二13)

在第四章討論世俗主義時,作者提醒這些沒有憐憫的富人:「嗐!你們有話說:今天明天我們要往某城裡去,在那裡住一年,做買賣得利.其實明天如何,你們還不知道.你們的生命是甚麼呢?你們原是一片雲霧,出現少時就不見了.你們祇當說:主若願意,我們就可以活著,也可以作這事,或做那事.現今你們竟張狂誇口；凡這樣誇口都是惡的.人若知道行善,卻不去行,這就是他的罪了.」(雅四13-17) 世俗主義不相信來世,祇有今生,而且把上帝排除在今生的規劃.但上帝可以在今生干預我們,不需等到來世才與我們算帳,我們要知道,謀事在人,成事在天.

如果大家留意,在報告事項中,九月二十三日下午在這裡有紀念中文大學音樂系主任麥嘉倫教授(Professor Mcclellen) 的追思會；麥教授為學生同事喜愛的教授,每年的畢業禮在這教堂作司琴,他突然因感冒菌入侵腦囪部,在家中倒下失救而離開的,終年五十三歲,非常意外,令我們悲傷惋惜.真的,「其實明天如何,我們還不知道.」「主若願意,我們就可以活著,也可以作這事,或做那事.」

變幻永恆-何輝的故事

我們誇耀的財富,才幹,地位,健康,甚至家庭幸福亦可一夜失去,這些看來閃爍的榮耀像肥皂泡沫,看似彩色繽紛,但要用力握著的時候,可能爆破,離我們而去.祇有永生,祇有作在主身上的才有永恆的價值.這星期收到校長傳給我一個電郵,推介一位十二月到訪中大的中學同學何輝,可否安排他到同學中間作見證,十二月聖誕期間同學多回家了,我預備十二月二十三日的崇拜講道前請他作見證,他的見證在恩雨之聲廣播過,該集取名〈變幻永恆-何輝的故事〉.

何輝是大學優才生,體育健將,美國跨國公司行政總裁,口才無礙,擁有二萬呎豪宅.但先後兩次無情的交通意外,不單奪去了何輝的事業和他的健康,更奪去了他摯愛的太太.信主前,何輝總是說說就淚流滿面,籍著牧師的探訪和在明白聖經,他重新建立他對人生的盼望,決意做一個殘而不廢,痛而不苦的開心快樂人.明白到曾經的擁有,物質的享受,都如美麗的肥皂泡一般,看起來五光十色,但觸摸起來,一個個都會破滅,無法在永生中被紀念.何輝用步行器學習走路,重新學講話,把他著聖經金句的衣服穿上.許多都會認真地讀這些句子,有的甚至停下來和他交談.他說:我以前只會自己顧自己,喜愛物質喜愛金錢.但是信耶穌後就不同了,什麼事都先惦記著別人.

結論:今天是選舉日,藉著選舉活動,我們關心這城市,但教會對世界最大的幫忙,就是提出不同的看法,與世界有分別就是對世界最大的貢獻.這段日子,反對國民教育運動弄得全城沸沸揚揚,我們懼怕子女被錯誤的思想洗腦；但作為信徒,我們有否警覺世俗主義-祇顧今生,不顧來世,把上主排出我們的規劃的世俗主義-已靜悄悄給我們洗腦呢?

讓我們看清楚真正的榮耀:服事他人,憐恤他人,追求公義,就是真正的榮耀.

讓我們看清楚真正的富足:信仰的富足,心靈的富足才是真正的富足.

讓我們看清楚甚麼是真正的永恆:在基督裡所作的才是永恆,才有永恆的價值.



\newpage

\section{天國的國民教育}
\label{sec:167}
\textbf{鄧瑞強博士 - 天國的國民教育}
\newline
\newline
\hyperref[sec:166]{< < < PREV SERMON < < <}
~
\hyperref[sec:toc]{[返目錄]}
~
\hyperref[sec:168]{> > > NEXT SERMON > > >}


馬可福音 8:27-38
\newline
\begin{longtable}{cl}
\hline
\hline
章節 & 經文 (和合本修訂版)\\
\hline
8:27 & \begin{tabularx}{0.7\textwidth}{X} 耶穌和門徒出去，往凱撒利亞‧腓立比附近的村莊去。在路上，他問門徒：「人們說我是誰？」 \end{tabularx} \\ \\ \relax
8:28 & \begin{tabularx}{0.7\textwidth}{X} 他們對他說：「是施洗的約翰；有人說是以利亞；又有人說是先知中的一位。」 \end{tabularx} \\ \\ \relax
8:29 & \begin{tabularx}{0.7\textwidth}{X} 他又問他們：「你們說我是誰？」彼得回答他：「你是基督。」 \end{tabularx} \\ \\ \relax
8:30 & \begin{tabularx}{0.7\textwidth}{X} 於是耶穌切切地囑咐他們不可對任何人說起他。 \end{tabularx} \\ \\ \relax
8:31 & \begin{tabularx}{0.7\textwidth}{X} 從此，他教導他們說：「人子必須受許多的苦，被長老、祭司長和文士棄絕，並且被殺，三天後復活。」 \end{tabularx} \\ \\ \relax
8:32 & \begin{tabularx}{0.7\textwidth}{X} 耶穌明白地說了這話，彼得就拉著他，責備他。 \end{tabularx} \\ \\ \relax
8:33 & \begin{tabularx}{0.7\textwidth}{X} 耶穌轉過來看著門徒，斥責彼得說：「撒但，退到我後邊去！因為你不體會神的心意，而是體會人的意思。」 \end{tabularx} \\ \\ \relax
8:34 & \begin{tabularx}{0.7\textwidth}{X} 於是他叫眾人和門徒來，對他們說：「若有人要跟從我，就當捨己，背起自己的十字架來跟從我。 \end{tabularx} \\ \\ \relax
8:35 & \begin{tabularx}{0.7\textwidth}{X} 因為凡要救自己生命的，必喪失生命；凡為我和福音喪失生命的，必救自己的生命。 \end{tabularx} \\ \\ \relax
8:36 & \begin{tabularx}{0.7\textwidth}{X} 人就是賺得全世界，賠上自己的生命，有甚麼益處呢？ \end{tabularx} \\ \\ \relax
8:37 & \begin{tabularx}{0.7\textwidth}{X} 人還能拿甚麼換生命呢？ \end{tabularx} \\ \\ \relax
8:38 & \begin{tabularx}{0.7\textwidth}{X} 凡在這淫亂罪惡的世代，把我和我的道當作可恥的，人子在他父的榮耀裡與聖天使一同來臨的時候，也要把那人當作可恥的。」 \end{tabularx} \\ \\
[1ex]
\hline
\hline
\end{longtable}
日期:2012年10月14日

主題:天國的國民教育

證道:鄧瑞強博士

經文:馬可福音 8:27-38

各位弟兄姊妹,主內平安.

近來,國民教育的問題,在香港引起軒然大波.國民教育,可以很複雜,但歸根究底,牽涉的,卻是我們生命裡的一些基本問題.諸如:我們的生命歸屬哪個社群?這個社群的界線何在?這個社群以什麼維繫?這個社群如何生活?我們在這社群裡要扮演什麼角色?這個社群終歸要如何走下去?其價值理想是什麼?回答了這些問題,才能講出要進行什麼形式和什麼內容的教育,好能在現在及將來,活出這個社群最美好的一面.

「我們的生命歸屬哪個社群?」是最基本的一個問題.社會上的權力結構會塑造人歸屬某種類型的社群.我們自己的成長背景可能會引導我們認同另一種類型的社群.每一個都需要認真思考,我們終歸以哪個身分作為自己的身分認同,這最終涉及我們將生命投身於哪種命途.

(1)生命誰屬問題?

有一次,在人生的路上,耶穌與他的門徒有一段對話,這對話終歸涉及到生命誰屬的問題.

馬可福音 8:27-29

27     耶穌和門徒出去,往該撒利亞腓立比的村莊去.在路上問門徒說,人說我是誰.     28     他們說,有人說,是施洗的約翰．有人說,是以利亞．又有人說,是先知裡的一位.     29     又問他們說,你們說我是誰．彼得回答說,你是基督.

門徒跟了耶穌一段日子,他們知道自己跟誰嗎?

上個禮拜日是選舉立法會議員的投票日.記者訪問一位婆婆,這婆婆的回答,大意是這樣的:「我不識字的,我隨便投的.」噢,她不知道她選擇了誰,她也不知道她將社群生活的重大決定權交給了誰.

門徒跟了耶穌很久,他們投票支持耶穌,但,他們知道他是誰嗎?

主耶穌問他的門徒,一般人認為他是誰?

有人認為他是施洗的約翰.施洗的約翰剛剛推展了一場宗教復興運動,他呼喚人悔改,迎接天國的來臨.他叫有食物的,分給沒有的.用現在的講法,即是捐食物給食物銀行,叫窮人有得吃.他強調一種正義的生活.他指出了君王的道德腐敗,甚至因此招致殺身之禍.主耶穌做的,與這施洗的約翰做的,的確有點相似,但卻多於施洗的約翰.

有人認為耶穌是以利亞,以利亞是一位預告新世代會來臨的先知.在以利亞的時代,社會上有太多假先知,引人走入迷途.以利亞不惜孤軍與社會佔大多數的假先知對決.後來,他不用經歷死亡,回到神的懷抱.主耶穌做的,與這以利亞所做的,的確有點相似,但卻多於以利亞.

有人認為耶穌是先知裡的一位,誠然,耶穌是一位偉大的先知,但耶穌不是一位平凡的先知.

耶穌再問門徒,「你們說我是誰?」

在現世,有施洗的約翰,以利亞,先知,不是足夠了嗎?有人追求社會的公義,有人抗衡世間的假先知,有人預告美好的明天,這不是足夠了嗎?

似乎,耶穌認為不夠.似乎,有些東西,比這一切更重要.似乎,有些東西,更值得我們將生命投身下去.

彼得回答說:「你是基督.」

在耶穌時代,「基督」這個稱號,是泛指一些從神而來的人物,這些人改變世間的一切權力結構,帶來永恆的公義與和平.耶穌之後,基督教的信仰群體,將「基督」這稱號用來專指耶穌,意味著他是從神而來的特別的一個人,意味著他是神自己的生命的顯露.基督讓我們看到神的真實性,體驗到神的愛,意識到神的永恆.

彼得回答說:「你是基督.」

在基督教信仰群體,我們大可將這答案理解為:「你是神.」

這答案是震撼的.當我們在塵世事物中尋找可依靠的人和物時,彼得直指問題核心,我們需要神.耶穌是神,我們渴望的,不是人間的種種,而是神自己.

在人世間,我們有很多追求,但我們追求神嗎?在人世間,我們有所歸屬.我們歸屬一個民族,歸屬一個社群,歸屬一個國家,但,我們意識到,我們歸屬神嗎?

與基督相遇,讓我們明白,我們要努力此生,但此生並不終極.我們要努力愛我們生活其中的社群,但這社群並不終極.塵世的一切,無論多好,都不是神.但歸屬神的人,會努力建造塵世,讓塵世更多屬於神.

在神的世界裡,獅子和羔羊是朋友,大地理應這樣和諧.但是,當塵世的人自以為是神的時候,災難便開始了.正如「動物農莊」(Animal Farm)一書所提示的,當塵世的豬認為自己是神,且教導一群走狗接受這道理時,人間的災難便開始了.神化自己,且教導人支持這神化,必會剷除異己,並且大義凜然地剷除異己.在這時,我們的確要問一問:我們屬誰?我們能否指出:耶穌是基督,唯有他是神.

曾經,在德國納粹獨裁的年代,他們的「國民教育」教育人民完全支持政府,很多基督徒也接受這種教育,以納粹為身分認同對象.但卻有一群神學家,發表＜巴門宣言＞(Barmen Declaration),指出唯有耶穌基督是主,任何世間權力要求人的絕對效忠,都只是一種偶像崇拜.這個宣言,在德國的信仰群體中,燃點起一枝燭光.

(2)天國的國民教育方向

歸屬耶穌基督的人,是屬於天國的.

天國有天國的國民教育,它指出天國的價值.

一個好的國民教育,不單教人維持現狀,也應教人實現種種價值,縱使這些價值在現實中仍未出現.

正如我們教一個小孩子,我們不單教他處理現狀,維持現況,也應教他如何在人生中開展真正的價值.好的教育,應向孩子展示價值的大路.沒有價值內涵,則沒有真正的明天.

馬可福音 8:31-33

31     從此他教訓他們說,人子必須受許多的苦,被長老祭司長和文士棄絕,並且被殺,過三天復活.     32     耶穌明明的說這話,彼得就拉著他,勸他.     33     耶穌轉過來,看著門徒,就責備彼得說,撒旦,退我後邊去罷．因為你不體貼　神的意思,只體貼人的意思.

耶穌對他的門徒,進行天國的「國民教育」.

屬於天國的人,其美善總會抗衡著塵世的醜惡.他追求公義,但公義不是神,不會假借公義之名而打擊對手.他追求和諧,但和諧不是神,不會假借和諧之名而消除異己.他追求社會發展,但社會發展不是神,不會假借社會發展之名而剷除追求不同價值的另類社群.他的生命歸屬神,不受塵世事物變成的神所掌控.如此,一個天國子民成為受苦的人.他必須受許多的苦,被維持世界權力結構的種種「長老,祭司長,和文士」棄絕.但天國之子將以愛還報暴力,以寬恕還報攻擊.在暴力下,他會被殺.但是,他會復活.

復活,不是「翻生」,而是生命與神結連.生命分沾了神的神聖,具備神聖的永恆和豐盛.復活的生命,永恆地發出人性和神聖的光芒.

彼得雖然承認耶穌為「基督」,但當牽涉到他的實際利益時,他即反面.馬可福音說:「彼得就拉著他,勸他」.「勸他」是翻譯得很客氣的用詞,《新漢語譯本》更準確地翻譯作「責備他」.彼得為了不想聽耶穌的看法,而責備耶穌.看來,彼得不願意與醜惡的世俗抗衡.看來,他寧願同流合污.看來,他渴望維持自己的既得利益.看來,他受「塵世國度」的國民教育太深.他沒有塵世以外的任何東西要追求.

彼得在生命的張力中.他跟從耶穌,卻不想跟從耶穌的生命方式.正如很多人移民外國,宣誓效忠那國家,卻從來沒有真正效忠過.彼得的身分認同,是跟從基督；但他的生命方式,卻是世俗的方式.

現在,輪到耶穌責備彼得.「責備」這個字與剛才那「勸」字,其實是同一個字.耶穌和彼得都立場鮮明.耶穌一定要彼得明白,不同的生命方式,有不同的命途.正如在分叉路口,有人問路,你必須清楚講明,這條路往哪裡去,那條路往哪裡去,一點不能含糊.因為走在不同路上,有不同結果.

耶穌責備彼得說:「撒旦,退我後邊去吧!因為你不體貼神的意思,只體貼人的意思.」

歸屬神的人,就是天國的子民.信耶穌的人,是進入一新的國度裡,他以神的心去思考,而不是按塵世的想法去思考.體貼人的意思,就是只想到現世的權力,利益,慾望的滿足.體貼神的意思,就是想到神對人類的渴望,想到愛,想到人性的尊貴,想到人性的未來.

(3)天國的國民教育的具體實踐

體貼神的意思,具體而言,是怎麼樣的呢?如何實踐呢?

馬可福音 8:34-38

34     於是叫眾人和門徒來,對他們說,若有人要跟從我,就當捨己,背起他的十字架來跟從我.     35     因為凡要救自己生命的,〔生命或作靈魂下同〕必喪掉生命．凡為我和福音喪掉生命的,必救了生命.     36     人就是賺得全世界,賠上自己的生命,有甚麼益處呢.     37     人還能拿什麼換生命呢.     38     凡在這淫亂罪惡的世代,把我和我的道當作可恥的,人子在他父的榮耀裡,同聖天使降臨的時候,也要把那人當作可恥的.

這裡有一個重要的

主題:「捨己」.

「捨己」之可能,必須先看到他人,並且為了維護他人的人性尊嚴,甘於放下自己的權利.

現在,在餐廳吃飯時,多了如下的場面.

侍應遞飲品時,一時大意,濺濕了客人.儘管那侍應即時道歉,作出種種合宜的補救措施.客人卻破口大罵,數落那侍應,又要召部長來,要部長對那侍應嚴重警戒一番.這客人有沒有得到合宜服務的權利,當然有.但是,當堅持這權利時,是否過份地羞辱了別人.「捨己」,是為了維護他人的人性尊嚴,甘於放下自己的權利.這是一種維護人性尊嚴的教育.

某些型態的國民教育,會造就狹隘的民族主義.過份的民族主義,會造成過份堅持自己的利益和權利,而造成排外.長此下去,對這群體毫無好處,也對這群體與別的群體的合諧共存毫無益處.好的教育,不應只顧及自己的存活,也應顧及地球上所有共同體的存活.合宜的民族主義是需要的,但更重要的是,要有超越民族主義的世界公民意識.甘於放下自己權利,而肯定他人尊嚴的「捨己」精神,從長遠而言,更能造福世人.

納粹般的獨裁政權也要求國民捨己,這當然不是天國子民的捨己對象.主耶穌說:「凡為我和福音喪掉生命的,必救了生命.」捨己的理由,必需出自基督和福音.基督和福音,總是對抗塵世的偶像和種種不義.為了讓他人得釋放,尊嚴被肯定,這才成為捨己的理由.唯有這樣,我們才能救自己的靈魂.所謂靈魂,即我們的內在生命.

(4)結語

讓我們再次肯定,我們的生命是歸屬於神的.

既然歸屬於神,便努力活出神的特質.

這種特質,是敢於為了建立別人的尊嚴和價值,獻出自己.

但願:榮耀歸於聖父,聖子,聖靈.阿們.



\newpage

\section{智者必懷仁}
\label{sec:168}
\textbf{林豪恩先生 - 智者必懷仁}
\newline
\newline
\hyperref[sec:167]{< < < PREV SERMON < < <}
~
\hyperref[sec:toc]{[返目錄]}
~
\hyperref[sec:169]{> > > NEXT SERMON > > >}


雅各書 3:13-4:3
\newline
\begin{longtable}{cl}
\hline
\hline
章節 & 經文 (和合本修訂版)\\
\hline
3:13 & \begin{tabularx}{0.7\textwidth}{X} 你們中間誰是有智慧有見識的呢？他就當在智慧的溫柔上顯出他的善行來。 \end{tabularx} \\ \\ \relax
3:14 & \begin{tabularx}{0.7\textwidth}{X} 你們心裡若懷著惡毒的嫉妒和自私，就不可自誇，不可說謊話抵擋真理。 \end{tabularx} \\ \\ \relax
3:15 & \begin{tabularx}{0.7\textwidth}{X} 這樣的智慧不是從上頭下來的，而是屬地上的，屬情慾的，屬鬼魔的。 \end{tabularx} \\ \\ \relax
3:16 & \begin{tabularx}{0.7\textwidth}{X} 在何處有嫉妒、自私，在何處就有動亂和各樣的壞事。 \end{tabularx} \\ \\ \relax
3:17 & \begin{tabularx}{0.7\textwidth}{X} 惟獨從上頭來的智慧，先是清潔，後是和平、溫良、柔順，滿有憐憫和美善的果子，沒有偏私，沒有虛偽。 \end{tabularx} \\ \\ \relax
3:18 & \begin{tabularx}{0.7\textwidth}{X} 正義的果實是為促進和平的人用和平栽種出來的。 \end{tabularx} \\ \\ \relax
4:1 & \begin{tabularx}{0.7\textwidth}{X} 你們中間的衝突是哪裡來的？爭執是哪裡來的？不是從你們肢體中交戰著的私慾來的嗎？ \end{tabularx} \\ \\ \relax
4:2 & \begin{tabularx}{0.7\textwidth}{X} 你們貪戀，得不著就殺人；你們嫉妒，不能得手就起爭執和衝突；你們得不著，是因為你們不求。 \end{tabularx} \\ \\ \relax
4:3 & \begin{tabularx}{0.7\textwidth}{X} 你們求也得不著，是因為你們妄求，為了要浪費在你們的宴樂中。 \end{tabularx} \\ \\
[1ex]
\hline
\hline
\end{longtable}
日期:2012年10月21日

主題:智者必懷仁

證道:林豪恩先生

經文:雅各書3:13-4:3

一．經文

13     你們中間誰是有智慧有見識的呢．他就當在智慧的溫柔上,顯出他的善行來.     14     你們心裡若懷著苦毒的嫉妒和分爭,就不可自誇,也不可說謊話抵擋真道.     15     這樣的智慧,不是從上頭來的,乃是屬地的,屬情慾的,屬鬼魔的.     16     在何處有嫉妒分爭,就在何處有擾亂,和各樣的壞事.     17     惟獨從上頭來的智慧,先是清潔,後是和平,溫良柔順,滿有憐憫,多結善果,沒有偏見,沒有假冒.     18     並且使人和平的,是用和平所栽種的義果.

二．釋經及應用

1．智者不是外加條件,而是有生活表現

3:13     你們中間誰是有智慧有見識的呢．他就當在智慧的溫柔上,顯出他的善行來.

雅各問:「你們中間誰是有智慧有見識的呢?」誰是智者呢?有人可能會想到住在深山寺院的高僧,或者隱居於郊外修院的修士.不過,這些遠離塵世的高僧和修士,都已經成為了故事中的人物,在書本中繼續流傳.在現代社會中,哪裡找到智者呢?智者在哪裡呢?遠在天邊,近在眼前,智者就藏在大學裡面,即是我們的禮拜堂所在之處.為甚麼這樣說呢?因為智慧這個字的希臘文Sophia,就藏在Philosophy這個字裡面.在大學校園,碰口碰面都是Master of Philosophy,或者Doctor of Philosophy,同時,大學每年也頒發不少Master of Philosophy,以及Doctor of Philosophy的證書.能夠Master(掌握)智慧的人,應該算是智者了吧!稱得上智慧的Doctor(老師),更加不可能不是智者.

有人認為,除了那些擁有Master of Philosophy和Doctor of Philosophy學位的人之外,社會中應該還有其他人都是有智慧有見識的.他們雖然沒有在學術上繼續鑽研,但卻具有專業資格,例如:律師,醫生,會計師．．．等等.又有人認為,有些人雖然沒有高等學歷,也沒有取得專業資格,但他們白手興家,建功立業,都應該算是有智慧有見識之列.再有人認為,讀萬卷書不如行萬里路,真正有智慧有見識的,應該是那些周遊列國,見聞識廣的人.

總括而言,現代人心目中有智慧有見識的,就是擁有較高學位的人,或者是擁有專業資格的人,或者是白手興家事業有成的人,又或者是曾經週遊列國的人.高等學位,專業資格,事業成就和周遊列國的經驗,都是外加在一個人身上的條件,這些外加的條件,都必須首先具有某些其他條件才能夠得到.在這樣的標準之下,沒有這些條件的人,就難以擠身於有智慧有見識之列了.

雅各回答「誰是有智慧有見識的呢?」這個問題的時候,提供了另外一個向度給我們參考.他不以學位,專業資格,事業成就,或者遊歷經驗這些外加的條件去判斷誰是有智慧有見識,而是以一個人處事為人的表現作為參考.他認為有智慧和見識的,「就當在智慧的溫柔上顯出他的善行來」.這句話直譯是「藉著美善的生活,在智慧的謙虛裡顯出他的行為.」在雅各心目中,有智慧有見識的,可能是有高等學位的人,可能是有專業資格的人,可能是事業有成的人,可能是曾經周遊列國的人,也可能是任何一個人,只要他能夠從這個人身上看見智慧的表現,就是謙虛地,美善地生活著.雅各的意思是,智者不需要宣傳自己的智商(IQ)有多高,不需要展示自己的學歷有多高,不需要炫耀自己的事業有多龐大,也不需要暗示自己踏足過世界上多少個國家,因為智者的智慧和見識自然會在日常生活和待人接物之中謙虛地流露出來.

當一個人謙虛地把他的智慧和見識透過生活行為表現出來的時候,會是怎樣的呢?又不會是怎樣的呢?接著,雅各從正反兩方面去說明.

2．不是從上頭來的智慧

3:14     你們心裡若懷著苦毒的嫉妒和分爭,就不可自誇,也不可說謊話抵擋真道. 3:15     這樣的智慧,不是從上頭來的,乃是屬地的,屬情慾的,屬鬼魔的.     3:16     在何處有嫉妒分爭,就在何處有擾亂,和各樣的壞事.

以疾妒來描述人與人之間的關係,通常都是負面的；這裡「紛爭」一字有自私,卑鄙和爭吵的意思,在亞里士多德的著作中曾用這字描述為自己的政治席位而不擇手段的爭取.雅各再以「苦毒」這個負面的字去形容嫉妒和紛爭,不是「負負得正」,而是更加負面.若一個人的心思意念被苦毒的嫉妒和紛爭所充滿著,他的智慧所表現出的行為會是怎樣的呢?會帶來甚麼結果呢?

一般來說,大家和平共處的時候,未必會被苦毒的嫉妒和紛爭佔據心靈.在甚麼情況之下,苦毒的嫉妒和紛爭最容易滋生和繁衍呢?就是在爭奪的時候.例如有兩位武士同時喜歡一位淑女,為了奪得美人歸,在戲劇中最常見的劇情發展就是決一死戰.又例如首長寶座只能一人登上,卻有幾個人去爭奪,於是掛在口邊的「君子之爭」有人認為變成了如動物農莊般狼豬互咬,又有人把這個過程描述為一場泥漿摔角,雙方都醜態百出.

我們坐在電視機面前,「食住花生等睇戲」,看著報紙高談闊論,當然可以大義凜然地指指點點,而且很容易令自己產生以為自己清高得出污泥而不染的幻象.其實,我們沒有泥足深陷,可能只是因為我們尚未披甲出戰,當我們進入戰場後,或許也清高不了多少.還記得當你看見朝思暮想的女同學竟然被你的情敵拖著手的時候,你的心裡有多嫉妒和苦毒嗎?還記得上次你申請升職,公司升了另一位同事沒有升你,你的心裡有多嫉妒和苦毒嗎?最令人義憤填膺的是,那位你朝思暮想的女同學竟然選擇那位沒有你那麼英俊瀟灑的情敵,於是你猜想你的情敵一定「出古惑」；最令你憤憤不平的是,公司升職竟然選擇那位沒有你表現得那麼好的同事,於是你猜想那位同事一定「有背景」或者「耍手段」.當一個人認為對手不擇手段的時候,這個人心裡苦毒的嫉妒和紛爭就有空間發揮,更有機會佔據這個人的心思意念,於是,既然對手已經做了「初一」,我做「十五」也不算太過份吧.結果,這個人把自己變成了他想像的對手一樣.

最近南非前大主教圖圖(Desmond Tutu)於英國《觀察家報》撰文,指摘美國前總統小布殊及前英揆貝理雅謊稱伊拉克擁有大殺傷力武器,於2003年入侵伊拉克,造成大量人命傷亡,更令全球陷入前所未有的亂局,認為兩人應為此接受海牙國際刑事法庭審訊.他指出:「問題不在薩達姆是善是惡,或他屠殺了多少人民,而是布殊先生和貝理雅先生不應讓自己墮落至如薩達姆那樣不道德.」

爭奪,不單單在情場上,商場上或戰場上出現,在最意想不到的場景也可以出現.雅各這些說話,就是對著信徒群體中想做領袖的人說的.這段說話的上文是3章1節,雅各如此說:「我的弟兄們,不要多人作師傅,因為曉得我們要受更重的判斷.」雅各指出,如果為了爭奪而讓苦毒的嫉妒和紛爭充滿自己的心,雖然仍然口若懸河,說得頭頭是道,不過,口裡所出的話,很可能只是反映說話者的動機和偏見,而未必是事實,也很可能話說間貶低別人,吹噓自己,而還以為自己在伸張正義.在信徒群體中作領袖,本來是宣講真理,然而,若內心被苦毒的嫉妒和紛爭自私所佔據,領袖本身就成為了謊言和虛假,與真理背道而馳,結果只會引致自己以及群體的紛亂和敗壞.

雅各指出,這種智慧不是從上面來的,而是屬於地下的,血氣的,鬼魔的.那麼,從上而來的智慧又是怎樣的呢?我們如何得知這種智慧是從上面來的呢?

3．是由上頭來的智慧

3:17     惟獨從上頭來的智慧,先是清潔,後是和平,溫良柔順,滿有憐憫,多結善果,沒有偏見,沒有假冒. 3:18     並且使人和平的,是用和平所栽種的義果.

保羅在加拉太書提出聖靈所結的果子有九方面的特質,「就是仁愛,喜樂,和平,忍耐,恩慈,良善,信實,溫柔,節制」(加拉太書5:22-23).雅各在這裡指出從上頭而來的智慧也有七方面的特質,包括:清潔,和平,溫良,柔順,滿有憐憫美善的果子,沒有偏見,沒有假冒.這些特質,就是有智慧有見識的人透過其生活行為和待人處事流露出來而被人看見.這七種特質的智慧,會指導以和平的方法,追求和平,最終會帶來和平.

雅各認為,從上而來的智慧七方面的特質排第一的是清潔.清潔的意思是心思和動機的純潔,這字在雅各書一章二十七節與「沒有玷污」同列,兩個詞互相補充.在運用智慧的過程中,心思和動機有多純潔,或者有多混雜,與和平有很大的關係.二十多年前,一百多萬香港人上街遊行,和平有序,沒有打鬥,沒有破壞,因為一百多萬人都把個人的利益和私慾暫時放下,大家的關心和著急都是一樣的.跟近,在政府總部外的廣場也有近十萬人的集會,就連中學生也有類似的發現.學民思潮的黃莉莉有這樣的分享:「在佔領開始時宣讀的絕食宣言中,我們曾表明有兩個目的,其一是希望大家放下慾望,和我們一同關注這件事,其二是要求政府撤回國民教育科.現在我們已經達到了第一個目的,12萬的市民放下自己生活,犧牲晚上的時間,來到政總和我們一同佔領.我記得12萬市民到來那晚,我們的音響設備十分不完善,在場外的市民根本無法參與我們的集會,但大家並沒有出現混亂,反而配合糾察的指示,安靜地坐下在地上.」再看看近來在香港附近的城市舉行的反日示威,那些破壞和混亂的場面令自己人感到非常尷尬,製造了大量被對手混淆視聽和指責的機會.有人認為,在示威隊伍中,很可能有些人懷著複雜和不可告人的動機而來,而很可能也有不少人除了反日之外,還乘機發洩各種不滿.如此混雜不純的心思和動機,所帶來的難免就是破壞和擾亂.

另外六方面的特質可以分為三組,第一組是「和平」,「溫良」,「柔順」,這三個字的原文都是以用一個字母開始的.原文溫良一字有溫柔和寬宏的意思,解作寬宏時,見用於指有權柄者的寬宏大量,用作溫柔時,有不以暴力去解決問題的指向.至於「柔順」,不是形容護髮素,而是有柔和,慷慨的意思,指在可能的情況下遷就他人,願意為整個群體而作出讓步.以上三種特質都是建立群體的人所要有的德性,和平是能與人和睦共處,寬宏是有容人的雅量,柔和代表願意遷就體諒,將群體的關注置於自己之先.

第二組是「滿有恩慈和善果」指慈悲為懷,樂善好施,與人為善.懷有這特質的智慧會產生有這樣美好果效的行動.

第三組是「沒有偏見和沒有假冒」.偏見,其中一個例子就是二章開始時所說的,按外貌待人,「有一個人帶著金戒指,穿著華美衣服,進你們的會堂去；又有一個窮人穿著骯髒的衣服也進去,你們就重看那穿華美衣服的人,說「請坐在這好位上」,又對那窮人說:「你站在那裡」,或「坐在我腳凳下」.虛偽就是裝腔作勢,假冒為善.雅各指出,從上而來的智慧沒有偏見和產偽這些成分.這些不只是個人所有的德性,也是一個和平的群體所必需有的特質.

三．總結

蒲松齡的《聊齋志異》中的有一名句這樣說:「智者不必仁,而仁者則必智」.智者是否不必仁呢?唐太宗賜給臣子蕭瑀的詩其中一句卻話:「智者必懷仁」.看來,雅各傾向認同後者.雅各指出,在群體中作為一個師傅,當然是一位有智慧有見識的人,不過,這智慧和見識不只是智識和能力,而同時也是品德和表現.

每逢選舉,無論是政府首長選舉或者是各級議會選舉,都有不少人走出來,聲稱自己有智慧有見識,有能力有抱負為大家解決困難,為社會帶來出路.他們應該參考雅各在這裡討論智慧所提出的,去除自己心裡苦毒的疾妒和紛爭,以免為他們領導的群體帶來更多的擾亂和壞事.同時,也需要培育從上而來的智慧的七種特質,時刻提醒自己,心中懷仁,眼中有人.若果能夠在發動任何運動的時候,在推出任何政策的時候,心中懷仁,眼中仍有人,那就可以與和平越來越接近了.

閱讀那位中學生投稿報章解釋撒離政府總部的文章,欣賞及敬佩之情油然而生.文章其中一段這樣寫:「韓老師,一位已經63歲的老師,他令我慚愧了,我對他深感敬佩亦感謝他.171小時的絕食,絕非輕易事情,但韓老師創造了歷史,他沒有說過辛苦,他尊重我們的決定,他堅持絕食到最後一刻.然而,我們並不願意看他再絕食下去,怎麼政府的不當要一位63歲的老人家來捱苦呢?心,很痛.所以必須完結佔領,完結絕食,讓老師好好休息!」無論我們對事件有甚麼立場,多多少少都認同這群中學生的智慧.他們在發動大規模行動的時候,在抗爭的時候,沒有被運動充昏頭腦,他們的眼裡仍然看見那位長者,他們的心仍然懂得為那位長者所受的苦而感到痛,他們仍能能夠拒絕以別人作為運動的代價,他們仍然能夠看見別人對他們的關心,仍然能夠以感恩的心去看待別人的關心.

猶太先賢說:「若人的善行比他的智慧多,他的智慧便持久；若他的智慧比他的善行多,他的智慧便不能持久.」(《先賢集3.10》)願我們社會中的智者,都懷著一棵仁心.



\newpage

\section{家的盼望}
\label{sec:169}
\textbf{賀志勇博士 - 家的盼望}
\newline
\newline
\hyperref[sec:168]{< < < PREV SERMON < < <}
~
\hyperref[sec:toc]{[返目錄]}
~
\hyperref[sec:170]{> > > NEXT SERMON > > >}


雅各書 5:13-18
\newline
\begin{longtable}{cl}
\hline
\hline
章節 & 經文 (和合本修訂版)\\
\hline
5:13 & \begin{tabularx}{0.7\textwidth}{X} 你們中間若有人受苦，他該禱告；有人喜樂，他該歌頌。 \end{tabularx} \\ \\ \relax
5:14 & \begin{tabularx}{0.7\textwidth}{X} 你們中間若有人病了，他該請教會的長老們來為他禱告，奉主的名為他抹油。 \end{tabularx} \\ \\ \relax
5:15 & \begin{tabularx}{0.7\textwidth}{X} 出於信心的祈禱必能救那病人，主必叫他起來；他若犯了罪，也必蒙赦免。 \end{tabularx} \\ \\ \relax
5:16 & \begin{tabularx}{0.7\textwidth}{X} 所以，你們要彼此認罪，互相代求，使你們得醫治。義人祈禱所發的力量是大有功效的。 \end{tabularx} \\ \\ \relax
5:17 & \begin{tabularx}{0.7\textwidth}{X} 以利亞與我們是同樣性情的人，他懇切地祈求不要下雨，地上就三年六個月沒有下雨。 \end{tabularx} \\ \\ \relax
5:18 & \begin{tabularx}{0.7\textwidth}{X} 他又禱告，天就降下雨來，地就有了出產。 \end{tabularx} \\ \\
[1ex]
\hline
\hline
\end{longtable}
日期:2012年10月28日

主題:家的盼望

證道:賀志勇博士

經文:雅各書 5:13-18

13     你們中間有受苦的呢,他就該禱告.有喜樂的呢,他就該歌頌.     14     你們中間有病了的呢,他就該請教會的長老來．他們可以奉主的名用油抹他,為他禱告．     15     出於信心的祈禱,要救那病人,主必叫他起來．他若犯了罪,也必蒙赦免.     16     所以你們要彼此認罪,互相代求,使你們可以得醫治.義人祈禱所發的力量,是大有功效的.     17     以利亞與我們是一樣性情的人,他懇切禱告,求不要下雨,雨就三年零六個月不下在地上.     18     他又禱告,天就降下雨來,地也生出土產.

一,Introduction 家的重要性

牧師曾經送過我一個大電視,大約有一尺這麼厚.品質很好,不過我擔心壞了之後就沒有地方可以換零件了.有一次,一個朋友來我家看到那個大傢夥,就跟我講:"你怎麼不換一個超薄的液晶電視掛牆上呢?"我跟他講我"戀舊".

現在一般的電器都講究"要新","要薄",要可以經常換.如果有什麼電器廣告講,我們的電器最大的特點就是"長壽",怎麼用都不壞,我想它一定不會好賣.你看,現在Iphone 5出來了吧,馬上大陸就有人出8000多港幣買新機.

家用電器可以經常換,人可以經常換麼?家可以經常換麼?據調查,2010年,北京市的離婚率高達39\%,上海38\%,香港也高達33.8\%,就是說香港每三對夫妻,就有一對離婚.所以有個香港學者講,今天中國最大的問題,不是政治的問題,不是經濟的問題,是家庭的問題.

家不僅僅是吃晚飯的地方,也不僅僅是睡覺的地方,它是讓我們的身體和心靈都可以得到休息的地方.按我媽的話講就是"家就是個加油站",如果在家裡都不能休息,你還有哪里可以休息呢?家太重要了.

二,解經

除了我們的小家,雅各書5:13-18就講到我們另外一個家--教會.教會是我們屬靈的家,在這裡我們尋求靈命的安息.所以一家教會的特點,決定了我們生命的狀況如何.在雅各的經文中,一家教會必須具有這三樣特點:

(1)關顧、醫治的家

教會必須是一個關顧、醫治的家.雅各在13節講"你們中間有受苦的,就該禱告", 14節講"你們中間有人患病嗎?他應該請教會的長老來,為他禱告,奉主的名給他抹油."  教會中有不少苦痛,    死亡、失去、分離、失敗,在教會中也不例外,連保羅都說"有誰軟弱,我不軟弱呢?".所以有時候我跟著伍牧師出去探訪病人,看著病人的樣子,我自己心裡都難過.可是教會絕不僅僅只有苦、只有痛,教會還有關顧和治療.雅各書告訴我們,面對這些苦痛,我們第一要禱告,第二要關顧,第三要用藥,合起來就是一個完整的治療.

以前猶太人沒有專門的醫生,祭司或者拉比還要負責看病；早期教會也是這樣子,長老不但會探望病人,替病人禱告,還負責治病.治病用什麼藥呢?用的就是橄欖油.現在有的基督徒鼓吹用禱告來治病,反對尋醫吃藥,其根據就是雅各的這句經文.他們以為這個抹油的抹字,是舊約裡頭"膏抹"的意思,是一種屬靈的宗教儀式.但實際上,膏抹的抹字,希臘文是"Chrio", 基督"Christ" 就是從這個字引申過來的.但雅各書裡抹油的抹字用的是"aleipho",不是"Chrio". 這個字在馬可福音6:13節用過,也是指一種治療方式.還有,大家記不記得,好撒瑪利亞人也是用油來治傷的.

所以說這些人誤解了"抹油"了這個字,他們以為抹油是個宗教儀式,因此他們片面強調單單靠禱告來治病療傷,好像禱告就是萬靈神丹一樣.他們以為依靠一點宗教儀式,或者說幾句屬靈話語,就能立即讓人脫胎換骨,藥到病除.但這樣的禱告,其實是想操控神,要神按我們的心意來抓藥!如果我們單單這樣來醫病的話,那教會就不是一個關顧、治療的家,教會可能會帶來更多傷害--這種傷害比在教會之外的傷害還可怕.

我有個朋友最近特別怕去教會,因為她一開口說自己的難處,還沒講幾句,就有好幾個人給她開藥方了:唉!你要多禱告!你要多倚靠神.  所以她一肚子苦水沒倒出來,還要向大家表示感謝,然後給大家一個標準答案:"是是是,我倚靠神不夠,我要多禱告,多倚靠神." 然後大家一聽她孺子可教,都欣然點頭,接著一起感謝神幫助了這位姊妹.

這個姊妹告訴我說,沒有人關心她到底發生了什麼事；沒有人給她抹油,大家只想給她開屬靈藥方,一些看似非常符合聖經、沒有信徒敢對它說不的藥方.但這樣的藥方她早知道,何必要別人再來教訓她呢?

城市大學的老師斯諾寫書說,他最怕的就是這種"屬靈官腔".你一分享難處,馬上得來的是"你要多禱告；你要依靠耶穌,不要依靠人,""你要看耶穌的面,不要看人的面."這種話聽起來冠冕堂皇,但聽了這話,你會不會跟充了電一樣,立刻就剛強壯膽、一無掛慮呢?

我們看看新約裡,耶穌無論是治療麻風病人,還是癱子；無論是面對撒瑪利亞婦人,還是面對行淫被抓住的婦人,他並沒有說:"你要多禱告,你要倚靠神,不要倚靠人." 他很少打屬靈官腔,他都是直接面對這些人的生命傷痛去勸勉、安慰、治療,所以雅各在這裡也提醒我們,別人受苦了,不但要勸人多禱告,還需要傾聽,還需要長老去抹油.長老是整個教會的代表,他去給病人抹油,就代表著整個教會是同一個肢體,一起關顧、安慰、治療.

雅各告訴我們,教會乃是醫治的家.我願我們這個家少一點"官腔",多一點禱告,也多一點抹油!

(2)歌頌的家

教會除了治療,教會還必須是一個歌頌神的家!13節雅各就說:"你們中間有喜樂的,就應當歌頌."歌頌是什麼意思?歌頌包括感恩、稱謝和敬拜.猶太人歌頌很多時候都是會唱出來的,你看詩篇就知道了.唱歌是一件很美妙的事,有益身心,哪怕唱卡拉OK都不錯啊.不過上次伍中恩教授就講,中國人裡頭,就屬我們漢民族最少唱歌了.幸好我們基督徒還是喜歡唱歌的,所以歡迎大家都來詩班唱一唱,我們聖誕節獻詩還需要大家支援了.

保羅曾經講過這樣一句發人深省的話:"我曉得怎樣處富貴,也曉得怎樣處貧賤."保羅和西拉在腓立比的監獄裡都唱歌讚美神.但很多時候,我們只曉得在傷痛時、有麻煩的時候禱告,但當富足和喜樂時,就把主忘得一乾二淨了.如果我們花很多時間去向主求這個求那個,但感恩歌頌卻只有寥寥幾句,我們這個屬靈的家會怎樣呢?

不過,歌頌神也要分時機場合,你身邊的人剛剛失業,請你千萬不要大聲歌頌神讓你買了一輛新車；你身邊的人剛剛喪失親人,請你千萬不要公開感謝神讓你"凡事充足,一無掛慮."你一無掛慮了,人家喪失親人,當作何感受呢?所以餘創豪就講:"這樣的歌頌感恩只適宜自己用,不適宜饋贈親友."

雅各在"你們當歌頌"之前加了這一句"你們中間有人受苦麼?",他不僅要人歌頌,他還要人觀察到別人的難處、軟弱去加以體諒.當你"以恩慈相待,存憐憫的心"的時候,你才能做到真正的感恩歌頌,教會也才真正成為了感恩歌頌的家.

(3)彼此認罪、代求的家

除此之外,16節講,教會還必須成為一個彼此認罪、彼此代求的家.沒有認罪,也就沒有真正的關顧治療,沒有真正的歌頌.

或許你要問:"我自己跟上帝認罪不就夠了麼?何必還要彼此認罪呢?"潘霍華是這樣解釋的:"神預備了弟兄來幫助我,他聽我認罪,就像基督聽我認罪一樣,他奉耶穌的名赦免了我的罪."不錯,在密室裡向神認罪,或許可以讓我不這麼尷尬,但這種認罪也不徹底.唯有向弟兄認罪,他又能帶領我到主的施恩寶座前,這樣就徹底了.

不過,彼此認罪容易不容易啊?一般來講,我們中國人很怕當面對質,因為當面對質會傷和氣.所以我爸媽經常跟我講:說話要婉轉一點,不要得罪人.婉轉是什麼意思啊?婉轉就是說話要轉幾個彎,你講了,對方有沒有收到,就看對方的領悟力了.所以你問中國人,這件事怎麼樣啊,我做得好不好啊? 一般都是回答:嗯,還好；嗯,還可以.

但這樣的話並不造就人.要彼此認罪,首先需要一個肯直接跟你講實話的弟兄或者姊妹.有的東西你沒有注意到,他幫你指出來了,你又肯認罪,那就是彼此得了一個弟兄了.箴言 27: 5說:"當面的責備,強過背地的愛情." 所以大家好好想想,如果你身邊沒有一個人敢同你提意見,那你就危險了；因為你沒有真正的朋友,沒有真正的弟兄姊妹.一間教會也如此,教會裡如果缺乏彼此認罪的弟兄姊妹,教會就只會維持表面的一團和氣,底下卻是冷漠疏離.當然,當面的責備不是事事當著所有人的面責備你,而是可以當著你的面,用愛心說誠實話.這一點也非常重要!

弟兄姊妹,為了我們這個家,為了我們自己身體與靈命的安息,讓我們彼此認罪,也彼此代求.但請記住,彼此認罪不是我跟你熟、關係好,我們兩個來認罪、代求一下,不熟的、關係不好的人就放在一邊.彼此認罪首先是向神認罪,其次是向教會的牧師、長老、執事認罪--他們是神的教會的代表,最後是向你得罪過的弟兄姊妹認罪,這樣我們才可以在神裡面彼此和好了.

結論:禱告的家

雅各這幾節經文,其實句句都跟禱告有關係.這讓我想起了一段往事.幾年前我去了溫州一家農村教會做探訪,那邊人都非常敬重傳道人.聽說來了個香港的傳道人,就來了好些老姊妹.她們有的是身體不好,有的是家裡關係處理不好,有的是孩子讀書有問題,大家都請我為她們禱告.她們說香港的傳道人是"義人", 比本地的傳道人要靈驗,禱告主一定會聽.我聽了哭笑不得,於是跟她們講"禱告的原義是"把我們的所想說給神聽",並沒有什麼魔力；而且彼得前書講"我們都是被揀選的族類,都是君尊的祭司."所以有心靈誠實,哪個基督徒禱告神都會聽的." 聽了我一番話,老姊妹們都點頭說:"賀傳道講得有道理!不過你這麼遠來,就不要浪費了,還是請你幫我們禱告一下吧!" 我還要解釋,旁邊的朋友朝我使了個眼神,我只好慢慢替她們一一禱告.

雅各告訴我們,教會是一個治療、歌頌、彼此認罪代求的家,是一個禱告的家.這一切都是因為我們有這樣的信心:我知道我的救贖主活著,末了必站立在地上.但這一信心不是出於教條,也不是出於"官腔",而是出於我們能真正理解並行出神的話語.就好像雅各書1:25節所說:"這人...乃是實在行出來,就在他所行的事上必然得福."

I have a dream,我盼望有一天我們都能行出神的教導,我願我們的教會有一日變得完全,變得像雅各所說的那樣,一個家,一個有情有義有真理的家.



\newpage

\section{不要那樣問而要}
\label{sec:170}
\textbf{鄧瑞強博士 - 不要那樣問,而要這樣想}
\newline
\newline
\hyperref[sec:169]{< < < PREV SERMON < < <}
~
\hyperref[sec:toc]{[返目錄]}
~
\hyperref[sec:171]{> > > NEXT SERMON > > >}


馬可福音 10:2-9
\newline
\begin{longtable}{cl}
\hline
\hline
章節 & 經文 (和合本修訂版)\\
\hline
10:2 & \begin{tabularx}{0.7\textwidth}{X} 有法利賽人來問他說：「男人休妻合不合法？」意思是要試探他。 \end{tabularx} \\ \\ \relax
10:3 & \begin{tabularx}{0.7\textwidth}{X} 耶穌回答他們說：「摩西吩咐你們的是甚麼？」 \end{tabularx} \\ \\ \relax
10:4 & \begin{tabularx}{0.7\textwidth}{X} 他們說：「摩西准許寫了休書就可以休妻。」 \end{tabularx} \\ \\ \relax
10:5 & \begin{tabularx}{0.7\textwidth}{X} 耶穌對他們說：「摩西因為你們的心硬，所以寫這誡命給你們。 \end{tabularx} \\ \\ \relax
10:6 & \begin{tabularx}{0.7\textwidth}{X} 但從起初創造的時候，神造人是造男造女。 \end{tabularx} \\ \\ \relax
10:7 & \begin{tabularx}{0.7\textwidth}{X} 因此，人要離開他的父母，與妻子結合， \end{tabularx} \\ \\ \relax
10:8 & \begin{tabularx}{0.7\textwidth}{X} 二人成為一體。既然如此，夫妻不再是兩個人，而是一體的了。 \end{tabularx} \\ \\ \relax
10:9 & \begin{tabularx}{0.7\textwidth}{X} 所以，神配合的，人不可分開。」 \end{tabularx} \\ \\
[1ex]
\hline
\hline
\end{longtable}
日期:2012年11月4日

主題:不要那樣問,而要這樣想

證道:鄧瑞強博士

經文:馬可福音 10:2-9

各位弟兄姊妹,主內平安.

在人類社會裡面,結婚和離婚常常是熱門的討論話題.有人說,婚姻是人類社會裡最奇怪的現象,未結婚的人和結了婚的人都同樣不開心,但他們不開心的理由卻剛好相反.未結婚的人,他們不開心,是因為他們未結婚.結了婚的人,他們不開心,正正是因為他們結了婚.

聽過一個故事是這樣的:

一對夫婦,正在度過結婚三十周年紀念,丈夫忽然悲從中來,異常傷心.妻子問他:「和我度過了三十年,你覺得很感動,是嗎?」丈夫說:「不是.我只是想起一件往事.」妻子問:「什麼往事?」丈夫說:「你記得嗎,三十年前,我欠了你父親一大筆錢,無法償還.你父親跟我說:『我的女兒無人要.你一是跟我的女兒結婚,一是準備坐三十年監.』我現在回想,若我當時選擇正確的話,到今日,我應該已經得到自由了.」

不知你在結婚三十周年紀念時,會否悲從中來?

有一個年青人,活得隨心所欲,去了很多想去的地方,做了很多想做的事,享受了很多想享受的生活樂趣.然後,他結婚了.結婚後,他滿懷感慨地對他的朋友說:「結婚,終於讓我明白到什麼是人生的真正幸福,但已經是太遲了.」

婚姻一般都不是一帆風順的.聽說,只有亞當和夏娃的婚姻才是完美的.因為,當亞當回家時,不用聽夏娃說別的男士對她如何好；而夏娃也聽不到亞當說他的母親燒菜如何棒.

若果沒有婚姻,人類便不會醒覺自己有「原罪」.若果沒有婚姻,人類還以為自己的愛的能力有多強大.若果沒有婚姻,人類很容易誤會自己已經是聖人.多虧婚姻的存在,我們才能遠離虛幻,面對現實,面對生命的軟弱,面對醜陋的人性.

遇著困難,不是人人都努力面對的,很多人寧願放棄.這邊廂有人努力結婚,那邊廂有人努力離婚.去愛一個人,已經不容易.但也有人花更大的氣力,去捨棄那曾經相愛的人.今日講道的內容,正與「離婚」有關.

今日的講道經文是:馬可福音 10:2-9

10:2     有法利賽人來問他說,人休妻可以不可以,意思要試探他.

法利賽人問耶穌:「休妻是否合乎律法規定?按信仰精神,可否離婚?」

在耶穌時代,猶太人一般接受休妻為合乎律法的行為.當時,有幾個重要的拉比.Rabbi Shammai 認為,若女方在性方面有不當的行為,丈夫便可以休棄她.Rabbi Hillel 比較寬鬆,他認為只要女方弄糟一碟菜,丈夫便可以休妻.Rabbi Akiba 最極端,他認為只要丈夫認為別的女人更美,更吸引,丈夫便可以休妻.

法利賽人問耶穌:「人休妻可以不可以?」

耶穌如何回答?

馬可福音 10:3-7

3     耶穌回答說,摩西吩咐你們的是甚麼.     4     他們說,摩西許人寫了休書便可以休妻.     5     耶穌說,摩西因為你們的心硬,所以寫這條例給你們.     6     但從起初創造的時候,　神造人是造男造女.     7     因此人要離開父母,與妻子連合,二人成為一體.

耶穌問那些法利賽人,摩西的吩咐是什麼?律法書的經文寫了什麼?

法利賽人回答說:「摩西許人寫了休書便可以休妻.」

的確,在舊約申命記 24:1明明地寫著:「人若娶妻以後,見她有什麼不合理的事,不喜悅她,就可以寫休書交在她手中,打發她離開夫家.」不同的猶太拉比,著眼點在於什麼叫「不合理的事」.他們的不同立場,源自他們對這經文的不同解釋.

主耶穌處理這經文時,不是著眼於什麼叫「不合理的事」,然後再多加一種解釋.主耶穌說,摩西這樣寫,是因為你們心硬.人面對婚姻的困難,不肯再付出,不肯提升自己的生命品質,不肯克服自己的罪性,堅持以離婚解決.如此心硬,摩西只能讓步.面對人性的軟弱的現實,摩西作出讓步.但主耶穌在這裡不肯讓步.他不贊同摩西這種權宜之策.人性是軟弱的,摩西很了解,故此,作出適切的讓步.人性是軟弱的,主耶穌也很了解,但是,他不從解決現實問題的角度去處理問題,他希望人看到更廣闊的世界,他希望人追求更高的價值,他希望人能活出更豐盛的人性.在這一點上,他不肯讓步.

《小王子》那本小故事書的作者Antoine de Saint Exup\'ery說:「假如你要建造一隻船,不要召集人去收集木頭,不要指派他們這樣任務那樣工作,而是要教導他們渴望大海的無邊廣闊.」( ``If you want to build a ship, don"t drum up people to collect wood and don"t assign them tasks and work, but rather teach them to long for the endless immensity of the sea." )

面對一個問題,很多時,重點不在做這做那,而在看到一個遠象(vision).有遠象,才有處理問題的真正方向.

主耶穌說:「但從起初創造的時候,神造人是造男造女.因此,人要離開父母,與妻子連合,二人成為一體.」

主耶穌將人的心靈打開,讓人看到真正遼闊的世界.一個比現實更真實的世界.

主耶穌叫法利賽人看,我們實在是活在這神聖世界中.婚姻的意義,不是由人去定義的.不是說,合則來,不合則去.這樣的婚姻只能滿足人的慾望,不能成全人性；只強化自私,不增加愛.主耶穌說,「從起初創造的時候」,神設立婚姻,在其中,人性得以完滿.這需要人學習自制,學習放下自私,學習為他人而活,學習為了他人而改變自己的品格德行,學習真正的愛.

在這學習愛的過程裡,人不是孤獨的,因為神在人間.主耶穌說,「神造人是造男造女」.神的創造力,在這男女之間顯明出來.人誠然要努力,但更要相信神.沒有對神的信心,不依靠神的真理,不仰賴神的幫助,人就很難在艱難的婚姻生活中走下去.

主耶穌說:「人要離開父母,與妻子連合,二人成為一體.」在婚姻裡,孤獨的,自私的,自利的個我,漸漸消失,而與愛人成就一愛的生命體.在這裡,人會發現,為了愛喪掉自己生命的,才得到生命.

主耶穌讓我們看到,我們不單只活在現實意義的世界裡,也活在神聖意義的世界裡.在這神聖意義的世界裡,人活著的意義,不單單建基於當下及將來,也建基於天地初開以來內注於人類生命裡的意義.在這神聖意義的世界裡,不單只是人類活著,神也活著,神活在人類之上,活在人類之中,也活在人類生命裡.在這神聖意義的世界裡,婚姻不單只是兩個人相愛相戀,或者說,不單只是「兩個客體的客觀距離接近了」,而是兩個人的生命互相滲入對方,以致你中有我,我中有你.我不再為自己而活,而為你而活；也因為有你在我裡面活著,我的生命才得以全面展現.在這種愛的交流裡,「神的形象」得以彰顯出來.

要建造一隻船,需要我們渴望大海的遼闊；同樣,要處理婚姻的艱難,需要我們想像人性存在的神聖深度.若果你不渴望大海,則在造船過程中,一塊木頭引起的難題,一顆釘子引起的痛苦,便叫你放棄.若果我們不渴望在婚姻裡展現生命的神聖深度,則一碟弄糟了的菜,旁邊一個漂亮的美女,便叫你離婚.

明顯,耶穌的重點不在解決「現實」的問題,他叫人看到現實背後的神聖意義.摩西的休妻律法,是要解決人的「現實」問題.法利賽人的問題焦點,是要解決人的「現實」問題.耶穌的焦點,卻不在這裡.在造船的過程中,若有人被一顆釘子戳穿了手指,這人問摩西:「不再造船,可以不可以?」摩西因為他放棄的心已決,只能說:「可以.」現實問題,現實解決.但在同樣情況下,這人問耶穌:「不再造船,可以不可以?」耶穌卻遠望大海,然後拍一拍他的肩膀,說:「看到嗎?多麼廣闊的大海.你願意和我出去看看嗎?」

我在教會(包括崇基禮拜堂)作了差不多二十年的小傳道,無數弟兄姊妹問過我無數「可不可以」的問題.小朋友問:「可不可以不返學?」再大一點的小朋友問:「可不可以不返教會?因為教會的聚會很悶.」青少年信徒問:「可不可以拍拖?」再大一點的青少年信徒問:「可不可以同時有兩個拍拖對象,一個正選,一個後備?因為導師教導我們,凡事都要有兩手準備.」青年信徒問:「可不可以和這人結婚?」再大一點的信徒問:「可不可以和這人離婚?」

面對無數的「可不可以」,在這裡,主耶穌教我們如何思考.這就是今日講道題目「不要那樣問,而要這樣想」的用意.不要常常問某一行為「可不可以」,而要多思考一下,產生這問題的事件,若放在一更廣闊的場景裡,將有何意義?而你在這意義的開展過程中,你將扮演什麼角色?

正如你在畫一幅畫,你問老師:「可不可以在這一點上,加一點顏色?」這個問題其實不能局部地看,這個問題牽涉的,不是「這一點」的問題.老師不能抽空地回答是否能在「這一點」上加一點顏色,因為「這一點」不只是關係到「這一點」,而是關係到全幅畫面.老師只能叫你思考,在全幅畫面裡,這一點扮演什麼角色?這一點的用意在哪裡?這一點色彩和畫面其他部分的色彩是否達致融合?這一點的用色是否帶出整個畫面的全面效果?故此,很難抽空地說,在這一點上,可不可以加一點顏色.

就著法利賽人的問題,主耶穌的確沒有處理行為的某一點上的「可以不可以」的問題.他沒有給出簡明的答案.他回答的,是教人得到答案的思考過程本身.他教人「不要那樣問,而要這樣想」.他教人不要將「離婚」問題只就「離婚」問題來處理.就著「離婚」問題,他教人思考人生在世的意義.這意義是源自自我的快樂追尋,抑或源自人的存在要開發的真愛?他教人思考,人是只向自己負責,抑或面向神而向神負責?他教人思考,人生的完滿,是來自自我的潛能本身,抑或我們需要一外在的他者?他教人思考,愛情是一自我感覺,抑或牽涉到人與他人的生命相連?

主耶穌的用意,不在於點出「可以不可以」,而在於指出,若要好好處理這問題,便需要這種思考過程.這種思考過程,需要你看到一個比你大得多的神聖世界.

主耶穌希望人學曉,不從「可以不可以」的方式處理人生,而要學曉看到生命的全幅神聖圖畫,追求一種配合這圖畫的生命質素,並倚靠神的恩典努力將這豐盛生命活出來.在這意義下,主耶穌講出以下一句看似「不可能」而需要神蹟才能將之變成「可能」的話.

馬可福音 10:8-9

8     既然如此,夫妻不再是兩個人,乃是一體的了.     9     所以　神配合的,人不可分開.

作為基督徒,我只能說,我相信恩典,也相信神蹟.

但願:榮耀歸於聖父,聖子,聖靈.阿們.



\newpage

\section{負傷的治療者}
\label{sec:171}
\textbf{伍渭文牧師 - 負傷的治療者}
\newline
\newline
\hyperref[sec:170]{< < < PREV SERMON < < <}
~
\hyperref[sec:toc]{[返目錄]}
~
\hyperref[sec:172]{> > > NEXT SERMON > > >}


以賽亞書 52:13-53:12
\newline
\begin{longtable}{cl}
\hline
\hline
章節 & 經文 (和合本修訂版)\\
\hline
52:13 & \begin{tabularx}{0.7\textwidth}{X} 看哪，我的僕人行事必有智慧，他必被高升，高舉，升到至高之處。 \end{tabularx} \\ \\ \relax
52:14 & \begin{tabularx}{0.7\textwidth}{X} 許多人因他驚奇---他的面貌比別人憔悴，他的外表比世人枯槁--- \end{tabularx} \\ \\ \relax
52:15 & \begin{tabularx}{0.7\textwidth}{X} 同樣，他也必使許多國家驚奇，君王要向他閉口。未曾傳給他們的，他們必看見；未曾聽見過的事，他們要明白。 \end{tabularx} \\ \\ \relax
53:1 & \begin{tabularx}{0.7\textwidth}{X} 我們所傳的有誰信呢？耶和華的膀臂向誰顯露呢？ \end{tabularx} \\ \\ \relax
53:2 & \begin{tabularx}{0.7\textwidth}{X} 他在耶和華面前生長如嫩芽，像根出於乾地。他無佳形美容使我們注視他，也無美貌使我們仰慕他。 \end{tabularx} \\ \\ \relax
53:3 & \begin{tabularx}{0.7\textwidth}{X} 他被藐視，被人厭棄；多受痛苦，常經憂患。他被藐視，好像被人掩面不看的一樣，我們也不尊重他。 \end{tabularx} \\ \\ \relax
53:4 & \begin{tabularx}{0.7\textwidth}{X} 他誠然擔當我們的憂患，背負我們的痛苦；我們卻以為他受責罰，是被神擊打苦待。 \end{tabularx} \\ \\ \relax
53:5 & \begin{tabularx}{0.7\textwidth}{X} 他為我們的過犯受害，為我們的罪孽被壓傷。因他受的懲罰，我們得平安；因他受的鞭傷，我們得醫治。 \end{tabularx} \\ \\ \relax
53:6 & \begin{tabularx}{0.7\textwidth}{X} 我們都如羊走迷，各人偏行己路；耶和華使我們眾人的罪孽都歸在他身上。 \end{tabularx} \\ \\ \relax
53:7 & \begin{tabularx}{0.7\textwidth}{X} 他被欺壓受苦，卻不開口；他像羔羊被牽去宰殺，又像羊在剪毛的人手下無聲，他也是這樣不開口。 \end{tabularx} \\ \\ \relax
53:8 & \begin{tabularx}{0.7\textwidth}{X} 因受欺壓和審判，他被奪去，誰能想到他的世代呢？因為他從活人之地被剪除，為我百姓的罪過他被帶到死裡。 \end{tabularx} \\ \\ \relax
53:9 & \begin{tabularx}{0.7\textwidth}{X} 他雖然未行殘暴，口中也沒有詭詐，人還使他與惡人同穴，與財主同墓。 \end{tabularx} \\ \\ \relax
53:10 & \begin{tabularx}{0.7\textwidth}{X} 耶和華的旨意要壓傷他，使他受苦。當他的生命作為贖罪祭時，他必看見後裔，他的年日必然長久。耶和華所喜悅的事，必在他手中亨通。 \end{tabularx} \\ \\ \relax
53:11 & \begin{tabularx}{0.7\textwidth}{X} 因自己的勞苦，他必看見光就心滿意足。因自己的認識，我的義僕使許多人得稱為義，他要擔當他們的罪孽。 \end{tabularx} \\ \\ \relax
53:12 & \begin{tabularx}{0.7\textwidth}{X} 因此，我要使他與位大的同份，與強盛的均分擄物。因為他傾倒自己的生命，以致於死，也列在罪犯之中。他卻擔當多人的罪，為他們的過犯代求。 \end{tabularx} \\ \\
[1ex]
\hline
\hline
\end{longtable}
日期:2012年11月11日

題目:負傷的治療者

經文:以賽亞書52:13-53:12

證道:伍渭文牧師

最近我們常常聽到這城市有病了,這城市是世界上貧富最懸殊的地方,我們社會存在不少深層次矛盾；近邊境的人賣不到奶粉,廣東道的名店不歡迎不懂說普通話的人.......若我們的城市有病,需要的不是斷症的人,我們需要治病的人.聯合國秘書長潘基文在前天「世界消除貧窮日」說:「要痛斥貧窮很容易,但要戰勝它卻很困難.遭受饑餓,匱乏和屈辱之苦的那些人需要的不僅僅是同情的言語,他們還需要具體支援.」(Poverty is easy to denounce but difficult to combat. Those suffering from hunger, want and indignity need more than sympathetic words; they need concrete support.) 城市的病,沒有妙藥靈丹,一服即愈(cure) ；我們需要守望這城市,關心這城市(care) ,要當城市的傷痛為自己的傷痛,才能為城市裹傷治病.

作為牧者,初到履新的教會,可能看到不少問題,洞察到多年累積下來的「病患」,愛護教會有經驗的牧者,不會祇給教會斷症,務求很快治療(cure) ,他會當這些負債是自己要償還的債務,如果說這是教會的傷痛,他就是負傷的牧者,與會友一起承擔,贏取會友的信任,同時進行治療,這牧者能夠有效在教會進行牧靈工作,他是一位負傷的治療者.

今天講道經文以賽亞書五十二13至五十三12就提到受苦的僕人,這位受苦的僕人就是一位「負傷的治療者」.

五十二13至五十三12是以賽亞書四首僕人之歌最後一首(42:1-7; 49: 1-9; 50:4-9; 52: 13-53: 12), 簡單來說以賽亞全書可分兩部分:1-39章說到上帝的審判,亞述國(伊朗)興起要南下攻打諸國,首當其衝的亞蘭(敘利亞)北方以色列國用軍事力量迫使南方的猶大結盟對付勢力強大的亞述:

「猶大王亞哈斯在位的時候.....有人告訴大衛家說:"亞蘭與以法蓮已經同盟." 王的心和百姓的心就都跳動,好像林中的樹被風吹動一樣.」(賽七1,2)以賽亞安慰亞哈斯不要懼怕他們,上帝會施行拯救,但不要和亞述國對抗.亞述國是上帝的工具,把拜偶像的北方以色列滅亡,他們是因拜偶像而遭滅亡的,而猶大若不猛然悔改,離開偶像,也會滅亡被擄.

下半部四十章至六十六章提到上帝的拯救,作者筆鋒改變,不作詳細的歷史描述,轉到上帝拯救的應許.先知受聖靈感動,彷彿看到二百年後將要發生的歷史,此時猶大亦因拜偶像被興起的波斯王國消滅,國人被擄巴比倫；但上帝大能的膀臂,會施行拯救,使他們歸回故土.四十1-3是神劇彌賽亞第一樂章的

主題:「安慰,安慰我的百姓」(Comfort ye, comfort ye my people) .「她爭戰的日子已滿；她的罪孽赦免了.」施洗約翰在曠野一開始就宣告這盼望,正如先知以賽亞書上記著說:

「一切山窪都要填滿,大小山岡都要削平高高低低的要改為平坦,崎崎嶇嶇的必成為平原.耶和華的榮耀必然顯現；凡有血氣的必一同看見,因為這是耶和親口說的.」(賽四十4, 5；參可一1, 2)

經文中受苦的僕人就是彌賽亞,但彌賽亞為何要受苦呢?為人贖罪祇需被釘在十字架,因為罪的工價就是死.但彌賽亞來不是去擺平罪(cure, quick fix)-處理罪的事件；他來是要關心罪人(care) ,他認同我們,承受罪帶來苦難.基督受苦,乃揭示罪如何扭曲人性,揭示罪如何破壞人與人的關係,破壞人與上主的關係:

「他被欺壓,在受苦的時候卻不開口；他像羊羔被牽到宰殺之地.」(賽五十三7)他生活在世上,承受扭曲人性的墮落世界帶給他的傷害,但靠著聖靈的能力,逆轉荒謬為意義,咀咒為讚美.

1.       他為罪被壓傷

「他在耶和華面前生長如嫩芽,像根出於乾地.他無佳形美容,我們看見他的時候,也無美貌使我們羨慕他.他被藐視,被人厭棄,多受痛苦,常經憂患.他被藐視,好像被人掩面不看的一樣,我們也不尊重他.他誠然擔當我們的憂患,背負我們的痛苦.」(五十三2-4a)

被藐視,被厭棄,不被尊重,是否是這城市某些人的感受呢?你們不屬於這裡!我們不歡迎你們!是否是這城市某些人的喊叫呢?

歧視:他被藐視,被人厭棄

在馬利亞的胎中,被聖靈感孕的基督就被岐視為非婚生之子illegitimate child,馬利亞也承受著未婚生子的污名,約瑟也曾經想暗暗的休了瑪利亞,直到耶穌公開出來傳道,行了第一宗神蹟,人們才明白和接受馬利亞是聖靈感孕,耶穌並非私生子(illegitimate child).

非婚生子女在香港沒有居留權大家都知道,他們等候被遺返內地前,不能在香港政府資助的學校讀書.以前「非婚生」有道德的價值判斷,但現在「非婚生」觀念重點不再是道德價值判斷,仍是既得益者用來保障自己的利益,非婚生不能合法分享權益,幸好有教會學校不理教育局,讓他們等候遺返內地期間接受教育.大量在內地城市工作的農民工,因為沒有戶口,他們的子女相對城市居民的子女是不合法的,不能享受均等教育機會,因為城市的居民,怕分薄了他們既得的教育資源.

2007年的除夕夜中央電視台節晚會,三十個農民工子女在詩歌朗誦《心裡話》,就道出了受歧視的心聲,令人感動:

``要問我是誰,過去,我總不願回答,因為我怕,我怕城裡的孩子笑話.

我們的校園很小,放不下一個鞍馬；我們的校舍簡陋,還經常搬家；

我們的教室很暗,燈光只有幾瓦；我們的桌椅很舊,坐上去吱吱啞啞.

但是我們作業工整,我們的成績不差.

要問我此刻最想說什麼,我愛我的媽媽,我愛我的爸爸.

因為-是媽媽把城市的馬路越掃越寬,

因為-是爸爸建起了新世紀的高樓大廈.

作文課上,我們寫下了這樣的話,別人與我比父母,我和別人比明天!

打工子弟和城裡的小朋友一樣,都是中國的娃,更是祖國的花."

被拒絕和出賣:好像被人掩面不看的一樣,我們也不尊重他

因著逃避希律王的刀劍,幼童耶穌成為政治難民,隨從約瑟馬利亞逃難埃及,遠離家鄉.代表羅馬嚴明法律的巡撫彼拉多,明知耶穌沒有犯罪,但屈從民粹壓力,在逾越節特赦釋放大盜巴拉巴,把耶穌釘在十字架,這是政治現實,耶穌早有心理預備；但令耶穌心碎的,是被跟隨他三年多的門徒出賣和否認.猶大以親嘴為記號,為了三十兩銀子,出賣自己的老師；曾公開對耶穌說:

「主啊,我們就是你下監,同你同死,也是甘心.」(路二十二33)的彼得,竟然三次不認主.

公義,忠誠,友愛因為罪被犧牲了,罪扭曲了人性.

這位受苦的僕人,在世上經歷罪帶來種種的傷害,所以他明白為何一位校長所扶掖作他副手的繼任人,上任後竟然拒絕栽培過自己的上司,不讓卸下行政責任的校長,留下來繼續教學；新校長告訴校董會,很難和離任的校長合作,但以前他們不是合作無間嗎?耶穌明白,因為他也曾被自己所愛的門徒出賣過,他可以醫治我們被罪壓傷的心靈.

「他也在各方面受過試探,與我們一樣,只是他沒有犯罪. 所以,我們只管坦然無懼地來到施恩的寶座前,為要得憐憫,蒙恩惠,作及時的幫助.」(來四15b, 16)

2. 為我們的罪孽被壓傷

「我們卻以為他受責罰,被神擊打苦待了.哪知他為我們的過犯受害,為我們的罪孽壓傷.因他受的刑罰,我們得平安；因他受的鞭傷,我們得醫治.我們都如羊走迷,各人偏行己路」(五三4b-6)

「負傷的治療者」是已故耶魯大學靈修學教授盧雲神父在1972年一本書的書名,這書風行一時,再版無數次,影響對治療者的角色和觀念.他指出人的疾病和心靈息息相關,亦和社會文化氛圍密密緊扣.一般醫學儀器祇能測試到表像,但深層的心靈傷痛就不容易量度.若我們明白情緒和身體息息相關,作為治療者就要進入對方心靈了解真相；但要對方開放心靈,治療者首先要開放自己.當我們面對人生的限制,感到無助,接受自己真正無助時,像耶穌在客西馬利園,面對十字架感到無助,才能依靠復活的盼望.一個不認識自己的限制,以為自已很有辦法和本領,不可能成為有效的治療者.當治療者不怕呈現自己的脆弱 ,就幫助對方不怕面對自己的脆弱,治療就成功了一半.耶穌豈不是說過:

「凡勞苦擔重擔的可以到我這裡來....我心裡柔和謙卑」(太11: 28, 29).

耶穌柔和謙卑,是他認識我們的軟弱,體貼我們,接納我們；而我們作為治療者心裡柔和謙卑,因為是認識自己的軟弱.

「我們卻以為他受責罰,被神擊打苦待了.哪知他為我們的過犯受害,為我們的罪孽壓傷.」是我們的罪使基督受傷,當我們看到城市的罪惡時,我們有否想到若不是主的恩典,我可能在其中也有份.

雅五19「我的弟兄們,你們中間若有人迷失真道的,有人使他回轉.這人該知道叫一個罪人從迷路轉回,便是救一個靈魂不死,並且遮蓋許多的罪.」信主不單使我們靈魂不死,而且有聖經作為腳前的燈,路上的光；有聖徒群體的伴隨和鼓勵,追求生命聖潔,免陷在罪惡中.我們可以想想,若不是信主後我們曉得不可作假見證,所以我們不敢與同事互租物業居住,作沒有利益衝突的虛假聲明,避了被廉正公署起訴的一劫.

Yehiel Dinur (耶喜.狄魯) 是納粹大屠殺(holocaust) 的幸存者,在審訊納粹頭號戰犯Adolf Eichmann(阿道夫.艾希孟)作按方証人.狄魯 直視被防彈玻璃後邊艾希孟,就像曾經被牽到宰殺之地的的羊面對屠夫.突然,狄魯飲泣,不久暈過去,不是因為傷痛或忿怒,他後來解釋:「那一刻我突然驚慄起來,我是懼怕自己,我可以成為他一樣,做出這樣可怕的事.這人是惡魔?是瘋子?還是令人感到心寒-是正常的?艾希孟就在我們的心中.」

當我們明白艾希孟也在我的心中時,我們心裡不得不柔和謙卑,我們就不會對人過份嚴苛,我們的心靈就能更加貼近被罪壓傷的人；我們就成為負傷的治療者.

3 暴力在他身上停住了

「他被欺壓,在受苦的時候卻不開口；他像羊羔被牽到宰殺之地.又像羊在剪毛的人手下無聲,他也是不開口.......耶和華卻定意將他壓傷,使他受痛苦；耶和華以他為贖罪祭.他必看見後裔,並且延長年日,耶和華所喜悅的事必在他手中亨通.他必看見自己的勞苦的功效,便心滿意足.」(賽五十三7, 10-11a)

經文說:

「他被欺壓,在受苦的時候卻不開口」,因此,暴力就在他身上停止了.以暴易暴祇產生暴力循環,這位受苦的僕人在受苦的時候不作反擊,暴力就消解了.

孟特拉因反對分離主義,被囚在惡劣的羅德島(Robben Island) 監獄二十七年,成為南菲總統後,因為基督的信仰,他沒有報復,反而成立由大主教圖圖(Tutu) 主持的『真相與和解委員會』 (Truth and Reconciliation Commission),讓受種族隔的受害者傾訴苦情,又讓施行罪行的人公開尋求饒恕,使對立的群族,得以和好.在他的自傳「自由的長征Long Walk to Freedom」,他指出因偏見和恐懼形成的隔離主義,對國家做成嚴重的傷害和撕裂,要幾代才能愈合,但人在那麼深的壓制下,才能造就出高貴的品格-饒恕,接納,和解的品格.他說:「我的國家地下蘊藏豐富的珍礦寶石,但我知道最珍貴的還是人民,他們比最純的鑽石更貴重和更真實」(My country is rich in the minerals and gems that lie beneath its soil, but I have always known that its greatest wealth is its people, finer and truer than the purest diamonds) .

社會需要治愈

以賽亞先知也引導被擄巴比倫的猶太人,看到苦難的積極意義,看見整個民族其實和人類的命運緊扣在一起.他們被擄,是被淨煉,徹底脫離偶像敬拜,預備彌賽亞的來臨,更新世界秩序,在新的秩序中:

「沒有數日夭亡的嬰孩,也沒有壽數不滿的老者......他們建造房屋,自己居住；栽種葡萄園,吃自己的果子.」(賽六十五20, 21)猶太人的神觀,藉被擄播遷異域擴張了,上帝是萬國之上主,上主關心萬族的福祉.沒有數日夭亡的嬰孩,因為政府注重公共衛生,讓最不能保護自己的嬰孩得到保護；沒有壽數不滿的老者,因為社會保障健全,讓老有所養；住自己建造的房屋,吃自己栽種的果子,因為公平貿易,沒有中間層層剁削.被擄回歸的上帝子民,要在全地傳揚上帝的救恩,而我們曾經如羊走迷,曾經被罪綑綁的人,就是上帝的子民,我們曾經被擄,為罪的奴僕,現在要被主差遣,安慰這城的百姓.

新時代已經降臨了! 「安慰,安慰我的百姓」(Comfort ye, comfort ye my people) .

若我們的城市有病,有很多深層次的矛盾,群族對立；這時候最需要的不是懂得疹斷的人,我們需要的負傷的治療者,守望這城市,關心在其中被罪壓傷的人,為他們包裹傷痛,背負他們的重擔,關懷他們.

負傷的治療者,曾像羊迷失過,現在不再迷失了；負傷的治療者,也許曾經不尊重主,出賣過自己的信仰,但當我們回轉悔改,像彼得一樣,得蒙赦免後,去安慰其他的人:「安慰,安慰我的百姓」(Comfort ye, comfort ye my people) .

因為這呼召,負傷的治療者,願意重新步履受苦的僕人的腳蹤,為這城市承擔憂患,化解冤仇.那受苦的僕人,已為我們走上了十字架；讓我們要為他走進人群中,成為負傷的治療者.阿們.



\newpage

\section{讓我們彼此相愛}
\label{sec:172}
\textbf{張伯笠牧師 - 讓我們彼此相愛}
\newline
\newline
\hyperref[sec:171]{< < < PREV SERMON < < <}
~
\hyperref[sec:toc]{[返目錄]}
~
\hyperref[sec:173]{> > > NEXT SERMON > > >}


日期:2012年11月18日

題目:讓我們彼此相愛

經文:約翰一書4:7

證道:張伯笠牧師

在馬可福音12章有一個文士和耶穌基督有一段對話.

有一個文士來,聽見他們辯論,曉得耶穌回答的好,就問他說,誡命中那是第一要緊的呢?耶穌回答說,第 一要緊的,就是說,以色列阿,你要聽.主我們神,是獨一的主.你要盡心,盡性,盡意,盡力,愛主你的神.其次,就是說,要愛人如己.再沒有比這兩條誡命更 大的了.那文士對耶穌說,夫子說,神是一位,實在不錯.除了他以外,再沒有別的神.並且盡心,盡智,盡力,愛他,又愛人如己,就比一切燔祭,和各樣祭祀, 好的多.耶穌見他回答的有智慧,就對他說,你離神的國不遠了.(可12:28-34)

約翰曾被稱為「雷子.」

約翰(John,亦譯「若罔」):西庇太的兒子,雅各的弟弟.耶穌最喜愛的門徒.與哥哥雅各一起被耶穌成為「雷子」.耶穌被捕後,曾跟到大祭司的院子裡去打聽消息.耶穌受難時十二使徒只有他一個人守在十字架下,耶穌就把母親瑪麗亞托付給他照看.

耶穌升天後,約翰和彼得搭伴積極進行傳道活動,被保羅稱為「教會的柱石」.約翰曾因傳道被羅馬政府放逐到小亞細亞以西的拔摩海島(亦譯「巴毛斯」或「帕特摩」),在那裡見了異象,寫成《啟示錄》.後來,約翰獲釋在以弗所傳道,並把瑪麗亞接到以弗所去奉養.在聖靈感動下,著有《約翰福音》與三封《約翰書信》.於公元104年去世.

先讓我們來看使徒約翰的自薦:使徒約翰在自我介紹中,一開始就謙卑地以普通弟兄的身份出現說:「我約翰就是你們的弟兄,和你們在耶穌的患難,國度,忍耐裡一同有分.」

我約翰就是你們的弟兄

請看,此時年老的使徒約翰和年輕時的雷子約翰已經完全變成兩人.那時他被稱為「雷子」,生性暴燥而又野心勃勃,當他作了主的門徒三年多時間後,還曾和他的哥哥雅各布一同通過他們的母親,要求主耶穌將來使他們兄弟倆人坐在主寶座的左右.

然而此時,被「稱為教會柱石」,德高望重,碩果僅存的年老使徒約翰,卻早已除盡了他年少時的暴烈性情和自大野心,而完全變成了他夫子的心地柔和謙卑的樣式,以致他現在甘願和弟兄們處於同等地位,自稱為「我約翰就是你們的弟兄」,並以能和弟兄們「在耶穌的患難,國度,忍耐裡一同有分」而感到歡喜快樂.

但也正在這時,這位謙卑柔和,愛主愛人,服事弟兄,服事眾人的主的忠僕,卻已在上帝的國中被升為至高,

這正如主所教導我們的:「因為自高的必降為卑,自卑的必升為高.」(路 14 : 11 ).因此,「你們裡頭為大的,倒要象年幼的,為首的,倒要象服事人的.是誰為大,是坐席的呢?是服事人的呢?然而我在你們中間如同服事人的.」(路 22 : 26 － 27 ).

另一次主也教導我們說:「你們中間誰願為大,就必作你們的用人,誰願為首,就必作你們的僕人.正如人子來,不是要受人的服事,乃是要服事人,並且要捨命,作多人的贖價.」(太 20 : 26 － 28 ).

從「雷子」約翰,轉變為大有愛心的心靈柔和謙卑的使徒約翰的事例上,也讓我們看到了福音改變人心大能的一個偉大的例證.

講到這裡,讓我們也省察一下自己:「這福音本是上帝的大能」,是否也在我們身上完成了一番奇妙的改變?

我們是否更謙卑,更溫柔,更聖潔,更愛主愛人,更完全獻身,更熱切地服務行善和傳道救靈?但願主救贖的大能也在我們每一個人的心內和身上完成這種偉大的改變.

有分於耶穌的患難

拔摩海島本是一個荒無人煙的海島,為當時流放犯人的地方.使徒約翰卻是為了主的傳道救靈聖工而被羅馬帝國充軍到這個荒島上.在這之前,羅馬當局曾將約翰 投人滾沸的油鍋中,妄圖殘害他.但約翰卻蒙上帝神奇的救護,毫無傷害.羅馬兵丁大衛驚奇,就將他從油鍋中拉上來.由於無法殺害約翰,但又不願讓他在弟兄姊妹和世人中間繼續發揮影響力,羅馬當局就下令將他放逐到拔摩海島上.

約翰先介紹他自己是誰,是在怎樣的情況下得到啟示.

A. 是弟兄

「我約翰就是你們的弟兄.」(v.9)

約翰只普通地稱自己是弟兄.他得了這麼重要的啟示, 一點也沒有要炫耀自己,覺得自己高人一等的意思.今天不少人事奉上有了些成就,恩賜上蒙神賜了些特別恩賜,就儼然以權威自居,仿似是先知,是聖靈受膏者,甚至是神自己.從約翰的自我介紹,就可分別出真假屬靈,和真假先知來.

B. 是為主受苦的

「和你們在耶穌的患難,國度,忍耐裡一同有分」(v.9)

約翰也介紹,他是和所有的信徒一樣,是因為耶穌的 國度而受苦的.為主受苦,是一句人不愛聽的話,但卻是當時許多信徒所遭遇的事,也是主再來的時候,將要發生在信徒身上的事.主再來前夕,不少愛主的人,將要為主

的原故,為神國度的原故而受苦.然而他們不是孤單的, 因為主耶穌自己走過,約翰也在這條路上走過.這帶給受 苦中人真是無比的安慰和鼓勵.

弟兄姊妹,主來的日子近了,真信主的人,日子將不會好過了.因為撒旦會動用一切可動用的力量來對付信徒,但我們不用怕,因為在我們前面有不少人都這樣走過去了. 而啟示錄就是向我們預告這些事,並宣告耶穌己經得勝.

C. 作見證的

「為神的道,並為耶穌的見證,曾在那名叫拔摩的海島上.」(v.9)

他為什麼會受苦?為什麼會被囚在拔摩島?不是他犯 了什麼罪,只因為他傳神的道.人若堅守神的道,持守聖 經的原則,一定會招來非議,甚至逼迫.相反,你若不堅持聖經原則,肯妥協,那你就能和其它的人,甚至是撒旦和敵基督都能相安無事,甚或大受歡迎.但約翰卻不願隨 污逐流,結果被放逐.

約翰為什麼會受苦?只因他為耶穌作見證.他若不為主作證,他就無事.他若不為主作證,他就可以和其它的 人一樣,過著舒適的生活,享受天倫.然而他揀選了要為 那愛他的主作見證,為那曾為他釘十架的主作證以至被放逐. 有分於耶穌的患難

拔摩海島本是一個荒無人煙的海島,為當時流放犯人的地方.使徒約翰卻是為了主的傳道救靈聖工而被羅馬帝國充軍到這個荒島上.在這之前,羅馬當局曾將約翰 投人滾沸的油鍋中,妄圖殘害他.但約翰卻蒙上帝神奇的救護,毫無傷害.羅馬兵丁大衛驚奇,就將他從油鍋中拉上來.由於無法殺害約翰,但又不願讓他在弟兄姊妹和世人中間繼續發揮影響力,羅馬當局就下令將他放逐到拔摩海島.

教會蒙福的條件之一就是愛,在愛的基礎上,教會得以建立；在愛的基礎上,教會得到興旺.教會中有彼此相愛的事實,這教會必定是蒙神悅納,蒙神祝福的.

有一本書是一個猶太的拉比寫的,裡面記載了一些聖經上所沒有寫的故事.是講到建造聖殿的那塊地的來歷.那塊地是耶布斯人阿珥楠的禾場,原來阿珥楠早年死了父親,他還有一個年幼的弟弟由阿珥楠撫養成人.他們兄弟二人非常相愛.後來弟弟成年了,按照傳統他們要分家產,在分家產的時候,有一塊地是多出來的,大家都不肯多要.哥哥要弟弟拿去,弟弟要哥哥拿去,讓來讓去,誰都不肯接受,每當耕種的時候,大家都搶著去種,每當收割的時候,大家都搶著去收割,為的是要把收割下來的穀物送到對方的家裡去.

有一年莊稼正熟了,兩個人都明白了對方是為什麼來的,大家擁抱在一起,熱淚滾滾地流下來.從此他們約定,這塊地既然誰都不肯要,就願意有一天得到機會把它獻給神用.後來大衛就得著了這塊地作建造聖殿之用.舊約中的聖殿就是建造在這塊標誌著弟兄相愛的基地上.

一, 彼此切實相愛(彼前1:22)

1, 「你們既因順從真理,潔淨了自己的心,以致愛弟兄沒有虛假,就當從心裡.」﹙彼得前書一章22節﹚切實相愛就是不虛偽.人生有兩種悲哀,其一 是你與人相處覺得像冰一樣,使你非常痛苦.其二使你更悲傷的是偽裝的愛.真正的愛是出於被真理潔淨的心,是單純的心.

保羅在《羅馬書》第13章8節到10節的三節經文中他講了這一段話:

凡事都不可虧欠人,惟有彼此相愛,要常以為虧欠.因為愛人的就完全了律法.像那不可姦淫,不可殺人,不可偷盜,不可貪婪,或有別的誡命,都包在愛人如己這一句話之內了.愛是不加害與人的,所以愛就完全了律法.(羅13:8-10)

我們在這裡看到保羅一直在講愛,講愛人如己,愛人如己這句話非常容易理解,什麼叫愛人如己呢?愛人如己就是愛別人就像愛自己一樣.

我們在羅馬書中看到保羅說,愛是什麼呢?愛是基督徒的行事為人的準則.他告訴我們說愛是成全了全部的律法.

首先,凡事都不可以虧欠人

保羅在這裡講,凡事都不可以虧欠人.保羅首先提醒我們說,凡事都不可以虧欠人作為我們處事為人的基本原則上.凡事都不可以虧欠人是說,任何事情都不要使別人吃虧. 我們若能遵守耶穌基督的吩咐,我就辦事不讓別人吃虧.如果你不能按照主的吩咐,那你還是要佔便宜.

主耶穌說:「所以無論何事,你們願意人怎樣待你們,你們也要怎樣待人.因為這就是律法和先知的道理.(太7:12)」

二,互相款待彼此服事﹙彼得前書第四章9-11節﹚

教會裡另一個問題是彼此比較,競爭.如果用自己的長處去量度別人的短處,會覺得人人都差勁.若拿自己的短處來量度別人的長處,你會自卑.彼此情景不同,但合起來卻是一幅美麗的圖畫.

我給你們講一個小故事,一個驢和馬的故事.這是伊索寓言,我們小時候常常讀伊索寓言:

一個商人有一 頭驢和一頭馬,馬是他騎的,驢是運貨的.有一天他去運貨去,驢擔當的貨物多了一點,他就把多的貨物放到馬的身上一點.馬就想:我是被騎的,不是拉貨的.它就很難過,它不敢對主人發牢騷,它怕主人用鞭子打它.它就向驢發牢騷,它說:「你瞧你,就那點力氣,把活都壓在我的身上.」

那天驢正巧感冒,不太舒服.驢就說:「馬大哥,我今天不太舒服,我現在擔不動了,你能不能在替我擔點?」馬就說:「這已經不少了,你還給我?原來我都一點 不擔的.」驢真的很難受,馬就一路埋怨,驢就一句話也不說,就往前走,可是走著走著,這驢就不行了,它就說:「馬兄啊!我快不行了.」馬說:「我不管你,你愛怎麼樣有怎麼樣.要死你就死.」這個驢就倒下了,累死了.

驢死了以後,商人就把驢身上的貨物都堆到馬的身上去了,這個馬就更後悔了,「剛才我要幫它擔一點,它活著,我們兩個總不這麼重吧!」而且連那個商人也騎到它面上去了,它更慘了.

在我們與人相處的過程當中,我們要學習彼此的配搭,學習同工們,團隊們在一起時候,能夠彼此的去建造.看到別人忙了就去幫一把手.我經常說我們的教會原則:第一次來是客人,第二次是主人,第三次來是僕人.這樣我們成長的更快.教會的弟兄姊妹在主的愛裡面要這樣的成長.

以弗所書2:19說的,我們「不再作外人和客 旅,是與聖徒同國,是神家裡的人了.」 一個美滿的家庭,一定有使家庭溫暖美好因素.那麼,在這個屬靈的家裡面,有什麼因素使之成為神的家?

1. 彼此相愛:在約翰福音15章裡面記載耶穌基督對門徒反覆講到你們要彼此相愛,(約13:35)也說 「你們若有彼此相愛的心,眾人因此就認出你們是我的門徒了.」

反之,如果我們有彼此相愛的心,就向世人表示,這是神的家,屬靈的家.使徒約翰在他的書信裡面,也好多次告訴我們,要有彼此相愛 的心,如果我們彼此相愛,神就在我們中間.

2. 彼此和睦:家裡若常爭吵,兄弟姐妹就會怕呆在家裡,反之,家裡若充滿和睦氣氛,大家都喜歡回家.

戴維在詩篇133:1 說:「看哪,弟兄和睦同居,是何等地善,何等地美!」

在詩篇34:12-14又說:「有何人喜好存活,愛慕長壽,得享美福,就要禁止舌頭不出惡言,嘴唇不 說詭詐的話.要離惡行善,尋求和睦,一心追趕.」

但是和睦并不容易得到,你想和睦,有的人不願跟你和睦,所以使徒保也說:「若是能行,總要盡力與眾人和睦 (羅12:18)」 .

1．為甚麼要相愛(3:13-16)

「我賜給你們一條新命令,乃是叫你們彼此相愛；我怎樣愛你們,你們也要怎樣(彼此)相愛.你們若有彼此相愛的心,眾人因此就認出你們是我的門徒了. 」

「你們若遵守我的命令,就常在我的愛裡 and \#8230; and \#8230;這些事 and \#8230; and \#8230;是要叫我的喜樂存在你們心裡......要彼此相愛,像我愛你們一樣,這就是我的命令.」 「我這樣吩咐你們,是要叫你們彼此相愛.」(約翰福音13章34-35節,15章10-12,17節)

Ａ．愛與恨是信徒跟世人的分別(3:13-15)

「沒有愛心的,仍住在死中」──這「死」當然不是指肉身之死,乃是指靈性的死.沒有愛心的人就是靈性死的世人.信徒必有愛心,否則仍未離開他舊日的「死行」(來6:1),

林前 13章 1-13節

我若能說萬人的方言,並天使的話語,卻沒有愛,我就成了鳴的鑼,響的鈸一般. 我若有先知講道之能,也明白各樣的奧祕,各樣的知識,而且有全備的信,叫我能夠移山,卻沒有愛,我就算不得什麼.我若將所有的賙濟窮人,又捨己身叫人焚燒,卻沒有愛,仍然於我無益.愛是恆久忍耐,又有恩慈；愛是不嫉妒；愛是不自誇,不張狂,不做害羞的事,不求自己的益處,不輕易發怒,不計算人的惡,不喜歡不義,只喜歡真理；凡事包容,凡事相信,凡事盼望,凡事忍耐.愛是永不止息.先知講道之能終必歸於無有；說方言之能終必停止；知識也終必歸於無有.我們現在所知道的有限,先知所講的也有限,等那完全的來到,這有限的必歸於無有了.我作孩子的時候,話語像孩子,心思像孩子,意念像孩子,既成了人,就把孩子的事丟棄了.我們如今彷彿對著鏡子觀看,糢糊不清(原文作:如同猜謎)；到那時就要面對面了.我如今所知道的有限,到那時就全知道,如同主知道我一樣.如今常存的有信,有望,有愛這三樣,其中最大的是愛.



\newpage

\section{守望你的鄰舍}
\label{sec:173}
\textbf{沈祖堯教授 - 守望你的鄰舍}
\newline
\newline
\hyperref[sec:172]{< < < PREV SERMON < < <}
~
\hyperref[sec:toc]{[返目錄]}
~
\hyperref[sec:174]{> > > NEXT SERMON > > >}


利未記 19:1-18
\newline
\begin{longtable}{cl}
\hline
\hline
章節 & 經文 (和合本修訂版)\\
\hline
19:1 & \begin{tabularx}{0.7\textwidth}{X} 耶和華吩咐摩西說： \end{tabularx} \\ \\ \relax
19:2 & \begin{tabularx}{0.7\textwidth}{X} 「你要吩咐以色列全會眾，對他們說：你們要成為聖，因為我耶和華－你們的神是神聖的。 \end{tabularx} \\ \\ \relax
19:3 & \begin{tabularx}{0.7\textwidth}{X} 你們各人都當孝敬父母，也要守我的安息日。我是耶和華－你們的神。 \end{tabularx} \\ \\ \relax
19:4 & \begin{tabularx}{0.7\textwidth}{X} 你們不可轉向虛無的神明，也不可為自己鑄造神像。我是耶和華－你們的神。 \end{tabularx} \\ \\ \relax
19:5 & \begin{tabularx}{0.7\textwidth}{X} 「你們宰殺祭牲獻平安祭給耶和華的時候，要獻得使你們可蒙悅納。 \end{tabularx} \\ \\ \relax
19:6 & \begin{tabularx}{0.7\textwidth}{X} 這祭物要在獻的當天或第二天吃；若有剩到第三天的，就要用火焚燒。 \end{tabularx} \\ \\ \relax
19:7 & \begin{tabularx}{0.7\textwidth}{X} 第三天若再吃，這祭物是不潔淨的，必不蒙悅納。 \end{tabularx} \\ \\ \relax
19:8 & \begin{tabularx}{0.7\textwidth}{X} 吃的人必擔當自己的罪孽，因為他褻瀆了耶和華的聖物，這人必從百姓中剪除。 \end{tabularx} \\ \\ \relax
19:9 & \begin{tabularx}{0.7\textwidth}{X} 「你們在自己的地收割莊稼時，不可割盡田的角落，也不可拾取莊稼所掉落的。 \end{tabularx} \\ \\ \relax
19:10 & \begin{tabularx}{0.7\textwidth}{X} 不可摘盡葡萄園的葡萄，也不可拾取葡萄園中掉落的葡萄，要把它們留給窮人和寄居的。我是耶和華－你們的神。 \end{tabularx} \\ \\ \relax
19:11 & \begin{tabularx}{0.7\textwidth}{X} 「你們不可偷盜，不可欺騙，也不可彼此說謊。 \end{tabularx} \\ \\ \relax
19:12 & \begin{tabularx}{0.7\textwidth}{X} 不可指著我的名起假誓，褻瀆你神的名。我是耶和華。 \end{tabularx} \\ \\ \relax
19:13 & \begin{tabularx}{0.7\textwidth}{X} 「不可欺壓你的鄰舍，也不可偷盜。雇工的工錢不可在你那裡過夜，留到早晨。 \end{tabularx} \\ \\ \relax
19:14 & \begin{tabularx}{0.7\textwidth}{X} 不可咒罵聾子，也不可將絆腳石放在盲人面前。你要敬畏你的神。我是耶和華。 \end{tabularx} \\ \\ \relax
19:15 & \begin{tabularx}{0.7\textwidth}{X} 「你們審判的時候，不可不公正；不可偏護貧窮人，也不可看重有權勢人的臉，總要公平審判你的鄰舍。 \end{tabularx} \\ \\ \relax
19:16 & \begin{tabularx}{0.7\textwidth}{X} 不可在百姓中到處搬弄是非，不可陷害鄰舍的性命。我是耶和華。 \end{tabularx} \\ \\ \relax
19:17 & \begin{tabularx}{0.7\textwidth}{X} 「不可心裡恨你的弟兄；要指摘你的鄰舍，免得因他承擔罪過。 \end{tabularx} \\ \\ \relax
19:18 & \begin{tabularx}{0.7\textwidth}{X} 不可報仇，也不可埋怨你本國的子民。你要愛鄰如己。我是耶和華。 \end{tabularx} \\ \\
[1ex]
\hline
\hline
\end{longtable}
日期:2012年11月25日

題目:守望你的鄰舍

經文:利未記19:1-18

證道:沈祖堯教授

在我接任校長後不久,牧師就跟我說,校長需要每年到崇基的崇拜講道一次,起初我也有些擔心,因為我並沒有接受過任何神學的訓練 and \#8230;

今次已經是我第三次來到這裡,我發覺每次我來到崇基的禮拜堂,坐在台上這個位置,處於這個莊嚴的風琴與牧師洪亮的聲音之間,會感覺到有一種潔淨的感覺；再聽到詩班優美的歌聲,有時候我們真的需要一些禮拜堂是這樣,尖尖的屋頂,才能聽到聲音有如在屋頂中旋轉,好像神在與我們說話.

所以很高興今日可以再與大家分享我最近從聖經裡的一些啟發,但我要再一次強調,我不懂希伯來文,也不懂拉丁文,所以我讀的聖經和你的都一樣,甚至我要看些更淺白的聖經版本,再想想我們今天生活在這個社會裡,這些說話對我們又有什麼意義呢?

今日與大家分享的題目是 －「守望我們的鄰舍」.大家一定聽過摩西在西乃山上接受耶和華給他的十誡,記載在出埃及記的第20章,當它指出我們不可有別的神及拜偶像之後,第一條對人的誡命中指出:要孝敬父母,使你們的日子能得以長久,然後再講出一系列的誡命,是人與人,我們與身邊的人有關的,大家也一定有聽過這些很短很短的誡命:不可以殺人, 不可以姦淫,不可以偷盜,不可作假見證陷害人,不可貪戀別人的房屋,妻子,僕俾,牛驢,並他一切所有的.

很簡單的五句話,很多時候我自己看過後都會覺得,與我何干?因為我又怎麼會殺人,怎麼會姦淫,怎麼會偷盜呢?所以覺得這幾句說話不太與我有關.但當我再翻開看後面幾章的聖經有更詳細解釋(利未記第19章──差不多再說一次十誡的內容),我便發現原來耶和華對摩西所說的這些話,其實與我們是很有關係的.特別它提到人與人之間的關係,都是用了許多的字眼且提出許多的細節,與我們今天的生活是非常貼切的.

我想讀出這段聖經:先讀中文,然後再讀英文翻譯版本.我們先打開利未記第19:1-18:

「耶和華對摩西說,你曉諭以色列全會眾說,你們要聖潔,因為我耶和華你們的神是聖潔的.你們各人都當孝敬父母．也要守我的安息日．我是耶和華你們的神.你們不可偏向虛無的神,也不可為自己鑄造神像．我是耶和華你們的神.你們獻平安祭給耶和華的時候,要獻得可蒙悅納. 這祭物要在獻的那一天,和第二天喫,若有剩到第三天的,就必用火焚燒.第三天若再喫,這就為可憎惡的,必不蒙悅納．凡喫的人,必擔當他的罪孽,因為他褻瀆了耶和華的聖物,那人必從民中剪除.」

好,我們留意,以下一段是說人與人之間的關係,我們慢些讀:「在你們的地收割莊稼,不可割盡田角,也不可拾取所遺落的.不可摘盡葡萄園的果子,也不可拾取葡萄園所掉的果子,要留給窮人和寄居的．我是耶和華你們的神.你們不可偷盜,不可欺騙,也不可彼此說謊.不可指著我的名起假誓,褻瀆你神的名．我是耶和華.不可欺壓你的鄰舍,也不可搶奪他的物．雇工人的工價,不可在你那裡過夜留到早晨.不可咒罵聾子．也不可將絆腳石放在瞎子面前．只要敬畏你的神．我是耶和華.你們施行審判,不可行不義,不可偏護窮人,也不可重看有勢力的人．只要按著公義審判你的鄰舍.不可在民中往來搬弄是非．也不可與鄰舍為敵,置之於死我是耶和華.不可心裡恨你的弟兄．總要指摘你的鄰舍,免得因他擔罪.不可報仇,也不可埋怨你本國的子民,卻要愛人如己.我是耶和華.」

讀完這段聖經後,我就發覺今天我們所做的事,與這段聖經是很有關係的.一開始它說:你要知道,我是耶和華,我是你們的神,是聖潔的,然後它再指示你們要如何行 and \#8230;因為我是神,所以你都要聖潔.在此給予我第一個教訓是Holiness(聖潔) and \#8230;

我們的聖潔並非只在星期日到禮拜堂唱完一首詩歌,讀了幾段聖經後,出去便成了聖潔,好像洗了澡一樣；我們的聖潔是行在每日的行為上,特別是你怎樣對待你的父母,以及你身邊的人.在講完這番話後,然後它說:你們要這樣做 and \#8230;因為我是耶和華,是聖潔的.

那到底要做些什麼呢?在這裡有很多對我們的提醒,我會把Eugene Peterson的英文翻譯版本,再配合這段中文來與大家分享 and \#8230;

第一,它說:當你耕田的時候,你收割你的莊稼,不要割得太盡.英文版這樣說:

9-10: When you harvest your land, don't harvest right up to the edges of your field or gather the gleanings from the harvest. Don't strip your vineyard bare or go back and pick up the fallen grapes. Leave them for the poor and the foreigner. I am GOD, your God.

它說,如果你有一塊田在耕種, 要收割,完全是正當的買賣收穫,是應得的；但不要割得太盡,要留一些邊不要割,好讓那些沒有東西吃的窮人,外地人靠著你留下的邊得以執拾；當你種葡萄的時候(大家到過很多葡萄園的話,看到那一行一行的葡萄園,大到不得了),你採摘完後,掉在地上的葡萄,請不要執起來,好讓那些沒有東西吃的窮人,外地人執拾得以為生.

這與我們今日生活的關係又是什麼呢?

各位,請你們想一想,不用我說你也知道是涉及我們當中怎樣對待貧窮的人,香港貧富懸殊的問題,我們仍然有很多人住在tong 房,板間房；我們仍然有很多人生活在貧窮線之下.幾年前,我在飛機上看到了無線電視的一集記錄片,片名叫「一百萬人的故事」.是什麼「一百萬人的故事」呢?原來香港還有一百萬人生活在貧窮線以下.

在第一集中,看到香港有很多人在晨早3,4點已要起床,到天橋底下賣一些你會覺得沒可能會賣出去的東西:破衣架,壞收音機,二手衣物... 他們何以要這樣做?因為他們沒有其他可以謀生的方法,他們也許是40,50歲失業的中年人,又或是外地來的人,支取綜援的人,但入不敷支,於是希望以這樣的方法能賺取1元數角,希望可以足夠讓他們到市場買些破爛但便宜的水果,蔬菜,作為糧食(這些景況很多就在我們身邊發生著).還有很多外地來的人 and \#8230;我在網上寫了一段說話,但帶來很多嚴厲的責罵,但我仍然覺得我們可否把內地來的移民是外來人呢?我們又該怎樣來對待他們呢?我們不但有來自國內的移民,也有很多南亞裔的移民,我曾經去過一間學校,全校都是印度,巴勒斯坦藉的學生,那裡的校長是我以前中學生物學的老師,他退休後到那間學校當校長,他說,他那間學校巳經是最好的了,名叫嘉多理..中學(我忘記了校名).但他說,在香港,有很多在尼泊爾,印度,巴勒斯坦來的人,他們可能已經來了香港很久,甚至生了孩子,但他們的孩子永遠也沒辦法讀到香港正規的大學,因為他們的孩子永遠也沒辦法得以學習聽,學習說,學習寫中文的,香港的學制沒有渠道讓他們就讀,因此,他們絕大部份只讀完中學後,便在重慶大廈賣東西,或在地盆當扎鐵,剷泥,幹粗活.

我們怎樣對待我們身邊的移民和難民呢?這處聖經說:當你有能力的時候,當你賺到錢甚至有過剩的時候,各位,請不要取得太盡,不要賺得太盡,不要不留餘地,不單是施捨,而是讓有需要的人可以執拾你過剩的東西,留下一些給他們.做生意並沒有錯,賺錢也不是錯,但這裡總提醒我們,不要取得太盡,要留有餘地.

第二,聖經說:不要說謊,不要偷盜,不要欺騙.馬克吐溫說:一個不說謊的人,不需要太好的記憶力,因為當你說了謊後,要記著以前說了什麼謊話,於是便需要說更多的謊話去遮掩之前的謊話,要想出一個謊話蓋過另一個謊話... 所以Mark Twain(馬克吐溫說):一個不說謊的人,不需要太好的記憶力.但你可能會說,我也不怎麼說謊,就算說謊也是無傷大雅的,我們那有什麼嚴重的謊話?各位,有些謊話真的是有嚴重後果的?你有沒有看過一套電影叫 Inside Job?內容是說在2009年的時候,由雷曼兄弟銀行倒閉這事開始,全球經濟衰退,何會以引至全世界經濟也陷入今天的困境呢?就是因為一些很有頭腦的人說了謊,如果沒才幹的人說了謊,破壞性不會那麼大,但很有本事的人製造了所謂衍生工具,或叫derivative(我也是看完這套電影才懂一點),是一些複雜得你根本不會明白的東西,看似買了之後會穩如泰山,不會虧損；但結果原來只是一個空殼,結局是將你的所有陷入一個萬劫不復的景況 and \#8230;

有才能的人, 讀過大學的人,甚至讀過博士的人,名校出身的人,有很多在goldman sachs, Lehman Brothers工作的人,都是最了不起的,Harvard Business School ,MIT,我不是刻意提這兩間學校,但來自最了不起的商學院,他們更懂得說一些對自己非常有利的說話,但可能害慘了全世界的慌言:

不可偷盜,不可欺騙.但可能你會說,我沒有他們那麼了不起!我只處身香港,也害不到全世界 and \#8230;.

過去的一個星期,每日的報紙也在報導DR. 事件(美容院醫生注射美容汁後出事),有些人在香港的人也在談論.如果我把你的血抽出來,洗一洗,再加一些不知名的東西然後再輸回入你的體內,你便會長命百歲,會更美麗,會更健康 and \#8230;.但最後造成了嚴重的後果,甚至喪失了生命...

同樣讓我看見:越讀書多的人,越有學位,越有專業資格的人,如果他要說謊,如果他沒有道德操守,如果他沒有看過這段的聖經,他所製造出來的危險及破壞,比一個未受過教育的人更嚴重.所以,不可偷盜,不可欺騙,不可說謊 and \#8230;其實不單只是說我們個人的操守會影響身邊兩三個人,可能會影響社會更多的人 and \#8230;

跟著,聖經一路說下去:不要欺壓你的鄰舍,不要搶別人的東西；雇工的人的工價,你不要留下過夜,不可留到隔天才發給人.即是說什麼呢?即是說,做老闆的,不要欠薪,不要扣工資,不要算盡得以付出最低的工資...機關算盡,以至自己能搾盡最大利益,搾乾工人.這裡說:不要利用工人,不要搶他們的東西,do not hold back the wages of hire hand overnight,過一個晚上也不能.聖經也說,牛在田裡耕地,它要吃的草你要給它吃,要渴的水你要讓它喝；如果我們當老闆的(不知當中有多少是老闆?)我們該怎樣去對待我們的工人呢?在這裡是有教導的:他也是你的鄰舍,也是你身邊的人.不可搶奪他的東西留到明日.還有,聖經說:不可欺負聾人,盲人；不可在路上放石頭,看他們會否被絆倒(不可沒有他們可用的廁所,不可沒有他們不能搭的電梯)；因為他是聾子,我不聘用他,或因為他有精神病的,憂鬱的,不可以歧視他們.因為聖經說:你把絆腳的石放在盲人面之前,你要想想我是耶和華,是你要敬畏的神!

我們在醫學院裡,曾經收了一個坐輪椅的同學,我們收錄了才發覺,要讓他從宿舍坐著輪椅走到醫院裡,要建一個斜坡,讓輪椅能走過,過馬路不用走樓梯.用了40萬!跟著要重建電梯,特別廁所,更發覺病床間的位置也太窄,容不下輪椅在病人床邊!所以差不多連病房也需改裝,還有手術室,要怎讓他能看到躺在病床上的病人?我們也要特別給他訂做合適的床!收了一個學生,我們花了數百萬改裝整間醫院!不過我覺得是值得的,因為我們是給與一個機會,讓一個本來不能讀醫的人,成為一個醫生.今年,他剛畢業,我們非常高興.雖然我們花了數百萬,但我相信這背後是非常有意義的,不單是用在一個學生的身上,可能將來可以用在有需要的學生身上,但一個更重要的意義是:我們不會因為你是盲人,不會因為你曾經患精神病,因為你腳不能走,而不給予同樣的機會.

今日,其實還有很多國家,我最近去到一個國家,我不說出國家名字,他們說他們國家裡會註明這份工作是不請女性的,我很愕然,因為香港有平等機會委員會,但他們說他們那裡沒有相關機構的,所以女人在他的國家裡,讀完大學後,找工作的機會比男人低得很.最近,我老師告訴我,她是一個女教授,我才知道在香港60年代的時候,男人與女人打同樣的一份工,但工資是大不相同的,她當時當醫生的人工只是男性的三份之一,而且假期福利也大不同.後來她與其他幾個女人一起爭取權益.對不起,我不是把女性與瞎子相提並論,只是想說出我們不可以有歧視,特別對有需要的殘疾人士.聖經說不可這樣行,因為這也是一個欺壓你鄰舍的行為.

還有什麼與我們有關呢?

聖經說:當施行審判的時候,對窮人或富人也要一致,不可偏護窮人也不可看重有勢力的人.這句說話, 在Eugene Peterson的英文譯本裡有小小不同的.這裡說你不要幫有錢人欺負窮人,應給予窮人一樣的審判結果,不要理會他有多少錢.好像有些說不要幫助有錢人,但在Eugene Peterson的譯本裡,其實是說 ... be fair to the poor and the rich...,其實我們也不應該欺負有錢人,也不應該欺負窮人,因為我們仇富的情緒會使我們欺負有錢人也說不定,可能那些有錢人無心地做了一些事,但由於我們因仇視他們,而過份地針對他們了.be fair to(公平對)窮人也be fair to有錢人.因此這裡說,當你審判的時候要本著公義,不可與之為敵,不可報仇,不可暗地裡恨他們,do not take revenge, do not secretly hate your neighbor,不可以眼還眼,不可以言語,以facebook,以任何媒體的方法「砌他」,欺負他.你可能會說,我這麼窮,怎麼欺負有錢人呢?但其實你可以的,現在我們的話很具殺傷力,所以我們要小心,因為上帝是窮人的上帝也是有錢人的上帝,叫太陽照好人,也照歹人.所以我們應該to be fair,在英文裡譯得非常清楚,不要stand by when your neighbor is in danger,當你鄰舍有危險的時候不要隔岸觀火,袖手旁觀.可能大家也記得,早前大陸有則新聞:有個電視鏡頭拍到,有個小女孩過馬路時被車撞倒,跟著很多的過路人也袖手旁觀,當沒事發生；很多車輛經過都視而不見,最後只有一位執垃圾的婆婆救起小孩,送去醫院,但小孩最後也救不了 and \#8230;

不可視若無睹,不可隔岸觀火.可能這些情況在媒體及facebook 的評論是說大陸是這樣的了,因為他們是冷血的.但我相信這情況並非只發生在大陸,不要讓這些事情發生在香港,不要讓這些事情發生在我們身邊,在你走過的地方.這個,也是你應該對待鄰舍的方法.

各位,我不想再說下去,但我看完這聖經後,我發覺這段一連串的教訓裡,其實是在教導我們一件事:我們不要做這樣,那樣做的背後原因,是你要有包容,有寬恕,你要有愛,你才能做得出那些.有些人實在令人討厭,或許他所說的話,或許他的意見,你並不同意,但你能否如孔子說:君子和而不同呢?能否有accommodation 呢?如果沒有包容,我們又怎會有言論自由呢?如果沒有包容,我們又怎能讓別人也有其權利呢?

昨晚在facebook上留存了這樣的一樣東西給我看:「中大大陸學生辱罵香港人」.我不知你有沒有看,昨晚瘋傳了一輪.因為曾經有一個在中大讀過的學生,也不知他是那一個系,畢業了沒有,他說香港人長得醜,論思考,說話也不經大腦,祟洋媚外,香港女孩不及他那裡的漂亮 and \#8230;總之他有幾篇東西在facebook瘋傳.跟著我也心知不妙,「中大大陸學生辱罵香港人」 and \#8230;接著我今天早上開電腦,已收到三個香港市民寫給我的電郵,叫我好好管教我的學生:作為大學校長,可否懲罰這些人,星加坡有個學生,罵星加坡人,就立即給人退學,你的大學又如何呢?多快速!如果昨晚我沒有看facebook,今早沒開電腦,我也不知道發生了什麼事!但已經有人把這件事放在我面前了 and \#8230;跟著我寫了一個電郵給我們的Dean of student(學生事務處長──他是我最要好的朋友又是我的擋箭牌),我對他說,請你查一查究竟發生了什麼事,這個是否我們學生,我們是否需要做些什麼?幸好,在我發出電郵發約一個小時之後,我收到一個電郵說這位同學又發出了一封電郵向全香港人道歉,他說昨晚的事他覺得有些過激的地方以及胡亂批評別人樣貌是不對的,對於香港人有極大的不尊重,他並非針對所有香港人,不過一時衝動,在此希望兩地團結,而大多數的香港人他是非常尊重的... 等等... (他終於寫了一篇這樣的文章).那我想我明天上班也不用那麼不知所措了,我真的可以在一日之內收到很多的要求,要求我們這樣做,那樣做,更連我也一起罵上了.

各位,這代表什麼呢?代表我們真的只有很小的包容.

剛才牧師說,今天的社會充滿了戾氣,充滿了互相指責,充滿了憤怒,只要一些雞毛蒜皮的事,看見別人在地鐵上吃東西,我便要立即拍下來,然後放上網,一起咒罵；只要有人在火車站擠迫到我,我便會立即還擊... 各位,我們想想,我們可否對我們身邊的人好一些呢?特別是對那些貧窮的,身體有缺憾的,軟弱的,或缺乏機會的.這並不是倫理道德課,今天我們說的是十誡裡上帝要我們做的,因為祂說:我是耶和華,我是聖潔的,所以你都要聖潔.聖潔的定義不代表你不去偷,去搶,你還要對身邊的人更好.還記得,在福音書裡有個比喻:有人死了後到了天堂,耶穌將他們分開兩邊,一邊是天堂,一邊地獄,耶穌對地獄的人說:你們何以在這裡呢?因為在我沒有東西吃的時候,你沒給我吃,在我沒衣服穿的時候,你沒有給我衣服穿,在我沒地方住的時候,你沒給我住.那些人就說,我何時遇見你沒東西吃,沒衣服穿,沒地方住呢?!耶穌說:那次你的鄰舍,你的弟兄中有這樣的需要,但你沒有伸出援手,這等如沒做在我的身上；跟著轉身向另一批人說:你們在天堂,是因為我沒東西吃的時候,你給我吃,我赤身露體的時候你給我衣服穿,我沒地方住的時候你給我住,這些做在微小的弟兄身上便如做在我的身上了.

對不起,我不是確切知道該段聖經實在說什麼,因為我真的沒讀過神學,但那句話的精髓是:做在你身邊最有需要的人身上,就是做在上帝的身上了!

Mother Theresa 說:我只是耶穌手上的一支小鉛筆,讓我不斷可以磨去自己,寫出上帝的愛是什麼.她所做的只是在街上收容那些快要死去的人,病到沒法醫治的人,或極其窮因的人,睡在溝渠邊沒飯吃的人 and \#8230;只是收容他們到她教堂裡,清潔他們的身體,給他們水喝,或許數小時侯他們就死去,但他們說這些人在短短的數小時改變很大,本來他們連動物也不如的身份,but they died like angels,死的時候好像天使一樣!也許是微小的東西,但就是可以做在我們鄰舍的身上!

希望今天這段說話可以讓我們每天反省一下,在我們在校園,在你工作的地方,甚至在家庭,在社交圈子裡,怎樣去關顧你的鄰舍呢?



\newpage

\section{不正義制度下的窮寡婦}
\label{sec:174}
\textbf{鄧瑞強博士 - 不正義制度下的窮寡婦}
\newline
\newline
\hyperref[sec:173]{< < < PREV SERMON < < <}
~
\hyperref[sec:toc]{[返目錄]}
~
\hyperref[sec:175]{> > > NEXT SERMON > > >}


馬可福音 12:38-44
\newline
\begin{longtable}{cl}
\hline
\hline
章節 & 經文 (和合本修訂版)\\
\hline
12:38 & \begin{tabularx}{0.7\textwidth}{X} 他在教導的時候，說：「你們要防備文士。他們好穿長袍走來走去，喜歡人們在街市上向他們問安， \end{tabularx} \\ \\ \relax
12:39 & \begin{tabularx}{0.7\textwidth}{X} 又喜愛會堂裡的高位，宴席上的首座。 \end{tabularx} \\ \\ \relax
12:40 & \begin{tabularx}{0.7\textwidth}{X} 他們侵吞寡婦的家產，假意作很長的禱告。這些人要受更重的懲罰！」 \end{tabularx} \\ \\ \relax
12:41 & \begin{tabularx}{0.7\textwidth}{X} 耶穌面向聖殿銀庫坐著，看眾人怎樣把錢投入銀庫。有好些財主投了許多錢。 \end{tabularx} \\ \\ \relax
12:42 & \begin{tabularx}{0.7\textwidth}{X} 有一個窮寡婦來，投了兩個小文錢，就是一個大文錢。 \end{tabularx} \\ \\ \relax
12:43 & \begin{tabularx}{0.7\textwidth}{X} 耶穌叫門徒來，對他們說：「我實在告訴你們，這窮寡婦投入銀庫裡的比眾人所投的更多。 \end{tabularx} \\ \\ \relax
12:44 & \begin{tabularx}{0.7\textwidth}{X} 因為，眾人都是拿有餘的捐獻，但這寡婦，雖然自己不足，卻把她一生所有的全都投進去了。」 \end{tabularx} \\ \\
[1ex]
\hline
\hline
\end{longtable}
日期:2012年12月2日

主題:不正義制度下的窮寡婦

證道:鄧瑞強博士

經文:馬可福音12:38-44

各位弟兄姊妹,主內平安.

大城巿裡面,有很多小故事.

某年某月,在一個被評為最適宜居住的大城巿裡,那個晚上,天氣寒冷,明哥回到深水Po 玉石巿場附近的天橋底.冷風吹來,只覺這個城巿異常冷漠.明哥有的,只是一張舖在地上的床褥,一張薄被,一兩袋雜物.在這個富裕的城巿裡,這裡是最不能沾上「富裕」這兩個字的地方.在豪宅璀璨的燈光下,這裡顯得異常黑暗.明哥是寂寞的人,但也不算寂寞,這裡有十多人,這晚一起露宿而睡.

夜了,明哥卷在薄被裡,矇矇矓矓地睡著了.或許,他夢見自己睡在有四堵牆的屋內.雖然不能說是好夢正酣,畢竟,或是睡熟了.忽然,有食環署的人來,也有警察來,叫醒他,將他的床褥,拖鞋,他僅有的一兩袋雜物,全當垃圾被扔到車上.他想取回自己僅有的東西,至少,取回自己的拖鞋,但是,全被阻止.那些人員說,現在是「洗街」,要將一切棄置街頭的垃圾清理掉.那些東西怎能算是垃圾,明哥心裡大叫,這就是他一切養生的,這就是他的所有.

沒有了,什麼也沒有了.明哥連同其他露宿者,透過小額錢債審裁處,向食環署追討賠償.有關官員卻一拖再拖.最後,當他追討的賠償還未有著落的時候,明哥患病離世了.當他死的時候,我想,他可能會問:「為何這個富裕的社會,富裕得容不下一點貧窮?」

這就是我們城巿的小故事.

為了令富裕的城巿顯得更富裕,一個一無所有的人,變得更加一無所有.

主耶穌講的道理,很多時透過城巿的小故事來表達.

今天,我們看看主耶穌見到的小故事.

講道

經文:馬可福音12:38－44

38     耶穌在教訓之間,說,你們要防備文士,他們好穿長衣遊行,喜愛人在街市上問他們的安,     39     又喜愛會堂裡的高位筵席上的首座．     40     他們侵吞寡婦的家產,假意作很長的禱告．這些人要受更重的刑罰.     41     耶穌對銀庫坐著,看眾人怎樣投錢入庫．有好些財主,往裡投了若干的錢.     42     有一個窮寡婦來,往裡投了兩個小錢,就是一個大錢.     43     耶穌叫門徒來,說,我實在告訴你們,這窮寡婦投入庫裡的,比眾人所投的更多．     44     因為他們都是自己有餘,拿出來投在裡頭．但這寡婦是自己不足,把他一切養生的都投上了.

有一次,主耶穌對他的城巿的現象,作了一番評論.

他說,你們要防備文士.文士,是律法的教師,是社會規範的制定者,是真理的代言人.你們看看他們.「他們好穿長衣遊行,喜愛人在街市上問他們的安,又喜愛會堂裡的高位,筵席上的首座.」他們愛好穿著顯示他們高貴身分的服飾,招搖過巿,熱衷於人家的奉承,愛好成為社會上的名流,喜歡在社交場合成為焦點人物.這些都是文士表面的生活,若他們只是如此,對他人也沒有什麼害處.但是,這些不只是一種表面現象.外表反映其內心.低俗的表面反映低俗的內心.他們的內心如何?主耶穌說:「他們侵吞寡婦的家產,假意做很長的禱告.這些人要受更重的刑罰!」他們覬覦寡婦的家產,他們剝削社會上最弱者的資源,他們對無助者落井下石.他們作長篇大論的祈禱,作長篇大論的演說,講出一大段表面愛護弱者實質剝削弱者的謊言.神判定他們的罪.

或許,我們會說,主耶穌講的,有何證據?

主耶穌便指出一個社會上的小故事,作為他立論的證據.

41     耶穌對銀庫坐著,看眾人怎樣投錢入庫．有好些財主,往裡投了若干的錢.     42     有一個窮寡婦來,往裡投了兩個小錢,就是一個大錢.

主耶穌坐著,觀察這個社會的運作.

在巴黎,有很多露天茶座.很多人坐下來,飲杯咖啡,看著走過的人群,細味城巿的底蘊.主耶穌坐著,細味當時社會的生活.

主耶穌在哪裡坐著?他不是坐在咖啡廳裡.他對著銀庫坐著.「銀庫」,很多英文聖經翻譯為 ``treasury".「銀庫」,涉及金錢,涉及人對金錢的態度,涉及這個社會制度如何處理人的金錢運用.主耶穌在這裡坐著,要看透這個社會在涉及金錢的事上對人有何要求?主耶穌對著的銀庫,是聖殿的銀庫,他在這裡坐著,要看看人以神聖的名義如何處理金錢事宜.主耶穌在這裡坐著,透過金錢與神聖,剖析人性.

主耶穌看到什麼?

他看見很多有錢人,投入了很多錢.為此,我們應該感激他們,也為他們感謝神.他們的付出,造福社群.

主耶穌看見一個「窮寡婦」.寡婦,已經是社會上的弱者.窮寡婦,是弱者中的至弱者.她投下了兩個小錢,就是一個大錢.「小錢」是當時流通貨幣裡的最小單位.所謂「大錢」,也只是很少的錢.讀到這裡,我記起我小時候返教會,崇拜時奉獻,因為貧窮,只能投下幾個毫子.幾個「毫子」,加起來就是一個「五毫子」.但是「五毫子」,也只是很少的錢.為何不投入一個「五毫子」?因為貧窮的人,在日常生活裡,只接觸「毫子」,連接觸「五毫子」的機會也很少.這個「窮寡婦」真的很窮.

面對這個情景,主耶穌作了如下的評論.

43     耶穌叫門徒來,說,我實在告訴你們,這窮寡婦投入庫裡的,比眾人所投的更多．     44     因為他們都是自己有餘,拿出來投在裡頭．但這寡婦是自己不足,把他一切養生的都投上了.

主耶穌說,這個「窮寡婦」付出的,比所有人加起來還多.按 Eugene Peterson 的 ``The Message" 的聖經翻譯,第44節是這樣說的:「其他人付出的,是他們不顧惜的；她的付出,是過份的,是她根本不能負擔的,她付上了她的一切.」( ``All the others gave what they"ll never miss; she gave extravagantly what she couldn"t afford -- she gave her all."- ``The Message")一個一無所有的人,變得更加一無所有.

主耶穌的評論,究竟有何所指?

首先,主耶穌是要指責當時的宗教制度,這種宗教制度也規範了社會的日常運作.他要指責這些制度的制定人.記得嗎?主耶穌曾說,那些文士,那些律法的解釋人,「他們侵吞寡婦的家產」.他們是明刀明槍地去搶「寡婦的家產」嗎?當然不會.這些社會上有頭有面的人,怎會作這明顯的壞事.他們透過「制度」－－神聖的制度－－去完成這件事.主耶穌坐著,觀察這「銀庫」運作,觀察這「銀庫」制度對一個「窮寡婦」的要求.然後,主耶穌說,你們看,這個女人,她窮得可憐,她一切都不足,她卻奉上了她一切所有,她奉上她賴以維生的一切.她將如何生活下去呢?

舊約申命記(如:申14:29)指出,人向神奉獻,其中一部分是要用來養活孤兒寡婦的.但是,現在,整個制度卻是讓貧者越貧,使這些貧窮者賴以維生的一切也被奪走.聖殿不是讓窮人在當中找到食物,讓他們空手而來,快樂地離去的嗎?是什麼改變了這一切?是什麼將窮人僅有的東西奪走?文士制定聖殿的制度,他們以神聖的理由規定這寡婦要作的奉獻.他們的生計依賴這些奉獻.神聖的制度保護了他們的收益.他們沒有想過,當那寡婦連她賴以存活的一切都奉上時,她將如何生活下去呢?

這些制度的設立,固然有其神聖的理由,固然有其宏大敘事,但是,主耶穌看到的,是這大制度下的小人物.他告訴我們,這個「窮寡婦」的故事,一個無人留意的淒慘故事.看,這個小人物,她告訴我們,這個「銀庫」制度出現了問題,這個使人找到生命寄托的聖殿出現了問題,那些長長的禱告出現了問題.

主耶穌的評論,不單對當時的宗教領袖說話,不單針對當時的「銀庫」制度,也針對跟從他的門徒.

跟從耶穌,到底想要什麼?

按馬可福音的記載,門徒跟從耶穌,要想得到今世的榮華富貴.

這是一件令人悲哀的事件,很多人跟從耶穌,只是想要今世的榮華富貴.

在馬可福音裡,主耶穌向門徒預告了三次他將要受死的訊息.第一次是在8:31-38,主耶穌的大弟子彼得即時的反應是「責備」耶穌.可能,彼得心想:「主,去受死,你是否傻了?」然後,主耶穌教訓他的門徒:「若有人要跟從我,就當捨己,背起他的十字架來跟從我.」第二次耶穌預告他的受死是在9:30-37,門徒的回應卻是,爭論誰為大.第三次耶穌預告他的受死是在10:32-45,主耶穌另外兩個大弟子雅各和約翰的回應,是求耶穌「賜我們在你的榮耀裡,一個坐在你右邊,一個坐在你左邊」.主耶穌的回應是:「人子來,並不是要受人的服事,乃是要服事人,並且要捨命作多人的贖價.」生命的神聖意義,在於尊重他人,愛護他人,甚至為了愛他人,放下自己的一切.

主耶穌的門徒明白嗎?他們做得到嗎?

一個人只顧自己,則做不成主的門徒.

主耶穌對著「銀庫」坐著,看見這窮寡婦,無論當中的制度合不合理,她因著神的名字,她默默地奉上了一切.經文說:「把她一切養生的都投上了」,這有一種「全然捨己」的意味,她甚至沒有顧念之後如何活下去的問題.

作為主的門徒,聽到這個窮寡婦的故事,會有什麼反應呢?或許,他們仍然會覺得,一個甘願捨己的人生抉擇是傻的.他們記掛的,仍是「誰最大」?誰能在今世的榮耀中得到最高的位置?

主耶穌坐著,說:「這窮寡婦投入庫裡的,比眾人所投的更多」.這是一個價值判斷.主耶穌教導他的門徒,什麼是「多」,什麼是「少」?什麼是「得」,什麼是「失」?什麼是「成」,什麼是「敗」?什麼是「豐盛」,什麼是「空洞」?什麼是「真情」,什麼是「假意」?什麼是「真信仰」,什麼是「假敬虔」?

很多人在這「銀庫」的故事面前,憂憂愁愁地離開.

這個「窮寡婦」的故事,預示了主耶穌的犧牲,因為講完這故事後,主耶穌便走向他的十架之旅了.

「窮寡婦」的故事,與主的十架故事一樣,一方面批判罪惡的人與制度,另方面卻指示真正的信仰人生.

今天,主耶穌對著我們的城巿的「銀庫」坐著,他要看看這個金錢世界是如何運作的,他要看看在金錢的事上我們活出了什麼樣的人性.

留心看看,這個在不正義制度下的窮寡婦.

但願:榮耀歸於聖父,聖子,聖靈.阿們.



\newpage

\section{爾國爾民}
\label{sec:175}
\textbf{林豪恩先生 - 爾國爾民}
\newline
\newline
\hyperref[sec:174]{< < < PREV SERMON < < <}
~
\hyperref[sec:toc]{[返目錄]}
~
\hyperref[sec:176]{> > > NEXT SERMON > > >}


約翰福音 18:33-37
\newline
\begin{longtable}{cl}
\hline
\hline
章節 & 經文 (和合本修訂版)\\
\hline
18:33 & \begin{tabularx}{0.7\textwidth}{X} 於是彼拉多又進了總督府，叫耶穌來，對他說：「你是猶太人的王嗎？」 \end{tabularx} \\ \\ \relax
18:34 & \begin{tabularx}{0.7\textwidth}{X} 耶穌回答：「這話是你說的，還是別人論到我時對你說的呢？」 \end{tabularx} \\ \\ \relax
18:35 & \begin{tabularx}{0.7\textwidth}{X} 彼拉多回答：「難道我是猶太人嗎？你的同胞和祭司長把你交給我。你做了甚麼事呢？」 \end{tabularx} \\ \\ \relax
18:36 & \begin{tabularx}{0.7\textwidth}{X} 耶穌回答：「我的國不屬於這世界；我的國若屬於這世界，我的部下就會為我戰鬥，使我不至於被交給猶太人。只是我的國不屬於這世界。」 \end{tabularx} \\ \\ \relax
18:37 & \begin{tabularx}{0.7\textwidth}{X} 於是彼拉多對他說：「那麼，你是王了？」耶穌回答：「是你說我是王。我為此而生，也為此來到世界，為了給真理作見證。凡屬真理的人都聽我的話。」 \end{tabularx} \\ \\
[1ex]
\hline
\hline
\end{longtable}
啟示錄 1:4-8
\newline
\begin{longtable}{cl}
\hline
\hline
章節 & 經文 (和合本修訂版)\\
\hline
1:4 & \begin{tabularx}{0.7\textwidth}{X} 約翰寫信給亞細亞的七個教會。願那位今在、昔在、以後永在的神，與他寶座前的七靈， \end{tabularx} \\ \\ \relax
1:5 & \begin{tabularx}{0.7\textwidth}{X} 和那忠信的見證者、從死人中復活的首生者、世上君王的元首耶穌基督，賜恩惠和平安給你們。他愛我們，用自己的血使我們從罪中得釋放 ， \end{tabularx} \\ \\ \relax
1:6 & \begin{tabularx}{0.7\textwidth}{X} 又使我們成為國度，作他父神的祭司。願榮耀、權能歸給他，直到永永遠遠。阿們！ \end{tabularx} \\ \\ \relax
1:7 & \begin{tabularx}{0.7\textwidth}{X} 「看哪，他駕雲降臨；眾目都要看見他，連刺他的人也要看見他；地上的萬族要因他哀哭。」這是真實的。阿們！ \end{tabularx} \\ \\ \relax
1:8 & \begin{tabularx}{0.7\textwidth}{X} 主神說：「我是阿拉法，我是俄梅戛 ，是今在、昔在、以後永在的全能者。」 \end{tabularx} \\ \\
[1ex]
\hline
\hline
\end{longtable}
日期:2012年12月9日

題目:爾國爾民

證道:林豪恩先生

經文:約翰福音18:33-37節;啟示錄1:4-8節

一．引言有人說,機場是一個經常有特別人士或者特別事件出現的地方,你有否在機場遇到這些特別的人士或特別的事情呢?

去年聖誕節期間,有一名婦女帶著兩個女兒從台灣出發乘飛機往英國,在香港機場轉機時被拒絕,強制送回台灣又無法入境,於是母女三人在香港機場滯留了十一日.

為甚麼這母女三人變成無地方收留的「人球」呢?原來這位婦女在台灣和英國籍的前夫打離婚官司,官司贏了,卻一直拿不到贍養費,對台灣司法灰心,就帶著2個女兒放棄中華民國國籍.當母女3人搭機到香港,準備轉機去英國找前夫,以為只要找到前夫,就可能拿到英國國籍,卻因為當下沒有任何國籍被強制回台,因為母女3人已經沒有中華民國國籍,在台也無法入境.(TVBS 20/12/2011)

世界那麼大,怎會無處容身呢?這樣奇怪的事情自古至今常有發生.當年,耶穌被押解到羅馬帝國派來的總督彼拉多面前,彼拉多問他:「你的同胞放棄你,把你交給我我這位殖民地官員處置.你到底做了甚麼事,連你自己的人也驅逐你呢?」其實,彼拉多似乎也不想處理這位被猶太人放棄的猶太人,於是向猶太人說:「我查不出他有甚麼罪來」,意圖把這個耶穌交回給猶太人.在此,耶穌彷彿成為了沒有歸屬的「人球」,在猶太人和彼拉多之間被「踢來踢去」.到底耶穌的歸屬在哪裡呢?他屬於哪一國呢?哪一國可讓他安身呢?我們又屬於哪裡的呢?哪裡可以讓我們安身呢?屬於哪裡對我們來說又有甚麼意義呢?

二．經文

18 : 33於是 彼拉多又進了總督府,叫耶穌來,對他說:「你是 猶太人的王嗎?」 18 : 34耶穌回答:「這話是你說的,還是別人論到我時對你說的呢?」 18 : 35彼拉多回答:「難道我是 猶太人嗎?你的同胞和祭司長把你交給我.你做了甚麼事呢?」 18 : 36耶穌回答:「我的國不屬於這世界；我的國若屬於這世界,我的部下就會為我戰鬥,使我不至於被交給 猶太人.只是我的國不屬於這世界.」 18 : 37於是 彼拉多對他說:「那麼,你是王了?」耶穌回答:「是你說我是王.我為此而生,也為此來到世界,為了給真理作見證.凡屬真理的人都聽我的話.」

三．釋經與應用

1．耶穌的國不屬於這個世界

從經文所見,一方面,耶穌見棄於世界,但同時,耶穌卻也宣稱「我的國不屬於這世界」.為甚麼耶穌說我的國不屬於這個世界呢?這個世界的國度有甚麼特色是與耶穌的國度不能相容的呢?從耶穌與彼拉多這段簡短的對話中,我們可以看見下列兩項.

第一,這世界的國度經常以民族及種族作為劃分.

18 : 33於是 彼拉多又進了總督府,叫耶穌來,對他說:「你是 猶太人的王嗎?」 18 : 34耶穌回答:「這話是你說的,還是別人論到我時對你說的呢?」 18 : 35彼拉多回答:「難道我是 猶太人嗎?你的同胞和祭司長把你交給我.你做了甚麼事呢?」

當彼拉多問耶穌是否猶太人的王的時候,耶穌的回答似乎想邀請他參與討論,詢問他的個人意見.然而,彼拉多卻拒絕邀請,拒絕進入耶穌的議題.彼拉多反問:「難道我是猶太人嗎?」他的意思是,你是猶太人的王這傳聞是你們猶太人自己的事情,我不是猶太人,因此,這事情與我無關.現在,這事情來到我面前,只因你們猶太人把你交給我而已,請你不要誤會以為我對你有興趣,以為我與你有甚麼關係.無論人們說你是甚麼身份也好,歸根究底,你是猶太人,我是羅馬官員,你們的事情與我無關.」

世界的國度常以民族作為定義,以民族為主體,這以乎是不證自明的道理.個人出生成為某個民族的一份子,對所屬民族產生感情,維持所屬民族的生存和利益,為所屬民族爭取榮譽,似乎也是天經地義的事情.至於其他民族,若與我們沒有利益衝突,無論他們打生打死,或者饑荒年年,這都是他們的「內政」,我們不會干涉別國的內政；若與我們有利益衝突,則把對方整個民族也視為敵人,輕則敵視,重則滅族,約七十年前納粹德國屠殺猶太人是一個近代的例子,約二十年前胡圖族人在三個月裡殺死100萬圖西族人是一個更近代的例子.

世界的國對非我族類的態度,有時連跟隨耶穌的門徒也難以擺脫.

有一次,耶穌經過撒馬利亞,和一個撒馬利亞婦人談話.耶穌的門徒看見,都為之而驚訝,他們心裡面想:「你為甚麼和他說話呢?」雖然大家是鄰居,但猶太人和撒馬利亞人並沒有來往.(約翰福音4章)門徒認為,作為猶太人的耶穌和撒馬利亞人不應該有甚麼關連的.

耶穌死而復活後,有一次,彼得得到異象去向外邦人哥利流傳福音,回到耶路撒冷後,立即被指控:「你進入未受割禮的人家和他們一同吃飯了.」他們認為耶穌是猶太人,耶穌的福音不應該和外邦人有甚麼關連的.曾幾何時,我們中國人也認為耶穌與我們無關,因為他是紅鬚綠眼的西方人,基督教是洋教,甚至有人慨嘆:「多一個基督徒,便少一個中國人」.

耶穌回答彼拉多說:「我的國不屬於這個世界」.如果耶穌的國屬於這個世界,就會被這個世界的民族觀念和種族觀念所限制.

第二,這世界的國以武力來維持,以戰鬥回應問題18 : 36耶穌回答:「我的國不屬於這世界；我的國若屬於這世界,我的部下就會為我戰鬥,使我不至於被交給 猶太人.只是我的國不屬於這世界.」

在彼拉多看來,耶穌被稱為猶太人的王簡真是一個笑話.如果耶穌真的是猶太人的王,猶太人就不會把他交給羅馬帝國的官員.彼拉多雖然官位不高,但作為一位羅馬帝國派往治理巴勒斯坦地區的總督,他自己也有聽命於他的軍隊,因此,世界上怎會有一位王是沒有部下的呢?彼拉多認定,耶穌不可能是王,也不可能擁有甚麼國度和權力.

耶穌的回答與令彼拉多大惑不解.彼拉多斷定耶穌是王的傳聞是笑話,耶穌卻暗示自己真的是王；彼拉多斷定耶穌手下無兵無將,絕對不成甚麼氣候；但耶穌卻明言自己有國,並且口口聲聲以「我的國」,「我的國」來回應.到底是耶穌在痴人說夢話,還是彼拉多判斷錯了呢?

其實,彼拉多和耶穌兩個人都知道,這個世界的國是以武力維持,並且以戰鬥回應問題的.因此,按照這樣的標準,彼拉多斷定面前這位無兵無將,甚至見棄於自己所屬的民族的耶穌不可能有國作王.另一方面,耶穌卻重覆說明自己的國不屬於這個世界,不以這個世界的國以武力維持及以戰鬥回應問題的方法運作.若果耶穌的國也以武力維持以及戰鬥回應問題,他的部下就會出來擊退對手,耶穌就不可能被猶太人交到彼拉多手下了.不以武力維持,不以戰鬥回應問題的國,到底是怎樣的國呢?這是彼拉多無法想像的,或者也是這個世界所有的國無法想像的.這個世界的國大多都信奉「寧要核子,不要褲子」的教條,這教條是指示這個世界的國經常保持戰鬥狀態及提昇戰鬥能力.

彼拉多想像不到不以武力解決問題的國到底是一個怎樣的國,因此,把耶穌視為一個笑話.有些時候,耶穌基督的教會也不太能領會怎樣是一個不屬於這個世界的國,因此,有時也採用這個世界的國以武力解決問題的方法去回應挑戰.

天主教教宗若望保祿二世在梵蒂岡聖伯多祿大教堂舉行「聖年千禧寬赦日」彌撒,發表了天主教會史無前例的公開告解談話,坦承教會在過去兩千年期間犯下的罪愆和過失並尋求天主寬赦.教宗提及罪愆其中兩項為:一,公元一 ○九五年法籍教宗伍爾班二世號召發動的十字軍東征,而十字軍東征戰役綿亙二個世紀,造成了基督教徒與回教徒之間迄今仍未能消弭的敵對仇恨；二,公元一二三一年教皇格列高利九世發佈通諭確立宗教裁判所體制,偵查,審判「異端份子」,而一四七八年西班牙的宗教裁判所開始運作後,梵蒂岡當局已授權施加酷刑迫使異端份子告解懺悔,截至一八○八年約有卅萬人遭火刑處死.

耶穌說他自己的確有國作王,他的國不以武力維持,也不以戰鬥回應問題,因此他的國不屬於這個世界.到底耶穌的國是一個怎樣的國呢?

2．耶穌的國是屬真理的群體

如果彼拉多有幸讀到今天的書信經課－啟示錄1:4-8節,他會從中找到啟迪.

1 : 4    約翰寫信給 亞細亞的七個教會.願那位今在,昔在,以後永在的上帝,與他寶座前的七靈,和那忠信的見證者,從死人中復活的首生者,世上君王的元首耶穌基督,賜恩惠和平安給你們. 他愛我們,用自己的血使我們從罪中得釋放, 1 : 5    併於上節1 : 6    又使我們成為國度,作他父上帝的祭司.願榮耀,權能歸給他,直到永永遠遠.阿們! 1 : 7    「看哪,他駕雲降臨； 眾目都要看見他, 連刺他的人也要看見他； 地上的萬族要因他哀哭.」 這是真實的.阿們! 1 : 8    主上帝說:「我是阿拉法,我是俄梅戛,是今在,昔在,以後永在的全能者.」

第一,相對於世界的國以戰鬥迎擊別人,耶穌的國度致力使人從罪中得釋放.啟示錄指出,耶穌為了使人從罪中得釋放,甚至要自己付出極大的代價,流血被殺.是甚麼力量推動著這樣的行動呢?如此付出為了甚麼呢?經文指出,因為「他愛我們」.愛,推動人為所愛的對象的福祉而行動,甚至不計較地付出.

東歐有對母女感情上有了裂痕,13歲的女兒一直認為是母親的卑微地位,使她在人前抬不起頭.母親終日忙碌辛苦,也不能使女兒快樂起來.

2002年2月,母親邀女兒去阿爾卑斯山滑雪,此舉當然想讓女兒高興.母子倆在滑雪中,由於缺乏經驗偏離滑雪道迷路了,又遭遇了可怕的雪崩.母女倆在雪山中掙扎兩天兩夜,幾次看見前來搜尋她們的直升飛機,都因她們身穿的是銀灰色滑雪裝,而難以被發現.

終於,女兒因體力不支昏迷過去,醒來時發現自己已經躺在醫院裡,而母親卻不在人世了.醫生告訴她,是她的母親用生命救了她.

原來,是母親割斷自己的動脈在雪地裡爬行,用自己的鮮血染紅一大片白雪,讓直升飛機因此發現目標.

第二,相對於這個世界的國以民族劃分,耶穌的國沒有民族之分.彼拉多說:「難道我是 猶太人嗎?你的同胞和祭司長把你交給我.你做了甚麼事呢?」彼拉多的視野有限,看不見啟示錄的作者所看見的.啟示錄的作者看見甚麼呢?「看哪,他駕雲降臨；眾目都要看見他,連刺他的人也要看見他；地上的萬族要因他哀哭.」彼拉多以為耶穌的國只與猶太人有關,與他這位羅馬官員無關,耶穌的門徒都以為耶穌的國只與猶太人有關,與其他民族無關,當基督教從西方傳過來時,中國人以為耶穌的國只與西方人有關,與中國人無關,啟示錄指出,耶穌的國度與所有人都有關係,無論是敵對他的人,還是跟從他的人.在耶穌的國度裡,民族或種族再不是分別的條件,在耶穌的國度裡,包括所有民族和種族的人.

隨著交通的方便及人口流動的頻繁,這個萬民萬族同時同地出現的景象,不需要等到耶穌再降臨,而就可能出現在我們身處的地方.加拿大的多倫多在二○○六年進行了人口普查,發現在248萬居民中一半出生於境外,全市居民使用140多種語言,有200多種人種.在這種生活狀況下,狹隘的民族或種族主義將無法讓大家和平共處,耶穌那個包含萬國萬族的國度可以成為我們的參考.

四．總結:我們生活於兩個國度之中耶穌說:「我的國不屬於這世界」,既然他的國不屬於這世界,他來這世界做甚麼呢?耶穌回答彼拉多時說:「我為此而生,也為此來到世界,為了給真理作見證.」耶穌來到這世界,是要為真理作見證,讓人知道,世界的國雖然實在,但不是唯一的；世界的國雖然強大,但不是永恆的.耶穌的國度為生活於世界的國的人們提供了另外一個更理想的選擇,也為生活於世界的國的人們提供一個更長遠的選擇.耶穌說:「我為此而生,也為此來到世界,為了給真理作見證.凡屬真理的人都聽我的話.」耶穌呼喚我們進入他的國度,進入真理的國度.啟示錄指出,耶穌這個國度,是他犧牲自己建立的,「他愛我們,用自己的血使我們從罪中得釋放,又使我們成為國度」.

世界的國的運作模式似乎強不可擋,民族之間,種族之間的隔閡與仇視揮之不去,人與人之間的,國與國之間常擺出「戰鬥格」,然而,就是在黑暗之處,光明更顯得重要和珍貴.南非實施種族隔離政策,造成種族之間更深的誤會和仇恨.時移世易,黑人終於上場做總統,他們要怎樣對待這些誤會和仇恨呢?

一九九四年,曼德拉成為南非第一位黑人總統,剛好南非被指定為世界欖球賽的主辦國.欖球運動向來只屬於白種人參與,而他們所成立球隊,自然不受黑人歡迎,每有國際比賽,黑人往往寧願為敵隊打氣.便有人主張,趁此機會把欖球隊成員全部換上黑人,這個議題不斷發酵,國內局勢也開始騷動不安.

此時,曼德拉召見了欖球隊隊長,隊長忐忑前往,不料曼德拉只是親切地問候:「聽說你的腳受傷了!好一點了嗎?」不僅沒有提及更換球員之事,更在他們練習時親自探望,不斷給予關懷鼓勵.

這一次的會面,使球隊放下心防,並樂意配合曼德拉的親民政策,到黑人社區教孩子們打欖球,漸漸地,兩個族群的人們,開始有了友善的對話.

決賽當天,欖球隊受到全場不分種族的民眾歡呼而士氣高昂,最後打敗勁旅紐西蘭隊,成為冠軍.(蒲公英希望月刊2012年11月)

民族或種族之間到底要戰鬥,還是握手,很大程度取決於我們用甚麼心態和方法去回應問題.

戰鬥意識支配著國與國的關係不足為奇,最奇怪的是,戰鬥意識竟然滲透到最親密的關係中.

一位男子到圖書館借書,他問圖書館的女職員:「請問《婚姻的幸福生活》這本書放在哪裡?」職員回答:「這是幻想小說,到右邊第三排櫃子去找.」那人再問:「「那麼《夫妻的相處之道》這本書又放在哪裡?」職員再回答:「那是武俠小說,到左邊第一排櫃子找吧!」

有一位網友回應武俠小說式的夫妻關係,有如此的感悟.

「一個懂得愛的人,寧可扮演輸家,也不去打敗自己的愛人.打敗了她,或者他,你想得到什麼呢．．．男人和女人,我們彼此不是戰爭!把槍放下,把子彈扔進湖水裡.誰也離不開誰,就好好珍惜吧,別在自己心愛的人面前逞能,較勁,爭長短．．．讓我們彼此依靠,一起旅行,一起看日出日落,聽海浪的聲音.」(節錄自網絡文章「女人讓一步男人讓兩步」)

世界的國以強大的吸力逼使我們以世界的方式生活和回應事情,若我們能夠以耶穌的國度的方式回應事情,我們就能夠在這世界之中為真理作見證.



\newpage

\section{塗抹與記念}
\label{sec:176}
\textbf{伍渭文牧師 - 塗抹與記念}
\newline
\newline
\hyperref[sec:175]{< < < PREV SERMON < < <}
~
\hyperref[sec:toc]{[返目錄]}
~
\hyperref[sec:177]{> > > NEXT SERMON > > >}


日期:2012年12月16日

題目:塗抹與記念

證道:伍渭文牧師

經文:詩25

塗抹與記念是今日的講道題目.我相信這題目也說明了這星期媒體所關注的事情.有人試圖塗抹做過的事情,因為對自己不利,但媒體老是記念他說過的每一句話,指出這人前言不對後語,愈是塗抹,愈顯出破綻.

我訂出這題目是因為中古一位著名修士愛克哈特大師(Meister Eckhard)一句說話:「祇有塗抹的手,才能寫出真實的東西.」塗抹的手,不是掩蓋做過的事,因為每人做過的事情,總會留下痕跡,若要人不知,除非己莫為.對愛克哈特(Eckhard)來說,塗抹的手,把自己曾經寫下來的拋棄,承認錯誤,重新改寫.能夠塗抹改寫,是勇於面對真相,這些人才能寫下真實的東西.詩篇,特別是認罪的詩篇,像今天詩篇經課二十五篇,就是塗抹的手所寫,塗抹不是文過飾非,塗抹是渴望重新改寫自己的生命.

這是一篇大衛所寫的或給大衛所寫的認罪詩,從內容來看,應該是大衛所寫的.為何在將臨期第一主日選讀認罪詩?因為將臨期主調是上帝將臨的審判,叫我們好好反思屬靈生命.為了幫助記憶,全詩運用離合詩體(acrostic) 寫成,每一節以一個希伯來字母開始,第五節用了兩個不同的字母,如此類推,二十一節用了所有希伯來文二十二個字母寫成.但為何全詩二十二節,經課祇讀出十節呢?因為這詩亦運用了交叉配列(chiastic) 方式,帶出重要的

主題:A頭(1-3, 等候仰望 ), B胸(4-7節,拯救赦免), C肚(8-10節,守約慈愛), D腰 (11), C"臀 (12-14節,將約指示), B"膝(15-18節,救我赦免), A"腳(19-21節,投靠等候) .如人體頭腳,胸膝,肚臀呼應,所以一至十節和十二至二十一節是平行對應的,二十二是後話(post script). 經課一至十節已把全詩的基本意思帶出來了,中間交叉是(腰)十一節,是全篇中心思想:「耶和華啊,求你因你的名赦免我的罪.」上主啊,求你記念你的慈愛,塗抹我的過犯,我要重新改寫生命.

等候仰望- 生命的低谷,靈魂的攀登(vv 1-3；19-21)

我們先看第一段1-3節,這裡說到大衛生命陷在低谷,遭仇敵追殺,生命堪虞.

耶和華啊,我的心仰望你. 我的上帝啊,我素來倚靠你； 求你不要叫我羞愧, 不要叫我的仇敵向我誇勝. 凡等候你的必不羞愧, 惟有那無故行奸詐的必要羞愧.再看配對的十九,二十節:「他們人多,並且痛痛地恨我......求你保護我的生命,搭救我」(vv.19, 20).

詩的處境:詩人對上主祈禱時,他知道不單上主聆聽,他的敵人也聽他的祈禱,並且希望上主沒有作為,使祈禱的人感到羞愧,上主不聽祈禱,上主無能為力.

這班人是誰呢?誰人要取大衛的生命呢?這班人被詩人形容為無故行奸詐,可繙作叛國的人,Eugene Peterson 繙譯:「 Massage, v v. 2, 3: I have thrown my lot with you; You won"t embarrass me, will you? Don"t embarrass any one of us who went out on a limb for you. It is the traitors who should be humiliated. 」第七節詩人求上主不要記念他幼年的罪愆,很有可能是大衛年長時,被自己的親生兒子押沙龍叛變追殺.(撒下15) 詩篇第三篇副題-大衛逃避他親生兒子押沙龍的時候所作的詩,1, 2節就有相同的話:「

耶和華啊,我的敵人何其加增!有許多人起來攻擊我；有許多人議論我說:"他得不著神的幫助."」

大衛的經歷告訴我們,在人世間,想不到的時候,所愛的人可以突然對自己不忠；在這社會,想不到的時候,推心置腹的朋友因利益可以突然反臉成為敵人.大衛被親生兒子背叛,篡位之餘,還想奪命,可以算是人生的溶毀(meltdown), 這溶毀把他牢牢黏貼在地上,不能動彈,但他記念過去上主的慈愛:「我素來倚靠你」(v.2) ,「求你記念你的憐憫和慈愛」(v.6a), 「耶和華啊,求你因你的恩惠,按你的慈愛記念我.」(v.7) ,上主就叫他離地而起,舉心向上,靈魂攀登,重新改寫自己的生命.

危難使人得到屬靈的覺醒,渴想回到上主的身邊；生命到了低谷的時候,也是詩人靈魂往上攀登的時候.「耶和華啊,我的心仰望你.」

但仰望不是被動的仰賴,仰望是主動的親近,英文:O Lord, I lift up my soul.你們當仰望主這句話,在所有完整聖餐禮儀中開始時以啟應出現.願主與你們同在,(答) 也與你同在.你們的心當仰望主；(答) 我們的心仰望主.我們當感謝我主上帝,(答)感謝上主是理所當然的.在禮儀學上,稱為「舉心向上」.仰望亦有廣東話所講「的起心肝」,認認真真的意思.

等候仰望-有目標的人生

等候需要忍耐,仰望是向高處行.運動員為了至高的目標,長年累月和家人隔離,接受嚴格的訓練,他們定睛至高的榮譽,別無旁騖.倫敦奧運剛完畢,優秀的運動員就計劃四年後的奧運會.我想,當一個運動家信主後,比較容易明白甚麼是專心跟隨主,在屬靈的操練上,也較一般人持之以恆.

1981奧斯卡最佳電影烈火戰車(Chariot of Fire) , 1981) ,描述的主角李愛銳(Eric Liddell) ,是1924年法國巴黎奧運會400公呎金牌得主,被譽為蘇格蘭的羚羊,受派往中國傳道,中日戰爭期間被囚在淪陷區的日軍集中營.後來日本和英國交換戰俘,他放棄機會,讓給一位懷孕的婦人,不久死在集中營中.這消息是2008年北京奧運會開幕前中國有關當局公開的.他的去世,頓時令整個營的年輕人感到震撼,充滿空虛感.-位宣教士轉述,李愛銳臨終最後一句話是:「全所有奉獻」(``It"s complete surrender'').

在電影中,李愛銳用賽跑來說明基督徒的信心,他說:「我可以用賽跑比喻信心.是艱辛的,需要集中意志,心靈用力.衝過彩帶時你樂透了,尤其你為勝出下了睹注.......但這能持續多久呢?我想給你們更恆久的東西,但祇能道出蹊徑.我沒有勝出方程式,每人都有各自的跑法.但直奔終點的力量從何而來呢?從裡面而來.耶穌說:看哪,「上主的國度就在你們的心中,你們若真正尋求我,就必尋見.」若你們獻身基督的愛,就跑出了精彩.

等候仰望-追求難趣 (difficult pleasure)

歷史上其中一個最偉大的登山運動員馬洛里(George Mallory)窮一生精力要征服天下第一峰珠穆朗瑪峰,最後死在登峰的途中,終年三十八歲.有記者問他為何要冒險犯難去闖絕頂,他答道:「因為它就在那裡」(Because it"s there).

然而,在消費至上的社會,易趣(easy pleasure) 垂手可得.有人用信用咭在網上成功訂購最新的iPhone,急不及待發短訊向所有人報喜,像父親喜獲麟兒.隨著通訊科技的發達,藝術成為一種隨手可得的即時消費.看世界名畫,可在手機或電腦的熒幕瀏覽,不需花錢花時間到藏館看真跡；聽音樂也如是,隨時下載,很多是免費的,透過耳機在地鐵,在巴士就聽到一流的管絃樂團.

我們漸漸失去等候和期待的經驗.到博物館看畫作真跡,是一個期待和邂逅的旅程.我們有等候的焦灼,見到作品時,有相逢的愉悅；面對作品,有凝視的崇敬；離開展館後,我們有分手後的回憶.但上網瀏覽太容易了,不需等候,也沒有回憶,因手一撥又回來了.漸漸,我們滿足於易趣,忘記他們是虛擬的,是浮光掠影,不是真正的實體.而且,我們也慢慢失掉去藏館看真跡和到音樂廳聽音樂的的欲望.愈容易得到的易趣,扼殺了我們內心深處熾烈的欲望,與享受難趣的能力.我們不會傾情拼命渴求得到一樣東西,只會漫不經心希望得到一百樣東西.

尋求赦免 vv. 5-7

耶和華啊,求你將你的道指示我,將你的路指教我!求你指教我,引導我進入你的真理,因為你是救我的上帝.我整日等候你.耶和華啊,求你記念你的憐憫和慈愛,因為這是亙古以來所常有的.求你不要記得我幼年的罪愆和我的過犯；耶和華啊,求你因你的良善,按你的慈愛記念我.

擺脫誘惑

當詩人仰望上主,向高處行時,他有一種力量,使他離地而起.渴想駕駛飛機的人,不是因為駕駛倉舒服,其實空間很小很窄,很不舒服；也不是沖上雲霄的剎那刺激,時間很短,而且是危險的時刻.他被吸引是在烏雲之上,崇嶺之上,那渺無邊際可以自由遨翔的萬頃窮碧.十九世紀著名蘇格蘭講道家湯馬斯.錢伯(Thomas Chamber)有一篇著名講章,講章著名不是因為內容,乃因為講章的名稱:The expulsive power of a new affection(新愛的排拒能力),經文是約壹2:15-17「

不要愛世界和世界上的東西,若有人愛世界,愛父的心就不在他裡面了.」反過來說,愛上帝,舉心向上,向高處行,我們就可以擺脫世界的誘惑了.

渴想上主就要討他喜悅,遵守他的真理,是屬靈生命的操練,是難趣(difficult pleasure) ,但這難趣有能力排拒其他眾多易趣的誘惑.如果說new affection新歡是上帝,我們可以說:新歡排舊愛.詩人說:「

耶和華啊,求你將你的道指示我,將你的路指教我!求你指教我,引導我進入你的真理,因為你是救我的上帝.」真理是靈魂的糧食,救我們擺脫試探的力量.

不再迷失

詩人求上主以真理引導他,因為他曾迷失,離開真道,陷在罪中.「

求你不要記得我幼年的罪愆和我的過犯」,大衛把自己幼年過犯連繫到現在的苦難.兒子背叛自己,但青年時自己何嘗不是背叛烏利亞嗎?烏利亞是自己的愛將,對自己忠心耿耿,但因他的妻子美貌,就設計殺了烏利亞,奪其妻子.他在年青時對烏利亞不忠,背叛信任,少年的兒子押沙龍看在眼裡,父親的罪在兒子的生命中,埋下了種子,「私慾既懷了胎,就生出罪來,罪既成長,就生出死來.」(雅1: 15)

想到這裡,大衛求上主赦罪他的罪過,不要記念他的過犯.

持守聖約 vv. 8-10, 12-14

第三段經文vv. 8-10描述與上主歡悅的關係.第一段,在苦難中有屬靈的覺醒,舉心向上,仰望等候.第二段追求真理,不再迷失,求上主用真理引導自己.但最後的一段經文中,調子改變了,從仰望等候,祈求引導,轉到讚美:

耶和華是良善正直的, 因此,他必教導罪人走正路.他要按公平引領謙卑人, 將他的道指教他們.凡遵守他的約和他法度的人, 耶和華都以慈愛信實待他.誰敬畏耶和華, 耶和華必教導他當選擇的道路.他必安然居住, 他的後裔必承受土地.耶和華與敬畏他的人親密,他要將自己的約指示他們.(vv. 8-14)

在這裡詩人強調與上主聖約的關係,他以上主為樂.

將臨期是等候耶穌的來臨,最反思靈命的時刻.為何耶要降臨,因為我們有罪,他來是拯救我們,我們要有悔罪的心等候迎接耶穌.

他來自至高之處,我們要仰望他,舉心向上切慕耶穌,追求屬靈生命的成長,向高處行,以他為生命的目標.

馬洛里(George Mallory)窮一生精力要征服天下第一峰,最後死在登峰的途中,因為他說:這峰就在那裡,向他招手,被它吸引.

我們呢?



\newpage

\section{豐盛生命}
\label{sec:177}
\textbf{張義南牧師 - 豐盛生命}
\newline
\newline
\hyperref[sec:176]{< < < PREV SERMON < < <}
~
\hyperref[sec:toc]{[返目錄]}
~
\hyperref[sec:178]{> > > NEXT SERMON > > >}


日期:2012年12月23日

題目:豐盛生命

證道:張義南牧師

主耶穌宣告說:「我來了,是要叫羊得生命,並且得的更豐盛.我是好牧人,好牧人為羊捨命.」(約翰福音 10:10-11)親愛的弟兄姊妹,在我們未信主以前,生命中還有很多的罪,以及污穢的東西纏累我們.但當我們歸入耶穌基督以後,主耶穌的生命就進入到我們的生命裡面.在聖經當中,「生命」這個詞語,在希臘原文的意思是「活著的」,與「死亡」是相對的.作為基督徒,他的生命是有光的,有愛的,有信心,有盼望的.這個信心是世人看不見的.然而我們作為一個基督徒的,從外表來看,與世人好像沒有太大區別,但我們的靈魂跟他們不同.我們基督徒的生命,是悔改的生命,當耶穌基督傳道的第一句話是說:「天國近了,你們要悔改.」所以當我們歸主之後,聖靈就會光照我們,要我們為罪,為義,為審判,自己責備自己(約翰福音16:7-8).所以我們的生命,因著耶穌的名,我們便得著生命.以前我們活在黑暗當中,我們屬靈的眼睛沒有打開.直至我們信主之後,主耶穌照明了我們心靈的眼睛,我們又明白了屬靈的奧秘,我們得到聖靈的光照,就會深切認罪悔改,當我們悔改之後,耶穌的寶血就把我們一個一個的罪塗抹了,我們就成為一個聖潔的生命,我們的生命裡有耶穌的真理.

耶穌基督宣告說:「我就是道路,真理,生命．若不藉著我,沒有人能到父那裡去.」(約翰福音 14:6)親愛的弟兄姊妹,在我宣揚生命之道的時候,我先講述一下耶穌揀選我的見證.因為聖經說:「我們既有這許多的見證人,如同雲彩圍著我們 and \#8230; 脫去容易纏累我們的罪 and \#8230;仰望為我們信心創始成終的耶穌．」(希伯來書 12:1-2).耶穌基督揀選我們,其實是在萬世以先,已預定我們成為他的信徒,他的兒女,聖靈與我們的心同證了我們是上帝的兒女.

感謝主,我是在1989年信主的,那時候,我是在現稱為南京理工大學的地方讀書,當時還沒有人跟我傳福音.我看到法國的一位作家－盧梭所寫的懺悔錄,書中的第一頁寫到「將來我們每一個人,都要受到　神的審判.」他挑戰我們,說,我們能在末日的時候,面對上帝的審判嗎?當時心裡真的非常震撼,因為我們在大陸學習的是唯物主義,是無神論的,我受了十幾年的教育,我學的是馬克斯主義哲學,馬克斯主義哲學認為基督教是人類的鴉片,基督教是帝國主義侵略中國的工具,在中國出版的小說及書籍,都是用來否定基督教信仰的.所以作為一個中國的大學生,要信耶穌基督不是一件容易的事.所以當我看到盧梭的懺悔錄時,我便開始時思考,人從那裡來?以後往那裡去?人是歸向那裡?我死了之後會怎麼樣?

之後在1988年,在中國的大學開始了自由化的思潮,我們並且對馬克斯主義和無神論主義開始不相信了,對傳統的文化也都不相信了.於是在南京理工大學讀書時,我走到當地的一所教堂,因為我以前學過哲學,文學,物理,科學,心理學,各樣學問.也學過歷史,連佛教,道教,伊斯蘭教,各樣宗教各樣文化,我都看過.包括氣功,我都有研究,另外占卜,算命,我都看過,但是我心靈感到沒有歸宿.當我走到南京莫愁路堂裡面,我坐在椅上,那裡跟這所教堂有點像,是由馮玉祥和宋子文幫助興建的教堂,是很有名的.我當時坐在下面,看到上面有人彈鋼琴,之後又由人講道,我心裡面就感到進入了心靈的港灣,進入了一個避風港,我心裡面有一個很強烈的意念,這個地方就是我的歸宿.

當時有一個學聯會,由一班老師籌辨,其中一個是東南大學的一個副教授,他是在1948 年加入大學生基督徒聯合會.他們認為當時中國的環境有一些寬鬆,而88,89年時候,是趙紫陽與胡耀邦治政的時候,政治環境的確是有一些寬鬆,所以他們就在那教堂開始了學生團契聚會.我當時參加了學生團契,而老師們籌辨學聯會也是沒有壓力的,都是為了傳福音,老師們亦為我們做見證.

在1948年,錢老師在上海交通大學參加了學聯會信主的,之後文化大革命開始,他講述當時紅衛兵怎樣批鬥又打他,要他放棄信仰,當中又怎樣受苦.在50年代,他在王明道先生於北京基督徒會堂有過半年時間的聚會,王先生說三自教會是不屬靈的,是有青年會的背景,但家庭教會是屬靈的,屬主的,而他又做了生命見證.當時我非常受感動,我想見了那麼多黨幹部,宣傳模範,因為我曾在大型的企業工作過,我們廠裡有勞動模範,有省勞動模範,有優秀黨員,但是我看到他們的品質,他們的生活方式以及各方面的誠實,對人的真誠,遠遠不及這個錢老師.我以前是做秘書的,是替那些勞動模範,表彰人物寫文章,又幫他們做宣傳的.因為在大陸有很多文章,書藉用來宣傳黨員有多麼好,多麼光榮偉大,但是實際上我並沒有找到這樣的一個人.但是我卻在錢老師身上找到這種誠實,他待人的真誠,又沒有謊言,他的仁愛深深的吸引了我.在我現在看來,他是一個有基督豐盛生命的基督徒.所以生命能夠生生命,一個好的生命能感染周圍的人,讓人能夠看到耶穌基督所散發出來的那種馨香氣息.

在大學校裡面,每逢星期六同學們都參加各樣的活動,例如跳舞之類.可是我對這些都沒有興趣,我寧願去學聯會參加團契聚會,儘然車程行遠.當時我信主的動機是我不用下地獄了,而且可以永遠住在天堂裡面,並沒有想到要去傳道.因為當時是八十年代末期,在中國大陸如果能做文評,又會寫作,而且我是報章的記者,也是大型企業的秘書.大家對我的前程都是看好的,他們都說以前的秘書當了副廠長,然後當了處長,之後更當了主任,我當時只有三十來歲,應該是很有前途的.但是有一天晚上的零晨時份,當時我在大學的宿舍裡,我睡在上鋪,聖靈來感動我,有一把聲音跟我說:「你要被聖靈充滿.」因為當時我是初信,又不明白什麼是被聖靈充滿,而且當時學校裡是很寧靜的,沒有什麼其他的聲音,我卻聽到這聲音,於是我便問:「怎麼被聖靈充滿呢?」之後,我便感到好像有電流一樣,流遍全身,從頭到腳就被聖靈充滿了,膏抹了.我覺得全身有無窮無盡的力量,心裡有一個意念,就是有了聖靈充滿的能力,什麼困難事情都可以克服完成.然後這聲音又繼續跟我說:「信上帝是你一生中最大的事情,信上帝是你一生中最大的事情,信上帝是你一生中最大的事情.」這聲音重覆了三遍這句話之後,我就爬起來,因為當時宿舍有些暗,我就打開燈,之後我把這些話記下來,「為此為大信仰上帝,信上帝是我一生中最大的事情.」這是上帝第一次對我的呼召,這次的呼召是要我傳道,信上帝是一生中最大的事情,不要為升官發財,成名成家,我過去的一些夢想,計劃,都要破碎.我要進入上帝給我的命定裡,要我去傳道的計劃當中.

後來我去了大型企業當秘書的工作,要給廠長,官員寫報告,可是在寫那些報告時,心裡不平安,因為在寫那些報告時,我要寫一些假話,並不是心裡想寫的.又有一些數據要造假的,我覺得這樣做基督徒不可以不誠實的,所以我把工作辭掉.當時身邊的幹部,工程師,副總,幹部,處長,主任等等,他們都感到不可思議,因為秘書這份工作多好,往上可以升官,其他人想掙也掙不到,而我為什麼會辭掉呢?所以他們都說我信上帝信迷了,變得不正常,但是我覺得是上帝呼召我.

當時我是在一所位於大山溝裡的軍工廠工作,我找到了一間教會,用來聚會的房子是非常簡陋,房子其實是養雞棚,連一般農舍都不如,只是在一個土堆上蓋一個茅舍.我從南京莫愁路堂一下子來到養雞棚的聚會點,這樣的落差太大了,在莫愁路堂裡聚會的,都是大學教授,大學同學,南京市民,身份好像都比較高.但是在偏遠的山區裡,那些基督徒都是蓬頭垢面,衣衫襤褸的,都是社會上邊緣的人.我記得有一次,有一位穿得破破爛爛的老姊妹,因為她的腿不好,又加上是下雨天,路上很滑,在聚會之後我扶著她,怕她滑倒.結果給工廠裡的職工看到了,他就對我說:「張義南,你那麼有文化,一直當秘書,又會寫作,各方面都好,大家都看好的.現在和沒文化的老太太她們混在一起,有什麼前途呢?」當時我太太依然在醫院裡當護士長,但是她對我也不理解,父母也對我不理解,朋友也對我不理解,為什麼信上帝信到要騎破車子去農村傳道?但是我當時只是想傳道,因為聖靈在我裡面燃燒著.感謝主,有一次我可以在容納一千多人的大教堂裡講道,我也和三自教會合作在農村裡建教堂.

後來在我家附近有一位姊妹,她是《迦南詩歌》創作者－小敏姊妹,他們的家庭教會傳道人來到我家.因為我們三自教會唱的是讚美詩,而他們家庭教會唱的是小敏姊妹的《迦南詩歌》.因為我從前是一名記者嘛,所以我也想去看一下,我就隨著農村來的傳道人參加他們的家庭教會.當時我到那些年青的農民,不是在三自教堂裡聚會,而是農村的院落裡聚會,他們唱著:「點燃,點燃福音的火,吹起,吹起福音號角」.他們唱的非常高亢,他們拍著手,臉上放著光,內心的復興的火,深深的打動我.我心裡說:「好呀,在這中華大地上,一場波瀾壯闊的福音運動就在農村升起啦.」我說我以後就要跟你們了.在他們的身上,我看到中國基督教會的希望.因為在三自教堂裡有很的限制,有三定政策和十不准,像不准接待海外的牧師,不准看海外的書籍,不准在教堂以外傳福音.另外,如受洗人數,按立牧師和長老都須得到批准.因為有很多限制,所以在傳福音時受到很的轄制,不能自由傳福音.而且由於從前栽培我的錢老師他們也認為,家庭教會是比較純正的教會.他們從前已經跟我這樣講,而我也親眼看到實際情況.當時三自教堂也有請我去講道,因為我有文化,有知識.但當三自教會知道我參加過家庭教會之後,都遠離我了,我想是因為他們怕被抓,怕罰款,怕受到連累,他們有很多的怕.所以之後我參加了家庭教會,開始傳道,之後更放下我的工作,全然奉獻來傳道.

在開始出去傳道的時候,我走到山上跟上帝作了一個禱告,我說:「主啊,你呼召我來傳道,我現在不敢工作,也沒有工資,那我的孩子跟妻子該怎麼辨?請你供應我的經濟需用,我才把工作全放下,全然奉獻來傳道.」可是上帝並沒有垂聽我這個禱告.後來我看到一本屬靈書籍,書中寫道要對上帝有信心,你要先憑信心邁出一步,上帝才會幫你開出路.所以要上帝供應你,先要憑信心把工作辭掉,出去傳道,然後看上帝供應不供應你.我想雖然你現在不供應我,但我全心相信將來衪必然供應我,我就把工作辭了.後來在我傳道半年之後,上帝就安排了一個人,主動地供應我的經濟方面的需要.所以主耶穌做就我們的生命是一個有信心的生命,這樣的一個信心不是人的理性所能想像的.我們都有很多自己的計劃,打算.但是上主要我們憑信心來仰望衪,來跟著衪走.這樣我們才能感受到,這是上主要我們做的,不是我們自己人的計劃所能做的事情.所以奉獻就是信心,而我就在這方面走了第一步.

之後我們家庭教會受到了逼迫,公安開始關注我們了.那時候我跟趙天恩牧師在河南鄉村傳道,在接觸他以後,主給了我一個意向,就是要我來寫作中國教會的歷史.因為在鄉村傳道的時候發現,有很的傳道人在文革政治活動中受苦,他們都有很多美好的生命見證,我覺得這些見證都被隱藏著,我認為要把他們都採訪和記載下來,讓後代的基督徒看到.於是我有了這負擔去採訪這批屬靈的老人.主耶穌真的非常奇妙,衪給了我們這個使命,就會給你開出路,讓你憑信心往前走.衪讓我遇見趙天恩牧師,他是研究中國教會的一個專家,因為他是英語思維,而我漢語基礎比較好,以前我當過高中的語文老師嘛,我們可以寫作,有互補作用.他不方便隨便在農村裡走,因為他來自香港,但我卻比較方便到各地去採訪見證.因為在九六年開始,官方對家庭教會有嚴打鬥爭,就是凡是不參加三自教會的家庭教會打成為邪教.這樣迫使到各地不同的家庭教會領袖聚在一起來禱告,當時我寫了一份中國家庭教會合一的論綱,就是要表明我們的信仰是純正的基督教信仰,不是邪教.我們家庭教會有很多不同的宗派,有基要福音派,有靈恩派,小群,奧秘派,改革宗長老會,一次得救永遠得救,耶穌家庭守真道教會,重生派等等.在我研究教會歷史,我發現這些農村裡的教派都被打成隔牆,互相不了解,都認為自己的教派最正確,而其他教派都不正確.所以作為大教派的家庭教會領袖,都去找各個教派的領袖來連結交通.於是我走到那些偏遠的山區,大家都約定好時間,在白天,在晚上,一起來討論應該怎樣合一,怎樣能夠表明我們的信仰是純正的.但在這個過程當中,一批一批的教會領袖都被抓起來.他們有的要挨打,有的去了勞改,有被罰款,有的被拘留.而我就一直在逃跑,有幾次差點被抓住,我被抄家了三次,其實我也知道在家裡住是有危險的.有一次,有六個公安翻牆跳進我家處抄家,不過剛好那晚我不在家裡,從此之後,從那以後,我開始在外面被追趕了五年半.在這五年半時間,我不能夠公開回去,因為如果公開回家,我會被公安派出所抓起來.我只能在夜深人靜的時間,周圍沒有線人,才能悄悄地回家.在那五年半的時間,我一直住在外邊的接待家庭,並一直在全國各地去傳道.無論在農村或是城市都跑去傳道,在家庭教會有兩個異象:「逃」與「傳」.

從前王明道與林獻羔都是坐在家裡等人來抓的,但到了八,九十年代,傳道人都學聰明了.因為如果被抓起來的話,要坐三至五年牢,那麼那幾年就不能夠傳道了,所以一被抓就要跑了.那麼一跑呢,無論是白天或晚上,就要天天傳道,因為已經不能幹社會的工作了.當一抄家以後,我們就加入家庭教會的團隊,家庭教會的領袖亦很高興,因為來了一個年青,有學識,又有文化的傳道人加入他們.當時我就是在這樣的一個生活光景.

我在農村教會時,接受到很多的訓練,又看到各種不同的團體.感謝主,因為祂造就我們要有豐盛的生命,衪要把我們的虛榮,驕傲,愛面子的心放下來,有個謙卑順服的心,要我們的生命有磨練,造就.在這個傳道的過程中,雖然我們看不到未來,一直在傳道,那以後怎麼辨呢?自己的孩子在上學,沒法管,我的妻子,也沒法管,我們就這樣傳道傳下去.我們聚會和傳道不像你們這邊可以有這麼安舒的環境,在我每次上講台講道的時候,腦裡常會想到,會不會在講道的時候,公安局突然衝進來?確實有一天晚上,當我們講完道之後,就有十八個公安把我們包圍了,他們用掍打我們.當時有弟兄姊妹往外跑,因為那時是晚上,有的被抓回來,有的跑掉了,而我就被抓到車上去,帶我到公安局那裡去.當車子駛到一個轉彎位時,車子就放慢行,我就趁機把車窗打開,就跳下去了.因為我知道被抓以後,一定會被關起來,關上多少年也說不定.公安之後向天開了一槍,意思是警告弟兄姊妹不要打算再逃了,而我在那次是逃脫了.

另外有一次,當我正在寫教會歷史的時候,有四名公安闖入,我就說了幾句話應付他們,但他們卻帶走了我一些寫作用資料.我知道他們回去細看那些資料之後,一定會回頭把我抓起來,所以我馬上跳牆出去,又把那些採訪的磁帶收藏到另一個地方.有時在聚會的時候,聽到警車開過來的聲音,我們都馬上收拾然後往後翻牆跑.就這樣跑了五年半之後,我被抓起來,就這樣被關起來兩年.那時白天晚上都被鎖上手扣,說要判罪我三年至無期徒刑.於是我禁食祈禱了八天,最後被判了兩年.

在這兩年監獄的時間裡面,我與流氓和黑社會,警察等人關在一起,我給他們唱詩,傳福音,有很多生命的造就.

為什麼主耶穌要我坐監?我感到因為是主耶穌愛我,衪才叫我坐監.但是妻子無法理解,她問我:「我家的張義南替上帝傳道,一邊跑要傳道,但為什麼卻要把他弄到監獄裡呢?」她跟上帝禱告,她跟上帝辯理,但後來聖靈感動跟她:「你要是想坐監,你還不配呢.」坐監是要有資格的,所以我在監獄裡很感謝主.在監牢的時候,主耶穌造就我的生命謙卑,順服,讓我們基督徒品格的流露,讓我們傳福音,讓我們參與屬靈爭戰,讓我們體會到中國各個方面的事情,我們可以有很好的禱告.我在受苦的過程,領受了更大的啟示,奧秘,異象.我回顧我走過的路,我感覺到主耶穌一步一步的造就我們在信心方面,把各種考驗面前,一次一次的更新我們的生命,一次一次的突破,一次一次提升我們的屬靈境界,提升我們的屬靈能力,提升我們的屬靈權柄.

有一個先知曾經這樣說:「如果一個基督徒,他的生命是一帆風順的,那上帝不會把大的使命託付給他.」一個傳道人如給他磨練多,給他造就多,給他苦難多,那給他的使命也大.所以中國家庭教會的傳道人,像我這樣經歷的有很多,其中有位同工被抓過十六次.好像一位家庭教會的領袖,他是一位老人家,他今年已經將近八十歲了,被判過兩次刑,一次被判了廿次,另外一次被判了十二年,他因為堅持家庭教會的立場,所以共被被判三十二年.然後他的兒子和其他家人,算一下,一共被判了六十年監.我在訪問期間,採訪過很多這樣催人淚下的見證,這就是主耶穌造就中國教會,愛中國教會的一個明證.因為上帝讓我們中國教會承受了許多苦難,就是要我們承接更大的使命.苦難愈深,你的使命就愈大.

所以感謝主,在我們面對困難的時候,在面對逆境的時候,我盼望各位弟兄姊妹,憑信心仰望那創始成終的耶穌基督.那時候不要抱怨上帝,要凡事感恩,上帝就會為你開出路,衪會把一個一個的難處,每過一關,就一個一個的把它挪開,祂就會造就你一個豐盛的生命.

來讓我總結一下,我們的基督徒的生命是有能力的生命,我們的生命是永恆的生命,聖潔的生命,有光,有愛的生命,我們能夠結出屬靈九種果子的生命,就是仁愛,喜樂,和平,忍耐,恩慈,良善,信實, 溫柔,節制的生命,那我們的生命經歷過考驗,磨練,然後得到得勝的生命.主耶穌曾告訴七教會,得勝的才能得到賞賜.當我們有豐盛的生命的時候,我們的內心才有真正的平安和真正的安舒,無論在順景,無論在逆境,你在監牢當中,你在豪華的飯店當中,你都是平安的,都是安舒的,將來的歸宿就是天上的賞賜,天上的冠冕.願　天父上帝祝福我們弟兄姊妹,一同得著豐盛的生命,都得到天國豐盛的賞賜.願　上帝祝福中文大學祟基學院禮拜堂,所有親愛的弟兄姊妹.

我們一同禱告,天父上帝我們感謝讚美你,主啊,感謝你使我能夠在這禮拜堂,在弟兄姊妹面前做見證.主啊,這見證的榮耀是屬於你的.主啊,保守我們不偷竊主的榮耀.主啊,感謝讚美你.主耶穌,你是羊的好牧人,要你羊得生命,並且得的更豐盛.主耶穌基督,你是豐盛的基督,上帝本性的一切豐盛,都有形有體的住在你裡面,你是讓我們在你裡面得豐盛,使我們的心能剛強起來.主啊,使我們勝過世界的誘惑,使我們勝過外在的逼迫.主啊,使我們堅定信心來跟隨你.主啊,你祝福弟兄姊妹的家庭.主啊,你祝福香港的眾教會,能夠彼此相愛,彼此連結,合而為一.主啊,讓我們能夠傳講個人得救的福音,重生的福音,悔改赦罪的真道.主啊,也能傳講的國度的福音,使香港成為一個基督化的城市,成為一個福音化的城市.主啊,請你在中文大學,合你心意的教會,傳講福音的教會,與我們的弟兄姊妹同在,使教會的人數天天加增,榮耀主你的聖名,高舉耶穌基督的旌旗,願你十字架的旌旗在這裡高高的飄揚.主與我們同在.祝福以下的時間,祝福我們領受主耶穌的聖潔和寶血,使我們的生命更加的豐盛,主啊!使我們的血液裡面,身體裡面,有主耶穌的生命,我們的血統是屬上帝的兒女的,神聖的血統的生命,主你與我們同在,保守我們信靠你到底,愛你到底,直到見你的面,永永遠遠都屬於你,禱告祈求感恩,奉主耶穌基督的名,阿們.



\newpage

\section{兩個不平凡的孕婦}
\label{sec:178}
\textbf{鄧瑞強博士(香港神學院) - 兩個不平凡的孕婦}
\newline
\newline
\hyperref[sec:177]{< < < PREV SERMON < < <}
~
\hyperref[sec:toc]{[返目錄]}
~
\hyperref[sec:179]{> > > NEXT SERMON > > >}


路加福音 1:39-45
\newline
\begin{longtable}{cl}
\hline
\hline
章節 & 經文 (和合本修訂版)\\
\hline
1:39 & \begin{tabularx}{0.7\textwidth}{X} 在那些日子，馬利亞起身，急忙前往山區，來到猶大的一座城， \end{tabularx} \\ \\ \relax
1:40 & \begin{tabularx}{0.7\textwidth}{X} 進了撒迦利亞的家，向伊利莎白問安。 \end{tabularx} \\ \\ \relax
1:41 & \begin{tabularx}{0.7\textwidth}{X} 伊利莎白一聽到馬利亞問安，所懷的胎就在腹裡跳動。伊利莎白被聖靈充滿， \end{tabularx} \\ \\ \relax
1:42 & \begin{tabularx}{0.7\textwidth}{X} 高聲喊著說：「你在婦女中是有福的！你所懷的胎也是有福的！ \end{tabularx} \\ \\ \relax
1:43 & \begin{tabularx}{0.7\textwidth}{X} 我主的母親到我這裡來，為何這事臨到我呢？ \end{tabularx} \\ \\ \relax
1:44 & \begin{tabularx}{0.7\textwidth}{X} 因為你問安的聲音一入我耳，我腹裡的胎就歡喜跳動。 \end{tabularx} \\ \\ \relax
1:45 & \begin{tabularx}{0.7\textwidth}{X} 這相信的女子是有福的！因為主對她所說的話都要應驗。」 \end{tabularx} \\ \\
[1ex]
\hline
\hline
\end{longtable}
日期:2012年12月30日

主題:兩個不平凡的孕婦

證道:鄧瑞強博士(香港神學院)

經文:路加福音1:39-45

各位弟兄姊妹,早安.聖誕節快到,願神賜予世人的喜樂和平安臨到你們.

今日福音書的經課,講到兩個女人閒話家常.兩個女人都「大了肚」,「大了肚」的女人總有很多話題,單單講一下孩子叫什麼名字也可談上半天.在香港,她們可能在討論應去哪間幼稚園排隊讓小孩子入讀.

我們讀一讀這段

經文:路加福音1:39－451:39　那時候,馬利亞起身,急忙往山地裡去,來到猶大的一座城；1:40　進了撒迦利亞的家,問伊利莎白安.1:41　伊利莎白一聽馬利亞問安,所懷的胎就在腹裡跳動.伊利莎白且被聖靈充滿,1:42　高聲喊著說:「你在婦女中是有福的!你所懷的胎也是有福的!1:43　我主的母到我這裡來,這是從哪裡得的呢?1:44　因為你問安的聲音一入我耳,我腹裡的胎就歡喜跳動.1:45　這相信的女子是有福的!因為主對她所說的話都要應驗.」

那時,馬利亞從聖靈感孕,「大了肚」.她立刻想到另一個也「大了肚」的親戚伊利莎白.馬利亞是否去尋求產前輔導?抑或,路加只是想表達,一個領受了聖靈指引的人,總是立刻有信仰的行動的.兩個懷孕的婦人相遇,應有很多話題.

這兩個懷孕婦人,她們的懷孕狀況都很不尋常.伊利莎白是一個老人家,應該已過了生育的年紀.她太老,不可能生育.老人家懷孕,若在今天,這位老婦人必定說:「Oh, My God! 糟糕!」另一位女子馬利亞,她太年輕,還沒有和男性有性接觸,不應該生育.這位未婚小女孩懷孕,若在今天,這位小女孩也必定說:「Oh, My God! 怎算?」一個太老,不可能生而神蹟地懷孕.一個太嫩,不應該生卻同樣神蹟地懷孕.

伊利莎白,她的生命尊嚴,全倚靠為夫家傳宗接代.這是一個古老的時代,女人的幸福要由男性去賦予的,女性的尊嚴要由男性去定義的.她等了很多年,她等待懷孕,她等待自己的生命有一點內容.等了很多年,卻一無所出.她失望嗎?肯定失望.她覺得人生有遺憾,生命有羞恥.當伊利莎白懷孕後,她說:「主在眷顧我的日子,這樣看待我,要把我在人間的羞恥除掉.」(路 1:25)這句話多麼沉重,可看出懷孕對這女人多麼重要.神看顧她,使她懷孕,她從這懷孕中找到生命的尊嚴.

馬利亞,她太年輕,她未結婚,她未親近男性,她還未長大.她生命的未來,有很多可能性.她從未想過懷孕,她不應該懷孕.對伊利莎白來說,懷孕會除去她的羞恥.但對馬利亞來說,懷孕只會帶來羞恥.

看來,懷孕對女性來說,是一件複雜的事.我不是指生理意義上的複雜性,而是指懷孕對女性身分的意義尋索上的複雜性.有時,懷孕除去羞恥；有時,懷孕卻帶來羞恥.能否懷孕?何時懷孕?牽涉複雜的權力結構.有些地方,女性是不能任意生育的,她的生育是有限額的.當然,錢多的人,限額會寬鬆一點.有些時代,女人的地位是由她生了多少男孩而決定的.當然,當女性沒法懷孕時,責任通常在女方.

伊利莎白的懷孕,反映了她的幸福全部由這懷孕賦予.這是一個舊時代.這時代將要結束了.

馬利亞的懷孕,是一種特別的懷孕,代表著一個新時代.馬利亞是「童女」懷孕.「童女」,或者說「處女」,當然是指未與男性有性關係的女人.神學家去注解「童女懷孕」這件事的信仰含意時,發現「童女懷孕」可以包含一種特別的神學信息,即:當女性的懷孕不必由男性去賦予時,則意指這女性的身分意義不必由男性去定義.懷孕,是最女性化的行為,唯有女性才能做,這是女性身分的關鍵特徵.若這身分的關鍵特徵不必男性去賦予,或者說,不必由男性去操控,則這是女性生命的一種解放.馬利亞的女性身體,全然是神創造裡的禮物.她的生命意義,不必由男性去定義,不必由男性去操控.沒有男性要求他生孩子以確立她在家庭中的地位,沒有一男性主導社會要求她生孩子以確定她的生存意義.在這裡,馬利亞的懷孕,呈現了她的女性特質,呈現了神的恩典,呈現了她的生命意義全由神的恩典去定義.在神的呼喚下,馬利亞可以自由地回應,以其女性生命,發揮女性生命的完滿意義.

這「懷孕」標誌著一個新時代,在這新時代裡,被權力社會壓制的人,能在神的恩典裡得到解救.又或者說,明白神的恩典的人,會努力在權力社會裡,去解救那些被壓制的人.明白這一點,又努力實踐這一點,基督教的福音才真是人類的好消息.明白這一點,又努力實踐這一點,聖誕節的故事,那童女懷孕生子的故事,才會對我們產生意義.明白這一點,又努力實踐這一點,我們才能了解馬利亞對自己的懷孕,為何有以下說法,她說(路1:46-53):

1:46　馬利亞說:我心尊主為大；1:47　我靈以神我的救主為樂；1:48　因為他顧念他使女的卑微；從今以後,萬代要稱我有福.1:49　那有權能的,為我成就了大事；他的名為聖.1:50　他憐憫敬畏他的人,直到世世代代.1:51　他用膀臂施展大能；那狂傲的人正心裡妄想就被他趕散了.1:52　他叫有權柄的失位,叫卑賤的升高；1:53　叫飢餓的得飽美食,叫富足的空手回去.

神藉馬利亞的懷孕,顛倒了人世間的權力秩序.

為免大家誤會我在這裡講一種不必男性的合作而女性自行懷孕的信息,我要鄭重說明,生理上的生育固然必須男女相愛相合才能發生.我在這裡的講法,是一種神學意義的論述,是對「童女懷孕」可能有的神學意義作一種開展而已.

馬利亞去探訪伊利莎白,向伊利莎白問安.這時,不單兩個孕婦相遇,也是兩個胎兒相遇.

路加福音的經文說:1:41　伊利莎白一聽馬利亞問安,所懷的胎就在腹裡跳動.伊利莎白且被聖靈充滿,1:42　高聲喊著說:「你在婦女中是有福的!你所懷的胎也是有福的!1:43　我主的母到我這裡來,這是從哪裡得的呢?1:44　因為你問安的聲音一入我耳,我腹裡的胎就歡喜跳動.

這是將施洗約翰和耶穌將來的使命,提前寫入他們仍是胎兒時相遇的情節裡.

胎中的施洗約翰,面對胎中的耶穌時,歡喜跳動,這是要表達施洗約翰將來的事奉,是迎接耶穌的到臨,並為他的救贖工程開路.

伊利莎白的懷孕,是老婦人的懷孕,是屬於舊時代的.馬利亞的懷孕,是童女懷孕,是一個新時代.伊利莎白肚中的胎兒,也感受到一個新時代的來臨.胎中的施洗約翰,歡喜跳動,是因為他感受到神的兒子的來世,是因為他體會到這是一個新時代的來臨.耶穌降生,的確開展出一個新的時代.在這時代裡,人能體會到神對人的愛有多深,深到神可以放下神聖而成為受苦的人.在這時代裡,人能看到歷史的遠象,一個以愛連繫,彼此分享的新世界.

有人說,現代生活是沉悶不堪的,這裡充滿機械的力量,充滿鐵籠一樣的結構.人在這裡,只感到壓力,需要不斷的刺激去尋找生活的樂趣.找不到遠象,找不到新意,找不到驚喜,故此,很難「歡喜跳動」.胎中的施洗約翰,他能「歡喜跳動」,因為他在耶穌那裡,找到生命不再一樣的路徑.

最後,我想多講一點馬利亞.

路加福音說:1:42　〔伊利莎白〕高聲喊著說:「你在婦女中是有福的!你所懷的胎也是有福的!1:43　我主的母到我這裡來,這是從哪裡得的呢?1:44　因為你問安的聲音一入我耳,我腹裡的胎就歡喜跳動.1:45　這相信的女子是有福的!因為主對她所說的話都要應驗.」

1:42及1:45都提到馬利亞是「有福的」,但兩節經文的原文用字不同.前者較傾向指「神對人的一般性的祝福」；後者較傾向指「因著人回應神的要求,這人成為特別蒙福的人」.舉例說,「為義受逼迫的人有福了,因為天國是他們的.」(太5:10)這裡「有福」的用字,便是後者.這人有福,是因為他「為義受逼迫」.因著他的行動,神祝福他.

路1:42指馬利亞的生命,蒙受神對人的一般祝福.縱使她是一個女子,一個未婚的,沒有男性作依傍,沒有孩子去增添她的尊榮的小女子,她仍是蒙受祝福的.

路1:45卻是指馬利亞的生命,因著她的抉擇,因著她對神的回應,她成為有福的人.

我們對這兩節經文看得再詳細一點.1:42　〔伊利莎白〕高聲喊著說:「你在婦女中是有福的!你所懷的胎也是有福的!」

一般的理解,是認為馬利亞蒙福,只是因為她懷了耶穌.但經文中的兩個子句「你在婦女中是有福的!(and)你所懷的胎也是有福的!」沒有因果結連,只有平行的連接.故此,有學者認為,馬利亞自己的蒙福,與她所懷的胎的蒙福,是獨立事件.縱使她沒有懷這胎,她也是蒙福的人.這是一個重大信息,一個女子可以不憑丈夫和兒子,縱使自身多麼微少,她仍可以頂天立地說,她蒙神祝福.

1:45　「這相信的女子是有福的!因為主對她所說的話都要應驗.」

她相信神對她生命的呼喚,她願意配合神的計劃,她的回應行動,使她成為特別蒙福的人.(可參考路11:27-28的說法.)

之前,我講過,伊利莎白的懷孕,會除去她的羞恥.但對馬利亞來說,懷孕只會帶來羞辱.馬利亞若要配合神的計劃,則她要承受的,是生命的艱難,而不是福樂.她要回應神,就不只是被動地回應,而是準備主動地承受莫大的痛苦.有福,不是指一帆風順,萬事如意.馬利亞是有福的,這是神對她的生命內涵的一種評價.在人生裡,她將踏進苦難,與其兒子耶穌,共行十架路.

在12月14日,美國發生校園槍擊案,死了20多人.這學校的校長,一位女士,也死了.這校長迎向槍手,試圖阻止他.我不知道,她當然想到什麼?她有什麼掙扎?有報導說,這女士如此做的時候,沒有半點猶疑.她的決定使她承受了莫大的痛苦,但她是有福的.

願耶穌降生的故事,讓你看到生命的遠象.在神恩典裡,這是一個新的時代,這應該讓你的心靈「歡喜跳動」.在這裡,神重視你的價值,無論你是一個多麼微少的人.也在這裡,你被呼召去參與耶穌降生的故事,以你的生命,讓他人得福,縱使這樣做的時候,你可能要作出某種犧牲.

願主的恩典臨到你們.

但願:榮耀歸於聖父,聖子,聖靈.阿們.

(本文的神學與釋經觀點,參考 James P. Martin, ``Luke 1:39-47," Interpretation 36 (October 1982): 394-399.)



\newpage

\section{孩子不見了}
\label{sec:179}
\textbf{林豪恩先生 - 孩子不見了}
\newline
\newline
\hyperref[sec:178]{< < < PREV SERMON < < <}
~
\hyperref[sec:toc]{[返目錄]}
~
\hyperref[sec:180]{> > > NEXT SERMON > > >}


路加福音 2:41-52
\newline
\begin{longtable}{cl}
\hline
\hline
章節 & 經文 (和合本修訂版)\\
\hline
2:41 & \begin{tabularx}{0.7\textwidth}{X} 每年逾越節，他父母都上耶路撒冷去。 \end{tabularx} \\ \\ \relax
2:42 & \begin{tabularx}{0.7\textwidth}{X} 當他十二歲的時候，他們按著過節的規矩上去。 \end{tabularx} \\ \\ \relax
2:43 & \begin{tabularx}{0.7\textwidth}{X} 守滿了節期，他們回去，孩童耶穌仍舊在耶路撒冷。他的父母並不知道， \end{tabularx} \\ \\ \relax
2:44 & \begin{tabularx}{0.7\textwidth}{X} 以為他在同行的人中間，走了一天的路程才在親屬和熟悉的人中找他， \end{tabularx} \\ \\ \relax
2:45 & \begin{tabularx}{0.7\textwidth}{X} 既找不著，就回耶路撒冷去找他。 \end{tabularx} \\ \\ \relax
2:46 & \begin{tabularx}{0.7\textwidth}{X} 過了三天，他們發現他在聖殿裡，坐在教師中間，一面聽，一面問。 \end{tabularx} \\ \\ \relax
2:47 & \begin{tabularx}{0.7\textwidth}{X} 凡聽見他的人都對他的聰明和應對感到驚奇。 \end{tabularx} \\ \\ \relax
2:48 & \begin{tabularx}{0.7\textwidth}{X} 他父母看見就很驚奇。他母親對他說：「我兒啊，為甚麼對我們這樣做呢？看哪，你父親和我很焦急，到處找你！」 \end{tabularx} \\ \\ \relax
2:49 & \begin{tabularx}{0.7\textwidth}{X} 耶穌對他們說：「為甚麼找我呢？難道你們不知道我應當在我父的家裡嗎？」 \end{tabularx} \\ \\ \relax
2:50 & \begin{tabularx}{0.7\textwidth}{X} 他所說的這話，他們不明白。 \end{tabularx} \\ \\ \relax
2:51 & \begin{tabularx}{0.7\textwidth}{X} 他就同他們下去，回到拿撒勒，並且順從他們。他母親把這一切的事都存在心裡。 \end{tabularx} \\ \\ \relax
2:52 & \begin{tabularx}{0.7\textwidth}{X} 耶穌的智慧和身量，並神和人喜愛他的心，都一齊增長。 \end{tabularx} \\ \\
[1ex]
\hline
\hline
\end{longtable}
講題:孩子不見了　The Child Is Missing

經文:路加福音2章41-52節

講員 :  林豪恩先生

場合:香港中文大學崇基學院禮拜堂

日期:2012年12月30日

一．引言

按路加福音的記載,耶穌主動地做的第一件事情,就是「玩失蹤」.「玩失蹤」,對孩子來說可以是沒有甚麼大不了的事情,但「孩子不見了」,對於父母來說,卻是驚心動魄的事情,特別是當社會充滿著拐子黨疑雲的日子.

2012年12月17日早上10時許,兩歲女童隨母親到伊利沙伯醫院看病期間,女童遭一名戴口罩男子強行抱走乘的士逃去,事隔半小時,女童被發現遭遺棄附近佐敦商廈梯間,未有受傷.警方調查發現事件疑屬大婆,二奶爭寵的家庭糾紛,同日拘捕56歲姓蕭婦人(大婆)及通緝在逃男子.【明報專訊18/12/2012】

2012年11月28日下午四時許,姓葉婦人帶同姓李兒子,到石硤尾街市對面巴域街買sung ,期間兒子站近,轉眼間失去蹤影,慌張尋子.當婦人四處張望之際,發現一名婦人拖著其子,於是急忙高呼愛兒名字,婦人即放手.葉婦追前截停婦人查問,對方稱:「仲以為係我gwo 仔添,拖錯jor ,拖錯jor !」隨即離去.葉婦疑遇上拐子婆決定報警.警員翻看附近店舖閉路電視,發現其子主動拉陌生女子,相信是誤會.【太陽報28/11/2012】

今天的福音經課也提到一件發生在二千多年前的「孩子不見了」事件.不見了的孩子名叫耶穌.他為甚麼失蹤呢?他的父母約瑟和馬利亞對於孩子不見了有甚麼反應呢?耶穌的失蹤以及他的父母的反應可以給我們甚麼啟迪呢?

二．釋經與應用

2 : 41每年逾越節,他父母都上 耶路撒冷去.

2 : 42當他十二歲的時候,他們按著過節的規矩上去.

2 : 43守滿了節期,他們回去,孩童耶穌仍舊在 耶路撒冷.他的父母並不知道,

2 : 44以為他在同行的人中間,走了一天的路程才在親屬和熟悉的人中找他,

2 : 45既找不著,就回 耶路撒冷去找他.

2 : 46過了三天,他們發現他在聖殿裡,坐在教師中間,一面聽,一面問.

2 : 47凡聽見他的人都對他的聰明和應對感到驚奇.

2 : 48他父母看見就很驚奇.他母親對他說:「我兒啊,為甚麼對我們這樣做呢?看哪,你父親和我很焦急,到處找你!」

2 : 49耶穌對他們說:「為甚麼找我呢?難道你們不知道我應當在我父的家裡嗎?」

2 : 50他所說的這話,他們不明白.

2 : 51他就同他們下去,回到 拿撒勒,並且順從他們.他母親把這一切的事都存在心裡.

2 : 52耶穌的智慧和身量,並上帝和人喜愛他的心,都一齊增長.

路加福音的作者透過上述經文把「耶穌不見了」的案情交待得清清楚楚.耶穌是如何失蹤的呢?他過完節後,留了在耶路撒冷而沒有跟隨鄉親父老一齊返回家鄉.馬利亞和約瑟在甚麼地方找到失蹤了的耶穌呢?他們在聖殿找到耶穌.耶穌為甚麼失蹤呢?他在聖殿與聖殿教師討論得非常投入而樂而忘返.失蹤孩子耶穌被找到之後有甚麼反應呢?他跟隨父母回家.

這件「孩子不見了」的事件到此完滿解決,作者的總結:「他就同他們下去,回到 拿撒勒,並且順從他們．．．耶穌的智慧和身量,並上帝和人喜愛他的心,都一齊增長」似乎給人大團圓結局的感覺.不過,當你細心留意孩子耶穌與父母的對話,以及母親馬利亞的反應,不知道有沒有給你一種特別的感覺?這種感覺不容易形容,如果說「得到佢個人,但得不到佢個心」又似乎誇張了一點.然而,這位不見了三天被尋回的兒子,難免令父母覺得有點不同,甚至感到有點陌生,有點驚奇,與失蹤前的兒子不再一樣了.

經文如此描述:「過了三天,他們發現他在聖殿裡,坐在教師中間,一面聽,一面問.凡聽見他的人都對他的聰明和應對感到驚奇.他父母看見就很驚奇.」其他人對耶穌感到驚奇不足為奇,因為他們不認識耶穌；然而,連耶穌的父母都感到驚奇,那就真的奇怪了.此時此刻的耶穌,似乎與父母一直認識的耶穌有所不同,這些不同,父母也未必能夠立即理解得到,因此只能感到驚奇.

另一點令父母感到不再一樣的地方,是耶穌的回應.馬利亞找到耶穌後,對他說:「我兒啊,為甚麼對我們這樣做呢?看哪,你父親和我很焦急,到處找你!」符合期望的回應可能是:「噢,對不起,令你們擔心.我與聖殿老師討論得太投入,忘記了跟你們回家.」然而,耶穌的回答竟然是:「為甚麼找我呢?難道你們不知道我應當在我父的家裡嗎?」耶穌的回答並沒有體諒父母心急如焚的心情,反而有點拙劣,然而,卻又令人感到煞有介事,不能隨便否定他.因此,作者指出:「他所說的這話,他們不明白」.

孩子找回來了,但似乎已經不是原來的孩子.面對著這位不再一樣的孩子,令父母感到驚奇和陌生的孩子,父母會否仍然有「孩子不見了」的感覺呢?

有一位小男孩升上小學一年級.媽媽帶他到課室,他不怕陌生,；禮貌地向老師行禮,然後快速地找到自己的座位,他好像還馬上跟旁邊的同學傾談,不費力便投入小學生的生活.他甚至好像忘了跟媽媽說「再見」,這讓媽媽有點失落.升讀小學不久,兒子好像變了樣子,他開始每事問.媽媽與他傾談,他舉一反三；媽媽教他做功課,他總是說不.他不是拒絕做功課,而是不聽媽媽的指示.他有一句口頭禪,愛說:「老師話唔係咁!」.有時,他會不滿意自己的字體,因為「老師話唔係咁」,然後他會用橡皮膠把寫好的字都擦掉.媽媽有點不耐煩,怎麼聽話的兒子這麼愛駁嘴?那麼反叛?媽媽開始擔心.(Smart Parents, 139)

這位媽媽可能也有「孩子不見了」的感覺.她不見了的,是從前那位凡事聽她指示,凡事以她為權威的孩子；她不見了的,是那位從前凡事依賴她,不會駁嘴的孩子.有些父母面對著孩子在心理上和知識上與自己的「分離」,感到擔心,就好像故事中的那位媽媽一樣.有些父母視孩子與自己不同的想法和意見為挑戰自己,因而感到憤怒；有些父母視孩子把父母的權威轉移給老師視為挫敗,因而在孩子面前批評老師,企圖奪回在孩子心中的權威.耶穌的父母如何回應這位對他說:「為甚麼找我呢?難道你們不知道我應當在我父的家裡嗎?」這些聽似有些大道理,但又令人感到被挑戰的說話呢?聖經如此寫道:「他所說的這話,他們不明白.」縱使父母不明白,但他們沒有因為自己的不明白而貶低孩子的想法,也沒有因為自己的不明白而否定孩子的想法.耶穌的媽媽的心胸似乎廣闊得丞相肚內可撐船,聖經說:「他母親把這一切的事都存在心裡」.馬利亞以廣闊的胸懷包容著他的兒子的改變,這些改變,有些父母或者會視之為反叛,有些父母或者會視之為成長.

孩子不見了,又豈只發生在兒童及青少年時期呢?往後,馬利亞繼續有「孩子不見了」的經歷,這些經歷一次比一次激烈和深刻.

馬太福音

12 : 46耶穌還在對眾人說話的時候,不料,他母親和他兄弟站在外邊想要跟他說話.

12 : 47有人告訴他:「看哪!你母親和你兄弟站在外邊,想要跟你說話.」

12 : 48他卻回答那對他說話的人,說:「誰是我的母親?誰是我的兄弟?」

12 : 49於是他伸手指著門徒,說:「看哪,我的母親,我的兄弟!

12 : 50凡遵行我天父旨意的人就是我的兄弟,姊妹和母親.」

少年時代的耶穌,在聖殿裡和聖殿教師談得興起,快樂不知時日過,成為了失蹤少年.成年後的耶穌,四出傳道,與家人聚少離多.這次,馬利亞帶同耶穌的兄弟到他工作的地方去找他,可能希望一家人有機會團聚.然而,耶穌的回應,可能比上次更加令馬利亞感覺到「孩子不見了」.上次,耶穌那句「為甚麼找我呢?難道你們不知道我應當在我父的家裡嗎?」令馬利亞疑惑,令馬利亞思前想後,令馬利亞感到兒子開始有自己的想法,再不是凡事以父母為標準的兒童.今次,耶穌的回應,有著重新定義家庭的意味,直接衝擊著馬利亞與耶穌的母子關係,因為耶穌說:「誰是我的母親?誰是我的兄弟?．．．看哪,我的母親,我的兄弟!凡遵行我天父旨意的人就是我的兄弟,姊妹和母親.」

如果你是馬利亞,你會有甚麼感受呢?「是誰懷胎十個月生你出來的呢?你是吃誰的奶長大的呢?是誰棉乾絮濕,含辛茹苦地把你養育成人呢?現在,你竟然問:誰是我的母親?你竟然說凡遵行天父旨意的人就是你的母親?」

耶穌重新定義家庭,重新定義誰是家人,不等於否認馬利亞是他的肉身母親,不等於其他女人可以取代馬利亞作為他的媽媽的關係,不過,這種家庭的重親定義,難免令耶穌的媽媽有失去和失落的感受.媽媽還是孩子心目中最重要的女人嗎?還是其實媽媽已經失去了孩子,孩子其實已經落入了其他人手上呢?孩子長大成人,是父母的期望；孩子成家立室,也是父母的安慰.不過,讓孩子成家立室時重新定義家庭和親密關係,卻不是容易的事情,甚至是痛苦的糾纏.

有一個男人名叫小李,他的母親承襲了一套家規,其中形響小李最大的有兩條:1.家務不是男性的工作,必須由女性負擔,因母親相信「下廚的男人沒出息」；2. 「自己人」與「外人」必須界限分明.自己人的意思是有血緣關係的人,彼此應團結愛護.媳婦嗎?屬外姓,因此是外人,不屬家庭一分子.母親叮囑小李不要,也不應對妻子好,因為「媳婦熬成婆,需要二十年」,入門不足二十年的媳婦,仍然「養吾熟」,一有機會,仍會溜.

當孩子出生後,小李的困擾出現了.小李很想照顧妻子和嬰兒,很想協助做家務,這是他對母子倆的愛的自然流露.然而,他又要面對自己母親的壓力.每次母親在家,都會制止他做家務,並因家務問題與他發生衝突,隨後更遷怒於媳婦.母親有時還會介入照顧和管教小李的兒子.經過一番商議之後,小李兩夫妻為減少與母親磨擦的機會,決定採取陽奉陰違的方式,在母親面前妻子扮演「小女人」的角色,小李亦絕不做任何家務；兩口子更日間盡量外出,晚上夜歸.他們認為,與母親的見面機會減少,衝突也會相應減少.但這樣做,他們心裡又感到內疚,對不起自己的孩子.況且,這樣安排,在母親居住時間短的情況下,還可以勉強應付.但是,母親與他們居住的時間卻愈來愈長,次數也愈來愈多,介入家庭生活也愈來愈深,連吃飯也不容許小李為忙著照顧小孩的妻子預留一點菜；只讓小李愛護母親和兒子,不讓他愛護妻子.最後,妻子終於忍無可忍,表示這樣繼續下去,會帶同兩個兒子離開李家.

小李尋求輔導員協助.輔導員告訴他不應太懦弱,應盡丈夫保護妻子的責任,直接告訴母親,指出她的家規如何不配合他自己對家庭的期望,且影響他的家庭生活.小李雖然痛恨母親,恨她不體諒他,甚至要控制他,破壞他的家庭幸福,但他仍然覺得這位輔導員太不了解他的處境了.作為孝順兒子,他怎能直接指出母親的不是呢!直接向母親說出心裡感受,不就等如指責和教訓母親嗎?同樣,作為孝順兒子,怎能提出請母親減少來家居住的次數和時間呢?(張包意琴:「變,從家規著手」)

耶穌的媽媽馬利亞如何回應她的兒子重新定義家庭呢?福音書沒有跟進這事件,因此我們沒有確實的資料回答.不過,從福音書記載耶穌往後的事蹟所見,馬利亞似乎沒有像上述的那位母親一樣介入兒子的生活,也沒有不能自制地操控著兒子的計劃.馬利亞當時未必感受良好,我們也不知道她如何能夠克服自己的擔心和和放下自己的期望,放手讓耶穌去接觸他所碰見的人,去與販夫走卒為伍,去走遍大城小鄉.放耶穌走,馬利亞不怕兒子一去不回頭嗎?容許耶穌相信「凡遵行我天父旨意的人就是我的兄弟,姊妹和母親」,馬利亞不怕兒子心中從此沒有了自己這位親生媽媽了嗎?讓我們繼續看第三個馬利亞不見了兒子的片段.

約翰福音

19 : 25站在耶穌十字架旁邊的,有他的母親,姨母, 革羅罷的妻子 馬利亞,和 抹大拉的 馬利亞.

19 : 26耶穌見母親和他所愛的那門徒站在旁邊,就對母親說:「母親,看,你的兒子!」

19 : 27又對那門徒說:「看,你的母親!」從那刻起,那門徒就接她到自己家裡去了.

從前,因為學業,孩子令父母有「孩子不見了」的感受；從前,因為工作和家庭,孩子又令父母有「孩子不見了」的感受；今次,孩子真的將會不見了.耶穌被釘在十字架上,馬利亞這位母親被賜下「白頭人送黑頭人」的苦杯.為甚麼馬利亞這位母親要喝這苦杯呢?因為她的孩子耶穌要活出他在世為人的意義,要實踐上主給他的使命.要阻止孩子去走他認為有意義的路嗎?要阻止孩子去實踐他的使命嗎?要阻止孩子去回應上主的召喚嗎?走他認為有意義的路,實踐他的使命,回應上主召喚,可能意味著孩子要走的路不是父母期望的路,可能不是平坦舒服的路,甚至可能是父母會擔憂的路.

還記得在中大讀書的時代,我們一班相熟的基督徒同學,常常走在一起事奉和分享.其中一個經常提及的話題,就是畢業後做甚麼事情回應上主的呼召.有一晚,其中一位同學告訴大家,他在某個聚會中感到上帝呼召他作傳道人,投身全職事奉.當我們為他感恩的時候,他卻為此而感到苦惱,因為他知道他的父母不會讚成他的決定.這情況使我們感到奇怪,因為我們知道他的父母也是基督徒,怎會不讚成他做傳道人呢?原來,他的父母早已告誡他,畢業後要找一份理想職業,要以事業為重,有時間的話在教會參與事奉也無妨.我們聽後,也不知道如何回應.當時,這位苦惱的同學拿起結他,唱起當年在教會界流行的齊唱新歌的其中一首,名為「媽媽不要哭」,歌詞是這樣的.

默默回望凝望著遠處　心中不能忘掉您的面

我為尋　我為尋覓曾話別遠去　跟您滿眼帶淚離別

傾出此生從未後悔過　心中激情惟願獻生命

我願能我願能盡獻心中所有　那懼前路坎坷

媽媽不要哭　我知您的心在痛　媽媽不要再流淚

知道令您甚掛慮　仍是要下決定　為了實現抱負使命

媽媽不要哭　可知我的心亦痛　媽媽深盼您明白

為尋求神國度去　衝破艱難

不怕累　不退後　不畏懼

媽媽會失去孩子嗎?孩子會忘記媽媽嗎?不會的.孩子在被創造的時候,已經透過臍帶和媽媽連繫著.孩子出生時,臍帶被剪斷,意味著孩子成為獨立的個體,然而,孩子與媽媽之間卻彷彿有一條無形的臍帶在連繫著.孩子無論成長到甚麼階段,無論往那裡去,心裡面仍然有媽媽,媽媽心裡面也仍然有孩子.耶穌在十字架上慷慨赴義的時候,他的眼睛仍然看著他的母親,他的心裡仍然掛念著母親.耶穌最後一項要親自辦理的事情,就是安排母親的生活,吩咐門徒好好照顧他的母親,對待他的母親好像對像自己的母親一樣.然後,耶穌才說出了最後的話:「成了!」才放心說:「父啊,我將我的靈魂交在你手裡.」馬利亞這位媽媽常常不見了自己的兒子耶穌,但馬利亞從來都沒有失去耶穌這位兒子.

三．總結

我們都必定是父母的孩子,有些同時也是孩子的父母,我們都體會過父母子女之間那種既緊密又微妙的關係.我們都在某程度上體會到孩子成長,邁向獨立的過程中所產生的與父母關係的張力.想到此,我們感激馬利亞,因為她廣闊的胸襟,接納著他的兒子耶穌在不同階段,不同程度的離開,容許他的兒子耶穌去追尋並實踐上帝在他身上的旨意.我們不能夠計算清楚馬利亞在我們的得救中的功積,但我們能夠確定馬利亞在耶穌生命中的參與,是與我們的得救有關的,是與我們得著的祝福有關.耶穌的母親馬利亞,請接受我們向你說聲感謝.但願天下為人父母的,因著我們在我們孩子生命中的參與,在將來的某一天,都能夠聽到受到我們的孩子祝福的人來向我們說聲多謝,感激我們讓我們的孩子成為他們的幫助.



\newpage

\section{奧秘揭開了}
\label{sec:180}
\textbf{鄧瑞強博士 - 奧秘揭開了}
\newline
\newline
\hyperref[sec:179]{< < < PREV SERMON < < <}
~
\hyperref[sec:toc]{[返目錄]}
~
\hyperref[sec:181]{> > > NEXT SERMON > > >}


以弗所書 3:1-12
\newline
\begin{longtable}{cl}
\hline
\hline
章節 & 經文 (和合本修訂版)\\
\hline
3:1 & \begin{tabularx}{0.7\textwidth}{X} 因此，我—保羅為你們外邦人作了基督耶穌囚徒的，替你們祈禱。 \end{tabularx} \\ \\ \relax
3:2 & \begin{tabularx}{0.7\textwidth}{X} 想必你們曾聽見神賜恩給我，把關切你們的職分託付我， \end{tabularx} \\ \\ \relax
3:3 & \begin{tabularx}{0.7\textwidth}{X} 用啟示讓我知道福音的奧祕，正如我以前略略寫過的。 \end{tabularx} \\ \\ \relax
3:4 & \begin{tabularx}{0.7\textwidth}{X} 你們讀了，就會知道我深深了解基督的奧祕； \end{tabularx} \\ \\ \relax
3:5 & \begin{tabularx}{0.7\textwidth}{X} 這奧祕在以前的世代沒有讓人知道，像如今藉著聖靈向他的聖使徒和先知啟示一樣， \end{tabularx} \\ \\ \relax
3:6 & \begin{tabularx}{0.7\textwidth}{X} 就是外邦人在基督耶穌裡，藉著福音，得以同為後嗣，同為一體，同為蒙應許的人。 \end{tabularx} \\ \\ \relax
3:7 & \begin{tabularx}{0.7\textwidth}{X} 我作了這福音的僕役，是照著神的恩賜，是照他運行的大能賜給我的。 \end{tabularx} \\ \\ \relax
3:8 & \begin{tabularx}{0.7\textwidth}{X} 雖然我比眾聖徒中最小的還小，他還賜我這恩典，讓我把基督那測不透的豐富傳給外邦人， \end{tabularx} \\ \\ \relax
3:9 & \begin{tabularx}{0.7\textwidth}{X} 又使眾人都明白甚麼是歷代以來隱藏在創造萬物之神裡的奧祕， \end{tabularx} \\ \\ \relax
3:10 & \begin{tabularx}{0.7\textwidth}{X} 為要在現今藉著教會使天上執政的、掌權的知道神百般的智慧。 \end{tabularx} \\ \\ \relax
3:11 & \begin{tabularx}{0.7\textwidth}{X} 這是照著神在我們主基督耶穌裡所完成的永恆的計劃。 \end{tabularx} \\ \\ \relax
3:12 & \begin{tabularx}{0.7\textwidth}{X} 我們因信耶穌，就在他裡面放膽無懼，滿有自信地進到神面前。 \end{tabularx} \\ \\
[1ex]
\hline
\hline
\end{longtable}
講題:奧秘揭開了　The Hidden Mystery Unveiled

經文:以弗所書3:1-12

講員:鄧瑞強博士(香港神學院)

場合:香港中文大學崇基學院禮拜堂

日期:2013年1月6日

各位弟兄姊妹,早安.

今日是新的一年的開始,希望各人在神的恩典裡,得享喜樂及平安.

大地運轉,冬去春來,大地從嚴寒回復溫暖,死寂回復生機,這容易引發人產生宗教情緒.宗教學者以利亞德(Mircea Eliade)說,新年是宗教節慶的重要時刻.人在新年裡,能宗教地體會萬物原初的狀態.在這原初狀態裡,神和人在和諧的關係裡,人從中獲得存在的意義.這意義給人勇氣與盼望,去克服現實生命中的種種逆境.新年的時候,我們喜歡許願.許願,可算是對這原初狀態的回憶,是失去樂園後對樂園的回憶,重尋樂園.在許願裡,我們祈求健康,這是對永生的一種渴望.祈求事事如意,這是對樂園裡萬物和諧的一種回憶.人會許願,是因為我們是宗教的動物.我們活在現實裡,但渴望超越現實的永恆.

在這新年期間,讓我們再思想一下,生命的真相是什麼?我們可以渴望什麼?我們在其中扮演的角色又是什麼?

今日講道的經文,是今日經課裡新約書信的經文:以弗所書3:1－12

3:1　因此,我─保羅為你們外邦人作了基督耶穌被囚的,替你們祈禱.

3:2　諒必你們曾聽見神賜恩給我,將關切你們的職分託付我,

3:3　用啟示使我知道福音的奧祕,正如我以前略略寫過的.

3:4　你們念了,就能曉得我深知基督的奧祕.

3:5　這奧祕在以前的世代沒有叫人知道,像如今藉著聖靈啟示他的聖使徒和先知一樣.

3:6　這奧祕就是外邦人在基督耶穌裡,藉著福音,得以同為後嗣,同為一體,同蒙應許.

3:7　我作了這福音的執事,是照神的恩賜,這恩賜是照他運行的大能賜給我的.

3:8　我本來比眾聖徒中最小的還小,然而他還賜我這恩典,叫我把基督那測不透的豐富傳給外邦人,

3:9　又使眾人都明白,這歷代以來隱藏在創造萬物之神裡的奧祕是如何安排的,

3:10　為要藉著教會使天上執政的,掌權的,現在得知神百般的智慧.

3:11　這是照神從萬世以前,在我們主基督耶穌裡所定的旨意.

3:12　我們因信耶穌,就在他裡面放膽無懼,篤信不疑地來到神面前.

經文的第一句,讓我想了很久.

「我──保羅為你們外邦人作了基督耶穌被囚的」.保羅,一個嚴守律法的猶太人.「為你們外邦人」,一班與自己全無關係的人,成了監獄裡的囚犯.一個人為了一些與自己無關的人坐監.

世上有些人進入監獄,是由於非一般的理由的.

最出名的例子,是古希臘的蘇格拉底,他被控「不敬虔」,並以新的道理敗壞年輕人,因而坐監.所謂不敬虔,就是不再信任當時社會的神明.這些神明,也可理解為社會的權威.蘇格拉底因為尋求真理而抵觸了這些權威,這是死罪.在監獄裡,他的朋友願意幫助他逃獄,但蘇格拉底堅持不逃獄.他認為身在雅典,便必須遵守雅典的法律.法律是正義的.他不能因為有些人誤用法律而違反法律.最後,飲下一杯毒酒,蘇格拉底便死在獄中.他是典型地為了正義的原則而死的人.

德國神學家潘霍華(Dietrich Bonhoeffer)因為與刺殺希特拉有關而入獄,在獄中寫下《監獄書信》,日後成了重要的神學文獻.他為何入獄?他認為作為基督徒,便要承擔民族及歷史的命運,為了他人的幸福,必須放下自己,盡一切力量,阻止悲劇.他最後死在集中營.

前南非總統曼德拉為了終止種族隔離政策,坐了27年監.他說:“For to be free is not merely to cast off one’s chains, but to live in a way that respects and enhances the freedom of others.”(「自由不是單單脫下鎖鏈,而是活出尊重及增進他人自由的方式.」)他就是為了別人的自由,他放下了自己的自由.

中國有劉曉波,他也在獄中.他的自辯書,題為＜我沒有敵人＞,裡面有句話說:「我期待,我將是中國綿綿不絕的文字獄的最後一個受害者,從此之後不再有人因言獲罪.」

世上還有很多這些人,他們被投在監裡,我在此不再列舉了.

他們都是體驗到宇宙人生裡某種真理的人.他們認為這真理是重要的.這真理牽涉人類存在的意義.他們堅持這真理,甚至不惜為此被囚.值得嗎?看到真相的人,會認為值得.他們的命運告訴我們,越接近真理,承擔揭露真相的責任越大.越願意承擔這責任,要承受的苦難便越大.要讓人得自由,有時必須失去自己的自由.要讓人得生命,自己必須進入死亡.

保羅,看到生命的真相,也願意為這真相被囚.他稱這被囚是「作了基督耶穌被囚的」(prisoner of Christ Jesus).看來,他視這被囚,是分享了基督的生命和使命.這不是失敗,而是光榮,顯示他正履行神給他的使命,成就神的計劃.一個人明白了真理,便應該承擔這真理的責任,勇往直前.耶穌在這罪惡的世界,活出了非暴力的愛,死在十架上.信仰耶穌的人,也當如此行.

保羅看到什麼真理,令他甘願被囚?

3:4　你們念了,就能曉得我深知基督的奧祕.

3:5　這奧祕在以前的世代沒有叫人知道,像如今藉著聖靈啟示他的聖使徒和先知一樣.

3:6　這奧祕就是外邦人在基督耶穌裡,藉著福音,得以同為後嗣,同為一體,同蒙應許.

這裡有一個重要的字眼:「奧秘」.因著神的啟示,保羅看到宇宙人生終極的奧秘.這奧秘不是什麼不為人知的秘密,這奧秘是神要世人明白的事情.所謂「奧秘」,就是神的心意,神的計劃.這心意,這計劃,不是隱藏的,不需要人去猜的.神是啟示的神,他會用盡千方百計讓我們明白他的心.神最終極的啟示,就是他來到人間,成為人.早前的聖誕節,便是記念神來到人間的事實.來到人間,為了承擔世人的罪,基督懷著愛與寬恕,他死在十字架上.他為了我們得到自由,他甘願失去自由.他為了我們明白寬恕,他忍受咒詛.他為了我們學曉真愛,他承受仇恨.他觸摸痲瘋病人,藉此告訴我們,這世界沒有人是不可觸摸的.他與稅吏吃飯,藉此告訴我們,這世界沒有人是斷然地被孤立的.他抱起小孩子,一個沒有地位沒有成就的小孩子,然後教導門徒,要學識尊重這樣的一個一無所有的人的尊嚴.他最後走向十架,藉此告訴我們,真正愛他人,是願意為他人犧牲的.這沒有什麼損失,這是生命的豐盛.死亡之後,便是復活.

保羅說:

3:5　這奧祕在以前的世代沒有叫人知道,像如今藉著聖靈啟示他的聖使徒和先知一樣.

3:6　這奧祕就是外邦人在基督耶穌裡,藉著福音,得以同為後嗣,同為一體,同蒙應許.

基督帶給我們的好信息,就是所有人(無論膚色,性別,人間的成就)都在神的恩典裡,同為後嗣,同為一體,同蒙應許.所有人都同為神的後嗣,所有人都可被稱為神的兒女.所有人同為一體,所有人情同手足.所有人同蒙應許,所有人都能享受天國的福樂.

請不要說,你的膚色跟我不同,我高級一點而你低級一點.請不要說,我做經理而你掃地,我高級一點而你低級一點.請不要說,我的財富比你多,我高級一點而你低級一點.

當神看到每個人的無限價值時,我們怎能輕視任何一個人?

當有人需要食物時,請不要說,這不關我的事.當你的同事需要幫助時,請不要說,這是你自己的問題.當你的伴侶需要你的愛護時,請不要說,你先好好檢討一下你自己.

當基督在我們還作罪人的時候便為我們死,好顯明什麼是愛,我們怎能不愛我們身邊的人?

面對基督,保羅頓然明白天地間最大的奧秘.在弗3:4裡,保羅稱之為「基督的奧祕」.原來,我們自己並不是宇宙的中心,我們只是神偉大創造的一個微少部分.原來,我們的生命被神的大愛托著.原來,當這大愛被我們實現出來時,人生的全幅意義才會顯明.

明白這奧秘的人,就不會時時以自己為中心,就不會時時只顧自己的事,就不會時時為自己的吃喝玩樂而煩惱.

有人說,態度決定高度.我說,視野決定深度.

你看得夠遠,你才明白你的位置在哪裡,你的天職在哪裡,你的深度在哪裡?

3:7　我作了這福音的執事,是照神的恩賜,這恩賜是照他運行的大能賜給我的.

3:8　我本來比眾聖徒中最小的還小,然而他還賜我這恩典,叫我把基督那測不透的豐富傳給外邦人,

3:9　又使眾人都明白,這歷代以來隱藏在創造萬物之神裡的奧祕是如何安排的,

3:10　為要藉著教會使天上執政的,掌權的,現在得知神百般的智慧.

3:11　這是照神從萬世以前,在我們主基督耶穌裡所定的旨意.

3:12　我們因信耶穌,就在他裡面放膽無懼,篤信不疑地來到神面前.

電影＜蜘蛛俠＞裡面,蜘蛛俠的叔叔Uncle Ben 對蜘蛛俠說:「偉大的能力伴隨著偉大的責任.」(With great power comes great responsibility)有能力的人,面對邪惡的世界,他的能力會要求他承擔更大的責任.當然,他也要排除萬難,願意承擔才行.

保羅在這裡講,「我作了這福音的執事(僕人),是照神的恩賜,這恩賜是照他運行的大能賜給我的.我本來比眾聖徒中最小的還小,然而他還賜我這恩典,叫我把基督那測不透的豐富傳給外邦人.」

保羅自認為卑微的人,但神的恩典臨到他,讓他明白神的奧秘,讓他看到生命的真相,讓他知道生命到底為何而活.恩典伴隨責任.真是天降大任,他成了福音的僕人.視野決定深度,他的視野要求他活出基督徒的深度.

從今以後,他就竭盡所能,促成不同群體在基督裡的和諧,讓不同的人在基督裡找到豐盛的生命,讓所有人明白什麼是生命的大愛.

看到了生命的奧秘,明白這奧秘的重要,他為此願意付出一切.他願意為此被囚嗎?他願意.

在這世上,很多人視野遠大,看到人生的真相,他們也願意為此被囚.他們的視野,決定了他們的深度.但願,這些人早日獲得釋放.

在這新年期間,讓我們看得更遠大,讓我們的新年願望配合這遠大的視野,而不只是「恭喜發財」,讓我們有勇氣為我們看見的奧秘盡上一分力量.

願主的恩典臨到你們.

但願:榮耀歸於聖父,聖子,聖靈.阿們.



\newpage

\section{主都知道}
\label{sec:181}
\textbf{陳偉光教授 - 主都知道}
\newline
\newline
\hyperref[sec:180]{< < < PREV SERMON < < <}
~
\hyperref[sec:toc]{[返目錄]}
~
\hyperref[sec:182]{> > > NEXT SERMON > > >}


日期:2013年1月13日

題目:主都知道

證道:陳偉光教授

今天好高興來到大家當中,我想問一問各位朋友,各位主內弟兄姊妹,你們都好嗎?是否都很好?今天很早就起來祈禱,讀聖經,今天我很早就來到這會堂,我坐在後面那裡,我抵達的時候只有詩班在練歌,並沒有太多朋友到臨,每一位進來的朋友,我都逐一為他們祈禱,我知道主有很特別的恩典要臨到在各人身上,何解? 有人說,今天是星期日,星期日是主日,許多時候,在崇拜的時候主席宣告神的話語,宣告聚會開始說:``This is the day that The Lord have made, we will rejoyce and be glad in it".今天是耶和華訂的日子,我們要在其中歡喜快樂,恍惚今天有什麼特別. 那今天是否特別?難道明天不是耶和華定的日子嗎?  昨天也是吧!一息間完了崇拜都是,其實每天都是一樣,所以每一刻都是那麼的珍貴,但我相信每一刻神都有特別,很意想不到的恩典,臨到我們每一個人身上,所以我們彼此代禱,求上帝憐憫我們這一群可憐的人.

我剛才說我來到是懷著很興奮,不單興奮,而且帶著一顆殷切的心情來跟大家分享上帝的話語,或是心裡的一些感動,我並非說客套說話,我真的很感動和殷切的,何解?因為我對這會堂有很特別的感覺,我在這裡唸書,在這裡工作,其實已是卅幾年了,星期天我不多在這裡,除了二十多年前我初來中大做事,我曾經在大學裡居住,有兩三年的時間,我曾經在大學做organist, 那是許久以前的事情,大概廿多年前了,在此之後基本上我都是平日星期一至星期五,又或甚至是星期六,才有機會來到這會堂,崇基週會都會在此舉行,感覺很特別,能夠在這個崗位上服事上帝,與大家彼此造就,是從來沒有試過,是我一生人的第一次,所以我懷著非常,非常殷切的心情,儘管是這樣說,但我必須要承認,其實我都懷著非常沉重和無奈的心情,今早來到這裡,何解?其實越殷切,便有越多的期望,我過去這幾個月時常都記著十一月十八日我便要來到這地方,和弟兄姊妹彼此交通,但苦於一直想著這件事情,想至昨天晚上,我一直都沒有機會好好地停下來,想一想,真是這樣,是一種很大很大的無奈! 如果用屬靈點的語氣來說,就是有許多攔阻,的而且確非常多的攔阻,我心裡有一番感動,大概是有個想法要與大家分享,事實上我亦寫了一段短短的文字,為了我們即時傳譯的同事,他看見我寫了一段文字,但很抱歉,一息間那段文字可能不會應用,因我大概知道有這些想法,但我未曾知道,老實說,我在等通知,等上帝通知我有什麼要說,但等待至昨晚我亦沒有機會好好地坐下來,認真的想一想,這是很真實,我在想,很奇怪! 我與大家分享的題目是「主都知道」,其實人生便是如此,往往都是事與願違,力不能勝,苦無出路,有種種原因,令你每日想做的事情,可能都沒有機會做,總的來說,幾十年的歲月,可能都是煩悶的感覺,悶得要吐,煩悶的感覺,總是沒有什麼出路,但正正就是經常落在這種狀態中,便知道你自己是何等需要上帝?!而且一切沒有辦法講得明白,非常複雜,自己解決不了,自己不能整理,非常凌亂的思緒,主都知道!我覺得好奇怪,讓我好好的體驗,今早來到述說一下我個人內心世界的分享,其實我自己都是經常地,甚至今次在這具體事情上,經歷著這樣的感覺,很多時候我們教書的人,一講便停不了,我現在說到引言.

其實本星期真是要好好的預備的時候,我有一個意想不到的事情發生,上星期六我到教會詩班練歌的時候,我們教會詩班的副指揮帶了一首歌來,叫 ``Ascribe to The Lord",其實今天夾在各位的程序表裡.這是位很好的指揮,一位年青的作曲家,他帶這首曲來唱.一唱,我心裡忽然間勾起一段回憶,很多年前的回憶,甚至可以形容是波濤洶湧的一種思緒,久久難平.這其實是 marana music 三十多年前在北美,怎至傳到香港,非常之流行的一個階段經常唱的歌,所以在我唱這首歌,除了上星期六我們的副指揮帶來唱的時候,我想至少有十五年我沒見過這首歌,所以當他一帶這歌來的時候,那感覺是很強.

我方才說我教會那指揮帶了那首歌過來,他是一個很好的作曲家,他完成了作曲博士學位,剛剛就巡迴演出了幾個大的作品,他很有心幫助我們教的詩班編新歌,他將那歌編成一個很精緻的合唱版本,是很美,比起大家看見這一頁的歌更美,他是我的學生,他排練完後問我說:「Victor, 你覺得這首歌怎麼樣?」我二話不說劈頭就說:「吶悶得要吐!」「 悶得要抽筋!」當然他嚇了一跳,我有心讓他一下子覺得尷尬!跟著我告訴他為什麼我覺得「悶到嘔」「悶得要抽筋」的因由,因為那是個人的感受,這首歌對我來說是很沉重的感覺.

大概十五,六年前,我的事業比較中段的階段,前段與中段中間,我很記得這首歌,每天的晚上,那時候我的女兒一兩歲左右,很難帶的小孩子,每晚都會哭到「吐」,我們夫婦倆都沒能睡得好,所以說倘若她能終於能睡著,是一件很痛快的事情,睡著的她也隨時會醒來,小朋友真是很奇怪,這``Ascribe to The Lord", 這首歌有一個很美的錄音版本,只要一聽這首歌,她就睡得不錯,於是我們兩夫婦便決定播給她聽她,一睡了便播,一播便要播上三,四十次,重複的 loop 著播,你知道這類歌,咕咕噥噥的,一個晚上要聽三,四十次,一聽就一兩年,你說想不想吐? 而且那些感覺都跑出來了,其實我自己是一個作曲的人,全家人睡了,才有機會定下來作曲,作曲並非一件很過癮的事,是不太浪漫的事情,是一件很孤獨,很沉悶,完全沒有助手,而且有很多細節要處理,有很多 tedious 的 details, endless tasks,永遠都完成不了的任務.不能睡,一邊聽著 ``Ascribe to The Lord",咕咕噥噥在隔壁的房間還播著,然後我在書房內作曲,一作便至天亮,那段日子,現在回想也頗淒涼,如此也有十年的光境,每星期平均睡三個晚上,其餘四個晚上工作到天亮,洗洗臉便上班去了,你說這是否十分離譜? 其實一點也不! 我有不少朋友也許沒有我當年的誇張,如此的撤底,但是都是人生的另一種基本狀態,這時代永遠都完成不了的任務(假如你曾經認真過的話),你有沒有這樣的經驗?或者沒有這樣的過分,但你總有永遠完成不了的任務,或者不容易過的日子?你們有點像不大願意回答的樣子!我想大家還沒預備好回答這個問題!假如你真是沒有關係,我鼓勵你回去好好的檢討,即便說你自己做人並不十分認真,一個跟隨基督的人,一個好的 Christ follower,你認認真真處理手上的任務,你總會有一些艱澀,有些不容易渡過的日子的感覺,總會如此,因為我們生活在這個時代,你要好好完成交到你手上的事情,但主的恩典好像不夠用,感覺是不容易,所以這好 repetitive 的歌,有一種沒有人明白的感覺.其實老實說,我怎會明白?我自己又搞不清楚自己,一塌糊塗的,還要等待有人明白,總希望有人會明白自己那十多年,或許那段日子會過去,但仍有另外一些日子我不容易明白的感受.

其實這個星期來之前,我們教會發生了一件事情,更加添了我一番情緒,我教會有一位年老的姊妹名叫「盛嫂」,她的先生名叫「盛哥」,盛嫂忽然間去世,一早起來洗過臉,跟他兒子說我今早有點不舒服,頭暈暈的,隨即叫救護車,救護車到來時她已經去世了.一位我們非常尊敬的老姊妹,故事背後很長,盛嫂的丈夫盛哥在教會服事了一生人的堂役,退休的時候,大家都說兩夫婦辛苦了一輩子,一生為教會工作,滿以為可以好好的休息,現在兒女也長大了.在他退休當日,就是他中風的那一天,中了風到現在廿年了.廿年來都躺在醫院,全身動彈不得,他每天都等盛嫂到都醫院裡,可以服侍他,照顧他的飲食,他很希望有一天他能早點讓天父接回天家,好讓盛嫂不用那麼辛苦,不過忽然之間盛嫂都離開了.我們大家都很憂心忡忡,帶著這噩耗到醫院告訴盛哥,牧師拿起詩篇23篇讀「耶和華是我的牧者」想安慰他,第一節還沒唸完半句,盛哥便大大聲聲的繼續唸下去,「耶和華是我的牧者,我必不至缺乏」他一直的唸下去,他唸的時候大家都覺得很慚愧,因為他心裡所相信的,他心裡那種平安,遠比想帶著安慰而來的人更加堅實.有時候我們帶著安慰的心,真淒涼!真可憐! 真不幸! 為什麼神的恩典是這樣在他們身上?我想我們真是要好好想想,基督徒要有亮麗的人生,有誰彼得上盛哥,盛嫂? 一生服侍,像是沒有看見什麼實質的回報,但他仍然在心裡堅實的依靠上帝,是否非常的亮麗呢?你夠膽量去安慰他?來不及告訴他們,這是我們無比欣羨的人生,主的恩典僅僅夠用,剛剛夠用,因為太多你會以為自己很了不起,就是在這樣日子裡,他堅心依靠上帝,不過畢竟人就是有很多限制,就像我們這些軟弱的人,我們便帶著安慰過去,然後我在想像,他的四個兒女,不斷有人來安慰他們,不知怎樣說,不知怎樣回應,其實心裡十分尊敬自己父母,我的父母有著最亮麗的基督徒的歲月,漂漂亮亮地走過來,可能人家安慰的時候你會道謝,其實是一些很複雜,很複雜的思緒,好 ``subtle",非常微妙的感覺.

活在這個世界上,如果你作為一個追隨上帝的基督徒,你心裡有一種苦澀,就是不會有太多人明白,甚至那些越是矢志要跟隨上帝的人,心裡越是艱澀,因為越少人明白當你跟隨主的時候,你過的日一點也不好過,但是沒人明白,不是太多人知道,唯有主知道! 我一直在想這些感覺,這星期發生的事情,唱一首舊的詩歌,見到一位年老我們尊敬的年老姊妹忽然間離開,都湧起這些思緒,心裡難以平復的情緒湧現.

這一個月我自己在讀傳道書,我在想聖經裡這個煩悶人生的感覺.傳道書可以說寫得最透徹,傳道書說虛空 ``smoke", ``spitting into the wind",完全是無聊的事情,很虛空.尤其是如果大家有興趣,可以打開傳道書第一章,傳道書的第一章,你讀一章便夠了,全章的精髓都在那裡!譬如傳道書的第一章,它說到第五節 「日頭出來 , 日頭落下 , 急歸所出之地 .」第六節「風往南颳 ,又向北轉 ,不住地旋轉 ,而且返回轉行原道 .」第八至十節「萬事令人厭煩, 人不能說盡 . 眼看 ,看不飽； 耳聽 , 聽不足 . 已有的事後必再有； 已行的事必再行 . 日光之下並無新事 .豈有一件 事人能指著說這是新的? 哪知 , 在我們以前的世代早已有了 .」你不需再唸下去,那種很 repetitive, 那 endless tasks,一切做完都是枉費的,都是白費,因為一早已有,你總不會好過你的列祖.以利亞極之沮喪的日子就是如此,還奢望自己好過之前的人,其實是不可能的事情,當你想想,一生人努力是為了什麼呢?一個煩悶的感覺.

其實傳道書在繼續的篇章中講出一件很重要的事情,人生的虛空,人生的煩悶,其實都是因為人類,我們說的是人類集體的貪婪,自私,彼此掠奪,結合下來,大家一起要承擔的結果,我想大家都很能看到這真實,特別活在香港,或是選擇在香港生活的我們這一群,香港是一個高度發展,極速進步的一個都市,我們的感受是非常之深,非常之強烈,如果將它再稍微收窄一下,回到我們身處的環境,大家在大學生活的這個社群,更加可見一斑,一個極速的城市,高度發展的時代步伐,從香港再縮影到大學,追求尖端科技,追求知識最卓越目標的一個群體來說,更加感受深刻,今天我們在大學裡,在社會中,一天都晚所高舉的是什麼? 持續增值,終身卓越的理想,請你告訴我,你聽見人家叫你持續增值,終身卓越,你是感到精神抖擻,無比亢奮,拚勁十足,還是惡夢連夜,神經衰弱,不想上班?你是那一種? 為什麼我們社會上那一群最勤奮,最上進,最自愛,最自重的人,都再提不起勁,難道這些全力以赴,做到最好,這些理想,不是應該的嗎?完全正面的目標,非常正確的人生態度,為何聽起來那麼討厭呢?那麼厭煩呢?的確是這樣,每一個人內心深處都藏著一份尊敬,一份對生命本質和它應有的價值的一份渴求,亦正是這份內心的信念告訴我們,如果我們能把手上的工作盡力做好,將人生之中種種不同的角色盡用心來演好的話,應該是一種很幸福的事情,是一個很高貴,很美麗的人生目標,我們會毫不猶豫,全情投入,但如果有人要求我們要全心,全力,全意,全人,全天候,全方位,全所有,整個人傾斜在某一些角色上,我想那絕對是一份過份苛求和苛索了.

我方才說彼此掠奪,永遠不滿足的一種貪婪,如果更有一天,我們撫心自問,已經竭盡所能,全人奉上,心力交瘁,然後那位神氣十足的領禱告訴你,「兄弟,光靠勤力是不夠的,新時代需要更加卓越的人材,要更深,更廣,更闊,更跨領域,跨媒介,集超人本領於一身的那種超人的表現,那你跟我也只得承認,這個時代的那種苛索,霸道,已經到了一種迷失和瘋狂的地步.眾人像我般一口氣說下去,還沒喘過氣來,那貪得毋厭的掠奪者,用文人寫的說話「人性的懶惰」是要有競爭,才有進步；有監察,才有自省!

然後永無休止的評審,這詞語大家很熟悉.評審,考核接踵而來,未曾達標的,請自動離場 and \#8230; 每天我們就是活在這一種不斷自我提升,拼命尋找原因去肯定自我價值的日子裡,所以除了贏過自己是不足夠的,還得勝過別人,總之要雙贏三贏,贏過全人.總之無論怎樣,要贏盡所有.不錯,我們心裡不安,因為藏在我們內心深處,對生命本質的信念通知我們,告訴我們這一切,剛才我所羅列的一切,都並非我們所認同的價值,我們從來沒有想過要賠上我們的生命去贏全世界,但說到這裡,談得如此興奮,如此振振有詞,我們每個人都不得不承認,不得不非常無奈的承認,這一切正正是我們生活的本質,生活的現實.有時候我們作為父母的看看下一代,作為老師的看看年青人,那種煩悶沮喪,怎至我形容是一種莫名的納悶,非常之悶的感覺,隨時會湧上心頭.

我還記得,小朋友無論是自己或是我自己的孩子,小時候天地曾經是非常廣闊,你可以隨意往來,不需既定的軌道,那時候作為小孩子,奔跑是我們的特權,自由是我們的專利,不過用不著等很久,這條道路不再是那麼寬敞,我們還沒有來得及看清楚自己那醜陋的模樣,細聽自己心中的想望,還沒有機會跨出我們生命歷奇的第一步,就已經...... 其實是出於好意,well intentioned 在熱情的守護之中,被四方的好人挾持我們,高聲喝停我們,罵我們這樣不可以做,那也不成,忙著要為我們引路,將我們的腳步擺在他們認為通達的軌道上,在這時代種種標榜下,恣意的跑出來攔路「小伙子,這不可作,那也不成」他們是這條路小朋友這樣不成,他們是well intentioned,是好心的,不過日子對小朋友來說,一點也不過癮,一點也不美麗,一點也不吸引.小朋友很快長大,到了少年時代,青年時代,幾經掙扎,大家總算找到自己志趣,好不容易一站又一站,那專業的裝備,即便說唸大學,再進修那些裝備,終於告一段落,好像覺得可以卸下小伙子的時候的束縛,迎向美好人生的旅途,憧憬自主的人生,心中確曾有過一份雀躍,你有沒有?我說的時候我都有丁點兒.我年青時候也想過,唸完書了,開始可以稍微自由自在一些,有一種憧憬,有一種雀躍,年輕的日子實在美好,我看著我的學生,男的有為上進,女的聰慧可人,感覺真是無敵,無可比擬,每一個都是俊男美女,看見他們拍拖,心裡多開心,感情路上大家開懷的盡情擁抱,盡情去愛,事業路途上,各人都全程投入,傾力以赴,大夥兒都帶著熱誠來出發,我不知道是否覺得,這多多少少是人生的寫照,我當時想起,這日子實在美好,能夠創出天地,走得痛快,活得精彩,然後有一天很有緣份找到了一個愛侶,找到自己愛的人,我們開始懂得自愛自重,因為為了愛人,你會更加見到自愛自重,更加勤奮上進.然後一天為人父母,更加顯得自己有種堅毅,有種承擔,好去做好兒女的榜樣.想一想,這些日子我們也曾經歷過,那時候你跟我也好像彷彿有用不完的精力,有花不盡的青春,一切我們好像都毫無保留,全程奉上,我們當下有如一個有為上進的青年人.想一想,如果你是這麼認真的話,你想到,我們真是不敢一刻忘記師長父輩的苦心,不想辜負眾人的期望,在家庭裡,在工作上,我們自覺盡心竭力,刻盡己任；在人生旅途上,我們已窮己之極,誠懇的做,認真的生活,我們傾情盡意,天地可鑑.所以便有個傢伙能十年來每星期都只睡三晚,在這種日子裡,大家是否覺得,不勝疲累?是否令人感到消磨意志呢?眼前的重擔,當下的難關,一項又一項的任務,一浪接一浪的挑戰,全足以叫人心力交瘁,我們那來的所謂敏銳的觸覺,清明的視野來繼續持守人生的方向?生命中精力最旺盛,體魄最強健的日子,我們已經不惜代價,全數奉上,活在這貪婪的世代,迎著這片土地上種種漂亮堂皇的目標,每天我們要與苛索霸道的時代理想正面交鋒,我那怕你有三頭六臂,任你有超人本領,最終都是一步一步走進深淵,落入迷惘,然後一天驚醒,不願意再投入的時候,你已經像我一般年紀了,人生已過了一大半.

我父親離世的日子,我量一量,已經過了四份三,所以我天天在倒數中,想一想那些日子感覺非常強烈,心裡湧起說不盡的漣漪與波濤,老實說,面對這種煩悶的,荒謬的人生,如果有些人在意的人在你身旁,如果有人在意的話,我自己覺得,那感覺是好丁點兒,整件事情來得比較不一樣.我想問問你,有否人 ``care"(關心)你心中的煩悶或無奈?有沒有人care?有沒有人很在意?我想大概都有一些,無論是深度的,又或是某程度的關懷在意,多多少少也會有的,不過我再問問你,那關懷的,在意的人,他是否明白,是否真正明白你的困難,你的處境?你覺得他明白嗎?可能你要想一想,有人會很關心你.我也很關心我太太,很關心我的女兒,老實說,若你問我全然明白與否?很明顯地你自己又明白自己幾多?你每天有這麼多的事情要去處理,到睡覺的那一刻也未必可以處理那些 unfinished business.一大堆事情,今天仍未搞妥,便趕著睡覺,閉上眼睛.我太太常說,我閉上眼睛不到半秒便已經睡著了,我真的是長期渴睡,在這狀態下我自己都被搞得一塌糊塗.我都沒想清楚那種煩悶是出於那裡?其實他也有問題,他也有困難,有一大堆事情情未搞定,他照顧自己都一塌糊塗,他又如何顧及你呢? 好!就算讓他有心,很想知道,他又有沒有能力知道?對,你自己也搞不定,你自己也搞不清雜亂的思緒.我要用上很長的時間整理自己的日記,寫 reflective journal,希望能透過書寫,透過文字整理自己複雜的,凌亂的思緒,更何況在你身邊,無論是親近的,無論是多關懷你的人,他所知悉的,的而且確是有限,的而且確是不容易讓人明白,但苦處在於在家庭裡,或是跟你最親的人,是你所在意的人,你對他的期望很高. 如果你最著緊的人是你太太的話,或是你的丈夫,你著緊他,你便緊張他緊不緊張你,所以你如果你很緊張他緊不緊張你,你便覺得我不用說,他也明白,但是你不說,他又如何明白呢?你對他期望很高,十分嚴格,於是問題便來了,家庭裡,眼見父母本來很好,為何有一天性情大變?到你開始有家庭有太太的時候,為什麼忽然間父母性情不同了,何解太太在拍拖時就很「棒」,現在總覺得不一樣了?真不明白!子女全部都是這種事情,是很難用三言兩語講個明白,你對他期望很高,你很想不用說他也明白,他也想明白,不過他明白不了全部,亦是你對他比較嚴苛了!你沒有好好的等待,給他空間,讓他明白.

我所說的這一切,其實想說出一件事情,這就是人生的現實,亦是人生的限制,世上或許真有非常美麗的性情,能甘願放下自己,完全從別人的角度去認識世界,去揣摩人生,去認識摰愛,但我們能夠相信世界上果然真的有些人可以不受自己的背景與經歷的限制,以致完完全全聽得懂我們內心的聲音嗎?縱是有願意關愛的人.人的關愛畢竟有限,你與我縱然遇上知音,但能夠彼此知道,亦不過是皮毛,古往今來最高尚的情操亦僅能及於傾力以赴,在崇基這地方常說,止於至善,崇基的校訓,我對止於至善有另一個角度加上的看法,我覺得止於至善這句話,揭視了聖哲深深明白到人的限制,縱願倒空自己,服役於人,亦無奈地得止步於我等這些蒼生微塵所歷知,所能及的那限度,便止在那兒,停在那裡,最多最盡倒是如此了,而且你知道永遠也去不到那位置,我們的而且確有這樣的限制.但很多人等著這些明白的聲音,很想知道,很想被知道,你別怪早幾年前那些人在旺角丟「通渠水」,我覺得我們的當時候的特首水平很低,他站出來非常嚴呵地斥責,不容許,絕對不能 tolerate,我讀這件事情的感覺就是有一個未被聽到的聲音,難道愛玩嗎?無聊乎?胡亂丟「通渠水」? 還有那些自焚的,還有要賠上自己,做人肉炸彈的呢?你以為在開玩笑嗎?傷害別人,亦傷害自己,就是因為聲音沒有被聽見! 如果有一個領袖他可以說:「老兄,我知道你有事情想說,可不可以讓我們聽清楚你想說的? 」這是一種不同領袖的水平,所以這時代所等待的,就是被聆聽的內心世界.

各位朋友,真是有這樣的人嗎? 能不受制於地域疆界,年齡,種族,能跨越星際時空,無遠弗屆能夠擁抱天涯地極,古往今來,能從萬物之始到宇宙之中,由你在母親的胎中未成形,直至年老衰敗那一天,走進人生盡頭的那一天,每時每刻,每分每秒,不分晝夜都能懷著無比的關切,走進人心靈極深的深處,完完全全,深深滲透在各人的心腸肺腑,認識你的思想意志,你覺得天地間,如果真有如此胸懷,是否超讚,很捧,很震撼?是驚天動地的信息?今天我和大家分享的就是這個訊息.

其實我差不多說完了,我用了這麼長時間來 cultivate,來形容,但我完全未曾好好地說一句上帝話語的分享,今天時間有限,但我覺得我需要用很長的時間去形容人生的煩悶,無奈,感受,好等大家把這些掘出來,別把它藏在枕頭之下,每天都不完成它便藏起它.我常跟我的學生說,如果有任何事情令你感到不安,感到unsettled 的話,解決不了的事情,其實大部份的人生的事情都是解決不了.請你牢牢的記住,好好地繼續想下去,記住它,不要讓它太容易的消滅,不要把它撲滅,不要收藏它,把它拿出來正面對著它,那怎麼辦?也沒得解決,我怎麼打算?把那些全都挖掘出來,好好的,正面的面對它,Squarely 的對著它, 然後你想想,跟著你對上帝說: 「你很厲害! 你是信靠得住的話,是否這樣一塌糊塗地走過人生,走過歲月?是否就是這樣? 如果不是的話,我們的信仰怎樣有說服力?你連自己也說服不了?你又如何說服別人,說明我們的信仰是如此驚天動地?」所以方才弟兄好好地,非常優雅的用普通話朗誦,將整篇詩篇一百卅九篇,他好好的朗誦出來,其實就是這信息,都不需要我說,上帝的說話已經說完了.你想想,有人關切你已是可遇不可求,人生極大的幸運,就是有人關心你,他所知的有限,他亦受制於自己的經歷背景,一切的限制,使他不能完全知道你,你仍然有種很落空,非常之失落,非常之虛空,誰明我心?誰能明白?

只要大家揭開詩篇一三九篇,稍稍的,輕輕的看看,你已經覺得無比的震撼,何等的驚天動地,詩篇一百三十九篇第一節說:「耶和華啊 ,你已經鑒察我, 認識我 .」祂看通看透,一切也逃不出祂的眼睛.第二節「我坐下 , 我起來, 你都曉得 ； 你從遠處知道我的意念 .」第三節「我行路 ,我躺臥 , 你都細察 ； 你也深知我一切所行的 .」意思即便說,無論你行路,你躺下,像看科幻片,在外太空一放大,超級望遠鏡一放大過來,全部看清,不單放大你的樣子,也放大你的心腸肺腑,你的一切思想,你想的每一個丁點的細節,你組織不來的,祂全都看透,真厲害啊!第四節「耶和華啊,我舌頭上的話,你沒有一句不知 道的 .」即便說你未開口,祂已知曉你想說些什麼,你什麼也還沒做,祂早已知道! 第五節「你在我前後環繞我, 按手在我身上 .」相當恐怖,即便說你探頭一看,祂在那兒,你回頭看,祂在,你向左看,祂在,你看武俠片才有,一轉身,全都在,多恐怖?但祂確實如此,祂到處存在.第六節「這樣的知識奇妙 ,是我不能測的 ,至高 ,是我不能及的 .」這就如我剛才說的,很震撼!對嗎? 很震撼,很厲害.讀過第一至五節,你覺得真的很震撼! 如果真的有這樣一個人有這般能力,你想想,非同小可!不知道你,或你身邊有多少人 and \#8230; 全香港,還有整個華南地區,全世界,祂不是今天,古往今來所有人,有沒有人有這般大的胸懷,全部包容其中,而且每一個細節,祂他都知道最軸心,你背後所有隱藏的動機,祂統統都知道,這是多震撼,多厲害?所以我將詩歌一百卅十九篇第一至六節我用一句說話歸納:「赤露敞開,無處掩藏」,對著全能的上主,我們在祂跟前是赤露敞開,無處掩藏.

這篇詩篇很容易讀,廿四節經文,你分開四段便成.六四廿四,六節一段,剛才我唸了六節,接著的六節是由第七至第十節,第二段我給它另外一個標題,我叫它「黑暗何懼,長夜不再」

139:7「我往哪里去躲避你的靈? 我往哪里逃,躲避你的面?」

139:8「我若升到天上,你在那裡；我若在陰間下榻,你也在那裡.」

139:9「我若展開清晨的翅膀,飛到海極居住」

139:10「就是在那裡,你的手必引導我；你的右手也必扶持我.」

第十節就是去到哪裡你的膀臂仍然扶持我. 這一句並非祂要尋找你,要毆打你,祂定必可把你找出來,但祂並非要打你,祂往那兒走你也不能逃避祂,祂往哪兒也能找著你.不過這是對那些自覺在漆黑死寂中,已經被忽略過的,已被遺忘的人－－便是我,是對這些人說的.主說:「我都看見,我都在那兒.」這句說話對於我來說,是無比的安慰.尤其是走過幾十年的人生,在這煩悶,沒有出路,沒辨法解決 and \#8230; 人生的事情那能解決?如果有個領袖跟你說:「我這?有一個解決方案,這一定是說謊者,沒有方案,那可解決?因為古往今來,所有問題都是大家一起自私,貪婪,彼此掠奪,集體的後果,集體承擔的後果,不要罵人罵得那麼絕對,你又有份兒,我們也是臭小子,有時候我們也佔人家便宜.早上駕車送我的女兒上學去,送了十六年由幼兒班開始,我看著一些家長怎樣插隊,我碰過大學的副校長,教會傳道人 and \#8230;　我也受不了,有沒有搞錯?你的女兒不遲到,難道我的女兒可以遲到麼?我只是隨便舉一個例子. 這些其實都是彼此掠奪,貪婪的結果,我們就想霸佔別人的領土,彼此互相傷害,全人類集體承擔彼此不負責的行為的一種後果.有否解決方案,如果有人跟我說能夠有解決方法,那人一定是瘋了,沒辦法解決!上帝說不打緊,就算不能解決,但最要緊的是祂知道,而且每一丁點都知道.你看傳道書,到那日主必審問,每一丁點你最隱藏的事,祂全都審問.所以我覺得這信息無比安慰,His laughing eyes are upon you.  祂的眼睛是看著你.這段經文139:11「我若說:黑暗必定遮蔽我,我周圍的亮光必成為黑夜；」139:12「黑暗也不能遮蔽我,使你不見,黑夜卻如白晝發亮.黑暗和光明,在你看都是一樣.」

我直至昨日我也沒空好好的想想寫講章,其實因我心裡有另外一個我感動,有一首歌我很喜愛,叫 ``No night there",那首歌我有很強烈的感覺,感覺強烈至我一定要在早上好好地將那首歌寫下來,變成一個好美的讚美詩,我用四個聲部,無伴奏要像黑人靈歌聲嘶力竭的喊,``No night there", 即那兒沒有夜晚,這首歌很勁就是說這句說話,黑暗和白晝在你看來都是一樣.那首歌形容著啟示錄廿一章所形容那新天新地的四方城,長闊高都是一樣三十公哩,那城很奇怪,形容那兒是沒有黑夜的,因為黑夜和白晝都是一樣的.為何黑夜白晝都是一樣?因為有基督的義袍的榮光一照,沒黑夜,那首歌詞實在非常震憾,and they count not time by years.原來哪兒的日子不是用年月日來計算的,不用計算的,完全超越了我們對時空的界限,超越常人被制約宇宙的秩序之中,是一個很震撼的事情,那城門從不關,因沒有敵人,沒人能傷害你,那?有基督在,這就是說我們要懂得擁抱那未來的家鄉,this world is not my home 不是一件消極的事情,而是我們知道我們等待的是一個遙遠的家鄉,我們看見那一刻,看見那種榮美,看見那種光明,看見上帝的大能的時候,我們就知道,今生的事情,那些你苦苦介懷的差異都是無關緊要的,是可以忽略的差異,對嗎?

我自己以前唸音樂的時候是很晚才開始的,剛開始時極度缺乏自信,我想我窮一生也趕不上,我在一,人家在十,跑盡了,我也只是從一跑到二而已,我已經doubled了自己,但人家一開始已經是九,那如何是好?覺得分別很大,但如果你一之後有十個零,跟十比,negligible, 你們明白我的意思嘛?所以我們的眼目實在太小了,看不見在上帝大能下這一切都是微不足道,但神所看的,耶和華不像人看人,人是看外貌,而耶和華是看內心,祂看見你內心的動機是多良善.像盛哥盛嫂,人家誤會他們,以為他們很可憐,要安慰他們,帶著安慰的訊息,其實他心裡知道.

我過了無比亮麗,很動人的一生,心裡很感謝上帝,用不著多說,因為主都知道,祂很明白,我沒有時間講完我的講章,牧師我們過了鐘,應該是,其實我這般的分享隨時可以停,我只是將一番情懷帶給大家,後面最精彩的經文還未講,第三段的詩篇一三九篇我最喜歡的一節是說,139:16「我未成形的體質,你的眼早已看見了；你所定的日子,我尚未度一日,你都寫在你的冊上了.」搞定了,對吧!其實不用想太多,其實大家定計劃,如果有一天真是有機會讓我做領袖,其實我又不是很討厭做,但因為我是一個完全沒有計劃的人,不需要的.你所定日子,老早就定下來,我還未出生,尚未到過一日,你全部寫在冊上.上帝有美好的計劃,有很好的心意,還差什麼?你不聽那你好等完蛋吧!因為你好好的聽信的說話,有那些人聽見,有那些人聽不到?是有分別的,說不定我就在這位置來結束,大家覺得有那些人可以聽到上帝的聲音,以致他懂得很努力,知道心裡仍然心裡煩悶和無奈的人生中,仍然心裡踴躍,靈?快樂. 有那些人是這樣的?他聽到上帝的聲音,每次他感到沮喪,感到無奈的時候,上帝說不用,OK的,我在,我全都明白,十分明白,那些丟「通渠水」的,自焚的,你不用做人肉炸彈,做黑寡婦,我知,我全知道!當你聽到這句的時候,你順服了,啊! 原來你明白,你是最關心,最仁慈.

大家許多時候好像不明白,為什麼我們要認識上帝的屬性,因為如果上帝是一個很淫猥的,很污穢的人,那便不得了!因為祂知但不知會祂會做甚麼,但是祂是很善良的,很愛你,祂會體察那些善良的,那些美好的,祂知道,那些全出於好的用意,即使沒有好的任何成就,祂知道,一個這樣善良的知道,不就成了?祂會告訴你怎樣才聽到,要一個好好的個人的選擇選擇,選擇在你,你要能聽到上帝的聲音,你要給上帝空間,你要提供空間給上帝,所以我說我今天到來和大家分享的時候,我一直都未有時間好好預備,但我不想找一些事來說,我直至今早祈禱的時候,看我具體要說什麼的時候,我心裡只有一個信息,就是你為今天的弟兄姊妹祈禱,所以今早我起來後,我沒機會再想想我要說什麼,不過我心裡就一直有個聲音,好好的為大家祈禱,因為這個群體都是很獨特,每一個都很不同,每一個都有他的需要,說不定我心中有過的掙扎,那些難處,一個如此軟弱不配,到如今都攪得一榻糊塗的人,那怕也可以當一個見證人.告訴大家,一個好比一榻糊塗的人,其實上帝都有夠用的恩典,是剛夠用,不太多,我仍然很淒涼,有很多不濟的地方,很多事都做得不好,好多人說:「Victor, 有沒搞錯,你這個不成,那又不是.」他都不喜歡,覺得我做得不好,或者覺得我不能滿足他的期望,但我仍然感到上帝有夠用的恩典,所以我覺得我的信息就到此為止,我們想想有那些人能夠在上帝的大能內,繼續很堅實的,心裡踴躍,靈?快樂的跟隨上帝,一定是那個跟上帝親近的人!

那你可要重組自己的優先次序,如果你又不看聖經,又沒默想,獨自安靜下來的時間,那會聽見?我們通常在辦公室只能聽到是冷氣的聲音,你知否冷氣是很嘈吵的? 關掉冷氣你便會開窗,開窗你便會聽到外面鳥兒的聲音,還有風聲,鳥聲風聲之外,還有草和樹葉衝擦的聲音,無語無語也無聲音可聽,它的量帶通遍天下,它的言語傳到地極.(詩篇十九篇 4 節)

不用多說,已經講完很多說話,只是全都是傻子,聽不到,這個宇宙本身就是很豐富,很全面的啟示,不過我們聽不到,give Him some space 好讓你可聽見祂的聲音. 如果你和上帝有一個好的關係,你發覺所有要解決的問題其實都有答案,而且祂還答應你,如果和上帝親近的人,他要像一棵樹裁在溪水旁,按時候結果子,葉子也不枯乾,凡他所作的盡都順利. 真的,真的要 take it seriously.  在如此煩悶,如此沒出路的人生底下,如果你看見有人仍然像一顆樹裁在溪水旁,他還可按時候結果子,他的葉子並未枯乾,他走了幾十年十分坎坷,十分難捱的日子,人將重擔加在他身上,世界叫他活得艱難,但他竟然矢志要別人活得輕省,這是否奇蹟?

這奇蹟我相信會在這群體中發生,因為我們是屬上帝的,我們聽到上帝的聲音,和祂 reassure 你,你是成的,這路是通的,是可信靠的,然後我們回去好好的認識上帝,事情開始有轉機,人生會不一樣.

我們一起祈禱:創造天地的主,賜寄我們生命的上帝,在我們心裡不斷運行的聖靈,救我們叫我們不用進入永死的耶穌,我們在這地方敬拜你,多謝你,我們眾人都是軟弱,在一個煩悶,沒出路的人生打滾,主,我們不感到意外,因為這本來就是人生的事實,當人的始祖犯了罪就最進入這狀態,我們沒法回轉,唯有等待救贖,一群蒙你恩,今早在這裡來全心敬拜,一心來擁抱你,主啊! 我們在這一刻求你三件事,我們求你幫助我們每一天更真切地來認識你,更加深深的來愛你,和更加緊緊地跟隨你,因為當我心深切地來認識你, 我們便會懂得傾情盡意來愛你；當我們傾情盡意來愛你的時候,我們懂得怎樣做基督的門徒,追隨你的腳步行每天的道路；當我刻懂得做一個基督的追隨者的時候,我們就能服侍我們的鄰舍,成為美好的見證他們,阿們.

我今天帶給大家這首歌,你們回去可以多唱唱,不錯的,尤其唱到中文版,這些如斯淒涼的人生,他竟然可以說,一生多稱頌,屈膝俯拜,所以我教會那副指揮,他譯了一個中文版本來唱,拿了中文來,我唱到這一句,眼淚掉下來,即便說,人生的而且確是這樣,但這一群人竟然sor sor 的,對,很淒涼,不過,我一生稱頌主,我還屈膝俯拜,我覺得這是難以想像的,最亮麗的見證,感謝上帝! 謝謝大家!



\newpage

\section{愛的洗禮}
\label{sec:182}
\textbf{陳榮超牧師 - 愛的洗禮}
\newline
\newline
\hyperref[sec:181]{< < < PREV SERMON < < <}
~
\hyperref[sec:toc]{[返目錄]}
~
\hyperref[sec:183]{> > > NEXT SERMON > > >}


講題:愛的洗禮　Baptism Of Love

講員:陳榮超牧師

場合:香港中文大學崇基學院禮拜堂

日期:2013年1月13日

聖誕節記念耶穌基督降生,福音書中從耶穌出生後很快便跳到他的福音事工.我們會發現在四福音中提及耶穌生平是很稀薄的.首先是在伯利恆出生,然後在兩歲的時候,約瑟及馬利亞帶著耶穌逃命,走到埃及,所以由兩歲至三十歲之間的耶穌生平是極之少.四福音的主要目的不是偉人生平傳記,並不是描述一個人的出生,青少年,怎樣學習讀書等等,這些資料很少,少的原故更見四福音的重點,耶穌基督工作的開端由洗禮開始,所以在紀念耶穌基督顯現的時候,福音的重點從洗禮開始,今天早上一開始的經課就是耶穌基督在約旦河河邊出現便是福音工作的開端,為什麼在談耶穌基督顯現會以這為經課?

因為耶穌基督的洗禮正是福音工作的開端,我們要思想福音開端是如何開展 - 耶穌基督的洗禮.我估計大部份坐在這裡的都是已受洗的信徒,亦有一組人已預備了七七八八的心情準備洗禮儀式,其中亦有少部份為慕道的,在心裡感到基督教信仰可以幫助他,協助他在生命上有所進一步.三類朋友對今天的經文可能有不同的感覺.當提到經歷洗禮的時候,我們當中有些已洗禮了幾十年,老信徒,盼望今天能去想想自己的洗禮和耶穌的洗禮的共同點,或認同點是如何?而其中一些而準備好,接受人生中非常隆重而嚴肅的儀式,在耶穌的洗禮上應有什麼回應.我們崇基禮拜堂已接受報名,大家也要趕快一點,洗禮快到了.另一班朋友對基督教信仰很感興趣,認為基督教信仰可以幫忙解決人生問題,或者家庭上的需要,究竟耶穌的洗禮對你們又有什麼共鳴?盼望今天的信息能讓信徒思想耶穌的洗禮與自己的洗禮的共同點和共鳴處是什麼?以致我們向前走時有一個開端的更新,在人生的起點中,讓耶穌的洗禮幫助我們.這是希望今天的信息中的目的.

路加福音第三章提到有一班猶太人(百姓)心裡盼望彌賽亞,救世主的降臨,當時有施洗約翰的出現,他們都想知道施洗約翰是否那位彌賽亞.但施洗約翰否認,並表示他雖然是以水替他們施洗,但有一位比他能力更強更大的,連提鞋也不配的,因為他是以聖靈和火替他們施洗,然後耶穌基督出現,耶穌也受了洗.正禱告的時候,天就開了,聖靈降臨在他身上,形狀彷彿鴿子；又有聲音從天上來,說:「你是我的愛子,我喜悅你.」耶穌基督受洗後有很特別的景象出現,有聲音說:「你是我的愛子,我喜悅你.」在信息中會作詳細解釋.首先回想耶穌基督來到約旦河受洗的情景.耶穌大約三十歲的一天,沒有指定那一天,他來到約旦河.那天,他決定要進行洗禮的儀式,那時已有很多人進行洗禮儀式,在那段日子中,常常有洗禮儀式,有些人因著施洗約翰的信息而來洗禮.施洗約翰比耶穌大六個月,並且有親戚關係,經文沒有提及時間,但那時正是施洗約翰傳講悔改信息,告訴百姓需要有一個悔改的洗禮,所以每天也有一大群人在約旦河邊接受洗禮儀式.關於施洗約翰的水禮儀式,猶太人自古已來已經有,他們到約旦河洗禮並不出奇.猶太人稱之為潔淨禮.在耶穌基督那個時代,有錢人家會在家中建一個水池,大約有半身的深度,而且盛滿水,若家中有人想行潔淨禮,便會走進其中,把人浸泡水中,然後起來抹淨,便是行潔淨禮.在舊約傳統,摩西律法書中,如果在街中踢到或接觸到屍體,便需要行潔淨禮.耶穌出世八天需要行割禮,但母親在過了四十天潔淨後,這是指猶太人的潔淨禮,就是在水中浸一下便能潔淨,有些人得罪了人,心裡很不舒服,便浸在水中行潔淨禮.所以這樣的水禮在當時是相當普遍的事情,有些人可能一天早,午,晚三次,因為經常犯錯,或者經常接觸到不潔淨的東西,有不潔淨的事在心裡,所以便行潔淨禮.在聖殿旁的牆邊也建了一些水池,幫助那些敬拜的人在敬拜之前先到水裡進行潔淨,覺得需要時便找聖殿的祭司,行潔淨禮後藉祭司有赦罪之權而相信罪得赦免.

施洗約翰到約旦河河邊,有些人覺得自己生命有不潔淨的地方,心裡覺得要到水裡得潔淨,能夠得到更新.在施洗約翰的時代,那些百姓到約旦河邊進行一般的潔淨禮,在施洗約翰的信息,他們感到有點不同,他們發現施洗約翰正是他們在心裡需求了很久的信息,他們心裡想施洗約翰是否他們所盼望的救世主和基督呢?經文說:「百姓指望基督來的時候,人都心裡猜疑,或者約翰是基督.」他們希望有新開始有新指望.為什麼令民眾帶來影響?因為信息中帶來兩個信息,在過去傳統聽不到.第一個是針對當時宗教人士在潔淨禮中佔了便宜,耶穌說:「安息日是為人設立的,人不是為安息日設立的.」的信息帶有相同意義,潔淨禮不是為人設立的,人不是為潔淨禮而設的.傳統讓我們為了儀式而儀式,忘記了儀式不是為人而做,人卻為儀式而做.所以當時的人為了這儀式,家家戶戶也建水池.祭司亦因為能夠赦罪而有貪污的情況,所以施洗約翰的信息含有針對當時的不公平的,把潔淨禮的意義淡薄了,只為了減輕心中的不安,因此施洗約翰把潔淨禮的意義從新更新,他並不是到人的家中,也不是到耶路撒冷聖殿前,而是在約旦河.約旦河代表活水泉,所以施洗約翰提到我們要以活水潔淨而不是一盆死水.第二要針對的是祭司赦罪的濫權.潔淨禮沒有赦罪的意義,卻要進行一個悔改的洗禮,不需要到祭司面前尋求赦免,而是到神面前悔改,這信息對耶路撒冷聖殿祭司有很大挑戰.所以我們能夠會意當時的百姓對施洗約翰的信息有這麼大的反應.因為對舊禮儀有新觀點,舊儀式有新看法,對一個習慣出現了新的動力.當百姓在約旦河進行洗禮,明白這跟往常的潔淨禮不一樣,因那是活水江河,可以洗淨我們的罪,因為大家存悔改的心,不為禮儀而進行禮儀,而是生命的更新,因此百姓聽了信息後,每天都有不少的人到約旦河進行洗禮,在那段日子有一天,耶穌基督就在約旦河接受洗禮.耶穌基督的洗禮帶給我們更深的層次.

在這裡先跟大家談談對洗禮的更深的看法,洗禮在希伯來文翻譯為潔淨禮或水禮,這個字本身是說把水聚集在一起的意思,指在水聚在一起時便可進行水禮.這字第一次出現在創世記,神創造天地時,天地混沌黑暗,神更在其中,如果我們與神一起時,便會知道那時周圍都是水,在第九節中,神把水聚在一起成為海,沒有水的地方是旱地,神完成這創造後,其他植物便開花結果,所以猶太人認為神把水聚在一起是一個生命的創造,因此每個進行水禮的人便會想起神創造生命的開始,把水聚集,神就讓生命開始.

在大學團契時,久不久便會到湖邊進行洗禮.記得曾經協助弟兄姊妹行進湖中較深水的地方,另外有兩位弟兄在其中協助信徒洗禮,那時在美國較北的地方,天氣很冷,有一位姊姊受洗後,從水中行上來,她讚嘆世界的美麗,鳥語花香,我四周觀看,還未到春天,天色灰灰暗暗,怎會有如此景象,所以到今天還很深印象,她看的跟我所看的完全不同.對我來說很深印象,她看到的我看不到,她聽到的我聽不到,神很恩待這位姊妹,她來自渺無人煙的地方,是美國土生華裔,中文並不通,來到大城市裡讀大學,我們有機會邀請她查經.有一天,她早上六點便拍我太太的門,太太在睡眼惺忪的時候開門,這位姊妹說她很想信耶穌,想受洗,神感動她要服侍人,很想先學中文,當時學中文並不容易,她的心意很大,她去當兵,因為那裡有言語學堂,她是一個很大心志的女孩子,信耶穌後,生命有新的開始,從窮窮鄉僻壤去讀大學,希望得到成就,神改變她的生命,讓她有新的開始,心裡有感動而願意服侍中國人,連當兵也願意,為了達成她的心志,所以每次去洗禮或者替人洗禮,便會想起1973年所認識的這位姊妹.

耶穌三十歲的時候,不知哪一天,來到約旦河,想到神的創造.根據考古學發現耶穌在約旦河受洗的地方,就是約書亞帶領以色列人渡過河到迦南地的地方,耶穌想到在這個地方,神的子民藉約書亞帶領到新的境界,就是以色列國在迦南地第一次得到的應許,是新的創造開始.施洗約翰把神新的權能國度權柄帶到以色列人當中,耶穌基督接受水禮是證明施洗約翰所說神的國度權能的開始,他實踐這開端,並把施洗約翰的信息在他的洗禮成為開始.水禮中耶穌基督提到人以悔改的心參與其中,生命不再像從前,人的生命已是神創造的,從前我在教會洗禮後會唱一首歌,音樂就是Happy Birthday,但歌詞不一樣,Happy birthday is two, only one will not be you, take Christ as your Savior, then you will have two. 所以接受了水禮的信徒,在那天神在你生命中是一個新開始,祂不是要替你算舊帳,生命是完全新的開始,在水中神創造我們的生命,祂把水聚集在一起時,當人在其中,一個新國度,一個新子民,一個新生命,一個新開端.在報名水禮時,是神在你生命中的一個新創造,如在徘徊在生命中,就不要錯過水禮帶來的應許.

耶穌的洗禮,他來到約旦河旁邊,慢慢走進水裡,沒有人帶他的,跟據過去的傳統,是施洗約翰陪著他,走進水裡替他施洗,但在猶太人的傳統,所有進行水禮的都不需要人相伴,各人憑自己心意走進水裡,在水中禱告,將心意交託,然後自己安靜地把心放在心上,慢慢蹲在水中,讓水浸過頭,然後回復自由.聖經記載耶穌從水中上來時,有聲音說:「你是我的愛子,我喜悅你.」聖靈降臨恍惚如鴿子,這是耶穌行水禮的大約情況.中文翻譯的受洗被認為是被動詞,表示需要有人幫忙的,這字翻譯出來有這意義,但實在不需要這被動的意義.耶穌的水禮是主動的,可能有人在旁邊自己進行水禮,就好像潔淨禮般,安靜在水裡,與神一起,向神禱告把心意交上,自己讓水蓋過頭髮,然後慢慢走上來,上來時帶來特別的能力,聖靈的降臨,有聲音從天而來,是外在的肯定,有聲音說:「你是我的愛子,我喜悅你.」這話讓我們明白水禮不但是潔淨過程,還是認罪得赦免的過程.潔淨禮本身不需要因有大罪而進行,耶穌基督的水禮帶來新信息,認同施洗約翰的信,這水禮帶來更深的意義,就是愛的洗禮,在其中耶穌基督展現了他來到世間的任務,建立愛神和互相相愛的群體,不單人的罪得赦免,這很重要,因為有罪便不能與神親近,我們需要潔淨,心裡很多害怕,人生帶著很多恐懼,需要神的憐憫,潔浸淨和幫助我們,解決我們的恐懼和重擔,耶穌完成這一切,但卻並不是他最重要的角色,今天罪得赦免,沒有罪在身又如何?各位信徒,神已赦免我們的罪,我們又作了些什麼?耶穌基督說不但我們罪得赦免,不但生命得潔淨,神的創造帶我們到愛的關係,耶穌基督有聖靈降臨在他身上,擁有內在能力,然後神說這是是我的愛子,我喜悅你!不是普通的人,是愛子,肯定你的地位,我是父親,是關係的建立,所以完成水禮後,便去召聚群體,呼喚十二門徒建立愛的群體,一個博愛的群體,神新的子民.新的子民即你和我,我們組合成愛的群體,把神愛的信息傳遞,認罪赦免是水禮的一面,成為愛的群體是另一面.沒有人在水禮後以為已經完成了便離開教會,以為有天堂的護照便不用作什麼.耶穌基督的洗禮是肯定從黑暗到光明,並把光明帶到群體中.初期信徒受洗後互相幫忙,甚至變賣田產,以致可幫助人,還在家中一同敬拜擘餅,讚美神,得到神的喜悅,神把人數日日加添.

今天結束時,重覆對三組人的鼓勵:已經接受水禮的弟兄姊妹,過去日子中回想耶穌基督的洗禮,要像耶穌基督的愛的洗禮,讓每個人連結一起去聚集,像神創造時把水聚集,使這創造的能力在我們身上發揮.第二組人,要快快接受水禮,你們有感動便要報名接受水禮,不要耽誤,這不是儀式,你有悔改的心,神把愛你的生命給你.最後一組在邊緣的人,一直在信仰邊緣徘徊,這信仰是很值得的,會為你帶來創新私改變,你會驚訝受洗後,無論任何年齡,你的一生將會有新的改變和體會,你會經歷聖靈降臨,有聲音說:「我所愛的,回來吧!」

耶穌的洗禮是一個愛的洗禮!

天父上帝,在你面前我們感謝你,你的創造藉耶穌基督的洗禮顯現在我們當中,讓我們看到的不是自以為是的儀式,我們在你面前帶領我們到約旦河中,開我們的眼睛,像在約旦河旁邊,我們有敬仰神的心,我們藉耶穌基督完成救恩工作,一個極大的應許,我們來到屬靈約旦河邊,我們清楚來到你面前,不用理會旁邊的人,不用理會其他聲音,我們好清楚明白在神面前接受水禮時就像耶穌基督受洗時有聖靈進入我們當中,有神的聲音對我說:「你是我所愛,我很愛你,我把你帶到愛中,藉著你把愛帶到群體當中,將福音傳開.」我們眾信徒這樣禱告,奉主名求.阿們!



\newpage

\section{公義如高山判斷如深淵}
\label{sec:183}
\textbf{伍渭文牧師 - 公義如高山,判斷如深淵}
\newline
\newline
\hyperref[sec:182]{< < < PREV SERMON < < <}
~
\hyperref[sec:toc]{[返目錄]}
~
\hyperref[sec:184]{> > > NEXT SERMON > > >}


講題:公義如高山,判斷如深海     Righteousness As Mighty Mountains Judgment As Great Deep

經文:詩篇36

講者:伍渭文牧師

場合:香港中文大學崇基學院禮拜堂

日期:2013年1月20日

最近因為報章深入偵查報導,揭發香港政府統計處前線調查員作假,他們到訪一層樓,發現內裡tong 開數間房就佯稱沒有人,因為上級視一個門牌一個個案,訪問數個戶主才當一個個案,不化算.統計數字不正確影響很大,因為很多政府的政策,規劃都基於統計數字；政策錯誤就浪費公帑,影響民生,延誤社會發展；若採用不真實的材料,偷工減料,豆腐渣工程,可引至塌樓斷橋,人命傷亡.

信實很重要,小朋友信任父母,當你答應過和他去街或賣東西給他們,他們就信以為真,期待你履行諾言,所以不要隨便答應他們.自然現象有可預計的規律,才有科學的研究和發展.從事商業活動,一定要有信譽,信守契約,不然經濟活動就癱瘓.我們需要有一個信實,可以信賴的世界.

今天的詩篇就強調上主的信實,但他的信實不是一個自然神論(Deism)的信實,像一個鐘表匠,上了發條,讓世界自然運作,遠離人世,對人的苦難無動於衷.詩篇三十六篇指出,上主的信實,就是他對人不變的慈愛.上主看到人在罪中的苦痛和苦難,與人立約,作他們的上主,施行拯救,被救贖的人要成為他的子民.上帝是信實的,會履行他的約,拯救保護他的子民,這是上主信實的慈愛.

詩篇三十六篇採用了常用的交叉配列(chiastic)方式:首尾呼應,中間對等,突出上主的慈愛和信實(ABB’A’).首段1-4節提到惡人「自誇自媚」(2節),「床上圖謀罪孽」(4)；尾段11,12節「不容驕傲人的腳踐踏我,不容惡人的手趕逐我.在那裡作孽的人,已經仆倒.」中間對上主的頌讚提到「慈愛」,祈求也提到上主的「慈愛」.很明顯,上主的慈愛對比惡人是詩篇三十六篇的主題.

聖經提到惡人時,一般和義人相比較,重點是他們的結局,而且用智慧傳統的議論文體來描述.智慧傅統(wisdom tradition)是用從創造論,不用救贖論的立場,以冷峻眼光來洞察世情,舊約箴言,傳道書,新約雅各書都是用智慧傳統的文體寫成.詩篇第一篇善惡的比較是典型的智慧傳統,惡人將像糠秕被風吹散；糠秕散發,看似金碧輝煌,但風吹過後無覓處,而義人像一棵樹,栽在溪水旁,按時候結果子,葉子也不枯乾.但詩篇三十六篇用他們與上主的關係作比較,而且運用不同的文學體裁:1-4節以第三身,用智慧體裁描述惡人；5-12節以第二身,用詩歌體裁頌讚上帝 (5-9) ；最後10-12節以第一身的祈禱結束.

1上主慈愛信實不變

信實的「切約」

慈愛這字貫穿整本聖經,指上主與人立約拯救世人的慈愛.這約不是平等協議,保護彼此利益的契約(contract), 是宗主國對其蕃屬宣布其主權和關係的約(covenant),主動來自宗主國,立約原文稱「切約」(cut the covenant), 主動立約的一方,把祭牲切開,從其中經過,寓意願意以生命作為保證,保護蕃屬；作蕃屬的也履行敬愛宗主的責任,立約強調彼此的親密的關係.上主與亞伯蘭立約,引領他離開吾耳(Ur),叫他成為大國,後代將如夜空星宿之多,亞伯蘭把牲口劈開,「日落天黑,不料有冒煙的爐和火把從那些肉塊中經過.」(創15:17)立約的一方以生命來保證,履行責任,如果爽約,就像祭牲被劈開,聖經中就常說耶和華以他的永生來起誓,保證他的慈愛信實不變.

救贖的慈愛

最早出現約的字是創6:18節上主與挪亞立方舟之約,以方舟拯救挪亞一家八口,洪水過後,並以彩虹為立約記號,永不用洪水審判世人.(創9:9-17) 從挪亞到亞伯拉罕,以撒,雅各,約瑟,摩西直到耶穌基督,上主都對他子民守約施慈愛.當以色列人在埃及法老治下為奴受苦,上主聽到他子民的哀聲,看到他們的困苦,興起摩西,帶領他們離開埃及,「讓我的百姓離開」(Let my people go).上主對摩西說將會「給他們(以色列人)興起一位先知像你.」(申18: 18)耶穌登山變像時,代表律法和先知的摩西,以利亞出現,談論耶穌去世的事(路9: 31) ,耶穌去世,原文就是「出埃及」,耶穌就是那位像摩西的先知.天使對約瑟說:「她(馬利亞) 要生一個兒子,你要給他起名叫耶穌,因他要將自己的百姓從罪惡裡教出來.」(太1: 21) 上主信實守約,拯救他的子民,這是他的慈愛,而惡人就是拒絕上主慈愛的人.

2. 惡人拒絕上帝的慈愛 vv. 1-4

「惡人的罪過在他心裡說:我眼中不怕上帝!他自誇自媚,以為他的罪孽終不顯露,不被恨惡.他口中的言語盡是罪孽詭詐；好他與智慧善行已經斷絕.他在床上圖謀不善道,不憎惡惡事.」

不同的懼怕

「我眼中不怕上帝.」(1節b) 不怕,英文可繙作不尊重,no regard.懼怕有多種:若你欠人錢拖延不還,見到債主你會懼怕,但這怕與你見到一獅子不同.又如你景仰德蘭修女,若她尚在人間,突然進入崇基禮拜堂,你有一種懼怕,這是一種敬畏,早期教會信徒彼此相愛,「恆心遵守使徒的教訓,彼此交接,擘餅,祈禱.」(使2:42) 眾人都懼怕.為何眾人都懼怕?因為他們看到人可以如斯共融相愛,他們看到一種神聖的景象,與平常割裂的景象,而懼怕.我們懼怕上帝,因為他神聖可畏,雖然可畏,但有一種令我們響往的吸引力,不同靈異的幽暗世界,令我們恐懼.

罪使人漸漸失去愛的能力

「惡人的罪過在他心裡說.」(1節a) 這裡說罪是一種屬靈的實體,像電影中的異形(alien) ,不屬於自己又屬於自己的實體.罪教唆惡人,不須敬畏上主,罪惡在惡人心中盤踞,毒害他的靈魂,使人心靈失去愛慕美善,聖潔的能力.與此相反,保羅提醒提摩太:「凡是真實的,公義的,清潔的,可愛的,有美名的,若有甚麼德行,若有甚麼稱讚,這些事你們都要思念.」(提腓4:8)

一個對上帝不虔不敬的人,對人也會不忠不義；失去對上主的愛,也會失去對人的愛.他白天說話盡是罪孽詭詐,睡覺時在床上圖謀罪孽.「他不行善事,不憎惡惡事.」(詩36: 4)為何他不憎惡惡事呢?為何他對不義的事不感到忿怒呢?因為他對罪惡的受害人的痛苦無動於衷.所有的罪惡,都有受害人(victim) .十誡中不可殺人,不可姦淫,不可偷盜,不可作假見證,不可貪戀人的財物,都是對人沒有愛心,所以耶穌說,誡命的總歸就是愛.

詩人描述惡人,以不懼怕上帝開始,以不憎惡惡事結束,說明惡人因為靈魂被罪侵蝕,失去愛的能力,失去愛慕真,善,美,聖潔的能力,也失去憎惡罪的敏銳,失去對被罪壓制的人的同情心.

3. 使我們飽足的慈愛vv. 5-9

從第三身冷峻的的智慧文體描述惡人,跟著筆鋒一變,作者用感情洋溢的詩歌體裁來頌讚上主和祈禱,對比惡人.

「你的慈愛上及諸天,你的信實達到穹蒼；你的公義好像高山,你的判斷如同深淵.上帝啊,你的慈愛何其寶貴!世人投靠在你翅膀下.他們必因你殿裡的肥甘得以飽足；你也必叫他們喝你樂河的水.因為,在你那裡有生命的源頭；在你光中,我們必得見光.」(vv. 5-9)

從慈愛到公義

「你的慈愛上及諸天,你的信實達到穹蒼；你的公義好像高山,你的判斷如同深淵.」(vv. 5, 6)

惡人自以為罪孽終不顯露,不被恨惡；詩人就頌讚上主公義像高山,上主的判斷如深淵；上帝的公義像高山立於前,毫不動搖；但上帝的作為像深淵,不是我們可以猜度,他何時審判,如何審判,不是我們可以和道的.慈愛和公義有何關係呢?公義是愛的均分,我們追求公義,因為愛每一個人,保護他們的權益,不要以為上主對罪孽袖手旁觀,對惡行好像默然無聲.

在你翅膀的蔭下-母鷹愛護雛鷹的關係

上帝啊,你的慈愛何其寶貴!世人投靠在你翅膀下.(v.7)

詩人跟著讚美上主信實的慈愛和保守.摩西在以色列人快要進入迦南美地時,在約但河東的摩押平原重申誡命(申命記Deuteronomy, second law),回顧歷史,說到上主在曠野引領保護以色列人:「耶和華遇見他在曠野荒涼野獸吼叫之地,就環繞他,看顧他,保護他,如同保護眼中的瞳人.又如鷹攪動巢窩,在雛鷹以上兩翅搧展,接取雛鷹,背在兩翼之上.」(申33: 11)

但有解經家說,這是聖殿敬拜的語境.「在你翅膀的蔭下」暗示在聖殿天使撒拉弗(seraphim)的個翅膀.以賽亞先知進入聖殿敬拜:「我見主坐在高高的寶座上．他的衣裳垂下,遮滿聖殿其上有撒拉弗侍立．各有六個翅膀．用兩個翅膀遮臉,兩個翅膀遮腳,兩個翅膀飛翔．彼此呼喊說,聖哉,聖哉,聖哉,萬軍之耶和華．他的榮光充滿全地.」(賽6: 1-3)

聖殿肥甘, 樂河之水-在敬拜中得著滿足

因為跟著第八節:「他們必因你殿裡的肥甘,得以飽足．你也必叫他們喝你樂河的水因為,在你那裡有生命的源頭.」.(8, 9a)

聖殿的裡的肥甘指獻祭祭牲的脂油,獻平安祭完畢,獻祭的人可分享祭品,形容與上主的親密關係,人在敬拜中得到最大的滿足.

今天是顯現節第二主日,詩篇經課提到喝你樂河的水,在經課的安排是有特別禮儀的意思.

樂河的水是生命的源頭,在創造的新世界,在伊甸園,有生命樹,有河從伊甸流出.(創2:9,10)這生命河在救贖的新天新地再度出現:街道當中一道生命水的河,明亮如水晶,從神和羔羊的寶座流出來.(啟22:1)

上主日講員提到耶穌受洗,受洗或受浸一字原來有把水聚在一起,創造天地時,地是空虛混沌,淵面黑暗,上帝把水分為上下,把水聚在一起,露出旱地,才生發生命,所以洗禮表示新生的開始.耶穌在約但河接受施洗約翰的洗禮,表示與人認同,站在人的地位,與人聯合,從水裡上來時,天開了,天父為耶穌作見證,有聲音從天上來,說:「你是我的愛子,我喜悅你.」(路3: 22) 「你是我的愛子,我喜悅你.」是彌賽亞的稱號,從此耶穌公開傳道,走上十字架.我們接受水禮的人,是與耶穌聯合,和他同死同埋葬同復活,和他聯合,得著新生,成為教會一份子,而教會是基督的新婦,所以今天福音書經課約翰福音21-11是一個婚宴,慶賀基督和教會的聯合.基督在婚宴中行了第一個神蹟,使水變酒,而且是最好的酒,表示基督是樂河的水,生命之源,祇有基督才能滿足我們.耶穌受洗後聖靈彷彿鴿子,降在耶穌身上,我們受洗的人,也領受聖靈的恩賜,互為肢體,彼此服事,這是今天書信經課林前12:1-11節的信息.

在你的光中,我們必得見光 v. 9b

在敬拜中,在聖言的亮光中,詩人有了真實的知識,不像惡人自誇自媚,自欺欺人.上主的慈愛,在基督裡拯救了我們,使我們經歷了真實的愛,也學曉愛人,因愛人追求公義.我們也有真的視野,認識自己的身份.是上主的子民,基督的身體,聖靈的團契,我們被召,是要建立基督的教會.

最後的祈禱vv.10-12

「願你常施慈愛給認識你的人．常以公義待心裡正直的人．不容驕傲人的腳踐踏我,不容兇惡人的手趕逐我.那裡作孽的人,已經仆倒．他們被推倒,不能再起來.」(vv.10-12)

一輪頌讚後,詩人作一簡短的祈禱,求上主敗壞惡人的計謀,保守自己脫離兇險.我們不用懼怕惡人,因為惡人已被徹底打敗,不能再起來.詩篇中的惡人,當然是實有其人,不懼怕上帝,也不憎惡惡事,不愛慕上主,也不憐愛人,無視罪惡帶給人的痛苦,聖經也是用惡人,表示罪惡的力量,「惡人的罪過在他心裡說.」(1節a) 罪就是拒絕愛,拒絕神聖和美善.但罪惡已經被上帝的愛所擊敗,基督為我們的罪,包括詩篇三十六篇所描述的惡人的罪,被釘十字架,上主像高山的公義滿足了,我們不明白為何上主竟然用這方法來拯救我們,他的判斷如深淵,超乎我們的想像.想到這裡,讓我們和保羅一起讚美(羅11: 33-36):

「深哉,上主的智慧深哉,　神豐富的智慧和知識.他的判斷,何其難測,他的蹤跡何其難尋,誰知道主的心,誰作過他的謀士呢,誰是先給了他,使他後來償還呢因為萬有都是本於他,倚靠他,歸於他．願榮耀歸給他,直到永遠.阿們.」

當我讀到「在你的光中,我們必得見光」時,我都想起早期北菲迦大基(Carthage,現今突尼斯)教父特土良(Tertullian,160-225 A.D.) .他是通曉希臘哲學和羅馬法律的著名學者,被聖經的陌生感覺所吸引,聖經像一度清泉,解除他長久以來希臘哲學不能給他的屬靈饑渴.他在光中看到光,信主以後,他把自己的學問獻上,成為知名的護教家和神學家,被譽被西方神學之父.

今天的詩篇告訴我們,世上最恆久不變的是上主信實的慈愛,這慈愛已由基督為我們釘死在十字架顯明了,我們也領受聖靈的恩賜,被差遣傳揚福音,建立上主的教會.阿們.



\newpage

\section{期待基督的喜悅}
\label{sec:184}
\textbf{陳曉東博士 - 期待基督的喜悅}
\newline
\newline
\hyperref[sec:183]{< < < PREV SERMON < < <}
~
\hyperref[sec:toc]{[返目錄]}
~
\hyperref[sec:185]{> > > NEXT SERMON > > >}


瑪拉基書 3:1-5
\newline
\begin{longtable}{cl}
\hline
\hline
章節 & 經文 (和合本修訂版)\\
\hline
3:1 & \begin{tabularx}{0.7\textwidth}{X} 萬軍之耶和華說：「看哪，我要差遣我的使者在我前面預備道路。你們所尋求的主必忽然來到他的殿；立約的使者，就是你們所仰慕的，看哪，快要來到。」 \end{tabularx} \\ \\ \relax
3:2 & \begin{tabularx}{0.7\textwidth}{X} 他來的日子，誰能當得起呢？他顯現的時候，誰能立得住呢？因為他如煉金匠的火，如漂洗者的鹼。 \end{tabularx} \\ \\ \relax
3:3 & \begin{tabularx}{0.7\textwidth}{X} 他必坐下如煉淨銀子的人，必潔淨利未人，熬煉他們像金銀一樣；他們就憑公義獻供物給耶和華。 \end{tabularx} \\ \\ \relax
3:4 & \begin{tabularx}{0.7\textwidth}{X} 那時，猶大和耶路撒冷所獻的供物必蒙耶和華悅納，彷彿古時之日、上古之年。 \end{tabularx} \\ \\ \relax
3:5 & \begin{tabularx}{0.7\textwidth}{X} 萬軍之耶和華說：「我必臨近你們，施行審判。我必速速作見證，警戒那些行邪術的、犯姦淫的、起假誓的、剝削雇工工錢的、欺壓孤兒寡婦的、屈枉寄居者的和不敬畏我的人。」 \end{tabularx} \\ \\
[1ex]
\hline
\hline
\end{longtable}
日期:2013年1月20日

經文:瑪拉基書3:1-5

證道:陳曉東博士

我們其實注意到在今天所讀的,舊約聖經瑪拉基書三章1-5節,有一個在人類歷史裡面,非常嚴峻的問題；也是非常讓人難以作答的問題.

以色列的百姓問先知「公義的　神在那裡?」我們知道,瑪拉基書接下來的回答,嘗試的是回答一個「神公義的問題」.也許,公義的　神,或者是在信仰裡面,　神義論的問題,是我們在整個的世代當中,常常面對,無從談起的一種的問題,曾經有人從邏輯上面來去挑戰教會的信仰.

如果,全知,全能,全善的真神,真是這位公義的神,為什麼世界有苦難?要嘛,祂就不是全知的神,所以,祂不知道這個世界有苦難；要嘛,祂就是全善的神,所以祂允許這個世界有苦難；我們知道這麼一種從人生所發出悲蒼的呼喊,實在是歷世代每一位追尋真理的人,心裡面一種非常痛苦的呼喊.

在我年輕的時候,我看過一部的戲劇,整個戲劇一開場的時候,有兩個流浪漢,一個流浪漢叫愛斯特拉岡(Estragon),另外一個流浪漢叫弗拉第米爾(Vladimir),他們站在一個空盪盪的舞台上面,整個的舞台只有一顆的枯樹,這2個流浪漢,在這顆枯樹下,等待著一位叫果陀(Godot and lrm;)的人出現,整個戲劇開場的第一句對白就是:.nothing to be done,沒啥好做的.接著演員就引用了箴言13章12節,一個前半段說:希望幻滅,心靈隨之破碎.

學習戲劇的人都知道,這是盛行在半個世紀前的戲劇《等待果陀》(法文:En attendant Godot and lrm;,英文:Waiting For Godot)的開場白.這一齣戲劇很正確的補捉了我們人類,從24世紀下半葉,到我們今天的經驗和感受,人生就好像一個沒有具體定義的地點,或是一條的鄉間小路,或者就是一個街上的拐角.

這齣戲用一條的鄉間小路,路旁的一顆枯樹,來暗示著我們所生活的世界,是那麼的荒蕪,整個的人類就好像是這兩個演員,無所事事的站在那邊打發時間,他們要等待,但是等待一個他們不認識的人,等待一個叫「果陀」的人.戲劇家很敏感的聽到這個名字背後,這「果陀Godot」這個字,和「God 神」這個發音有相似性.他們在等待一個叫「果陀Godot」的人,但是他們對果陀一無所知.雖然不知道是否真的有這麼一個人?而事實上,我們看過這部戲劇的人都知道,果陀在全個戲劇結束以後,也沒有出現過,這場等待,顯然是一個徒勞無功的等待.

常常有人心裡還沒有準備好,就去看這個戲劇的時候,都會氣得破口大罵,覺得自己所花得那百,十來塊錢,幾百塊錢,完全是浪費的一件事情.因為整個的戲劇,就是一個重覆兩幕的戲,每一次的戲劇結束以後,都會有一個據稱是果陀,所派遣的一個牧羊人,一個牧羊童出來通報說:果陀馬上就來了!大家忍耐一下.但是,從來沒有見過這樣的果陀!

至少這個時候,我們看到在戲劇上面,這兩位的演員,他們會懷疑,會討論,他們的處境是不是一種毫無義意的處境?他們的表演是不是一種毫無義意的表演?他們所站立的地方是不是一種毫無義意的所站立的地方?

我們在這個戲劇裡面,有好幾個精彩的片段.演員開始說:算了吧!咱們走.另外一個演員說:不行,我們不能走.愛斯特拉岡(Estragon)說:為什麼不能走!弗拉第米爾(Vladimir)說:因為我們在等待果陀啊!愛斯特拉岡(Estragon)說:啊!天哪!

整整的二幕長長的戲,劇情在原地踏步,毫無進展.演的人疲倦不堪,看得人也累的半死.但是這一切的折磨,就是為了和一種恐怖的寂靜,保持距離.

這個世界不能夠安靜下來,安靜下來實在太恐怖了,我們總是要做點什麼事情,讓我們在這樣一個無奈的世界裡面,來打發,打發時間.所以,兩個人言論滔滔不絕,透過一些很荒謬的話題,插葷打嗑不停的吵架又不停的合好；無奈的來去打遊戲,玩遊戲,去消費時間；沈悶到一個地步,甚至要去討論,怎麼樣的自殺方式才是最棒的自殺方式；到全戲結束的時候,都沒有搞清楚,誰是果陀?

為什麼要再這樣的一個詭異的地方,在上帝拋棄的地方,我們來去等祂?就算在自己的孤獨裡面,兩個人在那邊插葷打嗑,在這邊講笑話,在這邊玩遊戲,甚至在整個戲劇裡面,在加上兩個的角色出現,1個是地主波卓(Pozzo),還有他的奴隸幸運兒(Lucky)來進入到這樣一個的爭論裡面,他們的來到並沒有改變,這樣一個寂靜的可怕世界,他們的來到,只是讓我們這個世界更添煩,添亂.

當然,這樣一齣盛行於整個20世紀下半葉的戲劇,它的表演為什麼能夠長久不衰?因為整個的表演裡面,並沒有像我現在講的這樣的嚴肅,這麼的可怕,其實上整齣劇戲,在那種無奈裡面,是充滿了扭曲,充滿了誇張,充滿了好笑.

但是要回答一個人類不變的問題:人類的苦難,上帝的來臨,仁慈到底有什麼意義?人是不是生活在一種雜亂的,匆忙的,自言自語的世界當中?人類在這個世界上面,所用的一切高雅的語言,一切高級的語言,包括了神學的語言,講道的語言,藝術的語言,哲學的語言... 到底是不是一堆文化的垃圾?

事實上在這齣劇戲裡面,就是要告訴我們一個事實,或者是劇作家所理解的事實:這一切真的是垃圾!

所有的人都是那麼的無奈.我們都知道,人生存在這樣不定的把現實的幻像,但是我們不能拋棄這樣的幻像,我們只能夠在那無奈的等待裡面,等待著這位的果陀,等待著這位救世主.

薩繆爾-貝克特(Samuel Beckett)這位劇作家,讓劇中人變得荒誕可笑,讓劇中人變得那麼的悲傷.但是他把這種方式,整個的人類說:在當代所處的處境,原封不動的擺在人面前,讓人來狂笑,讓人來絕望,讓人來無奈.但是帶出一個非常難以表達的哀傷,人類還有希望嗎?人類還能夠為未來的世界展望些什麼嗎?

很多的信徒看這齣戲的時候,常常覺得很憤怒.因為覺得好像這齣荒謬戲劇,在消遣著基督教的信仰,特別是關於聖經裡有關盼望的信仰.如果我們帶著這樣的眼光,來去看這個世界的時候,也許我們今天坐在這邊,真是一場無意義的等待.也許我們點起這樣的第二根的蠟燭之後,這樣的蠟燭又是一對的,是劇作家所說的「文化的垃圾」.

但是我們從這樣一個沈澱澱的,這樣的一個理解當中,我們看到後面這個世道人心中,有一個不一樣的聲音,這個聲音在表達著一種深切的,人的關懷,人道的關懷.

我們知道在20世紀有一種的思想潮流叫「存在主義」,也許是一個非常紛雜的學派,但是在這樣的一個紛雜的學派後面,有一個共同的聲音,我們生活在一個難解的矛盾世界裡,在這樣一個矛盾世界裡面,我們怎樣活出人的尊嚴?我們怎樣活出人的自由?我們怎樣活出人的意義?

等待果陀這齣劇戲,寫作在1948年.也許大家知道,1948年正是世界人權宣言,所發表的時間,1948年人權宣言發表的時候,正是要指出,我們從一個戰亂的時間裡面站起來,在這樣一個戰亂的世界當中,我們站起來後,要怎麼樣活下去?整個的人類需不需要有尊嚴?整個的人類需不需要有自由?整個的人類需不需要有意義?

回望50多年前,這樣的一個聲音的吶喊,我們看我們今天21世紀,我們的世界跟等待果陀的世界,相差有多少?我們的世界是不是變得更幸福呢?更快樂呢?更美好呢?21世紀也許「商品拜物教」成了我們整個人類的主流信仰,我們還能夠在這樣一個的商品拜物世界裡面,在高喊著:「我購物,故我存在」的世界裡面,我們還能夠定義幸福嗎?我們還能夠活出人的尊嚴和自由嗎?

和薩繆爾-貝克特(Samuel Beckett)同樣的一位的文學大師阿爾貝-卡繆(Albert Camus),他曾經提出一個非常聰明的薛西弗斯(Sisyphis)理論,他藉用希臘一個古典的神話,提到這個薛西弗斯(Sisyphis),這個人類,因為他欺騙死神,所以被天神懲罰,做為一個永遠的懲罰,他的懲罰是非常荒謬的,為什麼這麼荒謬呢?這懲罰是:白天,他出來上班,把一塊的大石頭推到山的最高峰,這樣一個努力,勞力,是要用盡他一天所有的力氣,所有的能力,才能夠做到的事.但是到了黃昏的時候,這個大石頭就咕嚕咕嚕的滾到山的另外一邊.第二天,他的工作要重新的開始,在這樣一個重覆又重覆的工作裡,永遠的工作.但是,卡繆(Camus)只說一件事情.他說:正是這樣子,薛西弗斯(Sisyphis)是幸福的.為什麼他是幸福的?因為他在這樣一個無意義的工作裡面,他做出了一種的抵抗.他抵擋著,在他自己生存條件下,唯一可以反抗的事,就是活出他的尊嚴.

人生──卡繆是這樣說──是一片黑暗的明亮,這是一個不邏輯的講法,這是一個非邏輯的思學.但是他更能夠表達出我們人類的感受.我們今天的世界,包含了那麼多的對立,尖銳,讓人感受到存在的恐懼,甚至荒唐的感受.

有什麼感受,比我們看著電視,看著一位叫劉霞的女子,在電視機面前哭訴!八千萬優秀的黨員,代表著社會最先進的生產力,最先進的人類文化的精英政黨 and \#8230; 沒有任何的理由,把一個弱小的女子軟禁在家中二年.

有什麼圖畫比我們覺得人生更荒謬?那位為自己的民族,擺上自己所有生命的,一個弱女子──昂山素姬──要用那麼多的槍桿子,把她圍困在一個房間裡邊.這是我們人生中所看到的諸多荒謬事情中的,兩個小小的例子而已.

難怪卡繆說:我們的人生,每個人都是薛西弗斯(Sisyphis),差別只是我們有否認識到這點?他說:你起床,你坐車,你四個小時在辦公室,工廠工作,你吃飯,你坐車,你四個小時工作吃飯睡覺,星期一,星期二,星期三,星期四,星期五,星期六,大部分的時間,一天接一天,按著同樣的節奏,週而復始的如是,

直到某一天,你問:為什麼?當你問為什麼的時候,你開始有一份驚奇,你開始從一種的厭倦中,有了一個開始.但我們知道,人間的藝術家,文學家 and \#8230;沒有任何的答案,所以卡繆最後只能說:我們人只能是:有一種有靈性的警覺.

他的老搭檔尚．保羅．沙特(Jean-Paul Sartre)──後來跟他連友誼都破裂了──.完全不同意他的這樣子的想法,最後跟翻臉,絕交,為什麼呢?他說:卡繆,你錯了,我們要行動,我們要實踐.不幸的是,沙特也錯了.因為他的行動,是如風逝去的社會行動理論──馬克思主義──在他的法國社會裡面,終究沒有任何的影響力.

但是,我們今天所讀的福音經課,今天的福音書,不是荒誕的人生戲劇,而是一個救贖人類尊嚴,自由,以及幸福的偉大信息.我們從撒迦利亞所唱的詩歌裡面,我們看到了三個人物:一個叫:大衛；一個叫:亞伯拉罕；還有一個是他自己的親生的兒子:施洗約翰.

我們看到這樣一位的老先知,在他自己的生命讚歌裡面,在他生命的預言當中,他去稱頌上帝,但是這個稱頌,不是一種的無奈的等待,不是一個沒有根基的展望.他的稱頌把他在歷史的回望當中,看到了歷史當中,上帝在人間所施行的奇妙揀選.看見上帝在整個世界的起初,對亞伯拉罕的聖約!這樣一首的詩歌非常的有特色,因為他點出了基督信仰裡面,有二個非常重要的主題.

就像他說:以色列的　神,是應當稱頌的．因他眷顧他的百姓,為他們施行救贖． 在他僕人大衛家中,為我們興起了拯救的角, (正如主藉著從創世以來,聖先知的口所說的話.) 拯救我們脫離仇敵,和一切恨我們之人的手．(路加福音 1:68-71)

我們中文看不見這個時態,如果你看到,任何的有時態的語言,英文也好,德語也好,或者是它的原文,希臘文也好,你會看到當撒迦利亞或者是路加去記載著撒迦利亞這段詩歌的時候,他講到上帝的看顧,上帝的興起,上帝的演說,上帝的恩慈的顯示,都是用一個過去的時態.也就是說:我們的信心的基礎在那裡?我們盼望的基礎在那裡?不是在一個無根的世界,不是在一個被拋棄的世界,不是在一個不可理解的世界,而是在一個救恩的歷史裡面,神實實在在,在過去工作了.

所以祂在過去工作,我們也相信祂在今天工作,我們更相信祂在將來也繼續的工作!就是因為這樣,他告訴我們,把我們的心靈目光,去看整個救恩歷史裡面,上帝怎麼樣在歷史當中工作?怎麼樣在大衛家,這樣的一個興起,揀選事情!路加的眼光,不僅僅是讓我們看一個上帝揀選的一個國度的興起,路加透過撒迦利亞進一步的告訴我們,這樣的一位上帝,不僅僅是大衛的上帝,祂更是亞伯拉罕的上帝.

所以他在經文裡面,告訴我們,祂不僅僅選了大衛家,為我們興起了一位救贖主!不僅僅在大衛家為我們帶來一個揀選的革命,他更是在列祖的時代,向亞伯拉罕,向整個恩典的泉源,立下了祂的聖約!好叫整個世界,屬祂的百姓,能夠在這個聖約當中,被神來救贖!

正是這樣,我們看到的76節,為施洗約翰所唱的讚歌:孩子阿,你要稱為至高者的先知．因為你要行在主的前面,預備他的道路． 叫他的百姓因罪得赦,就知道救恩． 因我們　神憐憫的心腸,叫清晨的日光從高天臨到我們,要照亮坐在黑暗中死蔭裡的人．把我們的腳引到平安的路上.(路加福音 1:76-79)

我們從施洗約翰的使命當中,我們看到了一件非常重要的事情:人生也許像這些存在主義者那種敏感的心靈,所看到的是那麼的荒蕪,充滿了對立,無奈,衝突.但是在這樣的世界當中,我們不必插葷打嗑,我不必寄情於玩遊戲,也不必讓這樣的虛無的力量把生命完全的吞沒.

施洗約翰的蒙召告訴了我們,在黑暗的世界當中,生命是可以興起的,生命可以聆聽到上帝對我們的呼召,就是因為我們的生命能夠被呼召,所以施洗約翰去預備,準備,裝備,讓主的道在這樣一個黑暗的世界裡面,能夠展開.在這樣一個沙漠的世界裡面,能夠有綠洲.就是因為這樣的黑暗,所以施洗約翰被興起,他被興起去教導救恩的知識；就是因為黑暗,所以施洗約翰被呼召,被興起去照明,照明生命的死亡,並不是最終的恐懼.生命最大的恐懼,是我們遺忘了　神,是我們把　神用自己的驕傲殺死；生命的蒙召是能夠去引導,用生命來引導,乃至於生命可以走在平安的路上.

所以,「孩子阿,你要稱為至高者的先知．因為你要行在主的前面,預備他的道路．叫他的百姓因罪得赦,就知道救恩． 因我們　神憐憫的心腸,叫清晨的日光從高天臨到我們.」(路加福音 1:76-78)

公義的　神,透過祂所呼召的子民,在這個世界上為人帶來光照,為人帶來引導,為人帶來教導:「從來沒有人看見　神．只有在父懷裡的獨生子將他表明出來.」(約翰福音 1:18)

歷史的基督,在歷史的教會當中,被聆聽,被看見,被跟隨.耶穌應許:「有兩三個人奉我的名聚會,那裡就有我在他們中間.」(馬太福音 18:20)這偉大的信息是告訴我們:　神在我們中間!

在你的生命中,在我的生命中,在你和我的生命中,在我們彼此團契,連接成基督身體的見證中,在教會的教導裡面,告訴我們生命就是勝利.當我們願意把我們的生命,再一次的願意去潔淨；再一次來到主的聖壇前,領受祂的身體與寶血；再一次的在神面前有回轉,有悔放；在你的生活裡面,家庭當中,工作的場所,我們是「將臨期」被點起的蠟蠋,是這個世代的見證,是「神義論」的答案!

我們禱告:謝謝主呼召屬主的子民,這在一個「將臨期」的日子裡面,我們再一次的立定志向,求主你讓你的子民希望不絕,求主你讓你的子民在一個充滿著虛無,黑暗與混沌的歷史型態的世界當中,能夠如明光照耀,活出生命的勇氣,活出生命的信心,活出生命的盼望,奉靠耶穌基督聖名祈求,阿們!



\newpage

\section{奧秘揭開了}
\label{sec:185}
\textbf{鄧瑞強博士(香港神學院) - 奧秘揭開了}
\newline
\newline
\hyperref[sec:184]{< < < PREV SERMON < < <}
~
\hyperref[sec:toc]{[返目錄]}
~
\hyperref[sec:186]{> > > NEXT SERMON > > >}


以弗所書 3:1-12
\newline
\begin{longtable}{cl}
\hline
\hline
章節 & 經文 (和合本修訂版)\\
\hline
3:1 & \begin{tabularx}{0.7\textwidth}{X} 因此，我—保羅為你們外邦人作了基督耶穌囚徒的，替你們祈禱。 \end{tabularx} \\ \\ \relax
3:2 & \begin{tabularx}{0.7\textwidth}{X} 想必你們曾聽見神賜恩給我，把關切你們的職分託付我， \end{tabularx} \\ \\ \relax
3:3 & \begin{tabularx}{0.7\textwidth}{X} 用啟示讓我知道福音的奧祕，正如我以前略略寫過的。 \end{tabularx} \\ \\ \relax
3:4 & \begin{tabularx}{0.7\textwidth}{X} 你們讀了，就會知道我深深了解基督的奧祕； \end{tabularx} \\ \\ \relax
3:5 & \begin{tabularx}{0.7\textwidth}{X} 這奧祕在以前的世代沒有讓人知道，像如今藉著聖靈向他的聖使徒和先知啟示一樣， \end{tabularx} \\ \\ \relax
3:6 & \begin{tabularx}{0.7\textwidth}{X} 就是外邦人在基督耶穌裡，藉著福音，得以同為後嗣，同為一體，同為蒙應許的人。 \end{tabularx} \\ \\ \relax
3:7 & \begin{tabularx}{0.7\textwidth}{X} 我作了這福音的僕役，是照著神的恩賜，是照他運行的大能賜給我的。 \end{tabularx} \\ \\ \relax
3:8 & \begin{tabularx}{0.7\textwidth}{X} 雖然我比眾聖徒中最小的還小，他還賜我這恩典，讓我把基督那測不透的豐富傳給外邦人， \end{tabularx} \\ \\ \relax
3:9 & \begin{tabularx}{0.7\textwidth}{X} 又使眾人都明白甚麼是歷代以來隱藏在創造萬物之神裡的奧祕， \end{tabularx} \\ \\ \relax
3:10 & \begin{tabularx}{0.7\textwidth}{X} 為要在現今藉著教會使天上執政的、掌權的知道神百般的智慧。 \end{tabularx} \\ \\ \relax
3:11 & \begin{tabularx}{0.7\textwidth}{X} 這是照著神在我們主基督耶穌裡所完成的永恆的計劃。 \end{tabularx} \\ \\ \relax
3:12 & \begin{tabularx}{0.7\textwidth}{X} 我們因信耶穌，就在他裡面放膽無懼，滿有自信地進到神面前。 \end{tabularx} \\ \\
[1ex]
\hline
\hline
\end{longtable}
日期:2013年1月27日

主題:奧秘揭開了

證道:鄧瑞強博士(香港神學院)

經文:以弗所書3:1-12

各位弟兄姊妹,早安.

今日是新的一年的開始,希望各人在神的恩典裡,得享喜樂及平安.

大地運轉,冬去春來,大地從嚴寒回復溫暖,死寂回復生機,這容易引發人產生宗教情緒.宗教學者以利亞德(Mircea Eliade)說,新年是宗教節慶的重要時刻.人在新年裡,能宗教地體會萬物原初的狀態.在這原初狀態裡,神和人在和諧的關係裡,人從中獲得存在的意義.這意義給人勇氣與盼望,去克服現實生命中的種種逆境.新年的時候,我們喜歡許願.許願,可算是對這原初狀態的回憶,是失去樂園後對樂園的回憶,重尋樂園.在許願裡,我們祈求健康,這是對永生的一種渴望.祈求事事如意,這是對樂園裡萬物和諧的一種回憶.人會許願,是因為我們是宗教的動物.我們活在現實裡,但渴望超越現實的永恆.

在這新年期間,讓我們再思想一下,生命的真相是什麼?我們可以渴望什麼?我們在其中扮演的角色又是什麼?

今日講道的經文,是今日經課裡新約書信的

經文:以弗所書3:1－12

1     因此,我保羅為你們外邦人作了基督耶穌被囚的,替你們祈禱．〔此句乃照對十四節所加〕     2     諒必你們曾聽見　神賜恩給我,將關切你們的職分託付我,     3     用啟示使我知道福音的奧秘,正如我以前略略寫過的．     4     你們念了,就能曉得我深知基督的奧秘．     5     這奧秘在以前的世代,沒有叫人知道,像如今藉著聖靈啟示他的聖使徒和先知一樣．     6     這奧秘就是外邦人在基督耶穌裡,藉著福音,得以同為後嗣,同為一體,同蒙應許.     7     我作了這福音的執事,是照　神的恩賜．這恩賜是照他運行的大能賜給我的．     8     我本來比眾聖徒中最小的還小．然而他還賜我這恩典,叫我把基督那測不透的豐富,傳給外邦人．     9     又使眾人都明白,這歷代以來隱藏在創造萬物之　神裡的奧秘,是如何安排的．     10     為要藉著教會,使天上執政的,掌權的,現在得知　神百般的智慧．     11     這是照　神從萬世以前,在我們主基督耶穌裡所定的旨意.     12     我們因信耶穌,就在他裡面放膽無懼,篤信不疑的來到　神面前.

經文的第一句,讓我想了很久.

「我──保羅為你們外邦人作了基督耶穌被囚的」.保羅,一個嚴守律法的猶太人.「為你們外邦人」,一班與自己全無關係的人,成了監獄裡的囚犯.一個人為了一些與自己無關的人坐監.

世上有些人進入監獄,是由於非一般的理由的.

最出名的例子,是古希臘的蘇格拉底,他被控「不敬虔」,並以新的道理敗壞年輕人,因而坐監.所謂不敬虔,就是不再信任當時社會的神明.這些神明,也可理解為社會的權威.蘇格拉底因為尋求真理而抵觸了這些權威,這是死罪.在監獄裡,他的朋友願意幫助他逃獄,但蘇格拉底堅持不逃獄.他認為身在雅典,便必須遵守雅典的法律.法律是正義的.他不能因為有些人誤用法律而違反法律.最後,飲下一杯毒酒,蘇格拉底便死在獄中.他是典型地為了正義的原則而死的人.

德國神學家潘霍華(Dietrich Bonhoeffer)因為與刺殺希特拉有關而入獄,在獄中寫下《監獄書信》,日後成了重要的神學文獻.他為何入獄?他認為作為基督徒,便要承擔民族及歷史的命運,為了他人的幸福,必須放下自己,盡一切力量,阻止悲劇.他最後死在集中營.

前南非總統曼德拉為了終止種族隔離政策,坐了27年監.他說:"For to be free is not merely to cast off one"s chains, but to live in a way that respects and enhances the freedom of others."(「自由不是單單脫下鎖鏈,而是活出尊重及增進他人自由的方式.」)他就是為了別人的自由,他放下了自己的自由.

中國有劉曉波,他也在獄中.他的自辯書,題為＜我沒有敵人＞,裡面有句話說:「我期待,我將是中國綿綿不絕的文字獄的最後一個受害者,從此之後不再有人因言獲罪.」

世上還有很多這些人,他們被投在監裡,我在此不再列舉了.

他們都是體驗到宇宙人生裡某種真理的人.他們認為這真理是重要的.這真理牽涉人類存在的意義.他們堅持這真理,甚至不惜為此被囚.值得嗎?看到真相的人,會認為值得.他們的命運告訴我們,越接近真理,承擔揭露真相的責任越大.越願意承擔這責任,要承受的苦難便越大.要讓人得自由,有時必須失去自己的自由.要讓人得生命,自己必須進入死亡.

保羅,看到生命的真相,也願意為這真相被囚.他稱這被囚是「作了基督耶穌被囚的」(prisoner of Christ Jesus).看來,他視這被囚,是分享了基督的生命和使命.這不是失敗,而是光榮,顯示他正履行神給他的使命,成就神的計劃.一個人明白了真理,便應該承擔這真理的責任,勇往直前.耶穌在這罪惡的世界,活出了非暴力的愛,死在十架上.信仰耶穌的人,也當如此行.

保羅看到什麼真理,令他甘願被囚?

4     你們念了,就能曉得我深知基督的奧秘．     5     這奧秘在以前的世代,沒有叫人知道,像如今藉著聖靈啟示他的聖使徒和先知一樣．     6     這奧秘就是外邦人在基督耶穌裡,藉著福音,得以同為後嗣,同為一體,同蒙應許.

這裡有一個重要的字眼:「奧秘」.因著神的啟示,保羅看到宇宙人生終極的奧秘.這奧秘不是什麼不為人知的秘密,這奧秘是神要世人明白的事情.所謂「奧秘」,就是神的心意,神的計劃.這心意,這計劃,不是隱藏的,不需要人去猜的.神是啟示的神,他會用盡千方百計讓我們明白他的心.神最終極的啟示,就是他來到人間,成為人.早前的聖誕節,便是記念神來到人間的事實.來到人間,為了承擔世人的罪,基督懷著愛與寬恕,他死在十字架上.他為了我們得到自由,他甘願失去自由.他為了我們明白寬恕,他忍受咒詛.他為了我們學曉真愛,他承受仇恨.他觸摸痲瘋病人,藉此告訴我們,這世界沒有人是不可觸摸的.他與稅吏吃飯,藉此告訴我們,這世界沒有人是斷然地被孤立的.他抱起小孩子,一個沒有地位沒有成就的小孩子,然後教導門徒,要學識尊重這樣的一個一無所有的人的尊嚴.他最後走向十架,藉此告訴我們,真正愛他人,是願意為他人犧牲的.這沒有什麼損失,這是生命的豐盛.死亡之後,便是復活.

保羅說:

5     這奧秘在以前的世代,沒有叫人知道,像如今藉著聖靈啟示他的聖使徒和先知一樣．     6     這奧秘就是外邦人在基督耶穌裡,藉著福音,得以同為後嗣,同為一體,同蒙應許.

基督帶給我們的好信息,就是所有人(無論膚色,性別,人間的成就)都在神的恩典裡,同為後嗣,同為一體,同蒙應許.所有人都同為神的後嗣,所有人都可被稱為神的兒女.所有人同為一體,所有人情同手足.所有人同蒙應許,所有人都能享受天國的福樂.

請不要說,你的膚色跟我不同,我高級一點而你低級一點.請不要說,我做經理而你掃地,我高級一點而你低級一點.請不要說,我的財富比你多,我高級一點而你低級一點.

當神看到每個人的無限價值時,我們怎能輕視任何一個人?

當有人需要食物時,請不要說,這不關我的事.當你的同事需要幫助時,請不要說,這是你自己的問題.當你的伴侶需要你的愛護時,請不要說,你先好好檢討一下你自己.

當基督在我們還作罪人的時候便為我們死,好顯明什麼是愛,我們怎能不愛我們身邊的人?

面對基督,保羅頓然明白天地間最大的奧秘.在弗3:4裡,保羅稱之為「基督的奧祕」.原來,我們自己並不是宇宙的中心,我們只是神偉大創造的一個微少部分.原來,我們的生命被神的大愛托著.原來,當這大愛被我們實現出來時,人生的全幅意義才會顯明.

明白這奧秘的人,就不會時時以自己為中心,就不會時時只顧自己的事,就不會時時為自己的吃喝玩樂而煩惱.

有人說,態度決定高度.我說,視野決定深度.

你看得夠遠,你才明白你的位置在哪裡,你的天職在哪裡,你的深度在哪裡?

7     我作了這福音的執事,是照　神的恩賜．這恩賜是照他運行的大能賜給我的．     8     我本來比眾聖徒中最小的還小．然而他還賜我這恩典,叫我把基督那測不透的豐富,傳給外邦人．     9     又使眾人都明白,這歷代以來隱藏在創造萬物之　神裡的奧秘,是如何安排的．     10     為要藉著教會,使天上執政的,掌權的,現在得知　神百般的智慧．     11     這是照　神從萬世以前,在我們主基督耶穌裡所定的旨意.     12     我們因信耶穌,就在他裡面放膽無懼,篤信不疑的來到　神面前.

電影＜蜘蛛俠＞裡面,蜘蛛俠的叔叔Uncle Ben 對蜘蛛俠說:「偉大的能力伴隨著偉大的責任.」(With great power comes great responsibility)有能力的人,面對邪惡的世界,他的能力會要求他承擔更大的責任.當然,他也要排除萬難,願意承擔才行.

保羅在這裡講,「我作了這福音的執事(僕人),是照神的恩賜,這恩賜是照他運行的大能賜給我的.我本來比眾聖徒中最小的還小,然而他還賜我這恩典,叫我把基督那測不透的豐富傳給外邦人.」

保羅自認為卑微的人,但神的恩典臨到他,讓他明白神的奧秘,讓他看到生命的真相,讓他知道生命到底為何而活.恩典伴隨責任.真是天降大任,他成了福音的僕人.視野決定深度,他的視野要求他活出基督徒的深度.

從今以後,他就竭盡所能,促成不同群體在基督裡的和諧,讓不同的人在基督裡找到豐盛的生命,讓所有人明白什麼是生命的大愛.

看到了生命的奧秘,明白這奧秘的重要,他為此願意付出一切.他願意為此被囚嗎?他願意.

在這世上,很多人視野遠大,看到人生的真相,他們也願意為此被囚.他們的視野,決定了他們的深度.但願,這些人早日獲得釋放.

在這新年期間,讓我們看得更遠大,讓我們的新年願望配合這遠大的視野,而不只是「恭喜發財」,讓我們有勇氣為我們看見的奧秘盡上一分力量.

願主的恩典臨到你們.

但願:榮耀歸於聖父,聖子,聖靈.阿們.



\newpage

\section{狂人的福音}
\label{sec:186}
\textbf{陳恩明牧師 - 狂人的福音}
\newline
\newline
\hyperref[sec:185]{< < < PREV SERMON < < <}
~
\hyperref[sec:toc]{[返目錄]}
~
\hyperref[sec:187]{> > > NEXT SERMON > > >}


路加福音 4:14-21
\newline
\begin{longtable}{cl}
\hline
\hline
章節 & 經文 (和合本修訂版)\\
\hline
4:14 & \begin{tabularx}{0.7\textwidth}{X} 耶穌帶著聖靈的能力回到加利利，他的名聲傳遍了四方。 \end{tabularx} \\ \\ \relax
4:15 & \begin{tabularx}{0.7\textwidth}{X} 他在各會堂裡教導人，眾人都稱讚他。 \end{tabularx} \\ \\ \relax
4:16 & \begin{tabularx}{0.7\textwidth}{X} 耶穌來到拿撒勒，就是他長大的地方。在安息日，照他素常的規矩進了會堂，站起來要念聖經。 \end{tabularx} \\ \\ \relax
4:17 & \begin{tabularx}{0.7\textwidth}{X} 有人把以賽亞先知的書交給他，他就打開，找到一處寫著： \end{tabularx} \\ \\ \relax
4:18 & \begin{tabularx}{0.7\textwidth}{X} 「主的靈在我身上，因為他用膏膏我，叫我傳福音給貧窮的人；差遣我宣告：被擄的得釋放，失明的得看見，受壓迫的得自由， \end{tabularx} \\ \\ \relax
4:19 & \begin{tabularx}{0.7\textwidth}{X} 宣告神悅納人的禧年。」 \end{tabularx} \\ \\ \relax
4:20 & \begin{tabularx}{0.7\textwidth}{X} 於是他把書捲起來，交還給管理人，就坐下。會堂裡的人都定睛看他。 \end{tabularx} \\ \\ \relax
4:21 & \begin{tabularx}{0.7\textwidth}{X} 耶穌對他們說：「你們聽見的這段經文，今天已經應驗了。」 \end{tabularx} \\ \\
[1ex]
\hline
\hline
\end{longtable}
講題:狂人的福音　A Madmans Gospel

經文:路加福音4:14-21

講者:陳恩明牧師

場合:崇基禮拜堂

日期:2013年1月27日

各位早,謝謝伍牧師的介紹,我覺得好像參加安息禮拜一樣的. 亦覺得他彷彿像介紹一個不務正業的牧師似的,其實我真的是牧師,我是很喜歡分享福音,所以今天都是跟大家來分享福音,何解我叫它作《狂人福音》呢?

有時候我覺得我們教會,可能因為種種的原因,壓力,或許因為我們是歷史悠久之後,漸漸變得頗有體面,甚至是有財有勢,漸漸便變質了,是沒多少實質的力量,也沒什麼敢言的表達,也許甚至我們都有種信主日久,有了雞肋的味道,食之無味,棄之可惜；不上禮拜堂又怕給伍牧師罵,回來又不知道是為了什麼因由…

其實我們是很需要重拾那福音的火.一些沒有溫度,沒有熱力的東西便會硬化,僵化. 到了我這把年紀,我開始明白血液循環是十分重要的,為何你到一個階段會手腳冰冷?為何有時候,有些動作在你年青時候是很容易做到的,現在做了之後,很容易就會有後遺症,好像屈腰後就要按摩,要看看的,因為那循環的熱力不見了!

我希望主今天給我們恩典,雖然是很熟悉的東西,但是我們內裡真的需要重新燃點起自己內心的火熱,以致我們整個教會,仍然能在世上為主發光,發熱.

狂人?其實不用伍牧師介紹,我現在是來稍為借助一個人,作小小的自我介紹,如果你記得馬可福音第五章,有一個狂人,真的是頗瘋狂的,他是相當之有力量,而沒任何人能制服他,甚至人用鐵鏈捆鎖他的手,他的腳,他都能夠完全掙斷它. 我覺得這種 power 是非常的厲害,我看見這狂人的力量,雖然他有力量掙脫人們加在他身上規範,但他沒能力運用他的力量來保護自己,他反而會被自己的這種力量傷害,因為他也會拿起石頭砍自己,受傷以後,又再砍… 是很慘的一個情況,而且這個人－－他是自絕於人間,與家人脫離了關係,跟其他人沒任何的交往,他在墳墓堆裡出入,終日走來走去,狂叫大喊.我看見這人,我還看到一件事情,他是毫無遮攔的,他是衣不蔽體的,他是赤裸裸的,不知羞恥的,一個多麼破碎,多麼混亂,多麼墮落,多麼暴力的,甚至是沒什麼道德的一個人,這個人一點添加的價值都沒的,其實本身就已是在社會上被認為沒有價值的人,直到有一天,這受盡捆綁,破碎的人,見了耶穌基督,他說:「至高神的兒子,你現在來幹什麼? 你現在來幹什麼?!」 我覺得這仍然是一種不少人對耶穌的態度:「你現在想怎樣了? 我其實不需要你介入我的生命當中,你最好給我留多點空間!」 但是這位至高神的兒子,祂卻是繼續迫近這滿身邪惡力量捆綁的人,耶穌沒退縮,因為祂是愛；耶穌沒退縮,因為祂捨不得我. 耶穌恢復了我的身份,耶穌叫我知道我叫甚麼名字,耶穌叫我知道我真正的目的,我在世上生活的使命究竟是什麼?聖經描寫這馬可福音第五章的狂人,我覺得是我人生的寫照.這個人清醒過來,穿上衣服,坐在耶穌腳前.

也許我交代一下,自我介紹,我差不多是第三,第四代的信徒,但到了我青春期發動的時候,到我看見這社會上很多邪惡,很多解不開的苦罪的問題,當我看見人生最大的極限都是墳墓的時候,當我看見自己內心的黑暗,又在教會裡聽說了很多的規範,對我來說,那些好像是重擔；我看見我不斷地在逃避上帝,最終我真是徹底逃避祂,所有任何令男人快樂,威風的東西,我都覺得是我真正能夠用來麻醉自己生命的東西,我幾乎是毀掉自己的一生,但是耶穌愛我,使我清醒過來,穿上衣服,坐在祂的腳前.

再說這個狂人,這狂人離開了家,但被耶穌拯救之後,他想到處跟隨耶穌,耶穌卻叫他回家去.所以今天我到處跟男人說,不要不穿衣服,不要暴露,別傷害自己,別爭權奪勢,但那些權勢是你自己操控不了的,別被捆綁,然後我跟他們說: 「各位男仕,最要緊的是回家.」 當你回家看見你的所謂「黃臉婆」,原來是你幼年所娶妻子,她今天之所以變成這樣,完全是你一手造成,所以你要回家,擁抱她,摟她,愛她,就是這狂人,被耶穌改變成為一個清醒的人!

親愛的,你覺得這是否好福音? 其實這是不是我們每一個人在上主面前的一個縮影呢?原來這位耶穌,祂給這狂人知道,在耶穌眼中,這二千頭豬與這個在社會上一文不值的人來說,是一文不值的.其實福音是真的很奧妙,很精彩,這位耶穌肯為我們付出這代價,根本這是令我們沒辦法不五體投地來感激,一生一世都要來愛慕祂!

另外一個使我今天要分享這訊息的人,他也很瘋狂的,這人跟格拉森那狂人一點也不同,這人是有教養的,有唸書的,有地位的,甚至有道德的,有宗教的,不過他心中是非常狂熱的來追殺一些他認為犯錯的人.直到有一天,這位是如此苦苦追迫基督徒,將他們致諸死地的人,耶穌說「你用腳踢刺是難的!」使徒保羅問道:「你是誰?」耶穌說:「我就是你所逼迫的耶!」一位是狂熱地逼迫耶穌的人,後來他說了一段說話,說後就有羅馬的大官,甚至這個王,這位官長大人大聲的說:「保羅,你瘋了,你的學問太大,所以使你瘋狂了!」保羅說:「 非斯都大人,我不是癲狂,我說的乃是真實,是明白話. 」

各位,究竟誰對?使徒行傳第二十六章二十四節,二十五節,這如此簡短的對話,我覺得值得我們深化,想想今天我們和福音的關係是甚麼一回事,想像在我們生活中有沒有讓人覺得你都是有些另類的,有點兒瘋瘋癲癲的, 我不知道你聽後的感覺如何,即便說耶穌基督的教會在這世人的眼目中,覺得你們其實跟我們一式一樣,甚至比我們更差,而並非覺得你們這些人有些不同,不同之中甚至有點兒癲狂的.

你明白我在說什麼嗎?我很希望今天能引起一個討論,引你來想像一下,究竟要怎樣為主而瘋? 保羅說如果我們癲狂是為了主,原來是基督的愛激勵我們.如果我們清醒是為了你們,那瘋癲起來也還是有點兒分寸的.但基本上他講的,他信的,他作的,在某些人眼中實在有些瘋狂,人家說他瘋了,他說我沒瘋,那不就是典型精神病的情況. 他說我講的都是真理,其實這場辯論到今時今日還是進行中,那麼我們不妨溫習一下,究竟我們的這個福音是怎樣的一回事?

就在保羅被告在羅馬官的跟前,然後他堅持一定要上告到凱撒那兒,結果就有阿基帕和伯尼基兩人來聽保羅訴說他的心聲,整章聖經使徒行傳二十六章,我們可以慢慢細看,這讓我們看到教會今天所信的福音是甚麼的一回事.保羅說:「我說的其實一點也沒有新意,我沒做任何的壞事,不過他們定是要來告我,我之所以站在這裡受審,(使徒行傳二十六章六節:「是指望神向我們祖宗所應許的；)」這應許,是我們全以色列民族十二個支派 , 晝晝夜夜,引頸以待,我們晝晝夜夜,殷切事奉的神 , 向祂所祈盼的,就是因這指望. 那代表什麼? 原來保羅所說的,就是上帝的啟示和應許,指望著神應許給他們百姓,其實這指望亦都是給全阿當的族裔,阿當的族裔的絕望是什麼?  就是因與上帝隔絕而陷入死亡的捆鎖當中,而上帝賜予人類那一個最終極的盼望,親愛的,就是 “H O P E” 盼望,盼望,盼望,今天仍需要盼望.

二次大戰結束後,邱吉爾在他的辦公室裡向當時尚是年輕的葛培理提出一條問題:「青年人,請問那裡有盼望?」其實今時今日我們所持守的,我們所期待著的,仍然是同一個盼望,神叫死人復活,如果我們真明白叫死人復活內裡所包含的是甚麼,我相信我們的教會會有動力的,我們的生命會有能力,我怎麼會知道? 因為我因著這個盼望,我跟我二十二歲之前那種灰色絕望的人生,已經變得不同了,雖然我現在靠近我最後的日子,比我二十二歲的時候已經將近了四十年,但是我的盼望,卻是比我昔日是精彩得多!

你有沒有這盼望? 你是否傳講著這盼望?然後保羅他繼續講, 我在追殺基督徒的路上,就在這裡,那位耶穌基督向我顯現,祂說:「(徒26:16)我特意向你顯現,要派你作執事,作見證,將你所看見的事和我將要指示你的證明出來.」「我也要救你脫離百姓和外邦人的手. 我差你到他們那裡去,要叫他們的眼睛得開,從黑暗中歸向光明,從撒旦權下歸向上帝……」(徒26:17-18)耶穌基督是親身的顯現,然後向保羅說,我要派你,我要用你,我要救你,這是何等奇妙的工作.

耶穌的啟示,是整個教會的根基,今天有很多人還在研究,到底耶穌是誰? 研究來,研究去,自己拼湊出來的一個耶穌,始終沒有一個面貌,但唯有那一位親自復活,向門徒顯現,向使徒保羅顯現,並且吩咐他們,繼續在這世上為祂作見證,這一位上帝的兒子耶穌基督拿撒勒人耶穌,就是我們福音的核心.如果福音不是在說一位死而復活的拿撒勒人耶穌救主,並非十二使徒所見證的那位耶穌基督,那福音,無論有多少學問,都是一個沒有能力的福音.

一位活的耶穌基督,是今天仍然向人講說話,仍然向人來啟示,仍然可以被信服,被倚靠,仍然將祂的恩典,祂的能力賜予給所有信靠祂的人,然後我們看見這福音,除了是根據應許,是本著這活的啟示；還有一個很清晰的內容,福音是什麼?包括不少在教會很久的人,他都會說:其實也都是教人做好人而已,究竟我們有沒有唸書?究竟我們聽清楚了沒?是叫你眼睛得開,從黑暗轉向光明,從撒旦權下歸向神；是一個生命主權的回歸運動,是一個有罪的人和聖潔公義慈愛的上帝,一位活的神關係的恢復.  “Encounter between two living persons” 因為我們是被造,按著神的形像樣式造的,我們是一個活的人,是跟一位活的神來恢復關係,其實你知不知道福音最重要是什麼?

恢復上帝的榮耀,恢復上帝的榮耀,恢復這位創造主的權威和榮耀,離開這一切,我們所說的,其實都是不充足的.當然我們不是教人做壞人,但當我們不夠膽跟別人說,你是要跟這位創造主,而耶穌基督是那位創造未來的君王,你當如何面對祂?保羅告訴我們,是一個管治領域的改變,本來是被黑暗,被自我,被撒旦管轄的,但現在轉變了.

我想很多人都有置業的經驗,如果你簽了那張契約的時候,那麼你就擁有那樣東西. 福音就是將那張契約完全的轉名,轉回給原本有資格擁有我們的那一位,卸罪的恩典,所以保羅所說的,其實是整本聖經一直啟示的,是眾先知和摩西所講的說話,他也是說這些事,我的祖先所期望的死人復活,我的耶穌基督親自向我顯現,祂教我要做這些事,我亦真在大家眼前做了這些事,非斯都聽見後都說:「你瘋了!」 究竟為什麼?其實他是否聽明白了保羅所說的,只有一位創造主? 是否他聽明白了這拿撒勒人耶穌,是死而復活,是榮耀的主?是否他聽明白了保羅是有權威的,來為主作見證?是否他聽明白了那邪魔是失去勢力?是否他聽明白－－有一個信仰傳遍天下? 他的反應那麼大,我們何不想像一下是何解?他坐堂是有權威的,站在跟前的是一個可以聽後他發落的人,他是羅馬帝國中的人,坐在面前的是一個以色列亡國的人,你這亡國奴來傳揚一個上帝的國度?站在跟前的也是一個必死的人,你還沒曾死過,竟在傳揚復活的盼望?你說的是你看不見,摸不著的,非我經驗裡的東西,我都認為按我的理性被否定的,所以你是瘋的.那究竟誰瘋,誰不瘋?最終都是看耶穌基督究竟是否聖經所說的那一位.

所以我們回到最基本的一件事,就是我所信的福音是什麼一回事,我的福音難道只是叫我在這世上可以催吉避凶?平平安安嗎?我的福音是叫我死後榮登極樂嗎?儘管有東方和西方的分別.還是我們的福音,就是那一個讓耶穌基督真正進駐到我們的生命當中,讓上帝的榮耀,上帝的國度,上帝的主權湧入我們整個人生,然後說我一生無悔,是全心全意忠於你.

親愛的,讓這個上帝的榮耀,上帝管治的能力,上帝的釋放,上帝的卸罪,上帝使我們生命改變,保羅在他的見證中,他說我領受了這個職份,就是向猶太人和外幫人,任何人,無論男女老少,尊卑長幼,我跟他們講的訊息都是一樣,要悔改,要悔改,行事為人與所蒙的恩相稱,行事為人與悔改的心相稱,這就是與神和好的真理,榮耀復活,這個榮耀復活在聖經整體的啟示中乃是說,耶穌基督要作整個宇宙的王,所有一切都要在耶穌基督裡復興過來,成為美好,甚至可以說,比原來的創造更加精彩,因為是一個全新的創造,所以我們的福音,不是一個單單關乎個人生命的福音,是關乎整個人類宇宙的福音,而那核心是一位活的救主,無論如何我們要忠於祂,要見證祂.

保羅其實還沒說完的,不過非斯都也不想聽下去,不過保羅並沒有放棄,他說:「王啊!我知道你知道很多事情,或者在座也有很多人也都知道很多事情,但是好像這個阿基帕人一般,都好像在旁邊聽聽,沒有那個 commitment 承諾,我希望今天那位來尋找第一百隻羊的主,那一位來到請撒該從樹上下來,說我要進你家與你住的主,如果你看見你的生命裡仍然有陰暗,仍然有盲點,仍然有不信,仍然有邪惡的捆綁,不如你試著來找耶穌?」

如果你說我已經屬於耶穌,其實你知道在過去那麼多年當中,誰是真正的主,真正的,你內心最重要的是那一件事情,不如你今天說:「主啊!榮耀的主,你配得作我的主,作我的王,然後試一下,試試從今以後,你的人生是否真的有火?今天很多的人,不少的人被財,被色,被自我迷惑,他也說信耶穌,不過耶穌始終都是踏腳石,我今天的訊息就是,如果你說我是瘋的,我說我希望我是為耶穌而瘋的,耶穌從來不是一塊踏腳石,從來不是,耶穌從來都是至尊的宇宙的主宰,君王,祂能夠叫我們榮耀,改變得像祂,保羅在他的環境,你看看,辛辛苦苦面對著審判,他仍然向他的審判官來作出規勸,想人信主,你想不想?如果不想,你自己只顧自己的事,不多管他人,即便說我們已經正常化?已經不瘋了?即便說這些教會在等待著倒閉.但是如果我們說,主啊!你已經開了我的眼睛,你已經使我從黑暗進入光明,你已經使我主權易主,我覺得真是好!所以我無論往那



\newpage

\section{公義如高山判斷}
\label{sec:187}
\textbf{伍渭文牧師 - 公義如高山,判斷如深海}
\newline
\newline
\hyperref[sec:186]{< < < PREV SERMON < < <}
~
\hyperref[sec:toc]{[返目錄]}
~
\hyperref[sec:188]{> > > NEXT SERMON > > >}


日期:2013年2月3日

主題:公義如高山,判斷如深海

讀經:詩篇36

證道:伍渭文牧師

最近因為報章深入偵查報導,揭發香港政府統計處前線調查員作假,他們到訪一層樓,發現內裡tong 開數間房就佯稱沒有人,因為上級視一個門牌一個個案,訪問數個戶主才當一個個案,不化算.統計數字不正確影響很大,因為很多政府的政策,規劃都基於統計數字；政策錯誤就浪費公帑,影響民生,延誤社會發展；若採用不真實的材料,偷工減料,豆腐渣工程,可引至塌樓斷橋,人命傷亡.

信實很重要,小朋友信任父母,當你答應過和他去街或賣東西給他們,他們就信以為真,期待你履行諾言,所以不要隨便答應他們.自然現象有可預計的規律,才有科學的研究和發展.從事商業活動,一定要有信譽,信守契約,不然經濟活動就癱瘓.我們需要有一個信實,可以信賴的世界.

今天的詩篇就強調上主的信實,但他的信實不是一個自然神論(Deism)的信實,像一個鐘表匠,上了發條,讓世界自然運作,遠離人世,對人的苦難無動於衷.詩篇三十六篇指出,上主的信實,就是他對人不變的慈愛.上主看到人在罪中的苦痛和苦難,與人立約,作他們的上主,施行拯救,被救贖的人要成為他的子民.上帝是信實的,會履行他的約,拯救保護他的子民,這是上主信實的慈愛.

詩篇三十六篇採用了常用的交叉配列(chiastic)方式:首尾呼應,中間對等,突出上主的慈愛和信實(ABB"A").首段1-4節提到惡人「自誇自媚」(2節),「床上圖謀罪孽」(4)；尾段11,12節「不容驕傲人的腳踐踏我,不容惡人的手趕逐我.在那裡作孽的人,已經仆倒.」中間對上主的頌讚提到「慈愛」,祈求也提到上主的「慈愛」.很明顯,上主的慈愛對比惡人是詩篇三十六篇的主題.

聖經提到惡人時,一般和義人相比較,重點是他們的結局,而且用智慧傳統的議論文體來描述.智慧傅統(wisdom tradition)是用從創造論,不用救贖論的立場,以冷峻眼光來洞察世情,舊約箴言,傳道書,新約雅各書都是用智慧傳統的文體寫成.詩篇第一篇善惡的比較是典型的智慧傳統,惡人將像糠秕被風吹散；糠秕散發,看似金碧輝煌,但風吹過後無覓處,而義人像一棵樹,栽在溪水旁,按時候結果子,葉子也不枯乾.但詩篇三十六篇用他們與上主的關係作比較,而且運用不同的文學體裁:1-4節以第三身,用智慧體裁描述惡人；5-12節以第二身,用詩歌體裁頌讚上帝 (5-9) ；最後10-12節以第一身的祈禱結束.

1上主慈愛信實不變

信實的「切約」

慈愛這字貫穿整本聖經,指上主與人立約拯救世人的慈愛.這約不是平等協議,保護彼此利益的契約(contract), 是宗主國對其蕃屬宣布其主權和關係的約(covenant),主動來自宗主國,立約原文稱「切約」(cut the covenant), 主動立約的一方,把祭牲切開,從其中經過,寓意願意以生命作為保證,保護蕃屬；作蕃屬的也履行敬愛宗主的責任,立約強調彼此的親密的關係.上主與亞伯蘭立約,引領他離開吾耳(Ur),叫他成為大國,後代將如夜空星宿之多,亞伯蘭把牲口劈開,「日落天黑,不料有冒煙的爐和火把從那些肉塊中經過.」(創15:17)立約的一方以生命來保證,履行責任,如果爽約,就像祭牲被劈開,聖經中就常說耶和華以他的永生來起誓,保證他的慈愛信實不變.

救贖的慈愛

最早出現約的字是創6:18節上主與挪亞立方舟之約,以方舟拯救挪亞一家八口,洪水過後,並以彩虹為立約記號,永不用洪水審判世人.(創9:9-17) 從挪亞到亞伯拉罕,以撒,雅各,約瑟,摩西直到耶穌基督,上主都對他子民守約施慈愛.當以色列人在埃及法老治下為奴受苦,上主聽到他子民的哀聲,看到他們的困苦,興起摩西,帶領他們離開埃及,「讓我的百姓離開」(Let my people go).上主對摩西說將會「給他們(以色列人)興起一位先知像你.」(申18: 18)耶穌登山變像時,代表律法和先知的摩西,以利亞出現,談論耶穌去世的事(路9: 31) ,耶穌去世,原文就是「出埃及」,耶穌就是那位像摩西的先知.天使對約瑟說:「她(馬利亞) 要生一個兒子,你要給他起名叫耶穌,因他要將自己的百姓從罪惡裡教出來.」(太1: 21) 上主信實守約,拯救他的子民,這是他的慈愛,而惡人就是拒絕上主慈愛的人.

2. 惡人拒絕上帝的慈愛 vv. 1-4

「惡人的罪過在他心裡說:我眼中不怕上帝!他自誇自媚,以為他的罪孽終不顯露,不被恨惡.他口中的言語盡是罪孽詭詐；好他與智慧善行已經斷絕.他在床上圖謀不善道,不憎惡惡事.」

不同的懼怕

「我眼中不怕上帝.」(1節b) 不怕,英文可繙作不尊重,no regard.懼怕有多種:若你欠人錢拖延不還,見到債主你會懼怕,但這怕與你見到一獅子不同.又如你景仰德蘭修女,若她尚在人間,突然進入崇基禮拜堂,你有一種懼怕,這是一種敬畏,早期教會信徒彼此相愛,「恆心遵守使徒的教訓,彼此交接,擘餅,祈禱.」(使2:42) 眾人都懼怕.為何眾人都懼怕?因為他們看到人可以如斯共融相愛,他們看到一種神聖的景象,與平常割裂的景象,而懼怕.我們懼怕上帝,因為他神聖可畏,雖然可畏,但有一種令我們響往的吸引力,不同靈異的幽暗世界,令我們恐懼.罪使人漸漸失去愛的能力

「惡人的罪過在他心裡說.」(1節a) 這裡說罪是一種屬靈的實體,像電影中的異形(alien) ,不屬於自己又屬於自己的實體.罪教唆惡人,不須敬畏上主,罪惡在惡人心中盤踞,毒害他的靈魂,使人心靈失去愛慕美善,聖潔的能力.與此相反,保羅提醒提摩太:「凡是真實的,公義的,清潔的,可愛的,有美名的,若有甚麼德行,若有甚麼稱讚,這些事你們都要思念.」(提腓4:8)

一個對上帝不虔不敬的人,對人也會不忠不義；失去對上主的愛,也會失去對人的愛.他白天說話盡是罪孽詭詐,睡覺時在床上圖謀罪孽.「他不行善事,不憎惡惡事.」(詩36: 4)為何他不憎惡惡事呢?為何他對不義的事不感到忿怒呢?因為他對罪惡的受害人的痛苦無動於衷.所有的罪惡,都有受害人(victim) .十誡中不可殺人,不可姦淫,不可偷盜,不可作假見證,不可貪戀人的財物,都是對人沒有愛心,所以耶穌說,誡命的總歸就是愛.

詩人描述惡人,以不懼怕上帝開始,以不憎惡惡事結束,說明惡人因為靈魂被罪侵蝕,失去愛的能力,失去愛慕真,善,美,聖潔的能力,也失去憎惡罪的敏銳,失去對被罪壓制的人的同情心.

3. 使我們飽足的慈愛vv. 5-9

從第三身冷峻的的智慧文體描述惡人,跟著筆鋒一變,作者用感情洋溢的詩歌體裁來頌讚上主和祈禱,對比惡人.

「你的慈愛上及諸天,你的信實達到穹蒼；你的公義好像高山,你的判斷如同深淵.上帝啊,你的慈愛何其寶貴!世人投靠在你翅膀下.他們必因你殿裡的肥甘得以飽足；你也必叫他們喝你樂河的水.因為,在你那裡有生命的源頭；在你光中,我們必得見光.」(vv. 5-9)

從慈愛到公義

「你的慈愛上及諸天,你的信實達到穹蒼；你的公義好像高山,你的判斷如同深淵.(vv. 5, 6)

惡人自以為罪孽終不顯露,不被恨惡；詩人就頌讚上主公義像高山,上主的判斷如深淵；上帝的公義像高山立於前,毫不動搖；但上帝的作為像深淵,不是我們可以猜度,他何時審判,如何審判,不是我們可以和道的.慈愛和公義有何關係呢?公義是愛的均分,我們追求公義,因為愛每一個人,保護他們的權益,不要以為上主對罪孽袖手旁觀,對惡行好像默然無聲.

在你翅膀的蔭下-母鷹愛護雛鷹的關係

上帝啊,你的慈愛何其寶貴!世人投靠在你翅膀下.(v.7)

詩人跟著讚美上主信實的慈愛和保守.摩西在以色列人快要進入迦南美地時,在約但河東的摩押平原重申誡命(申命記Deuteronomy, second law),回顧歷史,說到上主在曠野引領保護以色列人:「耶和華遇見他在曠野荒涼野獸吼叫之地,就環繞他,看顧他,保護他,如同保護眼中的瞳人.又如鷹攪動巢窩,在雛鷹以上兩翅搧展,接取雛鷹,背在兩翼之上.」(申33: 11)

但有解經家說,這是聖殿敬拜的語境.「在你翅膀的蔭下」暗示在聖殿天使撒拉弗(seraphim)的個翅膀.以賽亞先知進入聖殿敬拜:「我見主坐在高高的寶座上．他的衣裳垂下,遮滿聖殿其上有撒拉弗侍立．各有六個翅膀．用兩個翅膀遮臉,兩個翅膀遮腳,兩個翅膀飛翔．彼此呼喊說,聖哉,聖哉,聖哉,萬軍之耶和華．他的榮光充滿全地.」(賽6: 1-3)

聖殿肥甘, 樂河之水-在敬拜中得著滿足

因為跟著第八節:「他們必因你殿裡的肥甘,得以飽足．你也必叫他們喝你樂河的水因為,在你那裡有生命的源頭.」.(8, 9a)

聖殿的裡的肥甘指獻祭祭牲的脂油,獻平安祭完畢,獻祭的人可分享祭品,形容與上主的親密關係,人在敬拜中得到最大的滿足.

今天是顯現節第二主日,詩篇經課提到喝你樂河的水,在經課的安排是有特別禮儀的意思.

樂河的水是生命的源頭,在創造的新世界,在伊甸園,有生命樹,有河從伊甸流出.(創2:9,10)這生命河在救贖的新天新地再度出現:街道當中一道生命水的河,明亮如水晶,從神和羔羊的寶座流出來.(啟22:1)

上主日講員提到耶穌受洗,受洗或受浸一字原來有把水聚在一起,創造天地時,地是空虛混沌,淵面黑暗,上帝把水分為上下,把水聚在一起,露出旱地,才生發生命,所以洗禮表示新生的開始.耶穌在約但河接受施洗約翰的洗禮,表示與人認同,站在人的地位,與人聯合,從水裡上來時,天開了,天父為耶穌作見證,有聲音從天上來,說:「你是我的愛子,我喜悅你.」(路3: 22) 「你是我的愛子,我喜悅你.」是彌賽亞的稱號,從此耶穌公開傳道,走上十字架.我們接受水禮的人,是與耶穌聯合,和他同死同埋葬同復活,和他聯合,得著新生,成為教會一份子,而教會是基督的新婦,所以今天福音書經課約翰福音21-11是一個婚宴,慶賀基督和教會的聯合.基督在婚宴中行了第一個神蹟,使水變酒,而且是最好的酒,表示基督是樂河的水,生命之源,祇有基督才能滿足我們.耶穌受洗後聖靈彷彿鴿子,降在耶穌身上,我們受洗的人,也領受聖靈的恩賜,互為肢體,彼此服事,這是今天書信經課林前12:1-11節的信息.

在你的光中,我們必得見光 v. 9b

在敬拜中,在聖言的亮光中,詩人有了真實的知識,不像惡人自誇自媚,自欺欺人.上主的慈愛,在基督裡拯救了我們,使我們經歷了真實的愛,也愛曉愛人,因愛人追求公義.我們也有真的視野,認識自己的身份.是上主的子民,基督的身體,聖靈的團契,我們被召,是要建立基督的教會.

最後的祈禱vv.10-12

「願你常施慈愛給認識你的人．常以公義待心裡正直的人．不容驕傲人的腳踐踏我,不容兇惡人的手趕逐我.那裡作孽的人,已經仆倒．他們被推倒,不能再起來.」(vv.10-12)

一輪頌讚後,詩人作一簡短的祈禱,求上主敗壞惡人的計謀,保守自己脫離兇險.我們不用懼怕惡人,因為惡人已被徹底打敗,不能再起來.詩篇中的惡人,當然是實有其人,不懼怕上帝,也不憎惡惡事,不愛慕上主,也不憐愛人,無視罪惡帶給人的痛苦,聖經也是用惡人,表示罪惡的力量,「惡人的罪過在他心裡說.」(1節a) 罪就是拒絕愛,拒絕神聖和美善.但罪惡已經被上帝的愛所擊敗,基督為我們的罪,包括詩篇三十六篇所描述的惡人的罪,被釘十字架,上主像高山的公義滿足了,我們不明白為何上主竟然用這方法來拯救我們,他的判斷如深淵,超乎我們的想像.想到這裡,讓我們和保羅一起讚美(羅11: 33-36):

「深哉,上主的智慧深哉,　神豐富的智慧和知識.他的判斷,何其難測,他的蹤跡何其難尋,誰知道主的心,誰作過他的謀士呢,誰是先給了他,使他後來償還呢因為萬有都是本於他,倚靠他,歸於他．願榮耀歸給他,直到永遠.阿們.」

當我讀到「在你的光中,我們必得見光」時,我都想起早期北菲迦大基(Carthage,現今突尼斯)教父特土良(Tertullian,160-225 A.D.) .他是通曉希臘哲學和羅馬法律的著名學者,被聖經的陌生感覺所吸引,聖經像一度清泉,解除他長久以來希臘哲學不能給他的屬靈饑渴.他在光中看到光,信主以後,他把自己的學問獻上,成為知名的護教家和神學家,被譽被西方神學之父.

今天的詩篇告訴我們,世上最恆久不變的是上主信實的慈愛,這慈愛已由基督為我們釘死在十字架顯明了,我們也領受聖靈的恩賜,被差遣傳揚福音,建立上主的教會.阿們.



\newpage

\section{義無反顧的獻身}
\label{sec:188}
\textbf{梁燕城博士 - 義無反顧的獻身}
\newline
\newline
\hyperref[sec:187]{< < < PREV SERMON < < <}
~
\hyperref[sec:toc]{[返目錄]}
~
\hyperref[sec:189]{> > > NEXT SERMON > > >}


日期:2013年2月10日

主題:義無反顧的獻身

證道:梁燕城博士

中國著名思想家劉再復,十九年前,他曾經在一篇文章寫道,有一次他走到德國的科隆大教堂(Cologne Cathedral),世上最大的教堂之一,那裡非常莊嚴寧靜,又聽到很美的詩歌,他便忽然有很大的感動,認定世上有一種神聖美善的價值存在.卻因為他曾經歷到中國的文化大革命,又經過中國改革的頭十年,當中有名的學者都經常受到批評和攻擊,在這種情況下,他說他們身上都有很濃的硝煙,很重的火藥味.但來到教堂裡,這種火藥味都被洗掉了.這世間應該是有美善神聖的世界,可惜我們都被困在現實的政治社會,種種複雜的處景中,失去了那份寧靜,那種美好.

大約在公元第四世紀的時候,在土耳其,當時的東羅馬帝國,有一個偉大的老修士,他原是一個富人,但當他經歷了上帝以後,便願意放下他所有金錢來幫助貧窮人.因為羅馬帝國當時的經濟狀況很差,所以當他把家財都捐出來之後,仍然有很多窮人,他發覺有很多小朋友沒有玩具,便自發做一些小玩具.每逢在聖誕的前一夜,他便背上一個袋,帶著那些小玩具送給貧窮的小朋友.趁他們睡著的時候,放到他們的旁邊.

這讓我記起,在中國剛開放的時候,還有很多沒鞋穿的小朋友,那時候我在想,如果可以讓每位小朋友都有一隻玩具熊,那多好呢.之後,中國的小朋友都幾乎可以擁有玩具熊.當時我有第二個夢,如果中國人的房子都有草地,可以有花園便好了.到現在,亦開始有很多人可以擁有了.不過我在想,如果中國下一步能有更多的公平,貧窮人得到更多的幫助,中國能發展成一個均富的社會,那便好了,這是我對習近平總理的夢.在江澤民總理時代,我希望每位小朋友都能擁有玩具熊；在胡錦濤總理時代,我希望中國人可以住得好房子；而現在我的夢想是均富,讓大家都分享到好處.

說回那位老修士,要改革社會是一件難事,但他把所擁有的獻出來,每逢在聖誕的前夕,他便背上袋,把玩具派發給小朋友.據說當時有三位小女孩,由於家貧又欠了責,家人要把她們賣了作妓女,她們既悲傷卻又很無耐.當老修士知道之後,他便想辦法籌了一筆錢,在聖誕前夕,偷偷地把那筆錢,從窗外扔進她們家裡.據說,當時她們家裡的火爐旁,掛了幾隻襪子正在烘乾,而那份禮物剛好掉進襪子裡.第二日,亦即是聖誕節那天,正準備被賣作妓女時,當她們穿上襪子,發覺裡面有筆錢,剛好足夠給她們贖身.她們便歡天喜地,便知道上帝幫助她們,後來都知道是那老修士幫助她們.直到老修士死後,有很多小朋友好像仍然見到他拿著一個袋,把袋內的玩具送給小朋友.相信大家聽到這裡,都知道老修士是誰,他名叫尼古拉斯 (Nicholas).後來流傳到荷蘭,變成 Santa Clause, Santa解作「聖人」,Clause則是Nicholas的簡稱.之後又去流傳到美國,「可口可樂」汽水在廣告中,把Santa Clause 畫成一位白髮白鬍鬚,穿上紅衣的胖子,便是今天我們看到聖誕老人的樣子.可能有很多人都不知道,聖誕老人原是一位虔誠的基督徒,他願意把自己的生命獻上給上帝,又樂意付出所有來幫助貧困人.他的生命來自他親身經歷的上帝,當他親歷上帝之後,便義無反顧地委身.

今日我有一段經文與大家分享,羅馬書十二章一至二節:<br /> 「所以弟兄們,我以　神的慈悲勸你們,將身體獻上,當作活祭,是聖潔的,是　神所喜悅的；你們如此事奉乃是理所當然的.不要效法這個世界,只要心意更新而變化,叫你們察驗何為　神的善良,純全,可喜悅的旨意.」

保羅在這裡勸勉基督徒,他說你要獻出你自己,成為一個活的祭,不單單是心靈而已,是要具體地把自己的整體獻出來,成為一個活的祭.在舊約時代裡,獻祭是要將無殘疾的牛羊獻在祭壇上,是要完全燒在壇上,為要完全獻上.但到了新約時代,就不是再用動物作獻祭,而是要把自己的生命作為祭,是一個聖潔的活祭,是　神所喜悅的.

我在想,在古時獻祭是要把牛羊切成一塊一塊的,然後完全地把它燃燒,獻在壇上,把它破碎.但在我接受耶穌基督的救恩之後,並沒有想過把自己生命獻上給　神.因為我是追求真理的人,十三歲開始閱讀莊子,十四歲讀尼采和柏拉圖,十五歲讀佛學,然後學習打坐,之後藉著打坐,開了天眼,看過另一邊的世界,當打坐到達終極的時候,我領悟到宇宙萬物有份慈悲仁愛,而這份慈悲仁愛是可以與人接通的,並且接通到宇宙最大的愛.

從中我有三重的領悟,第一個,我明白到宇宙是變幻無常的,一切皆空,這是佛教的學說.第二個領悟是,變中是有不變的,因為變化之中是有規律的,這是道家所說的道,同時亦是易經的學說.而第三個領悟,就是當人看到他人的痛苦,是會感到不安,所以人是有仁愛的心,而仁愛就是宇宙的本質.

宇宙──充天塞地都是偉大的仁愛,於是,我去尋索這份仁愛背後的最終安排者.究竟是誰安排宇宙萬物有規律?是誰使萬物由無變為有?雖然萬物是空,但它是有規律,而且其中又有仁愛放在人的人裡.這到底是誰?我知道這是中國古人所相信的「皇天上帝」,古人所謂的「天」,即是「人性的根源」.於是我去找尋這位「皇天上帝」,結果我找到基督教所稱的「上帝的啟示」,即是耶穌基督曾來到人間.於是我去考察他的可信性,之後我又發覺他藉著苦難,由一位苦難的上帝,來取代人類的痛苦和罪惡,從而啟發人類的身心靈.我最終尋到一位與人同甘共苦的上帝,所以明白到「上帝的啟示」就是耶穌基督的救恩,於是我接受耶穌基督,成為基督徒.

可是成為基督徒之後,我仍是有很重儒釋道的觀念.為什麼會信耶穌呢?只因為我明白他是真理,他能滿足我對真理的追求.在中文大學讀書時,我是一名怪人,因為要與其他新潮的年青人相反,追求「舊潮」,當時我愛穿長袍,幾乎要留辮子,又愛手拿寫有智慧詩句的扇子,而且又在大學的新亞校園裡打坐.當時有著名的教授牟宗三,同樣也是穿長袍的,其他人碰見我,以為我是另一位什麼教授,但其實我只不過是一名大一學生而已.我又刻意買來一把油紙傘,因為我要做一名古人,覺得現代人都是很膚淺的,立志回復古人的智慧.另外,我讀的是古書,而所買的是用線裝的那種,為要卷起來看,又拿紅筆來點句.之後大學搬到山上去,如果晚上穿上長袍的我走過,被其他人看到那被風吹起的長袍,著實被嚇壞.

當時我自命高人一等,因為我練過功,體驗過宇宙,又懂得古人的智慧,廿歲便講學儒釋道.當時其他人都質疑我,認為六十來歲才懂的,為何我會明白.又有人認同我的學問.亦有人以為是因為打坐的原故,使我本來由六十歲變成廿歲的外貌.當時校園裡有位有名的基督徒學生－溫偉耀,我告訴他,我終於相信耶穌基督,他都不大相信.因為之前我經常與他們辯論信仰問題,基督徒都不會儒釋道的學說,所以他們一定不能把我辯倒.溫偉耀是因為我的刺激之下,在及後的日子,由物理學改讀中國的儒釋道.說回當時那些基督徒帶我去查經班,誰知道我用佛教解釋聖經,把他們都嚇壞了.後來,他們邀請我去中文大學第一屆聯合夏令會,當時我大三,因為還沒有女朋友,我想可能在夏令會裡結識女孩子,覺得基督徒比較善良,而且在大學之後認識的女孩子,好像不太純樸,所以決定參加.

當時我帶著線裝的《朱子大全》參加夏令會,邊看邊用紅筆點句,有時又會拿扇子撥著,旁人還以為我是什麼講員.當時的夏令會的主題是「將你的生命主權交出來」,所研讀的是剛才的一節羅馬書.你的生命是否一個祭可以獻出來,是否一個活的祭獻給上帝.那時侯的講員是滕近輝牧師,滕牧師的外貌很屬靈,好像只要讓他站在講台上,便能夠得著復興,在他身上有聖靈,所以我很被感動,雖然現已忘了講道的內容,只記得他當時的容貌很像聖靈般.

接著的活動便是查經,所查的是那段是關於一位少年官去見耶穌,請教他怎麼才能得到永生,耶穌便叫他變賣所有的財產,並分給貧窮人,然後跟從祂.可是少年官感到很為難,於是便憂憂愁愁地離開.讀到這裡,我覺得我可以,只是變賣財產而已,因為我什麼都沒有,當然可以瀟脫地奉獻給　神.

有一個笑話,上帝問一個人,要他奉獻屋子,車子,他都能爽快地回答可以,但當上帝再問他,要他所擁有的一隻雞,他卻說不可以.於是上帝問他:「為什麼剛才屋子,車子都可以獻給我,為什麼一隻雞卻不可以呢?」他回答說:「因為我未有那些東西,只有那隻雞,所以不可以啦.」

當時的我同樣什麼都沒有,而且學過佛,對金錢與權力並不執著,所以我本來就可以放棄這些東西.當時領查經的是溫偉耀,他是一位很理性的,很有學識的基督徒,因為大家知道我是一個愛辯論的人,所以安排他作我的組長.溫偉耀解釋這段經文的意思不是在於變賣家財,而是在少年官而言,他最執著的便是他的財產.但在個人身上而言,可能是另外一些事情.今天,對你來說,什麼東西是讓你最放不下呢?在那一刻,我要面對幾個問題,什麼是委身,什麼是奉獻自己,什麼是活的祭?原來就是將自己心中最執迷的把它放下.

佛教都談放下執著,而基督教的核心也是說要放下執著,因為沒有執著才能坦承的交給　神,被衪使用.於是我心裡問上帝,我到底有什麼放不下,我明白到,是我的「學問」.我以為熟識儒家,佛家和道家的哲理,以及那些線裝書,便是我的財產.但如果上帝要我把這些都放下來,對我來說是很難為的.我一生的榮耀都是從這些東西得來的,我並沒有什麼,就只有這一些,這是我的驕傲.要我放下這一切來跟從你是很難的,而且我可以藉著這些學問,我也可以幫助神的工作,為什麼你要我放下呢?但當時上帝很清楚,要我放下最執著的.我以為最重要的學問,知識,自以為高人一等的,我自以為練功很高,這一切都要放下.可是一直想也想不通,直到最一天的晚上,我發覺與上帝之間被隔開了,好像尚欠一些才去到衪那裡,我很想與上帝有更深入的溝通.只是我知道不能依靠打坐去到上帝那裡,而是要先潔淨我的罪.可是我的罪雖已被耶穌基督潔淨了,但與上帝還有一段距離,因為我未把自己獻上.

當時溫偉耀也在夏令會的操場上,正思考自己的前途問題,我攔著他,因為想知道他是否能真實經歷　神,他回答是真的,是很真實的.我說:「為什麼我已經信了耶穌,卻沒有像你們那樣開心.你們基督徒個個都充滿喜樂,我卻與你們有差距呢,不能達到你們的喜樂.」於是溫偉耀回答說:

「當你想自己愈多,你便想上帝愈少.」

真是當頭棒喝的一句,原來這麼多年,我一直只是為了自己,我追求這麼多學問知識,只是為了炫耀自己.我信上帝,只是因為衪是真的,目的為要滿足自己對真理的追求.可是我卻沒有想過,信上帝是怎樣的一回事.雖然明白要喜樂,亦都感到一點的快樂,卻感覺仍有些距離,仍未達到上帝那兒.那一刻,我終於明白,上帝並非一盞神燈.我們祈禱,並不是擦一下神燈,裡頭走出一個大胖子,說著:「主人,有什麼吩咐.」我們總希望擁有一盞神燈,有位大神仙來幫助我們.所以我們很容易把最高的　神,當成為一個神仙來幫助自己.人們拜神,其實都是為著滿足自己,這些是虛假的民間神明.我們事實上要追求的,是創天造地,在最後掌管一齊的真神.所以並不是等我們去吩咐衪,而是我們去問祂:「上帝,你有什麼吩咐?」

「上帝,如果你真是存在,你是宇宙的主宰,你看透人間的悲歡離合,你明白世事有這麼多的艱難困苦.主,我想獻給你,但我相知道應該作什麼,我的人生應該怎麼樣,使你的愛能行出來.你的愛就是你的國度,當你的愛能行出來,就是你的國度在地上彰顯出來.」

那是我突然明白,信上帝並不是要來滿足我.原來是我要主動問衪,在世間有這麼多問題,上帝把我放在這裡,有什麼目的呢?衪賜給我這麼多知識與才能,到底是為什麼呢?我終於明白了,就在翌日的清晨,太陽初升的時候,我步出那個夏令營,在外面有一條小溪,我對著初升的太陽,禱告說:

「天父,我知道你真的存在,因為沒有你,宇宙就是偶然的出現.我現在願意將生命獻上給你.我不但接受耶穌基督成為我的救主,我要獻出我的生命,成為你可以用的人.

世界有著許多的苦難,(我當時是身處文化大革命,越戰的時代),很多殘忍的事,而我只是很卑微的一個大學生,只是比較讀多一點書而已,我不知道將來可以做什麼,我就是要把生命獻出來,我可是將所有的學問知識放棄,交給你,所有未來的前途,都由你帶領.使到我今日的存在,是可以讓這世界的黑暗,多一點的燭光,可以安慰多兩三個人.或者我的存在,可讓冰冷的世界多一點溫暖.天父,就讓我成為一個你可以用的人.」

我的內心本來有一道圍牆把自已圍著,就是以為自己很厲害,比很多人都高尚.但在那奇妙的一刻,圍牆突然溶解.原本的我是自我中心的,但突然之間溶解了,上帝進入我的生命,聖靈好像一道洪流湧進,那份溫暖慈愛就這樣湧進我的生命.就在那一刻,我明白上帝確然是真實的,怪不得,基督徒都因信神而瘋狂了,到處要傳揚衪,縱然被其他人辯倒,都不會認輸和動搖,因為他們真的經歷到上帝,就是那份慈愛,洪流以及神聖.我也經歷到了,明白基督徒為什麼為神而瘋狂,因為我也為之瘋狂了,而且從那時開始,一直瘋狂至今.我今天在此講道,並不是為了賺錢,因為講道其實蠻費精神體力的,那為什麼還要來呢?就是因為我為主瘋了的原故,是因為經歷到那份宇宙的仁愛,美善和寬恕,那份內心的觸摸.我就是將那份愛帶給世界,給予更多的人,將那份愛行出來,使這醜陋的世界加點美麗,使這冷漠的世界加點溫暖.就在那時候,我把生命都獻給　神.可是　神很奇妙,衪不但沒有叫我放棄學問,反而要我再讀多些,要我將來修讀儒家,道家和佛家至最高的層次,才能幫助到中國的.通常基督徒獻身後,都去修讀神學,但我獻身之後,明白上帝的指示,就是我不用放棄這些學問,而是修讀至最高的層次,將來我便會白箇中原因.所以我沒有選讀神學,而是去學易經,並成為博士,而論語,孔子都是修讀博士學位的研究內容.至於道家和佛家,我都繼續研讀,並且讀到最高峰,比如佛家的華嚴宗與禪宗,來了解他們最高的智慧.

最初獻身的時候,我想過立刻去到非洲宣教,在那裡的原野對著當地的土著,講述自己的見證,不過他們卻用矛把我殺掉,然後我便浪漫地倒在大原野上,為主而死.另外又想過,進入當時正在文化大革命中的中國,向紅衛乒傳福音,如他們要打我右邊,我連左邊也給他們打,及後被打至幾乎死去,我會臨死前說:「我會寬恕你們的.」我這份人,是有些英雄性格,常幻想英雄式般死去,因我最祟拜的就是文天祥,一邊赴刑場一邊唸詩,我也渴望這樣子死去.可是上帝感動,我還未開始替衪工作,怎可就此死去,去修讀好最高層的學問才是我的責任,所以我便一直讀下去.

之後人生也有不同的關口,在廿來歲的時候,獻身是件易事,因為什麼都沒有.但到了三十來歲,讀完博士學位之後,被一些教授看重；在社會上,又開始有了些名氣.當時心裡突然想到,現在已為成人了,可以作些詭詐的,為何不可去看看一,兩齣色情電影,為何要這麼純良,反正在美國沒人認得我.而且有時以真誠待人,他卻在背後傷害自己,發覺原來真誠待人是會受到傷害,那我是不是還要繼續以誠待人呢?

就在那段留美讀書的日子,有一晚,上帝在半夜把我弄醒過來,問道:「你在那裡?」就如昔日問亞當夏娃般,我回答說:「我正在往上爬,因我知道我的前途很好,我會很有名氣,很了不起,所以我要往上爬.」上帝再問道:「你可否放下這一切,來跟從我.」我立時再一次決定好,回答說:「主,你是這麼真實,我一定要放下這一切,來跟從你.所有一切的前途,我都不要了.如果能夠榮你的名,我就會去作；如果不能夠榮你的名,我就寧願捨棄.如果你要我不能完成博士學位,我也願意.只要我可使你的愛,送達至那些有需要的人.只要有人得著福音的好處,我是否可以成為博士,得到什麼社會地位,這都是次要的.」

就在那一次,我將生命再一之獻上.跟著,我又經歷到奇妙的事,就是一次厲害的趕鬼經驗.從此之後,我感受到上帝來臨之後,那一種非常龐大的能力.所以翌日早晨,我仍不斷想前晚到底發生什麼事,邪靈世界的事情和上帝的真實,於是我走到海邊,思想著上帝是這麼真實,我定要一生獻給祂.可是當我真誠對待上帝的時候,我不可以以虛假待人.我可否單純地待地人,但人是會傷害人,人是很醜惡的,有些人會在背後陷害自己,縱然如此,我可否仍舊保持著單純的心.

我知道上帝的能力強大,並且掌管一切,我交給衪是無問題的.所以我向上帝說:「偉大的天父,我寧願受傷害,也要以真誠待人,因我對你要真,對人亦要真,我要做一個單純的人.」

到了四十二歲的時候,我被人陷害得更厲害,那是我更大的學習.因為當單純的人,便會受到他人攻擊.最差的情況是,所遇到攻擊你的人,會隨著年齡增長而越厲害,他們都是好像好人的壞人；當你越有地位,陷害你的人,地位亦越高.當二,三十歲的時候,陷害你的人,算不得什麼.但到了四,五十歲時,陷害你的人,便變得越來越厲害,越來越奸詐.

我是否都要變成奸詐去應付這樣的人呢?但我仍選擇要堅持單純去待人.在四十二歲那次事件中,我要去學會如何去寬恕陷害自己的人,真的不容易.還沒有仇敵的時候,我也懂得說要去愛仇敵.但當有人傷害了自己,還可以去愛他嗎?真是非常困難,雖然我明白耶穌基督要我去愛他,但那個人真的壞透了,其實也想運用自己的能力去報復.因為我有一家的名氣,可利用傳播媒的影響力來攻擊他.但在最後我堅持不作出報復,因為我不希望上帝的名因此受到污染.

雖然決定不作出報復,但心內仍然氣憤難平.我雖不公開報仇,但可否私下報復呢?在加拿大那裡,冬天很寒冷,我大可在他家門前,潑下一盆水,因為太冷而在地面結冰,使他滑倒在地,又預先放下圖釘 and \#8230;人其實可以有很多毒辣的思想.

但我向上帝祈禱,說:「主,如果我心裡有憤恨,我是不可能服侍你的.當我要站上講台,講述你的愛的時候,怎能同時心存憤恨呢?可是每當想起這個人的時候,我的心便憤怒.」很奇妙,當時我有位好朋友,名叫楊慶球牧師,當時在剛剛在從香港教完神學,返回渥太華,路經多倫多來探望我家.那時剛好是感恩節,我們一起結伴外遊到海邊去靜思.到了星期日,雖然與楊牧師同行,是可以由他來主持一個小聚會,但聖經說:「你們不可停止聚會.」,所以還是決定到附近的西人教會有正式的聚集.那間教會聚會人數不多,牧師向我們說:「今天是感恩節,大家低頭想一想,今年你們有什麼是值得感謝的.」當我低下頭便想,今年沒有什麼值得感謝,全部都是被人陷害,被人貶低,就好像被人用腳踩在頭上,我還要舔他鞋底,這樣的遭遇,我怎能感謝上帝呢?但上帝在我心內說:「雖然如此,經歷了這麼多艱難痛苦,但你可不可以仍然感謝我?」在那時,我突然領悟到,當廿,三十歲被人陷害時,是沒有大問題的,因為自己和加害自己的人,都是普通人而已.當四十歲的時候,自己有了些名氣,加害自己的人,手段都會較之厲害,讓我們完全不能招架.但如果上帝是掌管天地,衪必然知道這一切,然而衪仍然容讓這些事發生,一定是有箇中意義的.

如果我沒有把生命獻給上帝,我是沒有權向上帝討論.但現在我已把生命獻給上帝,但衪仍然讓這些事發生,我便明白,這一定是好.只是我當時未明白是好的,所以我心才會憤怒.但我終於明白,衪原意是好的,我便不再憤怒,並且接納了這件事.原來我們的憤怒是因為我們不接納,就是不接納事情會這樣發生.我們不明白為什麼信任一個人,他竟然會在背後陷害自己,還故意佈局來苦害我.但我終於明白了:「主,只要你仍然在,這件事一定是好的,縱然我現在不明白,但我先感謝你.」

當我接納了這件事之後,便不再憤恨了,而且可以寬恕他,因為我知道一定會更好的.我不可再以仇恨來綑綁自己,我要再次站起來.我頓然明白,耶穌教我們愛和寬恕仇敵,原是要釋放我們內心的憤怒.原來是要釋放我們,讓我們再次站起來,讓我們能重新開始做我們該做的事,而不是便宜那害人的.

後來我遇到楊牧谷博士,他與我分享一節聖經,「在我敵人面前, 你為我擺設筵席」﹝詩23:5﹞正當打仗的時候,我們被敵人包圍著,為何上帝突然說要吃飯呢?以前不明所以,但現在我明白了.「我雖然行過死蔭的幽谷,也不怕遭害.」﹝詩23:4a﹞死蔭的幽谷不是單指在以色列的欣嫩子谷,那兒是有很盜賊.也可指在沙漠地帶,有些好像井一般的地形,因為下雨之後,那裡便成了好多水井,吸引很多羊來渴水,可是一不小心便很容易滑入水井之中,不能爬上來,死蔭的幽谷就是這一種的狀態.人生有時候就好像這樣子,掉入井裡,縱然努力掙扎,但總不能重新站起來,就在此時,牧羊人伸出他的手杖,心裡便立時很快慰.「因為你與我同在,你的杖,你的竿,都安慰我.」﹝詩23:4b﹞其實在我最艱難的時候,你是一直與我同行的,我發覺原來基督教最核心的,就是耶穌與我們同行這段路.他一直就在我們身旁,他的杖一直就在旁邊來安慰我們,所以縱然掉入井裡,我知道我一定能夠走出來；敵人來包圍的時候,你會擺設豐盛的筵席.

就在1993年,即是四十二歲那件事之後,上帝對我真正的使命終於來臨.上帝呼召我,要我回去中國,就是當全世界都在指責她的時候,批評她的時候,看不起她的時候.那時候因著89年民運之後,鄧小平南巡,大家對中國的前途充滿疑問,就在這時上帝對我作出呼召,這個民族受傷的心,她需要是醫治,無數的海外人士寫文章,不段地批評和咒罵,連連的示威,是對那兒的人沒有幫助的.對中國裡面千千萬萬受傷害,貧窮的人,你要去幫助他們.甚至共產黨都受到很多的傷害,你要用愛去醫治.但對我來說,有兩種人我是很難去愛的,一是共產黨,二是日本人.可是耶穌都愛他們,他們都需要被醫治和受安慰.最近我兒子教導我,要去愛日本人.

我那時明白,上帝的愛超過世界這一切的政治,經濟和社會的醜惡,是不帶有條件的.只要肯行出愛來,中國是會有改變的.於是我們便開始文化更新的使命,在中國則從九四年成立開始,到今天有十八年的工作,當我們進入國內,看到有很多的困難和需要,我們要用盡一切來幫助她.而第一個幫忙的項目是反腐敗,學廉潔.我們帶他們到加拿大了解,什麼是沒有腐敗的國家,他們學到什麼是審計系統,雖然中國腐敗的情況未能徹底解決,仍然存在很多問題,但至少學到利用審計來對付腐敗.

跟著的工作,是我們去到農村,看到很多貧窮的小孩子.那時上帝讓我看到,尼古拉斯 (Nicolas) 因看到有三個貧窮的小女孩快要被賣成為妓女,於是出手救她們出來.而當時中國剛剛興起,九十年代至二千年,很多少女要出城工作,做農民工,大約十五歲完成初中學業,經常被人欺凌,被人壓制.那些長得較為漂亮的,很容易會被引誘去賣淫,成為妓女.當我了解這情況之後,感到心很傷痛,因為想到自己的女兒同樣是十來歲,但因為那些女孩的家境貧窮,要出來謀生養家,或是為了供養弟妹讀書,而誤入歧途,若然換了是我的女兒,又會怎樣呢.國內千萬的父母,因為要出城工作,無暇教導兒女,造成很多問題.所以我當時去到北京見副總理兼民政部部長,他對我說:「我們沒法子幫助妓女,因為中國政府不承認國內有妓女,只有直至她們染上愛滋病或者毒癮之後,才可以幫助她們.」但是我覺得那時候已經太遲了,我想有沒有方法可以阻止她們陷入這處景.他說其實不太難,只要能幫助她們繼續上學便成.

於是我們在加拿大籌得二萬多,換成十四萬人民幣,當時幫助到一百四十個學生完成學業.只要能完成學業,他們便有前途,就這樣展開了關懷農村工作.除了到大學講學之外,我開始到農村了解他們的貧窮境況.後來有機會讓外交部的領使,得知我們對農村社會的改革建議,便介紹我到國務院反映意見,後來又介紹到黨中央,一步一步地,讓我能夠與政府直接講話,便將我見到的國家困難,民間的疾苦,貧窮人因缺乏教育與醫療而遇到很多苦難.跟著很快,國家便推行九年免費教育,之後醫療開始改革.又因著與政府表達過農民工的苦況之後,政府開始為他們建房子.我們盡努力為他們,到現在已幫助了六千人上學,一直至大學畢業.

現在我回望過去,我發現生命原來要不段獻上,當作活祭.正如聖經說:「不要效法這個世界,只要心意更新而變化,叫你們察驗何為神的善良,純全,可喜悅的旨意.」就是要心意更新變化,當我最初回國內時,只是希望透過文化,精神文明的建造,道德的價值與中國作出交流.之後發覺,道德原是要行出來,所以我去到農村幫助有需要的學生.之後發覺有機會向中國表達意見時,就要去盡量表達出來.所以這十幾年是我人生最快樂的時刻,縱然錢我是賺得最少.當時一九九三年,有一間電台機構,願以一百五十萬的年薪來聘請我作評論員,即價值現在約二,三百萬,對我來說實在很吸引.想了很久,但想到早國受了百多年的傷害,被欺凌侮辱,她實在需要愛的醫治,而且上帝的呼召比一切更重要,早國十三億的靈魂重建,都比一切更重要.所以我終於作了一個決定,獻身做這份正作.縱使薪酬遠低過一個教授,甚至牧師,但我仍然很豐足,因為上帝仍然給我充足的物質環境去生活,而且上帝給我的喜樂更是用錢也買不到的.我一生最大的願望是,能夠幫助到中國,這個曾經受苦的國家,而我看到她在轉好中,雖然仍然有很多問題和艱難.我有機會能夠參與陪伴同行這條路,將我們要說的告訴他們,盡了作中國人的責任.盡了能力來幫助六千人完成學業,又因為我們愛他們,他們懂得如何愛其他人.最近文化中心的學生拍了一段影片,他們是傳理系的大學生,拍下自己的影片來感謝我們.他們說自己從前只是曚曚懂懂的學童,因為貧窮而遇到很多困難,而沒有機會完成學業,卻多特你們這班民間的亞姨叔叔來幫助我們.

我們其實付出的金錢不是很多,卻因為我們陪他們一起走,對他們的擁抱,對他們的愛是最寶貴的.因為他們的父母出城工作,沒有空與他們一起；又因為父母文化水平不高,當他們讀大學時,不知道誰可以解答他們的問題.所以,他們把我們這群義工當作父母般,向我們問意見,我們又鼓勵他們前進,又在經濟上幫助他們.當中有些人,他們已經當上醫生,老師,甚至農村的人大主席,是經由直選出來的.看到他們不同的發展,我們都很高興.在他們最艱難的時候,我們伸出援手,就是這麼簡單,他們的一生就改變了,成為對社會有供獻的人.並且在中國最強大的時候,他們剛好出來,正好需要有愛心,道德,人格的時候,希望他們在這方面有所供獻.

他們學會無條件的獻身給上帝,學會無條件的愛人,而我們盡能力使這世界的冰冷小一點.

讓我們低頭祈禱:「親愛的天父,當我們看到天地何等的美麗,卻又看見人間這麼多的衝突,這麼多的可怕的事情.在美國,有人無故把小孩弄死,看到人心變得多麼黑暗,封閉又自我中心,天父啊,我們都這一切感到很悲痛.當我們又看到中國仍有千千萬萬的貧窮人,想奮鬥向前卻又缺乏幫助,我們感到很悲痛,可惜我們能力有限,我們不能幫助每一個人.

天父,我求你激發我們的心,讓我們完全的獻給你,讓我們成為活的祭,是燃燒的火,不會停止,讓我們供你永遠使用,讓我們生命流出的是你的愛,讓我們獻出心中最執著的,讓我們都有資源,能力,才華為你所用,無條件的奉獻給你,因此奉獻給有需要的人,求你的愛成為我們的能力,求你的愛帶領我們走每一步.

願主復興教會,願聖靈臨到我們當中,賜給我們能力,激發我們向上,讓我們明白到生命的重要和寶貴,讓我們的祭是活得有喜樂,充滿慈愛,我們感謝你.禱告奉主耶穌基督的名祈求,阿們.」



\newpage

\section{愛的洗禮}
\label{sec:189}
\textbf{陳榮超牧師 - 愛的洗禮}
\newline
\newline
\hyperref[sec:188]{< < < PREV SERMON < < <}
~
\hyperref[sec:toc]{[返目錄]}
~
\hyperref[sec:190]{> > > NEXT SERMON > > >}


日期:2013年2月17日

主題:愛的洗禮

證道:陳榮超牧師

聖誕節記念耶穌基督降生,福音書中從耶穌出生後很快便跳到他的福音事工.我們會發現在四福音中提及耶穌生平是很稀薄的.首先是在伯利恆出生,然後在兩歲的時候,約瑟及馬利亞帶著耶穌逃命,走到埃及,所以由兩歲至三十歲之間的耶穌生平是極之少.四福音的主要目的不是偉人生平傳記,並不是描述一個人的出生,青少年,怎樣學習讀書等等,這些資料很少,少的原故更見四福音的重點,耶穌基督工作的開端由洗禮開始,所以在紀念耶穌基督顯現的時候,福音的重點從洗禮開始,今天早上一開始的經課就是耶穌基督在約旦河河邊出現便是福音工作的開端,為什麼在談耶穌基督顯現會以這為經課?

因為耶穌基督的洗禮正是福音工作的開端,我們要思想福音開端是如何開展  and \#8211; 耶穌基督的洗禮.我估計大部份坐在這裡的都是已受洗的信徒,亦有一組人已預備了七七八八的心情準備洗禮儀式,其中亦有少部份為慕道的,在心裡感到基督教信仰可以幫助他,協助他在生命上有所進一步.三類朋友對今天的經文可能有不同的感覺.當提到經歷洗禮的時候,我們當中有些已洗禮了幾十年,老信徒,盼望今天能去想想自己的洗禮和耶穌的洗禮的共同點,或認同點是如何?而其中一些而準備好,接受人生中非常隆重而嚴肅的儀式,在耶穌的洗禮上應有什麼回應.我們崇基禮拜堂已接受報名,大家也要趕快一點,洗禮快到了.另一班朋友對基督教信仰很感興趣,認為基督教信仰可以幫忙解決人生問題,或者家庭上的需要,究竟耶穌的洗禮對你們又有什麼共鳴?盼望今天的信息能讓信徒思想耶穌的洗禮與自己的洗禮的共同點和共鳴處是什麼?以致我們向前走時有一個開端的更新,在人生的起點中,讓耶穌的洗禮幫助我們.這是希望今天的信息中的目的.

路加福音第三章提到有一班猶太人(百姓)心裡盼望彌賽亞,救世主的降臨,當時有施洗約翰的出現,他們都想知道施洗約翰是否那位彌賽亞.但施洗約翰否認,並表示他雖然是以水替他們施洗,但有一位比他能力更強更大的,連提鞋也不配的,因為他是以聖靈和火替他們施洗,然後耶穌基督出現,耶穌也受了洗.正禱告的時候,天就開了,聖靈降臨在他身上,形狀彷彿鴿子；又有聲音從天上來,說:「你是我的愛子,我喜悅你.」耶穌基督受洗後有很特別的景象出現,有聲音說:「你是我的愛子,我喜悅你.」在信息中會作詳細解釋.首先回想耶穌基督來到約旦河受洗的情景.耶穌大約三十歲的一天,沒有指定那一天,他來到約旦河.那天,他決定要進行洗禮的儀式,那時已有很多人進行洗禮儀式,在那段日子中,常常有洗禮儀式,有些人因著施洗約翰的信息而來洗禮.施洗約翰比耶穌大六個月,並且有親戚關係,經文沒有提及時間,但那時正是施洗約翰傳講悔改信息,告訴百姓需要有一個悔改的洗禮,所以每天也有一大群人在約旦河邊接受洗禮儀式.關於施洗約翰的水禮儀式,猶太人自古已來已經有,他們到約旦河洗禮並不出奇.猶太人稱之為潔淨禮.在耶穌基督那個時代,有錢人家會在家中建一個水池,大約有半身的深度,而且盛滿水,若家中有人想行潔淨禮,便會走進其中,把人浸泡水中,然後起來抹淨,便是行潔淨禮.在舊約傳統,摩西律法書中,如果在街中踢到或接觸到屍體,便需要行潔淨禮.耶穌出世八天需要行割禮,但母親在過了四十天潔淨後,這是指猶太人的潔淨禮,就是在水中浸一下便能潔淨,有些人得罪了人,心裡很不舒服,便浸在水中行潔淨禮.所以這樣的水禮在當時是相當普遍的事情,有些人可能一天早,午,晚三次,因為經常犯錯,或者經常接觸到不潔淨的東西,有不潔淨的事在心裡,所以便行潔淨禮.在聖殿旁的牆邊也建了一些水池,幫助那些敬拜的人在敬拜之前先到水裡進行潔淨,覺得需要時便找聖殿的祭司,行潔淨禮後藉祭司有赦罪之權而相信罪得赦免.

施洗約翰到約旦河河邊,有些人覺得自己生命有不潔淨的地方,心裡覺得要到水裡得潔淨,能夠得到更新.在施洗約翰的時代,那些百姓到約旦河邊進行一般的潔淨禮,在施洗約翰的信息,他們感到有點不同,他們發現施洗約翰正是他們在心裡需求了很久的信息,他們心裡想施洗約翰是否他們所盼望的救世主和基督呢?經文說:「百姓指望基督來的時候,人都心裡猜疑,或者約翰是基督.」他們希望有新開始有新指望.為什麼令民眾帶來影響?因為信息中帶來兩個信息,在過去傳統聽不到.第一個是針對當時宗教人士在潔淨禮中佔了便宜,耶穌說:「安息日是為人設立的,人不是為安息日設立的.」的信息帶有相同意義,潔淨禮不是為人設立的,人不是為潔淨禮而設的.傳統讓我們為了儀式而儀式,忘記了儀式不是為人而做,人卻為儀式而做.所以當時的人為了這儀式,家家戶戶也建水池.祭司亦因為能夠赦罪而有貪污的情況,所以施洗約翰的信息含有針對當時的不公平的,把潔淨禮的意義淡薄了,只為了減輕心中的不安,因此施洗約翰把潔淨禮的意義從新更新,他並不是到人的家中,也不是到耶路撒冷聖殿前,而是在約旦河.約旦河代表活水泉,所以施洗約翰提到我們要以活水潔淨而不是一盆死水.第二要針對的是祭司赦罪的濫權.潔淨禮沒有赦罪的意義,卻要進行一個悔改的洗禮,不需要到祭司面前尋求赦免,而是到神面前悔改,這信息對耶路撒冷聖殿祭司有很大挑戰.所以我們能夠會意當時的百姓對施洗約翰的信息有這麼大的反應.因為對舊禮儀有新觀點,舊儀式有新看法,對一個習慣出現了新的動力.當百姓在約旦河進行洗禮,明白這跟往常的潔淨禮不一樣,因那是活水江河,可以洗淨我們的罪,因為大家存悔改的心,不為禮儀而進行禮儀,而是生命的更新,因此百姓聽了信息後,每天都有不少的人到約旦河進行洗禮,在那段日子有一天,耶穌基督就在約旦河接受洗禮.耶穌基督的洗禮帶給我們更深的層次.

在這裡先跟大家談談對洗禮的更深的看法,洗禮在希伯來文翻譯為潔淨禮或水禮,這個字本身是說把水聚集在一起的意思,指在水聚在一起時便可進行水禮.這字第一次出現在創世記,神創造天地時,天地混沌黑暗,神更在其中,如果我們與神一起時,便會知道那時周圍都是水,在第九節中,神把水聚在一起成為海,沒有水的地方是旱地,神完成這創造後,其他植物便開花結果,所以猶太人認為神把水聚在一起是一個生命的創造,因此每個進行水禮的人便會想起神創造生命的開始,把水聚集,神就讓生命開始.

在大學團契時,久不久便會到湖邊進行洗禮.記得曾經協助弟兄姊妹行進湖中較深水的地方,另外有兩位弟兄在其中協助信徒洗禮,那時在美國較北的地方,天氣很冷,有一位姊姊受洗後,從水中行上來,她讚嘆世界的美麗,鳥語花香,我四周觀看,還未到春天,天色灰灰暗暗,怎會有如此景象,所以到今天還很深印象,她看的跟我所看的完全不同.對我來說很深印象,她看到的我看不到,她聽到的我聽不到,神很恩待這位姊妹,她來自渺無人煙的地方,是美國土生華裔,中文並不通,來到大城市裡讀大學,我們有機會邀請她查經.有一天,她早上六點便拍我太太的門,太太在睡眼惺忪的時候開門,這位姊妹說她很想信耶穌,想受洗,神感動她要服侍人,很想先學中文,當時學中文並不容易,她的心意很大,她去當兵,因為那裡有言語學堂,她是一個很大心志的女孩子,信耶穌後,生命有新的開始,從窮窮鄉僻壤去讀大學,希望得到成就,神改變她的生命,讓她有新的開始,心裡有感動而願意服侍中國人,連當兵也願意,為了達成她的心志,所以每次去洗禮或者替人洗禮,便會想起1973年所認識的這位姊妹.

耶穌三十歲的時候,不知哪一天,來到約旦河,想到神的創造.根據考古學發現耶穌在約旦河受洗的地方,就是約書亞帶領以色列人渡過河到迦南地的地方,耶穌想到在這個地方,神的子民藉約書亞帶領到新的境界,就是以色列國在迦南地第一次得到的應許,是新的創造開始.施洗約翰把神新的權能國度權柄帶到以色列人當中,耶穌基督接受水禮是證明施洗約翰所說神的國度權能的開始,他實踐這開端,並把施洗約翰的信息在他的洗禮成為開始.水禮中耶穌基督提到人以悔改的心參與其中,生命不再像從前,人的生命已是神創造的,從前我在教會洗禮後會唱一首歌,音樂就是Happy Birthday,但歌詞不一樣,Happy birthday is two, only one will not be you, take Christ as your Savior, then you will have two. 所以接受了水禮的信徒,在那天神在你生命中是一個新開始,祂不是要替你算舊帳,生命是完全新的開始,在水中神創造我們的生命,祂把水聚集在一起時,當人在其中,一個新國度,一個新子民,一個新生命,一個新開端.在報名水禮時,是神在你生命中的一個新創造,如在徘徊在生命中,就不要錯過水禮帶來的應許.

耶穌的洗禮,他來到約旦河旁邊,慢慢走進水裡,沒有人帶他的,跟據過去的傳統,是施洗約翰陪著他,走進水裡替他施洗,但在猶太人的傳統,所有進行水禮的都不需要人相伴,各人憑自己心意走進水裡,在水中禱告,將心意交託,然後自己安靜地把心放在心上,慢慢蹲在水中,讓水浸過頭,然後回復自由.聖經記載耶穌從水中上來時,有聲音說:「你是我的愛子,我喜悅你.」聖靈降臨恍惚如鴿子,這是耶穌行水禮的大約情況.中文翻譯的受洗被認為是被動詞,表示需要有人幫忙的,這字翻譯出來有這意義,但實在不需要這被動的意義.耶穌的水禮是主動的,可能有人在旁邊自己進行水禮,就好像潔淨禮般,安靜在水裡,與神一起,向神禱告把心意交上,自己讓水蓋過頭髮,然後慢慢走上來,上來時帶來特別的能力,聖靈的降臨,有聲音從天而來,是外在的肯定,有聲音說:「你是我的愛子,我喜悅你.」這話讓我們明白水禮不但是潔淨過程,還是認罪得赦免的過程.潔淨禮本身不需要因有大罪而進行,耶穌基督的水禮帶來新信息,認同施洗約翰的信,這水禮帶來更深的意義,就是愛的洗禮,在其中耶穌基督展現了他來到世間的任務,建立愛神和互相相愛的群體,不單人的罪得赦免,這很重要,因為有罪便不能與神親近,我們需要潔淨,心裡很多害怕,人生帶著很多恐懼,需要神的憐憫,潔浸淨和幫助我們,解決我們的恐懼和重擔,耶穌完成這一切,但卻並不是他最重要的角色,今天罪得赦免,沒有罪在身又如何?各位信徒,神已赦免我們的罪,我們又作了些什麼?耶穌基督說不但我們罪得赦免,不但生命得潔淨,神的創造帶我們到愛的關係,耶穌基督有聖靈降臨在他身上,擁有內在能力,然後神說這是是我的愛子,我喜悅你!不是普通的人,是愛子,肯定你的地位,我是父親,是關係的建立,所以完成水禮後,便去召聚群體,呼喚十二門徒建立愛的群體,一個博愛的群體,神新的子民.新的子民即你和我,我們組合成愛的群體,把神愛的信息傳遞,認罪赦免是水禮的一面,成為愛的群體是另一面.沒有人在水禮後以為已經完成了便離開教會,以為有天堂的護照便不用作什麼.耶穌基督的洗禮是肯定從黑暗到光明,並把光明帶到群體中.初期信徒受洗後互相幫忙,甚至變賣田產,以致可幫助人,還在家中一同敬拜擘餅,讚美神,得到神的喜悅,神把人數日日加添.

今天結束時,重覆對三組人的鼓勵:已經接受水禮的弟兄姊妹,過去日子中回想耶穌基督的洗禮,要像耶穌基督的愛的洗禮,讓每個人連結一起去聚集,像神創造時把水聚集,使這創造的能力在我們身上發揮.第二組人,要快快接受水禮,你們有感動便要報名接受水禮,不要耽誤,這不是儀式,你有悔改的心,神把愛你的生命給你.最後一組在邊緣的人,一直在信仰邊緣徘徊,這信仰是很值得的,會為你帶來創新私改變,你會驚訝受洗後,無論任何年齡,你的一生將會有新的改變和體會,你會經歷聖靈降臨,有聲音說:「我所愛的,回來吧!」

耶穌的洗禮是一個愛的洗禮!

天父上帝,在你面前我們感謝你,你的創造藉耶穌基督的洗禮顯現在我們當中,讓我們看到的不是自以為是的儀式,我們在你面前帶領我們到約旦河中,開我們的眼睛,像在約旦河旁邊,我們有敬仰神的心,我們藉耶穌基督完成救恩工作,一個極大的應許,我們來到屬靈約旦河邊,我們清楚來到你面前,不用理會旁邊的人,不用理會其他聲音,我們好清楚明白在神面前接受水禮時就像耶穌基督受洗時有聖靈進入我們當中,有神的聲音對我說:「你是我所愛,我很愛你,我把你帶到愛中,藉著你把愛帶到群體當中,將福音傳開.」我們眾信徒這樣禱告,奉主名求.阿們!



\newpage

\section{說一個感恩的故事}
\label{sec:190}
\textbf{林豪恩先生 - 說一個感恩的故事}
\newline
\newline
\hyperref[sec:189]{< < < PREV SERMON < < <}
~
\hyperref[sec:toc]{[返目錄]}
~
\hyperref[sec:191]{> > > NEXT SERMON > > >}


申命記 26:1-11
\newline
\begin{longtable}{cl}
\hline
\hline
章節 & 經文 (和合本修訂版)\\
\hline
26:1 & \begin{tabularx}{0.7\textwidth}{X} 「你進去得了耶和華－你神所賜你為業的地，並且居住在那裡的時候， \end{tabularx} \\ \\ \relax
26:2 & \begin{tabularx}{0.7\textwidth}{X} 就要從耶和華－你神所賜你的地上，將收成的各種初熟土產取一些來，盛在筐子裡，帶到那裡，耶和華－你神選擇作為他名居所的地方， \end{tabularx} \\ \\ \relax
26:3 & \begin{tabularx}{0.7\textwidth}{X} 到當時的祭司那裡，對他說：『我今日向耶和華－你的神宣認，我已來到耶和華向我們列祖起誓要賜給我們的地。』 \end{tabularx} \\ \\ \relax
26:4 & \begin{tabularx}{0.7\textwidth}{X} 祭司就從你手裡把筐子接過來，供在耶和華－你神的祭壇前。 \end{tabularx} \\ \\ \relax
26:5 & \begin{tabularx}{0.7\textwidth}{X} 你要在耶和華－你神面前告白說：『我的祖先原是一個流亡的亞蘭人，帶著稀少的人丁下到埃及寄居。在那裡，他卻成了又大又強、人數眾多的國。 \end{tabularx} \\ \\ \relax
26:6 & \begin{tabularx}{0.7\textwidth}{X} 埃及人惡待我們，迫害我們，將苦工加在我們身上。 \end{tabularx} \\ \\ \relax
26:7 & \begin{tabularx}{0.7\textwidth}{X} 於是我們哀求耶和華我們列祖的神。耶和華聽見我們的聲音，看見我們所受的困苦、勞役和欺壓， \end{tabularx} \\ \\ \relax
26:8 & \begin{tabularx}{0.7\textwidth}{X} 耶和華就用大能的手和伸出來的膀臂，以及大而可畏的事和神蹟奇事，領我們出了埃及， \end{tabularx} \\ \\ \relax
26:9 & \begin{tabularx}{0.7\textwidth}{X} 將我們領進這地方，把這流奶與蜜之地賜給我們。 \end{tabularx} \\ \\ \relax
26:10 & \begin{tabularx}{0.7\textwidth}{X} 耶和華啊，看哪，現在我把你所賜我地上初熟的土產供上。』隨後你要把筐子供在耶和華－你神面前，向耶和華－你的神下拜。 \end{tabularx} \\ \\ \relax
26:11 & \begin{tabularx}{0.7\textwidth}{X} 你和利未人，以及在你中間寄居的，要因耶和華－你神所賜你和你家的一切福分歡樂。 \end{tabularx} \\ \\
[1ex]
\hline
\hline
\end{longtable}
講題:說一個感恩的故事　To Tell A Thanksgiving Story

經文:申命記26章1-11節

講員:林豪恩先生

場合:香港中文大學崇基學院禮拜堂

日期:2013年2月17日

一.  引言

在去年中文大學畢業禮上,校長引用了一個耳熟能詳的半杯水故事.桌子上放著半杯水.有一個人說:「糟了!怎麼只剩下半杯水!」另一個人說:「真好!還有半杯水!」這兩個人描述相同的狀況,不過,卻用了兩種不同的態度.人生的經歷比半杯水的故事複雜得多,可以選取用作敘述的角度和主題也豐富得多.對於個人來說,你習慣以甚麼角度去敘述你的經驗呢?你的敘述突出甚麼主題呢?對於社會來說,我們的城市傾向以甚麼角度去敘述我們的故事呢?又突出甚麼主題呢?

今天的舊約經課,記載著摩西對離開埃及,將要進入迦南的以色列人的勸導,提及以色列人如何講述他們作為一個群體的歷史經驗.摩西建議他們用甚麼角度,選擇甚麼主題去講述他們的故事呢?選擇這些主題對他們又有甚麼意義呢?對我們又有甚麼啟發呢?

二.  經文

26 : 1「你進去得了耶和華－你神所賜你為業之地居住,

26 : 2就要從耶和華－你神賜你的地上將所收的各種初熟的土產取些來,盛在筐子裡,往耶和華－你神所選擇要立為他名的居所去,

26 : 3見當時作祭司的,對他說:『我今日向耶和華－你神明認,我已來到耶和華向我們列祖起誓應許賜給我們的地.』

26 : 4祭司就從你手裡取過筐子來,放在耶和華－你神的壇前.

26 : 5你要在耶和華－你神面前說:『我祖原是一個將亡的亞蘭人,下到埃及寄居.他人口稀少,在那裡卻成了又大又強,人數很多的國民.

26 : 6埃及人惡待我們,苦害我們,將苦工加在我們身上.

26 : 7於是我們哀求耶和華－我們列祖的神,耶和華聽見我們的聲音,看見我們所受的困苦,勞碌,欺壓,

26 : 8他就用大能的手和伸出來的膀臂,並大可畏的事與神蹟奇事,領我們出了埃及,

26 : 9將我們領進這地方,把這流奶與蜜之地賜給我們.

26 : 10耶和華啊,現在我把你所賜給我地上初熟的土產奉了來.』隨後你要把筐子放在耶和華－你神面前,向耶和華－你的神下拜.

26 : 11你和利未人,並在你們中間寄居的,要因耶和華－你神所賜你和你家的一切福分歡樂.

三.  釋經與應用

1．不全是困苦:困苦與祝福混雜

對以色列人說,埃及代表著壓迫,埃及經驗,代表著困苦,因此,出埃及,代表著釋放和盼望.在剛才誦經的經課中,他們是如此敘述埃及經驗的:

26 : 6埃及人惡待我們,苦害我們,將苦工加在我們身上.

26 : 7於是我們哀求耶和華－我們列祖的神,耶和華聽見我們的聲音,看見我們所受的困苦,勞碌,欺壓

以色列人以「惡待」,「苦害」,「將苦工加在我們身上」去描述埃及人對他們的態度,以「困苦」,「勞碌」,「欺壓」去表達自己的狀態.他們有足夠的歷史證據和親身經驗作這樣的描述.然而,這樣的描述卻不足以涵蓋以色列人在埃及的全部經驗.

有時候,我們都會感到自己被「惡待」,被「苦害」,感覺到別人「將苦工加在我們身上」,我們也會用「困苦」,「勞碌」,「欺壓」去描述自己的人生,然而,如果我們看得全面一些,這狀況未必是經驗的全部.二月初有一位歌手來在這個禮拜堂舉行的佈道會分享見證.她訴說在事業漸入佳境的時候,突然被某電視台封殺,外界更對她一沉百踩,散播很多對她很不利的消息.她感到自己受盡屈辱,飽嚐失敗的滋味,令她驚惶,害怕,擔心自己的歌唱事業能否繼續下去.對她來說,這些經驗和感覺到是真實的,不過,卻不足以描述她在樂壇的全部經驗.為甚麼不足以描述她在樂壇的全部經驗呢?因為,樂壇曾經讓她大放異彩,曾經把她捧上人生的高峰,曾經頒過不少獎項給她,更讓她有機會在紅磡體育館舉行個人演唱會.

摩西時代的以色列人在埃及的確經被惡待,然而,摩西建議他們避免以這階段的狀況去代表以色列人在埃及的全部經驗.因為,埃及曾經是以色列人的避難所,是他們休養生息的地方.

26 : 5．．．我祖原是一個將亡的亞蘭人,下到埃及寄居.他人口稀少,在那裡卻成了又大又強,人數很多的國民.

透過如此的敘述,以色列人承認,埃及曾經是以色列人的壓迫之地,也曾經是以色列人的發福之地.以色列人用「將亡」來描述他們來到埃及的狀況,在埃及,他們不但避過飢荒滅族的危機,更開枝散葉,繁衍後代.他們用「又大又強」,「人數很多」去形容當時的盛況.

在人生的某些時候,我們會容易一面倒地描述我們的經驗.例如有人把坐監和工作來個比較:「坐監大部分時間在一個 10 x 10 尺的空間,返工大部分時間在一個 8 x 8 尺的空間；坐監表現良好的話可以獲得假期,返工表現良好的話只會得到更多工作；坐監時守衛每天為你鎖好門窗,返工時你必須攜帶通行證並自己開啟門鎖；坐監你可以看電視和玩遊戲,返工時看電視或者玩遊戲的話將被解雇；坐監時你大部分時間在鐵欄後遙望外面,返工時你大部分時間都在窗花內想著外面；坐監時你必須與殘酷成聖的獄卒打交道,在辦公室他們的名字叫上司.」如果以比較客觀持平的態度去看,可能我們也會承認,工作其實正在滿足我們的某些需要.

有些時候我們會落在困苦壓迫的狀況之下,有些場境或人物會令我們感到被惡待和苦待,摩西提醒以色列人,這些狀況不是經驗的全部,因此避免以偏蓋全地讓自己跌進無奈和絕望的陷阱.如何能夠突破這種困惑呢?摩西建議以色列人以「寄居」的態度去看待埃及經驗,無論埃及如何令以色列人繁衍昌大,無論埃及如何令以色列人困苦壓迫,埃及不是永恆,埃及不是完美,埃及不是終極的依靠；無論樂壇如何把歌手捧上天空,又如何對同一位歌手一沉百踩,樂壇始終不是最終的依靠,也不是完美的地方；無論工作的地方和人事如何讓我們大展拳腳,或者令我們感到比坐監還辛苦,工作不是永恆的,上司和同事也不是完美的；就算是我們的家人,都有其不確定性和不完美性.正因如此,這些事物都可以帶給我們困苦的經驗,然而,卻又不盡是困苦,在某些時候,它彷彿又是祝福,正在滿足我們某些需要.

2．不全是歡愉:釋放路上挑戰滿途

如果埃及經驗的主調是困苦壓迫,那麼出埃及的經驗應該是釋放和歡愉.讓我們繼續看以色列人如何敘述他們的出埃及經驗.

26 : 7於是我們哀求耶和華－我們列祖的神,耶和華聽見我們的聲音,看見我們所受的困苦,勞碌,欺壓,

26 : 8他就用大能的手和伸出來的膀臂,並大可畏的事與神蹟奇事,領我們出了埃及,

26 : 9將我們領進這地方,把這流奶與蜜之地賜給我們.

以色列人如此概括地敘述出埃及的經驗,仍然沒有偏離事實,他們的確經驗過耶和華上帝施展神蹟奇事,在法老王和埃及軍隊面前拯救他們.不過,這也未必是他們出埃及的全部經驗.如果我們以為這就是以色列人出埃及的全部經驗,就等於以為童話故事結束的那一句:「從此以後,公主和王子就快快樂樂地住在一起了.」代表這對夫妻以後都沒有吵架,沒有衝突.無論出埃及故事如何強調拯救和釋放,這個大故事中仍然包含著許多其他與主旋律未必協調的小故事.這些小故事,展示著在釋放的路上仍然挑戰滿途.

第一個小故事記載在出埃及記14章,以色列人剛出發離開埃及,就懼怕得絕望.

14 : 10法老逼近的時候, 以色列人舉目,看哪, 埃及人追來了,就非常懼怕, 以色列人向耶和華哀求.

14 : 11他們對 摩西說:「難道 埃及沒有墳地,你要把我們帶來死在曠野嗎?你為甚麼這樣待我們,將我們從 埃及領出來呢?

14 : 12我們在 埃及豈沒有對你說過,不要攪擾我們,讓我們服事 埃及人嗎?因為服事 埃及人總比死在曠野好.」

第二個小故事記載在出埃及記16章,以色列人經驗過紅海神蹟之後,漸漸遠離埃及的時候,竟然緬懷埃及.

16 : 1以色列全會眾從 以琳起程,在出 埃及之後第二個月十五日到了 以琳和 西奈中間, 汛的曠野.

16 : 2以色列全會眾在曠野向 摩西和 亞倫發怨言.

16 : 3以色列人對他們說:「我們寧願在 埃及地死在耶和華手中!那時我們坐在肉鍋旁,吃餅得飽.你們卻將我們領出來,到這曠野,要叫這全會眾都餓死啊!」

以色列人不是以困苦和壓迫為主題去描述埃及經驗的嗎?為甚麼在出埃及的過程中卻不斷地出現緬懷埃及的片段呢?出埃及途中的以色列人缺水的時候,這種埃及情懷又出現了.(出埃及記)

17 : 1以色列全會眾遵照耶和華的吩咐,從 汛的曠野一段一段地往前行.他們在 利非訂安營,但百姓沒有水喝.

17 : 2百姓就與 摩西爭鬧,說:「給我們水喝吧!」 摩西對他們說:「你們為甚麼與我爭鬧呢?你們為甚麼試探耶和華呢?」

17 : 3百姓在那裡口渴要喝水,就向 摩西發怨言,說:「你為甚麼把我們從 埃及領出來,使我們和我們的兒女,以及牲畜都渴死呢?」

在第一個小故事中,以色列人對埃及的懷緬的想像是「我們坐在肉鍋旁,吃餅得飽」,到下面這個小故事,以色列人對埃及的懷緬後退到「我們寧願死在 埃及地」.這個小故事記載在民數記13章:

13 : 32於試探子中有人向 以色列人散佈有關所窺探之地的謠言,說:「我們所走過,所窺探之地是吞沒居民之地,並且我們在那裡所看見的百姓都身材高大.

13 : 33我們在那裡看見巨人,就是巨人中的 亞衲族人.我們在自己眼中像蚱蜢一樣,而在他們眼中,我們也確是這樣.」

14 : 1全會眾大聲喧嚷,那夜百姓哭號.

14 : 2以色列眾人向 摩西和 亞倫發怨言,全會眾對他們說:「我們寧願死在 埃及地,寧願死在這曠野!

14 : 3耶和華為甚麼要把我們領到那地,讓我們倒在刀下呢?我們的妻子和孩子必成為擄物.我們回 埃及去豈不更好嗎?」

埃及經驗,是否真的「我們坐在肉鍋旁,吃餅得飽」呢?埃及經驗,是否美好得「我們寧願死在 埃及地」呢?埃及經驗是甚麼,已經不再重要,重要的是以色列人當下的困難.當下,以色列人被埃及軍隊追捕；當下,以色列人缺水缺糧；當下,以色列人面對強敵當前,前路未卜.所有這些當下的處境,都是真實的,都是切身的.或者,並不是埃及經驗的美好令以色列人緬懷過去,而是他們當下的困境令他們美化埃及經驗,對埃及產生緬懷之情.

人們甚麼時候最緬懷初戀情人呢?可能是他／她對現在的配偶感到不滿的時候；我們這個城市的人們甚麼時候最緬懷過去呢?可能是走遍港九新界也買不到一鑵奶粉給自己的BB的時候,可能是這個城市的廉潔程度和競爭能力越來越被其他城市迎頭趕上的時候,可能是面對將來感到徬徨不安的時候.走在婚姻的路上,要做的不是回去尋找那位緬懷的初戀情人,而是想辦法增進改善目前的關係；對我們身處的城市來說,回應緬懷,豈不是也應該想辦法解決當下的困境嗎?

3．也無風雨也無晴,毋忘感恩

埃及經驗過漸漸遠去,有一天,以色列人也會走完曠野之路,他們如何看待「現在」呢?「現在」,是過去的苦難的陰霾嗎?「現在」,是個人勝利的成果嗎?摩西建議以色列人,不要忘記過去.記住過去,不是要令自己繼續生活在苦澀當中,而是記住在困苦的經歷中,在挑戰的路途上,現在能夠在這裡站立,有上帝的恩師及保守.若果以色列人能夠視「現在」為上帝的賜予,他們就會如此回應:

26 : 10耶和華啊,現在我把你所賜給我地上初熟的土產奉了來.』隨後你要把筐子放在耶和華－你神面前,向耶和華－你的神下拜.

26 : 11你和利未人,並在你們中間寄居的,要因耶和華－你神所賜你和你家的一切福分歡樂.

摩西知道有一天以色列人被引誘,不再以上帝的恩典為主題去講述自己的故事,卻以個人的勝利為主題去取代.因此,摩西在申命記第8章提醒他們:

8 : 11「你要謹慎,免得忘記耶和華─你的上帝,不守他的誡命,典章,律例,就是我今日吩咐你的.

8 : 12免得你吃得飽足,建造上好的房屋,住在其中,

8 : 13你的牛羊增多,你的金銀增多,你擁有的一切全都增多,

8 : 14於是你的心高傲,忘記耶和華─你的上帝.他曾將你從 埃及地為奴之家領出來,

8 : 15曾引領你經過那大而可怕的曠野,有火蛇,蠍子,乾旱無水之地.他也曾為你使水從堅硬的磐石中流出來,

8 : 16又在曠野將你列祖所不認識的嗎哪賜給你吃,為要磨煉你,考驗你,終久使你享福.

8 : 17你心裡說:『這財富是我的力量,我手的能力得來的.』

對於以色列人來說,這兩個敘述「現在」的故事互相競爭,採用第一個故事和第二個故事有甚麼不同呢?如果採用第一個故事,第一,他們會視現在為上帝的恩典,引發感恩的回應行動；第二,他們會記住自己曾經在埃及寄居,因此必須要與現在在他們當中寄居的人分享他們的恩典.這個故事很重要,因為建構著他們的社會生活,就是擁有生產資源的人有責任去讓主動放棄生產資源的人──在他們的社會中是利未人,以及沒有生產資源的人──在他們的社會中是寄居的,這兩批人都能夠生存於社會之中,成為社會的一個部份.

如果他們採用第二個故事,他們的信念就會變成:「這財富是我的力量,我手的能力得來的」,於是,他們只會「吃得飽足,建造上好的房屋,住在其中．．．牛羊增多．．．金銀增多,擁有的一切全都增多,於是心高氣傲．．．」這將會是另外一個社會狀況,另外一種生存狀況.

四.  總結

閱讀過摩西指導以色列人講述他們的故事,對此時此地的我們有甚麼值得參考呢?個人來說,你採用甚麼主題去講述你的故事呢?我們的社會又採用甚麼主題去講述我們的故事呢?

在我們的城市中,有一個人自嘲「自己的思想很不適合現代社會,可以說是一個怪人.」他自謙「只是小工業家,不能和地產大富豪相比」.不過,他的信念是「取之社會,用於社會」,因此,他自己過著簡樸的生活,把財富與社會分享.某年,他毅然賣掉居住了三十七年,位於九龍塘的豪宅,得款五千六百多萬元全數捐資助學,然後租住舊居旁邊的一個住宅單位.他的計算是這樣的:「雖然也有點不捨得,但打個譬喻,當房子的市價是六千萬元,每月的利息大概就有四十萬元,一想到自己睡一晚就不見了一萬三千三百三十三元,就覺得這樣太無謂,便賣了出去.」

現在,講述我們的城市的故事的版本眾多,哪些主題能夠引導我們個人和整個社會向健康和仁愛前進呢?願我們的故事裡面也有有拯救,有恩典,有感恩,有分享這些主題,我們當中的「利未人,並在我們當中寄居的,要因神所賜我們的一切福分歡樂」.



\newpage

\section{說一個感恩的故事}
\label{sec:191}
\textbf{林豪恩先生 - 說一個感恩的故事}
\newline
\newline
\hyperref[sec:190]{< < < PREV SERMON < < <}
~
\hyperref[sec:toc]{[返目錄]}
~
\hyperref[sec:192]{> > > NEXT SERMON > > >}


申命記 26:1-11
\newline
\begin{longtable}{cl}
\hline
\hline
章節 & 經文 (和合本修訂版)\\
\hline
26:1 & \begin{tabularx}{0.7\textwidth}{X} 「你進去得了耶和華－你神所賜你為業的地，並且居住在那裡的時候， \end{tabularx} \\ \\ \relax
26:2 & \begin{tabularx}{0.7\textwidth}{X} 就要從耶和華－你神所賜你的地上，將收成的各種初熟土產取一些來，盛在筐子裡，帶到那裡，耶和華－你神選擇作為他名居所的地方， \end{tabularx} \\ \\ \relax
26:3 & \begin{tabularx}{0.7\textwidth}{X} 到當時的祭司那裡，對他說：『我今日向耶和華－你的神宣認，我已來到耶和華向我們列祖起誓要賜給我們的地。』 \end{tabularx} \\ \\ \relax
26:4 & \begin{tabularx}{0.7\textwidth}{X} 祭司就從你手裡把筐子接過來，供在耶和華－你神的祭壇前。 \end{tabularx} \\ \\ \relax
26:5 & \begin{tabularx}{0.7\textwidth}{X} 你要在耶和華－你神面前告白說：『我的祖先原是一個流亡的亞蘭人，帶著稀少的人丁下到埃及寄居。在那裡，他卻成了又大又強、人數眾多的國。 \end{tabularx} \\ \\ \relax
26:6 & \begin{tabularx}{0.7\textwidth}{X} 埃及人惡待我們，迫害我們，將苦工加在我們身上。 \end{tabularx} \\ \\ \relax
26:7 & \begin{tabularx}{0.7\textwidth}{X} 於是我們哀求耶和華我們列祖的神。耶和華聽見我們的聲音，看見我們所受的困苦、勞役和欺壓， \end{tabularx} \\ \\ \relax
26:8 & \begin{tabularx}{0.7\textwidth}{X} 耶和華就用大能的手和伸出來的膀臂，以及大而可畏的事和神蹟奇事，領我們出了埃及， \end{tabularx} \\ \\ \relax
26:9 & \begin{tabularx}{0.7\textwidth}{X} 將我們領進這地方，把這流奶與蜜之地賜給我們。 \end{tabularx} \\ \\ \relax
26:10 & \begin{tabularx}{0.7\textwidth}{X} 耶和華啊，看哪，現在我把你所賜我地上初熟的土產供上。』隨後你要把筐子供在耶和華－你神面前，向耶和華－你的神下拜。 \end{tabularx} \\ \\ \relax
26:11 & \begin{tabularx}{0.7\textwidth}{X} 你和利未人，以及在你中間寄居的，要因耶和華－你神所賜你和你家的一切福分歡樂。 \end{tabularx} \\ \\
[1ex]
\hline
\hline
\end{longtable}
日期:2013年2月24日

題目:說一個感恩的故事

經文:申命記26章1-11節

證道:林豪恩先生

一.  引言

在去年中文大學畢業禮上,校長引用了一個耳熟能詳的半杯水故事.桌子上放著半杯水.有一個人說:「糟了!怎麼只剩下半杯水!」另一個人說:「真好!還有半杯水!」這兩個人描述相同的狀況,不過,卻用了兩種不同的態度.人生的經歷比半杯水的故事複雜得多,可以選取用作敘述的角度和主題也豐富得多.對於個人來說,你習慣以甚麼角度去敘述你的經驗呢?你的敘述突出甚麼主題呢?對於社會來說,我們的城市傾向以甚麼角度去敘述我們的故事呢?又突出甚麼主題呢?

今天的舊約經課,記載著摩西對離開埃及,將要進入迦南的以色列人的勸導,提及以色列人如何講述他們作為一個群體的歷史經驗.摩西建議他們用甚麼角度,選擇甚麼主題去講述他們的故事呢?選擇這些主題對他們又有甚麼意義呢?對我們又有甚麼啟發呢?

二.  經文

1     你進去得了耶和華你　神所賜你為業之地居住,     2     就要從耶和華你　神賜你的地上,將所收的各種初熟的土產,取些來,盛在筐子裡,往耶和華你　神所選擇要立為他名的居所去．     3     見當時作祭司的,對他說,我今日向耶和華你　神,明認我已來到耶和華向我們列祖起誓應許賜給我們的地.     4     祭司就從你手裡取過筐子來,放在耶和華你　神的壇前.     5     你要在耶和華你　神面前說,我祖原是一個將亡的亞蘭人,下到埃及寄居．他人口稀少,在那裡卻成了又大,又強,人數很多的國民.     6     埃及人惡待我們,苦害我們,將苦工加在我們身上．     7     於是我們哀求耶和華我們列祖的　神,耶和華聽見我們的聲音,看見我們所受的困苦,勞碌,欺壓,     8     他就用大能的手和伸出來的膀臂,並大可畏的事,與神蹟,奇事,領我們出了埃及．     9     將我們領進這地方,把這流奶與蜜之地賜給我們.     10     耶和華阿,現在我把你所賜給我地上初熟的土產奉了來．隨後你要把筐子放在耶和華你　神面前,向耶和華你的　神下拜.     11     你和利未人,並在你們．中間寄居的,要因耶和華你　神所賜你和你家的一切福分歡樂.

三.  釋經與應用

1．不全是困苦:困苦與祝福混雜

對以色列人說,埃及代表著壓迫,埃及經驗,代表著困苦,因此,出埃及,代表著釋放和盼望.在剛才誦經的經課中,他們是如此敘述埃及經驗的:

6     埃及人惡待我們,苦害我們,將苦工加在我們身上．     7     於是我們哀求耶和華我們列祖的　神,耶和華聽見我們的聲音,看見我們所受的困苦,勞碌,欺壓

以色列人以「惡待」,「苦害」,「將苦工加在我們身上」去描述埃及人對他們的態度,以「困苦」,「勞碌」,「欺壓」去表達自己的狀態.他們有足夠的歷史證據和親身經驗作這樣的描述.然而,這樣的描述卻不足以涵蓋以色列人在埃及的全部經驗.

有時候,我們都會感到自己被「惡待」,被「苦害」,感覺到別人「將苦工加在我們身上」,我們也會用「困苦」,「勞碌」,「欺壓」去描述自己的人生,然而,如果我們看得全面一些,這狀況未必是經驗的全部.二月初有一位歌手來在這個禮拜堂舉行的佈道會分享見證.她訴說在事業漸入佳境的時候,突然被某電視台封殺,外界更對她一沉百踩,散播很多對她很不利的消息.她感到自己受盡屈辱,飽嚐失敗的滋味,令她驚惶,害怕,擔心自己的歌唱事業能否繼續下去.對她來說,這些經驗和感覺到是真實的,不過,卻不足以描述她在樂壇的全部經驗.為甚麼不足以描述她在樂壇的全部經驗呢?因為,樂壇曾經讓她大放異彩,曾經把她捧上人生的高峰,曾經頒過不少獎項給她,更讓她有機會在紅磡體育館舉行個人演唱會.

摩西時代的以色列人在埃及的確經被惡待,然而,摩西建議他們避免以這階段的狀況去代表以色列人在埃及的全部經驗.因為,埃及曾經是以色列人的避難所,是他們休養生息的地方.

5

．．．我祖原是一個將亡的亞蘭人,下到埃及寄居．他人口稀少,在那裡卻成了又大,又強,人數很多的國民.

透過如此的敘述,以色列人承認,埃及曾經是以色列人的壓迫之地,也曾經是以色列人的發福之地.以色列人用「將亡」來描述他們來到埃及的狀況,在埃及,他們不但避過飢荒滅族的危機,更開枝散葉,繁衍後代.他們用「又大又強」,「人數很多」去形容當時的盛況.

在人生的某些時候,我們會容易一面倒地描述我們的經驗.例如有人把坐監和工作來個比較:「坐監大部分時間在一個 10 x 10 尺的空間,返工大部分時間在一個 8 x 8 尺的空間；坐監表現良好的話可以獲得假期,返工表現良好的話只會得到更多工作；坐監時守衛每天為你鎖好門窗,返工時你必須攜帶通行證並自己開啟門鎖；坐監你可以看電視和玩遊戲,返工時看電視或者玩遊戲的話將被解雇；坐監時你大部分時間在鐵欄後遙望外面,返工時你大部分時間都在窗花內想著外面；坐監時你必須與殘酷成聖的獄卒打交道,在辦公室他們的名字叫上司.」如果以比較客觀持平的態度去看,可能我們也會承認,工作其實正在滿足我們的某些需要.

有些時候我們會落在困苦壓迫的狀況之下,有些場境或人物會令我們感到被惡待和苦待,摩西提醒以色列人,這些狀況不是經驗的全部,因此避免以偏蓋全地讓自己跌進無奈和絕望的陷阱.如何能夠突破這種困惑呢?摩西建議以色列人以「寄居」的態度去看待埃及經驗,無論埃及如何令以色列人繁衍昌大,無論埃及如何令以色列人困苦壓迫,埃及不是永恆,埃及不是完美,埃及不是終極的依靠；無論樂壇如何把歌手捧上天空,又如何對同一位歌手一沉百踩,樂壇始終不是最終的依靠,也不是完美的地方；無論工作的地方和人事如何讓我們大展拳腳,或者令我們感到比坐監還辛苦,工作不是永恆的,上司和同事也不是完美的；就算是我們的家人,都有其不確定性和不完美性.正因如此,這些事物都可以帶給我們困苦的經驗,然而,卻又不盡是困苦,在某些時候,它彷彿又是祝福,正在滿足我們某些需要.

2．不全是歡愉:釋放路上挑戰滿途

如果埃及經驗的主調是困苦壓迫,那麼出埃及的經驗應該是釋放和歡愉.讓我們繼續看以色列人如何敘述他們的出埃及經驗.

7     於是我們哀求耶和華我們列祖的　神,耶和華聽見我們的聲音,看見我們所受的困苦,勞碌,欺壓,     8     他就用大能的手和伸出來的膀臂,並大可畏的事,與神蹟,奇事,領我們出了埃及．     9     將我們領進這地方,把這流奶與蜜之地賜給我們.

以色列人如此概括地敘述出埃及的經驗,仍然沒有偏離事實,他們的確經驗過耶和華上帝施展神蹟奇事,在法老王和埃及軍隊面前拯救他們.不過,這也未必是他們出埃及的全部經驗.如果我們以為這就是以色列人出埃及的全部經驗,就等於以為童話故事結束的那一句:「從此以後,公主和王子就快快樂樂地住在一起了.」代表這對夫妻以後都沒有吵架,沒有衝突.無論出埃及故事如何強調拯救和釋放,這個大故事中仍然包含著許多其他與主旋律未必協調的小故事.這些小故事,展示著在釋放的路上仍然挑戰滿途.

第一個小故事記載在出埃及記14章,以色列人剛出發離開埃及,就懼怕得絕望.

10     法老臨近的時候,以色列人舉目看見埃及人趕來,就甚懼怕,向耶和華哀求.     11     他們對摩西說,難道在埃及沒有墳地,你把我們帶來死在曠野麼,你為甚麼這樣待我們,將我們從埃及領出來呢.     12     我們在埃及豈沒有對你說過,不要攪擾我們,容我們服事埃及人麼,因為服事埃及人比死在曠野還好.

第二個小故事記載在出埃及記16章,以色列人經驗過紅海神蹟之後,漸漸遠離埃及的時候,竟然緬懷埃及.

1     以色列全會眾從以琳起行,在出埃及後第二個月十五日,到了以琳和西乃中間,汛的曠野.     2     以色列全會眾在曠野向摩西,亞倫發怨言,     3     說,巴不得我們早死在埃及地耶和華的手下,那時我們坐在肉鍋旁邊,喫得飽足．你們將我們領出來,到這曠野,是要叫這全會眾都餓死阿.

以色列人不是以困苦和壓迫為主題去描述埃及經驗的嗎?為甚麼在出埃及的過程中卻不斷地出現緬懷埃及的片段呢?出埃及途中的以色列人缺水的時候,這種埃及情懷又出現了.(出埃及記)

1     以色列全會眾都遵耶和華的吩咐,按著站口從汛的曠野往前行,在利非訂安營．百姓沒有水喝,     2     所以與摩西爭鬧,說,給我們水喝罷.摩西對他們說,你們為甚麼與我爭鬧,為甚麼試探耶和華呢.     3     百姓在那裡甚渴,要喝水,就向摩西發怨言,說,你為甚麼將我們從埃及領出來,使我們和我們的兒女並牲畜,都渴死呢.

在第一個小故事中,以色列人對埃及的懷緬的想像是「我們坐在肉鍋旁,吃餅得飽」,到下面這個小故事,以色列人對埃及的懷緬後退到「我們寧願死在 埃及地」.這個小故事記載在民數記13章:

32     探子中有人論到所窺探之地,向以色列人報惡信,說,我們所窺探經過之地,是吞喫居民之地,我們在那裡所看見的人民,都身量高大.     33     我們在那裡看見亞衲族人,就是偉人,他們是偉人的後裔,據我們看自己就如蚱蜢一樣,據他們看我們也是如此.

+2">14 1     當下全會眾大聲喧嚷,那夜百姓都哭號.     2     以色列眾人向摩西亞倫發怨言,全會眾對他們說,巴不得我們早死在埃及地,或是死在這曠野.     3     耶和華為甚麼把我們領到那地,使我們倒在刀下呢．我們的妻子和孩子,必被擄掠,我們回埃及去豈不好麼.

埃及經驗,是否真的「我們坐在肉鍋旁,吃餅得飽」呢?埃及經驗,是否美好得「我們寧願死在 埃及地」呢?埃及經驗是甚麼,已經不再重要,重要的是以色列人當下的困難.當下,以色列人被埃及軍隊追捕；當下,以色列人缺水缺糧；當下,以色列人面對強敵當前,前路未卜.所有這些當下的處境,都是真實的,都是切身的.或者,並不是埃及經驗的美好令以色列人緬懷過去,而是他們當下的困境令他們美化埃及經驗,對埃及產生緬懷之情.

人們甚麼時候最緬懷初戀情人呢?可能是他／她對現在的配偶感到不滿的時候；我們這個城市的人們甚麼時候最緬懷過去呢?可能是走遍港九新界也買不到一鑵奶粉給自己的BB的時候,可能是這個城市的廉潔程度和競爭能力越來越被其他城市迎頭趕上的時候,可能是面對將來感到徬徨不安的時候.走在婚姻的路上,要做的不是回去尋找那位緬懷的初戀情人,而是想辦法增進改善目前的關係；對我們身處的城市來說,回應緬懷,豈不是也應該想辦法解決當下的困境嗎?

3．也無風雨也無晴,毋忘感恩

埃及經驗過漸漸遠去,有一天,以色列人也會走完曠野之路,他們如何看待「現在」呢?「現在」,是過去的苦難的陰霾嗎?「現在」,是個人勝利的成果嗎?摩西建議以色列人,不要忘記過去.記住過去,不是要令自己繼續生活在苦澀當中,而是記住在困苦的經歷中,在挑戰的路途上,現在能夠在這裡站立,有上帝的恩師及保守.若果以色列人能夠視「現在」為上帝的賜予,他們就會如此回應:

10     耶和華阿,現在我把你所賜給我地上初熟的土產奉了來．隨後你要把筐子放在耶和華你　神面前,向耶和華你的　神下拜.     11     你和利未人,並在你們．中間寄居的,要因耶和華你　神所賜你和你家的一切福分歡樂.

摩西知道有一天以色列人被引誘,不再以上帝的恩典為主題去講述自己的故事,卻以個人的勝利為主題去取代.因此,摩西在申命記第8章提醒他們:

11     你要謹慎,免得忘記耶和華你的　神,不守他的誡命,典章,律例,就是我今日所吩咐你的.     12     恐怕你喫得飽足,建造美好的房屋居住．     13     你的牛羊加多,你的金銀增添,並你所有的全都加增,     14     你就心高氣傲,忘記耶和華你的　神,就是將你從埃及地為奴之家領出來的,     15     引你經過那大而可怕的曠野,那裡有火蛇,蠍子,乾旱無水之地．他曾為你使水從堅硬的磐石中流出來．     16     又在曠野,將你列祖所不認識的嗎哪賜給你喫,是要苦煉你,試驗你,叫你終久享福.     17     恐怕你心裡說,這貨財是我力量我能力得來的.

對於以色列人來說,這兩個敘述「現在」的故事互相競爭,採用第一個故事和第二個故事有甚麼不同呢?如果採用第一個故事,第一,他們會視現在為上帝的恩典,引發感恩的回應行動；第二,他們會記住自己曾經在埃及寄居,因此必須要與現在在他們當中寄居的人分享他們的恩典.這個故事很重要,因為建構著他們的社會生活,就是擁有生產資源的人有責任去讓主動放棄生產資源的人──在他們的社會中是利未人,以及沒有生產資源的人──在他們的社會中是寄居的,這兩批人都能夠生存於社會之中,成為社會的一個部份.

如果他們採用第二個故事,他們的信念就會變成:「這財富是我的力量,我手的能力得來的」,於是,他們只會「吃得飽足,建造上好的房屋,住在其中．．．牛羊增多．．．金銀增多,擁有的一切全都增多,於是心高氣傲．．．」這將會是另外一個社會狀況,另外一種生存狀況.

四.  總結

閱讀過摩西指導以色列人講述他們的故事,對此時此地的我們有甚麼值得參考呢?個人來說,你採用甚麼主題去講述你的故事呢?我們的社會又採用甚麼主題去講述我們的故事呢?

在我們的城市中,有一個人自嘲「自己的思想很不適合現代社會,可以說是一個怪人.」他自謙「只是小工業家,不能和地產大富豪相比」.不過,他的信念是「取之社會,用於社會」,因此,他自己過著簡樸的生活,把財富與社會分享.某年,他毅然賣掉居住了三十七年,位於九龍塘的豪宅,得款五千六百多萬元全數捐資助學,然後租住舊居旁邊的一個住宅單位.他的計算是這樣的:「雖然也有點不捨得,但打個譬喻,當房子的市價是六千萬元,每月的利息大概就有四十萬元,一想到自己睡一晚就不見了一萬三千三百三十三元,就覺得這樣太無謂,便賣了出去.」

現在,講述我們的城市的故事的版本眾多,哪些主題能夠引導我們個人和整個社會向健康和仁愛前進呢?願我們的故事裡面也有有拯救,有恩典,有感恩,有分享這些主題,我們當中的「利未人,並在我們當中寄居的,要因神所賜我們的一切福分歡樂」.



\newpage

\section{在曠野的一天}
\label{sec:192}
\textbf{伍渭文牧師 - 在曠野的一天}
\newline
\newline
\hyperref[sec:191]{< < < PREV SERMON < < <}
~
\hyperref[sec:toc]{[返目錄]}
~
\hyperref[sec:193]{> > > NEXT SERMON > > >}


講題:在曠野的一天　One Day at a Time in the Wilderness

經文:詩篇63

講員:伍渭文牧師

場合:香港中文大學崇基學院禮拜堂

日期:2013年3月3日

有一個人求上帝祝福他,給他蒙福的七天.上帝說:好,你可以選擇七天,我賜福給你!那人說:星期天,星期一,星期二……這七天裡你都賜福給我.上帝說:你很取巧,給你四天就算了.那人說:四天就四天吧；春天,夏天,秋天,冬天.上帝說:那不是一樣的嗎?減你一天,祇給你三天.那人說:好,三天就三天；我要昨天,今天,明天.上帝見這人很懂說話,心忖,給你一天,看你又如何.那人說:一天都行,一天裡面有白天和黑天,不論是白天和黑天,上帝你都要祝福我.

今天講道的題目「在曠野的一天」(One Day at a Time in the Wilderness) ,來自詩篇63篇標題:「大衛在曠野的時候,作了這詩.」大衛有兩次在曠野逃避敵人,第一次逃避掃羅王(Saul) :大衛殺死敵軍巨人哥利亞(Goliah) ,百姓擁戴頌讚他「掃羅殺死千千,大德殺死萬萬」,功高蓋主,掃羅王嫉妒,要追殺大衛(撒上23: 14,15) .另一次是逃避叛變的兒子押沙龍(撒下15, 16) .根據v. 16,「王必因神歡喜」,這次在曠野應該是逃避兒子押沙龍.

一位在曠野的在逃者,祇有一天的生活-one day at a time,七天的生活節奏,程序沒有了.以前可能星期二與核心內閣開行政會議,星期三,四批示奏摺,星期五上午會見祭司,傍晚安息日到會幕獻祭祈禱.曠野也沒有四季,在乾旱無水之地,植物稀少,沒有春天的發華,秋天的落葉.一位在曠野的在逃者,也失去昨天；昨天的權力,地位,財富,已成為過去,沒有昨天了.明天如何?會有明天嗎?明天是凶是吉還未知曉；缺水,野獸和追兵,都威脅著在逃者的生命,活到今天,已是恩典,明天如何還不知道；一覺醒來,要面對的是漫長的明天,不確實的明天.所以,在曠野的大衛,祇有今天-白晝和黑夜的一天,他祇能夠One day at a time.

在曠野,白晝和黑夜顯得非常真實；太陽下山後,四周愈來愈黑,舉目繁星閃閃,是那麼接近,但又是遙遠.這些星宿,億兆年已在那裡；望見的光,也可能來自億萬年前已消逝的星體.人世間的恩怨情仇,人世間的權力鬥爭,就像劃過黑夜的流星,轉眼即逝,何必爭過你死我活呢?

1. 在乾旱疲乏無水之地我渴想你vv . 1, 2

我的靈魂,我的肉身

詩篇63篇不是大衛逃離的歷史寫實記錄,是敬拜的公共禱文,詩篇的作者或編輯者,論及敬拜者要切慕上主,像大衛在曠野逃難時在乾旱無水之地,又饑又渴,感到疲乏,軟弱無力,渴想食物和水般渴慕上主.詩篇63: 1, 2節天主教思高譯本譯得很傳神:「天主(上主),你是我的天主(上主),我急切尋覓你；我的靈魂渴慕你,我的肉身切望你,我有如一塊乾涸無水田地.」

今天的的經課都提到食物:

賽55: 1, 2「你們一切乾渴的都當就近水來…..你們要留意聽我的話,就能吃那美物,得享肥甘,心裡喜樂.」真正滿足我們的是上主的話,上主的話是美物肥甘,帶來真正的喜樂和滿足,何必賣那不能使人飽足的食物呢?

林前13: 1-10提到百姓雖然「都從雲裡,海裡受洗歸了摩西」,他們離開埃及,經過紅海,在曠野蒙上主白天以雲柱,夜晚以火柱引路,吃了從天而降的嗎哪,但仍然念念不忘埃及的美食:「百姓坐下吃喝,起來玩耍.」(v. 7)

一般來說,我們尋覓上主,仰望他,是要舉心向上(lift up my soul) ,但這裡強調「我的肉身切望你」,「我還活著的時候要這樣稱頌你」(v. 4) ,這是一個在逃者,在曠野經歷生命朝不保夕時的深刻體會.肉體的私慾令我們忘記上主,但當我們明白肉體的限制時,亦是一個屬靈的轉機,使我們親近上主.

預苦期又稱為大齋期,我們不單用四十天預備心靈,迎接耶穌為我們受苦,被釘十字架.在這段時間,我們也效法耶穌,學習禁食,專注禱告,賙濟窮人(fast, prayer, alms giving) .操練禁食,聯想到饑餓貧窮的人的需要；很多教會,特別在歐洲,禁食同時舉行濟貧活動,藉此籌款,香港宣明會多年前就在大齋期舉行五餅二魚,饑饉三十(禁食三十小時)籌款活動.

禁食與禱告

但禁食基本的目的是用作專注禱告,不是為了禁食而禁食.禁食使我們身體疲乏,饑餓,感到自己的限制,明白到我們活著存留都在乎上主.在消費至上的社會,崇尚即時得到欲望的滿足,很容易做成一種罪性的自滿.罪的基本精神是脫離上主而存活,所以禁果又名分別善惡的果子,能夠分別善惡,就是能夠獨立自主.當南方猶太國面對北方以色列國和亞蘭王聯盟兵臨城下,以賽亞先知告訴猶太王亞哈斯:不要怕,上主會幫助你.「必有童女懷孕生子,給他起名叫以馬內利……在這孩子還不曉得棄惡擇善的時候,你所憎惡的那二王之地必得見棄.」(賽7: 14-16)

禁食可打破此種罪性的自滿.經過難受的初期,禁食者發覺頭腦愈來愈清醒,可以少睡一點,有一種輕盈,自由,喜樂的感覺,我們發覺根本不需要吃那麼多食物；我們愈來愈發現食得太多,缺少運動是引至各種疾病的根由；適量的節食可以使身體回復健康；放縱私欲,失去節制,就會傷害身體.

2. 我在聖所中曾如此瞻仰你v. 2, 3

「我在聖所中曾如此瞻仰你,為要見你的能力和你的榮耀.因你的慈愛比生命更好,我的嘴唇要頌讚你.」(詩63: 2, 3)

a. 因你的慈愛比生命更好v. 3

「我在聖所中曾經如此瞻仰你,為要見你的能力和榮耀.」(v.2) 能力和榮耀指聖殿的約櫃(ark) ,聖殿中最神聖的地方是至聖所,至聖所內面安放約櫃,裡面存放十誡兩塊法板.大祭司一年一次在贖罪日拿著祭牲的血,進入至聖所,把血灑在約櫃的頂部-施恩座(原文為遮蓋)-的四角,為人贖罪,施恩座有天使基路伯(serpherim) 守護.約櫃代表上主上主的臨在,代表他的拯救的能力和榮耀.當約櫃被非利士(Philistines)人擄走離開以色列時,聖經稱「榮耀離開以色列,因為神的約櫃被擄去了.」(撒上4: 21) 第七節「在你的翅膀下歡呼」是指守護約櫃的天使的翅膀.

以前,大衛王作為一國之尊,是最有榮耀和能力的人,但現在成為曠野的在逃者,失去了以前擁有的榮耀和能力,對上主的能力和榮耀有特別的體會.失戀的人,被信任的人所背叛的人,對上主不離不棄的愛特別感動.

詩人想起過去聖殿崇拜的情景,想起上主的能力和榮耀.此刻,天幕為蓋,野地為席,仰望頭頂繁星點點,上主何其親近；他切慕上主,像守夜的等候天明:「我在床上記念你,在夜更的時候思想你.」(vv., 5, 6) 此刻,心中就像飽足了的骨髓肥油,他不再懼怕死亡了.一個曠野的在逃者,沒有七天的生活程序,沒有昨天和明天,祇有今天,他最關注的是保全自己的生命,生命最重要.但他發現有一樣東西比生命更重要-與上主的親密關係比生命更重要,詩人說:「你的慈愛比生命更好」.(v., 3a)

b. 找到往死亡的路徑

剛過去的崇基週年教會研討會,主席介紹三位台座講員(panelists)其中一位李怡先生時,說他是「不默而生,寧嗚而死」；言論的自由,比生命更重要.當時我立刻想起一位在獄中的劉先生,他為了追求落實作為中國公民憲法應有的權利,發起零八憲章,也是不默而生,寧嗚而死.在中國大陸不少關心中國憲政發展的知識分子也簽署了這憲章,宣告自由比仕途的生命更好.一位我們熟悉的著名基督徒教授簽了這憲章後,大學校長跟他說:為何簽這份文件?他說:有人電郵給我,我看過後,其實內容很簡單,說一加一等如二,我很同意,所以簽了.之後,就不獲安排博士生跟他做研究.

生命比任何都重要,當人認為有一些東西比生命更重要時,他是有福的人,因為他找到了一樣東西,比生命更恆久,更有價值.不是「人一世,物一世」那麼短暫；沒有復活盼望的以弗所人,祇能說:「我們吃吃喝喝吧,因為明天要死了.」(林前15: 32) 孔子說,未知生,焉知死；但聯合國第二任秘書長韓馬紹(Dag Hammarskjold)說:若不認識我們會死,我們就不會好好去活.他說:「我們不需要尋求死亡,死亡會找上門的,但我們要找到往死亡的路,使死亡顯得有意義.」 (Do not seek death. Death will find you. But seek the road which makes death a fulfillment)

還有幾個月,在崇基工作轉瞬間就快十三年了,人生還有幾多個十三年呢?剛來時,胡秀英教授上斜路比我快,但看著她身體慢慢衰退,甚至因健康欠佳停止來教堂崇拜,最後安息主懷.我們真的要追求一些比生命要寶貴的東西,因為生命終有一天結束.有人晚上坐船看煙花的時候,船就沉了,生命突然終結；有人在白天在氣球升空的時候,生命突然終結.當我們面對死亡時,能否像詩人所說:上主啊,你的慈愛比生命更好.

「我們不需要尋找死亡,死亡自然會尋找我們,我們要的是尋找往死亡的路徑.」你找到往死亡的路徑嗎?詩人找到了,與上親密的關係,就是往死亡的路徑.聖經談到永生,原意就是與上主親密的關係,這關係連死亡都不能割裂,這就是永生.然而,與上主關係愈密切,也愈關心他者的需要.

3. 王必因上主歡喜 vv. 8-11

a.      生命顯揚上主的慈愛

「我在床上記念你,在夜更的時候思想你.」(vv., 5, 6) 大衛在床上記念甚麼呢?第八節:「我心緊緊跟隨你；你的右手扶持我.」我想大衛一定會想起過去的日子上主如何恩待他.

他本是牧羊小子,身形並不高大,但上主不是看外貌,看人的內心,上主透過撒母耳(Samuel) 到他父親家中找他,膏立他承繼掃羅王(Saul).(撒上16)他如何在掃羅王軍中打敗巨人哥利亞(Goliath) (撒上17) ,聲名大噪,引起掃羅嫉妒,要除之而後快,上主藉掃羅的兒子約拿單(Jonathan) 通風報訊,他聞風先遁,上主保守他避過掃羅的刀劍,並且立他為以色列的王.可是他心裡沒有緊緊跟隨上主,為了美色,殺害對自己忠誠的將軍烏利亞(Uriah) ,奪取他的妻子拔示巴(Bathsheba).(撒上11)我要回轉,我的心要緊緊跟隨上主,主的右手必會扶持我.

王位是上主給的,這位份要成為別人的祝福,我以後要好好利用這君王的身份,顯揚上主的慈愛.他祈求上主敗壞敵人的計謀,以致以後好好履行王的責任.「但那些尋索要滅我命的人,必往地底下去.他們必被刀劍所殺,被野狗所吃……因為說慌之人的口必被塞住.」(vv. 9-11)

b.      不同的身份,一樣的召喚

在聖殿的敬拜中,大衛讚美上主:「你的慈愛比生命更好」,但大衛並沒有因此而輕視生命.相反,他生活得更積極；他要透過這生命,把上主的慈愛告訴其他人.我們各人都有不同的身份,這些身份,是用來表明上主的慈愛.不同的身份接觸不同的人,不同的身份,履行不同的責任,服侍不同的群體.

上主的揀選我們是成為君尊的祭司,直接到他面前敬拜他.馬丁路德說:基督徒是最自由的君王,但同時又是最臣服的僕人.我們是最自由,因為罪都不能綑縛我們,甚麼能綑縛我們呢?因著愛心,我們甘心樂意去成為人的奴僕.身份是外在的,但自由是內在的,我們不要看自己是人的奴僕,我們是上主的奴僕,「你們是重價買來的,不要作人的奴僕」(林前7: 3)

沙士疫症的反思:曾經瀕臨死亡邊緣的人,一定會對自己的身份或職業有了不同的看法.今年是沙士(非典) 疫症十週年紀念,不少傳媒製作專輯,訪問多位抗疫醫護人員,重溫香港難忘的集體回憶.不少醫生因為嗅過死亡氣息,對醫生的身份跟以前看法不一樣了.很多受訪者分享,現在學曉醫人,不是單單醫病,他們關心病人的福祉,多於解決病患.香港中文大學校長沈祖堯說,自此以後,他更關心學生的成長和教育,最後承擔了校長的職任.有些醫生經此一役,更堅定留在公共醫療系統,不跟風為了較舒適的生活私人執業.

總結:「在曠野的一天」是大衛作為在曠野的逃者生命的寫照,聖經也用曠野的飄流比喻人生,我們的家不在曠野,乃在迦南美地,我們在世上是客旅,是寄居的,我們羨慕更美的家鄉.教宗本篤十六離任時說:他是一位朝聖者(pilgrim),步上最後的朝聖旅程.每一個基督徒都是朝聖者,行完人生的路,我們便歸家了,我們便迎見我們的上主了.今天的經文給我們一個應許,上主的慈愛,比生命更美好,上主也要使用我們的生命,去顯揚他的慈愛.阿們.



\newpage

\section{誰是那浪子}
\label{sec:193}
\textbf{盧龍光牧師 - 誰是那浪子?}
\newline
\newline
\hyperref[sec:192]{< < < PREV SERMON < < <}
~
\hyperref[sec:toc]{[返目錄]}
~
\hyperref[sec:194]{> > > NEXT SERMON > > >}


講題:誰是那浪子?　Who Is The Prodigal Son?

講員:盧龍光牧師

場合:香港中文大學崇基學院禮拜堂

日期:2013年3月10日

當我們知道耶穌基督在祂人生當中,最重要的起點,是祂的「受洗」及「試探」.

在那受試探的四十天裡面,撒旦要質疑祂的身份:「假如你是上帝的兒子,你能夠將石頭變成食物；假如你是上帝的兒子,你可以從聖殿的頂上跳下來,也不會受傷害；假如你是上帝的兒子…」.

今天我們同樣的要問我們自己:「我是誰?」然而今天我們說我們是基督徒,我們就是上帝的兒女,是真的嗎?可能在我們跌倒的時候,軟弱的時候,我們也會質疑自己.或者在我們成功的時候,我們也會質疑:我真的需要做一個基督徒嗎?其實,今天我們的人生當中,往往是受到我們對自己的認識,也決定我們的使命,工作,生命的方向.

今天的經文是選自一段大家都很熟悉的路加福音15章,有關浪子的比喻.這比喻是家喻戶曉,連一些不信主的人,大概也都聽過這個故事,甚至我們也說:浪子回頭,金不換.路加福音15:1－3節及 11節,這章裡頭有三個比喻:頭兩個是講到「迷羊」(路15:4-7)及「失去錢」(路15:8-10)的比喻,內容比較簡單:

1.    牧羊有100隻羊,迷失了一隻,他一隻也不能掉,就去尋找那一隻羊,找到了以後,就歡喜快樂；

2.    一個人丟了一塊錢,他會儘可能地去找,直到找著了,就很歡喜.

在路加福音15:1－3節這裡,好像有一個前言,就是法利賽人與文士,在質疑耶穌:祂是誰?這個人竟然接納罪人,又跟他們吃飯,祂不是老師嗎?在他們當時候,老師好像都是道貌岸然,不應該跟一些罪人,一些別人看不起的人在一起.但耶穌跟他們不一樣!

在這段經文裡面,都提到「迷失」的,能夠被「尋回」,是何等的寶貴.但是浪子的比喻卻篇幅最長的,而且裡面所講的內容,與前面兩個,好像有點不太一樣,但也有共同的地方.這段浪子的比喻,大家所熟悉的,就是那個小兒子,他好像就是主角,因為所提到有關他的事情,也是最長的.但我們知道這個故事裡有兩個兒子──小兒子與大兒子.但我們往往把注意力放在小兒子身上,而且覺得他才是浪子,因為他回頭了.

其實在這段經文裡面,「浪子」這個字,從來就沒有出現過,這個浪子使我們覺得他就是那個「浪子」,因為那個描寫,內容,在一開始的時候,他就主動要求他父親分家業,然後他也主動地離開父親,家庭往遠方去(vv.12-13a).但卻沒有描寫他為何要往遠方去?離開家,他放蕩去了──在家裡,我們都有不同性格的兒女,或許我們自己也是一個不同性格的人──這個小兒子,其實是代表一些比較主動的,有冒險精神的,他願意往遠方去,可能他的目的,並不是去放蕩的.他的目的,可能希望去開創一番事業,不要整天呆在家裡,他希望到外面去闖一下,希望能夠成功.但是人就是那麼脆弱的,人在外面的時候,就會遇到不同的試探,在這裡就記載了,他的結果是:任意放蕩,浪費錢財(v.13b),耗盡了一切就窮困起來,甚至去替人放豬(猶太人是不吃豬肉的,因為豬是不潔淨的,所以根本不會去碰豬的,不單單是骯髒,乃是猶太人不會去做的事,他也得去做),而他連豬吃的豆莢也得不到來充飢 (vv. 14-16),在這種光景底下,聖經很重要的信息,也是他主動的醒悟過來,對自己說:我父親有多少雇工,糧食有餘,我倒在這裡餓死嗎?他反省,恩考,對自己講話──這個也是人與其他動物最不一樣的地方.人在過程中會有變化的,會失落,跌倒,犯錯,但一方面我們有機會回頭的,會醒悟的,改變的.在這個故事裡面,小兒子主動的對自己說:我為什麼落在這種光景裡面?如同我方才說的,在試探當中,撒旦會對我們說話,有時候或許是我們對自己講話,又或許是上帝跟我們講話.但內容是什麼?撒旦永遠是質疑我們:唉,你這種光景也算是基督徒嗎?你這種德性也算是基督徒嗎?這講這種話,做這些事,也算是基督徒嗎?… 許多的試探,許多的質疑,甚至有許多的審判,定罪,自責.

保羅在羅馬書第八章告訴我們:上帝愛我們,捨棄祂的兒子,誰有權柄來控訴我們?誰有權柄來定我們的罪?只有撒旦,它才會控訴我們,定我們的罪.上帝只會責備我們,提醒我們,呼召我們,指出我們的錯,我們的罪,跌倒,過犯,但祂不會控訴我們,定我們的罪.

所以在這裡,這個浪子,在他生命的光景裡,他對自己說:我為什麼會這樣子?我為何要處在這樣的光景裡?他想起父親,家,他就決定:我要起來,回到我父親那裡去!而且也決定對父親說:父親,我得罪了天,又得罪了你,從今以後,我不配稱為你的兒子,把我當作一個雇工吧!這句話跟他剛離開家所講的話,完全不一樣.在12節那裡,小兒子怎麼對父親說話:「父親,請你把我應得的家業分給我」──我「應得」的,我是你的兒子,我是應該得到你的產業的,我應該的.

在我們生命裡,我們覺得有許多東西是我們「應得」的,我是上帝的兒女,你應該保守我；我是上帝的兒女,我怎麼會患上這種病啊?我是上帝的兒女,我怎麼會被人家害啊?我是上帝的兒女,我應該是一帆風順的,我們的生命裡,不應該有那麼多挑戰,苦難,我們總覺得我們是應該的….

但是這個小兒子,他醒悟過來──我不配──他願意回到家裡,不再是兒子,願意當一個雇工.他對自己的身份,自我的認識,看見他生命的經歷,過程及改變,他就有行動,不再喃喃自語了,他起來,往他父親那裡去.而父親每天就在那裡等候著孩子回來,所以遠遠看見他回來了,就動了慈心,跑去擁抱著他,親吻他,吩咐僕人把那上好的袍子給他穿,把那肥牛犢宰了,說:我們來吃喝慶祝,因為我這個兒子是死而復活,失而復得的!故事裡,一個孩子犯了這麼大的錯誤,回頭了,父親接納他.在這裡,浪子的故事可能已是高峰了,但是這裡有一句很重要的話說:我的兒子是死而復活,失而復得的(這段經文的最後,又出現了一次).

此時,大兒子回來,還沒進家門的時候,聽見慶祝的聲音(v25),就問:發生什麼事?僕人對他說:你弟弟回來了,你父親因為他無災無病地回來,把牛犢宰了.

是你,你會有什麼反應?大兒子就生氣,不肯進去,他父親出來勸他,他對父親說:你看,我服事你這麼多年,從來沒有違背過你的命令,而你從來沒有給我一只小山羊,叫我和朋友們一同快樂.

原來乖兒子心裡面有許多怨恨,心裡面有許多慾望得不到滿足,其實這個也是很正常,我們人總喜歡去比較,在家裡,朋友,同事中都會去比較,尤其是覺得父母親的偏心,老闆的偏心.大兒子可能不喜歡去冒險,很安定的,但是他心裡面藏了許多的埋怨,他跟弟弟有了比較,他知道弟弟在外頭的情況,他說,這個壞的弟弟,不值得我們去關心他,他活該,但他竟然回來了,回來時卻不懲罰他,趕他走,還要給他慶祝,這算什麼?

但父親卻勸他說:兒啊!你常和我同在,我一切的都是你的──父親已經把產業分給兩個兒子了,但他卻沒有一種歸屬感,他可能太過依賴這個家,依賴父親,所以並不覺得這個家是他的──父親提醒他說,弟弟已經拿走他的東西,這些都是你的,雖然爸爸還在,可惜他沒有對家的那種歸屬,所以縱使他常常在家裡,常與父親同在,但心卻不是如此,他的心不認為那是他的家,是他有份的,他不懂父親的心,他的心沒有與父親一起,他沒有父親那種懂得憐憫弟弟的心.

今天,我們大家都說,那個小兒子是浪子,因為他主動離開家,他要分家產,他跑到外面,壞事做盡… 還好,浪子回頭金不換,但從人情來說,這種人還是不好的,這種人到底還是應該要懲罰的.但事實上,他已受了他該受的懲罰了,他那種連豬都不如的生活,本是對他的一個審判,今天我們也常說,上帝要審判這個世界,什麼世界末日… 其實聖經裡說的審判是為了拯救,啟示錄最後幾章說,當新天新地的時候,是上帝的帳幕在人間,主耶穌要再一次降臨,並不是要毀滅這個世界,上帝乃是要更新這個世界,因為這世界是上帝所創造的,為什麼要毀滅它?就算在挪亞的時代,上帝用洪水審判這個世界,但並沒有毀滅這個世界,祂也說了:從今以後,不再有洪水毀壞地了.(創9:11)因為上帝審判的最終目的,乃是拯救.

今天你的心是怎麼想?你看到那個小兒子,他主動的離開家,而最後主動的醒悟過來,認罪,回家,父親的愛接納他了,大兒子雖然沒有離開家,但他的心並沒有與這個家連在一起；他雖然跟父親一起,他的心卻是離開父親的心很遠.究竟誰是那仍未回家的浪子?

大兒子在家門,但不願意進家門,他的心與父親的心相離很遠,是父親主動出來相勸,究竟今天誰是那浪子?小兒子當然是浪子,但大兒子可能也是個浪子,可能是更難回轉的浪子,因為他總不認為他有離開家,他不覺得他是一個浪子啊?他不覺得他有什麼不對啊?他也沒有感受到他與父親的心相離甚遠,所以他不願意進家門.

今天我們是不是也是不一樣的浪子?我們雖然好像沒有離開主,沒有離開教會,但你的心有與主連在一起嗎?你有發現你的心在與自己對話的時候,原來你還會叛逆上帝,不願意按照聖經的教導來過你的生活,按照聖經的價值標準來作為你生活的價值標準,你或許像那小兒子,承認自己也曾離開家,但你有真正地像那小兒子在聖經裡,重複說:我得罪了天,又得罪了你,我不配稱為你的兒子.

今天來到主耶穌的面前,我們是不配的.如果沒有這種「不配」的經歷的話,你還沒有機會去體會到做主兒女的甘甜.只有當你的經歷裡,生命裡,體驗到我是不配的,甚至夫妻之間──我不配做你的丈夫,太太,但你還是那麼愛我.今天我們這種不配的心,在聖經的信息裡,這是很重要的,我們常覺得「我應該的」,有這種心態,可能就是離開家的動力.唯有覺得自己不配,一個願意回轉的浪子,回到家裡；但那個大兒子很可惜,雖然常常在那裡,但他沒有真正的回家.

最後,故事結束了,也沒有告訴我們,大兒子最終有沒有進家門,你是哪一個兒子呢?你有沒有這種經驗呢?我希望我們不管是那種個性,情況,都願意回家,並發現自己的不配,來去經歷上帝的愛,就都願意回轉.在生命當中,不只是一次的回轉,你知道當主耶穌呼召門徒,彼得在約翰福音21章,耶穌問他:你愛我比這些更深嗎?這問題不是早就問過了嗎?在人群中,在路上時,耶穌不是早就說,你不愛我勝過這一切,你就不要來跟從我嗎?但彼得他們不是都一直跟著耶穌嗎?約翰福音21章,耶穌說,你來跟從我吧!

原來回轉,回應主耶穌的呼召,是否願意跟從耶穌,是一生不同階段的時候,都要做的決定,讓我們都做一個承認自己是浪子,不同情況的浪子,也同樣願意醒悟,認罪,回家.

請禱告:

慈愛天父:感謝你,因為你是我們生命的主,我們在你面前的時候,有人認為自己才是上帝,不願意承認你是上帝,今天在這裡參加崇拜的弟兄姊妹,我們的朋友,當我們願意來到這裡敬拜,參加的時候,都可能願意來思考,或許已承認你就是創造世界萬物,也創造我們生命的主,在我們父母生孩子的時候,或許有意外,但在你那裡,從來沒有意外.我們就是你所愛的,你所創造的,你給我們自由,我們可以承認你是我們的主,也可以拒絕你,我們今天可以回到你的家裡,也可以主動的離開!主啊,今天讓我們思考我們生命的經歷,無論是過去一年的,或是目前的景況,又或許是我們回頭看我們過去幾十年的人生的光景,有時候我們就發現,原來我們是一個浪子!不單單小兒子是浪子,其實大兒子原來也是個浪子,求你給我們在你面前領受你的恩典與話語,使我們有這種醒悟的能力,讓我們有謙卑,認罪的態度,也讓我們有個決心,用行動來回家.主,因為你在等候著我們,接納我們,邀請我們,願我們都回家,奉主耶穌名求,阿們.



\newpage

\section{伯大尼的筵席}
\label{sec:194}
\textbf{林豪恩先生 - 伯大尼的筵席}
\newline
\newline
\hyperref[sec:193]{< < < PREV SERMON < < <}
~
\hyperref[sec:toc]{[返目錄]}
~
\hyperref[sec:195]{> > > NEXT SERMON > > >}


約翰福音 12:1-8
\newline
\begin{longtable}{cl}
\hline
\hline
章節 & 經文 (和合本修訂版)\\
\hline
12:1 & \begin{tabularx}{0.7\textwidth}{X} 逾越節前六天，耶穌來到伯大尼，就是他使拉撒路從死人中復活的地方。 \end{tabularx} \\ \\ \relax
12:2 & \begin{tabularx}{0.7\textwidth}{X} 有人在那裡為耶穌預備宴席；馬大伺候，拉撒路也在同耶穌坐席的人中間。 \end{tabularx} \\ \\ \relax
12:3 & \begin{tabularx}{0.7\textwidth}{X} 馬利亞拿著一斤極貴的純哪噠香膏，抹耶穌的腳，又用自己頭髮去擦，屋裡充滿了膏的香氣。 \end{tabularx} \\ \\ \relax
12:4 & \begin{tabularx}{0.7\textwidth}{X} 有一個門徒，就是那將要出賣耶穌的加略人猶大，說： \end{tabularx} \\ \\ \relax
12:5 & \begin{tabularx}{0.7\textwidth}{X} 「為甚麼不把這香膏賣三百個銀幣去賙濟窮人呢？」 \end{tabularx} \\ \\ \relax
12:6 & \begin{tabularx}{0.7\textwidth}{X} 他說這話，並不是關心窮人，而是因為他是個賊，又管錢囊，常偷取錢囊中所存的。 \end{tabularx} \\ \\ \relax
12:7 & \begin{tabularx}{0.7\textwidth}{X} 耶穌說：「由她吧！她這香膏本是為我的安葬之日留著的。 \end{tabularx} \\ \\ \relax
12:8 & \begin{tabularx}{0.7\textwidth}{X} 因為常有窮人和你們在一起，但是你們不常有我。」 \end{tabularx} \\ \\
[1ex]
\hline
\hline
\end{longtable}
講題:伯大尼的筵席 A Meal at Bethany

經文:約翰福音12章1-8節

講員:林豪恩先生

場合:香港中文大學崇基學院禮拜堂

日期:2013年3月17日

一．引言

英國文豪狄更斯 (Charles Dickens) 以法國大革命為時代背景所撰之名著-《雙城記》開場之引言說:

這是最好的時代,也是最壞的時代;

這是智慧的時代,也是愚蠢的時代;

這是篤信的時代,也是疑慮的時代;

這是光明的季節,也是黑暗的季節;

這是希望的春天,也是絕望的冬天;

我們什麼都有,也什麼都沒有;

我們全都會上天堂,也全都會下地獄.

法國大革命這樣氣勢磅礡的大時代,與伯大尼小村莊的一頓晚飯有甚麼關係呢?相信沒有人認為有需要去考究它們之間的關係,其實很可能沒有甚麼歷史關係.不過,當我們用《雙城記》的引言去閱讀伯大尼的延席,會否也在小村莊的「飯局」中感受到大時代的氣勢磅礡呢?

二．經文

12 : 1逾越節前六日,耶穌來到伯大尼,就是他叫拉撒路從死裡復活之處.

12 : 2有人在那裡給耶穌預備筵席；馬大伺候,拉撒路也在那同耶穌坐席的人中.

12 : 3馬利亞就拿著一斤極貴的真哪噠香膏,抹耶穌的腳,又用自己頭髮去擦,屋裡就滿了膏的香氣.

12 : 4有一個門徒,就是那將要賣耶穌的加略人猶大,

12 : 5說:「這香膏為甚麼不賣三十兩銀子賙濟窮人呢?」

12 : 6他說這話,並不是掛念窮人,乃因他是個賊,又帶著錢囊,常取其中所存的.

12 : 7耶穌說:「由她吧!她是為我安葬之日存留的.

12 : 8因為常有窮人和你們同在,只是你們不常有我.」

三．釋經與應用

1．這是最好的時候,也是最壞的時候

伯大尼,本來是一個寂寂無名的小村莊,乏善可陳,因此,作者提到伯大尼,也沒有甚麼歷史背景或者風光人物介紹.那麼,作者如何介紹伯大尼呢?第一節如此描述:「耶穌來到伯大尼,就是他叫拉撒路從死裡復活之處」.對於作者來說,伯大尼這小村莊值得介紹的資料,就是耶穌曾經在這裡使拉撒路復活.對於伯大尼的村民來說,這村莊值得被提及的地方,很可能也是耶穌曾經在這裡使他們的小兄弟拉撒路復活.

讓我們想像,有一天,耶穌又途經伯大尼這小村莊,拉撒路的家庭,以及伯大尼的鄉親父老,會有甚麼反應呢?這位使拉撒路死而復生的救命恩人到訪,整個村莊都興奮起來,鄉親父老熱情地盡地主之誼接待耶穌,為耶穌和他的門徒預備筵席.拉撒路一家當然不會缺席,拉撒路的姐姐馬大繼續做她善長的接待工作,而拉撒路這位死而復生的焦點人物更與耶穌坐在一起.作者沒有記述耶穌和拉撒路的對話,或者,我們可以想像,拉撒路會對耶穌表示無限感激,而耶穌也可能詢問拉撒路最近的健康狀況.這是歡迎的筵席,伯大尼的鄉親歡迎他們的救命恩人耶穌到訪；這也是感謝的筵席,伯大尼的鄉親感謝耶穌把他們死了四天的小兄弟拉撒路救回來；這是慶祝的筵席,伯大尼的鄉親為拉撒路從死裡復活而舉杯.伯大尼的筵席,洋溢著感激之情；伯大尼的筵席,充滿著興奮；伯大尼的筵席,歡笑聲此起彼落.

如果伯大尼筵席的歡笑聲能夠傳到耶路撒冷,耶穌撒冷可能會這樣回應:「笑吧!吃吧!這是伯大尼村民和耶穌的最後晚餐了.難道你們不知道,這其實是耶穌的送喪筵席嗎?」為甚麼耶路撒冷會這樣說呢?因為在這筵席之前,他們已經決定要除去耶穌.約翰福音十一章記載耶穌叫拉撒路從死裡復活的事件傳到耶路撒冷,11章53節如此寫道:「從那日起,他們就商議要殺耶穌.」接著11章57節記載祭司長和法利賽人公開地發出對耶穌的通緝令:「若有人知道耶穌在哪裡,就要報明,好去拿他.」簡單純樸的伯大尼的村民當然不知道耶路撒冷的陰謀,他們怎會想到耶穌吃過這筵席之後就步向險境,就連耶穌的門徒也沒有察覺到他們的老師有殺身之禍.正在走向十字架的耶穌坐在充滿歡笑聲的筵席上,置身於慶祝被自己從死裡救回來的拉撒路的筵席,而自己卻步向死亡,耶穌的心情是怎樣的呢?伯大尼的筵席是一場喜宴,還是一場哀宴呢?這是一場喜劇,還是一場悲劇呢?

想起十多年前的意大利電影《一個快樂的傳說》,講述意大利一對猶太父子被送進納粹集中營,父親不忍年僅五歲的兒子飽受驚恐,利用自己豐富的想像力說他們正身處一個遊戲當中,必須接受集中營中種種規矩以換得分數贏取最後大獎.影片笑中有淚,將一個大時代小人物的故事,轉化為一個扣人心弦的悲喜劇.

人生的經歷,有時也是悲中有喜,喜中有悲.想起在網絡上流傳的一封信,是一位爸爸寫給他四歲女兒的最後一封信.

給可愛的女兒:

爸爸和妳玩了好多次躲迷藏,每次都一下子就被妳找出來.不過這一次,爸爸決定躲好久好久.

妳先不要找,等妳十四歲的時候(還要吃完十次蛋糕),再問媽咪爸爸躲在哪裡,好不好?爸爸要躲這麼久,妳一定會想念爸爸,對不對?不過,爸爸不能隨便跑出來,不然就輸了.

如果還是很想爸爸,爸爸就變魔法出現.因為是魔法,不是真的出現,所以不犯規,爸爸不算輸.

爸爸的魔法是:趁妳睡覺的時候,跑到妳夢裡大玩遊戲；在妳畫圖畫爸爸的時候,不管好不好看,妳覺得是爸爸,就是爸爸；當妳拿爸爸的照片看時,爸爸也在偷偷的看妳.

要記得,爸爸一直都陪著妳!妳已經是四歲的大姊姊了.爸爸要拜託妳一件事,要妳照顧和孝順爺爺,麻麻和媽咪,看妳是不是比爸爸以前做得好?有多好,媽咪會告訴你的.

爸爸猜想,我們這一次玩躲迷藏要玩這麼久,爺爺,麻麻,媽咪有時候看不到爸爸,他們一定會偷哭.偷哭就是犯規,就是失敗.

他們偷哭,妳就要逗他們笑,不然遊戲輸了以後,他們一定會哭得更厲害了.好不好,寶貝?我們是同一國的,來比賽看妳厲害,還是爸爸?準備好了嗎,比賽就要開始了!

伯大尼的筵席,是愉快的筵席；不過,筵席的外圍,卻是殺氣騰騰.這是最好的時候,也是最壞的時候,好壞交雜,好壞未定的狀況,似乎也是人類共同的經驗.

2．這是智慧的時候,也是愚蠢的時候

在這慶祝拉撒路復活的感恩筵席中,怎可能沒有拉撒路的另一位姐姐馬利亞在場呢?而且,按作者敘述這個故事的方式來說,伯大尼筵席的主角不是拉撒路,而是馬利亞.因為,馬利亞的行動吸引了所有在場人士的注意,也成為了討論的話題.她做了甚麼事情呢?經文如此描述:「馬利亞就拿著一斤極貴的真哪噠香膏,抹耶穌的腳,又用自己頭髮去擦,屋裡就滿了膏的香氣.」(12:3)她為甚麼做這樣的事情呢?作者沒有交待.然而,按著劇情理解,這很可能是馬利亞感激耶穌把他的弟弟從死裡救回來的表達.如果我們細心留意約翰的敘述,馬利亞在第十一章也曾經跪在耶穌的腳前:「馬大說:「夫子來了,叫你.」馬利亞聽見了,就急忙起來,到耶穌那裡去．．．看見他,就俯伏在他腳前,說:「主啊,你若早在這裡,我兄弟必不死.」(11:28-32)馬利亞那次俯伏在耶穌腳前,是表達悲傷,因為耶穌在拉撒路死了之後才來到,錯過了救活拉撒路的時機；今次馬利亞在筵席上俯伏在耶穌腳前,是表達感激,感激耶穌把她以為一去不返的兄弟復活過來.這份感激之大,令馬利亞把自己一點一滴累積多年的嫁妝也捨棄.或者,對馬利亞來說,傾出這真貴的香膏仍然不足表達對弟弟死而復生的感激.是的,就算是一斤極貴重的香膏,都不能和一個人的生命相比,何況這個人是自己的兄弟呢?

參加筵席的人士在飯桌上高談闊論之際,突然被一股香濃的氣味刺激.他們沿著香氣抬頭望過去,看見馬利亞的行動,就感到不高興.其中一個叫猶大的門徒出口批評:「愚蠢!這是愚蠢的行為!難道你不知道,這種香膏是很值錢的嗎?如果你不要,把它賣了,所得的錢可以用來賙濟窮人,為甚麼這樣浪費呢?」批評別人愚蠢的人,通常都認為自己有智慧.這位批評者如何有智慧呢?他的智慧在於他懂得曉以大義去滿足個人自私的慾望.這位批評者的「智慧」,連作者也看不過眼,給讀者補充一句說:「他說這話,並不是掛念窮人,乃因他是個賊,又帶著錢囊,常取其中所存的.」(12:6)遇到某些滿口仁義道德,說話胸還普世,但卻只是中飽私囊的人,怎不叫人氣憤呢!

耶穌的回應卻比作者溫和,他沒有當面揭露批評者的偽善,只是以勸解的口吻說:「由她吧!她是為我安葬之日存留的.因為常有窮人和你們同在,只是你們不常有我.」(12:7-8)賙濟窮人是美德,耶穌沒有因人廢言.不過,馬利亞此舉是否愚蠢呢?猶大的批評是否智慧呢?要講道理,當然是猶大有道理,他提出的是「為最多人帶來取大的好處」這個被廣泛接受的原則.不過,耶穌的回應,卻讓人感到馬利亞的行動才是令耶穌感到窩心的行動,因為她的行動正確地洞察了時機,也配合了時機.說對了道理,不等於做對了事情.對於幾天後就赴刑場,而跟隨自己出入的門徒竟也慒然不知的耶穌,馬利亞的行動彷彿知道自己的死亡,彷彿預備自己的安葬.講道理,馬利亞可能遠遠不及耶穌的門徒,然而,她卻在合適的時候,為耶穌做了合適的事情.

在合適的時候做合適的事情,往往更能打動人心.

有一個男士問他的女朋友:「如果有一位比我更好條件的男士追求你,你會走嗎?」那位女士回答:「不會.你在合適的時間出現,做了合適的事,在我來說已經是無可取替的好了.」

在伯大尼的筵席中,有人好像嬴了道理,卻暴露了愚蠢,有人看似欠缺道理,卻被耶穌悅納.社署的家庭教育宣傳海報寫著:「贏左場交 輸左個家 值得咩」.怎樣做才是有智慧,怎樣做才是愚蠢,值得我們再想一想.

3．這是光明的季節,也是黑暗的季節

杯盤狼藉之際,夜幕開始低垂,黑暗漸漸降臨大地.這天晚上,耶穌和門徒在伯大尼留宿一宵.眾門徒可能仍然為著馬利亞以香膏抹耶穌的腳的行動,以及猶大的批評而議論紛紛；同時,拉撒路姐弟和鄉親們可能忙著筵席的善後工作.就在此時,黑夜不動聲色地吞噬著大地,提示著耶穌第二天就要步入耶路撒冷這個刑場.縱使現場仍然眾聲喧嘩,耶穌難免泛起「一個人在途上」的孤單感.有些時候,感到孤單,不是因為身邊沒有人,而是身邊的人不能夠體會到你的狀況；就算有人能夠稍為明白到你的狀況,也不能分享你的命途.路,仍是自己一個人踏上；耶穌步向十字架的路,更只能是他自己一個人踏上.

黑暗隨者夜幕籠罩大地,不單止籠罩著清楚地知道自己命途的耶穌,黑暗不動聲色地伸延到拉撒路,這位無辜者的頭上.作者把鏡頭移離伯大尼的筵席,帶我們都耶路撒冷,在那裡,「祭司長商議連拉撒路也要殺了」.(12:10)為甚麼要連拉撒路也殺掉呢?約翰福音的作者在整卷福音書開首已經透露了玄機:「太初有道,道與上帝同在,道就是上帝.這道太初與上帝同在.萬物是藉著他造的；凡被造的,沒有一樣不是藉著他造的.生命在他裡頭,這生命就是人的光.光照在黑暗裡,黑暗卻不接受光.」耶穌是那光,黑暗不接受他,要消滅他；拉撒路是耶穌叫他從死裡復活的,分享著生命的光,見證著那光,因此,黑暗也容不下他,要吞噬他.

我們身邊的社會事件,政局時勢,又或者我們個人的遇遭,身體狀況,有時令我們感到黑暗覆蓋大地,好像處身於絕望的寒冬一樣,不見天日.不是嗎?連耶穌也被黑暗吞噬,為耶穌作見證的也要被黑暗吞噬,還有甚麼盼望呢?十年前的這個時候,我們這個城市被一種神秘的病毒襲擊,整個社會被恐慌籠罩.傳媒每天告訴市民,那座大廈有人病發,那個屋苑有人「中招」,人人自危.有一位醫院總監回顧說:「當時,醫院裡有一位醫生證實被感染,他召集員工會議.他在台上望下去,看見幾百位同事的眼神盡是擔憂和恐慌.這樣大規模的恐慌眼神,是我從沒見過的.」在彷彿被黑暗覆蓋的大地上,有一遠處傳來的聲音,隱約可聽.這是遠古詩人的歌聲,他唱著說:「我若說:「黑暗必定遮蔽我,我周圍的亮光必成為黑夜」,黑暗也不能遮蔽我使你不見,黑夜卻如白晝發亮,黑暗和光明在你看都是一樣.」(詩139:10-12)

耶穌沒有被黑暗嚇阻.第二天,他抖擻精神,繼續前行,以迎向黑暗的態度挑戰黑暗.耶穌知道,挑戰黑暗,就是光明.他說:「人子得榮耀的時候到了.我實實在在地告訴你們,一粒麥子不落在地裡死了,仍舊是一粒,若是死了,就結出許多子粒來.」(12:23)原來,黑暗並不能吞噬生命,黑暗只是生命的轉化過程.耶穌呼籲我們,不要被黑暗嚇倒,跟隨他的腳步,挑戰黑暗.十年前的這個時候,當我們的城市被黑暗籠罩的時候,縱使醫護人員也不能自保,但他們沒有被黑暗嚇阻,竭盡全力,迎向黑暗,挑戰黑暗.有醫學院教授在醫院發起祈禱會,以從上主而來的光明突破黑暗,以上主而來的盼望擊退絕望.光明,就在人們挑戰黑暗的行動中顯現,就在人們不被黑暗嚇倒的時候出現,也在人們仍然互相珍惜和鼓勵的狀況下出現.當時有一位醫科五年級生,在威爾斯親王醫院8A病房實習時染病,留醫期間父親在樓上九A病房去世,令他萬念俱灰.他的女友,現時的太太知道他的情況,哀求主診醫生及教授,只求見他一面.女友最後簽下生死狀,穿上保護衣便冒險入病房,突破當時不准探病的禁令.這位校友回憶說:「突然有個著保護衣ge 人搭我膊頭,我心諗邊個護士咁好.認得係佢,即刻喊jor 出黎.佢同我講,你要加油,你爸爸唔想睇到你lum !」醫護人員亦不斷加以鼓勵,那位教授更安排他穿上保護衣,到病房走廊送別父親.這位醫科學生十年前患病的時候目睹醫護人員無畏無懼,令他矢志成為呼吸科醫生,以得來不易的生命回饋病人.

黑暗籠罩伯大尼的筵席,但黑暗沒有嚇阻耶穌前進.從遠古傳來的歌聲激勵著耶穌前行:「黑暗也不能遮蔽我使你不見,黑夜卻如白晝發亮,黑暗和光明在你看都是一樣.」

四．總結

一個小村莊的筵席,透視著人世間好與壞,智與愚,光與暗的交雜,我們有沒有在旁觀這場筵席的時候看見自己身處的狀況呢?對我們來說,這是好的時候,還是壞的時候呢?這是智慧的時候,還是愚蠢的時候呢?這是光明的時候,還是黑暗的時候呢?當年,耶穌就在這好與壞,智與愚,光與暗混雜的世界中生活,在這個絕望與希望爭持的舞台上活出他的使命.今天,他呼籲我們跟隨他的腳蹤,在這舞台上「做場好戲」,好好生活.因為,好與壞,智與愚,光與暗的混雜,都不能遮蔽上主的眼睛令他看不見我們,就算我們處身這混雜的狀況,他仍然能夠引導我們,與我們同行.



\newpage

\section{真知道神救恩中的三}
\label{sec:195}
\textbf{陳天祥博士 - 真知道神救恩中的三重恩典}
\newline
\newline
\hyperref[sec:194]{< < < PREV SERMON < < <}
~
\hyperref[sec:toc]{[返目錄]}
~
\hyperref[sec:196]{> > > NEXT SERMON > > >}


日期:2013年3月17日

主題:真知道神救恩中的三重恩典

證道:陳天祥博士

我們今天要看的聖經是剛才頌讀的以弗所書經文.講道有很多種,其中一種是解經式的講道.無論聽道或解道,我也很喜歡解經的形式,因為我覺得這是神對事奉人員的要求,要分解真理的道,同時也是信徒的需要.聖經是新生嬰孩的靈奶,是能夠讓我們成聖的真道.當我們查經的時候,我們可以知道訊息的來龍去脈,看看是否合乎聖經的解釋.

我們先同心祈禱:「我們的天父,當我們聚集敬拜你的時候,我們有機會看你的話語,叫我們的心靈向著你,讓我們聽道時能分辦.看見你榮耀的時候,叫我們的內心能湧出敬拜和讚美.當看到我們需要負的責任,看到我們本相的時候,我們的內心要立志改變.求主祝福我們以下的時間.我們的祈禱是奉主耶穌基督的名.阿們.」

如果你問保羅,作為一個基督徒,最重要的是什麼?你想保羅怎樣回答?他的答案記載於腓立比書第3章.他說他以認識耶穌基督而自豪.這是保羅作為基督徒覺得最重要的事.如果你問保羅作為教會最重要的是做什麼?他的答案記載於我們剛才讀的聖經－以弗所書第1章.我們打開剛才已讀過的以弗所書第1章15至23節.我們最主要看三件事.第一件事是保羅對以弗所教會的信仰有一個反應.保羅對他們有一個要求.這要求是兩個知道.第17節:「求我們主耶穌基督的神,榮耀的父,將那賜人智慧和啟示的靈,賞給你們,使你們真知道他」.我讀的是舊譯本聖經,跟大家看的有一些不同；第二個知道在第18節:「並且照明你們心中的眼睛,使你們知道他的恩召有何等指望」.

我喜歡舊譯本聖經,覺得較貼近原來文字.原來聖經有兩個知道.第一個是真知道神.另一個是知道以下三樣東西,這三樣東西是並排下去的.這裡有兩個知道,而第二個知道中有三個不同的特質.這三個特質是什麼呢?在第18節:使你們知道他的恩召有何等的指望－恩照的指望.他在聖徒中得的基業有何等的豐盛－基業的榮耀.第三,並知道他向我們這信的人所顯的能力是何等浩大－是聖靈內住時所產生浩大的能力.這是三個不同的品質.

這裡有一個反應,兩個知道,三個品質.我們首先看反應.你看保羅的反應.在第15節:因此我既聽見你們順從主耶穌,親愛眾聖徒就為你們不住的感謝神.祈禱時常聽到你們求我們主耶穌基督的神,榮耀的父,將那賜人智慧和啟示的靈賜給你們,叫你們真知道他.大家覺得保羅說這些話是以什麼心境呢?是安心,是有很大要求,是擔心,還是興奮呢?大家覺得怎樣?他說因此我聽見你們順從了耶穌,我就會為你們不住的感謝神.是興奮嗎?是的.保羅在這裡表現有要求的興奮.很興奮,但有要求,好像你的兒子或親人被中大或港大收錄,你會很興奮.如果他們被耶魯或劍橋大學收錄,也會很興奮.你會有很多叮囑,告訴他們要怎樣努力讀書等等.這是有要求的興奮.為什麼值得興奮呢?正如第1章3至14節所講,神有莫大的恩典,而這莫大的恩典,這班信徒已經得著了,因此有莫大的興奮.

第3至14節中有三個稱頌.第一個稱頌在第6節:使他榮耀的恩典得著稱讚.第二個稱頌在第12節:叫他的榮耀從我們這首先在基督裡有盼望的人可以得著稱讚.第三個稱頌在第14節:這聖靈是我們得基業的憑據,直等到神之民被贖,使他的榮耀得著稱頌.這是很大的恩典.保羅提到三次稱頌,信徒已經得著了.是否很興奮呢?不單這樣,這班信徒確實得著救恩.很多人回教會,不過回教會人未必一定能信靠耶穌.但這班人真真正正能信靠耶穌基督,為什麼這樣說呢?

在第15節:因此我既聽見你們信靠主耶穌,親愛眾聖徒.他們有兩個表徵.第一是順從主耶穌.第二是對其他弟兄姊妹相愛.這兩個表徵證明他們是真正的基督徒.他們的信不只是參加聚會這般簡單,而是順從主的吩咐,按主的吩咐而行,順從主耶穌,而且對其他弟兄姊妹表現愛心.正如約翰一書所講,凡有愛心的都是由神而生的.這確實是一班愛主,真正的基督徒.他們真正有上文所講三個稱頌神的大恩典.是否應該很興奮呢?既然那麼興奮,為什麼有要求呢?因為我們得到神的恩典,不單只因信得救那麼簡單.保羅的要求是完全合理的.神的恩典是恩上加恩.神不單讓我們被稱為義,還要我們成聖.在以弗所書第1章4節所講,就如神在創立世界以前,在基督裡揀選了我們,使我們在他面前成為聖潔,無有瑕疵.這是神的心意,不單要我們稱義,而且讓我們成聖.

羅馬書第8章提到律法的義能成就在我們這些不隨從肉體,只隨從聖靈的人的身上.這是神的心意,不止稱義,而要成聖.不單單有生命,還要像約翰福音所講,要有豐盛的生命.不單單被召,而是要被選上.不要做被召但不被選上的人.不是要生命中受過試探,而是要我們成全完備地勝過試探.

神最終的目的是讓我們進到祂的榮耀,與祂分享祂的榮耀.正如羅馬書第8章35節所講,神預先定下的人就召他們來,召他們來的人便稱他們為義.稱他們為義的人就叫他們得榮耀.這樣才真正享受救恩的豐盛.大家看見這得豐盛的路嗎?保羅就是知道這得豐盛的路,因此對這班真正得救的以弗所信徒說:我對你們有要求.這班信徒像小孩子般要成長.我們也要成長,成長的重要一步是保羅在以弗所書所提到的兩個知道.這兩個知道是知道神,而第二個知道是知道人從神所得到豐盛的恩典.

這兩個知道是相連的.在舊譯本,第一個知道是真知道神,與第二個知道之間有一個「並」字.在新譯本,這個「並」字卻放在後面.其實在原來文字沒有這「並」字,但在舊譯本可看到兩個不同的知道.原本文字用了不同的字.原來文字的意思是因你對神的認識,所以你知道恩典對你來說有多大的豐盛.

我們或者可以從這兩個字明白它們的關連.這是兩個不同的字.第一個是真知道,在聖經曾出現二十次.第二個是知道,在聖經出現很多次,有653次.這兩個知道不是特別向著神,不是指奇怪的事或要超乎什麼才知道.很多時候指對普通事物的知道.在腓立比書第1章9節提及,你們的愛心,在知識的各樣見識上,多而又多.那個在愛心和知識的見識便是真知道的意思.至於出現了653次的第二個知道,在很多地方是講人的事情.例如以弗所書第6章提到,信徒推基古將保羅的事情告訴以弗所人,讓他們知道保羅的情況.這知道便是第二個知道.這是我們一般心理運作的認知.

我們可以問一個問題,我們可以明白怎樣知道神.你如何知道一件事?怎樣知道一個人呢?怎樣認識一個人呢?怎樣認識一件事情呢?開始時可能知道他的名字.如果你不知道孔牧師的名字,代表你不太認識孔牧師.你不知道一個人的姓氏怎樣寫,代表你不太認識他.此外,要認識他的事件和他以前做過什麼.好像剛才介紹我時也用一些事件,讓大家認識我.他的工作,作品及你與他的相交也有很大關連.例如要認識球,不只是知道它叫什麼或它是圓形.對於踢足球的人跟打籃球的人,他們對球的認識也不同.大家要有互換和相交才有認識.

你可能只透過報紙認識一個人,那你不算認識他.如果通過電郵,吃飯,見面,傾談或面書認識,可能會多些認識.更深入的是從這些表面事件明白這人的本質,例如他不變的性格和持守的原則,以進一步明白這人.我們明白神也是這樣.知道祂的名字和事蹟.我們祈禱與祂相交.我們看見祂如何引導和扶持我們.從這些事蹟,我們進一步明白祂的性情.以弗所書第1章提到,神在救恩中給我們恩召,給我們將來的基業,還有聖靈內住.這是神的工作.我們透過這些認識神,我們從這些表面的東西認識神不變的性格.

以弗所書第1章8節提到神的愛.我們體會神的愛,我們能體會神諸般的智慧和聰明.我們透過這些知道神原來那麼愛我們.神用智慧和聰明將救恩彰顯出來.我們能認識神,真知道神.第二個知道與第一個知道應該有相連的關係.

中英文聖經的看法很不一致.中文聖經譯作真知道,有些聖經譯作深深的知道,有些譯作充份知道.呂振中譯本和現代中文譯本則沒有加任何東西.英文聖經幾乎沒有加任何東西,除了NIV譯本譯作充份認識,因此很不一致.這個字是否有殊意思呢?根據字典,加了以上字首的字指特殊,有方向性,不是空泛的,是有焦點的.它也引用了其他人的說法,當中也有這意思.除了這意思外,認識這字也帶著經驗.不單指頭腦上的認識,還帶著經驗.在很多情況下,聖經人物也透過所參與的工作和事情明白神.例如亞伯拉罕,他稱神為耶和華以勒.大家記得那件事嗎?因為神在他預備獻上以撒前已為他預備一隻羊羔,所以他稱神為耶和華以勒,神必預備.例如摩西,神帶領他戰勝敵人時,他認識的神為耶和華尼西,即神作他的旌旗,戰勝的旗.當基甸遇見神,當時他可能喪命,但最後生命得以保存,於是他認識的神是賜他平安的神.因此,我們會從日常生活經歷中經歷神這個字的意思.綜合來說,真知道指透過與我們有特別關係和有份參與的事對神的認識.是什麼呢?

以弗所書第1章前半部分提到神給我們的救恩,我們真認識神,便認識神的恩召,認識將來一切在天上地下也在基督中同歸於一,於是我們也在基督裡面得到基業.我們對神的認識也包括神要聖靈住在我們的裡面,使我們有印記,使我們將來的基業有憑據.這些是神的工作.我們透過這些工作認識神.以下的字是要看見的意思,好像是自己與自己有特別的關係.我認為真知道神是神的工作.後面第二個知道是將神的工作和神的恩典內化,成為自己的特質.將神的工作內化,成為自己的特質,我覺得這是保羅對以弗所人的要求,將外面看到神的工作,內化成為你個人的特質.什麼是內化?如果外面有規矩,如果你覺得外面的規矩不是自己的,是被人強迫的,當內化後,這些規矩便成為你自己持守的原則.當外面的規矩成為你的原則,便是內化.小孩子成長時會看父母怎樣做.女孩會學媽媽,男孩會學爸爸.性別的身份從外在的做法變成內在的身份,便是內化,外面的會變成內面的.因此,我們可以說經文是神的工作的內化.經文中有清楚的支持.剛才提過神的工作,第一是恩召,第二是基業,第三是聖靈內住.保羅說:你們要知道,要內化,變成恩召的指望,基業的榮耀和聖靈內住的能力.將神的工作,神的恩召,神的基業,聖靈的內住變成恩召的指望,基業的榮耀和聖靈內住的能力.我們將神的工作內化.

無論是第一個還是第二個知道,也是非常困難.你看的事情只是表面,要進入事情的本質是很困難.神做很多啟示的工作.除了在以弗所書第1章提到的兩個知道外,還有第三個知道,在第9節:都是照他自己所預定的美意,叫我們知道他旨意的奧秘.這知道是使知道,是神使我們知道,給我們很多啟示,使我們知道.第二個支持我們知道的是聖靈,在第17節:將那使人智慧和那啟示的靈賞給你們.透過神的恩典我們能夠知道.第一,我們知道神的思典很浩大,有恩召,基業和聖靈內住,內化為我們的指望,榮耀,能力.這知道在信徒中有極大的需要.因為救恩除了是神的工作外,人也需要配合.你不知道,便不懂怎樣配合.

這也是讀科學時領會的原則.科學有四大目的,第一是措述／措寫／清楚描寫；第二是知道原理和理論；第三是預測；第四是操控或改變.同樣地,我們在信仰必須要預先知道.否則,我們無法配合得救的功夫,成為聖潔的子民.所以我們不要經常提倡只要有簡單的信心信靠神便足夠.我們要知道,要努力尋求和知道更多.這知道不只是頭腦上的知道,而是經歷上的知道.聖經希伯來書的作責備希伯來書的收信者,說他們不熟習仁義的道理,只能夠吃奶,不能吃乾糧,不能明白比較深奧的道理.我們要知道,當我們知道時,我們才能作出配合,作成得救的功課.

這三個內化對我們很重要.大家知道這三個內化在時態上有分別嗎?按照斯托德的註釋,恩召是從前的－創世以前神的恩召,神的揀選.基業是將來的－將來得基業.聖靈內住得能力是現在的,橫跨以前和將來.這三個不同的時態對現在也有影響.雖然恩召是以前的,對現在還有果效和影響力,影響指望.基業是將來的事.基業有兩個意思.第一,我們成為神的基業.第二,神成為我們的基業.兩個解釋在以弗所書第1章前半部分也有出現.在現在出現是什麼意思呢?解作我們是神的基業還是神是我們的基業呢?有些註釋書說前者,有些說後者.不排除兩個意思也有.無論如何,這是將來發生的事,現在來說構成榮耀的感受,是價值的感受.當然,現在的聖靈內住讓我們發揮聖靈給我們的能力.這三個特質讓我們現在嘗到信仰的好處.第一是恩召的指望.是指望,不是盼望.我喜歡聖經這個譯法,好像有期盼,現在所發生的.因為得到的恩典很大,所以你對信仰有期盼和指望,希望能在我們的身上發揮得淋漓盡致.

我們在聖經看到其他例子,可以表達這期待的心情.彼得後書第1章3至5節提到,神已經將關乎生命和敬虔的事賜給信徒,又說這是極大而寶貴的應許,賜給信徒,叫我們能脫離屬於情慾的敗壞,能與神的性情有份.彼得告訴我們,因此要分外殷勤.這份外的殷勤有期待的意思.你有否好像彼得所提到有份外的殷勤?如果是,你便能將恩召內化,成為你的期待.沒有期待的人,是憂鬱的人.憂鬱的人很愁苦,對所有事情沒有興趣和沒有期待,甚至會自殺,不想活下去.在越戰時,很多美軍被俘虜.他們被殘酷對待.他們回國後生病,變得沒有期待hopelessness.這些人受攻擊時不懂反抗,是完全沒有期待,生病的人.生命中當然有情況,屬靈亦然.一個在屬靈中沒有期待的人可能指停止聚會的人,對聖經和神沒有興趣,甚至離開教會,沒有期待,不能將神的恩典內化為指望.如果我們有期待,裡面充滿動力,會去認識,參與,成長.

不單這樣,我們內化時可以將基業變成榮耀的價值.什麼是榮耀?榮耀是光彩,榮譽.當然神才有真正的榮耀.人的榮耀由神而來.未信主之前,聖經告訴我們,人因自己所犯的罪虧缺神的榮耀.我們信主後便能分享神的榮耀.那種榮耀的感覺是什麼呢?心理角度看是自卑.覺得自己什麼也沒有,很自卑.一個信靠神的人應該在任何人面前也不會顯得自卑.我們在聖經看到一個明顯的例子－雅各.他在法老面前做了什麼呢?他為法老祝福.很奇怪吧!在法老面前,因為信仰的原故,雅各一點自卑也沒有.雖然在自己「沒有」的情況下,他仍然明的自己的「有」是何等的大.這便是我們的榮耀.

我們會否也這樣呢?追從聖經原來能幫助我們心理的健康.心理學上,我們追求兩樣關乎自我的東西,第一是自重self-esteem,第二是自賞self-regard.這兩樣東西對我們的心理健康很重要.原來我們將神的恩典內化的時候,我們便有這種特質.最後是聖靈內住的能力.NIV譯本將其譯作無比的能力incomparable power.<br />沒有東西能勝過的.沒有軟弱是不能讓這能力幫助.我們可能有什麼軟弱呢?我們可能有肉體的軟弱.我們面對世界,世界有敵擋神的風氣.我們在面對這世界時可能會軟弱.我們面對魔鬼的引誘時可能會軟弱.我們透過聖靈能發揮應有的能力,勝過我們的肉體,勝過世界,勝過魔鬼.不過,耶穌基督的恩典已經勝過所有,勝過肉體,勝過世界,勝過魔鬼.聖經中很多地方記載我們怎樣勝過肉體.在羅馬書第8章,肉體有軟弱,所以我們不能行律法的義.但我們在聖靈的工作裡面,聖靈賜給我們生命的糧,改變我們,使律法的意義成就在我們這些不隨從肉體,只隨從聖靈的人的身上.勝過肉體,勝過世界.約翰一書提到在我們裡面的比這世界還大.我們勝過世界.保羅說:就他而論,世界已經被釘在十字架上.就世界而論,他已經被釘在十字架上.是分割,沒有關係,是厭惡.如果有不良的風氣,神怎樣幫助我們?聖靈怎樣幫助我們?好讓我們厭惡這些東西?讓我們不會對這風氣就犯.神怎樣幫助我們勝過魔鬼呢?

希伯來書第2章說:耶穌基督的死,是敗壞那掌死權的魔鬼.因此,我們可以跟著耶穌說,當有試探的時候,你說撒旦,退在我後面吧.我們在聖靈裡有足夠的能力.這是心裡學所說的self-efficacy自我效能.這是很好的心理特質,使我們能健康地生活.當我們內化神的工作時,我們亦有自我效能.

各位弟兄姊妹,你希望成為心理健康的一份子嗎?原來內化神的工作可以達成這目的.當我們的心靈健康時,我們的心理便健康.不過,你一定要信,而且要做一個好基督徒.什麼是好基督徒?很多心理信仰的研究將基督徒分為兩種.一種是內在intrinsic的基督徒,另一種是外在extrinsic的基督徒.

什麼是外在的基督徒?他們會利用教會和信仰,達成自己地上的目的和好處.這是外在的基督徒.內在的基督徒會將神的生命和救恩活出來.這是內在,由裡面指揮,生命由裡面帶領,要活出信仰的真實.你們要信,要成為一個內在的信徒.這樣你們便能得到心靈的健康,當然也得到心理的健康.

我們祈禱:「求主幫助我們,讓我們看到神工作的偉大.當我們內化這一切事情的時候,能夠成為我們裡面的特質,成為一個在地上有能力的人.無論在什麼環境,我們也能夠做應該做的事,討你的喜悅,榮神益人,求主帶領我們.我們祈禱奉主耶穌基督的名.阿們.」



\newpage

\section{金色的麥穗是下垂的}
\label{sec:196}
\textbf{張闊先生 - 金色的麥穗是下垂的}
\newline
\newline
\hyperref[sec:195]{< < < PREV SERMON < < <}
~
\hyperref[sec:toc]{[返目錄]}
~
\hyperref[sec:197]{> > > NEXT SERMON > > >}


腓立比書 2:5-11
\newline
\begin{longtable}{cl}
\hline
\hline
章節 & 經文 (和合本修訂版)\\
\hline
2:5 & \begin{tabularx}{0.7\textwidth}{X} 你們當以基督耶穌的心為心： \end{tabularx} \\ \\ \relax
2:6 & \begin{tabularx}{0.7\textwidth}{X} 他本有神的形像，卻不堅持自己與神同等； \end{tabularx} \\ \\ \relax
2:7 & \begin{tabularx}{0.7\textwidth}{X} 反倒虛己，取了奴僕的形像，成為人的樣式；既有人的樣子， \end{tabularx} \\ \\ \relax
2:8 & \begin{tabularx}{0.7\textwidth}{X} 就謙卑自己，存心順服，以至於死，且死在十字架上。 \end{tabularx} \\ \\ \relax
2:9 & \begin{tabularx}{0.7\textwidth}{X} 所以神把他升為至高，又賜給他超乎萬名之上的名， \end{tabularx} \\ \\ \relax
2:10 & \begin{tabularx}{0.7\textwidth}{X} 使一切在天上的、地上的和地底下的，因耶穌的名，眾膝都要跪下， \end{tabularx} \\ \\ \relax
2:11 & \begin{tabularx}{0.7\textwidth}{X} 眾口都要宣認：耶穌基督是主，歸榮耀給父神。 \end{tabularx} \\ \\
[1ex]
\hline
\hline
\end{longtable}
講題:金色的麥穗是下垂的　A Golden Wheat Bows

經文:腓立比書2:5-11

講員:張闊先生

場合:崇基學院禮拜堂主日崇拜

日期:2013年3月24日

經文:腓力比書2:5-11

5:你們當以基督耶穌的心為心:

6:他本有　神的形像,不以自己與　神同等為強奪的；

7:反倒虛己,取了奴僕的形像,成為人的樣式；

8:既有人的樣子,就自己卑微,存心順服,以至於死,且死在十字架上.

9:所以,　神將他升為至高,又賜給他那超乎萬名之上的名,

10:叫一切在天上的,地上的,和地底下的,因耶穌的名無不屈膝,

11:無不口稱耶穌基督為主,使榮耀歸與父　神.

引言

記得在我高中的時候,書店裡擺出了很多關於成功學的書,這些書有的是教人如何包裝自己的外表來取得別人的好感和認同；有的是如何在語言上包裝,以達到一鳴驚人的效果；還有的是在動作上,細節上下功夫,讓自己的身上能夠散發出一種獨特的魅力,傳達出自己是擁有美好品格的資訊.

如今,與中國傳統的文化那種溫文爾雅,謙謙君子看似相悖的電視節目紅遍中國大江南北,各種選秀類,PK類的節目層出不窮,相親類的節目也是爭先恐後地出.

當我們看到這種現象的時候,不得不承認,全球化的影響已經深入到世界的許多地方.如今,如果你不懂得如何展現自己,不懂得如何表達自己的才幹,甚至不懂得如何誇大自己的能力,那麼你就會被認為沒有自信,被認為沒有實力,被認為沒有能力而最終被淘汰.因為這個世界的腳步太快了,有誰願意駐足停下來仔細的去品讀你呢?

那麼事實是否真的如同書中所說的那樣,只要你靠這些外在的包裝就一定能夠取得成功呢?就一定能夠讓人刮目相看你呢?就一定能夠展現出你超凡的魅力呢?是否不這樣就一定不會成功呢?

我想說,真正需要在意和包裝的,是你的內心,只有你由內而外發出的美好品質才會感染人,才會光照人.如果單看耶穌的外表並不一定出眾,他的衣著並不一定講究,他的舉止並不是最文雅的.但是他那顆舍己的心卻讓無數人的生命得到釋放,得到轉變.所以今天我們透過腓力比書2章,尤其是5-11節來看看,在耶穌身上所體現出的那種謙卑的心將如何指引我們現在的生活.

一,以基督的心為心

讓我們來到腓力比書2:5,我們也許會問這樣一個問題,什麼是耶穌基督的心呢?我們又要怎樣才能按照他的心去做呢?如果我們通讀5-11節,甚至是通讀第一章和第二章就會發現,基督耶穌的心就是甘願謙卑順服,甘願舍己犧牲的心.那麼我們該如何謙卑,如何舍己呢?我想在回答什麼是謙卑,如何謙卑之前,先從反面入手,看看什麼樣的人會驕傲.正所謂知己知彼,百戰不殆.通過驕傲來看謙卑.

二,從擁有來看驕傲

耶穌基督是神的兒子,有神的形象,而且他自己也就是神.這一點是關於三位一體教義,我不在此贅述了.我要帶出的資訊就是,上帝有什麼,他就有什麼.上帝的威嚴,榮耀,能力,權柄,冠冕,他都有.正如後面半句所說,這一切不是搶奪的,是本來就有的.

那麼我們擁有什麼呢?當我們的擁有有變化的時候,我們會不會產生驕傲呢?下面我們就從擁有的角度來分析三類容易驕傲的情況:

1,從沒有到有

在前一兩個月,我們普通話團契,安提阿團契一直在查創世記關於亞伯蘭的幾章經文,其中有一次查考的就是16章,亞伯蘭和夏甲同房生了以實馬利.在那段經文中記載,亞伯蘭當時已經八十多歲了,撒萊也已經年紀老邁了,不過兩個人還是沒有生子.上帝卻很多次向亞伯蘭應許,甚至是立約,說要賜給他一個後裔,可以成為他的繼承人,但是一直都沒有實現.亞伯蘭一直住在迦南,已經據他聽從上帝的呼召,相信上帝的應許,並出吾珥有十年了.但是在這十年裡,他沒有後代.最後撒萊坐不住了,於是把自己的使女夏甲交給亞伯蘭,讓她與亞伯蘭同房,以便得以生子,自己好被建立.亞伯蘭接受了這個建議,並且在同房之後,很快夏甲就懷孕了.這本來是一個皆大歡喜的結局,但是聖經沒有這樣結束,反而加入了一句話,使得整個故事再起波瀾.

聖經記載說:“她見自己有孕,就小看她的主母.”我想問,夏甲是哪裡人?對,是埃及人.但當時亞伯蘭住在的是迦南.夏甲是什麼地位呢?她是撒萊的使女.可見,夏甲背井離鄉,身份地位又不高,而且說讓她和誰同房就要同房.由此可見她當時是很卑賤的.那麼正是因為她的卑賤,正是她看到她肚子裡的孩子對於亞伯蘭和撒萊如此重要,所以她那顆卑微的心一下子反彈了,變得異常自信,得意,甚至驕傲,以至於輕看了她的祖母.

我稱這樣的心理為暴發戶的心理.那從因為之前窮,突然之間買彩票一夜中彩得了大獎的情況是會讓人驕傲,甚至是忘乎所以的.現實中,我們是否也有同樣的情況,作為學生,考試突然名列前茅,於是就高傲自滿；作為新科博士,覺得擁有了學位就高人一等了；作為教授,帶出來的學生個個都出人頭地,於是就身價大漲,資歷雄厚了呢?我爸爸曾告訴我,當你遇到好事,開心事時,可以得意,但是不要忘形.因為人要知道自己是誰,你的本分在哪裡.

2,一直擁有,且比別人的要多

其實這一點主要就是講那些在某些方面有天分的人.上帝創造我們,同時也賦予了我們天分.有的唱歌好,有的彈琴好,有的會做生意,有的會做文章.各有各的優點,各有各的長處.可是,就是這些優點和長處卻也容易讓我們驕傲.好一些的會在除了自己領域之外很謙卑,虛心求人指教.壞一些的呢,則是在任何領域都會擺出一幅驕傲的樣子來.前一種人,平時都和和氣氣,看起來非常謙卑,但是一旦到了自己的領域之後,則不容其他人置喙,因為自己是權威,並且在其他方面已經做出了妥協,這裡就沒有必要再和你客氣了.

但是神就是要擊打驕傲的人,賜福給謙卑的人.你越是在哪裡自信,在哪裡驕傲,神就是要在哪裡破碎你.雅各自認為心計很多,騙了長子名分,騙了長子的祝福,但是萬萬想不到卻被自己的舅舅所騙.記得曾有一位姊妹,自己做過這樣的見證.她說她唱歌很有天分,大家都經常誇她唱歌好.在一次聚會結束後,她提議每人唱一首最喜歡的讚美獻給神,因為她這次聚會很感動.大家一聽覺得是個好建議,所以都同意了.但是她又說要其他人先唱,她最後唱.這看起來是一種謙卑的表現,但是實際上卻暗藏了高傲在裡面,她想和別人有個強烈對比,好讓人誇讚她唱歌好.但是那次,除了她以外,所有人唱得都非常的感動,雖然聲音不是最美的,但是感情是最真摯的.而她,忘詞,破音,跑調,全都找上門來了.唱完以後,她覺得自己無地自容,因為上帝正是在她最驕傲的地方管教了她.正如經上所說:“驕傲來,羞恥也來.”(箴言11:2).

上帝給我們恩賜,是讓我們用來服侍他的,而我們千萬不要因為神給我們的恩賜,反而阻擋了我們與神之間的關係.不怕神不用,就怕不中用；不怕不中用,就怕神不用.

3,我們以為有,但實際上卻沒有.

曾有一位二十世紀的神學家叫尼布林,他有一本書叫《人的本性與命運》,在書中他闡述了他對罪的理解.我覺得可以給我們一個思路來思考驕傲的問題.首先從人本身來講,人是按照神的形象和樣式造的,那麼人在精神方面是有超越性的,但是人又受於自己身體的限制,是有有限性的.正是這種超越與有限,才導致了人的一種焦慮,而焦慮是罪的預備條件,最終會讓人產生罪.他反對那種由絕對真理而產生的那種排他性,認為這是最大的不寬容.

我們先不討論他的罪觀以及他認為的真理是什麼,我們倒是可以從他的態度來入手來看驕傲的問題.當我們用有先的身體和思維想要去把握無限的上帝和真理的時候,往往忘記了自己的有限性,反倒認為我們可以抓住真理,並且我抓住的真理才是真理,其他人的卻不是.在這樣一種絕對的真理下,人們很容易出現判斷以及審判.其實,我們可以堅持真理,我們可以為真理去犧牲一切.但是前提是我們做這麼大的犧牲前,是否有給自己留下一個看清真理的機會.

所以在教會裡,很多人覺得我理解的和我的經歷是相稱的,所以我看到的就是真理,那麼其他人說什麼都沒有用,別管你是牧師還是傳教士,別管你多麼有權威.

當我們用真理當做自己的武器,並且以好意去幫助別人斬斷罪惡時,這是最危險的時候.因為你以為你有,但是實際上你可能沒有.而你以高高在上的姿態去判斷別人,不一定是幫助了他,反而是傷害了他.上一次崇基學院邀請了甘小二導演來到中大進行演講,之後他又在神學院與我們一同觀看了他的一部電影.電影裡描述了一個女傳道,她和一個和尚有了孩子,教會不允許他們結婚,並且說她要走回正路.可是肚子裡的是她的孩子,身邊是她愛的人.教會只是用一個絕對真理式的指責告訴她不許怎麼樣,但是卻沒有人真正的聽她的訴說,聽她真實的想法和故事.這種公義大於慈愛的背後往往就是那種看似擁有卻沒有擁有真理的驕傲.

4,小結:

我們看到,驕傲會讓人得意忘形,讓本來沒有水準的人裝作有水準,正如那句老話,一瓶子不滿,半瓶子逛蕩.

驕傲會在我們最得意的地方遠離神.讓我們沒有進步的空間.

驕傲還會讓我們失去了真正看清真理的機會.

相反,我們看到耶穌,即使自己擁有一切,卻沒有驕傲,反而要謙卑順服.這真值得我們好好去反思.

三,如何謙卑

我們看到了都有什麼情況會驕傲,下面我們看看要如何來謙卑.

1,反倒虛己,取人於奴僕的樣式

剛才,其蘭姊妹用小孩子的聲音,小孩子的口氣,詞語,思維在和小朋友交流,耶穌基督也是一樣,用人的方式來讓我們認識到誰是上帝,我們又如何通過上帝的救恩得到拯救.

這種謙卑是很不容易的,他需要你放棄你擁有的一切,來迎合在你看來不如你的人.我的侄子來我家玩兒的時候,他拉著我讓我陪他玩兒捉迷藏,家裡不大,他藏哪兒我一眼就能看到,但是我卻要裝作沒看見,並且要裝作很費力的找他,如果我一下子就找到他,會讓他感覺很有挫敗感.如果沒有耐心的話,這樣的遷就都不會長久,何況耶穌基督降世為人三十多年呢?

曾有一位傳教士叫伯格裡,他創造了苗文,並且帶領大部分苗族歸向了基督.他放棄英國的好日子,反而來到最貧困的苗族,和他們吃的是一樣的,住的是一樣的,穿的是一樣的.最後又死在了異國他鄉.只有這種放下才能夠讓人感受到什麼才是真正的謙卑,才能感受到耶穌基督謙卑順服的力量.

2,看別人比自己強(2:3)

我們除了倒空自己,還要看別人比我們強.這一點真的很難做到.尤其是在自己做的真的很好的時候,還有是在別人真的比自己強的時候.

三國中,關雲長成名很早,他溫酒斬華雄,過五關斬六將,聲名遠播.敵人聽見其名無不聞風喪膽.也正是他有這樣的名聲,所以他越來越自信,以至於變成了自負和驕傲,不把任何人放在眼裡,最後白走麥城.

而三國中另外一個人,張飛卻很有特點.他魯莽時什麼事兒都能做出來,哥哥的話什麼都忘在了腦後,可是聰明的時候呢卻可以騙過所與人.在劉岱和王忠作為先鋒向駐守在徐州的劉備挑釁的時候,張飛就用了苦肉計生擒了劉岱,以至於讓劉備刮目相看.

所以保羅說看別人比自己強是有道理的,不管什麼人,他總有有優點的一面,總有比你強的一面.尤其是在教會裡,因為神是可以讓人改變的,你看曾經多麼大的罪人,如今都可以在教會裡為主事奉.而且他曾經失敗的經歷往往會轉變成為他今天事奉的財富.

我們也不要看不得別人比我們好,覺得別人好了,自己的優越感沒有了.當你真被別人超過去的時候,就要看看,他到底哪方面真的比自己的強,然後好學習以便今後的進步.

3,謙卑順服,以至於死,且死在十字架上

是不是說我們放下一切,就足夠謙卑了呢?是不是我們放下自己的面子肯去聽別人的建議就是完全謙卑了呢?我們再看經文,耶穌基督不但有人的樣式,而且最後是死了,且死在十字架上.死在十字架上要表達的是一種侮辱性的死亡.耶穌不光生下來卑微,在傳道時卑微,就是在死的時候也是以一種最帶有侮辱性方式死去的.這帶給我們什麼資訊呢?就是我們不光要謙卑地放下一切,而且還要有必“死”的信念,就是那種謙卑到底,以至於死.

有些人聽得進一時的建議,願意謙卑一時,但是一輩子聽建議,一輩子謙謙卑卑則受不了.因為會覺得,我已經改的很好了,我已經夠謙卑了.難道我就不能放縱一下嗎?

唐太宗李世民是少有的明君,這得力於他有一個直言進諫的宰相魏征.魏征的耿直是出了名的,有時也是讓唐太宗最頭疼的,在最激動的時候恨不得殺了魏征.但是冷靜下來後,他還是繼續聽取魏征的建議.當魏征死後,太宗說了這樣一段話:“夫以銅為鏡,可以正衣冠；以古為鏡,可以知興替；以人為鏡,可以知得失.我常保此三鏡,以防己過.今魏征殂逝,遂亡一鏡矣.

基督徒也要這樣,謙卑順服不是一時的功課,不是人前的功課,而是一生一世的功課.要帶有一種甘願完全治死老我,完全放下身段的謙卑,而不是假冒為善的裝裝樣子.不觸碰到底線的時候怎樣都無所謂,一旦觸碰到了底線就會毫不讓步.這個過程是痛的,包括我自己的經歷也都在內,但是這卻不是因為我們的痛就可以停下來的.

4,要有一樣的意念,愛心,心思.(2:2)

但是這些就夠了嗎?讓我們想想,是什麼能夠讓耶穌完全拋棄天上的一切來到人世間,甘願受死呢?正式因為他對我們的愛,他希望我們可以通過救恩,重新與上帝和好.若不是有這樣一個信念,耶穌是無法做出這麼大的犧牲,勇敢的邁向各各他的.看到電影中,耶穌的手被釘子一下一下的釘,每一下似乎真的釘在了我的心裡.

所以,當我們回到2:2節,我們看到,保羅希望眾人要有一樣的一年,愛心和心思.這是什麼意思呢?這就是合一,這也是我們為什麼謙卑順服的目的.如果沒有了這個信念和盼望,如果沒有了這個最終極的目的,我們是無法保持謙卑的心,即使有也可能是為了自己的利益,而不是整個教會.

但是很奇怪的就是,耶穌在十字架上,用生命來換取與神和好,在主裡肢體的合一,可是教會的歷史和現在卻經常上演著我們要拼了命的要與其他人分開,劃清界限.而這種情況是保羅,是耶穌都不想看到的.所以保羅強調要謙卑順服,看別人比自己強,最後才可以合一.

四,謙卑和榮耀

聖經非常奇妙,因為在其中,很多的記載都告訴我們一個道理,就是上帝做事與人不一樣.上帝會讓卑微的升高,讓痛苦的變喜樂.上帝總會在危難中,給人一種希望,總是讓整個事情會有一個巨大的翻轉.

耶穌是死了,但是正是因為他死了,所以他才復活,所以他才被升高,得到萬名之上的名,所以天上地下,以及地底下的都來跪拜稱頌他為主.

這樣來看,我們的謙卑不是自賤,我們的謙卑不是自卑,而是真的帶著痛悔的心,帶著愛心,帶著合一的心,帶著為了榮耀上帝的心才謙卑的,是為了在地上要彰顯主名,榮耀主名的謙卑.

結語

回到我們今天的題目,金色的麥穗是下垂的.我要特別感謝我們的伍牧師,因為這個題目是他給我題的建議,我覺得非常富有想像力以及詩情畫意.不但如此,因為這個題目還讓我想起了我奶奶與我之間的一段對話.

在我小的時候,有一次她來我家,跟我說她聽了一場關於謙卑的道,非常受益,那時,她就把其中聽到的一個比喻跟我講了一下.她說,你看麥子,它在未成熟的時候是直立著的,但是成熟了之後,就低下了頭.這說明,真正有水準的人,是非常謙卑的,因為他的學識,他的知識讓他懂得自己的不足.反而沒有內涵的人,成天才會趾高氣昂的不可一世.而今天,我想說,作為基督徒,我希望我們每個人都可以做那成熟的麥子,與之前的比喻有小小不同的是,這麥子裡的果實不是我們的學識,不是我們的知識和見識,也不是我們的經驗,而是我們清楚地知道我們是卑微不配的瓦器,是耶穌用重價贖回的失喪的靈魂.因為當我們認識到這一點,我們的才會真正的倒空謙卑自己,我們的上帝才會將更多的恩典,智慧和力量放入我們這群瓦器中,從而真正的成為沉甸甸的,成熟的麥子.好在耶穌再來的時候,可以被他收入庫中,而非像糠秕,或是被風吹散,或是丟在永火裡.

謙卑是一種美德,更是基督徒一種內在的,可以使我們與主耶穌基督相連的根本的需求!

願神祝福他自己的話語.



\newpage

\section{我必不至死仍要存活}
\label{sec:197}
\textbf{伍渭文牧師 - 我必不至死,仍要存活}
\newline
\newline
\hyperref[sec:196]{< < < PREV SERMON < < <}
~
\hyperref[sec:toc]{[返目錄]}
~
\hyperref[sec:198]{> > > NEXT SERMON > > >}


講題:我必不至死,仍要存活　I Shall Not Die, But I Shall Live

經文:詩篇118

講員:伍渭文牧師

場合:香港中文大學崇基學院禮拜堂

日期:2013年3月31日

「上主,給我們恩典,去寧靜地接受那不能改變的事物；給我們勇氣,去改變那應該改變的事物；給我們智慧,去識辨兩者的分別.」相信大家都聽過這篇求賜寧靜禱文(The Prayer of Serenity) ,禱文的作者就是著名美國神學家萊因霍爾德.尼布爾(Reinhold Neibuhr) .尼布爾說過,在復活節和聖誕節期間,他會選擇去一些高流教會(High Church),注重禮儀和聖樂的教會,少講道的教會-甚至沒有講道的教會,參加崇拜.作為牧師,我頗了解尼布爾的意思,因為在聖誕節或復活節,講道是不可能的使命(impossible task),基督的降生與基督的復活,是敬虔的奧秘,使我們默然無語.其實尼布爾不是不重視講道,他祇是同情要在復活節講道的講員.

基督已勝過死亡,讓我們來崇拜歌頌上主的作為,讓我們的心靈被提到天上.但復活的信息,對人世間的苦痛,有何適切作用呢?

當我預備今天復活節的信息時,我的思想離不開最近香港一連串自殺事件.其中有年青俊朗的醫生因患愛滋病輕生,一位專業的中年婦人,謀殺離異再復合的愛侶後,畏罪墮樓.他們都在生命絢爛,風華正茂時,突然離開.人選擇輕生,原因很多,其中不少是感到絕望,因為過去鑄成的大錯,心中愧疚自責,沒有顏面面對親友同事.亦有可能感到給人欺負愚弄,所以要報復,到弄出人命時,又後悔自責,不敢面對後果,無力面對控訴.仇恨好,自責好,控訴好,就像死亡,不斷侵蝕人的生命,使人感到沮喪,絕望,使人看不到出路,看不到拯救,但我們今天讀的詩篇給我們盼望,拯救來自上帝,我不至於死,仍要存活.

今天所讀的詩篇118篇,是慶典節期的感恩詩:「你們要稱謝耶和華,因他本為善,他的慈愛永遠長存!願以色列說:他的慈愛永遠長存.願亞倫的家說:他的慈愛長存.願敬畏耶和華的說:他的慈愛長存.」(vv. 1, 2, 3, 29),詩歌的開始和結束都說:你們要稱謝耶和華.

詩118是一組稱為Hallel-埃及讚美-詩113-詩118的詩章最後一首詩篇,描述上主大能的拯救；這也是在重要節期,來自各地的香客,由聖殿的詩班帶領,以啟應方式頌唱的進場詩歌:「給我敞開義門,我要進去稱謝耶和華.這是耶和華的門,義人要進去.」(詩118: 19, 20)

詩歌最初以第一身出現,是王帶領百姓作感恩祈禱,感謝上主保守國家勝過仇敵；因為用第一身,內容也可以作為個人的感恩詩歌.到了v.22節,提到被棄的石頭,在強大的入侵異邦眼中,以色列國無足輕重,像被人遺棄的石頭,才以第一身眾數出現:「匠人所棄的石頭,已成了房角的頭塊石頭.這是耶和華所作的,在我們眼中看為希奇.這是耶和華所定的日子,我們在其中要高興歡喜.」(vv. 22, 23, 24) 他們不單歌頌上帝的大能拯救,更獻上祭牲為感恩祭:「理當用繩索把祭牲拴住,牽到壇那裡.」(v.27) 耶穌設立聖餐後,聖經記載耶穌和門徒也唱了這詩:「他們唱了詩,就出來,往橄欖山去.」(可14: 26)

詩篇118篇說到詩人面對生命的巨大壓力,眾多敵人欲置他於死地,生命危在旦夕.「他們環繞我,圍困我……他們如同蜂子圍繞我,好像燒荊棘的火.」(詩118:11, 12) 但詩人尋求上主的拯救,上主答允拯救:「耶和華的右手高舉,耶和華的右手施展大能!我必不至死,仍要存活,並要傳揚耶和華的作為.」(詩118: ) 復活節的信息不單關於將來死後的復活,也關乎現今的拯救-拯救我們脫離各種死亡的力量.

1.  復活帶給我們平安,免去控訴和定罪,使我們不再懼怕.

「我在急難中求告耶和華,他就應允我,把我安置在寬闊之地.有耶和華幫助我,我必不懼怕.人能把我怎麼樣呢?」(詩118: 5, 6)

耶穌復活後第一次對十一位使徒顯現,他沒有說:當我被釘在十架時,你們在那裡?耶穌說:「願你們平安!」(路24: 36) 不用懼怕.耶穌沒有控訴他們,也沒有定他們的罪,他把我們安置在寬闊之地.相反,聖經中記述撒旦的工作是控訴,撒旦誘惑人犯罪,犯罪之後他又控訴人.「那在我們神面前,晝夜控訴我們弟兄的,己經被摔下去了.」(啟12: 10) 基督的復活,使我們不被定罪,不再被控訴.

逆境可怕,但懼怕逆境的懼怕更加厲害.沙士疫症十周年,回顧過去,我們發現懼怕沙士比沙士本身更可怕.教過書的人都曉得,教書之前我們很多焦慮,焦慮自己對所教的內容不熟掌握不好,焦慮自己在課堂表達能力不夠,不知學生是誰,對自己友善不友善,很多焦慮,這些焦慮令我們神經緊張,甚至睡得不好,但一開了學,心情就安定下好,焦慮退減了,壓力減少了,發現以前的顧慮是多餘的,杞人憂天.

詩人經過祈禱,上帝立刻賜他平靜,使他心靈寬闊,不再懼怕,有一種平安.復活沒有除去死亡,但除去對死亡力量的恐懼,復活除去控訴和定罪,復活安置我們在寬闊之地.基督的復活,打開了密封的世界,封閉墳墓的石頭己被滾開了,死亡的氣味已被驅散.

今天我不知道你面對甚麼急難,也不知道你有何懼怕,或被甚麼事情控訴,感到自責,但主說:願你們平安!我們也許會為自己過去的愚昧,放縱付上代價,受到應有的懲罰,但我們有生路,寬闊的路,不至死亡.「耶和華雖然嚴嚴的懲治我,卻未曾將我交於死亡.」(詩118: 18)

2.      復活給我們繼續事奉的力量-我們仍要存活, 並要傳揚耶和華的作為 v. 17

「匠人所棄的石頭,已成了房角的頭塊石頭.」v.22

在列國眼中以色列是被棄的石塊.

這是以色列人的感恩詩,在列邦眼中,以色列國微不足道,當外敵大軍壓境時,以色列人感到懼怕:「他們環繞我,圍困我……他們如同蜂子圍繞我,好像燒荊棘的火.」(詩118:11, 12) 在強大的異國眼中,以色列不過像一塊被棄的石塊,但上主要使他們成為基石,傳揚耶和華的作為.

在猶太人眼中,耶穌是被棄的石塊.

耶穌在馬可福音曾挪用詩118:22,說自己被猶太人拒絕,是一塊被人棄掉的石頭,但將成為重要的房角石:「匠人所棄的石頭已作了房角的頭塊石頭.這是主所作的,在我們眼中看為希奇.」(可12: 10, 11)

詩篇118被稱為彌賽亞詩篇(Messianic Psalm) ,預言基督勝過死亡,被棄的石塊,成為房角石 .耶穌復活後在以馬忤斯路上和兩位門徒解釋聖經的預言,基督如何按聖經所說,為我們的罪死了,而且埋藏了,第三天復活了；耶穌對十一位門徒顯現時,也對他們解釋聖經的預言:「這就是我從前與你們同在之時所告訴你們的話說:摩西的律法,先知的書和詩篇上所記的,凡指著我的話,都必須應驗.」(路24: 44) 律法,先知,詩篇(詩篇為首的聖寫Holy Writings-文學類) 是舊約的三部,耶穌應該對他們講解彌賽亞的詩篇,包括118篇.

這篇本來以色列國的感恩詩,說到上主拯救他的子民,不至於死；他們被列邦輕視,像一塊被遺棄的石塊,但成為重要的基石.從彌賽亞詩篇的角度詮釋這詩篇.就看到基督被釘十架,但不至於死,仍要存活；基督而且成為房角石,在其上,教會被建立起來.

彼得將成為教會的基石

當耶穌對十一位門徒講論基督這塊被棄的石頭,已成了房角石塊時,我想彼得特別有感觸,因為他的名字就是石頭(磯法) .在該撒利亞腓立比(Caesarea Philippi),彼得是第一個公開認信:你是基督,是永生神的兒子,跟著耶穌說:「你是彼得,我要將我的教會建立在這磐石上.」(太16: 18)

當全球上億的人都透過電視互聯網收看新任教宗方濟(Francis) 就職典禮；由於梵蒂岡是一個國家,所以全球有邦交的國家和地區都有使節觀禮,盛況空前.方濟在就職禮開始時首先進入聖彼得大教堂地下室,在傳說是羅馬首任主教使徒彼得的靈墓跪下祈禱.因為在天主教的傳統,彼得是基督地上的代表,教宗是彼得歷史的延續,各地主教是教宗空間的延伸和代表.

我們更正教(基督教) 則強調內在信仰的傳襲,彼得是第一位對耶穌位公開宣認:「你是基督,永生上主的兒子.」(太16: 16) 耶穌就以此認信為教會的根基:「你是彼得,我要把我的教會建造在這盤石上」(太16:18) 所有認信耶穌為基督的都是教會的一份子,彼得是會頭,我們在不同的時間,重覆了彼得的認信.

三次否認耶穌的基石

不過我們要知道,這位被稱為教會基石的彼得,曾三次否認耶穌:「我不認得他」,「我不是一党的」,「我不曉得你說的是甚麼」.(路22: 54-62) 但耶穌沒有放棄彼得這塊石頭,耶穌轉過身來看彼得,彼得就出去痛哭.耶穌復活後,對彼得三次說:「你愛我比這些更深嗎?」「你餵養我的小羊.」(約21: 15) 彼得三次不認主,耶穌三次把牧養教會的責任交給彼得.彼得跟隨主的心志,沒有因曾經不認主而死去:「我不至死,仍要存活,並要傳揚耶和華的作為.」

今天是復活節,願我們事奉上帝的心,得著挑旺,仍要存活,耶穌復活後並沒有責問門徒:當我釘十架時,你們在那裡?他說:願你們平安.主沒有問彼得:你為何三次不認我?你為何三次不認我?主說:你愛我比這些更深麼?你餵養我的小羊.

作為使徒之首,彼得有沒有意識到他就是上帝教會的基石?我想他會意識到,但他看這領導的地位不是榮耀,是責任.彼得曾到羅馬傳道並在尼祿王(Nero) 大迫害期間殉道,羅馬兵丁要把他釘在十字架,但他說不配和基督一樣,要求倒釘在十字架.

彼得不單自己沒有失去事奉上主的生命,這位被稱為石頭的彼得,更提醒信徒,我們就是上主的活石,我們被召,要建立靈宮.在彼得前書彼得引用詩篇118:22說,基督這塊被人所棄的房角石,已成為靈宮的頭塊石頭:「你們來到主面前,也像活石,被建造成為靈宮.作聖潔的祭司,藉基督奉獻神所悅納的靈祭.」(彼前2: 4, 5)

總結:我不至於死,仍要存活, 並要傳揚他.因為耶穌的復活,死亡的力量不能威脅我們的生命,願你們平安,我們不再被定罪,不再被控訴.我們也未被上主放棄,我們要完成上主交付給我們的工作,我們是活石,建造上主的靈宮.阿們.



\newpage

\section{信主之後}
\label{sec:198}
\textbf{鄭漢文博士 - 信主之後}
\newline
\newline
\hyperref[sec:197]{< < < PREV SERMON < < <}
~
\hyperref[sec:toc]{[返目錄]}
~
\hyperref[sec:199]{> > > NEXT SERMON > > >}


約翰福音 20:19-31
\newline
\begin{longtable}{cl}
\hline
\hline
章節 & 經文 (和合本修訂版)\\
\hline
20:19 & \begin{tabularx}{0.7\textwidth}{X} 那日（就是七日的第一日）晚上，門徒因怕猶太人，所在的地方門都關了。耶穌來，站在當中，對他們說：「願你們平安！」 \end{tabularx} \\ \\ \relax
20:20 & \begin{tabularx}{0.7\textwidth}{X} 說了這話，他把手和肋旁給他們看。門徒一看見主就喜樂了。 \end{tabularx} \\ \\ \relax
20:21 & \begin{tabularx}{0.7\textwidth}{X} 於是耶穌又對他們說：「願你們平安！父怎樣差遣了我，我也照樣差遣你們。」 \end{tabularx} \\ \\ \relax
20:22 & \begin{tabularx}{0.7\textwidth}{X} 說了這話，他向他們吹一口氣，說：「領受聖靈吧！ \end{tabularx} \\ \\ \relax
20:23 & \begin{tabularx}{0.7\textwidth}{X} 你們赦免誰的罪，誰的罪就得赦免；你們不赦免誰的罪，誰的罪就不得赦免。」 \end{tabularx} \\ \\ \relax
20:24 & \begin{tabularx}{0.7\textwidth}{X} 那十二使徒中，有個叫低土馬的多馬，耶穌來的時候，他沒有和他們在一起。 \end{tabularx} \\ \\ \relax
20:25 & \begin{tabularx}{0.7\textwidth}{X} 其他的門徒就對他說：「我們已經看見主了。」多馬卻對他們說：「除非我看見他手上的釘痕，用我的指頭探入那釘痕，用我的手探入他的肋旁，我絕不信。」 \end{tabularx} \\ \\ \relax
20:26 & \begin{tabularx}{0.7\textwidth}{X} 過了八日，門徒又在屋裡，多馬也和他們在一起。門都關了，耶穌來，站在當中，說：「願你們平安！」 \end{tabularx} \\ \\ \relax
20:27 & \begin{tabularx}{0.7\textwidth}{X} 然後他對多馬說：「把你的指頭伸到這裡來，看看我的手；把你的手伸過來，探入我的肋旁。不要疑惑，總要信！」 \end{tabularx} \\ \\ \relax
20:28 & \begin{tabularx}{0.7\textwidth}{X} 多馬回答，對他說：「我的主！我的神！」 \end{tabularx} \\ \\ \relax
20:29 & \begin{tabularx}{0.7\textwidth}{X} 耶穌對他說：「你因為看見了我才信嗎？那沒有看見卻信的有福了。」 \end{tabularx} \\ \\ \relax
20:30 & \begin{tabularx}{0.7\textwidth}{X} 耶穌在他門徒面前另外行了許多神蹟，沒有記錄在這書上。 \end{tabularx} \\ \\ \relax
20:31 & \begin{tabularx}{0.7\textwidth}{X} 但記載這些事是要使你們信耶穌是基督，是神的兒子，並且使你們信他，好因著他的名得生命。 \end{tabularx} \\ \\
[1ex]
\hline
\hline
\end{longtable}
講題:信主之後　After Faith

經文:約翰福音20:19-31

講員:鄭漢文博士

場合:香港中文大學崇基學院禮拜堂

日期:2013年4月7日

1.         信不信主之後都要死

1.1 各位弟兄姊妹平安.今日是復活節後的第二個主日.本來伍渭文牧師安排我上主日復活節講道,但我感到我是不配的,所以我跟伍牧師說:還是在復活節之後吧.今日我想同大家分享的講題是:「信主之後」,英文叫After Faith.

究竟信主是指甚麼意思的主? 是指:相信耶穌基督是復活主?是救主?還是生命的主?

1.2 首先讓我先聽一聽一個故事.從前,有一位媽媽,晚上在兒子睡覺前,會給他講故事.有一晚,這個媽媽講了這樣一個故事:

「從前有一個人,信主之後,……,最後死jor .」

在另一個晚上,這個媽媽又講了另一個故事:「從前有一個人,不相信主之後,……,最後死jor .」

1.3 信主不信主,其實最後都會死.究竟信與不信,又有甚麼分別?就看你所信的主是甚麼意思.

讓我們看一看聖經如何記載,當日的門徒在信主與不信主如何活出分別來.

2.         耶穌基督復活之後門徒見證所信的復活主

2.0 讓我們再讀一讀今天的經課,使徒行傳第五章27節到32節和修本這樣記載:這段經文,值得分解的有四項重點:是彼得和眾使徒這一群信徒在信主之後,當面臨信仰的質疑時,呈現出四種回應,或者是四層回應.

5 : 27 他們把使徒帶來了,就叫他們站在議會前.大祭司問他們,

5 : 28       說:「我們不是嚴嚴地禁止你們,不可奉這名教導人嗎?看,你們倒把你們的道理充滿了耶路撒冷,想要叫這人的血歸到我們身上!」

5 : 29       彼得和眾使徒回答:「我們必須順從上帝,勝於順從人.

5 : 30       你們掛在木頭上殺害的耶穌,我們祖宗的上帝已經使他復活了.

5 : 31       上帝把他高舉在自己的右邊,使他作元帥,作救主,使 以色列人得以悔改,並且罪得赦免.

5 : 32       我們是這些事的見證人；上帝賜給順從的人的聖靈也為這些事作見證.」

2.1 第一個回應在第29節,彼得和眾使徒回答:「我們必須順從上帝,勝於順從人.」

在這裡,我們看到信主的第一層意思是:要順從上帝,而不再順從人.

甚麼叫順從人?人包括:

(1)我自己　和

(2)他人,而他人有兩個層面,包括

(2a)一群人(存在主義者海德格在德語叫Das Man,廣東話即係:「人地」),和

(2b)掌握權力的人,有英譯本繙成human authority.

未信主之前,我們都生來是自我中心的,凡事都順從自己,按己意而行.

長大了,就開始認識到群眾壓力,身邊的人對你有期望,你又渴求他們對你的認同.給新入學的大學生的迎新營(O/Orientation Camp)每年都會出現,新生要依從老鬼設計的活動遊戲,要入群就難以不順從.

又或工作了一段時間的很快就明白工作間掌權者就是話事人,例如:碼頭工人在工資上和在工作條件上要求談判,但話事人總是不肯出來露面.其實黑社會也有話事人,香港政治也有話事人.

一句講完:人在江湖身不由己.

耶穌死後,復活升天.給這群信徒的影響是:他們不再順從己意,不再順從人群,不再順從人間權威,包括當時的宗教領袖如大祭司,而是單要順從上帝.我們信主之後,有沒有這種改變呢?

2.2 第二個回應在第30節,彼得和眾使徒回答:「你們掛在木頭上殺害的耶穌,我們祖宗的上帝已經使他復活了.」信主的第二個意思就是:相信復活的主,雖然死了,第三天從死人中復活.我們每周所一起頌讀的《使徒信經》是這樣認信的.上主日我們是這樣紀念復活的主的.

因他戰勝死亡,我們就不怕死了,我們死了,他也會叫我們復活.

彼得和眾使徒公開指責當時的權貴,認為是你們這群權貴用最殘暴,最羞辱人的方式──就是把耶穌掛在木頭上,殺害耶穌,然而我們歷代所信的上帝已經使他復活了.

彼得和眾使徒他們這樣說話,不怕自己也一樣地被這群權貴用最殘暴,最羞辱人的方式殺害嗎?用廣東話解釋:「一早預jor 啦.」

就好像香港近日出現「讓愛與和平佔領中環求真普選」的運動一樣,如果你問這一萬人,「咩你地唔怕要坐監?」用廣東話回答:「一早預jor 啦.」這就是我們基督徒十架的信仰,背後是一套十架的神學(Theology of the Cross). 其實我也在考慮要不要做這一萬人的其中一個,你們會不會支持我?我需要一百位選民支持我,五十人也不錯,至少二十三位認領我.

我在另一篇快將發表的文章會再詳細討論「讓公民授權萬人佔領中環為求還政於民」.

2.3 第三個回應在第31節,彼得和眾使徒回答:「上帝把他高舉在自己的右邊,使他作元帥,作救主,使 以色列人得以悔改,並且罪得赦免.」

信主的第三個意思就是:相信耶穌基督是救主,戰勝罪的權勢,成為君王,成為元帥,成為領主,讓他的選民有機會從罪的權勢得以解放出來,並且罪得赦免.然而,我們救主所要我們的是:我們先要悔改.

信主之後,就是要常存悔改的心,與主耶穌基督赦罪的救恩相稱.

正如保羅在羅馬書中陳明,上帝的救恩本來臨到以色列人,就是上帝的選民.可惜,有些猶太人錯過了,以致選民擴展到外邦人,包括在座你我的人.

有許多人會指責基督徒做了基督徒後一樣會犯罪,又有一些人認為既然信主後會犯罪就不如不信.他們指的是事實,就是事實上作為罪人我們信主的正正要公開承認我們的罪性,甚至會犯下罪行,然後靠著我們的耶穌基督十架上的代死,作了一次過獻上的贖罪祭.我們罪得赦免,並且得到從罪的權勢得以解放出來,得到一顆常存悔改的心.

2.4 最後第四個回應在第32節,彼得和眾使徒回答:「我們是這些事的見證人；上帝賜給順從的人的聖靈也為這些事作見證.」

信主的第四個意思就是:領受聖靈,為你的信仰作見證.

信主之後就是聖靈同在,同作見證.信主之後的經驗是很奇妙的循環:當信徒越順從上帝,上帝越賜給他們領受聖靈的恩臨；當信徒越領受聖靈的恩臨,他們就越能順從上帝.這個循環越來越可以構成見證.

這是靈性生活的核心部份.將來有機會再分享這部份. 一息間,我會分享我信主之後的相關經歷.

2.5 先返回到,信主之後原來還有第五個意思.

正如我們開首那個媽媽給孩子講故事時講了真相:信主之後以上四層的回應之後,還有第五個意思,最後要死. 你看一看第33節如何記載?

5 : 33       議會的人聽了極其惱怒,想要殺他們.

信主之後,隨時因主名而受死.

信徒冒死見證主,讓我自己在準備這次講道時也反省自己信主之後的人生.以下是我的見證.到今年七月,我信主就已經三十六年了.

我信主之後有四個不同階段的信仰追求

3.1 我是在1976年7月中旬信主的.

當時的我在中四升中五暑期參加了學校的暑假福音營,在營中信主的.地點是香港中文大學的聯合書院.從此我與中大結下一生情緣.

我讀的是男校,當時中四的班內分為兩半,一半是基督徒,一半不是.這一半常常向那一半傳福音.帶頭的基督徒班長挑戰未信主的暑假福音營,帶頭的非基督徒班長接受被傳福音的挑戰,但明言如果未信的在營內一個也沒有人信,以後請不可以再講耶穌.

四日三夜的福音營到了最後一夜,最強勁的佈道晚會完了後,我在思索明早出營前問心事環節一定會被公開地宣召,所以我就把我所聽過的道向自己勸說一遍,然後在子午時份自己引領自己祈禱信主.

到了明天,在問心事時,我班非基督徒班長一眾沒有一個站出來,基督徒班長那一半感到很失敗.誰知,我站了出來,但兩邊都不歡迎我信耶穌.從這刻開始,我知道我一生都會是活在兩造之間的邊緣人.

這是我信主之後的開始.

回首往來蕭瑟處,既多風雨既多晴.我的信仰成長,大概可以分為四個大階段:一是初信,二是尋覓,三是努力,四是復興.

3.2 首先是初信的四年.自信主後的中五開始跟同學返教會,當時家住銅鑼灣,中學在西環,返的教會是在灣仔的循道衛理聯合教會香港堂.我先返周六下午的少年團契,少年團契讓我成長得很好.當時導師教導我信主之後要在三年內把整本聖經六十六卷看完.我回應說:「在福音營裡傳福音的人又沒有講明信主之後要看這麼厚的書.」我其實在中學是不看書的.學校規定每年看六本英文課外書,然後測驗十題是非題或選擇題,我發現沒有看書的也猜中四題,測驗前聽班長講書就考對六題,所以不看也罷.

結果,順從導師教導,我用了一年半的時間,把新約二十七卷看完.新約聖經成為我中學唯一看過的書.導師也很滿意.然後我開始返主日學,返崇拜.信主後一年後的九月我十八年歲,我就去報洗禮預備班,在1977年聖誕節受洗.

初信之後,教會生活有不少恐懼.當時最怕查經,圍一個圈輪流讀,我怕讀錯所以我一早把我那節默唸,點知前面gwo 個讀多了一節讀埋我個節,我就無準備下讀下一節,次次都好緊張.所以,以後都及早準備,把要查的經文預習,和合本聖經其實有許多難讀難解的字,所以就要查經前及早查字典.  如此,教會生活成就了一些自學,學他的良好習慣,也因為要讀聖經而愛好了文字,因此而考進了大學.是香港中文大學,我還刻意選了崇奉基督的崇基學院.

3.3 1980年進入大學,開展了我信主之後的第二階段.我是先讀商科然後兩年後轉讀哲學為主修,成為我們崇基禮拜堂榮譽助理校牧鄧瑞強博士的同班同學.(他當時讀書讀得很好.)在講究利益和講究理性的年代成長,講利與講理二者成為我大學生時期信仰的挑戰.自入大學,我便每學期三五成群開信仰相關的讀書小組,三十年多來從未間斷.

在哲學系內得不少良師栽培,被他們鼓勵,在中大讀完碩士之後往英國牛津留學,成為研究生.我向神祈禱:「若你想我終身以哲學為信仰辯護,就讓我在牛津完成學業.」護教的心願雖然交上,但是護教的路不易走.好景不常,一語成讖,在牛津我見到巴別塔現象後,我也見到我的方向並不在此後,我選擇了離開,轉入教育,轉到倫敦大學教育研究所完成博士於1994年12月.

這一個十四年的學業,在尋覓中,經歷了困苦與復活.

3.4 下一個十四年的事業,在努力中,經歷悔改的心和赦罪的恩.

1995年1月,我返回香港中文大學母校的教育學院任教至今十八年.首十四年算作我信主之後的第三階段.

我自2006年領導價值教育和通識教育科兩個碩士,兼任課程協同總監,到了2012年7月底自動請辭.放下,回到基本點,開往水深處.從此,我專注開發轉化生命的學問.在大學教書不一定會讓我們的信仰得以攀升,反而容易靈性下沉.

大學是一個追求卓越的地方,所以追求卓越的人會有某一程度的狂妄,自負,不可一世.因為能在大學佔一席位的教授總在某一學問領域不單甚有造詣,而且可能在此領域已領風騷十數年.研究有成果,寫文夠多時就會升職,升上高位時就會兼任高職,兼職甚麼主任或咩長物長的領導人,就以為可以領導決策,位高權重,影響深遠.

來,讓我們建一個高塔,讓世人得見我們的榮耀.

誰知這是個象牙塔,是個巴別塔.

我留學到牛津時見到這個巴別塔,首先見到的是「窮智損德」的現象,就是追求智性的卓越至高之時,會漸漸失去德性的心.我回到中大後的十四年也見到這個巴別塔,這次見到的是「窮智損靈」的現象,就是因為智性難得卓越至高之時,會漸漸失去靈性的心.

在香港社會你們也可以見到這個巴別塔,是踏著千萬窮苦大眾的權貴高舉自己的成功.今天你們看看碼頭工人的苦況,真的好像《孤星淚》 Les Misérables裡香港版的悲慘世界,一個一個的貨櫃堆成無良僱主剝削工人的巴別塔.

在香港教會界我也見到這個巴別塔,是我們主所愛的弟兄姊妹高舉成功神學,主張因為神祝福我所以我名成利就的信仰解釋.

3.5 上主恩待我,讓我在這個美麗的校園成長.

崇基禮拜堂取代了我的巴別塔,成為我信仰不同時期的聖壇,我做大學生時在這裡上周會,我和內子畢業後回到這裡見證婚盟的,我做助教時曾在這裡守崇拜,我做老師時在這裡守望崇基精神,現在的我做崇基學院神學院神學生在這裡實習.

漸漸地,我信主之後轉入第四階段.

我是在2007年十月,在我自己教會(循道衛理聯合教會)的「以馬忤斯之旅」經歷信主之後的靈性復興,開始了另一段「學神」的人生旅程.在三天七十二小時的營會裡,再次經歷路加福音第24章門徒在往以馬忤斯路上重遇復活後的主耶穌基督,因為與主同行而得到靈性復興.今天主日,此時此刻,在大嶼山梅窩的衛理園裡,正在有四十多位姊妹在第67旅的旅程路上,得到主與她們同行的靈性更新.

你也渴望靈性復興嗎?

從此,我經歷著「進修中有退修,退修中有進修的靈智兼修」的日子.2008年9月我開始兼讀神學研究碩士,在三年後的2011年年中完成,隨即開始兼讀神道學學士.相信要兼讀四年到2015年才完成.三十多年來我是由學士讀碩士到完成博士,之後讀碩士後再讀學士.信仰好像攀山,有高有低.當中許多大小經歷,將來有機會再和大家分享.

今日講道我是用神學生身分來講道,今天的講道是我在這裡作神學生實習的一次講道.我自2012年1月開始在這裡不帶學分先實習六個月後,然後在暑假後再自去年9月初到今年4月底帶學分先實習,換言之,再見三個主日後,我就會開往水深另一處.

正如助理校牧林豪恩先生,Walter他也是被差來服侍大家的.順從主就是願意說:「我在這裡,請差遣我.」放心,我是會在5月19日回來看你們洗禮的,也會久不久回來探望你們.因為這裡是讓我信仰在不同階段曾經成長的屬靈的家.

3.6 正如我剛才說過我信主之後歷經四大階段:一是初信,二是尋覓,三是努力,四是復興.數算恩典,以上我生命的見證,奇妙地,每一階段大概十四年.

我從小學一年級讀到中學畢業共十四年,是第一個階段,前十年多次聽聞福音後四年得到初信造就.我從大學一年級讀到博士完成共十四年,是第二個階段,信仰轉入以小組研習為本的尋尋覓覓.

我從1995年重回母校任教後的十四年,是第三個階段,信仰重回體制式的教會參與,努力結合哲學,教育與信仰於實踐之中.

我從2008年讀神學啟動了另一個十四年,可能會是我最後的一個十四年,是回應上帝的呼喚,開始讀神學,學習牧養工作.

就是回應今天經文的主題:順從上帝,勝過順從人.

每一階段的轉變,都是在不同教會中周遊列國.在第一階段尾入大學前曾在本地教會到處短跑.

在第二階段尾到英國留學時在不同教會聚會.

如今這次是開始到不同教會實習,以及作講道和其他服侍.

實死,實死,人人必有一死.人之將死,其言也善.雖然不是生離死別,對於曾經相交一年半的崇基禮拜堂的弟兄姊妹,我還是臨別說真話好些.

4.    信主之後的四種信徒或四層回應

4.0 按我觀察,信主之後,其實有四種信徒.今日的經課,也標示出信主之後的四種回應或四層回應.

4.1 第一種是信主之後,無甚分別,在世從人.

他們在世界上所擁有的,依然擁有.信仰所加添的極其量是多了點宗教生活,例如上教堂.並未如使徒行傳第五章第29節,彼得和眾使徒回答:「我們必須順從上帝,勝於順從人.」

第一種信主之後,依然順從人,順從人間的權威,沒有立志順從上帝.在世成功,受人重視,依然是他們的價值.親愛的弟兄姊妹,你又如何?你信主之後,有沒有順從上帝?還是依然順從自己,順從人?

4.2 第二種是信主之後,平安喜樂,等待復活.

正如使徒行傳第五章第30節,彼得和眾使徒回答:「你們掛在木頭上殺害的耶穌,我們祖宗的上帝已經使他復活了.」這等信徒真的相信耶穌被掛在木頭上而為我們受害,受死,但最後復活了.所以,對於今生如何,他們都享受到等待復活的平安與喜樂.所以,正如當年的使徒和門徒,不怕死.

親愛的弟兄姊妹,你又如何?你信主之後,你平安嗎?你有沒有心裡面想爭取民主但又害怕權勢,不敢公開支持合理的訴求?

4.3 第三種是信主之後,高舉救主,以主為主.

正如使徒行傳第五章第31節,彼得和眾使徒回答:「上帝把他高舉在自己的右邊,使他作元帥,作救主,使 以色列人得以悔改,並且罪得赦免.」這等信徒真的悔改了,感受到罪得赦免的救恩. 相信因著耶穌基督戰勝罪的權勢,成為君王,成為元帥,成為領主,作為選民,我們可以從罪的權勢中得以解放出來.

罪的權勢包括政治上和經濟上.政治上不讓每位人民掌有等值的選舉權.經濟上不讓受僱工人得回勞力應得的工價.這些都因為罪的權勢得以橫行.親愛的弟兄姊妹,你又如何?你信主之後,你有沒有感動,讓勞苦工人和貧窮大眾從罪的權勢解放出來?

4.4 第四種是信主之後,領受聖靈,作主見證.

正如使徒行傳第五章第32節,彼得和眾使徒回答:「我們是這些事的見證人；上帝賜給順從的人的聖靈也為這些事作見證.」信主之後經歷到聖靈同在:越順從上帝的信徒上帝越賜給他們領受聖靈的恩臨,越領受聖靈的恩臨的信徒就越能順從上帝.

以此作為靈性生活的核心,有此經歷作為見證. 親愛的弟兄姊妹,你又如何?你信主之後,你感到聖靈同在嗎?你有沒有好像使徒一樣有信心,有勇氣為主作見證?

4.5 以上四種信徒可以出現在同一個生命的不同階段.

你信主之後,是那一種信徒?又或者,信主之後,你如何經歷以上不同階段?你渴慕剛才提及的那種信仰境界?

5. 信主之後

5.1 在構想我這次講道的題目時,我是在模仿一本影響倫理學界很深遠的書,叫做After Virtue《德性之後》,作者是叫Alastair MacIntyre麥金泰爾的哲學家.這部出版於上世紀八十年代的書,是一部西方倫理學的簡史,倫理學是講究美好人生如何得以活出的學問.

這部書影響深遠,復興倫理學的活力.其中一位受其影響的美國的神學家侯活士Stanley Hauerwas因而活化了基督教倫理學的討論.我們的賀志勇博士把他其中一本成名作Resident Aliens中譯成書,名為《僑居者》,快要出版了.

現在若要讀基督教倫理學或敘事倫理學,一定要先回到麥金泰爾After Virtue《德性之後》.這部《德性之後》書中描繪倫理學之重視德性,最早發凡於希臘時期,尤其以亞理士多德的哲學,然而經過中古時期和現代時期的理性高揚,當代的倫理學已經漸漸失落了德性的心.

作者主張返回失落了德性的心,重建德性為本的倫理學.

正如孟子說:學問無他,求其放心. 生命的學問,就是要找回失落了的心.德性之後,提醒我們返回失落了德性的心,重建德性為本的倫理學.信主之後,也同樣地提醒我們返回失落了信主的心,重建以主為本的信仰:順服上主勝過順服人.就是以信德作為基督徒倫理學的核心.

5.2 那裡尋回?

讓我們再回到新約聖經,從福音書到使徒行傳的重新閱讀,我們看到主死後復活所帶給早期教會信徒的啟發.讓我們回到當初被主耶穌基督十架救恩的感動的那一個神聖時刻,慢慢讓信主之後的日子依不同階段漸漸地成長.

讓我們每天每夜,安靜片刻,手握著主耶穌基督的十架,向我們的天父禱告,讓聖靈在你的心內向你說話.聆聽,就是神學.

5.3 讓我們一起低頭祈禱,一起用心回應:「親愛的上主,請你讓我們尋回我們信主之心.

讓我們再回到新約聖經,從福音書到使徒行傳的重新閱讀,教我們看到主死後復活所帶給早期教會信徒的啟發.今天同樣給我們信主之後有所啟發.

讓我們回到當初被主耶穌基督十架救恩的感動的那一個神聖時刻,慢慢讓信主之後的日子依不同階段漸漸地成長,求你與我們初信的同行,在尋覓中得見你,監察我們的努力,教我們靈性復興.今天同樣給我們信主之後得享平安,常存悔改的心,兼且有勇氣面對世上罪的權勢,尋求眾生困苦的解放.

讓我們每天每夜,都能得享安靜片刻,手握著主耶穌基督的十架,向我們的天父禱告,讓聖靈在我們的心內向我們說話.

讓我們聆聽,教我們學習聆聽的神學,教我們順服上主你,勝過順服人.

奉主耶穌基督的名祈禱,阿們.」

但願榮耀都歸給聖父,聖子,聖靈,永在的三一神.阿們.



\newpage

\section{以愛相繫}
\label{sec:199}
\textbf{陳耀鵬牧師 - 以愛相繫}
\newline
\newline
\hyperref[sec:198]{< < < PREV SERMON < < <}
~
\hyperref[sec:toc]{[返目錄]}
~
\hyperref[sec:200]{> > > NEXT SERMON > > >}


約翰福音 21:1-19
\newline
\begin{longtable}{cl}
\hline
\hline
章節 & 經文 (和合本修訂版)\\
\hline
21:1 & \begin{tabularx}{0.7\textwidth}{X} 這些事以後，耶穌在提比哩亞海邊又向門徒顯現。他怎樣顯現記在下面。 \end{tabularx} \\ \\ \relax
21:2 & \begin{tabularx}{0.7\textwidth}{X} 西門‧彼得、叫低土馬的多馬、加利利的迦拿人拿但業、西庇太的兩個兒子，和另外兩個門徒，都在一起。 \end{tabularx} \\ \\ \relax
21:3 & \begin{tabularx}{0.7\textwidth}{X} 西門‧彼得對他們說：「我打魚去。」他們對他說：「我們也和你一起去。」他們就出去，上了船；那一夜並沒有打著甚麼。 \end{tabularx} \\ \\ \relax
21:4 & \begin{tabularx}{0.7\textwidth}{X} 天剛亮的時候，耶穌站在岸上，門徒卻不知道他是耶穌。 \end{tabularx} \\ \\ \relax
21:5 & \begin{tabularx}{0.7\textwidth}{X} 耶穌就對他們說：「孩子們！你們有吃的沒有？」他們回答他：「沒有。」 \end{tabularx} \\ \\ \relax
21:6 & \begin{tabularx}{0.7\textwidth}{X} 耶穌對他們說：「你們把網撒在船的右邊，就會得到。」於是他們撒下網去，竟拉不上來了，因為魚很多。 \end{tabularx} \\ \\ \relax
21:7 & \begin{tabularx}{0.7\textwidth}{X} 耶穌所愛的那門徒對彼得說：「是主！」那時西門‧彼得赤著身子，一聽見是主，就束上外衣，跳進海裡。 \end{tabularx} \\ \\ \relax
21:8 & \begin{tabularx}{0.7\textwidth}{X} 其餘的門徒因離岸不遠，約有二百肘，就坐著小船把那網魚拉過來。 \end{tabularx} \\ \\ \relax
21:9 & \begin{tabularx}{0.7\textwidth}{X} 他們上了岸，看見那裡有炭火，上面有魚和餅。 \end{tabularx} \\ \\ \relax
21:10 & \begin{tabularx}{0.7\textwidth}{X} 耶穌對他們說：「把剛才打的魚拿幾條來。」 \end{tabularx} \\ \\ \relax
21:11 & \begin{tabularx}{0.7\textwidth}{X} 西門‧彼得就上船，把網拉到岸上，網裡滿了大魚，共一百五十三條；雖然魚這樣多，網卻沒有破。 \end{tabularx} \\ \\ \relax
21:12 & \begin{tabularx}{0.7\textwidth}{X} 耶穌對他們說：「你們來吃早飯。」門徒中沒有一個敢問他：「你是誰？」因為他們知道他是主。 \end{tabularx} \\ \\ \relax
21:13 & \begin{tabularx}{0.7\textwidth}{X} 耶穌走過來，拿餅給他們，也照樣拿魚給他們。 \end{tabularx} \\ \\ \relax
21:14 & \begin{tabularx}{0.7\textwidth}{X} 耶穌從死人中復活後向門徒顯現，這是第三次。 \end{tabularx} \\ \\ \relax
21:15 & \begin{tabularx}{0.7\textwidth}{X} 他們吃完了早飯，耶穌對西門‧彼得說：「約翰的兒子西門，你愛我比這些更深嗎？」彼得對他說：「主啊，是的，你知道我愛你。」耶穌說：「你餵養我的小羊。」 \end{tabularx} \\ \\ \relax
21:16 & \begin{tabularx}{0.7\textwidth}{X} 耶穌第二次又對他說：「約翰的兒子西門，你愛我嗎？」彼得對他說：「主啊，是的，你知道我愛你。」耶穌說：「你牧養我的羊。」 \end{tabularx} \\ \\ \relax
21:17 & \begin{tabularx}{0.7\textwidth}{X} 耶穌第三次對他說：「約翰的兒子西門，你愛我嗎？」彼得因為耶穌第三次對他說「你愛我嗎」，就憂愁，對耶穌說：「主啊，你無所不知，你知道我愛你。」耶穌說：「你餵養我的羊。 \end{tabularx} \\ \\ \relax
21:18 & \begin{tabularx}{0.7\textwidth}{X} 我實實在在地告訴你，你年輕的時候，自己束上帶子，隨意往來；但年老的時候，你要伸出手來，別人要把你束上，帶你到不願意去的地方。」 \end{tabularx} \\ \\ \relax
21:19 & \begin{tabularx}{0.7\textwidth}{X} 耶穌說這話，是指彼得會怎樣死來榮耀神。說了這話，耶穌對他說：「你跟從我吧！」 \end{tabularx} \\ \\
[1ex]
\hline
\hline
\end{longtable}
講題:以愛相繫　Bonded in Love

講員:陳耀鵬牧師

經文:約翰福音21:1-19

場合:香港中文大學崇基學院禮拜堂

日期:2013年4月14日

剛才我們聽了不同的經訓,其中一段記載在約翰福音21章1-19節,在這段經文裡,我們看見主耶穌基督,祂再一次地肯定彼得的牧職,祂知道彼得過往有三次不認祂,或許是這個原因,祂三次問彼得:「你愛我比這些更深麼?」彼得知道耶穌明白他的心意,他也有點擔心,因為他過往曾有三次不認主,所以他回答:「主阿,是的．你知道我愛你」.然後主耶穌就肯定他,叫他牧養祂的羊.

在這段經文裡,我們可以很簡單的看見在教會裡的牧養,甚至在教會裡不單單是牧者與會眾,甚至乎是會眾之間的關係是建基於愛的上面.我們怎樣活出這種愛呢?我們一群信徒回到教會裡,我們知道耶穌提醒他的門徒,也同樣地來提醒我們,我們如果彼此相愛,眾人就可以知道我們是耶穌基督的門徒.我們如何活出這種愛心呢?在腓立比書第二章1-8節,保羅在經文告訴我們,有一個很清楚的指引:

1 所以在基督裡若有甚麼勸勉,愛心有甚麼安慰,聖靈有甚麼交通,心中有甚麼慈悲憐憫, 2 你們就要意念相同,愛心相同,有一樣的心思,有一樣的意念,使我的喜樂可以滿足. 3 凡事不可結黨,不可貪圖虛浮的榮耀．只要存心謙卑,各人看別人比自己強. 4 各人不要單顧自己的事,也要顧別人的事. 5 你們當以基督耶穌的心為心. 6 他本有　神的形像,不以自己與　神同等為強奪的． 7 反倒虛己,取了奴僕的形像,成為人的樣式． 8 既有人的樣子,就自己卑微,存心順服,以至於死,且死在十字架上.

這段經文的中心,是保羅提醒腓立比教會的弟兄姊妹,需要以主耶穌基督的心為心(腓2:5),而腓立比書二章 6-8 告訴我們,究竟主耶穌基督的心是怎樣的一個心呢?首先在第6節那裡,很清楚地把主耶穌基督願意捨己的生命表達出來,不單如此,他還告訴我們主耶穌基督的心是願意進入人的生命裡頭的心──「既有人的樣子」.最後,主耶穌基督的生命更加是願意順服父神的生命.

因為在第8節:「就自己卑微,存心順服,以至於死,且死在十字架上」.祂的順服是卑微,徹底,甚至是一種羞愧的順服.因為大家都知道,當時候被釘十字架的罪犯,不像今天福音電影裡所描繪的,乃是全身一絲不掛地掛在木頭上面,所以主耶穌基督的生命是一個謙卑,順服,願意進入人的生命,這是我們該學效的.

在第3-4節裡,保羅列舉了五樣且重要,實際的提醒.當我們與人相處時,我們該如何「處事」呢?──「凡事不可結黨」；怎樣對待「名聲」呢?──「不可貪圖虛浮的榮耀」；怎樣「存心」呢?──「只要存心謙卑」；怎樣「視己」呢?──「各人看別人比自己強」；怎樣「待人」呢?──「各人不要單顧自己的事,也要顧別人的事」.

但在腓立比書第二章1-2節,更加清楚地列舉我們在教會裡面該如何切實地相愛而有的基礎,表現,後果:腓立比書是保羅在監獄裡寫給腓立比教會的弟兄姊妹們.腓立比是當時羅馬政府在馬其頓省的一個殖民地,腓立比教會也是保羅在他第二次傳道旅程裡所建立的教會,這教會所建立的時間不長,但當保羅寫信給他們的時候就稱讚他們:「從頭一天直到如今,你們是同心合意的興旺福音.」(腓1:5)換句話說,這教會是一間傳福音的教會.不單如此,這教會也是一所差傳的教會,因為他們用經濟來支持保羅的傳道旅程.在腓立比書4章16節,保羅感激他們,因為當他在帖撒羅尼迦開展工作時,腓立比書教會的弟兄姊妹支持這工作,他們不單差遣出去,還差人去到羅馬以協助事工(腓2:25-26).

他們差了以巴弗提去與保羅同工,所以腓立比書教會是在當地傳福音,向外地宣教的教會,所以難怪保羅縱然在監牢裡,但當寫腓立比書的時候,他有16次用「喜樂」,「歡歡喜喜」,「歡喜」等來形容他的心情,當中只有一次,保羅說:「使我的喜樂可以滿足」(腓2:2).在教會有什麼事,使弟兄姊妹做了後,教會傳道人的喜樂可以滿足呢?倘若我們再去推論一下,有什麼事弟兄姊妹做了後,使神心裡的喜樂可以滿足呢?第二章1-2節似乎就是那個答案了.答案就是能夠在教會裡過一個彼此相愛,以愛相繫的生活.這個豈不是也是約翰福音21章那裡,當耶穌基督再一次呼召,肯定使徒彼得的牧職時,來提醒他有關牧養的基礎,教會的基礎是建基於愛的上面.

以愛相繫的基礎是在基督裡面.在基督裡面又表示什麼呢?約翰福音一章14節:「道成了肉身,住在我們中間,充充滿滿的有恩典有真理.」換句話說,當我們的相愛是建基於基督裡面,所以才叫「在基督裡」,那個愛是建基於「恩典裡」.

中國人有一句話說:患難見真情.在患難,貧困,痛苦裡,更加顯出彼此關心的情懷,我們在主耶穌基督的恩典裡,所以我們可以彼此相愛.不單如此,我們彼此相愛的基礎也在主耶穌基督的真理裡面──道成了肉身,充充滿滿的有恩典有真理.

我們在教會裡面,我們講愛心,但這愛心不是沒有原則的愛心,因為教會也是真理的柱石,教會理所當然是相愛的地方,但教會更是一個敬拜這位聖潔,至高,至大,榮耀真神的地方,教會要保守神賜給我們那本聖經的真理.我們要用溫柔的心以愛相繫,也當有勇氣為了那交託給聖徒的真道,在需要的時候而竭力爭辯.像一些異端的勢力,求神幫助我們該如此來回應.另外,在教會裡,我們也該肯定愛的基礎是在基督的真理裡面.

正如作為父母親的,如果孩子要求一些他們不該擁有的東西時,父母會截鐵斬釘的拒絕他們.所以弟兄姊妹,如果教會的領袖很截鐵斬釘的告訴你,你有一些的行為,觀點不合聖經的原則,該悔改,改過,認罪,我們就要明白這愛是在真理中的基礎.

這愛也在基督的愛裡面,因為在英文的翻譯裡:if you have any encouragement from being united with Christ, if any comfort from his love──”from his love” ──是主耶穌基督的愛心.主耶穌基督的愛心是一個很廣泛的愛心,對不同的人,祂都願意施予愛的愛心,基督的愛的對象也是很廣泛,年長的,基督愛；年幼的,祂也愛.小孩子來到主耶穌的面前時,他們的父母很希望耶穌能祝福他們,但耶穌的門徒知道夫子很疲憊,但疲憊的耶穌仍然說,讓小孩子到我這裡來；有錢的人如同撒該一樣,耶穌會去他家坐席；窮人如同那只能獻上兩個小錢的寡婦,耶穌也稱讚她；當官的,耶穌看中他；一個少年官來到耶穌的面前,問耶穌:「良善的夫子,我當作甚麼事,纔可以承受永生.」聖經記載,耶穌就愛上他了；當乞丐的巴底買,耶穌也愛也醫治他；一些在社會上有地位的人尼哥底母,他能與耶穌抵足談心；羅馬的官長是耶穌的朋友；一些在社會上被唾棄的痲瘋病人,耶穌伸手摸他們並醫治；撒瑪利亞婦人也是耶穌傳福音的對象….耶穌不單與那些能與他相交的人來交談,耶穌愛的人很廣泛,求神幫助我們,使我們在這個群體裡,將我們愛的範圍,不斷地延伸.有人說,要一個新人在教會裡能留下來,他起碼要在這個群體裡找到七個有意義的關係,不一定是深交的,但是一些有意義的關係──微笑,問候,代禱支持,傷心的時候得到鼓勵.求神幫助我們使我們成為這七個關係中的一人.以愛相繫的基礎是在耶穌基督的恩典,真理及愛心之上.

接著看腓立比書2章1節,那裡有五個動作或動詞:「所以在基督裡若有甚麼勸勉,愛心有甚麼安慰,聖靈有甚麼交通,心中有甚麼慈悲憐憫」:

勸勉──是勸戒及勉力別人,是對人的；

安慰──如果對己太多的話,就有點阿Q精神；

交通──是人與人之間的交談與連絡；

慈悲──如果對己多於對人的話,就有一種不能自律的危險；

憐憫──對己過多的時候,就有一種顧影自憐,無病呻吟之嫌.

保羅提醒我們,在教會裡一些成熟的表現,乃不看自己的需要,而是強調別人的需要,所以他才說:「不要單顧自己的事,也要顧別人的事」.如果我們再看這五樣東西,似乎我們就能歸納這五樣東西為「施予多於接受」,保羅的語氣不是叫人人如此做,他乃叫你們這班腓立比的信徒(或今天的你和我)如此做.誠然我們不是說我們不能接受勸勉,安慰,慈悲,憐憫,當我們面對一些逆境,失意時,我們該回到教會,在牧者或弟兄姊妹身上,接受勸勉,安慰,慈悲,憐憫,但保羅在這裡所強調的,是我們要去付出勸勉,安慰,慈悲,憐憫.

今天美國總統是民主黨的奧巴馬,在過去幾十年以來,民主黨或共和黨的總統都各有功過,但如果你問一個美國人,在過往的四,五十年裡頭,他們最欣賞的民主黨總統是誰?我相信許多人都會說是約翰‧甘乃dik ,他在一九六三年十一月在他前往德州時,被人搶殺了.那時候,美國出現了一個問題,老一輩的美國人或許也會問這一個問題,如同在二○○一年九月十一日那時候,他們會問:「911的時候,你在那裡?」因此,當美國總統約翰‧甘乃dik 被刺殺的時候,美國人也問:「你在 JFK 約翰‧甘乃dik 被刺殺的時候,你在那裡?」他的死和911事件,是美國的一個分水嶺.他為何會那麼受歡迎呢?他不單單成功破解了古巴危機,他也公開答應,向全球宣告,要在十年之內要把人送上月球,縱然他自己無法親身看見,但卻是成功的.他所講過的話有很多,但其中一句是大家耳熟能詳的,就是當他就職上任時,在演說當中的一句話:「不要問國家可以為你做什麼,乃是要問,你可以為這個國家做什麼?」(ask not what your country can do for you - ask what you can do for your country).

弟兄姊妹當你回到教會的時候,當然我們是來敬拜神,但來到這裡的時候,不要問:神能為你做什麼?更加要問:你可以為神做什麼?強調施予多於接受.不要問教會可以為你做什麼?乃是問,你何以為教會做什麼?當你離開教會回到社區,學校,宿舍的時候,不要問你的室友可以為你做什麼?乃是問,你何以為你的室友做什麼?

求神幫助我們,使我們能夠成為光,鹽,更有影響力,就從我們這個群體開始.

最後,不單止要對人多於對己,不單強調施予多於接受,更加是強調感性多於理性.容許我再唸第二章,是以我個人的版本(非聖經的版本)來誦讀:所以在基督裡若有甚麼講道,愛心有甚麼研究,聖靈有甚麼發現,心中有甚麼理論,思想.

這裡是一個大學的群體,我們強調研究,發現,理論,思想,有沒有錯?誠然沒錯.我們在教會裡頭,也強調講道,是一種理性的交流,對嗎?但保羅似乎並不如此講,他所強講的是一種「情」的交流.

弟兄姊妹當我們回到教會時,除了我們是得到一種單向或理性的教導,我們更加需要一種「情」的交流──愛的交流,感受的體會,在需要時互相扶持／支持.

人生裡有許多起落,但回到教會的時候,是否能夠成為我們喜樂的泉源呢?人生裡有許多起伏,但回到教會的時候,弟兄姊妹可以用神的恩典──過往在痛苦,困苦,病痛,失意,失敗時的恩典──成為喜樂的提醒,泉流來面對人生的起伏,特別是一些波折的日子.

求神幫助我們,叫我們有「以愛相繫」的表現.

當我們有這個「以愛相繫」的基礎,並在主耶穌基督的恩典,真理,愛心裡,也當有以愛相繫的表現,要強調別人多於自己,注重施予賜於接受,針對情感,超越理性,如此,我們就帶來一個後果(腓2:2),而最終的目的,乃使神喜樂的心可以滿足.這裡告訴我們有一個合一的表現,與腓立比書教會的弟兄姊妹講合一是相當重要的,因為腓立比書教會(使徒行傳 16章)成立之時有三種人,他們的背景都非常懸殊的:

呂底亞:猶太人,有錢的女商人(賣紫色布疋的婦人)；

希臘的使女:被鬼附的,他跟隨保羅,多日喊叫,保羅就心中厭煩,把她身上的鬼趕了出來,她就信主了,這使女的主人,就拉保羅和他的同伴西拉去坐牢；

禁卒:在監牢裡,保羅和西拉在晚上禱告唱詩讚美神,就感動了禁卒及他全家一同信主.

有錢的女商人可能是一些富戶；這位希臘的使女是屬於一種勞動階層,甚至讓人鄙視的奴隸；而這個政府工職的禁卒,是屬於中產階級的人──他們的背景都很不一樣,很懸殊.所以保羅說要合一,這是相當重要的,而合一的後果有三個:意念相同,愛心相同,有一樣的心思(腓2:2).

當保羅提到的是「意念相同」,而不是意見相同.意念與意見有何不同?意念是指目標；而意見是達到目標的方法.我們的目標要一樣,但我們達到目標的方法可以不一樣──如同煮粥的方法有許多不同的方法,但目的都是一樣──所以,在教會裡,有誰可以說有一個至終的方法才是最好的呢?教會裡合一是重要的,神往往用一些出人意表的方法,甚至不是你我的方法,來成就祂的事工,但我們的意見可以不一樣,但目標一定要一致──意念相同,意見不同；愛心要相同,熱心可以不同.愛心是講對象,熱心是講程度,愛的對象必須是一樣,但人生裡有許多起跌,在跌的時候,可能我們的熱心就需要一些「喜」的人來扶持我們,使我們更加體會神喜樂的泉源.

或許在我們當中的弟兄姊妹,當你來到這裡聚會,參加查經班等等,吸引你繼續留下來的原因,可能不單是基督教的道理,反而是這裡美麗的環境,很吸引人的講道,師母親切的笑容,禮堂雄偉的建築,甚至是弟兄姊妹們的微笑,容貌等等的原因,但當你信主後,這一切就不該成為你的原因.當然那些可能是一些的誘因使你繼續留在這裡,但最重要的,那吸引你的,是那愛心的對象──耶穌基督.只因為耶穌基督的愛,為我們受死的緣故,所以我們仍然在這個地方.可是,熱心可以不一樣.

最後,我們的心思要相同,我們的心志可以不一樣.如同在教會裡,有許多的服事,像主席,詩班,印程序表的同工等等,我們的工作心志可以不一樣,但我們的心思乃是一樣的,就是要這個地方高舉耶穌基督的十字架,在這地方叫神的榮耀在這裡彰顯,更使我們的愛吸引更多人歸向主.求神幫助我們有這個心.讓我們低頭禱告:

親愛天父:我們謝謝你給我們今天這個時間,可以在這個悠久歷史的禮拜堂內敬拜你,看著十字架,窗外的景色,我們再一次記念耶穌基督為我們死以及復活.謝謝主,讓我們可以在教會裡,雖然在我們這裡有一些人是過客,或許我們畢業後就會離開,但我們在這個屬於你的教會裡,我們可以學習到,無論將來會在哪裡,都可以有一個彼此相愛的生命,為你作光,作鹽,求你讓我們記得腓立比書第二章1-2節的教導:所以在基督裡若有甚麼勸勉,愛心有甚麼安慰,聖靈有甚麼交通,心中有甚麼慈悲憐憫,我們就要意念相同,愛心相同,有一樣的心思,有一樣的意念,使神你的心可以滿足,奉主名求,阿們.



\newpage

\section{凱旋的群眾}
\label{sec:200}
\textbf{鄧瑞強博士 - 凱旋的群眾}
\newline
\newline
\hyperref[sec:199]{< < < PREV SERMON < < <}
~
\hyperref[sec:toc]{[返目錄]}
~
\hyperref[sec:201]{> > > NEXT SERMON > > >}


啟示錄 7:9-17
\newline
\begin{longtable}{cl}
\hline
\hline
章節 & 經文 (和合本修訂版)\\
\hline
7:9 & \begin{tabularx}{0.7\textwidth}{X} 此後，我觀看，看見有許多人，沒有人能計算，是從各邦國、各支派、各民族、各語言來的，站在寶座和羔羊面前，身穿白衣，手拿棕樹枝， \end{tabularx} \\ \\ \relax
7:10 & \begin{tabularx}{0.7\textwidth}{X} 大聲喊著說：「願救恩歸於坐在寶座上我們的神，也歸於羔羊！」 \end{tabularx} \\ \\ \relax
7:11 & \begin{tabularx}{0.7\textwidth}{X} 眾天使都站在寶座和眾長老，以及四個活物的周圍，俯伏在寶座前，敬拜神， \end{tabularx} \\ \\ \relax
7:12 & \begin{tabularx}{0.7\textwidth}{X} 說：「阿們！頌讚、榮耀、智慧、感謝、尊貴、權能、力量都歸於我們的神，直到永永遠遠。阿們！」 \end{tabularx} \\ \\ \relax
7:13 & \begin{tabularx}{0.7\textwidth}{X} 長老中有一位回應我說：「這些穿白衣的是誰？是從哪裡來的？」 \end{tabularx} \\ \\ \relax
7:14 & \begin{tabularx}{0.7\textwidth}{X} 我對他說：「我主啊，你是知道的。」他向我說：「這些人是從大患難中出來的，他們曾用羔羊的血把衣裳洗得潔白。 \end{tabularx} \\ \\ \relax
7:15 & \begin{tabularx}{0.7\textwidth}{X} 所以，他們在神寶座前，晝夜在他殿中事奉他；那坐在寶座上的要用帳幕覆庇他們。 \end{tabularx} \\ \\ \relax
7:16 & \begin{tabularx}{0.7\textwidth}{X} 他們不再飢，不再渴；太陽必不傷害他們，任何炎熱也不傷害他們， \end{tabularx} \\ \\ \relax
7:17 & \begin{tabularx}{0.7\textwidth}{X} 因為寶座中的羔羊必牧養他們，領他們到生命水的泉源；神必擦去他們一切的眼淚。」 \end{tabularx} \\ \\
[1ex]
\hline
\hline
\end{longtable}
講題:凱旋的群眾

經文:啟示錄7章9-17節

講員:鄧瑞強博士

經文:啟示錄7:9-17

場合:香港中文大學崇基學院禮拜堂

日期:2013年4月21日

各位弟兄姊妹,早安.

法國文豪雨果的《悲慘世界》(Victor Hugo: Les Misérables),上年又有新的電影版本.當視死如歸的群眾,唱著震人心弦的詩歌時,通常都是感人至深的.有一首群眾唱的歌,歌名叫＜Do You Hear the People Sing?＞很多人喜歡這首歌.這首歌的歌詞大意是這樣的:

「你聽到群眾的歌聲嗎?他們唱著忿怒的群眾的歌.這是人民的歌聲,他們不再甘心為奴.當你的心與戰鼓和應之時,新的生命便將來臨,就在明天.你會否加入這神聖的抗爭?誰能堅強地與我站在一起?在路障的彼岸,是否有一你渴望見到的世界?快來加入這場戰鬥,這給你自由的權利.…你願否奉獻一切你能奉獻的,以致我們的旗幟能前進?有些人會倒下,有些人會活下去.你會否站出來,把握你的命運?殉道者的血,會滋潤法國的土地.」

(Do you hear the people sing?     Singing a song of angry men?

It is the music of a people     Who will not be slaves again!

When the beating of your heart     Echoes the beating of the drums

There is a life about to start     When tomorrow comes!

Will you join in our crusade?     Who will be strong and stand with me?

Beyond the barricade     Is there a world you long to see?

Then join in the fight     That will give you the right to be free!…..

Will you give all you can give     So that our banner may advance

Some will fall and some will live     Will you stand up and take your chance?

The blood of the martyrs     Will water the meadows of France!)

網上有人說,聽到群眾唱這首歌,便想哭.又有人說,聽到這首歌,便想革命.

若群眾唱出了我們人性心底的渴望,我們也真的很難無動於衷.我們都不想做奴隸,我們都渴望自由.但是,要獲得這些看似理所當然的東西,也須要付出代價.我想起我們的國歌:「起來,不願做奴隸的人們!」但是要能這樣,必須「冒著敵人的炮火,前進!」這是驚心動魄的.

今日的講道經文,是今日的書信經課:啟示錄7:9－17.

這裡同樣提到群眾的歌聲,提到生命的艱難,提到要付出血的代價,提到我們所渴望的世界.這一切,都是驚心動魄的.

我們再讀這段經文,聽聽這些群眾的歌聲.

啟7:9　此後,我觀看,見有許多的人,沒有人能數過來,是從各國,各族,各民,各方來的,站在寶座和羔羊面前,身穿白衣,手拿棕樹枝,

啟7:10　大聲喊著說:「願救恩歸與坐在寶座上我們的神,也歸與羔羊!」

啟7:11　眾天使都站在寶座和眾長老並四活物的周圍,在寶座前,面伏於地,敬拜神,

啟7:12　說:阿們!頌讚,榮耀,智慧,感謝,尊貴,權柄,大力都歸與我們的神,直到永永遠遠.阿們!

啟7:13　長老中有一位問我說:「這些穿白衣的是誰?是從哪裡來的?」

啟7:14　我對他說:「我主,你知道.」他向我說:「這些人是從大患難中出來的,曾用羔羊的血把衣裳洗白淨了.

啟7:15　所以,他們在神寶座前,晝夜在他殿中事奉他.坐寶座的要用帳幕覆庇他們.

啟7:16　他們不再飢,不再渴；日頭和炎熱也必不傷害他們.

啟7:17　因為寶座中的羔羊必牧養他們,領他們到生命水的泉源；神也必擦去他們一切的眼淚.」

啟示錄的文體,屬於「天啟文學」(Apocalyptic Literature),作者「被聖靈感動」(1:10),進入一種特殊的信仰狀態,體悟到人生與歷史的終極真相,他以一種異常圖畫性的方式,向我們展現這真相.圖畫性的表達的好處,是活靈活現,令人印象深刻.我們讀的報刊,很多時都有反映社會現象或政治現實的漫畫,漫畫的構圖本身當然不是寫實的,但是,你有沒有發現,看似不寫實的漫畫,反而更深地展現出事實的底蘊.我看過一幅漫畫,香港的地鐵列車變成一隻食人的怪獸,各人爭相走避,但都難逃劫數.漫畫的標題是「鐵價不義」(「仁義」的「義」).從寫實的角度言,地鐵列車當然不是一隻怪獸；但從更深的角度言,這漫畫是否更真實地畫出交通開支對窮人造成的壓力?漫畫會出現「怪獸」,因這表達方式更能呈現出事情的真相.如此,你便會明白,啟示錄裡出現很多「怪獸」,也是要我們明白事情的真相.

回到今日的講道經文.

啟7:9　此後,我觀看,見有許多的人,沒有人能數過來,是從各國,各族,各民,各方來的,站在寶座和羔羊面前,身穿白衣,手拿棕樹枝.

場面異常浩大.從一種永恆角度觀看的世界裡,有無數人站在神面前,不只是得你一個人.中國宋朝哲人陸九淵講過:「東海有聖人出焉,此心同也,此理同也.西海有聖人出焉,此心同也,此理同也.」古往今來,同一天空下,總有不同的人領悟到同一真理.如孔子所說:「德不孤,必有鄰.」啟示錄鼓舞我們,若你看到羔羊耶穌生命裡展現的神性,若你心甘情願降服這位神,若你為這位神啟示的真理付出代價,你不必感到孤獨.古往今來,同一天空下,有無數人看到你所看見的真理,有無數人敬拜你所敬拜的神,有無數人願意付出你所願意付出的代價.你不是孤獨的,你與這些人,都是聖徒群體的一員.

這群聖徒「站在寶座和羔羊面前」.「寶座」和「羔羊」其實是對立的圖畫,猶如「王子」與「乞丐」一樣對立.「寶座」代表「權力」與「威榮」；「羔羊」代表「受害」與「卑下」.啟示錄用圖畫的方式,表達出「榮耀」與「十架」同在,「復活」與「死亡」聯合.這群聖徒「站在寶座和羔羊面前」,顯然,他們明白這同時強調十字架與復活的基督信仰.像羔羊那樣為僕,會散發出引導人心的領導力.敢於為神的真理而死,才算是永恆地活著.

這群人「身穿白衣,手拿棕樹枝」.按釋經書說,這是當時羅馬社會凱旋而歸的場面.勝利者身穿白衣,手拿代表勝利的棕樹枝.從事實而言,在羅馬社會裡,啟示錄講的這群群眾算是失敗者,像釘十架的耶穌一樣失敗.但從永恆的角度言,這群眾卻是勝利的群眾,像復活的基督一樣光明.我讀過一個報導,講到香港的一位年輕校長,月入八,九萬,卻放下這高薪厚職,開設免費的補習社,為那些中途輟學,有學習障礙,貧窮家庭,少數族裔等的青少年,提供一對一的免費補習服務.在他的感召下,有超過一千人甘願在他的補習社裡當義務教師.他們這個小群體,默默地唱出了生命之歌.在這個以錢來衡量成敗的社會,放下一百萬年薪的人,可算是失敗者.但從永恆的角度看,他是一個勝利的人.

在啟示錄裡,這群眾大聲呼喊:「願救恩歸與坐在寶座上我們的神,也歸與羔羊!」「救恩」可理解為「勝利」,「勝利」是屬於「在寶座上的神」及被殺的「羔羊」耶穌的.在他們看來,「勝利」的標準不由人定,只由神定.人間的英雄,會被浪花淘盡.但有些人,卻能在永恆裡,與神面對面.這些人有份於基督的凱旋.他們活在塵世,卻體會到塵世處處受罪惡污染而人在其中受苦.他們聆聽到神解救世人的呼喚,他們追隨上天的呼喚而行.這些人,擔負起生命的十字架.這些人,活在神面前.

群眾的呼喊引發天界的歌聲.

啟7:11　眾天使都站在寶座和眾長老並四活物的周圍,在寶座前,面伏於地,敬拜神,

啟7:12　說:阿們!頌讚,榮耀,智慧,感謝,尊貴,權柄,大力都歸與我們的神,直到永永遠遠.阿們!

聽說,最澎湃的合唱,要算是馬勒(Gustav Mahler)的第八號交響曲.首演時,合唱團人數是850人,獨唱者8人,管風琴家連同管弦樂團人數是171人,連同指揮,這首演動用了1030位音樂家合力演出.首演結束時,掌聲持續了20分鐘之久.馬勒寫信給一位指揮家Mengelberg,講到這第八號交響曲,馬勒說:「想像一下,整個宇宙開始鳴響.那不再是凡人的聲音,而是星球和太陽的運轉.」可以想像,這交響曲的演出,一定非常震撼.

但很明顯,馬勒沒有聽過啟示錄這裡講的天界合唱團的歌聲.當天界萬物都在頌唱時,馬勒那「千人交響曲」只能是小巫見大巫.

天界的歌聲,是對神的七頌讚美.「頌讚,榮耀,智慧,感謝,尊貴,權柄,大力」,是對神的榮耀和能力的無限讚美.

從終極角度言,勝利是屬於神的.

跟著的經文是一段對話,說明那群眾的身分.

啟7:13　長老中有一位問我說:「這些穿白衣的是誰?是從哪裡來的?」

啟7:14　我對他說:「我主,你知道.」他向我說:「這些人是從大患難中出來的,曾用羔羊的血把衣裳洗白淨了.」

「這些穿白衣的是誰?是從哪裡來的?」

「這些人是從大患難中出來的,曾用羔羊的血把衣裳洗白淨了.」

要成就生命的純潔,沒有捷徑可言.這凱旋的群眾,不是坐直升機升上來的,他們是從大患難中走過來的.我們不是活在完美的世界裡.期望世上沒有罪惡,沒有醜陋,沒有不義,是一種奢望.期望不用付代價,而經濟制度會變好,社會制度會變好,政治制度會變好,同樣是奢望.要以善勝惡,以公義取代不義,以仁愛克服仇恨,就必像羔羊面對豺狼,痛苦與艱難是少不免的了.

經文說,這群眾穿的白衣,是以羔羊的血把衣裳洗白的.這最少有兩個含意:一,這些聖徒不懼以善勝惡的艱難,他們參與了基督的十架行動.他們流出殉道者的血.他們不向強權低頭,因為他們相信神的大能.他們不向罪惡妥協,因為他們愛慕聖潔的光明.他們不會因自保而放棄真理,因為他們知道真理才是生命的泉源.

二,羔羊的血洗淨他們的衣裳,這意味他們不會因自己任何善行而自誇.人能面對神,全賴基督寶血的恩典.行善的人若自誇自己的善,他的善很快便會成為他行惡的理由和動力了.體會到生命全賴神的恩典,他的善就不會成為他的惡了.

這群聖徒活出了生命的神聖,也享受永恆的福樂.

啟7:15　所以,他們在神寶座前,晝夜在他殿中事奉他.坐寶座的要用帳幕覆庇他們.

啟7:16　他們不再飢,不再渴；日頭和炎熱也必不傷害他們.

啟7:17　因為寶座中的羔羊必牧養他們,領他們到生命水的泉源；神也必擦去他們一切的眼淚.」

坐寶座的神,賜給這群聖徒七重福氣.

坐寶座的「要用帳幕覆庇他們」,「他們不再飢」,「不再渴」,「日頭和炎熱也必不傷害他們」,「必牧養他們」,「領他們到生命水的泉源」,「也必擦去他們一切的眼淚」.

「坐寶座的要用帳幕覆庇他們」,神讓這群眾找到安身立命的家園.他們不再有眼淚,因為他們在人生的努力不會白費.他們會看到他們的義行的成果.基督徒在世,就要是接續基督的救贖工程.以服侍去扶助弱者,使世上再無弱者被欺凌.以寬恕讓人有新生,使世人不再以暴易暴.以愛去讓人感到生命的意義,以致人敢於去愛.在永恆裡,基督徒會看到這一切的實現.

總的來說,啟示錄的經文向我們展現終極的景象.誰是王?誰是勝利者?憑什麼得勝?經文對這一切都有所說明.經文提醒我們,生命有艱難,要勇敢面對.信仰的墮落,莫如假信仰之名,使自己在世俗的力量下,軟弱得成為世俗力量的奴隸.有些人的信仰,只是求神叫自己好運.祈禱的內容是:考試時能「貼」中「貼士」,買的股票能勁升而發一橫財,沒有付出努力卻能僥幸成功.這種信仰如果靈驗,只會叫人更軟弱,更不負責任,更走向塵世.啟示錄的經文講的信仰,卻是叫人認清人生的真相.它鼓勵我們,為了神,為了真理,為了正義,付上生命的代價.這才是信仰,這才是能與神面對面的生命.

上個禮拜,《時代雜誌》選出了今年度最有影響力的一百人.新上任的教宗方濟(Pope Francis)也榜上有名.介紹教宗方濟的網頁上寫著,他鼓勵我們:「不要害怕去愛!不要害怕成為溫柔的人!」(Do not be afraid to love! Do not be afraid to be tender!)對這簡單的呼籲,有紅衣主教說:「他說話像耶穌!」有一大主教回應道:「我想,這就是他的職務說明.」

我想,活像耶穌,也就是我們生命的內容說明.

但願:榮耀歸於聖父,聖子,聖靈.阿們.



\newpage

\section{使我成為別人的祝福}
\label{sec:201}
\textbf{許開明牧師 - 使我成為別人的祝福}
\newline
\newline
\hyperref[sec:200]{< < < PREV SERMON < < <}
~
\hyperref[sec:toc]{[返目錄]}
~
\hyperref[sec:202]{> > > NEXT SERMON > > >}


約翰福音 13:31-35
\newline
\begin{longtable}{cl}
\hline
\hline
章節 & 經文 (和合本修訂版)\\
\hline
13:31 & \begin{tabularx}{0.7\textwidth}{X} 猶大出去後，耶穌說：「如今人子得了榮耀，神在人子身上也得了榮耀。 \end{tabularx} \\ \\ \relax
13:32 & \begin{tabularx}{0.7\textwidth}{X} 如果神因人子得了榮耀，神也要因自己榮耀人子，並且要立刻榮耀他。 \end{tabularx} \\ \\ \relax
13:33 & \begin{tabularx}{0.7\textwidth}{X} 孩子們！我與你們同在的時候不多了；你們會找我，但我所去的地方，你們不能去。這話我曾對猶太人說過，現在也照樣對你們說。 \end{tabularx} \\ \\ \relax
13:34 & \begin{tabularx}{0.7\textwidth}{X} 我賜給你們一條新命令，乃是叫你們彼此相愛；我怎樣愛你們，你們也要怎樣彼此相愛。 \end{tabularx} \\ \\ \relax
13:35 & \begin{tabularx}{0.7\textwidth}{X} 你們若彼此相愛，眾人因此就認出你們是我的門徒了。」 \end{tabularx} \\ \\
[1ex]
\hline
\hline
\end{longtable}
講題:願我生命成為別人的祝福　Make Me a Blessing to Others

經文:約翰福音13章31-35節

講員:許開明牧師

場合:香港中文大學崇基學院禮拜堂

日期:2013年4月28日

我們一起來禱告:願僕人口中的言語,弟兄姊妹的心懷意念,都蒙上主的悅納,阿們.

在1986年10月25日,對我來說,這日子是相當難忘的.因為我當天是被按立為牧師.當時候大約有20-30位牧師的手,同時間按手在我頭上,我感到有一個很大的壓力,雖然20-30位的牧師都輕輕的按在我頭上,但加起來的力氣也相當之大的,那時候我就覺得,成為牧師將要面對著壓力,挑戰,但與之同時,我也願意成為一位好牧者,被上帝使用,在按立典禮之後,我收到許許多多的祝賀,其中最深刻的是一位姐妹,她特意製作了一張賀卡,卡上寫著:「願你的生命成為別人的祝福」.這句話成為我自己在按立之後,一個很大的心願.

當我唸到亞伯拉罕在創世記12章 - 13章2節,我在亞伯拉罕的身上看見他的生命是使別人得福的生命,也是我對自己的願望.今天,我們看見世界上存在著兩種的模式:

生命的存在是使他人受傷的──希特勒,他屠殺了六百萬的猶太人.

生命的存在是使他人得福的──亞伯拉罕的生命是受祝福的生命,並且他也要去祝福別人.

我們再來看創世記 12:1-3:

1 耶和華對亞伯蘭說,你要離開本地,本族,父家,往我所要指示你的地去. 2 我必叫你成為大國．我必賜福給你,叫你的名為大,你也要叫別人得福． 3 為你祝福的,我必賜福與他,那咒詛你的,我必咒詛他,地上的萬族都要因你得福.

亞伯蘭能夠得到上帝的賜福,是一個恩典,應許.但他如何蒙福呢?第一,他是因著順服上帝而蒙福──當上主要亞伯蘭離開本地,本族,父家的時候,「亞伯蘭就照著耶和華的吩咐去了」(創12:4)亞伯蘭有順服的心,這並不是簡單的事,因為:

他要離開本地──一個他所熟悉的地方；

離開本族──一個他族人所聚集的地方,一個富足之地；

離開父家──一個他有權力的的地方.

但上帝呼召他要去的地方,是他所不知道的地方,要他放下本地,本族,父家,這是一個相當大的挑戰,他需要冒險,信心.他能蒙福,除了是上帝對他的一個應許之外,但假使亞伯蘭不願意遵從上帝的吩咐的話,這祝福或許不會臨到的.

人類天生就有一種劣根性,反叛性,是不容易去順從的,上帝創造了亞當,夏娃時就說:「園中各樣樹上的果子,你可以隨意喫．只是分別善惡樹上的果子,你不可喫,因為你喫的日子必定死.」(創2:16-17)但結果他們還是吃了那善惡樹上的果子.因此,要順從並不是容易的,而順從上帝才能蒙福,這對基督徒來說,是一個挑戰.我個人相信,順從人不一定蒙福,但順從上帝的心意去做,你必定蒙福的.

從亞伯蘭的經驗裡,當上帝帶領他進入迦南地時(創12:7),他到了示劍,在那裡築壇,目的是感謝上帝的帶領,使他能平安抵達迦南示劍地.築壇也表示他願意親近上主.因此,築壇表達了他對上帝的感恩.亞伯蘭的生命能夠蒙福,因為他看見上帝的帶領及恩典.

在我自己牧會的過程之中,我深受著幾位牧者的影響.在我年青困苦的日子裡,這些牧者給予我許多的關愛及照顧,我感受到他們的生命是祝福著別人,所以我也願意踏上牧者的道路,我也盼望我的生命也是一個蒙福的生命,也能祝福別人.

亞伯蘭的蒙福是因為他看見上帝的祝福,當他領受了上帝的祝福後,他也去祝福別人.上帝帶領亞伯蘭到了迦南地之後,迦南地就遇上了飢荒,困難,在這些困難之中,讓亞伯蘭有另一層的體會,就是在危難時,看見上帝的祝福.聖經說,因著迦南地的飢荒,他們就下到埃及,在那裡解決他們飢荒的問題,但亞伯蘭也面對著另一個困難,就是他不能得罪法老王,可是因著他妻子撒萊的美貌,他擔心法老王會搶奪撒萊而把他殺掉,所以他就認撒萊為他的妹子,好保存性命.在這種困難裡,縱使飢荒的問題得以解決,但亞伯蘭卻陷入另一種的困局(因為撒謊是可以致命的),他滿以為用自己的方法可以解決,結果卻是遇上另一個更大的困難,就是法老王果然看中了撒萊,但上帝因著撒萊的緣故而擊打法老王及他全家,原本亞伯蘭是難逃一死的,他怎麼得以解脫呢?

亞伯蘭在這種的困局中,他是蒙受著另一個更大的祝福!法老王召喚亞伯蘭,一開口就問亞伯蘭:「你這向我作的是甚麼事呢?」(創12:18)接著法王又說:「為甚麼說,他是你的妹子?」(創12:19)本以為法老要向亞伯蘭治罪,但他卻說:「現在你的妻子在這裡,可以帶他走罷」.不單如此,法老「吩咐人將亞伯蘭和他妻子,並他所有的都送走了」也送了亞伯蘭許多的財物.為何會如此呢?

不曉得大家有沒有留意,當上帝祝福亞伯蘭的時候,上帝對亞伯蘭說:「為你祝福的,我必賜福與他,那咒詛你的,我必咒詛他」(創12:3).亞伯蘭作了這種錯事,理當處死,但因著上帝賜下了應許,這應許就是要祝福亞伯蘭,因此上帝介入此事,免去亞伯蘭的死.聖經並沒有提及撒萊的心境或表現,但聖經卻說:「耶和華因亞伯蘭妻子撒萊的緣故,降大災與法老和他的全家.」(創12:17)法老王怎麼會知道亞伯蘭撒謊呢?誰去告訴法老王呢?我相信是撒萊.對撒萊來說,她可能感到她的丈夫不能保護她,還把她送進王宮裡去,能保護她的,也只有天上的上帝.因此,撒萊極有可能大膽地去求見法老王,法老王就明白為何他及他的家族會被打擊,原來就是你的的上帝──耶和華,所以趕緊地把亞伯蘭送走.由於「為你祝福的,我必賜福與他,那咒詛你的,我必咒詛他」,法老就祝福亞伯蘭,送他許多的財物.創世記13:1-3那裡記載,亞伯蘭帶著許多從埃及而得的財物進入迦南地,因此在迦南地,他不用再受到飢荒之苦.

許多時候,我們往往是在困境中蒙上帝的祝福,在絕路中,我們體會到上帝的帶領,幫助,這種的祝福是更大的,上帝賜予的祝福不是說,你要什麼就給什麼的,或許,有時候你所求的,神不賜下；你不想要的,神卻賜下給你!──亞伯蘭就有這種的經歷.

如同我很喜歡唱歌,但在1985年,神把我的聲音取走,就連唱歌,講道也不行,後來神把我的聲音復還給我.經歷這件事之後,我很珍惜神給我的一切,在危難裡經歷到神的祝福,這祝福是又深又大的.

有一個寓言故事說到,一隻蚌對一隻蟹說:「我真是苦呀!我的身體裡面有一顆又大又圓的沙粒,使我很不舒服.」另一隻蚌聽後就說:「我很好,因為我裡裡外外都很舒服,沒有沙粒跑過來折騰我.」蟹就說:「你很舒服,沒有錯,但你知道你朋友所忍受的沙粒,將來會成為一粒的珍珠.」──有時當你經歷苦楚之後,你會得到一粒珍珠,我們必須有這一種的信心.有時候上帝的帶領,不是平平順順,乃是曲曲折折的,但當中卻隱藏著神的祝福.

上帝使亞伯蘭受祝福,也給了他一個使命──當他受到祝福後,他也要叫他周圍的人,以及萬族都要因他得祝福.另外,祝福亞伯蘭的,上帝也必祝福他.而且當亞伯蘭祝福他周遭的人時,上帝也會加倍的祝福亞伯蘭.

我原本的身體狀況不太好,但在我三十多年的牧會當中,神卻賜給我健康的身體,這就是祝福.如同李炳光牧師所言:我牧養群羊,群羊也牧養我.我去祝福,也得到別人的祝福,更得到上帝的祝福.因此,我的使命就是:讓我的生命成為別人的祝福.

什麼是福?從聖經上看,第一,是一個很好的關係,是平安,和諧的關係；第二,是幸福,快樂,喜樂；第三,是一個很好的狀態,是福氣；第四,是一個好的生命,有生命力的,有感染力的.當你祝福他人時,其他人也會再繼續祝福他們周遭的人.當我祝福他人時,上帝也繼續賜福予我,使我的生命可以繼續把這祝福傳遞下去.

盼望各位弟兄姊妹都能有這種心志,讓我們的生命成為別人的祝福.上帝把祝福賜下給我們,我們就把上帝給予的恩典與祝福繼續傳遞出去,因為神給我們的生命是可貴的!每天晚上,我們要懂得數算自己的日子,因為你人生的日子已少了一天,所以要珍惜上帝所給你的,也數算上帝給予的恩典.當你擁有一個體會上帝恩典的心,上帝祝福你生命的心,你就會更願意去祝福你身邊的人.有一位90多歲的老牧者說,每天早晨當我起來,我發現我還能起床,還能去飲茶,我就為此感謝神.你要體會神的恩典,好叫你能祝福你身邊的人.

有記者訪問愛因斯坦說:你是那麼有名的科學家,你認為你是不是一個成功的人?愛因斯坦回答說:我不想作一個成功的人,我想作一個有價值的人,因為成功是人家付予的,你總不能說自己很成功吧!而價值是你能夠給予的.

弟兄姊妹,你的生命是否能夠成為別人的祝福呢?

崇真會的萬牧師,當他的癌症復發時,一群老師,同學去醫院探望他,想安慰萬牧師,但他們去到那裡,反而受到萬牧師的安慰.萬牧師說過一段這樣子的話,他說:人在地上的帳篷拆毀了,還有天上的房屋存留著,人生好比蠟燭,蠟燭有長有短,而重要的不是那蠟燭有多長,有多短,重要的是這枝蠟燭曾經發光!

各位,你的生命有發光嗎?你的生命有否成為別人的祝福呢?有人燃燒自己,照亮自己；有人燃燒別人,照亮自己；有人燃燒自己,照亮別人,使別人也能得著祝福.祝福別人的生命,是很美麗的生命；祝福別人的生命,是我們體會到上帝的祝福.或許你不覺得神祝福你,但你仍然去祝福別人的,這就更好!

如果我們的生命能夠互相祝福,我們就能有更美好的生命!盼望你們能透過亞伯蘭的事件中,我們也能夠有這種的心志,正如我的心意一樣:願我的生命成為別人的祝福,阿們!



\newpage

\section{從水和聖靈生的}
\label{sec:202}
\textbf{伍渭文牧師 - 從水和聖靈生的}
\newline
\newline
\hyperref[sec:201]{< < < PREV SERMON < < <}
~
\hyperref[sec:toc]{[返目錄]}
~
\hyperref[sec:203]{> > > NEXT SERMON > > >}


講題:從水和聖靈生的　The Flame From Above

講員:伍渭文牧師

場合:香港中文大學崇基學院禮拜堂

日期:2013年5月19日

教會舉行水禮,大概有三個日子:聖誕節,復活節,聖靈降臨節.

聖誕節紀念耶穌的降生,聖誕之前一連四主日為將臨期,紀念上帝的兒子降世,將臨期的主調是耶穌來臨,原本上帝的來臨是審判,但上帝不願審判我們,以耶穌來臨作為拯救,但我們要預備迎見上帝.所以將臨期第二,三主日福音經課的主角不是耶穌,是耶穌的先鋒施洗約翰,他在曠野喊叫:天國近了,你們應當悔改,直到第四主日,福音經課才記載耶穌在伯利恆的降生.所以,在聖誕節安排水禮,有悔改歸向基督的氛圍,不是聖誕節高興,在普天同慶的日子領洗.

另一個日期是復活節.該天舉行水禮,充滿表明我們領洗的人與基督同死同復活.但復活節舉行水禮,亦因為復活節之前有四十天的預苦期.教會用四十天效法耶穌在曠野接受魔鬼的試探,靠著聖靈勝過試探,始祖在樂園受魔鬼試探失敗,耶穌在曠野得勝.在預苦期,我們多用時間躬身自省,禁食,祈禱,行善,預備心靈,紀念基督為我們受難,受死,埋葬,第三天復活.預苦期是悔罪迎向救贖的節期,在預苦期開始的聖灰日有塗灰禮,清楚說明悔罪:(主禮把灰塗在的額上或手掌上) 「你要緊記,你本出於塵土,乃要歸回塵土,你要離開罪惡,跟隨基督.」

聖誕節和復活節的水禮,基本上都有前期的時間,用來悔過反思生命；將臨期到聖誕節,觸動我們的情緒(appealing to our emotion) ；預苦期到受難和復活,獨動我們的意志(appealing to our will) .

第三個領洗的日期是聖靈降臨節.聖靈降臨是教會的開始,而領洗的人是歸入基督的身體-教會.每一位今天領洗的人,一如每一位曾經領洗的人,他們的信主,是經過不同人的關心.種子撒下了,可能在其他地方才長大,在另一個地方才結果.教會不同的信徒編造成一條線,牽引著他們,直到今天.

崇基教堂每一年半才舉行一次水禮,上次是聖誕,在年底；今次聖靈降節,在年中.

聖靈降臨時,聖經記載:「五旬節到了,門徒都聚集在一處.忽然,從天上在響聲下來,好像一陣大風吹過,充滿他們的屋子；又有舌頭如火焰顯現出來,公別落在他們的頭上.」 (使2:1,2) 聖靈如火降在信徒的頭頂,所以今天聖帷是象徵火的紅色,全年祇有今天的聖帷是紅色.但水禮的水和五旬節的火有何相關呢?我們豈不是說:水火本不相容嗎?

在聖經的教訓中,水和火都有審判和潔淨的兩種意思.在今天洗禮禮文中的釋義,提到舊約時期上帝與人立約:洗禮是挪亞之約,接受洗禮的人,就像昔日進入方舟,得到拯救.在這裡,洪水是審判,因為地上的人的罪甚大,上帝要用洪水審判人.方舟被水擊打,而裡面的人免受洪水淹沒.就如地耶穌在十字架上被釘,擁攬了我們的罪過,擋住上帝要落在我們身上的審判,卻以他的公義聖潔生命披在我們身上,像結婚的人分享對方的財富,也分擔了對方的負債.婚禮中帶戒指時新人彼此說:「我以我的身心尊敬你,我以我的財物,與你分享.」是歡悅的對換(joyous exchange) .在挪亞方舟之約,水是審判.

火也是審判,上帝曾因為所多瑪,蛾摩拉(Sodom and Gomorrah)的罪惡甚重,降火滅兩城.(創19)

但水也代表潔淨,在洗禮釋義中,保羅提及「弟兄們,我不願意你們不曉得,我們的祖宗從前都在雲下,都從海中經過,都在雲裡,海裡洗歸了摩西.」(林前10: 2) 這裡的水指紅海.領洗是過紅海,離開埃及,進入應許之地.他們離開圾及,為要來到西奈山,和上帝立約,作者上帝的子民,歸了摩西,就是接受摩西從西奈山從上帝頒下來的律法.

其實,五旬節的源起,是紀念是 以色列人出埃及後第五十天,上帝在西奈山頒給摩西《十誡》的日子.所以這水是分別的水,紅海把埃及,為奴之地,和流奶與蜜之迦南地分開.印度一位著名聖徒孫大信,他寫了一首歌紀念他的洗禮:十架在前頭,世界在背後,永不回頭,永不回首.The cross before me, the world behind me, no turning back, no turning back.

火也是潔淨.上帝首次向摩西顯現:「耶和華的使者從荊棘裡火焰中向摩西顯現.摩西觀看,不料,荊棘被火燒著,卻沒燒毀.摩西說:我要過去看這大異象,這荊棘為何沒有燒壞呢?」(出3:2) 這是潔淨的火,不是消滅的火.對真金來說,火使真全更純正,對雜質來說,火是消滅,對金來說,火使金更完備.

我們說浴火鳳凰,經過苦難的熬煉,生命改變,比以前更充實更美麗.所以水和火都有審判和潔淨的意義,在聖靈降臨節領洗是水火相容.

其實聖靈降臨節,並不是第一次上帝從天而降火焰,第一次燔祭的火是從天而降:「有火從耶和華面前而來,在壇上燒盡燔祭和脂油.」(利9: 24) 從天上來的火畫夜不會消滅,因為早晚祭司獻上祭牲,其脂油-成為燃料,而每次獻祭要在其上點火,否則就是獻凡火.亞倫的兒子拿答和亞比戶就因獻凡火被上主擊殺而死 (利10: 1.2) ；火從天而降,表示人不能靠自己親近上主,要靠從上主而來的拯救.

在聖靈降臨節洗禮,特別突顯水禮的教會性.我剛才已題及,今天每一位領洗或堅信,都是因著不少教會信徒的努力,一步一步把他們引到上帝面前.而聖靈的火焰,不單是浴火鳳凰的火,改變自己的生命,這火燃點我們,把我們差遣,帶著能力,為上帝作見證.耶穌升天時對門徒們說:「不要離開耶路撒冷,要等候父所應許的……你們要受聖靈的洗.……但聖靈降臨在你們身上,你們就必得著能力；並要在耶路撒冷,猶太全地,和撒馬利亞,作我的見証.」(使1:4 ,5,8)

我們每主日誦讀的信經,信網三提到聖靈時都和教會有關:「我信聖靈,聖而公的教會,聖徒相通,罪得赦免,身體復活,並得永生.」受洗歸入基督的人,同時也受洗歸入基督的身體,我們都領受聖靈給我們的恩賜,在教會配撘,各司其職.我們希望每一位領洗的,和每一位經常來崇基教堂聚會的,當崇基教堂是屬靈的家的會友,一起使教會成長,讓人來到這裡,找到基督的救恩.

耶穌是將臨期聖誕期的應許的光,在顯現節成為耀眼的明亮,經過受難和十架的苦痛復活,達到正午的日照,到五旬節成為火焰,降到眾人的頭上,陪伴眾聖徒,直到再度降臨,進入永遠的榮耀.(The light of Advent and Christmas becomes brilliance in Epiphany, turns into high noon sun at Easter, and becomes flame from above descending on the heads of believers at Pentecost and will return in eternal glory.)



\newpage

\section{引我們進入真理的聖靈}
\label{sec:203}
\textbf{張伯笠牧師 - 引我們進入真理的聖靈}
\newline
\newline
\hyperref[sec:202]{< < < PREV SERMON < < <}
~
\hyperref[sec:toc]{[返目錄]}
~
\hyperref[sec:204]{> > > NEXT SERMON > > >}


約翰福音 16:12-13
\newline
\begin{longtable}{cl}
\hline
\hline
章節 & 經文 (和合本修訂版)\\
\hline
16:12 & \begin{tabularx}{0.7\textwidth}{X} 「我還有好些事要告訴你們，但你們現在擔當不了。 \end{tabularx} \\ \\ \relax
16:13 & \begin{tabularx}{0.7\textwidth}{X} 但真理的靈來的時候，他要引導你們進入一切真理。因為他不是憑著自己說的，而是把他所聽見的都說出來，並且要把將要來的事向你們傳達。 \end{tabularx} \\ \\
[1ex]
\hline
\hline
\end{longtable}
講題:引我們進入真理的聖靈　The Spirit That Makes Us Understand

講員:張伯笠牧師

經文:約翰福音16:12-13

場合:香港中文大學崇基學院禮拜堂

日期:2013年5月26日

我們來看約翰福音十六章12-13節:「我還有好些事要告訴你們,但你們現在擔當不了.〔或作不能領會〕只等真理的聖靈來了,他要引導你們明白〔原文作進入〕一切的真理．因為他不是憑自己說的,乃是把他所聽見的都說出來．並要把將來的事告訴你們.」

在這段經文裡,耶穌告訴門徒說:讓我們進入真理,明白真理是通過聖靈.那麼真理是什麼?彼拉多在審判耶穌的時候,他問了耶穌一句:真理是什麼?

我在一個中國的報紙上說:共產主義是宇宙的真理.真理到底是什麼?今天我們要通過什麼來認識真理?耶穌基督說:我就是道路,真理,生命．(約14:6)我們如何來認識真理?靠我們自己的理性,感性是不能認識真理的.我們要認識真理,必需通過聖靈.

其實今天在許多教會裡,對聖靈工作的認識有許多的偏差.我在美國牧會十幾年了,我經常會去一些的教會,有的把聖靈工作看成非常重要,有的卻是不重要.我記得有一次我去了新澤西州講道,那同工主席開車到機楊接我說:「張牧師,你講聖靈要小心,我們教會不太強調聖靈的,我們教會的三位一體是──聖父,聖子,聖經.」「那你們教會應該是異端的教會.」

我們如何來認識聖靈?聖靈是真理的靈,聖靈是基督的靈,聖靈是神的靈,聖靈是全人的神──我們所信的神是三位一體的神──聖父,聖子,聖靈.如果我們不認識聖靈,我們就不認識神；如果我們不認識聖靈,我們更不可能稱耶穌為主；如果我們不認識聖靈,我們不可能承認我們是罪人.

聖靈是神的靈,聖靈是三位一體神的第三個位格.我們知道:

1.    上帝天父──設計救恩；

2.    聖子──道成肉身來到我們中間,完成救恩；

3.    聖靈──當耶穌基督升天的時候,祂留下保惠師給我們,聖靈在我們生命中來施行救恩.

聖靈是百分之百的神,如果你不信聖靈,你就不信神,一點也不信神.

今天我從三點跟弟兄姊妹來分享:

一,    聖靈光照人──認罪悔改是我們認識真理,得到上帝兒子的位分；

二,    聖靈引導我們──進入真理,明白真理,順服真理,結出豐盛的聖靈果實；

三,    聖靈給我們能力──使我們可以過一個全人不同的,全新的生活,這是聖靈給我們的一個地位.

(一)聖靈光照人

聖靈給我們認識神兒女的地位,剛剛我們唸了那段經文前面說,當我們認識聖靈的時候,耶穌就差下保惠師給我們,因著祂,我們就「為罪,為義,為審判,自己責備自己．為罪,是因他們不信我．為義,是因我往父那裡去,你們就不再見我．為審判,是因這世界的王受了審判.」(約16:8-11)

因此,當我們來到聖靈的面前,我們感謝主,我們有這個活著的靈在我們生命中,使我們得到完全不一樣的生命.所以當我們認識聖靈在我們生命的地位時,我們就會非常的感恩,因為過去我們是在亞當的地位,當聖靈引導我們明白真理,進入真理,我們得到一個在基督裡的地位.

亞當的地位是一個敗壞的生命,從祖先來的,自然的生命.我們有兩種生命:一種是自然的生命──從母腹中生出來,會隨從今世的風俗,空中掌權者的旨意,會隨著世界的潮流來過生活.另一種是屬靈的生命──當我們認識基督,在聖靈裡得到重生的時候,我們會過一個全新的生活,是在基督裡的生活.聖經說:「若有人在基督裡,他就是新造的人．舊事已過,都變成新的了.」(後5:17)是新人,新造的人,完全不同的人.

講一個例子,一個小偷,他有的是偷盜的生命,他每一天都要去偷東西,偷到了就喜樂,偷不到就很沮喪.有一天,他沒有偷到東西,他心裡很憂愁,晚上睡不著覺,就聽見遠處有人唱詩歌,像是有許多人在聚會,他心想,人多就有機會,他就去參加那個聚會去了.他一進去,發現牧師在講道,牧師看見他就指著他說,過去偷的,就不要再偷了.他嚇了一跳,心想,牧師怎麼知道我來了?所以,他就不敢偷了.結果,牧師一呼召,他就走上前信主了.

你看,這生命很不一樣了,信主以後,聖靈充滿他,他認識自己是罪人,他認識到自己過去的行為是壞的,所以他在主裡面開始成長,願意作一個新造的人.人的生命如果不靠聖靈的能力以及在主裡的追求,我們慢慢的就被拉走了.有一天,他下班回家,他肚子餓,他去包子舖,問老闆說:「這個包子多少錢?」老闆一邊包,一邊賣,沒有理睬他,他心想,我這樣子叫你,你都不理睬,過去我都不用叫你的,這是多好的機會啊!

人的靈很重要,人的思想決定人的行為,當他在想的時候,手就動了,一拿到包子就走回家了.到晚上吃飯的時候,他剛要謝飯禱告,突然想到這個包子是偷來的,他就禱告不下去了,他也睡不著覺了,因此,他半夜三點鐘打電話給牧師說:「我偷了包子,睡不著覺.」「那過去你偷了東西,能睡著嗎?」「過去我偷東西能睡著,偷不到就睡不著.為什麼現在我偷了東西,反而睡不著呢?」「你覺得是為什麼呢?因為你的生命不一樣了!你從一個自然的生命裡,成為一個屬靈的生命.」

我們可以為罪,為義,為審判自己責備自己,所以當我們認識聖靈,就先認識神兒女的地位──過去我們在亞當裡,今天我們在基督裡,成為新造的人,是在神的愛裡,在教會裡成長,更多的認識我們所信的耶穌基督是什麼.

(二)聖靈引導我們

人需要通過聖靈才能明白真理,認識真理,進入真理,聖經告訴我們要進入真理,如果沒有聖靈的工作,我們很難認識自己是罪人.因為我們一般人都希望能通過智慧來認識真理,但人的智慧卻是有限的,因為人的智慧是根據你所學習的背景,經驗,眼界,學習… 這些都是非常的有限,所以當人談論聖經時,他們會用智慧,理性為出發點,而非從聖靈的工作.

我認識耶穌基督不是通過智慧,完全是上帝奇妙的作為,是聖靈的工作.1989年我被中國政府通緝,這已是24年前的事了.當被中國政府通緝的時候,我逃到農村,向我傳福音的人,用人的眼光來看,沒有任何知識,學識,地位的人,就連自己的名字也不會寫的農村婦女.她幫助我,使我能住在她家裡面,並向我傳福音,我是北京大學中文系,一個不識字的人向我傳真理,你不覺得很滑稽嗎?我當時也覺得相當滑稽.她叫我給她讀一本書,因為她不識字,我打開那本書,叫《約翰福音》,是手抄的筆記本,我唸:「太初有道,道與　神同在,道就是　神.這道太初與　神同在.」我看看她的表情,心想,她能聽得懂嗎?我連讀也讀不懂,她沒有文化,她能聽得懂嗎?但她卻聽得懂,她不僅聽得懂,她反而是我第一個聖經老師,她告訴我:「道就是耶穌基督,沒有宇宙就有耶穌,祂成為肉身來到我們中間,為我們的罪釘死在十字架上,充充滿滿將真理和恩典向我們顯明,讓我們認識祂,使罪得以赦免,重擔可以脫落,憂愁可以變為喜樂…」那時候,我認識耶穌是通過一個沒有文化的人.

我記得有一個數學博士,他到教會聽道十年了,但都還沒信耶穌.有一天聖誕節,當我講完道,看見他,就問他說:「弟兄,你認識耶穌沒有?」「沒有」「為什麼呢?」「因為我學數學的,我怎麼算,都沒有神.你能夠算到神嗎?」「你若能把神算出來,你一定比神大,祂一直都在你的思維之中,因為祂是創造宇宙萬物的神,祂超越所有,你怎麼能夠用有限的理性把祂給算出來呢?我問你,你能不能算出一加一等於二?是為什麼?」「我算不出來,我不需要算,那是公理,我用就可以了.」「對啊!我們為什麼要算上帝呢?我們通過聖靈來相信祂,認識祂,就可以了,並活出祂在我們生命中的見證就行了.」那一天,他就決志信主了,他在我辦公室裡認罪,一直在流淚,他追求主十幾年,他算不出神,那天聖靈感動他,他認識耶穌基督是他生命的主.我們用我們的辦法無法算出神來的.

有一個小故事,一位教授過河,有一位划船的老漁夫在那裡,他問老漁夫:「老先生,你懂文學嗎?」「我不懂」「你懂詩歌嗎?」「不懂」「你懂音樂嗎?」「不懂」「你懂科學嗎?」「不懂」「那你這人活得有什麼意思?你什麼都不懂!」這時候,風浪來了,船要翻了,老漁夫問他:「先生,你會游泳嗎?」「我不會」「那你命沒了!」

我們通過知識不可能認識真理的,我們必須通過聖靈的感動來到上帝的面前認識祂.聖經說:「真理的聖靈來了,他要引導你們明白〔原文作進入〕一切的真理.」我覺得這裡用「進入」比「明白」更好.像我們明白聖經,但卻不能進入神的生命之中,我們明白上帝在我們身上的旨意,我們卻不能活出上帝在我們身上的旨意.所以,我們要作一個讓神的真理進入我們生命的人,這樣子,我們才可以得到重生,聖潔,洗淨,如此,才可以過一個新的,不同的生活.

約翰福音十四章16-17節:「我要求父,父就另外賜給你們一位保惠師,叫他永遠與你們同在,就是真理的聖靈,乃世人不能接受的.」

當我們進入真理聖靈的時候,我們首先是口中承認,心裡相信,耶穌基督是我們生命的主,我們得到上帝兒女的地位──這是聖靈引導我們的第一件工作:引導我們認識耶穌是神的兒子,引導我們承認自己是罪人,我們成為一個新造的人,得救的人,屬神的兒女,得屬神的位分.

彼得跟耶穌的對話中說:「西門彼得回答說,主阿,你有永生之道,我們還歸從誰呢.」(約6:68)當我們認識耶穌基督的時候,還要持續的跟隨耶穌基督.因為我們認識了真理,還不足夠,我們要進入真理,跟隨真理,活出真理.一個基督徒的生命是要活出來的,當他認識了真理,這真理要在他生活中彰顯出來的,不管是在他的工作中,家庭中,生活中,在面對著一個世界的挑戰中,他要活出真理,見證真理,這是非常重要的.

我們如何來活出真理?我們要求聖靈來幫助我們.我們不是一個宗教徒,乃是活生生上帝的兒女,因為神的靈是活潑的靈,祂的靈在我們生活中引導我們,我們靠著神的靈,在一切事上來到上帝的面前,聖靈的面前,我們便能得勝.

聖靈不但是位格,祂也是感動.祂住在我們生命裡,感動我們為罪,為義,為審判,可以活出耶穌基督給我們的生命.

當我信主以後,牧師告訴我,我當過一個全新的生活,主所喜悅的生活,就是每個禮拜要來作禮拜,每天要讀聖經,禱告,當你遇見別人委屈你,謾罵你的時候,你要用耶穌基督的愛,胸懷去包容他,要為你的仇敵禱告(有人打你左臉,把右臉也給他打,要你走一里路,你就走二里路),要忘記過去的仇恨… 這是很不容易的.

今天我們成為一個基督徒,在教會裡生活,如果沒有聖靈的感動,我們在教會裡的生活是枯燥的,對真理的認識是枯燥的.一旦有聖靈在我們裡面,我們會突然發現,我們能夠活出一個完全不一樣的生命.

美國一個曾經非常有名的首富洛克菲勒,他說:「我一生中持守了一個基督徒該有的生活,小時候母親說:『你已經認識真理了,你信了耶穌,就該活出基督徒該有的樣式,母親就給了他三個約定:過十一奉獻的生活,到教堂要坐第一排聽道,順從教會所有的事情,不要讓牧師操心.』這就是我成功的秘訣,我一生都遵守了母親給我的約定,我過了一個正常的基督徒生活.」難道這三點都達不到,就不是正常的基督徒生活嗎?當記者這樣子問他的時候,他回答說:「小時候我到教會來,母親都會給我20美分,並要我嚴格地守十一奉獻,20美分很容易,一點都不難,我就把2美分交給教會,可是我後來成為億萬富翁,這就比較難了!但我一生都遵行了母親給我的教導,從小到我離世之前,我皆奉獻超過十分之一.這訓練了我的胸懷,一個給予的胸懷.第二,母親說聽道要坐在第一排,所以我要提早45分鐘到達才能坐在第一排,這就訓練了我的紀律.第三,順從教會的決定,聽從牧師,雖然我覺得教會或牧師的決定,不一定都是對的,像對於奉獻的事,我都不去過問(註:他一生奉獻金錢建立了一千六百間教堂,一旦他奉獻了,他就知道這錢不再屬於他,這是上帝的錢,由教會來決定如何來使用),這訓練了我的順服.」

一個人在他的生命中,他有一個很重要的品德,就是聖靈引導他,進入一個正常的,神兒女的生活,來跟隨上帝,耶穌基督.

我感謝主,當我認識神以後,神改變了我.我小時候是一個非常叛逆的孩子,我的父親生了七個孩子,我是最不聽話的一個.我父親常說,這孩子早晚會出事,惹禍的,因為這孩子太不順服了.1989年我被中國政府通緝,我的通緝令透過電視台,向全國發放,我家人正在吃晚飯的時候,我的照片上了中央電視台,我媽媽一看見就哭了,而我父親沒有哭,他說:「你看,我當年說對了吧?我早晚就知道他會有這一天.」當我信了主之後,上帝改變了我,我從一個最不順服,最有自己想法的人,變為完全降服在聖靈的引導之下,凡事順服神的話語,聖靈的引導,好叫聖靈在我們生命中結出不一樣的果實來.耶穌基督說:「我是葡萄樹,你們是枝子.」(約15:5)我們需要在這葡萄樹上與耶穌連結,好結出耶穌基督給我們該有的聖靈果實,就是「仁愛,喜樂,和平,忍耐,恩慈,良善,信實,溫柔,節制．」(加5:22-23)如此豐盛的果實.

弟兄姊妹,今天你成為基督徒,你就要為主發光,尤其是在末後的時代.這時代,越來越不相信上帝的存在,不按上帝的旨意來生活.現今的世界已不用聖經來看世界,乃以世界的風俗來看世界,但我們作為神的兒女,如何為神來堅守真道呢?我們要用聖經的世界觀來看世界,只有聖經的世界觀是正確的.如果中國能走向文化基督化,中國的文化才有希望；中國只有走向教會的國度化,中國的教會才有希望；中國只有走向福音化,中國民族才有希望.今天中國人在想些什麼?中國現在有錢了,有錢了該做些什麼?中國人在世界上的地位越來越高了,我們該做些什麼呢?中國人在全世界被評價的指數是最低的,百分七十的國家,認識中國人沒有素質,為什麼呢?是因為我們沒有神,我們只靠我們的金錢.中國跑去全世界去消費,卻無法贏得人的尊重.是因為我們沒有生命,素質.這素質是從上帝那裡來的.

(三)聖靈給我們能力

一個人當他認識神,他就能活出真理──聖靈引導我們活出真理.使徒行傳第一章八節:「但聖靈降臨在你們身上,你們就必得著能力.」什麼樣的能力呢?就是「在耶路撒冷,猶太全地,和撒瑪利亞,直到地極,作我的見證.」──被聖靈引導的生活,就是要活出神兒女的見證.耶穌基督說:「你們就是這些事的見證.」(路24:48)

當耶穌基督要我們去傳福音的時候,祂沒有要我們作一個辯護律師,乃是作見證,這見證就是活出耶穌基督的見證.

當我從中國逃到美國的時候,我病了,住在美國醫院一年,台灣醫院一年,再回到美國,進進出出醫院有兩年的時間,因此,大概有四年是在疾病,死亡危險中渡過的,這是非常痛苦的四年.我在醫院時,常看到一些沒有信仰的人,他們在苦難中表現憂愁!感謝主,當時候我已信了耶穌基督,我深知我的生命在神的手中,即使我離開這個世界,也沒什麼可怕的,因為我已經有了生命得救的確據,我可以到上帝那裡去.但如果上帝讓我活下來的話,我該怎麼活呢?這四年,使我一切的夢想都破碎了,那時候,醫生告訴我,我只能活一年… 當苦難來臨的時候,你還能靠誰呢?病了一年後,朋友都沒了,那時候你就知道你是誰!所以我那時候就向神禱告:「如果你讓我活過去,醫治我的疾病,我就過另外一種生活,單單討你喜悅,一生服事你,榮耀你.」四年以後,我好了,進了神學院.感謝主,在神學院裡,學習破碎自己,倚靠聖靈來過一個全新的生活.那時候,生命就不太一樣了,感謝主!

聖靈是讓我們成為一個新造的人,也讓我們得著能力.一個可以過新生活的能力,傳揚福音的能力,做一個討耶穌基督喜悅的能力.這能力很重要,因為不是所有偉大的人物都有能力的,我們也不是每天都有能力的,我們需要聖靈的能力.即使像王明道這麼偉大的聖徒,他也有軟弱的時候──1955年,當他被抓時,解放軍的稽查員,從他家的牆上跳進去,拿著他對著他說:「不許動!」他嚇到倒在地上,也寫了認罪書,就參加三自,後來聖靈給他能力,他認識到他所做的,都不討神的喜悅,他就向公安人員說,要收回他的自白書,最後他被判了無期徒刑.

所以,當我們遇到難處,危險的時候,我們需要聖靈的能力,有了聖靈的能力,我們才可以過一個全新的生活.我們必需追求聖靈在我們生命中的充滿,引導,讓我們不僅僅進入真理,明白真理,也能活出真理.

在我教會有個曾經當解放軍團長的弟兄,來到美國在一家餐館打工,可是卻碰見當年打越南時的一個越南兵的班長,這班長是大廚,而這弟兄是服務員.他們常在餐館吵架,有一天,他們又吵起來,這越南兵在一米的距離內,開搶打在這弟兄的肚子上,造成這弟兄身上有十一個洞.到了醫院,他就想,我死了以後,去那裡?他突然想到,當他剛來到美國,去教會學英文,牧師告訴他,學英文不重要,認識真理才重要,如果你生命中有因難,並走到危險的境界時,就呼求耶穌… 他就喊出來了,心裡面說:「如果你存在,你救我!」一個月以後,他沒死,也沒信主,老婆也與他離婚,孩子不見了,除了兩佰塊錢,什麼都沒有,又不懂英文,就想死掉算了,反正活著也沒有意思,如果兩佰塊錢花完了,人也完了,人生的盼望在那裡?後來,他認識了一位牧師,那牧師帶領他信主,他也就受洗了.你看,聖靈引導我們,我們認識了真理,進入真理,我們便得了生命的能力!他說:每天生活過得真喜樂,神充滿我.

「若有人在基督裡,他就是新造的人．舊事已過,都變成新的了.」(林後5:17)

當我們要被聖靈充滿,首先我們要高舉耶穌基督,不是高舉自己；一個被聖靈充滿的人,是一個勇敢的傳揚福音；一個被聖靈充滿的人,是一個非常愛靈魂的人,願意把親朋好友帶到耶穌基督的面前；一個被聖靈充滿的人,要過聖潔的生活,因為聖靈是真理的靈,也是聖潔的靈；一個被聖靈充滿的人,一定會結出聖靈的果實,在這世界上能夠榮耀主的名,因著我們活出耶穌基督,我們把一切的榮耀歸給主.

弟兄姊妹,你願意被聖靈充滿嗎?我們來禱告:

天父,我感謝你,在這時刻,我把每個舉起手打開心的朋友交在你的手上,他們今天成為你的兒女,聖靈今天就進入他們生命當中,掌管他們的生命,引導他們的生命,祝福他們的生命,主!求你與他們同在.我們也把每一位來自中國大陸的學生,學者交在你的手上,求你引導他們進入真理,明白真理,成為上帝的兒女,無論走到那裡,都為你活出美好的見證,祝福我們的教會,與我們同在,奉主耶穌的名,阿們.



\newpage

\section{自由的滋味}
\label{sec:204}
\textbf{陳曉東博士 - 自由的滋味}
\newline
\newline
\hyperref[sec:203]{< < < PREV SERMON < < <}
~
\hyperref[sec:toc]{[返目錄]}
~
\hyperref[sec:205]{> > > NEXT SERMON > > >}


加拉太書 1:1-12
\newline
\begin{longtable}{cl}
\hline
\hline
章節 & 經文 (和合本修訂版)\\
\hline
1:1 & \begin{tabularx}{0.7\textwidth}{X} 我使徒保羅和所有跟我一起的弟兄，寫信給加拉太的眾教會。我作使徒不是由於人，也不是藉著人，而是藉著耶穌基督與使他從死人中復活的父神。 \end{tabularx} \\ \\ \relax
1:2 & \begin{tabularx}{0.7\textwidth}{X} 見上節 \end{tabularx} \\ \\ \relax
1:3 & \begin{tabularx}{0.7\textwidth}{X} 願恩惠、平安從我們的父神和主耶穌基督歸給你們！ \end{tabularx} \\ \\ \relax
1:4 & \begin{tabularx}{0.7\textwidth}{X} 基督照我們父神的旨意，為我們的罪捨己，要救我們脫離現今這罪惡的世代。 \end{tabularx} \\ \\ \relax
1:5 & \begin{tabularx}{0.7\textwidth}{X} 願榮耀歸給神，直到永永遠遠。阿們！ \end{tabularx} \\ \\ \relax
1:6 & \begin{tabularx}{0.7\textwidth}{X} 我很驚訝你們這麼快就離開那位藉著基督之恩呼召你們的神，而去隨從別的福音； \end{tabularx} \\ \\ \relax
1:7 & \begin{tabularx}{0.7\textwidth}{X} 其實並沒有另一個福音，不過有些人騷擾你們，要把基督的福音更改了。 \end{tabularx} \\ \\ \relax
1:8 & \begin{tabularx}{0.7\textwidth}{X} 但無論是我們或是天上來的使者，若傳福音給你們，與我們所傳給你們的不同，他該受詛咒！ \end{tabularx} \\ \\ \relax
1:9 & \begin{tabularx}{0.7\textwidth}{X} 我們已經說了，現在我再說，若有人傳福音給你們，與你們以往所領受的不同，他該受詛咒！ \end{tabularx} \\ \\ \relax
1:10 & \begin{tabularx}{0.7\textwidth}{X} 我現在是要得人的心，還是要得神的心呢？難道我在討人的喜歡嗎？我若仍舊想討人的喜歡，我就不是基督的僕人了。 \end{tabularx} \\ \\ \relax
1:11 & \begin{tabularx}{0.7\textwidth}{X} 弟兄們，我要你們知道，我所傳的福音不是按照人的意思； \end{tabularx} \\ \\ \relax
1:12 & \begin{tabularx}{0.7\textwidth}{X} 因為我不是從人領受的，也不是人教導我的，而是藉著耶穌基督的啟示而來。 \end{tabularx} \\ \\
[1ex]
\hline
\hline
\end{longtable}
講題:自由的滋味　A Taste Of Freedom

經文:加拉太書1章1-12節

講員:陳曉東博士

場合:香港中文大學崇基學院禮拜堂

日期:2013年6月2日

當我們的世界正處於在一種危險當中的時候,當我們今天所看見的世界,我們所生存的世界,濫用自由的現象,正在毀滅著人類的自由,我們需要回到神面前懺悔.

自由是一個大的題目,名詞.許多人想到自由的時候,都不知道從何說起.許多的哲學家,思想家,社會行動家,為自由下了不同的定義,當我們想到自由的時候,在你的腦袋裡,你想到什麼?當我坐到書桌前預備今天的信息時,許多思想在我腦袋裡出現──自由──我想到今天流行的自由主義,也想到我喜愛的一位法國哲學家笛卡爾 (René Descartes) 所說的:「我思故我在」.自由也讓我想到今天在許多人咀邊說的:「我的權益」.

一位英國牧師說:在我們今天的世界裡,我們真的窮得剩下 I-Me-Mine,所有的自由歸咎於個人的權益的思考裡.自由也使我們想起許多人所說:我就是我自己的主人:I am my master.如此,我們就看見在我們生活的世界裡,所有的反傳統,反權威,生命自覺,倫理的選擇… 都行使在自由的口號裡.

自由,自由…　有多少人用這樣一個名字來去行盡諸般的惡事,我們看見整個社會,人類文明的發展,都向著相對主義,價值判斷的中立來去推進.想到自由,就聯想到物競天擇,市場經濟,功利主義,放縱的情慾當中.或許有些人想到自由的時候,並不對這自由有太多的好感,但允許我們回到神的話語中,來為自由辯護.

今天的福音故事,讓我們在為惡,危急中得自由的故事.默想這之前,讓我先分享一個故事:

我有一次去歐洲的行程,特定繞了一個彎去了倫敦Westminster座堂看一些的雕像,其中一個雕像是一位中年的牧師,這牧師故事如下:

於1907的時候,在雲南武定縣有一個孩子出生,他出生的前一年剛好有宣教士在他所在的地方設立了教會,宣教士在那裡修建學校,使他們可以免於無知的自由；也設立醫院,使人可以脫離危急,病痛,使人得健康的自由.宣教士帶領苗族的百姓脫離愚昧,醫治肉體的痛苦.1921年這孩子在宣教士所辦的學校裡唸書,三年後又去了內地會所辦的學校進修,1926年,這孩子成為年青人,也成為教會,學校裡的老師,並受到教會的指派前往雲南嵩明,陸豐宣教,教學.1935年,他回到自己的家鄉,這時候抗戰爆發,而他仍然繼續教學,建造教會.中日戰爭時,許多的宣教士被迫的要跑到別的地方去逃難,但這年青人仍留在自己的家鄉,於撒普山總堂做傳道人.這年青人很有才幹,就被推薦為當時武定地區教會的會長.在沒有任何的幫助下,他辛勤地去牧養教會,福音在武定地區迅速的傳播,根據一些歷史的記載,一直到解放前的時候,在雲南武定這地方,已經有5500多的苗,彝,僳僳族的群眾在撒普山總教堂裡受洗,歸入主的名下.這年青人還編了苗文的頌主聖歌.1951年時,昆明教會按立他成為牧師,當時候,雲南已經成為解放區,他遵行聖經,順從掌權者的教導,和其他苗族教會的領袖在政府的要求下,簽署了三自宣言,但是他卻拒絕一件事,就是拒絕參加任何批鬥大會,他曾經說過:「我的手曾經為許多信徒施洗,不可用來沾染罪惡」.但好景不長,撒普山教堂,被政府關閉,這位牧師也不能繼續他的牧職,他被迫要去勞動生產,作為一位基督的僕人,他深知道一件彼得所教導的道理──順從神,不順從人的教導──所以當上帝的旨意和凱撒的命令發生矛盾的時候,他選擇了一件對他生命帶來挑戰的事情,他決定順服神的旨意,儘管聚會被禁止了,但是為了牧養羊群的需要,他帶領信徒在家中繼續地聚會,可惜在當時那種嚴酷的政治環境底下,他們的聚會被視之為反革命而被搗破了.在審判這位牧師的判決書裡說他:「死不悔改,繼續從事宗教間諜」.由於他在苗族裡擁有極高的威望,政府出於民族政策的需要,把他釋放了,且安頓在當時的楚雄州政協籌委會裡工作.1956年,他竟以牧師身份,代表苗族代表團副團長,上北京參加國慶典禮,受到當時的毛澤東主席的接見.可惜,文革開始以後,對基督徒的迫害更加的猛烈,於1969年他和21位的同工一同被捕,在獄中受到非常可怕的虐待,1973年12月29日在萬人的批鬥大會上,這位已經66歲的牧師,被判處死刑.見證者說,士兵用刺刀把他的舌頭當場攪碎,以防止他妖言惑眾,他用自己的血走完他在世上服事的歷程,在主裡安息了.在1998年時,倫敦 Westminster 決定在教堂的大門口,為24紀十位普世教會的殉道者樹立石像的時候,這一位寂寂無名的中國傳道人與馬丁路德(Martin Luther),潘霍華(Dietrich Bonhoeffer)等十位在世界上有名,無名的殉道者同列,他的名字叫王志明牧師.

當我在Westminster座堂前瞻仰這位與十位殉道者同列在大門口,頭上戴著苗族頭巾的雕像,我就默想自由的代價時,我看見與他同列的馬丁路德博士,並想起他說過的一句話:”Darkness cannot drive out darkness; only light can do that. Hate cannot drive out hate; only love can do that.” ──「黑暗不能驅逐黑暗,唯有光明可以做得到；仇恨不能驅逐仇恨,唯有愛可以成就」.

我們所生活的世界,是一個充滿了動盪,不安的世界,當我們從王志明牧師的生命見證中,來去默想什麼叫做「自由」的時候,我們也從保羅的書信中,看到一個非常重要的話語和信息,讓我們知道,在這個世界裡,真正能夠保障自由的不是我們人為所造出來的價值和觀念.耶穌基督才是自由的真正保障.

為什麼說「耶穌基督是自由真正保障」呢?因為耶穌基督就是自由,是自由的泉源,基礎,力量,生命.保羅說:「主的靈在那裡,那裡就得以自由.」(林後3:17)從這句精簡的話語裡,我們可以展開許許多多的默想,我們明白極其愛這個世界,當祂願意把救恩賞賜給人的時候,這救恩的內在實質,是神人和好的關係.而神人和好是帶來真正的自由,完全.為了保障這自由,聖經告訴我們,耶穌基督本來是與神同等,但祂不以這關係為強奪(腓2:6),他選擇了一件在人類歷史上從沒發生過的事,就是將道成為奴僕的肉身,木匠的肉身,在人類歷史中與人類同行,這並不是一個什麼浪漫的故事,而是非常血腥,可怕的故事,因為這個愛的故事,它的句號是羅馬的刑具,是十字架的等待,然而我們知道,教會並不是停留在十字架的哀號中,乃是與生命復活主相遇,並在聖靈降臨的大能中,對耶穌基督的再臨──教會的子民,在邊緣中的婦女,兒童,愚昧無知的小民──選擇了一個新的生命,生活,是以自由的群體來圍繞著十字架,主的身體與寶血來站立起來.無論這世上有多少的困阨,危難,這樣一個的自由群體,就成為了一個守護者,自由的實踐者,這才是真正的教會.

在聖經裡,我們看見這位自由的主,如何進入忍受困苦的生命,人在危難,急病中,把喜樂,醫治,關懷,愛帶進黑暗世界裡.當耶穌進入祂家鄉附近的一個迦百農村子時,路加如此記載說,村子裡的百夫長,是羅馬權利象徵的角色,他聽過耶穌的事,他家裡發生了一件事情──是他用盡他在世界上的所有的權利,知識,能力所不能解決的──他所愛,有用的僕人得了非常嚴重的疾病,所以他托當地的猶太人的長老,求耶穌去救他的僕人,當猶太人的長老到了耶穌那裡說:「你給他行這事,是他所配得的．因為他愛我們的百姓,給我們建造會堂.」(路7:4-5)耶穌聽了之後就與他們同去,離百夫長那家不遠時,聖經記載著:「百夫長託幾個朋友去見耶穌,對他說,主阿,不要勞動．因你到我舍下,我不敢當．我也自以為不配去見你,只要你說一句話,我的僕人就必好了.因為我在人的權下,也有兵在我以下,對這個說去,他就去．對那個說來,他就來．對我的僕人說,你作這事,他就去作.」(路7:6-8)他從推己及人的方式,看到這樣一位的先知,夫子,祂有醫治,能力的權柄,就毋需勞動到自己去做些什麼事.而耶穌對他這樣子的反應,聖經說:「耶穌聽見這話,就希奇他,轉身對跟隨的眾人說,我告訴你們,這麼大的信心,就是在以色列中我也沒有遇見過.」(路7:9)

從這樣子的一個故事裡,路加有幾個肯定,透過一個跟耶穌沒有直接關係的人,來指出一件非常重要的事:上帝的話語具有釋放人心,使人得自由的力量,這就是我們所說的福音.福音就是上帝的話語,它能釋放人心,使人的生命得自由,安康,平安.上帝的話語帶有能力,因為祂自己是能力的源頭,本體,自由,平安.

路加也肯定另外一件事:任何一個信心的操練,學習,不是停留在一些理論,討論裡.信心能夠發生,是需要有一個具體的環境.在我們每個人具體的環境裡,我們遇見自己,明白靠著自己無法解決不了的處境時,我們可以經歷到信心的臨在,同行,信心帶我們走出困阨的結果.

除此之外,路加也肯定說,當上帝的話臨到人的生命中時,我們的生命可以得著自癒,身心可以重獲自由,這是一個讓我們身心可以重新得以自由的「信」.

所以弟兄姊妹,當我們明白了我們剛才所說的,明白了當我們要品嚐自由的滋味時,我們要準備為自由付出代價,在處境中去體會得自由的喜樂.

今天,讓我們一起來敞開我們的心靈,願意在上帝的面前有一個禱告,一齊持守使人得自由的福音.

馬丁路德也講過有關基督徒的自由或基督徒如何看自由這事.馬丁路德有一個基本的看法.他認為每個人都是宗教人──有宗教信仰或沒有宗教信仰的人都是個宗教人──都有一種的傾向,希望透過自己的努力來達到一種價值觀的要求,好確認自己存在的價值,權勢,地位,名聲,金錢,使自己覺得自己是很重要的.當馬丁路德唸到加拉太書的時候,他很清楚的看見,保羅所指出的「因信稱義」是人類靠功德稱義的慾望,這在我們心裡是根深柢固的,就算我們已歸入主的名下,相信福音,跟隨福音的基督徒,許多時候也會故態復萌的,因為在我們內心,意識深處裡,還是希望靠自己某一種的功德,力量,能力,使我們能夠成就「神人和好」的得救關係.馬丁路德說:「人的本性,文化,在根本上是抵擋福音的；就以智慧和正直這兩種天賦為例,沒有基督,智慧是雙倍的愚蠢,正直是雙倍的罪惡,因為它們不僅不能認識到基督的智慧與義,還要阻攔,褻瀆基督的救恩.」保羅很公正的稱這個世界是罪惡和邪辟的,因為世界正處在它最好的時候,也是它最惡的時候,最嚴重的惡習和世界的智慧與義相比時,也不過是一個小小的過失而已.這些事阻攔人接受基督的福音,因為對於沒有基督的人而言,他們越好越有智慧,他不理睬,抵擋福音的可能就越大.這個世界,或許會看一些事是好的,不過,若沒有基督,它們都是錯的.

在人類的世界裡,我們認為自己有智慧,公義,理性,戰無不勝的思想武器,能夠愚公移山改變世界… 當我們靠著人類的智慧,人類所建造的價值,把自己的世界建造出來的時候,在皇冠上就是廣島的核子彈,國與國之間的軍事對峙,以聰明才智彼此恐嚇,威嚇….因此,從教會的先賢中,我們看到一個智慧；從聖經中,我們得到一個警覺的提醒,不要以自己有的智慧和義,來使自己成聖,完全的根據,因為那是上帝給我們的天賦,為了要成就一個更深的義和智慧型,就是當「神人和好」的時候,人在基督裡成就了真正的和平,自由.

保羅在加拉太書一章四節說:「基督照我們父　神的旨意為我們的罪捨己,要救我們脫離這罪惡的世代．」耶穌基督道成肉身是我們自由的保障；耶穌基督和平的福音才是生命得以完全的基礎──這也是保羅在加拉太書反覆論證的一個要點.在加拉太書五章一節說:「基督釋放了我們,叫我們得以自由.」

這個世界唯有回到耶穌基督那捨己,服事,與人和好的生命中,我們才可以得以自由──「基督釋放了我們,叫我們得以自由,所以要站立得穩,不要再被奴僕的軛挾制.」十三節說:「弟兄們,你們蒙召,是要得自由．只是不可將你們的自由當作放縱情慾的機會.」

對於自由的討論,許多人喜歡把它放在兩個範疇當中來說明:

一,    我們從怎樣的光景中得自由(Freedom from somewhere)；

二,    我們用自由進入到什麼光景中(Freedom to somewhere).

所以,從什麼環境中得自由,靠著自由進入什麼光景中,成為我們對自由有一個真實的把握.

我們今天的世界要從傳統的權威,價值,看法當中自由出來的時候,我們又靠著這自由進入一個怎麼樣的光景中呢?唯有我們看到,當人類擁有了他自己所以為擁有的自由時,他並不沒有好好把握這種的自由,他用自由進入一個保羅所說的,「放縱情慾」的自由.因此,我們要更多的思考耶穌基督所帶來的自由,這自由如何能幫助我們從生命裡那情慾的混亂中,得以被釋放出來.我們不再成為自己或世界慾望的奴僕,或權勢控制的奴僕,乃靠著耶穌,在基督裡,我們可以響往,追尋一種愛的自由,生命得潔淨的自由,追尋生命成聖的自由.

真正的自由是與聖靈同行的自由,是允許耶穌基督的靈在我們生命中結出豐盛的果實,而這果實可以奉獻於我們的生活,世界,宇宙,人的自由.保羅說:「聖靈所結的果子,就是仁愛,喜樂,和平,忍耐,恩慈,良善,信實,溫柔,節制．」(加5:22-23)

我們要安靜地敞開我們的生命,心靈,向神呼求從聖靈而來的自由:

主啊!你聽見我們的呼求,你聽見我們的心聲,在這樣一個看似平靜的世界中,隱藏著許許多多的風暴,我們所在的城市,面對著當怎樣行,無論是政府,民間,都有許多的討論,但是,我們祈求聖善的靈,給在這城市中那些屬你的子民,能夠舉起聖潔的手,為這個城市得著真正的自由來禱告,讓這城市裡有仁愛,喜樂,和平,忍耐,恩慈,良善,信實,溫柔,節制的自由,幫助我們有勇氣面對前面那些不可預測的結果,祈求自由的靈,在我們所愛的國家中,繼續地運行,引導許許多多的子民進入你自由,喜樂的福音當中,奉靠耶穌基督聖名祈求,阿們.



\newpage

\section{主的恩典夠用}
\label{sec:205}
\textbf{謝喜鵬先生 - 主的恩典夠用}
\newline
\newline
\hyperref[sec:204]{< < < PREV SERMON < < <}
~
\hyperref[sec:toc]{[返目錄]}
~
\hyperref[sec:206]{> > > NEXT SERMON > > >}


列王記下 4:1-7
\newline
\begin{longtable}{cl}
\hline
\hline
章節 & 經文 (和合本修訂版)\\
\hline
4:1 & \begin{tabularx}{0.7\textwidth}{X} 有個先知門徒的妻子哀求以利沙說：「你的僕人，我丈夫死了，他敬畏耶和華是你所知道的。現在有債主來，要帶走我的兩個孩子給他作奴隸。」 \end{tabularx} \\ \\ \relax
4:2 & \begin{tabularx}{0.7\textwidth}{X} 以利沙對她說：「我可以為你做甚麼呢？告訴我，你家裡有甚麼？」她說：「婢女家中除了一瓶油之外，甚麼也沒有。」 \end{tabularx} \\ \\ \relax
4:3 & \begin{tabularx}{0.7\textwidth}{X} 以利沙說：「你到外面去向所有的鄰舍借器皿，要空的器皿，不要少借。 \end{tabularx} \\ \\ \relax
4:4 & \begin{tabularx}{0.7\textwidth}{X} 然後你回家，關上門，你和你兒子在裡面把油倒在所有的器皿裡，倒滿了就放在一邊。」 \end{tabularx} \\ \\ \relax
4:5 & \begin{tabularx}{0.7\textwidth}{X} 於是婦人離開以利沙去了。她關上門，把自己和兒子關在家裡。他們把器皿拿給她，她就倒油。 \end{tabularx} \\ \\ \relax
4:6 & \begin{tabularx}{0.7\textwidth}{X} 器皿都滿了，她對兒子說：「再給我拿器皿來。」兒子對她說：「沒有器皿了。」油就止住了。 \end{tabularx} \\ \\ \relax
4:7 & \begin{tabularx}{0.7\textwidth}{X} 婦人去告訴神人，神人說：「你去賣了油還債，你和你兩個兒子可以靠著所剩的過活。」 \end{tabularx} \\ \\
[1ex]
\hline
\hline
\end{longtable}
講題:主的恩典夠用　God’s Grace Is Enough

講員:謝喜鵬先生

經文:列王記下4:1-7

場合:香港中文大學崇基學院禮拜堂

日期:2013年6月9日

親愛的弟兄姊妹,今天,我在這裡,用一段聖經跟弟兄姊妹一起勉勵!

請翻開列王記下4:1-7:

1 有一個先知門徒的妻,哀求以利沙說,你僕人我丈夫死了．他敬畏耶和華是你所知道的．現在有債主來,要取我兩個兒子作奴僕. 2 以利沙問他說,我可以為你作甚麼呢,你告訴我．你家裡有甚麼,他說,婢女家中,除了一瓶油之外,沒有甚麼. 3 以利沙說,你去,向你眾鄰舍借空器皿．不要少借, 4 回到家裡,關上門,你和你兒子在裡面,將油倒在所有的器皿裡,倒滿了的放在一邊. 5 於是婦人離開以利沙去了,關上門,自己和兒子在裡面．兒子把器皿拿來,他就倒油． 6 器皿都滿了,他對兒子說,再給我拿器皿來.兒子說,再沒有器皿了,油就止住了. 7 婦人去告訴神人．神人說,你去賣油還債,所剩的,你和你兒子可以靠著度日.

親愛的弟兄姊妹,我自己介紹我自己,1961年出生,我15歲蒙呼召.有一個老師在講道當中,我說:主啊!您賜給我有講道的恩賜.那時候我還不知道,2 年之後,17歲就走上講台,今年五十多歲,做了30多年.27,8歲時我藐視耶和華,輕視耶和華,厭棄耶和華,不要當醫生,不要當老闆,沒有技術,講聖經,要吃西北風,所以軟弱了五年,行屍走肉… 後來,我說:主啊!我真的怕了!我重新回到講台.

神是信實的!當我離開祂,祂仍然把講道的恩賜給我.今天能夠在這裡,分享主的話,我把每一場道當作最後的一堂傳道,我後來很認真,去理解聖經的話,我不是講者,我是聽者,這段聖經告訴我們一些什麼.

有個先生的婦人,他老公離開了世界,他老公年輕的時候就離開世界,別人離開世界留下兩棟房子給他太太,他留給他太太的是一本帳簿,人家是五百萬,他是負資產,這個婦人該怎麼辦?

孤兒寡婦,沒有還債的能力,有2個祈求:一,找上帝,二,上帝找不到,就找神的僕人.神的僕人啊!我老公是你的學生,常常這個禮拜堂講道,沒想到他怎麼這麼短命?離開世界了,現在我求你來幫我,幫我什麼?我欠人的錢,沒辦法還?她又做什麼?拿兒子抵債!

以利沙說:姐妹,你別緊張,喝杯水,慢慢說.你家裡有什麼?什麼都沒有!米缸,沒有米；水桶,沒有水；存摺,沒有錢,山窮水盡.

以利沙說:再想想,再想想-哦!我這裡除了一瓶油,一瓶小小的油,沒有什麼.你家裡有多少空器皿?窮人家庭沒有那麼多.

你回去,借空器皿,不要少借,要多借,關上門,母子把油倒出來,倒滿的放在一邊,倒滿的再放在一邊.兒子在倒油,媽媽拿器皿,器皿有沒有?器皿沒有了.油就止住了.他們就告訴以利沙:先生,家裡的器皿都滿了.他說,你拿回去市場去賣,把他換成美金,放在中國銀行,到世界哪個地方都可以刷卡.(你們有沒有中國銀行的卡呢?香港不說中國銀行,我們大陸就說中國銀行,中國銀行世界通).

親愛的弟兄姊妹們,我今天讀了這個聖經,有六句話:

一,    有還債的責任,沒有還債的能力.

你老公欠別人的錢,父債子還,夫債妻還,天公地義.

在四川,汶川地震的時候,有兩母子,那個樓房要倒榻的時候,媽媽的上半身被壓傷了,小孩子的手都給壓傷了,連動都動不了.孩子在裡面哭:媽媽,你救救我!媽媽聽到他孩子哭的聲音,眼淚往下直流下來,眼睛睛看著兒子走向死亡的邊緣,愛莫能助.

這個婦人,有人要叫他還錢,你是信耶穌的,講道理嗎?欠錢就還錢嘛?怎麼欠錢不還錢?她說我有還債的責任,但沒有還債的能力,家裡沒有錢.

今天我們基督徒都是這樣,有還債的責任,沒有還債的能力.

耶穌為我們死,我們要為主,耶穌赦免我們的債,我們要為主而活.有一句話,常常鼓勵我:耶穌為我們死,我們為祂而活,可是我們活不出基督,總是軟弱,總是憂愁,沒有能力去還債,還福音的債,還赦罪的債,引人歸主的債.

保羅說:我要誇我的軟弱,好叫基督的能力覆庇我.

二,    投靠耶和華,強似倚靠人!

我有還債的責任,沒有還債的能力.怎麼辦?找老闆,找神的僕人… 請他幫我禱告.

弟兄姊妹,20年前,有一個人來深圳,寫一封信回到家鄉,他說深圳的人情薄如紙,現在市場經濟,一切向錢看,沒有親情,沒有友情,沒有人情,凡事都講錢.這個寡婦去那裡找錢呢?我在想,人要找錢是難的.我跟女兒說,我們怎麼來發達呢?女兒說:每個人都想發達.

弟兄姊妹,我站講台站了三十多年,我對這世界還沒有完全失去,哇,這個寶馬,勞斯萊斯,高樓大廈,多少的吸引….

希伯來書有句話:我們都在肉身之內(來13:3).我們來表揚這個婦人,別的不求,別的不哭,找老師,找上司,找老闆.老闆啊!我的老公是你的學生,我老公是你教會的信徒,他為主做工,兩袖清風,現在負了這麼多的債,怎麼辦?你要幫幫我,你要幫幫我啊!

耶和華所賜的福沒有加上憂慮!

當你的人生,有困難的時候,可以找人幫你的忙,那個人幫到你的時候,這個人就成為你走天路的重擔和攔阻,你看到他就會覺得很累.寡婦倚靠 神,上帝幫助沒有負作用；　神的幫助不需要回報；你只要一顆感謝祂的心就夠了.

三,    欣賞擁有的!不埋怨缺乏的!

我後悔做個傳道人,後悔老公讀神學,後悔為主做工,後悔為他生二個孩子…婦人沒有後悔!神人啊!我欠你的錢,我老公一生愛主,你是知道的,現在貧窮,求你禱告　神幫幫我啊!以利沙問他:你家裡有什麼?他說家裡沒有甚麼?你想想.她說:哦!我家有小小一瓶油.以利沙說:你就拿這瓶油來感謝.

親愛的弟兄姊妹,我信主,從小就信主,事奉主35年,有很多抱怨,有很多苦毒,有很多不開心,有很多不滿意,有很多不公道,回顧我走過的一生.保羅說:愛不計算人的惡!

親愛的弟兄姊妹,有一次,聽個老師講道,他說:要欣賞你所擁有的,不想看你所缺乏的.我不知道擺在你們眼前的是什麼?也許你很富足,也許你很不開心,也許你四面受敵.保羅說:連活命的什麼都沒有了,心裡也是斷定必死的,然而欣賞你擁有的.當你欣賞你所擁有的,一定很開心.當你只看著你所缺乏的時候,你就埋怨,越埋怨,越埋怨.以色列走天路40年,最大的敵人就是埋怨,埃及有黃瓜,韭菜,房子,固定的…,在曠野什麼都不好,路也都不好,雲柱,火柱不好,衣服不會破不好,腳上的鞋不破不好,他們沒有欣賞擁有的!只看著自己所缺乏的.

有人問我,謝弟兄啊!你怎麼信　神?我說:我講不清楚.天堂我都沒去過,但是什麼是　神?我可以告訴你.那什麼是　神呢?　神就是很　神,你想不到,祂是　神；你做不到,祂是　神；復活顯現,祂是　神.

我們信主的人,有些人一生都很神奇,有很多恩典.今天,我們的主題:主的恩典夠用!主的恩典夠用!你欣賞你擁有的!你有一本聖經,我有一個女兒,有一個房產,有個教會,有個詩班,有鋼琴,要欣賞我們擁有的!　神會加恩!

四,    關上世界門,開啟恩典窗.

把世界的門關上,你回家把門關上,關外門,關裡門,與世隔絕,拿個瓶子來倒油.這個門關,那個窗開；那個窗關,這個門開,他們兩母子… 在外面車水馬龍,裡面是油廠的加工廠.

我在想… 這個油是怎麼倒出來的呢?應該有個無線的油接上去,越流越流不乾.

今天我們有困難,都在外面找老師,朋友,上司… 很少人回家,入內室關門.列王記下有句話:「以色列中沒有　神麼?」難道以色列中沒有先知嗎?為什麼你去找巴力西卜呢?

我們可以試一下,我們裡面有個小耶穌,一個屬靈的人,耶穌在我們裡面,一部份是自己的,一部份是耶穌的,當你自己工作的時候,耶穌在睡覺；當你叫醒耶穌的時候,我們的肉體會衰微.基督有豐盛的生命,關上門,開上窗,與世隔絕,叫醒我們裡面的小耶穌.當門徒在船上,狂風大浪的時候,說:我們喪命啦!他們忘記耶穌在船上,你去找狂風,你去找大浪,你越看越軟弱.

所羅門說:「看風的必不撒種．望雲的必不收割.」(傳11:4)今天,我們碰到什麼事情,老是看風,望雲.

五,    器皿沒有了,油就止住了.

器皿沒有了,油就止住了!油就止住了!器皿沒有了!我自己想,假如,這個婦人在倒油,孩子啊!拿器皿來,器皿都沒有了,油繼續在流,那這個油,會成為災難.

我跟你講兩件事情,91,92年,我在沙頭角開一間雜貨店,我們拿20箱可樂是一個價錢,拿50箱可樂是這個價錢,買多就會便宜!天氣熱,鐵皮房,可樂倒出來,哇!家中的床,被子,到處都充滿可樂,跟人家借了很多桶子來裝可樂,後來,聞到可樂就怕了,可樂本來是很好的嘛,太多了,成為災難了.我小的時候去釣青蛙,它很好吃,當埃及青蛙發難的時候,青蛙會是埃及的災難.

神的恩典,好好把握,這不是災難.有一次,我跟同工說:我們教會晚堂崇拜還沒坐滿,我們教會的禮拜,應該說有一千五百人左右,晚上也是差不多這樣的人,2,300人.我說弟兄啊!我們要發起福音運動,可以借器皿,空器皿怎麼借呢?這個地方這麼大,位子這麼多,還有很多空位,在十字路口,大街小巷,把人拉過來聽福音,裝滿油,借器皿；器皿沒有了,油就止住了；器皿有,油繼續流… 30人,300人,3000人,30000人.

親愛的弟兄姊妹,我只誇自己的軟弱,好叫基督的能力覆庇我,其他的別無可誇.我告訴你們,我懂的借空器皿.我95年來到藍山,找不到教會,弟兄,在自己的店裡,點支蠟燭….95年只有五個人做禮拜,晚上沒有電,就五個人.下個禮拜來了2個人共7個,下個禮拜9個,10個,30個,70個,120個,300個,700個,1500個,3000個.目前,我們教會信主的就超過2萬人.工作量很大.我說:主啊!有一天,你把一間教會交給我牧養,我們借空器皿,沒有場地,沒有教會,可以租.弟兄姊妹同心協力,我們有180個同工,七個唱詩班,有報紙,同聲傳釋,僕人之家,愛心大使….

我在想一件事情,今天上帝在中國做三件事情:

(1)預備人心.現在政府搞市場經濟,房地產,開發商.．．人富起來,但心靈空虛(你們有空就去深圳傳福音,現在是一個很好的機會).一百個家庭,當中有60個離婚,為什麼離婚呢?心中沒有　神,心中空虛,錢有了,愛心沒有了,人道沒有了.現在是傳福音的好機會.我們教會每個禮拜都有50個人來信耶穌…上帝在中國預備人心.

(2) 神在中國預備道路.唐崇榮,遠志明,張伯笠等人,中國政府不讓他們進去.但我在網路上,能聽的到唐博士的講道,人不能進,但網路進去了!現在網路是一個道路,人可以禁止,網路禁不了.

(3)借器皿.我一直在想,去禮拜堂,有很多困難,去禮拜堂,有很多攔阻.去體育館,運動場,沒有攔阻.因為那是政府工程.北京的鳥巢建好了,08年舉年了一次,我們再去北京看,鳥巢裡面是空的.我說:主啊!這個鳥巢拿來培靈真是太好了,在這裡講道太好了,那一天,你把全中國的運動場借給基督教,可以租,可以借….

弟兄姊妹,要借器皿,沒有場地就借場地；沒有人心,就借人心.

我跟你講一個人,大衛有一位祭司叫約押,後來所羅門起來要殺約押,在殺祭司的時候,所羅門說我不殺你,因為你當年幫我爸爸抬過約櫃,我爸爸借過你這個器皿,放你一條生路,我常常這樣子傳福音的.

我認識個老弟兄,78歲,他說,謝弟兄啊!我在家鄉建禮拜堂,禮拜堂建好之後,我就被人逼迫,誤會,冤枉,打擊,我就離家出走,來到深圳,一住就20多年,沒有感冒,生病,他說我一生做對一件事情,或許我衝動,任意…,我總把禮拜堂建起來,我們這個人,先借給 上帝用一下,裝滿油,耶和華因著那些油,會看顧那個人.

你們有沒有去過耶路撒冷,去耶路撒冷看,那裡有兩個城市,一個是耶路撒冷,一個是耶利哥.建造耶利哥的時候,耶利哥沒有聖殿,現在沒有耶利哥,移為平地,耶路撒冷有聖殿,大衛之星還在那裡.

弟兄姊妹,我們為主傳福音,借器皿,信主也好,沒有信主也好,我們請他幫我們做事,裝滿油. 神因著這個事工,忽然會揀選那個人.

六,    受聖靈充滿,不欠肉體的債.

油代表聖靈,這婦人把油倒滿了,聖靈充滿了,聖靈充滿了什麼都有.

一個姐妹,聽一個人傳福音,說:妳要信耶穌.她說:弟兄啊!信耶穌這麼好,信耶穌之後好不好打麻將?那個弟兄說:我去問牧師?過了六個月,這個姐妹受洗了.她對我說:謝弟兄啊!我當天問信耶穌之後好不好打麻將?過了六個月了,我麻將不打了,哦!那你連牧師也不需要問了.受聖靈充滿,就不欠肉體的債.

今天你沒辦法饒恕,放不下,你就欠饒恕的債；今天你苦毒,你沒辦法原諒,小氣,你欠苦毒的債；今天你撒謊,你欠撒謊的債；老是欺騙人,生活不榮耀神的名,欠債累累.五旬節聖靈充滿的時候,彼得站起來,滿有能力,不欠誰的債,我們欠債了嗎?我們求　神來預備,帶領我們,在神的愛中,受祂陶造,受祂恩惠,受祂祝福,受祂澆灌.

總結

今天,我們講主的恩典夠用,六句話:

一,    有還債的責任,沒有還債的能力.

二,    投靠耶和華,強似倚靠人!

三,    欣賞擁有的!不埋怨缺乏的!

四,    關上世界門,開啟恩典窗.

五,    器皿沒有了,油就止住了.

六,    受聖靈充滿,不欠肉體的債.

我講道有個特別,要講一個心得.今天,我們講寡婦這瓶油,有什麼心得?聽了這場道,有什麼體會?

今天在這裡有一句話,是我讀這聖經的心得.杯中是什麼,溢出是什麼.我自己在想,寡婦這瓶油又是什麼油?今天我真的很感謝　神!也很擔心,擔心什麼?擔心寡婦家的這瓶油.假設一瓶礦泉水,倒出來,不值錢.假如他家裡是瓶涼茶,倒出來,不是很好賣.假如家裡是一瓶垃圾,倒出來是垃圾.我問了解經家,他說這瓶是橄欖油.

親愛的弟兄姊妹,今天我們這個人,就是一個杯,你裡面裝了什麼?你裝了苦毒,倒出來就是埋怨.你心中充滿愁苦,倒出來是悲哀.

大衛心中有　神,說:耶和華是我的高臺!耶和華是我患難中的拯救!耶和華是我隨時的幫助!耶和華是我的讚美!我要讚美　神!

摩西四十年走天路,受盡委屈,摩西倒出來:「在　神的全家盡忠」(來3:2).

彼得三次不認主,後悔的時候,他講道30年,倒出來的,他說:「儆醒禱告」(彼前4:7).

我們用一段聖經作為結束.我們來讀哥林多後書4:7:「我們有這寶貝放在瓦器裡,要顯明這莫大的能力,是出於　神,不是出於我們.」

保羅說:我們這個杯,裝的是寶貝,碰出來的,就是寶貝.

弟兄姊妹,我們這個器皿,在新的一年,生命當中,多少碰碰撞撞,不滿意,缺乏,但是求主的寶貝放在瓦器裡,要顯明這莫大的能力.

你跟一個人,三天同吃同住,就知道這個人心中有什麼東西.有些人在服事當中,跟他在一起很煩,不是講東就是講西,不是講上就是講下.我說:弟兄啊!你能不能少說兩句呢?我跟你三天在一起,你杯中都是垃圾.信徒來教會,要聽寶貝,要受聖靈,在這裡不是聽垃圾.只有寶貝,才能甦醒人心!有寶貝,使人產生敬畏上帝的心!

寡婦家裡一瓶油,上帝藉著這瓶油來祝福她.使她這輩子無憂無慮.這瓶油,代表恩典；這瓶油,帶著寶貝；這個杯,代表從上天而來的恩澤.

弟兄姊妹,這瓶油就在我們裡面,願聖靈如膏油充滿我們每個弟兄姊妹.

親愛的主耶穌,我們感謝讚美您!主啊!話語是有限的,願主親自觸摸這個教會,使教會後來的榮耀,必大過先前的榮耀,主啊!我們倚靠你,你的恩典總是夠我們用,願主幫助我們在座的每個弟兄姊妹.耶穌,使這裡的牧者有力量來服事主,愛主更深,禱告祈求奉耶穌基督的名,阿們!



\newpage

\section{敬拜與讚美}
\label{sec:206}
\textbf{伍渭文牧師 - 敬拜與讚美}
\newline
\newline
\hyperref[sec:205]{< < < PREV SERMON < < <}
~
\hyperref[sec:toc]{[返目錄]}
~
\hyperref[sec:207]{> > > NEXT SERMON > > >}


講題:敬拜與讚美　Praise and Worship

經文:詩篇第22

講員:伍渭文牧師

場合:香港中文大學崇基學院禮拜堂

日期:2013年6月23日

有人說,有兩樣東西人是不能逃避的,就是你欠的稅 和死亡.欠稅可以有人代繳,或找稅貸緩急,找安信兄弟(財務公司廣告),但死亡則要自己單獨面對,別人不能幫忙.

傳說法國數學家,哲學家巴斯高(Pascal) 1662年臨終時說: 每人獨自死亡(One dies alone),願上帝不離棄我(May God never abandon me) .

被遺棄的孤單

詩篇22篇的作者論到他面臨死亡,要獨自面對死亡；在死亡面前,他感到被遺棄,被神遺棄,被人遺棄.上帝好像掩面不看他:「我的神,我的神!為甚麼離棄我?為甚麼遠離不救我,不聽我唉哼的言語?我的神啊,我白日呼求,你不應允；夜間呼求,並不住聲.」(vv.1, 2) 死亡是和生命的主隔絕,當死亡臨近時,生命的主竟然渺無蹤影!

他也被人所遺棄.惡人環繞他,嗤笑他:「但我是蟲,不是人,被眾人羞辱,被百姓藐視.凡看見我的都嗤笑我,他們撇嘴搖頭,說:『他把自己交托耶和華,耶和華可以救他吧!』」(詩22: 6-8)

中國人有一願望,希望落葉歸根.當年紀大的時候,很想回到自己長大的地方終老,我們很怕客死異鄉；臨走時,沒有熟悉的人伴隨,感到孤單.在以前或現在的農村地區,臨終時都在自己熟悉的家中,在長大的家中,由熟悉的人陪伴,渡過人生最後的時刻.

但隨著社會城市化,公共衛生水平提高,為了防禦疾病傳染,我們把病重的人安排在醫院留醫,希望得到最好和及時的照顧,但離開的人,要在陌生的環境離開.

其實當人面對死亡時,最需要是親友的陪伴.2002年在崇基學院教授哲學三十多年備受敬愛的陳特老師安息主懷,他寫的「生死徘徊十二年」說到親友的伴隨,最能幫助他面對頑疾的痛苦和對死亡的懼怕.

他說:「許多西方哲學家都說死亡是最個人也因此最寂寞最孤單的道路,但我至目前為止,卻總感到人間濃郁的情懷在陪伴著我,在死亡的路上,給了我許多勇氣和安慰.」「許多親友學生,不斷親自或以電話來問候,這些親友中有些是基督徒,為我禱告,我的妻子一年多來,一直陪伴在側.」

另一方便,他也感到上主的陪伴.「十二年來,面對著死亡與痛苦,我仍然努力使自己心靈平靜……我向上帝的禱告除了使身體恢復健康外,極重要的是使我心靈平安.」「我相信上帝在磨練我,不管肉體如何軟弱,他都在支持我……,如果不是他的恩典,我也不可能做到這些.」「就這樣,疏遠的上帝似乎又親近了,本來鬱悶的心情也開朗起來.」

死亡是一種隔絕

詩篇22篇給我們對死亡的透視:死亡是一種隔絕,和上帝隔絕,和人隔絕.而生命是與上帝和好,與人和好.

這詩的結構:分哀歌(lament)vv. 1-21b. ；應允(response), v. 21「你已經應允我,使我脫離野牛的角」及感恩 (thanksgiving), 受苦者宣告,vv.24-28；會眾宣告,29-32.

在詩篇的開始,詩人以哀歌體裁,向上主陳明自己的苦況:「我的神,我的神!為甚麼離棄我?為甚麼遠離不救我,不聽我唉哼的言語?我的神啊,我白日呼求,你不應允；夜間呼求,並不住聲.」(vv.1, 2)「我如水被倒出來,我的骨頭都脫了節,我的心在我裡面熔化.我的精力枯乾,如同瓦片；我的舌頭貼在我牙床上.你將我安置在死地的塵土中.犬類圍著我；惡党環繞我；他們扎了我的手,我的腳.」(vv. 14-16)

作者雖然四面楚歌,但沒有懷疑上帝賜給他生命的意義；詩人堅持生命是有意義:「但你是叫我出母腹的,我在母懷裡,你就使我有依靠你的心.我自出母胎就被交在你手裡,從我母親生我,你就是我的神.求你不要遠離我,因為急難臨近了,沒有人幫助我.」(vv. 9, 10. 11)

人若從現實世界看苦難,的確認為生命就是荒誕,約伯受苦時說:「願我生的那日和說懷了男胎那夜都滅沒.願那日變為黑暗,願神不從上面尋找它,願亮光不照於其上.願黑暗和死蔭索取那日,願密雲停在其上,願日蝕恐嚇它.願那夜被幽暗奪取,不在年中的日子同樂,也不入月中的數目」(伯3: 3-6)

但這詩篇是從信仰的角度看死亡的威脅,因詩22篇是公共禱文,是崇拜的禮文,所以在二十一節就有一個轉機:是應允-上帝宣告(oracle)- 「你已應允我,使我脫離野牛的角」,脫離前面所說:「有許多公牛四面困住我；它們向我張口,好像抓撕吼叫的獅子.」(v.12)

敬拜群體聆聽詩人的哀歌

跟著詩人感謝上主的拯救,詩篇從哀歌轉向感恩:「我要將你的名傳與我的弟兄,在會中我要讚美你!」(v. 22) 詩人要會眾和他一起讚美上主:「你們敬畏耶和華的人要讚美他!雅各的的後裔都要榮耀他!以色列的後裔都要懼怕他!」(v.23) 詩人得蒙上主拯救後,在大會中讚美:「我在大會中讚美你的話是從你而來的；我要在耶和華的人面前還我的願.」(v. 25)

會眾和應:「因為國權是耶和華的,他是管理萬國的.地上一切豐肥的人,必吃喝而敬拜；凡下到塵土中不能存活自己生命的人,都要在他面前下拜.」(vv. 28, 29)

我們留意,之前環繞他的對是一班譏笑他的人:「凡看見我的都嗤笑我,他們撇嘴搖頭,說:『他把自己交托耶和華,耶和華可以救他吧!』」(詩22: 7, 8) 但現在環繞他的,是一起讚美敬拜的會眾；這是一個大合唱,聲音洪亮.

但這班讚美敬拜的會眾,一直在會中保持靜默,專心在聆聽詩人的哀歌.事實上,在這詩篇,個人哀歌的篇幅比集體合唱更多,因為用愛心聆聽受苦的人在上帝面前的傾訴,就是對上帝的敬拜； 做在小子的身上就是做在上帝的身上.另一方面,詩人為何能夠盡情把苦況對上主陳明呢?因為他知道有一班愛他接納他的聆聽者.這是上主的會,這是耶和華的人.

敬拜群體和我們一起面對死亡的力量

肉體的死亡對我們一些人來說,可能比較遙遠,但死亡的力量,每日都壓迫著我們,我們每天都面對死線deadline.在這詩篇,詩人描述死亡的威脅非常生動:「有許多公牛四面困住我；它們向我張口,好像抓撕吼叫的獅子.」(v.12) 我們需要信仰群體一起支持我們去倆面對種種死亡的力量.

這星期我真的聽到獅子的吼叫.前幾天有機會和一對以前認識的夫婦協談,我徵得到他們的同意,分享他們正面對巴珊公牛的威脅.

約十天前凌晨一時,手機收到whatsapp一短信,有一位弟兄很想找我談話,我知這事情緊急,我立刻回覆:明天見吧；他即刻說:感謝上帝,認為是上帝聽他祈禱.作為丈夫父親的他最近在家中跌倒,腦頭受創,內有積血,要進行手術,手術後希望做針砭幫助復康,不料針砭之後手指感到有點麻痹,自然心慌,跟著失眠；沒有足夠的睡眠,健康愈來愈差.他朋友不多,也很久沒有返教會了,因為信任才找牧師談話,說出自己的懼怕.

一星期後他和太太一起見牧師.這對夫婦長期因為教導兒女方法不同,關係愈來愈差,沒有溝通,也不想溝通.要升中學的女兒成績不太好,掛心她不能找到好的中學,而且情緒不穩定,要看心理醫生.說到這裡,太太的眼睛潤濕.我看到巴珊大力的公牛,正張口正要把他們一家撕裂.

我和他們看這詩篇,叫他們不要懼怕死亡的力量,因為上主搭救我們,上主已應允:「使我脫離野牛的角.」上主可以使枯骨成為活人,當聖靈像風吹過,枯骨成為大有能力的軍隊.我們一起祈禱,也鼓勵他們一起讀經祈禱.面對死亡般的威脅最重要是信心:「從我母親生我,你就是我的神.求你不要遠離我,因為急難臨近了,沒有人幫助我.」

其實第一次見了這弟兄,我鼓勵他返教會崇拜,不可停止聚會,我們需要信仰群體的支持,激勵信心.跟著的主日,和太太見我之前,他回到住所鄰近的教會聚會,並發給我的崇拜程序表一首詩歌:我們愛,讓世界不一樣,用顏色劃下其中一句感動他的話「心再堅強也不要獨自飛翔」.是的,不可停止聚會,心再堅強也不要獨自飛翔.

崇基教堂因為在大學校園,一路以來的崇拜講道都注重信仰與知識的結合,信以求知(faith seeking understanding) ,1984-85年特別邀請學有專長的基督徒學者作專題講道,並結集出書,稱為大學講章.但我對崇基教堂有另一個期望,這教堂也是一個敬拜讚美的地方,特別我們教堂有一坐這麼美好的管風琴,提醒我們這裡是敬拜讚美的地方,在這裡,我們的信心得到堅立；在這裡我們可以對上帝傾訴心中的苦情,而有人聆聽,作為對上帝的敬拜；我們有人陪伴支持,一起面對人生種種橫逆.

其實活潑的敬拜有治療效用.在敬拜讚美中,當我們來到至高上主面前,我們從至高之處或在神舟太空艙俯瞰大地,地球不過是不同顏色的平面圖案,在地上我們彼此比較,爭過你死我活,但從至高上主的角度看云云眾生,都是按他的形像所造,同樣尊貴的人.

敬拜與讚美,甚具療效,藉著詩歌,我們誘發出深藏的情緒,認罪使我們面對自己,宣赦使我們得到接納,奉獻時教我們凡事感恩,代求使我們關懷,差遣讓我們有承擔,聽道使我們反思生命,得到聖經的智慧.

基督經歷死亡救我們脫離死亡

我想當我讀出詩人的哀歌時,心水清的人一定想起這詩篇就是耶穌的真實寫照.詩篇二十二篇的作者,呼救上帝免去他的死亡,上帝就撘救他.耶穌有沒有求上帝免去他的死呢?「我父啊!倘若可行,求你叫這杯離開我；然而不要照我的意思,祇要照你的意思.」(太26: 39) 耶穌沒有,他走上十字架.

「他把自己交托耶和華,耶和華可以救他吧.」(v.8) 這是祭司長,文士長老甚至同釘十架的強盜都譏笑耶穌的話.「他依靠神,神若喜悅他,現在可以救他.」(太27: 43) 「他們分我的外衣,為我的裡衣抽拈」(v. 18；約19: 24) 這是釘耶穌兵丁所為.「我的神!,我的神!為甚麼離棄我?」更是著名的十架七言.所以詩篇二十二篇又稱為第五福音書,詳盡預言耶穌如何被父上帝離棄,被人離棄,為人的罪經歷死亡.

基督救贖我們成為敬拜讚美的群體

但基督為何基督救贖我們?為何上主要摩西把百姓領出埃及?就是到西乃山領受律法作上帝的子民,作敬拜讚美的群體.馬丁路德說:得著上主,就是要敬拜上主.在敬拜讚美中,我們奉耶穌的名字來到父面前,我們奉耶穌的名開始聚會,奉基督的名字禱告；在基督裡,我們來到父上帝面前.基督「在聖所,就是真帳幕裡作執事」(來8:2) ,執事原文是liturgos(禮儀) .希伯來書作者說:「因那使人成聖的,和那些得以成聖的,都是出於-；所以他稱他們為弟兄,也不以為恥,說(引詩22: 22):「我要將你的名傳與我的弟兄,在會中我要讚美!」(來2: 12)

上帝的兒子成為人,以致人可以成為上帝的兒子；耶穌和我們是弟兄的關係,基督是長子,他首先復活,將來我們也會隨他一樣復活.所以基督在那裡,教會也應該在那裡,信仰群體就是基督的代表.若從這角度看緊隨的詩篇二十三篇,就會說:因著這個信仰群體的伴隨,我們不怕死亡力量的威脅,我們可以在敵人面前,擺設筵席,敬拜讚美.我們雖然行過死蔭的幽谷,也不怕遭害,因為信仰群體的伴隨我們,與我們同在.他們的杖,他們的杆,都安慰我們.阿們.



\newpage

\section{溫柔與嚴厲}
\label{sec:207}
\textbf{鄧瑞強博士 - 溫柔與嚴厲}
\newline
\newline
\hyperref[sec:206]{< < < PREV SERMON < < <}
~
\hyperref[sec:toc]{[返目錄]}
~
\hyperref[sec:208]{> > > NEXT SERMON > > >}


路加福音 9:51-62
\newline
\begin{longtable}{cl}
\hline
\hline
章節 & 經文 (和合本修訂版)\\
\hline
9:51 & \begin{tabularx}{0.7\textwidth}{X} 耶穌被接上升的日子將到，他決定面向耶路撒冷走去。 \end{tabularx} \\ \\ \relax
9:52 & \begin{tabularx}{0.7\textwidth}{X} 他打發使者在他前頭走；他們進了撒瑪利亞的一個村莊，要為他作準備。 \end{tabularx} \\ \\ \relax
9:53 & \begin{tabularx}{0.7\textwidth}{X} 那裡的人不接待他，因為他面向著耶路撒冷去。 \end{tabularx} \\ \\ \relax
9:54 & \begin{tabularx}{0.7\textwidth}{X} 他的門徒雅各和約翰看見了，就說：「主啊！你要我們吩咐火從天上降下來，燒滅他們嗎？」 \end{tabularx} \\ \\ \relax
9:55 & \begin{tabularx}{0.7\textwidth}{X} 耶穌轉身責備兩個門徒。 \end{tabularx} \\ \\ \relax
9:56 & \begin{tabularx}{0.7\textwidth}{X} 於是他們就往別的村莊去了。 \end{tabularx} \\ \\ \relax
9:57 & \begin{tabularx}{0.7\textwidth}{X} 他們在路上走的時候，有一個人對耶穌說：「你無論往哪裡去，我都要跟從你。」 \end{tabularx} \\ \\ \relax
9:58 & \begin{tabularx}{0.7\textwidth}{X} 耶穌對他說：「狐狸有洞，天空的飛鳥有窩，人子卻沒有枕頭的地方。」 \end{tabularx} \\ \\ \relax
9:59 & \begin{tabularx}{0.7\textwidth}{X} 他又對另一個人說：「來跟從我！」那人說：「主啊，容許我先回去埋葬我的父親。」 \end{tabularx} \\ \\ \relax
9:60 & \begin{tabularx}{0.7\textwidth}{X} 耶穌對他說：「讓死人埋葬他們的死人，你只管去傳講神的國。」 \end{tabularx} \\ \\ \relax
9:61 & \begin{tabularx}{0.7\textwidth}{X} 又有一人說：「主啊，我要跟從你，但容許我先去辭別我家裡的人。」 \end{tabularx} \\ \\ \relax
9:62 & \begin{tabularx}{0.7\textwidth}{X} 耶穌對他說：「手扶著犁向後看的人，不配進神的國。」 \end{tabularx} \\ \\
[1ex]
\hline
\hline
\end{longtable}
講題:溫柔與嚴厲　Tender and Straight

經文:路加福音9:51-62

講員:鄧瑞強博士

場合:香港中文大學崇基學院禮拜堂

日期:2013年6月30日

各位弟兄姊妹,早安.

魯迅的《自嘲》詩有兩句很出名的句子,「橫眉冷對千夫指,俯首甘為孺子牛」.若你身邊有一兩個人講你壞話,你已經煩惱不堪.若被千夫所指,則肯定苦不堪言.若要「橫眉冷對千夫指」,你想像一下,這需要多強的心力.面對千夫所指而仍勇往直前,必定是抓住一個重要的生命目標,或者說,被一個重要的生命目標抓住.這是一個生命的強者.生命的強者可成就大善,也可成就大惡,這視乎他對別人抱什麼心態.「俯首甘為孺子牛」就是對別人抱著良善的心.面對一個孩童,面對一個無知的弱者,竟能俯首,竟能俯就卑微,竟能像牛一樣讓他騎,竟能放下身段而為他洗腳,這種良善的心使強者成就大善.「橫眉冷對千夫指」而沒有善心的人,可以像坦克車,見人輾人.「橫眉冷對千夫指」卻能「俯首甘為孺子牛」的人,就像人力車,沒有殺傷性,卻能將你載到彼岸.

以上講的,是解釋今日的講題:「溫柔與嚴厲」.溫柔,指的是對別人的心.嚴厲,指的是對生命目標的堅定不移,具有「雖千萬人,吾往矣」的勇猛.

主耶穌正是一個溫柔而嚴厲的人.

今日,我們就是要看看主耶穌如何展示溫柔與嚴厲.我的講道經文,是今日的福音書經課:路加福音9:51－62.

路加福音9:51－62

9:51　耶穌被接上升的日子將到,他就定意向耶路撒冷去,

9:52　便打發使者在他前頭走.他們到了撒馬利亞的一個村莊,要為他預備.

9:53　那裡的人不接待他,因他面向耶路撒冷去.

9:54　他的門徒雅各,約翰看見了,就說:「主啊,你要我們吩咐火從天上降下來燒滅他們,像以利亞所做的嗎?」

9:55　耶穌轉身責備兩個門徒,說:「你們的心如何,你們並不知道.

9:56　人子來不是要滅人的性命,是要救人的性命.」說著就往別的村莊去了.

9:57　他們走路的時候,有一個人對耶穌說:「你無論往哪裡去,我要跟從你.」

9:58　耶穌說:「狐狸有洞,天空的飛鳥有窩,只是人子沒有枕頭的地方.」

9:59　又對一個人說:「跟從我來!」那人說:「主,容我先回去埋葬我的父親.」

9:60　耶穌說:「任憑死人埋葬他們的死人,你只管去傳揚神國的道.」

9:61　又有一人說:「主,我要跟從你,但容我先去辭別我家裡的人.」

9:62　耶穌說:「手扶著犁向後看的,不配進神的國.」

經文不長,卻可看到主耶穌的溫柔與嚴厲.

這段經文的開首,意義深長.

9:51　耶穌被接上升的日子將到,他就定意向耶路撒冷去.

「被接上升的日子」當然指到主耶穌復活升天的日子,也是指他的生命得享榮耀的日子.這也可以被理解成他的生命意義得以完滿實現的日子.主耶穌意識到人生在世是有一重大意義要實現的.這種意義感此刻壟斷了他的整個生命.這種意義感使他有一種明確的生命方向.這種意義感也使他冒險地奔向不測的命途.經文說,「他就定意向耶路撒冷去」.所謂「定意」的原意,大部分的英文聖經都翻譯出來,就是「堅決地固定自己的面」向著一個方向.用中文的講法,大致是「硬頸」地,「死牛一面頸」地向著一個方向.經文這裡說,主耶穌「死牛一面頸」地向著耶路撒冷去.「耶路撒冷」意味著十字架,受苦,死亡.主耶穌意識到生命的終極意義,這意義感令他有一種「雖千萬人,吾往矣」的勇猛.我們當中有多少人,能被生命的終極意義所抓住,以致能堅定地勇往直前,連死亡都不懼.

現代社會都訓練很多勇往直前的人.要發達的欲望,令很多人勇往直前,甚至像坦克一樣,將阻著前進的人輾在車底.人本心理學也強調人要勇往直前,實現自己的潛能.問題只是,這種自我實現會否神化了自己,消滅了他人?希特拉的自我實現,就消滅了很多人.

「死牛一面頸」地勇往直前,可以是極危險的.

主耶穌的勇往直前,是直向耶路撒冷,是直向十字架.基督教信仰的一個核心,就是:主耶穌並他釘十字架.若說這是一種自我實現,這卻是一種很特殊的自我實現.這種自我實現,是放下了自我,甚至讓自我甘願被消滅,為的是他人的生命意義得到實現,為的是他人得到寬恕和希望.

勇往直前的人是會面對別人的排斥和難阻的,總會面對「千夫指」.

9:52　便打發使者在他前頭走.他們到了撒馬利亞的一個村莊,要為他預備.

9:53　那裡的人不接待他,因他面向耶路撒冷去.

主耶穌面向耶路撒冷,就有人正正因他面向耶路撒冷而不接待他.「不接待」就是不接納,不同意,不歡迎.這裡提到「撒馬利亞」.「撒馬利亞」人與猶太人有一種緊張的關係.「撒馬利亞」人有自己的信仰習慣,有自己的聖經傳統,有自己的敬拜中心,這敬拜中心在基利心山上.若有人與你持不同真理觀點,不同生活方式,不同的信仰,則他很可能不接納你,不歡迎你.

面對這種不同,我們採取什麼態度?

魯迅說,「橫眉冷對千夫指」.主耶穌的門徒又抱什麼態度呢?

9:54　他的門徒雅各,約翰看見了,就說:「主啊,你要我們吩咐火從天上降下來燒滅他們,像以利亞所做的〔有古卷沒有「像以利亞所做的」這幾個字〕嗎?」

雅各和約翰的反應,是人類常見的反應.當認為自己擁有真理,又認為自己擁有神一樣的能力的時候,面對反對自己的人,最簡單的處理方法,就是殲滅他們.雅各和約翰像坦克車,要輾過一切擋在前面的人.

這是暴力政治.施行暴力政治的人,手裡有的是權力.有權力也不能隨意運用,必需有自己論述的真理去支撐權力的運用.有了真理的話語權,要消滅對方,只需要給對方一個污名,這樣,運用權力消滅對方便找到無懈可擊的理由.只要對方的污名是「邪惡軸心」,或「反革命」,或「不愛國」,對方便是屬於魔鬼的,屬於邪惡的,屬於反真理的.殲滅魔鬼,人人義不容辭.在這時,不努力殲滅所謂「邪惡」的人,就是屬於「邪惡」的.

何等的暴力,假真理的名而行.主耶穌的門徒,竟就是這群假真理之名而行的暴徒.我們生命的深處,是否有這種暴力的種子?

面對暴力的騷動,主耶穌如何回應呢?

9:55　耶穌轉身責備兩個門徒,說:「你們的心如何,你們並不知道.

9:56　人子來不是要滅人的性命,是要救人的性命.」說著就往別的村莊去了.

主耶穌責備這兩個門徒.「責備」這用字,與主耶穌「斥責」污鬼時「斥責」那個字是同一個字.主耶穌斥責他們,像斥責污鬼一樣.主責備他們:「你們的心如何,你們並不知道.」主的意思是說:「你們知不知道自己在做什麼?」「你們知不知道自己在想什麼?」「你們知不知道自己到底是誰?」「到底,你們知不知道什麼叫做基督徒?」

主耶穌「死牛一面頸」地去完成的目標,不是要消滅一切反對力量而達致的目標.主說:「人子來不是要滅人的性命,是要救人的性命.」一句簡單的話,點出了基督教信仰講的「救贖」.不是你擁有「真理」便可審判反對你的人「違背真理」,真正的真理不能被人擁有,真理只在你去學習愛一個反對你的人的時候呈現出來.不是你擁有「權力」便可以用權力去壓制反對你的人,真正的權力在於你放下權力,反倒虛己,為了他人的福祉而自己背負更大的責任.最大的力量就是「溫柔」.

上個月的《時代論壇》(2013年5月26日),社長李錦洪的＜心耕手記＞提到意大利哲學家博比奧(Norberto Bobbio)的一本小書《讚美溫柔》(In Praise of Meekness),這書講到:「人類的德行發展,都是陽者的強性道德,以征服,勝利與支配為主,這種德行觀其實是錯的,它應有一個對立面,那就是溫柔.她是一種非力量的道德,以助人,幫人的善為目標.」

基督教倫理學者侯活士(Stanley Hauerwas)及國際方舟團體的創辦人范尼雲(Jean Vanier)合寫了一本小書《暴力世界中的溫柔》(Living Gently in a Violent World).方舟團體是照顧智障人士的團體.一般人可能認為智障人士是異類.方舟團體認為,如何與智障人士一起生活正正考驗我們對待與我們不同的人的方式,也考驗整個社會的公義實踐.方舟團體「接待」智障人士,按他們的光景接納他們,與他們平等友愛地生活.《暴力世界中的溫柔》這本小書講到,世上好心腸的倫理學者及政治學者很努力的找一些理論,去肯定智障人士的權利,但是,找來找去,卻找不到一種理論,能說服人去和智障人士一起生活,讓他們在平等友愛的生活中體驗生命的完滿和救贖.智障人士最需要的,偏偏正是有人願意去和他們平等友愛地生活.這就是以溫柔成就的正義.

智障人士是與一般人不同的弱者,雅各和約翰說,殲滅他們.主耶穌卻責備他們.人來到這世界,到底是追求什麼?我們是否知道自己是誰?我們是否知道作一個基督徒意味什麼?主耶穌說:「人子來不是要滅人的性命,是要救人的性命.」主耶穌願意溫柔地與智障的人平等友愛地生活,主耶穌願意「俯首甘為孺子牛」.這是主耶穌溫柔的一面.

跟著,主耶穌與不同的人討論什麼叫做「跟從他」.

9:57　他們走路的時候,有一個人對耶穌說:「你無論往哪裡去,我要跟從你.」

9:58　耶穌說:「狐狸有洞,天空的飛鳥有窩,只是人子沒有枕頭的地方.」

9:59　又對一個人說:「跟從我來!」那人說:「主,容我先回去埋葬我的父親.」

9:60　耶穌說:「任憑死人埋葬他們的死人,你只管去傳揚神國的道.」

9:61　又有一人說:「主,我要跟從你,但容我先去辭別我家裡的人.」

9:62　耶穌說:「手扶著犁向後看的,不配進神的國.」

狐狸的洞,飛鳥的窩,是提供安全感的地方.但是,要達到主耶穌所講的那種生命目標,那種十字架的目標,那種救贖他人的愛的目標,就必須放下這種安全感了.坐在家中「歎冷氣,看電視」是安全的,但要為正義而行動,則便失去安全感了.

埋葬父親,是當行的義.主耶穌卻說:「任憑死人埋葬他們的死人,你只管去傳揚神國的道.」很多解經家為了使這句話變得合乎常理,便指所謂「死人」便是那些沒有永生的人,就讓那些沒有永生的人去安葬死去的父親吧.這種解釋是可能的.但主耶穌的意義可能很簡單,就是說:「死者已矣,你還是為天國的目標努力前進吧.」安葬父親,其實是天大的理由,讓人放下一切,儘管是天國的事情.但主耶穌卻沒有讓步,在世上,並沒有天大的理由,可以讓我們放下天國的事情.這是否不近人情?這是否很嚴厲?是的,耶穌是嚴厲的.有些事情,是需要我們「死牛一面頸」地去完成.否則,將是一事無成.

誠實地說,今日我們都輕忽信仰,視之為可有可無,遺忘了信仰的嚴厲要求.在工作場所裡,我們有天大理由不近人情.在對待我們不喜歡的人時,我們有天大理由發脾氣.在個人的追求上,我們有天大理由只顧自己,不理他人.我們有的是天大理由,這樣做是天公地道.但主耶穌說,這一切都不成理由.若你明白信仰為何,若你明白生命的終極意義,若你明白天父的心,你應明白,這一切理由都不成理由.沒有任何理由讓我們忘記基督徒的本分.

這個道理,初戀的人最明白.任何阻止你去見愛人的理由,都不成理由.你的意念只有一個,見你的愛人.難怪雅歌說:「愛情,如死之堅強」.主耶穌要求我們如此堅定地實踐我們的信仰.

有人說:「主,我要跟從你,但容我先去辭別我家裡的人.」主耶穌卻說:「手扶著犁向後看的,不配進神的國.」在農業社會,以犁耕地,生產糧食,是生命之所繫.主耶穌卻說,要進入神的國,你要有新的家庭關係,你要有新的生命所繫.你要視每個人為兄弟姊妹,視一切來自神的恩典.這樣,你才會學曉愛,學曉分享.

跟從主耶穌,是一件嚴肅的事.主的要求很嚴厲.

結語:在對待他人這事上,主耶穌是溫柔的,他渴望我們善待與我們為敵的人.在面對自己的天國責任時,主耶穌是嚴厲的,他渴望我們以專一剛毅的心,完成這信仰的人生.

但願:榮耀歸於聖父,聖子,聖靈.阿們.



\newpage

\section{機不可失}
\label{sec:208}
\textbf{賀志勇博士 - 機不可失}
\newline
\newline
\hyperref[sec:207]{< < < PREV SERMON < < <}
~
\hyperref[sec:toc]{[返目錄]}
~
\hyperref[sec:209]{> > > NEXT SERMON > > >}


路加福音 10:1-12
\newline
\begin{longtable}{cl}
\hline
\hline
章節 & 經文 (和合本修訂版)\\
\hline
10:1 & \begin{tabularx}{0.7\textwidth}{X} 這些事以後，主另外指定七十二個人，差遣他們兩個兩個地在他前面，往自己所要到的各城各地去。 \end{tabularx} \\ \\ \relax
10:2 & \begin{tabularx}{0.7\textwidth}{X} 他對他們說：「要收的莊稼多，做工的人少。所以，你們要求莊稼的主差遣做工的人出去收他的莊稼。 \end{tabularx} \\ \\ \relax
10:3 & \begin{tabularx}{0.7\textwidth}{X} 你們去吧！看！我差你們出去，如同羔羊進入狼群。 \end{tabularx} \\ \\ \relax
10:4 & \begin{tabularx}{0.7\textwidth}{X} 不要帶錢囊，不要帶行囊，不要帶鞋子；在路上也不要向人問安。 \end{tabularx} \\ \\ \relax
10:5 & \begin{tabularx}{0.7\textwidth}{X} 無論進哪一家，先要說：『願這一家平安。』 \end{tabularx} \\ \\ \relax
10:6 & \begin{tabularx}{0.7\textwidth}{X} 那裡若有當得平安的人，你們所求的平安就必臨到那家，不然，將歸還你們。 \end{tabularx} \\ \\ \relax
10:7 & \begin{tabularx}{0.7\textwidth}{X} 你們要住在那家，吃喝他們所供給的，因為工人得工錢是應當的；不要從這家搬到那家。 \end{tabularx} \\ \\ \relax
10:8 & \begin{tabularx}{0.7\textwidth}{X} 無論進哪一城，人若接待你們，給你們擺上甚麼食物，你們就吃甚麼。 \end{tabularx} \\ \\ \relax
10:9 & \begin{tabularx}{0.7\textwidth}{X} 要醫治那城裡的病人，對他們說：『神的國臨近你們了。』 \end{tabularx} \\ \\ \relax
10:10 & \begin{tabularx}{0.7\textwidth}{X} 無論進哪一城，人若不接待你們，你們就到大街上去，說： \end{tabularx} \\ \\ \relax
10:11 & \begin{tabularx}{0.7\textwidth}{X} 『就是你們城裡的塵土黏在我們的腳上，我們也當著你們擦去。但是，你們該知道神的國臨近了。』 \end{tabularx} \\ \\ \relax
10:12 & \begin{tabularx}{0.7\textwidth}{X} 我告訴你們，在那日子，所多瑪所受的，比那城還容易受呢！」 \end{tabularx} \\ \\
[1ex]
\hline
\hline
\end{longtable}
主題:機不可失　The Chance We Should Never Miss

經文:路加福音10章1-12節

講員:賀志勇博士

場合:香港中文大學崇基學院禮拜堂

日期:2013年7月7日

1 這事以後,主又設立七十個人,差遣他們兩個兩個的,在他前面往自己所要到的各城各地方去. 2就對他們說,要收的莊稼多,作工的人少．所以你們當求莊稼的主,打發工人出去收他的莊稼. 3 你們去罷．我差你們出去,如同羊羔進入狼群. 4 不要帶錢囊,不要帶口袋,不要帶鞋．在路上也不要問人的安. 5 無論進那一家,先要說,願這一家平安. 6 那裡若有當得平安的人,〔當得平安的人原文作平安之子〕你們所求的平安就必臨到那家,不然,就歸與你們了. 7 你們要住在那家,喫喝他們所供給的．因為工人得工價,是應當的.不要從這家搬到那家. 8 無論進那一城,人若接待你們,給你們擺上甚麼,你們就喫甚麼. 9 要醫治那城裡的病人,對他們說,　神的國臨近你們了. 10 無論進那一城,人若不接待你們,你們就到街上去, 11 說,就是你們城裡的塵土,粘在我們的腳上,我們也當著你們擦去．雖然如此,你們該知道　神的國臨近了. 12 我告訴你們,當審判的日子,所多瑪所受的,比那城還容易受呢.

一,引言:福音的稻田

12年前我來香港時,一間學校不到200個內地生,全港的大學估計不超過3000名內地生.今天,全港的內地生已經超過了5萬,光中文大學,內地生人數就已經超過4000.我是來到香港之後信主的,剛到香港,因為不會廣東話,就參加了校牧室舉辦的廣東話班,然後參加了校牧室和我的母會聯合舉行的查經班,終於在查經班歸信了主耶穌.到現在為止,我都清楚記得那段日子,也非常感恩那些幫助我信主的弟兄姊妹.

前兩個星期我去了南山教會講道,12年前我曾經去過那裡,那裡一片荒涼,還在開拓,現在是高樓林立,樹木成蔭,並且高樓中還隱藏著一家一家的教會.我當時真的非常激動,說「主啊,你真是在沙漠開江河,在曠野開道路的神!」

這是一個全新的時代,這是一個充滿著機會的時代,不僅僅是從經濟的浪潮來講,更是從福音的浪潮來講.今天的香港,已經成為了一個強有力的宣教基地,派出了一批一批的宣教士去歐洲,去大陸宣教,還開了不少培訓班去教導大陸的傳道人--崇基神學院就是一個非常好的例子；同時,今天的中國大陸,也是一個非常廣闊的禾場,方言一望,到處都是福音的莊稼,這能不讓我們大大地心動,大大地感動麼?

二,解經

一．莊稼之主的差遣(1-2節)

1-2節描繪了一幅圖畫,就是收割莊稼的圖畫.在舊約或者兩約之間的典籍,或者新約聖經中,收割莊稼大多數都是指末日的審判.比如何西阿書6:11節,以賽亞27章 11-12節,都是指莊稼成熟了,人們要把麥子跟稗子分開來,也就是說當神再來的時候,有一部分人會得救,蒙上帝祝福；有一部分人卻會被定罪.農夫要趕緊收割,是因為愛惜自己的勞動,愛惜自己的莊稼,生怕還沒收割,一場雨來了,就沒有收成了；而我們今天要去傳福音,同樣是愛惜我們的親人朋友和身邊的鄰舍,希望他們在審判來臨的時候,他們跟我們一樣,站立得住,可以享用天國的福氣.所以,如果你還有親戚朋友沒信主,而你又很愛他們,你會不會著急呢?

當年的耶路撒冷跟今天一樣,也有非常多的機會.當年的耶穌,也如同今天的我一樣,看到了一個非常大的禾場.但問題不在於硬件,而是工人!工人很重要,一間教會,一家大學,一間公司,哪個不需要合適的員工呢?耶穌向當年的門徒提出了一個挑戰,他打發他們趕緊去收割.熟悉農作物的朋友可能知道,收割是不等人的.莊稼都已經熟了,不等人的,你說「哎呀,我今天傷風了,下個禮拜再動手吧,」或者說「哎呀,我還要出去旅行!回來再收割啊!」那你可能血本無歸.尤其是種蘭花賣的人,等一個禮拜,等到花都謝了,那就真是哭都來不及了.

二．差遣的智慧

所以耶穌基督趁著時機,打發工人出去收割了.從主耶穌的安排中,我們看到了許多非常有智慧的教導:

第一,    1-2節講主耶穌設立70個人,打發他們兩個兩個出去.為什麼是70個人?有的抄本說是72個人,我就不多加分析了.之所以提到70個人,是因為在民數記十一章中,上帝吩咐摩西從以色列中選了70個長老分擔他的事奉.當時以色列民老是抱怨這,抱怨那,摩西不勝其煩,於是神就讓他選了70個長老,降了聖靈在他們身上,跟摩西一起事奉.那為什麼主耶穌要打發兩個兩個出去呢?因為一個人出去,沒有同工,就非常麻煩.我當年一個人坐火車出去旅行,連廁所都不敢上,因為怕我要去上廁所了,我的箱子被人拿走了怎麼辦!兩個人出去,可以互相守望,還可以互相提醒.遇到危險可以互相救助,遇到問題可以商量解決,有不對的地方有人指正,有低落的時候有人鼓勵--這就是同工.什麼事情都需要靠自己單打獨鬥,從來沒有人鼓勵你,也沒人批評指正你,那你就有禍了!做生意如此,做學問如此,做福音事工更是如此.有沒有單打獨鬥的事奉啊?覺得自己厲害,你們都不用來,我自己一個人就可以搞定了,如果是這樣子的事奉,一定會失敗.

其實無論是七十個人也好,還是兩個兩個派出去也好,都體現了主耶穌讓我們合一,讓我們一起同工的要求,並且他讓我們明白,主耶穌自己才是莊稼之主!現在有的教會或者有的教友,片面抬高某個牧師,好像這個牧師有什麼了不得的屬靈恩賜,比其他教友要高明很多,那就容易把收割工人當成莊稼之主了.記住,主耶穌才是莊稼之主,他讓我們一起事奉,兩兩出去,這是主耶穌的智慧.因為我們只有在彼此同工中,才能在屬靈上學習忍耐,學習恩慈,才能一起成長,一起成熟!

第二,3-4節,主耶穌引導我們進入第二個教導,他說「我差你們出去,如同羊羔進入狼群. 4 不要帶錢囊,不要帶口袋,不要帶鞋.在路上也不要問人的安.」一方面,這裡的教導來自舊約,主耶穌說不要帶鞋,是因為「逾越節的晚上,耶和華吩咐以色列人穿好鞋子吃羊羔,隨時準備離開埃及.」(出 12章)主耶穌說不要問人的安,是因為先知以利亞由於傳上帝話語的緊迫,吩咐僕人不要路上停留,向人問安.從這兩句可以看出,傳福音的使命是何等地緊迫啊!

但主耶穌的教導不止於此,他要告訴我們更深奧的屬靈的智慧.他告訴我們,當我們基督徒面對世界的時候,我們可能就像羊羔進入狼群一樣,我們很弱小,我們很危險.進入狼群的時候,除了主耶穌之外,我們別無倚靠!我們需要不斷提醒自己,神哪怕在最危險的時候都會眷顧我們.沒有這個盼望,我們就沒有勇氣!主耶穌還提醒我們,我們不要帶錢包,帶旅行箱,帶鞋子,因為我們基督徒面對這個世界時,不是去旅行,不是去觀光,我們是去探險的.這個世界充滿了誘惑和危險!狼群是外在的危險,但是我們內心對錢財˴對物質的倚靠,乃是內在的危險.信仰乃是一次未知的探險,隨時有地雷,有冷箭.但主耶穌告訴我們,在我們橫跨一生的信仰探險中,我們不能再倚靠錢包,行李箱,鞋子--那些從前我們以為是可以倚靠的東西,我們只能培養一種德性:對莊稼之主的倚靠.

以前我做學生事工累了的時候,我就想:要是我很有錢就好了,我可以請多幾個同工一起來事奉,請幾個工人做飯給學生吃,幫他們找工作,這樣學生一定會願意來聚會了.等我成熟一點,我就明白了:如果我真的這麼有錢,我還會做傳道人麼,那可不一定!而且,我真的提供飯食給學生,幫學生找工作,他們就一定來聚會了麼?不是!真正吸引他們來到這裡的是主的福音,是主賜下的平安!當我想倚靠那些東西來傳福音的時候,其實是我自己對福音的大能缺乏信心了.

第三,5-9節,耶穌不停提到一個字,就是「家」.耶穌的宣教策略,就是把事奉的中心從聖殿轉向了家庭,他要他的門徒先進入家庭,而不是進入會堂傳講神的平安.主耶穌是要提醒我們,神不僅僅是在教堂出現,他沒有被圍困在四面牆當中,,他更是我們的一家之主.

今天,我們在家裡面有沒有家庭崇拜呢?如果家裡還有沒信主的人,那是不是我們的莊稼禾場呢?我們的敬虔是不是僅僅表現在教堂裡,到了家裡我們就換了另一個樣子呢?這些問題都值得我們深思!

正因為家庭如此重要,所以主吩咐他的門徒,要一開始就進入家庭宣講福音,要祝福這一家的平安.平安就是人與神和好,人與人和好.也就是說,門徒很重要的任務,就是把和好帶到這個家庭,讓這個家庭從神那裡學會彼此和好,也跟神和好.這是非常不容易的工作.因為有時候原諒外人容易,原諒自己的家人卻不容易啊.中國人講「家和萬事興」,福音工作也是如此,如果福音能在一個一個家庭中傳播˴興旺,那就能真正興旺起來.

向家庭傳福音,為什麼耶穌又提到食物的問題呢?耶穌要門徒住到別人家裡去,別人擺上什麼,就吃什麼,這跟傳福音又有什麼關係呢?其實食物對於以色列人來講,是一個非常非常重要的問題.猶太人吃食物有潔淨的律例,比方說不吃雞,所以炸雞一定不吃；奶類的東西跟肉類一定不會一起吃,所以芝士漢堡包一定賣不出去.所以麥當勞在以色列就沒生意好做的.但是現在一個猶太人要去傳福音,很可能對方是希臘人,很可能提供的食物不符合潔淨的條例；你怎麼辦?當彼得傳去傳福音時,看到異象,要他吃一些猶太人不能吃的東西,他就說吃不得˴吃不得.但是耶穌對門徒說,你去傳福音,就要入鄉隨俗,人家擺什麼你就要吃什麼.

傳福音就一定要進入別人的生活,一定要有生活的交流.我們只有從日常生活中才能真正全面地瞭解一個人,我們必須進入一個人的生活,才能向他傳道.所以我們必須入鄉隨俗.你去四川傳福音,你不吃麻辣,你怎麼同當地人打交道呢?如果你同泰國人傳福音,你整天批評佛教,整天覺得他們貧窮落後,沒有法制,那怎麼向他們傳福音呢?

所以神向我們提到食物的問題,其實也就是「道成肉身」的問題.神親自成為人,進入一種文化,進入一種生活,入鄉隨俗,向人們展現天國具體的樣式；因此,他也要求門徒學他的榜樣,入鄉隨俗,進入別人家庭日常的生活去傳播福音.所以第9節主還特別提到了這一點:要醫治那城的病人.治病救人,不是最平常的生活麼?

三,結論 10-12節

不過福音的傳播不是一帆風順的事.10-12節,主要門徒向人們宣佈:神的國近了.什麼是神的國近了?神的國近了,就是指耶穌基督受難,復活的日子快要來到了.但是神的國不僅帶給我們祝福,帶給我們救贖,同時,他還帶給我們審判,就像所多瑪所受的審判那樣.我們是傳揚救贖的好消息的人,同時,我們也是宣佈神的審判將會臨到的人.當有人憑著信心領受神的道時,他們就得了神的平安,神的救贖；但是,當有人抗拒神的道時,我們也不必勉強,因為我們該當盡的責任已經盡了,我們擦掉腳上的塵土,往前方走就好了.

猶太人經過外邦人的地方時,就會把腳上的塵土擦下來才離開,表示對這個地方的咒詛.不過耶穌的教導已經賦予這個行動以新的含義.他是讓門徒做見證,去見證那要到來的審判是多麼可怕啊.就好像在黑夜裡,水壩要塌了,我們在黑夜裡大聲喊:「快起來啊,快起來啊,水壩要塌啦!」這種呼喊不是咒詛,而是愛.有的朋友說,你們基督教好霸道,審判這個,審判那個.我們可以告訴他們:我們不是為了審判而審判,我們只是那個在黑夜裡呼喊,叫醒大家的人.

神提醒我們,做為他的門徒,我們必須要在家庭中,在生活中傳福音.但是我們不是憑著自己的一腔熱情去傳福音就夠了.我們還要有屬神的智慧.有的弟兄姊妹傳福音不看對象,喋喋不休,那就不是傳福音,反而是阻擋福音了.傳福音一定要有以上所說的幾種智慧:第一,必須要有同工.出來行,行路難啊!我們不僅要有彼此配搭的幫手,而且我們也必須在彼此合作中磨練自己屬靈的能力.第二,我們要明白信仰不是旅行,乃試探險.因此我們不能再倚靠那些我們從前倚靠的東西,我們必須放下以前的安全感,全心全意去倚靠主.第三,我們要放下我們的驕傲,入鄉隨俗,真正進入對方的生活去傳播福音.弟兄姊妹,不輕看了自己.請觀看稻田,主的莊稼在面前,遼闊收成沒法估計!我們還在等什麼呢?



\newpage

\section{好撒馬利亞人是如何鍊成的}
\label{sec:209}
\textbf{林豪恩先生 - 好撒馬利亞人是如何鍊成的?}
\newline
\newline
\hyperref[sec:208]{< < < PREV SERMON < < <}
~
\hyperref[sec:toc]{[返目錄]}
~
\hyperref[sec:210]{> > > NEXT SERMON > > >}


路加福音 10:25-37
\newline
\begin{longtable}{cl}
\hline
\hline
章節 & 經文 (和合本修訂版)\\
\hline
10:25 & \begin{tabularx}{0.7\textwidth}{X} 有一個律法師起來試探耶穌，說：「老師！我該做甚麼才可以承受永生？」 \end{tabularx} \\ \\ \relax
10:26 & \begin{tabularx}{0.7\textwidth}{X} 耶穌對他說：「律法上寫的是甚麼？你是怎樣念的呢？」 \end{tabularx} \\ \\ \relax
10:27 & \begin{tabularx}{0.7\textwidth}{X} 他回答說：「你要盡心、盡性、盡力、盡意愛主—你的神，又要愛鄰如己。」 \end{tabularx} \\ \\ \relax
10:28 & \begin{tabularx}{0.7\textwidth}{X} 耶穌對他說：「你回答得正確，你這樣做就會得永生。」 \end{tabularx} \\ \\ \relax
10:29 & \begin{tabularx}{0.7\textwidth}{X} 那人要證明自己有理，就對耶穌說：「誰是我的鄰舍呢？」 \end{tabularx} \\ \\ \relax
10:30 & \begin{tabularx}{0.7\textwidth}{X} 耶穌回答：「有一個人從耶路撒冷下耶利哥去，落在強盜手中。他們剝去他的衣裳，把他打個半死，丟下他走了。 \end{tabularx} \\ \\ \relax
10:31 & \begin{tabularx}{0.7\textwidth}{X} 偶然有一個祭司從那條路下來，看見他就從另一邊過去了。 \end{tabularx} \\ \\ \relax
10:32 & \begin{tabularx}{0.7\textwidth}{X} 又有一個利未人來到那裡，看見他，也照樣從另一邊過去了。 \end{tabularx} \\ \\ \relax
10:33 & \begin{tabularx}{0.7\textwidth}{X} 可是，有一個撒瑪利亞人路過那裡，看見他就動了慈心， \end{tabularx} \\ \\ \relax
10:34 & \begin{tabularx}{0.7\textwidth}{X} 上前用油和酒倒在他的傷處，包裹好了，扶他騎上自己的牲口，帶他到旅店裡去，照應他。 \end{tabularx} \\ \\ \relax
10:35 & \begin{tabularx}{0.7\textwidth}{X} 第二天，他拿出兩個銀幣來，交給店主，說：『請你照應他，額外的費用，我回來時會還你。』 \end{tabularx} \\ \\ \relax
10:36 & \begin{tabularx}{0.7\textwidth}{X} 你想，這三個人哪一個是落在強盜手中那人的鄰舍呢？」 \end{tabularx} \\ \\ \relax
10:37 & \begin{tabularx}{0.7\textwidth}{X} 他說：「是憐憫他的。」耶穌對他說：「你去，照樣做吧！」 \end{tabularx} \\ \\
[1ex]
\hline
\hline
\end{longtable}
題目:好撒馬利亞人是如何鏈成的?

經文:路加福音10:25-37

場合:崇基學院禮拜堂主日崇拜(14/7/2013)

一．經文

10 : 25 有一個律法師起來試探耶穌,說:「老師!我該做甚麼才可以承受永生?」

10 : 26 耶穌對他說:「律法上寫的是甚麼?你是怎樣念的呢?」

10 : 27 他回答說:「你要盡心,盡性,盡力,盡意愛主-你的上帝,又要愛鄰如己.」

10 : 28 耶穌對他說:「你回答得正確,你這樣做就會得永生.」

10 : 29 那人要證明自己有理,就對耶穌說:「誰是我的鄰舍呢?」

10 : 30 耶穌回答:「有一個人從 耶路撒冷下 耶利哥去,落在強盜手中.他們剝去他的衣裳,把他打個半死,丟下他走了.

10 : 31 偶然有一個祭司從那條路下來,看見他就從另一邊過去了.

10 : 32 又有一個 利未人來到那裡,看見他,也照樣從另一邊過去了.

10 : 33 可是,有一個 撒瑪利亞人路過那裡,看見他就動了慈心,

10 : 34 上前用油和酒倒在他的傷處,包裹好了,扶他騎上自己的牲口,帶他到旅店裡去,照應他.

10 : 35 第二天,他拿出兩個銀幣來,交給店主,說:『請你照應他,額外的費用,我回來時會還你.』

10 : 36 你想,這三個人哪一個是落在強盜手中那人的鄰舍呢?」

10 : 37 他說:「是憐憫他的.」耶穌對他說:「你去,照樣做吧!」

二．好撒馬利亞人的故事及其無數現代版本

「有一個人從 耶路撒冷下 耶利哥去,落在強盜手中.他們剝去他的衣裳,把他打個半死,丟下他走了.  偶然有一個祭司從那條路下來,看見他就從另一邊過去了.又有一個 利未人來到那裡,看見他,也照樣從另一邊過去了.可是,有一個 撒瑪利亞人路過那裡,看見他就動了慈心,上前用油和酒倒在他的傷處,包裹好了,扶他騎上自己的牲口,帶他到旅店裡去,照應他.第二天,他拿出兩個銀幣來,交給店主,說:『請你照應他,額外的費用,我回來時會還你.』

好撒馬利亞人的其中一個現代版本

有一位婦女16年前由廣東惠州來港與丈夫團聚.自從丈夫10幾年前退休後,一家一度要靠領取綜援維生.為幫補家計,養大4名子女,這位婦人十幾年前開始在旺角花墟一帶撿拾紙皮.長年勞碌,年僅五旬的婦人看起來比實際年齡蒼老,頭髮灰白,腳趾更因長期在骯髒環境工作而潰爛.

2013年6月22日下午1時左右,這位拾紙皮婦人在洗衣街垃圾站旁的公廁外與外判清潔工起爭執,當曾婦如廁的時候,將滿載紙皮的手推車留在街上,工人遂以阻街為由將曾婦的手推車連紙皮「充公」,並用鐵鏈鎖在垃圾站內.曾婦其後返回欲取回手推車,但垃圾站工人卻說不能自行決定,要先報告上級.曾婦堅持取回手推車,並鑽進手推車與牆壁之間空隙大哭控訴常遭欺負,工人則報警指有人「企圖盜竊」.

商人黃先生路過目擊事件,留在現場安撫這位婦女,以及協助她向警察解釋事件.黃先生見曾婦被沒收紙皮,即掏出200元給曾婦,說:「當我幫你買紙皮.」之後更買水給她喝.後來這位婦女被送到醫院檢查,黃先生還到醫院探望她.(明報 2013年06月23日)

三．從道基督教誨的層次看好撒馬利亞人的故事

正如上述的新聞所展示,好撒馬利亞人故事中的戲情,有時也出現在我們現代人的日常生活之中.人們對於受苦者的回應,也好像故事中的人物一樣,有些人視若無睹,有些人則出手相助.

2011年10月13日傍晚,在廣東省佛山市一個五金城,一名2歲女童王悅(乳名「小悅悅」)在陰雨天黑夜獨自跑出家門百米外後,先後被兩輛汽車撞傷倒地,失救致死.閉路電視拍攝了小悅悅被車撞倒前後的過程:第一輛客貨車撞倒女童,前輪碾壓過胸腔和臉頰,司機稍作停留後,後輪繼續碾壓胸膛,隨後離去.其後3名途人路經事發現場,均沒有人救助受傷倒地的女童,隨後有另一輛小貨車駛過,再次輾壓女童雙腿.多名路人沒有伸出援手,而後來一名拾荒婦人將女童抱起到路邊,之後女童母親出現,迅速抱走女童.

這事件引起了鋪天蓋地的討論,包括在報章雜誌上,電視電台上,互聯網上.最多人提出的問題是:「為甚麼有些人看見受苦者視若無睹?」也有人從另一個角度問:「為甚麼有些人會出手相助呢?」

三．一．視若無睹總有理

其實,視若無睹不一定沒有理由的.在好撒馬利亞人故事中,祭司和利未人看見那位被躺臥在路上的對害苦者視若無睹,有人嘗試替他們的解釋:

第一種解釋是自己未必合適出手相助:祭司及利未人可能擔心這位受害者不是律法所界定的「鄰舍」,而現場狀況又不容許他們確認對方是否受割禮的人,因此還是留給其他人協助為妙.

可能有人感到這項理由似是而非,然而,受苦者是否值得我幫助?是否有人比我更合適去幫助這些受苦者呢?這些問題的確出現在我們的群體中,也影響著我們對受苦者的回應.六月初,我們一行二十人往柬埔寨首都金邊進行宣教服務體驗.某天,我們走進舊垃圾山旁的貧民區探討村民.同學們每人手上拿著最少兩袋物資,內有白米,即食面,食油,鑵頭,以及沐浴露等,逐家逐戶送給村民.一列列鐵皮屋,搭建在垃圾山的邊緣.同學們腳上踩踏著吸滿了雨水的垃圾,頭上頂著列日當空的太陽,然而,他們沒有埋怨,仍然透過翻譯耐心地與村民溝通,鼓勵他們,為他們禱告.探訪結束,我們離開貧民區,往金邊市中心去吃飯.當地的翻譯把我們帶到一個大型商場,同學們稱之為「金邊的又一城」.看見當地人在「金邊又一城」的消費水平,同學們受到不同程度的衝擊.他們發現,柬埔寨人不都是窮人,有些人比自己更有錢.於是他們問:「為甚麼柬埔寨的有錢人不去幫助他們的窮人,而要我們遠道而來去幫助他們呢?」回看好撒馬利人的故事,如果祭司和利未人這樣想,也未必是沒有道理的.不久之前四川發生地震,城中有言論說:「中國大陸錢多的是,何來需要我們的捐款!」這些言論,也引起不少人的共鳴呢!

第二種解釋是這樣說的:律法禁止祭司觸摸屍體,如果摸了,就是不潔淨.若沾染不潔,就不能主持獻祭,不能吃祭物,而且要行潔淨禮.祭司和利未人或者是為了潔淨的緣故而對路旁的不幸者避之則吉.

第三種解釋認為,祭司和利未人擔心劫匪還埋伏在附近.若停下來幫助受害者,自己也可能成為劫匪的目標.

上述兩項理由都與幫助受苦者的後果有關,若然路見不平拔刀相助,自己反而會被刺傷,這樣的後果難免阻礙人們對受苦者伸出援手.

人們討論小悅悅事件的時候,引用了之前發生的兩項案例,證明上述的擔憂是合理的.

2006年,一名南京青年彭宇在巴士站將一名跌倒在地的長者徐春蘭扶起,並將其送去醫院,豈料徐氏一口咬定是彭氏撞倒她,要求索償.南京鼓樓區法院接受徐氏的說法,在判詞指出,彭氏若 無撞人,應協助捉住涉事者,而不是先救人；可以通知徐氏的家屬,而無需將其送院；在醫院替徐氏支付200元,是其「知錯」而支付的「賠償」.結果幾經調解下,彭氏仍要向徐氏支付一萬多元.

2009年,許雲鶴駕車經過天津紅旗路時發現67歲的老人王秀芝,跨越馬路中間的護欄後摔倒,他隨即停車並召喚救護車將老人送往醫院,不料老人事後一口咬定是許雲鶴撞傷了自己.天津紅橋法院一審判決:許雲鶴承擔40\%的責任,賠償費10萬多元.

聽過這些說法,大家會否覺得祭司和利未人不去救助那位被劫匪打得半死的路人其實也有道理的呢?會否覺得他們視若無睹其實有合理解釋的呢?祭司和利未人的視若無睹有理由,那麼,撒馬利亞人又為甚麼出手相助呢?他們的分別在哪裡呢?

三．二．出手想助因慈心

故事中的祭司,利未人和撒馬利亞人所看見的情景都是一樣,但回應則不同,前者「看見就從另一邊過去」,後者「上前用油和酒倒在他的傷處,包裹好了,扶他騎上自己的牲口,帶他到旅店裡去,照應他」.他們的回應的不同在於祭司和利未人看見那位受苦者沒有動慈心,而那位撒馬利亞人卻動了慈心.慈心,或譯作憐憫之心,是路加福音經常用來形容耶穌的字眼.

7 : 11  過了不久,耶穌往一座城去,這城名叫 拿因,他的門徒和一大群人與他同行.

7 : 12  當他走近城門時,有一個死人被抬出來.這人是他母親獨生的兒子,而他母親又是寡婦.城裡的許多人與她一同送殯.

7 : 13  主看見那寡婦就憐憫她,對她說:「不要哭.」

7 : 14  於是耶穌進前來,按著槓,抬的人就站住了.耶穌說:「年輕人,我吩咐你,起來!」

7 : 15  那死人就坐了起來,開始說話,耶穌就把他交給他的母親.

耶穌有所行動,是因為他動了憐恤之心.那位撒馬利亞人對躺臥在路上的受害人有所行動,也是因為對他動了憐憫之心.

我們的先賢曾說過,惻隱之心,人皆有之.然而,生活經驗常常叫人收起惻隱之心.有人說,對敵人仁慈,即是對自己殘忍.敵人只是其中一項理由,正如上述所檢視,還有更多都說得通的理由讓人把憐憫收藏起來.不過,當人們越多說得通的理由去收起慈心,受苦的人就越受苦.在這些時候,我們的道理可能正在延長別人的痛苦.這些痛苦的人,可能是陌生的路人,也可能是我們的枕邊人.

有人認為,對很多夫妻而言,兩人發生衝突時,最直接的發洩方式並不是「廝殺」,而是展開冷戰.當兩人有了矛盾,一方率先擺出一副「我以後再也不理你」的姿態,同時另一方心想「不理就不理,誰怕誰啊」,雙方都認為自己有發難的理由,雙方都有理由收起對對方的慈心,於是,冷戰爆發了.其實,在背對著背躺在同一張床上的時候,雙方都知道對方徹夜難眠,雙方都感受到對方痛苦.這個時候,是最需要憐愛的時候,最需要恩慈時候.道理未必有效消減大家的痛苦,但憐愛和恩慈卻是最好的止痛藥.

憐憫之心在哪裡呢?其實離我們不遠.記者采訪扶起小悅悅到路邊的拾荒婦人時,問出了那句公眾最關心的話︰“孩子不知是死是活躺在地上,你不怕救了之後被人冤枉,被人坑嗎?”但她說自己什麼都沒想,下意識地上前救人,“這事很平常啊”.對受苦者生出慈心,對痛苦者伸出援手,但願真的成為我們日常生活很平常的事.

四．從信仰教導的層次看好撒馬利亞人的故事

耶穌講好撒馬利亞人故事,本來不是回答關於道基督教誨的問題,而是回答信仰問題.這個信仰問題是一位律法師向耶穌提出的:

10 : 25 有一個律法師起來試探耶穌,說:「老師!我該做甚麼才可以承受永生?」

「我該做甚麼才可以承受永生?」亳無疑問是一個信仰問題,而且是一個非常重要的信仰課題.

在好撒馬利亞人故事中,有三個主要人物,分別是祭司,利未人,和撒馬利亞人.這三個人物,其實可歸納為二種角色,第一種角色是代表著根正苗紅的聖殿神職人員的祭司和利未人,第二種角色是代表著正統猶太人血統以及正統猶太人信仰的邊緣者的撒馬利亞人.如果我們訪問耶穌時代的猶太人,這兩種角色哪一種與信仰比較有關係呢?當然是祭司和利未人,因為他們是根正苗紅的聖殿神職人員.再問,這兩種角色哪一種更接近上帝呢?當然是祭司和利未人,因為他們的任務就是代表人向上帝獻祭.再問,這兩種角色哪一種與承受永生更有關係呢?當然也是祭司和利未人,他們知道如何正確地敬拜上帝,也是正確地敬拜上帝的執行者.

然而,在耶穌所講的故事裡,這種理所當然的假設被挑戰了.我們作為聽故事的人,聽到耶穌如此講述的這個故事,你感覺到耶穌認為哪一種角色更接近信仰的精神?你感覺到耶穌認為哪一種角色與律法師起初的問題與承受永生更接近呢?

耶穌用來回答律法師關於永生的問題而講述的這個好撒馬利亞人的故事,會否引發我們思考,憐憫之心會否比身份和角色更接近永生呢?憐憫之心會否比正確而且專業地敬拜上帝更為有關係呢?

讓我們再退後一層,進入耶穌講這個好撒馬利亞人故事的場景.我們會看見有兩個主要人物,一位是前來向耶穌發問的律法師,另一位是被問的耶穌.按路加福音的作者的描述,這兩個角色也是很特別的,發出信仰問題的並不是一位慕道者,而是一位是律法師,即是當時的聖經教師.這個角色代表著信仰的知識份子,或者信仰的專家.這位信仰專家來問耶穌「如何承受永生」的問題,其實並不是他沒有答案,而是要來考核耶穌,或者挑戰耶穌,進而希望證明自己是對的,耶穌是錯的.耶穌的回答,卻巧妙地帶讀者去比較律法師所代表的信仰知識份子的角色以及撒馬利亞人所代表的信仰邊緣者的角色,兩者之間哪一種角色與承受永生較為有關係呢?從耶穌兩次肯定律法師所知道的,進而邀請他實踐,會否引發讀者去思考,只在知識上知道愛的道理,與承受永生仍有距離；唯有愛的實踐,把人帶到永生之中呢?

祭司並非等閒角色,正確的宗教禮儀也並非不重要,如果,祭司和他們主持的宗教禮儀朝向培育人們的憐憫之心,引發人們向受苦者伸出援手,相信耶穌對他們的評價就更正面了.聖經老師並非等閒角色,正確的信仰知識也並非不重要,如果,聖經老師和他們的信仰知識能夠帶引自己和人們把知識實踐,那麼,耶穌對他們的評價也會有所不同了.

五．總結

好撒馬利亞人的故事,一方面是道德層面的教誨,喚起人們「惻隱之心, 人皆有之」的本質.縱使現實上有很多合理的考慮足以阻礙人們展現憐憫的心,但其實人們心底裡仍盼望著一個有憐憫的世界.當有人因為救助別人反而被判有罪的事件出現後,在北京著名的三所大學所進行的民調,問「老人摔倒該不該扶」這問題,仍有高達64.8\%的受訪者認為應該扶起老人.

好撒馬利亞人的故事,也指向信仰的核心課題,提示著人們信仰和道德是不能分割的.在路加福音裡,耶穌曾經如此要求門徒仁慈,不是用經驗的原因,而是用信仰的要因.

6 : 36  你們要仁慈,像你們的父是仁慈的.」



\newpage

\section{重整生命焦點}
\label{sec:210}
\textbf{許開明牧師 - 重整生命焦點}
\newline
\newline
\hyperref[sec:209]{< < < PREV SERMON < < <}
~
\hyperref[sec:toc]{[返目錄]}
~
\hyperref[sec:211]{> > > NEXT SERMON > > >}


路加福音 10:38-42
\newline
\begin{longtable}{cl}
\hline
\hline
章節 & 經文 (和合本修訂版)\\
\hline
10:38 & \begin{tabularx}{0.7\textwidth}{X} 他們繼續前行，耶穌進了一個村莊。有一個女人，名叫馬大，接他到自己家裡。 \end{tabularx} \\ \\ \relax
10:39 & \begin{tabularx}{0.7\textwidth}{X} 她有一個妹妹，名叫馬利亞，在主的腳前坐著聽他的道。 \end{tabularx} \\ \\ \relax
10:40 & \begin{tabularx}{0.7\textwidth}{X} 馬大伺候的事多，心裡忙亂，進前來，說：「主啊，我的妹妹留下我一個人伺候，你不在意嗎？請吩咐她來幫助我。」 \end{tabularx} \\ \\ \relax
10:41 & \begin{tabularx}{0.7\textwidth}{X} 主回答說：「馬大，馬大，你為許多的事操心煩惱， \end{tabularx} \\ \\ \relax
10:42 & \begin{tabularx}{0.7\textwidth}{X} 但是不可少的只有一件。馬利亞已經選擇了那上好的福分，是沒有人能從她奪去的。」 \end{tabularx} \\ \\
[1ex]
\hline
\hline
\end{longtable}
講題:重整生命焦點　Reajusting To A Focused Life

經文:路加福音10章38-42節

講員:許開明牧師

場合:香港中文大學崇基學院禮拜堂

日期:2013年7月21日

引言:

◙那塊葉子著火起來了

小時,上《自然科》課,老師教我們做個簡單實驗；當時在炎炎夏日,烈日當空,我們把一片葉子放在地上,然後拿出放大鏡放在葉子上並向著太陽,等了一陣子,那塊葉子著火起來了,我們看見非常興奮.

老師解釋把陽光聚焦在放大鏡上,便能產生燃燒的效應.學者凱斯特鮑姆對人生的觀察也一樣,他說:「生命有焦點,讓人看來大有能力；欠缺生命焦點的人,會自視為效率低甚或無能,別人也這樣看他們……生命焦點,是真實的信心,信念和自信的源頭.」這是至理名言.

◙一個人集中精力做一件事,常會有最佳的成果.

正如激光集中光能斷刀槍.

◙相反,生命無焦點,就會忙亂!

經文

【路十38】「他們走路的時候,耶穌進了一個村莊；有一個女人名叫馬大,接祂到自己家裡.」

﹝原文字義﹞「馬大」太太,貴婦.寡婦

﹝文意註解﹞「耶穌進了一個村莊,」這村莊就是伯大尼(參約十一1).

【路十39】「她有一個妹子名叫馬利亞,在耶穌腳前坐著聽祂的道.」

﹝原文直譯﹞「…她『也』坐下來在耶穌腳前聽祂的話.」

﹝原文字義﹞「馬利亞」背叛者,苦.

﹝文意註解﹞按原文,本節有一個『也』(also)字,這個『也』字相當重要,表示她不是無所事事,光坐在那裡享受；而是在作完了一些事之後,『也』坐下來聽主講道.

【路十40】「馬大伺候的事多,心裡忙亂,就進前來說:『主阿,我的妹子留下我一個人伺候,你不在意麼?請吩咐她來幫助我.』」

﹝原文直譯﹞「但馬大被許多伺候的事分心了,就進前來說…」

﹝原文字義﹞「心裡忙亂」分心,困擾,負擔過重.被許多情操縱著,被困住.

【路十41】「耶穌回答說:『馬大,馬大,你為許多的事,思慮煩擾；」

﹝文意註解﹞「你為許多的事,思慮煩擾,」主不是說她『事情作得太多』,主乃是說她『思慮多到令她受到攪擾,苦惱』.主也不是說『作事不要用心思』,主乃是說『心思不要被事物弄得煩擾不安,失去在主裡該有的安寧』.

﹝話中之光﹞(一)問題不在於事多,事少,問題乃在於「思慮煩擾」；我們不作工則已,一作起工,就常常讓許多事物進到心裡面.

(二)我們要用心思指揮事情,不可讓事情指揮我們的心思.

【路十42】「但是不可少的只有一件；馬利亞已經選擇那上好的福分,是不能奪去的.』」

﹝原文字義﹞「不可少的」必須的；「上好的」美好的(good)；「福分」一分(part,portion),同得.

路加10:38- 42中,提及馬大「伺候的事多,心裡忙亂.」(路加10:40)套句現在的話說,馬大忙得心慌意亂.「心裡忙亂」就是分心,失去生命的焦點.

試想馬大的心多麼偉大!耶穌是貴賓,馬大心中第一要緊的事,就是讓祂受到最周全的招待.今日我們也曾見過這樣的婦人,她們出於愛心,為了使客人享受完美的招待,她們衝進衝出,忙個不停,好像有千頭萬緒的事等著幹的.如此重視客人,無可厚非；然而,馬大想用服事來表達她的愛心,卻收到了反效果.

因為她忙得心慌意亂,以致於忍耐不住,衝進客廳抱怨耶穌和她的妹妹.馬大對耶穌說,「你不在意嗎?」她的忙亂使她只看到自己,自己的工作,自己需要,自己的了不起；還有,只看到別人的不是,馬利亞不應坐著不幫手,耶穌不應坐視不理,任她一人做餐死,她在發怨言,發火和發炮,這樣招待客人真是弄巧反拙.

今日忙碌的馬大較多,坐主前的馬利亞較少.

現代人愛「忙碌」,甚至盲目崇拜「忙碌」

一個人自覺「忙碌」,自覺「重要」,「有價值」

一個人自覺「忙碌」,自覺「能幹,有魄力」

馬大集中於完成一件又一件的工作,並以完成為喜樂.

當忙死時,最看不過眼是有人繞著手,不來幫手,不幹.

或許,馬大的心思是好的,好好的弄一餐給耶穌和門徒,才聽道,她應成為好主人.

1在忙亂中,人眼花瞭亂,失去生命焦點

我們有許多事要做,要處理,許多人要聯絡,要見,

然而,馬大看重「緊急的事」,卻忽略「重要的事」.

2在忙亂中,失去自己

馬大可能也想煮好東西後,像馬利亞坐在耶穌腳前,陪伴耶穌,可惜她有許多事要做,許多責任要盡.她無法做到她想做的.我們也常如此,

等我退休,我會事奉；

等我有時間,我會多些陪伴子女；

等我開完會,我便來探你(幸好未完到訪)

3在忙亂中,看不到別人真正的需要

兩種不同的關懷

馬大著重耶穌外在的「吃」的需要,生理上的需要,一般人的需要.

耶魷來到伯大尼進入那家中,這是個大日子；而馬大衛了表示慶祝,盡他一己所能,把家中一切最好的,拿出來款待耶穌.

馬利亞著重耶穌內在的,心理上的需要,陪伴(精神).為了擺在面前的十字架,和內心的折騰,祂極盼望能擺脫不斷向祂需索的群眾；在伯大尼這個地方,找得一處安靜之所,即使是一兩個小時也是好的.

馬利亞的態度正切合耶穌的心意,而馬大卻在善意的盡力破壞.

耶穌向馬大提出『不可少的只有一件』──極可能是指『我所要的不是大擺筵席；一道菜,一點點簡單的食物,我便於願已足.』事情很簡單的,就是馬利亞明白耶穌的心意,而馬大卻不明白.

4在忙亂中,只看到自已的「是」,別人的「錯」

馬利亞在耶穌腳前坐聆聽他的道,而馬大則為伺候這位尊貴的上賓心裡忙亂.馬大希望主責備她妹妹沒有幫她的忙,但主耶穌溫和地責備了馬大的煩躁!

事奉可能會被驕傲和自以為是沾污；

讓祂佔住內心是不可少的一件事,是不能奪去的上好的福分.高斯評論說:「主想把我們由馬大轉變為馬利亞,正如祂想把我們從律法師轉變為鄰舍一樣.

5在忙亂中,人容易產生「埋怨」與「憤怒」

馬大「埋怨」與「憤怒」馬利亞及耶穌

過份的擔心與過份的分心,便需要處理.

在英國時,我們前後的花園都種了玫瑰花,師母亦對玫瑰花情有獨鍾,我經常在玫瑰園默想或唱歌；無意中發覺玫瑰花很香.有位西人朋友對我說:

「如果你不停下來,聞聞沿途上玫瑰花花的香味,你會感到人生通往天堂的路是佈滿沙石難行的路.」

“the way to heaven is a rough and rocky road” if we don’t stop and smell the roses along the way.

6忙亂中,使人懷疑上帝的看顧

【路十40】「馬大伺候的事多,心裡忙亂,就進前來說:『主阿,我的妹子留下我一個人伺候,你不在意麼?請吩咐她來幫助我.』」

耶穌若關心應該會在意「馬利亞」,叫馬利亞幫助她.

7在忙碌中,人會陷入思慮煩擾

【路十41】「耶穌回答說:『馬大,馬大,你為許多的事,思慮煩擾；」只坐無做會無能

只做無坐會無方向.

有坐有做產生能力與均衡.

富士德:「忙碌,好像魔鬼；不,它本身是魔鬼」

重整關係

【路十41】「耶穌回答說:『馬大,馬大,你為許多的事,思慮煩擾；」

【路十42】「但是不可少的只有一件；馬利亞已經選擇那上好的福分,是不能奪去的.』」

﹝原文字義﹞「不可少的」必須的；「上好的」美好的(good)；「福分」一分(part,portion),同得.

今日仍然有許多事,許多的煩擾,我們需要生命的焦點.只有一件事是最需要的?

馬利亞從忙亂中出來:親近主,敬拜主

馬利亞卻能,因為她在忙碌後,懂得坐下來

【路十39】「她有一個妹子名叫馬利亞,在耶穌腳前坐著聽祂的道.」

1在耶穌腳前:

▄在主腳前:謙卑,頭向下,敬拜主

忙碌,好像魔鬼；不,它本身是魔鬼

▄約櫃是耶和華的腳凳,上主坐在高高的寶座上,人類只能見到祂的腳凳.這表示人敬拜主,親近主.

▄馬利亞到了耶穌那裡,看見他,就俯伏在他跟前-(約11:32)

▄馬利亞就拿著一斤極貴的真拿達香膏抹耶穌的腳,又用自已的頭髮去擦…(約12:3)連續三節經文提及她是位在耶穌腳前的人.

2坐著:

不會周圍走,專一陪伴著主,這是一種禮儀,有朋友到,有人陪著他,與他傾談.

她知道如此,可能觸怒馬大；但她知道這是人生難得的機會,與主在一起.

3聽祂的道

▄ 聽,側耳聽,吸收,專注於

▄聽道:得著主的話,使人跳出忙碌的框框

▄馬利亞是四福音書中,被描寫得最多的.

我們發現她經常親近上主,生命充滿動力

▄馬大:活動型的:為耶穌做很多東西,使耶穌開心

馬利亞:靈修型:陪伴耶穌,聽他說話,逗耶穌開心

跳出框框

一位老師進入課室,拿出一張畫紙來,然後用黑色水筆,在紙上畫上一個黑色的圓圈,跟著問同學說:「你們看見什麼?」許多學生回應說:「我們看見一個黑色的圓圈」,老師再問:「還有什麼呢?」大家都靜下來,只有一位同學舉手說:「我看見沒有圓圈的部份」.

老師點點頭說:「還有什麼?」課室雅雀無聲良久,老師忍不住說:「你們看不見嗎?還有最重要的部份,就是這張紙哩!」

許多時候,我們總是跳不出眼界或思維的框框,以致生命被局限,無法伸展.未能跳出框框是現代人常見的情況,「人在江湖,身不由己」是我們常掛在口中的言辭.

路加福音記載伯大尼的一個家庭(路10:38-42)兩姊妹的表現.當中馬利亞在耶穌腳前坐聆聽他的道,而馬大則為侍候這位尊貴的上賓心裡忙亂.

馬大希望主責備她妹妹沒有幫她的忙,誰知主耶穌反而責備了馬大的煩躁.從中我們不難發現馬大在忙碌中未能跳出框框.她被許多要作的事,弄得心煩意亂.忙碌會使人進入「茫,忙,忘,亡」的四種狀態,

「茫」是沒有目標和方向；

「忙」是忙亂,不能自制；

「忘」是忘記自己真正需要；

「亡」就是死亡,與人與上帝的關係隔絕.

這四種情況都是忙碌可能引發的後果.

耶穌對馬大說:『馬大,馬大,你為許多的事,思慮煩擾」,馬大被許多的工作與責任纏繞著,他跳不出來.在今日注重工作,成效,表現的社會,我們常會陷入忙碌的困局中而不能自拔.

曾有釋經學者指出,馬大其實都想像馬利亞好好地服事耶穌後,便可以停下來專心聽主道,可惜是她要處理的事太多了,忙到團團轉,想停也停不了,這豈不就是今日大忙人的景況,總是跳不出自己的框框.

相反馬利亞卻能從框中跳出來.根據原文提及馬利亞時,應是她『也』坐下來在耶穌腳前聽祂的話,本節有一個『也』字,這個『也』字相當重要,表示馬利亞不是無所事事,光坐在那裡享受；而是在作完了一些事之後,『也』坐下來聽主講道,所以馬利亞能夠在忙碌的服事中抽身出來,她的成功之處也在此.

【】2001年可說是我人生事奉的高峰期,牧會近二十年,參與多間機構的義務工作,也完成了教牧學博士,教會重建亦快完成,當時滿有成功感,並想像新堂建成可以全家搬入基督教大樓,多麼美好呀.

然而繁忙的牧會工作,換來卻是茫,忙,忘,亡.我發覺我要像馬利亞從忙碌的事奉中停下來.

我毅然辭去教會的工作,全家前往英國進修,五年的進修重要的不是完成兩個高級學位,而是能夠接近大自然,親近上主與家人.

我離任應聘主因是要專注於人生的焦點,歲月不留人,我已是過五十的中年人了,離開六十退休之齡會愈來愈近,如果現在還不把握機會,集中生命焦點,恐怕錯失良機.

自從在英國大學完成了兩個高級學位,獲得優良成積後,指導教授不斷鼓勵我投身學術研究的工作,近年發表的學術論文也得到學術界的好評,在學術期刊出版,激發我的研究心志；因此我人生下半場已決定,除牧會外,在神學教育,學術研究及文字工作方面發展,這是一條窄路,鮮有牧者走的路,然而卻是應主呼召的路.

在忙碌中停下來作生命反思,然後朝向人生下半場再上路.

重整生命的焦點:

曾子有言:「吾日三省吾身」

從忙碌中,靜下來:

每天給自己些小時間,想一想:

1別人對我什樣?我對別人什樣?

2上主對我什樣?我對上主什樣?

3我對自己什樣?我有些什麼地方需要重整?

重整我對別人

重整我對上主

重整我對自己

「雖失望但不沮喪」

在美國職業高爾夫球賽中,名人錦標賽是最具盛名的比賽.2009年,肯尼·佩里在決賽時,原本一路領先,最後卻落居第二.《紐約時報》用「雖失望但不沮喪」這句話,來描述佩里輸掉比賽後的表現.

佩里說:「偶爾我會回想,如果我採取不同的方式會不會贏?但是我不會陷在懊悔當中.如果這是我人生中發生過最壞的事情,我的情況還不算太糟.我不會讓失望的情緒糾纏著我.人生中有太多更重要的事情……,我今天會回家跟家人團聚,享受天倫之樂.」

佩里感失望但不沮喪,因為他生命的焦點,不在於球賽的成敗,還有人生更重要的東西.

對基督徒而言,我們的焦點決定了我們如何面對人生的成敗.我們不應定睛在那些讓我們失望的事情上.

保羅說:「所以你們若真與基督一同復活,就當求在上面的事；那裡有基督坐在上帝的右邊.你們要思念上面的事,不要思念地上的事」(歌羅西書3章1-2節).

換言之,自我的肯定與生命的意義不在於自身的成就,而在於榮耀基督.倘能專注於此榮耀基督作為人生的焦點,你將會不枉此生.

跳出生命的框框,原來是可以創造更好的人生.耶穌稱讚馬利亞能夠跳出忙碌的框框來親近上主,她選擇了不可少的,重要的一件事,且是上好的福分.

在名著《天地一沙鷗》中,聰明的海鷗約拿單,發覺許多海鷗只懂得在海岸追追逐逐,一生就是以捕魚為樂,但牠發現自己飛得愈高,看得愈遠.

許多人的一生是馬大的人生,活在忙碌,忙亂,埋怨,煩擾的框框中,多麼的可惜呀!為何不向馬利亞學習,跳出框框,選擇活出自主,自決,自由,自在,蒙主耶穌稱讚的蒙福生命呢.記住,人生是可以作抉擇的,你可以選擇活在框中或者跳出框外.

默想:

你是馬大忙碌型或馬利亞安靜型,能跳出忙碌框框的人?

你一直以來生命的焦點是什麼?

你需要重整自己生命的焦點嗎?

在那方面你需要重整?



\newpage

\section{人生於世何所禱也}
\label{sec:211}
\textbf{鄧瑞強博士 - 人生於世,何所禱也?}
\newline
\newline
\hyperref[sec:210]{< < < PREV SERMON < < <}
~
\hyperref[sec:toc]{[返目錄]}
~
\hyperref[sec:212]{> > > NEXT SERMON > > >}


路加福音 11:1-4
\newline
\begin{longtable}{cl}
\hline
\hline
章節 & 經文 (和合本修訂版)\\
\hline
11:1 & \begin{tabularx}{0.7\textwidth}{X} 耶穌在一個地方禱告。禱告完了，有個門徒對他說：「主啊，求你教導我們禱告，像約翰教導他的門徒一樣。」 \end{tabularx} \\ \\ \relax
11:2 & \begin{tabularx}{0.7\textwidth}{X} 耶穌對他們說：「你們禱告的時候，要說：『父啊，願人都尊你的名為聖；願你的國降臨； \end{tabularx} \\ \\ \relax
11:3 & \begin{tabularx}{0.7\textwidth}{X} 我們日用的飲食，天天賜給我們。 \end{tabularx} \\ \\ \relax
11:4 & \begin{tabularx}{0.7\textwidth}{X} 赦免我們的罪，因為我們也赦免凡虧欠我們的人。不叫我們陷入試探。』」 \end{tabularx} \\ \\
[1ex]
\hline
\hline
\end{longtable}
講題:人生於世,何所禱也?　What Should We Pray for in Our Life

經文:路加福音11:1-4

講員:鄧瑞強博士

場合:香港中文大學崇基學院禮拜堂

日期:2013年7月28日

各位弟兄姊妹,早安.

今日講道的內容,是關於祈禱的.

懂得祈禱,應該是人與其他動物最大的分界線.理性與感情,不是人獸之別的終極界線.其他動物也有理性,也有感情.動物也能與人溝通,你呼喚牠的時候,牠會回應.但當有一天,你回家時,見到你的狗跪在床邊,為你祈禱,你肯定會懷疑,這還是一條狗嗎?祈禱,是人類特有的靈性行為.

如何祈禱?祈禱的內容是什麼?這是我們靈性生命最重要的學習內容.

有一次,有人問德蘭修女:「當你祈禱時,你向神說什麼?」

德蘭修女回答說:「我沒講什麼,我只是聽.」

那人再問:「那麼,當你祈禱時,神向你講什麼?」

德蘭修女回答說:「祂也沒講什麼,祂也只是聽.」

德蘭修女對祈禱有很深的體會,她講到一種留心對方,專注對方的默然交流.在這交流中,「我」變得不重要了.「我」的說話不重要了,「我」的渴求也不重要了.若果要講「我」,則是作為「神所聆聽的對象」這意義下的「我」.在神的聆聽裡,在神的愛護裡,在神的注目下,「我」才呈現出來.這是蒙神所愛的「我」,這是蒙神恩典眷顧下的「我」.離開了神的恩典,離開了神的聆聽,離開了神的愛,「我」便失去了重量,也失去其真實意義.

德蘭修女關於祈禱的講法,點出了祈禱對於認識「我是誰」,有莫大重要性.

回到今日的講道經文:路加福音11:1－4.

這是主耶教教導信徒祈禱的經文.

11:1　耶穌在一個地方禱告；禱告完了,有個門徒對他說:「求主教導我們禱告,像約翰教導他的門徒.」

11:2　耶穌說:「你們禱告的時候,要說:我們在天上的父〔有古卷:父啊〕:願人都尊你的名為聖.願你的國降臨；願你的旨意行在地上,如同行在天上〔有古卷沒有「願你的旨意行在地上,如同行在天上」的字眼〕.

11:3　我們日用的飲食,天天賜給我們.

11:4　赦免我們的罪,因為我們也赦免凡虧欠我們的人.不叫我們遇見試探；救我們脫離凶惡〔有古卷沒有「救我們脫離凶惡」的字眼〕.」

新約有兩個主禱文的版本.一個在馬太福音,這是我們崇拜時常用的版本.一個在路加福音,這版本的主禱文與我們習慣背誦的主禱文,有一些差異.我們今天要看的,就是路加福音的版本.

馬太福音的讀者,是猶太人,他們從少已習慣祈禱.但是過於習慣,便流於虛偽,變成法利賽式的祈禱.在馬太福音裡,耶穌主動教人以主禱文祈禱,以正確的祈禱內容,糾正錯誤的祈禱態度.

路加福音的讀者,是非猶太人,他們不懂耶穌那種祈禱方式.在路加福音裡,人問耶穌如何祈禱.耶穌教他們以主禱文祈禱,從而讓他們明白跟從基督的真正含意.

11:1　耶穌在一個地方禱告；禱告完了,有個門徒對他說:「求主教導我們禱告,像約翰教導他的門徒.」

每個宗教都有自己的祈禱內容,每個派別都有自己的祈禱方式.一種祈禱的教導,塑造一種信仰群體.施洗約翰教他的門徒一種祈禱方式,塑造他們成一獨特群體.門徒求問耶穌,他們是否也有一種屬於自己群體的祈禱.門徒這裡求問的,表面上是祈禱的內容,但在實質上,卻試探討信徒與神有何獨特關係,信徒有何特別的身分.主耶穌的回答,塑造一個獨特的信仰群體.主禱文教導我們,我們應以一種什麼樣的獨特方式,存活於世上.每唸一次主禱文,就是再一次確認自己的身分.

11:2　耶穌說:「你們禱告的時候,要說:我們在天上的父〔有古卷:父啊〕:願人都尊你的名為聖.願你的國降臨；願你的旨意行在地上,如同行在天上〔有古卷沒有「願你的旨意行在地上,如同行在天上」的字眼〕.

古代留傳下來的不同經文抄本,在這裡有不同寫法.有些抄本的經文短一點,有些抄本的經文長一點.我們可以理解,原先的經文應是「父啊」,但抄經的人也熟悉馬太的經文,受到馬太經文的影響,他將之改為「我們在天上的父」.若經文原先是「我們在天上的父」,則我們很難理解為何抄經的人將之改成「父啊」.關於經文校勘的問題,我不多講了.在新近出版的《和合本修訂版》裡,這節經文是:「父啊,願人都尊你的名為聖；願你的國降臨.」

路加福音的主禱文以一個呼喚作開始:「父啊」.簡單,直接,親切.相比之下,「我們在天上的父」較莊嚴,禮儀味道也較重.「父啊」是對於我們的生命寄托於何處的一個終極認定.在絕境裡,有人大聲呼喊:「錢啊!」在孤寂裡,有人大聲呼喊:「愛情啊!」可以想像,希臘的悲劇人物Oedipus在悲痛中會大聲呼喊:「無情的命運啊!」究竟是什麼東西承托或支配著我們的生命?是今世的財富?是感性的慰藉?抑或是無情的命運之力?

主禱文開宗明義提示我們,承托我們的,我們終極面對的,是「父」──一個接納我們,愛我們,對我們有期望的有情生命.路加福音獨家記載了主耶穌的一個比喻,去展示這位「父」的特質.就是那出名的「浪子回頭」的比喻.比喻裡講到的這位父,承受著小兒子的無情,大兒子的自義.他默默等待小兒子浪子回頭,循循善勸大兒子回心轉意.為的是讓他們體會何謂接納,何謂恩典,何謂愛.承托我們生命的,就是這位「父」.若我們將生命的意義建基於今生的成就,會有一種活法.若我們將生命的意義訴之無情的命運,會有另一種活法.若我們從心底呼喊「父啊」,將生命交給祂,我們的生命肯定會有一種新的活法.

「願人都尊你的名為聖.」中文這種譯法,給人一種感覺,就是神的名還未被尊為聖,我們現在只能渴望這事早日實現.這句的英文翻譯,一般是 “hallowed be thy name” (KJV, RSV, NIV),這更接近希臘文的結構,直譯是「被尊為聖的是你的名」.這是一種宣告與肯定,多於只是一種渴望.

啟示錄4:8提到有些神物不斷呼喊說:「聖哉!聖哉!聖哉!主上帝是昔在,今在,以後永在的全能者.」主禱文的「願人都尊你的名為聖」類近啟示錄的這個宣告,確認那顯出了大能,被尊為聖的,就是這位父,這位神.主禱文呼籲我們加入這個尊神的名為聖的行列.

尊一個人或一件物為聖,就是視之為終極價值.面對終極價值,其他的價值都可以捨去.人生有很多價值的取捨.孟子說:「生,亦我所欲也；義,亦我所欲也,二者不可得兼,舍生而取義者也.」孟子作了他的價值選擇,他尊「義」為聖,故甘願捨「生」取「義」.如今,我們的抉擇是:禮拜日,假期天,多睡一點我所欲也,崇拜神亦我所欲也.二者不可得兼時,我如何抉擇?

「被尊為聖的是你的名」,指示我們價值之所在,教導我們如何抉擇.

「願你的國降臨」,這句經文的結構與上句相似,用英文表達便是 “come thy kingdom”.同樣,與其說是一等待,不如說是宣告及肯定神的國降臨人間.啟示錄11:17提到天上的聖徒敬拜神說:「昔在,今在的主上帝－－全能者啊,我們感謝你!因你執掌大權作王了.」主禱文這句「願你的國降臨」類似啟示錄的這個讚頌,它呼喚我們看,神的國在人間.

當一個最卑污的罪人被寬恕,當一個最卑下的人被尊重,當一個最無助的人被關懷的時候,神的國便顯現出來了.當耶穌去探訪那些露宿者時,當耶穌去肯定那些中學文憑試不合格的人的價值時,當耶穌冒死去醫治那些沙士病人時,神的國便來臨了.神的國像一顆芥菜種一樣,毫不起眼,但確實在這裡.這是人間的希望所在.主禱文鼓勵我們去看,去肯定,耶穌昔日的救恩行動己將天國帶來人間,今日,天國仍在我們中間像種子一般發芽生長.

當你看到人間一切被慈愛的「父」承托著,當你確認只有神的名被尊為聖,當你看到天國在人間隱密地發展著,你的生命會全然不同.

有何不同?

主禱文跟著的幾句話,點出信徒與眾不同的地方.

11:3　我們日用的飲食,天天賜給我們.

這裡涉及生命最基本的需要──食物.我們每日需要的食物,天父「天天」賜給我們.我們習慣的馬太版本的經文,是「我們日用的飲食,今天賜給我們」.路加講「天天」,更凸顯天父持續的供應.不單「今天」供應,而是「天天」供應,更加叫我們不要為明天憂慮.

主耶穌教信徒這樣祈禱,是教我們一種面對物質世界的心態.

我們有物質性,但不全是物質,人有靈性.物質主義卻是要將一切約化為物質.物質主義會將價值約化為事物的「價格」,將真理約化為「有用」,將美約化為「追上時尚」.有一個小故事很能點出物質主義的精髓.《小王子》這本小書這樣寫:「如果你對大人說:『我看到一幢用玫瑰色的磚蓋成的漂亮的房子.它的窗戶上有天竺葵,屋頂上還有鴿子…』他們怎麼也想像不到這房子有多漂亮.你必須對他們說:『我看到一幢價值十萬法郎的房子.』這時,他們便會驚叫:『噢,多麼漂亮的房子啊!』」

當我們的心「物質化」了的時候,我們便會攫取,而不懂分享.只會計算一個孩子要花費多少萬,而不懂得與孩子一起捉迷藏.以消費的能力評估人的價值,而忘記人格.誠信變得無重量,懂得賺錢才最重要.

「我們日用的飲食,天天賜給我們」,教我們明白,來到我們手中的物質,全是天父的恩典.物質本身不是終極的,不是衡量一切的基礎.天父施恩典,人領受恩典,這才是生活的根本面貌.白白的得來,也當白白的捨去.

11:4　赦免我們的罪,因為我們也赦免凡虧欠我們的人.不叫我們遇見試探；救我們脫離凶惡〔有古卷沒有「救我們脫離凶惡」的字眼〕.」

「赦免我們的罪,因為我們也赦免凡虧欠我們的人」.這節經文在《新漢語譯本》譯作:「求你赦免我們的罪,正如我們也赦免所有虧負我們的人」.父赦免我們的罪,當然不是因為我們赦免了其他人.父赦免我們的罪,是恩典.我們是蒙恩的人,也當效法天父,去赦免其他人.

這裡涉及人與人之間最基本的相處之道:恕道.在路加福音6:27－29裡,主耶穌說:「你們的仇敵,要愛他!恨你們的,要待他好!咒詛你們的,要為他祝福!凌辱你們的,要為他禱告!有人打你這邊的臉,連那邊的臉也由他打.有人奪你的外衣,連裡衣也由他拿去.」這是一種不合常理的,以寬恕先行的存在方式.我想,這是跟從基督的人最極端的行為模式,卻就是基督徒群體的身分特徵.

克羅地亞裔的神學家沃弗(Miroslav Volf)寫了本巨著《擁抱神學》(Exclusion  and  Embrace),探討十字架在這仇恨世界裡的意義.克羅地亞人經歷過種族滅絕的痛苦,沃弗很明白什麼叫仇恨,但他也很明白寬恕的意義.他說:「十字架就是神放棄祂的自我,為的是不放棄人類；它〔十字架〕是神意欲在不用暴力下破除人類冤仇的力量,接納人類進入神聖共融中所帶來的結果.」(《擁抱神學》(台北:校園,2007),頁215)沒有這十字架的寬恕力量,人類就在仇恨的輪迴中打滾.

在寬恕者的眼中,看見敵人生命中的神的形象.在復仇者的眼中,只看到敵人就是魔鬼.沒有寬恕的心,只會覺得自己是神,仇敵是魔鬼.正邪不兩立,暴力的循環永不止息.要重建一切關係,先有寬恕的心.

大膽地放下仇恨,主動地寬恕,是基督徒的特質.念誦主禱文,就是願意活出這特質.人看見這種特質,便會看到神的國像微小的種子般在發芽.

「不叫我們遇見試探」,按釋經學者的解釋,這經文的重點不在「不遇見試探」,而在「遇見試探時不陷在其網羅裡」(參耶利米:《主禱文真義探索》(台北:使者出版社,1988),頁59.)其實,馬太福音的版文多加的一句「救我們脫離凶惡」,便解釋了禱文的重點不在沒有試探,而是面對考驗時不致跌入陷阱裡.這禱文涉及信仰的堅持.

人生在世,艱難重重.難道有少許挫折,便去求靠風水命理,放棄我們的天父.難道有少許不如意,便去訴之暴力,放棄十架的大愛.難道有少許失敗,便偏離正途,不再尊神的名為聖.

主禱文結尾的這一句,是求神保守我們,在人生的艱難中,不至失腳.我們的父是會不斷加我們力量的,問題只是我們自甘墮落.

父啊,被尊為聖的是你的名,來臨在世間的是你的國.我們日用的飲食,天天賜給我們.求你赦免我們的罪,正如我們也赦免所有虧負我們的人.不要叫我們在種種考驗中墮落.

但願:榮耀歸於聖父,聖子,聖靈.阿們.



\newpage

\section{泰然安心何處求}
\label{sec:212}
\textbf{林豪恩先生 - 泰然安心何處求}
\newline
\newline
\hyperref[sec:211]{< < < PREV SERMON < < <}
~
\hyperref[sec:toc]{[返目錄]}
~
\hyperref[sec:213]{> > > NEXT SERMON > > >}


路加福音 12:13-21
\newline
\begin{longtable}{cl}
\hline
\hline
章節 & 經文 (和合本修訂版)\\
\hline
12:13 & \begin{tabularx}{0.7\textwidth}{X} 人群中有一個人對耶穌說：「老師！請你吩咐我的兄弟和我分家產。」 \end{tabularx} \\ \\ \relax
12:14 & \begin{tabularx}{0.7\textwidth}{X} 耶穌對他說：「你這個人！誰立我作你們的判官，或給你們分家產的呢？」 \end{tabularx} \\ \\ \relax
12:15 & \begin{tabularx}{0.7\textwidth}{X} 於是他對他們說：「你們要謹慎自守，躲避一切的貪心，因為人的生命不在於家道豐富。」 \end{tabularx} \\ \\ \relax
12:16 & \begin{tabularx}{0.7\textwidth}{X} 然後他用比喻對他們說：「有一個財主，田地出產豐富。 \end{tabularx} \\ \\ \relax
12:17 & \begin{tabularx}{0.7\textwidth}{X} 他自己心裡想：『我的出產沒有地方儲藏，怎麼辦呢？』 \end{tabularx} \\ \\ \relax
12:18 & \begin{tabularx}{0.7\textwidth}{X} 就說：『我要這麼辦：要把我的倉庫拆了，另蓋更大的，在那裡好儲藏我一切的糧食和財物， \end{tabularx} \\ \\ \relax
12:19 & \begin{tabularx}{0.7\textwidth}{X} 然後要對我自己說：你這個人哪，你有許多財物積存，可供多年享用，只管安安逸逸吃喝快樂吧！』 \end{tabularx} \\ \\ \relax
12:20 & \begin{tabularx}{0.7\textwidth}{X} 神卻對他說：『無知的人哪！今夜就要你的性命，你所預備的要歸誰呢？』 \end{tabularx} \\ \\ \relax
12:21 & \begin{tabularx}{0.7\textwidth}{X} 凡為自己積財，在神面前卻不富足的，也是這樣。」 \end{tabularx} \\ \\
[1ex]
\hline
\hline
\end{longtable}
講題:泰然安心何處求　The Pursuit of the Worry Free Life

經文:路加福音12章13-21節

講員:林豪恩先生

場合:香港中文大學崇基學院禮拜堂

日期:2013年8月4日

一．引言

在香港生活,你覺得這個城市安全嗎?幾年前,有一對來大學做博士後研究的年青夫婦,認為香港是一個很危險的地方.為甚麼呢?原來,他們來香港之前把幾套「古惑仔」電影都看過,因此,以為黑社會橫行霸道,街頭械鬥,打家劫舍,是香港的常態.

事實上,和其他地區比較,香港算是安全的地方.根據政治及經濟風險顧問公司在2003年11月發表的調查,香港被評定為亞太地區最安全穩定的地方,例如,這裡少有集體暴力衝突,罪案率為亞洲最低.根據香港大學於2011年進行的公眾意見調查結果顯示,有95\%受訪者認為香港的日間環境是安全的,有76\%受訪者認為香港的夜間環境是安全的.(維基百科)

香港是安全的城市,不過,諷刺的是,生活在這個城市中的人,卻非常欠缺安全感,常常感到不安和焦慮.是甚麼令香港人如此缺乏安全感呢?1997年樓股皆升,經濟好景的時候,《社會指標》調查發現,有52.6\%表示經常憂慮會失去工作．另外,50\%認為香港僱員沒保障；65.6\%的受訪者認為難以信任他人；84.7\%感到現在擁有的東西將來將會失去.在經濟好景時,香港人都沒有安全感,更遑論金融風暴之後,泰然安心的日子更是可望而不可及.

雖然是難求,但無憂無慮的生活仍然是作為一個人與生俱來的期望,這期望也是古往今來人類的共同追求.今天的福音經課耶穌就講述了一個相關的故事,並且藉著故事的主角追求安逸快樂的過程,揭示出當中的迷思和陷阱.讓我們閱讀這個故事,以及耶穌後續的教導,尋找一點追求泰然安心的啟迪.

二．經文

12 : 13  人群中有一個人對耶穌說:「老師!請你吩咐我的兄弟和我分家產.」

12 : 14  耶穌對他說:「你這個人!誰立我作你們的判官,或給你們分家產的呢?」

12 : 15  於是他對他們說:「你們要謹慎自守,躲避一切的貪心,因為人的生命不在於家道豐富.」

12 : 16  然後他用比喻對他們說:「有一個財主,田地出產豐富.

12 : 17  他自己心裡想:『我的出產沒有地方儲藏,怎麼辦呢?』

12 : 18  就說:『我要這麼辦:要把我的倉庫拆了,另蓋更大的,在那裡好儲藏我一切的糧食和財物,

12 : 19  然後要對我自己說:你這個人哪,你有許多財物積存,可供多年享用,只管安安逸逸吃喝快樂吧!』

12 : 20  上帝卻對他說:『無知的人哪!今夜就要你的性命,你所預備的要歸誰呢?』

12 : 21  凡為自己積財,在上帝面前卻不富足的,也是這樣.」

三．釋經與應用

1．以偏概全

這個財主追求安逸快樂的第一個迷思是以偏概全,以為只要找到一條鎖匙就能開啟所有的門,找到一樣東西就能滿足所有的需要,解決所有的問題.

12 : 16  「有一個財主,田地出產豐富.

12 : 17  他自己心裡想:『我的出產沒有地方儲藏,怎麼辦呢?』

12 : 18  就說:『我要這麼辦:要把我的倉庫拆了,另蓋更大的,在那裡好儲藏我一切的糧食和財物,

12 : 19  然後要對我自己說:你這個人哪,你有許多財物積存,可供多年享用,只管安安逸逸吃喝快樂吧!』

12 : 20  上帝卻對他說:『無知的人哪!今夜就要你的性命,你所預備的要歸誰呢?』

12 : 21  凡為自己積財,在上帝面前卻不富足的,也是這樣.」

財物的確是維持生命的基本需要,溫飽不保,生存也受到威脅.財主的迷思,不在於他對財物的需求和儲備,而在於他以為財物能夠滿足作為一個人的所有需要.

生活在現代的城市人,其實非常了解這位財主的想法,一方面,我們更體會到錢財的重要,另一方面,我們更受到收入不穩定的威脅.職場的無情,專業科技變換的快速,與中年求職的艱難,警告著這一代的人,必須快速累積財富,否則難以安逸.因此,我們都與財主有著相同的渴望──就是在最短的時候之內可以找到足夠一生使用的錢.有了足夠的金錢,就不會再被焦慮和不安困擾.

最近有一位公眾人士在立法會發言說 :「四五十歲人都沒有物業,沒有土地,是廢人」.話雖然說得難聽,但不幸地這可能確是一眾男士的焦慮.期望快速地累積財富,不一定自私,作為男士,上有高堂,不有妻兒,我們只是盼望自己能夠盡責任,賺得足夠的金錢供養父母,養妻活兒而已.

將心比己,故事中的財主的想望其實乃人之常情,並不為過.不過,耶穌卻知道人很容易陷入替代的困境,就是當我們以為有一樣東西能夠解決問題,我們便會以這東西去替代甚他事物,縱使知識上我們明白甚他事物都有其重要性,但現實的結果往往就是忽略了其他事情.例如:一個人把控制體重的把盼望完全寄託於減肥藥的時候,他很可能會以藥物去代替控制飲食和做運動.因此,耶穌指出,財主的想法雖然看來合理,卻是危險的,因為他「為自己積財,在上帝面前卻不富足」.(V21)「一生何求」這首歌的其中兩句歌詞也道出了這種危機:「一生何求?常判決放棄與擁有,耗盡我這一生,觸不到已跑開,一生何求?迷惘裡永遠看不透,沒料到我所失的竟已是我的所有」.有些人拼命賺錢以改善生活,卻因此而賠上了健康,結果沒有機會享受預期金錢會帶來的快樂.有些人以為用盡心思和時間賺多點錢為家人帶來多一點幸福,不幸的是,家庭關係可能因為他的缺席而受到影響.

耶穌評論過財主的想望和行動後,接著以此教導他的門徒:

12 : 22  「所以,我告訴你們,不要為生命憂慮吃甚麼,為身體憂慮穿甚麼.

12 : 23  因為生命勝於飲食,身體勝於衣裳.

耶穌的請門徒們想一想真正重要的是甚麼?努力賺錢最終的目的是甚麼?下面這個故事可能也會給我們一點啟發.

有一個美國商人在墨西哥海邊一個小漁村的碼頭上,碰上一個打魚回來的漁夫.他問漁夫:「要多少時間才能抓這麼多?」

漁夫說:「才一會兒功夫就抓到了.」

美國人再問:「你為什麼不多抓一些魚?」

漁夫覺得不以為然:「這些魚已經足夠我一家人吃了!」

美國人又問:「那麼你一天剩下那麼多時間都在幹什麼?」

漁夫解釋:「跟孩子玩一玩,再跟老婆睡個午覺,晚上到村子裡跟兄弟們喝點小酒,玩玩吉他,我的日子可過得充實忙碌呢!」

美國人幫他出主意:「我是美國哈佛大學企管碩士,我倒是可以幫你忙.你應該每天多花一些時間去抓魚,到時候你就有錢去買條大一點的船,就可以抓更多魚,再買更多魚船,然後你就可以擁有一個漁船隊.到時候你就不必把魚賣給魚販子,而是直接賣給加工廠.然後你可以自己開一家罐頭工廠,控制整個生產,加工處理和行銷.然後你可以離開這個小魚村,搬到墨西哥城,再搬到洛杉磯,最後到紐約.在那裡經營你不斷擴充的事業.」

漁夫問:「這又花多少時間呢?」

美國人回答:「十五到二十年.」

「然後呢?」

美國人笑著說:「然後你就可以在家當皇帝啦!時機一到,你就可以宣布股票上市,把你的公司股份賣給投資大眾.到時候你就發啦!就可以幾億幾億地賺!」

「然後呢?」

美國人說:「到時候你就可以退休啦!你可以搬到海邊的小漁村去住.出海隨便抓幾條魚,跟孩子們玩一玩,再跟老婆睡個午覺,晚上到村子裡喝點小酒,跟哥兒們玩玩吉他囉!」

漁夫疑惑地說:「我現在不就是這樣了嗎?」

有一些事物很重要,但它可能只是成就目的的一種工具,耶穌提醒我們不要以此取代目的本身,不經意地取代其他事物,避免陷入財主「為自己積財,在上帝面前卻不富足」的危機.

2．自我掌控取代上天恩典

12 : 18  就說:『我要這麼辦:要把我的倉庫拆了,另蓋更大的,在那裡好儲藏我一切的糧食和財物,

12 : 19  然後要對我自己說:你這個人哪,你有許多財物積存,可供多年享用,只管安安逸逸吃喝快樂吧!』

12 : 20  上帝卻對他說:『無知的人哪!今夜就要你的性命,你所預備的要歸誰呢?』

財主的故事讓我們想到第二個尋找泰然安心的生活的迷思,就是要求自己掌控未來,要求自己解決自己的所有需要,不假外求.這樣的要求是一種迷思,也反映著財主「為自己積財,在上帝面前卻不富足」的無奈後果.如果不能信任上帝,甚至把上帝排拒於世界之外,人除了要求自己,還可以仰望甚麼呢?財主這方面的迷思,是可以明白的.自給自足令我們感到安全?還是要看別人面色感到安全呢?未來按自己的計劃實現令我們感到安全,還是受制於變幻莫測的未來安全呢?然而,財主的故事卻提醒我們,危機在於當我們越追求自給自足,越要求掌未來的時候,我們其實正在為自己建造一個封閉的世界,一個不受外力影響的世界.可惜這樣的世界只是我們的構想,現實並不存在.其中一個必然不能排拒的外力就是上帝,正如故事中上帝對財主說的話所揭示的:『無知的人哪!今夜就要你的性命,你所預備的要歸誰呢?』如果我們好像這位財主一樣,要求排拒所有外在的干擾,要求絕對的確定才感到安心,那麼,我們就只有繼續焦慮不安地生活.因此,耶穌請門徒反省:

12 : 25  你們哪一個能藉著憂慮使壽數多加一刻呢?

12 : 26  這最小的事你們尚且不能做,何必憂慮其餘的事呢?

人需要控制未來,才能夠安心,但又不可能控制,可以怎樣呢?耶穌鼓勵他的門徒相信宇宙間有一樣事情叫做恩典,仰望這種恩典去突破我們對人生的那種需要控制但又控制不了的困局.如何發現這種恩典呢?耶穌並不是對門徒保證凡事稱心滿意,而是邀請門徒把眼睛從自己的焦慮轉移望向廣闊的天空:

12 : 24  你們想一想烏鴉:牠們既不種也不收,既沒有倉又沒有庫,上帝尚且養活牠們.你們比飛鳥要貴重得多呢!

12 : 27  你們想一想百合花是怎麼長起來的:它也不勞動,也不紡線.然而我告訴你們,就是所羅門極榮華的時候,他所穿戴的還不如這些花的一朵呢!

12 : 28  你們這小信的人哪!野地裡的草今天還在,明天就丟在爐裡,上帝還給它這樣的妝飾,何況你們呢?

如果要求自我掌控以達至泰然安心是緣木求魚的話,接受世事難料,相信就算有我們不想發生的事情出現,上帝也會以恩典承托,會否是另一個更好的選擇呢?懷著這樣的信念,會否更能夠幫助我們面對順逆不定的人生呢?

在一個中國南方的農村,父母都以能夠為兒子建屋娶妻為己任和成就.一個貧窮的家庭要建一間屋,可以想像其中的沙石磚瓦需要多少年月去累積.為了兒子將來成家立室,這位媽媽在兒子斷奶後,就開始出海掏海.那一次,這位媽媽又再出海,小艇不知何故翻沉了.當時兒子尚幼,她等不到兒子叫一聲媽媽就離世了.

如果媽媽在天有靈,最掛心的,可能並不是兒子甚麼時候懂得叫媽媽,而是沒有父母愛護的兒子能否活下去?如何活下去?

開放的世界是會有意料不到的事情發生的,然而,正因為這世界是開放的,恩典就能夠進入這個世界,讓相信恩典的人能夠不被惶恐及不安吞噬.

這位媽媽的兒子雖然沒有自己的房屋,但仍然活了下來；就是現在,他仍然沒有屬於自己的居所,卻有人願意嫁給他,讓媽媽對他成家立室的心願實現；雖然他沒有經驗過父母的愛護,但沒有阻礙他學習成為一位負責任的父親.

如果這位媽媽在天有靈,她會看見上帝在他兒子的人生路上擺放了恩典,這不是她計劃的,也不是她能力所及的.當她的兒子讀小學的時候,老師發現這位男孩營養不良想幫助他,於是請這位男孩放學後留校協助老師,藉此帶他去吃營養豐富的午餐.

耶穌藉著這個財主的故事提醒我們,不用為自己建構一個完全掌控的世界,以獲取安心,相信上帝以恩典相待,並且方言觀看,發現恩典,才是通往安心的路.

3．焦慮自保取代人生追求

財主的故事揭示第三種追求安逸快樂的迷思,正正就是因為追求自己的安逸快樂.有些事情非常吊詭,一個以要求自我掌控的人,往往更容易被掌控不了的事情激怒；一個追求無憂無慮的人,往往有憂慮不盡的事情；一個以無病無痛為首要關注的人,往往經驗都身體不適帶給他的苦惱.故事中的財主花盡心思追求安安逸逸吃喝快樂的生活,始終都嘗不到安安逸逸吃喝快樂的滋味.這並不是因為財主沒有能力,以今天商業社會的角度來看,他非常有能力的生意人.我們試想想,這年農作物豐收,對於當時的農村經濟狀況有甚麼影響?按常理,糧食供應增加,價格便會下降,價格下降,那些平日餓著肚子的窮苦大眾,就可以吃多一點.可惜的是,村民和工人並沒有得享豐收的好處.為甚麼呢?因為財主囤積居奇,把糧食待價而沽,藉此賺取更高的利潤.12 : 20  上帝卻對他說:『無知的人哪!今夜就要你的性命,你所預備的要歸誰呢?』  這位財主不願意與村民分享豐收的成果,後果是他自己也享受不到豐收的成果.

耶穌接著教導他的門徒若要突破憂慮懼怕的奴役,就要重新審視人生的追求,調校生命的方向和目標.

12 : 29  你們不要求吃甚麼,喝甚麼,也不要掛慮.

12 : 30  這都是世上的外邦人所求的；你們需要這些東西,你們的父都知道.

12 : 31  你們只要求他的國,這些東西就必加給你們了.

12 : 32  你們這小群,不要懼怕,因為你們的父樂意把國賜給你們.

12 : 33  你們要變賣財產賙濟人,為自己預備永不壞的錢囊和用不盡的財寶在天上,就是賊不能近,蟲不能蛀的地方.

12 : 34  因為你們的財寶在哪裡,你們的心也在哪裡.」

有一位女士的經歷告訴人們,過份注視自己未必能夠幫助自己；反而把眼目關懷他人,把自己所有的,甚至是所失去的與人分享,竟然成了自己的拯救.

這位女士十多年前在兩年內先後失去了兒子和丈夫,她傷心欲絕,更懷疑患上了抑鬱症.為了不想女兒擔心,她於是主動向有關組織求助,希望能早日走出陰霾.經過兩年多的時間,內心仍會十分傷感.與其坐著甚麼都不幹,她接受邀請做義工關心跟自己同病相憐的喪親人士.當義工幫助別人,原以為自己只是施予,別人接受和得益.殊不知在幫助其他喪親人士的過程中,陳太覺得本身的傷痛同樣得到醫治.她眼見別人的情況跟自己的遭遇相似時,不是自憐自己的悽涼,反而覺得世界上還有其他人同樣受苦,自己並不孤單,生命的苦也本是平常.被陳太關懷的喪親人士,也覺得有過來人的明白和支持,便能夠重拾希望和信心,繼續努力活下去.她不是忘記了自己的悲傷,只是她將內心的悲傷化為幫人的力量,並為自己不幸的遭遇賦予新的意義,這樣,她跳出了生命的幽谷,活出人生的彩虹.(節錄自:黃志安「化悲傷為力量」)

四．總結

耶穌曾經對他的門徒說:「得著性命的,要喪失性命；為我喪失性命的,要得著性命.」(馬太福音10:30).通往泰然安心的道路,就是把自己交托給那位以恩典待我們的上主,願他的美意在我們身上成就.



\newpage

\section{警醒警醒再警醒}
\label{sec:213}
\textbf{賀志勇博士 - 警醒,警醒再警醒!}
\newline
\newline
\hyperref[sec:212]{< < < PREV SERMON < < <}
~
\hyperref[sec:toc]{[返目錄]}
~
\hyperref[sec:214]{> > > NEXT SERMON > > >}


路加福音 12:35-48
\newline
\begin{longtable}{cl}
\hline
\hline
章節 & 經文 (和合本修訂版)\\
\hline
12:35 & \begin{tabularx}{0.7\textwidth}{X} 「你們要束緊腰帶，燈也要點著， \end{tabularx} \\ \\ \relax
12:36 & \begin{tabularx}{0.7\textwidth}{X} 好像僕人等候自己的主人從婚宴上回來。他來叩門，就立刻給他開門。 \end{tabularx} \\ \\ \relax
12:37 & \begin{tabularx}{0.7\textwidth}{X} 主人來了，看見僕人警醒，那些僕人就有福了。我實在告訴你們，主人會叫他們坐席，自己束上腰帶，前來伺候他們。 \end{tabularx} \\ \\ \relax
12:38 & \begin{tabularx}{0.7\textwidth}{X} 他或是半夜來，或是天亮之前來，看見僕人這樣，那些僕人就有福了。 \end{tabularx} \\ \\ \relax
12:39 & \begin{tabularx}{0.7\textwidth}{X} 你們要知道，一家的主人若知道賊甚麼時候來，就不容賊挖穿房屋。 \end{tabularx} \\ \\ \relax
12:40 & \begin{tabularx}{0.7\textwidth}{X} 你們也要預備，因為在你們想不到的時候，人子就來了。」 \end{tabularx} \\ \\ \relax
12:41 & \begin{tabularx}{0.7\textwidth}{X} 彼得說：「主啊，這比喻是對我們說的呢？還是也對眾人呢？」 \end{tabularx} \\ \\ \relax
12:42 & \begin{tabularx}{0.7\textwidth}{X} 主說：「那麼，誰是那忠心又精明的管家，主人要派他管理自己的家僕，按時定量分糧給他們的呢？ \end{tabularx} \\ \\ \relax
12:43 & \begin{tabularx}{0.7\textwidth}{X} 主人來到，看見僕人這樣做，那僕人就有福了。 \end{tabularx} \\ \\ \relax
12:44 & \begin{tabularx}{0.7\textwidth}{X} 我實在告訴你們，主人要派他管理所有的財產。 \end{tabularx} \\ \\ \relax
12:45 & \begin{tabularx}{0.7\textwidth}{X} 如果那僕人心裡說『我的主人會來得遲』，就動手打僮僕和使女，並且吃喝醉酒， \end{tabularx} \\ \\ \relax
12:46 & \begin{tabularx}{0.7\textwidth}{X} 在想不到的日子，不知道的時候，那僕人的主人要來，重重地懲罰他，定他和不忠心的人同罪。 \end{tabularx} \\ \\ \relax
12:47 & \begin{tabularx}{0.7\textwidth}{X} 僕人知道主人的意思，卻沒預備，又未順他的意思做，那僕人要多受責打； \end{tabularx} \\ \\ \relax
12:48 & \begin{tabularx}{0.7\textwidth}{X} 至於那不知道而做了當受責打的事的，要少受責打。多給誰，就向誰多取；多託誰，就向誰多要。」 \end{tabularx} \\ \\
[1ex]
\hline
\hline
\end{longtable}
講題:警醒,警醒,再警醒　Be Watchful, Again and Again!

經文:路加福音12章35-48節

講員:賀志勇博士

場合:崇基學院禮拜堂主日崇拜

日期:2013年8月11日

12:35    你 們 腰 裡 要 束 上 帶 , 燈 也 要 點 著 ．  12:36    自 己 好 像 僕 人 等 候 主 人 , 從 婚 姻 的 筵 席 上 回 來 ． 他 來 到 叩 門 , 就 立 刻 給 他 開 門 .12:37    主 人 來 了 , 看 見 僕 人 儆 醒 , 那 僕 人 就 有 福 了 ． 我 實 在 告 訴 你 們 , 主 人 必 叫 他 們 坐 席 , 自 己 束 上 帶 , 進 前 伺 候 他 們 .  12:38    或 是 二 更 天 來 , 或 是 三 更 天 來 , 看 見 僕 人 這 樣 , 那 僕 人 就 有 福 了 .  12:39    家 主 若 知 道 賊 甚 麼 時 候 來 , 就 必 儆 醒 , 不 容 賊 挖 透 房 屋 , 這 是 你 們 所 知 道 的 .12:40    你 們 也 要 預 備 ． 因 為 你 們 想 不 到 的 時 候 , 人 子 就 來 了 .  12:41    彼 得 說 , 主 阿 , 這 比 喻 是 為 我 們 說 的 呢 , 還 是 為 眾 人 呢 .  12:42    主 說 , 誰 是 那 忠 心 有 見 識 的 管 家 , 主 人 派 他 管 理 家 裡 的 人 , 按 時 分 糧 給 他 們 呢 .  12:43    主 人 來 到 , 看 見 僕 人 這 樣 行 , 那 僕 人 就 有 福 了 .  12:44    我 實 在 告 訴 你 們 , 主 人 要 派 他 管 理 一 切 所 有 的 .  12:45    那 僕 人 若 心 裡 說 , 我 的 主 人 必 來 得 遲 ． 就 動 手 打 僕 人 和 使 女 , 並 且 喫 喝 醉 酒 ．  12:46    在 他 想 不 到 的 日 子 , 不 知 道 的 時 辰 , 那 僕 人 的 主 人 要 來 , 重 重 的 處 治 他 , 〔 或 作 把 他 腰 斬 了 〕 定 他 和 不 忠 心 的 人 同 罪 .  12:47    僕 人 知 道 主 人 的 意 思 , 卻 不 預 備 , 又 不 順 他 的 意 思 行 , 那 僕 人 必 多 受 責 打 ．  12:48    惟 有 那 不 知 道 的 , 作 了 當 受 責 打 的 事 , 必 少 受 責 打 ． 因 為 多 給 誰 , 就 向 誰 多 取 ． 多 託 誰 , 就 向 誰 多 要 .

引言  人生需要預備

最近跟幾個朋友一起吃飯,聊起想生孩子的事.那我就建議她們,想生孩子,要預備什麼什麼.她們就很懷疑地看著我,說‘你一個男人,又沒生過孩子,怎麼懂得生孩子的事’.我笑著說:“很多婦產科醫生都是男人,也沒生過孩子,還不是照樣可以幫人接生”.她們笑著點頭稱是.城市人速度太快,工作繁忙,食物又沒有以前健康,懷個孩子普遍比以前難,時間啊,營養啊,心情啊,身體啊,統統都要預備好才行.這個還算好,最怕是不想生的時候突然就懷上了,這下麻煩就更大了.最近有個朋友向我諮詢,她在國內大學教書,是公家人,不可以生第二胎,偏偏又不小心懷上了第二胎.作為基督徒,她不想墮胎；可作為大學老師,她如果把孩子生下來,就立即會被開除.她一點預備都沒有,問我該怎麼辦?弟兄姊妹,你說我該怎麼回答她好呢?

做人最怕什麼?最怕是請人吃飯自己卻沒帶錢,餐廳說“我們這裡不收卡”；最怕是要上飛機了,卻發現自己沒帶證件；要考試了,自己書還沒看幾頁.一個人沒有預備,就容易到處撞板；那一間教會呢,一間沒有預備的教會又會怎樣呢?我們今天來看一段經文,就是一段有關預備的經文,請大家一起翻到路加福音12:35-48節.

解經

一,             家主與僕人(35-38)

這是一段非常具有緊迫性的經文.35節一開始就提到 “你 們 腰 裡 要 束 上 帶 , 燈 也 要 點 著 ．” 這句話直接指向了出埃及記的逾越節的傳統.在出埃及記第十二章就提到“這夜是耶和華的夜”,以色列民要點著燈,束上腰帶,趕緊吃羊羔,因為耶和華隨時要帶領他們離開埃及.耶穌引用這句,是想告訴聽眾,救恩馬上就要完成了, 但是唯有警醒等待的人才可以得到.因為,如果耶和華來了,你還在睡覺的話,那你怎麼跟著他出埃及呢?

在35-38節這個比喻中,主人是從婚宴中回來的.婚宴,在路加福音裡一般是指彌賽亞的婚宴,也就是指天國來臨的意思.馬太福音中十個童女的故事也有這種意思.猶太人的婚宴要吃很多天,可以吃一天兩天,最多的可以吃滿七天(這種喜酒吃起來比較划算,值回了禮金).所以主人出去喝喜酒,僕人就不知道主人具體什麼時候回來了.  好的僕人,就會隨時預備,等候主人回來.

你有沒有試過夜晚回到家,忘記帶鑰匙,卻怎麼敲門都沒有人開門的經驗?或者你有沒有試過下班回到家裡,菲傭不知道哪裡去了,沒人煮飯,小朋友也不知道哪裡去了,打菲傭電話又不通? 如果你有過這樣的經驗,你就會理解主人的心情了.

再打個比方,有一天你很晚回家,又忘記帶鑰匙,你想‘回家敲門會不會吵醒家人’,可是一回到家,發現家裡燈還亮著,家裡人都沒睡,還在等著你,你會是什麼感受? 或者你辛苦工作了一天,回到家一敲門,傭人開門了,孩子撲過來抱住你,飯菜都已經準備好了,香噴噴的,你會是什麼感受? 說明白點,你會有一種回家真好的感受!看過電影《桃姐》的人就知道,桃姐幾十年盡心盡力,對於那個家庭來講不止是一個傭人,更是家庭不可缺少的一部分.家有老僕,是一種福氣!(家有一老,如有一寶)

一間教會也如此啊!如果你回到教會,發現司事人員一點不盡心,詩班唱詩也不排練,牧師講道都不寫講章,會眾唱詩歌都有氣無力,那你會把這間教會當做是自己屬靈的家麼,神會悅納這樣的敬拜麼?長此以往,這間教會會怎樣呢?所以一間教會有老會友,神的老僕人,也是一種福氣!

弟兄姊妹,這段經文之所以要我們警醒預備,不僅僅是為了主人,也是為了我們自己.因為主人和僕人構成了一個家的關係.張愛玲講“愛是需要經營的”,經營,就是打理生意的意思.意思是說,兩個人相愛,就像打理生意一樣,要小心謹慎,要有精明預備,要有長遠打算.其實,一個家何嘗不是如此呢?,家是幫我們擋風擋雨的地方.惟有我們時常警醒˴ 時刻預備,這個家才有統一的目標,才會彼此相愛,互相支持,才能平平安安˴ 幸幸福福地過下去.一間教會也是如此,保羅說‘我們是神的僕人,聚在一起就成了神的殿,神的家’.惟有我們都警醒˴ 預備好了,我們才能同心合意等候上主的再來,那主耶穌再來的時候,就會好好伺候我們了,就像故事裡家主伺候僕人一樣.

二,             賊人與家主(39-40)

39-40節的比喻是說主耶穌會像賊一樣到來,你說比喻什麼不好呢,主耶穌把自己比喻成賊.另外,巴勒斯坦的賊這麼厲害麼,偷東西把人家的房子都挖穿了?話說回來,巴勒斯坦的房子確實沒有香港這麼結實,大家記不記得路加5章提到有人拆房頂,放下癱子來請耶穌醫治的故事呢!    但就算這樣,我們中國人還是覺得,把主比喻成賊讓人不舒服!   實際上,這個比喻也是指向出埃及記的傳統.路加11:22節提到主的到來,就像打劫的壯士來了,把家裡搶了分贓一樣,都是指向同一個傳統.在出埃及記3:22-23節,上帝吩咐以色列人搶走埃及人的財物,因為埃及人長期欺壓以色列民,因此到出埃及的時候,上帝的審判就臨到埃及人了.也就是說,路加的兩個比喻,無論是壯士搶走人家的財產,還是賊偷走人家的財產,都指向上帝公義的審判.不過在這個比喻中,路加更強調:上帝的審判是在你意想不到的時候就來到了.

歷史學家湯因比說過“每個人都有一個錯覺,覺得自己總會有明天,但現實是,每個人有一天都會沒有明天.我們沒有一個人知道那天是哪一天,聰明的人就會懂得為那一天預備!” 我很小就怕死,只要下雨又天黑的時候,我就會很害怕,心想會不會有鬼來啊.

基督徒倒是不畏懼死亡,但基督徒可以學習一種態度,那就是每一天都活得像最後一天一樣.如果這是我們的最後一天,我們就不會花那麼多時間去看跟歷史完全沒關係的《大太監》啊,《宮心計》啊；就不會花那麼多時間去看股票,或者打麻將了.因為我們知道,這是最後一天了,我們還有很多事要做--如果我得罪了太太,我就趕快去和好；如果我欠了父母,就趕快去報恩；如果我想旅行,就趕快去訂票.人生很多無聊的應酬,虛偽的面具都快快丟掉,只有最後一天了嘛!如果我們有這種“只剩最後一天”的心態,我們就明白這第二個比喻的意思,因為你不知道上帝的審判究竟哪一天會到來,所以你每天都需要好好預備,為那個最後一天做準備!

有個朋友跟我講,你們基督教講“只要悔改就可以得救”,那我不如現在多做點壞事,多點享受,等我要死的時候悔改就可以了.我笑著對他說:“你是上帝麼?你能定自己什麼時候是最後一秒啊?你要能定,還需要上帝做什麼?”弟兄姊妹,大家說是不是呢?

三,             主人與管家(41-46)

那有人就會問了,“是不是警醒等候,等在那裡就可以了?”我回大陸的時候經常有人問我:  “你信主多少年了?”好像信主越久功力就越高,內力就越強一樣!但實際上信主時間越長,就越容易出一個問題:容易變成信仰上的老油條!來一天教會就撞一天鐘.我認得一個姊妹,她信主好些年了,我跟她講:“你真的是准點王,每次必定遲到一個小時,太準時了!”

41-46節這個比喻想告訴我們,警醒等候,並不是光等在那裡就可以了,因為等著等著我們可能就喪失最初的熱情,變成老油條了；但我們還有更重要的職分,就是在這地上做神的管家.在這個比喻中,主耶穌引用了一個忠心管家的楷模形象--約瑟.創世記39章講,約瑟很有才幹,他的主人波提乏把家中的事情都交給他打理；約瑟也很忠心,不論主人在家不在家,他都處理得很好.所以很明顯,忠心的管家會把 主人真正當主人,不但不會侵吞主人的財產,還會管理好家務,“按時分糧”；而“吃喝醉酒,還動手打下人”的僕人,其真正的心態是把自己當主人,不把主人放在眼裡.如果一家公司,一間教會,有這樣的經理或者管理人,是不是很危險?

況且,我們每個人都有一間更重要的公司,那就是我們自己!我們又是如何對待我們自己的恩賜,比方說才能,子女,甚至是時間呢?我們有沒有管理好上帝的這些禮物呢?我剛才提到的那位準時遲到一小時的姊妹,她在管理時間方面是不是有些“吃喝醉酒”呢?如果父母不好好教育子女,反而經常責打,是不是“動手打僕人和使女”呢?路加12章就講了一個“無知財主的比喻”,比喻中的那個財主有 很多很多糧食,多到沒有地方收藏了.他就告訴自己 ‘要建一個更大倉房收藏糧食’,然後對自己靈魂說,‘你有許多的財物積存,可作多年的費用,只管安安逸逸地吃喝快樂吧!’神卻對他說:“蠢才,今夜必要你的靈魂,你所預備的要歸誰呢?”這個財主是不忠心僕人的典型.他最可悲的是,他以為自己是財物的主人,結果卻不是.他不曉得自己僕人的身份,最後不但丟了財物,還丟了性命.

結論:警醒,警醒,再警醒(47-48)

話說到這裡,大家可能會想,天天要警醒警醒再警醒,天天要把這天當最後一天來過,做基督徒是不是很累?我們天天吃飯覺得累不累呢?其實,天天警醒,就跟天天跑步一樣,只不過天天跑步是鍛煉身體,天天警醒是鍛煉靈性.我神學院的老師孫寶玲牧師,天天跑步,都堅持幾十年了,馬拉松都跑了10次了.他不但不覺得天天跑步累,他還覺得一天不跑步就不舒服!所以說,凡事養成了習慣,你不但不覺得累,你反而能在其中得到享受.孔子說:“七十,從心所欲不逾矩!”也就是說,當他70歲的時候,他 可以在守規矩的情況下隨心所欲,而且還很舒服.所以,主耶穌的話不是要給我們捆綁,反而是要給我們自由.真正的自由乃是養成一種良好的習慣!警醒˴  警醒再警醒,不僅僅是為了主的再來,同時,也是為了鍛煉我們靈性的好習慣!

回到我講道開始時那個姊妹的疑問,她問我“她可不可以墮胎?”.我告訴她,你來問我這個問題,說明主的誡命你都知道了,但關鍵不是你知道不知道這個誡命,關鍵是你 是個什麼樣的人! 如果你沒有預備好承擔這份責任,那麼就算我再強調主的誡命,你也不會遵守；但如果你已經預備好了自己,無論孩子生下來會遭遇到怎樣的困苦,你都願意跟家人,跟弟兄姊妹一起去面對,那你自然會遵守主的誡命.所以關鍵不在於我的答案,關鍵在於你是怎樣的人!

弟兄姊妹,讓我們一起來讀47-48節,這就是我的結論!



\newpage

\section{在和平與紛爭的時代作信仰解讀}
\label{sec:214}
\textbf{鄭漢文博士 - 在和平與紛爭的時代作信仰解讀}
\newline
\newline
\hyperref[sec:213]{< < < PREV SERMON < < <}
~
\hyperref[sec:toc]{[返目錄]}
~
\hyperref[sec:215]{> > > NEXT SERMON > > >}


路加福音 12:49-56
\newline
\begin{longtable}{cl}
\hline
\hline
章節 & 經文 (和合本修訂版)\\
\hline
12:49 & \begin{tabularx}{0.7\textwidth}{X} 「我來是要把火丟在地上，假如已經燒起來，不也是我所希望的嗎？ \end{tabularx} \\ \\ \relax
12:50 & \begin{tabularx}{0.7\textwidth}{X} 我有當受的洗還沒有受，在這事完成之前，我是多麼地焦急！ \end{tabularx} \\ \\ \relax
12:51 & \begin{tabularx}{0.7\textwidth}{X} 你們以為我來是要使地上太平嗎？不！我告訴你們，是使人紛爭。 \end{tabularx} \\ \\ \relax
12:52 & \begin{tabularx}{0.7\textwidth}{X} 從今以後，一家五個人將要紛爭，三個和兩個相爭，兩個和三個相爭： \end{tabularx} \\ \\ \relax
12:53 & \begin{tabularx}{0.7\textwidth}{X} 父親和兒子相爭，兒子和父親相爭；母親和女兒相爭，女兒和母親相爭；婆婆和媳婦相爭，媳婦和婆婆相爭。」 \end{tabularx} \\ \\ \relax
12:54 & \begin{tabularx}{0.7\textwidth}{X} 耶穌又對眾人說：「你們看見西邊起了雲彩，就說：『要下大雨了』，果然就有； \end{tabularx} \\ \\ \relax
12:55 & \begin{tabularx}{0.7\textwidth}{X} 起了南風，你們就說：『要燥熱了』，也就有了。 \end{tabularx} \\ \\ \relax
12:56 & \begin{tabularx}{0.7\textwidth}{X} 假冒為善的人哪，你們知道分辨天地的氣象，怎麼不知道分辨這是甚麼時代呢？」 \end{tabularx} \\ \\
[1ex]
\hline
\hline
\end{longtable}
講題:在和平與紛爭的時代作信仰解讀　Christian reading of times in between peace and division

經文:路加福音12章49-56節

講員:鄭漢文博士

場合:香港中文大學崇基學院禮拜堂

日期:2013年8月21日

1. 誰的秩序?

1.1

從前有一個藝術家,他的創作室在別人眼中(包括他的管家眼中)看來是一遍混亂.

但是,對於這個藝術家來說,他的創作室是井然有序,各安其位.他要找的物品,他從來不會找不到.

有一天,在他外出時,管家請了個鐘點女傭來,這個女傭把這個創作室執拾得異常有秩序,不單如此,這個女傭也把她看來無用的物品全數棄掉.

藝術家回家後,發現了創作室一遍混亂,他要找的物品他找不到,他所收藏的藝術品全部被消失了.原先的秩序,被這位全無藝術修養的臨時工給打亂了.

請問大家,藝術家可以怎麼樣?

2.1

每個人都以為自己的秩序才是秩序,對神所創造本有的美好,對本有的秩序,未必認識,就算認識,也未必認同.

今日的講道,題目叫做:「在和平與紛爭的時代作信仰解讀」.

英文是:Christian reading of times in between peace and division.

今日的經課是來自路加福音第12章49至56節.今日我想用和合修訂本的聖經譯文,這段經文是這樣開始:

耶穌說:「我來是要把火丟在地上,假如已經燒起來,不也是我所希望的嗎?．．．你們以為我來是要使地上太平嗎?不!我告訴你們,是使人紛爭.」(路12:49; 51)

耶穌基督所傳的不是和平的福音嗎?點解會起火,點解要使人紛爭?

2. 聖父所創的天地,和平原是本性,紛爭是時代常態

2.1

2013年年頭香港有一位基督徒,提出「萬人佔領中環爭取真普選」的構想,一起火頭,紛爭至今.點解要起火?是甚麼因素使人紛爭?星星之火,可以燎原.

2013年年頭繼「和平佔中」運動的出現後,分出了三方取態:有支持的,有反對的,有態度未決的.

支持的又分為三方取態:有信任而跟從的,有爭取領導權的,有從中找機會加入自已取向的.

至於態度未決的又分為三方取態:有些想撘方面車的,有些不想太被認清立場的,也有找機會令自已陣營坐大的.

一石激起千重浪.

一路走來,這位戴耀延教授把所想所作的寫為文字,出版了一本書(舉起此書),書名叫:《佔領中環:和平抗爭心戰室》,書內是這樣介紹這位作者:「香港大學法律系副教授.主要研究範圍包括:香港特別行政區基本法,比較憲法,行政法,人權法,人權教育,公民教育,法律與宗教,法律與管治,法律與政治,法律與發展.公共服務方面,曾任基本法諮詢委員會委員,公民教育委員會委員,中央政策研究組非全職顧問.」

從這段履歷認識這教授,幾乎可以把他看成是建制權力的一部份,是既得利益的集團人士,至少是和平分子,維護人間的和諧共處,真的無法想像一個教法治的學者叫人犯法,主張公民抗命──幸好,他主張的是非暴力的和平抗爭.

2.2

這位書生,為何忽然論政?不單論政,為何忽然帶領一個政治運動?一介文人,為何忽然抗爭?他的教會如何教導他的,不是應該傳揚耶穌基督使人和平的福音的嗎?

火來了照耀人間,一耀不夠兩耀來.

他不單一人倡議,而且還找來朱耀明牧師,聯同香港中文大學教授陳健民,三人發起正式名為「讓愛與和平佔領中環」(Occupy Central with Love and Peace)的運動.

你有佔中三子,我有愛字輩人.

和平佔中運動是主張審議式民主(deliberative democracy)的,所以有必要主辦公開論壇.因而出現不同的反對者到場唱反調,也就出現「愛」字頭的群組不斷紛爭.愛群之後,出現新幫.上星期又出現「幫幫香港出聲行動」,以六人的陣勢對抗三子.

對抗的氣氛因為炎熱的夏天更形熾熱,七月中旬,在旺角街頭,一位林老師對警隊執法不公的激烈言論後,出現支持和反對的兩大勢力.政治力量無限上綱,為要打壓異己.

香港現正在進入紛爭的時代.

其實在香港一百五十多年以前,開埠以來都是紛爭的時代,其中有時稍有的和平歲月都是難能可貴的,就正如中國過去二百年,和平歲月都是難得的.

香港自九七後是全面紛爭的時代.至今只是紛爭的潛在激烈的漸漸浮現.

2.3

夠了夠了,我今天並不是香港政治講壇,不適宜把政治引入教會,教會最好與政治分家,教會傳的是耶穌基督和平的福音,政治難免講的是權與利的分爭.

我曾經被「提醒」(加引號)不要在講壇談政治,大概因為一談政治,信仰群體就會分化,生發紛爭.

Sorry,這裡是崇基禮拜堂,這是大學講章.信仰反思的層次並不一樣.

其實,這種去政治的教會觀從來是行不通的,因為教會存在世上,並不是存在一個真空之中,任何存在人世間的都難逃政治的.許多堂會式的教會內部,政治紛爭更加厲害.

教會的存在,就是分別為聖,與世間的權能有所不同,有所不認同,構成為另一種政治實體.

就算教會傳揚的是和平福音,種籽一落在地上,臨在紛爭時代中也會更添紛爭.

3. 耶穌傳揚的是和平的福音,臨在紛爭時代中也會更添紛爭

3.1

讓我們返回今天的經課,根據和合修訂本,路加福音第12章49至56節,記載耶穌親自說:

12 : 49	「我來是要把火丟在地上,假如已經燒起來,不也是我所希望的嗎?

12 : 50	我有當受的洗還沒有受,在這事完成之前,我是多麼地焦急!

12 : 51	你們以為我來是要使地上太平嗎?不!我告訴你們,是使人紛爭.

12 : 52	從今以後,一家五個人將要紛爭,三個和兩個相爭,兩個和三個相爭:

12 : 53	父親和兒子相爭,

兒子和父親相爭；

母親和女兒相爭,

女兒和母親相爭；

婆婆和媳婦相爭,

媳婦和婆婆相爭.」

12 : 54	耶穌又對眾人說:「你們看見西邊起了雲彩,就說:『要下大雨了』,果然就有；

12 : 55	起了南風,你們就說:『要燥熱了』,也就有了.

12 : 56	假冒為善的人哪,你們知道分辨天地的氣象,怎麼不知道分辨這是甚麼時代呢?」

路加福音傳的是和平的信息,為何耶穌要人紛爭?而且耶穌講明,紛爭不單出現在家的外面,還會出現在家的內面.

我的工作與教會服侍,長期與學校和家庭相關,專家發現青少年問題的成因,許多青少年的問題,有一半的解釋要追究到家庭的影響,另一半因素是因為父母家長的問題.

父慈子孝,母愛女順,向來只是人間的理想.家中跨代的融洽相處,向來都只是心之所願.至於婆媳相爭,基本上是大家庭常態.耶穌為何沒有提到夫婦兩邊的父母親相爭,因為在兒子與女兒在結婚那一天婚禮或婚宴之後,雙方家長都沒有再見面了.正所謂,眼不見為乾淨.

家人是最親密的,尚且難逃分爭.

最好的朋友一起經營生意,總會因利害計算而分手.

商界「比較好」,因為擺明車馬是為利益而相合,因利益而相分.商場沒有永遠的朋友,也沒有永遠的敵人.「比較好」的意思指的是比較真實地呈現分爭.

有人以為,學校是講究教育理想的地方,但是,教師工作產生的身心受傷的比率甚高.

大學呢?大學最講人類文明,最講天地立心.嘿嘿,對不起,我自1980年開始入讀中文大學,由崇基校園到牛津校園,經過倫敦校園再返回中大校園,教學十八年期間我再經歷過許多西歐,北美以及內地校園,我的結論是:大學也是一個最分爭的人間地獄.為甚麼?我不便在這講壇細訴真相,說人是非.

總之,正如耶穌指責:因為這裡有許多假冒為善的人哪!

3.2

在這段經文中耶穌用上了家作為例子,說明紛爭是人間的常態.

在這紛爭是常態的人間世上,耶穌基督就算傳的是叫人和平的福音,這種籽撒在這人間世上,都難逃劫數長出紛爭來.

因為人有罪性,而罪又入了人間世上,世間一切都變成污染的水一樣,多清的白雪飄落一池污水中,可以怎樣?

比如香港某教會邀請生下來就沒有手,沒有腳的Nick Vujicic今年六月四日到香港來分享信仰,坐得近坐得遠捐多少錢,這不也就產生紛爭來.就算六月四日不去點燃燭光,也不需要對著幹,豈不是無端生分爭嗎?

3.3

耶穌說:「一家五個人將要紛爭,三個和兩個相爭,兩個和三個相爭.」耶穌這三二對比應用在香港的政治,十分真實.

香港是我家,香港這個家也有分爭.五個香港人中,會出現三對二的紛爭.以一百為基數,六十對四十之比.在立法議會的投票中,正是這個比例:泛民派的票佔百分之六十的票,百分之四十為建制派的票.然而百分之六十票的泛民派只佔百分之四十議席,百分之六十票建制派卻佔百分之六十的議席.

議席不按民意分配,豈不是分爭嗎?如此怪現況,理由是因為政治體制不公義.

不公義的政治體制不會主動令得特區行政首長的選出會出現公義的程序的,所以有人主張以愛與和平佔領中環爭取一人一票選特首.

4. 如何依靠生命主,學會從紛爭中能作信仰解讀,得以回應?

4.1

紛爭是時代常態,是人間常態,我們信徒如何面對?如何分辨?

今日經課最後三節,路加福音12章54至56節這樣記述:

「耶穌又對眾人說:「你們看見西邊起了雲彩,就說:『要下大雨了』,果然就有；起了南風,你們就說:『要燥熱了』,也就有了.你們知道分辨天地的氣象,怎麼不知道分辨這是甚麼時代呢?」」

耶穌指責當時住在巴勒斯坦的猶太人,看見西邊起了雲彩,就是地中海的天空,就知道要下大雨了,果然正如所料.起了南風,就是沙漠往北吹來,就預測到要燥熱了,也果然正如所料.

對世界的現象懂得分辨,為何不懂得分辨時代的屬靈意義呢?

4.2

懂得分辨時代的屬靈意義,構成了信徒成長的重要栽培目標.(你們在這裡返教會,有這方面的成長嗎?)

我在自己的教會(馬鞍山堂),2009年1月至2011年12月這三年出任青少年部部長,就是青少年事工的信徒領袖.當時協助堂主任牧師的青少年牧者(Youth Pastor)忽然離任,唯有由我暫代其職,全盤衣帶領青少年事工.

我給青少年事工界定為:「以Christian readership的培養作為重心,培養青少年對生命及其所在的時代作基督信仰解讀.」

在我之前的青少年崇拜,信息分享偏向兩類講題,一是聖經經卷的解讀,二是信徒操守的教誨.對於教友子弟來說,他們的感受是,好的叫「舊調天天彈,不壓其煩」,不好的叫「你有你講大道理不難,我有我面對殘酷人間」.

所以,我引入我循道衛理聯合教會總議會倡議的兩個理念:「生命教育」(life education)和「信仰分享」(faith sharing),把信息訂為四種生命的信仰解讀(faith interpretation of four types of storied lives),包括(一)耶穌生平的生命解讀,(二)聖經人物的生命解讀,(三)信徒人物的生命解讀,(四)電影人物的生命解讀,所有信息最後都要回到青少年自身對自己的生命作信仰解讀.

我們並非放下聖經的教導,反而是通過聖經所呈現的信仰境界去閱讀我們信徒所身處的時代.學理地講試探用了敘事解釋學narrative hermeneutics,是文本text與脈絡context的相互詮釋.

每月四個主日,各配四種信息.當然每月第二週主日的電影主日特別多青少年來,連成年教友也借關心青少年而到臨,見證一下這較新派的信息分享.(你們想不想體驗下?)

4.3

而我是親自牽頭,開發電影人物的生命故事的信仰解讀這種青少年崇拜的形式.

例如,最近的一次在2013年7月14日,我是借《孤星淚》或叫《悲慘世界》(Les Misérables)這套音樂劇式的電影,抽出五個片段(共半小時),呈現「悲慘世界中人生五個關鍵時刻的五種祈禱」 (5 Prayers admist 5 Critical Life Moments in the Miserable World) .

故事主角Valjean為救一個他所愛的家姐女兒(甥女)而偷了一片麵包,以致入獄受苦十九年.出獄後,飽受歧視,無法維生,心生仇恨,直至遇上一位主教.Valjean 經歷種種苦困(受獄,受辱),本來充滿怨恨,但因主教的恩待所流露的基督精神,他悔改了,告別舊我,成為新人.曾經困苦,而又真心悔改的人,良知會時刻發用,同時自然地生出恩慈,關心愛顧身邊受困苦,有需要的人.

從悲慘世界中Valjean所經歷的人生,看到人生五個關鍵時刻,就是受苦,悔改,良知,恩慈,垂死,相應所需要這五種禱告:

一是受苦(者)的禱告suffering prayer,為呼喚創造主而祈禱

二是悔改(者)的禱告repenting prayer,為渴求拯救主而祈禱

三是良知(自身)的禱告conscientious prayer,為自己不埋沒良心而祈禱

四是恩慈(他者)的禱告charitable prayer,為他人能得恩待而祈禱

五是垂死(者)的禱告dying prayer,為安返天家而祈禱.

以上是為悲慘世界中的五種祈禱.

當然還有其他不同類型的祈禱,但借助電影中的生命故事,信徒讀者對現實人間的解讀能量得以敏感多了(capacity for faith interpretation being sensitized),以後也學會對著五種生命境況作相應的祈禱.

從路加福音第7章18-24節看到:

耶穌的憐憫,讓悲慘世界中的瞎眼的得看見,耳聾的得聽見,跛腿的得行走,貧窮的得福音.

這就是耶穌基督和平的福音,背景是人間紛爭是常態的悲慘世界.

5. 如何依靠聖靈,學會從紛爭中能作和平之子,得以復和?

5.1

如何面對這個紛爭是常態的悲慘世界?如何作和平之子?

解答就在今日經課的旨句──路加福音12章50節,耶穌說:「我有當受的洗還沒有受,在這事完成之前,我是多麼地焦急!.」

解經家對耶穌所說「我有當受的洗還沒有受」認為有兩種意思,或有兩重意思.

首先是指他在十架上為我們流血而死(Baptism of Passion),以至我們可以因他復活而戰勝死亡,因而得到心靈的和平(peace of soul).

其次是指聖靈的洗禮(Baptism of the Holy Spirit),也就是聖靈與你的生命同在,你所想所作都有聖靈的臨在與充滿(spirited soul).

5.2

經歷過基督和聖靈這兩重意義的受洗,信徒會如何在和平與紛爭的時代作信仰解讀?

我的體會有三重,一點連一點地實踐:

一是首先要認清世間的實況,正如輔導的第一步要認清現實,信徒也要明白,人生在世的人間世紛爭是常態,和平是難得的恩臨.在屬靈的爭戰中(spiritual warfare),作為信徒我們無法中立.我們無心不分爭,因為是世界不容我們.

二是要認同和平是基本的信念,正如「以愛與和平佔領中環運動」是一場非暴力的和平抗爭,但更重要的是心靈的和平.對於我來說,和平佔領中環運動就是要「讓中環價值不再佔領我們香港人的核心價值」,而我們信徒的核心價值就包括心靈的和平,但要能心靈和平就要回到最根本的愛.

三是在紛爭中要活出讓人感到和平的愛,正如在《悲慘世界》中垂死時的Valjean禱告時說:「彼此相愛就是讓人得見上主的面容」,英文是to love one another is to see the face of God.

因為我們信仰上帝的創造本有秩序,信仰基督救贖而更新,信仰聖靈的臨在而平和,就是讓我們在紛爭時代中,依然活出充滿和平信息的大愛.

但願榮耀都歸給聖父,聖子,聖靈,永在的三一神.阿們.

(鄭漢文博士為循道衛理聯合教會馬鞍山堂會友.

現職為香港中文大學教育行政與政策學系專業顧問,

同時在崇基學院神學院兼讀神道學學士課程.

曾於2012年1月至6月及2012年9月至2013年4月

於崇基禮拜堂作神學生實習.

將於2013年9月至2014年4月

到浸會大學校牧室實習.)



\newpage

\section{使人生轉彎為直}
\label{sec:215}
\textbf{許開明牧師 - 使人生轉彎為直}
\newline
\newline
\hyperref[sec:214]{< < < PREV SERMON < < <}
~
\hyperref[sec:toc]{[返目錄]}
~
\hyperref[sec:216]{> > > NEXT SERMON > > >}


路加福音 13:10-17
\newline
\begin{longtable}{cl}
\hline
\hline
章節 & 經文 (和合本修訂版)\\
\hline
13:10 & \begin{tabularx}{0.7\textwidth}{X} 安息日，耶穌在一個會堂裡教導人。 \end{tabularx} \\ \\ \relax
13:11 & \begin{tabularx}{0.7\textwidth}{X} 有一個女人被靈附身，病了十八年，腰彎得一點都直不起來。 \end{tabularx} \\ \\ \relax
13:12 & \begin{tabularx}{0.7\textwidth}{X} 耶穌看見，就叫她過來，對她說：「婦人，你的病好了！」 \end{tabularx} \\ \\ \relax
13:13 & \begin{tabularx}{0.7\textwidth}{X} 於是用雙手按著她，她立刻直起腰來，就歸榮耀給神。 \end{tabularx} \\ \\ \relax
13:14 & \begin{tabularx}{0.7\textwidth}{X} 會堂的主管因為耶穌在安息日治病，就很生氣，對眾人說：「有六天應當做工，那六天之內可以來求醫，在安息日卻不可。」 \end{tabularx} \\ \\ \relax
13:15 & \begin{tabularx}{0.7\textwidth}{X} 主回答他：「假冒為善的人哪，難道你們各人在安息日不解開槽上的牛和驢，牽去喝水嗎？ \end{tabularx} \\ \\ \relax
13:16 & \begin{tabularx}{0.7\textwidth}{X} 何況她本是亞伯拉罕的後裔，被撒但捆綁了十八年，不該在安息日這天解開她的綁嗎？」 \end{tabularx} \\ \\ \relax
13:17 & \begin{tabularx}{0.7\textwidth}{X} 耶穌說這些話，他的敵人都慚愧了；所有的人因他所做一切榮耀的事都很歡喜。 \end{tabularx} \\ \\
[1ex]
\hline
\hline
\end{longtable}
講題:轉彎為直的人生　Changing a Life from Bend to Straight

經文:路加福音13章10-17節

講員:許開明牧師

場合:香港中文大學崇基學院禮拜堂

日期:2013年8月25日

引言:

各位早晨:

我自得悉將成為崇基學院校牧,差不多每個月都有機會証道,感謝伍牧師的安排.

我很欣賞崇基禮拜堂每次皆有三段經課,並且講道以福音書為中心,使信徒得著有系統的講壇牧養.

今日的福音書的經文是路13:10-17,從中可探索何為健康的生命或人生?

各位,你自覺自己的生命健康嗎?

有人對你講過,你好健康嗎?

一個人健康是非常重要的,健康就是財富,健康就是力量,健康就是一切.

首先,何為健康?

一般來說,健康就是沒有疾病.不過,根據世界衛生組織(WHO,1948),“健康就是肉體,心思,精神及人際四方面皆處於良好的狀態,不僅是沒有不適和疾病.”

本段經文啟迪我們,主耶穌能轉彎為直使人得著真正的健康.

在我們生命中如何依靠耶穌,使生命轉彎為直,得著健康的生命?

信息:

今早的經課中,你會被那腰彎女子吸引住……

見經文

在今早福音經課中,我們可以看到四項的轉彎為直:

一 身體的轉彎為直

安息日耶穌在會堂裡証道,祂發現「有一個女人,被鬼附著病了十八年；腰彎得一點直不起來.」路十三11

這女子腰彎得非常嚴重,原文是彎在一起,雙重彎；

這不是普通的駝背,而是指她彎得胸部和腿部幾乎重疊在一起.你可以想像這女子的苦況:

她不能如一般人一樣,

不能抬起頭來做人；

不能舉目與人有眼神的接觸；

更不能如一般人般站立起來.

她可能是位年青貌美的姑娘,卻因嚴重的彎腰而儀容盡失,

愛美的女子的天性,她卻被迫失去；

她因這彎腰而不能結婚,失去家庭的支柱,常要寄人籬下,影響生計；

她更會因這病被人看為是她犯罪的結果,受到上主的責罰,因而被人歧視.

這位女子在安息日到會堂在什麼?

親近上主,更想脫離那鬼的困縛,被耶穌發現.

她所受的困苦,已十八年,每天都因這嚴重的彎腰,而大大受苦,耶穌看見她的苦況,明白她的苦況,主動要幫助她脫離那困局.

於是於耶穌用兩隻手按著那女子,她立刻直起腰來,就歸榮耀與上帝.

『按手』為何如此厲害呢?因為這是象徵傳遞祝福(參提前四14；提後一6),也象徵生命的聯合(參提前五22) ,主耶穌的生命與那女子生命聯合起來,就產生能力,使她能轉彎為直.

在我屬靈的生命中,常常發現我的能力再加上主耶穌的大能,便沒有難成的事了；這就正如保羅所說,「我靠著那加給我力量的,凡事都能作」(弗).

上主是位能破人生困局的主,並且是位主動破局者.從路13:10-17節,讓我們看見這樣的一位主.

在此,看見耶穌要使嚴重彎腰的女子轉彎為直,使她得著健康.

各位:

身體的健康非常重要,沒有健康的身體,你想做什麼都不行,我常求主給健康!

我四十多歲才到英國進修,得到的是健康,在那裡多運動,吃東西也簡單,經常要行.

還有,一位好弟兄的提醒:「小心,走人生下半場!」

二思想的轉彎為直

人的思想會有的不良,不正,偏見,需要轉彎為直.

管會堂的人存有自己的偏見,認為耶穌應遵守猶太律法,不應在安息日治病.這偏見也需要轉彎為直.

你不奇怪嗎?

管會堂的人並沒有因為那背彎十八年的女子蒙主醫治巳開心,反而怒氣沖沖,就是因為耶穌不應在安息日治病,這是違反猶太人安息日條例.

原來,按照猶太教拉比的教訓,除非人到要死的地步,否則不可在『安息日治病』,因認為治病是觸犯在安息日不得作工的規條.

他看重的是條例的實施多於人實際的需要.

在今日的社論,不難找到這種只懂拿著條例,面無人情,硬繃繃的人.

管會堂的,重視安息日的條例,並以為主上帝也看重律法的人.

耶穌提出兩個論據:

1耶穌在此指出人豈不比牲畜重要.

【路十三15】「主說:『假冒為善的人哪,難道你們各人在安息日不解開槽上的牛驢,牽去飲麼?」

﹝文意註解﹞

「假冒為善的人哪,」指他們假裝對律法熱心,其實對律法『愛人如己』的總綱(參太廿二39~40)卻毫不以為意.

「難道你們各人在安息日不解開槽上的牛驢,牽去飲麼?」這表示他們關心動物,勝過關心那更有需要的人.關心自己的家畜,而不關心別人的病痛,正表明他們愛己不愛人──不愛人如己.

人憐憫牛羊,怎不會憐憫人.

主耶穌強調強調在安息日醫治人,憐憫人比守安息曰的條例重要.

2這女人也是亞伯拉罕的後裔

【路十三16】「況且這女人本是亞伯拉罕的後裔,被撒旦捆綁了這十八年,不當在安息日解開她的綁麼?』」

﹝原文字義﹞「後裔」女兒.

﹝文意註解﹞「這女人本是亞伯拉罕的後裔,」意即她本是上帝選民中的一個.

﹝靈意註解﹞「不當在安息日解開她的綁麼?」指安息日是叫人得安息的(參創二3),不是叫人留在轄制底下,被捆綁,失去做人的能力.

管會堂的實在是曲解上主的意思,耶穌要轉他的思想為直,從律法主義者變為愛鄰舍者.

耶穌早已提出過:安息日是上主人而設,不是人為安息日而設.明顯地,人是比安息日重要.

其實,猶太人把數以千計的律法條文己已成為重擔,使人不得安息,成為重擔.

處理人事問題,到底按條例章則,或者憑愛心行事呢?

守緊章則,不講愛心,使人覺得無情!

只講愛心,沒有章則條例,沒有原則,會造成大細超,引致不公平!

當中如何平衡?公平公正處理?

耶穌要轉管會堂的為直.

上主基於愛,設立安息日,為使人在忙碌的工作中停下來休息.然而,滿有愛心的上主,怎能見到人在安息日受苦受捆縛而不拯救人,因此,基於愛,在守安息日的原則中仍有體諒.

各位:你曾否有些對人的偏見?需要轉彎為直?

你曾否有些對上主的偏見?需要轉彎為直?

讓上主光照你的思想,能助你轉彎為直:

每日定時靈修,或靜思時間,看看有些甚麼思想,需要更新?

一個傳統簡單的方式:

1讀經時,上主向我說話

2祈禱時,我向上主說話

三精神的轉彎為直

「被鬼附著病了十八年,」指她『被撒旦捆綁了十八年』(參16節).

「腰彎得一點直不起來,」象徵人被魔鬼轄制得只能俯視屬地的事物,而不能仰首望天.

不能望天,不能望人,是非常痛苦的事.

那麼,這位女子在安息日到會堂在什麼?親近上主,

更想脫離那鬼的困縛,被耶穌發現.

那女子被污鬼附身,使她受束縛,受痛苦

她成為污鬼的奴隸,要服從那惡者.

她成為污鬼的的囚犯,失去自由.

那惡者代表一股看不見的力量,操控著她的生命.

各位,這鬼是路加對當日的描述,用鬼作解釋.

在今日我相信鬼,不易上人身,也難以証實是否鬼附人身.

我認識兩位朋友,兩位都非常出名的精神科醫生,也是基督徒,甲相信有鬼的存在,乙不信認為這些都是精神問題導致.

我也相信有邪靈的存在,我牧會經常都驅魔:

在全完堂近四年,在四次,

在香港閩南堂,也有四次,

在望覺堂近二十,至少有十次,

到英國進修幾年,也有幾次.

每次那些見到鬼的,都講到鬼咁驚,鬼咁可怕,連我的毛管都動.

我每次都趕到鬼,暫時未失手過?

但是,是否每次都是鬼怪的騷擾,我也是弄不清的!

然而,對許多人說,一種難以解釋的精神困擾會操縱住人的心思意念,使人好像被困縛,好像失去自由般受盡苦楚.

這種精神的困縛可以是什麼呢?

懼怕

發脾氣

情緒失控

灰心喪志

抬不頭來做人

這種精神的困縛也可能是?

使我們生命的潛質受限制

使我們為人的上帝形象不能發揮起來

因為聖經描述的鬼是“Disabling Spirit”

「被鬼附著病了」:直譯是「有一個軟弱的靈(精神)」,意思是婦人病因是這個「軟弱的靈(精神)」.

耶穌對那女子說:“Woman you are freed from disability”                                    (ESV, V12)

因此,這精神的困擾,使人完全被制住,軟弱起來,失去做人的力量.

各位:

如果你發現你生命中,有一股使你軟弱的,令你墮落,使你生命會受損的靈,求主幫助你勝過它.

四 關係的轉彎為直

1那女子從彎腰中得醫治,感謝上帝,得與上帝和好.

2那女子從鬼的捆縛中得釋放,得自由,與自己和好,可以控制自己腰,歸榮耀與上帝.

3管會堂的,從偏見中感到羞愧,感到自己的不是.

4眾人感到歡欣,因那患病女子得醫治,耶穌又言之有理.

在福音書中提到人與神,人與人關係的建立,有一個重大的意義,就是「進入天國」,就是生命最幸福的境界.

耶穌要使身體彎曲的女子轉彎為直,健康起來,進入天國.

耶穌要使只懂得死守律法管會堂的人轉彎為直,脫離律法主義,成為有愛心的人,也進入天國.

耶穌要使被撒旦捆搏的人轉彎為直,得著自由,進入天國.

甚麼是健康呢?

【】一位九十多歲人記的弟兄患上癌症,問他的老友記,醫生朋友:「你看過我的醫療報告後,我好番的機會有多少?」

醫生朋友:「你好番想做甚麼?九十多歲人!」

這位長者突然醒悟:我要生命健康做什麼?許多壞蛋,惡人都健康去做盡壞事,我要健康,應有更高的意義!

甚麼是健康呢?

除了身,心,靈/精神及關係均健康,才是完整健康外,還有更深神學意義:

1神學家哈林 (Bernard Haring)提出一種看法:

對人類健康的全面理解,包括人全部力量和精力可達到的最和諧境界,賦予肉身實際情況最大的屬靈意義,加上將靈性最美好地體現出來.

真正的健康顯示於人的自我實現當中,就是人獲得自由管理自己所有天賦和能力,以盡做人的一切天職.

2巴特(Karl Barth) 根據這論點,認為健康是「人之所以為人的能力」.

3莫特曼(Molmann)說得更為清晰:「健康……是做人的力量」(加了強調效果).

Juergen Moltmann describe health:

“True health, is 真正的健康是

the strength to live, 有力量面對生活的一切

the strength to suffer, 有力量去面對苦難

and the strength to die. 有力量去面對死亡

Health is not the condition of my body,

健康並非人的真實境況,

it is the power of my soul to cope with the varying conditions of that body".

健康,就是指人具有能力去面對一切人生的境遇.

從神學範圍去看做人的意義:

a它暗示了不夠健康就是不夠力量完全成為人,也即暗示了只有健康的人才完全是人.

b這樣看來,疾病就破壞這種力量,令人變得虛弱,失去活動能力.

因此耶穌要醫治那女子,使她能夠恢復做人的力量.

假如健康是「做人的力量」,人對其健康或患病的狀態就可能有健康或病態的態度.

這顯示於人對快樂與苦難的承載力,及對生命中的喜樂和對死亡的哀傷的接受程度.

結語:

各位:

你是否健康?你在那方面有屈曲呢?

在身體,思想,精神及關係中,你那方面需要改善?

你身體有毛病嗎?

你思想積極嗎?

你精神活潑?

你與人關係良好?

真正的健康是能發揮做人的力量.

你能發揮做人的力量?

你能把人的潛質發揮出來?

願上主助你發揮做人的力量,成為健康的人.

耶穌也要使你彎彎曲曲的人生轉彎為直!請想想你人生的彎位在那裡?這個彎位是你生命的盲點,弱點,死穴需要奇妙的主,轉彎為直?



\newpage

\section{不怕兇惡的信息}
\label{sec:216}
\textbf{伍渭文牧師 - 不怕兇惡的信息}
\newline
\newline
\hyperref[sec:215]{< < < PREV SERMON < < <}
~
\hyperref[sec:toc]{[返目錄]}
~
\hyperref[sec:217]{> > > NEXT SERMON > > >}


講題:不怕兇惡的信息　Fear Not the Evil Tidings

經文:詩篇112

講員:伍渭文牧師

場合:香港中文大學崇基學院禮拜堂

日期:2013年9月1日

我一位親友在美國洛杉磯首次買了一房子,先生是印度人,入伙時請了當地印度教法師為他們的新居潔淨和祝福.他們曉得我是牧師,剛好在洛杉磯,也請我為新居祈禱祝福,並問我甚麼時候來,先來還是後來,還是同時來.我說我同時來,因為我很想認識印度教儀節,難得的機會.印度教儀式非常繁複,地上鋪上黃,紅顏色帷布作祭壇,放上偶像,香料水果,奉上代表金錢珠寶的米給偶像.主禮穿上祭袍,儀式歷時一個半小時；我盤膝而坐,弄到腿酸腳痺.之後,我叫他們站起來,圍成一圓圈為他們祈禱:感謝過去的保守,現在的祝福,將來的帶領,三分鐘就完了.

尋求祝福是所有人的願望,新居入伙祈福多多益善,管它甚麼宗教.

我們今天詩篇的經課詩112篇也論到有福的人,究竟有福的人是一個怎樣的人呢?

詩112開宗明義,蒙福的人不是你獻上甚麼豐厚的祭品,做了甚麼儀式,乃是(1)敬畏耶和華,甚喜愛他命令；(2) 施捨錢財,周濟貧窮；(3) 不怕兇惡,與邪惡爭戰.

詩篇是聖經聖殿集體敬拜用的公共禱文,禱文體裁多樣化:有個人情感抒發的祈求,讚美,感恩,哀歌；但有一種叫智慧文體,像詩篇112篇；用冷峻的眼光,洞察世情,透視人生,使人得到做人的智慧.作為公崇拜的禱文,智慧詩篇誘發屬靈生命的響往,模塑正確的人生觀,像北斗高懸,為天路客定向導航,不至迷失.崇拜不單有抒發作用,亦有誘發作用.

1. 敬畏耶和華,甚喜愛他命令 (Guardian)

喜愛上主的命令或律法是智慧傳統的主題,我們可以比較作為卷首篇詩第一篇和一一二篇,兩篇原文都以福哉斯人(Blessed is the man) 開始；詩篇第一篇有福的人是靜態的,田園的；像一棵樹,栽在溪水旁,按時候結果子,他晝夜默想上主的律法:「不從惡人的計謀,不站罪人的道路,不坐褻慢人的座位,惟喜愛耶和華的律法,晝夜思想,這人便為有福.(詩一1, 2)

詩112是動態的,很多經營活動.有商業活動:「家中有貨物,有錢財」；有家族親情:「他的後裔在世必強盛,正直人的後代必要蒙福」；有社會關懷:「他有恩惠,有憐憫,有公義.施恩與人,借貸與人的,這人事情順利.」他不怕惡勢力:「不怕兇惡的信息……他心堅定,總不懼怕……直到他看見敵人遭報.」這人有謀略智慧,疏財仗義,抱打不平.

信仰的承傳.詩篇提到「他的後代必要蒙福.他家中有貨物,有錢財；他的公義存到永遠!」(詩一2,3)表面看來,後代蒙福,家有錢財,是喜愛律法的回報；我們遵行律法,上帝就使我們事業成功,福蔭後人.但其實後代蒙福,因為他們幼承庭訓,喜愛律法.詩中提到「他的後代」,說明主角是一位父親和丈夫,他注重兒女的信仰和價值觀.後代蒙福,因為受了家庭宗教教育的薰陶,像父母親一樣遵守律法,因而蒙福,詩中主角是一位履行責任的父親和母親,不單是子女的監護人,也是信仰的監護人,他把信仰價值傳遞給兒女.非洲諺語:「青蔥歲月不扶直,乾涸日子難匡正.」(If you refuse to be made straight when you are green, you will not be made straight when are dry) .

大家都知道我年底將退休了,時光飛逝,以前知道,現在感覺到.師母常常問我:離世時我們遺下甚麼重要的東西給我們的兒子呢?人生苦短,很多事不是自己能夠完成的,事實上有很多事不是自己最適合去完成的,但你可以把價值信仰傳給其它人,特別是自己的下一代,影響他們的生命,我想這是我們給後代是寶貴的產業,我想這也是基督教高等教育的目的.「後代必要蒙福」並不是父蔭,仍因為後代從你的身上學曉蒙福的方程式:敬畏耶和華,甚喜愛他命令.

2. 施捨錢財,周濟貧窮 (Protector)

信仰的增值.蒙福者不單是家中的監護人,也是家中的供應者.他不單把信仰和價值教導兒女,他的信仰影響他如何待人接物,行事為人.因為實踐信仰,他耿直公正,與人和睦,樂於助人,得人信任,令人喜愛,所以事業亨通；「家中有貨物,有錢財」是躬身力行信仰得來的自然結果.信仰不單用來修身,亦可用來齊家治國,信仰為我們增值.

我1989至1991在台北教書時兼任一教會牧師,有一段時間在崇拜講解八福,有一位未信主的藥廠營銷經理經常參加崇拜,他的老板是教會長老,邀請他到教會聽道.他說:聖經的教訓對我的工作很有幫助,我經常用聖經的道理訓練同事,如何對待客戶,增加業績.

厚賜百物的上帝.舊約聖經常常以豐厚的財富和連綿的後代表示上主的祝福,新約聖經比較多強調在基督裡的屬靈福氣.因為相對於基督帶給我們屬靈的福氣,地上的財富就顯得不重要了,所以新約聖經比較少提到物質的豐富作為上帝的祝福,但聖經並沒有眨抑財富.保羅提醒提摩太,「不要依靠無定的錢財；祇要依靠那厚賜百物給我們享受的上帝.」(提前六17)上主仍然厚賜百物給我們享受,我們可以處卑賤,也可以處富有.

上帝也使用願意給上帝所用的財主.耶穌安葬的墓園是財主約瑟讓出來的,保羅在歐洲傳道第一站腓立彼得到賣紫布的呂底亞(Lydia)的經濟支持,打開歐洲宣教大門.馬丁路德(Martin Luther)1520年第一篇改教策論「致德意志諸侯書」(To the German Nobility),就呼籲作為有權力財富的諸侯,投身改教運動.

作先知亦作君王.有人認為教會不要靠攏權貴,擠進建制中心,應該在建制邊緣作為先知,發聲批評建制中心,因為權力容易使人腐化.但永遠的批評者,沒有執政的承擔的批評者,可能會嚴苛自大,鮮於諒解寬容,不懂讓步妥協.基督徒要效法基督的三重身份:先知,祭司,君王.先知表彰上主的公義,祭司顯明上主的憐恤,君王運用上主的智慧,在墮落的世界,落實天國的理想.若上主祝福你的事業,成為機構的行政主管,我們就有責任在自己影響的範疇,作好君王治理的責任.我們相信恩典可以治愈自然(grace heals nature) ,上主使用我們成為施恩的器皿,賜福他人.

權力在兩種:overpower and empower,宰制他人,和提升他人.約瑟在埃及王朝,但以理在巴比倫,尼希米在波斯的書珊(Susa)城,他們在異教帝國身居要位,同時活出信仰；信仰增值,幫助他們的事業,因為他們「敬畏耶和華,甚喜愛他命令」.

作為校牧,依照傳統我在周會講演前安排唱詩,讀經,祈禱；在學院的重要會議開會前和重要場合用膳前,常常要以校牧的身份祈禱謝恩,所以我要非常小心,怕忘記出席.我常常問自己:這些祈禱對參加者來說,是期待的資產(asset),還是要容忍的負債(liability) 呢?我們的信仰是否成為我們學業,事業的助力呢?我希望祈禱成為大學教育的增值,成為心靈教育的機會,我重視每一次領祈禱的機會,並寫周會禱文,讓同學藉禱文思考人生.

這人家中有貨物,有錢財,不是因為思想上主命令,背誦律法的回報,乃把律法貫徹在其商業活動,建立商譽,給人信賴,以至事業亨通.

保護者(Protector) .這位成功的商人,擁有令人欣羨的財富,但他沒有忘記貧窮的人,他施捨錢財,周濟貧窮,保護弱小,他是一位保護者(Protector) ,正如蜘蛛俠電影中有一句很有意思的話:「能力是用來保護弱小的」.我們留意當經文提到他「家中有貨物,有錢財」,同時提到他的公義存到永遠(v. 3),「義人被記念,直到永遠.」(v.6) 祇有公義存到永遠,家財貨物是暫時的,我們離世時不能帶走.然而,奇妙的事發生了,當我們把財富和有需要的窮人分享時,這短暫的財富便成為永遠的財富了,「他施捨錢財,周濟貧窮.他的仁義存到永遠.」(v.9)

3. 不怕兇惡信息的戰士(Warrior)

詩中主人行為正直,但卻被黑暗不容,他還被誣告帶到公堂被審.「正直人在黑暗中」(v.4),「他被審判的時候,要訴明自己的冤」.(v. 5) 他周濟窮人,是八福中憐恤人的人,卻要為義受逼迫,個受逼迫也是八福之一；但他沒有在退縮,他與不公義周旋到底.「他心確定,總不懼怕,直到他看見敵人遭報……他的角必被高舉,大有榮耀.」(v.8, 9b.)這裡提到敵人,提到(號)角,提到大有榮耀(凱旋),儼然是一個戰士得勝的意象.

挑戰黑暗.為何他惹來黑暗力量的惱恨呢?其中可能原因是,當我們關心弱者,窮困者的遭遇時,我們會進一步追問為何他們會如斯不幸?是否社會資源分配不公平?是不公的制度造成結構性貧窮呢?我們就由廣施憐愛進入尋求公義.很多醫療志願團體,開始時的扶傷治病,建造醫院；稍後就進入公共衛生領域,思考到資源分配是否公不公?是否環境因錢權尋租被污染,引至疾病?這樣就和利益集團起衝突.

揭示虛假.魯益斯(C.S. Lewis) 說:世界就是淪陷在敵人手中的領土 (Enemy-occupied territory-that is what this world is) .有時敵人裝扮成光明的天使,欺哄我們,迷惑我們,魔鬼就是說謊之父-世界給我們「虛假」(counterfeit).這段經文描述善人時,用了永遠一詞:「他的公義存到永遠」(v.3, 9) ；「他永不動搖.又被人紀念,直到永遠」(v.6),對比惡人的虛幻:他們「必咬牙而消化」,「心願要歸滅絕」.一個永遠,一個短暫,消化滅絕是它們的本質,但它們聲稱是永遠的,真實的.正如詩篇第一篇談到惡人的終局,如糠秕被風吹散；糠秕被風吹起時,給你金光閃耀的幻象,但經風一吹就消化滅絕,歸於無有.

魔鬼試探耶穌說:你若攀到聖殿至高的金頂,從那裡跳下來,上主的使者托住你,你不會受傷的.(太四5) 世界也叫我們攀上頂尖的地位,踏著同事的頭上位,不理會家庭,不理會身體,祇不斷攀爬,以為金頂使我們呼風喚雨,不會受傷.但我們到了金頂才發覺那裡是孤單和空虛,我們已失去朋友,失去家人,也失去健康,我們傷痕纍纍,那金頂是假冒counterfeit的能力,不是真實,也不永遠.

我們也常常假冒自己是上主:「今天,明天我們要往某城裡去,在那裡住一年,作買賣得利.其實明天如何,你們還不知道.你們的生命是甚麼呢?你們原來是一片雲霧,出現少時就不見了.你們當說,我們就可以活著,也可以作這事,或作那事.」(雅四13, 14)

世界也有它的虛假的軍隊,就是僱傭兵(mercenary) ,他們參與戰爭是為了金錢,不是為國家,那裡出得起錢僱用他們,他們就就到那裡打仗.其實人生就是戰場,我們都為一些東西來打仗,僱傭兵為自己爭戰,真正的戰士為國度而爭戰.今天我們為了甚應而戰爭呢?

今天的詩篇給我們看到一個真正乃是蒙福的人, 他是一個不怕凶惡信息的戰士.他首先為自己的兒女價值觀而爭戰,與世界爭奪其子女,他是真理的監護者(guardian)；作為戰士保護窮困弱者(protector) ,他施捨錢財,周濟貧窮；他是不怕威嚇的戰士(warrior),為公義的國度而爭戰.

戰事正酣,總要警醒.盟軍登陸法國諾曼第(Normandy) 是二戰的重要戰役.當傘兵從漆黑的高空跳下來時,生死未卜,他們要落未曾到過的地方,要遇到未曾見過的敵人,落地時可能飄離戰友,孤軍作戰.戰爭電影給我們的印象是這些空降師傘兵個個都士氣如虹,勇敢激昂,視死如歸.但根據一戰地實況報導,有些人落到地上就怕得捲伏在矮樹叢中,有些更睡著了；更匪夷所想的是,一班傘兵降落到一倉庫,發現裡面藏有烈酒,竟然喝得酊酩大醉,東斜西倒.(D-Day, quoted in John Eldredge, Wild At Heart. Discovering the Secret of A Man’s  Soul [Nashville, TN: Thomas Nelson, 2010], p. 160.)

這些人知道自己是軍人,也知道在處在戰爭中,但他們的行為表現,顯示他們好像不在戰爭中.我們呢?

不要醉酒,乃被聖靈充滿,讓我們放下各樣的重擔,脫去容易纏累我們的罪,拿起屬靈的兵器,不再作虛假的士兵,不再作僱傭兵為自己爭戰,讓我們為上主的國度爭戰,因為這國度是真實的,也是永恆的.阿們.



\newpage

\section{不斷思量的門生}
\label{sec:217}
\textbf{許開明牧師 - 不斷思量的門生}
\newline
\newline
\hyperref[sec:216]{< < < PREV SERMON < < <}
~
\hyperref[sec:toc]{[返目錄]}
~
\hyperref[sec:218]{> > > NEXT SERMON > > >}


路加福音 14:25-35
\newline
\begin{longtable}{cl}
\hline
\hline
章節 & 經文 (和合本修訂版)\\
\hline
14:25 & \begin{tabularx}{0.7\textwidth}{X} 有一大群人和耶穌同行。他轉過來對他們說： \end{tabularx} \\ \\ \relax
14:26 & \begin{tabularx}{0.7\textwidth}{X} 「無論甚麼人到我這裡來，若不愛我勝過愛自己的父母、妻子、兒女、兄弟、姊妹，甚至自己的性命，就不能作我的門徒。 \end{tabularx} \\ \\ \relax
14:27 & \begin{tabularx}{0.7\textwidth}{X} 凡不背著自己的十字架來跟從我的，也不能作我的門徒。 \end{tabularx} \\ \\ \relax
14:28 & \begin{tabularx}{0.7\textwidth}{X} 你們哪一個要蓋一座樓，不先坐下來計算費用，看能不能蓋成？ \end{tabularx} \\ \\ \relax
14:29 & \begin{tabularx}{0.7\textwidth}{X} 免得安了地基，不能蓋成，看見的人都笑話他，說： \end{tabularx} \\ \\ \relax
14:30 & \begin{tabularx}{0.7\textwidth}{X} 『這個人開了工，卻不能完工。』 \end{tabularx} \\ \\ \relax
14:31 & \begin{tabularx}{0.7\textwidth}{X} 或是一個王出去和別的王打仗，豈不先坐下來酌量，他能不能用一萬兵去抵抗那領二萬兵來攻打他的嗎？ \end{tabularx} \\ \\ \relax
14:32 & \begin{tabularx}{0.7\textwidth}{X} 若是不能，他就趁敵人還遠的時候，派使者去談和平的條件。 \end{tabularx} \\ \\ \relax
14:33 & \begin{tabularx}{0.7\textwidth}{X} 這樣，你們無論甚麼人，若不撇下一切所有的，就不能作我的門徒。」 \end{tabularx} \\ \\ \relax
14:34 & \begin{tabularx}{0.7\textwidth}{X} 「鹽本是好的；鹽若失了味，怎能叫它再鹹呢？ \end{tabularx} \\ \\ \relax
14:35 & \begin{tabularx}{0.7\textwidth}{X} 或用在田裡，或堆在糞裡，都不合適，只好丟在外面。有耳可聽的，就應當聽！」 \end{tabularx} \\ \\
[1ex]
\hline
\hline
\end{longtable}
講題:作個不斷思量的門生　Counting the Cost of Discipleship Always

經文:路加福音14章25-35節

講員:許開明牧師

場合:香港中文大學崇基學院禮拜堂

日期:2013年9月9日

引言

今日是2013年9月8日,是香港中文大學2013-2014年開學後,崇基禮拜堂於新年度,首次舉行的主日崇拜.今日特別歡迎第一次赴會的新朋友,請舉手,讓我認識你!

今日,對我來說,也是特別的日子,是我剛就任崇基學院校牧的第一篇講道.

按照教會年曆,講到福音書,路14:25-35,題目是「作個不斷思量的門生」,想不到這題目也非常配合開學的崇拜.

我自從三月確定擔任校牧後,四月便開始有機會講道,但是一直都未有機會聽伍牧師及其他助理校牧講道,從中得著學習,因過往在堂會事奉非常忙碌,特別是主日,因此無法來.從主日崇拜調查表的回應中,得知兄姊對講壇的評價高,期盼我們能得著鼓勵,精益求精.

各位:

有人問你說:「你有無智慧嗎?」或「你有無腦?」

有位兄弟,自覺自己好愚蠢,於是問他的朋友:「你覺得我蠢嗎?」他的朋友回應:「你竟然問我這問題,真是好蠢!」他為要解決愚蠢的問題,決定找醫生說:「我很蠢,怎麼辦?」醫生說:「入蠢無藥醫!不過我有些獨一秘方,你拿去試試吧!」

第一次,吃完,對醫生說:「我都係咁蠢!」

第二次,吃完,對醫生說:「我比前更蠢!」

第三次,吃完,對醫生說:「你騙我,那些根本不是聰明丸, 而是維他命丸」醫生對他說:「對了,你終於聰明了.」這位兄弟聰明了,因為他懂得思想.

信息

到底「有思想」是什麼?思想就是有觀點.當你對一件事物有觀點時,你就是有思想的.根據我的觀察,在未認識某人或某事之前,我們都是沒有觀點的.觀點越多且越接近本質,那麼你的思想也越是豐富和深刻.舉例來說,我見過又吃過蘋果,於是我對蘋果有下列的觀念:是心型的,是爽的,是有益的.

猶太入有句名言:“人類一思考,上帝就發笑” 這表示:上帝喜愛人思考,上帝因為人思考便高興起來,其實上帝也是個愛思考的上帝.

今早的經文提到「有極多的人和耶穌同行；祂轉過來對他們說:」【路十四25】 「有極多的人和耶穌同行,」同行的人雖然『極多』,但是主難得有真正的跟從者；這就是祂說出下面一段話:

一大群人為何跟著耶穌?因為有得聽,有得食,有得睇,有得傾,有病得醫治,有鬼得趕走…,跟住耶穌,真不錯.他們中,許多入跟耶穌可能只看到這些好處,忽略了跟從耶穌的人要成為耶穌的門徒:

要學耶穌,愛與犧牲(為人犧牲)

要跟耶穌,背十架走路(痛苦的窄路)

因此耶穌要他們思量:思想和計算代價.

第一點:思量對主的愛有多少

【路十四26】「『人到我這裡來,若不愛我勝過愛自己的父母,妻子,兒女,弟兄,姐妹,和自己的性命,就不能作我的門徒.」耶穌指出有意跟從耶穌的人,

愛我勝過愛其他人:愛我而使其他的失寵了:愛我勝過愛自己的父母,妻子,兒女,弟兄,姐妹,和自己的性命,這是說,愛了耶穌後,其他以前的至愛【父母,妻子,兒女,弟兄,姐妹,和自己的性命】便漸漸失寵了.的確,今日不少入,愛「貓,狗,雀」多過人的!日前,一位太太把她無限的大愛給三十隻貓,卻冷落丈夫,服事貓,多過丈夫,引起衝突,太太還打傷丈夫…

還有這裡的「愛」,原文是恨惡,可憎.恨:即不再愛了.所以「人到我這裡來,若不恨自己的父親,母親,妻子,兒女,弟兄,姊妹,甚至恨自己的生命,就不能作我的門徒.」

耶穌的意思是說,一個人有了最愛的,就會輕看其他的.如果你已是基督徒,基督是你的最愛,你便輕看其他入了.

愛主者,思量背主十架

愛主者,還有進一步的要求:背十架

「凡不背著自己十字架跟從我的,也不能作我的門徒.」【路十四27】「凡不背著自己十字架跟從我的,

(凡不背十架好像我一樣),」十字架在信徒身上,具有雙重的意義和目的:十字架乃是世人所加諸在主耶穌身上各種不合理的對待,和痛苦的頂點.所以背十字架跟從主,對信徒來說,是代表各種為主所忍受的1苦難,2羞辱與3犧牲.

第二點:思量自己能力有多少?

蓋一座樓

「你們那一個要蓋一座樓,不先坐下算計花費,能蓋成不能呢?」【路十四28】這一座樓,是在葡萄園中的高塔,用來監視賊入之用,因賊入會偷葡萄.人要計算財力,是否足夠建樓?

打一場仗

「或是一個王,出去和別的王打仗,豈不先坐下酌量,能用一萬兵,去敵那領二萬兵來攻打他的麼?」如果不行,就不要打了!「若是不能,就趁敵人還遠的時候,派使者去求和息的條款.」【路十四31-32】打仗要計兵力,一萬兵怎敵兩萬兵?

處事:要思量自己的實力或能力有多少?

第三點:思量能否撇下一切跟從主?

【路十四33】「這樣,你們無論甚麼人,若不撇下一切所有的,就不能作我的門徒.」

撇下一切所有的:是專一的跟從,這也是要付代價的.

彼得,安得烈,約翰,雅各,放下魚網,一切所有的都跟從主.

作主的門徒,要不斷思量:

1愛主有多少?

2實力有多少?

3能否撇下一切,跟從主?

小結

耶穌提出的這四個問題,不是說,他們想清楚後,才去跟從主,而是一路跟,一路思想這三個問題.

作中大的學生,你也要思想是要付代價的:

這不是指交學費的代價,付上讀書的代價.第一,勤力,不怕吃苦,因為精英雲集,入得中大,都是精英,或疑似精英的.你要努力,才可坐到船(Scholarship).如果你是只靠補習天王或天后教幾招,純粹靠掌握考試技巧變成摘星星人的,更要努力.(當然,補習社也有非常有實力的摘星人).第二,大學生命有寫意的一面,也有忙碌辛苦的一面.趕交報告,考試,論文,許多的聚會活動,你要估量時間,精神體力?第三,更重要的,作為大學生或老師,應該是個不斷思想的人.作為今日的老師同學也要思想:

作個不斷思想的人!

今日的社會,會使人失去思想的能力:

●就業的人,大量的工作,使人成為工作的機器,缺少思想的空間.

●學生,不停的大考,中考,小考；各種的測驗,使學生被迫成為書蟲,書獃子.缺少思想的空間.

●今日大量的資訊,要研究探討一個問題,不是沒有資料,而是資料太多,太濫,太眼花瞭亂了,資料塞住腦子,腦袋也塞住,腦筋不會動,只懂抄襲人家的思想看法,不懂分析,比較,批評,評論.我教中國教會歷史時,常提醒同學,不要給史料塞住個腦,多些用腦,我要你的腦汁,看到你的觀點.

因此,我們需要讓自己有多些思想的空間.特別是一班在學的老師同學.有思想的人,是個常常動腦筋的人,這種人是需要教育,訓練,並養成一種勤於思想的習慣:

記得18世紀法國思想家盧梭曾經說過:“無論就男性或女性來說,我認為實際上只能劃分為兩類人:有思想的人和沒有思想的人,其所以有這種區別,差不多完全要歸因於教育.”好的教育,讓人學會獨立思考,從而變得有思想；這獨立的思考,包括:1有自己的觀點,2主見,3懷疑,4批評.壞的教育,給人洗腦,從而變得沒有思想,只會一個傳聲筒,或者鸚鵡學舌,人云亦云.古希臘著名的思想家蘇格拉底說:“人之所以為人,就是因為他有思想,能獨立思考而不盲目相信流行的權威和說法.” 西塞羅更提醒世人說:「智者受理智的指導,常人受經驗的指導；而野獸受直覺的指導.」

一個有思想的人,站在佛家的觀點,就是一個已經覺醒的人.「佛」的意思,就是「已經覺悟的人」.覺悟,就是要認識到自己作為一個人的生命,掌握了三個問題:我來自哪裡?我現在如何?我將要去哪裡?各位讀者,你能夠常常自省這三個人生的基要問題,你會得著智慧面對人生.路加福音十五章所記載的浪子,他的生命得著轉機,就是因為他在豬場,走投無路,就醒悟過來,聰明的浪子也曾思想過這三條問題,一 我來自哪裡?來自有許多僱工,口糧有餘的父家；二 我現在如何?我現在恨不得拿豬所吃的豆莢充飢,也無人吃我,我怎能活下去!三 我將要去哪裡?我不能坐以待斃,我要回到父親那裡,向他認錯,並甘願作他的僱工,求得父親的收留(見路15:14-20節).結果回頭是岸的浪子,得到慈父的接納,人生悲劇改寫為喜劇.

著名哲學家康得在其名著《純粹理性批判》中提到過人類理性的三個永久性的人生問題,他認為以下簡單的問題,是我們要運用理性的腦袋思考的.這三個問題就是:一,我能夠知道什麼?二,我應當做什麼?三,我可以希望什麼?人一生中,其實也一直在這三個問題中掙扎.一個人活得精靈還是糊塗,就要對這三個問題進行審察,然後作出自己的回答.康得說,第一個問題是單純思辨的.第二個問題是單純實踐的.第三個問題,既是實踐的,又是理論的.以上三條問題,筆者經常問自己,從中得著智慧.舉例來說,我經過多年艱苦進修取得博士學位後,首先,我知道上主的心意,不會讓我人生下半場在一般堂會牧會；第二,我想我應該等候上主的安排在神學院或大學裡事奉；第三,我可以期望成為神學院講師或崇基的校牧,結果上主給我的是擔任校牧,牧養數以千計的大學生及教職員工,頗具挑戰性,卻是主所賜,我所愛的事奉,感謝上主.期盼你也能從以上三個問題,常作思考,得著人生的智慧.

環觀許多成功的人或對世界有貢獻的人都有相同的特徵,就是善於學習,勤於思考,用於實踐.學者蒙田有段教人深思的說話:「人要有三個頭腦,第一個,是天生的,上主給你的一個頭腦；第二個,是從書中得來的一個頭腦,你學到的東西；第三個,是從生活中得來的一個頭腦,你在生活中領悟到的智慧.倘若你已擁有這三個腦袋,真恭喜你了!如果有人問我:「許牧師,你有無腦?」我會這樣答他說:「我當然有腦,而且我有三個腦呢!」各位讀者,你又有多少個腦呢?

結語

無論在生活上,信仰上或求學上,願各位都能夠成為不斷思量的人.



\newpage

\section{人生中之圓與缺}
\label{sec:218}
\textbf{梁元生教授 - 人生中之圓與缺}
\newline
\newline
\hyperref[sec:217]{< < < PREV SERMON < < <}
~
\hyperref[sec:toc]{[返目錄]}
~
\hyperref[sec:219]{> > > NEXT SERMON > > >}


講題:人生的圓與缺　A Fulfilled or Broken Life

經文:路加福音15

講員:梁元生教授

場合:香港中文大學崇基學院禮拜堂

日期:2013年9月15日

一．前言

如果在座各位聽過我在主日崇拜中講道,大概會察覺到我所講的一切, 無論是舊約或是新約中的人和事,背後都環繞著一個主題: 那就是以聖經作為現代基督徒生活的參照.我們閱讀聖經,應盡量注意聖經話語與現實的關聯性 (contemporary relevance). 今天 我仍然打算用這樣的角度和大家看一段經文.

按照教會傳統年歷的安排,今天要講的經課是路加福音十五章,牧師告訴我,由於通知我的時間短促,我可以不依照年歷安排的經課去講,可以用另外一段經文,自定題目.我把那幾段經文看了多遍,決定還是依照安排以路加福音 十五章作為今日講道的主題,以之思考基督徒的人生.現在請大家翻開聖經新約路加福音第十五章.剛才讀經員替我們讀出路加十五章的上半段,但今天我們要看整章的聖經.包括下半段.

二．三個比喻

這章經文記載著耶穌基督向他的跟隨者所講的三個比喻.

第一個比喻:3-7節

「耶穌就用比喻說:你們中間,誰有一百隻羊,失去一隻,不把這九十九隻撇在曠野,去找那失去的羊,直到找著呢?找著了,就歡歡喜喜的扛在肩上,回到家裡；就請朋友鄰捨來,對他們說:我失去的羊已經找著了,你們和我一同歡喜吧!我告訴你們, 一個罪人悔改, 較 比為九十九個不用悔改的義人,歡喜更大.「

這是個失羊的故事.原來有一百隻羊,但一隻迷途,或者走失了,剩下九十九.

第二個比喻:8-10節

這是失錢的比喻:

「一個婦人有十塊錢,若失落一塊,豈不點上燈,打掃屋子,細細的找,直到找著嗎?找著了,就請朋友鄰捨來,對他們說 :我失落的那塊錢已經找著了,你們和我一同歡喜吧!我告訴你們, 一個罪人悔改,在神的使者面前,也是這樣為他歡喜.」

第三個比喻最為人所熟悉,11-32節.

這是浪子的比喻.由於時間所限,我不讀出整段的經文了,只請大家注意這家人有兩兄弟,小兒子離家出走,敗盡家財,窮途末路,挨飢忍餓,終於只有硬著頭皮回家,受到父親慷慨的款待,卻也引來長兄的嫉妒.

這三個比喻:迷失的羊,丟失的錢,和浪子回家,都是大家熟知的.我們往往會從傳揚福音的角度去瞭解這些比喻,強調耶穌對於失喪者的愛心:為那丟失的一隻迷羊和一塊大錢,如何放下群羊,加倍努力去尋回那失去的一隻；又為了那失去的一塊錢,而煞有介事的點燈照明,全屋打掃,亦表示其對失去的一塊重視之情.浪子的比喻更借父親之口說:這「失而又得」可以說是「死而復活」.

以上大概就是三個比喻的簡單介紹.

今天讓我們重新翻開這一章聖經,重新再閱讀和再思考這三個比喻,我希望從一個非一般的角度去看這三個比喻,看看可否得到新的亮光,和學到一些新功課.我建議的這個新的思考方法乃是從經文中的數字(numbers)入手.這個數字方法學(numerical approach)曾經被一些學者使用過,用來分析歷史文化.例如:研究中國歷史的一位專家黃仁宇教授,因不滿意傳統的歷史學家只集中其注意力於治亂,帝王將相,因而提出了數字管理和歷史發展的觀念,以財政史(數字管理)及社會史(大歷史或年鑒學派)的角度寫出《1587:The Year of No Significance》(即萬曆十五年)的史學名著.話說回來,在這段經文中和這些比喻里,明顯地,數字是很重要的構成部分.每個比喻里都含蘊著兩組數字:迷羊的比喻中,是一比九十九--一隻走失的羊, 而留在群中的則有九十九隻； 失錢的比喻中,是一比九--一塊錢丟失,還剩下九塊；而浪子比喻中,則是一比一,兄弟二人,一個去,一個留.如果以家庭來說,那就是一比三, 三個家庭成員,走失了一個.

其實,除了上面提及的兩組數字外 ,三個比喻中還有一個大數目(a big number),第一個比喻中是一百,第二個比喻中是十,第三個比喻中是三,而這個三也可以是一,一個家的一 .

先容許我從浪子比喻再詳細的解說:這組數字是一(兄)比一(弟),在加上父親便成為三,而這三人構成了另一個一, 這便是「 家」.浪子的故事,是一個「家」 中發生的一(父或兄)對一(浪子)的糾纏和衝突到復和的故事.至於迷羊比喻中,羊的總數是一百,走失的是一隻,還剩九十九,走失的是百分之一；失錢的比喻中錢的總數是十塊錢,丟掉了一塊,是十分之一.

這樣看來,一個有趣的現象就發生了:耶穌講這三個比喻時,似乎沒有把比例算准,或把差異當成一回事.九十九vs.一,或百份之一；和 十份之一,和二份之一,其重要性是相當的,是一樣的,竟然毫無分別!在我們來說,這當中有很大的差異,二份之一,十份之一,比之百份之一,其間差異甚大.我剛剛接到稅單,繳稅的數目是百份一或十份一,真是天淵之別.然而,這三個比喻中,有一點好像是相同的,就是「失去」的「一」--那一隻羊,那一塊錢,那一個人,那失去的」一」 才是最重要的,整章聖經的三個比喻都教導我們要用盡方法和費盡心思去尋回這個失去的一(只/塊/個).

三．人生有「缺」

從走失羊,丟掉錢,和浪子離家這一點,我要說到今天的主題了:缺.人生有缺!

缺即是失.也就是丟掉和不見了.若浪子沈淪,在外不歸,這個家就永遠有所缺,有所遺憾.這個家就不再是完整的家,圓滿的家.而在婦人眼中,那完整的大數是十塊錢,十是全部,是圓滿.(英文我們常說Perfect ten, ten is perfect.)一旦丟失了一塊,雖然只是十份之一,但代表了生命的不完整,代表了imperfect,代表了缺憾,代表了破碎.這缺憾和破碎的根源,可能來自自己的疏忽,自己的不小心；這「缺」或這「失」,是自己的過錯,或自己的犯罪和沈淪造成,也有可能來自上主的安排,或命運的戲弄.無論如何,造成了人生中的「缺」,有時是百份一的缺,有時是十份一的缺,有時缺憾極大,是二份之一.然而,缺無論是大或小,有缺則人生有憾,不可能成為完整,成為圓滿.因此之故,在我們的生命之中,最值得正視的問題,就是知到自己在哪裡丟失,在哪裡墮落,在哪裡有缺.對牧羊人而言,牽羊一百,在草原上放歌,山青水秀,雲蒸霞蔚,無拘無束,對生命的感覺是那麼的渾然一體,無掛無慮.但這種完整和完美的感覺,因一隻羊走失或丟掉而起了巨大變化,九十九不是一個完整的數目,從此我們得面對一個破碎的,有遺憾的人生.這就是牧人為甚麼要拼命尋找迷羊的原因,甚至要撇下那九十九去尋回那失去了的一.那是因為要「補缺」,彌補破碎的生命,尋找那整全的,圓滿的,完美的人生；不然,那九十九同樣也過著破碎和有遺憾的人生.換言之,牧羊人為了尋找那迷失的一,也在幫助那九十九回到原來完整和充實的生命.

那婦人點燈打掃,把原來平靜悠閒的生活改變成為動作多多,一片緊張的生活,原因也是一樣:尋回那失落的或不見了的生命中的那一塊.那是影響全局的一塊.失去了它,生命就不完整,就有缺憾.所以要盡力「補缺」,希望得到「救贖」,使生命回復完整和圓滿的感覺.

四．追求人生的圓滿

談過人生中的「缺「後,下面讓我在多講一下人生之圓,即人生中對圓滿和完善的追求.我們一般看這三個比喻,會把每一隻羊,每一塊錢,和每一個人,看成獨立的生命或個體.我在這裡建議大家換另一個角度去看:把一百隻羊看成一個整體,把十塊錢看成一個整體,也把父親及兒子三人組成的一個家看成一個整體.如此一來,我們就能較為容易明白「失去」一部分的痛苦和重要性.一百失去一,就不能圓滿,十失去了一,三人的家失去一人,同樣是殘缺不全的.如果我看這三個比喻,我會把完整及圓滿看成生命終極的關懷和目標,把一百隻羊,十塊錢,和三個人,看成不同的生命形態.換言之,有人的生命比較單純,或者比較短站,只有三或十；但也有人的生命較為複雜,經歷也多,牽連的個體和集體也多,甚至和別人的生命連帶一起,不易分開,因為那生命形態的構成部分,有一百之多.總之,我想說的:是生命形態的多元性.我們每個人的生命都不一樣,但這三個比喻給我們的啟示是:上帝的心意(也就是我們的終極關懷)是回到完整,達到圓滿的境地 .這三個比喻中,也把人生之無奈和困境說得很清楚:那就是比喻的核心信息:羊會迷途,錢會丟失,人會叛離.不顧一切的竭力尋找,以及慷慨寬容的接納和擁抱,就是整個「補贖」和「回歸」的努力和奮鬥,構成了人生最為重要的過程.在追求人生的圓滿和完善這一點上,基督教和儒家有著共通的地方.孔子說仁,有不同的解釋,背後就是要弟子找到自己的「缺「,然後努力改善,以期達致人生之完善和圓滿.我曾經寫過一篇小文章就此立論,於此不敘.

五．結語

我今天講的道可能有些人會覺得很玄,但人生本是如此,不易明白.因為有人的生命是一百,有人的生命是十,也有人是三；或者我們中間有人的生命是五,是六,七,八,九,或是六十,七十,八十,總之,我們要明白生命形態的多元性.但在我們的人生歷程中,一定會有迷失,陷落,甚至沈淪,即是離開上帝的安排,走上自我選擇或魔鬼引誘的歧路,成為迷羊,成為浪子.人生之至善,人生之大福,是甚麼呢?乃是尋找的尋見,分開的重聚,叛離的回歸,恢復那完美的整全的和圓滿的生命.

我預備講稿時,沒有注意時序.晚上出來,中大山頭明月當空,突然想起: 中秋已近.中秋節的意義,就是團圓,人的團圓,家的團圓.是人生中,人世間莫大的福分.願那失去的被尋回,丟失的被尋見,叛離的再回家.人人有團圓之福,有整全而不破碎的人生.



\newpage

\section{應使則使還是應使不使}
\label{sec:219}
\textbf{林豪恩先生 -「應使則使」 還是 「應使不使」}
\newline
\newline
\hyperref[sec:218]{< < < PREV SERMON < < <}
~
\hyperref[sec:toc]{[返目錄]}
~
\hyperref[sec:220]{> > > NEXT SERMON > > >}


路加福音 16:1-13
\newline
\begin{longtable}{cl}
\hline
\hline
章節 & 經文 (和合本修訂版)\\
\hline
16:1 & \begin{tabularx}{0.7\textwidth}{X} 耶穌又對門徒說：「某財主有一個管家，有人向主人告管家浪費他的財物。 \end{tabularx} \\ \\ \relax
16:2 & \begin{tabularx}{0.7\textwidth}{X} 主人叫他來，對他說：『我聽到了，你做的是甚麼事？把你所經管的交代清楚，你不能再作我的管家了。』 \end{tabularx} \\ \\ \relax
16:3 & \begin{tabularx}{0.7\textwidth}{X} 那管家心裡說：『主人辭我，不用我再作管家，我將來做甚麼呢？鋤地嘛，沒有力氣；討飯嘛，怕羞。 \end{tabularx} \\ \\ \relax
16:4 & \begin{tabularx}{0.7\textwidth}{X} 我知道怎麼做，好叫人們在我不作管家之後，接我到他們家裡去。』 \end{tabularx} \\ \\ \relax
16:5 & \begin{tabularx}{0.7\textwidth}{X} 於是他把欠他主人債的，一個一個地叫了來，問頭一個說：『你欠我主人多少？』 \end{tabularx} \\ \\ \relax
16:6 & \begin{tabularx}{0.7\textwidth}{X} 他說：『一百簍油。』管家對他說：『拿你的賬，快坐下，寫五十。』 \end{tabularx} \\ \\ \relax
16:7 & \begin{tabularx}{0.7\textwidth}{X} 他問另一個說：『你欠多少？』他說：『一百石麥子。』管家對他說：『拿你的賬，寫八十。』 \end{tabularx} \\ \\ \relax
16:8 & \begin{tabularx}{0.7\textwidth}{X} 主人就誇獎這不義的管家做事精明，因為今世之子應付自己的世代比光明之子更加精明。 \end{tabularx} \\ \\ \relax
16:9 & \begin{tabularx}{0.7\textwidth}{X} 我又告訴你們，要藉著那不義的錢財結交朋友，到了錢財無用的時候，他們可以接你們到永遠的住處去。 \end{tabularx} \\ \\ \relax
16:10 & \begin{tabularx}{0.7\textwidth}{X} 人在最小的事上忠心，在大事上也忠心；在最小的事上不義，在大事上也不義。 \end{tabularx} \\ \\ \relax
16:11 & \begin{tabularx}{0.7\textwidth}{X} 若是你們在不義的錢財上不忠心，誰還把那真實的錢財託付你們呢？ \end{tabularx} \\ \\ \relax
16:12 & \begin{tabularx}{0.7\textwidth}{X} 如果你們在別人的東西上不忠心，誰還把你們自己的東西給你們呢？ \end{tabularx} \\ \\ \relax
16:13 & \begin{tabularx}{0.7\textwidth}{X} 一個僕人不能服侍兩個主；他不是恨這個愛那個，就是重這個輕那個。你們不能又服侍神，又服侍瑪門。」 \end{tabularx} \\ \\
[1ex]
\hline
\hline
\end{longtable}
講題:「應使則使」還是「應使不使」　To Spend or Not to Spend

經文:路加福音16章1-13節

講員:林豪恩先生

場合:香港中文大學崇基學院禮拜堂

日期:2013年9月22日

一．引言

可以討論「應使則使」,還是「應使不使」,前題是「有錢使」.因此,有人稱這議題是快樂的爭論.對於「餐搵餐食餐餐清」的貧困人士,或者每逢月尾等出糧的「月光族」,他們對唐代詩人元稹的名句:「貧賤夫妻百事哀」感到更加切身.幾年前,明愛向晴軒進行了「婚姻衝突及和諧」研究調查 (2008),發現婚姻及家庭和諧與家庭總收入有密切關係.原來,家庭總收入增加,會促進婚姻及家庭和諧.報告指出,當香港貧富懸殊,通漲高企的時候,市民收入不足以抵銷百物騰貴,這種狀況對中低收入的家庭影響最為顯著.這些家庭的收入支柱往往都在做高工時,低工資的工作,因而缺乏時間與伴侶及家人相處,也缺乏時間管理子女及協助理家務,當然也缺乏資源和家人去旅行和消遣娛樂,因此間接影響婚姻及家庭關係.

無錢使可以令人苦惱,奇怪的是,有錢使也不一定令人感到輕鬆愉快.近年,香港政府財政年年有盈餘,香港財政的大管家「大把錢使」,比起那些面對赤字而要四出「撲水」的外國財金官員,他應該感到輕省暢快.然而,事實卻不是這樣,我們的大掌櫃不但沒有說自己輕省暢快,反而告訴公眾自己感到沉重,因為他被某人公開批評應使不使,不應使的卻使,並指責他是大罪人.

這個發生在此時此地的劇本,竟然與二千多年前耶穌所講述的一個故事遙遙相應,難免令此時此地的讀者產生聯想.故事的主角也是一位大掌櫃,他也被別人投訴,而投訴的罪名也是「應使不使」,浪費錢財.故事中的大掌櫃如何回應投訴呢?這投訴如何影響他的命運呢?當然,我們必需要問的是,耶穌藉著這個故事對門徒作甚麼教導呢?

二．經文

16 : 1  耶穌又對門徒說:「某財主有一個管家,有人向主人告管家浪費他的財物.

16 : 2  主人叫他來,對他說:『我聽到了,你做的是甚麼事?把你所經管的交代清楚,你不能再作我的管家了.』

16 : 3  那管家心裡說:『主人辭我,不用我再作管家,我將來做甚麼呢?鋤地嘛,沒有力氣；討飯嘛,怕羞.

16 : 4  我知道怎麼做,好叫人們在我不作管家之後,接我到他們家裡去.』

16 : 5  於是他把欠他主人債的,一個一個地叫了來,問頭一個說:『你欠我主人多少?』

16 : 6  他說:『一百簍油.』管家對他說:『拿你的賬,快坐下,寫五十.』

16 : 7  他問另一個說:『你欠多少?』他說:『一百石麥子.』管家對他說:『拿你的賬,寫八十.』

16 : 8  主人就誇獎這不義的管家做事精明,因為今世之子應付自己的世代比光明之子更加精明.

16 : 9  我又告訴你們,要藉著那不義的錢財結交朋友,到了錢財無用的時候,他們可以接你們到永遠的住處去.

16 : 10  人在最小的事上忠心,在大事上也忠心；在最小的事上不義,在大事上也不義.

16 : 11  若是你們在不義的錢財上不忠心,誰還把那真實的錢財託付你們呢?

16 : 12  如果你們在別人的東西上不忠心,誰還把你們自己的東西給你們呢?

16 : 13  一個僕人不能服侍兩個主；他不是恨這個愛那個,就是重這個輕那個.你們不能又服侍上帝,又服侍瑪門.」

三．釋經與應用

1．精明的管家可有值得學習之處嗎?

耶穌的故事是這樣說的:

16 : 1 「某財主有一個管家,有人向主人告管家浪費他的財物.

16 : 2  主人叫他來,對他說:『我聽到了,你做的是甚麼事?把你所經管的交代清楚,你不能再作我的管家了.』

16 : 3  那管家心裡說:『主人辭我,不用我再作管家,我將來做甚麼呢?鋤地嘛,沒有力氣；討飯嘛,怕羞.

16 : 4  我知道怎麼做,好叫人們在我不作管家之後,接我到他們家裡去.』

16 : 5  於是他把欠他主人債的,一個一個地叫了來,問頭一個說:『你欠我主人多少?』

16 : 6  他說:『一百簍油.(四千公升)』管家對他說:『拿你的賬,快坐下,寫五十.』

16 : 7  他問另一個說:『你欠多少?』他說:『一百石麥子.(四萬公升)』管家對他說:『拿你的賬,寫八十.』

16 : 8  主人就誇獎這不義的管家做事精明,因為今世之子應付自己的世代比光明之子更加精明.

耶穌對門徒說這個管家的故事,是否有甚麼要門徒向他學習呢?有人認為,雖然故事中的管家的行為極之具爭議性,但行事卻非常精明,用現代的語言說,他擁有高超的危機管理的能力,因此,倡議者列舉下列值得學習的地方給我們參考:

首先,管家洞識局勢,知道大勢己去,沒有對被解僱作不合理的幻想及無謂的糾纏；第二,他臨危不亂,沒有驚慌失措,仍然能夠冷靜地分析自己的處境；第三,除了分析處境,他也沒有忘記分析自己的能力,避免藥石亂投而走了不配合自己的路.第四,他不但能夠在危機中保持分析能力,更能仔細計劃,籌劃前路.最後,這位管家不是那些「講就天下無敵,做就有心無力」的人,他展現出執行計劃的能力,並且行動迅速,令計劃實現.

如果,我們把這位管家和上文十五章所述浪子的比喻中的小兒子比較,便更能發現這位管家的精明.浪子比喻中的小兒子也用金錢結交朋友,但他所花的金錢完全沒有回報,因為當這位小兒子饑寒交逼的時候,他朋友都離他而去,沒有一個人接濟他.十六章提及的管家則不同,他所使用的金錢是有目的及有回報的:「好叫人們在我不作管家之後,接我到他們家裡去」.更值得留意的地方是,浪子的比喻中的小兒子使用的是自己的錢,管家使用的是別人的錢,難怪管家「使錢」如此「疏爽」.把這兩位人物參照比較,精明程度高下立見.

不過,讓我們停一停,想一想.先不問上述的教導的信仰意義何在,就算不問信仰意義,這位管家虧空公款,以權謀私的精明果真值得人們學習嗎?這位管家的精明值得記在聖經留存給後世參考嗎?如果要這樣的精明故事,我們身邊絕對不缺.

在文匯報網站上,有一篇報導這樣寫著:2012年第一期某能源領域中央企業與省屬企業高層管理人員安全資格培訓班,5月末在北京西郊賓館悄然拉開帷幕.坐了一整天的「學員」們,入夜後出門沒幾步路,便三五成群湧入一間間挑著紅燈籠的居酒屋.一個小包間裡,「崔局」,「張總」,「蘇院長」,「小鍾」 等近十人開始了一場甚是喧嘩的飯局.「崔局」是飯局的核心人物,也是一行人中的「老幹部」,保持1.5cm的平頭是老崔當兵多年留下的印記.酒菜之間,他們聊起蘇家太太在加拿大生活如何,張家小孩兒在澳洲唸書怎樣,崔局兒子最近回國沒有……等長長的流水席結束,在一旁的記者吃驚地認識到,他們的配偶子女要麼均在國外,要麼一半在國外.他們在酒店裡的雄性吆喝引起了其他餐席的不滿,但沒人知道他們擁有一個共同的身份--「裸官」,「半裸官」.為甚麼要變成「裸官」呢?以下的故事足以說明原因.

1986年原河南省服裝進出口公司總經理在美國成立分公司,隨後把妻子從河南土產進出口公司調到美國分公司當經理.在國內的負責買賣配額和外匯額度;在國外的妻子買賣公司到美國和香港的貨物.1992年省裡換屆,董的靠山退休;1993年初,董明玉斷然出逃,一家三口在美國會合,成為改革開放後有史可查的 「裸官」外逃第一案.【文匯網訊2012-6-25】

故事中的管家瞞上欺下,違反誠信,不過,由於他處身的時空與我們太遠,也沒有對我們做成甚麼影響,所以我們比較容易抽空去談論他的處理危機如何精明.但是,如果類似人物出現在你的辦公室,或者類似的行為發生在我們的公職人員中,我們很可能會感到憤怒,而不會以他們的精明來作為學習素材了.那麼,耶穌為甚麼講述這位管家的故事呢?他要對門徒說甚麼呢?

2．也是神國的比喻

有一個常用而又簡單的方法能夠幫助我們明白耶穌說這個管家的故事及其後的評論的意思,就是閱讀這個段落的前文和後文,因為上下文往往提供伏線給我們理解目前的事情.從第十三章開始,讀者會發現越來越多關於神的國的講論,例如:「神的國好像甚麼?」(13:18)「我拿甚麼來比神的國呢?」(13:20)「從東,從西,從南,從北將有人來,在神的國裡坐席.只是有在後的,將要在前,有在前的,將要在後.」(13:29)接下來直至我們今天讀的這個段落之前的內容,都是講述為甚麼有些人在前,有些人在後,以及上主如何從東西南北找尋失落了的人來到神的國.後文19-31節財主和拉撒路的故事雖然沒有用上神的國這詞語,但故事的意境,亳無疑問是關於終未的,是關於彼岸的.因此,目前這個段落,管家的故事以及接著的評論,我們都可以視為是整個關於神的國的講論的一部份.

如果以神的國的脈絡去閱讀16章1-13節這個段落,大家就會發現,1-8節管家的故事和10-13節耶穌的評論這兩段看似不相干的部份,連結的地方就在第9節:「我又告訴你們,要藉著那不義的錢財結交朋友,到了錢財無用的時候,他們可以接你們到永遠的住處去.」這連結也提示了門徒應該以甚麼角度去理解管家的故事.當我們再閱讀管家的故事的時候,我們不妨再想一想,故事中的管家所求的是甚麼呢?

16 : 3  那管家心裡說:『主人辭我,不用我再作管家,我將來做甚麼呢?鋤地嘛,沒有力氣；討飯嘛,怕羞.

16 : 4  我知道怎麼做,好叫人們在我不作管家之後,接我到他們家裡去.』

管家在第4節說:「接我到他們家裡去」,和耶穌在第9節對門徒說:「他們可以接你們到永遠的住處去」互相呼應.如果我們再看前文浪子的比喻,剛好是一個浪子被父親迎接回家的故事,而後文財主和拉撒路的比喻,也是一個關於永恆的住處的故事,拉撒路回家了,安躺在亞伯拉罕的懷中,財主卻回不了家.從語境上的轉變來看,進入神的國,充滿著回家的意味.

3．神的國的價值不在於精明,而在於忠心

故事中的管家明白自己真正需要的不是錢財,而是一個接待他的家,因此,他精明地以錢財去換取一個接待他的家.耶穌提醒他的門徒,你們知道自己的需要是甚麼嗎?是錢財?還是一個永恆的家?如果門徒以為是錢財,而不是一個永恆的家,或者知道自己的需要是永恆的家,但不願意以錢財去換取,就落入耶穌在13節對門徒的警告:「一個僕人不能服侍兩個主；他不是恨這個愛那個,就是重這個輕那個.你們不能又服侍上帝,又服侍瑪門.」這就正如耶穌在第8節的評論:「今世之子應付自己的世代比光明之子更加精明」,因為今世之子精通今世的遊戲規則,現得精明出色,但光明之子卻不懂得神的國的遊戲規則,玩得進退失據,「兩頭不到岸」.

管家的故事暗示了今世的遊戲規則,也展示了那位管家玩得如何精明和出色.耶穌接下來的評論並非教導門徒去學今世之子在今世的遊戲中致勝,而是說明神的國的遊戲規則是怎樣的.如果,門徒要好像那位管家一樣,最終有家可歸,就要好像那位管家一樣熟識永恆的家的遊戲規則.

如果這個世界的致勝之道是欺騙,致勝的能力是精明,相對來說,神的國的遊戲規則是甚麼呢?致勝之道又是甚麼呢?讓我們以這個角度去聆聽耶穌對門徒的說話:

16 : 10  人在最小的事上忠心,在大事上也忠心；在最小的事上不義,在大事上也不義.

16 : 11  若是你們在不義的錢財上不忠心,誰還把那真實的錢財託付你們呢?

16 : 12  如果你們在別人的東西上不忠心,誰還把你們自己的東西給你們呢?

16 : 13  一個僕人不能服侍兩個主；他不是恨這個愛那個,就是重這個輕那個.你們不能又服侍上帝,又服侍瑪門.」

耶穌在這段說話中用了幾項對比:不義的錢財相對於永恆的住處；處事不義相對於處事忠心；不義的錢財相對於真實的錢財；別人的東西相對於自己的東西；瑪門相關於上帝.總括來說,最多出現的是不義和忠心.不義,可以指虛假和此世的意思,因此,耶穌稱那位精明的管家為不義的管家,即是虛假,不忠實的管家,不義的錢財,即是只能在這個世界使用而不能帶到永恆的錢財,因此在永恆的角度來看也是虛假和不真實的錢財.這位虛假不忠實的管家,用虛假不真實的錢財,精明地在這個虛假不真實的今世生存.

如果這個世界的遊戲規則是虛假,耶穌希望他的門徒知道,神的國是真實的,也要追求真實；如果玩這個世界的遊戲的致勝之道是精明和鑽空子,耶穌指出神的國獎賞的是忠心和可信賴.遊戲的規則不同,不過,遊戲的處境卻是相同,門徒和那位管家生活在相同的空間,使用相同的錢財.耶穌提醒門徒,生活在這個世界跟隨甚麼遊戲規則,如何使用這個世界的錢財,是進入神的國有關的,是和進入永恆的家有關的.原來,此岸和彼岸是有關的,現世和永恆是有關的,虛假和真實是有關的.如何相關於?我們可以參考耶穌在第9節的說話:「我又告訴你們,要藉著那不義的錢財結交朋友,到了錢財無用的時候,他們可以接你們到永遠的住處去.」這句說話是甚麼意思呢?接下來耶穌講述財主和拉撒路的故事可視為一項具體案例.

16 : 19 「有一個財主穿著紫色袍和細麻布衣服,天天奢華宴樂.

16 : 20 又有一個討飯的,名叫拉撒路,渾身長瘡,被人放在財主門口,

16 : 21 想得財主桌子上掉下來的碎食充飢,甚至還有狗來舔他的瘡.

16 : 22 後來那討飯的死了,被天使帶去放在亞伯拉罕的懷裡.財主也死了,並且埋葬了.

16 : 23 他在陰間受苦,舉目遠遠地望見亞伯拉罕,又望見拉撒路在他懷裡,

16 : 24 他就喊著說:『我祖亞伯拉罕哪,可憐我吧!請打發拉撒路來,用指頭尖蘸點水,涼涼我的舌頭,因為我在這火焰裡,極其痛苦.』

16 : 25 亞伯拉罕說:『孩子啊,你該回想你生前享過福,拉撒路也同樣受過苦,如今他在這裡得安慰,你卻受痛苦.

16 : 26 除此之外,在你們和我們之間,有深淵隔開,以致人要從這邊過到你們那邊是不可能的；要從那邊過到這邊也是不可能的.』

四．總結

如何運用短暫的錢財於永恆的國度呢?約翰衛斯理曾這樣說:「努力賺錢,努力儲蓄,努力奉獻.」如果這符合你的處境,不妨參考.

不過,錢財豈只鈔票呢?或者,最重要及最有價值的,不是鈔票,而是你這個人.大家又不妨參考下列的例子.

吳思源先生以「我出席了蔡廖水玉女士的追思禮拜」為題,在時代論壇發表了一篇文章,其中幾個段落這樣寫著:

想到七三年底創辦《突破》雜誌的蘇恩佩和蔡元雲,當時才三十歲出頭,在教會內外寂寂無名,竟夠膽辦一份雜誌,推動一個運動,殊不簡單.當時蔡廖水玉女士嫁蔡醫生不久,育有兩個小男孩,正如蔡醫為「突破」而放下他的行醫事業,蔡太也為了持家而放棄了她的護士專業.蔡太一生持守的信念是忠於上帝給她的召命.她常說:「我是個Home-Maker (建家的人)」.屈指一算,蔡太自一九七二年成為全職主婦,以香港專業護士的薪酬計算,她一共損失了接近一千萬港元,足夠在九七回歸前買三層樓,現在增值到最少兩千萬了,但她和蔡醫生從沒把金錢放在心上.八十年代「突破」曾經有一個口號:「抗衡文化,逆流而上」,劃時代地提倡「簡樸生活」.想起蔡醫和蔡太大半生,不就是身體力行的活出了這個信念……蔡醫一生從沒有擁有過一輛汽車,出入安步當車,以他在「突破」一份總幹事的薪水(八十年代初月薪大概是港幣三千元左右)養活一家幾口,蔡太持家有道實屬必然.記得之前有一個早上我剛巧到青年村開會,在村口的小巴車站看見蔡太打著傘子在烈日下候車出沙田,她言笑晏晏的跟我打招呼,這時候剛巧有幾輛汽車魚貫駛進青年村的停車場,是同工駕車上班,蔡太很熱情地向每一輛房車揮手,這情景至今我仍歷歷在目.

上帝賜給我們在這個世界生活的事物,包括金錢,財物,時間,能力,甚至身體,正如耶穌在第9節所說,總會有「無用的時候」,我們如何運用這些賜予呢?運用這些賜予來作甚麼呢?願上主幫助我們,應使則使.



\newpage

\section{為受屈的伸冤}
\label{sec:220}
\textbf{鄧瑞強博士 - 為受屈的伸冤}
\newline
\newline
\hyperref[sec:219]{< < < PREV SERMON < < <}
~
\hyperref[sec:toc]{[返目錄]}
~
\hyperref[sec:221]{> > > NEXT SERMON > > >}


講題:為受屈的伸冤　Justice for All Who Are Oppressed

經文:詩篇146

講員:鄧瑞強博士

場合:香港中文大學崇基學院禮拜堂

日期:2013年9月29日

各位弟兄姊妹,早安.

今日講道的經文,是詩篇146篇.

詩篇146篇至150篇,其實是整部詩篇結集的結語,是整卷詩篇的結案陳詞.

詩篇146篇至150篇這五篇詩篇,都以「你們要讚美耶和華」作開始,也以「你們要讚美耶和華」作結束.「你們要讚美耶和華」的希來伯文便是Hallelujah.這五首詩,被稱為「Hallelujah詩」.詩篇146篇,便是第一首「Hallelujah詩」.這首詩言簡意賅,道出一個讚美神的人,應當如何生活.

這首詩,說其簡單,是很簡單；說其複雜,卻也複雜.

簡單的,是其文字；複雜的,是其結構.

希伯來詩歌,以平行,對比,首尾呼應等手法,層層疊疊,向讀者展示一幅美麗的詩歌圖畫.

我按文字結構及意義結構,將這詩篇分成五個小段.

我花一點時間,講講這首詩篇的文字結構,然後,我回頭講其意義.

第一段是詩146: 1-2:第1節和第2節,是一平行,意義基本相同,只不過第2節深化了第1節講的.

「你們要讚美耶和華!」(第1節開首),是全詩的頭.

詩146:1　我的心哪,你要讚美耶和華!

詩146:2　我一生要讚美耶和華!我還活的時候要歌頌我的神!

第二段是詩146: 3-4:第3節和第4節,意義上也是一平行,第4節解釋第3節.

詩146:3　你們不要倚靠君王,不要倚靠世人；他一點不能幫助.

詩146:4　他的氣一斷,就歸回塵土；他所打算的,當日就消滅了.

第二段是第一段的對立.第一段講的,是「要」做的事,「要」讚美神.講的,是「還活的時候」,「我還活的時候要歌頌我的神」.第二段講的,是「不要」做的事,「不要」倚靠世人.講的,是「氣一斷」,「他的氣一斷,就歸回塵土」.

第三段是詩146: 5-8b(第8節的第二短句),這一段比較長,是全詩的中心.

詩146:5　以雅各的神為幫助,仰望耶和華─他神的,這人便為有福!

詩146:6　耶和華造天,地,海,和其中的萬物；他守誠實,直到永遠.

詩146:7　他為受屈的伸冤,賜食物與飢餓的.耶和華釋放被囚的；

詩146:8 (8ab)　耶和華開了瞎子的眼睛；耶和華扶起被壓下的人.

第5節,第6節:講的是神創造的行動.神是創造主.

第7節,第8節(至8b):講的是神救贖的行動.神是救贖主.

這中心段落的中心點,便是「他為受屈的伸冤」,這是全首詩篇的靈魂所在,也就是今日講道的題目.

第5節:講到「以雅各的神為幫助」的人,這些人是有福的.

這節經文明顯上承及對比之前第二段的第3節經文,「(世人)一點不能幫助」.倚靠世人的人,終歸一無所有.

第7節及第8節(8ab)講神的救贖,就下開後一段的文字,後一段文字也講救贖.

第四段是詩146: 8c-9:這是承接上一段的救贖主題.

詩146:8c　耶和華喜愛義人.

詩146:9　耶和華保護寄居的,扶持孤兒和寡婦,卻使惡人的道路彎曲.

第五段是詩146: 10:回應第1段對神的讚美.

詩146:10　耶和華要作王,直到永遠!錫安哪,你的神要作王,直到萬代!

之後,「你們要讚美耶和華!」是結尾,是呼應開頭的「你們要讚美耶和華!」

第一段和第五段互相呼應,互相平行,呼喚我們讚美神.

第二段和第四段,講的是讚美神的人的做人態度,講的是善惡的抉擇.學者稱這些主題為「智慧文學的主題」.

中間的第三段是中心思想:那位創造天地的創造主,正是救贖天地的救贖主.

如此,便是一典型的ABCBA的平行結構(chiastic structure).

以下,我們探討一下這首詩篇的意義內容.

這首詩首尾呼應,第一段和第五段平行.第一段第1節,提到「我的心哪,你要讚美耶和華!」對應第五段第10節的「錫安哪,你的神要作王」.

讚美神,就是接受神在你生命裡作王.

誰是你生命的王,是最終極的人生問題.在金融社會,有人說:Cash is King. 我聽過一個人說,當他面對衝突的處境時,他要「面子」的心,便會掌控他的一切神經.為了面子,會放棄真理,出賣朋友,埋沒良心.「面子」,是他生命的王.

讚美神,就是接受神在你生命裡作王.

第1節講「我的心哪」,第10節講「錫安哪」.這是從自己的生命開始,影響到群體的氣質.第2節講「我的神」,第10節講「你的神」.這是從「我」的信仰開始,讓「你」透過「我」的生命,願意接受「我的神」就是「你的神」.第2節講「一生」,「還活的時候」,第10節講「直到永遠」,「直到萬代」.這是以自己有限的一生,榮耀那永恆的神.

這樣的開頭,對比著這樣的結尾,是否意味著:

作為一個讚美神的人,我們的生命像「香薰」,不單自己馨香,也會令整個群體氛圍,變得馨香?

在生活的過程中,活出光明,以致「我的神」,被別人心悅誠服地接受為他的神?

在自己有限的歲月裡,讓他人體會到那位永恆的神的真實?

如此說來,敬拜神,就不單只是教堂內的活動,更是信徒生活的一種方式.何種生活方式,這便是詩篇146篇第二段及第四段討論的內容.

這詩篇第二段及第四段講的是人生的智慧.

這詩篇第二段講的,是人生的抉擇,讚美神的人,不要倚靠世人,縱或這是世人中最有權勢的君王,因為他們「一點不能幫助」.在現實中,人當然要互相幫助.這裡講的,不是這些,而是講人生的意義終歸何處的問題,講生命的王權終歸誰人的問題.

君王追求的是權,在現實政治裡,為了維護自己的統治權,君王可能漠視一切正義.世人追求的是自利,社會不再是彼此共在的基礎,而是不同利益團體角力的場所.一切行動,是為了替自己的利益團體瓜分到最大的利益.這裡講的「正義」,是自己團體能獲得應得的利益的「正義」,而不是為了最有需要者的利益而放下自己的利益的「愛的正義」.當每個人都追求自利,每個人都追求控制他人的權力時,我們就是在自己坐著的船的底部鑿洞,終將失去存活的基礎.

這詩篇第四段講的,與第二段講的一樣,是人生的抉擇.這裡陳示出義人和惡人的道路.言下之意,是呼喚人選擇義人的道路.我再讀這一段落的文字.

詩146:8c　耶和華喜愛義人.

詩146:9　耶和華保護寄居的,扶持孤兒和寡婦,卻使惡人的道路彎曲.

義人的道路和惡人的道路的對比,其實是詩篇第1篇的主題.詩篇這個開宗明義的主題,在詩集現在這結尾部分,像結案陳詞一樣,又再被提出來.

義人為神喜愛,進入神的生命裡；惡人為神所惡,被神離棄.在詩篇這段文字裡,「義人」和「惡人」這兩個用字,像開引號和關引號,引號中間講的東西,正是分判義人和惡人的基礎.引號中間有三個元素:「耶和華保護寄居的,扶持孤兒,扶持寡婦」.耶和華喜愛的人,是那些喜愛耶和華的道路的人.耶和華的道路,就是保護寄居者,扶持孤兒,扶持寡婦.

在舊約,寄居者,孤兒和寡婦,時常相提並論,三者都是弱勢者中的最弱勢者.寄居者,是活於自己不相屬的世界的人,他們腳下的土地,不是他們的土地,不會承受他們的重量.活在一個會傷害自己的主人家裡的外籍傭工,是寄居者.孤兒,是無父的人,沒有人供應他們的需要,沒有人教導他們人生的正途.寡婦,是社會的邊緣人,她們缺乏尊嚴,卻又無從建立尊嚴.在現代,她們能做的,是執紙皮,而執紙皮,無論如何也不能建立她們的尊嚴.

神做的,是保護他們,扶持他們.

無論這個世界如何輕視他們,放棄他們,神堅持保護他們,扶持他們.神喜愛那些與他同樣堅持,同樣行動的人.

最後,我要講這詩篇的第三段落,即這詩篇的核心段落.

這段落開首便說:「以雅各的神為幫助,仰望耶和華－－他神的,這人便為有福!」(詩146:5)

成為有福的人,人人都想.一般人說,擁有「健康與財富」(Health and Wealth),便是有福的.又有人說,能得到當下想要的享受,便是有福的.這些都是好的,但都只是從「擁有」的角度看幸福,而聖經的向度,總是「施比受更為有福」,是從「愛的分享」的角度看幸福的.神就是這樣做的.

「耶和華造天,地,海,和其中的萬物；他守誠實,直到永遠.」(詩146:6)神造萬物,一切都是他的.但他不是為所欲為,卻因愛護萬物而甘受束縛.「他守誠實」,即他要按他對萬物的愛的承諾而行.在舊約,神曾自行與人,與萬物訂立種種盟約,承諾按他的愛護而行.創造之主,甘受約束,按愛而行,成為救贖之主.

詩146:7　他為受屈的伸冤,賜食物與飢餓的.耶和華釋放被囚的；

詩146:8　耶和華開了瞎子的眼睛；耶和華扶起被壓下的人.

有學者認為「受屈的」(oppressed,「受欺壓的」)和「被壓下的」意義相近,這兩個字又構成文字上的開引號和關引號,引號中間,像第四段落一樣,又是三個元素:「賜食物與飢餓的,釋放被囚的,開了瞎子的眼睛」.

無食物,無自由,無光明,是人生的慘事.神要救助他們.

將這詩篇的第三段文字連同第四段文字看,詩人一共用了七個動詞,去描述神對弱者的救助:伸冤,賜食物,釋放,開瞎子的眼,扶起,保護,扶持.信息很清楚,神全面地,用不同行動,去救助世間不同的,有需要的人.

整篇詩篇的中心詩句,就是第7節的開頭:「他為受屈的伸冤」.＜現代中文譯本＞譯為:「他為被欺壓的人伸張正義」.神要求什麼?一言以蔽之,「正義」.

一首開頭和結尾都是讚美神的詩篇,中心講的卻是「正義」.讚美神的人,你們知道神的心在想什麼嗎?他想著「正義」.

這些「正義」是抽象的「程序正義」嗎?不是.神要求的,是實質的正義,是「愛的正義」.

「愛的正義」是茶餘飯後拿來作閒談的嗎?不是.「愛的正義」是要求我們實踐的.

為何人實踐「愛的正義」實踐得這麼差?

這是因為我們生命的王權置放錯了,也倚靠無用的力量,而追求幸福的方向也錯誤.

如何重入正軌,走上義路?

真正的讚美神.如這詩篇所言:「你們要讚美耶和華」.

最近聽我的學生的實習分享.

他在暑假,到了一個關懷貧窮人的機構作暑期實習.他和這機構的牧者去探訪住在「板間房」的人,這些板間房有很多木蚤,令住在其中的人受苦.我的學生和該機構的牧者說,機構可以免費為他們滅蚤.那裡的人說:為什麼在香港有這麼好的人?該機構的牧者簡單的說:「是耶穌教我們的.」

無錯,是耶穌教我們的.這真是一個懂得敬拜神的人.

但願:榮耀歸於聖父,聖子,聖靈.阿們.



\newpage

\section{計劃命運與人生}
\label{sec:221}
\textbf{沈祖堯教授 - 計劃, 命運與人生}
\newline
\newline
\hyperref[sec:220]{< < < PREV SERMON < < <}
~
\hyperref[sec:toc]{[返目錄]}
~
\hyperref[sec:222]{> > > NEXT SERMON > > >}


講題:計劃,命運與人生　Your Plan, Your Life and Your Destiny

證道:沈祖堯教授

場合:香港中文大學崇基學院禮拜堂

日期:2013年10月6日

我想在這個題目《計劃,命運與人生》中,與你分享並告訴你,你的計劃,命運與人生,它們彼此之間往往是沒有關係的,但也是出乎意料的.很多時候,你往往計劃了,以為你的一生必是如此,但到了最後,你的命運完全是在你的計劃之外,然而卻是在上帝的計劃之內.

所以,我們需要盤算一下,該如何計劃,面對我們餘下的日子,人生.如果事情發生了,與你之前所計劃的,所要走的路有分別時,你又該如何看呢?

我選擇了聖經裡的《以斯帖記》內的故事.《以斯帖記》記載到當猶太人被擄,在亞哈隨魯王的統治之下,發生了一件很特別的事,是與一位女子有關的.當我在唸這卷書的時候,我發現她與我有一些相似的地方,她的一生與我的前半生也有一些相似的地方.從這卷書中,我們也可以去想想,我們如何來計劃人生,面對生命中那些出乎意料的事,最後到達上帝要我們到達的目的地.

《以斯帖記》記載亞哈隨魯王的國家版圖很大,從中東到印度；他有許多的財富,統治許多的民族.一天,他在王宮裡設擺筵席,大宴群臣,展示他榮耀之國的豐富,和他美好威嚴的尊貴,以及美女,太太.當亞哈隨魯王飲酒心中快樂,就吩咐人,請容貌甚美的王后瓦實提出來.那時候,王后瓦實提也在自己的宮中大宴親朋,所以王后瓦實提卻不肯遵太監所傳的王命而來,所以王甚發怒,心如火燒,於是就廢除王后的位分.這事以後,亞哈隨魯王的忿怒止息,王的侍臣對王獻計,不如為王尋找美貌的處女.最後,他選了此卷書的主角──以斯帖.

以斯帖出身相當寒微,她父母死了,由她叔叔末底改撫養.末底改隱藏以斯帖的身份,囑咐她不要將她是猶太人的籍貫宗族告訴人,因為她是亡國奴,別讓人知曉她的出身是如斯卑微,直到她被選為王帝身邊的王后.

那我的出身是如何?我在灣仔酒吧區附近長大,家裡是以修理汽車為主的生意.我父親的出身是相當不錯的,因為祖父在上海做大生意──飯店,珠寶,股票市場,但由於戰爭,就逃來香港,住在外父車房的一個閣樓裡面.小時候,我在一所名為「中國兒童書院」裡就讀(現在已不存在了),是一所相當小的學校.所以,我的出身與以斯帖相若,我也不太想告訴人家,我是在灣仔酒吧區長大,唸的是一所已倒閉的小學….

人生的路,有時候就是很奇怪,走著,走著,走到一個你從未想像的地方.以斯帖走著,走著,走到王宮裡,因著王后失寵,她成為一位替代的王后,並且成為一人之下,萬人之上；而王愛以斯帖過於愛眾女,且說:「你要甚麼,我必賜給你,你求甚麼,就是國的一半,也必為你成就.」以斯帖從一個不讓人知曉的背景,高升到這種的地位,成為王心中最寵愛的王后,原因就只是單單因為王后瓦實提失寵所致.

至於我,我的事業,學業也是一條相當有趣的路.我唸書的時候,沒有考第一,在別人的眼中,我不是那種很出眾,很聰明的學生,但畢竟也在一九八三年於香港大學的醫學院畢業.畢業後,我無法進入我嚮往的瑪麗醫院裡,跟隨著我很崇拜的教授們一起工作,就去了一間相當遙遠的,新開的沙田威爾斯親王醫院裡工作.當我剛開始工作時,是有點不愉快的際遇,到了 2003 年是我人生中最重要的轉捩點.

2003年的沙士,是一個災難,一個前所未有的困難──經濟,民生,生命.但就在這個看似「危」的裡頭,對我來說,卻成了一個「機會」──讓香港許多人感覺到,原來香港的醫護人員也有可取之處,他們在危難之中,不會選擇退避──當時,就掀起我的人生的另一頁.接下來的十年,所發生的事,更加是曲折離奇:從一個對中文大學毫無認識的,到對大學教育,方向,學生,都越來越感到興趣,接下來又出現另一個相當大的轉變.

這個轉變,在《以斯帖記》裡是非常重要的.當時候,有一位奸臣哈曼,他要所有經過朝門的一切臣僕,都向他跪拜,惟獨末底改不跪不拜,因為他是屬於耶和華,只會拜他的上帝.這事不單為末底改招來殺身之禍,所有的猶大人,就是末底改的本族也受到牽連.哈曼對亞哈隨魯王提儀,推行一個政策:「有一種民,散居在王國各省的民中,他們的律例,與萬民的律例不同,也不守王的律例,所以容留他們與王無益.」亞哈隨魯王也沒有去細查事情,就從自己手上摘下戒指,給哈曼用王的戒指蓋印,說:「這民也交給你,你可以隨意待他們.」哈曼就定意要在幾天之內,全然剪除,殺戮滅絕通國的猶太人.這時候,末底改知道,要救猶太人,只有一個辦法,就是找亞哈隨魯王身邊最寵愛的人──以斯帖.

讓我們來唸一下以斯帖記 4:12-17:「12 人就把以斯帖這話告訴末底改. 13 末底改託人回覆以斯帖說,你莫想在王宮裡強過一切猶大人,得免這禍． 14 此時你若閉口不言,猶大人必從別處得解脫,蒙拯救,你和你父家,必致滅亡．焉知你得了王后的位分,不是為現今的機會麼. 15 以斯帖就吩咐人回報末底改說, 16 你當去招聚書珊城所有的猶大人,為我禁食三晝三夜,不喫不喝．我和我的宮女,也要這樣禁食．然後我違例進去見王．我若死就死罷. 17 於是末底改照以斯帖一切所吩咐的去行.」

以斯帖這段話,是一個很重要的人生決定.王不是隨便可以去見的,以斯帖就顯得相當左右為難,末底改說,你別以為你住在宮中,得免這禍,你若閉口不言,耶和華也會用另外一個方法救我們,猶大人必從別處得解脫,不過,你和你父家,你必致滅亡.以斯帖就吩咐人回報,招聚全國的猶大人為她禱告,她就冒生命危險去行,她說了一句很坦誠的話:「我若死就死罷」(if I die, I die).故事的發展,大家都想像得到,是一個完美的大結局──王向她伸出手中的金杖,她向王說了她要說的,最後所有的猶太人得蒙救贖,並將哈曼和他十個兒子都一同掛在木架上.

對我來說,我又有什麼樣的體會呢?

在三年半以前,當我還很高興埋首在我的醫學研究時,無憂無慮地做逸夫書院的院長,也享受每天上班,寫報告,出國講課… 突然間,一個末底改說,他想提名我當下一屆的中文大學校長.我的反應是,不是啊!我不懂得做這些事,這不是我的專長和專業.但上帝卻說:「你別以為你在醫學院安舒,做你想做的,就可以解決現今社會,學校的問題,你不做,我可以找別人來做.」──對我來說,這是一個命運大轉彎,我有點措手不及的.

其實,在聖經裡,上帝也揀選了一些人,讓那些感到非常的驚訝!當摩西上到何烈山上時,看見焚而不毀的荊棘,耶和華就在那裡打發他,將以色列百姓從埃及為奴之地領出來.摩西說,什麼?你找我?我是甚麼人?我素日不是能言的人,我又無兵,無武器,我只不過是一個牧羊人,你一定是在開玩笑!

摩西就是用這種的心態來帶領以色列人出埃及.

到了約書亞記那裡,縱使約書亞在曠野走了四十年,也做了摩西的副手四十年,他說,你叫我帶領這群烏合之眾的以色列人過到河的對岸,與那些高大的民爭戰,奪取那塊流奶與蜜之地?你開玩笑吧!而且還不需要去爭戰,只要繞城,呼喊,城就會倒塌?換了今天,如果神叫我去某個地方,繞城,呼喊,樓房就會倒塌?有沒有搞錯?這怎麼可能?對不對?你一定在開玩笑.

但神常常打發人去做一些事,是人未曾想過的,甚至有一些的先知被打發的時候,他說,我不去.神就把那個人丟進海裡,並安排一條大魚吞了那人,三天之後,魚就把他吐在旱地上,他就去了尼尼微城那裡.

Life is always different from your plan, and destiny can be totally unexpected - 你的人生可能計劃了要做許多的事,可能你計劃去某個城作買賣,明天我移民去那裡,在那裡置房產,做生意,生孩子等等,過一個很安舒的生活.你可以有許多,許多的計劃.

我可以告訴你,當我還沒有接受這個校長的任命時,我也有一個退休的計劃,甚至我已經在珠海買了房子,好預備退休後住在那裡的,沒想到,我接受了這個任命.

你的計劃與你要走的路,以及最後的目的地,往往是不一樣的.在這些的故事裡頭,在我的故事裡頭,或許在你的故事裡頭,你該如何的來面對你人生的計劃呢?計劃又有何用?我要向東,祂又要我向西,或許與我所接受的訓練,所唸的科系,所做的工作,完全無關,那我還需要計劃嗎?我該怎樣去面對人生下一個… 新的一頁呢?在我生命當中,上帝有什麼安排是我所不知道的呢?

在這裡我想與你們分享自已的體會.

第一,上帝在歷史裡──不管在香港,中國,世界的歷史裡──祂仍掌管著.祂掌管著我們未來五十年,一百年,一萬年要發生的事. 在這個歷史的過程當中,上帝會興起一些人,或許是你,是他,你的兒女,我… 做一些祂要在歷史裡完成的事,像摩西就是一個例子.上帝掌管歷史,不過,祂是用人來實現祂要在歷史中完成的事. 紅海絕少會自己分開的,或出現一個很大的自然環境,而改變歷史,往往是用人來改變歷史.

第二,當神興起這些人的時候,這些人往往都覺得自己還沒準備好──not prepared, not qualify, not ready.不是每個人一出生,就想當校長,特首,或任何一個領導人.神常常在人未曾想到的事興起一些人(像孫中山先生剛開始也沒想過要去當一個革命家,他只想當個醫生而已),不曉得何故,神總是揀選了這些人做一些的事.當你去問問這些人的經歷,他們會告訴你,他們是沒預備好的,沒想過的,甚至不一定是夠資格的,他們在那一刻,是否已經預備好去做了呢?未必的.

因此,我們的命運可能是一個驚奇.問題就在於,你有沒有這個信心來接受這個任命呢? 你有沒有這個信心把你的生命交給祂,讓祂來掌管呢?

第三,我特別在以斯帖記裡,學到一個功課,就是方才我們所唸的那段話,末底改說:你莫想在王宮裡就安全,人若知道你是猶太人,你也遭災.這句話使我想到,無論我們今天在那裡工作,幹哪一種行業… 像在大學裡,應該是最受到保護,最自由的.又或許我今天的工作已經很穩定了,有房子,車子,什麼都有了,社會上那些貧富懸殊的問題,也不能影嚮我,我毋須與人去爭奪什麼?

但聖經告訴我們,我們往往不能獨善其身.當歷史的洪流臨到,沒有一個人可以獨善其身,如果我們的社會,國家出現問題,不穩定時,你我不能說,這不關我的事.因為那些事總會跑到你這裡來,你總不能獨善其身.你和你的家人,都會受到影響──這是我在以斯帖記裡所看見的教訓.

另一個教訓,就是當未底改說, 你若閉口不言,上帝也會從別處拯救猶太人,你不做,也會有其他人去做, 但你和你父家,必致滅亡.這裡給我一個很大的提醒,就是:上帝仍然可以成就祂的計劃,因你不聽話,祂的計劃就得稍微有點更改.

這一點很重要!我不能因為我的叛逆,不聽話,而改變歷史的進程.在神的安排,計劃裡,沒有難成的事,也沒有驚奇的事.可能,不是由你來完成這事,而是由其他人來擔起這擔子.這是聖經給我們一個很清楚的教訓.如果,你不做,別人去做,這就不再是你的呼召,異象,也不是與神有一個很緊密的合作關係,你的命運就改變了,你人生最後的目的地就完全不一樣了.我不曉得這會是一個災難,還是一個假的祝福.但肯定的是,你人生的命運改變了,因為你沒有聽從神的話.

各位,我希望這故事對你來說,有獨特的意思,因為我並不想單單和你分享我的故事,我乃想用以斯帖記這個故事,來提醒我們每一個人:上帝在你的生命當中,是有一個計劃,這計劃不一定與你所計劃的相若,甚或是完全的不一樣.但上帝的計劃是要完成祂在歷史裡,那個宏觀的計劃──對於香港,中國,世界,人類的歷史.

你可以參與在這個計劃裡,是一個大的參與,或小的參與,或許你的家人,子女也會參與這計劃….你可以拒絕,說,我不要,我不想作這些事,我不想有這些的責任,不過,這個計劃絕對不會因為我們的不聽從而改變的,這計劃還有其他人可以去完成,只不過,我們的生命,命運,與上帝之間的關係就改變了.

事實上,如果三年半前,我沒有接受這個校長的任命,我的生命也會很不一樣.我可以告訴你,可能我與上帝之間的關係距離是遠了許多,因為我不覺得我每一天都需要倚靠祂的能力來渡過到今天.現在,我天天都需要上帝的能力.

在香港,許多事情都變得那麼政治化,但在這種困難的環境裡,作為一個基督徒,我們應該躲在王宮裡,還是與我們的民族,有血源關係的人一起走過這些艱難的道路,一齊見證上帝要在香港,中國的歷史中要完成的事呢?

我沒有一個答案給你,我希望你自己去找到屬於你自己的人生,願上帝祝福我們每一個人!謝謝!



\newpage

\section{做個感恩快樂人}
\label{sec:222}
\textbf{許開明牧師 - 做個感恩快樂人}
\newline
\newline
\hyperref[sec:221]{< < < PREV SERMON < < <}
~
\hyperref[sec:toc]{[返目錄]}
~
\hyperref[sec:223]{> > > NEXT SERMON > > >}


講題:做個感恩快樂人　Being A Happy Man with A Thankful Heart

講員:許開明牧師

場合:香港中文大學崇基學院禮拜堂

日期:2013年 10月13日

經文:

【路11:17-19】十個麻瘋病人,一個感恩

【林後九15】「感謝神,因祂有說不盡的恩賜!」

【帖前5:16-18】「要常常喜樂,不住的禱告,凡事謝恩.」

祈禱:

引言:

各位,如果有人問你,你會怎樣回答?

今天有無一些值得感恩的事?我活著起來

今年有無一些值得感恩的事?我成為崇基的校牧

人生有無一些值得感恩的事?我成為基督徒

如果你能很快地回答我這三個問題,你就是一位感恩快樂人!

信息

【】一位小姐正在超級市場推廣雞仔餅,凡經過的人都派一件餅.

一位跳皮的小女孩看見,立即上前,取了一個.在旁的母親問她:『阿女,你要講什麼?』.

小女孩搔搔頭,想了一陣,然後對那小姐說:『

給我多件雞仔餅?』

【】一位父親謝飯祈禱後不久,說:『今日這些飯為何那麼硬?菜為何那麼林?魚為何那麼腥?』孩子就問父親:『阿爸,上帝應該聽你前面的感恩,還是後面的埋怨?』

多少時候,我們好像那小孩不懂得說:『多謝!』

今天從十個痲瘋病人得醫治事跡,我們要學習感恩的功課:

路加十九章的十個痲瘋病人得醫治事跡中,我們可分5個重點瞭解:

1痲瘋者的苦況,2痲瘋者的祈求,3痲瘋者得醫治,4痲瘋者痊癒的反應,5耶穌的回應.

信息

一.     痲瘋者的苦況

十個長大痲瘋的病人(12)

當耶穌進入那村子,他遇見十個長大痲瘋的病人,遠遠的站著.(12)

遠遠的站著(12)

＃為何要『遠遠的站著』?

→因為患者按摩西律法,不能與人接近,最少要相距五十碼.

大痲瘋的病人,要與家人,親人,朋友,村民,所有人隔離,甚至會堂都拒絕他們進入,獨自過著孤單的生活,身體,心靈,靈性都受傷.

如果一定要走向人群,他們遠遠便要大叫:「不潔,不

潔!」人們一聽就會走開.你試想他們的感受?

○路加記載是耶穌主動進入這個村子,他明知那村子有患大痲瘋的病人,祂也主動進入,去接近他們.沒有家人,親人,朋友,甚至會堂裡的人願意接近患大痲瘋的病人,給與他們關懷或幫助,唯有耶穌可以.

→我們有時或許也有此情況,感到求助無門,無人可依靠,唯有主耶穌是你的幫助.

二.     痲瘋者的祈求

這十個患痲瘋病的病人高聲說:『耶穌,夫子,可憐我們吧!』(13)

\#為何不是求醫好而是可憐他們.

這十個患痲瘋病的病人長期生活在被隔離,孤單無助的生活裡,他們逼切地希望得醫治.他們相信耶穌是他們唯一的救星.他們相信耶穌是夫子,是老師,是先知,是救星,甚至相信耶穌是神,因痲瘋病只有神才能醫治.

○他們十個人一同高聲說,一定很響亮(像狗吠聲),亦很感人,耶穌怎會不理會呢?

○他們說:『可憐我們吧!』像是向上帝求助的呼聲,是毫無辨法,向人求乞的聲音.上帝若憐憫人,人便可以得拯救,得醫治,得自由.

→我們得著恩典也是由於神的憐憫,神的開恩!不是因為你是什麼人?好人,尊貴人,有權勢的人.

三.     痲瘋者得醫治

耶穌看見,就對他們說,你們去把身體給祭司察看,他們去的時候就潔淨了.(14)

\#為何對他們說,你們去把身體給祭司察看?

痲瘋者求醫治,耶穌並沒有說:「我可憐你們,現在你們立即得醫治!」這對耶穌來說是不難的說話.

但,奇怪?耶穌只告訴他們『去把身體給祭司察看』,為什麼要這樣?

因為要他們有信心,相信耶穌能醫治他們!

他們可以懷疑耶穌的能力:

-因為耶穌是遠遠的對他們說話.

-因為耶穌沒有用手踫他們.

-因為他們有十個人,以前他們可能知道耶穌醫過某個痲瘋病人,  但從未見過耶穌一醫可醫十人.

○他們需要有信心,才能得醫治.果然『他們去的時候就潔淨了.』

他們去,就是信心的表現,他們便真的得醫治了!

○  『去把身體給祭司察看』是痲痲病人自覺得醫治後,為要得証明,需要的行動.如果他們沒有信心自己會得醫治,他們是不敢到祭司那裡求証的.

○  十個痲痲病人都有信心.

四.     痲瘋者得醫治後的反應

內中有一個見自己已經好了,就回來大聲歸榮耀與上帝.(15)

又俯伏在耶穌腳前感謝他.這人就是撒瑪利亞人.(16)

十個祈求醫治,十個都好了!只有一個回來感恩,為何如此?

只有一個感恩,九個在那裡?

撒瑪利亞人的感恩表現

1﹒看見自已好了:當他離開耶穌時,他突然看見身體皮膚白色的鱗片消失,再經過祭司的證明,他確實完全痊癒了.

他內心充滿感恩,是耶穌的醫治,上帝的憐憫,使他從此可以過新生活,不再被隔離,過孤單的生活,他十分的開心.他知道不要忘恩,不要忽略感謝上帝.

他看見過往患病的痛苦,今日得醫治的喜樂,這是恩典.

各位,我們要懂得感恩,首先也要懂得看見上帝的恩典,看見別人對你的恩惠.

我們要常常去看見,去發現上帝的恩典.

他不是猶太人,是撒瑪利亞人,上帝也醫治他!

2﹒回來:回來感謝耶穌不是易事,十個只一人做到.因為當他們痊癒後,十居其九立刻想到的是回家告訴家人,kiss 親親太太,hug抱抱兒女；

或者通知各親友他已痊癒了；

或者像廣東人湯雞還神,開過感恩會大吃一頓.

但那九個,那大部份的人忘記了使他們痊癒的是神,是耶穌的憐憫.他們最值得做,最首先的應是『感恩』.

只有一個想到感謝上帝,多謝耶穌,撒瑪利亞人原本可以跟那九個繼續向前走,不必回來 (turn back) ,但他想到感恩的重要,他做對了,他不單接受恩典(Gift),也沒有忘記施恩的人(Giver).

相反那九個只想得恩典(Gift),卻忘記恩典的源頭,施恩的神(Giver).

只有這一個得到耶穌的欣賞,得到歷代信徒的欣賞,更立下感恩的好榜樣.他回來與主建立好關係

3﹒大聲歸榮耀與上帝

為何那個撒瑪俐亞人要大聲歸榮耀與上帝,因為他曾經大聲的求救:『耶穌,夫子,請可憐我!』,他既然大聲求救,也要大聲感恩.

多少時刻我們熱切地求主幫助,卻沒有熱切地感恩讚美,歸榮耀與上帝.你有否大聲唱詩感謝神!我唱詩都是大大聲的!

-他與耶穌距離最少五十碼,竟然也能得醫治,他要感恩!

-他不是猶太人,按照律法是不可以得到醫治的恩典,他原本是不配得恩典的.但上帝卻沒有掉棄他,他實在十分的感動,所以他要大聲歸榮耀與上帝,感謝上帝!

大聲的讚美也是要人聽到,向人見證耶穌是他的大恩人,更是向人傳福音:耶穌是神!祂連當時束手無策的病也能醫治.

4﹒俯伏在耶穌腳前感謝他

那人做了一個很大的動作來感謝耶穌:俯伏在耶穌的腳前,這是表示極大的尊敬,極大的感恩,不單把耶穌當作夫子,更把祂視為上帝,因只有上帝才能醫麻瘋.

祂願意一生以耶穌為主為王,這是一種委身的行動.

【】例:非洲的感謝傳統

西非有種感謝傳統,受過恩惠的人向著恩人下跪,頭伏在地上,表示極大的感謝,令人感動.

人懂得感恩,人就會謙卑,人就會記得神,人就與神建立良好的關系.

這個人真不簡單,他知恩,感恩,報恩:把自己的生命擺上檯.

古時以色列民感恩是獻上牛羊,他沒有牛羊,獻上他自己最寶貝的東西,就是他的『生命』.

五. 耶穌的反應

耶穌說,潔淨了的不是十個人麼?那九個在那裡呢?17

除了這外族人,再沒有別人回來歸榮耀與上帝麼?18

就對那人說,起來,走罷,你的信救了你了.19

耶穌問『那九個在那裡』?

明顯,耶穌對那九個是不滿,身為猶太人也不懂感恩:不懂得向上帝感恩.

那九個為何不來感恩呢?可能由於下列的理由:

第一個說,感恩!容我先與家人見面,遲些才去感謝耶穌.

第二個說,感恩!容我觀察一下,我是否真正康復?

第三個說,感恩!不知這康復期可維持多久?會否復發?

第四個說,感恩!呀,根本無患過如此醜死人病.

第五個說,感恩!呀,我開心到根本忘記感恩的事.

第六個說,感恩!我已經在祭司面前感恩了.獻上肥牛,盡上責任,我個人照足本子辦事!

第七個說,感恩!其實耶穌無做過什麼!只叫我見見祭司.

第八個說,感恩!我的病可能是自己痊癒的!我最近一直有好轉.

第九個說,感恩!耶穌不用那麼拘禮!那麼客套,大家心照嘛!

→我們也可以有許多不感恩的籍口!九比一也提醒我們大部份人都忘記感恩的功課.

○耶穌對他說:『你的信救了你了』your  faith has made you whole:他的信心,他對神的感恩,使他得完全.

▲   一個信徒未懂得感恩,他便算不得完全!

▲那九個沒有感恩的,只得身體的醫治；

那懂得感恩的,不獨得身體的醫治.還得精神的,靈性的醫治!所以是完全的醫治.其實,

醫治的原意是「使之完全!」Healing means to be made whole.

▲   按原文的意思是-耶穌對他說:『起來,跟我走吧,你的信心,你的感恩,使你得到完全的醫治!』

▲   起來!從此他可以振作地起來,健康地起來,不再閃閃縮縮過日子.給耶穌欣賞,給耶穌讚,更是快快樂樂的起來!

身體站起來,心境也積極起來

▲   跟我走,跟從耶穌,做主的門徒,去傳揚福音,多麼有意義呀!

精神煥發起來,做耶穌活的人版見證!

更重要的,這人感謝耶穌,得耶穌的喜悅和欣賞

→那九個沒有感恩的,只得到身體的醫治,身體的健康.這人得到完全的健康:

1身體的健康；得醫治

2心靈的健康:得站起來,積極做人

3精神的健康:得跟從耶穌,做主的門徒

4關係的健康:得到耶穌喜悅和欣賞

懂得感恩得到完整的健康!

經文:

【林後九15】「感謝神,因祂有說不盡的恩賜!」

【帖前5:16-18】「要常常喜樂,不住的禱告,凡事謝恩.」喜樂與感恩有密切關係!

一個人常有二個的思想:上主有說不盡的恩賜(講不盡的恩)

凡事都可謝恩:好事,不好的事皆可謝恩.

感恩的操練

感恩是需要操練養成良好的習慣的:以下是一些建議:

1.     為『得到的東西』感恩

這是最容易感恩的:我得學位,得成就,得健康,得到東西,自然會感恩.神一定有給你的恩典!

-我看見在小學,中學,大學,研究院的日子,有不小值得感恩的事情:

-我有得食,住,有溫書地方,

-我得認識神,得到好配偶,得三個可愛的兒女,得到健康,得到許多讀書的機會.

得到的東西,不一定是物質的,也可以是非物質的,是精神的,例如:

平安的感覺,滿足感,安全感,

甚至生命一種醒悟,醒覺.

2.     為『已擁有的或曾經擁有的東西』的感恩

許多時候,我們都沒有為自己擁有的東西而感恩,卻常為到沒有的東西而埋怨.

【】例:我還有一條好腿

一位工人一條腿癱瘓了,被解僱.朋友問候他說:『現在怎樣呀?一隻腳不能動了!』

他說:「你不要忘記上帝仍然給我有條好腿」學習去看你擁有的東西許多時,我們只看我們失去的,我們沒有的東西.

【】曾經有人做過以下的調查:

○如果你有食物在雪櫃,有衣服在衣櫃,有瓦遮頭,有地方往,你比世上75\%的人富有和幸福.

○如果你銀行有儲蓄,銀包常有錢,你已可列入世上8\%最富有人之一.

○如果你今早起床感到十分健康,無病無痛,你已每日比數以百萬面臨死亡的又患病的人幸福.

○   如果你沒有有遇上戰爭,被囚禁,被虐待,你已比世上數以千萬人幸福.

○   如果你父母仍然健在,夫婦沒有離異,兒女常在身邊,你是今日香港和英美國家少有的幸福之人.

【】霍金也感恩

無意中看到這篇文章,提及霍金也感恩,令人感動.當代科學大師霍金的學術演講剛結束,一位年輕記者走到霍金面前,提出了一個十分不解的問題:『霍金先生,你的病將你永遠地固定在輪椅上,難道你沒有為自己失去了太多而悲傷過嗎?』霍金臉上掛著微笑,用不太靈活的手指,艱難地敲打著胸前的鍵盤,螢幕顯示出以下的幾行文字:

我的手指還能活動,

我的大腦還能思考,

我有終生追求的理想,

有我愛和愛我的親人和朋友,

最重要的是我還有一顆感恩的心……．

驟然間,肅穆的會場上再次響起了如雷的掌聲．年輕的記者的心靈被震撼了,望著霍金那並不高大的身軀,恍惚讀懂了一個十分重要的人生課題,那就是「感恩」.處於傷健景況下的霍金,有許多的限制,許多的困難,許多的無助,而且他的身體情況是會愈來愈差的,他人生的破局良方,就是不看他所無的,而是能夠多看他所有的.

3.     為『當然的事』感恩

我常提醒自己『幸福不是必然的!』『健康快樂也不是必然!』

【】為能自由呼吸而感謝神

一位九十歲的信徒剛動過手術,她見證說:「我從來沒有想到必須為一件小事向神感恩.九十年來我都不曉得必須為自己自由呼吸而心存感謝.現在我的鼻子暢通了,這真是極大的恩典.

當醫生把管子插進我的鼻孔,那真是難受…現在我為了能自由呼吸,整天感謝神.」

【】阿身體,我要感謝你!

這位朋友快離世時,請牧師到他病床前.

朋友說:「我快要見主面了,我想多謝許多在世時,曾經幫助過我的朋友,其中一位就是你!」

牧師很感動,但更感動的是教友說:「牧師:我也要感謝我的身體,包括我的眼,耳,口,鼻,手手,腳腳,它們服事了我九十年了,我感謝我的身體,是我的好朋友!」

各位,你有否感謝天父給你的身體呢?

當我懂得感謝我的身體時,我會愛我的身體.

【】一位年青軍官在戰爭中被殺,並死於一間教堂.他的父母為紀念愛子,戰爭紀念日捐了一筆錢給教會作紀念.牧師請這軍官的父母捐獻前講幾句話,讓會眾瞭解他們的心願,善用這些錢.

有對父母聽完他們的分享,十分感動.在回家途中,做父親的說:「我們也為我們的孩子捐筆錢給教會,好嗎?」母親很不高興回答:「怎可以!我們的孩子並沒有死.」

父答:『就是因為我們的孩子沒有死,我們更應該感恩呀!』

4.     為『失去的,痛苦的事』感恩

「感謝神」這首詩,我很欣賞其中二句:『感謝神賜路旁玫瑰,感謝神玫瑰有刺.』連『刺』也感謝,確不是容易的事.

【】Fanny Crosby是位盲人聖樂家.她是因為醫生用錯藥而盲.

有人問她:『如果她遇到那位醫生,會否責怪他?』

她說:『我不會責怪他,因為他雖然做錯事,但上帝成就美事!我反而要感謝神,因我盲了,更能專心寫作詩歌,頌讚神.』

你可曾想過你失去的東西,可能也值得感恩!

保羅說:「凡事謝恩,…」(帖前5:18)這是信仰的美德.

在「密室」(The Hiding Place)一書中,嘉莉(寇裡) (Corrie Boom)提到自己與姐姐貝西(Betsie),在第二次大戰時被送到另一所集中營,那兒情況更糟,既小且髒,跳蚤又多.

一天早晨,她們讀經讀到帖前五章十六至十八節「要常常喜樂,不住的禱告,凡事謝恩」時,貝西叫停,要為集中營的一切感恩.

起初,嘉莉(寇裡)怎樣也不肯為這麼多跳蚤感謝神,但她勝不過姐姐的堅持,就為集中營的每一樣事,包括跳蚤,向神感恩.

不久,她們竟能在集中營中公開舉行查經班,而不受德國士兵的干涉,原因是德國士兵為了跳蚤太多,不願前來查看.借著這件事,嘉莉學習了「凡事謝恩」這一課.

著名的聖詩〈感謝神〉是筆者很喜歡唱的一首詩歌,雖然這首詩原是來自一首瑞典詩歌,先翻成英文,再而翻成中文,然而歌詞仍是很有味道:

「感謝神賜溫暖春天,感謝神淒涼秋景；感謝神禱告蒙應允,感謝神未曾垂聽,感謝神路旁有玫瑰,感謝神玫瑰有刺……」

當中二句:「感謝神路旁有玫瑰,感謝神玫瑰有刺」特別令人回味.基督徒作家杏林子寫著,她每次唱到這二句時,就有一股溫熱的淚水忍不住要從眼穴裡翻湧而出.

在《感謝玫瑰有刺》一書中記載,年少時,杏林子不明白,「為什麼感謝了溫暖的春天,還要感謝淒涼秋景；禱告蒙應允,應該要謝天謝地；然而沒有應允,不該是忿忿不平,怨聲載道嗎?為什麼要感謝?何必要感謝?玫瑰花的芬芳美麗,人人喜愛,但誰會喜歡它的刺呢?」年紀稍長,杏林子才感悟「玫瑰有刺」的意思,因著她那病痛的刺,激發她寫了許多勵志的書,例如:《生之歌》,《生之頌》,《杏林小記》等書；如果她當初沒有得病的話,也不會有那麼多的書,而且也有可能活得沒有那麼積極,在《生之歌》中,她將生命比喻成一首歌,歌詠其中的奧秘.她透露,在國小六年級時,因她患了類風濕關節炎,這根刺使她身心飽受病痛的煎熬,畢業後被迫家中自修,並逐漸培養出寫作的興趣；雖然常受病痛影響,然而作者還是創作不輟,他的作品充滿為生命奮鬥的意念.除此之外,更可貴的是他能將意志力化作行動力創立伊甸基金會,關顧弱勢,積極鼓舞,協助他們奮發向上.

5感恩原來是具有生命的正能量

這是個十分強大的力量.而且就跟情緒一樣,這能量也會吸引與其共鳴的能量.

換句話說,當你心存感恩,就會引發更多值得感恩的事物.接受別人的付出或服事,不只要心存感謝,還要把你對他的「感謝」親口說出來,用文字寫出來.

美國加州大學戴維斯(Davis)分校教授羅伯特(Robert),長期致力於研究積極思考的力量對人的益處.他發現一周感恩五次的人,會活得更健康,更善於應付壓力.

感恩會使體內的血清素(serotonin)和多巴鞍含量升高,這兩種化學物質可以令人產生愉悅感和滿足感.懂得感謝,將會讓你的心靈與宇宙的造物主聯繫,產生人生的正能量呢.

保羅教導我們說:「……凡事謝恩；因為這是神在基督耶穌裡向你們所定的旨意. 」 (帖前 :五 16~18)

凡事謝恩,是指在任何境況下,都能感恩.事實上,感恩是一種生活態度,是一種能力善於發現生活中的感恩的事項,並能承受甚至享受一切人生的順或逆的境況.一個常常能夠感謝父母,感謝家人,感謝朋友,感謝上帝,過感謝的生活肯定是一個快樂的人.

英國作家薩克雷(William Makepeace Thackeray)說:「生活就是一面鏡子,你笑,它也笑；你哭,它也哭.」

如果我們時時能用感恩的心來看這個世界,便會發現這個世界其實很可愛!

樹上小鳥的輕唱,

太陽無私的光明與熱能,

路旁花朵的芬芳,都會令你感到心曠神怡.

學會培養感恩的心,無論是順境或逆境也能冷靜下來,學習數算上帝的恩典,

凡事謝恩,我們就會發覺一切幸福也不是必然.

每天清晨和睡前想一想,上帝在這天的奇妙恩典,獻上感恩的禱告；每天記下值得自己感恩的事情,讓每天活在感恩中,你會發現幸福人生在你手呢.

感恩創造神蹟

其實感恩是最重要的人生態度.因著我們心存感恩我們可以成就許多事情.在「五餅二魚」的故事,我們讀到當耶穌拿起了小童所獻上的五個餅和二條魚(根據約翰的記載),祝謝了,便餵飽五千人,而且還剩下了十二個籃子的零碎.感恩可以創造神蹟!

美國上世紀其中一位最偉大的教育家是海倫凱勒女士(Helen Keller).海倫凱勒本是一位傷殘人士,又盲,又聾,又啞.然而她卻成就了大事.

在她的傳記中,她寫上:「我感謝上主賜給我這些身體的殘缺；我更感謝上主,衪賜我有能力與勇氣去克服這一切的障礙,從而看到了這個美麗的世界!」因著會感恩,不可能發生的事終會發生.

今天的成年人常指摘年青一代不會感恩,常認為他們所享有的是理所當然.

許多時我們自己又何嘗持有「凡事謝恩」(帖前五18)的態度呢?在日常生活中,我們是否抱怨多過感恩；因為我們怨恨太多而謝恩太少,所以我們錯失了許多寶貴的體驗,錯失了享受生命的機會.

作為基督徒,我們應有「萬物皆是從主而來,我們當把從主而來的獻給主.」(代上二十九14下)的持守.我們的生命是上主所賜,我們當把我們的生命獻給主為主使用.具體地,就是盡力去服侍上主所創造,所愛顧的,尤其是有特別需求的人群.

小結:

1為「得到的」

2為「擁有的」

3為「當然的」

4為「失去的」也能感恩,就能「凡事謝恩」

感恩,產生正能量,創造神蹟

做個感恩快樂人

肯定感恩的價值

【】一句謝謝你

一句謝謝你,可使陌生人變成朋友.

一句謝謝你,可使你受人注意.

一句謝謝你,可使人樂於效勞.

一句謝謝你,可顯出你優美的人格和教養.

請不要吝嗇『一句謝謝你』

常常感恩,積極做人

【】假如

一位信徒甲常常喜樂說:『哈利路亞,感謝主!』

信徒乙終日毫無喜樂.

乙見甲又感謝主,便問甲說:『你常常都感謝主,為什麼我就不能?』

乙:『假如你無飯食,假如你無工開,假如你無地方住,假如你無錢,假如你還能感謝主嗎?』

乙:『你太多假如了,所以你活得不開心,也沒友感恩!』

你不能感恩就是因太多『  假如』?

各位:把感恩成為你生活的操練吧!

總結

一位神僕說:『每人每日都應為自己能夠活著而感恩!』

心會跳,鼻會吸氣,口會講,手會動,腳會行,眼會看,頭腦會思想,這一切都是恩典.『我要感恩,至少因為我活著!』這已是上帝的恩典了!

各位我們經常有感恩的心,感恩的行動,我們也會活得有衝勁,活得快樂!記住今日的講題:「做個感恩快樂人!」

感恩使你生命完整,有健康的身體,心靈,精神,關係,靈性,這是人生最大的幸福!

●給你一個小小的使命:

想想上帝給你的恩典,向上帝說聲:『多謝!』

想想家人(父母,愛人或子女)給你的恩惠,向他們說聲:『多謝!』

想想你的朋友,曾經給你的幫忙,說聲:『多謝!』

『只一個感恩,那九個在那裡?』

你是那一個,還是那九個呢?

一個懂得感恩的人是一個積極的人,

一個懂得感恩的是一個充滿活力的,

一個懂得感恩的更是一個快樂的人.

你是這種人嗎?

祈禱:

主耶穌,許多時候,好像那不懂感恩的麻瘋病人,我們不懂感恩的人,反而有不少的埋怨,不滿；常常忘記向主感恩,向人講多謝,求主赦免.

願主幫助我們常常感恩:常發現神在我們身上奇妙作為,為到得到的感恩,為到擁有的感恩,甚至為到你不想我們有的,失去的,也能感恩.靠主幫助,使我們成為一個完全健康的感恩快樂的人.

阿們.



\newpage

\section{我要向山舉目}
\label{sec:223}
\textbf{伍渭文牧師 - 我要向山舉目}
\newline
\newline
\hyperref[sec:222]{< < < PREV SERMON < < <}
~
\hyperref[sec:toc]{[返目錄]}
~
\hyperref[sec:224]{> > > NEXT SERMON > > >}


詩篇 121:1-8
\newline
\begin{longtable}{cl}
\hline
\hline
章節 & 經文 (和合本修訂版)\\
\hline
121:1 & \begin{tabularx}{0.7\textwidth}{X} 我要向山舉目，我的幫助從何而來？ \end{tabularx} \\ \\ \relax
121:2 & \begin{tabularx}{0.7\textwidth}{X} 我的幫助從造天地的耶和華而來。 \end{tabularx} \\ \\ \relax
121:3 & \begin{tabularx}{0.7\textwidth}{X} 他不叫你的腳搖動，保護你的必不打盹！ \end{tabularx} \\ \\ \relax
121:4 & \begin{tabularx}{0.7\textwidth}{X} 保護以色列的必不打盹，也不睡覺。 \end{tabularx} \\ \\ \relax
121:5 & \begin{tabularx}{0.7\textwidth}{X} 保護你的是耶和華，耶和華在你右邊蔭庇你。 \end{tabularx} \\ \\ \relax
121:6 & \begin{tabularx}{0.7\textwidth}{X} 白日，太陽必不傷你；夜間，月亮也不害你。 \end{tabularx} \\ \\ \relax
121:7 & \begin{tabularx}{0.7\textwidth}{X} 耶和華要保護你，免受一切的災害，他要保護你的性命。 \end{tabularx} \\ \\ \relax
121:8 & \begin{tabularx}{0.7\textwidth}{X} 你出你入，耶和華要保護你，從今時直到永遠。 \end{tabularx} \\ \\
[1ex]
\hline
\hline
\end{longtable}
講題:我要向山舉目(I Lift Up My Eyes to the Hills)

講員:伍渭文牧師

經文:詩篇121:1-8

場合:香港中文大學崇基學院禮拜堂

日期:2013年10月20日

海明威小說The Snows of Kilimanjaro(乞力馬扎羅的雪,或譯雪山盟) 一開始記載了坦桑尼亞的非洲第一高峰-乞力馬扎羅山 (菲洲語「神的殿堂」)山頂附近,有一隻凍僵的豹子屍體,出現一個懸念:豹子到這樣高寒的地方來尋找什麼?

小說的結束給了我們答案.

故事主角是一位美國的知名小說作家,去非洲狩獵時不慎染病瀕臨死亡,他妻子陪著他,等候救援.小說用倒敘法讓小說家開始回憶這一生,包括戰爭,破裂的婚姻生活…….臨終前他感嘆過去曾經才華橫溢,也有很多可歌可泣的經歷,但此刻卻無法付諸於文字了.最後,救援人員來到,把他接上飛機；在飛機上他看到前方潔白壯麗的乞力馬扎羅山(Kilimanjaro).不過,這一切只不過是他的幻覺,妻子發現他已氣絕身亡,小說到此結束.

雪山飛行象徵主角死後靈魂上升到神的殿堂,乞力馬扎羅山就是神的殿堂.

今天詩篇經課121篇詩題稱為上行之詩,詩人要向山舉目,因為這山有上帝的聖殿,這是錫安山,聖殿雄立山上,他往高處行到殿敬拜上主.詩篇119是最長的詩篇,以後從120到134是一組短短的詩,共十五首詩篇,像從雄偉的大陸來到美麗的海島.這十五首詩稱為上行之詩,是信徒進入聖殿時所唱的進堂詩.

每個成年的猶太人,一年三次都會到耶路撒冷的聖殿,一次是逾越節,紀念他們離開埃及；一次是住棚節,紀念他們在西乃山領受律法；另一次是修殿節 (Feast of Dedication),紀念馬迦比(Maccabees)家族公元前 165 年從希臘人手中光復聖殿.統治巴勒斯坦安提阿古四世(Antiochus IV ) 公元前 167 年率軍入到耶路撒冷,強行推行希臘化,聖殿安放宙斯神像,並在壇上獻猶太人認為不潔淨的豬.耶穌曾在修殿節上耶路撒冷,約十22, 23:「在耶路撒冷有修殿節,是冬天的時候,耶穌在殿裡所羅門的廊下行走.」

這十五首上行之詩不單是進場的禮儀詩歌,也是朝聖往耶路撒冷聖殿漫長的旅途中,朝聖者用來彼此鼓勵所唱的詩歌.這些上行之詩也是每一位朝聖者-基督徒-行走天路的屬靈寫照,每首詩歌都提到天路客的屬靈爭扎,試探,和凱歌.

誰是你的靠山 (What are your mountains)?

今日詩篇經課詩121指出天路客面對的一種試探,把上主所造的世界和從世界得到的祝福等同上主.「我要向山舉目,我的幫助從何而來?我的幫助從造天地的耶和華而來.」(vv.1, 2)Eugene Peterson, The Message 繙譯第一,二節很傳神:‘I look up to the mountains; does my strength come from my mountains? No, my strength comes from God, who made heaven, and earth, and mountains.’( 我要向山舉目,我的幫助從山而來嗎?不,我的幫助是從上主而來,他造天造地,也造諸山.)

峰頂經驗不等同上主

十多年前我第一次有機會到歐洲旅行,在從萊比錫(Leizip)到柏林(Berlin)的火車途中認識一位年青的畫家,他正趕往柏林舉行畫展,他的太太是作家和評論家,夫婦都是知識份子,很有主見和識見,我問他的宗教信仰,有沒有到教堂?他說:我的宗教就是藝術,我從繪畫和觀賞畫作得到與宗教信仰一樣的滿足,他說不需要宗教信仰.

八零年代初,我在美國深造神學時,參加美國教會的主日學,有一位美國同學很喜歡跑步,每年都參加馬拉松,他說他喜愛跑步到了一個地步令他懼怕,因為從跑步的過程,他感到有一種類宗教的滿足,他慎戒這滿足代替上主的敬拜,所以他盡量主日不跑步,勉力自己無論如何忙碌也參加主日崇拜.

究竟雪山盟小說的主角,攀登雪山去經歷上帝,還是到雪山尋找上帝呢?小說沒有說明,留給讀者自己思考.很多人在城市居住的人,特別喜歡到西藏旅行,目的就是去經歷體會一種類似宗教的精神洗禮,西藏旅行的過程就是目的,不是手段；他們在西藏不是希望找到比西藏更大,更終極的實體.在美國荷理活不少富有及知名度高的明星,特別喜歡西藏,並熱心支持藏民的人權.他們追求高檔的品味,同時追求不同宗教的屬靈經驗,並從這些得到滿足.追尋世界美好品味,追求世界不同的屬靈經驗(spirituality),對一些人來說,都是追求個人滿足的消費行為.

但詩篇121篇提醒我們,無論是物質或精神的峰頂經驗:最高檔的酒店,最壯麗的風景,最神聖的朝聖旅程,都不能真正滿足我們,因它們來自這被造的世界,祇有創造世界的上主才能滿足我們.

你的靠山可靠嗎?

“I look up to the mountains”,「我要向山舉目」,每一個人心目中,都有一個高山作為攀爬的目標.登高臨絕頂,俯看群山小,站上高崗,大地就在我的腳下,滿有成就感.信主的人也會求上主祝福他們,得到所追求的高山,作為我們的成就和依靠,但詩篇121篇提醒我們,創造天地的主,安放高山的主,他自己才是我們的幫助.

1997年七月發生亞洲金融風暴,韓國不少信奉成功神學的教會陷入信仰危機,成功神學就是一種看法:信仰為我們增值,給我們正能量,祇要我們祈求,上主就慷慨賜我們財富,健康.當時不少信徒破產,生活陷入困境,健康每況愈下,甚至自尋短見；上主好像不理我們,離棄我們.

這詩篇告訴我們,我們的上主大過我們的靠山.他創造天地,也賜給我們靠山,但他高過靠山,他給我們財富,給我們健康,給我們和睦相愛的家庭,但這些可以發生改變:突然我們失去工作,染上頑疾,被家人背叛.正在上映的電影「引力邊緣」(Gravity) 有一段黑色的幽默:一位在太空工作,每逢太空艙經過地上的家的上空,都和地上太太送上飛吻,但五星期回到地上,發現太太和一位律師逃離他了.但上主沒有改變,而且這時候成為我們的保護,使我們的腳不動搖.當我們憂心忡忡,失眠不能入睡,上主也陪著我們:「保護以色列的,也不打盹也不睡覺.」(v. 4)

上主的保護-生命的脆弱

我們觀察到,貫穿詩篇121篇整篇有一個詞語,就是上帝的保護.我們的腳會動搖的,祇有上主才能使我們的腳不動搖 (v. 3) .我們也容易受傷害,往聖殿經過必經的曠野,烈日暴曬,夜間風寒,容易染病；野獸出沒,蛇蟲毒害,危機四伏,朝聖者體會到旅途的兇險,生命的脆弱.你想一想過去的日子你有否遇到生命的危險,上主保護你.我預備這講道時想起兩宗事情,過去在講台我比較少分享個人的經歷,因為信仰建立在上主的話語,不在個人特別的經歷.

中學五年級時有一次我在西貢大清水游泳,望過對面的小島看似很近,想是二千呎吧,就下水游過去,但愈游會發覺不對勁,三十分鐘還未到岸,而且水流愈來愈急,開始著有點懼怕,但己經過半,回頭也有一段路程,就祈禱求上主保守,希望快快游到岸；我對上帝說,你若保我生命,我願意獻身作傳道.還有幾百呎就到岸了,但突然腳抽筋,水流更急,心知不妙,突然間有一漁船經過,把我救起.那船剛剛在那時間出現,不遲不早,我知是上帝拯救了我.

另一次二十多年前在台灣,一位同工駕車送我從台北到台中一神學院講早會,車剛離開超級公路收費站不遠,駕盤不能動,輪軸斷了,是非常罕見的意外,要是在高速行使中,後果不堪設想,我們事後多天想起來猶有餘悸,也感恩不已,是上帝保護了我們的生命.

這詩提到上主的保護,就說明生命本身是脆弱的,我們不可自以為是,可以掌控明天.但俗世的精神(secularism)就是想掌控明天,聖經常常提醒我們,其實明天如何我們還不知道.「嗐!你們有話說:” 今天,明天我們要往某城一年,做賣買得利.” 其實明天如何,你們還不知道.你們的生命是甚麼呢?你們原來是一片雲霧,出現少時就不見了.」(雅五13-15)詩篇用了禮儀祝福的格式,說明我們需要上主的保護.

祝福的格式-進入世界的依靠

詩篇121篇很像民數記祭司的祝福:「願耶和華賜福給你,保護你.願耶和華使他的臉光照你,賜恩給你.願耶和華向你仰臉,賜你平安.」(民六24-26)詩篇開始是敬拜者自問自答:「我要向山舉目,我的幫助從何而來?我的幫助從做天地的耶和華而來.」(vv. 1, 2) 跟著是祭司祝福的話:「他必不叫你的腳搖動,保護你的必不打盹……你出你入,耶和華要保護你,從今時直到永遠.」(vv. 3-8) 這是祭司在聖殿敬拜結束時對行將離開的人祝福的話.

甚麼是祝福呢?祝福當然是好的話,英文祝福(Benediction)文義「美善的話語」.在崇拜結束時,當我們進入世界,面對不能掌控的明天,也面對自己生命的脆弱,明白到需要依賴上主,過每天的生活.

接受祝福是接受自己的限制和脆弱,接受上帝的權柄,和信賴他的慈愛,這和墮落剛剛相反.罪的基本精神是不信賴上主,拒絕上主的愛.上主說:「園中各樣的果子,你可以隨意吃,祇是分別善惡樹上的果子,你不可吃,因為你吃的日子必定死.」(創二16, 17) 禁果又稱分別善惡的果子,能夠分別善惡表示可以獨立自主,不需要上主.不可吃禁果祇是陳述事實,善意的提醒和警告,像提醒枝子不能脫離樹幹,離開就枯毀了.吃分別惡的果子是撕裂和上主長久信任無間的關係,拒絕上主的愛,接受祝福表示信賴和依靠上主的愛.

墮落犯罪後,上主來到樂園說:「你在那裡?」(創三9) 始祖聽到上主的聲音,就藏在園的樹木中躲起來, 但我們崇拜首項宣召(call to worship) 發出 ,會眾就和以歡樂的詩歌.崇拜結束時我們接受上主的祝福-美善的話語,表示對上主無限信賴.我們的禮序剛好是失樂園的逆轉,崇拜就像重拾童真,回到樂園信靠天父,接受他的祝福.

祇要我們信賴他,就得到上主的祝福,我們不需登上高山,才尋見上主；事實上不是每一個人有體力登上乞力馬扎羅山(Kilimangaro,上主的殿堂) ,也不是每一個人有經濟能力到耶路撒冷的聖殿朝拜.然而,至高的上主竟然來到我們中間親近我們,「耶和華在你右邊蔭庇你.你出你入,耶和華要保護你,從今時直到永遠.」(v. 8)

悠然見鞍山 (Re-Enchantment of the Mundane)- 生活的驚喜

詩篇121把上主的超越性和內在性並列.上主是造天造地,同時上主不睡覺不打盹保護我們；在你右邊保護你,你出你入保護你.美國洛杉磯時報今年九月十二日報導,美國航天局(NASA)證實,無人外太空探測器航海家1號(Voyager 1)在發射36年後,已經真的離開了我們太陽系,進入星際空間已超過一年了,它的蹤影無人知道.

這新聞叫我認識到宇宙的浩瀚,航海家1號是存在的,但遠離人間,在人不能測度的地方,不能知曉的地方.今天很少人否認有造物主,但造物主與我們有關係嗎?他造了世界後就遠離我們,我們怎可以接觸造天地的主呢?當我們強調上主的超越性到了極限時,就把上主像航海家1號推出太陽系之外的星際空間,不能測度,不能知曉,和我們日常生活沒有關係,我們稱這看法為自然神論(Deism) .

上主就像一個鐘表匠,為鐘表上了發條就隱退了,也許上主的存在,可作為道德倫理的基礎,德國哲學家康德(Kant) 就說:「頭頂的星空,內裡微少的聲音,使我驚懼.」上主是存在的,否則良心直覺的倫理要求,就沒有超越的基礎了.但上主與我們平常生活毫無關係.世界按自然規律運行,人與人按倫理規范相處,經濟按無形的手運作,規現矩矩,一切自然運行,自動駕駛 (auto pilot) ,沒有驚喜.

但詩篇121篇提到上主既造天地,但同時和我們親近；睡覺時看著我們,出入伴隨著我們,生活中實有上主的蹤影,祇是我們沒有察覺,與他失諸交臂；生活其實潛藏了很多驚喜和意義,對上主來說沒有一宗事情是意外的,祇是我們暫時不明白.

可能有一部分弟兄姊妹都知道,有幾位會友,在我年底御任任職十三年半崇基校牧職務時,把我的講章,禱文,每周靜思選一部分結集成書出版,作為送給我的禮物,我感謝他們的厚意,深感不配.當他們問我為這書起一個名字,作為內容的和主題時,我立時想到用「悠然見峖山」(Re-Enchantment of the Mundane) .因為我們崇基教堂主日崇拜禮序中,認罪,宣赦之後唱詩時,會把帷幕徐徐揭開,象徵蒙赦免後,撥開陰霾,天光臨照,重新上路；穿過玻璃幕牆的十字架,悠然見到蒼蒼峖山,它就在那裡!穩固不動絲毫,猶如上主的慈愛,沒有改變,永不動搖.「悠然見峖山」其實修改自陶淵明詩:採菊東籬下,悠然見南山,在不經意間發現驚喜和樂趣.

上主若不和我們同行,生活就沒有驚喜,失去應有的魅力；生活是精彩的,因為上主出入伴隨,保護我們,我們睡著時,上主也不打盹看顧著我們.若我們把上主排出生活之外,基督祇成為一種精神理念,他的教訓成為普世價值,但和我們日常的生活沒有關係.

詩篇121篇告訴我們,上主是超越偉大,他是造天地的主,但同時又是內在的,他貫乎眾人之中,又住在眾人之內,和我們親近.

上主如何保護我們呢?上主透過朝聖群體的伴隨,保護我們,使我們的腳不動搖.不要忘記,這是上行之詩,是朝聖者上耶路撒冷途中彼此鼓舞一起唱的詩歌,是敬拜群體進入聖殿敬拜時唱的進堂詩.

前幾天我們以前牧養過的教會一位老姊妹離世,終年八十五歲,臨終時三十多位教會弟兄姊妹伴隨著她,師母也在其中,在眾聖徒唱詩聲,她寧靜安詳安息主懷.這姊妹住在政府屋村,有一隻腳微跛,但熱心探訪,常常帶著自費烹調好的雞,送給需要探訪的家庭.她愛教會的會友,會友也愛她.在這位老姊妹彌留之際,上主派聖徒伴隨保護她:「耶和華要保護你,從今時直到永遠.」(v. 8)

我在想,崇基教堂如何能彰顯基督的精神呢?在這裡,除了大學講章,優美的音樂,莊嚴的禮儀,會不會是信仰群體彼此的愛呢?上行之詩提醒我們,在朝聖的旅途上,上主賜下聖徒群體,伴隨我們,使我們的腳不動搖.

登山上行

多瑪斯．牟敦梅頓(Thomas Merton)是近代 著名靈修大師,他的自傳《七重山》(The Seven Storey Mountain) 自一九四八年出版以來即風行全球,譯成二十多種文字.《紐約時報》將之譽為二十世紀的《懺悔錄》,當代最具影響力與重要性的靈性經典,感動,改變了數百萬人的生命.因為這書,美國二戰後退伍軍人,大學畢業生進入修道院,追求獨處靈修生活.

《七重山》的名字取自但丁(Dante)《神曲》(Divine Comedy)的淨界山(Purgatorio) 的七重山嶺,天路客向上攀登途中,遇到層層試探,七宗罪交相誘惑,經過淨界的生命淨化,便到達頂峰的伊甸園,進入天界(Paradiso).這書描述牟敦在經歷了感情,智識與靈性上的掙扎和折磨後,從荒唐無知走向真理與永恆的生命歷程.

在書的結語,牟敦說:「我喜愛獨處唯一的原因-在被造的東西中消失,向它們死去,不知道它們,因為它們提醒我離你多麼遠.它們提醒我:你和這些相距甚遠,雖然你在其中.你創造它們,它們因你而存在,它們把你遮蓋起來,讓我看不見你.我將獨自生活,離開它們.福哉獨處!」(That is the only reason why I desire solitude-to be lost to all created things, to die to them and to the knowledge of them, for they remind me of my distance from You. They tell me something about You: that You are far from them, even though You are in them. You have made them and Your presence sustains their being, and they hide You from me. And I would live alone, and out of them. O beata solitudo! ”) (Thomas Merton, The Seven Storey Mountain. An Autobiography of Faith, 1998: 461) .

我要向山舉目,我的幫助從何而來呢?你向山舉目,鎖定一些目標,因為你認為這些可以給你峰頂的經驗,但其實這些令人欣羨的峰頂經驗祇是過眼雲煙,不能真正和長久滿足我們,因為它們來自被造的世界,祇有創造世界的上主才能滿足我們.

甚麼是你的靠山呢?你的靠山真的可靠嗎?請緊記,最可靠的靠山也會塌下來,當靠山塌下時我們就會跨倒,這時候,上主扶持我們,叫我們的腳不動搖.

我要向山舉目,我的幫助從山而來嗎?不,我的幫助是從上主而來,他造天造地,也造諸山.阿們.



\newpage

\section{這是甚麼態度}
\label{sec:224}
\textbf{鄧瑞強博士 - 這是甚麼態度?}
\newline
\newline
\hyperref[sec:223]{< < < PREV SERMON < < <}
~
\hyperref[sec:toc]{[返目錄]}
~
\hyperref[sec:225]{> > > NEXT SERMON > > >}


路加福音 18:9-14
\newline
\begin{longtable}{cl}
\hline
\hline
章節 & 經文 (和合本修訂版)\\
\hline
18:9 & \begin{tabularx}{0.7\textwidth}{X} 耶穌向那些自以為義而藐視別人的人講了這比喻： \end{tabularx} \\ \\ \relax
18:10 & \begin{tabularx}{0.7\textwidth}{X} 「有兩個人上聖殿去禱告，一個是法利賽人，一個是稅吏。 \end{tabularx} \\ \\ \relax
18:11 & \begin{tabularx}{0.7\textwidth}{X} 法利賽人獨自站著，自言自語地禱告說：『神啊，我感謝你，我不像別人勒索、不義、姦淫，也不像這個稅吏。 \end{tabularx} \\ \\ \relax
18:12 & \begin{tabularx}{0.7\textwidth}{X} 我每週禁食兩次，凡我所得的都獻上十分之一。』 \end{tabularx} \\ \\ \relax
18:13 & \begin{tabularx}{0.7\textwidth}{X} 那稅吏遠遠地站著，連舉目望天也不敢，只捶著胸，說：『神啊，開恩可憐我這個罪人！』 \end{tabularx} \\ \\ \relax
18:14 & \begin{tabularx}{0.7\textwidth}{X} 我告訴你們，這人回家去比那人倒算為義了。因為凡自高的，必降為卑；自甘卑微的，必升為高。」 \end{tabularx} \\ \\
[1ex]
\hline
\hline
\end{longtable}
講題:這是什麼態度!?What Attitude Is This ?

經文:路加福音18:9-14

講員:鄧瑞強博士

場合:香港中文大學崇基學院禮拜堂

日期:2013年10月27日

各位弟兄姊妹,早安.

近日,YouTube有一條片非常流行.片中有一男士當街跪在地上,他面前的,是他的女友.他的女友扯著他的頭髮,不斷掌摑他.這女士罵這男士,竟敢帶其他女人回家.這女士問這男士:「你有沒有當我是你的女友?」

這段片叫全城譁然,何等女人,如此男人.

我看後,第一個反應是:「當街當巷如此羞辱一個人,這是什麼態度!」

心裡有個問題:「你有沒有當他是你的男友?」

我衷心希望大家不要遇上這種女人,也不要作這種女人.

我想,大家應該沒有親身經驗過以上經歷.

但,是否也多少體會以下的經歷呢?

有一個理直氣壯的宗教人士來到你面前,在精神意義上,公開地扯著你的頭髮.在精神意義上,公開地不斷掌摑你.他責罵你,說:「你竟敢作出這種事.」然後,他問:「你有沒有當你自己是一個基督徒?」

我衷心希望大家不要遇上這種宗教人士,也不要作這種宗教人士.

當我們做錯事時,我們渴望有人溫柔地接納我們的認錯,我們渴望有人說:「平平安安地回去吧,你的罪被赦免了.」

今日的講道經文,是今日的福音書經課:路加福音18: 9-14.

路18:9　耶穌向那些仗著自己是義人,藐視別人的,設一個比喻,

路18:10　說:「有兩個人上殿裡去禱告:一個是法利賽人,一個是稅吏.

路18:11　法利賽人站著,自言自語地禱告說:『神啊,我感謝你,我不像別人勒索,不義,姦淫,也不像這個稅吏.

路18:12　我一個禮拜禁食兩次,凡我所得的都捐上十分之一.』

路18:13　那稅吏遠遠地站著,連舉目望天也不敢,只捶著胸說:『神啊,開恩可憐我這個罪人!』

路18:14　我告訴你們,這人回家去比那人倒算為義了；因為,凡自高的,必降為卑；自卑的,必升為高.」

這位法利賽人若用現代的形象表達,大致是這樣的.

他穿著最名貴的西裝,手裡拿著i-phone,翹起二郎腿,坐在教堂最寬敞的位置上.他身邊剛好有一個人在認罪祈禱.他抖著腳,看著這個信徒,然後自言自語地說些東西,看似是祈禱,但更似是對那認罪的人說的.

他自言自語說:「神啊,我真是感謝你.我有很多錢,不用上班,常來教堂親近你,不似得有些人時時說要上班,很久也不來教堂一趟.我一生走運,這樣輕易便成功了,不似得有些人要耍些什麼手段博上位.我有當模特兒的妻子,不似得有些人由朝到晚花盡心思去博取異性歡心,這簡直浪費時間.神啊,我真是感謝你,我可以不似身邊的這個人.還有,我一年去兩次聖地旅行,去朝拜你,去親近你,不似得那些人只顧營營役役,不去親近你.我上個禮拜去了一次慈善晚會,慷慨地捐了錢,還上台跟主人家拍了照.神啊,我真是感謝你.」

他繼續抖著腳,心裡有無比滿足.

若你見到這種人,你會如何想?

我的想法很簡單:「這是什麼態度!」

我希望大家不要遇上這種宗教人士,也不要作這種宗教人士.

回到聖經,看看這個法利賽人.經文說:他站在那裡,自言自語.

他像是在禱告,但又不像是在禱告.他是自言自語,講自己的「威水史」,在「晒命」.

他禱文的第一句是:「神啊,我感謝你.」

很好,他懂得感謝神,這是做人的首要原則.

他禱文的第二句是:「我不像別人勒索,不義,姦淫,也不像這個稅吏.」

這句話便很不好了.他站上道德高地,宗教高地,輕視其他人.

「我不像別人」這講法便將人二分.我站得高一點,他們全在我下面.他們作這壞事,作那壞事,我與他們完全不同.我是好人,其他人都是壞人.我與其他人不一樣.世人都是罪人,我卻是例外.

我們不喜歡這種態度.這算是什麼態度!

我們不喜歡一個人站在道德高地或宗教高地,去指控其他人不夠道德,不夠敬虔.

尤其是當這個法利賽人說「也不像這個稅吏」時,我們的反感更大.這個稅吏就在他的身邊,是他的鄰舍.他這樣說,就是刻意傷害他,刻意刺痛他,刻意羞辱他.在某個意義上,這個法利賽人是抽著這稅吏的頭髮,一巴一巴的打過去.讓他被羞辱,讓自己的義被高舉.

法利賽人是宗教上的敬虔人物,他持守著宗教傳統,對神要求的正義有一敏銳的心.這一切,本來都是好的.只是一念之差,一個宗教人將神的一切恩典,看為是自己擁有的偉大數質,將自己推上道德與宗教的高地,視其他人為低下的人.就是這一念之差,他展示了人類生命最大的邪惡.

主耶穌為何講這比喻呢?這比喻針對什麼人呢?

路18:9　耶穌向那些仗著自己是義人,藐視別人的,設一個比喻.

主耶穌講得很明白,這比喻是講給那些自認為自己比別人更道德,更有宗教性,因而藐視別人,輕看別人,甚至羞辱別人的人說的.

你是這樣的人嗎?

人很容易自以為自己偉大,而別人一文不值.

那位女士對著那位跪在地上的男士說:「你有沒有當我是你的女友?」

她認為那位男士帶女人回家是絕對邪惡的,面對這絕對邪惡,她自己就是絕對正義的.她站在絕對正義的位置上,她認為可以公然羞辱這個男士.難道,她真是絕對正義嗎?當她公然羞辱另一個人時,難道不展示了她的邪惡?

有一次,在餐廳吃東西,有一個侍應不小心弄髒了一個客人的衣服,那客人破口大罵,很多難聽的話都說出口了,那侍應也哭了.這客人站在「顧客就是上帝」的高地上,羞辱一個服事他的人.這樣,不單他的衣裳髒了,他的人格也髒了.

有些老師,對著學生說:「爛泥附不上壁.真蠢!」有些父母,指著兒女說:「生舊叉燒好過生你!」有些雇主對員工說:「看你一世都是打份牛工的了!」這些老師,父母,雇主的態度,算是什麼態度!

老師,其實蒙了上天的恩典,有機會吸收知識,去教育學生.父母,是蒙了上天的恩典,有機會養育兒女.雇主,是蒙了上天的恩典,能賺到足夠的錢,請員工來幫自己的手,或服事自己.他們每個人都領受了恩典.但一個不好的老師,不好的父母,不好的雇主,便藉著這些恩典,將自己放在高位上,去輕視不合自己標準的人.他們不就是主耶穌口中的法利賽人,他們利用了神給他們的恩典,利用了他們的優勢,輕視了他們身邊的人.

若主耶穌講完那個法利賽人的比喻後,要我們競猜:誰是我們生命中的法利賽人,我們猜會是誰?

你可能想:有很多法利賽人呢!這個是,那個也是.

或者想:身邊全部都是.

你想:世人都是罪人,我卻是例外.

或許,當我們認為身邊很多人都是法利賽人時,這想法本身便清楚顯明,我們本身也就是一個法利賽人.當你說,「我不像他們」.這時,你會發現,你的講法和法利賽人的講法,是那麼相似.「我不像他們」,我多麼正義,他們多麼不義.這不就是法利賽人的想法嗎?

今日,我站在這裡,指指點點,也是利用了神給我的恩典,站在一個高位上,說了一些可能輕視了你們的話.我也是一個法利賽人.

在主耶穌的比喻裡,有一個稅吏.

經文說:「那稅吏遠遠地站著,連舉目望天也不敢,只捶著胸說:『神啊,開恩可憐我這個罪人!』」

一個卑微的靈魂,瑟縮在聖殿最隱閉的角落裡.

我最感興趣的是,為何這個人會到聖殿祈禱.他明知別人歧視他,他自己也歧視自己.別人指控他的罪,他自己也指控自己的罪.他遠遠地站著,他自覺要跟那些聖潔的人保持距離.他連頭也不敢抬,他自覺和神有一種距離.他捶著胸,他自覺和自己的生命也有一種距離.他在這裡,要忍受旁人的冷言冷語,要忍受神的目光,要忍受內心的控訴.他不可能在聖殿裡開心地唱歌,他不可能聚會後與牧師握手,他不可能聚會後和別人一起聚餐.他在這裡,是痛苦的.他,為何要回來呢?或許,回到聖殿,是自少養成的一種宗教習慣.或許,在這裡找到一點僅餘的生存寄托.或許,在這裡還找到一點被寬恕的希望.

他像一個快要跌落懸崖的人,緊緊抓住崖邊的一根小草.

法利賽人和這個稅吏,都是可憐的人.但法利賽人看不到自己的可憐,他意識不到自己生命的邪惡.這稅吏明白自己的可憐,也意識自己生命的不義.法利賽人利用了宗教,利用了恩典,利用了神,令自己可以高人一等,可以藐視別人.稅吏在宗教裡,明白自己多麼需要恩典,多麼需要神的寬恕.宗教對他而言,是最後的救命符.

主耶穌這個比喻,讓我們看見,我們與法利賽人,相去不遠.當我們意識到自己與法利賽人相去不遠時,也會發現,我們與那稅吏也相去不遠.

這人性的真面目,是令人失望的.

但是,福音之為福音,就是讓我們在人性的絕望中,因為神的接納,再次燃起希望.

誰明白自己的醜陋,誰願意不再踐踏別人的醜陋,誰願意謙卑下來,這人會發現,神看顧卑微的人.他比那些虛偽自大的人,更接近神.

這比喻的結尾說:「我告訴你們,這人回家去比那人倒算為義了；因為,凡自高的,必降為卑；自卑的,必升為高.」

聽說,在一間教會,牧師用主耶穌這比喻講完一篇道後,他請一個會友作結束祈禱.這會友大聲說:「神啊,我感謝你,我不像這法利賽人.」

或許,我們應該祈禱說:「神啊,我感謝你,你憐憫像稅吏一樣的我.」

但願:榮耀歸於聖父,聖子,聖靈.阿們.



\newpage

\section{拿單為公義發聲}
\label{sec:225}
\textbf{彭順強博士 - 拿單:為公義發聲}
\newline
\newline
\hyperref[sec:224]{< < < PREV SERMON < < <}
~
\hyperref[sec:toc]{[返目錄]}
~
\hyperref[sec:226]{> > > NEXT SERMON > > >}


講題:拿單:為公義發聲　Nathan: Voicing-Out-for-Justice

經文:撒母耳記下11:2-5,14-17,12:1-7

講員:彭順強博士

場合:香港中文大學崇基學院禮拜堂

日期:2013年11月3日

背景

大衛:

l   正當在軍事上的成就和開拓疆土達至高峰的時候.

拿單:

l   一位早期先知,和大衛的諫臣.

l   大衛的連年爭戰將結束的時候,他告訴拿單很想為耶穌華建造居所.但拿單阻止他.拿單預言所羅門會建聖殿,和一個王朝,及大衛的家系將會有一位的彌賽亞到來(撒下七1-7；代上十七1-15).

l   拿單斥責大衛殺烏利亞娶拔士巴,並預言大衛家將要承受可怕的後果(撒下十二10-12):大衛家中發生強姦亂倫之事,他3個兒子他相繼死去,並有內戰(撒下十三至十八；王上一).拔士巴與大衛苟合所懷的孩子也不能存活(撒下十二14).

l   大衛快要離世時,他的兒子亞多尼雅欲篡奪王位(王上一1,5).拿單促拔示巴提醒大衛先前曾應許給所羅門繼位之事,並及時的支持她(王上一10-27).大衛就授權拿單等人,為所羅門加冕(王上一28-53).

l   拿單是一個重要的史家(代上二十九29；代下九29).

l   拿單與大衛一起發展聖殿崇拜的音樂事工,有十分重要的貢獻(代下廿九25).

(一 )大衛向拔示巴 (Bathsheba) 犯姦淫

撒下十一2-5

1 過了一年,到列王出戰的時候,大衛又差派約押(Joab)率領臣僕和以色列眾人出戰.他們就打敗亞捫人(Ammonites),圍攻拉巴(Rabbah).大衛仍住在耶路撒冷.

2 一日,太陽平西,大衛從床上起來,在王宮的平頂上遊行,看見一個婦人沐浴,容貌甚美.3 大衛就差人打聽那婦人是誰.有人說:「她是以連的女兒,赫人烏利亞的妻拔示巴.」4 大衛差人去,將婦人接來.那時她的月經才得潔淨.她來了,大衛與她同房,她就回家去了.5 於是她懷了孕,打發人去告訴大衛說:「我懷了孕.」

解經:

1 過了一年,到列王出戰的時候,大衛又差派約押(Joab)率領臣僕和以色列眾人出戰.他們就打敗亞捫人(Ammonites),圍攻拉巴(Rabbah).大衛仍住在耶路撒冷.

l   拉巴(Rabbah):(meaning “the great one”) was the capital city of the Ammonites.

2 一日,太陽平西,大衛從床上起來,在王宮的平頂上遊行,看見一個婦人沐浴,容貌甚美.

l   David may have had his bed on the roof (see 1 Sam 9:25).

3 大衛就差人打聽那婦人是誰.有人說:「她是以連的女兒,赫人烏利亞的妻拔示巴.」

l   The “wife of Uriah” is mentioned also in the genealogy of Jesus. (Matt 1:6)

l   Uriah probably belonged to Jerusalem’s nobility which may have had Hittite (赫人) associations (cf.Exek 16:3).

l   “Uriah” means “Yahweh is my light.”

4 大衛差人去,將婦人接來.那時她的月經才得潔淨.她來了,大衛與她同房,她就回家去了.

l   Bathsheba’s purification: implies that David had disregarded also the ritual law (see Lev 15:24) since in the ancient Near East menstruation usually entailed cultic uncleanness.

l   Intended to stress that it was a favorable time for conception.

l   After this affair Bathsheba returned to her house, and as far as David was concerned, the relationship was, apparently, at an end. However, very soon there arose far-reaching consequences which influenced many lives and, perhaps, altered the course of Israelite history.

(二) 大衛設計殺烏利亞 (Uriah)

撒下十一14-17

14 次日早晨,大衛寫信與約押,交烏利亞隨手帶去.15 信內寫著說:「要派烏利亞前進,到陣勢極險之處,你們便退後,使他被殺.」16 約押圍城的時候,知道敵人那裡有勇士,便將烏利亞派在那裡.17 城裡的人出來和約押打仗.大衛的僕人中有幾個被殺的,赫人烏利亞也死了.

例證 (Ilustration):

1.      《星夢傳奇》/「香港電視」

2.      「抵制蒙哥馬利巴士」(Montgomery Bus Boycott)

不可容忍不公義之法

1955年3月,一位名叫葛狄詩‧歌雲(Claudette Colvin)的15歲黑人少女拒絕遵守「吉姆‧克羅法」(Jim Crow laws),在巴士上讓坐給一位白人.同年12月1日,露莎‧柏詩(Rosa Parks)同樣地拒絕讓坐給白人,故因而被捕.於是依嘉‧力遜(Edgar Nixon)便策劃「抵制蒙哥馬利巴士」(Montgomery Bus Boycott),並由馬丁路德‧金來帶領這運動.這抵制運動共維持了385天.因著形勢的緊張,馬丁路德‧金的居所被炸彈所炸.馬丁路德‧金也在這個運動期間被捕.最終「聯邦地方法院」(United States District Court)下令結束蒙哥馬利市所有公共巴士的隔離政策.

遊行和抵制為要打倒不公義

馬丁路德‧金在「為自由而遊行」(“Walk for Freedom”)的文章中指出,蒙哥馬利市的抗議活動重點,其實並不在於黑人與白人之爭,而是在於「公義」與「不公義」之爭:「真正的緊張並不是在蒙哥馬利的黑人市民和白人市民之爭,其實是公義與不公義之爭,是光明勢力與黑暗勢力之爭.若是會有一個的勝利(當然將會有一個的勝利),那麼這將不單只是黑人市民勝過白人市民,而是公義的得勝,和不公義的失敗.」

要堅持到底

在蒙哥馬利市的抵制運動中,有76位的神職人員被捕,並約100個市民被起訴.但馬丁路德‧金勸慰人們要敢於犧牲,並要堅持到底,以非暴力和愛來面對:「但我們從起初便明白到,我們確實要面對重大的犧牲,和這都是叫人不舒服的經歷.我們要立定決心去堅持到底,這都是我們所堅持去作的.就算是在被起訴期間,我們仍然堅持非暴力的態度,和以愛是最重要的.」

(三) 拿單(Nathan)斥責大衛罪行

撒下十二1-7

1 耶和華差遣拿單去見大衛.拿單到了大衛那裡,對他說:「在一座城裡有兩個人:一個是富戶,一個是窮人. 2 富戶有許多牛群羊群； 3 窮人除了所買來養活的一隻小母羊羔之外,別無所有.羊羔在他家裡和他兒女一同長大,吃他所吃的,喝他所喝的,睡在他懷中,在他看來如同女兒一樣.

4 「有一客人來到這富戶家裡,富戶捨不得從自己的牛群羊群中取一隻預備給客人吃,卻取了那窮人的羊羔,預備給客人吃.」

5 大衛就甚惱怒那人,對拿單說:「我指著永生的耶和華起誓,行這事的人該死! 6 他必償還羊羔四倍,因為他行這事,沒有憐恤的心.」

7 拿單對大衛說:「你就是那人!耶和華以色列的神如此說:『我膏你作以色列的王,救你脫離掃羅的手； 8 我將你主人的家業賜給你,將你主人的妻交在你懷裡,又將以色列和猶大家賜給你；你若還以為不足,我早就加倍地賜給你. 9 你為甚麼藐視耶和華的命令,行他眼中看為惡的事呢?你借亞捫人的刀殺害赫人烏利亞,又娶了他的妻為妻. 10 你既藐視我,娶了赫人烏利亞的妻為妻,所以刀劍必永不離開你的家.』

解經:

1 耶和華差遣拿單去見大衛.拿單到了大衛那裡,對他說:「在一座城裡有兩個人:一個是富戶,一個是窮人. 2 富戶有許多牛群羊群； 3 窮人除了所買來養活的一隻小母羊羔之外,別無所有.羊羔在他家裡和他兒女一同長大,吃他所吃的,喝他所喝的,睡在他懷中,在他看來如同女兒一樣.

4 「有一客人來到這富戶家裡,富戶捨不得從自己的牛群羊群中取一隻預備給客人吃,卻取了那窮人的羊羔,預備給客人吃.」

l   We have a parable, not an allegory (寓言) or fable (神話傳說).

l   It clearly stresses the enormity of the nature of the offense by means of a contrast of opposites: rich-poor; many flocks and herds-single pet lamb like one’s own child.

l   This is further emphasized by yet another contrast; the rich man acts as an outwardly generous and considerate host yet he is entertaining his visitor with fare obtained by the most despicable means.

l   Involved not only in stealing but also in adultery and murder.

5 大衛就甚惱怒那人,對拿單說:「我指著永生的耶和華起誓,行這事的人該死! 6 他必償還羊羔四倍,因為他行這事,沒有憐恤的心.」

l   該死:lit. “son of death”, usual rendering has been “deserves to die”.

l   他必償還羊羔四倍,因為他行這事,沒有憐恤的心:It is possible that v6a, or the whole verse, is a later explanatory insertion by a scribe (cf. Seebass, ZAW 86 [1974] 204) who was aware that the law required only a fourfold restitution (see Exod 22:1 [MT 21:37]; cf. also Luke 19:8).

例證 (Illustrations):

1. Martin Luther King: article “Suffering and Faith”-

「因著我為我的人民爭取自由,我在過去的幾年間都很少有安靜的日子.我曾五次被捕,並因囚於阿拉巴馬州監獄裡面.我的屋曾兩次被炸.我家人和我經常在死亡的威嚇之下渡日.我曾在刺刀下成為死亡邊沿的受害者.因此,我確實曾被迫害的風暴不斷襲擊.我必須承認,我曾多次感到我不再能夠忍受這個如此大的重擔,並且曾嘗試退隱到更為寧靜和平靜的生活去.但每當這引誘出現時,卻又會有一些東西增強和支持我的決心.我現在學會了,當我們背負上主的軛時,祂的擔子其實是輕散的.」

孤注一擲:退縮或面對

「我個人的試煉教曉了我,別人加給我的苦難之價值.當苦難到來時,我不久便明白到,我可以以兩種的方法來回應我的處境:可以以痛苦來回應,或可以將苦難轉化為創造性的力量.我決意選擇後者.」

試煉:更新自己,學懂助人

「因著明白到受苦的必須,我於是便嘗試將它化為一樣美德.為了要將自己從痛苦中拯救過來,我曾嘗試看我的個人嚴酷考驗為一個機會,以更新我自己,和醫治那些在不幸處境中的人.我在過去的幾年中,都不斷活出一個信念:外間加在我身上的痛苦,是有其救贖的作用!」

更能經歷耶穌的苦和神的真實

「有些人仍然以『十字架』為絆腳石,並且其他人看它為愚拙的,但我卻前所未有地堅信,『十字架』是上帝對於社會和個人拯救所施的大能.因此,我現在能夠如使徒保羅一樣,謙卑而驕傲地說:『我身上背負著主耶穌的傷痕.』我在過去幾年所曾經歷過的受苦和痛苦時刻,卻能夠使我更加靠近神.我是前所未有地更加堅信,有一位個人的神真實地存在著.」

2. 潘霍華 (Dietrich Bonhoeffer; 1906-45):

2.1 教會的定義:「只有真正為他人而存在的教會,才是真正的教會」.

2.2  教會和信徒要關心社會:“Christ for others,” “Church for others,” “Religionless Christianity”



\newpage

\section{她是誰的妻}
\label{sec:226}
\textbf{林豪恩先生 - 她是誰的妻}
\newline
\newline
\hyperref[sec:225]{< < < PREV SERMON < < <}
~
\hyperref[sec:toc]{[返目錄]}
~
\hyperref[sec:227]{> > > NEXT SERMON > > >}


路加福音 20:27-38
\newline
\begin{longtable}{cl}
\hline
\hline
章節 & 經文 (和合本修訂版)\\
\hline
20:27 & \begin{tabularx}{0.7\textwidth}{X} 有些撒都該人來見耶穌。他們說沒有復活這回事，於是問耶穌： \end{tabularx} \\ \\ \relax
20:28 & \begin{tabularx}{0.7\textwidth}{X} 「老師，摩西為我們寫下這話：『某人的哥哥若死了，有妻無子，他該娶哥哥的妻子，為哥哥生子立後。』 \end{tabularx} \\ \\ \relax
20:29 & \begin{tabularx}{0.7\textwidth}{X} 那麼，有兄弟七人，第一個娶了妻，沒有孩子死了。 \end{tabularx} \\ \\ \relax
20:30 & \begin{tabularx}{0.7\textwidth}{X} 第二個、 \end{tabularx} \\ \\ \relax
20:31 & \begin{tabularx}{0.7\textwidth}{X} 第三個也娶過她；同樣地，七個人都娶過她，沒有留下孩子就死了。 \end{tabularx} \\ \\ \relax
20:32 & \begin{tabularx}{0.7\textwidth}{X} 後來，那婦人也死了。 \end{tabularx} \\ \\ \relax
20:33 & \begin{tabularx}{0.7\textwidth}{X} 那麼，在復活的時候，那婦人是哪一個的妻子呢？因為他們七個人都娶過她。」 \end{tabularx} \\ \\ \relax
20:34 & \begin{tabularx}{0.7\textwidth}{X} 耶穌對他們說：「這世代的人有娶有嫁， \end{tabularx} \\ \\ \relax
20:35 & \begin{tabularx}{0.7\textwidth}{X} 惟有配得那要來的世代和從死人中復活的人不娶也不嫁。 \end{tabularx} \\ \\ \relax
20:36 & \begin{tabularx}{0.7\textwidth}{X} 因為他們不能再死，和天使一樣；既然是復活的人，他們就是神的兒子。 \end{tabularx} \\ \\ \relax
20:37 & \begin{tabularx}{0.7\textwidth}{X} 至於死人復活，摩西在《荊棘篇》上就指明了，他稱主是亞伯拉罕的神，以撒的神，雅各的神。 \end{tabularx} \\ \\ \relax
20:38 & \begin{tabularx}{0.7\textwidth}{X} 神不是死人的神，而是活人的神，因為對他來說，人都是活的。」 \end{tabularx} \\ \\
[1ex]
\hline
\hline
\end{longtable}
講題:她是誰之妻?　Whose Wife Shall She Be?

經文:路加福音20:27-38

講員:林豪恩先生

場合:崇基學院禮拜堂主日崇拜

日期:2013年11月10日

一．引言

學生團契的團友發出聖經IQ題問:「在聖經中,哪一種人最「抵死」?」答案是撒都該人,因為他們「殺都該」(被殺都應該).撒都該人是否「抵死」,我們容後再討論.不過,要殺撒都該人,應該不是容易的事.在所羅門作以色列王的時代,有一位大祭司名叫撒都該,相傳「撒都該人」就是大祭司撒都該的後人.有一位聖經學者如此描述撒都該人:「撒都該人在初時不是一個宗教的黨派,也不是一個政治團體,而只是一個社團吧了.在數字上,他們比法利賽人少很多,而且大部分的人是屬於猶太國的貴族,和祭司家族中非富則貴的人家.他們的領袖都是議會的議員長老,軍官,政客,和參與公眾事務管理的官員.他們只有一個野心,就是使自己成為在位的君主所不可或缺的人,好讓他們能照自己的見地來左右甚至指揮國家政府.」今天的福音經課提到,這些高貴顯赫的撒都該人也來會見耶穌,他們的出動所為何事呢?他們與耶穌的討論留給我們甚麼啟發呢?

二．經文

20 : 27  有些撒都該人來見耶穌.他們說沒有復活這回事,於是問耶穌:

20 : 28  「老師, 摩西為我們寫下這話:『某人的哥哥若死了,有妻無子,他該娶哥哥的妻子,為哥哥生子立後.』

20 : 29  那麼,有兄弟七人,第一個娶了妻,沒有孩子死了.

20 : 30  第二個,

20 : 31  第三個也娶過她；同樣地,七個人都娶過她,沒有留下孩子就死了.

20 : 32  後來,那婦人也死了.

20 : 33  那麼,在復活的時候,那婦人是哪一個的妻子呢?因為他們七個人都娶過她.」

20 : 34  耶穌對他們說:「這世代的人有娶有嫁,

20 : 35  惟有配得那要來的世代和從死人中復活的人不娶也不嫁.

20 : 36  因為他們不能再死,和天使一樣；既然是復活的人,他們就是上帝的兒子.

20 : 37  至於死人復活, 摩西在《荊棘篇》上就指明了,他稱主是 亞伯拉罕的上帝, 以撒的上帝, 雅各的上帝.

20 : 38  上帝不是死人的上帝,而是活人的上帝,因為對他來說,人都是活的.」

三．釋經與應用

1．男人的妻子還是上帝的兒女

用現代人的視野來看,撒都該人真是「抵死」,他們「抵死」之處,就是在甚麼時候,甚麼地方見到女人,都是用同一種眼光來看女人.當他們在社群中看到女人,他們會問:「她是誰的妻子?」就算死後復活在天堂看見女人,他們仍然是問:「她是誰的妻子?」我們可以想像,就算他們見到天真活潑的小女孩,也會問她將來是誰的妻子,議論她是否「好生養」.對撒都該人來說,女人只有一個身份,就是男人的妻子.除此之外,別無其他.

為甚麼撒都該人的社會如此把女人的可能性壓縮為只作為男人的妻子這身份呢?為甚麼他們把女人的價值削減為只是作為男人的妻子呢?是否因為當時的男人都太過浪漫,整天掛念著創世記所說的「骨中的骨,肉中的肉」那種親密的伴侶關係呢?讓我們看看經文如何說:(20 : 28) 「某人的哥哥若死了,有妻無子,他該娶哥哥的妻子,為哥哥生子立後.」原來,他們視女人的唯一身份為男人的妻子,並非追求親密的伴侶關係,也不是要為兄弟照顧他的寡婦,而是女人為男人生子續後,讓男人生能傳宗接代這種功能.話說回頭,為甚麼撒都該人如此重視生子續後,存宗接待呢?讓我們從他們對事件的描述中尋找一點原因.

20 : 27  有些撒都該人來見耶穌.他們說沒有復活這回事,於是問耶穌:

20 : 28  「老師, 摩西為我們寫下這話:『某人的哥哥若死了,有妻無子,他該娶哥哥的妻子,為哥哥生子立後.』

20 : 29  那麼,有兄弟七人,第一個娶了妻,沒有孩子死了.

20 : 30  第二個,

20 : 31  第三個也娶過她；同樣地,七個人都娶過她,沒有留下孩子就死了.

20 : 32  後來,那婦人也死了.

20 : 33  那麼,在復活的時候,那婦人是哪一個的妻子呢?因為他們七個人都娶過她.」

在他們間短的敘述中,重覆地說了很多次同一個字,就是「死」這個字:v28「某人的哥哥若死了」；v29「沒有孩子死了」；v31「沒有留下孩子死了」；v32「那婦人也死了」.死,對撒都該人來說,是完全的終結,是完全的消失,甚至是完全的絕望,因為他們不相信復活,他們的世界沒有復活,所以如果他們要克服死亡帶來的中斷,要延續生存,就必需要在現世中生子續後,使這位男人的延續性不致消失.因此,生子續後,成為了當時的男人,甚至是整個社會的焦慮,也成為了他們的劇本的主題和內容.

讓我們想一想,在以生子續後,傳宗接代是的劇本裡,女性可以有甚麼角色和任務呢?又讓我們再想一想,在這套生子續,傳宗接代的人生遊戲中,沒有機會參與其中的女性,或者因為種種原因不能參與其中的女性,他們會有甚麼遭遇呢?她們的生存狀況會是怎樣的呢?社會如何看待她們呢?她們如何體驗自己的價值呢?讓我們從幾千年前的希伯來民族,穿越時空,回到現在.

此時此刻座在這裡參加崇拜的男士,相信都已經克服了生子續後,傳宗接代的焦慮,因此,出生在你們家庭的女孩子,都非常幸福和幸運,因為你們待她們如公主一樣珍而重之.然而,試想想,如果你的女兒「投胎」在香港以北的地方,可能會得到怎樣的對待呢?重慶晚報的報導稱,「重慶市出生人口男女比例嚴重失調,去年11月市統計局的千分之一人口抽樣調查,重慶出生嬰兒男女性別比達139:100(正常男女嬰的出生比例為102:107).重慶市計生委副主任潘曉陽說,在偏遠鄉鎮二胎,三胎的男女嬰比例失衡情形更嚴重,個別地方高達五,六百甚至七百比一百.男女嬰失衡現象在其他城市也同樣普遍,人口專家預測,照此下去,16年後,全大陸將有三千多萬光棍,重慶至少有一百萬以上的光棍.」

這篇報導指出正常男女嬰的出生比例為102:107,而重慶市的比例是139:100,有些地方甚至五百比一百,我們不禁要問,女嬰去了哪裡呢?女嬰去了哪裡,是我們的關心,但卻未必是他們的問題,這篇報導提出的問題,是「16年後,全大陸將有三千多萬光棍」.他們關注的,仍然是男人,他們擔心的是男人娶不到妻子,而不是被消失的女嬰.

在西南部的某個農村,青年男女結婚後,第一個重要任務就是生子.麗珍婚後不久就傳來喜訊,全家上下都懷著「一索得男」的期望,麗珍也希望自己能夠第一時間「完成任務」.等了十個月,嬰兒終於出生,可惜,卻令所有人都失望,因為是女嬰,家人唯有寄望麗珍的第二胎.二年後,麗珍再度懷孕.在農村,雖然可以生第二胎,但這也是最後一次機會,不容有失.麗珍戰戰兢兢,患得患失地等待了十個月,嬰兒終於出生.這晚,家中各人都垂頭喪氣,同時也煩躁不安.最後,女嬰幸運地被容許生存下來,連夜被送往居住於深山的親戚家中.女孩在親戚家中牙牙學語,但沒有人讓她叫爸爸或媽媽.有時,女孩要求零食,親戚會對她說:「你爸爸媽媽這個月還沒有送錢過來,不要說零食,飯也沒得吃.」寄人籬下,女孩感受到自己不受歡迎,沒有人要,彷彿是多餘的一個.

某天,這位女孩終於可以回家了,因為她的媽媽生了一個男嬰.回家,原來不等於回到父母的懷抱.女孩發現,爸爸媽媽的懷抱,只向弟弟開放；爸爸媽媽的笑臉,只向弟弟展現.女孩再度陷入迷惑,到底這兩位她稱為爸爸媽媽的成人,是否她的父母呢?如果是,為甚麼他們只抱弟弟,而不抱自己呢?為甚麼他們只擁著弟弟睡,而拒自己於千里之外呢?為甚麼弟弟撒嬌得到父母的呵護,而自己撒嬌卻被打罵呢?有時,女孩有意無意地做出一些吸引父母注意的行為,希望父母關注她,可惜換來的只是更多的責罵,有時甚至被逐出家門.女孩認定,無論去到那裡,自己都是沒有人要的,都是多餘的,都是沒有價值的.

這位女孩做錯了甚麼事情,以至落入如此艱難的景況呢?很可能不是因為她做了甚麼,或者沒有做甚麼,只是因為她在男性生子續後的劇本裡找不到角色,甚至扮演了阻礙男性存宗接代的負面角色,這成為了她的「原罪」,成為了她一生舉步維艱的沉重包袱.

撒都該人理直氣壯地質問耶穌:「在復活的時候,那婦人是哪一個的妻子呢?」(20:33).耶穌指出撒都該人的想法是錯誤的,女人的基本身份並不是男人的妻子,女人的價值也不是為男人傳宗接代,女人如男人一樣,都是上帝的兒女.他說:「既然是復活的人,他們就是上帝的兒子.」(20:36)

在為男人存宗接代的劇本中,女人被物化了,價值被貶低了.然而,最令女人傷痛的,不是這些抽象的概念,而是失去或者扭曲了最基本而又最需要的親密關係,就是失去父母無條件的愛和重視.無論是男人或是女人,都需要親密關係.誰是我的父親?誰是我的母親?我是屬於誰的呢?有一位福音歌手曾經為這些狀況而深感傷痛,自卑自憐,憤世嫉俗.有一天,她終於知道自己是誰的女兒,終於知道自己屬於誰,終於在父親的愛護與重親中找到自己,體會到自己的身份及價值.她以「911的女孩」這首歌來唱出自己的故事.

在九月十一的早晨,

有一位女孩誕生在一個沒有平安的世界,

只好帶著無奈

她的媽媽早已離開,

爸爸又終日喝醉回來對她的精神虐待,

叫她抬不起頭來.

自卑的心靈陪伴著她孤獨長大,

也許她從來不懂甚麼是愛

我就是那九一一的女孩,我還有我的夢和未來,

天父的介入　讓我明白在愛裡沒有失敗

我喜歡做九一一的女孩,每一天都有天父同在

因為祂的寬恕與愛,讓我自由也釋懷

在九月十一的夜晚,生日蛋糕沒有人分享

她默默向上帝祈禱,許下一個願望

有一個聲音呼喚她要把握現在

只要勇敢去愛不怕受傷害

多年以後再次擁抱,眼淚忍不住掉下來,

自從爸爸接受耶穌　他變得很可愛

破鏡重圓是因為愛,仇恨的心已經拋開,

感謝宇宙萬物的主宰,垂聽我的祈禱到現在!

在某種歷史時空之中,女人被社會局限於有限的身份,甚至失去身份,令她們感到迷失和徬徨.讓我們把目光望離一些,超越特定歷史的時空,看見我們原來是天父的兒女,是被重視和被愛的兒女.

2．把現狀延伸至永恆還是以永恆來改進現狀

那位「911的女孩」可以認識及體會作為天父的兒女的身份,突破了成長背景的限制,修補了扭曲關係的傷痛,除了她的個人際遇和努力之外,她處身的社會能否提供這樣的空間也是重要的條件.如果,她處身的社會由承續撒都該人理念的人掌控,她就不會有這樣的機會了.撒都該人設想了一個這樣的狀況去考問耶穌:(20 : 29 - 33)「 有兄弟七人,第一個娶了妻,沒有孩子死了.第二個,第三個也娶過她；同樣地,七個人都娶過她,沒有留下孩子就死了.後來,那婦人也死了.那麼,在復活的時候,那婦人是哪一個的妻子呢?因為他們七個人都娶過她.」你看,在撒都該人的世界裡,永恆只是現在的延續.女人作為男人的妻子,為男人存宗接代,在撒都該人的世界裡,女人這樣的生存狀態是不會改變的,他們甚至認為這種狀況應該維持永恆－如果有永恆的話.撒都該人對於女人的身份和角色不能改變的堅持,甚至成為了他們反駁復活的論證.

耶穌提議撒都該人不要被他們當下的社會狀況局限了男人和女人的關係的可能性.耶穌指出:(20 : 34 - 36)「這世代的人有娶有嫁,惟有配得那要來的世代和從死人中復活的人不娶也不嫁.因為他們不能再死,和天使一樣」.

撒都該人認為,耶穌的說法是荒謬的,因為他們根本就不相信天使和復活這些事情,他們的信念中沒有耶穌所描繪的永恆遠象,也不相信有這永恆的真實,對他們來說,真實就是現在,現在是唯一的真實；耶穌反駁撒都該人,指出他們才是荒謬.撒都該人的荒謬在於相信「實然」就是「應然」,如果有永恆的真實,永恆的真實都只不過是眼前偶然的歷史處境的延伸,就是他們現在的利益和地位的延伸.耶穌提出,歷史的偶然性不是絕對的,也不是終極的真實.耶穌呼籲人們不要把偶然的現狀當作永恆的真實,反而應該以永恆的真實去檢視偶然的現狀,以永恆的真實作為遠像去改進現狀.

撒都該人和耶穌在這場辯論裡突顯的衝突,並沒有在當時結束.有人延續撒都該人的信念,把偶然的歷史現狀當作永恆的真實,極力去把偶然的現狀維持到千秋萬代；也有人接受耶穌的呼喚,認識到現狀並不一定等於永恆,而以永恆的真實去檢視現狀和改進現狀.這兩種態度,在歷史進程中互相視對方為荒謬,互相角力.

我們的四大民間傳說之一,悲壯的愛情故事,「梁山伯與祝英台」中的祝英台,為甚麼要女扮男裝到杭州去求學呢?因為她處身的社會是不容許女性出門學習的,也沒有教導女生的學校.公開的學習和公開的考試都只為男性而設,是男性的專利.直至晚清,基督教傳教士來華,開設女校,招收女生,當時的鄉親父老,甚至姨媽姑姐的反應都是:荒謬!大家閏秀是三步不出閏門的,女子人家怎能上學校讀書!紳親父老們的腦海中,只有一代一代承繼下來男尊女卑,男外女內的傳統,無法想像女性如男性一樣讀書的圖畫,也沒有男女平等的遠像.對他們來說,實然就是應然,現在就是永遠.傳教士對於鄉親父老的抗拒讓女孩讀書,也感到荒謬,他們看見不一樣的圖畫,在他們的圖畫裡,男人和女人都是天父的兒女,男人和女人在天父的樂園無分高低,不分貴賤地平等存在.他們知道眼前的狀況並非如此,但他們不停留於現狀,他們被遠像所推動,以遠像去改變現狀.

有時,現狀是令人留戀的,尤其是當持份者是現狀的概得利益者的時候,就算有遠像,都可以為了眼前的利益而放棄遠像,把眼前的扭曲和醜惡維持下去.歐洲人在美洲大地開發家園,追求自由平等的新世界,立約說:「凡人生而平等,秉造物者之賜,擁諸無可轉讓之權利,包含生命權,自由權,與追尋幸福之權.」當有人呼籲奴隸主實踐憲法,釋放黑奴的時候,奴隸主的反應也是:「荒謬!釋放黑奴,就等於損失50萬的勞動人口,誰人去承擔生產棉花,煙草,稻米的成本?釋放黑奴,必定拖垮經濟,損害國家長久以來累積的資本.」

然而,有一群人,他們沒有忘記上主給他們的遠像,鍥而不捨地努力令現實接近遠像.有一次,代表整個芝加哥牧師的一個團體攜著名為「上帝的即時釋奴令」來到白宮,要求林肯總統維護及執行憲法.林肯總統沒有積極回應.牧師團其中一位名叫荷瑞斯格裡萊,對林肯的消極態度感到氣憤,寫了一篇《兩千萬人的祈禱》的公開信來炮轟總統.不久之後,黑人被釋放,從奴隸做回自由人.

可惜,百多年多,不少人仍然只是看見眼前的差異和利益:他們的皮膚是黑色的,我們的皮膚是白色的,容許他們搬進來會污染我們的社區,容許他們上車會弄髒我們的車廂,容許他們進入這音樂廳會破壞氣氛,容許他們進這學校會教壞我們的孩子,有些人甚至想,容許他們投票會破壞我們的政治生態.

不過,也有另外一些人,心中仍然懷著從上而來的遠像,他們的眼睛彷彿看見天堂的景象,並且被這景象吸引和推動去改變現狀.他如此描述他所看見的:「當我們讓自由之聲響起來,讓自由之聲從每一個大小村莊,每一個州和每一個城市響起來時,我們將能夠加速這一天的到來,那時,上帝的所有兒女,黑人和白人,猶太教徒和非猶太教徒,耶穌教徒和天主教徒,都將手攜手,合唱一首古老的黑人靈歌:「終於自由啦!終於自由啦!感謝全能的上帝,我們終於自由啦!」

四．總結

耶穌說:「配得那要來的世代和從死人中復活的人」,都是上帝的兒女.個人來說,我們需要這遠像.因為我們可能處身於壓縮我們的身份和價值的經濟和社會制度,可能遇到被貶抑和被輕視的處境.我們需要這遠象來確認自己的身份,肯定自己終極的價值.我們的社會也需要這遠象.最近,我們的大學有一位教授政治學的老師,發表文章提醒我們:「香港社會正處於一個在歷史上少見的分裂局面.土地發展,房屋以超越當下利益分配的零和遊戲.政策,基建設施,全民退休保障,環境保護,舊區重建,文化保育,中港關係,普選特首,佔領中環,對於以上種種的問題,香港內部均是充滿了爭議和對立.香港社會似乎已失去了一種達至共識的能力」,他呼籲香港人:「透過愛和理想,突破由物質差異而引致的身分,地位和意識形態的不同,團結一致,重建香港的未來.」(黃偉豪:撕裂香港－走向兩極化對立的危機)我們需要耶穌所啟示的永恆的真實,並且以這永恆的視野檢視現在,塑造每一刻的現在.



\newpage

\section{苦難是玩意或愛意}
\label{sec:227}
\textbf{許開明牧師 - 苦難是玩意或愛意?}
\newline
\newline
\hyperref[sec:226]{< < < PREV SERMON < < <}
~
\hyperref[sec:toc]{[返目錄]}
~
\hyperref[sec:228]{> > > NEXT SERMON > > >}


路加福音 21:5-19
\newline
\begin{longtable}{cl}
\hline
\hline
章節 & 經文 (和合本修訂版)\\
\hline
21:5 & \begin{tabularx}{0.7\textwidth}{X} 有人談論聖殿是用美石和供物裝飾的，耶穌就說： \end{tabularx} \\ \\ \relax
21:6 & \begin{tabularx}{0.7\textwidth}{X} 「你們所看見的這一切，日子將到，沒有一塊石頭會留在另一塊石頭上而不被拆毀的。」 \end{tabularx} \\ \\ \relax
21:7 & \begin{tabularx}{0.7\textwidth}{X} 他們問他：「老師，甚麼時候有這些事呢？這些事將臨到的時候有甚麼預兆呢？」 \end{tabularx} \\ \\ \relax
21:8 & \begin{tabularx}{0.7\textwidth}{X} 耶穌說：「你們要謹慎，不要受迷惑，因為將有好些人冒我的名來，說『我是基督』，又說『時候近了』，你們不要跟從他們！ \end{tabularx} \\ \\ \relax
21:9 & \begin{tabularx}{0.7\textwidth}{X} 當你們聽見打仗和動亂的事，不要驚惶；因為這些事必須先發生，但終結不會立刻就到。」 \end{tabularx} \\ \\ \relax
21:10 & \begin{tabularx}{0.7\textwidth}{X} 於是耶穌對他們說：「民要攻打民，國要攻打國， \end{tabularx} \\ \\ \relax
21:11 & \begin{tabularx}{0.7\textwidth}{X} 將有大地震，多處必有饑荒、瘟疫，又有可怕的異象和大神蹟從天上顯現。 \end{tabularx} \\ \\ \relax
21:12 & \begin{tabularx}{0.7\textwidth}{X} 但這一切的事以前，有人要下手拿你們，迫害你們，把你們交給會堂，並且關在監裡，又為我名的緣故拉你們到君王和統治者面前。 \end{tabularx} \\ \\ \relax
21:13 & \begin{tabularx}{0.7\textwidth}{X} 但這些事終必成為你們作見證的機會。 \end{tabularx} \\ \\ \relax
21:14 & \begin{tabularx}{0.7\textwidth}{X} 所以，你們要立定心意，不要預先考慮怎樣申辯； \end{tabularx} \\ \\ \relax
21:15 & \begin{tabularx}{0.7\textwidth}{X} 因為我必賜你們口才和智慧，是你們一切敵人所敵不住、駁不倒的。 \end{tabularx} \\ \\ \relax
21:16 & \begin{tabularx}{0.7\textwidth}{X} 連你們的父母、兄弟、親族、朋友也要把你們交給官府；你們中間也將有被他們害死的。 \end{tabularx} \\ \\ \relax
21:17 & \begin{tabularx}{0.7\textwidth}{X} 你們要為我的名被眾人憎恨。 \end{tabularx} \\ \\ \relax
21:18 & \begin{tabularx}{0.7\textwidth}{X} 然而，你們連一根頭髮也不會損失。 \end{tabularx} \\ \\ \relax
21:19 & \begin{tabularx}{0.7\textwidth}{X} 你們憑著堅忍，就必保全性命。」 \end{tabularx} \\ \\
[1ex]
\hline
\hline
\end{longtable}
講題:到底苦難是上帝的玩意,無意,還是愛意? 　Is “Suffering” God’s Stuff or Love?

經文:路加福音21:5-19

講員:許開明牧師

場合:香港中文大學崇基學院禮拜堂

日期:2013年11月17日

1引言:

苦難,跟人類有不解之緣.

苦難,也是人類一個不解之謎.

十年前沙士疫症,年前南丫島海難事件,多個月前埃及熱氣球爆墜事件,最近菲律賓的海燕大風災,千萬人死亡,……人生活在許多苦難裡.

多年前,帶領一個大學生團契研經班,參加者都是成績相當優異的理科生.在研經中提及到苦難問題.當時有位青年問我:「明哥(當年的青年都這樣喚我的),點解上帝讓這個世界有苦難?點解善良的上帝造了美麗伊甸園後,又要造分別善惡樹和蛇,埋下痛苦的根源?點解無所不知的上帝,預知人會犯罪,犯罪後會受罰受苦,卻容許連串的痛苦發生,讓痛苦進入世界.

首先,亞當受汗流滿面得糊口之苦；夏娃受生產兒女之苦；蛇受以肚行路和吃土之苦；跟著罪入世界,該隱殺人所受的刑罰和痛苦更甚.

上帝既然是全知全善全能的主,為何不阻止苦難的發生,甚至消滅苦難?上帝是否玩野?愛眼睜睜的看著人類受苦；然後,又差派祂的獨生子耶穌來到替人受苦,看著自己的兒子,令自已心痛呢!」

一連串的問題,剛神學院畢業的我,當時真是無言以對?當其時,我也清楚這位青年正在失戀中,受著年青人最深的痛苦,最大的打擊.

【】其實,當年剛出道也有同樣疑問,為何堂主任何牧師是那麼好人,一位愛神,愛人,愛教會的良牧,卻要受病痛煎熬呢?踏入退休之年,先是愛妻匆匆離他而去,使他心如刀割；喪妻之痛還在,又發現自己患上頑疾,無愛妻的陪伴,獨自面對痛苦難當的末期癌症,真是痛上加痛.我是見他最後一面的人看著一位好牧者,原本是像雄糾糾的軍人,飽經苦痛瘦削不堪的遺體,心中酸溜溜的感覺.

2信息:

講道的命題:到底苦難是什麼?

是上帝的玩意,無意/有意,還是愛意呢?

本段經文:路21:5-19節,耶穌談到世界末日前,人要面對多種苦難?

按照猶太人的末世觀,也是他們的世界觀,包括

現在的日子→主的日子→主再來的日子.

現在的日子:走向衰敗,滅亡的日子,會愈來愈差!

主的日子:受苦的日子,有不少災難臨到的日子.

主再來的日子,榮耀得勝的日子!

2.1聖殿的被毀,失去聖殿之苦,精神上的苦

根據猶太歷史家約瑟夫的記述,整個聖殿用四十六年才完工.這殿建築宏偉,堂皇壯麗,建殿的大理石,有的甚至長達十公尺以上,另有巨型金飾,真是金碧輝煌,光輝耀目.

聖殿是上主的榮耀,也是猶太人國家民族的榮耀!

【路廿一6】「耶穌就說:『論到你們所看見的這一切,將來日子到了,在這裡沒有一塊石頭留在石頭上,不被拆毀了.』」

結果:在主後七十年按字義完全應驗了,當時羅馬人在提多率領下,徹底毀滅了耶路撒冷及聖殿建築物,甚至石頭也被撬開,以取得當聖殿被燒燬時從屋頂上熔流其間的金箔.也有人說,當時尚存留有一些石頭未被移動.

2.2人禍之苦:

【路廿一10】「當時耶穌對他們說:

『民要攻打民,國要攻打國；」

「民要攻打民,」指大至民族,小至部落村莊之間的戰爭.

「國要攻打國,」指國際間的戰爭.

結果:從主升天以後直到如今,歷世歷代烽火連緬不絕.

2.3 天災之苦:

【路廿一11】「地要大大震動,多處必有饑荒瘟疫；又有可怕的異象,和大神蹟,從天上顯現.」

震動:使人無法安居至死

饑荒:使人饑餓至死

瘟疫:使人患病至死

可怕的異象:使人恐懼至死

2.4為義受逼迫至死

【路廿一12】「但這一切的事以先,人要下手拿住你們,逼迫你們,把你們交給會堂,並且收在監裡,又為我的名拉你們到君王諸侯面前.」

信徒因信主耶穌要受苦:

被拿,被逼迫,收收在監裡,被審,受凌辱之苦!

【路廿一16】「連你們的父母,弟兄,親族,朋友,也要把你們交官；你們也有被他們害死的.」

被至親所害,好友所害,何等的苦!

【路廿一17】「你們要為我的名,被眾人恨惡.」

被眾人人恨惡,非常難受!

從中可見,有個人的苦,群體的苦,肉體的苦,精神的苦,

人際的苦.

以上的苦是人逃避的了的!

以上的苦是將要發生的!

以上的苦未發生,已令人恐懼不安!

預言苦難:令人擔憂!

一位兄弟去睇相,相士告知:「你有一有十年的惡化運!那十年後?又有十年的惡運；跟著你仍然是惡運；以後?不過你會慢慢適應惡運.」

在這世界觀或末世觀,苦難乃是上帝的計劃中,這些苦難是上帝所知道,容許的,甚至是來自上帝的.

那麼,就令人頭痛了!

苦難,可以是來自上帝

路得記1:20 -22拿俄米認為上帝也會容許災禍臨到人的身上.當拿俄米及路得回到家鄉伯利恆時,她對城中的人說:「不要叫我拿俄米,因那是甜的意思,要叫瑪拉,因那是苦的意思.因全能者使我受了大苦.」20

「我滿滿的出去,耶和華使我空空的回來.耶和華降禍與我,全能者使我受苦.」21

一般人的思想,上帝只是賜福的上帝.拿俄米卻有新的體會:耶和華是會降禍,耶和華是會使人受苦.人生的禍福是來自上帝,上帝是掌權生命的神.

慈愛全能的神,怎可以容讓世間有苦難?

聖經啟示神存在,且說這位神是一位慈愛和全能的神.於是我們說:

「如果神是慈愛但不是全能的,祂無法幫助我們脫離苦難,這是我們接納的.

如果神是全能但不是慈愛的,祂不會幫助我們脫離苦難,這也是我們理解的.

但如果神是慈愛又是全能的,祂不去除掉世間的苦難,尤其是連祂的兒女也要受苦,那是不可以接受的.

既然世間有苦難,連基督徒也有苦難,那就證明基督教所講的慈愛全能的神是不存在的.」

這是個很嚴肅的問題.

若然如此,我們可能會問:

1苦難,到底是否上帝的玩意?上帝有意捉弄人!看著人受苦,然後伸手救人,或由得人自生自滅.

2苦難,到底是否上帝的無意或上帝有意的,有目的!

3苦難,到底是否上帝的愛意?上帝容許苦,甚至給人苦難,是出於對人的愛!使人經歷生命的成長,經歷生命的關係.

這個答案,我最後才回答!

3定義探討

中文裡「苦難」一詞,可字面解釋為「困苦與災難」.

在英文裡,相關的 字彙則有 “suffering”(受苦,源自拉丁文 sufferre,「在重壓下忍受」) , “distress”(憂苦),”agony”(痛楚)等等.所以,苦難通常只能是被動 的忍受,如同突遭重壓一般,往往令人感到萬分痛苦,喘不過氣,以致喪膽, 灰心,失望.而耶穌自願走向十字架,默然接受眾人置之於死的苦難,實非「 常人」行徑.

「苦難」可能是身體的,心靈的,情緒的,人際的,環境的.

「苦難」一詞,字面上解釋為「困苦與災難」.

在聖經中,「苦難」則有不同的定義,舊約希伯來文的用詞是「痛苦」或「疼痛」(ke’ev).而在新約希臘文用詞中是「苦楚」(Pathema) [1] {path’-ay-mah}

1) 指人所正遭受或曾受過的苦難

1a) 指外來的, 一個受苦的經歷, 不幸的事情, 慘禍, 災禍, 痛苦憂傷

1a1) 指基督所曾受的苦難

1a2) 亦指基督徒所必須受的苦難, 為著基督也曾同樣忍受過

1b) 指內心所經歷的憂傷, 衝擊, 犧牲

2) 也是「患難」(thlipsis)的意思.

2a) 壓制, 受苦, 迫害, 磨難

2b) 內心的愁煩, 苦惱

我發現人生一個苦會引發另一個苦,甚至更多更重的苦.一個人的苦,又會引致周圍的人受苦,並且一直傳遞下去,可以是沒完沒了的苦.

例如一位兄弟患病,會失去工作能力,因失業而脾氣暴躁,因常發脾氣暴躁引致家庭不和,全家受苦；一個苦引發一個又一個的苦,殃及池魚,屢見不鮮.

苦難難耐,到底上帝在那裡呢?為何上帝容許苦難呢?這是古今中外,人不斷會尋問,思想,掙扎,又好像無法得到答案.

苦難是人生的事實

上月,聽到王維琪得知政府不發免費牌照給香港電視時,概歎說:「…人生痛苦十常八九,只要活著,我仍會繼續砌下去!」苦難確是是人生的實況.

米蘭‧昆拉(Milan Kundera,1929年4月1日－)德曾說:「我痛苦,故我存在.」人的存在,就會經歷苦難.真樹與假樹之別.真樹(有生命)會生病(苦難),假樹(無生命)不會有蛀蟲.

人在世上是勞苦愁煩,病痛,衰老,死亡,在苦海的洶湧波濤中.

苦難好像空氣佈滿全世界,無論什麼地方都會有苦難,無論什麼人都會遭遇苦難,苦難從人生的開始,便與人形影不離；這正如雅各說:「我平生的年日又少,又苦.」(創世記47:9).

誰不怕苦難呢?或許會有；然而,大部份人面對苦難難免會懼怕,擔心,焦慮不安.

人生自嬰兒呱呱落地,一出母腹就會哭,因為他進人苦海般的世界.倘出生的孩子不會哭,醫生就會打他,令他受痛,放聲大哭.為什麼要如此對待初生嬰兒?因為他會哭,哭了才會呼吸,才會使肺開展,這樣嬰兒的生命才得以保存；所以醫生打出生不會哭的嬰兒,令他受苦,並不是虐待,乃是要救他的命.

所以,當孩子哭時,父母便笑了!

不過,當我個仔出生時,他一哭,好大聲,我嚇到手都震!上主給我們人生的苦難,也許有他一定的意義.

4 苦難像把刀

智者說:「苦難好像一把菜刀,可以拿來切菜,也可以拿來殺人,就看你怎樣拿法.」你是拿刀刃還是拿刀柄呢?

你若是拿刀刃就會被割傷；

拿刀柄,你就會控制它.

苦難來了,你只看苦難,不看苦難的背後的涵義.單單看苦難,就是抓刀刃,結果你被苦難所傷；面對苦難,你不看苦難,看苦難後面還有位上帝,這就是抓住刀柄,你就利用刀作事.

當一粒沙進到蚌的身體裡以後,蚌是很痛苦的,可是它不斷把一種內分泌途在傷口上,不久就出現一顆珍珠.我們的苦難也能化為祝福.

4.1苦難會催毀人

苦難,使人遠離神

某君:一次失戀,就不談戀愛

一次失戀,就不返教會

一次失戀,就認定失敗者

一次失戀,否定上帝!

苦難常常是人們抗議,起訴上帝,否定上帝甚或詛咒上帝的理由.因為“萬能”的上帝“應該”為人被遺棄在苦難之中負責.

存在主義者卡繆((Albert Camus, 1913~1960))曾說:「即使接受上帝存在,受苦的伊凡(托斯妥耶夫斯基著作「卡拉馬助夫兄弟們」中的主角之一)也不會在人類遭受的苦難面前向上帝屈服,反而在對人類遭受的苦難進行了長時間的思考以後,怒火越燒越旺,最後把「即使你存在」變成「你根本不值得存在」,甚至變成「你就是虛無」.」

因為苦難得不到說明,人類的受苦似乎就成為拒絕上帝的理所當然的理由.

4.2 從苦難中活出意義來

維克多‧ 弗蘭克(Viktor Emil Frankl 1905.03.26－1997.09.02)是意義治療法的創立者,他的理論已成為佛洛德,阿德勒之後維也納精神治療法的第三學派.第二次世界大戰期間,他曾被關進奧斯維辛集中營,受盡非人的折磨,九死一生,只是僥幸地活了下來,他在集中營之所以仍能活著,是因為有「活出意義來」的信念.

在他的著作《活出意義來》這本小書中,他特別強調意義的尋求,是人的最基本的需要.當這種需要找不到明確的指向時,人就會感到精神空虛,弗蘭克稱之為存在的空虛.這種情形普遍地存在於當今西方的「富裕社會」.

當這種需要有明確的指向卻不可能實現時,人就會有受挫之感,弗蘭克稱之為「存在的挫折」.這種情形發生在人生的各種逆境或困境之中.

在苦難中,人們一般是靠希望活著的,即相信或至少說服自己相信厄運終將過去,然後又能過一種有意義的生活.

因此,面對苦難重重的人生,人當如何破苦局呢?就是「從苦難中活出意義來!」

4.3 苦難具積極意義

上主給我這病定有其用意

曾在新聞節目,聽到一則,美國科舉家的新發明,傷健人士可以直接用思想控制的精密機械臂,當時找來一位因病而全身癱瘓多年的中年女士進行實驗,在其腦內植入感應器,讓她能自如地憑思想控制機械臂.

這位女士於成功通過實驗後,接受傳媒訪問,感慨地說出了這樣說:「多年以來,我一直認為上主給我這病定有其用意,現在我終於找到答案了,哈利路亞.」這位女士能夠堅強地活下去,就是因為在苦難中仍不失盼望,找到受苦的人生意義.

a苦難給人動力

以色列人在埃及若沒有受逼迫,若沒有遭受大的苦難,誰也不能把他們從埃及領出來,由於上帝藉著苦難催逼著他們,他們才願意離開埃及歸回迦南.

浪子若不遭受天災人禍,遭遇大饑荒；若不是窮苦起來,走投無路,浪子不會回家.

人在苦難中,發現自己的有限,生命的難以操控,因而轉向一位無限的上帝；或者到處尋求友朋的匡助,患難見真情,經歷人間的溫暖.

人的苦難中,千方百計要逃脫苦困,因而潛質盡顯,潛能大發；今日許多的發明,均由於人要逃脫苦難,勝過苦難而來.所以,苦難磨練人格,激發人的潛能,建立入神關係,促進社會科學文明進步.

請留意,人類吃了禁果雖然犯罪受苦,卻得到分辨善惡樹的智慧能力,眼睛因而明亮起來了.

舒伯特說:「我的音樂是才華和厄運之子.很奇怪地,眾人似乎最欣賞我在困厄之際所蘊釀的音樂.」

b苦難使人生命茁壯

正如一位哲人所說:「生命中的每個苦難,每個傷痛,每個打擊,都有它的意義.」

生命,總是在挫折和磨難中茁壯；

思想,總是在徘徊和失意中成熟；

意志,總是在殘酷和無情中堅強.

苦難是人認識社會,理解人生的生動教材,苦難是人成熟的機會,苦難是競爭社會中人面臨的必然挑戰.

著名成功學大師卡內基說:「苦難是人生最好的教育.」古今中外大量事實說明,偉大的人格無法在平庸中養成,只有經歷熔煉和磨難,願望才會激發,視野才會開闊,靈魂才會昇華,人生才會走向成功.

c苦難使人成長

主借著苦難催迫我們信他,也借著苦難叫我們的靈命長進.申命記32:11說:“又如鷹攪動巢窩,在雛鷹以上兩翅扇展,接取雛鷹,背在兩翼之上.”這是母鷹教雛鷹舉飛的方法,是攪動巢窩,使巢窩不舒服,叫雛鷹不得不離開,這樣雛鷹才會學飛.當雛鷹無力要跌下的時候,母鷹就扇展兩翅接取雛鷹,背在兩翼之上.上帝就是這樣對待我們,他攪動我們舒服的巢窩,把我們推出窩外去,使我們被迫使用我們的兩翼,就是信靠和順服,叫我們的靈命長進.

上帝一方面叫我受苦難,一方面保守我們.彼得前書5:10說:“那賜諸般恩典的上帝,曾在基督裡召你們,得享他永遠的榮耀,等你們暫受苦難之後,必要親自成全你們,堅固你們,賜力量給你們.”主借著苦難叫我們靈命長進,好像花草樹木借著太陽之曬而長進一樣.約伯遭遇苦難以後說:“我從前風聞有你,現在親眼看見你.”(約伯記42:5)苦難好像天使領我們與主愈近.詩人說:“我未受苦以先,走迷了路；現在卻遵守你的話.”“我受苦是與我有益,為要使我學習你的律例.”(詩篇119:67,71)“看哪,我受大苦,本為使我得平安.”(以賽亞書38:17)

d苦難使我們更深的認識神

從路21:5-19的苦難,如何面對?耶穌的指示如何?

1不要驚?

【路21:9】「你們聽見打仗和擾亂的事,不要驚惶；因為這些事必須先有,只是末期不能立時就到.』」

2不用準備怎樣做?

【路21:14】「所以你們當立定心意,不要豫先思想怎樣分訴；」當我們遇見逼迫和反對的時候,千萬不要憑自己說話.

3主必同在:

【路21:15】「因為我必賜你們口才智慧,是你們一切敵人所敵不住,駁不倒的.」

【路21:17】「你們要為我的名,被眾人恨惡.」

【路21:18】「然而你們連一根頭髮,也必不損壞.」

【路21:19】「你們常存忍耐,就必保全靈魂(或作必得生命).」

結語:人在苦難中,主必同在

對受苦的人來說,上帝不是從天上給他們安慰,而是安靜地參與在受苦者之中.對受苦者來說,最重要的是有沒有人瞭解他的痛苦,認同他的處境,與 他一起吶喊出心中的抗議.上帝沒有應許信祂的人會一帆風順,卻應許「與他 們同在」,這種「同在」就是最大的安慰了；每個人都需要安慰與鼓勵,上帝 永遠是扮演這角色.

上帝並沒有叫我們一定要解釋苦難的原因,其實祂自己也不常解釋苦難存在的原因；上帝只要我們與哀哭的人同哀哭,與喜樂的人同喜樂,要記得與受 苦的人同在,因為我們的愛白白得來,也要白白地分享給別人.解釋受苦的原 因未必能安慰受苦者,惟有認同,體會,接納才能使悲傷的人得著安慰!

記得羅曼羅蘭(Romain Rolland,1866年1月29日－1944年12月30日)有句名言:「只要看見一雙哭泣的眼睛,就值得我為生命而受苦.」

耶穌的另一個名字「以馬內利」是主對人類的承諾.他是與人同在的神.基督教認愛,神不只安慰受苦之人,更是參與在人們的苦難之中.對受苦的人來說,最重要的是有沒有人了解他的苦,認同他的處境.

作為上主的兒女,只要你知道主耶穌也曾「哭」,為已離世的拉撒路哭了(見約19:33),祂讓我們認識到道成肉身的主,來到我們中間,祂知道我們的苦,經歷我們的苦,並且能分擔我們的苦,並且勝過人間的苦,耶穌說:「…要叫你們在我裡面有平安,在世上你們有苦難,但你們可以放心,我已經勝了世界.」(約16:33)

莫特曼(Jürgen Moltmann,1926年－)強調人是透過苦難認識上帝的愛.人是從「苦難」和十字架來認識上帝,這樣的啟示是隱藏而間接的.莫特曼的苦難神學指出:在人類的痛苦中,上帝沒有隱藏,祂總是與受苦的人同在,就如聖經以賽亞書六十三章9結所說:「他們在一切苦難中,祂也同受苦難」.因此,沒有苦難就沒有愛!

「耶和華靠近傷心的人,拯救靈性痛悔的人.義人多有苦難,但耶和華救他脫離這一切.」詩篇第三十四篇18,19節

若然如此,我們可以說:

1苦難,不會是上帝的玩意,因為上帝是天父,看顧人的上的神.

2苦難,是上帝所容許苦難的,是上帝有意的,藉此要對人有某種意義!

3 苦難,是上帝的愛意?上帝容許苦,甚至給人苦難,是出於對人的愛!使人經歷生命的成長,經歷生命的關係.

所以上帝容許苦難,不是祂的玩意,而是衪的愛意,使人更深體會祂的真愛.

天地是有缺陷

【】台灣著名作家張曉風女士:『天地是有缺陷的,但缺陷造成了縐折,縐折造成了奇峰幽谷之美,月亮是不是常圓的,人生不如意事十常八九.當我們心平意和地承認這一切缺陷的時候,我們忽然發現沒有甚麼是不可接受的.…』

她又說:「有時,人生的殘缺不但是可以接受,而且是美麗的.」

e從苦難中,可看到一位吊詭的上帝:

上帝會賜福,上帝也會賜禍,給人苦難

上帝可以理解,上帝又不可以理解

苦難是公平的:人人有份

苦難是不公平:好人受苦,壞人亨福,大,小苦不均

苦難是好的:建立人；苦難不好的:催毀人

苦難是痛苦,苦難於雨過天晴是喜樂

上帝設立分別善惡樹是不好,會死!

上帝設立分別善惡樹是好的,會得智慧,分辨對錯!

■   在她所著的雨後彩虹一書中,她如此說:『五年前,我的丈福音腦癌去世,自己被迫踏上單親之路,與唯一女兒相依而活.一年前,女兒亦突感染病病毒引至腦炎及急性腎衰竭,在不足二天之內離世了.我的經歷,很多人都認為神對我實在太不公平和殘酷,甚至--,一些長輩更認為我是命苦的人.』

『但我亦想,不幸的事既然已經發生了,苦難已經臨到我身上,是否我要在哀傷自憐或埋怨中渡日?--』

她說:『我是個有信仰的人,我知道自己的生命是主救贖過來的；昔日約的伯受苦,神的兒子何嘗不是要受苦甚至受死.我受的苦,我很相信神必有祂的旨意,因此我不問神為何我要受苦,反之是不斷求問神,究竟要我在苦難學習甚麼的功課,我深信主會加力於我,使我行過這流淚谷.』[見《雨後彩虹》,頁132-133]

梅太勇敢地站起來,在贐明會服事關心一些極度傷心的人.她的生命成為他人的祝福.

上帝僕人說:「基督教信仰不是為解釋苦難,而是為承載苦難.」

經歷傷痛的人,才了接受傷人的苦,才懂得安慰受傷的人.

我遇過家人的離世:叔叔的離世,兄長的離世,同工牧者的離世.主持安息禮拜時,更能體察教友的哀傷.我試過考不進大學,體會過年青人讀書失敗的滋味,失戀,失意.但現在總算苦難不多,現在也一路順風,生活順利如意,你便要感恩,為到上帝給你的生命也感謝.

無論如何,在我的體驗中,苦難是滿有上帝的愛意!



\newpage

\section{天堂與地獄}
\label{sec:228}
\textbf{鄧瑞強博士 - 天堂與地獄}
\newline
\newline
\hyperref[sec:227]{< < < PREV SERMON < < <}
~
\hyperref[sec:toc]{[返目錄]}
~
\hyperref[sec:229]{> > > NEXT SERMON > > >}


路加福音 23:33-43
\newline
\begin{longtable}{cl}
\hline
\hline
章節 & 經文 (和合本修訂版)\\
\hline
23:33 & \begin{tabularx}{0.7\textwidth}{X} 到了一個地方，名叫髑髏地，他們就在那裡把耶穌釘在十字架上，又釘了兩個犯人：一個在右邊，一個在左邊。 \end{tabularx} \\ \\ \relax
23:34 & \begin{tabularx}{0.7\textwidth}{X} 〔 這時，耶穌說：「父啊！赦免他們，因為他們所做的，他們不知道。」〕士兵就抽籤分他的衣服。 \end{tabularx} \\ \\ \relax
23:35 & \begin{tabularx}{0.7\textwidth}{X} 百姓站在那裡觀看。官長也嘲笑他，說：「他救了別人，他若是基督，是神所揀選的，救救他自己吧！」 \end{tabularx} \\ \\ \relax
23:36 & \begin{tabularx}{0.7\textwidth}{X} 士兵也戲弄他，上前拿醋送給他喝， \end{tabularx} \\ \\ \relax
23:37 & \begin{tabularx}{0.7\textwidth}{X} 說：「你若是猶太人的王，救救你自己吧！」 \end{tabularx} \\ \\ \relax
23:38 & \begin{tabularx}{0.7\textwidth}{X} 在耶穌上方有一個牌子寫著：「這是猶太人的王。」 \end{tabularx} \\ \\ \relax
23:39 & \begin{tabularx}{0.7\textwidth}{X} 同釘的犯人中有一個譏笑他，說：「你不是基督嗎？救救你自己和我們吧！」 \end{tabularx} \\ \\ \relax
23:40 & \begin{tabularx}{0.7\textwidth}{X} 另一個就應聲責備他，說：「你是一樣受刑的，還不怕神嗎？ \end{tabularx} \\ \\ \relax
23:41 & \begin{tabularx}{0.7\textwidth}{X} 我們是應得的，因為我們是自作自受，但這個人沒有做過一件不對的事。」 \end{tabularx} \\ \\ \relax
23:42 & \begin{tabularx}{0.7\textwidth}{X} 他對耶穌說：「耶穌啊，你進入你國的時候，求你記念我。」 \end{tabularx} \\ \\ \relax
23:43 & \begin{tabularx}{0.7\textwidth}{X} 耶穌對他說：「我實在告訴你，今日你要同我在樂園裡了。」 \end{tabularx} \\ \\
[1ex]
\hline
\hline
\end{longtable}
講題:天堂與地獄　Heaven and Hell

經文:路加福音23:33-43

講員:鄧瑞強博士

場合:香港中文大學崇基學院禮拜堂

日期:2013年11月24日

各位弟兄姊妹,早安.

大家有否曾經活在地獄裡?

地獄是苦難和絕望的深淵.

在歷史上,人類製造了很多地獄.

在中國近代史上,有一個時期,人為了自保,會出賣家人,愛人,朋友.將政治的罪名加給對方,使自己免於這個罪名.為了免於這個罪名,一個人便無可避免地跌入罪的網羅裡.一個人為了自保,將別人釘上十字架,他的生命很難不被十架的血濺污.這是人間地獄.

人間地獄的恐怖,莫如納粹德國的死亡集中營.集中營的恐怖,又以「奧斯維辛」(Auschwitz)為最.進了這個集中營,要出去的話,只能通個焚化爐的煙囪,化作一縷輕煙,飄送出去.

很多電影講「奧斯維辛」的經驗,其中一齣電影叫《蘇菲的抉擇》(Sophie’s Choice, 1982).主角蘇菲是一個波蘭人,本身不是一個反納粹的人,並千方百計避免捲入反納粹的事情裡,這樣的一個無辜的人,仍不能倖免,被送入「奧斯維辛」.為了自保,她出賣道德原則,出賣肉體.在保留兒子的生命,抑或保留女兒的生命之間,她被迫作出抉擇,作出著名的「蘇菲的抉擇」.她選擇放棄女兒,即送女兒去死.一個無辜的人,作出無可奈何的抉擇,使自己不再無辜.這種深深的負罪感,使她覺得神放棄了她,而她也放棄神.這種深深的負罪感,使她不能選擇日後的幸福,而只能選擇與一個常常虐待她的人走在一起,最終走向死亡.一個無辜的人無法再是無辜,一個幸福的人無法進入幸福,一個信神的人無法再相信神.這是人間地獄.

哲學家阿多諾(Theodor Adorno)說:「奧斯維辛之後,再沒有詩.(No poetry after Auschwitz)」

今日講的道,是講神走在人間,進入人間地獄,慘死在十字架上.我們大概可以說:「神死在十架上之後,再沒有語言.」

不過,為了講道,我仍需要講一點東西,言說那不能言說的.

今日的講道經文,是今日的福音書經課:路加福音23:33-43.

23:33　到了一個地方,名叫「髑髏地」,就在那裡把耶穌釘在十字架上,又釘了兩個犯人:一個在左邊,一個在右邊.

23:34　當下耶穌說:「父啊!赦免他們；因為他們所做的,他們不曉得.」兵丁就拈鬮分他的衣服.

23:35　百姓站在那裡觀看.官府也嗤笑他,說:「他救了別人；他若是基督,神所揀選的,可以救自己吧!」

23:36　兵丁也戲弄他,上前拿醋送給他喝,

23:37　說:「你若是猶太人的王,可以救自己吧!」

23:38　在耶穌以上有一個牌子〔有古卷加:用希臘,羅馬,希伯來的文字〕寫著:「這是猶太人的王.」

23:39　那同釘的兩個犯人有一個譏笑他,說:「你不是基督嗎?可以救自己和我們吧!」

23:40　那一個就應聲責備他,說:「你既是一樣受刑的,還不怕神嗎?

23:41　我們是應該的,因我們所受的與我們所做的相稱,但這個人沒有作過一件不好的事.」

23:42　就說:「耶穌啊,你得國降臨的時候,求你記念我!」

23:43　耶穌對他說:「我實在告訴你,今日你要同我在樂園裡了.」

事情發生的地點,叫「髑髏地」.「髑髏」,是一個「死亡」的標誌.在這裡,耶穌被釘在十字架上,他的左邊右邊,釘有兩個犯人.耶穌,一個暴力下的無辜者.宗教的暴力,政治的暴力,群眾的暴力,集合成瘋狂的力量,將這無辜的人釘在十架上.宗教提供了這暴力所需要的神聖理由,政治提供了這暴力所需要的客觀條件,群眾的助力而不是抵抗,令這暴力的發生變得出奇的自然.在「髑髏地」發生的這件事件,不是一樁罪犯行刑的一般事件,不是一樁犯罪的人得到應有的公義裁決的一般事件,這是一樁無論你有罪抑或無辜,都難逃劫數的暴力事件.一個人,無論你有罪,抑或無辜,都在十架上受煎熬,等待死亡.這是人間地獄.

這不是特例,這不是孤立事件,這是人類存在的,制造地獄的邪惡結構的一次展露.所有人都牽涉其中.有些人是處心積慮造成這暴力事件,有些人是沉默地讓這事發生.「有罪的人」被牽涉其中,「無辜的人」也不能倖免.這件事之後,不能沒有後遺症.不純粹是因果的後遺症,而是存在論上(ontological)的後遺症.人間地獄之後,人還能相信上帝嗎?否定上帝,抑或仍能相信上帝?拒絕上帝,抑或仍能將生命交在上帝手裡?堅持絕望,抑或仍能保持希望?

上帝是否死在人類的十字架上?若果是,這意味什麼?若果上帝死了,人類還有希望嗎?這真是一個大問題.

路加福音的經文說:

23:35　百姓站在那裡觀看.官府也嗤笑他,說:「他救了別人；他若是基督,神所揀選的,可以救自己吧!」

23:36　兵丁也戲弄他,上前拿醋送給他喝,

23:37　說:「你若是猶太人的王,可以救自己吧!」

23:38　在耶穌以上有一個牌子寫著:「這是猶太人的王.」

23:39　那同釘的兩個犯人有一個譏笑他,說:「你不是基督嗎?可以救自己和我們吧!」

這些問題和呼喊,可以翻譯成人類在人間地獄裡的存在叩問.

神啊,人間地獄的存在,是否顯示你不存在?神啊,若要人相信你,你得做點什麼.你做點什麼,救救自己吧!

神啊,為何在人間的災難中,你顯得那麼無能?神啊,你做點什麼,救救自己吧!

神啊,若你存在,拯救救這個地獄一般的世界吧!

耶穌沉默無語.上帝沉默無語.

天地間,回盪著人類的「天問」.

在《蘇菲的抉擇》那齣戲裡,蘇菲在地獄一般的集中營中,放棄了信仰.當她離開集中營後,她在教堂裡試圖自殺,她要死在神的面前,以死去控訴神.地獄的生活,令人絕望.絕望的人,以絕望的手法去表達絕望.若上帝存在,她就要用絕望去否定這上帝.地獄一般的生活,黑得令人性的一點光明都沾滿黑暗.這是徹底的絕望.

中世紀的詩人但丁(Dante),在其《神曲》(Divine Comedy)的＜地獄篇＞(Inferno)裡說,地獄的門口寫著:「所有進去的人,要放棄一切希望.」在那條分隔生者和死者的「亞開龍河」(River Acheron)河邊,那些準備進入地獄的靈魂,不斷發出咒詛.他們咒詛上帝,咒詛父母,咒詛人類,咒詛自己的出生.地獄與絕望與咒詛,連成一體.地獄一樣的生活,令人放棄希望.但只有那些真正全然放棄希望的人,才能真正進入地獄.絕望的人,心中充滿咒詛.若果只是滿懷咒詛,就只會陷入更深的絕望中.

在十架的地獄裡,上帝在哪裡?

耶穌說:「父啊!赦免他們；因為他們所做的,他們不曉得.」(23:34)

這句話透露了反抗地獄力量的另類力量,神是這力量的根源.

我知道,更多人的禱告是:「父啊!不要赦免他們；因為他們所做的,他們曉得.」若果這禱告的背後,是埋怨與仇恨,則已是地獄的種子.一個地獄過後,另一個地獄便會出現.

「父啊!赦免他們.」這裡牽涉一種人性的變革.沒有這種人性的變革,地獄便會層出不窮地出現.

我知道,單講一種表面的寬恕,一種廉價的寬恕,可能會使人變成沉默的群眾,助紂為虐,成為製造地獄的幫凶.

我知道,基督教信仰不只講內心的變革,也要求種種制度的變革,讓天國降臨人間.

我只是想講,一切制度的變革,一切反抗邪惡的行動,若背後沒有一種人性的變革,可能都只是在「地獄的超穩定結構」裡周而復始地運轉.

主耶穌不會不明白制度改變的重要性,主耶穌不會不明白沉默可能變成邪惡的幫凶.主耶穌卻直指制度改變所需的人性基礎,直指反抗地獄所需的精神力量.在現實上,我們有時需要用一把刀去阻擋另一把刀,但是,若我們不相信,不渴望,不夢想一個不用刀去威脅人的世界,這個世界就不可能出現.

上帝在哪裡?

當耶穌向我們展示一個新世界,一個神聖的世界,並透過自己在受苦中的寬恕去展示這個新世界實現的可能時,上帝反抗地獄的心,便向我們展示出來.上帝就在這裡.

面對神在人間地獄裡的沉默,我們可以如何回應?

耶穌十架旁的另一個犯人,作了一個示範.

23:40　那一個〔犯人〕就應聲責備他〔剛才譏諷耶穌的犯人〕,說:「你既是一樣受刑的,還不怕神嗎?

23:41　我們是應該的,因我們所受的與我們所做的相稱,但這個人沒有作過一件不好的事.」

23:42　就說:「耶穌啊,你得國降臨的時候,求你記念我!」

這個犯人不因為自己身陷地獄一樣的苦難中,而反抗上帝.他視這為上帝對人類的罪的一種揭露,一種審判.他承認自己的罪.當塵世出現人間地獄時,有誰還能說自己是無罪的呢?

這個犯人肯定耶穌是無辜的,「這個人沒有作過一件不好的事」,他看見耶穌的良善.在十架的黑暗中,這個犯人在十架的耶穌那裡發現光明.

這個犯人作出一個突破性的祈求.他說:「耶穌啊,你得國降臨的時候,求你記念我!」他在耶穌那裡看到一個新世界,新國度.他求耶穌記念他.在猶太文化裡,一個將死的人,會求神記念他,這就是求神將他列在義人的行列中,能永活在神的面前.世上的偉人,很多人為他們立碑立傳,就是記念他們,在記念中讓他們進入永恆.這個犯人求耶穌記念他,像神記念他一樣.他相信,在這個地獄一樣的世界之外,有一個天堂一樣的世界.他相信,在這個毫無道理的悲慘世界背後,仍有一個正義的世界,他渴望自己能像義人一樣,活在其中.他相信,上帝在人間地獄的沉默,是要讓人類看見自己的罪惡,讓人在沉默的耶穌那裡看見愛的可能性,激發人渴望一個來臨中的國度,呼喚人為這國度奉獻自己,好叫自己成為神記念的義人.耶穌作為上帝,死在十架上.這個犯人在這即將死亡的上帝那裡,看見人類的希望.

這個犯人的祈求,好像呼喚我們,不要在人間地獄中怨恨,咒罵,絕望.要信仰上帝,要承認自己的邪惡,要渴望自己有分於那即將來臨的天國.

23:43　耶穌對他說:「我實在告訴你,今日你要同我在樂園裡了.」

在人間地獄裡,耶穌像火把,將黑暗驅散,讓人看到光明.

耶穌堅定地回應這個犯人的渴求.他的渴望會被實現,並且不用等到遙遠的未來,而是近在咫尺,今日,當下.若果人類的歷史,是被逐出樂園的歷史.在這裡,人當下找到回歸樂園的大門.

在十架上,耶穌沒有神蹟,上帝沒有作為,但在人的絕望深處,燃點起滿載恩典的希望.上帝就在這裡,向人顯現.基督教信仰,就是在十架上死亡的上帝那裡,捉摸到上帝的存在.

在十架上,有耶穌,有受苦的犯人在耶穌左右.受苦的你,你的反應會像哪個犯人?

但願:榮耀歸於聖父,聖子,聖靈.阿們.



\newpage

\section{眾裡尋他千百度}
\label{sec:229}
\textbf{鄧瑞強博士 - 眾裡尋他千百度}
\newline
\newline
\hyperref[sec:228]{< < < PREV SERMON < < <}
~
\hyperref[sec:toc]{[返目錄]}
~
\hyperref[sec:230]{> > > NEXT SERMON > > >}


馬太福音 11:2-11
\newline
\begin{longtable}{cl}
\hline
\hline
章節 & 經文 (和合本修訂版)\\
\hline
11:2 & \begin{tabularx}{0.7\textwidth}{X} 約翰在監獄裡聽見基督所做的事，就派他的門徒去， \end{tabularx} \\ \\ \relax
11:3 & \begin{tabularx}{0.7\textwidth}{X} 問耶穌：「將要來的那位就是你嗎？還是我們要等候另一位呢？」 \end{tabularx} \\ \\ \relax
11:4 & \begin{tabularx}{0.7\textwidth}{X} 耶穌回答他們：「你們去，把所聽見、所看見的告訴約翰： \end{tabularx} \\ \\ \relax
11:5 & \begin{tabularx}{0.7\textwidth}{X} 就是盲人看見，瘸子行走，痲瘋病人得潔淨，聾子聽見，死人復活，窮人聽到福音。 \end{tabularx} \\ \\ \relax
11:6 & \begin{tabularx}{0.7\textwidth}{X} 凡不因我跌倒的有福了！」 \end{tabularx} \\ \\ \relax
11:7 & \begin{tabularx}{0.7\textwidth}{X} 他們一走，耶穌就對眾人談到約翰，說：「你們從前到曠野去，是要看甚麼呢？看風吹動的蘆葦嗎？ \end{tabularx} \\ \\ \relax
11:8 & \begin{tabularx}{0.7\textwidth}{X} 你們出去到底是要看甚麼？看穿細軟衣服的人嗎？那穿細軟衣服的人是在王宮裡。 \end{tabularx} \\ \\ \relax
11:9 & \begin{tabularx}{0.7\textwidth}{X} 你們出去究竟是要看甚麼？是先知嗎？是的，我告訴你們，他比先知大多了。 \end{tabularx} \\ \\ \relax
11:10 & \begin{tabularx}{0.7\textwidth}{X} 這個人就是經上所說的：『看哪，我要差遣我的使者在你面前，他要在你前面為你預備道路。』 \end{tabularx} \\ \\ \relax
11:11 & \begin{tabularx}{0.7\textwidth}{X} 我實在告訴你們，凡女子所生的，沒有一個比施洗約翰大；但在天國裡，最小的比他還大。 \end{tabularx} \\ \\
[1ex]
\hline
\hline
\end{longtable}
講題:眾裡尋他千百度　Here Comes the Chosen One

經文:馬太福音11:2-11

講員:鄧瑞強博士

場合:香港中文大學崇基學院禮拜堂

日期:2013年12月15日

各位弟兄姊妹,早安.

現在是「將臨期」,教會在這段期間,記念主耶穌的降生,也盼望主耶穌的再來.

主耶穌曾經來世,世人拒絕他.

今日,主耶穌若以某種形式來到我們中間,我想,世人也會再次拒絕他.

俄國大文豪陀思妥耶夫斯基(Dostoyevsky)寫了一本巨著:《卡拉馬佐夫兄弟》(The Brothers Karamazov),其中有一章成了宗教文學的經典,這一章稱為「宗教大法官」(The Grand Inquisitor).話說,在十六世紀的某天,耶穌再次來到人間.在教堂前,他叫一個女孩子復活.那時,一個紅衣主教見到,這個紅衣主教是宗教裁判所的宗教大法官.他下令將耶穌捉住.

在宗教法庭的監獄裡,宗教大法官對耶穌說(經過改寫,大意如下):「真是你?真是你麼?…你為什麼到這裡來妨礙我們?…為了人性的自由,你拒絕將石頭變麵包,但是一般人只是平凡的人,不理解自由.對他們而言,再沒有比自由更難以忍受的東西了.他們寧願要麵包.他們會對我們說:『你們儘管奴役我們吧,只要給我們食物吃.』…耶穌,你答應給他們天上的糧,但是我再重複一次,在軟弱而永遠敗德不義的人類眼裡,還是地上的麵包好得多.…我們給他們麵包,他們將自由給我們.…他們對我們這些教士驚歎,把我們看作神,因為我們作為他們的領袖,竟甘願把他們所懼怕的自由承擔下來而統治他們,他們覺得放棄麵包而做自由人太可怕了.…耶穌,你知道嗎?倘若人類真的有了自由,他們也要找尋一個敬拜的對象.他們會互相廝殺,迫令對方敬拜自己的神.最後,什麼能令所有人放下爭鬥,成為眾人無可爭辯的敬拜對象呢?是麵包.但是,耶穌,你為了人的自由,卻拒絕了麵包.你應許的自由,讓平凡的人煩惱不堪.…為了這些平凡的人,我們建築了宗教大廈,以權威領導他們,讓他們交出自由,解除他們的人性負擔.…沒有了自由,他們不會煩惱,也不會自相殘殺.…耶穌,你說,這有什麼不對?…耶穌,你走吧!不要再來了!永遠不要再來了.」

於是,耶穌走了.

若耶穌今日再來,他遇到的排斥可能更大了.

因為現代人不單要麵包,也要自由.但是,我們不要耶穌講的那種自由,耶穌講的那種自由,是放下自己的自由而甘願去愛他人的自由.我們不要這種自由,我們要那種「想如何便如何」的自由.我們不要那種「受愛束綁」的自由,我們要那種「逼令別人愛我們」的自由.

由於我們又要麵包,又要這種自由,我們有雙重理由拒絕耶穌.

「將臨期」不單記念耶穌的降生,不單盼望耶穌的再來,也是反省的時刻,反省我們對耶穌的拒絕.

「眾裡尋他千百度」,我們尋尋覓覓,要尋找耶穌,但我們卻又用雙手,蒙著眼睛,讓我們看不見他.

今日的講道經文,是今日的福音書經課:馬太福音11: 2-11.

在這福音故事裡,一個重要人物問了一個重要問題:「耶穌,是你嗎?」

我們看看這段經文(馬太福音11: 2-11):

11:2　約翰在監裡聽見基督所做的事,就打發兩個門徒去,

11:3　問他說:「那將要來的是你嗎?還是我們等候別人呢?」

11:4　耶穌回答說:「你們去,把所聽見,所看見的事告訴約翰.

11:5　就是瞎子看見,瘸子行走,長大痲瘋的潔淨,聾子聽見,死人復活,窮人有福音傳給他們.

11:6　凡不因我跌倒的就有福了!」

11:7　他們走的時候,耶穌就對眾人講論約翰說:「你們從前出到曠野是要看什麼呢?要看風吹動的蘆葦嗎?

11:8　你們出去到底是要看什麼?要看穿細軟衣服的人嗎?那穿細軟衣服的人是在王宮裡.

11:9　你們出去究竟是為什麼?是要看先知嗎?我告訴你們,是的,他比先知大多了.

11:10　經上記著說:『我要差遣我的使者在你前面預備道路』,所說的就是這個人.

11:11　我實在告訴你們,凡婦人所生的,沒有一個興起來大過施洗約翰的；然而天國裡最小的比他還大.」

這個福音故事,是以一個人心中的疑惑作開始的.這個人叫約翰,由於他在約旦河為人施洗,所以又叫「施洗約翰」.施洗約翰是一個又敬虔,又正直的人物.他指斥希律王的錯,被希律王關在牢裡.他時時渴望上天差派救主來這世間,打救他的同胞,糾正世間的不義.在他心中,耶穌應是這人物.

施洗約翰曾經公開向人說,他只是用水為人施洗,但耶穌的能力比他更大,耶穌要用聖靈與火為人施洗.耶穌手裡拿著簸箕,處理田上的麥粒,把麥子收在倉裡,把無用的糠用不滅的火燒盡.這些講法,帶有末日審判的意味.按施洗約翰的估計,耶穌會如烈火一般,大力將正義的天國帶來,義人被肯定,惡人被消滅.耶穌將會有一番大作為.

但是,耶穌的行為,與施洗約翰心中的理想,看來有一段距離.

施洗約翰是一個克苦嚴肅的人,常常禁食禱告,用現在的講法,不抽煙不飲酒,不會「夜蒲」,不會吃喝玩樂.他呼喚人淨化心靈,悔改歸神.耶穌的行為卻又如何?按馬太福音11:19 所講的:「〔耶穌〕也吃也暍,人又說他是貪食好酒的人,是稅吏和罪人的朋友.」施洗約翰得到的消息,耶穌就是這樣的人.

若有人說你的老公或你的男朋友,是「又吃又暍,是貪食好酒的人,是黑道人士和三教九流的人的朋友,還時時和妓女在一起」.你驚不驚?擔心不擔心?

在施洗約翰的心中,耶穌似花花公子多一點.

於是,施洗約翰派兩個門徒去問耶穌:「是你嗎?我等的人是你嗎?我期待的救主是你嗎?」

當我們對神的期待,與神的表現,有一段距離時,這也是我們的問題:「是你嗎?你是真神嗎?抑或,我要投靠另一個?」

當神不靈驗時,當我們要求麵包而神給我們自由時,當我們要求成功而神給我們十字架時,這時,我們會問:「是你嗎?是你嗎?抑或,我搞錯了.」

耶穌沒有直接回答這問題,他講出很多事實.

11:4　耶穌回答說:「你們去,把所聽見,所看見的事告訴約翰.

11:5　就是瞎子看見,瘸子行走,長大痲瘋的潔淨,聾子聽見,死人復活,窮人有福音傳給他們.

11:6　凡不因我跌倒的就有福了!」

耶穌所展示的大能,與施洗約翰期望的「烈火」一樣的大能,有很大落差.

耶穌沒有差遣天軍來消滅羅馬士兵,耶穌沒有興起群眾來推翻殘暴的希律王,耶穌也沒有叫火從天降下來淨化這個不義的世界.耶穌卻是與世上最弱的弱者在一起,耶穌扶起他們,耶穌在他們生命裡重燃希望.耶穌開展了一場人性的革命.

盲的,破的,患重病的,聾的,全是社會上的弱者,他們需要幫助,但他們更需要一個尊重和接納他們的社會.死了的人,不可能再有希望.耶穌復活他們,就是叫最無希望的人再有希望.死人復活,應是最大的事,但在死人復活之後,還有一項「窮人有福音傳給他們」.看來,「窮人有福音傳給他們」應是以上事情的一個扼要說明.

「窮人」,是「彎腰乞討」的人,他們無法挺起胸膛做人.但是,在耶穌裡,這些人能聽到有福的信息,能體會到幸福.這是一種翻天覆地的改變.

我讀過一本書,講到一個啞的人,在夢中遇見耶穌,耶穌用流利的手語和她交談.(參考:侯活士,范尼雲:《暴力世界中的溫柔》)耶穌沒有醫治她,但按她的本相接納她.她在耶穌面前,不用擔心耶穌不明白她,不用擔心耶穌不接納她,當然更不用擔心耶穌會歧視她.在耶穌面前,她安然地活出自己,安然地表達自己,安然地愛與被愛.這個「窮人」,能體會到幸福.她不需別人施捨,她不用擁有什麼,她卻能有尊嚴地活著,單單這樣,她已感到無比幸福.因著耶穌,窮人能體會何謂「有福」的生命.

耶穌做的,和施洗約翰期望的,有一段距離.我不知道,施洗約翰聽完這番話後,能否明白耶穌就是他等待的人.

我不知道,我們聽完耶穌這番話後,能否明白耶穌就是我們等待的人.

耶穌最後加上一句話,非常耐人尋味.

11:6　「凡不因我跌倒的就有福了!」

看來,在耶穌心中,很多人會因他所作的跌倒；或者說,因他所作的事,反對他.

畢竟,無論耶穌何時來到人間,世人多是拒絕他的.

試想想,聖誕節的故事,是講什麼的?

是講天上的神,降生在馬槽裡.單單這一點,我們都已排斥他了.在香港,我們講的是「贏在起跑線」上；但耶穌做的,是選擇「輸在起跑線」上.關於這點,連東方的博士都搞錯.他們以為耶穌會生在皇宮裡,他們去找希律王.耶穌卻是生在馬槽裡.「贏在起跑線」上的心態,會扭曲一切.遊戲的童年充塞著各種大人的指標.同伴變成競爭的對手,使人失去友誼.人生的成就只在乎勝負,使人失去永恆的價值追求.得失的計量,使人物化,失去人性的尊嚴.耶穌生在馬槽裡,「輸在起跑線」上.耶穌明白,俯就卑微,比強者更堅強；彎腰為人洗腳,比偉人更偉大；親吻窮人,比一切富人更富有.

耶穌「輸在起跑線」上,因他願意成為弱者,與弱者同行,他珍惜每個人的人性尊嚴.他渴望,每個人都能活出神聖的光明.「凡不因耶穌的作為跌倒的,就有福了!」但是,「贏在起跑線」上的世界,會排斥「輸在起跑線」上的生命.因耶穌的作為,排斥耶穌的,畢竟佔多數.

主耶穌講完了他做的事後,他反問身邊的人一個問題,也反問我們.

11:7　他們走的時候,耶穌就對眾人講論約翰說:「你們從前出到曠野是要看什麼呢?要看風吹動的蘆葦嗎?

11:8　你們出去到底是要看什麼?要看穿細軟衣服的人嗎?那穿細軟衣服的人是在王宮裡.

11:9　你們出去究竟是為什麼?是要看先知嗎?我告訴你們,是的,他比先知大多了.

11:10　經上記著說:『我要差遣我的使者在你前面預備道路』,所說的就是這個人.」

耶穌問他們:你們去曠野,想找什麼?

若在今日,耶穌會問我們:「你們返教會,到底想尋找什麼?到底想得到什麼?」

我們不要只問耶穌:「耶穌,是你嗎?」也要反過來問自己:「我們找什麼?我們是找耶穌嗎?我們是找那個我們塑造出來的耶穌嗎?我們是找那個能幫我們『贏在起跑線』上的耶穌嗎?抑或,我們是找那個甘願『輸在起跑線』上的耶穌呢?」我們在這裡,到底要找什麼?凡不因耶穌所做的而排斥耶穌的人,有福了.

當時的人,到曠野去,是要尋索先知,是要聆聽神的要求.

耶穌說,施洗約翰的確是先知,並且是一個偉大的先知.但是,這仍不是人終極要尋求的.

耶穌繼續說:

11:11　「我實在告訴你們,凡婦人所生的,沒有一個興起來大過施洗約翰的；然而天國裡最小的比他還大.」

施洗約翰是很偉大的,但畢竟仍只是指向耶穌的路標.一個餐牌標示出各種食物的名字,甚至寫出各款sung 菜如何美味,仍只是餐牌,不是那些美味的食物本身.一個指示我們去遊樂場的路標,無論設計得多精美,也只是路標,不是遊樂場本身.在曠野的呼聲,無論多嘹亮,無論對天國發表了多偉大的觀點,仍只是呼聲,不是天國本身.一個正在吃美味食物的乞丐,強過那些只能看餐牌的富豪.一個在遊樂場玩耍的小孩,強過那些只是看路標的人.一個在天國裡的人,當然強過那些只是聆聽先知,那怕是很用心地聆聽先知的人.

有一個故事是這樣的:

很久很久以前,有一個廚師,能用平凡的材料,煮出各樣人間美食.村民們很崇拜他,常常圍住他,問他各種關於煮食的學問.後來,這個廚師死了.村民們為他造了一間廟,尊他為「廚神」.為了要讓廚神講過的東西永留萬世,便按他生前的講論寫了很多食譜.為了精確解釋這些食譜的內容,發展出不同的學派.不同學派之間,為了爭奪解釋的正統性,互相廝殺,很多人為此死了.多年後,這「廚神」的食譜已發展成精細的學問.求這些學問的人,絡繹不絕.但,奇怪的是,從來沒有人好奇地問:按「廚神」食譜煮出來的菜,到底是什麼味道的呢?

今日,我們對耶穌有很多談論,滿足於這些談論的人多,真正與耶穌有生命連繫的人少.我們對耶穌講的天國有很多描述,但願意按耶穌的指示進入天國的人少.我們尋求神,擁有關於神的很多知識,但能順服神,敬拜神的人少.

主耶穌說:「凡婦人所生的,沒有一個興起來大過施洗約翰的；然而,天國裡最小的,比他還大.」

誠心渴望你是一個在天國裡的人.

但願:榮耀歸於聖父,聖子,聖靈.阿們.



\newpage

\section{生命因何活得更美好}
\label{sec:230}
\textbf{鄭漢文博士 - 生命因何活得更美好?}
\newline
\newline
\hyperref[sec:229]{< < < PREV SERMON < < <}
~
\hyperref[sec:toc]{[返目錄]}
~
\hyperref[sec:231]{> > > NEXT SERMON > > >}


馬太福音 2:13-23
\newline
\begin{longtable}{cl}
\hline
\hline
章節 & 經文 (和合本修訂版)\\
\hline
2:13 & \begin{tabularx}{0.7\textwidth}{X} 他們走後，忽然主的使者在約瑟夢中向他顯現，說：「起來！帶著小孩子和他母親逃往埃及，住在那裡，等我的指示；因為希律要搜尋那小孩子來殺害他。」 \end{tabularx} \\ \\ \relax
2:14 & \begin{tabularx}{0.7\textwidth}{X} 約瑟就起來，連夜帶著小孩子和他母親往埃及去， \end{tabularx} \\ \\ \relax
2:15 & \begin{tabularx}{0.7\textwidth}{X} 住在那裡，直到希律死了。這是要應驗主藉先知所說的話：「我從埃及召我的兒子出來。」 \end{tabularx} \\ \\ \relax
2:16 & \begin{tabularx}{0.7\textwidth}{X} 希律見自己被博學之士愚弄，極其憤怒，差人將伯利恆城裡和四境所有的男孩，根據他向博學之士仔細查問到的時間，凡兩歲以內的，都殺盡了。 \end{tabularx} \\ \\ \relax
2:17 & \begin{tabularx}{0.7\textwidth}{X} 這就應驗了耶利米先知所說的話： \end{tabularx} \\ \\ \relax
2:18 & \begin{tabularx}{0.7\textwidth}{X} 「在拉瑪聽見號咷大哭的聲音，是拉結哭她兒女；她不肯受安慰，因為他們都不在了。」 \end{tabularx} \\ \\ \relax
2:19 & \begin{tabularx}{0.7\textwidth}{X} 希律死了以後，在埃及，忽然主的使者在約瑟夢中向他顯現， \end{tabularx} \\ \\ \relax
2:20 & \begin{tabularx}{0.7\textwidth}{X} 說：「起來，帶著小孩子和他母親回以色列地去！因為要殺害這小孩子的人已經死了。」 \end{tabularx} \\ \\ \relax
2:21 & \begin{tabularx}{0.7\textwidth}{X} 約瑟就起來，帶著小孩子和他母親進入以色列地去。 \end{tabularx} \\ \\ \relax
2:22 & \begin{tabularx}{0.7\textwidth}{X} 但是他因聽見亞基老繼承他父親希律作了猶太王，怕到那裡去；又在夢中得到主的指示，就往加利利境內去了。 \end{tabularx} \\ \\ \relax
2:23 & \begin{tabularx}{0.7\textwidth}{X} 他們到了一座城，名叫拿撒勒，就住在那裡。這是要應驗先知所說的話：「他將稱為拿撒勒人。」 \end{tabularx} \\ \\
[1ex]
\hline
\hline
\end{longtable}
講題:生命因何活得更美好?　What Makes Life Being Well?

經課:馬太福音2章13-23節

講員:鄭漢文(rogercheng@cuhk.edu.hk)

場合:香港中文大學崇基學院禮拜堂

日期:2013年12月29日

1. 一個不曾發生的傳說

今天是聖誕節後的第一個主日,也是2013年最後一個主日,我很高興得到兩位在任的牧師──伍渭文牧師以及許開明牧師──的邀請,讓我好像回鄉探親一樣回來崇基禮拜堂講道,我今日講道的主題是:生命因何活得更美好?

首先,讓我分享一個傳說,故事是這樣記載的:

「今天是2047年7月1日,香港終於實現雙普選,即是行政首長和所有立法議會議員都是由香港選民提名並且一人一票等值地,直接地選出.

今天,在一份叫《崇基日報》的頭條新聞這樣報導:

『在2013年12月29日的這一天,有數個博士預言,香港回歸五十周年的2047年經由普選產生的行政首長會是:一位64歲的男士,他在香港中文大學畢業,是位曾經在崇基禮拜堂聚會的基督徒.

香港島在中環和西環外另有個地方叫東環,是一個被廢置的工廠,當年有一群人在那裡結集,為要搗截這位直選產生的行政首長,決定要暗殺所有在2014年內附合以上所描述特徵的30歲男士.

(在座有無附合以上所描述特徵的30歲男士?)

幸好,在2013年12月29日的同一天,上主的使者告訴這位男士,讓他暫時離開香港一段時間.直到這群殺手被捕,這個東環工廠集團被瓦解,故事中的男士才返回香港,並且在2047年前的雙普選中,成功地由香港選民提名,由全港合資格的市民一人一票等值地,直接地選出.』」

以上是一個不曾發生的傳說,純屬虛構,用來警世的,但好像以上的恐怖故事,其實早在二千多年前一個嬰孩誕生時就發生了,而且死的無辜者更多(死得人多).

2. 殺人生命的政治:馬太福音第二章經文敘事分析

今日的經課取自馬太福音第二章13至23節,是故事的後半部,整整的第二章是記載一個嬰孩誕生時發生禍及無辜死得人多的恐怖故事.請大家打開聖經,是在新約的最開首.

這是一個個案,個案有「時,地,人,事」四方面,讓我們分開何事What,何地Where,何時When,何人Who四方面一一分析這個殺人事件:

首先是何事?What happened?故事在開頭的第二章第1-2節說:

1	當希律王的時候,耶穌生在猶太的伯利恆．有幾個博士從東方來到耶路撒冷,說

2	那生下來作猶太人之王的在那裡.我們在東方看見他的星,特來拜他.

猶太人之王要降世,滋事體大,值得重視,但是在第3節說:

3	希律王聽見了,就心裡不安．耶路撒冷合城的人,也都不安.

希律王和耶路撒冷合城的人為甚麼不安?是他們的名位,權力或利益會因此而受影響嗎?

正如希律王要依靠背後的羅馬政權的喜惡來管治猶太地,香港的政治領袖也是要依靠背後的「阿爺」的喜惡來管治這個特區.忽然傳出另有一個真命天子,當然會令當權派不安!

下一步會做甚麼?就是通過時間和地點去指認這個人出來,英文叫做to identify the one.

從第4節開始,是要找查此人生在何地?Where?

4	他就召齊了祭司長和民間的文士,問他們說,基督當生在何處.

5	他們回答說,在猶太的伯利恆．因為有先知記著說,

6	『猶大地的伯利恆阿,你在猶大諸城中,並不是最小的．因為將來有一位君王,要從你那裡出來,牧養我以色列民.』

地方確定了,下一步是時間,此人生在何時?When?第7節說:

7	當下希律暗暗的召了博士來,細問那星是甚麼時候出現的.

8	就差他們往伯利恆去,說,你們去仔細尋訪那小孩子．尋到了,就來報信,我也好去拜他.

這個人在此時此地,還只是一個嬰孩.但為免博士懷疑,政治掌權者說了句語言偽術:

「你們去仔細尋訪那小孩子．尋到了,就來報信,我也好去拜他.」

故事下面的發展,顯示很難不相信:政治掌權者是不會說真實話的.

故事轉到博士的經歷,第10節開始記載:

10	他們看見那星,就大大的歡喜.

11	進了房子,看見小孩子和他母親馬利亞,就俯伏拜那小孩子,揭開寶盒,拿黃金乳香沒藥為禮物獻給他.

12	博士因為在夢中被主指示,不要回去見希律,就從別的路回本地去了.

故事轉到約瑟的經歷,第13節開始記載:

13	他們去後,有主的使者向約瑟夢中顯現,說,起來,帶著小孩子同他母親,逃往埃及,住在那裡,等我吩咐你．因為希律必尋找小孩子要除滅他.

14	約瑟就起來,夜間帶著小孩子和他母親往埃及去．

15	住在那裡,直到希律死了．這是要應驗主藉先知所說的話,說,『我從埃及召出我的兒子來.』

故事轉到希律王的反應,第16節開始記載:

16	希律見自己被博士愚弄,就大大發怒,差人將伯利恆城裡,並四境所有的男孩,照著他向博士仔細查問的時候,凡兩歲以裡的,都殺盡了.

17	這就應了先知耶利米的話,說,

18	『在拉瑪聽見號咷大哭的聲音,是拉結哭他兒女,不肯受安慰,因為他們都不在了.』

經文記載:「凡兩歲以裡的,都殺盡了.」希律王感到被愚弄而發怒,殺盡兩歲以下的無辜嬰兒.找不到這個誰人Who,就把所有近似的人殺盡.

要不要這麼殘暴?廣東話叫這種做法:「有殺錯,無放過.」

政治掌權者要殺敵人,是為了自利而清除異己,清算他人.

要殺敵人時兼濫殺無辜,是殘暴不仁.濫殺無辜時卻殺不到敵人,是為之無能.

正如香港政治的鬱悶,未能謀求全民的幸福的政治,真的是乏善足陳.

政治的黑暗還未完結,殘暴的政治掌權者的後人上位,亞基老接著他父親希律作了猶太王,依然繼續令人害怕.

第19節到23節完結了這個故事:

19	希律死了以後,有主的使者,在埃及向約瑟夢中顯現,說,

20	起來,帶著小孩子和他母親往以色列地去．因為要害小孩子性命的人已經死了.

21	約瑟就起來,把小孩子和他母親帶到以色列地去.

22	只因聽見亞基老接著他父親希律作了猶太王,就怕往那裡去．又在夢中被主指示,便往加利利境內去了.

23	到了一座城,名叫拿撒勒,就住在那裡．這是要應驗先知所說,他將稱為拿撒勒人的話了.

馬太福音第二章的故事,出現起承轉合的四部曲:

一是起於猶太地伯利恆中有人想殺人,

二是承於因避禍而離鄉別井遠走到埃及,

三是轉於重返以色列境內進入加利利,

四是合於留在拿撒勒成大長人.

馬太福音第二章故事,人物眾多,主角當然是傳說中將為猶太王的嬰孩耶穌,但是,影響情節變化相關的有四種人,包括:

(一)政治掌權的希律王,

(二)博學求真的博士,

(三)初為父親的約瑟,

(四)主的使者.

3. 養人生命的學問:博學求真,心繫永恆,謙卑尋道,記掛蒼生的博學之士

生命因何而活得美好?剛才經文解讀時,我們已經呈現政治掌權者可以草菅人命.

生存的最基本條件,就是免受死亡的威嚇.這是政治最低限度要為人民謀求的.

聖誕節慶祝耶穌的誕生,馬太福音第二章告訴我們:不是因為不回頭傳話給殘暴希律王的博士,不是因為在適當時傳話的主的使者,不是因為非親生的爸爸但願意聽從而披星戴月的約瑟,耶穌一出世就被政治殺害了.

所以,今天是聖誕節後的第一個主日,就是要讓我們反思,我們在出世後:有那些恩人你們要感謝─你的生命才能存活到今天?有那些恩人你們要感謝─你的生命才活得這麼美好?

馬太福音第二章記載的恩人,包括:東方來的博士.

究竟有多少個博士?傳統說有三個.

其實聖經沒有明言有三個,而是有幾個博士從東方來到耶路撒冷.因為送上黃金,乳香,沒藥為三種禮物,所以就好似一人各自呈獻一種,共有三個送禮人.

另外有一些相關文獻記載,其實是一群朝聖者,朝聖博士群中為首的有三位,三位都是有名有姓的.(叫做:美爾加斯Melchus,卡斯珀Casper和法迪沙達Phadizarda.)

有些後期抄本,指出這些智者是星象學家.他們好博學,能夠觀天象而通世情.

和修本的聖經譯為:有幾個博學之士.

所以,這幾個朝聖的博學之士,並不是我們現今大學讀完碩士再讀博士的學位持有者.

他們也不是現今大學中的教授級博士,拿到研究資助經費,然後成立小組,進行實地考察.

他們也不是已經在大學任教,做研究出論文可以上位,權傾大學任主管的教授.

他們也不是捐了金錢給大學,大學回報給予榮譽博士,無做學問卻又想有名譽的社會人士.

他們也不是學有專精的專家學者,但卻才力狹窄時心胸又窄的所謂博士.

聖經稱這些智者為博士,當然並非現今大學的博士,而是博學求真,心繫永恆,謙卑尋道,記掛蒼生的智者.

因為通天文,知地理,鑑人情,所以老遠從東方跑來,帶同黃金,乳香,沒藥作為三種禮物,為求得見拯救世人的君王.

黃金,乳香,沒藥這三種禮物代表甚麼呢?有解經家認為:黃金代表敬重王者的尊貴,乳香代表尊崇神的獻祭,沒藥代表膏抹救世主的犧牲.

尋求人生可以活得更美好,是人的本性.

政治的本質,就是讓所有人都能活出美好的人生.從西方希臘哲學家亞理士多德論城邦民主,和中國古代孟子論仁政,到當今的自由主義和社會主義的意識型態相爭,都認為政治是要建立讓人活得美好的美好社會.

我個人近年正在經營這樣的「轉化生命的學問」,就試探討讓人生活得更美好,讓社會政治更美好的學問,是保護弱勢社群的學問,是讓人盡其才,活出真我的學問.

有經營這種「轉化生命的學問」的博士嗎?

有誰願意說:「主啊!我在這裡,請差遣我.」

4. 育人生命的親人:為子女活得更好而奔波的父母親約瑟和瑪利亞

政治與學問是世間的兩種權能,活出美好人生需要超阻力,借助力:

一是超克殺人生命的政治,化黑暗阻力為燃點光明；

二是借助養人生命的學問,但願上天早生聖人；

幸好還有第三種人,就是育人生命的親人,人人皆有的家人.

在上主日助理校牧Walter林豪恩先生,以「欲休還娶」為題,跟大家講解過馬太福音第一章,其實聖靈感孕的耶穌,令得馬利亞幾乎被約瑟休了.第二章就看見,有多少次因為約瑟聽從主的使者所吩咐,與剛生育不久的馬利亞到處奔波.因為約瑟,不為己出的耶穌,才可以婚生,才可以出生,才可以長大成人.

剛才我們從馬太福音第二經文的地理轉變,就體會到這個父親在耶穌出生後的起承轉合:

一是起於猶太地伯利恆中有人想殺人,

二是承於因避禍而離鄉別井遠走到埃及,

三是轉於重返以色列境內進入加利利,

四是合於留在拿撒勒成大長人.

正如古代有孟母三遷,是一個單親家庭中的母親三次遷家,為的是讓兒子成長得更好.

回想現今,你想一想非婚生子女的一對父母,帶同初生嬰孩,有些在中國境內跑到香港,正如我們有些香港人曾跑到英國或者加拿大,再跑回香港.

在香港,他們分享了一點社會保障的資源,呼吸一下自由的空氣,是因為人的本性在於盼望活得更美好.

我們可否讓香港的政治容得下尋求美好生活的過客?

正如這個崇基禮拜堂容納了不少曾經是過客的你們,保護過心靈曾經軟弱的你們,成為收留信仰孤兒的家園.

崇基禮拜堂好像一個家,伍渭文牧師好像我們的爸爸,每主日分享屬靈食物餵養大家.

是時候,談一談2013年年底退休的伍渭文牧師.

5. 得聽主心的使者:讓人人成為上主使者的牧養

伍渭文牧師在崇基出任校牧十三年半,任內建立崇基禮拜堂.

他管治清明,得享清風甚少紛爭.

他貴為博士,講章禱文文采出色.

他情如父親,牧養靈內子女無數.

他的同在,如主的使者,讓人得知上主的心意.

他的牧養,讓弟兄姊妹也成為主的使者.在別人生命路上,互相提醒,得聽上主的心意.

生命因何活得更美好?今天的講道我們得見四大恩情:

在家庭有父母親人,在學校有生命學問,在社會有政治清明,一生一路走來有主的使者.

余杰,王怡兩位在2010年出版的《一生一世的仰望》中,訪問八位牧者,其中一位是伍渭文牧師,他們為伍牧師的訪問起了這樣的題目:「引導大學師生追求人生的真諦」,認為校牧是大學校園裡的”靈魂的守望者”,認為校牧在大學中提供終極關懷的維度和”增值”的教育,認為崇基禮拜堂仿如暗夜的燭光.

崇基禮拜堂教友為伍牧師榮休而出版了《悠然見鞍山》,這本書只是十三年半牧養書信冰山的一角,還有尚未窺探的深藏不露.

香港中文大學校長沈組堯教授這樣作序:「伍牧師學問淵博,靈感充沛；發揮成為鏗鏘可誦的文章,擲地有聲而感人至深.」

崇基學院院長兼香港中文大學文學院院長梁元生教授這樣作序:「他的禱文往往都文辭優美,並且回應時事,激動人心,容易引起迴響和共鳴.在這樣忙碌的工作環境下,最難的是懂得欣賞悠然,一種放得開的性格,和入世而又出世的生活態度.」

其實,我今天能站在這裡講道,也是伍牧師起用新人,有容人之量,我覺得他其實是慧眼識英雄(自許一笑).是兩年前的2012年1月開始,伍牧師讓我在這裡作神學生實習,到今年2013年4月底完結.之後,他邀請我每年四次的第五個主日到此講道,並得許開明牧師同意.

今天是2013年最後一次講道,也是伍渭文牧師作為校牧榮休的最後一次崇拜.我本一介文人,身無長物,我謹以這篇講章作為賀禮,當然不及黃金,乳香,沒藥的貴重,不過也算是一名博士在經營「轉化生命的學問」的一篇成品.(伍牧師,你放心投資在我身上!)

再者,各位不要見笑,我用了「伍渭文牧師」的人名以及《悠然見鞍山》的書名,作了七言絕詩一首,呈獻給上主,謹向伍牧師──我們的生命師傅──榮休致敬:

「幔開鞍山見堅穩,牧情悠然現胸襟,今雖退伍師常在,渭河長流主禱文.」

讓我們一起低頭禱告:

「上主啊,我們行經時間,你賜下甘霖.

回首往來蕭瑟處,也無風雨也無晴.我們感恩.

上主啊,求你教我們回首,細想我們的人生,是因何而活過來?

誰曾為我付出甚麼?我們卻不曾說過多謝.懇請上主你,祝福這個誰.

上主啊,求你教我們感恩,細想我們人生歷過困苦鞎難時,是誰幫助我們渡過?

是育我生命的親人,是養我生命的學問,是活我生命的政治,是主你所差來引導我的使者?

上主啊,求你教我們盼望,讓我們成為他人生命中的主的使者,助他人活得更好.

上主啊,求你悅納伍渭文牧師在這裡十三年多的服侍,求你恩待立意要事奉你的伍牧師和他的一家,求你繼續祝福他和他一家.

奉主耶穌基督的名祈禱,阿們.」

但願榮耀都歸給聖父,聖子,聖靈,永在的三一神.阿們.

(鄭漢文博士為循道衛理聯合教會馬鞍山堂會友.

現職香港中文大學教育行政與政策學系專業顧問,

同時在崇基學院神學院兼讀神道學學士課程.



\newpage

\section{年有道與祂共渡}
\label{sec:231}
\textbf{許開明牧師 - 新年有「道」,與「祂」共渡}
\newline
\newline
\hyperref[sec:230]{< < < PREV SERMON < < <}
~
\hyperref[sec:toc]{[返目錄]}
~
\hyperref[sec:232]{> > > NEXT SERMON > > >}


約翰福音 1:1-18
\newline
\begin{longtable}{cl}
\hline
\hline
章節 & 經文 (和合本修訂版)\\
\hline
1:1 & \begin{tabularx}{0.7\textwidth}{X} 太初有道，道與神同在，道就是神。 \end{tabularx} \\ \\ \relax
1:2 & \begin{tabularx}{0.7\textwidth}{X} 這道太初與神同在。 \end{tabularx} \\ \\ \relax
1:3 & \begin{tabularx}{0.7\textwidth}{X} 萬物都是藉著他造的，沒有一樣不是藉著他造的。凡被造的， \end{tabularx} \\ \\ \relax
1:4 & \begin{tabularx}{0.7\textwidth}{X} 在他裡面有生命，這生命就是人的光。 \end{tabularx} \\ \\ \relax
1:5 & \begin{tabularx}{0.7\textwidth}{X} 光照在黑暗裡，黑暗卻沒有勝過光。 \end{tabularx} \\ \\ \relax
1:6 & \begin{tabularx}{0.7\textwidth}{X} 有一個人，是從神那裡差來的，名叫約翰。 \end{tabularx} \\ \\ \relax
1:7 & \begin{tabularx}{0.7\textwidth}{X} 這人來是為了作見證，是為那光作見證，要使眾人藉著他而信。 \end{tabularx} \\ \\ \relax
1:8 & \begin{tabularx}{0.7\textwidth}{X} 他不是那光，而是要為那光作見證。 \end{tabularx} \\ \\ \relax
1:9 & \begin{tabularx}{0.7\textwidth}{X} 那光是真光，來到世上，照亮所有的人。 \end{tabularx} \\ \\ \relax
1:10 & \begin{tabularx}{0.7\textwidth}{X} 他在世界，世界是藉著他造的，世界卻不認識他。 \end{tabularx} \\ \\ \relax
1:11 & \begin{tabularx}{0.7\textwidth}{X} 他來到自己的地方，自己的人並不接納他。 \end{tabularx} \\ \\ \relax
1:12 & \begin{tabularx}{0.7\textwidth}{X} 凡接納他的，就是信他名的人，他就賜他們權柄作神的兒女。 \end{tabularx} \\ \\ \relax
1:13 & \begin{tabularx}{0.7\textwidth}{X} 這些人不是從血生的，不是從情慾生的，也不是從人的意願生的，而是從神生的。 \end{tabularx} \\ \\ \relax
1:14 & \begin{tabularx}{0.7\textwidth}{X} 道成了肉身，住在我們中間，充充滿滿地有恩典有真理，我們也見過他的榮光，正是父獨一兒子的榮光。 \end{tabularx} \\ \\ \relax
1:15 & \begin{tabularx}{0.7\textwidth}{X} 約翰為他作見證，喊著說：「這就是我曾說：『那在我以後來的先於我，因為在我以前，他已經存在。』」 \end{tabularx} \\ \\ \relax
1:16 & \begin{tabularx}{0.7\textwidth}{X} 從他的豐富裡，我們都領受了恩典，而且恩上加恩。 \end{tabularx} \\ \\ \relax
1:17 & \begin{tabularx}{0.7\textwidth}{X} 律法是藉著摩西頒佈的；恩典和真理卻是由耶穌基督來的。 \end{tabularx} \\ \\ \relax
1:18 & \begin{tabularx}{0.7\textwidth}{X} 從來沒有人見過神，只有在父懷裡獨一的兒子將他表明出來。 \end{tabularx} \\ \\
[1ex]
\hline
\hline
\end{longtable}
講題:新年有「道」,與「祂」共渡　Living With “Logos”In The New Year

講員:許開明牧師

經文:約翰福音1:1-18

場合:香港中文大學崇基學院禮拜堂

日期:2014年1月5日

引言

今日是2014第一個主日.恭祝大家新年進步.

我們透過今日的經文:約1:1-18節來思想,「新年有「道」,與「祂」共渡」.

我們都喜歡新年,新年給人一種新的感覺,從新開始的機會.

一首詩如此寫:

「舊的過去,又來了新的一年,

每個人又有新的機會,新的希望,新的夢想.

過去的門已關了起來,

我們抬起頭來迎接新的一年.-」-Hilda Butler Farr

每年都有開始,都有新年,有無意思?

因著新的開始,新的更新,新的希望,我們要慶祝並存著感恩面對未來.

然而,當想到今年2014年,你會點?

你會發現,有許多不知道的挑戰,困難?因此有疑惑,有憂慮!

但你也會有些知道?

你會老一點!不單是樣子老,思想老,心態也會老,老化,衰老!

你的限制多一點,身體的限制,工作的限制,發展的限制!

人對未來,未能把握,或人怕未來:會病,會死,會失敗,會受苦,因而期望有一位有能力又掌握未來的上帝出現.

然而,是否有這樣的上帝?

如果有的話,當你相信上帝時?是否病,死,失敗,受苦便不再出現?

如果我們期望的上帝,相信祂之後,我們仍然面對人生的種種困難,會死,會失敗,會受苦.

我們的上帝不能使我們趨吉避凶,排除萬難.反而信耶穌要受逼迫,遭患難,背十架,你為何還要信祂,還要跟從祂?

再簡單點說,你向上帝祈求十樣東西,上帝只應允你一樣,或一樣都無,你還要相信祂嗎?依靠祂嗎?

我今日的題目是「與祂(上帝)同行」,如果祂好像一無可為,那你為何仍要與祂同行?

再問,我們祈求與上帝同行,是有目的,是存有一種功利的心態,這又是否正確呢?

這一連串的問題,希望引發大家思想,信仰的真義.

今日的經課約1:1-18節是約翰福音書的序言.

約翰寫作約翰福音,目的是要使當代羅馬人,希臘人和受希臘文化影響的猶太人接受耶穌是基督,是上帝.他介紹上帝就是道,上帝與道同在,這道就是耶穌基督.

約翰綜合當代希臘與希伯來人的思想,建立「上帝觀」:

太初有道,(未有世界之前,已有道,已有言語)

道創造天地,(道創造這世界,言語創造世界,上帝說:有就有)

道進入世界,(道創造完世界後,衪進入世界,是看不見的靈(言語))

道成肉身,(道為要使人看得衪,穿上了外衣,讓人看見)

這住在人的中間.(道住在入群的中間,使人聽見說話,看見外衣)

上帝及基督都是「道」Logos

「道」是由赫拉克利特(Heraclitus,西元前535-475年)首先提出來的,他認為 λόγος 為萬物的起源,而萬物也有一部分的 λόγος 在其中；他又認為「道」 λόγος 是神聖之火,是神.火產生熱力和光能.「道」是「動性法則／理性」,宇宙和一切存在所依賴的.

到了西元前300年左右,斯多亞學派(Stoics)的始創者季蒂昂的芝諾(Zeno of Citium)提出λόγος是充滿宇宙和一切存在所依賴的「動性法則／理性」(Active Reason),人類也擁有一部分的「道」.

與約翰差不多時期的,有亞歷山太的斐羅(Philo of Alexandria,西元前20至西元後40年),他認為「道」是創造的法則(Creative Principle).「道」就是一個仲介體:神,人世界的存在物.

承接柏拉圖有關創造的思想,世界是由「仲介體」所創造,而「仲介體」未必等於至高的神,斐羅認為絕對神聖的神不會與不完美的物直接接觸,而「道」就是一個仲介體,負責創造的工作,把先存而沒有形質的物轉化成四種基本元素.

因此,在他的思想體系中,「道」並不是那位絕對超然的神,只是仲介,僅次於神的聖言,既非造成亦非生成,是處於神和人,宇宙之間的存在.

總括而言,希臘哲學思想認為:

1「道」λόγος 並不是那位絕對神聖的神,只是神創造世界的仲介體,不是真正的源頭.2「道」並沒有位格,更不會與人接觸.

然而,約翰福音一章1至18節,清楚回應創造世界的是神,而λόγος,神的第二位格,並不是次等,祂與神的本質和來源相同,在時間和空間被造之前,已經與神同在.約翰福音一章1至3節以重複和正反角度去描述道是神,是創造者(不是被造,不是次等).

約翰更指出道成了肉身(σὰρξ ἐγένετο,約一14),說明耶穌並不是幻影,也不是附在一個人的肉體當中來執行任務,教導世人正確認識有關神的真理.保羅稱這為敬虔的奧祕(提前三16).這位絕對神聖超然的創造者,竟然還會住在我們中間(ἐσκήνωσεν ἐν ἡμῖν ,約一14)

住(σκηνόω)與會幕(σκηνὴ；出三十三7)的字根相同,而會幕(或聖所)代表神可以住在他們中間(出二十五8),意即「神與人同在」.道與世人同住,8對人類來說真的是不可思議的福氣和恩典,也反映了神是何等愛祂所創造的人!

「道」,它關聯到一個動詞,這動詞在英文繙譯裡一點也不顯眼,但在希臘文裡是十分突出的,但不在其動詞本身,而在於它使用的時態.

我們先來看這個名詞「道」.

它是甚麼意思?在希臘哲學中logos這字是常用的字.

約翰使用這個字時,是受到希臘哲學的影響,也受的影響倒是來自希伯來哲學,這是很重要的區別,我建議每一個學生都看伯尼博士(Dr. Burney)的書《第四福音之亞蘭原文》(The Aramaic Origin　of the Fourth Gospel),在某些點上,我並不同意伯尼博士的見解.但無疑的,他證明瞭約翰福音的思想是屬於亞蘭文(當時希伯來人通用的語言)方式的,而不是希臘文方式的；並且他發現這本書必定是一個巴勒斯坦猶太人的著作；並且作者寫此書時決不會遲於主後七十五年.因此他主張此書作者的想法是屬於希伯來人的.

當約翰使用希臘文裡的logos時,他必定是得自希伯來文Memra這字的啟發.這個字具有希臘哲學的意義,但更是合乎希伯來哲學.那麼,希臘文的這個字是甚麼意思呢?

它所指的乃是整個思想領域,是隱藏在每一具體東西背後之抽象觀念.

其觀念或許可用「智慧」這個字來表達.希臘哲學家認為智慧是在一切工作之先,且是一切現象原先的本體.

希伯來哲學家說:事物是來自思想的,思想的成果.希伯來哲學家說:事物是來自思想的.只要有事物,就證明有思想.

希伯來哲學家進一步又說,「如果事物是來自思想,那麼思想就來自思想者.」因此,這宇宙會有位最偉大的思想者,就是耶和華.

希臘哲學家並沒有說得這麼多.

希臘哲學家說,「在一切事情背後,必定有思想,但思想是抽象的(看不見,神秘).」希伯來哲學家說,「除非有思想者(上帝),否則就不能有抽象的思想.」因此,希伯來哲學家說,「敬畏耶和華是智慧的開端.」希伯來哲學家說,「在這宇宙中沒有至終無法解決的問題.」它是在思想者的心思中得著解決的.「起初神創造天地.」

神就是思想者,然後有思想.而一切事物,乃是這思想者思想和行動所造成的後果.關於第1節中的名詞就說到此為止.我們再來看其中的動詞,「太初有道,道與神同在,道就是神.」正如前面說過的,這樣的繙譯,對我們來說並沒有甚麼奇特之處,因為英文並不像希臘文那樣有變化.

希臘哲學家重理性的,抽象的思維:

希伯來人的思想重超理性,超自然的思維,上帝是超然的；還有希伯來人在艱苦中成長的民族,重真實的經歷,上帝是來到人們中間的上帝,耶和華是進入人類時空歷史的上帝,是亞伯拉穻的,以撒的,雅名的上帝.

「道」有豐富的含意,從中可見基督教上帝的偉大:

1言語(Word)

這字「道」在約出現了三十九次,基本音思是所說的「話」.

創世記強調道是「言語」,上帝藉言語創造萬物.

「上帝說,萬物便造成」詩篇三十三:六

言語是帶有能力的,耶廿三:二十九「我(上帝)的話豈不像火,又能打碎磐石的大錘麼?」

2原理(Order):宇宙的秩序與法則,創造萬物的原理,因此,「道」在萬物之中,萬物都分享上帝神性的部份.主前六世紀,希臘哲學家赫拉克利特說:「道是永恆的,整個宇宙藉它而存在.」

3智慧(Wisdom):上帝的智慧.

箴言三:十九「耶和華以智慧立地,以聰明定天.」

智慧是有位的存在,並且早於萬物未受造之前已經存在.

4先存性(Pre-existence)/永恆性(Eternity)

5有位格(Personality)

上帝是道,耶穌也是道.

我們所信的上帝是創造的主,生命的主,永恆的主:祂具有下列特點:

1言語,2原理,3智慧,4永恆,5有位格的主(有位格,有性格,有獨立的存在,才能道成肉身).這是值得信靠的上帝；同時也是非常偉大的上帝,超過人的思想和想像.

Is your God big enough?

你的上帝夠大嗎?

約翰眼中的上帝是道(Logos)是非常偉大!更偉大的是道成肉身!

道(話)成肉身(約一14~18) ,    神的存在是客觀的真理；但神是住在人不能靠近的光裡,是人未曾看見,也是不能看見的(提前六16).雖然說,神的事情,人所能知道的,原顯明在人心裡；因為神已經給他們顯明.自從造天地以來,神的永能和神性是明明可知的,雖是眼不能見,但藉著所造之物,就可以曉得,叫人無可推諉(羅一19~20).

然而不能看見,不能觸摸,總是不太具體,缺少實感.神愛世人,祂樂意看見世人得著恩典,脫離罪,死,魔的轄制,得著永生,藉著耶穌基督得兒子的名分,並且活在祂面前成為聖潔,無有瑕疵；使祂榮耀的恩典得著稱讚(參照弗一1~14).

今日許多人有意無意中把上帝局限,他們所信的是微小的上帝,從而對上帝不敬也不自知,有下列情況,我們要留心:

1人把上帝私有化

過份強調我的上帝,我的主!他心中的上帝,眼中只有他,為他服務.上帝對我特別好!上帝一定幫我!上帝不會虧待我!

你看!我常上教堂,我全家歸主,我今年又奉獻了一大筆錢給教會,我又熱心事奉和傳福音,上帝不會也不應虧待我.

有些人注重自己的功德,祈求得到上帝的喜悅.把上帝太過擬人化,私有化了.

早期以色列民也有這種狹窄的上帝選民的思想,以為上帝會特別優待亞伯拉穻的子民,上帝的選民.事實上,以色列亡國,被擄,被大屠殺,是苦難重重的一族.

請不要忘記上帝是我的上帝,你的上帝,他的上帝,我們的上帝,普世的上帝.

我求上帝賜我健康,平安,喜樂,無可厚非!但也不要忽略為鄰舍向上主祈求,為社會向上主祈求,為國家向上主祈求,為世界上的人向上主祈求.

2人想做上帝

人建造巴別塔,人想做上帝,人想得到榮耀!

人以為可以控制一切,人只看到自已的成就.

小山晃佑(Koyama Kosuke)寫道:

「上帝固然給我們陣陣涼風,但科技卻給我冷氣機；

上帝固然創造天上的飛鳥,但科技卻給我們噴射的飛機；

上帝沒有醫治病人,醫生卻用科技來治療病人.

今日每一種發明,實用都不歸諸上帝而歸於人；

愈多人手的科技,愈少神秘的上帝；

上帝,你可以退休,在幽靜的天堂裡休養.」

現代人已不再說,「我們在天上的父,而是我們在地上的科技(電腦,電話,電子產品)」人自視太高,人自以為可以做上帝.

3 人想利用神,人要上帝辦事

受後現代思潮的影響:

今日不少的教會強調要享受神,享受與上帝一齊的喜樂!

神是給人享受的嗎?

與神一起,不一定是喜樂的,神會鞭撻人,令人不平的.

我在一間教會牧會,逢周六祈禱會,姊妹們上午七時便到堂祈禱,一位姊妹每次祈禱都列出一系列要上帝作工的工作清單:

上帝我個細仔,快考試,給他得到好成績；上帝我個女已三十歲了,給她配偶；

上帝我個大仔失業了,給他工作；上帝我經常胃痛,醫好我!

上帝是人的奴僕嗎?上帝受命幫人做野?如果上帝照人的吩咐做事?

上帝還是上帝嗎?不是你下命令,上帝就照做!

你有信心,上帝會成就?

你無信心,上帝要成就的事,祂也會成就.

巴特(Karl Bath):讓上帝成為上帝(Let  God be God)；

上帝是主,自主,自由,自我.

與「道」同行

1) 不一定得福的.2) 不一定一帆風順的.3) 不一定因你的信心而應允.

與上帝同行,是與上帝經歷人生的甜酸苦苦辣!

一位弟兄說:「我所信的上帝,不是聽我話的上帝,而是我要聽祂話的上帝.」

所以與上帝同行,不是說:上帝一定聽你禱告,這是迷信上帝!

上帝一定在危難時伸手幫你,這是迷信上帝!

上帝一定不會讓你跌倒!這是迷信上帝!

與上帝同行,與主同在是:相信上帝一定在我身旁,我的人生有了「道」

我不再孤單,我有信靠的對象,我有上帝,就有了「道」,有智慧,

我在一切的困境中中仍有希望.

一位在德國集中營的弟兄所寫的詩:雖然我看不見藍天,我仍相信有藍天；雖然我見不到陽光,我仍然相信有太陽；雖然我看不到上帝,但我仍然相信有上帝.

存在主義者尼采(德語:Friedrich Wilhelm Nietzsche,1844年10月15日－1900年8月25日) 認為上帝創造世界後,衪就離開這世界,讓世界自然運作,因為上帝不獨創造萬物,而且創造了原則和規律. 正如造鐘錶的工匠,造好後便任由那鐘錶自由運作,不用再理會它. 所以道創造萬物,道出世界,道不再與人同在同行.

人的存在是充滿危險的,他的一句名言:“人生像走兩個山峰間的鋼線,前進是危險的,退後更危險,站著是最危險的,總之人生就充滿了危險!”沒有上帝同行的人生是危險而無助的. 有上帝同行的人生,雖然仍有危險,但不是無奈的,是充滿希望和力量.

我們不知,今年2014是怎樣的一年?

願我們信靠上帝,勇往直前.阿們.



\newpage

\section{祈禱進階2天意從來高難問}
\label{sec:232}
\textbf{賀志勇博士 - 祈禱進階II:天意從來高難問?}
\newline
\newline
\hyperref[sec:231]{< < < PREV SERMON < < <}
~
\hyperref[sec:toc]{[返目錄]}
~
\hyperref[sec:233]{> > > NEXT SERMON > > >}


講題:祈禱進階II:天意從來高難問?　Is It Hard to Know What God’s Will Is?

講員:賀志勇博士

場合:香港中文大學崇基學院禮拜堂

日期:2014年1月12日

一,引言  天意從來高難問?

最近有朋友向我抱怨,說她兩個兒子自己家的飯就不想吃,一聽說要吃麥當勞就歡呼.她說她搞不懂孩子怎麼會這樣了.另一個朋友插嘴說:那是因為小孩子去麥當勞就當是去過節,是去慶祝的,那當然要歡呼啦!

原來如此!去麥當勞就像去過節,是會讓孩子歡呼的,那來教堂做禮拜呢?來教堂做禮拜是不是像過節,是不是會讓人歡呼呢?我們中間有沒有人每個星期天一起床,會大聲說:“太好了,我又可以回教會做禮拜啦!”       我們來教會是抱著這種心情呢,還是說,我們來教會是因為我們有一個大老闆,乃是在天上的,所以每週我們不得不來,雖然偶爾我們也會遲到一下,但這個老闆總是有恩典˴有憐憫的,沒有什麼可以讓我們與祂的愛隔絕的.是不是這樣呢?

如果說做禮拜不是那麼讓人激動人心,那麼歡呼雀躍,那我們的信仰生活--或者說,我們基督徒的生活,是不是也存在著同樣的處境呢?

我有一些朋友就認為信仰生活並不是那麼讓人興奮,因為他們覺得,他們搞不懂上帝到底要他們做些什麼,會給他們些什麼?

比方說,A女士常常問上帝到底是劉先生還是張先生更適合做她的老公,更符合上帝的心意,可是上帝總是保持沉默；而B先生就想知道他是不是該換份工作,他本來是個麻醉醫生,最近突然想轉行去做時尚設計(fashion design),但他不知道上帝是不是喜悅他這種轉變,或者說上帝會不會祝福他四十歲才開始打拼的人生夢想；而C女士最近在為她的兩個兒子發愁,兩個兒子越來越瘦,卻總是不肯好好吃飯,甚至為了搶遊戲機差點打起來了,她搞不懂上帝為什麼要這樣對她呢,她不是很按時上禮拜堂禮拜,按時奉獻,按時上查經班麼,怎麼還這麼不蒙福呢?

他們都常常禱告,可是上帝似乎總是保持沉默--這就是現代基督徒的困惑,雖然他們不認同尼采說的“上帝已死”,但他們懷疑上帝是不是常常在睡覺?有時候我們是不是也有這種遭遇--我們凡事常常禱告,但卻總是沒有回音；還有另一種情況是,當我們聽到別的信徒分享,說神聽了她的禱告,給了她很多祝福,或者讓她知道該怎麼樣怎麼樣做出選擇的時候,我們腦海裡就會浮現一個大大的問號(question mark),我們會想:是不是真的啊?是不是你的幻覺啊?

二,解經

1,可以得知神的旨意

這兩種情況都體現了一點,我們對於我們能夠明白神的旨意這一點不太有把握.那麼,究竟我們能不能夠明白神的旨意?若是可以,我們怎樣才能明白呢?

以弗所書5:17節說“不要做糊塗人,要明白主的旨意如何”,這說明主的旨意是可以明白的,而詩篇32:8節更清楚顯明了這一點,在這一節裡,耶和華對大衛說:“我要教導你,指示你當行的路.”這就告訴我們,尋求神的旨意並不是捉迷藏,神並不是讓我們獨自面對世界,祂會教導指示我們.

這就產生了一個問題,是不是我們的任何事,神事先都一一幫我們計畫好了,我們只需要求問神的旨意,然後奉天承運,按旨行事,就一定會蒙主的祝福了；否則如果我們沒有神的明確指示就自己行動,那我們是不是就自由散漫,偏行己路呢?問題的答案就在於,神有沒有做出這樣的教導和指示.

2,聖經中顯明神的旨意

帖撒羅尼迦前書4:3-4節很清楚地寫道:“神的旨意就是要你們成為聖潔,遠避淫行；要你們各人曉得怎樣用聖潔˴尊貴守著自己的身體.”另外,在5:18節中又寫道:“凡事謝恩,因為這是神在基督耶穌裡向你們所定的旨意.”

這兩句經文清楚表明,神的旨意,或者說神對我們做事的指示主要歸結為這兩點:第一,成為聖潔；第二,凡事謝恩.

聖潔原文的意思就是分別出來,所以聖經上有一句話說“分別為聖”,也就是說神的心意,是要我們基督徒行事為人跟這個世界的人有分別.如果別人買六合彩我們也買六合彩,如果別人背後講是非我們也學人講是非,如果別人眼裡只有穿細軟衣服的,沒有稅吏和兵丁,我們眼裡也只有穿名牌的,沒有弱小孤苦的人,那我們跟這個世界有什麼分別?如果我們與這世界行事沒有什麼分別,我們當然無法為神做出自己的見證.

所以當有的教友問我,什麼樣的選擇合神的心意時?我會告訴他,只要他的選擇是聖潔的,那麼他的任何一個選擇都是合神心意的.

如果說聖潔是針對行動來說的話,那凡事謝恩是針對我們的心態來說的.也就是說,神的旨意是讓我們在聖潔的同時,還抱著一種凡事謝恩的心態.凡事謝恩,用約伯的話來說就是“賞賜的是耶和華,收取的也是耶和華,耶和華的名是應當稱頌的.”用保羅的話說就是“我靠主大大的喜樂......我知道怎樣處卑賤,也知道怎樣處豐富.”(腓立比書4章)

凡事謝恩這一旨意告訴我們,神希望我們成為這樣的人:在我們成功的時候仍舊懂得謙卑順服,而在我們失意的時候並不會消沉不振.

綜合起來,神的旨意就是要我們在行事為人上聖潔,可以做世人的榜樣；同時要我們無論在順境逆境,都要懂得神才是天地的主,我們無須太過懼怕,都要懂得感謝他.

3,如何實行神的旨意

就這麼簡單?有的教友不太放心,可能還會追問一句,怎麼樣才叫做成為聖潔?首先想說的是,成為聖潔的確有一些法則,正如詩篇40:8節所說:“我的神啊,我樂意照你的旨意行；你的律法在我心裡.”這節經文就告訴我們,神的旨意可以在神的律法中去尋求,去體會.其次,對於基督徒來講,這律法不再是舊約的律法,而是主基督為我們所立的“新”律法.耶穌基督並沒有改變律法或“補充”律法,他只是對律法作了更忠實,更深刻的解釋.比方說,他在登山寶訓中說:“你們聽見有吩咐古人說‘不可殺人’,又說‘凡殺人的,難免受審判.’只是我告訴你們‘凡向弟兄動怒的,難免受審判”.這不僅解釋了律法,也更深化了律法.

所以成為聖潔,凡事謝恩看起來簡單,實際行動,還必須紮根到聖經中去,去體味聖經的敘事,才能更好地活出這兩個原則.

或許你又要問了:“我連聖經都沒有看完,有的章節,更是一拿起來看就想睡覺,那怎麼去活出神的旨意呢?”不用擔心!聖經早替你想到了答案.以弗所書2:13節就講到:“因為你們立志行事都是神在你們心裡運行,為要成就祂的美意.”也就是說,神的旨意這麼高難度的任務並不是讓我們單獨承擔的,神的靈會在我們心裡運行,會與我們同在,一起去活出神的旨意.

但是我們並不能躺在那裡,單單等著神在我們心裡運行.正如羅馬書12:2節所說:“你們不要效法這個世界,只要心意更新而變化,叫你們察驗何為神的善良˴純全˴可喜悅的旨意.”察驗這個字,在原文就有體驗並且證實的意思,也就是說,你不能躺在那裡等著神的旨意,人有人的責任,神的旨意,是需要我們在每天的生活中去體驗並且去證實的.

不僅僅如此,保羅在歌羅西1:9節還說:“願你們在一切屬靈的智慧悟性上,滿心知道神的旨意.”一提到悟性,有的朋友就想,從小老師就說我悟性不夠,那怎麼辦呢!其實悟性這個字翻譯得不好,原文的意思乃是理智的意思.也就是說,神希望我們用理智去思考,去行動,去好好把握祂的旨意.

總之,用廣東話說,信仰是要用腦的.神的旨意並不神秘,它只是需要我們好好動用我們的理智,在我們的生活中好好去體驗,並通過我們的聖潔˴感恩的生活,去證明“神真是又真又活的神”.

當然,我不排除在我們生命中某些特定事件,神可能會親自來指引,但我想那可能是跟我們的救贖和成聖相關的事吧.也就是說,上帝有沒有替我們預定一個太太或老公,起床的時候穿鞋要先穿左腳還是右腳,我到底要不要換一份工作這一類事情,上帝有可能替你指定了答案,但也很可能上帝並不回答這類問題；在這些事上,祂可能更希望我們自己動動腦筋,用理智去打理好.祂更關心的是,我們在做選擇時,能否成為聖潔,並凡事謝恩--至少從聖經的敘事來看是如此.

三,結論 天意並不高難問

當我這樣分析完,有個朋友就跟我講:“這樣太沒有安全感了!本來以為我什麼事神都安排好了,我只要好好去明白他的心意,按照他的旨意辦,就事事都好辦了.結果你告訴我,神讓我們很多事自己做決定.”沒錯,神希望我們有事去尋求祂的心意,因為那是人謙卑順服的責任；但是傳道書3:10-11節也提醒我們:“我見神叫世人勞苦,使他們在其中受經練.......然而神從始至終的作為,人不能參透.”

經練就是磨練.很多時候,上帝就像父親,祂並不總是拉著我們的手,綁定我們這一步,下一步.他給了我們自由!更多時候,他像腳前的燈,路上的光,照亮我們前行的方向,卻放手讓我們去運用理智,並磨練我們的品格和能力,去選擇活出聖潔.

而信仰生活就像探險,探險,就充滿了許許多多的未知,我們不知道我們會碰上什麼危險,也不知道我們將會有什麼驚喜,我們唯一知道的是:上帝與我們同在!

而正是因為上帝與我們同在,那些生命中的小事--結婚生子,探訪病人,同別人和好--都獲得了指引,都獲得了救贖；而也正是通過這些信仰生活中的小事,我們鍛煉了我們的信仰品格,學習凡事感恩,並且活出神的旨意.如果搞懂了上帝的這些旨意,我們還覺得信仰生活不是一件讓人興奮的事麼?

天意並不高難問!



\newpage

\section{成長的階梯}
\label{sec:233}
\textbf{鄧瑞強博士 - 成長的階梯}
\newline
\newline
\hyperref[sec:232]{< < < PREV SERMON < < <}
~
\hyperref[sec:toc]{[返目錄]}
~
\hyperref[sec:234]{> > > NEXT SERMON > > >}


約翰福音 1:35-39
\newline
\begin{longtable}{cl}
\hline
\hline
章節 & 經文 (和合本修訂版)\\
\hline
1:35 & \begin{tabularx}{0.7\textwidth}{X} 又過了一天，約翰同兩個門徒站在那裡。 \end{tabularx} \\ \\ \relax
1:36 & \begin{tabularx}{0.7\textwidth}{X} 他見耶穌走過，就說：「看哪，神的羔羊！」 \end{tabularx} \\ \\ \relax
1:37 & \begin{tabularx}{0.7\textwidth}{X} 兩個門徒聽見他的話，就跟從了耶穌。 \end{tabularx} \\ \\ \relax
1:38 & \begin{tabularx}{0.7\textwidth}{X} 耶穌轉過身來，看見他們跟著，就問他們說：「你們要甚麼？」他們對他說：「拉比，你在哪裡住？」（「拉比」翻出來就是老師。） \end{tabularx} \\ \\ \relax
1:39 & \begin{tabularx}{0.7\textwidth}{X} 耶穌說：「你們來看。」他們就去看他在哪裡住。這一天他們就跟他同住；那時大約是下午四點鐘。 \end{tabularx} \\ \\
[1ex]
\hline
\hline
\end{longtable}
講題:成長的階梯　Stages of Growth

經文:約翰福音1:35-39

講員:鄧瑞強博士

場合:香港中文大學崇基學院禮拜堂

日期:2014年1月19日

各位弟兄姊妹,早安.在這新的一年裡,深願神的恩典顯明在你們身上.

佛教有一種講法,講出人認識真相的不同境界.「見山是山,見水是水.見山不是山,見水不是水.見山還是山,見水還是水.」

「見山是山,見水是水」,是第一境界.能見山是山,見水是水,是能按感官照應這個世界,也能明白不同「名相」之別.這是認識的開始.

「見山不是山,見水不是水」,是第二境界.能見到「山」依賴各種條件而成,見到「水」依賴各種條件而成,山與水,本來無一物,本是「空」.既然是「空」,也就無所謂「山」與「水」.同時,也明白所謂「山」,所謂「水」,只是由人的心識而來,一切由心所造,「山」與「水」只是心識所造成的不同「名相」而已.從認識現象,進而認識現象之為現象背後的構成條件,是深一步的認識.

「見山還是山,見水還是水」,是第三境界.不執著世間的「色」,也不執著「空」.如《心經》所言:「色不異空,空不異色.色即是空,空即是色.」「山」與「水」固然是「空」,但「空」也透過「山」與「水」呈現出來.「物」不離「心」,「心」也不離「物」,彼此在依存關係中.由於不執著「色」,也不執著「空」,故「見山還是山,見水還是水」.這是自覺的活在真相中,這是最深的認識.

從一種客觀的,外在之「知」,慢慢體會到「物我交融」的「體會」之知,很多時,都是信仰成長的一種標記.

今日的講道經文是:約翰福音1:35－39.

這是一個小故事,是主耶穌與跟從者的簡單對答.這些簡單的對答,透露出生命成長的階梯.

約翰福音1:35－39.

1:35　再次日,約翰同兩個門徒站在那裡.

1:36　他見耶穌行走,就說:「看哪!這是神的羔羊!」

1:37　兩個門徒聽見他的話,就跟從了耶穌.

1:38　耶穌轉過身來,看見他們跟著,就問他們說:「你們要什麼?」他們說:「拉比,在哪裡住?」(拉比翻出來就是夫子.)

1:39　耶穌說:「你們來看.」他們就去看他在哪裡住,這一天便與他同住；那時約有申正了.

(1)成長的第一階段:外在之知

1:35　再次日,約翰同兩個門徒站在那裡.

1:36　他見耶穌行走,就說:「看哪!這是神的羔羊!」

1:37　兩個門徒聽見他的話,就跟從了耶穌.

這裡有一個信仰前輩,他有兩個門徒.

這信仰前輩就是施洗約翰,他一出生,就注定成為曠野中的先知.

街道上,熙來攘往.但他一眼看出,有神聖的人物在其中.

這有點似一個武林高手,單單聽人群中的腳步聲,便能指出有超級輕功的絕頂高手.

施洗約翰對他的兩個門徒說:「看哪!這是神的羔羊!」

「看啊」:這用字提醒我們校正焦點.

我們看不清,是因為焦點模糊.人生無焦點,是找不到真理的.

施洗約翰呼籲我們校正焦點去看.

看什麼呢?看「神的羔羊」.

「羊」通常不是人的焦點,因為「羊」代表弱者.

在香港,最威水的地方,如:匯豐銀行,門口是不會放隻羊的,他們會放獅子.這才夠威,夠煞氣.

世人看世界時,是看獅子.但施洗約翰教他的門徒看世界時,是看羊.

看到「羊」的價值是極難的.

在這個努力賺取的世界裡,看到「羊」的「犧牲」的價值是極難的.

在這個獅子與狼橫行的世界裡,看到「被殺的羔羊」的價值是極難的.

施洗約翰的偉大,在於他有一種信仰洞見,看到世人看不到的真價值.

他的那兩個門徒聽見老師的指點,便跟從了耶穌.

一般而言,「羊」很難吸引人跟從.孔武有力的人,較似「大佬」,容易吸引人跟從.若你要跟「大佬」,會否跟一個好似「羊」一樣的「大佬」?

很明顯,這兩個門徒信任老師的指點,他們願意跟從那「羔羊」.

他們對耶穌的認識,是初步的,是外在的,我稱這為「外在之知」.

有人指出真理,他們便按真理的方向走去,這是認識真理的開始.

很多人返教會,信耶穌,也是經人指點,對耶穌有一種「外在之知」.我們接受別人的指導,便跟從了耶穌.這是信仰的開始.

(2)成長的第二階段:內省之知

1:38　耶穌轉過身來,看見他們跟著,就問他們說:「你們要什麼?」

兩個門徒跟著耶穌,耶穌停下腳步,轉過身來,望向他們.

主耶穌向他們問了一個關鍵的問題.主耶穌問他們:「你們要什麼?」

這句話是主耶穌在約翰福音裡講的第一句話.

在約翰福音的神學裡,耶穌是「神的道」,「神的說話」.耶穌開口講的第一句話,應有特殊的神學用意.

主耶穌第一句說的話,是問跟從他的人,「你們要什麼?」

這裡的重點,不在我們認識的「外在」對象究竟是什麼?重點在於追問,我們「內裡」想在這些外在對象身上得到什麼?這不是「外在之知」的追問,而是一種內心的探索,內向的反省.

哲學有一門分支,叫「知識論」,專門探討有關「知識」的各種問題.較早階段的「知識論」,探討我們如何將一外在對象變成我們的知識.但自啟蒙運動的哲學家康德(Kant)後,我們便意識到我們人性的內在條件,會決定我們如何去構作我們的知識.簡單說,用一個花形的「模」去造曲奇餅,造出來的曲奇餅就是花形的.我們人性裡面其實有一建構知識的「模」,用這「模」建構出來的知識,其實都有這個「模」的形狀.自康德起,我們不單討論「外在之知」,也內向地反省我們內裡的條件如何建構我們對外物的認識.

這個認識的「模」,就是我們建構知識的「範疇」(categories).後來的知識論,進一步顯示,我們的語言本身也是一個「模」.以不同語言去表述同一事物,也會構成不同的認識.一個語言,造就一個世界.

再進一步,我們發覺我們對事物的不同興趣(interests),也會將事物變形,以套在我們的興趣的「模」裡.舉例說,當看見一隻狗時,愛好吃狗肉的人,會看見食物；做狗隻配種生意的人,會看見商機；尋找寵物的人,會看見可以作伴的生命.

你尋找什麼?你要什麼?會影響我們對外物的認識.

在這裡,關鍵不是一隻外在的狗,而是我們內心想以這隻狗來作什麼?

主耶穌問那兩個跟從者的問題:「你們要什麼?」重點不在其「外在之知」,而在其「內省之知」.他們是否知道他們跟從主的內在目的?

各位弟兄姊妹,你們要什麼?

一個追求成功的人,會將耶穌塑造成一個賜人成功的主,去滿足其渴望.

一個追求快樂的人,會將耶穌了解為一個賜人快樂的主,去滿足其追求.

每一個人,看來都因著自己的渴望,塑造一個合乎自己心意的耶穌.

主耶穌問:「你們要什麼?」

這是一個很好的提問,讓我們去反省,我們有什麼人生目標?我們要追尋什麼樣的價值?我們渴望活出什麼意義?

看來,主耶穌要我們正視,我們是如何理解他的.主耶穌要我們正視,我們是懷著什麼目的去跟從他的.主耶穌要我們正視,我們如何建構我們對他的認識.主耶穌要求我們,有一種「內省之知」.有了這種「內省之知」,我們才會明白,我們對耶穌的認識,多麼受我們自身限制.明白這種限制,我們才有機會從這限制走出來.這是一種「自知之明」,這是信仰成長的第二步.

(3)成長的第三階段:體會之知

1:38　耶穌轉過身來,看見他們跟著,就問他們說:「你們要什麼?」他們說:「拉比,在哪裡住?」(拉比翻出來就是夫子.)

1:39　耶穌說:「你們來看.」他們就去看他在哪裡住,這一天便與他同住；那時約有申正了.

主耶穌問那兩個跟從者,「你們要什麼?」

那兩人答得也妙.他們說:「拉比,在哪裡住?」

主耶穌邀請他們去看.

他們看後,便當下與主耶穌同住.

「住」代表生命的同在,親近,結連.

「住」代表生命的交流,生命的開放,生命進入愛的關係中而放下自我.

約翰福音在另外兩個重要地方提到「住」,雖然《和合本》的中文翻譯沒有顯示出來,但「住」的意思是明確的.

第一處經文是約翰福音6:56,是耶穌的話.《呂振中譯本》這樣翻譯:「嚼我肉,喝我血的人住在我裡面,我也住在他裡面.」耶穌在這裡講到自己是生命之糧.以他為糧的人,生命與他結連.這人住在耶穌裡,耶穌也住在這人裡.

第二處經文是約翰福音15:5,也是耶穌的話.《呂振中譯本》這樣翻譯:「我是葡萄樹,你們是枝子.住在我裡面,而我住在他裡面的,這種人才能多結果子；因為離了我,你們就不能作什麼.」

明顯,約翰福音有一套「住」的神學.信徒與基督同住,就是信徒住在基督裡,基督住在信徒裡.在這裡,信徒不再是抽離的,旁觀的按自己的框架去理解耶穌.當耶穌住在我們裡面時,耶穌會挑戰我們內在的理解框架,會改變我們的理解框架,以致活著的不再是我,乃是基督在我裡面活著.當基督在我們裡面活著的時候,當我們活著就是基督時,便不是我們帶著種種塵世的目的去理解基督,而是基督透過我們「如其所如」地呈顯自己.

或許,現代社會的問題,是每個人都是一個密不透風的「自我」(ego),一個排外的「主體」(subject).每個人都有一種凌駕他人的觀點,每個人都要滿足自身要求的目的.你若問:「你們要什麼?」十個人就有十種不同的要求,並且各不相讓.

當我們願意讓基督住在我們裡面時,我們才可放下我們裡面那獨大的「自我」.當我們願意住在基督裡面時,我們便可放下塵世的種種目的,油然地回到愛的家園裡.

在希伯來人的文化裡,「認識」和「性的連合」是同一個字.對他們來說,要有「生命的連合」,才算是真正的「認識」.以弗所書4:15說,《和合本》的翻譯是:「用愛心說誠實話」.這經文其實也可理解成:「在愛裡活出真理」.原來,真正要認識真理,是要透過愛的連繫去將真理活出來.當我們的生命與基督結連時,當我們與他人在基督的愛裡結連時,真理就在我們的生活裡流露出來.我們不再對基督只有「外在之知」,也不再懷著自己的目的去塑造一個我們想要的基督,而是在與基督的生命連合裡,真實地體會基督的生命,也讓基督的生命油然地在我們的生命中流露出來.這是對真理的最高認識,這是真正明白基督,這也是信徒成長的最高階段.

耶穌邀請那兩個門徒去「看」.

你去「看」,才能進入基督的故事裡.

門徒「看」了,當下便與基督同住.

你又是否願意與基督同住,讓基督活在你裡面,而你又活在基督裡面呢?

但願:榮耀歸於聖父,聖子,聖靈.阿們.



\newpage

\section{誰能住你的聖山}
\label{sec:234}
\textbf{伍渭文牧師 - 誰能住你的聖山?}
\newline
\newline
\hyperref[sec:233]{< < < PREV SERMON < < <}
~
\hyperref[sec:toc]{[返目錄]}
~
\hyperref[sec:235]{> > > NEXT SERMON > > >}


講題:誰能住在你的聖山?　Who May Live on Your Holy Hill?

詩篇:詩篇15

講員:伍渭文牧師

場合:香港中文大學崇基學院禮拜堂

日期:2014年2月2日

今天是年初三,俗稱赤口,謂見面時容易引起口角,不宜拜年,是求神問卜最好時刻.每年大初三,沙田車公廟因為交通方便,香火最盛,不少人為著能夠第一個進入廟堂上第一柱香而輪候通宵.

但有人為表虔誠,新春期間特別到遠離市區的大嶼山寶蓮寺祈福；認為愈遠離塵囂的廟宇,愈是靈驗,神聖就是脫俗出凡.不過自從有架空觀光纜車接駁的寶蓮寺成為旅遊景點後,現在也成為市集了.我在山東省泰山觀光時,就見過坡度陡峭的「天梯」,導遊告訴我們,有人半行半跪上山參拜祈福.

今天詩篇經課也提到參拜上主.「耶和華啊:誰能寄居你的帳幕?誰能住在你的聖山?」(詩十五1) 帳幕指在曠野飄流時期的會幕,聖山指錫安山,在錫安山聖殿雄立. 跟著幾節經文列出參拜上主的十項資格:

積極來說:(1) 行為正直,(2)作事公義,(3)心裡說實話.(二節)

消極來說:(4)不以舌頭讒謗人,(5)不惡待朋友,(6)不隨伙毀謗鄰里.(三節)

第七,八項提及其社交朋友:「眼中藐視匪類」是遠離惡人」,「尊重那敬畏耶和華的人」是樂交看重屬靈事情的朋友,見賢思齊.(四節)

第九和第十項最為具體:不能乘人之危高息借貸,不能受賄賂以害無辜.(五節)

以上十項都涉及個人操守和人際關係.參拜上主不在遠離塵囂,甚至苦待自己表示虔誠.詩篇列出的十點要求,正是雅各書所教導:真正的敬虔乃是看顧患難中的孤兒寡婦,並保守自己不沾染世俗.這些人才能住上主的聖山,才能參拜上主.

進殿禮儀或訓育手冊

格式:研究詩篇的學者對詩十五篇有兩種看法.從格式上看,這詩似是崇拜進堂禮儀.當敬拜者來到神聖空間的門檻,祭司以問答方式要敬拜者說出誰能進入這神聖空間,敬拜的人一口氣唸出十項資格,然後祭司說:「行這事的人必永不動搖.」(五節下) 詩的答問格式儼然是進入聖殿的始禮儀節.

內容:但亦有學者指出,從內容來看,是一首智慧詩,像詩篇第一篇談及甚麼是蒙福的人生.全詩列舉十項行事為人指引,積極消極並舉,正面反面參證,重複叮嚀,是訓育手冊.家長自幼以此教導孩童,要正心誠意,說話真實,行事公義,遠離惡行,扶貧濟世,才能敬拜上主.詩篇的十項要求,小孩子可用雙手十隻手指逐項細數,十是完全的數目,像完整的十誡.其實,這詩可作訓育,亦是進堂禮儀.

在基督信仰,崇拜禮儀和倫理訓誨不可分割.(註 )在崇拜用以公禱的詩篇,不少就是智慧詩篇,教導信徒如何行事為人.我們每主日誦讀的經課,書信部分就更多行事為人的教導:夫婦如何相待,主僕如何相處,要視工作為召命…….進入聖殿的神聖空間,首先躬身反思和人的關係,先與人和好,才能到聖殿獻禮物.(太五23, 24).

詩篇十五篇應該是進堂禮儀,同時也是教導孩童成為敬拜者的德育手冊,積極消極,正面反面,重複教導,如何才得進入上主的臉光中,敬拜上主.

進入新年,永不動搖

今天是馬年第一個主日,我們進入了新的一年,這篇進堂詩篇是極好的提醒和安慰,因為它給我們一個應許-永不動搖.站在聖殿門口的祭司在詩篇開始時發問:「誰能住在你的聖山?」敬拜者跟著唸出十項條件,聽畢祭司在詩篇結束時就祝福應許:「行這事的人必永不動搖.」(詩十五5)

在新春期間,世人都從俗到廟宇祈福,因為對前路感到沒有把握,誠惶誠恐.年紀愈大的人,經歷愈多,更加明白到生命的脆弱和限制.有一年,一位平時鮮談信仰的大學同事告訴我,他初三第一次到車公廟上香祈福.我聽後甚驚訝,因為平時的他講證據,強調凡事要有數據,要科學分析才能作評論.該年大年初三前幾個月,家人因重病離開,自己又血壓突然高昇,感到人生實在不可恃.他說:我所尊敬的老師以前跟我說過,當你年紀愈大,經歷愈多時,你就會愈相信命數了.謀事在人,成事在天,不由你不信.不單年長的如此相信,不少年輕人也相信.最近一位畢業不久的同學,也特別回來跟我約談,問我從信仰角度如何瞭誡命數,他愈來愈發現,人生冥冥中像有主宰.

當然我們希望新的一年風調雨順,萬事勝意,但實際很多時候事與願違.迷信的人,祈福抽到下下籤時怎麼辦?感謝上主,詩篇在終結時給我們應許:「行這事的人必永不動搖」,沒有任何事叫我們跌倒.但這句話也指出我們會遇到衝擊,產生震蕩,甚至危危欲墜,但上主保守我們永不動搖.因此我們可以昂首進入新的一年,不會患得患失.不過,上主的保守有一個前設:行這事的人.現在讓我們看看他是一個怎樣的人.

行為正直

首先他是行為正直的人,英文NIV是用whose walk is blameless-行為無可指摘,指一個人的品格,他的行事為人.這人的品格可從其外在表現看出來-行事公義,說話誠實,但看得到的行事公義和說話誠實,來自更深層的內在看不到的清潔的心-心裡說實話.跟著詩人用三個消極的說法反複論說:行事公義的人,不會惡待朋友；說話誠實的人,不會以舌頭讒謗人,亦不會隨伙毀謗鄰里.

以舌頭讒謗人是主動,隨伙是被動共謀.有時我們不敢拿起第一塊石頭,但也不敢力排眾議,制止是非的傳播,為弟兄的名譽申訴,我們太多時候屈從群眾的壓力,保持緘默,隨伙毀謗鄰里.在上主眼中,主犯和共謀都沒有說實話.在波士頓有一紀念碑,紀念在納粹時期在集中營遭毒氣殺害的,有一段發人深省的刻文,指出在不公義面前保持緘默就是共謀,就是隨伙毀謗鄰里,作者馬丁牧師寫道:

最初他們抓猶太人,我不說話,因為我不是猶太人；

然後他們抓工會會員,我不說話,因為我不是工會會員；

後來他們抓新教徒,我不說話,因為我不是新教徒；

最後他們向我走過來,已經沒人為我說話了.

4.    追求良伴

人看見的行為公義,說話誠實,來自人看不見的正直內心.清潔的良心,如同泉眼,自然湧出清潔的水.雅各書的作者談到舌頭的毒害,我們不能同時讚美上主又咒詛人,就用泉眼的比喻:「泉源從一個泉眼裡能發出甜苦兩樣的水嗎?」(雅三11)

但如何能保持顆一清潔的良心呢?就是「眼中藐視匪類」,遠離惡人」；「尊重那敬畏耶和華的人.」(四節) 惡人或敬畏耶和華的人,代表和我們交往的朋友或事物,我們吸收的價值觀念.我們看些甚麼書,結識甚麼人,誰是我們的英雄,也會影響我們成為怎樣的人.

滕近輝牧師安息了.大家可能都知道備受信徒愛戴的滕近輝牧師最近安息主懷.過去星期一我和師母出席了中國神學研究院紀念他的感恩崇拜,我在「中神」當學生和教授時他是院長.席中多人述說滕牧師如何影響他的生命,藉著聽他講道,看他的寫作,改變觀念.他應允義務作中國神學研究院的院長,使人看重神學研究.他到印尼宣教,推動普世宣教運動.他注重祈禱,躬身力行,每天早上五時為教會祈禱.他追求和睦,默默工作,雖成為國際知名講員,但行事低調,不張揚,不論斷他人.

我們食甚麼食物,就有甚麼的身體；吸收甚麼觀念,就形成甚麼的品格.我們可以藉著一個人的行為和言語,窺探到他們看過的書和所交的朋友.遠離惡人,親近尊重敬畏上主的人,幫助我們塑造正直的品格.其實這也孟母為幼兒孟子三遷的原因.我們今天會三遷入名校校區,但教會忽略孩子看甚麼電視節目,也沒有著意培育他們的閱讀習慣和內容.其實,觀念或價值慢慢模塑我們的品格.

電影鐵娘子(The Iron Lady) 有一段對白對輕視理念,對隨從起伏多變的感覺,屈從民粹的從政者,確是警語箴言,對崇尚感覺的年青一代甚有啟發.

戴卓爾夫人落任後常常懷念亡夫,甚至幻想他仍在家中,有一天醫生到診.

醫生:傷痛是非常自然的狀況.

戴卓爾夫人:我的丈夫離世多年,是癌病.

醫生:它仍然帶給你一些迷惘,你仍然被情感所困.

戴卓爾夫人:我仍然被情感所困?人現在不再思考了,他們憑感覺.

「啊,對那些我感到不安.」「啊,我非常抱歉,我們這班人覺得…….」

你知否,現今最大的問題是我們被一班關注感覺多於思想和概念的人管治.

思想和概念,引起我的興趣.問我想些甚麼,戴卓爾夫人你正想些甚麼呢?

小心你的思想,因它成為你的言語.小心你的言語,因它成為你的行為.

小心你的行為,因它成為你的習慣.小心你的習慣,因它成為你的品格.

小心你的品格,因它成為你的命運.我們如何思想,就會成為如何的人.

我父親常常以此教導我.

在新的一年, 讓我們追求敬畏上帝的思想,遠離令我們偏離德行的影響.若能養育自己的靈魂,我們就有免疫能力,抵禦貪婪,借貸給有需要的人而不會收取高息,更不會為錢財而受賄賂以害無辜,因為物必先腐而後蟲生.箴言教導我們:「你要保守你的心,勝過保守一切,因為一生的果效,是從心發出.」(箴四23)

我一月開始御任校牧職務,轉到崇基神學院兼課,有一天我專心備課,整天在家書房埋首伏案閱讀,專注思考,沒有開電視,沒有聽收音廣播,沒有開手機電話,沒有看鐘,非常安靜專注,直到日落天黑,起來進食.突然感到思想的洗滌,心靈的滿足.我有所感悟,若不經歷三天埋在地下,就沒有復活.平時起居作息,我多是跟時鐘推移,整點的新聞報告不會放過；隨身智能手機,讓電郵和短訊即看即覆,保持和外界的即時(real time) 接觸.但這次埋首「無待」,我為主場的深入閱讀,不被外在零碎的資訊分心,使我心曠神怡,我認為這才是真實時間(real time)-深入細致,沈澱情緒,過濾思考,面對自我的真實時間.

我想,大家也可以在新的一年每一星期,抽出時間,關掉電視機,關掉手機,拿起一本書,一本讓你更敬畏上主的書,好好深入閱讀,充實自己的靈命.

上主引領同行

耶穌和門徒上了山.誰能住在你的聖山?這進堂詩篇和進入禮儀是律法還是福音?當祭司問誰能住上主的聖山時是攔阻我們,還是邀請我們?這十項要求使我們灰心絕望,還是使我們感恩?

請注意今天的福音經課-八福.「耶穌看見這許多人,就上了山,既己坐下,門徒到他跟前來.他就開口教訓他們.」(太五1)是耶穌帶領門徒上了山,祇有一位能住在上主的聖山,就是耶穌.他行為正直,作事公義,心裡說實話,他不以舌頭讒謗人,不惡待朋友,不隨伙毀謗鄰里.我們在耶穌裡,也一同登上聖山.八福是豐盛生命旅程的邀請,八福之始就提及「虛心的人有福了」或譯「自覺靈命貧窮的人有福了」.因為認識自己的缺乏,進而依靠耶穌.

升天主日的詩篇.根據教會年而編寫的經課,詩篇十五篇是升天日必誦讀的詩篇,每年升天主日的晚禱的詩篇,就是詩篇十五篇.耶穌為我們受難,埋葬,復活,升天,進入(enter)父神的榮耀.而每主日我們也奉耶穌的名,藉耶穌的拯救,進入天父面前敬拜他.

結論:詩篇十五篇是進堂詩Entrance,是進入神聖空間的進入禮儀.今天我們也進入新的一年,是上主給我們的日子,因為來自上主,也是進入神聖的境域.我們並不孤單,因為基督與我們同行.靠著基督給我們的力量,我們也可以行為正直,作事公義,心裡說實話.因為遠離惡事,親近敬畏上主的,我們也有力量,對困苦者施濟,不受賄賂以害無辜.

這樣人無論遇到任何風浪,蒙上主保守,永不動搖.阿們.

註:在基督信仰,崇拜禮儀和倫理訓誨不可分割.

崇拜是與上主交談對話,起點是人能明白的類比,用人世間可瞭解的事情來論述上主；採用類比才能明白上主；終點是人不能明白的隱密.但在整個崇拜過程,悟性和隱密並未分割,有別於奧圖Otto, Rudolf的宗教經驗論述.奧圖強調,在敬拜中,悟性和非悟性互不交集.他在1917年出版討論崇拜的經典著作《神聖的概念》(English translation, 1923, The Holy - On the Irrational in the Idea of the Divine and its Relation to the Rational) ,探討當人進入神聖空間,與聖者相遇時的象外經驗(Numinous).他說當我們進入神聖空間,與那聖者相遇時,就會產生不能言說的神秘感(mysterium),但同時感到渺小,引發敬畏驚懼(tremendum) 和欣悅狂喜(fascinans) .

奧圖研究神聖的觀念,基於一般宗教,一般宗教的神聖都和凡俗割裂的:空間的佈置,儀節的鋪排,主禮者的服飾,其口中唸唸有詞,都使參與禮拜的人,廁身在一個和凡俗,慣常割裂的環境,催生神聖感.此刻,我們遇見那超越悟性和德性的聖者.奧圖的書英文直譯就是「神聖-上主概念的非悟性面向和與悟性的關係」.

然而,誠如一位作者所說:若那聖者在心靈留下痕跡,接受者可以明白和作回應.」The divinest goodness can make no impression on a human soul, unless the soul is able to express that same goodness in  response (Clarke, Liturgy and Worship, p. 15) .談到敬拜的事情,耶穌豈不是對那撒瑪利亞婦人說:「你們所拜的,你們不知道；我們所拜的,我們知道.」(約四22) 悟性在敬拜中並未隱退,因而敬拜的人,可以明白上主對我們的倫理要求,甘心順服,接受差遣,回應世界的需要.在聖殿的神聖空間,以賽亞回應上主:「我在這裡,請差遣我.」(賽六8)



\newpage

\section{活出鹽與光的生命}
\label{sec:235}
\textbf{許開明牧師 - 活出鹽與光的生命}
\newline
\newline
\hyperref[sec:234]{< < < PREV SERMON < < <}
~
\hyperref[sec:toc]{[返目錄]}
~
\hyperref[sec:236]{> > > NEXT SERMON > > >}


馬太福音 5:13-20
\newline
\begin{longtable}{cl}
\hline
\hline
章節 & 經文 (和合本修訂版)\\
\hline
5:13 & \begin{tabularx}{0.7\textwidth}{X} 「你們是地上的鹽。鹽若失了味，怎能叫它再鹹呢？它不再有用，只好被丟在外面，任人踐踏。 \end{tabularx} \\ \\ \relax
5:14 & \begin{tabularx}{0.7\textwidth}{X} 你們是世上的光。城造在山上是不能隱藏的。 \end{tabularx} \\ \\ \relax
5:15 & \begin{tabularx}{0.7\textwidth}{X} 人點燈，不放在斗底下，而是放在燈臺上，就照亮一家的人。 \end{tabularx} \\ \\ \relax
5:16 & \begin{tabularx}{0.7\textwidth}{X} 你們的光也要這樣照在人前，叫他們看見你們的好行為，把榮耀歸給你們在天上的父。」 \end{tabularx} \\ \\ \relax
5:17 & \begin{tabularx}{0.7\textwidth}{X} 「不要以為我來是要廢掉律法和先知。我來不是要廢掉，而是要成全。 \end{tabularx} \\ \\ \relax
5:18 & \begin{tabularx}{0.7\textwidth}{X} 我實在告訴你們，就是到天地都廢去，律法的一點一畫也不能廢去，直到一切都實現。 \end{tabularx} \\ \\ \relax
5:19 & \begin{tabularx}{0.7\textwidth}{X} 所以，無論誰廢掉這誡命中最小的一條，又教導人也這樣做，他在天國裡要稱為最小的。但無論誰遵行並如此教導人的，他在天國裡要稱為大。 \end{tabularx} \\ \\ \relax
5:20 & \begin{tabularx}{0.7\textwidth}{X} 我告訴你們，你們的義若不勝過文士和法利賽人的義，絕不能進天國。」 \end{tabularx} \\ \\
[1ex]
\hline
\hline
\end{longtable}
講題:發出鹽與光的生命　Living out a Life of Salt and Light

經文:馬太福音5:13-20

講員:許開明牧師

場合:香港中文大學崇基學院禮拜堂

日期:2014年2月9日

祈禱:

引言

今日是農曆年初十,仍是充滿新春的氣氛.

恭祝大家新年如意,身體健康,福杯滿溢!

本段經文大要分析

今日的經課是屬於耶穌最著名的講章,名為登山寶訓.

最初一篇,開始的一篇,是怎樣做個有福的人?

耶穌連續講了八個褔,也有認為是九個福.

這八個褔,虛心的人,哀動的入,溫柔的人第均與基督徒品格有關.

本段兩大重點:

第一個:跟著耶穌提到門徒要努力成為鹽和光.耶穌一定想過只是八福對某些人來說,還未夠的.於是提到兩大方面.

鹽是內心的修維.

光是外在的表現.

孔子有言「內聖外王」的意思,不謀而合.

第二個:再而提到成為鹽和光還不足夠,仍需守律法.

耶穌強調,他來不是廢除律法,而是成全律法,並要勝過律法的義.

一門徒為鹽

【】中國人讚人愛說「你個樣好甜!」

不會說:「你個樣好咸,或你好咸!」

對廣東人來說,好咸,不是順耳的,也另有別意的.

但是耶穌卻鼓勵門徒要做「鹽,要夠咸」!

他另有真義在其中.

在古代的世界,鹽和光具有極崇高的價值.羅馬人說:『沒有別的東西比太陽和鹽更有用處.』(Nil untilius sole et sale).

希臘字稱鹽為:divine(theion),神性.

有一句諺語,以拉丁文來念是一個押韻的句子,

在耶穌那個時代,在人們的心目中,鹽與三種品質有關:

(一)保持純潔的作用.鹽與純潔有關.它白色的光彩,無疑地使這聯係很容易建立.羅馬人說:鹽是萬物中最潔淨的,因它是從最純淨的兩種東西──太陽與海水而來.

鹽的確是首先奉獻給上帝的祭物,猶太人在一日的末了獻鹽為祭.因此基督徒若要作世上的鹽,他必須是純潔的榜樣.

【】今日世界給我們的挑戰是請把標準低落甚至

「阿sor ,隻眼開隻眼閉」

我們生活的世界有一個特徵,就是標準低落:

人格的標準:

誠實的標準,

工作勤奮的標準,

正直的標準,

道德的標準,都在走下坡；

基督徒必須維持崇高的標準,在言語上,行為上,甚至思想上,保持絕對的純潔.

(二)鹽防腐的作用.古代的世界,鹽是最普通的防腐劑,用以保存東西免致腐壞,排除腐爛與朽壞.

蒲魯他克(Plutarch)用一種奇怪的方式來說明鹽的功用,他說:肉就是死去的身體或屍體的一部分,如果聽憑自然,定將歸於腐朽；但鹽好像一個新的靈魂介入一個死去的身體一般,保持了它的新鮮,使之免於腐壞.

保持咸味,不是目的,而是保持了它的新鮮,有魚味,菜味,肉味.

鹽保存了食物,使之免於腐朽.基督徒要作世上的鹽,也就必須對生命有防腐性的影響力,保持基督徒的味道.

我們都知道近朱赤,近墨者黑的道理.跟某些人在一起易於為善,跟另一些人在一起易於使標準鬆懈.

因此,基督徒需要經常聚在一起崇拜,團契,彼此的守望.

(三)鹽調味的作用.鹽最廣大最明顯的特質,是可以調味.沒有鹽的食物就沒有滋味,甚至成了令人生厭的東西.

基督教對於生命,正如鹽對於食物一般.可悲的是人們常把基督教與此相反的作用相提並論.

V13.　鹽的主要用途是調味和防腐,門徒如果真正應了神的召喚,就能將世界變得更為純潔,美好.可是,若不能保存他們與眾不同的特質,他們還能做什麼呢?

鹽不獨能防腐,保持魚的鮮味與魚味.

更能把鹽放在魚中,肉中,菜中,便能調味,使食物更有味道!

【】有味的弟兄

在公元二千年,中華基督教會差派我往英國伯明翰大學進修.我們一家五口前往,去到異地遇上許多的困難.住屋給人騷擾,孩子讀書給人欺侮,駕車又不知所蹤?在一次講道中,認識一位教會的李弟兄,關心我們一家給我們許多指引和幫助.

有人情味的弟兄,他也如此關顧留學生,引領多人信主.

注意:鹽是會失去鹹味!

【太五13】「你們是世上的鹽；鹽若失了味,怎能叫它再鹹呢?以後無用,不過丟在外面,被人踐踏了.」

鹽沒有鹹味就失去了價值.嚴格說來,純鹽是不可能失去鹹味的,可是從死海灘上掘出的不純「鹽」,一旦氯化鈉分解了,慢慢也就不鹹了.

古時猶太地方和中東各地吃的鹽,多是『礦鹽』或『井鹽』,那些鹽塊從礦場或井裡挖出來時,就像拳頭大的石頭一般.在這石塊的外層包著一層鹽.

猶太人在吃飯時,右手拿著菜,左手握著鹽塊,用舌頭舐鹽以取味.家中每個人都有一塊鹽放在廚房中,吃飯時就各用自己的鹽塊.等到鹽塊的鹹味逐漸減少,至完全失味時,原先的鹽塊就成為石頭了,於是隨手棄置,這些鹽份用盡的石頭就任人踐踏了.

在古老的傳統中,如果猶太人背叛了信仰,然後再回頭時,在蒙接納回到會堂之前,必須躺在會堂門口,請走進會堂的人用腳踩在他身上.有些基督教團體也繼承了這一個傳統,凡被教會法規驅逐的基督徒,在他蒙接納歸回以前,也要被迫躺在教堂門口,對走進來的人說:『踐踏在我身上,因我是那失了味的鹽.』

給入踐的感覺,一定是不好受的,自願給入踐,或者不怕給入踐,不會是容易的事.接受自己犯錯後,應受罰.

許多人給人指責,可以知錯,但往往想避開受罰.或者明知做錯,卻不肯認錯,因為認錯,就要受罰.

【】菲律賓總統不認錯,因認錯,要道歉,就要受罰,要賠償.

【】日本不承認南京大屠殺,因認錯,要道歉,就要受罰,要賠償.

【】六四事件為何不可平反,你我也明白.

有鹽的信徒也應是敢於認錯,失去鹽味,接受踐踏的.

1不失去純潔的理想

2不失去調和作用

3不失去防腐作用

4不失去味道,失去智慧

猶太的拉比常用鹽來比喻智慧(參,西四6),所以希臘語代表失去味道的字實際上又有「變得愚蠢」的意思(耶穌無疑用的正是亞蘭語裡有這兩個意思的詞 ta-pe-l).一個愚蠢的門徒對世界是沒有影響的.

→基督徒要小心,你會失去鹽味,咸味的.

你會失去智慧的,因此要小心.

【】多年前,我應紀律部隊團契的邀請講道.

講題是「弟兄跌倒了,怎麼辦?」

我是第一次應邀,記得一位認識的,非常熱心的弟兄,很想見到他.可惜,走到團契見不到,非常失望!

講完道後,於是我問請我的弟兄,為何不見某某弟兄?

那弟兄靜靜告訴我,「許牧師,實不相瞞,某某弟兄也跌倒了!你這篇也是適合他的,可惜他不來?」

後來,我知道這弟兄跌倒,連信仰也失去,我感到很難過,一位曾經熱心愛主的弟兄,也中招!

一位打擊罪惡的人,竟然為罪惡所勝!

我們是失去鹽,失去咸味的!我也提醒自己!

三 門徒為光

14-16.　光,像鹽一樣,因為與暗區隔,才能對環境有所影響.門徒只有明顯的與眾不同,才能影響周遭的人.

但門徒為光,不是自己發光,不是發自己的光.

門徒好行為的動機,不是為了炫耀自己的優良品質,而是為了把眾人的注意力引到召喚他們的神這裡來.只有這樣,門徒才能照亮所有的人(參,腓二15).

光的作用,不是光照自己,甚至不是光照他人；而是引導他人,去就真光(耶穌).

讓基督的光,透過我們的生命,所言所行而顯明出來.

【】基督徒是光,我們要問:我所言所行的,會否令自己,別人跌倒,造成對別人的傷害呢?

門徒的光來自基督的光

耶穌是世界的光(約八12),是以賽亞曾預言的神的僕人(賽四十二6,四十九6),但是這個任務已交給了祂的門徒(參,徒十三47).在談光的時候,突然提出造在山上的城,顯然有些不自然,其實是為了強調與世明顯不同的重要性107.

世上的光(五14,15)

耶穌說:『我在世上的時候,是世上的光』(約九3).

祂吩咐祂的跟隨者要作世上的光,祂所要求的,是他們應當像祂一樣.

耶穌所採用的這種表達方式,正是初次聽見的猶太人十分熟悉的表達方式.

以色列人,自稱耶路撒冷是『外邦人的光』,有一位著名的拉比被稱為:『以色列的燈』.因著猶太人所使用的這種表達方式,使我們可以知道耶穌採用它的關鍵.

有一件事是猶太人十分清楚的,沒有人能夠點燃他自己的燈.

耶路撒冷的確是外邦人的光,但是上帝點燃了以色列的燈.無論是國家或上帝的僕人所發出來的光,都是借來的.

基督徒的光也是如此.耶穌所要求我們的,並不是我們應當發出我們自己的光來.我們必須發出反映祂的光輝.基督徒所發的光,是因為基督他心裡而映照出來的亮光.我們常說『容光煥發的新娘』,她所發出的光輝,是從她內心的愛所產生的.

耶穌說基督徒要作世上的光,這句話是甚麼意思呢?

(一)光輝最主要的是被人看見.

【】巴勒斯坦的房屋十分黑暗,只有一個小小圓形的窗戶,可能直徑只有十八英吋,燈是船形,好像一隻醬油碟子那樣,裡面盛著油,浮著燈蕊.

在火柴尚未發明的年代,要重燃一盞燈並非易事.

那時通常燈是放在一根粗糙的木頭作成的燈架上,人們外出的時候,為安全起見,把燈從架子上拿下來,放在陶土做成的蒲式耳底下(每蒲式耳約合中國的三斗餘,為東方人的一種量穀的容器──譯註),使點燃的燈不致在他們回來以前發生危險.燈光原來的功用是給人看見.

因此基督教就是要給人看見,有人說得好:『不可能有秘密門徒的身分存在,因為不是秘密破壞了門徒的身份,就是門徒的身份破壞了秘密.』一個人的基督信仰應該給所有的人看得見.

【】盲人點燈,為何?你看不見?

因為我要讓人看見,使人不會撞倒!.

【】基督徒是光,我們要問:我所言所行的,會否令自己,別人跌倒,造成對別人的傷害呢?

我們怎樣對待店舖裡的伙記,怎樣在飯店裡點菜,怎樣對待我們的僱工或服侍我們的僱主,怎樣玩遊戲或駕駛汽車,還有我們日常所說的言語及所念的書籍；基督徒無論論在工廠,商店,造船廠,礦場,學校,手術室,廚房,高爾夫球場,都要好像在教堂一樣.耶穌並沒有說:『你們是教堂的光……』

祂是說:『你們是世上的光』.人生存在世界上一天,他的信仰就應當向所有的人證實出來.

(二)光是一種指引.在任何江河的港口,我們都可以看見有一列燈光在水道上作記號,使船隻得以安全的駛入,即使在城市的街道上,若沒有燈光,我們都知道將會遭遇多大的困難.光可以使路途清晰.

(三)光時常可以當作警戒主的記號.當前面有危險的時候,它警告我們停止.

有時候基督徒的責任就是把必需的警戒帶給他的同胞.這常是很困難的,要做到使對方得益處而不給他少許的傷害,尤其不容易.可是人生中最慘痛的悲劇是一個人,尤其是一個青年人,來對我們說:『如果及時勸阻,我決不會落到今日的地步.』

一位老師與校長奧爾雄(Florence Allshorn)說只要她有機會責備學生,她總沒有忘記同時『我用口責備他們的不是,同時用手臂環抱他們』.如果我們所給予的警戒不是出於怒氣,憤懣,批評,定罪或有意加以傷害；而是出於愛心,一定會生效力.

能夠使人看得見的光,警戒人的光,引導人的光,就是基督徒所必須作的光.

為上帝發光(五16)

這裡有兩件最重要的事:

(一)讓人們看見我們的好行為.在希臘文中有兩個字代表『好』,agathos單單是指一件東西有良好的質地.

kalos的意思不單是一件東西好,而且也美麗迷人,引人入勝.這裡所用的『好』字,就是kalos.基督徒的善行,不僅是良善而已,它們必須要吸引人.

【】德蘭修女,把她大半生放在服事印度最窮困,患絕症的人,有次遇上一位非常有錢的人.那人為表示對德修的欣賞,對她說:「你就算給我一百萬,我也不會做你的工作.」

德蘭回應說:「我若為一百萬也不會做這樣的工作.」

德蘭修女一生致力讓人身上看見的,不是自己而是基督的光.

(二)更要注意的是我們的好行為所引人注意的,並不是我們自己而是上帝.

【】穆迪(D. L. Moody)參加一次聚會,有幾位對信仰十分認真的青年人,也來參加聚會.一天晚上,他們舉行了一個通宵的禱告會,第二天早上,他們正要離開的時候,遇見了穆迪.他問他們在作甚麼?

他們說明原委之後,接下去說:『穆迪先生,請看我們的臉是多麼光亮.』

穆迪很有禮貌的回答說:『摩西並不知道他的臉在發光.』自己感覺到的良善,把注意力集中在本身的良善,這不是基督徒的良善.

人總愛別人留意他的光,他對社會的貢獻,他光茫.

至於作光這一面的責任,王宣佈了光的雙重影響──城與燈.

a「城造在山上,是不能隱藏的,」

基督徒的地位如同建造在山上的一座城,是顯明在世人眼前的,所以是不能隱藏的.

「城」是由許多建築物集合而成的,故它象徵教會集體的見證,表明教會惟有建造在一起,才能大大地發揮光的功能.

祂說,「城造在山上是不能隱藏的.」「城造在山上」,是指教會作為光,去產生影響力.沒有一個基督徒能單獨造成這種影響.一個人可能成為山頂的火炬,但不能成為一座城.要集合在國度裡的一群人,纔能照亮廣袤無邊的原野.

一個座落在山上的城,人們從遙遠的山谷,就能望見它,它所發出的光輝,遠近都能看到.這幅圖畫描寫了教會的影響力,但今日教會幾乎已經完全喪失了這種影響力.

它描繪出教會的社會法則如何照耀著它的那個世代.這比喻不是隨便選的.

這座城的統治者是神,在城中沒有任何玷污,可憎,虛謊的事.城中集合了一切美麗,整齊,光亮與喜悅的事.這一切正是教會應有的光景.因此教會將守護這城的眾門,不容許任何敗壞和傷害侵入.若那城達到這境地,在它的邊境之內,社會的理想得以實現,以致一人哀哭,眾人同流淚；一人喜樂,眾人同歡笑；對貧苦的聖徒,供以今世所需；對軟弱的人,提供屬靈的幫助；沒有一人說,他的東西有一樣是自己的,他們都是凡物公用當這一切實現時,教會就成了造在山上的城,所發的光就能照耀到地極.

b光的另一層影響是燈.

王從一座城轉而論到一所房屋.「照亮一家的人.」倘若城的比喻,是說到那照亮廣大穹蒼之光,那麼燈的比喻自然是說到在個人生活中所發的光.

一個人雖然不能作造在山上的城,卻可以作照亮一家人的燈.這是光另一層的影響.一個家若有一個基督徒,全家就都得著亮光,因為在這家中有一盞為耶穌發光的燈.

「人點燈,不放在斗底下,是放在燈台上,就照亮一家的人.」不要忘記這裡所說,別放燈在斗底下的話.

麥拉連博士(Dr. Alexander Maclaren)對這一節聖經的註解道盡了它的精華.他說,「沒有人點了燈就放在斗底下,乃是放在燈台上,但是倘若有人要把燈放在斗底下,結果將如何呢?

只有兩個可能:若不是斗將燈熄滅,就是燈使斗起火焚毀.」

把斗(計糧的容器,大約盛九公升)扣在油燈上,就可能把燈熄滅,含義恐怕是,點燃燈不是只等著熄滅用的(Jeremias, PJ, pp.120-121).怎會熄滅燈?

「放在斗底下,」斗是木造的量器,用以量取糧食,故它象徵生活的掛慮；燈放在斗底下,意思是信徒被生活的問題所扣住,以致遮蔽了亮光.

信徒被自己的焦慮,憂心,懼怕所累,會失去光茫!

「是放在燈臺上,」燈台是放置燈的正當器物,它象徵見證主的生活.一個隱蔽的門徒活在世上,和一個失去特色的門徒一樣無大用處(13節).你們在天上的父是馬太很愛用的表達法(參五45,六1,9,七11等),反映著耶穌教訓裡一個主要重點.在較早期猶太人的心目中,以為神是整個以色列民族的而不是哪個人的父親,儘管這句稱謂的後一種意義到公元一世紀就已開始使用了(見 ICC, p.44).

【】持燈者成一燈路:英國作家拉斯金(John Ruskin),與一位弟兄在一起,談:「怎樣做個基督徒?」坐在家裡,望見窗外遠處,有位點燈的人,手拿著火把,在山坡上點燈.當時天色昏暗,看不見點燈的人本身,但從那些燈火可知他的存在.

拉斯金突然有感說:「做一個基督徒就是這樣,你可能看不見他,但從他的好行為,知道他的的存在.因為在他所經過的地方,都點著發燈,都發了光.」

三門徒當守律法

進入上帝的國:

太五章1至12節強調信徒只要懂得八福之道:

虛心的人,哀 慟 的 人,溫 柔 的 人,飢 渴 慕 義 的 人,憐 恤 人 的 人,清心的人,使 人 和 睦 的 人,

為 義 受 逼 迫 的 人.

人做好自己,人靠自己,人由心發出,便行了!

人不需要規範,人不要律法,誡命,可以嗎?

自古到今都有人提出,我們做好自己,一切由心發出,人能夠做好自己,根本不需要任何規矩與律例.

耶穌說:「不行!人需要律法,耶穌來,不是廢除律法而是成全律法!」

現在,讓我們把第17-20節的幾個要點歸納如下:

「我不是來摒棄舊約,而是來成全舊約所指示的.舊約的任何部分都不能摒棄,都要成全(正如我現在的傳道和教訓所行的成全那樣).

拒不接受舊約任何部分的基督徒是一個糟糕的基督徒；堅定的基督徒定要接受舊約的引領,也要教訓其他人這樣做.

但是,一個真正基督徒的態度可不同於文士和法利賽人的律法主義,而應是一種更深的奉獻,奉獻自己來行神的旨意,自第21節以下將對之進行闡述.」

律法的精華(五17-20)

耶穌對於律法,作何解釋?祂說祂來不是要廢掉律法,乃是成全它.

這是說祂來是要使律法真正的意義清楚表現出來.律法真正的意義又是甚麼呢?

整個律法後面真正的原則是甚麼?耶穌來成全的是甚麼原則?祂來向我們顯明的原則,真正的意義何在?

十誡原是一切律法的基礎和本質,我們可以在其中看到它整個的意義能夠總括在一個字裡面──『尊重』,或者更進一步的『敬畏』.

『敬畏』與「尊重」

敬畏上帝和上帝的名,敬畏上帝的日子,尊重父母,尊重生命,尊重財產,尊重人格,尊重真理與別人的好名聲,尊重自己,使我們不會被錯誤的慾望所控制── 這就是十誡的基本原則

十誡的基本原則是敬畏上帝,尊重同胞和我們自己,這也人是愛的表現..

沒有這基本的敬畏與尊重,就不可能有律法,一切的律法都是以此為基礎.

耶穌來成全的就是這種敬畏和尊重,祂在實際的生活中指示人敬畏上帝,並尊重所有的人.

敬畏與尊重並不在於遵守一大堆瑣碎的規則與條例,並不在於祭祀,卻在於憐憫；並不在於律法,卻在於愛心；

並不是在禁令中強迫人不做一些甚麼事情,而是在誡命中告訴他們,將他們的生命模塑成主動的,愛人的生命.

敬畏和尊重是十誡的基礎,是永遠也不能廢去的,它們是維護人與上帝以及人與人之間關係的永遠的要素.

門徒勝過法利賽人的義

5:20    我 告 訴 你 們 , 你 們 的 義 , 若 不 勝 於 文 士 和 法 利 賽 人 的 義 , 斷 不 能 進 天 國 .

基於「敬畏」「尊重」,和「愛」

不獨做到,且是走多一步,做得更好.

有道德智慧的人,往往都比人想多一步,看深一點,不要亂衝亂撞.(智慧的牧者,是浙江省省長的機要秘書.神學生跟實習,常跟我交談,把他的牧會智慧傳授)

5:21    你 們 聽 見 有 吩 咐 古 人 的 話 , 說 , 『 不 可 殺 人 , 』 又 說 , 『 凡 殺 人 的 , 難 免 受 審 判 . 』

5:22    只 是 我 告 訴 你 們 , 凡 向 弟 兄 動 怒 的 , 難 免 受 審 判 . 〔 有 古 卷 在 凡 字 下 添 無 緣 無 故 的 五 字 〕 凡 罵 弟 兄 是 拉 加 的 , 難 免 公 會 的 審 斷 ． 凡 罵 弟 兄 是 魔 利 的 , 難 免 地 獄 的 火 .

5:27    你 們 聽 見 有 話 說 , 『 不 可 姦 淫 . 』

5:28    只 是 我 告 訴 你 們 , 凡 看 見 婦 女 就 動 淫 念 的 , 這 人 心 裡 已 經 與 他 犯 姦 淫 了 .

5:38    你 們 聽 見 有 話 說 , 『 以 眼 還 眼 , 以 牙 還 牙 . 』

5:39    只 是 我 告 訴 你 們 , 不 要 與 惡 人 作 對 ． 有 人 打 你 的 右 臉 , 連 左 臉 也 轉 過 來 由 他 打 .

5:40    有 人 想 要 告 你 , 要 拿 你 的 裡 衣 , 連 外 衣 也 由 他 拿 去 .

基督的跟隨者,也要遵守律法,明白這是敬畏上帝,尊重人,愛人的表現.

結語

【】有次,聽一位年過九十,上帝所重用的僕人跟我們一班後輩分享牧會心得,其中一句,我十分深刻:

他說:「教會不是無我唔得,只是有我會更好!」

雖然老牧師說話有點自負,也好像有點驕傲,過份自信.

但回想到這位牧者確實在教會做過許多有意義的工作,使許多人的生命得到建立,九十多歲的長者驕傲得起.

我不敢好像這位牧者說出這句話.

然而,若果我要你鹽的話,我問自己一個問題:

「這個世界有沒有因為我的存在,變得好?」

「我的家庭,有沒有因為我的存在,變得好?」

「我的工作地方,有沒有因為我的存在,變得好?」

「我的社區,甚至社會,國家,有沒有因為我的存在,變得好?」

如果我是鹽和光,我的存在,我活著便不是可有可無.

在本段經文,耶穌教導我們要成為鹽和光,需守律法,並勝過律法的義.



\newpage

\section{心懷天道人自正}
\label{sec:236}
\textbf{鄧瑞強博士 - 心懷天道人自正}
\newline
\newline
\hyperref[sec:235]{< < < PREV SERMON < < <}
~
\hyperref[sec:toc]{[返目錄]}
~
\hyperref[sec:237]{> > > NEXT SERMON > > >}


詩篇 119:33-40
\newline
\begin{longtable}{cl}
\hline
\hline
章節 & 經文 (和合本修訂版)\\
\hline
119:33 & \begin{tabularx}{0.7\textwidth}{X} 耶和華啊，求你將你的律例指教我，我必遵守到底！ \end{tabularx} \\ \\ \relax
119:34 & \begin{tabularx}{0.7\textwidth}{X} 求你賜我悟性，我就遵守你的律法，且要一心遵守。 \end{tabularx} \\ \\ \relax
119:35 & \begin{tabularx}{0.7\textwidth}{X} 求你叫我遵行你的命令，因為這是我所喜愛的。 \end{tabularx} \\ \\ \relax
119:36 & \begin{tabularx}{0.7\textwidth}{X} 求你使我的心趨向你的法度，不趨向不義之財。 \end{tabularx} \\ \\ \relax
119:37 & \begin{tabularx}{0.7\textwidth}{X} 求你叫我轉眼不看虛假，使我活在你的道路中。 \end{tabularx} \\ \\ \relax
119:38 & \begin{tabularx}{0.7\textwidth}{X} 求你向敬畏你的僕人堅守你的話！ \end{tabularx} \\ \\ \relax
119:39 & \begin{tabularx}{0.7\textwidth}{X} 求你使我所懼怕的羞辱遠離我，因你的典章本為美。 \end{tabularx} \\ \\ \relax
119:40 & \begin{tabularx}{0.7\textwidth}{X} 看哪，我切慕你的訓詞，求你因你的公義賜我生命！ \end{tabularx} \\ \\
[1ex]
\hline
\hline
\end{longtable}
主題:心懷天道人自正　Walk Straight in the Word of God

經文:詩篇119:33-40

講員:鄧瑞強博士

場合:香港中文大學崇基學院禮拜堂

日期:2014年2月23日

各位弟兄姊妹,早安.

北宋黃庭堅說過:「士大夫三日不讀書,則義理不交於胸中,對鏡覺面目可憎,向人亦語言無味.」這句話後來簡化成:「三日不讀書,便覺語言無味,面目可憎」.

猶如我們的身體,每天都需要補充養份；我們內在的生命,也需要時時獲取智慧與啟迪.身體沒有養份,便會虛弱；內在的生命缺乏智慧,便會語言空洞,缺乏內涵.空洞的語言,猶如青蛙的叫聲,響亮卻是單一,越聽越覺得厭煩,失去一種吸引人的品質.靈修學者傅士德(Richard Foster)的《屬靈操練禮讚》(Celebration of Discipline)一書第一章第一句便說:「我們這個時代的禍因是淺薄.事事尋求立時的滿足乃是一個基本的屬靈問題.今天最迫切的需要不是要有大量聰明能幹的人,或者大有恩賜的人,乃是有深度的人.」

我想起蘇東坡詩裡的一句話:「腹有詩書氣自華」.一個人若有內涵,他的生命自會散發出生命的光采.基督徒應是散發出生命光采的人.今天我講道的經文,是詩篇119篇33-40節.詩篇119篇整篇詩篇,講出詩人的一種體會,他的體會與蘇東坡的體會相似,都是一種「腹有詩書氣自華」的體會.不過,聖經裡的詩人,他的「詩書」是神的話,他講的「氣自華」是屬神生命的光采.一個人的生命若浸淫在神的話語裡,則他的生命自有一種屬天的光華.為此,我起的講道題目是:「心懷天道人自正」.

未讀詩篇119: 33-40前,我講一講詩篇119篇的特色.

希伯來文有22個字母,詩篇119篇按這些字母分為22段,每段8句.第1段的8句詩句,開首的字母全是希伯來文的第1個字母.第2段的8句詩句,開首的字母全是希伯來文的第2個字母.餘此類推.即是說:希伯來文的第1個字母是“aleph”,這詩篇的頭8節詩句,開頭全是“aleph”.希伯來文的第2個字母是“bet”,這詩篇第9-16節的詩句,開頭全是“bet”.餘此類推.這造成一種「押頭韻」的效果,方便背誦.全詩22段,每段8節,全詩便是176節.

詩篇119篇的第二個特色,是用八個同義詞來指稱神的律法(torah),分別是:律法(law),法度(testimonies),訓詞(precepts),律例(statutes),命令(commandments),典章(ordinances),神的話(God’s word),應許(promise).而差不多每一節經文,都出現其中一個字眼.即是說:「神的律法」滲透在整篇詩篇的詩句裡,滲透在整篇詩篇的思想裡.背誦這詩篇的人,神的律法會縈繞在他的心靈中.

就這兩個特點,詩篇119篇透過其文學手法傳遞著一些重要信息.

這詩篇有意識地順序使用希伯來文的22個字母,亦有意識地將「神的律法」這用字安置入每句詩句中.其用意是,將每個字母,及由這些字母所造成的句子,並由這些句子所表達的思想,都被「神的律法」滲透,被「神的律法」洗禮.背誦這詩篇的人,會形成一習慣,會造成一意識,就是他用的每個字母,講的每句句子,想的每個思念,都被「神的律法」指導,被「神的律法」駕馭,被「神的律法」充滿.

「腹有詩書氣自華」,「心懷天道人自正」.若果「腹有詩書」能叫人氣度不凡的話,則「心懷天道」自能校正人要走的道路.

若果「三日不讀書,便覺語言無味,面目可憎」的話,則日日思念神的律法的人,語言應有內容,面目應有光采.

讓我們讀兩段詩篇119篇的經文,試試體會一下詩人的感受.當「神的律法」充滿詩人的思想時,他的思念,他說的話,他面對逆境時的感受,他在人生中追求的意義,都有全新的面貌.

詩篇119篇的第13段(每句經文以字母Mem開首)

詩119:97　我何等愛慕你的律法,終日不住地思想.

詩119:98　你的命令常存在我心裡,使我比仇敵有智慧.

詩119:99　我比我的師傅更通達,因我思想你的法度.

詩119:100　我比年老的更明白,因我守了你的訓詞.

詩119:101　我禁止我腳走一切的邪路,為要遵守你的話.

詩119:102　我沒有偏離你的典章,因為你教訓了我.

詩119:103　你的言語在我上膛何等甘美,在我口中比蜜更甜!

詩119:104　我藉著你的訓詞得以明白,所以我恨一切的假道.

詩篇119篇的第14段(每句經文以字母Nun開首)

詩119:105　你的話是我腳前的燈,是我路上的光.

詩119:106　你公義的典章,我曾起誓遵守,我必按誓而行.

詩119:107　我甚是受苦；耶和華啊,求你照你的話將我救活!

詩119:108　耶和華啊,求你悅納我口中的讚美為供物,又將你的典章教訓我!

詩119:109　我的性命常在危險之中,我卻不忘記你的律法.

詩119:110　惡人為我設下網羅,我卻沒有偏離你的訓詞.

詩119:111　我以你的法度為永遠的產業,因這是我心中所喜愛的.

詩119:112　我的心專向你的律例,永遠遵行,一直到底.

當「神的律法」滲透著詩人的生命時,他明白什麼才是智慧,他看到什麼才是正路,他知道憑什麼去克服人間的艱難.

詩人願意讓「神的律法」去滲透他的生命,我們又讓什麼滲透我們的生命呢?

在香港,無盡的廣告用盡各種方式滲透我們,無盡的消費品呼喚我們,不停的娛樂刺激我們.廣告向我們展示這是一個消費的世界,你的價值被你的消費能力定義,你的目標是追求更高的消費能力,你的幸福就是當下的消費滿足.手機科技製造了不間斷的娛樂的可能.每個人都埋頭苦幹地望著自己的手機.不停的看視像,不停的玩遊戲,不停的與人連線.我們不能忍受一刻沒有刺激.沒有視像或沒有遊戲刺激的一刻,那怕是一刻,都是莫大的苦悶.有了手機,我們不能等待別人慢一點的反應或回覆,那怕是慢了一刻,都是莫大的折磨.

當一切都只是消費時,便沒有價值,只有價格.當我們習慣了不停的刺激,不能靜待一刻時,我們便失去了沉澱的力量.水太急,則什麼都沖走.生命太急,也等不到有重量的東西沉澱下來.留下的,只是一沒有重量的,淺薄的生命.

詩篇119篇的詩人,讓「神的律法」滲透他的字母,他的句子,他的思想,以致他在人生種種逆境中,仍屹立不倒.我們卻容讓什麼滲透我們的生命呢?

讓我們看看今日的講道經文:詩篇119: 33-40

詩119:33　耶和華啊,求你將你的律例指教我,我必遵守到底!

詩119:34　求你賜我悟性,我便遵守你的律法,且要一心遵守.

詩119:35　求你叫我遵行你的命令,因為這是我所喜樂的.

詩119:36　求你使我的心趨向你的法度,不趨向非義之財.

詩119:37　求你叫我轉眼不看虛假,又叫我在你的道中生活.

詩119:38　你向敬畏你的人所應許的話,求你向僕人堅定!

詩119:39　求你使我所怕的羞辱遠離我,因你的典章本為美.

詩119:40　我羨慕你的訓詞；求你使我在你的公義上生活!

我們將這段經文兩節兩節的看.

詩119:33　耶和華啊,求你將你的律例指教我,我必遵守到底!

詩119:34　求你賜我悟性,我便遵守你的律法,且要一心遵守.

詩人尋求一些值得一生堅持的東西.

但要找到這些東西,需要神的恩典賜予.人有很多盲點,就算正路在面前,我們仍然可能看不見.我們需要神賜予悟性,賜予明白的心.

前幾天,我有一個學生和我分享其行山經驗.行山期間,忽遇暴風雨,看不到下山的路.後來他發現,下山的路其實就在前面,但看不見就是看不見.

神讓我們看見,堅定我們生命的「錨」在哪裡.神的律法,是我們生命的「錨」.

詩119:35　求你叫我遵行你的命令,因為這是我所喜樂的.

詩119:36　求你使我的心趨向你的法度,不趨向非義之財.

在消費的世界裡,人以能消費為所喜樂的.若以消費為樂,便很難不追求擁有更高的消費能力.這種追求,很難不貪錢,很難不向錢看,很難不為了錢而放棄某些重要的價值.在消費的世界裡,什麼都有一個價格,也只是一個價格.聽說,有些少女願意出賣自己肉體,只是為了去買一個名牌手袋.不同的人出賣不同的有價值的東西,只是為了一個有價格的東西.

詩人說,他甘願在神的引導下,聽從神的命令.他認為,這是他所喜樂的.聽從神的律法,不為什麼,只因為這本身是好的.他找到一樣真正有價值的東西,既然找到了,就不會被不義之財所迷惑.誘惑,對迷途的人,最有吸引力.對於一個生命飽足的人來說,引誘的「餌」是沒有吸引力的.

詩119:37　求你叫我轉眼不看虛假,又叫我在你的道中生活.

詩119:38　你向敬畏你的人所應許的話,求你向僕人堅定!

第38節可以簡單地譯為:「你向僕人堅定你的應許,使他敬畏你.」

在現代世界,不看「虛假」是很難的,因為我們已進入「虛擬實在」的時代.「虛假」的東西比真實的東西看來更真實.手袋的「牌子」比手袋本身更被重視.虛擬的“Bitcoin”比真實的貨幣更引人投資.一個人的實質學問並不比他擁有的學位重要.很多時,虛擬的信仰比真實的信仰更能吸引人.

在這「虛擬世界」裡,有很多不同的虛擬的潮流,有很多虛擬的價值,有很多虛擬的上帝.能否看出其虛假?遠離其虛假?

詩人說,只有神的律法,能叫人不看虛假,進入真正的「生命」裡.只有神能實現他的應許,只有神能將看來不現實的應許變成真實.虛假的東西只引誘我們進入虛假,在虛假裡,是沒有真實的.

詩119:39　求你使我所怕的羞辱遠離我,因你的典章本為美.

詩119:40　我羨慕你的訓詞；求你使我在你的公義上生活!

不按世界的標準而話,是引來羞辱的.但神的律法卻對敬畏神的人,有另一種評價.有了神的律法,我們才能從另一個角度看自己,評價自己.神的律法支撐著敬畏神的人,使敬畏神的人能知道自己走的路是正路.

詩篇119篇的經文,鼓勵我們,以「神的律法」滲透生命,才能抗衡世間的貪婪與虛假,才能找到生命.你不以「神的律法」滲透生命,其他的力量便會滲透你,消費的力量,娛樂的力量會滲透你,你的生命便容易走入迷途了.

想像一下,我們是失明的人,走在一充滿陷阱的路上.我們需要一條穩當的路軌,引導我們走到彼岸.「神的律法」就是這條路軌.



\newpage

\section{天降大任}
\label{sec:237}
\textbf{鄧瑞強博士 - 天降大任}
\newline
\newline
\hyperref[sec:236]{< < < PREV SERMON < < <}
~
\hyperref[sec:toc]{[返目錄]}
~
\hyperref[sec:238]{> > > NEXT SERMON > > >}


創世記 12:1-4
\newline
\begin{longtable}{cl}
\hline
\hline
章節 & 經文 (和合本修訂版)\\
\hline
12:1 & \begin{tabularx}{0.7\textwidth}{X} 耶和華對亞伯蘭說：「你要離開本地、本族、父家，往我所要指示你的地去。 \end{tabularx} \\ \\ \relax
12:2 & \begin{tabularx}{0.7\textwidth}{X} 我必使你成為大國，我必賜福給你，使你的名為大；你要使別人得福。 \end{tabularx} \\ \\ \relax
12:3 & \begin{tabularx}{0.7\textwidth}{X} 為你祝福的，我必賜福給他；詛咒你的，我必詛咒他。地上的萬族都必因你得福。」 \end{tabularx} \\ \\ \relax
12:4 & \begin{tabularx}{0.7\textwidth}{X} 亞伯蘭就遵照耶和華的吩咐去了；羅得也和他同去。亞伯蘭離開哈蘭的時候年七十五歲。 \end{tabularx} \\ \\
[1ex]
\hline
\hline
\end{longtable}
講題:天降大任　Called by God

經文:創世記12:1-4

講員:鄧瑞強博士

場合:香港中文大學崇基學院禮拜堂

日期:2014年3月16日

各位弟兄姊妹,早安.

或多或少,我們都曾經想做一個偉大的人,但當機會真的來到的時候,我們通常是退縮,甘於平凡.侏儒渴望做巨人,但當真的可以做巨人的時候,他卻因不習慣這巨大而害怕得要命,寧願做侏儒.我們會談論劉曉波,但一想到他在獄中,一想到他的妻子在受苦,我們便會與他保持一種心理距離,一種適合於談論而不是投入的距離.我們喜歡談論偉大的人物,因為我們或多或少渴望偉大.但我們會保持距離,因為我們抓緊的,只是一塊渺少之地.

我們是動物與天使之間的生命,縱使我們可能大部分時間活得似動物,但總會在某些時刻靈光一閃,聽到一種呼喚,呼喚我們活得比動物多一點,似天使多一點.縱使是一隻井底之蛙,仍會看到藍天.總有一天,牠會聽到藍天的呼喚,召喚牠跳出井來.我們每天都接到很多“call”,有些來自朋友,有些來自家人,更多的是「白撞」的.但我們也會接收到另一些“call”,這些“call”來自我們的良心,來自我們生命的深處,來自我們在天上的父.這些呼喚,是很微弱的,也是震耳欲聾的.這些呼喚,微弱得你幾乎聽不到,像黃葉飄落地面的聲音,你一分心,便聽不見了.這些呼喚,卻也是震耳欲聾的,像母親在街頭大聲呼喚在街尾玩耍的仔仔回家吃飯一樣,街知巷聞,你不能充耳不聞.

不知各位聽過這種呼喚沒有?

我們今日的講道內容,會講亞伯拉罕,他聽到這種呼喚.

沒有這些呼喚,人的生命就像一潭死水.呼喚,叫這潭水充滿生機.

我們可能試過去旅行時,去過一些洞穴,洞穴裡有一潭水.偶然,由高處落下一滴水珠,打在水裡,在極寧靜處,我們會聽到很明亮的「滴答」,整潭水像得到了靈魂,活了起來.來自神的呼喚,同樣會叫我們的生命活起來.

很多現代人的生命悶得發慌,很想聽到某種呼喚,令自己感到某種意義.很多書,很多講座,很多生命導師,教人尋找「生命的呼喚」.我看過一些資料,教人尋索這些「生命的呼喚」.這些資料說,首先,你要渴望做一些不平凡的事,或者渴望做一個不平凡的人.給自己一點勇氣,想像自己能改變世界.要敢膽去想像.你的想像力有多大,你的世界就有多大.最好的開始,是從你的興趣入手.若你喜歡動物的話,你有沒有想過去拯救某些瀕危物種?若你喜歡音樂的話,你有沒有想過開一場音樂會?若你喜歡吃東西的話,你有沒有想過開一間有特色的餐廳?最重要的是喚醒你的潛能,實現你的夢想.這就是你內心向你發出的呼喚.

對於一個不知自己想要什麼的人,這些講法也有點用處,起碼可以叫他去做點事情,好過坐在那裡「發霉」.但是,這些講法充其量只是鼓勵人去實現自己的夢想,未必能讓人更接近永恆.一隻井底的青蛙,縱使在井底開過音樂會,呱啦呱啦地唱,牠其實仍覺得欠了點什麼.牠頭上的藍天,仍在呼喚著牠.

一個聆聽到上天呼喚的人,他若真的聆聽,真的回應,他的生命就會經歷一種死而復生.過去的種種,像昨天的流水,一去不返.今日的生命,像是剛來到這世界一樣,是全新的.「舊事已過,都變成新的了.」

我想起陶淵明.

相傳他不肯為「五斗米」,逢迎低劣的官員.於是辭官歸故里,寫下《歸去來兮》這首詞:

「歸去來兮!田園將蕪胡不歸?既自以心為形役,奚惆悵而獨悲?悟已往之不諫,知來者之可追；實迷途其未遠,覺今是而昨非.」

「覺今是而昨非」,聆聽某種永恆呼喚的人,很多時有這種體會.Say “No” to the past, say “Yes” to a new life. 就好像一個奴隸忽然間變成一個自由人一樣,過去的一切全過去了,現在是一個新的開始.又好像一個失明的人,做過手術後,視力恢復了.黑暗過去了,現在活在光明中.

偉大的人物,多多少少,都經歷過這種「死而復生」.

今日的講道經文是:創世記12:1－4上

亞伯拉罕聆聽到神的呼喚,他將經驗「死而復生」.

創世記 12:1-4a

12:1　耶和華對亞伯蘭說:「你要離開本地,本族,父家,往我所要指示你的地去.

12:2　我必叫你成為大國.我必賜福給你,叫你的名為大；你也要叫別人得福.

12:3　為你祝福的,我必賜福與他；那咒詛你的,我必咒詛他.地上的萬族都要因你得福.」

12:4a　亞伯蘭就照著耶和華的吩咐去了.

經文一開始,是神對一個人的呼召.

這是一個偉大的時刻.一個卑微的人,面對偉大的神,聆聽神的呼喚,被委派去成就大事.

這是叫人戰慄的時刻.神的呼喚,就像基督的呼喚一樣,他叫人背起人類的十架.

這是叫人退縮的時刻,因為我們害怕十架,害怕苦難,害怕死亡.

耶和華對亞伯蘭說:「你要離開本地,本族,父家,往我所要指示你的地去.」

亞伯蘭後來改名叫亞伯拉罕,新的名字代表新的人生.

神吩咐亞伯蘭,離開本地,本族,父家.離開故土,人便會變成大地的流浪客.對流浪客而言,一切都不再固定,一切被可能隨時失去. “Fiddler on the Roof” (這片名在香港譯作《錦繡良緣》)這齣講述猶太人的電影,最能講出「離開本地」的慘況.俄國沙皇一聲令下,這群人便要收拾行裝,遠赴他方.流浪客的人生,像是片中那在屋頂上拉小提琴的琴手,他不能在屋裡拉,因他沒有屋可進,他沒有家園可言.他在屋頂這危險之地,奮力拉出生命之歌.

亞伯蘭要離開本族,離開親人.離開親人,一個人將一無所依.在荒涼的大地上,他將孤獨前行.我想起陳子昂的＜登幽州台歌＞:「前不見古人,後不見來者.念天地之悠悠,獨愴然而涕下.」亞伯蘭之所以沒有「獨愴然而涕下」,是因為他仰望天上的神.

亞伯蘭也要離開父家.父家,是他的身分所在.一個人放下過去的身分,是為了承擔新的身分.有一次,有人向耶穌說:「你的母親和你弟兄在外邊找你.」耶穌卻說:「誰是我的母親?誰是我的弟兄?凡遵行神旨意的人就是我的弟兄姊妹和母親了.」(參可3:32-35)耶穌放下了他過去的身分,因他要以新的身分,改變人的關係.人間不再強調你的家族關係,我的家族背景,人與人平等相待.

要離開本地,本族,父家,是對過去的一切道別,迎向一未知的將來.

這將要面對徹底的孤獨,面對赤裸的自己,面對歷史的十架.

聽說,在美國內戰時,有人問林肯總統:「神會在你這一邊嗎?」支持解放黑奴的林肯答道:「我關心的,不是神是否在我們這一邊,而是我們是否在神那一邊.」

「神在我們這一邊」,這種想法是叫人安心的,不再驚慌的.但「我們是否在神那一邊」,這問題叫人戰慄不已.這問題迫使林肯離開本地,本族,父家.他不能像他的父輩那樣處理黑奴的問題.午夜夢迴,林肯一定非常孤獨,他承擔著歷史的十架.

亞伯蘭要離開過去的一切,去哪裡呢?

往神所要指示他去的地方去.

這是什麼地方?

不知道.

這不是什麼生涯規劃,這是生命的歷險.這不是先有一個人生目標,然後計算一下自己要付什麼代價,而是我們放下自己的想法,也不知要付的代價有多高,只是單單順從神的呼喚.這不是想要完成某種夢想,然後求神祝福,而是將生命交給神,讓神帶領我們走到我們意想不到的地方.

南非的杜圖大主教(Archbishop Desmond Tutu)寫了一本書:《上帝有一個夢》(God has a Dream).他講到,很多人聽從了神的呼喚,離開了本地,本族,父家,承擔了歷史的十架,然後,緩慢地,一個強行種族隔離政策的南非改變了.書裡講到第一次民主選舉時的情況.杜圖這樣寫:

「人們依序走進投票亭,當他們從另一端出來之後,就完全變成新人了.黑人帶著過去被踐踏的尊嚴,被非人方式對待的苦楚走進這裡,隨後投下了票,然後說:『嘿,我自由了──我已重拾尊嚴,也被當成人來看待,我自由了!』…白人走進投票亭時,因著過去享受的種種不義特權而背負了沉重的罪咎,但從投票亭出來後,他們也成為新人:『嘿,我自由了,肩上的重擔已經卸下,我自由了!』」(杜圖,亞伯藍斯合著:《上帝有一個夢》(雅歌),頁34-35.)

這不是聽從神的呼喚的人夢想過的,而是他們從來沒有夢想過的.他們從來沒有夢想過坐了27年監的曼德拉,轉眼便成為了民選總統.他們從來沒有夢想過充滿種族仇恨的地方,能成就復和.

神發出了呼喚,要求人放下過去的一切,走向一未知的將來.有人回應了,有人行動了,然後,新的歷史便開始了.

12:2　「我必叫你成為大國.我必賜福給你,叫你的名為大；你也要叫別人得福.

12:3　為你祝福的,我必賜福與他；那咒詛你的,我必咒詛他.地上的萬族都要因你得福.」

12:4a　亞伯蘭就照著耶和華的吩咐去了.

神呼喚亞伯蘭,不是要成就他的個人夢想,而是要他成為祝福萬民的僕人.亞伯蘭的個人生命消失了,他成為神祝福萬民的一個具體標記.世上有某些人,是某種東西的具體標記.有些人是民主的標記,有些人是愛的標記,有些人是謊言的標記.擁抱這些人,就是擁抱他們代表的事物.亞伯拉罕是聆聽神的呼喚的人的一個模範.擁抱他生命的人,就是願意聆聽神的呼喚的人,他們也會成為別人的祝福.反對他的人,就是拒絕聆聽神的呼喚的人,他們成為人間的禍患.

靈性上的偉大人物,他們都是屬於同一類別的人,他們站在亞伯拉罕身邊,都很熟悉神的呼喚,都很明白其中的孤獨,都很明白當中的十架.

德蘭修女有一天,聽到神的呼喚,她這樣寫:「我要離開〔現在〕的修院,在窮人中間工作,在他們中間生活.這是一道命令.我知道我屬於那裡,但我不知道如何去到那裡.」(參\url{http://www.peacefulearth.com/magazine/motherteresa.html})

是否有點熟悉?那位呼召亞伯拉罕的神,也呼召她.神呼喚德蘭修女,離開本地,本族,父家,往神所要指示她的地去.她不知應如何做,但她知道這是不可抗拒的命令.她要成為別人的祝福,同時,她要背起人類的十架.她成了回應神的呼喚的人的又一個具體模範.

今日,神又再向我們發出他的呼喚.

我們會如何回應?

若果神用電話call我們,他大致會聽到這樣的回覆:「我現在很忙,請遲一點再call過.」

神一定會再call我們的.



\newpage

\section{容易迷失的歸途}
\label{sec:238}
\textbf{伍渭文牧師 - 容易迷失的歸途}
\newline
\newline
\hyperref[sec:237]{< < < PREV SERMON < < <}
~
\hyperref[sec:toc]{[返目錄]}
~
\hyperref[sec:239]{> > > NEXT SERMON > > >}


講題:容易迷失的歸途　Going Astray on Home Coming Journey

經文:詩篇95

講員:伍渭文牧師

場合:香港中文大學崇基學院禮拜堂

日期:2014年3月23日

我今天的題目改了,但經文仍是原來詩篇九十五篇.我改動原來的題目因為最近馬來西亞航機HM370的迷航,對一百五十一位中國乘客來說,他們在是歸家途中突然離奇迷失.事情發展到現在,我們知道航機極可能是被騎劫,偏離歸家航道.因為有人心裡迷糊,為了某種原因,以數百人的生命作籌碼,偏離歸途. 今天舊約經課出埃及記17:1-7,也是詩95引用的歷史殷鑑,也提到迷失的歸途:以色列人因為不信,心裡剛硬,在曠野迷失飄流,延宕日子,消耗生命,失去四十年的光陰,最後一個也不能進入迦南.詩95: 9-11:「那時,你們的祖宗試探我,並且觀看我的作為.四十年之久,我厭煩那世代,說:這是心裡迷糊的百姓,竟不曉得我的作為!我在怒中起誓說:他們斷不可進入我的安息」

這星期,有關馬航的消息鋪天蓋地,我們和乘客的家屬都一樣心裡焦灼,要是客機真的被騎劫,不是墮機,我們還懷有一線希望.在眾多新聞中,美國CNN引用一位航空專家的一句話,很有啟發性:「所有的飛機駕駛員的職責是啟航,航行和通訊,按這樣的次序」(The duty of all pilots is to aviate, navigate and communicate, in that order) . 舊約的以色列人和今天的信徒也是在歸家途中的天路客,曠野就是寄居的世上,曠野飄流,就是天路客的旅途寫實.所有天路客也有啟航,航行和通訊.過紅海是起航,洗禮就是過紅海:「弟兄們,我們不願意你們不曉得,我們的祖宗以前在雲下,都從海中經過,都在雲裡,海裡受洗歸了摩西,並且都吃了一樣的靈食,也喝了一樣的靈水,所喝的是出於他們的靈磐石,那磐石就是基督.但他們中間多半是神不喜歡的人,所以在曠野倒斃.」(林前10: 1-5)我們領受聖洗禮的一刻,就開始了天路客的旅程(Aviating) ；我們在地上寄居的日子是航行(Navigating) .

在天上航行,要與地上的指揮中心保持聯系(Communicating),聽從他們的指引.一架飛機起航時要聽從機場指揮塔,祇有他們才知道何時可進入跑道,下降也如是.航行過程中若經過某國家航空識別區,要聽從當地指揮中心,上報有關航機資料,否則被視為敵對飛機,再越過所屬領空,便有可能被炮彈擊落.1983年一架南韓007號客機,就因為誤闖前蘇聯領空,沒有和地面保持通訊,被導彈擊落,機上269人全部罹難. 保持和地面通訊是空中航行不可或缺的,因為祇有地上的雷達才看得透天上的環境.信徒在地上行經時間,同樣需要天上的智慧.在曠野旅行的客旅都知道,因為沒有明顯的地標,容易迷失,要仰視天上的北斗來辨別方向.我們每主日參加崇拜,就和上主溝通,跟上主說話,也聆聽上主,聽他說話；讓生命從新聚焦,調較方向,免得心裡迷糊.詩篇九十五篇的結構就有這雙向的溝通(Communication).

結構: 第一至七節第一段,是敬拜者以啟應對答方式頌揚上主；第二段八至十一節,是上主藉祭司以宣講方式訓示會眾；中間的連結是「惟願你們今天聽他的話」,這連結也是詩篇的中心概念.我們為何要聽上主的話呢?詩篇提出兩個正面和反面的原因:1. 我們屬於他,所以要聽他的話；2.  不聽他話的人,不能進入應許之地. 第一段-人和上主說話-上主是我們的創造主和救贖主,所以要聽他的話: 這一段有三個邀請.第一,第二邀請是平行體,可由主禮和會眾應答讀出,表示來聖殿敬拜的歡騰,是外顯的,禁捺不住的興奮.「來啊,我我們要向耶和華歌唱；向拯救我們的磐石歡呼!我們要來感謝他,用詩歌向他歡呼.」(vv. 1, 2)歡呼是激越澎湃的情緒,Eugene Peterson 英譯本 The Massage頗能譯出經文的神韻:來啊,讓我們喊出讚美,把屋頂升起；讓我們操進他跟前,以讚美托高橫樑(Come, let’s shout praises to God, raise up the roof….Let’s march into his presence singing praises lifting the rafters with our hymns.) .

跟著指出歡呼的原因:「耶和華為大神,為大王,超乎萬神之上.地的深處在他手中,山的高峰也屬他.海洋屬他,是他造的,旱地也是他手造成的.」(vv.3-5) 他是創造主,賜我們生命,他極偉大,超乎一切,可敬可畏. 第三個邀請「來啊,我們要屈身敬拜,在我們的耶和華面前跪下.因為他是我們的神,我們是他草場的羊,是他手下的民.」(6, 7) 是內歛靜態,莊嚴肅穆,以肢體的俯伏和下跪,表達對上主的敬畏.「耶和華的面前」指進到至聖所的前面,面對何等的神聖空間,我們要俯伏,我們要下跪.崇拜從開始的歡騰到後來的靜默肅穆,有點像我們宣召後第一首詩是拼發歡欣的頌揚,到聽完道的回應詩歌較緩慢和安靜,讓我們反思.

崇拜方式多元豐富,歡騰莊嚴並舉,應答直述交替；但激越中不失莊重,興奮中不失謙恭,熱切中不失分寸,以聲音和肢體語言來表達各種情緒. 第三個邀請強調上主是立約的救贖主:「我們是他草場的羊,是他手下的民,惟願你們今天聽他的話.」這話是在西乃山上主和以色列人立約時說的.他們離開埃及,目的是來到西乃山與上主立約,接受上主藉摩西頒下的十誡.出埃及記19: 4-6:「我向埃及人所行的事,你們都看見了,且看見我如鷹將你們背在翅膀上,帶來歸我.如今你們若實在聽從我的話,遵守我的約,就要在萬民中作屬我的子民,因為全地都是我的.你們要歸我作祭司的國度,為聖潔的國民.」

第二段vv.8-11上主和人說話,以歷史作為殷鑑來說明為何要聽他的話: 第一段是敬拜者以第一身說話:「來啊,讓我們」頌讚感恩上主；但vv. 8-11的第二段像祭司的講道,會眾以第二身聆聽: 你們不可硬著心,像當日在米利巴,就是在曠野的瑪撒.那時,你們的祖宗試探我,並且觀看我的作為.四十年之久,我厭煩那世代,說:「這是心裡迷糊的百姓,竟不曉得我的作為!」(vv. 8-10) 詩篇九十五篇前半段是敬拜者對上主說話,歌唱感謝,以聲音和動作,用應答和直述；後半段是聆聽教誨,是上主對人說話.其實,在崇拜中,上主也透過我們和他說話提醒我們,因為詩歌,禱文都蘊含聖經的真理,激勵我們.我們給上主的,上主加倍歸還給我們,這是崇拜的誘發作用.在崇拜中上主除了直接透過講員的宣道與我們說話,也利用我們對他所呈獻的詩歌和禱告對我們說話,誘發我們應有的態度和信心.我們現在看看這詩篇九十五篇的信息.

1. (不能逆轉的啟航) 沒有返回點的啟航 (Aviation Without A Point of No Return) vv. 1-7

詩篇用創造和救贖一起說明我們和上帝的關係,兩者有分別但不能分離,因為救贖是在創造的場景進行.奴役以色列人的埃及政權是現實的,我們可以在現今不少專制政權中,找到他們的基因.要越過的紅海現在還在地球上；被救贖的選民要返回真實的世界,在萬民中作屬上主的子民. 上主是創造主又是救贖主,說明我們沒有辦法逃避上主,沒有信仰的人也要面對生命的主.上主是創造的主,我們同屬被造之物,同屬生物界.生物有開始,也有結束；生物有出生,也有死亡.在哲學的世界,可以談論不朽,無限,完美.但在生物世界,需要生生不息,新陳代謝；人有一生,必有一死.杜甫《登高》詩:「無邊落木蕭蕭下,不盡長江滾滾來」；落葉飄零,無邊無際,蕭蕭而下；不盡長江,洶涌澎湃,奔騰而來.長江後浪推前浪,花果轉眼成污泥.

所有飛機必須有燃料才能起飛,在燃料的標板,有一臨界點-「不能返回點」,PNR (Point of No Return),提醒機師燃料剛剩下一半,此刻尚有可能返回起飛的機場,過後就沒有機會了.但生命的航程沒有「不能返回點」,人可以重返出生的地點,但不能返回過去的歲月,因為作為生物,一開始就不能返回起始點了.

2. 容易迷失的航程(Easily Astray Navigation) vv.8-10

你們不可硬著心,像當日在米利巴,就是在曠野的瑪撒.那時,你們的祖宗試探我,並且觀看我的作為.四十年之久,我厭煩那世代,說:「這是心裡迷糊的百姓.竟不曉得我的作為!」(vv. 8-10) 所有機師都知道,他們才不怕跨越「不能返回點」,因為已經預備了一個機場,可以安全降落,而且有人在那裡等候歡迎.那安全降落的機場就是我們的目標,也是我們的歸屬,我們的家,我們一生都問:我們情歸何處呢?甚麼是值得我們付托終身的呢? 詩篇第九節提到百姓心裡迷糊,一遇到困難就埋怨為甚離開埃及,來到曠野荒蕪之地,他們對情歸何處感到迷惘.

這段歷史就是今天的舊約經課出17: 1-7,地點就在利非訂(Rephidim)-紅海到西乃山最後的駐營地；拔營出發,下一站就是西乃山了.西乃在望了,他們等候很久的聖山就快到了,上主呼召摩西把他們領出來,就是要到這山和上主立約. 但黎明前的一刻總是最黑暗的,他們在最後的一站的利非訂發現沒有水源,怎麼辦呢!我們會因缺水而渴死的.「百姓在那裡甚渴,要喝水,就向摩西發怨言,說:你們為甚麼將我們從埃及領出來,使我們和我們的兒女並牲畜都渴死呢?」(出17: 3)  他們的要求合理嗎?這是生死攸關的缺水,沒有水怎能到西乃山敬拜呢?我們命都不保了!要求是對的,但態度不對.他們懷疑上主不理他們,放棄他們,並懷疑一路持守的信仰；過紅海是極大的錯誤,留在埃及多麼安全. 其實不是第一次發怨言了,出16章他們埋怨沒有食物,「巴不得我們早死在埃及地耶和華的手下!那時我們坐在肉鍋旁邊,吃得飽足,到這曠野,是要叫我們這全會眾都餓死啊!」(出16: 3) 民數記14: 20說百姓有十次之多試探上主.

但恩慈的上主知道他們的需要,天降嗎哪,驅使鵪鶉(quail) 作為他們的食物. 他們也不可能不知道上主會及時給他們水喝的.因為在不久之前,上主將不能喝的瑪拉苦水變甜.「百姓就向摩西發怨言,說:我們喝甚麼呢?耶和華指示他一棵樹,他把樹掉在水裡,水就變甜了.」(出15: 24, 25)

甚麼是試探上主呢?試探上主是不相信上主,希望上主按我們的意思做事.出17: 7「因以色列人爭鬧,又因他們試探耶和華說:耶和華是在我們中間不是?」,詩95: 8「你們不可硬著心,像當日在米利巴(爭鬧),就是在曠野的瑪撒(試探).」他們不相信上主,除非要看到證據.耶穌在曠野受魔鬼試探:「你若是上帝的兒子,可以跳下去,因為經上記著說:主要吩咐他的使者,用手托著你,免得你的腳碰在石頭上.」(太4: 6) .不錯,這是上帝的應許,上帝會保護我們,但不是按我們的時間,我們的方法.這是把上主「雙規」:在規定的時間,規定的方法,證明上主是在我們中間.我們不用買醫療保險,因上主會保護我們不生病；病了不用看醫生吃藥,上主會醫治我.當病情愈來愈重時,又會說是我們不夠依靠上主,沒有信心祈禱.試探上主核心問題是對他不夠信心,要「雙規」他,要看看「耶和華在我們中間不是.」

試探上主是控制上主,按我們的時間,依從我們的方法,為我們辦事.詩篇說:他們「並且觀看我的作為」,但經結果是「竟不曉得我的作為」.試探上主是靠眼見,不憑信心；但沒有信心,就不能看見上主的作為. 我們觀察到,每一次百姓遇到困難,就想回到過去埃及的舒適安全環境,埋怨摩西帶他們來到曠野.他們最大的試探是覺得還有「不能返回點」,還來得及返回埃及.但航機已啟航,飛機已越過「不能回頭點」,我們的飛機正飛向新機場降落.

針對這身份危機,詩篇強調以色列人的歸屬:「他是我們的神,我們是他草場的羊,是他手下的民.」(詩95: 7) 荷馬史詩奧德賽(Odyssey) 記述了主角俄底修斯(Odysseus, Uylsses) 情歸何處的旅途.希臘伊薩卡王(Ithaca) 俄底修斯運用計謀,木馬屠城攻陷特洛伊(Troy) .這位英雄十年離家後思故鄉,因為惹怒海神,歸途風急浪高,漫長迂迴,險阻層層,試探重重. 因為人的愚昧,到達了海港但不能上岸.在歸途中,俄底修斯得風之主埃俄洛斯(Aeolus) 的款待,他把所有的風收在袋裡,並發動西風送行,因而俄底修斯很快就到了家鄉的海港,但同行的人以為袋中收藏了禮物,俄底修斯獨占不給他們,就偷偷打開看看.袋口一開,風就出來把他們吹回原處.他們再三求埃俄洛斯再送西風時被拒,這是愚昧造成的迷途. 俄底修斯也試過一晌貪歡,在溫柔鄉中忘記了家中等候的愛妻.但俄底修斯始終沒有忘記自己的身份,他要歸家.延宕十年,最後回到妻兒身邊.奧德賽是在字典中已成為情歸何處的尋覓,每個人都尋找那安全降落的機場,因為我們一出生就越過了「不能返回點」.

3 . 命運攸關的訊息 (The Communication of Destiny) v. 11

所以我在怒中起誓說:「他們斷不可進入我的安息」(v.11) 這是極嚴厲的警告,我們不能對上主的忍耐視若無睹,硬著心腸拒絕聖靈的感動,否則就失去屬靈的敏銳,愈來會麻木,最後失去悔改的機會.我很喜歡富蘭克林(Benjamin Franklin) 的名言:To err is human, to repent divine, to persist devilish, 犯錯皆凡人,悔改靠神恩,沉溺是魔障.俄底修斯曾經迷途,但沒有忘記他的家；馬可跟保羅到歐洲傳道,因為種種原因中途放棄,但在巴拿巴鼓勵培育下,最後成為保羅激賞的同工. 這些百姓最後一個也不能進入迦南,死在曠野；更可惜是,在一生的生命中,看不到上主的作為:「這是心裡迷糊的百姓,竟不曉得我的作為!」這些死在曠野的百姓,看不到後來約書亞如何把耶利歌城攻陷.他們不明白,人看來苦澀的生命,可以轉化成祝福他人的生命,像瑪拉的苦水變成甜美,供人飲用.他們也不會明白,上主會為跟隨他的人,天降嗎哪,總有及時出人意表的供應.

你知道嗎?你的生命已啟航了,我們生命的起點就是不能回頭點.在地上航行,我們需要聆聽天上的智慧,才不至迷失.今天你情歸何處呢?你要降落的機場在那裡呢?你預備好降落的機場嗎?你有機場給你降落嗎?

今天的詩篇九十五篇告訴我們:上主說:聆聽我話的,必曉得我的作為；聆聽我話的,必得著所賜的安息.

聆聽我話的,必曉得我的作為；聆聽我話的,必得著所賜的安息.



\newpage

\section{信仰讓生命不再一樣}
\label{sec:239}
\textbf{鄭漢文博士 - 信仰讓生命不再一樣}
\newline
\newline
\hyperref[sec:238]{< < < PREV SERMON < < <}
~
\hyperref[sec:toc]{[返目錄]}
~
\hyperref[sec:240]{> > > NEXT SERMON > > >}


約翰福音 9:1-41
\newline
\begin{longtable}{cl}
\hline
\hline
章節 & 經文 (和合本修訂版)\\
\hline
9:1 & \begin{tabularx}{0.7\textwidth}{X} 耶穌往前走的時候，看見一個生來就失明的人。 \end{tabularx} \\ \\ \relax
9:2 & \begin{tabularx}{0.7\textwidth}{X} 門徒問耶穌：「拉比，這人生來失明，是誰犯了罪？是這人還是他的父母呢？」 \end{tabularx} \\ \\ \relax
9:3 & \begin{tabularx}{0.7\textwidth}{X} 耶穌回答：「既不是這人犯了罪，也不是他的父母，而是要在他身上顯出神的作為來。 \end{tabularx} \\ \\ \relax
9:4 & \begin{tabularx}{0.7\textwidth}{X} 趁著白日，我們必須做差我來的那位的工；黑夜來到，就沒有人能做工了。 \end{tabularx} \\ \\ \relax
9:5 & \begin{tabularx}{0.7\textwidth}{X} 我在世上的時候，是世上的光。」 \end{tabularx} \\ \\ \relax
9:6 & \begin{tabularx}{0.7\textwidth}{X} 耶穌說了這些話，就吐唾沫在地上，用唾沫和了泥抹在盲人的眼睛上， \end{tabularx} \\ \\ \relax
9:7 & \begin{tabularx}{0.7\textwidth}{X} 對他說：「你到西羅亞池子裡去洗。」（西羅亞翻出來就是「奉差遣」。）於是他去，洗了，回來就看見了。 \end{tabularx} \\ \\ \relax
9:8 & \begin{tabularx}{0.7\textwidth}{X} 他的鄰舍和素常見他討飯的人，就說：「這不是那從前坐著討飯的人嗎？」 \end{tabularx} \\ \\ \relax
9:9 & \begin{tabularx}{0.7\textwidth}{X} 有的說：「是他」；又有的說：「不是，卻是像他。」他自己說：「是我。」 \end{tabularx} \\ \\ \relax
9:10 & \begin{tabularx}{0.7\textwidth}{X} 於是他們對他說：「你的眼睛是怎麼開的呢？」 \end{tabularx} \\ \\ \relax
9:11 & \begin{tabularx}{0.7\textwidth}{X} 那人回答：「有一個名叫耶穌的，他和了泥抹我的眼睛，對我說：『你到西羅亞池子去洗。』我去一洗，就看見了。」 \end{tabularx} \\ \\ \relax
9:12 & \begin{tabularx}{0.7\textwidth}{X} 他們對他說：「那個人在哪裡？」他說：「我不知道。」 \end{tabularx} \\ \\ \relax
9:13 & \begin{tabularx}{0.7\textwidth}{X} 他們把以前失明的那個人帶到法利賽人那裡。 \end{tabularx} \\ \\ \relax
9:14 & \begin{tabularx}{0.7\textwidth}{X} 耶穌和泥開他眼睛的那一天是安息日。 \end{tabularx} \\ \\ \relax
9:15 & \begin{tabularx}{0.7\textwidth}{X} 法利賽人又問他是怎麼得看見的。他對他們說：「他把泥抹在我的眼睛上，我一洗，就看見了。」 \end{tabularx} \\ \\ \relax
9:16 & \begin{tabularx}{0.7\textwidth}{X} 於是法利賽人中有的說：「這個人不是從神來的，因為他不守安息日。」另有的說：「一個罪人怎能行這樣的神蹟呢？」他們之間就產生了分裂。 \end{tabularx} \\ \\ \relax
9:17 & \begin{tabularx}{0.7\textwidth}{X} 於是他們又對那盲人說：「他開了你的眼睛，你說他是怎樣的人呢？」他說：「他是個先知。」 \end{tabularx} \\ \\ \relax
9:18 & \begin{tabularx}{0.7\textwidth}{X} 猶太人不信他以前是失明，後來能看見的，等到叫了他的父母來， \end{tabularx} \\ \\ \relax
9:19 & \begin{tabularx}{0.7\textwidth}{X} 問他們說：「這是你們的兒子嗎？你們說他生來是失明的，現在怎麼看見了呢？」 \end{tabularx} \\ \\ \relax
9:20 & \begin{tabularx}{0.7\textwidth}{X} 他的父母就回答說：「他是我們的兒子，生來就失明，這是我們知道的。 \end{tabularx} \\ \\ \relax
9:21 & \begin{tabularx}{0.7\textwidth}{X} 至於他現在怎麼能看見，我們卻不知道；是誰開了他的眼睛，我們也不知道。他已經是成人，你們問他吧，他自己會說。」 \end{tabularx} \\ \\ \relax
9:22 & \begin{tabularx}{0.7\textwidth}{X} 他父母說這些話，是怕猶太人，因為猶太人已經商定，若有宣認耶穌是基督的，要把他趕出會堂。 \end{tabularx} \\ \\ \relax
9:23 & \begin{tabularx}{0.7\textwidth}{X} 因此他父母說「他已經是成人，你們問他吧」。 \end{tabularx} \\ \\ \relax
9:24 & \begin{tabularx}{0.7\textwidth}{X} 於是法利賽人第二次叫了那以前失明的人來，對他說：「你要將榮耀歸給神，我們知道這人是個罪人。」 \end{tabularx} \\ \\ \relax
9:25 & \begin{tabularx}{0.7\textwidth}{X} 那人就回答：「他是不是個罪人，我不知道；有一件事我知道，我本來是失明的，現在我看見了。」 \end{tabularx} \\ \\ \relax
9:26 & \begin{tabularx}{0.7\textwidth}{X} 他們就問他：「他給你做了甚麼？是怎麼開了你的眼睛？」 \end{tabularx} \\ \\ \relax
9:27 & \begin{tabularx}{0.7\textwidth}{X} 他回答他們：「我已經告訴你們了，你們不聽，為甚麼又要聽呢？難道你們也要作他的門徒嗎？」 \end{tabularx} \\ \\ \relax
9:28 & \begin{tabularx}{0.7\textwidth}{X} 他們就罵他：「你是他的門徒，而我們是摩西的門徒。 \end{tabularx} \\ \\ \relax
9:29 & \begin{tabularx}{0.7\textwidth}{X} 神對摩西說話是我們知道的，可是這個人，我們不知道他從哪裡來。」 \end{tabularx} \\ \\ \relax
9:30 & \begin{tabularx}{0.7\textwidth}{X} 那人回答，對他們說：「他開了我的眼睛，你們竟不知道他從哪裡來，這真是奇怪！ \end{tabularx} \\ \\ \relax
9:31 & \begin{tabularx}{0.7\textwidth}{X} 我們知道神不聽罪人，惟有敬奉神、遵行他旨意的，神才聽他。 \end{tabularx} \\ \\ \relax
9:32 & \begin{tabularx}{0.7\textwidth}{X} 從創世以來，未曾聽見有人開了生來就失明的人的眼睛。 \end{tabularx} \\ \\ \relax
9:33 & \begin{tabularx}{0.7\textwidth}{X} 這人若不是從神來的，甚麼也不能做。」 \end{tabularx} \\ \\ \relax
9:34 & \begin{tabularx}{0.7\textwidth}{X} 他們回答他說：「你完全是生在罪中的，還要來教訓我們嗎？」於是他們把他趕出去了。 \end{tabularx} \\ \\ \relax
9:35 & \begin{tabularx}{0.7\textwidth}{X} 耶穌聽說他們把他趕出去，就找到他，說：「你信人子嗎？」 \end{tabularx} \\ \\ \relax
9:36 & \begin{tabularx}{0.7\textwidth}{X} 那人回答說：「主啊，人子是誰？告訴我，好讓我信他。」 \end{tabularx} \\ \\ \relax
9:37 & \begin{tabularx}{0.7\textwidth}{X} 耶穌對他說：「你已經看見他，現在和你說話的就是他。」 \end{tabularx} \\ \\ \relax
9:38 & \begin{tabularx}{0.7\textwidth}{X} 他說：「主啊，我信！」他就拜耶穌。 \end{tabularx} \\ \\ \relax
9:39 & \begin{tabularx}{0.7\textwidth}{X} 耶穌說：「我為審判到這世上來，使不能看見的看見，能看見的反而失明。」 \end{tabularx} \\ \\ \relax
9:40 & \begin{tabularx}{0.7\textwidth}{X} 同他在那裡的法利賽人聽見這些話，就對他說：「難道我們也失明了嗎？」 \end{tabularx} \\ \\ \relax
9:41 & \begin{tabularx}{0.7\textwidth}{X} 耶穌對他們說：「你們若是失明的，就沒有罪了；但現在你們說『我們能看見』，你們的罪還在。」 \end{tabularx} \\ \\
[1ex]
\hline
\hline
\end{longtable}
詩篇 23:1-6
\newline
\begin{longtable}{cl}
\hline
\hline
章節 & 經文 (和合本修訂版)\\
\hline
23:1 & \begin{tabularx}{0.7\textwidth}{X} 耶和華是我的牧者，我必不致缺乏。 \end{tabularx} \\ \\ \relax
23:2 & \begin{tabularx}{0.7\textwidth}{X} 他使我躺臥在青草地上，領我在可安歇的水邊。 \end{tabularx} \\ \\ \relax
23:3 & \begin{tabularx}{0.7\textwidth}{X} 他使我的靈魂甦醒，為自己的名引導我走義路。 \end{tabularx} \\ \\ \relax
23:4 & \begin{tabularx}{0.7\textwidth}{X} 我雖然行過死蔭的幽谷，也不怕遭害，因為你與我同在；你的杖、你的竿，都安慰我。 \end{tabularx} \\ \\ \relax
23:5 & \begin{tabularx}{0.7\textwidth}{X} 在我敵人面前，你為我擺設筵席；你用油膏了我的頭，使我的福杯滿溢。 \end{tabularx} \\ \\ \relax
23:6 & \begin{tabularx}{0.7\textwidth}{X} 我一生一世必有恩惠慈愛隨著我；我且要住在耶和華的殿中，直到永遠。 \end{tabularx} \\ \\
[1ex]
\hline
\hline
\end{longtable}
講題:信仰讓生命不再一樣　Faith Makes Life Different

經文:約翰福音9:1-41;詩篇23:1-6

講員:鄭漢文博士(rogercheng@cuhk.edu.hk)

場合:香港中文大學崇基學院禮拜堂

日期:2014年3月30日

1. 四個失明人士的苦樂故事

弟兄姊妹,平安.

今天,我好像回鄉探親一樣,返回崇基禮拜堂這個信仰大家庭.

我沒有甚麼手信,只有一組信息,一首詩歌,和一段哀愁.

首先讓我分享一個故事.

從前有四個失明的人士,一同討論:誰最苦?誰最幸福?

第一個是一生都有視力的人,覺得自己最幸福,因為一生可以看見大自然的美麗,人間的燦爛,而且許多世間的歡愉都跟視覺有關的,例如看書,看電影,看美人.但是,這一種人好容易患上一種叫「屬靈視障」(spiritual visual-handicap)的苦困而不知,失去的是一點靈明.

第二個是到晚年才失明的人,覺得跟第一個人的幸福相若,只欠到晚年有點遺憾,無法在離世前把所愛的人好好地看一眼.在失去視力的日子,無法享受跟視覺有關的世間歡愉.但是瞎眼之時,心靈反得看見一些之前看不到的東西.會是甚麼東西呢?

第三個是生來瞎眼到死依然失明的人,一生無法看見.他的世界,失去色彩.生活不便,間受歧視.然而,他的時間不會受視覺有關的世間歡愉所困,他的聽覺更靈敏,他的心靈比前兩個人更寧靜.

第四個是生來瞎眼到後來得醫治而開眼看見的人,他經歷前三個人的生命感,而且他看見施恩的上帝.

若然上帝問你,以上四種人生,你會選擇哪一種?

我們都渴望從未失明:因為相信一個終身開眼(能看見)的人生比較幸福.

我們總不想永久失明:因為相信一個終身瞎眼(不能看見)的人生比較痛苦.

今天的信息,是環繞這四種失明人士.

2. 親眼看見的屬靈視野就是證明:約翰福音第九章經文敘事分析

今日經課約翰福音第九章第1至41節這段經文可以分為四部份,就是分為起承轉合.

2.1  起:治好生來瞎眼的

(請打開聖經)和合本聖經(約翰福音第九章第1至7節)是這樣記載這個故事的起頭:

1耶穌過去的時候,看見一個人生來是瞎眼的.

2門徒問耶穌說:「拉比,這人生來是瞎眼的,是誰犯了罪?是這人呢?是他父母呢?」

3耶穌回答說:「也不是這人犯了罪,也不是他父母犯了罪,是要在他身上顯出神的作為來.

4趁著白日,我們必須做那差我來者的工；黑夜將到,就沒有人能做工了.

5我在世上的時候,是世上的光.」

6耶穌說了這話,就吐唾沫在地上,用唾沫和泥抹在瞎子的眼睛上,

7對他說:「你往西羅亞池子裡去洗.」(西羅亞翻出來就是「奉差遣」.) 他去一洗,回頭就看見了.

約翰敘述這個故事,先從耶穌與門徒的共同經歷入手,門徒跟隨這個師傅,路上一邊走動一邊學習,耶穌隨走隨傳.我將耶穌這種使徒教學稱為「邊走邊談 Walk and Talk教學法」(英文可以稱之為Walkie-Talkie Pedagogy).

所以,門徒在路上見到甚麼人生現象,就引發相關的信仰反思.

這回所見,是城中皆知的這個人是生來是瞎眼的.

門徒趁耶穌注視這瞎子的機緣,提問了一個自古以來存在的「苦罪懸謎」:「人之受苦是否因為人的罪所產生?上帝為甚麼容許苦難存在?」

舊約最出名的案例,大概可算是《約伯記》.

回到約翰描繪的當時,就大家所知,眼前這個人不單止是瞎眼的,而且生來是瞎眼的.為甚麼是生來就瞎眼?

第一種解釋是:是這個人犯了罪?是疚由自取?但是,他豈不是在母腹中就犯了罪?

第二種解釋是:是這個人的父母犯了罪?但是,這個人豈不是在無辜受害?

門徒以為只得這兩種解釋,這種兩難式的思維,其實有不少假設.

首先假設:生來是瞎眼是一種苦痛,不是好事.

其次假設:生來是瞎眼這類苦事,是有得解釋其原因的,即必然有因才有此果.

第三假設:生來是瞎眼這類苦事,是某人犯了罪,起了這種起因,才有此苦果.

第四假設……算了算了,大家知我以前是讀哲學的,從哲學分析下去還可以再細分許多假設.

耶穌不是門徒的哲學tutor,而至少是信仰上的屬靈導師(spiritual mentor),他直接否定門徒數不清的假設,且直接渡引他們得見屬靈的視野(spiritual vision),就是通過醫治「一個人人以為是天生終身肉眼失明者」而同時醫治門徒的「屬靈視障」(spiritual visual-handicap).

耶穌借助醫治這個生來瞎子(born blind),同時打開門徒的心眼,讓他們看見眼前的師傅其實就是世上的光,眼前所作的同時是上帝的作為,就是在趁白日做著的是「那差我來者的工」.

大家有沒有留意,故事起首這一段的敘述下,這個生來失明的瞎子完全是被動的.

一,他未知道眼前看不見的耶穌是那一位.(廣東話:你邊位?)

二,他未必明白這位耶穌在跟一班同行的過路人在講甚麼.(你地講咩?)

三,他沒有邀請這位耶穌去醫治自己生來的目盲.(可能從來無信仰過有此被醫好的可能性.)

四,他沒有預期,耶穌會吐唾沫在地上,用唾沫和泥抹在瞎子的眼睛上.(你有無攪錯?咁污穢!)

五,但是,他按耶穌吩咐前往那個是「奉差遣」意思的西羅亞池子裡去洗.(人叫佢又去做?)

六,最後的事實,他回頭就看見了.(瞎眼忽然看見時,耶穌反而不見了.)

生命不可能出現的好轉,好的改變,本來最值得慶祝的.

一個人人以為生來失明的瞎子,竟然可以重見光明.豈非令人振奮?

非也,非也.重見光明的人立即面對的是黑暗,而不是身邊人的祝賀.

首先面對的是鄰舍的質疑,繼而是法利賽人的質詢,還有親人的置身度外.

從前風聞人間有惡,現在親眼看見.

這竟然是這個肉身視障的人的第一層屬靈經歷.故事發展下去,轉入他的第一身,同時呈現他的屬靈經歷.

2.2  承:生來瞎眼被治好後首先面對的鄰舍

故事第二段是這樣的:生來瞎眼被治好後,首先要面對的是鄰舍這班群眾.

第8至12節是這樣記載的:

8他的鄰舍和那素常見他是討飯的,就說:「這不是那從前坐著討飯的人嗎?」

9有人說:「是他」；又有人說:「不是,卻是像他.」他自己說:「是我.」

10他們對他說:「你的眼睛是怎麼開的呢?」

11他回答說:「有一個人,名叫耶穌,他和泥抹我的眼睛,對我說:『你往西羅亞池子去洗.』我去一洗,就看見了.」

12他們說:「那個人在哪裡?」他說:「我不知道.」

每個人都渴望有關心自己的鄰舍,可以為自己的苦事而分憂,為自己的樂事而開心.

生來失明的瞎子也一樣需要這樣的鄰舍.

可惜,他所面對的鄰舍,一早把他定型:那素常見他是討飯的.即係:乞丐.

長期被藐視,很難重新獲得尊重.

鄰舍都無預期到這個乞丐會重見光明的,所以不知如何處理,但又不可以不理,因為這件「生命不再一樣」的事涉及大家的信仰:若果生來失明的原因在於有人犯罪,這人天生瞎眼如今得見,就無法不涉及上帝,例如涉及:罪得赦免的恩典.

所以,他們把他帶到法利賽人那裡.

2.3  轉:法利賽人盤問醫治的事

故事轉入第三段.這部份從第13節到34節,稍長的情節又可以分為三個分節.

第一節,法利賽人先盤問當事人──即係這個已經開眼得見光明的瞎子.

第二節,法利賽人再盤問當事人的家人.

第三節,當事人在法利賽人面前答辯.

第一節,法利賽人先盤問當事人.

13他們把從前瞎眼的人帶到法利賽人那裡.

14耶穌和泥開他眼睛的日子是安息日.

15法利賽人也問他是怎麼得看見的.瞎子對他們說:「他把泥抹在我的眼睛上,我去一洗,就看見了.」

16法利賽人中有的說:「這個人不是從神來的,因為他不守安息日.」又有人說:「一個罪人怎能行這樣的神蹟呢?」他們就起了紛爭.

17他們又對瞎子說:「他既然開了你的眼睛,你說他是怎樣的人呢?」他說:「是個先知.」

第二節,法利賽人再盤問當事人的家人.

18猶太人不信他從前是瞎眼,後來能看見的,等到叫了他的父母來,

19問他們說:「這是你們的兒子嗎?你們說他生來是瞎眼的,如今怎麼能看見了呢?」

20他父母回答說:「他是我們的兒子,生來就瞎眼,這是我們知道的.

21至於他如今怎麼能看見,我們卻不知道；是誰開了他的眼睛,我們也不知道.他已經成了人,你們問他吧,他自己必能說.」

22他父母說這話,是怕猶太人；因為猶太人已經商議定了,若有認耶穌是基督的,要把他趕出會堂.

23因此他父母說:「他已經成了人,你們問他吧.」

第三節,當事人在法利賽人面前答辯.

24所以法利賽人第二次叫了那從前瞎眼的人來,對他說:「你該將榮耀歸給神,我們知道這人是個罪人.」

25他說:「他是個罪人不是,我不知道；有一件事我知道,從前我是眼瞎的,如今能看見了.」

26他們就問他說:「他向你做甚麼?是怎麼開了你的眼睛呢?」

27他回答說:「我方才告訴你們,你們不聽,為甚麼又要聽呢?莫非你們也要作他的門徒嗎?」

28他們就罵他說:「你是他的門徒；我們是摩西的門徒.

29神對摩西說話是我們知道的；只是這個人,我們不知道他從哪裡來!」

30那人回答說:「他開了我的眼睛,你們竟不知道他從哪裡來,這真是奇怪!

31我們知道神不聽罪人,惟有敬奉神,遵行他旨意的,神才聽他.

32從創世以來,未曾聽見有人把生來是瞎子的眼睛開了.

33這人若不是從神來的,甚麼也不能做.」

34他們回答說:「你全然生在罪孽中,還要教訓我們嗎?」於是把他趕出去了.

原來最能令這個蒙恩的瞎子明白自己生命的不再一樣,竟然是不信此醫治瞎眼是神蹟的法利賽人,盤問的過程同時成為瞎子「屬靈視障」的超越.

法利賽人越質疑,瞎子越相信,越發發現第一身經歷的屬靈意義,大概經歷三個層次:

他由第一層的單純的肉眼看見,相信醫治他的這個未見到的人「(耶穌)是個先知.」

轉入第二層的理解,宣認自身的信仰:「能開我眼的一定是神的使者.」

再進入第三層,這個開我眼目的(耶穌)不可能是個罪人,「因為我們都知道神不聽罪人,惟有敬奉神,遵行他旨意的,神才聽他.」

甚麼叫見證(eye witness)?這段對答就是見證的好例子:這個本來失明的瞎子能夠看見,由此信仰這改變是神的作為,在種種質疑下用自己的第一身經驗加以證明這是神的作為.

這個瞎子見證說:「有一件事我知道,從前我是眼瞎的,如今能看見了.」

你呢?你可以見證甚麼?

你有無至少一件事你是知道的,從前你是眼瞎看不到神,如今能看見了?

你何時曾經超越你的「屬靈視障」?

2.4  合:靈性的盲目

最後一層的信仰體認,還是要返回到神那裡去.

讓我們再閱讀第35至41節這個故事的結束.

35耶穌聽說他們把他趕出去,後來遇見他,就說:「你信神的兒子嗎?」

36他回答說:「主啊,誰是神的兒子,叫我信他呢?」

37耶穌說:「你已經看見他,現在和你說話的就是他.」

38他說:「主啊,我信!」就拜耶穌.

這個被醫治被質詢的眼明新人,終於再見恩人.

耶穌親自渡引他的屬靈視力達到最高峰:「你現在親眼所見的就是神的兒子,現在直接和你說話.」

再見耶穌的重遇,就是超越「屬靈視障」後的生命高峰.

39耶穌說:「我為審判到這世上來,叫不能看見的,可以看見；能看見的,反瞎了眼.」

40同他在那裡的法利賽人聽見這話,就說:「難道我們也瞎了眼嗎?」

41耶穌對他們說:「你們若瞎了眼,就沒有罪了；但如今你們說『我們能看見』,所以你們的罪還在.」

自以為看見的屬靈視障,其實才真的是人的罪性.罪人,其實是屬靈失明人士.

明知有屬靈視障的虛心,反而不受罪性所盲目所障閉.救恩,其實是靈眼重見光明.

信仰讓生命不再一樣.

經歷過信仰讓生命不再一樣的你們,都是活在「靈眼得見上帝的光明之子」,得見上帝在世間事顯現的榮耀,所以,看世界已經不再一樣.

信仰就是「瞎眼今得看見」,生命不再一樣.

約翰福音第九章,講的是以為永久失明的絕望,變成早年失明但終於開眼看見的感恩.

約翰福音第九章,講的不單是一個瞎子重見光明的故事,而且還是一群自以為是的失明故事.

約翰福音第九章,講的不單是超越「屬靈視障」的信仰歷程,還有「屬靈視障」的失明故事.

約翰福音第九章,講的是第一身經驗建立信仰,一個我親身「見到」「證明」的「見證」(eye witness),親眼看見的屬靈視野就是證明.

3.  前我失明,今得看見:奇異恩典

歷代不止約翰福音第九章的這個瞎子見證說:「我只知道,我以前失明,今我看見.」

大家記不記得,有那一首歌的第一段是唱出這種信仰體驗?

對,是《奇異恩典》(Amazing Grace)的第一段,中文其中一個譯本是這樣唱的:

「奇異恩典,何等甘甜,我罪已得赦免；

前我失喪,今被尋回,瞎眼今得看見.」

Amazing grace! How sweet the sound.

That sav’d a wretch like me!

I once was lost, but now am found,

Was blind, but now I see.

這首大家愛唱的《奇異恩典》的詩詞原作者叫:John Newton(約翰·牛頓),寫於1772年.你若在網上尋索,例如打“奇異恩典”,維基(\url{http://zh.wikipedia.org/wiki/})就會出現以下的介紹:

「《奇異恩典》(Amazing Grace)是最著名的基督教聖詩之一.歌詞為約翰·牛頓所填,出現在威廉·科伯(William Cowper)及其他作曲家創作的讚美詩集Olney Hymns的一部分.

約翰·牛頓(1725年-1807年)曾經是一艘奴隸船的船長.在1748年5月10日歸家途中遇到了暴風雨,他經歷了“偉大的判決”.在他的航海日誌中他寫道:他的船處在即將沉沒的重大危險中,他喊道“主憐憫我們”.他的內心逐漸走向光明…….

在非洲一個港口等待奴隸裝船的時候,牛頓寫下了歌曲《耶穌之名如此美妙》.以後,他和奴隸貿易切斷了關係,當上了一名牧師,並且加入了威廉·威伯福斯領導的反奴隸運動.

這首讚美詩也在美國黑人基督教社會流傳.最初這首讚美詩選用過20多種旋律,到1835年的時候,借用了《新不列顛》的曲調,廣泛流傳直到今天.」

以下附錄全首歌詞,以及兩個不同的漢譯／修訂版本.

Amazing Grace《奇異恩典》

Amazing grace! (how sweet the sound)

That sav’d a wretch like me!

I once was lost, but now am found,

Was blind, but now I see.	奇異恩典,何等甘甜,

我罪已得赦免；

前我失喪,今被尋回,

瞎眼今得看見.

奇異恩典,何等甘甜,

我罪已得赦免；

前我失喪,今被尋回,

瞎眼今得看見.

‘Twas grace that taught my heart to fear,

And grace my fears reliev’d;

How precious did that grace appear,

The hour I first believ’d!	如此恩典,使我敬畏,

使我心得安慰；

初信之時,即蒙恩惠,

真是何等寶貴!

如此恩典,使我敬畏,

使我心得安慰；

初信之時,即蒙恩惠,

真是何等寶貴.

Thro’ many dangers, toils and snares,

I have already come;

‘Tis grace has brought me safe thus far,

And grace will lead me home.	許多危險,試煉網羅,

我已安然經過；

靠主恩典,安全不怕,

更引導我歸家.

經過許多,危險網羅,

飽受人間苦楚,

此恩領我,平安渡過,

他日歸回天府.

The Lord has promis’d good to me,

His word my hope secures;

He will my shield and portion be,

As long as life endures.

Yes, when this flesh and heart shall fail,

And mortal life shall cease;

I shall possess, within the veil,

A life of joy and peace.	將來禧年,聖徒歡聚,

恩光愛誼千年；

喜樂頌讚,在父座前,

深望那日快現

將來在天,安居萬年,

恩光如日普照,

好像最初,蒙恩景況,

讚美永不減少.

The earth shall soon dissolve like snow,

The sun forbear to shine;

But God, who call’d me here below,

Will be forever mine

以上引述到領導反奴隸運動的威廉·威伯福斯(William Wilberforce),是受到John Newton的信仰生命所影響,所感召.他們的故事曾經在2007年拍成電影,電影也是命名為《奇異恩典》(Amazing Grace).(可上相關網站:\url{http://www.amazinggracemovie.com}.)

故事有兩個片段非常感人.有一段是John Newton晚年漸漸失明,William Wilberforce探望他時,John Newton反而分享他越來越看得清楚上帝.這片段詮釋了這樣的屬靈視覺:「原以為看見實瞎眼的前生,因神的恩待,如今瞎眼反得看見」.

另一段是William Wilberforce在紳士俱樂部高歌聖詩《奇異恩典》以明心志,與戰友一同廢除奴隸貿易,讓其他人生命不再一樣.

信仰讓生命不再一樣,同樣的基督信仰,解放了John Newton 的靈性失明,解放了William Wilberforce職志猶豫,解放了英國的奴隸貿易,也解放了全球的黑奴.

4. 有牧者的生命多麼不一樣:詩篇23篇

當我準備今次講道時,發現約翰福音第9章和詩篇23篇兩段經課都很值得講.

詩篇23篇這段經文,是我初信之時,便已經非常喜歡的舊約經文.

剛才因為已經用了大篇幅講解約翰福音第9章,現在用另一個方式去思考詩篇23篇.

請大家一同打開聖經到詩篇23篇.

現在不用講解方式,讓我們一起用一種像在山谷中三重迴響式的朗讀方式,一同用廣東話讀出和合本的經文,讀法是我讀一句,弟兄重複讀同一句,姊妹再重複讀同一句:

23:1	耶和華是我的牧者．我必不至缺乏.

23:2	他使我躺臥在青草地上,領我在可安歇的水邊.

23:3	他使我的靈魂蘇醒,為自己的名引導我走義路.

23:4	我雖然行過死蔭的幽谷,也不怕遭害．因為你與我同在．你的杖,你的竿,都安慰我.

23:5	在我敵人面前,你為我擺設筵席．你用油膏了我的頭,使我的福杯滿溢.

23:6	我一生一世必有恩惠慈愛隨著我．我且要住在耶和華的殿中,直到永遠.

這段經文在讀出時,你有沒有體會一種「我有牧者真好」的生命感?

被好牧者所牧養,羊就不再受自己的盲目所困,不會走迷.

不單一個一個的生命需要牧者,整個香港社會都需要牧者,因為香港正是一個像群羊在走迷的社會,充滿著不同型態的「屬靈視障」.

沒有信仰,人可以如何面對困苦的人生境況?沒有信仰,如何面對殘酷的社會現實?

劉進圖先生被刀手所斬.在急救床上,重新通過信仰,反思生命.

劉進圖這位弟兄可以相信甚麼?我們可以相信甚麼?

劉弟兄在他【劉進圖感言之六:受苦的意義 生命的反思】中分享到:

(source: \url{http://inews.mingpao.com/htm/Inews/20140317/aa11008y.htm})

「一個新聞工作者無辜遇襲,引起了社會巨大迴響,這是完全出乎我意料的.…………

除了協助調查,我在病榻上比較上心的,是求問上帝,容許我遇上這一劫,到底有什麼深意?我和太太都是基督徒,相信「萬事互相效力,叫愛神的人得益處」,但無端被斬傷,要花很多個月才能重新行走,這樣的事對我有什麼益處?盧龍光牧師是我和太太的證婚牧師,他來病房為我祈禱時說了一句話,希望上帝讓我在病房內學會人生更深刻的一課,這句話讓我反覆咀嚼,正是我內心深處的渴求.…………

我從少年時期就返教會,學會人生目標應是榮神益人,努力事奉.過去我一直認為,事奉是奉獻自己的才智,能力和時間,做有益於社會的事情,自然就能榮神益人.例如,我努力讀書,考取好成績,便是向世人作好見證.又例如,我捨棄較賺錢的法律專業,為興趣和使命跑去當記者,還盡量抽時間教學及服公職,這就是努力事奉,榮神益人.

如果我承認今次我遇襲產生了積極的社會意義,那就意味我是通過一個非自願的,被動的,與個人才智能力全不相干的途徑,單純靠承受傷害和痛楚來造就別人,這完全違反了我一貫的概念.其實,我這個人頗有點精英主義,上司曾批評我恃才傲物,我從來沒有想過,放下聰明才智,不再勞心勞力,人還能做什麼?我從來不曾明白,所謂無權勢者的權力,受苦的事奉,到底是什麼意思?我大半生走過的所謂榮神益人努力事奉的路,到底當中有多少真實而深刻的意義?有多少其實是為了滿足自我的虛榮?這一連串的問題,在靜夜無眠的晚上,折磨著我的心.

我沒有答案.但我學會一件事,我需要謙卑下來,重新思考人生的方向和意義.」

順著今天講道的關懷,這位走過死蔭幽谷也有點怕遭害的羊──劉弟兄,在引領我們一同展開跨越「屬靈視障」的上行,盼望看到更深更闊的視野.

新聞人的視野是新聞自由的信仰.全香港人都會相信這是我們的核心價值.

核心價值還包括:全民普選行政首長和立法議員的民主.

究竟雖千萬人吾往矣的「佔中三子」信仰甚麼?

朱耀明牧師當然相信上帝,並且願意一生成為上帝的僕人,牧養迷失的羊群.

戴耀廷教授也相信上帝,相信公義臨到人間,自身甘願成為群羊之首.

雖然長年在基督教信群中生活,陳健民教授不承認自己相信基督宗教.

當這個學期有次他在這崇基禮拜堂作週會講者時,有人問到:「若然有關權貴要傷害你時,你還會相信甚麼?」陳健民教授答:「我會回頭,望向這個十字架.」

請大家一同,望向崇基禮拜堂這個十字架,對你來說,代表了甚麼?

對你來說,十字架,代表了甚麼?

5.　地上的爸爸,再見,在天父的家

對於我,十架回望,我總想嚎啕大哭,我總想恩臨大地,我總想為神奮戰.

回望我的人生,我經驗過三次學業的重大挫折,三次事業的非常挫敗,三次健康的深切困纏,到近兩年工作也特別地不暢快.

我跟內子和兒子曾經講過:「經歷苦困,遭逢惡待,若然我沒有信仰,我早已變成大惡人.」

經歷人間苦困,正是信仰生命的功課.

近月,我在學習告別的功課.

首先,在告別一種時代的氣氛,感到難以繼續堅持.

其次,在告別我所創建的兩個碩士,被解甲,不再由我主領.

更傷的是,自從十日前開始,在告別父親,學習在一種哀愁中生活.

我返回中大在教育學院教書,已經進入第二十年.

我的教學精神是:「每次教課,可能最後一課.」所以,我會盡心盡力教好每一課,讓學生從我教學生命所得到的會是最好的.

正如我的人生態度是:「每次相見,可能最後一見.」所以,難離難捨,總有一些.

上周三到醫院探望病床上的父親時,竟然是最後一次相見.一天之後,他便離世了.

近年家父身體抱恙多時,3月20日晚上病發辭世,主懷安息,平安作別,享年近89歲.

我需要在這裡作一次告解,一次懺悔.

家父姓鄭名錦潮,廣東省潮陽縣人.他與林碧蘭女士結婚,育有5子1女,我排行第五.

我是全家第一個信耶穌的,在1976年信主時我剛升讀中五,1977年一過十八歲我便洗禮,為我施洗的是當時循道衛理聯合教會香港堂的主任牧師李炳光牧師.

除了數年後我四哥也信主,我對家人在信仰上的影響,可謂乏善足陳.

2007年十月,我入了教會一個叫以馬忤斯之旅的三天信仰營後,經歷「屬靈視障」的醫治──重新親眼遇見耶穌.營內我立定了數個心志,其中一個是我不再成為家人信主的障礙──特別是不希望我的第一個信耶穌的言行反而構成家父不信耶穌的障礙.

出營後,我立即和內子以及兩個兒子商量,再與四哥四嫂及其兩兒子協作,再加上其信耶穌的第三代,一行十多位基督徒,在父親十一月生日那天,邀請家父以及全家未信耶穌的十多位,一同前赴政府大球場(所在山谷的正民村其實是我們的故鄉),參加2007年12月「香港葛福臨佈道大會」.

有趣的是,在步行進入政府大球場的路上,遇上了李炳光牧師,李牧師立即為在路旁的家父祈禱.在是次「香港葛福臨佈道大會」,家父走落草地,跟大會祈禱,決志信主.信主後,我們一家現在所聚會的馬鞍山堂主任牧師陳建基牧師曾經親身探望他,確定他信仰並與他一同祈禱.去年十二月底,在退休前的前任校牧伍渭文牧師為家父的長孫女Angie和長孫婿Sam在這個崇基禮拜堂主持婚禮.

所以,4月13日晚上的安息禮,我們邀請了這三位牧師(主持,講道,祝福),與我們一起與家父告別.

有牧者的人生多好,有信仰的生命不再一樣.

信仰讓生命不再一樣.

除了對人生的苦難困阨有不一樣的屬靈視野外,還對生命的所持有的感到上帝的恩待.

信仰讓我們反思上帝給我們的恩賜,你有何恩賜?

我相信在天上的父親給我講笑,講學,講道這三講能力.不知何時有機會能夠「三講合一」?這其實是我讀神學的心志(開講有宣教意義的教育棟篤笑).

但是,在地上的父親從來未聽過我講學,也未聽過我講道.

留下可堪追憶的,只有他曾經聽過我講笑的片段,留下的,是他在笑的記憶.

今天講道事前,我得到許開明校牧容許,讓我在這次講道念及家父.

我謹願呈獻此講章給天父,作家父返回天家的印記:「地上的爸爸,再見,在天父的家.」

但願榮耀都歸給聖父,聖子,聖靈,永在的三一神.阿們.

(鄭漢文博士為循道衛理聯合教會馬鞍山堂會友.

現職香港中文大學教育行政與政策學系專業顧問,

同時在崇基學院神學院兼讀神道學學士課程.

曾於2012年1月至6月及2012年9月至2013年4月於崇基禮拜堂作神學生實習.

現於2013年9月至2014年4月在浸會大學校牧處作第二次神學生實習.)



\newpage

\section{作個有用的人}
\label{sec:240}
\textbf{許開明牧師 - 作個有用的人}
\newline
\newline
\hyperref[sec:239]{< < < PREV SERMON < < <}
~
\hyperref[sec:toc]{[返目錄]}
~
\hyperref[sec:241]{> > > NEXT SERMON > > >}


講題:作個有用的人　Being a Useful Person

講員:許開明牧師

場合:香港中文大學崇基學院禮拜堂

日期:2014 年4月13日

AUDIO

引言

今日是棕枝主日,紀念二千多年前,主耶穌乘驢進入聖城耶路撒冷.群眾拿著棕枝歡呼 ,歡迎,合城的人都驚動起來.這事件四福音均記載,可見是非常重要的事跡.今天是紀念主耶穌進入聖城的主日,願我們一同回顧和思想.

【】在四十多年前,我住在官塘最旺的一條街,住的大廈對面有間著名百貨公司,記得開幕時,請了一位明星﹝江漢﹞來展彩,十分震動,許多圍著中匯,我你和許多人都打開窗一賭明星的風采,年少的我看到這明星,十分開心!好像今日的韓星「金秀賢」般.

你想像一下,耶穌可說是當代的名人,大英雄像「明星」般進入耶京,多麼轟動,如果耶穌基督到訪崇基禮拜堂,場面也會很偉大!

然而,在這段經文中,也有下列迷思?我們需要了解,才能明白這段經文.

經文的迷思:

1為何耶穌騎驢進城,不驢馬?

騎馬,不是更威風,更引人注目?

2為何耶穌騎驢駒,不騎母驢?

耶穌那麼重,壓扁驢仔!

3為何耶穌高調進入耶京?

耶穌不是告訴人,不要透露他彌賽亞的身份?耶穌作風不是一向低調的嗎?

4為何群眾如此興奮,熱烈的歡迎,合城震動起來!後來卻要釘死耶穌?

由非常歡迎到非常厭惡拒絕,真是匪夷所思,令人難明?

經文的迷思,我只能作簡單的解釋:

1為何耶穌騎驢進城,不驢馬?

騎馬,上戰場,代表戰爭!

騎驢,載貨物,代表回鄉或進城!

『驢』是溫馴的動物,古猶太人在未普遍使用馬之前,曾被用作君王的座騎(參士十4；十二14；撒下十六2).由於牠不像戰嗎哪樣威風,所以聖經用驢和馬作對比,以表明騎驢者的謙和,溫雅(參亞九9~10).

這也是應聖經的預言說:

【太廿一5】「『要對錫安的居民(原文作女子)說:

“看哪,你的王來到你這裡,是溫柔的,又騎著驢,就是騎著驢駒子.”』」

2為何耶穌騎驢駒,不騎母驢?

驢駒是驢仔,,母驢仍牠身旁的,不能承擔太重,耶穌應有五呎十寸(180),170磅,不輕的.

原因是:耶穌騎著驢駒進入耶路撒冷,這隻從來沒有人騎過的驢駒,特別適合於神聖的用途.

在潔淨禮中所使用的純紅的母牛犢,一定是『未曾負軛』的(民十九2,申廿一3)；載負耶和華約櫃的牛車必須是從來未作別用的牛車(撒上六7).

這場合特別神聖的氣氛,藉著這從來未有人騎過的驢駒表明出來.

3為何耶穌高調進入耶京?

耶穌一向低調的,這次騎驢進入,跟他的低調的作風不太一樣.

耶穌選擇在人多熱鬧之時進城,這與祂平常謙虛,低調的作風不太一樣；主耶穌似乎精心策劃進城一事.他有意接受群眾們用最熱情,最尊崇的方式迎接祂,歡迎祂作王,耶穌清楚表明他自己是彌賽亞,是基督.是以色列人期待的那一.

耶穌清楚向猶太人表明身份,不要說:耶穌塑有表明,至於接受與否?決定在群眾.

然而,主耶穌進城不是要作世上的王,而是要受苦,被釘死在十字架上,完成祂救贖世人的使命.

主耶穌在世上最後的幾天,應驗舊約先知的話騎驢進城,象徵祂有王權,祂不採用武力,而是謙謙和和地騎著驢,帶著愛進城,並且勇敢地走向各各他十字架犧牲的道路.祂所要建立是和平,仁愛的國度,這就是主耶穌傳播的資訊.

4為何群眾如此興奮,熱烈的歡迎,合城震動起來!後來卻要釘死耶穌,除去他們曾是最擁戴,尊敬的人呢?

由非常歡迎到非常厭惡拒絕,真是匪夷所思,令人難明?

十分明顯,因為耶穌不是他們心中的王?

猶太人期待的是一位西門馬加比般的耶穌.

群眾揮舞棕樹枝迎接英雄凱旋歸亦發生於西門馬加比進人耶路撒冷時.

當時羅馬將軍安提阿哥佔領了耶路撒冷城,決意要除滅猶太教,將希臘的生活方式與崇拜帶進巴勒斯坦.

他故意汙穢聖殿,在祭壇上獻豬肉,向奧林皮烏山的丟斯神獻祭,甚至把聖殿的會所改作公共的妓院.

後來馬加比起來反抗,收復了他們的國土.

耶穌也潔淨聖殿,卻指責在聖殿以不法手段,擺賣圖利的人.耶穌希望聖殿是禱告的殿,是敬拜上帝,冒犯了許既得利益者或持份者的猶太人.他們心中的彌賽亞不是那位侵犯他們利益的人,而是要使他們獲取更多利益,打退羅馬人,不用立稅給該撒的人.

民眾需要的也是像馬加比的武功高度強,又能招兵買馬,動員群眾,恢復以色列河山的人.耶穌應該比馬加比更厲害,因為祂是神的兒子,又能行神蹟.耶穌來到聖城可以大顯身手,把羅馬兵打到落花流水,然後恢復我山河!

怎料耶穌來到聖殿是推倒買賣的人,與法利賽人辯論,跟被捉拿,被戲弄,一點兒馬加比的氣質都沒有!耶穌看來沒有用.

因此,由擁戴耶穌,變成釘死耶穌.

在本段經文:探討「作一個有用的人」

你有用嗎?

我有用

我都有D用

我無用喇!

我真無用!

1一般人理解「有無用」,在於「我自己有無用」?

我有什麼才幹?我有什麼能力?我有什麼的成就?

從不同的角度,去探索成為有用的人

哲學家說「認識你自己」,知道你自已的強弱點,特點,盲點,做有用的人

心理學家說「接受你自己」,接納自己是個怎樣人,自我限制,做有用的人

教育學家說「肯定你自己」,肯定自我的價值,自我的潛能,做有用的人

宗教家說「超越你自己」,超越自我品格,能力,智慧限制,做有用的人

基督教的信仰說「捨棄你自己」,為主所用(太十六24)不靠自己,仰望上主,成為有用的人

近代學者認為:成為有用的人是一位能夠「實現自我的人」

心理學家羅吉斯(Carl Rogers)相信:「自我實現的傾向」促使每個人都有能力發展自己的潛能,以維持或改善自己的生活,這是我們與生俱來就有追求成長的基本需求.

另一位心理學家馬斯洛(Abraham Maslow),以提出生理,安全,歸屬與接納,自尊,自我實現五個階層漸進的「需求層級理論」著稱,他主張這些逐步滿足的需求能建立良好的基礎,最終成為一個潛能完全發揮,人格協調一致的「自我實現的人」(self-actualized individual).

一個「自我實現的人」,在日常生活中展現出來的,也就是一個「有自信的人」,自覺「有用的人」.

自信顯示個人對於自我概念與外在環境有正向客觀的評價.一個有自信的人對於自己有充份的自覺,信心和勇氣,敢於嘗試並接受挑戰,對於外在的突發狀況能夠很快抉擇反應,自我修正,適當處理,每次妥善處理完一個挑戰之後,會有成長的滿足感,反饋給自己而良性循環得到更有自信的感覺.

2我會被人用?人家覺得我有用嗎?

所謂「有用」,不獨是「我有無用?」

人家認為我有用嗎?我感到「人家認為我有用?」

【】老板告知你:「你好幫得我手,無你不行!」

(你老板看來會用到你盡!不過你也會很開心!)

許多人的自我價值,有無用決定於人家的看法?

愚昧人用我,可能不識用我,當我也無用!

聰明人用我,可能用到盡；或覺得我無用,不用我!

懶人用我,可能用盡我!自已繼續懶下去!

糊塗人用我,不懂我的職責或才能,亂用的我!激死我!

人看人,有無用?基於他對你的期望.

一個人有無用,建基於別人的看法,是很危險的.

3基督徒問:我肯被上帝用?耶穌怎樣用我?

在本段經文,我們看到:

1)門徒:聽耶穌話,成為有用的人

【太廿一2】「耶穌就打發兩個門徒,對他們說:『你們往對面村子裡去,必看見一匹驢拴在那裡,還有驢駒同在一處；你們解開牽到我這裡來.」

【太廿一3】「若有人對你們說甚麼,你們就說:“主要用牠.”那人必立時讓你們牽來.』」

【太廿一6】「門徒就照耶穌所吩咐的去行,」

【太廿一7】「牽了驢和驢駒來,把自己的衣服搭在上面,耶穌就騎上.」

「耶穌就騎上,」按照原文字面,好像耶穌是騎在兩匹驢(牠們)上面,其實一個人不可能同時騎上兩匹驢,祂應是騎上驢駒(參可十一7),而那匹驢駒因從未被人騎過(參可十一2),須有其母陪伴在旁,亦步亦趨.

我們聽從主,成為有用的人.

【慕迪Moody】二十世紀,上帝所重用的佈道家.

他在一次青年退修會,聽到上帝的呼召,起來向大家說:

「我身上288磅肉,磅磅都是主的.」因而委身宣教.

2)驢的主人:聽門徒的話,把有的驢及驢駒給主用.

這主人也相當慷慨,把兩驢拿出來,相當貴重的東西,卻成就美事.

把有的東西,有錢出錢,有力出力,校牧室需要大家奉獻支持:

崇拜事工:講道,音樂,程序,主日學,查經班

學生事工:本科生,研究生.團契,輔導,查經小組

內地生:本科生,研究生.團契,輔導,查經小組

要有效做好牧養,需要聘請人手,歡迎奉獻.

【】一杯涼水

對於旅行者,非常頸渴的人,非常重要,可救一命!

要上山,從山溪水

要到井邊打水,取井水飲用.

3)驢駒:肯讓耶穌坐上,成為有用人

一般人的理解:

這驢駒更是從來沒有人騎過的,似乎是沒有用的.

「驢駒」是指尚未長成的小驢子,提到的這匹小驢駒,更是從來沒有人騎過的,似乎是沒有用的,但是「主要用牠」.現在小驢駒承擔很重要的任務--背負著主耶穌,成就神的工作.

從這匹小驢駒的身上,我看到主要用的人,是:肯被主用的人,平凡普通的人,聖潔謙卑的人,以及願交主權的人.

這隻小驢駒肯被主耶穌使用,順服地讓主騎乘.主不看我們能不能,是問我們肯不肯.

這匹小驢駒是平凡,不顯眼的,好像沒什麼價值,但是當主使用牠,牠的價值就不一樣了.

主若肯用我們,而我們又願意給主使用時,我們的價值就能顯露出來.

驢子沒什麼大本事,身高沒有馬高,速度沒有馬快,載運的貨物沒有馬,駱駝多,就是這樣平平凡凡的驢生,卻因著找運者是這位奉主名來的王是應當稱頌的,而得了當得的榮耀.

上帝常使用年青的,像青的無花果(伯法其』原文字義是『尚未成熟的無花果之家』)

驢駒:小驢仔,無人騎過的,較可愛.

4)棕枝:有多種的用途.這代表能幹的人,滿有恩賜才能的人.

▲   棕樹的全身360种功用,它的果子,葉子,枝子,汁,根,,,,,皆有用.

▲   生長在巴勒斯坦棕樹的果子,可以食,對人體有無比的價值與功能.

1棕樹他會結棗子,這一個棗子甜甜的,他不是結一點點,而是結一串串很多的棗子,

2棕樹有一些特點,它是長在沙漠的地方,它的樹形很美麗,聖殿中間,非常多的裝飾,是用棕樹作裝飾,木板上,門板上都刻著是棕樹,你看到中東一帶,一些古老的建築,柱子最上面刻著就是棕樹,

3棕樹叫「他馬」他們有一些女孩子的名子取名叫他馬,不是她媽,聖經中猶大的媳婦叫做他馬,雅各有個女兒也叫他馬,大衛有個女兒叫做他馬,他馬是什麼意思,他馬就是像那個棕樹一樣,那樣子的美麗,婀娜多姿,棕樹樹形,在中東一帶是美麗的代表.

4它結的果實非常的多,它不是一點點果實,它有一點像結檳榔,一串一串,棕樹結的棗子甜的不得了,那個棗子做成棗子醬,就跟蜜一樣的甜,把花苞刺破就流出許多的汁漿來,那一個汁漿做出來就像蜂蜜一樣,幾乎所有的解經家都相信,摩西所講的帶你們進入流乃與蜜之地,這一個蜜就指的是這一個,因為他就像蜜棗一樣甜的很,這就是棕樹,

5它的樹齡很長,可以活到二百年,蘋果樹大概三十年,香蕉是結了一年就要砍掉,我不知道芒果是多少年,葡萄樹也是三十年,就不能再長了.

棕樹可以種二百年,所以你種一顆棕樹,世世代代有的吃,聖經中說到了年老還結果子,葉子也不枯乾,汁漿非常的豐富,我們一起看詩篇92:14 他們年老的時候仍要結果子,要滿了汁漿而常發青,就是只著這一個,這一個棕樹.

【】有些像棕樹,多才多藝為主大大的使用.

5耶穌:願獻己身為上帝所用,成為極有用的人

耶穌是個有用的人:有計劃的人

(1)把握時間:在逾越節,最多入的時間入城.

(2)把握地點:在聖殿,最多人娶的地方.

(3)把握歷史的場景:學效馬加比進城的故事

(4)把握舊約的預言:

以賽亞和撒迦利亞的預言:

要對錫安的女子說:

「看哪,你的王來到你這裡!

是溫柔的,又騎著驢,

就是騎著驢的駒子.」

(5)把握時間事前作安排:

兩個門徒,驢主人,驢駒

(6)把握群眾的舉動,顯示祂彌賽亞的真實身份

二一7,8 主騎驢進城,踏在鋪著衣服和棕樹枝的路上,

群眾的歡呼聲不絕縈繞耳際.至少這刻,人們公認他是王.

二一9 眾人喊著說:「和散那歸於大衛的子孫,奉主名來的,是應當稱頌的.」

這節是引自詩篇一百一十八篇25至26節,明顯指到歡迎彌賽亞的到來.

和散那原本的意思是「現在就拯救」；也許眾人的心意是:「拯救我們脫離羅馬欺壓的政權吧.」下半節是頌讚感歎.

「大衛的子孫」和「奉主名來的是應當稱頌的」清楚指出耶穌被公認為彌賽亞,是應當稱頌的那位──他行使耶和華的權柄,完成神的意旨.

(7)有用的人,經得起挫敗的人

耶穌經歷民眾歡呼過後,卻是苦難的十字架?由高潮走向低潮?

一個人最大的打擊是從上跌下來

從光榮到羞辱

從成功到失敗

從歡呼到喝采

【】記得多年前,我完成教堂重建,由四層變十四層,我向區會請求想休息,並想接受高等學位的資助往英國大學進修,當時已有三間英國大學會收我讀書.我於2001年到伯明翰.全家陪阿爸讀書,去到英國突然有種挫敗感:

堂主任,校監,多個基督教機構,部門主席.

講一聲,事就這樣成了!

教會重建,困難重重,意見分歧,也逐一解決.

相當有自豪感；生養也不錯3個,出書也不錯10本.

到英國讀書,突然,覺得自己好無用,無能!

太太:不要讀書大晒,無人服事你呀!

考車一次一次失敗；成功後,仔女話不放心!

我:只讀書,教授嚴格,難通過!

太太:返工,做家務,園丁:很受仔女尊重

大女:阿爸,你只會讀書!

細女:為何你不用返工?日日都是讀書!

主耶穌經得起挫折,從折折中站起來!

經歷十架的痛苦和羞辱,死亡的無奈；才有復活的人喜樂和榮耀.

至於我們應該怎麼做個有用的呢?

1肯定上帝做我,我有存在的價值

【】一位少年人想學牛頓,在蘋果樹下默想,

等了很久,都沒有甚麼悟到?

於是他望著蘋果樹很大,很高,但蘋果很細小,為什麼?

於是他又周圍望,望到園中有個大西瓜,卻長在小小的藤上及地上；當中一定有原因.但他想了久都不通.突然有個蘋果掉下來,落在他的頭上,他立時醒悟了.如果蘋果太大掉下來會傷人的頭；,西瓜太重應生在地上.

上帝創造的萬物,都有一定的價值和理由.

因此,他也有價值.

2肯定自己,立定目標.我們要做個有用的人,或為別人所用,為上帝所用,做個有用的人,就要先認清「自己」,肯定「自我」的力量,建立起自信心.然後,再立定志向,想想以後要做什麼,發奮追求理想,努力自我實現.

3肯定起來行動,勇敢去作.

要實際行動.俗語說:「坐而言,不如起而行.」

當我們認清人生的目標後,就要靠自己不斷的努力,努力成為有用的人.

1不要怕.懼怕,會變成無用的人.

生手怕熟手,熟手怕高手,高手怕失手.

床頭驚鬼,床尾驚賊!(去香港,去英國)

2有計劃:耶穌進城滿有計劃:在逾越節,進城門,準備驢,門徒,主人等.

3懂用人:門徒,主人,驢駒,棕樹,群眾.

結語

【】人生像三種的船

棄船:無用的船,停在港口,阻礙其他的船行走.

帆船:順風而行,沒有自己.

電船:船有摩打,有目的而行,好好的順從舵手行.

耶穌自己是一個很有用的人,他也善於用人,並且甘願把自已獻給上帝所用；

我們要成為真正有用的人,就是願意被上帝所用.

馬偕博士有一句著名的話「寧為燒盡,不願鏽壞」,願為主所盡用!

【】最後以一個故事作結束,

一個弟兄願為神所用,但他剛考上大學,於是與神討價還價要多裝備等畢業再來事奉.

畢業後,主來到他面前要差派他,他則想先讀完碩士再出來事奉.主耶穌沒說一句話,拿起一朵玫瑰花撥掉幾片花瓣.

這弟兄受完訓碩士,主來到他面前,他卻想等事業有成再事奉；主沒說一句話,拿起那朵花再撥掉幾片花瓣.

等他事業有成,他又想等兩個女兒大學畢業,無家庭顧慮再出來事奉；主拿起那朵花再撥掉幾片花瓣.

過幾年女兒大學畢業,他也老了,他又想過兩年退休後再出來事奉；主拿起那朵花撥掉最後的幾片花瓣.

他退休後很高興地來到主面前,說:“我可以把一生擺上了.”主拿出那朵花說:“這花代表你的生命,他已經不是原來的花,你還能給我什麼?”

人生有限,時間不等人,當還有時間應把握機會,快來事奉主.

最近我發現香港的神學院,好像我當年廿歲出頭去讀神學的愈來愈少,神學院學生老化,四五十歲才讀的愈來愈多,這也不是不好.只是如果是受職的同工,做不久便退休,或者今日最需要的兒童及青少年工作,找個年紀太大的也不大合適且疲於命.

【】週報文章:「生命,讓上帝叫ｃｕｔ!」

很有意思,不要自己ＣＵＴ自己,讓上帝來ＣＵＴ!

導演在拍戲時,叫一聲ＣＵＴ,就立刻停下來.

上帝對我們講聲ＣＵＴ,生命就完結!我們不想上帝講ＣＵＴ,誰都怕這聲ＣＵＴ,然而人生始終有日上帝會對我講:ＣＵＴ.



\newpage

\section{以馬忤斯晚餐的光影}
\label{sec:241}
\textbf{伍渭文牧師 - 以馬忤斯晚餐的光影}
\newline
\newline
\hyperref[sec:240]{< < < PREV SERMON < < <}
~
\hyperref[sec:toc]{[返目錄]}
~
\hyperref[sec:242]{> > > NEXT SERMON > > >}


路加福音 24:13-24
\newline
\begin{longtable}{cl}
\hline
\hline
章節 & 經文 (和合本修訂版)\\
\hline
24:13 & \begin{tabularx}{0.7\textwidth}{X} 同一天，門徒中有兩個人往一個村子去；這村子名叫以馬忤斯，離耶路撒冷約有二十五里。 \end{tabularx} \\ \\ \relax
24:14 & \begin{tabularx}{0.7\textwidth}{X} 他們彼此談論所發生的這一切事。 \end{tabularx} \\ \\ \relax
24:15 & \begin{tabularx}{0.7\textwidth}{X} 正交談議論的時候，耶穌親自走近他們，和他們同行， \end{tabularx} \\ \\ \relax
24:16 & \begin{tabularx}{0.7\textwidth}{X} 可是他們的眼睛模糊了，沒認出他。 \end{tabularx} \\ \\ \relax
24:17 & \begin{tabularx}{0.7\textwidth}{X} 耶穌對他們說：「你們一邊走一邊談，彼此談論的是甚麼事呢？」他們就站住，臉上帶著愁容。 \end{tabularx} \\ \\ \relax
24:18 & \begin{tabularx}{0.7\textwidth}{X} 兩人中有一個名叫革流巴的回答：「你是在耶路撒冷的旅客中，惟一還不知道這幾天在那裡發生了甚麼事的人嗎？」 \end{tabularx} \\ \\ \relax
24:19 & \begin{tabularx}{0.7\textwidth}{X} 耶穌對他們說：「甚麼事呢？」他們對他說：「就是拿撒勒人耶穌的事。他是個先知，在神和眾百姓面前，說話行事都大有能力。 \end{tabularx} \\ \\ \relax
24:20 & \begin{tabularx}{0.7\textwidth}{X} 祭司長們和我們的官長竟把他解去，定了死罪，釘在十字架上。 \end{tabularx} \\ \\ \relax
24:21 & \begin{tabularx}{0.7\textwidth}{X} 但我們素來所盼望要救贖以色列民的就是他。不但如此，這些事發生到現在已經三天了。 \end{tabularx} \\ \\ \relax
24:22 & \begin{tabularx}{0.7\textwidth}{X} 還有，我們中間的幾個婦女使我們驚奇：她們清早去了墳墓， \end{tabularx} \\ \\ \relax
24:23 & \begin{tabularx}{0.7\textwidth}{X} 不見他的身體，就回來告訴我們，說她們看見了天使顯現，說他活了。 \end{tabularx} \\ \\ \relax
24:24 & \begin{tabularx}{0.7\textwidth}{X} 又有我們的幾個人往墳墓那裡去，所發現的正如婦女們所說的，只是沒有看見他。」 \end{tabularx} \\ \\
[1ex]
\hline
\hline
\end{longtable}
約翰福音 20:21-22
\newline
\begin{longtable}{cl}
\hline
\hline
章節 & 經文 (和合本修訂版)\\
\hline
20:21 & \begin{tabularx}{0.7\textwidth}{X} 於是耶穌又對他們說：「願你們平安！父怎樣差遣了我，我也照樣差遣你們。」 \end{tabularx} \\ \\ \relax
20:22 & \begin{tabularx}{0.7\textwidth}{X} 說了這話，他向他們吹一口氣，說：「領受聖靈吧！ \end{tabularx} \\ \\
[1ex]
\hline
\hline
\end{longtable}
講題:以馬忤斯晚餐的光影-餐桌,小崗,與仄徑The Light  and Shadows of the Supper at Emmaus-The Table, the Small Hill, and the Narrow Road

經文:路加福音24:13-24;約翰福音20:21-22

講員:伍渭文牧師

場合:香港中文大學崇基學院禮拜堂

日期:2014年4月20日

大家看看這兩幅卡拉瓦喬(Caravaggio) 的名畫-以馬忤斯的晚餐Supper at Emmaus.卡拉瓦喬1601年這幅是倫敦大英博物館鎮館之作；1606年這張藏於義大利米蘭布雷拉畫廊(Pinacoteca di Brera),由香港賽馬會贊助港幣482萬最近在香港展出,展出場地是金鐘亞洲協會香港中心中心 (Asia Society Hong Kong Center),保險費達六億港元.

大家看看兩幅同一主題的畫作有何不同?

卡拉瓦喬是緊接文藝復興後巴洛克(Baroque)時期的大師,文藝復興(Renaissance) 時期畫作強調人類理想的完美,人物看來優雅脫俗,文藝復興響往古希臘和古羅馬對人性的理想化,文藝復興的人板如達芬奇,樣樣皆能,是畫家,雕刻家,懂醫術,亦是數學家,天文學家；是巧手的工匠,又是美學大師,思想深邃,學識淵博.

但巴洛克時期比較重真實的人性,不再是理想的投射,是現實的描繪,有血有肉.卡拉瓦喬善於運用黑暗,托出明亮,畫的背景被黑暗遮蓋,使人集中注視人物的神態和臉相,但第一幅畫也有文藝復興時期的優雅臉相.

贖罪的畫作

但第二幅耶穌臉相顯得疲憊,憔悴.再看第一幅的桌上擺設,食物比較豐富,第二幅貧乏,而且氣氛灰暗,善於利用光影的卡拉瓦喬要在第二幅畫作表達甚麼呢?

前一段時間我在電視第一次看到有關展覽的介紹,說這是作者表達尋求上帝救贖的作品,就想到用「以馬忤斯晚餐的光影」作復活主日的講道題目,因為復活就是救贖的確據；但後來發現展覽的主題也是這題目,所以我今天的題目加了副題-餐桌,小崗,與仄徑.在這這畫作,作者把自己的心靈疲憊,投射在基督身上,基督負負我們的重擔.不錯,聖經豈不說過:「他被掛在木頭上,親身擔當了我們的罪.」(彼前2: 24)繪畫第二幅畫作時,卡拉瓦喬是一位被追緝的殺人犯.[i]

以馬忤斯餐桌的團契

卡拉瓦喬和以後受他影響的畫家很多都以以馬忤斯的晚餐為主題作畫,一個原因是這主題是復活的基督和兩位不見經傳的平民百姓共進晚餐,其中一位叫革流巴(Cleopas) ,祇出現過一次,另一位名叫甚麼聖經都沒有交待.而為使徒們作畫則要非常小心,第一次他被委約畫使徒馬太就被退回,教會不能接受他把使徒馬太還原為有血有肉的稅吏；但他重繪的馬太,腳跟還是帶點骯髒,作為抗議.

另一個採取以馬忤斯為畫作的原因,是畫家們都接受信仰傳統,認為這不是普通的晚餐,是耶穌和兩位門徒共享聖餐.聖經記載:「到了坐席的時候,耶穌拿起餅來,祝謝了,擘開,遞給他們.他們眼睛明亮了,這才認出他來.」(路24: 30, 31)從耶穌拿起餅來,祝謝,擘開,遞給門徒的動作,兩位門徒就立時認出這是耶穌,因為這是耶穌設立聖餐的同一動作.直到今天,不同地區,不同文化族群舉行聖餐時,都是同樣的動作.正如保羅所說:「我當日傳給你們的,原是從主領受的,就是主耶穌被賣的那一夜,拿起餅來,祝謝了,就擘開,說:這是我的身體,為你們捨的,你們應當如此行,為的是紀念我……」(林前11: 23-26)

復活的基督在以馬忤斯路上,教導門徒聖經和與門徒分享餅和酒,就象徵復活的基督藉聖道和聖禮活在教會中.[ii]卡拉瓦喬以以馬忤斯的晚餐畫作,來救贖自己的罪行,因為基督的死,就是為我們贖罪,領受他的聖體寶血,就是接受他的救贖.

我想卡拉瓦喬不是懼怕被捕,甚至承受刑責,乃想到過去所作所為給人帶來的傷害,這個不是牢獄可以償還給受害人或他們的家人的；那位棄船先逃的南韓船長,無論他如何公開道歉,或將來被判終身監禁,都不能除去內心的罪疚.

卡拉瓦喬六歲時歐洲遭遇鼠疫(bubonic  plague) ,家人包括父親因疫病離世,心靈因而重創,失去親人的少年卡瓦喬自始變得沒有安全感,性情暴躁,流連市集結,與同輩結成繪畫的童黨,身上常配利器,惹事生非,他們的口號:「無懼亦無望」( nec  spe, nec metu, without hope, without fear) .很自然,這樣的生活很容易闖出禍來.結果他在羅馬為了女人殺了她的情人,逃到馬爾他(Malta)島又槍傷人,被單獨關在高度設防監牢中.但他越獄；他下二百呎峭壁,泅泳三英哩,直到找到一條船,逃到西西里(Sicily).死時不夠四十歲,是因為中了油彩中的鉛毒.一位傳記作者說:「差不多可以說,他是不能避免犯險悖逆.當他一被權威人士接納,被教宗歡迎,被馬爾他候爵歡迎,旋踵間他就要做出岔子,弄壞事情,可以說是悲劇性的宿命.」

「悲劇性的宿命」是傳記作者的講法,但對尋求救贖的畫家,就沒有諉過成長的坎坷,他需要內心被上主赦免的平安,這位光影大師能夠逃獄,身體得自由,但卻不能擺脫良心被控訴的綑鎖,他要藉繪畫以馬忤斯的晚餐,走出罪的陰影.而復活節的信息就是赦罪的平安.聖餐時,我們我引述耶穌的話,說:這是我的身體為你們捨的,這是我的寶血為你們而流的.稍後,基督復活對門徒顯現第一句話就是:「願你們平安!」(路24: 36) 這是罪得赦免的平安,不是一般問候語,因為在約翰福音20: 21, 22,復活的基督對門徒說:「願你們平安!父怎麼差遣了我,我也照樣差遣你們.說了這話,就向他們吹一口氣,說:你們受聖靈.你們赦免誰的罪,誰的罪就赦免了；你們留下誰的罪,誰的罪就留下了.」

復活後耶穌說「願你們平安」的話,不單是猶太人見面時的問候語,以前的「願你們平安」是期待上帝神國要帶來的平安,復活後是期待的成就.就如我們一路為抵押給銀行的樓房供款,直到有一天供完了(maturity),雖然住在同一房子,但感覺不同,這是真正屬於自己的房子了.復活後耶穌說:「願你們平安!」就是上帝已成就了平安!

今天,復活的主對你說:願你們平安,我們不再受罪的綑縛,走出罪的控訴,come out of the shadows of sin走出罪的光影!

各各他小崗十架的孤寂

卡拉瓦喬善用光影,繪畫人物,他畫的黑暗很有層次,論者都同意這方面是他是受到教會禮儀的啟發.每年聖週(Holy Week) 教會舉行聖周受難經禮儀 (Tenebrae [darkness] Offices)三天,第一天在星期三晚上舉行,其餘兩次在受難日早上和星期六早上,崇拜的特色是聚會都在黑暗中舉行,教堂安放受難經燭台,放置十五枝蠟燭,每唱一段詩篇,吹熄一枝,剩下中間一枝,崇拜結束時也熄滅了,象徵基督進入受難的黑暗,等候復活帶來的光明.

跟隨十架的祈禱

其實,不單每年的聖週以黑暗來表達基督為我們受難,教會在每一天的日課(Daily Offices, Daily Prayers) 紀念基督的受難.按著耶穌在各各他小崗十架受難的時刻,信徒要有三次禱告時間:羅馬人從晨曦早上六時開始計算的第三,六,九個時辰(Terce,Sext, None) ,耶穌是被掛的午前禱(早上九時) ,天地開始昏暗的午禱(午正十二時) 和斷氣的申初午後禱(下午三時) ,加上起床的晨禱(Laud) ,晚餐時的晚禱(Vesper) ,睡前的寢禱(Compline) 已經是六次祈禱了,但還有一次祈禱,差不多在黑暗中,或微弱亮光進行的早禱.早禱是一天最早的祈禱,在午夜到晨曦任何時間進行的祈禱.有些人在半夜進行早禱,因為耶穌在童女的比喻中說,等候新郎的童女的燈要預備好油,不要打盹睡著了,那怕新郎半夜來了,所以你們要儆醒.(太25: 6, 13) 當然我們很難半夜起來禱告,很早時早禱(Matin)和晨禱(Laud)就合併,起來第一個禱告就是早禱.

黑暗中的早禱

但在修道院,在午夜後晨曦前,仍舉行早禱.我有機會參加大嶼山神樂院的早禱,印象深刻,對黑暗有嶄新的體會.大嶼山神樂院是屬熙篤會(Trappist) ,修士的職志就是祈禱,為世界祈禱,安靜寡言,他們到現在仍嚴守一天七次祈禱的日課(Daily Offices).早禱不單在深夜進行,而且祇有極微弱的燭光,崇拜最後,擊打聲一響,緊有的燭光立時熄滅,修士停止頌唱,頓時光影消逝,音響滅絕,靜默幾分鐘後,修士們隨隨引退.獨坐教堂,感到黑暗撲面而來,宇宙此刻為你停止,此震撼難忘的經歷,到此刻印象猶新,我經歷了絕響和絕光絕影.

黑暗開出光明

早禱在黑暗中進行,早禱是一天的開始,表徵從黑暗中開出希望,在黑暗中看到光明.聖經描述耶穌在十架的受難:「從午正到申初,遍地都黑暗了.約在申初,耶穌大聲喊著說:以利!以利!拉馬撒巴各大尼?就是說:我的神!我的神!為甚麼離棄我?」(太27: 45, 46) 遍地都黑暗,表示耶穌承受極大的苦難,天地變色,但這眼見的受難,被釘受辱,卻帶來實在的得勝,the victim is the victor那被殺的羔羊同時是掌管萬有的主.

受難者也是得勝者

我們每主日背誦信經:基督在本丟彼拉多(Pilate) 手下遇難,被釘十架,埋葬,降到陰間…….神學家對降到陰間有兩種解釋:一是極其降卑；二是降到陰間對舊約聖徒凱旋的宣告,罪被打敗,他們等候的救贖已經來到,救主已一役竟其功.所以馬太跟著說:「忽然,殿裡的幔子從上到下裂為兩半,地也震動,磐石也崩裂,墳墓也開了,已睡聖徒的身體,多有起來的.到耶穌復活以後,他們從墳墓裡出來,進了聖城,向許多人顯現.」(太27 : 51-53)

人生光影的逆轉

十架帶來自然光影的變化,「從午正到申初,遍地都黑暗了.」；也帶來對人生光影不同的詮釋.前明報主編劉進圖弟兄街頭被斬,對他和家人來說是生命極黑暗的時刻,但此事卻把香港撕裂的社會難得的團結起來,在立法局不同政見的議員都同聲譴責,上萬人自發上街守護香港新聞自由,在校園學生貼出標語:「我們都是劉進圖」；”They can’t kill us all” (他們不能殺清我們) .跟著劉進圖弟兄在明報發表系列感想,其中一篇這樣說:

「我從少年時期就返教會,學會人生目標應是榮神益人,努力事奉.過去我一直認為,事奉是奉獻自己的才智,能力和時間,做有益於社會的事情,自然就能榮神益人.例如,我努力讀書,考取好成績,便是向世人作好見證.又例如,我捨棄較賺錢的法律專業,為興趣和使命跑去當記者,還盡量抽時間教學及服公職,這就是努力事奉,榮神益人.

如果我承認今次我遇襲產生了積極的社會意義,那就意味我是通過一個非自願的,被動的,與個人才智能力全不相干的途徑,單純靠承受傷害和痛楚來造就別人,這完全違反了我一貫的概念.其實,我這個人頗有點精英主義,上司曾批評我恃才傲物,我從來沒有想過,放下聰明才智,不再勞心勞力,人還能做什麼?我從來不曾明白,所謂無權勢者的權力,受苦的事奉,到底是什麼意思?我大半生走過的所謂榮神益人努力事奉的路,到底當中有多少真實而深刻的意義?有多少其實是為了滿足自我的虛榮?這一連串的問題,在靜夜無眠的晚上,折磨著我的心.我沒有答案.但我學會一件事,我需要謙卑下來,重新思考人生的方向和意義.」

作為少年劉進圖的牧者,我頗了解他所說的榮神益人,就是讓光照在人前.作基督徒若有好的行為,令人欣羨欣的成就,就是好的見證.但自己生命陷在黑暗中,可以榮神益人嗎?這次黑暗的經歷,把他過去以為照在人前的光完全掩蓋了,他成為軟弱無助,躺在深切病房的受害者.

但在十架事件中,人看來是失敗和羞辱,卻是得勝,他明白基督藉自己身上承受的看得見的黑暗,打敗了那看不見的黑暗,甚至有時偽裝光明正義的黑暗.中文黑色叫玄,人生的黑暗藏有上帝的玄妙和玄機,上帝在人認為失敗的基督,帶出凱旋.上帝也會在我們認為黑暗的日子,扭轉我們的生命,就如在以馬忤斯路上的兩位失意的門徒.

仄徑的伴隨

以馬忤斯的故事說明,人生很多時候就像這兩位門徒,突然失去素來所盼望的.「但我們素來所盼望,要贖以色列民的就是他!」以馬忤斯在耶路撒冷西北偏西十一公哩,行到村口,日頭已經平西了,他們應該是下午起行,避過正午的烈日,開始時彼此談論耶穌的墳墓空了,屍首不見了,滿懷憂愁,太陽在前面,愈行愈覺太陽下沉,四周愈來愈昏暗,就像他們的心情.

隱藏在同路人背後的基督

但此時耶穌以陌生的同路人的身份,「親自就近他們,和他們同行」,耶穌說:「你們走路彼此談論的是麼事呢?」我們看到,是耶穌主動就近他,陪同他們,並聆聽他們,忍耐教導他們,引導他們慢慢進入真理.今天你是否感到素來所盼望的失去了,覺得上帝離你很遠?其實上帝沒有離開,祇是他隱藏在同路人的背後,一直陪伴著我們,祇是我們眼睛迷糊看不出來.「眼睛迷糊」是上帝的工作, They were kept from recognizing him (NIV) ,因為上帝隱秘行事不為人所知.

十架令光影轉移,陽光普照的正午突然遍地昏暗,聖父掩面不看聖子,因為此刻聖子背負了世人的罪孽.在人生的路途上,我們也會經驗突然風雲色變,此時候我們不要忘記,上差遣眾多守護天使,貼近我們,陪伴我們,用聖經開導我們,走出陰霾.上主不單以聖道和聖靈伴隨我們行走人生路,更以聖徒陪伴我們,互相鼓勵.

行道才能經驗真實的信仰

最後,我們看看兩位門徒甚麼時候才認識耶穌呢?不在聆聽聖經時,仍實踐聖經接待客旅時.單單學習聖經而不遵守聖經,我們永遠不會經驗信仰的真實.不錯,接待客旅是律法的要求,日已平西,前面愈來愈黑了,不過耶穌好像還要往前行,夤夜趕路,若是這樣,他們可以順水推舟,不須強留耶穌.

但他剛剛聽完耶穌解釋聖經,心裡火熱:「他們彼此說:在路上,他和我們講解聖經的時候,我們的心豈不是心裡火熱的嗎?」(路24: 32) 他們沒有停在聽道的層面,他們要行道,為這位陌生人設身處地著想,前面那麼黑,怎可以餓著肚子趕路呢?聖經說「他們卻強留他」,這表示他們真的關心陌生客旅.當他們接待陌生人,分享食物和居所時,普通的晚餐,就被祝聖為聖餐,他們就認出這人就是耶穌.善行使信仰鮮活(faith made active in charity).

結論:以馬忤斯的故事開始時是滿臉愁容,舉步維艱的仄徑,難行的窄路；故事結束時是歡樂的筵席,興奮的心情,因為聖道的啟蒙,同路人的伴隨.這故事提醒我們,要在人生的路上,代表基督陪伴失意的人,耐心聆聽他們的苦痛,贏取他們的信任,用聖經開導他們；也要對有需要的人,給予援手.

這是以馬忤斯晚餐的光影,也是復活節的色彩.

[i]一位傳記作者說:「差不多可以說,他是不能避免犯險悖逆.當他一被權威人士接納,被教宗歡迎,被馬爾他候爵歡迎,旋踵間他就要做出岔子,弄壞事情,可以說是悲劇性的宿命.」(“It almost like he cannot avoid transgressing.  As  soon as he’s welcomed by authority, welcomed by the Pope, welcomed by the Knight of Malta, he has to do something to screw it up. It’s almost like a fatal flaw”, Andrew Graham -Dixon, Sacred and Profane, 2011)

[ii]馬丁路德改教時期,引起教會本質的反思.甚麼才是真正的教會呢?改教家都有共識:那裡正確宣講聖道,那裡正確施行聖禮,那裡就是真正的教會,至於教制組織和禮儀,各地可以不盡相同.(奧斯堡信條七論教會)



\newpage

\section{當世界變得越來越得人驚的時候}
\label{sec:242}
\textbf{鄧瑞強博士 - 當世界變得越來越得人驚的時候}
\newline
\newline
\hyperref[sec:241]{< < < PREV SERMON < < <}
~
\hyperref[sec:toc]{[返目錄]}
~
\hyperref[sec:243]{> > > NEXT SERMON > > >}


約翰福音 20:19-23
\newline
\begin{longtable}{cl}
\hline
\hline
章節 & 經文 (和合本修訂版)\\
\hline
20:19 & \begin{tabularx}{0.7\textwidth}{X} 那日（就是七日的第一日）晚上，門徒因怕猶太人，所在的地方門都關了。耶穌來，站在當中，對他們說：「願你們平安！」 \end{tabularx} \\ \\ \relax
20:20 & \begin{tabularx}{0.7\textwidth}{X} 說了這話，他把手和肋旁給他們看。門徒一看見主就喜樂了。 \end{tabularx} \\ \\ \relax
20:21 & \begin{tabularx}{0.7\textwidth}{X} 於是耶穌又對他們說：「願你們平安！父怎樣差遣了我，我也照樣差遣你們。」 \end{tabularx} \\ \\ \relax
20:22 & \begin{tabularx}{0.7\textwidth}{X} 說了這話，他向他們吹一口氣，說：「領受聖靈吧！ \end{tabularx} \\ \\ \relax
20:23 & \begin{tabularx}{0.7\textwidth}{X} 你們赦免誰的罪，誰的罪就得赦免；你們不赦免誰的罪，誰的罪就不得赦免。」 \end{tabularx} \\ \\
[1ex]
\hline
\hline
\end{longtable}
講題:當世界變得越來越得人驚的時候　 When the World Becomes Increasingly Scary

經文:約翰福音20:19-23

講員:鄧瑞強博士

場合:香港中文大學崇基學院禮拜堂

日期:2014年4月27日

各位弟兄姊妹,早安.

覺不覺得香港變得越來越得人驚呢?

我不知道是否世界上每個地方都能講＜國王的新衣＞這個故事.若果能夠講這故事的話,我想,每個聽眾都會讚賞那個說真話的小孩子.但是,如果有一個地方,那個講真話的小孩子會被人斬斷腳的話,那麼,這是否很得人驚呢?

我不知道是否世界上每個地方都能讀《動物農莊》這本小書.若果能夠讀這本小書的話,我想,每位讀者都會對「所有動物一律平等,但有些動物比另一些動物更為平等」這想法反感.但是,若是有一個地方,那些對這個想法口口聲聲反感的人(或動物),卻在事實上實踐這個想法,並且視追求真正平等的人為違反這「動物農莊」的基本的法律的話,那麼,這是否很得人驚呢?

人類歷史已經很長了,痛苦的悲劇已經夠多了.從歷史的教訓裡,我們已經領悟到一些道理.有些道理,現在已近乎常識,如:權力使人腐化,必須加以監管,政權的興亡視乎是否尊重人民的意願,強權不敵真理….這些都是人人掛在口邊的常識,但奇怪的是,竟有人違反常識.尊重一下常識,難道真是很難嗎?很多人誠心自問,都知道那是一條錯路,卻為了私利走下去,這是否很得人驚呢?

耶穌的故事,發生在一個十分得人驚的世界裡.

正正這世界越來越得人驚,我們更需要耶穌的故事.

今日的講道經文是:約翰福音20: 19-23

20:19　那日(就是七日的第一日)晚上,門徒所在的地方,因怕猶太人,門都關了.耶穌來,站在當中,對他們說:「願你們平安!」

20:20　說了這話,就把手和肋旁指給他們看.門徒看見主,就喜樂了.

20:21　耶穌又對他們說:「願你們平安!父怎樣差遣了我,我也照樣差遣你們.」

20:22　說了這話,就向他們吹一口氣,說:「你們受聖靈!

20:23　你們赦免誰的罪,誰的罪就赦免了；你們留下誰的罪,誰的罪就留下了.」

故事的開始是講某個晚上,門徒因害怕猶太人,門都關了.

在約翰福音的神學裡,「猶太人」扮演著一個「神學角色」,他們是反對神,反對真理,反對光明的人.由於他們的存在,外面的世界變得黑暗.由於他們的存在,手無寸鐵的人被迫隱匿在黑暗的角落裡.

在這個故事裡,不久前發生了一件很得人驚的事.一個叫耶穌的人,他叫貧窮的人有麵包可吃,他叫被社會遺棄的痲瘋病人獲得愛護,當然,他也叫暴力抗爭的人放下刀劍.他宣講天國,這是久違的信息,也是人人渴望的信息.他的行動,他的信息,讓人體驗神真實的慈愛.這種慈愛,多年來被宗教的沉重的規條埋葬了,已沒有多少人體驗過.他的行動,他的信息,勾起人們對救贖的渴望.當時的人,被羅馬軍隊欺凌,被沉重的稅項剝削,生命像將殘的燈火.耶穌來了,他顧惜這些將殘的燈火,他在人心中,燃起對正義與和平的渴望.

就當人的心越來越渴望真理,神的愛,正義的天國時,耶穌被人殘酷地殺死了.他顧惜世上將殘的燈火,但邪惡卻撲滅了他這正義的火.那位下死刑判決的彼拉多,明知這判決是錯的,也違反良知,殺死這無辜的耶穌.

一個人純粹因為他尊重弱者,渴望正義,宣講真理而被殺,這真是得人驚呀!

他的門徒很害怕.他們沒有帶著寫上「They can’t kill us all」的橫幅上街示威.他們真的害怕.

我們只是平凡的人,像野地裡的一朵小花,今日還在,明天便不見了.我們怎能以肉身擋著刀劍,以身軀擋著坦克,以螳臂之力托著巨大的十架?

門徒關上門,因為他們害怕.他們的恐懼驅使他們抽離現實.恐懼令他們不敢面對世界.他們不想面對外面的不義,他們只想安安分分,活在密室中.恐懼讓人活得似蟑螂,只能在幽暗的角落爬行.

約20:19　耶穌來,站在當中,對他們說:「願你們平安!」

門雖然關了,但不妨礙耶穌的同在.外面的恐怖叫人懦弱,但這並不阻礙神在人心裡創造勇氣.死亡的威脅令人卻步,但這並不阻礙基督向人展示其復活的力量.有些人,是天生的勇者,如荊軻,可以高歌:「風蕭蕭兮易水寒,壯士一去兮不復還!」但更多人只是軟弱的人.軟弱的人若能生起面對逆境的勇氣,則更顯出這是一大神蹟,這是一種恩典下的勇氣.這種勇氣,不是他們固有的,而是神在他們心裡,一步一步鼓勵他們,而他們一步一步培養出來的.這有點似天生雙腳無力的人,因著其他人的幫助及鼓勵,自己的鍛練及努力,終於能站起來一樣.這種站起來是非一般的站起來,足能贏得真正的掌聲.神要我們經驗生命的神蹟,賺取這種掌聲.

復活的耶穌,向恐懼的門徒顯現,對他們說:「願你們平安!」這句話的原意,基本上不是一個願望,而是陳述一個事實:「平安就在你們當中」.你們不需要等待平安,而是認清楚你們根本就在平安之中.

「平安」是聖經的一個基本用字,代表生命的豐盛與完滿,代表生命在光明與正義中.「平安」不是指無風無浪,無驚無險,而是指風浪與恐怖都奪不走也掩蓋不到的生命的豐盛內容.當代最經典的一幅表達「平安」的畫面,是一個手無寸鐵,孤立無援的人,站立在一排坦克之前.這是一個很得人驚的處境,但無損那人展示「平安」.何謂「平安」?就是在恐怖的巨大力量前,顯示一個卑微的人的生命裡的無限豐盛.這種「平安」,展示人性深不可測的深度.這個人,揭露一個坦克以外的世界,展示一種坦克以外的力量.他以卑微的生命,迎向巨大的「反生命」的力量,宣告人在極恐怖的處境裡,仍可展露恩典下的平安.

這不是那種事事如意,我們「被動地」享受的平安.而是那種在逆境面前,在得人驚的事情面前,需要我們「主動地」活出的平安,那種被神的恩典激活的平安.

復活的耶穌,進入關了門的內室,向恐懼的人說:「平安就在你們的生命中.」「生命的豐盛意義就在你們的靈魂深處.」

20:20　說了這話,就把手和肋旁指給他們看.門徒看見主,就喜樂了.

在耶穌的手和肋旁看到的,就是他受傷受苦的記號.

受傷受苦受死的人,如今成了一個復活的人.傷痕依舊,但無損其展示生命的榮耀和永恆.這種展示是很重要的,會幫助人面對一個得人驚的世界.

近來,向我們展示傷口的,是明報的前總編劉進圖先生.他發表了若干病床上的＜感言＞.他說:「真理在胸筆在手,無私無畏即自由.」他說:「我仍會密切留意,未來數年幾家大學傳理系的收生水平會否下跌,以及畢業生加入新聞行的比例會否降低,這是觀察香港新聞自由前景的重要指標.」他說:「有同學在留言中委婉相勸,想我趨吉避凶,意思是離開傳媒這個高風險行業,甚或離開香港這個板盪之城.我完全明白同學的善意,親友中也有好些人有這樣的看法.不過,我這一生走的是窄路,路上同行的人雖然較少,但風景獨好,一路過來心安意暢,且深信上帝與我同行,縱有死蔭幽谷,也必蒙恩受庇.思量再三,我在同學群組上留言回應,表明心跡:『窄路險途無怨悔,恩澤滿懷寸心知』.」

帶著傷痕,他思念著真理,思念著一個城巿的前景,思念著無人願意走的窄路.

一個身懷受苦記號的人告訴我們,受苦並不表示失去了平安,受苦並沒有取消生命的豐盛,受苦反而讓人看清楚什麼是生命中的渣滓,什麼是生命中的精華.他的話對恐懼的人,有莫大的鼓舞作用.

帶著受苦記號的人,叫其他人不要在艱難中失去勇氣,他的話最有說服力.

經文說,復活的耶穌,把手和肋旁的傷痕給門徒看.門徒看見這身懷受苦記號的主,就喜樂了.

20:21　耶穌又對他們說:「願你們平安!父怎樣差遣了我,我也照樣差遣你們.」

主耶穌再重申一次,「平安就在你們當中」.這是一個事實,無人能奪走你的平安,因為「平安」是從神而來的禮物.

面對越來越得人驚的世界,不是人人能體會這種「平安」的,也不是人人能從受苦者的記號中支取勇氣的.但總有些人能明白,無論夜多深沉,恍惚整個宇宙的黑暗都壓下來,仍無法掩蓋一枝火柴的光明.總有一些人明白,在叫人恐懼的力量背後,有一種給人勇氣的力量.在死亡的力量背後,有復活的力量.在看似無意義的世界背後,有一位上帝.

總有人看到真相的.這些看見的人,領受著天命,帶著一種使命活著.

主耶穌對這些人說:「父怎樣差遣了我,我也照樣差遣你們.」

天父差遣耶穌,耶穌差遣那些看見他復活的人.這是一種使命的傳承.

主耶穌是一個領受天命的人,他有使命在身.天父差遣他,他要實踐天父的旨意.

在約翰福音裡,有兩節經文扼要地交代天父差遣耶穌來世的用意.約翰福音3: 16-17:「神愛世人,甚至將他的獨生子賜給他們,叫一切信他的,不至滅亡,反得永生.因為神差他的兒子降世,不是要定世人的罪,乃是要叫世人因他得救.」

天父有「好生之德」,他愛我們,他希望我們找到生命的豐盛內容,即找到「永生」.要得到「永生」,便要透過耶穌.耶穌帶著使命來生活.他在敵擋真理的地方,活出真理.當時的行政長官彼拉多問他:「真理是什麼?」耶穌就用他的一生去回答這個問題.這就是:在黑暗瀰漫的日子,他堅持光明的生命.當人人都甘願被謊言蒙蔽時,他卻大聲說:「真理必叫你們得以自由.」(約8:32)當人人都因恐懼而關上門時,他卻以有釘痕的手鼓勵人勇敢面對.我們是要透過耶穌這種生命,去體會「永生」的內容.

他差遣那些明白他的人,接續他的人生.在這個越來越得人驚的世界,要堅持真理,要有勇氣,要活出「永生」.

20:22　說了這話,就向他們吹一口氣,說:「你們受聖靈!

20:23　你們赦免誰的罪,誰的罪就赦免了；你們留下誰的罪,誰的罪就留下了.」

要完成天父的使命,需要從天而來的力量.主耶穌向接續他的生命的人說:「你們受聖靈」.

一講到力量,我們想起金錢,警棍,核彈.但聖靈這種力量,卻像春天的風,輕輕地吹,吹醒大地,呼喚青草,綻放花朵.聖靈是一種「無力之力」,就像劉進圖先生受傷後,「無力地」躺在病床上時,卻呼喚起成千上萬的人起來,為言說真理的自由而努力.這是一種「無力之力」,一種「十架的大能」.聖靈會給人力量的,但這不是神蹟的力量,而是受苦的力量.正如聖靈給耶穌力量去背起人類罪惡的十架,藉此顯出在得人驚的世界裡,平安是可能的,永生是可能的,豐盛的生命是可能的.

在聖靈裡的人,成了人類的鏡子,照出世人的真面目.「你們赦免誰的罪,誰的罪就赦免了；你們留下誰的罪,誰的罪就留下了.」這裡講的「罪」,不是指道德的罪,而是指人拒絕耶穌.這節經文是說,一個有聖靈的人,一個能承擔十架的人,一個接續耶穌生命的人,他能判定誰是跟從耶穌的,誰是反對耶穌的.或者說,一個講真話的人,會顯出誰是講假話.一個有真勇氣的人,會顯出誰是因懦弱而不敢放下刀劍.一個有永生的人,會顯出誰的生命只是過眼雲煙.

看來,這個世界會變得越來越得人驚的.但在心靈的最深處,復活的主會向我們顯現的.他會向我們展示他的傷口,他會告訴我們這些傷口不會妨礙生命的平安與豐盛.若你看到這位復活主的話,他會差遣你的生命的.你的生命就會帶著天父的使命.聖靈會給你力量的,但這是一種「無力之力」.然後,你會活出像基督一樣的生命.這種生命像鏡一樣,照出人類的善惡.

你看見復活的主嗎?你願意領受他的差遣嗎?

但願:榮耀歸於聖父,聖子,聖靈.阿們.



\newpage

\section{晨光第一線復活主再現}
\label{sec:243}
\textbf{許開明牧師 - 晨光第一線,復活主再現}
\newline
\newline
\hyperref[sec:242]{< < < PREV SERMON < < <}
~
\hyperref[sec:toc]{[返目錄]}
~
\hyperref[sec:244]{> > > NEXT SERMON > > >}


講題:晨光第一線,復活主再現　The Morning was Broken with The Risen Lord

講員:許開明牧師

經文:約翰福音21

場合:香港中文大學崇基學院禮拜堂

日期:2014年5月4日

引言:

今日是復活後第三主日,我們要思想在約21章中「主復活的大能」.主復活,乃顯示耶穌是神,顯示耶穌的大能.然而,復活的大能在那方面呢?從神學的角度來看,應有兩方面:一方面是從死亡的權勢下復活再生,死,原本是生命的終點,存在的終點；死,使生命不存在,使人失去生命力；但,復活的基督,使死亡,不再是生命的終點,不再是存在的終點；這顯出一位大能的基督.第二方面,復活勝過死亡的權勢,死亡的權勢是使人無生命力,引致毀滅,消失；但,復活的基督,使死亡的權勢失去效能,失去生命力,步向毀滅,消失；這顯出一位得勝的基督.第三方面,復活主,身體改變了,能突然出現,忽然消失,穿牆過壁,能不被限制,不被約束,不被阻攔,不被擊倒,超越時間與空間；還有,復活的基督耶穌,是有身體的,又好像無身體,這是位榮耀的基督.然而,這樣的基督對我們有何意義?當然有喇!大能的基督,使我們相信祂,仰賴祂.榮耀的基督,使我們敬拜祂,事奉祂.得勝的基督,使我們也得勝罪惡,得勝死亡,不斷更新改變!

然而,如果基督的復活的意義只停留在此,這個基督有兩個問題,值得深思.第一個問題:這復活的基督,會否成為「利我」的基督?對我有好處,使我得勝罪惡,得勝死亡,使我得力勝過困境,使我活得更好,使我得到庇佑,使我得到祝福.信主復活得復活,得永生,得好處,得益處等,無可厚非.然而,這樣的話,這跟一般人求神拜佛,祈求趨吉避凶,其實相似,只是換了個信仰的對像:更高,更有份量,更勁的復活主.這是不壞,但未算好,也與其他宗教信仰有沒太大的分別?

第二個問題:這復活的基督會否變成「高不可樊的基督」?大能的基督,勝利的基督,榮耀的基督.非常偉大,非常值得羨慕；但是,與我們這軟弱,卑微,有限的人類真有一大距離!這也是不壞的,但也未算夠好的基督.

然而,約翰福音筆下的復活的基督,與一般教會單單強調的榮耀基督,有所不同.耶穌是位道成肉身的復活主,也是位不同凡響的基督,是位落地的(Down to Earth)的基督!這位落地的,既是群體所信奉?也是信眾個人的復活主,能使人得力面對人生,這是約翰福音給我們信息..

讓我們進入約翰福音21章中,思想這位落地的復活主.

「這些事以後,耶穌在提比哩亞海邊,又向門徒顯現；祂怎樣顯現記在下面.有西門彼得,和稱為低土馬的多馬,並加利利的迦拿人拿但業,還有西庇太的兩個兒子,又有兩個門徒,都在一處.」【約廿一1-2】這裡所舉一共有七個門徒；『七』在聖經裡象徵今世的完全,故『七個門徒』代表全體主的門徒,甚至也包括我們後來信主的人.1約翰福音有七個神蹟,2耶穌有七個自承,3七樣東西去形容耶穌:道,生命,人的光,真光,恩典,真理,獨生子.

「西門彼得對他們說:『我打魚去.』他們說:『我們也和你同去.』他們就出去,上了船,那一夜並沒有打著甚麼.」【約廿一3】彼得說:「我打魚去,」『打魚』就是重操故業(參太四18~22).門徒彼得可能這樣想,耶穌復活了!耶穌的事可以自己辦!若果耶穌周圍顯現,一定有許多人相信祂,現在用不著我,也用不著門徒了.到底跟著耶穌還有何意義呢?事實上復活了的耶穌都不易跟嘛,祂可以突然出現又突然會消失,真了不起呀!但跟我們這班打漁佬真是有點距離了!現在的復活主的神性好勁,但人性好弱,好係唔係人咁品了!點跟他,事實上,耶穌復活後,還沒有叫我們跟他呢!

門徒中彼得大哥為人最爽,先提出不如大家重操故業,「搵食,搵食」「唔搵怎會有得食?」不如大家重操故業,做番老本行,有專業加上經驗,搵得到,有得食!昔日跟隨耶穌的理想大志,看來都是無法實現的了!帶著無奈與挫敗,彼得先說:「我打魚去!」其他六位門徒也應:「我們也和你同去」於是他們帶著挫敗的心,打魚去了!

面對人生的挫敗,人便會想到避開他,不想繼續面對挫敗,免得心靈再受傷害,這是人心理的自衛術.例如:認識的朋友中,有因為教書受不起學生的氣,或與同事不和,或與校長的期望有出入,或給教育局朝令夕改的政策弄到筋疾力盡而帶著挫敗感離職；從此不再教書,因不想再面對挫敗.

然而,當人感到事與願違時,挫敗感更大.「那一夜並沒有打著甚麼,」古時漁夫喜歡在晚上打魚(參路五5),據說夜晚正是打魚最好的時候.彼得是個老練的漁夫,他應該知道在何時,何處可以打到魚；但他們竟然整夜打不著魚!這真是令人奇怪!六個人跟著他,一無所獲,大家都要餓死麼?自己失敗,是很難受的事.到底發生什麼事?七個懂捕魚的能手,天時地利人和是成功的要素,整晚上竟會「一條魚都沒有!」真是大失預算.在我們人生的經歷中,也常有大失預算,大失所望的時候.

就在黑暗中和門徒的挫敗中,復活主就在此時出現了.「天將亮的時候,耶穌站在岸上；門徒卻不知道是耶穌.」【約廿一4】晨光第一線,復活主出現,天色也愈來愈光亮了!約翰生動地引出復活主的顯現,將有重要事情發生.然而,在海岸仍然大霧,所以他們七個漁夫都看不清楚岸上的人是誰?復活主若隱若現,帶點神秘感,懸疑的味道,約翰寫得引人入勝.事實上,耶穌在我們的生命中,也會如此.你不一定看得祂很清楚,有時你甚至懷疑復活主的存在及主的帶領!

猶太人對日子的觀念跟我們不同,乃是由黑暗開始,逐漸迎向光明.人生也是由痛苦開始,逐漸迎向喜樂!事實上,一個人出世,都是由黑暗痛苦開始,不獨是媽媽受生產之痛,連胎兒也是.快出生的胎兒原本生活在黑暗平靜的子宮中,出生之際,熱烈的燈光便會使胎兒受驚並感到不舒服.猶太人從畫夜中得著的人生智慧,苦盡甘來,夜盡天明；掌管畫夜的上主,使人生總有希望,黑暗痛苦總會過去,就是這種思想使猶太人成為非常堅強的民族!從約翰福音的記載,復活的主是位臨在和內蘊()的基督有下列特點:

一 復活主關心他們「有吃的沒有」

「耶穌就對他們說:『小子(朋友,老友),你們有喫的沒有?』他們回答說:『沒有.』」【約廿一5】.門徒整晚一無所獲,七個人身體都非常的疲乏,心靈感非常失望,正帶著無奈的心情準備回岸上休息.怎知岸上有位疑似是漁夫的人,問候他們說:「老友,你們有吃的沒有?」或「兄弟,搵到食嗎?」當你帶著挫敗的心情時,有人懂得適時給你關懷和,這是多麼的窩心!復活主是好牧人,祂關心我們人生最基本的生理需要說:

「小子,你們有吃的沒有?」祂主動來到我們中間,祂關心我們的景況,祂關心我們最關切的需求.

二 復活主指導他們的工作:

當門徒回應岸上人說:「沒有」,耶穌說::『你們把網撒在船的右邊,就必得著.』他們便撒下網去,竟拉不上來了,因為魚甚多.」【約廿一   5-6】門徒坦誠告知那陌生人:「沒有!」他們承認自己的失敗!人面對挫敗,要肯承認並接受現實,才有改善的可能!人生的轉機往往在於人肯承認錯誤,並接受別人的建議和改正.復活主指出他們失敗的原因是弄錯了方向,並指導他們正確之道.

人生的失敗是由於擺錯位.彼得和門徒都是漁夫,且是具有豐富的漁夫,打左打右應該是不會錯的!岸上這位人兄是誰?為何要聽他的話?如果門徒不聽,不信,不願放下己見,他們便繼續網不到魚,繼續挫敗下去!

肯謙虛的承認錯誤,承認不足,接受意見,人才能脫離挫敗,得著成長.古時,漁夫會彼此提醒,在岸上的漁夫會提醒船中的漁夫魚網應撒在那裡?做人處事,當局者迷,我們需要別人的提醒.接受別人善意的提醒,使人得到改進.許多人永遠活在挫敗中,因為不能接受別人的提醒!

三復活主供應他們的需要

「他們上了岸,就看見那裡有炭火,上面有魚,又有餅.」【約廿一9】「就看見那裡有炭火,」注意,岸上的炭火,正與院子裡的炭火(參十八19)遙相對應；那夜彼得是在炭火旁失敗了,現在主是在炭火旁供應了他.「上面有魚,」這是說他們根本無須到海上去打魚,因為岸上已經有魚了.「上面有魚,又有餅,」『魚』代表海中的豐富；『餅』代表陸上的豐富.「耶穌對他們說:『把剛才打的魚,拿幾條來.』」【約廿一10】耶穌已預備了食物了,為何要加送?吃他們網到的魚,是大魚,是他們的戰利品,吃得特別滋味!

「西門彼得就去(或作上船),把網拉到岸上,那網滿了大魚,共一百五十三條；魚雖然這樣多,網卻沒有破.」【約廿一11】這是彼得奇妙的觀察,那網滿了大魚,竟沒有破.意即所網到的魚,一條也沒有漏掉,是一個神蹟.然而這神蹟,還有更深的神學意義.希伯來人相信上主掌管大海的主,大海是邪惡勢力的所在.古時的人以為怪獸來自大海,不信,請看但以理書和啟示錄.以色列的神只是陸地的神,民族的神,牧羊人的神,大海不是他所管轄的範疇,復活主是天地的主,陸地及大海之主,大小魚的主!復活主顯出祂的無比大能,祂知道那裡有魚；祂不讓魚群靠近門徒的船,進入他們的網中；祂可以使大量的魚在網中卻不會破,所以一條魚也不失掉.更了不起的,是可以控制魚的數量進入網中,153條魚有豐富含義.最簡單的解釋是由耶柔米(Jerome)提出來的.他說加利利海裡有一百五十三種不同的魚；這次捕到的魚包括了每一個種類.所以這個數目象徵:有一天萬國萬民都要在耶穌基督面前聚集的事實.另一個奧古斯丁(A)的發現,若由一加二加三加四……到十七,總數是153；當中十是十誡,七是恩典:上帝賜予世人雙寶.這神蹟,還有一非常重要之處,門徒只認識那位使魚網豐盛的主,卻可能看不到那管制魚群之主,祂能使魚群不進入門徒的網.大海要聽他的話,大魚也要受他指揮!還有,復活主能使人由無到有,又由到豐盛和有餘!

「耶穌說:『你們來喫早飯.』門徒中沒有一個敢問祂:『你是誰?』因為知道是主.耶穌就來拿餅和魚給他們.」【約廿一12-13】耶穌服侍門徒,因為知道他們又饑又餓又冷!但復活的主的偉大,不在祂的外在能力上,更在於祂內在生命的力上.復活主顯出祂好牧人的樣式,祂口中有問候,並有落實的行動,預備早餐團契,在團契中有人與人之間的愛與關懷.

四  復活主給他們使命與挑戰

人要脫離挫敗,人生的低谷,需要「生命的目標與理想」.一位雙失的青年,失戀與失學,寫下遺書後跳海.跳下大海後,努力掙扎,大叫救命,結果給救起來!水警問那位跳海的兄弟:「你既然寫下遺書決定死了,為何還要大叫和掙扎?」這位青年說:「當我跳下海後,在大風大浪中,忽然間領悟人生是要搏鬥的,跟風浪搏鬥,時間搏鬥,死亡搏鬥,環境搏鬥,命運搏鬥,疾病搏鬥!」這位跳海不死的兄弟,從此人生積極起來了.

復活主給彼得生命的目標方向與挑戰.「他們喫完了早飯,耶穌對西門彼得說:『你愛我比這些更深麼?』

彼得說:『主阿,是的；你知道我愛你.』耶穌對他說:『你餧養我的小羊.』「耶穌第二次又對他說:『約翰的兒子西門,你愛我麼?』彼得說:『主阿,是的；你知道我愛你.』耶穌說:『你牧養我的羊.』」第三次對他說:『約翰的兒子西門,你愛我麼?』彼得因為耶穌第三次對他說,你愛我麼,就憂愁,對耶穌說:『主阿,你是無所不知的,你知道我愛你.』耶穌說:『你餧養我的羊.』【約廿一15-17】

在本段經文復活主有三大挑戰:首先從失敗中站起,愛復活主而餵養群羊,愛復活主而受苦受死的挑戰:

1從失敗中站起的挑戰

彼得曾三次不認主,於是耶穌也問彼得三次.耶穌要彼得從失敗的地方站起來!一個人要面對自己的失敗,再站起並不容易.但事實上許多成功的人都經過失敗,有些人一生被失敗所勝,有些人卻從失敗中站起來.其實生命中的困難挫折,從積極一面去看是有它的意義.我國讀書人有言:「刀要石磨,人要事磨；刀不磨不快,人不磨不精.」  有位神僕人一句話常留在我心中:「脫離困難的最好方法,就是要學習在困難中站穩.」然而,令人感慨的是,許多人的一生都是無法從使他失敗的弱點,盲點中站起來,以致一生都無法翻身.許多人一生活在挫敗中,一直是沿地踏步,逃不出昨日的壞行為,壞習慣,壞的生活方式；被過去的罪咎纏繞,過去的,昨日的老我一直控制著他,直到如今,還會到永遠.危險在:這種人不會進化,反而會退化,僵化,硬化.

我在堂會工作,請工友的時候,有次一位女士見工,一進我辦公室,我聞到煙味,我坦誠問她,你吸煙嗎?她想了一陣答:是」我說:「我們教會,除教堂不准吸煙外,與兩間相連的學校相連也不可吸煙?」她說:「我是不會戒煙的,我真的無法勝過煙的吸引!我很享受食煙的樂趣」我應:「我欣賞你的坦誠,但無人會享受聞煙的樂趣!真不好意思,我們不會考慮你!」她很不高興的離去,我看著這位女士的背影,有點感觸!不是因為找工會有困難,而是這煙可以傷她身體,影響與家人的健康與關係.

2 愛復活主而餵養群羊的挑戰

耶穌要彼得愛主並愛的主的羊.耶穌那句話:你愛我嗎?是很力!我試過!大家不妨試試.但請留意,耶穌問彼得說:「你愛我嗎?」不是為自已的好處,而是為教會羊群的好處,並且耶穌以身作則,把一個好牧人的樣式表現出來,諸如關心問候,引導,供應所需等,讓彼得得學習.

3愛復活主而受苦殉道(martyrdom)的挑戰

復活主對彼得說:「我實實在在的告訴你,你年少的時候,自己束上帶子,隨意往來；但年老的時候,你要伸出手來(被捉拿),別人要把你束上(無自由),帶你到不願意去(無保障)的地方.」」「耶穌說這話,是指著彼得要怎樣死榮耀神.說了這話,就對他說:『你跟從我罷.』」【約廿一18-19】

在火爐旁,彼得彷佛是那條「在火中烤的魚」的魚,在火中燒著,發出魚香,給主給吃,為主所用!是會受苦,受傷,甚至受死!復活主挑戰彼得從失敗中站起來,面對苦難的十架,死亡!痛苦的死,恐怖的死!根據古代教會的傳說,使徒彼得不但是為主被釘十字架而殉道,並且還要求行刑的人將他倒釘,因為他自覺不配與主同有一樣的死法. 請注意,耶穌要彼得跟他走條死路,彼得竟無拒絕,勇敢為主殉道.

回顧過去的2013年的香港,對香港市民來說會有份無奈感.我頗認同一位現代人的描述:「在高大堅硬的牆和雞蛋之間,我永遠站在雞蛋那邊.」這句話是出自作家村上春樹(Murakami Haruki)在以色列接受耶路撒冷文學獎時的得獎感言,引發人思考微不足道的個人,對抗既有權力和體制的深層意義.

高牆與雞蛋可以說是今日香港的寫照.過去2013年的香港人彷彿是一隻雞蛋,在很多議題上面對猶如高牆的當權者,都像在以卵擊石.踏入2014年,香港人彷彿仍是雞蛋,只是擲到高牆的蛋愈來愈多；而我們面對的高牆卻愈來愈高和硬.無論如何,我欣賞那些真為香港爭取自由,民主,普選的人；為理想而不斷努力的人.這些人好像彼得一樣,明知前面是艱難之路,十字架的路,仍然勇往直前.

結語

晨光第一線,復活主再現；經過黑夜有平安,旭日初升,又是新的一天.各位,我們每過人都會經經歷過做錯事,以致有挫敗感和失落感.許多人從此一敗不起,抬不起頭來!南韓歲月號的船長遇難時作了錯誤的事,叫學生留在船艙,自己卻走人,以致使全船幾百人死亡；特別心痛的時是看到原本穿上救生衣的學生,可以早走的話,便有多人可以獲救,不致全部人罹難.面對這錯敗,船長如何站起來!

復活的基督,永遠帶來盼望.復活主會扶起挫敗的人,使到彼得和其餘的門徒從挫敗中站起來,接受使命,勇往直前.從這段經文,讓我們清楚認識到二大方面:一方面,復活主不是高高在上,而是臨在人間的主,且是好牧人.另方面復活主的大能不獨在從死人中復活,使絕望轉為希望,悲傷轉為喜樂,也在於衪牧養關顧之大能,當中包括:1祂關心問候挫敗的人,2祂指引挫敗的人走正路,3祂為挫敗的人預備食物,4祂給挫敗的人使命與挑戰戰；祂使挫敗的人站起來,充滿生命力,勇敢面對未來艱難的挑戰.

近代神學家John Hunter曾說:「每位信徒因自省:基督的復活對你有何意義﹖它對你日常生活有何影響﹖

基督復活是歷史的事跡,還是你信仰的內容.你的生命有無因為復活主而更新﹖你有否讓復活主進入你生活中.現在,許多信徒仍然沒有信靠復活主,好像基督沒有復活,因此絲毫沒有感到基督復活的大能,何等可惜.」

願你經歷復活主的大能,也能從挫敗中站起!阿們.



\newpage

\section{我必不致缺乏}
\label{sec:244}
\textbf{伍渭文牧師 - 我必不致缺乏}
\newline
\newline
\hyperref[sec:243]{< < < PREV SERMON < < <}
~
\hyperref[sec:toc]{[返目錄]}
~
\hyperref[sec:245]{> > > NEXT SERMON > > >}


詩篇 23:1-6
\newline
\begin{longtable}{cl}
\hline
\hline
章節 & 經文 (和合本修訂版)\\
\hline
23:1 & \begin{tabularx}{0.7\textwidth}{X} 耶和華是我的牧者，我必不致缺乏。 \end{tabularx} \\ \\ \relax
23:2 & \begin{tabularx}{0.7\textwidth}{X} 他使我躺臥在青草地上，領我在可安歇的水邊。 \end{tabularx} \\ \\ \relax
23:3 & \begin{tabularx}{0.7\textwidth}{X} 他使我的靈魂甦醒，為自己的名引導我走義路。 \end{tabularx} \\ \\ \relax
23:4 & \begin{tabularx}{0.7\textwidth}{X} 我雖然行過死蔭的幽谷，也不怕遭害，因為你與我同在；你的杖、你的竿，都安慰我。 \end{tabularx} \\ \\ \relax
23:5 & \begin{tabularx}{0.7\textwidth}{X} 在我敵人面前，你為我擺設筵席；你用油膏了我的頭，使我的福杯滿溢。 \end{tabularx} \\ \\ \relax
23:6 & \begin{tabularx}{0.7\textwidth}{X} 我一生一世必有恩惠慈愛隨著我；我且要住在耶和華的殿中，直到永遠。 \end{tabularx} \\ \\
[1ex]
\hline
\hline
\end{longtable}
講題:我必不至缺乏　I Should Not Be in Want

經文:詩篇23:1-6

講員:伍渭文牧師

場合:香港中文大學崇基學院禮拜堂

2014年5月11日

英文simple and naïve, and simple and profound 是兩個完全不同的觀念.Simple and naïve ,簡單天真,廣東話可繙作「簡單無腦」是揶揄人的話；simple簡單,有內容單薄,理據貧乏的意思.但simple and profound,簡明深入是讚賞,因為思想通透,使人容易明白；但意義深長,需要再三玩味,始能窮其究竟.詩篇23篇就是一篇簡明深邃的詩篇；容易了解,但意義深遠.

這詩篇主題清晰,文字優美,用田園的意象,表達上主和我們的親密關係:上主是我們的牧者,我們是他的羊,還有青草地,溪水旁,這是聖經中的貝多芬的第六交響曲-田園交響曲.我們彷彿聽到溪水流淌,雀鳥鳴叫,看到雜樹交蔭,綠草如茵,感到涼風拂臉.羊群散開,安詳躺在青草地,享受暖和的日照；有些在溪水旁喝水.這些意象都是當時的農業社會容易明白的.但以牧人和羊比喻上主和人的關係,就有深邃的神學意義了,耶穌在約翰福音就說過:「我是好牧人；好牧人為羊捨命.」(約10: 11)讓我們在田園交響曲的背景音樂中再讀一次詩篇二十三篇.

今天是母親節,上主愛我們,他賜我們父母養育我們,我們也有作父母的天職,牧養上主所賜的兒女,或上主交付我們的屬靈兒女,讓我們從這詩篇看看上主如何牧養我們,也學曉如何作好上主交付我們的牧養的工作.

先求上主甦醒自己的靈魂 v.3

「他使我的靈魂甦醒,為自己的名引導我走義路.」(v. 3)

青草地溪水旁,可指上主供應我們肉身的需要,一無所缺；亦可指在上主面前的安竭,就如主日到教堂崇拜,讓心靈安息下來,領受靈魂的糧食,竭息安靜可讓生命有重量的東西沉澱下來.作為父母,我們不單讓孩子們得到肉身的青草地和溪水旁,這個在今時富裕的香港,不難做到.我們的孩子最重要的,在追求人頭地時,不要失去自己的靈魂,常常保持靈魂清醒,走在正確的義路上,因為人像羊一樣,容易迷失.

我們是容易迷失的羊

羊的性情溫順,中文的善字從羊而來；羊是合群動物,中文群字也從羊而來.但羊容易迷路,迷途羔羊；羊群心理也指出羊本身不能獨立判斷,受其他羊群影響.亞理士多德在其的動物誌指出,羊是最愚蠢的四腳哺乳動物,一旦遇到大風雪,如果牧人不驅趕牠,牠就站在裡不動,直至凍死,我們是容易迷失的羊,這個和知識教養沒有關係,因為利令智昏,人怎能憑自己勝過撒旦的誘惑呢?

2008年美國時代周刊專題報導金融海嘯,封面用了「貪婪的代價The Price of Greed」,這些金融界的精英,都是著名大學的優秀畢業生,薪俸豐厚,但因為貪心,用種種財技謀求更大的利潤,引致全球經濟蕭條,多人失去工作,不少人更因破產弄致神經錯亂,甚至輕生.失去靈魂不止「利」字當頭的金融界,在「名」字當頭的高等教育界,也響起失去靈魂的警號.

2007年在哈佛大學教學三十年,作教務長八年的Harry Lewis,寫了一本書:沒有靈魂的卓越-博雅教育還有將來嗎? (Excellence Without  A Soul-Does Liberal Education Have A Future?)他在哈佛大學親身經驗觀察到,像哈佛,耶魯,史坦福等著名大學,慢慢最為追求排名而背棄最初辦學的理想.這些著名大學最初都是由教會或基督徒教育家所創立,透過博雅教育,培育學生有正確的價值觀,耿直的品格,服務社會的熱情.作者認為以人為本的教育網址.edu (education),現在應該改作.com (company) .

中大沈祖堯校長在去年九月對新生入學禮致辭,希望同學明白儉樸生活亦可以過得幸福,寄語學生公平待人,誠實,維護公義,大學的成功不能以排名,建築物多寡,籌得的款項來衡量,而是在於能否培養學生擁有高尚情操,心存謙卑,能過儉樸卻快樂的生活.

這段話在網上熱傳多時,成為一道清泉,因為很多內地高校校長都誇耀畢業同學一出來的薪金多高,硬體設施多好；這也難怪,因為地方政府都是以GDP(國民生產總值)作為升遷的政績.

我們作父母都希望兒女成長,有正確的人生觀,不要迷失本性,失去靈魂.我想麥克亞瑟將軍為兒子的禱告也是天下為父牧者的禱告.我們看看這篇禱文.

在今天的程序表有一禱文:麥克亞瑟將軍為兒子的禱告

上主阿!建立我的兒子.

使他剛強以致能知道何時為軟弱.使他勇敢以致能面對心中所畏懼.

輸得要光明磊落,不折不撓,保有榮譽感.在得勝中,保持謙遜溫柔.

建立我的兒子,使他有高瞻的理想,一位認識你的兒子…….

不要讓他走在安逸,舒適的坦途上,

牽引他在困境,壓力和挑戰的路上,穩步前行.

讓他敢於面對風浪；讓他對失敗的人滿有憐恤.

建立我的兒子,讓他手潔心清,目標遠大.

讓我的兒子首先克勝自己才去領導他人；

懂得歡笑,但不忘記如何哭泣；記取過去的教訓,又能擁抱將來.

我祈求上主成就這一切後,又讓他學會足夠的幽默感,

以致他辦事認真,但不會對自己過份嚴苛.賜他謙虛,以致他記得偉大中的簡樸,

真智慧的開放兼聽,而真剛強也有溫柔的一面.

如蒙垂念,作為一位父親,容我再低聲細稟:「我沒有枉度此生了.」

- General Douglas MacArthur 麥克亞瑟將軍 (伍渭文譯)

每年我都以這禱文送給崇基學院畢業同學-我們的兒女,因為為師者要有為父母的心.

麥克亞瑟將軍是美國二次世界大戰盟軍統帥,韓戰時任聯合國三軍總司令,1942年舉家遷往菲律賓履行二戰責任,當年獲選為全美「年度模範父親」(National Father of The Year) .他說:「我是職業軍人並以此為傲,但我更以作為一位父親為傲,無尚的驕傲.軍人先得拆毀才能建立,但父親祇建立不拆毀.一個有潛在的死亡,另一個隱含創造和生命.死亡軍旅強大,但生命兵團更強大.當我離世後,我盼望我的兒子想起我時,不是在戰場的父親,乃在家中和他一起經常誦唸「我們在天上的父」那一位.(By profession I am a soldier and take pride in that fact. But I am prouder-infinitely prouder-to be a father. A soldier destroys in order to build; the father only builds, never destroys. The one has the potentiality of death; the other embodies creation and life. And while the hordes of death are mighty, the battalions of life are mightier still. It is my hope that my son, when I am gone, will remember me not from the battle but in the home repeating with him our simple daily prayer, ‘Our Father Who Art in Heaven.’)

使我成為一個父親

《建立我的兒子》這篇禱文,感動不少父母,叫他們的靈魂甦醒過來.一位美國教會機構領袖榮納唐(Thom S. Rainer)和兒子1985年 出席全美國青年大會講演時,分享這禱文如何使他的靈魂甦醒過來.他全職讀神學時,因為家庭的需要,除了擔任一小教會牧師外,更需每周在銀行兼職三十個小時,太太就在家中看顧三個五歲,三歲和五個月的兒子,自己因為學業和工作,忽略了作父親和丈夫的責任.有一晚他在一本雜誌看到麥克亞瑟將軍《建立我的兒子》禱告,心中有一種罕有不能言喻的觸動,認為首先不是為兒子禱告,乃為自己禱告,因為父母的生命,深深影響的兒女的生命.

他停下神學院的功課,拿出打字機,改寫這禱文,題目為:「禱告:一位不完全的父親為三位寶貝兒子禱告」(A Prayer: From an Imperfect Father for His Three Priceless Sons) ,禱文首句:「使我成為一個父親,讓我的兒子看到我面對軟弱的能力；面對懼怕的勇氣；接受失敗的恩典；勝出時的謙遜溫柔……Make me the father, O Lord, who will show my sons the strength to face weakness; the courage to face fear; the grace to accept honest defeat; and the humility and gentleness to accept victor…..」

走出死蔭的幽谷vv. 4-6

經文到了四節,主詞從「他」(上主)第三身轉到「我」第一身,而上主稱「你」而不稱「他」了,雖然牧人羊的意象沒有改變,但作者的情緒不同了,「我雖然行過死蔭的幽谷」也不怕遭害.突然,寧靜的田園風雲驟變,日照隱退,烏雲蓋頂,跟著雷電交加,好像貝多芬第六交響曲(田園交響曲)中間的「雷雨Thunderstorm」樂章,一時鼓動鑼鳴,情況危急,險象橫生.

有釋經學者試圖在巴勒斯坦近耶路撒冷附近,尋找那裡是死蔭的幽谷呢?如果我們稍稍了解聖經地理,詩中的青草地不是大草原,像中國青海風高草低見牛羊一望無垠的大草原.詩篇二十三篇的青草地,祇是有水源的小塊草地,行過一個山頭,沒有溪流,就是曠野.而曠野就是逃避敵人,逃避追殺,逃避死亡的地方,是死蔭之地,像美國西部電影,西部曠野是逃命之地,因為敵人會懸紅搜尋你,職業獵犯者也會為了獎金會追殺你.

大衛曾在放牧時和獅子和熊搏鬥,因為牠們把羊羔啣走.「我就追趕牠,擊打牠,將羊羔從口中救出來……我揪著牠的鬍子,將牠打死.」(撒上17 :34)行過死蔭的幽谷不單指面對猛獸的危險,「行過」這詞,應指作為在逃者的一段經歷.他曾在曠野兩度成為在逃者(fugitive),大衛五年在曠野逃避野上司掃羅追捕,稍後又逃避兒子押沙龍的追殺；作為在逃者,死亡如影隨形,所經之處都可能是葬生之地,都是死蔭的幽谷.

這死蔭不單是外在的,也是內心的；獅子來襲,可全心應付,但被曾賞識自己,曾經交心的上司攻擊,被自己養大的兒子追殺,就令到自己心力交瘁.兒子押沙龍對自己不忠,大衛一定想起自己的劣行,如何像種子撒在押沙龍心中；因為別示巴(Bathsheba)的美色,他用計謀殺了烏利亞(Uriah),一位對自己忠心耿耿的將軍,奪去其妻,兒子一定看在眼裡,藏在心內.(撒下11)

同一個辦公室,今天是青草地溪水旁,明天因為人際關係的破裂,可能成為死蔭的幽谷.同一個家庭,因為背叛不忠誠,亦會成為死蔭的空間.有一位大公司的總裁,精明能幹,日理萬機,甚麼人都不怕,但總是怕回家見自己的太太,車開到門口不敢進去,在家附近繞去繞去,被內心的死蔭纏裹,揮之不去.

今天我們可能沒有死亡的威脅,但每一個人都會被一些陰影shadow纏累我們,甚至控訴我們,叫我們失去志氣,失去力量去履行責任.但經文給我們盼望:「你的杖,你的竿,都安慰我,在我敵人面前,你為我擺設筵席,你用油膏了我的頭,使我的福杯滿溢.」(詩23: 4, 5)在危難中,這位牧者保護我們,與我們同在；過去靈魂沉睡時,我們做了深感後悔的事,而被良心和我們的仇敵魔鬼控訴,活在過犯的陰影下,這位牧者也伴隨我們,作我們的中保,因著赦免走出死蔭的陰霾,得著赦罪的平安和喜樂,像路福音十五章11-32節浪子的比喻,父親為悔改歸家的兒子擺設筵席.

而且,他不單塗抹我們的過去,也給我們事奉他的將來:「你用油膏了我的頭」.祇有祭司,先知或君王才被膏油,象徵上主聖靈的同在和祝福,受差遣履行重要的職責.我們可能因犯罪受到擊打,但擊打的目的叫我們醒悟過來,所以「你的杖,你的竿,都安慰我」.杖-短棒,教訓桀敖不馴的羊；你的竿-頭曲的長棍,用來牽引挽回離群的羊,目的使他們走在正確的路上.

在開始時給我們看到青草地,溪水旁田園的寧靜,在結束時我們來到聖殿敬拜的榮耀:「我一生一世必有恩惠隨著我；我且要住在耶和華的殿中,直到永遠.」(v. 6)

直到永遠表示不能奪去的福氣,因這些福氣是牧人犧牲了自己的生命作保證的.

認信耶和華是我的牧者 v. 1

耶穌在約翰福音第十章為詩篇二十三篇這位牧人作了詮釋,耶穌說:「我是好牧人；好牧人為羊捨命.若是雇工,不是牧人,羊也不是他自己的,他看見狼來,就撇下羊逃走；狼抓住羊,趕散了羊群.雇工逃走,因他是雇工,並不顧念羊.我是好牧人；我認識我的羊,我的羊也認識我,正如父認識我,我也認識父一樣；並且我為羊捨命.」(約10: 11, 14)

這位牧人用自己的生命救贖羊群,並愛羊群到底,保守他們,走在義路,直到永遠.「為自己的名引導我走義路」,為自己的名表示我們是屬於他的,我們受洗歸入他的名下,因為這牧人用生命賣贖我們.這詩篇緊隨二十二篇非常有意思,二十二篇是牧人為羊捨命的詩篇:「我的神,我的神!為甚麼離棄我?」(v. 1)「凡看見我的都嗤笑我；他們撇嘴搖頭,說:他把自己交托耶和華,耶和華可以救他吧!耶和華既喜悅他,可以救他吧!」(vv. 7, 8)「犬類圍繞我；他們扎了我的手,我的腳……他們分我的外衣,為我的裡衣拈鬮.」(vv. 16, 18)

經過詩篇二十二篇血腥奮力搏鬥,這位為羊捨命的牧人,配得宣告他在我們生命中的絕對主權,我們是屬於他的,他是屬於我的.詩篇二十二篇就像貝多芬第五交響樂-命運交響樂,激昂奮勇,常用作描述戰爭和凱旋,表達人為生死存活,為命運爭戰；經過奮勇的爭戰,才進入第六交印響曲-田園交響曲的憩靜甯逸,詳和舒泰.是十字架的淒厲流血,天昏地暗,磐石崩裂,宇宙搖盪,帶來赦罪的平安.

耶和華是我的牧者,我必不致缺乏；得著他作為牧者,我們一無所缺.祇有接受他的主權,我們才真正得著滿足.但認了上主作為我的牧者,就不能承認其他作為我的牧者了,主權祇有一個歸屬.

所以,「耶和華是我的牧者」是信仰的宣告.我們可能迷失,靈魂也會沉睡,但不能沒有信仰,放棄信仰.我們可以像羊有時頑梗,桀敖不遜,不祇桀敖不馴,對上主無禮,但總要不放棄信仰就有回轉的機會,再被上帝膏立,被他使用,靈魂可以甦醒過來.

但有一些人是不能叫醒的,就是假裝睡著的人,假裝睡著,就是自己欺騙自己.

李怡先生有一次在報刊專欄中提到一篇「不能叫醒裝睡的人」的文章,在反國民教育時期在網上流傳.作者是一位中文大學同學,來自鄰近地區,她看到公民廣場前「鐵屋吶喊」的橫幅幾乎流淚受觸動寫了這篇長文.她說在鄰近地區,教的老師,聽的學生,幾乎無人相信那些歌頌黨國的鬼話,但人人照說照學照表態,並憑此取得利益和權力.交章說:「國民教育要告訴你的,不是一個假的歷史真相,一個假的真實世界,一個假的良心……只是要告訴你,這個世界,真實不重要,良心不重要,所謂歷史真相也不重要,重要的,只是利益與權力……沉睡的人,你總有辦法去喚醒.但是,你永遠沒有辦法喚醒一個裝睡的人!」

詩篇二十三篇告訴我們一個真理,也是歷史真相:耶和華是我的牧者,我必不至缺乏.這位牧人甚愛我們,為我們的過犯,死在十字架,他要引導我們走有義路,一生以恩惠慈愛伴隨我們,直到永遠.要是作父母要對兒女牢記一句影響他們一生的說話,就是:耶和華是我的牧者,我必必不至缺乏,這是至重要的真理,我們如何回應這真理呢?

詩人回應說:是的,耶和華是我的牧者,我必不致缺乏.這是信仰的宣告,也是主權的交付,承認生命的歸屬.當我們這樣對上主說,你是我的牧者時,我們就回到青草地,溪水旁,安竭下來,靈魂得著甦醒,再次蒙上主膏立,繼續牧養交付給我們的人.阿們.



\newpage

\section{生路一條}
\label{sec:245}
\textbf{賀志勇博士 - 生路一條}
\newline
\newline
\hyperref[sec:244]{< < < PREV SERMON < < <}
~
\hyperref[sec:toc]{[返目錄]}
~
\hyperref[sec:246]{> > > NEXT SERMON > > >}


約翰福音 14:1-11
\newline
\begin{longtable}{cl}
\hline
\hline
章節 & 經文 (和合本修訂版)\\
\hline
14:1 & \begin{tabularx}{0.7\textwidth}{X} 「你們心裡不要憂愁；你們信神，也當信我。 \end{tabularx} \\ \\ \relax
14:2 & \begin{tabularx}{0.7\textwidth}{X} 在我父的家裡有許多住處；若是沒有，我就早已告訴你們了。我去原是為你們預備地方去。 \end{tabularx} \\ \\ \relax
14:3 & \begin{tabularx}{0.7\textwidth}{X} 我若去為你們預備了地方，就必再來接你們到我那裡去，我在哪裡，叫你們也在哪裡。 \end{tabularx} \\ \\ \relax
14:4 & \begin{tabularx}{0.7\textwidth}{X} 我往哪裡去，你們知道那條路。」 \end{tabularx} \\ \\ \relax
14:5 & \begin{tabularx}{0.7\textwidth}{X} 多馬對他說：「主啊，我們不知道你去哪裡，怎麼能知道那條路呢？」 \end{tabularx} \\ \\ \relax
14:6 & \begin{tabularx}{0.7\textwidth}{X} 耶穌對他說：「我就是道路、真理、生命；若不藉著我，沒有人能到父那裡去。 \end{tabularx} \\ \\ \relax
14:7 & \begin{tabularx}{0.7\textwidth}{X} 既然你們認識了我，也會認識我的父。從今以後，你們就認識他，並且已經看見他了。」 \end{tabularx} \\ \\ \relax
14:8 & \begin{tabularx}{0.7\textwidth}{X} 腓力對他說：「主啊，將父顯給我們看，我們就知足了。」 \end{tabularx} \\ \\ \relax
14:9 & \begin{tabularx}{0.7\textwidth}{X} 耶穌對他說：「腓力，我與你們在一起這麼久了，你還不認識我嗎？看見我的就是看見了父，你怎麼還說『將父顯給我們看』呢？ \end{tabularx} \\ \\ \relax
14:10 & \begin{tabularx}{0.7\textwidth}{X} 我在父裡面，父在我裡面，你不信嗎？我對你們所說的話不是憑著自己說的，而是住在我裡面的父在做他的工作。 \end{tabularx} \\ \\ \relax
14:11 & \begin{tabularx}{0.7\textwidth}{X} 你們要信我，我在父裡面，父在我裡面；即使不信，也要因我所做的工作信我。 \end{tabularx} \\ \\
[1ex]
\hline
\hline
\end{longtable}
講題:生路一條　A Way Letting Us Live

講員:賀志勇博士

經文:約翰福音14:1-11

場合:香港中文大學崇基學院禮拜堂

日期:2014 年5月18日

14:1 耶穌又對他們說:“你們心裡不要愁煩；要信上帝,也要信我.14:2 在我父親家裡有許多住的地方,我去是為你們預備地方；若不是這樣,我就不說這話.14:3 我去為你們預備地方以後,要再回來,接你們到我那裡去,為要使你們跟我同在一個地方.14:4 我要去的地方,那條路你們是知道的.”14:5 多馬對他說:“主啊,我們不知道你要到哪裡去,怎麼會知道那條路呢?”14:6 耶穌說:“我就是道路,真理,生命；要不是藉著我,沒有人能到父親那裡去.14:7 你們既然認識我,也會認識我父親的.從此你們認識他,而且已經看見他了.”14:8 腓力對耶穌說:“主啊,把父親顯示給我們,我們就滿足了.”14:9 耶穌回答:“腓力,我和你們在一起這麼久了,你還不認識我嗎?誰看見我就是看見父親.為什麼你還說‘把父親顯示給我們’呢?14:10 我在父親的生命裡,父親在我的生命裡,你不信嗎?我對你們說的話不是出於我自己,而是在我生命裡的父親親自做他的工作.14:11 你們要信我,我在父親的生命裡,父親在我的生命裡；如果不信這話,也要因我的工作而信我.

一,前言

主前500多年,有一個青年被波斯王扔在獅子坑裡,王對他說:“你所常事奉的神,他必救你.”最後,他果然安然無恙的從獅子坑裡出來了.這個人就是但以理!我們現代人會怎麼看這個故事呢?神話,傳說,還是放在一邊,看看就好了?誰會管一個遙遠的獅子坑的故事呢?反正香港連獅子都沒有!去海洋公園不用擔心獅子,也不用擔心鯊魚,因為海族館的玻璃還是足夠結實的!

這是不是真的只是一個神話故事呢?我記得我剛來香港時去了一次晨曦島,島上住的都是吸毒犯,倚靠福音來戒毒.接待的牧師告訴我,他以前就是個大毒梟,在歐洲和香港販毒,結果自己也染上毒癮,怎麼戒也戒不掉.他被抓剛來島上的時候,毒癮發作,難受得在海邊沙灘打滾,恨不得把自己的皮都抓掉.他以為自己沒救了,結果福音救了他.病好之後,他決定奉獻自己做牧師,用福音救更多像他這樣的人.他告訴我說:“對於吸毒太深的人,其他辦法都沒有用了,耶穌就是唯一的生路.”我當時聽了心裡很震動.臨走的時候,他知道我來自大陸,就要我跟他一起唱《我的中國心》,他說他在歐洲販毒的時候,就經常唱這首歌來解鄉愁.看著他喜樂的樣子,我覺得他真的已經被救上岸了.我之後再沒有見過他,雖然當時我還沒信主,但他的話一直在我腦海裡--耶穌就是唯一的生路.但現在耶穌已不在人世,這條路我們該怎麼走呢?或許約翰福音14章1-11節會給我們一些提醒.

二,解經

【14 :1-4】14:1 耶穌又對他們說:“你們心裡不要愁煩；要信上帝,也要信我.14:2 在我父親家裡有許多住的地方,我去是為你們預備地方；若不是這樣,我就不說這話.14:3 我去為你們預備地方以後,要再回來,接你們到我那裡去,為要使你們跟我同在一個地方.14:4 我要去的地方,那條路你們是知道的.”

在約翰福音第 13 章,耶穌給出了一條新命令“彼此洗腳,彼此相愛,跟隨耶穌”.“洗腳”和“跟隨”是這一新命令的核心信息. 但是,耶穌反復強調他要離開門徒去天父那裡,沒有人能夠直接跟著他去.這種強調在門徒的心中激起了極大不安.因為如果耶穌不離開,這些門徒就不用擔心,這樣的話 ,耶穌到哪裡,他們就跟到哪裡；耶穌要他們做什麼,他們就去做什麼.但是,耶穌說他要離開了.

耶穌清楚知道門徒的困惑,所以他安慰他們不要憂傷 .但怎樣才能去掉這些困惑和憂傷呢?耶穌在1-4節強調 : “要信靠上帝,信靠我.”跟隨耶穌的關鍵就是:信靠上帝,信靠耶穌的道路.信靠不僅僅是頭腦裡的相信,信靠是願意根據耶穌的指示去實踐.有許多人知道很多聖經的知識,但卻行不出來,說穿了,原因就是還是不能徹底信靠耶穌.所以,我聽到最多的口頭禪是:“我知道這很好,但是我就是做不到.”

我想到了我老爸,他心臟做了兩次搭橋手術,現在已經有12年了.跟他同時做心臟手術的病友都離世了,但是他還在.為什麼?因為他非常聽醫生的話,這十幾年,他每天走路三個小時,減肥減了40磅；早睡早起,每天的菜大多是素,而且很少很少油鹽.我老爸信靠醫生的話,所以他還活著；而信靠耶穌話的人,比我老爸命更好,他能跟隨耶穌得永生,因為耶穌就是那條生路!.

(14:5 -7) 14:5 多馬對他說:“主啊,我們不知道你要到哪裡去,怎麼會知道那條路呢?”14:6 耶穌說:“我就是道路,真理,生命；要不是藉著我,沒有人能到父親那裡去.14:7 你們既然認識我,也會認識我父親的.

但多馬還是有疑問,5-7節,多馬問道:“主啊,我們不知道你要到哪裡去,怎麼會知道那條路呢?”多馬不知道,耶穌離去,乃是要去上十字架,所以多馬覺得我們需要有知識,如果我們有了“主去的地方的知識”,我們就知道該走什麼路自己到天父那裡去.也就是說,多馬覺得那條道路是可以通過我們自己的知識和計畫去獲得的,哪怕耶穌不在了,我們照樣可以靠著自己的知識和能力跟上去.

但是,耶穌卻認為,道路問題不是一個知識問題,而是一個信心問題.作為門徒,最重要的是憑信心跟隨他.只要門徒相信耶穌的話,他們就應該知道:耶穌本人,或者說耶穌展示的身份˴品格與行動,就正正是通往天國的的道路.在耶穌看來,如果他們相信他的獨子身份,他們就可以自然而然地推導出“去的道路”.但是,門徒推不出在耶穌看來是顯而易見的結論.問題出在哪裡?不在知識,乃在信心.

我們今天有很多的知識,可能比使徒知道得還多得多；以前的人或者沒有聖經,我們今天卻很容易得到一本聖經,甚至有幾本聖經；可是什麼阻攔了我們跟隨耶穌呢?可能的理由是“我工作很忙,”“我要去這城那城做生意,”“我家裡要裝修,”“我親戚朋友要陪伴”；其實真正的原因,是我們對於耶穌的身份還認識不清,我們還沒有全然相信耶穌就是那“道路,真理,生命”,所以也就沒有信心去全心全意跟隨耶穌.

“道路,真理和生命”並不是並排的關係,而是一種遞進關係--那就是,通過基督的道路認識真理,通過真理去獲得生命. 什麼是道路?道路就是你行動的方向.舉個例子來講,我們為自己的生命規劃了許多條道路--怎樣作論文,怎樣去進修,怎樣去升職,怎樣去買樓,我們把那些當做是生命最重要的行動方向,但是通往真理的道路只有一條--跟隨耶穌!正如馬太福音所說“人活著不是單靠食物,而是靠神口裡所出的一切話.”神的話就是真理,而真理就展現在耶穌基督三十三年的生命裡,“謙卑,柔和,善良,憐憫……”,當我們跟隨耶穌的生命這樣去行動時,我們發現我們最終也會具有像基督那樣的性情--“謙卑,柔和,善良,憐憫”,這些我們可以說是“聖靈的果子”,也可以說是從基督而來的生命.但是,很可惜,有時候這些性情離我們還很遠.所以侯活士警告我們說:“有時候,有的基督徒可能會成為別人信基督的最後一個障礙!因為他們彰顯不了耶穌基督的性情!”要警醒!因為耶穌基督的性情,才是真正生命的彰顯,也唯有透過它,我們才真正有生路.

(14: 8-11) . 14:8 腓力對耶穌說:“主啊,把父顯示給我們,我們就滿足了.”14:9 耶穌回答:“腓力,我和你們在一起這麼久了,你還不認識我嗎?誰看見我就是看見父.為什麼你還說‘把父顯示給我們’呢?14:10 我在父的生命裡,父在我的生命裡,你不信嗎?我對你們說的話不是出於我自己,而是在我生命裡的父親自做他的工作.14:11 你們要信我,我在父的生命裡,父在我的生命裡；如果不信這話,也要因我的工作而信我.

雖然耶穌說得很清楚了,但從8-11節來看,耶穌關於信心的一番話並沒有滿足門徒的知識渴望 .老實說,就算耶穌說得再多,門徒也無法完全理解--理解依靠經驗 ,因此,除非門徒已經在信心中跟隨耶穌˴並在跟隨中和真理建立了直接的關係,從而對憑信心跟隨有了經驗,他們才可能真正理解˴真正經歷上帝,也就是“看見”天父 .

我們需要注意腓力說話的口氣:“讓我們看看天父,滿足我們吧!”原文這裡的 “看看”用的是命令語氣,直譯這句話是:“把天父顯現給我們看吧!” 腓力的語氣就證明他還沒有真正理解耶穌,因為如果他真正理解了耶穌的身份˴權柄的話,他豈敢這樣對耶穌說話?

耶穌反問腓力是帶著責備的:“腓力,你還不認識我嗎?” 腓力跟隨耶穌也有多時了,可以說很認識耶穌.但古人說:“傾蓋如故,白頭如新.”意思是說,有的人在車蓋下交談兩句,就已經是老朋友了；有的人認識了一輩子,還像剛剛認識一樣.用廣東話說就是“腓力與耶穌不合嘴型.”為什麼會這樣呢?我想,大概是因為腓力跟耶穌建立的關係還不夠深刻徹底,用廣東話來講就是“感情不夠深”!

如果問大家:什麼樣的人算是好朋友呢?有的人可能會回答:有事肯借錢給我的那個人!還有的人可能會回答:有人講我壞話,我又不在場,肯出頭支持我,撐我的人! 總之,交情不一般!他們為什麼肯這樣做?答案就是:信得過你!所以,當耶穌責備腓力不認識他時,實際上是說:“腓力,你還信不過我麼?你跟了我三年,看見我這麼多的行事為人,你還信不過我麼?” 耶穌三年所做的事,充分彰顯了天父的尊容,門徒只需要跟著做就能不斷領受天父的恩典.但問題在於門徒信不信得過他? 耶穌要腓力“如果不信這話,也要因我的工作而信我.”他其實是告訴腓力:“腓力,你該好好觀察我,觀察我的生命,觀察我的行事為人,如果你真的用心去觀察了這一切,你會相信我的身份和我的宣告的.”

同樣,今日我們在教會裡不斷領受著主的話語,我們信不信得過聖經所見證的那位拿撒勒人耶穌呢?我們有沒有用心觀察過其他基督徒的生命行事,從而看到他背後彰顯的主耶穌呢?這不僅是對腓力,對我們同樣是一個挑戰.

三,結論

現在車子的GPS很發達,你不知道路,GPS可以幫你導航.有一次跟一個朋友開車出去,快到目的地了,GPS老是提醒我們“前方100米左轉”；我們轉了好幾圈,發現總是在原地打轉.下車去問,才知道中間有條路在修路,堵起來了,而GPS又沒有更新,所以就一直提醒我們“前方100米左轉”.人的知識總容易出錯,信錯了東西,就永遠到不了目的地!

在本段經文中,對於“怎樣跟隨耶穌”這個問題,不難看到,門徒的思考方式又回到他們的老路上去了--人總是容易重複自己的弱點:他們覺得,如果自己能夠擁有關於“耶穌去的地方”的知識,即使耶穌不在他們面前,他們也能跟上來.但他們沒有發現,這樣的“跟上”,是依靠自己的知識或能力 “跟上來”,卻不是靠耶穌帶領著“跟上來”; 他們也並不知道,耶穌的離去,乃是要上十字架,這樣一條道路,單靠我們的知識與能力能夠跟上麼?

知識與能力是我們已經曉得,已經掌握在手裡的東西；我們覺得把一切抓在手裡才有安全感,但是“信是所望之事的實底,是未見之事的確據.”我想到了年青而倔強的但以理,波斯王不允許人敬拜除他以外的人或神,違者必扔在獅子坑中；但以理雖然知道這個禁令,但還是像往常一樣,一日三次在家中敬拜耶和華.他不是不曉得危險,但是,他只順從真理,哪怕付上生命--真正的信靠,乃是不根據外在的環境來決定自己應有的行動,是要付上代價的.

親愛的弟兄姊妹,出來走路有危險,不能單靠買保險!“林中有兩條路,”但生路只有一條,你選擇了哪一條呢?



\newpage

\section{基督在我們裡面}
\label{sec:246}
\textbf{鄧瑞強博士 - 基督在我們裡面}
\newline
\newline
\hyperref[sec:245]{< < < PREV SERMON < < <}
~
\hyperref[sec:toc]{[返目錄]}
~
\hyperref[sec:247]{> > > NEXT SERMON > > >}


約翰福音 14:15-24
\newline
\begin{longtable}{cl}
\hline
\hline
章節 & 經文 (和合本修訂版)\\
\hline
14:15 & \begin{tabularx}{0.7\textwidth}{X} 「你們若愛我，就會遵守我的命令。 \end{tabularx} \\ \\ \relax
14:16 & \begin{tabularx}{0.7\textwidth}{X} 我要求父，父就賜給你們另外一位保惠師，使他永遠與你們同在。 \end{tabularx} \\ \\ \relax
14:17 & \begin{tabularx}{0.7\textwidth}{X} 他就是真理的靈，是世人不能接受的。因為他們既看不見他，也不認識他；你們卻認識他，因他常與你們同在，也要在你們裡面。 \end{tabularx} \\ \\ \relax
14:18 & \begin{tabularx}{0.7\textwidth}{X} 我不會撇下你們為孤兒，我必到你們這裡來。 \end{tabularx} \\ \\ \relax
14:19 & \begin{tabularx}{0.7\textwidth}{X} 再過不久，世人不再看見我，你們卻會看見我，因為我活著，你們也要活著。 \end{tabularx} \\ \\ \relax
14:20 & \begin{tabularx}{0.7\textwidth}{X} 到那日，你們就會知道我在父裡面，你們在我裡面，我也在你們裡面。 \end{tabularx} \\ \\ \relax
14:21 & \begin{tabularx}{0.7\textwidth}{X} 有了我的命令而又遵守的人，就是愛我的；愛我的人，我父要愛他，我也要愛他，並且要親自向他顯現。」 \end{tabularx} \\ \\ \relax
14:22 & \begin{tabularx}{0.7\textwidth}{X} 猶大（不是加略人猶大）問耶穌：「主啊，為甚麼親自向我們顯現，而不向世人顯現呢？」 \end{tabularx} \\ \\ \relax
14:23 & \begin{tabularx}{0.7\textwidth}{X} 耶穌回答他說：「凡愛我的人就會遵守我的道，我父也會愛他，並且我們要到他那裡去，與他同住。 \end{tabularx} \\ \\ \relax
14:24 & \begin{tabularx}{0.7\textwidth}{X} 不愛我的人就不遵守我的道。你們所聽見的道不是我的，而是差我來之父的。 \end{tabularx} \\ \\
[1ex]
\hline
\hline
\end{longtable}
講題:基督在我們裡面　Christ is In Us

經文:約翰福音14:15-24

講員:鄧瑞強博士

場合:香港中文大學崇基學院禮拜堂

日期:2014年5月25日

各位弟兄姊妹,早安.

中文大學新亞書院有一個「合一亭」,是記念錢穆先生的.錢穆先生以九十六歲高齡,寫下他人生最後一篇文章:<中國文化對人類未來可有的貢獻>.這篇文章有時稱為＜天人合一論＞.錢穆先生認為,「天人合一」論是中國文化對人類作出的最大貢獻.

讓我們簡約地讀讀這篇重要文章.

錢穆先生說:「從來世界人類最初碰到的困難問題,便是有關天的問題.…西方人喜歡把『天』與『人』離開分別來講.換句話說,他們是離開了人來講天.…中國人是把『天』與『人』和合起來看.中國人認為『天命』就表露在『人生』上,離開『人生』,也就無從來講『天命』.離開『天命』,也就無從來講『人生』.…西方人常把『天命』與『人生』劃分為二,他們認為人生之外別有天命,顯然把『天命』與『人生』分作兩個層次,兩次場面來講.如此乃是天命,如此乃是人生.…所以西方文化顯然需要另有天命的宗教信仰,來作他們討論人生的前提.…總之,中國古代人,可稱為抱有一種『天即是人,人即是天,一切人生盡是天命的天人合一觀』.…即如孔子的一生,便全由天命,細讀《論語》便知.…倘孔子一生全可由孔子自己一人作主宰,不關天命,則孔子的天命和他的人生便分為二.…就此而言,孔子的人生即是天命,天命也即是人生,雙方意義價值無窮.換言之,亦可說,人生離去了天命,便全無意義價值可言.….孔子為儒家所奉稱最知天命者,其他自顏淵以下,其人品德性之高下,即各以其離於天命遠近為分別.」(參《新亞生活》36:4 (2008年12月))

以我有限的見識來說,以上引文有幾個重點.

一,人類其中一個大問題,就是有關「天」的問題.我十分同意,「天」的問題,是人生的一大問題.人類生命有限,但有一種追求永恆的天性.人類仰賴物質而活,但尋求一種物質以外的價值.追求「天」,追求「超越」,追求「神」,是人之常情.

二,中國文化的特質,是將「天」與「人」合起來說,即所謂「天人合一」.「天命」內貫於「人生」裡.的確,沒有「天命」的「人生」,是淺陋不堪的.遠離「人生」的「天命」,則無處可尋.

三,單純講「天人合一」,畢竟仍屬抽象之論,必須具體地說.錢穆先生以孔子作為「天人合一」的具體說明.孔子作為一完美人格的代表,他的生命全依「天命」.「天命」透過孔子這具體的人,展現其內容.沒有孔子,天命無從揭露.沒有天命,孔子只是凡夫.當孔子以其「人生」展露「天命」為何時,「天人合一」便有了具體說明.「天命」流行於孔子生命裡,他成了有別於他人的「聖人」.他的生命成了一標準,判定其他人的「人生」價值.

四,相對於中國文化言,西方文化將「天」與「人」分開,在「人生之外別有天命」.「天命」既在「人生」以外,尋索「天命」,便「需要另有天命的宗教信仰」了.按錢穆先生的講法,西方的宗教自有一超越的理,卻與「人生」兩相分離,不能在「人生」中體悟「天道」.

錢穆先生的文章,點出了幾個人生大問題.

一,「天」,透過什麼途徑展露其內容?

二,「天」,是一個什麼樣的天?

三,「天」,與一般世人關係如何?

今日的講道,將以這幾個問題為中心.

在錢穆先生眼中,基督教算是「西方」的宗教.基督教講的「天」,一定是與「人」分隔的了.這是事實嗎?

有了以上資料作基礎,我們看看今天的講道經文:約14: 15-24

約14:15　「你們若愛我,就必就遵守我的命令.

約14:16　我要求父,父就另外賜給你們一位保惠師〔或譯:訓慰師〕,叫他永遠與你們同在,

約14:17　就是真理的聖靈,乃世人不能接受的；因為不見他,也不認識他.你們卻認識他,因他常與你們同在,也要在你們裡面.

約14:18　我不撇下你們為孤兒,我必到你們這裡來.

約14:19　還有不多的時候,世人不再看見我,你們卻看見我；因為我活著,你們也要活著.

約14:20　到那日,你們就知道我在父裡面,你們在我裡面,我也在你們裡面.

約14:21　有了我的命令又遵守的,這人就是愛我的；愛我的必蒙我父愛他,我也要愛他,並且要向他顯現.」

約14:22　猶大(不是加略人猶大)問耶穌說:「主啊,為甚麼要向我們顯現,不向世人顯現呢?」

約14:23　耶穌回答說:「人若愛我,就必遵守我的道；我父也必愛他,並且我們要到他那裡去,與他同住.

約14:24　不愛我的人就不遵守我的道.你們所聽見的道不是我的,乃是差我來之父的道.

首先,我們討論一下,「天」透過什麼途徑展露其內容?

孔子曰:「吾十有五而志於學,三十而立,四十而不惑,五十而知天命,六十而耳順,七十而從心所欲,不踰矩.」

一個人努力求道,努力體會生命意義何所寄託,努力實現價值,則去到某個階段,便能「知天命」,明白「天」的心意.明白過後,便會按天的心意去理解萬事萬物,不再因別人的批評,謾罵,惡言惡語而改變自己的人生方向,此之謂「耳順」.然後,不為名利所誘,不為世物所纏,不為逆境所困,而能大膽地實現「天」的價值,此之謂「從心所欲,不踰矩」.就是透過這種聖人人格的實踐,「天」展露其內容.聖人不出,天道不顯.聖人出,天道朗現.這就是錢穆先生所言,「天命」不離「聖人」的「人生」.

今日的講道經文出自約翰福音.約翰福音很重視「天」與「人」的關係.

約翰福音開首時說:「太初有道,道與神同在,道就是神.」(約1:1)天道即神聖.然後,「道成了肉身,住在我們中間.」(約1:14)約翰說:「從來沒有人看見神,只有在父懷裡的獨生子將他表明出來.」(約1:18)

天父無形,以基督為形.耶穌將天父的生命表達出來.

耶穌說:「你們所聽見的道不是我的,乃是差我來之父的道.」(約14:24)

耶穌一如孔子,內心全無自己的心念,只有「天命」,「天道」,「天意」.從耶穌生命流露出來的,正是「天」的內容.耶穌體驗這上「天」,為一「有情之天」,他稱之為「天父」.

有一次,耶穌的門徒腓力對耶穌說:「求主將父顯給我們看,我們就知足了.」耶穌對他說:「腓力,我與你們同在這樣長久,你還不認識我嗎?人看見了我,就是看見了父；你怎麼說『將父顯給我們看』呢?」(約14:8-9)

天父透過耶穌展露其「容貌」,耶穌透過順從天父旨意去活出生命的意義.基督教的「天命」與「人生」,並不如錢穆先生講的,分隔得很遠.

現在,我們討論一下,這個「天」,是一個什麼樣的天?

按錢穆先生的講法,「天」與「人」要和合起來說,互為表裡.順這思路而下,便很容易得到「天道即人心」的想法.

孟子說:「盡其心者,知其性也；知其性,則知天矣.」人回復自己的本性,便知道什麼是「天」了.孟子又說:「萬物皆備於我矣.反身而誠,樂莫大焉.強恕而行,求仁莫近焉.」天之道,全然在我心內.回到自己,則有大樂.努力接納他人,「仁」便得到實現了.人的「仁心」擴而充之,便是天道大行人間.天道無他,唯「仁」而已.

在這一點上,基督教的確與儒家有別.

面對上天,耶穌呼喚其為「天父」.天父不離基督,卻與基督有別.天父透過基督去揭露自己,基督實現天父的旨意去完成自己人間的使命.天父與基督,關係密切.但基督教有一「三位一體」的教義,去確保天父與基督有別.

關於耶穌,基督教以「神人二性」的教義,去描述耶穌的生命.耶穌兼具「神人二性」,表明「天」不離「人」,「天命」不離「人生」,神性與人性連在一起.關於基督的「神人二性」到底如何,神學家於主後451年,開了一次神學大會,確立一「迦克頓信經」(Creed of Chalcedon).這信經確立基督教對基督神人二性的正統理解.這信經說:基督的神性和人性,兩者「不相混亂」(inconfusedly),「不相交換」(unchangeably),「不能分開」(indivisibly),「不能離散」(inseparably).簡單說,在基督生命裡,神性不離人性,人性不離神性；沒有離開神性的人性,也沒有離開人性的神性；但神性依然不同於人性,神性是神性,人性是人性,兩者不會交換性質.

這「迦克頓信經」,一方面說「天命」不離「人生」,另方面卻要保證,天人有別,「人」不能與「神」混淆.這教義表明,不能直接將「人心」等同「天命」.

從這點看,基督教是有「天人合一」的一面,但同時,卻又指出「人」不是「天」,免得有任何人神化自己.基督教對人的「狂妄自大」很敏感.人最大的罪,莫如以為自己是「神」.

「天」不是「人心」的無限擴充,那麼,「天」是什麼?

面對上天,耶穌卑微地呼喚「天父」.這不是人內在的道德心的外化,而是一有情的大生命.人面對「天父」,不是一個肉軀的小我,面對自己的道德心靈外化的「大我」.天父是人之外,人之上的大生命.

耶穌的一生,順服天父的心意.除了天父的心意,他沒有自己的心意了.從這一點講,他與父原為一(約10:30).從另一方面講,天父是耶穌之外的偉大生命,臨到他,呼喚他,要求他.就這而論,父是比他大的(約14:28).

天父,不是我們的自我的投射.天父是偉大無比的生命,他臨到我們,叫我們謙虛,叫我們走出自己去回應身外的呼喚,叫我們思想自身生命之外的偉大世界.

倫理學者利維納斯(Emmanuel Levinas)指出,我們身邊的弱少者向我們發出的求救呼喚,可看成是無限偉大的上帝對我們的呼喚.他們呼喚我們走出自大的自己,走出自滿自足的自己,走出虛偽而不可一世的自己.

最後,我要講一下這個「天」與一般世人關係如何?

公都子問孟子:「鈞是人也,或為大人,或為小人,何也?」孟子曰:「從其大體為大人,從其小體為小人.」

從儒家看,有些人能像孔子那樣,回到本心,以「人生」展現「天命」,則成為「大人」.另一些人,「耳目之官不思,而蔽於物」,而成為「小人」.關鍵在於,是否能展現人心裡的一點靈明,一點仁性?

耶穌說:「你們若愛我,就必就遵守我的命令.」(約14:15)愛,並不是那種主觀的感覺.愛,是遵守耶穌的命令.遵守耶穌的命令,便是遵守耶穌背後的天父的命令.這天父會將真理的聖靈賜給我們,這聖靈會在我們生命裡(約14:17).如此,天父的心意會在我們裡面,而實踐天父心意的力量也在我們裡面.這是世人不認識的,猶如「小人」不明白「大人」.不願意用自己的生命實現天父旨意的人,當然沒有天父的真理在心中,也不會明白那是什麼東西.

耶穌說:「還有不多的時候,世人不再看見我,你們卻看見我；因為我活著,你們也要活著.」(約14:19)耶穌乃是展現天父心意的人,若我們誠心實現天父心意,會在自己生命裡,發現基督生命的痕跡.我們會與基督的生命,連在一起,兩者都是以卑微的人生,活出天父的大愛.如耶穌所言:「到那日,你們就知道我在父裡面,你們在我裡面,我也在你們裡面.」(約14:20)

天父的心,貫通耶穌的生命.我們因著愛耶穌,甘願實現天父的心,則天父的心及耶穌的心,會在我們生命裡,永遠與我們同在.

看來,基督教的「天命」,都是不離「人生」的.「天命」透過「人生」而彰顯.但為了避免「人生」誤入歧途,基督教仍堅持「天命」高於人,讓人謙卑,讓人細心聆聽自身之外的他人靈魂的呼喊.能聆聽的,甘願回應的,便像基督一樣,以生命揭示天父.這樣的一個人,與基督同活,被聖靈充滿,去展示天父的無限仁慈.

每一個人都被呼召進入這樣的生命裡?你願意聆聽這呼召嗎?

但願:榮耀歸於聖父,聖子,聖靈.阿們.



\newpage

\section{澆灌前夕的等候}
\label{sec:247}
\textbf{伍渭文牧師 - 澆灌前夕的等候}
\newline
\newline
\hyperref[sec:246]{< < < PREV SERMON < < <}
~
\hyperref[sec:toc]{[返目錄]}
~
\hyperref[sec:248]{> > > NEXT SERMON > > >}


使徒行傳 1:6-14
\newline
\begin{longtable}{cl}
\hline
\hline
章節 & 經文 (和合本修訂版)\\
\hline
1:6 & \begin{tabularx}{0.7\textwidth}{X} 他們聚集的時候，問耶穌：「主啊，你就要在這時候復興以色列國嗎？」 \end{tabularx} \\ \\ \relax
1:7 & \begin{tabularx}{0.7\textwidth}{X} 耶穌對他們說：「父憑著自己的權柄所定的時候和日期，不是你們可以知道的。 \end{tabularx} \\ \\ \relax
1:8 & \begin{tabularx}{0.7\textwidth}{X} 但聖靈降臨在你們身上，你們就必得著能力，並要在耶路撒冷、猶太全地和撒瑪利亞，直到地極，作我的見證。」 \end{tabularx} \\ \\ \relax
1:9 & \begin{tabularx}{0.7\textwidth}{X} 說了這些話，他們正看的時候，他被接上升，有一朵雲彩從他們眼前把他接去。 \end{tabularx} \\ \\ \relax
1:10 & \begin{tabularx}{0.7\textwidth}{X} 他升上去的時候，他們定睛望天，看哪，有兩個人身穿白衣站在他們旁邊， \end{tabularx} \\ \\ \relax
1:11 & \begin{tabularx}{0.7\textwidth}{X} 說：「加利利人哪，你們為甚麼站著望天呢？這離開你們被接升天的耶穌，你們見他怎樣升上天去，他也要怎樣來臨。」 \end{tabularx} \\ \\ \relax
1:12 & \begin{tabularx}{0.7\textwidth}{X} 有一座山，名叫橄欖山，離耶路撒冷不遠，有安息日可行走的路程。那時，門徒從那裡回耶路撒冷去， \end{tabularx} \\ \\ \relax
1:13 & \begin{tabularx}{0.7\textwidth}{X} 他們一進城，就上了所住的樓房；在那裡有彼得、約翰、雅各、安得烈、腓力、多馬、巴多羅買、馬太、亞勒腓的兒子雅各、激進黨的西門，和雅各的兒子猶大。 \end{tabularx} \\ \\ \relax
1:14 & \begin{tabularx}{0.7\textwidth}{X} 這些人和幾個婦人，包括耶穌的母親馬利亞，和耶穌的兄弟，都同心合意地恆切禱告。 \end{tabularx} \\ \\
[1ex]
\hline
\hline
\end{longtable}
講題:澆灌前夕的等候　The Quite Awaiting before The Grand Downpour

經文:使徒行傳1:6-14

講員:伍渭文牧師

場合:香港中文大學崇基學院禮拜堂

日期:2014年6月1日

今天是教會年的升天主日,耶穌復活後四十日升天,應該是星期四,教會以之後的主日為升天主日Ascension Sunday.再過一星期是復活後五十日的五旬節或聖靈降臨(Pentecost)主日.對當時的門徒來說,耶穌升到天上是喜也是憂.喜是耶穌升天坐在父上帝的右邊,安坐王位(Enthronement),表明他的救贖工作被父上帝肯定,與父同掌王權,這是耶穌救贖工作的一部份,信經這樣表述:「在本丟彼拉多手下受難,被釘十字架,死了,葬了；下到陰間,第等三天從死裡復活；後升天,坐在全能父上帝的右邊；將來必從那裡降臨,審判活人死人.」

但另一方面,他們要適應耶穌和他們同在的方式了.使徒行傳告訴我們,耶穌復活後,「用許多的憑據將自己活活地顯給使徒看,四十天之久向他們顯現,講解神國的事.」(使1: 3)他們已習慣耶穌親自在他們中間教導他們,他們親眼看到他的樣子,親耳聽到他的聲音,但耶穌升天後,此情不再了,他們會很不習慣,甚至產生極大的焦慮.

1. 不再憑眼見,要憑信心

上帝在世界缺席嗎?

一般教會會都紀念聖誕和復活節,聖靈降臨節也有紀念但不多,升天節就更少了.但升天節的信息,對今日教會有適切的提醒.升天是執掌王權,但似乎我們眼睛見到的,上帝在世界缺席,上主好像不在場.就如我們有爸爸或媽媽,不是孤兒,但爸爸媽媽總是忙這個忙那個,最需要他們時他們都不出現,缺席的父母.今天詩篇經課提到上帝是駕行諸天,但同時看顧地上的孤兒寡婦.「神在他的聖所作孤兒的父,作寡婦的伸冤者.神叫孤獨的有家,使被囚的出來享福；惟有悖逆的住在乾燥之地.」(詩68: 5,6)但實際上,有時面對社會種種不義,惡人的囂張,公義得不到伸張,上主好像缺席,不理我們.

看不到上帝,上帝就不理我們嗎?

我們很想上主給我們解釋,很想像耶穌復活後升天前的四十天,耶穌就在我們眼前,有疑問就問過明白,無論答案如何,祇要耶穌親口答我們,我們就安心.但耶穌升天後,我們不能很快問過清楚明白.耶穌要我們學習不需憑眼見,也曉得上主與我們同在.就像小孩子進幼兒院離開父母,看不到父母,以為父母不理他們.父母就是不讓孩子看到自己,他們躲在旁邊,眼目沒有離開過孩子,有時和當席老師談論孩子,也不讓孩子們知道,目的讓孩子們成長學習,和同學們一起成長學習.孩子要明白他們可能看不到父母,但父母總會在放學時出現,接他們回家,或安排人接他們回家.父母關心他們在學校的生活,他們要學習信靠父母.

要明白神國的隱藏性

四十天耶穌對門徒顯現,講論甚麼呢?講論神的國.而神的國不是憑眼見,神的國是憑信心.

門徒在耶穌快要升天之際,還期望升天所代表的掌權,就是上帝國度明顯的展現:「主啊,你復興以色列國就在這時候嗎?」(使1: 6) 耶穌說:「父憑自己的權柄所定的時候,日期,不是你可以知道的.」升天代表得著所有的權柄,但這權柄不是眼睛看得到的,要用信心體會的權柄.復興以色列國,在地上完全實現上帝的公義,時候還沒有來,還需要忍耐等候.

十字架已一役竟其功,成了!

基督被釘在十架,人看來是失敗,但憑信心已一役竟其功.當耶穌在十架喊叫,成了!撒旦就徹底被打敗,現在祇是「執埋手尾」,清理殘餘.當然,撒旦還在欺騙我們,裝模作樣,好像甚麼都沒有發生,好像還不知鹿死誰手.復活對憑眼見的人來說,祇是墓石滾開,墳墓空了,屍首不見了,裹屍布像人形完完整整放在一起,是一宗懸疑不解的事.不錯,表面看來,墓園的溫度和濕度沒有變化,巴勒斯坦仍然由羅馬中央派來的巡撫彼拉多統治,而且因為耶穌的屍首不見了,當權者還加派羅馬兵丁維穩,羅馬政權仍然穩如泰山.但憑信心我們已曉得世界已不一樣,我們的罪已被赦免,我們窺探到死亡不是生命的終局,復活如嶄新的靈風,吹進世界每一角落,更新萬有,復活打開了密封的世界.

應用:門徒要打破舊有的框框,來思考和耶穌的關係.初信主時我們祈禱比較個人,而且覺得很蒙上主垂聽,極有可能上主因為我們初信,以此給我經歷他的真實,但當我們靈命成熟時,我們就要學習多為其他人代禱,有時禱告還遲遲未蒙應允,學習等候上主自己的時間.

今天也許在我們的生命有很多疑團,很想上主給我們答案,但到現在上主仍然默然不語,也許直到見主面時才完全明白.上主過去如何引導我們,但將來不一定用相同方法引導我們.過去上主給你很容易相處的上司,將來的上司不一定那麼容易相處.但要緊記,雖然我們眼睛看不到上帝,但上帝的眼睛沒有離開我們.

2. 我們繼續耶穌的故事Jesus’ Story Goes on In Us

「從今以後,我不在世上,他們卻在世上；我往你那裡去.聖父啊!求你因你所賜給我的名保守他們,叫他們合而為一,像我們一樣.」(約17: 11)「你怎麼差我到世上,我也照樣差他們到世上.」(約17: 18)

馬太,馬可,路加福音都以耶穌升天為壓卷的事蹟,但耶穌的救贖工作才剛開始,我們要受聖靈的澆灌,得著能力,把福音傳遍地極.「不要離開耶路撒冷,要等候父所應許的,就是你們聽見我所說過的.約翰是用水施洗,但不多幾日,你們要受聖靈的洗.」(使1: 4)路加視福音書為前書:「提阿非羅啊,我已經作了前書,論到耶穌開頭一切所行所教訓的.」(使1: 1) 序幕剛揭起,好戲在後頭.

1942年英國軍隊在埃及Alamein首次戰勝德國「沙漠之狐, Desert Fox」隆美爾Rommel,之前沒有勝過,之後也沒有敗過.當時英軍士氣大振,以為敗隆美爾一役後,大戰很快就結束,時任首相邱吉爾(Winston Churchchill) 發表演說:」This is not the end, it is not even the beginning of the end, but it is perhaps the end of the beginning.」( 這不是結束,也不是結束的開始,也許可以說,剛剛開了個頭)

如果新約聖經以四卷福音書為高潮,之後的使徒行傳和書信便是反高潮；但福音書不是記錄一粒麥子如何死,乃說這粒死了的麥子如何結出許多的子粒來.福音書記錄耶穌在世上一個非常重要的工作-呼召和訓練門徒-胚胎的教會.福音書是序幕,使徒行傳才是戲肉.耶穌升天,聖靈降臨,教會成立,福音傳遍天下,耶穌再來,建立新天新地.

今天福音經課:「從今以後,我不在世上,他們卻在世上；我往你那裡去.聖父啊!求你因你所賜給我的名保守他們,叫他們合而為一,像我們一樣.」(約17: 11)耶穌升天不在世上,我們卻在世上,稍後耶穌禱告:「你怎麼差我到世上,我也照樣差他們到世上.」(約17: 18)

我們怎知道上帝沒有離開這世界呢?我們如何知道上帝是真實的呢?就是因為教會-基督的身體.今天強調個人的喜好,我行我素,對教會嘲諷多於愛護,升天節提醒我們,教會在地上要繼續耶穌的故事.

a. 不濟事的教會存留到現在,有上主的看顧

薄伽丘(Giovanni Boccaccio,1315-1375)《十日談》(il Decameron) 是義大利文藝復興時代的著名短篇小說,第二天提到一個故事,說到一位信主的猶太商人楊諾·德·雪維尼傳福音給一位名叫亞伯拉罕的猶太教朋友,也是商人.這位未信主的商人跟傳福音給他的基督徒說:

「楊諾,你一心要我改信基督教,現在我也同意了,不過還得先讓我到羅馬去一遭,瞻仰一下你所謂天主派遣到世上來的‘代表’,看看紅衣主教的作為和氣派.使我有所感悟,否則我還是信我的猶太教.」楊諾聽他說要到羅馬看看教會,心裡焦急,因為想他看到教會的實在情況,一定不會改信基督教.從羅馬回來後,亞伯拉罕真的歸歸依基督信仰了.原因呢?

亞伯拉罕說:「我只覺得羅馬不是一個‘神聖的京城’而是一個容納一切罪惡的大溶爐:照我看,你那位高高在上的神職人員,本該做教會的支柱和基礎,卻正日日夜夜,用盡心血,千方百計,要叫教會早些垮台,但教會還是屹然不動,倒反而日益發揚光大,這使我認為一定有聖靈在給它做支柱,做基石.我信了,給我行洗禮吧.」

b. 信徒的生命見證,顯出教會是屬靈的生命有機體

教會是屬靈生命的有機體,不是因為不被腐敗侵蝕,乃有一班活出信仰的信徒,保持著信仰的貞潔和承傳.比薄伽丘早一世紀的十三世紀,歐洲城市興起,商貿發展,大學湧現,道明會就成立為「傳道者修會」(Order of Preacher),他們其中知名的領袖是多瑪斯.亞奎那,銳意向智識份子解釋信仰.城市興起也帶來貧富懸殊,方濟會成立,扶傷濟貧,修士諾守神貧,生活簡樸,以愛心幫助被邊緣化的族群,創會人為聖方濟或聖法蘭西斯(St. Francis of Assisi) ,著名的聖法蘭西斯禱文相傳是他所寫,其實此禱文最早見於1912年十二月巴黎一份名叫La Clochette (小鐘League) 的靈修刊物,出版者為神聖彌撒聯會(La Ligue de la Sainte-Messe) ,作為彌撒(崇拜) 的禱告.因為禱文內容,會通聖法蘭西斯的精神,故托其名.

使我作你和平之子,在憎恨之處播下你的愛；在傷痕之處播下你寬恕；在懷疑之處播下信心.

使我作你和平之子,在絕望之處播下你盼望；在幽暗之處播下你光明；在憂愁之處播下歡愉.

使我作你和平之子,在赦免時我們便蒙赦免；在捨去時我們便有所得；迎接死亡時我們便進入永生.主啊使我少為自己求,少求受安慰,但求安慰人；少求被瞭解,但求瞭解人；少求愛,但求全心付出愛.

二次大戰後,這禱文被大量印刷傳閱,對如何面對後戰爭的仇恨,虛無,悔疚,如何復和解怨起了治癒作用.

上帝為甚麼讓苦難臨到我們?若你是馬航HM370乘客的親人,或南韓歲月號沈船269遇難者的家屬,你都會問為甚麼發生這些災難?為何是他們的親人?也許我們說這是人為災難,與上主扯不上關係,但2003年南亞海嘯,2008年汶川大地震,這些天災又如何解釋呢?十三年半作為校牧期間,遇到重大災難時,不少學生和同事都會問我,為何如此呢?

其實,作為牧師,我也會問上主,為何如此呢?為何上主好像缺席呢?

我答:我不知道為何有這些苦難,上主行事隱秘難明；但上主的慈愛顯而易見,因為耶穌為我們釘十架.

我常引用申命記 29:29,「隱秘的事是屬耶和華我們神的；惟有明顯的事是永遠屬我們和我們子孫的,好叫我們遵行這律法上的一切話.」沙士疫症,海嘯,地震是隱秘的事,屬於上帝；但救援工作是明顯的事,屬於我們,愛人如己的是誡命.

當上帝透過被聖靈感召的人,參與苦難的救援工作,上帝就在當中,沒有缺席.何處有苦難,何處就看到上帝的眼淚；何處有苦難,何處就聽到基督在哭泣.耶穌在馬太福音二十五章告訴我們,在未日審判時,能進神的國「因為我餓了,你們給我吃,渴了,你們給我喝；我作客旅,你們留我住；我赤身露體,你們給我穿；我病了,你們看顧我；我在監裡,你們來看我.」(太25: 35, 36)上帝沒有缺席,祇不是按我們的方法和方式出席.

應用:耶穌升天,教會在地,我們繼續耶穌的故事.我們被他差遣,帶著使命和能力,在地上作明顯的事,遵行上帝的誡命,叫人曉得,隱秘行事的上帝,並沒有缺席.

3. 同心等候,回應時代的召命

耶穌升天的信息教導我們:不要憑眼見,要憑信心；也提醒我們,教會在世界是代表上帝,繼續耶穌在的工作.最後,我們看看這班門行如何等候聖靈的降臨.

「當下,門徒從那裡回耶路撒冷去,進了城,就上了所住的一間樓房,在那裡有彼得,約翰,雅各,安得烈,腓力,多馬,巴多羅買,馬太,亞勒腓的兒子雅各,奮銳黨的西門和雅各的兒子猶大.這些人同著幾個婦人和耶穌的母親馬利亞,並耶穌的弟兄,都同心合意地恆切禱告.」(使1: 13)他們同心禱告,等候上帝的作為,靜心等候時代的召命.

這個簡單的樓房,是普世宣教的發源地.在禱告中,他們眾人如同一人,他們靜靜等候上帝.等候不是簡單的事,特別當我們對現實不滿時；以為祇要行動,事情就可以解決.耶穌給他好的榜樣,用三十年預備三年半的工作,復活後用四十天好好教導門徒有關天國的真理,並要他們不要分散,一起禱告等候.他們要好好預備,不能靠血氣,要靠聖靈的能力；不能單打獨鬥,要同心合意.

十天之後,聖靈降臨,「忽然,從天上有聲下來,好像一陣大風吹過,充滿了他們所坐的屋子,又有舌頭如火焰顯現出來,分開落在他們的頭上.他們就都被聖靈充滿,按著聖靈所賜的口才說起別國的話來.」(使2: 2-4)

等候,在禱告集體等候,使我們得到共識,同感一靈,看準上帝在此時此刻要我們參與的工作,回應時代的召命.

馬丁路德(Martin Luther)和伊斯拉姆(Erasmus) 都是歐洲文化史上曠世奇才,兩人學識淵博,大將風範.伊斯拉姆是希臘文專家,1516年出版新約希臘文；路德1517年張貼九十五條論綱(Ninety Five Theses) ,指出當時教會背離聖經的教訓,開始了改教運動(Reformation 樣Movement) .兩人對當時教會的腐敗,同樣心痛欲絕.有人評論兩人:伊斯拉姆準確看出那一棵樹幹在朽爛,但路德就舉起斧頭砍去那腐朽的樹幹.路德批評伊斯拉姆:他應該為教會落淚時,竟然嘲諷教會.

在開始時,伊斯拉姆名氣很大,歐洲的王侯將相都和他通訊,請益治國之道.相比之下,路德不過是一名修士,大學的聖經教授；但伊斯拉姆比較個人主義,沒有路德的勇氣.路德聽到時代的召命,和他的改教團隊,回應歷史,浩浩蕩蕩,逆轉教會的路線,回到聖經.但伊斯拉姆自改教運動開始後,就漸漸沈寂下來,歷史缺席,非常可惜.

我常常覺得,因為種種歷史原因,崇基學院能夠在香港中文大學擁有校牧室和神學教育,是上帝給我們的機會,而崇基禮拜堂的主日崇拜,也有其獨特的時代使命,這是大學的教會,我們在這裡聚會並非偶然.讓我們在升天節的主日,像昔日的門徒同心禱告,眾人如同一人,一起等候上帝,願意被上帝使用,一起建立這主日崇拜,藉此推動校園福音工作.

總結:耶穌升天,我們從此不再憑眼見,要憑信心,信心使我們看到肉眼看不到的上帝；這上帝並沒有在世界缺席,透過他所差遣的教會,我們顯明上主並未遺忘這世界；而當教會同心等候,一起為福音齊心努力,我們就和道上帝在這時此地,要我們一起為他作的工,回應時代的召喚.



\newpage

\section{領受靈氣平安伴隨}
\label{sec:248}
\textbf{許開明牧師 - 領受靈氣,平安伴隨}
\newline
\newline
\hyperref[sec:247]{< < < PREV SERMON < < <}
~
\hyperref[sec:toc]{[返目錄]}
~
\hyperref[sec:249]{> > > NEXT SERMON > > >}


講題:領受靈氣,平安伴隨　Receiving the Holy Ghost, Peace Be With You All

講員:許開明牧師

場合:香港中文大學崇基學院禮拜堂

日期:2014年6月8日

引言

人生命中最強大的敵人,最常出現,最具威力的是什麼呢?就是『懼怕』 .宗教的產生,起源於人對大自然的恐懼,對存在的恐懼,對死亡的恐懼.基督教的福音,也是「使人不懼怕」的福音.請留意,基督的降生,天使的喜訊是「不要懼怕!」復活主常說的一句話是「願你們平安!」事實上,許多人是因為「怕」而尋找神,依靠神.

基督教的信仰所關注或強調的,不止是「除去懼怕」,而且是更進一步「勇敢的面對懼怕」,以平安入世,勇敢為人的態度,面對人生.為何如此說?在今日的講道中,會嘗試加以解釋.

近代心理學家發現,在人類眾多的懼怕中,基本上有三類的懼怕:第一類是怕跌倒,第二類是怕巨響,第三類是怕被遺棄.

第一類,怕跌倒:包括怕失敗,怕出錯,怕站不住,怕不合格,怕失去:失去健康,失去錢,失去生命等.第二類,怕巨響:包括怕突如其人的災難,風災,雨災,火災,空難,海難,地震等.第三類,怕被遺棄:包括怕失戀,怕被父母子女遺棄,怕被出賣,怕被人陷害,怕被人取笑,怕不被重視等.因此,誰可以 一無所懼呢?

再者,以上三種懼怕是會互有關連和和相通的,即是第一類懼怕,會引發第二類的懼怕,再引發第三類的懼怕.使到懼怕加劇,令人無法勝過,抵受不住.

所以懼怕是人生最大的敵人.它足以快速地嚇死,也可以慢慢地把人折磨至死.

從約20章的記述,首先,門徒把門關上,因怕猶太人會突然來提拿他們；這是屬第二類,怕巨響的懼怕.繼而,門徒怕一旦被捉拿,便會被監禁,會失自由；會被打,會失去健康；會被殺,失去生命, 這是屬第一類,怕跌倒,怕失去的懼怕.進而,門徒怕信耶穌後,被自已同胞猶太人,甚至家人親友的遺棄.這明顯是屬第三類,怕被遺棄.

倘若我們明白這點,你便會明白為何見過復活主的門徒仍然會活在恐懼之中.而且他們懼怕的感覺慢慢地加劇.縱然他們見過復活主仍未能脫離懼怕.直到耶穌說,「願你們平安! 」再加上向他們吹一口氣,情況才慢慢好起來!

今日,同樣許多相信耶穌的人仍然會有懼怕?

許多人有一種錯誤的觀念,一個信耶穌的人是不應該懼怕的,懼怕是無信心,靈性軟弱的表現.我並不同意上面的看法,信耶穌的人仍然有懼怕的.

某教會的牧者與信徒前往探望一位教會的長老.這位長老知道在世的日子無多,他坦誠的表示他知道快離世很不開心,因為捨不得太太和子女；他也自覺在事奉上有虧欠,怕見主面.一位老姊妹聽見很快回應說:「長老,你是不應該不開心和怕死的,快見主面不是好得無比的嗎!點解你會病到無晒信心?」牧師見狀,即時說:「姊妹,不如你出一出去,我想個別跟長老交通一下!」這位直腸直肚的姊妹有一種錯覺,長老是不應有恐懼不安的.事實上,信主的人也是人,面對死亡也會有不願意,不開心,他在牧師面前,坦誠說出來,好過臨終還放在心裡,得不到釋放.

想問,恐懼是否毫無價值?它既然是所上帝容許的,它是否也有意義?從另一角度看,人的憂慮和恐懼是正常的,本身確有意義,是好的.人怕火,怕水,怕電,使人避開危險.我們不會讓年幼的子女獨留家中,因為他們「唔醒驚」,所以有危險!一位精神科醫生說:「憂慮和恐懼好像血壓,沒有不能活,過高使人受損.事實上它們是我們生存不可缺的一種情緒.」從信仰層面看,人因懼怕,尋找信仰,依靠上帝.人要懂得懼怕上帝,才能敬畏(Fear)上帝.

問題是一個人常常被「懼怕」纏繞,就活得不快樂；懼怕會使畏縮,失常,失去做人的力量；更可怕的是活在懼怕中的人,往往會成為可怕的人,使周圍的人不安.人在可怕的環境中,會懼怕!但是有時在安舒的環境中,也會自己嚇自己.真奇怪,今日許多人是活在自己嚇自己的情況中.

曾經在台灣聽過一個故事:傳說在基隆的某路段常有鬼魂出沒,許多人經過那路段時心裡總是不安.有一天,一個爸爸載兒子去上幼稚園,途經那路段時,孩子突然手指著電線桿說:「爸爸,你看電線桿上有兩個人!」老爸抬頭一看什麼也沒有,心想:「該不會是碰上那種東西吧?還是快走為妙!」老爸就飛快地離開了.

到了傍晚,媽媽去接兒子放學,又經過那路段時,孩子又突然手指著電線桿說:「媽媽,你看電線桿上有兩個人!」這回真把媽媽嚇死了.

這件事也不知怎麼傳開的,連電視台都來採訪,還架好攝影機,對準那支電線桿,然後請那孩子過來,準備開拍.工作人員就問:「你說的兩個人在哪裡?」孩子手指著電線桿說:「就是那裡啊,那裡有兩個人啊!」大家順著孩子的指示往上一看,都呆了!

真的有兩個人:「交通安全, ”人人”有責」!

小心,不要被「懼怕的幻覺」操縱著你的生命,自己嚇自己.

人遇到突如其來事,總會被懼怕操縱,想不好的事!所以匪徒有機可乘.不知,你有否接聽過這樣的電話「你個仔在我們手中,快拿贖款來!」電台不停提醒,仍有人繼續未醒,受騙!

耶穌賜我們平安,勝過懼怕

1「願你們平安,跳出懼怕!

這是在耶穌復活的第一天晚上,門徒因為害怕猶太人,就關門待在家中.耶穌來到他們中間說:「願你們平安」,又把釘痕給門徒看,他們就歡喜.耶穌將平安與喜樂帶到門徒中間.【約廿19】

「門徒因為害怕猶太人,就關門待在家中」, 懼怕,使人關起門,困在屋中,不敢見人,不敢作事,失去做人的能力.

耶穌來到他們中間說:「願你們平安」.注意,這是一種在逼迫恐懼中的平安.

耶穌是穿牆過壁,突然在他們中間,這顯出復活主的大能；同時,顯出祂的愛.耶穌不叩門,不用開門,免得他們受驚!當人關上門,怕受逼害時,最驚是開門!耶穌明白他們的需要,不使他們受驚,在平安的環境中,把平安賜他們.然而,領受了「平安」,便要勇敢的,受差遣到不安的環境中,傳楊平安的福音.

還有,一個奇妙的地方,在最近讀經時才發現,耶穌要讓門徒看見,他是曾經經歷過「不安」的主.「說了這話,就把手和肋旁,指給他們看；門徒看見主,就喜樂了.」【約廿20】耶穌就把手和肋旁,指給他們看,手的釘痕,肋旁的傷痕:這是那從死復活的主,是擁有身體的主,不是人的錯覺,更不是幻覺.原文「看見」”cido”:不止是看到,且是看清楚,看得明白.他們就喜樂了!這復活主也是受過試探,受過苦,受過痛,受過死的主,因此他能明白我們,知道我們的傷和痛的主.

但如今這復活主,復活了! 在門徒眼中,復活主是相當奇妙的,也是奧妙難明的.復活主是有身體,卻不被身體限制,仍然能穿牆過壁!超越肉身,超越牆壁,超越距離,超越時間,超越的上帝.祂不被限制,也不被阻攔!

復活主也能超越自己有身體的局限,人最偉大的地方是能超越自己的局限.舉例來說,無手的人仍能畫畫,用腳畫畫；無腳的人仍能行路,用手行路；無眼的人仍能講到書,用手讀書.超越自己的無能和有限!是需要一股強大的心力和能力.

從神學上來看耶穌基督,特別是從約翰的神學來看耶穌基督:超越的道,甘願進入有限人身(肉身),受人身的限制,但祂也能脫離人身,超越有限,進入無限.既是「全能」也可以「不能」,又由 「不能」回到「全能」 .這是奧秘,超越人的理性,邏輯與想像.

耶穌既可以是神也可以是人,既有神性也有人性,且神性與人性共存,互不排斥,互不衝突,互(相)得益彰.如果耶穌只有神性,無入性,就高高在上,不可親近!如果耶穌只有人性,無神性,就軟弱不堪,不能相助!耶穌既是神,有神性,就能相助!

耶穌既是人,有人性,就可親近,並能經歷人間的不安,痛苦,能體恤我們的軟弱!這就是祂的偉大之處!自古到今,神學家,可能各走極端.只強調耶穌的神性時,好像中古世紀神學家,過份強調耶穌的神性,只有榮耀的基督,萬主之主,萬王之王.這耶穌高高在上,不食人間煙火,不知世人痛苦.這個基督不是聖經裡的基督.

另方面,只強調耶穌的人性時,好像近代的一些神學家,受到人文主義的影響,把基督過份人性化,再加上去神話化,耶穌與人一樣受苦受壓,是被釘的十架,這位耶穌與我們非常接近,甚至沒有分別,少講復活的基督,就算復活都只是精神復活.這樣過份人性化的基督,永遠都像把耶穌釘在十字架受苦,也不是聖經裡的基督.

對我這研究聖經和神學多年的人,我仍然堅持相信一位擁有神性和人性,既可「全能」也可「不能」,既「無限」也可「受限」,既「能屈又能伸」 的耶穌基督,這麼可愛又可敬的基督,在我研究過宗教後,發覺只有基督教的耶穌基督督是最明顯的.

願你們平安,接受差遣

接受差遣,領受聖靈,進入懼怕,勝過懼怕

接著,又對他們說:「願你們平安」,就差遣他們,賜他們聖靈與赦罪的權柄.使門徒在受差遣的路上,有聖靈同在的平安與能力.

門徒得到主賜的平安後,主就差遣他們,賜下聖靈,進入世界.從此他們勝過懼怕,打開門,進入世界,去到人群.他們不再怕猶太人,羅馬人,因有聖靈的同在!

第一個:「願你們平安」,使他們跳出恐懼

第二個:「願你們平安」,並接受聖靈,使他們勝過恐懼,是更高的層次的平安.

到底平安是什麼?

「平安」是一句奇妙的話,它不只是一句貼心的問候語,更重要的是它背後的信念,在它背後的信仰.

為何耶穌在復活後,多次見到門徒總是問候他們,「願你們平安」.「平安」這句話不只是一個習慣性的問候語,且有更豐富的意義.

當中有兩二大特色:

第一大特色:總的來說,「平安」乃顯出美好 .我們常用「平安」互相請安,因為我們明白「平安」的意義.在希臘文中「平安 (eirene)」的字義是「平安,和平,和睦」.在希伯來文中是shalowm,意思是「平安,快樂的,友好的,福祉,健康,興盛」. 「願你們平安」,「願上帝將各樣美好/美善都賜給你!」

第二大特色:平安兩字在不同的場景,有不同的狀態.

a國與國際間:平安是和好,和平沒有戰爭

路加福音14:31-32提到,「一個王出去和別的王打仗,豈不先坐下酌量,能用一萬兵去敵那領二萬兵來攻打他的嗎?若是不能,就趁敵人還遠的時候,派使者去求和息的條款.」其中這個「和息」就是「平安」的意思.「和息」達到「平安」的目的,創造「和平」的關係.

b在人際間:平安是和睦的關係

保羅在哥林多前書 7:15b說:「神召我們原是要我們『和睦』」.這是指著人與人之間的關係說的,這和睦就是「平安」.又在寫給以弗所信徒的信中指出,耶穌基督以自己的身體廢掉猶太人與外邦人間隔離的牆,使人與人之間,在基督的生命裡合而為一.保羅說:「如此便成就了『和睦』」(弗2:15b),這和睦就是「平安」.

c 在崇拜中:平安是有秩序,和諧

保羅在哥林多前書14:33說:「因為神不是叫人混亂,乃是叫人『安靜』」.這個「安靜」是「混亂」的相反詞,是「有秩序」的意思.所以,保羅的話在指出,神所要給人的不是混亂,乃是「秩序」.這「秩序」就是「平安」.

D在人與人的道別中:「平安」乃表達祝福

「平安」在希臘文中,「 eirene」,希伯來文shalowm的意義.「平安」如果用在臨別贈語時,意思等於「祝你健康」的意思.保羅和西拉在傳道的途中,曾一度被囚在腓立比,他們把福音傳給了看守他們的獄卒,當官長們決定要釋放保羅和西拉時,獄卒就很高興的對他們說:「官長打發人來釋放你們,如今可以出監,平平安安地去吧.」(徒 16:36).這個「平平安安」有「平安」和「祝你健康」的含意.

「願你們平安」(Peace be with you)

約翰記載耶穌顯現的事件中,他連續三次提到耶穌說:「願你們平安」. 「願你們平安」是「願平安與你同在!」無論你去到那裡,平安都跟隨你.還有意思是「願上帝將各樣的美善賜給你!」

這「平安」是一種禮物,豐富的恩典.

這「平安」是一種好的願望.

這「平安」是一種力量/能力.

3 領受聖靈

『(耶穌)說了這話,就向他們吹一口氣,說:你們受聖靈!』(約二十:22).

今日是教會年曆:聖靈降臨節主日,願我們對聖靈有所認識.聖靈(Holy Spirit/ Holy Ghost),靈原文是“氣”,氣,令你想起什麼?

1氣是看不到,2生命攸關,無氣就無命,空氣,

3力量有關,有氣有力.生氣:破壞力量,正氣:做正事的力量,勇氣:不怕困難的力量.

聖靈是什麼?

3.1氣與「氣息」

『(耶穌)說了這話,就向他們吹一口氣,說:你們受聖靈!』(約二十:22).耶穌為何要給門徒吹氣?有象徵意義嗎?有何神學意義?

從創造的故事看,神將生氣吹在亞當鼻孔裡,亞當成為活人.為何神要給亞當吹氣,神說:「我們要照著我們的形像,按著我們的樣式造人,使他們管理海裡的魚,空中的鳥,地上的牲畜,和全地,並地上所爬的一切昆蟲」(創一:26).神要人有與祂相似的形像,因此吹一口氣在亞當的鼻孔裡,就是要人與神有相類似的形像.從中可見,幾項重要的思想:

a神要人能管理/治理全地,(能力,權柄)

b生養眾多(數量)

c神賜給他有智慧可以為各種動物命名(智慧).

但是,因著亞當夏娃犯罪(失去人的能力,智慧,德性),亞當夏娃就失去屬靈的能力,失去做人的能力.耶和華說:「人既屬乎血氣,我的靈就不永遠住在他裡面；然而他的日子還可到一百二十年」(創六:3)這一節經文是要表明人屬乎血氣的話,神的靈就無法繼續住在人的裡面.

3.2聖靈是「手」

神讓以西結看見遍滿骸骨的平原耶和華的靈(注:原文作“手”)降在我身上,耶和華藉他的靈帶我出去,將我放在平原中,這平原遍滿骸骨(結三十七:1).於是,我遵命說預言.正說預言的時候,不料,有響聲,有地震；骨與骨互相聯絡.我觀看,見骸骨上有筋,也長了肉,又有皮遮蔽其上,只是還沒有氣息(結三十七:7,8).

骸骨可以復活嗎? 氣息 rûach(風,靈),有神的氣息就有新的活力.當以西結向風發預言,氣息就進入骸骨,骸骨就活了,並且站起來,成為極大的軍隊. (結三十七:10).

上帝給亞當吹一口氣,成為有靈的活人

以西結給枯骨吹口氣,使枯骨生出筋和骨,合起來成為強大的軍隊.

耶穌給門徒吹一口氣,門徒受聖靈,建立教會,掌握天國的匙,赦免誰的罪,誰的罪得赦.

門徒受靈氣,重新再造,勝過懼怕,負起責任和使命,勇往直前.門徒受聖靈,仍然會面對恐懼不安,只是能夠勝過恐懼,在恐懼中前進.

聖靈的同在,帶來平安,平安帶出一種定力.一位兄弟向雜技員學行銅線.第一天開始,先學行麻繩,且離地一尺,學足一年,都是如此!直到第二年,才行鋼線.為什麼?因為要克服恐懼,人 有定力,才能勝過懼怕.聖靈是氣,給人加力,得著氣力；還有聖靈是推動力,使人勇敢向前,不為懼怕所勝!

反思:六四事件中的恐懼

1989年,當權者因恐懼而屠殺大派手無寸鐵的大學生.被殺學生的爸母站起來,有良知的人起來,要平反六四,當權更害怕起來了!他們的懼怕,使他們繼續打壓,使周圍的人也恐慌起來!人在恐懼中,會被懼怕操縱,以致失去理智,做出愚蠢的事.

我們要跳出恐懼,敢說應該說的事,敢做應該做的事!我十分欣賞德國的潘霍華牧師(Rev. Bonhoeffer),因他能不被恐懼所操縱,冒死反對納粹德國希特拉(Hitler)的統治!因他的死也激勵多人勇敢的站起來!潘霍華牧師說:「當基督呼召一個人的時候,祂要這個人跟從祂,為祂捨命.」”When Christ calls a man, He bids him come and die.’’ 為基督喪掉生命,不退縮,不跑掉! 我們需要主賜的平安,不是無災無難,無風無雨的平安,而是在人間的風雨中仍能前進的平安.

我也很喜歡一首聖詩:《堅固保障》A Mighty Fortress is Our Lord . 想起一位偉大的宗教改革家,馬丁路得(Martin Luther 1483- 1546),他因天主教教廷的打擊而逃往德國的一古堡名(瓦爾特堡, Wartburg),寫作的一首詩:「上主是我堅固保障,永不動搖之保障…」我曾到過Wartburg,並在那裡默念Martin Luther.一份正氣凜然的感覺湧現.

回想今日在打壓下的民運人土,仍敢站出來的民運人士,流亡海外,經過25年可以安安逸逸過活,仍敢站出來說話!繼續為民主而努力的人士.十八萬人在維園漆黑中,燃起燭光的的參加者,渴求民主自由的燭光,遠望一大片火光,產生了一股極大的力量,我看見了在黑暗中的希望.

結語

在講道開始時,我說:「基督教的信仰所關注或強調,不止是「除去懼怕」,而且是更進一步「勇敢的面對懼怕」,以平安入世,勇敢為人的態度,面對人生.」

耶穌應允我們,賜給我們兩項寶貝:1是「平安」的力量,使人得到「定力」,面對懼怕.2是「聖靈」 的同在,使人加力,並有推動力,受差遣,進入懼怕中,勝過懼怕.

各位,你當下有些什麼懼怕嗎?

那日晚上,耶穌站在門徒中間說:「願你們平安!」跟著向他們吹氣說:「你們受聖靈」.今日復活主對每個真誠相信祂的人也說「願你們平安!願你接受聖靈的同在,不要懼怕,平安入世,勇敢為人!」



\newpage

\section{美好}
\label{sec:249}
\textbf{林豪恩先生 - 美好}
\newline
\newline
\hyperref[sec:248]{< < < PREV SERMON < < <}
~
\hyperref[sec:toc]{[返目錄]}
~
\hyperref[sec:250]{> > > NEXT SERMON > > >}


創世記 1:2-4
\newline
\begin{longtable}{cl}
\hline
\hline
章節 & 經文 (和合本修訂版)\\
\hline
1:2 & \begin{tabularx}{0.7\textwidth}{X} 地是空虛混沌，深淵上面一片黑暗；神的靈運行在水面上。 \end{tabularx} \\ \\ \relax
1:3 & \begin{tabularx}{0.7\textwidth}{X} 神說：「要有光」，就有了光。 \end{tabularx} \\ \\ \relax
1:4 & \begin{tabularx}{0.7\textwidth}{X} 神看光是好的，於是神就把光和暗分開。 \end{tabularx} \\ \\
[1ex]
\hline
\hline
\end{longtable}
講題:美好　Beautiful and Good

經文:創世記1章2-4節

講員:林豪恩先生

場合:崇基學院禮拜堂主日崇拜

日期:2014年6月15日

一．引言

那些年,日本有一位享譽盛名的歌手,名叫五輪真弓；香港有一位出色填詞人,名叫鄭國江,他把五輪真弓的其中一首歌填上廣東話歌詞,名叫「為什麼」,其中第一段如下:

「為什麼生世間上　此間許多哀與傷

為什麼爭鬥不絕　歡欣不永享

問為何人存隔膜　顏面無真相

問那天可找得到　理想中的烏托幫」

世事紛擾,憂患不絕,還有甚麼可以用上「美好」去述說的呢?這些年,流行的描述是:「This city is dying, you know?」因此,有網友認為這是一個不適合生孩子的世界:「 我們的下一代會在甚麼環境下成長? 色情,毒品,校園欺凌,數唔晒地充斥著,總之社會風氣敗壞……過呢個世界好彩好命ge 人一定係少數,唔好彩唔好命ge 人一定係大多數,就好似貧富越來越懸殊,所以下一代九成都係延續上一代ge 不幸……空氣有毒,水又有毒,食物又有毒.你話,生個仔出來係咪即係害左佢?」

「美好」這個詞彙,似乎已經離開人們的嘴唇很遠,或者也漸漸從人們的腦海消失.更不幸的是,「美好」這個詞彙被使用時,例如某人的名句:「香港好,祖國好,祖國好,香港更好」不但沒有令人覺得美好,反而令人心情複雜.

縱使如此,美好並沒有消失.今天的舊約經課,上帝一連多次對著我們的世界讚賞說美好.在上帝的眼中,世界如何好,如何美呢?上帝的評價對我們有甚麼啟迪呢?

二．釋經和應用

A. 上帝看著是好的

1. 有光有暗是好的

1:1起初,上帝創造天地.

1:2地是空虛混沌,淵面黑暗；上帝的靈運行在水面上.

1:3上帝說:「要有光」,就有了光.

1:4上帝看光是好的,就把光暗分開了.

1:5上帝稱光為「晝」,稱暗為「夜」.

天地仍是一片漆黑,甚麼也看不見的時候,稱不上好.突破黑暗,需要有光.當光劃破黑暗,帶來光明,上帝認為這樣的狀態是好的.「上帝稱光為「晝」,稱暗為「夜」.有晚上,有早晨」(1:5).有一套電影名為「黑夜話事人We Own the Night 」,他們以為,白日是警察管治,晚上是黑幫的世界.人們也有一種錯覺,以為白日是屬於正義的,晚上是屬於邪惡的.也有人以為白日是屬於上帝的,晚上是屬於魔鬼的.因此,有些人害怕夜晚.晚上是否黑暗無光呢?晚上是否屬於邪惡或者魔鬼的呢?第14-18節有這樣的補充:

1:14上帝說:「天上要有光體,可以分晝夜,作記號,定節令,日子,年歲,

1:15並要發光在天空,普照在地上.」事就這樣成了.

1:16於是上帝造了兩個大光,大的管晝,小的管夜,又造眾星,

1:17就把這些光擺列在天空,普照在地上,

1:18管理晝夜,分別明暗.上帝看著是好的.

原來,夜晚沒有甚麼值得害怕,因為夜晚並不是完全黑暗.廣闊的夜空中躺著皎潔的明月,掛著閃閃的星星,安靜地,輕柔地照著大地.夜晚不是屬於邪惡的,也不是屬於魔鬼的,如白日一樣,夜晚是屬於上帝所創造的光體所管理的.白日陽光燦爛,夜晚月色皎潔,晝夜各有精彩,上帝認為這樣的循環是好的.詩人曾經有這樣的經驗和發現:

詩139:11我若說:黑暗必定遮蔽我,我周圍的亮光必成為黑夜；

詩139:12黑暗也不能遮蔽我,使你不見,黑夜卻如白晝發亮.黑暗和光明,在你看都是一樣.

原來,有日有夜,有光有暗,如此普通和平常的循環,已經是美好的.

2. 有地有海是好的

1:6上帝說:「諸水之間要有空氣,將水分為上下.」

1:7上帝就造出空氣,將空氣以下的水,空氣以上的水分開了.事就這樣成了.

1:8上帝稱空氣為「天」.

1:9上帝說:「天下面的水要聚集在一處,使乾地露出來.」事就這樣成了.

1:10上帝稱乾地為「地」,稱聚集在一起的水為「海」.上帝看為好的.

當天地仍然是汪洋一片,空虛混沌,大地被水覆蓋的狀況下,稱不上好.上帝對空虛混沌的狀況作出規劃,創造空間把水上下分開；然後,上帝又把下面的水在平面上分開,陸地出現了.有天空,有海洋,有陸地的景致,上帝認為是好的.喜歡遠足的朋友很可能有這樣的體會,在藍天白雲之下,站在青翠的山崗上,眺望碧海青天,這才是真正的心曠神怡,而不是住在日出康城.

3. 地上有植物是好的

1:11上帝說:「地要長出植物,就是含種子的五穀菜蔬,和會結果子,果子裡有種子的樹,在地上各從其類.」事就這樣成了.

1:12於是地長出了植物:含種子的五穀菜蔬,各從其類；會結果子,果子裡有種子的樹,各從其類.上帝看為好的.

陸地出現了,地面上怎樣的狀態是美好的呢?原來,上帝認為大地上有植物,而且是會自我繁殖,結出果子的植物,是好的.為甚麼呢?因為有植物,動物才有生存的條件；因為有果子,人類的生命才得到延續.

4. 空中有飛鳥,水裡有生物是好的

1:20上帝說:「水要滋生眾多有生命之物；要有鳥飛在地面以上,天空之中.」

1:21上帝就創造了大魚和在水裡滋生的各樣活動的生物,各從其類,以及各樣有翅膀的鳥,各從其類.上帝看為好的.

大地有植物是美好的,那麼,天空和海洋怎樣的狀態是美好的呢?上帝認為天空有鳥飛翔,海洋有魚游動的狀況是好的.海中的大魚,有時被古人誤作海怪,甚至被視為邪惡,但其實都只是上帝的創造之一,沒有甚麼值得懼怕.

5. 地上有走獸是好的

1:24上帝說:「地要生出有生命之物,各從其類,就是牲畜,爬行動物,地上的走獸,各從其類.」事就這樣成了.

1:25於是上帝造了地上的走獸,各從其類；牲畜,各從其類；和地上一切的爬行動物,各從其類.上帝看為好的.

天空有飛鳥在飛翔,海洋有魚兒在滶游,陸地上似乎仍然很平靜.創造者的畫筆輕掃大地,大地變得生氣勃勃,熱鬧起來.畫面上有走獸在奔跑,有牲畜在行動,有昆蟲在爬行.看著這圖畫,上帝也不禁讚嘆美好.

6. 在其中有人類,就非常好

1:26上帝說:「我們要照著我們的形像,按著我們的樣式造人,使他們管理海裡的魚,天空的鳥,地上的牲畜和全地,以及地上爬的一切爬行動物.」

1:27上帝就照著他的形像創造人,照著上帝的形像創造他們；他創造了他們,有男有女.

1:28上帝賜福給他們,上帝對他們說:「要生養眾多,遍滿這地,治理它；要管理海裡的魚,天空的鳥和地上各樣活動的生物.」

1:29上帝說:「看哪,我把全地一切含種子的五穀菜蔬和一切會結果子,果子裡有種子的樹,都賜給你們；這些都可作食物.

1:30至於地上一切的走獸,天空一切的飛鳥,並一切在地上爬行的,有生命的動物,我把綠色植物賜給牠們作食物.」事就這樣成了.

1:31上帝看一切所造的,看哪,都非常好.

有光有暗是好的,有海有地是好的,地上有植物是好的,天空有飛鳥是好的,海中有魚兒是好的,地上有動物是好的,怎樣可以令這幅圖畫更美好呢?上帝最後加上一筆,人類就出現在這些美好的狀況之中.在一切的創造之中加上人類,人類負責「管理海裡的魚,天空的鳥,地上的牲畜和全地,以及地上爬的一切爬行動物」,這樣的狀態,上帝認為比之前更好,因此說:非常好!

B. 人類看為是好的

創造者認為美好的世界,何以和今天人們的體驗如此不同呢?在上帝的創造中,光不會對創造者說:這樣不好.地不會對創造者說:這樣不好.植物不會對創造者說:這樣不好.飛鳥,魚兒,動物都不會對創造者說:這樣不好(至於那條會和人說話的蛇,應該不算入動物的類別吧!).只有人類會對創造者說:我認為你這樣判斷不好,如果我的判斷不比你更好,也應該和你一樣好.

3:1耶和華上帝所造的,惟有蛇比田野一切的活物更狡猾.蛇對女人說:「上帝豈是真說不許你們吃園中所有樹上的果子嗎?」

3:2女人對蛇說:「園中樹上的果子,我們可以吃,

3:3惟有園當中那棵樹上的果子,上帝曾說:『你們不可吃,也不可摸,免得你們死.』」

3:4蛇對女人說:「你們不一定死；

3:5因為上帝知道,你們吃的日子眼睛就明亮了,你們便如上帝能知道善惡.」

3:6於是女人見那棵樹的果子好作食物,也悅人的眼目,且是可喜愛的,能使人有智慧,就摘下果子來吃了,又給她丈夫,她丈夫也吃了.

第一章上帝看著是「好」這個「好」字,在第三章這個段落兩次了出現,其一次是「你們便如上帝能知道善惡」的「善」字,其二次是「於是女人見那棵樹的果子好作食物」的「好」字.蛇建議人挑戰創造者,要自己和祂一樣判斷甚麼是好,甚麼是不好.於是,人把自己當成上帝.在第一章經常出現的句子是:「上帝看著某某是好的」,現在變成:「女人見那棵樹的果子好作食物」.人們要自己決定何為好,並且努力去改變天地,使之體現人們認為的美好.

C. 是好的和還不好的

1. 創造者認為有日有夜是好的,人們卻認為「黑夜如白晝發亮」更好,於是,有些人飽嚐光污染的惡果.香港璀璨的夜景一向為港人所自豪,但在漂亮的裝飾背後,我們要付出代價──光污染除了使星星被遮掩,更對生態,環境和市民的健康造成影響.有鬧市業主被閃動招牌害得「神經衰弱」；有居民亦慘受室外強光之擾,沒有一覺好睡.

2. 創造者認為有海洋有陸地是好的,不過,有些人認為把海洋變成陸地更好.吊詭的是,當某些人努力移山填海,增加陸地的時候,另外一些人卻連自己站腳之地也保不住.有否記得,有本地動話電影中一隻小豬仔說:「我聽日飛啦!」「馬爾代夫--藍天白雲水清沙幼椰林樹影…… 」不過,一個殘酷的事實是,由於全球暖化,馬爾代夫即將消失.如果聯合國計算準確的話,最快一個世紀這些島嶼將被海水逐一吞噬.

3. 創造者認為大地被植物覆蓋是好的,不過,有些人卻認為砍伐樹本更好.聯合國專家估計,全球每分鐘就有60英畝熱帶雨林被砍伐.這意味著每年有5萬平方英里(相當於大半個英國)的熱帶雨林消失.大地被人類「增值」之後,土地出產增加,不過,所生產出來的,卻又未必是人們想要的.據官方報告顯示,中國有近五分之一的耕地受到污染.2013年4月,幾批來自主要大米出產省份湖南的大米中被發現含有大量鎘.這一發現引發了外界對耕地的擔憂,並一度使得湖南大米價格下跌達14\%.(cn.wsj.com/big5/20140417/bch173038.asp)

4. 人類按其認為好的方式去改造大地,大地漸漸變成抬頭望不到飛鳥,俯首看不到魚兒,地上失去走獸的蹤影.據聯合國有關報告顯示,自工業革命開始,全球已進入最大規模的物種滅絕時代.目前,全世界每天都有上百個物種消失.人類在這場危機裡扮演了舉足輕重的角色-因其自身的活動,例如大批的森林,草原,河流消失了,而把其他物種的自然滅絕速度提高了100至1000倍,而當地球生態系統被損害到一定程度時,必將導致人類賴以生存的體系崩潰. (\url{http://hk.crntt.com/doc/1005/7/0/7/100570759.html})

5. 人類也不覺得自己很好

諷刺的是,人類能夠按自己認為好的方式去改造世界的時候,人類自己也未必互相認為好.反而,人們在多數情況下,都評價對方不好,例如這個民族認為那個民族不好,這個人認為那個人不好.John Gottman是一位研究溝通的學者,尤其集中於夫婦的溝通.他發現四樣破壞婚姻的溝通方式,稱之為「未日四騎士」,因為他發現,若夫婦相處上大量出現這四種行為,他們的婚姻離未日不遠.這四位騎士的相同地方都是指責對方不好.

第一,批評:批評有別於投訴,投訴是對個別處境和為的不滿,批評是對一個人的性格廣泛的抱怨.「你又出錯了,你有甚麼事情做得令我稱心的呢?」說這類話對關係的破壞性頗高,被批評的一方會感到被人身攻擊.

第二,蔑視:當一些負面情緒未能化解,我們會對另一方的負面想法和評價增加,甚至升級到要羞辱和輕視對方.例如:「這種事你都做得出,你好有家教啊!」這種輕視和羞辱對方的話,自然會引動第三位騎士出現.

反擊和自衛:受襲的一方,感到受傷之餘,會將內心的憤怒化為自衛和還擊.例如:「你不比我好得多呢!你看你這副怪相,哪有人會喜歡.」

當一個人認為另一個人不好的時候,這個人本身會很不快樂,同時也會製造痛苦給另一個人.父母認為「生舊叉燒」好過自己的子女,太太埋怨自己嫁錯郎,丈夫後悔自己娶錯妻……如是者,大家一齊陷入痛苦的深淵.

四．總結

今天的經課以詩篇第八篇回應創世記第一章:

8:1 耶和華我們的主阿,

你的名在全地何其美．

你將你的榮耀彰顯於天.

8:2 你因敵人的緣故,

從嬰孩和喫奶的口中,建立了能力,

使仇敵和報仇的,閉口無言.

8:3 我觀看你指頭所造的天,

並你所陳設的月亮星宿,

8:4 便說,人算甚麼,你竟顧念他．

世人算甚麼,你竟眷顧他.

8:5 你叫他比天使微小一點,

並賜他榮耀尊貴為冠冕.

8:6-8 你派他管理你手所造的,使萬物,

就是一切的牛羊,田野的獸,

空中的鳥,

海裡的魚,

凡經行海道的,

都服在他的腳下.

8:9 耶和華我們的主阿,

你的名在全地何其美.

如果詩人果真是大衛,他一定深明世途險惡,也處身於世事紛亂之中.他必定發現,圍繞著他的人,敵我難分,就算是自己的兒子,也是謀朝篡位之輩.其實,在夜深人靜之時,仰望長空,撫心自問,大衛也會看見自己比他們好不了多少.詩篇第八篇是讚美創造主的詩,不過,同時也是一首自我反省的詩.如果詩人的眼目仍然被明爭暗鬥所控制,他便不會有機會「觀看你指頭所造的天,並你所陳設的月亮星宿」；如果沒有機會仰望星空,他便不會有「人算甚麼,你竟顧念他．世人算甚麼,你竟眷顧他」的自省和感悟.諸葛亮在前出師表說:「臣本布衣,躬耕南陽,苟全性命于亂世,不求聞達于諸侯.先帝不以臣卑鄙,猥自枉屈,三顧臣于草廬之中,諮臣以當世之事,由是感激,遂許先帝以驅馳.」當人類自覺「人算甚麼,世人算甚麼」的時候,竟然發現創造主仍然委以重任「派他管理你手所造的,使萬物,就是一切的牛羊,田野的獸,空中的鳥,海裡的魚,凡經行海道的,都服在他的腳下.」詩篇第八篇作為回應,呼籲人們,體認自己在天地宇宙之間其實是微不足道的,也難以為宇宙萬物設立好的標準,人們應該知恥近乎勇,學習發掘創造的美好,欣賞創造的美好,致力讓創造的美好得以被體現.



\newpage

\section{時間的旅人盼望之旅}
\label{sec:250}
\textbf{陳劍小姐 - 時間的旅人--盼望之旅}
\newline
\newline
\hyperref[sec:249]{< < < PREV SERMON < < <}
~
\hyperref[sec:toc]{[返目錄]}
~
\hyperref[sec:251]{> > > NEXT SERMON > > >}


講題:時間的旅人--盼望之旅　The Traveler of Time - For the Hope

講員:陳劍小姐

場合:香港中文大學崇基學院禮拜堂

日期:2014年6月22日

親愛的弟兄姊妹,平安.今天是我人生第一次在主日崇拜講道,願上帝祝福講的聽的都一同得造就.讓我們一起來做一個簡短的禱告.

親愛的三位一體的上帝,因您的臨在,而使我們的聚集有意義；因您親自賜下話語,好使這講臺成為您自己的講臺；因您體恤我們所有人的需要,好使我們都在您的裡面得到餵養和飽足.奉耶穌基督的名,阿們.

我現在是非常緊張的.但是我這個緊張的心情和保羅當年拿著外邦信徒的奉獻踏上耶路撒冷之旅的緊張是不一樣的.今天,我就用保羅的耶路撒冷盼望之旅,來跟大家分享這個主題--盼望.

我們主要會看三段經文:羅馬書的15章,使徒行傳21章,以及馬太福音第10章.

首先我會跟大家一起看看當時出於什麼樣的處境,保羅作為一個普通人,何不獨善其身就好,為何不禱告求上帝挪去一切艱難險阻,反而要選擇“明知山有虎,偏向虎山行”?其次,我會跟大家一起看看,原來保羅走上這條艱難的道路是經過深思熟慮,是充滿勇氣和盼望的.接著我們會從保羅的經歷中明白,真正的盼望有兩個要點:一,真盼望是來自與神真正相愛的關係；二,擁有這種真正的盼望,才能夠活出真正的平安.

首先我們要知道,保羅不只是一個被高度符號化了的聖經人物,他更是一個活生生的真人.就好像你現在身邊的弟兄姊妹一樣的真實.從脾氣性格上來說,他並不比我們高尚到哪裡去,從生活境遇上來說,他跟我們一樣有每日煩惱.然而,在他人生的一個關鍵的時刻,他做了一個非同一般的選擇.

1,保羅的困境:

保羅用許多的時間和精力,費盡各種周折,好不容易收集到了一筆外邦信徒捐助的款項,打算要親自拿去給耶路撒冷的教會.這事情記錄在羅馬書15:25-28,還有使徒行傳21章裡.關於這件事有幾個背景要交代:

1),大家可以猜一下保羅拿著這筆錢,是以什麼心態去耶路撒冷的?是得意洋洋的,還是忐忑不安的?

其實保羅所收的奉獻有限,遠低於他的預想.保羅當初奮力開拓的外邦基督徒的疆土在當時已經遭到破壞,比如在加拉太書裡面提到的加拉太衝突等,所以羅馬書15章寫捐獻教會的時候,沒有提到哥林多和加拉太的名字.所以保羅並不是以一個成功的外邦人使徒的姿態去耶路撒冷的,相反,他的處境其實相當可憐.可以說他根本沒有談判的資本.

2),同時保羅知道在羅馬已經流傳著關於他所傳福音不純的流言.他明白,雖然他決心將西班牙當作新的福音使命的目標,但是他很有可能不能夠獲得羅馬信徒群體全力的支持.這也是他很擔憂的一件事.可謂是內憂外患.

3),最後,根據聖經提供的資料,保羅真的在耶路撒冷被抓了,大家還記不記得是出於什麼原因?對,是因為他帶著外邦人進了聖殿,這可是觸犯了猶太人的天條,所以他就被抓了,還被押送回羅馬去受審坐監.

在這樣一種局面下,保羅面臨二個選擇:A繼續他的獨立福音使命.但這個策略已經被證明是行不通的,雅各幫派的人以及反對保羅的猶太信徒破壞了保羅幾十年來苦心經營的外邦教會福音疆域；B 承認他策略的失敗.讓他發展出的外邦信徒轉而服從摩西律法,也就意味著,眼睜睜看著他從福音侍工一開始所為之奮鬥的一切前功盡棄.

這兩個選擇,都不符合保羅一直以來為之奮鬥的信仰.因為在保羅看來,“合一”才是基督徒最應該追求的價值.保羅堅持:福音大能帶來的合一,是對任何隔閡的打破,不管是種族的,性別的,階級的,文化的任何一種隔閡,合一就是要突破這些隔閡.所以,福音怎麼可以是猶太基督徒的私有財產呢?外邦信徒又何必要行律法成為猶太人之後,才可以成為基督徒呢?

因此保羅一定要把這筆外邦信徒捐獻的款項帶到耶路撒冷,唯有這樣,才有可能建立一個外邦教會和中心教會彼此合一的新關係,但究竟是否可以如此成就,保羅也不知道,他是憑著對主耶穌的愛,來踏上這凶多吉少的旅程的.有學者認為羅馬書是保羅設想在耶路撒冷可能會的迫害,而為自己預先寫的辯護詞,在羅馬書中他飽含深情地說到:“誰能使我們與基督的愛隔絕呢?是患難嗎?是困苦嗎?是迫害嗎?是饑餓嗎?是赤身露體嗎?是危險嗎?是刀劍嗎?”

所以路加在使徒行傳中記錄到大家紛紛勸保羅不要前往耶路撒冷,保羅是怎麼回答的呢?讓我們一起來看使徒行傳21:13:“你們為甚麼哭,使我心碎呢?我為主耶穌的名,不但被捆綁,就算死在耶路撒冷我也都準備好了.”

這樣一段話不禁讓我想起著名文學家魯迅的名言:“真的勇士,敢於直面慘澹的人生,敢於正視淋漓的鮮血.”毫無疑問,保羅是真的勇士,他的勇氣來自於“我為主耶穌的名”這句話,他盼望的是主耶穌要在他身上成就的心意.正因為此盼望,兇險的耶路撒冷之旅成為了保羅的盼望之旅.

真正的盼望

然而是不是所有基督徒都像保羅這樣地盼望呢?

在弟兄姊妹的交流中,我們常常聽到許多見證,有的是疾病得醫治,有的是孩子考上理想的學校,有的是工作得到升遷,有的是人際關係得到好轉,等等,這些都是值得高興和感恩的事；但是我們也會聽到教會裡有一些弟兄姊妹說,他們覺得上帝不聽自己的禱告,因而他們覺得上帝不愛自己；或者,他們覺得人生充滿各種難題,於是懷疑上帝到底有沒有在掌權,有沒有在帶領.上帝有沒有按照我們所求的賜給我們,事情有沒有按照我們所希望的那樣發展,這些已經在不知不覺間變成了我們人生的盼望.這其實是很可怕的,因為我們可能會不停地在失望又希望,希望又失望之間打轉,最後突然遇到一件特別困難的事,然後就陷入絕望.

這就像舊約的以色列民一樣,原本是跟迦南地周圍的民族去比誰的神打仗更厲害,慢慢適應了農耕畜牧業文化後,發現迦南地信奉的土地神畜牧業神好像蠻靈的,於是他們放棄了耶和華神,而去崇拜當地的各種偶像,以祈求五穀豐登,有個好收成.結果呢,他們被擄失散,在他們心裡沒有平安,只有無盡的概歎.這是因為他們沒有在愛當中與上帝連接,再多的宗教禮儀又有什麼用呢?上帝在以賽亞書當中痛心疾首地說:“天哪,要聽!地啊,側耳而聽!因為耶和華說:‘我養育兒女,將他們養大,他們竟悖逆我.牛認識主人,驢認識主人的槽；以色列卻不認識我,我的子民不明白我.’”

上帝是有位格的上帝,是主體性的存在,與上帝建立彼此相愛的關係,才是真盼望的根源.耶穌也說,誡命中最大的一條就是:“盡心,盡性,盡意,盡力,愛主你的神.”這正是我們要說的第一點:盼望,正是建立在這個字--“愛”上面.

1,盼望是來自與神的相愛.

有一首詩歌叫做《我知誰掌管明天》,不知道有沒有弟兄姊妹會唱?這首詩歌非常美麗,它唱到:“我不知明天將如何,每一天只為主活,有許多未來的事情,我現在不能識透,但我知誰掌管明天,我也知誰牽我手.”

這首詩歌是一位團契的弟兄介紹的,他是一個有趣的人.他最害怕的就是在眾人面前禱告,別人唱讚美詩時流眼淚他問你們為什麼要哭?他在團契默默服侍了差不多7年,都只是做一些後勤的工作,並沒有所謂成長成一名精兵.他的生活也跟你我的生活一樣,並非一帆風順,有一次他家族企業裡面的一個事故,讓他一直到現在都仍然有明顯的創傷後遺症.他的信仰沒有大道理,沒有口號,你叫他講聖經故事,他可能會憋紅了臉,最後可能說我還是去買生果給大家吃,就趕快跑掉.就是這樣的一個弟兄,有一次在一個聚會上,大家問彼此最喜歡的讚美詩是什麼.他說是這一首.原因是什麼呢?他說:我不懂得講為什麼,但是我就是知道,上帝是我的阿爸父,祂掌管我的一切.我不知道明天會怎樣,但是我知道誰掌管明天!

這位弟兄的盼望,不在於他人生的現狀如何,也不在於他是否心想事成,更不在於他覺得他找到了一個天靈地靈世上最靈的神,而是在於他明確地知道,自己是在一段關係裡,就是與上帝相愛的關係裡,他對上帝的愛的理解,不是概念式的,而是像孩子一樣單純地切身經歷的.你問一個小孩子,為什麼你愛你爸爸?他可能沒有辦法給你總結出十個中心三條理論,但是他會告訴你,如果爸爸帶他去哪裡,不需要問,他也會乖乖牽著爸爸的手一起去.再用兩個人拍拖的比喻來說,真愛是不需要每分鐘說一百遍我愛你,再從一百個角度分析和解釋為什麼愛,也不需要跑遍整個香港島為了買一份對方最愛吃的甜點來證明你的愛.兩個人只需要想到彼此,就知道愛是真實存在,真實得超過一切眼睛可以看見,手可以觸摸的真實.盼望正是來自於這種真實的愛.

然而你可能會問,我是真的很愛主,我也跟保羅一樣凡事盼望著主耶穌的旨意成就,並深深相信主耶穌的旨意就是最好的.那麼我會得到什麼?

2,真盼望才能帶來真平安.

我告訴你,你會得到真正的平安.

這就是關於盼望的第二點,真盼望才能帶來真平安.

在今天宣讀的經文中,馬太福音10章的內容,是主耶穌關於“discipleship”門徒職份的教導.主耶穌先是告訴他們傳道的注意事項,和可能會遇到的各種患難危險,但是叫他們不要害怕,他們的一切都在天父的看顧之中,最後,主耶穌應許為他失喪生命的,將要得著生命.我們一起來讀10:38-39:“不背著他的十字架跟從我的,也不配做我的門徒.得著生命的,將要失喪生命；為我失喪生命的,將要得著生命.”這生命不是一般的生命,乃是豐盛的生命,平安的生命,永遠的生命.

保羅的耶路撒冷之行,仿佛是這段教導的一個活生生的例子,面對許多困難甚至最後被抓被拉去坐監,保羅仍顯得那麼心有平安如江河,還在法庭上坦然為自己的信仰來辯護.正如他之前想到自己作為主的門徒會受到許多磨難時,他說到:“但靠著愛我們的那一位,我們在這一切事上就得勝有餘了.”保羅所謂的得勝,是他終於靠著禱告成功去除掉一切的患難了嗎?不,而是他得到了主耶穌基督在愛裡面所賜的真正的平安.就像著名神學家莫特曼說的,真正的生命的健康,不是在於除去生命中一切的苦難,而是在於你有力量與這些苦難共舞.約翰福音中記錄了耶穌對於門徒的一個重要的應許,約翰福音14:27:“我留下平安給你們,我把自己的平安賜給你們；我給你們的,不像世界所給的.你們心裡不要難過,也不要恐懼.”

這就是主所賜的真平安.這平安在於耶穌基督同門徒之間愛的關係,因而門徒可以放心將一切交托給主,就好像剛才我們說的,小孩子對愛自己的父親,完全的信任和交托；真愛的兩個人之間,對彼此完全的信任和依賴,然而上帝的愛還要超過以上這兩個比喻的愛,因為上帝是用永恆的大能的慈愛來愛著我們,上帝說:就算你父母遺棄你,我也不離開你.這才是值得我們去盼望的平安.因而是晴是雨,是患難還是順境,是讚譽還是詆毀,是白天還是黑夜,無論什麼樣的生活遭遇,都不會動搖我們的勇氣,力量和信仰.保羅,當他拿著那註定被人拒絕被人漠視的奉獻款項,走向註定要被人傷害的耶路撒冷之旅時,他不是一個隨著命運漂泊的可憐人,也不是因為境遇而無法選擇自己的命運,相反,他滿懷著盼望,心有平安如江河,迎向主耶穌的榮耀.

正如保羅在羅馬書5章中說的:“就是在患難中也是歡歡喜喜的.因為知道患難生忍耐,忍耐生老練,老練生盼望,盼望不至於羞恥,因為所賜給我們的聖靈將神的愛澆灌在我們心裡.”願我們大家在各自的生活中,都可以活出對上帝的愛,從而在這份安穩的愛中產生盼望,最後獲得真正的平安.



\newpage

\section{讓愛與和平佔領人間世}
\label{sec:251}
\textbf{鄭漢文博士 - 讓愛與和平佔領人間世}
\newline
\newline
\hyperref[sec:250]{< < < PREV SERMON < < <}
~
\hyperref[sec:toc]{[返目錄]}
~
\hyperref[sec:252]{> > > NEXT SERMON > > >}


羅馬書 6:12-23
\newline
\begin{longtable}{cl}
\hline
\hline
章節 & 經文 (和合本修訂版)\\
\hline
6:12 & \begin{tabularx}{0.7\textwidth}{X} 所以，不要讓罪在你們必死的身上掌權，使你們順從身體的私慾。 \end{tabularx} \\ \\ \relax
6:13 & \begin{tabularx}{0.7\textwidth}{X} 也不要把你們的肢體獻給罪作不義的工具，倒要像從死人中活著的人，把自己獻給神，並把你們的肢體獻給神作義的工具。 \end{tabularx} \\ \\ \relax
6:14 & \begin{tabularx}{0.7\textwidth}{X} 罪必不能作你們的主，因你們不在律法之下，而是在恩典之下。 \end{tabularx} \\ \\ \relax
6:15 & \begin{tabularx}{0.7\textwidth}{X} 那又怎麼樣呢？我們在恩典之下，不在律法之下，就可以犯罪嗎？絕對不可！ \end{tabularx} \\ \\ \relax
6:16 & \begin{tabularx}{0.7\textwidth}{X} 難道你們不知道，你們獻自己作奴僕，順從誰就作誰的奴僕嗎？或作罪的奴隸，以至於死；或作順服的奴僕，以至於成義。 \end{tabularx} \\ \\ \relax
6:17 & \begin{tabularx}{0.7\textwidth}{X} 感謝神！因為你們從前雖然作罪的奴隸，現在卻從心裡順服了所傳給你們教導的典範。 \end{tabularx} \\ \\ \relax
6:18 & \begin{tabularx}{0.7\textwidth}{X} 你們既從罪裡得了釋放，就作了義的奴僕。 \end{tabularx} \\ \\ \relax
6:19 & \begin{tabularx}{0.7\textwidth}{X} 我因你們肉體的軟弱，就以人的觀點來說。你們從前怎樣把肢體獻給不潔不法作奴隸，以至於不法；現在也要照樣將肢體獻給義作奴僕，以至於成聖。 \end{tabularx} \\ \\ \relax
6:20 & \begin{tabularx}{0.7\textwidth}{X} 因為你們作罪的奴隸時，不被義所約束。 \end{tabularx} \\ \\ \relax
6:21 & \begin{tabularx}{0.7\textwidth}{X} 那麼，你們現在所看為羞恥的事，當時有甚麼果子呢？那些事的結局就是死。 \end{tabularx} \\ \\ \relax
6:22 & \begin{tabularx}{0.7\textwidth}{X} 但如今，你們既從罪裡得了釋放，作了神的奴僕，就結出果子，以至於成聖，那結局就是永生。 \end{tabularx} \\ \\ \relax
6:23 & \begin{tabularx}{0.7\textwidth}{X} 因為罪的工價乃是死；惟有神的恩賜，在我們的主基督耶穌裡，乃是永生。 \end{tabularx} \\ \\
[1ex]
\hline
\hline
\end{longtable}
講題:讓愛與和平佔領人間世　Let Love and Peace Occupy the World

經文:羅馬書6章12-23節

講員:鄭漢文(rogercheng@cuhk.edu.hk)

場合:香港中文大學崇基學院禮拜堂

日期: 2014年6月29日

1. 「護體膏」與「解靈氣」的故事

弟兄姊妹,平安.

今天,我又再回鄉探親一樣,返回崇基禮拜堂這個信仰大家庭.

我的家人學生和好友,今天都共聚在此,盼望聖道的臨在.

華人社會非常喜歡武俠小說,現在首先讓我分享一個武俠小說中的故事.

正所謂功夫深道理真.話說有個江湖傳聞,存在兩套令得生命可以「超勁」(起死回生)的武學妙方,功夫(a martial art of practice):一種叫「護體膏」(Body Protection Lotion),另一種叫「解靈氣」(Spiritual Renewal Respirator).

「護體膏」是一種神奇的膏藥,聞說是要從六種死物加上六種毒液經六六三十六個月才煉成.早晚搽在身上,走近身邊的別人的生命能量就會被吸收過來,彼消我長,別人生命力會消退,而搽此護體膏的我會生命力加強.稍有小成,變成生命力增長的手法,成為人間世上讓其他人敬畏而遠之或取之而後快的神功.

「解靈氣」倒是另一種修煉(spiritual practice),聞說是要早午晚三個時段,吸收日月精華:一是就是在日出東方的時分向天上的光芒靜靜祈禱一分鐘,二是在日正當空的時分向一個有需要的人子提供幫助,三是在日落西方的時分當著大地吹過的風深深吸入呼出吐納三回.

解靈氣會把每天沾在身上的毒素化解成生命能量變為靈氣,更奇妙的是,你的解靈氣會讓與你親近的人分享到你的化解毒素而成為靈氣.稍有小成,人與人的心漸漸相連,心連心相連於天地之間.

究竟一個搽滿「護體膏」的武夫與一個煉就「解靈氣」的俠士相遇,二人如要交手,交手之下,結果會是如何?

究竟會是前者把後者的生命力全吸走而成為走肉行屍,還是後者化走前者體上所有的毒而成為一個有靈的活人?

人生在世,難為正邪定分界.二者擇一,正邪之間,你會選擇哪一種?

個體上從罪的奴僕解放為義的奴僕:羅馬書6:12-23解經

今天講道,我們會用一二三的結構,我寫了一首打油詩來表示:「一段經文總體讀,兩條線索交織解,三面體制同反思,實踐理解你我他.」

今日經課是羅馬書第六章第12至23節.讓我們一同誦讀這段經文.

請打開聖經到《羅馬書》第六章第12節讀到第23節,讓我建議一種三把聲的輪讀一段經文,這種三一讀法是:我讀第12/15/18/21節,姊妹讀在我讀完後開聲讀13/16/19/22節,弟兄在姊妹讀完後開聲讀14/17/20/23節.讓我試一試效果.

(6: 12) 所以,不要容罪在你們必死的身上作王,使你們順從身子的私慾.

(6: 13) 也不要將你們的肢體獻給罪作不義的器具；倒要像從死裡復活的人,將自己獻給神,並將肢體作義的器具獻給神.

(6: 14) 罪必不能作你們的主,因你們不在律法之下,乃在恩典之下.

(6: 15) 這卻怎麼樣呢?我們在恩典之下,不在律法之下,就可以犯罪嗎?斷乎不可!

(6: 16) 豈不曉得你們獻上自己作奴僕,順從誰,就作誰的奴僕嗎?或作罪的奴僕,以至於死；或作順命的奴僕,以至稱義.

(6: 17) 感謝神!因為你們從前雖然作罪的奴僕,現今卻從心裡順服了所傳給你們道理的模範.

(6: 18) 你們既從罪裡得了釋放,就作了義的奴僕.

(6: 19) 我因你們肉體的軟弱,就照人的常話對你們說.你們從前怎樣將肢體獻給不潔不法作奴僕,以至於不法；現今也要照樣將肢體獻給義作奴僕,以至於成聖.

(6: 20) 因為你們作罪之奴僕的時候,就不被義約束了.

(6: 21) 你們現今所看為羞恥的事,當日有什麼果子呢?那些事的結局就是死.

(6: 22) 但現今,你們既從罪裡得了釋放,作了神的奴僕,就有成聖的果子,那結局就是永生.

(6: 23) 因為罪的工價乃是死；惟有神的恩賜,在我們的主基督耶穌裡,乃是永生.

這段經文其實由兩條線索所交織而成,保羅是在對比兩種生命的景況,簡單地分別是:

S線:罪的奴僕(servant to sin)；和

R線:義的奴僕(servant to righteousness).

借助這S和R兩條線索好像照明燈(torch),現在我們重新閱讀這段經文,讓我們引用這樣的一種從字裡行間(highlighting)顯影的技巧:

先把其中一條線索勾出來,突顯這條線索所呈現的生命景況,

然後重施故技,再把另外一條線索勾出來,突顯另一條線索所呈現的另一種生命景況,

再而比較一下,你會你會選擇哪一種?

首先是S線:個體肉身成為罪的奴僕的光景(predicament of servanthood to sinfulness)──

(6: 12) 所以,不要容罪在你們必死的身上作王,使你們順從身子的私慾.

(6: 13) 也不要將你們的肢體獻給罪作不義的器具；

(6: 14) 若你們還在律法之下罪,必能作你們的主.

(6: 16) 豈不曉得你們獻上自己作奴僕,順從誰,就作誰的奴僕嗎?或作罪的奴僕,以至於死；

(6: 17) 你們從前…作罪的奴僕,

(6: 19) 我因你們肉體的軟弱,就照人的常話對你們說.你們從前怎樣將肢體獻給不潔不法作奴僕,以至於不法；

(6: 20) 因為你們作罪之奴僕的時候,就不被義約束了.

(6: 21) 你們現今所看為羞恥的事,當日有什麼果子呢?那些事的結局就是死.

(6: 23) 因為罪的工價乃是死.

當我們還是罪的奴僕的時候,我們順從身子的私慾,我們的肢體獻給罪作不義的器具,肉體的軟弱令我們將肢體獻給不潔不法作奴僕,行不義,行羞恥事.結果呢,結局就是死.因為罪的工價乃是死.這是S線下的生命光景.

這樣的生命景況實在不是我們想望的.有另類的選擇嗎?有,保羅描繪給我們知道.

好在還有的是R線:個體肉身成為義的奴僕的境界(predicament of servanthood to righteousness)──

(6: 13) 倒要像從死裡復活的人,將自己獻給神,並將肢體作義的器具獻給神.

(6: 14) 罪必不能作你們的主,因你們在恩典之下.

(6: 15) 這卻怎麼樣呢?我們在恩典之下,不在律法之下,就可以犯罪嗎?斷乎不可!

(6: 16) 豈不曉得你們獻上自己作奴僕,順從誰,就作誰的奴僕嗎?或作順命的奴僕,以至稱義.

(6: 17) 感謝神!因為你們從前雖然作罪的奴僕,現今卻從心裡順服了所傳給你們道理的模範.

(6: 18) 你們既從罪裡得了釋放,就作了義的奴僕.

(6: 19) 我因你們肉體的軟弱,就照人的常話對你們說.…現今也要照樣將肢體獻給義作奴僕,以至於成聖.

(6: 22) 但現今,你們既從罪裡得了釋放,作了神的奴僕,就有成聖的果子,那結局就是永生.

(6: 23) 惟有神的恩賜,在我們的主基督耶穌裡,乃是永生.

因為在主基督耶穌裡,因為神的恩典,我們從罪裡得了釋放,就作了義的奴僕.我們像從死裡復活的人,將自己獻給神,並將肢體作義的器具獻給神.作了神的奴僕,就有成聖的果子,那結局就是永生.這是R線下的生命光景.

我們S線走向R線,從罪的奴僕(servant to sin)釋放出來,成為義的奴僕(servant to righteousness).由無法有行義的自由,成為可以行義的自由,不單行義,還可以成聖,得享永生.一來恢復神創造我們成為有靈的活人,而且靈氣可以逼人,讓其他人都可以把毒化解,不再受罪所捆綁.

我們信耶穌基督的救恩的,都明白以上的道理.

都我們活在人間世上,依然感受到被奴役,所捆綁,無法明白我們所還在經歷的苦況.

今天我想把以上的道理再向前再推進一步:把身體所負載的罪性由個體的層次推向群體的層面,探討我們如何面對罪的權勢(sinfulness as principalities and powers),以及「結構的罪」(structural sin)下我們都變成了「被罪者」(sinned against).

3. 群體上從罪的奴僕解放為義的奴僕:經濟體制內出現結構性的惡也不羞恥

保羅在羅馬書提供了一個理解這個人間世上有惡的神學觀點.

例如,在羅馬書第五章16節提到:罪入了世界,罪佔領了人間世.我們活在人間世界上,我們個別的身體成為群體的部份.

罪佔領了的人間世化為經濟,政治與社會的軀體,毒素在蠶食我們的體制,出現不同的體制出現不公義欠仁愛而不感羞恥.我們如何從取死的體制中解放出來,把我們在人間世上的生命變化成獻給神的器皿?

羅馬書6章21節提到:罪佔領了人間世,人在其中,會行羞恥的事,會行惡事而不感到羞恥.我們現在活在人間,世上有何現象是理該令人感到羞恥,但卻不斷令其出現的人卻不感到羞恥?罪佔領了人間世,在其中形成更大的身體,叫做體制,這個身體可以令在其中的人異化,會行可恥的事,會行惡事而不感到羞恥..

先說經濟所出現的體制不公義,量度貧富懸殊的堅尼斯系數,香港的已超過0.5,按過去的研究,超過了0.4的社會其實已經會出現革命的了.

在香港,極少數的有錢人掌握大多數的財富,跡近一成人擁有百分之九十的財富,餘下的九成人只享有一成的財富,這公義嗎?他們不感到可恥的嗎?而這群掌有大量生產資料的資本家,不斷掌控生產關係,對勞動階層不斷剝削,結果是貧者越貧,富者越富.這樣公義嗎?

近年通過的最低工資,只是杯水車薪.

最近反東北撥款的抗爭,背後是立法議會出現體制不公義被懷疑當中有官商勾結.當中或有掌握地權者,通過發展自己有份的土地或者有群帶關係的業權人,因發展而得利益互相施肥.

犧牲享有自己生活方式的住民等等利益.犧牲香港擁有可以自足的農田耕作土地與文化.

作為土地正義聯盟執委以及東北三區村民組織支援者的朱凱迪在〈東北613行動的前因後果〉這樣解說:

「議員和村民互相支撐的那幾次財委會,不應稱為「拉布」,因為保皇黨提倡的所謂「有效率的審議」,實際上就是不去挑戰議會的不平等結構,官員和議員跟議案之間的利益衝突,以及不去深究議案本身的細節.

這幾次財委會的審議裡,公眾可以知道,雖然政府不斷吹噓新界東北是解決未來住屋問題的關鍵,可是公營房屋佔地僅三十六公頃,即整個發展區範圍的6\%；公眾可以知道,政府官員在立法會內公然說謊,否認曾經進行東北區內的業權調查,反而更加清楚,政府設定持有四千平方米以上的地主可以就地發展的政策,是諮詢地產業界的委員會後得出的結果,證實了透過發展向地產商輸送利益的指控；公眾可以知道,當計劃徹底破壞原有社區和環境時,政府原來沒有區內長者數目的資料,所以根本不打算照顧他們的特殊需要,沒有想好令失去農地的農友得以復耕的方案,連村民的「公屋原區安置」也是假的,因為當第一期公屋落成之日,村民早已全部被逼遷四散.

另外,面對香港這個半民主半威權的政體,每次在立法會的認真審議都必然上升至對議會不民主體制及制度暴力的醒悟.比起四年前反高鐵,香港特區政府和立法會的狀況如今更加不堪.」

大家對是否值得棄守東北可以不同的觀點.但周五立法議會下財委會這樣通過首期發展撥款的方式,根本就是議會暴政,大家不覺得羞恥的嗎?

因為有「護體膏」早晚搽在身上的資本家和政客,不會感到羞恥的.煉得越上乘,把走近身邊的別人的生命能量就會被吸收過來,彼消我長,別人生命力會消退,而搽此護體膏的人會生命力加強.就算香港整個大東北,他們都可以吸為己有,毫不羞恥.

民資民膏,盡情收挖.「護膏」神功,走火入魔.

保羅在羅馬書第六章再三提醒:當我們還是罪的奴僕的時候,我們順從群體中掌權者的私慾,我們的群體獻給罪作不義的器具,個別肉體的軟弱令我們將群體獻給不潔不法作奴僕,行不義,行羞恥事.結果呢,結局就是死.因為罪的工價乃是死.

結構的罪的工價也是死,在結構的罪下受苦的我們是:我們做到死,生不如死.

我們都是「被罪者(sinned against)」.

4. 讓愛與和平佔領人間世:化除政治體制中軀體的毒素得以釋放為自主人

罪佔領了的人間世,化為經濟,政治與社會的軀體,毒素在蠶食我們的體制,出現不同的體制不公義,政策欠關愛,關係難溝通.

罪佔領了人間世,在其中形成更大的身體,叫做體制,這個身體可以令在其中的人異化,繼續主宰人.經濟體制不公義,被包庇在政治體制的不公義.利益集團一直掌權,人民民不聊生.現在讓我們進入「三面體制同反思」的第二項反思:政治體制.

經濟以外,香港人核心價值還包括:全民普選行政首長和立法議員的全面民主化.

大家一定想起我今天會談及:「讓愛與和平佔領中環」這個運動.

是的,看今天講道的題目就會預期到,但談的是另一個層次,就是從信仰角度解讀心靈佔領.今年3月30日我在家父剛離開人間世後,在這裡講道,題目為信仰讓生命不再一樣,曾經提及以下一段話:

「究竟雖千萬人吾往矣的「佔中三子」信仰甚麼?朱耀明牧師當然相信上帝,並且願意一生成為上帝的僕人,牧養迷失的羊群.戴耀廷教授也相信上帝,相信公義臨到人間,自身甘願成為群羊之首.

雖然長年在基督教信群中生活,陳健民教授不承認自己相信基督宗教.

當這個學期有次他在這崇基禮拜堂作週會講者時,有人問到:「若然有關權貴要傷害你時,你還會相信甚麼?」陳健民教授答:「我會回頭,望向這個十字架.」

請大家一同,望向崇基禮拜堂這個十字架,對你來說,代表了甚麼?

對你來說,十字架,代表了甚麼?」

是否代表:讓愛與和平佔領人間世?

是的,罪佔領了的人間世,我們要反佔領.用甚麼來反佔領?

罪佔領了的人間世,化為更大的身體,名字叫體制,不斷產生異化力量,形成了「結構的罪」(structural sin),讓體制下的你我再次成為罪的奴隸,好人都會變成惡人,惡人變成大惡人.我們都成了「被罪者(sinned against)」,在體制下受苦.所以,我們無法不在體制內反佔領.

讓我們耶穌基督給我們的愛與和平,來反佔領我們所身在的人間世.

羅馬書6章17-18節提到:從前我們是罪的奴僕,如今釋放了,從心裡順服了所傳給你道理的模範.道理就是讓我們從罪的奴役中得到釋放出來,藉著聖靈更新我們的生命而讓我們漸漸肉身成道,把耶穌基督的愛活現在人間世,作為上帝與人復和的和平之子.

近日活在近似專制的暴政下,大家感到是被宰制被綑綁,還是自由自在?

有人說:權力腐化人,絕對的權力絕對腐化人.(Power corrupts; absolute power absolutely corrupt)

罪佔領了人間世,其中形成更大的身體,叫做體制,這個身體可以令在其中的人異化,順從身子的私慾,作不義的器皿.

一個不是人民選出來的政府,會因其中份子的私慾,行不仁不義的惡事.所以為何要行政長官與立法議員雙普選?因為被選出來的行政的人民公僕和立法的人民代表,若有因私慾而行羞恥的事,可以在下次改選時被人民撤換.

我們身體需要定期檢查,不好的身體部份需要維修或撤換.機器如是,政治體制也如是,普選就是定期檢查,改選就是定期維修,更換組件,非常正常.

拒絕普選的人,反而不正常,若是在體制內的當權者就是因為謀私(而不願放權),在體制外的可投票者就是因為(思想)還未得到釋放(而未能自主).

面對這個時代,我又如何?

約書亞的一句話,就是我的回應:「至於我和和家,定必事奉耶和華.」

每次我回來這裡,總有教友問我關於和平佔中運動的事.我最近也被《時代論壇》(Christian Times)訪問了,問我作為一位基督徒如何回應佔中?我的看法是:這是一個讓中環價值不再佔領我們心靈的公民覺醒運動.這個訪問剛巧是刊登在今天的1400期的頭版(你們崇拜完後可以立即購買,先睹為快.)我在這裡不多談了.

今天是620至629公民投票的最後一天,七十多萬人已經投了,你呢?

有人說,這是佔中投票而非公民投票,這是(無法制力的)民間運動,而非政府法制行為.

你們說的是,但正正是政府不讓公民投票的真正自主力量得以運作,唯有用民間方式讓公民自主的力量得以釋放出來.

今天朝十到晚十的一整天,在香港中文大學的康本園增設了一個實票投票站.

不要單單聽完道,還要行道.

一路之近,可以非常遙遠,只因各有各的道,遙遠是人的心靈與心靈的距離.

人心距離,關係撕裂.

先讓讓愛與和平佔領自己所身在的人間世:仁愛關係重建的呈獻

罪佔領了的人間世,化為經濟,政治與社會的軀體,毒素在蠶食我們的體制,人心異化後,出現不同的體制不公義,政策欠關愛,關係難溝通.

剛才講過經驗與政治的體制異化,現在讓我們進入現「三面體制同反思」的第三項反思:罪入了人間世下的社會體制這一面.

在崇基學院神學院出任院長十九年而在今年七月底退休的盧龍光牧師,在2014年6 月1 日循道衛理聯合教會主日靈音稿這樣描述近年的中港關係以及社會撕裂:

「這廿五年,不少香港人從愛中國變得厭惡中國；青年人的本土意識高漲而排拒中國!在面對香港民主化進程受北京粗暴地干預的情況下,一些知識分子發起和平佔中運動!香港人和內地同胞追求更大的民主不但是世界的大勢所趨,更合乎普世價值和聖經所揭示的價值觀,使人的尊嚴和上帝所賦予的人權得以建立和保障；香港人以行動爭取民主是本身的責任和權利,支援內地同胞作同樣的追求是理所當然的.

香港本土意職的膨脹而排拒中國是妄顧客觀處境和自我為中心的思想,不符合上帝的福音,是要我們打破人類族群間的阻隔和對立.

香港人和內地人在生活習慣和文明程度的差距是社會發展的差異,香港人只不過是在不同的政治,社會,經濟發展下形成今日的香港文化,我們不應看不起同樣為上帝所創造的中國人.」

關係,總是由近至遠.近年的本土運動不是沒有國家的關懷,而是國家沒有守望我們本土的關懷.唯有自求多福,不再做主子的奴才.

人心距離,關係撕裂.我這近半年感受得最深.

知我近況的友好,都關心我近日被捲入自己工作間的體制不公義抗爭運動之中,成了指控其中數項體制不公義政策欠關愛的吹哨者(whistle blower).

正如我的講道學老師教落:不要把講台變成炮台.我在這裡不便談及細節及當中人物.關心我和我所在的體制的,總會有方法得知詳情.請為我代祈禱.

簡單地說:我今年三月20日父親過世後一直因多事困纏而身心抱恙,我還自討苦吃:打從六月四日開始,寫了一連串公開信向院長及院方發出「體制不公義,政策欠關愛」(Unjust Institution, Uncaring Policies)的四項抗議.這次我稱為「六月風暴」的結果是:事功未成,我自己早吃了苦果──經歷了種種反噬及人情冷暖──弄得自己身心社靈受創,大傷元氣(還可能要面對法律及大學的其他追究),甚需要下半年恢復安康.

我需要復修,需要的正是江湖傳聞的「解靈氣」.

按我們的信仰,甚麼是我們的「解靈氣」?

大家從今天講道主題大概都會猜到答案了.對,就是仁愛與和平.「解靈氣」就是我們把活在上帝的仁愛,以及因耶穌基督所得的恩典而有的和平,成為化解體制內的毒素的行者(agency),讓體制靈氣重新活現.「解靈氣」會把每天沾在身上的毒素化解成生命能量變為靈氣,更奇妙的是,你的「解靈氣」會讓與你親近的人分享到你的化解毒素而成為靈氣.稍有小成,人與人的心漸漸相連,心連心相連於天地之間.

讓我們在經濟體制,政治體制以及社會體制之中,化解當中的罪的奴役,成為公義的體現,成就仁愛的關係,讓愛與和平佔領人間世.在不得已之時,也許體制不公義,政策缺關愛是需要抗議,需要抗爭的.但依然可以用和平的方式抗爭,在抗爭路上建構仁愛的關係.

我們還活在世上有何意義?正如羅馬書6章13節提醒:讓我們肉身成道,把我們在人間世上的生命變化成獻給神的器皿,變成神在人間世上的「解靈氣」,成為一份呈獻(dedication).

讓我們一起低頭禱告:

「親愛的上主啊,

求你賜我們「解靈氣」,好像你在以色列出埃及離開奴役但還流落在曠野中的嗎哪(manna).

讓我們學會這樣的一種屬靈修煉(spiritual practice),讓我們早午晚三個時段都能從你而得著活在人間世的精神力量.

讓我們每天日出東方的時分,向天上的光芒靜靜祈禱一分鐘,讓聖父爸爸你差遣我們每一天的存活在人間世上.

讓我們每天日正當空的時分,向人間世上有需要的人子提供幫助,讓聖子耶穌基督你的憐憫臨在.

讓我們每天日落西方的時分,向著大地吹過的風深深吸入呼出吐納三回,讓聖靈你清理我們罪入人間世所染上的雜質,清洗我們的心靈.

但願上主你能給我們的「解靈氣」,能把每天沾在我們身上和我們體制上的毒素化成生命能量,還更奇妙地讓親近我們的人分享當中的化解毒素成為靈氣.

求你變化我們的生命成為你在人間世上的「解靈氣」,讓我們每天存活都讓愛與和平重新佔領人間世,讓體制內的不仁不義漸漸消除,讓人人在體制內不再異化,成為有靈的活人.」

但願榮耀都歸給聖父,聖子,聖靈,永在的三一神.阿們.

(鄭漢文博士現職香港中文大學教育行政與政策學系專業顧問,

同時在崇基學院神學院兼讀神道學學士課程.

曾於2012年1月至6月及2012年9月至2013年4月於崇基禮拜堂作校牧實習.)



\newpage

\section{國度國土國民}
\label{sec:252}
\textbf{伍渭文牧師 - 國度,國土,國民}
\newline
\newline
\hyperref[sec:251]{< < < PREV SERMON < < <}
~
\hyperref[sec:toc]{[返目錄]}
~
\hyperref[sec:253]{> > > NEXT SERMON > > >}


講題:國度,國土,國民　The Kingdom, the Land, and the People, Psalm

經文:詩篇145

講員:伍渭文牧師

場合:香港中文大學崇基學院禮拜堂

日期:2014年7月6日

七一在烈日,大雨下,遊行渲洩對政府不滿的人數,竟然在歷年前列,難怪陳佐耳先生說:領土回歸十七年,但人心尚未回歸.不論是本地或國際媒體,今日講道題目「國度,國土,國民」都是最近沸沸揚揚的頭條新聞,像東海,南海島嶼主權爭議,伊拉克遜尼派反政府武裝力量宣稱已建立「黎凡特伊斯蘭國」(Levant Islam State).

今天是聖靈降臨後第四主日,根據教會年(Church Year)的設計,上半年關於基督生平:將臨期(Advent),聖誕期(Christmastide),顯現期(Epiphany),預苦期(Lent),聖周(Holy Week),復活期(Eastertide),升天日(Ascension),直到聖靈降臨日(Pentecost).聖靈降臨後開始有關基督徒的下半年,稱為常年期(ordinary time).耶穌復活和門徒顯現四十天升天日,談論有關天國的事情,並叫門徒在耶路撒冷等候聖靈降臨,得著能力,把天國的福音傳遍世界.今天的經課就談論神的國度的建立,國民的使命和國土的擴展.

認識詩篇.今天的講道的經文選自詩篇145: 8-14節.我們要承認,一般信徒對詩篇印象並不深刻,除了幾篇著名的詩篇如23篇「耶和華是我的牧者,我必不至缺乏」,19篇「諸天述說神的榮耀,窮蒼傳揚他的手段」, 當然大部份人都曉得最長的詩篇是119篇.我們知道詩篇是讚美詩和禱告,認為讚美禱告貴乎虔誠,忽略詩篇其實是公共崇拜禮儀中的讚美和禱告,跟私人隨意禱告讚美不同；每篇詩篇都經過深思熟慮,合乎公共崇拜體統寫成的.在印刷術發明之先,我們是聆聽詩篇,沒有機會閱讀詩篇,為了容易記誦,詩篇採用了一些文學手法來演繹.

詩145就採用了離合詩體(Acrostic) ,每節第一個字依次採用二十二個希伯來文字母開始,第十三節用了兩個字母,共二十一節.全詩分四段,重覆兩個主題:上主偉大,vv.1-3,vv. 4-6；上主慈愛,vv.7-9.上主偉大,vv. 10-13；上主慈愛14-21.

第一段「上主偉大」,vv.1-3: 「我的神我的王啊,我要尊崇你!我要永永遠遠稱頌你的名!我要天天稱頌你,也要永永遠遠讚美你的名!耶和華本為大,該受大讚美；其大無法測度.」這段文字用了「尊崇你」,「本為大」,「該受大讚美」,「其大無法測度」描述上主的偉大.Vv. 4-6繼續上主偉大的主題,還加入了第三身來讚美:「這代要對那代頌讚你的作為,也要傳揚你的大能.」「人要傳說你可畏之事的能力」.

第二段「上主慈愛,」vv. 7-9: 「他們記念你的大恩就要傳出來,並要歌唱你的公義.耶和華有恩惠,有憐憫,不輕易發怒,大有慈愛.耶和華善待萬民；他的慈悲覆庇他一切所造的.」這段文字出自西乃山立約時的宣告:「耶和華在雲中降臨,和摩西一同站在那裡,宣告耶和華的名.耶和華,耶和華,是有憐憫有恩典的神,不輕易發怒,並有豐盛的慈愛揚給和誠實.」(出34:5,6) 這慈愛先從約民親身經歷,並傳揚給萬民.

第三段「上主偉大」,vv. 10-13: 「耶和華啊,你一切所造的都要稱謝你；你的聖民也要稱謝頌你,傳說你國的榮耀,談論你的大能,好叫世人知道你大能的作為,並你國度威嚴的榮耀.你的國是永遠的國!你執掌的權柄存到萬代!」第三段重覆上主偉大的主題:「談論你的大能」,「知道你大能的作為」.與第一段不同的,這裡出現了國度的觀念:「你的聖民也要稱謝頌你,傳說你國的榮耀……你的國是永遠的國.」

第四段「上主慈愛」,vv. 14-20:「凡跌倒的,耶和華將他們扶持；凡被壓下的,將他們扶起.萬民都舉目仰望你；你隨時給他們食物.你張手,使有生氣的都隨願飽足.耶和華在他一切所行的,無不公義；在他一切所做的都有慈愛.凡求告耶和華的,就是誠心求告他的,耶和華便與他們相近.敬畏他的,他必成就他們的心願,也必聽他們的呼求,拯救他們.耶和華保護一切愛他的人,卻要滅絕一切的惡人.」

這詩篇除了採用離合詩體,一個字母開始一節,還採用了平行體；四段兩主題,重覆上主偉大,上主慈愛,A, B ,A’, B’.今天讀經員讀出的詩篇經課就是四段中間的二,三段,而平行體中間的內容,也是全詩的中心思想.

在心中建立的慈愛和公義的國度vv. 7, 8

「他們記念你的大恩就要傳出來,並要歌唱你的公義.耶和華有恩惠,有憐憫,不輕易發怒,大有慈愛.」這兩節經文強調認信,信仰的表白:記念你的大恩,歌唱你的公義.神國沒有看得見的領土,祇有印證於心中的信仰.

甚麼是神聖的領土.構成國家有三個元素:主權,領土和人民.聖經的國度重點在上主的主權,當人的心被上主的慈愛和公義管治時,上帝的國就在我們心中.俗世的國度,依靠武力建立和維持,槍杆子出政權.此刻我們看到伊拉克遜尼派武裝力量,如何靠武力和殺戮建立「黎凡特伊斯蘭國」(Levant Islam State).強大的國家也須靠大量經費維持軍隊,保衛國土；甚至用超過國防經費,維持境內的穩定,保住政權.但上主的國在人的心中,其神聖不可侵犯的國土,就是順服上主的心.

以色列人曾經以為約櫃神聖不可侵犯,因裡面安放了兩塊法版,也以為聖殿神聖不可侵犯,因為這裡是敬拜上主的地方；有了約櫃和聖殿,上主一定祝福我們.但當他們拜偶像,心中遠離上主時,約櫃可以被擄走,聖殿可以被燒毀,祇有我們的心,才是上主眼中的神聖空間,聖靈在其中居住.

被憐憫征服的心.當我們的心被上主的慈愛和憐憫所征服,上主的國就出現在我們心中.一般提到上主國度,我們強調公義,但上帝國度建立在他對人的憐愛,第八節:「耶和華有恩惠,有憐憫,不輕易發怒,大有慈愛.」這是聖經特別用語,是上主和他子民在西乃山立約頒十誡時的宣告:「耶和華在雲中降臨,和摩西一同站在那裡,宣告耶和華的名.耶和華,耶和華,是有憐憫有恩典的神,不輕易發怒,並有豐盛的慈愛和誠實.」(出34:5,6)

這宣告是以色列人拜金牛犢失敗後,上主給予的第二次機會.摩西上西乃山,準預領受頒下的十誡,遲遲不見人下來,以色列人就要求亞倫讓他們造金牛犢代表上主,敬拜金牛犢.摩西下山聽到人聲喧嘩,看到百姓放肆跳舞,膜拜金牛犢,就把兩塊法版扔在山下摔碎了.拜偶像惹動上主的怒氣,凡不悔改的遭利未人殺戮,一天死了三千人.摩西求情:「這百姓犯了大罪,為自己造了金像.倘或你肯赦免他們的罪……不然,求你從你所寫的冊上塗抹我的名.」(出33 : 31, 32)上主就憐憫赦免百姓,給予第二次機會.摩西是基督的預表,耶穌是第二個摩西,登山垂示寶訓,像摩西從西乃山頒下律法,因為基督的原故,上主給以色列人第二塊法版.

公義和憐憫.公義和憐憫不同,公義是按應得,而憐憫按需要.上主按我們的需要,憐憫我們,不按我們應得的審判我們,憐憫原是向審判誇勝.「耶和華有恩惠,有憐憫,不輕易發怒,大有慈愛.」公義由愛衍生,因為愛每一個人,希望每個人都有同等機會,同等權利,分享到經濟發展的成果,我們追求更合理的社會資源分配.我們關心每一個人的健康,所以關心公共衛生分配到的資源是否足夠.

手段正確.所以我們若明白以愛和憐憫追求公義,就小心不可污名化,妖魔化對手達至社會公義的目的,因為我們也要愛我們的對手.舊約經課撒迦利亞書:「看哪,你的王來到你這裡……謙謙和和地騎著驢駒子.我必除掉以法連的戰車和耶路撒冷的戰馬；爭戰的弓也必除滅,他必向列國講和平.」(亞 9: 9, 10)

動機純良.書信經課羅馬書七章也提醒我們,追求公義時,小心自己的動機是否純良.「我覺得有個律,就是我願意為善的時候,便有惡與我同在.,因為按著我裡面的意思,我是喜歡神的律；但我覺得肢體有個律和我取心中的律交戰,把我擄去,叫我附從那肢體中犯罪的律.」(羅 7: 21-23)在追求公義時,我們要注意自己的動機,是否有無愧的良心和無偽的愛心.

要承受地土,非攻佔地土.上主國度的領土在我們心中,不在外在的土地,提醒我們溫柔的人,才能承受地土；真正的滿足在於心靈的充實,不在擁有更多和更多.有些人心靈貧乏,窮得祇剩下錢.非常諷刺,不少重視領土完整,強調神聖靈土一吋也不讓的國家,投放大量經費保護疆界,部署內部維穩工程,卻不重視已擁有的領土,任由土地長期環境污染,水土流失.同樣,現代人不斷擁張地界,攻佔地土,累積財富,卻不會去承受享自已已擁有的地土.為了見識見識,到處遊覽,卻對自己的家人,自己本身真正需要,很陌生.八福中的第三福-溫柔的人有福了,因為他們必承受地土.溫柔的人是謙柔知足的人,不被不止息的欲望所驅策,不斷爭取擁有更多又更多,沒有時間建立友誼,鞏固親情,沒有獨處默想,最後失去朋友,失去家人,也失去健康,也迷失了自我.

小結:神國就在我們心中,要讓上主的慈愛和憐憫佔領我們的內心,才有真正的滿足.我們的心在上主眼中,是最神聖的國土,因為在這裡,上主掌權.

善待萬民眾生的國度 vv.9, 14-17

「耶和華善待萬民；他的慈悲覆庇他一切所造的.」在平行的第四段,提到上主慈愛時,上主關心社會上被欺壓的族群:「凡跌倒的,耶和華將他們扶持；凡被壓下的,將他們扶起.」(v. 14)上主也關心地球上自然生態,養育眾生:「萬民都舉目仰望你；你隨時給他們食物.你張手,使有生氣的都隨願飽足.」(v.15, 16) 並為同一地球所有受造聖靈,主持公道:「 耶和華在他一切所行的,無不公義；在他一切所做的都有慈愛.」(v. 17)

公義和慈愛不可割裂,追求慈愛的平均分配,就是追求公義.

權力有兩種,一是壓制他人的權力(overpower),一是建立他人的權力(empower) ,上主的主權是釋放和給力: 「耶和華善待萬民；他的慈悲覆庇他一切所造的.」香港人心並未回歸,那麼多人七一上街表達不同的訴求,其中不少是抗議最近公佈的一國兩制白皮書,擴大一國主權(overpower) ,收緊兩制的自治(empower) .

「好國家指數」.六月二十五日英國公佈獨立機構調查「好國家指數」,愛爾蘭名列首位.指數利用來自聯合國等國際組織的數據來量度一個國家在科技,文化,國際和平與安全,世界秩序,地球與氣候,繁榮與平等,健康與福祉七個範疇共35項指標的表現,得出哪個國家貢獻最大,哪個國家最「自私」的結論.在125個入選國家與地區當中,中國排名107,美國排名第21位,英國排名第7位.

調查機構負責人安霍爾特(Simon Anholt)解釋說,這項指數只是他要在近期啟動的眾多研究項目之一,旨在探討國家的真正意義.他說:「他們的存在是否純粹為了滿足其政客,商家與市民的利益,還是為了全人類和整個地球活躍地做貢獻?」

好國家指數的理想和目標,正是詩篇145篇神國的目標.「耶和華善待萬民；他的慈悲覆庇他一切所造的」是一種國際胸懷,追求族群共融,社會公義.過去香港社會,政府有一種統治策略,叫行政吸納,大學畢業生的精英,都被高薪的政務官僚職系所吸納,成為統治班子,把反政府的力量消弭於未然.

因為公民社會的發展,我觀察到不少優秀的畢業生選擇投身社會壓力團體,國際非政府組織.這班冒起的年輕人,他們的理想,不是高薪安定的生活,乃要建設更理想的社會,認為土地規劃利用,要和當地區民充份溝通,追求社區人性化,族群共融,追求意義多於追求富裕.我們應該為著這班滿有理想的年輕人打氣,鼓勵他們,因為他們提醒我們可以有另類的選項,另一種生活模式,雖然我們不一定同意,但他們迫我們去反思種種現象,不能視而不見.

成為世人的祝福的聖民vv. 10-13

「耶和華啊,你一切所造的都要稱謝你；你的聖民也要稱謝頌你,傳說你國的榮耀,談論你的大能,好叫世人知道你大能的作為,並你國度威嚴的榮耀.你的國是永遠的國!你執掌的權柄存到萬代!」

好叫世人知道.聖民的聖是指屬於上主,被選揀分別為聖.聖民被揀選,目的是叫世人認識上主國度的榮耀,明白上主拯救的大能,這是聖靈降臨成立教會的目的,把福音傳遍世界.我們留意平行的第一段開始時用「我要尊崇」,「我要永永遠遠稱頌你的名!」(vv. 1, 2) 跟著第三身「人要傳說你可畏之事的能力」(v. 6) 到第十節出現「你的聖民也要稱頌你」.上主國度永永遠遠,我作為百代之過客怎可以永永遠遠稱頌上主呢?所以在第四節出現「這代要對那代頌讚你的作為」,一個人會過去,但聖教會繼續,離世的聖民叫得勝的教會(church triumphant) ,留在地上的教會是作戰的教會(church militant) .

我上次在升天主日崇拜講道中,提到耶穌升天,坐在父上主的右邊,表示他掌握了天上地下所有的權柄,但在世上我們看到好像上主不在現場.耶穌升天表示我們進入了不再憑眼見,要憑信心的教會時代；耶穌升天後,聖靈才會降臨.耶穌的故事並未完結,也不是中場休息,乃是好戲剛剛開始,教會走上了舞台,上主的國民走上了舞台.神國的真實性,神國的慈愛和公義,就要由被聖靈充滿的信徒顯明了.

神國有隱藏性.所有權柄來自上主,異教的帝國是上主在歷史默默行事的器皿,在其中更有敬畏上主的信徒參與.約瑟在埃及帝國,但以理在巴比倫帝國,尼希米在波斯帝國,都被上主大大使用,沒有這些信徒參與異教帝國的治,可能人民生活更痛苦.所有權力來自上主的授命,「王的心在耶和華手中,好像隴溝的水隨意流轉.」(箴 21:1)我們要為掌權者禱告,叫他們有兼聽的智慧,善待異見者的寬容.

使人和睦的教會.伊拉克遜尼派與什葉派的教派武裝衝突和仇恨,更使我想起數年前在西敏寺大教堂召開全球聖公會大學校董會.因為同是聖公宗,我們能夠在教會裡面很有歷史價值的「耶路撒冷會議廳」(Jerusalem  Chamber)開會, 1643年英國國會被清教徒大多數把持,在西敏寺大教堂召開神學會議,制訂英國與蘇格蘭教會共同信仰綱領,著名的西敏寺信條Westminster Confession就在那裡簽署的.開會兩個多星期,就在那裡教會安排了伊拉克的什葉派和敵對的遜尼派舉行和談,伊拉克敵對派系信任教會作調停人,多於信任英國政府,教會成了使人和睦的使者,在現今的香港社會,教會使人和睦的角色更為重要.

上帝之城.公元410年,北方蠻族入侵,羅馬陷落,網紀崩毀,人心惶惶,奧古斯丁寫了名著「上帝之城」,比較人的國度和上帝的國度,在紛亂中叫人定睛上帝屬靈永恆的國度,天上的新耶路撒冷.因為自公元313年米蘭諭令(Edict of Milan)頒佈,教會地位合法,及稍後皇帝君士坦丁領洗,教會和地上俗世權勢走得愈來愈近,依附俗世權力的教會要醒悟過來；人的國度會過去,但上帝的國度永遠長存.

有耳可聽就可聽.因為上主的國度是屬靈的國度,祇有信心才能掌握到.當人被上帝的慈愛和憐憫及公義管治時,就經歷到神國的平安.但正因為是隱藏的,所以被人忽略,也被人故意用藉口推辭.福音書經課: 「我們向你吹笛,你們不跳舞；我們向你舉哀,你們不捶胸.約翰來了,也不吃也不喝,人就說他是被鬼附的；人子來了,也吃也喝,人又說他是貪食好酒鬼人,是稅吏和罪人的朋友.但智慧之子總以智慧為是.」(太11: 17-19)

總結:你知道嗎?上帝的國度就是上帝的主權,也是他的慈愛憐憫和公義；他要的領土就是你的心.讓我們的心被上帝的慈愛和公義所充滿,讓我們被他差遣,對世界傳揚他的偉大和慈愛.



\newpage

\section{你聽見嗎}
\label{sec:253}
\textbf{許開明牧師 - 你聽見嗎?}
\newline
\newline
\hyperref[sec:252]{< < < PREV SERMON < < <}
~
\hyperref[sec:toc]{[返目錄]}
~
\hyperref[sec:254]{> > > NEXT SERMON > > >}


馬太福音 13:1-13,18-23
\newline
\begin{longtable}{cl}
\hline
\hline
章節 & 經文 (和合本修訂版)\\
\hline
13:1 & \begin{tabularx}{0.7\textwidth}{X} 就在那天，耶穌從房子裡出來，坐在海邊。 \end{tabularx} \\ \\ \relax
13:2 & \begin{tabularx}{0.7\textwidth}{X} 有一大群人到他那裡聚集，他只好上船坐下，眾人都站在岸上。 \end{tabularx} \\ \\ \relax
13:3 & \begin{tabularx}{0.7\textwidth}{X} 他用比喻對他們講了許多話。他說：「有一個撒種的出去撒種。 \end{tabularx} \\ \\ \relax
13:4 & \begin{tabularx}{0.7\textwidth}{X} 他撒的時候，有的落在路旁，飛鳥來把它們吃掉了。 \end{tabularx} \\ \\ \relax
13:5 & \begin{tabularx}{0.7\textwidth}{X} 有的落在土淺的石頭地上，因為土不深，很快就長出苗來， \end{tabularx} \\ \\ \relax
13:6 & \begin{tabularx}{0.7\textwidth}{X} 太陽出來一曬，因為沒有根就枯乾了。 \end{tabularx} \\ \\ \relax
13:7 & \begin{tabularx}{0.7\textwidth}{X} 有的落在荊棘裡，荊棘長起來，把它擠住了。 \end{tabularx} \\ \\ \relax
13:8 & \begin{tabularx}{0.7\textwidth}{X} 又有的落在好土裡，就結出果實，有一百倍的，有六十倍的，有三十倍的。 \end{tabularx} \\ \\ \relax
13:9 & \begin{tabularx}{0.7\textwidth}{X} 有耳的，就應當聽！」 \end{tabularx} \\ \\ \relax
13:10 & \begin{tabularx}{0.7\textwidth}{X} 門徒進前來問耶穌：「對眾人講話，為甚麼用比喻呢？」 \end{tabularx} \\ \\ \relax
13:11 & \begin{tabularx}{0.7\textwidth}{X} 耶穌回答他們說：「因為天國的奧祕只讓你們知道，不讓他們知道。 \end{tabularx} \\ \\ \relax
13:12 & \begin{tabularx}{0.7\textwidth}{X} 凡有的，還要給他，讓他有餘；凡沒有的，連他所有的也要奪去。 \end{tabularx} \\ \\ \relax
13:13 & \begin{tabularx}{0.7\textwidth}{X} 我之所以用比喻對他們講，是因為他們看卻看不清，聽卻聽不見，也不明白。 \end{tabularx} \\ \\ \relax
13:14 & \begin{tabularx}{0.7\textwidth}{X} 在他們身上，正應驗了以賽亞的預言：『你們聽了又聽，卻不明白，看了又看，卻看不清。 \end{tabularx} \\ \\ \relax
13:15 & \begin{tabularx}{0.7\textwidth}{X} 因為這百姓的心麻木，耳朵發沉，眼睛閉著，免得眼睛看見，耳朵聽見，心裡明白，回轉過來，我會醫治他們。』 \end{tabularx} \\ \\ \relax
13:16 & \begin{tabularx}{0.7\textwidth}{X} 但你們的眼睛是有福的，因為看得見；你們的耳朵也是有福的，因為聽得見。 \end{tabularx} \\ \\ \relax
13:17 & \begin{tabularx}{0.7\textwidth}{X} 我實在告訴你們，從前有許多先知和義人要看你們所看的，卻沒有看見；要聽你們所聽的，卻沒有聽見。」 \end{tabularx} \\ \\ \relax
13:18 & \begin{tabularx}{0.7\textwidth}{X} 「所以，你們要聽這撒種的比喻。 \end{tabularx} \\ \\ \relax
13:19 & \begin{tabularx}{0.7\textwidth}{X} 凡聽見天國的道而不明白的，那惡者就來，把撒在他心裡的奪了去；這就是撒在路旁的了。 \end{tabularx} \\ \\ \relax
13:20 & \begin{tabularx}{0.7\textwidth}{X} 撒在石頭地上的，就是人聽了道，立刻歡喜領受， \end{tabularx} \\ \\ \relax
13:21 & \begin{tabularx}{0.7\textwidth}{X} 只因心裡沒有根，不過是暫時的，一旦為道遭受患難或迫害，立刻就跌倒。 \end{tabularx} \\ \\ \relax
13:22 & \begin{tabularx}{0.7\textwidth}{X} 撒在荊棘裡的，就是人聽了道，後來有世上的憂慮、錢財的迷惑把道擠住了，結不出果實。 \end{tabularx} \\ \\ \relax
13:23 & \begin{tabularx}{0.7\textwidth}{X} 撒在好土裡的，就是人聽了道，明白了，後來結了果實，有一百倍的，有六十倍的，有三十倍的。」 \end{tabularx} \\ \\
[1ex]
\hline
\hline
\end{longtable}
講題:你聽見嗎?　Are you listening?

經文:馬太福音13:1-13,18-23

講員:許開明牧師

場合:香港中文大學崇基學院禮拜堂

日期:2014年7月13日

引言:

「你聽見嗎?」這句話,你說過沒有?

一位丈夫回家,推開門,見太太背著門正在看書,好像不察覺他已回家,於是便嘗試由細聲至大聲,喊「太太」三聲. 太太不耐煩回答說:「我已回答了三次,你聽到嗎?」很多時候,我們以為別人聽不見,其實卻是自己聽不見!是自己耳朵有問題!聆聽的障礙自古到今都非常普遍,常見的毛病有聽不清楚,聽不明白,聽不入耳,聽錯等等.

聽覺是人最早接觸外界的器官,早在母腹中胎兒已有聽覺,聽見媽媽的聲音.人離開世界時,最遲消失的感官也是聽覺.聽覺對人非常重要.有人聽道,常常有領受和得著,生命得到成長.有人聽道多年,生命卻毫無改變與成長.為什麼? 原因在於聆聽的障礙,這問題在耶穌時代曾經發生,今日依然出現,也因此引發耶穌說出著名的「撒種比喻」.

經文結構

全章結構安排已經把聽眾分為兩類:第一類是聽見耶穌公開講論撒種比喻的「許多人」(1-3節)；另一類是聽見耶穌後來單獨. 秘密解釋該比喻涵義,以及為什麼要用比喻教訓人的「門徒」(10-11節).

從比喻的內容也能看見區分和對立.通常人稱這些比喻為「神國的比喻」,因為它們確實每個都談及「天國」的事情.13章的比喻談到:1)沃土和瘠土的區別,2)稗子和麥子的區別,3)天國之子與那惡者之子的區別.這些比喻的主題都與人在聆聽神國度理時的態度有關.

用比喻的原因:

人不愛畫公仔畫出腸,比喻的目的正是要人想一想!比喻需要闡釋,讓人發掘其中的意思,原來較說者直接表達的更多；比喻有奧秘成份,不是每個人都聽得明白,當中卻蘊藏深意.大眾聽不明白,跟從主的門徒卻會明白.

天國的奧秘

新約裡的「奧祕」是指人從未知曉的事實.這些事若不靠屬神的啟示(揭開),是永遠無法明白的.不過,神現在卻把這些事向人啟示.13章12節耶穌把這些奧祕向門徒啟示,卻把群眾摒於門外,表面似乎是專斷的做法.可是主解釋說:「凡有的,還要加給他,叫他有餘.凡沒有的,連他所有的,也要奪去.」因為門徒相信主耶穌,所以得著更多.凡有(知識)的 ,還要加給他(知識),叫他有餘.凡有(財富)的,還要加給他(財富),叫他有餘.凡有(智慧)的,還要加給他(智慧)的,叫他有餘.有的,愈有；無的,愈無,連僅有的也奪去!愈懂聽道,愈有…

天國的力量不斷增長(31-33節),天國的「奧祕」(11節)之所以是「奧祕」,只因不是人盡皆知.門徒明白,得到道種,接受上帝的管治,就進入上帝國度.

比喻的結構

1         撒種的比喻,向大眾講論(1-9節)

2a         運用比喻的目的,向門徒說明(10-13節)

2b         解明撒種的比喻,向門徒說明(18-23節)

比喻大要 (太13:4-8)

土壤	後果

種子落在堅硬的路旁

種子落在土淺石頭地上

種子落在佈滿荊棘的地.

種子落在好土.

道種被飛鳥吃掉

道種很快便發芽,但沒有根；日頭出來,就枯乾了.

道種雖然發芽,但因為荊棘的緣故,沒法長起來.

種子發芽,長大結實,收成有一百倍的,有六十倍的,有三十倍的.

撒種的方式

甲,「道」在路旁(太13:4):

如果種子撒落在這樣的地方 - 其實有些種子不論是用甚麼方法撒下去的,都會落到這種地方 - 它是決不會有機會進入泥土之內.聽了道,卻不明白的人,就是「聽見,不接受」「不明白」不單指頭腦的不領會,不瞭解,也有「不接受」,「不相信」的意思.聽見而不理會對方所講,或者是不重視,也可能是對講者或那道有成見.種子一旦落在路旁就不會被人理會,唯有空中的飛鳥會發現它的存在,然後吃掉.失去道!

乙,「道」落在土淺石頭地(太13:5-6)

就是「聽見,接受,但沒有根基」.在表面上只有一層薄薄的泥土,下面卻是石灰石的岩層.在石頭上面的泥土可能只有幾寸,落在這種土地上的種子,一定能夠發芽,而且芽發得佷快,因為這種土地很快就吸收到太陽的熱氣.可是泥土不深,種子的根往下尋找養料和水分的時候,所碰到的只是石頭,使種子飢渴而死,經不起陽光的照射.

『心裡沒有根』,,無法長久存在人的心裡.人不能只是喜歡聽道,卻沒有容許道長久存在心裡的決心.人就承受不起為義受到患難,逼迫,心中的道就會枯乾,凋萎,立刻就跌倒了.

丙,「道」在荊棘裡(太13:7,22)

世上的思慮,錢財的迷惑,把道擠住了,就是「聽見接受卻受世俗的誘惑影響」.荊棘的土地就是欺哄人的土地.撒種的人在撒種的時候,看見地表好像很乾淨.可是在地裡仍有茅草與蔓藤的鬆根,還有將要復生的雜草,以及四季生長的有害植物. 當雜草欣欣向榮的時候,好種就死去了.世上的思慮,錢財的迷惑,就如同荊棘一般,會把人心中屬天的道擠住,叫道沒有健全生長的空間,以致於空有枝葉不能結實.

丁,「道」在好土裡(太13:23)

「道」在好土裡,人聽道明白了,後來結實,有一百倍的,有六十倍的,有三十倍的.    好土的泥層深厚,純淨並鬆軟,使種子可以深入,獲得它的養料,使它的生長不受阻礙,在好土裡的種子,會結出豐盛的子粒.

請問你的道種又如何?

突破聆聽的障礙

聽道者有時只注意講道者,忽略了道本身,怎樣講?是否有準備?是否動聽?是否有內涵?

聽道者往往忽略耶穌所說的,要準備好自己成為好的土,盛載上帝的道.接受上帝的道有好幾種方式,產生的結果是在於接受者內心的情況.任何話語的命運維繫於聽者,正如有人曾說:『一則笑話的成功,並不是靠說笑者的舌頭,而是在於聽笑話者的耳朵.』一則笑話如果說給一個具有幽默感,並且準備微笑的人聽,一定會達到目的.如果一則笑話說給一個無幽默感,不苟言笑或決不肯自娛的人聽,一定會招致失敗.

誰是這比喻中所警戒的聽者呢?

就是道種落在路旁的聽者.上帝的道無法進入某一種心靈之中,這就好像種子落在曾被許多腳跡踏過的硬土上一樣.

道在路旁的聽者.就是人不重視道,人拒絕所聽的道.

偏見可以使人對於每件他所不願意看見的東西瞎眼

不受教的靈能夠樹立一道不易攻破的障礙.

不肯謙虛受教的心.

管理大師彼得‧杜拉克(Peter Ferdinand Drucker)曾說:“管理問題百分之六十來自溝通不良,溝通不良大半來自於缺乏聆聽.”「拒絕聆聽」容易在人際之間製造距離.想想下列情況:

別人還沒說完,就下了結論,

還沒聽到結局,就急著發表見解,

還沒聽懂,就要別人聽你說,

還沒有聽完,就先回答,結果就是顯出自己的愚昧,自取羞辱.

道在淺土的聽者.有一種聽者的心靈猶如土淺的土地.他是屬於一種未曾把事情經過深刻透徹考慮的人.

對於聽得的道,沒有晝夜思想,沒有根,得不到生長,慢慢枯萎.現代人對於每一種新流行的玩意兒都表示高興,很快就接受,不過也很快就丟棄.

道在荊棘的聽者.有一種聽者對於生活中許多的事物都發生興趣,可是時常把最重要的擠住了.現代生活越來越擁擠,越來越迅速,特徵是人逐漸忙得無暇祈禱.許多事佔據了他的時間,使他無暇研究上帝的道.他要參加許多會議,有價值的工作與慈善服務,使他留不出時間獻給愛心與服務的源頭.他的工作緊緊地抓住他,使他疲累得想不到其他的事.引起危險的並不是那些明顯的壞事,而是好事.因為『次好的永遠是最好之勁敵.』現代人有許多思慮,角色,身份…工作抓住人,嗜好抓住人,學業抓住人.我們聽到道,卻無時間去消化去實踐.

道在好土的聽者.還有一種人如同好土一般.他接受上帝的道有四個階段.

像好土那樣,他的心靈是開啟的,他聽見,聽明上帝的道.

他準備聽從.他決不會因為過分驕傲或過分忙碌,以致不能聽見.有許多人如果能夠簡單地停下來,聽從智慧朋友的忠告或上帝的聲音,一定不會遇到那些令人痛心的事.

他明白.他徹底想過整個的事情,並且知道這事對於他個人的意義,準備好要接受.

他把所聽見的道轉變為行動,使好的種子結出美好的子粒.真正的聽者,就是那聽見,瞭解並遵從的人.

怎樣成為好土?

我們需要突破聆聽的障礙!耶穌說:有耳可聽的,就應當聽(太13:9)

聆聽,好像是一個很簡單的舉動,但正如「聆聽」二字的寫法,「耳」朵,只是其中一個部份,還需要「令」人覺得被聆聽,及「一心」一意地聆聽別人.聆聽並非只用耳朵,更要投入,要用心,用眼,要有同情才起真正的作用.總之,聆聽是一種藝術.

不要常常說:「我明!」

幾年以前,有位智者說:「當有人來給你訴說他的問題,千萬不要隨口說:『我明白,我明白.』對人對神也不要如此!當然,我們對他頻頻說:「我明白」,是有意使他放鬆,消除他的緊張,但是很多時候結果是適得其反,使對方生氣,令他失望.聆聽是重要的,因為聆聽的好壞,所產生的後果很大,為訴說的人,為聆聽的人,都有很大的影響.

聆聽是一種藝術.

所謂藝術,就是做事的適當方法和熟練技巧.有所謂繪畫的藝術,唱歌的藝術,寫作的藝術,彈奏鋼琴的藝術.藝術包含某一程度的熟練,這熟練是由小心的學習,恆久的操練,不斷的改善,加上熱烈的願望而逐漸積聚而成的.聆聽作為一種藝術,有很高的要求.諺語有云:「訥於發言而敏於聆聽.」

聽比講重要

一般人會欣賞一個有口才,說話滔滔不絕的人,而不會欣賞一位懂得聆聽的人.其實上帝造人一個口,兩個耳,可見聽比講來得重要.一般人都覺得發言比聆聽要容易得多,因為聆聽需要克制自我.

聆聽必須把注意力從自我身上轉移到他人身上,以便進入他人之內,而注意他人常是一件吃力的難事:不睬不理,不聽不聞,才是最容易的.此外,我們大家都免不了有成見偏見,以及先入的觀念,這些都是聆聽藝術的勁敵.

一個懂得聆聽的人容易被上帝的道感動；反過來這人也會感動家人,親友,他的社會,國家,全世界.

「聆聽」和「聽聞」不同

聽聞比較被動,是單向的接受,聆聽則是有意識的保持警覺,努力注意訴說者的言詞,和有意表達的意思.所以,聽道是「聆聽」.

聆聽是一項發現或發掘

當我們聽道,或聽演說,或學術會議,或與人飲茶,或看新聞報導,或聽音樂,或聽古典音樂,流水聲,雀仔的叫聲,行雷閃電,我們常能有所發現或發掘寶藏.

但是, 言詞不一定是一句有聲的話語,而可能是一個顧盼,一個手勢,一脈眼神,甚至少許的靜默.如果我們只留意發聲的言詞,我們也許只能得到一些片面的零星資料,甚至或許根本沒有接觸到對方的真正問題.因此,除了言詞的意義之外,還必須注意言詞的聲調,面部的表情,說話的快慢,以及在訴說途中的停頓和靜默.基本上,聆聽就是關愛.

聆聽是多方面的使用

我們能夠用耳朵,用頭腦,用情感,用心靈聆聽.

有時,我只用耳朵來聽,只感到言語的聲波敲擊我們的耳膜.有時這些聲音是有意義的,這時,我們當然會高興聽下去 .

我們也能夠只用頭腦來聽,這時,我們很容易覺到對方來意不善,有意向我們挑戰.當對方提出一種想法,或一個觀念,我們就在心中盤算怎樣招架,甚至如何反擊；我們一面聽他講述,一面在心中搜尋質疑,安排論證；並等候機會殺出,指證他大錯特錯.有時我們只用頭腦來聽,討論一下子就變成辯論,對話很快就變對抗.聆聽,不要只偏好「用腦」.聆聽三步曲:1用耳去聽,聽清楚每句話,2用腦去聽,分析所講的是否合理?3用心去聽,就是試著進入他的內心

用心去聽,就是試著進入他的內心,「穿著對方的鞋子行走」；就是以他的觀點看事；在適當的時機,提出適當的問題,為了要明白他所說的,就全神灌注地靜聽,確認當前的問題；顯示我們見到實有一談的價值,並表示衷心的感激,因為他信任我們,給我們分享他個人的,內心的感受.這樣的聆聽,才可以說:「聆聽就是愛.」

如果頭腦和心能夠學習去聆聽他人的思想和心靈,愛護,和關懷的花朵就會盛開,並使人們和睦相處合而為一.只有真正的聆聽可以彌補裂痕,摧毀藩籬.拒絕聆聽,則製造冷膜,離間人心,產生仇恨.

聆聽是一種治療

交談不一定有治療力!講話能夠喚醒他人,也能夠激怒他人,但本身不一定治療的能力.

開始聆聽,便要開始治療的進程.嘗試學習聆聽風,聆聽樹,聆聽鳥,聆聽自己,聆聽天鵝(媽媽,太太,仔女),聆聽上主.在明華堂,常常可聽見鳥兒求偶的聲音.聆聽一個人,每天聆聽,即使每天五分鐘也好.如果每個人都這樣學習聆聽,大家都會更加健康.當聆聽是一種出於真心的關懷的時候,它就帶著醫治的能力,能給憂傷的心靈帶來極大的安慰.

聆聽需要專注,同時要主動,也必須用心.專心聆聽對方的分享,傾聽他人的心事,聽他訴苦,瞭解他,而不急著給予建議,也不急著告訴他下一步該怎樣做.因為訴苦的人在當下所渴望的就是有人肯瞭解他,關心他,願意多聽他,讓他能宣洩自己積壓的情緒,所以「快快地聽」,「專注地聽」是很要緊的.當你進入對方的心靈世界,明白他的感受,為對方的難處流淚,為他禱告,這最能感動人心.

被動和主動的聆聽

容我以另一個角度來說:聆聽能有兩種情況,被動的聆聽和主動的聆聽.這裡,我們不是在講論無心的聽聞,而是有意識的,集中注意力的聆聽.然而這樣的聆聽也可以分為被動的,或是主動的.

所謂被動的聆聽,是以接受的心境去聆聽,而主動的聆聽則是同時以分析和批判的態度去聆聽.如果我們要聆聽全部的意義,我們應努力以被動,即接受的心境去聆聽,為了聽到對方所訴說的一切,而不是我所希望聽到的,或要求聽到的東西.

許多做父母問子女問題, 子女的回答常令他們不開心.因為不是他們他們的期望.(例:父母期望好成績的兒子去學醫,兒子卻選擇讀「音樂」)今日香港社會有個不好的現象,有些政府官員說話,或內地官員說話,無論他們說甚麼?某些對他有成見的人都聽不入耳,以致造成社會矛盾,連正話,好話都失去.真正被動式的聆聽,我們對自己不作什麼佈局,而是讓對方的言詞自然地流入我的耳朵之內.唯有在這種情況下,言詞的意義,其他的表達,才能原原本本地吸收消化.不然的話,我們聽到的,可能不完全是訴說者的原來的意思,而是我們自己的理論和期望的反映.

實際上,良好聆聽的標準就是:被動的和接受的心境.存有這種心境的聆聽者,在他看來,每一個人都是獨一無二的,都有個人不能轉移的權利.每一個人都應尊重他所說的,他絕不把前來訴說的人,分門別類的,配置在鴿籠裡.「如果我要人尊重我的話,我也要尊重別人的話,聆聽就是尊重!」

耹聽,需要快過講話

雅各書1:19經文說:「你們各人要快快地聽,慢慢地的說,慢慢地動怒.」今日的政教問題,政府白皮書等各人要快快地聽,慢慢地的說,慢慢地動怒.父說:「請講慢些,孩子,你講得好快!」兒子回應:「阿爸,不是我講得快,而是你聽得太慢!」讓我們多用理解的耳朵聆聽,用充滿愛的語言說.「快快地聽,慢慢地說.」多聽,少論斷,多理解,少批判.多聽別人怎麼說,不要動輒指摘別人.而要耐心聽,聽清楚對方在說什麼,想要表達的是什麼,讓我們多一點體恤和耐心,不要心浮氣躁.

聆聽,要專心與尊重

一個人向你說話,只希望你能聆聽.其實,你只需要聽,可以不說話的.我卻見過有子女原本想把某件事告訴父母親,最終卻變成父母向兒女訓話,雙方不歡而散!聽人說話的時候要專心聽,不要顯出「無聊」的樣子,或露出一副「我早知道了」的表情.不要打斷對方的談話,要讓對方完整地把話講完,即使他有很多抱怨,也要耐心聽完,因為「傾聽就是尊重」.如果對方還沒有說完你就打斷,會讓對方感受到他所說的毫無價值,不值得你留意,這樣一來就顯出你對他的輕視.希波的奧古斯丁(Augustine of Hippo)有言「請你,我,大家一起把傲慢放在一邊.請我們中任何一人千萬不要自以為擁有全部的真理.讓我們尋找真理,猶如尋找什麼未知的東西.讓我們以熱愛,以真誠尋找真理.我們中誰也不魯莽,誰也不冒昧,自以為找到了所有的真理,如果我質問你太多,也請你讓我聆聽你.請讓我同你交談對話,因為在我方面,我並不冒昧地認為,我已經完全瞭解了你.」

我們需要一種開放的氛圍和胸襟.無論是你,或者是我,都不以為自己擁有全部的真理.我們每一個人都獨一無二的,包括我們的過去,我們的現在,我們的將來.我們需要彼此聆聽.我們需要聆聽生命,聆聽經驗,聆聽他人,也必須答覆他人.懷抱他人,把他溶入我內,為了完全瞭解他人而努力不懈.我們每一個人都有聆聽的能力,也有作出回應的能力,以及建立富有意義的人際關係的能力.

在我講道生涯中,曾經遇過一位對我很有鼓勵的弟兄(現已離世),他每次都愛坐在最前,很用心的聽道,總帶著微笑向著我,有時會點頭示意.若有不同看法時,他會約我食飯提出來討論.這位弟兄,有事跟他商討時,他會細心聆聽:當人跟他說話的時候,他會暫時放下手邊的工作,放下書報,離開電腦,放下手機,看著對方的眼睛.他很用心地聽我講話,我就感到很溫暖,被尊重.他多聽別人說話,自己肝卻不多說話.

聆聽要用情

這裡有兩個名詞值得我們思索一下它們的意義,一個是同情(sympathy),另一個是移情 (empathy).

同情是有同樣的感受；這是一種心境,雙方互相分享對方的情緒,感覺,情景,形同身受.

移情是一種投射的能力,把一個人的性格投射到另一個人身上,把後者成為前者的化身.我們愛誰,就生活在他身上.偉大的演員就有這種能力,能夠投入劇中人的性格和行動舉止.

1聆聽三要:專注,接納,不批判

如何要「令」人覺得被聆聽?最重要的是聆聽者在過程中保持三個基本態度:「專注」,「接納」及「不批判」.

分享者可以在一個舒適,平和的氣氛下,放鬆地表達自己的想法與感受,使他們覺得安全,被尊重及無壓力.除了這三個基本態度外,聆聽者更需要放下自己的經歷,概念及想法,才能專注地聆聽分享者說話的內容.最近讀過一篇文章,提醒我聆聽的技巧:

第一,聆聽者要如鏡子般「一心」的如實反映分享者的情況,在傾訴的過程中,分享者可以對自己的情況有更完整的檢視,從中找出另一個角度的看法來瞭解自己當時的情況或情緒.

第二,聆聽者應該如影子般「一心」跟隨及陪伴分享者,隨著分享者的步伐,讓他們主導談話.值得注意的是聆聽者有時會側重於自己「想聽」的範疇及方向,不自覺地主導及帶領整個分享,因此應該讓分享者表達想說和想分享甚麼.

第三,聆聽者應該如棋子般「一心」跟隨棋手,面對困局及勝局.

聆聽,要先忘我

我們常常只把焦點放在自己身上,而沒有為別人多想,結果當別人說話的時候我們沒有留心聽.箴言18:13 經文說:「未曾聽完先回答的,便是他的愚昧和羞辱.」很多時候我們只想著自己心裡要講的,而沒有耐心好好聽完對方所講的話.我們容易在聽人講話的時候心不在焉,一面聽,一面心裡在想著自己的事情.

有一次,集盲,聾,啞殘缺於一身的美國女作家海倫凱勒在一所大學演講,會後一位同學走到她前面問她:“海倫凱勒女士,請問一個人要如何才能獲得最大的快樂?”這位同學大概是從海倫凱勒身上得到了許多啟發和感觸,所以問了這麼樣一個問題.

對這個問題,海倫凱勒(Helen Adams Keller)的回答很簡潔,很有力,她只簡單地說了兩個字:“忘我.”Forget yourself. 忘掉你自己吧!很多時候我們太多顧念自己,無論做什麼都只想到自己的利益,顧念自己的感受.

當一切都是以「自我」為中心時,我們活著就很痛苦.人際關係的突破始於「忘我」.若是你過去只在意自我的表然而忽略與人的關係,沒有忘我地投入,讓我們有一個全新的開始.

聆聽,開放心靈讓上帝說話

蘇格拉底告訴我們,自從孩童的時候開始,他不時聽到一個聲音,他稱之為「神聖的記號」,這是指從人的良知發出的聲音.這聲音禁止他做某些事情.有時候,這聲音可能是一個明確的指示；有時候可能只是預示未來的好運或惡運.

現代心理學家稱這種現象為「內在聽聞」,並指出這種情況並不少見.法國的聖女貞德在十三歲那年,感受到一種內在的催促,形容為「祂(上主)的聲音」,催促她自敵人的手中拯救法國.這些來自上天的聲音只是發自他∕她們良心的呼聲嗎?還是另外的一件事情?弗洛伊德的門生心理學家特奧多爾‧列伊克(Theodor Reik)出版了一本書,書名是《用第三隻耳朵聆聽》.這第三隻耳,就是上主的聲音.上主的兒女,多在安靜耹聽上帝的聲音,使你得著智慧,更新,和進步.

結語:各位,今日的講道,你聽見嗎? 願你聽見,聽明,也聽從!



\newpage

\section{一生之久的愛}
\label{sec:254}
\textbf{林豪恩先生 - 一生之久的愛}
\newline
\newline
\hyperref[sec:253]{< < < PREV SERMON < < <}
~
\hyperref[sec:toc]{[返目錄]}
~
\hyperref[sec:255]{> > > NEXT SERMON > > >}


哥林多前書 13:4-8
\newline
\begin{longtable}{cl}
\hline
\hline
章節 & 經文 (和合本修訂版)\\
\hline
13:4 & \begin{tabularx}{0.7\textwidth}{X} 愛是恆久忍耐；又有恩慈；愛是不嫉妒；愛是不自誇，不張狂， \end{tabularx} \\ \\ \relax
13:5 & \begin{tabularx}{0.7\textwidth}{X} 不做害羞的事，不求自己的益處，不輕易發怒，不計算人的惡， \end{tabularx} \\ \\ \relax
13:6 & \begin{tabularx}{0.7\textwidth}{X} 不喜歡不義，只喜歡真理； \end{tabularx} \\ \\ \relax
13:7 & \begin{tabularx}{0.7\textwidth}{X} 凡事包容，凡事相信，凡事盼望，凡事忍耐。 \end{tabularx} \\ \\ \relax
13:8 & \begin{tabularx}{0.7\textwidth}{X} 愛是永不止息。先知講道之能終必歸於無有；說方言之能終必停止；知識也終必歸於無有。 \end{tabularx} \\ \\
[1ex]
\hline
\hline
\end{longtable}
講題:一生之久的愛　The Love to Last a Lifetime

經文:哥林多前書13:4-8

講員:林豪恩先生

場合:香港中文大學崇基學院禮拜堂

日期:2014年7月27日

一. 引言

作為助理校牧這身分的最後一次講道,我想送一段歷久彌新的經文給大家.我相信大家非常熟識這段經文,其中一個原因是這段經文在婚禮的「曝光率」非常高,差點令人以為這是婚禮專用經文.我最近再讀這段經文的時候,有一點個人的體會.經文中提及的愛,不單單在婚禮適用,原來對我們整個人生也適用.這愛,彷彿應對著我們人生的每一個階段.

二. 經文

13:4愛是恆久忍耐,又有恩慈；愛是不嫉妒；愛是不自誇,不張狂,

13:5不做害羞的事,不求自己的益處,不輕易發怒,不計算人的惡,

13:6不喜歡不義,只喜歡真理；

13:7凡事包容,凡事相信,凡事盼望,凡事忍耐.

13:8愛是永不止息.

三. 信息

1. 愛是恆久忍耐；又有恩慈

第一句送給父母以及分擔父母責任的朋友.忍耐這個字的意思是有耐心地等待,而且不氣餒,甚至在被頂撞和對抗的時候,也不反擊,不報復.恩慈即是仁慈,憐憫和慈愛.恆久忍耐,不放棄,不攻擊,還要疼惜慈憐,這樣的愛最適合送給要養育子女的父母.

作為父母的,仍記得這樣的經驗嗎?有一位新任媽媽在網絡上求助說:「我想問下各位媽咪,妳地有冇BB夜晚成晚喊既問題?我個BB出生左10日咁上下,佢夜晚只係識得喊又唔肯訓覺,而且喊到拆晒天,搞到我好困腦.我又餵過奶,又幫佢換過尿片,都唔知佢想點?好多時食奶佢又食得唔多,唔夠一陣佢又要再食過,又掃唔到風,有時會嘔返D奶出黎!簡直係全晚俾晒佢.我覺得好辛苦,要等幾幾多個月BB夜晚先至唔喊呢?」

終於等到嬰兒變成一位滿有主見的「小成人」,這個時候,父母所要付出的忍耐和恩慈,可否減少呢?有一位家長分享說:「我家的老大是一個小五學生,是不是這個年齡就開始反叛呢?舉例:請他去洗澡,叫了五,六次,還不動眼盯電視.關電視,他不高興,說不想洗澡就不洗澡之類的話.昨天更糟,和他去逛街.告訴他要回家時,他的語氣及臉色極差.問他想怎樣?他不回答.老公開車來接時,告訴老公,老公責備他,問他是不是不想回家.他竟然說是.老公生氣就叫他下車,他真的下車了.我跟老公實在不知如何.於是車開走繞一圈又回來.我因擔心故拉回他.回家後又好像甚麼事也沒有發生似的.」

孩子到甚麼時候才不需要父母操心呢?俗語說:「養兒一百歲,長憂九十九」.據英國《每日郵報》引述調查報道,在孩子年滿30歲之前,英國父母平均經歷了超過8年時間的不眠之夜.哪一個年齡階段最令父母擔心呢?原來是18-30歲階段是孩子,在這期間父母會因為牽腸掛肚而歷經多達1872個難眠夜晚(例如:兒女去夜街,父母在家等).可見子女就算長大成人,這恆久忍耐又有恩慈仍然是父母需要的.

2. 愛是不嫉妒

這一句送給識行識走的寶寶,以及剛有弟弟或妹妹的幼兒.

有沒有見過這些幼兒看到媽媽在照顧別人時,就會表現奇怪,可能會吸吮大拇指,玩弄頭髮而大哭大喊.因為他們覺得媽媽被佔有,覺得自己被冷落了.因此,他們的奇怪行為是希望奪回媽媽的注意力.有些小朋友在弟弟或妹妹出生後,在他的小小世界裡甚至連結到絕望的想法:「爸爸媽媽都只關心BB,我要被拋棄,不再被愛了!」

有一位聰明的媽媽,懷著小娃娃時,就教大兒子常在媽媽的肚子上告訴娃娃說:「我是哥哥,你在裡面,可要乖乖的呀!」臉上充滿了當哥哥的喜悅.可惜,當媽媽肚中的娃娃出生後,娃娃的一顰一笑都成為大家注意的焦點,每個人都搶著抱,搶著逗小娃娃,而哥哥卻被遺忘了.哥哥好像也變得不再可愛了,常會央求媽媽抱他,要不就大哭大鬧；哥哥也會吵著要喝牛奶(要用奶瓶裝的才算數)；哥哥開始像個娃娃似的,動不動就哭,動不動就尿褲子；哥哥不再懂事聽話了,有時還似個小惡棍般地把小娃娃弄哭,甚至要媽媽把娃娃抱去丟掉.

嫉妒,似乎是感到失去愛的反應.聖經說,愛是不嫉妒,會不會愛就是治療嫉妒的良方呢?小朋友需要爸爸媽媽重覆地作出愛的保證:「弟弟出生了,但是我們還是很愛你!」雖然嫉妒是正常的情緒,但是大人還是要善用方法,以愛把這此負面情緒轉化為幫助他們成長的動力.

3. 愛是不自誇,不張狂

自跨,俗語即是「認叻」,張狂這個字的字根是「吹」,有吹噓的意思,我們俗語說「吹水」剛巧也有「吹」的意思.「愛是不自誇,不張狂」這句,我們可以考慮送給開始上學的兒童.

小朋友喜歡「認叻」,因此往往成為笑話的主題.例如:

有三個小男孩在自誇自己的爸爸.其中一位說:「我爸跑得飛快,可以在弓箭射出時起跑,然後箭還沒到時,人就先到了終點.」 另一位男孩說:「那沒什麼,我爸可以在開槍時起跑,子彈還沒到終點他就先到終點了.」 最後一個說:「那沒什麼,我爸在公家機關上班,他五點下班,但還不到四點就到家了.」

另外三個小朋友比較誰的爸爸比較偉大. 第一位說:洞庭湖的洞是我爸爸挖的.第二位 說:喜馬拉雅山是我爸爸建造的.第三位說:死海是我爸爸殺死的.

小朋友「認叻」有時會向另一個竟想不到的方向:

第一個小朋友:我昨天嘔了一次.

第二個小朋友:我仲慘啦,個日我嘔左十次.

第三個小朋友:你地都唔夠我慘,我一日嘔左一百次啊!

曾聽過有父母驚訝孩子上學後竟然學懂說謊,說自己做了班長,查實並無此事.孩子這些吹噓和言過其實的行為,有時令父母既憤怒又擔心.為甚麼這個時期的小朋友喜歡自誇?有時,小孩子「認叻」是要尋求認同和自我肯定；有時是為了博取成人的讚賞及注意；有時只是把願望和想像當做真實發生的事說出來而已.

聖經說,「愛是不自跨,不張狂」.愛,會否也是消除自誇和吹噓的元素呢?美國孩子研究協會講師格路恩堡．馬茲諾夫人反對用揭穿孩子的謊話,做為去除孩子說謊方法.她說:「無論孩子說謊的理由是什麼,父母親們必須認為這只是她們成長過程中的一種現象而已；務必使孩子明白自己為父母所信任不疑.」她認為說謊並不是意味著孩子開始變壞,而是在親子關係中或許有那麼一點溝通上的不良所導致的一個行為產生而已；耐心的傾聽與良好的溝通,正是最好的教養原則.

但願孩子經驗得越多成人透過接納和信任表達的愛,就越能夠減少自誇和張狂的表現.

4. 不做害羞的事,不求自己的益處

這一句送給青少年朋友.「害羞的事」的意思是不得體和不適當的行為,「求自己的益處」的意思是一定要得到自己想要得到的.

青少年朋友的特徵是改變,這包括發育時期身體的改變,情緒的改變,興趣和活動範圍的改變,對各種關係的重要性的改變,這些改變互相牽連,互相影響.這些改變足以產生巨大的動力,推動他們自我發展,邁向成熟,也可以驅使他們年少輕狂,誤入歧途.有些青少年以拒絕父母的勸告,來肯定自己的獨立自主,雖然父母苦口婆心,他們亦知道父母是對的,但總是給朋輩的外力所吸引,反叛是證明自己長大 的呼喊.

有一位母親,育有一女一子,女兒本來是她心目中的天使,怎料踏入少年期之後,就變成家中的小魔鬼.她的女兒自少天資聰敏,移民澳洲後成績一直都是Ａ級學生,她還彈得一手好琴,年紀輕輕就取得演奏級的水平.她自幼返教會,個人盡責守規矩.但踏入十四歲,經過一次失戀和交了一些貪玩的朋友後,她就來一個大變身,開始夜歸,逃學,飲酒,不再學琴,打扮入時,拒絕與父母交談,動怒時還會說一兩句粗話,她的轉變叫母親十分擔心.(區祥江:「暴風少年的沖擊」)

青少年在尋找自己的過程中, 如何能夠「不做害羞的事,不求自己的益處」呢?成人如何能夠幫助處身於轉變中,以及尋求突破自我的青少年朋友減少不適當的行為,以及懂得考慮身邊的人的意見呢?聖經說,愛是「不做害羞的事,不求自己的益處」,會否提示我們,青少年在尋找自我的過程中仍然十分需要愛呢?

我曾經認識一位朋友,年少時誤入歧途,染上毒癮.儘管他也想孝順父母,但毒癮已控制到他連父母都不認,傷透父母的心.直到有一次,他被黑社會追斬後,奄奄一息,傷痕纍纍地躺在醫院病床上,醒來後第一眼看到的,是媽媽跪在床前為他禱告.他說:「當時,我終被感動了,原來媽媽沒有放棄我,主耶穌也沒有放棄我.」他接受媽媽的安排,到香港晨曦會去戒毒.

各位青少年朋友,如果我們知道並不需要完美,也不需要特別,也有人愛自己；如果我們在軟弱和迷惑的時候,仍然知道有人對自己不離不棄,可能,我們就沒有那麼大的引誘去做害羞的事,也未必會被追求自己的益處囚禁而能夠看見身邊的人的需要.

5. 不輕易發怒,不計算人的惡

這句是送給青年人的.有人說:「三十歲之前而不是左派,這個人一定沒有靈魂.」青年人滿懷理想,並且以理想去審視世界；青年人也充滿動力,並且以動力去改變世界.青年人的理想和動力,可以為更新世界,推動世界前進.

有時,我們會發現理想和現實之間的距離.這種距離可以令人埋怨,為甚麼我爸爸不是李嘉誠；可以令人憤怒,為甚麼這些人「阻住地球轉」；可以令人產生破壞和報復的想法.

有一位青年曾經有這樣的經驗:「步入大學,意味著可以自由一些了.我可以隨意挑選我喜歡的書閱讀,可以選修我喜歡的課程,可以騎著自行車四處去聽講座.我遇到了幾位令人尊敬的師長,他們為人正直,教會我應該獨立思考而不再盲從,教導我應該有批判的精神,教育我學習辨別真理和謊言.我像一頭飢餓覓食的小獸,搜尋著過去不曾見過的文字,傾聽著不同的聲音.可接觸越多,痛苦反而越大.我常常站在先哲巨人們的肩膀上,指責前人的錯誤竟然如此愚昧；我常常看著反思文革和八九學潮的資料,追問「歷史怎麼會是這個樣子」?我尤其搞不懂的是,為什麼我這樣一個懵懂的學生對此都有所察覺,而我們的主流聲音卻隻字不提?我不再相信歌舞昇平──那些不過是粉飾太平.我變成一個「憤青」,憤怒且嫉世.曾經的滿懷希望轉變為滿腔憤怒,如同槍炮,向周圍的一切開火.沒有光的日子,無疑是黑暗而且痛苦的,我常常歎息「黑夜給了我黑色的眼睛」,而我已經沒有勇氣再繼續「尋找光明」了.我開始聽一些頹廢的地下音樂；我開始對自殘和自虐的行為津津樂道；我會躲在黑暗的角落裡；我對著天橋上來來往往的車流喊:愚蠢的人們,你們的行為同自殘並無區別.我有時也想,幹嘛這樣作賤自己?孫鵬飛:「走出「憤青」」)

世界上很多事情可以令人憤怒,也有很多對我們不好的人以及他們的惡行可以去數算.聖經提醒我們,如果我們只是看見別人對自己的不好,總是認為身邊的人都虧欠自己,因此而合理化自己經常埋怨,隨便發憤,隨時反擊,那麼,這有可能是欠缺愛的表現.聖經說,愛是不輕易發怒,不計算人的惡.

甘地稱他主張的非暴力,不合作運動為Satyagraha.所謂Satyagraha,是握住真理,愛與虔敬的意思.甘地說,實踐Satyagraha的人,要放下憤怒和報復,相信人最根本的善,以自己最平靜的受苦,喚醒和邀請掌權者和所有人的心的轉變,一起合作走向真理這條道路.

6. 不喜歡不義,只喜歡真理；

這句送給在社會打滾多年的成年人.人生已經走了一大半,如果要建功立業,應該已經略有所成；如果屢遇挫折,一事無成,難免感到焦慮和急躁.眼看自己時日無多,而且長江後浪推前浪,可能成為挺而走險的推動力.

眼見有人「轉一轉手就賺到盆滿缽滿」；又見有人「行行企企食飯幾味」；聽說有人「飲杯芧台就官運亨通」,不禁令人慨嘆:「識多個人好過識多隻字」；甚麼有種「忠忠直直終須乞食,奸奸姣姣朝煎晚炒」的苦澀.《河南商報》刊登一篇檢察官撰寫的文章,指出內地接近八成被查處的省部級高官在腐敗期間獲得晉升.愈腐敗愈升官的原因是這些貪官編織了兩張網,一張是官場關係網,一張是上司保護網.無官不貪的後果是劣幣驅逐良幣,狗官備受重用,清官慘遭排斥.(彈指春秋:「品德不壞 致富不快」)

聖經說,愛是「不喜歡不義,只喜歡真理」,對於感到「人在江湖身不由己」的朋友是一項重要的提醒.一事無成,事日無多,的確令人焦急.然而,古人有云:「是非成敗轉頭空,青山依舊在, 幾度夕陽紅.」原來,在扭曲的世代,仍然堅持有所為,有所不為,這也是愛的表現.孟子在與別人辯論什麼叫做「大丈夫」時提出做人的準則:富貴不能淫,貧賤不能移,威武不能屈.

7. 凡事包容,凡事相信,凡事盼望,凡事忍耐.

這一句我們尊敬的前輩們.

有人認為,當人踏入人生最後歷程,是一個把以往成敗總結的時候.故此,如果以往的發展是完美的,有成就的,在這個階段便會享受其成果,心理上便有圓滿的感覺,不然,就會在心理上積存失望感.這種圓滿包括接受一生的功業,而無後悔自己之作為.

然而,回望過去,誰人能夠完全沒有遺憾呢?愛是凡事包容:包括接受自己以及別人不同的生活方式及其貢獻；包括在經濟及環境威脅下所作之努力.遠眺將來,我們能夠操控的越來越少了,學習放下吧!如何能夠放得下呢?愛是凡事相信,凡事盼望.那麼,如何面對現在呢?竟然又回到開始－－都是忍耐.從前,我們忍耐,等待孩子成長,現在,我們要忍耐的是自己.眼睛沒有從前般銳利了,行動沒有從前般快捷了,身體沒有從前般強壯了,甚至精神也沒有從前般旺盛了.這些改變,都需要我們以忍耐待之.

保羅說:「所以我們不喪膽．外體雖然毀壞,內心卻一天新似一天.」(林前4:16)這一天新似一天的體會,其中一個可能是從整合和發現自己人生的意義而來的.有一位女士有這樣的體悟:

「坐在梳妝台前梳著頭,我的心情越來越沉重.兩鬢的白髮似乎越來越茂密了,我不禁嘆了一口氣.當初的我,深怕拔了會長更多,所以都不敢拔掉它.但是,隨著年幾越來越大,就算不去拔,白髮仍不斷的冒出來.對我而言,那不是智慧的結晶,而是煩惱的總數.這晚睡前,我忍不住向丈夫抱怨.丈夫看著我的頭髮,接著認真地說:「這一根白髮,是你天天料理三餐,為了家務操勞的結果；這一根,是你在孩子生病時,細心照顧的結果；這一根,是我們在難過和沮喪的時候,你陪伴我們的結果......謝謝你,有你真好.」聽到這些話,我感動得笑了.原來,白髮也是愛的記號,就像是家人為我戴上的美麗冠冕.(蒲公英月刊:「白髮」)

8. 愛是永不止息

最後一句,也作為總結:「愛是永不止息」,送給每一位朋友.正如約伯說:「我赤身出於母胎,也必赤身歸回」(伯1:21)在人生之中,在天地之間,有甚麼是真實而恆久的呢?當世上一切都帶不走的時候,有甚麼是可以永遠相隨的呢?當沒有人能夠和自己一齊穿越生命的盡頭的時候,還有甚麼能夠陪伴呢?保羅有這樣的體會:

「誰能使我們與基督的愛隔絕呢?難道是患難嗎?是困苦嗎?是逼迫嗎?是飢餓嗎?是赤身露體嗎?是危險嗎?是刀劍嗎?」(羅8:35)

「因為我深信無論是死,是生,是天使,是掌權的,是有能的,是現在的事,是將來的事,是高處的,是低處的,是別的受造之物,都不能叫我們與上帝的愛隔絕；這愛是在我們的主基督耶穌裡的.」(羅8:38-8:39)

愛是永不止息.願這真實而恆久的愛伴隨我們走過人生的每一程,帶領我們進入永恆.



\newpage

\section{耶穌頭腦簡單麼}
\label{sec:255}
\textbf{許開明牧師 - 耶穌頭腦簡單麼?}
\newline
\newline
\hyperref[sec:254]{< < < PREV SERMON < < <}
~
\hyperref[sec:toc]{[返目錄]}
~
\hyperref[sec:256]{> > > NEXT SERMON > > >}


馬太福音 14:13-21
\newline
\begin{longtable}{cl}
\hline
\hline
章節 & 經文 (和合本修訂版)\\
\hline
14:13 & \begin{tabularx}{0.7\textwidth}{X} 耶穌聽到了，就從那裡上船，私下退到荒野的地方去。眾人聽到後，從各城來，步行跟隨他。 \end{tabularx} \\ \\ \relax
14:14 & \begin{tabularx}{0.7\textwidth}{X} 耶穌出來，見有一大群人，就憐憫他們，治好了他們的病人。 \end{tabularx} \\ \\ \relax
14:15 & \begin{tabularx}{0.7\textwidth}{X} 傍晚的時候，門徒進前來，說：「這地方偏僻，而且時候已經晚了，請叫眾人散去，他們好進村子，自己買些食物。」 \end{tabularx} \\ \\ \relax
14:16 & \begin{tabularx}{0.7\textwidth}{X} 耶穌對他們說：「不用他們去，你們給他們吃吧！」 \end{tabularx} \\ \\ \relax
14:17 & \begin{tabularx}{0.7\textwidth}{X} 門徒說：「我們這裡只有五個餅、兩條魚。」 \end{tabularx} \\ \\ \relax
14:18 & \begin{tabularx}{0.7\textwidth}{X} 耶穌說：「拿過來給我。」 \end{tabularx} \\ \\ \relax
14:19 & \begin{tabularx}{0.7\textwidth}{X} 於是他吩咐眾人坐在草地上，就拿著這五個餅和兩條魚，望著天祝福，擘開餅，遞給門徒，門徒又遞給眾人。 \end{tabularx} \\ \\ \relax
14:20 & \begin{tabularx}{0.7\textwidth}{X} 他們都吃，並且吃飽了。門徒把剩下的碎屑收拾起來，裝滿了十二個籃子。 \end{tabularx} \\ \\ \relax
14:21 & \begin{tabularx}{0.7\textwidth}{X} 吃的人中，男的約有五千，還不算婦女和孩子。 \end{tabularx} \\ \\
[1ex]
\hline
\hline
\end{longtable}
講題:耶穌頭腦簡單麼?　Did Jesus Have a Hole in the Head?

經文:馬太福音14:13-21

(參考:馬可福音6章35-44節；路加福音9章10-17節；約翰福音6章1-13節)

講員:許開明牧師

場合:香港中文大學崇基學院禮拜堂

日期:2014年8月3日

引言

小時候,曾經有人問過我:「你信耶穌,耶穌會給你飯吃嗎?」 這一個問題,簡單,直接.但是背後卻包含嚴峻的質疑:「信耶穌有何實質作用?」信耶穌對我有何意義?信耶穌是否實際?

信息

在高小主日學,我無意中聽見「五餅二魚」的神蹟,好像解答了我心中的疑惑.信耶穌是真的可以得飯吃!「五餅二魚」餵飽五千人的神蹟,真利害!耶穌竟然能使小變大,從少變多,跟隨耶穌有飯吃!還可以吃飽,吃剩!超棒!信耶穌有飯吃,食耶穌得永生,信耶穌什麼都有!信耶穌就這麼簡單,這麼有利,這麼清楚,就好了!

這個「五餅二魚」神蹟是否在解決這個信仰的迷思?

我相信不是!因為耶穌在曠野受魔鬼的試探時,魔鬼曾經建議耶穌把石頭變成餅,當時耶穌拒絕說:「人活著不是單靠食物,乃是靠上帝口裡所出的一切話」(太4:4).假如每個人來到耶穌基督跟前,耶穌就把石頭變成餅,一定會有許多人為食而來,為食而信,為了某些「好處」來相信耶穌.耶穌若可如此輕易吸引人歸向祂,祂不必背負釘十架的痛苦了.魔鬼正是希望耶穌這樣踩入陷阱,可是耶穌怎會如此輕易上當?還有,信耶穌若祇求吃飽,祇求神蹟,純為滿足精神上好奇心,這樣的信仰豈非很低層次?祇求滿足個人食慾,豈非與動物,野獸無異!

如果五餅二魚的神蹟,不是要突顯信耶穌使人食飽,使人得利得福?它背後要突顯什麼些呢?

本段經文可見,耶穌經過一天的工作,醫病,趕鬼,傳道,直到晚上,出來看見群眾餓了.

門徒的建議:天將晚的時候,門徒進前來說,這是野地,時候已經過了,請叫眾人散開,他們好往村子裡去,自己買喫的.(14:15)

耶穌的回應:耶穌說,不用他們去,你們給他們喫罷.(14:16)

門徒的反應:門徒說,我們這裡只有五個餅,兩條魚.(14:17)

耶穌的回應:耶穌說,拿過來給我.

門徒的方法:他們去,自己搵食；只有五餅二魚(小小的,不夠吃的),無法解決五千人的饑餓問題.

耶穌的回答是「他們沒有吃的,你給他們吃吧!」

門徒回答:「我們只有五餅二魚!」意思是:即使我們拿出所有,祇有五餅二魚!怎能餵飽數以千萬計的群眾.耶穌這種答案,是否太簡單?如果今日我們這樣的回應別人的需要:

「他們沒有人去,那你去吧!」

「他們沒有人做,那你做吧!」

「他們沒有錢還債,那你替他還吧!」

那是否太簡單呢?太兒戲?

當超過千萬計的群眾在捱餓?

第一種門徒方法:叫人去『自己買吃的』(15節),鼓勵自力更新!這是律法的原則:自己解決(Do It Yourself).是否明智?相信大多數現代人也會支持!門徒當時很明智啊!看見已是夜深,人在曠野,超過五千人的饑餓問題,門徒怎能解決!委實是愛莫能助!祇好每個人各自面對,把問題交回各人,雖然並非最好的處理方法,卻是許多人都會採用的處理方式.

第二種耶穌方法:「你們給他們吃罷!」這是恩典的原則,不靠自己,而是依靠白白的恩典.這是向人施恩.對許多現代人來說,非常不智,易養懶人!不過,當時耶穌看見的現實是:「天黑了!群眾既疲倦又饑餓,人數又眾多,數以千萬計的人群湧入村莊,怎能應付?甚至可能引起暴亂.這個超過五千人要開飯的大問題,實在是他們自己無法解決,別人也不能幫忙的實際難題,一定要幫他們, 幫到多少就多少.」

細心思想,這兩種想法有何不同?

第一種門徒方法:知道別人的問題與需要,鼓勵當事人自己想辦法,自己解決問題.你的老師,你的上司可能都是這樣教導你!

第二種耶穌方法:知道別人的問題與需要,先體恤他們的景況,想想自己可以怎樣幫忙?這可能是自尋麻煩的方法,令你一身蟻!

這卻是耶穌的方法!

兩種方法,也反映人面對難題的兩種態度:

門徒消極型:看到的是困難:天黑(難行),曠野(難找食物),超過五千群眾開飯問題(難以滿足),祇有五餅二魚(不夠那麼多人吃).消極型的人,看到什麼都是困難,無出路.眼前縱然並無困難,都會想像困難出來!所以,永遠都在困難中.這種人 一日到黑都是「我死啦!(廣東人),我苦啊!(閩南人)」

耶穌積極型:看到的除了困難,還有小小的五餅二魚；門徒當中,有一位有信心的安德烈(約6:1-13),也看見這個可能.只要「門徒肯給他們吃!」,就有希望!只要肯向前行,縱然艱難,或者緩慢,也會到達目的地.一位智者說『我這個人不會站著不動.動,雖然撞傷,會跌倒,只要我仍然向前行,我便會接近目的地.』當門徒說:「只有五餅二魚」,不夠啊!耶穌!耶穌說:「拿過來給我.」

「耶穌說:『拿過來給我.』」(太14:18)你們拿來給我(你要相信我),當你們肯拿來給我,當你相信我,神蹟就發生!這神蹟就是少少的變成多多.在你手上的少少,交到耶穌手裡會變成多多的.

約翰福音6:1-13記載:

門徒腓力:是現實的,理性者,計算過,就算二十兩銀子的餅(相當於當時一個人半年的工資),叫他們祇吃一點(約6:7),也無法滿足五千人.

門徒安得烈:仍把小孩「僅有的五餅二魚」帶到耶穌跟前,他是有信心者.他盡了自己的責任,把餘下的交給主.

五餅二魚是什麼?能夠餵飽五千人,可能是特別製造的?

餅是大麥餅:當代最平價的餅,被視為野獸吃的,是窮人吃的餅.魚是細小的像沙甸魚,且是腌製的.魚乾使人能吃下又乾又粗的大麥麵包.五餅二魚只夠一個大男人一餐,小孩或夠一日的糧.

許多人輕視「小小的五餅二魚」,「小小的資源」,「小小的金錢」,「小小的能力」.門徒說:「我們只有小小的能力,所以是無能為力!」耶穌卻說:「你們給他們吃,盡管只有「小小的五餅二魚」也是行!」上主看的,不是「你能給多少」,而是「你的信心有多大!」.五餅二魚,小小的食物餵飽多多的人,超過五千人,真是不可思議!小小的請求:你們給他們吃,就成就大大的事!小小的人物:門徒或小孩,卻在主內成就偉大的事.初期教會門徒集結在小小的馬可樓祈禱,結果帶來聖靈降臨的大復興!

耶穌拿起來!

於是吩咐眾人坐在草地上,就拿著這五個餅,兩條魚,望著天,祝福,擘開餅,遞給門徒,門徒又遞給眾人.(太14:19)

1「於是吩咐眾人坐在草地上

「坐」斜臥,斜靠.「坐在草地上,」得著休息,又饑又渴的群眾經過一天追隨耶穌,十分的疲倦!安息於主的腳前(參路十39).等候主的作為,現代人太急,太快要反應!人面對困難,不是像門徒所建議的「散開」,反而是聚在一起,並非離開主.耶穌不單叫他們坐,且有秩序地坐.馬可6:39節,「耶穌吩咐他們,一小組一小組坐在青草地,一排一排的,有一百一排,有五十一排的.」分排而坐,使人有秩序,有規矩,才能安靜等候,以便傳遞,不會遺漏,不會混亂,搶食!

2就拿著這五個餅,兩條魚,望著天,祝福,擘開餅,遞給門徒；門徒又遞給眾人.這個神蹟當中有五個「給」字(give or gave)

1『你們給他們吃罷』you yourselves give them something to eat (16節),耶穌鼓勵門徒要「給」!

2『拿過來給我』bring them here and give them to me(18節),首先給主使用!

3「望著天(祈禱「祝福」 原文是「謝恩」gave thanks to God),表明神是一切供應的源頭.主耶穌給父神謝恩.

4「遞給門徒」gave them to disciples

5「門徒又遞給眾人」Disciples gave them to the people

這五個「給」的次序太好了,豐滿供應的路在於此.

給與的人生:就算是小小的給,也會有奇妙的事發生.

分享的精神

擘開就是分享,是痛苦的,正如主耶穌擘開自已,分享的生命.分享食物,分享時間,分享資源.擘開是痛苦的,要犧牲,犧牲時間,金錢,奉獻會有痛的感覺.這段經文的耶穌就是如此:

耶穌聽見了,就上船從那裡獨自退到野地裡去．眾人聽見,就從各城裡步行跟隨他.13:13

耶穌出來,見有許多的人,就憐憫他們,治好了他們的病人.14:14

【耶穌聽見了】,所指的是施洗約翰(祂的先鋒,好朋友)被人斬頭的事(見前太14:1-12).在此刻,耶穌需要休息,聽到施洗約翰被殺,他的心情低落,他日以繼夜的講道醫病,十分的疲乏,不想做事,因此退到曠野.同時,耶穌需要時間祈禱,他要準備步上十字架,面對群眾,面對死亡.但他的腳蹤卻被加利利人發現,各城的人出來要聽道,見他行神蹟,趕鬼等.

【耶穌憐憫他們】他看見許多病人,許多人流離失所,許多人饑餓,沒有東西吃,許多疲倦無力的人.耶穌看見人的需要.耶穌原想退到曠野休息,祂需要休息,他也想休息；但數以千萬計的群眾卻不讓他休息,想得到他的道,他的醫治,他給予他們生命的飽足.

耶穌看見人的需要,就憐憫他們,放下自已的獨處,休息.關注群眾的需要.把他的時間給人,與人分享.今日我們大部份都是忙碌的人,需要休息的人.但我們仍需要看見周圍的人的需要,仍需要不失去憐憫人的心.許多時候,我們只看見自己的需要,自己的事情,而忽略了鄰舍的需要.

基督教的上帝觀:強調上帝是一位樂意分享,樂意給予的神,信徒要學習怎樣突破人的自私自我,願意給予與分享.整本聖經就是告訴我們,上帝如何與人分享祂的恩典與產業,也如何把自己給了我們.「神既不愛惜自己的兒子,為我們眾人捨了,豈不也把萬物和他一同白白的賜給我們麼?」(羅八 32 )許多人的人生,主要是要「取」,「得」,「搏」.基督徒的人生,卻是要效法主耶穌的「給與」,「分享」.

這神蹟的結果:

他們都喫,並且喫飽了．把剩下的零碎收拾起來,裝滿了十二個籃子.(太14:20)

喫的人除了婦女孩子, 約有五千.(太14:21)

這是一幅美麗的圖畫.人人都吃,人人都吃飽,並有廚餘12個籃子.這表示這神蹟使人人都得吃,有公平,人人都吃飽,有滿足,還有12籃子的剩餘,即有盈餘.

耶穌能使少變多,缺乏的變充足,甚至有餘.不單滿足現在,且能應付未來!耶穌顯示了人類所期待的上帝國,在這神蹟中實現了,多麼的興奮呀!上帝國是理想的國,有公平,公正,有愛,和平,滿足,平安,有人人的分享和參與等.

主耶穌在這神蹟中,顯示了一個人際的經濟定律:

主給我,我給需要的人,那需要的人又給更有需要的.主跟我分享,我跟有缺乏的人分享,那有缺乏的人又與比他更有缺乏的人分享.在這世界,人與人懂得給與分享,這世界人人都得飽足!無論是食物,財富,資源等.

面對世界有限的資源,我們需要學效耶穌的「給與」和「分享」的精神.這個世界的資源便能好好的運用.上帝的國就得到實現.

把零碎捨起來!

眾人吃飽後,耶穌吩咐門徒拾起零碎的.為什麼?因為猶太人在聚餐中通常的習慣要留下一部份食物給服事他們進餐的人.零碎載滿12個籃子,是留給12門徒的,他們都帶著籃子,是瓶形的,表示少少的會聚成多多.

猶太人外出遠行時都帶著它.籃子可放自己的食物及物件.或把乞食得到的,放在籃子裡.每個門徒都得到零碎食物在籃子裡,這是他們獎賞.我們學習不浪費,善用零碎的食物(廚餘:給飢餓人吃),零碎的金錢(像五餅二魚的奉獻),零碎的時間,用來看書,寫文章,或做有意義的事情.現代人浪費食物,浪費資源,浪費時間,浪費能源.善用零碎,不浪費是耶穌的智慧.

本段經文是神蹟嗎?

神蹟在那裡看到?

1神蹟在於用「五餅二魚」給五千人吃飽,加上婦女及兒童,可能近萬人.

2神蹟在於不單叫千萬人吃飽,還有12籃子的剩餘.

3神蹟在於耶穌拿起餅,祈禱,祝謝,擘開,傳遞下去!

4神蹟在於人人都要參與,神人協力!以往的神蹟只是耶穌一人幹,醫病,趕鬼,使瞎子看見,使痂腿的行走,使死人復活等.從耶穌而來的餅和魚,提醒我們都要學耶穌拿起,擘開,傳下去!

u這是個特別神蹟:

1特別在於不是立刻看見神蹟,不是清楚看見,是慢慢的,最後才發現超過餵飽了五千人,還有足夠的盈餘,不多不少,12籃子(受到耶穌的控制,這也是神蹟!).

2 特別在於這好像是小神蹟,又好像是大神蹟.好像是小神蹟:因這不像醫好瞎子,趕出污鬼,那麼使人驚訝!但又好像是大神蹟:

耶穌的神蹟,一對一的多,醫好一個瞎子,醫好一位血漏病人,醫好十個長大麻瘋,……今次神蹟是用五餅二魚給超過五千人食飽,還有12籃的剩餘.用五餅二魚(一個人一餐的食量) 竟然能給五千人吃飽,解決千萬人饑餓的問題,生存的問題,多大的神蹟,耶穌在世最大,最犀利的一次過最多人受惠的神蹟.

3特別在於這好像不是神蹟,又好像是神蹟!

是神蹟!耶穌為這廣大的群眾變出神蹟來,耶穌可把石頭變食物的.這時的群眾不可能都帶有食物,只可能有少部份人有帶食物,就算少部份人拿出食物,也不可能餵飽近萬人,所以這件一定是神蹟.

好像不是神蹟:可能是小孩或門徒拿出五餅二魚後,感動了許多的人也取出他們的五餅二魚,不再自私,所以人人都吃飽,耶穌並沒有做什麼,是人為的,不算是神蹟!

其實,這不是一般神蹟?

這大群人都帶著食物的,因為他們走曠野,那麼多的路.只是他們自私,不肯拿出來,直到門徒或小孩拿出「五餅二魚」,他們才受感動效法門徒.因此人人拿出,人人吃飽,還有剩餘.所以,這神蹟更深去看:是耶穌把自私自利的人改變為慷慨助人的人,這也是一種神蹟.

許多時候,我們太集中於看耶穌的神蹟:是改變環境,改變身體的缺陷；我們忽略了耶穌改變心境,改善人的思想也是一種神蹟,甚至可能是最大的神蹟!猶太人有一種看法,思想是先於一切的存在.耶和華是智慧者,思想者,所以祂是創造者.改變一個人的思想,就可能改變一群人,進而一個社會,一個國家,甚至一個時代.這改變的力量,豈不是最大的神蹟!

我認為這神蹟的看法是兩種講法都對:耶穌施展了使「五餅二魚」由少變多的能力!同時施展了「改變人心由自私變慷慨的能力」.

4這是神人協力的神蹟.耶穌把「五餅二魚」先謝恩給神,然後分給門徒,門徒分給眾人,眾人分給眾人,如此合作,產生神蹟.主耶穌有恩典和能力,門徒有生命時間,互相合作,同心協力,產生無限動力.各人都像耶穌般拿起,擘開,傳下去.擘開後,一半放入口,另一半傳下去,要顧自己吃也要顧別人吃.如果到了不能再擘開,傳下去時,那人可能拿出他的五餅二魚繼續傳下去,這種不自私的行為不斷的發生,五餅二魚餵飽超過五千人是可以發生的!

u結語

當我們細心探討,你會看見耶穌的作為,他的頭腦簡單麼?相信你早已猜中我的答案.「絕不簡單!」.如果我們認同耶穌頭腦不簡單,那些不簡單在那些地方?

面對超過五千人的饑餓難題,魔鬼要迫耶穌出招,就把曠野的石頭變作食物吧!單靠五餅二魚,怎可以餵飽五千,甚至近萬人? 耶穌在此困局,只有出招,行一個神蹟吧!那就等於聽從魔鬼的話,「把石頭變作餅」,不用上十架,不用死,也不會復活,不要顯示出耶穌是那位能從死裡復活的真神.魔鬼一定知道,耶穌剛聽到施洗約翰被人所害,突然遭斬頭.耶穌也是人,有人的弱點,就在耶穌也有心情低落的時候,把他擊敗!就在耶穌需要休息的時候,利用群眾的需要,強迫耶穌行「把石變作餅」的大神蹟,成功的話,魔鬼就勝過上帝.還有人們來找耶穌並非為聽耶穌的道,而是為了食物,信耶穌有食物,永遠崇拜肚腹,以肚腹為神,人類變成動物,甚至禽獸為食而生存,信耶穌也只為食.

此際,耶穌墮入非常的危險的景況!然而在此情況下,耶穌卻能轉危為安,大勝一場,他的頭腦絕不簡單,充滿高度智慧!

1首先,耶穌沒有把石頭變作餅,因為石頭是石頭,餅是餅；把不是餅的石頭變作餅,是違反理性,且是虛偽,把不是的變成真實,信實的上帝是不可以如此的.

2再而,耶穌是用小小的「五餅二魚」變成多多的「五餅二魚」,由少變成多,由缺乏的變成豐富有餘.這神蹟是靜靜的進行.既似神蹟又不似神蹟.使人不會單為看神蹟而來.

3耶穌更透過「五餅二魚」的神蹟來講道.從中可見他的智慧:

a他吩咐眾人坐在草地上:安靜下來

b排列好:方便傳遞

c就拿著這五個餅,兩條魚,

d望著天,祝福:給上帝謝恩

e擘開餅,遞給門徒．作榜樣

f門徒又遞給眾人.門徒照著做.

這些動作中有上帝的道.要人懂得關顧他人,信靠上帝,給與,分享的精神,不要浪費等.是一篇活潑有力,互動的講道,使人透過這神蹟而認識上帝.

4 耶穌這「五餅二魚」的神蹟,是神人協力的結果,讓眾人有機會參與,還有耶穌這神蹟不獨要改變環境,改變饑餓為飽足,更重要的「改變人心」,由自私自利變為慷慨樂於助人.不單單滿足有些人的好奇心,變戲法,或滿足肚腹.

五餅二魚的神蹟更不簡單,當中提示我們:

1不要忽略“小”的神蹟,盡所能,不看“不能“

2不要忽略“給“,不要一味“取”,”得”

3不要忽略“分享“, 不要“獨有“

4不要忽略“節儉“,不浪費

從五餅二魚的神蹟,讓我們看到耶穌是關顧人,滿有能力,且是智慧高超的主,願我們多親近主,多得著人生的智慧.阿們.



\newpage

\section{堅強與軟弱}
\label{sec:256}
\textbf{張美珍博士 - 堅強與軟弱}
\newline
\newline
\hyperref[sec:255]{< < < PREV SERMON < < <}
~
\hyperref[sec:toc]{[返目錄]}
~
\hyperref[sec:257]{> > > NEXT SERMON > > >}


列王記上 19:9-18
\newline
\begin{longtable}{cl}
\hline
\hline
章節 & 經文 (和合本修訂版)\\
\hline
19:9 & \begin{tabularx}{0.7\textwidth}{X} 他在那裡進了一個洞，在洞中過夜。看哪，耶和華的話臨到他，說：「以利亞，你在這裡做甚麼？」 \end{tabularx} \\ \\ \relax
19:10 & \begin{tabularx}{0.7\textwidth}{X} 他說：「我為耶和華－萬軍之神大發熱心，因為以色列人背棄了你的約，毀壞了你的壇，用刀殺了你的先知，只剩下我一人，他們還要追殺我。」 \end{tabularx} \\ \\ \relax
19:11 & \begin{tabularx}{0.7\textwidth}{X} 耶和華說：「你出來站在山上，在耶和華面前。」看哪，耶和華從那裡經過。在耶和華面前有烈風大作，山崩石裂，耶和華卻不在風中；風後有地震，耶和華也不在其中； \end{tabularx} \\ \\ \relax
19:12 & \begin{tabularx}{0.7\textwidth}{X} 地震後有火，耶和華也不在火中；火以後，有輕微細小的聲音。 \end{tabularx} \\ \\ \relax
19:13 & \begin{tabularx}{0.7\textwidth}{X} 以利亞聽見，就用外衣蒙臉，出來站在洞口。聽啊，有聲音向他說：「以利亞，你在這裡做甚麼？」 \end{tabularx} \\ \\ \relax
19:14 & \begin{tabularx}{0.7\textwidth}{X} 他說：「我為耶和華－萬軍之神大發熱心，因為以色列人背棄了你的約，毀壞了你的壇，用刀殺了你的先知，只剩下我一人，他們還要追殺我。」 \end{tabularx} \\ \\ \relax
19:15 & \begin{tabularx}{0.7\textwidth}{X} 耶和華對他說：「去吧，從原路回去，往大馬士革的曠野去。到了那裡，你要膏哈薛作亞蘭王， \end{tabularx} \\ \\ \relax
19:16 & \begin{tabularx}{0.7\textwidth}{X} 又膏寧示的孫子耶戶作以色列王，並膏亞伯‧米何拉人沙法的兒子以利沙作先知接續你。 \end{tabularx} \\ \\ \relax
19:17 & \begin{tabularx}{0.7\textwidth}{X} 將來逃過哈薛之刀的，必被耶戶所殺；逃過耶戶之刀的，必被以利沙所殺。 \end{tabularx} \\ \\ \relax
19:18 & \begin{tabularx}{0.7\textwidth}{X} 但我在以色列中留下七千人，是未曾向巴力屈膝，未曾親吻巴力的。」 \end{tabularx} \\ \\
[1ex]
\hline
\hline
\end{longtable}
詩篇 85:8-13
\newline
\begin{longtable}{cl}
\hline
\hline
章節 & 經文 (和合本修訂版)\\
\hline
85:8 & \begin{tabularx}{0.7\textwidth}{X} 我要聽神－耶和華所說的話，因為他必應許賜平安給他的百姓，就是他的聖民；他們卻不可再轉向愚昧。 \end{tabularx} \\ \\ \relax
85:9 & \begin{tabularx}{0.7\textwidth}{X} 他的救恩誠然與敬畏他的人相近，使榮耀住在我們的地上。 \end{tabularx} \\ \\ \relax
85:10 & \begin{tabularx}{0.7\textwidth}{X} 慈愛和誠實彼此相遇，公義與和平彼此相親。 \end{tabularx} \\ \\ \relax
85:11 & \begin{tabularx}{0.7\textwidth}{X} 誠實從地而生，公義從天而現。 \end{tabularx} \\ \\ \relax
85:12 & \begin{tabularx}{0.7\textwidth}{X} 耶和華必賜福氣給我們；我們的地也要出土產。 \end{tabularx} \\ \\ \relax
85:13 & \begin{tabularx}{0.7\textwidth}{X} 公義要行在他面前，使他的腳蹤有可走之路。 \end{tabularx} \\ \\
[1ex]
\hline
\hline
\end{longtable}
羅馬書 10:5-15
\newline
\begin{longtable}{cl}
\hline
\hline
章節 & 經文 (和合本修訂版)\\
\hline
10:5 & \begin{tabularx}{0.7\textwidth}{X} 論到出於律法的義，摩西寫著：「行這些事的人，就必因此得生。」 \end{tabularx} \\ \\ \relax
10:6 & \begin{tabularx}{0.7\textwidth}{X} 但出於信的義卻如此說：「你不要心裡說：誰要升到天上去呢？（就是說，把基督領下來。） \end{tabularx} \\ \\ \relax
10:7 & \begin{tabularx}{0.7\textwidth}{X} 或說：誰要下到陰間去呢？（就是說，把基督從死人中領上來。）」 \end{tabularx} \\ \\ \relax
10:8 & \begin{tabularx}{0.7\textwidth}{X} 他到底怎麼說呢？「這話語就離你近，就在你口中，在你心裡，」（就是說，我們傳揚所信的話語。） \end{tabularx} \\ \\ \relax
10:9 & \begin{tabularx}{0.7\textwidth}{X} 你若口裡宣認耶穌為主，心裡信神叫他從死人中復活，就必得救。 \end{tabularx} \\ \\ \relax
10:10 & \begin{tabularx}{0.7\textwidth}{X} 因為，人心裡信就可以稱義，口裡宣認就可以得救。 \end{tabularx} \\ \\ \relax
10:11 & \begin{tabularx}{0.7\textwidth}{X} 經上說：「凡信靠他的人必不蒙羞。」 \end{tabularx} \\ \\ \relax
10:12 & \begin{tabularx}{0.7\textwidth}{X} 猶太人和希臘人並沒有分別，因為人人都有同一位主，他也厚待求告他的每一個人。 \end{tabularx} \\ \\ \relax
10:13 & \begin{tabularx}{0.7\textwidth}{X} 因為「凡求告主名的就必得救」。 \end{tabularx} \\ \\ \relax
10:14 & \begin{tabularx}{0.7\textwidth}{X} 然而，人未曾信他，怎能求告他呢？未曾聽見他，怎能信他呢？沒有傳道的，怎能聽見呢？ \end{tabularx} \\ \\ \relax
10:15 & \begin{tabularx}{0.7\textwidth}{X} 若沒有奉差遣，怎能傳道呢？如經上所記：「報福音、傳喜信的人，他們的腳蹤何等佳美！」 \end{tabularx} \\ \\
[1ex]
\hline
\hline
\end{longtable}
馬太福音 14:22-33
\newline
\begin{longtable}{cl}
\hline
\hline
章節 & 經文 (和合本修訂版)\\
\hline
14:22 & \begin{tabularx}{0.7\textwidth}{X} 耶穌隨即催門徒上船，先渡到對岸，等他叫眾人散去。 \end{tabularx} \\ \\ \relax
14:23 & \begin{tabularx}{0.7\textwidth}{X} 疏散了眾人以後，他獨自上山去禱告。到了晚上，只有他一人在那裡。 \end{tabularx} \\ \\ \relax
14:24 & \begin{tabularx}{0.7\textwidth}{X} 那時船已離岸好幾里 ，因風不順，被浪顛簸。 \end{tabularx} \\ \\ \relax
14:25 & \begin{tabularx}{0.7\textwidth}{X} 天快亮的時候，耶穌在海面上走，往門徒那裡去。 \end{tabularx} \\ \\ \relax
14:26 & \begin{tabularx}{0.7\textwidth}{X} 但門徒看見他在海面上走，就驚慌了，說：「是個鬼怪！」他們害怕得喊叫起來。 \end{tabularx} \\ \\ \relax
14:27 & \begin{tabularx}{0.7\textwidth}{X} 耶穌連忙對他們說：「放心！是我，不要怕！」 \end{tabularx} \\ \\ \relax
14:28 & \begin{tabularx}{0.7\textwidth}{X} 彼得回答他說：「主啊，如果是你，請叫我從水面上走到你那裡去。」 \end{tabularx} \\ \\ \relax
14:29 & \begin{tabularx}{0.7\textwidth}{X} 耶穌說：「你來吧！」彼得就從船上下去，在水面上走，往耶穌那裡去； \end{tabularx} \\ \\ \relax
14:30 & \begin{tabularx}{0.7\textwidth}{X} 只因見風很強，害怕起來，將要沉下去，就喊著說：「主啊，救我！」 \end{tabularx} \\ \\ \relax
14:31 & \begin{tabularx}{0.7\textwidth}{X} 耶穌立刻伸手拉住他，說：「你這小信的人哪，為甚麼疑惑呢？」 \end{tabularx} \\ \\ \relax
14:32 & \begin{tabularx}{0.7\textwidth}{X} 他們一上船，風就停了。 \end{tabularx} \\ \\ \relax
14:33 & \begin{tabularx}{0.7\textwidth}{X} 在船上的人都拜他，說：「你真是神的兒子。」 \end{tabularx} \\ \\
[1ex]
\hline
\hline
\end{longtable}
講題:堅強與軟弱　Strong vs Weak

經課:列王記上19:9-18;詩篇85:8-13;羅馬書10:5-15;馬太福音14:22-33

講員:張美珍博士

場合:香港中文大學崇基學院禮拜堂

日期:2014年8月10日

各位弟兄姊妹,主內平安.

請問各位:你認為自己是一個堅強的人,還是軟弱的人呢?

有沒有人在孩童時代很堅強,但長大後卻變得軟弱?抑或在求學時期很軟弱,但長大後出來工作卻顯得很堅強?又或者, 上班的時候很堅強,下了班卻變得軟弱?在朋友面前表現得很堅強,但對著太太或丈夫又變得很軟弱?

堅強或軟弱的性格,是與生俱來,還是後天造成?

上述這些疑問,其實早在十多年前,當我還在加拿大卡加利(Calgary, Canada)唸書時曾經出現過的.一直在我心目中非常堅強的一位至親,在移民卡加利不到一年竟然變得很軟弱.我心中的疑團沒法解開,適逢那時我為當地一份華人報章兼職寫《移民故事》專欄,我便以這位至親的故事為題材,寫了一篇文章,採用的題目正是《堅強與軟弱》.想不到,文章刊登後,在卡加利華人社群竟然引起不少迴響.有些移民致電報社約我見面,與我分享他們的經歷(當中有來自國內弟兄姊妹醫學研究人才,也有台灣美麗島事件的政治人物),企圖為我解開這個疑團,結果他們的故事成為了我《堅強與軟弱》一文的續篇.當時的結論是,堅強與軟弱並非先天的性格使然,它是一種內心的感覺或狀態,亦可以說是一種情緒反應,在不同的時空下交替出現.

令我想起十多年前這樁小事,乃因今次在預備講章時,在今天的經課中,看見先知以利亞的經歷,再度勾起我對「堅強與軟弱」的反思,並有了新的體會.以利亞是上帝所重用的先知,在很多的處境下都顯得很堅強；然而,今天的經課卻描繪了他軟弱的一面.

王上十七章一節首次出現他的名字──「基列寄居的提斯比人以利亞」.他的身世頗傳奇,背景沒人知曉；翻查聖經地圖,「提斯比」這個地方都無法找到.這位彷彿從天而降的人物,在其後兩章的經文,記載了他如何靠著耶和神上帝的名呼風喚雨,令死人復活,又令火從天降.

他跑去對當時以色列的王亞哈(他拜偶像,又行耶和華眼中看為惡的事,娶了拜巴力的耶洗別為妻)說:「這幾年我若不禱告,必不降露不下雨.」(王上十七:1)說完這話,就照耶和華的吩咐,去住在約但河東邊的基立溪旁,喝那裡的水,有烏鴉早晚給他叼餅和肉來給他吃,直至基立溪的水乾了,又照耶和華的吩咐往西頓的撒勒法去投靠一個寡婦.

到了那個寡婦家,才發現原來這家庭非常貧窮,只剩一把麵及一點油,食完這一餐都不知何時才有下一餐吃.但以利亞吩咐寡婦給他作一個小餅,然後再為自己及兒子作餅,罈內的麵及瓶內的油一直源源不斷地供應他們,沒有短缺.他們三人就靠這些過日子,直到雨降地上令糧食出產為止.在這期間,以利亞又令那婦人的兒子從死裡復活.

三年後,以利亞再見亞哈,相約在迦密山比試,看先知的真偽.巴力的先知以牛犢為祭,多次發功,甚至自殘,天上都沒有反應,但以利亞一祈禱,火便從天降下,燒盡祭壇上的一切,令圍觀民眾俯伏在地,敬拜耶和華.以利亞趁機叫民眾捉住巴力的先知(有四百五十人),到基順河邊把他們殺掉；之後再上迦密山喚雨,結束旱災及饑荒.

當亞哈回家,將一切告訴妻子耶洗別後,耶洗別發出了追殺令,以利亞逃命,去了別是巴,在曠野走了一日的路程,在一棵羅騰樹下尋死.有天使供應他水及餅,再走了四十晝夜,來到何烈山(何烈山又名叫西乃山,這兩個名字常常互用；摩西頒十誡就在這個山上.)

在王上19:9-18,即今日的舊約經課,一開始便描繪以利亞在何烈山上,因極度害怕而躲在山洞中的經歷.如果大家看過之前兩章的記載,真的無法解釋為何一個為奉耶和華之名行了那麼多神蹟,得上帝如此重用的先知,會突然變得那麼軟弱可憐.

他聽到耶和華以微小的聲音問他:「你在這裡作什麼?」以利亞這樣回答:「我為耶和華萬軍之神大發熱心.因為以色列人背棄了你的約,毀壞了你的壇,用刀殺了你的先知,只剩下我一個人,他們還要尋索我的命.」這聲音出現過兩次,而兩次以利亞的答案都一樣.很明顯,他的內心正是被這思慮纏擾.以利亞是因為怕死而變得軟弱嗎?當然不是；若他怕死,就不會只走了一天的路程,就要在羅騰樹下尋死.

那是什麼原因令他變得軟弱?大家有沒有發現,在他兩次重覆的說話中,反映了他的思維,而他的思維,反映了他的假設──他為耶和華大發熱心,不應得到這個結局──只剩下他一人,而且正被人追殺.他發晦氣要尋死,感到不開心,最後躲了起來,是因為事情的發展,與他的期望出現了很大的落差,令他情緒低落.

今日我們環顧身邊,又或者打開報章,不難發現有很多人與以利亞這個經驗很相似.很多在熱戀中的男女,事業有成的人,成績不錯的學生,當遇到挫折時,或失戀,或失業,或考試失手,竟然會倒地不起,甚至自尋短見,實在令人感到惋息及難過.這些可能在別人眼中曾經堅強的人,為何會變得軟弱,無力面對生活?

相信在座中不少弟兄姊妹,都曾遇過事與願違的情況,你的反應是如何?古語有云:「不如意事十常八九」,當你發現事情的發生,並不如你主觀的想像(或願望,或假設)般出現時,你會變得軟弱,發晦氣,逃避,放棄?還是變得越來越堅強,積極承受及面對事情所給你的挑戰?不要少看自己的假設,我認識不少朋友,因為自己一廂情願的假設而令婚姻破裂,工作碰壁,朋友反臉,也有信徒因此離開上帝…….

以利亞的經驗,對我們信徒來說,是一個很重要的信仰經驗:以利亞這般堅強的人,都會有軟弱的時刻,更何況我們這些平凡人呢!這種軟弱感讓我們知道自己無論自己多能幹,都無法控制一切所發生的事,亦無法掌管明天.這讓我們意識到,人,原來不過很渺少,很脆弱,亦沒有什麼值得誇口.在我們最軟弱的時刻,原來就是讓我們最接近上帝的時刻.

以利亞最後如何重新振作起來呢?乃因上帝親自尋找他,讓他細訴心裡的恐懼,用祂的大能堅固他的信心,又讓他知道是他的思維出了問題,粉碎了他的假設──

1．以利亞以為上帝總是以大能者的姿態出現,不明白自己為何為上帝大發熱心後,結局竟然是給人追殺.弟兄姊妹們,大家有沒有相似的經驗?有沒有問上帝,我如此熱心為你,為何在職場上獲得升職的不是我?為何考試成績未如理想?然而,上帝真的每次都以大能者的姿態出現嗎?「耶和華說,你出來站在山上,在我面前.那時耶和華從那裡經過,在他面前有烈風大作,崩山碎石,耶和華卻不在風中；風後地震,耶和華卻不在其中；地震後有火,耶和華也不在火中；火後有微小的聲音.」(王上十九:11-12)耶和華有時刻意用不同的姿態出現,考驗我們對祂的認識及信心.

2．以利亞以為自己是惟一的倖存者,孤軍作戰:「我為耶和華萬軍之神大發熱心.因為以色列人背棄了你的約,毀壞了你的壇,用刀殺了你的先知,只剩下我一個人,他們還要尋索我的命.」但耶和華告訴他:「我為自己留下七千人,是未曾向巴力屈膝的.」(王上十九:18)祂提醒以利亞:不願與世俗之王妥協的,尚有很多人；你,並不孤單.

以利亞的思維得到改變後重新上路,耶和華亦為他預備了一位拍檔,即是他徒弟以利沙,與他一起並肩作戰.最後以利亞不但沒有被耶洗別所殺,反而他預言亞哈及耶洗別兩夫婦因作惡多端而不得好死卻應驗了,之後以利亞更沒有經過死亡,乘旋風昇天.

以利亞的經驗告訴我們,最軟弱的時刻,亦是上帝恩典最豐富的時刻.使徒保羅亦提到他曾三次祈求上主為他拔走肉體上的一根刺(林後十二:7-8).但主耶穌對他說:「我的恩典夠你用的,因為我的能力是在人的軟弱上顯得完全.」所以保羅更喜歡誇自己的軟弱,因他什麼時候軟弱,什麼時候就剛強了(林後十二:9).

當我們軟弱的時候,別忘記上帝恩典夠我們用的.還有,當我們以為自己孤軍作戰之時,別忘上帝會為我們預備同行者,就好像我們回來崇基禮拜堂不單是為敬拜──根據約翰福音四章24節所說:「上帝是個靈,所以拜祂的,要用心靈和誠實來拜祂」,敬拜上帝不一定局限在教會,但大家一起回來這裡崇拜,就是讓我們大家都知道,在這世上願意一起追求並實現上帝國的,大有人在.你和我,並不孤單.



\newpage

\section{驕傲與偏見}
\label{sec:257}
\textbf{賀志勇博士 - 驕傲與偏見}
\newline
\newline
\hyperref[sec:256]{< < < PREV SERMON < < <}
~
\hyperref[sec:toc]{[返目錄]}
~
\hyperref[sec:258]{> > > NEXT SERMON > > >}


馬太福音 15:21-28
\newline
\begin{longtable}{cl}
\hline
\hline
章節 & 經文 (和合本修訂版)\\
\hline
15:21 & \begin{tabularx}{0.7\textwidth}{X} 耶穌離開那裡，退到推羅、西頓境內。 \end{tabularx} \\ \\ \relax
15:22 & \begin{tabularx}{0.7\textwidth}{X} 有一個迦南婦人從那地方出來，喊著說：「主啊，大衛之子，可憐我！我女兒被鬼纏得很苦。」 \end{tabularx} \\ \\ \relax
15:23 & \begin{tabularx}{0.7\textwidth}{X} 耶穌卻一言不答。門徒進前來，求他說：「這婦人在我們後頭喊叫，請打發她走吧。」 \end{tabularx} \\ \\ \relax
15:24 & \begin{tabularx}{0.7\textwidth}{X} 耶穌回答：「我奉差遣只到以色列家迷失的羊那裡去。」 \end{tabularx} \\ \\ \relax
15:25 & \begin{tabularx}{0.7\textwidth}{X} 那婦人來拜他，說：「主啊，幫幫我！」 \end{tabularx} \\ \\ \relax
15:26 & \begin{tabularx}{0.7\textwidth}{X} 他回答：「拿孩子的餅丟給小狗吃是不妥的。」 \end{tabularx} \\ \\ \relax
15:27 & \begin{tabularx}{0.7\textwidth}{X} 婦人說：「主啊，不錯，可是小狗也吃牠主人桌上掉下來的碎屑。」 \end{tabularx} \\ \\ \relax
15:28 & \begin{tabularx}{0.7\textwidth}{X} 於是耶穌回答她說：「婦人，你的信心很大！照你所要的成全你吧。」從那時起，她的女兒就好了。 \end{tabularx} \\ \\
[1ex]
\hline
\hline
\end{longtable}
講題:驕傲與偏見　Pride and Prejudice

講員:賀志勇博士

經文:馬太福音15:21-28

場合:香港中文大學崇基學院禮拜堂

日期:2014年8月17日

馬太福音15章21-28節  傲慢與偏見

15:21 耶穌離開那裡,退到推羅西頓的境內去.  15:22 有一個迦南婦人,從那地方出來,喊著說,主阿,大衛的子孫,可憐我．我女兒被鬼附得甚苦.

15:23 耶穌卻一言不答.門徒進前來,求他說,這婦人在我們後頭喊叫．請打發他走吧.15:24 耶穌說,我奉差遣,不過是到以色列家迷失的羊那裡去.

15:25 那婦人來拜他,說,主阿,幫助我. 15:26 他回答說,不好拿兒女的餅,丟給狗吃. 15:27 婦人說,主阿,不錯．但是狗也吃他主人桌子上掉下來的碎渣兒. 15:28 耶穌說,婦人,你的信心是大的．照你所要的,給你成全了吧.從那時候,他女兒就好了.

一,引言  傳統的缺點

印度有活人殉葬的傳統,就是丈夫死了之後,妻子要自焚陪葬,如果妻子不肯,就會被人逼著去自焚.英國人佔領印度的時候,覺得這個傳統太不人道,決定要廢除這個傳統.英國總督就找來印度教貴族,跟他們講:“我們不能允許逼女人陪葬這種事,這是殺人!我們不允許殺人,我們的傳統是誰殺人,我們就處死誰.”印度教貴族當然不願意放棄,他們說:“但陪葬是我們的傳統,不可以改變.”於是英國總督說:“好吧!那你們堅持你們的傳統,我們堅持我們的傳統,大家走著瞧!”最後印度教貴族怕死,放棄了活人陪葬的傳統.

每個人都活在傳統當中,傳統塑造了我們的身份和特點.就以食物來講,我有個香港朋友,吃過很多地方的面,最後感歎說:“吃了這麼多種面,還是我們香港的鮮蝦雲吞面最好吃.”傳統這個東西有許多優點,比方說,中華民族的傳統就是熱情好客˴勤勞善良啊；不過話又說回來,傳統可能最突出的兩個缺點就是:第一,偏見；第二,傲慢.從剛才的故事來看,印度人的殉葬傳統代表著醜陋的偏見--女人就是男人的財產,可以任意處置；而英國人的傳統代表著武力的傲慢--你不聽話,我就弄死你.這兩種特點,至今還源源不斷的發揮破壞力,最近的巴以衝突傷亡已經過萬,就是很明顯的例子了.

今天要講的經文,馬太福音15章21-28節,就是一段有關傳統的故事,它講一個迦南女人哀求耶穌治療她女兒被鬼附的病,但耶穌因為她不是猶太人,所以起先不肯醫治(這是猶太人的傳統),而且說不要拿兒女的餅喂狗這樣的話刺激她,但這個女人說狗也吃餅渣啊,最後結果很戲劇性,耶穌反而誇獎這個女人信心大,幫她女兒醫治了.    就這樣看這個故事,這跟傳統有什麼關係呢?要明白這個故事,讓我們來看看這個故事之前發生了什麼事.

二,解經

1,故事之前的經文 15:1-20節

馬太15章1-20節,講一些法利賽人和文士,特意從耶路撒冷大老遠跑來見耶穌,指責他的門徒違反了傳統--吃飯之前不洗手.耶穌回敬他們說“吃進口的食物不能汙穢人,但是從口裡˴心裡發出的東西才會汙穢人”,批評法利賽人是借著傳統來違反神的誡命,是瞎子領瞎子.

問題在於“吃飯之前洗手”是不是神的誡命?如果是,那耶穌豈不是公然教人違反誡命麼?於是我們把摩西五經翻來翻去,發現律法書裡面並沒有這一條.利未記裡提到潔淨的禮儀,是指人如果被汙穢之後,一定要潔淨；但並沒有說人在飯前一定要洗手.再追蹤一下,原來飯前洗手這一條出現在“米示拿”裡,是猶太拉比對於律法的解釋.這就跟人大釋法一樣,只可能解釋得更緊˴更嚴格,不可能更寬鬆.拉比解釋“潔淨律法”這一條,就加了不少自己更緊的地方,這就演變成了猶太人洗手的傳統.

在這一段中,耶穌強調的是:真正敬拜神,乃是內心中出於愛的敬拜,而不僅僅是出於宗教傳統的敬拜. 其實飯前洗手並沒有什麼不好,起碼講衛生啊.這個傳統沒什麼不好的,但如果傳統不是出於內心對神˴對人的愛,而是出於人為的“律法加強版”,那傳統就變成了壓制人的藉口,就成了偏見,讓人的生命不得自由.

2,21-24節 猶太人與外邦人

接下去,在21-24節中,我們發現耶穌第一次,也是唯一一次到了外邦地區--推羅˴西頓,這是很違反傳統的事,而且這一次還遇到了一個求助的迦南女人,這更有點奇怪了.熟悉舊約歷史的人就知道,當年約書亞的任務就是要滅掉迦南人,而約書亞這個名字是希伯來文發音,如果用希臘文發音就是“耶穌”.迦南人跟以色列人本是世代的仇敵,但是今天,以色列人的敵人居然來到主的面前,還大喊道耶穌“主啊,大衛的子孫,請可憐我吧.”當法利賽人死守傳統,拒絕主耶穌的時候；以色列的敵人反而認出耶穌是主,這是不是有趣呢? 故事到這裡卻有了個奇怪的轉折,主耶穌似乎並沒有感動,他反而開始守傳統了,他說:“我奉差遣,不過是到以色列家迷失的羊那裡去.”

有時候,傳統會表現得非常冷冰冰--因為迦南人是仇敵,因為這個婦人不是猶太人,所以根據傳統她不能得到恩典和憐憫.我們的社會是不是也有這樣的傳統呢?比方有人堅持“父母都不是香港人,他們在香港生的孩子不應該享受香港的恩典和福利”；還有中國幾千年重男輕女的傳統,是不是都有這樣的特點呢?但在這裡,請不要怪主耶穌堅持傳統,因為他深知道在那個時候他自己的使命是什麼--向以色列人傳福音.在這種對話氣氛中,這故事該如何推動下去呢?

3,25-28節  狗與碎渣

在25-28節,主耶穌再次重申了傳統,他說:“不好拿兒女的餅丟給狗吃,”這句話會讓人很不舒服,這不是罵人是狗麼?無論在中國文化還是當時的文化,說人家是狗都很不尊重.主耶穌怎麼可以如此無禮,他不是一向謙卑溫柔的麼?士可殺不可辱.但這個女人居然繼續請求說:“主阿,不錯．但是狗也吃他主人桌子上掉下來的碎渣兒.”哎呀,何等的作踐自己,何等的丟臉啊!

如果我們這樣來看這段經文,我們會覺得氣憤,起碼會覺得很怪異--主耶穌怎麼會這樣講話呢?但如果對照馬可福音7章27節,我們發現主耶穌說的是:“讓兒女們先吃飽,不好拿兒女的餅丟給狗吃.”這句話如果仔細再聽,我們發現裡面藏有玄機,第一,主耶穌強調的是讓兒女先吃飽,但並沒有說之後一點不能分給“狗”吃；第二,這裡的狗字,其實是“小狗”的意思 ,不是外面到處亂跑的野狗,而是自己家裡養的小狗.

如果根據我們的文化傳統,我們覺得主耶穌的話在罵人,是傲慢,是侮辱,是拒絕；因為我們的傳統是不能丟面子,勿衰得,勿輸得,所以我們不會這麼低聲下氣去求人.但這個女人卻能夠聽出主耶穌的言外之意,繼續祈求,這是為什麼?是不是因為她已經絕望了,是不是因為她知道她唯一的希望就是主耶穌? 事實也是如此.當我們覺得自己還有點什麼,自己還可以搞定的時候,我們當然無法謙卑俯伏在主耶穌面前去祈求,所以我們也聽不明白主的話.真正的謙卑,來自我們意識到其實我們沒有別的路可以走,“貧窮的人有福了,因為天國是他們的.” 這個迦南女人正是意識到自己已經無路可走了,所以她繼續謙卑,也毫不退縮地跟主說:“是的,主啊,但狗也吃主人桌子上掉下的碎渣兒.”當她這樣祈求,主就成全了她的信心.

每次我讀經讀到這一段,我心裡就覺得非常感動.我想起自己信主的經過,何嘗不是這樣俯伏在主的面前,才得憐恤,蒙恩典的!我本科學的是電子自動化,當時覺得一輩子做電子,生命沒啥意思；為了找到些意思,我就改學了哲學,結果哲學沒有解決我的生命問題,反而帶來了更多問題,於是我又去學基督教.當時去查聖經,主要是去找麻煩的；想用科學來教育一下基督徒,告訴他們信這本聖經有多荒謬啊!查經查了一年半,每個星期都去吃教會的免費午餐,但是人家問:“你來了這麼久,信了麼?”我都搖頭:“真不好意思,我還是信不了.”直到有一次,我的好朋友從浸會大學20層樓上墜下來,一下子,我的生命也失重了,因為我想不通,像他這麼優秀的一個人,怎麼一下子就死了呢?在空虛與傷痛了幾個月之後,有一天,我俯伏在主面前,承認自己有罪,才真正得到了釋放.從那時起,我再看聖經,發現以前很多的疑惑不再是疑惑,主的話語,終於能一點一點看懂了.當我意識到,其實我一無所有的時候,我才發現,我的生命有了一個新的起點.我從前的傳統,慢慢地被聖經的傳統影響˴改變；我從前生命中的驕傲與偏見,慢慢地被融化,我的生命也凝結成新的樣式.

三,結語

回頭來看,這兩段經文都提到了食物--飲食口味是人最難磨滅的傳統,我到現在還喜歡吃麻辣口味.接著,馬太由食物再連接到了猶太傳統--飯前要洗手啊,恩典不能流向外邦人啊.猶太人非常守傳統,有一次去以色列,我們的導遊是一位師母,她說她去一個猶太朋友家,那家的男人看到她來了,都面朝牆壁,屁股對著她,因為她不是本家的女人,他們不願意看到她的臉.

屁股對著人家的傳統無可厚非,但傳統如果演變成了一種偏見,或者是一種傲慢,那就危害大了.耶穌的偉大正在於此,他不是要毀滅傳統,他只是要改變傳統中的形式主義,因為他深深明白傳統形式的背後應該承載什麼樣的內容.形式不是最重要的,重要的是裡面的內容.所以洗手不洗手不是最重要的,身份也不是最重要的,而一個罪人悔改認信,得到愛與憐憫才是最重要的.主耶穌深深明白這個道理:一成不變只會流於狹隘,於是他勇敢地跨越了傳統.

方舟子,一個內地名人曾經這樣批評過這段經文,他說:“我們無需去當吃碎渣兒的狗.”他之所以這樣說,是因為他的傳統的偏見.他不明白,那些碎渣並不是可有可無的東西,那其實就是主耶穌的身體；為了我們,主耶穌自己成了碎渣.我們今天的聖餐禮,不是一同擘餅,一同領受這些碎渣麼?但是,聖經的話“向聰明通達的人就藏起來,向嬰孩就顯出來”,因為聰明通達的人,很容易就滿有傲慢與偏見.

領受恩典的是那些像嬰孩一樣的人,就像那個謙卑但又風趣的迦南婦人.表面上看,是她的女兒得了醫治,但這婦人豈不是一樣麼?有誰感受到這個婦人的焦慮,恐懼,擔憂與絕望呢?可是當她鍥而不捨地俯伏在主面前時,主耶穌跟她開了個小玩笑,卻最終大大地成全了她,也醫治了她.經歷過主醫治的人會明白那種感受,主的愛帶來不僅僅是身體的醫治,更帶來了生命的醫治.就像現在一樣,主已經宣告:死亡不是我們最終的結局,我們將掙脫死亡的毒鉤,進入更美好的家鄉.

弟兄姊妹,我們能不能夠俯伏在主的跟前,跟主說:“主啊,分一點您的碎渣給我吧.”如果不打破我們過往世俗氣的傳統,那恩典 又如何能臨到呢?



\newpage

\section{有不認識約瑟的新王起來}
\label{sec:258}
\textbf{林豪恩先生 - 有不認識約瑟的新王起來}
\newline
\newline
\hyperref[sec:257]{< < < PREV SERMON < < <}
~
\hyperref[sec:toc]{[返目錄]}
~
\hyperref[sec:259]{> > > NEXT SERMON > > >}


出埃及記 1:8-2:10
\newline
\begin{longtable}{cl}
\hline
\hline
章節 & 經文 (和合本修訂版)\\
\hline
1:8 & \begin{tabularx}{0.7\textwidth}{X} 有一位不認識約瑟的新王興起，統治埃及。 \end{tabularx} \\ \\ \relax
1:9 & \begin{tabularx}{0.7\textwidth}{X} 他對自己的百姓說：「看哪，以色列人的百姓比我們還多，又比我們強盛。 \end{tabularx} \\ \\ \relax
1:10 & \begin{tabularx}{0.7\textwidth}{X} 來吧，讓我們機巧地待他們，恐怕他們增多起來，將來若有戰爭，他們就聯合我們的仇敵來攻擊我們，然後離開這地去了。」 \end{tabularx} \\ \\ \relax
1:11 & \begin{tabularx}{0.7\textwidth}{X} 於是埃及人派監工管轄他們，用勞役苦待他們。他們為法老建造儲貨城，就是比東和蘭塞。 \end{tabularx} \\ \\ \relax
1:12 & \begin{tabularx}{0.7\textwidth}{X} 可是越苦待他們，他們就越發增多，更加繁衍，埃及人就因以色列人愁煩。 \end{tabularx} \\ \\ \relax
1:13 & \begin{tabularx}{0.7\textwidth}{X} 埃及人嚴厲地強迫以色列人做工， \end{tabularx} \\ \\ \relax
1:14 & \begin{tabularx}{0.7\textwidth}{X} 使他們因苦工而生活痛苦；無論是和泥，是做磚，是做田間各樣的工，一切的工埃及人都嚴厲地對待他們。 \end{tabularx} \\ \\ \relax
1:15 & \begin{tabularx}{0.7\textwidth}{X} 埃及王又對希伯來的接生婆，一個名叫施弗拉，另一個名叫普阿的說： \end{tabularx} \\ \\ \relax
1:16 & \begin{tabularx}{0.7\textwidth}{X} 「你們為希伯來婦人接生，臨盆的時候要注意，若是男的，就把他殺了，若是女的，就讓她活。」 \end{tabularx} \\ \\ \relax
1:17 & \begin{tabularx}{0.7\textwidth}{X} 但是接生婆敬畏神，不照埃及王的吩咐去做，卻讓男孩活著。 \end{tabularx} \\ \\ \relax
1:18 & \begin{tabularx}{0.7\textwidth}{X} 埃及王召了接生婆來，對她們說：「你們為甚麼做這事，讓男孩活著呢？」 \end{tabularx} \\ \\ \relax
1:19 & \begin{tabularx}{0.7\textwidth}{X} 接生婆對法老說：「因為希伯來婦人與埃及婦人不同；希伯來婦人健壯，接生婆還沒有到，她們已經生產了。」 \end{tabularx} \\ \\ \relax
1:20 & \begin{tabularx}{0.7\textwidth}{X} 神恩待接生婆；以色列人增多起來，極其強盛。 \end{tabularx} \\ \\ \relax
1:21 & \begin{tabularx}{0.7\textwidth}{X} 接生婆因為敬畏神，神就叫她們成立家室。 \end{tabularx} \\ \\ \relax
1:22 & \begin{tabularx}{0.7\textwidth}{X} 法老吩咐他的眾百姓說：「把所生的每一個男孩都丟到尼羅河裡去，讓所有的女孩存活。」 \end{tabularx} \\ \\ \relax
2:1 & \begin{tabularx}{0.7\textwidth}{X} 有一個利未家的人娶了一個利未女子為妻。 \end{tabularx} \\ \\ \relax
2:2 & \begin{tabularx}{0.7\textwidth}{X} 那女人懷孕，生了一個兒子，見他俊美，就把他藏了三個月， \end{tabularx} \\ \\ \relax
2:3 & \begin{tabularx}{0.7\textwidth}{X} 後來不能再藏，就取了一個蒲草箱，抹上柏油和樹脂，將孩子放在裡面，把箱子擱在尼羅河邊的蘆葦中。 \end{tabularx} \\ \\ \relax
2:4 & \begin{tabularx}{0.7\textwidth}{X} 孩子的姊姊遠遠站著，要知道他究竟會怎樣。 \end{tabularx} \\ \\ \relax
2:5 & \begin{tabularx}{0.7\textwidth}{X} 法老的女兒來到尼羅河邊洗澡，她的女僕們在河邊行走。她看見在蘆葦中的箱子，就派一個使女把它拿來。 \end{tabularx} \\ \\ \relax
2:6 & \begin{tabularx}{0.7\textwidth}{X} 她打開箱子，看見那孩子。看哪，男孩在哭，她就可憐他，說：「這是希伯來人的一個孩子。」 \end{tabularx} \\ \\ \relax
2:7 & \begin{tabularx}{0.7\textwidth}{X} 孩子的姊姊對法老的女兒說：「我去叫一個希伯來婦人來作奶媽，替你乳養這孩子，好嗎？」 \end{tabularx} \\ \\ \relax
2:8 & \begin{tabularx}{0.7\textwidth}{X} 法老的女兒對她說：「去吧！」那女孩就去叫了孩子的母親來。 \end{tabularx} \\ \\ \relax
2:9 & \begin{tabularx}{0.7\textwidth}{X} 法老的女兒對她說：「你把這孩子抱去，替我乳養這孩子，我必給你工錢。」那婦人就把孩子接過來，乳養他。 \end{tabularx} \\ \\ \relax
2:10 & \begin{tabularx}{0.7\textwidth}{X} 孩子長大了，婦人把他帶到法老的女兒那裡，就作了她的兒子。她給孩子起名叫摩西，說：「因我把他從水裡拉出來。」 \end{tabularx} \\ \\
[1ex]
\hline
\hline
\end{longtable}
講題:有不認識約瑟的新王興起　A New King Arose Who Did Not Know Joseph

經文:出埃及記1:8-2:10

講員:林豪恩先生

場合:香港中文大學崇基學院禮拜堂

日期:2014年8月24日

揭幕:新時代的開始

1:8有不認識約瑟的新王起來,治理埃及,

故事的首句-「有不認識約瑟的新王起來,治理埃及」,標示著一個新時代的開始.古人說:「一朝天子一朝臣」,權力出現交替,權力的保護傘和關係網也會隨之而改變.又有人說:「新人事新作風」,新當權者對事情的判斷,施政理念和施政方法,都有可能改弦易轍,推倒重來.試想想,如果你是約瑟的後人,現在「有不認識約瑟的新王起來,治理埃及」,你會有甚麼反應呢?相信很多人會心感不妙,靈敏的會嗅到危機的味道.如果有選擇的空間,有人可能會考慮移民他方.大家都知道,新王「不認識約瑟」,不是「不知道有約瑟」這個人那麼簡單.「不認識約瑟」,其實代表著新王否認埃及和約瑟的關係,改變埃及對約瑟後人的政策.例如:有一位富翁百年歸老,富翁的兄弟來向富翁的兒子求助,如果富翁的兒子對這些叔伯說:「我不認識你們」,大家就明白是甚麼意思了.「有不認識約瑟的新王起來,治理埃及」,將會是一個怎樣的時代呢?約瑟的後人將會有甚麼遭遇呢?他們又會怎樣面對呢?

在繼續閱讀下去之前,大家不妨想一想,我們個人的際遇,城市的際遇,甚或是世界上不同地方教會的際遇,會否曾經或者正在面對「有不認識約瑟的新王起來,治理埃及」這不妙的狀況呢?閱讀這個故事的時候,歡迎發揮聯想,對號入座.

一.  看見你對我的威脅

1:9對他的百姓說:「看哪,這以色列民比我們還多,又比我們強盛.

1:10來吧,我們不如用巧計待他們,恐怕他們多起來,日後若遇甚麼爭戰的事,就連合我們的仇敵攻擊我們,離開這地去了.」

這位不認識約瑟的新王,他看約瑟後代的眼光和先王不同了.他選擇「看」一項而「不看」兩項.第一,他看約瑟的家人,即是經文所說的以色列人為「威脅」.這位新王創造了「在埃及的以色列人威脅論」,把這論述提升至民族層面:「看哪,這以色列民比我們還多,又比我們強盛」.這樣的論述煽動埃及人的民族情緒,製造群族對立.同時,他把「以色列人威脅論」提升至國家安全層面:「恐怕他們多起來,日後若遇甚麼爭戰的事,就連合我們的仇敵攻擊我們」.新王妖魔化約瑟後人的形象,把他們描繪為裡通外敵的危險人物.其實,約瑟的後人有否對新王構成威脅呢?他們能否威脅埃及的國家安全呢?從這位新王的說話可知事實未必如此.第10節說:「來吧,我們不如用巧計待他們,恐怕他們多起來,日後若遇甚麼爭戰的事,就連合我們的仇敵攻擊我們,離開這地去了.」新王關注的,是約瑟的後人「離開這地去了」.為甚麼不容讓他們「離開這地去」呢?如果約瑟的後人離開,埃及就失去一大批廉價勞工.這位不認識約瑟的新王是聰明的,他製造「以色列人威脅論」以爭取埃及人支持他改變對待約瑟後人的政策.

二.  不看你過去的貢獻

新王看約瑟後人為「威脅」,同時卻掩面不看另外兩項.第一,他掩面不看約瑟過去對埃及的貢獻.創世記41章記載,約瑟曾經是埃及的恩人,幫助埃及渡過多年的飢荒.故事是這樣的:

41:1過了兩年,法老做夢,

41:8到了早晨,法老心裡不安,就差人召了埃及所有的術士和博士來；法老就把所做的夢告訴他們,卻沒有人能給法老圓解.

41:25約瑟對法老說:「法老的夢乃是一個.上帝已將所要做的事指示法老了.

41:26七隻好母牛是七年,七個好穗子也是七年；這夢乃是一個.

41:27那隨後上來的七隻又乾瘦又醜陋的母牛是七年,那七個虛空,被東風吹焦的穗子也是七年,都是七個荒年.

41:28這就是我對法老所說,上帝已將所要做的事顯明給法老了.

41:29埃及遍地必來七個大豐年,

41:30隨後又要來七個荒年,甚至在埃及地都忘了先前的豐收,全地必被饑荒所滅.

41:31因那以後的饑荒甚大,便不覺得先前的豐收了.

41:34法老當這樣行,又派官員管理這地.當七個豐年的時候,征收埃及地的五分之一,

41:35叫他們把將來豐年一切的糧食聚斂起來,積蓄五穀,收存在各城裡做食物,歸於法老的手下.

41:36所積蓄的糧食可以防備埃及地將來的七個荒年,免得這地被饑荒所滅.」

41:37法老和他一切臣僕都以這事為妙.

41:39法老對約瑟說:「上帝既將這事都指示你,可見沒有人像你這樣有聰明有智慧.

41:40你可以掌管我的家；我的民都必聽從你的話.惟獨在寶座上我比你大.」

41:41法老又對約瑟說:「我派你治理埃及全地.」

41:48約瑟聚斂埃及地七個豐年一切的糧食,把糧食積存在各城裡；各城周圍田地的糧食都積存在本城裡.

41:53埃及地的七個豐年一完,

41:54七個荒年就來了.正如約瑟所說的,各地都有饑荒；惟獨埃及全地有糧食.

41:55及至埃及全地有了饑荒,眾民向法老哀求糧食,法老對他們說:「你們往約瑟那裡去,凡他所說的,你們都要做.」

41:56當時饑荒遍滿天下,約瑟開了各處的倉,糶糧給埃及人；在埃及地饑荒甚大.

41:57各地的人都往埃及去,到約瑟那裡糴糧,因為天下的饑荒甚大.

約瑟洞悉天機,出謀獻策,並且親自督導政策,救埃及人於大飢荒之中,也令埃及成為鄰近地區的救助.因此,若數典不忘恩,約瑟的家族,不應該被描繪為威脅,而應該是感激的對象.不過,對那位不認識約瑟的新王來說,為了他的統治目的,他選擇掩眼不看約瑟對埃及的貢獻這段歷史,否定約瑟對埃及的恩情.

三.  不看曾經對你的承諾

第二個層面的「不看」,是不看埃及曾經對約瑟的家族作為的承諾.因為約瑟成為了埃及的救命恩人,法老曾經對約瑟的家人所出如此的承諾:(創世記)

45:16這風聲傳到法老的宮裡,說:「約瑟的弟兄們來了.」法老和他的臣僕都很喜歡.

45:17法老對約瑟說:「你吩咐你的弟兄們說:『你們要這樣行:把馱子抬在牲口上,起身往迦南地去.

45:18將你們的父親和你們的眷屬都搬到我這裡來,我要把埃及地的美物賜給你們,你們也要吃這地肥美的出產.

45:19現在我吩咐你們要這樣行:從埃及地帶著車輛去,把你們的孩子和妻子,並你們的父親都搬來.

45:20你們眼中不要愛惜你們的家具,因為埃及全地的美物都是你們的.』」

47:5法老對約瑟說:「你父親和你弟兄到你這裡來了,

47:6埃及地都在你面前,只管叫你父親和你弟兄住在國中最好的地；他們可以住在歌珊地.你若知道他們中間有甚麼能人,就派他們看管我的牲畜.」

曾幾何時,埃及王把約瑟的家人視為貴賓,並且承諾歡迎他們搬遷來埃及定居,讓他們安居樂業.這位不認識約瑟的新王,掩面不看埃及對約瑟家族的承諾.這些承諾,在新王接任後變成了「逝去的諾言」.

四.  控制你對我的威脅:經濟政策

既然視你為威脅,接下來的行動當然就是要增強控制,甚至出手打壓,以減低你的威脅.這位不認識約瑟的新王回應「以色列人威脅論」主要有以下兩方面,這兩方面的控制方法似乎在現代社會仍然常見.第一方面和經濟生活有關:

1:11於是埃及人派督工的轄制他們,加重擔苦害他們.他們為法老建造兩座積貨城,就是比東和蘭塞.

1:12只是越發苦害他們,他們越發多起來,越發蔓延；埃及人就因以色列人愁煩.

1:13埃及人嚴嚴地使以色列人做工,

1:14使他們因做苦工覺得命苦；無論是和泥,是做磚,是做田間各樣的工,在一切的工上都嚴嚴地待他們.

統治者要控制人民,其中一種做法是令他們成為純粹的經濟動物.比較文明的社會,統治者鼓勵人民做高級經濟動物,集中所有精力去追求經驗增長,累積更多的財富；而可憐的經濟動物,只能為了糊口而生存.「日日做到無停手」,「手停口停」,人們便不會有空間去理會和追求其他事情了.如果人們問為甚麼日做夜做,一生辛勞?埃及王建議他們埋怨自己命苦.是自己的命生得不好,不是誰人的錯,也不是政策的錯.

不過,有時統治者的如意算盤未必打得嚮.埃及王以為,透過勞役耗盡約瑟後人的精力,他們便不會有閒情逸緻去生兒育女；把以色列人折磨得半死,令他們感到命苦,便不會想到開枝散葉.然而,在這個故事中,埃及王預測對了一半,卻又錯了一半,對的是被他的新政所針對的人的確是筋疲力盡,埋怨命苦；錯的一半是,他們並沒有因此而減少生育,反而更加生養眾多.

五.  控制你對我的威脅:人口政策

經濟手段未能達到目的,埃及王運用第二個層面的策略,就是人口政策.埃及王的人口政策的目的,是要透過此消彼長,去把埃及帝國內的以色列人的比例降底.現代人運用人口政策相對比較文明,但這位古埃及王的人口政策是相當直接及殘酷的.

1:15有希伯來的兩個收生婆,一名施弗拉,一名普阿；埃及王對她們說:

1:16「你們為希伯來婦人收生,看她們臨盆的時候,若是男孩,就把他殺了；若是女孩,就留她存活.」

1:17但是收生婆敬畏上帝,不照埃及王的吩咐行,竟存留男孩的性命.

1:18埃及王召了收生婆來,說:「你們為甚麼做這事,存留男孩的性命呢?」

1:19收生婆對法老說:「因為希伯來婦人與埃及婦人不同；希伯來婦人本是健壯的,收生婆還沒有到,她們已經生產了.」

1:20上帝厚待收生婆.以色列人多起來,極其強盛.

1:21收生婆因為敬畏上帝,上帝便叫她們成立家室.

1:22法老吩咐他的眾民說:「以色列人所生的男孩,你們都要丟在河裡；一切的女孩,你們要存留她的性命.」

殺希伯來人的男嬰,是一項殘酷的暴政.在這暴政之下,場面應該是非常恐怖,呼天搶地的.奇怪的是,聖經作者的敘述,卻聽不到哀哭聲,也看不到恐怖場面.在這暴政之下,被記述的是古代的醫護人員如何勇敢地抗拒毀滅生命的命令,堅守保護生命的神聖使命,並且智慧地與埃及王週旋.專業人員未能有效完成埃及新王的命令,他竟然發動群眾運動,動員埃及人以實際行動去對付以色列人(1:22法老吩咐他的眾民說:「以色列人所生的男孩,你們都要丟在河裡；一切的女孩,你們要存留她的性命.」).如果埃及人的手果真的染滿了以色列男嬰的鮮血,埃及人和以色列人必然立刻變成敵人,兩族人民之間的仇恨必然世代相傳.由此可見,這位新王的新政如何惡毒,何等殘酷!故事沒有跟進到底新王這項殺以色列男嬰的命令在民間有沒有被貫徹執行.怨有頭,債有主,男嬰被殺的哀號沒有在被奴役者的家中出現,在甚麼地方出現呢?在殘酷的統治者自己的家中出現.在甚麼時候被聽到呢?在埃及王幾十年來不斷瘋狂地壓迫之後,在埃及王莫視上帝提出的九個嚴厲警告仍然不悔改之後.

12:29到了半夜,耶和華把埃及地所有的長子,就是從坐寶座的法老,直到被擄囚在監裡之人的長子,以及一切頭生的牲畜,盡都殺了.

12:30法老和一切臣僕,並埃及眾人,夜間都起來了.在埃及有大哀號,無一家不死一個人的.

新王的暴政,把約瑟的後人被推進黑暗時代.在經濟上,他們越來越艱苦而無所收獲；在生活上,他們被監控得越來越嚴厲；在安全上,他們後代的生命受到越來越廣泛的威脅；在社會上,他們被視為叛國賊,變成被仇視的對象.被壓逼者可能仍未能弄清楚為甚麼自己有這樣不幸的遭遇,他們只知道自己被擲入痛苦的深淵中,一片黑暗,無奈又無助.

這位不認識約瑟的新王已經完全控制局勢,在他隻手可以遮天的政治舞台上,已經沒有甚麼對手戲可以看.在這樣的狀況下,作者把鏡頭轉向一個掙扎求存的小家庭,一個大時代的小人物.

六.  大時代的小人物

2:1有一個利未家的人娶了一個利未女子為妻.

2:2那女人懷孕,生一個兒子,見他俊美,就藏了他三個月,

2:3後來不能再藏,就取了一個蒲草箱,抹上石漆和石油,將孩子放在裡頭,把箱子擱在河邊的蘆荻中.

2:4孩子的姊姊遠遠站著,要知道他究竟怎麼樣.

2:5法老的女兒來到河邊洗澡,她的使女們在河邊行走.她看見箱子在蘆荻中,就打發一個婢女拿來.

2:6她打開箱子,看見那孩子.孩子哭了,她就可憐他,說:「這是希伯來人的一個孩子.」

2:7孩子的姊姊對法老的女兒說:「我去在希伯來婦人中叫一個奶媽來,為你奶這孩子,可以不可以?」

2:8法老的女兒說:「可以.」童女就去叫了孩子的母親來.

2:9法老的女兒對她說:「你把這孩子抱去,為我奶他,我必給你工價.」婦人就抱了孩子去奶他.

2:10孩子漸長,婦人把他帶到法老的女兒那裡,就作了她的兒子.她給孩子起名叫摩西,意思說:「因我把他從水裡拉出來.」

今天的舊約經課結束在這裡,我們要述說的故事也暫停在這裡.熟識這個故事的讀者會知道,這位摩西就是帶領以色列人脫離奴隸身分,從壓迫之中得到解放的偉大人物,不過,這已經又是幾十年之後的事情了.在此之前的這幾十年的日子,約瑟的後人是怎樣渡過的呢?讀者無從得知,因為故事接著的下一句,已經是摩西長大成人的時代了.

那些日子,他們看不見上帝親自降臨改變歷史,也看不見上帝差派能人異士施行拯救.然而,他們沒有自暴自棄,在惡劣的環境中盡量發揮他們的智慧努力生存.他們沒有能力逼使埃及王改變,但從這個故事中,讀者可以看見他們以何等的智慧和勇氣,精確的計劃和細心的部署,幫助一個小生命生存下來.或者,在這樣處境之下的其他小家庭,也用盡辦法去保存生命.哪位嬰孩是將來的盼望?故事之外的讀者知道,但故事之內的所有人都不知道.甚麼時候會有改變?故事之外的讀者知道,但故事之內的所有人都不知道.誰能帶來盼望?誰能帶來改變?可能是這位,也可能是那位,因此,努力幫助每一位都能夠生存下去.盼望從哪裡來?改變從哪裡來?可能是從一個小家庭而來,可能是從河邊而來,可能是從王宮而來.無論從哪裡來,無論甚麼時候來,他們沒有放棄,他們仍懷著盼望.

七.  有不認識……的新王起來,治理……

如果把第一句經文變成一項填充題:「有不認識……的新王起來,治理……」,你會填入甚麼呢?你會想到聯想到甚麼呢?想到你的自己的個人際遇?想到教會在內地某些城市的的遭遇,或者在伊拉克北部的遭遇?想到我們這個城市的過去和將來?無論想到甚麼,願這個故事成為我們的鼓勵,啟發我們帶著盼望,以堅忍和智慧繼續前行.



\newpage

\section{我們在東亞處境下怎麼背起我們的十字架}
\label{sec:259}
\textbf{松谷曄介牧師 - 我們在東亞處境下怎麼背起我們的十字架?}
\newline
\newline
\hyperref[sec:258]{< < < PREV SERMON < < <}
~
\hyperref[sec:toc]{[返目錄]}
~
\hyperref[sec:260]{> > > NEXT SERMON > > >}


講題:我們在東亞處境下怎麼背起我們的十字架? 　How are we going to take up our cross in the East Asian context?

講員:松谷曄介牧師

場合:香港中文大學崇基學院禮拜堂

日期:2014年8月31日

首先,我們一起來做一個禱告 .

我們的主,我們感謝您今天帶領我們到這裡.特別地感謝您給我們聽您的話語的時間.您的話語是我心中的歡喜和快樂(耶利米書15:16).求您差遣聖靈進入我們的心裡,讓我們能夠明白和接受您的旨意.奉耶穌基督的名禱告,阿們.

＜何謂“自己的十字架”?＞

我們做基督徒,意味著做基督的門徒而跟從基督.我們都覺得這是理所當然的事.但是,如果我們看今天的經文,十二個門徒當中的領袖彼得,當時他還不太明白“跟從耶穌基督”到底是什麼樣的事情.當耶穌指示門徒們他將來要受難和復活的時候,彼得對他說:“主啊,萬不可如此!這事必不臨到你身上”.耶穌就對彼得說:“撒旦,退我後邊去吧!你是絆我腳的；因為你不體貼神的意思,只體貼人的意思”.他好像按照自己的想法而 誤會主耶穌的話語.

＜十字架-為了“和好”＞

然後,耶穌對門徒們說: “若有人要跟從我,就當舍己,背起他的十字架來跟從我”.我們跟從耶穌的時候,並不是沒做什麼而跟從他的.反而,他吩咐我們“背起自己的十字架”而跟從他.那麼,“背起自己的十字架”到底是什麼意思呢?

“背起十字架”並不是我們也應該被釘在十字架的意思.也不是如字面那樣背起很大的木頭 .有人以為,十字架簡直是苦難的象徵,所以我們常常聽見的用法是“那個人現在面臨著很大的挑戰和艱苦,他背 著十字架”.我覺得這個用法不是完全不對的,但是這是有一點私人化(個人化)的理解.從整本聖經來看 ,“十字架”有更大的意義.我們首先要了解主耶穌的十字架到底是什麼意思 ,然後才能明白“我們背起自己的十字架”是什麼意思.

十字架本來是羅馬帝國時代的一個很殘酷的死刑工具,所以本來有否定的含義 .但是,現在很多教堂把十字架挂在教堂的前面或者屋子上.有一些基督徒帶著有十字架的服飾品.為什麼呢?原因是 ,耶穌基督復活了之後,門徒們知道了,主耶穌的十字架原來是讓人得救的事情.後來,十字架從“否定的含義”變成了“肯定的 含義”.現在,我們基督教把十字架看成上帝賜給我們的“贖罪,拯救,愛,公義,和平,和好 ” 等等的象徵.

尤其是保羅,他在書信裡面,把耶穌基督的十字架看成上帝賜給我們的“和好(reconciliation)”的象徵.羅馬書第5章10節:“因為我們作仇敵的時候,且籍著神兒子的死,得與神和好”.對保羅 來說,罪的根本 問題就是上帝與我們人之間的關係破裂 .所以,我們常說的“贖罪,拯救,得救”意味著上帝與人之間的關係恢復 ,換句話來說,這就是“和好”,或者“和睦”,“和解”.

我們按照這樣對十字架的理解 來一起學習,“背起自己的十字架”是什麼意思.

＜猶太人教會和外邦人教會的和好＞

耶穌的十字架主要是為了恢復 上帝與人之間的關係.其實,他的十字架還有另外的意義,就是成就人與人之間的和好.

以弗所書 第2章13-16節:“13你們從前遠離神的人,如今卻在基督耶穌裡,靠著他的血,已經得親近了.14因他使我們和睦,將兩下合而為一,拆毀了中間隔斷的牆；15而且以自己的身體廢掉冤仇,就是那記在律法上的規條,為要將兩下籍著自己造成一個新人,如此便成就了和睦.16既在十字架上滅了冤仇,便籍這十字架使兩下歸為一體,與神和好了.”

這裡所說的“你們”包括 猶太人和外邦人.在新約聖經裡,外邦人主要指的是 希臘人.初期教會的一個很嚴重的問題是猶太人教會和外邦人教會之間的關係,以及 同一個教會裡的猶太人基督徒和所謂外邦人基督徒之間的關係.猶太人基督徒和外邦人基督徒都接受了耶穌基督,他們每一個人都籍著耶穌基督的十字架而 與上帝和好.雖然他們都是基督徒,猶太人和外邦人之間仍然有 各種各樣的不同點,包括民族 ,語言,文化,生活習慣等等,雙方都難以接納對方.他們中間還有難以克服的“隔斷的牆”.因此,保羅在這裡解釋十字架的時候非常強調 :“因他使我們和睦,將兩下合而為一,拆毀了中間隔斷的牆”.

保羅在不同的書信裡 反複地說:“並不分猶太人,希臘人”,是因為他相信猶太人基督徒和外邦人基督徒都籍著十字架的和好來 可以有親密的關係.但是,很遺憾的是 ,當時實際上還沒有實現那樣的關係.所以保羅一直為了成就猶太人教會和外邦人教會之間的和好而貢獻力量.

＜和好的職分＞

保羅的熱情在哥林多後書裡很明顯地表現出來.哥林多後書第5章18-19節:“18一切都是出于神；他籍著基督使我們與他和好,又將勸人與他和好的職分賜給我們.19這就是神在基督裡,叫世人與自己和好,不將他們的過犯歸到他們身上,並且將這和好的道理托付了我們.”

他在這裡說,上帝賜給我們勸人和上帝和好的職分,而且上帝把這和好的道理(話語)托付給 我們.意思就是,上帝呼召我們每一個人,吩咐和鼓勵我們為了傳揚“和好的福音”而奮鬥!

如果我們可以理解耶穌的十字架是為了成就“上帝與人之間的和好”以及 成就“人與人之間的和好”,那麼,我們現在可以明白,“背起我們的十字架而跟從耶穌基督”是什麼意思.“背起自己的十字架”就意味著我們也為了傳揚“和好的福音”而貢獻力量.

如前所說,耶穌基督的十字架有兩個方面,就是上帝與人之間的和好,以及 人與人之間的和好.因此,傳揚“和好的福音”有兩個方面.第一個是告訴別人耶穌基督的十字架的意義而勸他們與上帝建立和好的關係.這是我們基督徒和教會的一個很重要的大使命.

傳揚“和好的福音”的第二個方面就是為了人與人之間的和好或者教會與教會之間的和好而努力 .為了更深入地了解這一點,接下來我們看看保羅所做的具體的“和好的職分”.看看他怎樣背起他自己的十字架.

＜和好的職分－保羅把外邦人教會的奉獻 送到耶路撒冷＞

保羅寫信給在以弗所和哥林多的外邦人教會的時候,保羅要求他們為了耶路撒冷的猶太人教會奉獻.保羅打算把那些奉獻送到耶路撒冷的猶太人教會.當時,耶路撒冷是民族主義最高漲 的地方.對民族主義者來說,像保羅所說的“不分猶太人,不分外邦人.他們都是神的子女”這種想法是接受不了的事.如果保羅去到耶路撒冷的話,被抓起來或者被殺的可能性是很大 的.所以,他還是要去耶路撒冷,就是一個冒死的決定.

為什麼保羅拼命要把外邦人教會的奉獻送到耶路撒冷的猶太人教會呢?這不僅是為了物質上的支援,而且是為了把雙方的教會結合起來,合而為一.保羅相信,雙方的教會互相沒有關係是不應該的狀態.如果外邦人教會跟耶路撒冷的猶太人教會沒有關係的話,那麼外邦人教會會失去從舊約聖經時代以來的上帝的救贖歷史 ,外邦人教會的上帝跟亞伯拉罕,以撒,雅各的上帝也會沒有關係了.反過來,如果耶路撒冷的猶太人教會跟外邦人教會沒有關係的話,猶太人教會變成一個民族主義的教會,不能維持教會的普世性.如果雙方的教會不能建立主內的和好關係,教會就不是 真正的教會了.因此,為了讓 教會成為真正的教會,保羅冒險將 外邦人教會的奉獻送到耶路撒冷的猶太人教會.

我們應該記著,這並不是保羅一個人做的私人的使命 .他的背後有支持他的不同教會和許多弟兄姊妹的存在.他不是孤單地背起十字架,反而很多人為了建立和好的關係一起背起保羅的十字架.所以可以說,保羅的使命不是私人的使命,反而是教會托付他的使命.

＜與喜樂的人要同樂,與哀哭的人要同哭＞

今天的另外一處經文 是羅馬書第12章9-21節.保羅得知在羅馬教會內部的猶太人基督徒和外邦人 基督徒之間有一些矛盾,他因為盼望著雙方的和好而 寫了這封書信.10節:“愛弟兄,要彼此親熱；恭敬人,要彼此推讓”,16節:“要彼此同心”,保羅在這裡勸他們“彼此相愛,彼此和好”.

“彼此相愛,彼此和好”是什麼樣的狀態呢?15節:“與喜樂的人要同樂,與哀哭的人要同哭”.無論屬于什麼民族,無論有什麼語言,文化,在耶穌基督裡一起分享喜樂,一起承擔負擔,這就是耶穌基督籍著 他的十字架賜給我們的“和好的關係”.

＜我們在東亞 的和好的職分＞

大家會感覺,今天的講道題目“我們在東亞 處境下怎麼背起我們的十字架?”有一點誇大.  當然,自己的信仰生活,自己的教會,本地的教會,都是重要.但是,基督教信仰和基督的教會並不是一個私人的事情,而是一個普世性的事情.所以,我們應該多多關心世界各地的教會和弟兄姊妹而為他們禱告.尤其是日本,香港,中國內地,台灣,韓國,北朝鮮等等,這些所謂東亞 的地域互相是鄰居 ,所以互相關心和 禱告是很重要 的.

今天我想跟大家一起想一想,我們東亞 的教會和教會之間有沒有“和好”的關係?比如說,日本教會與香港教會,有沒有和好?有人會這樣想 :“唉,雙方之間沒有什麼仇恨,怎麼需要和好呢?” 　對,可能沒有仇恨.但是,日本教會和香港教會是不是已經成為像聖經所說的“與喜樂的人要同樂,與哀哭的人要同哭”這樣的關係呢 ?

現在,很多日本人喜歡來香港旅遊,喜歡買東西,吃點心等等.很多香港人也喜歡去日本旅遊,很愛吃日本菜,很愛看日本的卡通和漫畫.但是,很可惜的是,日本基督徒對香港的教會幾乎沒有什麼關心和興趣.香港基督徒對日本教會的理解也是不太多.這樣的狀態 是不是“和好”的關係呢 ?日本教會和內地教會之間的情況也是一樣.

明年是二戰 70周年,日本與中國,韓國之間還有矛盾和衝突,還沒有真正的和好.在這個時代,至少各地的教會應該成為和好的榜樣,但是實際上,還是像保羅時代的猶太人教會和外邦人教會的關係一樣,日本教會和其他東亞 的教會之間還沒有成就和好的關係.

＜我為什麼來了香港＞

我們每一個人都在東亞 的各地做主耶穌的門徒.我們應該為了東亞 各地的教會之間的和好而貢獻力量.大家會問問 :“具體地來說,我們應該怎麼辦?你自己為了這件事情做了什麼 呢?”那麼,我想簡單地分享我在這裡做的研究活動.

我在日本的教會裡牧會了六年 之後,去年秋天來到香港.我現在在崇基學院神學院做訪問研究員.主要的研究題目是二十世紀以來的中國教會歷史.我為什麼研究這個題目呢?是因為,如果日本教會將來能夠跟內地的教會自由地交流的話,我們日本教會需要了解對方的歷史.我一直希望,日本的弟兄姊妹和內地的弟兄姊妹將來有更多的交流,而且有“一起喜樂,一起哀哭”這樣和好的關係.所以我 現在的研究是為了將來的和好關係做預備 .

其實,來了香港之後,上帝還 給我增加另外兩個研究題目.第一,“香港教會的研究”, 所以我慢慢開始讀有關香港教會的資料,探訪不同的本地教會.第二,學習廣東話(粵語).幸好,中文大學裡面有很好的語言學校,我就開始上學,學了大概半年廣東話.有些朋友問我:“你為什麼還想要研究香港教會啊?已經會說普通話,為什麼還需要學習廣東話?”理由是很簡單的.是因為,我想多了解香港教會,想要跟香港的弟兄姊妹一起交流.是因為,我盼望日本教會和香港教會將來 也有密切的關係.是因為,我希望日本的基督徒和香港的基督徒的關係成為“一起喜樂,一起哀哭”這樣的關係.

我來香港留學,一方面是我個人自己決定要來 的,但是另一方面,日本的一批弟兄姊妹也了解和贊成我的異象(vision),而一直為我禱告 .他們也是跟我一起背起為了和好的十字架.

＜我們在東亞處境下怎麼背起我們的十字架?＞

為了建立東亞 教會的和好關係,我們具體地可以做什麼呢?其實,應該做的事和可以做的事非常多!

第一:試試開始學日語.如果大學生有機會去國外留學,去日本吧!然後,重要的是 ,參加本地的日本教會, 跟日本的弟兄姊妹一起崇拜,一起禱告,一起傳福音,一起喜樂,一起哀哭.

第二:有些人會想:“我已經不是可以學外國語的年齡,而且已經工作,不能留學”.但是,如果有機會 去日本旅遊的話,禮拜天也留在日本,然後參加本地的日語崇拜.不懂日語不是重要的問題.雖然語言不同,但是日本教會的弟兄姊妹也念一樣的主禱文,使用一樣的聖經.把自己放在那種環境的時候可以感覺到,我們無論是哪裡人,無論講什麼語言,我們都是在主內合一的,上帝的和好就在那裡.

第三:如果自己不能學習外國語,自己不能去國外,還有可以在這裡做到的事情,就是鼓勵和支持身邊的人學習和去國外.尤其是支持下一代的年輕人也是很重要 的.自己做不到的vision托付他們,鼓勵他們,為著他們禱告,這些都是參與上帝的和好工作.

在東亞 ,基督徒的比例只有韓國高一點,但是整體來說,我們基督徒還是少數派.日本一億兩千萬的人口當中,基督徒(包括天主教徒)的比例是百分之一以下,差不多一百萬.無論是日本,無論是其他的地方,亞州地域的基督教宣教有不同的困難.但是我還相信,如果東亞 的各地教會之間能夠建立和好的關係,我們就能夠把“和好的福音”更好地傳揚給這個世界,可以做美好的見證.

今天,我特別感謝上帝,他給我一個機會在這裡 講道.感謝崇基學院禮拜堂邀請我這一個 從日本過來的牧師.可能,大家都是第一次聽日本牧師的講道 .其實,我是第一次在香港的本地教會裡講道.而且在我的人生當中,這次就是第一次用中文來講道 .跟大家一起崇拜的這一天,我永遠不會忘記.我想把崇基學院禮拜堂的事情告訴日本的弟兄姊妹.日本教會也會漸漸地開始為著香港教會,為著內地教會,為著東亞 教會禱告.

請大家告訴你身邊的朋友和家人:今天一個日本牧師過來講道 了!雖然日本教會的規模很小 ,但是在日本確實也有教會和弟兄姊妹.請大家以後為著日本教會禱告.

耶穌基督在我們當中,在我們面前已經賜給我們很寶貴的“和好”.從今天起,從現在起,我們一起背起為了和好的十字架吧!我們都是跟從基督的門徒,也是一起喜樂和一起哀哭的上帝的家庭!我們一起把“和好的福音”傳揚下去吧!

我們低頭禱告.

親愛的上主,今天我們都聽到了您對我們的呼召.

求您差遣我們去這個世界.

求您為了傳揚您的和好福音而使用我們.

在東亞為了成就在主內的和好,我們願意把我們的時間和能力都獻上給您.

求您祝福我們每一個人,求您賜給我們需要的力量.

求您幫助我們能夠為了東亞教會之間的和好而貢獻力量.

奉主耶穌的名求,阿們.



\newpage

\section{饒恕別人就是釋放一個囚犯}
\label{sec:260}
\textbf{伍渭文牧師 - 饒恕別人就是釋放一個囚犯}
\newline
\newline
\hyperref[sec:259]{< < < PREV SERMON < < <}
~
\hyperref[sec:toc]{[返目錄]}
~
\hyperref[sec:261]{> > > NEXT SERMON > > >}


創世記 50:15-21
\newline
\begin{longtable}{cl}
\hline
\hline
章節 & 經文 (和合本修訂版)\\
\hline
50:15 & \begin{tabularx}{0.7\textwidth}{X} 約瑟的哥哥們見父親死了，就說：「也許約瑟仍然懷恨我們，會照我們從前待他一切的惡，重重報復我們。」 \end{tabularx} \\ \\ \relax
50:16 & \begin{tabularx}{0.7\textwidth}{X} 他們就傳口信給約瑟說：「你父親未死之前曾吩咐說： \end{tabularx} \\ \\ \relax
50:17 & \begin{tabularx}{0.7\textwidth}{X} 『你們要對約瑟這樣說：從前你哥哥們惡待你，你要饒恕他們的過犯和罪惡。』現在求你饒恕你父親的神之僕人們的過犯。」他們對約瑟說了這話，約瑟就哭了。 \end{tabularx} \\ \\ \relax
50:18 & \begin{tabularx}{0.7\textwidth}{X} 他的哥哥們又來俯伏在他面前，說：「看哪，我們是你的奴隸。」 \end{tabularx} \\ \\ \relax
50:19 & \begin{tabularx}{0.7\textwidth}{X} 約瑟對他們說：「不要怕，我豈能代替神呢？ \end{tabularx} \\ \\ \relax
50:20 & \begin{tabularx}{0.7\textwidth}{X} 從前你們的意思是要害我，但神的意思原是好的，要使許多百姓得以存活，成就今日的光景。 \end{tabularx} \\ \\ \relax
50:21 & \begin{tabularx}{0.7\textwidth}{X} 現在你們不要害怕，我必養活你們和你們的孩子。」於是約瑟安慰他們，講了使他們安心的話。 \end{tabularx} \\ \\
[1ex]
\hline
\hline
\end{longtable}
講題:饒恕別人就是釋放一個囚犯-上主的主權與人的責任　To Forgive Is To Set Free A Prisoner

經文:創世記50:15-21;37:18-27

講員:伍渭文牧師

場合:香港中文大學崇基學院禮拜堂

日期:2014年9月14日

今天題目:饒恕別人就是釋放一個囚犯,並不是公權力的特赦.美國政府面對眾多非法移民,在適當的時候也會宣布特赦,把問題浮上面,易於管理；香港成立廉公署時也有特赦,既往不究.特赦的對象,其實是政府；若政府執法得力,就不須特赦.今天題目:饒恕別人就是釋放一個囚犯,仍是一個受傷害的人,饒恕傷害你的人.相信大家看到這題目時,可能想到一句潛台詞:這個囚犯就是你自己.我們一生在過去總會有意或無意傷害了人,也會感到被人所傷害,要是心存怨憤,不肯饒恕傷害我們的人,自己就被怨憤所囚禁,甚至被報復的情緒驅策著我們,找機會報復.

被罪控訴的囚犯

但今天舊約經課讀出的創世記五十章末段的話,囚犯不是被傷害的約瑟(Joseph),乃傷害約瑟的兄弟.約瑟為雅各(Jacob)晚年所出,特別得寵.雅各共有十二個兒子,即後來以色列十二支派,從拉結生的祇有約瑟和便雅憫.雅各為了娶得拉結(Rachel),甘心在舅父拉班(Laban)手下工作七年,竟被騙娶得眼睛沒有神氣的妹妹利亞(Leah),拉班看準雅各會再為漂亮的拉結工作多七年.「以色列原來愛約瑟過於愛他的眾子,因為約瑟是他年老生的；他給約瑟做了一件彩衣.約瑟的哥哥們見父親愛約瑟過於愛他們,就恨約瑟,不與他說和睦的話.」(創三十七3, 4)父母偏心,做成兄弟姊妹之間的嫉妒,是很非常深刻難忘的,因他們自幼一起生活.小孩子是敏銳的觀察者,差勁的解釋者；他們不會明白雅各用了十四年才娶得拉結,她所生的老么當然當是爸爸的掌上明珠.更令哥哥們不快的,是約瑟把他們的惡行報告給父親知道；穿彩衣者,就是打小報告的人.

約瑟自少常見異夢,而且圓夢的能力,與日俱長,後來為法老王圓夢備受重用.他十七歲時發了一個夢,見到自己在田裡所捆的禾稼站在中間,兄長們的禾稼圍著自己的禾稼下拜；另一次在夢中見到太陽,月亮,與十一個星向自己下拜.(創三十七6-11)當約瑟告訴兄長這兩個夢時,他們更加恨約瑟.有一次趁在外邊放牧時,他們想用計謀把約瑟殺掉.但流便(Reuben)不忍流弟弟的血,建議扔他在坑裡自生自滅；猶大(Judah)更仁慈,說服眾人把約瑟賣往埃及作奴隸.約瑟曾求兄長們放過自己,但他們硬了心要除掉約瑟.到約瑟成為埃及宰相,在七個豐年時就做好防饑的儲糧工程,到了全地連續七年饑荒,在迦南地的雅各叫眾兒子南下到埃及買糧食.

約瑟想起以前作的夢:十一個星向自己下拜,就說:你們說家中有十二人,為何祇十個來?你們是否是奸細?立即回去帶來兩位幼弟,證明你們是誠實人.兄長們知道家中祇有最幼的便雅憫留在家中,約瑟已經被賣掉了,那裡可能找回他呢?此時他們想起賣弟弟的惡行.「他們彼此說:我們在兄弟身上實在有罪,他哀求我們的時候,我們見他們心中的愁苦,卻不肯聽,所以這場苦難臨到我們身上.」(創四十二21)

那麼多年來,那位穿彩衣弟弟哀求的聲音揮之不去,謀害約瑟的兄弟,多年來被罪債追討；就算受害者約瑟不追討,上帝也會追討,因為他們流了無辜者的血.

創世記:墮落開始,恩典結束

然而,今天舊約經課的重點,不是指出我們犯了罪,傷害了人,被罪追討；仍告訴我們,上主有赦罪之恩.創世記開始時,描述罪帶來墮落,失去樂園,兄弟仇殺:該隱殺亞伯,雅各以紅豆湯迫使以掃交出長子名份,並假扮以掃騙取父親的祝福；兄弟合謀賣掉約瑟.

但創世記結束時,看到上主對罪人的憐憫.約瑟對兄長們說:「從前你們的意思是要害我,但神的意思是好的,要保全許多人的性命,成就今日的光景.不要害怕,我豈能代替神呢?……現在你們不要害怕,我必養活你們和你們的婦人孩子.」(創五十20, 21)約瑟自己看所受的苦難是上主的拯救計劃,不少解經家都認為約瑟就是基督的預表(type):為父所愛,但被兄弟們出賣,經歷種種苦難,卻帶來拯救和生命；一如耶穌被出賣,掛在十字架,拯救世人.

我豈可代替上主呢?

「我豈能代替神呢?」這句話指出,傷害別人,不單是受害人和施害人之間的私人恩怨,犯罪的人要面對上主的審判.我們都知道,刑事案不是原告和被告的私人恩怨.觸犯刑法,主控官就代表司法追究罪犯,原告祇是法庭上的證人,罪犯要面對法律制裁.我們得罪人時也得罪上主,上主會憑公義按各人所行的報應各人,約瑟把申冤的事,交給公義的上主；沒有人可以逃罪,但上主有赦免的恩典.赦免是上主的本性,他以憐憫罪人為他的榮耀.讓我們看看今天的詩篇經課.

詩篇經課一0三8-13

這詩是感恩詩,感謝上主赦免過犯,饒恕我們的罪過.這段詩篇經課重覆又重覆三同義詞:憐憫,恩典,慈愛.慈愛hesed,指立約的愛,上主和選民立的約不是商業協議,利益交換,乃是宗主國對其蕃屬,或貴族對所擁有的百姓,宣布其主權:你們是屬於我的,我會愛你們,保護你們.

背景是出埃及記三十四6 頒下十誡時立約的宣告:「耶和華,耶和華,是有憐憫有恩典的神,不輕易發怒,並有豐盛的慈愛和誠實.」從上主本身性情看,是慈愛；從罪人本身的需要看,是憐憫和恩典；誠實faithfulness,指這慈愛恆久不變.我們要注意,立約頒下的是第二塊法版,第一塊因為拜金牛犢摔碎了.摩西上西乃山,久久不下來,他們按捺不住,造金牛犢代替上主敬拜.摩西下山聽到百姓圍繞金牛犢跳舞唱歌,忘形放肆,「便發烈怒,把兩塊版扔在山下摔碎了」.(出三十二19)

人不能憑自己遵守律法,摩西求上主赦免百姓的罪:「不然,求你從你所寫的冊上塗抹我的名.」(v. 32)摩西知道百姓不容易帶領,求上主同行,並顯出他的榮耀作為保證:「求你顯出你的榮耀給我看.」(出三十三18)「耶和華說:『我要顯出我一切的恩慈,在你面前經過,宣告我的名.我要恩待誰就恩待誰；要憐憫誰就憐憫誰』；又說:『你們不能看見我的面,因為人見我面不能存活.』耶和華說:『看哪,在我這裡有地方,你要站在磐石上.我榮耀經過的時候,我必將你放在磐石穴中,用我的手遮掩你,等我過去,然後我要將我的手收回,你就得見我的背,卻不得見我的面.』」(出三十三21-23)

你就得見我的背

摩西求見上主的榮耀,上主顯示他的憐憫,好像牛頭不搭馬嘴,其實憐憫就是上主的榮耀.上主叫我們看看基督在十字架的死,就是上主顯出他榮耀的時候.當上主的榮耀經過時「你就得見我的背」是甚麼意思呢?為何見到上主的背,就是見到他的榮耀?因為沒有人能見到上主的本體,我們一進入上主赤裸裸的聖潔,就消滅了；就如細菌見到激光極大的熱能,紛飛煙滅.上主的背就是耶穌,從來沒有人見過上主,祇有在父懷裡的獨生子將他表明出來.基督就是上主的憐憫,他的饒恕是永遠的.

永遠的赦免

「耶和華有憐憫,有恩典,他不長久責備,也不永遠懷怒.」(詩103:8)「他赦免你的一切罪孽,醫治你的一切疾病.」(v.3) 一切,表示完全；因為完全赦免,所以忘記了,不會再題起,所以他不長久責備,也不永遠懷怒.上主不同人,就算彼此相愛的配偶,我們激氣時,都會題起對方過去的不良記錄,很難洗底.一切,也包括了將來,是無條件的接納.反而常常題起的,是他信實不改變的慈愛.

聖餐中叫我們常常想起基督犧牲的大愛,不是我們的過犯.聖餐紀念anamnesis一字,有拿起來的意思；聖餐的紀念不是回到過去,哀慟基督的死；紀念是把過去提到現在,使現在產生意義.像結了婚的人,提起過去簽了婚書的事實,作為現在的憑據,將來的盼望.單單回顧婚禮,看看錄影有甚麼意義呢?不是愐懷過去,乃彼此題醒,我們已經結了婚,現在還有了兒女,就好好相愛,一起面對明天,我們己活在彼此的盟誓和承諾中.聖餐的紀念,題醒我們活在上主的慈愛中和憐憫中,不被控訴和定罪了,所以「他赦免你的一切罪孽,醫治你的一切疾病.」

憐憫按需要

詩篇一0三11-14跟著解釋甚麼是憐憫.公義按我們應得的,憐憫按我們需要的.很多法庭前面都有蒙眼代表公義的塑像,手裡拿著天秤,不看人面,祇看客觀的斤兩.憐憫按我們的需要:「他知道我們的本體,思念我們不過是塵土.」(v. 14) 所以「他沒有按我們的罪過待我們,也沒有照我們的罪孽報應我們.」(v. 10)

饒恕使我們敬畏

「天離地何等的高,他的慈愛向敬畏他的人也是何等的大!東離西有多遠,他叫我們的罪離我們也有多遠.」(vv. 11, 12)敬畏上主,覺得虧欠上主,認罪悔改的人,才到上主的赦免.同時心被恩感的人也會敬畏上主,愛上主愛人.怎能知道我們蒙饒恕呢?就是我們饒恕人.今天福音書經課就提到被主人饒恕欠一千萬銀子的人,不肯饒恕欠他十兩銀子,主人「便揪著他,掐住他的喉嚨……把他下在監裡.」(太十八28-30)主人就大怒,把他交給掌刑的,等他還清了所欠的債.

這人不肯饒恕欠自己十兩銀子,揭示了一個事實-他沒有真正悔改,尋求上主饒恕；他的生命並沒有改變,他沒有相應的行為顯出他的信心,因為信心沒有行為是死的.正如雅各書所說:「因為那不憐憫人的,也要受無憐憫的審判；憐憫原是向審判誇勝.(雅二13)

饒恕的信息對今天有何意義呢?

饒恕是理解

作為蒙上主憐憫的教會,明白憐憫原是向審判誇勝,憐憫優勝過審判.公義,審判按我們應得的；憐憫,饒恕按我們的需要.政治權力,我們強調應得,人大常委的框架是中央應有的權力,不能更改,但若說政治智慧,就要考慮香港人的需要；易地而處,我們也要理解掌有權力的人的需要,不能一味說這是我應得的,政治智慧就是想到對手的需要.

饒恕是理解,按對方需要理解對方.今天書信經課提及「信心軟弱的,你們要接納,但不要辨論所疑惑的事.」(羅十四1)羅馬書不錯指的是教會信徒之間如何接納,但這教訓對人性有洞察力.解決問題,除了辨論說理之外,人與人之間的信心和信任都很重要.香港和中央之間信心和信任不夠,單憑辯論是不行的.我們蒙上主饒恕的人,明白到憐憫是按需要,要推動彼此理解,追求和睦,建立信任.

歷史的上主

今天舊約經課和詩篇經課都談及饒恕,創世記約瑟饒恕兄長,詩篇一0三8-13西乃山立約上主饒恕百姓；兩棕事件遙遙相應.出埃及目的是來到西乃山與上主立約作他的子民:「你和以色列的長老要去見埃及王,對他們說:耶和華-希伯來人的神遇見了我們,現在求你容我們往曠野去,走三天的路程,為要祭祀耶和華-我們的神.」(出三18)

上主運籌帷幄,出埃及事件的出現,因為七年大饑荒,雅各及其家族逃避饑荒來到埃及,因為約瑟貴為宰相,法老王撥出尼羅河三角洲東部的肥沃的歌珊地(Goshen) 生聚,到摩西出生時已住了350年,摩西八十歲時帶領以色列人出埃及.

約瑟被賣到埃及,在護衛長波提乏家中被主母引誘不逐,遭誣告坐獄,但因為表現良好行事有智慧竟然被司獄賞識,把所有囚犯交給約瑟管理.約瑟在獄中為下囚的法老的膳長和酒政解夢,之後酒政出獄復職,法老王發夢需要智慧的人解夢時,膳長推薦約瑟,約瑟為法老解夢得寵,貴為宰相,在七個豐年期間進行儲糧工程,預備七個荒年來臨.

上主使用埃及帝國,儲糧拯救人民,我們不要妖魔化任何政權,也許是異教或無神論的政權,祇要這政權行公義,憐愛人民,也是上主手中的器皿,因為上主是歷史的主.

毋違天上來的異像

約瑟所處的世界,異教埃及帝國當權,與父親所教導的一神信仰迥異；曾祖父亞伯拉罕離開吾耳蒙應許的國度,了無蹤影；因為全地七年大饑荒,上帝的選民要向異教帝國乞求糧食.當他遇見賣自己的兄長們,不是想起報復,乃想起上主十七歲給他的異夢,為何祇得十個人?應該十一個星下拜才對.在人看來他事業非常成功,但約瑟看重的是上主在他生命中的計劃.

當然他看不到四百多年後要發生的出埃及史詩,就從被兄長敵視謀害,被賣到遠方就開始譜寫了.他也看不到傲視一切的政權的強大軍隊,一瞬間就被紅海沒頂,所有埃及長子一夜之間喪命.約瑟祇求做好當前的角色,不忘記少年時候上主給他的異夢,不斷思考如何落實領受的托付,毋違從上而來的遠象.他不斷操練圓夢的恩賜,恩賜需要不斷操練,如火挑旺.十七歲時他圓夢時沒有智慧,對兄長們的感受不夠敏銳.但經過磨練,在獄中沒有停止運用圓夢的恩賜,最後在朝廷被法老王所重用.恩賜是上主所賜的,異像或遠象也是上主所賜的.約瑟沒有因為世上事業的繁重,忘記上主在年幼時給自己的感動.

人當盡的責任

約瑟被賣到埃及,貴為宰相,儲糧拯救生命,除看到上主冥冥中的主權,也看到人當盡的責任.正當哥哥們硬了心要殺約瑟時,流便(呂便)力排眾議,不要流弟弟的血,扔在坑裡讓約瑟自生自滅就算了.為何不提議釋放約瑟呢?流便這提議是在最少的空間,作出最寬大的可能了；釋放約瑟是不可能的,這是哥哥們定下的框架,但可以不流約瑟的血,讓他在坑中自生自滅.

當流便爭取到讓約瑟活下來的時候,猶大就有機會提出更進取的意見,見到往埃及的以實馬利人經過,把握機會建議賣約瑟到圾及為奴隸,讓約瑟生命得以存活.不然,在坑裡缺水缺糧,面對遇野獸,也會兇多吉少.流便敢言,鼓勵猶大跟進,猶大最初沒有附和流便的.我們可能不是約瑟的材料,約瑟一表人材,相貌俊美,能幹機智,辦事能力極強,是帝國宰相人材,但可以作流便和猶大,祇要我們把握機會,對不義的事情敢言敢有作為.

結論:饒恕原是向審判誇勝.在今天我們需要學習饒恕的理解,讓上主給我們政治智慧.

饒恕原是向審判誇勝.憐憫最終勝出,不是劍拔弩張的對立,敵視,暴力,極權勝出.面對此刻天雨欲來風滿城,黑雲壓城城欲摧的香港,我靈鎮靜.

最後,當人盡上責任,敢於指出不義,終會帶來改變,逆轉死亡為生命.



\newpage

\section{厚此薄彼豈有此理}
\label{sec:261}
\textbf{許開明牧師 - 厚此薄彼,豈有此理!}
\newline
\newline
\hyperref[sec:260]{< < < PREV SERMON < < <}
~
\hyperref[sec:toc]{[返目錄]}
~
\hyperref[sec:262]{> > > NEXT SERMON > > >}


馬太福音 20:1-16
\newline
\begin{longtable}{cl}
\hline
\hline
章節 & 經文 (和合本修訂版)\\
\hline
20:1 & \begin{tabularx}{0.7\textwidth}{X} 「因為天國好比一家的主人清早去雇人進他的葡萄園做工。 \end{tabularx} \\ \\ \relax
20:2 & \begin{tabularx}{0.7\textwidth}{X} 他和工人講定一天一個銀幣，就打發他們進葡萄園去。 \end{tabularx} \\ \\ \relax
20:3 & \begin{tabularx}{0.7\textwidth}{X} 約在上午九點鐘出去，看見市場上還有閒站的人， \end{tabularx} \\ \\ \relax
20:4 & \begin{tabularx}{0.7\textwidth}{X} 就對那些人說：『你們也進葡萄園去，我會給你們合理的工錢。』 \end{tabularx} \\ \\ \relax
20:5 & \begin{tabularx}{0.7\textwidth}{X} 他們也進去了。約在正午和下午三點鐘又出去，他也是這麼做。 \end{tabularx} \\ \\ \relax
20:6 & \begin{tabularx}{0.7\textwidth}{X} 約在下午五點鐘出去，他看見還有人站在那裡，就問他們：『你們為甚麼整天在這裡閒站呢？』 \end{tabularx} \\ \\ \relax
20:7 & \begin{tabularx}{0.7\textwidth}{X} 他們說：『因為沒有人雇我們。』他說：『你們也進葡萄園去。』 \end{tabularx} \\ \\ \relax
20:8 & \begin{tabularx}{0.7\textwidth}{X} 到了晚上，園主對工頭說：『叫工人都來，給他們工錢，從後來的起，到先來的為止。』 \end{tabularx} \\ \\ \relax
20:9 & \begin{tabularx}{0.7\textwidth}{X} 約在下午五點鐘雇的人來了，各人領了一個銀幣。 \end{tabularx} \\ \\ \relax
20:10 & \begin{tabularx}{0.7\textwidth}{X} 那些最先雇的來了，以為可以多領，誰知也是各領一個銀幣。 \end{tabularx} \\ \\ \relax
20:11 & \begin{tabularx}{0.7\textwidth}{X} 他們領了工錢，就埋怨那家的主人說： \end{tabularx} \\ \\ \relax
20:12 & \begin{tabularx}{0.7\textwidth}{X} 『我們整天勞苦受熱，那些後來的只做了一小時，你竟待他們和我們一樣嗎？』 \end{tabularx} \\ \\ \relax
20:13 & \begin{tabularx}{0.7\textwidth}{X} 主人回答其中的一人說：『朋友，我沒虧待你，你與我講定的不是一個銀幣嗎？ \end{tabularx} \\ \\ \relax
20:14 & \begin{tabularx}{0.7\textwidth}{X} 拿你的錢走吧！我樂意給那後來的和給你的一樣， \end{tabularx} \\ \\ \relax
20:15 & \begin{tabularx}{0.7\textwidth}{X} 難道我的東西不可隨我的意思用嗎？因為我作好人，你就眼紅了嗎？』 \end{tabularx} \\ \\ \relax
20:16 & \begin{tabularx}{0.7\textwidth}{X} 這樣，那在後的，將要在前；在前的，將要在後了。」 \end{tabularx} \\ \\
[1ex]
\hline
\hline
\end{longtable}
講題:厚此薄彼?豈有此理!　What an Outrageous Manner of Lord Showing Favoritism!

經文:馬太福音20章1-16節

講員:許開明牧師

場合:香港中文大學崇基學院禮拜堂

日期:2014年9月21日

引言

各位,在你人生的經歷中,有否遇上厚此薄彼的家人,老師,老板?

倘若遇上厚此薄彼的對待,你的感受又如何呢?

信息

從現代人角度看,這個聖經中的家主,真是厚此薄彼!

1)先僱的工人,整天(即12小時)勞碌受熱,每人只得一錢.那只做一小時的,也得一錢.多勞不多得!不公平!

2)那遲開工的只做了1小時,工作12時便終止.遲來工作的,卻先得到工錢；先被僱的,遲遲才得工錢.在前的反在後!在後的反在前,不合理!

3)家主跟先僱的工人,講定一天一錢銀子；跟後面的其他的僱工說:「所當給的,我必給你們!」這是雙重的標準制度或多重的標準制度,大細超,並不一致.

須處理三個問題:

1)那麼,多勞多得是錯的嗎?

2)做少給多,做多給少,公平麼?

3)人受到不公平的對待,同一公司對員工行兩種制度,厚此薄彼應該嗎?第一批工人不平則鳴難道不該麼?

今天所講的「葡萄園比喻」,是耶穌基督所講的連串天國比喻中的第九個；也是馬太福音書所獨有的資料.

一位傳道人在教會查看這段經文時,向一些大學生發問:「如果是你聽完這比喻, 你將如何把這個故事接續寫下去?」學生們很有想像力,他們有的說要這樣描述下去:「這些感覺受到不公平待遇的人,他們要去拉白布條抗議老闆不公平!」有的人說:「要描寫這些工人準備向勞工處申訴!」；有的則說；「要召開記者會說明老闆沒有公平對待所有勞工.」有一個同學的答案最使人印象深刻,他說:「翌日,這個老闆繼續去街上找工人,結果大家都在閃躲著老闆,他們希望在太陽快下山的時候,才出現給老闆看到,然後被邀請進入葡萄園工作,這樣就可以像前一天最後進入葡萄園的工作者一樣,工作最短時間,卻能領到和最早進入葡萄園工作者一樣的工錢.」

各位,你也會覺得這家主厚此薄彼嗎?

聖經學者蒙特佛雷(C.G.Montefiore)把這個比喻稱作:『所有的比喻中,最偉大,最榮耀的一個比喻.』初次聽起來,你未必同意,但細心看下去,卻發現了這比喻另有故事,另有意義!厚此薄彼是真的嗎?

從家主的表現,好像是真的!

葡萄園收成之時,正是炎熱如火的季節.炎熱過後,有傾盤大雨和狂風,若不盡早採摘,葡萄會被狂風打壞及被暴雨吹走!因此家主急於找工人!

工人期待被徵用,得到工酬,以養活家人.他們清早到市場,期望得到工作.還有,按理愈早找到工作,愈得到多些工資.所以家主需要工人,工人也需要家主,大家都需要對方,在今日的社會或上帝國裡,我們都需要對方.

聘請五批的工人:

比喻雖然以平常生活為素材,它的記載卻非常特別,特別之處就是僱用工人的次序與發工錢的次序是顛倒的.比喻中園主僱用了五批工人:第一批工人,看1節,是清早時候僱用的,也就是早上六點,全日做12小時；第二批人,看3節,是從已初的時候請的,就是早上九點,做9小時；第三批工人,看5節,是從午正的時候僱用的,就是中午十二點,做6小時；第四批工人,看5節,是從申初的時候僱用的,也就是下午三點,做3小時；第五批工人,看6節,是從酉初的時候僱用的,就是傍晚五點,做1小時.

發五批工人的工資:

當發工錢的時候,看8節「到了晚上(就是工作結束,六點鐘),園主對管事的說:叫工人都來,給他們工錢,從後來的起,到先來的為止.」第一批領工錢的人是下午5點鐘進葡萄園工作的工人,他們只工作了一個小時；然後是第四批工人,然後是第三批工人領工錢,然後是第二批,最後領工錢的是第一批進去工作的工人.假如你是第一批進去的工人,也就是從早上六點做到晚上六點的工人,你做了十二個小時.你會怎樣想?

當你看到最後一批進去的工人,進去領工資,然後笑容滿面的出來,因為他們雖然只工作了一個小時,卻竟然領到一錢工資的時候,你會怎樣想?耶穌竟然說:「這葡萄園就是預表上帝國!」你覺得怎樣?為何會這樣呢?

上帝國與地上國的比較

耶穌在這比喻中,顯示上帝國與地上國的差異如下:

1)兩種不同的待遇

這個比喻中的工人分作兩種待遇.

v2:「講定一天一銀錢」.第一種是與主人談妥價錢訂立契約的第一批工人,他們說:『如果你給我們多少工資,我們就來工作.』他們是為工資而作,他們的行為,表示他們所關心和看重的是工作所得酬報.

第二種是那些後來被雇用的人,雙方並沒有契約.

v.3 「--所當得的,我必給你.」他們所要的祇是工作的機會,他們不曾提到工資或契約.他們所要的祇是工作,他們願意讓主人來決定酬報的數額,因為他們看重所事奉的是主,並相信家主不會虧待他們.

第一種是計較酬報的人,也是今日許多人的寫照:在這世界上許多得到酬報的人,在天國裡祇有很低的地位,因為他的思想是集中在他的酬報上面.第二種是不計酬報的人,這是天國裡的人.在這世界上,有許多在人看來是酬報很低的窮人,他們在天國裡卻要為大,因為他不想到酬報多少,祇是為了工作的事奉.這種人以事奉主為重.不計較酬報,得到酬報感恩更多.如果你是老板,你欣賞那一類的工人?

2)律法與恩典

上帝既講律法規則；又講恩典!在天國更強調恩典.

律法:重視個人利益,講定工價多少,就給多少, 看重善行功德.

恩典:重視憐憫與慷慨的給予, 給與有需要,貧苦的人,無能力去守律法的人.父母教導子女或者上司面對下屬:若用律法原則,便講規矩原則；若用恩典原則,便是接受一些人有自已的困難, 給予愛心體諒；未能完全按照原則或規矩處理.上帝的國度顯然是著重恩典的原則.

3)在前的在後,在後的在前

一般人看法是先到先得的原則,家主在上帝國卻不依從這原則.天國強調恩典的原則,超越這個觀點.守法,是為好的管治,求公平,是群體的願望.恩典,是為「需要」,不可忽視「恩典」,顧及一些無能力守法的人.正如耶穌的時代,法利賽人己定了許多的律法,耶穌釋放律法,給予他們恩典. 今日已有許多規條. 法制,原則捆綁著個人或群體,我們更需要恩典.

家主,憐憫人的主,明白最後一批的苦況:1)他們由朝等到晚,等了很久都無人徵用他們,等候,焦急,是非常辛苦的事.2)他們要養妻活兒,一家人等著他賺錢開飯.3)這批工人只有一小時,若仍然無人聘任,就糟了,他們極希望還有一線希望.

家主知道最早受僱的第一批人很幸福:1)他們是第一批被選中的工人,肯定有工開.不用徬徨終日,早已知道將獲發工資.2)家主早已應允給予一銀錢銀子,工資與羅馬兵的工資相等,相當不錯,預先講定的工資必定到手,不用擔心,毋須憂慮.所以,家主無論是對第五批,四,三,二批的人都有恩典.甚至是第一批,也是早早的給予他們恩典,早已約定一錢銀.所以,第一批的計較,埋怨是有問題的.

回應

第一批工人的問題:埋怨

工人埋怨的產生:

1)比較:工作12小時,得一錢銀子；工作1小時也得一錢銀子.

2)期望:以為必要多得!

3)不能接受:不公平感覺,不滿足感覺!

4)埋怨發怒:自己不開心,引起主人及周圍人不開心.

主人回應:

1)早已有約定:一天一錢銀子.

2)我給你的和給那後來的是我的錢,怎樣運用是我的權利!

3)我早於19:30節說明天國的原則:「許多在前要在後,在後的要在前.」

4)第一批工人的埋怨是「不滿足」,「貪心」,「妒忌」(紅了眼!)我們也要小心,不要跌入這埋怨的陷阱中!

不滿與埋怨

第一批工人的不滿,是逐漸加深的!他們先看到不公平的家主:

1)最遲來做工的,卻先得工錢?最先的,卻最後得工錢?

2)工錢待遇非常不合理:只做一小時的,得一銀錢?做三小時的,得一銀錢?做六小時的,得一銀錢?做九小時的,得一銀錢?做十二小時的,得一銀錢?

他們是愈看愈不忿的,愈火滾的!

3)他們勞碌整天,且是火熱的時份,那麼辛苦,當時地上無火,但心中有火!

計較→埋怨,這是不好嗎?合理的,不一定「不好」!

人會計較,懂得數算,是聰明人.人懂得分辨公平與不公平,是有品的人.人埋怨,抱不平,不平則鳴,是勇敢的人.計較→埋怨,似是負面的情緒,卻有一定的作用與意義.

他們的計較與埋怨有什麼問題?

這批工人的問題是只從自已角度出發:

1)自已最早返工,最遲放工!別人則工時不足!

2)自已最苦,長期在炎熱天氣下.別人下午三時才來未嘗過炎熱.

3)自己做得最多,最長,最辛苦,但得到的工資最少,是五批工人中最慘,最不公道!今日的公司或機構裡,也常見到有這種的工人.只見自已,不見他人.只見自己的「無」,不見自已的「有」

憤怒的工人

其中一位工人,一定是惱怒得把工錢擲在地下,所以家主說:「拿起你的工錢,走罷!」(14)(Take your Pay and Go!),這是一種命令的語句.其他的工人卻是取了工錢,有錢落袋,才埋怨,才跟家主講數.講成功可以再袋多些,講不成功也有一日的工資在袋了!真是智慧之子!

這位憤怒的工人,不肯「袋住先!」寧願不要,也 「不接受不公平的對待」.他不肯「袋住先」,他敢頂撞老板,自已把錢擲在地上.好可能會被老板取回,因為是這位兄弟自已不要的!不過仁慈的老板,沒有這樣做!今日也有不少不肯「袋住先」的香港人!不過今日不肯「袋住先!」的不是葡萄園的工人,而是在一國兩制下的香港人,各有不同.

第一批工人抒發不滿

埋怨的產生往往是:「期望與事實有出入」期望愈大→ 失望愈大→埋怨愈多.有些人對政改仍有期望,不肯「袋住先」,有些人認為「袋住先」,慢慢講數,都可能有期望.然而這些期望,合理嗎?對第一批的工人,有他們的理.對家主來說,也有他的理.各有自己的道理.然而,我有理,也要明白對方的理!蠻不講理的人,是只講自己的理,不理別人的理,這是非常危險的,會造成衝突!聽過家主的解釋後,我相信那憤怒的工人,會拿起他的錢,離開! 在爭拗中,讓我們自己講理,也謙卑聽聽別人的理!這樣埋怨會減少,甚至消除!

第二批至最後一批的問題:閒站的問題

閒站,Idleness,是指失業,無所事事.第一批工人看到自已的需要,卻看不到其他人的需要:主人卻能看見那些閒站人,無事做的人的需要:1)閒站,沒有工作,沒有收入,不能養活家人.2)閒站,使他們便會感到痛苦,需要工作,卻沒有工作,看著一個跟一個的工人都被召入葡萄園,自己卻一小時又一小時的等,多麼焦急,多麼的痛苦呀!3)長期的閒站,令自已的自尊受損,覺得自己一無是處,痛苦至極.4)閒站的人,容易受試探跌倒!無所事事,容易生事,甚至出事!貪窮的人閒站,易犯罪,為了自己或家人而鋌而走險!神學家加爾文(John Calvin)盡量令自已不會「閒站」,他說:「閒站!我連一分鐘的都不想!」當他因病臥床之時,朋友叫他放下手上的工作,他卻回應說:「莫非你想主回來時,看見我正閒站著,無所事事麼?」」座上眾退休人士,小心閒站!

各位,家主同情「閒站人的苦況」,因此特別關顧他們!給他們工作,解他們的苦惱!給他們一錢銀,使他們可養妻活兒!

什麼是公平?

厚此薄彼是不公平麼?

到天黑的時候,主人分工錢,由晚來的人先分,下午五點來的人得到一個他連得,輪到那些早上就進來的工人,他們心想一定會得到更多,不過也是每人一個他連得.

領了之後,就埋怨起家主來了:「我們整天勞苦受熱,那後來的只做了一小時,你竟叫他們和我們一樣嗎?」

做十二小時和做一個小時領一樣的錢,這樣公平嗎?這是一個倫理的問題.

「公平」是一個全人類社會普遍都有的道德觀念,但是要定義什麼是公平卻不是一件簡單的事.

簡單的就是「對待所有同樣的情形,人,或事情,要用同樣的方法和態度.

對待不同的情況,人或事情,則按其不同,而有所不同.」

這個原則聽起來簡單,但是問題很多.「什麼是同樣的情形?什麼是同樣的人和事情?」

在後現代,我們常聽到有人主張,要尊重每個人的個別差異,因為每個人生來就不同.〔有性別差別的,體能差別的,有智商差別,情緒差別等〕.

但是,在差異當中,一群人要生活在一起,我們就要問有什麼是大家共同具有的共通性,我們可以照著這個共通性制定一個大家都要遵守的標準呢?

若要尊重每個人的不同,又如何按其不同情況,而有所不同呢?由此可見,「公平」不是一個簡單的概念.許多時候,我們聽見有人喊:「我們要尊重『他者』或弱勢者」,也就是那些跟我們不一樣的人.但是如果更仔細聽,可能會發現,所謂「尊重他者」的倫理堅持,事實上,是要求別人「尊重他」,因為他和別人不同,所以他也不必遵守社會所共同制定的規則,他就是他口中的「他者」.「我們整天勞苦受熱,那後來的只做了一小時,你竟叫他們和我們一樣嗎?」

這些清早就入園工作的工人,向葡萄園主人抗議,聽起來很有道理.但是主人對抗議工人的回答有兩個很清楚的理由.第一,主人和工人本來就約定了工資是一個他連得,主人有照他的約定；第二,主人給那些晚來的人也是一個他連得,這錢是由主人的口袋中拿出來的,並沒有讓早來的工人吃虧.所以主人講:「因為我做好人,你就紅了眼嗎?」這句話也可以譯成,「因為我慷慨,你就嫉妒嗎?」(NRSV,“Or are you envious because I am generous?”)

公平, 我有講到一個定義就是,「對待所有同樣的情形,人,或事情,要用同樣的方法和態度.對待不同情況的,就有所不同.」但是公平還有另外一個意思,就是照原來的約定,主人和工人講定一個他連得,主人就付一個他連得,主人有照這個約定,這就是公平.

清早入園的工人會感覺不公平,是因為看見晚來的工人得到一個他連得的工資,所以暗中在想,自己可能會領更多,但是後來的結果並不是照他的期待來發生,令他們大失所望.所以並不是主人對早來的人失約,不公平,而是早來的人有錯誤的期待,在失望時,就生「遭受不公平對待」的感覺.

試想老板對你特別好,賞識你,你很開心!這位老板對你的同事特別好,你就問:「為何要對他那麼好呢?」同樣,你對老板好,這是「忠心,殷勤」表現.人家對老板好,是「擦鞋」,「拍馬屁」.

上帝給你許多恩典, 你會不停感恩,欣然接受!上帝給你的不喜歡人(仇人)許多恩典,你也可能心裡問:「為何這樣的人,上帝也給他那麼多恩典?」

朋友,稱呼?

「朋友啊,朋友!」(etaire)葡萄園的主人稱呼這些早進來,正在埋怨的工人.在新約聖經「朋友」這個詞只有出現三次,三次都是在馬太(跟約翰福音的「朋友」不同),一次是太22:12,國王稱那些沒有穿禮服來赴宴席的人為「朋友」.另外一次是太26:50,猶大帶祭司長來捉耶穌時,耶穌稱猶大衛「朋友」.朋友,在當時的話,是向陌生人用尊敬的口氣,來稱呼不認識的人.好像廣東人稱呼人「老友」,不是「老朋友」,一點也不「老」.各位,你願意因為計較報償,卻被耶穌稱你作「朋友」嗎?

在馬太福音的教會,可能有猶太血統的基督徒,後來發現,沒受割禮的外邦人,不但成為基督徒,甚至還漸漸成為教會的領袖.猶太基督徒心中就產生疑問:「主啊!這樣公平嗎?」主人慷慨,難道不是一件好事嗎?你有感受到上帝的恩典,上帝對你的愛沒有?如果有,就不會再嫉妒了.

我們的公平觀念正確嗎?

從耶穌的這個比喻,我們還可以再反省另外一件事.由早上就進入葡萄園做工的人來看,可知道我們對「公平」的感覺,很多時候可能會出錯,所以我們需要經過再三反省.許多時候,我們以為自己的「公平公義」才是對的,殊不知其實裡面有我們自己看不見的盲點 {好像第一批的工人}.

我們唯有在上帝的話語面前,謙卑自己,不斷反省,才可能避免錯誤的「正義感」.葡萄園工人的比喻,讓我們認識到上帝是一位慷慨的主人,每一個在祂園中工作的,都受祂的恩典.我們若清楚祂對我們的恩典,在上帝面前就沒有什麼好埋怨的了.當心有不平時,要仔細在上帝的面前省察自己,就會知道是自己的自私,驕傲；或是真正有不公義的事.願上帝賞賜我們每一個人有感恩和智慧的心.

結語

家主厚待後來的工人,好像薄待第一批的工人,厚此薄彼,豈有此理?真有此理!上帝國的理,與一般人的理有不同!如果我們明白家主的厚此薄彼,是基於對最後一批及其他工人的關愛.也知道與第一批工人是講定價才開工的,並無虧待他們.你又知道家主是主人,他擁有使用權,賜予權.你們便會反過來欣賞這位主人!

天國(Kingdom of God)像葡萄園,是上帝管理的地方,天國觀與世界觀不同:

1)世俗觀:著重法律規條,講功德,報酬.天國觀:看重恩典及個人的需要.

2)世俗觀:多勞多得,先到先得,主人賺多,工人賺少→比較,批評

天國觀:多勞不多得；先到不先得,主人賺少,工人賺多→恩典

3)世俗觀:看你的表現,能力與 專長,貢獻,看你的存在價值→比較,批評.天國觀:上帝看你是一個人,是他的子民,他的兒女,看你的需要,看你的強弱,你的缺乏,你的不足→ 你需要的恩典.

願我們相信上帝,進入上帝國,並且把上帝的國切實地行出來!

親愛的弟兄姊妹,今天我們的社會或教會,聰明人太多；老實人太少.感覺自己拼命做事的人多；知道自己軟弱的人少.看到自已需要的人多；顧及別人需要的人少.看到對自已不公平的多；發現對別人不公平的人少.講求規矩條例的人的多；酌情給強勢者的人少.比較,抱怨的人多；慷慨施予的人少.容易眼紅,計較報償的人太多；感恩,默默服事的人太少.這可能是馬太福音所描述的初期教會,也有可能是我們今天的教會與社會景況,值得我們反思.



\newpage

\section{你仗著什麼權柄做這些事}
\label{sec:262}
\textbf{鄧瑞強博士 - 你仗著什麼權柄做這些事?}
\newline
\newline
\hyperref[sec:261]{< < < PREV SERMON < < <}
~
\hyperref[sec:toc]{[返目錄]}
~
\hyperref[sec:263]{> > > NEXT SERMON > > >}


馬太福音 21:23-27
\newline
\begin{longtable}{cl}
\hline
\hline
章節 & 經文 (和合本修訂版)\\
\hline
21:23 & \begin{tabularx}{0.7\textwidth}{X} 耶穌進了聖殿，正教導人的時候，祭司長和百姓的長老來問他：「你仗著甚麼權柄做這些事？給你這權柄的是誰呢？」 \end{tabularx} \\ \\ \relax
21:24 & \begin{tabularx}{0.7\textwidth}{X} 耶穌回答他們說：「我也要問你們一句話，你們若告訴我，我就告訴你們我仗著甚麼權柄做這些事。 \end{tabularx} \\ \\ \relax
21:25 & \begin{tabularx}{0.7\textwidth}{X} 約翰的洗禮是從哪裡來的？是從天上來的，還是從人間來的呢？」他們彼此商議說：「我們若說『從天上來的』，他會對我們說：『這樣，你們為甚麼不信他呢？』 \end{tabularx} \\ \\ \relax
21:26 & \begin{tabularx}{0.7\textwidth}{X} 若說『從人間來的』，我們又怕眾人，因為大家都認為約翰是先知。」 \end{tabularx} \\ \\ \relax
21:27 & \begin{tabularx}{0.7\textwidth}{X} 於是他們回答耶穌：「我們不知道。」耶穌也對他們說：「我也不告訴你們，我仗著甚麼權柄做這些事。」 \end{tabularx} \\ \\
[1ex]
\hline
\hline
\end{longtable}
講題:你仗著什麼權柄做這些事?  By What Authority Are You Doing These Things?

經文:馬太福音21:23-27

講員:鄧瑞強博士

場合:香港中文大學崇基學院禮拜堂

日期:2014年9月28日

各位弟兄姊妹,早安.

若果我在這裡講道,忽然間,有人擲上一條香蕉.這時,我首先要搞清楚,這條香蕉究竟是用來擲我的,還是用來擲牧師的.若這條蕉是用來擲我的,則我會看看這條蕉是否一條好蕉.若是,則我便可以安心地吃掉它.然後,我會問擲蕉的人:「你為什麼向我擲蕉呢?」我想,擲蕉的人應該是不太滿意我在這裡講道.他的潛台詞是:「你憑什麼在這裡講道?」面對這問題,我大致的回答是:「是這禮拜堂的牧師請我來的.」他若進一步問:「他憑什麼權柄這樣做?」我大致的回答是:「你又憑什麼權柄質疑他這樣做?」這種質疑來來回回,最後便要提出大家所憑藉的權柄的終極根據.我大致訴之於制度,這裡的制度是向來如此的,禮拜堂的牧師有權柄請他想請的人來講道.這個人若要反對這制度,則要提出一些偉大的理由.例如:訴之於他的「聽道權」,他有自由去選擇聽誰或不聽誰.這種「聽道權」又根植於他那不可磨滅的自主權,這種自主權是神聖不可侵犯的,也是他的終極理由了.為了表達這種自主權,他向我擲了一條蕉.為了表達我有我的存在理由,我吃了這隻蕉.

兩種行動,我在這裡講道,他在這裡擲蕉.各有其理由,各按其理由賦予的權柄行動,彼此問對方:「你仗著什麼權柄做這些事?」

在一個沒有共識的社會,在一個社群撕裂的社會,在一個不同群體追求不同價值的社會裡,追問對方:「你憑什麼這樣做?你仗著什麼權柄做這些事?」是經常發生的.這是權力的較力,或者說,是不同的意義追求的比拼.不同的「意義追求」,終極的說,其實就是追問「誰是上帝」.

在一個眾人共同生活的社會裡,一個公共行動,不能是完全任意的.這個公共行動要有意義,背後必須有一意義結構.這意義結構支撐著這公共行動,賦權給行動者.若一個人要作一些重大的公共行動,其背後的意義結構必須偉大得能賦予他行動所需的十足權柄.普通事物,很難賦予人行動的十足權柄.一般而言,要讓人行動有十足權柄,要將這意義結構提升到宗教層面才可.唯有神聖的意義結構,才能給人無限的權柄,去實踐重大的事情.

無論如何,唯有神聖之力,才能叫一個人成為不顧一切的「聖戰」戰士.唯有神聖之力,才能叫一架坦克毫不猶疑地衝向一個手無寸鐵的身軀.當然,也唯有神聖之力,才能叫人奮不顧身,面向坦克,或者,走向十架.

就算在一個世俗化的社會裡,不同的,重大的公共行動,背後的意義結構都得披上宗教的外衣,才能讓其授權的行動,帶有神聖的光環.沒有神聖光環,其授權的行動就沒有說服力.帶有神聖光環的東西,都沾上「神聖不可侵犯」的宗教性質.例如:神聖不可侵犯的領土,神聖不可侵犯的國家主權,神聖不可侵犯的人權,神聖不可侵犯的自由,神聖不可侵犯的民族精神.每一樣都是「神聖不可侵犯」的,每一樣都是終極而不需要再解釋的,每一樣都能賦予行動的人無限的行動權柄.只要一帶上神聖的光環,行動者就有無限的權柄去行動,義無反顧,勇往直前.

「神聖不可侵犯」的領域,就是宗教的範圍.宗教涉及的,便是「誰是上帝」的問題.有些上帝,會要求人獻人為祭.有些上帝,會要求人放下屠刀.在社會上,不同社群彼此追問對方:「你仗著什麼權柄做這些事?」終歸而言,就是追問什麼是對方背後的「神聖不可侵犯」的「上帝」.問題是:這個「上帝」是要求人獻人為祭,泯滅人性呢?抑或是要求人放下屠刀,重視人的尊嚴呢?

在耶穌時代,有一天,當時的權力核心問了耶穌一個重要問題:「你仗著什麼權柄做這些事?給你這權柄的是誰呢?」

主耶穌做了什麼事,令這些權威人士提出這質問呢?

讓我們打開聖經:馬太福音21: 12-16.這段經文先於我們剛才讀的福音書經課,給我們提供一些背景資料.我們看看主耶穌在耶路撒冷的聖殿,做了些什麼?

馬太福音21: 12-16

21:12　耶穌進了神的殿,趕出殿裡一切做買賣的人,推倒兌換銀錢之人的桌子,和賣鴿子之人的凳子；

21:13　對他們說:「經上記著說:我的殿必稱為禱告的殿,你們倒使它成為賊窩了.」

21:14　在殿裡有瞎子,瘸子到耶穌跟前,他就治好了他們.

21:15　祭司長和文士看見耶穌所行的奇事,又見小孩子在殿裡喊著說:「和散那歸於大衛的子孫!」就甚惱怒,

21:16　對他說:「這些人所說的,你聽見了嗎?」耶穌說:「是的.經上說『你從嬰孩和吃奶的口中完全了讚美』的話,你們沒有念過嗎?」

當時,主耶穌在耶路撒冷,準備守逾越節.

耶路撒冷當時和現在都是一樣,是一個政治敏感而容易引起衝突的地方.這是當時的強國人羅馬人統治的殖民地,但這塊殖民地的人的自決心很強.他們日日等待彌賽亞,等待一個救主般的人物,渴望民族的自決.這裡有很多革命思想,也有很多群眾運動.

逾越節,是一個敏感時刻.這是猶太民族的一個重要集體回憶的時刻.他們在這日記念神對他們的救贖,記念神解放他們脫離奴役.這日,猶太民族心靈亢奮.這日,羅馬士兵嚴陣以待.

敏感的地方!敏感的時刻!

主耶穌就在這時,進入聖殿.聖殿是宗教,經濟,社會權力的核心.有點似香港的中環.中環,我們這個城巿的聖殿.

主耶穌首先趕走一切在這裡做買賣的人.這些買賣是合法的,也是聖殿運作上必需的.獻祭者要在這裡買祭牲,捐獻者要在這裡兌換聖殿專用的錢幣,一些沒有頭像的錢幣.聖殿的運作,需要這些買賣活動.但在主耶穌眼中,在主耶穌的意義框架裡,這一切活動都是「賊」的活動.用我們現在的講法說,這一切活動只展示「經濟霸權」,是合法的,但不合人性.主耶穌引用神聖的經文說:「我的殿必稱為禱告的殿」.一個城巿的真正中心,應是「禱告」的地方,應是人謙卑下來而尋求永恆意義的地方.

「禱告」,是「對話」,是人與天的對話,是卑微的人領悟永恆的天意的對話.保持這種對話精神,一個人才不會將另一個人視為「經濟剝削」的對象,而視為共同生活,體現天意的夥伴.

在聖殿裡,在這個城巿的中心,主耶穌醫好了瞎子和瘸子.很奇怪,這裡有瞎子和瘸子,因為瞎子和瘸子是不容許進入聖殿範圍的.就像我們的城巿中心,我們中環的聖殿,是不歡迎社會上的弱者的.這些人,應該是跟隨主耶穌,走進聖殿的.主耶穌沒有趕走他們,沒有譴責他們,主耶穌醫好他們.

在那裡,還有一群小孩子喧喧鬧鬧,他們讚頌耶穌.小朋友,在當時社會裡,不被視為一個整全的人.他們的存在,是沒有重要性的.一班毫無見識的小朋友,看來參與了某種群眾運動.聖經說,祭司長和文士(法律專家)對這些小孩子的舉動「甚惱怒」.主耶穌的回應是:「嬰孩的口,也能讚頌真神的」,或者說:「連最無知的小孩,也是懂得說國王是沒有穿衣的」.

明顯,主耶穌在敏感的城巿,在敏感的時刻,身處城巿的中心,阻礙了城巿的運作.他擾亂了當地的經濟運作,也批判了這運作的「盜賊」本性.他接納那些不准進入這城巿中心的弱者,醫好他們,重建他們的人性.他引發了小孩子的呼喊,這些看似無知的人講出了真相,但權力核心不容許這真相.

從維穩的角度看,主耶穌引發一場混亂.盜賊,不喜歡被人揭露為盜賊.淘汰弱者的社會,不喜歡瞎子和瘸子.習慣虛偽的社會,不容許揭露真相.

於是,耶路撒冷的權力核心正面質問主耶穌.

我們看看今日的講道經文:馬太福音21: 23-27.

21:23　耶穌進了殿,正教訓人的時候,祭司長和民間的長老來問他說:「你仗著什麼權柄做這些事?給你這權柄的是誰呢?」

21:24　耶穌回答說:「我也要問你們一句話,你們若告訴我,我就告訴你們我仗著什麼權柄做這些事.

21:25　約翰的洗禮是從哪裡來的?是從天上來的?是從人間來的呢?」他們彼此商議說:「我們若說『從天上來,』他必對我們說:『這樣,你們為什麼不信他呢?』

21:26　若說『從人間來』,我們又怕百姓,因為他們都以約翰為先知.」

21:27　於是回答耶穌說:「我們不知道.」耶穌說:「我也不告訴你們我仗著什麼權柄做這些事.」

祭司長,是這城巿的權力核心,他維繫著城巿的經濟運作,平衡著政治權力,而這一切又靠賴他的宗教身分,他是神聖的祭司長.他身邊是「民間的長老」,這是當時的民意代表.這一個「經濟,政治,民意,宗教」權力的集合體,在那個風雨飄搖的城巿裡,我想,他們最關心的,是如何維持安定,繁榮.那管經濟是剝削的,政治是壓迫的.他們的宗教授予他們這樣做的權力.在這權力裡,他們質問耶穌:「你仗著什麼權柄做這些事?給你這權柄的是誰呢?」

這真是一個好問題.主耶穌憑什麼在城巿的中心引起混亂.

當然,在主耶穌的角度看,他們又憑什麼權柄,維持城巿這樣運作下去.

權力的背後,是某種意義結構.意義結構的背後,是某種「神聖不可侵犯」的東西.這「神聖不可侵犯」的東西,就是人賴以存在的宗教.

追問權力來源,終極而言,就是追問「你的上帝是誰?」你效忠誰?

主耶穌以一個問題,去回答這個問題:「約翰的洗禮是從哪裡來的?是從天上來的?是從人間來的呢?」

約翰的洗禮,是一種叫人悔改的洗禮.這悔改是要求人有行動的.約翰對要受洗的人說:「有兩件衣裳的,就分給那沒有的.」這是愛.對要受洗的稅吏說:「除了例定的數目,不要多取.」這是公義.對要受洗的兵丁說:「不要以強暴待人,也不要訛詐人.」這是和平.簡單而言,約翰的洗禮,要求人放棄種種墮落的行為,放棄對人的漠視,放棄貪污,放棄以政治特權欺凌無權者,實踐愛,公義,和平.

主耶穌問:「這一切要求,背後的權柄,來自天上,抑來自人間?」

或者我這樣表達:「這一切要求,是否有神聖性,抑或只是地底泥?」

這群權力核心彼此商議:「我們若說『從天上來,』他必對我們說:『這樣,你們為什麼不信他〔指約翰〕呢?』」

即是說,若他們認為那些東西是神聖的,那麼,他們為何反對提倡這些東西的人呢?

他們又說:「若說『從人間來』,我們又怕百姓,因為他們都以約翰為先知.」

即是說:若他們認為那些東西只是地底泥,一文不值,他們又怕「民意」反對他們,因為民意明顯地認為這些東西是神聖的.

這群權力精英不敢承認「愛,公義,和平」,是有神聖意義的.但他們又不敢否認這些,因為他們怕被民意唾棄.

他們回答耶穌說:「我們不知道.」

這群權力核心是痛苦的,他們若有良知,會明白真理在哪裡.他們不敢承認真理,也不敢否定真理.他們生命的神,令他們的眼光注視塵世.

主耶穌展示了他相信的神,這個神給他行動的權柄.這位神曾經賦權給約翰去施洗.賦權給約翰,使他要求受洗的人,要徹底改變生命,要實踐愛,公義,和平.這位神也賦權給主耶穌,使他批判那「盜賊」一樣的經濟制度,使他重建弱者的生命,使他讓小孩子講出事情的真相.

可能,有一天,當你實踐愛,公義,和平時,有人會問你:「你仗著什麼權柄做這些事?給你這權柄的是誰呢?」

你會如何回答呢?

但願:榮耀歸於聖父,聖子,聖靈.阿們.



\newpage

\section{當枯萎葡萄吐新葉}
\label{sec:263}
\textbf{伍渭文牧師 - 當枯萎葡萄吐新葉}
\newline
\newline
\hyperref[sec:262]{< < < PREV SERMON < < <}
~
\hyperref[sec:toc]{[返目錄]}
~
\hyperref[sec:264]{> > > NEXT SERMON > > >}


講題:當枯萎葡萄吐新葉　When the Faltering Vine Rebounds

經文:詩篇80

講員:伍渭文牧師

場合:香港中文大學崇基學院禮拜堂

日期:2014年10月5日

引言

「那是最美好的時代,那是最糟糕的時代；那是智慧的年頭,那是愚昧的年頭；那是信仰的時期,那是懷疑的時期；那是光明的季節,那是黑暗的季節；那是希望的春天,那是失望的冬天；我們全都在直奔天堂,我們全都在直奔相反的方向.」

這是狄更斯(Charles Dickens) 名著雙城記膾炙人口的開場白,上一次以此開場白為講章的引言,已經是二十五年前的六月了.雖然這段開場白也適切描述香港當前的處境,但我是用來描述今天的講道經文詩篇八十篇的主題-復興.

復興副歌.復興是今日講道經文詩篇八十篇的主題,出現三次的副歌(refrain) 都為了復興向上主祈禱 :「神啊!求你使我們回轉(復興, Restore us, O God, NIV),使你的臉發光,我們便要得救.」(v.3)「萬軍之神啊!求你使我們回轉(復興)……」(v.7)「耶和華萬軍之神啊!求你使我們回轉(復興)……」(v. 19)而且透過上主愈來愈莊嚴的稱謂加強表達他們的懇求.

詩篇的背景提到以色列人被外敵欺凌,國破家亡,在黑暗的處境,和充滿眼淚的日子,他們呼求上主的拯救,反思民族的命運,幡然醒悟,明白到苦難的根由是由於叛離上主,背棄真道,因而回轉,得到復興.

復興渴想.復興(restoration) 是今天常常被人談論的詞句.在香港,我們談論要復興獅子山下的拼搏精神；現今香港人太好逸惡勞了,香港若有明天,就要復興獅子山下的拼搏精神.在中國大陸,習近平主席一上台就以復興中國為夢想:「中華民族偉大復興,實現國家富強,民族振興,人民幸福.」因為中國已大國崛起,財力雄厚,羽翼已豐.但聖經的復興,始於屬靈的覺醒,感到自己生命脆弱無助,哀求上主憐憫救拔.

大時代中的覺醒      vv. 1-7

雙城記的開場白排行最佳小說開場語首位,描述法國大革命前夕新舊交替的大時代.大時代之稱為大時代因為它擱置了慣常思想方法,衝擊我們平時的行為習慣,讓出空間,去思考人生和價值意義.大時代令人不安,因為出現了和慣常不同的世界,使我們失去了安全感.崇基神學院師生這段時間嚮應罷課,讓我們擱置慣常生活節奏,反思人生的次優和所處的社會,這點有積極的教育意義.

詩篇八十篇也談及大時代:「耶和華─萬軍之神啊,你向你百姓的禱告發怒,要到幾時呢? 你以眼淚當食物給他們吃,又多量出眼淚給他們喝.你使鄰邦因我們紛爭；我們的仇敵彼此戲笑.」(vv. 4-6)「你為何拆毀這樹的籬笆,任憑一切過路的人摘取?林中出來的野豬把它糟蹋,野地的走獸拿它當食物」(vv. 12, 13)「這樹已經被火焚燒,被刀砍伐.」(v. 16上)國破家亡,聖殿被毀,被擄異域,外族鐵腕統治；慣常所見的一切不單擱置,而且徹底改變,一去不返；人生何等無常,生命何等脆弱,此刻他們心靈騰出空間,反思與上主的關系,因而得到屬靈的覺醒.我們看看這詩篇所描述的歷史背景.

詩篇的處境.「領約瑟如領羊群之以色列的牧者啊,求你留心聽!坐在基路伯的啊,求你發出光來!在以法蓮,便雅憫,嗎哪西前面施展你的大能,來救我們.」(詩八十1, 2)

「在以法蓮,便雅憫,嗎哪西前面施展你的大能」指以色列國(北國)的陷落.雅各有十二兒子,成為以色列十二支派,之後分成南北兩國；北國十支派以以法蓮為首,南國猶大和西緬兩支派,先後因為拜偶像被上主管教.北國首都撒瑪利亞在元前722年被亞述國(Assyrians)(現伊拉克境)攻陷,國民被擄,留在國土的被迫和異族混婚,成為後來的撒瑪利亞人.南國也在公元586年聖殿被毀,百姓被擄巴比倫(現伊拉克境).直到波斯帝國(現伊朗境)興起,公元536年古列(Cyrus)王讓被擄的猶太人回歸故土,修築城牆,學習律法,重建聖殿.這詩篇是集體的哀歌,也是個別信徒面危難時要學習的禱告.

危難的禱告.第一,二節用了三個生動的祈願:「求你留心聽!」上主好像充耳不聞；「求你發出光來!」我們在黑暗中不知下一步如何,找那一個醫生提供第二個意見才決定做不做手術呢?「求你施展大能」(NIV, Awaken your might. Eugene Peterson: Get out of the bed-you ‘ve slept long enough! Come on the run before it’s too late. ),上主好像睡著了,快醒來,再遲來不及了.像福書中門徒的船遇到風暴,船滿了水,耶穌在船尾枕著枕頭睡覺,門徒叫醒了他:「夫子!我們喪命,你不顧嗎?」(可四37-39)

親密的關系.詩人也用了兩個形容詞說明與上主關系親密:他是「帶領約瑟的牧者」的拯救者.約瑟在埃及作宰相,在豐年期間儲糧,好在饑荒年間救了雅各十二支派.也是「坐在基路伯」的施恩者.至聖所內置約櫃,存放兩塊法版,約櫃頂稱施恩座,有神獸基路伯守護,施恩座原文是遮蓋(cover) ,因每年大祭司用祭牲的血灑蓋其上為人贖罪.

因為關系親密,詩人知道上主一定垂念我們的苦況:「你以眼淚當食物給他們吃,又量出眼淚給他們喝.」(v. 5)希伯來書作者就用「坐在基路伯」的意象,鼓勵我們禱告:「所以我們祇管坦然無懼地來到施恩的寶座前,為要得憐恤,蒙恩惠,作隨時的幫助.」(來四16)

復興不是因為強盛,乃認識生命的脆弱,遇到不能承受生命的重.有一次崇基周會安排了龍應台演講有關她的新書「大江大海」-民國時期的內戰,講演前讀經選讀登山寶訓的八福.龍應台一開始就說:要是我們中國人懂得為戰爭的苦難哀慟和對難民憐恤,中國就有福了.是的,在戰爭的大時代中,我們對福氣有不同的定義.在九一一恐怖襲擊中,紐約世貿雙塔被烈火焚燒快要倒下時,頻臨死亡的人撥出最後的一個電話,就是給家人說:我愛你們.劫後餘生的人,生命優次不再一樣了.

生命意義的求索 vv.8-15

在哀求上主憐憫他們當下的苦情,施行拯救後,詩人回顧過去民族的歷史.

8.你從埃及挪出一棵葡萄樹,趕出外邦人,把這樹栽上.

9.你在這樹根前預備了地方,他就深深扎根,爬滿了地.

10.他的影子遮滿了山,枝子好像佳美的香柏樹.

11.他發出枝子,長到大海,發出蔓子,延到大河.

12.你為何拆毀這樹的籬笆,任憑一切過路的人摘取?

13.林中出來的野豬把他糟踏；野地的走獸拿他當食物.

14.萬軍之神啊,求你回轉!從天上垂看,眷顧這葡萄樹,

15.保護你右手所栽的和你為自己所堅固的枝子.

像一葡萄樹被上主特意栽種,上主的選民何等蒙福.它的枝葉曾經豐茂,他們經歷過出埃及過紅海的悲壯,進入迦南因有神助徵戰所向披靡,他們目睹過大衛所羅門王朝盛世,示巴女王也慕名遠道求見,現在何竟凋萎如斯?被他們視為污穢的外族如野豬任意糟蹋,情何以堪!情何以堪!

聖經常用葡萄樹比喻上主的選民,今天舊約經課以賽亞書五1-7提到本來上主苦心經營的葡萄園,現在竟然結出野葡萄,所以園主把葡萄園拆毀,重新建造:「我必撤去籬笆,使它被吞滅；拆毀牆垣,使它被踐踏.我必使它荒廢,不再修理,不再鋤刨,荊棘蒺藜倒要生長,我也必命雲不降雨在其上.」(賽五5, 6)福音書經課太二十一33-46,以上主為擁有葡你園的園主,造好葡萄園所有設施遠行,交給受託的園戶打理,派僕人按時候收果子時,卻被兇惡的園戶打傷,打死,最後差遣兒子來也被殺死,猶太人就是兇惡的園戶.

葡萄樹的比喻使我們想起詩篇第一篇的義人:「他要像一棵樹,栽在溪水旁,按時候結果子,葉子也不枯乾.」「詩一3」我們不是野樹,是上主有目的栽種的樹,所需的空氣,雨露上主供應,結出的果子和茂葉,供他人果腹和遮蔭.用葡萄樹比喻選民很有意思,葡萄樹本身沒有美麗的花朵吸引人,它不需要很多種植面積也可以生長,因它向上生長.繞著一枝柱子,一面牆,或一道籬笆就可以枝葉茂密,很多不起眼的籬笆和欄桿,都以葡萄枝葉來美化；不好看的場景,也利用豐茂的葡萄樹濃蔭遮擋.基督徒要關心被人忽略,邊緣化的弱勢族群,給他們冠冕,得到作為人的尊嚴.因為作在弟兄中一個最小的身上,就是作在基督的身上.(太二十五40)

復興不是追求個人靈性的滿足,復興者會尋求生命的意義,和在所處時代可有的角色.

復興始於和上主修和          vv. 14- 18; 參八十一10-13

「萬軍之神啊,求你回轉,從天上垂看,眷顧這葡萄樹,保護你右手所栽的和你為自己所堅固的枝子.這樹已經被火焚燒,被刀砍伐；他們因你臉上的怒容就滅亡了.願你的手扶持你右邊的人,就是你為自己所堅固的人子(子民,思高版).這樣,我們便不退後離開你；求你救活我們,我們就要求告你的名.」(詩八十14-18)

從當前充滿眼淚的痛苦,回想過去上主的看顧和恩典,最後詩人進入問題的核心:國破家亡是上帝的審判,因為我們背離上主,上主藉外族強權管教我們.「這樹已經被火焚燒,被刀砍伐；他們因你臉上的怒容就滅亡了.」我們的敵人是上主,我們與上主為敵.祇要我們回轉歸向上主,求告上主,那些壓迫我們的外族也速速滅亡.

從最初盼望上主復興以色列國,從眼見的外在政治國度,詩人帶我們進到個人心靈的覺醒.祇要每個選民心靈覺醒,回歸上主,就得到復興；當每一個人復興時,整個國家也會復興起來,這是屬靈的軟實力.詩人關注的是子民的靈性復興,我們從詩篇的編輯也看到這信息.

詩篇的排列.打開聖經,中間就是詩篇；打開詩篇,中間是第三卷,詩篇的編輯者仿效摩西五經把詩篇分成五卷；打開第三卷,有三組詩篇,今天詩篇經課,就取自中間的一組-七十九,八十,八十一,八十二,八十三篇,談到以色列人國破家亡,被邪惡強權欺凌的苦痛,祈求好像不理他們的上主回轉拯救他們,是典型的哀歌詩篇 (lament psalms).這五首詩V型排列,首尾呼應,中間平行,突出單獨的81篇作為主要信息,是詩篇常用的文學手法.

被列國蹂躪.首尾七十九篇和八十三篇提到列國蹂躪他們:「神啊,外邦人進入你的產業,污穢你的聖殿,使耶路撒冷變成荒堆,把你僕人的屍首交與天空的飛鳥為食,把你聖民的肉交與地上的野獸,在耶路撒冷周圍流他們的血如水,無人葬埋,我們成為鄰國的羞辱,成為我們四圍人的嗤笑譏刺.」(詩七十九1-4)詩八十三3, 4:「他們同謀奸詐要害你的百姓,彼此商議要害你所隱藏的人.他們說:來吧,我們談及將他們剪滅,使他們不再成國!使以色列人不再被人記念!」

祈求復興.中間八十,八十二篇分述求上主復興被踐踏的家國及求審判施暴的強權.八十篇:「施展你大能,來救我們,神啊,求你使我們回轉(或譯復興)……萬軍之神啊,求你使我們回轉(復興)……耶和華萬軍之神啊,求你使我們回轉(復興).」(詩八十2, 7, 19) ；八十二篇:「神站在有權力者的會中,在諸神中行審判……你們要死,與世人一樣,要僕倒,像王子中的一位.神啊,求你起來審判世界,因為你要得萬邦為業.」(詩八十二1, 7, 8)

聽從真道.但經文的重點不是強權被審判,也不是國家在政治上的復興,乃在以色列人作為上主子民的靈性上醒覺.整組的五篇詩篇像V字排列,突顯中間的八十一篇,其鑰節就在十三,十四節:「甚願我的民肯聽從我,以色列肯行我的道,我便速速治服他們的仇敵,反手攻擊他們的敵人.」(詩八十一13,14)

是人的回轉.詩篇八十篇復興一字可譯作回轉與復興,開始時詩人哀求上主回轉,拯救他們脫離異族的壓迫,滅消敵人,復興失去的榮耀.慢慢地他們愈來愈明白,回轉的不是上主,回轉的是子民,祇要他們歸回上主,靈性覺醒,就得到復興.

不斷的回轉.復興不是參加一個晚會,有令人振奮的音樂和講道,one shot就搞掂.枯萎的葡萄由來有自,為何曾經枝葉茂盛的葡萄會枯萎呢?可能茂盛一時的表象,樹幹裡面已埋下害蟲的卵子,可能土壤沙石太多,不夠肥沃,不能持續茂盛下去,可能沒有經常灌溉,可能旁邊雜草太多,還有荊棘的種子.植物不會突然長起來,也不會突然枯萎的.要枯萎的葡萄再吐新葉,就要常常除去野草蒺藜,鬆土施肥,勤於灌溉.

蜘蛛的網羅.有人形容沈溺罪惡就像昆蟲陷進被蜘蛛的網羅,是慣常自願投進去的.但房間能夠讓蜘蛛結網有很多條件的:有一定的濕度,一定的幽暗,一定數量的塵埃,做成這種客觀條件是疏於打理房子,還有一隻蜘蛛,等候捕捉昆蟲的蜘蛛,牠是主謀.保持清潔的房子,蜘蛛就沒有方法結出蜘蛛網了.

豐茂葡萄也枯萎.曾經豐茂的葡萄樹也會漸漸枯萎,提醒我們,大衛所羅門的盛世王朝,需要有大衛和所羅門敬畏上主的人當君王,才能持續不衰.不是制度的完善,乃執行制度的人.亞洲週刊2009年十一月二十六最有一篇文章報導天鵝絨革命二十週年的捷克近況,該文題目為:「政治被黑道劫持--捷克天鵝絨革命二十年的悲歌」.捷克街頭賭場被黑幫用來洗錢,政治被黑道劫持,經濟繁榮背後,捷克成為歐洲貪污重災區,人人懷念哈維爾.

建制核心的見證.反之,在最困難的異教帝國的權力核心,但以理在巴比倫,尼希米在波斯書珊城(Susa),他們日理萬機,但敬畏上主,懷念選民的命運福祉.尼希米得到波斯王祝福,帶領被擄百姓回歸耶路撒冷,重建破爛倒塌的城牆.但以理是大時代被擄到巴比倫在宮廷受訓,供侍朝廷,但堅持不拜偶像,被扔在火裡,蒙上主保護不傷絲毫,之後為王解忘記了的夢, 被王賞識重用.「王阿,你夢見一個大像,這像甚高,極其光耀,站在你面前,形狀甚是可怕.這像的頭是精金的,胸膛和膀臂是銀的,肚腹和腰是銅的,腿是鐵的,腳是半鐵半泥的.你觀看,見有一塊非人手鑿出來的石頭打在這像半鐵半泥的腳上,把腳砸碎.」(但二31-34)但以理看出精金頭甚是可怕的像,其腳是泥,一砸即碎.

在嚴重撕裂的社會,祇有被相方所信任的人,才能推動修和解怨,祇有與上主修和,我們才能有力量推動修和.香港此刻正出現大時代,特別是青少年經這星期所發生的事的洗禮,跟以前不再一樣,人心所需要是靈性的復興,面對這挑戰,我們預備好了嗎?各位,也許此刻我們感到能力不足,甚至靈性枯萎,讓我們回想過去,就如詩人回想過去民族蒙恩的歷史,上帝沒有改變,祇要我們回轉,枯萎的葡萄也會吐新葉.今天書信經課腓三4b-14, 提到一位便雅憫人使徒保羅,因著他回轉悔改,和其他回轉的使徒,就把羅馬帝國顛覆了,不是靠武力,乃見靠福音的大能,以愛勝過仇恨.

求上主復興我們,讓我們枯萎的葡萄,再吐新葉.啊們.



\newpage

\section{生命是一場筵席}
\label{sec:264}
\textbf{鄧瑞強博士 - 生命是一場筵席}
\newline
\newline
\hyperref[sec:263]{< < < PREV SERMON < < <}
~
\hyperref[sec:toc]{[返目錄]}
~
\hyperref[sec:265]{> > > NEXT SERMON > > >}


馬太福音 22:1-14
\newline
\begin{longtable}{cl}
\hline
\hline
章節 & 經文 (和合本修訂版)\\
\hline
22:1 & \begin{tabularx}{0.7\textwidth}{X} 耶穌又用比喻對他們說： \end{tabularx} \\ \\ \relax
22:2 & \begin{tabularx}{0.7\textwidth}{X} 「天國好比一個王為他兒子擺設娶親的宴席。 \end{tabularx} \\ \\ \relax
22:3 & \begin{tabularx}{0.7\textwidth}{X} 他打發僕人去，請那些被邀的人來赴宴，他們卻不肯來。 \end{tabularx} \\ \\ \relax
22:4 & \begin{tabularx}{0.7\textwidth}{X} 王又打發別的僕人，說：『你們去告訴那被邀的人，我的宴席已經預備好了，牛和肥畜已經宰了，各樣都齊備，請你們來赴宴。』 \end{tabularx} \\ \\ \relax
22:5 & \begin{tabularx}{0.7\textwidth}{X} 那些人不理就走了，一個到自己田裡去，一個做買賣去。 \end{tabularx} \\ \\ \relax
22:6 & \begin{tabularx}{0.7\textwidth}{X} 其餘的抓住僕人，凌辱他們，把他們殺了。 \end{tabularx} \\ \\ \relax
22:7 & \begin{tabularx}{0.7\textwidth}{X} 王就大怒，發兵除滅那些兇手，燒燬他們的城。 \end{tabularx} \\ \\ \relax
22:8 & \begin{tabularx}{0.7\textwidth}{X} 於是王對僕人說：『喜宴已經齊備，只是所邀的人不配。 \end{tabularx} \\ \\ \relax
22:9 & \begin{tabularx}{0.7\textwidth}{X} 所以你們要往岔路口上去，凡遇見的，都邀來赴宴。』 \end{tabularx} \\ \\ \relax
22:10 & \begin{tabularx}{0.7\textwidth}{X} 那些僕人就出去，到大路上，凡遇見的，不論善惡都招聚了來，宴席上就坐滿了客人。 \end{tabularx} \\ \\ \relax
22:11 & \begin{tabularx}{0.7\textwidth}{X} 王進來見賓客，看到那裡有一個沒有穿禮服的， \end{tabularx} \\ \\ \relax
22:12 & \begin{tabularx}{0.7\textwidth}{X} 就對他說：『朋友，你到這裡來怎麼不穿禮服呢？』那人無言可答。 \end{tabularx} \\ \\ \relax
22:13 & \begin{tabularx}{0.7\textwidth}{X} 於是王對侍從說：『捆起他的手腳，把他扔在外邊的黑暗裡；在那裡他要哀哭切齒了。』 \end{tabularx} \\ \\ \relax
22:14 & \begin{tabularx}{0.7\textwidth}{X} 因為被召的人多，選上的人少。」 \end{tabularx} \\ \\
[1ex]
\hline
\hline
\end{longtable}
講題:生命是一場筵席　Life is a banquet

經文:馬太福音22:1-14

講員:鄧瑞強博士

場合:香港中文大學崇基學院禮拜堂

日期:2014年10月12日

各位弟兄姊妹,早安.

若果有一天,有人邀請你,「去飲啦」,你會有什麼反應?

生活裡,我們遇到大大小小的邀請,邀請你去赴宴.有些筵席是增進友誼的,有些筵席是浪費時間的,有些筵席卻是與生命的意義息息相關的.有些筵席你不想去,有些筵席你不應去,有些筵席去不去卻是需要你再三思量的.赴不赴宴,是一個複雜的問題.

中國歷史上有很多著名的筵席,其中一場便是著名的「鴻門宴」.話說當時項羽大軍壓境,劉邦不得不赴一場「鴻門宴」.項羽的謀臣范增力勸項羽趁此良機殺了劉邦.在宴會裡,范增傳召項莊,叫項莊在席間舞劍助慶,趁機會刺殺劉邦.正所謂「項莊舞劍,意在沛公」.劉邦那邊,卻有個項伯,也隨即拔劍起舞,以身體阻擋項莊,使其無法接近劉邦.劍在舞,心在跳.生與死在筵席中相遇,禍與福在酒杯間穿插.每個在席的人,都各有盤算,各有行動,張力大得不得了.後來,劉邦借尿遁,才逃過一劫.

一場筵席,涉及主人,涉及客人,也涉及一個主場.誰是主場的主人,有時是顯而易見的,有時卻是惹起爭辯的.誰是客人,這人算是客人抑或算是敵人,有時是顯而易見的,有時卻是惹起爭辯的.若香港是一個主場,若香港人是主人,那麼,大陸自由行的人,是我們歡迎的客人,抑或他們是不受歡迎的人?很多國際難民在香港,我們要設筵席接待他們,抑或告訴他們,這裡沒有他們的筵席?這是一個複雜的問題.

如何善待過客,如何款待客人,在聖經裡,常是一個複雜的神學問題.

所以,當有人說,一場筵席已經擺好了,「去飲啦」,當中涉及的問題,複雜無比.

今日的講道經文,是今日的福音書經課:馬太福音22: 1-14.

這段經文討論筵席的問題,而今日講道的主題是:生命是一場筵席.

太22:1　耶穌又用比喻對他們說:

太22:2　「天國好比一個王為他兒子擺設娶親的筵席,

太22:3　就打發僕人去,請那些被召的人來赴席,他們卻不肯來.

太22:4　王又打發別的僕人,說:『你們告訴那被召的人,我的筵席已經預備好了,牛和肥畜已經宰了,各樣都齊備,請你們來赴席.』

太22:5　那些人不理就走了；一個到自己田裡去；一個做買賣去；

太22:6　其餘的拿住僕人,凌辱他們,把他們殺了.

太22:7　王就大怒,發兵除滅那些兇手,燒燬他們的城.

太22:8　於是對僕人說:『喜筵已經齊備,只是所召的人不配.

太22:9　所以你們要往岔路口上去,凡遇見的,都召來赴席.』

太22:10　那些僕人就出去,到大路上,凡遇見的,不論善惡都召聚了來,筵席上就坐滿了客.

太22:11　王進來觀看賓客,見那裡有一個沒有穿禮服的,

太22:12　就對他說:『朋友,你到這裡來怎麼不穿禮服呢?』那人無言可答.

太22:13　於是王對使喚的人說:『捆起他的手腳來,把他丟在外邊的黑暗裡；在那裡必要哀哭切齒了.』

太22:14　因為被召的人多,選上的人少.」

這段經文表面看,有很多不合理的地方.「不去飲」,幹嗎是殺人放火.一個人沒有穿禮服,幹嗎要將他的手腳捆起,丟他到外邊的黑暗裡去.表面看,有點不合常理,不近人情.

經文討論的,不是日常的筵席,而是筵席的神學問題.筵席,從深層次的意義上說,常常是涉及生與死的.像那場「鴻門宴」,生與死在席間擦身而過.「去飲啦」,有時是一種生死的抉擇.

這段經文,直接來說,是討論「誰能參與天國的筵席」這樣的一個神學問題.它提供的答案,也很簡單.有些人原先是有參與資格的,但他們放棄了,這招致他們的禍患.有些人原先是沒有參與資格的,但他們被恩待,得以參加.簡單的說,這故事的教訓是:珍惜恩典.

現在,我們進一步看看這個比喻.

經文說,「天國好比一個王為他兒子擺設娶親的筵席」,這不是一場普通的筵席,這是一場皇室婚宴.這是王子的婚宴.這個王子,很可能是這個王的繼承人.透過這場婚宴,這個王可能讓來賓認識一下他的王國的繼承人.

王邀請的客人,應該也不是善男信女.按慣例,筵席擺好了,主人便叫僕人去請賓客來.一請再請,但客人毫不領情.經文說:「那些人不理就走了；一個到自己田裡去；一個做買賣去；其餘的拿住僕人,凌辱他們,把他們殺了.」

這有點違反常理.不想赴宴,是不用殺人的.但是,這場筵席可能不是一場普通的筵席,這場筵席的背後,可能是一場權力的角力.我們看看王的反應,王的反應是異乎尋常的.經文說:「王就大怒,發兵除滅那些兇手,燒燬他們的城」.請飲的,與被請的,都異常暴力.

聽說,近日,有人在Facebook表達自己對當下時事的看法,邀請朋友閱讀.朋友完全不接受其觀點,不只是 “Unlike” 他的文章,還 “Unfriend” 他.於是,在他的朋友的Facebook世界裡,他的存在「被消失」了.同時,在他的Facebook世界裡,他的朋友的存在也自動消失.昨天還是朋友,今天便勢成水火.「去不去飲」,真的是會導致某些人在某個世界消失的.

如何理解上述經文,不同學者有不同看法.有學者說,這是指到以色列人不回應耶穌的福音,故此,後來,因他們的罪,他們的聖城耶路撒冷被毀,聖殿被焚燒.這種看法很正路.但也有學者認為,這段經文背後,是涉及兩個城邦之間的權力角力.一個城邦向另一城邦宣示主權,我是主,你是客,但另一城邦不同意,作出反抗.而這又招致這城邦的鎮壓,殺人放火.若果不能理解這場筵席背後的權力角力的話,就不能理解這場筵席為何這麼血腥.

人與人相遇相處,開展一場生命的筵席.但一場筵席,可以同時是決戰的戰場.在這筵席裡,有人是主人,他的權力使他成為主人.他有權設定這筵席的主題.這主題的設定,隱晦地宣示他的主權.他有權挑選他的客人.某某人滿足了他的某種條件,才有資格成為他的客人.主人和客人之間,沒有恩情,有的只是像無情的僱主和無情的僱員之間的合約關係.

在這合約關係裡,一個主人有義務邀請合乎條件的人成為他的客人,進入一場筵席.他的客人也有義務作出某種回應,否則,這個合約關係便破裂了.在這合約裡,在主人眼中,客人永遠是外人,而不是家裡的人.主人的家是封閉的,主人的心是封閉的.主人與客人之間,處在對立的兩極,生命從沒有交流.這場筵席的背後,隱藏著「我要控制你」及「我不想被你控制」的張力.當這張力爆發開來時,筵席便成為生死決鬥的場所.

誇張一點來說,這種筵席類似兩幫黑幫在講數.表面上有酒有肉,稱兄道弟,但台底下卻擺著刀鎗,如箭在弦.

這種筵席不是天國的筵席,這種關係模式也不屬於天國.明顯,馬太福音全然否定這種生命的筵席.

馬太福音講另一種筵席,這筵席有另一種請人赴宴的模式.

太22:8　於是對僕人說:『喜筵已經齊備,只是所召的人不配.

太22:9　所以你們要往岔路口上去,凡遇見的,都召來赴席.』

太22:10　那些僕人就出去,到大路上,凡遇見的,不論善惡都召聚了來,筵席上就坐滿了客.

這是一種「無條件」的邀請.這種邀請放下了所有先決條件,不是他滿足了我的某種要求,我才邀請他.不是因為他好,我才邀請他.不是因為他對我有用,我才邀請他.我邀請他,先於一切條件.任何人都被邀請.甚至經文說:「不論善惡都召聚了來」.

「不論善惡都召聚了來」,這是否很危險?

是很危險的.當你主動伸出手去建立友誼時,是很危險的,因為別人可以拒絕你,或者,他可能拉跌你.當你真誠談戀愛時,是很危險的,因為你將打開心門,失去了保護,你的愛人可以直闖你的內心,搞亂一切.當你和一些不知名的生命打交道時,是很危險的,因為這生命可能是恐佈的Alien(異形),也可能是趣怪的ＥＴ.

當一個醫生,去西非照顧感染伊波拉病毒的病人時,他是很危險的.他其實是邀請別人進入他的生命中,無論是善是惡,是有病毒抑或無病毒,他不計較這些人是誰.他在那裡,展開了一場生命的筵席.

當神邀請人進入這筵席時,神是很危險的.聖餐禮就是這場神人筵席的再現.聖餐禮,源自耶穌與門徒的最後筵席.在席間,耶穌說,我的血為你們而流,我的身體為你們而捨去.為了邀請人進入一場筵席,要了耶穌的命.

當一個主人放下了他的身段,放下了他的既有想法,放下了他對客人的篩選,放下了他對權力的執著,這時,主人和客人之間的權力鬥爭才會停止.因著主人無條件的邀請和接納,客人得以和他建立分享的關係,而不只是合約的關係.筵席變成是友誼的歡聚,分享生命的時刻.唯有在這種歡聚中,客人不再是主人掌控下的客人,而是一個人,一個與他人分享生命的人.

近來,電視上常看到一齣政府的宣傳短片.一位大叔,在一粥檔裡,舀粥給老人家吃.他說:「我不覺得自己在做善事.我有一碗粥便與別人分享,就這麼簡單.我們可以做些事情嘛.小人物就小人物吧.小人物做小事.」看來,他是免費送粥給老人家吃.他再說:「這個年代,『互助』這個詞,好像忘記了.」然後他再說:「有能力幫助別人,是上天賜予的福氣.」的確,我們很多時忘記了,生命本是一場友愛的筵席.這場筵席,本來就是恩典.恩典的筵席,才是天國的筵席.

說到這裡,要處理最後一段經文了.

太22:11　王進來觀看賓客,見那裡有一個沒有穿禮服的,

太22:12　就對他說:『朋友,你到這裡來怎麼不穿禮服呢?』那人無言可答.

太22:13　於是王對使喚的人說:『捆起他的手腳來,把他丟在外邊的黑暗裡；在那裡必要哀哭切齒了.』

這段經文,引起很多爭論.有人說,怎能要求一個在街上召來的人穿禮服?這人可能無時間穿,又或,他可能窮得無禮服可穿.就算不穿禮服,趕走他便算,幹嗎要將他手腳捆起,丟在外邊的黑暗裡.

近來有研究顯示,這段經文的文字與觀念和兩本「天啟文學」(Apocalyptic Literature)的內容相關.這兩本「天啟文學」就是《以諾一書》(I Enoch)及《亞伯拉罕啟示錄》(Apocalypse of Abraham).這兩卷書卷都不在我們的聖經裡,但馬太福音的經文和它們的內容共同分享著當時猶太人的宗教觀念.透過這些聖經外的資料,我們現在較明白馬太福音的文字的含意.長話短說,我將學者的結論講出來.學者認為,這裡講的「禮服」,不是參加筵席的條件,而是一個好好地參與筵席的人會有的成果.這正如一個小學生,入讀小一時是不需有小學畢業時的學問的,但到他畢業時,他應有小六程度的學問.小六程度的學問,不是小一入學時的條件,但應是一個小學生好好學習後的成果.進入恩典的筵席,是不用禮服的,穿禮服不是入席的條件.但在恩典的筵席裡,人應慢慢穿上禮服,好與蒙受的恩典相稱.聖經最後的一卷書《啟示錄》,是一卷具有「天啟文學」性質的書,也有類似觀念.其中講到信徒參與天上的一場筵席(「羔羊的婚宴」),在19:8那裡說:「就蒙恩得穿光明潔白的細麻衣.這細麻衣就是聖徒所行的義.」恩典是無條件的,但信徒也要在恩典裡努力結出仁義的果子,這成為信徒生命的禮服.這不是入席的條件,而是坐席後應有的成果.(這裡的釋經資料參考:David C. Sim: “Matthew 22.13a And 1 Enoch 10.4a: A Case of Literary Dependence?” Journal for the Study of the New Testament No.47 (1992): 3-19)

問題是:很多人在筵席中白坐,一生白活,就像學生們在課室裡望天打卦,虛度光陰一樣.一生在恩典裡,卻毫無蒙恩的人應有的表現.粥水賺了很多,卻沒有舀一碗粥水給人.金錢賺了很多,卻沒有好好施予.才華橫溢,卻沒有服務過一個人.就因為這樣,主問:「你的禮服呢?」或者問:你在恩典裡過了一生,為何半點感恩的行動都沒有?

至於「捆起他的手腳來,把他丟在外邊的黑暗裡；在那裡必要哀哭切齒了」,這是「天啟文學」慣用的手法,去表達他遠離神,遠離神的福樂,遠離神的光明.生命在無意義的黑洞裡打滾.

生命是一場筵席,只要在這裡仍然有人連一碗粥水都沒有,則這場筵席仍有太多不正義的權力,仍有人被阻攔以致不能去領受恩典,而我們的「禮服」仍未著得夠好.

今日講道經文的最後一句說:「因為被召的人多,選上的人少.」我們每一個人都被邀去參與一場生命的盛宴,但有多少人在這生命筵席中能夠真誠的與人暢飲高歌,一同為神的恩典歡慶?

但願:榮耀歸於聖父,聖子,聖靈.阿們.



\newpage

\section{這個上帝不夠棒}
\label{sec:265}
\textbf{陳偉光教授 - 這個上帝不夠棒!}
\newline
\newline
\hyperref[sec:264]{< < < PREV SERMON < < <}
~
\hyperref[sec:toc]{[返目錄]}
~
\hyperref[sec:266]{> > > NEXT SERMON > > >}


講題:這個上帝不夠棒!　This God is not smart enough

講員:陳偉光教授

場合:香港中文大學崇基學院禮拜堂

日期:2014年10月19日

畢業典禮上,請來了一位教會領袖當主禮嘉賓.一個不尋常的安排,加上一篇充滿主觀宗教經驗的演詞,叫所有坐在這莊嚴學術殿堂上的人,無不落在意外與觀望的矛盾心情中.我還來不及為這位宗派領袖的勇氣而雀躍,心裡已湧起一份難以言喻的不安.演說由勸勉年青人謙卑尋求上帝開始,進而談個人的信仰堅持,就這樣,這位長者驕傲地回顧了他的一生,堂而皇之地「見証」服侍上帝的人如何格外「蒙福」.聽他誇口少時讀書雖不太用功,卻往往在考試時成績過人；工作上雖不及同事般能幹,卻又因機緣巧合而獲得上師提拔；人生中一個又一個本來不屬於他的機會,都滾滾地隨他而來；從金錢,聲譽,到家庭,後代,無一不因上帝特別眷佑,而凡事亨通,羨煞旁人.興奮的詞情,隨著講者高漲的情緒,聲線分外嘹亮,充滿著會堂的每一角落；只是,他熱情地傳達的那份「高尚」情操,此刻卻顯得異常冰冷.眼前的「尊貴人」雖近在咫尺,我的心卻早已飛到遠處,遠在那片不屬此世的靈空中,靜觀典禮廳如何繼續被這份自覺高人一等的基督徒優越感所霸佔.有人以福音為榮,因為福音叫他在人前很有體面.我所信的福音與他的不同,我的上帝也肯定沒有像他的「棒」.莫明的納悶忽然從內心升起,我們這些「不屬此世」的人[1],一時間竟有點不知如何自處.

「請把施洗約翰的頭放在盤子裡,拿來給我.[2]」有誰能比這位「預備主的路,修平祂的道[3]」的天國先鋒,死得更不體面?小女孩的一句戲言,施洗約翰就這樣死無全屍,草草埋葬[4].有沒有想過身為父母的撒迦利亞與以利沙伯,此刻如何在人前見証上帝的眷佑?天使不是曾應許過,這兒子「將要為大[5]」嗎?孩子出生的時候,不是曾明明地為父母「除掉人間的羞恥[6]」嗎?不錯,原來上帝所器重的人,其際遇往往連最低微的也不如[7],這是蒙愛的人的專利,是上帝特別留給他們的「荊棘」冠冕,最不體面卻最是光榮.天國的奧秘,在人看為愚拙,在我們卻為上帝的大能.面對這超世的價值,淺薄無知的世代,只懂說某人定必是「受責罰,被神擊打苦待了」[8],耶穌卻權威地宣告:「凡婦人所生的,沒有一個興起來大過施洗約翰的.[9]」

這份被世人輕視的光榮,叫多少曾信誓旦旦地立志跟隨主的人,都慚愧得無地自容.回頭看自己那顛簸跌蕩的基督徒歲月,甚麼信心的危機,信仰的低谷,豈不都是因為所等待的那個上帝,根本就不存在.那不過是一個自己所編造的上帝,一個只許配合我的需要而行事的上帝,還不過是那昔日的老我,假借上帝之名,找個漂亮體面的新名字,好繼續走舊日的路.就這樣,一個又一個稱心如意的計劃,都貼上了「靠主引領」的標籤,插上「因主之名」的旗幟,心裡平安,臉上光彩,自欺欺人.「我好自專,隨意自定程途,直到如今![10]」這John Newman著名的詩句,成了多少基督徒一生的寫照,走足一輩子,到頭來還是像那少年的官一樣,「憂憂愁愁地走了」[11].大概我們從沒認真地想過,要「變賣[12]」些甚麼來跟隨主,且付上如此大的代價後,最終上帝為我們預備的,不過是一份被世人輕視的光榮.

在現實的生活裡,人要為自己的選擇與行為,承擔應有的後果,這是必然的事.今天睡得晚了,明天自然精神不振；這邊我得了額外的好處,那邊就有人蒙受不公的虧損.明明是自己霸佔了別人的領土,卻狡猾地說成,這裡有個三贏的方案,睜著眼把上帝請出來,為自己的貪婪與懶惰添上光環.不是嗎?平日沒有用心學習,竟敢求上帝在重要關頭格外施恩保守；平日生活不像基督徒的父母,竟奢望孩子會乖巧善良,誠實友愛.教會甚麼時候給編造了一個教人「不勞而獲」的上帝,讓我們可以隨心所欲地如願以償?莫說人要為一己的選擇承擔後果,就算事情全屬別人的錯失,我們也從來無權置身事外.為彼此墮落的行為,集體付上代價,這正是人類歷史的縮影.狡猾的人,把自己一手造成的困局,都歸到上帝的賬上,還挑戰上帝「豈不能[13]」大行奇事,如此這般地除去人間的痛苦?連信徒也在等上帝「出手」,自己好光榮體面而全無代價地,由上帝來給他收拾這個爛攤子.

誰會被下在監裡,被置諸死地,而仍舊不以福音為恥[14]?難道會是個終日鼓著一肚子氣,等待上帝為他在人前一雪前恥,揚眉吐氣的人嗎?我們自稱天國兒女,引以為傲,怎可能光把人帶進神的家[15],卻從沒向他們展現如何委身,走窄路,如何預備接受這份世人看為愚拙的光榮?我們一方面為自己定下鴻圖大志,向別人作出種種承諾,然後再把上帝拖下水,說甚麼全交主手,等待祂為我們成就大事；而另一方面,我們又怕自己開出的支票,上帝無法兌現,想著,想著,想到一個不聽禱告的上帝,會叫人多丟臉,也就趕快找來了一大堆漂亮的屬靈藉口,設法為這不濟事的上帝「補鑊」,好意地弄個下台階來替祂開脫.

「禍哉!我滅亡了![16]」每次禱告時,想到這昔日的光景,舊日的愚昧,只有停不了的眼淚,再無半句可說的話.一個以福音為恥的人,既「得罪了天,又得罪了你」[17],能不像那稅吏般,只敢遠遠的站著,連連的捶胸,求上帝開恩可憐嗎?[18]

多少基督徒的人生,就是如此載浮載沉地,走過了一輩子.到教會來「崇拜」的人,口裡雖說等候上帝,心底裡卻對這個信仰感到失望.他們期待著一個大能的上帝,為他們解決人生的種種問題,讓他們可以心悅誠服地「尊崇」,「拜服」.可是,失望得很,眼前的這個上帝實在不夠棒.人會有這種想法,亦非今天的事,主耶穌在世的日子,不也是同樣被藐視,唾棄嗎?然而,更可悲的是,有人想這個上帝既不能成事,何不自己動手為祂添姿抹粉,甚或改頭換面,索性把祂換成一個要多體面,有多體面的上帝,然後自己好一臉尊貴地,當上了代言人.今天,這代言人繼續在推銷著一個現代人夢想中的上帝(dream God),一個我從不認識的上帝.今天,或許教會也得認真思考,如何面對這個dream God的挑戰.我們要教會的「崇拜」禮台,繼續被一個虛擬的上帝所霸佔,甘願由他肆意地吸引人,來「崇尚」,「拜望」一個人手所造的偶像,還是勇敢地把一個世人看為愚拙,我們卻看為神的大能[19]的選擇,擺在眾人面前,忠心地引導他們去領受一份被世人輕視的光榮?至於那些自以為聰明的代言人,他們「以敬虔為得利的門路」[20],早已得了自己的獎賞.

真正的「崇拜」,叫我們無法逃避信心與主權取捨.我們真心相信上帝在我們身上的計劃,是信得過,靠得住的嗎?我們真的願意把這個選擇權,完全交到一個肉眼看不見的上帝手上嗎?我們這群過來人,跌跌撞撞地走到今天,是多麼渴望那真正的「崇拜」,讓我們在其中,也就是在今天我們服侍的教會中,彼此守護心靈,並肩「重建」那已降臨人間的天國.

曾經流蕩的人,今天都齊集在崇拜的禮台前,心裡懊悔不已.想到昔日的自己,「驕癡無忌,舊事乞莫重提![21]」此刻,我們仍舊「遠遠的站著,連舉目望天也不敢[22]」,心裡只連連的說:「向來未曾如此虛心求主,導我前行![23]」

[1]    約十五:19

[2]    太十四:8

[3]    賽四十:3

[4]    太十四:12

[5]    路一:15

[6]    路一:25

[7]    太十一:11

[8]  賽五十三:4

[9]    太十一:11

[10]    參聖詩【導我前行】(又名【慈光歌】,見「頌主聖歌」第240首第二節中段之翻譯).

[11]    太十九:22

[12]    同上

[13]    約十一:37

[14]    提後一:8

[15]    太二十三:15

[16]    賽六:5

[17]    路十五:21

[18]    路十八:13

[19]    羅一:16

[20]    提前六:5

[21]    見註10(參第二節末段之翻譯).

[22]    路十八:13

[23]    見註10(參第二節首段之翻譯).



\newpage

\section{講得好好行得離譜}
\label{sec:266}
\textbf{許開明牧師 - 講得好好,行得離譜!}
\newline
\newline
\hyperref[sec:265]{< < < PREV SERMON < < <}
~
\hyperref[sec:toc]{[返目錄]}
~
\hyperref[sec:267]{> > > NEXT SERMON > > >}


講題:「講得好好,行得離譜」!　Speak Well, Act Badly

講員:許開明牧師

場合:香港中文大學崇基學院禮拜堂

日期:2014年11月2日

從古到今,不難發現有一種人:「講一套,做一套」或「言行不一致」.套用廣東人的俗語,就是「講得好好,行得離譜!」

今日福音經課的法利賽人正是這樣的人,是耶穌所不滿的人.

法利賽人使我想到一個小故事.

一位住在世俗中的神父,看見住在對面的妓女,每天家裡都有許多男士出入.神父就責備妓女:「你是大罪人,每天都得罪上帝,也得罪了那些男人背後的女人.」

那妓女聽到神父的話,感到很不開心,她真誠地向上帝認罪!祈求寬恕,並祈求上帝給她其他謀生之道.

但她找不到其他工作,餓了一星期後,只好重操故業.只是每次想起神父嚴厲的樣子,責罵的聲音,知道自己有罪,都向上帝求赦免.

神父眼見妓女沒有聽他的話,就生氣得甚麼也不做,每天監視著妓女的房子,每見一個男人進入,就放一塊石頭.過了一段日子,就對妓女說:「你看見這堆石頭嗎?每塊石頭都代表你一宗罪過,我要再提醒你,小心你的壞行為,將受到上帝的懲罰.」

妓女聽見後,既自責又恐懼,她流淚禱告:「主阿,求你憐憫我,幫我從不幸的人生中解脫吧!」

上帝聽了她的呼求.當夜取去她的生命,也同時帶走了神父.奇怪的是,神父的靈魂去了地獄,妓女的卻上了天堂.

神父心中不忿,說:「我一直都持守公義,行事不偏不倚,竟然要下地獄?!那個妓女天天犯罪,你卻讓她上天堂!」

上帝回答說:「你天天代我作審判官和代言人,加重這女人愁苦,激發我的怒氣,應下地獄；相反那女人日夜哭泣向我求恕,感動了我憐憫的心,讓她的靈魂上天堂.」

在上帝的規條中,妓女的行業是不正當的.但是,神父只會執掌規條,只會痛罵,指責,沒有體諒妓女不作此業難以維生的殘酷事實.神父只怪責她犯罪,卻沒有幫助她脫貧,教導她正當謀生之道.l這顯得他毫無同情心,兼妄作上帝的代言人.神父口中的規條,只成為束縛他思想行為的枷鎖.

《馬太福音》第二十三章,分為三段.表面好像是針對法利賽人.其實,這些話之目的並非針對他們,主要是要「警戒門徒」,不要效法他們的行為(3).

耶穌所作的事大膽又危險,因為法利賽人是當時最受歡迎和尊敬的人,真是「未死過!」.

效法:是一種學習的方式.唱歌,彈琴,寫字都是透過效法老師或別人而獲得知識或技能.這是一種「有意識」的「效法」.在成長過程中,我發覺我說話速度很慢,原來無意中我效法了「父親」的說話方式.後來,我又發覺我是很勤力的人,原來也是無意中學效了母親做事方式.

法利賽人是拉比(即老師),是許多人效法的對像!因此耶穌特意提醒門徒,不要效法他們不當之處.

第一段從第一節到十二節,記述主耶穌對門徒和眾人的警戒,叫他們不要效法文士和法利賽人的虛假.

第二段從第十三節一直到三十六節,是主耶穌對法利賽人的責備,述說法利賽人的七禍,如果加上第十五節的小字就是八禍.這是遙遙對應主耶穌剛出來傳道時,在登山寶訓裡講的八福.主在地上的職事快結束時,提到文士法利賽人的八禍,這對我們來說是刻骨銘心的教導.

第三段從第三十七到三十九節,主為耶路撒冷憂傷,歎息.

主對門徒與眾人的警戒(太23:1-12)

法利賽人的三個嚴重問題(1-7)

主耶穌在這裡勸導門徒和眾人,千萬不要效法文士和法利賽人的行事為人.主提到他們在人生的態度和表現上有三個很嚴重的問題.

第一個問題,他們能說不能行(“實踐”他們所宣講的)“they do not practice what they preach, NIV”.

1廣東話:“講就天下無敵,做就有心無力”

或“講就咁講,做就人在江湖,身不由己?”

文士和法利賽人:對摩西許多律例規條啷啷上口,自己卻不照著行.

主說他們是能說不能行；只會要求別人,卻不要求自己,這是「寬己嚴人」.

這涉及四種的人生態度:是1)「寬己嚴人」,2)「嚴己寬人」,3)「嚴己嚴人」或 4)「寬己寬人」?

你是屬於哪一類人?哪一種較佳:「嚴己寬人」?即比較言行一致的人.

2他們把難擔的重擔捆起來,捆起來擱在人的肩上.

第四節說,他們把難擔的重擔捆起來.換句話說,不只一點點,而是把很多的擔子捆在一起,擱在別人的肩上,自己卻連一根指頭也不肯動.

政改框架落閘,政府把難擔的重擔捆起來放在香港人的肩頭上,然後說:「袋住先,可能重有得傾,不過框架就無得傾.」

佔中運動,演變成雨傘運動,無收傘的期限,無限期佔中,也一樣是把難擔的重擔捆起來(不是一樣的重擔,是一大堆的重擔)放在政府,警隊及香港市民的身上.

各不退讓,大家「企硬」,面對面,總會期望對方退讓；大家都阻住大家,大家都無路行,大家都無法向前,最痛苦的是大家都不想見到對方；到見面時,你要聽我講, 各自表述.

今日的法利賽人也造成了困局!

許多時候,我們有意或無意中會把難擔的重擔放在別人的肩頭上.

當我們只顧自己的看法,只顧一些規矩原則,甚至自己的理想時；

當我們未能進入別人的處境或想法時,或者是不願意進入別人的處境或想法時,

這種情況往往出現,我們有意無意成為法利賽人的翻版.

好像法利賽人或文士,這種在宗教,社會上有權威的人,很容易把自己當成上帝的代言人或代理人:「上帝,這樣說…」或者「上帝不會這樣做…」

我們看一個人,究竟他所說的和他所行的,是否可以分開?

講一套,行一套的人,為什麼我們不能接受呢?因為言行不一?

主耶穌在這段聖經的開頭指,文士和法利賽人坐在摩西的位置上.是否重要的是他所說的,卻不是他所行!

一個犯了罪的牧師,他講的道,你是否也應聽從呢?

一個不貞潔的牧師,他在講貞潔之道?你是否也應當聽?

一個犯了罪的神父,他施聖餐是否仍然有效呢?

若然,他是坐在摩西的椅上,似乎你就要聽,因為你不是尊重講者,而是尊重他的身份及上主所頒的律法!這個觀點你是否同意?

如果你認為可以言行可分開,不理會他是誰?只要他一坐在摩西的位上,或站上講台,你就要聽?!

今天在猶太人的會堂仍然有一張特別的椅子.聽說有的還是用石頭雕刻出來的,特別讓律法師,文士和法利賽人坐在上面,傳遞摩西的律法,並教導人.

第二個問題,他們一切所做的都要叫人看見.

換句話說,這些人非常會做表面工夫,門面功夫.

他們配戴的經文帶子做得特別寬,衣裳的繸子做得特別長,都是為了叫人看見他們有敬虔的外表.

做給人看,做場戲!

換來好聲譽,好敬虔,好名聲!但這不是堅實的東西!

希望別人注意自己,以提高自己的身價,以搏取別人的欣賞,是今日現代人所推崇的,也是無可厚非的.

今日許多的商品廣告,目的也是叫人看見.今日許多的成功人士,未成功前也要人看見,成功後更要人看見!

如果動機正確,不是為炫耀自已,而是為了推廣.

小朋友穿了件新衣,問:「我靚唔靚?」我回答:「靚!」,他整天就開心了!

一位小學生,說他「醜」,他便不開心一生!

如果人的價值由身材決定,由衣服決定,由樣子決定!多麼悲哀!

人為了叫人看見你,就重視你,尊重你,欣賞你,讚美你.

你的價值只建立在別人所見所聞所言,這種人不會活得開心.

今天在以色列地,特別在哭牆附近,我們仍然可以看得到一些人,他們將經文做成帶子或是一個小盒子,綁在手上,甚至將經文放在一個小黑盒子裡,綁在額頭上,表示他們是非常敬虔的人,注重神的律法和神的話.

我曾認識住在隔鄰的一位老婆婆,經常在門口看聖經,每天看超過十小時,早,午,晚都在看,我們見到她,都讚她,她就很開心.但後來才發現,原來她不懂文字.

第三個問題,他們喜歡得著人的尊崇和稱讚.

主耶穌從很多方面談論這問題.例如:

1)主說他們喜歡筵席上的首位(社會隆重的讌會 ),

2)喜歡會堂裡的高位(宗教集會的場合).

筵席所指的是世上的生活,他們喜歡在那裡坐首位.

會堂是宗教的中心,他們喜歡在那裡坐高位.

3)他們走在路上,街上(公開公眾的場合)時,喜歡別人向他們問安.

他們喜歡被人稱為夫子,師尊,稱為父.

現代人跟法利賽人一樣,喜歡別人注意他,欣賞他,讚揚他.

特別是一些自覺自己很有身份,地位,或很有成就的人,更常有這種渴求.

主責備他們,也教導門徒們不可效法他們.

坐在某個位上?

得到某個名銜?拉比(老師),師尊(導師)

若一個人常常要人擦鞋,欣賞,稱讚,才會開心,這樣的人一定會失望.

沒有人擦,捧,讚時,就失去生命動力.他要依靠別人,才能活得起勁,讓別人去決定他的生命,多麼可憐!

法利賽人是誰?

1)名稱:法利賽人 (Pharisees, 希伯來文 perushim, 字根是parash,意思是「分離」) ,指一些為保持純潔而與俗世保持距離的人[1],與撒都該人追求俗世的權力及物慾相對.

2)組織:是公元前二世紀猶太教中的一個派系,主要由文士,律法師組成.

3)性質:政治上他們代表著一個保守的宗教集團,反對希臘化,擁護馬卡庇(Maccabi)起義(公元前168年),在羅馬管治時期,對於羅馬政府雖然無積極反抗,但採取不合作態度,由於他們最關心的是宗教而非政治,所以祇要羅馬政權不干涉其宗教生活,雙方可以相安無事.

4)信仰:他們很強調維護猶太教傳統與猶太人生活規範,祇與同聲同氣的人交往.又相信靈魂不死,肉身復活,神罰,天使,精靈等.

5) 律法:法利賽人認為上帝交托摩西的律法分為「成文」與「口傳」兩種:

成文律法《妥拉》是對人的神聖啟示,但同時要靠先知的教誨同歷代相傳之口傳律法補充,所以人應該用上帝賦予的理性去解釋《妥拉》,而不應該盲從《妥拉》裡面的文字,在解釋律法的時候,會著眼於當中的精神,所以他們很懂得把《妥拉》裡面的精神配合時代需要,賦予全新意義.

6)影響:法利賽人群體大約在猶太人被擄到巴比倫的時候就已出現；到第二聖殿被毀以後,法利賽教派重新建立為拉比猶太教 ── 成為了日後傳統猶太教的典範,今日猶太教所有教派的根本.

從上可知,法利賽人,他們是受猶太人尊重的人,維護猶太人律法傳統的人.認識聖經與律法,很可能是第一批聖經的詮釋者；強調律法應注重與地區文化的結合,可能是本土神學家.同時,法利賽人注重傳統與現實的結合,並非完全死硬派.

4 法利賽人的根本問題是驕傲.

耶穌說:「凡自高的必降為卑,自卑的必升為高」(太23:12)

驕傲(英語:Pride),又寫為憍傲,是一種內在的情緒狀態,一般而言,有兩大類常見的意涵.

作為負面的意思,驕傲是指一種對於個人的地位或成就的自我膨脹與炫耀,通常與傲慢(hubris)是同義詞.

作為正面的意思,驕傲是一種對於達成目標,或是對於某個選擇或行動的結果,感到滿意的情緒.

它也可以是一種受到讚賞之後產生的情緒,對自我獨立,堅強的形象感到滿足,成為某個社會群體一份子後產生的歸屬感與認同,並因此衍生出滿意的情緒.是對自我認識的一種表現.

「我比別人強」的感覺和表現,那是因為你只看了自己比別人好的一面,還有很多時候,其實你並不是那麼優秀的.所以,正確的認識自己,才不會隨意的比較而驕傲的自以為是. 自大的表現:

我是對的,絕對的對!

我是聰明

我是長輩,你應聽話

驕傲:高高在上的．別人說的話是不一定．他說的話就是肯定!太自我與自大了．一副我是王,我最厲害的樣子．有不可一世的傲慢態度!這是形容負面的驕傲.

「我為你感到驕傲」:這類驕傲是正面的．因為你的超優表現讓別人以你為榮．感覺有你這個朋友或這個小孩感到很棒．是一句肯定的話語!

猶太人認為:世界上有許多不美麗的東西,其中一項是「自大」.

猶太古語有句:「自大的人呀,沒有你,太陽照樣東升,西下?」

猶太法典認為自大是罪,是愚昧.

驕傲的人是自大的人,猶太人說:「如果一個人內心已被自己佔滿時,就不會有留給神和人住的地方了!」

所以謙虛就是「信靠上帝,給上帝地方住,讓上帝有所作為.」

驕傲的人是沒有上帝的!

猶太法典說:「謙虛的人是盡量隱藏自己的優點和功績,就像你隱藏自己作過的壞事或錯事一般.」

謙虛的人同時能夠從積極面來看任何人.

「真正謙虛的人,就是不論遇到什麼人,都認為別人有優於自己的地方,並表示尊敬.

假如遇到的人比自己年長,

就認為他比過去的我優越,因為對方的善行比自己多,因此要表示尊敬.

假如遇到的人比自己年輕,

就認為他比過去的我優越,因為他所犯的罪比自己少,應當尊敬他.

假如遇到的人生活比自己優裕,

就認為這是他比自己更努力而得到的成果,因此表示尊敬.

假如遇到的人生活比自己貧困,

就認為他比自己更堅強面對生活的挑戰,並表示尊敬.

假如遇到的人比自己聰明,

就對他的智慧表示敬佩.

假如遇到的人比不上自己聰明,

就認為他所犯的錯誤比自己少,因為聰明總被聰明誤,因此並表示尊敬.」

猶太古語認為「謙卑像葡萄,愈是結實纍纍的葡萄,愈會垂下頭來；所以愈是偉大的人,就愈平易近人!」

猶太法典:「謙卑人像流動的水是由高處向下流的；這種水保持清潔和美麗:不動的水是死水,會堆積髒東西,或會滿瀉,沾污別人.」

結語

【發牌的是上帝】

美國總統艾森豪威爾(Dwight David Eisenhower,1890年10月14日－1969年3月28日),是著名的總統和軍人,大戰英雄.年青時常跟家人一起玩紙牌.

有次,他連續執了幾副很壞的牌,很不開心,抱怨自己運氣很差.

他媽媽停下來,嚴正地說:「如果你要玩,就必須用手上的牌玩下去,不管你的牌怎樣?」停了一陣又說:「做人也一樣,發牌的是上帝,不管是怎樣的牌,你都要拿著,你能做的就是盡你的全力,求得最好的結果.」

這些母親的話語,艾森豪威爾一直都記著,無論環境順逆,他都以積極的態度去迎接生命的挑戰,盡力做好每一件事.

發牌的是上帝,我們無法改變這人生的牌,當虛心的承受,但盡力運用智慧使它不致太壞,甚至反敗為勝!

各位,你人生的牌又如何?

多數人可能會說,我的不好!你的比較好!

無論如何,謙卑的接受,因為發牌的是上帝,試想你的生與死都不在自己手呢!就連中間那段作息存留也在乎祂呢!

對上帝謙卑,就是相信上帝,依靠上帝做人；對人謙卑是放下自己,欣賞別人,肯定每個人都有存在權利與價值.

不要效法法利賽人四樣東西:

1)能說不能行

2)一切所行的,都要叫人看見

3)喜歡得著人的尊崇和稱讚

4)驕傲自高與自大.

短禱:

天父上帝,求使我們不要學效法利賽人的虛偽,言行不一；更不要變成驕傲自大的法利賽人.靠主耶穌的名求.阿們.

註:(Dwight David Eisenhower,1890年10月14日－1969年3月28日),是美國陸軍五星上將和第34任總統(1953年－1961年).第二次世界大戰期間,他擔任盟軍在歐洲的最高指揮官；二戰時期,軍中的士兵都稱呼他為艾克(Ike).負責計劃和執行監督1944年至1945年裡,進攻維希法國和納粹德國的行動.1951年又出任北大西洋公約組織武裝力量最高司令



\newpage

\section{等待不是靜看流水落花}
\label{sec:267}
\textbf{鄧瑞強博士 - 等待,不是靜看流水落花}
\newline
\newline
\hyperref[sec:266]{< < < PREV SERMON < < <}
~
\hyperref[sec:toc]{[返目錄]}
~
\hyperref[sec:268]{> > > NEXT SERMON > > >}


馬太福音 25:1-13
\newline
\begin{longtable}{cl}
\hline
\hline
章節 & 經文 (和合本修訂版)\\
\hline
25:1 & \begin{tabularx}{0.7\textwidth}{X} 「那時，天國好比十個童女拿著燈出去迎接新郎。 \end{tabularx} \\ \\ \relax
25:2 & \begin{tabularx}{0.7\textwidth}{X} 其中有五個是愚拙的，五個是聰明的。 \end{tabularx} \\ \\ \relax
25:3 & \begin{tabularx}{0.7\textwidth}{X} 愚拙的拿著燈，卻沒有帶油； \end{tabularx} \\ \\ \relax
25:4 & \begin{tabularx}{0.7\textwidth}{X} 聰明的拿著燈，又盛了油在器皿裡。 \end{tabularx} \\ \\ \relax
25:5 & \begin{tabularx}{0.7\textwidth}{X} 新郎遲延的時候，她們都打盹，睡著了。 \end{tabularx} \\ \\ \relax
25:6 & \begin{tabularx}{0.7\textwidth}{X} 半夜有人喊：『看，新郎來了，你們出來迎接他。』 \end{tabularx} \\ \\ \relax
25:7 & \begin{tabularx}{0.7\textwidth}{X} 那些童女就都起來挑亮她們的燈。 \end{tabularx} \\ \\ \relax
25:8 & \begin{tabularx}{0.7\textwidth}{X} 愚拙的對聰明的說：『請分點油給我們，因為我們的燈要滅了。』 \end{tabularx} \\ \\ \relax
25:9 & \begin{tabularx}{0.7\textwidth}{X} 聰明的回答：『恐怕不夠你我用的；你們還是自己到賣油的那裡去買吧。』 \end{tabularx} \\ \\ \relax
25:10 & \begin{tabularx}{0.7\textwidth}{X} 她們去買的時候，新郎到了。那預備好了的，與他進去共赴婚宴，門就關了。 \end{tabularx} \\ \\ \relax
25:11 & \begin{tabularx}{0.7\textwidth}{X} 其餘的童女隨後也來了，說：『主啊，主啊，給我們開門！』 \end{tabularx} \\ \\ \relax
25:12 & \begin{tabularx}{0.7\textwidth}{X} 他卻回答：『我實在告訴你們，我不認識你們。』 \end{tabularx} \\ \\ \relax
25:13 & \begin{tabularx}{0.7\textwidth}{X} 所以，你們要警醒，因為那日子，那時辰，你們不知道。」 \end{tabularx} \\ \\
[1ex]
\hline
\hline
\end{longtable}
講題:等待,不是靜看流水落花　Awaiting Does Not Mean Sitting Still nor Letting Go

經文:馬太福音25:1-13

講員:鄧瑞強博士

場合:香港中文大學崇基學院禮拜堂

日期:2014年11月9日

各位弟兄姊妹,早安.

「等待」抑或「不再等待」,是生命裡的一個重大問題.

有人問我:「我應該不應該和這人拍拖呢?這人不太理想,但我又渴望拍拖.我是否應該多等一會,看看有否更合適的人出現呢?」等待抑或不再等待,一個牽涉幸福的重大問題!

有一次我困在一部升降機裡,那升降機的門半開半掩,但開的程度未夠一人通過.腦海裡閃出兩幅畫面.一幅畫面是:那部升降機忽然下墜.另一畫面是:我拉開那道門,正走過的時候,那部升降機又忽然開動了.我應該什麼也不做,等人救援,抑或想辦法走出去呢?等待抑或不再等待,一個牽涉生存的重大問題!

二十世紀中葉,有一齣話劇叫《等待果陀》(Samuel Beckett, Waiting for Godot,1953年首演),以「等待抑或不再等待」去呈示人生的存在光景.劇中有兩個主角.兩個主角看來都不知應該做什麼?也不知自己在做什麼?甲(Estragon)說:「我們現在幹什麼呢?」乙(Vladimir)答:「我不知道.」甲說:「我們走吧!」乙說:「我們不能走.」甲問:「為什麼不能走?」乙答:「我們等待果陀.」甲有點絕望,然後說:「呀!」

這兩個人在等待一個名叫「果陀」的人.「果陀」是誰,他們不太知道.「果陀」幾時來,他們也不知道.他們等呀等呀,不知等到何時,也不知應做什麼.「果陀」差了一個男孩告訴他們,他今晚不會來,但明天會來.但明天又如何呢?明天「果陀」也沒有來.於是,那段對話又重覆了.甲說:「我們走吧!」乙說:「我們不能走.」甲問:「為什麼不能走?」乙答:「我們等待果陀.」

這是生命的寫照.

我們在等待,等待什麼?等待某些事情的來臨,等待某些理想的實現,等待某個偉大的時刻.這可能是一個漫長的等待,等待的過程可能是單調的.但等呀等呀,「果陀」還沒有來,理想還未實現,偉大的時刻仍未來臨.有人說:「我們走吧!」另一人說:「我們不能走.」「為什麼不能走?」「我們等待果陀.」

然後,這個「等待」抑或「不再等待」的循環又再開始.這個「等待」抑或「不再等待」的疑問又重覆出現.

有人說:「人一出生,便是等待死亡.」

有人等得不耐煩,選擇了斷生命.

這個主題也出現在《等待果陀》這齣話劇裡.那兩位主角實在等得太無聊了.甲(Estragon)說:「我們上吊吧!」乙(Vladimir)說:「多等一會,看看他會說什麼?」甲說:「誰?」乙說:「果陀.」

我們會絕望,是因為我們曾經有希望.從不希望的人,無所謂絕望.但當絕望出現的時候,保持「希望」卻是唯一解藥.

問題是:當等待,等待,再等待,仍等不到「果陀」時,仍要等待嗎?

「果陀」真的會來嗎?

「等待」抑或「不再等待」,是生命裡的一個重大問題.

關於「等待」這個神學主題,主耶穌講了一個故事.這就是我們今天的福音書經課:馬太福音25:1－13.這也是我今天的講道經文.

馬太福音25:1－13:

25:1　「那時,天國好比十個童女拿著燈出去迎接新郎.

25:2　其中有五個是愚拙的,五個是聰明的.

25:3　愚拙的拿著燈,卻不預備油；

25:4　聰明的拿著燈,又預備油在器皿裡.

25:5　新郎遲延的時候,她們都打盹,睡著了.

25:6　半夜有人喊著說:『新郎來了,你們出來迎接他!』

25:7　那些童女就都起來收拾燈.

25:8　愚拙的對聰明的說:『請分點油給我們,因為我們的燈要滅了.』

25:9　聰明的回答說:『恐怕不夠你我用的；不如你們自己到賣油的那裡去買吧.』

25:10　她們去買的時候,新郎到了.那預備好了的,同他進去坐席,門就關了.

25:11　其餘的童女隨後也來了,說:『主啊,主啊,給我們開門!』

25:12　他卻回答說:『我實在告訴你們,我不認識你們.』

25:13　所以,你們要警醒；因為那日子,那時辰,你們不知道.」

主耶穌這個故事,處理「等待」這個主題,不是處理人與人相處的問題,故此,請不要去想:「為何那些聰明的童女不早點提醒那些愚拙的童女多預備油?」「那些聰明的童女為何這般自私,不願意分享?」「那個新郎為何這麼絕情,不開門給那些愚拙的童女?」我們放下這些問題.

這個故事設定一個問題:「在一個必需等待的人生裡,你是如何等待的?」一個明白天國是什麼的人,他會正確地等待.一個正確地等待的人,他活在天國裡.

在主耶穌的故事裡,「聰明的童女」與「愚拙的童女」都是在等待,等待那位新郎,等待一場生命的筵席,等待「果陀」.這場等待很漫長,不知等到何時.黑夜來臨了,她們「都打盹,睡著了」.這場等待可能很無聊,令人很疲累.她們的感覺是一樣的,反應是一樣的.人累了,便睡覺；餓了,便吃飯.無論聰明與愚拙,這都是一樣的.

「聰明」與「愚拙」的分別,在哪裡呢?

經文說:「愚拙的拿著燈,卻不預備油」(25: 3)

為什麼她們不預備油呢?

最簡單的解釋是:她們以為這只是一場短暫的等待.

將一件需要長久等待的事情看成只需要片刻等待,是「愚拙」的.

減肥,需要漫長的運動操練,但我們不能等待,只想快速減肥,吃藥也不夠快,便去做手術「抽脂」,有人因此喪命.這是「愚拙」的,也是危險的.

有些宗教教派,不能忍受天國太遲來臨.他們會用極端手法,即時實現他們的「天國」.像近日的「伊斯蘭國」,便用極端手法,建立他們心中的天國.任何他們看為是「反天國」的東西,都被消滅.這是「愚拙」的,也是危險的.

不能等待的人,容易採取暴力.從效果來說,暴力是最快達到目的的手段.若一個極權的國家對著人民說,若你們不信奉「ＸＸ主義」,便被處死,則這國家的人民很快便會表明信奉「ＸＸ主義」,當然,我們不會知道這是出於真心抑或假意.若一個人不能等待慢慢工作,慢慢儲錢,他用暴力去搶,則表面上他很快達到賺錢的目的.若一個父親沒有耐性去教導他的孩子,他一巴掌打過去,他的孩子便即時聽話了.不能等待的人,容易採取暴力.

在主耶穌的故事裡,那些愚拙人之所以愚拙,不是她們不等待,而是她們以為這等待只是片刻的等待.她們以為只需等待片刻,她們要得到的東西便會得到了.以為可以即時得到要得到的東西,不用等待,不願意等待,不值得去等待,這些想法是「愚拙」的,也是危險的.因為只有使用暴力,才能達到這種即時效果.

我們活在一個高速的時代,一切效果都是即時的.我們有即食麵,有即沖咖啡,有高速的上網手機,可即時接通世界.我們的耐性越來越低了.上網的速度慢了一點,我們已不能忍受了.若上網的速度只慢了一點點,我們都不能忍受,我們怎能忍受反應慢半拍的人?我們越來越喜歡反應快的人,越來越不喜歡停下來思考的人.動作片越來越受歡迎,要思考的東西越來越受冷落.與其講道理,不如「曬馬」.與其漫長地等待社會改變,不如即時攤牌.簡而言之,我們越來越不能等待,不能忍受等待.這是「愚拙」的,也是危險的.

主耶穌故事裡,那些聰明的童女,她們預備了足夠的油.她們明白,人生需要等待,事情的改變需要等待,她們更明白這等待可能需要漫長的時間.

我想講一個關於「等待」的故事.

在基督教貴格會(Quakers)的歷史裡,有一個很重要的人物,他的名字叫John Woolman(武爾曼),生平年代是1720-1772,他生活於新澤西的殖民地,當時美國仍未立國.他是裁縫,在貴格會的農民和商人中間生活.在那個年代,奴隸為主人帶來很多財富.但武爾曼得到神的啟示,認為奴隸制度是不合乎神的心意的.

長達20年,他努力向貴格會的會友,宣揚解放奴隸的信息.他去到偏遠的牧場,分享他得到的啟示.他寧願禁食,都不吃由奴隸預備或奉上的食物.如果他無意中得到奴隸的幫助,他堅持付錢給那奴隸.

很容易想到,不是人人都歡迎武爾曼的信息.但武爾曼從一個地方到另一個地方,從一個會議到另一個會議,宣講他得到的啟示.他的宣講,造成了會友間的張力,當時的現實是貴格會的會友很多是有奴隸的,但另一方面會友們的心中是有上帝的.現實與理想,利益和上帝,人對真理的順服和反抗,在會友間形成了莫大的張力.貴格會的人是不會採取投票形式去解決問題的,他們認為投票縱使能很快有答案,但這只會令落敗的人收聲,並不能帶來社群裡的真正和平.當然,他們是和平主義者,不會因為某些人有某些看法,暴力地將之逐出教會.長達20年,他們活在這張力裡,他們沒有取消這張力,也沒有以即時有效的方式解決這張力.20年後,他們慢慢地達到共識,就是要取消奴隸制度.20年後,現實與理想的張力,將人心綻開,使人投向正義和真理.這時刻,他們各人是心甘情願地取消奴隸制度的.

要達到這個共識,20年是否太長?

或許,是很長,但貴格會的群體是美國最早解放奴隸的宗教群體.他們解放奴隸的行動,比起美國那場與解放奴隸有關的內戰早了差不多80年.而80年後的那場戰爭,解決這問題的方式是急促而暴力的,製造了多少傷亡,帶來了多少怨恨和遺憾.

(註:上述武爾曼的故事,參閱帕爾默(Parker Palmer):《隱藏的整全》(香港:基道),頁199-201.)

武爾曼的故事,是容讓生命在脆弱的雞蛋裡有足夠時間慢慢孵化.這是需要等待的,這是需要漫長的等待的.太急促的動作,或過早的破開蛋殼,生命便完結了.

愚拙的童女,沒有想過要長期的等待,她們過早破開蛋殼,生命因得不到足夠的孕育時間而結束了.聰明的童女,她們明白等待,是需要時間的.

我們的社會此時此刻在極大的張力中,張力的雙方都在「等待」抑或「不再等待」的抉擇中.當一些需要長期等待的東西被看成是不用等待而當下可得時,可能是危險的,愚拙的.

很多人明白生命在雞蛋中孕育孵化是需要時間的,他們願意等待.但是等呀等呀,等了很久,生命仍未出現,「果陀」還未來,「理想的訴求」還未實現,他們在漫長的等待中灰心了,失落了.我們願意等,但真的等到什麼嗎?

主耶穌的故事,最後一個重點是:請不要灰心,新郎是會來的.

主耶穌說:「所以,你們要警醒；因為那日子,那時辰,你們不知道.」(25:13)

各位,請不要灰心,天國是會實現的.縱使長路漫漫,天國是會降臨的.

各位,在這個看似沒有上帝的世界,請保持對上帝的盼望.在盼望的漫長歲月裡,請堅持實踐上帝要求的正義和仁愛.天國的奧秘,就在這生活中.

我引用關於德蘭修女的一本書中的一段話作結束:

「人們不講道理,思想謬誤,自我中心,不管怎樣,總是愛他們；

如果你做善事,人們說你自私自利,別有用心,不管怎樣,總是要做善事；

誠實與坦率使你易受攻擊,不管怎樣,總是要誠實與坦率；

你耗費數年所建設的可能毀於一旦,不管怎樣,總是要建設；

人們確實需要幫助,然而如果你幫助他們,卻可能遭到攻擊,不管怎樣,總是要幫助.」(引用自Lucinda Vardey:《一條簡單的道路》(新店:立緒文化),頁166-167.)

但願:榮耀歸於聖父,聖子,聖靈.阿們.



\newpage

\section{數算自己的日子}
\label{sec:268}
\textbf{伍渭文牧師 - 數算自己的日子}
\newline
\newline
\hyperref[sec:267]{< < < PREV SERMON < < <}
~
\hyperref[sec:toc]{[返目錄]}
~
\hyperref[sec:269]{> > > NEXT SERMON > > >}


講題:數算自己的日子　Counting Our Days

經文:詩篇90

講員:伍渭文牧師

場合:香港中文大學崇基學院禮拜堂

日期:2014年11月16日

雨傘運動開始至今剛五十天了,反對佔領道路的聲音愈來愈烈.有人觀察很多學生和群眾運動最後的失敗,因為一開始時沒有設定撤離機制,對何時結束運動沒有共識,至運動高潮後形成膠著狀態,不知如何結束,也不想結束.《風俗通》有一個故事:「殺君馬者道旁兒」,說到一匹好馬跑得很快,但路邊看客不停地鼓掌,馬兒於是不停地加速,結果不知不覺地累死了.馬兒是死在路邊看客的掌聲之中,忘記自己的方向和終點,也忘記自己是血肉之驅. 我們在人生舞台奔跑,有想過退場的機制嗎?還是我們也被路旁兒掌聖所驅策,不停往前衝,不顧身體健康,不顧和家人朋友的關係,這些路旁兒可能來自大眾媒體,它們塑造世界成功人士的標準,鼓動我們對號入座,並為此不停奔跑,奔跑. 今天詩篇經課九十篇就提醒我們,人生短促,要有智慧做好退場機制.講道題目取自詩篇九十11, 12:「誰曉得你怒氣的權勢?誰按著你該受的敬畏曉得你的忿怒呢?求你指教我們怎樣數算自己的日子,好叫我們得著智慧的心.」 這兩節經文把整篇詩篇前後兩段連接起來.一至十節論及人生短促,但可植根永恆；十三至十七節論及在短促人生所作的工作,可被上帝堅立,行經時間,留下痕跡.我們先看第一段.

植根永恆-跨越朽壞 vv 1-10

詩篇九十是安息禮拜常用的經文,經文有律法的嚴厲,迫使人面對死亡的事實,說到人生短促,人有一生必有一死,死亡是眾人的命運,好好反思生命.但細看經文,也有福音的安慰:死亡不是人生的歸宿,上帝才是人生的歸宿,英文:We die, but we do not die out! 人死如燈滅是die out,就如旅客乘機場鐵路到機場,機場就是旅程的終站,生命完結就完結了.但旅客往機場還要登機,飛到另一站；機場不是終站 ,是進入另一個世界的門檻.   永恆的靠山, vv 1, 2   「主啊,你世世代代作我們的居所,諸山未曾生出,地與世界你未曾造成,從亙古到永遠,你是神!」(vv  1, 2)   永恆不朽是人對上帝的期望,如果有上帝,他應該是永恆不朽的.相反,人有生老病死.但這詩重點不是歌頒上帝的永恆,也不是嗟嘆人生何其短促,乃強調短促的生命,植根永恆的上帝:「主啊,你世世代代作我們的居所」.天主教思高譯本直接譯作「作了我們的靠山」:「吾主,從永遠到永遠,你作了我們的靠山,群山尚未形成,寰球尚未生出,從永遠直到永遠,你就已經是天主.」上主是我們永遠的靠山,沒有一樣東西攔阻這靠山,過去如此,將來也如此.   「從亙古到永遠,你是上帝」也指出我們短暫人生的出現,是上主冥冥中的計劃.上帝不是臨時接手,將來也不會有任何力量挑戰他的主權,沒有人可以更改上主為我們一生寫下的劇本.我們植根於永恆的上帝,也參與上帝永恆的計劃,在上帝永恆的國度有份.   神聖的墓穴:把永恆不朽的上帝和短暫朽壞的世界聯系在一起,在耶穌時期的泛希臘羅馬世界,是嶄新的信仰,產生前所未有的宗教現象-聖徒的崇敬,甚至敬拜(The cult of saints).「從亙古到永遠,你是上帝」是當時耳熟能詳的觀念,人們認為祇有神祇才是永恆不朽,凡人終歸塵土,人死後靈魂便脫離朽壞的皮囊,升到星河(Milky Way) .星河是古代人可目測到宇宙的至高處,那裡繁星顆顆,千萬年就停在穹蒼之頂,永恆不動絲毫,那裡平靜如鏡,不沾凡塵.   我們肉體在月亮之下流轉變幻,最後朽爛成一堆黃土；肉體的歸宿是地上,靈魂的歸宿在天上,各有所屬,不能逾越.肉體是靈魂的牢獄,死亡使靈魂得釋放,祇有羅馬皇帝死後連同肉體羽化升天(apotheosis) ,證明他們有神性,是神祇地上的代表,強化其神權為外衣的極權統治,但皇帝死後連同肉體升天也是一廂情願的想法.   因為沒有復活的盼望,保羅引用哥林多的人說:「我們就吃吃喝喝吧!因為明天死了」.(林前十五32)   若河漢是吾家,我們連同復活不朽的身體,也居住其間；那躺在地上聖徒的墳墓,就成為天堂人間的交匯處,這是前所未有的宗教現象,因為古代泛希臘羅馬文化貶抑朽壞的肉身,不相信身體復活；認為短暫朽壞之軀,不可承受永恆之福.   但早期教會對聖徒的墳墓極其崇敬,他們多是殉道者-見證人.不論是以血澆奠作紅色的見證,或是禁欲作白色的見證,他們的墓地成了朝聖者參觀祈禱的地方；他們的行狀,也成為信眾的榜樣,他們朽壞的身體,已連上不朽的上帝,分享上主永恆的生命.公元三世紀開始,很多教會裡面都安放了聖徒的墓穴,其上安放大型的燭台,靜靜燭光,射在神鑫的頂部,儼然像穹蒼的星河(Milky Way) .這是天堂,也是人間,聖徒的墓穴就是天堂人間的交匯.   天主教會一路來都保持這傳統,在教堂內安放聖徒的靈寢,強調其使徒的統緒和傳襲.現任教宗方濟在梵帝崗聖彼得大教堂就職前,就進到教堂地下室,在安放傳說為使徒彼得遺體的地方,先行祈禱,然後就職.   上主世世代代作我們的居所,上主是我們的靠山,我們跨越朽壞,植根永恆.早期教會非常重視聖餐,認為吃了基督永恆的身體,我們也分享永恆的生命；馬丁路德也說,吃了上帝羔羊的生命,本是狼的生命也漸漸變成羊的生命,我們食甚麼就成為甚麼,we are what we eat.   詩篇跟著解釋生命短暫的原因,說明為何從永恆之主而來的人,去得如此匆匆.     罪帶來死亡, vv 3, 5, 7, 8   「你使人歸於塵土……你叫他們如水沖去……我們因你的怒氣而消滅,因你的忿怒而驚惶是人惹動上主的怒氣而消滅,因他的忿怒而驚惶」(vv 3, 5, 7, 8)   是人惹動上主的怒氣而消滅,因為罪人失去永恆的生命:「你使人歸於塵土」(v 3), 「你叫他們如水沖去」(v 5)  歸回塵土是審判,大齋期開始的聖灰日的塗灰禮就說明這點.主禮把去年棕枝主日的棕枝,燒成灰燼,以十架記號塗在領灰者的額上或掌心說:「父上帝,你用地上的塵土創造我們,求你施恩,使這些灰燼向我們揭示生命的有限,並提醒我們必須悔罪,因為只有仰賴你的恩典,我們才能領受在我們的救主耶穌基督裡頭的永生.」眾答:阿們.在塗灰時,主禮向領灰者額上劃十字並說:「你要記住,你本是塵土,仍要歸於塵土.當遠離罪,效忠基督.」   罪使生命顯得虛空,  v 10   「我們一生的年日是七十歲,若是強壯可到八十歲；但其中所矜誇的不過是勞苦愁煩,轉眼成空,我們便如飛而去.」(v  10)   人如水沖去,人生如睡一覺,不單說到人生的短促,也說到人生的虛幻.人世間種種繁華世界,像黃梁一夢,鏡花水月,並不真實；沖去的水留不住,沒有永恆的價值.這種沒有上主的空虛感,催迫才華橫溢的詩人,政治評論家多瑪斯．牟敦梅頓(Thomas Merton) 追求修道院獨處生活.他的自傳《七重山》(The Seven Storey Mountain) 自一九四八年出版以來即風行全球,譯成二十多種文字.《紐約時報》將之譽為二十世紀的《懺悔錄》.   在書的結語,牟敦說:「我喜愛獨處唯一的原因-在被造的東西中消失,向它們死去,不知道它們,因為它們提醒我離你多麼遠.它們提醒我:你和這些相距甚遠,雖然你在其中.你創造它們,它們因你而存在,它們把你遮蓋起來,讓我看不見你.我將獨自生活,離開它們.福哉獨處!」(That is the only reason why I desire solitude-to be lost to all created things, to die to them and to the knowledge of them, for they remind me of my distance from You. They tell me something about You: that You are far from them, even though You are in them. You have made them and Your presence sustains their being, and they hide You from me. And I would live alone, and out of them. O beata solitudo! ”) (Thomas Merton, The Seven Storey Mountain. An Autobiography of Faith, 1998: 461) .   上帝是我們的靠山,我們超越朽壞,超越勞苦愁煩,也超越罪帶來的苦難.

植根永恆-跨越苦難 vv 13-15「耶和華啊,我們要等到幾時呢?求你轉回,為你的僕人後悔,求你使我們早早飽得你的慈愛,好叫我們一生一世歡呼喜樂.求你照著你使我們受苦的日子和我們遭難的年歲,叫我們喜樂.」(vv 13-15)第十三節:「求你回轉,為你的僕人後悔」不是上主後悔做錯了；思高譯本作:「上主,求你歸來,尚待何時?求你快來憐恤你的僕役!」英文新修訂本(New Revised Standard Version):Turn, O Lord! How long! Have compassion on your servants! 用了三個感嘆號,表達祈求上主憐憫之切.「我們經過的日子都在你震怒之下,我們渡盡的年歲好像一聲嘆息.我們一生的年日是七十歲,若是強壯可到八十歲」香港平均壽數女性八十六,男性八十,壽數之高在世界前列.我們常常用平均壽數,減去自在的年齡,推算自己還有多少日子,但其實明天如何我們實在不知道,我們想也想不到一位滿懷希望的內地二十歲尖子來香港作交短期換生,竟然因為成績的壓力,墮樓輕生?也有人正值壯年想也想不到突然身惟頑疾,這星期先後兩日收到消息,認識的兩位芳華正茂的中年信徒腫瘤復發和新發.

托爾斯泰(Leo Tolstoy)名著安娜.卡列寧(Anna Karenina) 開場白說得很對:「幸福的家庭都彼此相似,不幸的家庭各有各的不幸.」(Happy families are all alike; every unhappy family is unhappy in its own way. ) 人生不單短暫,而且充滿各種不同的不幸.信靠耶穌不是免去各種不同的不幸,乃在遭難的日子,仍然有喜樂,仍然飽得上主的慈愛.上主沒有在人生航道中除去礁石,這些礁石叫我們的船隻破裂沉沒,上主的恩典像海水上漲,把船升高,迎礁而上,漫過礁石. 我們會面對至親離開的苦痛,但上主看聖民之死,極其寶貴；我們會流淚,但上主用袋裝.上帝的兒子耶穌已來到世界,不單解決罪 (fix sin) ,也關顧被罪壓傷的人(care about sinners).上主在諸般患難中看顧我們,撫平我們的傷痛.

植根永恆-跨越短暫 vv 16, 17

「願你的作為向你僕人顯現；願你的榮耀向他們子孫顯明.願主我們神的榮美歸於我們身上.願你堅立我們手所作的工；我們手所作的工,願你堅立.」(vv 15, 16) 人生如此短暫,想做的事又如此的多,怎麼辦? 最近網上熱傳一視頻,一位六十歲名叫方勵先生電影人和一班三十歲以下的年青人分享人生,講演題目:惜命不是養生,是用來折騰.講者說他身邊很少同齡人,因為他們都養生去了,他們喜歡打哥爾夫球. 一開始就說他是很成功的人,因為他很快樂,因為他道知人是甚麼；接著他給大家看美國亞利桑亞州的巴林杰隕石坑(Barringer  Meteor Crater) 的圖片,直徑1.2公哩,深175米,約五萬年前由隕石撞擊地球而成,說明人的晚出,力量微小,生命短暫,要珍惜生命；我們要在最短的時間完成自己喜歡做的事,人生一定要好玩.在事業高峰期他辭掉了高薪厚職,自己創業,目的就想做自己喜歡做的工作.他想接觸不同的人,認為朋友很重要,和朋友相處給他很大的滿足感.他不斷開始不同的新工作,做金融投資,人才培訓,最後還拍電影,同時照顧四十多個項目,都很成功.做那麼多來折騰自己因為和時間競賽,睡眠時間不多,因為人生苦短,要活得多姿多彩. 這視頻非常勵志,激動人心.但講者做得到,不是每個人做得到.講者誠然天資優異,目標明確,努力不懈,加上上天眷顧,才有今天的成就.不過做了那麼多的工作,真的感到不枉此生嗎?當我們離開世界,這些工作會持續不衰嗎?我們的手作作的工,能否堅立不墮呢? 喬布斯(Steve Jobs)傳記作者 (Walter Isaacson) 與喬布斯作了超過四十次訪談,最後一次在他家中的後花園進行,訪問是在陽光普照的下午進行.訪問時喬布斯身體感到很不舒服,其間喬布斯把話題轉到宗教和死亡.華達.以撤遜問喬布斯是否相信「死後生命」(afterlife),是否相信有上帝時,他說:「五十五十.這是大奧秘,人永遠不知道,但我希望死不是生命的終結,我但願累積的智慧不會隨死亡歸於無有,可以永續.但他跟著說:但可能它真的像開關制,click一下,你就離開了；(停頓一秒後他笑說) 所以我們蘋果產品沒有開關制的.」 短暫人生所做的一切,詩人看來最重要是被永恆的上主所紀念,被他所肯定,而不是被人所紀念.「願你堅立我們手所作的工；我們手所作的工,願你堅立.」(v  17)我們手所作的,影嚮我們子孫的生命,而他們也影響他們子孫的生命,代代相傳,有持續的果效.「願你的作為向你僕人顯現,願你的榮耀向他們的子孫顯明.」(v  16) 詩人用僕人一詞,說明短暫的人生,也可參與上主永恆的計劃,成為他的器皿,被他使用,所做的工被他堅立.這些工作不在乎人看來是否豐功偉業,上主求於管家是忠心,今天福音經課太二十五21:「你這又良善又忠心的僕人,你在不多的事上有忠心,我要把許多的事派你管理.」不論我們的秉賦才幹多少,一千好,二千好,五千好,重要的是忠心做好管家的職任,用愛心做好小事情；動機純良,忠於所托,就是「又良善又忠心」的僕人,在他永恆國度裡有份. 幾年前我在一本雜誌讀到有關一位年長姊妹的見證,她結婚不久丈夫就因病去世,獨力養大一對兒女,用聖經真理建立他們的價值觀.到了大兒子生了孩兒,她又忙幫湊孫,也熱心協助教會探訪.後來婆婆患了重病,身體漸漸衰弱,很多人探望她因為愛她.婆婆很少叫人為自己身體祈禱,反而關心其它患病的人,她最後的遺言是甚麼呢?是叫孫兒寫下一些recipe烹飪配方,她關心離開之後兒媳孫兒能否吃到他們喜歡的菜譜.她認為食快餐是生存,但自己去買菜下廚弄菜是生活,吃得健康又省錢,在餐桌前上一家人又吃又談,她更用廚藝在家款待朋友. 遺言是recipe烹飪方,很有意思.她臨終前還掛心兒孫吃得好.她預備好見上主,為何那麼年輕丈夫就離開?為何自己到享晚年福時染上頑疾?為何是她?她沒有問上主,因她知道上主愛她,她祇關心其它人的需要.帶著滿足和感恩,她知這上主在那邊等她.上主給她的托付是獨力養大兩兒女,有正確的價值觀.她是又良善又忠心的僕人.她沒能力同時照顧四十多個項目,也不覺得生命短暫,惜命就要折騰自己,在最短的時做是多的事. 心靈的覺醒:今天是教會年最後的第二個主日,還有一個主日,就結束教會年了,再經過為為期四主日的將臨期,就是聖誕.教會年的結束,和人生命的結束一樣,叫我們想起上帝的審判.教會年最後數星期的經課,都涉及審判的主題.將臨期也緊扣審判的主題,明白審判的嚴厲和真實,誘發我們對救贖的熱切盼望,因基督的來臨使我們免去上帝的審判.   教會年快將結束時,一連三個主日的福音書經課馬太福音二十五章,都有審判相關的主題:上主日主題是「儆醒」,要學聰明的童女,預備好有油的燈,迎接新郎來臨；下主日主題是「審判」,分辨綿羊和山羊；今天主題是「交賬」,要作個忠心有見識的管家,善用上主交付我們的秉賦.   植根永恆的盼望:但審判的基調不是驚懼惶恐,乃是歡欣喜悅.童女興奮等候新郎的出現和緊隨的婚宴；綿羊等待分辨的程序,之後可以進天國；忠心有見識的管家勝卷在握,不會患得患失,等候被主人稱讚嘉許.「主啊,你世世代代作我們的居所.」永恆的上帝是我們永遠的靠山,使我們以喜樂面世間種種橫逆.詩篇九十篇重點不在嗟嘆人生苦短,嘲諷人生虛空,乃宣告植根永恆,堅立人生.   機場鐵路的起始是香港站,跟著是九龍站,青衣站,赤立角機場,但機場不是我們的終站,我們還是登機,沖上雲霄,翱翔萬頃窮碧,飛向那永久的城邑,我國永遠的家鄉,坐在駕駛艙的,是永恆的上帝.阿們.



\newpage

\section{公義和平安互相親嘴}
\label{sec:269}
\textbf{鄧瑞強博士 - 公義和平安互相親嘴}
\newline
\newline
\hyperref[sec:268]{< < < PREV SERMON < < <}
~
\hyperref[sec:toc]{[返目錄]}
~
\hyperref[sec:270]{> > > NEXT SERMON > > >}


講題:公義與平安互相親嘴　Righteousness and Peace Kiss Each Other

經文:詩篇85

講員:鄧瑞強博士

場合:香港中文大學崇基學院禮拜堂

日期:2014年12月7日

各位弟兄姊妹,早安.

今日有弟兄姊妹領受洗禮,故我的講道會比平常的講道短一點.

我今日的講道經文,是今日的詩篇經課:詩篇85篇.

由於時間關係,我不再讀經文一次了.

首先,讓我們簡單的看看這詩篇的結構及佈局.

第1-3節,是一種回憶,描述神對他的子民「過去」的恩情,用的是過去的時態,神的子民是以「第三身」的「他們」去表達的.

昔日是美好的,言下之意,今日問題多多.

第4-6節,是求神「回轉」(第4節),是渴望神的解救,神的子民是以「第一身」的「我們」來表達的.

第7節是總結上文的祈求,卻又特別點出神的子民的希望所在.神是「慈愛」的,他的「救恩」會賜給我們的.

第8節,詩人以「第一身」的「我」發言,對於人生的困局,他像先知一樣,講出其神學洞見.

第9-12節,就是詩人的神學洞見.第9節與第12節平行,第10節與第11節平行,這四節經文組成A-B-B’-A’的平行.

第13節,詩人回望了歷史,在信仰中對前景懷著盼望,如今,面對當下的困局,他安靜地對世人講出正道何在.

我們看第一段:第1-3節

詩85:1　〔可拉後裔的詩,交與伶長.〕耶和華啊,你已經向你的地施恩,救回被擄的雅各.(聖經＜新譯本＞將「救回被擄的雅各」譯為「恢復了雅各的產業」(NIV:“restored the fortunes of Jacob”),這翻譯可能好一點,因這詩篇未必與以色列人的「被擄歸回」事件有關.)

詩85:2　你赦免了你百姓的罪孽,遮蓋了他們一切的過犯.〔細拉〕

詩85:3　你收轉了所發的忿怒和你猛烈的怒氣.

「可拉後裔的詩,交與伶長」.不知大家對「可拉」有什麼印象?

「可拉」是利未支派的人,在摩西帶領神的子民出埃及時,可拉帶頭叛變.民數記16:31-32如此記述:「他們腳下的地就開了口,把他們和他們的家眷,並一切屬可拉的人丁,財物,都吞下去.」

可拉因背叛而被「大地」吞噬,現在,他的後裔卻蒙憐憫,成了聖殿的詩班員.他們唱的這首詩篇,充滿「地」的意象.不知他們何所思,何所感?我想,他們會用心唱這首詩,向參與崇拜的人,訴說神對不義者的懲罰,訴說神對自己的恩情,訴說行在正路上的重要性.

第1-3節短短三節經文,回顧了神對他的子民施予的厚恩.看看那些用字,第1節:「施恩」,「救回」,第2節:「赦免」,「遮蓋過犯」,第3節:「收轉了所發的忿怒」.在神的眷顧裡,一切如斯美好.

但是,好景不常,神收轉了的忿怒,在當下,又再爆發出來.「忿怒」,成了下一段的主題.

第二段經文:第4-6節

詩85:4　拯救我們的神啊,求你使我們回轉,叫你的惱恨向我們止息.

詩85:5　你要向我們發怒到永遠嗎?你要將你的怒氣延留到萬代嗎?

詩85:6　你不再將我們救活,使你的百姓靠你歡喜嗎?

這一段經文,與上一段經文,氣氛相差很遠.「惱恨」,「發怒」,「怒氣」充滿這段經文.過去的平安,今日不再復見.人在其中,感到神審判的列怒,渴求神回復昔日的恩典.

「回復」(return, restore)這個詞,貫串著這詩篇的頭六節經文.第1節:「恢復了雅各的產業」,即return雅各的產業給雅各.同一個字,出現在第3節:「你收轉了所發的忿怒」,即你的忿怒return.同一個字,又出現在第4節:「求你使我們回轉」,即make us return,或者說return to us.又是同一個字,出現在第6節:「你不再將我們救活…」,字面直譯即will you not return our lives.神的恩典,表明在神向我們返回,回收他的怒氣,回復我們的福氣.如今,神好像發怒了,人的呼求是:神啊,求你回轉,回復我們的生命與福氣.

這裡,道出了人類的「救贖和墮落」的周而復始.神的恩典回來了,人很開心,然後,人反叛,人墮落.然後,神的怒氣回來了,神的審判回來了.再然後,人悔改,人呼求,呼求神return,return our fortunes,return our lives.這種「救贖和墮落」的循環,能否終止呢?

這詩篇還有最後一個return,這個return講出終止「救贖和墮落」的循環的關鍵.第8節:「我要聽神─耶和華所說的話；因為他必應許將平安賜給他的百姓─他的聖民；他們卻不可再轉去妄行(NIV: “Let them not return to folly.”).」神的忿怒會不斷return的,直到人不再return to folly.

人常常求神回轉,自己卻常常回轉到歧途那裡去.這樣,不是神回不回轉的問題,而是你自己回不回轉的問題.直到你不再轉去惡行,你就不再需要求神回轉了.究竟,真正需要return的,是神,抑或是我們呢?直到當我們不再轉去邪惡,人世間的「救贖和墮落」的循環便終止了.

回頭看第7節:「耶和華啊,求你使我們得見你的慈愛,又將你的救恩賜給我們.」這是總結第4-6節的祈求.這裡引入兩個用字:「慈愛」和「救恩」.這裡提到「救恩」,是上接第4節:「拯救我們的神啊」,下開第9節:「他的救恩誠然與敬畏他的人相近」.而人能經驗神的救恩,全因神的「慈愛」.

神的慈愛,是人經歷救恩的基礎.但單單有這基礎也未足夠,人要有反應,要走向慈愛才行.經文第8節的重心就是在此.人要轉離妄行,轉向正路,才能經歷救恩.

這就是第9-12節這一段經文的主題.

詩85:9　他的救恩誠然與敬畏他的人相近,叫榮耀住在我們的地上.

詩85:10　慈愛和誠實彼此相遇；公義和平安彼此相親.

詩85:11　誠實從地而生；公義從天而現.

詩85:12　耶和華必將好處賜給我們；我們的地也要多出土產.

神會向我們回轉嗎?神的救恩會臨到我們嗎?

第9節指出:神的救恩誠然與人相近,但是,是與那些「敬畏他的人相近」.你要神向你轉回,你也要向神轉回才是.當人犯罪時,感到神離開了.好像神向左走,你向右走.你求神轉回,這時,神就算轉回,也只是在你後面追著你.直到,你也轉回,面向他,這時,兩者才能正面相遇.就像在沙灘上,一男一女向著對方跑過去,相遇,相擁,親吻.這就是第10-11節所講的.

「慈愛和誠實彼此相遇；公義和平安彼此相親.」慈愛,是神的本性.「誠實」(truth),是轉離「妄行」,敬畏神的人應有的品質.當神轉向人,當人轉向神的時候,兩者便相遇了.「公義」,是神的行事方式.「和平」(Shalom),是人生命的完滿,是人福樂的極致.當神要求的公義在地上被實踐的時候,人生的完滿便在神的慈愛裡展露.「公義與平安互相親嘴」,they kiss each other.

「誠實從地而生；公義從天而現.」這對應第10節.過一「誠實」的生活,是敬畏神的人應有的責任.當人盡其性之時,公義的天道便展現人間.慈愛的天與精誠的人接通,互相親吻,天道在人間.這就是第9節所講的:「神的榮耀住在我們的地上」,神在人間,以敬畏他的人為家.然後,大地才能回復秩序,人才能享受「和平」(Shalom),一種豐盛的人生.這就是第12節講的:「耶和華必將好處賜給我們；我們的地也要多出土產.」

第1節講到「地」,在結尾之時,第12節又回到「地」這個主題.「地」是人類的家園,同時,也是神施恩的場所.第1節講到神向「地」施恩.然後,好像一切混亂了.這時,若然人轉離妄行,轉離虛假,敬畏真神,言行誠實,則「地」便再蒙恩,能「多出土產」.

整首詩,似乎暗指當時大地乾旱,第3節講的神「猛烈的怒氣」暗指炙熱的太陽.神的子民向神求雨.雨下來,像「公義從天而現」.大地有了雨,便長出植物,「誠實從地而生」.詩人從大自然的災難,聯想到人間的荒涼.與其問:「有沒有足夠的雨水去生長植物」,不如問:「在神充沛的慈愛裡,人能否活出誠實的果子來」?人若結不出誠實的果子,大地會依舊荒涼下去.人若結出誠實的果子,則神的公義會彰顯人間,人間便成為讓人享平安的樂土.

這首詩的結尾,第13節:「公義要行在他面前,叫他的腳蹤成為可走的路.」在這裡,「公義」像開路使者,它走在神面前,開出神的大道,而這大道留在人間,也成為我們當走的大道.這結尾很簡短,卻呼喚我們走在這公義的正路上.對於今日要受洗的弟兄姊妹來說,這首詩對你的提醒是:轉離你過去的妄行,走上神開出的公義大道上.

社會會出現混亂,你的人生也會出現混亂.不要只求神回轉,我們也要回轉.神的義路在此,讓我們走在其中.願神的公義與人間的平安,互相親嘴.

但願:榮耀歸於聖父,聖子,聖靈.阿們.



\newpage

\section{天使的宣告}
\label{sec:270}
\textbf{伍渭文牧師 - 天使的宣告}
\newline
\newline
\hyperref[sec:269]{< < < PREV SERMON < < <}
~
\hyperref[sec:toc]{[返目錄]}
~
\hyperref[sec:271]{> > > NEXT SERMON > > >}


路加福音 1:26-38
\newline
\begin{longtable}{cl}
\hline
\hline
章節 & 經文 (和合本修訂版)\\
\hline
1:26 & \begin{tabularx}{0.7\textwidth}{X} 到了第六個月，天使加百列奉神的差遣往加利利的一座城去，這城名叫拿撒勒， \end{tabularx} \\ \\ \relax
1:27 & \begin{tabularx}{0.7\textwidth}{X} 到一個童女那裡，她已經許配大衛家的一個人，名叫約瑟；童女的名字叫馬利亞。 \end{tabularx} \\ \\ \relax
1:28 & \begin{tabularx}{0.7\textwidth}{X} 天使進去，對她說：「蒙大恩的女子，你好，主和你同在！」 \end{tabularx} \\ \\ \relax
1:29 & \begin{tabularx}{0.7\textwidth}{X} 馬利亞因這話就很驚慌，又反覆思考這樣問候是甚麼意思。 \end{tabularx} \\ \\ \relax
1:30 & \begin{tabularx}{0.7\textwidth}{X} 天使對她說：「馬利亞，不要怕，你在神面前已經蒙恩了。 \end{tabularx} \\ \\ \relax
1:31 & \begin{tabularx}{0.7\textwidth}{X} 你要懷孕生子，要給他起名叫耶穌。 \end{tabularx} \\ \\ \relax
1:32 & \begin{tabularx}{0.7\textwidth}{X} 他將要為大，稱為至高者的兒子；主神要把他祖先大衛的王位給他。 \end{tabularx} \\ \\ \relax
1:33 & \begin{tabularx}{0.7\textwidth}{X} 他要作雅各家的王，直到永遠；他的國沒有窮盡。」 \end{tabularx} \\ \\ \relax
1:34 & \begin{tabularx}{0.7\textwidth}{X} 馬利亞對天使說：「我沒有出嫁，怎麼會有這事呢？」 \end{tabularx} \\ \\ \relax
1:35 & \begin{tabularx}{0.7\textwidth}{X} 天使回答她說：「聖靈要臨到你身上；至高者的能力要庇蔭你，因此，那要出生的聖者要稱為神的兒子。 \end{tabularx} \\ \\ \relax
1:36 & \begin{tabularx}{0.7\textwidth}{X} 況且，你的親戚伊利莎白，就是那素來稱為不生育的，在年老的時候也懷了男胎，現在懷孕六個月了。 \end{tabularx} \\ \\ \relax
1:37 & \begin{tabularx}{0.7\textwidth}{X} 因為，出於神的話，沒有一句不帶能力的。」 \end{tabularx} \\ \\ \relax
1:38 & \begin{tabularx}{0.7\textwidth}{X} 馬利亞說：「我是主的使女，願意照你的話實現在我身上。」於是天使離開她去了。 \end{tabularx} \\ \\
[1ex]
\hline
\hline
\end{longtable}
講題:天使的宣佈　The Angelic Announcement

經文:路加福音1:26-38;46-48

講員:伍渭文牧師

場合:香港中文大學崇基學院禮拜堂

日期:2014年12月21日

引言:聖誕節與報佳音是分不開的-聖誕報佳音.這天晚上,那應許明亮的星,已劃破伯利恆的夜空,並傳來天使歌聲:「在至高之處榮耀歸與神!在地上平安歸與所喜悅的人!」(路二14)但聖經第一次報佳音,並不是對野地的牧羊人,乃是對童女馬利亞.對一個母親來說,真正的喜訊是她發現有身孕那一刻:「恭喜!恭喜!你有喜了.」嬰孩誕生下來祇是時間的問題,花落蒂熟,沒有懸念.以前還有是男是女的懸念,現在很早就知道了.在甚麼時間,甚麼地方誕生也沒有懸念,時辰八字都可以藉剖腹生子決定.

聖誕真正的喜訊,是天使對馬利說:「蒙大恩的女子,我向你問安……馬利亞,不要怕!你在神面前已經蒙恩了.你要懷孕生子,可以給他起名叫耶穌……聖靈要臨到你身上,至高者的能力要蔭庇你,因此所要生的聖者必稱為神的兒子.」(路一30-35)

一  道成肉身是最大的神蹟-密縫的世界被捅開了

已故著名聖經學者 何思虔爵士( Sir Edwyn Hoskyns ) 說:「道成肉身就像一匕首,把密縫的世界歷史捅開了」(The incarnation is like a dagger thrust into the weft of human history) .永恆不朽的上帝進入人類歷史,從此世界不再一樣了,耶穌的降生成為歷史的分水嶺:耶穌出生前稱公元前(B.C., Before Christ),出生後稱公元後 (A.D. Anno Domino, 主後) .

耶穌的復活,打開了封閉的世界,就像密封墳墓的墓石被推開,復活就像嶄新的靈風,吹進密封的世界,更新萬有.其實早在馬利亞聖靈感孕時,密封的世界已被捅開了.魯益斯(C.S.Lewis) 說得好:最大的神蹟,不是復活,是道成肉身；如果耶穌道成肉身,復活絕對不難,使水變酒也頂簡單.

早期教會面對是大的異端是諾斯底主義或靈智派(Gnostics),他們否認耶穌是道成肉身,使徒約翰說:「一切的靈,你們不可都信,…..凡靈認耶穌是成了肉身來的,就是出於神的.」(約壹四1, 2)首位被教會判為異端的馬吉安(Marcion) 在其福音書(Marcion Gospel)首句說:「該撒提比流十五年,彼拉多的時候,耶穌從天而降到加利利的迦伯農,在會堂施教」(In the fifteen year of Tiberious Caesar, in the time of Pontius Pilate, Jesus descended from heaven at Capharnaum, a town of Galilee, and taught there in the synagogue).

基督是靈體,附在一個叫耶穌的人的肉身,當耶穌這肉身被釘時,基督的靈體已離開.上主不會成為肉身,也不會被釘死在十字架上.因此早期教會崇拜頌讚的一個主題,就是道成肉身:「大哉,敬虔的奧秘!無人不以為然,就是神在肉身顯現,被聖靈稱義,被天使看見,被傳於外邦,被世人信服,被接在榮耀裡.」(提前三16)

不錯,夏蟲怎能與語冰呢?夏天的蟲,還沒有見冰就凍死了；水中的游魚,怎能和空中的飛鳥談話呢?魚一離開水面就缺氧死亡.生活在同一世界不同空間的生物尚且不能溝通,何況至聖,至善,永恆不朽的上帝,如何同朽壞滿有罪污的人溝通呢?

大哉,敬虔的奧秘!無人不以為然,就是神在肉身顯現.

最近上映的電影「星際啟示錄」(Interstellar) 對不同維度的人的溝通,甚有啟發.

時間不同

電影說到地球遭遇到不明的沙塵風暴,農作物又遭受蟲害,地球不適合人居住.地球本身沒有辦法自救,美國太空總署兩手準備,派人到外太空尋找救治地球的方法和合適人居住的星球,把人類遷移.身為父親的男主角是太空穿梭機駕駛員,離開時跟十歲的女兒相約一定回來,但回來時父女將同齡,因為所往的星球地心吸力不同,在那裡一小時已是地球上七年了.因為意外,父親在探索的星球逗留時多了時間,回到太空站已過了二十多年了；而且通訊設施壞了,祇能收到地球傳來的消息,不能發出訊息.

在熒幕上父親把二十多年地面傳來的訊息細看,見到兒子大學畢業,結婚,生子,樣子隨年齡變化.十歲的女孩也長大了成為中年出色的科學家,但怪自己爽約,在同年齡的生日當天,沒有出現,很氣很失望很難過,因為離開二十多年都收不到一個訊息,以為父親放棄了他們.但在另一端,父親因不能發信息報平安感到沮喪；兩人處於不引力的空間,父女情深,奈何河漢阻隔,不能溝通.這時候全場靜默,多人眼眶潤濕.

維度不同

為父的終在一星球找到救治地球的方法,但坐穿梭機回地球恐為時已晚了.地球的沙塵暴還在肆虐,疫症漫延,家人,全球人類危在旦夕.他想出辦法,冒險進入宇宙的蟲洞(wormhole) 和黑洞(dark hole) .蟲洞可縮短空間的距離:一張紙的兩端是是長距離,但對摺就在隔離,進入蟲洞可縮短返回地球的距離；跟著進入黑洞,更可以回到過去,因為黑洞吸力很大,甚至超過光速,追回很多時間,但入蟲洞和黑洞溫度很高,隨時喪命.不過到了黑洞另一邊,可以把拯救地球的方法,傳給自己的科學家的女兒,公佈世界,拯救地球.

穿過黑洞後,父親進入另一第五空間新維度,回到過去,看到過去十歲時的女兒和年輕的自己,但不能和女兒溝通,因為處於不同的維度,祇有一個溝通的方法-透過重力,把書房的書推跌,做成密碼,十歲的女兒知道外界的「鬼」(ghost)想和自己溝通,密碼寫成文字是「STAY」(逗留).十歲女兒叫父親留下來,不要往外太空,但父親硬著心離開,因他要找辦法拯救地球和他的家人.

長大成為科學家的女兒思念父親,回到幼時的書房,翻看密碼記錄,又看父親臨走時給自己的腕表的秒針不定跳動,知道是摩士密碼,是父親從另一個維度傳給自己的信息,快快讀取,找到解決沙塵暴和疫症的方法.

電影充滿扣人心絃的張力:地球快將沉淪,人類快將毀滅,父女分隔,父親如何把救贖的方法,從一個維度傳遞到另一個維度的女兒呢?不同維度的人尚且這麼難溝通,何況上主和人呢?

地球不能自救,來自地球之外的救法已藉密碼傳給我們,這聖經就是拯救的密碼,祇要用心解讀就會明白.

「道成肉身就像一匕首,把密縫的世界歷史捅開了」.

二  舉自卑微,顧念貧窮 (讓我們先看第一段經文)vv. 26-29

拯救世人就要道成肉身,但選擇何人成孕呢?耶穌由聖靈感孕,為童女馬利亞所生,顯出上帝的降卑為人,其實耶穌可以寄胎在大祭司女兒的腹中.路加運用天使加百列兩次不同的宣告作對比,突顯童女馬利亞的卑微.

兩個天使的宣告

第一次天使宣告的對象是祭司撒迦利亞,地點是耶路撒冷的聖殿,時間是撒迦利亞在聖殿供職獻祭的時候(註一)；魏峨的聖殿,尊貴的祭司,神聖的時刻,天使宣告他不育的太太將懷孕生子叫約翰:「他要使許多以色列人回轉,歸於主他們的神.他必有以利亞的心志能力,行在主的前面.」(路一16, 17)

第二次天使的宣告在寂寂無名的拿撒勒城一個清貧的家庭,我說清貧因為耶穌出生八天到聖殿行割禮,約瑟獻上的是斑鳩和雛鴿,沒有錢買羊羔(路二24) ；對象是一約十三歲,沒有甚麼履歷的鄉村女孩.當時十二,三歲就配婚,再過一,兩年就過門了.這宣告確是臨到Nobody in nowhere for no reason.

天使宣告說:「蒙大恩的女子,我問你安,主和你同在了!……你要懷孕生子,可以給他起名叫耶穌.他要為大,稱為至高者的兒子,主神要把他祖大衛的位給他.他要作雅各家的王,直到永遠,他的國沒有窮盡.」(路一28, 31-33)馬利亞聽後非常驚慌,反覆思想這樣問安有甚麼意思呢?誠然,馬利亞丈夫約瑟是大衛的後人,但大衛的後人很多,不一定要選中約瑟.這是上帝的主權,他定意揀選卑微的,叫尊貴的羞愧,揀選人看來是愚拙的,叫自以為聰明的蒙羞.

應用:我們從馬利亞被揀選,知道上主特別眷顧卑微的人.聖誕節的意義不在耀眼閃爍的燈飾,奢華的消費,乃在顧念被邊緣化生活窮困的人.

但馬利亞確有一些屬靈氣質令我們欣賞,作為我們的榜樣.

三  跟進感動,反覆思想,尋找互勉 (讓我們看看第二段經文) vv 39-45

首先,馬利亞對這事非常認真,她跟進天使的話:「那時候馬利亞起身,急忙往山地裡去,來到猶大的一座城.進了撒迦利亞的家,問以利莎伯的安.」(路39)為甚麼她這麼心急找以利莎伯呢?她很想分享這經歷,也想聽以利莎伯分享她的經歷,了解「那素來稱為不能生育的」如何成孕.

當馬利亞問安時,以利莎伯就被聖靈感動,高聲說:「你在婦女中是有福的!你所懷的胎也是有福的!我主的母到我這裡來,這是從哪裡得的呢?因為你問安的聲音一入我耳,我腹裡的胎就歡喜跳動!因為主對她所說的話都要應驗.」(路一42-45)

馬利亞探望以利莎伯,使後者得著激勵,因為一聽到問安,胎兒就跳動,受聖靈感動說出讚美的話.而馬利亞亦因以利莎伯的話,信心更加堅定,因為此刻天使的話已成就-她已經成孕了.以利莎伯六個月的胎兒,向自己剛成孕的胎兒歡喜跳動.巴哈(J.S. Bach) 更為了胎兒的跳動歡欣,寫了千古名曲-作品編號147-耶穌我靈所樂Jesu, Joy of Man’s Desiring.

馬丁路德說信徒的互勉(mutual consolation)就是教會的記號.我們崇拜時豈不是誦讀使徒信經宣告:我信聖靈,聖而公的教會,聖徒相通.教會也以此探訪為「聖母訪親日」慶節(Feast of Visitation) .因為這聖徒互勉,馬利亞唱出尊主頌(Magnificent):「我心尊主為大；我靈以神我的救主為樂.因為他顧念他使女的卑微,從今以後,萬代稱我有福…….」

反覆思想聖經

請注意路加第一章描述以利莎伯的讚美和稍後她丈夫撒迦利亞的頌讚,都提及是聖靈感動之下的讚美,惟獨馬利亞的尊主頌沒有提到是聖靈感動,為甚麼?因為尊主頌的概念是馬利亞很熟悉的,這著名的尊主頌,取材自撒母耳母親哈拿(Hanah) 求得兒子,蒙上主答允有孕的頌歌.(撒上二1-11)頌歌首句:「我的心因耶和華快活,我的角因耶和華高舉,我的口向仇敵張開,我因耶和華的救恩歡欣.」

可見馬利亞熟讀聖經,她也常反覆思想聖經.當聖嬰誕生下來,牧羊人尋打聖嬰,並把天使的話告訴約瑟和馬利亞.「凡聽見的,就詫異牧羊之人對他們所說的話.馬利亞卻把這一切的事存在心裡,反覆思想.」(路二18, 19)其它人詫異所聽到的,是對靈異的事情有好奇心,但馬利亞卻把這一切存在心裡,反覆思想.

四  尊主頌  (讓我們看看第三段經文) vv 46-56

尊主出少年

尊主頌是教會聖樂重要篇章,歷代不少著名作曲家為歌詞作曲,但這頌歌居然出自十三,四歲的少女,令人震驚,可能嗎?也許因為馬拉拉的出現,我們相信十三,四歲的少女,真的可以改變世界:”One child, one teacher, one book, one pen can change the world( -個小孩,一個老師,一本書,一枝筆,可以改變世界) “.2014年諾貝爾和平獎十七歲共同得主馬拉拉(Malala Yousafzai) 十二歲就應邀在BBC寫網誌,揭露在塔利班治下的巴基斯坦女生沒有機會受教育,兩年前被塔利班槍擊,險些喪命.

家庭的重要:顯然,馬拉拉有良好的家庭教育,她在家看到的父母,不是老是盯實電視機或手機熒幕,他們喜歡閱讀,也影響馬拉拉喜歡閱讀.

馬利亞的父母應該是虔敬的人,影響馬利亞,正如哈拿影響撒母耳,以利莎伯和撒迦利亞影響施洗約翰.聖誕節提到一個家庭-約瑟,馬利亞,耶穌；家庭是上主重要的畏器皿,成就他的救贖計劃.

清心的重要:尊主出少年.理論上,年齡愈大讀書愈多,經歷愈多應該認識上帝更多,對的-假如我們保持內心清明.但隨著年歲增加,我們罪過也增加,不夠清心；才智和財富增加,計算和牽挂也增加,不夠信心.但聖經說清心的人有福了,因為他們必得見上主.耶穌說:進天國的要像小孩子,因為小孩自己甚麼都尚未擁有,就單信賴供應他們的父母,他們有信心.

童女馬利亞蒙大恩,因為她清心；也有信心,承受恩典的容器夠大.保羅鼓勵年輕的提摩太:「不要叫人小看你年輕,總要在言語,行為,愛心,信心,清潔上,都作信徒的榜樣.」(提前四12)又說「人若自潔,脫離卑賤的事,就必作貴重的器皿,成為聖潔,合乎主用,預備行各樣的善事.」(提後二21)

尊主到老年

童女懷孕是一生的責任.馬利亞因為耶穌的原故,與約瑟和幼孩耶穌避難埃及,逃躲希律王的刀劍,因為希律維穩,把新生王滅於萌芽之際,凡兩歲以下的孩子都要殺掉.到耶穌十二歲到聖殿守節,回程時才發現不見耶穌,趕回聖殿尋找耶穌.「我兒,為甚麼向我們這樣行呢?看哪,你父親和我傷心來找你!」(路二48)這是育有少年子女(teenage children )的父母心聲,老是擔心他們的行縱.馬利亞最後親眼見到耶穌被釘十字架,痛不欲生.聖靈感孕是一生的承諾,為了責任不拒絕活下去的承擔.

我很多時候在婚禮提醒新人,不要忘記他們曾經是單細胞.其實每個人都是從單細胞走過來,是上主藉父母之間的愛情把我們帶到世上,他們如何在我們不能保護自己時保護我們,養育我們,直到今天,結婚是父母一個偉大工程的完成,但婚後他們還會想念我們,掛念我們,所謂一生兒女債.

基督徒懷有基督:馬利亞聖靈感孕,懷有基督；我們也被聖靈感動,領受了聖靈,懷有基督的生命.信仰是一生的跟隨,無論如何困難,為了信仰,不拒絕活下去.

教宗約翰XXIII.他七十六高齡當選教宗,被看為過渡教宗,看守教宗.有一次他說:「你看在我桌面文件一大堆,要做的事很多很多,我不懼怕死亡,但不拒絕活下去,祇求完成上主交付給我的事情.」任內約翰XXIII推行羅馬天主教的革新,召開梵蒂崗第二大公會議(1962-65),推動讀經運動,禮儀更新,肯定平信徒的職份,促進天主教會現代化,影響深遠.

不拒絕活下去有時比死亡更困難.當馬利亞知道這嬰孩名叫耶穌-救主,她知道她不單此刻背負未婚懷孕的羞辱,也會慢慢明白,將來她也得面對兒子因擔當世人罪孽被釘十架的苦痛.

能夠不拒絕活下去需要很大的恩典.不拒絕活下去是願意承擔責任；在工作環境裡,家庭姻親中,同事朋友間,有時會面對使人窒息的關系,看這些關系的已死去容易:當這些關係死了就可以擺脫這關係,但不拒絕這關系,盼望關系的新生,就需要更大的恩典和承擔.

結論和思考:上星期當我預備這講章時,剛發生令人難過的新聞,一位十五歲懷疑援交私影的少女被殺,棄屍垃圾房,那麼年輕就離開了.當我們思想童女馬利亞時,也要記念社會中不少身處兇險的少女:有些是自己迷失,有些是生活所迫,有些是遭人威迫；和等候他們歸家心裡焦灼的父母.她們和馬利亞差不多的年齡,也有上主的形像,也是基督尋找的對象,為上帝所愛.

懷有基督生命的信徒,不拒絕為福音原故活下去,不拒絕面對這些社會黝黑面,因為耶穌正為此而生,道成肉身.

註一:這是祭司一生工作的高潮.耶路撒冷當時共有祭司一萬八千人,分二十四班.撒迦利亞屬亞比雅第八班次,輪到他們班次的要抽簽,抽到的可有兩次一星期在聖殿供職,早上九時及下午三時獻一歲的羊羔為祭.因為人數眾多,抽中的才有機會獻祭,一生祇一次獻祭.



\newpage

\section{報到}
\label{sec:271}
\textbf{林豪恩先生 - 報到}
\newline
\newline
\hyperref[sec:270]{< < < PREV SERMON < < <}
~
\hyperref[sec:toc]{[返目錄]}
~
\hyperref[sec:272]{> > > NEXT SERMON > > >}


路加福音 2:22-40
\newline
\begin{longtable}{cl}
\hline
\hline
章節 & 經文 (和合本修訂版)\\
\hline
2:22 & \begin{tabularx}{0.7\textwidth}{X} 按摩西律法滿了潔淨的日子，他們就帶著孩子上耶路撒冷去，要把他獻給主。 \end{tabularx} \\ \\ \relax
2:23 & \begin{tabularx}{0.7\textwidth}{X} 正如主的律法上所記：「凡頭生的男子必歸主為聖」； \end{tabularx} \\ \\ \relax
2:24 & \begin{tabularx}{0.7\textwidth}{X} 又要照主的律法上所說，用一對斑鳩，或用兩隻雛鴿獻祭。 \end{tabularx} \\ \\ \relax
2:25 & \begin{tabularx}{0.7\textwidth}{X} 那時，在耶路撒冷有一個人，名叫西面；這人又公義又虔誠，素常盼望以色列的安慰者來到，又有聖靈在他身上。 \end{tabularx} \\ \\ \relax
2:26 & \begin{tabularx}{0.7\textwidth}{X} 他得了聖靈的啟示，知道自己未死以前必看見主所立的基督。 \end{tabularx} \\ \\ \relax
2:27 & \begin{tabularx}{0.7\textwidth}{X} 他受了聖靈的感動，進入聖殿，正遇見耶穌的父母抱著孩子進來，要照律法的規矩而行。 \end{tabularx} \\ \\ \relax
2:28 & \begin{tabularx}{0.7\textwidth}{X} 西面就把他抱過來，稱頌神說： \end{tabularx} \\ \\ \relax
2:29 & \begin{tabularx}{0.7\textwidth}{X} 「主啊，如今可以照你的話，容你的僕人安然去世； \end{tabularx} \\ \\ \relax
2:30 & \begin{tabularx}{0.7\textwidth}{X} 因為我的眼睛已經看見你的救恩， \end{tabularx} \\ \\ \relax
2:31 & \begin{tabularx}{0.7\textwidth}{X} 就是你在萬民面前所預備的： \end{tabularx} \\ \\ \relax
2:32 & \begin{tabularx}{0.7\textwidth}{X} 是啟示外邦人的光，是你民以色列的榮耀。」 \end{tabularx} \\ \\ \relax
2:33 & \begin{tabularx}{0.7\textwidth}{X} 孩子的父母因論耶穌的這些話就驚訝。 \end{tabularx} \\ \\ \relax
2:34 & \begin{tabularx}{0.7\textwidth}{X} 西面給他們祝福，又對孩子的母親馬利亞說：「這孩子被立，是要叫以色列中許多人跌倒，許多人興起；又要成為毀謗的對象， \end{tabularx} \\ \\ \relax
2:35 & \begin{tabularx}{0.7\textwidth}{X} 叫許多人心裡的意念顯露出來；你自己的心也要被劍刺透。」 \end{tabularx} \\ \\ \relax
2:36 & \begin{tabularx}{0.7\textwidth}{X} 又有位女先知，名叫亞拿，是亞設支派法內力的女兒，年紀已經老邁，從童女出嫁，同丈夫住了七年， \end{tabularx} \\ \\ \relax
2:37 & \begin{tabularx}{0.7\textwidth}{X} 就寡居了，現在已經八十四歲。她不離開聖殿，禁食祈求，晝夜事奉神。 \end{tabularx} \\ \\ \relax
2:38 & \begin{tabularx}{0.7\textwidth}{X} 正當那時，她進前來感謝神，對一切盼望耶路撒冷得救贖的人講論這孩子的事。 \end{tabularx} \\ \\ \relax
2:39 & \begin{tabularx}{0.7\textwidth}{X} 約瑟和馬利亞照主的律法辦完了一切的事，就回加利利，到自己的城拿撒勒去了。 \end{tabularx} \\ \\ \relax
2:40 & \begin{tabularx}{0.7\textwidth}{X} 孩子漸漸長大，強健起來，充滿智慧，又有神的恩典在他身上。 \end{tabularx} \\ \\
[1ex]
\hline
\hline
\end{longtable}
講題:報到　Register

經文:路加福音2章22-40節

講員:林豪恩先生

場合:香港中文大學崇基學院禮拜堂

日期:2014年12月28日

引言

「報到」這個口語在日常生活中用途十分廣泛,相信人人都有「報到」的經驗,甚至一天之內要到不同的地方向不同的人「報到」.還記你第一次「報到」是向誰「報到」嗎?你一定不會記得,因為這應該是別人替你做的.根據香港法例,新生嬰兒必須在出生後42天內,到出生登記處「報到」.在香港,這首次的「報到」,嬰兒是不需要出現的.

有些「報到」人們是樂意做的,有些「報到」是勉為其難的；有些「報到」的地方是人們不願意去的,有些「報到」的地方卻是非常受歡迎的.香港,就是一個受歡迎的嬰兒「報到」地方,如果沒有限制,鄰近地區的媽媽都想在香港分娩,目的是要讓嬰兒取得香港出生證明書.原來,在哪裡「報到」,向誰「報到」和取得甚麼身分有關,而這個身分又可能和一個人的命運有關.

如果沒有人為孩子「報到」,會發生甚麼事情呢?某天半夜的台北街上,警察截查一位少女,發現她沒有身分證,翻查警政以及戶政的電腦系統,竟然也找不到她的身分資料.這位少女會否是非法入境者呢?要不要遞解出境呢?遞解往哪裡去呢?後來,少女的父親出現,疑問才水落石出.原來這女孩出生後第三天,還來不及辦出生登記,母親就不辭而別,音訊全無.因為沒有向有關當局「報到」,這位女孩從小就沒有醫療保健,不能上學,也沒有辦法找工作,只能在社會上漂泊,還染上毒癮.那次被警察發現後,少女的父親希望當局幫幫忙,還她的女兒一個身分.(tvbs.com)

今天的福音經課,也可以看作一個嬰兒「報到」的故事.在故事中,誰人往甚麼地方去「報到」呢?這樣的「報到」有甚麼意義呢?對我們又有甚麼啟發呢?

經文

2:22按摩西律法滿了潔淨的日子,他們帶著孩子上耶路撒冷去,要把他獻與主,

2:23(正如主的律法上所記:「凡頭生的男子必稱聖歸主.」)

2:24又要照主的律法上所說,或用一對斑鳩,或用兩隻雛鴿獻祭.

2:25在耶路撒冷有一個人,名叫西面；這人又公義又虔誠,素常盼望以色列的安慰者來到,又有聖靈在他身上.

2:26他得了聖靈的啟示,知道自己未死以前,必看見主所立的基督.

2:27他受了聖靈的感動,進入聖殿,正遇見耶穌的父母抱著孩子進來,要照律法的規矩辦理.

2:28西面就用手接過他來,稱頌上帝說:

2:29主啊!如今可以照你的話,釋放僕人安然去世；

2:30因為我的眼睛已經看見你的救恩－

2:31就是你在萬民面前所預備的:

2:32是照亮外邦人的光,又是你民以色列的榮耀.

2:33孩子的父母因這論耶穌的話就希奇.

2:34-35西面給他們祝福,又對孩子的母親馬利亞說:「這孩子被立,是要叫以色列中許多人跌倒,許多人興起；又要作毀謗的話柄,叫許多人心裡的意念顯露出來；你自己的心也要被刀刺透.」

2:36又有女先知,名叫亞拿,是亞設支派法內力的女兒,年紀已經老邁,從作童女出嫁的時候,同丈夫住了七年就寡居了,

2:37現在已經八十四歲,並不離開聖殿,禁食祈求,晝夜事奉上帝.

2:38正當那時,她進前來稱謝上帝,將孩子的事對一切盼望耶路撒冷得救贖的人講說.

2:39約瑟和馬利亞照主的律法辦完了一切的事,就回加利利,到自己的城拿撒勒去了.

2:40孩子漸漸長大,強健起來,充滿智慧,又有上帝的恩在他身上.

釋經與應用

如果我們留意上文,就會發現耶穌往耶路撒冷聖殿去「報到」,並不是孤立的事件,因為在講述耶穌「報到」的故事之前,作者在上文記述了另一個大規模,勞師動眾的「報到」事件.

2:1當那些日子,凱撒奧古斯都有旨意下來,叫天下人民都報名上冊.

2:2這是居里扭作敘利亞巡撫的時候,頭一次行報名上冊的事.

2:3眾人各歸各城,報名上冊.

在歷史上來說,那次報名上冊其實是古代人口普查的典範.當時羅馬帝國的凱撒奧古斯都為了充分掌握國情,進行了三次人口普查.奧古斯都的人口普查很有系統和組織,他任命20名調查官員分赴各行省,與行省總督和地方官員共同組成「人口普查工作組」,對本行省的所有的居民及其財產進行調查.帝國的每一位居民都要宣誓不說一句假話,然後逐項申報個人財產及家庭狀況.(宮秀華:「奧古斯都與羅馬帝國初期的人口普查制度」)

今天,當我們讀路加福音2章下半章關於耶穌往耶路撒冷聖殿「報到」的故事時,讓我們同時參照2章上半章那個耶穌父母被凱撒奧古斯都要求「報到」的故事,在互相比對之下,看看有甚麼發現.

向誰報到:誰是生命的主?

第一,這兩次「報到」的對象是誰呢?又有甚麼不同的含意呢?讓我們再讀第一個「報到」故事:

2:1當那些日子,凱撒奧古斯都有旨意下來,叫天下人民都報名上冊.

2:2這是居里扭作敘利亞巡撫的時候,頭一次行報名上冊的事.

2:3眾人各歸各城,報名上冊.

2:4約瑟也從加利利的拿撒勒城上猶太去,到了大衛的城,名叫伯利恆,因他本是大衛一族一家的人,

2:5要和他所聘之妻馬利亞一同報名上冊.

按路加福音記載,這是凱撒的命令,要天下人民都向他「報到」.凱撒奧古斯都要求天下人民向他「報到」,所謂何事呢?有人認為,凱撒這樣做,除了了解國情之外,宣揚統治權力的意義更為重大,他在告訴報到者:「你們已經被我征服了,我的權力可以進到你們每一家每一戶每一個人,從今以後,我是你們的主宰.」統治者希望自己成為人民的主宰這種慾望,在古代社會是為普遍,就是在現代,在某些極權國家仍然可見.

在亞洲東北某國,去年的流行歌曲是《除了他我們誰也不認》,其中一段歌詞如下:

就算風雲變幻,逆風吹來,我們的心中只有您一人,永永遠遠生死與共,只擁護愛戴您唯一的領導.偉大的金正恩同志,我們除了他誰都不認!偉大的金正恩同志,向您宣誓忠誠!(維基百科)

可以想像,凱撒奧古斯都要求人民向他「報名上冊」,或許也希望人民「心中只有他一人」,也希望人民「除了他之外誰都不認」,希望人民「向他宣誓忠誠」.事實上,當時的確有把奧古斯都神化,稱他為「主」,為「拯救者」的行動.

不過,路加福音報導這件「報名上冊」的帝國大事後,接著報導一個民間「報名上冊」的小故事.這個民間小故事和那件帝國大事的意識形態並不一致,因為當羅馬帝國企圖透過要求人民報名上冊來滲透凱撒是主,是拯救者這信息的時候,民間的小故事揭示真正的主和拯救者另有其人.首先,我們看看在民間的「報到」故事中如何回應「凱撒是主」這宣稱,下文再探討如何回應「凱撒是拯救者」的另一個宣稱.

2:22按摩西律法滿了潔淨的日子,他們帶著孩子上耶路撒冷去,要把他獻與主,

這好像是平淡的行動,不過,意義卻深遠.試想想前文鋪陳了一種怎樣的政治氣氣?在凱撒要求效忠的大氣候之下,在「凱撒是主」的政治宣傳之下,約瑟和馬利亞按著他們的信仰把孩子帶到耶路撒冷聖殿去,要把他獻給凱撒以外的另一位主.這樣的行動,似乎帶有顛覆意味.且慢,這並不表示約瑟和馬利亞在搞革命,也不能指控他們背叛凱撒.前文指出,他們是遵從凱撒命令的百姓,馬利亞縱使腹大便便,也長途跋涉,按著要求往約瑟的老家猶大去報名上冊.不過,他們把兒子帶到耶路撒冷獻給主,即是耶和華上帝,又似乎顯示,在他們心中,「凱撒是主」的政治宣傳,並沒有取代「耶和華是主」的信仰.他們被要求向凱撒「報到」,但他們仍然繼續向耶和華上帝「報到」.他們明白,凱撒是當時的統治者,但並不是嬰孩生命的賜予者.凱撒的力量遍及每一家每一戶每一個人,但不是他們的兒子的最終主宰.約瑟和馬利亞這對父母能夠分辨得到,應該把甚麼獻給凱撒,把甚麼獻給上帝；應該教導兒子和凱撒建立怎樣的關係,又應該教導兒子和上帝建立甚麼關係.凱撒的力量雖然強大,但約瑟和馬利亞知道,凱撒不是人生的全部,若要把兒子獻上,對象不是凱撒,而是上帝.

今天,社會彌漫著的競爭壓力也非常強大,恐嚇著父母不能讓自己的孩子「輸在起跑線上」,壓迫父母認同成功就是一切.在這樣強大的力量支配之下,父母把孩子獻給誰呢?帶孩子到甚麼地方「報到」呢?心理學家大衛•艾爾金德在他的著作中提出「忙碌兒童」這個概念.後來更有人提出「忙碌兒童綜合症」(Hurried child syndrome)這種現象,所指的是家長將孩子的生活安排得過滿,強迫孩子不斷參加各種學習班和培訓班,以幫助他們為入學做好準備,極力要求孩子在學業上取得佳績,並且希望孩子像小大人一樣行為和做事.

今天,父母不會忘記為孩子去出生登記處「報到」,但會否忘記帶孩子到上帝面前「報到」呢?當父母忙著把孩子交給這位老師,交給那位教練的時候,會否忘記把孩子交給上帝呢?約瑟和馬利亞這對父母提示我們,雖然很多力量企圖主宰孩子的人生,企圖要成為孩子的主,然而,誰才是孩子生命的賜予者呢?誰才配得作為孩子的主呢?若要把孩子交托,除了上帝,還有誰配得呢?

向誰報到:誰是拯救者?

凱撒命令人民報名上冊,讓人民意識到他的主權,宣揚他的拯救.然而,如果凱撒是拯救者,他的拯救是怎樣一回事呢?會不會只是把人民由一位統治者「過戶」到另一位統治者的手上呢?人民被奴役的狀況會改變嗎?在凱撒奧古斯都那個時代,統治者要求人民報名上冊不外乎兩個目的,一為徵兵,二為徵稅.有人說:「古來征戰幾人回」,徵兵,差不多等於要了你的命；又有人說:「苛稅猛於虎」,徵稅,雖然沒有直接要你的命,但交重稅的感覺也是很要命.如果凱撒果真是拯救者,他要求人民向他「報到」,在某種意義上等於要了人民的命.

真正的拯救者是怎樣的呢?路加福音在接著的「報到」故事作了一個對比,澄清了誰是真正的拯救者,以及真正的拯救是怎樣的.

2:22按摩西律法滿了潔淨的日子,他們帶著孩子上耶路撒冷去,要把他獻與主,

為甚麼進行這樣的行動呢?路加馬上加以解釋:

2:23(正如主的律法上所記:「凡頭生的男子必稱聖歸主.」)

原來,這樣的行動是和一個偉大的拯救故事有關的.在出埃及記13章14-16節如此記載:「日後,你的兒子問你說:『這是甚麼意思?』你就說:『耶和華用大能的手將我們從埃及為奴之家領出來.那時法老幾乎不容我們去,耶和華就把埃及地所有頭生的,無論是人是牲畜,都殺了.因此,我把一切頭生的公牲畜獻給耶和華為祭,但將頭生的兒子都贖出來.這要在你手上作記號,在你額上作經文,因為耶和華用大能的手將我們從埃及領出來.』」約瑟和馬利亞把耶穌帶到聖殿向耶和華上帝「報到」,就是要把他們的孩子帶進這個拯救的故事之中,讓他成為這個拯救故事的一部份,延續著這個拯救故事.

昔日,拯救者從強大的埃及法老手上,把約瑟和馬利亞的祖先拯救離開為奴之地,不再做奴隸.路加福音藉著這個「報到」故事指出,這位拯救者已經臨到人間,他不是凱撒奧古斯都,他是誰呢?有沒有人誰認出他呢?路加告訴我們最少有兩個人認出真正的拯救者,第一位名叫西面:

2:28西面就用手接過他來,稱頌上帝說:

2:29主啊!如今可以照你的話,釋放僕人安然去世；

2:30因為我的眼睛已經看見你的救恩－

2:31就是你在萬民面前所預備的:

2:32是照亮外邦人的光,又是你民以色列的榮耀.

第二位是女先知亞拿,

2:38正當那時,她進前來稱謝上帝,將孩子的事對一切盼望耶路撒冷得救贖的人講說.

能夠看見真正的拯救者,為甚麼是這兩位長者呢?當時那些宗教上的專家和領袖們呢?為甚麼他們看不見呢?對於第一個問題,路加福音提供了以下資料讓讀者細味:

2:25在耶路撒冷有一個人,名叫西面；這人又公義又虔誠,素常盼望以色列的安慰者來到,又有聖靈在他身上.

2:36又有女先知,名叫亞拿,是亞設支派法內力的女兒,年紀已經老邁,從作童女出嫁的時候,同丈夫住了七年就寡居了,

2:37現在已經八十四歲,並不離開聖殿,禁食祈求,晝夜事奉上帝.

路加福音指出,西面和亞拿是虔誠的人,是真正的信仰者.耶穌說:「清心的人有福了!因為他們必得見上帝.」(馬太福音5:8)那麼,當時的宗教專家和領袖又為甚麼看不見真正的拯救者呢?約翰福音記載了他們的想法讓讀者參考:

11:47祭司長和法利賽人聚集公會,說:「這人行好些神蹟,我們怎麼辦呢?

11:48若這樣由著他,人人都要信他,羅馬人也要來奪我們的地土和我們的百姓.」

11:49內中有一個人,名叫該亞法,本年作大祭司,對他們說:「你們不知道甚麼.

11:50獨不想一個人替百姓死,免得通國滅亡,就是你們的益處.」

原來,祭司長和法利賽人心中想著的是凱撒,他們以凱撒的權力角度去看待耶穌,因此,嬰兒耶穌在他們眼前,只是一個無權又脆弱的嬰孩；就是面對成人耶穌所行的神蹟,他們都看不到拯救者已經臨到,只看到凱撒的威脅.對他們來說,凱撒的力量實在太強大了,就算他們不願意認同凱撒是拯救者,凱撒也令他們看不到他們盼望的拯救者.今天,會不會也有一些強大的力量,令我們誤以為它是拯救者呢?或者令我們看不見真正的拯救者呢?

真正的拯救者和凱撒這冒認的拯救者有甚麼分別呢?凱撒奧古斯都這類靠著軍事力量的冒充拯救者,他美其名解放天下人民,其實只是把天下人民變為他的奴僕.今天的書信經課(加拉太書)指出,真正的拯救者是把人從奴僕釋放變為兒子:

4:4及至時候滿足,上帝就差遣他的兒子,為女子所生,且生在律法以下,

4:5要把律法以下的人贖出來,叫我們得著兒子的名分.

4:6你們既為兒子,上帝就差他兒子的靈進入你們的心,呼叫:「阿爸!父!」

4:7可見,從此以後,你不是奴僕,乃是兒子了；既是兒子,就靠著上帝為後嗣.

路加福音敘述這兩個「報到」的故事,指出縱使在迷惑和壓力之下,仍要懂得分辨誰是主,誰是拯救者.今天,我們能夠撥開我們時代的迷霧,抗衡我們時代的壓力,分辨得到誰是主,誰是拯救者嗎?

總結

人生難免要「報到」,向誰「報到」可以產生很不同的意義.醫生每天回醫院「報到」,會否同時向創造生命的上主「報到」呢?執法人員每天向上司「報到」,會否同時向公義的上主「報到」呢?到最後,每個人都要作相同的「報到」.到甚麼地方「報到」呢?不錯,是死亡登記處.不過,人人都知道,死亡登記處只能發出死亡證,並不是離世者的歸宿,人最終仍然要向生命的主及拯救的主「報到」.你看得出誰是生命的主嗎?分辨得出誰是拯救之主嗎?



\newpage

\section{洞悉人生視野真實現實超現實的世界}
\label{sec:272}
\textbf{許開明牧師 - 洞悉人生視野:真實,現實,超現實的世界}
\newline
\newline
\hyperref[sec:271]{< < < PREV SERMON < < <}
~
\hyperref[sec:toc]{[返目錄]}
~
\hyperref[sec:273]{> > > NEXT SERMON > > >}


講題:大開眼界:透視人生真實,現實,超現實的世界　The Visual Field of Life: The World of Truth Reality and Virtual Reality

講員:許開明牧師

場合:香港中文大學崇基學院禮拜堂

日期:2015年1月4日

引言

今天是2015年的首個主日,新年伊始,恭祝各位新年進步!

今天的講題是「大開眼界:透視人生真實,現實,超現實的世界」

新的一年,我們將怎樣渡過?

怎樣才可以讓有限的日子過得好?

我相信關於這兩條問題的答案應該與我們如何看待這個世界有關!

成見

有一則故事或許大家曾聽過,或者就曾發生在你身上.

一天晚上,有位小姐在赤蠟角機場等候登機,她慣常地找個地方坐下,拿出一包薯片,一本書來,因為她喜歡一邊看書,一邊吃薯片,別有一番滋味.當她沉醉於看書時,她無意中發現,隔鄰的男士,竟然未經許可伸手來取她的薯片吃,若無其事的,一片片的隨意的放進口裡吃.

開始,她假裝若無其事,不想因偷吃薯片小事,在大庭廣眾,大吵大鬧.她繼續看書,並使勁地咀嚼薯片,讓對方知道她的反應.那偷薯片賊仍然不為所動,繼續不停的吃,時間一分一秒的過去,她也愈來愈憤怒,是否應該啞忍呢?

她察覺這位男士全身肌肉,相當高大,好女不吃眼前虧呀!還是忍下去!但想深一層,小女子怎能因他大強壯,就怕她,真是忍無可忍了!正想鼓起勇氣破口大罵時,發現薯片袋只餘最後一片了.

她心想:「給你最後一次機會,如果你讓給我吃,那就算了吧!」怎知那男士把最後一片也拿起,對著小姐露出好像善良的微笑,然後很小心地把薯片分開一人一半,遞給她半片,然後很開心的吃著自己的半片!

這位小姐心想:「天啊!這傢伙臉皮也太厚了吧!吃了人家半包薯片,一聲多謝都沒有,還不知羞的對我微笑,簡直變態!」她帶著怒氣登上機,當飛機飛定後,她把手伸入手袋取她的書,無意中摸到一樣東西,是什麼?

是一包未開封的薯片!原來,她一直吃的薯片,是那位男士的.那男士才有風度,她才是偷薯片賊!不問自取的不是那男士,竟然錯怪了人家!

各位,這位小姐,她在現實的世界裡有成見,這成見使她活在自已的心靈世界,這世界與事實不符.這種情況經常會發生.我們不留心的話就會產生誤會,可能釀成大錯!

我們的經課是約1:14-34節.當中,有個重要的信息:

約翰福音的希臘思想

約翰福音成書時間是在第二世紀,此時基督教的信仰已傳到希臘,進入世界.

作者當時居於以弗所,他是一位熟識猶太教及希臘文化的人,他非常關注如何把基督教信仰傳給希臘人,傳給全世界.他發現希臘人的神Logos與猶太人的神相當相似,這便是他找到的「共通語言」.

希臘人有兩個偉大的觀念:

洛格斯(Logos)的觀念.

Logos創造世界.希臘文的洛格斯有兩個意義,一個是言語,一個是理性.就猶太人而言,他們完全熟悉那滿有能力的上帝之言.「上帝說,要有光,就有了光」(創一3).就希臘人而言,則完全熟悉理性的思想,理性表現於秩序與規律.他們看世界是一個偉大而又可靠的秩序:

晝夜永不落空地相繼；

一年四季,交替不息；

以從不改變的軌跡運行；大自然有它的永恆的法則.

這秩序是怎樣產生的呢?希臘人會毫不猶疑的回答,是上帝的洛格斯,上帝的心思(理性)負責世界莊嚴的秩序.

是誰使人類有能力去思想,去推理,去知道呢?希臘人又會毫不猶疑地回答,是上帝的洛格斯,上帝的心思在人心中使人成為有思想,會推理的動物.

因此,約翰福音作者就把握住這一點,並用這種思想方法來談論耶穌.

他對希臘人說:「上帝統領一切,掌管一切的偉大心思,吸引住他們的一生.現在上帝的心思已藉耶穌這個人來到世間.只要看看他,你們就見到了上帝的心思意念.」約翰發現了一個新的範疇,有助於希臘人去思想,耶穌就是上帝的洛格斯以人的樣式行動.

兩個世界觀.

希臘人常常認為有兩個世界.現實的世界就是我們活在其中的世界.它本身雖有其美好之處,但它只是一個影子,複印本,一個不真實的世界.另一個是真實的世界,在這個世界中,偉大而實際存在的事物要存到永遠.

柏拉圖(Plato)曾將這些觀念整理成為他的思想體系.他認定在那個不可見的世界(真實世界)中有每一件事物的完美樣式,而這個世界(現實世界)的事物,只不過是這些永恆樣式的影子般的複印本而已.

簡單說,柏拉圖認為真實世界裡,某處有一張桌子的完美樣式,我們這世界所有的桌子都是不完美的複製品；某處有善與美的完整樣式,所有世間的善與美都不過是不完整的複製品.

而那偉大的實在,那最高的思想；一切格式的格式,一切樣本的樣本就是上帝.問題是,怎樣方能進入那真實的世界?怎樣方能擺脫這世界的影子而進入永恆的真實?

耶穌就是來到世間的實在!

約翰福音作者宣稱說:這正是耶穌使我們能夠做到的事.耶穌就是來到世間的實在.在希臘文,真正的(real)這詞語為ale-thinos,與ale-thes很接近,

後者是真正的意思,ale-theia,意即真理(truth).

欽定本和標準修訂本都把ale-thinos譯為true(真實的)；如果譯作real(真正的)就會好得多.

耶穌是真正的光(約一9)；

耶穌是真正的糧(約六32)；

耶穌是真正的葡萄樹(約十五1)；

在我們這個如影子的,不完美的世界中唯有耶穌才是真正的實在.

(按:true及real──的中譯非常接近.中文聖經和合,呂氏,現代等譯本都譯為真(的).英文K.J.V.及R.S.V. 譯為True,T.E.V.則譯為real).

神蹟(Semeia)的意思是窗戶.

耶穌的每一個行動(神跡)不單是歷史中某一時間的一項簡單的作為,而是一扇窗戶,讓我們經由它見到「真實」.透過神跡這扇窗戶引伸出別的「真實」來.這就是約翰福音作者稱耶穌的奇蹟為神蹟(Semeia)的原因.耶穌所行的神蹟不但使人驚奇,更是窗戶,使人見到真實的上帝.

透過耶穌基督的生命,我們可以看到:

耶穌是真實的,因為他是神,是真光,是真理;

耶穌是現實的,因他道成肉身(人身),進入這個現實的世界,人可見他,摸他,跟他接觸…因耶穌引出上帝的真實,約翰就引出耶穌是基督.

耶穌是超現實的,因他從主裡復活,永遠活著,升天掌權!

耶穌以這三度空間,讓我們可以更深的認識祂和思想祂,這是我修讀神學多年的發現,將會放入我著作中.

約翰福音論述耶穌

約翰福音作者引出耶穌,使人大開眼界.

(施洗)約翰為耶穌作見證,口喊著說:「修直主的道!」(約1:15- 18)

施洗約翰當時推介耶穌,只靠自己的「呼喊」.沒有報紙,刊物,電台,電視,網絡!

約翰福音作者引出耶穌基督幾個特點:

耶穌在約翰之後,也在之前.他可見的身體在約翰之後,他的靈體在約翰之前!基督是先存(Pre-Existence),現存(Existence),永存(Eternity).

耶穌是豐盛的.包括擁有上帝的智慧(Wisdom),能力(Power),愛德(Love),和有關上帝的一切知識(Knowledge of God).

耶穌擁有恩典(Grace)和真理(Truth)!摩西的律法,是使人知道規矩,懂得怎樣做,做不到時所應受怎樣的懲罰；耶穌的恩典,給行不到的人有悔改的機會,改正歸真.

耶穌引出上帝來,使人可以接觸父上帝.

猶太的當權者派人來,盤問施洗約翰:「你是誰?」(約一19)

各位,你曾否試過在街上走,突然間有人過來問:「你是誰?」嗎?我年少住官塘雞寮時,曾有街童問我.你猜我怎樣回答呢?看到對方不懷好意,我不敢告訴他我的姓名.但又不能不回答,所以靈機一動說:「我係我!」(我是我)結果被打了幾拳,Why?不知!施洗約翰也曾被人追問:「你是誰?」

第一批是從首都耶路撒冷來(即代表今日的中央):聖殿的祭司和利未人集團

問:兄弟,你是誰?你整日大聲呼喊,吸引一大群人,聽你講,又領洗,跟著你!(現代版)

答:我不是你猜想的那位基督,我不是猶太人等候的那位基督!

問:你是否以利亞大先知呢?

答:不是!

問:你是那位大家講的那位「好堅」的先知(卓越的先知)嗎?

答:不是.

問:你到底是誰?你自己說:你是誰?等我向「大佬」有個交待!

答:我就是聖經上以賽亞先知所說:「在曠野有人聲喊著說『修直主的道路』.」(古時,王者經過的地方,找人先開路,使君王在平安穩妥,平坦的路上經過.)我只是一個聲音,讓人準備自己等候基督的來到!

第二批是地方集團派來的,律法:法利賽人集團,作進一步的盤問.

問:「你既不是基督,又不是以利亞,你既不是那先知,為何幫人施洗呢?(你憑著什麼身份?什麼權柄為人施洗呢?)」

答:「我是用水施洗!(約翰沒有回答他為何洗,憑什麼權柄去洗!)但有一位將會用聖靈為人施洗；這人在我之後來的,比我年輕；但其實他在我之前.(可以想像聽者覺得他語無倫次了!)這人在你們當中,不過你們還沒有認識他；他是非常尊貴的,我連幫他解鞋帶洗腳也不配!」(但這二批人對耶穌未有興趣,只對約翰有!)

到底約翰是誰?猶太高層要動用兩批人盤問他?

出身在祭司家庭,是貴族,卻不做祭司,到曠野和高山做先知.

表現出位,受人擁戴:根據《聖經》的描述,施洗者約翰長大後在曠野工作.他身穿的駱駝毛的衣服,吃的是蝗蟲野蜜,過著最簡單的生活,是不貪圖地位,權勢,具有淡泊心志之人.

傳悔改的洗禮.許多人(數以千,以萬計)跟隨著約翰,使當權者不得不注視他!大膽指出人的錯,連當權者希律,也直言不諱.施洗約翰在傳教的過程中勇敢指出人的錯誤:有很多人聽了他的講道後接受了施洗者約翰的洗禮,施洗約翰在傳教的過程中並不是以甜言蜜語來籠絡這些人,而是嚴肅地指出他們的罪.當時的王希律·安提帕斯((古希臘語:Ἡρῴδης Ἀντίπατρος,Hērǭdēs Antipatros,英語:Herod Antipas,簡稱:Antipatros))做了一件不合猶太人律法的事,施洗約翰不畏王權與自身的安危,勇敢指出希律·安提帕斯的罪,因此他被希律·安提帕斯下令逮捕,但希律·安提帕斯顧忌他的威望,一直不敢殺他.約翰敢於發正義之聲,義正言辭的指出人的錯處,為民請命,受群眾歡迎,是當權者最忌的.因此,先派祭司集團來盤問他,要摸底,古時的人已懂得摸底!到底他是誰?他有何後台?

約翰給他們的答案是:

他不是基督(不擁有反羅馬,反希律的勢力)

他不是以利亞(不擁有先知的勢力,群眾擁戴的勢力)

他不是那先知(並非像摩西般的民族英雄,頒律法者,政治與宗教的勢力)

他只是耶穌基督的開路先鋒,猶太集團不用擔心,只是小人物一個!

然而,猶太領袖仍然懷疑約翰,於是派出德高望重,猶太會堂的主持,受人尊重的法利賽人來問:一針要見血,單刀直入:「你既不是基督,你又不是以利亞,你既不是那先知, 你為何施洗呢?(你憑著什麼身份?憑什麼權柄為人施洗呢?)」

假如有人在主日走上聖壇說:「誰想洗禮?上來洗禮!」,作為牧師的我自然會問:「你是誰啊?」如果約翰未能回答,他憑著什麼身份,他憑什麼權柄為人施洗,法利賽人和祭司聯合起來,可捉拿,甚至可處死他!

當時,約翰是所答非所問的!為何又安然無事呢?唯一的解釋是當時無數人跟隨著約翰,他們動手的話,恐怕會引起騷動!反而,約翰藉此機會,引出耶穌基督來!約翰巧妙地經過當時的政治及宗教者的盤問,還引出了耶穌基督,真是令人大開眼界!避重就輕,反敗為勝!

到第二天(次日),跟隨著約翰的群眾,一定很早到約旦河邊,看看還有什麼令人驚奇的事發生.果然,昨天談及的主角耶穌基督出現了!

當約翰見到耶穌時,說:「看啊,上帝的羔羊!」(29)

耶穌像上帝的羔羊,在聖殿中,把自己的生命獻上,替人贖罪!耶穌像逾越節的羔羊,犧牲生命,流出寶血,灑在門楣,使人不致被殺害!在約翰的心中,想到耶穌是這樣的救贖主!

約翰另一個看見,是看見「聖靈彷佛像鴿子(神聖鳥,不可捕捉和殺害；上帝的靈像鴿子),從天而降,降在他(耶穌)身上!」(32)因此耶穌乃用聖靈為人施洗的.證明他是上帝的兒子,是基督.

從約翰的信息中,我們可透視到三樣東西:

真實的耶穌:耶穌是真理,豐盛,先存,現存,永存,有價值的.

現實的耶穌:耶穌可被看見,被我看見(約1:29-34:29,32,33,34)共5次.

超現實的耶穌:上帝的羔羊,聖靈像鴿子在耶穌身上.

它們為要使看見的人明白其所代表的,並且皆指向另一個事物.羔羊,指向上帝的羔羊

鴿子,指向聖靈的鴿子.這三個透視,對我們的生命深具意義

根據柏拉圖的觀點:

真實的世界:看不見的世界,真實的,理想的,完美的,有價值的,永恆的.

現實的世界:看得見的世界,實際的,非完美,變化中的,短暫的.

超現實的世界:不真實的,但又是現實的延伸,心靈的,想像的,虛擬的世界,看不見的超現實的世界, 超越現實的,但可以當成真的世界.

這三個世界的意義是什麼?你活在怎樣的世界當中?你應該活在怎樣的世界當中?

真實世界:看不見的,最真實的,最完美的,最有價值的,最理想的世界.這是我們要追求的方向,人太現實時,往往忘記追求那些「真實」的東西.我對於大學生追求真普選,過去七十多天鍥而不捨追求理想的精神是欣賞的.然而是否用佔中這種方式,我是有保留的.

現實世界:接受現實,不追求完美主義.追求完美的人在現實世界裡是痛苦的,容易患憂鬱症.「不完美主義」是日本近幾年心理學上非常重要觀念,為拯救日本抑鬱症及自殺率居第一社會,更需要「不完美主義」的想法.日本學者齋藤茂太提倡「人生百分之八十主義」,因為要達到「百份之百」很易受挫折.現實的世界,不完美是自然的,追求完美反而失去自我,使自己陷入困境.

面對逆境,不向現實低頭

有一句英文諺語:When life gives you lemons ,make lemonade!字面翻譯就是:當生命帶給你檸檬時,拿它們做檸檬水吧!

檸檬,肉汁酸,皮又苦,沒有人喜歡吃.但檸檬水味酸帶甜,很多人可以接受.尤其是冰凍的檸檬水,在炎夏更是生津止渴的上佳飲品.

那句諺語其實有一個引伸的意思,就是:當生命中出現不如意,又酸又苦的經歷(檸檬)時,你可轉念把它轉化成不一樣的好事(又甜又好喝的檸檬水).

一位小女孩 Alexandra Scott,她剛滿一歲不久,就診斷出患有神經母細胞瘤(Neuroblastoma), 一種極難治愈的兒童癌症.Alex的健康嚴重受影響,但她很堅強努力,到兩歲可靠腳架站起來.跟著下來,連續不斷接受治療,可以開始學走路.她的健康仍是不斷惡化.到四歲時,Alex 跟父母起了個念頭,在自己的家前園開個小攤子擺賣檸檬水籌款,把籌到的\$2000捐給她的醫院做研究費去幫助其他患癌的兒童.接著下來,每年Alex和家人仍繼續他們的檸檬水小攤檔 (Lemonade Stand)籌款,而且感動了各地的人支持,紛紛加入開設檸檬水小攤募款的行列.2004年 Alext 8歲時去世,當時的 Alex’s Lemonade Stand 活動已籌得超過一百萬的研究費.

Alex的生命很真實地註解和活出了這句諺語:『當生命帶給你檸檬時,拿它們做檸檬水吧!』她把生命中又酸又苦的檸檬,變成無數患癌的兒童帶來醫治的希望.

超現實的世界 (包括心靈的世界／想像的世界／虛擬的世界):

時間一天一天的過,

過著那個屬於我的虛擬和平世界,

過著那個屬於我的虛擬快樂生活,

躲在我的小小虛擬世界,看著我的螢幕…

陪伴著我的心情,一路的走下去…

網絡雖是「虛擬」世界,但是,卻指向「現實」世界中的實體,因此與「現實」世界是有互動性,關係性及影響性的.如此,我們便不能夠就此一刀切地,把所謂虛擬的「網絡世界」抽離於現實,並將它與現實對立起來.因此,人或許,從一開始便過份的誇大了「虛擬」世界與「現實」世界的差異,把「虛擬」世界中的一切從「現實」抽離出來.以致,這前設成為了人在網上任意妄為的藉口.

其實,網絡是真實世界的一個組成部分.二者緊密聯繫,不可分割.世界越是步向真正的網絡化時,人的生活便越離不開網絡.「虛擬」與「現實」之間的差異縮小,從而與真實社會同在,並在許多領域上形成疊合.

如此,網絡雖是「虛擬」,但卻不是脫離「現實」而獨立存在的另一空間.毫無疑問,人才是「虛擬」與「現實」的主角.

「人既是現實社會與虛擬社會溝通的媒體,又是現實社會與虛擬社會的主體.」因此,人便為二者中「橋」的角色,一方面將現實世界中的「感情」投射進虛擬世界,另一方面,也把虛擬世界中的「感情」帶於現實世界的裡面.

可以這樣說,網絡雖然是虛擬,但人從網絡所經驗到的卻是真實.其經歷彷彿有如活現在現實世界的當中.如此,無論「虛擬」與否,有血有肉的人仍是當然的「主角」.

結語

人生經常面對著這三個世界:

「我要真普選!」這是追求長遠價值,長期的努力.

「我要真出糧!」這是現實的追求.

「我要真網友!」這是在網上,超現實的追求.

真實:是最真,最美,最好的(看不見)

現實:是真實的影子,引出真實(看得見)

超現實:是現實影子,現實的延伸(虛擬的世界,以假為真,也看得見,又像看不見)

我常提醒自己,這三個世界都各有長短:

真實的世界固然有意義,人不應忽視,但一味追求完美真實的世界,人的生活會變得很辛苦!

現實的世界來得實際,但人太現實,毫無理想,人不會有進步:

超現實的世界,使人追求虛的,浮的,使人脫離現實,要加以小心!

但不懂超現實的人太乏味,只會原地踏步!e.g 有時發點夢,像小鳥飛!

【兩億三千萬】

日前,收到朋友送來兩億三千萬,未敢獨自享用,要與大家一齊分享!

第一個憶,是回憶:

回憶自己的一生,有多少值得珍惜的親情,友愛；值得感恩的人生快事；值得自豪的大小成功；值得欣慰的賞心樂事；儘管並非事事如意,但能平平安安地活到今天,留下不少快樂回憶,足以老懷安慰.

第二個憶,是失憶:

把過去種種不愉快的經驗,耿耿於懷的傷害,別人對自己的不夠好,吃虧與不幸的遭遇,通通拋向九霄雲外,人就分外地輕省與年青.

第一個千萬,千萬要開開心心過日子:

凡事一笑置之,用正面積極的角度面對,在任何環境下都懂得自得其樂,笑口常開.

第二個千萬,千萬要保重身體:

小心飲食,少肉多菜,恆常保持輕巧的運動,身心康泰.

第三個千萬,千萬要開放心靈,接受新事物:

不要以年老固步自封,和世界保持接觸,聆聽不同的聲音,讓今天的我比昨天進步,與人分享分擔.

我們都要向著這兩億三千萬邁進,快快抓緊它啊!(不是從押款車掉下來的,請放心!

期盼眾弟兄姊妹都得到這二億三千萬,這筆人生巨款,稅局都不能向你抽稅的.

求主賜福各位!



\newpage

\section{多馬的認信}
\label{sec:273}
\textbf{梁明浩牧師 - 多馬的認信}
\newline
\newline
\hyperref[sec:272]{< < < PREV SERMON < < <}
~
\hyperref[sec:toc]{[返目錄]}
~
\hyperref[sec:274]{> > > NEXT SERMON > > >}


約翰福音 20:24-29
\newline
\begin{longtable}{cl}
\hline
\hline
章節 & 經文 (和合本修訂版)\\
\hline
20:24 & \begin{tabularx}{0.7\textwidth}{X} 那十二使徒中，有個叫低土馬的多馬，耶穌來的時候，他沒有和他們在一起。 \end{tabularx} \\ \\ \relax
20:25 & \begin{tabularx}{0.7\textwidth}{X} 其他的門徒就對他說：「我們已經看見主了。」多馬卻對他們說：「除非我看見他手上的釘痕，用我的指頭探入那釘痕，用我的手探入他的肋旁，我絕不信。」 \end{tabularx} \\ \\ \relax
20:26 & \begin{tabularx}{0.7\textwidth}{X} 過了八日，門徒又在屋裡，多馬也和他們在一起。門都關了，耶穌來，站在當中，說：「願你們平安！」 \end{tabularx} \\ \\ \relax
20:27 & \begin{tabularx}{0.7\textwidth}{X} 然後他對多馬說：「把你的指頭伸到這裡來，看看我的手；把你的手伸過來，探入我的肋旁。不要疑惑，總要信！」 \end{tabularx} \\ \\ \relax
20:28 & \begin{tabularx}{0.7\textwidth}{X} 多馬回答，對他說：「我的主！我的神！」 \end{tabularx} \\ \\ \relax
20:29 & \begin{tabularx}{0.7\textwidth}{X} 耶穌對他說：「你因為看見了我才信嗎？那沒有看見卻信的有福了。」 \end{tabularx} \\ \\
[1ex]
\hline
\hline
\end{longtable}
講題:多馬的認信　Thomas’ Confession

經文:約翰福音20:24-29

講員:梁明浩牧師

場合:香港中文大學崇基學院禮拜堂

日期:2015年1月18日

Have you ever had any nick names? Many people have their nick names. Some nick names are pleasant, some are very funny, and still others are a little bit harsh to hurt people’s heart. There was a man who had a dishonorable nick name. His name is Thomas. People often call him ‘doubting Thomas.’ From the Gospel of John, which we read this morning, we can guess how Thomas got this nick name.

All the disciples except Thomas had seen the risen Lord. And they told Thomas that they had seen Jesus. Then Thomas said that he would not believe until he really sees Jesus and touches him. Because of this episode, the later generations began to call him ‘doubting Thomas.’

But, using that yard stick, that nick name has to be given not only to Thomas but also to all other disciples. Before Jesus was crucified, he told them, “I am going to be killed soon but after three days I will rise again. After I have risen, I will go ahead of you into Galilee.” Now, what should they do after Jesus died? For three days, they should have eagerly waited with excitement, expecting the coming victorious resurrection. But that was not the case.

At the day of resurrection, a couple of women broke the news to the disciples that the Lord has risen. When they heard that news, where do you think they should have gone? 1) To their houses? 2) To the tomb, 3) to Galilee? Yes, they should go to Galilee to meet Jesus. But they couldn’t believe the news. So they ran into the tomb to check. Do you remember what they did after they saw the empty tomb? Luke records Peter’s response in this way, “…then he went home, amazed at what had happened.” Peter didn’t go to Galilee. And that night, Jesus appeared where the disciples gathered and showed them his body, but they didn’t believe it. So Jesus said to them, “Why are you frightened, and why do doubts arise in your hearts?”

Therefore, the nick name “doubting” should be given to all disciples.

In fact, Thomas was deserved to be called ‘a man of loyalty.’ Do you know why? One day when Jesus heard that Lazarus died, he said to his disciples, “Let’s go back to Judea.” Then, all other disciples got scared and said, “The Jews there tried to stone you, and yet you are going back?” But Thomas said, “Let us go with him so that we may die with him.” Thomas was a kind of person who was willing to risk his own life to follow Jesus.

But for whatever reason, he became a person who would never believe Jesus’ resurrection. We don’t know why. But one thing we can say is that even a person like Thomas who once had a strong faith and passion can lose a faith and doubt.

How about you? Have you ever doubted before, because of fear or anxiety or disappointment?

Thomas did not believe Jesus’ resurrection. He wouldn’t change his mind even after he was told the news from his fellow disciples.

Then, John begins verse 26 in this way, “Eight days later…”

John says that ‘Eight days’ passed by without anything happening to Thomas. How do you think he felt during those eight days? What do you guess he thought for those eight days? For him, those eight days might have felt longer than eight years. He must have experienced painful loneliness. Jesus appeared to all disciples but him. Jesus met them. But Jesus never appeared to him for those long eight days. For Jesus, it is not a big deal to show himself to Thomas.

Thomas could have thought like these…:

“If it is true that the Lord rose again, how could He treat me like this? I am the same disciple as others and have followed him faithfully with my best. Did I put my life in worthless and hopeless things? Am I not his disciple? Am I not a child of God? All the other disciples seem to find favor in God’s eye except me. Am I abandoned?”

Those painful thoughts must have kept crossing his mind. Thomas must have prayed and beseeched wholeheartedly, and painfully waited and waited for Jesus to appear. Does our Lord know his mind and care what Thomas is going through? Those long eight days are just passing by without any answer to his earnest prayers.

Is there anyone in this room who is currently going through the similar eight days as Thomas had to experience, those eight days for which Jesus seemed not to care about Thomas and just kept silent to him? Is there anyone who is feeling that your longstanding prayers are not answered, your prayers for healing do not work, and problems in your family and finance and other problems in your life are still pending unresolved? And are you now passing through the darkest time in your life? Are you exhausted not able to pray anymore? Even then, do not despair. You will meet God’s timing. Many people of faith in the Bible are those who passed through long periods of time when God seemed not to care for them and kept silent. Some of them had to spend those times under a false accusation. Please do not give up. Even though you feel that God’s timing seems to be delayed, it will arrive. Our expected time will not always come into line with the time in God’s mind, but it will surely arrive. Therefore, our time of waiting for God’s timing is not a time of our being abandoned. It’s a time of grace when God is still working in our life.

Some time after Jesus was resurrected, Peter, John, and several disciples decided to go out to fish. As we all know, Peter is a veteran fisherman, and he knows very well about the Sea of Galilee. But that night, it seemed that fish were not cooperative with Peter. What a miracle for this skilled fisherman to catch nothing! Yes, it could be called a miracle. Somebody guessed, that night a miracle might have happened this way. When Peter cast his net to the right, Jesus drove all fish to the left. And Peter cast a net to the left, and then Jesus sent a signal to all fish to move to the right. If it was not the case, how come no fish had been caught all night?

While Peter and his friends were experiencing continuous failure in fishing, the Lord was preparing to meet those disappointed disciples. What they really needed was Jesus, not fish. And after passing through that long night, early in the morning they met Jesus. Jesus was waiting them on the shore. He shouted them, “Children, you have no fish, have you?” He already knew their circumstances. Then, he said to them, “Cast the net to the right side of the boat, and you will find some.” And hearing Jesus’ words, fish might have begun to move to the right side of the boat. Fishermen cast a net and caught large number of fish and were unable to haul the net into the boat.

There are times in our life that we have to go through, catching no fish all night. But even then, every moment of our life is under God’s grace because God is in our life always. Even the time of silence at which nothing happens is also the time of grace at which God is still working.

God was with Thomas during those eight days of dark silence.

After eight days passed, the Lord came again where his disciples gathered. This time, he seemed to come for Thomas. As soon as Jesus came in, he approached Thomas and said, “Thomas, see my hands. Put your finger here and touch this hole. Reach out your hand and put it into my side. Stop doubting and believe.” What Jesus said were the exact same thoughts as Thomas said to the other disciples eight days ago. That means… Jesus was listening to what Thomas was saying to his fellow disciples! That means… Jesus was listening when Thomas was praying. That means… Jesus knows what’s in Thomas’ heart, all his tears, sighs, and groans.

Psalm 139(vv.1-4) says, “You have searched me, Lord, and you know me.  You know when I sit and when I rise;    you perceive my thoughts from afar. You discern my going out and my lying down; you are familiar with all my ways.  Before a word is on my tongue, you, Lord, know it completely.”

Do you feel that the Lord is silent to you? Then, please remember! God knows our thoughts! God knows us. And God cherishes each one of us.

Thomas didn’t touch Jesus’ hands. He couldn’t do it. In fact, he didn’t need to touch him. Why? Because Jesus already knew all of his circumstances, heart, conditions, fears, disappointment, and groans. Jesus knew all about him. That’s more than enough for Thomas.

Now, I want all of you to remember what Jesus said to Thomas. Jesus told him, “Stop doubting and believe.” What is seen is not everything. As Paul says, “What is seen is temporary, but what is unseen is eternal.”(2 Cor. 4:18) We, who only live extremely limited to one generation, cannot fathom the will and works of God which are being fulfilled throughout all generations. We can only have a glimpse of a tiny part of them which emerge in this one generation we are living.

A long time ago, I heard a story. There was a frog. One day early in the morning, it met a dayfly, and they became good friends. In the evening, as it was getting darker, the frog told the dayfly, “Hey! It’s getting dark now, let’s play tomorrow.” And the dayfly said, “Tomorrow? What is tomorrow?” And the frog answered, “You don’t know what tomorrow is? It’s going to be very dark soon, and then you cannot play anymore. At night, we have to go to bed. We will get up in the morning. That’s tomorrow.” And the dayfly said, “Night? Tomorrow? Are you kidding? I don’t get it. There is no night, no tomorrow.” The next day, the frog couldn’t find the dayfly. A couple of days later, he found a new friend. It was a grasshopper. They had a very good time every day. The frog was very happy because the grasshopper understood when he said “See you tomorrow.” Several weeks later, it was getting cold. One day, the frog told the grasshopper, “Hey! I think it’s time to prepare for a long winter’s sleep. We can’t play in the winter. Let’s meet next year.” With amazement, the grasshopper asked him, “What is a winter? What do you mean by ‘next year’? I haven’t heard about that before.” The frog was amazed too, but kindly explained to him what winter is. But the grasshopper said, “Snow? Winter? Next year? I don’t understand what you are talking about. Hey, friend, there is nothing like winter, nothing like next year.” When the frog woke up next year, he couldn’t find that grasshopper.

The wisdom of a dayfly which lives one day is not like that of a grasshopper. A grasshopper which lives one season cannot understand the thought of a frog. We cannot fathom nor understand our creator’s wisdom, knowledge, and thought. I am sure that we all remember what God said through Isaiah. God said, “My thoughts are not like yours, and my ways are different from yours. As high as the heavens are above the earth, so high are my ways and thoughts above yours.” (55:9)

That God knows us. That God cares us. That God is with us. And that God helps us. Therefore, in the midst of every circumstance that surrounds us, let us hold on to the promise of God and shout to the Lord in our confession, “You are my Lord and my God.” And by faith, let’s live a victorious life, echoing this confession in every moment of our life…. until our Lord comes again.



\newpage

\section{堅石與空氣}
\label{sec:274}
\textbf{鄧瑞強博士 - 堅石與空氣}
\newline
\newline
\hyperref[sec:273]{< < < PREV SERMON < < <}
~
\hyperref[sec:toc]{[返目錄]}
~
\hyperref[sec:275]{> > > NEXT SERMON > > >}


講題:堅石與空氣　A Rock and A Breath

經文:詩篇62

講員:鄧瑞強博士

場合:香港中文大學崇基學院禮拜堂

日期:2015年1月25日

各位弟兄姊妹,早安.

今日的講道題目,是「堅石與空氣」.堅石,拿上手,是有質感的,有重量的；空氣,要捉也捉不住,虛無飄渺.一個重,一個輕.重與輕,是生命的一大問題.司馬遷說:「人固有一死,或重於泰山,或輕於鴻毛,用之所趨異也.」人總有一死,這死是重是輕,在乎這死所追求的目的是什麼.捷克作家昆德拉(Milan Kundera)寫了本《生命中不能承受之輕》去探討生命的重與輕的問題.這部小說後來被拍成電影.主角托馬斯(Tomas)將「愛」與「性」分開,對他而言,「愛」過於沉重,他只想享受「性」的輕狂.但是,一個純潔而痴情的特麗莎(Tereza),偶然地闖入了他的生命.這個「偶然」而來的女人,卻給托馬斯的生命帶來「必然」.特麗莎教曉托馬斯欣賞貝多芬最後一首四重奏曲的最後一樂章,這樂章的主題是:「非如此不可!非如此不可!」(Es muss sein)從此,這樂章的主題滲入托馬斯的生命裡,揮之不去.面對人生某些沉重的責任,如婚姻,如不向獨裁者低頭,托馬斯學曉了某種承擔,認為「非如此不可」.當然,一個人的生活,不能每樣事情都要「非如此不可」的,這將會太沉重,但是,當所有事情都失去「非如此不可」的理由時,我們便找不到做任何事的必然意義.一切都變得「輕省」,但也成為一種「不能承受之輕」,因為這種輕,輕得毫無內容,毫不實在.無怪乎昆德拉說:「我們也或多或少地贊同:我們相信正是人能像阿特拉斯(Atlas)頂天一樣地承受著命運,才會有人的偉大.貝多芬的英雄,就是能頂起形而上重負的人.」(昆德拉:《生命中不能承受之輕》,時報文化出版,頁51.)

我今日的講道經文,是今日的詩篇經課:詩篇62篇.透過這詩篇,我們一起思考生命的「重」與「輕」,「堅石」與「空氣」.

這首詩篇,在希伯來文裡,有一個明顯的特徵,就是在這詩篇12節的經文裡,有6節是以 ʼak 這個字作開頭的.ʼak 這個字,可譯成「only」(單單),也可譯成「truly, indeed」(事實就是這樣).作者如此用這個字,是引人注意,叫人留意事情的真相和底蘊. ʼak 這個字出現在以下的經文裡:第1節的「專等候神」的「專」字；第2節的「惟獨他是我的磐石」的「惟獨」；第4節的「專要從他的尊位上把他推下」的「專」；第5節的「專等候神」的「專」字；第6節的「惟獨他是我的磐石」的「惟獨」；第9節也有這個字,「下流人真是虛空」的「真」.

我好想將詩人的用字凸現出來,好顯明他的用意,故此我將這詩篇自行作一個直譯,我將 ʼak 這個字全部譯成「係呀,單單…」.

詩篇62篇:

係呀,單單在神裡面,我的生命,平靜下來.從他而來,我的拯救.

係呀,單單他是我的磐石,我的拯救,我的高臺,我不會被大大動搖.

要到幾時呢?你們攻擊一個人,要殺死他,你們每一個人.像歪斜的牆,將倒的壁.

係呀,單單要從他的高位上,他們商議,將他拉下來.他們的樂趣在於謊話.他們口講祝福的話,內心卻是咒詛.

係呀,單單在神裡面,要平靜啊,我的生命啊,因為從他而來,我的盼望.

係呀,單單他是我的磐石,我的拯救,我的高臺,我不會被動搖.

以神為基礎,我的拯救和我的榮耀.我力量的磐石,我的避難所,在神裡面.

當時時倚靠他,你們眾人.在他面前,將心傾倒出來.神是我們的避難所.

係呀,單單只是一息空氣,這就是世人.一句謊話,這就是人類.放在天平裡,所有人加起來,比一息空氣還輕.

不要倚靠欺壓,不要對掠奪懷有幻想.就算這樣能加增財富,你的心不要想著它.

神說了一件事,我聽到兩件:能力屬乎神.

你,主啊,就是慈愛,因為你照著人的行為報應他.

這詩篇很簡單,我將之分成兩部分.第一部分是1至7節,詩人在表達信心,但在表達信心的文字中間,卻是人生的凶險景象.第二部分是8至12節,詩人向世人發出呼聲,呼籲世人認清生命的「輕」和「重」,回歸正途.

係呀,單單在神裡面,我的生命,平靜下來.從他而來,我的拯救.

係呀,單單他是我的磐石,我的拯救,我的高臺,我不會被大大動搖.

這兩句經文都是以 ʼak 作開首,有學者認為, ʼak 這個用字,有力排眾議的用意.當人人都抱有一種看法時, ʼak 這個用字,就是表達另一種看法,一種被認為是更真實的看法.

當世人在不同的地方尋找生命的平靜時,當世人不認為神是生命的磐石時,當世人離棄神而尋找不同的拯救時,詩人站起來,力排眾議.他說:「係呀,單單在神裡面,我的生命,平靜下來.從他而來,我的拯救.係呀,單單他是我的磐石,我的拯救,我的高臺,我不會被大大動搖.」他要指出,世人的想法是錯的,真正的真相是:唯有神才是我們生命平靜的根源,唯有神是讓生命穩妥的磐石.

當他表達出唯有在神那裡才能找到生命的平靜和穩妥時,他同時便指出世人的錯誤想法所造成的不平靜和不穩妥.

要到幾時呢?你們攻擊一個人,要殺死他,你們每一個人.像歪斜的牆,將倒的壁.

係呀,單單要從他的高位上,他們商議,將他拉下來.他們的樂趣在於謊話.他們口講祝福的話,內心卻是咒詛.

單單讀經文,便感覺到生命的不平靜和不穩妥.在這裡,安穩只是表面的,背後是殺機.祝福只是表面的,背後是咒詛.世人的樂趣,竟然是空無內容的謊話.

第3節寫「像歪斜的牆,將倒的壁」,究竟是指想殺人的人,像「歪斜的牆,將倒的壁」那樣危險,隨時壓死人；抑或是指那位受害者的生命,像「歪斜的牆,將倒的壁」那樣搖搖欲墜,我們不能確定.詩歌體的文字,只留下意象,讓人自行想像.但這文字的形象很鮮明,離開神的人生和世界,「危危乎」,搖搖欲墜.

在剛過去的聖誕假期裡,跟我任教的神學院的師生去了一趟聖地.

在約旦,有個「皮特拉」(Petra)古城.這古城的特色是,在懸崖上,有很多住人的山洞,也有很多防禦敵人或攻擊敵人的制高點.在以色列,有一個在懸崖上的堡壘馬撒大(Masada),依靠天險,這堡壘可謂固若金湯.詩篇第4節講到的「高位」,可能就是指這些在高點上的,易於防守,也易於攻擊的戰略位置,但是,無論這些位置看來如何穩妥,在凶險的世人手裡,也只是如「歪斜的牆,將倒的壁」那樣「危危乎」.所以第4節說:「係呀,單單要從他的高位上,他們商議,將他拉下來.」世人以為穩妥之地,事實是毫不穩妥的.

這詩篇的第1至2節,是客觀地描寫信仰的事實,人在神那裡可以找到平靜和穩妥的保護.第3至4節,是客觀地描寫人間的凶險,世人依靠的保障,是毫不安穩的.面對人間的凶險,信神的人也可能會軟弱,必須重新堅定自己倚靠神的心.第5-6節的經文,讀起來好像是重複第1至2節,但5-6節不是客觀事實的描寫,而是主觀地喚醒自己,叫自己堅定地進入那信仰的客觀事實裡.留意第5節講的「要平靜啊」,用的是「祈使語氣」.

係呀,單單在神裡面,要平靜啊,我的生命啊,因為從他而來,我的盼望.

係呀,單單他是我的磐石,我的拯救,我的高臺,我不會被動搖.

若果你明白「皮特拉」或「馬撒大」等位於高點的堡壘對於保護當時人的生命的重要性,你就會明白神作為我們的磐石,高臺的拯救作用.「皮特拉」或「馬撒大」,最終都被羅馬人攻破.真正的磐石,高臺,拯救,不是「皮特拉」或「馬撒大」,而是神.「係呀,單單只有神,是我們的磐石,我們的拯救,我們的高臺」.

第7節是以上一段文字的結論:

以神為基礎,我的拯救和我的榮耀.我力量的磐石,我的避難所,在神裡面.

面對人生的不平安和不穩妥,詩人不單呼喚自己回歸神,也呼喚世人回歸神.這開展了這詩歌的第二段文字.

當時時倚靠他,你們眾人.在他面前,將心傾倒出來.神是我們的避難所.

經文的對象,轉換為「你們眾人」.詩人呼喚世人回歸神.要在神面前,將心傾倒出來.這是說,在神面前,要裡外一致的以真面目示人.不要再虛虛偽偽,不要再樂於說謊,不要再口裡祝福而心裡咒詛.並且,不要再為自己找種種東西去安頓自己,要以神為安頓生命的避難所.

ʼak 這個字又出現了,詩人想世人認清人生的真相.

係呀,單單只是一息空氣,這就是世人.一句謊話,這就是人類.放在天平裡,所有人加起來,比一息空氣還輕.

不要倚靠欺壓,不要對掠奪懷有幻想.就算這樣能加增財富,你的心不要想著它.

世人以為自己很利害,很重要,很有重量.但實情是如何的?詩人說:「世人如煙,係呀,世人只是一息空氣,也如一句謊話,裡面毫無內容.」如空氣一樣輕飄飄的人類,全部加起來,也不比一息空氣重.

人生這般輕,是因為人的心記掛一些毫無重量的事情.「不要倚靠欺壓,不要對掠奪懷有幻想」.「幻想」這個字,在希伯來文裡,與「空氣」一字有相同字根(hbl).什麼是「幻想」,就是如空氣般不實在的想法.人以虛偽的心,虛假的話,去賺取不義之財,以為這樣能令自己穩妥.但從來用這方法令自己穩妥的人,都心知肚明,這只帶來恐懼與不安,從不穩妥.

神說了一件事,我聽到兩件:能力屬乎神.

「神說了一件事,我聽到兩件」,有點似廣東話的「話頭醒尾」,這是希伯來文的強調手法.神講什麼呢?「能力屬乎神」.「能力」這個字,有「堅固」的含意,也指向「安穩」.「能力」當然與「力」有關,從物理學而言,「重量」就是「力」.「能力屬乎神」,就是說,神是「堅」的,有重量的.在他裡面,找到真正的「重量」,也找到安穩.

這詩篇的結語是:

你,主啊,就是慈愛,因為你照著人的行為報應他.

神就是愛.他是人間的愛者,向世人呼喚.當愛者遇上愛者,便有關係,責任,十架.有了關係,責任,十架,重量便產生了,力量便產生了,意義便產生了.愛人者,揭露了世人的有情抑或無情,暴露了世人的重與輕.

神判斷人生命的輕與重.輕如空氣的人,就讓他離地飄起,消失於虛無飄渺間,不留下半點歷史痕跡.在神那裡找到重量的人,就讓他立足大地,在永恆裡留下努力活過的記號.

詩歌如茶,讓人細味.這詩篇讓人體會生命的輕與重.為何生命會輕得毫無重量?為何一句關懷的話說得這般假?為何交一個朋友時只想到他的利用價值?為何處理事情時只想到如何可以不負責任?詩人希望我們多想一想:為何生命會這樣輕?

這詩篇形容神的用字,是有重量的.神是磐石,是高臺,是「能力」.這一切都是「堅」的,有內容的,有重量的.最後一節經文,講到神是愛,這「愛」也是重得驚人的.這是守約的愛,不離不棄的愛,是犧牲自己而成全他人的愛.在神那裡,人才能找到重量.有了重量,才能安穩.

斗轉星移,浪花淘盡英雄.人與事,如輕煙飄去.但是,在我們心底裡,是否仍渴望某種重量?

堅石與空氣,重與輕,永恆與短暫,兩者的張力會常常出現在我們生命中.無論世人如何看,詩人卻說:「係呀,單單神是我的磐石,我的拯救,我的高臺,我不會被動搖.」

但願:榮耀歸於聖父,聖子,聖靈.阿們.



\newpage

\section{權柄與權力}
\label{sec:275}
\textbf{伍渭文牧師 - 權柄與權力}
\newline
\newline
\hyperref[sec:274]{< < < PREV SERMON < < <}
~
\hyperref[sec:toc]{[返目錄]}
~
\hyperref[sec:276]{> > > NEXT SERMON > > >}


馬可福音 1:21-28
\newline
\begin{longtable}{cl}
\hline
\hline
章節 & 經文 (和合本修訂版)\\
\hline
1:21 & \begin{tabularx}{0.7\textwidth}{X} 他們到了迦百農，耶穌就在安息日進了會堂教導人。 \end{tabularx} \\ \\ \relax
1:22 & \begin{tabularx}{0.7\textwidth}{X} 他們對他的教導感到很驚奇，因為他教導他們正像有權柄的人，不像文士。 \end{tabularx} \\ \\ \relax
1:23 & \begin{tabularx}{0.7\textwidth}{X} 當時，會堂裡有一個污靈附身的人，他在喊叫， \end{tabularx} \\ \\ \relax
1:24 & \begin{tabularx}{0.7\textwidth}{X} 說：「拿撒勒人耶穌，你為甚麼干擾我們？你來消滅我們嗎？我知道你是誰，你是神的聖者。」 \end{tabularx} \\ \\ \relax
1:25 & \begin{tabularx}{0.7\textwidth}{X} 耶穌斥責他說：「不要作聲，從這人身上出來吧！」 \end{tabularx} \\ \\ \relax
1:26 & \begin{tabularx}{0.7\textwidth}{X} 污靈使那人抽了一陣風，大聲喊叫，就出來了。 \end{tabularx} \\ \\ \relax
1:27 & \begin{tabularx}{0.7\textwidth}{X} 眾人都驚訝，以致彼此對問：「這是甚麼事？是個新的教導啊！他用權柄命令污靈，連污靈也聽從了他。」 \end{tabularx} \\ \\ \relax
1:28 & \begin{tabularx}{0.7\textwidth}{X} 於是耶穌的名聲立刻傳遍了全加利利周圍地區。 \end{tabularx} \\ \\
[1ex]
\hline
\hline
\end{longtable}
講題:權柄與權力　Authority and Power

經文:馬可福音1:21-28;詩篇110

講員:伍渭文牧師

場合:香港中文大學崇基學院禮拜堂

日期:2015年2月1日

歷時七十多天的佔領運動過去了,香港跟以前不一樣了.佔領者本是違法,政府多次勸籲佔領者都不肯離開,但政府又不直接執法,祇等待民事訴訟,借用機會,協助執達吏執法,乘機清場.有人因此說,政府已失去應有的權威或權柄Authority, 現在祇憑權力Power統治.   權柄與權力有何分別呢?權柄本身有一樣吸引力,使人樂於順從；權力是一種外在的制約,不得不遵守.當法國遺派核子動力航母戴高樂號前往以拉克對付以斯蘭國,作為對巴黎漫畫雜誌社恐怖襲擊的回應,這是權力威攝-軍事權力的威攝.但教宗前一段時間在菲律賓訪問,民眾夾道歡迎,熱烈自發的親炙他,民眾是被一種屬靈的權柄所吸引.印象中年輕人一般是反權威的,其實不是反權威,乃反對權威背後的虛偽；每次天主教的世界青年大會,我觀察到來自世界各地包括北美,歐洲的年輕人,對教宗一樣熱情,對屬靈的事情充滿熱枕,說明了人不是不想獻祭,乃在渴望尋找可獻祭的祭壇.   今天福音經課就提到耶穌講道帶有一種權柄,不像文士,使聽道的人非常驚訝.「眾人很希奇他的教訓；因為他教訓他們,正像有權柄的人,不像文士.」(可一22,23)聖經權柄一字(exousian) 可以翻作能力,權柄帶有能力,有一種權能.耶穌講道和文士不同,使聽的人感到驚訝；而且耶穌運用這權能驅逐污鬼.財富可以購買很多東西,權位也可以帶來不少影響力,但財富和權位不能帶給人尊重和生命的吸引力.   報載鄰近地區的一個豪華旅行團在歐洲旅行,用餐時叫滿一了滿檯食物, 結果吃不完,非常浪費,結賬時又沒有絲毫歉意.餐廳侍應提議下次可少叫一些,免得浪費,但惹來豪客不滿,輕蔑地說:「錢是我們的,我們花得起,你不用管!」侍應不亢不卑說:「對,錢是你們的,但這些是我們的出產,我們不想見到我們的資源被浪費.」   耶穌出來傳道,就帶著不一樣的權柄,使人驚訝,吸引著聽道的人,這樣柄權是在非常卑微的迦伯農會堂開始的.   1   權柄出卑微 (可一 21)   「到了迦百農,耶穌就在安息日進了會堂教訓人.」(可一21)短短的一句,馬可表達了耶穌的主動.之前,耶穌到約但河接受施洗約翰的洗禮；行到海邊,呼召打魚為生的西門,安得烈,雅各,約翰為門徒；進入迦百農的會堂,在安息日教訓人.他總是去到我們中間,接觸我們,服侍我們.迦百農的會堂,人數應該不會很多,參加的人,背景也不會顯赫,但耶穌的偉大事工,就以卑微的迦百農會堂開始.   「耶穌就在安息日」指出耶穌順著當時的習慣,規章和限制來傳道,彰顯福音的能力.我記起唸神學的時候,當時的院長滕近輝牧師講了一句非常受用的話:教會任何不完善的制度,也不會防礙傳道的工作.任何的一個制度總有改善的空間,但滕牧師的意思叫我們在諸般限制中做好自己,不要逃避困難,諉過環境而裹足不前,磋跎歲月.   甚為可惜,現在不少年輕傳道人因為教會人事張力,傳統的制肘,頻頻轉換工場,奮然走上十字架了結自己,希望在另一個新工場得到復活.頻頻轉換工場,對自己和教會都不是好事.事實上,困難的磨練,使人成長.   在洛杉磯探親時有機會和當地一位牧師午餐,他是印尼華僑,在洛杉磯牧會四十年同一間教會,現在教會主日普通話上千人聚會,英語也六百多人,廣東話二百多人,接近二千人聚會的教會都是從很卑微的開荒工作一步一步建立起來.四十年在同一間教會忠誠的服務,這牧者贏取了會友的愛戴,他的說話也帶有令會友聽從的權柄.   2   權柄來自順服 (可一22, 太二十三2,3)   文士能說不能行.「眾人很希奇他的教訓；因為他教訓他們,正像有權柄的人,不像文士.」(可一22)究竟文士講解聖經和耶穌有何不同呢?同樣是講解聖經,為何文士的教訓沒有權柄呢?原因是他們沒有遵行聖經.耶穌說過:「文士和法利賽人坐在摩西的位上,凡他們所吩咐你們的,你們都要謹守遵行；但不要效法他們的行為；因為他們能說,不能行.」(太二十三2,3)   當新生王要降生時,希律王非常不安,深受威脅,詢問文士:「基督當生在何處?」(太二4)他們回答說:「在猶太的伯利恆.最為有先知記著,說:猶太地的伯利恆啊,你在猶大諸城中並不是最小的；因為將來有一位君王要從你那裡出來,牧養我以色列民.」(太二6)但他們並沒有像東方千里迢迢的賢士尋覓基督.   文士熟悉聖經,但失去對真理的催迫感,也不遵行聖經.所以耶穌說,聽他們的講論但不要效法他們的行為.當時的祭司長和文士擁有解釋聖經的權力(power) ,隨著這權力也得到不少利益,成為耶路撒冷聖殿城的管治階層；他們的解釋也可以作為官方的標準.但這些權力和利益使他們的屬靈觸角慢慢麻木起來,最後失去了使人心歸附的權柄.   生命虔敬的能力.過去在美國洛杉磯探親期間,有機會在張伯笠牧師教會講道,看到在教會入口處,張牧師為來訪新朋友購買自己寫的書「在逃者」簽名.訪客多來自中國大陸,他們大多數聽過張牧師的見證.「在逃者」記述作為八九民運領袖的張伯笠,在中俄邊疆逃避通緝的過程中,得到幾位農村家庭教會的姊妹的幫助和分享福音而信主的.她們讀書不多,但信仰堅定和熟讀聖經,也遵守聖經教訓生活.她們有一種特別的氣質,吸引著這位在逃者,這是在天安門廣場的學生領袖所沒有的.這些學生領袖很有理想,也有學識,組織力強,詞鋒犀利,辯才無礙,但就沒有這幾位姊妹的生命虔敬和愛心,張伯笠在她們身上,窺探到一個真實的屬靈世界.   無論我們身處於人看來是卑微或限制的處境,祇要我們聆聽上主的話語,躬身力行,就有一種令人響往的屬靈權柄,吸引人歸向基督.

挑戰邪惡力量的權柄 vv. 23-26

挑戰邪惡的力量.「在會堂裡有一個人被污鬼附著.他喊叫說:拿撒勒人耶穌,我們與你有甚麼相干?你來滅我們嗎?我知道你是誰,乃是神的聖者.」(可一23-24)我們不知道這被污鬼附身的人,會否平時聚會也不斷喊叫,抑或這人平時很安靜,一如常人,因為文士的教導沒有權柄和能力,所以大家相安無事.   外在的綑綁.但耶穌的教導有權能,這權能挑釁潛伏的邪惡力量.這事使我想起數年前天主教某敢言神甫,嚴厲批評某地產商的商業霸權無孔不入,包攬所有衣,食,住,行；而且分毫必賺,點滴不漏,與聖經規定收割麥穗,不可拾遺,留給窮困者充飢的教訓相去甚遠,儼然是魔鬼附身的行為,耶穌先祖路得曾生活窮困,因為拾麥穗才存活下來.(路得記二)   耶穌驅走污鬼:「耶穌責備他說:不要作聲!從這人身上出來吧.」(可一25)污鬼喊叫,因他知道,耶穌會驅趕牠,他認識耶穌是上帝的聖者.聖邪豈能兩立呢?我們也要效法耶穌,直視並驅走在這社會壓制人的邪魔鬼道,恢復人的尊嚴和人性.耶穌驅走污鬼,使眾人驚訝:「這是甚麼事?是個新道理啊!他用權柄吩咐污鬼,連污鬼也聽從了他.」(可一27)   內在的綑綁.「他用權柄吩咐污鬼,連污鬼也聽從了他」這話給我們很大的鼓勵,無論綑綁我們的力量有多大,多次嘗試都不能擺脫,這些綑綁甚至使我們失去自信,使我們絕望.今天的經文提醒我們,但祇要我們對上主的話有信心,祇要我們願意順服上主的話,我們就可以得著釋放,我們也可以經歷上主的權能,驅走隱伏內裡的心魔,淨化內心的邪情私慾,經歷心靈的釋放.   順服是經驗上主話語的權柄和權能的必需條件.文士懂說但不行,講道沒有權柄；我們若懂聽而不行,也不能經驗上主話語的權柄.

上主的權柄在人心中,詩110: 2, 3

上主的國度.耶穌驅逐污鬼是標誌性事件,表示上主的國度已來臨了,今天詩篇經課詩篇一一0是新約聖經引用最多的詩篇；直接引用七次,間接十五次的彌賽亞詩篇,這詩篇說到國度和權力的運用.詩篇以大衛的登基,預表大衛之後耶穌作為彌賽亞掌權登基.大衛登基,執掌王權,管治國民,但他祇是代替上主,君王有崩逝,政黨有輪替,朝代有興衰,但上主永遠坐著為王:「我主(上帝)對我主(大衛)說,你坐在我的右邊,等我使你仇敵作你的腳凳.」(v.1)昔日耶路撒冷由麥基洗德統治,今天由大衛統治,他們都是替上主管治,君王也要被上帝管治,才能秉公義施仁政.   建立他人的權力.麥基洗德是善用上主托付權力的好榜樣,他以公義治理耶路撒冷,麥基洗德就是公義的王的意思.他愛好和平,所以又稱為撒冷王,撒冷salem就是和平salom.他守住先祖傳下來的信仰:女人的後裔要傷蛇的頭,等待上主應許的救贖,帶領子民敬畏上主,他是一國的祭司；集君王,祭司於一身,善用上帝給他的權力.非常可惜,同時代鄰近的所多瑪和俄摩拉王,沒有善用上主給予的權力,把百姓陷在邪淫的罪中,惹來上主的忿怒和審判；慾火焚身,天火焚城.   創世記十四章記載,當亞伯蘭戰勝東方四王,救回侄兒羅得經過耶路撒冷時,麥基洗德以是至高上主的祭司的身份,帶著酒和餅出來迎接亞伯蘭,為他祝福；亞伯蘭把奪取的戰利品十分一獻給他.「你是照著麥基洗德的等次永遠為祭司」(4節下)是預表基督,希伯可書說:「他無父,無母,無族譜,無生之始,無命之終,乃是與上主的兒子相似.」(來7: 3)一位無顯赫背景,甚至無父無母,身世坎坷的人,因為善用權力,勤於職責,篤信力行,. 蒙上主揀選為基督的預表.   被愛約制的權力.這詩篇也給我們窺探到,上主最終極和絕對的權力不是唯我獨尊的極權(autocratic),是有愛的交流和契通:「我主(上帝)對我主(基督)說」.在馬可福音書,耶穌曾公開問猶太人:「文士怎麼說基督是大衛的子孫呢?大衛被聖靈感動,說:主(上主)對我主(彌賽亞)說:你坐在我的右邊,等我仇敵作你的腳凳.大衛既自己稱他為主,他怎麼又是大衛的子孫呢?」耶穌把上主對(我主)大衛的話,轉為上主對彌賽亞自己的話.   在這上帝本體三一神的對話中,我們看到至高至終極的權力,是被愛和團契約制的.基督信仰挑戰極權,因為極權沒有他人,祇有自我；沒有契通,祇有獨斷,所以權力很容易成為壓制(overpower) 他人,不是建立(empower) 他人.   隱藏的國度.詩篇談到這國度是隱藏的:「耶和華必使你從錫安伸出能力的杖來；你要在你仇敵中掌權.」(v. 3)驟眼看來,四周還是仇敵盤踞.今年波蘭奧威辛解放七十年紀念,不少猶太人,特別是教導的拉比目睹同胞一車一車運往集中營,甚至知道是毒氣室時都問:上主,你在那裡?教宗訪能菲律賓期間,一位十三歲女童問他:為何現在還有像她那麼年輕的女孩子被拐賣當雛妓呢?上主為何不理會呢?但詩篇告訴我們,從信仰的角度來看,他們已成為上帝的腳凳(v. 1),這位彌賽亞君王正抬起頭來(v.7),昂首前進.詩篇用了首尾相應的文學手法,說明上帝權柄永續,他的國愈來愈壯大.當他怒的日子,他打傷列王,他要在列邦中刑罰惡人,屍首就遍滿各處,他要在許多國中打破仇敵的頭.(vv. 5-6)這是天父世界,罪惡雖然好像得勝,天父卻仍掌管.這是天父世界,我心不必憂傷,上帝是主是王,天地同唱,歌聲充滿萬方.   上帝的國度是隱藏的,但透過他子民的生命顯為真實.當教會在地上為不公義發聲,為所不公義所壓傷的人裹傷時,上主的國度就顯為真實.   歸附的子民.要成為一個國家或國度要有三個元素:國土,國民和主權.佔領事件發生,令一位前任負責香港事務的中央官員感嘆:九七領土回歸了,但人心到今天尚未回歸.彌賽亞的國度有主權,也有子民,但沒有看得見的國土.但沒有國土怎能成為一個國家呢?其實上主的國度是有國土的,這國土就是國民的心.當彌賽亞掌權時,他的子民就誠心歸附:「當你掌權的日子,要以聖潔的妝飾為衣,甘願犧牲自己；你的民多如清晨的甘露.」(2節)一日晨曦之始,一生青春之端 (Join you at the fresh break of day, join you with all vigor of youth), 上主的子民,獻上自己,欣然歸附.   潤物細無聲.「清晨的甘露」使我們想起晨曦前夕的黑夜.在黑夜裡,植物停止光合作用,大地張開雙手,放鬆自己,靜靜接受空氣的覆蓋和水氣的滋潤,杜甫詩《春夜喜雨》:「隨風潛入夜,潤物細無聲.」我們也要回歸平靜安穩,學曉默默無聲,等候上主.   最近看到一位紐約報專欄作家極力推薦的書:My Bright Abyss-Meditation of A Modern Believer(我的光明深淵-一個現代信徒的思考),作者Christian Wiman為美國期刊「詩社」(Poetry) 的主編.作為詩人,他很少和讀者互動,因為詩人重視聆聽來自某處的聲音(somewhere else) ,要一首詩成為詩,你要有對詩的來源有一種修士般的虔敬(monkish devotion) ,沒有信仰的詩人也是這樣.但七年多前他生命陷入嚴重的低潮,不能寫作,墮入愛河,但同時獲知得到不治的頑疾,生命感到被撕開,他今年才四十多哩.他開始返回上帝旁邊,認真思考和上帝的關系,寫了七頁短文,投稿到一間很小的雜誌發表,結果引起很大的回響,他開始認真藉文字和人溝通.   他發現其實現代人很想了解信仰,但厭煩套話浮淺；我們現代人說話太多太快,思考默想太少太短.他寫作非常認真,作為詩人,他聆聽每一個字的音節,輕重和色澤,並等候聆聽某處的聲音才動筆.漫長的黑夜,水氣靜靜的濕潤,才孵育出清晨的甘露；我心默默無聲,才進到上主的真實.   今天經文談到一種權柄,吸引著人使人樂於歸附的權柄,這權柄可以改變人的生命.無論在多麼卑微的處境,祇要我們聆聽聖經,順服聖經,我們也可以得到這種權柄,經歷生命的改變,也要被上主使用,令世界變好一些.



\newpage

\section{往聖城路上的城與鄉}
\label{sec:276}
\textbf{梁元生教授 - 往聖城路上的城與鄉}
\newline
\newline
\hyperref[sec:275]{< < < PREV SERMON < < <}
~
\hyperref[sec:toc]{[返目錄]}
~
\hyperref[sec:277]{> > > NEXT SERMON > > >}


路加福音 9:21-56
\newline
\begin{longtable}{cl}
\hline
\hline
章節 & 經文 (和合本修訂版)\\
\hline
9:21 & \begin{tabularx}{0.7\textwidth}{X} 耶穌切切吩咐他們，命令他們不可把這事告訴任何人； \end{tabularx} \\ \\ \relax
9:22 & \begin{tabularx}{0.7\textwidth}{X} 又說：「人子必須受許多的苦，被長老、祭司長和文士棄絕，並且被殺，第三天復活。」 \end{tabularx} \\ \\ \relax
9:23 & \begin{tabularx}{0.7\textwidth}{X} 耶穌又對眾人說：「若有人要跟從我，就當捨己，天天背起自己的十字架來跟從我。 \end{tabularx} \\ \\ \relax
9:24 & \begin{tabularx}{0.7\textwidth}{X} 因為凡要救自己生命的，必喪失生命；凡為我喪失生命的，他必救自己的生命。 \end{tabularx} \\ \\ \relax
9:25 & \begin{tabularx}{0.7\textwidth}{X} 人就是賺得全世界，卻喪失了自己，或賠上自己，有甚麼益處呢？ \end{tabularx} \\ \\ \relax
9:26 & \begin{tabularx}{0.7\textwidth}{X} 凡把我和我的道當作可恥的，人子在自己的榮耀裡，和天父與聖天使的榮耀裡來臨的時候，也要把那人當作可恥的。 \end{tabularx} \\ \\ \relax
9:27 & \begin{tabularx}{0.7\textwidth}{X} 我實在告訴你們，站在這裡的，有人在沒經歷死亡以前，必定看見神的國。」 \end{tabularx} \\ \\ \relax
9:28 & \begin{tabularx}{0.7\textwidth}{X} 說了這些話以後約有八天，耶穌帶著彼得、約翰、雅各上山去禱告。 \end{tabularx} \\ \\ \relax
9:29 & \begin{tabularx}{0.7\textwidth}{X} 正禱告的時候，他的面貌改變了，衣服潔白放光。 \end{tabularx} \\ \\ \relax
9:30 & \begin{tabularx}{0.7\textwidth}{X} 忽然有摩西和以利亞兩個人同耶穌說話； \end{tabularx} \\ \\ \relax
9:31 & \begin{tabularx}{0.7\textwidth}{X} 他們在榮光裡顯現，談論耶穌去世的事，就是他在耶路撒冷將要完成的事。 \end{tabularx} \\ \\ \relax
9:32 & \begin{tabularx}{0.7\textwidth}{X} 彼得和他的同伴都打盹，但一清醒，就看見耶穌的榮光和與他一起站著的那兩個人。 \end{tabularx} \\ \\ \relax
9:33 & \begin{tabularx}{0.7\textwidth}{X} 二人正要和耶穌分離的時候，彼得對耶穌說：「老師，我們在這裡真好！我們來搭三座棚，一座為你，一座為摩西，一座為以利亞。」他卻不知道自己在說些甚麼。 \end{tabularx} \\ \\ \relax
9:34 & \begin{tabularx}{0.7\textwidth}{X} 說這些話的時候，有一朵雲彩來遮蓋他們；他們一進入雲彩就很懼怕。 \end{tabularx} \\ \\ \relax
9:35 & \begin{tabularx}{0.7\textwidth}{X} 有聲音從雲彩裡出來，說：「這是我的兒子，我所揀選的。你們要聽從他！」 \end{tabularx} \\ \\ \relax
9:36 & \begin{tabularx}{0.7\textwidth}{X} 聲音停止後，只見耶穌獨自一人。當那些日子，門徒保持沉默，不把所看見的事告訴任何人。 \end{tabularx} \\ \\ \relax
9:37 & \begin{tabularx}{0.7\textwidth}{X} 第二天，他們下了山，有一大群人來迎見耶穌。 \end{tabularx} \\ \\ \relax
9:38 & \begin{tabularx}{0.7\textwidth}{X} 其中有一人喊著說：「老師！求你看看我的兒子，因為他是我的獨子。 \end{tabularx} \\ \\ \relax
9:39 & \begin{tabularx}{0.7\textwidth}{X} 他被靈拿住就突然喊叫，那靈又使他抽風，口吐白沫，並且重重地傷害他，不輕易放過他。 \end{tabularx} \\ \\ \relax
9:40 & \begin{tabularx}{0.7\textwidth}{X} 我求過你的門徒把那靈趕出去，他們卻不能。」 \end{tabularx} \\ \\ \relax
9:41 & \begin{tabularx}{0.7\textwidth}{X} 耶穌回答：「唉！這又不信又悖謬的世代啊，我和你們在一起，忍耐你們，要到幾時呢？把你的兒子帶到這裡來！」 \end{tabularx} \\ \\ \relax
9:42 & \begin{tabularx}{0.7\textwidth}{X} 他正來的時候，那鬼把他摔倒，使他重重地抽風。耶穌斥責那污靈，把孩子治好了，交給他父親。 \end{tabularx} \\ \\ \relax
9:43 & \begin{tabularx}{0.7\textwidth}{X} 眾人都詫異神的大能。 \end{tabularx} \\ \\ \relax
9:44 & \begin{tabularx}{0.7\textwidth}{X} 「你們要把這些話聽進去，因為人子將要被交在人手裡。」 \end{tabularx} \\ \\ \relax
9:45 & \begin{tabularx}{0.7\textwidth}{X} 門徒卻不明白這話，其中的意思對他們隱藏著，使他們不能明白，他們也不敢問這話的意思。 \end{tabularx} \\ \\ \relax
9:46 & \begin{tabularx}{0.7\textwidth}{X} 門徒互相議論，他們中間誰最大。 \end{tabularx} \\ \\ \relax
9:47 & \begin{tabularx}{0.7\textwidth}{X} 耶穌看出他們心中的議論，就領一個小孩子來，叫他站在自己旁邊， \end{tabularx} \\ \\ \relax
9:48 & \begin{tabularx}{0.7\textwidth}{X} 對他們說：「凡為我的名接納這小孩子的，就是接納我；凡接納我的，就是接納那差我來的。你們中間最小的，他就是最大的。」 \end{tabularx} \\ \\ \relax
9:49 & \begin{tabularx}{0.7\textwidth}{X} 約翰回應說：「老師，我們看見一個人奉你的名趕鬼，我們就阻止他，因為他不與我們一同跟從你。」 \end{tabularx} \\ \\ \relax
9:50 & \begin{tabularx}{0.7\textwidth}{X} 耶穌對他說：「不要阻止他，因為不抵擋你們的，就是幫助你們的。」 \end{tabularx} \\ \\ \relax
9:51 & \begin{tabularx}{0.7\textwidth}{X} 耶穌被接上升的日子將到，他決定面向耶路撒冷走去。 \end{tabularx} \\ \\ \relax
9:52 & \begin{tabularx}{0.7\textwidth}{X} 他打發使者在他前頭走；他們進了撒瑪利亞的一個村莊，要為他作準備。 \end{tabularx} \\ \\ \relax
9:53 & \begin{tabularx}{0.7\textwidth}{X} 那裡的人不接待他，因為他面向著耶路撒冷去。 \end{tabularx} \\ \\ \relax
9:54 & \begin{tabularx}{0.7\textwidth}{X} 他的門徒雅各和約翰看見了，就說：「主啊！你要我們吩咐火從天上降下來，燒滅他們嗎？」 \end{tabularx} \\ \\ \relax
9:55 & \begin{tabularx}{0.7\textwidth}{X} 耶穌轉身責備兩個門徒。 \end{tabularx} \\ \\ \relax
9:56 & \begin{tabularx}{0.7\textwidth}{X} 於是他們就往別的村莊去了。 \end{tabularx} \\ \\
[1ex]
\hline
\hline
\end{longtable}
講題:“往聖城路上的城與鄉:不接待主的鄉村, 不悔改的城市”　Villages of Non-hospitality and Cities of Non-repentance: On the Road to the Holy City

經文:路加福音9:21-56;10:13-16

講員:梁元生教授

場合:香港中文大學崇基學院禮拜堂

日期:2015年2月8日

前言:城與鄉

不接待的村莊:撒瑪利亞

不悔改的城市:哥拉汛,伯賽大,迦百農,西頓,泰爾

結語；朝向聖城之路

前言:城與鄉

各位弟兄姊妹,各位朋友:早安.

首先,我要感謝上主再一次給我機會在這裡和大家分享,同時,我也在此多謝校牧許牧師的邀請和伍牧師的調換時間.本來許牧師要我在上個主日講道,但由於大學的退思會由上星期六開到星期日中午才結束,所以要求調期,多謝伍牧師答允.在此再次謝謝.

今天我要和大家思想的聖經,記載在新約路加福音第九章和第十章,是關於城市和鄉村的.

我想先說明一點:我是學歷史的,而專門研究的範圍是城市史和移民史.所以每當我讀聖經時,總會對經文中提到的城市及地區之間人口的遷移和流動,格外留意.在這個講臺上, 我記得也曾和大家講過幾聖經城市的故事,包括希臘大城雅典,馬其頓的邊城腓立比等.今天我和大家分享的一段聖經, 也環繞著城市和鄉村,以及從此地區到彼地區移動的主題.現在讓我們來讀聖經中路加福音第九章的最後一段(九:21-56 節)這裡提到一條村莊:一條坐落於撒瑪利亞地區的村莊,一條不接待主耶穌和祂門徒的村莊； 之後,我們跳開幾段再讀下去,來到路加福音第十章12-15節,那裡提到好幾個城市的名字, 它們是:哥拉汛,伯賽大,迦伯農, 以及推羅和西頓,前面三個個是新約時期亦即是耶穌當年常經過而且十分熟悉的城市,後兩個是有千年歷史的腓尼基古城, 以商業聞名, 主前三百年已經衰頹,不復當年舊貌. 在這段經文中,當耶穌提及這些城市時,祂用了聖經中舊約的一座充滿罪惡的大城來比附,就是所多瑪.(十:12-15)這些城市,用耶穌自己的話來說, 是“不悔改”的城市.所以,我今天講道的題目雖然叫做“城與鄉”,但更貼切的題目應該是:“不接待主的鄉村, 不悔改的城市.”而遇到這兩個“不”--不接待的鄉村和不悔改的城市--的經驗,發生在門徒跟隨耶穌前往聖城耶路撒冷的路途上. 我想在這裡事先強調一點,這樣的經驗,同樣會發生在每個信徒的人生“旅途”上,包括你和我的人生.

換言之, 在我們的人生旅途或信仰旅途上, 都會有經過鄉村和走過城市的經驗, 而這些經驗又並不都是完全愉快的, 更多時候,這些經驗是讓我們沮喪和使我們感到失望和困惑的.從上面的經文總結這種經驗, 可以用兩個“不”字來形容, 就是“不接待”和“不悔改”( 英文我就用兩個Non 字, 即 ‘Non-hospitality and Non-repentance’),不接待的村莊,就是village of non-hospitality ,不悔改的城市,我就翻譯作 cities or towns of non-repentance .

不接待的村莊 : 撒馬利亞

讓我們先來看那個「不接待主的村莊」.這裡指的是撒馬利亞. 路加九章21節說:「耶穌被接上升的日子將到,他就定意向耶路撒冷去,便打發使者在他前頭走.他們到了撒瑪利亞的一個村莊,要為他預備.那裡的人不接待他,因他面向耶路撒冷去.他的門徒雅各約翰看見了,就說:“主啊,你要我們吩咐火從天上降下來,燒滅他們,像以利亞所作的嗎?“ 耶穌轉身責備兩個門徒, 說:“你們的心如何, 你們並不知道.人子來不是要滅人的性命,是要救人的性命.”說著 就往別的村莊去了.」

撒馬利亞這條村裡的人,不要接待主耶穌和他的門徒.首先, 讓我們來看看聖經地理,好瞭解事情發生的背景.事情發生的當兒, 正值耶穌帶著門徒由加利利去耶路撒冷的途中. 換言之, 耶穌當時正由北向南走,從加利利海向亞拉巴海(即鹽海)的方向走, 目的地是耶路撒冷.由北向南走,也就是從巴勒斯坦向猶太的方向走,沿途必定需要經過撒馬利亞, 那是人煙較少的鄉村地方.這裡記載:撒馬利亞村莊裡的人都不願意接待耶穌和門徒.換言之,耶穌和門徒來到撒馬利亞的村莊, 受到冷待,不受歡迎.撒瑪利亞是一條不接待主,不歡迎主,甚至是拒絕主的村莊.那麼,我們要問: 撒瑪利亞人對主及門徒的拒絕,原因何在?聖經說:是因為他們要向耶路撒冷去.這句簡單的話, 背後所包含的民族間的恩怨情仇,以及文化和地域間的模糊界線,十分複雜,一時間不容易說得清楚.要明白撒馬利亞人之所以不歡迎主耶穌,就要瞭解猶太人和撒馬利亞人的種族關係和歷史恩怨. 據撒馬利亞人說,他們才是亞伯拉罕真正的子孫,是約瑟的後裔,是屬於以法蓮和嗎哪西支派的.到所羅門王死後,兒子羅波安繼位,但群眾叛亂,迦南地區分成南北二國,北國叫以色列,擁戴尼八的兒子耶羅波安為王, 統治地方包括了撒馬利亞全境；南方的叫猶大,奉所羅門王的兒子羅波安為王,以耶路撒冷為首都.兩國分立多年,也時有爭戰.兩國雖然系出同源,但為爭正統而屢屢鬧得不個可收拾.撒馬利亞人說他們敬拜的上帝才是以色列的真神,他們指斥猶太人經歷了後來的巴比倫的擄掠和流徙,民族受到外邦人的沾汙,不再是真正清白的血統；而猶太人則指責撒馬利亞人為古達(古代伊拉克)的移民,辭中帶著貶義,意即“賤民”.再者,猶太人以耶路撒冷為聖城,在那裡敬拜上帝, 但撒馬利亞人則以為基利心山才是聖地, 在那裡建立敬拜上帝的殿宇. 這些因宗教和種族上的不同意見造成了猶太人和撒馬利亞人相互的排擠和歧視,到耶穌基督的時代還是如此.

我不厭其煩地把猶太人和撒瑪利亞人之間的複雜關係說了個梗概,其中最重要的一點是:他們是同種同文,有著共同的祖先,也有著共同的歷史,和共同的傳統,但卻都執著一方的正統性和排擠對方的異質性和混雜性.這種執著與堅持,造成相互的排擠,歧視,甚至敵視和武力鬥爭.不尋求共融,而追求獨大,至尊,這種自以為是,坐井觀天的態度,我叫它做“鄉村的心態”,或“地方主義”(provincialism),重視傳統,強調正統,排擠他者,視野狹隘,是一條極深刻的區分內外的民族,文化和地域的界線.環顧今日,這個全球化的世界仍然佈滿了鄉村的界線! 而在香港這個城市中,鄉村心態,排擠他者,自以為是的撒瑪利亞鄉村的界線和痕跡,仍然是處處可見.但令人栗然的是:在跟隨主耶穌的人群中,也有這樣的“村民”, 門徒中就有人說:主啊, 請用天火把村莊燒掉!.後來. 耶穌用好撒瑪利亞人的比喻, 告訴門徒:祭司和猶太人, 並不是都帶著慈悲心和仁愛心的, 反不如異族的撒瑪利亞人.而且, 他也同情打水的撒瑪利亞婦人, 和醫治撒瑪利亞的麻瘋病患者.

不悔改的城市:哥拉汛,伯賽大,迦百農,泰爾(推羅),西頓

接著下來, 耶穌提醒門徒說, 他們不但在鄉村遭到拒絕,而且在城市中也不受歡迎.他們從加利利海邊出發, 要把福音帶到各個城鎮.雖然有些人信了福音, 也來跟從主； 但許多人都硬著心, 不肯悔改. 耶穌特別點名指出幾個城市,就是:哥拉汛,伯賽大和迦百農.說:「哥拉汛啊,你有禍了.」又說:「伯賽大啊, 你有禍了.」「迦百農啊,你已經升到天上, 將來必推下陰間.」

聖經的編者在這段經文前面, 加上小標題: 不悔改的城市. 背景是耶穌設立七十個人, 並且派遣他們出去, 傳平安的福音.他的一番話就是針對那些不接待門徒和棄絕福音的城市說的, 特別點名的三個不悔改的城市 (cities/towns of non-repentance),就是上面提到的哥拉汛(Chorazin),伯賽大(Bethsaida)和迦百農(Capernaum).和這些不知悔改的城市比較, 耶穌提到另外的幾個城市, 就是舊約時代的所多瑪(Sodom),和當時位於巴勒斯坦地中海沿岸的兩個海港:推羅(Tyre泰爾)和西頓(Sidon).這幾個都是著名的商業大城,是繁榮的海港, 人煙稠密, 商業茂盛.所多瑪的故事我們應是耳熟能詳的了.由於城市中的罪惡滿盈, 乃招天火焚城之殃, 只有羅得倖免於難.至於推羅和西頓, 早在舊約時代已經是著名的港口, 例如舊約提到推羅和西頓, 都盛讚其地交通發達,商業繁盛,物產豐富.毫無疑問,推羅(泰爾)和西頓,是典型的商港,舊約聖經曾經稱它們為“大城”,為“商港”,為“列國的大碼頭”,又因居民但與逸樂,追求享受,故此又叫它們做“歡樂的城”.(賽二十三章) 推羅(泰爾)和西頓這兩個大商港,位於地中海海邊,自舊約時期到主前三百年,一直都是商業繁盛,海舶雲集,人口眾多的城市.(在這樣的商業大城中,人來人往,交通便利,但多數人營營役役,茲茲求利,生活萎靡,燈紅酒錄,夜夜笙歌, 驕縱奢華,亂拜偶像.和古代罪惡之城所多瑪和俄摩拉,同屬相同類型的城市；而在現代社會中,許多工商業為主的大城市,如紐約,芝加哥,倫敦,洛杉磯,乃至當前中國的上海,深圳,以及我們身處的香港,都是相類的商業城市.在近代城市發展歷史上,這些工商業城市興起時,吸引大量的外地和農村移民來做工,又多是單身男士, 他們聚居在骯髒狹窄的貧民窟或工廠宿舍,以煙酒,賭博和嫖妓消耗勞力賺來的錢財和生命,至於生意人,則追求財富,出賣靈魂的,比比皆是.故此,許多人都把這些商業大城標籤做“罪惡之城”(wicked city), 歎息它們對傳統價值觀的破壞,和對青春和生命的糟蹋.推羅,西頓,所多瑪,都是這種唯利是圖,追求奢華生活和物欲享受,及以商業主義(commercialism)作為核心價值的“不悔改”的城市.這些罪惡城市,當然後來都遭到上帝的懲罰.它們的悲慘結局,我們較能理解, 因為城中的人民背叛上帝,丟棄敬虔,以金錢為宗教,以物欲及人為的成就作為膜拜的偶像.

然而,耶穌所說的“不悔改”的城市,除了上面所說的商業大城推羅,西頓,所多瑪之外,也有哥拉汛,伯賽大和迦百農.這些城市是不同類型的,我叫他們做“家鄉之城”, 因為耶穌自小生活在這裡,而大部分門徒也都生於斯和長於斯.這幾個城市座落在加利利湖的沿岸,可以說是耶穌和門徒的家鄉.住在這裡的人,絕大部分都是認識耶穌和他家庭的人,當然也是熟悉耶穌門徒的人.但他們因熟悉而拒絕,因認識而不信,對耶穌和門徒傳播的福音不肯接受,反呲之以鼻,譏笑耶穌及門徒說:本地能出甚麼先知嗎?

我們認為我們認識一個人,主要是靠著熟悉或熟絡,而熟悉的根據又往往由於掌握很多資料:口傳的,文獻的,或在現代來說是從電腦來的,是谷歌(google)得來的,但這種並不真實的掌握資料所帶來的熟悉,造成了我們的妄自尊大和自以為是.誰人比我--一個香港仔--更熟悉香港和瞭解香港?誰人比我--一個中國人--更認識中國和更明白中國?這就是“家鄉之城”的一種思維方式,一種近乎“民粹主義”和“本土主義”的思維和心態.縱然耶穌在迦百農,伯賽大和迦百農行過神跡,施行教導, 廣傳福音,和醫病趕鬼,但在“家鄉之城”卻受到拒絕!百姓不信,因為他們自以為熟悉.誰是耶穌?他不是木匠的兒子嗎?拿撒勒能出甚麼人來呢? 熟悉產生懷疑和不信, 自以為掌握實情和真相,卻是對自己的瞞騙.這就是我對哥拉汛,伯賽大和迦百農這幾個“家鄉之城”被主認為“不悔改城市”的解釋.

結語:朝向聖城之路

耶穌和門徒是在往耶路撒冷的途中,碰到“不接待的鄉村”和“不悔改的城市”的.我們跟隨主,走人生的道路,也是一條走向耶路撒冷之路,一條朝往聖城之路, 一條通向天國之城的路. 在我們的人生路上,在我們的信仰路上,或在我們的心路歷程上,我們也必然會和主的門徒一樣,會經過城與鄉,也必然會受到人的冷遇,也必然會有兩個“不”,即“不被接待”(non-hospitality)和“不願悔改”(non-repentance)的經驗.我們的朝向聖城之路,我們的信仰之路,我們作基督徒跟隨主的路,我再強調地說:路途並不是一帆風順的, 途中充滿羈絆,頗多荊棘,今日只提出路加福音提及的兩個 “不”(不接待及不悔改,亦即是拒絕和遺憾),和大家一起檢討與反思.這兩個“不”及其所代表的“拒絕”與“遺憾”,及其植根的土壤--“狹隘的地方主義”和“極端的民族主義”,物欲掛帥的“商業主義”和自我中心的“知識主義”,在信主之後的人生路上,仍然必會遇到,仍然會時刻伴隨著我們.然而,值得我們信靠和安慰的,是有主同在；而有主同行,這樣我們就可以一直走到聖城--天國之城--那是我們和一眾信徒的目的地,也是基督徒人生之路的終點.



\newpage

\section{你是我的愛子}
\label{sec:277}
\textbf{鄧瑞強博士 - 你是我的愛子}
\newline
\newline
\hyperref[sec:276]{< < < PREV SERMON < < <}
~
\hyperref[sec:toc]{[返目錄]}
~
\hyperref[sec:278]{> > > NEXT SERMON > > >}


馬可福音 1:9-15
\newline
\begin{longtable}{cl}
\hline
\hline
章節 & 經文 (和合本修訂版)\\
\hline
1:9 & \begin{tabularx}{0.7\textwidth}{X} 那時，耶穌從加利利的拿撒勒來，在約旦河裡受了約翰的洗。 \end{tabularx} \\ \\ \relax
1:10 & \begin{tabularx}{0.7\textwidth}{X} 他從水裡一上來，就看見天裂開了，聖靈彷彿鴿子降在他身上。 \end{tabularx} \\ \\ \relax
1:11 & \begin{tabularx}{0.7\textwidth}{X} 又有聲音從天上來，說：「你是我的愛子，我喜愛你。」 \end{tabularx} \\ \\ \relax
1:12 & \begin{tabularx}{0.7\textwidth}{X} 聖靈立刻把耶穌催促到曠野裡去。 \end{tabularx} \\ \\ \relax
1:13 & \begin{tabularx}{0.7\textwidth}{X} 他在曠野四十天，受撒但的試探，並與野獸同在一起，且有天使來伺候他。 \end{tabularx} \\ \\ \relax
1:14 & \begin{tabularx}{0.7\textwidth}{X} 約翰下監以後，耶穌來到加利利，宣講神的福音， \end{tabularx} \\ \\ \relax
1:15 & \begin{tabularx}{0.7\textwidth}{X} 說：「日期滿了，神的國近了。你們要悔改，信福音！」 \end{tabularx} \\ \\
[1ex]
\hline
\hline
\end{longtable}
講題:「你是我的愛子」　You are my Beloved Son

經文:馬可福音1:9-15

講員:鄧瑞強博士

場合:香港中文大學崇基學院禮拜堂

日期:2015年2月22日

各位弟兄姊妹,早安.

新春期間,祝願各人蒙神賜福,生命行走在正義的路上.

今日的講道經文,是今日的福音書經課:馬可福音1: 9-15.

這段經文很短,用三個情節,交代一個人物的出場,交代他的身分,交代他要經歷的事情,交代他要成就的事業.簡言之,交代他的存在意義,也交代他對我們的存在所產生的意義.

馬可福音1: 9-15

可1:9　那時,耶穌從加利利的拿撒勒來,在約旦河裡受了約翰的洗.

可1:10　他從水裡一上來,就看見天裂開了,聖靈彷彿鴿子,降在他身上.

可1:11　又有聲音從天上來,說:「你是我的愛子,我喜悅你.」

可1:12　聖靈就把耶穌催到曠野裡去.

可1:13　他在曠野四十天,受撒旦的試探,並與野獸同在一處,且有天使來伺候他.

可1:14　約翰下監以後,耶穌來到加利利,宣傳神的福音,

可1:15　說:「日期滿了,神的國近了.你們當悔改,信福音!」

耶穌是如何走進歷史舞台的.

(第一幕:受洗)

馬可福音沒有馬太福音和路加福音那些我們熟悉的聖誕故事,沒有天使報訊,沒有馬利亞的未婚懷孕,沒有約瑟的徬徨,沒有希律王的恐懼,也沒有遠道而來朝拜的博士.耶穌一出場,就是他受洗的場面.

他,來自加利利的拿撒勒,來自窮鄉僻壤,來自社會的底層.一個不起眼的人,去約翰那裡,領受悔改的洗.以這一舉動,他與所有需要悔改認罪的世人站在一起.從社會層面言,他是卑微者.從宗教角度言,他與卑微者站在一起.

他來到約旦河,受約翰的洗禮.「洗禮」這個字,在馬可福音的用字裡,也代表著苦難.在馬可福音後來的經文裡,當耶穌走向十架苦難時,他問他身邊的門徒:「我所喝的杯,你們能喝嗎?我所受的洗,你們能受嗎?」(可10:38)故此,耶穌的受洗,預告著他背負人類的苦難.這是一個注定要承受人類苦難的人物.

可1:10　「他從水裡一上來,就看見天裂開了,聖靈彷彿鴿子,降在他身上.」

「天裂開了」,是頭等大事.舊約先知以賽亞講到:「願你裂天而降」(賽64:1),天裂開了,就是神要在人間行事的時候.同時,聖靈以可見的形象降在他身上,他的生命,將是天命貫注的生命,也是天命流瀉出來的生命.

可1:11　「又有聲音從天上來,說:『你是我的愛子,我喜悅你.』」

天上的聲音是肯定耶穌的獨特身分.他雖然是卑微者的一員,但他卻也是被神特別肯定的一個人.記載這同一件事情時,馬太福音的寫法是:「從天上有聲音說:『這是我的愛子,我所喜悅的.』(太3:17)」馬太的寫法聽起來,是神向眾人介紹耶穌.馬可福音的寫法,卻側重這是神向耶穌作的一種私隱認定.對象是「你」.面向天父,耶穌體會到一種獨特關係,他是「神的愛子」,神喜悅他.

「你是我的愛子」,這經文源自詩篇2:7,在詩篇那裡,這用字是指向彌賽亞君王.「我所喜悅的」,這經文源自以賽亞書42:1,經文這裡是指向神的卑微僕人.在洗禮裡,在耶穌心中,天父與他有一種獨特交流,他完全明白他在天父心中的角色.他是以色列人期待的王者,但同時是神差來人間的最卑微的僕人.

耶穌,一個卑微人物,天因他而裂開,天命將透過他而彰顯大地.他,一個卑微人物,聖靈之力卻在他生命中燥動.他,一個卑微人物,卻聆聽到天父對他的生命的肯定.他,生命的尊貴像王者,卻是要以背起人類重負的方式,一種奴僕的方式,去展示這王者的生命.洗禮,表示他與人類的苦難認同,同時,顯明他以自甘受苦的方式去承托受苦的人類.

馬可的經文很簡短,卻充滿無比張力.耶穌像是一顆細小的種子,但這卻是一奇特的種子,當它長成後,巨大無比,其根可穿透整個地球.

(第二幕:受考驗)

一個人明白自己的身分,便要接受這個身分的考驗.

可1:12　「聖靈就把耶穌催到曠野裡去.」

這裡的「催」字,有「強力催迫」的含意.內貫於耶穌生命的聖靈,催迫耶穌去曠野.耶穌不能「慢吞吞」的,擇好時機才去面對身分的考驗.

聖靈催迫耶穌去到曠野.曠野遠離城巿,與人類日常的城巿生活有區隔.在曠野,一切不再一樣,生活再不能理所當然.不能一打開冰箱,就有食物.不能一開空調,便有冷暖氣.進入曠野,一個人便不能按自己的意欲而生活,一切將仰賴天父的安排.這是一個失去自己而卻又尋回自己的地方.

可1:13　「他在曠野四十天,受撒旦的試探,並與野獸同在一處,且有天使來伺候他.」

在曠野四十天,當然指向過去以色列人出埃及後在曠野四十年的經歷.這當然也是今日開始的四十天「大齋期」(Lent)的依據.四十天,受撒旦的試探考驗.考驗的內容只有一樣,你會否因著塵世的種種原因(如:權力,利益,情慾的快樂等),抗拒聖靈的感召,遺忘天父的召命,而捨棄自己作為神的兒女的身分.

「與野獸同在一處,且有天使來伺候他」.「與野獸同在一處」,應作何解.學者有兩種基本看法.一批學者對「野獸」抱正面看法,他們認為這是回到伊甸園的情景,人和動物和諧共處,耶穌的角色是作為「新亞當」,示範如何回應試探.或者,這是預示「新天新地」裡人和動物和諧共處的局面,在此,耶穌的角色便是「神自己」,是神將「新天新地」引來(參:賽43:20).

另一批學者對「野獸」抱負面看法,野獸代表危險與邪惡,與撒旦一樣.

在我看來,在這裡,「野獸」和「天使」正好作一對比.就算活在野獸群中,就算豺狼當道,仍不失自己的身分,仍相信上天的使者的保護.

短短兩節經文,描述了耶穌面對的考驗的艱難.在曠野,一切習以為常的生活都放下了.周圍都是野獸,一切安全平靜的局面都失去了.撒旦無處不在的試探,要將一個人生命裡點點滴滴的神聖性都要除滅得一乾二淨.神聖性失去了,人便只是行屍走肉.

耶穌頂住這些考驗,他仰望蒼天,守住生命的靈性,信賴天使的幫助.

面對這般考驗,仍不折腰,他真是一個人,他真配稱為「神的愛子」.

(第三幕:公開發言)

一個人明白自己神聖的身分,一個人頂得住考驗,這個人才有資格向這世界發言.

有些人一發言,全世界便作嘔,因為他們毫無人性,向撒旦低頭,活得似一頭野獸.

可1:14　「約翰下監以後,耶穌來到加利利,宣傳神的福音,

可1:15　說:『日期滿了,神的國近了.你們當悔改,信福音!』」

耶穌向世人宣講一個不一樣的信息.這不是股巿上落的消息,不是樓價升跌的預測,不是成功之道.他宣講來自神的信息,來自神的好消息,來自神的福音.

「日期滿了」,這裡「日期」的聖經用字是 “kairos”,這個字強調的是「時間的質」,而不是「時間的量」.從「時間的量」而言,我們今日坐在這裡崇拜,各人花的「時間的量」是一樣的.但從「時間的質」而言,若你在崇拜中睡著了,則你度過的時間只有「量」而沒有「質」；若你在崇拜中遇見神,聆聽到他的呼聲,則你度過的時間的「量」可能只是一個鐘頭,但時間的「質」卻有永恆的重量.對於一個活得毫無神聖的人而言,縱使他很長命,他只是在「時間的量」中流逝,毫無質感.對於一個活出自己的神聖身分的人而言,縱使他很命短,他卻在「時間的量」中留下永恆的足跡.

耶穌說:「日期滿了」,是時候各人面對自己生命的「時間的質」了.你的時間,只是「量的流逝」,抑或滿有「神聖的質」?中國人說,「虛度光陰」,就是光陰走過後,留下的,只是虛空一片,沒有任何實質.

「神的國」即「天國」,「天國」不是「天堂」,天國是指神在這個塵世的實質管治.天國近了,「近了」即是「不遠」,但「近了」也是「未到」.若即若離,似遠還近.用神學的講法是:「既濟」與「未濟」(already but not yet).天國的「既濟」與「未濟」,看來是神已完成他要走的那段路,現在是等接棒的我們走完我們的路.我們還在賽道中,觀眾緊張莫明,靜待事情的變化,因我們會跌棒,會跌倒,會因種種原因放棄.

天國大於人心,但天國的實現,卻有待人心的改變.「你們當悔改,信福音!」所謂「悔改」,所謂「信」,要求人整個生命的改變,「覺今是而昨非」.真正的悔改,是經得起考驗,活出你的生命的神聖性.

耶穌的受洗,受考驗,公開宣講這些事蹟,對我們有何意義呢?

先講一下「洗禮」.「洗禮」這個字,已進入日常語言裡.我們有很多「洗禮」.有「成功的洗禮」,有「失敗的洗禮」,有各種各樣的「洗禮」.

為了探索香港人經歷過什麼「洗禮」,和經歷這些「洗禮」後人有什麼變化,我好奇地在google輸入「雨傘運動的洗禮」作搜索,走出來的第一篇文章是許寶強教授寫的＜毋忘初衷＞(\url{http://www.inmediahk.net/node/1029592}).文章的一段文字這樣寫:「本來以為,亂世中的教育工作,會舉步為艱.然而,雨傘運動的洗禮,卻開啟了過去難以想象的可能性.過去的二個多月,不少以各種方式參與這場運動的學生青年,快速成長,在不同程度拓展了對社會的認識,改變或鬆動了單一狹隘的價值觀念,進一步掌握組織活動照顧自己和他人的能力,部分同學的表現更叫人刮目相看,甚至令我們這些在大學任教的老師懷疑,過去和現在給學生的三,四年大學課堂教育,成效是否還不及接受七十多日的廣場薰陶?教育不就是啟發知識,孕育價值和學習技能的工作嗎?」

一場社會運動的「洗禮」,會改造某些人的生命.這些人的改變,又會進一步改變另一些人對這些人甚至對這世界的看法.這些改變,當然又會進一步引發其他改變.

我們的一生,會經歷很多不同的「洗禮」,這些「洗禮」可能令我們走上「正途」多一點,也可能令我們走上「歧途」多一點.耶穌受洗後,神面對面的跟他說:「你是我的愛子,我喜悅你.」這可說是展示了一種終極洗禮,耶穌在這洗禮中呈現他的終極身分.我們又如何呢?我們有否在人生種種洗禮中經驗這種終極洗禮?有否在其中領悟到,我們的終極身分,不是什麼「成功之子」,「財金之子」,「革命之子」,而是「神的愛子」.你是神的愛子.當然,耶穌作為「神的愛子」和我們作為「神的愛子」的意義不一樣.耶穌作為「神的愛子」,是他永恆的本性就是如此.我們作為「神的愛子」,是出於神的恩典和救贖.在神的恩典和救贖中,你是神的愛子.你有否忘記你這身分?你的生命有一種來自神的「神聖性」,來自神的靈明,來自神的永恆之光.在這個充滿獸性的世界裡,重拾你的神聖身分,是至關重要的.你是神的愛子,別忘記你這個身分.

若你沒有忘記「你是神的愛子」這身分,考驗便來了.在一個極權世界,你堅持生命的神聖,麻煩便來了,撒旦便來了,與野獸一樣的人便來了.他們要你忘記身分,忘記生命的神聖,最終,要你忘記活在世上的意義.當你一切都忘記時,「時間的質」便消失了,你只在「時間的量」中悄悄流逝.曠野是凶險的.請記住你的身分,猶如記住你的名字.你是神的愛子.你的身分意味著一種特權,在神的恩典裡,作為神的兒女,你們配得天使的幫忙.

最後,我要說的是,那些守住自己「神的兒女」身分的人,那些頂住世界考驗的人,他們的生命,對世界而言,就是一個聲音,一個呼喚,一個宣告.他們是歷史的發言人,他們是人性的嚮導,是神的聲音.

你是神的愛子.你有沒有忘記這身分?

但願:榮耀歸於聖父,聖子,聖靈.阿們.



\newpage

\section{彌賽亞的命運交響曲}
\label{sec:278}
\textbf{伍渭文牧師 - 彌賽亞的命運交響曲}
\newline
\newline
\hyperref[sec:277]{< < < PREV SERMON < < <}
~
\hyperref[sec:toc]{[返目錄]}
~
\hyperref[sec:279]{> > > NEXT SERMON > > >}


講題:彌賽亞的命運交響曲　The Fate Symphony of the Messiah

經文:詩篇二十二

講員:伍渭文牧師

場合:香港中文大學崇基學院禮拜堂

日期:2015年3月1日

孤獨面對信仰.自從有了孫女後,對一句話:「上主祇有兒女,沒有孫兒女」特別有感覺.我們的父母是基督徒,但我們不一定是基督徒,成為基督徒是自已非常孤獨的抉擇.我記起中學信主後非常熱心傳福音,有一次邀請了一位小學同學到教會聽佈道會,到講員呼召時,我見同學沒有舉手決志,我一時心急用雙手托起他的手；當然這是犯規,無效.

獨孤面對死亡.面對死亡,人也是孤獨的.我們可以在病榻陪伴,安慰鼓勵,但每一個人都要獨自經歷死亡,一如獨自經歷信仰.信主和死亡都是回到生命的主那裡,是頂個人的事,沒有人可以代替.

今天講道經文取自詩篇經課詩篇二十二篇,作者面臨死亡,他祈求上主的拯救.從經文內容來看,死亡的威脅可能來自疾病:「我如水被倒出來,我的骨頭都脫了節,我心在我裡面熔化.我的精力枯乾,如同瓦片；我我的舌頭貼在我牙床上.你將我安置在死地的塵土中.」(vv. 14, 15) 但詩人憑信心知道上主已聽他的禱告:「你已經應允我(v. 21下),轉哀歌為讚美感恩.

大齋期的詩篇.今天是大齋期(Lent)第二主日,詩篇經課選讀詩篇二十二篇是有原因的.大齋期是復活前四十天的預備期,從前周三聖灰日開始計數,惕除歡樂的主日,剛好是四十天,效法耶穌在曠野禁食祈禱.

耶穌降生聖誕節的預備期是將臨期(Advent),重點是情緒的牽動,將臨期燭光愈來愈明亮,我們渴慕基督的來臨,此期間是黑暗和光明的糾纏,我們願意驅走心中的黑暗,迎接基督的降生.而大齋期的預備期重點是意志的磨練,大齋期是死亡與生命的爭戰.在這四十天,耶穌在曠野思考他的命運.他將公開傳揚福音,呼召和訓練門徒,教導他們有關天國的奧秘,告訴他們他要赴耶路撒冷受死,藉著死打敗撒旦,然後復活升天,差遣聖靈澆灌門徒,他們將得著能力把福音傳遍天下,然後他會再來,審判世界.

四十天在曠野耶穌禁食祈禱,思考他自己的命運,在世上要完成的使命.

命運交響曲.在一次的講道中,我說到詩篇二十三篇「耶和華是我的牧者,我必不致缺乏,他使我躺臥在青草地上,領我在可安竭的水邊……」像是貝多芬的第六交響曲-田園交響曲.但這位牧者,帶領我們進入恬靜田園之前,曾為我們的靈魂奮勇爭戰,犧牲自己,他是配得我們跟隨的牧者.詩篇二十二篇像是貝多芬的命運交響曲,第一樂章明亮的快板,激昂奮進,最後凱旋得勝.

第五卷福音書.詩篇二十二篇詳細描述耶穌在十字架被釘的過程,甚至被稱為第五福音書.福音書的作者多次引用詩篇二十二篇的話,描述被釘十架的耶穌:「我的上帝,我的上帝!為何離棄我?……他們扎了我的手,我的腳.……他們分我的外衣,為我的裡衣拈鬮」,「他把自交給耶和華,耶和華可以救他吧!」

我們看看全詩的結構:

敬拜禮儀:首先我們要明白,這是一首崇拜禮儀詩歌,轉哀歌為感恩讚美的關鍵就在上主的應允:「你已經應允我」.(v. 21下)「謙卑的人必吃得飽足」(v. 26),「地上一切豐肥的人必吃喝而敬拜」(v. 29)指敬拜者和祭司一起吃獻上的祭物.

哀歌傾訴:詩歌一開始是哀歌,詩人藉哀歌傾訴苦情,vv 1-21；前段 vv. 1-10描述他感到被上主遺棄,第十一節是第一次祈求:「求你不要遠離我,因為急難臨近了,沒有人幫助我.」後段vv. 12-18描述四面楚歌:「有許多公牛圍繞我,巴珊大力的公牛四面困住我,它們向我張口,好像撕裂吼叫的獅子.」(v. 12)「犬類圍著我,惡党環繞我；他們扎了我的手,我的腳.」(v. 16)接著vv.19-21上是第二次的祈求:「耶和華啊,求你不要遠離我!我的救主啊,求你快來幫助我!求你救我的靈魂脫離刀劍,救我的生命脫離犬類,救我脫離獅子的口；使我脫離野牛的角(放在你已經應允我前面,平行前句).」哀歌過後是神諭(o racle),上主藉祭司的宣告,這裡是暗示的神諭,沒有直接說「神如此說.」

回應(暗示):「你已經應允我」(v. 21下).

讚美感謝:二十二節至三十一節是讚美和感謝,先是祈求的人自己讚美感恩:「我要將你的名傳與我的弟兄……我在大會中讚美你的話是從你而來；我要在敬畏耶和華的人面前還我的願.謙卑的人必吃得飽足,尋求耶和華的人必讚美他.願你們的心永遠活著!」(vv. 22-26).

跟著會眾一起讚美感恩:「地的四極都要思念耶和華,並且歸順他；列國的萬族都要在你面前敬拜.因為國權是耶和華的；他是管理萬國的.……」(27-31)會眾一直在會中細心聆聽面對死亡威脅的信徒的哀歌傾訴,然後和他一起感恩讚美.

看完全詩的結構後,我們看看這詩對我們的現實意義.

獨處與和群(Solitude and Community)

在信仰的經歷上,我們也會面有時感到被上主遺棄.現在的遭遇和我們的神學和過去的經歷完全割裂,完全不同.對舊約的聖徒來說,我們和上主的關係牢不可破,連死亡也不間斷.列祖雖然離開世界,但他們和上主的關係沒有終止；上主是亞伯拉罕的神,雅各的神,以撒的神.甚麼是永生?永生就是和上主緊密的關係,這關係連死亡也不能間斷.「我們的祖宗倚靠你；他們倚靠你,你便解救他們.」(v. 4)但此刻詩人覺得上主不理他,遠離他.「我的上帝,我的上帝!為何離棄我?為甚麼遠離不救我?不聽我唉哼的言語?」(v. 1)

他現在的經歷和信仰是衝突的,也和過去的經驗相左的.「但你是叫我出母腹的；我在母懷裡,你就使我有依靠的心.我自出母胎就被交在你手裡；從我母親生我,你就是我的上帝.」(vv. 9, 10 )

在信仰的征途中,有時我們也會像詩人一樣,感到迷惘,我們經歷過上主的看顧,我們的禱告也蒙應允,上主顯得很真實可靠,但有時上主好像隱退了,我們感到被上主遺棄,上主好像跟我們變臉,我們感到極度孤單.

被人遺棄.更令他難受的是病重的他,得不到朋友的支援,反而被嗤笑,說他祈禱沒有用了,說上主放棄他了,可能他們認為他的病就是上主對他的懲罰.或者,這些根本不是他的朋友,是敵人,幸災樂禍.他深深感到被徹底遺棄,獨自面對死亡的臨近.

我以前遇到一個老人家,女兒和太太都來教會聚會,但自己不來,有一天他病了要進醫院留醫.一位認識他已退休的傳道人對他說:「就是因為你不去聚會,上主讓你生病.」他知道要進行手術,但卻拒絕進手術室.後來我探望他,告訴他上主愛他,叫他信靠上主接受手術.我跟他說不知道你為何生病?也許人年紀大了,生病的機會增多,但我們的生命在愛我們的上主手中,若他不同意,一根頭髮也不會掉下來,來上主賜下平靜安穩,面對一切.我說完後,他改變了態度,幾天後就接受手術.

以前他怕進手術室,因他以為上主會為難他,是他的敵人,躺在手術床太危險了!

群體的承托

但詩人並沒有放棄,他到上主面前,也在信仰群體面前,藉哀歌傾訴,而眾聖徒靜心聆聽,支持著他.上主沒有離開他,信仰群體也沒有離開他.公共崇拜要莊重敬虔,但不能流於形式,禱告要有血有肉,上主顧念我們真實的需要,最個人的也是最普及(the most personal is the most universal) .在聚會中,眾人靜默細聆聽一個人的哀歌時,他們也聽到自己的聲音,透過這人的禱辭,他們也向上主禱告.

彌賽亞被棄絕.在彌賽亞身上,我們也看到孤獨,甚至被棄絕,同時看到群體和契通,solitude and community .好像一個要生孩子的母親,沒有人可幫助她去承受生產之痛苦,她要孤獨經歷,但支持她的是將要誕生的孩子.她一想到腹中的骨肉要誕生,她看到嬰兒將面世,她忍受生產之苦楚,單獨承受,她帶著喜悅去承受這苦難.因基督被棄絕,帶來了修和和友愛群體的誕生:「我要將你的名傳與我的弟兄,在會中我要讚美你.」(v. 22) 希伯來書把這句話用在教會身上,但在前面加了:「原來那為萬物所屬,為萬物所本的,要領許多的兒子進榮耀裡去……所以,他稱為他們為弟兄也不以為恥,說:我要將你的名傳與我的弟兄,在會中我要頌揚你.」(來二10-12)

希伯來書的作者著跟說:「兒女既同有血肉之體,他也照樣親自成了血肉之體,特要藉著死敗壞那掌死權的,就是魔鬼,並要釋放一生因怕死而為奴僕的人.」(來二14, 15)基督的爭戰,藉著死使我們危去對死亡的懼怕.

中古神學家說基督的死,像魚餌,表面被吃掉,但魚被釣住了.死亡吃掉基督,但基督藉著死敗壞那掌死權的魔鬼,基督死後第三天復活就把死的毒釣除掉了,此後我們不再懼怕死忙了.

不再懼怕死亡 (Fear Death No More)

上主的應允

轉哀歌為感恩讚美的關鍵就在上主的應允:「你已經應允我」.(v. 21下)這是整個禮儀中的轉捩點.但請注意,這裡沒有明顯的神諭(oracle) , 說他的病將被治好,那麼詩人得到甚麼應允呢?他認為上主應允不會離開他,上主的群體也不會離開他,他們正在會中聆聽他的哀歌哩.上主不離開就不懼怕死亡.詩人第一個祈求:「求你不要遠離我!因為急難臨近了,沒有人幫助我.」(v. 11) 第二個祈求:「耶和華啊,求你不要遠離我!我的救主啊,來你快來幫助我!求你救我的靈魂脫離刀劍.」(vv 19, 20)目的是上主不要離開,不是醫好他的病,他最關注的是若重病帶他離開世界,上帝不和他同在時,這才是最令他害怕,「急難臨近」極可能是指死亡臨近.到了死亡的門檻,親友止步,除了上主,誰能幫助我們呢?

一位中學校長獻身作傳道,分享在醫院牧靈經歷.在醫院垂危病人的病榻前,他發現他很重要；當醫務人員束手無策,親友充滿焦慮,眾人都沉默不語時,他的臨在陪同很有象徵信仰的盼望,上帝的安慰,他說這身份是當中學校長時所沒有的.在死亡面前,我們沒有任何的幫助,祇有上帝,上帝應許不離開,他就轉哀歌為讚美.

若上主醫好他,他要感恩,抑或不然,他也釋然,因為上主沒有離開他,死亡也不能間斷他和上主的關係.

十字架除了對死亡的懼怕

這詩描述詩人因病患弄至身體的枯槁:「我如水被倒出來,我的骨頭都脫了節,我心在我裡面熔化.我的精力枯乾,如同瓦片；我我的舌頭貼在我牙床上.你將我安置在死地的塵土中.」(vv. 14, 15) 但基督是被酷刑催殘至身體的枯槁.

受難曲.十字架的刑罰是羅馬人非常不人道的刑具,有羅馬公民權的免釘十字架.二零零四年二月二十五日聖灰日米路吉遜(Mil Gelson)導演的《受難曲》(The Passion of the Christ) 上影,因為太血腥定為三級片,其實導演祇是寫實,十架的刑罰就是那麼血腥殘酷.

《馬利亞的泣訴》.最近香港電台一分鐘閱讀介紹當代愛爾蘭文學巨人柯姆．托賓(Colm Tóibín)所著《馬利亞的泣訴》(The Testament of Mary).全書以第一身書寫,這位全程目睹兒子慘痛遭遇讓她痛不欲生,她的聲音終於被聽見:「我看見我的兒子被綁起來而且流著血,我聽見他的哭喊……」當長釘刺穿耶穌手臂時,兩千年來不曾有人聽過她(馬利亞)的一聲尖叫.聖母的沉默,至今才有人書寫她的泣訴.

殘酷施虐因為懼怕死亡.公開殘酷虐打並加上一些儀節,可以合理化殘酷,使施刑者覺得是替天使罰.當旁觀者從其中得到快樂時,他們就愈來愈大胃口,要求更大的殘酷,最後受虐者便被處死了；公開處死受刑者,像基督被釘十字架,其實是舒減旁觀者和施暴者對死亡的害怕.

現代殉道士.公開對不能反抗的人施虐,因宗教信仰而被處決,更是極大的邪惡.我想我們都因為二十一名在利比亞遭俘虜的埃及科普特(Coptic)基督教徒被公開處決,心裡感到忿怒和不安.

這些人和家人分隔,遠赴利比亞謀生,因為愛；他們不放棄信仰,寧死也不背叛基督,因為信.他們肯定會為了脫離危難而祈禱,但上主沒有救他們脫離敵人的手,但應允不離開他們.在熒幕上,我看到他們臉沒懼色,安詳跪在地上,我想他們下跪時在敬拜上帝；他們的血與基督連在一起,成為可敬佩的殉道士(martyr)-見證人.

3        透識真相,為信仰而活.

我在想,要是我是那二十一位埃及基督徒,我會為了生命背棄信仰嗎?我們也許沒有這福氣做流血的殉道士.但我想,一個平時不會為信仰而活的人,他會在關鍵時刻,為信仰而死嗎?我想起在更早時候在熒幕上也被公開處決的日本記者後藤健二(Kenji Goto) .後藤是因為前往敘利亞試圖解救另一名日本人質湯川遙菜(Haruna Yukawa)而被捕的.為何後藤冒生命的危險去救朋友呢?

我想因為信仰,後藤是基督徒,他不單關心朋友,也關心在戰地生活的人.他說過:「那些我去採訪報道的地方遭遇著巨大的困難,但在那裡人們也每天生活著……那些人總有些話要說,有信息要傳達.如果我能幫助他們將信息傳遞給世界,那麼就可能促成某種解決方法.」後藤因為使其他人活得更好而失去生命,他也是殉道者-見證人.他的血與基督連在一起,為基督作了美好的見證.

我想,若我們能為信仰而活,就比較有可能為信仰而死.

大齋期開始的塗灰禮:當主禮把灰塗在額上時說:「你來自塵土,也要歸回塵土,要遠離罪惡.」明白我們來自塵土,也要歸回塵土,死亡就不是那麼懼怕了.

聯合國第二任秘書長(1925-300韓馬紹(Das Hammarskkjold  ) 對死亡和命運有睿見: 「我們不需要尋求死亡,死亡會找上門的.但我們要找到往死亡的路,使死亡顯得有意義……命運本身不能逃避,也不需響往,它不是沒有原委的隱密,因為它說明這世界和人類歷史的跡,都有意義.」(Do not seek death. Death will find you. But seek the road which makes death a fulfillment…Destiny is something not to be desired and not to be avoided, a mystery not contrary to reason, for it implies that the world, and the course of human history, have meaning.)

沒有人響往死亡,除了基督；因為十字架是他的命運交響曲,他在其上為我們奮勇爭戰,打敗死亡.「兒女既同有血肉之體,他也照樣親自成了血肉之體,特要藉著死敗壞那掌死權的,就是魔鬼,並要釋放一生因怕死而為奴僕的人.」



\newpage

\section{行真理者必來就光}
\label{sec:279}
\textbf{鄧瑞強博士 - 行真理者必來就光}
\newline
\newline
\hyperref[sec:278]{< < < PREV SERMON < < <}
~
\hyperref[sec:toc]{[返目錄]}
~
\hyperref[sec:280]{> > > NEXT SERMON > > >}


約翰福音 3:14-21
\newline
\begin{longtable}{cl}
\hline
\hline
章節 & 經文 (和合本修訂版)\\
\hline
3:14 & \begin{tabularx}{0.7\textwidth}{X} 摩西在曠野怎樣舉蛇，人子也必須照樣被舉起來， \end{tabularx} \\ \\ \relax
3:15 & \begin{tabularx}{0.7\textwidth}{X} 要使一切信他的人都得永生。 \end{tabularx} \\ \\ \relax
3:16 & \begin{tabularx}{0.7\textwidth}{X} 「神愛世人，甚至將他獨一的兒子賜給他們，叫一切信他的人不致滅亡，反得永生。 \end{tabularx} \\ \\ \relax
3:17 & \begin{tabularx}{0.7\textwidth}{X} 因為神差他的兒子到世上來，不是要定世人的罪，而是要使世人因他得救。 \end{tabularx} \\ \\ \relax
3:18 & \begin{tabularx}{0.7\textwidth}{X} 信他的人不被定罪；不信的人已經被定罪了，因為他不信神獨一兒子的名。 \end{tabularx} \\ \\ \relax
3:19 & \begin{tabularx}{0.7\textwidth}{X} 光來到世上，世人因自己的行為是惡的，不愛光，倒愛黑暗，這就定了他們的罪。 \end{tabularx} \\ \\ \relax
3:20 & \begin{tabularx}{0.7\textwidth}{X} 凡作惡的人都恨惡光，不來接近光，恐怕他的行為被暴露。 \end{tabularx} \\ \\ \relax
3:21 & \begin{tabularx}{0.7\textwidth}{X} 但實行真理的人就來接近光，為要顯明他的行為是靠神而行的。」 \end{tabularx} \\ \\
[1ex]
\hline
\hline
\end{longtable}
講題:「行真理的必來就光」　Those Who Do What is True Come to the Light

經文:約翰福音3:14-21

講員:鄧瑞強博士

場合:香港中文大學崇基學院禮拜堂

日期:2015年3月15日

各位弟兄姊妹,早安.

香港社會流傳著一個神話,叫做「獅子山下」的神話.這個神話的內容,不是很精確的,但這神話很好用.當香港的經濟不好時,有人大聲喊:要發揮「獅子山下」精神.當香港的經濟好到不得了,但低下階層生活得越來越艱難時,有人大聲喊:要發揮「獅子山下」精神.當香港社會向上流動的路越來越窄時,有人大聲喊:要發揮「獅子山下」精神.當年青人無論如何努力工作都很難有安居之處時,有人大聲喊:要發揮「獅子山下」精神.無論這個城巿遇到什麼問題,看來,只要一喊「獅子山下」精神,便有解決希望.漸漸,這個所謂「獅子山下」的神話,變成一個體制「神話」,用來支撐現存體制.在一個世俗化的世界,政治體制容易扮演某種宗教的角色.宗教需要神話,使不滿的人在神話裡找到面對艱難的力量和希望.「獅子山下」的神話,安撫人心,讓生活在社會底層的人懷有某種希望.

後來,有人試圖賦予這個神話一點新的元素,便在獅子山頭,掛上一條直幡.這是一個安插在舊神話上的新符號,意味著一種新的神話精神.當然,維護體制的神話,怎能容許新的元素.很快,這直幡被移走了.

新的神話符號,其實是代表著一群人,他們在舊的神話裡,找到新的意義.但是,持守著舊的神話的人,怎能容許這種新意義.維護著舊的神話的人,是以這社會神話,守護著他們既有的一切.對他們而言,所謂「獅子山下」神話,就是無論社會變成怎樣,總希望香港人繼續默默耕耘,默默付出,力保社會的底層邏輯一切依舊.但對另一些人而言,所謂「獅子山下」精神,就是集合螞蟻的力量,去推倒大象.這兩群人對這神話,有不同的解讀.

原先,大家都是活在獅子山下,都是同一群人.但漸漸,對於如何活在獅子山下,如何發揮「獅子山下」精神,有不同解讀,便變成兩類人.這兩類人漸行漸遠,便形成正邪對立,我是光明的,你是黑暗的.這種對立撕裂,令人心酸,因為畢竟所有人都是活在獅子山下,都是在這裡一起長大的人.如何能在舊神話裡找到一些共同的身分回憶,以致能有一合作的空間；又如何在這神話的新元素裡看到新的可能,以致有變革的餘地.這真是一個大問題.

以上講的,其實是要交代約翰福音成書的一些背景.我今日講道的經文,是今日福音書經課裡的約翰福音的經文.

一般學者認為,約翰福音是一卷較後期的作品.書中的內容,反映成書時期猶太人和猶太基督徒的緊張關係.簡單來說,約翰福音的內容,一方面記述耶穌的事蹟,另方面卻同時反映耶穌復活升天幾十年後的時代形勢.

早期的基督徒,大多是猶太人.他們與猶太人,同一根源,活在同一天空下.在寫作約翰福音時,耶路撒冷的聖殿已被毀,人心虛怯,聖殿這宗教中心被摧殘,意味著心靈的無家可歸.若在香港,這差不多就等於終審法院及其代表的法治制度的瓦解,依賴法治而活的香港人,所感到的惶恐,應與當時的猶太人差不多.人心虛怯,便需要抓住一些東西來安頓自己.信仰的傳統,是人心虛怯時的救命符.當時的猶太信仰領袖,緊張兮兮地力保這傳統的純正.這不只是信仰,更是他們的身分認同.當這個世界什麼都瓦解時,能做的,便是力保自己的身分.而身分的核心,就是信仰.失去這一切,將失去自己及民族的將來.

就在這混亂時刻,猶太信仰的一個新興分支,塑造著一種新的猶太人身分,期盼著一種新的歷史前景.這就是基督徒群體的興起.漸漸,猶太人意識到基督徒的與別不同.縱使這些基督徒也是猶太人,但基督徒對耶穌的了解,與一般猶太人截然不同.同一天空下,這兩種猶太人分享著很多相同的東西.但在一個關節位上,他們有不同的理解.這關節位當然涉及上帝.上帝愛猶太人,但上帝能否愛猶太人以外的人?上帝關心猶太人的歷史前程,但上帝是否更關心整個世界的歷史前程呢?猶太基督徒是一批對耶穌所產生的意義極其敏感的人.他們在耶穌裡,看見神,看見神的新展現,看見神的新作為.他們在耶穌那裡,找到一種新的身分,找到新的安身立命的基礎,找到新的歷史使命.他們堅信,在耶穌裡,才有真正的明天.但是,一般猶太人看不到耶穌這種新意,他們一口咬定,耶穌是褻瀆神的人.他們認為一定要清除耶穌及其跟從者,才能保住傳統的信仰.沒有這傳統信仰,他們將失去身分,失去自己.

到了約翰福音的寫作時期,猶太人對基督徒的敵對,看來,到了一個高峰,到了勢不兩立的局面.在約翰福音裡,猶太人趕基督徒離開會堂(參:約9:22, 12:42, 16:2),這是耶穌生前未有過的事.對猶太人而言,被趕離會堂,意味著他們在家鄉被視為「外人」,在自己的土地被視為不宜在此居住.在社會衝突裡,有時,會有一些人叫另一些人移居外地,他們忘記了,這裡是這些衝突裡的所有人的家園.

這是叫人心酸的局面,社群撕裂,猶太基督徒被猶太同胞排斥,其排斥程度已到「正邪」對立的形勢,各自都為建立自己的身分而努力,各自都在混亂的歷史情境中找出路.這就是約翰福音寫作的背景.

有了以上背景,我們看看今日的講道經文:約翰福音3:14－21.

這段經文的上文,是講到一個叫尼哥德慕的法利賽人向耶穌問道.法利賽人,代表猶太信仰的古板傳統.約翰福音3:2那裡說,「這人夜裡來見耶穌」.這個古板傳統是屬於黑暗的,他不能接受耶穌顯明的新事物.他來求問「重生」的事情,求問如何能「從上而生」,求問如何能有一新的生命,新的身分.約翰福音3:10,耶穌「窒」他說:「你是以色列人的先生,還不明白這事嗎?」簡單的經文,表達了約翰福音寫作時的背景.猶太的宗教權威,不明白耶穌.這種不明白,在當時引發猶太人對基督徒的打壓,形成一種「黑暗與光明」的對立.當然,從基督徒的角度來看,猶太人是屬於「黑暗」的.

在這背景裡,耶穌講出約翰福音3:14－21的話.

約3:14　「摩西在曠野怎樣舉蛇,人子也必照樣被舉起來,

約3:15　叫一切信他的都得永生.

約3:16　神愛世人,甚至將他的獨生子賜給他們,叫一切信他的,不至滅亡,反得永生.

約3:17　因為神差他的兒子降世,不是要定世人的罪,乃是要叫世人因他得救.

約3:18　信他的人,不被定罪；不信的人,罪已經定了,因為他不信神獨生子的名.

約3:19　光來到世間,世人因自己的行為是惡的,不愛光,倒愛黑暗,定他們的罪就是在此.

約3:20　凡作惡的便恨光,並不來就光,恐怕他的行為受責備.

約3:21　但行真理的必來就光,要顯明他所行的是靠神而行.」

約3:14　摩西在曠野怎樣舉蛇,人子也必照樣被舉起來.

這裡涉及一個舊約典故,記載在民數記21:4－9.當時,以色列人離開埃及的奴隸生活,出了埃及後,走在曠野的路上.他們覺得,這條通往自由的路,太艱辛,太不值得.路上無糧無水,相比奴隸的生活,現在更不如.他們埋怨.神就使毒蛇出沒在他們中間.

他們哀求.神就叫摩西造一條銅蛇,掛在杆上.誰人望這銅蛇,誰就活了.通往自由的路,是艱難的,神沒有免去這艱難,神只醫治在這艱難的路上遇到的毒蛇的毒.約翰福音說:「人子也必照樣被舉起來」.神沒有拿走通往自由的路上遇到的毒液,但神派耶穌來,他可醫治這些毒液.

「被舉起來」這個「舉」字,有兩個含意.一是指耶穌被高舉在十字架上,受死.一是指耶穌被高舉在萬有之上,獲得無限的光榮.從約翰的神學而論,為了榮耀神而受死,為了解救他人而受死,這「受死」本身就是榮耀.

約3:15　叫一切信他的都得永生.

所謂「信」,就是參與這種生命.所謂「永生」,就是生活在神的臨在裡.這種生命,在耶穌的身上展現出來.

約3:16　神愛世人,甚至將他的獨生子賜給他們,叫一切信他的,不至滅亡,反得永生.

在NET Bible裡,這句經文的開頭是:“For this is the way God loved the world…”,然後,這聖經譯本加上一個注腳,這是指神愛世人的方式,同時也是指神愛世人的深度.神的方式是:捨去自己的至愛,去表達愛的純淨與深度.

在這裡,猶太基督徒與一般的猶太人對耶穌有不同看法.在一般猶太人眼中,耶穌在安息日醫病,是破壞宗教律法；與外邦人接觸,混亂了民族的界線；寬恕罪人,破壞了義與不義的劃分.而這一切,都是猶太人確立自己身分的關鍵因素.在混亂的時代,猶太人很緊張自己的獨特身分,他們自視為正統的,神的選民.他們以一種內向的,排外的,自義的,自覺純潔的,不容挑戰的,死守律法的方法,去建立他們的身分認同.耶穌的方式是他們很不習慣的.耶穌向他人開放,向異己開放,向仇敵開放.耶穌不用「伊斯蘭國」那種排外的仇恨去確定自己獨特的身分,耶穌以愛仇敵的方式去展示自己的身分.為了人與人之間的平等的愛,他甚至犧牲自己.而猶太基督徒,在耶穌那裡,看見他們傳統的信仰的一種新的展示,或者說,是他們傳統信仰的真實內容.

「神愛世人」是法利賽人極不習慣的想法.他們要不斷定人的罪,去肯定自己的義.他們死守規條到一個地步,看來是有違人情與常理的,藉此去保持自己的純潔.他們排斥外人,去保全自己「神的選民」的身分.用現今的話來說,就是為了保持自己的根正苗紅,為了保全自己「被上天選召」的身分,過一種排除異己的生活.他們認為,在混亂時代,這是他們能混下去的最佳身分.但基督徒在耶穌身上,看到愛,看到「神愛普世人類」,看到千差萬別的人能在歷史長留下去的必須要有的素質,即人與異己之間的互相尊重,寬容和仁愛.

耶穌來到世界,是否製造另一身分,製造另一標準,製造生死不相往來的「敵我」矛盾呢?不是的.經文說:

約3:17　因為神差他的兒子降世,不是要定世人的罪,乃是要叫世人因他得救.

耶穌不是要製造一標準,藉此去定人的罪,耶穌只是想展示,真正的生命可以是如何的.他展示一個在神的臨在下的生命,可以是怎樣的.這是活出「神愛世人」的精神的生命,這是走向異己,伸出友誼的手,以仁愛激發仁愛的生命.

約3:18　信他的人,不被定罪；不信的人,罪已經定了,因為他不信神獨生子的名.

咦,剛才不是說,耶穌降世,不是要來定人的罪的,現在為何又講到「不信的人,罪已經定了」呢?

耶穌的確不是要定人的罪,是人自己定自己的罪的.

當人人都和平溝通時,你忽然出手打人,則你的打人便顯示出你是打人的.在和平的人面前,你是打人的.你自己定了自己的罪.

約3:19　光來到世間,世人因自己的行為是惡的,不愛光,倒愛黑暗,定他們的罪就是在此.

約3:20　凡作惡的便恨光,並不來就光,恐怕他的行為受責備.

打人的,總在暗角裡,因為他們知道這不是光明的行為.經文說:「不愛光,倒愛黑暗,定他們的罪就是在此.」耶穌是光明,他沒有定黑暗的罪,他只是讓黑暗顯為黑暗.在黑暗中的人,他們的行為顯明他們是屬黑暗的,他們自己的行為定自己的罪.

在耶穌時代,「光明與黑暗」的對立,是猶太文化裡常見的元素,但不同派系有不同的詮釋.上一世紀有一重大考古發現,就是發現「死海古卷」.這些古卷可能是屬於猶太信仰的一個隱蔽社團「愛色尼派」(Essenes)的,其中一卷叫＜光明之子與黑暗之子的戰爭＞(The War of the Sons of Light Against the Sons of Darkness).死海古卷的宗教群體,將「光明」與「黑暗」提升到「形而上」的層次.「光明」與「黑暗」的戰爭,是一種宇宙戰爭.

耶穌時代,社會上有一群人,叫「奮銳黨」(Zealots)人,他們視「光明」與「黑暗」的對立,為政治上的「被壓迫者」與「壓迫者」的對立.他們殲滅黑暗,為的是重見光明.

耶穌在這文化裡,不採取以上看法,卻將「光明」與「黑暗」的對立視為一種道德契機.「黑暗」中的人,會否因有「光明」的人作對照,而知道自己活在「黑暗」中?在這裡,「光明」與「黑暗」的對比,不是宇宙大戰,不是政治的殲滅戰,而是點燃燭光,因有光明,而看出黑暗.

有時,看見一個人讓座給另一個人,我便看見自己在黑暗中.沒有人審判我,只是因為對照著光明,我知道自己在黑暗中.我自己定自己的罪.

約3:21　但行真理的必來就光,要顯明他所行的是靠神而行.

「神愛世人」,意謂神藉耶穌的生命,走在人間,走出自己的安全領域,去到異己當中,以仁愛激發仁愛.這就是「光明」.若這不是「光明」,我就不知光明是什麼了.靠神而行的人,就是活出這種仁愛的人.這些人願意在光天化日之下,而不是在暗角裡,展示這種生命.「行真理的必來就光」,活出生命大道的人,不怕光明,渴望光明.他們的行為與神的行為相稱,他們的行為顯明他們是屬於神的.

在一個混亂時代,容易造成正邪對立,黑暗與光明對峙.在耶穌時代及「後耶穌時代」,當時的宗教領袖為了保持一種「上天揀選了我」的身分,走向自閉,更大的排斥異己,置耶穌及基督徒於死地.神能用這種人,開創歷史的明天嗎?基督徒,在耶穌身上,看見神的仁愛,看見「神愛世人」,他們以仁愛去開拓自己的身分,這身分不建基於種族,社會身分或政治特權,凡在基督裡的,凡活出神的仁愛的,無論你是誰,都是兄弟姊妹.這是人類的一種全新的身分.人類的前景,若要有光明,不正是需要多一些這樣的人嗎?

「但行真理的必來就光,要顯明他所行的是靠神而行.」

但願:榮耀歸於聖父,聖子,聖靈.阿們.



\newpage

\section{人生關鍵時刻時候到了}
\label{sec:280}
\textbf{許開明牧師 - 人生關鍵時刻:「時候到了!」}
\newline
\newline
\hyperref[sec:279]{< < < PREV SERMON < < <}
~
\hyperref[sec:toc]{[返目錄]}
~
\hyperref[sec:281]{> > > NEXT SERMON > > >}


約翰福音 12:20-36
\newline
\begin{longtable}{cl}
\hline
\hline
章節 & 經文 (和合本修訂版)\\
\hline
12:20 & \begin{tabularx}{0.7\textwidth}{X} 那時，上來過節禮拜的人中，有幾個希臘人。 \end{tabularx} \\ \\ \relax
12:21 & \begin{tabularx}{0.7\textwidth}{X} 他們來見加利利的伯賽大人腓力，請求他說：「先生，我們想見耶穌。」 \end{tabularx} \\ \\ \relax
12:22 & \begin{tabularx}{0.7\textwidth}{X} 腓力去告訴安得烈，然後安得烈同腓力去告訴耶穌。 \end{tabularx} \\ \\ \relax
12:23 & \begin{tabularx}{0.7\textwidth}{X} 耶穌回答他們說：「人子得榮耀的時候到了。 \end{tabularx} \\ \\ \relax
12:24 & \begin{tabularx}{0.7\textwidth}{X} 我實實在在地告訴你們，一粒麥子不落在地裡死了，仍舊是一粒；若是死了，就結出許多子粒來。 \end{tabularx} \\ \\ \relax
12:25 & \begin{tabularx}{0.7\textwidth}{X} 愛惜自己性命的，就喪失性命；那恨惡自己在這世上的性命的，要保全性命到永生。 \end{tabularx} \\ \\ \relax
12:26 & \begin{tabularx}{0.7\textwidth}{X} 若有人服事我，就當跟從我；我在哪裡，服事我的人也要在哪裡；若有人服事我，我父必尊重他。」 \end{tabularx} \\ \\ \relax
12:27 & \begin{tabularx}{0.7\textwidth}{X} 「我現在心裡憂愁，我說甚麼才好呢？說『父啊，救我脫離這時候』嗎？但我正是為這時候來的。 \end{tabularx} \\ \\ \relax
12:28 & \begin{tabularx}{0.7\textwidth}{X} 父啊，願你榮耀你的名！」於是有聲音從天上來，說：「我已經榮耀了我的名，還要再榮耀。」 \end{tabularx} \\ \\ \relax
12:29 & \begin{tabularx}{0.7\textwidth}{X} 站在旁邊的眾人聽見，就說：「打雷了。」另有的說：「有天使對他說話。」 \end{tabularx} \\ \\ \relax
12:30 & \begin{tabularx}{0.7\textwidth}{X} 耶穌回答說：「這聲音不是為我，而是為你們來的。 \end{tabularx} \\ \\ \relax
12:31 & \begin{tabularx}{0.7\textwidth}{X} 現在正是這世界受審判的時候；現在這世界的統治者要被趕出去。 \end{tabularx} \\ \\ \relax
12:32 & \begin{tabularx}{0.7\textwidth}{X} 我從地上被舉起來的時候，我要吸引萬人來歸我。」 \end{tabularx} \\ \\ \relax
12:33 & \begin{tabularx}{0.7\textwidth}{X} 耶穌這話是指自己將要怎樣死說的。 \end{tabularx} \\ \\ \relax
12:34 & \begin{tabularx}{0.7\textwidth}{X} 眾人就回答他：「我們聽見律法書上說，基督是永存的；你怎麼說，人子必須被舉起來呢？這人子是誰呢？」 \end{tabularx} \\ \\ \relax
12:35 & \begin{tabularx}{0.7\textwidth}{X} 耶穌對他們說：「光在你們中間為時不多了，應該趁著有光的時候行走，免得黑暗臨到你們；那在黑暗裡行走的，不知道往何處去。 \end{tabularx} \\ \\ \relax
12:36 & \begin{tabularx}{0.7\textwidth}{X} 你們趁著有光，要信從這光，使你們成為光明之子。」 \end{tabularx} \\ \\
[1ex]
\hline
\hline
\end{longtable}
講題:生命的關鍵時刻:「時候到了!」　The Critical Time of Life: “The Hour Has Come Now”

經文:約翰福音12章20-36節

講員:許開明牧師

場合:香港中文大學崇基學院禮拜堂

日期:2015年3月22日

大能臨頭

引言

本主日是預苦期或大齋期第五主日,耶穌進入耶路撒冷,正是猶太人的逾越節,他知道他生命的關鍵時刻:「死亡的時候到了!」

今天與各位思想的題目,是每個人都會遇到的,生命的關鍵時刻:「時候到了!」

祈禱:

【】香港人都會記得粵語殘片的橋段:

一位犯人手腳被搏著,垂低著頭,另一位肥大的壯男人拿著大刀.

只聽見有把聲大叫「時辰已到!處斬!」

那肥壯男士舉起大刀,向著犯人的頸部準備手起刀落,快人頭落地.

犯人也合上眼,知道快被處決,可能吸一口氣,然後接受死亡.

就在這關鍵時刻!

突然聽見,聖旨到.「刀下留人!」

在在這刻,一切都改變了,犯人不但不用死,而且無罪釋放,峰迴路轉,場面感人.

人生總會遇上「時候到了!」的關鍵時刻!

然而,以上峰迴路轉的情景,是甚少的出現的!

反而,這些時刻,往往是驚心動怕,令人擔憂的!

病人,聽醫生的報告或做手術,「時候到了!」

運動員,決賽的「時候到了!」

人生,迎向死亡的「時候到了!」

你的感覺如何?緊張!恐懼!不知所措!

【搶救那一幅畫】

一百年前,法國一份報章向公眾提出了一個很有趣的問題,引起許多人的回應,這個問題至今仍然值得我們反思:

「假如法國最大的博物館羅浮宮失火了,在這關鍵的時刻,只允許你搶救出一幅畫,你會搶救那一幅?」

羅浮宮存放著許多世界知名的油畫,許多人就提出答案,其中有三幅是最多人選的,包括:「蒙娜麗莎」,「拿破崙一世加冕大典」,「自由神領導人民前進!」你估下,那一幅有可能呢?

在成千上萬的答案中,獲獎的答案是:?

「我搶救離出口最近的那一幅!」

答案原因:「那一幅是我最有把握搶救到的,雖然它可能比不上「蒙娜麗莎」,但放得在羅浮宮,一定是名貴的畫.」

答案的啟示:「在人生的關鍵時刻,好像失火的緊張時刻,要搶救一幅名畫,目標不是最有價值的那一幅,而是最有可能實現的那一個!」

【】在海潚中,一位父親和兒女跌下水,父親是救生員,在大風大浪,救了女兒,卻不救兒子.有何問:你只救女兒,不救兒子呢?是否你愛女情深!

他回答:「在危難中,我根本無法決定救仔或女,只能救最接近我的女兒!」

在危救,目標不是應該救那位,而是最有可能救到的那一個!

經文大要

本段的經文,耶穌基督也正面對著這幾個關鍵時刻!

我們來看耶穌如何面對,當中對我們有何啟迪?

1希臘人到訪人子的時刻到了:基督教信仰進入世界時刻到了

2人子得榮耀的時刻到了

3人子面對死亡的時刻到了

4人子讓上主舉起的時刻到了

5門徒趁著有光的時刻到了

耶穌面對著幾個關鍵的時刻

1希臘人到訪人子的時刻到了:基督教信仰進入世界的時刻到了

信仰傳到地極,到世界.法利賽人說:「…世上的人都隨從祂了.」(約11:19)在這時刻,福音廣傳,正進入世界了!

信仰進入知識,理性(抽象),圖像(想像):

福音給貧苦的人,給漁夫,給稅吏,給財主,但福音也給希臘人:有知識,有思想,有文化的人；好像保羅,路加醫生,約翰等人.

希臘人,代表求知欲強的人,代表愛旅行周圍探索知識的人.

中國人也早有此智慧:「讀萬卷書不如行萬里路」

希臘人是:讀萬卷書十行萬里路!(基督徒去過聖地後,睇聖經通了很多,我80年,93年去過,一次比一次深刻.)

基督福音要深入民間,應有其實用性；感受主的愛.

要深入人心,必有哲理性:救恩的原理.深入人心才能持久,直到永恆.

基督教進入人的腦(理性),也進入人的心(感情),人的身(經歷).

約翰福音顯示基督教的信仰受希臘人的重視,其信仰也隨之而豐富起來了!

希臘人到訪人子的時刻,這不獨是福音傳到地極的時刻,更重要的是福音得以進入人心,當代的人的心,千萬代人的心,直到永恆.

希臘人求見,門徒問耶穌:「他們想見你,可以嗎?」耶穌有無回答?好像「無」.

其實是「有!」；但,不是直接的回答,衪說:「若有人服事我,就當跟從我,我在那裡,服事我的人,也要在那裡.若有人服事我,我父必尊重他.」(約12:26)

耶穌知道希臘人求見,為要求知識,想獲得真理.

他提出三個人生關懷:榮與辱的人生,生與死,愛惜或不愛惜生命?

然而,耶穌更進一步顯示:你們見我,是出於好奇,我要求的卻是:不獨見我,跟隨我,服事我!還要委身於我,還要與我建立關係,進入希伯來所強調的人與上帝的I and Thou的親密的關係.

2人子得榮耀的時刻到了

人子釘十架是得榮耀,還是得羞辱呢?

當時的十架是羅馬人捉到叛國賊,被判有罪的刑罰,死也死在十架被曬,無埋葬之地,在人的眼中,是極羞辱的記號.

耶穌卻認為是榮耀的記號,因沒有死,便沒有復活,沒有勝過死亡,沒有榮耀.

在眾人以為羞辱的十架,耶穌卻堅持為榮耀的記號,更是不容易.他說:「人子得榮耀的時刻到了」(約13:24)

在社會上走群眾走的路線是寬路,堅持己見往往是一條窄路,並且別人否定你的觀點時,擇善固執,更是難得.

人子的榮耀觀與世人不同:

世人是要征服:勝過別人,得榮耀!耶穌在十架卻是失敗者,仍得榮耀.

世人要擁有權力,超越別人,得榮耀!耶穌放下權力,讓別人超越他,

卻得榮耀.

世人要擁有統治權,凌駕別人,得榮耀!耶穌甘心失去治權,給人凌辱,

卻得榮耀.

請注意耶穌所說的人子得榮耀,不獨要使天父得榮耀,且最終為要天父得榮耀.

並非為自已,而是為父神.他說:「父親阿,願你榮耀你的名.」(約13:28)

許多信徒求得榮耀,是個人得榮耀, 而非父上帝得榮耀.耶穌強調的,是「榮歸天父」.

3人子面對死亡的時刻到了

人子的抉擇:逃避死亡或接受死亡?

逃避死亡:不用再恐懼!

接受死亡:經歷恐懼及不可知的未來!

耶穌像一般人一樣,他一樣有如下的掙扎:

(約13:27)「」

但是接受死亡如種子,埋在地下結出子粒,得更多的生命!

從死復活,得長遠的,永恆的生命!

自然界的法則:「生命從死而來」- 置之死地而後生.

麥粒落在地裡死了,才能結出許多子粒來!

耶穌一人的捨命,而叫多人得著永生.

從自然界的法則:「生命從死而來」,引伸出「放棄自私」的人生法則.

人子的抉擇:愛惜生命或不愛惜生命(犧牲生命)?

錫身,為自已,滿足自己,保護自已,只為自己而活,使別人生命.

唔錫身,犧牲自我,成全他人,為別人而活!

錫身,為罪惡(金錢,賭博,情欲)而活,為無意義的東西而活.

4人子讓上主舉起的時刻到了

人子的抉擇:被上帝舉起(被釘十架)或不被舉起(不被釘十架)

但舉起是痛苦的!是面向死亡的,恐懼,不安,是走向死亡的.

然而舉起後吸引萬民!舉起也同樣經過死亡,步向復活.

在面對十架要被舉起時,耶穌也有掙扎,迎向痛苦十架,還是逃避十架?

一般人面對痛苦,羞辱,死亡採取什麼態度?逃避它們!

耶穌卻選擇面對,勇敢面對,因這是祂的使命,祂生命的終極目的.

面對困難,往往不是逃避它,而是勇敢地接受它,承認它的存在,然後勝過它!

【】正如一個人發現患癒症,逃避不是辦法,而是接受它,並且依靠神,藉著治療去勝過癌細胞.

【】面對自已的缺點,自己的臭脾氣,自己的弱點,盲點,是接受?勝過?

討論:人的命運,人的時日光陰到底在上帝的手,人的手,還是魔鬼的手?

對於人的時間,有時像在神的手,有時像在人的手,有時像在撒旦的手?

這最終決定於人對人生態度,或人的信念.

按照詩人的看法:「我終身的事都在你手裡」(詩31:15)

“My times are in your hand”(Psalm 31:15)[我所有的時間都在你手裡!]

若然時間在上主的手,生命也在上帝的手,人就應信靠上帝.

【繩子82】

一位登山者努力向上爬,為要趕盡快到目的地,不斷向上爬,爬到日落西山,在黑夜非常暗,伸手不見五指,禁不住驚慌起來,一不小心,快速地不斷向下跌,他想我死定了!

誰知以為會死,卻不死,腰間的繩拉住他,沒有跌落到谷底或岩石上,他被懸挂在空中,不斷有山上寒風吹來,繩子不斷搖晃,令人提心吊膽,他本能地向天求救:「上帝呀,救命呀!」

有聲音說:「你相信我會救你嗎?」

他說:「我相信,絕對相信!」

那聲音說:「那你就剪斷你條繩吧!」

那登山者一直沉默著,他無法放棄那救他命的那條繩.

他相信那條繩,多過那聲音!

結果,第二天,登山隊發現,有一具凍僵了的屍體懸掛在半空中,

但令人惋惜的是,這屍體離地面只有十尺.

如果這發山者相信那聲音,大膽地剪斷繩子,跳下來便不會死了!

信仰要敢於冒險,這是許多基督徒的經歷.

一位神的僕人說:「我們要大著膽相信耶穌,大著膽跟隨耶穌,大著膽做人.」

5門徒趁著有光的時刻到了

趁著有光,免得黑暗勝過你們.

古時的人,無有發明電燈照明,晚間基本上在黑暗中不能工作.

他們要趁著白畫,有光的時候,有機會的時候就要工作.

生命要趁著有光,趁著有生命的機會,趁著可工作的機會,

一個生命的成功與否,在於他能否懂得趁著生命的機會,善用生命的機會?

你是否趁著在大學讀書的機會,讀好書,識多D朋友?過好大學的生活?

你是否趁著有機會事奉上帝,作好事奉?

作父母的,你是否趁著孩子還會跟你的機會,一起玩耍,一起去旅行,一起上禮拜堂.

耶穌基督趁著在世上生命的機會,傳講上帝的國,建立上帝的國,使人得真理,得永生,得救恩!

黑暗勝過人,黑暗是一種勢力, 使人看不見,使人恐懼,使人無知,使人不敢向前,甚至使人被操縱.

【】中國的貪污問題嚴重,現政府打蒼蠅也打老虎.其中軍隊的高層的腐貪所作所為, 對軍隊造成了很大內傷,內地軍事專家,退役少將羅援近日接受內地傳媒訪問時,說了句值得中國人民反思的話:他說,「目前幾乎沒有哪個國家有能力和膽量打敗中國,唯一能打敗我們(中國)的就是我們自己；腐敗不除,未戰先敗,腐敗是軍隊戰鬥力的第一號殺手(頭號的敵人)」.打擊貪腐,也是打擊黑暗,向光而行!

趁著有光,把握機會:不知道機會?求主賜智慧,順從之!

關鍵時刻有三個特點:

1臨界點

a是由生到死的臨界點!

或由死到生!

b是羞辱到榮耀的臨界點

或由近到遠

c由舉起到埋下的臨界點

2吊詭與抉擇

生與死的抉擇

苦與樂的抉擇

近與遠的抉擇

成與敗的抉擇

3態度問題?

你人生的態度,決定一切:正面或負面?

a榮與辱?

b接與拒?

正面態度!勇敢面對!

死就事死罷!

不怕死,不怕失敗!

[]呆若木雞”dumb as a wooden chicken”,好嗎?

前以為是不好的,「呆」像「木雞」,怎會好?

三個境界:

《莊子,達生》,記一名紀渻子,為齊王養鬥雞:

a十天後,齊王問:鬥雞養好了嗎?

紀說:還沒有,所養鬥雞盛氣凌人!見雞就打!極為慌張!不能鬥!

b再十天後,齊王問:鬥雞養好了嗎?

紀說:還沒有,鬥雞只要聽到雞嗚,雞影,就要鬥,神經緊張.還不能鬥!

c再過十天,齊王問:鬥雞養好了嗎?

紀說:可以了!他見雞不動,見雞影不動,雞鳴也不動,不為所動,簡直是[呆若木雞].充滿自信,對其他鬥雞,視若無賭,勇往直前,無懼成與敗,其他雞一見這鬥雞都不敢向前,反而逃跑.

耶穌面對十字架,時候到了的時刻,就是到這種「呆若木雞」的地步了!沉著,靜默,鎮定,不驚不慌,勇往直前,面向十架!面向人生的關鍵時刻!

各位:你又如何?

可曾想過如果有一天臥在床上,知道生命的燭光快要熄滅了,你會如何?

結束祈禱:

掌管生命的主宰,人生總有關鍵的時刻,時辰到了的時候!願我們能學效主耶穌,沉著面對,作出合宜的抉擇,並以積極的態度面對,憑著信心前進.奉主名求.阿們.



\newpage

\section{騎驢進京殺機四伏}
\label{sec:281}
\textbf{賀志勇博士 - 騎驢進京,殺機四伏}
\newline
\newline
\hyperref[sec:280]{< < < PREV SERMON < < <}
~
\hyperref[sec:toc]{[返目錄]}
~
\hyperref[sec:282]{> > > NEXT SERMON > > >}


馬可福音 11:1-11,15-18
\newline
\begin{longtable}{cl}
\hline
\hline
章節 & 經文 (和合本修訂版)\\
\hline
11:1 & \begin{tabularx}{0.7\textwidth}{X} 耶穌和門徒快到耶路撒冷，來到伯法其和伯大尼，在橄欖山那裡。耶穌打發兩個門徒， \end{tabularx} \\ \\ \relax
11:2 & \begin{tabularx}{0.7\textwidth}{X} 對他們說：「你們往對面村子裡去，一進去的時候會看見一匹驢駒拴在那裡，是從來沒有人騎過的，把牠解開，牽來。 \end{tabularx} \\ \\ \relax
11:3 & \begin{tabularx}{0.7\textwidth}{X} 若有人對你們說：『為甚麼做這事？』你們就說：『主要用牠，但會立刻把牠牽回到這裡來。』」 \end{tabularx} \\ \\ \relax
11:4 & \begin{tabularx}{0.7\textwidth}{X} 他們去了，看見一匹驢駒拴在門外街道上，就把牠解開。 \end{tabularx} \\ \\ \relax
11:5 & \begin{tabularx}{0.7\textwidth}{X} 在那裡站著的人，有幾個說：「你們解開驢駒做甚麼？」 \end{tabularx} \\ \\ \relax
11:6 & \begin{tabularx}{0.7\textwidth}{X} 門徒照著耶穌的話說，那些人就任憑他們牽去了。 \end{tabularx} \\ \\ \relax
11:7 & \begin{tabularx}{0.7\textwidth}{X} 他們把驢駒牽到耶穌那裡，把自己的衣服搭在上面，耶穌就騎上。 \end{tabularx} \\ \\ \relax
11:8 & \begin{tabularx}{0.7\textwidth}{X} 有許多人把衣服鋪在路上，還有人把田間的樹枝砍下來鋪上。 \end{tabularx} \\ \\ \relax
11:9 & \begin{tabularx}{0.7\textwidth}{X} 前呼後擁的人都喊著說：「和散那！奉主名來的是應當稱頌的！ \end{tabularx} \\ \\ \relax
11:10 & \begin{tabularx}{0.7\textwidth}{X} 那將要來的我祖大衛之國是應當稱頌的！至高無上的，和散那！」 \end{tabularx} \\ \\ \relax
11:11 & \begin{tabularx}{0.7\textwidth}{X} 耶穌到了耶路撒冷，進入聖殿，看了周圍的一切。天色已晚，他就和十二使徒出城，往伯大尼去。 \end{tabularx} \\ \\ \relax
11:12 & \begin{tabularx}{0.7\textwidth}{X} 第二天，他們從伯大尼出來，耶穌餓了。 \end{tabularx} \\ \\ \relax
11:13 & \begin{tabularx}{0.7\textwidth}{X} 他遠遠地看見一棵無花果樹，樹上有葉子，就過去，看是不是在樹上可以找到甚麼。他到了樹下，竟找不到甚麼，只有葉子，因為不是無花果的季節。 \end{tabularx} \\ \\ \relax
11:14 & \begin{tabularx}{0.7\textwidth}{X} 耶穌就對樹說：「從今以後，永沒有人吃你的果子。」他的門徒都聽到了。 \end{tabularx} \\ \\ \relax
11:15 & \begin{tabularx}{0.7\textwidth}{X} 他們來到耶路撒冷。耶穌一進聖殿，就趕出在聖殿裡做買賣的人，推倒兌換銀錢之人的桌子和賣鴿子之人的凳子； \end{tabularx} \\ \\ \relax
11:16 & \begin{tabularx}{0.7\textwidth}{X} 也不許人拿著器具從聖殿裡經過。 \end{tabularx} \\ \\ \relax
11:17 & \begin{tabularx}{0.7\textwidth}{X} 他教導他們說：「經上不是記著：『我的殿要稱為萬國禱告的殿嗎？你們倒使它成為賊窩了。』」 \end{tabularx} \\ \\ \relax
11:18 & \begin{tabularx}{0.7\textwidth}{X} 祭司長和文士聽見這話，就想法子要除掉耶穌，卻又怕他，因為眾人都對他的教導感到驚奇。 \end{tabularx} \\ \\
[1ex]
\hline
\hline
\end{longtable}
講題:騎驢進京,殺機四伏　Riding into the Holy City, Surrounded by Dangers

經文:馬可福音11:1-11,15-18

講員:賀志勇博士

場合:香港中文大學崇基學院禮拜堂

日期:2015年3月29日

1耶穌和門徒將近耶路撒冷,到了伯法其和伯大尼,在橄欖山那裡.耶穌就打發兩個門徒,2對他們說,你們往對面村子裡去.一進去的時候,必看見一匹驢駒拴在那裡,是從來沒有人騎過的.可以解開牽來.3若有人對你們說,為什麼作這事.你們就說,主要用它.那人必立時讓你們牽來.4他們去了,便看見一匹驢駒,拴在門外街道上,就把它解開.5在那裡站著的人,有幾個說,你們解驢駒作什麼.6門徒照著耶穌所說的回答,那些人就任憑他們牽去了.7他們把驢駒牽到耶穌那裡,把自己的衣服搭在上面,耶穌就騎上. 8有許多人,把衣服鋪在路上,也有人把田間的樹枝砍下來,鋪在路上.9前行後隨的人,都喊著說,和散那,(和散那原有求救的意思,在此乃是稱頌的話)奉主名來的,是應當稱頌的.10那將要來的我祖大衛之國,是應當稱頌的.高高在上,和散那.11耶穌進了耶路撒冷,入了聖殿,周圍看了各樣物件.天色已晚,就和十二個門徒出城往伯大尼去了.

第一 引言:騎什麼很重要麼?

我小時候喜歡看動畫片(Cartoon),記憶很深的就是哆啦A夢(大陸又叫機器貓),大家有沒有看過啊?裡面有可愛又懶惰的野比大雄(大陸又叫康夫),有美麗又善良的靜香(大陸又叫小靜).大雄暗戀靜香,但是大雄很膽小,經常要哆啦A夢拿出法寶來幫忙追靜香,所以經常鬧笑話.有一次,大雄想到靜香的夢裡面看看靜香對自己印象怎麼樣；靜香正夢見有怪獸來抓自己,然後就見大雄(夢裡面就是王子)騎著一匹東西來救她.靜香正在高興,但是,不對不對,突然間靜香吃驚地喊道:“大雄,你怎麼騎著豬來啦?你難道沒有白馬麼?”“太生氣啦!騎著豬還敢來救我!”最後,靜香在自己的夢裡拋下大雄,跟一個騎著白馬的王子走了.

這個故事告訴我們,騎什麼是很重要的,騎錯了東西後果很嚴重!今天我們要看的經文是馬可福音11:1-11節,講的就是耶穌基督騎驢進京的故事.

第二,解經

列位,如果你們查看其他幾卷福音書,你們會發現,其實主耶穌耶路撒冷已經來過好幾次了,都是悄悄地來的,那這次,他為什麼要大張旗鼓呢?讓我們像看電影一樣,再一次重溫這一段騎驢旅程吧.

1,  經文:1節

從第1節,讓我們看看這次進京的出發地點--伯法其和伯大尼,伯大尼有什麼特別之處呢?哦!有伯大尼的馬大和瑪利亞,這兩位姊妹有個非常有名的弟弟--死裡復活的拉撒路.在約翰福音中,拉撒路復活是耶穌基督在進入耶路撒冷之前的最後一個神跡,但約翰福音12章提到,祭司們不僅要殺耶穌,連拉撒路也要殺,“因有好些猶太人為了拉撒路的緣故,回去信了耶穌”.(約12:11)這樣看來,馬可不是隨隨便便提到伯大尼,在馬可的敘述背後,不動聲色地隱藏著巨大的殺機；也就是說,馬可在這裡暗示了:主耶穌其實知道,這將是他最後一次踏入這城門,而之後,他將被作為一個“罪犯”,在自己的母親˴弟子,和世人面前被鞭打,被羞辱,直到被釘死.

這次進京的時間選在逾越節前5天,在最後的晚餐之前4天.據一位羅馬總督的資料講,那段時間耶路撒冷聚集了大概250萬從世界各地去的人,要宰殺大概25萬頭羊.每次我想到殺羊的場面:羊被捆起來,頭被抬起來,一刀下去,鮮血飛起來,那個場面太血腥太殘忍,何況要殺的是25萬頭羊!而主進城的那一周,殺羊在不斷進行.每一次鮮血濺起,都預表著更血腥的事件--耶穌基督,逾越節的羔羊,被釘死,被高高掛在十字架上. 如果這是一幕電影,那我們會從這個時間地點嗅到很濃重的血腥味--殺機四伏!

這一幕一幕,讓我想起了亞伯拉罕獻以撒的故事.獻以撒的地方,正正就是當年的耶路撒冷,所羅門聖殿的中心.當年,亞伯拉罕聽從耶和華的命令,把以撒帶到這裡,捆起來,在孩子的哀呼聲中,手起,刀就要落......何等的驚心動魄,何等相似的場景呢!

也許信仰就是這樣.有時候,信仰並不僅僅是“在美麗的教堂做做禮拜,聽聽美妙的音樂,哈,我心一片平安”這麼簡單；相信一件事情是要付代價的；我們不會一邊說“醫生,我相信你醫術很好”,一邊拒絕這個醫生給自己做手術.相信醫生,就要把自己託付給醫生,上手術臺,哪怕我們知道手術有風險!託付給醫生就是代價.同樣,信耶穌也是要付代價的；這個代價有時候甚至會血淋淋!就像亞伯拉罕獻以撒,就像耶穌基督這一次進城背後的故事...

2,2-7節 預備

撒迦利亞曾經預言,主將站在橄欖山上打敗那些攻打耶路撒冷的人.(亞14:1-5) 但在2-7節我們看到,我們的主耶穌到橄欖山並沒有打仗,卻是來找驢駒.如果這是一幕電影,你會發現整個畫面非常寧靜.耶穌吩咐門徒去取驢駒過來,門徒按著吩咐去做了,那些人任憑他們牽走了驢駒.主就上了驢駒.整個畫面很日常,沒有波瀾,沒有起伏.但熟悉電影的人往往知道,平靜的畫面往往預示了更大的不平靜.

驢駒就是驢的崽子,不大.有人問我,主為什麼要找一匹驢駒呢?人家這麼小,太辛苦了吧!他雖然不認為主耶穌虐待動物,但是他奇怪主耶穌為什麼不騎大驢子呢?從馬太福音看,明明驢駒的驢媽媽就在旁邊嘛!據撒母耳記6:7節說明:背負耶和華東西的必須是未曾負軛的,原來耶穌選驢駒是為了完成神聖潔淨的禮儀；更重要的是,先知曾經說過:“看哪,你的王來到你這裡, 是溫柔的,又騎著驢,就是騎著驢駒子.”(賽62:11)原來主耶穌這樣做是要實現先知的預言.他不但完成了律法的要求,他更完成了天父的心意.在神導演的拯救世人的大戲中,我們的主角耶穌基督一絲不苟地完成了天父的計畫；他從橄欖山出發,用他的順服,對耶路撒冷發動了攻擊.

在這場大戲中還有一個小人物--驢駒的主人.或許他不明白主為什麼要這匹驢駒,但主吩咐了,他就這樣做了.信仰就是這樣,信仰就是一種順服.撒母耳說:聽命勝於獻祭,順從勝於公羊的脂油.(撒上15:22) 驢駒的主人並不明白主究竟要驢駒做什麼,但是他以他安靜的順服來實踐他的信仰；而主耶穌非常明白前路的可怕殺機,但他以他安靜的順服來完成天父的安排-- 所以說,信心沒有什麼花巧可言,無論最後你明白或不明白,無論代價有多大,因為主這樣吩咐,所以你就順服了,這就是信心.

3,8-11節進城

好!馬上就要進城了,就要進城了,正在進城了.鏡頭一點點拉近.耶穌基督這次騎驢進京,看起來有些聲勢浩大,有點像七一遊行.你看8-10節寫著:隨行的人“有的把衣服鋪在路上,有的鋪樹枝在路上,還有人高喊‘高高在上和散那!’”

如果你是圍觀的猶太百姓,你可能會議論:“啊,這有點像馬加比起義勝利之後榮耀進城的場面,我的先祖經歷過哦!但是,不對,他怎麼騎著這麼小的驢子!你看,這像啥樣子啊,驢子太小啦,腳都點著地啦!”

如果你是羅馬士兵,你可能會說:“猶太人搞什麼鬼,騎著這麼小的驢子不說,還又搖樹枝又鋪衣服的.(馬太福音21章)這哪能跟我們羅馬人進城比.咱們是騎著高頭大馬,鋪著紅地毯,兩邊吹著號迎接,多帥氣,多威武!”

如果你是祭司,你可能會說:“還奉主名來了,這麼小的驢子!讓你得意,讓你煽動百姓.讓你不得好死!”

騎什麼很重要!你騎著高頭大馬跟騎著小毛驢,檔次差得不是一點兩點.這就跟你是開著計程車,還是開著勞斯萊斯,差別大得很啊!世界看重的是白馬王子,騎著豬來的怎麼會是王子呢?同樣,騎著驢駒來的怎麼會是王呢? 可是聖經說:“你的王來到你這裡,是溫柔的,又騎著驢,......”(亞9:9)先知預言了,主耶穌將會徹底顛覆我們對榮耀的理解!榮入聖城,不在於主的聲勢,不在於華貴風光,而在於他的謙卑˴溫柔,和明知要死卻堅定赴死的順服!這次進城,耶穌想用他的舉動,最後一次努力去喚醒從世界各地來的人--去想想先知的話語吧,去明白上帝的旨意吧!

可是主做的這一切似乎有點徒勞.不過幾天,這群高聲呼喊著“和散那”的人就會對他失望,離他而去--因為主給不了他們想要的東西.今天呼喊“和散那”的人,很有可能就是過幾天喊著“釘死他,釘死他”的人.因為他們失望了--這個王沒有給他們帶來權力˴富貴,或者是自由˴民主˴解放；顯給他們看的卻是無能˴順服與屈辱,這樣的人怎麼會是“奉主名來的”,這樣的人怎可能是彌賽亞?釘死他,他褻瀆了上帝!

之後的故事我們都很清楚,耶穌潔淨聖殿,祭司長和文士就想法子除滅耶穌--殺機四伏.

結論:殺機中的生機

受難節你會想到什麼?每一次受難節,我都會想到這段榮入聖城.當主耶穌溫柔地騎著驢駒進入城門時,歡迎主耶穌的人高喊著“和散那”,“和散那”就是“求主拯救”的意思,但是他們不知道,他們必須真正委身於主,他們必須肯為這個呼喊付出代價,這一呼喊才會蒙應允.主耶穌常常避開群眾,他不信任他們.他知道有人在高呼讚美的時候,轉過頭就會離棄他.

弟兄姊妹們,我們會不會在教堂裡熱情洋溢地敬拜主,出了這門就忘了祂呢? 我們現在以輕鬆的心情閱讀這段經文,可能因為我們並不覺得這個世界殺機四伏.我們沒有遭遇過辦公室政治,我們可能也無需在星期天上班,日子過得也還不錯,所以我們不太會覺得殺機四伏.我們喜歡的名利˴地位,或者工作˴成就,甚至是民主˴自由,這些都是好東西,但如果我們不能用主耶穌的方式來對待它們,它們就慢慢佔據了我們的時間˴熱情˴智慧.這就是殺機,因為它奪取的不是肉體,卻是忠誠--對主的忠誠.在香港這個都市,最可怕的殺機其實是忙碌,它是一種慢性毒藥,奪取我們對主的忠誠,我們還安慰自己說:“其實我也很愛主,只是太忙了.....”

但這段經文再一次提醒我們,面對溫柔一刀的世界,我們沒有別的辦法,我們只有順服並委身于神,付上代價,我們才能夠得拯救.正如主耶穌所說的:“ 凡要救自己生命的,必喪掉生命；凡為我和福音喪掉生命的,必救了生命.”(馬可8:35-37)

主耶穌這句話絕不是空口無憑,因為復活節也不遠了.如果我們不願意做出改變,我們就無從體驗復活.而對於真正的門徒來講,殺機四伏並不可怕,正如詩人荷爾德林所說的:“哪裡有危險,哪裡就生長著解救之力.”



\newpage

\section{願你們平安}
\label{sec:282}
\textbf{邢福增教授 - 願你們平安}
\newline
\newline
\hyperref[sec:281]{< < < PREV SERMON < < <}
~
\hyperref[sec:toc]{[返目錄]}
~
\hyperref[sec:283]{> > > NEXT SERMON > > >}


約翰福音 20:19-31
\newline
\begin{longtable}{cl}
\hline
\hline
章節 & 經文 (和合本修訂版)\\
\hline
20:19 & \begin{tabularx}{0.7\textwidth}{X} 那日（就是七日的第一日）晚上，門徒因怕猶太人，所在的地方門都關了。耶穌來，站在當中，對他們說：「願你們平安！」 \end{tabularx} \\ \\ \relax
20:20 & \begin{tabularx}{0.7\textwidth}{X} 說了這話，他把手和肋旁給他們看。門徒一看見主就喜樂了。 \end{tabularx} \\ \\ \relax
20:21 & \begin{tabularx}{0.7\textwidth}{X} 於是耶穌又對他們說：「願你們平安！父怎樣差遣了我，我也照樣差遣你們。」 \end{tabularx} \\ \\ \relax
20:22 & \begin{tabularx}{0.7\textwidth}{X} 說了這話，他向他們吹一口氣，說：「領受聖靈吧！ \end{tabularx} \\ \\ \relax
20:23 & \begin{tabularx}{0.7\textwidth}{X} 你們赦免誰的罪，誰的罪就得赦免；你們不赦免誰的罪，誰的罪就不得赦免。」 \end{tabularx} \\ \\ \relax
20:24 & \begin{tabularx}{0.7\textwidth}{X} 那十二使徒中，有個叫低土馬的多馬，耶穌來的時候，他沒有和他們在一起。 \end{tabularx} \\ \\ \relax
20:25 & \begin{tabularx}{0.7\textwidth}{X} 其他的門徒就對他說：「我們已經看見主了。」多馬卻對他們說：「除非我看見他手上的釘痕，用我的指頭探入那釘痕，用我的手探入他的肋旁，我絕不信。」 \end{tabularx} \\ \\ \relax
20:26 & \begin{tabularx}{0.7\textwidth}{X} 過了八日，門徒又在屋裡，多馬也和他們在一起。門都關了，耶穌來，站在當中，說：「願你們平安！」 \end{tabularx} \\ \\ \relax
20:27 & \begin{tabularx}{0.7\textwidth}{X} 然後他對多馬說：「把你的指頭伸到這裡來，看看我的手；把你的手伸過來，探入我的肋旁。不要疑惑，總要信！」 \end{tabularx} \\ \\ \relax
20:28 & \begin{tabularx}{0.7\textwidth}{X} 多馬回答，對他說：「我的主！我的神！」 \end{tabularx} \\ \\ \relax
20:29 & \begin{tabularx}{0.7\textwidth}{X} 耶穌對他說：「你因為看見了我才信嗎？那沒有看見卻信的有福了。」 \end{tabularx} \\ \\ \relax
20:30 & \begin{tabularx}{0.7\textwidth}{X} 耶穌在他門徒面前另外行了許多神蹟，沒有記錄在這書上。 \end{tabularx} \\ \\ \relax
20:31 & \begin{tabularx}{0.7\textwidth}{X} 但記載這些事是要使你們信耶穌是基督，是神的兒子，並且使你們信他，好因著他的名得生命。 \end{tabularx} \\ \\
[1ex]
\hline
\hline
\end{longtable}
講題:願你們平安　Peace Be With You

經文:約翰福音20章19-31節

講員:邢福增教授

場合:香港中文大學崇基學院禮拜堂

日期:2015年4月12日

最近亞視不獲續牌的消息,令不少人開始懷緬這些電視台.我記得在麗的時代,劇集「人在江湖」有膾炙人口的插曲「人在旅途灑淚時」.為何這插曲那麼受歡迎?除了因為歌曲悅耳外,相信歌詞也道出了許多人的心聲.「明白世途多麼險阻,令你此時三心兩意,看遠路,正漫漫,誰謂抉擇最容易.前路渺茫請君三思,問你可曾心生悔意,要退後,也恨遲,人在旅途灑淚時.」許多人都用旅途來比喻人生,這條路從出生走到死亡,旅途上充滿各種的挑戰.有些人儘管經歷各種難處,但生命仍活得精彩.也有些人在旅途上迷失了方向,不知前面該走的路如何,不得不灑下眼淚.特別是當他們曾致力奮鬥的理想,或者是持守的信念與價值幻滅的時候,生命頓然變得沒有意義.正因如此,有人只能活在自己所搭建的謊言世界之內,並將之視作真理,就好像國王的新衣中那位國王,活在「謊言」之中,總比被那個天真的小孩道出真相為好.又有人裝睡,無論如何,都不想被別人叫醒.因為面對現實,承受真相往往是個教人感到痛苦不堪的歷程.

今天,我們會興奮地記念主耶穌的復活.但大家有沒有想過,在第一個復活節當日,「復活」卻只不過是個「謊言」而已.對那些要把耶穌置諸死地的祭司

和文士而言,復活無疑是個謊言(他們擔心的是門徒把屍體偷去〔太27:64〕).但對於跟從耶穌三年多的門徒而言,這何嘗又不是一個幻想?我甚至懷疑,這些

門徒是否仍記得主耶穌曾預言人子被釘在十字架,第三日要復活(太20:18-19)這句話?坦白說,門徒對耶穌的復活,根本沒有一絲奢望,他們甚至沒有企圖借

借「復活」的信念來凝聚自己,例如真的計劃偷走耶穌的屍體,編一個關於復活的謊言,藉著謊言來延續自己的運動(他們沒有這個膽量,也沒有這種領袖才能),反倒是處於迷失之中,沒有方向,只有恐懼.約翰福音記載,當時門徒懼怕猶太人想斬草除根,進一步逼害他們,所以一直躲在屋內把門緊緊關上.其實,祭司長等人根本沒有把這些門徒放在眼內,十二個門徒中,先是猶大出賣了耶穌,接著彼得又公開否認認識耶穌.當耶穌走向各各他十字架之路途中(教會傳統視稱作「苦路」),我們看不見門徒,那時代替耶穌背負十字架的是古利奈人西門,而為耶穌哀哭的,只有一些婦女(路23:26-28).其中一站苦路是說有一名女子幫助耶穌擦去臉上的血.耶穌死後,鼓起勇氣去領回耶穌屍體,是亞利馬太來的約瑟.至於把耶穌按著猶太人規矩來安葬的,是尼哥德慕.甚至首先去墳墓的,是婦女(抹大拉的馬利亞等),而不是門徒.敢問:耶穌的門徒在哪裡?唯一的答案是:在耶穌被捕以後,門徒早就潰不成軍了!

門徒畏縮地躲在一起,除了害怕受耶穌的牽連外,也是極其難過的.馬可福音記載,當婦女告訴門徒耶穌復活的消息時,他們「正哀慟哭泣」(可16:10).

這也難怪,那位呼召他們的主已被釘死了.他們當然會為耶穌的死而難過,畢竟大家相處三年,這是人之常情.也許,又會有人深切為自己在這星期的表現,特別是表現得如斯不濟而難過,因而進一步懷疑自己,否定自己.但他們內心也是充滿著許多的疑問,甚至我們不排除有人會覺得被耶穌欺騙與出賣.三年來,他們「撇下一切」,跟從耶穌,作他的門徒,最後換來甚麼?換了是你,內心也難掩這種不滿的情緒.耶穌的死,意味著過去三年的日子是走錯了路,押錯了注,真是悔不當初.門徒不僅要面對是否要否定自己的過去,更對自己的現在及將來充滿疑惑!他們迷失了方向,沒有了意義,人生的痛苦與悲哀莫過於此.所以,當他們聽到耶穌復活的消息時,福音書描寫他們最直接的反應就是:「以為是胡言,就不相信」(路24:11),「可是不信」(可16:11).就在這所滿佈著恐懼,迷失與疑惑的房子裡,主耶穌向門徒顯現,並向他們說:「願你們平安」.「願你們平安」是主顯現時最常說的一句話(太28:9,路24:36,約20:19,21,26).還記得在數天前,就是逾越節晚餐後,主耶穌曾跟門徒說:「我留下平安給你們,我把我的平安賜給你們.我所賜給你們的,不像世人所賜的；你們心裡不要憂愁,也不要膽怯.」(約14:27)在耶穌顯現前,如果門徒仍記得耶穌所說的話,一定感到莫大的諷刺,主真的留下平安嗎?沒有!他們所感受到的,只有憂愁與膽怯!那麼,對於當時那群面對著恐懼,迷失與疑惑的門徒而言,顯現後的耶穌,「平安」到底有何意義.第一,平安讓門徒能夠走出恐懼.當主耶穌在這間充滿恐懼的房子出現時,說了第一句「願你們平安」,然後就把手和肋旁的傷痕指給門徒看.為甚麼復活的主還留著釘十字架時的傷痕?理論上祂絕對可以變得完美(祂可以不受時空間的限制,突然出現,又突然消失),但是主卻選擇可以看得見,甚至是摸得著的傷痕作為祂復活的記號,最簡單的原因,就是主耶穌要證明那位曾被釘在十字架上受苦至死的人子,如今已勝過死亡,從死裡復活了──「傷痕」就是最好的證明.面對復活的主,門徒「就喜樂了」(約20:20).門徒從恐懼到喜樂,關鍵就是主真的復活,這位死而復活的主也就是平安與喜樂的源頭.耶穌身上的傷痕,無疑是耶穌被釘死的記號,但如今,又成為他復活的標記.

詩歌《因祂活著》其中一段歌詞是:「因衪活著,我能面對明天,因衪活著,不再懼怕；我深知道,他掌管明天,生命充滿了希望,只因衪活著.」對,若有主耶穌與我們同在,我們就不必害怕!因為整個世界都改變,我們的主永不改變.基督是復活的主,他永遠活著.主已經復活了!祂第一句跟門徒說的話就是「願你們平安」.祂賜給我們世人不能給予的平安,就是復活的生命.第二,是平安與使命的問題.主耶穌向門徒第二次說「願你們平安」後,接著說:「父怎樣差遺了我,我也照樣差遺你們」(約20:21).這話其實耶穌在客西馬尼園的「分離禱告」中也曾向門徒揭示:「你怎樣差我到世上,我也照樣差他們到世上」(約17:18).耶穌把一幅父－子－門徒的圖畫呈現出來,現在子已經完成父所交託的使命,門徒就是子的接棒者,主耶穌向門徒顯現,就是差遺他們,繼續完成上帝救贖人類的工作.不是吧!這群已經潰不成軍,活在恐懼與迷失的門徒,竟然是基督的接棒者,主耶穌是否所託非人呢?其實,門徒的迷失,是因為他們在耶穌死後失去了方向,看不到前面的路.如今復活的主向他們託付,正是要為門徒的未來賦予使命與意義.如果門徒在耶穌死後懷疑甚至否定過去三年,甚至認為自己原來是活在謊言之中時,復活的主正要揭示,祂就是基督,上帝的兒子.主耶穌一直堅守自己的使命,沒有把父差派的任務丟棄,如今復活的主就要把這重大的責任交給門徒.今後,門徒的生命不再是活在謊言與迷失之中,而是要把自己的生命與主所託付的責任與使命結合.主耶穌差遺他們,又向他們吹了一口氣,說:「領受聖靈吧.你們赦免誰的罪,誰的罪就得赦免；你們不赦免誰的罪,誰的罪就不得赦免」(約20:22-23)一方面,耶穌回應了逾越節晚餐應許會賜下聖靈一事,聖靈是耶穌職事的繼承者.耶穌向門徒吹氣,就好像上帝創造時向亞當吹了一口氣一樣,讓門徒能得著新生命.另方面,耶穌授予門徒有赦罪的權柄,這正是門徒承接使命後,按照福音的大能而獲得的權柄.平安不僅是安慰,也是順服接受差遺與使命,並且授予權柄.這一群退縮的門徒,因著主的復活,重新堅定信心,並將繼續為福音使命而奮鬥.我們看見,躲在門後迷失的門徒,後來在使徒行傳中,成為一群勇敢見證,傳揚福音的人,甚至不懼死亡,以生命來見證復活的主.是甚麼讓他們改變過來?關鍵就是從這位復活的主,賜予平安的主,領受使命與權柄.今天,福音使命的棒已通過歷世歷代的信徒傳到我們手裡,我們是否因著不同的原因而感到迷失與乏力呢?在這些難關之中,復活的主所賜下的平安,能否讓我們從憂愁與膽怯中得到勇氣與力量?也許,我們會看見許多代表軟弱與挫敗的「傷痕」,但同樣是這些傷痕,可以讓我們看見復活主的臨在.

第三,是平安與相信.原來耶穌這次顯現時,多馬並不在場.當門徒向多馬說:「我們已經看見主了」,並向他陳述有關情況,特別是主把手和肋旁的傷痕都指給門徒看.多馬的反應是疑惑與不信:「除非我看見他手上的釘痕,用我的指頭探入那釘痕,用我的手探入他的肋旁,我絕不信」(約20:25).就在八日後,耶穌再次向門徒顯現,這次多馬也在場了.門也是關上,但主耶穌突然出現,向眾人說:「願你們平安!」然後他特別向多馬說:「把你的指頭伸到這裡來,看看我的手；把你的手伸過來,探入我的肋旁.不要疑惑,總要信」.不過,在復活的主面前,摸不摸主的傷痕已經不再重要了.多馬高呼:「我的主,我的上帝!」多馬顯然徹底降服下來了!接著主耶穌說:「你因看見了我才信嗎?那沒有看見卻信的有福了」.這句話是甚麼意思呢?其實,多馬並不比眾門徒差,因為其他門徒第一次聽到婦女說主復活時,也是抱著懷疑的態度.第一批見證主復活者,都是因看見才相信的.所以,耶穌這話的目的,不是要斥責多馬的小信.僅得留意,主耶穌復活後的顯現,都是向著他的跟隨者與門徒.為何耶穌不向更多原來就反對他的人顯現,例如他可向彼拉多,或到公會,在定他罪的人面前大顯威風.但是他並沒有這樣作,福音書從沒有記載耶穌向不信他的人顯現.顯然,耶穌從復活到升天的六個星期內,最大的目的是要建立門徒,能夠親眼看見復活的主,實在是莫大的福份.他深悉自己升天後,便不再顯現,要令那些「沒有看見的人」相信他,就只能靠這群見證人.保羅在哥林多前書十五章討論基督的復活時這樣說:「我當日所領受又傳給你們的,最重要的就是:照聖經所說,基督為我們的罪死了,而且埋葬了；又照聖經所說,第三天復活了,還顯給磯法看,又顯給十二使徒看,後來一次顯給五百多弟兄看,其中一大半到現在還在,卻也有已經睡了的.以後他顯給雅各看,再顯給眾使徒看,最後也顯給我看；我如同未到產期而生的人一般.我原是使徒中最小的,不配稱為使徒,因為我曾迫害過上帝的教會.然而,由於上帝的恩典,我才成了今日的我,並且他所賜給我的恩典不是徒然的.我比眾使徒格外勞苦；其實不是我,而是上帝的恩典與我同在.無論是我或是其他使徒,我們都如此傳,你們也都如此信了.」(林前15:3-11)「我們都如此傳,你們也都如此信了」──保羅這句話正是歷世歷代每位基督徒的寫照.不過,無論是有機會親眼看見復活的主的,或是因著這些見證人所傳而相信的,都必須與主有第一手相遇的經歷.就好像多馬一樣,因著與主相遇而放下自我──「我的主,我的上帝!」我們如果能夠經歷主,就不容易因眼前的疑惑而否定主的真實,進而能夠感受到祂同在的「平安」.弟兄姊妹,你的信仰的生命中,是否也曾因與主的相遇而呼高「我的主,我的上帝」呢?與主相遇的經驗,可以幫助我們在不如意的環境中,在看不見上帝的黑暗日子裡,甚至是懷疑上帝的真實的困境中,仍能堅持信仰的真實.記得有一次在國內認識了一位教會的長老.他帶領我們到山區的教會探訪信徒時,在途上特別指著一座橋給我們看.這座橋在他的生命中,扮演著十分重要的角色.原來在文革時期,他被下放勞動,負責在山邊修建這座橋樑.不知是甚麼原因,他在工作期間突然失足墜下,當他以為自己必定喪命時,身體卻被另一些東西承託住.他深信是上帝拯救他的性命,於是他立志餘生要好好的為主而活.八十年代,教會恢復活動,他便積極參與教會事奉,每次當他在事奉過程中面對困難,想要放棄的時候,他就回到這座橋邊,再次思想上帝如何拯救他的生命,他與上帝所立的約,然後便堅持下去.各位弟兄姊妹,我們的生命會否同樣有這樣一座橋,讓我們在面對困難的時候,可以提醒我們不要忘記上帝的恩典?這座生命之橋,既是我們生命中最低谷的地方,但也同時可以是見證上帝與我們同在的地方.如果每個人的生命都是由生走向死亡的旅程,那麼,基督徒在人生路上奔跑,究竟與其他人有甚麼不同?我們是因著信仰而走一條容易及平坦的康莊大道嗎?或是好像中國人駕車時懸掛那個「一路平安」的小符,路路平安,沒有意外?我相信,基督徒生命的不同,是表現於:一,我們是否因著相信復活的主,而相信復活的生命,並因末後的盼望而活得更有勇氣；二,我們是否願意接受復活的主所交託予我們的使命,以此作為我們奔跑人生路的方向；三,我們是否能夠與復活的主相遇,經驗祂的真實.求主幫助我們,讓我們的生命真的能夠因信仰而與別不同.因為復活的主已已祂平安賜下,這是世人不能給予的平安,讓我們的生命也能作復活主的見證.



\newpage

\section{床前明月光}
\label{sec:283}
\textbf{鄧瑞強博士 - 床前明月光}
\newline
\newline
\hyperref[sec:282]{< < < PREV SERMON < < <}
~
\hyperref[sec:toc]{[返目錄]}
~
\hyperref[sec:284]{> > > NEXT SERMON > > >}


講題:床前明月光　A Sleepless Night

經文:詩篇4

講員:鄧瑞強博士

場合:香港中文大學崇基學院禮拜堂

日期:2015年4月19日

各位弟兄姊妹,早安.

當人睡覺時,腦海裡最不願意去想的問題,便是:「為何還未睡著?」

這問題一是不出現,一出現了,便會緊緊地纏著你,像怨女遇上痴男,糾纏不休.看看台頭的鐘,已是深夜一點鐘了,心中問:「為何還未睡著?」再看,已是深夜兩點鐘了,心中問:「為何還未睡著?」再看,已是深夜三點鐘了,心中仍在問:「為何還未睡著?」這問題一纏著你,便永不放手.

聽說,人類其中一個最艱深的哲學問題,便是:「夜這麼深了,為何還有一條狗在外面吠?」

在親密的人際關係裡,其中一種最大的張力是:與你一起睡覺的人,他已睡到「雷聲大作」,而你卻在無奈地聽他的「雷聲」.在這處境裡,一個充滿張力的問題油然而生.究竟是讓他繼續睡覺而我醒著,抑或是讓他醒著而我能睡覺?

有人說,失眠的人,應有正面思想.最正面的思想,莫如自己想:「當世人沉睡時,我還清醒.」

又或者,失眠的人應自負地對早睡的人說:「你們哀歎時光飛逝吧,我們的時光是不會飛逝的.我們過的一個晚上,時間比你們的一輩子還要長.」

對於基督徒來說,最有趣的問題是:「為何聽道時,一下子便睡著了.但到了晚上,卻睡來睡去也睡不著.」

聽說,有一個人,為了克服「失眠」對人的折磨,下了一個宏願,說:「在找到克服失眠的方法前,誓不睡覺!」有一天,他的朋友對他說:「我找到克服失眠的方法了,就是睡覺.」

由於「失眠」很普遍,很多文學作品都與「失眠」有關,這差不多都可以構成一種獨特的文學範疇.蘇東坡的＜水調歌頭＞是一個例子,上半段一開首便說:「明月幾時有?把酒問青天.」下半段開首說:「轉朱閣,低綺戶,照無眠.」月光流轉,照著一個失眠的人.

失眠的人是孤寂的,或者說,孤寂的人是容易失眠的.李白的＜靜夜思＞便是講一個孤寂的失眠人.「床前明月光,疑是地上霜；舉頭望明月,低頭思故鄉.」在床上,卻睡不著.望著地上的月色,疑是地上的霜.睡不著,是否因為太冷呢?再望明月,便想到故鄉.噢,原來他遠離親人,一個人在途上.孤苦伶仃,感到大地和人事的淒冷.或許,就是因為這份淒冷,久久不能眠.

有人說,失眠的人是「眼光光」的,但就因他「眼光光」,他才能看到人生的實況.不過又有人說,就是因為看到人生的實況,一個人才失眠.

失眠與生命的實況,看來緊密相連.某些人,看到某些生命實況,他們失眠.某些人,看到另一些生命實況,他們能安然入睡.

失眠與對抗失眠,是人生的一大問題.失眠從何而來,如何對抗失眠,又與我們如何理解生命的實況有關.當夜闌人靜,看著「床前明月光」時,你的信仰能否讓你看到生命的實況,而能安然入睡呢?

詩篇第四篇的作者,看來與李白和蘇東坡一樣,都是一個失眠的人.不過,他以信仰的心,去克服失眠.今日,我們看看他的見解,能否幫助我們克服失眠?

我講道的經文,是今日的詩篇經課:詩篇第四篇

詩4:1　〔大衛的詩,交與伶長.用絲弦的樂器.〕顯我為義的神啊,我呼籲的時候,求你應允我!我在困苦中,你曾使我寬廣；現在求你憐恤我,聽我的禱告!

詩4:2　你們這上流人哪,你們將我的尊榮變為羞辱要到幾時呢?你們喜愛虛妄,尋找虛假,要到幾時呢?〔細拉〕

詩4:3　你們要知道,耶和華已經分別虔誠人歸他自己；我求告耶和華,他必聽我.

詩4:4　你們應當畏懼,不可犯罪；在床上的時候,要心裡思想,並要肅靜.〔細拉〕

詩4:5　當獻上公義的祭,又當倚靠耶和華.

詩4:6　有許多人說:誰能指示我們甚麼好處?耶和華啊,求你仰起臉來,光照我們.

詩4:7　你使我心裡快樂,勝過那豐收五穀新酒的人.

詩4:8　我必安然躺下睡覺,因為獨有你─耶和華使我安然居住.

我將這詩篇的第1節及第8節,看為是一個禱文的開頭及結尾.然後,第2至3節,4至5節,6至7節分成三個小段落.

(1)第1節:公義的神聆聽人的呼求

夜闌人靜,明月當空.詩人思潮起伏,久久不能眠.他向夜空的主祈禱.

詩4:1　顯我為義的神啊,我呼籲的時候,求你應允我!我在困苦中,你曾使我寬廣；現在求你憐恤我,聽我的禱告!

「顯我為義的神啊」可以簡單地譯為「我的公義的神啊」.

詩人感到人生有欠缺,世事有不足.欠缺的東西,核心是公義.詩人耿耿於懷的,不只是自己的事,更是國事與天下事.世間沒有公義.但公義不是一種空想,公義源自一公義的神.公義的神看著世人,也審判著世人.

真正的祈禱,不是黑夜中獨自哀鳴,而是向神呼求.這種呼求,建基於神的性情.若神是公義的,則我們遇上不義時,便應向神求公義的伸張,而不是胡亂的呼天搶地.

「我的公義的神啊,我呼籲的時候,求你應允我!」

詩人信任神,也信任從神而來的公義.

「我在困苦中,你曾使我寬廣」:這經文用了「窄」與「寬」的意象.我在「窄」中,你使我「寬」了.這可以引伸很多解讀.如:我在死巷中,你使我走上康莊大道.有人這樣解讀:我的喉嚨被人卡著,欲語不能；你卻鬆開這捆綁,讓我開口發聲.

詩人的經驗,曾到窄巷死地,神卻解救他,使他走上生命的大道上.

如今,他又遇上不義的事,被「滅聲」,但他偏偏要大聲呼求.他知道,世間有公義,因為公義的主存在.

(2)第2至3節:面對喜愛虛妄的人

詩4:2　你們這上流人哪,你們將我的尊榮變為羞辱要到幾時呢?你們喜愛虛妄,尋找虛假,要到幾時呢?〔細拉〕

「你們這上流人哪」,現在一般只譯為「世人哪」(sons of men).這班人向詩人做了什麼呢?

他們侮辱人的尊嚴,損害人的聲譽.他們喜愛虛妄,作不誠實的指控.他們弄虛作假,大話連篇.或者說,他們指鹿為馬,顛倒是非；以骯髒的手段,將正義的人貶損為不義的人；以謊話技術,將「愛人的人」污衊成「害人的人」.

面對這樣的世界,詩人的反應,異常平靜.

詩4:3　你們要知道,耶和華已經分別虔誠人歸他自己；我求告耶和華,他必聽我.

詩人看到這個世界是分為兩類人的.一類是「虔誠人」,即信仰忠貞的人,即堅持真理的人,他們是屬於神的.神將他們從世人中分別出來,他們不屬世界,是屬於神的.另一類人當然是「非虔誠人」,他們將假話當真話,他們離遠神.「虔誠人」是屬神的,他們屬於永恆的層次.「非虔誠人」與他們的話一樣,是虛空的,轉眼間便灰飛煙滅.

詩人看自己是「虔誠人」,是屬於神的人.他渴求正義的禱告,正是神的心意.故此,他相信,他求告神的時候,神必聽他.

詩人看來,由於人分兩種,「虔誠人」與「非虔誠人」,而「非虔誠人」顛倒是非,陷害「虔誠人」,這是「非虔誠人」的本性使然,故此,「虔誠人」被屈,不必大驚小怪.世人的奸險,只是顯出,誰屬上帝,誰敵對上帝.屬神的人,應明白必蒙神的保守,故也不必過份驚惶.

(3)第4至5節:面對心懷忿怒的人

詩4:4　你們應當畏懼,不可犯罪；在床上的時候,要心裡思想,並要肅靜.〔細拉〕

「你們應當畏懼,不可犯罪」,＜聖經新譯本＞將之譯為「你們生氣,卻不可犯罪」.這裡講出世人常常失眠的原因.他們在床上,應平靜的時候卻不平靜,他們常常怒火中燒.

詩人徹夜未眠,是因為他想到,世人在夜裡,仍是充滿忿怒,充滿殺機.或許,在黑夜裡,這些人正密謀陷害人的大計.

詩人的反應,是在禱告中規勸他們.勸他們不要犯罪,勸他們平靜下來.夜深了,平靜吧.

詩4:5　當獻上公義的祭,又當倚靠耶和華.

這仍是詩人的規勸.他呼籲心懷忿怒,夜裡未睡的人,當獻上公義的祭,又當倚靠耶和華.神是公義的神,「獻上公義的祭」,就是獻上合乎神性情的禮物,即:作出合乎神心意的行動.這其實就是呼籲「非虔誠人」從惡道中回轉,走回正道.

在一個充滿不義的世界,在隱藏著種種殺機的黑夜裡,詩人默默地祈禱,祈求人熄滅心中的怒火而能安眠,祈求人放下犯罪的念頭而能平靜下來.

(4)第6至7節:面對放棄希望的人

詩4:6　有許多人說:誰能指示我們甚麼好處?耶和華啊,求你仰起臉來,光照我們.

當世人弄虛作假,顛倒是非時；當世人心懷怨忿,充滿殺機時；另一些人就覺得前途暗淡,放棄希望.他們說:「誰能指示我們甚麼好處?」或者說:「還有誰能給我們一個美好的明天呢?」哀莫大於心死,這些人放棄了對明天的想像.沒有明天了,明天已經死了.沒有誰可以救我們了.放棄吧.

憂慮與絕望,都叫人無法安睡.當沒有明天時,一個人怎能睡覺.在絕望的心靈裡,黑夜的空洞,使人找不到落腳的地方.沒有落腳之處,人怎能安心,怎能安然睡覺?

詩人的祈禱是:「耶和華啊,求你仰起臉來,光照我們.」這經文源自民數記6: 24-26的祭司禱文,全文是:「願耶和華賜福給你,保護你.願耶和華使他的臉光照你,賜恩給你.願耶和華向你仰臉,賜你平安.」這禱文講到從神而來的光明,這是在世界的黑暗裡讓人感到平安的光明.

詩4:7　你使我心裡快樂,勝過那豐收五穀新酒的人.

世人的快樂,是外在的,物質的快樂.當衣食豐足時,世人才感快樂.詩人說,有另一種更深層次的快樂.這是「內在的快樂」,「心裡的快樂」,這是勝過「外在的,物質的」快樂的.這種快樂,不受外在環境影響,而源自創造明天的神.不要為明天憂慮,不要以為沒有明天了.明天是屬於神的,人不能從神手中奪去明天.

(5)第8節:神使人安然睡覺

詩人的結語是:

詩4:8　我必安然躺下睡覺,因為獨有你─耶和華使我安然居住.

「必」,即一躺下便睡著.經文出現兩個「安然」.第一個「安然」,指生命的圓滿,無缺無憾.第二個「安然」,指生命的平安,無畏無懼.詩人在無缺無憾,無畏無懼的平安中,一躺下便睡著.這平安的源頭,唯獨是神.

詩人面對的世界是叫人不安的.謊話連篇,陷義人於不義,充滿忿怒,心懷惡念,對明天絕望,充滿憂慮不安.這是一個叫人失眠的世界.

詩人的回應是:一,更肯定自己的身分,自己是屬於神的,是屬於公義的.這不是自我麻醉,自我清高,而是提醒自己,要行事公義,與神的性情一致.二,默默為世人祈禱,希望他們悔改.三,鼓勵絕望的人,幫助他們在神的光明中找到希望.

詩人沒有苦毒,沒有怨言,沒有仇恨,沒有失望.他只是默默地守望,默默地愛,堅定地信靠神.

最後,他以「一躺下便睡著」,去回應世人的失眠.他能「一躺下便睡著」,是因為神是真的.信靠神的人,能在他賜的平安中克服失眠.

你願意以信靠神的心去克服失眠嗎?

但願:榮耀歸於聖父,聖子,聖靈.阿們.



\newpage

\section{好牧人}
\label{sec:284}
\textbf{戴浩輝牧師 - 好牧人}
\newline
\newline
\hyperref[sec:283]{< < < PREV SERMON < < <}
~
\hyperref[sec:toc]{[返目錄]}
~
\hyperref[sec:285]{> > > NEXT SERMON > > >}


詩篇 23:1-6
\newline
\begin{longtable}{cl}
\hline
\hline
章節 & 經文 (和合本修訂版)\\
\hline
23:1 & \begin{tabularx}{0.7\textwidth}{X} 耶和華是我的牧者，我必不致缺乏。 \end{tabularx} \\ \\ \relax
23:2 & \begin{tabularx}{0.7\textwidth}{X} 他使我躺臥在青草地上，領我在可安歇的水邊。 \end{tabularx} \\ \\ \relax
23:3 & \begin{tabularx}{0.7\textwidth}{X} 他使我的靈魂甦醒，為自己的名引導我走義路。 \end{tabularx} \\ \\ \relax
23:4 & \begin{tabularx}{0.7\textwidth}{X} 我雖然行過死蔭的幽谷，也不怕遭害，因為你與我同在；你的杖、你的竿，都安慰我。 \end{tabularx} \\ \\ \relax
23:5 & \begin{tabularx}{0.7\textwidth}{X} 在我敵人面前，你為我擺設筵席；你用油膏了我的頭，使我的福杯滿溢。 \end{tabularx} \\ \\ \relax
23:6 & \begin{tabularx}{0.7\textwidth}{X} 我一生一世必有恩惠慈愛隨著我；我且要住在耶和華的殿中，直到永遠。 \end{tabularx} \\ \\
[1ex]
\hline
\hline
\end{longtable}
使徒行傳 4:5-12
\newline
\begin{longtable}{cl}
\hline
\hline
章節 & 經文 (和合本修訂版)\\
\hline
4:5 & \begin{tabularx}{0.7\textwidth}{X} 第二天，官長、長老和文士在耶路撒冷聚集， \end{tabularx} \\ \\ \relax
4:6 & \begin{tabularx}{0.7\textwidth}{X} 又有亞那大祭司、該亞法、約翰、亞歷山大，和大祭司的親族都在那裡。 \end{tabularx} \\ \\ \relax
4:7 & \begin{tabularx}{0.7\textwidth}{X} 他們叫使徒站在中間，問他們：「你們憑甚麼能力，奉誰的名做這事呢？」 \end{tabularx} \\ \\ \relax
4:8 & \begin{tabularx}{0.7\textwidth}{X} 那時，彼得被聖靈充滿，對他們說：「民間的官長和長老啊， \end{tabularx} \\ \\ \relax
4:9 & \begin{tabularx}{0.7\textwidth}{X} 倘若今日我們被查問是因為在殘障的人身上所行的善事，就是這人怎麼得了痊癒， \end{tabularx} \\ \\ \relax
4:10 & \begin{tabularx}{0.7\textwidth}{X} 那麼，你們大家和以色列全民都當知道，站在你們面前的這人得痊癒，是因你們所釘在十字架、神使他從死人中復活的拿撒勒人耶穌基督的名。 \end{tabularx} \\ \\ \relax
4:11 & \begin{tabularx}{0.7\textwidth}{X} 這位耶穌是：『你們匠人所丟棄的石頭，已成了房角的頭塊石頭。』 \end{tabularx} \\ \\ \relax
4:12 & \begin{tabularx}{0.7\textwidth}{X} 除他以外，別無拯救，因為在天下人間，沒有賜下別的名，我們可以靠著得救。」 \end{tabularx} \\ \\
[1ex]
\hline
\hline
\end{longtable}
約翰一書 3:16-24
\newline
\begin{longtable}{cl}
\hline
\hline
章節 & 經文 (和合本修訂版)\\
\hline
3:16 & \begin{tabularx}{0.7\textwidth}{X} 基督為我們捨命，我們從此就知道何為愛；我們也當為弟兄捨命。 \end{tabularx} \\ \\ \relax
3:17 & \begin{tabularx}{0.7\textwidth}{X} 凡有世上財物的，看見弟兄缺乏，卻關閉了惻隱的心，神的愛怎能住在他裡面呢？ \end{tabularx} \\ \\ \relax
3:18 & \begin{tabularx}{0.7\textwidth}{X} 孩子們哪，我們相愛，不要只在言語或舌頭上，總要以行為和真誠表現出來。 \end{tabularx} \\ \\ \relax
3:19 & \begin{tabularx}{0.7\textwidth}{X} 從這一點，我們會知道，我們是出於真理的，並且我們在神面前可以安心， \end{tabularx} \\ \\ \relax
3:20 & \begin{tabularx}{0.7\textwidth}{X} 即使我們的心責備自己，神比我們的心大，他知道一切。 \end{tabularx} \\ \\ \relax
3:21 & \begin{tabularx}{0.7\textwidth}{X} 親愛的，我們的心若不責備我們，在神面前就可以坦然無懼了。 \end{tabularx} \\ \\ \relax
3:22 & \begin{tabularx}{0.7\textwidth}{X} 我們一切所求的，就從他得著，因為我們遵守他的命令，行他所喜悅的事。 \end{tabularx} \\ \\ \relax
3:23 & \begin{tabularx}{0.7\textwidth}{X} 神的命令就是：我們要信他兒子耶穌基督的名，並且照他所賜給我們的命令彼此相愛。 \end{tabularx} \\ \\ \relax
3:24 & \begin{tabularx}{0.7\textwidth}{X} 遵守神命令的，住在神裡面，而神也住在他裡面。從這一點，我們知道神住在我們裡面，這是由於他所賜給我們的聖靈。 \end{tabularx} \\ \\
[1ex]
\hline
\hline
\end{longtable}
約翰福音 10:11-18
\newline
\begin{longtable}{cl}
\hline
\hline
章節 & 經文 (和合本修訂版)\\
\hline
10:11 & \begin{tabularx}{0.7\textwidth}{X} 「我是好牧人，好牧人為羊捨命。 \end{tabularx} \\ \\ \relax
10:12 & \begin{tabularx}{0.7\textwidth}{X} 雇工不是牧人，羊不是他自己的，他一看見狼來，就撇下羊群逃跑；狼抓住羊，把牠們趕散。 \end{tabularx} \\ \\ \relax
10:13 & \begin{tabularx}{0.7\textwidth}{X} 雇工逃走，因為他是雇工，對羊毫不關心。 \end{tabularx} \\ \\ \relax
10:14 & \begin{tabularx}{0.7\textwidth}{X} 我是好牧人；我認識我的羊，我的羊也認識我， \end{tabularx} \\ \\ \relax
10:15 & \begin{tabularx}{0.7\textwidth}{X} 正如父認識我，我也認識父一樣；並且我為羊捨命。 \end{tabularx} \\ \\ \relax
10:16 & \begin{tabularx}{0.7\textwidth}{X} 我另外有羊，不屬這圈裡的，我必須領牠們來，牠們也要聽我的聲音，並且要合成一群，歸一個牧人。 \end{tabularx} \\ \\ \relax
10:17 & \begin{tabularx}{0.7\textwidth}{X} 為此，我父愛我，因為我把命捨去，好再取回來。 \end{tabularx} \\ \\ \relax
10:18 & \begin{tabularx}{0.7\textwidth}{X} 沒有人奪去我的命，是我自己捨的；我有權捨棄，也有權再取回。這是我從我父所受的命令。」 \end{tabularx} \\ \\
[1ex]
\hline
\hline
\end{longtable}
講題:好牧人(信靠的人生)　Good Shepherd

經文:詩篇23:1-6;使徒行傳4:5-12;約翰一書3:16-24;約翰福音10:11-18

講員:戴浩輝牧師

場合:香港中文大學崇基學院禮拜堂

日期:2015年4月26日

引言

有一次見到一個可能大家都經歷過的現象.有一天,當我在一餐廳喝咖啡並看著街上的人來人往時,看到有幾個人好像很辛苦的抬著東西慢慢的走動,但看不到他們在抬甚麼.於是我心裡想,這些人為甚麼好像很辛苦的抬東西,但又看不出他們抬甚麼?莫非他們在假裝辛苦,好讓老板以為他們不是在偷懶嗎?當他們繼續慢慢的前行,我清楚看到他們在抬甚麼.原來他們是在抬一大塊玻璃,我初時在餐廳看見這些人的角度時,剛好看不到這透明的玻璃,但當換上另一個角度時,太陽的照射,這重擔就清楚在目了.各位弟兄姊妹,人生豈不也是這樣嗎?有多少時候我們所背負的重擔,在人家看來,好像是透明似的,真是有苦自己知.

然而,人生又是否那樣無奈?對於我們信主的人來說,人們不明白我們的苦況,與我們的信仰有何關係?我的意思是上帝在這一點上有何位置?上帝對我來說有甚麼意思?如果是有意思的話,是怎樣的呢?我們的信仰是七天中的一個多小時,抑或是我們生命的中心?要回答以上的問題,我認為可以從詩篇二十三篇找到答案.讓我們在這復活節的期間,從這首詩中,檢視我們的信仰和人生,好使我們信靠主,有力量刻勝生活種種的重擔.詩二十三就是一首既熟悉又不太清楚的詩歌,讓我們從這首歌中了解信靠的人生是怎樣的.

一  在信靠中更新的滿足

傳統認為這首詩的作者是大衛,他從一位牧童,成長為以色列最偉大的君王.聖經把他的崛起和成為君王的事蹟都清楚記載,有苦自己知的經歷也常出現在這些記載的字裡行間.而今天這首信心的詩歌可以說是大衛從信心看人生,是他解決有苦自己知的秘訣.

詩人一開始表達「耶和華是我的牧者,我必不致缺乏(我一無所缺).」他清楚交代了他和耶和華的關係,是以「我的」這個屬格來表示.有上帝就是我的「滿足」,一無所缺.這個世界如果人人都感到一無所缺,從來沒有人的慾望會怎樣?我相信有兩個可能:第一,人類仍然會住在山洞或完全滅亡,因為人將會不求進步.第二,可能世界上再沒有罪惡.有渴求有慾望是人的本性,但甚麼才使我們滿足?詩人在這裡表達:耶和華明白他,他的禱告耶和華會聆聽.詩人是受格,是上帝所關心的對像,詩人繼續以牧人和羊群的意象表達這滿足的原因,「祂使我躺臥在青草地上,領我在可安歇的水邊.」當我們聽到這如畫的詩句時,可能我們不其然的想著在歐洲某湖邊的大片草原.但要知道作者不是歐洲人,而所描繪的背景當然也不是歐洲,而是又旱有多沙石的猶大山區,清草地不是容易找到的,而有水源安息之處更不容易找到.然而不論青草地或是近水邊安息之處,在原文都是眾數.換言之,不單是一次的引領,而是常常的引導和供應.我們再從以色列的牧羊經歷看,我們不難發現,這裡所描述的是十分安慰人的.青草地和安竭的水邊不僅是不容易找到,往往也是充滿危險的地方,因為兇猛的野獸通常就在這些地方等待獵物出現和飲水,就在不防備之際突襲.因此,詩人形容上帝所引導的地方是「可安竭」的水,不是處處危險的水邊.

這滿足的表達也有其歷史的寫照,以色列在曠野四十年之久正如申二7所述:「…你走這曠野,祂都知道了.這四十年耶和華你的上帝常與你同在,故此你一無所缺.」不論以色列民或是詩篇的作者都讓我們清楚一件事,這個世界沒有使我們滿足的,惟有是那位牧者豐盛的帶領.

祂不單使祂所牧養的羊群滿足,這位牧者更更新他的羊群.經過祂的帶領,經過祂所給予的水和安息,羊群得著更新.「祂使我的生命重整,為祂的名的緣故,祂引導我走仁義的路徑.」和合本的「甦醒靈魂」可以繙作「生命重整」就好像廢墟或以不在使用的路徑重新使用一樣.「生命重整」可以十分具體的回復生命力,也可以有深層的了解,就是一個一塌糊塗的生命完全恢復過來.在我們眼前的,有不少這樣的見證,就是一生受毒品所控制的,因耶穌的緣故,得著釋放,完全恢復過來.

各位弟兄姊妹,有苦自己知的滋味可能我們和詩人都嚐過；然而,今天詩人要給我們另一種體驗,不論我們走得多疲憊,所面對的難處和挑戰有多真實,針對我們的人如何兇猛；讓我們信靠那位帶我們走過這一切的主,在祂裡面,得著生命的重整和更新.

二  在信靠中現實的生活

「耶和華是我的牧者」這句話,也成為不少攻擊基督教的人的口實.他們認為,這句話已顯示了基督徒的逃避現實的心理.弗洛依德在他的書「未來的幻想」中說:「宗教使人類維持在一個嬰孩期的安全狀態和行為中,所以人類要在這種宗教變態的心理中釋放出來.」有不少人也以為這首詩就是弗氏所指的宗教信仰,但我們要從詩的背後了解詩人所經歷的,我們就知道他們的指摘是沒有根據的.

詩篇二十三篇並不是指出一個信仰的逃避地方,在第1-3節所指出的美麗帶領中, 我們不難看到詩人也有死蔭幽谷的經歷,正如第4節指出:「我雖然行過死蔭的幽谷,也不怕遭害,因為你與我同在；你的杖,你的竿,都安慰我.」在人生的歷程中,死蔭的幽谷是真實而不可避免的, 但也可能因著這種遭遇, 我們更能體會耶和華是我的牧者的意思.有一次當我和妻子並一位同工駕車經大埔道返回沙田時,突然間下了大雨,那時我正下斜坡,限速是七十公里,我大概是五十多六十公里時速,正要減速,就在一個彎路失控,車子三百六十度旋轉,雨水太大,我看不到前路,這時好像死蔭的幽谷是真實,只要一株樹或是一排防撞欄,都可以使我們死亡或重傷.但祂的同在比現實更真,感謝上帝,我們最後停在一個草地上,我下車檢查,只見前車輪破裂了,於是把後備輪換上,到了汔車維修站,發現第二個車輪也壞了.如在路上兩車輪齊破裂,我們無法把車駕駛到維修站,上帝的恩典,兩個車輪雖是同時因撞擊破裂,但奇妙的不是一起泄氣；上帝在一切危險中與我們同在.當然,這不是每人都可經歷的上帝保守,但感謝主,我經歷了.

當前三節以「我與祂」來表達作者與神的關係時,後三節以「我與你」來表達神人關係.

詩人繼續以牧者的意象來表達耶和華同在的意思, 「你的杖你的竿都安慰我」以色列的牧人常配備兩件牧人用具,一是杖或作棒,好像現代的棒球的棒一樣,是用來防禦野獸或是不友善甚或偷羊的人來攻擊時自衛使用的.第二件是一枝很長的竿,用來一方面引領小羊, 另一方面可以分開公山羊間的爭鬥.當詩人說:「你的杖你的竿都安慰我」時,代表了上帝對他並不單是看著,而是以上帝保護的權能和拯救的能力來安慰他.弟兄姊妹,讓我們體會祂保護的杖和引領的竿如何在我們周圍.現實生活縱然殘酷,但祂保護引領是同樣的真實.

「在我敵人面前,你為我擺設筵席；」我們不僅經歷生活的困難,人生也少不免有很多與我們作對的人.我們的上帝如何幫助我們?詩人形容耶和華就像請客的家主,詩人經過1-3節的生活,有安息,有吃的,有喝的.詩人不再以羊作表徵,而是以耶和華筵席的座上客作表徵,耶和華就是那位請客的家主 (詩七十八19-20；52-53).通常得到貴賓式的招待會有油膏抹,例如:耶穌被馬利亞膏抹.賓客也是家主所保護的,例如:創世記十九章關於羅得的故事.最重要的是這筵席是祝捷的筵席,是仇敵還未攻擊前；耶和華已為詩人作了勝利的祝捷.當我們感到四面受敵被圍攻時,我們可以作甚麼?詩篇二十三篇要成為我們的禱告,讓我們在這些現實的境遇中信靠那位為我們祝捷的上帝.

三  在信靠中積極的人生

在現實的人生中,我們生活在一個差不多甚麼多可以替代的世代.科學的發展使很多以前不能做到的東西都可以更換替代,例如:換心,換腎等手術,前幾天在一個崇基學院的飯局中更有人指出很快可以換頭.以遺傳工程技術複製人體任何器官,甚至另一個你.另外,我們的工作崗位可以被替代,甚至我們的婚姻配偶,也可能會給替代了.然而,上帝與我們的關係是一個一生親密的關係,這關係是不能替代的.我不能替代我的妻子與上帝有這親密的關係,我的兒女也不能因我與上帝的關係,而得著與上帝的關係,雖然我們不論在禱告,教育和行事為人上,都盡力帶領兒女到上帝面前,但他們與神的關係是獨特的.各位弟兄姊妹,上帝藉主耶穌拯救了我們,祂與我們每一個個人的關係是無人能代替的.我們在祂面前,每一個人都是寶貴的,因耶穌那位好牧人為羊(為我們)捨命.有了這樣的關係,我們的人生會怎樣?我們的人生就是屬於主的人生,是一個積極的人生,是一個一生一世都有上帝恩惠和慈愛相隨的人生.換言之,從主那裡我們的人生是一個豐盛的人生和一個信靠的人生,我們怎樣運用它?詩人提醒我們,他要「住」在上帝的殿中直到永遠,意思是他要一生事奉上帝,你呢?你的一生如何運用?很多的人生是一個虛空的一生,因為他沒有所屬,而且常感不足,常常抱怨.詩人給我們一個方法可以勝過這一切,那就是「住在耶和華殿中直到永遠.」換言之,那就是肯定一生都事奉主,一生都為主使用,那你便不枉此生.

結語

「生活有壓力,有苦自己知」這兩句話不能應用在我們這些有好牧人耶穌救贖了的人身上,當我們面對死蔭幽谷或對頭時,讓我們信靠這位好牧人的帶領保護,從祂那裡支取力量,努力事奉祂.



\newpage

\section{繫連妙結}
\label{sec:285}
\textbf{許開明牧師 - 繫連妙結}
\newline
\newline
\hyperref[sec:284]{< < < PREV SERMON < < <}
~
\hyperref[sec:toc]{[返目錄]}
~
\hyperref[sec:286]{> > > NEXT SERMON > > >}


講題:繫連妙結　The Tie That Binds

講員:許開明牧師

場合:香港中文大學崇基學院禮拜堂

日期:2015年5月10日

宇宙間有股強大的力量,要是失去的話,就會產生極大的問題:要是失去了它,一切物件可能會消失；要是失去了它,所有人類也許會滅亡!這股力量,沒有信仰的人可能會認為它來自大自然；有信仰的人,則會認為它來自上帝 - 上帝創造了這力量,發揮著這力量.但同時間撒旦卻不斷破壞這力量,阻擋這力量.到底這股力量是甚麼?有著這麼大的威力?答案就在今日的講題「繫連妙結」之中.

這力量就是「連結」的力量.(power of connection)

小時候,無意中發現自己有個肚臍,年少時總覺得這個肚臍沒有用,覺得上帝造這東西真失敗,又不好看,又不知所謂.但長大後才知道它的貢獻原來非常大,造這肚臍的上帝真偉大!壞看越覺得它美麗.

肚臍是我們曾經有過「臍帶」的證明:當我們還在母親的肚裡,我們就是靠著臍帶與母親連結,取得從母親而來的供應,滋潤.臍帶正是母子間的「能量連結線」.沒有了它,我們也就不會存在,也根本不會出生,不會成長.

臍帶被剪掉後,就留下了肚臍 - 也就是我們離開母腹,進入世界,再清楚不過的記號.正是它令我們知道,我們與母親有著緊密的關係,這也是今日母親節我想要提醒大家的事.

臍帶提示我們母子以臍帶結連,彼此的關係亦靠此逐漸建立起來:我們與他人建立親密關係的需求；滋養自己和他人,培育他人；情感的輸送；還有心靈相通的能力.雖然我們與母親的臍帶從出生那刻已被剪掉,但與母親,父親,兄弟姊妹,配偶,朋友,同事的關係卻保持著連結.

身為基督徒,就需要與上主有美好的結連.正因為衪是創造我們父母親的主,我們也就是來自上主的了.從神學上來說,我們的始祖由上帝所造,上帝就是始祖之祖,是太祖,是最先知祖,所以我們要與衪有美好的連結.連結從古到今都非常重要,人與人要連結；人與上帝要連結；人與自然也要連結.但要是連結本身發生了問題,更多的問題也就接踵而來了.

香港市民之間也極需要一個良好的連結,有好了的連結才能夠有成效,政府官員也察覺到這點,最近亦經常落區,想要團結市民,希望爭取他們支持政府推出的普選方案,這是市民有目共睹的.

近來關於「連結」的問題,不只在香港,在外國也出現了.那是關於一個「移花接木的手術」的.

試想像,有一個A君,他的頭腦靈活,身體卻已衰殘,癱瘓；又有一個Ｂ君,他的身體靈活,臨終前說想要捐出自己的身體.

說起捐出身體,我們會想到器官移植,也就是將一具健康肉體的器官,「轉移」到另一具需要該器官的肉體裡,讓後者的生命得以延續.到目前為此,人們接受的往往是心臟,肝臟,腎臟等的器官移植.那麼移植整個腦袋呢?恐怕很多人都會搖頭認為是不可能的任務.

但意大利醫生塞爾吉奧 ‧ 卡納維洛卻聲稱,「全身移植」的醫療技移在兩年內可望成真,即是像剛才把Ａ君的頭顱,移植到Ｂ君身上般的手術,這當然引來了極大的爭議.從生命哲學來分析,這個手術涉及了三方面的問題.

首先是「能或不能」的問題:全身移植手移的可行性,其實仍有不少爭論.儘管卡納維洛指他們已經攻克了脊髓轉接最後的一部分,但直到今日,頭部移植手術都只局限於動物試驗.而且,神經科學學者Jerry Silver亦不相信這項在老鼠,豬或者狗身上成功過的試驗可以應用到人身上.他認為切斷頭部之後,要耦合破損組織間的神經軸突簡直是異想天開.

復旦大學附屬中山醫院副院長,上海市器官移植重點實驗室主任,朱同玉認為,現在進行頭部移植手術的時機並不成熟,而其中最大的困難就是「脊髓連結」,它是一個非常困難的連結.而且要是把遺傳基因不同的Ａ的頭與Ｂ的身體強制結合,非常有可能會排斥,連作不順暢,這個問題難以解決.

其次就是「應或不應該」的問題:道德上這個手術非常有爭議性.科學給人們生活帶來的影響有著兩面性,而那些在歷史上徘徊於正確和錯誤,是與非之間的科學試驗,就被稱為「黑暗科學」.「身首異處」頭身分離可否存,「換身」手術引來的巨大道德爭議,整具全新的軀體,都會為病人帶來龐大的心理壓力.醫學道德委員會亦幾乎不會托准這種測試.「如果頭部移植手術成功,那麼人的完整性遭到破壞.」西安交通大學醫學院教授王明旭告訴《中國科學報》記者「而且即便移植成功,存活下來的到底是A頭部的提供者還是B身體的提供者?」

另外,如果更長遠地看這個手術的話,自然會問到有關病人後代的問題.即使手術成功,病人的後代遺傳屬性究竟由誰來決定呢?而更令人擔心的是,會不會有不法份子用這個手術來謀取私利,或者有人想自己長生不老,無限延續壽命而剝奪他人的生命?

第三就是「有價值或無價值」的問題:卡納維洛指全身移植手術可以延長絕症病人的壽命,幫助癱瘓病人,以及將身首異處的頭和身聯繫起來,達到不死的效果:人頭由舊的身體出離出來後,換上新的身體就可以不斷存活下去,讓生命延續千年,甚至不死.試想像如果一些特別聰明的人,如愛因斯坦,如果得到全身移植的話,會有怎樣的事發生.

其實,不只是身體的連結,所有的連結都涉及上面三方面的探討.它們可以被歸納為能力及科學的問題(能或不能),倫理和宗教的問題(應該或不應該),以及人生哲學的問題(有價值或無價值).連結並不是全都屬於正向的,反向的連結亦不難找到,甚至是正反同時存在的連結,正如卡納維洛的全身移植手術.正向的連結可以為我們帶來快樂,相反地,反向的自然是破壞的連結,帶來反效果.

讓我們稍為回到聖經的內容,因為聖經之中其實也有不少的連結存在,當中耶穌就用過「葡萄樹與枝子」來形容門徒與衪的連結 - 耶穌是葡萄樹,當然門徒就是枝子.枝子總是常在葡萄樹上,常在的意思是「remain」,是保持著的意思,香港人會說「Keep住」.

「常在」這個詞語的意義,也由約翰用來勉勵門徒和耶穌的關係:主的話語要「常在」門徒心裡(約15:7)；主的仁愛要「常在」門徒的心裡(約15:9)；主的喜樂要「常在」門徒的心裡(約15:11)；與主的關係(朋友,師生,差者與使者)也要「常在」門徒的心裡(約15:13-16).只要門徒的關係好,自然就會結出果子,有美好的成果(約15:16).

聖經的內容,以及剛才提到的「連結」,原來與今天 - 母親節主日,也有很大的關係.約翰提出的四個連結原則,用在父母與子女,家人之間的關係,都十分有用.

聖經提到,「你們若在我裡面,我的話也常在你們裡面…」我們要怎樣保持常在主的裡面?就是讓主的話常在我們心裡.

小時候我參加主日學的時候,每週都要學一句金句,我覺得很有用,所以常常將它們記在心裡:「虛心的人有福了,因為天國是他們的.」(太5:7)；「你若能信,在信的人,凡事都能!」(可9:23)；「你們是世上的鹽,鹽若失了味,怎能叫它再鹹呢?」(太5:13)；「我(耶穌)留下平安給你們,我將我的平安賜給你們…你們心裡不要憂愁,也不要膽怯.」(約14:27)信徒要是把主的話常放在心中,也就常在主裡了.

所以說,每一句說話的威力其實都不少,正如丈夫也要時常記得太太說過她最喜歡的,最害怕的,還有就是最想丈夫做的.有些太太更會不時反問丈夫她說過的話,測試丈夫聽力之餘,也測試記憶力.所以常聽話的丈夫,肯定是常把太太的話放在心中的.

不過,常把家人負面的說話放在心中,卻會破壞彼此的連結,像是父母的責罵,子女的取笑和彈劾,當然也是少記為妙.有時我們把父母,家人不經意的說話,袋一世,自然會痛苦一生,破壞彼此的連結.

要減少這種情況的出現,當然就是要減少負面的說話和破壞性的連結了.「乜你咁蠢,咁懶」「自私,霸道」,「無用」,「無良」,「糊塗」這些指責,標籤性的說話,正是我們需要減少的.少指責,多講自己的感受,多表達關心就更好.兒子冬天穿得少衣服的話,即時的責罵只會令他更拒絕添衣,反而用關心的口吻「我驚你凍親,著多件衫啦」卻會更有效.

除此之外,不用「你」而用「我」來表達,亦是減少破壞性連結的方法之一.以「我」作為出發點,可以讓人覺得平等；相反以「你」為出發點,就會讓人感覺是「由上對下」的命令,難以使人順服.嘗試不說「你激死我」,改為說「我好激氣」；嘗試不說「你唔愛我」,改為說「我覺得無被愛」；嘗試不說「你真係蠢」,改為說「我覺得你呢件事處理得唔夠聰明」.多用「我」,少用「你」,尊重由此而生.

另外,善用「點樣」而不用「點解」亦可以帶在正面的連結.「點解」是不滿的感歎號；「點樣」是尋求答案的問號.前者只會質疑對方的智慧,令對方有輸的感覺；後者卻是邀請對方一同解決難題,是讓對方貢獻的好機會.

以上都是在家庭裡,建立好的「言語連結」的方法.除了「言語連結」,還有「仁愛連結」.

聖經裡提到「我(耶穌)愛你們,正如父愛我一樣；你們要常在我的愛裡」(約15:9),又提到「你們要彼此相愛,像我愛你們一樣,這就是我(耶穌)我命令」(約15:12),還有「人為朋友捨命,人的愛心沒有比這個大的.」(約15:13)所以我們都要常在主的愛裡,大家彼此相愛.主的愛是捨命的愛,犧牲為人的愛,我們正是用這種愛去愛家人.關於父母對子女的愛,有這樣的一個故事:

從前有一個名叫「小強」的少年,他經常為母親分擔家務,不過有一天,他給母親寫了一張「小強賬單」:小強幫媽媽到超級市場買東西,媽媽應付酬勞20元；小強自己起床摺被,媽媽應付酬勞20元；小強幫媽媽洗碗碟,倒垃圾,媽媽應付酬勞30元；小強幫媽媽掃地洗地,抹檯,媽媽應付酬勞30元；媽媽合共需付100元.小強把這張賬單寫好,放在餐桌回房間溫習.媽媽閱過賬單,寫了回覆給小強.讀者們可以猜猜,媽媽有沒有給小強100元?

媽媽給小強的回覆是這樣寫的:媽媽過去15年含辛茹苦養大小強,小強應付0元；媽媽每天為小強準備早午晚餐,茶點,宵夜,小強應付0元；媽媽每週陪小強學琴,學游泳,學打球,小強應付0元；以上合計,小強應付0元.這張「媽媽的賬單」,小強一直留在手裡,因為它令小強明白母親的愛是無條件的,是不計回報的.若是每個家庭裡,每個人都表現到這種愛,家庭的種種問題也就解決了.

有一本名為《世紀末城市心》的書,作者麥成輝形容後現代社會出現了兩大「愛的疾病」,一是「愛滋」,二是「愛無能」.前者的出現,是由於婚前自由亂愛,不認真忠誠地愛,結了婚仍有多餘的愛去愛情人,包二奶,卻不認真愛枕誘人,是亂愛的後果.常說男人愛新鮮,女人愛保鮮；男人愛憑視覺,女人愛憑觸覺；男人愛面子,要受人尊重,女人愛被錫,要常常被愛.後者則是沒有能力愛,不懂怎樣去愛的同時,卻想要得到別人的尊重,愛護.

更錯誤的是,現代人有不少對愛的錯誤觀念,衍生出錯誤的行為.愛無能的人,只求自私自利,只是個人的好處作出發點.「結婚是錯誤,失孩子是失誤,離婚是覺悟,包二奶是醒悟,再婚是執迷不悟.」這重荒天下之大謬的說法,不重視婚姻之餘,更放縱了自己,破壞了美好的家庭,傷透了家人的心.

怎樣的愛才是可取的愛呢?我們可以看看外國人與中國人的對比.外國人是「愛在心裡口不停」,中國人卻是「愛在心裡口難開」 - 中國人總是不擅於具體表達對家人的愛.有一本叫《親情》的書,作者說過「父親去世後,我才發現他實在愛我,但以前的我竟然不知.」真令人傷感.父親母親們何不早些讓子女感受到你們的愛,也讓伴侶感受到你的愛?

中國媽與英國媽的例子,更清楚地說明愛的表達方式應該為何:中國媽一見一見孩子踏進家門,就說:「快D洗手,有點心食,食完先練琴,然後沖涼做功課.」；英國媽一見孩子入門,就說:「我的心肝寶貝,您返來.我好想念你呀,快些來,錫一啖,講講今日在學校的事給阿媽聽.」各位喜歡中國媽還是英國媽呢?我覺得英國媽有可取之處.但當然,最好是中國媽和英國媽的混合:媽媽先說我愛你,然後有點心食,真開心.

更多關於愛的方法,心理教育家Dr. Rose Campbell提出了家人需要多層次的接觸,它們就是眼睛的接觸(關注,尊重,反意),身體的接觸(愛錫,擁抱),耳朵的接觸(留意子女所講,對父母言語的反應)和心靈的接觸(彼此了解,關怴,思想意見的交流.時代周刊特別指出一個常被父母照顧,擁抱,撫摸,尤其是父親,的孩子,長大後會較能忍受壓力,較能與陌生人相處,較快速地學習,以及有較強自信心和安全感.相反在單親家庭,沒有父親或母親的愛護,對子女的損害則較大.

以上都是家人對於「仁愛連結」的建立方法,以下讓我們看看家人間的「情緒連結」:

對於「情緒連結」,耶穌說過「這些事我已經對你們說了,是要叫我的喜樂,存在你們心裡,並叫你們的喜樂可以滿足.」(約15:11)主的喜樂,是真正的喜樂,是沒有憂慮,煩惱的喜樂.它存在於我們心裡,就是要叫我們的喜樂可以滿足.人與人的喜樂是可以傳遞的,家人相處就要更加留意.曾經有一位弟兄被激怒,回到家裡說了一句「我比人激親,唔好掂我!」怎料家人反而去碰他,就是因為他說了一句「唔好掂我」.

家庭時間內有喜樂,歡樂的話,家人方能團結,大家才會喜歡待在一起.一起講笑話,放聲大笑,大叫,甚至一起大哭；或者是來個家庭之夜,大家輪流表演唱歌,彈琴,跳舞等等……這些都是可以讓大家分享歡樂的好方法.當大家都有歡樂的時間,情感戶口便會有所進帳了.

崇拜及祈禱,更加可以令家人在靈性上,精神上和情感上都有所得益.家人一起敬拜,祈禱,可以令家庭更健康,團結家人,互相了解,幫助.有一位弟兄說過「一起禱告的家庭,一定連結在一起,禱告如同一條繩,把家人聯繫起來.

小時候,我在福建的農村長大,每晚都由嬤嬤帶領,到天台一起家庭崇拜.我們望著月亮,看著星星,一起禱告,一起為著我孱弱的身體禱告,一起為希望有肉可以吃而禱告,一起為到香港見爸爸一事禱告.很慶幸當時的禱告,大部分都成就了.又有一次我在英國進修,每天上學,上路開車時我們一家都會禱告,求主眷顧,開路.有一天由我駕車,我的孩子祈禱說:「求主幫助阿爸開車不會撞車,如果撞車不會撞到頭,如果撞到頭不會撞死或撞蠢!」大家聽到都笑了.

正面的情緒互相分享,可以令大家共享喜樂；反面的情緒卻會破壞關係,因為控制不好「怒氣」正是導致家庭不和的死穴.發怒時有個危機:「罪伏門前,他必戀慕你(創4:7)」.所謂的「罪伏門前」,原文就是「罪像守候大門口的惡魔,專吃家裡的人.」；而「戀慕你」則是指「控制,黏著你」.該隱容許脾氣來控制他,結果他就一怒之下殺死了亞伯.

當然,發怒是人之常情,但人不能被怒火所控制.為人父母亦經常會被子女激怒,正所謂「佛都有火」.其實牧師都有會有火,有時更不止是有火,而是火滾,怒火中燒:有一次我被兒子激怒,我警告他說「我Ｄ火升緊」,他也回我一句「阿爸你小心,我Ｄ火都升緊!」所以大家都要注意,怒火會破壞關係.

提到關係,家庭之中亦有所謂「關係的連結」.「你們若遵行我所吩咐的,就是我的朋友了.以後我不再稱你們為僕人；因僕人不知道主人所作的事；我乃稱你們為朋友；因我從我父所聽見的,已經都告訴你們了.」(約15:14-15)耶穌這樣跟門徒說,而朋友一詞正是代表「平等的,互敬的」,朋友間會彼此認識,了解,更會無所不談.孩子長大後,父母更應以「朋友」來對待子女,大家有商有量,給予意見,甚至是提醒.

「不是你們揀選了我,是我揀選了你們,並且分派你們去結果子,叫你們的果子常存.使你們奉我的名,無論向父求甚麼,祂就賜給你們.」(約15:16)我們與主的關係亦像「師生」,因為是被揀選了的門徒.同時,這段經文也表達了差者,使者的本質,例如父母分派子女家務,這樣的工作分派,正是為了養成對家庭的歸屬感和責任.

孩子對父母的聽從程度,最少有三個表現階段:「細時聽話」,「少年時聽Ｄ話」,「青年時則唔聽話」.父母有時候會很「激氣」,因為孩子越大,就越不聽話,越愛「駁嘴」.其實越是懂得「駁嘴」,代表孩子的智商高,口材好,駁贏了父母,某程度上亦是「叻仔」的表現.正如丈夫渴望太太的尊重,父母親也渴求子女的尊重,這也是再正常不過的.但很多時候父母都要放下自尊,接受孩子的批評.

我的三個孩子都長大了,對我和另一半都有很多批評,我也明白是不容易接受的:「阿爸你今日講度果陣講錯野,你zar 車果陣撞樹…」吃一頓飯,被孩子不停地「數不是」,甚至連媽媽也加入批判行列.我的兒子通常對我比較仁慈,會維護我,但可惜家中男女比例是2:3,真慘.

剛才提到的三個聽話程度,可以用這三個比喻來代表:小孩子是鴨子階段,像鴨仔一樣常跟著父母,十分聽話,這時候可以有主僕的關係,可分派他們工作；少年時是鴿子階段,總愛與父母討論,聽取父母的意見,這正是師生的關係；到了青年階段,孩子則見面更少,甚至可能在外築巢,父母的意見更只是參考,親子間就像是朋友的關係.清楚了解這些不同階段,不同的關係,扮演不同的角色,正是做個醒目父母的關鍵.

家庭的連結不像臍帶或是葡萄樹與枝子,是看不見的連結.大家都要切記,無論是看得見或看不見的連結,要是正向的話,就會有正能量,反向的話就只會產生負能量,不斷破壞彼此的關係.英因循道會的創辦人衛斯理(John Wesley)先生受父母的感受最大:他的母親共有十九個兒女,有人詫異衛斯理母親為什麼能在這麼多子女身上有如此大的感動力,她就回答說並沒有什麼秘訣,只是星期一晚找美莉,獨自帶她到房中跟她談道和禱告；星期二約翰,星期三南基,星期四雅各…如此類推,每晚就一個人.衛斯理夫人就是這樣把神介紹給孩子.

還記得在一個追思崇拜中,一位兒子講述父親的一生,最後一段給我很大的提醒.「我的父親是位好Daddy,又是位好丈夫.我需要人時,他關心我；我痛苦時,他安慰我；我走錯路時,他糾正我；我需要禱告時,他與我禱告.我感到自豪因我有位好父親.」

各位,假如有一日,我也離開這世界,我的兒子卻這樣說:「我父親是位名人,是教會長老執事,但卻常常不在家,是社會上成功人士,卻不是個好爸爸,因為他沒有關心我,陪伴我.」我當然不想這種事情發生,因此我總是提醒自己,我要當一個像我父親一樣的爸爸,他是我獨一無二的父親,我也要像他般,做個「似樣」的爸爸.

在家庭的各種連結,也就是言語,仁愛,情緒,關係的連結,大家能做到嗎?如果被問及「能不能連結?」,「應不應連結?」,「是不是有意義的連結」,大家的答案又會是甚麼呢?要是身為父母的話,更應該常常反問自己家人的連結怎樣?是不是理想的連結?我們也要多留意,在上面的四個連結裡,哪些方面要特別多留意?家人之間,又有沒有負面的連結?有的話又應該怎樣消除,化解呢?

求主賜福各位的家庭都能「繫連妙結」!阿們!



\newpage

\section{天上的信息地上的使者}
\label{sec:286}
\textbf{伍渭文牧師 - 天上的信息,地上的使者}
\newline
\newline
\hyperref[sec:285]{< < < PREV SERMON < < <}
~
\hyperref[sec:toc]{[返目錄]}
~
\hyperref[sec:287]{> > > NEXT SERMON > > >}


提摩太後書 4:9-22
\newline
\begin{longtable}{cl}
\hline
\hline
章節 & 經文 (和合本修訂版)\\
\hline
4:9 & \begin{tabularx}{0.7\textwidth}{X} 你要趕緊到我這裡來。 \end{tabularx} \\ \\ \relax
4:10 & \begin{tabularx}{0.7\textwidth}{X} 因為底馬貪愛現今的世界，已經離棄我，往帖撒羅尼迦去了；革勒士往加拉太去；提多往撻馬太去； \end{tabularx} \\ \\ \relax
4:11 & \begin{tabularx}{0.7\textwidth}{X} 只有路加在我這裡。你來的時候把馬可帶來，因為他在服事上於我有益。 \end{tabularx} \\ \\ \relax
4:12 & \begin{tabularx}{0.7\textwidth}{X} 我已經打發推基古往以弗所去。 \end{tabularx} \\ \\ \relax
4:13 & \begin{tabularx}{0.7\textwidth}{X} 我在特羅亞留給加布的那件外衣，你來的時候要帶來，那些書也帶來，特別是那幾卷羊皮的書。 \end{tabularx} \\ \\ \relax
4:14 & \begin{tabularx}{0.7\textwidth}{X} 銅匠亞歷山大多方害我；主必照他所行的報應他。 \end{tabularx} \\ \\ \relax
4:15 & \begin{tabularx}{0.7\textwidth}{X} 你也要防備他，因為他極力抗拒我們的話。 \end{tabularx} \\ \\ \relax
4:16 & \begin{tabularx}{0.7\textwidth}{X} 我初次上訴時，沒有人前來幫助，竟都離棄了我，但願這罪不歸在他們身上。 \end{tabularx} \\ \\ \relax
4:17 & \begin{tabularx}{0.7\textwidth}{X} 惟有主站在我身邊，加給我力量，使我能把福音完整地傳開，讓所有的外邦人都聽見；我也從獅子口裡被救出來。 \end{tabularx} \\ \\ \relax
4:18 & \begin{tabularx}{0.7\textwidth}{X} 主必救我脫離一切的兇惡，也必救我進他的天國。願榮耀歸給他，直到永永遠遠。阿們！ \end{tabularx} \\ \\ \relax
4:19 & \begin{tabularx}{0.7\textwidth}{X} 請向百基拉、亞居拉和阿尼色弗一家的人問安。 \end{tabularx} \\ \\ \relax
4:20 & \begin{tabularx}{0.7\textwidth}{X} 以拉都在哥林多住下了。特羅非摩病了，我把他留在米利都。 \end{tabularx} \\ \\ \relax
4:21 & \begin{tabularx}{0.7\textwidth}{X} 你要趕緊在冬天以前到我這裡來。友布羅、布田、利奴、革老底亞和眾弟兄都向你問安。 \end{tabularx} \\ \\ \relax
4:22 & \begin{tabularx}{0.7\textwidth}{X} 願主與你的靈同在！願恩惠與你們同在！ \end{tabularx} \\ \\
[1ex]
\hline
\hline
\end{longtable}
講題:天上的信息,地上的使者-這個冬天不太冷　Heavenly Message, Earthly Messenger

經文:提摩太後書4:9-22

講員:伍渭文牧師

場合:香港中文大學崇基學院禮拜堂

日期:2015年5月17日

代言人是近年新興的推銷商品方法.前段韓劇「來自星星的你」風靡港,台,澳,大陸,走在大商場,盡是是男女主角金秀賢,全智賢的肖像.但商品的代言人,跟本不認識商品,用過商品；不少消費者就因為鍾情代言人,透過購買商品.偶像用這化妝品,我也用這化妝品.賣了不良商品,損失金錢事小,但吃了不良的藥品就影響健康了.

最近國內通過法案,規定代言人一定要親身用過和認識所代言的商品,代言人應該是見証人,否則就是作出虛假聲明.

耶穌的福音也有代言人.今天是耶穌升天主日,下周是聖靈降臨節.耶穌升天時對門徒說:「聖靈降臨在你們身上,你們就必得著能力,並要在耶路撒冷,猶太全地,和撒馬利亞,直到地極,你我的見証.」(使一8)

在商業消費社會,知名的代言人抬捧商品,使商品增值；但福音的代言人,是天上的福音使我們尊貴,我們不過是地上的使者.我們來自塵土,出身卑微,而且生命軟弱,容易迷失,但天上的信息已托負給我們了.為甚麼?因為我們是見證人.

一   地上的使者-我們的本性容易迷失, 使一16,1 7；約十三25, 26

培根有句名言:「犯錯因凡人,悔改靠神恩,沈溺是魔障」(To err is human, to repent divine, to persist devilish) .地上的使者,指出使者生命的軟弱,容易懼怕退縮,容易迷失貪愛世界,容易橫梗悖逆.

猶大被召成為使徒,但因貪財出賣耶穌,失去了使徒的職份.「他本來列在我們數中,並且在使徒的職任上得了一份.這人用他作惡的工價買了一塊田,以後身子僕倒,肚腹崩裂,腸子都流出來.」(使一17, 18)

猶大的角色令人費解,甚至引起爭議.猶大是否按著上帝寫好的劇本,演出交付給他歹角的角色；他的使徒職份是偽裝的,根本是撒旦的臥底；甚至有人說猶大的功勞不少,沒有他耶穌怎樣上十字架?

耶穌沒有被猶大出賣也可被捉拿,被釘十字架.耶穌被門徒出賣,旨在說明他嚐盡人間罪惡帶來的痛苦.莎士比亞名劇「凱撒大帝」(Julius Casesar)中的一幕,凱撒步上元老院階級時,被其親信刺了二十多刀,其中至命的來自最信任的布魯圖:Et tu, Brute? (You too, Brutus? 布魯圖,你都有份嗎?) ,跟著凱撒倒在地下.

猶大的使徒的職份也是真的.為此彼得進行了認真的補選程序,而不是宣佈猶大使徒職份無效.耶穌呼召猶大進入使徒之列,在猶大快要出賣他時想輓回他,懸崖勒馬.「耶穌就蘸了一點餅,遞給加略人西門的兒子猶大.」(約十三26)蘸一點餅遞給猶大,是表示一種關愛.猶太人有一習俗,請客的席中安排上好的餅,由主人親自遞給所敬愛的客人中國主人奉上雞脾.但猶大並未領情,沒有甚麼反應,他硬了心,拒絕耶穌的愛.「他吃了以後,撒旦就進入了他的心.」(約十三27) 並不是餅有甚麼魔力,吃了撒旦就進入心中,乃當人拒絕上主的愛時,撒旦便有機可承了.

我們不知道甚麼時候,猶大衛了金錢出賣耶穌:「這人用他作惡的工價買了一塊田」.是否當猶大看準了這塊地,想擁有這塊地,而出賣自己使徒的身份呢?像以掃為了一碗紅豆湯,出賣了長子的名份.

在諸多被揭發的內地貪官案例中,很多官員初時是勤儉愛民,心底善良,但因為制度不健全,受到金錢的誘惑,德行被腐蝕,變成另一個人.報載一位軍隊貪官被查抄時:「在豪宅內的地下室發現有一噸多重的現金,而宅內各種金銀珠寶更是不可勝數,要臨時叫來十幾輛軍用車才 全部運走.」

誠然,猶大是從恩典中墮落了,因為貪婪出賣耶穌,出賣了使徒的身份.另一個在恩典中墮落的例子是底馬.

保羅第二次下在監前,初次被提訊時,沒有人前來幫助他；他期望站在他旁邊的人,竟然離開他,特別是底馬:「因為底馬貪愛現今的世界,就離開我往帖撒羅尼迦去了.」(四6)底馬是誰?為何保羅感嘆底馬離他而去呢?底馬就是保羅在腓利門書所提及曾經和他出生入死的同工.

世界的魅力.現今的世界,永遠閃爍著迷惑,叫天路客放軟了腳步,模糊了遠象,偏離了方向.貪愛一字,顯示了世界永恆的魅力-眼目的情慾,肉體的情慾,和今生的驕傲.有一首歌這樣描述現今的世界:我們活在這時代,除了追求金錢,成功,名譽,炫耀外,追求其他的價值會被嘲弄,甚至遭消滅；唯有金錢,成功,名譽,炫耀,因為我們活在唯有物資的世界.(We are living…in the age, in which the pursuit of all values other than money, success, fame, glamour have either been discredited or destroyed. Money, success, fame, glamour, for we are living in the age of the thing.)

貪愛一字,也表明信徒不是經過誠實的思考而放棄信仰,乃是被吸引隨流飄離信仰.魯益師(C. S. Lewis)說:「我們要常常被提醒所信為何,沒有一樣信仰會自動在我們的頭腦中茁壯成長的,它需要培育滋養.事實上,若你好好查驗一百位離棄基督信仰的人,我懷疑多少是經過誠實思考而離開信仰,很多都是隨流飄走,不是嗎?」(We have to be continually reminded of what we believe. Neither this belief nor any other will automatically remain alive in the mind. It must be fed. And as a matter of fact, if you examined a hundred people who had lost their faith in Christianity, I wonder how many of them would turn out to have reasoned out of it by honest argument? Do not most people simply drift away?- C.S. Lewis, Mere Christianity) .

你有沒有好友,多年和你在團契一起事奉,為信仰大發熱心,甚至是領你信主,但現在己沒有返教會,他沒有放棄信仰,但金劍已沈埋,壯志蒿萊.崇基神學院邢福增院長在最後的週四晚會寄語畢業同學時提到:他是安素堂長大的,作為歷史學者,他特別安素堂的歷史有興趣.安素堂是紀會前循道衛理聯合教會前會督黃安素牧師的,台灣東吳大學也在校園紀念他的安素堂.離世前,黃安素牧師按立了三位他信賴的同工,一位是梁林開牧師,另外兩位現在已離開教會了.

信仰不會自動在我們心中茁壯成長的,它需要培育滋養,特別是信徒互相激勵；一堆炭燒得熊,一顆炭容易熄滅.保羅沒有因為底馬離開而灰心,因他看到另一群熱心的同工.

我們信福音和傳播福音信息的人,是福音帶給我們尊貴的身份,盡管我們是軟弱的.信息是屬天的,傳信息的使者是屬地的,這是保羅在提後四9-22給我們的信息.

二   地上的使者-生命的限制,在限制中做好自己,提後四13

「我在特羅亞留於加布的那件外衣,你來的時侯要以帶來,那些書也要帶來,更要緊的是那些皮卷.」(提後四13)

今天使徒行傳經課,選立代替猶太的使徒,結果馬提太遞補空缺,事件給我的氣氛,教會像旭日初升,朝氣勃勃；使徒將分散到各地傳揚福音.這景象好像電影「七劍下天山」的開始,你到那處,他到這處…….但讓我們進到使徒行傳的後面,看到各使身體隨著時間的流逝,開始身體已不及當年勇了,因為我們是地上的使者；有些福音的使者更因為因傳福音而坐牢,身體衰老得更快.

提後四9-12記述一位年過六十的老人福音使者,上了枷鎖,單獨被囚禁在地下只有四立方公尺的斗室中,陽光透過窄小的頂部的窗戶,灑在頭髮稀疏的老人的頭頂.然而,這位老人沒有自怨自艾,這位老人就是使徒保羅,他是因傳揚福音而下獄.面對同工的變節和離棄,他心中也沒有任何憤懟.

冬天很快就來到,因為地中海氣候,冬雨夏乾,屆時囚室四壁潮濕,陣陣冷風,寒氣迫人,而且冬天過後,保羅自己預感將面對身首異處的處決:「我現在被澆奠,離世的日子到了.」(提後四6)隨著年日的流逝,他身體開始衰弱,他需要麼加衣禦寒:「我在特羅亞留於加布的外衣,你來的時候可以帶來.」(四13)在人看來,這是一個寒冷的冬天,但對保羅來說,這個冬天不太冷.

保羅第一次被囚是軟禁,可見客人,但這次上了枷鎖,而且身體已不及當年勇了.然而,冬天寒凍,有衣物禦寒；禁見訪客,可以閱讀,亦可寫作.「我在特羅亞留於迦布的那件外衣可以帶來,那些書也要帶來,更要緊的是那些皮卷.」(四13)我們若接受自己的限制,在制約的環境中,也可以從容自若,盡量做好自己.保羅曾經年輕過,那時他拿著大祭司的文書,進入大馬士革,搜捕基督徒,那時候保羅何其強壯.史提反被擲石塊,行刑的人,都把衣服放在保羅的腳前,那時他何其威武；壯年時,保羅曾風塵僕僕,遠赴歐洲傳道,未嘗言倦.他曾靠主行神蹟,使死人復生,病人得愈,驅魔趕鬼,如今天氣寒冷,難道不可以求上主賜他金剛不壞之身?

保羅明白,人的限制和自然規律,是上主為我們立下的界限,不能逾越,他甘心順服,要求提摩太攜來外衣,抵禦風寒.其實很早的時候,他曾三次求主醫治他身體的軟弱-拿走身上的的刺,這刺很可能指眼疾,加六11:「請看我親手寫給你們的字是何等的大呢?」但上主不允,他就順服,毫無怨言,藉此認識到自己的限制,學到謙遜.在不同的人生季節,都有不同的限制,萬物各按其時,在諸般的限制中,盡量做好自己.

這個冬天不太冷,因他有外衣蔽體,有書卷閱讀.

三   地上的使者-生命像泥土,可被陶造成大器

這個冬天不算冷,也因為他看見許多人被上主使用.

窩心暖意.保羅沒有覺得這是冬天太嚴寒,因為很多事使他感到窩心.不錯,他看到人的軟弱,離棄信仰,但同時有更多的人,生命被陶造,持守真理,勇往向前.保羅不會因底馬貪愛世界離開他而意志消沈,因他看到更多的人被主的愛所吸引,愛慕主的顯現多於眼目的情慾,羨慕將來公義的冠冕多於現今的名和利.

他看到革勒士往加拉太,是為事工的需要,提多往撻馬太也是為了事工,而且路加一路都陪伴左右,還有馬可,他就是馬可福音所記載,耶穌被賣時,有一少年人被兵丁追問:你可否認識耶穌?他急忙把身上的布拋開逃跑,他是一位十分怛怯的人.「有一個少年人,赤身披著一塊麻布,跟隨耶穌,眾人就捉拿他.他卻丟了麻布,赤身逃走了.」(可十四51)第一次保羅與巴拿巴到歐洲傳福音時,他中途怕太艱辛而離隊,所以第二次保羅到歐洲傳道時,堅持不攜帶馬可.(徒十三13；十五36~38)但勸慰子巴拿巴慧眼識英雄,有培訓的恩賜,曾經栽培保羅,在他悉心提攜下,馬可得到接納和磨練,最終成為幹練的宣教士:「在傳道的事上於我有益.」

保羅的眼目,不單看到令他失望的底馬,他看到更多被主改變生命的同工.在教會,祇要我們細看,我們也會看到上帝在不少人的生命中工作,像泥土在陶匠手中被陶造,生命得著激勵.

冬天己經來到了,春天還會遙遠嗎?

保羅很早就知道,人生季節最後的一站是冬天,所以他很早就栽培提摩太為接班人,他對提摩太說:「但你所學習的,所確信的,要存在心裡；為你知道是跟誰學的.」(三14)保羅知道他快要澆奠離世了,巴拿巴怎樣訓練自己,培育自己,他也要如何訓練提摩太,接續他未完成的工作.人要過去,但上主的工作會持續發展.保羅不是為了事工而做事工,他是透過不同人的生命,去推動事工,所以保羅在人生的晚年,特別珍惜同工,除了我們提過的提摩太,馬可,路加,還有許多人名:百基拉亞居拉夫婦,亞尼色弗一家人等等.

總結:地上的使者,可以成器,也可以被棄.

家學淵源提摩太.在不同的人物類型中,我們看到一路順暢,預備接班的提摩太,他年少時己受教於敬虔的袓母羅以和母親友妮基.「想到你心裡無偽之信,這信是先在你外祖母羅以和你母親友妮基心裡的,我深信也在你的心裡.」(一5)我們看到家庭教育對人材的培育何其重要.羅以-友妮基-提摩太,提摩太家學淵源,得天獨厚,保羅來料加工,便成大器.盼望我們也重視家庭宗教教育,食飯前謝恩,臨睡前講一簡短聖經故事,然後合手謝謝天父睡覺.

大器晚成馬可.另一位是起步緩慢的馬可,馬可也有敬虔母親馬利亞,她開放家庭作為禱告的地方-馬可樓(徒十二12).可能長於富裕家庭,馬可難耐艱辛,他也生性怯懦,但經過生命導師巴拿巴悉心訓雕琢,鼓勵誘導,給他第二次機會實踐磿練,最終也能大器晚成,在傳道的事上於保羅有益.

摸不透的底馬.最後是摸不透的底馬.底馬隨流飄去,究竟是外在的壓力令他退縮,或是靈魂深處的潛伏的牽引使他迷失呢?作為馬其頓省省會的帖撒羅尼迦的大都會生活,何時開始對他有特別的吸引力呢?能與保羅,路加等使徒同工並受他們肯定,底馬表現令人激賞,但在靈魂深處,卻孕育了靈性的危機,底馬已被現今世界的閃爍炫耀所吸引.

究竟底馬何時讓私慾在心中成孕呢?保羅不知道,但底馬自己應該知道,我們的靈魂渴慕甚麼,其他人可能不知道,但我們自己應該知道.保羅開始不用馬可,但後來愈來愈欣賞馬可,我想馬可有一顆羨慕-聖工的心,他的靈魂深處,渴想事奉上帝,被主的愛所吸引,愛慕主的顯現多於眼目的情慾,羨慕將來公義的冠冕多於現今的名和利.

天上的信息,已交付給我們這些地上的使者了,我們來自地上的塵土,也有一天歸回塵土,但因著天上的信息,我們有永恆的生命,我們和上帝的關系,連死亡也不能間斷.

我們地上的使者,誠然是瓦器,但因為天上的信息,我們懷有莫大的能力.我們的心可以被迷惑,但也可以被上主的愛所吸引,愛慕主的顯現,羨慕將來公義的冠冕.



\newpage

\section{起來不是奴隸的你}
\label{sec:287}
\textbf{鄧瑞強博士 - 起來,不是奴隸的你}
\newline
\newline
\hyperref[sec:286]{< < < PREV SERMON < < <}
~
\hyperref[sec:toc]{[返目錄]}
~
\hyperref[sec:288]{> > > NEXT SERMON > > >}


羅馬書 8:12-17
\newline
\begin{longtable}{cl}
\hline
\hline
章節 & 經文 (和合本修訂版)\\
\hline
8:12 & \begin{tabularx}{0.7\textwidth}{X} 弟兄們，這樣看來，我們不是欠肉體的債去順從肉體而活。 \end{tabularx} \\ \\ \relax
8:13 & \begin{tabularx}{0.7\textwidth}{X} 你們若順從肉體活著，必定會死；若靠著聖靈把身體的惡行處死，就必存活。 \end{tabularx} \\ \\ \relax
8:14 & \begin{tabularx}{0.7\textwidth}{X} 因為凡被神的靈引導的都是神的兒子。 \end{tabularx} \\ \\ \relax
8:15 & \begin{tabularx}{0.7\textwidth}{X} 你們所領受的不是奴僕的靈，仍舊害怕；所領受的是兒子名分的靈，因此我們呼叫：「阿爸，父！」 \end{tabularx} \\ \\ \relax
8:16 & \begin{tabularx}{0.7\textwidth}{X} 聖靈自己與我們的靈一同見證我們是神的兒女。 \end{tabularx} \\ \\ \relax
8:17 & \begin{tabularx}{0.7\textwidth}{X} 若是兒女，就是後嗣，是神的後嗣，和基督同作後嗣。如果我們和他一同受苦，是要我們和他一同得榮耀。 \end{tabularx} \\ \\
[1ex]
\hline
\hline
\end{longtable}
主題:起來,不是奴隸的你　Not Be a Slave

經文:羅馬書8:12-17

講員:鄧瑞強博士

場合:香港中文大學崇基學院禮拜堂

日期:2015年5月31日

各位弟兄姊妹,早安.

我們國家的國歌,第一句便說:「起來,不願做奴隸的人們.」

每天,電視定時定刻播放這國歌,提示我們,「起來,不願做奴隸的人們.」

這其實是一個很好的提示:我們都是奴隸.

問題只是,有多少人懂得「起來」,不再沉醉在奴隸之鄉.

做奴隸是容易的.人叫你吃,你便吃.人叫你睡,你便睡.無需運用自由,無需看清真相,無需思考價值.

要釋放奴隸,是艱難的.要解開他的手腳,這還容易.要釋放他的靈魂,則甚艱難.清朝完結後,民國時代的人,仍對著末代皇帝,自稱奴才,俯首下拜.革命打倒了體制上的專制,但靈魂卻仍不懂得飛翔.

人類文化很早已留意到人類自身的奴役狀況.柏拉圖著名的「洞穴比喻」,講到人困在洞穴裡,背向真實的世界,自以為見到的影像是真實,不相信在影像世界之外有一真實的世界.井底之蛙,不相信井外另有天地.這些人是井底的奴隸.柏拉圖的講法雖然古老,卻十分現代,在這個手機橫行的年代,多少人沉醉在影像的世界裡.

來到近代,有著名的電影「Matrix」,同樣的主題,展示人的奴隸性.不過,今次,人是科技的奴隸.在超級人工智能力量下,人只是被利用的工具,用完即棄.

基督教信仰傳統,由始至終,都是講奴隸得釋放的故事.舊約的出埃及故事,講的是以色列人脫離埃及法老的奴役,成為一群屬於神國度的新子民.新約的故事一脈相承,講的是基督徒脫離「罪」的奴役,成為天國的子民.

基督徒多多少少仍是「罪」的奴隸,雖然我們腳上的枷鎖已被解開.基督徒實質上是自由人,雖然我們很多時仍未懂得如何運用信仰的自由.

今日的講道,內容環繞著奴役與自由.

講道經文是:羅馬書8:12－17(我採用《新漢語譯本》的譯文)

羅8:12　弟兄們,這樣看來,我們都有虧欠,但不是虧欠肉體,去隨從肉體而活.

羅8:13　因為你們若隨從肉體而活,就必定會死；但你們若靠著聖靈,把身體的惡行置於死地,就必活著.

羅8:14　因為凡是被神的靈引領的,都是神的兒子.

羅8:15　事實上,你們所領受的,不是奴役的靈,使你們再有懼怕；你們所領受的,是兒子名分的靈；靠著這靈,我們喊叫:「阿爸!父啊!」

羅8:16　聖靈親自與我們的靈一同見證我們是神的兒女.

羅8:17　既是兒女,就是承受產業的人,是承受神產業的人,且與基督一同承受產業──只要我們與他一同受苦,與他一同得榮耀.

這段經文開宗明義講到人不應該仍在奴隸的狀態中.

羅8:12　弟兄們,這樣看來,我們都有虧欠,但不是虧欠肉體,去隨從肉體而活.

羅8:13　因為你們若隨從肉體而活,就必定會死……

「我們都有虧欠」,意即我們都不是自滿自足的,我們都「有所待」,不同的力量牽引著我們的生命.

「但(我們)不是虧欠肉體,去隨從肉體而活」,意即我們沒有必要成為「肉體」的奴隸.這種奴隸沒有好結果,至終只是無意義的死亡.

保羅渴望我們正視生命的光景.一,各種力量努力主宰我們；二,現實上,「肉體」的力量主宰我們；三,我們實在沒有必要被「肉體」主宰；四,被「肉體」主宰,只有死路一條.

這裡所講的「肉體」,不是指肉體的正常欲望.保羅口中的「肉體」,是一個神學用字,指「與神為敵」的一種存活態度.身體的「物質性」沒有本然的邪惡,邪惡的是人敵對神的態度.一個敵對神的人,他那沒有「物質性」的「心靈」都是「屬肉體」的.「隨從肉體」的存活態度,叫人滿足於成為一隻快樂的動物,而不是具有神形象的人.我這種講法對動物沒有惡意,但人淪為動物,在街頭做愛(註:這是近日在香港發生的新聞),則失去人性應有的尊嚴.

保羅口中的「肉體」,又可指看重一個人的肉身特徵,而忽略一個人內在的靈性內容.保羅為何這麼反對「割禮」,就是因為一個人容易因為這個「肉體」的特徵而以為自己便是屬神的人,而忘記一個人屬不屬神,重點不在肉體的特徵,而在乎內在的靈性內容.一個人倚靠自己的「肉體特徵」而以為自己是屬神的,則這人是「屬肉體」的.同樣,一個人擁有外在的財富或美貌,便以為自己的生命豐盛或優美,這也是「屬肉體」的.

一個在「肉體」的存活態度裡被奴役的人,他只有現世性,現實性,物質性.他像一個朝生暮死的人,千方百計去奪取下一刻要吃的東西,想像不到明天,也沒有靈魂.

羅8:15　事實上,你們所領受的,不是奴役的靈,使你們再有懼怕……

這經文說得很白,活在「肉體」中,就是活在「奴役的靈」之下,這是一種將我們變成奴隸的力量.這種力量控制我們的板斧,是使我們懼怕.我是一個畏高的人,這種懼怕,令我不能走上險峰.有很多人,因為沒有靈魂,也畏懼人性的高度.

緬甸的昂山素姬(Aung Sang Suu Kyi)講過:「使人腐化的,不是權力,而是恐懼.支配權力的人,因恐懼失去權力而腐化；被權力支配的人,因恐懼權力所帶來的苦難而腐化.」她又說,恐懼所引致的腐化,是一切腐化的根源.這很容易理解,若我們恐懼自己擁有得不夠多,我們可能生發貪念.若我們恐懼別人勝過我們,我們可能生發敵意.北韓的領導人恐怕手下不聽話,便用高射炮射他.一種因懼怕而生的絕對腐化.

一個沒有靈魂的人,容易恐懼.一個恐懼的人,又容易出賣他的靈魂.

在香港,政治素質的變化,社會的矛盾撕裂,經濟的不正常發展,在在都使人覺得,令人恐懼的元素越來越多,看來,叫人出賣靈魂的機會也將越來越大.

基督徒在世上,所事何為?

羅8:13　因為你們若隨從肉體而活,就必定會死；但你們若靠著聖靈,把身體的惡行置於死地,就必活著.

羅8:14　因為凡是被神的靈引領的,都是神的兒子.

羅8:15　事實上,你們所領受的,不是奴役的靈,使你們再有懼怕；你們所領受的,是兒子名分的靈；靠著這靈,我們喊叫:「阿爸!父啊!」

羅8:16　聖靈親自與我們的靈一同見證我們是神的兒女.

信耶穌,不只是接受一些信念,而是改變整個存活的方式.

信耶穌,意味著讓聖靈主宰我們的生命,放棄「肉體」的生活方式.這意味著:除了欲望的滿足,更重要的,還有心靈的自由；除了當下的利益,還有普世的價值；除了竭力的攫取,更重要的,還有真誠的饋贈；除了表面的風光,更重要的,還有內心的豐盛；除了肉體,還有靈魂；除了今生,還有來世.

人生是有選擇的,只要你不被「奴役的靈」所嚇怕.「奴役的靈」會恐嚇你:若你不收聲,則「無飯」吃.若你不乖乖唸這台詞讓我拍段短片放上網,你就失去免費遊學的機會.若你走出來爭取正義,你的前途將會被不正義地打壓.

人生是有選擇的,只要你不懼怕.基督徒不懼怕是有理由的,因為基督徒相信有神,相信人有靈魂,相信永生.

人生是有選擇的,我們沒有必要成為「肉體」的奴隸,我們可以選擇成為「神的兒女」,成為真正的自由人.

在過去的日子裡,只有皇帝才被稱為「天子」－－神的兒子.現在,在基督的救贖下,我們基督徒也被稱為「神的兒女」.

羅8:15　……你們所領受的,是兒子名分的靈；靠著這靈,我們喊叫:「阿爸!父啊!」

羅8:16　聖靈親自與我們的靈一同見證我們是神的兒女.

想想,「神的兒女」,多麼尊貴的身分.我們可向人說,萬有的主宰,是我們的「阿爸,父」.想想這身分,應該有何等的氣度!怎能做人自私自利,有失這身分.怎能行事不明不白,有失這身分.怎能思想卑劣低下,有失這身分.

神的兒女,行事光明正大.保羅在以弗所書5:8－9這樣說:

弗5:8　從前你們是暗昧的,但如今在主裡面是光明的,行事為人就當像光明的子女.

弗5:9　光明所結的果子就是一切良善,公義,誠實(即:仁愛,正義,真理).

基督徒成為神的兒女,是因為神收養我們.被收養之後,我們的身分便被神所定義了.舊的已經過去,一切都更新了.

基督徒見證著,在基督裡,在聖靈的引領下,成為「天父的兒女」是一件真實的事情.作為「天父的兒女」,我們的確能活像基督.

如何活像基督?

羅8:17　既是兒女,就是承受產業的人,是承受神產業的人,且與基督一同承受產業──只要我們與他一同受苦,與他一同得榮耀.

活像基督的秘訣:「我們與他一同受苦」.這是最無懼地,最大膽地運用自己的自由.現代社會一般都重視個體的自由,所謂個體的自由,就是只要不傷害他人,不妨礙他人的自由,則個體可自由抉擇他想要的東西.按這原理,終極而言,他可選擇讓座或不讓座給一個有需要的人.但十架下的自由,就是自由地放下自己的權利,讓座給一個有需要的人,因為關愛這個有需要的人.在自由的社會裡,沒有法律可規範我們去愛一個人.但與基督一同受苦的人,卻因感受到基督的大愛,而自由地,甘心的去關愛他人.

愛到什麼程度呢?

德蘭修女訂出一個簡單的原則:「愛,直到你為此受到損傷(Love until it hurts).」這就是與基督一同受苦.

德蘭修女又說:「我體會到一種弔詭,就是當你去愛,直到你為此受到損傷時,你不會感到損傷,只體會到更多的愛.」(I have found the paradox, that if you love until it hurts, there can be no more hurt, only more love.)

這就是包藏在「十字架」裡面的弔詭.一個人學曉捨去生命時,才會得著生命.一個人因愛而放下自由時,他才活出真正的自由.一個人跟從基督走向死亡時,他才進入永生.最終極的弔詭是:十字架的羞辱,正是神至大的榮耀.

這就是經文說的:「只要我們與他一同受苦,與他一同得榮耀」.所謂「榮耀」,就是當歷史終結時,當一切真相大白時,當天國降臨時,基督的神性全然顯露.而在那時,我們卻有幸去分享基督的榮耀,因為在塵世的歷史中,我們分享了基督的生命,承擔了十架的擔子,參與了天國降臨的進程.這份「榮耀」,就是我們作為神的兒女,能承受的寶貴的產業.

人生活著,總有一種存活態度.隨從「肉體」的,敵擋神的,看似前途似錦,卻是死路一條.看似生命豐盛,卻只是奴隸一名.隨從「聖靈」的,活像基督的,看似十架死路,卻是生命大道.看似一無所有,卻能承受神無限的產業.你如何抉擇呢?起來,不是奴隸的你.

願你在基督裡,在聖靈的引導中,成為天父的兒女,得享真正的自由.

但願:榮耀歸於聖父,聖子,聖靈.阿們.



\newpage

\section{捉迷藏}
\label{sec:288}
\textbf{鄧瑞強博士 - 捉迷藏}
\newline
\newline
\hyperref[sec:287]{< < < PREV SERMON < < <}
~
\hyperref[sec:toc]{[返目錄]}
~
\hyperref[sec:289]{> > > NEXT SERMON > > >}


創世記 3:8-13
\newline
\begin{longtable}{cl}
\hline
\hline
章節 & 經文 (和合本修訂版)\\
\hline
3:8 & \begin{tabularx}{0.7\textwidth}{X} 天起了涼風，那人和他妻子聽見耶和華神在園中來回行走的聲音，就藏在園裡的樹木中，躲避耶和華神的面。 \end{tabularx} \\ \\ \relax
3:9 & \begin{tabularx}{0.7\textwidth}{X} 耶和華神呼喚那人，對他說：「你在哪裡？」 \end{tabularx} \\ \\ \relax
3:10 & \begin{tabularx}{0.7\textwidth}{X} 他說：「我在園中聽見你的聲音，我就害怕；因為我赤身露體，我就藏了起來。」 \end{tabularx} \\ \\ \relax
3:11 & \begin{tabularx}{0.7\textwidth}{X} 耶和華神說：「誰告訴你，你是赤身露體呢？莫非你吃了那樹上所出的，就是我吩咐你不可吃的嗎？」 \end{tabularx} \\ \\ \relax
3:12 & \begin{tabularx}{0.7\textwidth}{X} 那人說：「你賜給我、與我一起的女人，是她把那樹上所出的給我，我就吃了。」 \end{tabularx} \\ \\ \relax
3:13 & \begin{tabularx}{0.7\textwidth}{X} 耶和華神對女人說：「你怎麼會做這種事呢？」女人說：「那蛇引誘我，我就吃了。」 \end{tabularx} \\ \\
[1ex]
\hline
\hline
\end{longtable}
講題:捉迷藏　Hide and Seek

經文:創世記3:8-13

講員:鄧瑞強博士

場合:香港中文大學崇基學院禮拜堂

日期:2015年6月7日

各位弟兄姊妹,早安.

「捉迷藏」應該是一個很古老的遊戲.

「捉迷藏」的廣州話是「捉伊人」.在廣州話裡,「人」(耶魯注音:yan4)有時會唸作「因」(耶魯注音:yan1)(如:晚上不要一個「人」〔唸成「因」〕四圍走).「捉伊人」,一個古老遊戲的古老叫法.

在聖經裡,人跟神玩的第一個「遊戲」,便是「捉迷藏」,「捉伊人」.

創世記第三章一開頭,講到人犯第一個罪.

只要犯了第一次罪,便會一不做二不休,繼續犯罪下去.犯罪,變成人的某種天性.

人類的第一宗犯罪事件,便是「蛇」引誘人食「禁果」的故事.這犯罪事件,聽起來,不合乎常理,卻合乎人之常情.蛇引導人去留意「人神之間」有一條不應踰越的界線,進一步,蛇慫恿人越過這界線去得到某些人類現在沒有的東西.

人越過某些界線去得到某些自己現在沒有的東西,很合乎人之常情.但神既然創造人並為人設定界線,又講明這一切安排都「甚好」,人不相信神的安排是「好」的,並認為自己做得比神更好,這想法不合常理.

創世記說,「蛇」是狡猾的動物.牠講的話半真半假,混淆真相和謊言.牠有超高的說話技巧,引誘人以錯誤的方式去獲取虛假的成功.人渴望某種成功,因而被蛇的言辭迷惑,是合乎常情的.但人竟然不信任神,而信任比人還低下的蛇,這是不合常理的.人不信任神的真理,而去信任「蛇」一樣的低等生物的說話,是不合常理的.可悲的是,人就是寧願聽「蛇」的謊言.

人聽了「蛇」的說話,破壞了保護人性的圍欄,吃了「禁果」.人性變了,大地變了,歷史的軌跡也變了.

人不敢面對神,而玩起「捉伊人」的遊戲.這正如一個小孩子犯錯後,會「匿埋」,不願意父母找到他.

今日的講道經文:創世記3:8－13

創 3:8　天起了涼風,耶和華神在園中行走.那人和他妻子聽見神的聲音,就藏在園裡的樹木中,躲避耶和華神的面.

創3:9　耶和華神呼喚那人,對他說:「你在哪裡?」

創3:10　他說:「我在園中聽見你的聲音,我就害怕；因為我赤身露體,我便藏了.」

創3:11　耶和華說:「誰告訴你赤身露體呢?莫非你吃了我吩咐你不可吃的那樹上的果子嗎?」

創3:12　那人說:「你所賜給我,與我同居的女人,她把那樹上的果子給我,我就吃了.」

創3:13　耶和華神對女人說:「你做的是甚麼事呢?」女人說:「那蛇引誘我,我就吃了.」

人犯了罪,聽見神的臨近,便躲藏起來.

神問了一個問題:「你在哪裡?」這是神在聖經裡問的第一個問題.

在聖經希伯來文裡,「你在哪裡」只是一個字「ayeka」.這一個字,是神對人類發出的永恆詢問.這個詢問,是說明:神主動尋找人.

猶太教神學家Abraham Heschel寫了一本書,叫 “God in Search of Man”(中文譯本:《覓人的上帝》,山東大學出版社),作者說,大多數宗教講的上帝,都是沉默的,隱藏的,這個上帝並不關心人對他的尋索.但聖經講的上帝,卻是一個尋覓人的上帝.聖經講的全部人類歷史,一言以蔽之,就是「上帝尋覓人」的故事.神主動尋索迷途的人,是基督教最動人的主題之一.

神關心我們對他的尋索.我們的祈禱,是我們尋索神的一種行動,神關心我們對他的尋索,故神也在意我們的祈禱.甚至,在我們祈禱以先,神已著緊我們的需要.神尋覓人,甚至到了一個地步,他來到人間,尋找我們.

Abraham Heschel這樣寫:「當亞當和夏娃從神的面前躲藏時,主呼喚道:『你在哪裡?』這是一個反復發出的呼喚.它是一個細小聲音的微弱回聲,不是以言詞說出來的,不是用思想概念來傳達的,而是不可說且神秘的,就如同充滿全世界的榮光一樣不可說而神秘.它被包裹在沉默之中,隱晦而微弱,彷彿萬事萬物都是這個問題凝固了的回聲:『你在哪裡?』」

「信仰來自敬畏,來自對我們暴露在他面前所感到的畏懼,來自對回答上帝的挑戰的焦慮,來自意識到我們被召喚.宗教是由上帝的問題和人的回答構成的.……除非上帝提問,否則,我們的一切求索都是枉費心機.」【Abraham Heschel(海舍爾):《覓人的上帝》(濟南:山東大學出版社,2003),頁128】

神不斷地問人:「你在哪裡?」直到人回答神「我在這裡」的時候,人才從存在的隱蔽之處中走出來,成為一個真正的人.

我們活在絕對的漆黑之中,不知身在何處,不知要做什麼,不知往何處去.直到那一刻,上帝在呼喚:「你在哪裡?」就在那一刻,我們才找到座標,才知道方向,才明白有人在關心我們.當我們說「我在這裡」時,我們會重新找到希望,找到前進的方向,找到存活的意義.

上帝發問,我們被詢問,我們被尋索,我們被呼喚.在我們回應這詢問,尋索,呼喚時,我們會被召去正視我們存在的真正意義,被召去放下我們自身的個人考慮,被召去回答比我們大得多的創造主.

神的問題「你在哪裡」,向我們每個人發出.你活了這麼久了,你的生命去到什麼地步?你的心思在何處?你在人生的什麼領域,搞什麼東西?

想想:你現在在什麼地方?你走的路是一條什麼路?你繼續走下去,會去到什麼光景?你的善性還在嗎?你持守的真理還在嗎?你的靈魂還在嗎?

整個天地間,都充滿神的詢問:「你在哪裡」?

我們每時每刻,只要你的心專注,都會聽到神的詢問.

神固然知道我們在那裡,甚至他走在我們面前問我們.他之所以要問,是因為若他不問,我們會迷失.需要這「問」的,不是神,是我們.

神的「問」,迫亞當從隱藏之處走出來,從虛無的黑暗走出來面對神聖的光明,從罪的深淵走出來去正視罪的真相.

亞當從隱蔽之處走出來.

創3:10　他說:「我在園中聽見你的聲音,我就害怕；因為我赤身露體,我便藏了.」

亞當沒有正面回答神的問題,他解釋他為何隱藏.

他隱藏,是因為他害怕.他害怕,是因為他赤身露體.

其實,自亞當被造以來,他都是赤身露體的.以前他赤身露體,卻不感到害怕.現在他赤身露體,卻感到害怕,是因為他對身體的態度改變了.

身體,其實是人最原始的一個具體界限.我們不是一道靈氣,不是一個無軀殼的意識,我們是有身體的生命.神創造有身體的我們,讓我們以具體的方式將神的形象展現人間.當然,將神的形象全然展現人間的具體生命,就是耶穌的「道成肉身」.

身體,是人的界限,是人的限制,也顯示人的脆弱無能.強暴的人,就是不重視這個界限,越過這個界限,侵佔他人的身體.獨裁的暴君,就是不承認人性的限制,自視為神,要求人膜拜.損人利己的人,就是忘記自己的脆弱無能,以為憑著自己的手段,可令自己幸福.

蛇的誘惑,就是慫恿人越過人性應該守著的界限,用自己的手段去克服脆弱無能,以提升自己去到神的高度.蛇誘惑人誤用自己的身體,漠視身體所標示的界線,限制和脆弱無能.人意圖將自己改造成一個好像不再受身體限制的神明,一種無限的,沒法攔阻的,強大的純力量.

就在這時,神呼喚人,神問:「你在哪裡?」

神的臨在,意味著一個人不能取消的界限的臨在.亞當看看自己,他仍只是一個一無所有的軀體.他曾試圖超越人性的限制,去成為一個像神一樣的人.但現在看來,這是徒然的.但這徒然的嘗試,卻徹底改變了他的身體的態度.以前,他是赤身露體的,但他並不覺得這有什麼不好,這是守護人性的籬笆.但是,蛇誘惑了他,他打破這個籬笆.如今,這個籬笆仍在,這個人性的限制仍在,他卻看之為不好的.在神面前,他的整個赤露的身體展示出來了.這身體代表著限制與脆弱,這卻是他現在不願意接受的.一個人不願意接受也不願意正視的限制與脆弱被暴露,只會引起不自在的感受.這種不自在的感受,是一種「羞恥」感.羞恥的人不願意正視人,面對人,在人面前敞開自己.亞當不願意面對神,他隱藏起來.如今,神的詢問迫他走出來正視自己,正視自己應守護著的界線和限制.

創3:11　耶和華說:「誰告訴你赤身露體呢?莫非你吃了我吩咐你不可吃的那樹上的果子嗎?」

神詢問亞當:「誰告訴你,你是赤身露體的?」其實,亞當一直都是赤身露體的,故此,神的問題是追問:「是什麼令你覺得赤身露體是不好的?」或者說:「是什麼令你覺得人性的界限,限制,脆弱,是不好的?」神繼續說:「莫非你吃了我吩咐你不可吃的那樹上的果子嗎?」或者說:「莫非你不信任我,要越過界線,看看是否能令自己更好?」事實上,聽神的說話,守著自己的界線,接納自己的有限,有什麼不好呢?但是,人總是好奇,很想知道,在種種禁果背後,是否藏著某種神沒有賜給我們的快樂?

創3:12　那人說:「你所賜給我,與我同居的女人,她把那樹上的果子給我,我就吃了.」

亞當仍然沒有正面回答神的問題,他將矛頭指向其他人.

他埋怨女人,也埋怨將女人給他的神.這個女人,是神「給」我的.那個果子,是女人「給」我的.不同的人將不同的東西「給」我,我只是一個被動接受的人.亞當扮演受害者的角色.一切都與他無關.

亞當將自己隱藏得更深了.他沒有頂天立地的保護女人,他也沒有為神賜的這個伴侶而感謝神,當然,他也沒有為自己做的事而負上一個人的責任.這一切,都因為人不信任神,而信任一條比人還不如的蛇；因為人不願意守著界限而意想成為無限；因為人認為神做得不夠好,而認為自己做得比神好.女人,現在不是伴侶,而淪為代罪羔羊.上帝,不再是感恩的對象,而是人類禍患的根源.唯有自己是對的,不用向任何人負責.

這是不是很悲哀呢?

蛇每天都誘惑我們,而我們每天都試圖越過一切界線,然後,不負責任地說,這一切都不關我的事.

這是很悲哀的,這是我們的悲哀.

創3:13　耶和華神對女人說:「你做的是甚麼事呢?」女人說:「那蛇引誘我,我就吃了.」

女人的回答,明顯比男人的回答好得多.女人沒有埋怨神,沒有埋怨伴侶,沒有說一切東西都是別人給她的,沒有暗示自己只是受害者.女人很爽快地承認,「那蛇引誘我,我就吃了.」

習慣上,我們將所謂「原罪」的責任,加諸女性,但很明顯,這只是男人不負責任而諉過於女人的結果.面對神的詢問,女人正面地回答了,她從隱藏之處站了出來.而男人呢,他仍隱藏著,他隱藏在自大的自我背後.

創世記的故事,每天都在發生.人不斷反抗身體所代表的界限,限制與脆弱.權力的追逐,暴力的使用,自我的神化,每天都在發生.神像尋找浪子的父親,發出呼喚,追問:「你在哪裡」?我們卻像亞當一樣,與神玩「捉伊人」的遊戲,將自己隱藏起來.

今日,神也向你發出這詢問:「你在哪裡」?正面回應這詢問的人,會深刻領悟自己存在的意義.

但願:榮耀歸於聖父,聖子,聖靈.阿們.



\newpage

\section{我的女兒病重了}
\label{sec:289}
\textbf{伍渭文牧師 - 我的女兒病重了}
\newline
\newline
\hyperref[sec:288]{< < < PREV SERMON < < <}
~
\hyperref[sec:toc]{[返目錄]}
~
\hyperref[sec:290]{> > > NEXT SERMON > > >}


馬可福音 5:21-43
\newline
\begin{longtable}{cl}
\hline
\hline
章節 & 經文 (和合本修訂版)\\
\hline
5:21 & \begin{tabularx}{0.7\textwidth}{X} 耶穌又坐船渡到對岸，有一大群人聚集到他身邊；他正在海邊。 \end{tabularx} \\ \\ \relax
5:22 & \begin{tabularx}{0.7\textwidth}{X} 有一個會堂主管，名叫葉魯，也來了，一見到耶穌，就俯伏在他腳前， \end{tabularx} \\ \\ \relax
5:23 & \begin{tabularx}{0.7\textwidth}{X} 再三求他，說：「我的小女兒快要死了，求你去為她按手，使她痊癒，可以活下去。」 \end{tabularx} \\ \\ \relax
5:24 & \begin{tabularx}{0.7\textwidth}{X} 耶穌就和他同去。有一大群人跟隨他，擁擠著他。 \end{tabularx} \\ \\ \relax
5:25 & \begin{tabularx}{0.7\textwidth}{X} 有一個女人，患了經血不止的病有十二年， \end{tabularx} \\ \\ \relax
5:26 & \begin{tabularx}{0.7\textwidth}{X} 在好多醫生手裡受了許多苦，又花盡了她所有的，一點也不見好，反而更重了。 \end{tabularx} \\ \\ \relax
5:27 & \begin{tabularx}{0.7\textwidth}{X} 她聽見耶穌的事，就夾在眾人中間，從後面來摸耶穌的衣裳， \end{tabularx} \\ \\ \relax
5:28 & \begin{tabularx}{0.7\textwidth}{X} 因她想：「我只摸到他的衣裳，就會痊癒。」 \end{tabularx} \\ \\ \relax
5:29 & \begin{tabularx}{0.7\textwidth}{X} 於是她的流血立刻止住，她覺得身上的疾病好了。 \end{tabularx} \\ \\ \relax
5:30 & \begin{tabularx}{0.7\textwidth}{X} 耶穌頓時心裡覺得有能力從自己身上出去，就在眾人中間轉過來，說：「誰摸我的衣裳？」 \end{tabularx} \\ \\ \relax
5:31 & \begin{tabularx}{0.7\textwidth}{X} 門徒對他說：「你看眾人擁擠著你，還說『誰摸我』呢？」 \end{tabularx} \\ \\ \relax
5:32 & \begin{tabularx}{0.7\textwidth}{X} 耶穌周圍觀看，要見做這事的女人。 \end{tabularx} \\ \\ \relax
5:33 & \begin{tabularx}{0.7\textwidth}{X} 那女人知道在自己身上所成的事，就恐懼戰兢，來俯伏在耶穌跟前，將實情全告訴他。 \end{tabularx} \\ \\ \relax
5:34 & \begin{tabularx}{0.7\textwidth}{X} 耶穌對她說：「女兒，你的信救了你，平安地回去吧！你的疾病痊癒了。」 \end{tabularx} \\ \\ \relax
5:35 & \begin{tabularx}{0.7\textwidth}{X} 耶穌還在說話的時候，有人從會堂主管的家裡來，說：「你的女兒死了，何必還勞駕老師呢？」 \end{tabularx} \\ \\ \relax
5:36 & \begin{tabularx}{0.7\textwidth}{X} 耶穌不理會他們所說的話，就對會堂主管說：「不要怕，只要信！」 \end{tabularx} \\ \\ \relax
5:37 & \begin{tabularx}{0.7\textwidth}{X} 於是他帶著彼得、雅各和雅各的弟弟約翰同去，不許別人跟著他。 \end{tabularx} \\ \\ \relax
5:38 & \begin{tabularx}{0.7\textwidth}{X} 他們來到會堂主管的家裡，耶穌看到一片吵鬧，並有人大聲哭泣哀號， \end{tabularx} \\ \\ \relax
5:39 & \begin{tabularx}{0.7\textwidth}{X} 就進到裡面，對他們說：「為甚麼大吵大哭呢？孩子不是死了，是睡著了。」 \end{tabularx} \\ \\ \relax
5:40 & \begin{tabularx}{0.7\textwidth}{X} 他們就嘲笑耶穌。耶穌把他們都趕出去，帶著孩子的父母和跟隨的人進了孩子所在的地方， \end{tabularx} \\ \\ \relax
5:41 & \begin{tabularx}{0.7\textwidth}{X} 就拉著孩子的手，對她說：「大利大，古米！」翻出來就是說：「女孩，我吩咐你，起來！」 \end{tabularx} \\ \\ \relax
5:42 & \begin{tabularx}{0.7\textwidth}{X} 那女孩子立刻起來走動—她已經十二歲了；他們就非常驚奇。 \end{tabularx} \\ \\ \relax
5:43 & \begin{tabularx}{0.7\textwidth}{X} 耶穌切切地囑咐他們，不要讓人知道這事，又吩咐給她東西吃。 \end{tabularx} \\ \\
[1ex]
\hline
\hline
\end{longtable}
講題:我的女兒病重了　My Daughter Falls Severely Sick

經文:馬可福音5:21-43

講員:伍渭文牧師

場合:香港中文大學崇基學院禮拜堂

日期:2015年6月28日

你過去有沒有在別面前下跪的經驗?是否想到求婚那一刻,男士們?那是浪漫的回憶,恐怕是惟一的一次吧.你過去有沒有人,突然在你面前下跪,心情非常焦灼,懇求你幫助他們呢?我有兩年時間在台北一間基督教高校工作,其中一年作教務長；有一位來自南部的媽媽一進入我的辦公室就跪在我面前,希望我取錄她的女兒；她的突然行動,令我當時不知所措,我體會到母親對女兒的愛心.我當然沒有因為她下跪就改變決定,但給她打了電話聯絡了另一間學校,錄取了她的女兒,她女兒大學畢業後還特別前來道謝.

如果一個地位崇高的父親,能夠放下身段,在公眾地方眾目睽睽之下,向一個人下跪,懇求幫助,心中一定有極大的催迫,今天的福音經課就提到這樣的一個父親.「我的女兒快要死了,求你去按手在她身上,使她痊愈,得以活了.」(可五23)我的女兒病重了,快來救她!

今年經課選讀的符類福音是馬可,去年是馬太,明年是路加,約翰每年分散放進去,之後又是新的循環.馬可福音是公認最早的福音書,大概公元65年完書,90\%的馬可內容在馬太找到,50\%在路加找到；即馬太和路加參考馬可,加上其它資料寫成他們的福音書.馬可福音的特色是明快,耶穌很多行動,強調耶穌的權能:「上主的兒子,耶穌基督福音的起頭.」(可一1)看馬可如同看實況記錄,很有電影感.馬可被認為是根據彼得口授筆錄寫成的,彼前五13「我的兒子馬可也問你們安.」今天的福音經課是典型馬可的實況記錄,短短二十二節經文,生動記錄了耶穌答允父親的懇求,醫她他的病重或準確的說,使剛死去的十二歲女兒復生；在往女兒家的途中,也醫好一位患了十二年血漏病的婦人.我們先看這位愛女情深的父親睚魯.

1              焦灼的父親

管會堂睚魯:睚魯是管會堂的,是非常有地位名望的人,不是看守會堂,開門關門.管會堂是治理會堂長老會的主席,監督會堂崇拜禮儀的進行,相當現在香港猶太會堂的管理委員會主席,是香港猶太社群的意見領袖,地位崇高.作為猶太人的意見領袖,睚魯一定非常關注耶穌所作所為,格拉森二千隻豬被耶穌驅逐的群鬼附身,投涯墮海淹死,被鬼附身的人清醒過來,留在社區並在附近的鄉鎮傳講耶穌的大能,睚魯一定聽過,甚至見過這人,格拉森就在鄰近,「耶穌坐船又渡到那邊去」(可五21).

睚魯是經過一番掙扎才找耶穌的,身為管會堂的猶太人意見領袖,公開請求耶穌到家中醫治女兒,是政治不正確的,因為猶太人拒絕耶穌.約翰福音九章生來瞎眼的瞎子被耶穌醫好,他的父母也怕公開支持耶穌,說:「他是我們的兒子,生來是瞎眼的,如今怎麼能看見,我們卻不知道；是誰開了他的眼睛,我們也不知道.他已經成了人,你們問他吧,他自己必能說.」(約九20, 21)聖經說:「他父母說這話,是怕猶太人,因為猶太已經商議定了,若有人認耶穌是基督的,要把他趕出會堂.」(約九22)但管會堂的睚魯也是父親睚魯.

父親睚魯:一個有女兒的父親,是認識了兩個和自己生理,心理完全不同的女人.因為愛,他跟第一個女人-妻子-走在一起；也因為愛,讓第二個女人-女兒-來到世界.耶穌稱女兒為閏女:「『閏女,我吩咐你起來!』那閏女立時起來走,他們就大大驚奇.閏女已經十二歲了.」(可五41-43)猶太女孩到了十二歲就被視作成人了,半年後就配婚,再過兩年,就出嫁了.十二歲的女孩還是待字閏中,還有半年時間,就有不同的身份了.

這時候的閏女,父母最愛,因為很快女兒就要離開了.這時候,若女兒是獨女,在她心中,父親還是惟一的男人,從父親相處過程中,女兒也認識和自己生理,心理不同的男人,她也將會一生和自己喜歡的男人-丈夫-生活.

這位閏女是頂幸福的女孩,不單是父母親的掌上明珠,可能不少家庭正盯著父親是管會堂的女孩,想自己的兒子娶過來作妻子,但她的重病,令多個家庭心情焦慮,這閏女正命懸一線；多人的盼望,特別是父親睚魯,就在耶穌身上.此刻耶穌不再是睚魯注意的新興宗教運動的領袖,是拯救他女兒的惟一盼望了.耶穌也答應了睚魯的懇求,跟隨他前往他的家,因他說:「求你按手在她身上,使她痊愈,得以活了.」

2              坎坷的女兒

但往睚魯的家途中,出現另一個「女兒」,孤單的女兒,生命坎坷的女兒.經文說:「有一個女人,患了十二年的血漏,在好些醫生手裡受了許多的苦,又花盡了一生所有的,一點也不見好,病勢反倒更重了.」(可五25, 26)經文說有一個女人,因為她父母不在場,也沒有說她有丈夫.「有一個女人」的稱呼,是上主給予的身份；女兒,妻子是和人的關系.但上主給她的身份,帶給她特有的婦科病血漏,若是生為男生就沒有這病了.她不單身體痛楚,而且給信賴的醫生欺騙,花盡所有的錢,但病情愈來愈嚴重.

人群擁擠著耶穌,看看睚魯的掌上明珠如何將被耶穌醫好；睚魯有社會地位,但這婦人卻花盡所有的,病情卻每況愈下,但誰來為自己申訴呢?睚魯的女兒一出生,自己的病就開始；睚魯的女兒漸漸長大,自己的病也漸漸嚴重；睚魯的女兒愈來愈得人喜愛,自己愈來愈被多人拒絕,因為患血漏在當時被視為不潔,根據利未記十五章,連坐過的椅子也被視為不潔,不能參加會堂的崇拜,不能進入聖殿獻祭,被社會隔離.

她惟一的盼望也是耶穌:「她聽見耶穌的事,就從後頭來,雜在眾人中間,摸耶穌的衣裳,意思說:『我祇摸她的衣裳,就必全痊愈.』」(可五28)她知道不可能像睚魯的女兒,期望耶穌會按手在自己身上,根本人太多了,埋不到身.她想,祇要摸到耶穌衣裳就可以了.她拼命擠進人群,摸到耶穌衣裳的襚子；立時,血漏的源頭乾了,她感到十二年的病,霍然而愈!

公開見證

耶穌頓時覺得有能力從自己身上出去,就在眾人中間轉過來,說:「誰摸我的衣裳?」門徒很詫異耶穌這樣問,對耶穌說:「你看眾人擁擠你,還說『誰摸我』嗎?」(可五31)耶穌知不知道誰摸他的衣裳呢?耶穌應該知道,三十二節說:「耶穌周圍觀看,要見作這事的女人.」

耶穌問這問題,有幾個目的:

想這得醫治的女人明白,是她對耶穌的信心,不是耶穌的衣裳醫好她；我們不是追求上主的權能,乃信靠滿有權能的基督.

想和她建立個人的關系:「女兒,你的信救了你.」她公開稱她為女兒,想她知道,睚魯如何愛她患重病的女兒,上主也一樣愛長期患血漏的她,她是上主的女兒,不單是一個女人.

讓她公開被接納為潔淨的人,重新歸入社群；從此可以進入會堂敬拜,進入聖殿獻祭.

堅固睚魯的信心,因為跟著有從家裡來的人報訊:「你的女兒死了,何必還勞動先生呢?」(可五35)眶魯聽了,甚麼話都說不出來,懼怕頓時牢牢抓著他,耶穌對他說:「不要怕,祇要信!」這婦人因為信得到醫治,你女兒雖然死了,你若信,也可以使她復生.

不要怕,祇要信!

睚魯怕甚麼呢?他怕失去愛女,回到家中,女兒已不能呼吸,手腳冰冷.死亡令人懼怕,因為它奪去我們所愛的人,沒有女兒的生活會如何呢?此刻,睚魯的情緒由焦灼轉為懼怕,耶穌看出他情緒的改變,就把眾人回絕,不許別人跟隨他.之前聖經說:「有許多人跟隨,擁擠他」,現在祇帶從彼得,雅各和約翰,其它的門徒也一概留下來,不許跟隨前往睚魯的家.耶穌了解睚魯的心情,此刻睚魯需要安靜,去消化女兒去世的消息,他的信心正經歷極大的試煉.耶穌和三個門徒陪著作父親的,疾步走向他的家.

上主日我應邀到夏門一間教會講道,崇拜完畢,一位太太留下來飲泣,要我為她祈禱.她第一次來那間教會聚會,因為當天突然感到很悲傷,前所未有的悲傷,想起一個朋友就來了.她說一個月前一次車禍,丈夫,弟弟,舅舅都離開了,她要回家鄉生活,陪伴兩個小孩子,但擔心家姑家公不疼愛她.聽她這樣一說,我突然沒言以對,緊握著她的手去感受她的恐懼和悲傷,半分鐘後跟她說:上主明白我們的感受,面對這樣的遭遇,悲傷很正常,就讓自己放懷悲傷,不用壓抑,耶穌已接收了屬於他的子民的靈魂,聖經說:「我流淚,你用皮袋裝.」(詩五十六8)最重要回到家鄉尋找教會聚會,得到支持.我們幾個人圍著她祈禱,把她交托上主.

耶穌陪同睚魯到了他的家,看到人們面對死亡,顯得無助,祇得亂嚷哭泣哀號,耶穌說:「孩子不是死了,是睡著了.」他們就嗤笑耶穌,耶穌把他們都趕出去,祇帶父母和三位門徒進到孩子所在之處,「就拉著孩子的手,對她說:『大利大古米!』(翻出來就是說:閏女,我吩咐你起來!)那閏女立時起來走,他們就大大地惊奇.」(可五41, 42)耶穌使死去的女兒活過來.

睚魯面對女兒的死亡,他很懼怕,耶穌陪著他走,並使死去的女兒復生.今天我們也會因為關系的死亡變得懼怕,最近收到一位信徒給我的信,我徵求的的同意,和大家分享:

有一個問題想知你怎看……

自從經過苦難後,我發現自己的人生時鐘較快了許多,總覺得自己好快離開世界,唔敢慢下來,最嚴重的,是覺得很多要用時間處理的事情,我看得很灰,包括建立關係,培養興趣等……重新起步的難處,就是要重新相信,生命可以再來……真的嗎?還有本錢再輸嗎?輸感情,輸失敗……突然對際遇有很多unknown…..究竟約伯是如何重新開始?

十年感情突然夭折,將近五年了,還是不敢和人建立關系,懼怕花時間心血建立起的東西,也會突然失去.我回信說:你的傷痛未癒,求上主醫治你；關系如人,死了就要埋葬,開始新生.

應用:今天你有甚麼懼怕?耶穌說:「不要怕,祇要信!」耶穌陪著我們走,醫治我們,使我們經歷新生.

我的女兒病重了

今天講道題目:我的女兒病重了.誰是病重的女兒?當群眾擁擠耶穌,當群眾跟隨睚魯和耶穌走向睚魯的大宅,這女兒在人心中毫無懸念是那臥病在床,頻臨死亡十二歲的閏女,這女兒為父母所疼愛,為眾多家庭所賞識,欲收為媳婦的女孩,她是眾人心中的女兒,她的病牽動不少的家庭,她是眾人的焦點.用今天的情況類比,這閏女因為父親是管會堂的,常常成為媒體談論的對像:那一家人希望娶他為媳婦,那一個男孩才是可能的白馬王子；因為家庭背景,她自幼就站在不同的起跑線,是城中名媛,連病重了都成為媒體關注的焦點.也許,這閏女已被媒體看成是這城市的女兒了.

但在耶穌眼中,被人忽略,被頑疾折磨了十二年,花盡財富,血漏病反而愈來愈差的婦人,也是女兒,雖然沒有人訪問過她,她沒有新聞價值,不曾出現在媒體.但耶穌對這可憐的婦人說:「女兒,你的信救了你!」閏女是上主的女兒,所以耶穌去她的家醫治她；血漏病的婦人也是上主的女兒,所以在途中耶穌停下來醫好她.

應用:所有需要幫忙的人都是上主疼愛的兒女,那些被社會遺忘的人,承受人生諸般苦難的人,上主憐愛他們猶如憐愛自己的兒女,我們要關心他們.

眾人擁擠,惟她觸摸

今天,誰是上主的女兒?誰能得到上主的醫治?群眾擁擠耶穌,群眾中可有人有疾病,身體或心靈的疾病?他們有否想得從耶穌身上得到醫治,得到屬靈的祝福?他們擁擠著耶穌,那麼接近耶穌了；但聖經記載,祇有這位患血漏病的婦人得到醫治,因為她渴慕得醫治,她渴望得到痊愈.

應用:耶穌說:「女兒,你的信救了你.」今天,祇要我們有信心,認識自己的身體,靈魂的需要,也可以得到醫治.

八福第一福是虛心的人有福了,因為天國是他們的；次福為哀慟的人有福了,因為他們必得安慰.在上主面前虛心,知道虧欠,進而有屬靈的覺醒,為自己的虧欠哀慟,這是蒙福的起始點.有多大的器皿,就能盛載多大的恩典.眾人擁擠,惟她觸摸,她感到需要上主,信靠上帝,就經歷上主的能力,得到痊愈.



\newpage

\section{我們被差遣}
\label{sec:290}
\textbf{賀志勇博士 - 我們,被差遣!}
\newline
\newline
\hyperref[sec:289]{< < < PREV SERMON < < <}
~
\hyperref[sec:toc]{[返目錄]}
~
\hyperref[sec:291]{> > > NEXT SERMON > > >}


馬可福音 6:7-13,30
\newline
\begin{longtable}{cl}
\hline
\hline
章節 & 經文 (和合本修訂版)\\
\hline
6:7 & \begin{tabularx}{0.7\textwidth}{X} 他叫了十二個使徒來，差遣他們兩個兩個地出去，也賜給他們權柄制伏污靈， \end{tabularx} \\ \\ \relax
6:8 & \begin{tabularx}{0.7\textwidth}{X} 並且吩咐他們：途中不要帶食物和行囊，腰袋裡也不要帶錢，除了手杖以外，甚麼都不要帶； \end{tabularx} \\ \\ \relax
6:9 & \begin{tabularx}{0.7\textwidth}{X} 只要穿鞋子，也不要穿兩件內衣。 \end{tabularx} \\ \\ \relax
6:10 & \begin{tabularx}{0.7\textwidth}{X} 他又對他們說：「你們無論到何處，進哪家，就住在哪裡，直到離開那地方。 \end{tabularx} \\ \\ \relax
6:11 & \begin{tabularx}{0.7\textwidth}{X} 若有甚麼地方的人不接待你們，不聽你們，你們離開那裡的時候，要跺掉你們腳上的塵土，證明他們的不是。」 \end{tabularx} \\ \\ \relax
6:12 & \begin{tabularx}{0.7\textwidth}{X} 使徒就出去傳道，叫人悔改， \end{tabularx} \\ \\ \relax
6:13 & \begin{tabularx}{0.7\textwidth}{X} 又趕出許多鬼，用油抹了許多病人，治好他們。 \end{tabularx} \\ \\ \relax
6:14 & \begin{tabularx}{0.7\textwidth}{X} 耶穌的名聲傳開了，希律王也聽見。有人說：「施洗的約翰從死人中復活了，因此才有這些異能在他裡面運行。」 \end{tabularx} \\ \\ \relax
6:15 & \begin{tabularx}{0.7\textwidth}{X} 但別人說：「他是以利亞。」又有人說：「是先知，正如先知中的一位。」 \end{tabularx} \\ \\ \relax
6:16 & \begin{tabularx}{0.7\textwidth}{X} 希律聽見卻說：「是我所斬的約翰，他復活了。」 \end{tabularx} \\ \\ \relax
6:17 & \begin{tabularx}{0.7\textwidth}{X} 原來，希律為他兄弟腓力的妻子希羅底的緣故，派人去抓了約翰，把他綁了在監獄裡，因為希律已經娶了那婦人。 \end{tabularx} \\ \\ \relax
6:18 & \begin{tabularx}{0.7\textwidth}{X} 約翰曾對希律說：「你佔有你兄弟的妻子是不合法的。」 \end{tabularx} \\ \\ \relax
6:19 & \begin{tabularx}{0.7\textwidth}{X} 於是希羅底懷恨他，想要殺他，只是不能。 \end{tabularx} \\ \\ \relax
6:20 & \begin{tabularx}{0.7\textwidth}{X} 因為希律怕約翰，知道他是義人，是聖人，所以就保護他，雖然聽了他的講論十分困惑，仍然樂意聽他。 \end{tabularx} \\ \\ \relax
6:21 & \begin{tabularx}{0.7\textwidth}{X} 有一天，恰巧是希律的生日，希律擺設宴席，請了大臣、千夫長和加利利的領袖。 \end{tabularx} \\ \\ \relax
6:22 & \begin{tabularx}{0.7\textwidth}{X} 他的女兒希羅底進來跳舞，使希律和同席的人都很高興。王就對女孩說：「無論你要甚麼，向我求，我都會給你」； \end{tabularx} \\ \\ \relax
6:23 & \begin{tabularx}{0.7\textwidth}{X} 又對她多次起誓說：「無論你向我求甚麼，就是我國家的一半，我也會給你。」 \end{tabularx} \\ \\ \relax
6:24 & \begin{tabularx}{0.7\textwidth}{X} 她就出去對她母親說：「我該求甚麼呢？」她母親說：「施洗約翰的頭。」 \end{tabularx} \\ \\ \relax
6:25 & \begin{tabularx}{0.7\textwidth}{X} 她就急忙進去見王，求他說：「我願王立刻把施洗約翰的頭放在盤子裡給我。」 \end{tabularx} \\ \\ \relax
6:26 & \begin{tabularx}{0.7\textwidth}{X} 王就很憂愁，然而因他所發的誓，又因同席的人，不願食言， \end{tabularx} \\ \\ \relax
6:27 & \begin{tabularx}{0.7\textwidth}{X} 就立刻派一個衛兵，吩咐拿約翰的頭來。衛兵就去，在監獄裡斬了約翰， \end{tabularx} \\ \\ \relax
6:28 & \begin{tabularx}{0.7\textwidth}{X} 把頭放在盤子裡，拿來給那女孩，她就給她母親。 \end{tabularx} \\ \\ \relax
6:29 & \begin{tabularx}{0.7\textwidth}{X} 約翰的門徒聽到了，就來把他的屍體領去，放在墳墓裡。 \end{tabularx} \\ \\ \relax
6:30 & \begin{tabularx}{0.7\textwidth}{X} 使徒們聚集到耶穌那裡，把一切所做的事、所傳的道全告訴他。 \end{tabularx} \\ \\
[1ex]
\hline
\hline
\end{longtable}
講題:我們,被差遣　We Are Sent Forth!

經文:馬可福音6:7-13,30

講員:賀志勇博士

場合:香港中文大學崇基學院禮拜堂

日期:2015年7月5日

6:7          耶穌叫了十二個門徒來,差遣他們兩個兩個的出去．也賜給他們權柄,制伏污鬼．

6:8          並且囑咐他們,行路的時候,不要帶食物和口袋,腰袋里也不要帶錢,除了拐杖以外,甚麼都不要帶．              6:9                只要穿鞋．也不要穿兩件掛子. 6:10     又對他們說,你們無論到何處,進了人的家,就住在那裡,直到離開那地方. 6:11             何處的人,不接待你們,不聽你們,你們離開那裡的時候,就把腳上的塵土跺下去,對他們作見證. 6:12                門徒就出去,傳道叫人悔改.6:13           又趕出許多的鬼,用油抹了許多病人,治好他們.

6:30        使徒聚集到耶穌那裡,將一切所作的事,所傳的道,全告訴他.

引言 被差遣的馬丁路德金

1950年代的美國有許多不可思議的事情,比方說黑人與白人不能在同一家餐廳吃飯,不能在同一所公園玩耍；甚至在公車上,黑人與白人也要坐不同的座位.1955年,當一位黑人申請入讀密西西比大學時,有關當局把他關進了精神病院:因為只有精神有問題的黑鬼才會想到密西西比大學來讀書. 羅莎夫人的出現改變了這一切,1955年,羅莎夫人拒絕在公車上給白人讓座,她已經厭倦了這種毫無自尊的屈服,她因此被告上法庭.而之後被民權運動推舉出來幫助羅莎夫人的人叫做馬丁.路德.金,他是一名新任的浸信會牧師,當時才26歲,看起來像個男孩多過像個男人,當他被推舉出來的時候,他有點不知所措,因為他實在不覺得自己可以擔當如此大任.

剛當選沒幾天,三K黨就宣佈要對付他,而他也莫名其妙因為開車超速的原因被扔進了監獄.在被放出來那天夜晚,他在廚房裡想自己是不是要找些理由辭職--還是專心去做學術吧.幾分鐘之後他接到一個電話,要他趕緊滾出這個城市,否則就要炸掉他的房子.他很驚恐,在禱告當中,他對主講:「主啊,我必須承認我現在很軟弱,我正在顫抖,我正在失去我的勇氣.」但之後他聽見了耶穌的聲音,那聲音告訴他說,要繼續徵戰,主應許永遠不會離開他,永遠不會只剩下他一個.三天後,一個炸彈落在馬丁家裡,但他平靜面對這一切,他說:「我三天前的宗教經驗讓我知道如何平靜.」 之後十三年當中,馬丁路德曾有過其他宗教經驗,但沒有一次能跟那個夜晚的經驗相比,因為在那個夜晚,他感受到:主差遣了他,主一直與他同在.

解經

這一段驚心動魄的歷史讓我想起另外一段歷史,那段歷史發生在2000年前,那一次,被主差遣的是12使徒.主在自己的家鄉被拒絕之後,他開始把門徒差遣出去,向更廣闊的地方去宣告福音.

一,主差遣:應許與團契

6:7          耶穌叫了十二個門徒來,差遣他們兩個兩個的出去．也賜給他們權柄,制伏污鬼．

馬可6章七節說主賜給門徒制服污鬼的權柄,如果你細心的話,你會發現三章15節,主已經應許要賜給門徒這一權柄.第7節和12節以及13節說明:主的應許從來沒有落空過,當主應許了,他就真的會給你!亞伯拉罕之所以70歲了還要舉家離開吾珥,摩西八十歲的時候還帶著以色列人離開埃及,是因為他們相信主的應許:神為他們預備了更美好的家鄉,並且,神與他們同在!

我常常在想,如果亞伯拉罕有高血壓,摩西有心臟病,他們還會不會接受差遣.聖經沒有告訴我們他們的體檢報告,但聖經卻告訴我們他們一路遭遇的艱難.如果不相信神的應許從不落空的話,他們這一把年紀,這麼拼命是做什麼?聖經的故事說明,被差遣絕對不是一件浪漫的事--但做基督徒不就是這樣麼!做一名基督徒,不就是被主呼召和差遣,不怕艱難在這個世界為他做見證麼?當馬丁路德金的頭顱被打過十多次,當他宣揚和平的抗爭卻被黑人自己趕下台時,當他堅持對神公義的信心卻一直看不到結果時,他卻依然堅信: 「主差遣了我,祂的應許永不落空」,正是這樣一種信心,支持馬丁路德金跨越最黑暗的十三年,直到黑人真正獲得民權!

不過,主呼召並差遣使徒,並沒有讓他們一個人上路,他給使徒預備了同伴,讓他們兩兩出去.我認識一些禱告的同伴,他們常常跟我分享他們的禱告,雖然他們不是經常經歷神跡,主更不會每次都答應他們的請求,但主經常從不同角度,讓他們經歷到主意想不到的帶領.每次聽他們講到他們如何禱告,神如何回應和帶領,他們熱情洋溢的信心就讓我深深被激勵.弟兄姊妹,你有沒有這樣的同伴呢?

二,主差遣:趕緊上路,滿有信心

6:8 並且囑咐他們,行路的時候,不要帶食物和口袋,腰袋里也不要帶錢,除了拐杖以外,甚麼都不要帶．6:9只要穿鞋．也不要穿兩件掛子.

使徒們就要上路了,或多或少有點激動吧.但看到8-9節中主耶穌吩咐他們要帶的東西,估計他們就激動不起來了.我想起我這次去美國訪學,連日搬家實在太辛苦啦,所以我發誓出門盡量少買東西,少帶東西.正如朋友講的,帶上Credit卡和美金就好了.

可是當年的使徒沒有Credit卡啊,而且主耶穌還吩咐他們不要帶錢；不要帶食物,只帶拐杖,鞋子,那他們出去吃什麼,用什麼呢?只帶一件衣服,怎麼換衣服啊?主派這些人出去為他爭戰,卻連生活必需品也不給他們預備,你會不會奇怪這是什麼樣的一段旅程啊.可是當年以色列人的先祖也經歷了這樣的旅程,出埃及記12章11節這樣記載了耶和華的命令:「你們吃羊羔當腰間束帶,腳上穿鞋,手中拿杖,趕緊地吃.」 為什麼耶和華讓以色列人只穿鞋和拿拐杖?因為以色列人要馬上趕路,因為埃及人會反悔,會很快追上來.正如12章39節說的「因為他們被催逼離開埃及不能拖延,也沒有為自己預備什麼食物.」

哦,原來是這樣!主耶穌的吩咐跟當年耶和華的吩咐是一樣的,當年以色列人緊急離開了埃及,這說明現在門徒傳福音的使命也何等的緊迫.主耶穌在世的時間沒有多久了,因此,主要差派他們趕緊出去--主耶穌其實想告訴他們:」莊稼熟了,不去收割就會爛掉啦!而且,趁著我還在的時間,你們趕緊去磨練自己吧,這樣哪怕你們有跌倒和疑惑,我還可以親自指點你們,帶領你們.你們只有在事奉中才能學習成長,才能親身體會我的應許是真的.這樣,哪怕有一天我不在了,你們受過了磨練,就可以把我的使命繼續下去.「

我常常聽人說,「現在我很忙,等我退休了再來事奉吧.」但這幾年的事奉讓我明白,事奉不能拖延.因為只有當我們落手落腳,獨立去事奉主,我們才能磨練自己的信仰,而只有我們的信仰經過了磨練,我們才會真正信靠「主給我們的能力是真的,主說的話不會落空.」 正如馬可6章30節所說,使徒們傳道回來,都很興奮,將一切所做的事,所傳的道都說給主耶穌聽.我能夠明白那樣一種興奮,每一次我看到有人決志信主時,我也感嘆主的大能真是不可思議.

主的能力何等奇妙!當以色列人兩手空空地出埃及時,他們不曉得將來要如何活下去；可是主的大能養活了他們四十年；同樣,使徒們出去時,他們也不知道靠什麼生活,主耶穌要讓他們明白,門徒沒有Credit卡,沒有美金可以倚靠,主的能力將會在旅程中保護他們,養活他們.

三,主差遣:款待與離開

6:10          又對他們說,你們無論到何處,進了人的家,就住在那裡,直到離開那地方. 6:11何處的人,不接待你們,不聽你們,你們離開那裡的時候,就把腳上的塵土跺下去,對他們作見證.

馬丁路德金在他生命中最後十三年,就深深體驗了那奇妙的能力.雖然主並不保證他道路平坦,無災無難,可是主安排了許多接待他的人.有很多人向他伸出了友情之手,不僅有黑人,還有很多白人,其中有不少是無神論者.馬丁最親密的顧問中就有共產黨員.他接受了各種各樣的款待,他並不去區分所謂的政治觀點,只要那些人肯認同真理與公義,他就接受他們的款待.不抵擋真理的就是朋友.

連主自己揀選的門徒,也各有各的政治觀點了.當我們太過於堅持自己的身份,政治觀點時,我們反而忽略了主給我們的任務是勸人與主和好.所以第10節主告誡門徒,傳福音的過程中不要挑三揀四,不要在乎對方是不是一定是以色列人.如果有人願意接待,願意聽他們講真理,他們就需要在那裡住下.

當然,我們要預料到,肯定有不肯接待我們的人.馬丁就曾經痛心地說過:「當我們的運動向南方白人牧師發出直接的呼籲時,他們大部分都架起雙手--有些甚至起來反對我們.」 我也常常聽到有人痛苦地說,自己傳福音給別人,別人卻遲遲不肯接受.這時,我就會想起主耶穌在11節的教導.當以色列人從外邦回來時,就要先跺掉腳上的塵土,一方面表明自己潔淨了,另一方面也傳達了主的審判.主其實告訴我們,傳福音也好,事奉也好,無需糾結於個人的得失成敗,只要我們已經忠心盡力,其餘的交給主就好了.我們只需跺掉腳上的泥土,告訴他們主的旨意,主自己會施行他的公義,因為」申冤在我,我必報應.」

結論:我們,被差遣!

在馬丁路德金被暗殺之前的最後一個晚上,他做了這樣一個演講,他說:「就像其他人一樣,我也想長命百歲.長壽有它的位置,但我現在已不再在乎.我只想完成上帝的旨意,而祂容許我上到山上,從那裡俯視,讓我看見那應許之地.......所以我今晚很高興.我心中一無掛慮.我不懼怕任何人,因為我的眼睛已經看見上主榮耀的來臨.」

馬丁用他一生和平的態度,堅定的信心,告訴我們上主的應許是真的.他是這現代世界的使徒.弟兄姊妹,我們不也一樣,我們不也是耶穌基督的精兵,不也是蒙差遣的麼?我相信,如果我們不是僅僅享受教會之內的溫情,而是落手落腳地去事奉,去見證福音,我們就會更深刻地體會主的帶領,主的應許.「出於神口裡的話,沒有一句是不具有能力的.」 而當我們經歷過差遣之後,我們才會像馬丁路德金那樣大聲告訴自己:「我心中一無掛慮,我不懼怕任何人,因為我的眼睛已經看見上主榮耀的來臨.」

所以,請牢記,我們並非無名小卒,我們乃是被差遣的精兵.



\newpage

\section{上帝你還在嗎}
\label{sec:291}
\textbf{張美珍博士 - 上帝,你還在嗎?}
\newline
\newline
\hyperref[sec:290]{< < < PREV SERMON < < <}
~
\hyperref[sec:toc]{[返目錄]}
~
\hyperref[sec:292]{> > > NEXT SERMON > > >}


講題:「上帝,你還在嗎?」　God, Are You Still There?

講員:張美珍博士

場合:香港中文大學崇基學院禮拜堂

日期:2015年7月12日

各位弟兄姊妹早晨!主內平安!

今日是禮拜堂自四月底關閉裝修後,首次再在此舉行崇拜.請勿誤會,我的講題,並非因為我們離開了兩個多月,而要問上帝還在這禮拜堂嗎；因為敬拜上帝,正如在約翰福音所說,不在山上,也不在耶路撒冷,不受時間空間限制,而是在乎心靈與誠實.

今天這個講題:「上帝,你還在嗎?」對很多信徒來說,這可能是一個最簡單不過的問題；大家既然來這裡參加崇拜,很明顯是相信上帝的存在；但是,這問題於我來說,這卻是一個相當沉重的問題.

今日所讀的福音經課:可六:14-29,記述了為耶穌開路的先知──施洗約翰,亦是一位上帝所重用的僕人,被希律王因為一句戲言而斬首,以一個悲慘的結局為他的人生劃上句號.一位不畏強權,為耶穌鋪路的先知,怎麼會落得如此悲慘的下場?上帝,你還在嗎?

過去這一年來,伊斯蘭國這個極端宗教組織,綁架了不少基督徒,有記者,遊客,以及平民,公然斬首,並在處決前拍下片段在全世界播放；每看見這些新聞,我都想起施洗約翰.二千多年前這種野蠻的殺人方式,竟然在文明社會重現.每見到這些片段,我都不期然會問:上帝,你還在嗎?

這問題反映了我背後的假設,就是假如上帝還在的話,這些可怕的事件,是否就不會發生?每天打開報章,所看的國際新聞都離不開天災人禍,充斥著血腥與暴力的照片；本地新聞,若沒有重大意外事件發生的話,就離不開議事廳的爭拗,對政治的批判及對官員的謾罵.為何我們的世界,我們生活的社會,會變成這樣?上帝,你還在嗎?

人生,有時都幾荒謬.我們活在一個不可理喻的世界.

大約一個月前,多謝白耀燦師兄的邀請,看了由他主演的話劇﹝終局﹞,在導賞手冊中有這樣的一段說話:

「第二次世界大戰粉碎了多少美好的家庭,也粉碎了人類可以幸福地生存在世上的夢想；美夢變成夢囈,人間變成另一個地獄……,人本身就是一個地獄.

人們開始重新檢視自身的生存狀態,重新評估一切價值,他們質疑著:神去了哪裡?為何造物主要把人捨棄在這個千瘡百孔的世界裡?人存在有什麼意義?人應往何處去?」

我本人沒有經歷過世界大戰,但先父在國內經歷過抗日戰爭,是抗日英雄；其後經歷國共內戰,在四九年逃難至香港時,由家境富裕驟變至一貧如洗,由天堂跌落地獄,令我對於《終局》這段說話產生強烈共鳴.是的,上帝去了哪裡?為何造物主要把人捨棄在這個千瘡百孔的世界裡?人存在有什麼意義?人應往何處去?

其實這些問題,早在我三十多年前進神學院念書時,已不斷在我腦海中盤旋.可能少時嚐過貧窮的滋味,經歷過顛沛流離的生活,這些問題在我年紀還很輕時便已出現,令我比同齡人變得早熟.嘗試在神學院找尋答案的當兒,影響我最深的,是潘霍華(Dietrich Bonhoeffer)的神學,以及《對神誠實》(Honest to God)這本書.此書的作者是當時英國倫敦Woolwich區聖公會主教魯賓遜(John A.T. Robinson),他將當代三位歐洲基督教神學家的思想結合起來,包括有潘霍華的「去宗教化的基督信仰」(Religionless Christianity),田立克(Paul Tillich)的「存有的基礎及存在的深處」(Ground of beings and Depth of existence),布特曼(Rudolf Bultmann)的「去神話化」(Demythologization)等觀點融會貫通,對傳統基督教及上帝存在的問題,進行徹底的反思,讓我們在俗世中更認識我們的上帝.

《對神誠實》一書認為傳統基督教信仰,將上帝放在天上(up there)或外面(out there),是不正確的,"God is no longer a daddy in the sky". 上帝應該是活在人中間或內心的深處,與我們並存的.魯賓遜主教認為傳統的基督教信仰已經過時,對上帝的觀念應有新的看法,

潘霍華(1906-1945)是一位德國牧師及神學家,由於反抗希特拉政權而被捕,最後以絞刑被處決,離世時還不到40歲.他在獄中兩年多所寫的書信,被輯成一本書,中文譯名為《獄中書簡》,當中表達了他對上帝存在的一個很突破性的看法:

「上帝正在教導我們,我們生活為人,可以不要靠賴他.與我們同在的上帝,就是離棄我們的上帝.……這就是基督教與其他宗教主要不同之處.人的宗教意識使他在痛苦時去仰賴世上有力的上帝,以上帝為救星.聖經卻指引他去尋找一個無能無力和受痛苦的上帝.唯有受痛苦的上帝才能幫助人.我們所描述的已成長了的世界,是在於放棄了對上帝的錯誤觀念,且準備為聖經所指示的上帝而戰鬥；上帝以祂的軟弱去征服世界的強權.」《獄中書簡》第141頁.

其實,每次當我問「上帝,你還在嗎?」,只是出於對人生的感慨；我絕對知道我問錯了問題.當潘霍華提醒我們,世界已經成長,上帝不再是以強而有力的第三者出現,上帝存在於人的內心,而當我們成長後,我們絕對應該要向自己所作的負責.世界與上帝,就好像子女與父母的關係一樣；子女年幼時,事事靠父母照顧；但當子女成長後,他應該有獨立自主的能力,並且要向自己的所作所為負責任.

是的,世界已經成長；今日我們鮮見上帝以超自然的力量來干預世界的運作.當我們看見世上那麼多天災人禍,要問的,不是上帝還在嗎；而是要問人類,問我們自己──作為世界的管家,為何會把這個世界弄得那麼糟?

我最近看了一條短片,在此與大家分享:《一名法學院學生的演講》

我每次看完北大法學生這個演講片段,我都有股想哭的衝動；上月學院職員退修會上,我與同事分享這條短片後,也無法控制自己的眼淚.這條片給我最深刻的,是這四個字:「不要變壞」!今日我們的祖國神舟十號己經升空,國民生產總值成為全球最大經濟體之一,但為何在國際上卻仍然無法贏得到別人的尊敬,甚至經常成為笑柄?中國食品沒人敢買,中國製造等於血汗工場!這正正是因為很多人變壞了,為了追求經濟利益不擇手段,罔顧人命.

創世記第一章,記述了上帝創造天地的故事,世界本來就是美好的──「上帝看著是好的」,後來祂把管理權交到按祂形象所造的人類手上.今日,我們把世界管得如何?不要變壞──說得真好!今天,並非上帝已經不存在,要是我們變壞了,埋沒了自己的良心,再沒有位置留給上帝.上帝依然存在,只是並非以我們期望的方式來存在.每見人間罪惡時,所問問題的對象,不應是上帝,而是「人,你變壞了麼?」

一位我非常敬重的老師,亦曾任崇基學院院長的沈宣仁教授(1931-2004),在他離世前一年,我往美國洛衫磯探望他；當我們的話題轉到他的病情以及人間苦難時,他提到近年來的一些體會:

「我發覺自己已經不再為那些人間罪惡,天災人禍而流淚；相反地,我只會為一些感人的事蹟,人性的真善美而流淚.苦難在人世間來說,本就是常態,是與生俱來的；唯有人性的真善美,才是不尋常.很多人遭遇苦難時都會怨天尤人,這是不應該的,因為這個世界本來就是這樣.」《沈宣仁教授追思禮拜文集第73頁》

變壞了的人,變了質的世界,原來已變了常態.要活出上帝形象,有如要違反地心吸力,反其道而行,是需要不斷努力的；我們有沒有勇氣,如那位女學生所說,對那些已變壞的人說:我跟你不一樣!因我們是有上帝形象的!

上帝還在,祂正活在我們的心中!願我們一起努力,活出上帝的真善美!

~~~~~~~~~~~~~~~~~~~~~~~~~~~~~~~~~~~~~~~~~~~~~~~~~~~~~~

附件:一名法學院學生的演講

我是一名法學院的學生.我的每一門課的教授,都曾經在他的課堂上講過這麼一句話.他們常常說:「法律是這麼規定的,但是現實生活中(卻不是).」

現實生活,是一種很神奇的生活.在現實生活中,那些尊重規則的老實人,往往一輩子都默默無聞；反倒是那些弄虛作假的人,到最後會名利雙收.於是乎像我這樣的年輕人,就經常有那些看著很有經驗的前輩,過來拍拍你(我)的肩膀跟你(我)說「年輕人你還不懂」.

我想問的是:「我們年輕人,你(我們)能為這個世界做甚麼?」

總有一天,銀行行長會是九零後(九十後)；企業家會是九零後(九十後)；甚至國家主席都會是九零後(九十後).當全社會都被九零後(九十後)佔領的時候,我想問你們九零後們(九十後們):「大家想把這個社會變成怎樣?」

我知道,不是每一個人,他都能夠成為那種站在風口浪尖上去把握國家命運的人物.你,我,都是再普通不過的升斗小民.是這個龐大的社會機器上,一顆小小的螺絲釘.

讀書的時候,每天都被父母耳提面命,說你幹啥都不要給我耽誤學習；畢業的時候,到處投簡歷(履歷),凄凄惶惶的等一家企業收留自己；逢年過節被逼婚；結婚買了房子,要花自己年輕的時候的,最好的二十年來償還貸款.

(這個社會)讓每一個年輕人都忙著生活,而沒有夢想,沒有時間關心政治,沒有時間關心環境,沒有時間關心國家的命運…還有哪甚麼精力去為這社會做甚麼?

但是後來我發現,還是有一件事情,你跟我都可以做到.這件事情就是:我們這一代人,在我們老去的路上,一定一定不要變壞 - 不要變成你年輕的時候,最痛恨,最厭惡的那種成年人.

如果將來,你去路邊擺攤,你就不要賣地溝油小吃,不要缺斤短兩；你將來開了工廠當了老闆,你不要偷工減料生產一些次品.

每一個普通人,他在自己普通的崗位上做一個好人,是有非常非常嚴重(重要)的意義的 - 因為我們每一個人,生下來都注定會改變世界.

我是一個學法律的,如果我將來是一個公正嚴明的法官,那麼這個社會就因為多了一個好法官,而變好了一點點.我希望大家都記住,即使給了你十萬個理由讓你去作惡,你都要保持自己的操守跟底線 - 僅僅就因為一個理由,這個理由就是,你不是一個(一隻)禽獸,你是一個人.

我更希望我們所有的九零後們(九十後們),你們(我們)都能成為那種難能可貴的年輕人,一輩子都嫉惡如仇,絕不隨波逐流；你絕不趨炎附勢；你絕不摧眉折腰；你絕不放棄自己的原則；你絕不絕不失望於人性.

所以,我親愛的九零後們(九十後們),如果將來再是有那些人跟你說,年輕人你不要看不慣,你要適應這個社會,這時候你就應該像一個真正的勇士一樣直面他,你告訴他:「我跟你不一樣,我不是來適應社會的 - 我是來改變社會的!」

謝謝.



\newpage

\section{教會沒有畢業典禮的語言學校}
\label{sec:292}
\textbf{松谷曄介牧師 - 教會:沒有畢業典禮的語言學校!?}
\newline
\newline
\hyperref[sec:291]{< < < PREV SERMON < < <}
~
\hyperref[sec:toc]{[返目錄]}
~
\hyperref[sec:293]{> > > NEXT SERMON > > >}


講題:“教會:沒有畢業典禮的語言學校!?”　    The Church : A language school without graduation ceremony!?

經文:申命記6章4-9節/ 詩篇131章1-3節/ 羅馬福音10章14-17節/ 馬太福音19章13-15節

講員:松谷曄介牧師

場合:香港中文大學崇基學院禮拜堂

日期:2015年08月02日

申命記6:4-9/

4「以色列啊,你要聽!耶和華－我們神是獨一的主.5你要盡心,盡性,盡力愛耶和華－你的神.6我今日所吩咐你的話都要記在心上,7也要殷勤教訓你的兒女.無論你坐在家裡,行在路上,躺下,起來,都要談論.8也要繫在手上為記號,戴在額上為經文；9又要寫在你房屋的門框上,並你的城門上.」

詩篇131:1-3/

1耶和華啊,我的心不狂傲,我的眼不高大；

重大和測不透的事,我也不敢行.

2我的心平穩安靜,好像斷過奶的孩子在他母親的懷中；

我的心在我裡面真像斷過奶的孩子.

3以色列啊,你當仰望耶和華,從今時直到永遠!

羅馬10:14-17/

14然而,人未曾信他,怎能求他呢?未曾聽見他,怎能信他呢?沒有傳道的,怎能聽見呢?15若沒有奉差遣,怎能傳道呢?如經上所記:「報福音,傳喜信的人,他們的腳蹤何等佳美! 」16只是人沒有都聽從福音,因為以賽亞說:「主啊,我們所傳的有誰信呢?」17可見,信道是從聽道來的,聽道是從基督的話來的.

馬太19:13-15

13那時,有人帶著小孩子來見耶穌,要耶穌給他們按手禱告,門徒就責備那些人.14耶穌說:「讓小孩子到我這裡來,不要禁止他們；因為在天國的,正是這樣的人.」15耶穌給他們按手,就離開那地方去了.

1

兩個星期前,我在另外一間教會裡第一次用廣東話講道,所以這一次是我人生當中第二次用廣東話講道.去年我在這裡第一次用普通話講道的時候,崇拜後有一位朋友鼓勵我說；“下一次用廣東話講道吧!”當時,我還想不到自己真的有一天能用廣東話講道.

2013年10月我來到香港的時候,基本上不會講廣東話,所以我跟神學院的老師和同學,教會的教友講話的時候一般用英文或者普通話來溝通.之後,我在一間語言學校裡,禮拜一到五,早上九點到十二點,學了半年的廣東話.老師教書的時候基本上只用廣東話.剛開始的時候,聽得不太明白,但是天天在廣東話的環境裡,慢慢就聽得明白了.

學了半年廣東話後,有一天,我在大學裡面遇到了一位本地的神學生.之前,我主要是用普通話和英文來跟他溝通.那個時候,我開始嘗試用廣東話聊天.雖然現在不記得講了什麼,但是我還記得,他很驚喜我可以用廣東話與他交流.我也很高興能夠用他的語言來跟他溝通.我發覺,語言的溝通就是心靈的溝通.

2

大家知不知道,我們基督徒有一個共同的語言?當然不是廣東話,也不是普通話,也不是日本話,那是“上帝話”!無論是日本人還是香港人,中國人,韓國人,台灣人,歐美國家的人,我們都可以一起崇拜,讚美上帝,同心禱告,並且我們的心也可以互相溝通,這是因為我們都有一個共同的語言“上帝話”.基督徒之間之所以有喜樂和安慰的交流,也都是因為我們能透過“上帝話”來溝通.

那麼我們怎樣學習上帝話呢?我認為,學習上帝話有點像學習外語.我們第一次聽聖經故事或者剛開始聽講道的時候,或許不太明白到底是什麼意思.雖然知道每個詞,每句話的意思,但是不一定能夠接受上帝的存在,不一定能夠明白上帝對我們的愛.但是,繼續參加教會的聚會,繼續聽聖經的話語,上帝的話語,我們慢慢就會明白上帝的話語了.剛開始的時候不太明白的聖經話語,漸漸地就明白了.當我們能聽明白上帝話語的時候,我們會感覺非常高興,是因為我們的心開始能夠跟上帝溝通了.

3

學語言的時候,聽說讀寫都很重要,但是語言原來是“聲音”,所以“聽”和“講”更重要.那麼,聽和講,哪一個更重要呢?我認為 ,我們好好聽才能夠好好講,所以,學語言的時候最重視的應該是“聽”.

今天的舊約經文有一節比較有名:“以色列啊,你要聽!耶和華－我們神是獨一的主”(申命記6:4).現在我們有了這種一本書的聖經,但是古代的上帝的子民當然沒有聖經.聖經這本書是後來為了記錄上帝子民的歷史慢慢寫出來的.摩西帶領以色列人出埃及的時候,最重要的事還是“聽”上帝的話語.先知摩西把他親自聽到的上帝話語轉告給以色列人,要他們好好聽!當然摩西本身不是上帝,所以他們不應該崇拜他.他們在摩西的話語裡面聽到了上帝的話語.他們的信心透過“聽”上帝的話語來慢慢成長.

我們觀察小朋友成長過程的時候,可以發現他們能很快吸收很多語言.剛出生的寶貝無論懂不懂語言,爸爸媽媽還是抱著寶貝不停地跟他講話.這樣,小朋友一直聽家長和周圍的其他大人的講話,漸漸地自然地學會了語言.

詩篇131篇這樣描述著,像寶貝在母親的懷抱裡面一樣,我們人在上帝的懷抱裡面,依靠祂.上帝用祂的手抱著我們,無論我們懂不懂他的話語,祂不會放棄我們,一直在跟我們講話.我們在上帝的懷抱裡長大,慢慢地能聽明白上帝的話語了.

這兩年,我觀察了我的兩個女兒在香港學習外語的過程.大女兒雖然在日本人學校唸書,但是她經常在我的旁邊聽我跟其他人講的英文和中文,逐漸地能聽明白一些或者能推測大概的內容了.小女兒在香港本地幼兒園上學.她剛開始上學的時候,廣東話一點都不會講.我和太太,幼兒園的老師都有一點擔心,她能不能習慣本地幼兒園.最初的兩三個月,她在語言溝通方面有一些困難.有一次,她好像做了不應該做的事情,老師勸告她那樣做是不行的,但是她不懂老師講什麼,就哭起來了.但是,過了半年,她已經聽得明白老師和同學講的廣東話了,而且她自己也慢慢開始主動地講話了!她從來沒有學過什麼廣東話課本等等,她不過是在百分之百的廣東話環境裡,一直被周圍老師和朋友的言語所熏陶,後來很自然地學會了廣東話.

4

語言環境是很重要的.我們需要一直留在主耶穌基督那裡,繼續聽祂的話語.主耶穌說:“讓小孩子到我這裡來,不要禁止他們；因為在天國的,正是這樣的人”.耶穌吩咐大人,他們應該像小孩子一樣.我們要成為小孩子,這裡應該有几個不同的含義.但是,今天我想解釋,像小孩子一樣也意味著“像小孩子一樣地學習語言”.

小孩子出生後,成長的過程當中,一直聽爸爸媽媽講的話來學習語言.同樣,我們大人都像小孩子一樣要聽上帝的話語.我們大人往往認為,我們已經學了很多,所以以後可以不用那麼努力地繼續學習了.

但是,羅馬書10:17說:“信道是從聽道來的,聽道是從基督的話來的.” 這個“道”可以說是“上帝的話(the Word of God)”.相信上帝的話是從聽上帝的話開始的,這也意味著,為了繼續相信上帝的話語,我們都要繼續聽上帝的話語,即基督的話語.無論是牧師還是大學老師,主日學老師,無論是接受洗禮一年的人,還是五十年的人,我們都一樣要繼續努力學習上帝的話語.

5

我們學習“上帝話”的時候,環境是很重要的.一般來說,我們一個人的話,不能學好語言,一定要選擇語言課程或者要去語言學校.因為學語言的時候一定要聽老師講的正確發音和語法.而且,如果有一些同學可以一起聊聊天,練習講話,互相鼓勵,一起努力學習的話,進步會更快.

我們學習“上帝話”的時候也需要合適的環境.為了能夠明白上帝的話語,我們需要去可以學習“上帝話”的學校.我們需要加入可以繼續聽“上帝話”的共同體.那就是教會.在教會這個共同體裡面,我們才能夠聽“上帝語”.這個“上帝話”不僅僅是講道,我們在詩歌,禱告,讀經裡面也可以聽“上帝話”.而且,崇拜後跟弟兄姊妹們一起分享信仰故事的時候,跟弟兄姊妹聊聊天的時候,也可以聽他們講“上帝話”.我們在崇拜時一起聽“上帝話”,一起用“上帝話”來唱詩,禱告,用“上帝話”一起分享,彼此鼓勵,互相安慰.

我們在教會這個環境當中不僅可以 學習“上帝語”,練習“上帝語”,而且也可以教別人“上帝語”.在這個意義上,我們可以說,教會是學習“上帝語”的一種語言學校.

6

學習語言的時候,離不開環境.如果我們離開一個很好的環境,語言能力會馬上下降.如果離開教會這個環境,我們的“上帝話”的能力將會越來越不好,[聽得不好了,講得不好了 聽力和口語都會下降].我以前服侍的日本教會有很多不來崇拜的信徒.一些人由於家庭的事情,一些人由於工作的事情,慢慢不來崇拜,一年只有一兩次來而已.明顯的是他們講的“上帝話”很差了,是因為他們好久沒有聽上帝的話語,好久沒有與教會的弟兄姊妹用“上帝話”溝通了.他們當中也有曾經擔任長老和主日學老師,但是他們的“上帝話”的能力的確下降了.所以,我牧會的時候,一直在想怎樣可以讓他們重新學好“上帝話”,怎樣可以讓他們重新想起來他們原來講的“上帝話”.

如果我的兩個女兒離開香港這個語言環境的話,她們就很快就會忘記在這裡學會的語言.回日本以後,我想盡量給她們一些可以接觸中文和英文的環境,但是周圍的環境基本上還是日語,所以維持外語能力不容易.盡管如此,我希望她們長大後再次回來香港,學中文和英文.我相信,她們到那時候會很快重新學好語言的.當她們回到這個環境後,她們會想起來很小的時候自己能講的語言,自己能聽的語言.她們過去兩年在香港學的語言,雖然回國後會忘記,但是我相信她們將來還可以想起來.

“上帝話”也是一樣.忘記“上帝話”的人,如果重新回到可以聽和可以講“上帝話”的一個合適的環境,那些人也可以慢慢想起來“上帝話”,慢慢聽得懂,慢慢能講出來.這個合適的環境就是教會.

7

教會是我們學習“上帝話“的語言學校.老師是誰?當然是耶穌基督.教會的牧師並不是老師,只是班長而已,牧師也需要跟其他的弟兄姊妹們一起從耶穌基督的話語裡來學習“上帝話”.

學習“上帝話”的學校和一般的學校有一個重要的不同點.教會,這個學校沒有畢業典禮.一般學校,學生的目的是好好學習後順利畢業,但是學習“上帝話”的學校是一輩子持續學習的地方.我們人是學不完“上帝話”的,因為,上帝的話語是無限的,永遠的.

教會,這個學校除了沒有畢業典禮之外,還有一些特點,就是沒有不及格和退學.離開教會一段時間的人,自己以為“我已經不及格,已經退學了(退教會),沒有資格再回教會了,但是,主耶穌基督所有的名單裡面還記錄了他們的名字.基督一直等著他們再回來“復學”,曾經一起學習“上帝話”的同學,一起崇拜的弟兄姊妹都一直等著那些人的“復學”.

8

“上帝話”是上帝國度的語言.我們無論是哪裡人,我們都是上帝國度的子民,我們無論講日文還是中文,我們都是在講上帝國度的“國語”.

我和家人後天要回日本去了.我自己一月份再回來,但是明年春天將要回日本.我們一家人在香港的這兩年學了廣東話,普通話,英文,當然我們都講得不太地道,但是還可以跟崇基禮拜堂的大家交流,超越語音,文化,民族,歷史等等,我們都成為了主內的朋友和家人.為什麼可以這樣呢?都是因為在香港的大家和從日本來的我們原來有一個共同語言,即“上帝話”.

無論在日本還是香港,願我們今後也在不同地方的教會裡繼續聽上帝的話語,學習“上帝話”吧!我們下次見面的時候,我和我太太,兩個女兒的廣東話,普通話,英文或許都不好了,但是我家人以後也會繼續留在教會裡,不會忘記“上帝話”.請大家也繼續留在教會裡,繼續聽和學習“上帝話”.如果我家人和大家都沒有忘記“上帝話”,我們將來再次見面的時候,還可以用“上帝話”來溝通,一起崇拜,一起讚美神,一起禱告,一起交流.

「以色列啊,你要聽!耶和華－我們神是獨一的主.你要盡心,盡性,盡力愛耶和華－你的神」

崇基禮拜堂,上帝的子民,弟兄姊妹們啊!你們要盡心,盡性,盡力聽上帝的話語!



\newpage

\section{語出驚人}
\label{sec:293}
\textbf{許開明牧師 - 語出驚人}
\newline
\newline
\hyperref[sec:292]{< < < PREV SERMON < < <}
~
\hyperref[sec:toc]{[返目錄]}
~
\hyperref[sec:294]{> > > NEXT SERMON > > >}


約翰福音 6:35,41-51
\newline
\begin{longtable}{cl}
\hline
\hline
章節 & 經文 (和合本修訂版)\\
\hline
6:35 & \begin{tabularx}{0.7\textwidth}{X} 耶穌對他們說：「我就是生命的糧。到我這裡來的，絕不飢餓；信我的，永不乾渴。 \end{tabularx} \\ \\ \relax
6:36 & \begin{tabularx}{0.7\textwidth}{X} 可是，我告訴過你們，你們已經看見我，還是不信。 \end{tabularx} \\ \\ \relax
6:37 & \begin{tabularx}{0.7\textwidth}{X} 凡父所賜給我的人，必到我這裡來；到我這裡來的，我總不丟棄他。 \end{tabularx} \\ \\ \relax
6:38 & \begin{tabularx}{0.7\textwidth}{X} 因為我從天上降下來，不是要按自己的意願行，而是要遵行差我來那位的旨意。 \end{tabularx} \\ \\ \relax
6:39 & \begin{tabularx}{0.7\textwidth}{X} 差我來那位的旨意就是：他所賜給我的，要我一個也不失落，並且在末日使他復活。 \end{tabularx} \\ \\ \relax
6:40 & \begin{tabularx}{0.7\textwidth}{X} 因為我父的旨意是要使每一個見了子而信的人得永生，並且在末日我要使他復活。」 \end{tabularx} \\ \\ \relax
6:41 & \begin{tabularx}{0.7\textwidth}{X} 猶太人因為耶穌說「我是從天上降下來的糧」，就私下議論他， \end{tabularx} \\ \\ \relax
6:42 & \begin{tabularx}{0.7\textwidth}{X} 說：「這不是約瑟的兒子耶穌嗎？我們豈不認得他的父母嗎？現在他怎麼說『我是從天上降下來的』呢？」 \end{tabularx} \\ \\ \relax
6:43 & \begin{tabularx}{0.7\textwidth}{X} 耶穌回答，對他們說：「你們不要彼此私下議論。 \end{tabularx} \\ \\ \relax
6:44 & \begin{tabularx}{0.7\textwidth}{X} 若不是差我來的父吸引人，就沒有人能到我這裡來；到我這裡來的，在末日我要使他復活。 \end{tabularx} \\ \\ \relax
6:45 & \begin{tabularx}{0.7\textwidth}{X} 在先知書上寫著：『他們都要蒙神教導。』凡聽了父的教導而學習的，都到我這裡來。 \end{tabularx} \\ \\ \relax
6:46 & \begin{tabularx}{0.7\textwidth}{X} 這不是說有人看見過父，惟獨從神來的，他才看見過父。 \end{tabularx} \\ \\ \relax
6:47 & \begin{tabularx}{0.7\textwidth}{X} 我實實在在地告訴你們，信的人有永生。 \end{tabularx} \\ \\ \relax
6:48 & \begin{tabularx}{0.7\textwidth}{X} 我就是生命的糧。 \end{tabularx} \\ \\ \relax
6:49 & \begin{tabularx}{0.7\textwidth}{X} 你們的祖宗在曠野吃過嗎哪，還是死了。 \end{tabularx} \\ \\ \relax
6:50 & \begin{tabularx}{0.7\textwidth}{X} 這是從天上降下來的糧，使人吃了就不死。 \end{tabularx} \\ \\ \relax
6:51 & \begin{tabularx}{0.7\textwidth}{X} 我就是從天上降下來生命的糧；人若吃這糧，必永遠活著。我為世人的生命所賜下的糧就是我的肉。」 \end{tabularx} \\ \\
[1ex]
\hline
\hline
\end{longtable}
講題:語出驚人　The Speech Made Scared

講員:許開明牧師

經文:約翰福音6:35,41-51

場合:香港中文大學崇基學院禮拜堂

日期:2015年8月9日

「語出驚人,口出狂言」

本節講章引用經文出自約翰福音第六章第三十五節:

「35耶穌說:『我就是生命的糧.到我這裡來的,必定不餓；信我的,永遠不渴.』

及約翰福音第六章第四十一至五十一節:

「41猶太人因為耶穌說『我是從天上降下來的糧』,就私下議論他,42說:『這不是約瑟的兒子耶穌嗎?他的父母我們豈不認得嗎?他如今怎麼說「我是從天上降下來的」呢?』43耶穌回答說:『你們不要大家議論.44若不是差我來的父吸引人,就沒有能到我這裡來的；到我這裡來的,在末日我要叫他復活.45在先知書上寫著說:「他們都要蒙上帝的教訓.」凡聽見父之教訓又學習的,就到我這裡來.46這不是說有人看見過父,惟獨從上帝來的,他看見過父.47我實實在在地告訴你們,信的人有永生.48我就是生命的糧.49你們的祖宗在曠野吃過嗎哪,還是死了.50這是從天上降下來的糧,叫人吃了就不死.51我是從天上降下來生命的糧；人若吃這糧,就必永遠活著.我所要賜的糧就是我的肉,為世人之生命所賜的.』」

引言

試想像一下,有一天你沿著崇基禮拜堂旁的斜路往禮拜堂正門方向走,突然有位兄弟大叫:「我是從天而降的食物,我的身體是可以吃的,吃過我的肉永遠不餓；喝過我的血永遠不會渴.來,快吃我的肉,喝我的血吧!」當時你會有甚麼反應呢?

各位,今天我們讀到的經文,在二千年前由耶穌向群眾演講過.你想,當年的群眾有怎樣的反應?

耶穌為甚麼如此「語出驚人」呢?可能是希望引人注意；可能是想刺激思考；也可能是想令人回味[1]!那麼,為何說耶穌「語出驚人」呢?「驚人」的地方主要在於這幾方面:

約翰福音第六章第四十一節中,耶穌說祂是從天而降的身體,更可將肉,血分給眾人吃,喝.到了第五十二節,猶太人們便問,耶穌怎能把祂的肉分給他們吃.而且耶穌說「我就是生命的糧」所用的「我就是」(Ego eimi)譯為中文就是「我是是」,是神的自稱.

所以難怪有人會認為「耶穌是傻子,智慧不足,IQ 低」,甚至「是瘋子,有精神病,神經痛,『發癲』」,甚或是粗俗卻又帶點雅氣和理性的「『癡線』」.猶太人亦認為耶穌不是神,只是人 - 人怎會把自己的身體給別人吃?這觀點不對嗎?這是個理性而基於事實的觀點,亦是十分自然的.

耶穌之所以「語出驚人」,甚至是「口出狂言」,是不是因為祂頭腦有問題呢?這場爭論又是為何會開始的呢?

耶穌說祂是天上的糧,於是猶太人之間就彼此爭論,但是私下的,不公開的,卻是一面倒的,也不只屬於猶太人的!因為他們的對手是耶穌,這是耶穌與猶太人間的爭論 - 猶太人抱有的共同看法就是「否定耶穌從天而降」,也「否定耶穌是生命的糧」.他們的爭論是沒有出路的,為什麼?這是因為「雞同鴨講」,完全不對「嘴型」:猶太人說的是地上的糧；耶穌說的卻是天上的糧.

約翰福音之有著希臘哲學家的思想背景,而柏拉圖提出的數個重要概念,足以影響當代以及現代世界:

在量度上「以有限想像無限」；在空間上「以世上的想像討論天上的事物」；在感官上「以眼目想像非眼見」；在觀念上「以主觀想像客觀」.

從地上的糧到天上的糧,會有以下的分別:

地上的糧上可見的,天上的糧則是看不見的；地上的糧會腐壞,天上的糧卻是永恆的,吃完的人就得永生；地上的糧,吃完了還有餓,天上的糧則是吃了就不會再餓的.由此我們便可以知道,天上的糧對人生來說,深具意義!

試想想,我們每日都要「搵食」,更是要「搵三餐食」或是「煮三餐食」,吃飽了又要洗碗洗碟,清潔餐桌,更要消化,排洩,還要處理掉那些排洩物!光是花在進食的時間已經差不多三小時,「又煮又食又要清潔,好麻煩呀!」

但如果有食物,吃過一次就以後都不用再吃；吃過一次說以後都不會肚餓的話,不就可以節省許多時間,可以減少進食的麻煩,更可以不再受飢餓的苦嗎?這對熱愛工作的你來說,是多麼美好的事啊!「生命的糧」究竟是怎樣的一回事?

在第三十五節耶穌宣告說:「我就是生命的糧(這讓群眾得希望).到我這裡來的(這讓群眾得出路),必定不餓(這讓群眾得飽足)；信我的,永遠不渴(這讓群眾得解渴).」

這樣的「天糧」,你想吃嗎?這種「天糧簡直是人類的夢想,是人心深處的渴望!一位古代學者說過「人的欲望無盡,只有一位無限於神,才能使人滿足.」觀乎耶穌的「天糧」,即可知耶穌就是這樣能夠滿足人心的神.

猶太人與耶穌的對立,對質

在第四十二至四十三節,猶太人反問「這不是約瑟的兒子耶穌嗎?他的父母我們豈不認得嗎?他如今怎麼說我是從天上降下來的呢?」

猶太人之所以否定耶穌是從天而降,理據在於第四十二節的一句「這不是約瑟的兒子耶穌嗎」- 猶太人認識耶穌的父親.猶太人認為,「耶穌是有父親的,所以他不是從天而降的」；猶太人又認為,「因為耶穌在地上有父親,說明耶穌就是為人所生的,不是從天而降的」.

可是,耶穌提的「我是生命的糧」,並不指祂的肉身,而是生命的要素,是靈糧,是超越物質的,是精神上的,是靈性的.

第四十二節又提到猶太人反問「他(耶穌)的父母我們豈不認得嗎」.這是說明猶太人覺得他們認識耶穌的父母,所以就認為耶穌不是生命的糧.這不是「牛頭不搭馬嘴」嗎?這也讓人想到猶太人瞧不起耶穌.

經常都有人說:「佢?我都識啦,好叻咩?」,「果個牧師,我都識啦,好叻咩?」因為你認識他,所以他就「唔叻」,這完全是不合邏輯的推論!猶太人的一番話,也見不得與「耶穌是從天而降的糧」有何關係.

耶穌說祂是生命的糧,主要基於以下的理據:

「耶穌回答說:「你們不要大家議論.(第四十三節)若不是差我來的父吸引人,就沒有能到我這裡來的；到我這裡來的,在末日我要叫他復活.(第四十四節)在先知書上寫著說:他們都要蒙神的教訓.凡聽見父之教訓又學習的,就到我這裡來.(第四十五節)」

從上面的文字,我們可知(1)是父吸引人到耶穌這裡來,得「生命的糧」,(2)是聽見父的教訓和學習的,才說耶穌這裡來得「生命的糧」,(3)是相信耶穌的人才到祂這裡來得「生命之糧」的.

另外,從約翰福音第六章的整體來說,我們可進一步綜合耶穌是「生命的糧」的理據如下:

首先「生命的糧」就是維持生命的要素,從人的精神,靈性層面來說,那就是上帝的話(因為「人活著不是單靠食物的」).耶穌作為「道」,「言語」,「上帝的話」,就正是人類生命的要素；耶穌也是「神」,是「生命的源頭」,人與「生命的源頭」建立關係,就可得「供應」,「飽足」和「力量」.

至於得耶穌＝「生命之糧」的方法,從約翰福音第六章第三十四至四十節,我們只要見到耶穌,繼而走向耶穌,相信耶穌,跟他握手,建立關係,耶穌就給予我們平安與永生.

人面對信仰(天糧,真理)時,就是受著上帝的吸引.本段經文中有個特別的字「吸引(希臘文為helkuein)」- 人會被上帝吸引,似磁石一樣走向上帝,但同時又會抗拒上帝的吸引,於是人就常在這抗拒與吸引之下,處於兩股拉力當中.這就好像漁夫拉起魚網,漁夫受著「魚穫」吸引,卻又受「魚的重量」所阻礙.

人會受真理,正義吸引,也會在面對強權時出現拉力.六四的時候,我聽過兩句話,印象深刻:第一句是「你有坦克,但我有真理」；第二句則是「我有真理,但你有坦克」.兩句話的字數,用字都一樣,但意思卻大衛不同.面對坦克,追求真理的人更需要力量,需要勇氣!

耶穌所賜的「天糧」,就是要解決人的這個難題:在古時獻祭中有種信仰,當獻祭之後上帝就會進入祭肉中,人吃過就會充滿力量,並著勇氣勝過罪惡,實踐真理.約翰福音的作者透過第五十一節(「『我是從天上降下來生命的糧；人若喫這糧,就必永遠活著；我所要賜的糧,就是我的肉,為世人之生命所賜的.』」),告訴基督徒,耶穌基督在聖餐禮中邀請信徒吃那代表基督身體的無酵餅,喝那代表基督寶血的葡萄酒,當信徒憑信心吃下,喝下,就表示基督的生命已經進入信徒的生命中,而這聖餐就正正代表「耶穌基督生命的糧」.這更是承傳了以往獻祭的傳統,祭司獻祭祝聖之後,上主就從天而降到祭物之中,而牧師就正好在聖餐禮中擔當祭司的角色,祝聖後耶穌基督就從天而降,臨到祭品(聖餐的餅和酒)中,信徒吃下這些耶穌基督進入了的聖餐後,便得著生命的糧,得著力量,得著勇氣,面對罪惡的世界.

一個人有病,除了服藥打針直接對付病菌,更重要的是要強身健體,也就是要透過運動,食療,睡眠等等來增強體力,勝過罪惡.

耶穌跟猶太人對立

在對立之中,基於大家提出的論據,我們可知耶穌的理據明顯比較充分.

在人生中,人也少不免會面對種種對立,遭否定的時刻,我們或會問天父,為何你要容許人與人之間有對立的情況?要使人不和,使人受苦呢?在人的世界,不論是學生跟老師,同事們之間,兄弟姊妹間,上下兩代人之間,也有對立；在聖經裡,也有地糧跟天糧,物質與精神,眼見與非眼見,可量度與不可量度,短暫與永恆之間的對立；在宇宙萬物之間,也充滿著善和惡,正與邪,陰陽兩者的對立.

面對這些世間常見的對立,我們是怎樣應對的?一般來說,人是傾向「勢不兩立」,「消滅對立」,「聯手對立」,還是「消滅對立」?

試想像這樣的一個故事:某君遭上司無理的責罵,幾位同事為此抱不平,想要反駁上司,事後上司卻召那位遭無理取鬧的下屬,教訓他說:「你們好大的膽子,竟敢聯手反我!」火冒三尺的上司繼續罵,「改天給你們統統幹掉!」

某君回家就告訴太太這事,太太說跟某君說:「這樣不好,不要讓老闆以為你有意與他對立,以為你想要聯同其他同事跟他作對.你從今應避免跟那些同事太接近,免得你成為第一個遭殃的人!」

某君就聽太太的話,一改常態,避開同事之餘又對上司誠惶誠恐似地,恭敬得不得了.「看他今天一副沒有骨氣,對上司唯命是從的樣子!小心這傢伙出賣我們!」同事這樣說.於是同事也學他一樣,討好上司,對立結束,辦公室也回復平靜,好像一團和氣的樣子.

故事好像就此結束,但一天某君接到上司傳召約見,某君就心想,一定是老闆看到我近來的好表現,一改常態想要對我另眼相看!

果然像某君所想,上司笑著對他說:「感謝你多年來的服務!這是給你的信!」某君打開信封,卻嚇了一大跳 - 信上的字寫著「解僱,請另謀高就!此信即日生效!」某君只好收拾東西,直到他離開,也沒有人安慰他.

某君才走出大門,上司就笑盈盈地對其他人說:「我早就想幹掉他,看在你們支持他的份上,不貿然行事.最近看到你們對他也不怎樣,我就決定一腳把他踢走!好,今天午飯由我請客!」可是到了午飯的時候,誰都不發一語,有位比較敢言的同事忍不住問:「老闆,你下個想踢走的是哪個?」老闆沒有回答,可是他的笑容卻暗示了他心中有數…

這個故事,值得我們深思一個問題:「對立一定就是不好的嗎?」若問約翰福音的作者,我可以肯定他會說:「對立是好的.」約翰想突出的是光暗之間,善惡之間,上帝與撒旦之間,地糧與天糧之間的對立.

《猶太法典》之中,有一條自問自答的問題:人的眼睛由黑白兩部份組成,但神為何要讓我們只透過黑的部份去看東西?《猶太法典》說這是因為人生必須透過黑暗,才可以看到光明.

希伯來人這個字,本來的意思就是「站在(河的)對岸的人」,他們被稱為智慧的民族,也跟他們的稱謂有關係:站在對岸觀察,就可以看得更廣,更清楚,也就是說人要常常突破環境的框框,群組的框框,思維的框框,學問的框框,工作的框框等等.除此之外,他們是站在對岸,跟別人對立的,他們不畏懼,該反對時就反對,但也會原諒別人不贊成自己.假若世界是一致的,根本不可能有進步；假若明知是罪惡也不反對,罪惡就會得勝,會橫行霸道；假若罪惡不被視為「不對」,世界更會一片混亂!《猶太法典》又說:「假如世界像一艘船,所有人都向同一個方向走,這個世界就必將傾覆.」

對立,正如前面所述,是無處不在,自然存在的.在家庭中會有父母,夫婦,父母子女之間的對立；在社會中會有階級間的,貧富之間的對立；在政治中會有黨派,政見,執政與在野黨之間的對立……面對對立,人只想消滅對立,但其實即使沒有對立,它也會逐漸出現.

我們就是要肯定對立,容許對立!尤其面對不公義,罪惡之時,更要堅持對立.而且對立的存在,也不必然是壞事,正如一個商場裡也會有一般的店舖(士多,超市),也會有專門的店舖(眼鏡店,書店)等等,各自各精彩,顧客也各取所需.若論我們應完全消除對立,或是接受,容忍對立,與它共存?有些對立如貧富懸殊,我們應當消滅；有些對立如人體內的細菌,我們束手無策；有些對立如政治世界的左,右派及唐代經常進諫的魏徵,有助進步.在權力存在的同時,我們需要可以制衡權力的「反派」,否則權力只會不斷膨脹,甚至腐化!

信仰的接納與懷疑

耶穌宣告:「我是生命的糧；到我這裡來的(出路),必定不餓(真理)；信我的,永遠(永生)不渴.」這是約翰福音第六章三十五節的經文,其中耶穌行了個大神蹟,就是用五餅二魚餵飽五千人.耶穌就是有著使人得飽足的能力.於是許多群眾都跟隨耶穌,想要強迫耶穌作王,耶穌卻獨自退到山上去.到了晚上門徒下海,在小船上卻遇大風浪,害怕得很.害怕之中,卻見耶穌在海上行走,更覺害怕,「人怎能在海上行走?」,「是幽靈嗎?是鬼?是神?」

然後耶穌到了另一個地方,群眾又跟著他,耶穌就教訓他們說(第二十六節):「我實實在在的告訴你們,你們找我,並不是因見了神蹟,乃是因吃餅得飽.(第二十七節)不要為那必壞的食物勞力,要為那存到永生的食物勞力,就是人子要賜給你們的.」

也就是說耶穌要群眾從追求「地糧」,改為追求「天糧」.當時群眾的反應如何呢?對於耶穌的宣告,從古到今都有兩種不同的反應,要不接納,要不懷疑!我們在面對真理的時候,也總是有這兩種不同的反應!一般來說我們都認同「接納真理」的人,不認同「懷疑真理」的人,但細心思考一下與眾不同的約翰福音,我們必會有新的發現:

有些群眾接納「耶穌是生命糧」,正如耶穌在第三十二,三十三節,耶穌說「我實實在在的告訴你們,那從天上來的糧不是摩西賜給你們的,乃是我父將天上來的真糧賜給你們.因為神的糧就是那從天上降下來,賜生命給世界的」之後,第三十四節中群眾就說:「主啊,常將這糧賜給我!」

有些猶太人懷疑耶穌,正如第四十一,四十二節中提到「猶太人因為耶穌說我是從天上降下來的糧,就私下議論他,說:『這不是約瑟的兒子耶穌嗎?他的父母我們豈不認得嗎?他如今怎麼說我是從天上降下來的呢?』」這正是懷疑論(希臘文為σκέπτομαι,skeptomai,意指思考,探索,考慮或研究)涉及的普遍地向知識,事件,意見或信仰持有疑問的態度,又或是懷疑一些理所當然的主張的表現.懷疑論就是一種要求所有論點或資訊都要充分地被理據所支持的整全方式.傳統哲學懷疑論衍生自「skeptomai」,指的是向「任何事物都不加以斷言」的概念,按照邁克爾‧舍默的說法,懷疑論主要是一個發現真理的過程,多於純粹地,普遍地否定一必.基於這個原因,宗教懷疑論者其實亦或會相信耶穌的存在,但同時質疑他是否彌賽亞或會否創造奇蹟(可以參照Historicity of Jesus).通過「系統性的懷疑」及「不斷的測試」以獲得知識,正好可以解釋猶太人從否定中進而繼續否定,甚至是絕對的否定,盲目的否定.

在今時今日,有些人會否定上帝的存在,這可能是因為他們假設「上帝存在的話,上帝就一定要能被我們看見」,或是「如果上帝存在的話,苦難就不應存在」.於是前者就看不見上帝,他們否定上帝,並且絕對地否定；後者則仍然看見苦難的存在,否定上帝,並絕對地否定上帝.這正是一種從懷疑變成否定的過程,造成封閉,固執,非理性,也就猶太人的信仰態度.

當然,話分兩頭說,懷疑信仰的人也可以變成確信者的.正如多馬本來懷疑耶穌的復活及其真實性,結果耶穌就向他顯現,他就從懷疑變成確信；保羅亦同是本來懷疑,否定耶穌的,後來在大馬色路上遇見復活主,又是由懷疑變成確信.今天的信徒,有的由懷疑,到相信,到確信；也有的由相信,到懷疑,到確信.這兩種「信」,相比「一講就信」,更好!皆因經過懷疑,追尋,信仰就會有更鞏固的根基.

結語

在今次的講章中,首先耶穌宣告「我是生命的糧」,語出驚人,甚至是「口出狂言」的程度.若問耶穌頭腦有問題嗎,我們可以從耶穌簡短的回應中,找到充分的理據,理論,證明耶穌豈會是傻子,瘋子,祂簡直是智慧過人.若問耶穌說的「生命糧」為何物,有何意義,我們從經文可知「生命糧」就是信耶穌的人所吃的聖餐,是使人得力量勇敢面對罪惡的世界的恩物.若問耶穌的一生之中,上帝容許「耶穌面對猶太人與其的對立,對質」的場面出現,箇中意義何在,我們可以回答說「對立」可以使人在辯護的過程中申明觀點,並思考他人的觀點,這樣世界,社會的權力就可得到平衡,思想不會一面倒,權力亦不會絕對化,甚至是腐化.那麼,人對信仰有接納的同時,上主又會否容許我們對信仰有懷疑呢?答案是可以的,因為懷疑不是罪,甚至有好處.對信仰有所懷疑,便可以透過懷疑來得到更多,更新,更深入的認識；透過懷疑,更可使我們確信.我們作為受過高等教育的基督徒,更不應停止懷疑及探索.

最後給大家來個小習作,大家能否懷疑一下,自己是存在或是不存在?你在做夢,還是仍然處於清醒狀態呢?剛才聽道睡著的人,請醒來吧.如果你能夠清醒的話,就證明你是存在的.其實剛剛你也存在過,只是「人在心不在」而已!

[1] 筆者亦曾以「唧唧耶穌」,「上帝頂呱呱」,「牧師呱呱叫」和「誰是謀害上帝的兇手?」等為題,寫過數篇文章.



\newpage

\section{吃肉喝血的群體}
\label{sec:294}
\textbf{關瑞文教授 - 吃肉喝血的群體!}
\newline
\newline
\hyperref[sec:293]{< < < PREV SERMON < < <}
~
\hyperref[sec:toc]{[返目錄]}
~
\hyperref[sec:295]{> > > NEXT SERMON > > >}


約翰福音 6:51-58
\newline
\begin{longtable}{cl}
\hline
\hline
章節 & 經文 (和合本修訂版)\\
\hline
6:51 & \begin{tabularx}{0.7\textwidth}{X} 我就是從天上降下來生命的糧；人若吃這糧，必永遠活著。我為世人的生命所賜下的糧就是我的肉。」 \end{tabularx} \\ \\ \relax
6:52 & \begin{tabularx}{0.7\textwidth}{X} 因此，猶太人彼此爭論說：「這個人怎能把他的肉給我們吃呢？」 \end{tabularx} \\ \\ \relax
6:53 & \begin{tabularx}{0.7\textwidth}{X} 耶穌對他們說：「我實實在在地告訴你們，你們若不吃人子的肉，不喝人子的血，在你們裡面就沒有生命。 \end{tabularx} \\ \\ \relax
6:54 & \begin{tabularx}{0.7\textwidth}{X} 吃我肉、喝我血的人就有永生，並且在末日我要使他復活。 \end{tabularx} \\ \\ \relax
6:55 & \begin{tabularx}{0.7\textwidth}{X} 我的肉是真正可吃的；我的血是真正可喝的。 \end{tabularx} \\ \\ \relax
6:56 & \begin{tabularx}{0.7\textwidth}{X} 吃我肉、喝我血的人常在我裡面，我也常在他裡面。 \end{tabularx} \\ \\ \relax
6:57 & \begin{tabularx}{0.7\textwidth}{X} 永生的父怎樣差我來，我又怎樣因父活著，照樣，吃我肉的人也要因我活著。 \end{tabularx} \\ \\ \relax
6:58 & \begin{tabularx}{0.7\textwidth}{X} 這是從天上降下來的糧，不像你們的祖宗吃過嗎哪還是死了；吃這糧的人將永遠活著。」 \end{tabularx} \\ \\
[1ex]
\hline
\hline
\end{longtable}
講題:吃肉喝血的群體　The Cannibals!

講員:關瑞文教授

經文:約翰福音6:51-58

場合:香港中文大學崇基學院禮拜堂

日期:2015年8月16日

很久以前,在一個城鎮裡,有一所教會,教友的人數不多也不少.雖經牧師多番努力,人數仍未見大增.不久,有一所俱樂部,在附近興建了起來.俱樂部的CEO頭腦精明,推出了很多迎合客人的設施,不管你是男女老幼,總能在當中找到你的樂園.因此,開業不久,俱樂部已招徠不少會員.牧師甚為好奇:「俱樂部何來如此的吸引力?」 他決定走訪該CEO.CEO言精意簡,指出秘訣很簡單:「投其所好,讓客人賓至如歸.」牧師聽後,恍然大悟,立即開始更新教會:革新了兒童主日學,以遊戲活動教導聖經,慕求小朋友快快樂樂；崇拜,就以最大眾化的形式進行,慕求人人享受；每周多設幾個祈禱會,為眾人祈福；設立講道評分意見箱,讓各人有冤有路訴等等.沒錯,首半年,人數真的增加了.可是,半年後,教會狀況打回原形.再過半年,人數甚至急速下滑.牧師大惑不解:「我已做到投其所好,讓教友賓至如歸.何解?」他再次請教該CEO.CEO說:「當然吧! 我的是俱樂部,天職是效忠客人.你的是教會,天職是效忠上帝.我的吸引力在於俗世至上,你的吸力在於展示有別於俗世的另類人生.你的聖經在使徒行傳2.38ff不是已經說了嗎?這裡清楚說明教會之所以有吸引力的原由:38彼得說:『你們各人要悔改 . . .』 40彼得 … 勸勉他們說:『你們當救自己脫離這彎曲的世代.』41於是 . . . 那一天,門徒約添了三千人,42都恆心遵守使徒的教訓 … 43眾人都懼怕 . . . 44信的人都在一處,凡物公用,45並且賣了田產,家業,照各人所需用的分給各人.46他們天天同心合意恆切地在殿裡,且在家中擘餅,存著歡喜,誠實的心用飯,47讚美上帝,得眾民的喜愛.主將得救的人天天加給他們.不是嗎?我的吸引力,在於我能融入俗世；而你的吸引力,該是因為你與眾不同,甚至另類得使『眾人懼怕』.」

弟兄姊妹,今天我們讀的約翰福音6:51-58,正是要提醒主的跟隨者,教會是另類群體.

約翰福音6:51-58是一段很受注目的經文,聖經學者花了不少心血研究這組經文.該怎樣讀這段經文呢? 不同人有不同的讀法.一些基督徒,讀這段經文,覺得是耶穌向信他的人宣佈一個福氣,就是他是生命之糧,信他的人有永生.讀畢,連忙感謝上帝,因為知道自己已經獲得了天國入場券.毫無疑問,信耶穌的人真的有永生.然而,這並非這段經文的重點.

曾接受少許神學訓練的人,讀到這段經文時,就常常把閱讀焦點放在聖餐神學上,欲透過這段經文,斷定究竟聖餐的餅和酒,是記念耶穌的受死,還是經祝聖後變質真的成為耶穌基督的聖體寶血,又還是餅和酒與基督真正的身體和血是同時存在於聖餐中.這個焦點,也許不錯,卻似乎沒有照顧到約翰福音經文本身的文本脈絡及當日處境.

今天,讓我們一起重讀這段寶貴經文.

話說,當天,耶穌談論到自己是生命之糧,並說自己的肉和血可吃可喝.不知怎的,這句話卻惹起了猶太人的爭論:「52這個人怎能把他的肉給我們吃呢?」53耶穌聽到了,就說:「我實實在在地告訴你們,你們若不吃人子的肉,不喝人子的血,就沒有生命在你們裡面.54吃我肉,喝我血的人就有永生,在末日我要叫他復活.」說完了這數句,耶穌好像覺得自己說得還未足夠到位,因此,他再以更強的語氣,進一步澄清說,「55我的肉真是可吃的,我的血真是可喝的.」

弟兄姊妹,你猜耶穌這樣澄清,會是讓群眾釋疑,還是讓人厭棄他呢?答案是,他這樣澄清後,反而讓不少人更厭棄他:60他的門徒中有好些人聽見了,就說:「這話甚難,誰能聽呢?」61耶穌心裡知道門徒為這話議論,就對他們說:「這話是叫你們厭棄…」66從此,他門徒中多有退去的,不再和他同行.

讀這段經文,我們要問兩個問題:1.怎麼當耶穌說自己的肉和血可吃可喝,竟然會惹來如此爭論?2.怎麼耶穌的澄清和解釋,看來變得弄巧成拙?連本來跟隨他的門徒中都多有退去?

首先,讓我們探索第一個問題:1.怎麼當耶穌說自己的肉和血可吃可喝,竟然會惹來如此爭論?好像很莫名其妙,是嗎?然而,若我們翻查當時的希羅文獻,我們就很容易明白當中理由.原來,「吃人肉,喝人血」這句話,在當日是很流行用來罵人的.罵什麼?罵你是反對派,叛徒,是我們中間的奸細,是出賣同胞的非我族類. 並且,既然是非我族類,其心必異,必定是危險人物.這個罵人用語,在當日的文獻是不罕見的,在荷馬史詩裡,在修西得底斯的敘事裡,在約瑟夫的歷史裡等等,都可找到例子.

因此,當耶穌說:「54吃我肉,喝我血」時,聽在猶太人的耳裡,無疑是耶穌在說,「我是你們中間的非我族類,我心必異於你心.我帶來的世界,與你們的世界,迥然有異.我為世界定下來的議程,非你們所期望的.」

也許,在聽眾當中的耶穌跟隨者,還以為耶穌這句「吃我肉,喝我血」,不是這個意思.他們可能在想,「吃我肉,喝我血」這句話,一定是象徵用語或者是比喻修辭.它一定有言外之音!總之,耶穌不可能說自己是世界的反叛者,不可能說自己是來搞亂世界的秩序.

正等待耶穌澄清之際,耶穌卻進一步說:53「我實實在在地告訴你們,你們若不吃人子的肉,不喝人子的血,就沒有生命在你們裡面.54吃我肉,喝我血的人就有永生…55我的肉真是可吃的,我的血真是可喝的.56吃我肉,喝我血的人常在我裡面,我也常在他裡面.57…吃我肉的人也要因我活著.」

聽到這個所謂澄清,耶穌的跟隨者都嚇呆了!為什麼?看看耶穌如何在字裡行間改變了他的用字,你便會明白.

起初,耶穌在v51的「吃」字,原文是phago.可是,當他進一步澄清時,在 v54開始,他改用了 trogo 這個字.Phago 與 trogo,中文聖經都翻譯為「吃」, 有分別嗎?有.Phago 是隻一般而簡單的字,平平凡凡是「吃」的意思.Trogo卻是隻很細膩的字,很繪聲繪色地形容「吃」的當下動作,過程型態等等,大概是廣東話的「嗒」,「啅」,「梅」,「含」,「吮」,「咬」之類.

這就是說,耶穌最初是說 「51我是從天上降下來生命的糧；人若吃這糧,就必永遠活著.」可是,當跟隨他的人心懷好意地想這個 「吃」字是比喻用語時,耶穌立刻澄清:「我不是在說笑,也不是在講比喻,我真是說我的肉是可以『嗒』,『啅』,『梅』,『含』,『吮』,『咬』的.」因此,為了強調這個「吃」的真實性,在 v55, 耶穌就索性非常直接的說,「我的肉真是可吃的,我的血真是可喝的.」耶穌要清清楚楚讓人知道,「我是你們中間的 『非我族類』,我心異於你心.我帶來的世界,與你們的世界,迥然有異.我為世界定下來的議程,也不是你們所期望的.」

還有,在言談間,耶穌又說,53「我實實在在地告訴你們,你們若不吃人子的肉,不喝人子的血,就沒有生命在你們裡面.」這無疑在說,「跟隨我的人,是被呼召加入一個新的群體,並且有義務去建造這個群體.而且,對於世界而言,這個新群體是 『非我族類,其心必異』的另類群體!」

這也回答了我們剛才發問的第二個問題:怎麼耶穌的澄清和解釋,看似變得弄巧成拙?連本來跟隨他的門徒,都多有退去?這顯然是因為,耶穌向跟隨他的人所提出來的挑戰,大得他們不願意接受. 他們從來沒有想過,他們被呼召,是要被呼召成為世界的異己,世界的非我族類!

弟兄姊妹,耶穌說:「在俗世社會中,我們是非我族類,我們的心必異於俗世的心.我們的理想世界,與俗世的理想世界,迥然有異.跟隨主的人,為世界定下來的議程,不是俗世所期望的.」

弟兄姊妹,教會是另類群體.教會看人生,看世界,看社會的角度,是迥異於俗世的.教會有自己獨特的語言,獨特的世界觀,獨特的關懷,獨特的議程,獨特的觀點.

近年來,香港人很煩惱.政治和經濟的生態急速地在變化.教會也很煩惱:雨傘,政改,本土運動,中港矛盾,同性婚姻,甚至「等埋首副」事件,教會似乎都不能也大概不應該置身事外.頭痛的是,面對這些爭論,教會該站在甚麼立場?到底教會是追尋民主,還是人民民主專政,還是君權神授?到底教會該是建制還是反建制?今天的經文說, 這些都不是教會的立場.教會是另類群體:教會看人生,看世界,看社會的角度,是迥異於俗世的.教會有自己獨特的語言,獨特的世界觀,獨特的關懷,獨特的議程,獨特的觀點,不需要也不會在俗世意識型態中選擇其一.耶穌基督的跟隨者只懂唯獨忠心於主耶穌基督的教訓,最多只會在適當時式有需要時與俗世意識形態短暫共事 (ad hoc alliance),並不會與俗世意識形態建立永恆契約 (eternal covenant).同樣,究竟基督徒該是左傾還是右傾?是自由主義者,新自由主義者,資本主義的倡導者,還是馬克思主義的擁護者?今天的經文告訴我們, 這些都不是教會的立場.教會是另類群體:教會看人生,看世界,看社會的角度,是迥異於俗世的.教會有自己獨特的語言,獨特的世界觀,獨特的關懷,獨特的議程,獨特的觀點,不需要也不會在俗世意識型態中選擇其一.跟隨耶穌基督的人會說:「我唯獨忠心於主耶穌基督,並不會與任何俗世意識形態建立永恆契約.」

耶穌說,我們是吃肉喝血的群體,我們的身份,價值觀念,世界觀,與世人有異,與世人認為好的,妙的理想世界以至行動方案,議程,有質的分別.

且慢!那基督徒就該不問世事,獨善其身,就該只躲在四面牆壁裡崇拜嗎?絕對不是.聖子道成肉身,沒有想過要獲取世間權力,沒有想過要搞亂世界秩序,他只執意實踐天父的旨意而堅持向世人展示天國的格局―他傳福音給貧窮的人,報告被擄的得釋放,瞎眼的得看見,叫那受壓制的得自由,報告上帝悅納人的禧年(路加福音 4:18).結果,耶穌離不開政治,為了實踐天父所差遣的,他得罪了幾乎所有權力群體:(1)羅馬政府懷疑他；(2) 親建制的撒都該人和文士設法要陷害他；(3) 猶太「本土派」法利賽人視他為異己；(4)「左翼」革命份子奮銳黨人看他不順眼；(5) 甚至群眾到頭來都高呼「釘他十字架!釘他十字架!」.在意識型態上, 耶穌非左非右非中間,並且左右不逢源.最後, 耶穌就是死在政治刑具十字架之上.不過, 他卻改變了世界.初期教會亦然.她因為誓死效忠於主耶穌基督的教訓,結果被當時不少猶太人指控為攪亂天下份子.記載在使徒行傳17.6裡,猶太人對初期教會的指控是這樣的: “These people who have been turning the world upside down…”(NRSV).還有,一本很精彩的書,書名為 Christianity and the State in Asia,　當中研究了一些亞洲教會個案,發現當中不少地區的信徒群體,因為立心實踐主耶穌基督的教訓而不按人的話而生活,結果給好些專權的政府帶來了改革的壓力.他們都不是甚麼左傾或右傾群體,也沒有任何遠大的政治藍圖,他們只是一心一意作主的道的見證人,忠於神言,而聽信人言,結果卻在混濁的社會,世代中,成了暗路明燈.

按耶穌基督生命之糧的教導,教會是另類群體.另類群體,英文是 alternative community.這個英文字 alternative很有意思.Alternative,既有另類之意,亦有選擇 (choice) 之意.沒錯,只有另類,才是真正的選擇.弟兄姊妹,你們記得嗎?麥兜故事有這樣一則.某天,麥兜與媽媽到茶餐廳進餐:

兜:唔該,我要一個常餐呀.

太:常餐﹖常餐有咩食架﹖

伙:同特餐一樣囉.

太:咁特餐偽咩黎架﹖

伙:同快餐咁上嚇囉.

太:咁快餐又係咩黎架﹖

伙:車,快餐咪唶係午餐.

太:午餐食咩架﹖

伙:午餐同晚餐一樣架.

太:咁晚餐又食D 咩呀﹖

伙:晚餐咪即係常餐囉.

太:咁呀,我要兩個常餐呀!

沒有alternative-另類, 就沒有 alternative-選擇.

有另類,有不一樣的,世人才真有選擇.上帝的教會欲要成為世人的真正選擇,就要保持其在主耶穌基督裡的另類特質-在世而不融入於世,發聲卻不說俗世意識形態之言.這樣,教會才能成為俗世的暗路明燈.阿們.



\newpage

\section{信仰路上的四個重要時刻}
\label{sec:295}
\textbf{林豪恩先生 - 信仰路上的四個重要時刻}
\newline
\newline
\hyperref[sec:294]{< < < PREV SERMON < < <}
~
\hyperref[sec:toc]{[返目錄]}
~
\hyperref[sec:296]{> > > NEXT SERMON > > >}


約書亞記 24:1-18
\newline
\begin{longtable}{cl}
\hline
\hline
章節 & 經文 (和合本修訂版)\\
\hline
24:1 & \begin{tabularx}{0.7\textwidth}{X} 約書亞召集以色列的眾支派到示劍，他召了以色列的長老、領袖、審判官和官長來；他們都站在神面前。 \end{tabularx} \\ \\ \relax
24:2 & \begin{tabularx}{0.7\textwidth}{X} 約書亞對眾百姓說：「耶和華－以色列的神如此說：『古時你們的列祖，就是亞伯拉罕和拿鶴的父親他拉，住在大河那邊事奉別神。 \end{tabularx} \\ \\ \relax
24:3 & \begin{tabularx}{0.7\textwidth}{X} 我將你們的祖宗亞伯拉罕從大河那邊帶出來，領他走遍迦南全地，又使他的子孫眾多。我把以撒賜給他， \end{tabularx} \\ \\ \relax
24:4 & \begin{tabularx}{0.7\textwidth}{X} 我又把雅各和以掃賜給以撒，將西珥山賜給以掃為業。但雅各和他的子孫下到埃及去了。 \end{tabularx} \\ \\ \relax
24:5 & \begin{tabularx}{0.7\textwidth}{X} 我差遣摩西和亞倫，照我在埃及中間所做的，降災與埃及，然後把你們領出來。 \end{tabularx} \\ \\ \relax
24:6 & \begin{tabularx}{0.7\textwidth}{X} 我領你們的祖宗出埃及，你們就到了紅海。埃及人帶領戰車騎兵，追趕你們的祖宗到紅海。 \end{tabularx} \\ \\ \relax
24:7 & \begin{tabularx}{0.7\textwidth}{X} 你們的祖宗哀求耶和華，他就用黑暗把你們和埃及人隔開了，又使海水衝向埃及人，淹沒他們。我在埃及所做的，你們都親眼見過。你們在曠野住了很多日子。 \end{tabularx} \\ \\ \relax
24:8 & \begin{tabularx}{0.7\textwidth}{X} 我領你們到約旦河東亞摩利人所住之地。他們與你們爭戰，我把他們交在你們手中，你們就得了他們的地為業。我也在你們面前滅絕他們。 \end{tabularx} \\ \\ \relax
24:9 & \begin{tabularx}{0.7\textwidth}{X} 那時，摩押王西撥的兒子巴勒起來攻擊以色列人，派人去召比珥的兒子巴蘭來詛咒你們。 \end{tabularx} \\ \\ \relax
24:10 & \begin{tabularx}{0.7\textwidth}{X} 但我不願聽巴蘭，所以他反而為你們連連祝福。這樣，我救了你們脫離他的手。 \end{tabularx} \\ \\ \relax
24:11 & \begin{tabularx}{0.7\textwidth}{X} 你們過了約旦河，來到耶利哥。耶利哥人、亞摩利人、比利洗人、迦南人、赫人、革迦撒人、希未人、耶布斯人都與你們爭戰，我卻把他們交在你們手裡。 \end{tabularx} \\ \\ \relax
24:12 & \begin{tabularx}{0.7\textwidth}{X} 我派遣瘟疫在你們前面，將亞摩利人的兩個王從你們面前趕出，並不是用你的刀，也不是用你的弓。 \end{tabularx} \\ \\ \relax
24:13 & \begin{tabularx}{0.7\textwidth}{X} 我賜給你們的地，不是你們開墾的；我賜給你們的城鎮，不是你們建造的。你們卻住在其中，又得吃那不是你們栽植的葡萄園和橄欖園的果子。』 \end{tabularx} \\ \\ \relax
24:14 & \begin{tabularx}{0.7\textwidth}{X} 「現在你們要敬畏耶和華，誠心誠意事奉他，除掉你們列祖在大河那邊和在埃及事奉的神明，事奉耶和華。 \end{tabularx} \\ \\ \relax
24:15 & \begin{tabularx}{0.7\textwidth}{X} 若你們認為事奉耶和華不好，今日就可以選擇所要事奉的：是你們列祖在大河那邊所事奉的神明，或是你們所住這地亞摩利人的神明呢？至於我和我家，我們必定事奉耶和華。」 \end{tabularx} \\ \\ \relax
24:16 & \begin{tabularx}{0.7\textwidth}{X} 百姓回答說：「我們絕不離棄耶和華去事奉別神。 \end{tabularx} \\ \\ \relax
24:17 & \begin{tabularx}{0.7\textwidth}{X} 因為耶和華－我們的神曾領我們和我們的祖宗從埃及地為奴之家出來，在我們眼前行了那些大神蹟，並在我們所行的一切路上，和所經過的各民族中保護了我們。 \end{tabularx} \\ \\ \relax
24:18 & \begin{tabularx}{0.7\textwidth}{X} 耶和華又把各民族和住此地的亞摩利人都從我們面前趕出去。所以，我們也必事奉耶和華，因為他是我們的神。」 \end{tabularx} \\ \\
[1ex]
\hline
\hline
\end{longtable}
以弗所書 6:10-20
\newline
\begin{longtable}{cl}
\hline
\hline
章節 & 經文 (和合本修訂版)\\
\hline
6:10 & \begin{tabularx}{0.7\textwidth}{X} 最後，你們要靠著主，依賴他的大能大力作剛強的人。 \end{tabularx} \\ \\ \relax
6:11 & \begin{tabularx}{0.7\textwidth}{X} 要穿戴神所賜的全副軍裝，好抵擋魔鬼的詭計。 \end{tabularx} \\ \\ \relax
6:12 & \begin{tabularx}{0.7\textwidth}{X} 因為我們的爭戰並不是對抗有血有肉的人，而是對抗那些執政的、掌權的、管轄這幽暗世界的，以及天空靈界的惡魔。 \end{tabularx} \\ \\ \relax
6:13 & \begin{tabularx}{0.7\textwidth}{X} 所以，要拿起神所賜的全副軍裝，好在邪惡的日子能抵擋仇敵，並且完成了一切後還能站立得住。 \end{tabularx} \\ \\ \relax
6:14 & \begin{tabularx}{0.7\textwidth}{X} 所以，要站穩了，用真理當作帶子束腰，用公義當作護心鏡遮胸， \end{tabularx} \\ \\ \relax
6:15 & \begin{tabularx}{0.7\textwidth}{X} 又用和平的福音當作預備走路的鞋穿在腳上。 \end{tabularx} \\ \\ \relax
6:16 & \begin{tabularx}{0.7\textwidth}{X} 此外，要拿信德當作盾牌，用來撲滅那惡者一切燒著的箭。 \end{tabularx} \\ \\ \relax
6:17 & \begin{tabularx}{0.7\textwidth}{X} 要戴上救恩的頭盔，拿著聖靈的寶劍—就是神的道。 \end{tabularx} \\ \\ \relax
6:18 & \begin{tabularx}{0.7\textwidth}{X} 要靠著聖靈，隨時多方禱告祈求，並要為此警醒不倦，為眾聖徒祈求。 \end{tabularx} \\ \\ \relax
6:19 & \begin{tabularx}{0.7\textwidth}{X} 也要為我祈求，讓我有口才，能放膽開口講明福音的奧祕， \end{tabularx} \\ \\ \relax
6:20 & \begin{tabularx}{0.7\textwidth}{X} 我為這福音的奧祕作了帶鐵鏈的使者，讓我能照著當盡的本分放膽宣講。 \end{tabularx} \\ \\
[1ex]
\hline
\hline
\end{longtable}
約翰福音 6:56-69
\newline
\begin{longtable}{cl}
\hline
\hline
章節 & 經文 (和合本修訂版)\\
\hline
6:56 & \begin{tabularx}{0.7\textwidth}{X} 吃我肉、喝我血的人常在我裡面，我也常在他裡面。 \end{tabularx} \\ \\ \relax
6:57 & \begin{tabularx}{0.7\textwidth}{X} 永生的父怎樣差我來，我又怎樣因父活著，照樣，吃我肉的人也要因我活著。 \end{tabularx} \\ \\ \relax
6:58 & \begin{tabularx}{0.7\textwidth}{X} 這是從天上降下來的糧，不像你們的祖宗吃過嗎哪還是死了；吃這糧的人將永遠活著。」 \end{tabularx} \\ \\ \relax
6:59 & \begin{tabularx}{0.7\textwidth}{X} 這些話是耶穌在迦百農會堂裡教導人的時候說的。 \end{tabularx} \\ \\ \relax
6:60 & \begin{tabularx}{0.7\textwidth}{X} 他的門徒中有好些人聽見了，就說：「這話很難，誰聽得進呢？」 \end{tabularx} \\ \\ \relax
6:61 & \begin{tabularx}{0.7\textwidth}{X} 耶穌心裡知道門徒為這話私下議論，就對他們說：「這話成了你們的絆腳石嗎？ \end{tabularx} \\ \\ \relax
6:62 & \begin{tabularx}{0.7\textwidth}{X} 如果你們看見人子升到他原來所在之處，會怎麼樣呢？ \end{tabularx} \\ \\ \relax
6:63 & \begin{tabularx}{0.7\textwidth}{X} 聖靈賜人生命，肉體毫無用處。我對你們所說的話就是靈，就是生命。 \end{tabularx} \\ \\ \relax
6:64 & \begin{tabularx}{0.7\textwidth}{X} 可是你們中間有些人不信。」耶穌起初就知道哪些人不信他，哪一個要出賣他。 \end{tabularx} \\ \\ \relax
6:65 & \begin{tabularx}{0.7\textwidth}{X} 於是耶穌說：「所以，我對你們說過，若不是蒙我父的恩賜，沒有人能到我這裡來。」 \end{tabularx} \\ \\ \relax
6:66 & \begin{tabularx}{0.7\textwidth}{X} 從此，他門徒中有很多退卻了，不再和他同行。 \end{tabularx} \\ \\ \relax
6:67 & \begin{tabularx}{0.7\textwidth}{X} 耶穌就對那十二使徒說：「你們也要離開嗎？」 \end{tabularx} \\ \\ \relax
6:68 & \begin{tabularx}{0.7\textwidth}{X} 西門‧彼得回答他：「主啊，你有永生之道，我們還跟從誰呢？ \end{tabularx} \\ \\ \relax
6:69 & \begin{tabularx}{0.7\textwidth}{X} 我們已經信了，又知道你是神的聖者。」 \end{tabularx} \\ \\
[1ex]
\hline
\hline
\end{longtable}
講題:信仰路上的四個重要時刻　Four Critical Moments of a Pilgrimage

經文:約書亞記24:1-18;詩篇34;以弗所書6:10-20;約翰福音6:56-69

講員:林豪恩先生

場合:香港中文大學崇基學院禮拜堂

日期:2015年8月23日

引言

眾所周知,維持一段關係很不容易,由拍拖到結婚,難免想過是否要分手；結婚之後,以為「公主和王子從此以後就快快樂樂地住在一起」,怎知道,離婚的念頭竟然偷偷地潛入不少夫妻的腦海中.有一位作者說,若是有人作調查,應該不難發現,目前活在地球上的已婚人士,從沒想過「我要離婚」的人大概不多!

基督徒維持和上主的關係,也是一件不容易的事.相信為數不少的基督徒在信仰路上都掙扎過是否要繼續前行,甚至想過離開.今天經課的四段經文讓我們意識到,信仰路上的確有不少挑戰,這是考驗信仰的時刻,也是重新確認信仰的時刻.

一．約書亞記24:1-18:進入新階段的信仰重新確認

24:1約書亞召集以色列的眾支派到示劍,他召了以色列的長老,領袖,審判官和官長來；他們都站在上帝面前.

以色列人出埃及之後,在曠野漂流了四十年,終於來到應許之地──迦南.進入迦南,代表著以色列人進入新的階段.新的階段是充滿變數的,可能有無限的機遇,也可能有無盡的挑戰,甚至可能是迷失的開始.有一位讀幼稚園的小孩子問他的父母:「為甚麼我要努力讀書?」他的父母告訴他:「因為要讀好的小學.」當他讀小學時又問父母:「為甚麼我仍要努力讀書?」「因為要讀好的中學.」當他讀中學的時候又問父母:「為甚麼我仍要努力讀書?」「因為要進入大學.」當他進入大學後,他感到興奮,但漸漸又感到迷茫.以色列人幾十年來,跨越兩代所追尋的應許之地,已經到達了,然後又如何呢?在這新的階段,應該如何生活呢?帶領以色列人進入迦南的領袖約書亞特別關注人民的信仰狀況.在新的階段,放在以色列人面前的,有眾多信仰選擇.應該如何面對這些選擇呢?

24:2約書亞對眾百姓說:

24:14「現在你們要敬畏耶和華,誠心誠意事奉他,除掉你們列祖在大河那邊和在埃及事奉的神明,事奉耶和華.

24:15若你們認為事奉耶和華不好,今日就可以選擇所要事奉的:是你們列祖在大河那邊所事奉的神明,或是你們所住這地亞摩利人的神明呢?至於我和我家,我們必定事奉耶和華.」

約書亞指出,「市面上」有很多信仰選擇,就以色列人所熟知的,包括有:「在埃及事奉的神明」,「你們列祖在大河那邊所事奉的神明」(即是亞述和巴比倫等地的信仰),「你們所住這地亞摩利人的神明」.從現實上考慮,埃及是南面的強國,大河那邊是北方強盛的勢力,選擇他們的信仰,有靠背強國的優勢.另一方面,若從當下處境考慮,「這地亞摩利人的神明」也是一個明哲保身的選擇,有利得到當地人的接納和幫助.如此多有優勢的選擇,應該如何看待他們本身對耶和華上帝的信仰呢?

香港教會更新運動早前發表了《2014香港教會普查簡報》,特別指示大專及職青信徒的流失狀況.報告顯示,1999年時,教會佔有這組年齡層41.6\%,至2014年,教會於這組年齡層跌至27.4\% (全港為31.0\%).很多教會表示,近年來,大專信徒至初職10年以內信徒的流失嚴重.(\url{http://www.hkchurch.org/GenericStyles/Content.asp?ID=11650 and PaperID=0010})

有人形容,基督徒中學生進入大專校園,或者基督徒大學生進入工作世界,好像小朋友進入主題公園一樣,樂而忘返.又有人形容,這些基督徒進入一個新階段,好像羊進入狼群一樣,有些變成了獵物,有些變成了狼.新的階段,有很多新的選擇,這些選擇,好像比原來的更吸引；新的階段,有很多新的挑戰,原本的信仰,好像已經應付不了新的處境.

進入迦南開始新的階段的以色列人,面對眾多看似更有優勢的選擇,如何回應呢?

24:16百姓回答說:「我們絕不離棄耶和華去事奉別神.

在新的階段,面對新的挑戰,他們重新確認對耶和華的信仰,以及對信仰的委身.是甚麼令他們對信仰作出重新確認和委身呢?

24:17因為耶和華－我們的上帝曾領我們和我們的祖宗從埃及地為奴之家出來,在我們眼前行了那些大神蹟,並在我們所行的一切路上,和所經過的各民族中保護了我們.

24:18耶和華又把各民族和住此地的亞摩利人都從我們面前趕出去.所以,我們也必事奉耶和華,因為他是我們的上帝.」

他們對耶和華的信仰,是他們的親身經歷.對別人而言,或者這些只是「事件」,但對當事人而言,這是「經歷」；別人可以把你的「經歷」用各種學說解釋掉,但對當事人來說,這些經歷是人生的一部份,不容否定的.踏入新的階段,面對新的挑戰,縱使未必能肯定本來的信仰仍否管用.然而,回望過去,某時某地某事件,那次的信仰經歷仍是真實的.這些真實的信仰經歷片段,是信徒進入新的人生階段再一次確認信仰的基礎.

二．詩篇34:善惡交戰中的信仰確認

34:15耶和華的眼目看顧義人,他的耳朵聽他們的呼求.

34:16耶和華向行惡的人變臉,要從地上除滅他們的名字.

34:17義人呼求,耶和華聽見了,就拯救他們脫離一切患難.

34:18耶和華靠近傷心的人,拯救心靈痛悔的人.

34:19義人多有苦難,但耶和華救他脫離這一切,

34:20又保護他全身的骨頭,連一根也不折斷.

34:21惡必害死惡人,恨惡義人的,必被定罪.

34:22耶和華救贖他僕人的性命,凡投靠他的,必不致定罪.

此詩背景

21:10那日大衛起來,躲避掃羅,逃到迦特王亞吉那裡.

21:11亞吉的臣僕對他說:「這不是那地的國王大衛嗎?那裡的人跳舞唱和:『掃羅殺死千千,大衛殺死萬萬』,不是指著他說的嗎?」

21:12大衛把這些話放在心裡,就很懼怕迦特王亞吉.

21:13於是他在眾人眼前一反常態,在他們中間裝瘋作癲,在城門的門扇上胡寫亂畫,任由唾沫流在鬍子上.

21:14亞吉對臣僕說:「看哪,你們看這人瘋了,為甚麼帶他到我這裡來呢?

21:15我豈缺少瘋子,你們竟然帶這人到我面前瘋癲嗎?這個人可以進我的家嗎?」

就是因為有人宣傳:「掃羅殺死千千,大衛殺死萬萬」,令大衛被掃羅追殺,四處逃走,流離失所.在受苦中之中的人很可能會問:「到底我做了甚麼事情令我得到如此的對待呢?」大衛堅持自己對掃羅沒有惡意,也沒有做過對不起他的事情,他被追殺,只因為掃羅對他的嫉妒.對一個信仰者來說,上帝應該是賞善罰惡的,惡人當道,自己卻無辜受苦,這會令人對信仰提出質疑,動搖信仰的基礎.

約伯無辜受苦,約伯的妻子對他說:2:9他的妻子對他說:「你仍然持守你的純正嗎?你背棄上帝,死了吧!」人們都相信,上帝是「德福一致」的,也就是說,有德者必能蒙福.在約伯的觀念裡,只要敬畏上帝,上帝必能保住他的產業且不止息地賜福他,而只要做到「無可指謫」,上帝就無降禍的理由.約伯的痛苦,不單單是遭遇災難,更在於作為虔誠的信徒,陷入善惡混雜,禍福不分的失序世界.

34:17義人呼求,耶和華聽見了,就拯救他們脫離一切患難.

詩人表示,義人是可以有患難的,患難臨到,並不等於你不是義人.上帝的公義,不在於不讓患難出現在義人身上,而在於對身處患難中的義人予以肯定,並且關注和行動.

34:19義人多有苦難,但耶和華救他脫離這一切

義人是可以有苦難的,甚至是「多有苦難」,就正如耶穌說:「為義受逼迫」.詩人再次肯定,上帝的公義,在於顧念遇到苦難的義人,並且出手相救.

舊約中不少先知也因為看見惡人當道,義人受難而詰問上帝.我們作為平凡的基督徒,更難免因為自己的不如意遭遇,但又明明不是自己做錯了甚麼而得來的,而質疑上帝是否公義.詩人提醒我們,這些時候是再一次確認信仰的時候.歷世歷代不少在社會中試圖實踐「行公義,好憐憫」的信徒,看見世界惡人當道,公義不彰的時候,都曾經陷入信仰的掙扎.當他們再次確認信仰後,他們更有基礎及盼望在世上實踐「行公義,好憐憫」的使命.

三．以弗所書6:10-20:為信仰而戰的信仰確認

經過了新階段的衝擊,又經過了正邪混雜的考驗,信仰已經相當穩定.這個時候,可能滿懷使命,為上帝作見證,為信仰作出貢獻.在這個似乎是最不需要再作出信仰確認的時候,保羅認為又是要作出信仰確認的時候.

6:10最後,你們要靠著主,依賴他的大能大力作剛強的人.

6:11要穿戴上帝所賜的全副軍裝,好抵擋魔鬼的詭計.

6:12因為我們的爭戰並不是對抗有血有肉的人,而是對抗那些執政的,掌權的,管轄這幽暗世界的,以及天空靈界的惡魔.

6:13所以,要拿起上帝所賜的全副軍裝,好在邪惡的日子能抵擋仇敵,並且完成了一切後還能站立得住.

6:14所以,要站穩了,用真理當作帶子束腰,用公義當作護心鏡遮胸,

6:15又用和平的福音當作預備走路的鞋穿在腳上.

6:16此外,要拿信德當作盾牌,用來撲滅那惡者一切燒著的箭.

6:17要戴上救恩的頭盔,拿著聖靈的寶劍-就是上帝的道.

去年六月底,新加坡警察商業調查局正式起訴新加坡最大基督教會的辦創人康希牧師與另外五位教會領袖; 據稱商業事務局與新加坡慈善總監己調查已有兩年時間, 起訴指控財務違規, 他們私自盜用教會建堂基金用作康牧師師母何耀珊的音樂事業上. 新加坡慈善總監稱,從2010年5月開始的調查發現了財務違規行為,該教會被盜用的慈善資金總額至少為2,300萬新加坡元(1,800萬美元).(jgospel.net/news/christian-news)

這是甚麼一回事呢?這不是一般為了私人利益的虧空公款,而是源於一個見證動機,這是城市豐收教會的「跨界計劃(Crossover Project)」.基督徒音樂評論人如此評論這事件:2002年康希夫婦開始這個事工,目的是要「使用何耀珊的歌唱和音樂,來接觸從來沒有聽過福音的人和地」.何耀珊想要在短時間內快速崛起,非要借力使力不可 ,她向荷里活投石問路的首支單曲”Where Did Love Go” 找來天王級製作人David Foster操刀.Foster是誰呢?他曾製作過Whitney Houston,Madonna,Michael Jackson等超級巨星的唱片,是當前美國音樂界最火的人物.我十分佩服何耀珊的野心和格局,試問亞洲有哪一位歌手能請得到這樣的製作人?(讀者大概能嗅出這一千九百萬美元的跨界基金的一部份用到哪裡去了.)接著的十年,我們看到「跨界計劃」的執行團隊,非常努力的打造何耀珊成為流行巨星,這些年何的作品在Billboard Hot Dance Chart(美國告示牌舞曲榜)屢獲佳績,從創意發想,行銷包裝,到公關造勢,處處看得出跨界團隊的專業和用心.然而,令人遺憾的是,他們忽略了「跨界計劃」最重要的核心價值──用音樂做為媒介來接觸未信的朋友.網友們對何耀珊這些作品的看法,大部份是負面的,大家都很納悶,都在問,這樣的音樂要如何「跨界」?用什麼信息內容與未信的人「對話」?從這個角度來看,「跨界計劃」的音樂創作是不及格的,從China Wine到Mr. Bill,從Where Did Love Go到Fancy Free,容我用比較重的語氣說,這些作品反映的是時下被罪沾污的世界觀,看不到希望,看不到出路,也完全看不到被基督更新的生命.跨界團隊對音樂題材的選取上,似乎只考慮到行銷和賣點,完全不顧內容是否與所宣揚的信仰有所抵觸.何在創作這些作品的同時,大概已經忘了,她是基督徒,是師母,而且正在執行「跨界計劃」,然而諷刺的是,聽眾沒忘,網友也沒忘.在宣教的領域裡,我欣喜跨界計劃正在大膽摸索一條不太有人敢走的新路,他們勇於接觸一般教會無法接近的群體,但這個夢,碎了.這一切,使得那飽受質疑的一千九百萬美元跨界宣教基金,在用途的正當性上找不到具有說服力的理由；而那個「為義受逼迫,被人誤解」的故事劇本也變得難以自圓其說,整個事件淪為一場鬧劇,也讓「跨界」這個立意良好的概念蒙上一層陰影.(王星然 : 「康希」來了－談何耀珊的跨界音樂)

在這個「顧客就是上帝」,「市場就是真理」的時代,凡是懷著使命要在這個世代為信仰作見證的基督徒,都可以在保羅的這段教導中得提醒,並且再一次對信仰作出確認.

6:12因為我們的爭戰並不是對抗有血有肉的人(生得見的人),而是對抗那些執政的,掌權的,管轄這幽暗世界的,以及天空靈界的惡魔(看不見的權勢).

6:13「穿戴上帝所賜的全副軍裝」還是「世界所宣揚的全套市場策略」

6:14「用真理當作帶子束腰」還是「用包裝當作帶子束腰」,「用公義當作護心鏡遮胸」還是「用公關當作護心鏡遮胸」

6:15「用和平的福音當作預備走路的鞋穿在腳上」還是「用財富與健康的福音當作預備走路的鞋穿在腳上」

6:16「拿信德當作盾牌」還是「拿市場口味當作盾牌」

6:17「戴上救恩的頭盔」還是「戴上名氣作高帽」,「拿著聖靈的寶劍-就是上帝的道」還是「拿著市場的潛規則」.

這好像最不需要再次確認信仰的時候,卻是最需要重新確認信仰的時候.

最後,如果信仰沒有市場,那又當如何面對自己,面對上帝呢?

四．約翰福音6:56-69:多人退去時的信仰確認

6:57永生的父怎樣差我來,我又怎樣因父活著,照樣,吃我肉的人也要因我活著.

6:58這是從天上降下來的糧,不像你們的祖宗吃過嗎哪還是死了；吃這糧的人將永遠活著.」

6:59這些話是耶穌在迦百農會堂裡教導人的時候說的.

永生的話

6:60他的門徒中有好些人聽見了,就說:「這話很難,誰聽得進呢?」

6:61耶穌心裡知道門徒為這話私下議論,就對他們說:「這話成了你們的絆腳石嗎?

6:62如果你們看見人子升到他原來所在之處,會怎麼樣呢?

6:63聖靈賜人生命,肉體毫無用處.我對你們所說的話就是靈,就是生命.

6:64可是你們中間有些人不信.」耶穌起初就知道哪些人不信他,哪一個要出賣他.

6:65於是耶穌說:「所以,我對你們說過,若不是蒙我父的恩賜,沒有人能到我這裡來.」

6:66從此,他門徒中有很多退卻了,不再和他同行.

6:67耶穌就對那十二使徒說:「你們也要離開嗎?」

6:68西門‧彼得回答他:「主啊,你有永生之道,我們還跟從誰呢?

6:69我們已經信了,又知道你是上帝的聖者.」

二○○九至二○一一年間,華人信念發展研究中的「信徒靈性發展研究」以網上問卷調查了數千名基督徒.他們的年齡由16至67歲,平均是29歲.究竟離開信仰的人在一年前有甚麼特徵,是與那些仍持守信仰的人有所不同?分析顯示,他們有較大比例是:(一)全時間的大學生(大學生活)；(二)忙(世務纏身)；(三)不重視傳統價值觀念；(四)當時沒有參加教會聚會(缺乏教會生活).(許志超,劉永發,劉月瑩,張樹輝:「離開信仰的一些先兆」)

在信仰的路上,總有人因不同的原因而離開的.耶穌並非是報喜不報憂的宣傳機器,有人離他而去,耶穌坦然面對,聖經如實記載.

6:66從此,他門徒中有很多退卻了,不再和他同行.

6:67耶穌就對那十二使徒說:「你們也要離開嗎?」

你們也要離開嗎?為甚麼留下來呢?

葛凡西／法蘭西斯·柯林斯(Francis S. Collins),是一位遺傳學家,領導人類基因組計劃,並發現了多種疾病基因.1989年與徐立之共同發現囊性纖維化的致病基因(通稱CF gene).葛氏讀大學時物理和化學科的成績出眾,並在一九七四年輕易地完成了他在耶魯大學的博士學位.讀書期間,他都以無神論者自稱,拿到博士學位後,他看見醫學界中遺傳因子學的挑戰,於是進修醫科,並於一九七七年拿到他的醫學學位.之後他在一所醫院工作並研究遺傳因子學.葛氏在醫院實習時,常看見許多人在垂死邊緣,然而他發現那些自稱是基督徒的,在面對死亡時,是多麼有平安,跟一般人不一樣.一天,有個每天受盡病痛折磨,卻因信基督而滿有平安的老婦問他說:「你信甚麼?」那一刻來得尷尬極了,因為他只能回答:「我不太肯定.」當時廿六歲的他,已當醫生,擁有兩個博士學位,對生命,疾病和死亡應已有答案；然而當他被問及有關他自己永恆歸宿這最基本的問題時,卻是啞口無言.(伍煒國:天外有天)

很多年之後,他寫了一本書,名為《上帝的語言》(The Language of God).這本書寫到,生物學,天體物理學和心理學都支持對上帝的信仰.他表示:「一旦正確地被理解和尊重,科學與信仰其實是可以相互豐富和補充的」,「有一些科學無法回答的非常重要的問題……我發現信仰為找到答案提供了更好的途徑.」

6:67耶穌就對那十二使徒說:「你們也要離開嗎?」

6:68西門‧彼得回答他:「主啊,你有永生之道,我們還跟從誰呢?」

你呢?你還要跟從誰呢?這是否也是你要作出信仰確認的時刻呢?



\newpage

\section{從餐桌掉下來的餅碎}
\label{sec:296}
\textbf{伍渭文牧師 - 從餐桌掉下來的餅碎}
\newline
\newline
\hyperref[sec:295]{< < < PREV SERMON < < <}
~
\hyperref[sec:toc]{[返目錄]}
~
\hyperref[sec:297]{> > > NEXT SERMON > > >}


馬可福音 7:24-30
\newline
\begin{longtable}{cl}
\hline
\hline
章節 & 經文 (和合本修訂版)\\
\hline
7:24 & \begin{tabularx}{0.7\textwidth}{X} 耶穌從那裡起身，往推羅境內去，進了一家，他不願意人知道，卻隱藏不住。 \end{tabularx} \\ \\ \relax
7:25 & \begin{tabularx}{0.7\textwidth}{X} 立刻有一個婦人，她的小女兒被污靈附著，一聽見耶穌的事，就來俯伏在他腳前。 \end{tabularx} \\ \\ \relax
7:26 & \begin{tabularx}{0.7\textwidth}{X} 這婦人是希臘人，屬敘利亞的腓尼基族。她求耶穌從她女兒身上趕出那鬼。 \end{tabularx} \\ \\ \relax
7:27 & \begin{tabularx}{0.7\textwidth}{X} 耶穌對她說：「讓孩子們先吃飽，拿孩子的餅丟給小狗吃是不妥的。」 \end{tabularx} \\ \\ \relax
7:28 & \begin{tabularx}{0.7\textwidth}{X} 婦人回答：「主啊，桌子底下的小狗也吃小孩子的碎屑呀！」 \end{tabularx} \\ \\ \relax
7:29 & \begin{tabularx}{0.7\textwidth}{X} 耶穌對她說：「憑著這句話，你回去吧，鬼已經離開你的女兒了。」 \end{tabularx} \\ \\ \relax
7:30 & \begin{tabularx}{0.7\textwidth}{X} 她就回家去，見小孩子躺在床上，鬼已經出去了。 \end{tabularx} \\ \\
[1ex]
\hline
\hline
\end{longtable}
講題:從餐桌子掉下來的餅碎　The Scraps from the Table

經文:馬可福音7章24-30節

講員:伍渭文牧師

場合:香港中文大學崇基學院禮拜堂

日期:2015年9月6日

猶太人有一個傳統:收割時不要拾起掉下來的麥穗,讓窮人和孤兒寡婦拿起來充飢.路得記(Ruth )第二章記載外邦摩押(Moabite)女子耶穌先祖路得在田間拾麥穗,被地主波亞斯(Boaz)看見,開始發展一個 浪漫愛情故事.今天的福音書也說到一位外邦敘利腓利基族(Syrian Phoenicia)女子,乞求從桌子上掉下來的食物,餵養她的女兒.她的女兒被污鬼附身,她俯伏在耶穌面前求耶穌趕走污鬼.

推羅西頓是加利利北方沿海的外邦人港口,在現代黎巴嫩南,北兩端.耶穌進入推羅西頓境內原本是避靜,不想人和道.耶穌在世首先在猶太人中間工作,他們是上帝的選民,他就是選民所期待的彌賽亞.所以耶穌對她說:『讓兒女先吃飽,不好拿兒女的餅丟給(小)狗吃.』這話表面聽來很刺耳,耶穌好像很不近人情,而且語帶刻薄.但這裡的重點,不在兒女和 和狗的分別,乃在得到餅的先後次序.耶穌用當時社會用膳時熟悉的圖畫,說明救恩先臨到猶太人,後才來到外邦人,就像在餐桌上與父親共席的兒女先得餅食,在桌下等待的狗才得飽足.在當時社會,狗也是家中要員；狗不是寵物,給我們呵護.狗是守護者,忠誠的朋友,我們幼小兒女的玩伴,晚上睡覺時保護我們,防備盜賊和野獸入侵；有些狗(牧羊犬)更幫助我們牧羊.耶穌用小狗的字,強調狗和家人的親密關系.所以,用膳時狗也圍繞在桌下,並期待分享應有的一份.

為何福音從猶太人開始呢?上帝的國先在預備好了的土壤,才能生根發芽長大.我們都曉得,種稻米的人,先把種子下在秧田(seedbed) 培育 ,讓它發芽成為秧苗,才在田裡插秧.猶太人是選民,有上帝的律法,但因為拜偶像,經歷諸般苦難.他們被亞述國人侵,強行雜婚,混亂種族宗教；又被擄 巴比倫,數百年顛沛流離,痛定思痛,才能澈底棄絕偶像,信奉獨一上主.也因為被擄巴比倫,耶路撒冷聖殿被毀,崇拜地點從聖殿轉到遍地開花的會堂,崇拜禮儀重點從獻祭轉向律法的誦讀和講解.

崇拜誦讀經課的傳統就是從被擄時期得到強化和完善,每安息日首輪經課誦讀律法書(創,出,利,民,申五經),一年讀完；次輪讀大,小先知書,作為律法書的解釋.耶穌第一次進入迦百農會堂讀出以賽亞書,應該是次輪經課.教會承繼了這傳統,一年讀完一卷福音書,並以舊約,詩篇,書信來解釋福音書；鮮有宗教像基督信仰那麼注重經典的研讀.被擄回歸耶路撒冷後,又多年遭希臘,羅馬鐵蹄殖民統治,更使他們盼望彌賽亞的降臨. 福音種子很難在多神的異教土壤生長.

我們再看可七25-30: 當下,有一個婦人,他的小女兒被污鬼附著,聽見耶穌的事,就來俯伏在他腳前.這婦人是希利尼人,屬敘利非尼基族.她求耶穌趕出那鬼離開她的女兒.耶穌對她說:「讓兒女們先吃飽,不好拿兒女的餅丟給狗吃.」婦人回答說:「主啊,不錯；但是狗在桌子底下也吃孩子們的碎渣兒.」耶穌對她說:「因這句話,你回去吧；鬼已經離開你的女兒了.」她就回家去,見小孩子躺在床上,鬼已經出去了.

一 信心使我們得著上主的餅

禮失而求諸野.耶穌同意福音先臨到猶太人,以後才傳到外邦人,但猶太人拒絕自己的彌賽亞,祝福就臨到這外邦的女子.「讓兒女先吃飽」是原則,但當兒女拒絕餅時,餅就掉落餐桌下,給饑餓等待食物的家犬.馬可引入這外邦婦人求餐桌掉下來的碎渣兒,來對比作為先得餅的兒女猶太人,禮失而求諸野.

可六55, 56總結當時耶穌的工作:「凡耶穌所到的地方,或村中,或城裡,或鄉間,他們都將病人放在街市上,求耶穌祇容他們摸他的衣裳繸子；凡摸著的人就都好了.」太四24更說:「那裡的人把一切害病的,就是害各樣疾病,各樣疼痛的和被鬼附的,癲癇的,癱瘓的,都帶了來,耶穌就治好了他們.」

兩批不信的猶太人.但可六1-6指出,最認識耶穌的家鄉猶太人竟然對這些神蹟亳無感動,反而厭煩他(They took offence at him,可六3).甚麼是厭煩呢?厭煩是沒有理性的敵視情緒,就是不喜歡,唔liki, 先入為主的偏見:「這人從那裡有這些事呢?所賜給他是甚麼智慧?他手所做的是何等的異能呢?這不是木匠嗎?不是馬利亞的兒子雅各,約西,猶太,西門的長兄嗎?他妹妹們不也是在我們這裡嗎?」(可六2, 3)不錯,先知在本地一般不受尊敬,但他們特別的硬心和不信令耶穌也覺詫異.請注意:因為他們的不信,耶穌在那裡不得行甚麼異能,而且離開他們往周圍鄉村教訓人去了.要對上主有信心,才能經歷上主的能力,看見上主的作為,這是一個屬靈原則.

另一批拒絕耶穌的是法利賽人和文士,聖經特別題及他們「從耶路撒冷來,到耶穌那裡聚集.」極有可能他們來實地瞭解耶穌在加利利一帶所作所為.作為文士和法利賽人,他們熟悉聖經,也教導聖經,是猶太人的宗教領袖,他們是上帝兒女中的兒女,餐桌上的餅首先給他們的,但他們拒絕這些餅.他們批評耶穌的門徒不按傳統完成洗手的禮儀就吃飯; 堅持從市集俗世之地回來,謝飯前先要洗手脫俗；而對剛發生五餅二魚吃飽五千人,和耶穌在海面行走的震撼性神蹟視若無賭,沒有興趣；祇談吃飯前洗俗手,他們轉移視線,模糊焦點.耶穌責備他們:「你們是離棄神的誡命,拘守人的遺傳.」(可七8)

偏見與扭曲.為何耶穌嚴嚴責備他們,而對家鄉的猶太人的不信表示詫異呢?因為後者不信是偏見,但前者是扭曲；耶穌趕鬼,他們可以扭曲說耶穌是靠鬼王趕鬼.作為律法的守護者和教化者扭曲律法的原意,破壞性更大更危險,因為牽涉到是非對錯.睡著的人可以叫醒,但裝睡的人永遠叫不醒,裝睡的人是有意識地拒絕清醒；偏見是可以慢慢糾正的無知,扭曲是對抗真理的邪惡；偏見是明顯的拒絕,扭曲是暗裡的顛覆.

相比之下,這外邦婦人有令耶穌嘉許的信心.她接受自己身份的限制,耶穌是猶太人所盼望的彌賽亞,她不是猶太人,所以她沒有像其它猶太人把女兒帶出來,放在耶穌所經之處,摸著他衣裳的繸子得醫治.她也接受耶穌的話:「讓兒女們先吃飽,不好拿兒女的餅丟給狗吃.」但她抓著上主的憐憫:「主啊,不錯；但是狗在桌子底下也吃孩子們的碎渣兒.」

屬靈的洞見,經課的主題.這句話有透穿上主心意的屬靈洞見:上主愛每一個人,不分種族,不分階級,特別社會上最無依靠的人,上主更加憐愛.今天誦讀其它三段經課也說明這點.以賽亞三十五4-7上:彌賽亞來的時候,「要使軟弱的手堅壯,無力的膝穩固.對膽怯的人說:你們要剛強,不要懼怕.」詩篇一四六7-9「他為受屈的伸冤,賜食物與饑餓的.耶和華釋放被囚的.耶和華開了瞎子的眼睛.耶和華扶起被壓下的人.耶和華喜愛義人.耶和華保護寄居的,扶持孤兒和寡婦,卻使惡人的道路彎曲.」；雅各書二5,13「神豈不是揀選了世上的貧窮人,叫他們在信上富足.……那不憐憫人的,也要受無憐憫的審判,憐憫原是向審判誇勝.」

我們崇拜讀出四段經課,彼此詮釋,互通款曲,幫助我們更明白當天的信息,很明顯,關心憐恤周遭有需要的人是經課的共通主題.

信心來自對自己的認識.為何婦人有如斯的信心呢?因為她覺得需要,她小女兒被污鬼附身,在異教的城邑,可能已試盡不少神靈,但無濟於事.推羅是古代地中海著名海港,不少神祇隨商貿傳入推羅,推羅曾有大祭司謁巴力(Ethbaal)成為國王,女兒耶洗別(Jezebel)嫁給亞哈王(Ahab)為皇后,把以色列陷在雨神巴力的偶像敬拜中.對比第五章睚魯求耶穌醫好自己的女兒,婦人信心的跨度是大的.睚魯身為管會堂的猶太人意見領袖,求耶穌醫好自己女兒要力排眾議,睚魯地位崇高,挺得住.但婦人要跨越異教多神信仰跨度很大,當時宗教是全面滲透到各生活層面,離棄異教信靠耶穌需要很大的決心.相對睚魯,她勢孤力弱.睚魯親自求耶穌,因他家中有婢僕看守重病的女兒；這婦人自己俯伏在耶穌面前,為何不是作父親求耶穌呢?可能父親在家看守著女兒,可能她丈夫早已離世,她是寡婦,和小女兒相依為命.

她的信心來自認識自己的無助,她需要耶穌.家鄉的猶太人和來自耶路撒冷的文士和法利賽人呢?他們自滿,他們認為全都知道:我們看著耶穌長大的!我們認識律法及傳統!因此他們覺得不需要耶穌.

發現地心吸力的以撒牛頓(Isaac Newton)是一位虔誠的基督徒,他對上主的信心來自對自己的認識.臨終時他說:我像一個小孩在茫茫大海的海邊,偶然在沙灘發現一些可愛的貝殼和小滑石,令我欣喜不已；但我所不知的,就像那茫茫無邊際的大海.所以牛頓謙遜,牛頓篤信.信心是認識自己的不足,認識自己的需要.

這婦人謙卑等候桌上掉下來碎渣兒,她等候上主的憐憫,但需憐憫的人竟然有這深邃的屬靈洞察力；她是今天雅各書經課所說的貧乏人,但上主使她信上富足.八福中說:憐恤人的人有福了,因為他們必蒙憐恤.憐恤人時,上主也憐恤我們；憐恤人時,接受我們幫助的人也會教曉我們明白生命的意義.

二 愛心使我們分享上主的餅

「主啊,不錯；但是狗在桌子底下也吃孩子們的碎渣兒.」因這句話,你回去吧；鬼已經離開你的女兒了.」可七9.

有一位在智障學校工作的校長說,我的學生就是我的天使,我的屬靈導師,教曉我何為單純的信心,何為知足.靈修神學大師盧雲(Neuwen)神父辭掉耶魯大學的教席,到多倫多一間收留發展障礙的兒童院「黎明之家」(L.’Arches Daybreak)工作,寫出洞悉人性,滿有屬靈睿智的感人著作.黃仁龍前律政司長有一次在崇基學院的周會分享他出身基層,做了大律師後保持每月探訪露宿者,不單為露宿者,也是為了自己；不要在名氣,富裕中迷失了自我.因為探訪露宿者,以後看到一大疊文件時,不是見到一堆文字和數字,乃看到這些文件所代表眾多的臉龐,他們尋求公義,等候申訴.

最近在微信朋友圈看到強力推薦一篇感人短文「善良無須考核」,有關1998年巴西電影「中央車站」選角的故事,此片當年獲柏林電影節金熊獎,好萊烏外國記者協會的金球獎最佳外語片.導演沃爾特(Martine de Clemont-Tonnerre)在選角時挑選了很多藝校學生,但都不夠滿意.

一天,為此一籌莫展的沃爾特在西郊火車站的站前廣場上遇到了一個十多歲的擦鞋小男孩.小男孩問道:「先生,您需要擦鞋嗎?」沃爾特低頭看看自己腳上剛剛擦過不久的皮鞋,搖搖頭拒絕了.那個小男孩紅著臉追上來:「先生,我整整一天沒吃東西了,您能借給我幾個錢嗎?我從明天開始多多努力擦鞋,保證一周後把錢還給您!」 沃爾特看著面前這個衣衫襤褸,肚子乾癟的小男孩,不由的動了憐憫之心,就掏出幾枚硬幣遞到小男孩手裡.

半個月後,他又一次經過西郊火車站時,突然看到一個瘦小的身影離的老遠就向他招手喊道:「先生,請等一等!我在這裡等您很久了,今天總算把錢還給您了!

沃爾特忽然發現這小孩很符合自己腦海中構想的小男孩角色,他把幾枚硬幣還給小男孩:「這點零錢是我誠心誠意給你的,就不用還了,明天你到市中心的影業公司導演辦公室來找我,我會給你一個很大的驚喜.」

第二天一大早,門衛就告訴沃爾特,說外面來了一大群孩子.他詫異的出去一看,就見那個小男孩興奮的跑過來,一臉天真的說:「先生,這些孩子都是同我一樣沒有父母的流浪孩子,他們也渴望有驚喜!」.

導演沒想到一個窮困流浪的孩子竟會有一顆如此善良的心!通過反復觀察和篩選,確實有幾個比小男孩更機靈,更適合出演劇本中的小主人公,但他最後還是只把小男孩留了下來,在錄用合同的免試原因一欄中只寫了幾個字:善良無須考核!因為他覺的小孩很善良,在自己面臨困境的時候,卻把本屬於自己一個人的希望,無私的分享給別人,而電影中的孩子,正是這樣一個善良,博大,無私的人.

這個小男孩叫文尼斯基(Vinicius de Oliveira),若干年後成為一家影視文化公司董事長的文尼斯基寫了一部自傳《我的演藝生涯》.在書的扉頁上面,是沃爾特的親筆題字:善良無須考核.下面則是他給予文尼斯基的評價:「是善良,曾經讓他把機遇讓給別的孩子；同樣也是善良,讓人生的機遇沒有錯過他!」

三 守住身份,履行責任,轉限制為機會

相對猶太人,生在敘利腓利基的希利尼人是外邦人,較遲接受福音,但這限制亦是機會.我們可以想想,當她回到腓利基,看見污鬼離開女兒,看到鄰居親友如何為她高興和驚訝,她可以為耶穌作美好的見証,也為推羅,西頓地中海古代著名海港的福音傳播,撒下種子.

打開眼睛.上主賜我們在人生不同的季節有不同的身份,目的使我們在不同的時間,接觸到不同的人,分享上主給我們的從天上有來的餅.我們要關心在我們周圍有需要 的人,他們就在我們桌子附近,桌子下邊,在他們有需要時,給予援手.其實他們祇需要掉下來的碎渣兒,就是那麼多.撥一電話問候 一下,停下腳步聽聽他們,寫一個電郵關心,寄一張問候卡致意,花一些時間探訪他們.面對大量亞拉伯難民,奧大利和德國贏得我們的尊重,因為他們分享桌上的餅,給難民蔭庇,他們的領土稱得上是神聖的領土.

愛心使信心保持活潑(faith active in charity)

對上主的信心和對人的愛心相輔相承,信是對上主白白恩典的回應,藉信稱義,別無其它方法.愛,對人的愛心,是對上主恩典的回應,心被恩感而愛上主所愛的人.沒有對人的愛對上主的信不會生根,愛心使信心活潑不墮,faith active in charity.

馬丁路德說:(因著信)基督徒是最自由的君王,(因著愛)同時也是最順服的僕人.沒有信心的愛容易使我們因現實的不公平而氣餒,所以聖經把愛心和最後審判連繫起來( 太二十五章).沒有愛心的信將信仰私人化,信仰成為個人的宗教消費品；不愛看得見的弟兄,怎可以說愛看不見的上主呢?

把握機會.我們不同的身份是是上主給我們的機會,但機會需要及時把握的,機會不會永遠等待你.機會就像一個人從前面走來,前額有頭髮,可以抓住,後面是兀的,過後就捉不住了.在上帝的國度開展過程,上主有他的計劃,也給我們不同的身份去完成他的計劃.我們小時候是兒女,將來是父母,再遲一些是祖父母.耶穌肯定猶太人和外邦人的界限,提醒我們守住身份,把握機會,履行應盡的本份.這些界限有時限的,離開工作崗位, 便不能與我們的同事共事了；兒女長大後,到就不能常伴左右了；大學畢業後,便沒有在學時那麼多時間讀書了.

信心超越限制.不同的身份也有不同的界限和限制,但像這婦人,信心超越外在身份的限制.保羅也教導我們:『並不分是猶太人,希利尼人,自主的,為奴的,或男或女,因為你們在基督裡都成為一了.』(加3;28) 從信心的角度看,限制亦是機會；每一張放食物的桌子都有它安放的空間,有它的桌下和它的毗鄰,而食物總有可分享的碎渣兒,祇要我們留心,總有可以服事的對像.

總結:信心使我們得到上主所賜的餅；愛心使我們分享上主的餅；不同身份是上主給我們的機遇,雖然有限制,但對上主的信心和對人的愛心,總會轉限制為機會.



\newpage

\section{你們在路上議論的是什麼}
\label{sec:297}
\textbf{鄧瑞強博士 - 你們在路上議論的是什麼?}
\newline
\newline
\hyperref[sec:296]{< < < PREV SERMON < < <}
~
\hyperref[sec:toc]{[返目錄]}
~
\hyperref[sec:298]{> > > NEXT SERMON > > >}


馬可福音 9:30-37
\newline
\begin{longtable}{cl}
\hline
\hline
章節 & 經文 (和合本修訂版)\\
\hline
9:30 & \begin{tabularx}{0.7\textwidth}{X} 他們離開那地方，經過加利利；耶穌不願意人知道， \end{tabularx} \\ \\ \relax
9:31 & \begin{tabularx}{0.7\textwidth}{X} 因為他正教導門徒說：「人子將要被交在人手裡，他們要殺害他；被殺以後，三天後他要復活。」 \end{tabularx} \\ \\ \relax
9:32 & \begin{tabularx}{0.7\textwidth}{X} 門徒卻不明白這話，又不敢問他。 \end{tabularx} \\ \\ \relax
9:33 & \begin{tabularx}{0.7\textwidth}{X} 他們來到迦百農。耶穌在屋裡問門徒說：「你們在路上議論的是甚麼？」 \end{tabularx} \\ \\ \relax
9:34 & \begin{tabularx}{0.7\textwidth}{X} 門徒不作聲，因為他們在路上彼此爭論誰最大。 \end{tabularx} \\ \\ \relax
9:35 & \begin{tabularx}{0.7\textwidth}{X} 耶穌坐下，叫十二個使徒來，說：「若有人願意為首，他要作眾人之後，作眾人的用人。」 \end{tabularx} \\ \\ \relax
9:36 & \begin{tabularx}{0.7\textwidth}{X} 於是耶穌領一個小孩過來，讓他站在門徒當中，又抱起他來，對他們說： \end{tabularx} \\ \\ \relax
9:37 & \begin{tabularx}{0.7\textwidth}{X} 「凡為我的名接納一個像這小孩子的，就是接納我；凡接納我的，不是接納我，而是接納那差我來的。」 \end{tabularx} \\ \\
[1ex]
\hline
\hline
\end{longtable}
主題:「你們在路上議論的是什麼?」　What Were You Discussing on The Way?

經文:馬可福音9:30-37

講員:鄧瑞強博士

場合:香港中文大學崇基學院禮拜堂

日期:2015年9月20日

各位弟兄姊妹,早安.

今日,在返教會的途中,若你身邊有伴,「你們在路上議論的是什麼?」

在人生的旅途上,你和最好的朋友,把酒談心的時候,你們議論的是什麼?

我們談論的事,多得很,天文地理,文學歷史,經濟政治,信仰理想.談論的內容,是我們人生用心之所在.想想,平時你多談論什麼?想想,你生命的真正用心,在什麼地方?

有一次,耶穌和門徒在途上.耶穌心中思念著生命的大事,門徒卻熱烈地談論別的大事.如此,引發了一場對話.

我們看看這故事,就是今日的講道經文:馬可福音9:30－37.

可9:30　他們離開那地方,經過加利利；耶穌不願意人知道.

可9:31　於是教訓門徒,說:「人子將要被交在人手裡,他們要殺害他；被殺以後,過三天他要復活.」

可9:32　門徒卻不明白這話,又不敢問他.

可9:33　他們來到迦百農,耶穌在屋裡問門徒說:「你們在路上議論的是什麼?」

可9:34　門徒不做聲,因為他們在路上彼此爭論誰為大.

可9:35　耶穌坐下,叫十二門徒來,說:「若有人願意做首先的,他必做眾人末後的,作眾人的用人.」

可9:36　於是領過一個小孩子來,叫他站在門徒中間,又抱起他來,對他們說:

可9:37　「凡為我名接待一個像這小孩子的,就是接待我；凡接待我的,不是接待我,乃是接待那差我來的.」

這個故事的核心,是耶穌問門徒一個問題:「你們在路上議論的是什麼?」

他們默不作聲,因為他們在路上彼此爭論的事情,連他們也知道,是遠離耶穌很遠的事情.

這是一群跟著耶穌的人,一群最接近耶穌的人,一群理應最明白耶穌使命的人,卻是一群距離耶穌很遠的人.

這是天地間的一大悲哀.

追隨真理的人,卻思想著騙人的謊言.在教會裡崇拜上帝的人,卻在人生的路上崇拜著世間的榮華富貴；信仰耶穌的基督徒,卻信仰著世間的權力.最接近神的人,卻是最遠離神.最接近天國的人,卻一步一步走向地獄.這是驚心動魄的.這是天地間的一大悲哀.

可9:35　耶穌坐下,叫十二門徒來,說:「若有人願意做首先的,他必做眾人末後的,作眾人的用人.」

門徒爭論誰為大,耶穌卻教導他們作別人的僕人.門徒的視線,在世間生存背後的權力結構；耶穌的視線,卻在神對世人生命的終極判斷.

門徒在路上的爭論,其實道出了世間運作的真相,一種關於權力的真相,一種權力無所不在的真相.他們明白世間的運作,從根本上說,就是一場權力的角力.他們沒有跳出權力運作的框架,只希望自己在這框架下得到最大的好處.他們爭論誰為大.

有一個意大利的哲學家葛蘭西(Antonio Gramsci),講「文化霸權」(cultural hegemony),講到權力其實不單在有形的體制中運作,其實也在文化的各個領域裡展示其力量.在各文化領域裡潛藏的霸權,建構我們各種想法.文化霸權建構出「什麼是有價值」,「什麼是好」,甚至建構出「什麼是真實」的想法.

寫作《世界是平的》(The World is Flat)這本出名的書的作者湯瑪斯．弗里曼(Thomas Friedman)同時寫了另一本沒那麼出名的書,書名叫《從貝魯特到耶路撒冷》(From Beirut to Jerusalem).書中講到,當以色列士兵射殺巴勒斯坦人時,他們是見不到有無辜的平民遭到殺害,也聽不見有天真的孩童發出尖叫.他們看見的,是「恐怖分子」遭到「清除」,「邪惡」得到消滅.從另一面講,當巴勒斯坦人襲擊以色列的無辜百姓時,他們也沒有什麼良心不安,他們看對方只是像狗一樣的「恐怖分子」,殺害他們是執行上天旨意.

他們每一個人,都只看到權力建構出來的所謂「真實」.對方到底是「什麼」,要視乎在權力角逐的遊戲中誰勝誰負.強權者可以定義對方是「什麼」,這便成為對方的「真相」.在國際舞台上,很多人爭論誰為大.

爭論誰為大,會扭曲人性.

有一齣電影,片名是《The Mission》(香港譯作《戰火浮生》),講到有幾個耶穌會的神父,去到南美洲位處巴西,阿根廷,巴拉圭交界的伊瓜蘇瀑布(Iguazu Falls)上游,向當地的土著傳福音.他們的佈道很成功,建立了一個基督化的社區,但這卻觸動了西班牙和葡萄牙這些殖民者的神經.因為在殖民者眼中,這些土著是一種商品.殖民者「捕捉」這些土著作奴隸,將他們拿去賣錢.在神父眼中,這些土著是神眼中寶貴的兒女.在殖民者眼中,這些土著是發財的工具.神父抵制著殖民者,於是,殖民者便竭力銷毀神父建立的事工.

教廷派了一位主教去處理這事.在主教面前,耶穌會的神父派了一位土著小孩,唱出美妙的讚美詩.之後,這位主教向在場的殖民者權力代表說,「你們怎麼可能認為這孩子只是一頭動物?」一位殖民者的權力代表回應說:「只要教鸚鵡唱歌,鸚鵡也會唱歌的.」主教說:「呀,是的,但怎能教鸚鵡唱得如此美妙?」那位殖民權力的代表說:「這孩子是森林裡長大的,是一頭懂人話的動物.…必須要用劍將他們制服,必須要用鞭使他們成為有用的勞動力.(This is a child of the jungle, an animal with a human voice….They will have to be subdued by the sword and brought to profitable labor by the whip.)」

殖民者是世上爭論誰為大的典範人物,他們的爭論會將一個人看成為「動物」,看成為奴隸,看成為可被利用的勞動力.

最後,殖民者進行大屠殺,摧毀整個土著社區.

「爭論誰為大」,是人類最關鍵,最核心的爭論.按法國哲學家福柯(Michel Foucault)的講法,「權力無所不在」.權力會製造出真理的體制,而真理的體制會為權力提供可用的知識.權力令它製造出來的真理看來可信；這真理系統又會令這權力的運用變得合法.如此,「爭論誰為大」,就不單只是爭坐上王位,也是爭取真理論述的主宰權.

「爭論誰為大」這框架下的真理論述,會二分這個世界.有權力者,是人.無權力者,是動物,是奴隸,是物.在「爭論誰為大」的框架裡,有權力者可以建構一種對失敗者的論述.這種論述涉及很多方面,而最基本的,可以有兩方面.首先,是「本體論」的論述(ontological discourse),即論述對方是「什麼」.其次,便是「倫理學」的論述,即論述如何對待對方.

巴勒斯坦裔學者薩依德(Edward Said)寫了一本書,書名是《東方主義》(Orientalism).書中講到,西方霸權以自己的利益觀點,建構出一種對「東方」的論述.在西方霸權的權力眼鏡下,在「本體論」層次而言,「東方」是非理性的,墮落的,幼稚的,低下的；相對而言,「西方」是理性的,道德的,成熟的,優越的.正因為這種「本體論」上的差異,西方人征服,管治,開發「東方」,在倫理上便是合理的,這是使「東方」進步的合理途徑.這本書展示了權力的邏輯:權力產生真理的論述,真理的論述展示一種「本體論」,「本體論」合理化一種「倫理學」.若對方被論述為只是一頭「動物」,則馴養牠,驅策牠,並利用牠,在倫理上便變成合法.

門徒在路上爭論誰為大,他們的爭論是人類根深蒂固的爭論,也是我們心底潛藏的爭論.這爭論涉及我們在真理上,在本體論上,在倫理學上,如何處置這個世界.耶穌沒有在他們的爭論框架下打轉,耶穌展示一個與權力世界完全不一樣的神聖世界,耶穌稱這世界為「天國」.

耶穌說:「若有人願意做首先的,他必做眾人末後的,作眾人的用人.」

有人以為,耶穌在這裡講的是爬上高位的策略.先扮懵,先做「細」,實質是扮豬食老虎,看準時機,一步登天.這種看法將耶穌看成是最成功的「世界仔」,深明如何「上位」的竅門.基督徒若有這種想法,則是跟著耶穌而卻是最遠離耶穌的人.

耶穌是要顛覆整個「爭論誰為大」的思想框架.「爭論誰為大」是以權力作為一切生活的基礎.擁有權力,便擁有產生真理論述的力量.這真理的論述,會在「本體論」層次將所有人的本性及價值作一劃分.劃分出:誰是優越的,有價值的；誰是低下的,可放棄的.然後,以這「本體論」觀點合理化倫理行動.對於被認定為優越的,我們親近,讚揚.對於那些被認定為低下的,我們遠離,放棄.

耶穌說:「若有人願意做首先的,他必做眾人末後的,作眾人的用人.」

「爭論誰為大」是迫一些無權力的人作「下人」,耶穌的世界卻是甘願成為別人的「下人」.背後的邏輯,不是權力結構,而是愛與關懷.為了愛,可以多走一步.為了愛,可以將座位讓給人.為了愛,可以將錢捐給人.為了愛,可以不是為自己而活,而是放下自己,承托他人.

這是一種新的真理論述.這真理論述讓我們看到神眼中看到的世界.在這世界裡,我們每個人都是僕人,都是神的僕人,被差派到世上,為了神的正義,去服事身邊的人.身邊的人不是「爭論誰為大」的對手,而是我們懷著愛去服事的對象.

一種新的對人的「本體論」的看法,便會產生一種新的倫理行動.

耶穌作出一個示範.

可9:36　於是領過一個小孩子來,叫他站在門徒中間,又抱起他來,對他們說:

可9:37　「凡為我名接待一個像這小孩子的,就是接待我；凡接待我的,不是接待我,乃是接待那差我來的.」

「小孩子」,在法律上,社會上,宗教上,都是無地位的.He is just a nobody. 他們無權主宰自己的命運,他們是屬於大人的.在「爭論誰為大」的框架裡,這些一無是處的人,正好是欺負的對象.耶穌卻抱起他們,沒有在權力爭奪的框架下擠壓他們.對於這些無權無勢的人,耶穌視他們為保護的對象,像母親懷抱兒女一樣.

接著,耶穌說:「凡為我名接待一個像這小孩子的,就是接待我；凡接待我的,不是接待我,乃是接待那差我來的.」

這是一種身分的認同.若這個人是一個弱者的話,則耶穌便同樣視自己為一個弱者.在「本體論」層次而言,耶穌沒有視自己為高於弱者的人,而是視自己為與弱者站在同一位置的人.接待他,就是接待我,因為我與他認同.這種真理的論述,會將抱起弱者的生命視為倫理上的無上命令.這是神的心意.

耶穌的說話和行動,是在一大背景下說和做的.我們回到經文的起首,看看這故事的背景.

可9:31　耶穌教訓門徒,說:「人子將要被交在人手裡,他們要殺害他；被殺以後,過三天他要復活.」

在路上,耶穌和門徒談論的,是他的死亡和復活.這位從天降世的神,認同了卑微的世人,變得一無所有.十架上的耶穌,在法律上,社會上,宗教上,都變得毫無地位.He is just a nobody. 在十架上,耶穌成了一個一無是處的小孩子.耶穌與小孩子認同,不是因為小孩子得意,可愛,純真,而是因為他們是無權無勢的代表,是自己走上十架時的模樣.

接待小孩子,就是接待十架上的耶穌.接待十架上的耶穌,一個像小孩子那樣無權無勢的耶穌,就是接待那差耶穌來的天父.

若果我們只顧「爭論誰為大」,便不可能接近十架上的耶穌.事實上,當耶穌被釘在十架上的時候,這些「爭論誰為大」的門徒是一哄而散的.

耶穌走上十架,他所作的,就如抱起小孩子.他以極端的愛,認同弱者,並扶起弱者.他以極端的愛,去對抗權力世界將弱者壓下去的做法.他以極端的愛,去對抗這個「爭論誰為大」的結構.

信耶穌,到底是什麼意思?

放下「爭論誰為大」的思維,顛覆權力結構,不以權力迫人做「下人」,而是因著愛而甘願成為服務別人的人.視一無所有的人為兄弟,接待他們.接待他們,就是接待與他們認同的耶穌.接待耶穌,就是活在天父的心意裡.活在天父的心意裡,我們便重生了.所謂「重生」,就是神聖的豐盛生命.

聽聽耶穌的問題:「你在路上議論的是什麼?」你會如何回答呢?

但願:榮耀歸於聖父,聖子,聖靈.阿們.



\newpage

\section{三個屬靈成長的秘訣}
\label{sec:298}
\textbf{蒲錦昌牧師 - 三個屬靈成長的秘訣}
\newline
\newline
\hyperref[sec:297]{< < < PREV SERMON < < <}
~
\hyperref[sec:toc]{[返目錄]}
~
\hyperref[sec:299]{> > > NEXT SERMON > > >}


提摩太後書 2:1-7
\newline
\begin{longtable}{cl}
\hline
\hline
章節 & 經文 (和合本修訂版)\\
\hline
2:1 & \begin{tabularx}{0.7\textwidth}{X} 我兒啊，你要在基督耶穌的恩典上剛強起來。 \end{tabularx} \\ \\ \relax
2:2 & \begin{tabularx}{0.7\textwidth}{X} 你在許多見證人面前聽見我所教導的，也要交託給那忠心而又能教導別人的人。 \end{tabularx} \\ \\ \relax
2:3 & \begin{tabularx}{0.7\textwidth}{X} 你要和我同受苦難，作基督耶穌的精兵。 \end{tabularx} \\ \\ \relax
2:4 & \begin{tabularx}{0.7\textwidth}{X} 凡當兵的，不讓世務纏身，好使那招他當兵的人喜悅。 \end{tabularx} \\ \\ \relax
2:5 & \begin{tabularx}{0.7\textwidth}{X} 運動員在比賽的時候，不按規則就不能得冠冕。 \end{tabularx} \\ \\ \relax
2:6 & \begin{tabularx}{0.7\textwidth}{X} 勤勞的農夫理當先得糧食。 \end{tabularx} \\ \\ \relax
2:7 & \begin{tabularx}{0.7\textwidth}{X} 我所說的話，你要考慮，因為主必在凡事上給你聰明。 \end{tabularx} \\ \\
[1ex]
\hline
\hline
\end{longtable}
講題:三個屬靈成長的秘訣　Three Secrets of Spiritual Growth

經文:提摩太後書2章1-7節

講員:蒲錦昌牧師

場合:香港中文大學崇基學院禮拜堂

日期:2015年9月27日

一, 前言:

各位弟兄姊妹平安!

我們世界的確變了,而且變了許多.

我們可以從一個二十歲年青人的角度,想像一下他/她的生活,大概就會了解世界變了多少.

像平常一樣,昨晚也是半夜一,兩點才睡,和朋友在網上閒聊,不知不覺時間就過去了,放下電話,倒頭便睡.今早起床,已是十時,幸好大學沒課,否則又要走堂了.先檢查一下手機的信息,才洗面刷牙.爸媽也上班了.趕著出門口,搭地鐵戴著耳機一面聽音樂,一面看手機熒幕上彈出來的新聞──唉,啟晴村食水含鉛事件越弄越大件事,竟然出動到法官成立專責調查委員會,政府又一次表現差勁!都差不多食午飯時間,what’sapp同學吃飯吹吹水,上網找個地方吃東西.啊,差點忘記今天要交功課,找個地方坐下,用i-pad整理好未完成的功課電郵給老師吧!噢!有同學what’sapp我今晚政府總部又有活動.去不去呢?

當然,以上的故事只是靠想像的香港青年人生活片段.但是,只要細心一想,我們便會發現,青年人今日所生活的世界和我們五十歲以上(三十年前是年青人)的很不同,資訊科技的發展正大幅度改變著我們的生活.

如果我們不嘗試走近年青人的世界,了解他們多一點,可能我們便像生活在同一空間但卻是兩個世界,彼此之間沒有交集.我們看報紙電視,他們看網上新聞；我們在熟悉的朋友和生活圈子中活動,他們相識遍天下,滿是陌生人；我們和他們傾偈,他們卻低頭做著另一些事情；我們是他們最親的親人,但是,他們卻視手機為最佳伴侶.我們認為他們是小孩,他們認為自己是大人.

信仰對21世紀資訊科技的年青一代具有什麼意義呢?能否成為他們人生的支柱呢?教會怎樣幫助他們成長呢?保羅二千年前對年青傳道人提摩太的勸勉,又是否可以跨越時空成為他們的幫助呢?

二, 主題:

保羅在寫提摩太後書時,已接近他事奉及人生的終點.他將盼望都寄託在年青一代身上,其中一個就是提摩太.

在提後2:1說:「我兒啊,你要在基督耶穌的恩典上剛強起來.」又在提後1:7說:「因為上帝賜給我們,不是膽怯的心,乃是剛強,仁愛,謹守的心.」保羅希望提摩太夠堅強,可以承擔福音使命的重任.

保羅還說(提後2:2):「你在許多見証人面前聽見我所教訓的,也要交託那忠心能教導別人的人.」他盼望靠著他訓練出來的提摩太,可以把福音的棒一代一代傳下去.

教會要傳承下去,不能不培養年青人,否則就沒有下一代的領袖.看看初期教會的歷史,不是每個年青人的成長都是一帆風順的,個性都是完美的,提摩太個性柔弱,馬可怕艱難.但是,教會要信任他們,給他們機會.

例:青少年主日幾個年青人上講台講道,每人五分鐘,都是有心事奉的年青人.

小朋友學行總會跌跌撞撞,跌下再爬起來,哭完再走.年青人也在成長中學習,不能只成功不犯錯.教會也要包容.

對於提摩太的成長,保羅給了他三個成長的錦囊.這些錦囊,不單青年人合用,其他年紀的其實也很值得參考.他以三個身份,工作,事業來描述這三個錦囊是什麼,就是軍人,運動員和農夫.

1. 軍人──委身(訓身),專注

提後2:3-4:「你要和我同受苦難,好像基督耶穌的精兵.凡在軍中當兵的,不將世務纏身,好叫那招他當兵的人喜悅.」

保羅時代的軍人,是有僱傭兵的,即專業軍人,打仗就是他們的職業,有點像私人屋苑的保安,保障住客安全,軍人則保護國家的安全.當一個人準備接受軍人的職務時,他第一個要做的就是放下其他的工作,事務,才能做個專業軍人.他不能一邊當兵,一邊耕田,又或者一邊做生意買賣.否則,那隊軍隊只是一班雜牌軍.

現在社會上常講multi-task,工作不會一生人做一個職業,不可能留在一間公司終老,隨時沒工做,要懂得多些技能才可轉行.年青人要周身刀,小學已一大堆証書.

資訊科技時代年年青人懂的東西很多,也會同時做幾件事.

例:看過講員一邊講一邊注視手機,試過崇拜手機響,不尊重上帝.不能停一下嗎?真的那麼重要嗎?你是醫生24小時on call等著救人一命嗎?

的士有六,七部手機在前面,很擔心自身的平安.

保羅第一個錦囊,是要我們學專注,分得清楚什麼最重要,什麼是次要,什麼是絲毫不重要的.一個人其實人生有限,太多群組太多信息根本處理不到,要學專注.

軍人當兵,是叫招他當兵的喜悅,基督徒在世生活,是以上帝的喜悅作最終標準.我們也以此為目標嗎?也以此衡量自己所做一切嗎?目標清晰的話,一切就豁然開朗,否則,在大海飄浮,始終沒有方向,不知要去那兒.

2. 運動員──紀律,鍛鍊

提後2:5 「人若在場上比武,非按規矩,就不能得冠冕.」

提後2:5(和合本修訂版):「運動員在比賽的時候,不按規則就不能得冠冕.」

古希臘的奧林匹克競技和比賽,就有非常之嚴格的訓練,包括食物的控制和體能訓練.現在奧運的運動員,不單只有營養師照顧飲食,有物理治療師,職業治療師協助身體不同部分的鍛鍊,心理學家調控比賽前的心理狀態,還要嚴格控制賽前所用過的藥物,以避免違反賽事的規定.十多二十年訓練,只為了一,兩分鐘的比賽.「台上一分鐘,台下十年功.」真不簡單.在準備期間,運動員許多時候要抵受一日復一日枯燥無味的訓練,家人的分離,強度大的體能運動量,疼痛的身體,不合口味的食物,受傷的危險等.要成功,要得獎,這一切都得忍受.

一個人要做好一件事,也要這種決心和紀律,也要鍛鍊自己應付人生的賽事.

例:全完堂周刊文章的經驗.要鍛鍊,要紀律自己.

資本主義社會教導人自我滿足,以自己的欲望和喜好選擇和決定,身處這樣的世界,更需要紀律來限制自己,懷抱遠大目標,鍛鍊自己,全力以赴,才會有成功一天.

3. 農夫──勤懇

提後2:6 「勞力的農夫,理堂先得糧食.」

農夫在人類歷史的不同時代,都是艱辛的工作.翻土,播種,除草,除蟲,照顧植物度過不同的天氣,在炎熱的,寒冷的季節下田,一點也不好受.凡是經歷過這些階段的,也會懂得珍惜食物.有不少家長陪小朋友到新界元朗,錦田租一片小小的地來耕種,也是希望他們懂得珍惜所有的豐足.

保羅講的道理很簡單,肯堅持下去,肯付出汗水的農夫,才能得到收穫,才    有機會先嘗到收成.

香港社會素來不講一分耕耘,一分收穫,而講走捷徑.資訊科技世界不講長時間的努力,而講即時的反應,短時期的點擊率和點讚(like).

葛拉威爾:《異數:超凡與平凡的界線在哪裡?》,談披頭四(Beatles)及比爾‧蓋茨(Bill Gates)的成功,有黃金一萬小時之說.不是靠天才,而是努力達至成功.

保羅繼續說:「我所說的話你要思想,因為凡事主必給你聰明.」我們要對社會,世界的潮流以信仰的角度去反思,對別人的想法用自己的腦袋去思考,才不至於隨波逐流,失去基督徒見證的能力.

三, 結論:

做基督徒不是一件容易的事,要承擔傳福音,牧養世界的重任更加不是.盼望我們,特別是我們中間的年青人,都有軍人一樣的專注,運動員一樣的紀律與鍛鍊,農夫一樣的勤懇,使我們在屬靈生命上不斷成長,在人生中有所成就,貢獻社會,榮耀上帝,見証上帝的美名,阿們!



\newpage

\section{上帝與金錢是誰在掌權}
\label{sec:299}
\textbf{許開明牧師 - 上帝與金錢,是誰在掌權?}
\newline
\newline
\hyperref[sec:298]{< < < PREV SERMON < < <}
~
\hyperref[sec:toc]{[返目錄]}
~
\hyperref[sec:300]{> > > NEXT SERMON > > >}


馬可福音 10:17-37
\newline
\begin{longtable}{cl}
\hline
\hline
章節 & 經文 (和合本修訂版)\\
\hline
10:17 & \begin{tabularx}{0.7\textwidth}{X} 耶穌剛上路的時候，有一個人跑來，跪在他面前，問他：「善良的老師，我該做甚麼事才能承受永生？」 \end{tabularx} \\ \\ \relax
10:18 & \begin{tabularx}{0.7\textwidth}{X} 耶穌對他說：「你為甚麼稱我是善良的？除了神一位之外，再沒有善良的。 \end{tabularx} \\ \\ \relax
10:19 & \begin{tabularx}{0.7\textwidth}{X} 誡命你是知道的：『不可殺人；不可姦淫；不可偷盜；不可作假見證；不可虧負人；當孝敬父母。』」 \end{tabularx} \\ \\ \relax
10:20 & \begin{tabularx}{0.7\textwidth}{X} 他對耶穌說：「老師，這一切我從小都遵守了。」 \end{tabularx} \\ \\ \relax
10:21 & \begin{tabularx}{0.7\textwidth}{X} 耶穌看著他，就愛他，對他說：「你還缺少一件：去變賣你所有的，分給窮人，就必有財寶在天上；然後來跟從我。」 \end{tabularx} \\ \\ \relax
10:22 & \begin{tabularx}{0.7\textwidth}{X} 他聽見這話，臉就變了色，憂憂愁愁地走了，因為他的產業很多。 \end{tabularx} \\ \\ \relax
10:23 & \begin{tabularx}{0.7\textwidth}{X} 耶穌看了看周圍，對門徒說：「有錢財的人進神的國是何等的難哪！」 \end{tabularx} \\ \\ \relax
10:24 & \begin{tabularx}{0.7\textwidth}{X} 門徒對他的話非常驚奇。耶穌又對他們說：「孩子們，要進神的國是何等的難哪！ \end{tabularx} \\ \\ \relax
10:25 & \begin{tabularx}{0.7\textwidth}{X} 駱駝穿過針眼比財主進神的國還容易呢！」 \end{tabularx} \\ \\ \relax
10:26 & \begin{tabularx}{0.7\textwidth}{X} 門徒就更為驚訝，彼此對問：「這樣，誰能得救呢？」 \end{tabularx} \\ \\ \relax
10:27 & \begin{tabularx}{0.7\textwidth}{X} 耶穌看著他們，說：「在人不能，在神卻不然，因為在神凡事都能。」 \end{tabularx} \\ \\ \relax
10:28 & \begin{tabularx}{0.7\textwidth}{X} 彼得就對他說：「看哪，我們已經撇下一切跟從你了。」 \end{tabularx} \\ \\ \relax
10:29 & \begin{tabularx}{0.7\textwidth}{X} 耶穌說：「我實在告訴你們，凡為我和福音撇下房屋，或是兄弟、姊妹、父親、母親、兒女、田地， \end{tabularx} \\ \\ \relax
10:30 & \begin{tabularx}{0.7\textwidth}{X} 沒有不在今世得百倍的，就是房屋、兄弟、姊妹、母親、兒女、田地，並且要受迫害，在來世得永生。 \end{tabularx} \\ \\ \relax
10:31 & \begin{tabularx}{0.7\textwidth}{X} 然而，有許多在前的，將要在後；在後的，將要在前。」 \end{tabularx} \\ \\ \relax
10:32 & \begin{tabularx}{0.7\textwidth}{X} 他們行路上耶路撒冷去。耶穌在前頭走，他們很驚訝，跟從的人也害怕。耶穌又叫十二使徒來，把自己將要遭遇的事告訴他們， \end{tabularx} \\ \\ \relax
10:33 & \begin{tabularx}{0.7\textwidth}{X} 說：「看哪，我們上耶路撒冷去，人子將被交給祭司長和文士；他們要定他死罪，又交給外邦人。 \end{tabularx} \\ \\ \relax
10:34 & \begin{tabularx}{0.7\textwidth}{X} 他們要戲弄他，向他吐唾沫，鞭打他，殺害他；三天後，他要復活。」 \end{tabularx} \\ \\ \relax
10:35 & \begin{tabularx}{0.7\textwidth}{X} 西庇太的兒子雅各和約翰進前來，對耶穌說：「老師，我們無論求你甚麼，願你為我們做。」 \end{tabularx} \\ \\ \relax
10:36 & \begin{tabularx}{0.7\textwidth}{X} 耶穌對他們說：「要我為你們做甚麼？」 \end{tabularx} \\ \\ \relax
10:37 & \begin{tabularx}{0.7\textwidth}{X} 他們對他說：「在你的榮耀裡，請賜我們一個坐在你右邊，一個坐在你左邊。」 \end{tabularx} \\ \\
[1ex]
\hline
\hline
\end{longtable}
講題:上帝與金錢,是誰在掌權?　God or Money, Who is the Boss?

講員:許開明牧師

經文:馬可福音10:17-37

場合:香港中文大學崇基學院禮拜堂

日期:2015年10月11日

今天,我們會看馬可福音第10章第17節至第27節,因為今天我想分享的講題就是「上帝與金錢,是誰在掌權?」(God or Money, Who is the Boss?)你是否已有答案?或許你估到我所要講的答案,但最重要的,是你的答案!

經文簡述

先讓我們看看馬可福音10:17-37.

第一幕:有一個人跑來找耶穌

第十章第十七節:「耶穌出來行路的時候,有一個人跑來,跪在他面前,問他說:『良善的夫子,我當做甚麼事才可以承受永生?』」這個「人」,是一個年青人的官,又是財主,他「跑來」,為何他如此匆忙?是否一時衝動?然後他在主耶穌面前「跪下」,這麼有身份的人,竟跪在一位中年人,沒有甚麼前途；沒有社會地位,不受猶太上層,羅馬權貴所尊重的人物；且是窮困的拿撒勒人耶穌.第十八節:「耶穌對他說:『你為甚麼稱我是良善的? 除了神一位之外,再沒有良善的.」耶穌暗示你想做到「良善」不是容易的事!

第二幕:詢問那人

從耶穌的問題,我們可以了解這人達到多少良善呢?第十九節:「誡命你是曉得的:不可殺人,不可姦淫,不可偷盜,不可作假見證,不可虧負人,當孝敬父母.你到底達到多少?」在這誡命中的七條,你做到嗎?第二十節:「他對耶穌說:『夫子,這一切我從小都遵守了.』」

第三幕:耶穌指出:他還缺少了一件事

第二十一節:「耶穌看著他,就愛他,對他說,你還缺少一件．去變賣你所有的,分給窮人,就必有財寶在天上．你還要來跟從我.」耶穌的要求,可能對他來說,太高太難了!還要是一個比一個難!第一難是變賣所有；第二難是把財產分給窮人；第三難是跟從耶穌.但那人的屬靈洞察力並未配上預備委身(包括犧牲在內)的態度,所以他憂憂愁愁的離去了.耶穌的挑戰就是:你還有一件事,你還欠缺這件事!

而面對這個重大的挑戰,你會選擇做財富的主人嗎?如果你做到這三件事:即是去變賣你所有的,分給窮人,就必有財寶在天上,你 還要來跟從耶穌.你會選擇「耶穌做主人」,還是「財富做你的主人」?而那個人就認為財富是他的一切!他不變賣,不分給窮人,不跟從主!第二十二節:「他聽見這話,臉上就變了色,憂憂愁愁的走了．因為他的產業很多.」根據原文字義,「變了色」同時解作「陰了天,發黑」的意思.這位年輕人最終憂憂愁愁的走了.這個令耶穌失望的反應,很清楚地說明了他還沒有一個正確的觀念.因為他不願意變賣那沒有永恆價值的產業,去換取上帝所賜予人在天上永存的財產--永生.這也是整個事件中最終要帶出來的信息和挑戰.

第四幕:耶穌指出那人的問題所在

第二十三節:「耶穌周圍一看,對門徒說,有錢財的人進神的國是何等的難哪.」第二十四節:「門徒希奇他的話.耶穌又對他們說:『小子,倚靠錢財的人進　神的國,是何等的難哪.」為甚麼耶穌的話讓門徒都希奇呢?

因為有錢的人,是蒙福的人,財富是良善人的報應!詩篇第三十七章第二十五節說:「義人不討飯」為甚麼進上帝國卻會有困難?當時猶太人普通的道德觀念,認為人的享福,純粹是一個良善的人應有的報應,所以財富是美善的人格和上帝賜福的象徵.但詩篇第三十七章的詩人一言以蔽之,第二十五節是這樣寫的:「我從前年幼,現在年老,卻未見過義人被棄,也未見過他的後裔討飯.」這樣就難怪門徒當時如此的驚奇了!他們一向以為人財富越多,進天國的機會越高.於是耶穌第二次稍為改變祂的語調,重申祂的理論「耶穌又對他們說:『小子,倚靠錢財的人進神的國,是何等的難哪.」亦即是說「人若把自己的信靠,放在財富上,是很難進上帝的國的.」第二十五至二十七節:「駱駝穿過鍼的眼,比財主進神的國,還容易呢.門徒就分外希奇,對他說:『這樣誰能得救呢.耶穌看著他們說,在人是不能,在神卻不然,因為神凡事都能.」

第五幕:門徒彼得做得到!

第二十八節至三十二節這樣說:「彼得就對他說,看哪,我們已經撇下所有的跟從你了.耶穌說,我實在告訴你們,人為我和福音,撇下房屋,或是弟兄,姐妹,父母,兒女,田地.沒有不在今世得百倍的,就是房屋,弟兄,姐妹,母親,兒女,田地,並且要受逼迫．在來世必得永生.然而有許多在前的將要在後,在後的,將要在前.他們行路上耶路撒冷去．耶穌在前頭走,門徒就希奇,跟從的人也害怕．耶穌又叫過十二個門徒來,把自己將要遭遇的事,告訴他們說.」

猶太人的金錢觀

讀畢這段經文,就讓我們來看看猶太人的金錢觀.猶太人是很會賺錢的民族,他們對錢有正面的態度:聖書給人光明,金錢給人溫暖；金錢是提供機會的美物.他們相信完全沒有財富或金錢的人,憂愁更多.錢是貨幣,也是一個人擁有財富的標記,所以它具有價值.他們還認為對金錢不感興趣的人賺不到錢；然而把金錢看得太重,金錢就變成包袱.對他們而言,世界有三大的事情,就是金錢,金錢和金錢.

耶穌的金錢觀

不過,耶穌的金錢觀又是怎樣的呢?

財富是短暫的,不能永存

第一,耶穌認為:財富是短暫的,不能永存.馬太福音第六章第十九至二十節這樣說:「不要為自己積儹財寶在地上,地上有蟲子咬,能鏽壞,也有賊挖窟窿來偷．只要積儹財寶在天上,天上沒有蟲子咬,不能鏽壞,也沒有賊挖窟窿來偷．」

錢會受損

由這段經文,我們可以看到錢的四個特點.第一個特點是:錢會受損.在耶穌時代,人們的財富常由身上的衣著顯示.最昂貴的是羊毛織的衣服,而蟲子最喜歡咬這類衣服,在當時人們是不知如何來防止蟲子對他們財物上的破壞.耶穌接著指出地上的財寶也會鏽壞,「鏽壞」會讓我們想到當時的錢幣所用的金屬是會被氧化生鏽的,因而遭到破壞.但這「鏽壞」兩個字在新約的其他提到的地方都和食物和吃有關,也許這裡是指當時的財富也由所貯藏的穀物表示,耶穌指出這地上的財寶是一樣地會被老鼠或各種蟲子吃光.

錢會被人偷

錢的第二個特點是錢會被人偷.在耶穌時代,盜賊會挖窟窿來偷,他們總是挖穿當時以泥土做的牆來偷其他人家中的財物,其實今天照樣存在這樣的小偷強盜,只是今天這個世代還有一些更聰明的人,藉著他們在電腦上的知識和技術,盜用冒充你的名字,把你在銀行的錢輕輕鬆鬆地就轉走了.

錢不能被帶進墳墓

錢的第三個特點就是錢不能被帶進墳墓.

錢會飛走

而最後一個錢的特點就是錢會飛走,箴言第二十三章第五節這樣說:「你豈要定睛在虛無的錢財上嗎?因錢財必長翅膀,如鷹向天飛去.」巴爾扎克也這樣講過:「對於浪費的人,金錢是圓的,可是對於節儉的人,金錢是扁平的,是可以一塊塊堆積起來的.」而在史祈生寫的《五十靈筵》中記錄過:「一九五三年美國科學促進會年會上,波士頓的一位心理學醫生朗誦論文說:『文明時代的金錢病就是失眠,胃潰瘍,頭痛,神經失常.』這些病不是財富多少影響的,而是對金錢的概念影響的.」

所以有錢財的人,難進上帝的國

所以有錢財的人,難進上帝的國,相反,無錢財的人易進上帝的國.馬可福音第十章第二十五節裡提過「駱駝穿過鍼的眼」,這是古時猶太人的城門,往往是在正常的門上另開一個小門,此小門名叫「鍼眼門」.白天時大開城門,供人們和貨物進出；每到黃昏,便關閉城門,而僅打開那扇小鍼眼門.那門小到只可容「人」進出,故若遇有馱貨物的駱駝時,須先把貨物從牠身上卸下來,然後叫駱駝屈著身子匍匐,再加上人力的推和拉,才能通過.所以「駱駝穿過鍼的眼」遂成了一句俗諺,以形容事情非常困難,但也並不是完全不可能之意.耶穌用這來形容年青,富有的官,他的財物太多了,他不能卸下他的所有,好像馱負貨物的駱駝一樣穿不過鍼的眼!

富裕有如肥胖

其實富裕有如肥胖,上世紀初有一個肥胖的演員,叫利奧‧羅斯頓.他體重三百八十五磅,腰圍六‧二英尺.一九三六年他在英國演出時,因心肌衰竭被送進湯普遜醫院急救,搶救人員用了當時最先進的設備.仍沒法挽回他的生命.臨終時,羅斯頓喃喃自語:你的身軀很龐大,但你的生命需要的只是一顆心臟.

羅斯頓這句話,觸動了湯普遜醫院的院長,他讓人把羅斯頓這句遺言刻在醫院的大樓上.一九八三年,一位叫默爾的美國石油大亨住進了湯普遜醫院.他因為石油生意,往返歐,亞,美之間,終因勞累而出現了心肌衰竭,在英國時住進這間醫院.他在醫院包了一層樓,增設五部電話和兩部傳真機,遙控他的業務.當時有報紙形容這裡是「湯普遜──美洲的石油中心.」

默爾的心臟手術成功,在這裡住了五個月就出院了.不過他迅即結束了他的大生意,在蘇格蘭鄉村買下一棟別墅住了下來.一九八八年,湯普遜醫院百年慶典,邀請默爾參加,記者問他為甚麼賣掉自己的公司,他指出了醫院大樓上的那一行字,說:利奧‧羅斯頓.

記者有點不明所以,媒體對他也沒有報道.後來,有一個人在默爾的傳記中發現了這麼一句話:富裕和肥胖沒有兩樣,也不過是擁有超過自己需要的東西吧了.

這就是他賣掉公司的答案.你的身軀很龐大,但你的生活需要的僅僅是一顆心臟.你的財富很豐厚,但你所需要的僅僅是你生活中最核心的部份而已.最核心部份,是最基本的衣食住行.其他部份,包括銀行中的多幾個零的數字,包括名車,名牌,豪宅的炫耀,實際上和人身體的肥胖沒有兩樣,都是生命中多餘的東西.不,甚至是負累.

羅斯頓一句話,讓默爾看破了人生,從而作出拋棄在名利場追逐而享受生命的選擇.追求更好的人生.

由這個故事中,我們可以看到,富裕和肥胖沒有兩樣,都是擁有超過自己需要的東西.人們都知道肥胖是多餘的,卻疏忽了豐富的財富也是多餘的.我們所需要的財富,僅僅是生活中最核心的部分而已.

錢是危險的

耶穌還認為,錢是危險的.耶穌對金錢似乎較負面,他認為錢是危險的,因為錢會阻礙人進入上帝的國.錢亦容易成為上帝,成為人生命中的主人.錢也容易成為人的偶像.因為人容易眼中只有錢,見錢開眼!更可能會看錢的份上,就由錢話事,人易被錢牽著走.從來沒有人比耶穌更清楚的指出,物質和財富所帶來的危險,它們到底有些甚麼危險呢?

第一,物質財富容易引導人的心專注這個物質的世界,因而把大部份精力和興趣集中在它上面,於是很難看見物質之外的世界了.人最感困難的一件事就是與財富分手.莊信博士講過一個故事,有一次他被邀請到一個十分華麗的城堡和每一處都很美麗的花圃參觀,當他遊覽完畢之後,便轉向他的朋友說:『人擁有這些華麗堂皇的地方,實在很難心甘情願地死去!』所以人有太多財富的話,自然會把自己的思想和興趣放在現存的物質世界上面.

第二,倘若人的唯一興趣是放在物質的事情上面,他的思想便容易以價錢去衡量一切事物了.他的思想就是凡事都計算用多少錢去買,而忘記了在這個世界裡,有些東西是用錢買不到的,有些東西是沒有標價的,而世界上真正貴重的東西是不能用錢去買的.人若開始習慣思想:凡事可以用錢去買時,這就是的生活的致命傷了.

「這東西值多少錢?」,「這人帶的手袋值多少錢?」,「這人每月入多少錢?」等.當國家或社會,當權者以經濟掛帥,非以民生及民主為先,容讓富者愈富,貧者愈貧時,社會會變成怎樣呢?當一間公司老板只會向錢看,不顧員工福利,剝削薪酬,這老板變成金錢的奴隸,這公司會變成怎樣呢?當教會推動事工時,只懂看錢份上,無錢不行,這教會又會變成怎樣呢?普希金曾說過:「世界的設計創造應以人為中心,而不是以謀取金錢,人並非以金錢為對象而生活,人的對象往往是人.」其實,上帝在創造世界時也是以萬物之靈--人類,為中心!

富翁之死

接下來,我想分享一個富翁的故事.

有個富翁在急流中翻了船,爬到溪間的石頭上大喊救命.一個年輕人奮不顧身地盪舟去救,但是由於山洪下瀉而漸張的湍流,使他的船進行非常緩慢.「快呀!」富翁高喊:「如果你救了我,我送你一千塊!」船仍然移動緩慢.「用力划啊!如果你划到,我給你兩千塊!」青年奮力地划著,但是既要向前,又要抗拒水流的阻力,船速仍然難以加快!「水在漲,你用力呀!」富翁嘶聲喊著:「我給你五千塊!」說時洪流已經快淹到他站立的地方.青年的船緩緩靠近,但仍然嫌慢.「我給你一萬塊,拚命用力呀!」富翁的腳已經淹在水中了.但是船速反倒愈慢了.「我給你五萬…….」富翁的話還沒說完,已經被一個大浪打下岩石,轉眼捲入洪流,失去了蹤影.

青年頹喪地回到岸上,蒙頭痛哭:「我當初只想到救他一命,但是他卻說要送我錢,而且一次又一次地增加.我心想,只要划慢一點點,就可能多幾萬塊的收入,哪裡知道,就因為慢了這麼一下,使他被水沖走,是我害死了他啊!(也失去了錢!)」青年捶著頭:「但是當我心裡只有義,而沒想到利的時候,他為什麼要說給我錢呢?」

從上可知:人不小心,會失去人格,很易墮入錢財的陷阱,害人害己.

金錢 vs 永生

在金錢與永生的兩者中,我們要知道愛神愛財是不可共存的.馬太福音第六章第二十四節說:「一個人不能事奉兩個主．不是惡這個愛那個,就是重這個輕那個．你們不能又事奉神,又事奉瑪門.〔瑪門:財利的意思〕」

讓我先講講「金錢」,金錢古時的作用是用於量度貨物的價值.金錢可以購買貨物,金錢越多,購買的貨物越多,人擁有貨物越多,便愈富有.得到金錢,可使人的基本生命的需要得到解決,得到安全感；繼續擁有財富,使人可以享受生命,因而有快感.當人擁有更多的金錢,更多的財物,他可能也擁有更多的空間地方,土地,物業,勞工.並且擁有權力,名聲,更多的財利.

我們再看看「永生」,永生是屬於上帝的,是代表上帝,因為只有上帝擁有永生,永遠的生命,永恆的存在.當中還有更深的意思,就是永恆的價值,永恆的意義,永恆的真理和永恆的良善.上帝是耶和華,耶和華又譯作永恆主,具永恆的真實生命.

讓我們再回歸今天讀過的馬可福音第十章第十七至二十七節裡面所提到的「這人」,「這人」是一個年青,有錢的官；他擁有年青有力的生命；他富有:有不少田地；他是位官,擁有權力.他還有一樣更好的東西,就是他擁有理想,要追求良善.這是個很好的,很有潛質的,很有為的青年,可惜在面對耶穌的挑戰時,在金錢與上帝,他選擇了金錢!財富與永生之間,他選擇了財富!所以結果就是,他憂憂愁愁的走了!面對在金錢與上帝,財富與永生:我們也容易像這位有財有勢的年青人,作出一樣的選擇!

在這個故事中,它提示我們常常面對的幾個掙扎.

「理想」與「現實」

第一個掙扎就是「理想」與「現實」,現實就是「金錢或財富」,面理想就是「永生」,這人得到金錢或財富了,還想得到永生!魚與熊掌均想兩者兼得!事實上,這是十分的困難!人生常要在兩者中選取其一!人若把「金錢或財富」 變成「終極理想」,那就糟了!「金錢或財富」 就在人的生命中掌權!

「現今短暫」與「永恆價值」

第二個掙扎就是「現今短暫」與「永恆價值」.這人追求的金錢,「現今短暫」是可見的,抓得住的；而「永恆價值」是看不見的,抓不到的,但是有價值的.耶穌挑戰這人:追求的永生.我們要積財在地還是積財在天呢?地財會變壞,但天財不會變壞,並可以存到永恆!

「達半」與「達成,完整」

第三個掙扎是「達半」與「達成,完整」.這人做到十誡中大部份,守住六誡,但並不是全部,他有一種不完全的感覺.耶穌挑戰這人:要他達到完全!有人的生命,求求奇奇,不求完全!

「無做壞事」與「去做好事」

第四個掙扎是「無做壞事」與「去做好事」.這人滿足於無做違規的,不該的,傷人的事!他沒有違反律法.他實在說:「我一生未嘗傷害過別人.」他講的是鐵一般的事實.耶穌卻挑戰他,「你究竟做了那件好事呢?」繼續的一句話好像是要挖苦他--「你有這麼多的財產,這麼富有,你對別人做了那些善事呢?」我想講一個「人見天使」的故事:有位信徒到天堂門口,正想進入.

守門天使不淮入!信徒問:「為什麼?」天使說:「因為天堂沒有你名字!」信徒說:「我從本無做壞事!」天使說:「但,你從來也無做過好事!天堂是給做好事的人進入的!」各位,不少人常常滿足於無做錯事,壞事,或傷人的事,卻忽略了耶穌要人「愛你的鄰舍」,去做善事,而且是盡所能的去做善事!

「獲得」與「捨去」

第五個掙扎是「獲得」與「捨去」.財主已擁有地上的財富了,他突然想到也應擁有天上的財富!這人追求「獲得金錢」,自私的擁有,耶穌就挑戰這人:要「捨去」金錢 ,與人分享的,使窮人有出路!所以耶穌說:「你還缺少一件」,第一是「去變賣你所有的」,表示他的「所有」,要他去「變賣」所有的,轉化成金錢和不愛錢的人.第二是「分給窮人」,表現出馬太福音第十九章第二十節說的「愛人如己」,而「就必有財寶在天上」表示他現在所有的,不過是地上的,他在『天上』仍舊一無所有.第三是「你還要來跟從我」,表示他即使將所有的分給人,在他就好像馬太福音第十四章第7節裡所說,他的心裡仍可能只有窮人而沒有主,所以他還要來「跟從主」,亦即要像馬太福音第十章第三十七至三十八節:要愛主過於愛一切,連生命獻上!這樣,他才可以變成錢財的主人了!

因此,人要懂得用錢,做錢的主人.卡萊爾說過:「金錢(有時萬能)能做很多事,但它(有時無能)不能做一切事.我們應該知道它的領域,並且把它限制在那裡；當它想進一步發展時,甚至要把它們踢回去.」

舉一個例子,就是我們想買書的時候,我們要自己問三個原因,第一是「為甚麼花錢」；第二是「有沒有必要花錢?」要是可買可不買,可花可不花的錢就盡量節省下來.第三是「有沒有花錢的意義和價值?」我們要學會賺錢,學會控制錢.約翰衛斯理說過:「盡你所能去賺取,盡你所能去儲蓄,盡你所能去施捨.」學會控制錢後,把錢用在施捨,聖工,捐獻上,積財寶在天了!

結語

培根曾說:「金錢是好的僕人,卻是壞的主人.」馬克·吐溫也曾說:「如果你懂得使用,金錢是一個好奴僕,如果你不懂得使用,它就變成你的主人.」弗爾丁也說:「如果你把金錢當成上帝,它便會像魔鬼一樣地折磨你. 」上帝與金錢,是誰在掌權?God or Money, Who is the Boss?

這引伸出你生命常常面對幾個重要的選擇,財富與永生:第一是「現實與理想」；第二是「短暫與永恆」；第三是「達半與達成」；第四是「無做壞事與去做善事」；第五是「獲得與捨去」.我們的終極(Ultimate Concern),到底是積財於地還是積財於天?

昔日耶穌問那年青富有的官:「財富與永生,你的選擇是甚麼?」他沒有回答,卻憂憂愁愁的(黑著面,滿面愁煩)走了!從前大部份人是憂愁地來見耶穌,但開開心心的離去!今日,如果耶穌問你:

「財富與永生?」你的選擇是甚麼呢?「上帝與金錢?Who is your Boss?」你的回答是甚麼?

結束前,我們一起祈禱.



\newpage

\section{前來與被召}
\label{sec:300}
\textbf{林豪恩先生 - 前來與被召}
\newline
\newline
\hyperref[sec:299]{< < < PREV SERMON < < <}
~
\hyperref[sec:toc]{[返目錄]}
~
\hyperref[sec:301]{> > > NEXT SERMON > > >}


馬可福音 10:35-45
\newline
\begin{longtable}{cl}
\hline
\hline
章節 & 經文 (和合本修訂版)\\
\hline
10:35 & \begin{tabularx}{0.7\textwidth}{X} 西庇太的兒子雅各和約翰進前來，對耶穌說：「老師，我們無論求你甚麼，願你為我們做。」 \end{tabularx} \\ \\ \relax
10:36 & \begin{tabularx}{0.7\textwidth}{X} 耶穌對他們說：「要我為你們做甚麼？」 \end{tabularx} \\ \\ \relax
10:37 & \begin{tabularx}{0.7\textwidth}{X} 他們對他說：「在你的榮耀裡，請賜我們一個坐在你右邊，一個坐在你左邊。」 \end{tabularx} \\ \\ \relax
10:38 & \begin{tabularx}{0.7\textwidth}{X} 耶穌對他們說：「你們不知道所求的是甚麼。我所喝的杯，你們能喝嗎？我所受的洗，你們能受嗎？」 \end{tabularx} \\ \\ \relax
10:39 & \begin{tabularx}{0.7\textwidth}{X} 他們對他說：「我們能。」耶穌對他們說：「我所喝的杯，你們要喝；我所受的洗，你們也要受。 \end{tabularx} \\ \\ \relax
10:40 & \begin{tabularx}{0.7\textwidth}{X} 可是坐在我的左右，不是我可以賜的，而是為誰預備就賜給誰。」 \end{tabularx} \\ \\ \relax
10:41 & \begin{tabularx}{0.7\textwidth}{X} 其餘十個門徒聽見，就對雅各和約翰很生氣。 \end{tabularx} \\ \\ \relax
10:42 & \begin{tabularx}{0.7\textwidth}{X} 耶穌叫了他們來，對他們說：「你們知道，外邦人有君王作主治理他們，有大臣操權管轄他們。 \end{tabularx} \\ \\ \relax
10:43 & \begin{tabularx}{0.7\textwidth}{X} 但是在你們中間，不可這樣。你們中間誰願為大，就要作你們的用人； \end{tabularx} \\ \\ \relax
10:44 & \begin{tabularx}{0.7\textwidth}{X} 在你們中間誰願為首，就要作眾人的僕人。 \end{tabularx} \\ \\ \relax
10:45 & \begin{tabularx}{0.7\textwidth}{X} 因為人子來，並不是要受人的服事，乃是要服事人，並且要捨命作多人的贖價。」 \end{tabularx} \\ \\
[1ex]
\hline
\hline
\end{longtable}
講題:前來與被召　Come Forward and Being Called

經文:馬可福音10章35-45節

講員:林豪恩先生

場合:香港中文大學崇基學院禮拜堂

日期:2015年10月18日

引言

馬可福音第九章記載,耶穌主動向門徒揭示自己的遭遇,告訴他們說:「人子將要被交在人手裡,他們要殺害他；被殺以後,三天後他要復活.」(9:31)門徒聽到這消息後,他們之間就出現了爭論．門徒的這場爭論,引起了耶穌的注意,因此他問門徒:「你們在路上議論的是甚麼?」(9:33)門徒似乎知道他們所爭論的並不是甚麼光明正大的事,因此不敢告訴耶穌,但作者告訴我們,他們在路上"彼此爭論誰最大."(9:34)於是,耶穌勸導他們說:「若有人願意為首,他要作眾人之後,作眾人的用人.」(9:35)

令人驚訝的是,轉過頭來,當耶穌再次向門徒說:「看哪,我們上耶路撒冷去,人子將被交給祭司長和文士；他們要定他死罪,又交給外邦人.他們要戲弄他,向他吐唾沫,鞭打他,殺害他；三天後,他要復活.」(10:33-34)門徒再次聽到這消息後,竟然又引起類似的反應．門徒今次如何反應呢?耶穌又如何回應呢?讓我們細閱今天誦讀的福音經課──馬可福音10:35-45節．

經文

10:35西庇太的兒子雅各,約翰進前來,對耶穌說:「夫子,我們無論求你甚麼,願你給我們做.」

10:36耶穌說:「要我給你們做甚麼?」

10:37他們說:「賜我們在你的榮耀裡,一個坐在你右邊,一個坐在你左邊.」

10:38耶穌說:「你們不知道所求的是甚麼.我所喝的杯,你們能喝嗎?我所受的洗,你們能受嗎?」

10:39他們說:「我們能.」耶穌說:「我所喝的杯,你們也要喝；我所受的洗,你們也要受；

10:40只是坐在我的左右,不是我可以賜的,乃是為誰預備的,就賜給誰.」

10:41那十個門徒聽見,就惱怒雅各,約翰.

10:42耶穌叫他們來,對他們說:「你們知道,外邦人有尊為君王的,治理他們,有大臣操權管束他們.

10:43只是在你們中間,不是這樣.你們中間,誰願為大,就必作你們的用人；

10:44在你們中間,誰願為首,就必作眾人的僕人.

10:45因為人子來,並不是要受人的服事,乃是要服事人,並且要捨命作多人的贖價.」

釋經與應用

10:35西庇太的兒子雅各,約翰進前來,對耶穌說:「夫子,我們無論求你甚麼,願你給我們做.」

10:36耶穌說:「要我給你們做甚麼?」

10:37他們說:「賜我們在你的榮耀裡,一個坐在你右邊,一個坐在你左邊.」

上一次,十二門徒試圖自己協調排名,但看來不成功．協調,從來都是不容易的事情．因此,有人想到,與其花時間去和競爭對手磋商,不如另覓他途,撇下對手,直達天庭．雅各和約翰這對兄弟班似乎也採取類似的策略,企圖先發制人,捷足先登．

雅各和約翰如此進取,個別進前來找耶穌,其他門徒會如何反應呢?他們會不會遂個來敲耶穌的門呢?耶穌又如何回應呢?

10:42耶穌叫他們來,對他們說:「你們知道,外邦人有尊為君王的,治理他們,有大臣操權管束他們.

10:43只是在你們中間,不是這樣.你們中間,誰願為大,就必作你們的用人；

10:44在你們中間,誰願為首,就必作眾人的僕人.

10:45因為人子來,並不是要受人的服事,乃是要服事人,並且要捨命作多人的贖價.」

門徒個別進前來找耶穌,私下要求；耶穌的做法是召集他們一齊到他面前來,公開回應．第一種情景是門徒私下前往找耶穌,第二種情景是耶穌召集門徒,這兩種情景對照,就是今天的題目"前來與被召"的對比．

雅各和約翰的行動,相對於耶穌的回應,讓我們聯想到前來與被召這兩種不同的信仰態度．作為耶穌的跟隨著,我們是前來還是被召的呢?還是兩種狀態交雜並存呢?從這段經文所敘述的事件來看,懷著前來的心態和處於被召的狀態可能有甚麼不同呢?

1. 進前來的門徒懷著操控的心態；被召的門徒懷著聆聽的心態

10:35西庇太的兒子雅各,約翰進前來,對耶穌說:「夫子,我們無論求你甚麼,願你給我們做.」

雅各和約翰進前來,並非直接向耶穌提出要求,而是先對耶穌說了上述的話．他們為甚麼要先對耶穌說這些話呢?聽到這些話的人都會預期,對方接著就會提出要求；聽到這些話的人同時也會感受到,對方說這些話是意圖減低要求被拒絕機會,也會感受到對方隱約地在操控自己要滿足他們所提出的要求．

10:42耶穌叫他們來,對他們說:

相對來說,耶穌召集門徒來到他面前,期望他們放下操控的執著,放下要求的欲望,進入虛心聆聽的心境．在這裡,出現了操控的心態和聆聽的狀態的不同和對比．聆聽,從來都不是容易的事．這裡啟發我們反思,會不會聆聽和操控有著此消彼長的關係呢?越想操控的時候,就不願意聆聽；越害怕被操控的時候,也越抗拒進入聆聽的狀態．

來到上主的面前,我們懷著怎樣的心態呢?我們企圖以要求和投訴去操控上主達成我們的期望?還是安靜地聆聽他對我們說話呢?還是,我們已經知道上帝不會被我們操控,但又不想被上帝干預我們的生活和行事為人,所以寧願與上帝維持距離,不讓自己有機會進入聆聽上帝說話的狀態呢?

2.進前的門徒追求別人服事；被召的門徒服事別人

10:36耶穌說:「要我給你們做甚麼?」

10:37他們說:「賜我們在你的榮耀裡,一個坐在你右邊,一個坐在你左邊.」

如果我們聽得清楚門徒的說話,他們並不是要求任何時候都在耶穌的左右,例如並不是說當耶穌釘十字架的的時候,我們一個在你的右邊,一個在你的左邊；他們的要求是有等定條件的,就是當耶穌在榮耀裡的時候,一個在右,一個在左．他們的想像,可能是王帝上朝的景象,一個坐在王帝的右邊,一個坐在左邊,下面有群臣聽從他們的指令,外面有百姓為他們效力．

如果有人問,耶穌的門徒為甚麼那麼熱烈地爭論排名?為甚麼如此努力地爭奪位置?從耶穌對他們的教導來看,門徒很可能念念不忘世界的權力架構和遊戲規則．

「你們知道,外邦人有尊為君王的,治理他們,有大臣操權管束他們.(42)

雖然門徒他們放下一切跟隨耶穌,但世間的權力架構和遊戲規則太過根深蒂固．他們進入耶穌的國度,但所看重的和所追求的,仍然是世間的權力模式,縱使這種權力架構帶來的是管束和剝削,他們自己和他們的社會正在經歷這種羅馬政權的管束和剝削,他們就是無法跳出這種框框．可怕的是,他們把這種權力觀念帶進耶穌的國度,因此,他們爭排名,爭位罝．這些門徒的跟隨者在歷世歷代的教會中爭名譽,爭地位,爭話事權．有些信徒在這些狀況中受傷,或者厭倦這些狀況,因此,四處尋覓,只想尋找一片能夠清心地敬拜上帝的靜土,一處能夠專心地聆聽上帝話語的地方．

耶穌召集門徒來,直接告訴他們,

10:43只是在你們中間,不是這樣.

不是這樣,是怎樣的呢?耶穌接著說:

你們中間,誰願為大,就必作你們的用人；10:44在你們中間,誰願為首,就必作眾人的僕人.10:45因為人子來,並不是要受人的服事,乃是要服事人,並且要捨命作多人的贖價.」

耶穌似乎明白門徒的想望,追求地位和權力來把其他人變為他們的僕人,因此,向他所呼召來的人們澄清,並不是要受人的服事,乃是要服事人,甚至為了別人而犧牲自己．

我們都未必想過要把別人變成自己的僕人,但現實卻可以是另一回事．崇基學院院長陳偉光教授以"獵人與掠物"在本學年第一次學院週會發表了發人心醒又感人至深的演講,其中一段如此說:

"一年後的今天,我心裡更焦急,因為我不知道還有多少日子,可以向你訴說我心裡最關切的事情.不久前,我確診患上了癌症,雖然在暑假期間接受了治療,但無人能夠確知前境如何.我回顧自己一生,由年青時候的尋夢者,跟你們一樣坐在這個崇基的禮堂裡,憧憬著美好的明天,由不甘平庸,矢志努力不懈,全力爭勝,到一天看著自己努力的成果,唯恐失去,於是竭盡所能,全力守護.原來窮一生之力,也不過是另一個由擺脫被掠奪,到自己當上獵人的故事."

耶穌的國度和世界的權力架構是不同的．進前來的門徒效法世界的權力遊戲,為大的治理其他人,為首的操權管束其他人；被耶穌呼召的人效法耶穌,為大的和為首的都是服事其他人．世界的權力架構和遊戲規則,是務求把其他人變成自己的僕人,耶穌的國度,是務求令自己成為其他人的僕人；世界的遊戲規則,是運用其他人來幫助自己達成理想,此為之成功；耶穌的國度,是承存別人,視別人的福祉為自己的成功．

3.從前來到被召:受苦就是服事,也是通往服事的路

耶穌的門徒都是被召的,也是前來的,前來的心態與被召的狀態交雜．如何讓前來的心態逐漸淡出,讓被召的心態植根成長呢?耶穌給予門徒的邀請竟然是分享他的苦難．

10:38耶穌說:「你們不知道所求的是甚麼.我所喝的杯,你們能喝嗎?我所受的洗,你們能受嗎?」

10:39他們說:「我們能.」耶穌說:「我所喝的杯,你們也要喝；我所受的洗,你們也要受；

10:40只是坐在我的左右,不是我可以賜的,乃是為誰預備的,就賜給誰.」

杯和洗,代表審判而來的災殃與苦難,甚或被災難或危險所淹沒．耶穌以此代表為拯救而承受的苦難:

約翰福音18:10西門‧彼得帶著一把刀,就拔出來,把大祭司的僕人砍了一刀,削掉了他的右耳,那僕人名叫馬勒古.18:11於是耶穌對彼得說:「收刀入鞘吧!我父給我的杯,我豈可不喝呢?」

路加福音12:50我有當受的洗還沒有受,在這事完成之前,我是多麼地焦急!

或者,跟隨耶穌的門徒在當時還沒有想到,受苦竟然就是服事,也是通往服事的路．對耶穌而言,這是真確的．

以賽亞書53:5他為我們的過犯受害,為我們的罪孽被壓傷.因他受的懲罰,我們得平安；因他受的鞭傷,我們得醫治.

53:6我們都如羊走迷,各人偏行己路；耶和華使我們眾人的罪孽都歸在他身上.

受苦竟然就是服事?這也是新聞工作者劉進圖先生受到襲擊後的反思:

"到了周末,來探病的朋友告訴我,周日有遊行,左中右的新聞團體破天荒聯合行動,一起反暴力.又有朋友說,這次事件喚醒了港人,關注新聞及言論自由,他認為我這一次受苦有很大的積極意義,可能比我過去的新聞工作貢獻更大.我理性上知道這位朋友為什麼這樣說,他的看法有一定的客觀基礎,但情感上我很難接受,因為他這說法顛覆了我一貫的人生觀.

我從少年時期就返教會,學會人生目標應是榮神益人,努力事奉.過去我一直認為,事奉是奉獻自己的才智,能力和時間,做有益於社會的事情,自然就能榮神益人.例如,我努力讀書,考取好成績,便是向世人作好見證.又例如,我捨棄較賺錢的法律專業,為興趣和使命跑去當記者,還盡量抽時間教學及服公職,這就是努力事奉,榮神益人.

如果我承認今次我遇襲產生了積極的社會意義,那就意味我是通過一個非自願的,被動的,與個人才智能力全不相干的途徑,單純靠承受傷害和痛楚來造就別人,這完全違反了我一貫的概念.我從來不曾明白,所謂無權勢者的權力,受苦的事奉,到底是什麼意思?我大半生走過的所謂榮神益人努力事奉的路,到底當中有多少真實而深刻的意義?有多少其實是為了滿足自我的虛榮?我沒有答案.但我學會一件事,我需要謙卑下來,重新思考人生的方向和意義.(【劉進圖感言之六】受苦的意義 生命的反思)

總結

我們都是進前來尋找上主,也是被上主呼召,因此前來的心態與被召的心態交雜並存．歸根究底,其實我們都是被呼召的．呼召我們的上主盼望我們逐漸地以聆聽祂取代向祂抱怨和索求；逐漸地建立以服事他人取代運用他人的價值觀．最後,讓我們以陳偉光教授的演講的另一個段落作結:

“各位,今天站在你面前的崇基院長,只是一個平庸不過的小人物.但這小人物卻有幸,曾遇過世上的巨人,就是那些一生被好處拿盡,便宜佔盡,落得兩手空空,但仍然傾情盡意,服侍鄰舍,照顧患難中的孤兒寡婦的人；就是那些不懂改變世界,只知要改變自己,沒想過要「砌低」別人,只一心想要「砌低」自己,叫世上從今可以少一個帶來別人不幸的因素的人.世界上沒有多少人知道他們的名字,他們死後亦大概無人記念,但他們在所行經的路上,早已留下了馨香的氣味,佳美的腳蹤.他們的人生,令我無比欣羨!”(陳偉光:獵人與掠物)



\newpage

\section{我們好像做夢的人}
\label{sec:301}
\textbf{鄧瑞強博士 - 我們好像做夢的人}
\newline
\newline
\hyperref[sec:300]{< < < PREV SERMON < < <}
~
\hyperref[sec:toc]{[返目錄]}
~
\hyperref[sec:302]{> > > NEXT SERMON > > >}


講題:「我們好像做夢的人」　We Were Like Dreamers

經文:詩篇126

講員:鄧瑞強博士

場合:香港中文大學崇基學院禮拜堂

日期:2015年10月25日

各位弟兄姊妹,早安.

蘇東坡說:「世事一場大夢」.又說:「萬事到頭都是夢,休休.明日黃花蝶也愁.」成敗得失,轉眼成空.聰明如蘇東坡,人生的起起跌跌,令他自嘆:「我被聰明誤一生」.在不如意的人生中,能做的,就是對月獨飲.「人生如夢,一尊還酹江月」.

夢,看來就是一種不真實,像鏡花水月,不能把握,不能掌控.蘇東坡認為,這就是人生.

心理分析大師弗洛伊德(Freud)卻認為,夢就是願望的實現(wish fulfillment).當心靈從外在世界的壓力中釋放出來,當心靈逃脫了物質的枷鎖,這就是夢(The dream is the liberation of the spirit from the pressure of external nature, a detachment of the soul from the fetters of matter).

看來,日有所思,夜有所夢.在日間,我們有所思,有所期望,但卻被現實之力捆綁著,壓迫著.到夜間,在夢鄉,當一切現實枷鎖被擺脫時,人才能在夢中實現其願望.夢,告訴我們心底的渴望.

面對不能掌控的人生,或許,我們會如蘇東坡般,視之為夢,然後去飲酒消愁.又或者,面對不能掌控的人生,我們仍心中有夢.夢想著那不可能實現的夢境的實現.

夢,總是突破現實的種種限制.夢,是心靈自由的人的專利.夢,是人與神交通的媒介.

我們面對著種種現實枷鎖,政治的,經濟的,人心的.若你說:「現在的這一切是會改變的」.或許有人會說:「你發夢」.這些真的是夢.問題只是:這夢會有實現的一天嗎?

若你問蘇東坡,這夢會有實現的一天嗎?他可能會答,飲酒去罷.

若你問弗洛伊德,這夢會有實現的一天嗎?他可能會答,這夢只能實現在夢裡.

若你問詩篇的作者,這夢會有實現的一天嗎?詩人會答,人間真正的夢,其實也是上帝的夢.在人看來,這夢不可能實現.但在那位「從無造有」,使不可能變成可能的上帝而言,這夢必有實現的一天.

今天講道的經文,是今日的詩篇經課:詩篇126篇.

這詩篇講到夢,講到無法掌控的人生,講到出人意外的夢境實現.

詩126:1　〔上行之詩.〕當耶和華將那些被擄的帶回錫安的時候,我們好像做夢的人.

詩126:2　我們滿口喜笑,滿舌歡呼的時候,外邦中就有人說:耶和華為他們行了大事!

詩126:3　耶和華果然為我們行了大事,我們就歡喜.

詩126:4　耶和華啊,求你使我們被擄的人歸回,好像南地的河水復流.

詩126:5　流淚撒種的,必歡呼收割!

詩126:6　那帶種流淚出去的,必要歡歡樂樂地帶禾捆回來!

這首詩篇十分簡單,人生的逆境是真實的,而夢境的實現是出乎意外的.面對這一切,人能做的,便是大笑一場.

這首詩是猶太人過節時,吃完飯後唱的詩,表達出對前景的盼望,縱使當下的處境,是艱難,是困苦,是叫人流淚的.

藉著這首詩,猶太人宣告,縱使人生黑得見不到一點光,他們仍然相信光明是存在的.光明最後會勝過黑暗.他們認為,越是黑暗,越是要相信光明.沒有這種信念,黑暗便不費吹灰之力將我們吞滅.

這首詩的開頭,對這首詩的性質有一種稱呼,稱之為「上行之詩」.「上行之詩」的真正意思,已經失傳了.但一般學者認為,這是朝聖者一步一步接近耶路撒冷時唱的詩歌.若果這講法是真的,則這首詩表達出一個敬拜神的人,當他的心靈越來越接近上帝時,上帝在他的心靈中興起的盼望.

詩126:1　當耶和華將那些被擄的帶回錫安的時候,我們好像做夢的人.

詩歌的表達,有時留下詮釋的空間.「當耶和華將那些被擄的帶回錫安的時候」,可以理解為,「當耶和華復興錫安的時候」,或可以理解為,「當耶和華轉面面向錫安的時候」.無論如何,這經文表達一種轉變,逆境被扭轉,失望被消除,黑暗被驅散.原先,神好像背向祂的子民.如今,神面向祂的子民.當神露面的時候,一切被改變了.

當這一刻出現時,詩人說:「我們好像做夢的人」.這是否發夢?詩人說:是啊,是啊,我們正在發夢.或許有人說:「起來吧,不要再發夢了」.詩人會說:「我不是睡著,這是真的,這是夢境成真」.

「當耶和華將那些被擄的帶回錫安的時候」,這如何可能?被擄的人,怎可能回歸故土?作奴隸的人,怎可能得到自由?不可能有希望的人,怎可能重拾希望?

這是因為神的面轉向他們,這是因為神實現祂的救贖應許,這是因為神根本反對世上一切的欺壓和強暴.神定意要實現人類真正的夢想.

對蘇東坡而言,夢只是表達世事的虛幻.對弗洛伊德而言,夢是人無法實現卻也無法放棄的心底渴求.但在聖經的世界,夢是神和人溝通的媒介.透過夢,神向人啟示祂的救贖計劃.在舊約的創世記,約瑟的夢暗示了神的子民的救贖.在新約的馬太福音,另一個約瑟的夢預示著救主耶穌的降生.在聖經裡,夢不是人無法實現的渴望,夢是神定意要實現的作為.夢,是先知的遠象.先知看到神的手在行動.先知在人生的迷惘裡,叫人因著神,不要絕望.

猶太教有一種講法,人死後,會被審問.而被問的第一個問題,便是:在人生中,你是否對救贖有所期待?或者說:在黑暗中,你是否相信光明?

有一個故事是這樣的.

在基督教世界的某個角落,有一條小小的村莊.這條村莊務農為生,但遭逢大旱.村裡的牧師呼籲所有村民到教堂祈禱求雨.所有人陸陸續續走向教堂,在人群中,有一個小孩子,也唯有這個小孩子,帶著一把傘,一把黃色的傘.當所有人都只是兩手空空地去求雨時,有一個小孩子帶著一把傘,一把黃色的傘.

可以想像,當雨真的落下來時,唯有這個小孩子,能舉起雨傘.當上主問這群人,「你們是否對救贖有所期待?」唯有這個小孩子,能大聲答:「我期待」.

當耶和華將自由賜予祂的子民的時候,我們好像做夢的人.這怎麼可能?這是否只是虛幻的夢境?在神的大能裡,人心底的夢境是會成真的.

詩126:2　我們滿口喜笑,滿舌歡呼的時候,外邦中就有人說:耶和華為他們行了大事!

詩126:3　耶和華果然為我們行了大事,我們就歡喜.

這不是夢,這是真的,連不信的人都驚訝,神竟然行了如此的大事.

信仰的夢,總是關乎現實世界的改變的.為了這改變,我們先改變自己.如何改變?我們堅信上帝會按祂的恩典與救贖而行.這種信念,會在信仰者身上流露出來.有這信念的人,行事公正,毫無詭詐,尊重他人的尊嚴,對天國的實現滿懷希望.

上帝會實現人間的救贖,等候上帝的人,總不失望.

詩人說:「耶和華果然為我們行了大事,我們就歡喜」.這是真正的笑,上帝抹去人的眼淚.

詩126:4　耶和華啊,求你使我們被擄的人歸回,好像南地的河水復流.

這是詩歌的第二個段落,再次提到「被擄歸回」.

但在這裡,「被擄歸回」不是一個現實,而是祈求的內容.

看來,在人世間,人經歷到神的救贖後,不是一步到天堂.恩典之後,是生命的另一艱難.過去的恩典,是讓我們對神滿懷信心.憑著這份信心,面對另一個逆境.

詩人求神的恩典臨到,「好像南地的河水復流」.憑這詩句,我們推斷他們當時面對的逆境是大旱.無水,便無法耕種.無法耕種,便無法生活.這裡講的「南地」,便是巴勒斯坦地的南部,那裡是「內蓋夫沙漠」(Negev Desert).雨水稀少.但是,若下起雨來,只花幾個鐘,河水便氾濫.於是,「南地的河水復流」,是指雨水來得急,來得猛,來得叫人措手不及.而雨水,代表著神的恩典.而神的恩典臨到,就像南地的雨水一樣,叫人預計不到,叫人驚喜莫明.

詩126:5　流淚撒種的,必歡呼收割!

若果他們真的遇到旱情,則他們看著他們已經撒下的種,必然會流淚.

無論是否遇著旱情,撒種的人總是流著淚撒的.撒種,是經營生命,是對生命的發芽生長有期待,但同時對生命的發芽生長毫無把握.這是叫人流淚的.為人父牧者對子女的成長,為人師表者對學生的栽培,社會運動者對社會改善的期望,這一切,都是叫人流淚的.因為,要成就這些事情,要下的苦功多,但對於是否真有收成,卻毫無把握.

「必歡呼收割」,這是流淚撒種的人的期盼.這期盼會否實現?一定會的.因為,在過去,神叫「被擄的人歸回」,這已表明,神的恩典是真實的,久旱逢甘露是真實的,在神的手裡,人類的善良的夢是會實現的.

於是,詩人大膽地宣告他的結論:

詩126:6　那帶種流淚出去的,必要歡歡樂樂地帶禾捆回來!

一般而言,詩歌都是首尾呼應的.第一節講「夢」,如今,結尾講「禾捆」.這如何首尾呼應呢?學者發現,舊約聖經裡的一個重要的夢,便是約瑟的夢,而約瑟夢見的,正正是禾捆.後來,約瑟做了埃及的宰相,而他的家人卻在迦南地斷糧.最後,各人在夢見禾捆的約瑟手下得到解救.

這首詩歌道盡了人間的起跌.悲哀的奴役之後,有叫人歡笑的釋放.開心的背後,又有憂愁.憂愁之時,卻仍要盼望喜樂.沒有怨天尤人,沒有絕望.只是把持著神的應許,走在歷史艱難的路上,深信在歷史盡頭,神的應許會實現.

在這苦難的世界,我們都是發夢的人.

在二十世紀,有一個人發了一個千秋大夢.美國的馬丁路德金牧師說:「我有一個夢.我夢見有一天,這個國家會起來實現它的信條的真正意義:『我們認為這些真理是自明的,這真理就是人人被創造而為平等的』.我夢見有一天,在喬治亞的紅山上,昔日奴隸的兒子,和昔日奴隸主的兒子,能坐下來分享手足之情.…我夢見有一天,我的四個孩子能活在一個不以他們的膚色而以他們的品德來評價他們的國家裡.…今天,我有一個夢.我夢見有一天,一切深谷都要填滿,一切山岡都要削平；險峻的要改為平地,崎嶇的要改為平原.耶和華的榮耀必要顯現,所有的人都必一同看見.〔引自賽40: 4-5,這裡譯文採自＜新譯本＞〕」

在一個不正義的世界,馬丁路德金按神的應許,看到歷史的前景,然後講出他的夢.今日,他的夢實現了.流淚撒種的,今日能歡呼收割!

或許,我們應該發夢.讓我們夢見,有一天,嬰孩能在愛的環境中成長,不用吃假的奶粉；年輕人能在大人中學到正路與真理,不用假話上位；中年人有工可作,有屋可住,不怕因講一點理想而被捕；老人家能不用「嫌命長」,有一個尊重他們的環境,有一個養育他們的社會.

這些夢會實現嗎?詩人說,這些夢會實現的.

但願:榮耀歸於聖父,聖子,聖靈.阿們.



\newpage

\section{誡與命愛與怕}
\label{sec:302}
\textbf{許開明牧師 - 誡與命,愛與怕}
\newline
\newline
\hyperref[sec:301]{< < < PREV SERMON < < <}
~
\hyperref[sec:toc]{[返目錄]}
~
\hyperref[sec:303]{> > > NEXT SERMON > > >}


講題:誡與命,愛與怕　Commandment and Life, Love and Fear

講員:許開明牧師

場合:香港中文大學崇基學院禮拜堂

日期:2015年11月1日

你有自由嗎?

各位,你有自由嗎?現在的你有行動的自由?如果你無行動的自由,你會怎樣?現在的你有說話的自由?如果你無說話的自由,你又會怎樣?現在的你有選擇做或不做的自由嗎?如果你沒有選擇做或不做的自由,你會怎樣?現在的你有信仰自由?如果你沒有信仰的自由,你會變成怎樣?

自由好不好?

自由是好或是不好?大家都會認為自由是「好」的,不自由是「不好的」.一般人對自由最普遍的觀念是:人人生而自由,享有使用各種權利之自由,擁有言論,信仰,宗教的自由,只要不侵犯他人,現代人都非常自由.但是過去的人並非如此,例如:有些人生下來就是奴隸,而他所生的小孩也是奴隸,世世代代都是沒有自由的奴隸；直至今天,在印度的印度教徒仍然在遵守這個制度.在過去的時代也沒有言論自由,沒有受教育的自由,沒有信仰的自由,沒有戀愛的自由,沒有參與政治的自由,沒有爭取權力的自由.以前的人並不是想做麼就能做什麼.但是經過14,15世紀文藝復興的人文主義運動,16世紀的宗教改革,17世紀的啟蒙運動,到18世紀法國大革命,19世紀的南北戰爭,以及20世紀的各種自由運動之後,這個時代可以說是有史以來最自由的時代,現代人愛做什麼就能做什麼.我們想去哪裡就可以去那裡,想買什麼就可以買什麼,想吃什麼就可以吃什麼,想做什麼就可以做什麼.現代人有了自由,感到自在,方感到人生幸福.

但是,如果人有無限自由,絕對的自由,又是一件好事嗎?如果甚麼都可以做,如果甚麼都可以不做,這個人會快樂嗎?答案當然是「不會」.

人有無限和絕對的自由嗎?

人沒有絕對的,無限的自由!因為人類受造的本體是「自由中的不自由」,因我們在本質上有許多限制,如:受創造律,自然律,罪權與死權的限制；受限於時限與年限(時間,年齡),和所住的疆界(空間).人活在造物主創造的定律裡,在凡事有定期,定時,定理下運作,人無能逆轉.(傳3:1-8, 8:6-8)還有因著始祖犯罪,人類就世世代代伏在罪權轄制之下立志為善由得我們,只是行出來由不得我.(羅7:18下)我們成了罪奴,不得自由,是何等痛苦.就在這種種限制中,逼使我們尋求造物主,或者可以揣摩而得到神(徒17:27).在自己無法逆轉的情況下,我們應轉而敬畏仰望這位設定天下萬事萬務的上主(傳3:14).

誰有絕對的自由?

那就是基督教的上帝-耶和華,祂無形無體,卻是抽象的,超越的觀念.在出埃及記3章13至15節,耶和華說:「我是『自有永有者』.」「自有永有者」的意思,就是神是自己存在,不依靠任何別的存在而存在,他是永恆的存在,他不是被造的.而「我是」就是「我存在」的意思；耶和華是「我是」的上帝,強調自我自由的上帝.「我是」-源自希伯來文的連繫動詞Hayah (是)；「我是自有的」意思是有兩個,第一個是“I am what I am”或“I am who I am”或“I am that I am”,而第二個就是“I cause to be”.所以,在信仰方面來看,耶和華是自我的,自由的,自在的上帝.祂擁有真自由,並能賜人真自由；人須要與上主有恆常美好的連結,方能成為真正有自由的人.

一些未有信仰的朋友,他們不肯成為基督徒,因怕「失去自由」.一些朋友說信耶穌很好,基督徒也不錯,但是因為成為基督徒後就不能抽煙,不能醉酒,不能說謊；星期天還要到教會,要奉獻,要禱告；即是很多他們本來喜歡的事都不能做,許多他們不喜歡的事卻必須做,因而認為成為基督徒是非常不自由.他們喜歡無拘無束,所以他們不想成為基督徒.似乎在上帝律法及誡命的約束下,基督徒沒有什麼自由,這是一個不爭的事實.信耶穌,得真理,真理不是能使人得自由嗎?為什麼信了耶穌,就不准這樣又不准那樣,反而要受束縛呢?

聖經的自由

然而,聖經不是強調「識真理,得自由」嗎?認識耶穌,不是能使人得自由嗎?怎會反而失去自由呢?

到底聖經所反映的自由是什麼的自由呢?什麼都可以做就是自由嗎?被上帝的命令約束就等於不自由嗎?答案是恰好相反.那麼,真正的自由是什麼?自由的定義到底是什麼?有人說:「自由就是有選擇」,可是怎樣才能有選擇呢?唯有在有些能做,有些不能做的情況下,才會產生選擇的機會,然後再選擇當中的,才有真正的自由.也就是說:「有限制才有選擇,有選擇才有自由」.就好像學生考試中的是非題,考試的方式限制你只能打V或打X,雖然有限制,但也因此才有選擇,在選擇中才有真正的自由.

當一艘船決定要到美國時,就不能去英國了,這艘船的乘客就失去了隨便航行自由了,因為他們被限制一定要去夏威夷這個目的地.但實際上,乘客卻在限制中感到快樂.這說明了什麼?當我們以為沒有限制就是很自由,可以隨心所欲的自由航行,但當眾人不知要到哪裡去,反而心被懸在那裡,這是虛假的自由.可是當我們有一個目的地,表面上雖然使這艘船的航行,受到了這個目的地的限制,因而不能隨心所欲的自由航行,但實際上,正正因為有限制,反而使船有了方向,使眾人的心有所歸依,從而得到了真正的自由.這就叫做「有限制才有自由,顯出自由的可貴!」.

在選擇中的自由,上帝的確是可以造一個沒有分別善惡樹的伊甸園,也不需要讓人有限制,讓人不需要遵守不可吃分辨善惡果子的命令,這樣,亞當和夏娃,不就可以無拘無束,自由的生活在沒有命令,而且所有果子都可吃的伊甸園中嗎?

但是,若是如此,那麼伊甸園才真正是一個「沒有自由」的環境,亞當夏娃就是最沒有自由的人. 因為所謂的「自由」,就是要在環境中提供一個讓人能夠有選擇的機會.若是人只能遵守上帝命令,而沒有違背上帝命令的機會,那麼就表示沒有選擇,那麼也意味著人沒有自由.只有當人可以選擇吃或不吃,遵守或違背命令時,才算是有真正的自由.

神不是創造一個完全只能接受祂的愛,聽從其命令的機械人,若是這樣,人就更可悲,伊甸園就很可怕了!因為上帝乃是創造一個擁有自由意識,可自行作選擇的「人」.所以,伊甸園中有不能吃的果子,以及後來人會違背上帝的命令,這些都證明了人是有自由的人.伊甸園中有不能吃的果子正是人有自由的記號.這就叫做「有選擇才有自由,有限制才有選舉,有選擇才有自由.」

對耶穌的考問

回到今日的經課:馬可福音12章,首先是有關納稅給凱撒的問題.(太22．15－22；路20．19－26)由第13節開始:後來,他們打發幾個法利賽人和希律黨人到耶穌那裡,要用他自己的話陷害他.他們來了,就對他說:「老師,我們知道你是誠實的,無論誰你都一視同仁；因為你不看人的面子,而是誠誠實實傳神的道.他們打發幾個法利賽人和希律黨人到耶穌那裡,問他納稅給凱撒合不合法.15我們該不該納?」法利賽人:反對納稅給該撒；希律黨人支持納稅給該撒.耶穌知道他們的虛偽,就對他們說:「你們為甚麼試探我?拿一個銀幣來給我看.」16他們就拿了來.耶穌問他們:「這像和這名號是誰的?」他們對他說:「是凱撒的.」17耶穌對他們說:「凱撒的歸凱撒；神的歸神.」他們對他非常驚訝.

在這裡,耶穌的回答是:凱撒的歸凱撒；神的歸神.耶穌的回答使兩方都滿意和接受!耶穌答得好,勝了!

復活的問題

(太22．23－33；路20．27－40)

馬可福音還有講復活的問題,由第18節開始:撒都該人來見耶穌.他們說沒有復活這回事,於是問耶穌:19「老師,摩西為我們寫下這話:『某人的哥哥若死了,撇下妻子,沒有孩子,他該娶哥哥的妻子,為哥哥生子立後.』20那麼,有兄弟七人,第一個娶了妻,死了,沒有留下孩子.21第二個娶了她,也死了,沒有留下孩子.第三個也是這樣.22那七個人都沒有留下孩子.最後,那婦人也死了.23在復活的時候,她是哪一個的妻子呢?因為他們七個人都娶過她.」24耶穌說:「你們錯了,不正是因為不明白聖經,也不知道神的大能嗎?25當人從死人中復活後,也不娶也不嫁,而是像天上的天使一樣.26論到死人復活,你們沒有念過摩西書中《荊棘篇》上所記載的嗎?神對摩西說:『我是亞伯拉罕的神,以撒的神,雅各的神.』27神不是死人的神,而是活人的神.你們是大錯了.」

在這裡,耶穌的回答是:「人死後無嫁娶事,誰是那女人的丈夫,根本不成問題.還有那句:神不是死人的神,而是活人的神.」又是雙方都討好!耶穌答得好,勝了!

最大的誡命

(太22．34－40；路10．25－28)

第三位來考問耶穌的是文士,是高人,身份高,知慧高,比前面的兩位:撒都該人,更厲害的一位?耶穌在這次答辯中,勝了嗎?

馬可福音記載:那文士來,聽見他們的辯論,知道耶穌回答得好,

就問耶穌說:「誡命中哪一條是第一呢?」29耶穌回答:「第一是:『以色列啊,你要聽,主-我們的神是獨一的主.30你要盡心,盡性,盡意,盡力愛主-你的神.』31第二是:『要愛鄰如己.』再沒有比這兩條誡命更大的了.」32那文士對耶穌說:「好,老師,你說得對,神是一位,除了他以外,再沒有別的了；33並且盡心,盡智,盡力愛他,又愛鄰如己,要比一切燔祭和祭祀好得多.」34耶穌見他回答得有智慧,就對他說:「你離神的國不遠了.」從此以後,沒有人敢再問他甚麼.

在我看來,這文士表面在讚賞耶穌,耶穌回答撒都該人和法利賽人均好,而且他本人也非常讚同耶穌的觀點!然而,就在同一段的經文,在馬可福音12章38至40節,耶穌特別提到門徒要防備文士的偽善,並具體指出他們是如何的偽善.

如果我們了解,當時背景可見其中有不簡單的地方:耶穌的年代已有兩種發展的傾向,猶太主義者在解釋上有一種雙重的傾向,他們一方面把律例條文無限地微分成數百千條的細則；另一方面則把繁雜的條文縮短,以一句概括的話表達出來.有一次,一位仰慕猶太教的人向偉大的拉比希利請教,怎樣才可以把一切律法用一句話作出總綱呢?森黎(Sammlai)曾說過,摩西從西乃山上接受了613 條訓令；當中365 條是照一年的日曆記錄出來,餘下的248條是照人的世代次數而定的；後來到了大衛,他把613條縮短,在這事上,人們大致分為兩派,第一派是繁文褥節派,強調律法,並以守法為榮,文士主要屬此派.第二派是簡化派,反對律法主義,愛精簡.

在這兩派爭辯中:耶穌傾向簡化派,這肯定引起了繁文褥節派的不滿.那文士成功地令耶穌表達他是簡化派的；使另一方人反感.當然耶穌應會知道,有此後果的.

在誡命下的自由

上帝在造人之後,安排有分別善惡樹,人不可吃,這是上帝定下的規矩.在摩西時代,上帝頒布十誡,目的有兩個,一是要人知道什麼是做得對,什麼是做得不對.二是要人知道上帝的要求標準,人才能懂得達標!詩篇119篇44至45節說到:「我要常守你的律法,直到永永遠遠.我要自由而行,因我素來考究你的訓詞.」英文版的聖經是說:「I will walk about in freedom, for I have sought out your precepts. (我要尋求禰的律法,使我能行在自由之中).」這裡說到:人在上帝的律法及誡命中才能得到自由.這是不是搞錯了!照道理說,人在律法及誡命的約束中,應當是不自由的,怎麼會是自由的呢?這真是令人感到矛盾又困惑,為什麼聖經中的自由觀與我們一般人所想的自由觀不同呢?我們今天就要來思考「自由」的這個主題.

沒有法律限制,沒有道德觀念,什麼都可以做,是不是自由呢?這樣反而是沒有自由,因為什麼都可以做,表示沒有限制,沒有限制表示沒有選擇,沒有選擇就沒有自由.所以,現在我們發現詩篇119篇44至45節,這一個經節,最少隱含了4個真理:

一,上帝的律法,誡命或法則,並不表示使人失去自由,恰好是使得到自由.因為有限制就表示有選擇,有選擇表示有自由.

二,上帝的法則使讓人的行為有一個準則,使人心有所往,使生命有方向,因此使生命有了真正的自由.

三,上帝的誡命使人知道什麼是應該做的及不應該做的,使人確知上帝的標準.

四,上帝的律法使人的行為受到約束,使人在律法之下有所保障；就像十誡的後六誡:當孝敬父母,不可殺人,不可姦淫,不可偷盜,不可陷害人,不可貪戀他人的所有.你不要以為這限制你的自由,其實是保障你的自由.

五,上帝的律法使人有一條蒙福的管道,因為上帝有應許,凡是遵行的人就蒙福.

這就是上帝律法中所隱含的偉大真理,它將我們原本錯誤的自由觀念完全扭轉過來.

誡與律是聖經中一大教訓,上帝藉此教導我們應當怎樣為人,並且指示我們當作和不當作的.(馬丁路德)上帝的律法,立下行為的準則.上帝說了,人跟著做!誡與律限制約束人的行為,使人行在正途上.上帝的律法,早在上帝造人時,已刻在人的心板上；然而,因人的軟弱,上帝也刻在石板上.

在倫理選擇下的自由

讓我們看看保羅自由觀的原則和選擇的標準:

一,善與惡

保羅在羅馬書7章14至20節,他提出在他生命中有兩個律,一個是肉體的惡；另一個是心裡的善.他常要作選擇,然而,許多時他卻感力有不逮!

二,益與害

哥林多前書6章12節說到:「凡事我都可行．但不都有益處.凡事我都可行,但無論那一件,我總不受他的轄制.」哥林多前書10章23節也說:「凡事都可行．但不都有益處.凡事都可行．但不都造就人.」聖經只有一個原則,凡事都可行,但不都有益處.上帝給我們自由；上帝允許我們去做；上帝任憑我們做,但是不是每件事都對我們有益處,而且我們還需要對我們所做的負責任,有時還要付出代價.

三,有能力做與不做

因此現在我們瞭解到:自由不是可以做你想做的；而是選擇不做你可以做的.當父母,牧師,傳道人說:不能做這件事,所以你就不去做,這個叫「順服」.因為你發現做這件事對你不好,而你不願去做,這個叫「自由」.因為你可以做,有機會能做,但你知道你是上帝揀選的,主重價買贖回來的基督徒,是上帝的兒女,所以願意不做,這才叫「自由」.

愛與怕

在有限的自由中,我們守律法誡命是基於什麼?基於怕?我怕殺人要填命,要判死,所以我不會殺人.我怕偷東西,會被拉,被判坐監,所以我不會偷.我怕遲返家罰無零用錢駛,所以我早於12時就返家!因怕才去守律法,不錯,但不算好,因是消極的,無奈的,非純出於真心的!

我們守律法誡命應該是基於愛!因為愛上帝,你遵守上帝的誡命!因為愛人,你不會殺人和傷害人!愛產生了推動力!在耶穌來說,一個能愛神愛人的人,他便能守十誡,甚至數以百計的律例!不過,你也要留意:耶穌所言的愛,不是賣口乖,而是要付上一切的愛!馬可福音12章30節說:「你要盡心,盡性,盡意,盡力,愛主你的神.」

「盡」的原文字義是全部,全然的意思；「愛」的原文字義是神聖的愛(agapao).「盡心」(all your heart)即是你的感情；「盡性」(all your soul)即是你的意志；「盡意」(all your mind)即是你的思想；「盡力」(all your strength)就是你的力量.「愛主你的神」意即以整個人,全生命的貫注去愛神；在此四重的說法旨在表達強調之意.

馬可福音12章31節這樣說:「其次,就是說,要愛人(鄰舍)如己(Love your neighbor as yourself).再沒有比這兩條誡命更大的了.」愛別人好像愛自己,你怎樣愛自己?其實就是說你要:「盡心(all your heart)你的感情,盡性(all your soul)你的意志,盡意(all your mind)你的思想,盡力(all your strength)你的力量」來愛自己!願我們是基於愛去守誡命,行律法!

因著愛,願犧牲自己的自由,使人得真自由.

聖經中彰顯出我們所信的上帝,乃是一位願意犧牲自己自由,使我們能獲得真正自由的上帝.我們所信的主,愛我們愛到一個地步,祂願意犧牲自己的生命,使我們能得到真正的自由.這也是基督信仰與其他宗教信仰大衛不同的地方.其他宗教中的神明大都有完全的自由,然後要求信徒要去虔誠修道,以便能獲得真正自由.但是我們所信的神,卻是願意犧牲自由,而把真正的自由給我們.

因著愛,上帝甘願失去自由.

基督徒所信的全能上帝當然是最自由的.但是祂因為愛我們的緣故而願意約束祂自己,願意失去自由.

為什麼上帝會失去自由呢? 因為當上帝和我們立約的時候,上帝就失去了完全的自由.因為這個時候,上帝必須要遵照祂和我們所立的約定來行事.當我們遵守約定時,上帝就必須遵守,祂沒有自由可以違背這個約定.他在與人的約定中,祂失去了自己的自由.「約定」就是:簽約的雙方過渡了一定的權力,自願被約定內容所限制.所以「新約」及「舊約」中那些人與上帝的約定,都是上帝失去自由的「記號」.比如說,「彩虹」就是上帝失去自由的記號.因為創世記第9章那裡說到上帝與人之間的那個「彩虹之約」.同樣地,聖經中的「預言」及「應許」,也都是上帝失去自由的記號.而「先知」則是限制上帝自由的人,因為先知所傳講的預言及上帝的又許限制了上帝的作為.上帝必須按照預言而行事,必須使應許實現.因此,「先知」也就是使上帝失去自由的證人.

主耶穌也為我們失去自由

不只是我們所信的這位上帝願意為我們而失去自己的自由,我們的主耶穌也是如此.主耶穌原本是自由的,但因為愛我們的緣故,願意降世成為不自由的人；並且為了使我們能脫離罪惡的轄制,與上帝合好,主願意代替我們的罪犧牲在十架上.

「馬槽」與「十字架」都是主耶穌失去自由的記號.主耶穌為了使我們能獲得真正的自由,而不受罪的轄制,因此他願意失去祂的自由,使我們能得到自由.所以,真正的自由並不是我們能擁有和我們配得擁有的,但是主卻願意透過在十字架上的犧牲而將真正的自由賜給我們.他乃是用他的身,以及他的血來使我們獲得真自由.

到這裡我們發現:真正有自由的神,為了使我們得到自由,卻願意為我們而犧牲自己自由.耶穌並非選擇自由,而是選擇犧牲自由,選擇受苦；這就是主對我們的愛.因為主愛我們,祂願意失去自己的自由,而使人獲得真正的自由.這才是從亙古到永遠,自由當中最偉大的真理.

我們愛神愛人也值得愛到失去自由

假如基督徒宣稱我們是主的門徒,那麼我們應當要效法我們的主,願意犧牲自己的自由,而使他人獲得真正的自由.我們雖然在主裡面已經擁有真正的自由,但是我們卻願意犧牲自己的自由,為要使他人獲得在主裡面的真自由.因此,在哥林多前書9章19節,保羅說:「我雖是自由的,無人轄管；然而我甘心作了眾人的僕人,為要多得人.」為使他人得到自由,我們應該願意犧牲自己自由,這也就是基督應該走的十架道路.

願我們是基於愛去守誡命,行律法!並愛到一個地步為神為人犧牲自己的自由,使別人得釋放得自由!耶穌做得到,你做得到嗎?



\newpage

\section{信仰中的有餘多餘與不足充足}
\label{sec:303}
\textbf{鄭漢文博士 - 信仰中的有餘(多餘)與不足(充足)}
\newline
\newline
\hyperref[sec:302]{< < < PREV SERMON < < <}
~
\hyperref[sec:toc]{[返目錄]}
~
\hyperref[sec:304]{> > > NEXT SERMON > > >}


馬可福音 12:38-44
\newline
\begin{longtable}{cl}
\hline
\hline
章節 & 經文 (和合本修訂版)\\
\hline
12:38 & \begin{tabularx}{0.7\textwidth}{X} 他在教導的時候，說：「你們要防備文士。他們好穿長袍走來走去，喜歡人們在街市上向他們問安， \end{tabularx} \\ \\ \relax
12:39 & \begin{tabularx}{0.7\textwidth}{X} 又喜愛會堂裡的高位，宴席上的首座。 \end{tabularx} \\ \\ \relax
12:40 & \begin{tabularx}{0.7\textwidth}{X} 他們侵吞寡婦的家產，假意作很長的禱告。這些人要受更重的懲罰！」 \end{tabularx} \\ \\ \relax
12:41 & \begin{tabularx}{0.7\textwidth}{X} 耶穌面向聖殿銀庫坐著，看眾人怎樣把錢投入銀庫。有好些財主投了許多錢。 \end{tabularx} \\ \\ \relax
12:42 & \begin{tabularx}{0.7\textwidth}{X} 有一個窮寡婦來，投了兩個小文錢，就是一個大文錢。 \end{tabularx} \\ \\ \relax
12:43 & \begin{tabularx}{0.7\textwidth}{X} 耶穌叫門徒來，對他們說：「我實在告訴你們，這窮寡婦投入銀庫裡的比眾人所投的更多。 \end{tabularx} \\ \\ \relax
12:44 & \begin{tabularx}{0.7\textwidth}{X} 因為，眾人都是拿有餘的捐獻，但這寡婦，雖然自己不足，卻把她一生所有的全都投進去了。」 \end{tabularx} \\ \\
[1ex]
\hline
\hline
\end{longtable}
列王記上 17:8-16
\newline
\begin{longtable}{cl}
\hline
\hline
章節 & 經文 (和合本修訂版)\\
\hline
17:8 & \begin{tabularx}{0.7\textwidth}{X} 耶和華的話臨到他，說： \end{tabularx} \\ \\ \relax
17:9 & \begin{tabularx}{0.7\textwidth}{X} 「你起身到西頓的撒勒法去，住在那裡，看哪，我已吩咐那裡的一個寡婦供養你。」 \end{tabularx} \\ \\ \relax
17:10 & \begin{tabularx}{0.7\textwidth}{X} 以利亞就起身往撒勒法去。他到了城門，看哪，有一個寡婦在那裡撿柴。以利亞呼喚她說：「請你用器皿取點水來給我喝。」 \end{tabularx} \\ \\ \relax
17:11 & \begin{tabularx}{0.7\textwidth}{X} 她去取水的時候，以利亞又呼喚她說：「請你手裡也拿點餅來給我。」 \end{tabularx} \\ \\ \relax
17:12 & \begin{tabularx}{0.7\textwidth}{X} 她說：「我指著永生的耶和華－你的神起誓，我沒有餅，罈內只有一把麵，瓶裡只有一點油。看哪，我去找兩根柴，帶回家為我和我兒子做餅。我們吃了，就等死吧！」 \end{tabularx} \\ \\ \relax
17:13 & \begin{tabularx}{0.7\textwidth}{X} 以利亞對她說：「不要怕！你去照你所說的做吧！只要先為我做一個小餅，拿來給我，然後為你和你的兒子做餅； \end{tabularx} \\ \\ \relax
17:14 & \begin{tabularx}{0.7\textwidth}{X} 因為耶和華－以色列的神如此說：『罈內的麵必不用盡，瓶裡的油必不短缺，直到耶和華使雨降在地上的日子。』」 \end{tabularx} \\ \\ \relax
17:15 & \begin{tabularx}{0.7\textwidth}{X} 婦人就照以利亞的話去做。她和以利亞，以及她家中的人，吃了許多日子。 \end{tabularx} \\ \\ \relax
17:16 & \begin{tabularx}{0.7\textwidth}{X} 罈內的麵果然沒有用盡，瓶裡的油也不短缺，正如耶和華藉以利亞所說的話。 \end{tabularx} \\ \\
[1ex]
\hline
\hline
\end{longtable}
詩篇 146:1-10
\newline
\begin{longtable}{cl}
\hline
\hline
章節 & 經文 (和合本修訂版)\\
\hline
146:1 & \begin{tabularx}{0.7\textwidth}{X} 哈利路亞！我的心哪，你要讚美耶和華！ \end{tabularx} \\ \\ \relax
146:2 & \begin{tabularx}{0.7\textwidth}{X} 我一生要讚美耶和華！我還活著的時候要歌頌我的神！ \end{tabularx} \\ \\ \relax
146:3 & \begin{tabularx}{0.7\textwidth}{X} 你們不要倚靠君王，不要倚靠世人，他一點也不能幫助。 \end{tabularx} \\ \\ \relax
146:4 & \begin{tabularx}{0.7\textwidth}{X} 他的氣一斷，就歸回塵土，他所打算的，當日就消滅了。 \end{tabularx} \\ \\ \relax
146:5 & \begin{tabularx}{0.7\textwidth}{X} 以雅各的神為幫助、仰望耶和華－他神的，這人有福了！ \end{tabularx} \\ \\ \relax
146:6 & \begin{tabularx}{0.7\textwidth}{X} 耶和華造天、地、海和其中的萬物，他守信實，直到永遠。 \end{tabularx} \\ \\ \relax
146:7 & \begin{tabularx}{0.7\textwidth}{X} 他為受欺壓的伸冤，賜食物給飢餓的人。耶和華釋放被囚的， \end{tabularx} \\ \\ \relax
146:8 & \begin{tabularx}{0.7\textwidth}{X} 耶和華開了盲人的眼睛，耶和華扶起被壓下的人，耶和華喜愛義人。 \end{tabularx} \\ \\ \relax
146:9 & \begin{tabularx}{0.7\textwidth}{X} 耶和華保護寄居的，扶持孤兒和寡婦，卻使惡人的道路彎曲。 \end{tabularx} \\ \\ \relax
146:10 & \begin{tabularx}{0.7\textwidth}{X} 耶和華要作王，直到永遠！錫安哪，你的神要作王，直到萬代！哈利路亞！ \end{tabularx} \\ \\
[1ex]
\hline
\hline
\end{longtable}
講題:信仰中的有餘(多餘)與不足(充足)　(Un)necessariness and (In)sufficiency in Faith

經課:馬可福音12:38-44;列王記上17:8-16;詩篇146:1-10

講員:鄭漢文博士

場合:香港中文大學崇基學院禮拜堂

日期:2015年11月8日

開始(In the beginning)

從前有一個人,自年青時已開始形成一種習慣,就是當親朋戚友生日時,當他要送禮物時,都會在家裡找,找自己不再有用的物品,即是對他來說是多餘的東西,包裝好後,作為生日禮物送出.如是者,年復一年,年復一年.

直至有一天,他退休兼六十大壽,他筵開百席,所有他曾出席生日會的親朋戚友都被邀請出席贈慶,各位來賓都送上生日禮物.

筵席散去後,回到家,他把一份又一份的禮物打開,他發現一份又一份曾經從家裡送出的物品,一一都返回到自己的家裡.

大家好,我今日返回到這裡──好像我的家的崇基禮拜堂,闊別了這個講壇近一年半.今天要同大家分享的經課記載在馬可福音第12章38至44節.

讓我們從三個對比來解讀這段經文:

第一個對比是:譴責(blameworthy)與讚賞(praiseworthy)

第二個對比是:有餘(residual/unnecessary)與不足(insufficient)

第三個對比是:多餘(redundant/unnecessary)與充足(sufficient)

第一個對比是:譴責文士與讚賞寡婦

(First contrast: blameworthy scribes vs praiseworthy widow)

這段經文,耶穌究竟在譴責誰?在讚賞誰?

耶穌明顯地在譴責文士,讚賞寡婦.他為何譴責文士?他為何讚賞寡婦?

首先,讓我們看看耶穌如何譴責文士(另可參看:太23:1-36；路20:45-47):

「38他在教導的時候,說:“你們要防備文士.他們好穿長袍走來走去,喜歡人們在街市上向他們問安,

39又喜愛會堂裡的高位,宴席上的首座.

40他們侵吞寡婦的家產,假意作很長的禱告.這些人要受更重的懲罰!”」

如果放在今天的香港,你會想起誰好像這些文士?

若根據耶穌這譴責文士的批判精神,你在今天的香港會認為誰會當受譴責?

當時,耶穌對文士譴責時正好出現一個情景,就是寡婦投錢作奉獻(另可參看:路21:1-4),耶穌善借機緣講道:

「41耶穌面向聖殿銀庫坐著,看眾人怎樣把錢投入銀庫.有好些財主投了許多錢.

42有一個窮寡婦來,投了兩個小文錢,就是一個大文錢.

43耶穌叫門徒來,對他們說:“我實在告訴你們,這窮寡婦投入銀庫裡的比眾人所投的更多.

44因為,眾人都是拿有餘的捐獻,但這寡婦,雖然自己不足,卻把她一生所有的全都投進去了.”」

如果放在今天的香港,你會想起誰好像這些寡婦?

若根據耶穌這讚賞寡婦的憐憫精神,你在今天的香港會認為誰會當受讚賞?

我在香港中文大學教育學院作教師培育工作快二十一年,九年前開始在崇基學院神學院作兼任老師,為強化基督教教育而開課.

我希望將來有一天,能為神學院開一門的課,稱為「耶穌教學法」(Pedagogy of Jesus).就是借用四部福音書裡如何記載耶穌的言行而提煉出一套教學方法,相信「耶穌教學法」會非常精彩,也值得我們教師,家長,信徒領袖以及教牧同工效法.

將來的事,將來再說.

從馬可福音第12章38至44節這段經文,我們已經可以看到耶穌教學法的精神(spirit)與技法(technique):心與力.

首先,道的教學是無法中立的,耶穌是不會中立的,因為耶穌要問的是你的心,包括你的信心與愛心.

你的心若是惡的,耶穌會直斥其非,而且是公開的,不留情面的.耶穌會直斥你的惡行,就算是尊貴的階級比如文士,虛偽的惡行他不留情面地公開批評.

若根據耶穌這譴責文士的批判精神,你在今天的香港會認為誰的虛偽惡行當受不留情面地公開批評?

其次,道的教學是不會離地的,耶穌是直接引例,借眼前的當下機緣而說明要講的道.當下的境況:

「41耶穌面向聖殿銀庫坐著,看眾人怎樣把錢投入銀庫.有好些財主投了許多錢.42有一個窮寡婦來,投了兩個小文錢,就是一個大文錢.」

從這個實況,耶穌引出他要講的一個道理:「有餘的捐獻」與「不足也全投上」.

第二個對比是:有餘的捐獻與不足也全投上

(Second contrast: residual offering vs offering of necessity)

從這個實況,耶穌要說明一個道理,信徒有兩種奉獻的心態:

有些人只是把有餘的捐獻；

有些人雖然不足也全投上.

為讓大家更容易明白,讓我用量化的方式,有點跨張的方式,來放大這個道理.

(一)有餘奉獻者的心態是好像:無限大減去有限細,所餘的依然是無限大.

(二)不足奉獻者的心態是好像:有限細減去有限細,所餘的變成一無所有.

用百份比數一數,前者,是少於千萬份之一；後者,是近乎百份之百.

因此耶穌說:後者比前者奉獻更多:「43b這窮寡婦投入銀庫裡的比眾人所投的更多」.

對於一個千億萬元身家的富豪,要他捐一萬幾千元,這樣的慈善,連利息都不及,其實無損於他的財富.無限大減去有限細,所餘的依然是無限大.

對於一個在貧窮線下,依靠綜合援助維生的赤貧,要他捐一千幾百元,這樣的善行,連生計都不能維持,等同他的財富全部.有限細減去有限細,所餘的變成一無所有.

在座的,是比上不足,比下有餘,你奉獻時,捐獻時是那一種心態?

曾經有一個大學教師,他常與學生飲茶食飯,邊食邊教,堅持在生活中進行活的教學.

每次「埋單找數」,都由他結賬.有些同學「大安旨意」,有些同學表白不安(「唔好意思」).

這位老師堅持做法.一來是他的收入可觀,請客尚算有餘；二來他感到學生或有不足,不想學生因怕付費而卻步,失去學習機會.

然而,請客多了,漸漸地,他也感到支出吃力.繼續呢?還是停止呢?

每次因付出而「肉痛」,他都回想起讀大學時的一位老師.一次一位,一次痛想起一位老師,如是者,他回想起所有恩師.

每段回想,都重新發現當時自己做大學生時都是老師請客.原來,他就是在老師恩待之下邊食邊學地成長,才會有今天.

何以報恩?大恩難言謝,利息滿人間.

所以,他從此立願,把收入的十份之一,用在學生身上,作每月的十一奉獻.

你呢?你回想一下:

你是在怎樣的被恩待下成長,才有今天?

你會怎樣地恩待他人,作為對上帝的回應?

第三個對比是:「多餘的信仰」與「信仰的充足」

(Third contrast: redundant faith vs faith of sufficiency)

知我者,知道我從大學本科到博士完成都是修讀哲學的.十四年的哲學訓練,發現其中有兩個觀念是非常基本的,就是「必要」(necessary)和「充足」(sufficient).

例如,在邏輯學裡,我們會說:A是B的必要條件.C是D的充足條件.

若然A是B的必要條件,那麼,沒有了A,B就無法成立.

若然C是D的充足條件,那麼,有了C,D就足以成立.

讓我舉例說明一下:

若然空氣是我們生存的必要條件,那麼,沒有了空氣,我們的生命就無法維持.

若然清水是我們生存的必要條件,那麼,沒有了清水,我們的生命就無法維持.

但汽水呢?

汽水是不是我們生存的必要條件?若不是,那麼,沒有了汽水,我們依然可以生存.

汽水不單不是我們生存的必要條件,而且喝多了有害於我們的生命.

金錢財富呢?

金錢財富是不是我們生存的必要條件?

若是,那麼,沒有了金錢財富,我們就無法生存.

金錢財富會不會好像汽水一樣,其實儲存多了有害於我們的生命?

除非你沒有清水可喝,唯有以汽水暫代.

除非你無法生存,金錢財富也就是:夠,就夠.

馬可福音第12章38至44節會否在告誡我們:

「41b…….有好些財主投了許多錢.…….43b這窮寡婦投入銀庫裡的比眾人所投的更多.44a因為,眾人都是拿有餘的捐獻……」

若然把有餘的捐獻,所捐獻的這些金錢財富其實是你不必要的(unnecessary),即「毫不入肉的」(不令人肉痛的),那就是多餘的(residual).

這樣的信仰,在你的人生也是多餘的(redundant),是可有可無的.

(廣東話曰:「佢信耶穌?sai 氣!」)

反之,若然你把不足的(insufficient)也投入,所投入的這些金錢財富其實是你必要的(necessary),即「十分食入肉的」(令人肉痛的),那就是所捐獻的實在會影響生活,甚至會影響生存.

這樣的信仰,在你的人生才真正地充足,因為你的信仰是生存的充足條件,那麼,沒有了這些捐獻的金錢財富,生命依然足以成立.

若然基督信仰是我們生命的充足條件,那麼,有了上帝的信實,我們人生在世就足以成立.

充足,其他就不必要,就算於卑賤,於貧窮,於困乏,我因信我的神我就可以安身立命(my faith suffices).相關的神學觀念是Providence of God,即神會供應足夠給我們.

這就是耶穌讚賞窮寡婦的道理:「雖然自己不足,卻把她一生所有的全都投進去了.」

正如另一段經課列王記上17:8-16提到,以利亞對撒勒法的寡婦表示:

「14因為耶和華－以色列的神如此說:『罈內的麵必不用盡,瓶裡的油必不短缺,直到耶和華使雨降在地上的日子.』」

15婦人就照以利亞的話去做.她和以利亞,以及她家中的人,吃了許多日子.

16罈內的麵果然沒有用盡,瓶裡的油也不短缺,正如耶和華藉以利亞所說的話.」

今日我想同大家分享的,是:「有餘與不足,照見信仰的多餘與充足」.

(Unnecessariness and insufficiency reveal the unnecessariness and sufficiency of our Faith.)

祈禱結束 (Ending prayer)

讓我借詩篇146章1至10節作為禱文,讓我們一起低頭祈禱結束:

1哈利路亞!

我的心哪,你要讚美耶和華!

2我一生要讚美耶和華!

我還活著的時候要歌頌我的神!

3你們不要倚靠君王,不要倚靠世人,

他一點也不能幫助.

4他的氣一斷,就歸回塵土,

他所打算的,當日就消滅了.

5以雅各的神為幫助,

仰望耶和華－他神的,這人有福了!

6耶和華造天,地,海和其中的萬物,

他守信實,直到永遠.

7他為受欺壓的伸冤,

賜食物給饑餓的人.

耶和華釋放被囚的,

8耶和華開了盲人的眼睛,

耶和華扶起被壓下的人,

耶和華喜愛義人.

9耶和華保護寄居的,扶持孤兒和寡婦,

卻使惡人的道路彎曲.

10耶和華要作王,直到永遠!

錫安哪,你的神要作王,直到萬代!

哈利路亞!

但願:榮耀歸於聖父,聖子,聖靈,三一的神.阿們.



\newpage

\section{那日子臨近}
\label{sec:304}
\textbf{林豪恩宣教師 - 那日子臨近}
\newline
\newline
\hyperref[sec:303]{< < < PREV SERMON < < <}
~
\hyperref[sec:toc]{[返目錄]}
~
\hyperref[sec:305]{> > > NEXT SERMON > > >}


希伯來書 10:11-14,19-25
\newline
\begin{longtable}{cl}
\hline
\hline
章節 & 經文 (和合本修訂版)\\
\hline
10:11 & \begin{tabularx}{0.7\textwidth}{X} 所有的祭司天天站著事奉神，屢次獻上一樣的祭物，這祭物永不能除罪。 \end{tabularx} \\ \\ \relax
10:12 & \begin{tabularx}{0.7\textwidth}{X} 但基督獻了一次永遠有效的贖罪祭，就坐在神的右邊， \end{tabularx} \\ \\ \relax
10:13 & \begin{tabularx}{0.7\textwidth}{X} 從此等候他的仇敵成為他的腳凳。 \end{tabularx} \\ \\ \relax
10:14 & \begin{tabularx}{0.7\textwidth}{X} 因為他僅只一次獻祭，就使那些得以成聖的人永遠完全。 \end{tabularx} \\ \\ \relax
10:15 & \begin{tabularx}{0.7\textwidth}{X} 聖靈也對我們作證，因為他說過： \end{tabularx} \\ \\ \relax
10:16 & \begin{tabularx}{0.7\textwidth}{X} 「主說：那些日子以後，我與他們所立的約是這樣的：我要將我的律法放在他們的心上，又要寫在他們的心思裡。」 \end{tabularx} \\ \\ \relax
10:17 & \begin{tabularx}{0.7\textwidth}{X} 並說：「他們的罪惡和他們的過犯，我絕不再記得。」 \end{tabularx} \\ \\ \relax
10:18 & \begin{tabularx}{0.7\textwidth}{X} 這些罪過既已蒙赦免，就不用再為罪獻祭了。 \end{tabularx} \\ \\ \relax
10:19 & \begin{tabularx}{0.7\textwidth}{X} 所以，弟兄們，既然我們靠著耶穌的血得以坦然進入至聖所， \end{tabularx} \\ \\ \relax
10:20 & \begin{tabularx}{0.7\textwidth}{X} 是藉著他給我們開了一條又新又活的路，從幔子經過，這幔子就是他的身體。 \end{tabularx} \\ \\ \relax
10:21 & \begin{tabularx}{0.7\textwidth}{X} 既然我們有一位偉大祭司治理神的家， \end{tabularx} \\ \\ \relax
10:22 & \begin{tabularx}{0.7\textwidth}{X} 那麼，我們該用誠心和充足的信心，同已蒙潔淨、無虧的良心，和清水洗淨了的身體來親近神。 \end{tabularx} \\ \\ \relax
10:23 & \begin{tabularx}{0.7\textwidth}{X} 我們要堅守所宣認的指望，毫不動搖，因為應許我們的那位是信實的。 \end{tabularx} \\ \\ \relax
10:24 & \begin{tabularx}{0.7\textwidth}{X} 我們要彼此相顧，激發愛心，勉勵行善； \end{tabularx} \\ \\ \relax
10:25 & \begin{tabularx}{0.7\textwidth}{X} 不可停止聚會，好像那些停止慣了的人，倒要彼此勸勉，既然知道那日子臨近，就更當如此。 \end{tabularx} \\ \\
[1ex]
\hline
\hline
\end{longtable}
馬可福音 13:1-8
\newline
\begin{longtable}{cl}
\hline
\hline
章節 & 經文 (和合本修訂版)\\
\hline
13:1 & \begin{tabularx}{0.7\textwidth}{X} 耶穌從聖殿裡出來的時候，有一個門徒對他說：「老師，請看，這是多麼了不起的石頭！多麼了不起的建築！」 \end{tabularx} \\ \\ \relax
13:2 & \begin{tabularx}{0.7\textwidth}{X} 耶穌對他說：「你看見這些宏偉的建築嗎？這裡將沒有一塊石頭會留在另一塊石頭上而不被拆毀的。」 \end{tabularx} \\ \\ \relax
13:3 & \begin{tabularx}{0.7\textwidth}{X} 耶穌在橄欖山上，面向聖殿坐著；彼得、雅各、約翰和安得烈私下問他說： \end{tabularx} \\ \\ \relax
13:4 & \begin{tabularx}{0.7\textwidth}{X} 「請告訴我們，甚麼時候有這些事呢？這一切事將成的時候有甚麼預兆呢？」 \end{tabularx} \\ \\ \relax
13:5 & \begin{tabularx}{0.7\textwidth}{X} 耶穌說：「你們要謹慎，免得有人迷惑你們。 \end{tabularx} \\ \\ \relax
13:6 & \begin{tabularx}{0.7\textwidth}{X} 將有好些人冒我的名來，說『我是基督』，並且要迷惑許多人。 \end{tabularx} \\ \\ \relax
13:7 & \begin{tabularx}{0.7\textwidth}{X} 當你們聽見打仗和打仗的風聲，不要驚慌；這些事必須發生，但這還不是終結。 \end{tabularx} \\ \\ \relax
13:8 & \begin{tabularx}{0.7\textwidth}{X} 民要攻打民，國要攻打國，多處必有地震、饑荒。這都是災難的起頭。 \end{tabularx} \\ \\
[1ex]
\hline
\hline
\end{longtable}
講題:那日子臨近The Day Approaching

經文:希伯來書10:11-14,19-25;馬可福音13:1-8

講員:林豪恩先生

場合:崇基學院禮拜堂主日崇拜

日期:2015年11月15日

一．引言

那日子臨近,你會有甚麼反應呢?

2011年8月15日(星期一)下午二時,記者目擊約五十名信徒在元朗西鐵站大堂內集合,大部分都是十至二十歲的年輕人,他們明顯有備而來,都揹了個爬山大背囊,外加一張露營地蓆.他們大多都穿戴印有「Jesus」,「耶穌」等字樣的T恤或帽子.約二時半,眾人分別乘坐巴士或的士前往大帽山郊野公園,記者混入其中,發現他們下車後,便有默契地在停車場集合.當時已有近百名信徒在場,其後不少人陸續到達,有人拖篋,有人帶同愛犬走難,部分人更搬了幾箱礦泉水和即食麵上山.原來,當天凌晨時份,他們陸續收到教會高層的電郵或手機短訊緊急通知,隸屬不同小組的信徒,應依照之前所訂下的路線,逃到大帽山或飛鵝山頂,甚至可避走東莞或廣州.(東週刊)

今天的福音經課,耶穌和門徒也談及那日子臨近的話題.那日子臨近,應該如何反應呢?

二．釋經與應用

那日子的迷惑(馬可福音)

13:1耶穌從殿裡出來的時候,有一個門徒對他說:「夫子,請看,這是何等的石頭!何等的殿宇!」

13:2耶穌對他說:「你看見這大殿宇嗎?將來在這裡沒有一塊石頭留在石頭上,不被拆毀了.」

13:3耶穌在橄欖山上對聖殿而坐.彼得,雅各,約翰,和安得烈暗暗地問他說:

13:4「請告訴我們,甚麼時候有這些事呢?這一切事將成的時候有甚麼預兆呢?」

13:5耶穌說:「你們要謹慎,免得有人迷惑你們.

13:6將來有好些人冒我的名來,說:『我是基督』,並且要迷惑許多人.

13:7你們聽見打仗和打仗的風聲,不要驚慌.這些事是必須有的,只是末期還沒有到.

13:8民要攻打民,國要攻打國；多處必有地震,饑荒.這都是災難的起頭.就更當如此.

這些事是必須有的(this is routine history)

門徒問耶穌:

13:4「請告訴我們,甚麼時候有這些事呢?這一切事將成的時候有甚麼預兆呢?」

耶穌回答所提及的現象:戰爭,地震,餓荒,成為了後世討論「那日子臨近」的熱門指標.你感到戰爭越來越多嗎?你感到天災越來越頻密嗎?你感到「那日子」越來越接近了嗎?想到這些現象,你有甚麼反應嗎?你不相信戰爭和天災越來越多嗎?且聽熱心人士所做的資料搜集:

戰爭

踏入二十世紀,戰爭越來越嚴重.第一次世界大戰據報動員了六千五百萬人,死一千七百萬,傷殘三千萬,財物損失三百億美元.在過去的歷史上,你何曾看見這種情況?不久之後,第二次世界大戰爆發,全球都捲入這場空前的大戰中,人命和財物損失更倍數於第一次大戰.單單平民就死了一千八百萬,而金錢損失,最保守的估計是一千億美元.原子彈的威力本應喚醒人心,可是,第二次大戰之後,人還是瘋狂的發動戰爭:五十年代的韓戰,六十年代的越戰,中東,阿富汗,伊拉克等地區的戰事更無日無之．今天呈現在我們眼前的戰爭告訴我們,世界已來到盡頭,主耶穌快再來.(Pastorlau.org)

地震

耶穌說,祂再來之前世界必越來越多地震.那麼這預言是真的嗎?世界果真如此嗎?我們不要說別的,單就二0一0年,世界六級以上的地震就至少有二十四次.我曾請教過幾位地震專家,地震是否越來越頻密,他們的答案都是「是的」.四級地震在這十年中增加了367次 ; 三級的增加了2664次; 二級的增加了3427次.(Pastorlau.org)

聽到這些資料,感到驚慌嗎?感到驚慌就對了,或者正是分享這些資料的朋友所期望的反應.如果人們不驚慌,怎會有危機感呢?怎會產生回應的行動呢?

不過,奇怪的是,當時耶穌提及這些現象時,卻對聽眾說:「不要驚慌」(13:7)．為甚麼不必驚慌呢?他說,「因為這些事是必須有的」(13:7)．「這些事是必須有的」這一句,Eugene Peterson在The Message Bible譯作「this is routine history」.耶穌似乎在說,戰爭,地震,餓荒這些天災人禍,都是歷史的一部份,並非罕世奇事,他呼籲門徒不用驚慌．

1.2. 要迷惑許多人

面對天災人禍這些既混亂又危急的狀況感到驚慌是人之常情吧,耶穌對門徒說,不要驚慌,那又應該如何反應呢?

13:5耶穌說:「你們要謹慎,免得有人迷惑你們.

13:6將來有好些人冒我的名來,說:『我是基督』,並且要迷惑許多人.」

關於「那日子」的論述,我們以為相關的反應應該是「驚慌」,因此,我們自己驚慌,也懷著使命把驚慌傳給其他人.耶穌吩咐門徒的反應是「謹慎」,因為「有人迷惑你們」.迷惑,令人感到模糊混亂,真假難分,對錯難辨．那些迷惑人的人說「我是基督」是甚麼意思呢?就是要那些原本跟隨基督的,或者尋找基督的,轉而跟隨他們．人們跟隨了他們,就等於離開了原來所跟隨的,放棄了原本要走的路,用近年流行的說話,就是忘記了初衷．「那日子臨近」的威脅,不是天災人禍,而是艱難的狀況下所產生那種令人改變,放棄初衷的迷惑.

很多人懷著初衷去開始,這些初衷可能是理想,可能是價值,可能是所追尋的美善．當兩個人結婚的時候,總會是懷著彼此相愛,建立家庭的初衷,不會懷著彼此傷害,拆散家庭的惡念?投身教育工作的人,都懷著作育英才的初衷,而不會是虐待學童的惡念?投身醫護工作的人,都本著濟世為懷的初衷,而不會懷著草菅人命的惡念?忘記初衷,未必是刻意的,而是處身於艱難境況,在時不與我的狀況之下被迷惑的反應.在家庭中,如果不是他／她對待得那那麼差,我才不會對他／她那麼差．在學校中,如果不是學生那麼頑劣,我才不會如此失控．在醫院中,如果不是工作環境那麼惡劣,我才不會如此對待病人．耶穌知道他的門徒和眾生一樣,也是容易被艱難的景況迷惑的,在面對危難的時候,或者挫敗的時候,就是最容易忘記初衷的時候．非常時期用非常手段．是的,本來沒有打算如此,也不想如此,但為勢所逼,唯有如此．

福音書記載了耶穌和門徒一項具體地面對迷惑的事件,路加福音的版本是這樣的:

22:39耶穌出來,照常往橄欖山去,門徒也跟隨他.

22:40到了那地方,就對他們說:「你們要禱告,免得入了迷惑.」

22:41於是離開他們約有扔一塊石頭那麼遠,跪下禱告,

22:42說:「父啊!你若願意,就把這杯撤去；然而,不要成就我的意思,只要成就你的意思.」

當艱難還沒有臨近,人們還可以高言大志.艱難當前,耶穌也想到「把這杯撤去」的可能性.但他沒有入了迷惑,沒有忘記自己的道成內身來到世間的初衷,堅持「不要成就我的意思,只要成就你的意思」．他也如此勸勉跟隨他的人:

22:45禱告完了,就起來,到門徒那裡,見他們因為憂愁都睡著了,

22:46就對他們說:「你們為甚麼睡覺呢?起來禱告,免得入了迷惑!」

22:47說話之間,來了許多人.那十二個門徒裡名叫猶大的,走在前頭,就近耶穌,要與他親嘴.

22:48耶穌對他說:「猶大!你用親嘴的暗號賣人子嗎?」

22:52耶穌對那些來拿他的祭司長和守殿官並長老說:「你們帶著刀棒出來拿我,如同拿強盜嗎?

22:53我天天同你們在殿裡,你們不下手拿我.現在卻是你們的時候,黑暗掌權了.」

黑暗掌權了,門徒能夠不陷入迷惑嗎?對手使用暴力,難道門徒仍要和平理性嗎?對手在暗角行事,難道門徒仍要光明磊落嗎?

22:49左右的人見光景不好,就說:「主啊!我們拿刀砍可以不可以?」

正如上文所說,人們多是懷著良好的初衷,但見光景不好,就難免感到迷惑.要堅持初衷嗎?還是「你做初一我做十五」呢?

22:50內中有一個人把大祭司的僕人砍了一刀,削掉了他的右耳.

面對危難,在光景不好的時候,有人以行動取代思考,以暴易暴．

22:51耶穌說:「到了這個地步,由他們吧!」就摸那人的耳朵,把他治好了.

也是在危機之中,在光景不好的時候,有人毋忘初衷．沒有被暴力把自己變成暴力,沒有被黑暗把自己變為黑暗．以平靜面對威嚇,以和平面對暴力,以慈悲回應傷害．

那日子臨近,威脅不是天災人禍,天災人禍是歷史的一部份；威脅是人們面對光景不好而驚慌失措,入了迷惑.迷惑令人信心動搖,令人看不見前景,令人與人之間失去互信,互相傷害．

二．既知道那日子臨近,就更當如此(希伯來書)

今天的書信經課,希伯來書10章25節說:「既知道那日子臨近,就更當如此」．「更當如此」,即是如何呢?

10:11凡祭司天天站著事奉上帝,屢次獻上一樣的祭物,這祭物永不能除罪.

10:12但基督獻了一次永遠的贖罪祭,就在上帝的右邊坐下了.

10:13從此,等候他仇敵成了他的腳凳.

10:14因為他一次獻祭,便叫那得以成聖的人永遠完全.

10:19弟兄們,我們既因耶穌的血得以坦然進入至聖所,

10:20是藉著他給我們開了一條又新又活的路,從幔子經過,這幔子就是他的身體.

10:21又有一位大祭司治理上帝的家,

10:22並我們心中天良的虧欠已經灑去,身體用清水洗淨了,就當存著誠心和充足的信心來到上帝面前；

10:23也要堅守我們所承認的指望,不致搖動,因為那應許我們的是信實的.

10:24又要彼此相顧,激發愛心,勉勵行善.

10:25你們不可停止聚會,好像那些停止慣了的人,倒要彼此勸勉,既知道那日子臨近,就更當如此．

2.1. 信心

10:22並我們心中天良的虧欠已經灑去,身體用清水洗淨了,就當存著誠心和充足的信心來到上帝面前；

艱難的處境令人或到迷惑,令人信心動搖.就算是本來懷著堅強信念的人,也會陷入疑惑迷亂之中.德蘭修女是1979年諾貝爾和平獎得主,遠道在加爾各答貧民區中服侍最窮苦的人,世人一直以為她是信心巨人.然而,在照顧貧民的時候,許多人毫無希望地死在她手中,她向神禱告,但大部份的禱告是不成就的.她感到非常孤單,精神飽受折磨,甚至無法禱告下去.她在一篇對耶穌的懺悔中說:「我呼求,我緊抓,我要求,卻沒有一人回答,沒有一個我能抓住.沒有,一個也沒有,只有孤獨.我的信心在哪裡? 就算我的內心深處,也是什麼都沒有,只有空虛和黑暗.我的神,這莫名的痛苦真是好大啊!我沒有信心.」

是甚麼狀況動搖你的信心呢?是甚麼事情令你如德蘭修女一樣發出心靈的呼喊:"我沒有信心"呢?是感到和伴侶相處不來的時候嗎?是學業或工作遇到挫折的時候嗎?是局勢不好,時不與我的時候嗎?德蘭修女在信心動搖的時候向上主說:「我沒有信心」,向上主訴說內心深處的困惑和痛苦,反映出她仍信任上主,如果她的內心早已麻木的話,早就放棄向神呼求了.

希伯來書的作者呼籲我們,那日子臨近,就更當「存著誠心和充足的信心來到上帝面前」,這能幫助我們面對艱難,免得入了迷惑.

2.2. 盼望

10:23也要堅守我們所承認的指望,不致搖動,因為那應許我們的是信實的.

縱使是懷著偉大理想的,懷著美好心意的人,在艱難的時候,都會感到迷惑和失望.馬丁路德金牧師懷著為有色人種爭取平等的理想而奔波勞碌,但卻遇到眾多的阻攔．攻擊不單來自反對者,不單來自極端組織對他自身及家人的死亡威嚇,更有來同路人的分裂和對抗．是甚麼支持他堅持下去呢?因為,他有一個夢,這個「到了那一天,上帝的一切兒女,無論黑人,白人,猶太人,非猶太人,基督教徒和天主教徒,都能夠攜手一同慶祝得到自由的遠象．」就是這份盼望,支撐著馬丁路德金牧師,因此,他說:「朋友們,今天我要告訴你們,儘管當前還有許多困難挫折,我仍然懷有一個夢想.」因著有這盼望,「我們能夠把絕望的大山鑿成一塊希望的磐石」.

在過去的幾年間,越有理想的朋友,好像越感到迷茫和挫敗.絕望的大山一座一座地出現在眼前,寸步難行.希伯來書呼籲我們,在光景不好的日子,更當「堅守我們所承認的指望」,因為美善的上主最終都會肯定美善.

2.3. 愛心

10:24又要彼此相顧,激發愛心,勉勵行善.

馬太福音24章12節耶穌又論都那日子的時候說:「因為不法的事增多,許多人的愛心就冷淡了」.當社會失去信任,極端地對立的時候,人們傾向互相攻擊,有我無你.

馬丁路德金牧師在「我有一個夢」演講中如此勉勵民眾:「我們不要用仇恨的苦酒來緩解熱望自由的乾渴.我們必須永遠站在高處,便我們的鬥爭方式保持尊嚴,堅守紀律.我們一定不能使富有創造性的抗爭淪為使用暴力的低下行動.我們必須努力不懈地站在以靈魂力量來對付肉體力量的神聖高度.已經席捲黑人社會的戰鬥氣氛決不要導致我們對一切白種人的不信任.因為我們許多白人兄弟今天到這裡來集會,就證明了,他們已經認識到他們的命運和我們的命運緊緊聯結在一起,他們的自由和我們的自由完全分不開.」

那日子臨近,就更當「彼此相顧,激發愛心,勉勵行善」.無論我們認為自己有多正確,對方有多不該,我們仍然是生活在同一天空之下.我們不想對方把我們趕盡殺絕,我們也不能把對方排除出去.

2011突尼斯發生「茉莉花革命」,變革之火燃燒到北非,各地的人民先後對其政府發出要求改革開放的怒吼.然而,經歷「阿拉伯之春」的一些國家,卻陷入各種力量互相仇視,非要消滅對方或者控制對方不可的惡鬥泥沼.不過,突尼斯卻脫穎而出,取得成功．他們的「全國對話四方組織」因而獲得本年度的諾貝爾和平獎.其實,突尼斯也如其他北非國家一樣,一度陷入內戰的邊緣.為挽救突尼斯陷於內戰,突尼斯總工會倡議組成全國對話機制,邀請其他三個有影響力的民間組織加入:突尼斯工商聯盟,人權聯盟和律師聯合會.難能可貴的是,總工會將突尼斯工商聯盟拉進對話.他們原本是天生對頭--工會代表底層工人,而工商聯盟由他們的僱主組成.這對冤家聯手,四個組織代表不同利益人群,顯示了民間團結的決心.同為宗教政黨,突尼斯復興運動黨也不如埃及穆兄會極端.他們願意交出權力,共接受選舉落敗的結果.黨領袖下台後,四處受邀講演,而短暫出任埃及總統的穆兄會成員穆爾西,卻成為軍方階下囚.

希伯來書作者呼籲人們,在艱難迷惑的時候,更當「彼此相顧,激發愛心,勉勵行善」.在同一天之下,有你,有我,也有他．

三．總結

「那日子臨近」雖然有天災人禍的亂象,然而,歷史上何時沒有天災人禍呢?因此,「那日子臨近」,不是驚慌的時候,而是謹慎不被迷惑的時候；不是放棄的時候,而是更需要堅持信望愛的時候.



\newpage

\section{當挺身昂首}
\label{sec:305}
\textbf{鄧瑞強博士 - 當挺身昂首}
\newline
\newline
\hyperref[sec:304]{< < < PREV SERMON < < <}
~
\hyperref[sec:toc]{[返目錄]}
~
\hyperref[sec:306]{> > > NEXT SERMON > > >}


路加福音 21:25-36
\newline
\begin{longtable}{cl}
\hline
\hline
章節 & 經文 (和合本修訂版)\\
\hline
21:25 & \begin{tabularx}{0.7\textwidth}{X} 「日月星辰要顯出預兆，地上的邦國也有困苦，因海中波浪的響聲而惶惶不安。 \end{tabularx} \\ \\ \relax
21:26 & \begin{tabularx}{0.7\textwidth}{X} 人想到那要臨到世界的事，就都嚇得魂不附體，因為天上的萬象都要震動。 \end{tabularx} \\ \\ \relax
21:27 & \begin{tabularx}{0.7\textwidth}{X} 那時，他們要看見人子帶著能力和大榮耀駕雲來臨。 \end{tabularx} \\ \\ \relax
21:28 & \begin{tabularx}{0.7\textwidth}{X} 一有這些事，你們就當挺身昂首，因為你們得救贖的日子近了。」 \end{tabularx} \\ \\ \relax
21:29 & \begin{tabularx}{0.7\textwidth}{X} 耶穌對他們講了一個比喻說：「你們看無花果樹和各樣的樹， \end{tabularx} \\ \\ \relax
21:30 & \begin{tabularx}{0.7\textwidth}{X} 樹葉一長出來，你們看了自然就知道夏天近了。 \end{tabularx} \\ \\ \relax
21:31 & \begin{tabularx}{0.7\textwidth}{X} 同樣，當你們看見這些事發生，就知道神的國近了。 \end{tabularx} \\ \\ \relax
21:32 & \begin{tabularx}{0.7\textwidth}{X} 我實在告訴你們，這世代還沒有過去，一切都要發生。 \end{tabularx} \\ \\ \relax
21:33 & \begin{tabularx}{0.7\textwidth}{X} 天地要廢去，我的話卻絕不廢去。」 \end{tabularx} \\ \\ \relax
21:34 & \begin{tabularx}{0.7\textwidth}{X} 「你們要謹慎，免得被貪食、醉酒和今生的憂慮壓住你們的心，那日子就忽然臨到你們， \end{tabularx} \\ \\ \relax
21:35 & \begin{tabularx}{0.7\textwidth}{X} 如同羅網一樣，因為那日子要臨到所有居住在地面上的人。 \end{tabularx} \\ \\ \relax
21:36 & \begin{tabularx}{0.7\textwidth}{X} 你們要時時警醒，常常祈求，使你們能逃避這一切要來的事，得以站立在人子面前。」 \end{tabularx} \\ \\
[1ex]
\hline
\hline
\end{longtable}
講題:「當挺身昂首」　Stand Up and Raise your Heads

經文:路加福音21章25-36節

講員:鄧瑞強博士

場合:香港中文大學崇基學院禮拜堂

日期:2015年11月29日

各位弟兄姊妹,早安.

今日是將臨期的第一個主日.將臨期提示我們,有一異常重要的事,即將來臨.這事重要到一個地步,值得我們花一段日子去預備.將臨期提示我們,記念基督降生的日子快到了.這個日子很重要,基督的降生,改變了人類的命運.不單如此,將臨期也提示我們,要認真等待基督再來的事實.

在日常生活裡,我們都試過,等待某些重要日子的來臨.找工作的人,在面試的日子臨到之前,會多番準備自己,好讓自己在面試時有好的表現.第一次約會心儀對象的人,會在第一次約會日子前,多番準備,好能在對方心裡,留下美好印象.準備結婚的人,可能在結婚的日子一年之前,已作出種種準備.要減肥,要瘦身,要讓身體在最佳狀態,要訂造最漂亮的禮服,要在適當時間通知各方人士,要安排婚禮如何進行,要預計婚宴時的細節,等等等等.這是對重要事件的真正預備.面對重要日子,我們很少毫無準備,讓那重要的日子忽然臨到似的.在準備自己的時候,將來的事件會變成現在生活的內容.現在每一刻的生活,都滲透著將來的事件.每一刻的思想,都被將來的事件牽動.好像當下的每刻,就是為了預備一切,好迎向那重要的一刻.

將臨期,就是提示我們,有些重要事情即將出現.這事情很重要,重要到值得我們傾盡全力去作準備.耶穌的再來,就是這件重要的事.

今日講道的經文,是今日的福音書經課:路加福音21:25－36.

路21:25　「日,月,星辰要顯出異兆,地上的邦國也有困苦；因海中波浪的響聲,就慌慌不定.

路21:26　天勢都要震動,人想起那將要臨到世界的事,就都嚇得魂不附體.

路21:27　那時,他們要看見人子有能力,有大榮耀駕雲降臨.

路21:28　一有這些事,你們就當挺身昂首,因為你們得贖的日子近了.」

路21:29-30　耶穌又設比喻對他們說:「你們看無花果樹和各樣的樹；它發芽的時候,你們一看見,自然曉得夏天近了.

路21:31　這樣,你們看見這些事漸漸地成就,也該曉得神的國近了.

路21:32　我實在告訴你們,這世代還沒有過去,這些事都要成就.

路21:33　天地要廢去,我的話卻不能廢去.」

路21:34　「你們要謹慎,恐怕因貪食,醉酒,並今生的思慮累住你們的心,那日子就如同網羅忽然臨到你們；

路21:35　因為那日子要這樣臨到全地上一切居住的人.

路21:36　你們要時時警醒,常常祈求,使你們能逃避這一切要來的事,得以站立在人子面前.」

(1)這世界搖搖欲墜

舊約時代和新約時代中間,有一段約400年的時間.在這「兩約中間」時期,發展了一種特殊文體,一種稱為「天啟文學」(apocalyptic literature)的文體.這種文體的其中一個特徵,就是將人間世界放在一宇宙視域中.這宇宙視域揭露神在萬有之上.在這宇宙視域下,人間世界顯得渺小,短暫,沒有終極意義.

有了以上的理解基礎,我們看看剛才讀的經文的頭兩節.

路21:25　「日,月,星辰要顯出異兆,地上的邦國也有困苦；因海中波浪的響聲,就慌慌不定.

路21:26　天勢都要震動,人想起那將要臨到世界的事,就都嚇得魂不附體.

經文希望我們在一更大的視域中,看到這個世界的真相.

我們以為,人間的權力結構是穩固的.我們以為,經濟的成果有終極意義.我們以為,世上的功業是永恆的.但經文說,在一更大視域下,便看到這一切毫不穩固.日,月,星辰都要發出信息,天地間的一切東西都站不住腳.「天勢都要震動」意指天上的各種勢力都被動搖.世上的列國,不能因自己的強勢而站穩.因著天地間一切被視為穩當的東西都被推倒,列國就「慌慌不定」.那要臨到世界的事,就是宇宙的真相的揭露,神的終極臨在.面對這終極真相,人便「嚇得魂不附體」.

我們面對的這個世界,很不完美.你派人肉炸彈來殺我,我便派戰機去炸你.難民走過來,我便築起圍牆來擋住.你試圖為經濟體系製定新的規則,我便用強大的經濟體系來圍堵你.你打了我一巴掌,我要了你條命.每個國家都試圖稱霸,每個人都試圖稱王,每個人每日的生活都在競爭之中.每個人都從過去的戰爭中吸取了沉重的教訓,每個人卻沉重地預備著另一場戰爭.

主耶穌的說話,是指出我們這個世界不是終極的,世間的權力不是終極的,人間的痛苦也不是終極的.將我們的世界放在宇宙的視域裡,便會發覺我們的人世渺少得可憐.無論我們如何努力將自己的歷史看為是歷史的中心,在字宙的視域裡,我們的所謂偉大,只是輕煙一縷,微不足道.

人的眼光若短淺,則會在塵世的名利追逐中,虛度一生.

比人間世界更高層次的天,比人間權力更高層次的天界力量,也在顫抖,也在崩潰,何況是人間世界呢?

那令天地震動的東西,是什麼呢?

(2)人子大有能力地降臨

路21:27　那時,他們要看見人子有能力,有大榮耀駕雲降臨.

路21:28　一有這些事,你們就當挺身昂首,因為你們得贖的日子近了.」

令天地震動的,是一將來事件.

這將來事件,也就是將臨期提示我們等待的事件.

這將來事件,就是「人子」大有能力,大有榮耀地駕雲降臨.

「人子」駕雲降臨這意象,來自舊約但以理書7章13節.這意象也是天啟文學常有的意象.「人子」是神性人物,他從天而來,有神的能力,去主持人間正義.他要消滅邪惡,解放被欺壓的人.他來了,要揭露宇宙人生的真相.顯明人間邪惡的虛妄,顯明神的真實.真理揭露之時,也是萬物被神的正義審判之時.對邪惡的人來說,「人子」是他們的判官.對受苦的人來說,「人子」是他們朝夕等待的救主.

這裡的「人子」,明顯是指耶穌.但重點不在當時仍在世上的耶穌,而是指那死後復活而升入高天的耶穌.

耶穌死而復活,是表明神那審判罪惡並勝過罪惡的力量.耶穌的復活,是展示與這個世界不一樣的另一個世界,一個從神而來的世界.「人子」駕雲降臨,意味著我們今天要等待的明天,要等待的事件,是一件對抗我們現今世界一切邏輯的事件.我們的世界的邏輯,是死去了便是死去了,失去了便是失去了,失敗了便是失敗了.故此,每個人都會竭力爭勝,置對方於死地.若被對方擊敗了,便會失去明天.但是,「人子」駕雲降臨,正正說明,那被釘十架的主,被神復活了.他站在人類歷史的終點,帶著審判的大能,向我們走過來.在他的復活裡,顯明在人類的生死之上,有一主持正義的神.為義受死者,他的死亡不是終點,不是失敗,因此也不必絕望.

正因為「人子」駕雲降臨,我們明白,人類歷史的明天,不是由拿著殺人武器殺盡一切人的勝利者確立的,而是將這暴力擊倒的神聖正義力量確立的.新的世界,不是由將人釘十架的力量帶來的,而是由肯定被釘十架者的復活力量確立的.

若果沒有這個世界以外的另一個世界,若果沒有這個世界的生活邏輯以外的另一種邏輯,若果這個世界就是一切而沒有另一個具有神聖意義的世界,則我們只是這個搖搖欲墜的世界的囚徒,只能將生命壓注在這個世界的遊戲規則裡.「將臨期」是叫我們等待那即將駕雲降臨的「人子」,他給我們帶來一個新的世界.就因為有這新的世界,我們可以審視這個世界的邪惡,拒絕這個世界的潛規則,甚至對抗這個世界.這就是「活在世界而不屬世界」的意義.

在世上,信徒是一群特殊的人.我們在今日,期待著這件重要的事情發生.我們向前展望,在信心的眼睛裡,看見那即將來臨的主.我們是迎向前去,迎接救主的再臨.

任何人,只要認真地期待一件重要的事情來臨,他的日子便會和這件事情息息相關.若你計劃某日結婚,你便會不斷計劃,在什麼日子之前要找牧師作婚前輔導,在什麼日子之前要減去多少磅,在什麼日子之前要變成什麼模樣.每日的生活,看來和這件重要的事,是無法分開的.

認真等待主耶穌再來的人,也應如此.生活的每一日,都和這等待有關.生命的每個環節,都滲透著人子再來的氣息.人子再來,是帶來一全新的世界.在等待時,信徒會特別留意到這個世界的腐敗和墮落.人子再來,是將神的正義帶來.在等待時,信徒也會感覺到,這世界的不義和神的正義,有一強烈對比.人子再來,向我們展示一個比現世大得多的神聖世界.等待主再來的信徒,心靈也會因這等待而擴闊,能明白天國的意義.

等待人子再來的信徒,會為了抵抗世界的沉淪而有一忠勇奮鬥的精神.但要奮勇上進,迎向基督,是艱險莫名的.世上所有令人沉淪的力量,威嚇人的力量,置人於死的力量,所有釘耶穌上十字架上的力量,全都起來反對你,威嚇你,誘惑你.世界會告訴你,這個宇宙是沒有神的,強權就是正義,暴力就是道理.你在這遊戲中失敗了,便只有死路一條.

期望主的再來,令信徒生命改變.這生命改變,是不甘於同流合污,不甘於成為釘耶穌上十架的幫兇.如此的生命,是艱難的.主耶穌說:「一有這些事,你們就當挺身昂首,因為你們得贖的日子近了.」在艱難時刻,要想起主的復活,並復活主的再來.在復活主的力量下,挺身昂首,站起來,堅持仁愛和正義.不要怕,因為得贖的日子近了.

這樣的等待,會否落空?

會不會好像「等埋發叔」那樣,到頭來「得個桔」?

主耶穌說:「天地要廢去,我的話卻不能廢去」.當然,主耶穌以他的復活來證明他有超越於天地的大能.

(3)拒絕這世界的生活方式

等待,不是坐著等,而是將等待的事情變成生活的方式.若你等待復活的主再來,那麼,你的生活便要有新生命的氣息.若你等待的主是審判罪惡的主,那麼,你的生活便要光明正大.若你等待的,是一個新世界的展現,則你的生活便要與這世界的生活方式不同.

主耶穌說:

路21:34　「你們要謹慎,恐怕因貪食,醉酒,並今生的思慮累住你們的心,那日子就如同網羅忽然臨到你們；

路21:35　因為那日子要這樣臨到全地上一切居住的人.

路21:36　你們要時時警醒,常常祈求,使你們能逃避這一切要來的事,得以站立在人子面前.」

所謂「貪食」,即縱慾；「醉酒」,即朝生暮死而失去自控；「今生的思慮」,即種種令我們對神的真理沒有興趣的東西.這三者合起來,便是塵世生命的基本內容.

這一切,都與「人子」所帶來的新生命和新世界格格不入的.

所謂「警醒」,就是意識到自己不屬這個世界,而屬於那從天而來的世界的.而「常常祈求」,就是心中時常以神的事為念.

若真的等待主的再來,便要讓再來的主充滿我們的生命,以致我們有力量放下這世界的生活方式.

我們的生活方式,告訴我們是否真正等待主的再來.

等候主再來的人,他活在再來的主的那個世界裡.

小時候,父母在外面工作時,我們兄弟姊妹在家裡,便會搬動傢具,移空場地來玩耍.玩耍之餘,也常打架.當父母回來的時間近了,我們便會還原一切.若是打過架了,便會抹乾自己和對方的眼淚,和對方和解,表現出親切的模樣.父母回來了,見一切如常,各人友愛,便很欣慰.

意識到父母回來,我們改變了,變得友愛.

意識到主耶穌的再來,我們又會作什麼改變呢?

但願:榮耀歸於聖父,聖子,聖靈.阿們.



\newpage

\section{應當一無掛慮}
\label{sec:306}
\textbf{鄧瑞強博士 - 應當一無掛慮}
\newline
\newline
\hyperref[sec:305]{< < < PREV SERMON < < <}
~
\hyperref[sec:toc]{[返目錄]}
~
\hyperref[sec:307]{> > > NEXT SERMON > > >}


腓立比書 4:4-7
\newline
\begin{longtable}{cl}
\hline
\hline
章節 & 經文 (和合本修訂版)\\
\hline
4:4 & \begin{tabularx}{0.7\textwidth}{X} 你們要靠主常常喜樂。我再說，你們要喜樂。 \end{tabularx} \\ \\ \relax
4:5 & \begin{tabularx}{0.7\textwidth}{X} 要讓眾人知道你們謙讓的心。主已經近了。 \end{tabularx} \\ \\ \relax
4:6 & \begin{tabularx}{0.7\textwidth}{X} 應當一無掛慮，只要凡事藉著禱告、祈求和感謝，將你們所要的告訴神。 \end{tabularx} \\ \\ \relax
4:7 & \begin{tabularx}{0.7\textwidth}{X} 神所賜那超越人所能了解的平安，必在基督耶穌裡，保守你們的心懷意念。 \end{tabularx} \\ \\
[1ex]
\hline
\hline
\end{longtable}
講題:「應當一無掛慮」　Be Anxious for Nothing

經文:腓立比書4:4-7

講員:鄧瑞強博士

場合:香港中文大學崇基學院禮拜堂

日期:2015年12月13日

各位弟兄姊妹,早安.

將臨期有四個主日,一般而言,第一周的主題是「盼望」,第二周的主題是「平安」,第三周的主題是「喜樂」,第四周的主題是「愛」.今周的主題是「喜樂」.

談論「喜樂」,越來越艱難,不同人指出,我們這個時代的特徵是「憂慮」.聖經說:「不要為明天憂慮」.但現實是:我們每刻都為明天憂慮.

教宗方濟指出,巴黎的恐怖襲擊可看為是第三次世界大戰的一小部分.在他看來,第三次世界大戰的形式是漫長地,零碎地打,縱使你只是一個正在聽演唱會的平民,仍會牽涉其中.巴黎的恐怖襲擊預示著這場戰爭的開始.若真是這樣,則無論你身處何方,都有可能成為一種新型態的戰爭的受害者.世界如此,能不憂慮嗎?

英美詩人奧登(Wystan Hugh Auden),寫了首長詩＜憂慮的時代＞(The Age of Anxiety),宣告我們進入了一個「憂慮的時代」.有幾行詩句如此說:

我們寧願毀滅,都不甘願被改變,

我們寧願死在自己的憂慮中,

也不願意屈身爬過時代的轉折,

而讓我們的幻想死去.

We would rather be ruined than changed

We would rather die in our dread

Than climb the cross of the moment

And let our illusions die.

這是令人驚奇的,面對戰爭,我們都不願意改變自己,放下仇恨,我們寧願同歸於盡.我們不願意放下驕傲的幻想,寧願死在自己製造出來的憂慮裡.

現代社會的生活方式,也可能造成種種個人憂慮.我們越來越個人化,越來越強調自我的存在意義.這原是好的,讓我們每個人都體會存在的價值.但很多時,強調自我便抗拒社群.社群是價值體系的傳道體,離開社群,便失去意義的承托.沒有意義的承托,面對生命危機時,便容易憂慮.

現代社會的生活,強調控制.我們要控制面上皺紋第一次出現的時間,我們要嚴格計劃明天.主耶穌講過一個故事.他說:「有一個財主,心裡說:『要把我的倉房拆了,另蓋更大的,在那裡好收藏我一切的糧食和財物』.這財主對自己的靈魂說:『靈魂哪,你有許多財物積存,可作多年的費用,只管安安逸逸的吃喝快樂吧!』神卻對他說:『無知的人哪,今夜必要你的靈魂；你所預備的要歸誰呢?』」(參路12: 16-20)要控制,便要控制明天.但明天之為明天,就是因為明天仍未出現,在我們的控制之外.你要試圖控制一種在你控制力以外的東西,你總會憂慮.

神從「無」(nothingness)中「創造」萬有,也從「無」中創造人.人的存在深處,有一種虛無性(non-being).神學家田立克(Paul Tillich)說:當我們體會到這種「虛無」力量時,我們便會「憂慮」.在生命存在的層面上,這「虛無」以死亡之力威脅我們.在心靈意義的層面上,這「虛無」以「空虛」及「意義的失落」來威脅我們.在道德成就的層面上,這「虛無」以「罪無可恕」(condemnation)的道德重量威脅我們.

在茫茫天地間,人感到自身內裡的「虛無」力量,去取消我們的道德努力,意義尋索,最後取消我們的生命存在.人力圖克服這「虛無」之力,但有限的人只會力不從心,憂慮不已.

神以創造的大能,從「虛無」中創造我們,也只能憑神的創造大能,去克服這「虛無」的威脅,驅散憂慮.離開神,我們自身的「虛無」,又會吞噬我們.

今日講道的經文,是:腓立比書4:4－7.主題是「應當一無掛慮」.

腓4:4　你們要靠主〔或譯:在主裡〕常常喜樂.我再說,你們要喜樂.

腓4:5　當叫眾人知道你們謙讓的心.主已經近了.

腓4:6　應當一無掛慮,只要凡事藉著禱告,祈求,和感謝,將你們所要的告訴神.

腓4:7　神所賜,出人意外的平安必在基督耶穌裡保守你們的心懷意念.

在這憂慮的時代,這經文作出很好的提示.「應當一無掛慮」,就是:面對任何事都不要憂慮.我們留意到,這段經文的每一節,都顯示和神的緊密關係.沒有神,我們很難面對憂慮.

這段經文的開始,其實是一個「命令」.

腓4:4　你們要靠主〔或譯:在主裡〕常常喜樂.我再說,你們要喜樂.

「你們要喜樂」,「再說,你們要喜樂」,兩次的動詞,用的都是命令語氣.你們務要喜樂,this is an order,不得猶疑,不得討價還價.在憂慮的年代,人自自然然地便會陷入憂慮,要你不憂慮,要你喜樂,是反常態,必須下命令.你可能說,人有悲歡離合,總有傷感時刻,命令人要喜樂,是否不近人情?聖經當然不會控制人的感情流露.這裡講的喜樂,建基於一種生命最根本的信念,這信念就是:縱然人間「虛無」的力量不斷威脅我的存在,我仍然堅信神創造性的大能與恩典托著我.死亡的「虛無」力量威脅我,叫我出賣靈魂,但我仍然堅信,神復活的力量能克服我的死亡.「空虛」及「意義的失落」的「虛無」力量威脅我,叫我放棄價值的追尋,但我仍然堅信,神是意義和價值的泉源.「罪無可恕」的道德重量威脅我,叫我放棄自己,朝生暮死,但我仍然堅信,神會以祂的無限寬恕去接納我這個罪大惡極的人.就因為這信念,就因為這信念背後的那位無限慈愛的神,我們無理由在「虛無」面前屈服,反而在「虛無」裡,深深地依靠著那使「無」變「有」的神.這種堅信,這種生命,就是「喜樂」.保羅命令我們,要有這種生命內容.其實,這種生命內容,就是基督徒應有的生命內容.故此,保羅的意思是:要喜樂,因為這根本就是基督徒的生命樣式.

保羅命令我們,要喜樂,還要「常常」喜樂.要能做到這樣,必需生命的轉化.要能克服憂慮,要能克服「虛無」,不能靠自己,因為「虛無」存在於我們生命的深處.要能「常常」喜樂,必需生命的轉化.轉化的方向是,活「在主裡」.你活著,不再是自己活著,而是基督在你裡面活著.要「在主裡」,才能常常喜樂.「在主裡」,就不是在自己裡.以主耶穌對天父的信靠,作為自己對天父的信靠.以主耶穌那捨己精神,來鼓勵自己不要常常想著自己.以主耶穌對人的關懷,來強化我們對他人的愛.生命有所依,放下自我,心中有愛,才能有真喜樂.

一個人若有真喜樂,其他人是感受得到的.

腓4:5　當叫眾人知道你們謙讓的心.主已經近了(或者譯:主近了).

喜樂的人最凸出的特質,是「謙讓」.什麼是「謙讓」(希臘文是epieikes)?聖經學者巴克萊(William Barclay)說,這個字是其中一個最難翻譯的聖經用字.這個字包含忍耐,克制,寬容,豁達,不執著於刻板的規條,公正卻又仁慈等意思.這個字用來描繪神對罪人的寬容,是最合適的.神不輕易發怒,且有豐盛的慈愛,這就是「謙讓」.我將「謙讓」翻譯為「寬宏大量」.一個喜樂的人,一個放下憂慮的人,他對自己或他人,總是「寬宏大量」的.

能「寬宏大量」,是因為「主已經近了」.「主已經近了」,可以就空間說,也可以就時間說.我們能寬宏大量,因為主就在我們身旁.我們能寬宏大量,因為主耶穌即將再來.人生有憂慮而不得喜樂,是因為我們面對「虛無」的威脅.「虛無」可能來自道德上的無能為力.我們體會自身的無藥可救,他人的罪無可恕.但主耶穌就在身旁,他會接納我們,也鼓勵我們接納他人.並且,既然主再來的日子近了,那就將審判大權交還上帝,讓神去審判而不是我們自己去審判.「主已經近了」,讓我們從自己或他人的道德墮落中,找到恩典的根源,從而有喜樂的可能.又或者,「虛無」來自生命存在的即將逝去,死亡的黑暗威脅著我們.在死亡面前,人一切的名利功業,都即將歸回空無.死亡的憂慮,會叫人瘋狂.「主已經近了」,是克服死亡的空無的最大力量.從空間上說,主就在身旁,是最大的解救.就像一個人墮入空無的深淵時,上帝之手在身旁抓住他.從時間上說,主即將來臨,是最大的盼望,因為主再來時,會將復活的新生命,賜給我們這群必死的人.人若能在再臨的上帝中,找到復活的希望,人是會放下憂慮,重拾喜樂的.

腓4:5　當叫眾人知道你們的「寬宏大量」,主近了.

如何能使自己常常「在主裡」,以致能「常常喜樂」呢?

腓4:6　應當一無掛慮,只要凡事藉著禱告,祈求,和感謝,將你們所要的告訴神.

「應當一無掛慮」,再一次,是一個命令.放下一切憂慮,this is an order.

如何做呢?

保羅說:「只要凡事藉著禱告,祈求,和感謝,將你們所要的告訴神」.

凡事都放在禱告中.思想的時候,在禱告中去思想.感受的時候,在禱告中去感受.行動的時候,在禱告中去行動.不要離開禱告,也即不要離開神.當凡事都在禱告中進行時,則每件事都滲透著神的恩典,神的慈愛,神的渴望.當一塊麵包滲滿著奶油時,這麵包便變成一塊美味的「奶油多」,味道不再一樣.同樣,當神的恩典滲滿我們面對的事情時,事情的味道也不再一樣.不再是叫人憂慮的,而是叫人喜樂的.

凡事都放在禱告裡,那麼,應如何禱告呢?

要「藉著禱告,祈求,和感謝,將你們所要的告訴神」,這句話,我理解為「禱告和祈求時,要伴以感謝,要在感謝裡禱告和祈求,將你們所要的告訴神」.感謝,是因為我們知道萬事從神而來.就像我們崇拜奉獻後所唱的獻禮文:「萬物都是從主而來,我們把從主而來的獻給主.」這感謝的心,就是一切禱告背後的心情.

以感謝的心,將一切告訴神.無論是生活裡一般的交托,抑或是急迫時的懇求,或禱告,或祈求,都在感謝的心情裡,帶到神面前.在生命平順時,在禱告裡,感謝神.在生命艱難時,在祈求裡,求神幫助,但仍感謝神,生命的艱難在神手裡總是恩典.這就是「在主裡」的生命.主耶穌在平順時,在一花一木裡,指出天父的美意.主耶穌在十架苦難前,在懇切的祈求裡,求天父拿走苦難,但心裡仍堅信,若天父不拿走苦難,則這苦難總有天父的美意.主耶穌明白,苦難與愛,難捨難離.要愛,就有苦難.承擔苦難,愛就流露出來了.主耶穌以感謝的心,接納生活的全部,視之為天父的心的展示.

當我們放下憂慮,凡事在感謝中向神禱告時,我們會得到什麼?

腓4:7　神所賜,出人意外的平安必在基督耶穌裡保守你們的心懷意念.

喜樂的人迎接的,是從神而來的平安.這平安是「出人意外」的,是超乎我們理解力的,是我們連想像也想像不到的.人能理解的,是塵世事物本有的邏輯,但神要賜給我們的,是塵世事物本性外的新天新地,是我們的今天不能推論的明天,是一全新的生命.「真正的新」,總是超乎人類的理解力的.我以「死人復活」為例.死人的「死」,已終結一切可能.「死」本身的邏輯,並不提供復活的可能.「復活」來自一「死亡」以外的全新領域,這是神的領域.唯有神,能將這個塵世沒有的「全新的明天」帶給我們這群必死的人.這不能理解,不能想像的「新的生命」,就是經文這裡所講的「人意外的平安」.這平安必在基督耶穌裡保守我們的心懷意念.

在這個憂慮的時代裡,「將臨期」對我們作出重大提醒.

「將臨期」提示我們這群基督徒,要懷有盼望,因為「主近了」.主近了,可以指空間上的近,主就在我們身旁；也可以指時間上的近,祂再來的氣息已滲滿我們現今的時間.主耶穌要再來,祂的再來會帶來一全新的明天.這全新的明天,充滿著神那種從「無」創造「有」的創造能力.在這創造能力裡,將要克服那些迫使我們憂慮的「虛無」力量,並在「虛無」之上創造一不再憂慮的喜樂人生.

我最後要說的是:你們要在主裡常常喜樂.我再說,你們要喜樂.這不是一般勸勉,this is an order.

但願:榮耀歸於聖父,聖子,聖靈.阿們.



\newpage

\section{成年耶穌也問題多多麼}
\label{sec:307}
\textbf{許開明牧師 - 成年耶穌也問題多多麼?}
\newline
\newline
\hyperref[sec:306]{< < < PREV SERMON < < <}
~
\hyperref[sec:toc]{[返目錄]}
~
\hyperref[sec:308]{> > > NEXT SERMON > > >}


路加福音 2:41-52
\newline
\begin{longtable}{cl}
\hline
\hline
章節 & 經文 (和合本修訂版)\\
\hline
2:41 & \begin{tabularx}{0.7\textwidth}{X} 每年逾越節，他父母都上耶路撒冷去。 \end{tabularx} \\ \\ \relax
2:42 & \begin{tabularx}{0.7\textwidth}{X} 當他十二歲的時候，他們按著過節的規矩上去。 \end{tabularx} \\ \\ \relax
2:43 & \begin{tabularx}{0.7\textwidth}{X} 守滿了節期，他們回去，孩童耶穌仍舊在耶路撒冷。他的父母並不知道， \end{tabularx} \\ \\ \relax
2:44 & \begin{tabularx}{0.7\textwidth}{X} 以為他在同行的人中間，走了一天的路程才在親屬和熟悉的人中找他， \end{tabularx} \\ \\ \relax
2:45 & \begin{tabularx}{0.7\textwidth}{X} 既找不著，就回耶路撒冷去找他。 \end{tabularx} \\ \\ \relax
2:46 & \begin{tabularx}{0.7\textwidth}{X} 過了三天，他們發現他在聖殿裡，坐在教師中間，一面聽，一面問。 \end{tabularx} \\ \\ \relax
2:47 & \begin{tabularx}{0.7\textwidth}{X} 凡聽見他的人都對他的聰明和應對感到驚奇。 \end{tabularx} \\ \\ \relax
2:48 & \begin{tabularx}{0.7\textwidth}{X} 他父母看見就很驚奇。他母親對他說：「我兒啊，為甚麼對我們這樣做呢？看哪，你父親和我很焦急，到處找你！」 \end{tabularx} \\ \\ \relax
2:49 & \begin{tabularx}{0.7\textwidth}{X} 耶穌對他們說：「為甚麼找我呢？難道你們不知道我應當在我父的家裡嗎？」 \end{tabularx} \\ \\ \relax
2:50 & \begin{tabularx}{0.7\textwidth}{X} 他所說的這話，他們不明白。 \end{tabularx} \\ \\ \relax
2:51 & \begin{tabularx}{0.7\textwidth}{X} 他就同他們下去，回到拿撒勒，並且順從他們。他母親把這一切的事都存在心裡。 \end{tabularx} \\ \\ \relax
2:52 & \begin{tabularx}{0.7\textwidth}{X} 耶穌的智慧和身量，並神和人喜愛他的心，都一齊增長。 \end{tabularx} \\ \\
[1ex]
\hline
\hline
\end{longtable}
講題:《成人耶穌也問題多多麼?》　Was Jesus once a Troubled Youth?

講員:許開明牧師

經文:路加福音2:41-52

場合:香港中文大學崇基學院禮拜堂

日期:2015年12月27日

引言:

關於耶穌的降生,馬太福音和路加福音都有詳細的記載；而關於耶穌出來傳道的事跡,四福音都記載.但是,關於耶穌的成年,只有路加福音2:40- 52有記載,這段是相當寶貴的經文,讓我們先讀一遍經文.路加福音2:41-52

經文

「41每年逾越節,他父母都上耶路撒冷去.42當他十二歲的時候,他們按著過節的規矩上去.43守滿了節期,他們回去,孩童耶穌仍舊在耶路撒冷.他的父母並不知道,44以為他在同行的人中間,走了一天的路程才在親屬和熟悉的人中找他,45既找不著,就回耶路撒冷去找他.46過了三天,他們發現他在聖殿裡,坐在教師中間,一面聽,一面問.47凡聽見他的人都對他的聰明和應對感到驚奇.48他父母看見就很驚奇.他母親對他說:『我兒啊,為甚麼對我們這樣做呢?看哪,你父親和我很焦急,到處找你!』49耶穌對他們說:『為甚麼找我呢?難道你們不知道我應當在我父的家裡嗎?』50他所說的這話,他們不明白.51他就同他們下去,回到拿撒勒,並且順從他們.他母親把這一切的事都存在心裡.52耶穌的智慧和身量,並上帝和人喜愛他的心,都一齊增長.」

成年耶穌有問題麼?

2.1　不聽從父母?

父母帶著12歲快成人的耶穌到耶路撒冷守節,守滿了節期,耶穌卻沒有跟著父母走,結果走失了,而父母也不知.耶穌已到成人,應懂守規矩,既然不跟著去,應該早些通知父母!免父母擔心才對!成人耶穌跟時下許多青少年一樣,真沒有責任心!

2.2　不體諒父母的苦心?

父母為了找成人耶穌,在某處折返,跑回耶路撒冷共用了三天時間找耶穌,不見了耶穌三天,這三天,真不易過:第一天,父母離開耶路撒冷,以為耶穌在同行的男孩中,找了一天,找不到.第二天,他們回頭到耶路撒冷,在回程中再找,找了一天,找不到.第三天:他們到耶路撒冷,在城中又找了一天,最後才在聖殿找到耶穌.父母一定很擔心,掛慮,寢食不安!後來在聖殿裡找到耶穌.他母親對他說:「我兒啊,為甚麼對我們這樣做呢?看哪,你父親和我很焦急,到處找你!」成人耶穌的回答竟然是「為甚麼找我呢?難道你們不知道我應當在我父的家裡嗎?」成人耶穌明明知道父母愛他才會找他,竟然還回應:「為甚麼找我呢? 」這樣太不體諒父母,太不近人情了,太無禮貌了!還有,那句:「難道你們不知道我應當在我父的家裡嗎?」這句話令馬利亞摸不著頭腦,也令人摸不著頭腦:「耶穌的家不是在拿撒勒,約瑟跟馬利亞的家嗎?」如果耶穌認為他應當以聖殿為家的話,他應早說明,並告知父母,他會留在聖殿.成人耶穌不應以責怪的口吻回應馬利亞的,太不敬了.

3　無時間觀念?

成人耶穌在殿裡,坐在教師中間,一面聽,一面問,耶穌應該知道,何時要回家的,聽與問雖然有意義,也應有時限,有節制.耶穌應該知道,父母在找他的.若然他父母不找他,他是否一直問下去,毫無限制呢?

2.4　不懂向父母道歉?

從上的分析,耶穌不是全部不對,也有不對的,他竟然沒有「向父母道歉」.雖然耶穌在聖殿,他的智慧受到人的讚賞,但是他有智慧卻沒有品德!不過,耶穌應該很有品的嗎?

成人的耶穌,或者說仍是少年的耶穌,也有許多的問題?這些問題好像也是今日青少年常犯的問題:1不聽從父母,2不體諒父母的苦心,3 無時間觀念,4不懂向人道歉等.

然而,對於成人耶穌這些表現,是否真有問題?

3.1　不聽從父母,不通知父母不回家

原來猶太人守節,都是成群結隊步行,婦女小孩走前面,男人及大男孩走後面,一面走一面吟唱著上行之詩,或談論耶路撒冷城的雄偉,聖殿的見聞……熱鬧而有趣,到晚上要住宿過夜時,男人們才散去找妻小,女人們也各自散去找家人,一家團聚.

就在回程的第一天晚上,約瑟夫婦發現耶穌不見了,一夜不知問了多少親朋好友,就是沒有下落.約瑟夫婦判斷兒子可能沒跟上回程的隊伍,所以只好回頭再走一天,第三天終於又累又急的回到聖殿.

耶穌可能是跟著大男孩,並沒有與父母同行,他根本是無法通知到父母,他的行蹤的,他不想回家的消息,想在聖殿聽書,根本無法傳達.

說他不聽從父母,不通知父母不回家,這指控不成立.還有耶穌行了成人禮,便被當成成人了,他自己回家也是可以的.

3.2　不體諒父母的苦心

「48他父母看見就很驚奇.他母親對他說:『我兒啊,為甚麼對我們這樣做呢?看哪,你父親和我很焦急,到處找你!』」

「驚奇」這兩字,我們很多的時候會對此產生誤解,以為是耶穌的父母因為兒子所表現出來的聰明和智慧而驚喜高興；其實不然,在原文中明確指明他們因為看見「他」,感到驚訝和不知所措,這就是說他母親的那句話「我兒,為什麼向我們這樣行呢?看哪,你父親和我傷心來找你.」是有點責備的語氣,對耶穌的這種行為有些不滿.

「49耶穌對他們說:『為甚麼找我呢?難道你們不知道我應當在我父的家裡嗎?』」

單憑耶穌一句話「為甚麼找我呢?難道你們不知道我應當在我父的家裡嗎?」,並不能確定此觀點就是不體諒父母的苦心.不過,耶穌直率,真實地表達他的看法,這是年青人常會有表現.剛成人的耶穌可能覺得這是他首次來聖殿,這是他難得的機會,與一班滿有學識的宗教領袖在一起,是很好學習的機會,且他已行了成人禮,已是成人了,他應走自己的路,更應找上主要他走的路,體貼天父的心意,住在天父的家裡才是人生最重要的事情.他體貼天父的心,重要過體貼人的心,在耶穌道成肉身中,一般神學家看到的是魔鬼與上帝間的張力,卻忽視天上的父與地上的父母的張力,這張力在耶穌的一生中都存在著.這張力也存在於跟隨耶穌的基督徒,不容忽視!

3.3　無時間觀念

成人耶穌在殿裡,坐在教師中間,一面聽,一面問,耶穌太專注於聆聽和發問,年青人往往因太自我與自私,太專注自己的事情,也忽視了別人,忽視了父母會因為不見他而擔心起來,沒有好好管理好自己的時間.

然而,我們也要知道,由於時間的有限,我們也要知道時間須要分先後優次.成人的耶穌在聖殿中,他醒覺到自己與上主的獨特關係,他醒覺到自己「應以父的事為念」,這是他人生的重要取向.

耶穌非常珍惜與聖殿中的教師會的學習與研討,這不獨是難得機會,也是他人生中優先要抓住的機會,他懂得把握,抓住是人生的智慧.今日難得年青人在時間上有分先後優次的概念,也有一樣人生目標和他熱切的追求.

3.4　不懂向父母道歉

如果從耶穌的回答:「為甚麼找我呢?難道你們不知道我應當在我父的家裡嗎?」他不覺得自己做錯甚麼,當然不用道歉.

從宗教的角度,如果耶穌是上帝獨生子,是基督,他所作的都應是對的,根本不用道歉!反而人應該向成年的耶穌學習呢?成人了!耶穌的邁向成人的表現,對我們有何提示?今日我們大部份都是成年人了,但我們長大了麼?

小結:

簡言之,耶穌並沒有時下般的青少年問題,反而有成年人值得學習之處.

讓我們先認識成人禮及其意義.十二歲生日對猶太男孩是個重要的日子,他要參加一項特別的典禮(bar-mitzvah,意思是律法之子),

從此他由一個孩子的身分轉變為成人,要接受社會的權利和義務,自己負擔起宗教上的責任.十二歲以前的孩子是由父母代替他作決定,引導他認識神,但從此以後,他要與天父建立個人的關係了.一般虔誠的猶太人父母,他們的兒子在這個年齡都要接受Bar-mitzvah的儀式,Bar-mitzvah的意思就是「律法之子」.

耶穌可能是去聖殿接受這個儀式,祂也可能已經在拿撒勒的會堂中接受了這個儀式.聖經學者認為,有一點可以確定,祂進聖殿的時候,是帶著律法之子的權利進去的.在那之前,一個孩子的父母必須擔負這孩子在宗教方面的責任,但受過了儀式之後,他就必須自己負責.耶穌在這時,已經是律法之子了.

在猶太人的社會中,十三歲的孩子被視為由一個人生階段踏入另一個人生階段的開始,在這個過程中要做好充足的準備;十三歲的孩子被正式稱為「聖約之子」,他們需要自我承擔責任；而十二歲則是由孩童成為成年人的重要標誌,所以他們父母把他帶到聖殿來學習.

成人耶穌的表現

4.1　尋求獨立,自主,對自己所作的負責

成人禮後,一個孩子的身分轉變為成人,要接受社會的權利和義務,自己負擔起宗教上與處事為人的責任.第一,耶穌不跟同伴走,他決定留在聖殿跟拉比學習.第二,耶穌認為他當以父的事為念,而非人(包括父母)的事為念.

耶和華是最獨立自主的主

當時以色列人在埃及做奴隸,神呼召摩西,讓他帶領以色列人出埃及.摩西對神說:「我到以色列人那裡,對他們說:『你們祖宗的神差遣我到你們這裡來』,他們必問我:『他叫甚麼名字?』我要對他們說甚麼呢?」神回答摩西:「我是『自有永有者』.」

這裡的「我是」和「自有永有者」其實是一個詞:「是」的第一人稱形式.希伯來文的的「是」意義很豐富,有「存在」的意思,「我是」就是「我存在」,這個存在是一個現在持續式,意思是神是持續的永恆的存在.中文(和合本/新譯本等)意譯為「自有永有者」.神有兩個名字:第一個名字說明了他是在他自己裡面存在,「我是那自有永有者」,《七十士譯本》譯為ειμι ο ων「我是那存在者」,就是說,我是那個實質的,獨立的,不變的,必要的存在,祂現在是,以前是,將來也是.

他是自己存在(self-existent),他是屬於他自己的存在,不依靠於其他任何東西,而作為自己存在,他必然是自己充分的(self-sufficient),因此是完全充分的(all-sufficient),也是一切存在和祝福的不竭的源泉.在成人禮中,猶太人藉此要求人要獨立!在婚禮中,猶太人在此提到人要離開父母,成為獨立的人,二人要成為一體.

4.2　從聽問中,尋求知識,認識,了解

成人耶穌,問問題多多,藉著聆聽與發問,尋求知識:

「46過了三天,他們發現他在聖殿裡,坐在教師中間,一面聽,一面問.47凡聽見他的人都對他的聰明和應對感到驚奇.48他父母看見就很驚奇.」

猶太人注重「發問的精神」,開啟知識寶庫的鑰匙就是「發問的精神」.行過成人禮的耶穌,是律法之子,可進入聖殿,聆聽和發問,從而得著知識.

第一,發問是思想的初步,研究的動機.一切知識的獲得,大都從發問而來；第二,新發明,新創造也常常由發問而開其端.勇於發問,勤於發問的人,頭腦自然會日益豐盈,眼光自然會日益敏銳.從發問中,獲得知識；從知識中,獲得新的亮光；從新的亮光中,獲得覺悟.

舉一些例子,壺水開沸,誰不常常見到?卻只有瓦特把它當作問題研究,因而發明了蒸汽機,使人類至今蒙受其利.生,老,病,死,都是極普遍的人生現象.釋迦牟尼偏偏要尋根究柢,求個解答.他因此拋棄王位和家庭,獨自去潛修靜想,終於創立了佛教.我們中國的孔子,也是個好問的人,他說過「不恥下問」的話,他到太廟裡去的時候,看見每樣事物都要問.韓愈《師說》:「師者,所以傳道授業解惑也.」意思是說一個好教師的職責不僅是要傳授知識,技能,也要隨時為學生解答一切的疑問；反過來說,一個好學生不僅要在教師的門下認真學習知識,技能,同時也要主動的請求教師解決自己所遇到的一切疑惑,少年耶穌就是這方面的典範.

耶穌在聖殿中,聆聽拉比的講解,並與拉比的對談中,他得著醒悟:第一,他與上主有獨立的父子關係；第二,他一生應與天父的事為念.第三,他的家,他的國,不在地上而在天上.

一個成長的人,應經常發問,反省,不斷更新,不斷進步.作為大學生,大學的同事,作為一位認真的信徒,也當如此!

4.3　尋求人生理想方向

雖然耶穌行過成人禮後,便是成年,但實際上,12歲的耶穌,仍是青少年.他熱衷於尋找人生的方向和理想!中一時,我大概12或13歲,老師的作文題目:「我的志願」.我在中學教書,教中一班的作文題目也是「我的志願」.

耶穌的家:居住的家與屬靈的家；「生父的家」與「天父的家」；現實的家與理想的家.耶穌和家人到耶路撒冷,守完節後,他應該跟從父母回家去!但耶穌期盼,在此成年之時,尋找人生的方向與上帝的心意,結果他留下來在聖殿中聆聽與發問,尋索生命的方向,尋求上帝的心意.對上帝的追尋,以父的事為念(Or be about my Father’s business),不以自己的事為念,以地上的父為念與天上的父為念成為了耶穌的掙扎,恰好說明耶穌是一位有理想,有異象之人!當耶穌以天父的事為念,走上十字架的路,他便全力以赴!人生須要有方向和目標,才能活得積極.

中學時,我很喜歡寫作,作詩作詞,雖然我沒有成為作家(然而,個人專著15,與學者或牧者合著超過15本),但我活得很積極!1969年,我已寫了一本個人詩集,用文言寫的,其後轉寫新詩.

大學時,我很喜歡唱歌,因同房是音樂系,在房間放一些歌唱家的錄音帶等,引到我大唱特唱,但我也沒有成為歌唱家.

牧會期間,除牧養工作,上主給我建堂的使命,在望覺堂,香港閩南堂,全完堂都有建堂,忙碌得很；然而,我也留下美好的建堂腳蹤.

上主呼召我到崇基學院來,中華基督教會借調我到這裡,我也有雄心壯志做好校牧的工作.

4.4　尋求智慧作出抉擇

經過聆聽,發問後,成人耶穌要作出順從的抉擇:第一,身為律法之子,他要認識律法,順從律法,按規矩而行!第二,身為人之子,他要順從地上的父母!第三,身為神之子,他要順從上主!成人的耶穌要作出智慧的抉擇,不獨是在對與錯的抉擇,而是順從誰的抉擇?

「50他所說的這話,他們不明白.父母不明白,甚至不認同耶穌的使命,耶穌有位天上的父!51他就同他們下去,回到拿撒勒,並且順從他們:耶穌遵守上主的誡命,順從父母,孝敬父母,跟他們回家盡人的責任.」

縱然父母未能明白主,成年的主,不是用對抗方式,而是順從的方式.事實上,12歲的青少年,在世只有12年,入世未深,順從也是需要的!今日有些青少年常跟父母對立,甚至對抗,或父母基於自尊採取對立, 結果弄致家庭永無寧日!

舉一個例子:我叫:「阿女,幫我拿一枝筆來!」

幼稚園時:女兒總是聽的,很少發問；聽見,就會說:「好呀!」

小學時:女兒都會聽,但也會問,回應說:「等一下!」因她要完成手上的工作,或是會問:「為何只叫我,不叫家姐?」,這是一個公平的問題.

中學時:女兒就少聽,卻變得多問,反問:「你自己拿!」,因為你不應該叫我,自己的事自己做.又或者會十問九唔應,十問九唔做!總會質疑,拿一枝筆,為何不自己去做?於是建立了一個觀念:「自己的事自己做」

高中時:最近,女兒大了一年,講了句「自己的事自己做」後,不久,又靜靜拿枝筆給我,說:「下次自己做!」

我也欣然接受,那是因為「愛」,使她最終順從爸爸!耶穌也一樣,耶穌也愛地上父母,雖然父母不明白他,他因為愛而順從.

4.5　尋求完美的增長指標

「52耶穌的智慧和身量,並上帝和人喜愛他的心,都一齊增長.」

第一,祂心智的成長──智慧增長.

第二,祂身體的成長──身量增長.

第三,祂屬靈的成長──神喜愛他.

第四,祂社交的成長──人喜愛祂.

52節告訴我們:成人耶穌的四樣東西「都一齊增長.」聖經學者有以下的看法:發拉爾(Furrar)說,「這裡用的增長一詞是從『先鋒』一詞來的,意指軍隊一面砍伐樹木開路,一面向前挺進.」

柴義爾(Thayer)是位元偉大的文字學家說:「增長的意思是用鎚擊打使其變長,好像鐵匠打鐵一樣.」兩者的定義都有相同的概念.我們若細察「增長」這詞的構造,它真正的意思是向前砍,向前打,不停的劈砍.因此這詞是指一種不屈不撓的行動,而不是被動的發展.

這孩子耶穌在前十二年中漸漸長大,是被動的發展,沒有任何的責任.成人禮後,祂成了律法之子,要守規!祂的生命中有了責任,要承擔!祂在未來不只是被動的發展,並且要學會控制整個生命,要主動.祂的一切都要增長:身體,智慧,靈性,人際關係,要付出代價!增長,成長都是痛苦的,不是輕易的事!

結語

成年耶穌有問題麼?當然無問題!反而,成年的耶穌,值得我們這班成年人,甚至已邁向老年人反思:今日作為成人,我們要問自己五個「有或沒有?」.

第一,有沒有尋求獨立自主?

第二,有沒有從聽與問中尋求知識?

第三,有沒有尋求智慧作出順從誰的抉擇?

第四,有沒有尋求理想的人生方向?

第五,有沒有尋求完美的增長指標?

成年耶穌的表現,引出許多值得我們深思的表現,願共勉.



\newpage

\section{戀戀山城}
\label{sec:308}
\textbf{林豪恩先生 - 戀戀山城}
\newline
\newline
\hyperref[sec:307]{< < < PREV SERMON < < <}
~
\hyperref[sec:toc]{[返目錄]}
~
\hyperref[sec:309]{> > > NEXT SERMON > > >}


耶利米書 31:7-14
\newline
\begin{longtable}{cl}
\hline
\hline
章節 & 經文 (和合本修訂版)\\
\hline
31:7 & \begin{tabularx}{0.7\textwidth}{X} 耶和華如此說：「你們當為雅各歡樂歌唱，為萬國中為首的歡呼。當傳揚，頌讚說：『耶和華啊，求你拯救你的百姓，拯救以色列的餘民。』 \end{tabularx} \\ \\ \relax
31:8 & \begin{tabularx}{0.7\textwidth}{X} 看哪，我必將他們從北方之地領來，從地極召集而來；同他們來的有盲人、瘸子、孕婦、產婦；他們必成群結隊回到這裡。 \end{tabularx} \\ \\ \relax
31:9 & \begin{tabularx}{0.7\textwidth}{X} 他們要哭泣而來。我要照他們懇求的引導他們，使他們在河水旁行走正直的路，他們在其上必不致絆跌；因為我是以色列的父，以法蓮是我的長子。 \end{tabularx} \\ \\ \relax
31:10 & \begin{tabularx}{0.7\textwidth}{X} 列國啊，要聽耶和華的話，要在遠方的海島傳揚，說：「趕散以色列的必召集他，看守他，如牧人看守羊群。」 \end{tabularx} \\ \\ \relax
31:11 & \begin{tabularx}{0.7\textwidth}{X} 因為耶和華救贖了雅各，救贖他脫離比他更強之人的手。 \end{tabularx} \\ \\ \relax
31:12 & \begin{tabularx}{0.7\textwidth}{X} 他們來到錫安的高處歌唱，因耶和華的宏恩而喜樂洋溢，就是五穀、新酒和新的油，並羔羊和牛犢。他們必像有水澆灌的園子，一點也不再有愁煩。 \end{tabularx} \\ \\ \relax
31:13 & \begin{tabularx}{0.7\textwidth}{X} 那時，少女必歡樂跳舞；年輕的、年老的，都一同歡樂；因為我要使他們的悲哀變為歡喜，並要安慰他們，使他們的愁煩轉為喜樂。 \end{tabularx} \\ \\ \relax
31:14 & \begin{tabularx}{0.7\textwidth}{X} 我必以肥油使祭司的心滿足，我的百姓也要因我的恩惠知足。這是耶和華說的。 \end{tabularx} \\ \\
[1ex]
\hline
\hline
\end{longtable}
詩篇 147:12-20
\newline
\begin{longtable}{cl}
\hline
\hline
章節 & 經文 (和合本修訂版)\\
\hline
147:12 & \begin{tabularx}{0.7\textwidth}{X} 耶路撒冷啊，要頌讚耶和華！錫安哪，要讚美你的神！ \end{tabularx} \\ \\ \relax
147:13 & \begin{tabularx}{0.7\textwidth}{X} 因為他堅固了你的門閂，賜福給你中間的兒女。 \end{tabularx} \\ \\ \relax
147:14 & \begin{tabularx}{0.7\textwidth}{X} 他使你境內平安，用上好的麥子使你滿足。 \end{tabularx} \\ \\ \relax
147:15 & \begin{tabularx}{0.7\textwidth}{X} 他向大地發出命令，他的話速速頒行。 \end{tabularx} \\ \\ \relax
147:16 & \begin{tabularx}{0.7\textwidth}{X} 他降雪如羊毛，撒霜如灰燼。 \end{tabularx} \\ \\ \relax
147:17 & \begin{tabularx}{0.7\textwidth}{X} 他擲下冰雹如碎渣，他發出寒冷，誰能當得起呢？ \end{tabularx} \\ \\ \relax
147:18 & \begin{tabularx}{0.7\textwidth}{X} 他一出令，這些就都融化，他使風颳起，水便流動。 \end{tabularx} \\ \\ \relax
147:19 & \begin{tabularx}{0.7\textwidth}{X} 他將他的道指示雅各，將他的律例典章指示以色列。 \end{tabularx} \\ \\ \relax
147:20 & \begin{tabularx}{0.7\textwidth}{X} 他未曾這樣對待別國，至於他的典章，他們向來都不知道。哈利路亞！ \end{tabularx} \\ \\
[1ex]
\hline
\hline
\end{longtable}
約翰福音 1:10-18
\newline
\begin{longtable}{cl}
\hline
\hline
章節 & 經文 (和合本修訂版)\\
\hline
1:10 & \begin{tabularx}{0.7\textwidth}{X} 他在世界，世界是藉著他造的，世界卻不認識他。 \end{tabularx} \\ \\ \relax
1:11 & \begin{tabularx}{0.7\textwidth}{X} 他來到自己的地方，自己的人並不接納他。 \end{tabularx} \\ \\ \relax
1:12 & \begin{tabularx}{0.7\textwidth}{X} 凡接納他的，就是信他名的人，他就賜他們權柄作神的兒女。 \end{tabularx} \\ \\ \relax
1:13 & \begin{tabularx}{0.7\textwidth}{X} 這些人不是從血生的，不是從情慾生的，也不是從人的意願生的，而是從神生的。 \end{tabularx} \\ \\ \relax
1:14 & \begin{tabularx}{0.7\textwidth}{X} 道成了肉身，住在我們中間，充充滿滿地有恩典有真理，我們也見過他的榮光，正是父獨一兒子的榮光。 \end{tabularx} \\ \\ \relax
1:15 & \begin{tabularx}{0.7\textwidth}{X} 約翰為他作見證，喊著說：「這就是我曾說：『那在我以後來的先於我，因為在我以前，他已經存在。』」 \end{tabularx} \\ \\ \relax
1:16 & \begin{tabularx}{0.7\textwidth}{X} 從他的豐富裡，我們都領受了恩典，而且恩上加恩。 \end{tabularx} \\ \\ \relax
1:17 & \begin{tabularx}{0.7\textwidth}{X} 律法是藉著摩西頒佈的；恩典和真理卻是由耶穌基督來的。 \end{tabularx} \\ \\ \relax
1:18 & \begin{tabularx}{0.7\textwidth}{X} 從來沒有人見過神，只有在父懷裡獨一的兒子將他表明出來。 \end{tabularx} \\ \\
[1ex]
\hline
\hline
\end{longtable}
講題:戀戀山城　Love in Jerusalem

經文:耶利米書31:7-14;詩篇147:12-20;約翰福音1:10-18

講員:林豪恩先生

場合:香港中文大學崇基學院禮拜堂

日期:2016年1月3日

一．詩篇147:12-20保護

147:12耶路撒冷啊,你要頌讚耶和華!錫安哪,你要讚美你的上帝!

147:13因為他堅固了你的門閂,賜福給你中間的兒女.

147:14他使你境內平安,用上好的麥子使你滿足.

147:15他發命在地；他的話頒行最快.

147:16他降雪如羊毛,撒霜如爐灰.

147:17他擲下冰雹如碎渣；他發出寒冷,誰能當得起呢?

147:18他一出令,這些就都消化；他使風颳起,水便流動.

147:19他將他的道指示雅各,將他的律例典章指示以色列.

147:20別國他都沒有這樣待過；至於他的典章,他們向來沒有知道.你們要讚美耶和華!

147:12耶路撒冷啊,你要頌讚耶和華!錫安哪,你要讚美你的上帝!

耶路撒冷是一座山城,位於以色列中央山脈海拔約七百五十公尺的平臺上,東距約但河約三十公里,西距地中海約五十公里．耶路撒冷是一座舉世聞名的歷史古城,距今已有5000多年的歷史.自3000年前被大衛王攻下並將其定為首都,故又稱「大衛之城」或稱「錫安」；其後,所羅門王在這裡建造第一座聖殿,使耶路撒冷成為猶太人的宗教和權力中心.

為甚麼大衛王選擇耶路撒冷這座山城作為首都呢?

重慶是中國著名的山城．抗日戰爭時期,國民政府預料南京將會失守,就在日軍佔領南京前,趕及遷都重慶．選擇重慶,就是取其易守難攻的地理優勢．在1937年到1944年間,來自浙江,湖南,安徽,湖北,江蘇,上海等長江中下游地區的百萬移民遷往重慶及其周邊地區.同時,數以萬計的企業,學校,近8萬噸黃金都搬遷至重慶,因此重慶成為抗日戰爭時期中國的政治,經濟,文化中心,以及第二次世界大戰的遠東指揮中心．

耶路撒冷這座山城三面溪谷之底與城落差達四十至七十公尺,又非常陡峭險峻,也是易守難攻的天險之地．多少年來,不少以色列人以為山城就是他們的保障,然而當這座他們以為堅不可摧的山城在主前586年被巴比倫人徹底拆毀後,他們才如夢初醒,驚覺原來易守難攻的山城也有靠不住的時候.他們不禁疑惑,如果山城也不能保護他們,到底甚麼才是可靠的保障呢?當時以色列人的疑惑,不禁令人聯想到一九九七年的香港,當所有人都以搶購「磚頭」為保障的時候,一年之後,這些保障竟然變成了負資產.二零零八年之前,當不少人以為購買債券為低風險的安全保障的時候,突然間金融海嘯掩至,債券一夜之間變成了廢紙.

詩人在另一首詩說:「我要向山舉目」,但原來山城也可以失陷,詩人不禁疑問:「我的幫助從何而來?」

147:12耶路撒冷啊,你要頌讚耶和華!錫安哪,你要讚美你的上帝!

147:13因為他堅固了你的門閂,賜福給你中間的兒女.

147:14他使你境內平安,用上好的麥子使你滿足.

當那時候,人們終於醒覺,幫助不是從山城地理優勢而來,原來,”我的幫助從造天地的耶和華而來”．透過描述耶和華上帝對天地的掌管,詩人指出耶和華是有能力的保護者.

147:15他發命在地；他的話頒行最快.

147:16他降雪如羊毛,撒霜如爐灰.

147:17他擲下冰雹如碎渣；他發出寒冷,誰能當得起呢?

147:18他一出令,這些就都消化；他使風颳起,水便流動.

耶和華這位有能力的保護者是如何保護祂的子民呢?如果我們以當下父母保護子女的表現去期望耶和華的保護,就可能會感到失望.循道衛理楊震社會服務處曾經訪問育有六至十一歲子女的父母,發現近七成家長過分保護子女.鄰近地區有極端的例子是這樣的:

2015年11月19日下午放學,溫嶺市城西街道某民工子弟學校三年級某班班長小鄒因為值日問題與同班同學小文發生爭執,推撞中小文的指甲劃到了小鄒的臉.小鄒的父親要幫女兒討回「公道」.第二天,他替女兒出頭衝到教室打了同學小文,又把來勸阻的老師打成「熊貓眼」.(《錢江晚報》)

耶和華不是溫嶺市小鄒的父親,也不是保定市某公安局的局長李剛,祂不會縱容祂的子民以祂為靠山仗勢唬人,以勢壓人.出乎意料之外,耶和華這位大能者對耶路撒冷的保護,對某些人來說,可能被看作束縛和限制.詩人如此說:

147:19他將他的道指示雅各,將他的律例典章指示以色列.

147:20別國他都沒有這樣待過；至於他的典章,他們向來沒有知道.

以色列人以為,保護來自山城；無獨有偶,現今官場中不少人也以為,保護來自背後的大靠山.然而,哪一座山可靠呢?是元老之子,中央政治局委員薄熙來嗎?是經歷五十年軍旅生涯,官至中央軍委副主席徐才厚嗎?"山"真的可以靠嗎?有一首詩歌如此唱:「大山可以挪開,小山可以遷移,但主的慈愛永不離開你.」

就正如孩子向父母投訴:「為甚麼你總是對我有要求,而不要求鄰家的小朋友?」父母語重心長地說:「因為我們愛護你,才要求你,而不要求鄰家的小朋友.」詩人指出,耶和華給予律例典章,是祂對子民的眷顧.真正的保護來自哪裡呢?詩人認為是遵守耶和華的律例典章.在以色列第一位君王掃羅被立的場合如此記載:"撒母耳將君王的典章對百姓說明,又記在書上,放在耶和華面前．"(撒上10:25)律例典章的目的是要規範君王與百姓的關係,讓社會運作及人民生活在一種協定的信任之上．詩篇72篇如此說:

72:1上帝啊,求你將判斷的權柄賜給王,將公義賜給王的兒子.

72:2他要按公義審判你的民,按公平審判你的困苦人.

72:3大山小山都要因公義使民得享平安.

72:4他必為民中的困苦人伸冤,拯救窮乏之輩,壓碎那欺壓人的.

俗語說,物必先腐後而蟲生,山城何以敗落?耶路撒冷何以被拆毀?以賽亞先知如此說:

1:21可歎,忠信的城變為妓女!從前充滿了公平,公義居在其中,現今卻有兇手居住.

1:22你的銀子變為渣滓；你的酒用水攙對.

1:23你的官長居心悖逆,與盜賊作伴,各都喜愛賄賂,追求贓私.他們不為孤兒伸冤；寡婦的案件也不得呈到他們面前.

「子貢問政.子曰:足食,足兵,民信之矣.子貢曰:必不得已而去,於斯三者何先?曰:去兵.子貢曰:必不得已而去,於斯二者何先?曰:去食.自古皆有死,民無信不立.」(子貢問治理國家.孔子說:「糧食充足,軍備充足,人民信任.」子貢說:「如果不得不去掉一項,那麼這三項應先去掉哪一項?」孔子說:「去掉軍備.」子貢說:「如果不得不再去掉一項,那麼這兩項應去掉哪一項?」孔子說:「去掉糧食.自古人終究是要死的,如果失去人民的信任,那就沒有立國之本了.」)

"大山小山都要因公義使民得享平安"．這才是真正的保障．

戀戀山城,以色列人失去山城,才發現甚麼才是保障,也看見原來有人一直在眷戀他們．

二．耶利米書 31:7-14管教

31:7耶和華如此說:你們當為雅各歡樂歌唱,因萬國中為首的歡呼.當傳揚頌讚說:耶和華啊,求你拯救你的百姓以色列所剩下的人.

31:8我必將他們從北方領來,從地極招聚；同著他們來的有瞎子,瘸子,孕婦,產婦；他們必成為大幫回到這裡來.

31:9他們要哭泣而來.我要照他們懇求的引導他們,使他們在河水旁走正直的路,在其上不致絆跌；因為我是以色列的父,以法蓮是我的長子.

31:10列國啊,要聽耶和華的話,傳揚在遠處的海島說:趕散以色列的必招聚他,又看守他,好像牧人看守羊群.

31:11因耶和華救贖了雅各,救贖他脫離比他更強之人的手.

31:12他們要來到錫安的高處歌唱,又流歸耶和華施恩之地,就是有五穀,新酒,和油,並羊羔,牛犢之地.他們的心必像澆灌的園子；他們也不再有一點愁煩.

31:13那時,處女必歡樂跳舞；年少的,年老的,也必一同歡樂；因為我要使他們的悲哀變為歡喜,並要安慰他們,使他們的愁煩轉為快樂.

31:14我必以肥油使祭司的心滿足；我的百姓也要因我的恩惠知足.這是耶和華說的.

以賽亞先知曾說:

49:14錫安說:耶和華離棄了我；主忘記了我.

49:15婦人焉能忘記她吃奶的嬰孩,不憐恤她所生的兒子?即或有忘記的,我卻不忘記你.

山城雖然被毀,但耶和華沒有忘記祂的子民．耶路撒冷雖然不保,詩人指出耶和華並非是沒有能力的保護者．

31:10列國啊,要聽耶和華的話,傳揚在遠處的海島說:趕散以色列的必招聚他,又看守他,好像牧人看守羊群.

正如希伯來書所說:

12:6因為主所愛的,他必管教,又鞭打凡所收納的兒子.

12:7你們所忍受的,是上帝管教你們,待你們如同待兒子.焉有兒子不被父親管教的呢?

上帝的子民可曾想到,耶和華把他們當作兒子般對待?縱使他們失敗跌倒,耶和華仍不放棄他們．如果他們好像年幼的冰心向她的媽媽發問:"媽媽,你為甚麼愛我?"31章9節的話正是耶和華的回答:"因為我是以色列的父,以法蓮是我的長子."

戀戀山城,以色列人對山城的眷戀,遠遠不及耶和華對祂的子民的眷戀那麼長闊高深．先知呼喚以色列人慶祝回歸,然而回歸不單單是耶路撒冷這山城,而是回歸一個按照耶和華的律例典章而建立的公義和公平的社會,這未必是一個強國,而是一個連最弱小的也有位置的社會:

31:8我必將他們從北方領來,從地極招聚；同著他們來的有瞎子,瘸子,孕婦,產婦；他們必成為大幫回到這裡來.

31:13那時,處女必歡樂跳舞；年少的,年老的,也必一同歡樂；

四．約翰福音1: 10-18臨在

耶穌在馬太福音23:37曾說:「耶路撒冷啊,耶路撒冷啊,你常殺害先知,又用石頭打死那奉差遣到你這裡來的人.我多次願意聚集你的兒女,好像母雞把小雞聚集在翅膀底下,只是你們不願意.上帝的子民可能無法想像,有一天,上帝自己道成肉身親臨這山城,讓人們都看見山城是山城,上帝是上帝,山城不是上帝,上帝不是山城．

1:10他在世界,世界也是藉著他造的,世界卻不認識他.

1:11他到自己的地方來,自己的人倒不接待他.

1:12凡接待他的,就是信他名的人,他就賜他們權柄作上帝的兒女.

1:13這等人不是從血氣生的,不是從情慾生的,也不是從人意生的,乃是從上帝生的.

1:14道成了肉身,住在我們中間,充充滿滿地有恩典有真理.我們也見過他的榮光,正是父獨生子的榮光.

1:15約翰為他作見證,喊著說:「這就是我曾說:『那在我以後來的,反成了在我以前的,因他本來在我以前.』」

1:16從他豐滿的恩典裡,我們都領受了,而且恩上加恩.

1:17律法本是藉著摩西傳的；恩典和真理都是由耶穌基督來的.

1:18從來沒有人看見上帝,只有在父懷裡的獨生子將他表明出來.

曾經,上帝透過律例典章的保護去表達祂對子民的眷愛；曾經,上帝透過管教挽回去表達祂對子民的眷愛；現在,上帝親臨人間,所說的仍是那句情話:"賜他們權柄作上帝的兒女"(1:12)．有一位教神學的老師說過以下的故事:

很久以前有一位神醫藥王,他超卓的醫術活人無數,但有一天他的仇家在他三歲的兒子身上下了劇毒,這本是無藥可解的,但藥王苦思三天,終於想到一個方法,是要先給他兒子吃一種極苦的藥,再用金針渡穴大法刺遍他身上各大穴道,然後將他放在蒸籠七七四十九天,將毒逼出．但他兒子吃了若藥後開始嚎啕大哭,再聽說父親要用金針,又要用蒸籠"對付"他,心中不禁淒苦．藥王無論怎樣解釋,他不是聽不明白就是聽不進去．藥王問兒子:"你不相信我是為了你好嗎?"兒子回答:"我一直相信你對我是好的,但我眼前看到的是麼可怕的針和籠,叫我怎樣相信呢?"藥王最後說:"好吧!我刺你一針,也刺自己一針；你在蒸籠一日,我也在蒸籠一日．"兒子眨眨眼,半信半疑:"真的?但你內功那麼深厚,可能根本不怕痛?"藥王歎了一口氣,接著閃電般在自己身上點了一百指:"寶貝兒子,為了讓你相信我,我作甚麼也願意．我現在已封了我的內功,你滿意了吧!時間緊逼,我們現在就開始痛苦的療程．"(關啟文:我信故我思)

聖子耶穌基督甘願道成肉身,擁抱及品嚐人間的苦怒哀樂,在承受十架苦難時,在說出“我的神!我的神!為甚麼離棄我?” (太二七:46)之時,仍能稱上帝為父,說:‘父啊,我將我的靈魂交在你手裡!’ (路二三:46)上帝多次以父母的心賜去對待祂的子民,稱他的子民為兒女；相對來說,祂的子民卻經常感到困難去視上帝為父母,道成肉身來到人間的聖子耶穌基督展示了無論在怎樣的處境之下,視上帝為父母是可能的．

五．總結

著名作家魯益師在《四種愛》一書中說:「去愛,就會受傷害.只要你愛什麼,你心中就會有苦惱,甚至極為傷心.如果你真要保持你的心不受傷害,你必須什麼都不愛.你可以……避免所有情感上的牽累,把你的愛安全地封鎖在你的自我中心裡…….除天堂之外,唯一能使你安全,避免所有關於愛的危險的地方－那就是地獄.」

眷戀山城的人,就是希望靠著山城的保護而不會受到傷害．眷戀生活在山城中的人的上帝,卻因為其眷戀而經受傷痛．弟兄姊妹,當我們想到那位受傷的上帝的時候,願我們也感受到祂對我們的眷戀．



\newpage

\section{教會論之下的人觀}
\label{sec:309}
\textbf{鄧瑞強博士 - 教會論之下的人觀}
\newline
\newline
\hyperref[sec:308]{< < < PREV SERMON < < <}
~
\hyperref[sec:toc]{[返目錄]}
~
\hyperref[sec:310]{> > > NEXT SERMON > > >}


哥林多前書 12:12-27
\newline
\begin{longtable}{cl}
\hline
\hline
章節 & 經文 (和合本修訂版)\\
\hline
12:12 & \begin{tabularx}{0.7\textwidth}{X} 就如身體是一個，卻有許多肢體，身體的肢體雖多，仍是一個身體；基督也是這樣。 \end{tabularx} \\ \\ \relax
12:13 & \begin{tabularx}{0.7\textwidth}{X} 我們無論是猶太人是希臘人，是為奴的是自主的，都從一位聖靈受洗成了一個身體，並且共享這位聖靈。 \end{tabularx} \\ \\ \relax
12:14 & \begin{tabularx}{0.7\textwidth}{X} 身體原不只是一個肢體，而是許多肢體。 \end{tabularx} \\ \\ \relax
12:15 & \begin{tabularx}{0.7\textwidth}{X} 假如腳說：「我不是手，所以不屬於身體」，它不能因此就不屬於身體。 \end{tabularx} \\ \\ \relax
12:16 & \begin{tabularx}{0.7\textwidth}{X} 假如耳朵說：「我不是眼睛，所以不屬於身體」，它也不能因此就不屬於身體。 \end{tabularx} \\ \\ \relax
12:17 & \begin{tabularx}{0.7\textwidth}{X} 假如全身是眼睛，聽覺在哪裡呢？假如全身是耳朵，嗅覺在哪裡呢？ \end{tabularx} \\ \\ \relax
12:18 & \begin{tabularx}{0.7\textwidth}{X} 但現在神隨自己的意思把肢體一一安置在身體上了。 \end{tabularx} \\ \\ \relax
12:19 & \begin{tabularx}{0.7\textwidth}{X} 假如全都是一個肢體，身體在哪裡呢？ \end{tabularx} \\ \\ \relax
12:20 & \begin{tabularx}{0.7\textwidth}{X} 但現在肢體雖多，身體還是一個。 \end{tabularx} \\ \\ \relax
12:21 & \begin{tabularx}{0.7\textwidth}{X} 眼睛不能對手說：「我用不著你。」頭也不能對腳說：「我用不著你。」 \end{tabularx} \\ \\ \relax
12:22 & \begin{tabularx}{0.7\textwidth}{X} 不但如此，身上的肢體，人以為軟弱的，更是不可缺少的； \end{tabularx} \\ \\ \relax
12:23 & \begin{tabularx}{0.7\textwidth}{X} 身上的肢體，我們認為不體面的，越發給它加上體面；我們不雅觀的，越發裝飾得雅觀。 \end{tabularx} \\ \\ \relax
12:24 & \begin{tabularx}{0.7\textwidth}{X} 我們雅觀的肢體自然用不著裝飾；但神配搭這身子，把加倍的體面給那有缺欠的肢體， \end{tabularx} \\ \\ \relax
12:25 & \begin{tabularx}{0.7\textwidth}{X} 免得身體不協調，總要肢體彼此照顧。 \end{tabularx} \\ \\ \relax
12:26 & \begin{tabularx}{0.7\textwidth}{X} 假如一個肢體受苦，所有的肢體就一同受苦；假如一個肢體得光榮，所有的肢體就一同快樂。 \end{tabularx} \\ \\ \relax
12:27 & \begin{tabularx}{0.7\textwidth}{X} 你們是基督的身體，並且各自都是肢體。 \end{tabularx} \\ \\
[1ex]
\hline
\hline
\end{longtable}
講題:「教會論之下的人論」　Churchly Anthropology

經文:哥林多前書12:12-27

講員:鄧瑞強博士

場合:香港中文大學崇基學院禮拜堂

日期:2016年1月24日

各位弟兄姊妹,早安.

啟蒙時期的大哲學家康德(Kant)在他的巨著《純粹理性批判》(Critique of Pure Reason)裡指出,他的哲學甚或是人類的哲學都是為了解決三大問題:一,「我能知道什麼?」二,「我應做什麼?」三,「我可以盼望什麼?」而這三大問題,又可以用一個終極問題去總括,或者說,是為了回應一個終極問題,這問題就是:「人是什麼?」

「我能知道什麼?」這問題涉及形而上學及知識論.

「我應做什麼?」這問題涉及倫理學.

「我可以盼望什麼?」這問題涉及神學.

就康德來說,這些問題之所以有意義,是因為這些問題反映人性.

中國人說:「狗眼看人低」.若你眼中看到的人,都是低一點的,則可能反映你的人性裡,多少有一點狗性.我們認識到的世界,其實是反映我們的人性.知識論指向人性論.

若你的倫理學的首要原則是「擊敗身邊所有人」,則多少反映你孤獨的人性,因為在你身邊的人,不是朋友,全是對手,而你每刻的思想,便是以千方百計使他們敗在你的腳下.倫理學揭露人性.

若你的盼望是心想事成,而神是你心想事成的保證,則你的盼望反映了你的人性,你的人性是看自己為神.神學反映人性.

人是什麼?這可能是一關鍵問題.

弄不清人性,則帶來錯誤之知,錯誤之行,錯誤之期望.

常言道,最怕是「看錯人」,或許,更怕的是「看錯了人性」.

中國學者劉再復寫了一本書,書名叫《人論二十五種》,寫出千奇百怪的人性.書中討論的第一種人性,他名之為「傀儡人」.他寫道:「在表面上,高等傀儡人和真正的人已沒有差別,但在深層的精神面上,他們與真正的人仍有巨大的區別,這就是它沒有人之所以成為人所必需的獨立人格和表現這種人格的屬於自身的語言,即雖會說話,但不會說屬於自己的話.綜上所述,我們便可以知道所謂傀儡人,乃是被他人所操縱所掌握而沒有自己的靈魂和沒有自己的話語的人.」(中信出版社,2010年,頁4)

劉再復又講述一種「陰人」,所謂「陰人」,「也就是陽光世界下那種帶有鬼氣和鬼質的人」.「陰人在人鬼之間,是一種喪失生命激情的人.沒有血氣,沒有生命激情,萬物萬事都已看透,這也罷；但他們又偏偏不甘寂寞,仇視人間熱情.…倘若看到同行或看到其他人取得什麼成就,他們一定要從背後放一冷箭,高明地給予中傷.」(頁116, 117)

劉再復刻劃出眾生怪相.這種刻劃,差不多就是書寫人間悲劇了.

人性是什麼,不同思想家有不同說法.

馬克思會說,根本沒有人類共通的人性.資產階級和無產階級是對立的,兩種階級分屬兩個黨.人只有黨性,黨性大於任何意義的人性.當無產階級被統治時,唯一的目標是推翻資產階級的統治.當無產階級取得統治權時,其唯一渴望是無產階級專政.「無產階級」與「非無產階級」,看來是沒有共通的人性的.

我們習慣的人性論,是自由主義的人性論,這種人性論是以個體自由來定義的.自由主義的奠基人之一穆勒(John Stuart Mill)寫了本《論自由》(On Liberty),裡面講出一個關於自由的經典定義,穆勒說:「惟一名副其實的自由,是以我們自己的方式追求我們自己的善,只要我們不試圖剝奪別人的自由,或者妨礙他們獲得自由的努力.」(“The only freedom which deserves the name is that of pursuing our own good in our own way, so long as we do not attempt to deprive others of theirs, or impede their efforts to obtain it.”)

簡言之,人之為人,就是他有自由.他有自由追求自己的善,他有自由行自己的路,他有自由活自己想要的人生.只要不干涉別人,一個人是可以完全按自己的想法去運用自己的自由的.自由主義的人性,就是完全自由的個體.

然而,只要你覺得自己的自由是最重要的,則你總會覺得其他人多少是你自由的限制.你想坐得鬆動一點,緊緊坐在你身旁的人,便限制了你的自由.你想在地鐵裡有多一點站立的空間,站在你身邊的人,便限制著你的自由.你想安靜一點睡覺,但別人看電視的自由便限制了你的自由.在墮胎的議題上,高舉自由的 “Pro- Choice” 陣營,強調母親的自由抉擇,若母親認為胎兒妨礙了她的自由,她便可選擇清除這障礙.在離婚的議題上,若一個人覺得身邊的人已變成負累,在高舉個體自由的原則下,則很難說服他維持婚姻.自由主義高舉個體自由,他人的價值,視乎這個他人是增加我的自由抑或弱化我的自由.若他增加我的自由,則他多少已變成我的僕役.若他弱化我的自由,則他是我必須清除的障礙.若個體自由是最高原則,則身邊的每個人都只會是潛在的競爭對手,很難肯定他人的價值.一個要求絕對自由的人,多是孤獨的.

現在歐洲出現無數難民,若歐洲人每個人都只講個體自由,則沒有人能提出任何理由,要去接待難民的.純粹個體自由原則是不包含憐憫的.

德國大量接收難民,德國人質疑,這會否令德國有過多的回教徒.德國總理默克爾(Angela Merkel)回應得好,她說:「我們不是有太多回教徒,而是有太少基督徒.我們太少討論基督教的人觀.」

按默克爾的講法,拒絕難民,是遺忘了基督教對人的看法.

每一種思想都有一種人觀.基督教是如何看待「人」的呢?

今日,我們以哥林多前書12:12－27的經文,討論基督教的人觀.

(1)基督論,教會論,人論

12:12　就如身子是一個,卻有許多肢體；而且肢體雖多,仍是一個身子；基督也是這樣.

12:13　我們不拘是猶太人,是希臘人,是為奴的,是自主的,都從一位聖靈受洗,成了一個身體,飲於一位聖靈.……

12:27　你們就是基督的身子,並且各自作肢體.

經文討論人是什麼,討論的時候,不是以自由的個體的角度去討論的,而是以一群人的角度去討論的.討論這群人時,又是從基督的角度去討論的.用神學的講法是,討論「人論」,不離「教會論」；討論「教會論」,不離「基督論」.所以,今日講道的主題是:從基督論底下的教會論去討論人論.

基督是誰?基督是最自由的人,但他自由地選擇放下自己的自由,去為他人而活.自由主義的自由,是「free from others」；基督的自由,是「free for others」.基督在愛裡,將他人的生命承擔在自己的肩上.他人不是限制他自由的障礙,他人是他愛的對象,是他負責任的對象.為了他人,基督甘心放下自由.這是十架裡的真自由.

看到這份愛,明白這份愛,回應這份愛的人,會在愛中與基督連成一體,組成「基督的身體」,即「教會」.教會不是一般的社團,而是基督生命在世間的延續.接受「洗禮」的人,是在「洗禮」中與基督同死同活.同領「聖餐」的人,是分享基督的血和肉,以致生命在基督中轉化.

與基督連合的人,他的生命不再是為自己而活的,而是為他人而活的.因著基督,他與其他人的生命也連在一起.無論「是猶太人,是希臘人,是為奴的,是自主的」,不分種族,不問身分,都在基督裡合而為一.不同的人能夠合而為一,是因為他們都不再為自己而活,都不再高舉自己那唯我獨尊的自由,都因基督的捨己而不再堅持自我.所有被基督的愛感召的人,都在基督裡合而為一,成為一個被稱為「基督的身體」的信仰群體,這就是教會.

要各自為政,各有「排他的」自我的人合而為一,是一艱巨的工程,故經文說,各人「都從一位聖靈受洗,成了一個身體,飲於一位聖靈」,這即是說,要以聖靈之力,要以天地間的靈氣,要以神的大能,方能成就這事.

一個以聖靈之力建立的教會群體,會展現出一種什麼樣的人性?

(2)重視多元,但不是各自孤立的多元

12:14　身子原不是一個肢體,乃是許多肢體.

12:15　設若腳說:「我不是手,所以不屬乎身子；」它不能因此就不屬乎身子.

12:16　設若耳說:「我不是眼,所以不屬乎身子；」它也不能因此就不屬乎身子.

12:17　若全身是眼,從哪裡聽聲呢?若全身是耳,從哪裡聞味呢?

12:18　但如今,神自己的意思把肢體俱各安排在身上了.

教會群體重視生命的多元性.「若全身是眼,從哪裡聽聲呢?若全身是耳,從哪裡聞味呢?」教會群體反抗無味的單一性.獨裁的權力能夠以暴力製造和諧的大一統體制.大一統體制的缺失,是失去人的獨特性,即失去人性.教會卻使信徒在群體中保持其獨特性.「身子原不是一個肢體,乃是許多肢體」.然而,雖然各人有其獨特性,卻不會因此孤芳自賞,將他人視為潛在對手,反因基督的愛,與他人連成一體.正如經文說:「設若腳說:『我不是手,所以不屬乎身子』；它不能因此就不屬乎身子.設若耳說:『我不是眼,所以不屬乎身子』；它也不能因此就不屬乎身子.」

生命各有特色,但不是各自精彩,而是將生命的火堆在一起,讓火燒得更旺.

(3)重視多元,但在基督裡合一

剛才的經文強調生命的多元,以下的經文強調多元中的合一.

12:19　若都是一個肢體,身子在哪裡呢?

12:20　但如今肢體是多的,身子卻是一個.

12:21　眼不能對手說:「我用不著你；」頭也不能對腳說:「我用不著你.」

生命在基督裡是互相承擔,互相配搭,互相扶持的.

生命的互相結連,互相扶持,少不了對弱者的加倍關懷,提攜,鼓勵.

12:22　不但如此,身上肢體人以為軟弱的,更是不可少的.

12:23　身上肢體,我們看為不體面的,越發給它加上體面；不俊美的,越發得著俊美.……

12:25　免得身上分門別類,總要肢體彼此相顧.

生命的互相結連,要去到情感相通的地步.他人的成功,即我的成功.他人的傷痛,即我的傷痛.

12:26　若一個肢體受苦,所有的肢體就一同受苦；若一個肢體得榮耀,所有的肢體就一同快樂.

情感結連,才表明生命的彼此相屬.外在有交往,情感不結連,仍是將他人視為他人,他人終歸而言,仍是與我無關的.

(4)結語

基督招聚他的信徒組成教會,這不單是聚會的場所,而是活出新的人性的地方.信仰基督,就是在基督裡改造人性.這人性可以在真正的教會生活中體現出來.

這種新的人性,就其本體論結構而言,是放棄那種「獨立的」,「孤離的」,「我與你有別」而「對立於你」的,高舉個體自由的人性.而建立那種「社群的」,「互相扣連的」,「我與你有別」卻又「我中有你,你中有我」的,以基督的捨己的愛為生命內容的人性.

教會裡的新人性,就其倫理而言,是放棄自私的,自我的,我要比各人強的勝利者思維,而是活出那種互相扶持的,互相尊重的,關懷弱者的,為他人而活的,與人同情共感的人性.

我們能盼望什麼?我們盼望基督創造的新人性,能在教會中活現出來.並希望教會展示的新人性,能影響世人,使世人因著教會的生命見證,也羨慕這人性,也願意活出這人性.或許,這就是天國.

但願:榮耀歸於聖父,聖子,聖靈.阿們.



\newpage

\section{漠視真話封口滅口}
\label{sec:310}
\textbf{許開明牧師 - 漠視真話,封口滅口!}
\newline
\newline
\hyperref[sec:309]{< < < PREV SERMON < < <}
~
\hyperref[sec:toc]{[返目錄]}
~
\hyperref[sec:311]{> > > NEXT SERMON > > >}


路加福音 4:14-30
\newline
\begin{longtable}{cl}
\hline
\hline
章節 & 經文 (和合本修訂版)\\
\hline
4:14 & \begin{tabularx}{0.7\textwidth}{X} 耶穌帶著聖靈的能力回到加利利，他的名聲傳遍了四方。 \end{tabularx} \\ \\ \relax
4:15 & \begin{tabularx}{0.7\textwidth}{X} 他在各會堂裡教導人，眾人都稱讚他。 \end{tabularx} \\ \\ \relax
4:16 & \begin{tabularx}{0.7\textwidth}{X} 耶穌來到拿撒勒，就是他長大的地方。在安息日，照他素常的規矩進了會堂，站起來要念聖經。 \end{tabularx} \\ \\ \relax
4:17 & \begin{tabularx}{0.7\textwidth}{X} 有人把以賽亞先知的書交給他，他就打開，找到一處寫著： \end{tabularx} \\ \\ \relax
4:18 & \begin{tabularx}{0.7\textwidth}{X} 「主的靈在我身上，因為他用膏膏我，叫我傳福音給貧窮的人；差遣我宣告：被擄的得釋放，失明的得看見，受壓迫的得自由， \end{tabularx} \\ \\ \relax
4:19 & \begin{tabularx}{0.7\textwidth}{X} 宣告神悅納人的禧年。」 \end{tabularx} \\ \\ \relax
4:20 & \begin{tabularx}{0.7\textwidth}{X} 於是他把書捲起來，交還給管理人，就坐下。會堂裡的人都定睛看他。 \end{tabularx} \\ \\ \relax
4:21 & \begin{tabularx}{0.7\textwidth}{X} 耶穌對他們說：「你們聽見的這段經文，今天已經應驗了。」 \end{tabularx} \\ \\ \relax
4:22 & \begin{tabularx}{0.7\textwidth}{X} 眾人都稱讚他，並對他口中所出的恩言感到驚訝；他們說：「這不是約瑟的兒子嗎？」 \end{tabularx} \\ \\ \relax
4:23 & \begin{tabularx}{0.7\textwidth}{X} 耶穌對他們說：「你們一定會用這俗語向我說：『醫生，你醫治自己吧！我們聽見你在迦百農所做的事，也該在你自己的家鄉做吧。』」 \end{tabularx} \\ \\ \relax
4:24 & \begin{tabularx}{0.7\textwidth}{X} 他又說：「我實在告訴你們，沒有先知在自己家鄉被人接納的。 \end{tabularx} \\ \\ \relax
4:25 & \begin{tabularx}{0.7\textwidth}{X} 我對你們說實話，在以利亞的時候，天閉塞了三年六個月，遍地有大饑荒，那時，以色列中有許多寡婦， \end{tabularx} \\ \\ \relax
4:26 & \begin{tabularx}{0.7\textwidth}{X} 以利亞並沒有奉差往她們中任何一個人那裡去，只奉差往西頓的撒勒法一個寡婦那裡去。 \end{tabularx} \\ \\ \relax
4:27 & \begin{tabularx}{0.7\textwidth}{X} 在以利沙先知的時候，以色列中有許多痲瘋病人，但除了敘利亞的乃縵，沒有一個得潔淨的。」 \end{tabularx} \\ \\ \relax
4:28 & \begin{tabularx}{0.7\textwidth}{X} 會堂裡的人聽見這些話，都怒氣填胸， \end{tabularx} \\ \\ \relax
4:29 & \begin{tabularx}{0.7\textwidth}{X} 就起來趕他出城。他們的城造在山上；他們帶他到山崖，要把他推下去。 \end{tabularx} \\ \\ \relax
4:30 & \begin{tabularx}{0.7\textwidth}{X} 他卻從他們中間穿過去，走了。 \end{tabularx} \\ \\
[1ex]
\hline
\hline
\end{longtable}
講題:《漠視真話,滅聲滅口》　Neglect the Truth, Kill to Silence!

經文:路加福音4章14-30節

講員:許開明牧師

Edited by Jojo Leung

場合:香港中文大學崇基學院禮拜堂

日期:2016年1月31日

引言

看完本段經文,你可能跟我一樣,會問以下幾個問題.第一,為何耶穌在自己家鄉傳道不受歡迎,卻在外邦人中大受歡迎?第二,猶太人不是非常強調「本土主義」(簡稱自己人)的嗎?為何耶穌會反其道而行?第三,為何猶太人對耶穌會由欣賞,瞧不起,厭惡,憎恨到除滅耶穌?第四,到底是什麼引致耶穌殺身之禍?呵!原來耶穌講了二個以色列歷史事實,為何如此?第五,為何猶太人把耶穌由山頂推下去,耶穌卻無事,竟能從中間直行過去?

路加福音四章十三節:「魔鬼用完了各樣的試探,就暫時離開耶穌.」由這節經文可見,魔鬼不走掉,隨時出現,並與耶穌角力!十四節:「耶穌滿有聖靈的能力回到加利利,祂的名聲就傳遍了四方.」這裡的「耶穌滿有聖靈的能力」,意指聖靈是祂工作有能力的來源；「能力」在原文是指如同炸藥爆炸一般的有力量(dynamite).最初的「傳聞」,已經變成了「大喧噪」,到處的人都在談論祂.「回到加利利」的主耶穌在勝過魔鬼的試探之後,曾在猶太地傳道約有一年之久(參約二1－五47).這裡是記錄祂開始第二年的公開事奉.

從中可見,魔鬼與聖靈是兩股不同的力量:魔鬼,利用會堂的人與群眾；聖靈,與耶穌同在,給耶穌力量!

魔鬼利用三種破壞性的思維

第一種是「朝不起人」的思維.「這不是約瑟的兒子麼?」其意思就像「本地薑(約瑟的兒子麼?)唔竦」.在二十二節中:「眾人都稱讚祂,並希奇祂口中所出的恩言；又說:『這不是約瑟的兒子麼?』」自從主耶穌出來傳道之後,此處是惟一提到祂的養父「約瑟」名字的經節.一般認為約瑟很早就過世,主耶穌身為長子,在祂出來傳道之前,曾有一段年日接續養父約瑟的職業,作過木匠(參可六3),幫助母親撫養弟妹們.原本拿撒勒人對主的恩言不過帶著欣賞,評論的態度來「稱讚」,「希奇」,突然卻因主的身份低微──「約瑟的兒子」,就不肯好好領受,最後把祂攆出去(參29節).這原本集中耶穌的優良表現,後來轉移到耶穌家底.家庭或背景不顯赫,是否就否定人家的表現或能力?

我經常與中大同學,特別是明華宿生交談,發現許多成績好的學生,都不是出於什麼名校?就算當年1978我考入中大,成績也不錯,我讀的那間中學在觀塘都無人識.最近人氣急升的某公眾入物曾說:「某人無博士學學位,無學術貢獻?不能擔任港大副校長」,令人感到好像朝不起無博士學位的學者.日前又說「大學生應有理性,卻變得無理性,無知易被政黨利用」等,又引起軒然大波.

第二種思維是「只顧自己人」的思維,亦即是平日我們所說的「肥水不流別人田」,或者經文中俗語的「醫生,你醫治自己罷」.經文二十三節中:「耶穌對他們說:『你們必引這俗語向我說:“醫生,你醫治自己罷.我們聽見你在迦百農所行的事,也當行在你自己家鄉裡.”』」「醫生,你醫治自己罷」意即『你先幫助自己人罷』.感覺就等同「耶穌,你先在家鄉行神蹟,幫助自己人吧.」若然,耶穌在家鄉施行神蹟；對猶太人來說,最好的神蹟就是用來消滅羅馬的政權,這就引起猶太入跟羅馬人的對立,造成大亂.如此,耶穌的神蹟便被猶太人所利用,也被魔鬼所利用.

猶太人的自己人

猶太人重視自己人,背後有歷史的因素:亡國,被擄,被欺壓,自然地會求自保!耶穌的故鄉是拿撒勒,猶太人處於的拿撒勒是大城市,耶穌也具國際的視野.

一,大城市,國際都市

拿撒勒是耶穌的家鄉,也是祂最早期工作的地方之一.拿撒勒並不是一條村莊,而是一個Polis,意即一個市鎮或城市；當時人口已達二萬人之眾,人數算多.

二,在高山下

它位於加利利小山斜坡下一低陷之處,靠近耶斯列平原.小孩子只需要爬上城上的山頂,便可以見到方圓數哩驚人的全景.

三,歷史名城

喬治·史密夫爵士(Sir George Adam Smith)從山巔俯覽全地,以色列的歷史展現在觀看者眼前.這裡的以斯德倫(Esdraelon)平原便是:1底波拉和巴拉爭戰之所；2基甸在這裡取得他的勝利；3掃羅在此遭遇敗亡；4約西亞在此戰死沙場；5這裡有拿伯的葡萄園,6有耶戶殺死耶洗別的地方；7有以利沙住過的書念；8有以利亞和巴力的先知莊麗爭戰的迦密山.遠方的藍天深處,躺著地中海和其中的島嶼.

四,交通要塞

這通往世界各地的門戶亦由拿撒勒城山巔敞開,沿著山勢有三條大路.第一,從南方伸展直抵耶路撒冷,是朝聖者必經之路；第二,另外一條是海道,由埃及通達大馬士革,滿載貨物的行商坐賈,穿梭往來其間；第三,有一條大路向東伸展,成為亞拉伯商人必經之路；而羅馬的軍團亦由此而邁向東邊的疆界.

如果我們以為耶穌長成於窮鄉僻壤,那就是一個錯誤的想法.相反的,祂乃是在一個交通發達和歷史地位極重要的城鎮中成長.因此耶穌有廣闊的世界觀.猶太人的神觀,從自己人的神,演變成普世的神:首先,是亞伯拉穻的家族之神,摩西的以色民族的神,詩人曠野之神(山之神),大衛之以色列國的神,被擄後猶太人的神,耶穌的神是天上的父:普世之神.

猶太人對耶穌的不滿

這樣就奇怪了,民族性變得愈來愈本位主義.猶太人的神觀,民族化,私有化,成為了祝福自己人,排斥外人的神!在「自己人」的觀念下,猶太人對耶穌有以下的不滿.

第一,自己人應幫助自己人,所以耶穌這自己人應行神蹟幫助自己人；耶穌卻在故鄉不行神蹟!猶太人期望耶穌行神蹟,利用神蹟去消滅羅馬的政權,猶太人想利用耶穌達聖民眾的願望.

第二,自己人不應罵自己人,耶穌卻指出自己人是「不信」的,表現不佳.猶太人受不起耶穌的批評,不能面對歷史的事實.

第三,自己人不應讚賞恩待外邦人,耶穌人卻恩待外邦人,猶太人不認同耶穌,他的做法跟他們不一樣,猶太人不能接受自己人這樣的做法.

多元化的社會；多元的思想；多元的聲音；多元的角度.強調「自己人」或「本土主義」今日的自己人,小圈子,在家庭,社會,國家都不是好事.自己人對待自己人的時候,或許會自己人瞧不起自己人(本地薑唔辣),自己人罵自己人罵得更厲害,自己人不敬重自己人等.猶太人對自己人耶穌,更是反轉豬肚:罵他,出賣他,釘他十字架……

近代本土主義(localism)即以本土文化為其主要意識,主張本地優先的政策, 抱有本土主義思想的人士被稱為「本土派」.地方主義(localism),又稱本土主義,地方保護主義,地域主義,指一個地區為維護本土社會利益而實行文化及人口保護,防止外來事物可能衝擊或侵蝕本土生態.近年香港有不少「本土主義」冒起.

第一,強調保護本地人利益.從爭學位,搶床位,到購奶粉,買樓房,均涉及本地市民與內地居民之間的利益和資源分配問題.

第二,常以「香港的核心價值」將本港與內地資識別.香港市民崇尚人權,自由和擁護法治精神,而相比內地和世界各地,香港居民有幸享有高度的言論,結社,集會等各種人身自由,並受發展已上百年的良好司法制度保障.

第三,當中的「本土」以港式文化,價值及港人生活方式為基礎.一般而言,本土文化創意產品乃以廣東話為主的流行歌曲,電影,電視劇為主.

第四,著力保護「具本港歷史意義」建築,這是這思潮的一大現象.部分建築物確具其歷史價值及建築特色,值得保育.

然而,凡事自我中心,過度自我膨脹,愈走愈極端,完全不顧香港鄰近地區的發展,亦不願與內地省市合作交流,「本土主義」便會變成了Parochialism,即思想狹隘,堅拒對外開放的本位主義,非常危險.這種「自己顧自己」(inward-looking),或「自己人只顧自己人」,過度高估自我重要性的態度和主義實在與「全球融合」的潮流相違背,也與一向思想開放,樂於與外溝通的香港人格格不入.

第三種思維就是堅持己見,不容異己的人,引致「想殺死人」的思維.「就起來攆祂出城」:「借刀殺人」二十九節:「就起來攆祂出城,他們的城造在山上,他們帶祂到山崖,要把祂推下去.」為何會引發殺機?因為耶穌指出事實.第一個史實在第二十五節:「我對你們說實話,當以利亞的時候,天閉塞了三年零六個月,遍地有大饑荒,那時,以色列中有許多寡婦；以利亞並沒有奉差往他們一個人那裡去,只奉差往西頓的撒勒法,一個寡婦(無社會地位的人)那裡去.」(參王上十七1~15)第二個史實是第二十七節:「先知以利沙的時候,以色列中有許多長大痲瘋的；但內中除了敘利亞國(外邦敵人)的乃縵,沒有一個得潔淨的.』」(參王下五1~14.)

主在二十五至二十七節所舉神救助兩個外邦人的實例,其重點是說,過去以色列人拒絕了為神傳救贖信息的使者,神就差他們轉向外邦人.如今主耶穌就是這時代的救贖主,他們若拒絕祂,歷史還會重演(參十13~15；羅九至十一章).

第二十八節:「會堂裡的人聽見這話,都怒氣滿胸；」猶太人自認是神的選民,得天獨厚,擁有特權可以享受神的恩典,如今聽見耶穌似乎偏袒外邦人的話,就都甚為氣忿.這是排外的思想,並以為上帝是特別祝福他們,因上帝祝福亞伯拉罕.所以就會出現第二十九節:「就起來攆祂出城,他們的城造在山上,他們帶祂到山崖,要把祂推下去.」因為古時猶太對待假先知,乃是從數十呎的懸崖上將他推下去,並向他扔石頭.拿撒勒人從稱讚耶穌的態度變成恨惡,要把祂置於死地,乃因祂:沒有在那裡行過任何神蹟(參23節),達不到他們的期望；責備他們不尊敬本地的先知(參24節)；指出他們不是之處；向他們宣告以色列人未蒙恩,但外邦人反而蒙恩(參25~27節).故此,他們斷定祂是假先知,定祂死罪.傷害民族的情感:自己人竟不蒙恩.耶穌被扣帽子:假先知!古時,被扣為間諜,漢奸；現今就稱為恐怖份子,危險人物.他們認為耶穌把「肥水流向別人田」是:思想不正確,散播思想不正確的言論,以及這些思想不正確的言論會構成顛覆.保羅在加拉太時,也被稱為「那翻轉國家的傢伙」,他也自知在羅馬帝國,時日無多了!

這裡看到猶太人借刀殺人(殺自己人),加上魔鬼與耶穌間的角力.撒旦借會堂裡的人殺耶穌；會堂裡的人,又借群眾的力量殺耶穌；撒旦在背後發功,極其陰險!你曾陷過這種困局嗎?

借刀殺人目的是:滅聲.影片中,手槍裝上減聲器後,殺人不動聲色,一個個倒下,最使人驚慌的是:手槍無聲,連倒下的人都無聲,因為他們在短時間就死去.在今日社會上,常有滅聲出現,恐怖得很!

到底有誰想滅聲?第一是當權者,尤其是不容其他聲音的當權者；第二是敵對者,即不容異己者；第三是犯事者,因他們有不可告人之秘密怕被洩露.

常用滅聲的方法有很多.第一是禁聲,用盡方法,禁止人出聲.第二是滅聲,行酷刑,切去舌頭,拉出人的舌頭,甚至斬去人頭,即是滅口.第三是使人無聲無色消失,突然人間蒸發,不知所蹤.第四就是消滅入的思想,用不斷的洗腦方式.

到底耶穌如何面對滅聲,滅口呢?第三十節:「祂卻從他們中間直行,過去了.」原文的「直行」,可能是指行過,走過；「過去了」即離開了的意思.本節原文意思未明,有兩種可能:第一,主耶穌在他們面前顯出神蹟,行空走過懸崖,而未墜落崖下；第二,主耶穌的出聲的威儀鎮嚇住他們,令他們鬆手,就在他們鬆手之際,從他們中間泰然自若地走開了.

這次會堂裡的人滅不到口,當然也滅不了聲；滅不了歷史的事實；滅不了真理；滅不了耶穌,更滅不了上帝!縱然耶穌最後也死在十字架,卻能從死人中復活,耶穌並沒有被滅,反而數以億萬計的奉耶穌的名受洗；無數人向祂發聲,為祂發聲,代祂發聲.耶穌當日的發聲,被紀錄下來的變成文字,變成人聲.總言之,耶穌沒有因為會堂裡的人要滅聲而停止發聲.祂繼續發出正義之聲!

後現代的虛無力量

虛無主義作為哲學意義,為懷疑主義的極致形式.認為世界,生命(特別是人類)的存在是沒有客觀意義,目的以及可以理解的真相.虛無,這個詞最早來源於拉丁語中的「nihil」,意為「什麼都沒有」.尼采認為虛無主義是人們意識到「上帝已死」所導致的,並強調人們應通過重新賦予生活意義來克服它.‘Out, out, brief candle!’ 「熄滅吧,熄滅吧,短命的蠟燭!」‘Life’s but a walking shadow, a poor player ’「生命不過是一個行走的影子, 一個可憐的演員.」

現今都市人正面對後現代主義的一股「虛無力量」,這力量相當厲害,宜慎防之.筆者指出下列九大點,值得都市人深思,試述如下:

一. 無時間-使人失去生命力

二. 無空間-使人失去承載力

三. 無心間-使人失去心力

四. 無真理-使人失去方向

五. 無真我-使人失去本性

六. 無感動-使人變得麻木

七. 無大我-使人變得自私

八. 無整全觀念-使人自光短小

九. 無至善-使人變得腐朽

「虛無」的力量是撒旦的力量,它想人把「有」的變成「無」 .把無的當作「有」,這與上主的創造正正對立.「虛無」是使人真假難分,是非難分,光暗不明,立場不定,更危險是向罪惡妥協,明知故犯,並光明正大作壞事!都市人要活出豐盛的人生,實在要對付這「九大無」,又可稱為「九大鬼」,這些鬼經常在人的生命中作怪.

都市人要留意這三「無」,加以防範的那就是:

1對過往人生的錯誤,切勿說「無問題」；

2對現在人生的方向,切勿說「無所謂」；

3對未來人生的理想,切勿說「無可能」.

如果你陷入這「三無」,你很快也會消失於人間!

經文提醒我們要小心三種引致虛無的思維:

第一,        瞧不起人,導致不尊重人；

第二,        只顧自己人,導致排斥人；導致分黨派；導致小圈子,製造對立,對抗；

第三,堅持己見,不容異己的人,導致想殺死人,或導致想害人,想傷人,那些人不對咀型,不同意見,不同看法的人.這種人,會變成恐怖份子,傷害無辜的群眾,跟他們無仇無怨的群眾,特別是是與他們有仇的,他們卻殺不到,以傷害大量的人群來洩憤,來表達仇恨!

各位,聖經所提的這三類人,可能在你公司,機構,學校,甚至教會裡,我也發現這三類.然而,我們更要小心的,在我們在不知不覺中,也會被惡者所勝,也變成這樣的人!真要小心,那虛無的力量.



\newpage

\section{堅離地道}
\label{sec:311}
\textbf{林豪恩先生 - 「堅離地」道}
\newline
\newline
\hyperref[sec:310]{< < < PREV SERMON < < <}
~
\hyperref[sec:toc]{[返目錄]}
~
\hyperref[sec:312]{> > > NEXT SERMON > > >}


路加福音 9:28-36
\newline
\begin{longtable}{cl}
\hline
\hline
章節 & 經文 (和合本修訂版)\\
\hline
9:28 & \begin{tabularx}{0.7\textwidth}{X} 說了這些話以後約有八天，耶穌帶著彼得、約翰、雅各上山去禱告。 \end{tabularx} \\ \\ \relax
9:29 & \begin{tabularx}{0.7\textwidth}{X} 正禱告的時候，他的面貌改變了，衣服潔白放光。 \end{tabularx} \\ \\ \relax
9:30 & \begin{tabularx}{0.7\textwidth}{X} 忽然有摩西和以利亞兩個人同耶穌說話； \end{tabularx} \\ \\ \relax
9:31 & \begin{tabularx}{0.7\textwidth}{X} 他們在榮光裡顯現，談論耶穌去世的事，就是他在耶路撒冷將要完成的事。 \end{tabularx} \\ \\ \relax
9:32 & \begin{tabularx}{0.7\textwidth}{X} 彼得和他的同伴都打盹，但一清醒，就看見耶穌的榮光和與他一起站著的那兩個人。 \end{tabularx} \\ \\ \relax
9:33 & \begin{tabularx}{0.7\textwidth}{X} 二人正要和耶穌分離的時候，彼得對耶穌說：「老師，我們在這裡真好！我們來搭三座棚，一座為你，一座為摩西，一座為以利亞。」他卻不知道自己在說些甚麼。 \end{tabularx} \\ \\ \relax
9:34 & \begin{tabularx}{0.7\textwidth}{X} 說這些話的時候，有一朵雲彩來遮蓋他們；他們一進入雲彩就很懼怕。 \end{tabularx} \\ \\ \relax
9:35 & \begin{tabularx}{0.7\textwidth}{X} 有聲音從雲彩裡出來，說：「這是我的兒子，我所揀選的。你們要聽從他！」 \end{tabularx} \\ \\ \relax
9:36 & \begin{tabularx}{0.7\textwidth}{X} 聲音停止後，只見耶穌獨自一人。當那些日子，門徒保持沉默，不把所看見的事告訴任何人。 \end{tabularx} \\ \\
[1ex]
\hline
\hline
\end{longtable}
講題:「堅離地」道　Faith On a Cloud

經文:路加福音9章28-36節

講員:林豪恩先生

場合:香港中文大學崇基學院禮拜堂

日期:2016年2月7日

一．引言

有考究認為,「離地」一詞應該源於2012至2013年度財政預算案中,財政司司長曾俊華自稱中產的事件.自此之後「離地」成為近年流行的一個潮語,通常用作形容詞,形容某些人的想法甚至行為脫離普羅大眾,不食民間煙火,不知民間疾苦.他們雙腳離地,活於天際,高高在上.(香港網絡大典)

例句一:「就如一D 月薪三十幾萬ge 高官,都話自己係中產,連一碟粟米石斑飯都唔知幾錢ge 人,就絕對「離地」.」(SUN奇古怪:【離地】lei4 dei6)

例句二:【明報2015年4月26日】「港府昨啟動大型地區政改宣傳活動,主理政改的政務司長林鄭月娥率領30多名港府官員,乘坐掛上「2017,一定要得」橫額的開篷巴士,在港九新界多區巡遊宣傳政改,歷時約3小時.巴士以「快閃」方式駛過,官員全程「落區不落地」,只在車上揮手和嗌咪,並未落車接觸市民.泛民主派狠批官員落區的方式「離地」．」

「離地」這個潮語非常好用,當Ａ組別的人形容Ｂ組別的人離地的時候,Ｂ組別的人也認為Ａ組別的人離地:

例句三:「有班中產政客,個個大律師,完全唔知人間何世,諗埋D ye 都係意識形態居多,曲高和寡,唔好話牛頭角順嫂,分分鐘連中環白領都唔知佢地講乜,gwo D 咪「離地中產」囉!(SUN奇古怪:【離地】lei4 dei6)

當看見人們用「離地」或「堅離地」這潮語來互相描繪的時候,我們可能感到有趣．不過,你有沒有看見在網絡上有人用「離地」或「堅離地」來形容基督徒或教會呢?看到人們如此形容基督徒及基督徒群體時,你有甚麼感覺呢?

早前有人以「離地「耶撚」見證盡失」為題在網上發表文章,引起迴響．文中幾句說:「教徒們越被稱為「耶撚」,就越離地,不去好好回到地面聽聽人家的想法,最後則有更多人稱基督徒作「耶撚」,成了無限循環.」

今天的題目「「堅離地」道」的道,不是馬路,是使徒行傳中常說的「這道」,例如:「10:36上帝藉著耶穌基督-他是萬有的主-傳和平的福音,把這道傳給以色列人；17:11這地方的猶太人比帖撒羅尼迦的人開明,熱心領受這道,天天查考聖經,要知道這道是否真實.」這道,是「離地」的,還是不「離地」的呢?

二．經文

9:28說了這話以後約有八天,耶穌帶著彼得,約翰,雅各上山去禱告.

9:29正禱告的時候,他的面貌就改變了,衣服潔白放光.

9:30忽然有摩西,以利亞兩個人同耶穌說話；

9:31他們在榮光裡顯現,談論耶穌去世的事,就是他在耶路撒冷將要成的事.

9:32彼得和他的同伴都打盹,既清醒了,就看見耶穌的榮光,並同他站著的那兩個人.

9:33二人正要和耶穌分離的時候,彼得對耶穌說:「夫子,我們在這裡真好!可以搭三座棚,一座為你,一座為摩西,一座為以利亞.」他卻不知道所說的是甚麼.

9:34說這話的時候,有一朵雲彩來遮蓋他們；他們進入雲彩裡就懼怕.

9:35有聲音從雲彩裡出來,說:「這是我的兒子,我所揀選的,你們要聽他.」

9:36聲音住了,只見耶穌一人在那裡.當那些日子,門徒不提所看見的事,一樣也不告訴人.

三．釋經與應用

人們多在甚麼時候尋找宗教信仰呢?或許你也聽聞過,或者我們自己也是,多在面對生死,困頓,彷徨不安的時候．在失落的時候,我們尋找寄托,在那一刻,宗教信仰成為了我們心靈上的慰藉．正如有些人在帶領聚會時呼籲信徒:弟兄姊妹,現在我們來敬拜上帝,放下外面所有的紛亂和煩擾,放下個人的困難和重擔,放下人際關係的糾纏和恩怨,放下前景不明朗的焦慮．．．正如彼得所說:「我們在這裡真好!」既然「我們在這裡真好」,當然不想離開．離開,豈不是就要再次面對社會的紛亂和煩擾?再次擔起個人的困難和重擔?再次投入人際關係的糾纏和恩怨?再次感受前景不明朗的焦慮?因此,彼得說出了不少信徒的心聲:「可以搭三座棚,一座為你,一座為摩西,一座為以利亞.」我們便可以永遠停留在這「真好」的境界裡面．雖然作者認為彼得「不知道所說的是甚麼」,但是,從上下文來看當時的處境,彼得的建議是吸引的．

離地的平安與在地的危機(面對政治勢力的威脅)

9:7希律分封王聽見耶穌所做的一切事,就困惑起來,因為有人說:「約翰從死人中復活了.」

9:8又有人說:「以利亞顯現了.」還有人說:「古時的一個先知又活了.」

9:9希律說:「約翰我已經斬了,但這是甚麼人?關於他,我竟聽到這樣的事!」於是希律想要見他.

山下,是複雜的政治狀況．猶大地是羅馬帝國一個行省,由分封的諸侯希律王管治．為了穩固自己的利益及羅馬皇帝的政權,希律的管轄要在政治正確性極高的狀況下進行,因此,他要除去一切不穩定的因素．9:9希律說:「約翰我已經斬了」,現在再來一個耶穌,引起希律不安地問:「這是甚麼人?」下文(路加福音13章31節)提到,希律已經有了定案,繼施洗約翰之後,連耶穌也要殺掉.

另一邊的政治力量,是猶太人復國運動.在猶太人而言,耶穌是先知；就他的言行而言,他滿有能力；在群眾力量而言,他有號召力.因此,福音書提到,「眾人要來強逼他作王」． (約翰福音6:15節)

面對兩面的政治壓力,當時的猶太宗教會成為耶穌的保護嗎?事實是不會,因為猶太宗教領袖懼怕羅馬政治權力,因此也決定去耶穌而後快,免得添煩添亂.

與山上的寧靜祥和相比,山下充滿殺機．如果作為神祇的目的是接受敬拜,如果作為信徒的目的是獲得心靈慰藉,彼得那「夫子,我們在這裡真好!可以搭三座棚,一座為你,一座為摩西,一座為以利亞」的建議真是相當安全．

離地的自足與在地的不能(面對百姓生活的壓力)

9:12太陽快要下山,十二使徒進前來對他說:「請叫眾人散去,他們好往四面村莊鄉鎮裡去借宿和找吃的,因為我們這裡地方偏僻.」

9:13耶穌對他們說:「你們給他們吃吧!」他們說:「我們不過有五個餅,兩條魚,若不去為這許多人買食物就不夠.」

9:14那時,男人約有五千.

如果耶穌和門徒能夠逃避各方政治勢力施與的壓力,也總不能夠閉目不看民間疾苦．五餅二魚的事件,正正反映出百姓溫飽的需要．除了溫飽的需要,耶穌在路加福音第四章發表他的「使命宣言」時他提及:「貧窮人,被擄的人,失明的人,及受壓迫的人」．人間的需要彷彿無窮無盡,就正如掃落葉一樣,掃去一片,又來一片．門徒自知沒有能力為那五千人提供一餐,更何況長期而又多樣的各種需要呢?

與在山上的自足相比,山下令人感到無能為力．如果要避免「講就天下無敵,做就有心無力」的尷尬,彼得那留在山上的建議是自我感覺良好的．

離地的和平與在地的邪惡(面對邪惡力量的壓力)

9:37第二天,他們下了山,就有許多人迎見耶穌.

9:38其中有一人喊叫說:「夫子!求你看顧我的兒子,因為他是我的獨生子.

9:39他被鬼抓住就忽然喊叫；鬼又叫他抽瘋,口中流沫,並且重重地傷害他,難以離開他.

9:40我求過你的門徒,把鬼趕出去,他們卻是不能.」

山上與山下的對比實在太強烈了．山上的境況,彷如超人迪加主題曲的描述:「銀河唯一的秘密,天際最強人物,正氣朋友,性格忠實,英勇未變質．」摩西,以利亞,加上耶穌,天地正氣都集於一山．山下則是一個正邪混雜的地方．邪惡力量「如同吼叫的獅子,遍地遊行,尋找可吞吃的人.」(彼得前書5:8)「9:40我求過你的門徒,把鬼趕出去,他們卻是不能.」正義朋友對著邪惡力量無能為力,情何以堪呢?留在山上,就不需要面對這種尷尬及不安．

離地是吸引的,然而,人們對基督徒群體的指責,正是這種自我感覺良好的離地心態和表現．話說回頭,門徒這次離地經驗,其實是耶穌製造給他們的,是耶穌帶他們上山離地．耶穌製造離地經驗給門徒,到底所為何事呢?

離地的道

上文提到希律王問耶穌是誰,其實所有人都在問耶穌是誰．有人猜他是以利亞,有人猜他是摩西,還有人猜他是施洗約翰．耶穌和他的門徒也聽聞這些傳言,他也和門徒談及這項話題:

9:18耶穌自己禱告的時候,門徒也同他在那裡.耶穌問他們說:「眾人說我是誰?」

9:19他們說:「有人說是施洗的約翰；有人說是以利亞；還有人說是古時的一個先知又活了.」

9:20耶穌說:「你們說我是誰?」彼得回答說:「是上帝所立的基督.」

在山上,門徒親眼見到摩西和以利亞:

9:30忽然有摩西,以利亞兩個人同耶穌說話；

摩西,以利亞和耶穌三者同時出現,解答了上文所有人的疑問．耶穌不是摩西,不是以利亞,耶穌就是耶穌．那麼,對於門徒來說,耶穌是誰呢?

9:34說這話的時候,有一朵雲彩來遮蓋他們；他們進入雲彩裡就懼怕.

9:35有聲音從雲彩裡出來,說:「這是我的兒子,我所揀選的,你們要聽他.」

耶穌帶門徒離地,並不是帶他們去拜見摩西和以利亞這些英雄人物,也不是帶他們去避世,而是帶他們去退修／進修．「這是我的兒子,我所揀選的」,在山上,門徒親耳聽到從天上來關於耶穌的身份的宣告；更重要的,是關於門徒與耶穌的關係的指示:「你們要聽他」．這次離地的經驗對門徒來說是重要的,「從前風聞有你,現在親眼看見你」．這次離地的經驗,是門徒信仰進深的機會,是門徒領受使命的機會．

對耶穌而言,這次離地,也是他再次確認使命的時刻:

9:31他們在榮光裡顯現,談論耶穌去世的事,就是他在耶路撒冷將要成的事.

面對離地的指責,反駁之前,讓我們回首再問,我們信奉的這道是離地的嗎?答案同時是是,也是不是．這道本不在地,就此而言,這道是離地的:

1:1太初有道,道與上帝同在,道就是上帝.

1:2這道太初與上帝同在.(約翰福音)

然而,這道卻臨到大地,成為「在地」的道．

1:14道成了肉身,住在我們中間,充充滿滿地有恩典有真理,我們也見過他的榮光,正是父獨一兒子的榮光.(約翰福音)

四．總結

上山是為下山．耶穌帶領門徒離地,其實是為了在地．對耶穌而言,他絕對沒有離地還是在地的困擾,離地與在地之間的爭論只出現於跟隨他的群體．其實,這也沒有甚麼好爭論,地是被道創造的,這道當然不等同於地,也不出於地；然而這道的跟隨者在地生活,也在地見證,因此,不可能離地．問題是「地」並非固定不變,我們所站的「地」在時空之下不停地改變．基督徒群體在此時此地被指責離地,未必因為我們不實踐信仰,可能只是實踐信仰的觀念和行動停留了在過去的某一時空而已．網絡本是暢所欲言之地,面對絡網言論,有時「認真就輸了」．不過,連溫文而且忠誠於教會的神學院老師也對基督徒群體提出離地的指責,就值得我們留心及反省了．這位老師的慷慨之言其中一段如此說:

「對很多仍有理想,不願跪拜和臣服強權的人來說,近年香港眾多教牧,信徒領袖,教會和神學院界的表現讓人心痛和失望.香港自梁振英上台後,隨著中國整體左風興起,並中國因強大而來的驕傲,虛偽,霸道和扭曲事實的面孔頻密地出現在我們眼前,眾多年輕人和有心人充滿鬱悶,感到被欺凌,感到無奈,感到沒有出路．．．雨傘運動前後,香港教會的主流走向了「被慎言,被閉口,怕得罪,怕爭議,怕激進,怕政治,怕面對現實」的光景.教會甘心「離場,離群,離德」的後果,比可能產生的撕裂的影響更為嚴峻．．．當香港人和香港基督徒被大氣候和政局問題深深困擾時,講台上,眾教會和領袖卻多選擇「離場,離群,離題」,甚至「離德」.在最困擾的議題上,非常困擾的市民和信徒得不到牧養,得不到真正多元,多角度,但重視真理,審視時代的思辨.」

這位老師提醒我們,「地」改變了．就算沒有他的提醒,大家都有親身體會,我們的城市不停在改變,有些人甚至說:這不再是我們熟識的香港．「地」改變了,我們當如何在這地跟隨耶穌呢?

在文章結束時,這位老師向此時此地的信徒群體發出呼籲:

「愈是試煉的日子,愈是顯出我們內心素質的日子.讓我們將這最壞的時候,轉化成最美麗的時候.香港教會,中國大陸教會,台灣教會,海內外的華人教會,在這個世界關注華人的時代,讓我們的美善,勇氣,直言,愛心,對神的信靠和忠誠,能在這個充滿灰心,擔心,困擾的時代裡作光作鹽,成為神與人喜悅的見證.」(蔡少琪:基督徒只能有理想,不能甘心成為強權的奴隸)



\newpage

\section{人禍意外天譴}
\label{sec:312}
\textbf{鄧瑞強博士 - 人禍,意外,天譴}
\newline
\newline
\hyperref[sec:311]{< < < PREV SERMON < < <}
~
\hyperref[sec:toc]{[返目錄]}
~
\hyperref[sec:313]{> > > NEXT SERMON > > >}


路加福音 13:1-9
\newline
\begin{longtable}{cl}
\hline
\hline
章節 & 經文 (和合本修訂版)\\
\hline
13:1 & \begin{tabularx}{0.7\textwidth}{X} 正當那時，有些在場的人把彼拉多使加利利人的血攙雜在他們祭物中的事，告訴耶穌。 \end{tabularx} \\ \\ \relax
13:2 & \begin{tabularx}{0.7\textwidth}{X} 耶穌對他們說：「你們以為這些加利利人比其他的加利利人更有罪，所以受這害嗎？ \end{tabularx} \\ \\ \relax
13:3 & \begin{tabularx}{0.7\textwidth}{X} 我告訴你們，不是的！你們若不悔改，都同樣要滅亡！ \end{tabularx} \\ \\ \relax
13:4 & \begin{tabularx}{0.7\textwidth}{X} 從前西羅亞樓倒塌，壓死了十八個人，你們以為那些人比一切住在耶路撒冷的人更有罪嗎？ \end{tabularx} \\ \\ \relax
13:5 & \begin{tabularx}{0.7\textwidth}{X} 我告訴你們，不是的！你們若不悔改，都照樣要滅亡！」 \end{tabularx} \\ \\ \relax
13:6 & \begin{tabularx}{0.7\textwidth}{X} 於是，耶穌用比喻說：「有一個人在葡萄園裡栽了一棵無花果樹。他前來在樹上找果子，卻找不到， \end{tabularx} \\ \\ \relax
13:7 & \begin{tabularx}{0.7\textwidth}{X} 就對園丁說：『看哪，我這三年來到這棵無花果樹前找果子，竟找不到。把它砍了吧，何必白佔土地呢？』 \end{tabularx} \\ \\ \relax
13:8 & \begin{tabularx}{0.7\textwidth}{X} 園丁回答：『主啊，今年且留著，等我在樹周圍掘開土，加上肥料， \end{tabularx} \\ \\ \relax
13:9 & \begin{tabularx}{0.7\textwidth}{X} 以後若結果子便罷，不然再把它砍了。』」 \end{tabularx} \\ \\
[1ex]
\hline
\hline
\end{longtable}
講題:「人禍,意外,天譴」　Man-Made Disaster, Accident, God’s Punishment

經文:路加福音13:1-9

講員:鄧瑞強博士

場合:香港中文大學崇基學院禮拜堂

日期:2016年2月28日

各位弟兄姊妹,早安.

今年農曆新年前後,發生了兩場地震.

新年前,發生了台灣高雄地區的206地震.年廿八,清晨時分,一場地震,震碎了多少人的美夢.原本是團圓的日子,對很多人來說,卻是永別.

大家都知道,地震是由於不同的地殼板塊的摩擦,擠壓,相撞而引起的.新年後,發生了另一場地震.這次地震,卻不是由於「地殼板塊」的運動造成,而是由於「政治板塊」的摩擦,擠壓,相撞而造成.這次地震是年初二清晨時分香港的旺角暴力事件.原本是喜慶的日子,但很多人卻頭破血流.

地震的恐怖,在於我們立足的大地－－我們以為堅固的大地,在裂開,在崩潰,在瓦解.當我們立足的大地在瓦解時,我們便不再知道應立足於何處.人在地震中,依附的任何物件,都在震動.你無法把持一些什麼,可令自己穩定下來.你腳下的大地,在顛簸,在搖動,在怒吼.「一切堅固的東西都煙消雲散了」.這句子原是＜共產黨宣言＞裡的一句句子,正好用來形容地震所帶來的分崩離析.

1995年1月,日本發生阪神大地震(或稱「神戶大地震」),受災地區滿目瘡痍,像大戰過後的情景.日本作家村上春樹,對這大地震作了深度沉思,寫出了一系列的短篇小說,後來這些小說結集成《神的孩子全跳舞》(上海:上海譯文出版社,2009)一書.在書中,村上春樹藉小說主人公說:「可也真是不可思議,地震這東西.我們從來都深信腳下地面是牢固不動的,甚至有『腳踏大地』這樣的說法.想不到有一天事情突然變得不是這樣.本應堅固的地面,岩石竟變得液體一樣軟軟乎乎.在電視新聞上聽到了,是叫『液狀化』吧?」(《神的孩子全跳舞》,頁76)正是「一切堅固的東西都煙消雲散了」.在這小說裡,各人都因地震而改變生命.有人因為地震,不再接受現在的,不假思索的,了無新趣的生命.有人因為地震,不再想及明天,今朝有酒今朝醉,誰知地震明天不會在自己腳下發生呢?有人因為地震,常常回憶悲慘的畫面,無法投入生活.也有人因為地震,積極地尋索生命的意義,積極地尋索當大地分崩離析後雙腳仍可立足之處.簡言之,地震震碎了他們固有的生命,他們要在地震的廢墟中重新站起來.

台灣的高雄地震,救災的情景,令人動容.社會整體,無分藍綠,每人都盡上自己的本分,每人都獻上最好的禮物.有人免費送上飯食,有人免費提供住宿.沒有失蹤者被遺棄,那怕她只是一個精神失常的流浪者.每個離世者的遺體都被重建,而獲得應有的人性尊嚴.希望這次地震,能修補台灣的社會裂痕.願天父看顧台灣.

至於香港因不同政治板塊撞擊而造成的流血事件,則令人心傷.板塊相撞而形成的壓力,總是需要宣洩的.令人心傷的是,我們的社會最終竟然只能以這種方式宣洩政治壓力.這次政治性的,社會的地震,震碎了香港社會一直以來以為是堅固的社會地基.我們以為,就算一百萬人上街,都不會砸碎一個櫥窗.忽然,地震了,我們立足的大地,頓然崩潰了,留下了無底的深淵.這個深淵能否填補,現在不能知道.願天父看顧香港.

若我們身處大自然的地震的災場,看見生命的失去,我們將有何話可說?若我們身處政治的,社會的地震,看見有人失去生命,我們將有什麼想法?

若我們將這些事告訴耶穌,耶穌會如何回應呢?

讓我們看今日的講道經文:路加福音13:1－9.

路13:1　正當那時,有人將彼拉多使加利利人的血攙雜在他們祭物中的事告訴耶穌.

路13:2　耶穌說；「你們以為這些加利利人比眾加利利人更有罪,所以受這害嗎?

路13:3　我告訴你們,不是的!你們若不悔改,都要如此滅亡!

路13:4　從前西羅亞樓倒塌了,壓死十八個人；你們以為那些人比一切住在耶路撒冷的人更有罪嗎?

路13:5　我告訴你們,不是的!你們若不悔改,都要如此滅亡!」

路13:6　於是用比喻說:「一個人有一棵無花果樹栽在葡萄園裡.他來到樹前找果子,卻找不著.

路13:7　就對管園的說:『看哪,我這三年來到這無花果樹前找果子,竟找不著.把它砍了吧,何必白佔地土呢!』

路13:8　管園的說:『主啊,今年且留著,等我周圍掘開土,加上糞；

路13:9　以後若結果子便罷,不然再把它砍了.』」

那個時候,有人將一件事情告訴耶穌,經文說:「有人將彼拉多使加利利人的血攙雜在他們祭物中的事告訴耶穌.」《新漢語譯本》的翻譯較簡明,說:「有一些人來告訴耶穌,有些加利利人在獻祭的時候被彼拉多殺了.」

加利利人,不是耶路撒冷人.耶路撒冷人較接近羅馬的權力中心,為了生存,已學曉與權力打交道.他們未必是建制派,卻與權力建制保持恰當關係.加利利人,相對耶路撒冷人而言,是草莽之人.草莽之徒,活在社會邊緣,容易成為革命分子.加利利,是革命分子的溫床.耶穌被稱為「加利利人耶穌」,這稱呼頗堪玩味.

經文提及,有一群加利利人去獻祭.在當時,獻祭之地就是耶路撒冷的聖殿.這群加利利人往耶路撒冷獻祭時,被當時統治耶路撒冷的殖民官彼拉多懷疑為革命分子.看來,這些人當場被殺了,他們的血和他們的祭物攙在一起.

這是一場人間悲劇,是一場社會地震.一群人,一群加利利人,一群社會的邊緣人,他們或許是平民,或許是革命分子,或許是抱有革命想法的平民,無論如何,他們當時的舉動,只是獻祭,只是一種宗教上的行動,只是表達生命的敬虔.這殖民官是殘酷的,或許他認定他們是革命分子,或許他不確定他們是革命分子卻要先下手為強,或許他根本只是「有殺錯無放過」而藉此去震攝他人.這是政治板塊的撞擊,一場造成流血的社會地震.

耶穌對這事有何看法?

當時,有另一件悲劇發生.經文提及,「西羅亞樓倒塌了,壓死十八個人」.若不計這是否涉及「豆腐渣工程」,這事純屬意外.台灣地震的傷亡,主要也是來自塌樓.塌樓造成人命傷亡,無論是意外,抑或天災,總令人向神發出疑問.

耶穌,這些人禍,意外,算是一種天譴嗎?

這些人禍,意外,甚或天災,耶穌,你有一種說法嗎?

我們聽聽主耶穌的回答.

路13:2　耶穌說；「你們以為這些加利利人比眾加利利人更有罪,所以受這害嗎?

路13:3　我告訴你們,不是的!你們若不悔改,都要如此滅亡!

路13:4　從前西羅亞樓倒塌了,壓死十八個人；你們以為那些人比一切住在耶路撒冷的人更有罪嗎?

路13:5　我告訴你們,不是的!你們若不悔改,都要如此滅亡!」

首先,耶穌為那些人禍,意外,天災的受難者的人格尊嚴,討回一種公道的看法.或者說,耶穌為這些已逝者洗刷污名,洗刷旁觀者可能給他們的種種污名.在耶穌看來,這些受難者不比其他人更有罪,不比其他人更值得死,更應該死,他們不是遭受天譴.耶穌保護了這些已逝者的人格尊嚴.

其次,耶穌的回答,改變我們處理問題的位置.面對人禍,意外,天災,我們可能會扮演審判官的角色,評論那些受害人.或許,我們會說,他們身為被懷疑的加利利人,為何這麼蠢,走去耶路撒冷呢?明知那是豺狼彼拉多的地盤,走過去不就是送死嗎?他們怎麼還不檢點,行動惹彼拉多的懷疑呢?這些人在人禍,意外,天災中傷亡,是純屬不幸,咎由自取,抑蠢鈍過人,抵死有餘呢?他們是否比人差了點什麼?或者終極地說,他們是否比別人更有罪?耶穌的回答,拒絕這一切問題,拒絕我們扮演審判官的角色.

與其問別人是否更愚蠢,更抵死,更有罪,不如看看自己到底活出個什麼生命來.在人禍,意外,天災中,有人死了,他們去到了人生的終點.耶穌以此來問我們,我們如何面對我們自己的終點.在生命的終點裡,我們將面對那位對生命作全盤審問的上帝.與其評論別人,審判別人,不如面對上帝的審問,檢視自己的人生.與其作一個審判官,不如作一個被審判者.這不是說,不問世事,走向內心,找一內在的安頓之處；而是說,將天下人的生死,看作是自己生命中的大事.我們不能作「花生友」,只作旁觀者,對他人指指點點,而是視他人的生死,與自己生命之所當行,所當活,憂戚與共.無視他人的生死,令生命蒼白.在他人的生與死的提示下,面向上帝對生命的審問,活出一種不會空白地生,空白地死的生命,才是正道.

面對人禍,意外,耶穌說:「你們若不悔改,都要如此滅亡!」

「都要如此滅亡」,耶穌這樣說,重點不是要去評價人禍與意外中受害者的人格與價值,而是要去指出死亡本身可能揭露的意義.死,可以輕於鴻毛,可以重於泰山.死,可以控訴生命的空洞,也可述說生命的永恆內容.那些受害人遭受的死,揭露了生命的空洞的一面.但是,生命可以不是這樣的,死亡也可以不是這樣的.人,總有一死.面對死,面對大限,也就是面對那位對生命作全盤審問的上帝.死亡面前,如何好好生活,以致在死亡臨到的那刻,不是感到生命的無限空虛,而是體會活過的這一生的無限意義呢?

如何能在大限臨到前好活一場呢?

耶穌說,要悔改.悔改,是生命的轉向.從面向塵世之人和塵世之物,轉向萬物以外的上帝.這種轉向,對不同的人有不同的意義.對有權力的人來說,就是不要倚仗權力去成就自己的詭計,要明白權力來自神,要向神負責.死亡一到,你的權力和計謀,對你有何益處呢?對擁有財富的人而言,要明白財富是為你和他人的幸福而擁有的,財富來自神的祝福,要向神負責.以金錢幫助有需要的人,就是眼目從金錢轉向神.對心中憂傷的人而言,心中不要緊緊抓住過去的傷痛,要仰望神,要將那些傷痛交給神,要以神的寬恕去面對令你傷痛的人.你的人生要向神負責,若你只注目自己過去的種種,怎能好好地大活一場呢?

要悔改呀,否則人生就是白活一場,死亡臨到時,只會感到空虛一片.

面向上帝而活的人,面對人禍,意外,天災的受害者,不會對死者指指點點,反而會負起生命的責任,在能力範圍內,力圖使以後的人不需要再經歷那種無意義的死亡.若不義叫人受苦,就要抗拒不義.若意外叫人受傷,就要力圖查找意外的根源,將意外的根源根治.若天災叫人受害,則要確保一切設施合乎規格,減低天災帶來的傷亡.

這是一個地震後的世界,一切都不再一樣了.你在咖啡館飲咖啡時,有人會向你開槍；你在劇院裡聽音樂時,炸彈會在你身旁爆炸；你在超巿買日常用品時,你會忽然一命嗚呼.耶穌說,要悔改,要面向神,要在死亡臨到前活一種有永恆意義的生命.

路13:6　於是用比喻說:「一個人有一棵無花果樹栽在葡萄園裡.他來到樹前找果子,卻找不著.

路13:7　就對管園的說:『看哪,我這三年來到這無花果樹前找果子,竟找不著.把它砍了吧,何必白佔地土呢!』

路13:8　管園的說:『主啊,今年且留著,等我周圍掘開土,加上糞；

路13:9　以後若結果子便罷,不然再把它砍了.』」

「無花果」與「葡萄園」,都代表神的子民.三年,代表生命已有足夠的成長時間.你已夠成長了,但是,作為神的子民,你的生命仍未悔改,仍眷戀塵世,仍只活出空洞的生命.簡言之,生命仍沒有應有的果子.園主說:「把它砍了吧,何必白佔地土呢!」既然這生命無用,讓他終結吧.這是對我們生命的審判.然而,管園的說:「主啊,今年且留著,等我周圍掘開土,加上糞；以後若結果子便罷,不然再把它砍了.」且慢,多留一會吧!這是神的恩典,忍耐,憐憫.

每一刻,都有人經驗人禍,意外,天災；每一刻,我們都面對死亡；每一刻,我們都面對神對我們生命的審問；每一刻,神都宣判我們的大限已到；每一刻,神又以祂的恩典和憐憫寬容我們,托著我們.

弟兄姊妹,在大限忽然臨到前,悔改吧.

最後,我以劉進圖弟兄的一段話作結.劉進圖弟兄,前明報總編輯,於2014年2月被凶徒斬了六刀,生命忽然徘徊於死亡的邊緣.他在《迎鋒而立－－刀後「劉」言》(香港:明報出版社,2015)一書裡這樣說:「除了身體的運作規律改變,住院的日子還帶來了思想,情緒和心靈深處的變化,我的感覺就像進了修道院,每天都在默想沉思,在探索一些過去因太過忙碌而無法了解的事情.如果朋友問我,有甚麼得著?我會說,得著很多,一是原來過去日忙夜忙的事情,有許多並不是真正重要的；二是原來過去習以為常的生活方式,並不是正常的；三是原來每天生活裡有很多值得感恩的人和事；四是原來這個社會有許多人比自己不幸,需要關懷.你也許會說,這些都是顯而易見的道理吧.是的,可是我過去一直看不見.」(頁143)

但願:榮耀歸於聖父,聖子,聖靈.阿們.



\newpage

\section{歸去來兮吾歸何處}
\label{sec:313}
\textbf{鄧瑞強博士 - 歸去來兮,吾歸何處}
\newline
\newline
\hyperref[sec:312]{< < < PREV SERMON < < <}
~
\hyperref[sec:toc]{[返目錄]}
~
\hyperref[sec:314]{> > > NEXT SERMON > > >}


路加福音 15:11-32
\newline
\begin{longtable}{cl}
\hline
\hline
章節 & 經文 (和合本修訂版)\\
\hline
15:11 & \begin{tabularx}{0.7\textwidth}{X} 耶穌又說：「一個人有兩個兒子。 \end{tabularx} \\ \\ \relax
15:12 & \begin{tabularx}{0.7\textwidth}{X} 小兒子對父親說：『父親，請你把我應得的家業分給我。』他父親就把財產分給他們。 \end{tabularx} \\ \\ \relax
15:13 & \begin{tabularx}{0.7\textwidth}{X} 過了不多幾天，小兒子把他一切所有的都收拾起來，往遠方去了。在那裡，他任意放蕩，浪費錢財。 \end{tabularx} \\ \\ \relax
15:14 & \begin{tabularx}{0.7\textwidth}{X} 他耗盡了一切所有的，又恰逢那地方有大饑荒，就窮困起來。 \end{tabularx} \\ \\ \relax
15:15 & \begin{tabularx}{0.7\textwidth}{X} 於是他去投靠當地的一個居民，那人打發他到田裡去放豬。 \end{tabularx} \\ \\ \relax
15:16 & \begin{tabularx}{0.7\textwidth}{X} 他恨不得拿豬所吃的豆莢充飢，也沒有人給他甚麼吃的。 \end{tabularx} \\ \\ \relax
15:17 & \begin{tabularx}{0.7\textwidth}{X} 他醒悟過來，就說：『我父親有多少雇工，糧食有餘，我倒在這裡餓死嗎？ \end{tabularx} \\ \\ \relax
15:18 & \begin{tabularx}{0.7\textwidth}{X} 我要起來，到我父親那裡去，對他說：父親！我得罪了天，又得罪了你， \end{tabularx} \\ \\ \relax
15:19 & \begin{tabularx}{0.7\textwidth}{X} 從今以後，我不配稱為你的兒子，把我當作一個雇工吧。』 \end{tabularx} \\ \\ \relax
15:20 & \begin{tabularx}{0.7\textwidth}{X} 於是他起來，往他父親那裡去。相離還遠，他父親看見，就動了慈心，跑去擁抱著他，連連親他。 \end{tabularx} \\ \\ \relax
15:21 & \begin{tabularx}{0.7\textwidth}{X} 兒子對他說：『父親！我得罪了天，又得罪了你，從今以後，我不配稱為你的兒子。』 \end{tabularx} \\ \\ \relax
15:22 & \begin{tabularx}{0.7\textwidth}{X} 父親卻吩咐僕人：『快把那上好的袍子拿出來給他穿，把戒指戴在他指頭上，把鞋穿在他腳上， \end{tabularx} \\ \\ \relax
15:23 & \begin{tabularx}{0.7\textwidth}{X} 把那肥牛犢牽來宰了，我們來吃喝慶祝； \end{tabularx} \\ \\ \relax
15:24 & \begin{tabularx}{0.7\textwidth}{X} 因為我這個兒子是死而復活，失而復得的。』他們就開始慶祝。 \end{tabularx} \\ \\ \relax
15:25 & \begin{tabularx}{0.7\textwidth}{X} 「那時，大兒子正在田裡。他回來，離家不遠時，聽見奏樂跳舞的聲音， \end{tabularx} \\ \\ \relax
15:26 & \begin{tabularx}{0.7\textwidth}{X} 就叫一個僮僕來，問是甚麼事。 \end{tabularx} \\ \\ \relax
15:27 & \begin{tabularx}{0.7\textwidth}{X} 僮僕對他說：『你弟弟回來了，你父親因為他無災無病地回來，把肥牛犢宰了。』 \end{tabularx} \\ \\ \relax
15:28 & \begin{tabularx}{0.7\textwidth}{X} 大兒子就生氣，不肯進去，他父親出來勸他。 \end{tabularx} \\ \\ \relax
15:29 & \begin{tabularx}{0.7\textwidth}{X} 他對父親說：『你看，我服侍你這麼多年，從來沒有違背過你的命令，而你從來沒有給我一隻小山羊，叫我和朋友們一同快樂。 \end{tabularx} \\ \\ \relax
15:30 & \begin{tabularx}{0.7\textwidth}{X} 但你這個兒子和娼妓吃光了你的財產，他一回來，你倒為他宰了肥牛犢。』 \end{tabularx} \\ \\ \relax
15:31 & \begin{tabularx}{0.7\textwidth}{X} 父親對他說：『兒啊！你常和我同在，我所有的一切都是你的； \end{tabularx} \\ \\ \relax
15:32 & \begin{tabularx}{0.7\textwidth}{X} 可是你這個弟弟是死而復活，失而復得的，所以我們理當歡喜慶祝。』」 \end{tabularx} \\ \\
[1ex]
\hline
\hline
\end{longtable}
主題:「歸去來兮,吾歸何處」　Home Sweet Home

經文:路加福音15:11-32

講員:鄧瑞強博士

場合:香港中文大學崇基學院禮拜堂

日期:2016年3月6日

各位弟兄姊妹,早安.

每個人都有歸家的渴望,每個人卻也有離家的衝動.我們的靈魂需要安頓,我們的肉體卻尋找放縱.家人的擁抱是溫暖的,但不夠新鮮刺激.

於是,各人以不同方式離開家園,闖蕩江湖.

今日的講道內容,是離家與歸家的故事.

講道經文:路加福音15:11－32

路15:11　耶穌又說:「一個人有兩個兒子.

路15:12　小兒子對父親說:『父親,請你把我應得的家業分給我.』他父親就把產業分給他們.

路15:13　過了不多幾日,小兒子就把他一切所有的都收拾起來,往遠方去了.在那裡任意放蕩,浪費資財.

路15:14　既耗盡了一切所有的,又遇著那地方大遭饑荒,就窮苦起來.

路15:15　於是去投靠那地方的一個人；那人打發他到田裡去放豬.

路15:16　他恨不得拿豬所吃的豆莢充飢,也沒有人給他.

路15:17　他醒悟過來,就說:『我父親有多少的雇工,口糧有餘,我倒在這裡餓死嗎?

路15:18　我要起來,到我父親那裡去,向他說:父親!我得罪了天,又得罪了你；

路15:19　從今以後,我不配稱為你的兒子,把我當作一個雇工吧!』

路15:20　於是起來,往他父親那裡去.相離還遠,他父親看見,就動了慈心,跑去抱著他的頸項,連連與他親嘴.

路15:21　兒子說:『父親!我得罪了天,又得罪了你；從今以後,我不配稱為你的兒子.』

路15:22　父親卻吩咐僕人說:『把那上好的袍子快拿出來給他穿；把戒指戴在他指頭上；把鞋穿在他腳上；

路15:23　把那肥牛犢牽來宰了,我們可以吃喝快樂；

路15:24　因為我這個兒子是死而復活,失而又得的.』他們就快樂起來.

路15:25　那時,大兒子正在田裡.他回來,離家不遠,聽見作樂跳舞的聲音,

路15:26　便叫過一個僕人來,問是甚麼事.

路15:27　僕人說:『你兄弟來了；你父親因為得他無災無病地回來,把肥牛犢宰了.』

路15:28　大兒子卻生氣,不肯進去；他父親就出來勸他.

路15:29　他對父親說:『我服事你這多年,從來沒有違背過你的命,你並沒有給我一隻山羊羔,叫我和朋友一同快樂.

路15:30　但你這個兒子和娼妓吞盡了你的產業,他一來了,你倒為他宰了肥牛犢.』

路15:31　父親對他說:『兒啊!你常和我同在,我一切所有的都是你的；

路15:32　只是你這個兄弟是死而復活,失而又得的,所以我們理當歡喜快樂.』」

一段相當長的經文,描繪出一個相當破碎的家.小兒子離家出走了,大兒子雖然身在家園,卻覺得自己像不受重視的奴僕,他否定那歸家的弟弟,也否定父親對自己的愛,他身在家,卻離家甚遠,無家可歸.維繫著整個家庭的,是父親,是懷著無比大愛的父親,是承受著羞辱與責難的父親.沒有這父親,這個家早已散了.

今日,我不討論大兒子了,只集中在小兒子身上.

路15:12　小兒子對父親說:『父親,請你把我應得的家業分給我.』他父親就把產業分給他們.

小兒子離家的心很強,他正尋找他心中的家園.他要求即時取得父親遺產裡屬於他的那一份.這要求是很過分的,因為父親根本未死,怎能在父親還在的時候分遺產.再者,父親還在,一切產業都是屬於父親的,小兒子怎能說其中一份是屬於他的.小兒子的這種做法,等同對父親說,我想你早點死,甚或說,我根本當你已死.在父親還在時,取走一部分父親的財富當成是自己的,基本上就是當父親「無到」,與父親劃清界線,從今以後,我行我路.

父親忍受了羞辱,容讓小兒子的離家.父親知道,強迫與禁止,不能約束兒子的心,也無法叫他體會到家的可貴.父親將產業分給這小兒子.

路15:13　過了不多幾日,小兒子就把他一切所有的都收拾起來,往遠方去了.在那裡任意放蕩,浪費資財.

急不及待,小兒子拿了一切屬自己的,往遠方去了.「有咁遠,走咁遠」.看來,他要割斷過去的一切,遺棄一切的親情,不再忍受這個家的半點牽連.或許,這個家給他太多規矩,他要狠狠地遠離這些規矩,要盡情放肆.或許,這個父親太多叮嚀,他不想再聽到父親的聲音,他要做回自己,一個想如何便如何的自己.或許,這個花天酒地的世界實在太吸引了,人不好好墮落一下,實在對不起自己.他遠離家園,遠離規矩,遠離限制,要多風流有多風流,要多揮霍有多揮霍.只要有錢可花,要多快樂有多快樂.不要談論明天,今朝有酒今朝醉便是了.

久不久,那聲音會來造訪我們.那聲音說:「放棄舊時的信仰吧!放棄那些婆婆媽媽的道德枷鎖吧!放棄那些所謂正直的想法吧!看看這個世界,這個花花綠綠的世界,我們就只有這個世界的了.有錢可以花,就盡情花吧!有放縱的機會,就盡情放縱吧!那管天塌下來,盡情尋歡作樂吧!」我們很難抗拒這聲音,因為這聲音講的比我在這裡講的,聽起來更真實更動聽.

上帝應許的福樂,對我們來說,沒有多大吸引力.世界應許的名成利就,對我們來說,吸引力大得多.世人大聲說,來吧,不要稀罕天父的慈愛,而要用自己的手,確立自己的價值.來吧,向世界證明,你是一個有價值的人,然後取走你應得的獎品.若世界認為美貌是好的,那你便用盡一切方法令自己美貌一點,無論那方法是塗些什麼在面上,吃些什麼,甚或是整容.若世界認為成績是重要的,那你便用盡一切方法令自己有多點成績,無論那方法是多補一點習,多取一個學位,甚或是「出貓」.小兒子告訴父親,我好想你早點死,我要用自己方式,過一個世人認為快樂的人生.或許,我們也想告訴天父,天父,請你遠離我,我想用自己方式,過一個世人認為快樂的人生.我們越聽從世界的召喚,我們離開心靈的家越遠.

路15:14　既耗盡了一切所有的,又遇著那地方大遭饑荒,就窮苦起來.

路15:15　於是去投靠那地方的一個人；那人打發他到田裡去放豬.

路15:16　他恨不得拿豬所吃的豆莢充飢,也沒有人給他.

若你想用自己的手去創造自己的幸福,去掌控自己的快樂,很好,但你很快發現,你總會有疲累的一刻,你總會有掌控不到的一刻,你總會有失手的一刻.你會發現,要用自己的手去掌握自己的命運,是很難的.你想走在成功的頂端,但原來要走到這一步,你要常常處在戰鬥狀態,將對手擊敗,要超越昨日的自己,又要防範不知從何而來的攻擊.你若將幸福與快樂押注在你的成功之上,則你只會進入一無盡的漩渦裡.

小兒子終於用盡他的財富了,也遇到他掌控不到的逆境.或許,他仍想用自己的手去抓住快樂,但時不我與,無情的命運打擊他.世界的法則是,你成功時,保證人人擁護你.當你失敗時,人人離棄你,這也是世界的法則.你自恃擁有某些東西,用這些東西可換取快樂,好.但當你自恃的東西失去時,別人看你連牲畜也不如.小兒子想吃豬食的東西,人也不給他.

他可能想著,有朝東山再起.或許,他要捱過這艱難,而看看命運會否再次眷顧他.或許,他就這樣過一生.什麼也有可能.

路15:17　他醒悟過來,就說:『我父親有多少的雇工,口糧有餘,我倒在這裡餓死嗎?

路15:18　我要起來,到我父親那裡去,向他說:父親!我得罪了天,又得罪了你；

路15:19　從今以後,我不配稱為你的兒子,把我當作一個雇工吧!』

在這人生的關鍵時刻,小兒子「醒悟過來」,有些英文聖經譯作「He came to himself」(KJV, RSV,ASV),這就是新約希臘文的基本意思,「他回到他自己」.他心中問:「我到底是誰?」這真是一個好問題,我到底是誰.小兒子開始醒覺,我不是由世界的成敗去界定的,我不是由別人的看法去界定的,我甚至不是由我對自己的評價去界定的,我是父親的兒子.無論我如今去到什麼地步,我仍是有一個父親的,我是他的兒子.這一種對生命的重新認識,使小兒子重新思考,家在何處.

小兒子打算歸家了,有趣的是,他在這裡作台詞.面對父親,他應該說什麼呢?他在盤算,應說什麼才好.很明顯,他根本不認識他的父親.他對父親的想法,全是根據世界上「等價交易」,「成王敗寇」的邏輯想像出來的.或許,他在外面闖蕩久了,以為家園裡的想法和在江湖上的想法是一樣的.輸了,便只能當一個奴隸.小兒子不相信無條件的接納,也不期望無條件的寬恕.失敗了,只好「死死氣」回去「扮死狗」,任由父親處置.

台詞想好了,便踏上歸途.

路15:20　於是起來,往他父親那裡去.相離還遠,他父親看見,就動了慈心,跑去抱著他的頸項,連連與他親嘴.

一切都出乎意外.小兒子想像的,與真實的情景相差甚大.

主動的,是父親,他主動的跑了過來,見到這衣衫襤褸的小兒子,還未待這小兒子講出他早已想好的台詞,父親已擁抱著他,親吻他.

接納克服了內咎,恩典拒絕了世間的計較思量,仁慈照亮了人性的黑暗,家的溫暖安頓了疲累的心靈.

這是神聖光輝永恆照耀的一幕.我們看見,神那無限的仁慈.神將開雙手,接納每個願意回家的人,以祂為家.

我相信,很多人一生想望的,就是讓父親這樣溫柔的一抱.自以為是走過的錯路,羞辱父親換來的內咎,自把自為帶來的生命恥辱,我行我路導致的無比失落,就在這一抱裡,化作無比溫暖.我們多麼渴望父親的這一抱.

荷蘭畫家林布蘭(Rembrandt van Rijn, 1606-1669)以「浪子回頭」這故事畫了一張巨畫,去捕捉父親擁抱小兒子這神聖一刻.靈修神學家盧雲(Henri Nouwen, 1932-1996)按著林布蘭這幅畫的提示,寫了一本書,叫《浪子回頭》.每個人都渴望被父親這麼一抱,然而,盧雲卻說,歸家的路甚難.

為什麼浪子歸家的路甚難?

盧雲在書中這樣寫:「我經年在屬靈生活的各方面指導學生,幫助他們明白屬靈生活的重要.可是,我自己可曾敢走進中央〔按:指林布蘭的畫的中央,父親擁抱小兒子的那一幕〕,跪下,讓饒恕的神擁我入懷?」(頁18)為何這麼難呢?盧雲說:「我很想控制自己的靈程,至少仍舊能預估一些成果；撇棄旁觀的安全,成為軟弱的歸家兒子,簡直太不可能了.」(頁19)原來,盧雲發覺,無論如何,他仍想以自己的方式操控自己的命運.他繼續說:「趨向畫作中央的每一小步,都像是達不到的一項要求:再次要求我放開掌控的意欲,再次撇棄預知生命的渴望,再次捻息不知這一切導向何方的恐懼,再次降服於不止息的愛.」(頁21)盧雲再說:「我的心思,感受,情緒,熱望,常常遠離神選擇為家的地方.回家,留在神的居所,傾聽真理與愛的話語,的確是我最怕的旅程,因為我知道神是忌邪的戀人,祂要的是我整個人.」(頁24)一個有所保留的人,就無法回家了.

我們渴望家,卻又會離家.離家後,會經歷浪子的起起跌跌.聖經故事裡的浪子,經歷生命的淒慘後,他想起自己是誰,他歸家去.但是,現實上,我們就算意識自己是浪子,也未必會回家,理由是我們不願意放下自己一手建立的成就,我們不相信我們生命的尊貴建基於天父的恩典,我們不敢不依賴自己的計算而只依賴天父的仁慈.因此,世上有無數疲累不堪而不回家的浪子.

路15:21　兒子說:『父親!我得罪了天,又得罪了你；從今以後,我不配稱為你的兒子.』

路15:22　父親卻吩咐僕人說:『把那上好的袍子快拿出來給他穿；把戒指戴在他指頭上；把鞋穿在他腳上；

路15:23　把那肥牛犢牽來宰了,我們可以吃喝快樂；

路15:24　因為我這個兒子是死而復活,失而又得的.』他們就快樂起來.

父親根本不在乎小兒子的台詞講什麼,因為他的心早已擁抱著這小兒子.無論這小兒子過去對父親是愛是恨,現在講的台詞是真心抑假意,他在世上是成功或失敗,對父親來說,這一切都不重要.在父親眼中,這個就是他的兒子.要將上好的袍子給他穿,恢復他的尊嚴.把戒指戴在他指頭上,恢復他的身分.把鞋穿在他腳上,因他不是赤腳的僕人.吃一頓美味的「肥牛肉」,這意味著一種類似天國的福樂.

你可能不相信,有這樣的父親.你可能不相信,你能成為這父親的兒女.是的,難以置信.但福音之為福音,就是有點難以置信的.現在要做的,是領悟自己的身分,努力回家.

但願:榮耀歸於聖父,聖子,聖靈.阿們.



\newpage

\section{何處覓蓬萊}
\label{sec:314}
\textbf{陳偉光教授 - 何處覓蓬萊}
\newline
\newline
\hyperref[sec:313]{< < < PREV SERMON < < <}
~
\hyperref[sec:toc]{[返目錄]}
~
\hyperref[sec:315]{> > > NEXT SERMON > > >}


講題:何處覓蓬萊　　Where Is That Fabled Fairy Isle?

講員:陳偉光教授

場合:香港中文大學崇基學院禮拜堂

日期:2016年3月13日

我今天的題目是《何處覓蓬萊》.「蓬萊」大概是一個夢想中的樂境,世人嚮往的烏托邦,但這地方究竟在那兒?「何處覓蓬萊?」一個充滿唏噓,無奈的問號,既問自己,又問蒼天.我第一次聽到這個問題,是少年時代,在許冠傑的歌裡面.歌詞不過是寥寥數十字,卻挑起了很多人心中的感慨,說出了雖然人人都「心嚮往蓬萊,想身處蓬萊」,但有誰曉得「蓬萊所在」?要到那天它才不再「似被浮雲掩蓋」?語調中的愁緒,勾起了大眾久別的當年,寫出了我們心中殷切期待的明天；既帶點苦味,又滿懷憧憬.我心裡在想,世人如此,那麼我們屬主的人又如何?我們的世界究竟是怎樣的?我們所盼望,所等待的,究竟是甚麼?我們所宣稱的國度,所尋找的家鄉,真的能夠說服自己,又說服別人嗎?

這些問題,我想了很多年,幾乎陪伴我走過整個信仰的歷程.我曾經把其中的一些感受,寫在兩年前崇基網頁上,同樣以《何處覓蓬萊》為標題的一篇我的網誌中.不過當日所寫的,只能說是這一份感動的起點.隨著日子的過去,隨著其間所發生的種種事情,我想或許是時候,讓我在信仰的反思中,把心中這份從未減退的感動,在今早重新整理一下.

網誌的開始,是憶述新聞報導中的一幕.容許我用一點時間,讀出這片段.

「噩耗傳來,父親在地盤意外中不幸離世,才十二歲的小女孩,勇敢地伴著母親,向那圖以數萬元草草了事的承建商,爭取應得的賠償[1].鏡頭下輕輕掠過的一幕,雖不過數秒的新聞片段,卻真使人心酸,教人不忍.小小的年紀,在殘酷的際遇中,再無暇等待別人的安慰與憐恤,只知道要堅強地保護媽媽,不讓人家欺負她.堅強,成熟,懂事,原來都是有代價的.不錯,這代價確實很大；然而,造物主是善良的,公義的上天給他們預備的賞賜既豐且隆.我碰過多少滿懷熱誠與真情的年青人,他們原來都是在苦難的灌溉,困乏的滋潤中茁壯成長.看他們笑得燦爛時仍稚氣未減,但每當遇上關乎信念與價值的挑戰,他們都會沉著機智而毫不妥協地堅守壁壘.這些都是我所教過最優秀的學生,男的有為上進,女的聰慧可人,不少更以一等榮譽畢業,再當選全校優異畢業生代表.如此出色的一群,大可為自己選擇一條更寬闊的路,可是,我眼見今天的他們,還是在自己所認定的使命中,奮力掙扎,編寫著不一樣的故事.他們中有專注保護本土基層權益的法律工作者,也有選擇到第三世界國家服務的醫療人員.少時艱難的歲月,不幸的際遇,不但沒有為他們帶來詛咒,反成就了意想不到的祝福.這種福氣,是上天為那配受苦難的人預備的專利[2],也是那些生活在無驚無憂,無風無浪中的孩子,所無緣領略的.」

重讀當日所寫的,我心中依然激動.世間的天國,人間的樂境,在一群不向命運低頭的孩子身上,忽然顯得輪廓鮮明,指日可待.但與此同時,我心裡更湧起一種莫名的納悶,究竟當日我所看見的祝福,就是那份化身作逆境的福氣,是否必然地能夠表現在如此高尚的情操上?昔日瑰麗的情操,今天是否已日漸褪色?各位,你有沒有留意到,曾經戰勝窘厄,走出陰霾的人,總喜歡相信,昔日能熬過非常的日子,跨越漫長的流淚谷,全憑自己心中的一份信念,相信世界不論變到那個樣子,人心那怕墮落到更壞的境況,詭詐還是永遠不能勝過良善,狂傲也定要在謙卑前降服,終有一天,真理與公義必要被高舉.各位,是不是聽起來非常動人?但微妙地,可怕地,只需要一點點的偏差,這信念足以叫我們的世界再容不下別人,這信念甚至足以令人相信,只要一天我們曾拒絕屈服強權,我們就可登上正義的高台,當上真理的代言人,高呼命運只許掌握在自己手中,人要對抗逆境,擺脫厄運,就得勇猛駕馭,強勢出擊,憑著一己堅定不屈的意志,在渺無生機的絕境中,創造明天,建設理想中的烏托邦.

魔鬼的詭計實在無處不在,就是這一絲微妙而隱晦的分別,足以叫我當日所讚嘆的華麗,所歌頌的「逆境中的堅毅,艱苦中的奮進」,在不經意間,變成了自義,自恃.可知道,這種變了質的「堅毅奮進」,今天已被包裝成漂亮的理想,就如廣告標語中所推銷的價值觀:「我們明知會輸,都一定要贏」,又有「命運在我,自己想點,自己話事」[3],更有政治領袖一邊呼籲群眾要「謙卑,謙卑,再謙卑」,另一邊又宣告「從今以後,再沒有人需要為自己的認同而道歉!」全都是非常動聽的說話,非常吸引的想法,都是我們期待已久的「肯定」,是公開而正面地,肯定我們的「自我中心」.我們感覺良好,心中暢快,但所為的,已經不再是清澈如昔的理想,而是一份自覺良善的優越.一線之差,天地之別.

我所描述的,其實是一個好像我這種一生走走跌跌,載浮載沉,心靈願意,肉體軟弱的信徒的故事.就如只要某天,我曾經有過一絲的真情和熱誠,我的焦點很快就會轉到別人的冷酷與無情,恨不得世上懂得「行公義,好憐憫」的,就只剩下我一人.多少時候,就在這一線之間,就在這種「眾人皆醉我獨省」的感覺裡,我成了魔鬼的俘虜.

還有更多時候,魔鬼叫我們心裡不甘,眼見別人做得轟轟烈烈,我的上帝卻不曾華麗登場,這樣等下去再不是辦法,唯有決定自己出手.看,我們真的需要很大的恩典和憐憫,讓我們在跟世人一起擁抱真理,愛慕良善,追求美麗光明的事物的同時,能夠展現出一種與別不同的生命素質,向世人見證我們所認識的「堅毅奮進」,是建基於我們知道唯有上帝是一切的主宰,萬物眾生不過是祂手中的工作.我們對抗不義,或可到流血的地步；我們堅守正直,或可不顧性命,但能掌管明天的至終只有上帝.當此人生種種,吊詭的是我們雖得全力以赴,卻深知事情的關鍵全不在我.如此,就任它風雨飄搖,境遇變遷,那怕純真被輕看,友愛遭踐踏,美意給佔盡便宜,善工落得全盤失敗,我們仍然可以心裡平安,「與素常一樣」[4],過著渴慕天父性情的生活,滿懷感激地在祂為我預備的挫敗中繼續奮進,仍舊做著該做的事情.

世人堅毅,我也奮進,逆流之中,似是同舟.彼我之別,只在追隨基督的人,原來從不相信自己,卻唯獨依賴上帝,「信任祂安排」[5].有沒有留意到,在剛才我為大家選唱的詩歌 “I Lift My Eyes”《舉目遙觀》中,就有這麼一句:「信任祂安排,蒙祂保守心安定」,一份在信徒生命中,也是我在聖詩文學中所找到,最高貴的氣質和情操.詩篇的作者說:「我曾耐心等候耶和華」[6] (詩四十:1),所描繪的,也正是這種的情操,就是在聽不見祂聲音,看不見祂蹤影,無法確定祂是否還在我們中間的日子,我仍要等候祂.「等候」乃基督的追隨者一種很特別的氣質,我們等候,卻不無奈；遠景模糊,心仍踏實.不論是那種「天天有人來幫助大衛[7]」的光境 (歷上十二:20),還是那些「耶和華的言語稀少[8]」的日子 (撒上三:1),或在平順的旅途中,或在磨難的歲月裡,我們依然耐心等候祂.世人說「信心」難以捉摸,我們就拿「等候」來向他闡明；世人指「盼望」不過自欺,我們就拿「等候」來給他示範；就是世人對「愛仇敵」感到大惑不解時,那「曾耐心等候耶和華」的人亦只會說,這一切都不過是上帝的工作,仇敵既是上帝所賜的,也就是我的鄰舍.這樣看來,那不幸中仍能奮進,被置於低微而沒有心裡不甘,終久不見成果而不至心存苦毒的,都非甚麼人間異類,他們只因知道上帝的道路更高,就樂意一心順服,在等候中繼續奮力前行.而這一切,原來都是在「信,望,愛」裡,才能得以圓滿.

渴慕天父性情的人在那裡,那裡就有「真,善,美」；何處得見等候耶和華的人,「信,望,愛」就在其中.「真,善,美」使人堅毅奮進,「信,望,愛」教我任祂安排.如此,在順服裡堅剛,在靜候中奮進,實乃真理的兩面,猶如劍之兩刃,二而為一.然而,令人遺憾的是,在今天的教會裡,這忠勇堅毅,熱誠奮進的行列,或許仍然聲勢浩大,旗幟飄揚,但那甘心順服,耐心等候的門徒,卻已不見幾人.

問題的癥結,在於人們始終執意相信,若真理仍在,浩氣尚存,那麼公義就必須在地上確實地顯明自己,讓人在日光之下,眾目之前,得以同證它的真實.且看今天,多少人已再沒有那種耐性,等待甚麼遙遠的家鄉.他們只知道,今生的選擇,理應在今生得見果效,得償報應.今世的種種,若誠然皆在上主的許可中而成,那末祂不也應該就在今世,為祂如此種種的判斷,來給我們一個圓滿的交代嗎?

剛才我們讀過聖經約翰福音第八章的一段記載,正好把我們帶到耶穌時代一個相似的場景.當日的文士和法利賽人,從一開始就是我們眼中的壞人,但換一個角度看,很多時候他們不過是勇於主持正義,公開指斥律法道德所不容的事情.我想像他們可能都真誠地相信自己正在替天行道,全心全意地,徹徹底底地,做著當做的事情.他們的熱誠,他們的義憤,大概都是值得嘉許的,但非常可惜,就是因為這一線之差,就是因為這一點點隱晦而微妙的分別,他們都由「忠心的僕人」,變成了「毒蛇的種類」[9].聖經說:「耶和華不像人看人,人是看外貌,耶和華是看內心」[10].原來關鍵全在我們內心的態度,全在我們在「不喜歡不義,只喜歡真理」的時候,是否已經忘記了自己也是一個罪人,忘記了「在我裡頭沒有良善」[11],以至我們在揭破別人罪行,當上仲裁者的一刻,自己竟然已經做了上帝.耶穌當日對這群捉拿淫婦的人所說的話,其實是對我們說的:「你們中間誰是沒有罪的,誰就可以先拿石頭打她」[12].看來我們也應該「從老到少,一個一個的都出去」[13]!請留意,這時候,那位真正無罪的耶穌,對我們這些可憐的罪人,只說了一句話,一句驚天動地,改寫你我終極命運話:祂說:「我也不定你的罪!去吧,從此不要再犯罪了」[14].

回過頭來,檢視我們做主門徒的歲月,且別說曾否為公義出頭,以至在不經意間,變得自義,自恃,自覺優越.就是在個人的層面上,在自覺「艱苦奮進」的日子中,我們何嘗不是苦苦糾纏,死不放手.試想,有誰甘心在自己有生之年,讓一切苦了一輩子的事情,都可就此作罷?世事縱然不能盡如人意,但對於那決意「堅忍耐煩,任勞不避[15]」的人來說,期望努力終能得到一點回報,也實不為過.無怪乎人都羨慕Thomas More筆下的烏托邦,因為它為悲苦的人生創造了一個輪廓鮮明的夢想,而這夢想更可望有一天,能真正得以實現人間.

在信主的人來說,我們也曾艱苦經營,用心栽種,窮一生所有,極力之能及,天天做個忠心的門徒.我們屬主的人,當然更有理由相信天下萬物既然由祂統管,努力進入天國的人,必能得見天國實現眼前.可是,在現實的經驗中,我們的道路卻沒見得比別人的通達,日子也未嘗比別人的好過.尤有甚者,在那些我們最為關切的事情上,就好比我們的孩子,愛侶,前途,事業等,往往投入的心力越大,換來的無奈也更深.

是天父對我特別嚴苛,要狠狠地把我教好,讓我痛定思痛,還是祂對我特別器重,要把我進行地獄式訓練,苦其心志,勞其筋骨?

記得電影【加菲貓】中的名句嗎?”Tomorrow will be a better day!” 勵志之語,想必無人異議.可是若把它當真,以為人只要努力耕耘,天國必可實現人間,那恐怕我們都要注定失望.危險的是,世人總愛憑一己良好的意願,繼續這樣想下去.看英國大文豪William Blake如何歌頌那「在地獄似的絕境中,重建人間的天國[16]」的使命,就可見類似的情懷,在世人的心裡,已不成疑問地,被視作人間最高貴的情操.

一線之差,足叫人迷惘一生,魔鬼的詭計本來就是這樣.原來上帝既非把我嚴懲,亦非給我特訓,祂只是不忍看著我糊里糊塗地跟世人一樣,真的以為天國終會在人間實現!

這是何等的智慧,讓我們在這不會有完美答案的寄居之地,細味甚麼是沒法持久的健康,苦無出路的困局,不常相伴的笑臉,難望復和的關係,還有那更多更多終久難圓的夢,我們好在其中領略上帝所賜的「豐盛生命」,是多麼與世人所祈望的截然不同.

也就是這種超然的「豐盛」,每天都提醒著我們,「這世界非我家」[17],我心仍要等候耶和華.

八十年前,美國印第安那州一個寧靜的旁晚,作曲家 Virgil Brock 相約朋友在湖邊聚舊,漫步暢談.這時候忽然有人發現遠處風雲色變,看來暴風雨頃刻將至.大家心裡都有點惆悵,恐怕眼前又是一個雷電交加的晚上.不過很奇怪,這時候,各人的目光忽然又回到眼前的景象,原來憂思愁緒或可一刻擁上心頭,卻不曾擾動這一幅優美寧靜的圖畫,湖水依然平靜如鏡,在夕陽照射下,金光閃耀,仿佛上帝臨在,親手保護；想到世事多變,儘管厄運當前,不論面對任何境遇,屬主的人大可無懼凶險,深知是誰鋪張穹蒼,統管萬有.

大家正沉醉於此刻美景,當中一位自小雙目失明的朋友,忽然非常感慨,禁不住大聲的說:「我從未見過像眼前般美麗的日落!」大家當然知道他其實沒可能真的看見此番景緻,所以對他當下如此雀躍,都感到十分詫異.作曲家 Virgil Brock 連忙回應說:「這句話出自你的口,實在太奇妙!」”People are always amazed when you talk about seeing.” 此人雖眼不能見,心中卻如此激動,確實令在場各人都讚嘆不已.想不到這位自小雙目失明的朋友,這時候竟然很認真地跟 Virgil Brock 說:「我真是看見的!」“Yes, I can see! I see through other people’s eyes.” 「是透過你們的眼睛,讓我看見的.」“And I think I often see more.” 「我深至覺得,說不定我看見的更多.」”I see beyond the sunset!” 不錯,我們都看見「日落」,但我們不是每一個都看見「日落的那邊」.

Virgil Brock就在當天晚上,創作了這一曲 “Beyond the Sunset”,也就是剛才由 Jeannette 和 Francis所唱出的《斜陽日影去》.這個新的音樂版本,是我今早特意為各位編寫的.我希望八十年前那動人的一刻,今早能叫我們再次心裡熾熱；我期待那昔日曾令人心醉的清純皎潔,今天可重臨這片「混濁之地」[18]；我深願歌曲中每一句曾打動我的詩詞,此刻再向眾人說話.

昔日我也曾「心嚮往蓬萊,想身處蓬萊」[19],今日方知此處本非蓬萊所在.曾經無奈,人才懂得耐心等待,這才恍然驚覺,原來除祂以外,我心別無所愛.[20]

[1]     意外發生於二零一一年七月八日,事件在翌日被傳媒廣泛報導.

[2]     這裡所指的,是那些不會在艱難的歲月裡,終日怨天尤人,自暴自棄者,亦可比作聖經中那「忍受苦難的人」.

[3]     廣東話口語用詞.

[4]     但六:10

[5] 「信任祂安排,蒙祂保守心安定.」這種基督追隨者的氣質,在Timothy Dudley-Smith 的詩詞中表現得淋漓盡致,其中最廣為人知的 “I Lift My Eyes” 《舉目遙觀》一詩,可見於「恩頌聖歌」第408首.

[6]     詩四十:1

[7]     歷上十二:20

[8]     撒上三:1

[9]     太十二:34

[10]   撒上十六:7

[11]   羅七:18

[12]   約八:7

[13]   約八:9

[14]   約八:11

[15]   節錄自國父孫中山勸勉年青學子的名句:「人生的事,其不如意者,固十常八九.總要在堅忍耐煩,任勞不避,才能祈於有成」.

[16]   “Love seeketh not Itself to please, Nor for itself hath any care; But for another gives its ease, And builds a Heaven in Hells despair.” (William Blake) 『「愛」是不求自己的益處,不戀慕在施贈中得享優越；是「愛」使人甘作他人的僕役,在地獄似的絕境中,重建人間的天國.』(陳偉光譯)

[17]   “This World Is Not My Home” 這是一首著名的黑人靈歌,歌詞可見於「生命聖詩」第485首.

[18]   見「生命聖詩」第485首《舉目遙觀》第一段末.

[19] 節錄自許冠傑的名曲【何處覓蓬萊】,首段歌詞如下:「心嚮往蓬萊,想身處蓬萊,仰望蓬萊所在,似被浮雲掩蓋.」

[20]   「除你以外,在天上我有誰呢?除你以外,在地上我也沒有所愛慕的.」(詩七十三:25)



\newpage

\section{復活節信息生命的意想不到}
\label{sec:315}
\textbf{許開明牧師 - 復活節信息:生命的 “意想不到!”}
\newline
\newline
\hyperref[sec:314]{< < < PREV SERMON < < <}
~
\hyperref[sec:toc]{[返目錄]}
~
\hyperref[sec:316]{> > > NEXT SERMON > > >}


路加福音 24:1-12
\newline
\begin{longtable}{cl}
\hline
\hline
章節 & 經文 (和合本修訂版)\\
\hline
24:1 & \begin{tabularx}{0.7\textwidth}{X} 七日的第一日，黎明的時候，那些婦女帶著所預備的香料來到墳墓那裡， \end{tabularx} \\ \\ \relax
24:2 & \begin{tabularx}{0.7\textwidth}{X} 發現石頭已經從墳墓滾開了， \end{tabularx} \\ \\ \relax
24:3 & \begin{tabularx}{0.7\textwidth}{X} 她們就進去，只是不見主耶穌的身體。 \end{tabularx} \\ \\ \relax
24:4 & \begin{tabularx}{0.7\textwidth}{X} 正在為這事困惑的時候，忽然有兩個人站在旁邊，衣服放光。 \end{tabularx} \\ \\ \relax
24:5 & \begin{tabularx}{0.7\textwidth}{X} 婦女們非常害怕，就俯伏在地上。那兩個人對她們說：「為甚麼在死人中找活人呢？ \end{tabularx} \\ \\ \relax
24:6 & \begin{tabularx}{0.7\textwidth}{X} 他不在這裡，已經復活了。要記得他還在加利利的時候怎樣告訴你們的， \end{tabularx} \\ \\ \relax
24:7 & \begin{tabularx}{0.7\textwidth}{X} 他說：『人子必須被交在罪人手裡，釘在十字架上，第三天復活。』」 \end{tabularx} \\ \\ \relax
24:8 & \begin{tabularx}{0.7\textwidth}{X} 她們就想起耶穌的話來。 \end{tabularx} \\ \\ \relax
24:9 & \begin{tabularx}{0.7\textwidth}{X} 於是她們從墳墓那裡回去，把這一切事告訴十一個使徒和其餘的人。 \end{tabularx} \\ \\ \relax
24:10 & \begin{tabularx}{0.7\textwidth}{X} 把這些事告訴使徒的有抹大拉的馬利亞、約亞拿，和雅各的母親馬利亞，還有跟她們在一起的婦女。 \end{tabularx} \\ \\ \relax
24:11 & \begin{tabularx}{0.7\textwidth}{X} 她們這些話，使徒以為是胡言，就不相信。 \end{tabularx} \\ \\ \relax
24:12 & \begin{tabularx}{0.7\textwidth}{X} 彼得起來，跑到墳墓前，俯身往裡看，只見細麻布，就回去了，因所發生的事而心裡驚訝。 \end{tabularx} \\ \\
[1ex]
\hline
\hline
\end{longtable}
講題:《生命的意想不到!》　Easter Message: “The Unexpected” in Our Lives

經文:路加福音24章1-12節

講員:許開明牧師

場合:香港中文大學崇基學院禮拜堂

日期:2016年3月27日

「1七日的頭一日,黎明的時候,那些婦女帶著所預備的香料,來到墳墓前.2看見石頭已經從墳墓滾開了．3他們就進去,只是不見主耶穌的身體.4正在猜疑之間,忽然有兩個人站在旁邊．衣服放光.5婦女們驚怕,將臉伏地．那兩個人就對他們說,為甚麼在死人中找活人呢.6他不在這裡,已經復活了．當記念他還在加利利的時候,怎樣告訴你們,7說,人子必須被交在罪人手裡,釘在十字架上,第三日復活.8他們就想起耶穌的話來,9便從墳墓那裡回去,把這一切事告訴十一個使徒和其餘的人.10那告訴使徒的,就是抹大拉的馬利亞,和約亞拿,並雅各的母親馬利亞,還有與他們在一處的婦女.11他們這些話,使徒以為是胡言,就不相信.12彼得起來,跑到墳墓前,低頭往裡看,見細麻布獨在一處,就回去了,心裡希奇所成的事.」

今天是2016年的復活節,對於21世紀的你,「復活」是否一件意想不到的事呢?你有否懷疑過耶穌復活的真確性?在大學禮拜堂,在座的大部份都是大學生,或是校友,或是愛思想的人士.你也肯定有疑惑過!沒有?希望你今天也質疑一下,思想一下:到底耶穌的復活可信嗎?耶穌的復活對今天仍然有意義嗎?

無論是信徒或是非信徒,我們都會問這個問題.對於非信徒來說,耶穌的復活是匪夷所思,難以接受的.但是,對於信徒來說,特別是信了幾十年的信徒,他們大多會深信不疑,並拿出一些證據來:包括空墳墓,門徒見證,天使見證,甚至裹屍布等,以上都是聖經所記載的.當然,空墳墓仍在,不過不只一個.在1980年時,我以神學生的身份到訪當地,那時有2個空墳墓；在1993年,我代一位學者帶隊去,又聽說有3個空墳墓.我們如何確定那個空墳是耶穌的呢?信徒所提的證據,對非信徒來說,仍欠缺說服力,因為只要問一句:「你親眼見過復活的耶穌?聖經這樣說,你就這樣信?」

從現實世界的角度來討論,以上帶出了兩大基督教所關注的問題:到底耶穌的復活可信嗎?耶穌的復活對今日仍然有意義嗎?其實這兩條問題很難找到一個滿意的答案,甚至找不到合理的答案,除非我們進入聖經的世界或信仰的世界,要憑信心接受才可以.

近代聖經學者或神學家,例如Karl Bath,一般都會認為,我們處身於兩個世界.第一個是現實的世界,是客觀存在的世界,而且是人可以接觸,理解的世界(Plato於公元前五世紀引出理念世界).第二個是聖經的世界或信仰的世界,聖經中的人物所記述,所經歷,所信仰的世界.因為你接受基督教信仰,就肯定聖經的信仰權威,你也相信聖經的世界能銜接到今日的世界,並聖經能向你講話.那麼,你也應該相信耶穌已復活,相信耶穌的復活對你仍有意義!所以我們是從聖經的世界或信仰的世界,是以相信聖經世界中的耶穌是復活主,以及祂的復活對我們也具有意義為前提.

對於聖經,初期教會,「耶穌復活,也是一件意想不到的事!」你曾否有「意想不到」的經驗?例如,立法會對2015年政改投票之結果,是令人意想不到；那麼多學生受到生活各種壓力而輕生,都可能是意想不到；最近的李波事件,所謂真相,真令人意想不到,聽得愈多愈樸索迷離,愈令人感到疑惑.

路加福音第二十四章中講述復活節的清晨,婦女準備去膏抹耶穌的身體,經歷到許多的意想不到.古時猶太人的殯葬規矩是用香膏膏死人的身體,並用細麻布加上香料裹上屍體,這是對死者的尊敬.婦女為表達對耶穌的愛和尊敬,她們準備了香料,並第一時間(耶穌周五下午離世至周六是安息日,婦女不能辦事,只有等主日清晨才能)來到準備膏抹耶穌,按常理死人不會走的,但真想不到會見不到耶穌!第一個意想不到是原本封住墳門的大石頭為什麼會滾開了?怎麼可能滾開了?第二個意想不到是她們入內竟然看不見耶穌的身體,可能給人搬走了?第三個意想不到是她們忽然看見兩個人,站在旁邊,衣服放光,明顯是天使.第四個意想不到是天使問她們說:「為何在死人中找活人?」當然,他們不是來墳墓裡找「活人」.她們可能仍不明天使的意思,所以天使再明言說:「他不在這裡,已經復活了!」

這種「意想不到」是越來越強烈的!第一,石頭輥開,可能有人搬開!第二,不見耶穌的身體,可能有人偷屍!第三,看見天使,可能是幻像!第四,天使問:「為何在死人中找活人?」第五,天使告知耶穌已復活了!

「意想不到」的事

「意想不到」的事,是在人未想到而發生的,就如婦女不想耶穌死,不想在墳墓中找不到耶穌.這意想不到的事發生後,當時人的反應各有不同,婦女相信天使所言,相信耶穌復活這「意想不到」的喜事,十分開心,並去把喜訊告訴11位門徒.但大部份門徒未經查證,便認為是胡言(nonsense),否定,不合理,不相信,耶穌對他們而言仍然是死的,因為門徒對「意想不到」的事,盲目否定,所以他們得不到喜樂.另外,彼得卻起來,去探考,去求真,彼得對「意想不到」的喜事,加以調查,具求真精神,結果他發現耶穌身體真的不見了,他感到驚奇.

各位,你面對「意想不到」的消息,你的態度會如何?人對於不同「意想不到」的事會有不同反應,面對「意想不到」的壞事,人可能特別的傷心,難過,失望,或憂慮,也可能令人不敢相信；面對「意想不到」的好事,人可能會特別的興奮,高興,也可能令人不敢相信.

「意想不到」的事,有好事與壞事,善事與惡事,正面與負面的事.現代人也會追求「意想不到」的好事發生,以顯示自己的智慧與能力.當一件又一件「意識不到」的好事發生,例如從飛鳥想像到飛機實現了；從手動的扇想像到電動的風扇發明了；從只聽到聲音的收音機到看到視像的電視機發明了!當人類有越來越多的成功發明,越有自信時,人可能會走向極端,以為自己是上帝,能發明並使以前看來「意想不到的事」得到實現.人想更控制一切,甚至自以為自己可以成為上帝般無所不能,控制一切,如時間,速度,空間,環境,生命,健康,心靈,經濟,社會,甚至他人,現在的一切,未來的一切.現代人也知自己的有限,但他相信自己有無限的潛力.

現代人要避免「意想不到」的壞事發生,因此努力控制一切.人想令一切都在人的操控之中,不會有失誤,不會有失望,這是好的意願,因為這種意願更促進了社會的文明進步.然而,當人過份追求「控制一切」,並且越能成功,越有自信時,人便會去向極端,以為自己是上帝,人就會有意無意地以為自己能操控一切.

不過,人也無法避免「意想不到」的壞事發生的,甚至是避無可避的,這正顯示主耶穌的大能.基督徒由耶穌「從死中復活」,更可以看到上主能轉化「十字架與死亡」這不好的事為美事!跟隨門徒的門徒,大部份認為耶穌是彌賽亞,想不到耶穌會被釘十字架死亡.然而,更出乎意料的事是,羞辱的十架會變成榮耀,及耶穌能從死亡中「復活」.如果耶穌沒有釘的十架痛苦,便得不到榮耀,人也得不到救恩.如果耶穌沒有死亡,失去生命,耶穌便得不到復活和生命.人「意想不到」的壞事,在上帝的手裡能轉化成好事或美事.

信奉上帝的人,常有這種經歷.《路得記》記載拿俄米和丈福音以色列遭饑荒,一家人遷到摩押地,在那裡十年,組織小康之家,落地生根,可能全家以後都定居在摩押地.怎料,拿俄米先死丈夫,隨後兩個兒子也相繼死亡.在男權社會裡,她和兩位媳婦,即時失去了

經濟支柱,無依無靠.這是「意想不到」的不幸,拿俄米就帶著不肯離她而去的媳婦路德回鄉,看能否投靠鄉里.回鄉時丈夫,兒子皆死,她實在是個苦命人,她不再叫拿俄米(甜),改叫瑪拉(苦).這「意想不到」的不幸,卻因上帝奇妙的帶領,讓可憐路得在田間認識拿俄米的親戚波阿斯財主,二人後來結為夫婦,有二項結果,第一是使路得和婆婆拿俄米今後衣食無休；第二是路得婚後不久便生子,這兒子便作拿俄米養子,名俄備得,俄備得是耶西之父,大衛的祖父.人以為意想不到的壞事,在上帝手裡也會變成美事.

還記得1978年的時候,我想考入中大,因為我的成績不錯,特別是中文和歷史科,我自信地認為自己一定有機會,怎料事與願違.當時這事情對我的打擊非常大,我走到教堂鐘樓祈禱,在禱告中,上主給了我一句安慰的話:「他們不收你,是他們的損失,但我要用你!」在鐘樓上,我已立志獻身,將來要修讀神學.其實我並沒有損失,因預科的成績好,我受聘於一間中學教書,並且晚上在香港聖樂學院修讀聖樂,讀到我當時很喜歡的教會聖樂科.於第二年,我報讀了崇基宗教系的神學.在神學領域中,我讀得很開心,並在牧會,教學,研究與寫作上,有更多的發展機會.

然而,我們也得小心,意想不到的惡事也可能會產生,我們應加以警惕!例如近年伊斯蘭國(Islamic State)發動接二連三的恐怖襲擊,日前比利時機場及地鐵遭恐怖襲擊,30多人死,近200人傷,令人心寒.在21世紀,我們發現有「意想不到」的惡存在,這惡一直都存在,只是未見這惡如此恐佈,這惡雖然與現代社會格格不入,卻在不斷的增長中.這惡就是恐佈主義,今天這主義的主要策動者就是伊斯蘭國,他們自覺是受壓迫的一群,想要擊敗敵人,而他們的敵人有兩方,一是美國及西方國家；二是伊斯林什葉派.

在政治上,他們伊斯蘭的遜尼派要建立政教合一的伊斯蘭國,復興昔日伊斯蘭帝國.在宗教上,他們利用群眾宗教的熱情,產生宗教合一的力量,特別是傳統的「聖戰」的觀念.「聖戰」,原非神聖的爭戰,而是自己與自己的爭戰,戰勝自己的私欲與心魔,其後被改成為阿拉而戰,乃為達成阿拉的心願而戰.然而,如果阿拉的心願實在難明,變成政客或野心家的心願,那就糟了!

在手段上,因他們認為打「聖戰」為達成阿拉(某狂魔)的心願,可用一切的手段,包括武力,因此他們以「恐怖襲擊」為手段,即以殺死大量人,包括無辜的人,無武力的非執法人員,殺人者成事後自殺可得到上天堂為酬.他們的目的就是要製造恐慌,引起憂慮,使更多人受其控制,然後逐步建立他們的王國.即使建立不到王國,他們也要削弱敵人的力量,並達到報仇的目的.

恐怖主義者,根據學者的分析,可分為軟件及硬件兩種.軟件上,他們帶著多年積壓的仇恨,掛著聖戰宗教名義與熱誠,存著攻擊與破壞的不當思維,懷著政教合一建立王國的理想,以及群體力量的支持.硬件上,他們有石油,得到波斯灣巨大的捐款,也有大量的先進海陸空,先進武器,眾多的軍隊,包括青少年,兒童,婦女也加入其中,他們亦向全世界伊斯蘭信徒招募,當兵也有收入,也可生活無休.

你我意想不到的惡,從以巴的衝突,千年世仇,現在若然演變基督教國家與伊斯蘭國家的世仇,甚至中東阿拉伯國家與西方及遠東國家的世仇,人類大撕殺,第三次世界大戰出現,人類便會終結,這正中撒旦的詭計.然而,我們也得小心,意想不到的惡事,也可能會產生,加以警惕!可能意想不到的惡會出現,所以香港人要發聲,因為要是這次我們對「李波事件」不出聲,就會否有人破壞《基本法》,甚至產生一地兩檢?真是越想越多可能衍生的問題!

意想不到的上帝

以我多年讀神學與教神學的經驗,我發現我們常常忽略我們所信的上帝是位「意想不到的神」.從創造論看,上帝從無造有,是令人「意想不到」!從基督論看,上帝的道成肉身,基督的神人兩性,是令人「意想不到」!從救贖論看,上帝啟迪保羅因信稱義,基督代贖,救恩論等,是令人「意想不到」!從末世論看,上帝給人永生,世界末日,新天新地是令人「意想不到」!從上帝論看,上帝的意念高過人的意念；上帝是無限的,自由自在的,無所不在的,無所不能的,無所不知的等.如此厲害的上帝,這些事也是令人「意想不到」,因此,上帝可以成就「意想不到」的事,成為意想不到的神!

意想不到的人生觀

從基督的十架與復活顯示出一套意想不到的人生觀.在一週將盡,邁向安息日,基督被釘十架受盡人間的苦難,凌辱,無助,身心靈均受痛楚,痛至死在十架,被埋葬；在主日清晨,一週剛開始,基督復活帶來新的開始,新的一天,新的希望......這顯示了一種從十架的苦,走向復活的人生,這是一種「意想不到」的人生.

從創造論中,上帝從黑暗中造出光明,人生是從黑暗的時刻逐漸步向光明.一天是始自日落,這跟東方的一天始自日出不同,拉比經典中說「與其明亮的開始,黑暗的結束；不如黑暗開始,光明結束.」來得令人快慰,人生也應如此.

從上帝做人的眼睛來看,猶太的拉比在《猶太法典》中自問自答,

問道:「為何上帝造人的眼睛有黑和白兩部份?又為何人只能透過黑的部份才能看東西?」又回答說:「因為人生要經歷(黑)黑暗和(白)光明,而且人生必須經過黑暗,才能看到光明.」

猶太人像中國人一樣經歷許多苦難,在苦難中得著一個重要的意識,就是「苦難意識」.這意識在日常的生活,節期,婚禮中,不斷提醒人們,人生是苦的,並需常常存有「苦難」意識.每個周五的黑夜是安息日的開始,全家人一起祈禱,一起迎接黑夜來臨,相信經過黑夜有光明和平安,相信經過苦難有甘甜,使他們勇敢面對未來.另一個例子是每年新年,過逾越節,他們會穿上華麗衣服,並預備豐盛的食物,當中一定要吃粗麵包和苦菜,因為它們象徵屈辱,苦難和失敗,在這節期,吃粗麵包和苦菜就是要提醒他們曾在埃及為奴,做苦工,受盡苦楚和凌辱!在新婚喜慶的日子裡,猶太婚禮儀節中,規定新人不要把酒飲完,然後放入盤子裡,而要新人喝完酒後,把酒杯摔破,這行為表示兩人願意同甘共苦同過艱難的一生,並提醒婚後不會風平浪靜,雙方總會有衝突,仍需共同渡過,卻不要如公主,王子結合永遠幸福快樂的神話.

猶太人認為猶太歷史是苦難史,人生也是苦難人生,生活的磨練,身體的疾病,,天災人禍,流離失所等,但他們非悲觀主義者,因著苦難意識,他們做人更積極進取,不輕言放棄.在迦太基著名的博物館有幅名畫「將軍」,象徵苦難的魔鬼與人下棋.在棋局中,人累積過去一切的智慧,經驗,戰略,大戰魔鬼；畫中魔鬼快把人將軍,勝過人,人快敵不過,但有位棋高手看到這畫,大聲叫:「還有一著,人是不會輸!」

自稱為上帝子民的猶太人自出生,開始準備與苦難爭戰.自稱為基督門徒的基督徒,你相信耶穌,就準備要跟從耶穌背起自已的十架,與苦難和罪惡爭戰.從耶穌的被釘十架,死亡於禮拜五下午三時,是安息日的開始,大地進入黑暗；耶穌的復活則在安息日過後,主日的黎明,天色愈來愈光亮,耶穌復活了!信耶穌復活的人的生命,也是由黑暗步向光明.

復活事件的轉化

復活,這意想不到的事件,會引致生命的轉化(Transformation),第一種轉化是從死裡復活,勝過死亡,使人由懼怕到不怕死亡.第二種轉化從無可能(死),轉化為可能(復活).第三種轉化是從每天開始黑暗到光明,勝過黑暗,使人充滿希望；在看來越無希望的地方,出現了希望；在墳墓中可找到活人,那從死裡復活的人.第四種轉化是婦女聽完天使的報訊後,從失憶到憶起耶穌的話:「第三天要復活」,婦女的心靈得到轉化,由無知到開竅(心靈的甦醒),婦女想起主的話,從失憶中甦醒起來,回憶能使故事,人物活現；在現實中與過去相遇,回憶所愛的人,也會引發人的感情.上主要以色列民常有記憶,記住上帝的作為:逾越節,受苦節,復活節.記憶很重要,因為忘記密碼,就進不到屋,登入不到電郵,開不到夾萬(古老夾萬,需要一組組字輸入).你記起密碼,就會很快樂；記不起,無記性,那就會很麻煩了!第五種轉化是婦女從接受信息者(receiver)轉化成傳遞信息的人(messenger).第六種轉化是婦女見不到耶穌的身體,她們的情緒從失望,恐懼,驚奇,轉化到充滿平安,喜樂,和動力去傳福音.

從十架的痛苦轉化成復活的喜樂,使人由軟弱變成剛強,猶太人在「苦難意識」加上對「彌賽亞」 (基督徒:耶穌基督)的盼望,產生一種轉化力量,面對人生的逆境!

結語

這股生命的力量是今日的學生需要的,也是中年,年長的人所需要的.今日我強調從聖經的世界/信仰的世界去看「復活主」,你會看到一位意想不到的復活主,在人意想不到中顯出祂的大能,祂的大能能轉化生命,從死亡轉化為再生,復活；從苦難轉化為平安,力量；從懦弱轉化為剛強；從黑暗轉化為光明；從惡事轉化為美事.

祈禱

復活主,願我們都相信你是從死人中復活,勝過黑暗,使人意想不到的主,並能轉化生命,賜人力量,叫我們得力走崎嶇的人生路.阿們.



\newpage

\section{從死裡首先復活的耶穌}
\label{sec:316}
\textbf{鄧瑞強博士 - 從死裡首先復活的耶穌}
\newline
\newline
\hyperref[sec:315]{< < < PREV SERMON < < <}
~
\hyperref[sec:toc]{[返目錄]}
~
\hyperref[sec:317]{> > > NEXT SERMON > > >}


啟示錄 1:4-8
\newline
\begin{longtable}{cl}
\hline
\hline
章節 & 經文 (和合本修訂版)\\
\hline
1:4 & \begin{tabularx}{0.7\textwidth}{X} 約翰寫信給亞細亞的七個教會。願那位今在、昔在、以後永在的神，與他寶座前的七靈， \end{tabularx} \\ \\ \relax
1:5 & \begin{tabularx}{0.7\textwidth}{X} 和那忠信的見證者、從死人中復活的首生者、世上君王的元首耶穌基督，賜恩惠和平安給你們。他愛我們，用自己的血使我們從罪中得釋放 ， \end{tabularx} \\ \\ \relax
1:6 & \begin{tabularx}{0.7\textwidth}{X} 又使我們成為國度，作他父神的祭司。願榮耀、權能歸給他，直到永永遠遠。阿們！ \end{tabularx} \\ \\ \relax
1:7 & \begin{tabularx}{0.7\textwidth}{X} 「看哪，他駕雲降臨；眾目都要看見他，連刺他的人也要看見他；地上的萬族要因他哀哭。」這是真實的。阿們！ \end{tabularx} \\ \\ \relax
1:8 & \begin{tabularx}{0.7\textwidth}{X} 主神說：「我是阿拉法，我是俄梅戛 ，是今在、昔在、以後永在的全能者。」 \end{tabularx} \\ \\
[1ex]
\hline
\hline
\end{longtable}
講題:「從死裡首先復活的耶穌」　　Jesus, the Firstborn from the Dead

經文:啟示錄1:4-8

講員:鄧瑞強博士

場合:香港中文大學崇基學院禮拜堂

日期:2016年4月3日

各位弟兄姊妹,早安.

今日的講道,會講啟示錄的經文.

一般相信,啟示錄成書的日子,約為主後90年至95年左右,這是羅馬皇帝圖密善(Domitian,或譯作「多米田」)統治的後期.很多羅馬皇帝都會被人崇拜,凱撒大帝死後,被人尊稱為神.他的養子及繼承者凱撒奧古斯都順理成章稱自己為「神的兒子」,又認為自己是宙斯的化身.後來的尼祿,認為自己是亞波羅的化身.圖密善更進一步,在自己生前,自稱是「主及神」(lord and god),要他的臣民下拜.一系列的造神運動.成為神的好處,是以後不再需要和人類講道理.一代又一代的造神運動,一代又一代的不講道理.

生在神化了的皇帝的年代,你必須要學曉低頭下跪.頭低下來,你便看不清真相.又或者說,你要看不清真相,才能低下頭來.低下頭,閉上咀,生命留了下來,靈魂卻失去了.

那個時代,是要你向一個人下拜的時代.那個時代,是阿諛奉承的人大展所長的時代.那個時代,是考驗人的正直的時代.

就在那個時代,神的僕人約翰,領受神的啟示,向神的子民發出鼓勵的信息.他提示我們,毋忘天上的神.誰遺忘天上的神,誰就在混亂的人間失去方寸.

我們看看神的僕人約翰的說話:啟示錄1:4－8

啟1:4　約翰寫信給亞西亞的七個教會.但願從那昔在,今在,以後永在的神,和他寶座前的七靈,

啟1:5　並那誠實作見證的,從死裡首先復活,為世上君王元首的耶穌基督,有恩惠,平安歸與你們!他愛我們,用自己的血使我們脫離罪惡,

啟1:6　又使我們成為國民,作他父神的祭司.但願榮耀,權能歸給他,直到永永遠遠.阿們!

啟1:7　看哪,他駕雲降臨!眾目要看見他,連刺他的人也要看見他；地上的萬族都要因他哀哭.這話是真實的.阿們!

啟1:8　主神說:「我是阿拉法,我是俄梅戛,是昔在,今在,以後永在的全能者.」

我將這經文分成三個段落.

第一段落是第4節至第5節上半節.這段落告訴我們,真神是誰.

啟1:4　約翰寫信給亞西亞的七個教會.但願從那昔在,今在,以後永在的神,和他寶座前的七靈,

啟1:5上　並那誠實作見證的,從死裡首先復活,為世上君王元首的耶穌基督,有恩惠,平安歸與你們!

「約翰寫信給亞西亞的七個教會」,這七教會都在現今的土耳其.「七」這個數字,代表完全,完整,故「七教會」也可理解為地上所有教會.啟示錄是約翰寫給七教會的信.作為一封信,信的內容一定與收信人的生命息息相關,也是收信人可以明白的.故此,啟示錄講的東西,無論在我們看來多古怪,也一定是收信人明白的,也對他們有莫大的鼓舞作用.

造神運動在地上越是猖狂,「有一位神在天上」這信息越是重要.若沒有一位神在天上,就很難想像人能不向地上的神低頭.

這封信的開首,便提及天上的真正的神,這神有一種三一性.

我們看看與聖父有關的部分.經文說:「但願從那昔在,今在,以後永在的神」.這是與聖父有關的經文.這經文譯得不好,《新漢語譯本》的翻譯是「願今在,昔在,來臨中的那一位」,英文翻譯一般是 “from Him who is and who was and who is to come” (NIV, NASB, RSV).經文最先強調的,是聖父的「今在」,然後才是「昔在」.再者,聖父不是靜態的「以後永在」,而是動態的「在來臨中」.

當地上變得黑暗,人心變得虛怯的時候,人會懷疑,神還在嗎?約翰說,神在,神此刻在,神今在.神的今在,猶如祂的昔在.昔日,祂為受欺壓的人伸冤；今日,祂同樣如是.昔日,祂懲治法老王；今日,祂同樣如是.

聖父不是靜態的,在天上坐著搖搖椅,而是動態的,「在來臨中」.聖父有情,當人在痛苦呼喊時,不會坐著不動.聖父看見,聖父聽見,聖父在來臨中.

關於三一神的聖靈的部分,經文說:「和祂寶座前的七靈」.「七靈」這講法的典故,來自舊約撒迦利亞書,詳細的考據不能在這裡交代了.簡言之,「七靈」代表神的大能.要面對世界的艱難,必須要倚靠神的靈.當然,「七」代表完全,完整,「七靈」或者也可代表神那完備的靈.

關於聖子的部分,經文講得很詳細.啟1:5說:「並那誠實作見證的,從死裡首先復活,為世上君王元首的耶穌基督」.這裡用三重描述來談論聖子耶穌.這三重描述,描繪耶穌現在的身分.

耶穌是「那誠實作見證的」,他是誠實的證人,為真理作證.當地上的君王神化自己時,地上一定充滿謊話.因為唯有謊話,才會指一個人是神.耶穌是真理的見證人,啟示錄也鼓勵讀這信的人,要堅定自己,作一個誠實的真理見證人.

接著,經文說:耶穌是「從死裡首先復活」的.在虛假的時代,堅持真理,只有死路一條.或許,這條死路離我們也不太遠.耶穌走上了這條死路,之後,他在真神的大能下復活了.這就顯出生命的微妙局面.那些為了生路而繼續謊話的人,至終會滅亡.那些為了真理而走入死路的人,他們至終會享受永恆的生命.

這正是耶穌說過的:「凡要救自己生命的,必喪掉生命；凡為我喪掉生命的,必得著生命.」(太16:25)

關於聖子,經文最後說:他是「世上君王的元首」.世人只看到地上的君王,服膺他們的權力；但基督徒應看到,在這一切地上君王之上的真正君王,及祂的真正權力.地上的君王,有生殺大權；但天上的基督君王,卻有賜予永生的大能.或許,我們真的要想想,死亡的威脅大,抑或永生的吸引力大?

約翰祝福收信者,從三一真神而來的,是「恩惠,平安」.恩惠,即禮物,是白白賜予的,這不同於在威脅下乞求得來的苟廷殘喘.平安,應是當時四處可見的政治標語,就是那在強大的軍事力量保證下的羅馬和平(Pax Romana).三一神賜的和平,不是那種和平,而是一種真理下的和平,一種為正義受苦的和平,一種展示永恆生命的和平.在塵世,我們期望哪種和平?

我們去到經文的第二個段落.

第二段落是第5節下半節至第6節.這段落告訴我們,我們是誰.

啟1:5下　他愛我們,用自己的血使我們脫離罪惡,

啟1:6　又使我們成為國民,作他父神的祭司.但願榮耀,權能歸給他,直到永永遠遠.阿們!

這裡用三重描述去展示基督在信徒生命裡的作為.這三重描述,描繪基督過去為我們做的事.基督的作為,改變了我們的生命面貌.

如何改變呢?

經文開宗明義說,「他愛我們」,這就是基督在我們身上恆常作的事.基督永遠愛我們,我們生命永遠有所屬.一個人有所屬,就不用四處去找點什麼來讓自己安心.我們需要有人愛我們,是這份愛,讓我們在這個混亂而虛偽的世界,仍能堅持下去.基督愛我們,這份愛多麼重要.

基督愛我們,愛到什麼地步?經文說:他「用自己的血使我們脫離罪惡」.真正的愛,是叫人受傷的.要救人脫離罪惡,是要犧牲的.救贖,要流的血,是自己的血,不是流他人的血.唯有流自己的血,才能流露愛.唯有愛,才有足夠大的力量,救人離開深淵.因著基督,信徒能經驗新生.這種新生命,是知道誰是真正的神,知道生命應向誰交賬.

基督愛我們,將我們的生命帶回正途,便改變了我們的終極身分.經文說,他「使我們成為國民,作他父神的祭司」.作為天父的祭司,生命要自重,在人生中,要放棄一切偶像,唯獨敬拜一位真神.地上神化了自己的皇帝,會要求你當他是神一樣下拜.強大的武力,會逼你低頭.人性的罪,會誘惑你出賣真理.這個時候,請謹記你作為天父祭司的身分.這是你的終極身分.要在黑暗的世界,指出光明.

我們去到經文的第三個段落.

第三段落是第7節.這段落告訴我們,我們面對一個怎麼樣的明天.

啟1:7　看哪,他駕雲降臨!眾目要看見他,連刺他的人也要看見他；地上的萬族都要因他哀哭.這話是真實的.阿們!

也是三重描述,描繪基督將要成就的事.

我們講過,神是「今在,昔在,來臨中的那一位」,這「今在,昔在,來臨中」是透過基督的「今在,昔在,來臨中」來表現.我們已講過基督的今在和昔在,現在重點是基督的「來臨中」.

首先,基督將「駕雲降臨」,「雲彩」代表神聖,基督將以神的身分來臨.我們不是面對虛無的明天,明天是有內容的,明天被基督的再來所充滿.

「眾目要看見他,連刺他的人也要看見他」,再一次提到基督的受苦,但如今,受苦的基督要展示公開的勝利.世上一切偶像與邪惡,會見到基督的得勝.

關於將來,經文最後說,「地上的萬族都要因他哀哭」,這是因為萬民要面對基督的審判,要為自己的罪惡而懺悔.我們可大膽地預期這明天,真神在上,正義總會彰顯,真理終會獲勝,惡人必被懲治.

以上三大段經文的結尾是(第8節):「主神說:我是阿拉法,我是俄梅戛,是今在,昔在,來臨中的全能者.」

若你看東正教的聖像畫,很多時便是按這經文的意思入畫.畫中央是基督,上方左右寫上“Α”(alpha,阿拉法)及“Ω”(omega,俄梅戛),在頭上的光環裡寫上「今在」,然後,在不同地方寫著“pantokrator”(全能者).

若當時的皇帝圖密善(Domitian)自稱是「主及神」(lord and god),則啟示錄在這裡要一正視聽,誰是主,誰是神,一目了然,就是那位alpha及omega,那位創始成終者.祂是萬有之源,也是萬有的成全者.祂是alpha,是萬有存在的根基.祂是omega,對萬有歷史有最後話事權.作為萬有存在的根基,祂能叫死人復活.作為歷史的最後話事人,祂對萬有作出最後審判.在神那裡,萬物的存在,應有其意義.正義與邪惡,最後也總有所分判.然後經文說,祂是「全能者」.這講法有點似舊約的「萬軍之耶和華」,強調祂的大能大力.在光明與黑暗的戰鬥中,祂穩操勝券.

啟示錄用短短的經文,畫出一巨型圖畫.古往今來,邪惡無盡,黑暗無邊,痛苦莫明.然而,在這一切之上,有真理正義,有光明生命,有恩惠平安.這圖畫是否足以安慰你,提示你,鼓舞你呢?

但願:榮耀歸於聖父,聖子,聖靈.阿們.



\newpage

\section{民以食為先}
\label{sec:317}
\textbf{許開明牧師 - 民以食為先}
\newline
\newline
\hyperref[sec:316]{< < < PREV SERMON < < <}
~
\hyperref[sec:toc]{[返目錄]}
~
\hyperref[sec:318]{> > > NEXT SERMON > > >}


約翰福音 21:1-19
\newline
\begin{longtable}{cl}
\hline
\hline
章節 & 經文 (和合本修訂版)\\
\hline
21:1 & \begin{tabularx}{0.7\textwidth}{X} 這些事以後，耶穌在提比哩亞海邊又向門徒顯現。他怎樣顯現記在下面。 \end{tabularx} \\ \\ \relax
21:2 & \begin{tabularx}{0.7\textwidth}{X} 西門‧彼得、叫低土馬的多馬、加利利的迦拿人拿但業、西庇太的兩個兒子，和另外兩個門徒，都在一起。 \end{tabularx} \\ \\ \relax
21:3 & \begin{tabularx}{0.7\textwidth}{X} 西門‧彼得對他們說：「我打魚去。」他們對他說：「我們也和你一起去。」他們就出去，上了船；那一夜並沒有打著甚麼。 \end{tabularx} \\ \\ \relax
21:4 & \begin{tabularx}{0.7\textwidth}{X} 天剛亮的時候，耶穌站在岸上，門徒卻不知道他是耶穌。 \end{tabularx} \\ \\ \relax
21:5 & \begin{tabularx}{0.7\textwidth}{X} 耶穌就對他們說：「孩子們！你們有吃的沒有？」他們回答他：「沒有。」 \end{tabularx} \\ \\ \relax
21:6 & \begin{tabularx}{0.7\textwidth}{X} 耶穌對他們說：「你們把網撒在船的右邊，就會得到。」於是他們撒下網去，竟拉不上來了，因為魚很多。 \end{tabularx} \\ \\ \relax
21:7 & \begin{tabularx}{0.7\textwidth}{X} 耶穌所愛的那門徒對彼得說：「是主！」那時西門‧彼得赤著身子，一聽見是主，就束上外衣，跳進海裡。 \end{tabularx} \\ \\ \relax
21:8 & \begin{tabularx}{0.7\textwidth}{X} 其餘的門徒因離岸不遠，約有二百肘，就坐著小船把那網魚拉過來。 \end{tabularx} \\ \\ \relax
21:9 & \begin{tabularx}{0.7\textwidth}{X} 他們上了岸，看見那裡有炭火，上面有魚和餅。 \end{tabularx} \\ \\ \relax
21:10 & \begin{tabularx}{0.7\textwidth}{X} 耶穌對他們說：「把剛才打的魚拿幾條來。」 \end{tabularx} \\ \\ \relax
21:11 & \begin{tabularx}{0.7\textwidth}{X} 西門‧彼得就上船，把網拉到岸上，網裡滿了大魚，共一百五十三條；雖然魚這樣多，網卻沒有破。 \end{tabularx} \\ \\ \relax
21:12 & \begin{tabularx}{0.7\textwidth}{X} 耶穌對他們說：「你們來吃早飯。」門徒中沒有一個敢問他：「你是誰？」因為他們知道他是主。 \end{tabularx} \\ \\ \relax
21:13 & \begin{tabularx}{0.7\textwidth}{X} 耶穌走過來，拿餅給他們，也照樣拿魚給他們。 \end{tabularx} \\ \\ \relax
21:14 & \begin{tabularx}{0.7\textwidth}{X} 耶穌從死人中復活後向門徒顯現，這是第三次。 \end{tabularx} \\ \\ \relax
21:15 & \begin{tabularx}{0.7\textwidth}{X} 他們吃完了早飯，耶穌對西門‧彼得說：「約翰的兒子西門，你愛我比這些更深嗎？」彼得對他說：「主啊，是的，你知道我愛你。」耶穌說：「你餵養我的小羊。」 \end{tabularx} \\ \\ \relax
21:16 & \begin{tabularx}{0.7\textwidth}{X} 耶穌第二次又對他說：「約翰的兒子西門，你愛我嗎？」彼得對他說：「主啊，是的，你知道我愛你。」耶穌說：「你牧養我的羊。」 \end{tabularx} \\ \\ \relax
21:17 & \begin{tabularx}{0.7\textwidth}{X} 耶穌第三次對他說：「約翰的兒子西門，你愛我嗎？」彼得因為耶穌第三次對他說「你愛我嗎」，就憂愁，對耶穌說：「主啊，你無所不知，你知道我愛你。」耶穌說：「你餵養我的羊。 \end{tabularx} \\ \\ \relax
21:18 & \begin{tabularx}{0.7\textwidth}{X} 我實實在在地告訴你，你年輕的時候，自己束上帶子，隨意往來；但年老的時候，你要伸出手來，別人要把你束上，帶你到不願意去的地方。」 \end{tabularx} \\ \\ \relax
21:19 & \begin{tabularx}{0.7\textwidth}{X} 耶穌說這話，是指彼得會怎樣死來榮耀神。說了這話，耶穌對他說：「你跟從我吧！」 \end{tabularx} \\ \\
[1ex]
\hline
\hline
\end{longtable}
講題:《民以食為先》For people to eat first

經文:約翰福音21章1-19節

講員:許開明牧師

場合:香港中文大學崇基學院禮拜堂

日期:2016年4月10日

1這些事以後,耶穌在提比哩亞海邊又向門徒顯現.他怎樣顯現記在下面:2有西門．彼得和稱為低土馬的多馬,並加利利的迦拿人拿但業,還有西庇太的兩個兒子,又有兩個門徒,都在一處.3西門．彼得對他們說:「我打魚去.」他們說:「我們也和你同去.」他們就出去,上了船；那一夜並沒有打著甚麼.4天將亮的時候,耶穌站在岸上,門徒卻不知道是耶穌.5耶穌就對他們說:「小子!你們有吃的沒有?」他們回答說:「沒有.」6耶穌說:「你們把網撒在船的右邊,就必得著.」他們便撒下網去,竟拉不上來了,因為魚甚多.7耶穌所愛的門徒對彼得說:「是主!」那時西門．彼得赤著身子,一聽見是主,就束上一件外衣,跳在海裡.8其餘的門徒離岸不遠,約有二百肘,就在小船上把那網魚拉過來.9他們上了岸,就看見那裡有炭火,上面有魚,又有餅.10耶穌對他們說:「把剛才打的魚拿幾條來.」11西門．彼得就去,把網拉到岸上.那網滿了大魚,共一百五十三條；魚雖這麼多,網卻沒有破.12耶穌說:「你們來吃早飯.」門徒中沒有一個敢問他:「你是誰?」因為知道是主.13耶穌就來拿餅和魚給他們.14耶穌從死裡復活以後,向門徒顯現,這是第三次.15他們吃完了早飯,耶穌對西門．彼得說:「約翰的兒子西門,你愛我比這些更深嗎?」彼得說:「主啊,是的,你知道的愛你.」耶穌對他說:「你餵養我的小羊.」16耶穌第二次又對他說:「約翰的兒子西門,你愛我嗎?」彼得說:「主啊,是的,你知道我愛你.」耶穌說:「你牧養我的羊.」17第三次對他說:「約翰的兒子西門,你愛我嗎?」彼得因為耶穌第三次對他說「你愛我嗎」,就憂愁,對耶穌說:「主啊,你是無所不知的；你知道我愛你.」耶穌說:「你餵養我的羊.」18我實實在在地告訴你,你年少的時候,自己束上帶子,隨意往來；但年老的時候,你要伸手出來,別人要把你束上,帶你到不願意去的地方.」19耶穌說這話是指著彼得要怎樣死,榮耀神.說了這話,就對他說:「你跟從我吧!」

引言

中國人有句話:民以食為先(天),你同意嗎?每天你也要吃飯,吃對人生意義重大,但意義何在?你會怎樣回答?食,是以食為先,或以食為天?聖經如何看「食」這個問題?本主日的經課是約翰福音21章1至19節,可能有助你思想這問題.

自古到今,「食」是人類非常關注的課題.「食」有很多原因或目的,可以是為存在,工作,享受,健康,父母要食,嘴饞等.我們需要定時食,八時吃早餐；十二時吃午餐；七時吃晚餐.飲食,是人的本性,《禮記》寫道:「飲食男女,人之大慾存焉」.在這物質富裕的時代,人類在生活上的進步就是從有用,到合用,然後到好用；同樣地,人類在飲食上的進步是從有食物可食,到合符口味,再到吃美味好吃的食物,要求色,香,味全,甚至要吃得健康,吃得環保,吃得道德,吃得創新等.

飲飲食食

世上好些主要宗教,包括佛教,伊斯蘭教,猶太教等,都對飲食有嚴格規範.然而,大部份現代基督新教信徒,都不認為飲食的選擇跟信仰原則有什麼關係,傾向相信凡物皆可吃,因沒有律法禁止；又謂上帝的國不在乎吃喝,所以吃什麼穿什麼,跟信仰完全無關.然而事實是否只是如此?

從基督信仰角度,我們應重新思索飲食在信仰中的位置,以及信仰在飲食中的位置.當上帝創造人後,不久,人就面對「食」的誘惑,並在「食」的問題上跌倒,他們因食了上帝禁止的食物而受罰,那日子必死,失去了永生,但卻得到智慧和發現了自由的意志.到底「食」是帶來祝福或咒詛呢?食可使人得飽足,食得好使人感覺幸福；但食得太多,太肥,太甜,太鹹,太放縱都不行,食錯了,更足以致命.

「食」是人的必須品,是人人賴以為生的東西.「食」不止是「食到」,「食飽」,或「食得好不好」的問題.「食」與人生的各方面都息息相關,是看似簡單,卻又複雜的問題.「食」不止食這麼簡單,藉著食,還可以建立關係,聯絡感情,談情說愛,談生意,算帳,拜託別人幫忙等.

食與人生,社會民生皆有關連

社會大眾最關切的不是「民主」而是「民以食為天」,因為彼此問候時會問:「吃過了沒有?」更有許多與「吃」有關的詞語用法,例:吃苦,吃虧,吃教,吃人,吃不消,吃得開,吃裡扒外,食古不化等.維生之計是「搵食」,「為個口」,「糊口」,而大同理想是「民生」,「皆有所養」,但食是民生的首要問題等.

猶太人食的智慧

第一是飲食有度,如廁有度.飲食有度,即只「吃滿三分一胃的容量,喝三分一,留下三分一空」；而如廁有度,即不站著吃,吃不言,枕不語.

第二是飲食合宜,每天只吃早晚兩餐,晚上才飲酒(早上飲是石頭,午間飲是紅銅,晚上飲是白銀；三天一喝是黃金),鼓勵素食,特別是蔬果.

第三是飲食小心,要分潔淨與不潔淨的食物,凡不潔淨的食物不接觸,偶蹄動物即羊,牛,鹿,是潔淨的,所以可吃；而奇蹄動物即駱駝,兔子,豬,是不潔淨的,所以不可吃；水中動物而有鱗有鰭的是潔淨的,所以可吃；而無鱗無鰭的,例如鰻魚,鱔,貝類等是不潔淨,不可吃.另外也不吃對身體有害的食物,如果吃了,即等於輕視自己,會得不到祝福.還需要分有益與無益的食物,例如魚(小魚),蛋,蜂蜜,餅,麵包穀物等,小魚能使人體生殖力強勁,而這早餐能為大腦提供營養.耶穌為門徒預備的是營養早餐,比麥當奴和大家樂的早餐還要好.猶太人甘願受約朿卻能得健康的身體,這是有價值的.

小結

從一個人如何吃,能否食之有道,能否自控,我們可以看出他是個怎樣的人.今早,我們要透過這段經文,從復活主給門徒準備的精美早餐來思想.

食的問題與始祖有關

上帝創造人後,把他安置在伊甸園,並賜人各類果子作食物,只是分辨善惡樹上的果子不可吃,吃了的話那日子必死.(創2:15-17)可惜,當女人聽見蛇說:「你們不一定死,反而吃了會使眼睛明亮,便能知道善惡.」於是女人見那棵樹好作食物,也悅人眼目,且是可喜愛的,能使人有智慧,就摘下果子來吃了,又給他丈夫,她丈夫也吃了,結果人類要受罰.可見人受罰也與食有關,於是在尋找食物上,成為一種艱辛,因為人類要汗流滿面才得糊口,如果始祖早知這點,可能不會如此衝動食禁果.

無論如何,人類始祖已開始面對食的誘惑.食禁果在理性上,給人聰明智慧；食禁果在生理上,滿足人的慾望；但食禁果在靈性上,就是不聽從上帝的話,將失去永生.當中已引出一個問題:你的抉擇會是什麼?你會選擇食禁果,滿足即時的食慾,短暫的快樂,慢慢失去生命嗎?還是只聽上帝話,埋沒自由意志,不食禁果,得永遠居住在伊甸園和永遠不死?食與人生都充滿誘惑和疑惑,因為食為先,食就好像人生中最首要的.不相信嗎?試想想,要是吃不夠,不可吃,就會飢餓受苦,會病,會死.說到這裡,你可能又會說:食不可能是人生最重要的,有人為得自由,得公平,得生命更高的理想,就會絕食；年青男女為身型苗條會節食；敬虔信徒會禁食.

【耶穌出道前面對食的試探】

魔鬼叫饑餓了四十天的耶穌把石頭變作食物,當時耶穌餓了四十天,也就是人性最軟弱的時候,祂面對人生中最大的敵人,也最需要吃飯來得到力量,才可面對強敵.如果耶穌打輸了這一仗,人類便得不到翻身,得不到救恩了.耶穌毅然拒絕了魔鬼,並引出一句震驚古今的名言:「人活著不是單靠食物,乃是靠上帝的話」,這裡引出上帝的話或道也是人生命所需.在耶穌來說,遵守上帝的話比滿足「食」的慾望重要.很明顯地,對耶穌而言食不是為先,祂在亞當跌倒的地方站起來.

當中有兩個重要的思想.第一,平衡理解食物的涵義.當耶穌提到「人活著不是單靠食物,乃是靠神口裡所出的一切話.」(路4:4)正指出身體的食物與靈性的食物是缺一不可的,當要耶穌在兩者作抉擇時,耶穌會選擇「靈性的食物」.

第二,生命的養分與目標不僅是食物,所以,耶穌說:「我的食物就是遵行差我來者的旨意,做成他的工.」(約4:34)鼓勵我們透過順服及事奉,得著生命養分,這個目標遠高於食物本身.

【耶穌最偉大的,最大型的神蹟:以五餅二魚餵飽五千人】

耶穌看見聽道的群眾得到上帝的道,卻在飢餓中,他就施行餵飽五千人的神蹟,使人經歷上主的大能與大愛.得道與得食二者皆成,這是天國的臨到,人類嘗到在地若天的滋味.所以,如果天堂沒有可吃的,對許多嘴饞的人是一種損失.

耶穌復活後,與門徒相聚,享用鮮魚早餐

「有(1)西門彼得,和稱為(2)低土馬的多馬,並加利利的迦拿人(3)拿但業,還有西庇太的兩個兒子(4)(5),又有兩個門徒(6)(7),都在一處.」(約21:2)彼得連同六位門徒重操故業「打魚去」,自己的專業可使自己得食,也使人得食,更可以賺錢養活,在當代人來說,確是明智之舉.一起尋找食物,成捕魚團隊,互相幫助,彼此扶持,更明智.可惜,他們一起的時候仍找不到食物,因為「那一夜並沒有打著甚麼」.古時漁夫喜歡在晚上打魚(參路5:5),據說夜晚正是打魚最好的時候,彼得是個老練的漁夫,他知道在何時,何處,何種方法可以打到魚；但他們竟然整夜打不著魚,多麼令人感到失望和挫敗；對彼得這老漁夫來說,更是深受打擊.

在挫敗中,有人來問:「『小子,你們有魚吃的沒有?』他們回答說:『沒有.』」(約21:5)各位,當一個人考試成績不好時,最怕人問:「你的考試成績如何?」同樣地,漁夫捕不到魚,最怕就是有人問:「你捕到魚嗎?」或「你有魚可賣嗎?」他們回答說:「沒有.」這是一種對自己的失望,也會使別人失望.彼得和門徒跟了耶穌數年,彼得甚至把船和漁具全部拍賣來跟從主,結果最後夫子被釘十字架死了,真令人失望.停了數年捕魚生涯,為了生計,他首次重操故業,以為可以東山再起,再展雄風,怎料,整夜勞碌都只是「得個吉」,對彼得來說,打擊是何等沈重.

人的不能,正顯出復活主的大能

「耶穌說:『你們把網撒在船的右邊,就必得著.』他們便撒下網去,竟拉不上來了,因為魚甚多.」(約21:6)當時門徒並未知那在岸上的是耶穌.那麼為何門徒會聽話?可能是這位陌生人是用了肯定的,命令的方式說:「你們把網撒在船的右邊,就必得著.」而不是:「你們試試把網撒在船的右邊,就必得著.」或「可能把網撒在船的右邊,便會得著.」現代人一般都不敢說,肯定的,命令式的話.結果,「他們便撒下網去,竟拉不上來了,因為魚甚多.」(約21:6)從無魚到多魚,耶穌問:你們有魚吃沒有?「沒有」應是令人失望的,但復活主能扭轉乾坤,「把網撒在船的右邊,必得著」,結果網拉不上來,因為魚甚多.從無到有,從有到豐盛；從捕不到魚,到捕大魚,捕大量的大魚；從靠自己的經驗與能力而失敗；到信靠復活主的指示而成功,這些都可看見復活主的能力.

復活主的能力

第一是轉化能力(Transformation),祂能把無轉化成有；把黑暗轉化成光明；把死亡轉化成復活；把不好轉化成好；把醜陋轉化成美麗等.

第二是優化能力(Optimization),祂能使擁有優化到擁有很多；使光明優化到非常光明；使再生優化到永生；使好優化到極好；使美優化到極美.

第三是從轉化能力到優化能力(From Transformation to Optimization),從無到有,從擁有到擁有很多；從黑暗到光明,從光明到非常光明；從死亡到再生,從再生到永生；從不好到好,從好到到極好；從醜到美,從美到極美.這種復活主轉化與優化的能力,使人生愈來愈好.

【認出復活主】

「耶穌所愛的那門徒對彼得說:『是主.』那時西門彼得赤著身子,一聽見是主,就束上一件外衣,跳在海裡.」(約21:7)這是意想不到的收獲!使那不知名的門徒,就是耶穌所愛的門徒說:「這是主!」這話使彼得醒悟,以往的彼得是最聰明的,最先認得主的,為何他會遲鈍起來?可能是彼得曾三次不認主,這屬靈上的失敗,使他振作不起；更可能是彼得提出重操故業,結果整晚無所獲,更振作不起；更大可能是,彼得這位大漁夫整晚打不到魚,但這個陌生人竟能如此輕易指出右邊有魚,且是非常多的魚,使彼得在想這個人是誰?

他不會想起是主,因為理性上,他不能接受已死的耶穌能復活.

「這是主」是什麼意思呢?主是神的意思,他曾經稱呼耶穌為主.那門徒的「這是主」喚醒了他的信心的眼睛,「這是主」才能行出這種神蹟；「這是復活的主!」才能在現在行出這種神蹟；「主,沒有死,主從死裡復活!且在眼前!」

「耶穌所愛的那門徒對彼得說:『是主.』那時西門彼得赤著身子,一聽見是主,就束上一件外衣,跳在海裡.」(約21:7)彼得開心得跳到水裡去,束上上衣即是他放棄眼前的魚獲,要上岸,想要看清楚復活的主；跳下海,即他不顧這樣會把上衣弄濕,因為見復活主比魚獲更為重要.「其餘的門徒,(離岸不遠,約有二百肘〔古時以肘為尺,一肘約有今時尺半〕),就在小船上把那網魚拉過來.」(約21:8)「約有二百肘」,「肘」為古時長度單位,每肘約為十八英吋,故當時離岸約有一百呎.但其餘的六位門徒卻沒有甚麼反應,與彼得的興奮截然不同,可能除了那門徒外,其餘他們不確定祂是復活主.

復活主為門徒預備早餐

彼得上岸不久,他們坐小船上了岸,因只離約一百呎,不算太遠,當門徒上岸後,發現那裡有炭火,上面有魚,又有餅.「上面有魚」即是說他們根本無須到海中去打魚,因為岸上已經有魚了.「上面有魚,又有餅」,「魚」代表海中的豐富,而「餅」代表陸上的豐富.上主早已為他們預備了魚和餅,這是他們最好的食物.既然已有餅和魚,為何耶穌要門徒把剛才打的魚,拿幾條來?(約21:10)是因為不夠吃嗎?他們共有八人,耶穌沒有預計這麼多人吃嗎?但祂是復活主,是神,怎會不知道呢?或是剛捕到的魚比較新鮮味美?因為耶穌的魚是變出來的或隔夜的嗎?原因是耶穌要彼得發現這神蹟的奇妙.耶穌一說,大哥彼得就行動,把網拉到岸上,那網滿了肥美和價值高的大魚,共有一百五十三條,是個很大和驚人的數目,這數字更是有意思或意義,這是上主也數過的,是所有魚的分類數目,全共界民族的數目.還有,魚雖這麼多,網卻沒有破,沒有一條魚失掉,是上主的保守的大能,這網也可能是剛好只能承受一百五十三條魚,可見到上主能控制一切.

「耶穌說:『你們來吃早飯.』門徒中沒有一個敢問祂:『你是誰?』因為知道是主.」(約21:12)精美的早餐,從猶太人的食物來看,就是魚,蛋,蜂蜜,餅,麵包穀物,耶穌為門徒預備的是營養早餐,比麥當奴,大家樂的早餐還要好.在飢餓,失望,疲乏中吃鮮魚早餐,特別能夠感覺主的大愛.主在適當的時候,做適當的事,令人難忘,所以愛要懂表達.

復活主於鮮魚早餐後的考問

吃完精美的早餐後,復活主給彼得愛的召命.耶穌的論點是,若彼得愛主,就當餵養主的羊.吃並不是最終目的,吃飽後一起談心事,考問和託付使命才是最終目的.

這裡有天時,地利,人和.天時就是早上,捕魚後,彼得和其他門徒從無到有所獲,心情好的時候；地利就是在加利利海邊,有海風在吹,是個很舒服的環境；人和就是吃完早餐,吃飽後,他們都感到很滿足.

復活主與門徒相處交流的方法有很多,祂不用命令式,教訓式,恐嚇式,利誘式,但會用對話式,互動式,親切式,激發方式等,但更常用的是愛比較式,愛深入式,因愛而為,為的就是預告,祂將會變得孤單,又會失去自由,受苦並走上不願意的路.於是,復活主要考問及託付使命給彼得.當時的彼得肯定記得,數日前在大祭司的園子裡,一班僕人和差役生了炭火,有人問他是否跟從耶穌的人,彼得曾三次不認主,彼得當時心情矛盾,既開心見到復活主,卻又想避開復活主,最怕耶穌舊事重提,在眾多門徒面前,他便顏臉無存了.

這餐不是常餐,是非常餐,既舒服溫馨,又帶點難受,帶些悲情,帶點師生味,帶點愛情；當中有些恐懼,又點意外的驚喜,百感交集.耶穌並沒有問彼得:「彼得,你記得,你曾經三次不認我嗎?」這令人有秋後算帳的感覺!也沒有問:「彼得,你為何三次不認我?」這更是質問,把對方迫到牆角,沒有退路.體貼人的主,不會這樣問,因為這樣問已沒有意思了,只會弄巧反拙.耶穌三次問的方式與彼得三次的回答更是「高手過招」,值得深究.

【第一次考問】

「他們喫完了早飯,耶穌對西門彼得說:『約翰的兒子西門,你愛我比這些更深麼?』彼得說:『主阿,是的；你知道我愛你.』耶穌對他說:『你餵養我的小羊(初信或年少的信徒).』」(約21:15)

〔原文〕「愛」是神聖的愛,犧牲的愛,深謀遠慮的愛(主耶穌問話裡所用的字原文是agapao)；友誼的愛,回應的愛,情感的愛(彼得答話裡所用的字原文是phileo)      ,另外還有比較的愛,優次的愛.

「你愛我比這些更深麼?」「這些」若作主詞解,則本句意即:

「你愛我比這些門徒愛我更深麼?」「這些」若作受詞解,則可有兩種解法:(1)「你愛我比愛這些東西(指船,漁網,魚,餅,炭火等)更深麼?」(2)「你愛我比愛這些人(指其他門徒)更深麼?」

【第二次考問】

「耶穌第二次又對他說:『約翰的兒子西門,你愛我麼?』彼得說:『主阿,是的；你知道我愛你.』耶穌說:『你牧養我的羊.』」(約21:16)

﹝原文﹞「愛」神聖的愛,犧牲的愛,深謀遠慮的愛(主耶穌問話裡所用的字,agapao)；友誼的愛,回應的愛,情感的愛(彼得答話裡所用的字,phileo).

﹝文意﹞「你牧養我的羊」,「牧養」不僅僅是餵食物,還包括看顧,監察,引導和保護等；「羊」指一般的羊,即信主多年或中年的人,不再是小羊.

【第三次考問】

「第三次對他說:『約翰的兒子西門,你愛我麼?』彼得因為耶穌第三次對他說,你愛我麼,就憂愁,對耶穌說:『主阿,你是無所不知的,你知道我愛你.』耶穌說:『你餵養我的羊.』」(約21:17)

﹝原文﹞「愛」友誼的愛,回應的愛,情感的愛(包括主耶穌的問話和彼得的答話裡所用的字,phileo).「知」是指裡面主觀的知覺；「知道」是外面客觀的認識.

分析

復活主從愛心早餐考問中,轉化彼得,將他由失落變到回轉.復活主以愛為動力,彼得若愛主,就會愛教會中的羊,愛的方式是餵養,即以上主的話,餵養門徒.以愛的行動:關懷,照顧門徒的身心需要.

復活主要彼得回答,你愛我麼?不是復活主不知,而是信仰要認信,要讓其他人也知他對復活主的愛.

默觀式靈修:想像復活主對我說話

許多信徒以這段經文,即耶穌對彼得的考問作靈修,我也曾試過.在靈修中,我想像復活主,如此問我:「許開明,你愛我比這些更深麼?」這些可能是:工作,學業,家庭,喜好.我說:「我愛不愛你,你是知道的!」我每天在教會工作超過14小時,常常探病和探訪教友,又經常講道,並且用了很長時間準備講道,亦勤於行政工作,推動教會事工等.復活主回答說:

(1)「你若愛我,你要好好的照顧我給你的身體,給身體多休息!」

因為我曾經很用心事奉,進修,教學,寫作或寫論文,參與學術會議,每天只睡四小時.

(2)「你若愛我,你要好好的照顧我給你的身體,你要約束你的飲食.」

因為我自小很喜歡吃,並且胃口很好.牧會期間,會友很喜歡吃,特別是在跑馬地牧會,查完經後,就去找美食,吃到身體也有不良反應,體重不斷加增.

(3)「你若愛我,除了牧養教會這大家庭,也不要忽略了我給你的家庭,你的家人,也是我給你的羊,你也要以上主的話餵養他們,和關顧他們.」

我曾經因為教會建堂,日間除一般牧會工作外,也有許多重建教堂的事務要負責,晚上有無數的會議,因而忽略家人.

以上三次的靈修默想,都有價值,使我在身體得到調整和與家人的關係得到改善.你也可以方式默想.

愛主失自由,走苦路

「我實實在在的告訴你,你年少的時候,自己束上帶子,隨意往來,但年老的時候,你要伸出手來,別人要把你束上,帶你到不願意去的地方.」(約21:18)

【年少有自由】

「你年少的時候,自己束上帶子,隨意往來」,象徵當彼得未認識主時,作事隨己意,不被約束,但生命並無目標,自由自在地做自己想做的.這種自由只滿足自己的慾望,以「搵食」為中心.

【年老失自由】

「但年老的時候,你要伸出手來,別人要把你束上,帶你到不願意去的地方.」這象徵了當彼得相信復活主,就能甘心忍受別人所給予的不便與束縛.年老的彼得身不由己,失去自由.這種無自由,是為主的緣故,甘願犧牲自己的自由.

「耶穌說這話,是指著彼得要怎樣死來榮耀神(被判死時,要求倒轉釘十架).說了這話,就對他說:『你跟從我罷.』」(約21:19)就是要他跟耶穌死.耶穌向彼得提出,並挑戰彼得重新來跟從主,因彼得曾不認主,離開主,現在悔改再跟從主.彼得有自由去選擇:跟從或不跟從耶穌,他選擇了跟耶穌,甘願跟從主而失去自由.彼得找到比餬口更重要的,就是完成基督的使命.

結語

復活主在加利利海邊特別為門徒設立的愛心鮮魚早餐,是一個早餐的團契,是擘餅的聖餐聚會.這一餐是改變生命的聖餐,使彼得及其他門徒生命得回轉(Turing back).當耶穌面向十字架步向死亡,彼得曾一次又一次的不認主,不再相信和跟隨主.當婦女傳播耶穌復活的好消息後,彼得仍然無動於衷,為要餬口,便重操故業,打魚去.

復活主在愛的早餐中,使彼得回轉.第一,主先讓他看見他的大能,從無到有,再到有餘,顯出了轉化和優化的大能；第二,主挑起他的愛心,要他愛主也愛主的羊；第三,主給予他餵養群羊使命,鼓勵他回轉重新跟從主,結果彼得成為上主重用的門徒,更被封為第一位教皇.在信仰靈程中,我們也會失落遠離主,如果有的話,願我們得到回轉歸向復活主.

在這愛心早餐中,彼得看到復活主耶穌比他所捕的魚獲或財富重要,甚至比他的自由和生命重要,而你又如何呢?彼得不會以食為先,因他回轉後的生命乃以「主」為先了.

祈禱

天父上帝,感謝你賜人食物,使人得以生存.然而,「食」只是人生的一部份.求主教導我們懂得食,不被食慾操控,反能操控食慾,藉著食物得到健康力量,多作主工；藉著一起食,增長人神與之間團契.阿們.



\newpage

\section{被殺的羔羊與掌管萬有的主}
\label{sec:318}
\textbf{伍渭文牧師 - 被殺的羔羊與掌管萬有的主}
\newline
\newline
\hyperref[sec:317]{< < < PREV SERMON < < <}
~
\hyperref[sec:toc]{[返目錄]}
~
\hyperref[sec:319]{> > > NEXT SERMON > > >}


啟示錄 7:9-17
\newline
\begin{longtable}{cl}
\hline
\hline
章節 & 經文 (和合本修訂版)\\
\hline
7:9 & \begin{tabularx}{0.7\textwidth}{X} 此後，我觀看，看見有許多人，沒有人能計算，是從各邦國、各支派、各民族、各語言來的，站在寶座和羔羊面前，身穿白衣，手拿棕樹枝， \end{tabularx} \\ \\ \relax
7:10 & \begin{tabularx}{0.7\textwidth}{X} 大聲喊著說：「願救恩歸於坐在寶座上我們的神，也歸於羔羊！」 \end{tabularx} \\ \\ \relax
7:11 & \begin{tabularx}{0.7\textwidth}{X} 眾天使都站在寶座和眾長老，以及四個活物的周圍，俯伏在寶座前，敬拜神， \end{tabularx} \\ \\ \relax
7:12 & \begin{tabularx}{0.7\textwidth}{X} 說：「阿們！頌讚、榮耀、智慧、感謝、尊貴、權能、力量都歸於我們的神，直到永永遠遠。阿們！」 \end{tabularx} \\ \\ \relax
7:13 & \begin{tabularx}{0.7\textwidth}{X} 長老中有一位回應我說：「這些穿白衣的是誰？是從哪裡來的？」 \end{tabularx} \\ \\ \relax
7:14 & \begin{tabularx}{0.7\textwidth}{X} 我對他說：「我主啊，你是知道的。」他向我說：「這些人是從大患難中出來的，他們曾用羔羊的血把衣裳洗得潔白。 \end{tabularx} \\ \\ \relax
7:15 & \begin{tabularx}{0.7\textwidth}{X} 所以，他們在神寶座前，晝夜在他殿中事奉他；那坐在寶座上的要用帳幕覆庇他們。 \end{tabularx} \\ \\ \relax
7:16 & \begin{tabularx}{0.7\textwidth}{X} 他們不再飢，不再渴；太陽必不傷害他們，任何炎熱也不傷害他們， \end{tabularx} \\ \\ \relax
7:17 & \begin{tabularx}{0.7\textwidth}{X} 因為寶座中的羔羊必牧養他們，領他們到生命水的泉源；神必擦去他們一切的眼淚。」 \end{tabularx} \\ \\
[1ex]
\hline
\hline
\end{longtable}
講題:被殺的羔羊與掌管萬有的主　The Lamb of God and Ruler of All (Pantocrator)

經文:啟示錄7:9-17

講員:伍渭文牧師

場合:香港中文大學崇基學院禮拜堂

日期:2016年4月17日

「此後,我觀看有許多的人,沒有人能數過來,是從各國各族各民各方來的,站在坐在寶座和羔羊面前,身穿白衣,手拿棕樹枝,大聲喊著說:願救恩歸與坐在寶座上我們的神,也歸與羔羊.眾天使都站在寶座上我們的神,也歸與羔羊.眾天使都站在寶座和眾長老並四活物的周圍,在寶座前,面伏於地,敬拜神,說阿們,頌讚,榮耀,智慧,感謝,尊貴,權柄,大力,都歸與我們的神,直到永永遠遠.阿們!……這些人從大患難中出來的……他們不再饑,不再渴；太陽必不傷害他們,任何炎熱也不傷害他們,因為寶座中的羔羊必牧養他們,領他們到生命的泉源；上帝必擦去他們的眼淚.」(啟七9-17)

IMG-20160415-WA0001

引言:島子「過紅海」聖水墨說起.2012第五屆崇基基督教文化節藝術家是清華大學島子(王敏)教授, 接受基督教媒體訪問,與時任校牧對談基督教藝術創作的社會性,談到參展作品「過紅海」-一幅沒有在畫冊出現的作品.因為策展人一看很感動,還沒有拍攝就放進玻璃畫框.島子說創作這畫時,有一種很沉重的感動,他說:「中國的近代的歷史,我們中國還沒有過紅海.……國父在辛亥革命之前就是基督徒,他要引領國人「走向共和」.……這個「共和」的道路是沒有完成的.我們現在還是一個極權社會……80年代鄧小平打開了計畫經濟的一個缺口,引入資本主義的經濟進來,但是很快這種資本主義又強化了極權統治.正因這個感動,所以畫了這幅《過紅海》,那個是鋪天蓋地的巨浪壓過來,霎時把法老的追兵全部都撲到底下,它是用大幅度的潑墨和破墨完成.然後,手持神杖的摩西在這邊已經上岸了,在呼召他的子民趕快走,在畫面底部用很小筆觸,用焦墨畫的.你仔細看,這個巨浪是一個字母 “v”,victor.」

我回應說:「上面的巨浪壓下來,與另一幅畫「逾越節」的受害者(victim)-被殺的羔羊是一樣的,耶穌是得勝者(victor) 與受害者(victim),耶穌已經被殺了,同時因他被殺,我們也可以得勝,把追趕我們的罪壓下去,過紅海.……得勝者(victor)同時要成為受害者(victim),才能成就救贖.」 過紅海事件中,救贖者是得勝者(victor),使海水分開成坦途讓選民走過,又使海水復合淹滅埃及追兵,使選民脫離為奴之地.同時這救贖者也是受害者(victim).耶穌就是那逾越節的羔羊,這羔羊的血塗在門框,門楣,審判天使見到代贖的血,就越過這家庭.事實上,耶穌在登山變時,摩西跟以利亞談論耶穌「去世」之事,原文就是出埃及(路九31).

耶穌復活後同時是掌管萬有的主,和被殺的羔羊,這是今天啟示錄經課給我們的信息.

1 復活的吊詭性-同時是被殺羔羊和掌管萬有的主The Victim and Victor at the same time

「「此後,我觀看有許多的人,沒有人能數過來,是從各國各族各民各方來的,站在坐在寶座和羔羊面前,身穿白衣,手拿棕樹枝,大聲喊著說:願救恩歸與坐在寶座上我們的神,也歸與羔羊.眾天使都站在寶座和眾長老並四活物的周圍,在寶座之前,面伏於地,敬拜神,說阿們.頌讚,榮耀,智慧,感謝,尊貴,權柄,大力,都歸與我們的神,直到永永遠遠.」(啟七9-12)

教導神國的道理:今天是復活後第四主日,復活使耶穌從被殺的羔羊成為掌管萬有的主.教會年的復活期共五十天,直到聖靈降臨節,開始教會年的下半年,這時期我們稱國度期.教會年上半年關基督的生平,將臨期是他降生的預備期,跟著聖誕期,顯現期(約但河受洗開始,登出變像結束).耶穌的死的預備期是四十天的預苦或大齋期,其高潮是聖周的受難和死後三天的復活.

復活後耶穌四十天對門徒顯現,教導他們天國的道理.耶穌正要升天之時門徒問他:「主阿,你復興以色列國,就在這時嗎?」(使－6)若耶穌擁有天上地下所有的權柄,是掌管萬有的主,是得勝者Victor,就應該顯明可見的國度,推翻羅馬鐵蹄統治,復興以色列國,還要再等多久?耶穌對他們說:「父憑著自己的權柄,所定的時候日期,不是你們可以知的.」(徒－6下) 耶穌指出上主的國度是隱藏,我們不會完全明白,就像耶穌復活的清晨..

隱藏的國度: 耶穌降生之時,有異星橫空,吸引遠方的賢士千里迢迢來到耶路撒冷,尋找新生王.耶穌誕生當晚,伯利恆上景星明亮似晝,天使歌聲響澈雲霄:在至高之處榮耀歸於上主,在地上平安歸於他所喜悅的人.耶穌死的時候,從午正至申的三個小時,遍地為黑喑所吞噬,日頭變黑；氣斷時,聖殿分隔神人的幔子裂開,地大震動,磐石也崩裂,墳墓開了.但復活的清晨,溫度,氣壓如常,墓園的草如常沾滿了露水,雀鳥鳴叫,好像甚麼事都沒有發生過,祗是墓石滾開了,若探頭往墓裡頭看看,才見到裹屍布像人形放在石棺上,耶穌屍體不見了.

耶穌作為被殺的羔羊的救贖歷史,眼睛可見,有跡可尋.但復活升天後被殺的羔羊成為掌管萬有的主,就不是憑眼見,需要憑信心.經過四十天的教導,門徒仍以為上帝的國是眼睛看得到的得勝者(victor).其實上帝的得勝就在十字架的道理:甘心成為受害者(victim),犧牲自己,以善勝惡,以愛勝恨.

願意受傷(vulnerable):就是在此刻,耶穌也像昔日被人拒絕輕視,包括信徒.耶穌不想以掌管萬有者的身份去強制我們跟隨他,要我們甘心順服,所以接受被拒絕和被傷害(vulnerable).耶穌升天在我們的視線隱退,他要我們運用信心不憑眼見,曉得他並沒有離開我們,他藉聖道,聖禮和聖徒親近我們,甚至透過不完全甚至軟弱的人代表他傳講他的話語,接受被拒絕和不順服的可能.有時我們厭煩他管得我們太嚴,輕視他的教訓,甚至重釘他十字架,耶穌是我們愚頑悖逆的受害者(victim).

耶穌升天說他會再來,但數十年了,信徒還被羅馬鐵蹄統治,地震天災饑饉延年,更因為拒絕羅馬皇帝該撒的造神運動,堅持祇有一位上帝,基督是萬王之王,受到空前的逼害,猶其是尼祿(Nero, 54-68)和多馬田(Domitian, 81-96)時期.約翰寫啟示錄的目的,指出上主是掌管歷史,正在運籌維幄,讓歷史成就他隱密的旨意,上主不單將來掌權,現在他也坐在寶座上,教會要定睛終末的盼望,最後審批,和新天新地.

啟示錄書中的七印就說明當時的情況.我們要知道,雖然啟示錄有很多我們現在看不明白的象徵,但作為書信,收信的人是明白的.我們很簡單講講啟示錄的結構,特別是今天講道經文的上文脈絡.

一至三章:地上的七教會,他們面對不同的屬靈試探,其中我們最熟悉的老底嘉教會,不冷不熱.七教會不同的試探也是歷代教會的試探,每一位信徒的試探,我們要勝過試探,成為得勝者.教會是那位稱為阿們,誠信真實的見證人(啟三14),抗衡地上國度的虛假和不義.

第四章: 天上的敬拜.四活物:代表受造界飛得最高的鷹,最有智慧的人,野生最強的獅和受飼養最強的牛.二十四位長老:代表信仰群體的舊約十二支派和新約十二使徒,他們和四活物都在敬拜坐在寶座的上主.四活物遍體長滿眼睛,表示地上發生的事上主完全知道,寶座前的玻璃海表示連象徵邪惡的海都在上帝旨意之內,雖然有獸會從海裡上來,進行殺戮(十三章),但最後新天新將出現,海也不再有了(二十一章).

第五章: 被殺羔羊.祇有他才能打開封了印的書卷,讓我們窺探到歷史的意義和人類的命運.

第六,七章: 打開六印.啟示錄七9-17是揭開六印後的凱歌,現在讓我們一起打開書卷的六個印,了解約翰當時面對的世界.

第一至四印是一整體,提到四匹不同顏色的馬.一印白馬,騎士拿弓象徵征戰,征戰原因可能是擴張領土,宣揚羅馬和平的大夢,但征戰經年造成國庫虛空.二印紅馬,騎士拿刀奪去太平是內戰.三印黑馬,騎士手執天平,量度一升麥子,配給一個人一天的糧食；但天平另一邊是酒和油,指出路旁有饑俘,朱門酒肉臭.四印灰馬,騎士是死亡,手是空的,是人自己引至死亡.因耀武揚威,四出征戰稱霸；因為內部權力鬥爭,失去彼此的信任和太平,造成社會不穩；大量資源用作維穩,引至經濟下滑,田地荒廢,產生饑荒.四印是當時社會的真實寫照,是信徒的親身經驗,祇是他們不太明白,為何上帝讓這些事發生.我們現在看看當時的歷史記載.

歷史記載,公元60年,發生地震；公元62,羅馬軍隊在東面被波斯(現伊朗)軍重挫,死傷無數；公元64羅馬大火,尼祿王嫁罪於基督徒,基督徒被扔進獅子坑,或被活活燒死；其實祇要他們被捕後踐踏十字架,稱尼祿為主就可以得釋放.公元70年,羅馬提多將軍鐵腕結束四年猶太人的獨立運動,聖殿夷為平地,六千猶太人被殺.耶穌已在福音書預言,因為猶太人拒絕彌賽亞,「這一切的罪,都要歸到這世代了.」(太二十三30)因為你們殺害先知,你們的家,要成為荒場.(太二十三37-39)公元68,尼祿崩,羅馬四將領爭奪皇位,延年內戰,民不聊生,經濟下滑,糧食短缺.公元68年,義大利拿不列斯(Naples)威蘇爾(Vesuvius)活火山爆發,整個龐貝城(Pompeii)二萬人剎那被埋.四印代表人世間的征戰,內戰,饑饉和死亡.

約翰時代的信徒問:上主真的掌管萬有嗎?上帝的國何時降臨呢?

應用:回歸以來我們不斷談論主權的問題,我們關心十年後香港變成如何呢?我們感到這城市好像愈來跟以前不同了.現在是否為時已晚呢?啟示錄七章的凱歌提醒我們,上帝不祇是終末的上帝,也是掌管萬有的主,歷史正一步一步成就上主的計劃.伊斯蘭國的恐怖襲擊引起全球恐慌,但有從事回教國家宣教士報告,正因為這些行為,在伊拉克地區,大批回教徒離開原來的信仰,歸向基督信仰,這是我們想也未想過的.

一本探究18世紀法國大革命的著作《舊制度與大革命》,在國內知識界引起廣泛關注,王岐山強力推薦,這書是警示當前中國社會的處境與法國大革命前期有相似性.當前中國與大革命前的法國都處於最繁榮的時期,物質財富的增加促使人民權利意識的覺醒,對特權,腐敗,不公正的容忍度更低.中國經歷三十多年的改革開放,經濟快速發展後,出現了利益分化嚴重,貧富差距加大,社會不公等現象.是中國最好的時候,也是中國最危險的時刻.

最近網上瘋傳中國政法大學叢日雲教授在畢業典禮上的演講,他指出未來社會將產生巨變.目前局勢風雲變幻,暗潮湧動,前途莫測.面對這巨變,勇敢者與膽小者的做法,關鍵是前者不讓自己的人性在隨波逐流中泯滅.即使不敢做抗爭者,也不要去對抗爭者背後放冷箭,助紂為虐.他希望自己的學生在大潮襲來時,選擇站在理性一邊,站在人民一邊.

不讓人性隨波逐流,為真理作見證,就是啟示錄七章所說的殉道者.

2誠信真實的國度

「長老中有一位問我說,這些穿白衣的是誰,是從那裡來的.我對他說,我主,你知道.他向我說,這些人是從大患難中出來的,曾用羔羊的血,把衣裳洗白淨了.所以他們在神寶座前,晝夜在他殿中事奉他.坐寶座的要用帳幕覆庇他們.他們不再饑,不再渴.日頭和炎熱,也必不傷害他們.因為寶座中的羔羊必牧養他們,領他們到生命的泉源.神也必擦去他們的眼淚.」(啟七13-17)

信徒都是殉道者-見證人.穿白衣的人,在揭開第五印時出現過.「在祭壇底下,有為神的道,並作見證,被殺之人的靈魂大聲喊著說,聖潔真實的主阿,你不審判住在地上的人給我們伸流血的冤,要等到幾時呢?」(啟六9) 這些被殺的殉道者,他們的血成印了天上祭壇的祭品,為上帝所悅納.約翰肯定想到尼祿和多馬田大迫害中的大量殉道者.殉道原文作見證者,他們「為神的道,並作見證」而被殺.殉道者原文就是見證人,他們見證上帝國的真實,揭示地上國度的虛妄；誠信真實和說謊虛妄,在啟示錄是常用的對比概念.

上主國度是真實誠信,對這些迷惑眾人的謊話不能保持箴默,要用真理揭露謊言.羔羊是「那位阿們,誠信真實的見證者」(啟三14),有七角七眼(啟五6)角表示力量,眼表示智慧,羔羊不是因軟弱而被殺,是因為世界拒絕他所說的真理而被殺.說謊者容不下真理,撒旦就是說謊之父,一直「迷惑住在地上的人」的獸(啟十三14),在最後審判,一切虛謊的人都不能進入新天新地.(啟二十一8, 27, 二十二15)

穿白衣者已被上主稱為義,他們為神的道,為誠信真實作見證而失去生命.尼祿,多馬田不是神,寧可被扔進鬥獸場面對餓獅,也不能褻瀆上帝.殉道者的誠信真實,揭示當時社會的虛妄.

地上國度的虛妄:要揭露虛妄, 就要指出甚麼是誠信真實,活出不一樣的價值觀生活.當每一個人都唯利是圖,有人開茶餐廳免費招待窮人；當做業主的不斷加租,把空間改裝又改裝,割成tong 房,有人願意以低於市值租給這位善心茶餐廳老板,因為原來的餐廳要加租了.這種願意吃虧的自我犧牲行為就說明,除了利潤,還有更美好的東西.這是一種救贖的行為,免得我們因金錢而迷失,甚至沉淪.

3 仁愛的國度

被殺羔羊永遠提醒我們信仰的核心是愛.凱歌一開始,我們看到一個宏大普世的群體,因為羔羊犧牲的愛被召喚成為一個群體.「此後,我觀看,有許多人,沒有人能計算,是從各邦國,各支族,各語言來的,站在寶座和羔羊面前,身穿白衣,手拿棕樹枝.」(啟七9)

羅馬軍團強大,但愛的團契更強大；軍團用來鎮壓殺戮,摧毀生命；愛用來是重建盼望,帶來新生,就像今天使徒行傳經課的多加(Dorcas)和她的寡婦團隊.

彼此相愛:今天首輪經課使徒行傳九36-43記載:「在約帕(Joppa,現在的台拉維夫)有一個女徒,名叫大比大(Tabitha),繙希利尼話就是多加(「多加」就是「羚羊」的意思)；她廣行善事,多施賙濟.當時,她患病而死,有人把她洗了,停在樓上.呂大與約帕相近；門徒聽見彼得在那裡,就打發兩個人去見他,央求他說:「快到我們那裡去,不要耽延.」(使九36-38) 認識多加的信徒很愛多加,派人找彼得來施行神蹟使多加復生.「眾寡婦都站在彼得旁邊哭,拿多加與他們同在時,所做的裹衣給他看.」(使九39)眾寡婦失去丈夫,無依無靠,感戴多加幫助.多加本身不富有,所以自己縫紉裡衣給眾寡婦.她們不捨多加離開,希望彼得使她復活,這是一個多麼感人愛的團契.彼得沒有金子,銀子,但有上帝給他的權能,也愛寡婦們和多加,他跪下禱告,求上主使多加復生.

神的國在世上沒有羅馬帝國的財富和軍隊,但有愛的團契教會和屬靈的權能.今天詩篇經課詩篇二十三篇:「耶和華是我的牧者,我必不致缺乏」,上帝藉聖徒這愛的團契,使我們不致缺乏.在復活期升天之後,耶穌藉聖道,聖禮,和聖徒與我們同在,他並沒有離開我們.

多加患病,最後死去,但多加不懼怕死亡,因為上帝的羔羊已經被殺,而且第三天從死裡復活,成為掌管萬有的主.使徒行傳多加的復生,也因為一群愛她的人不捨,找緊叫人請彼得來,當然她們也自己為多加禱告；因為她們的愛,上主藉彼得施行神蹟.多加復生的神蹟,始於多加本人關心人的需要,雖然自己緊絀,但用手為寡婦一針一針縫紉裡衣,神蹟就這樣開始.失去丈夫的寡婦因為多加的愛,告別苦澀,驅走焦慮,也開始關心其它的人,成為服事的團隊.聖經說「眾寡婦都站在彼得旁邊」,這圖畫使人感動.

多加生病離世,聖經沒有提到她的家人.管會堂睚魯的閨女病重,作為父親的眶魯跪在人前,懇求耶穌醫治；拉撒路死了,我們見到馬大,馬利亞趕快請耶穌來,並為拉撒路離開而痛哭.但多加生病離開,始終沒有看見她的親人出現.可能多加本身是孤兒,守獨身的孤兒；可能是寡婦,丈夫早逝而沒有兒女的寡婦.但多加並不孤單,她視教會為她的家,專心幫助有需要的人,特別是這班失去依靠,孤單的寡婦；不單幫助他們,且訓練她們成為事奉的團隊.

事實上,若我看整章使徒行傳,我們發現這愛的團契的身影.保羅在大馬色悔改後失去視力,亞拿尼亞受差遣使他看見並接待保羅.眾人懼怕保羅,巴拿巴引他見眾使徒.彼得周游傳道,呂大的聖徒接待他,多加教會代表派出兩個代表到呂大找彼得.耶穌升天,聖靈降臨,教會誕生,教會是聖靈的團契,也是愛的團契.

被殺羔羊被棄絕,帶來教會的誕生:在十字架上,耶穌好像被父神遺棄,極其孤獨,他引用詩篇二十二,喊叫:「我的神,我的神,為何離棄我,為甚麼遠離不救我.」(詩二十二1)但同時耶穌感到不孤單,因他看見藉著他的死,教會將會誕生,像母親孤獨承受生產的痛苦,因看見快要來臨的嬰孩.基督被棄絕,帶來了修和和教會友愛群體的誕生:「我要將你的名傳與我的弟兄,在會中我要讚美你.」(詩二十二22) 希伯來書把這句話用在教會身上,但在前面加了:「原來那為萬物所屬,為萬物所本的,要領許多的兒子進榮耀裡去……所以,他稱為他們為弟兄也不以為恥,說:我要將你的名傳與我的弟兄,在會中我要頌揚你.」(來二10-12)希伯來書的作者著跟說:「兒女既同有血肉之體,他也照樣親自成了血肉之體,特要藉著死敗壞那掌死權的,就是魔鬼,並要釋放一生因怕死而為奴僕的人.」(來二14, 15)

結論:賺回你的生命Earn It

掌管萬有的主,不單掌管歷史,也掌管我們的一生.我們不能打開生命的書卷,祇有被殺的羔羊才能打開；十字架使我們明白生命的奧秘,步履耶穌的腳蹤,才活出有意義的生命.

被殺的羔羊能打開你封閉的書卷.也許你生命的書卷封閉了而且打上印,不是你不明白,乃是你太明白,因為不敢面對,把書卷封閉.今天耶穌可以把書卷打開,他要赦免你一切的過犯,因為這些過犯,耶穌被釘在十字架上.我們需要耶穌,因他有能力.為何需要耶穌,因為我們人太敗壞,罪孽太深重,才需要耶穌.今天耶穌賜你自由,可以因他的接納,接納你的過去,讓被殺的羔羊,打開你封閉的書卷.

被殺的羔羊,永遠提醒我們,他是愛我,為我捨己.說到付重價的拯救,使我想起電影「雷霆救兵」(或譯「拯救大兵瑞恩」,Saving Private Ryan).1944年諾曼第戰事中,三位親兄弟戰死沙場,作為寡婦的媽媽傷心欲絕,但還有一子瑞恩(Ryan)在前線生死未卜.美101空降師派出隊長約翰.米勒(John Miller) (湯漢斯飾)組成拯救隊尋找Ryan,送他返家,以慰其母親連喪三子之痛.拯救隊損折兩位隊員後終於找到Ryan,他正在一橋頭頑抗挺進的德軍.Ryan不忍離棄戰友,拒絕回家,他說: 「這是我緊有的兄弟(The only brothers I have left).」最後米勒的拯救隊與Ryan 並肩作戰,終於保橋成功；不過,差不多所有拯救隊成員都犧牲了-為了一個人而犧牲.隊長米勒臨終說:Earn it!很難翻譯,直譯是「賺回來吧!」其實意思是:「過一個與蒙拯救代價相稱的生活,證明你值得付出如此的代價得回你的生命.」

各位,我們正站在大時代的門檻,這是黑暗的時代,也是光明的時代,讓我們因著對世人的愛,做一個見證人,為誠信真實的耶穌作見證,因為他是被殺的羔羊,也是掌管萬有的主.阿們.



\newpage

\section{平安主賜勇往直前}
\label{sec:319}
\textbf{許開明牧師 - 平安主賜,勇往直前}
\newline
\newline
\hyperref[sec:318]{< < < PREV SERMON < < <}
~
\hyperref[sec:toc]{[返目錄]}
~
\hyperref[sec:320]{> > > NEXT SERMON > > >}


約翰福音 14:21-31
\newline
\begin{longtable}{cl}
\hline
\hline
章節 & 經文 (和合本修訂版)\\
\hline
14:21 & \begin{tabularx}{0.7\textwidth}{X} 有了我的命令而又遵守的人，就是愛我的；愛我的人，我父要愛他，我也要愛他，並且要親自向他顯現。」 \end{tabularx} \\ \\ \relax
14:22 & \begin{tabularx}{0.7\textwidth}{X} 猶大（不是加略人猶大）問耶穌：「主啊，為甚麼親自向我們顯現，而不向世人顯現呢？」 \end{tabularx} \\ \\ \relax
14:23 & \begin{tabularx}{0.7\textwidth}{X} 耶穌回答他說：「凡愛我的人就會遵守我的道，我父也會愛他，並且我們要到他那裡去，與他同住。 \end{tabularx} \\ \\ \relax
14:24 & \begin{tabularx}{0.7\textwidth}{X} 不愛我的人就不遵守我的道。你們所聽見的道不是我的，而是差我來之父的。 \end{tabularx} \\ \\ \relax
14:25 & \begin{tabularx}{0.7\textwidth}{X} 「我還與你們在一起的時候，已對你們說了這些事。 \end{tabularx} \\ \\ \relax
14:26 & \begin{tabularx}{0.7\textwidth}{X} 但保惠師，就是父因我的名所要差來的聖靈，他要把一切的事教導你們，並且要使你們想起我對你們所說的一切話。 \end{tabularx} \\ \\ \relax
14:27 & \begin{tabularx}{0.7\textwidth}{X} 我留下平安給你們，我把我的平安賜給你們。我所賜給你們的，不像世人所賜的。你們心裡不要憂愁，也不要膽怯。 \end{tabularx} \\ \\ \relax
14:28 & \begin{tabularx}{0.7\textwidth}{X} 你們聽見我對你們說過，我去了還要回到你們這裡來。你們若愛我，就會因我到父那裡去而喜樂，因為父比我大。 \end{tabularx} \\ \\ \relax
14:29 & \begin{tabularx}{0.7\textwidth}{X} 現在事情還沒有發生，我預先告訴你們，使你們在事情發生的時候會信。 \end{tabularx} \\ \\ \relax
14:30 & \begin{tabularx}{0.7\textwidth}{X} 我不再和你們多說了，因為這世界的統治者將到，他在我身上一無所能。 \end{tabularx} \\ \\ \relax
14:31 & \begin{tabularx}{0.7\textwidth}{X} 我這麼做是照著父命令我的，為了讓世人知道我愛父。起來，我們走吧！」 \end{tabularx} \\ \\
[1ex]
\hline
\hline
\end{longtable}
講題:《平安主賜,勇往直前》　Peace Lord Gives, Let Us Leave!

經文:約翰福音14章21-31節

講員:許開明牧師

場合:香港中文大學崇基學院禮拜堂

日期:2016年5月1日

「21有了我的命令又遵守的,這人就是愛我的．愛我的必蒙我父愛他,我也要愛他,並且要向他顯現.22猶大(不是加略人猶大)問耶穌說,主阿,為甚麼要向我們顯現,不向世人顯現呢.23耶穌回答說,人若愛我,就必遵守我的道．我父也必愛他,並且我們要到他那裡去,與他同住.24不愛我的人就不遵守我的道．你們所聽見的道不是我的,乃是差我來之父的道.25我還與你們同住的時候,已將這些話對你們說了.26但保惠師,就是父因我的名所要差來的聖靈,他要將一切的事,指教你們,並且要叫你們想起我對你們所說的一切話.27我留下平安給你們,我將我的平安賜給你們．我所賜的,不像世人所賜的．你們心裡不要憂愁,也不要膽怯.28你們聽見我對你們說了,我去還要到你們這裡來.你們若愛我,因我到父那裡去,就必喜樂,因為父是比我大的.29現在事情還沒有成就,我預先告訴你們,叫你們到事情成就的時候,就可以信.30以後我不再和你們多說話,因為這世界的王將到．他在我裡面是毫無所有．31但要叫世人知道我愛父,並且父怎樣吩咐我,我就怎樣行.起來,我們走罷.」

引言

假如有一日,一個人走到了人生的盡頭,他會對配偶,子女,家人說些什麼話呢?每個人總會有這一天的,所以這是值得每個人都思想的問題.人到臨終,其言也善,也有意思!約翰福音14章21-3節,說到當耶穌走向耶路撒冷,步向死亡的十字架,耶穌說了一些值得當日門徒深思的話.本主日經文,我們也要思想一下.

祈禱

上主,你的話語能甦醒人心,使愚人有智慧；你的訓詞正直,能快活人心,明亮人的眼目.祈求上主,祈求上主,使我們用耳和用心來聆聽.阿們.

在今日的福音經課,約翰福音14章21-31節中,耶穌講出一段語重心祥的臨終感言.

1耶穌臨離世時提出:去愛與被愛,不冷感!

「21有了我的命令又遵守的,這人就是愛我的．愛我的必蒙我父愛他,我也要愛他,並且要向他顯現.」

愛我的,我也愛他,我父也愛他,我要向他顯現,我和父要到他那裡,我們與他同住!人需要去愛及被愛,愛使關係建立；愛增進關係；愛更使關係維持.耶穌指出:去愛與與被愛的關係,這種關係是重要的,使人不冷感,愛可以超越生死.活到超過半百的人,開始發現有些朋友跟你道別,有些你所愛的朋友,仍在你的心中,愛使人得力量!

2耶穌臨離世時時提出:同在感

「23耶穌回答說,人若愛我,就必遵守我的道．我父也必愛他,並且我們要到他那裡去,與他同住.」

人需要同在感,使人不孤單.在人心深處,最可怕的是孤單,孤獨.聽坐監的人說,最可怕的不是失去自由,而是孤單,無人關心和傾訴,生不如死.

3耶穌臨離世時時提出:平安

「26但保惠師(幫助者,安慰者),就是父因我的名所要差來的聖靈,他要將一切的事,指教你們,並且要叫你們想起我對你們所說的一切話.(指教者,提醒者 )」

3.1 父派保惠師

保惠師即是幫助者(the Helper)和安慰者(the Comforter),是要指教你們一切的事,提醒你們主所說的一切話.

3.2 主賜的平安

「27我留下平安給你們,我將我的平安賜給你們．我所賜的,不像世人所賜的．你們心裡不要憂愁,也不要膽怯.(留下平安,給予平安,平安同在或有平安伴隨)」

從耶穌的神性來看,耶穌的平安是屬主的平安,無憂無慮的平安,所以祂能賜人平安.從耶穌的人性來看,耶穌的平安是屬人的平安,耶穌正面向痛苦十字架,墳墓與死亡.耶穌正面對著不安,一個正經歷不安的人,去叫人勝過不平,意義何在?祂明白你的不安.例如一個經歷過牙痛的牙醫比一個無有經歷過牙痛的牙醫不同；一個未曾患病的醫生和一個有醫病的醫生,對病人都有不同.

主這樣說:「我留下平安,給予平安,平安同在或有平安伴隨,我賜的是不一樣的平安.」你不要憂愁!你不要被憂愁所充滿!Do not let your heart be troubled!你不要膽怯,你不要被膽怯所充滿!Do not let your heart be afraid!許多時候,我們落入憂愁與恐懼中,因我們容讓它們存在我們心中,並被它們所操控,一個人被憂愁和恐懼操控,就活得不開心.

4耶穌臨離世時說出:他們將會喜樂

「28你們聽見我對你們說了,我去還要到你們這裡來.你們若愛我,因我到父那裡去,就必喜樂(Rejoice),因為父是比我大的.」

If ye loved me, ye would rejoice(喜樂,歡欣)因為我去,還要到你們這裡來.I go away and come to you again. 這說出了主還會回來.

5 耶穌臨離開時說出:他們的未來,會面對不安

「29現在事情還沒有成就,我預先告訴你們,叫你們到事情成就的時候,就可以信.」

在這一節,耶穌預先說了,讓他們可以預先知道,使他們回應,並能作好準備,面對未來,不會慌張.人在無把握的時候,會感不安或不快樂,失去安全感.

6 耶穌臨離開時說出:世界的王要來!

但他不能對我作什麼!他不能控制我,阻礙我!He has no hold on me!他不能在我身上作什麼.He has no hold on me.他沒有能力在我身上作什麼.He has no power on me. 耶穌的目的是,要使門徒安心,因為撒旦不能勝過他.

今日許多人向罪惡,強權,撒旦低頭,以致失去安全感,活在痛苦中.今日的人懼怕撒旦的能力多於懼怕上帝,敬畏上帝,以致罪惡橫行.自古到今,撒旦想人懼怕他,向他讓步,他就可以向前進一步；撒旦想你不讓步,表現緊張,做錯事,然後中他的計.有一位丈夫,他的妻子被另一位男士勾引,丈夫一時衝動之下,可能做出四種傻事,第一是傷害或殺害那位對不起他的男士；第二是傷害或殺害那位對不起他的妻子；第三傷害或殺害那位對不起他的男士及妻子；第四是自殘或自殺,因為他自覺自己是失敗者.

7 耶穌臨離開時說出:我們要起來,向前,迎戰!

「31但要叫世人知道我愛父,並且父怎樣吩咐我(the Father gave me commandment我父給我命令),我就怎樣行.起來,我們走罷(Get up, and let us go).」

耶穌鼓勵門徒起來,懷著替天行道的心,是勇敢面對世界的王.世界的王來了,我們一起向迎戰!

認識兩個世界觀

自古到今,人類早已存在了兩個世界觀:現實與理想；兩股力量:正與邪；兩種存在:天堂與地獄；兩種狀態:快樂與痛苦；兩種抉擇:走正與走歪；兩項勢力:建立與毀滅等.

現實的世界是希羅的思想,也是今日人的思想.現實世界是物質的世界,步向滅亡的世界；現實世界是罪惡的世界；現實世界是正,邪對立的世界,邪惡當道.反之,理想的世界是現實世界的理想和目標.米歇爾·德·蒙(法語:Michel de Montaigne；1533－1592)是法國在北方文藝復興時期最有標誌性的哲學家,指出人需要理想和目標才能好好的存在,不是因為人堅強,反而是因人太脆弱,無目標就無力活下去.布萊茲‧帕斯卡(Blaise Pascal；1623－1662),是法國神學家,哲學家,數學家,物理學家,化學家,音樂家,教育家,氣象學,他提出人是會思考的蘆葦,蘆葦是脆弱的,但思考卻使人堅強,人是愈思考愈強的動物.基督教認為理想世界是再創造；是救贖；是恢復原來上主創造的原狀,一個超越物質的世界,正義能必勝,上帝國能彰顯,新天新地來臨就是上帝掌權的世界.

現世的世界是物質的世界,是定型的世界,是走向腐敗,滅亡,沒有明天,絕望的世界.反之,基督的世界是屬靈的世界,是完美的世界,因上帝或基督是不屬於現世,是永生,是喜樂,是希望!柏拉圖(Plato)提出:兩個世界的融合,如何把理想的世界,引進現實世界?理想的世界是完美的善；而現實的需要改變.柏拉圖的理想國是烏托邦,所以把理想國實現在地上是可能的.神學家加爾文(John Calvin)也想過如何把基督的世界,上帝國,引進現今的世界,因為基督世界是完美的,而現世能引進基督的世界,就能改變.神學家巴特(Karl Barth)認為教會是一班跟隨基督的人,應當在世上扮演這角色,拉近這兩個世界,並把天國實現於地上.教會的角色和使命就是把基督的世界,上帝國引入教會,教會把這國度推展向社會,這應是教會的使命.

現時,大部份教會卻未能把上帝國引入教會,反而有些教會把現世的國度引入教會,教會內爭權奪利,明爭暗鬥,與世俗機構無異.

當教會推展事工,他們只向錢看,不是前看,更遑論向上帝的國度看,這可能把「經濟掛帥」的世俗思維引入教會.當教會的崇拜,只求滿足信徒的需求,要聽動聽好笑的道,要有動人的音樂,要有所得著,要有宏偉的聖堂或十架,這可能把「消費主義」帶進教會,信徒變成了客戶.當教會的領導明爭暗鬥,爭權奪利,只講求階級制度,無意或無意講求服從人多於神,明白某些人的心意多於尋求上主的旨意,

這可能已把現今世界的模式帶進教會.當教會以自身的發展作為存在的目的,不是朝向上帝國的方式,只追求增長,增多,增強,要人多,錢多,堂多,鬥大,門名,鬥人；自建教堂,自建神學院,培養自己特色的傳道人；自建書室,只賣自己批准的書籍；本派主義,本堂主義,本位主義等等.這種自我,自大的氣派,把現世大企業經營模式,一條龍方式引入教會,然後沾沾自喜.當教會只注重自身的發展,而非上帝國的拓展.教會要把那看不見的,理想上帝國或基督世界引進世界,有賴信徒有此異象,教會也應有此教導.

在約翰福音14:21-31節的經文,是耶穌上十架前,對門徒的臨終贈言,在苦難不安的世代,信徒如何自處,作出美好的提醒!

耶穌臨終遺言

耶穌以上的七項遺言,是理想的人生所需.第一需要是平安感,使人不憂不懼；第二需要是同在感,使人不孤單；第三需要是去愛與被愛感,使人不冷感；第四需要是喜樂,使人不憂慮；第五需要是知悉前路,使人不慌張；第六需要是盼望和機會,使人不失望；第七需要是勇敢,使人不怕站起,迎向敵人.

耶穌以上的七項遺言,主要是帶出「安全感」.心理家認為,人最需要的是「安全感」,人有安全感就能面對一切.得平安,就能以平安去回應神.

有一位殘暴的將軍,去到那裡都帶來不安,某村莊的百姓聽見這惡將軍的消息,只有逃走,但有一人不走,還敢面對將軍,將軍對他說:「你知我是誰?我可以眼都不眨,一劍穿心,把你殺死!」你猜那人怎樣回答,那人竟回答:「我知你是誰,我也可以告訴你,我也可以眼都不眨,接你一劍穿心!」那殘暴的將軍聽了,嚇了一跳,放走了這個敢對著他「眼都不眨」的人.那人以和平方式面對暴力,這不是怯懦,反而是英雄,這就是耶穌所教導的平安.以平安,去面對不安,反而得到平安.製造不安,去對抗令人不安者,只有引發更大更多的平安,而今日的恐佈主義者,就是以製造不安來威脅人.被威脅者以製造不安來反擊也只有產生不安,絕不能平息不安.耶穌在面對不安時,不是用暴力制止不安,而是使自己成為-位平安者,勇敢面對不安的製造者.平安是會產生一種力量,去制止「不安」.如果你想要平安,就為平安而努力,以平安去回應,從而得到平安.

相反,你以令人不安的手法,並不能帶來平安.恐怖主義者在周圍製造不安,如果英美各國也以令他們不安的手法回擊,只有帶來更大的不安,變成以巴衝突的翻版,使各國與恐怖襲擊者成為世仇,永無寧日,這正是撒旦的心意.

平安主賜,勇往直前

面對今日的香港社會,曾有人寫了下列的感言:

《十年就是目前(第三版本)》

一國兩制,互相哄騙,

說的是片刻承諾,

聽的是愛恨交纏,

裝睡的不能叫醒,

醒來的沒法入眠,

港人治港非自治,

爭取普選進寒天.

舞照跳,馬照跑,

紙醉金迷到白天,

富有的離岸持股,

耕種的照樣沒田,

空置的豪宅依舊,

tong 房的住客又添,

城市燈光耀目,

天橋底下辛酸.

失蹤的笑謂自由,

自由的戴上鎖鏈,

報章琳瑯滿目,

編輯部門翻天,

軍隊隱身城內,

警察就在身邊,

人人受到保護,

恐懼不斷蔓延.

人生一瞬,

轉眼十年,

秉行公義,

去惡求善,

願主使用,

活我今天.

當我讀完此感言,也有點無奈,但作者有句話,令我深刻:「人生一瞬,轉眼十年,…,願主使用,活我今天.」

有位朋友看完醫生後,哭著打電話給太太.太太問:「你怎樣回事?」他說:「醫生給我藥吃,今後每天一定要吃一粒!」太太問:「那你為何那麼難過?」他說:「醫生只給我四粒!」假如你是那人,只有四天可活,你們會如何過你的日子?當我們過了某個歲數,當我們去殯儀館多年了,告別一個又一個的好朋友,你感到生命的有限.

在有限的生命裡,通常有兩種表現出現.第一個是把握有限的生命,完成自己的使命或有意義的事情.有記者訪問一位前鋒球員說:「如果你走近龍門,想把球踼入時,有天使突然出現對你說:『你一入球,就會心臟病發死亡,你會怎樣做?』」那球員會如何回答呢?他說:「我會把球踼入,死就死吧,因這是我的使命.」這人就好像耶穌,他勇敢面向十字架,準備犧牲生命,完成救恩的工作；這是生命最有價值的事.

第二個是把握有限的生命,享受生命,做自己喜歡的事.讀過印度一個具人生哲理的小故事.一個人被老虎追,一見古井就跳下去,被掛在井內樹枝上.往上看,有老虎在井口；往下看,有一些毒蛇在井底.他想今次死定了,無意中卻發現有些可愛的水果,肚子亦餓,吃完才死吧.

人生張力就是完成使命還是享受生命.人生最重要的追求是甚麼呢?錢?成功?安全感?答案是安全感.在耶穌的臨終贈言,耶穌給了門徒給七樣東西,歸納而言就是安全感.第一需要是平安感,不憂不懼；第二需要是同在感,不孤單；第三需要去愛與被愛感,不冷感

第四需要喜樂,不憂慮；第五需要知悉前路,不慌張；第六需要盼望和機會,不失望；第七需要是勇敢,不怕站起,迎向敵人.

今日不少人陷入高危一族,我敢相信這類族人在上列需要上出現問題.我們不妨自我測試,自己問自己上列的問題.

你或許似曾聽過,某師傅快走了.臨走前,他約了全體徒弟相聚,告知他要走了.邪派魔頭快到,將會作王,師傅留下秘笈及護身符.但跟他們走一段路後,師傅要走去見師傅的師公(阿爸父),耶穌師傅鼓勵他們起來,勇敢向前!你是門徒,師傅將離世,留下七項秘笈,將與邪魔交手,你會害怕嗎?

主賜平安,勇往直前!許多時候,主耶穌會對我們說:「起來,我們走吧!」Get up, and let us go! (留意是一起走!)



\newpage

\section{監獄風雲}
\label{sec:320}
\textbf{林豪恩先生 - 監獄風雲}
\newline
\newline
\hyperref[sec:319]{< < < PREV SERMON < < <}
~
\hyperref[sec:toc]{[返目錄]}
~
\hyperref[sec:321]{> > > NEXT SERMON > > >}


使徒行傳 16:16-39
\newline
\begin{longtable}{cl}
\hline
\hline
章節 & 經文 (和合本修訂版)\\
\hline
16:16 & \begin{tabularx}{0.7\textwidth}{X} 後來，我們往那禱告的地方去時，有一個被占卜的靈附身的使女迎面走來，她使用法術使她的主人們發了大財。 \end{tabularx} \\ \\ \relax
16:17 & \begin{tabularx}{0.7\textwidth}{X} 她跟隨保羅和我們，喊著說：「這些人是至高神的僕人，對你們傳講救人的道路。」 \end{tabularx} \\ \\ \relax
16:18 & \begin{tabularx}{0.7\textwidth}{X} 她一連好幾天這樣喊叫，保羅就心中厭煩，轉身對那靈說：「我奉耶穌基督的名吩咐你從她身上出來！」那靈立刻出來了。 \end{tabularx} \\ \\ \relax
16:19 & \begin{tabularx}{0.7\textwidth}{X} 使女的主人們見發財的指望沒有了，就揪住保羅和西拉，拉他們到市上去見官； \end{tabularx} \\ \\ \relax
16:20 & \begin{tabularx}{0.7\textwidth}{X} 又帶他們到行政官長們面前，說：「這些騷擾我們城的，他們是猶太人， \end{tabularx} \\ \\ \relax
16:21 & \begin{tabularx}{0.7\textwidth}{X} 竟傳佈我們羅馬人所不可接受、不可遵守的規矩。」 \end{tabularx} \\ \\ \relax
16:22 & \begin{tabularx}{0.7\textwidth}{X} 群眾就一齊起來攻擊他們。官長們吩咐撕開他們的衣裳，用棍子打； \end{tabularx} \\ \\ \relax
16:23 & \begin{tabularx}{0.7\textwidth}{X} 打了許多棍，就把他們下在監裡，囑咐獄警嚴緊看守。 \end{tabularx} \\ \\ \relax
16:24 & \begin{tabularx}{0.7\textwidth}{X} 獄警領了這樣的命令，就把他們下在內監，兩腳拴在木架上。 \end{tabularx} \\ \\ \relax
16:25 & \begin{tabularx}{0.7\textwidth}{X} 約在半夜，保羅和西拉正在禱告，唱詩讚美神，眾囚犯也側耳聽著的時候， \end{tabularx} \\ \\ \relax
16:26 & \begin{tabularx}{0.7\textwidth}{X} 忽然，地大震動，甚至監牢的地基都搖動了，監門立刻全開，眾囚犯的鎖鏈也都解開了。 \end{tabularx} \\ \\ \relax
16:27 & \begin{tabularx}{0.7\textwidth}{X} 獄警一醒，看見監門全開，以為囚犯已經逃走，就拔刀要自殺。 \end{tabularx} \\ \\ \relax
16:28 & \begin{tabularx}{0.7\textwidth}{X} 保羅大聲呼叫：「不要傷害自己！我們都在這裡。」 \end{tabularx} \\ \\ \relax
16:29 & \begin{tabularx}{0.7\textwidth}{X} 獄警叫人拿燈來，就衝進去，戰戰兢兢地俯伏在保羅和西拉面前。 \end{tabularx} \\ \\ \relax
16:30 & \begin{tabularx}{0.7\textwidth}{X} 然後獄警領他們出來，說：「二位先生，我必須做甚麼才可以得救？」 \end{tabularx} \\ \\ \relax
16:31 & \begin{tabularx}{0.7\textwidth}{X} 他們說：「當信主耶穌，你和你一家都必得救。」 \end{tabularx} \\ \\ \relax
16:32 & \begin{tabularx}{0.7\textwidth}{X} 他們就把主的道講給他和他全家的人聽。 \end{tabularx} \\ \\ \relax
16:33 & \begin{tabularx}{0.7\textwidth}{X} 當夜，就在那時候，獄警把他們帶去，洗他們的傷；他和他所有的家人立刻都受了洗。 \end{tabularx} \\ \\ \relax
16:34 & \begin{tabularx}{0.7\textwidth}{X} 於是獄警領他們上自己的家裡去，給他們擺上飯。他和全家的人，因為信了神，都滿心喜樂。 \end{tabularx} \\ \\ \relax
16:35 & \begin{tabularx}{0.7\textwidth}{X} 到了天亮，官長們打發差役來，說：「釋放那兩個人吧。」 \end{tabularx} \\ \\ \relax
16:36 & \begin{tabularx}{0.7\textwidth}{X} 獄警就把這些話告訴保羅：「官長們打發人來，要釋放你們，現在可以出監，平平安安去吧。」 \end{tabularx} \\ \\ \relax
16:37 & \begin{tabularx}{0.7\textwidth}{X} 保羅卻說：「我們是羅馬人，並沒有定罪，他們竟在公眾面前打了我們，又把我們下在監裡；現在要私下趕我們出去嗎？這不行！叫他們自己來領我們出去吧！」 \end{tabularx} \\ \\ \relax
16:38 & \begin{tabularx}{0.7\textwidth}{X} 差役把這些話回稟官長們；官長們聽見他們是羅馬人，就害怕了， \end{tabularx} \\ \\ \relax
16:39 & \begin{tabularx}{0.7\textwidth}{X} 於是來勸他們，領他們出來，請他們離開那城。 \end{tabularx} \\ \\
[1ex]
\hline
\hline
\end{longtable}
講題:監獄風雲　Prison On Fire

經文:使徒行傳16章16-39節

講員:林豪恩先生

場合:香港中文大學崇基學院禮拜堂

日期:2016年5月8日

一．引言

今天的新約經課－使徒行傳16章16-34節當中有一句耳熟能詳的經文,有時在街上也可看見,因為有些教會把這「金句」掛於教堂外牆,向公眾展示．這句被教會認為足以展示基督教信仰的經文記載在16章31節:「當信主耶穌,你和你一家都必得救.」這「金句」出現在甚麼事件之中的呢?除了「你和你一家都必得救」,這事件還能帶給我們甚麼盼望呢?

二．經文

16:16後來,我們往那禱告的地方去.有一個使女迎著面來,她被巫鬼所附,用法術,叫她主人們大得財利.

16:17她跟隨保羅和我們,喊著說:「這些人是至高上帝的僕人,對你們傳說救人的道.」

16:18她一連多日這樣喊叫,保羅就心中厭煩,轉身對那鬼說:「我奉耶穌基督的名,吩咐你從她身上出來!」那鬼當時就出來了.

16:19使女的主人們見得利的指望沒有了,便揪住保羅和西拉,拉他們到市上去見首領；

16:20又帶到官長面前說:「這些人原是猶太人,竟騷擾我們的城,

16:21傳我們羅馬人所不可受不可行的規矩.」

16:22眾人就一同起來攻擊他們.官長吩咐剝了他們的衣裳,用棍打；

16:23打了許多棍,便將他們下在監裡,囑咐禁卒嚴緊看守.

16:24禁卒領了這樣的命,就把他們下在內監裡,兩腳上了木狗.

16:25約在半夜,保羅和西拉禱告,唱詩讚美上帝,眾囚犯也側耳而聽.

16:26忽然,地大震動,甚至監牢的地基都搖動了,監門立刻全開,眾囚犯的鎖鍊也都鬆開了.

16:27禁卒一醒,看見監門全開,以為囚犯已經逃走,就拔刀要自殺.

16:28保羅大聲呼叫說:「不要傷害自己!我們都在這裡.」

16:29禁卒叫人拿燈來,就跳進去,戰戰兢兢地俯伏在保羅,西拉面前；

16:30又領他們出來,說:「二位先生,我當怎樣行才可以得救?」

16:31他們說:「當信主耶穌,你和你一家都必得救.」

16:32他們就把主的道講給他和他全家的人聽.

16:33當夜,就在那時候,禁卒把他們帶去,洗他們的傷；他和屬乎他的人立時都受了洗.

16:34於是禁卒領他們上自己家裡去,給他們擺上飯.他和全家,因為信了上帝,都很喜樂.

16:35到了天亮,官長打發差役來,說:「釋放那兩個人吧.」

16:36禁卒就把這話告訴保羅說:「官長打發人來叫釋放你們,如今可以出監,平平安安地去吧.」

16:37保羅卻說:「我們是羅馬人,並沒有定罪,他們就在眾人面前打了我們,又把我們下在監裡,現在要私下攆我們出去嗎?這是不行的.叫他們自己來領我們出去吧!」

16:38差役把這話回稟官長.官長聽見他們是羅馬人,就害怕了,

16:39於是來勸他們,領他們出來,請他們離開那城.

三．釋經與應用

1．獄中歌唱的盼望

你有沒有被關進監獄的經驗?我相信你絕對不希望有．三月的其中一次讀書會,有一位朋友來分享他在大陸被關進監獄的經歷:

「我被戴上手扣,在兩名管教的押解下,被推進一間囚室．進門後,厚重的鋼門砰然一響地關上,此一刻,我的心碎了!我整個人愣住了!呆了!整個人精神處於崩潰邊緣,好一會都無法清醒過來,無法接受這個事實．羞愧與屈辱,挫拆與自卑．從此我就被終身打上一個「罪犯」的烙印,想到這個烙印我就無地自容,討厭自己．我喪失了自我,迷失了自我．」(程翔:千日無悔)

使徒行傳16:22-24節記載了耶穌基督的使徒保羅和他的宣教同伴被關進監獄的情況:

16:22眾人就一同起來攻擊他們.官長吩咐剝了他們的衣裳,用棍打；

16:23打了許多棍,便將他們下在監裡,囑咐禁卒嚴緊看守.

16:24禁卒領了這樣的命,就把他們下在內監裡,兩腳上了木狗.

如果你是保羅和西拉,你會否感到驚惶,憤怒,和疑惑嗎?我們沒有放火,沒有殺人,到底犯了甚麼大罪,被公開羞辱和毒打?被視為重犯關在監獄的內室?還要在兩腳鎖上木狗呢?我們到底做錯了甚麼,以致落得如此糟糕的下場呢?

那位來和我們分享的朋友憶述:

「我質問自己,我的人生道路為甚麼會走到監獄的死胡同?是命運使然,抑或是我錯誤的選擇?如果是命運,上天為甚麼要薄待我?如果是自己錯誤的選擇,實在太過冤枉了,因為我自問一生的所有重大決定,從來都是把國家的利益放在第一位來考慮的．」(程翔:千日無悔)

他以陷入「死蔭的幽谷」來形容自己當時的狀況．在「死蔭的幽谷」中,他感到死亡擦身而過,陷入精神崩潰的邊緣．

希望你沒有這種被關進監獄的經驗．不過,人生中突如其來的打擊卻足以令人陷入「死蔭幽谷」．如那位突然被捕的朋友一樣,在「死蔭幽谷」的人都禁不住問為甚麼?為甚麼是我?然而,有些時候,根不找不到合理解釋；或者,就算合理解釋也不能救人離開困局．使徒行傳記載了保羅和西拉被關進監獄之前所發生的事情:

16:16後來,我們往那禱告的地方去.有一個使女迎著面來,她被巫鬼所附,用法術,叫她主人們大得財利.

16:17她跟隨保羅和我們,喊著說:「這些人是至高上帝的僕人,對你們傳說救人的道.」

16:18她一連多日這樣喊叫,保羅就心中厭煩,轉身對那鬼說:「我奉耶穌基督的名,吩咐你從她身上出來!」那鬼當時就出來了.

16:19使女的主人們見得利的指望沒有了,便揪住保羅和西拉,拉他們到市上去見首領；

16:20又帶到官長面前說:「這些人原是猶太人,竟騷擾我們的城,

16:21傳我們羅馬人所不可受不可行的規矩.」

按那位被巫鬼附身的使女說,保羅所做的是「傳說救人的道」；對那位使女的主人來說,保羅解救了使女是「斷了他們的財路」．本質上,這是私人恩怨,

對香港人來說:「這是錢可以解決的問題」．然而,某些人另有用心,就把這「搵食」的問題上綱上線變成公共安全問題,變成威脅羅馬管治的政治問題來挑動羅馬官員的神經．身陷牢獄,不一定是罪有應得的．墜入「死蔭幽谷」,也不一定找到合理原因．挫敗,疑惑,驚惶,憤怒,埋怨,羞辱,無奈……都是可理解的反應．除了怨天尤人,坐以待斃之外,還可以怎樣呢?保羅和西拉被投入黑獄這事件,啟發我們可以有獄中歌唱的盼望．

16:25約在半夜,保羅和西拉禱告,唱詩讚美上帝,眾囚犯也側耳而聽.

這種獄中歌唱的盼望,並不只是出現在保羅和西拉身上,也出現在那位來和我們分享的朋友的生命中．

「我在看守所期間,就已經開始禱告．當時為甚麼會禱告起來?我想,在那個完全無助的狀態下,是痛苦,無奈和求助的心態促使我走向主．」

「我的第一次禱告是這樣說的:主啊!求你幫助我,使我能夠堅定走我的愛國度路,為中國的改革開放,民主自由,和平統一而奮鬥．」

「在監獄期間,我繼續通過禱告去戰勝困難,重新出發．我的每一天,都是從這個禱告開始:感謝主賜我一顆喜樂的心,使我時刻能夠:以微笑面對痛苦,以禮貌面對冷酷,以寬恕面對迫害,以謙卑面對狂妄,以關懷面對冷漠,以激勵面對踐踏．特別是即使身處逆境,仍能常存感恩,惜福,隨緣,知足的心．」(程翔:千日無悔)

你感到自己處身於怎樣的光景呢?有人在獄中仍向上帝祈禱,唱歌讚美上帝；有人在「死蔭幽谷」仍不怕遭害；有人在黑暗之中仍相信:「黑暗也不能遮蔽我使你不見,黑夜卻如白晝發亮,黑暗和光明在你看都是一樣.」(詩139:12)

有人說,未來十年可能烏雲蓋頂．今天經課記載的事件啟發我們,就算身陷牢獄,我們仍可以有獄中歌唱的盼望．

2．化敵對為友伴的盼望

16:23打了許多棍,便將他們下在監裡,囑咐禁卒嚴緊看守.

16:24禁卒領了這樣的命,就把他們下在內監裡,兩腳上了木狗.

獄警和囚犯是一種怎樣的關係呢?保羅和西拉所犯的事是否嚴重到要囚在內監裡?兩腳上木狗會否傷害被囚者及為他們帶來痛楚?如果問禁卒,他可能回答說:「我只是執行命令．」

那位朋友說:「我第一次領會到漢字「獄」字為甚麼是從爪,從言,從犬,就是因為監獄生活把「人」貶為一隻「會說話」的「狗」．」「有一次我看到獄警召集各囚室的頭犯開會,他坐在大班椅上,六七個犯人蹲在他前面,他用皮鞋尖逐一對著每個人,吩咐每個人要做的事．又有一次我看到獄警召見一個爺爺輩的囚犯,爺爺輩在孫子輩的面前蹲下,孫子輩的用皮鞋在他面前打劃著說話．」(程翔:千日無悔)

獄警和囚友必然敵對的嗎?此時此地的市民,親身經歷,或者在傳媒上看見金鐘,旺角,銅鑼灣街頭的那79 天,以及農曆新年旺角那一夜,不禁對執法者和示威者對抗的場面感到震驚．執法者和示威者都是這個城市的一份子,甚至可能是鄰居,同學或朋友,是甚麼令他們變成敵對呢?是誰人令他們變成敵對呢?這是不少人心中的疑惑,也引起很多人內心不安．

菲利普‧津巴多(Philip Zimbardo)教授在1971年做過一項出名的“史丹佛監獄實驗(Stanford Prison Experiment )”．實驗發現,人們會被角色支配,並按角色來行動．扮演守衛和囚犯的本來都是正常人.然而,在囚犯和守衛對抗的形勢之下,僅僅幾天的時間,守衛就變成了暴徒,而犯人則完全精神崩潰了．

佔中之後,有一位警務人員和一位示威者應邀對談,那位警員說:「警察身為執法者,by nature 要維繫社會制度穩定.如果有人要衝擊這個制度,其實必然會跟警察正面衝突,無可避免.」(當警察遇上示威者 - 陳sir VS. 葉寶琳(立場新聞2015/9/2))執法者和被執法者之間確實有角色的差異,除了對立之外,還有其他可能性嗎?或者說,對立是我們之間唯一的關係嗎?

16:26忽然,地大震動,甚至監牢的地基都搖動了,監門立刻全開,眾囚犯的鎖鍊也都鬆開了.

16:27禁卒一醒,看見監門全開,以為囚犯已經逃走,就拔刀要自殺.

16:28保羅大聲呼叫說:「不要傷害自己!我們都在這裡.」

16:29禁卒叫人拿燈來,就跳進去,戰戰兢兢地俯伏在保羅,西拉面前；

16:30又領他們出來,說:「二位先生,我當怎樣行才可以得救?」

16:31他們說:「當信主耶穌,你和你一家都必得救.」

16:32他們就把主的道講給他和他全家的人聽.

16:33當夜,就在那時候,禁卒把他們帶去,洗他們的傷；他和屬乎他的人立時都受了洗.

16:34於是禁卒領他們上自己家裡去,給他們擺上飯.他和全家,因為信了上帝,都很喜樂.

作者花了相當的篇幅去敘述保羅與禁卒的互動．在沒有改變角色,即是禁卒仍是禁卒,囚犯仍是囚犯的身份之下,可以創造另類的互動狀況和關係質數．他們打破了禁卒與囚犯兩者角色的定型,超越了對抗性的互動,出現了憂戚與共,患難相助的關係．囚犯的福祉依賴禁卒,禁卒的福祉,也受囚犯影響； 16章33節展示,禁卒為囚犯清洗傷口,囚犯為禁卒施行洗禮．原來禁卒和囚犯兩者是相依共存,可以互相幫助的．

禁卒不能改變囚犯的身份,犯囚也不能要求禁卒反守則偷偷放犯囚走．禁囚和犯囚在甚麼基礎之上發展另類互動呢?很可能是因為保羅看對方超越了角色,而看到對方本身是一個人,因此以人性來對待對方．這種人性的對待也喚起了禁卒的人性,視對方為一個人,而不單單是一個囚犯,去照顧作為一個人的基本需要．在甚麼基礎之上幫助雙方視對方為一個人呢?不是角色,如果是角色,當滿足不到角色的要求,禁卒就會「拔刀要自殺」(v27)；也不是權威,如果是權威,禁卒就會「戰戰兢兢地俯伏在保羅,西拉面前」．那基礎是耶穌基督的福音,在上帝裡面,他們能夠平等相待,化焦慮為喜樂(v34)．

其實,執法者和被執法者都是人．執法者的人性,有時只是被制度扭曲和掩蓋而已．程翔曾談辦理他的案件的執法人員時如此透露:

「經過幾次安排我家人來會見後,我的兄弟妹們與大哥口中的國安局Ｘ先生變得熟絡起來．有一次,妹妹瑞儀就大膽地試問他:「你接觸我二哥這麼長時間,又審訊了他,你覺得他是一個怎樣的人?」聽妹妹說,這位Ｘ先生很認真地思考了一下,然後正色地說:「他是一個好人!」我妹妹認為他回答時很認真的,一點都不是敷衍或客套．他沒有必要討好我,事實上他的紀律使他連名字都不能給我家人,所以即使討好了也沒有甚麼用．他這樣回答,肯定是發自內心．我聽了這句話,心中很感慨,也很安慰．感慨的是:一個連審訊者都認為是「好人」人,審訊者卻不得不設法羅織罪名給他．安慰的是,一個人真正的道德品格是裝不出來的,也是掩蓋不了的．即使在遭遇困難險阻的時候,人是被投入黑獄裡了,但可幸我真正的本性仍然呈現出來,而呈現出來的本性,足以使一個奉命加害我的人也承認他要處理的「疑犯」其實是一個好人!

只要人性不滅,化敵對為友伴的盼望仍然可期待．只要那位創造人性的上帝仍在,人性不滅仍然是可期待的盼望．

3．冤假錯案得平反的盼望

16:35到了天亮,官長打發差役來,說:「釋放那兩個人吧.」

16:36禁卒就把這話告訴保羅說:「官長打發人來叫釋放你們,如今可以出監,平平安安地去吧.」

出獄當然重要,然而,如何出獄,可能更重要．

【20160104on.cc東網專訊】 巴西伯南布哥州昨午發生集體越獄．報道指,當時正值家屬探監時間,兩人把一輛汽車停在監獄外,其後引爆炸藥,炸開監獄外牆,逾100名囚犯乘機逃獄.包括今次,伯南布哥州在一周內已發生兩宗越獄,日前有53名囚犯挖地道逃獄.

半夜監門打開的時候,保羅他們不願背負逃犯之名而逃走．早上獲得釋放,他們竟然還是不走．保羅要求甚麼呢?有甚麼比被釋放出獄更重要呢?

16:37保羅卻說:「我們是羅馬人,並沒有定罪,他們就在眾人面前打了我們,又把我們下在監裡,現在要私下攆我們出去嗎?這是不行的.叫他們自己來領我們出去吧!」

16:38差役把這話回稟官長.官長聽見他們是羅馬人,就害怕了,

16:39於是來勸他們,領他們出來,請他們離開那城.

有過來人士表示,如果監獄對犯人從身體上和精神上深深打下「罪犯」,「監躉」的烙印,長遠來說,會對罪犯構成心理上和精神上永久性的傷害,妨礙他們在出獄後融入社會,而且這永久性的傷害也是有些人再犯罪的因素．(程翔:千日無悔)

保羅他們被釋放,但他們更需要的是平反．社會不接受把有罪的當作無罪,更不應該把無罪的當作有罪．在這事件中,保羅他們得回公道,尊嚴獲得歸還．最近,有一件在1989年發生的悲劇終於尋冤得雪:

調查1989年英國希斯堡球場人踩人慘劇的死因陪審團,2016年4月26日裁定警方當年安排出錯及疏忽導致慘劇,96名喪生的利物浦球迷是「非法被殺」.英國首相卡梅倫稱,裁決還了死者一個遲來的公道,意義重大.涉事南約克郡警隊現任主管就警方當年的「災難性錯誤」作「毫無保留道歉」.事發於1989年4月15日,利物浦與諾定咸森林在南約克郡設菲爾德市的希斯堡球場(Hillsborough Stadium)進行英格蘭足總盃準決賽,吸引數萬計球迷到場.為疏導人潮,賽事的警方指揮官達肯菲爾德(David Duckenfield)下令打開一個出口閘門,容許大批利物浦球迷湧入看台,結果觸發人踩人,造成96人死600多人傷.事後球迷被指生事導致慘劇.2009年慘劇發生20周年,當局成立獨立小組調查慘劇起因.2012年發表報告,指警方涉嫌隱瞞真相和篡改文件,務求諉過於無辜球迷.案件其後重審.陪審團自2014年開始聽取約1000名證人作證,經過兩年審訊,裁定球迷在事發當日的行為並沒有釀成慘劇.警方指揮官達肯菲爾德去年出席聆訊時承認自己當年未能封閉球場一條隧道是造成人命傷亡的直接原因.(明報2016年4月27日)

今天是母親節,有一群母親二十多年來一直盼望,盼望那件也在1989年發生的慘劇尋冤得雪,盼望他們在生的以及已經離世的家人的污名得到平反,他們的盼望有實現的一天嗎?世界上仍然有不少人和他們一齊盼望著這一天的來臨．

四．總結

今天經課講述保羅和他的宣教隊伍的短暫牢獄經歷,啟發我們對盼望的眾多想像．除了「當信主耶穌,你和你一家都必得救」之外,還有「獄中歌唱的盼望」,「化敵對為友伴的盼望」,以及「冤假錯案得平反的盼望」．



\newpage

\section{誰是我的父親}
\label{sec:321}
\textbf{鄧瑞強博士 - 誰是我的父親?}
\newline
\newline
\hyperref[sec:320]{< < < PREV SERMON < < <}
~
\hyperref[sec:toc]{[返目錄]}
~
\hyperref[sec:322]{> > > NEXT SERMON > > >}


羅馬書 8:12-17
\newline
\begin{longtable}{cl}
\hline
\hline
章節 & 經文 (和合本修訂版)\\
\hline
8:12 & \begin{tabularx}{0.7\textwidth}{X} 弟兄們，這樣看來，我們不是欠肉體的債去順從肉體而活。 \end{tabularx} \\ \\ \relax
8:13 & \begin{tabularx}{0.7\textwidth}{X} 你們若順從肉體活著，必定會死；若靠著聖靈把身體的惡行處死，就必存活。 \end{tabularx} \\ \\ \relax
8:14 & \begin{tabularx}{0.7\textwidth}{X} 因為凡被神的靈引導的都是神的兒子。 \end{tabularx} \\ \\ \relax
8:15 & \begin{tabularx}{0.7\textwidth}{X} 你們所領受的不是奴僕的靈，仍舊害怕；所領受的是兒子名分的靈，因此我們呼叫：「阿爸，父！」 \end{tabularx} \\ \\ \relax
8:16 & \begin{tabularx}{0.7\textwidth}{X} 聖靈自己與我們的靈一同見證我們是神的兒女。 \end{tabularx} \\ \\ \relax
8:17 & \begin{tabularx}{0.7\textwidth}{X} 若是兒女，就是後嗣，是神的後嗣，和基督同作後嗣。如果我們和他一同受苦，是要我們和他一同得榮耀。 \end{tabularx} \\ \\
[1ex]
\hline
\hline
\end{longtable}
講題:「誰是我的父親?」　Who Is My Father?

經文:羅馬書8:12-17

講員:鄧瑞強博士

場合:香港中文大學崇基學院禮拜堂

日期:2016年5月15日(聖靈降臨節)

各位弟兄姊妹,早安.

「爸爸」這種人物是很奇特的,你在不同年紀看他,會有不同感覺.你讀小學的時候,會覺得爸爸無所不知,無所不能.他竟然知道地球是圓的,又懂得弄馬桶.你讀中學的時候,會覺得他「識少少,扮代表」.你讀大學的時候,忽然覺得他簡直無知.然後,你一直長大,不知到了什麼年紀,有一天,你忽然覺得爸爸是如此智慧,如此高明.

馬克吐溫講了一段很出名的話,他說:「當我十四歲時,我的父親無知至極,我不能忍受這老人家在我周圍.但當我二十一歲時,我驚訝這老人家在這七年來竟學到了如此多的東西.」( “When I was a boy of 14, my father was so ignorant I could hardly stand to have the old man around. But when I got to be 21, I was astonished at how much the old man had learned in seven years.”- Mark Twain )

在西方,爸爸的命運頗算悲慘.聽心理學家弗洛伊德(Freud)說,西方男子大抵有一「殺父戀母情結」(Oedipus Complex).在西方,一個男孩子漸漸成長,要建立一獨立的個體人格.在這時,他會發覺父親正是在他之上而限制他成長的人,為了發展個體人格,西方男子有一去除父親影響力的潛在心理,這就是所謂「殺父情結」.對西方男子而言,個性獨立的要求,大於家庭紐帶的力量.

要求個體人格的獨立自主,已是個人主義社會的普遍現象,現在已無分東方西方了.在我們的社會,當年輕人拿起一塊磚頭擲向警察時,「殺父情結」長出的苗已清楚可見.所謂父親,所謂權威,已被看成為壓制個人自主和自由的強權,與之衝突,變成無可避免.

靈修學者盧雲(Nouwen)說,這是「無父的一代」.年輕人越來越抗拒父親,越來越抗拒所謂權威,越來越抗拒既定的體制.他們不聽大人們講的理想,因為大人們本身也沒有理想.他們不想融入大人們所設計的社會裡,因為就是這些大人們將社會搞到一塌糊塗.他們不再珍惜大人們的嘉許,因為大人們根本不明白他們的期望.盧雲說:「他們覺得一切都不對勁,可惜找不著可行的取代方法.所以,他們沮喪失意,通常的發洩方法是漫無目的之暴力,只有破壞卻動機不明；他們或是自暴自棄,與世隔絕.」(《負傷的治療者》,中譯本,頁38)

在美國,很多社會學者稱美國的社會為「無父的社會」,這裡有另一類的「無父的一代」.在這裡,因著種種原因,如:離婚等,造成很多無父的家庭.按某些統計,有43\%的美國孩童沒有父親在身邊,這造成很多問題.90\%離家出走的孩童,71\%的青少年懷孕,63\%的青少年自殺,85\%行為失常的兒童….等等,皆來自「無父的家庭」.(參:\url{http://fatherhoodfactor.com/us-fatherless-statistics/})

有些人反抗父親,有些人被父親離棄,都成了無父的人.這些人都是孤獨的.孤獨地摸索前路的人,多麼渴望有父親在身邊提著燈,幫助他們看明前路.

其實,我們每個人都需要父親,但不是那種強橫,暴力,剛愎自用的父親.我們需要父親去指示生命真正的高度,我們需要父親去幫助我們揭開存在的奧秘,我們需要父親成為我們背後的鼓舞力量.

問題是:誰是我們的父親?這父親在哪裡?

無父的人在徬徨,在尋覓,生命的靈吹拂大地,讓人體會深藏於天地間的父性.這是天父世界,這世界滲透著天父的大仁大愛.仁慈的天父,正是我們眾人的父親.有了這父親,生命才有真正的靠山.

我們看看今日的講道經文:羅馬書8:12－17

羅8:12　弟兄們,這樣看來,我們並不是欠肉體的債去順從肉體活著.

羅8:13　你們若順從肉體活著,必要死；若靠著聖靈治死身體的惡行,必要活著.

羅8:14　因為凡被神的靈引導的,都是神的兒子.

羅8:15　你們所受的,不是奴僕的心,仍舊害怕；所受的,乃是兒子的心,因此我們呼叫:「阿爸!父!」

羅8:16　聖靈與我們的心同證我們是神的兒女；

羅8:17　既是兒女,便是後嗣,就是神的後嗣,和基督同作後嗣.如果我們和他一同受苦,也必和他一同得榮耀.

這段經文一開始,講到生命的張力.

羅8:12　弟兄們,這樣看來,我們並不是欠肉體的債去順從肉體活著.

羅8:13　你們若順從肉體活著,必要死；若靠著聖靈治死身體的惡行,必要活著.

羅馬書的作者保羅說,生命有所待,有所欠.我們好像等著什麼似的,我們好像被什麼牽引著的.保羅說,我們不是欠了肉體一些什麼,我們沒有必要跟著肉體的指揮棒而活.所謂「肉體」,在保羅神學裡,其實是指人性的脆弱,這脆弱源自人具有必朽壞的肉體,或者說,伴隨人的必死性而有的脆弱.這種脆弱令人尋求自保,這自保的心催迫人去滿足肉體的慾望和需要.這自保的心過強,便會令人成為自我的奴隸.一個自我的奴隸,只看到自己,只想到自己,只感受到自己,只為自己而活.在「我」之外,再看不到什麼.這個「超大的我」,沒有他人,沒有上帝.這是一個拒絕一切父親的孤兒.他是一個永遠浪跡天涯的浪子,他不會回家,因為他看不到家.他不會回歸父親的懷抱,因為他不承認任何父親.保羅說,我們沒有欠這個「超大的我」什麼,我們不用跟著他走.保羅說,「你們若順從肉體活著,必要死」.一種無父,無權威,無他人的生命,在保羅看來,只是一條死路.

保羅講另一種生命,他說:「若靠著聖靈治死身體的惡行,必要活著.」聖靈,天地間促發生命的靈氣.這是吹入沒有氧氣的密室的生氣.聖靈引導我們離開死路,而接上生命的源頭.這有點似戒毒,要將一個人從被控制的生命中拉回來,他好像經歷某種死亡,然後,他才能開展新生.

保羅繼續說:

羅8:14　因為凡被神的靈引導的,都是神的兒子.

羅8:15　你們所受的,不是奴僕的心,仍舊害怕；所受的,乃是兒子的心,因此我們呼叫:「阿爸!父!」

上天有眼,他看見人努力展示生命的內容,看見人拒絕父親和權威,看見人被父親離棄,看見人一個又一個離家遠行,看見人用自己的腳踏出不同的路,看見人走入荒原,看見人因拒絕任何權威或失去權威而不知方向,看見人在孤寂中不斷消耗自己的生命,看見人在死亡的面前恐懼,看見人因自保而越發孤寂,看見人瘋狂,看見人在絕望中死去.上天有眼,他看見人間的一切.上天有情,他不願人間如此.

天上的靈風吹遍大地,生命的靈四處去觸動人的心靈.來吧,走出孤寂.來吧,不要再這麼自我.來吧,對這人間多一點信任.靈風吹遍,讓人體會到一種生命的力量.這種生命的力量,源自天地間最大的一顆愛心.在宇宙間,有一天父,他不忍我們在無父的迷惘和孤寂中死去.在聖靈裡體會這天父的心的人,只要他願意與這天父聯繫,他就成為神的兒女.

這是在生命之靈的引導下的一種新的自覺,原來,就算人間再沒有父親,天地間仍有天父,我們的生命仍有所屬,我們仍可以有回家的渴望.這位天父,仁愛為懷.他不會以大欺小,在基督裡,我們看到他甘願俯就卑微.他不會為了自保而欺壓我們,在基督裡,我們看到他甘願捨棄自己而成全我們.他不會因我們不聽話而離棄我們,在基督裡,我們看到他在我們還作罪人的時候千方百計地拯救我們.縱使人間的父親離棄我們,從使我們反抗人間的父親,仍有一天上的父,他溫柔地看著我們,指引著我們.

保羅說,「凡被神的靈引導的,都是神的兒子.你們所受的,不是奴僕的心,仍舊害怕；所受的,乃是兒子的心,因此我們呼叫:阿爸!父!」這是一種新的生命體會.人生在世,不是無父,而是有父.我們是天父的兒女,不是不自由的奴僕.不是世界各種勢力的奴僕,也不是天父之下無地位的奴僕,而是天父所收納的兒女.奴僕,沒有自由,沒有前景,不能得到任何東西,有的只是恐懼的心.但身為兒子,有身分,有尊嚴,將來要承受父親的產業.在這「無父的一代」,基督徒要成為「有父的一群」.這是一新的人生.

有了新的領悟,人生就有新的責任,保羅說:

羅8:16　聖靈與我們的心同證我們是神的兒女；

羅8:17　既是兒女,便是後嗣,就是神的後嗣,和基督同作後嗣.如果我們和他一同受苦,也必和他一同得榮耀.

我們的心,與生命的靈拍和,就是「聖靈與我們的心同證我們是神的兒女」.經文說,神的兒女,就是神的後嗣,是和基督同作後嗣.基督是誰,基督是天父的心在人間的彰顯.基督彰顯天父的父性,彰顯這父性的大愛,彰顯這大愛的深度.十字架的大愛,就是天父的心.

天地間的生命之靈,感化我們去體會天父的心.天父的心,落實在基督的生命裡.我們若認天父為父,若接受天父的兒女的身分,則我們是過一種新的生活,領受一種新的責任.這種生活,就是像基督一樣生活.這種責任,就是承擔起基督所承擔的責任.

「無父的一代」,拒絕一切外加的責任,視一切外加的責任為枷鎖.聖靈感化人成為天父的兒女,就是讓人從心裡感應天父的心,從而甘心承擔天父的愛的責任.沒有這種人性的轉化,人很難不按自我的私慾而行,也很難像基督一樣背負愛的十架.

基督的受苦,其實展示天父的父性.基督以十架的苦,十架的愛,展示真正的父性.這種父性,其權威的建立,不是以暴力,而是以甘心犧牲的大愛來建立,以愛來感召人順服.這種父性,是明白人的反叛,而在人反叛時,已預備了回家的路.這種父性,展示了人性應有的高度,這高度的極致,就是以生命活出天父的心.

生命之靈感化我們,使我們成為神的兒女,展示像基督一樣的生命,這生命就是為父的生命,展露天父一樣的心靈.在這「無父的一代」,多麼需要有人甘願為父.

為父,是艱巨的工程,要背起十架.在這裡,保羅鼓勵我們,「如果我們和他一同受苦,也必和他一同得榮耀」.為父者,像復活的基督,當完成生命的工程時,能回到天父那裡去.這是我們的盼望.

網上有篇文章:“What It’s Like To Be A Fatherless Daughter”(一個無父的女兒是什麼一回事),文章說:「在我經驗裡,無父的女兒是最孤寂的生命.超乎想像的哀愁,空洞,破碎.總覺得有些東西從我們生命中遺失了.只感到十分寂寞,十分不安全.只感到不值得被愛.」

(參:\url{http://thoughtcatalog.com/shannon-bendell/2015/03/what-its-like-to-be-a-fatherless-daughter/})

或許,我們就是這樣的一個人,願生命的靈帶領我們,去認識天上的父.或許,我們已體會自己是天父的兒女,這時,讓我們效法基督的愛,去成為他人的父親.這個「無父的一代」,需要更多真正的父親.

但願:榮耀歸於聖父,聖子,聖靈.阿們.



\newpage

\section{與神摔跤破碎然而蒙福}
\label{sec:322}
\textbf{朱建忠宣教師 - 與神摔跤:破碎然而蒙福!}
\newline
\newline
\hyperref[sec:321]{< < < PREV SERMON < < <}
~
\hyperref[sec:toc]{[返目錄]}
~
\hyperref[sec:323]{> > > NEXT SERMON > > >}


創世記 32:22-32
\newline
\begin{longtable}{cl}
\hline
\hline
章節 & 經文 (和合本修訂版)\\
\hline
32:22 & \begin{tabularx}{0.7\textwidth}{X} 他夜間起來，帶著兩個妻子，兩個婢女和十一個孩子，過了雅博渡口。 \end{tabularx} \\ \\ \relax
32:23 & \begin{tabularx}{0.7\textwidth}{X} 他帶著他們，送他們過河，他所有的一切也都過去， \end{tabularx} \\ \\ \relax
32:24 & \begin{tabularx}{0.7\textwidth}{X} 只剩下雅各一人。有一個人來和他摔跤，直到黎明。 \end{tabularx} \\ \\ \relax
32:25 & \begin{tabularx}{0.7\textwidth}{X} 那人見自己勝不過他，就摸了他的大腿窩一下。雅各的大腿窩就在和那人摔跤的時候扭了。 \end{tabularx} \\ \\ \relax
32:26 & \begin{tabularx}{0.7\textwidth}{X} 那人說：「天快亮了，讓我走吧！」雅各說：「你不給我祝福，我就不讓你走。」 \end{tabularx} \\ \\ \relax
32:27 & \begin{tabularx}{0.7\textwidth}{X} 那人說：「你叫甚麼名字？」他說：「雅各。」 \end{tabularx} \\ \\ \relax
32:28 & \begin{tabularx}{0.7\textwidth}{X} 那人說：「你的名字不要再叫雅各，要叫以色列，因為你與神和人較力，都得勝了。」 \end{tabularx} \\ \\ \relax
32:29 & \begin{tabularx}{0.7\textwidth}{X} 雅各問他說：「請告訴我你的名字。」那人說：「何必問我的名字呢？」於是他在那裡為雅各祝福。 \end{tabularx} \\ \\ \relax
32:30 & \begin{tabularx}{0.7\textwidth}{X} 雅各就給那地方起名叫毗努伊勒 ，說：「我面對面見了神，我的性命仍得保全。」 \end{tabularx} \\ \\ \relax
32:31 & \begin{tabularx}{0.7\textwidth}{X} 太陽剛出來的時候，雅各經過毗努伊勒，他的大腿就瘸了。 \end{tabularx} \\ \\ \relax
32:32 & \begin{tabularx}{0.7\textwidth}{X} 因此，以色列人不吃大腿窩的筋，直到今日，因為那人摸了雅各大腿窩的筋。 \end{tabularx} \\ \\
[1ex]
\hline
\hline
\end{longtable}
講題:與神摔跤:破碎然而蒙福!　Wrestling with God: Broken, But Blessed!

講員:朱建忠宣教師

經文:創世記32:22-32

場合:香港中文大學崇基學院禮拜堂

日期:2016年5月22日

引言:弟兄姊妹,主內平安!今天邀請大家一起來思想,作為基督徒,在我們的屬靈歷程當中,到底什麼是我們應該努力抓取的,什麼又是我們應該放下的.我們一起來看聖經.

讀經:創32:22－32

22 他夜間起來,帶著兩個妻子,兩個使女並十一個兒子,都過了雅博渡口. 23 先打發他們過河,又打發所有的都過去, 24 只剩下雅各一人.有一個人來和他摔跤,直到黎明. 25 那人見自己勝不過他,就將他的大腿窩摸了一把,雅各的大腿窩正在摔跤的時候就扭了. 26 那人說:「天黎明瞭,容我去吧.」雅各說:「你不給我祝福,我就不容你去.」27 那人說:「你名叫什麼?」他說:「我名叫雅各.」 28 那人說:「你的名不要再叫雅各,要叫以色列,因為你與神與人較力,都得了勝.」 29 雅各問他說:「請將你的名告訴我.」那人說:「何必問我的名?」於是在那裡給雅各祝福. 30 雅各便給那地方起名叫毗努伊勒 b,意思說:「我面對面見了神,我的性命仍得保全.」 31 日頭剛出來的時候,雅各經過毗努伊勒,他的大腿就瘸了. 32 故此,以色列人不吃大腿窩的筋,直到今日,因為那人摸了雅各大腿窩的筋.

我們剛才所讀聖經是創世記第32章,22節到32節,我們的題目是與神摔跤--破碎然而蒙福.為了讓大家更好的來理解這段經文,我們需要簡單的介紹一下這段經文的背景.我們都知道我們的這位朋友雅各因為聯合他母親利百佳用計騙得父親的長子祝福之後,以掃就怨恨雅各,並且暗下決心,心裡說:「為我父親居喪的日子近了,到那時候,我要殺我的兄弟雅各.」(創21:41) 於是,雅各無奈之下,聽從母親安排,離家投奔母舅拉班住處.為了拉班的兩個女兒及羊群雅各服事了二十年後,看到母舅的氣色也大不如前時,就悄悄的起來,帶著兒子,妻子並巴旦亞蘭所得的一切牲畜和財物離開了拉班.到了第三日,有人告訴拉班,雅各逃跑了,拉班便帶了他的眾兄弟去追趕,一連七日,終於在基列山追上了.但是神插手,阻止了拉班加害於雅各.

所以當我們讀到第32章的時候,雅各是剛剛躲過凶險的一關,別過了前來殺他的母舅之後繼續前行.此時的雅各,舉目望去,故土依希就在前方,但這一路是吉是凶依然沒有定數,二十年前兄長以掃那一句:「等父親居喪的日子,我便殺了雅各.」彷彿又在耳邊想起,以至於神的使者在他的前面路過也不能使他安心.倘若兄長以掃若帶著人馬過來追討二十年前的怨恨,自己妻兒老少這麼多人又如何躲得過這一劫難.想到這裡,雅各打發人先去以東地,探探兄長以掃對自己此行的態度如何?焦急的等待一直持續到了傍晚時分,探子終於回到了雅各那裡,但報來的消息卻是:「以掃帶了四百人,正迎著你而來.」聖經在第七節,第八節記載說:雅各就甚懼怕,而且愁煩,便把與他同在的人口和羊群,牛群,駱駝分作兩隊以便在以掃擊殺這一隊的時候或許可以逃避另一隊.然後開始提醒耶和華神,(9節,11節)「耶和華我祖亞伯拉罕的神,我父親以撒的神啊,你曾對我說:『回你本地本族去,我要厚待你.』  求你救我脫離我哥哥以掃的手,因為我怕他來殺我,連妻子帶兒女一同殺了.」

但是,禱告之後雅各還是不放心.緊接著,在夜色之中,火光之下,雅各急忙為迎接以掃預備禮物(14-16節):母山羊二百隻,公山羊二十隻,母綿羊二百隻,公綿羊二十隻,奶崽子的駱駝三十隻各帶著崽子；母牛四十隻,公牛十隻,母驢二十匹,驢駒十匹.每樣各分一群,而且群群有相隔的距離.雅各心裡想:「我借著在我前頭去的禮物解他的恨,然後再見他的面,或者他容納我.於是禮物先過去了；雅各便在隊中住宿,只是一躺下驚恐就漫過了他的心:「萬一以掃殺了全家連同妻子兒女,再奪去全部財產怎麼辦?」「以掃向來擅於獵殺技能,此次又帶了四百人過來,一定是來者不善.」雅博渡口已經再也不能入眠.

我們一起來看聖經22,23節:雅各夜裡起來,帶著他的兩個妻子,兩個婢女和十一個兒子,都過了雅博渡口,我們不知道雅各為什麼突然改變主意,連夜起來帶他們過河,但不管是他有意還是無意,神已經奇妙地預備了與他單獨相處的時間,24節說「只剩下雅各一人,有一個人來與他摔跤,直到黎明.」弟兄姊妹,或許你我都遭遇過艱難,但當我們在艱難之際,是否有預備好與上主獨處的時間呢?

對雅各來說,二十年前,他一個人離鄉背井,在伯特利的時候,主向他顯現.現在,在他最需要的時候,神再次施恩來見他.弟兄姊妹們,神總是在他的兒女最需要幫助的時候出現,但是顯現的方式卻有所不同,向天路客亞伯拉罕,神就做客旅；對將軍約書亞時,神就做元帥,向雅各,神就以摔跤手來到,因為,雅各整個壯年的人生都在與人角力,以掃,以撒,拉班等都是他角力的對象.這裡聖經明說:是神主動來與他摔跤,說明神何等熱切的要除去我們信心生活中的阻礙.

我們來看聖經25節,「那人見自己勝不過他」.起初雅各拼命抵抗,雅各為他心愛的拉結將石頭轉離飲羊井的那股力量依然存在!當年雅各隻身投奔舅父拉班,在巴旦亞蘭的飲羊井旁,與拉結相遇,不顧牧人們的反對,執意搬去壓在井口的大石,為拉結之羊飲水.此次,他也無意示弱.弟兄姊妹,其實我們也一樣.神在我們裡面摔跤,哪一次我們願意認輸?江守道:「在你軟弱之時,請別憂傷；當你剛強之時,卻要小心跌倒.」V25「那人見自己勝不過他,就將他的大腿窩摸了一把,雅各的大腿窩正在摔跤的時候就扭了.」聖經在這裡讓我們看到,對方只在他大腿窩上打了一下,雅各就脫了臼.神要提醒我們:人的計謀,人的強壯,人的財力等這些都靠不住,越以此自得,你就跌的越慘:多年前,雅各用一碗紅豆湯換得長子名分,卻為此種下了兄弟仇恨的種子,父親年邁的時候,雅各連同母親騙得長子的祝福,最後卻背井離鄉,在外飄蕩二十年,在母舅那裡用計致富,最後卻差點惹下殺身之禍.

例1 我爺爺的經歷:人生依靠自己,終必落空.

聖經馬太福音16章26節,「人若賺得全世界,賠上自己的生命,有什麼益處呢?人還能拿什麼換生命呢?」弟兄姊妹,我們強壯的大腿窩需要被神以溫柔的手來觸摸,我們才懂得屈身下跪,求神開路.

V26,「那人說,天黎明瞭,容我去吧!」既然勝負已分,是否可以容神離去?既然知道我們自己的無有,是否只有自我消沈?「不!」雅各說:「你若不給我祝福,我就不容你走!」(V27)我們要知道與我們摔跤的是神自己,我們就不能因失敗而離去,而應從反抗轉向倚靠.我們需要的是緊抓不放,有如一個淹在水中的人,抓住一根救命稻草一般.邁爾博士在此形容道:我們可以看到這一位跛腳的先祖是如何拼命的依賴神來提供幫助和拯救.他已經放棄了一切人為的抵抗,緊緊的抓住那人,就像是受驚的孩子抱住自己父親的頸項一樣.弟兄姊妹,你知道嗎?這才是蒙福的態度,也是得力的途徑.

我們知道這一次是神自己刺激了這場鬥爭,好叫衪以自己似乎被打敗的方式而征服我們.我們手上的武器乃是眼淚和懇求.這次的祝福不是給行善受報的人,而是給拼命求取的人；不是給意氣奮發的人而是給精疲力竭的人.弟兄姊妹,你是否也有這樣一個絕境之時,正像泰戈爾詩里所描繪的:我們的儲糧已絕,前途已盡,旅程已終,精力已竭.要知道在我們信仰的歷程當中有時正是:人的盡頭,神的起頭.我想我們中間有些人是從內地過來讀書；也有些人是已經在香港工作；再或者你們中間,因著神的感動,也在教會或機構中做著義工服事.如果你的學習或者工作順利,那你要為此感謝神.但如果你正面臨艱難,你更要抓住神不放,因為這或許正是神要你學習信心與仰望的功課. 詩127:1 「若不是耶和華建造房屋,建造的人就枉然勞力；若不是耶和華看守城池,看守的人就枉然警醒.」讓我們在各樣的境遇中,都來學習依靠神.

我們再看27節.那人問他「你叫什麼名字?」他說:「我名叫雅各.」那人肯定知道他的名字叫雅各,那為什麼還要問他呢?我們依稀記得,曾經也有人問他這個問題,他卻撒了慌!在以撒年紀老邁,眼睛昏花的時候,雅各手臂上裹著皮毛就近他的父親.以撒在床頭問他:「我兒,你是誰?」他對父親說:「我是你的長子以掃.」雅各因為擅於使用計謀或早已臭名遠揚.主當然知道他的名字,所以當主問他,你是誰的時候,乃是要雅各承認自己是誰,然後交給主來改變.

V28.那人說:「你的名不要再叫雅各,要叫以色列.」以色列名字的意思乃是「神的王子」,在這裡聖經又補充到:因為你與神與人較力都得了勝.雅各他打輸了,為什麼神卻說他得了勝?原來,這裡聖經告訴我們:當人自以為剛強,成功,有能力時,他在神面前不會被神看為得勝者；相反,當一個人常常認識到自己軟弱,無能,而專心尋求神的恩典時,他在神面前才被視為得勝.所以,使徒保羅說:「如果誇耀,我就誇耀我的軟弱,好叫基督的能力覆庇我.因為神的能力在人的軟弱上顯的完全」.(林後12:9)

例2:自己這近三年深圳服侍的經歷.

縱觀雅各的一生,實在讓我們看到神的賜福與護佑在他身上彰顯的淋灕盡致.1,未出生時,神應許:將來大的要服事小的.2,被迫離家時,神在伯特利顯現:「我是亞伯拉罕的神,也是以撒的神……我與你同在,你無論往哪裡去,我必保佑你,領你歸回這地,總不離開你棄你,直到我成全了向你所應許的.」3,拉班追趕了七天,終於將他追上時,拉班自己說:「在我手上本有能力害你,但是你的神阻止了我.」4,雅博渡口,雅各七次下跪前行,憂心忡忡,最終以掃前仇不計,流淚相見.5,兒子在示劍惹禍,當地的人聚集起來要追殺雅各全家時,聖經記載(創35:5)「神使那周圍城邑的人都甚驚懼,就不追趕.」

希伯來書在11章21節記載雅各臨終之前乃是「扶著杖頭敬拜神」,這句話不但提示他對神的信心,為著神將會借著神自己的應許賜福給他的子孫而感謝神,更表示他不忘在雅博渡口瘸了腿的經歷.為這瘸腿的經歷,他特要在生命臨近終點之時,起來扶著杖頭來敬拜神,獻上他最後的感恩,承認神在這件事上奇妙的作為,乃是他一生之中最值得記住,感恩的大事.

雅各在雅博渡口與神摔跤.摔跤在希伯來文里與「雅各」,與雅博都是半諧音,記錄者有意作這樣的處理一定有他的特殊含義,總是讓我們記住這一刻的地點與這一刻的事件.讓我們也記住這是雅各一生的轉折點.在這裡我們需要思想和回應的有:

1,這個和雅各摔跤的人可曾找過你?

2,你是否覺得長期以來,你一直追尋,珍愛的某些東西現在必須放棄?

3,你是否明明知道,應當全人順服在神面前,卻在你裡面有一股強烈的悖逆的意願使你無法向神降服?

4,你是否願意,被神以他奇妙的手觸摸你的強壯之處,使你成為柔軟,並且蒙福?

另外,我們也必須認為到:屬靈生命的成長需要經歷「雅博渡口」,但它的刺激並非僅僅在於轉瞬之間,乃是一生之久.雅各腿瘸了,並且整個以色列民都不吃大腿窩的筋,正是作為不斷的提醒.實際上你會發現經歷了這次事件之後雅各並非就是痛改前非,煥然一新.在他的身上我們可以看到人性的幽暗面及軟弱面.但也正是因為如此,更能讓我們這些軟弱的人去體認,去接近這一位最終的得勝者:神的王子-雅各.

弟兄姊妹,或許這一個受傷的大腿窩在你的身上也發生過,請你不必刻意的求神挪去!儘管它有時會有不便,痛苦,甚至煎熬,但那也是激勵,輝煌,與得勝.但願弟兄姊妹都能經歷並記念雅博渡口那一夜的摔跤；但願弟兄姊妹在與神摔跤之時,也能像先祖雅各一樣在絕境當中,被神以祂奇妙,慈愛的手翻轉,使我們裡面抓取世界的雅各能真正成為神的王子以色列!



\newpage

\section{一生之久的恩典}
\label{sec:323}
\textbf{鄧瑞強博士 - 一生之久的恩典}
\newline
\newline
\hyperref[sec:322]{< < < PREV SERMON < < <}
~
\hyperref[sec:toc]{[返目錄]}
~
\hyperref[sec:324]{> > > NEXT SERMON > > >}


詩篇 30:1-12
\newline
\begin{longtable}{cl}
\hline
\hline
章節 & 經文 (和合本修訂版)\\
\hline
30:1 & \begin{tabularx}{0.7\textwidth}{X} 耶和華啊，我要尊崇你，因為你救了我，不讓仇敵向我誇耀。 \end{tabularx} \\ \\ \relax
30:2 & \begin{tabularx}{0.7\textwidth}{X} 耶和華－我的神啊，我呼求你，你醫治了我。 \end{tabularx} \\ \\ \relax
30:3 & \begin{tabularx}{0.7\textwidth}{X} 耶和華啊，你救我的性命脫離陰間，使我存活，不至於下入地府。 \end{tabularx} \\ \\ \relax
30:4 & \begin{tabularx}{0.7\textwidth}{X} 耶和華的聖民哪，你們要歌頌他，要頌揚他神聖的名字。 \end{tabularx} \\ \\ \relax
30:5 & \begin{tabularx}{0.7\textwidth}{X} 因為，他的怒氣不過是轉眼之間，他的恩典乃是一生之久；一宿雖然有哭泣，早晨便必歡呼。 \end{tabularx} \\ \\ \relax
30:6 & \begin{tabularx}{0.7\textwidth}{X} 至於我，我凡事順利，就說：「我永不動搖。」 \end{tabularx} \\ \\ \relax
30:7 & \begin{tabularx}{0.7\textwidth}{X} 耶和華啊，你曾施恩，使我穩固如山；你轉臉不顧，我就驚惶。 \end{tabularx} \\ \\ \relax
30:8 & \begin{tabularx}{0.7\textwidth}{X} 耶和華啊，我曾求告你；我向耶和華懇求： \end{tabularx} \\ \\ \relax
30:9 & \begin{tabularx}{0.7\textwidth}{X} 「我被害流血，下到地府，有何益處呢？塵土豈能稱謝你、傳揚你的信實嗎？ \end{tabularx} \\ \\ \relax
30:10 & \begin{tabularx}{0.7\textwidth}{X} 耶和華啊，求你應允我，憐憫我！耶和華啊，求你幫助我！」 \end{tabularx} \\ \\ \relax
30:11 & \begin{tabularx}{0.7\textwidth}{X} 你將我的哀哭變為跳舞，脫去我的麻衣，為我披上喜樂， \end{tabularx} \\ \\ \relax
30:12 & \begin{tabularx}{0.7\textwidth}{X} 使我的靈歌頌你，不致緘默。耶和華－我的神啊，我要稱謝你，直到永遠！ \end{tabularx} \\ \\
[1ex]
\hline
\hline
\end{longtable}
講題:「一生之久的恩典」　Lifetime Favour

經文:詩篇30:1-12

講員:鄧瑞強博士

場合:香港中文大學崇基學院禮拜堂

日期:2016年6月5日

各位弟兄姊妹,早安.

我們常說:「天跌下來當被cum 」,但當天花板跌下來時,就不能當被cum 了.很難想像,在我們的城巿,有現代建築物的天花板整幅塌下來.看來安全的地方,也可以不堪一擊.(按:2016年5月20日,香港城巿大學某運動中心的天花板塌了下來.)

在人生裡,我們都會為自己製造不少「生命的天花板」,去保護自己,使生命不用日曬雨淋.我們會努力賺錢,我們會努力成就某些事業,我們會努力達致成功.這些都是「生命的天花板」,讓我們免受別人的嘲笑,免受社會的風雨的欺凌,讓我們安心.但是,怎能保證這「生命的天花板」不會一下子塌下來.

對於生命的無常,詩人都是很敏感的.今日,我們看一篇詩篇,看看詩人對無常生命有何洞見.

今日講道的經文:詩篇30篇

詩30:1　〔大衛在獻殿的時候,作這詩歌.〕耶和華啊,我要尊崇你,因為你曾提拔我,不叫仇敵向我誇耀(或譯:不讓仇敵因勝過我而喜樂).

詩30:2　耶和華──我的神啊,我曾呼求你,你醫治了我.

詩30:3　耶和華啊,你曾把我的靈魂從陰間救上來,使我存活,不至於下坑(或譯:下到深淵).

詩30:4　耶和華的聖民哪,你們要歌頌他,稱讚他可記念的聖名.

詩30:5　因為,他的怒氣不過是轉眼之間；他的恩典乃是一生之久.一宿雖然有哭泣,早晨便必歡呼.

詩30:6　至於我,我凡事平順便說:我永不動搖.

詩30:7　耶和華啊,你曾施恩,叫我的江山穩固(在和合本聖經裡,「江」字旁邊有幾點,表示原文沒有這字.這句話可譯為:「叫我穩固如山」,或者譯:「叫我穩如泰山」)；你掩了面,我就驚惶.

詩30:8　耶和華啊,我曾求告你；我向耶和華懇求,說:

詩30:9　我被害流血,下到坑中,有甚麼益處呢?塵土豈能稱讚你,傳說你的誠實嗎?

詩30:10　耶和華啊,求你應允我,憐恤我!耶和華啊,求你幫助我!

詩30:11　你已將我的哀哭變為跳舞,將我的麻衣脫去,給我披上喜樂,

詩30:12　好叫我的靈歌頌你,並不住聲.耶和華──我的神啊,我要稱謝你,直到永遠!

這首詩第一個值得我們留意的地方,就是這首詩的題目.這首詩的題目是:「大衛在獻殿的時候,作這詩歌」,原文直譯是:「詩,獻殿之歌,關乎大衛的」.這詩不一定是大衛作的,也可能是為大衛而寫的,總之是關乎大衛的.值得留意的,是大衛不可能獻殿,大衛時代根本沒有聖殿.聖殿是他的兒子所羅門建的.

若果說,這詩的主題與大衛及聖殿有關,則舊約裡最相關的其中一個故事便是歷代志上21章的故事.那裡講到,大衛戰無不勝,便心高氣傲,數點戰士,以顯兵威.有點似現代的閱兵,為了宣揚國家的實力.神不喜悅大衛耀武揚威,便降瘟疫,死了七萬人.滅命的天使,站在耶路撒冷一個原居民阿珥楠(Ornan the Jebusite)的禾場上.大衛見到,即時向神認罪,並在這禾場上築了一座壇,在其上獻祭.神的怒氣就止住了.這築壇的地方,就是他的兒子日後建聖殿的地方.

詩篇30篇的主題出現「大衛」和「聖殿」的字眼,若這主題真與上述事件有關,則詩篇30篇便是講人的驕傲與神的懲罰,滅命天使所站之地正是親近神的聖殿之所在,認罪的時刻正是經歷神的憐憫的時刻.

希伯來詩歌的特色是對比,平行.詩篇30篇也不例外,有大量對比,平行.

第1節:「不讓仇敵因勝過我而喜樂」,對應第11節:「你將我的麻衣脫去,給我披上喜樂」.兩次的「喜樂」原文字根相同,前者是動詞,後者是名詞.

第2節:「耶和華──我的神啊,我曾呼求你,你醫治了我」,對應第10節:「耶和華啊,求你應允我,憐恤我!」

第3節講「下坑」,第9節也講「下坑」,下到生命的深淵,無助的深淵.

第4節講「神的聖民要歌頌神」,第8節講「我向耶和華懇求」.

第5節講「他的怒氣不過是轉眼之間；他的恩典乃是一生之久」,第7節講「你曾施恩,叫我穩如泰山；你掩了面,我就驚惶」.

這首詩的中間就是第6節:「至於我,我凡事平順便說:我永不動搖」.人的這種自大,正是一切問題的根源.

這首詩的結尾(第12節),講到「歌頌」神,「稱謝」神,這其實是對比這首詩的題目:「獻殿之歌」.

我們看看第1－3節.

詩30:1　耶和華啊,我要尊崇你,因為你曾提拔我,不讓仇敵因勝過我而喜樂.

詩30:2　耶和華──我的神啊,我曾呼求你,你醫治了我.

詩30:3　耶和華啊,你曾把我的靈魂從陰間救上來,使我存活,不至於下坑.

詩人的人生經驗,他下到「坑」中,下到人生的低谷,下到生命的深淵.詩人形容這是「陰間」,一個死地,一種絕境.叫天不應,叫地不聞.人在這裡,好像只能靜待死亡.詩人在深淵中,向神求救.然後,第1節講「你曾提拔我,不讓仇敵因勝過我而喜樂」.「提拔」這個字,希伯來文是dalah,指從井裡打水.人在深淵中,如在井中,神像打水一樣將我們打上來.生命的喜樂,一定不是屬於仇敵的,而是屬於我們的.因著這一切經歷過的經驗,詩人讚美神.

這一段落裡,出現很多十分形象化的對比.我在深淵,神將我拉上來；我在陰間,神將我救上來；我在死亡邊緣,神使我存活；我呼求,神醫治.簡言之,我在絕境中,神卻將我救回.

我們看這詩篇的第二個段落:第4－5節.

詩30:4　耶和華的聖民哪,你們要歌頌他,稱讚他可記念的聖名.

詩30:5　因為,他的怒氣不過是轉眼之間；他的恩典乃是一生之久.一宿雖然有哭泣,早晨便必歡呼.

這兩節經文不是向神呼籲的,而是呼籲我們這些信神的人的.詩人呼籲我們進入他的經歷中,這不獨是他的經歷,這也是我們的經歷.

詩人用強烈對比的字眼去表達他的生命體會.他經驗到神的怒氣,但也經驗到神的恩典,神的悅納.他經驗到神的怒氣只不過是轉眼之間,神的悅納卻是一生之久.哭泣是有的,可能持續一整晚；但早晨一到,內心便爆發出歡樂.

以基督徒的經驗來說,我們經歷罪的深重,像被壓在深淵裡,但也經驗被神寬恕的輕省,像被拉出深淵.我們經驗人性的軟弱,但也經驗在基督裡的剛強.我們經驗恐懼,也經驗天賜的平安.我們會經驗死亡,也會經驗復活.

詩人提示我們,在神的恩典裡,生命的新開始是可能的.「神的怒氣不過是轉眼之間,他的恩典乃是一生之久」.

詩篇30篇的中心,是第6節.

詩30:6　至於我,我凡事平順便說:我永不動搖.

這便是一切問題的核心.「我凡事平順便說:我永不動搖」.這經文用別的說法來表達,便是:我發達的時候,我覺得大地在我腳下；我成功的時候,我認為永遠不會失敗；我健康的時候,我認為自己永不會病；我「生生猛猛」的時候,我認為自己不會死.

或許,這就是人類墮落背後的平凡想法.我要用我手,創造我的人生.昔日,亞當夏娃在伊甸園,吃禁果時,可能都是想,我們不用神指引我們了,我們自己尋求永生的大道便是了.昔日,當德國人擁護希特拉上台時,他躊躇滿志,意氣風發,一步一步地用自己的手創造國家的明天.這時,他最能講出這經文:「至於我,我凡事平順便說:我永不動搖」.講這經文的人,是一個自大的人.自大的人,毫無感恩的心.他不會敬畏天,也不會尊重人.這心態引導人一步一步的墮落,一步一步地接近黑暗的深淵.

詩人將這經文放在這詩篇的中間,讓我們好好面對自己心底裡的這種想法.

詩篇30篇的第四個段落:第7－8節,這回應第二段落的內容.

詩30:7　耶和華啊,你曾施恩,叫我穩如泰山；你掩了面,我就驚惶.

詩30:8　耶和華啊,我曾求告你；我向耶和華懇求.

放下自大,詩人認識生命的真相.詩人用一些對比去呈現這真相.神施恩時,面向我們時,我們能安穩如山；神掩面時,不看我們時,我們便驚惶失措.生命的一切,都在神的手中,而不是在我們手中.

或許,有一天,我們會投資失誤而破產,會忽然中風而行動不便,會患重病而徘徊在死亡邊緣.這就是詩人講的「下坑」,墮入深淵,生命無光.「生命的天花板」忽然塌了下來,我們以為穩當的保護,忽然消失了.

我們來到這詩篇最後一個段落:第9－12節,這回應第一段落的內容.

詩30:9　我被害流血,下到坑中,有甚麼益處呢?塵土豈能稱讚你,傳說你的誠實嗎?

詩30:10　耶和華啊,求你應允我,憐恤我!耶和華啊,求你幫助我!

詩30:11　你已將我的哀哭變為跳舞,將我的麻衣脫去,給我披上喜樂.

詩30:12　好叫我的靈歌頌你,並不住聲.耶和華──我的神啊,我要稱謝你,直到永遠!

第9節看來是詩人和神之間的討價還價.我們要留意,這不是平等的雙方在討價還價,不是那種「神你若救我,我便tong 雞還神」的討價還價,這討價還價只是一種「目中無神」的叫價.詩人自認自己只不過是一個敬拜神的,像塵土一樣微小的人.他若死去,神便會損失一位敬拜者.若神不解救詩人,他死去,神當然會失去一位敬拜者,這是塵世的一個事實,詩人絕無威脅神之意.與其說這是討價還價,不如說詩人藉此將自己放在神之下,表明自己敬拜神的心.

詩人十分明白神的本性,「神的怒氣不過是轉眼之間,他的恩典乃是一生之久」.在這裡,詩人再次表達蒙神恩典的經驗,又是用一些對比性的語言.「哀哭」與「麻衣」,多麼充滿死亡的意味.在神的恩典裡,這變為「跳舞」和「喜樂」.

第12節對應這首詩的「題目」,也同時作為一總結,詩人呼籲我們歌頌神.

有個電視節目,叫《跟住矛盾去旅行》,內容大家都知道了.詩篇30篇可以叫《跟住矛盾去生活》,內中呈現生活裡各種事情的兩極:「完全的自大」與「像塵土一樣的卑微」,「認為自己穩如泰山」與「忽然落入陰間的深淵中」,「如在深坑裡」與「被神提上來」,「神的怒氣,轉眼之間」與「神的恩典,一生之久」,「晚間的哭泣」與「早晨的歡呼」,「神施恩」與「神掩面」,「哀哭,麻衣」與「跳舞,喜樂」等.充滿對立,但這就是生活.詩人在訴說這些對立時,卻肯定神施恩是多於神掩面的,生命的喜樂是多於哀哭的,神的怒氣不過是轉眼之間而他的恩典卻是一生之久.在這不對稱的對立裡,你會否和詩人一樣,稱謝神,直到永遠.

但願:榮耀歸於聖父,聖子,聖靈.阿們.



\newpage

\section{獨排眾議舉動觸目}
\label{sec:324}
\textbf{許開明牧師 - 獨排眾議,舉動觸目}
\newline
\newline
\hyperref[sec:323]{< < < PREV SERMON < < <}
~
\hyperref[sec:toc]{[返目錄]}
~
\hyperref[sec:325]{> > > NEXT SERMON > > >}


路加福音 7:36-50
\newline
\begin{longtable}{cl}
\hline
\hline
章節 & 經文 (和合本修訂版)\\
\hline
7:36 & \begin{tabularx}{0.7\textwidth}{X} 有一個法利賽人請耶穌和他吃飯，耶穌就到那法利賽人家裡去坐席。 \end{tabularx} \\ \\ \relax
7:37 & \begin{tabularx}{0.7\textwidth}{X} 那城裡有一個女人，是個罪人，知道耶穌在法利賽人家裡坐席，就拿著盛滿香膏的玉瓶， \end{tabularx} \\ \\ \relax
7:38 & \begin{tabularx}{0.7\textwidth}{X} 站在耶穌背後，挨著他的腳哭，眼淚滴濕了耶穌的腳，就用自己的頭髮擦乾，又用嘴連連親他的腳，把香膏抹上。 \end{tabularx} \\ \\ \relax
7:39 & \begin{tabularx}{0.7\textwidth}{X} 請耶穌的法利賽人看見這事，心裡說：「這人若是先知，一定知道摸他的是誰，是個怎樣的女人；她是個罪人哪！」 \end{tabularx} \\ \\ \relax
7:40 & \begin{tabularx}{0.7\textwidth}{X} 耶穌回應他說：「西門，我有話要對你說。」西門說：「老師，請說。」 \end{tabularx} \\ \\ \relax
7:41 & \begin{tabularx}{0.7\textwidth}{X} 耶穌說：「有兩個人欠了某一個債主的錢，一個欠五百個銀幣，一個欠五十個銀幣。 \end{tabularx} \\ \\ \relax
7:42 & \begin{tabularx}{0.7\textwidth}{X} 因為他們無力償還，債主就開恩赦免了他們兩個人的債。那麼，這兩個人哪一個更愛他呢？」 \end{tabularx} \\ \\ \relax
7:43 & \begin{tabularx}{0.7\textwidth}{X} 西門回答：「我想是那多得赦免的人。」耶穌對他說：「你的判斷不錯。」 \end{tabularx} \\ \\ \relax
7:44 & \begin{tabularx}{0.7\textwidth}{X} 於是他轉過來向著那女人，對西門說：「你看見這女人嗎？我進了你的家，你沒有給我水洗腳，但這女人用眼淚滴濕了我的腳，又用頭髮擦乾。 \end{tabularx} \\ \\ \relax
7:45 & \begin{tabularx}{0.7\textwidth}{X} 你沒有親我，但這女人從我進來就不住地親我的腳。 \end{tabularx} \\ \\ \relax
7:46 & \begin{tabularx}{0.7\textwidth}{X} 你沒有用油抹我的頭，但這女人用香膏抹我的腳。 \end{tabularx} \\ \\ \relax
7:47 & \begin{tabularx}{0.7\textwidth}{X} 所以我告訴你，她許多的罪都赦免了，因為她愛的多；而那少得赦免的，愛的就少。」 \end{tabularx} \\ \\ \relax
7:48 & \begin{tabularx}{0.7\textwidth}{X} 於是耶穌對那女人說：「你的罪都赦免了。」 \end{tabularx} \\ \\ \relax
7:49 & \begin{tabularx}{0.7\textwidth}{X} 同席的人心裡說：「這是甚麼人，竟赦免人的罪呢？」 \end{tabularx} \\ \\ \relax
7:50 & \begin{tabularx}{0.7\textwidth}{X} 耶穌對那女人說：「你的信救了你，平安地回去吧！」 \end{tabularx} \\ \\
[1ex]
\hline
\hline
\end{longtable}
講題:獨排眾議,舉動觸目　Override All Objections Alone, Attract the Gaze of People

講員:許開明牧師

場合:香港中文大學崇基學院禮拜堂

日期:2016年6月12日

經文:路加福音7:36-50

「36有一個法利賽人,請耶穌和他喫飯．耶穌就到法利賽人家裡去坐席.37那城裡有一個女人,是個罪人．知道耶穌在法利賽人家裡坐席,就拿著盛香膏的玉瓶,38站在耶穌背後,挨著他的腳哭,眼淚溼了耶穌的腳,就用自己的頭髮擦乾,又用嘴連連親他的腳,把香膏抹上.39請耶穌的法利賽人看見這事,心裡說,這人若是先知,必知道摸他的是誰,是個怎樣的女人,乃是個罪人.40耶穌對他說,西門．我有句話要對你說．西門說,夫子,請說.41耶穌說,一個債主,有兩個人欠他的債．一個欠五十兩銀子,一個欠五兩銀子．42因為他們無力償還,債主就開恩免了他們兩個人的債.這兩個人那一個更愛他呢.43西門回答說,我想是那多得恩免的人.耶穌說,你斷的不錯.44於是轉過來向著那女人,便對西門說,你看見這女人麼．我進了你的家,你沒有給我水洗腳．但這女人用眼淚溼了我的腳,用頭髮擦乾.45你沒有與我親嘴,但這女人從我進來的時候,就不住的用嘴親我的腳.46你沒有用油抹我的頭,但這女人用香膏抹我的腳.47所以我告訴你,他許多的罪都赦免了．因為他的愛多．但那赦免少的,他的愛就少.48於是對那女人說,你的罪赦免了.49同席的人心裡說,這是甚麼人,竟赦免人的罪呢.50耶穌對那女人說,你的信救了你,平平安安的回去罷.」

引言

各位,就今天所讀的經課,想同大家,從一個比較少人提及的角度,就是「獨排眾議」去思想.

自古到今,每個人在社群中都會掙扎著,這個問題:做一個隨波逐流的人,或是做個有性格的人?我們應走群眾路線,還是走自己的路?或再有人說:我只是個小人物,可能走群眾路線較安全；但如果我是大人物,領袖,我就應獨排眾議,走自己的路.然而,走群眾路線,往往往是跟風,沒有個人的風格,沒有個性,個人常被忽略.走群眾路線的人,可能失去自己.但走自己路的人,可能不被接受,失去群眾.很明顯,機構主管常須掙扎:堅持走自己的路,還是顧及下屬的看法?我們處事為人在這兩極端間應如何作取捨,如何平衡?期盼這段經文給我們一些亮光和啟迪!

第一幕:耶穌與有罪女子,在法利賽人家裡,有罪女子出現

「36有一個法利賽人,請耶穌和他喫飯．耶穌就到法利賽人家裡去坐席.37那城裡有一個女人,是個罪人．知道耶穌在法利賽人家裡坐席,就拿著盛香膏的玉瓶,38站在耶穌背後,挨著他的腳哭,眼淚溼了耶穌的腳,就用自己的頭髮擦乾,又用嘴連連親他的腳,把香膏抹上.」

第二幕:耶穌與法利賽人

1法利賽人不滿那有罪女人；耶穌說故事:多得主恩免,多愛主

「39請耶穌的法利賽人看見這事,心裡說,這人若是先知,必知道摸他的是誰,是個怎樣的女人,乃是個罪人.40耶穌對他說,西門．我有句話要對你說．西門說,夫子,請說.41耶穌說,一個債主,有兩個人欠他的債．一個欠五十兩銀子,一個欠五兩銀子．42因為他們無力償還,債主就開恩免了他們兩個人的債.這兩個人那一個更愛他呢.43西門回答說,我想是那多得恩免的人.耶穌說,你斷的不錯.」

2耶穌指出法利賽人的虧欠,有罪女人補償他的虧欠

「44於是轉過來向著那女人,便對西門說,你看見這女人麼．我進了你的家,你沒有給我水洗腳．但這女人用眼淚溼了我的腳,用頭髮擦乾.45你沒有與我親嘴,但這女人從我進來的時候,就不住的用嘴親我的腳.46你沒有用油抹我的頭,但這女人用香膏抹我的腳.47所以我告訴你,他許多的罪都赦免了．因為他的愛多．但那赦免少的,他的愛就少.48於是對那女人說,你的罪赦免了.」

第三幕:耶穌與同席的人同席的人議論

「49同席的人心裡說,這是甚麼人,竟赦免人的罪呢.50耶穌對那女人說,你的信救了你,平平安安的回去罷.」

在本段經文,我們會看見三位「獨排眾議」的人.首先,「獨排眾議」是什麼意思?我的定義是:堅持自己的主張,排斥或不理會眾人或人家的看法.

三位獨排眾議

獨排眾議的女子

這個有罪女人有下列獨特的表現,一定引起議論的.猶太女子不可接近男子,壞女人更不應,更不可接觸男人的身體.此女子對著男子流淚,滴在男子腳上,又披頭散髮,用髮抹腳,再用香油去抹.這有罪女人的所為,一般人不會贊同,法利賽人更覺「反胃」,但耶穌不反對,反而讚賞,令當時人也感到希奇.

獨排眾議的主人家:法利賽人

那法利賽人既然是主人家,是他主動邀請耶穌的,他卻虧待耶穌.在客人進入屋子裡頭,作主人的有三件事要做.第一是見面禮,主人會搭手在客人肩上,給他平安的一吻.這一吻是表示尊敬的記號；作為一個傑出的拉比,這是永遠不會省略的.第二是洗腳禮,路上佈滿泥塵,腳上的鞋履亦只不過是用帶子把鞋底綁在足上,因此作主人會把冷水傾潑在客人的雙足,使客人感到潔淨涼快.第三是抹油禮,最後,主人會點起一撮味道香甜的香料,或是在客人頭上抹上數滴玫瑰香油,使客人聞到陣陣香味,感到愉快.這三件事是款客之道,然而今作主人的一件也沒有辦到.

如果你是耶穌,你會有何感覺?不被重視,不被尊重?這位主人家西門是法利賽人,應是最認識禮儀和律法的人.法利賽人,這個名稱是從希伯來文來的,意思是「分開」,「分別出來」,「與別不同」.法利賽人要監察:猶太人要從世俗中,藉著信奉上帝,分開出來；

猶太人從世俗中,藉著守律法,分開出來.這身份使法利賽人,會有「獨排眾議」的行為.

獨排眾議的耶穌

當主人家及眾人不滿這有罪女人的行為時,耶穌卻容許她接近,並且稱讚這女子的所為,又赦免這有罪女人的罪,這是法利賽人及各人所不能接受的,因為只有上帝才能赦罪,西門請吃飯,耶穌不但沒有多謝主人招待,反而指責他招待不周.雖然這西門值得被罵,但一般人不會公開直指其非,耶穌卻表現直接.這就是獨排眾議的耶穌.

五餅二魚神蹟後,群眾希望耶穌帶領猶太人推翻羅馬政權,耶穌卻不是他們期望的那位.當那穌釘在十字架上,有人挑戰道:「如果你是上帝的兒子,可以從十架上跳下來!」魔鬼也想耶穌走群眾路線,從聖殿跳下,耶穌又獨排眾議.耶穌餵飽五千人後,眾人想擁他為王,耶穌就退到山上去.

為何他們三位要獨排眾議,表現獨特?

獨排眾議的女子

我先說一個獨排眾議的老人的故事,這是一個丹麥作家寫的,是一個很特別的故事,小朋友不易明,但大人會會心微笑.

一個農村老人離家去趕集,他要賣掉家裡的一匹馬.他跟老婆告別後,往市集去.老人見到一個人拉著一頭牛,那個人問:「老人願不願把他的馬換那人的牛?」老人願意便得了頭乳牛,老太婆每天早上有新牛奶飲,成交!不久,老人見到有人趕著一頭豬,趕豬人對老人說:「你願不願把你的一頭牛換我的一頭豬?」老人想大豬生小豬,老太婆見到小豬那麼可愛,一定開心,成交!再向前走,見到一個人抱著一隻雞賣.那人問要不要以豬換雞?老人心想,老太婆天天有蛋吃,於是用豬換了雞.再向前,他見到有個帶著一袋欄蘋果,那人問老人願否以雞換爛蘋果?老人想,爛蘋果埋在院子裡,可以長出蘋果樹,老太婆最愛食蘋果,一定高興,成交!三個商人朋友見到老人拿著一袋爛蘋果,感到奇怪,老人就告知他怎樣由一匹馬換到一袋爛蘋果.商人大笑並說:「老人,老太婆一定罵你蠢,並賞你一巴掌.」老人回答說:「可以打賭,老太婆不會罵我,打我,甚至會親吻我.」結果,老人回家並告知老太婆換物的經過,想不到,老人每講一個交易,老太婆就大笑,並讚他做得好,最後忍不住還錫他一口.三個商人都輸了,老人得了三個金幣.老太婆為何欣賞老人?因為老人每個交易都想著老婆,想到老婆的需要,這就是愛,這愛使他表現獨特,得到老太婆的愛.

對那女子來說,她表達感謝耶穌,敬愛耶穌的好機會；對她那女子來說,是她補償主人家對耶穌待客不周,不足的機會.這個女人帶一個盛香膏的白玉瓶,當耶穌斜倚在長凳上吃飯,頭挨近桌子那邊時,她站在耶穌背後,用眼淚洗耶穌的腳,並用自己的頭髮擦乾,又用嘴連連親他的腳,然後把昂貴的香膏抹上.如此的敬拜和奉獻,顯示她深信沒有東西是好得連耶穌也受不起的.這女子因著對耶穌的愛,勇敢接近耶穌,不理會別人的閒言閒語,她計劃接觸耶穌的腳,腳可能是人體最污穢的部份,因為易沾上路上泥濘；腳也可能是最臭的部份,因此猶太人進食前要洗腳,這女子是十分機智:沒有水洗腳,她是用眼淚水；沒有毛巾,她用頭髮抹腳,是用女性最高,最尊貴部份人抹最低,最污穢部份；沒有吻臉,她用口親腳,以示乾淨和尊敬,沒有抹油,她用她頸上的香膏抹腳,她頸上掛著一小瓶的香膏,香膏的名字叫「雪花香膏」(alabasters),是非常昂貴,她想把這些香膏,即是她所有的全部,傾注在耶穌的腳上.這用行動表達出她對耶穌的愛與尊重,也就是人的兩項基本的需要.

這女子是有備而來的,有錢人的家居都有一個偌大的天井,天井有時會有一個噴水池,天氣暖和之際,更可在此進膳.當時的習慣是當拉比在屋子裡進食,任何人都可以自由進入,以聽取他滿有智慧的金玉良言,這樣解釋了這個婦人為甚麼會在此出現.那女子把握機會,表達對耶穌的愛.

獨排眾議的西門

他對耶穌的無禮,這種奇特的表現是無意或故意?應是故意的,因為這種見面的傳統,怎可忘記呢!但經文未有提供線索,我們推想他是法利賽人,常站在與耶穌對立那一方,如果他對耶穌表現過份的親善,可能引起同道的不滿.這種偽善的表現,受到耶穌的指責,也使那女子感到西門對耶穌招呼不周,這種獨排眾議的表現是不好的,引起別人不滿.

法利賽人眼中的壞女人不應接近耶穌,因為她們不配得；法利賽人眼中的耶穌應拒絕壞女人,因為這壞女人會使耶穌身敗名裂；眾人眼中西門不應對耶穌太親切,因為他是位問題人物.許多時候,我們對別人都有些自以為是的看法,對家人,對朋友,對上司下屬；這些未必是正確,未必就是上主的心意.

耶穌引發西門作獨立的思維,西門是法利賽人,自高自大,覺得自己勤守律法,就是無罪的人,所以看不起那有罪的女人,目睹她犯倫理的規則,主動接近男人,並有不恰當的行為.耶穌採用的方法是相當智慧的方法,不直斥其非,不向他發表偉論.只向他講個比喻,然後要他自己悟出答案來.

「41耶穌說:『一個債主,有兩個人欠他的債,一個欠五十兩銀子,一個欠五兩銀子,42因為他們無力償還,債主就開恩免了他們兩個人的債.這兩個人那一個更愛他呢?』43西門回答說:『我想是那多得恩免的人.』耶穌說:『你斷的不錯.』」

在此,你見耶穌對這種自以為是,高人一等的人,跟他講故事,讓他懂得獨立思考,並講出他認為對的答案,不是拾人牙慧的.西門說對了,耶穌就讚他,令他開心.大家建立友善的關係後,耶穌才直斥其非,指出他的無禮,是很有智慧的做法.

獨排眾議的耶穌

在此耶穌提出一個獨特的看法:罪等於欠上帝的債.人的債即是人的罪；欠債要還債就是人犯罪要得赦免,每個人欠上帝的債,所以每個人都要還債.對於無力還債的人需要上帝的赦免,得恩典也小.欠債少的,得赦免少,得愛少,他愛上帝的心也少；相反,欠債多的,得赦免也大,得愛也大,他愛上帝的心也大,得恩典也大!因此耶穌要我們不可批評或歧視那些欠債多的人.

這個獨特觀點,打破法人及眾人的思維,原本以為上帝特別愛那些遵守律法的人,愛上帝的子民,愛義人,愛男人多於女人.上帝愛義人,但上帝也愛罪人.不過,不論上帝多愛有重罪的人,犯重罪的人的罪果仍是要自己承擔,只是上帝不會不愛有罪的人,犯重罪的人,因為上帝愛世人,愛好人,壞人,罪人.

耶穌把握此機會講出:恩免愈多,愛也愈多.耶穌稱讚那女子,並赦了她的罪,使她得平安.「50耶穌對那女人說:『你的信救了你,平平安安的回去罷.』」原文直譯就是:「…你的信救了你,在平安裡走罷.」

獨特/獨立思維的挑戰

社會跟風隨波的氣氛

你玩過「請你跟我這樣做」的遊戲嗎?許多人都想「請你跟我這樣做」,要求人們有同一思想,同一動作,同一眼神,同一語句,就如當權者插手影響人的思維,請你跟我或跟某人這樣做.年少時跟毛主席做,毛主席說:「要學雷鋒.」今日仍然有些國家令人感到可惜的是,要人跟他的一套:打壓,令人受驚,妥協；軟禁,令人失去自由；禁聲,不可出聲,不可唱歌；洗腦,不是要令人思想空白,而是只有一種思想,真可怕.伏爾泰(Voltaire)說過:「人類永遠是瘋狂的,而那些自以為可以治療別人瘋症的人則是瘋狂之最(最瘋狂的人).」最可怕是,人們在不知不覺中,產生羊群的心態,好像喝了迷湯一樣,跟著別人走.

個人太注重別人的眼光和期望

女人不可主動接近男人,有罪的女人更不可主動接近受尊敬的夫子.那有罪的女人卻獨排眾議,不理會人家的閒言閒語.今天的人有兩極,不是太自我,就是太顧及他人的眼光,敢仰起鼻子走的人少,被人牽著鼻子走的人多；敢於擇善固執,堅持自已的理想,走自己的路的人少!

我們常發現自己要活在別人的訴求和期望中:太太對我有訴求或期望, 是子女的訴求/期望；上司或下屬對我有訴求,想我快完成他交辦的工作,最怕拖泥帶水；學生對我有期望,希望我能不偏不倚,公平,公正,但對他寬鬆些；作為公民,國家對你有期望,就是順服當權者,和諧,努力建設繁榮安定.我活在這麼多人的要求中,要是我達不到他們的要求,我就失望,連要求我的人也失望.結果,我只能選擇性達到要求,去完成我真的認同的要求,我達不到某些人的要求,不要自責.更重要的是清楚我自己有何要求和期望,因為自己的往往被忽視,例如自己的休息,健康,抱負,時間,機會,強項有否充份的發揮等.我如何滿足眾多人訴求和期望?我若被人訴求和期盼拉著走我便失去自由,失去自我.

人失去獨特,獨立的思維

在傳統教會長大,常聽說:「聖經這樣說,我們便這樣信」,「牧師(爸爸)這樣說,就是這樣」,「不要疑惑,只要信,單純的信,就有福!」這種信仰的態度,可能對某些人是有意義的,但對於一些真誠的,理性的,愛探求的信徒,是不足的.耶穌對那有罪女人的欣賞,關懷,寬恕,是獨特的,是法利賽人不會認同和接受的.作為一位大學生,或是從事研究的學者更需要有這種獨立的,獨特的思維,才不會被人牽著鼻子走,被人蒙蔽,失去自己.

獨特的人生觀-被咬了一口的蘋果

蘋果公司的標識為什麼是個被咬了一口的蘋果?引起注意和討論!一直好奇:為什麼蘋果公司的標識,是個被咬了一口的蘋果,直至今天,偶然看到,才懂得,被咬了一口的蘋果,有個淒涼的故事.第一,是被人咬了一口；第二,是被蟲咬了一口；第三,是被上帝咬了一口

被人咬了一口

被人咬,即是被人欺負或傷害過嗎?英國劍橋大學有一位年輕畢業生圖靈(Turing),以他為首,建造了世界上第一台真正的電腦,從而破獲了德國自以為無法被破解的密碼,完全掌握了德國潛艇部隊的秘密,最終滅了德國海軍,為提前打敗德國,做出了決定性的貢獻.但是艾倫.圖靈Alan Mathison Turing的命運,很悲慘,戰爭結後,他回到劍橋大學,過著同性戀的隱秘生活,當時的同性戀,在英國被視為重罪,痛苦非常.1954年6月8日,圖靈42歲,正逢進入他生命中,最輝煌的創造頂峰,一天早晨,女管家走進他的臥室,發現檯燈還亮著,床頭上還有個有毒的蘋果,只咬了一小半,圖靈沉睡在床上,一切都和往常一樣,但圖靈永遠地睡著了,不會再醒來……人的生命會被咬,被傷害的.

被上帝咬了一口

這是一個很獨特的,積極正向的故事.有一個人從小雙目失明,懂事後,他為此深感煩惱,認定是老天懲罰他,覺得這輩子都完了.後來,一位牧師對他說:「世界上每個人都是被上帝咬過一口的蘋果,都是有缺陷的,有的人缺陷比較大,是因為上帝特別喜歡他的芬芳.」他很受這段話的鼓舞,從此把失明看作是上帝的特殊鍾愛,開始振作起來；若干年後,成為一位有品有有牌的專業技師-盲人推拿師.

講這句的牧師加上說:上帝知道這件事後,笑道:「我很喜歡這個美麗而睿智的比喻,但要聲明一點,所謂缺陷是生理上的；那些有道德缺陷的人是爛蘋果,不是我咬的,是蟲蛀的.」

這個故事告訴我們,身體的缺陷,美麗或醜陋,甚至人生不好的遭遇,並不是老天的懲罰,因為我們都是上帝最鍾愛的人；不要因為你的某種缺憾而懊惱.這是一種獨特的,正向的思維.

楊絳的百歲感言也表現了一種獨立與獨特的思想,他的題目是《世界是自己的,與他人毫無關係》,裡面是這樣寫的:「我今年一百歲,已經走到了人生的邊緣,我無法確知自己還能走多遠,壽命是不由自主的,但我很清楚我快『回家』了.我得洗淨這一百年沾染的污穢回家.我沒有『登泰山而小天下』之感,只在自己的小天地裡過平靜的生活.細想至此,我心靜如水,我該平和地迎接每一天,準備回家.上蒼不會讓所有幸福集中到某個人身上,得到愛情未必擁有金錢；擁有金錢未必得到快樂；得到快樂未必擁有健康；擁有健康未必一切都會如願以償.保持知足常樂的心態才是淬煉心智,淨化心靈的最佳途徑.一切快樂的享受都屬於精神,這種快樂把忍受變為享受,是精神對於物質的勝利,這便是(我的)人生哲學.」

結語

獨排眾議就是堅持自已的主張看法,不理會或甚至排斥他人的意見或看法.這種態度有好亦有壞,問題是我們怎樣做.法利賽人獨排眾議,沒有按規矩接待耶穌,引起耶穌和那女子的不滿,這不能接受的表現,但那有罪女人卻善用這機會,有罪女人雖然有罪,她勇敢接受耶穌,補償西門的無禮,充份表達對耶穌的愛,卻成為佳話.

又有一個關於貓和豬的故事,牠們是好朋友.一天貓掉進大坑,豬拿來繩子,貓叫豬把繩子扔下來,結果它整捆扔了下去.貓很鬱悶的說:「這樣扔下來,怎麼拉我上去? 」豬說:「不然該怎麼做?」貓說:「你應該拉住一頭繩子啊!  」豬就跳下去,拿了繩子的一頭,說:「現在可以了!」貓哭了,哭得很幸福.

有的人不是很聰明,卻值得你終生擁有.豬雖然很笨,但豬聽貓話,豬愛貓,使貓亦愛笨豬,並感到很幸福.我有時也像豬,但笨一點也無妨,只要令身邊的人感到幸福!

結束祈禱

上主,在今天多元化的社會裡,講求齊一,共識,共融,埋堆,合作.祈求上主,叫我們仍能執善固執,不怕獨排眾議,分別出來,建立獨立與獨特氣質,阿們.



\newpage

\section{或去或留}
\label{sec:325}
\textbf{林豪恩先生 - 或去或留}
\newline
\newline
\hyperref[sec:324]{< < < PREV SERMON < < <}
~
\hyperref[sec:toc]{[返目錄]}
~
\hyperref[sec:326]{> > > NEXT SERMON > > >}


路加福音 8:26-39
\newline
\begin{longtable}{cl}
\hline
\hline
章節 & 經文 (和合本修訂版)\\
\hline
8:26 & \begin{tabularx}{0.7\textwidth}{X} 他們到了格拉森人的地區，就在加利利的對面。 \end{tabularx} \\ \\ \relax
8:27 & \begin{tabularx}{0.7\textwidth}{X} 耶穌上了岸，就有城裡一個被鬼附的人迎著他走來。這個人好久不穿衣服，不住在屋子裡，而住在墳墓裡。 \end{tabularx} \\ \\ \relax
8:28 & \begin{tabularx}{0.7\textwidth}{X} 他看見耶穌，就喊叫著俯伏在他面前，大聲說：「至高神的兒子耶穌，你為甚麼干擾我？我求你，不要叫我受苦！」 \end{tabularx} \\ \\ \relax
8:29 & \begin{tabularx}{0.7\textwidth}{X} 這是因耶穌曾吩咐污靈從這人身上出來。原來這污靈屢次抓住他；他常被人看守，又被鐵鏈和腳鐐捆鎖，他竟把鎖鏈掙斷，被鬼趕到曠野去。 \end{tabularx} \\ \\ \relax
8:30 & \begin{tabularx}{0.7\textwidth}{X} 耶穌問他：「你的名字叫甚麼？」他說：「群」；這是因為附著他的鬼多。 \end{tabularx} \\ \\ \relax
8:31 & \begin{tabularx}{0.7\textwidth}{X} 鬼就央求耶穌不要命令他們到無底坑裡去。 \end{tabularx} \\ \\ \relax
8:32 & \begin{tabularx}{0.7\textwidth}{X} 那裡有一大群豬正在山坡上吃食，鬼央求耶穌准他們進入豬裡；耶穌准了他們。 \end{tabularx} \\ \\ \relax
8:33 & \begin{tabularx}{0.7\textwidth}{X} 於是鬼從那人出來，進入豬裡，那群豬就闖下山崖，投進湖裡，淹死了。 \end{tabularx} \\ \\ \relax
8:34 & \begin{tabularx}{0.7\textwidth}{X} 放豬的看見這事就逃跑了，去告訴城裡和鄉下的人。 \end{tabularx} \\ \\ \relax
8:35 & \begin{tabularx}{0.7\textwidth}{X} 眾人出來，要看發生了甚麼事；到了耶穌那裡，發現那人坐在耶穌腳前，鬼已離開了他，穿著衣服，神智清醒，他們就害怕。 \end{tabularx} \\ \\ \relax
8:36 & \begin{tabularx}{0.7\textwidth}{X} 看見這事的人把被鬼附的人怎麼得醫治的事告訴他們。 \end{tabularx} \\ \\ \relax
8:37 & \begin{tabularx}{0.7\textwidth}{X} 格拉森周圍地區的人，因為害怕得很，都求耶穌離開他們；耶穌就上船回去了。 \end{tabularx} \\ \\ \relax
8:38 & \begin{tabularx}{0.7\textwidth}{X} 鬼已從身上出去的那人懇求要和耶穌在一起，耶穌卻打發他回去，說： \end{tabularx} \\ \\ \relax
8:39 & \begin{tabularx}{0.7\textwidth}{X} 「你回家去，傳講神為你做了多麼大的事。」他就走遍全城，傳揚耶穌為他做了多麼大的事。 \end{tabularx} \\ \\
[1ex]
\hline
\hline
\end{longtable}
講題:或去或留　To Go or to Stay

經文:路加福音8章26-39節

講員:林豪恩先生

場合:香港中文大學崇基學院禮拜堂

日期:2016年6月19日

一． 引言

或去或留,可能是英國人近日最艱難的抉擇之一.幾天之後(6月23日),英國人將要投票選擇脫離歐盟,還是留在歐盟.我們的財政司司長曾俊華加入評論,認為美國加息與否,也會待英國脫歐公投後才作決定.他還指出,香港政府也留意英國脫歐的影響,已進行一些估計及壓力測試,他認為會間接對環球金融有一定的衝激.

或者或留,也是香港人不時的掙扎.那些年,當我們還是中學生的時候,一年之中難免要唱幾次這首惜別歌:

人生於世上有幾個知己　多少友誼能長存

今日別離共你雙雙兩握手　友誼常在你我心裡

今天且有暫別　他朝也定能聚首　縱使不能會面　始終也是朋友

說有萬里山　隔阻兩地遙　不須見面　心中也知曉　友誼改不了

「今天且有暫別,他朝也定能聚首」,拜WhatsApp所賜,二十多年後,我們那一屆的中學同學,果然重新聚首一Group,那是一個WhatsApp Group.「說有萬里山,隔阻兩地遙,不須見面,心中也知曉」,我們的WhatsApp group每天二十四小時不停有人交談,當身處香港的同學睡覺的時候,歐美的同學就接力,到歐美的同學說Good night的時候,亞洲地區的同學就紛紛「蒲頭」.

更意想不到的是,最近有人突然召集小學同學聚會,所為何事呢?原來有同學在六月中舉家移居離開香港.加拿大移民部最新統計資料顯示,今年首3月抵Po 加拿大的香港移民較去年同期增加85\%,而同期香港人提出的申請人數亦增加17\%.【明報專訊6/6】

此時此地,或去或留這個話題,又再熱鬧起來.

可能,或去或留這個問題,根本就是人生的一部份,伴隨一生.有些人可能在感情關係上掙扎著或去或留,有些人可能在工作上思考著或去或留,有些人甚至在信仰生活上徘徊在去與留之間.

今天福音經課裡的人物,也面對著或去或留的處境.讓我們進入他們的故事,尋找一點啟發,豐富我們關於去留的思考.

二．經文

8 : 26 他們到了格拉森人的地方,就是加利利的對面.

8 : 27 耶穌上了岸,就有城裡一個被鬼附著的人迎面而來.這個人許久不穿衣服,不住房子,只住在墳塋裡.

8 : 28 他見了耶穌,就俯伏在他面前,大聲喊叫,說:「至高神的兒子耶穌,我與你有甚麼相干?求你不要叫我受苦!」

8 : 29 是因耶穌曾吩咐污鬼從那人身上出來.原來這鬼屢次抓住他；他常被人看守,又被鐵鍊和腳鐐捆鎖,他竟把鎖鍊掙斷,被鬼趕到曠野去.

8 : 30 耶穌問他說:「你名叫甚麼?」他說:「我名叫『群』」；這是因為附著他的鬼多.

8 : 31 鬼就央求耶穌,不要吩咐他們到無底坑裡去.

8 : 32 那裡有一大群豬在山上吃食.鬼央求耶穌,准他們進入豬裡去.耶穌准了他們,

8 : 33 鬼就從那人出來,進入豬裡去.於是那群豬闖下山崖,投在湖裡淹死了.

8 : 34 放豬的看見這事就逃跑了,去告訴城裡和鄉下的人.

8 : 35 眾人出來要看是甚麼事；到了耶穌那裡,看見鬼所離開的那人,坐在耶穌腳前,穿著衣服,心裡明白過來,他們就害怕.

8 : 36 看見這事的便將被鬼附著的人怎麼得救告訴他們.

8 : 37 格拉森四圍的人,因為害怕得很,都求耶穌離開他們；耶穌就上船回去了.

8 : 38 鬼所離開的那人懇求和耶穌同在；耶穌卻打發他回去,說:

8 : 39 「你回家去,傳說神為你做了何等大的事.」他就去,滿城裡傳揚耶穌為他做了何等大的事.

三．釋經與應用

1. 離去是為了「被離去」的人

8 : 26 他們到了格拉森人的地方,就是加利利的對面.

上文(8:22)提到,耶穌帶著門徒上船,對門徒說:「我們可以渡到湖那邊去.」那邊,就是格拉林,即是加利利的對面.對面,不單單是地理上的對面,也是文化和宗教上的對面,對這些加利利猶太人來說,對面是一個異域:

8 : 27a 耶穌上了岸,就有城裡一個被鬼附著的人迎面而來.

那裡是一個有污鬼的地方.

8:28b 這個人許久不穿衣服,不住房子,只住在墳塋裡.

那裡是一個墓地.

8 : 32 那裡有一大群豬在山上吃食.

8 : 34 放豬的看見這事就逃跑了

那裡有豬,以及有養豬的人的地方.

按作者描述,這個地方的每一項條件都是不潔的,加起來更加不潔,因此,應該是猶太人避之則吉的地方.那麼多地方不去,耶穌為甚麼帶著門徒離開加利利自己的家鄉,去這陌生的地方沾污自己呢?

而且,這次旅程並不好玩.如果要用玩來形容,這是絕對是「玩命之旅」.不是嗎?且看門徒的經歷:

8:23船行的時候,耶穌睡著了.湖上忽然起了狂風,船將灌滿了水,很危險.

8:24門徒去叫醒他,說:「老師!老師!我們快沒命啦!」

門徒可能不知道為甚麼要去那樣的地方,就算知道,可能也未必同意如此驚險的行程.那異域到底有甚麼吸引耶穌離家而過去呢?

8 : 27 耶穌上了岸,就有城裡一個被鬼附著的人迎面而來.這個人許久不穿衣服,不住房子,只住在墳塋裡.

原來,在那異域,有一個人,這是一個「被離開」的人.為甚麼說這是一個「被離開」的人呢?經文指出,這個人本來是屬於那個城市,但現在卻在墳地裡,因此,他很可能是被他本來屬於那個城市要求離開的.我們可以想像,這個人被城市驅逐,可能流落到鄉村；又被鄉村驅逐,可能流落到田野；又被農民驅逐,最後留落到墳場.

耶穌離開自己的地方,會不會是為了這個「被離開」自己地方的人呢?

這個人不但被要求離開自己的地方,更被逼離開他自己.他的表現,已經不是他原來的自己:

8:28b 這個人許久不穿衣服

即是說,他本來是穿衣服的.

8 : 29 他常被人看守,又被鐵鍊和腳鐐捆鎖,他竟把鎖鍊掙斷

即是說,他本來是自由的,也是有自制能力的.

8 : 28 他見了耶穌,就俯伏在他面前,大聲喊叫,說:「至高神的兒子耶穌,我與你有甚麼相干?求你不要叫我受苦!」

即是說,當時活著的已經不是他,乃是污鬼在他裡面活著,並且奪去了他的自主及自尊.

耶穌離開自己的地方,會不會是為了這個「被離開」自己的人呢?

這是合理的猜想,因為耶穌過去格拉森這個「異域」,所產生的效果就是令這個人找回自己,以及回到自己的地方.

8 : 33 鬼就從那人出來

然後:

8 : 35 眾人出來要看是甚麼事；到了耶穌那裡,看見鬼所離開的那人,坐在耶穌腳前,穿著衣服,心裡明白過來

如果門徒問耶穌:為甚麼離開加利利自己的家鄉,冒著海上的風浪,過去一個外邦不潔的地方呢?耶穌可能回答:離家,是為了「被離家」的人；前去,因為那裡有「被離家者」.

有一位校友在醫院裡做護士.當夜班的日子,她感到最辛苦.每次她要離家上班,她的兩歲兒子都哭個不停.關門前,她安慰兒子不要哭；關門後,她自己忍不住流淚.她等待著兒子有一天會明白,媽媽要離家,因為醫院裡有些人不能回家,要留院治療.離家,是為了那些不能回家的人.

有十多位崇基學生將於下星期五離港前往泰國進行短宣體驗,接待我們的,是海外基督使團,從前稱為中國內地會.中國內地會的出現,是因為一百五十多年前,有一位英國人名叫戴德生(James Hudson Taylor)聽到中國有很多人在天父上帝的大家庭之外,於是他離開英國的家,遠赴東方,希望幫助中國人回家.直到今天,這個宣教團體差派了超過1400位來自全球30多個國家的宣教士在亞洲18個地區,超過100個群體中服事.這些人離家,都是為了協助離家的人回家.

去,總有一個理由.差異不在於所去之處或遠或近,而在於是否帶著使命.

2. 留下,是為了留下的人

或去或留,有時未必是自己的選擇.就正如故事中的格拉森人,他被逼離開自己的城市,「無得留低」.耶穌又如何呢?那個地方的人一樣驅逐他,耶穌都是「無得留低」.

8 : 37 格拉森四圍的人,因為害怕得很,都求耶穌離開他們；

耶穌如何回應呢?經文繼續說:「耶穌就上船回去了.」(8 : 37下)面對被驅逐,耶穌似乎沒有甚麼大反應,一切順其自然.這令人想到 蘇軾的名句:

回首向來蕭瑟處,歸去,也無風雨也無晴.

要去則去,要留則留,沒有甚麼大不了.所謂安然離去,原來「安然」的心境,是可以不被別人的讚譽或毀謗決定的,做了要做的事情,就可以心安理得了.

或去或留,何時會是絕對自主的呢?從前,那位格拉森人被驅逐離開他的城市,「無得留低」；現在,他自己選擇不「留低」:

8 : 38 鬼所離開的那人懇求和耶穌同在；

然而,「耶穌卻打發他回去」(8 : 38)

曾經有人說,離開容易,留下更需要勇氣.對那位格拉森人來說,離開不容易,留下的確更需要勇氣.這個城市對他來說,既熟識,又陌生.或許,他是那些街道,仍是那些建築物,但在這個城市所發生的事情,以及這個城市對待他的態度和方法,都令他感到陌生,也令他感到擔憂和恐懼.有些時候,或去或留,都不是容易的事情.

馬傑偉教授兩個月前在【人生進退】專欄發表了名為「移民避港?心在,不在乎去留」一文,文中有幾段或者也說出了不少香港人的心情:

「那應該是2012年的事了.個人的低潮也反映於書房之外,香港正一步一步掉入困境.現在回想,曾蔭權,湯顯明涉貪只是毒瘡而非毒瘤,接續而來的連番亂象與亂治,才是港人集體抑鬱的開始.自梁振英上任以來,接二連三的明爭暗鬥,幾乎打擊香港每一個引以自豪的開明傳統.法治,自由,人權的底線,一次又一次被衝擊.堂皇官話再掩蓋不住謊言與偽善.幾年前,我們還是每一次香港跌破低價時,都萬分肉緊的盤點損失；如今憾事太多,大家頹坐愁城,疲累不堪,只能招架,守得一時得一時.

2012年底我最感窒息的時侯,朋友提出遠走台灣「避避風頭」……大家都認真起來,有一種逼虎跳牆,「搵條生路給自己」的決絕心態.此地不容人,自有留人處.

那個聖誕,我們入住宜蘭充滿藝文色彩的民宿.午夜無眠,我在眾人睡去的深夜,在幽暗的陽台,外望雨後大片大片的水田,靜水寂寂,倒影雨後雲天.偶然一隻睡醒了的千羽鶴,展開巨翅,劃過水田,觸爪留下一圈一圈無聲漣漪,如此美靜水鄉,怎會不是安息的居所?

我以為,落腳台南,暫離香港,生活可以由灰黑的寒冬,一下子轉到陽光充沛的盛夏.這倒不是妄想,多少移民故事可以證實,在彼邦休養生息,再建新門牆,遠勝舊家鄉.我甚至跟家人駁辯,堅持己見,肯定地說:「就算找不到合適可買下來的農地,也要租一個小小的apartment住下來.」

當然,在台灣不必掛慮「搵食」,才可以享受彼岸風華.要為口奔馳討生活,美夢很快消失.若要經營生意(如咖啡室),你就必須處理當地的複雜「政治」.在台北開咖啡店的港人Raymond說,台灣只是個用繁體字的中國大陸社會.人治,關係,貪污,港人不一定可以受得住.」

馬傑偉教授思量者:「去或留?移不移民台灣?離不離開香港這個傾倒中的城市?」

那位格拉森人主動求去,向耶穌提出申請,可惜,「那機會不屬於他的」.耶穌拒絕帶他走,反而把任務交給他:

8 : 39 「你回家去,傳說神為你做了何等大的事.」

耶穌對那位格拉森人的吩咐,會否太過強他所難呢?這個城市曾經拒絕他,也拒絕耶穌,難道會歡迎他的見證嗎?見證,可能從來都不是錦上添花,而只是在沒有花朵的地方播種:

哪裡有仇恨,就讓我播種仁愛

哪裡有殘害,就讓我播種寬恕

哪裡有猜疑,就讓我播種信仰

哪裡有絕望,就讓我播種希望

哪裡有黑暗,就讓我播種光明

哪裡有憂苦,就讓我播種喜樂　(聖法蘭西斯禱文)

留下,是一種狀況.這種狀況可能是當事人的選擇,也可能是別無他選.然而,以甚麼態度生活在「留下」這種狀況之中,或者是可以選擇的.那位格拉森人選擇了.看看他的活力和熱誠:

8 : 39b 他就去,滿城裡傳揚耶穌為他做了何等大的事.

耶穌離家,是為了那「被離家」的人；那位格拉森人留下,看來卻是為了留下在格拉森的其他人.

馬傑偉教授在剛才引述的那篇文章繼續說:

「2014年夏天,我面對一生人未遇過的危機,幾經掙扎,問題解決了,我重新認識自己.從幽谷走出來,才知道那幾年灰暗得不能自拔,害了自己也苦了家人.老婆大人還願意「陪我癲」,在台灣奔走了一年.心結解開來,心裡最重要的追求,不外乎是家人朋友,不外乎在困局之中保持良知,在醜惡的世界中選擇善良.在香港如是,在台灣如是,在世界遠方那個你我心目中的桃園,也不過如是.今天回看,心靈安頓了,台灣那個計劃變得很遙遠.香港還是一個悶局；但放開心結,此地還是充滿親情與友情.」

馬傑偉教授總結說:「心在,不在乎去留,也不在乎身在何方.」(馬傑偉【人生進退】移民避港?心在,不在乎去留)

四.總結

耶穌和格拉森人的故事,有去有留.或去或留,總有因由.使徒保羅曾在腓立比書說:

1:20沒有一事能使我羞愧；反倒凡事坦然無懼,無論是生是死,總要讓基督在我身上照常顯大.

1:21因為我活著就是基督,死了就有益處.

1:22但是,我在肉身活著,若能有工作的成果,我就不知道該挑選甚麼.

1:23我處在兩難之間:我情願離世與基督同在,因為這是好得無比的；

1:24然而,我為你們肉身活著更加要緊.

1:25既然我這樣深信,就知道仍要留在世間,且與你們眾人一起存留,使你們在所信的道上又長進又喜樂,

願上主幫助我們,或去或留,「總要讓基督在我身上照常顯大」.



\newpage

\section{尋根歷史的叩問}
\label{sec:326}
\textbf{梁元生教授 - 尋根:歷史的叩問}
\newline
\newline
\hyperref[sec:325]{< < < PREV SERMON < < <}
~
\hyperref[sec:toc]{[返目錄]}
~
\hyperref[sec:327]{> > > NEXT SERMON > > >}


路加福音 8:11-15
\newline
\begin{longtable}{cl}
\hline
\hline
章節 & 經文 (和合本修訂版)\\
\hline
8:11 & \begin{tabularx}{0.7\textwidth}{X} 「這比喻是這樣的：種子就是神的道。 \end{tabularx} \\ \\ \relax
8:12 & \begin{tabularx}{0.7\textwidth}{X} 那些在路旁的，就是人聽了道，隨後魔鬼來，從他們心裡把道奪去，以免他們信了得救。 \end{tabularx} \\ \\ \relax
8:13 & \begin{tabularx}{0.7\textwidth}{X} 那些在磐石上的，就是人聽道，歡喜領受，但沒有根，不過暫時相信，等到碰上試煉就退後了。 \end{tabularx} \\ \\ \relax
8:14 & \begin{tabularx}{0.7\textwidth}{X} 那落在荊棘裡的，就是人聽了道，走開以後，被今生的憂慮、錢財、宴樂擠住了，結不出成熟的子粒來。 \end{tabularx} \\ \\ \relax
8:15 & \begin{tabularx}{0.7\textwidth}{X} 那落在好土裡的，就是人聽了道，並用純真善良的心持守它，耐心等候結果實。」 \end{tabularx} \\ \\
[1ex]
\hline
\hline
\end{longtable}
講題:尋根:歷史的叩問　In Search of Roots:Knowing our Past

經文:使徒行傳21:37-40,22:1-5;路加福音8:11-15

講員:梁元生教授

場合:香港中文大學崇基學院禮拜堂

日期:2016年6月26日

引言

今天要和大家思考的題目是“尋根”,意思是尋找歷史,尋找自身的歷史.緣起是校牧許開明牧師帶領的“尋根團”正在華東地區訪問尋根.許牧師和團友們大概明天才會回港,原來他想邀請崇基神學院院長邢福增教授作為今次主日崇拜講員,可惜邢教授不能抽身,故此問我能否擔此重任.臨危受命或隨時替補,以往在伍渭文牧師擔任校牧時期,我也是常常會做的.時間雖短,我也欣然答應了.我要講的題目也是從許牧師這次帶領的“尋根團”或“尋根之旅”變化出來的.我期待許牧師會在不久之將來向我們報告這個“尋根之旅”的經過和感受.今天我只是想借著這次“教會大學尋根之旅”作為引子來談談“尋根”對基督徒的意義.今次的“教會大學尋根之旅”由崇基校牧帶同一批崇基人--崇基校友,崇基教員,崇基學生及崇基禮拜堂會友,前往蘇州,南京和上海去考察,主要是重訪往昔幾間教會大學的遺址,包括以前是東吳大學的蘇州大學,以前是金陵大學的南京大學,前是金陵女子大學的南京師範大學,他們到了上海,也希望多了姐一下以前是上海聖約翰大學和上海滬江大學而現在是華東政法大學及華東理工大學的一些情況.他們訪問這些大學,尋找舊時教會大學的遺跡,主要的原因是因為崇基學院和這些教會大學有著很深的淵源和發展的關係.一九五一年崇基學院在香港創辦,至今已有六十五年.但它的開辦,是和整個中國動盪的政局以及大學和高教的改組息息相關的.四九年新中國成立,共產黨執政,對西方國家,對資本主義,對傳統文化,對基督教都嚴厲批評和表示反對.原來由基督教創辦的十三所教會大學都面臨解散或改組的命運.許多大學後來併入國立或省立的大學,以新的面貌重新出現,如燕京大學併入了北大,嶺南大學併入了中山大學,金陵大學改組成為南京大學, 金陵女子學院改為南京師範大學等等.

而一些教會大學的校董,教授,老師和職工,在那革命和動亂的年代,也只有離校他往,另覓生計,或者遷居港臺,以避共產統治.崇基學院就是在這樣的背景下創辦的.三位主要的創辦人:當時的聖公會主教何明華會督(Bishop Ronald Hall)是其中之一,其餘兩人:崇基創校校長李應林博士,是廣州嶺南大學解放前最後的一位校長,崇基校董會首任主席歐偉國,曾經肆業於上海聖約翰大學,是上海先施公司總經理歐彬之子, 回到香港之前曾經擔任上海聖約翰大學的校董.1949年後崇基學院在香港的籌備和成立,是這三位元老和其他的同道人,特別是香港基督教會領袖,協力合作的成果.許開明校牧和其他崇基人前往上海,蘇州和南京去考察和尋訪,是希望透過歷史記憶,透過重新認定崇基的“機構身份”(institutional identity)或“團體身份”(corporate identity)的過程, 去肯定以往的中國教會大學對崇基學院的影響和貢獻,明瞭崇基歷史的源頭和根由.一般認為,如果我們對一家公司,一所學校,或一個團體的歷史知道得越多,凝聚力和向心力就會越強；故此,現代許多大企業也都會設立公司檔案部門,希望員工從歷史和過去中培養對公司或企業的忠誠和熱愛.我們回望過去,叩問歷史,為的是從過去中明白現在的由來,對自己的身份多一些瞭解.然而,作為基督徒,我們卻常常對回望過去和叩問歷史存著一種懷疑的態度, 甚至有些學者說基督教有一種“非歷史”(ahistorical),甚至是“反歷史”(anti-historical)的看法.

誤解的由來

基督教是“非歷史”或“反歷史”的嗎?當然不是.不然,從基督教角度寫的歷

史書何其多,或者單單寫基督教歷史的書也是汗牛充棟,多如牛毛.你只要到崇基牟路思怡圖書館去轉一轉,就知我所言非虛.問題是:許多基督教開辦的教會學校, 甚至許多基督教會,大都不知其自身的歷史,大部分信徒也對其身處的學校或教會的由來,也所知甚少,或一概不知.除了缺少歷史感,也少有探問歷史的精神.對個人,對群體,對單位,對機構的過去,沒有探知的衝動.

這種態度可能源於對基督教神學的一點誤解.在新約聖經,特別是保羅書信中,明顯地強調下面的幾個資訊:

首先,信徒對信仰生活的態度,是“向前走”和“莫回頭”的: 在路加福音第 九章 六節中, 路加記下耶穌基督說過的這句話:“凡是手拿著犁往後看的, 不配進上帝的國.”保羅在寫信給以弗所及腓立比教會的基督徒時,也多次強調:“要忘記過去,努力面前,向著標杆直跑,為要得著前面的獎賞.”再者,保羅又用了“老我“和“新我”(或“舊人”和“新人”)的比喻,說明“悔改”與“重生”對基督徒人生的重要性.保羅說:“要脫去舊人,穿上新人”, 生命因而改變,“舊事已過,都變成新的了!”過去,舊事,歷史-似乎不值得留戀,基督徒的人生,應該是迎接新生,努力前行,不由舊事舊人拖累的.這種對立性的思維,把“新與舊”,“過去和現在”作出截然二分的對立性看法,做成了一般基督徒對過去,對歷史,等同或相近於 “舊人”或“老我”,從而引起一種負面的或冷淡的態度,不願探知,也不願處理,只求棄舊而迎新,只想前瞻而不屑回顧.在我們學術圈中,曾經對“現代化”的概念及其與“傳統”的關係,也有過同樣的糾結.有一陣子幾乎所有研究現代化的學者都把“現代”和“傳統”置於對立面,認為傳統是負累,需要拋開,需要割切,需要揚棄,才可以走向現代,走向未來.但到了今時今日,學術圈內已經有了很大的改變,認為傳統和現代不應對立,也不需割切,兩者是立足在同一條線上的,時間上有前後,階段上有次序,但二者仍有緊密的關係,傳統中帶著現代,現代中留有傳統,二者是“發展”而非“對立”的關係.就如同我們的生命,它的過去,現在和將來,也是在同一條發展線上,環環相扣,節節相連,沒有已往,也就沒有今日,遑論未來.今日讓我們重新檢視我們對歷史的態度,對往昔和將來,對傳統和現代,對新人和舊人的關係,作一反思.

從保羅的歷史說起

如果基督教神學真的有一點“非歷史”或“反歷史”的傾向,而那傾向是受了保羅的

的影響,那麼,容許我們從保羅自身的歷史看看他對這個問題的看法. 的確,保羅十分重視“悔改”和“重生”的經驗,也鼓勵信徒“忘記過去,努力面前”,“脫去舊人,穿上新人”.然而,他從來沒有忘記自身的歷史.他對自己的出身,自己的過去,自己的往跡,自己的家族, 自己的身份,都有著深刻的認識和瞭解.正因為他對自己的過去有深刻的認識,才會對自己的悔改與更新,對自己的選擇和委身,有無悔的勇氣,有理性的解釋,有極大的承擔.

我們知道:保羅是個偉大的使徒和宣教師,是初期教會和基督教神學的奠基者. 一般人講他的生命和工作,都會從大馬色(大馬士革)路上的遭遇和改變為開端:以往的事是掃羅的事,新的一切從保羅的瞎眼和重光開始.但保羅真的撇掉過去的一切嗎?他真的完全拋開自身的歷史嗎?不是的.

讓我們聽保羅說自身的歷史:他說:“我原是猶太人,生在基利家的大數,長在這城裡,在迦馬列的門下,按著我們祖宗嚴謹的律法受教,熱心事奉神,像你們今日一樣.”(徒22:3)在另外一處,他又提到:他是法利賽人,而且是便雅憫族人,父親是是羅馬人, 很可能是羅馬的公民和官員.

保羅並不強調他的優秀學歷和文化背景.但他對自身的歷史知悉甚詳,對家族的來源也有深厚的認識.從他的自述,我們可以清楚認識到保羅的背景:

他是便雅憫支派的猶太人,如果你對猶太民族略有認識的話,你會知道便雅憫支派源於雅各十二個兒子中的幼子,是約瑟的弟弟,是以色列最為疼愛的.十二支派中,最受寵愛和器重.

他生於大數, 大數是位於敘利亞濱海的大城；保羅少時在這城市長大,至少生活過六年之久,再到耶路撒冷受教育.在大數他會見到來往的商旅,地中海沿岸的不同民族,不同文化背景的人,故此他可以說希伯來話,也能用當時流行的亞蘭語和希臘文；到了耶路撒冷後,他受教於猶太名師迦伯列,學的是經典和法律.在猶太人中他是高學識一族,是法利賽人,相信神跡,相信死後復活,而不是撒都該派的.他年輕時滿懷基要主義的宗教激情,逼害他認為新教或異教的基督徒,結果是在追逐基督徒時在大馬色的路上遇到大光和主的使者,因而眼睛失明,生命改變,成為主的跟隨者,後來更成為初期教會中最出色的神學家和宣教師.

“尋根”的思考

我們尋找保羅的家世,叩問保羅的歷史,為的是從保羅的“根”談到纏繞我們當代基督徒面對的幾個大問題:這幾個問題包括:

血緣與民族:

從血緣或血脈的根來看保羅家世,毫無疑問,他是個正統的具備支派優勢的猶太人.他又接受最嚴謹的宗教信仰的薰陶和訓練,對種族和傳統有深刻的認知和瞭解,能夠用希伯來話講道和誦讀摩西五經, 連當時在亞歷山大的著名猶太學者斐羅(Philo)也無法做到的,所以說保羅是個學童純正,背負道統,如假包換的猶太人.從然在他接受基督之後,猶太的原教旨主義者皆認為他已經背棄了猶太的正統,但他並沒有離棄他的民族,反之,他對他的民族有真正的承擔, 努力讓他們認識基督彌賽亞與其民族傳統的關係.最為重要的是,這個血統純粹和受到傳統深刻訓練的猶太人,卻成了為主在外邦人中工作最得力的宣教師,顯出上帝並不把“正統”和“外邦”放置在一個對立面上.請告訴我,誰是猶太人?誰是便雅憫的後代?誰是希伯來文化的傳人?請告訴我,誰是中國人,誰是皇帝的後裔?請告訴我,誰是香港人?五千年來你的家是否一直居住在新界?請告訴我,悄悄地告訴我,不要喧嘩.

社會和國家:

從家國情懷或國族的根去看保羅的國民身份,情況就有點曖昧.因為在他出生之時,猶太人已經沒有自己的國家.它環保持的是分散在各地的猶太人社會或社區.至於保羅,他的國民身份是羅馬人,羅馬的公民.他的父親曾經擔任過羅馬的官員,而這個身份使到保羅後來在被關押和審訊時得到一些優待:羅馬公民的身份是不少被統治的猶太人所欣羡的,正如香港在長期英國殖民地統治的時代,也有許多中國籍的港人用盡方法,希望入籍成為英國公民一樣.

文化和教育:

保羅是羅馬帝國的公民,但不以自己是猶太人為恥,他的文化的根,牢牢地建立在猶太文化的土壤上,身為羅馬公民,但他反而努力學習猶太人的文化,篤信猶太人的宗教,成為一名有教養,有學識,有抱負的人.他對語言的運用(希伯來,希臘,亞蘭)有傑出的能力,他對法律,條文,規矩,傳統,又嚴格的訓練和深刻的認識,口才便給,論辯滔滔,是耶路撒冷名師迦伯烈的弟子.在香港,我們也有同樣的機會,從大數到耶路撒冷,從香港到另一個大城(倫敦,新加坡,北京,上海)去學習,去接受嚴謹的文化和學術訓練,為的是充實自己,豐富學識,加強能力等,都是可以的.然而,五倫國籍為何,也可以保持自己中國人,甚至香港人的身份,成為自己文化的代言人,像保羅一樣.

信仰與生命:

尋根,保羅從便雅憫支派和猶太人的歷史,而我們則從中國個各民族源流中去尋找我們的血脈之根；,保羅的文化之根, 可以從大數,從耶路撒冷,從迦伯列的教誨去知悉,而我們的文化,則可以從中國的儒家傳統,也可以從香港的社會教育等方面去瞭解一二,但尋根中至為重要者,是尋找生命的根,即是找到可以委身和投放生命的信仰.保羅在大馬色的路上找到了.從此生命有根,不再是浮萍無寄,也不是永遠在追逐的途上.路加福音八章那裡,耶穌說過一個大家熟悉的比喻,就是撒種的比喻.我讀到八個比喻之時,最為震撼的是一句話:“沒有根.”沒有根,神的道也就不會生長,也不會結實.生命沒有根,沒有基督在裡面,一切的在外面的尋找, 包括在血脈中,在文化中及在歷史中的尋找,最終也歸徒然.生命之根,應是我們人生中最首要的選擇.



\newpage

\section{生命同行的鄰舍}
\label{sec:327}
\textbf{鄭漢文博士 - 生命同行的鄰舍}
\newline
\newline
\hyperref[sec:326]{< < < PREV SERMON < < <}
~
\hyperref[sec:toc]{[返目錄]}
~
\hyperref[sec:328]{> > > NEXT SERMON > > >}


路加福音 10:25-37
\newline
\begin{longtable}{cl}
\hline
\hline
章節 & 經文 (和合本修訂版)\\
\hline
10:25 & \begin{tabularx}{0.7\textwidth}{X} 有一個律法師起來試探耶穌，說：「老師！我該做甚麼才可以承受永生？」 \end{tabularx} \\ \\ \relax
10:26 & \begin{tabularx}{0.7\textwidth}{X} 耶穌對他說：「律法上寫的是甚麼？你是怎樣念的呢？」 \end{tabularx} \\ \\ \relax
10:27 & \begin{tabularx}{0.7\textwidth}{X} 他回答說：「你要盡心、盡性、盡力、盡意愛主—你的神，又要愛鄰如己。」 \end{tabularx} \\ \\ \relax
10:28 & \begin{tabularx}{0.7\textwidth}{X} 耶穌對他說：「你回答得正確，你這樣做就會得永生。」 \end{tabularx} \\ \\ \relax
10:29 & \begin{tabularx}{0.7\textwidth}{X} 那人要證明自己有理，就對耶穌說：「誰是我的鄰舍呢？」 \end{tabularx} \\ \\ \relax
10:30 & \begin{tabularx}{0.7\textwidth}{X} 耶穌回答：「有一個人從耶路撒冷下耶利哥去，落在強盜手中。他們剝去他的衣裳，把他打個半死，丟下他走了。 \end{tabularx} \\ \\ \relax
10:31 & \begin{tabularx}{0.7\textwidth}{X} 偶然有一個祭司從那條路下來，看見他就從另一邊過去了。 \end{tabularx} \\ \\ \relax
10:32 & \begin{tabularx}{0.7\textwidth}{X} 又有一個利未人來到那裡，看見他，也照樣從另一邊過去了。 \end{tabularx} \\ \\ \relax
10:33 & \begin{tabularx}{0.7\textwidth}{X} 可是，有一個撒瑪利亞人路過那裡，看見他就動了慈心， \end{tabularx} \\ \\ \relax
10:34 & \begin{tabularx}{0.7\textwidth}{X} 上前用油和酒倒在他的傷處，包裹好了，扶他騎上自己的牲口，帶他到旅店裡去，照應他。 \end{tabularx} \\ \\ \relax
10:35 & \begin{tabularx}{0.7\textwidth}{X} 第二天，他拿出兩個銀幣來，交給店主，說：『請你照應他，額外的費用，我回來時會還你。』 \end{tabularx} \\ \\ \relax
10:36 & \begin{tabularx}{0.7\textwidth}{X} 你想，這三個人哪一個是落在強盜手中那人的鄰舍呢？」 \end{tabularx} \\ \\ \relax
10:37 & \begin{tabularx}{0.7\textwidth}{X} 他說：「是憐憫他的。」耶穌對他說：「你去，照樣做吧！」 \end{tabularx} \\ \\
[1ex]
\hline
\hline
\end{longtable}
講題:生命同行的鄰舍　Life CoWalking Neighbour

經文:路加福音10章25-37節

講員:鄭漢文博士

場合:香港中文大學崇基學院禮拜堂

日期:2016年7月10日

1. 撒馬利亞人的慈心與自義

各位弟兄姊妹,平安.

今日,我又再回鄉探親一樣,返回崇基禮拜堂這個信仰大家庭.

我的家人,學生和好友,今天都共聚在此,盼望聖道的臨在.

讓我先以一個自以為是的好撒馬利亞人故事作為開始,故事是這樣的:

「從前有一個人,落在強盜手中,被剝去衣裳打個半死.有一個撒馬利亞人經過,看見那人,就動了慈心,上前替他包裹傷處,還扶他帶到店裡照顧他,第二天還替他結帳.

然後,到了第三天,那人又落在強盜手中,強盜又把他打個半死,好撒馬利亞人又經過,動了慈心出手幫助.

如是者,每天不斷重覆「好撒瑪利亞人」的故事.

強盜每天出動,那人每天被害,好撒馬利亞人每天動慈心出手幫助.

不過,後來發現,原來好撒馬利亞人正住在強盜的隔壁--好撒馬利亞人每一天都目睹強盜出門搶劫,卻因著「非暴力」與「不取代上帝工作」的理由沒有上前阻止,並繼續每天事後扮演「好鄰舍」的角色.」

以上的故事,是引述自建道神學院神學系助理教授陳韋安教授(2016年6月30日)寫的〈 好撒瑪利亞人的暴力〉這篇文章的開首,究竟今天我們基督徒是如何回應城中的挑戰?

2.　好撒馬利亞人的喻道故事:路加福音10:23-37解經

今天讓我們看一看,耶穌當年是如何面對挑戰?

今日經課是取自路加福音10章25至37節.經文明顯分為三段:

第一段25至29節,是:耶穌面對別人挑戰的處境

第二段30至35節,是:耶穌講了一個故事作比喻

第三段36至37節,是:耶穌引伸了一個教導作回應

讓我們一同先誦讀這段經文.

第一段25至29節,是:耶穌面對律法師探問永生這個處境的挑戰.先請弟兄一起讀出:

25有一個律法師起來試探耶穌,說:「老師!我該做甚麼才可以承受永生?」

26耶穌對他說:「律法上寫的是甚麼?你是怎樣念的呢?」

27他回答說:「你要盡心,盡性,盡力,盡意愛主-你的上帝,又要愛鄰如己.」

28耶穌對他說:「你回答得正確,你這樣做就會得永生.」

29那人要證明自己有理,就對耶穌說:「誰是我的鄰舍呢?」

第二段30至35節,是:耶穌講了一個故事作比喻,後世稱為好撒馬利亞人.請姐妹一起讀出:

30耶穌回答:「有一個人從耶路撒冷下耶利哥去,落在強盜手中.他們剝去他的衣裳,把他打個半死,丟下他走了.

31偶然有一個祭司從那條路下來,看見他就從另一邊過去了.

32又有一個利未人來到那裡,看見他,也照樣從另一邊過去了.

33可是,有一個撒瑪利亞人路過那裡,看見他就動了慈心,

34上前用油和酒倒在他的傷處,包裹好了,扶他騎上自己的牲口,帶他到旅店裡去,照應他.

35第二天,他拿出兩個銀幣來,交給店主,說:『請你照應他,額外的費用,我回來時會還你.』

第三段36至37節,是:耶穌引伸了一個教導作回應,重新詮釋愛鄰如己.讓我讀出:

36你想,這三個人哪一個是落在強盜手中那人的鄰舍呢?」

37他說:「是憐憫他的.」耶穌對他說:「你去,照樣做吧!」

這個好撒馬利亞人的故事家傳戶曉,耳熟能詳.

不單成為教會教導的經典,也成為公民社會中慈惠人士的一個象徵.

我還有甚麼可以講呢?

例如,三年前林豪恩先生就以〈好撒馬利亞人是如何鍊成的?〉作題目在這裡崇基禮拜堂講道.我重新多次聽了他那次講道,參考了他的講章,刻意地要有所不同.

我還可以講些甚麼呢?

3. 回首往來四十年:生命同行的啟示

為了準備這一次講道,為了找尋一個推陳出新的觀點,我回顧了我從四十年前信主到近年接受神學訓練.當中有些心得,得出三個重點.

四十年前的七月中,我在中四升中五的學校暑期福音營中信耶穌,今天剛好信主四十年.

第一個重點是:毋忘初心.

回想我信主的那一剎那間,感到上帝親自來找尋我,從此得到上主與我同行,四十年來,經歷過三次學業低谷,三次事業低谷,三次健康低谷,兩年前經歷三合一的第十次低谷,上主依然不離不棄,與我生命同行.

1980年,我考入中大讀書,一年後,愛上了一位愛好文學的少女,她也是位基督徒,大家相戀五年後的1986年,我們在這所崇基禮拜堂進行了結婚祝福禮,育有兩兒,至今「夫婦同行三十年」.今天下午的讀書會,我們會再分享.

近年我在崇基學院神學院接受神學訓練時,我的心志是結合講學,講笑與講道這三者作服事.而我的講道學是學自三位老師,分別是幫我洗禮的李炳光牧師,教我宣講福音邀請不可推的吳振智牧師,指導我在崇基禮拜堂神學生實習的伍渭文牧師.

每次講道,我都會重溫三位老師所教,心中所學,與我生命同行的種種美好回憶.

李炳光牧師講道,起題精要,旁徵博引,文采飛揚.

十分記得課上他說過一個笑話:許多次講道完畢,教友會向他致謝並表示得益良多.但當李牧師追問他們聽到甚麼時,他們又說不出所以言來.所以,李牧師教導我們,要把講道歸結為讓聽道者縈繞在心的一句話.

至於吳振智牧師,他的講道法,是起於釋經,終於基督,言必歸十架.換言之,講道必須由經文的釋道開始,最後要把聽道者引領回到基督,從alpha到omega,都言必歸向十架精神.

伍渭文牧師,大家都十分認識.他教導我三代經課講道,聖道必須通過經文聖言而呈現,避免講述可能純是個人偏見的所謂經歷,但最後聽道者必須有信仰的應用.

珠玉在前,一位一位在我生命比鄰的同行者,伴我上路,伴我講道.

我今天在以上三種傳統下,添加了我的「敘事想像」的方法(method of narrative imagination),希望有所推陳出新.

「生命同行」(Life CoWalk),就是我今天解釋這段經文的總念.

讓上主的鄰憫,通過我們聖靈感動我們的心回應受困苦者,與他們的生命同行一段路,以體現愛鄰如己,這就是永生生命的存活,這就是受苦的主耶穌的基督精神──苦者尤比鄰.

4. 生命同行的個體性:以受困苦者的心為心

路加福音第10章這段經文是順序的,現在讓我們先把經文的次序倒轉讀,就是從耶穌其實想教導我們甚麼的向度重讀.

耶穌以故事回應挑戰的最後反問律法師:

36你想,這三個人哪一個是落在強盜手中那人的鄰舍呢?」

37他說:「是憐憫他的.」

耶穌逆轉了律法師的問法,由「恩惠慈善的施予者」(charitable)的角度,轉為「蒙恩受惠的受苦者」(sufferer)的本位,來詮釋愛鄰如己.

為讓大家可以生發同理心,代入案主當事人的視點,讓我們把故事轉為第一身的自述,好似警局報案一樣,受害者這樣落案:

「我有一天從耶路撒冷下耶利哥去,忽然有幾個強盜跳出來手中.他們剝去我的衣裳,把我打個半死,丟下我走了.

偶然有一個祭司從那條路下來,看見我就從另一邊過去了.(我心好痛:乜佢係神職人員嗎?仲咩唔幫我?)

又有一個利未人來到那裡,看見我,也照樣從另一邊過去了.(我心好痛:乜佢係公職人員嗎?仲咩唔幫我?)

可是,有一個撒瑪利亞人路過那裡,看見我就動了慈心,他上前用油和酒倒在我的傷處,包裹好了,扶我騎上自己的牲口,帶我到旅店裡去,照應我.

第二天,他拿出兩個銀幣來,交給店主,叮囑店主,請他照應我,還承諾,額外的費用,回來時會還給店主.」

如果案主是你,誰是與你生命同行的鄰舍?

就是能明白你的困苦,伴你生命同行一段苦路的同行者.

你能夠保證生命不會遭逢不測,被某一種意義的強盜所搶劫?

讓我們現在靜下來,低頭,閉目,回想人生:人生中曾幾何時,你生命遇上傷害你的人時,又有誰在你身旁,給你無限憐憫,以你的困苦的心為心?

憐憫的憐,就是心字旁,與鄰為伴.

你我都可以成為受困苦者的鄰舍,好好保守你可以對受苦者動慈心.

5.　生命同行的群體性:杜絕強盜的惡行

但是,以上所講,只停留在一種個人化的進退取捨.故事太快終結於:做個好人.

讓我們進深一點,故事可不可以回到今日講道的開首:就是杜絕強盜的惡行.如果及早杜絕強盜的惡行,就不會出現受苦的案主.

剛才講道開頭,我引用了陳韋安教授寫的〈 好撒瑪利亞人的暴力〉這篇文章的喻道故事,究竟今天我們基督徒是如何回應城中的挑戰?

「從前有一個人,落在強盜手中,被剝去衣裳打個半死.有一個撒馬利亞人經過,看見那人,就動了慈心,上前替他包裹傷處,還扶他帶到店裡照顧他,第二天還替他結帳.

然後,到了第三天,那人又落在強盜手中,強盜又把他打個半死,好撒馬利亞人又經過,動了慈心出手幫助.如是者,每天不斷重覆「好撒瑪利亞人」的故事.

強盜每天出動,那人每天被害,好撒馬利亞人每天動慈心出手幫助.

不過,後來發現,原來好撒馬利亞人正住在強盜的隔壁--好撒馬利亞人每一天都目睹強盜出門搶劫,卻因著「非暴力」與「不取代上帝工作」的理由沒有上前阻止,並繼續每天事後扮演「好鄰舍」的角色.」

陳韋安教授這樣評述:

「今日的香港,實體的強盜「不存在」--因為強盜無處不在,強盜成為一個客體,強盜成了一個結構.因此,教會活在強盜系統中,唯有從這系統中醒覺,對抗這客體,才能成為真正的「好撒瑪利亞人」.不然,不停在結構暴力中貼膠布只是一種「好撒瑪利亞人的暴力」.不是嗎?一邊弄傷你的傷口,一邊為你貼膠布.很變態的啊.」

所以,生命同行的鄰舍,不單在個人與個人的交往上,還要在社群的倫理結構上改變.

例如,如何杜絕強盜的惡行?

例如,如何建設收容受困苦者的旅店?

讓我們再把路加福音第十章耶穌講的這個故事本土化:

「有一天有一個香港人從深水Po 搭地鐵返工到中環,忽然有幾個強盜在中環跳出來.

他們當中,有交通公司濫收車費,有老闆剝削他的工資,有上司要佢工作超時到半死,最後丟下退休無保障的他走了.(強盜當中,有人把人移民去了,或有人把資產移到別國去了.)

偶然有一個貌似行政長官的人從那條路下來,看見他,就從另一邊過去了.(都唔係我有份民選出來的,係無咩良心!)

又有一個好像高級的公務員來到那裡,看見他,也照樣從另一邊過去了.(都係超離地的所謂公僕,係唔會明白民間疾苦!)

可是,有一個流浪漢路過那裡,看見他就動了慈心,他上前用油和酒倒在他的傷處,包裹好了,扶他坐上自己的手推車,帶他到教會裡去,照應他.

第二天,他拿出兩個銀幣來,交給牧師,叮囑牧師,請他照應我,還承諾,額外的費用,回來時會還給教會.」

在香港,如何杜絕強盜的惡行?如何形成鄰舍守望(Neighborhood Watch),守望受困苦者不再被強盜搶奪.

在香港,如何建設收容受困苦者的旅店?如何形成受困苦的安居(Resting Place),得以養傷復康.

在香港,信徒做了甚麼?教會又做了甚麼?

我們當然要有好撒馬利亞人的慈心,但很容易流於自義,成為自以為是的「港豬」.

我們還要建立守望鄰舍的倫理,杜絕強盜的惡行與權勢.

今次的講道,在這裡就此打住,再多講,又會出事.

在最近這兩三年間,香港的公共領域十分動盪,我自己也疲於奔命參與其中.

有一回被問及我的政治取態,追問我代表甚麼party,我答曰:Roger represents suffering party.

我鄭漢文,若要代表,就是代表受困苦者這苦群(Suffering Party).

四十年追隨基督,最後明白耶穌精神就是以受苦者的心為心,英文叫:option for the suffering.

讓我們一起低頭禱告:

「親愛的上主啊,

求你教我們有以受苦者的心為心,學習成為「生命同行」的鄰舍.

求你教我們如何杜絕強盜的惡行?建設收容受困苦者的旅店?

求上主讓你的鄰憫,通過們聖靈感動我們的心回應受困苦者,與他們的生命同行一段路,以體現愛鄰如己,這就是永生生命的存活,這就是受苦的主耶穌教導我們的基督精神──苦者尤比鄰.」

但願榮耀都歸給聖父,聖子,聖靈,永在的三一神.阿們.

(鄭漢文博士現職香港中文大學教育行政與政策學系專業顧問,

同時兼任崇基學院神學院優質生命教育中心名譽副主任)



\newpage

\section{馬大世界馬利亞心}
\label{sec:328}
\textbf{林豪恩先生 - 馬大世界馬利亞心}
\newline
\newline
\hyperref[sec:327]{< < < PREV SERMON < < <}
~
\hyperref[sec:toc]{[返目錄]}
~
\hyperref[sec:329]{> > > NEXT SERMON > > >}


路加福音 10:38-42
\newline
\begin{longtable}{cl}
\hline
\hline
章節 & 經文 (和合本修訂版)\\
\hline
10:38 & \begin{tabularx}{0.7\textwidth}{X} 他們繼續前行，耶穌進了一個村莊。有一個女人，名叫馬大，接他到自己家裡。 \end{tabularx} \\ \\ \relax
10:39 & \begin{tabularx}{0.7\textwidth}{X} 她有一個妹妹，名叫馬利亞，在主的腳前坐著聽他的道。 \end{tabularx} \\ \\ \relax
10:40 & \begin{tabularx}{0.7\textwidth}{X} 馬大伺候的事多，心裡忙亂，進前來，說：「主啊，我的妹妹留下我一個人伺候，你不在意嗎？請吩咐她來幫助我。」 \end{tabularx} \\ \\ \relax
10:41 & \begin{tabularx}{0.7\textwidth}{X} 主回答說：「馬大，馬大，你為許多的事操心煩惱， \end{tabularx} \\ \\ \relax
10:42 & \begin{tabularx}{0.7\textwidth}{X} 但是不可少的只有一件。馬利亞已經選擇了那上好的福分，是沒有人能從她奪去的。」 \end{tabularx} \\ \\
[1ex]
\hline
\hline
\end{longtable}
講題:馬大世界馬利亞心　Having a Mary Heart in a Martha World

經文:路加福音10章38-42節

講員:林豪恩先生

場合:香港中文大學崇基學院禮拜堂

日期:2016年7月17日

一. 引言

基督徒作者韋喬安(Joanna Weaver)以馬大,馬利亞和拉撒路這三姐弟的家庭寫了三本書,其中一本名為《馬大世界馬利亞心》(Having a Mary Heart in a Martha World),此書所採用的就是今天的福音經課--路加福音10章38-42節的故事.

作者以馬大和馬利亞的性格差異來解讀這個故事:

「馬利亞如果像一道陽光,馬大就像一聲雷鳴；如果馬大是車頭,馬利亞就是跟在後頭的車廂.馬利亞能夠優遊自得地漫步在生命的花園中,為了嗅聞一朵玫瑰而駐足留連；至於馬大,她會明快地將玫瑰花採下,俐落地修剪花枝,然後搭配著滿天星和蕨葉插在花瓶裡.」

以性格差異來解讀,很自然會得出以下的結論:

「這是兩種不同的性格表現,沒有誰對誰錯.我們每個人都不一樣,神造我們就是這樣.每一種性格,每一種天賦,都有它的過人之處和弱點,它們會獲得應有的讚賞,也會遭遇不同的試探……這意味著馬利亞和馬大都可以獲得這上好的福分,也意味著不論我們的個性和天賦為何,每個人都可以獲得這上好的福分,就看我們怎麼抉擇.」

作者提出證據:

「當耶穌糾正馬大的錯誤時,祂並不是說:「為什麼妳不能學學妳妹妹馬利亞呢?」祂知道馬大永遠不會變成馬利亞,馬利亞也永遠不會變成馬大.不過,當她們面對同一個選擇--該去工作還是去崇拜?-耶穌的回應是:「馬利亞選擇了上好的福分.」

基於對這個故事的解讀,作者提出以下的教導和應用:

「今天,在教會裡多的是馬大這樣的人,但少有像馬利亞這樣的人!我們的生命裡,充滿了太多次要的事情,心中塞滿了沒必要的事物!甚至連我們對主耶穌的服事,也被一大堆的事擠滿了,連主耶穌本身都容不下!信徒在世,應當努力工作,完成神所託付的責任──家庭,社會,教會,三者同等重要,必須平衡兼顧,不可偏廢.但在這一切之先,應和主耶穌建立個人,親密,健全的關係.每天在工作盡責任之前,先專心安靜親近敬拜主耶穌,渴慕得著祂話語的滋潤和飽足,享受和祂同在的甜美喜樂平安,明白祂的心意,從祂支取恩典能力!如此,在照顧家庭,社會工作,教會事奉上,都必蒙福,工作有果效,且能榮神益人.」

本書其中一句最能解釋本書名稱,也是本書的主旨所在:

「祂可以讓我們生活在馬大的世界中,卻賜我們馬利亞的內心.」

作者在本書中的提醒,的確對我們很有幫助.我們都生活在馬大的世界中,正如作者所說,無論是工作,家庭和教會,都諸事繁多,忙個不停.我們的確需要從上主而來的平安以安定我們的心,從上主而來的力量來應付沒完沒了的工作.

本主日的經課編排,提示我們可以從另一個角度來閱讀這個故事.今天的舊約講述亞伯拉罕接待上主的使者(或者是上主自己)到訪的故事；詩篇的回應是與上主同住的人的特質；書信闡釋上主道成肉身的神學意義；而福音書正正是道成內身的上主--耶穌基督到訪一個家庭的具體故事.由此觀之,這個馬大和馬利亞的故事,是一個關於有客到訪及接待客人的故事,更是一個關於耶穌到訪及接待耶穌作為客人的故事.上主道成肉身臨到世界,世界可否仍照原來的樣子去迎接祂呢?道成肉身的耶穌臨到我們,我們適合否仍活在原來的世界中呢?耶穌親臨,我們應該如何接待祂呢?

二. 經文

10:38他們繼續前行,耶穌進了一個村莊.有一個女人,名叫馬大,接他到自己家裡.

10:39她有一個妹妹,名叫馬利亞,在主的腳前坐著聽他的道.

10:40馬大伺候的事多,心裡忙亂,進前來,說:「主啊,我的妹妹留下我一個人伺候,你不在意嗎?請吩咐她來幫助我.」

10:41主回答說:「馬大,馬大,你為許多的事操心煩惱,

10:42但是不可少的只有一件.馬利亞已經選擇了那上好的福分,是沒有人能從她奪去的.」

三. 釋經與應用

1. 馬大世界的待客之道

猶太人非常重視待客之道,最著名的好客事蹟,就是出自他們的先祖亞伯拉罕身上.他們相信,亞伯拉罕之所以百歲得子,與他的好客有直接關係.他把耶和華派來的三個遠道而來客人請進了帳篷,給他們水和食物.經文中說,亞伯拉罕不僅跑去迎接,還「俯伏在地」,三人臨走時他去遠送行.然後,三人中的一位說:「到明年這時候,我必要回到你這裡,你的妻子撒拉必生一個兒子.」

有一位朋友曾經在猶太人家中作客,大開眼界,記下了他的見聞:「在特拉維夫,我住在公司給他租的一間小房子裡,有個晚上他約見一個有求於他的年輕人,把我也帶上.他們談完話後,那年輕人擋住他的手,要把我們三人的賬一起付掉.

「他是我的客人.」丹恩指著我對他講.

年輕人還是堅持要一起付帳.他倆又嘰嘰咕咕討論起來.最後丹恩說:「OK,你買吧,但他是我的客人.」

現在想起來,還是感到這事很有趣味.猶太人好像對host(主人)的身份有特殊的理解,host不止是一個暫時的身份,還關係到一種他們很在乎的名分.猶太傳統的待客之道完全浸潤在他的習慣和言行裡.他給我上了一課,他如此在乎我們之間的身份關係:「他是我的客人」,不可以隨便被人奪走(哪怕人家並無這樣的想法),也不能讓他產生被我隨意處置之感,哪怕只是一杯飲料.」(多聯新聞網:「猶太人的至誠待客之道」)

如果說馬大的世界是一個為了不必要的事情而忙亂的世界,馬大可能會感到委屈,馬大其實是在努力實踐待客之道.在他們的世界中,盡最大的努力去款待客人,是當盡之義,是傳統,是習俗,也是美德.

馬大的心情也是可理解的.經文雖然只突出耶穌一人,然而,這突然到訪的客人,合理假設應該不止耶穌一人,至少還有他的門徒吧!突然之間要接待十幾位大漢,單單是飲食,都足夠馬大「心裡煩亂」了,對嗎?從這樣的處境去看,如果我們硬要說她發出抱怨,她的抱怨也是事實:「我的妹妹留下我一個人伺候」.因此,在馬大的世界裡,她提出的要求:「請吩咐她來幫助我」也是合理吧!

在馬大的世界中,她恰如其分地做好主人接待客人的角色,反而是馬利亞好像走錯了位置.從畫面所見,她並非和姊姊馬大同一位置去為客人預備這樣,預備那樣,卻「在主的腳前坐著聽他的道」,把自己當作耶穌的門徒一樣.如果鄰居三姑和六婆看見,可能會議論說:「女仔人家,尚未出嫁,走入男人堆幹甚麼?」

2. 馬大世界的待客之道不合時宜

這個故事最特別的地方,是耶穌這位客人對主人家的評論:

10:41「馬大,馬大,你為許多的事操心煩惱,

10:42但是不可少的只有一件.馬利亞已經選擇了那上好的福分,是沒有人能從她奪去的.」

為了做好主人的角色,為許多事操心煩惱是可理解的,但耶穌卻說似乎在說這種操心是不必要的；反過來說,耶穌卻嘉許馬利亞這位按當時的標準沒有扮演好主人角色的表現,並且明言她「已經選擇了那上好的福分」.這是故事的張力所在,也是令人驚訝的地方.

「那上好的福分」是甚麼呢?上文記載了耶穌說過這樣的話:

10:21「父啊,天地的主,我感謝你!因為你把這些事向聰明智慧的人隱藏起來,而向嬰孩啟示出來.父啊,是的,因為你的美意本是如此.

10:22一切都是我父交給我的.除了父,沒有人知道子是誰；除了子和子所願意啟示的人,沒有人知道父是誰.」

10:23耶穌轉身私下對門徒說:「看見你們所看見的,那眼睛有福了.

10:24我告訴你們,從前有許多先知和君王要看你們所看的,卻沒有看見,要聽你們所聽的,卻沒有聽見.」

「眾里尋他千百度,驀然回首,那人卻在燈火闌珊處.」馬大按著她的世界的運作模式,符合世界的要求,卻無暇回首,因此未能發現「在燈火闌珊處」的「那人」,甚至也不覺得馬利亞「選擇了那上好的福分」.

福終於來到,福又擦身而過,叫人倍感可惜.內地網民曾經熱烈地討論一個廣告,這個廣告大約是這樣的:

一位媽媽對懷中的嬰孩說:等你長大了,媽媽就享福了；後來,對長成大男孩的兒子說:等你考上大學,媽媽就享福了；後來,對大學畢業的兒子說:等你結完婚有了孩子,媽媽就享福了；後來,兒子生了個女兒,她對懷中的孫女說:等你長大了,奶奶就享福了!

或者,福已經臨到,然而,我們的世界令我們察覺不到,於是,福白白地溜走了,我們繼續按我們的世界的指導去追尋福.

我們上一代的父母,為家庭的基本生存條件而掙扎,哪有閒情去享兒女之福或家庭之樂呢?現在生活環境可能改善了,但父母也不一定有閒情去享受兒女之福或家庭之樂.如馬大一樣,有些父母「為許多的事操心煩惱」,四出為兒女找尋最好的玩具,驀然回首,兒女最好的玩具可能就是他們的父母；另一些父母「為許多的事操心煩惱」,四出為兒女找最好的圖書,驀然回首,最好的圖書可能就是父母與孩子一齊讀的圖書；又有一些父母四出為兒女找最好的課程,驀然回首,孩子最不能被取代的老師可能就是父母自己；有些父母四出去為兒女找最好的旅遊路線,其實,最快樂的旅程,就是和你最喜愛的人一齊共渡的旅程.兒女們長大成人,成家立室,年長的父母最期待的,豈不是一家團聚,享兒孫之福嗎?奇怪的情況是,當兒孩齊齊整整聚首一堂時,兩老卻在廚房裡忙個不停,整頓晚飯,根本沒有坐下共膳,到大家都差不多吃完了,才放下圍裙出來.作為兒女的,要兩老服侍自己,最後吃菜頭菜尾,情何以堪呢?

福,可能觸手可及,然而,如果當事人仍然堅持舊有的世界--縱使那世界在過去某時是真實的,福始終又擦身而過.

正如那位前述的猶太人一樣,他自覺是主人,對方是客人.在他的世界中,主客分明,他要扮演好主人的角色,禮待客人,照顧客人,保護客人.在這樣主客分明的世界中,如此盡責的主人難免為許多事操心煩惱.在他的世界之外,會否有其他可能性呢?例如主客互動?又例如反客為主?

3. 耶穌是主而不是客

10:38他們繼續前行,耶穌進了一個村莊.有一個女人,名叫馬大,接他到自己家裡.

10:39她有一個妹妹,名叫馬利亞,在主的腳前坐著聽他的道.

10:40馬大伺候的事多,心裡忙亂,進前來,說:「主啊,我的妹妹留下我一個人伺候,你不在意嗎?請吩咐她來幫助我.」

道成內身的上主到訪,誰是主,誰是客呢?馬大接他到自己家裡,馬大當然是主,耶穌當然是客.所以,馬大要求耶穌吩咐馬利亞去幫助她.有趣的地方是,馬大稱呼耶穌為主,但她卻堅持自己是主,因為她要求耶穌按她的議程去行事.

10:41主回答說:「馬大,馬大,你為許多的事操心煩惱,

10:42但是不可少的只有一件.馬利亞已經選擇了那上好的福分,是沒有人能從她奪去的.」

耶穌雖然來作客,但他回應馬大的要求時,卻彷如以主人的身份回答.這位客人,其實是名正言順的主.試想想,耶穌臨到你,你把他當作客人來招呼,還是以他為主呢?耶穌臨在我們的禮拜堂,我們把他視為客人來把最好的表演給祂觀賞,還是把他當作主人呢?如果我們把他當作客人去招呼,就難免諸事繁忙,為許多事操心煩惱,然而,如果我們還原祂為主人的角色,就可以如馬利亞般,「在主的腳前坐著聽他的道」.《溏心風暴之家好月圓》中荷媽(李司琪飾)有一句金句:「大家唔使驚,荷媽係大廳!」道成內身的上主到訪,我們會否感到:「大家唔使驚,耶穌係大廳」呢?

上主臨到人間,人們不需要再徬徨等待,最好的已到,按著祂的道去重構世界的人,就會發現祂；祂來不是接受我們的吩咐去實現我們的議程,而是來呼喚我們聽祂的道,以祂的道來調整我們的議程,以祂道改變我們對事物的態度,重塑我們的世界.如果馬大以耶穌為主去重塑她的世界,她本來那煩躁不安世界就可以如馬利亞的世界一樣安穩了.

四. 總結

馬大和馬利亞的故事,給予我們的挑戰未必是「應該去工作還是去崇拜」的選擇,而是有沒有發現耶穌這位客人其實是主人；同時,也未必是考驗我們有否以馬利亞的心去生活在馬大的世界中的能力,而是提示我們,耶穌是主而不是客,我們「在主的腳前坐著聽他的道」是適合的回應,也應以祂的道來重構我們的世界.以馬利亞的心去活在馬大的世界,有時是不可能的,因為馬大的世界並不容許馬利亞去聆聽並實踐上主的道.例如善良的人不忍看見飢荒地區的人白白餓死,也不忍看見飢荒地區的兒童營不良,然而,有些世界的遊戲規則無法令這些善良的人的善心得到實現.例如,有新聞報道指出:「美國農業部的資料顯示,2015年首5個月美國東北地區已倒掉3100萬磅鮮乳,比去年同期增加31\%,為有史以來最多.」(自由時報2015-07-13)又有報道說:「2015年元旦起廣東各地紛紛出現酪農倒奶的現象,酪農每天倒掉的鮮奶超過20噸.」(蘋果日報2015年01月09日)除了牛奶及牛奶製品,你還發現「送去垃圾堆也不可以送給有需要的人」的其他類似情況嗎?這些現象說明,我們的世界需要改變,否則是會阻礙人們追求更美好的事物,阻礙人們聆聽和實踐上主的道.

現代世界中所有人都「為許多事操心煩惱」,連教會也諸事繁多,彷彿「忙」是「有用」和「成功」的代名詞一樣,因此沒有人敢說自己「得閒」,免得被人認為「懶惰」和「無用」.然而,我們都不能不反思一下忙忙碌碌所為何事,尤其是那在地上為主作見證的教會.「著名天主教靈修學作者 Thomas Green指出「為主作工」與「作主工作」兩者有分別；前者是按個人認為主所喜愛的,其實是自己所愛的,為主而作的事工；後者則是順服主的心意,按照祂喜愛而作的事奉.Green以送禮來說明,「為主作工」是送禮者單方面認為對方會喜歡而送的禮物；「作主工作」則問了對方,然後按對方心意致送的禮物.」(胡志偉:教牧與退修)



\newpage

\section{求主教導我們禱告}
\label{sec:329}
\textbf{蘇成溢牧師 - 求主教導我們禱告}
\newline
\newline
\hyperref[sec:328]{< < < PREV SERMON < < <}
~
\hyperref[sec:toc]{[返目錄]}
~
\hyperref[sec:330]{> > > NEXT SERMON > > >}


何西阿書 1:2-10
\newline
\begin{longtable}{cl}
\hline
\hline
章節 & 經文 (和合本修訂版)\\
\hline
1:2 & \begin{tabularx}{0.7\textwidth}{X} 耶和華初次向何西阿說話。耶和華對他說：「你去娶一個淫蕩的女子為妻，收那從淫亂所生的兒女；因為這地行大淫亂，離棄耶和華。」 \end{tabularx} \\ \\ \relax
1:3 & \begin{tabularx}{0.7\textwidth}{X} 於是，何西阿去娶了滴拉音的女兒歌篾。她就懷孕，為何西阿生了一個兒子。 \end{tabularx} \\ \\ \relax
1:4 & \begin{tabularx}{0.7\textwidth}{X} 耶和華對何西阿說：「給他起名叫耶斯列；因為再過片時，我要懲罰耶戶家在耶斯列流人血的罪，也必終結以色列家的王朝。 \end{tabularx} \\ \\ \relax
1:5 & \begin{tabularx}{0.7\textwidth}{X} 到那日，我必在耶斯列平原折斷以色列的弓。」 \end{tabularx} \\ \\ \relax
1:6 & \begin{tabularx}{0.7\textwidth}{X} 歌篾又懷孕，生了一個女兒，耶和華對何西阿說：「給她起名叫羅‧路哈瑪；因為我必不再憐憫以色列家，絕不赦免他們。 \end{tabularx} \\ \\ \relax
1:7 & \begin{tabularx}{0.7\textwidth}{X} 我卻要憐憫猶大家，使他們靠耶和華－他們的神得救；我必不讓他們靠弓、刀、戰爭、馬匹與騎兵得救。」 \end{tabularx} \\ \\ \relax
1:8 & \begin{tabularx}{0.7\textwidth}{X} 歌篾在羅‧路哈瑪斷奶以後，又懷孕生了一個兒子。 \end{tabularx} \\ \\ \relax
1:9 & \begin{tabularx}{0.7\textwidth}{X} 耶和華說：「給他起名叫羅‧阿米；因為你們不是我的子民，我也不是你們的神。」 \end{tabularx} \\ \\ \relax
1:10 & \begin{tabularx}{0.7\textwidth}{X} 然而，以色列的人數必多如海沙，不可量，不可數。從前在甚麼地方對他們說「你們不是我的子民」，將來就在那裡稱他們為「永生神的兒子」。 \end{tabularx} \\ \\
[1ex]
\hline
\hline
\end{longtable}
歌羅西書 2:6-15
\newline
\begin{longtable}{cl}
\hline
\hline
章節 & 經文 (和合本修訂版)\\
\hline
2:6 & \begin{tabularx}{0.7\textwidth}{X} 既然你們接受了主基督耶穌，就要靠著他而生活， \end{tabularx} \\ \\ \relax
2:7 & \begin{tabularx}{0.7\textwidth}{X} 照著你們所領受的教導，在他裡面生根建造，信心堅固，充滿著感謝的心。 \end{tabularx} \\ \\ \relax
2:8 & \begin{tabularx}{0.7\textwidth}{X} 你們要謹慎，免得有人用他的哲學和虛空的廢話，不照著基督，而是照人間的傳統和世上粗淺的學說，把你們擄去。 \end{tabularx} \\ \\ \relax
2:9 & \begin{tabularx}{0.7\textwidth}{X} 因為神本性一切的豐盛都有形有體地居住在基督裡面； \end{tabularx} \\ \\ \relax
2:10 & \begin{tabularx}{0.7\textwidth}{X} 你們在他裡面也已經成為豐盛。他是所有執政掌權者的元首。 \end{tabularx} \\ \\ \relax
2:11 & \begin{tabularx}{0.7\textwidth}{X} 你們也在他裡面受了不是人手所行的割禮，而是使你們脫去肉體情慾的基督的割禮。 \end{tabularx} \\ \\ \relax
2:12 & \begin{tabularx}{0.7\textwidth}{X} 你們既受洗與他一同埋葬，也就在此禮上，因信那使他從死人中復活的神的作為跟他一同復活。 \end{tabularx} \\ \\ \relax
2:13 & \begin{tabularx}{0.7\textwidth}{X} 你們從前在過犯和未受割禮的肉體中死了，神卻赦免了你們一切的過犯，使你們與基督一同活過來， \end{tabularx} \\ \\ \relax
2:14 & \begin{tabularx}{0.7\textwidth}{X} 塗去了在律例上所寫、敵對我們、束縛我們的字據，把它撤去，釘在十字架上。 \end{tabularx} \\ \\ \relax
2:15 & \begin{tabularx}{0.7\textwidth}{X} 基督既將一切執政者、掌權者的權勢解除了，就在凱旋的行列中，將他們公開示眾，仗著十字架誇勝。 \end{tabularx} \\ \\
[1ex]
\hline
\hline
\end{longtable}
路加福音 11:1-13
\newline
\begin{longtable}{cl}
\hline
\hline
章節 & 經文 (和合本修訂版)\\
\hline
11:1 & \begin{tabularx}{0.7\textwidth}{X} 耶穌在一個地方禱告。禱告完了，有個門徒對他說：「主啊，求你教導我們禱告，像約翰教導他的門徒一樣。」 \end{tabularx} \\ \\ \relax
11:2 & \begin{tabularx}{0.7\textwidth}{X} 耶穌對他們說：「你們禱告的時候，要說：『父啊，願人都尊你的名為聖；願你的國降臨； \end{tabularx} \\ \\ \relax
11:3 & \begin{tabularx}{0.7\textwidth}{X} 我們日用的飲食，天天賜給我們。 \end{tabularx} \\ \\ \relax
11:4 & \begin{tabularx}{0.7\textwidth}{X} 赦免我們的罪，因為我們也赦免凡虧欠我們的人。不叫我們陷入試探。』」 \end{tabularx} \\ \\ \relax
11:5 & \begin{tabularx}{0.7\textwidth}{X} 耶穌又對他們說：「你們中間誰有一個朋友半夜到他那裡去，對他說：『朋友！請借給我三個餅； \end{tabularx} \\ \\ \relax
11:6 & \begin{tabularx}{0.7\textwidth}{X} 因為我有一個朋友旅途中來到我這裡，我沒有東西招待他。』 \end{tabularx} \\ \\ \relax
11:7 & \begin{tabularx}{0.7\textwidth}{X} 那人在裡面回答：『不要打擾我，門已經關了，孩子們也同我在床上了，我不能起來給你。』 \end{tabularx} \\ \\ \relax
11:8 & \begin{tabularx}{0.7\textwidth}{X} 我告訴你們，雖不因他是朋友起來給他，也會因他不顧面子地直求，起來照他所需要的給他。 \end{tabularx} \\ \\ \relax
11:9 & \begin{tabularx}{0.7\textwidth}{X} 我又告訴你們，祈求，就給你們；尋找，就找到；叩門，就給你們開門。 \end{tabularx} \\ \\ \relax
11:10 & \begin{tabularx}{0.7\textwidth}{X} 因為凡祈求的，就得著；尋找的，就找到；叩門的，就給他開門。 \end{tabularx} \\ \\ \relax
11:11 & \begin{tabularx}{0.7\textwidth}{X} 你們中間作父親的，誰有兒子求魚，反拿蛇當魚給他呢？ \end{tabularx} \\ \\ \relax
11:12 & \begin{tabularx}{0.7\textwidth}{X} 求雞蛋，反給他蠍子呢？ \end{tabularx} \\ \\ \relax
11:13 & \begin{tabularx}{0.7\textwidth}{X} 你們雖然不好，尚且知道拿好東西給兒女，何況天父，他豈不更要把聖靈賜給求他的人嗎？」 \end{tabularx} \\ \\
[1ex]
\hline
\hline
\end{longtable}
講題:求主教導我們禱告　Lord, teach us to pray

經文:何西阿書1:2-10;歌羅西書2:6-15;路加福音11:1-13

講員:蘇成溢牧師

場合:香港中文大學崇基學院禮拜堂

日期:2016年7月24日

1 引言

1.1 路加11:1-13與馬太6:9-15 (主禱文)；7:7-12 (祈求,尋找,叩門)

在馬太6:6-15耶穌教導門徒如何祈禱,並提供一個禱文範本,是因為主要提醒他們有一個正確的祈禱態度 (5-8節)；至於在7:7-12 耶穌論及祈求,尋找,叩門,卻是與待人的態度有關 (12節)

1.2 路11:1-13與太6:9-15；7:7-12約的比較

1.2.1主禱文內容  路加版本較短,包括 i) “父啊/亞爸” (馬太:我們在天上的父)  ii) 兩個願:主名為聖,主國降臨  (馬太:三個願:主名為聖,主國降臨,主旨得成)  iii) 天天賜日用糧食 (馬太:今天賜日用糧食)  iv) 兩個求:罪得赦免,不入試探 (馬太:三個求:罪得赦免,不入試探,脫離惡者)

1.2.2論及祈求,尋找,叩門的生活例子,用求魚求蛋,豈會給蛇給蠍子 (馬太:求餅求魚,豈會給石頭給蛇) 各例子都是人的食物,但有些遠看形狀相似卻對身體有害之物

2  路11:1-13有關禱告的教導

2.1 禱告的態度

2.1.1 真摯而坦率.主禱文的用語“父啊/亞爸” (2節),提醒了向天父禱告,要像兒女向父母傾訴一樣,其親切程度如同父母對子女的暱稱,子女對父母也如是.

2.1.2 主權在天父.9-13節共出現五次 “求” 字,表明祂是主,我們是祈求者；絕非我們是吩咐者,祂是受命完成任務者.

2.2 禱告的學習

2.2.1 尊上主旨意.禱告主名為聖,主國降臨,表明在祈禱時,要先承認上主的神聖與掌權；然後求衪供應,幫助和引領我們生命中肉體和心靈所需. (正如 西2:8提及的來自世界的各類試探和挑戰；詩19:12-14也告誡基督徒留意自己的隱惡,任意妄為,言行和意念.)

2.2.2 確信上主美意.9-10節的重複句及三個被動式動詞 “就給…. 就找到…. 就給開門” ,強調天父是隨時隨地回應禱告的；此外,魚和蛋是日用糧食/餅 (參5節) 以外的副食品,用此例子表明天父對禱告的回應,必更美好.

3 總結與反思

3.1 確信天父

每當禱告時,要確信天父是聽禱告並且會回應的. (有猶太拉比論祈禱說:i) 上主接近祂所造的,如同耳與口般接近  (參詩138:6 “耶和華雖崇高,卻看顧卑微的人.”)  ii) 人不能同時聽兩個人說話,但全世界的人若同一時間向上主呼求,上主仍聽到他們的呼聲  iii) 人可能因朋友的請求感到煩惱,但對上主來說,只要有人將他的需要和請求放在衪面前,祂只會更愛他們.

3.2 禱告實踐

3.3.1 真摯而坦白.甘易蓬神父的《祈禱與現代人》提到:祈禱時人應毫無保留地將真實的自己放在耶穌面前,是最基本的態度.也即是說,在祈禱時能將生命完全剖白的人,才是最能與上主交心的禱告者.

3.3.2 專一倚靠.向上主祈禱,切勿有測試,循例,無可奈何的想法,乃要像小孩子一樣,在生命各個需要的時刻,第一時間想到的,就是最親愛的人.



\newpage

\section{從差異與分歧到合一與共融}
\label{sec:330}
\textbf{范晉豪座堂牧師 - 從差異與分歧到合一與共融}
\newline
\newline
\hyperref[sec:329]{< < < PREV SERMON < < <}
~
\hyperref[sec:toc]{[返目錄]}
~
\hyperref[sec:331]{> > > NEXT SERMON > > >}


講題:從差異與分歧到合一與共融

講員:范晉豪座堂牧師

場合:香港中文大學崇基學院禮拜堂

日期:2016年7月31日

在云云美國漫畫英雄電影中,我最喜愛的始終是《X-men》,它不會像超人與蝙蝠俠,美國隊長與Ironman般堆砌劇情,「為打而打」,X-men內兩派變種人之間的張力,源自彼此分歧的信念.在面對社會主流歧視的當下,到底要勇武抗爭,還是以非暴力精神,相信人與之間在差異歧見中,仍能共融共存?

《X-men》漫畫於1963年面世,X教授與磁力王兩派領袖的創作靈感,取材自60年中兩位政治立場南轅北轍的美國黑人民權運動領袖.他們就是馬丁路德金和麥爾坎·X(Malcolm X).馬丁路德金是Ｘ教授的原型,他主張非暴力抗爭,不願挑起黑人與白人之間的仇恨,要針對是不平等的制度,要教育的是下一代的黑人與白人,放低彼此的歧見共融共存；相反,麥爾坎·X是磁力王的原型.他主張勇武抗爭,彼此的矛盾不能透過和平理性的對話去達成共識,這是戰爭,只有透過暴力,對抗白人,為黑人申張正義,取回他們應得的一切.

故事中擁有異能的變種人因為跟普羅大眾有差異,遭受社會排斥與歧視,在現實世界也略見不鮮.歸根究底,這源於有限的人類對自己沒有無限的自信,面對和自己差異的族群感覺受威脅,不安,恐懼與焦慮油然而生.在歷史裡,我們看見社會建制的主流容易透過權力控制大局,剷取異己.政治如是,宗教亦如是,但如此狀況實在與基督信仰有著根本的矛盾.

在新約教會時代,外邦使徒保羅因為向外邦人宣教,其使徒身分被當時教會主流的猶太基督徒所質疑,他所傳的福音更被指為討外邦人喜歡,得人心的技倆(加1:10) .保羅對這指控十分憤怒(加1:8-9) ,他強調沒有不同版本的福音,可供體貼猶太人和外邦人兩種福音,任君選擇.「沒有另一個福音.」(加1:7) 我們要傳的絕非單單猶太民族信仰系統,而是要傳那位捨已身救贖歷世歷代的人脫離罪惡世代的基督(1:4) .基督來,不單是為猶太人,亦不只為基督徒,而是為歷世歷代所有人宣告上帝的福音.前南非杜圖大主教其中一本著作《上帝不是一個基督徒》(God is not a Christian) ,正要表達無限的上帝不像有限的人類般思想狹隘,基督徒不應以「護教」之名,限制上帝那無邊界的愛,就如昔日守割禮的基督徒,把信仰控制在猶太人的安全區域以內.

事實證明,基督所傳的上帝是向整個世界不同信念,政見,種族,宗教的人張開雙手,歡迎他們回到衪家裡共融的上帝.不單新約,我們從舊約也可看見那位向人張開聖臂的上帝.上帝揀選亞伯蘭,以撒和雅各這一族,從其後裔建立以色列上帝的選民,絕不是排拒外邦人,而是讓這民族成為萬民得福的器皿(創12:12:3, 22:18, 26:5, 28:14) .在聯合王國最興盛的所羅王時代,所羅門向天舉手,奉獻聖殿禱告上主說:「求你在天上你的居所垂聽,照著外邦人所祈求的而行,使天下萬民都認識你的名,敬畏你,像你的民以色列一樣」(列上8:43) 詩篇96篇也和應對上主的讚美是不分國界與種族,列國要述說他的榮耀,萬民要述說他的奇事.

面對如此一位上帝,吸引歷世歷代萬族萬民,進入其無邊際的愛裡,難道我們還用恐懼因為人與人之間彼此的差異,視彼此為敵,建立圍牆,透過種種權力遊戲,控制與排除異己嗎?耶穌對歷代的門徒回答說:「不要怕!只要信!"(可5:16)

福音書記載令人驚訝的大多是耶穌基督的言論與作為,但路7:1-10記載了一個人卻反過來令耶穌驚訝與讚嘆.他不是猶太人,也不是耶穌的門徒,他是羅馬駐守的百夫長.對國族意識強烈的猶太人來說,統治他們的羅馬人軍官絕對是仇人.但這令耶穌驚訝的人竟然能打破種族間的國仇家恨,託猶太長老請耶穌醫治他頻死的僕人.當耶穌往那家裡去,他自覺不配見主,託朋友請耶穌說一句話,他深信僕人必能得救,就好像軍隊聽從上級命令般.

一隊低級的雜牌兵只著眼自己的利益與安危,他們眼中的戰友與敵人沒有太大的分別,只是替代自己送死的工具,如此互不信任的軍隊,只是一盤散沙,必敗無疑.相反,一隊精銳部隊如同一體,上司是值得信任的大腦,服從的下屬信任並將自己生命交與上司的指揮,為是要帶來整體的勝利.有限的人類真的不可信,且充滿不安與焦慮,若我們甘為雜牌軍,只定晴專注自己與他人之間的利益差異,這只會帶來鬥爭與仇恨,彼此盤算如何控制更多資源,如何打倒對手,排除異己.

人縱然不可信,但無限的上帝是可信的,人應專注在自己與上帝的關係中建立信心.在這全然不同於我們的上帝(Wholly Other)裡,我們能從與衪共融中,體會到彼此差異可以不帶侵略性,共融不是要侵蝕了人的個性,反而是豐富完全了我們的生命,打破了人與人,我與你之間的敵我鴻溝,讓共融,信任,盼望,人類整體福祉成為可能.在上帝裡面我們不只看見自己,愈來愈看見和我們不同的異已,其實也是同一屋簷下,上主家裡的人.

這位沒有露面的百夫長對耶穌充滿無比的信心,因為這份簡單卻偉大的信心,他願意相信神蹟,即上帝的作為能臨現這充滿仇恨,脆弱,不安的人間.耶穌稱讚他這麼大的信心實在前所未見.他希望歷代的跟隨者能對上主建立這份信心,用生命延續傳揚自舊約至新約,上帝要使萬民得福的福者,叫我們在這時代為這紛亂分歧的世間締造合一共融的神蹟.



\newpage

\section{異鄉人在故鄉}
\label{sec:331}
\textbf{賴品超教授 - 異鄉人在故鄉}
\newline
\newline
\hyperref[sec:330]{< < < PREV SERMON < < <}
~
\hyperref[sec:toc]{[返目錄]}
~
\hyperref[sec:332]{> > > NEXT SERMON > > >}


彼得前書 2:9-17
\newline
\begin{longtable}{cl}
\hline
\hline
章節 & 經文 (和合本修訂版)\\
\hline
2:9 & \begin{tabularx}{0.7\textwidth}{X} 不過，你們是被揀選的一族，是君尊的祭司，是神聖的國度，是屬神的子民，要使你們宣揚那召你們出黑暗入奇妙光明者的美德。 \end{tabularx} \\ \\ \relax
2:10 & \begin{tabularx}{0.7\textwidth}{X} 「你們從前不是子民，現在卻成了神的子民；從前未曾蒙憐憫，現在卻蒙了憐憫。」 \end{tabularx} \\ \\ \relax
2:11 & \begin{tabularx}{0.7\textwidth}{X} 親愛的，你們是客旅，是寄居的，我勸你們要禁戒肉體的情慾；這情慾是與靈魂爭戰的。 \end{tabularx} \\ \\ \relax
2:12 & \begin{tabularx}{0.7\textwidth}{X} 你們在外邦人中要品行端正，好讓那些人，雖然毀謗你們是作惡的，會因看見你們的好行為而在鑒察的日子歸榮耀給神。 \end{tabularx} \\ \\ \relax
2:13 & \begin{tabularx}{0.7\textwidth}{X} 你們為主的緣故要順服人的一切制度，或是在上的君王， \end{tabularx} \\ \\ \relax
2:14 & \begin{tabularx}{0.7\textwidth}{X} 或是君王所派懲惡賞善的官員。 \end{tabularx} \\ \\ \relax
2:15 & \begin{tabularx}{0.7\textwidth}{X} 因為神的旨意原是要你們以行善來堵住糊塗無知人的口。 \end{tabularx} \\ \\ \relax
2:16 & \begin{tabularx}{0.7\textwidth}{X} 雖然你們是自由的，卻不可藉著自由遮蓋惡毒，總要作神的僕人。 \end{tabularx} \\ \\ \relax
2:17 & \begin{tabularx}{0.7\textwidth}{X} 務要尊重眾人；要敬愛教中的弟兄姊妹；要敬畏神；要尊敬君王。 \end{tabularx} \\ \\
[1ex]
\hline
\hline
\end{longtable}
講題:異鄉人在故鄉　Aliens in Homeland

經文:彼得前書2:9-17節

講員:賴品超教授

場合:香港中文大學崇基學院禮拜堂

日期:2016年8月7日

甚麼是福音?＝信耶穌得(死後)永生,人生的最後／真正歸宿；人生＝寄居的客旅,訃文“在世寄居n載".

上一代非常流行的詩歌:這世界非我家,我無一定住處,天堂若非我家,我必流離無依, …. 故我不再貪愛這世界為我家.對世界無眷戀,但也沒有甚麼責任,除了傳福音,無甚有價值的事可做,等上天堂.看似消極厭世.

有人批評基督教逃避社會責任,甚至不環保,因污染也無妨.不一定!遊客可以不留下也不帶走甚麼(環保遊),甚至有可能把一個更好的地方留給下一手.(曾租屋住,搬走時乾淨過入住之前,業主大讚.皆因家母愛潔.)

今天的經文,先講福音:上帝在基督裡的工作,拯救,揀選,更新,呼召(2:9-10),信徒因而有獨特身份,就是作為“被揀選的族類,是有君尊的祭司,是聖潔的國度,是屬上帝的子民”(彼前2:9)在此天國公民身份的基礎上,經文再講另一個身份,就是寄居的,作客旅的.

為何說是寄居的,作客旅?按照經文所說,此書信是給在多個不同地方的離散寄居者(παρεπιδήμοις διασπορα),包括“本都,加拉太,加帕多家,亞西亞,庇推尼”等地的信徒(彼前1:1).讀者很可能是在社會／政治上的意義上是寄居者,不是住在家鄉／本國的猶太人.在耶穌與保羅的時代,不少猶太人散居各地,似中國人多華僑.收信者可能主要是離散(diaspora)的猶太僑民,但也不排除也有外邦人(學者對此有爭論).無論如何,由於天國公民的身份才是primary /首要／基本／必然,身住那一國家則是secondary／次要／偶然(可變),就有如寄居.在世寄居的天國子民,應如何生活?

彼得對作客旅／寄居者的勸告,信仰上敬畏上帝,在教會生活上親愛弟兄,在個人生活上:(1)積極而言,品行端正(2:12),行善(2:15)；(2)消極而言,要禁戒肉體的私欲(2:11),不要“藉自由遮蓋惡毒” (2:16)；(3)在社會／政治生活上:要順服君王,長官,以至既有制度.

(1)品行端正,行善－任何情況都要做,但作為僑居異地者更為重要,因可給當地人良好印象,改善與鄰舍關係.

(2)孤身到外,沒有家庭朋輩壓力,覺自由,想做就做.對當地無感情,只有sense of insecurity,總是想多賺個錢旁身／回鄉,在家鄉不敢做的,在外地可放縱地做,尤其如果當地不犯法.小心,若讓當地人覺這些外來者很自私,貪婪,甚至aggressive,要來掠奪甚麼,反惹仇恨.

(3)人到外地,對外地的制度有時覺得不習慣,不合理或對自己有損,覺得還是習慣故鄉的一套.(太煩了!我鄉下不是如此!).很容易不尊重當地文化風俗甚至法律.(通常當地人都不喜歡治外法權/Islamic law in Britain)小心,入境問禁,尊重當地文化,較易受當地人照受,減少被排斥.

彼得的勸告,似乎就是勸信徒做奉公守法的良好市民.不要讓人覺得我們是自私自利,不理公益的外來者,不應將個人私欲凌駕公共秩序利益,多做有益於公眾的事情(奉公)；也不要讓人覺得我們,破壞既有秩序(守法).這些勸告對於個人的好處(保身)固然重要,但對經文來說(明節)更為重要:我們如此行,目的是宣揚上帝的美德.(正如我們的學生去做交換生,總是選成績好,品行好的,為免破壞CUHK的名聲,因為這些學生在外地的行為,影響校譽／image.)【曾遇中國留學生戲謔,在外國做壞事,記得大聲講日文!】

信徒如此生活,非為自保,而是出於對上帝的信仰.我們信徒本身是“自由的”(彼前2:16),但也是“作上帝的僕人”(彼前2:16).如此生活,作這一切,不是無可選擇,被迫如此,也不是因為潮流興／入鄉隨俗,也不是因為君王的命令,而是為了“主的緣故” ,為了“歸榮耀給上帝”(彼前2:13),塞無知人的口(2:15),不要給人話柄.

對君王,臣宰以及人的一切制度的順服,不是基於甚麼愛國情操(we are not their nationals / citizens),不是對他們有何好感,更不是因為政策或官員們對我們有優惠,利益輸送,非因私慾而有所偏愛／恨惡,非因懼怕君王的權力(羅馬書的講法是:他們不是空空配劍的),更非因對權力的崇拜,而是為了彰顯上帝藉基督的救贖所顯出來的美德.

正因如此,順服君王／一切制度,不是一種毫無原則的順服(因敬畏上帝才順服君王)；對君王／政府,有一個恰當的態度,一種恰如其份的尊敬.該段經文的小結是“務要尊敬眾人,親愛教中的弟兄,敬畏上帝,尊敬君王” (彼前2:17),當中肯定尊敬君王只是信徒的眾多不同的責任的其中之一而已,所採取的態度也有分別.對上帝的是敬畏(φοβεισθε) (1:17),對教內的弟兄才是親愛(αγαπατε)(1:22),對(教外)的眾人是尊敬(τιμήσατε),而對君王也是同一字根的尊敬(τιματε).

換言之,這既不是用對上帝的敬畏去對敬畏君王【這是拜偶像!斷乎不可!】,也不是用對教內弟兄的親愛去親愛君王【有實際上的困難,不單可能有國仇家恨,更是親疏有別,對方不一定領情,只會讓人以為獻媚】,而是用對待一般教外眾人的尊敬,去尊敬君王而已.(不少人是either hate or love君王,但我們不巴結,也不卑視或敵視【這其實並不容易】,而是予以尊敬,正如尊敬其他人一樣,一視同仁,不因政治地位而有異.)【據說,大學高層曾討論可否不收大額捐助大學者診金；最後接受廉署的忠告:任何人,包括首富／特首 (校監) 來看病也要照收診金,一視同仁,當然會善待,正如善待一般病人一樣.】

信徒即使在被邊緣化的環境,甚至受迫逼／敵視 ,也要堅持信仰原則,一方面拒絕被同化,另一方面努力行善,即使不被欣賞,甚至要為行善而受苦,仍要堅持,better to suffer than to sin,正如基督立下榜樣(2:18-25) .

早期教會的實踐

早期教會的信徒,不僅猶太裔的基督徒如此,就連本來在種族上不是猶太人的信徒,也因為信仰的原因及所帶來的價值觀／生活態度,變得與原居地的人有所疏離；就是說,他們雖然是羅馬人並身處本國本土,但卻被視為異鄉人,被邊緣化,甚至誤解和歧視.此即所謂:異鄉人在故鄉.

早期教會的一份文獻(約為二至三世紀)正是如此描述基督徒:“他們雖住在他們自己的祖國,但好像是國中的寄寓者,他們作為公民,參預一切事務,但又像外人樣,忍受一切痛苦.他們以海外為家,每一國外都變成他們的外國.”

在教外人眼中信徒是怪人,是異類 (有如電影中的aliens 異形).早期的信徒不時受誣告,誹謗(無神,淫亂,吃人肉)；但他們仍按照《聖經》的教導(耶29:7 “因為那城得平安,你們也隨著得平安”),在具體的行動上,除了保持貞潔外,更致力為所住的城市求平安,做不少公益的事,照顧孤兒寡婦,貧苦無依者,甚至得到一些異教者的稱許.

當時社會:似很矛盾,崇尚自由(個人享樂)但又有政治壓制.皇帝／政府有權取締某些宗教,甚至有權要求人去拜官方所定的“神”,包括皇帝.日後,羅馬帝國大搞對君王的崇拜,要求信徒獻祭,不少信徒拒絕,大受壓逼.他們無法順服,因這是違反聖經的教導.原本就是因為敬畏神才尊敬君王,如硬要迫他們二揀一,最後唯有“順從神,不順從人,是應當的"(徒5:29).

香港的處境

我們的處境有點似早期教會面對的處境:似很矛盾,個人自由vs. 政制壓制.

我們生活在非常自我的世代,高舉自我的文化風尚,me generation,玩自拍(selfie 列入世紀重要關鍵詞), 高舉權利／自由不講義務／責任, 很容易“藉自由遮蓋惡毒” (2:16),尤其反映在消費文化(只要付得起,不用想環保,公平貿易),性文化(只要喜歡可以很自由,想做就做,只要不影響別人)(對社會真的無影響?)

香港人一直有很多自由,缺乏民主.最近好似自由也少了,激烈批評政府的人無端端被辭職或調職,失去贊助,又講不出為何,暫時無法證實有任何的政治迫害,但又隱隱讓人覺得好像是有某種政治的目的,收窄言論自由,誅滅異己的聲音.

與社會主流文化仍有不少重大的分歧,政治制度也不一定是我們所喜歡.面對這種處境,信徒應如何生活,如何彰顯神的美德?今天經文一些啟迪.

個人:小心私慾.在我們的社會,我們的工作環境,與我們的信念,價值觀可以很不一樣,朋輩social pressure不容易resist.(合資買六合彩,好少錢,又可聯絡感情,輸了當做善事,why not!? - 留意, 這有可能是“藉善事遮蓋貪婪”.)文化的滲透更難防範(傳媒／廣告:人生要講享受,甚麼都可以被人當作是基本人權的一部份),甚麼都可當作一種消費,甚至以消費的心態去教會.需要好好的思考, 如何用僅有的自由去作榮耀神的事, 而不是“遮蓋貪婪”.

教會:我們的社會人與人之間十分疏離,更需要有教內弟兄姊妹之間的親愛；但小心,這不是藉口以公益為名以謀私,將教會變成小圈子利益集團.香港教會人數也不少,社會有地位的也不少(佈道會專找高官名人作見證),參與很多教育慈惠工作,甚至有人覺得代表精英階層(有人甚至為入名校而入教).值得思考,教會會否以傳福音為理由,實則擴展版圖/enterprise, 反而犧牲公眾利益(少批評政府兩句,政治上多點識做,讓你開多間學校).

政治:順服君王,接受一切制度,似與民主／政制發展相違背,實則不然.我們沒有君王,仍對為官的仍給予應有的尊重,no more and no less,更期盼能有人與人之間(包括君王或長官)有在平等的基礎上互相尊重；你有一票我也有一票,不應有某些人因有某種社會地位而有多一票或兩票的特權.換言之,在民主制度中,有更大的可能去實現經文所表達的理想,就是信徒可自由地表達信仰,尊敬眾人也受眾人尊重 (有言論,結社以至宗教信仰的自由,也是基本人權) ,更不會因行善或爭取公義而受政治上的打壓,可以更自由地去彰顯上帝的美德──就是公義,光明,仁愛.

香港的信徒,對國家民族的歸屬感與認同也許相對模糊,反而對流散的處境和感覺卻並不陌生,不少人曾移民或寄居異地(手上或有外國護照)；但香港教會一直堅持的信念是“為這城求平安",而這甚至為持不同政治取向的基督徒所接受.這正是 “在世寄居的天國子民”,作為僑居的異類為身份認同所應有的態度.我們雖因信仰而有如異鄉人在故鄉,但仍應為這城求公義,民主,人權,這不僅有利於信徒信仰的實踐 ,也有利於社會大眾.

我們要堅持,也許要付出一定的代價(文化上,社會上被孤立,甚至政治上被打壓),但仍要追隨基督的十架.求主幫助我們,在我們寄居在世時,不單在個人生活,也有社會生活上,能宣揚衪美德使主得著榮耀,更希望可以在我們離世之時留下一個更美好,而不是更壞的世界給下一代.



\newpage

\section{看得見的抗爭看不見的和平}
\label{sec:332}
\textbf{鄧瑞強博士 - 看得見的抗爭．看不見的和平}
\newline
\newline
\hyperref[sec:331]{< < < PREV SERMON < < <}
~
\hyperref[sec:toc]{[返目錄]}
~
\hyperref[sec:333]{> > > NEXT SERMON > > >}


路加福音 12:49-59
\newline
\begin{longtable}{cl}
\hline
\hline
章節 & 經文 (和合本修訂版)\\
\hline
12:49 & \begin{tabularx}{0.7\textwidth}{X} 「我來是要把火丟在地上，假如已經燒起來，不也是我所希望的嗎？ \end{tabularx} \\ \\ \relax
12:50 & \begin{tabularx}{0.7\textwidth}{X} 我有當受的洗還沒有受，在這事完成之前，我是多麼地焦急！ \end{tabularx} \\ \\ \relax
12:51 & \begin{tabularx}{0.7\textwidth}{X} 你們以為我來是要使地上太平嗎？不！我告訴你們，是使人紛爭。 \end{tabularx} \\ \\ \relax
12:52 & \begin{tabularx}{0.7\textwidth}{X} 從今以後，一家五個人將要紛爭，三個和兩個相爭，兩個和三個相爭： \end{tabularx} \\ \\ \relax
12:53 & \begin{tabularx}{0.7\textwidth}{X} 父親和兒子相爭，兒子和父親相爭；母親和女兒相爭，女兒和母親相爭；婆婆和媳婦相爭，媳婦和婆婆相爭。」 \end{tabularx} \\ \\ \relax
12:54 & \begin{tabularx}{0.7\textwidth}{X} 耶穌又對眾人說：「你們看見西邊起了雲彩，就說：『要下大雨了』，果然就有； \end{tabularx} \\ \\ \relax
12:55 & \begin{tabularx}{0.7\textwidth}{X} 起了南風，你們就說：『要燥熱了』，也就有了。 \end{tabularx} \\ \\ \relax
12:56 & \begin{tabularx}{0.7\textwidth}{X} 假冒為善的人哪，你們知道分辨天地的氣象，怎麼不知道分辨這是甚麼時代呢？」 \end{tabularx} \\ \\ \relax
12:57 & \begin{tabularx}{0.7\textwidth}{X} 「你們又為何不自己判斷甚麼是合理的呢？ \end{tabularx} \\ \\ \relax
12:58 & \begin{tabularx}{0.7\textwidth}{X} 你同告你的冤家去見官，還在路上，要盡力跟他和解，免得他拉你到法官面前，法官把你交給法警，法警把你下在監裡。 \end{tabularx} \\ \\ \relax
12:59 & \begin{tabularx}{0.7\textwidth}{X} 我告訴你，就是最後一小文錢還沒有還清，你也絕不能從那裡出來。」 \end{tabularx} \\ \\
[1ex]
\hline
\hline
\end{longtable}
講題:「看得見的抗爭,看不見的和平」　Fight for Peace

經文:路加福音12:49-59

講員:鄧瑞強博士

場合:香港中文大學崇基禮拜堂

日期:2016年8月14日

各位弟兄姊妹,早安.

今天天氣好嗎?昨晚睡得好嗎?吃過早餐沒有?附近有什麼吃早餐的好地方?暑假去了哪裡玩?玩過的地方,哪裡最美?哪裡最抵玩?……

若我只懂問這些問題,只是講這些話題,只對這些東西有興趣,我想,你大大概概會認為我是個「貪玩識食」的人,你大大概概會喜歡和我「吹水」,你大大概概會覺得我平易近人.若你同樣喜歡這些話題,又喜歡和我「吹水」,則我們會相處融洽.就算交不成朋友,我以上講的東西,也應不會惹起你的反感,不會讓你受威脅,更不應引發你的敵意.

我們都希望耶穌多講些吃喝玩樂的東西,若他能多講一點投資必勝術,更好；若能教我們如何美容瘦身容光煥發,則perfect.若晚上回家,打開電視,見到耶穌介紹加利利的風光,多好.想像一下,耶穌教我們在春天應走那條路線,去欣賞加利利田野遍地的野花；在夏天,他教我們坐著遊船,在加利利湖垂釣；在秋天,跑到迦密山(Mount Carmel)上看日落；在冬天,則上黑門山(Mount Hermon)滑雪.何等寫意!看完旅遊節目,耶穌繼續主持飲食節目,教我們如何「今晚食好Ｄ」,教我們如何慢火烤一條加利利湖的「聖彼得魚」,教我們如何品嚐一支戈蘭高地的紅酒.何等快活!耶穌若是如此,則會通街都是Fans,我們會愛死他.他簡直就是高品味的優雅典範.

耶穌沒有主持旅遊節目,也無主持飲食節目,他主要的話題,是傳講「天國近了,你們應當悔改」.他因此失去很多Fans.天國,是一個新的世界,是一個與這個世界全然不同的新世界.在這個新世界裡,「飢渴慕義的人是有福的,因為他們必得飽足」.在這裡,人不是渴慕美食,而是渴慕神的正義.他說:「有人打你的右臉,連左臉也轉過來由他打」,又說:「要愛你們的仇敵,為那逼迫你們的人禱告」；又說:「要愛你們的仇敵,為那逼迫你們的禱告.這樣就可以作你們天父的兒女,因為祂叫日頭照好人,也照歹人；降雨給義人,也給不義的人」.這種生活形態與我們慣常的生活形態很不同.我們慣常的生活形態,受我們的欲望牽引.耶穌講的生活形態,受天父的愛主導.天國的生活形態,與我們慣常的生活形態,有一種張力.耶穌與我們,有一種對抗性.

若果耶穌只是問我們:「昨晚睡得好嗎?」我們會感到他的關懷.但耶穌若指著你,說:「昨夜你為何不睡,想著那些邪惡念頭」,或說:「由頭到腳,你只是一個自私的人」,或說:「你來這裡崇拜,內心半點敬虔也沒有」,則我們便可能不大喜歡耶穌了.若耶穌指著你,當眾指著你,公然揭露你的罪,你和他之間,便不可能和平共處了.不要說是耶穌指著你,若是現在我指著你,你也要殺了我.

耶穌將火丟在地上.

我們看看今日的福音書經課:路加福音12:49－59

路12:49　「我來要把火丟在地上,倘若已經著起來,不也是我所願意的嗎?

路12:50　我有當受的洗還沒有成就,我是何等地迫切呢?

路12:51　你們以為我來,是叫地上太平嗎?我告訴你們,不是,乃是叫人紛爭.

路12:52　從今以後,一家五個人將要紛爭:三個人和兩個人相爭,兩個人和三個人相爭；

路12:53　父親和兒子相爭,兒子和父親相爭；母親和女兒相爭,女兒和母親相爭；婆婆和媳婦相爭,媳婦和婆婆相爭.」

路12:54　耶穌又對眾人說:「你們看見西邊起了雲彩,就說:『要下一陣雨』；果然就有.

路12:55　起了南風,就說:『將要燥熱』；也就有了.

路12:56　假冒為善的人哪,你們知道分辨天地的氣色,怎麼不知道分辨這時候呢?」

路12:57　「你們又為何不自己審量甚麼是合理的呢?

路12:58　你同告你的對頭去見官,還在路上,務要盡力地和他了結；恐怕他拉你到官面前,官交付差役,差役把你下在監裡.

路12:59　我告訴你,若有半文錢沒有還清,你斷不能從那裡出來.」

耶穌不是講飲講食,耶穌直指人的罪.

他將火丟在地上.

「火」:在聖經裡,代表神的審判,代表神要煉淨一切渣滓,代表消除一切不義.

耶穌不是來到我們當中,風花雪月,嬉嬉哈哈,他是將火丟在地上,他的說話要評斷人性的是非曲直,他的行動要揭露世人的善惡正邪.

晨曦未露,曙光未顯之時,大地仍是暗淡矇矓,一切都隱藏在黑暗之中,正邪不辨.當朝陽乍現,晨光展露之時,鼠輩的行蹤便會敗露,真相會大白.耶穌是真理的光,他的來到,是照亮一切.獅子被顯為獅子,豺狼被顯為豺狼,一切都無法隱藏.當一切在黑暗裡時,混然不分,有一種黑暗中的寧靜.但當正邪分明,邪惡被指明,不義被揭露時,正邪便無可避免地有一決戰了.如講題所示,此之謂「看得見的抗爭」.

路12:49　「我來要把火丟在地上,倘若已經著起來,不也是我所願意的嗎?

路12:50　我有當受的洗還沒有成就,我是何等地迫切呢?

這兩節經文,聖經《呂振中》譯本譯作:「我來把火丟在地上；我是多麼迫切願它著起來阿!我有該受的洗,我是多麼迫切,直到完成阿!」

耶穌將火丟在地上,意味他要審判,分判正邪,他是多麼願意這審判快些完成,天國快些降臨.要完成這審判,耶穌必須經歷一個「洗禮」,那就是指「十字架」的苦難洗禮.十字架,完全呈現世人的罪,也完全呈現神的大愛.

神的大愛不會姑息罪惡,卻要揭露人性裡隱藏著的邪惡.邪惡的力量不會取消神的大愛,而終會被神的大愛所克服.十字架的苦難,是由於人類的罪惡造成.耶穌不顯露於世,世人的罪深藏心底,世間會有片刻黑暗中的寧靜.耶穌出現,如真光照耀,邪惡會被揭露.當時的政治力量,社會力量,宗教力量,為了自身的利益,激發了民眾的恐懼,假上帝之名,將耶穌推上十架上.

對耶穌而言,這個十架洗禮,正好讓每個人都眼白白地看見自身的不義和罪惡.光天化日,每個人都因著自身的軟弱,自私,仇恨,恐懼,參與一場殺戮.對抗這十架的殺戮,是神的大愛.有一種人類測不透的大愛,在十架上顯露,靜靜地告訴世人,只有這份愛,才能在邪惡的世界裡有正義的可能.

耶穌期待審判,也期待十架所揭露的救贖的完成.

「不義」被揭露後,便會頑強地和「正義」正面交鋒.

耶穌揭露「不義」的後果,便是引發正邪交戰.

路12:51　「你們以為我來,是叫地上太平嗎?我告訴你們,不是,乃是叫人紛爭.

路12:52　從今以後,一家五個人將要紛爭:三個人和兩個人相爭,兩個人和三個人相爭；

路12:53　父親和兒子相爭,兒子和父親相爭；母親和女兒相爭,女兒和母親相爭；婆婆和媳婦相爭,媳婦和婆婆相爭.」

光明顯露,就不能回到黑暗中的寧靜了.

相傳,東方有個小島,島上的居民,曾經過著和諧的生活,他們和諧,是因他們只有一個共同方向,就是努力向上.然後,來了一個人,這個人史稱「一男子」,他來了,將火丟在地上,然後,小島上每個人都在爭吵.原本沉悶的顏色,開始分裂出不同色彩,黃色,藍色,紅色,綠色,黑色,多姿多采.黃色和藍相爭,紅色和綠色相爭,紅色加黑色和綠色相爭,綠色和黃色又相爭.

能否回到起初的和諧?不可能了.

因為小島上的人,不再只向上,也學曉向左,向右,向前,向後.

爭吵,是因為有不義和正義的爭論,有人心底的邪惡和由神而來的良善的對抗.不義被揭露之後,不義會掙扎,正義會抵抗,不可能回到黑暗裡的平靜了.

這是人要懂得抉擇的時刻.耶穌說:

路12:54　「你們看見西邊起了雲彩,就說:『要下一陣雨』；果然就有.

路12:55　起了南風,就說:『將要燥熱』；也就有了.

路12:56　假冒為善的人哪,你們知道分辨天地的氣色,怎麼不知道分辨這時候呢?」

以色列地的西邊,是地中海.西面吹來的風,帶有濕氣,會下雨.以色列地的南邊,是曠野.南面吹來的風,帶著曠野的熱氣,帶來燥熱.那時的人,能看天做人.耶穌問:你們能預測天氣,怎麼不能分辨世情呢?

是的,我們面對的時勢是什麼呢?要如何抉擇呢?

我想,耶穌的想法是清楚的,看看罪在那裡被揭露,看看是那種力量將他人推向十架,看看是誰摧毀天國的美善,然後,面向神的審判,作自己的抉擇.

最後,耶穌以一法庭的情景結束這番話.

路12:57　「你們又為何不自己審量甚麼是合理的呢?

路12:58　你同告你的對頭去見官,還在路上,務要盡力地和他了結；恐怕他拉你到官面前,官交付差役,差役把你下在監裡.

路12:59　我告訴你,若有半文錢沒有還清,你斷不能從那裡出來.」

根據上下文,這段經文應是提示我們,神的審判將臨,時間緊迫,你不能再漫不經心地生活,要作出及時的抉擇了.耶穌將我們的生命,置放在神的審判前.就在當下,神要審判你,你如何面對?

我將這段經文的意義改變一下,將之應用在當下的衝突處境上.

面對社會上的各種衝突,我想,每個人都難免要有一立場.立場既定,便難免與另一立場有張力,有對抗.如耶穌所說,「父親和兒子相爭,兒子和父親相爭；母親和女兒相爭,女兒和母親相爭」.家庭的上一代和下一代有張力,教會裡不同顏色的陣營也有裂痕.這是無可避免的.在我看來,這仍未到上帝與魔鬼對決的地步.

如何處理這些衝突呢?

由於這不是上帝與魔鬼的對決,我想,或許,最好的方法,是讓自己及對方都能以罪人的身分在神面前認罪求赦.我是罪人,故此,就算我認為自己全對,我也無權任意處置對方.他是罪人,故此,就算他認為自己全對,我也體諒他會忽略他的盲點.就因為我們都是罪人,我們不能因彼此不同而排斥對方,而只能共同仰賴神的憐憫.這是將生命放在神之下而建造起來的和平,一種心靈裡的和平,一種「看不見的和平」.我們彼此可以仍屬不同陣營,但我不是審判官,我無權審判你.我們只能在神的審判下,謙和地以各自的立場,去造福社會,去祝福你的生命.衝突要不至釀成悲劇,我想,必須將自己置放在神的審判之下.

一對雙打的運動員落場打比場,遇著對手多變的戰術,或許,兩個人有不同的反應和策略,或許,兩個人的想法各有道理.這時,最好的做法,不是自己判斷,不是堅持己見,而是問問教練,讓教練指導.神總是我們的教練.

但願:榮耀歸於聖父,聖子,聖靈.阿們.



\newpage

\section{皆大歡喜}
\label{sec:333}
\textbf{林豪恩先生 - 皆大歡喜}
\newline
\newline
\hyperref[sec:332]{< < < PREV SERMON < < <}
~
\hyperref[sec:toc]{[返目錄]}
~
\hyperref[sec:334]{> > > NEXT SERMON > > >}


路加福音 13:10-17
\newline
\begin{longtable}{cl}
\hline
\hline
章節 & 經文 (和合本修訂版)\\
\hline
13:10 & \begin{tabularx}{0.7\textwidth}{X} 安息日，耶穌在一個會堂裡教導人。 \end{tabularx} \\ \\ \relax
13:11 & \begin{tabularx}{0.7\textwidth}{X} 有一個女人被靈附身，病了十八年，腰彎得一點都直不起來。 \end{tabularx} \\ \\ \relax
13:12 & \begin{tabularx}{0.7\textwidth}{X} 耶穌看見，就叫她過來，對她說：「婦人，你的病好了！」 \end{tabularx} \\ \\ \relax
13:13 & \begin{tabularx}{0.7\textwidth}{X} 於是用雙手按著她，她立刻直起腰來，就歸榮耀給神。 \end{tabularx} \\ \\ \relax
13:14 & \begin{tabularx}{0.7\textwidth}{X} 會堂的主管因為耶穌在安息日治病，就很生氣，對眾人說：「有六天應當做工，那六天之內可以來求醫，在安息日卻不可。」 \end{tabularx} \\ \\ \relax
13:15 & \begin{tabularx}{0.7\textwidth}{X} 主回答他：「假冒為善的人哪，難道你們各人在安息日不解開槽上的牛和驢，牽去喝水嗎？ \end{tabularx} \\ \\ \relax
13:16 & \begin{tabularx}{0.7\textwidth}{X} 何況她本是亞伯拉罕的後裔，被撒但捆綁了十八年，不該在安息日這天解開她的綁嗎？」 \end{tabularx} \\ \\ \relax
13:17 & \begin{tabularx}{0.7\textwidth}{X} 耶穌說這些話，他的敵人都慚愧了；所有的人因他所做一切榮耀的事都很歡喜。 \end{tabularx} \\ \\
[1ex]
\hline
\hline
\end{longtable}
講題:皆大歡喜　The Entire Crowd was Rejoicing

經文:路加福音13章10-17節

講員:林豪恩先生

場合:香港中文大學崇基學院禮拜堂

日期:2016年8月21日

當下的皆大歡喜

中文有一成語說:「人同此心,心同此理」,意指對合情合理的事情,大家都有一致的感受與看法.然而,在世道迷離,又動輒得咎的社會中,人們學會了明哲保身,這種「人同此心,心同此理」的表現未必容易被看得見.最近,有一群孩子讓我看見這種赤子之心的表現,置身其中,觸動心靈.

有一天,我陪女兒去觀看免費音樂劇,名叫《開心博士森林記》.開心博士是一個周身法寶的旅行家,他周遊列國,所到之處都會帶給小朋友無限歡樂.在一次探險旅程中,博士經過「寧靜森林」,發現一隻大河馬被困在兩棵巨樹之間的夾縫中,無法動彈.博士相信友誼無分國界,只要有一顆善良的心,幫助別人跨過難關,美好的事便會隨之而來.大河馬在博士出手相助下重獲自由,大夥兒翩翩起舞,「寧靜森林」一時間好不熱鬧!

這樣,問題就來了.原來「寧靜森林」的國王最討厭快樂,他設計了一個「快樂指數」的量度工具,來控制「寧靜森林」的快樂狀況.森林的成員每一次載歌載舞的時候,快樂指數就會上升；他們興奮地玩遊戲的時候,快樂指數又會上升.結果,上升的指數驚動了國王,追查之下,發現為「寧靜森林」帶來快樂的「罪魁禍首」,就是開心博士.

觸動心靈的場面不是演員的表現,而是觀眾的反應.國王命令士兵追捕開心博士,開心博士逃離舞台,走入觀眾席,你們知道觀眾席的孩子們如何反應嗎?孩子們馬上起哄:有些叫開心博士來他們那裡躲藏；有些在士兵查問開心博士的行蹤時指手劃腳,打發士兵離開.最後,當士兵發現開心博士的位置而前來拘捕時,你們能想像孩子們有甚麼行動嗎?在附近的孩子走出座位,站在路上攔截士兵前進.

孩子們在守護甚麼呢?或許,在他們的年紀,他們可能會回答:保護開心博士.日後,他們回想這段經驗,他們就知道,其實他們是在守護一些「人同此心,心同此理」的精神或價值,當中可能包括快樂,和平,正義,友愛等.

在孩子們的保護之下,士兵無功而還,開心博士繼續傳播快樂,全場歡呼,皆大歡喜.孩子們的歡樂,不單單是開心博士的人身安全得到保障,更是他們赤子之心所懷抱的價值得到保存及肯定,也是因為他們能夠體現自己在保護這些心底裡的渴望的付出及努力的成果.

耶穌時代的皆大歡喜

13:10安息日,耶穌在會堂裡教訓人.

13:11有一個女人被鬼附著,病了十八年,腰彎得一點直不起來.

13:12耶穌看見,便叫過她來,對她說:「女人,你脫離這病了!」

13:13於是用兩隻手按著她；她立刻直起腰來,就歸榮耀與上帝.

13:14管會堂的因為耶穌在安息日治病,就氣忿忿地對眾人說:「有六日應當做工；那六日之內可以來求醫,在安息日卻不可.」

13:15主說:「假冒為善的人哪,難道你們各人在安息日不解開槽上的牛,驢,牽去飲嗎?

13:16況且這女人本是亞伯拉罕的後裔,被撒旦捆綁了這十八年,不當在安息日解開她的綁嗎?」

13:17耶穌說這話,他的敵人都慚愧了；眾人因他所行一切榮耀的事,就都歡喜了.

現在,讓我們聆聽另一個故事.

從加利利到耶路撒冷的路上,每逢安息日,耶穌和跟隨者如常參加當地會堂的聚會.被稱為夫子的耶穌,有時被邀在聚會中講道.

安息日聚會有人講道,是恆常事情,恆常事情,通常都不會引起注意.就在大家都以為「無驚無險」又過一天的時候,舞台上突然出現了另一個人.這個人的出現,令恆常的安息日聚會變得不恆常；令昏昏欲睡的聽眾,都抬起頭來；令工作人員,可能包括聚會的司事,會堂的長老等,都緊張起來.

大家不要想像來者是一位手持武器的恐怖份子,相反,被舞台燈射著的,只是一位「手無殺雞之力」的女子.這位女子絕對不似會帶來威脅,她很可能沒有直視群眾,群眾也看不清楚她的面孔,因為她是一位駝著背的病者,駝背,令她長期都抬不起頭來.

這一刻,舞台射燈聚焦在耶穌和這位駝背女子身上.耶穌望著她,她繼續低下頭.在場所有人都屏息以待.

就在這一刻,耶穌出手了.耶穌一出,就是兩隻手,按在那位駝背女子的肩膀上.神奇的事情出現了,那位駝背女子不但沒有被耶穌的雙手按倒,相反,她立刻企直,本來彎曲的脊骨畢直起來.那位十八年來都駝背的女子,終於可以抬起頭來,她第一個反應,就是讚美上帝,感謝上帝.

群眾有甚麼反應呢?有些人可能仍然在驚愕之中,有些人可能在思考應否拍掌,就在大家都未及反應的時候,有人衝上舞台.

這位怒氣沖沖的人是誰呢?原來是會堂的主管.會堂主管,負責監督敬拜聚會的進行,使之不違猶太教的傳統,也有權邀請嘉賓來講道.當時,他親眼目睹嘉賓講員違反律法,感到義憤填膺:「有六日應當做工；那六日之內可以來求醫,在安息日卻不可.」雖然憤怒,會堂主管並不是直接向耶穌說話,而是向群眾說話；他並不是批評耶穌的行動,而是批評那位駝背女子的行動.或許,會堂主管想避免和耶穌直接衝突,只想向他的群眾糾正耶穌的錯.

如果你是在場的群眾,聽過會堂主管的教訓會有甚麼反應呢?有些人可能慶幸自己沒有太早拍掌；有些人可能心裡感到耶穌對那位女士做的是好事,但理性上又同意這事情的確違反安息日的律法.作為普通平民百姓,怎可能有知識和權力去討論那些複雜的律法條文呢?正如城中某高官曾說話:「唔通的士司機,酒樓侍應,麥當勞服務員會〔就二十三條草案〕逐條同我討論?草案都係專家去睇；好似立法會議員,法律界同埋學者.」無奈的現實是,群眾雖然沒有知識和權力去參與律例的討論和制定,但卻可能是最受到律例管控,最沒有能力鑽法律圈子的一群.誰來保護這些平民百姓呢?

現在,舞台上的氣氛相當緊張.會堂主管既然搬律法出來否定耶穌的行動,耶穌就和他進行律法辯論.參加這此安息日聚會的群眾意想不到有這額外的節目,也從未見過有人敢膽公開挑戰會堂主管對律法的詮釋.

13:15「假冒為善的人哪,難道你們各人在安息日不解開槽上的牛,驢,牽去飲嗎?

13:16況且這女人本是亞伯拉罕的後裔,被撒旦捆綁了這十八年,不當在安息日解開她的綁嗎?」

會堂主管和群眾都聽得明白耶穌對安息日律法的詮釋,相信今時今日的讀者也讀得明白.耶穌用了一對重要的詞彙去解釋安息日,其一是解開,其二是捆綁.安息日的精神是解開,而不是捆綁.耶穌畫龍點睛,大家馬上明白,皆大歡喜.

皆大歡喜,是替那位被耶穌醫好的女子高興嗎?可能是,但也不單只這樣?皆大歡喜,很可能是因為群眾心裡的渴望獲得回應,很可能是因為有群眾心底裡的共同價值獲得肯定.或許,群眾心裡面一直疑惑:為甚麼安息日的律例如此繁多,甚至如此令人窒息呢?或許,群眾有時也會反省:耶和華不是有恩典有憐憫的上帝嗎?為甚麼安息日如此不近人情呢?不過,縱使心裡有疑惑,縱使感到不妥,作為群眾,在宗教專家及強大的宗教建制的管理之下,又可以怎樣呢?今天,群眾親眼看見安息日是有恩典和有憐憫的,安息日是讓人自由地去讚美上帝的,群眾心的疑惑終於撥開雲霧見青天,群眾感到的捆綁終於得到釋放.

今天,當你閱讀這個故事的時候,你會認同及投入哪個角色呢?可能有人認為自己是那位長期患病的女士等待醫治；可能有人覺得自己是那位會堂主管堅負著維護傳統,執行律法的重任；可能有人期許自己是耶穌,挑戰僵化了的律法詮釋,還原律法的美善本意.不過,正如故事中的情況一樣,我們當中大多數人都不是站在舞台上,而只是一般群眾,隨著劇情的發展而心情起伏.作為沒有名字,沒有面孔的群眾,我們是否沒有角色呢?例如在安息日的異化過程中,一般平民百姓會否也有份參與其中呢?讓我們一齊閱讀另一個更遠古的故事.

皆大歡喜的異化

按出埃及記的記載,在安息日「你和你的兒女,奴僕,婢女,牲畜,以及你城裡寄居的客旅,都不可做任何的工.」(20:10)作者給與一個神學的理由:「因為六日之內,耶和華造天,地,海和其中的萬物,第七日就安息了；所以耶和華賜福與安息日,定為聖日.」(20:11)在申命記中,安息日被賦予多一種歷史意義:「你要記念你在埃及地作過奴僕,耶和華－你的上帝用大能的手和伸出來的膀臂領你從那裡出來.因此,耶和華－你的上帝吩咐你守安息日.(5:15)因此,安息日的精神,是上帝的對生命的一種祝福與賞賜,也是對人和整個生存環境拯救的恩典.

上帝對生命的祝福與賞賜,以及對人和生存環境拯救的恩典,都是美好的價值.如何把這些美好的價值變成具體行為,並成為群體的規範,問題就不簡單了.有時,我們不經意地製造了作繭自縛的困局.

讓我們設想一群小朋友上課的處境:

老師說,各位同學,上課時要保持安靜.

過了一會,阿一同學叫道:老師,阿二和和阿三沒有保持安靜,他們輕聲談話.

阿二回應:老師,阿一也沒有保持安靜,他沒有舉手就叫出來.

於是老師說:保持安靜就是不可以談話,也不可以沒有舉手就叫出來.

又過了一會,又聽到阿一的聲音:老師,阿四和阿五……

阿一沒有說完,老師就接過來:阿一,我不是說過說話前要先舉手的嗎?

阿一委屈地說:我有舉手,只是你看不見而已.

老師唯有進一步指示:各位同學,舉了手之後,還要等老師批准了才可以說話.

又過了一會,阿二舉手,老師示意阿二可以說話:老師,阿六沒有保持安靜,他老是把間尺擦桌面.

老師回應說:阿六,上課要保持安靜,不可以製造聲音.

又過一會,阿六舉手,獲批說話:老師,阿七剛才打呵欠,製造聲音……

如此類推,有人向老師揭發有同學的椅子發出聲音,有人告訴老師有同學的筆盒發出聲音,有人……

每一項被揭發出來的事情,都的確有發出聲音,打破安靜,那麼,老師就只能不停地說:不可以,不可以,不可以……

後來,課堂上突然傳來一聲巨嚮:呠…呠…呠…所有同學都屏息以待,有一位同學知道藏不了,非常尷尬地說:老師,我知道上課時不可以發出聲音,但我實在忍不住了……

以下劇情純屬虛構,如有雷同,實屬巧合.不過,類似的巧合的劇情有時發生在現實世界中(例如法國人爭論應該穿甚麼泳衣去沙灘),也發生在教會中(例如爭論參加崇拜時應該穿甚麼衣服).

在很久很久之前,當以色列人聽到在安息日甚麼工都不可以作的規定後,就發生了以下這個記載在民數記的故事,我們可以想像,在安息日律例的發展歷史中,這樣的故事可能只是冰山一角.

15:32以色列人還在曠野的時候,發現有一個人在安息日撿柴.

15:33發現他撿柴的人把他帶到摩西,亞倫以及全會眾那裡.

15:34他們把他收押在監裡,因為還不知道要怎樣懲罰他.

15:35耶和華吩咐摩西說:「這人應當處死；全會眾要在營外用石頭打死他.」

15:36於是全會眾把他帶到營外,用石頭打死他,是照耶和華所吩咐摩西的.

美好事物的變質過程,總會有很多因素.作為群眾的我們,如果以為自己只是受害者,可能忽略了自己的參與及角色.

最近,有新聞報導說:

大欖郊野公園風景怡人,不少市民都會帶同子女前來遠足和燒烤,親親大自然,而一直在路邊生長,一觸碰即合上的含羞草群,既為小童們的「寵兒」,同時也是教育子女的靈丹妙藥.

有報章記者發現,漁護署把含羞草群連根拔起,移到別處重新種植,並裝上圍欄.然而,被移植的含羞草可能水士不服,幾乎全部死去.

常識告訴大家,漁護署職員怎麼會主動走去大欖郊野公園清除含羞草呢?原來,漁護署收到家長投訴,指其子女遭含羞草刺傷,故署方將該處的含羞草連根拔起,並於公園另覓地點重新種植.惟移植後設有圍欄,防止再有兒童被含羞草刺傷.(都市日報3/8/2016)

漁護處的處理方法是否恰當,當然可以討論.然而,市民遊覽大欖郊野公園的時候再看不見含羞草,這樣的結果,我們這些平民百姓是有角色在其中的.

總結

下次去大欖郊野公園還能否看到含羞草,尚算小事.對耶穌來說,基督教信仰還有沒有「上帝對生命的祝福與賞賜,對人和生存環境拯救的恩典」這些美好價值,還是已經異化成為沒有人,也沒有上帝,只有一堆捆綁人的律例規範,這是他不能忘懷的.

在往耶路撒冷的路上,耶穌透過在安息日醫治駝背女子的行動,還原了安息日的美意,皆大歡喜.原來,在人們的心底裡,在最赤誠的地方,仍然渴望美善,也對美善有回應.就正如想觀看《開心博士森林記》的那群孩子,本能地挺身而出保護開心博士,守護快樂,和平,正義,友愛等美好價值.但願作為成人的我們好好守護他們,好讓他們可能看見世界紛亂,仍能保存赤子之心.



\newpage

\section{潤物細無聲}
\label{sec:334}
\textbf{林豪恩先生 - 潤物細無聲}
\newline
\newline
\hyperref[sec:333]{< < < PREV SERMON < < <}
~
\hyperref[sec:toc]{[返目錄]}
~
\hyperref[sec:335]{> > > NEXT SERMON > > >}


講題:潤物細無聲　Silent and Soft, Life to Rise

經文:詩篇1

講員:林豪恩先生

場合:香港中文大學崇基學院禮拜堂

日期:2016年9月4日

一． 引言

詩篇第一篇是比較特別的,它不像其他多數詩篇一樣,在崇拜或典禮中作為儀文來運用,它被放在篇首,作為整部結集的引言.這首詩的第一個字可譯為「幸福」或「快樂」,這個字在第二篇最後一句也出現.因此,有學者認為詩篇第一篇及第二篇合起來是詩篇的雙重引言.第一篇以智慧文學的形式書寫,是一首智慧詩,第二篇以先知文學的傳統表達,共同的地方是「幸福／快樂」這主題.

二． 釋經及應用

1．虛則實之

我們可以把這首詩分為三個部分來閱讀,第一部分是第一及第二節:

1:1不從惡人的計謀,不站罪人的道路,不坐褻慢人的座位,

1:2惟喜愛耶和華的律法,晝夜思想,這人便為有福!

本詩原文的第一個字並非「不」字,而是和合本譯本在第二節的最後一個字「有福」.新譯本把原文的第一個字放回原位,就容易讀出詩歌的題旨.

1有福的人:不從惡人的計謀,不站罪人的道路,不坐好譏笑的人的座位.

2他喜愛的是耶和華的律法,他晝夜默誦的也是耶和華的律法.

有福的人是怎樣的呢?或者,如何得到幸福和快樂呢?幸福和快樂好像是抽象的東西,捉不實,抓不緊,因此,人們通常都透過擁有一些捉得到,又抓得緊的東西來來體驗幸福快樂.

詩人提出「喜愛耶和華的律法,晝夜思想」,就會得到幸福和快樂,這樣的建議,難免令人感到虛幻不實,難以掌握.相對來說,作者所否定的「惡人的計謀,罪人的道路,褻慢人的座位」則實際得多,有計劃,有方法,可掌握.

詩歌常用平衡句及對比句,惡人,罪人和褻慢人三者雖然有差異,但讀者明白所指的都是和下句「喜愛的是耶和華的律法」的人相對的；計謀,道路和座位三者的涉入程度越來越深,但讀者也明白所說的大概也是與下句「晝夜思想」耶和華律法相反對狀況.至於耶和華的律法,狹義指舊約聖經被稱為律法書或摩西五經的首五卷,廣義可說是上主創造天地萬物的原意,以及交托給人類的恩賜和使命.有人稱之為「天道」,「天道」在不同的傳統有不同的內涵,但大體上都是表示終極真理,本原,本體,規律,原理,境界等等,指運作永恆一切的道.香港有一間基督教出版社名為「天道書樓」,所指應該就是上帝的道.

理想地說,人們都是擇善去惡的.例如人們期許:宗教都是導人向善的；教育也是推崇「德智體群美」的.不過,在現實操作上,有時感到「天道」,「善」,「義」這些指導未能幫助我們達到目標,未能令我們得到我們想要的東西,也未能令我們稱雄稱霸；或者,最深層次的猶豫,是「天道」,「善」,「義」這些指導未能保障我們的生存,因為它們太清高,太抽象,太虛幻,總意就是「不能當飯吃」.為了生存,我們總要捉老鼠,因此,不理白貓還是黑貓,總之捉到老鼠的就是好貓.對於我們來說,如果這是童話世界,我們一定選擇「耶和華的律法」,而不會是「惡人的計謀,罪人的道路,好譏笑的人的座位」,然而,當我們要戰鬥,要打敗對手,要達到目的,要遂其心願的話,我們就可能「非常時期用非常方法」了.

普通話團契八月的其中一個聚會看了一套電影,名叫Chappie,中文有譯作「超能查派」,又有譯為「成人世界」.故事說有一間公司為南非某城市生產了機械人警察,成功協助該城市收復失控的犯罪狀況.設計機械人的工程師再進一步設計出全世界第一個有意識及情感的人工智慧,然後試驗性地放進其中一個警察機器人中,這個有意的機械人被稱為Chappie.因緣際會之下Chappie落入罪犯手中,這個像新生嬰兒一樣對世界感到畏懼而好奇的人工智慧不得不在周圍成人的複雜世界中找出自己生存的方式.創造他的工程師教導他要有禮貌,也不能犯罪,不能殺人,所以,當罪犯教他使用槍械射擊紙板人的時候,他拒絕.後來,罪犯稍為改變方法,在他面前示範使用各種武器,攻擊各種非紙板人目標,他看到罪犯們如何威風,也跟著學,看到自己成功擊中目標,他感到興奮及威風.下一個鏡頭,Chappie已經跟足罪犯的打扮和言語,身上配備武器,大搖大擺地隨罪犯出發.罪犯如何能夠令Chappie協助他們去犯案呢?因為Chappie的電力供應有時限,要生存下去,就必需要解決這個問題.罪犯告訴他,如果要解決這個生存下去的問題,就要有大量金錢.面對生存問題,設計師交給Chappie的原則-不可以犯罪,不可以殺人,就顯得軟弱無力,「惡人的計謀,罪人的道路,好譏笑的人的座位」似乎是較為現實的可行出路.

美善的原則,價值,精神,是否太虛弱呢?「惡人的計謀,罪人的道路,好譏笑的人的座位」會否真正實在呢?

2．實則虛之

第二部份是第三至第五節,作者運用比喻,比較上文所說有福的人和惡人結局的差異:

1:3他要像一棵樹栽在溪水旁,按時候結果子,葉子也不枯乾.凡他所做的盡都順利.

1:4惡人並不是這樣,乃像糠秕被風吹散.

1:5因此,當審判的時候,惡人必站立不住；罪人在義人的會中也是如此.

有謂:實則虛之,虛則實之.作者描繪了兩種景象.一種是第四節「糠秕被風吹散」的景象.「惡人的計謀,罪人的道路,好譏笑的人的座位」曾經好像很具體,很實際,很威風,但如此實在的東西,作者指出其實是虛幻的.另一種景象是第三節「一棵樹栽在溪水旁,按時候結果子,葉子也不枯乾」,用來代表上文那些「喜愛耶和華的律法,晝夜思想」的人.「盡都順利」是指那棵樹,自然生長,果實櫐櫐,而不是指人一帆風順和不會遇到困難.風吹雨打在所難免,那棵樹仍然能夠自然生長,按時結果,作者暗示了其中一項重要條件,就是「栽在溪水旁」.溪水緩緩地流經,「無言無語」,甚至「也無聲音可聽」,四季不息,默默地滋潤著那棵樹,讓那棵樹活出它被創造的美好.

按詩意,滋潤著樹的溪水,可聯想到上文「耶和華的律法」.「晝夜思想」,扎根在神話語的泉源中的果效,就正如那樹長年累月被溪水潤澤所得到成長.

第六節是結束,為上述的差異提出了一個信仰上的理由.

1:6因為耶和華知道義人的道路；惡人的道路卻必滅亡.

作者提示,按天道而行,天會看顧及確認其道,同時否定惡人的路.按耶和華上帝創造天地萬物的原意,以及交托給人類的恩賜和使命而行,創造者會看顧並確認其創造的美意,而否定與其相反的惡念和惡行.

3．重思獲得幸福／快樂的途徑

讓我們回到詩篇第一篇的主旨,如何獲得幸福／快樂呢?作者在詩中暗示,幸福／快樂不是上帝賜予的,而是人在生活中經驗的.怎樣的生活是幸福快樂的呢?就是「不從惡人的計謀,不站罪人的道路,不坐好譏笑的人的座位」,取而代之的是「喜愛耶和華的律法,晝夜思想」.作者指出,耶和華的律法如溪水滋養樹木成長一樣,潤澤我們的生命.不過,人們的困難往往輕視對生命的滋養潤澤,另一方面,心裡的焦慮及外在的聲音常常鞭策我們,若要幸福快樂,就要強大,要富有,要操控.所以,有些國家告訴她的人民,不惜破壞自然環境的代價都要追求成為經濟大國；「不要褲子,也要核子」也要追求成為軍事強國.他們要人民相信並接受:國家強大,人民才會幸福快樂.事實上,世界上最強大的國家,他們的人民是否最幸福快樂呢?

有一位德國人,期望奉獻一己之力,讓所有的人們都能同享快樂的祕訣,於二○一三年走訪不同年度世界最幸福／快樂國家調查中排首位的十三個國家,實地訪談當地居民,以及各國研究幸福相關議題的權威學者.寫成了一本書,名為:「為甚麼我們最幸福-最幸福的十三個國家」.作者發現,第一,那些世界上數一數二的經濟大國和軍事強國,都不是最幸福或快樂的國家；第二,令人感到幸福或快樂的,都不是因為經濟和軍事強大,而是因為價值和觀念這些聽似虛無,摸不到,捉不實的東西.例如:

看不見的富裕祕境──挪威

幸福關鍵字:信任

我們的社會政策試著把所有人都納入考慮,希望藉此縮小貧富差距.這樣的方式帶給人與人之間充分的信任感.在這裡,你可以相信其他人是可靠的,不必隨時提心吊膽,擔心被騙,這樣真好!一個人幸福與否和身邊的人都有關係,如果你信任別人,就可以活得更自在.

從性別到階級都平等的世界──瑞典

幸福關鍵字:凡事都要剛剛好

瑞典語中,「Lagom」意謂著凡事不要太匆忙,但也不是過度緩慢；沒有過多階級概念,可是每個人都知道,最終是由來做決定的.或許這裡多一點女性溫柔,那裡多一點陽剛之氣,但都是有分寸地行事.

還有一些治安不是那麼好的國家,例如墨西哥；甚至經濟不甚發達的國家,例如哥斯大黎加等.如果不是強大,是甚麼滋養潤澤著人們的生命,使人感到幸福快樂呢?令挪威人感到幸福快樂的是「信任」；令瑞典人感到幸福和快樂的是「凡事都要剛剛好」；令墨西哥人感到幸福快樂的是「人生有意義」.原來,幸福快樂與否,與該群體的核心價值有關,有些國家的核心價值是令人們疏離,對立,焦慮和不安的,這些核心價值摧殘生命；另一些國家的核心價值令人與互相信任,令社群互相包容,互相接納,互相幫助,這些核心價值滋養潤澤著人們的生命,令人感到幸福快樂.有些地方的核心價值是自求多福,每個人都爭取資源,務求使自己脫離公共醫療體系,務求使自己的子女逃出本地的教育

系統,務求使自己超越食物安全製度……結果,無論是逃得脫的人,或者是逃不脫的人,都生活在不安和不滿之中.

三． 總結

那棵樹,在溪水的滋潤中生長結果.我們的心田,又有甚麼滋潤著呢?我們的群體,有甚麼核心價值滋潤著呢?詩人說:「喜愛耶和華的律法,晝夜思想,這人便為有福」.創造主的創造美意,能否成為潤澤我們的心田的良方呢?能否成為我們群體的核心價值呢?耶穌說,律法的總綱是愛,讓我們看看愛如何滋養人生:

著名作學冰心的文字滋潤著萬千的心靈,她的心靈又如何被滋養的呢?在「母親」一文中,她如此說:

有一次,幼小的我,忽然走到母親面前,仰著臉問說:「媽媽,你到底為甚麼愛我?」母親放下針線,用她的面頰,抵住我的前額,溫柔地,不遲疑地說:「不為甚麼,──只因你是我的女兒!」

這句說話,滋養著冰心幼小的心靈,她心裡有一項確認,這項確認又繼續成為她一生的滋潤:

假使我走至幕後,將我二十年的歷史和一切都變更了,再走出到她面前,世界上縱沒有一個人認識我,只要我仍是她的女兒,她就仍用她堅強無盡的愛來包圍我.她愛我的肉體,她愛我的靈魂,她愛我前後左右,過去,將來,現在的一切!

這樣的體會,讓冰心對「世界是如何建造起來的」這樣一個深奧的神哲學問題有如下的總結:

她的愛不但包圍我,而且普遍的包圍著一切愛我的人；而且因著愛我,她也愛了天下的兒女,她更愛了天下的母親.小朋友!告訴你一句小孩子以為是極淺顯,而大人們以為是極高深的話:「世界便是這樣的建造起來的!」

或者有人認為世界是由物質建造的,或者有人認為世界是由經濟力量建造的,或者有人認為世界是由權力建造的,在母親的愛的滋養之下,冰心認定世界是由愛建造起來的,是那種滋養潤澤生命的愛.

或許,我們父母在各種限制之下,在我們的成長過程中未能如冰心的母親一樣滋潤我們的心田；或者,當我們細看自己的已經長大成人的心田,仍然發現乾涸,甚至龜裂?有甚麼能滋潤我們的心田呢?使徒約翰說:

4:16我們知道並且深信上帝是愛我們的.上帝就是愛,住在愛裡面的就是住在上帝裡面；上帝也住在他裡面.(約翰一書)

詩篇第一篇提示我們,上帝美善的心意如溪水滋養樹木一樣,潤澤著人們通向幸福快樂的旅途.詩篇十九篇的作者進一步指出:上帝的話語:「比蜜甘甜,比蜂房下滴的蜜甘甜」(19:10).願上主恩待我們,讓我們在祂的話語中嚐到甘甜,心靈得到滋潤.



\newpage

\section{人飢餓非因無餅}
\label{sec:335}
\textbf{鄧瑞強博士 - 人飢餓非因無餅}
\newline
\newline
\hyperref[sec:334]{< < < PREV SERMON < < <}
~
\hyperref[sec:toc]{[返目錄]}
~
\hyperref[sec:336]{> > > NEXT SERMON > > >}


阿摩司書 8:1-12
\newline
\begin{longtable}{cl}
\hline
\hline
章節 & 經文 (和合本修訂版)\\
\hline
8:1 & \begin{tabularx}{0.7\textwidth}{X} 主耶和華又指示我一件事，看哪，有一筐夏天的果子。 \end{tabularx} \\ \\ \relax
8:2 & \begin{tabularx}{0.7\textwidth}{X} 他說：「阿摩司，你看見甚麼？」我說：「一筐夏天的果子。」耶和華對我說：「我子民以色列的結局到了，我必不再寬恕他們。 \end{tabularx} \\ \\ \relax
8:3 & \begin{tabularx}{0.7\textwidth}{X} 那日，宮殿裡的詩歌要變為哀號；必有許多屍首拋在各處，安靜無聲。」這是主耶和華說的。 \end{tabularx} \\ \\ \relax
8:4 & \begin{tabularx}{0.7\textwidth}{X} 你們這些踐踏貧窮人、使這地困苦人衰敗的，當聽這話！ \end{tabularx} \\ \\ \relax
8:5 & \begin{tabularx}{0.7\textwidth}{X} 你們說：「初一幾時過去，我們好賣糧；安息日幾時過去，我們好擺開穀物；我們要把伊法變小，把舍客勒變大，以詭詐的天平欺哄人， \end{tabularx} \\ \\ \relax
8:6 & \begin{tabularx}{0.7\textwidth}{X} 用銀子買貧寒人，以一雙鞋換貧窮人，把壞的穀物賣給人。」 \end{tabularx} \\ \\ \relax
8:7 & \begin{tabularx}{0.7\textwidth}{X} 耶和華指著雅各的驕傲起誓說：「他們這一切的行為，我必永遠不忘。 \end{tabularx} \\ \\ \relax
8:8 & \begin{tabularx}{0.7\textwidth}{X} 地豈不因這事震動？其中的居民豈不悲哀嗎？全地必如尼羅河漲起，如埃及的尼羅河湧起退落。 \end{tabularx} \\ \\ \relax
8:9 & \begin{tabularx}{0.7\textwidth}{X} 「到那日，我要使太陽在正午落下，使這地在白晝黑暗。」這是主耶和華說的。 \end{tabularx} \\ \\ \relax
8:10 & \begin{tabularx}{0.7\textwidth}{X} 「我要使你們的節期變為悲哀，你們一切的歌曲變為哀歌；我要使眾人腰束麻布，頭上光禿；我要使這悲哀如喪獨子，其結局如悲痛的日子。 \end{tabularx} \\ \\ \relax
8:11 & \begin{tabularx}{0.7\textwidth}{X} 「看哪，日子將到，我必命饑荒降在地上；人飢餓非因無餅，乾渴非因無水，而是因不聽耶和華的話。」這是主耶和華說的。 \end{tabularx} \\ \\ \relax
8:12 & \begin{tabularx}{0.7\textwidth}{X} 他們必飄流，從這海到那海，從北邊到東邊，往來奔跑，尋求耶和華的話，卻尋不著。 \end{tabularx} \\ \\
[1ex]
\hline
\hline
\end{longtable}
講題:人飢餓非因無餅　Not a Famine of Bread

經文:阿摩司書8:1-12

講員:鄧瑞強博士

場合:香港中文大學崇基學院禮拜堂

日期:2016年9月18日

各位弟兄姊妹,早安.

中秋節雖然過了,仍祝願各人家人常團聚,生命得圓滿.

蘋果電腦創辦人喬布斯(Steve Jobs)有一次在史丹福大學的畢業禮上演講,講到人不應盲從別人的想法,不應活別人的生活,而應敢於去實現自己的夢想.這演講最後兩句話是 “Stay Hungry. Stay Foolish.” 若將這兩句話抽離喬布斯的演講,則這兩句話差不多變得無意義.因為根本不用Stay Foolish,很多人本身就是Foolish；也根本不用Stay Hungry,很多人本性已是飢餓無比,甚至飢不擇食.世上最差的人,可能就是某些Hungry and Foolish的人,他們對權力飢渴,卻又愚不可及,最後禍己害人,遺臭萬年.

我們會否是飢餓而愚蠢的人?

人有多層次,多種類的飢餓,渴望.我們會肚餓,會口渴,想呼吸,這些都是身體上的飢餓.我們不想被人無端端拘捕,不想打仗,不想死,不想辛苦賺來的錢被人搶去,不想有病,這些是我們對安全感的渴求.我們渴望朋友,渴望有所歸屬,渴望擁抱,渴望有人一起賞月,這些是我們對關係的渴望,對歸屬的渴望.失去一切關係,我們可能會餓得精神失常.我們也渴求別人的讚賞,你唱歌時渴望有知音人,你講的話希望有人用心聆聽,這是我們對榮譽的渴望,對認同的渴望.還有,我們渴望實現自己的夢想,有人花八年造一隻會飛的飛機而感到無比滿足,有人花十年苦練而將世界紀錄推前零點一秒而興奮莫名.有人願意走過千山萬水,就是希望這種人性飢餓得到滿足.

人有種種飢餓,問題只是,我們是否以愚蠢的方式去滿足它們?

我們看看阿摩司先知對我們的提醒.

今日的講道經文是:阿摩司書8:1－12

摩8:1　主耶和華又指示我一件事:我看見一筐夏天的果子.

摩8:2　他說:「阿摩司啊,你看見甚麼?」我說:「看見一筐夏天的果子.」耶和華說:「我民以色列的結局到了,我必不再寬恕他們.」

摩8:3　主耶和華說:「那日,殿中的詩歌變為哀號；必有許多屍首在各處拋棄,無人作聲.」

摩8:4　你們這些要吞吃窮乏人,使困苦人衰敗的,當聽我的話!

摩8:5　你們說:月朔幾時過去,我們好賣糧；安息日幾時過去,我們好擺開麥子；賣出用小升斗,收銀用大戥子,用詭詐的天平欺哄人,

摩8:6　好用銀子買貧寒人,用一雙鞋換窮乏人,將壞了的麥子賣給人.

摩8:7　耶和華指著雅各的榮耀起誓說:他們的一切行為,我必永遠不忘.

摩8:8　地豈不因這事震動?其上的居民不也悲哀嗎?地必全然像尼羅河漲起,如同埃及河湧上落下.

摩8:9　主耶和華說:到那日,我必使日頭在午間落下,使地在白晝黑暗.

摩8:10　我必使你們的節期變為悲哀,歌曲變為哀歌.眾人腰束麻布,頭上光禿,使這場悲哀如喪獨生子,至終如痛苦的日子一樣.

摩8:11　主耶和華說:日子將到,我必命饑荒降在地上.人飢餓非因無餅,乾渴非因無水,乃因不聽耶和華的話.

摩8:12　他們必飄流,從這海到那海,從北邊到東邊,往來奔跑,尋求耶和華的話,卻尋不著.

阿摩司先知生活的時代,北國以色列在耶羅波安二世統治下,進入以色列國少有的強盛年代.政治穩定,經濟繁榮,軍事力量強大,國際貿易頻繁.這造就了很多有錢人,有錢人過著奢華生活.為了維持奢華生活,這些有錢人強力地欺壓窮人,藉此獲得厚利.社會越富裕,窮人越可憐,正所謂「朱門酒肉臭,路有凍死骨」.就在這光景裡,上帝發出祂的聲音.

剛才讀的那段經文,以明確的宣告作開始,以色列的結局到了.

摩8:1　主耶和華又指示我一件事:我看見一筐夏天的果子.

摩8:2　他說:「阿摩司啊,你看見甚麼?」我說:「看見一筐夏天的果子.」耶和華說:「我民以色列的結局到了,我必不再寬恕他們.」

摩8:3　主耶和華說:「那日,殿中的詩歌變為哀號；必有許多屍首在各處拋棄,無人作聲.」

「夏天的果子」(希伯來文:qayits)和「結局」(希伯來文:qets)這兩個字,在希伯來文的讀音相近.上帝叫先知看,先知看見「一筐夏天的果子」,這意味「結局」.上帝說:「我民以色列的結局到了,我必不再寬恕他們.」何等斬釘截鐵,沒有寬容餘地.一般而言,上帝總是寬容再寬容的,如此決絕,神的子民一定是作了大逆不道的事了.上帝繼續說:「那日,殿中的詩歌變為哀號；必有許多屍首在各處拋棄,無人作聲.」「無人作聲」是感嘆詞,譯得傳神一點,可以譯為「噓,噓,不要作聲!」上帝的審判已決,祂的行動已定,還叫先知不要出聲,以免打草驚蛇.

以色列人作了什麼大逆不道的事?

摩8:4　你們這些要吞吃窮乏人,使困苦人衰敗的,當聽我的話!

摩8:5　你們說:月朔幾時過去,我們好賣糧；安息日幾時過去,我們好擺開麥子；賣出用小升斗,收銀用大戥子,用詭詐的天平欺哄人,

摩8:6　好用銀子買貧寒人,用一雙鞋換窮乏人,將壞了的麥子賣給人.

摩8:7　耶和華指著雅各的榮耀起誓說:他們的一切行為,我必永遠不忘.

以色列人吞吃窮人,殘踏弱者.

他們說:「月朔幾時過去,我們好賣糧；安息日幾時過去,我們好擺開麥子」.

月朔即初一,初一和安息日,都是宗教節期,是獻祭的日子.但那些以色列人,守節的時候,心中想著的,卻是生意,賺錢.他們沒有與民眾一同在上帝面前歡慶生命,而是想著將民眾變成自己魚肉的對象.親近上帝一刻也嫌太久,生意做少一點卻嫌太多.

虛偽的信仰,表現在不正義的賺錢態度上.

「賣出用小升斗,收銀用大戥子,用詭詐的天平欺哄人」,「將壞了的麥子賣給人」,他們一點也不公正.賣出的時候,將量器弄小；收錢的時候,卻用加重了的法碼.天平是假的,麥子也摻雜次貨.收的,是人參的價錢；給你的,卻只是普通的樹根.人有賺錢的渴望,這是正常的,但不能用弄虛作假的方式.

賺了錢後,他們有更大資源強化欺壓人的制度.

摩8:6　好用銀子買貧寒人,用一雙鞋換窮乏人.

人窮,就需要借貸.很多時,小小債項,就叫窮人不能翻生.這裡,有錢人買起窮人作奴隸.他們付出的,只是少許金錢.這是以自己的成功,剝削他人的尊嚴.沒有人與人之間的平等,而是強化階級之間的分野.不是分享富裕,而是盡情搜刮.巧取豪奪地將別人變成奴隸,而不是渴望別人自由.將他人商品化,「物化」,而不是賦予他人神聖的價值.在這個社會,沒有人與人之間的友愛,只有競爭,將他人踐踏在腳下.

有一個致力消除現代奴隸制度的全球組織,叫 “Walk Free Foundation”,在今年五月發表報告,估計全球約有4千5百萬人陷在現代的奴役中.香港約有29,500人算是現代奴隸,按人口比例而言,我們的情況好過埃塞俄比亞和肯雅,但差過菲律賓和西非的加納.(參:\url{http://www.globalslaveryindex.org/findings/})這些人多數是被嚴重剝削的外籍傭工或被騙來港賣淫的人.在香港,久不久會有虐待外傭的新聞,這是將人「物化」的極端表現.

當時的以色列人,過著虛假的信仰生活,做著騙人的買賣,維持著非人化的社會經濟關係.宗教的,經濟的,社會政治的面貌,都全面腐化.對上帝,沒有半點敬畏.對公平,沒有半點尊重.對窮人,沒有半點仁慈.這簡直就是一個無法無天的世界.

這種腐敗情況,在歷史裡周而復始地出現,像夏天的蚊一樣,老是纏著人不放.這種腐敗情況,驅之不散,因為它的根源,來自人性的飢餓.我們想食飯,這是很自然的.吃飽後,我們想吃得好過人.我們不想只吃白飯,我們想吃鮑魚.要吃鮑魚,我們得賺錢多過人.賺錢多過人,又會被別人讚許.被別人稱讚,又是我們的一種基本渴望.被別人讚許後,我們想贏盡所有掌聲.贏盡所有掌聲,我們才多少感到實現了自我.實現自我,是我們心底最深的渴求.人性的飢餓驅使我們一步一步的向前走.當這種飢餓以愚蠢的方式去滿足時,當又hungry又foolish時,便會視自己為神,而視其他人為奴隸,造成一個無法無天的世界.

上帝無動於衷嗎?

摩8:7　耶和華指著雅各的榮耀起誓說:「他們的一切行為,我必永遠不忘.」

實在太好了,神在看,神在審判,神在行動.

這是受苦者最後的希望.

上帝與人立約,以天地作證.當人不仁不義時,天地也為之震怒.

摩8:8　地豈不因這事震動?其上的居民不也悲哀嗎?地必全然像尼羅河漲起,如同埃及河湧上落下.

摩8:9　主耶和華說:到那日,我必使日頭在午間落下,使地在白晝黑暗.

摩8:10　我必使你們的節期變為悲哀,歌曲變為哀歌.眾人腰束麻布,頭上光禿,使這場悲哀如喪獨生子,至終如痛苦的日子一樣.

地大震動,河流翻騰,白日無光,是因為我們顛倒了愛與正義的秩序.上帝的審判速速降臨,對受壓的貧窮人而言,這便是福音.

摩8:11　主耶和華說:「日子將到,我必命饑荒降在地上.人飢餓非因無餅,乾渴非因無水,乃因不聽耶和華的話(更好譯為:乃因聽不到耶和華的話).

摩8:12　他們必飄流,從這海到那海,從北邊到東邊,往來奔跑,尋求耶和華的話,卻尋不著.」

人有種種飢餓,這是真的,但人真正的飢餓,非因無餅；人真正的乾渴,非因無水.人飢餓乾渴,是因為我們聽不到神的說話,神的真理.我們聽不到,不是因為上帝沒有說,而是我們不再願意聽.現代人像戴著各種聽筒,聽自己喜歡聽的東西,除此以外,再聽不入其他聲音了,遑論上帝的聲音了.

上帝造人,有種種飢渴.C.S.Lewis說:人最大的渴望,是像神一樣.人渴望在不需要上帝的情況下,為自己爭取最大的幸福.金錢,野心,戰爭,階級,帝國,奴隸制度等等,陳列在人類的歷史上,展示我們是怎樣在上帝以外去創造自己的幸福的.上帝造我們的時候,是以祂自己來滿足我們的飢渴的.在上帝以外,根本沒有令人飽足的東西的存在.(參:C.S.Lewis: Mere Christianity, Section 2.3, “The Shocking Alternative”,任何版本.)

耶羅波安二世的統治年代,約為公元前786-746.阿摩司先知提及的大地震,約發生在公元前760年,約在耶羅波安二世統治的中期.地震的警告,顯然對以色列國沒有產生警示作用.地震後約40年,以色列於公元前721年滅國.

人總有種種飢餓,人總會千方百計滿足自己的飢餓.人總是hungry的,但希望我們不要foolish.先知說:人真正的飢餓,非因無餅；人真正的乾渴,非因無水.人飢餓乾渴,是因為我們聽不到神的說話,神的真理.

要明白我們為何飢餓,也要明白胡亂地滿足種種飢餓是何等愚昧.我的結語是:Why Hungry ? How Foolish !

但願:榮耀歸於聖父,聖子,聖靈.阿們.

(後話:本人將早期的崇基講章結集為書,書名《敢於跟隨主》,基道出版,有興趣者可購買.)



\newpage

\section{不做壞事無做好事}
\label{sec:336}
\textbf{許開明牧師 - 不做壞事, 無做好事?}
\newline
\newline
\hyperref[sec:335]{< < < PREV SERMON < < <}
~
\hyperref[sec:toc]{[返目錄]}
~
\hyperref[sec:337]{> > > NEXT SERMON > > >}


講題:不作壞事,亦無做好事　Did Nothing Evil, Do No Good!

講員:許開明牧師

場合:香港中文大學崇基學院禮拜堂

日期:2016年9月25日

引言

許多人有這種人生觀:認為我「不會作壞事,亦無做好事」 也可以接受.好像聖經中,這個財主並不是一個窮兇惡極,欺壓窮人的人,他容許一個滿身是瘡的乞丐在他的門口討食,他無作壞事啊,只不過,他也無做善事或好事,使討飯的得飯吃,得衣服穿,得入屋取暖.這種人,不算好,也不壞啊!

然而,中國人有 「止於至善」之追尋,「止於至善」出自儒家之主要典籍《大學》.《大學》乃「初學入德之門」(北宋儒者程顥語),開首即揭示德性之源及其展開與圓滿之義,即無休止的進步,成德之高,成德之堅永,盡在這「圓」與「未圓」之間.「圓」與「未圓」互相映發,自知「未圓」,故能虛懷若谷；力求「圓滿」,故得至健如天(「至健」語出《易經》).「止於至善」蘊含至健與虛懷,同時兼有「不斷追求圓滿」的自強不息精神以及「不自以為圓滿」的開放精神,這與聖經教導不謀而合.

天堂門口

一位弟兄離世後,上到天堂門口,正想進入,門口天使叫他等等,要先看看進天堂人士名單,發現並沒有弟兄的名字.這弟兄本以為自己從來無做壞事,一定可以進入天堂,就詢問天使原因,天使查完說:「你無做壞事,但亦無做好事,仍然不符合資格,進不到天堂.」

1. 地上(陽間 )兩種人的情景

1.1. 財主

「有一個財主 (『戴維斯』(Dives,源於拉丁文『富豪』一詞)),穿著紫色袍和細麻布衣服,天天奢華宴樂(大魚大肉,吃喝玩樂).」(路十六19)

由「財主」一詞代表這個人有財富；「穿著紫色袍和細麻布衣服」即代表這個人有身份,因為「紫色袍」為王室和尊貴人的標誌 (參斯八15,箴卅一22,啟十八12),而「細麻布衣服」亦是用細麻織成的布料所製作的名貴衣服；「天天奢華宴樂」即代表這個人有很多時間 ,「奢華宴樂」在原文含有「不但吃的是山珍海味,並且講究華麗的排場,極盡享樂之能事」的意思.這個人穿得好,食得好,玩得好,有的是一種自我專注,自得其樂,自私自利的人生.自古到今,無數人追求的就是這種人生.猶太人更有一種人生哲學,一個人有錢,成為富翁是上主的祝福.

1.2. 拉撒路

「又有一個討飯的,名叫拉撒路,渾身生瘡(Covered with sores全身被瘡所遮蓋),被人放在(他不能行進去) 財主門口,要得財主桌子上掉下來的零碎充饑；並且狗來舔他的瘡.」(路十六20－21)

「討飯的」,代表拉撒路活在飢餓中；「拉撒路」的意思是神所幫助的,他是一個在苦難中等候神和人幫助的人.「渾身生瘡」即代表他活在羞辱中,外表難看加上肉體的痛楚,「渾身生瘡」這句話的希臘文是常用的醫學名詞,新約聖經只在此處出現過.「渾身生瘡(Covered with sores)」和「穿著紫色袍和細麻布衣服(dressed in purple and fine linen)」形成了一個很強烈的對比.

「被人放在財主門口」說明了他活在等候食物中,「要得財主桌子上掉下來的零碎充饑」,據說古時猶太人吃飯時不用刀叉,也不用圍巾,人們用雙手取食物吃；而在富有人家,常以麵包來擦手,擦過手的麵包就隨手丟棄,拉撒路和狗在門口所等待的,便是這類麵包.這裡就顯示了財主天天奢華宴樂與拉撒路天天等食餅碎,食不飽的情況,有著強烈的對比.

「並且狗來舔他的瘡」,這句話有兩個意思,第一,猶太人視他們所輕看的人如狗,故指拉撒路的處境如狗,其痛苦得不著人的同情；第二,瘡會流血,所以狗舔瘡就是舔血,意味著他活活的被狗折磨而死.(參王上廿一19)再一次對比了財主和拉撒路的情況:財主有工人服侍,拉撒路卻有狗和他爭食,並舔他的瘡.

財主錯在那?就是他不作壞事,亦沒有做好事 .這個財主並不是一個窮兇惡極,欺壓窮人的人,否則,他絕不會容許一個乞丐在他的門口討食.他的問題是眼中只有美食,沒有窮人；他沒有善用錢財,辜負了神託付他那麼多錢財的用心,上帝給他的錢財只為滿足自己,令自己快樂.財主的罪是什麼?財主只顧自己,忽視別人(即拉撒路)的需要,感受和痛苦.

同樣地,現今的地產商可能只顧自己的利益,而不擇手段.恐怖主義者可能只顧自己的信念,而不斷進行恐怖襲擊.怪獸家長為達致把自己孩子培育成龍成鳳的期望,而忽略孩子的個人感受及漠視孩子自我的發展.

財主,他失去一種對人(即拉撒路)的認同感,對拉撒路的飢餓,疾病,被狗欺負的情況都全不理會.

心理學家弗洛姆(Eric Fromm)發現,認同感是人人普遍的原望,它能使個體的生命突破人我的隔離力量,得到他人的認同感,使人脫離孤立與隔離狀態,把他和他人結合起來.哲學家叔本華 (Arthur Schopenhauer)提出人與人的相處 「同情」很重要, 同情是 「同感」或 「共感」,能對他人的痛苦感同身受,進而達至 「共苦」,與他人分擔痛苦,這是更進一步.所以人不被認同是痛苦的.拉撒路不被認同,他長久活在痛苦裡.他眼看自己吃不飽,穿不暖,不能走路,無法翻身,沒有自尊的情況,卻無法改變,那種痛苦生不如死.

一個現代人落到如此地步,就會感到生不如死,他可能就會自殺和輕生.最近中大有班同學研究自殺的問題,我發現:一個像拉撒路如此,不是沒有東西吃而餓死,也可能有得吃但不食而死!財主眼中,沒有拉撒路,沒有看見飢餓,患病,無助的人,只顧自己天天奢華宴樂,大魚大肉,吃喝玩樂.

2. 地下(陰間)兩種情況

兩種不同的命, 卻有著相同的結局:死亡, 但他們進入陰間有著兩種不同的待遇.

「後來那討飯的死了 (是餓死或病死,自殺死?),被天使帶去放在亞伯拉罕的懷裡.財主也死了,並且埋葬了.」(路十六22)

拉撒路有天使的同在和帶領,更被放在亞伯拉罕的懷裡,就如有著上帝的看重和兼顧,「懷裡」一詞取材於猶太人赴筵坐席時的情景,即每個坐席的人都憑著自己的肘側身半躺著,有如處身在他後側之人的懷裡(參約十三23),所以「懷裡」含有「與那人同赴席,為人所愛的意思」.

後來財主也死了,可能是因是暴飲暴食而死,然後就被人埋葬了.財主生前非常顧惜他的身體,穿的是紫色袍和細麻布衣服,吃的是珍饈美味,為他的身體耗費錢財,但是死後也不過是把身體埋葬了,他更要下到陰間受苦.

「他在陰間受痛苦,舉目遠遠的望見亞伯拉罕,又望見拉撒路在他懷裡,」(路十六23)

「陰間」是人死後靈魂暫時停留,等候末日受神審判的地方(參徒二27,31).它分成兩部分,一部分收留義人的靈魂,又稱樂園(參路廿三43)；另一部分收留罪人的靈魂(參伯廿四19),中間有深淵相隔,可望而不可及(參26節).事情並未了結,財主他的靈魂,那有知覺的自我,下了陰間.陰間是舊約的用語是希臘文Sheol(音譯索奧),意思為「逝去之靈的居所」.

拉撒路死後,被天使帶領到天父懷裡,得安慰；而財主死後,則被埋葬,並到地獄去受苦.

3. 財主在陰間的表現

(財主)「就喊著說:『我祖亞伯拉罕哪,可憐我罷,打發拉撒路來,用指頭尖蘸點水,涼涼我的舌頭；因為我在這火焰裡,極其痛苦.』亞伯拉罕說:『兒阿,你該回想你生前享過福,拉撒路也受過苦(拉撒路受苦時,你無作善事,由得他受苦) ；如今他在這裡得安慰 (因為他受盡苦,上帝幫助他脫離苦了),你倒受痛苦 (因為你見苦不救,你要受懲罰).」(路十六24－25)

財主的苦可分成兩方面:

第一方面是身體的苦,即火焰的苦(參24節)；第二方面是「回想」的苦,包括他悔恨的苦和祈求不允的苦(參27~31節),其苦真是永世沒有休止的苦.今生的選擇決定死後永恆的結局,一旦死亡,結局就限定了,享福或受苦,再也無法彼此對調了(參26節).

拉撒路「得安慰」

「得安慰」 與「保惠師」(參約十四16)在原文同屬一個字群,信徒生前有聖靈在旁安慰與引導,死後仍得被呼召靠近所信的主,即與主同在(參腓一23).

4. 有深淵的限定

「不但這樣,並且在你我之間,有深淵限定,以致人要從這邊過到你們那邊,是不能的；要從那邊過到我們這邊,也是不能的.”」(路十六26)

4.1. 有深淵限定:空間

陰間與樂園之間有深淵相隔,致財主求一滴冷水以涼舌頭也不可得.這時,他想起父家中尚有五位兄弟未得救,於是他期盼拉撒路能被差遣去傳福音給他們,可惜已經太遲了.

所以你想得人幫助也要把握時間和機會,就算亞伯拉罕,拉撒路想幫也因有深淵而不行!人應當把握機會做善事,好事,有益於人的事, 死後就沒有機會了.就如財主得不到水喝,得不到人(拉撒路)的幫助,得不到超越界限,這是不能改變的規限!

人生有許多規限,所以你要認識你的規限.小時你不好好讀語文,長大後才學就會有限制,有機會讀書時你不讀,過了就再沒有機會.趁孩子還小就多陪他；待他長大後你陪他,他可能會走開.明白與接受聖經所言:「生有時,死有時......在崇基神拜堂的日子有時,凡事都有定時」的人,才是智慧人.

「27財主說:『我祖阿,既是這樣,求你打發拉撒路到我父家去；28因為我還有五個弟兄；他可以對他們作見證,免得他們也來到這痛苦的地方.』29亞伯拉罕說:『他們有摩西和先知的話,可以聽從.』30他說:『我祖亞伯拉罕哪,不是的；若有一個從死裡復活的,到他們那裡去的,他們必要悔改.』31亞伯拉罕說:『若不聽從摩西和先知的話,就是有一個從死裡復活的,他們也是不聽勸.』」(路十六27－31)

財主要求亞伯拉罕差派拉撒路沾水給他,這是為自己的表現；而他要求亞伯拉罕差派拉撒路去為五兄弟作見證,不要到地獄 ,這則是為家人的表現,可惜機過已過,財主已後悔莫及.

4.2. 有生死的規限:時間

有限的時間即死亡和死線(Dealine) 極具意義.我們重視也應重視死亡及死線.財主無死亡意識,不懂好好管理自己的時間,每天奢華宴樂,浪費時間；拉撒路在困苦中,連生命時間也沒有的,只能努力存活,且隨時會死亡,不能掌握時間.

海德格(Martin Heidegger) 提出人生的「大限」,即死亡,深具意義,因為只有人意識到 「死亡」才能認真地活著,並且要有使命的活著,好好活此生.

路加福音在這故事中,更提出了「死亡」在上主的手裡,可以扭轉乾坤,就如在陽間的時候,財主很快樂,但拉撒路很痛苦；在陰間的時候,財主很痛苦,拉撒路卻很得安慰,上帝在他世幫助拉撒路脫貧脫苦.

時間的長短,人不可確定,但在有限時間的運用,人可以作主.過去所做的也會影響現在,因財主在世時不顧拉撒路,結果死後受罰.同樣地,現在所做的將會影響未來,拉撒路現在受苦,但他仍然忍受痛苦直到死,結果後來得到天使引路到樂園睡在主懷中,這顯明了上月的講道主題:「人在做,天在看!」人要知道,過去所做的,影響現在,現在所做的影響將來,這是人生大道理；人也要把握現在的機會,過去了就機會不再.

海西阿德說過:「善於掌握自己時間的人,是真正偉大的人. 」為好好管理時間,不同的學者有以下三大貢獻:

第一,把時間分為年,月,日,時,分,秒；

第二,每天定為二十四小時,讓人可以把時間分成三份:八小時工作,八小時睡眠及八小時自己安排的時間,如吃飯,如廁,交通,玩樂,學習,與家人一起,交際等；

第三,「朝九晚五」成為今天全世界「公共辦公時間」,讓人盡量在此時間內完成工作,建立精確的時間運用.

為了探索時間與人生,學者們又有幾個生命時間的觀念:

第一,時間長度以及長壽的概念.

第二,時間闊度,讓人多方面的認識與掌握,例如一個讀書人,他擁有書,也會看書和寫書,又例如我愛讀神學,哲學,歷史文化,聲學,心理學等,這些都是闊度的例子.

第三,時間密度,今天大部份人追求時間的密度,或者是被迫追求,例如是寬頻上網,速度當然愈來愈快,但我們的忍耐力卻愈來愈差；另一個例子是即食文化普及,人們震吃得快和好味；還有人們追求商業成本效益,要賺快錢,在短時間得大利益；更有死線的催逼,使人精神緊張,容易暴燥,工作或讀書為首位,使他們目中只有自己,沒有別人.現代人常追求時間的密度和速度,人都巔(整天要達到巔峰) ,使人失去長度,闊度和深度.

第四,時間深度, 一種 「內在時間 (immanent Time)」,即有意義的時間,例如男女一起吃燭光晚餐,家人一起的歡樂時光,主日崇拜,與親人臨終的告別等.

在本段經文引發思想同情別人,與人有認同感及樂於為善的時間.你還有時間,你怎樣善用時間呢?

昨天我往老人院探訪媽媽,聽到一位老婆婆在彈鋼琴,她已九十多歲仍彈得一手悅耳的聖詩,我也忍不住口跟著唱,聲音響遍老人院,不過大部份老人對我的歌聲都無反態應,因為他們都是耳聾的,包括我的媽媽,只有兩位婆婆豎起手指公,我就知道他們沒有耳聾.我也買了蛋卷給媽媽,她吃得很開心,與九十多歲的婆婆 (最年長Pianist)合作,她彈我唱,兩位也近百歲的婆婆無耳聾,聽我唱詩聽得很開心,成為我的聽眾,我實在也作了美事,善事.作善事,不一定要驚天動地的偉大事業,日行一善,多作好事,不要像財主一樣只顧自己,不顧門外的拉撒路.

總結

朱自清這樣寫過:「洗手的時候,日子從水盆裡過去；吃飯的時候,日子從飯碗裡過去；默默時,便從凝然的雙眼過去.我覺察他去得匆匆了,伸出手遮挽時,他又從遮挽著的手邊過去.天黑時,我躺在床上,他便伶伶俐俐地從我身邊跨過,從我腳邊飛去了.等我睜開眼和太陽再見,這算又溜走了一日,我掩著面歎息.但是新來的日子的影兒又開始在歎息裡閃過了. 」

有智者這樣說過:「告訴我你怎樣運用你的時間,我就告訴你,你是個怎樣的人!」各位弟兄姊妹,到底你又是一個怎樣的人呢?希望大家從今起會不做壞事,常常做好事,善用你的生命.



\newpage

\section{天行健君子以自強不息}
\label{sec:337}
\textbf{鄧瑞強博士 - 天行健,君子以自強不息}
\newline
\newline
\hyperref[sec:336]{< < < PREV SERMON < < <}
~
\hyperref[sec:toc]{[返目錄]}
~
\hyperref[sec:338]{> > > NEXT SERMON > > >}


提摩太後書 1:1-14
\newline
\begin{longtable}{cl}
\hline
\hline
章節 & 經文 (和合本修訂版)\\
\hline
1:1 & \begin{tabularx}{0.7\textwidth}{X} 奉神旨意，按照基督耶穌裡所應許的生命，作基督耶穌使徒的保羅， \end{tabularx} \\ \\ \relax
1:2 & \begin{tabularx}{0.7\textwidth}{X} 寫信給我親愛的兒子提摩太。願恩惠、憐憫、平安從父神和我們的主基督耶穌歸給你！ \end{tabularx} \\ \\ \relax
1:3 & \begin{tabularx}{0.7\textwidth}{X} 我感謝神，就是我接續祖先用純潔的良心所事奉的神，在祈禱中晝夜不停地想念你。 \end{tabularx} \\ \\ \relax
1:4 & \begin{tabularx}{0.7\textwidth}{X} 我一想起你的眼淚，就急切想見你，好讓我滿心快樂。 \end{tabularx} \\ \\ \relax
1:5 & \begin{tabularx}{0.7\textwidth}{X} 我記得你無偽的信心，這信心先存在你外祖母羅以和你母親友妮基的心裡，我深信也存在你的心裡。 \end{tabularx} \\ \\ \relax
1:6 & \begin{tabularx}{0.7\textwidth}{X} 為這緣故，我提醒你要把神藉著我按手所給你的恩賜再如火挑旺起來。 \end{tabularx} \\ \\ \relax
1:7 & \begin{tabularx}{0.7\textwidth}{X} 因為神賜給我們的不是膽怯的心，而是剛強、仁愛、自制的心。 \end{tabularx} \\ \\ \relax
1:8 & \begin{tabularx}{0.7\textwidth}{X} 所以，不要以給我們的主作見證為恥，也不要以我這為主被囚的為恥；總要靠著神的大能，與我為福音同受苦難。 \end{tabularx} \\ \\ \relax
1:9 & \begin{tabularx}{0.7\textwidth}{X} 神救了我們，以聖召召我們，不是按我們的行為，而是按他的旨意和恩典；這恩典是萬古之先在基督耶穌裡賜給我們的， \end{tabularx} \\ \\ \relax
1:10 & \begin{tabularx}{0.7\textwidth}{X} 但如今藉著我們的救主基督耶穌的顯現已經表明出來；他把死廢去，藉著福音，將不朽的生命彰顯出來。 \end{tabularx} \\ \\ \relax
1:11 & \begin{tabularx}{0.7\textwidth}{X} 我為這福音奉派作傳道，作使徒，作教師。 \end{tabularx} \\ \\ \relax
1:12 & \begin{tabularx}{0.7\textwidth}{X} 為這緣故，我也受這些苦難。然而，我不以為恥，因為我知道我所信的是誰，也深信他能保全他所交託我的，直到那日。 \end{tabularx} \\ \\ \relax
1:13 & \begin{tabularx}{0.7\textwidth}{X} 你從我聽到那健全的言論，要用在基督耶穌裡的信心和愛心常常守著，作為規範。 \end{tabularx} \\ \\ \relax
1:14 & \begin{tabularx}{0.7\textwidth}{X} 你要靠著那住在我們裡面的聖靈，牢牢守住所交託給你那美好的事。 \end{tabularx} \\ \\
[1ex]
\hline
\hline
\end{longtable}
講題:天行健,君子以自強不息　Be Persevering on the Way of the Cross

講員:鄧瑞強博士

經文:提摩太後書1:1-14

場合:香港中文大學崇基學院禮拜堂

日期:2016年10月2日

各位弟兄姊妹,早安.

《易經》的第一卦,是「乾」卦.《象傳》對「乾」卦作了一注解,說:「天行健,君子以自強不息」.對於這句話,最簡單的理解是:「天」即大自然的天,天行健,即大自然運轉不停,人在此體會到自己應努力不懈,自強不息.或者,我們可以將「天」理解成萬物背後冥冥中的神聖力量,這神聖力量不停地,默默地以其生命力托住蒼生,作為人類的我們,感悟到這份神聖力量的奧妙,從而甘心的透過自身的努力實踐,將這神聖生命力展露出來,去滋潤萬物.在這實踐過程中,君子自覺自己正是實現上天的心意.君子的自強不息,正是活現上天在人間的剛健運行.「天行健,君子以自強不息」,人堅毅努力地實現美善,正是因為上天正是這樣堅毅努力地實現美善.

我們相信,萬物背後,有一位承托著萬物的天父,祂的心意表明在人間的耶穌身上.有一次耶穌說:「我父做事直到如今,我也作事」.這句話可譯為:「天行健,我也如此」.基督徒是誰,是明白耶穌的人,也是透過耶穌去明白天父的人.基督徒是一群甘心讓自己的生命向基督轉化的人,向基督轉化,就是活似基督.活似基督,就是讓天父的心在自己的生命中展現.甘心如此,是因為體會天父那不止息的仁愛,受這仁愛感召,而甘願參與在仁愛的行動中.

保羅,對耶穌的生命,有透切的理解.因這理解,他參與了基督的生命,像基督一樣生活.到了某一刻,他感到自己的生命,已走到盡頭.他渴望將自己的生命領悟,傳給自己的入室弟子.他寫了《提摩太後書》,給他的徒弟提摩太.若我有徒弟,我都很想像保羅一樣,講這番說話.

今日講道的經文,是今日的書信經課:提摩太後書1:1－14

《提摩太後書》是一位老牧者向一位年輕牧者講的最後一番話,交代一點後事,作一點叮嚀.老牧者要將信仰使命的棒交給下一代,他向他的徒弟講解這棒背後的意義.老牧者如此著緊那使命,是因為那使命承載著天父的心.這天父的心,是萬物的意義所在.

若你知道一救世秘密,你會懷著什麼心情,將這秘密傳下去呢?

在《倚天屠龍記》這本小說裡,倚天劍和屠龍刀背後,正正隱藏著一個救黎民於水火的秘密.小說的細節不說了.只道,這秘密最後傳到滅絕師太那裡.話說當時,滅絕師太被困於萬安寺,心想必死無疑,便將倚天劍和屠龍刀的秘密傳給周芷若.並為了保住周芷若的性命,滅絕師太抱著周芷若從高塔跳下,快落地時,運勁將周芷若托起,以保她的性命,也就能保存那個救世秘密.

若倚天劍和屠龍刀的秘密,傳了給你,而你又行將就木,你會如何?

我們看看保羅如何將生命的奧秘傳給他的徒弟提摩太.

提後1:1　奉神旨意,照著在基督耶穌裡生命的應許,作基督耶穌使徒的保羅,

提後1:2　寫信給我親愛的兒子提摩太.願恩惠,憐憫,平安從父神和我們主基督耶穌歸與你!

保羅,身分是「基督耶穌的使徒」,這身分出自「神的旨意」(天父的天意).這身分的使命是什麼?是將天意的秘密說清楚,這秘密就是天父透過耶穌向世人應許了永生.人憑自己是不可能知道這秘密的,能知道這秘密,只能出自天父的旨意.

不明白這秘密的人,會按自己的理解,去構想這秘密.正如江湖上,很多人去搶屠龍刀,死傷無數,但其實他們根本不曉得箇中秘密.福音的秘密,也是這樣.很多人相信福音,但可能不曉得當中的秘密.要明白福音的奧秘,必須嚴格依從福音的傳承,由知道真相的人,一代一代的傳下來.

提後1:3　我感謝神,就是我接續祖先用清潔的良心所事奉的神.……

提後1:5　想到你心裡無偽之信,這信是先在你外祖母羅以和你母親友妮基心裡的,我深信也在你的心裡.

保羅提到,他是「接續祖先」的信仰的,並以「清潔的良心」去事奉神的.他的信仰是沿著真正的傳統的,他的人生是純潔地按這傳統而活的.他用心地保存著信仰的秘密.

提摩太,是一個混血兒,他父親是希臘人,母親是猶太人.明顯,提摩太的母親及其外祖母在他的信仰教育中發揮了主導角色,提摩太從她們那裡承繼了先祖的信仰.保羅鼓勵他,對這傳統信仰保持「無偽之信」.

傳統是信仰的根.根,吸收養分,穩定樹幹.根,看不見,但對樹木而言,重要得很.一棵樹表面看可能很高大,但若底下的根不夠深,不夠健康,則長得越高,倒下的機會便越大.

聽一些基督徒說,他們教會裡的男女關係較不清不楚.這可能是現代潮流,但肯定不合信仰的傳統.傳統可能很老套,但卻是保護我們的籬笆.信仰要有根,必須要有傳統.問題只是,這傳統如何傳承下去.

若你手握著倚天劍和屠龍刀,你驚不驚?若你知道救世的秘密,你驚不驚?

我想,驚是常情.提摩太面對保羅的交托,一定很驚.

提後1:6　為此我提醒你,使你將神藉我按手所給你的恩賜再如火再挑旺起來.

提後1:7　因為神賜給我們,不是膽怯的心(原文是「靈」字),乃是剛強,仁愛,謹守的心.

保羅鼓勵接棒者,不要膽怯,不要驚.這裡提到保羅按手在提摩太身上.「按手」,是師父將能力和責任傳給徒弟,像滅絕師太將掌門指環戴在周芷若的手上.接過指環,你就是掌門了.保羅的身分,責任與事奉能力,來自天父.藉著按手,提摩太承接了這從天父而來的身分,責任,能力.

我們只是軟弱的人,根本承受不起天命責成的重任.保羅鼓勵提摩太,成就大任的,不是我們的「洪荒之力」,而是天父賜的聖靈之力.

神賜給我們的聖靈,不是膽怯的靈,乃是剛強,仁愛,謹守的靈.「剛強」者,完成使命的「勇力」；「仁愛」者,為了他人的好處而甘心放下自己；「謹守」者,明辨是非的「智慧」.「剛強,仁愛,謹守」可說就是中國傳統講的「智,仁,勇」,這不是人可修煉而有的特質,而是天父的靈的特質.天父的靈將成就天父使命所需的「智,仁,勇」賜給願意承擔的人.

有齣電影,叫《Amazing Grace》(奇異恩典),故事主角叫威伯福斯(William Wilberforce),一個有錢的年輕人,英國下議院的議員,名成利就.但他就是對生命懷有一悲憫之情.後來,因著不同的人的鼓勵,在1789年提出禁止奴隸買賣的議案.當時,奴隸制度,牽涉很多人的財富,甚至牽涉國家的強大.這議案當然被人否決.很多人反對威伯福斯,但威伯福斯鍥而不捨,經過將近20年的努力,這議案在1807年終於獲得通過.然後,再經過26年的努力,才能成功爭取通過奴隸解放的法案.這法案通過後三天,威伯福斯便與世長辭了.

威伯福斯為何花一生去解放奴隸,當然,是驚訝於人類竟能如此「非人性」地對待他人,是對他人有深切的憐憫.有次,他的僕人問他:「是你找到上帝嗎?」威伯福斯的回答是:「不,是上帝找到了我.」

真的,天父找到了他,是「天命」臨到他,並將「智,仁,勇」的聖靈賜給他,讓他完成非常人能完成的事情.

最後,我要講的,是保羅將一個什麼樣的人生領悟,一個什麼樣的信仰使命,交給了提摩太.

提後1:8　你不要以給我們的主作見證為恥,也不要以我這為主被囚的為恥；總要按神的能力,與我為福音同受苦難.

提後1:9　神救了我們,以聖召召我們,不是按我們的行為,乃是按他的旨意和恩典；這恩典是萬古之先,在基督耶穌裡賜給我們的,

提後1:10　但如今藉著我們救主基督耶穌的顯現才表明出來了.他已經把死廢去,藉著福音,將不能壞的生命彰顯出來.

提後1:11　我為這福音奉派作傳道的,作使徒,作師傅.

提後1:12　為這緣故,我也受這些苦難.然而我不以為恥；因為知道我所信的是誰,也深信他能保全我所交付他的〔或譯:他所交託我的(按:我選擇這後者的翻譯)〕,直到那日.

提後1:13　你從我聽的那純正話語的規模,要用在基督耶穌裡的信心和愛心,常常守著.

提後1:14　從前所交託你的善道,你要靠著那住在我們裡面的聖靈牢牢的守著.

先看第10節,這就是保羅所持守的,要傳遞的終極奧秘.

提後1:10　但如今藉著我們救主基督耶穌的顯現才表明出來了.他已經把死廢去,藉著福音,將不能壞的生命彰顯出來.

天父的旨意,顯明在耶穌的生命裡.耶穌以「死亡」廢去「死亡」,以「苦難」取消「苦難」,以「絕望」拒絕「絕望」,這就是福音,這就是生命中最大的奧秘.在「看來沒有上帝之處」顯示「上帝的真實臨在」,在「被人棄絕的死亡」裡顯出「永不朽壞的光明」,在「為了解放奴隸而傾盡一生」裡顯明這不是「浪費生命」而是「填滿生命」.這是至仁,至勇,至智的生命,這是聖靈下的「智,仁,勇」的生命.

提後1:8　你不要以給我們的主作見證為恥,也不要以我這為主被囚的為恥；總要按神的能力,與我為福音同受苦難.

保羅,就是為這奧秘而被囚,受苦,受辱.他呼籲他的徒弟,也要藉聖靈之力,剛健勇毅地持守這信仰奧秘,也要為這信仰而承受人生的艱難.

當時的人,重視「榮譽」,遠離「羞辱」,這跟中國人有點相似.但什麼是榮譽,什麼是羞辱,一定要搞清楚.羞辱者,莫過於耶穌釘十架.榮譽者,亦莫如耶穌的復活.世人厭棄的十架,其實是神的榮譽的前奏.保羅看透耶穌所啟示的真理,他看清天父的心,故此,他不以當下的苦為苦,不以當下的恥為恥.

提後1:9　神救了我們,以聖召召我們,不是按我們的行為,乃是按他的旨意和恩典；這恩典是萬古之先,在基督耶穌裡賜給我們的,

保羅明白,在「萬古之先」,在永恆裡,有一托住萬物的天心,即天父的心,這心在基督的十架與復活裡顯明.保羅被這天父的心所感動,所感召.進入這心,即與基督認同,以受苦的愛去顯明何為復活的光明.

提後1:11　我為這福音奉派作傳道的,作使徒,作師傅.

提後1:12　為這緣故,我也受這些苦難.然而我不以為恥；因為知道我所信的是誰,也深信他能保全他所交託我的,直到那日.

保羅領受使命,為此承受恥辱,他相信,一切源自天父.天父交託的,只要我們努力承擔,天父總會完成的.「君子自強不息」,是因為他深深體會「天行健」.

面對莫大的艱難,面對無告的羞辱,面對不被理解的使命,提摩太的壓力,可想而知.這也是每個基督徒的壓力,壓力之大,可想而知.

最後,請聽保羅安靜的鼓勵與安慰:

提後1:13　你從我聽的那純正話語的規模,要用在基督耶穌裡的信心和愛心,常常守著.

提後1:14　從前所交託你的善道,你要靠著那住在我們裡面的聖靈牢牢的守著.

但願:榮耀歸於聖父,聖子,聖靈.阿們.

(後話:本人將早期的崇基講章結集為書,書名《敢於跟隨主》,基道出版,有興趣者可購買.)



\newpage

\section{後現代性與基督教信仰}
\label{sec:338}
\textbf{謝任生牧師 - 後現代性與基督教信仰}
\newline
\newline
\hyperref[sec:337]{< < < PREV SERMON < < <}
~
\hyperref[sec:toc]{[返目錄]}
~
\hyperref[sec:339]{> > > NEXT SERMON > > >}


講題:後現代性與基督教信仰　Christian Faith and Post- Modernity

講員:謝任生牧師

場合:香港中文大學崇基學院禮拜堂

日期:2016年10月9日

深哉,上帝的豐富,智慧和知識!他的判斷何其難測!他的蹤跡何其難尋!誰知道主的心?誰作過他的謀士?誰先給了他,使他後來償還呢?因為萬有都是本於他,倚靠他,歸於他.願榮耀歸給他,直到永遠.阿們!

所以,弟兄們,我以上帝的慈悲勸你們,將身體獻上當作活祭,是聖潔的,是上帝所喜悅的,你們如此事奉乃是理所當然的.不要效法這個世界,只要心意更新而變化,叫你們察驗何為上帝的善良,純全,可喜悅的旨意.

引  言

今天早上我講道的題目是後現代性與基督信仰.首先,我想稍為解釋甚麼是後現代性(Post Modernity).講後現代就是有現代和前現代.我不打算詳細講解現代和前現代,只是簡述一下,而集中在講後現代性.其實,我們也正是在後現代性時期.那麼甚麼是後現代性?從內容上看,後現代性大致體現為三種比較鮮明的現象 :

1/ 否定性(解構性)的後現代性    所謂否定性(解構性)的後現代性,就是要對現代性的基礎,即理性主義進行徹底顛覆,特別是否定現代性的所有二元結構(諸如 : 主/客,真/假,本質/現象,事實/價值,等等).即是說,後現代性的看法是,根本就不存在甚麼「絕對的真理」,「普遍的規律」,「永恆的結構」等.這種現代性的結構無論看上去多麼合理,多麼複雜,實際上都是虛假的和不正確,因為一切都是相對的. 因此,我們不是經由現代性的絕對主義和個人主義來建構一個完滿的人文主義的社會.相反,只有通過瓦解這種捆綁人的現代性,才能實現人的自由和真正解放.還有,後現代主義強調消費主義.如此,人才是自己的呎度.這就是後現代的袪絕對主義,否定理性上所主張的真理.

2/ 建設性(建構性)的後現代性    所謂建設性(建構性)的後現代性,就是試圖在批判和反思現代性的同時,重建一種超越現代性的價值觀,以「無根基」的後現代性完全取代現代性. 也就是說,要擺脫現代性普遍模式,通過建立人與自然,人與社會,人與人的新型關係,才能真正達到目的.這就是後現代性的個人主義,也是環保運動,綠色革命的動力.

3/   虛假性(迪士尼式)的後現代性

所謂虛假的後現代性就是建構一種迪士尼樂園式的社會,就是把現實生活放入迪士尼樂園式的虛幻中,再成為好像是真實的現實世界.人在迪士尼樂園中的世界是充滿快樂的,但又不是真實生活上的快樂.迪士尼樂園主要是要針對人如何摧毀現代性中種種理性對人的控制與壓抑,以祈尋求思想和生活上自由的最大化,從而最大限度地釋放出被現代性所壓抑的創造性.這也就是迪士尼樂園成功的地方,叫人可以暫時享受在虛擬中的自由與快樂,可是,這也是令人最空虛的地方,因為這個迪士尼樂園的世界,我們樂在其中時是快樂的,但又不是恆久的,既是真實的但又是虛幻的.

一/ 後現代性與羅馬書十一章33節至十二章2節的關係

保羅為在羅馬書中也曾提及現代性社會及後現代性社會的問顯.羅馬書十一章33節至十二章2節具有一種橫向意義(Lateral Meaning).第一章至第十一章是保羅以柏拉圖哲學的理性思維去闡釋因信耶穌基督而稱義的教義,這是一種現代性陳述上帝的因信稱義,而第十二章3節開始講及信徒生活的倫理問題,是一種在因信稱義的基礎上去建構一種基督宗教的人文主義,教導信徒要當怎樣行事為人,叫人可以知道如何察驗何為上帝的善良,純全可喜悅的旨意.這可以說是後現代的建設性(建構性)的人文主義,是保羅把當時羅馬帝國太平盛世的腐化道德重新的建構成為基督宗教的神聖倫道理道德的社會觀.然而,在這兩大段落中,保羅加入了一段詩一般的插曲:

深哉,上帝的豐富,智慧和知識!他的判斷何其難測!他的蹤跡何其難尋!誰知道主的心?誰作過他的謀士?誰先給了他,使他後來償還呢?因為萬有都是本於他,倚靠他,歸於他.願榮耀歸給他,直到永遠.阿們!

所以,弟兄們,我以上帝的慈悲勸你們,將身體獻上當作活祭,是聖潔的,是上帝所喜悅的,你們如此事奉乃是理所當然的.不要效法這個世界,只要心意更新而變化,叫你們察驗何為上帝的善良,純全,可喜悅的旨意.

這段經文的意義是甚麼?就是保羅一方面要以希腊哲學式來闡明因信稱義的合理性,而另一方面保羅也知道理性不是解決人際間關係的問題,而必須要從一種超越當時社會的人倫墮落的新方向,就是在耶穌基督裡的新人際關係來建立一種新的文人社會.就如著日基督教的出現,要去取代著日羅馬帝國日落西山的腐敗社會一樣.而其中的重點不在理性主義,也不在人文主義,而是以上帝的奧秘和大愛來建立一個新的國度,就是基督的國度.

那麼這段經文對今日的後華人教會時期有何關係?這正是我想和大家分享的主旨.巴克萊(William Barclay)對這段經文解釋說 :

保羅從未寫過比這一段更具獨特性的經文.在這裡我們看到由追求真理的頭腦轉而為敬拜讚美的心靈.最後,一切都必過去進入奧秘之中.這是說,在人不能純以理性去了解上帝時,只要心中卻是愛著上帝.這樣的人又能夠說萬有都是本於祂,倚靠祂,歸於祂,那麼還有甚麼沒有說的呢?在人的情況中,有某種的兩極性.上帝給人頭腦,人的責任就是盡量的用它去思想.不過有的時候,頭腦也有到了極限的時候.在那時,人只有接受上帝的恩典和讚美上帝.

巴克萊在以上的解釋說明了信仰的兩條腳:人理性的信仰和人渴愛上帝的奧秘.人可以很理性的相信上帝,但卻不一定需要去愛慕上帝.這是說,現代性所主張的理性主義與人心靈中原有的感性的愛是有分別的.現代性要求人以理性為依歸來解釋上帝和聖經.可是,愛是很難以理性來體現的,因為人在愛慕上帝之中,必須是以感性來建立與上帝的關係.人對上帝有信心不一定就會必然去愛上帝,但是人若是真心愛上帝,就必然會去相信上帝.所以保羅在十二章1-2節說,所以,弟兄們,我以上帝的慈悲勸你們,將身體獻上當作活祭,是聖潔的,是上帝所喜悅的,你們如此事奉乃是理所當然的.不要效法這個世界,只要心意更新而變化,叫你們察驗何為上帝的善良,純全,可喜悅的旨意.

這段插曲可以說是全卷羅馬書的軸心思想,就是:上帝的奧秘和祂的大愛,並且把前後兩大部份的羅馬書,透過上帝的奧秘和大愛聯繫起來.只是,我們不竟是活在一個後現代的時代中,既有前現代主義要求人對權威和傳統的尊重和服從,但又會有執著對現代性的理性主義,凡事講求理性和科學,神學更是如此.另一方面,我們又不能逃避後現代主義所要求的講個人自由,多元主義,強調消費就是富有和成功的表記.只是,我們必須明白,在現代性的理性主義是有一定的限度,因為有些事情是超越理性所能解答的.例如,理性不能給上帝一個完滿的稱號.就是「上帝」一詞也不一定就是上帝的真實名稱,而只是一個普通名詞(common noun).因為上帝是可以有稱謂但又可以沒有稱謂的,祂原是超越兩者的.所以,理性解答不到的東西就不能說是錯的.相反,我們應當說,若是理性不能解答的問題,我們就當承認理性是有一定的限度.我們就要從另一角度來尋找答案.這就是後現代性的要求:要「終結」現代性所有的傳統.但是,對基督宗教而言,這就是在人在上帝的愛中先學習相信上帝和追尋祂的奧秘.這不是靠理性的求知,而是靠本性上對上帝的愛慕.

神學家巴特(Karl Barth)在他的巨著《羅馬書註釋》論及這兩段經文時解釋說,這位以因信稱義,使人因信耶穌基督裡能認識這位「非直觀性」的上帝.他原是一位「隱而不露」,「深下可測」的上帝.這正是基督宗教在後現代性對上帝應有的認識,去體驗驗上帝愛和奧秘.

但是,奧秘之為奧秘,就不是人能以人的頭腦和理性所能把握的.巴特更認為:人要「直接認識上帝?不可能!參與上帝的決定?不可能 !控制他,約束他,使他承擔義務,去和他建立一種相互關係?不可能!沒有甚麼「聯邦神學」!他是上帝,他自身是上帝,唯獨他是上帝.這就是羅馬書中肯定的『是!』」  那是說,神學與理性只能作非常有限度的解釋上帝,有時更是自作主張的解釋上帝.所以巴特認為 : 人若「要認識上帝,就得靜靜地面對他祈禱,而他(是)在人難以企及的光芒籠罩之中.」   這就是保羅在十二章開始所說,所以,弟兄們,我以上帝的慈悲勸你們,將身體獻上當作活祭,是聖潔的,是上帝所喜悅的,你們如此事奉乃是理所當然的.不要效法這個世界,只要心意更新而變化,叫你們察驗何為上帝的善良,純全,可喜悅的旨意.

巴特同時更是指出,信仰上帝與生活上帝是不能分割的.那是說,我們對上帝的認識必須基建於我們真實的生活上,特別是文化上,而非單靠理性思維去分析.他說 : 「為了理解《羅馬書》,建議大家趕緊去讀各種純屬世俗的文獻,尤其是報紙.因為,思想如果是真實的思想,就是生活的思想,因而也就是上帝的思想(意指從上帝而來思想).」  所以,人可以用理性來認識上帝,但他同時更需要對上帝存有體認與愛慕,並且達到以基督的心為心的心性.《使徒信經》是大公教會所認同的信經,但是,我們不但要確信上帝,更應當是我愛上帝,又是我們愛上帝.以愛來認信上帝才是我們對信仰的起點,並且可以統懾人的理性,感性,心性,從而進入上帝的神性中,正如保羅所說的:「現在活著的不再是我,乃是基督在我裡面活著.」

可是,我們在一般的主日拜中,許多時是以理性來崇拜,而非以心性來敬畏上帝.余時英先生正看出這一點.他說人應當以心學去代替神學.這正是一語道破今日基督徒靈性的弱點,太過理性化了上帝,而缺乏心性上去經驗上帝.

但是,在某些教會的崇拜中,我們又會以非理性的情緒化崇拜,特別是在詩歌頌唱的環節上.我們忽略了虔敬下拜的敬畏上帝.我曾經問過一位東正教的姊妹怎樣認讖東正教的神學.她直接了當的對我說,首先你不用去研讀甚麼東正教的神學書藉,而是來參加我們東正教的主日崇拜.因為東正教的神學是注重眼睛的凝視,和靈性思念上帝的心學,這不是人頭腦的理學.只要你以敬虔之心去默觀東正教教堂內一切的陳設,就會感受到對上帝的認識和與你同在.這也許是我們所當反省的地方.

我要講的要點是,今日華人教會仍是死守著著日前理性主義的五十年不變的教條或及不明文的教規,而忽略了我們在信仰上是要不要效法這個後現代的世界,只要心意更新而變化,為後現代建立一個新的社群,就是上帝的教會,叫你們察驗何為上帝的善良,純全,可喜悅的旨意,好在這個好像沒有方向的後現代,指出上帝完美的國度.

(二) 目前香港華人教會的處境 : 對後理性時代的後知後覺

另一方面,我們在教會中有著太多的教條主義,宗派主義,神學學派主義和個人英雄主義,把我們在的耶穌基督裡所釋放的自由,變成被人間所謂屬靈主義所束縛,結果是 : 我們信得很無奈,但又不敢放棄.雞肋信仰!信之無力,棄之可惜!可是,信仰不竟不是教條.教條是不能改動的,但信仰卻是活生生的,是心意更新而變化的.Karen Armstrong 在《神的歷史》一書中一開始就說過類似的話.她說:「信仰不是教條,教條是不可以改變的,但信仰可是也是必然會的,因為信仰是隨著的我們生命成長,使我們對信仰,對事物,對人物,對觀念等而有所改變或改正的.這是生命的成長.」  所以,保羅說,信仰是會更新而變化的,察驗何為上帝的善良,純全,可喜悅的旨意.這是為何我們要在後現代文化中建構一種真正自由去進入上帝的奧秘和愛慕的關係.

可是照目前所見,一般教會多仍然活在前現代的傳統與權柄之下.缺乏「心意更新而變化,察驗何為上帝的善良,純全,可喜悅的旨意.」於是我們常常把上帝雪藏了,不許上帝在我們的生命中,在祂的教會裡作出任何有生命力的改變.林語堂是牧師之子,早年看透了教會的後知後覺和自我疆化.他曾經一度背棄基督教信仰而尋佛問道,最後晚年才回歸基督教中.他說 :

基督教甚至是把信仰放進雪格裡把之疆化,成為雪條又成為教條,不得有任何的變化.世上一切宗教差不多都是如此,而特別是基督教,它已經僵硬,放入箱裡,且加上防腐劑,它不容許任何討論.很奇怪,在這件關於宗教的事情上,每個人都似乎認為他所擁有的是獨一無二的真理.

另一方面,教會又常常自我主張的改變,並且變得令人十分感慨.林語堂又感慨說,上帝所能做的事情多過人的宗教所夢想的,但是在宗教中也有許多歸之於上帝的事情,上帝連做夢也沒有想到要去做.   這會不會只是地上教會自我主張,而非經由上帝的慈悲的改變?我在此想說的是 : 我們必須認識清楚 : 教會不等於上帝,聖經也不等於上帝,而早期的《使徒信經》及第一版的《尼西阿信經》更是從無提及過聖經的.那是說,我們信仰的核心不是聖經,而是在耶穌基督裡的三位一體的上帝.只是這些話不是一般教牧敢說出來的,因為一旦說出就會被標籛為不信派,自由派,甚至連朋友也沒有.事實上,我們就是活在所謂的屬靈的聖經主義中,忘記了聖經是為要見證上帝,教會也是為了服侍上帝.東正教主教Kallistos Ware說過以下一段引人反思的話 :

閱讀聖經的目的是要培育(cultivate)人靈性中的智慧,而非單單去增加人對聖經的知識.不錯,聖經知識能引導我們對聖經有新的發現,但是智慧卻是要引導我們看見上帝的榮美.所以,「基督宗教的任務不是為每個問題提供答案,而是要促使我們向上地意識到一種神秘感(即是奧秘).上帝並非是我們知識中的對象,好像是我們對神秘事物的誘因.」

不幸的是,一般信徒多會把主日歸給教會,把周一至周六歸給自己.信仰必須是建立在愛慕上帝上,而非建立在經文主義的教條中.可惜的是,今天我們的信仰又是怎樣的?是否

1. 信仰沒有靈性的深度,

2. 信仰也沒有文化的深度,

3. 一般的講道也沒有深度,

4. 崇拜的禮儀又是沒有深度.

我們今日大多數的教會崇拜,多是演唱會式的崇拜,有的只是清談節目或楝篤笑的歡樂今朝.一位某間大教會的執事感嘆地向我說過,「我們某某宗派的教會大是過上帝的!」這話是出自一位教會執事又是著名高等學府中的教授.這真可不是笑話!這更是一項致命的諷刺!這就是今日後現代性的基督教信仰的哀歌.

三/ 信徒離場!

以上的處境造成教會中一種牧者及教會領袖不敢正視的現象,就是 : 信徒離場(Church Leaving)!有人向我說過,按教會更新運動的一份統計指出,每主日到教會崇拜的基督新教信徒,佔全港人口有5\%的.但是,實質上另外有15\%的基督新教教徒是不去主日崇拜的.  這表示甚麼?值得思考.我在\url{http://exchristian.hk/wiki/doku.php/}這個網頁中讀到以下的名單 :

■黃霑 ■黃子華■黃耀明 ■劉敏儀(程翔妻子)■雲海 ■杜葉錫恩 ■胡恩威 ■周兆祥 ■康熙 ■蘇守忠 ■林夕

■倫永亮  ■梁文道 ■謝君豪 ■張小嫻 ■陳潔靈 ■楊傳廣 ■張柏芝 ■鄧麗君 ■孫中山 ■曹宏威 ■胡適 ■杜汶澤

以上只是一個網絡中提及一些離場教會: (本地)的名單\url{http://exchristian.hk/wiki/doku.php/}  當然,離場教會的人還多得很多.只是他們離場教會不等於他們離開上帝,也不等於放棄信仰.在這些人中,有些仍是篤信上帝.只是許多時教會崇拜的講道中,因為道也不在場,而崇拜上帝往往成了追隨名牧的崇拜.真心追尋上帝的人又怎能受得起呢!我想我們要反省.套用英國詩人E.T. Elliot 的一句話 :「一種宗教不僅需要有一群懂得自己正在做甚麼的牧師,更需要有一群懂得他們正在做甚麼的崇拜者.」 .這表示甚麼?這就是一般信徒太受後現代主義的影響,但是教會領袖卻是太受前現代主義及現代主義的影響.

後現代的基督徒知識份子對地方教會的離場,最少顯示出兩個現象:一是教會的後知後覺,教會領袖仍是持守聖經主義,凡事以經文為主導,而非以基督之心和上帝的愛來作為主導.那是說,當有問題發生時,大多去尋求聖經經文的支持,而非以基督之心去體諒對方.二是神學教育的後知後覺.所謂神學上的後知後覺是指到我們太受西方神學模式和學科所主導,凡事都向西方神學看.我們好像仍未找到屬於自己的神學教育方向和軟體(soft ware).在此我們要衷心感謝西方教會和神學教育,使我們對基督信仰有更新和更深的認識.只是,我們不應當就此滿足.我們應當是時候去建構出屬於自己文化的學科制度和神學概念.這才是緮現代性所賦於我們的良機!

香港華人教會應有的醒覺

從以上離場信徒名單觀察,他們大多都是知識份子,或是高級知識份子,又或是公共知識份子.換言之,他們都是有過一定文化教養的人,或是對文學,文藝有素養的人.這是否教會的後知後覺?是否我們把文化與文藝文學推之離場?教會是否只會活在前現代主義的屬靈權威主義,聖經主義或是仍停留在現代主義的唯獨理性主義,聖經主義或神學獨尊的教條中?這是現實,教會可以我行我素,但是一些有識之士的信徒因為教會把文化與文藝去場化,他們也離場教會.難道我們就可以如此下去,自嗚清高?我們究竟是生命為了生活,還是生活是去深化生命?教會的存在是為了奉獻和人數的數字,還是為了上帝的榮耀?

結  論

讓我引用我保留一份多年前的短文作結束.我多年前一位尋道者向明報投寫的文章,雖然時變勢遷,但至今天還是落地鏗鏘有聲.容我在此作為我的論.這位署名李秀明讀者對當時一位教會領袖的言說感到憤慨.她說 :

記得一位朋友說過,相信上帝的日子就是擁著一團幸福,因為一切是非善惡都有了判斷的標準,生命的來與去已知道根源和歸宿；對於不明所以的事有了交托,毋須徬徨.外邊世界擾亂紛爭,可以探頭出去觀望,然後安全地關上窗戶,燃點室內的燭光,不怕給人扯走.以穩定的信念觀看生活中的浮光掠影,不單處變不驚,且樂趣無窮.如此這般,我們就平平安安,喜喜樂樂,帶著感恩生活下去….然而,現實生活果真如此平面和簡單嗎?我們是否在自欺欺人?

作為一個尋道者,我相信人生的終極意義或基要價值可以從聖經啟示中窺探；但作為一個平凡人,我亦深信真誠地面對現實,有助洞悉聖經啟示的意義和價值.…..可是多少時候,基督徒的信仰表達,聚會形式,以至術語,心態,見識,確實令人難以信服.究其原因,就是我們對信仰的詮釋缺乏現實基礎；尤更甚者,大部份基督徒對信仰的處理,只作抽象空泛的思維概念,在信徒自己的小圈內兜兜轉轉,互相確認(reinforcement).於是,在宗教程序,術語,心態上不斷重複一些消化,體會或省思不足的教條,儼如活在信仰的幻象中.

其實,基督徒勇敢地回應世界並非媚俗,因為惟有誠實地演繹現實的生活,才能慢慢體悟人生終極意義,甚至有天超越現實的局限.可是,很多時候,當教會面對一些極具爭議的課題時,一是迴避,一是搬出一大走空泛教條；然而,卻沒有認真考慮問題的複雜性和個別性.…… 教會這樣閉門造車的膚淺教條,只會慢慢去向自欺欺人的胡同,甚至有天可能會作出否定現實的荒謬行為或反智取向.結果,信徒有時不得不將信仰「暫停營業」,按常理來應付某些生命的抉擇；或是將信仰變成一種虛偽的生活裝飾品,閒來粉飾一下不中用的生活；焉有理智而又具勇氣的人,誠實而沉默選擇自己應走的道路,奔向那思想自由的天地.

「我們必須明白,信仰並非一張行為清單,亦非且一列規範措施,更非一套賞罰機制.信仰是人類生活的經驗,藉此而令人真實地體現上帝的存在.當我們認真地去生活,去愛,並且有勇氣而誠實地承擔生活每種經驗,我們便已投入人生的敬拜之中,消除種種障礙,真真正正地演繹我們的信仰人生.在演繹過程中,每一個人的方式,快慢,進退乃同起跌都不同,正是這份自由和不同的經驗,才顯出創造者的奇工.深切希望基督教群體可以(用)認真而有智慧的批判態度,正視一個新時代的信仰,而代替一種盲目的篤信；以愛來代替審判,關心生命多於宗教的規條.」

這位在明報2000年一月的明報副刊中表白心聲的尋道者提問,帶出了羅馬書十一章至十二音2節的一個核心問題 : 靈性與文化的問題.豈不正是保羅在羅馬書十二章說 : 所以弟兄們,我以神的慈悲勸你們,將身體獻上,當作活祭,是聖潔的,是神所喜悅的；你們如此事奉乃是理所當然的. 不要效法這個世界(這是走向凋謝的信仰),只要心意更新而變化,叫你們察驗何為上帝的善良,純全,可喜悅的旨意.(這才是成長的信仰).

由此說來,教會必須重新調校前現代性的文化,特別是傳統主導及權柄主義.在現代性的心態上,我們需要重新糎訂理性主義的意義和限制性,而非甚麼都以理性來作斷定.我們也需要重新檢討後現代性的個人主義,自由主義及消費主義,進而進入一種新的,後後現代的,屬於基督的人神開係和人與人之的關係.

後現代性時期必然也有它的極限,而且正在尋求新的方向.這就是如何建構一個後後現代性的社會 : 重新建立新的價值觀,新的人際關係,更是新的神人關係.也就是今日教會的新契機.這就是把信仰重新融入上帝的奧秘和大愛中.以愛來表達我們的信仰,以上帝的奧秘作為我們的認信.這就是保羅所強調的上帝的奧秘和上帝的大愛.這才是今日教會當持守的信仰.



\newpage

\section{宣教的行動}
\label{sec:339}
\textbf{林崇智牧師 - 宣教的行動}
\newline
\newline
\hyperref[sec:338]{< < < PREV SERMON < < <}
~
\hyperref[sec:toc]{[返目錄]}
~
\hyperref[sec:340]{> > > NEXT SERMON > > >}


講題:宣教的行動　Mission in Action

講員:林崇智牧師

場合:香港中文大學崇基學院禮拜堂

日期:2016年10月16日

西方教會失去生命力,人數不斷減少,已是不爭的事實,對我來說近年有兩個很埋身的體驗,一個是去年我參加英國循道會週年大會,他所顯示的教會是一個面向世界,服事世界的教會.會議討論公平貿易,以巴和平,世界的貧窮,饑荒,環境污染,人權,自由等重要問題.但最重要,甚至是致命問題大會逃避去討論,就是過往十年,英國循道會教友人數少了十萬人,人數由原來的30萬減至20萬,很多禮拜堂關閉.討論聖公會合一,其中一個原因是同一教區內,大家教友人數都很少,教友若集中在幾所禮拜堂就可減少營運支出.在教會聚會中絕大部份是長者和幼童,若未來十年持續減少十萬,英國循道會便減少了一半教友.教會自顧不下,那位重要和偉大的社會關懷看來都無以為繼.

今年是美國衛理公會創辦人亞斯理會督逝世200周年,我們往美國參加世界循道衛理宗協會後,有幾位牧師繼續進行9日尋根之旅.回來後,同工們說:「那些他們訪問的教會很大很美,但崇拜人數只有幾十人,美國衛理公會本士教友人數持續下跌,西方教會人數不斷減少.有很多著作出了研究和解釋,但可以肯定的是教會沒有全面地實踐宣教使命.今日教會沒有初期教會被聖靈所催逼,被宣教的火焰燃燒,這肯定是其中一個重要原因,教會不積極委身宣教,教會就會奄奄一息,死亡是遲早的事.

著名教會歷史學家約翰佛斯德,寫了一本重要的小書叫「教會存在的意義」(Why the Church),他在第一章結尾說:「基督徒懷抱堅強信心,認為他們以教會和上帝選民的身份存在於世界,是能夠給與世界和人生重大意義,上帝創世的原意已被罪惡切底破壞,由於耶穌基督的降生,現在人類可以重新開始得著改變,上帝的計劃已經在主耶穌裡完成.今日主耶穌仍然在世界上,祂仍然工作,祂仍然救贖世人,只是,他不再只有一雙手了,而是千千萬萬的人準備與他同工,信徒的聲音對各國各族的人傳出上帝的信息,教會是耶穌基督的身體,代表上帝的靈仍留在這世界上」.這是真實嗎?今日教會是一群蒙上帝呼召的人,藉著宣教的生命和行動,帶領人到上帝面前,在世人面前榮耀上帝.所以教會是為著委身宣教而存在,每一個基督徒的生命都必需要成為一個宣教的記號.

宣教可以用不同方式表達,今天我將宣教分作三個重點表達形式來分享,第一,宣教就是宣講.每一個基督徒是「要奉主的名傳悔改赦罪的道,直到萬邦」.他們要像聖奧古斯丁說,「要在所有可能的地方,盡你一切可能方法,向所有人宣講天國的福音」,所以基督徒要照保羅對提摩太的教導「務要傳道,不論得時與不得時」,基督徒永不忘記自己有一個傳道者的身份.使徒行傳告訴我們,那些因司提反被迫害的人,他們隨走隨傳,天國福音就甚至隨著基督徒受迫害而四散,傳到世界各地.信徒無論在任何情況下,順境或逆境,都宣講天國福音,是初期教會快速發展的原因.

當耶穌在世上日子快要結束,祂只有幾個星期,將他在世上未完成的工作交給門徒去完成,他要門徒先親身經歷宣講天國福音是甚麼回事,於是他差派70位門徒,兩個兩個出去到山區,向未信的人去宣講福音,他告訴門徒「要收的莊稼,作工的人少」,但這是不一項容易的工作(今日仍是一樣),因為「我差你們去如同羊羔進入狼群」,你們會被人拒絕,被人反對,甚至被人排斥攻擊,但你們要知道,你們只是我的代表,不要太介懷別人的反應!「聽從你們的就是聽從我,拒絕你們的就是拒絕我」.但你們不要擔心害怕,因為當你們宣講時,必定有聖靈能力與你們同在.聖經說這七十人回來匯報他們的宣講工作,在聖靈大能下有很好的成果.聖經記載耶穌聽到門徒的匯報後,充滿極大快樂.宣講福音能討耶穌喜樂.

今日彼得後書提醒我們基督徒,要為宣講福音常作準備,「你們要心裡尊主基督為聖,有人問你們心中盼望的緣由,就要常作準備,以溫柔的心回答各人.」問題是我們對上帝及有關上帝的事,有深入認識嗎?我們對耶穌基督的真理認識嗎?我們在信仰上有沒有把握,爭取學習和進深的機會?還是我們雖然表面上信主多年,但仍如聖經批評彼得在主耶穌登山變象時的表現,「他根本不知道自己在說甚麼」!我們面對的世界充滿挑戰,膚淺的信仰是不能回應世俗的挑戰.我們若不在信仰知識和見識上進深,我們的宣講變成鳴的鑼,響的跋不能回應人的需要.

宣教第二個表達重點,宣教就是見證,路24:49耶穌對門徒說:「你們就是這些事的見證」,我們若勉強的區別,宣講的重點是你對信仰的認識,但見證不再是指頭腦上對上帝和有關上帝的信仰認知,不再是分析和理解問題,而是我們生命中有沒有經歷過上帝的愛,我們有沒有經歷聖靈的大能,在我們生命中帶來的拯救和改變.因為見證人是去見證你所看到,所聽到,所經歷到的生命,將上帝的公義,仁愛,平安和喜樂在你的生命中活出來.作見證是聖靈透過我們生命的工作,所以使徒行傳說:「門徒等候直到五旬節藉著聖靈降臨,他們才能成為耶穌的見證,沒有聖靈大能就沒有見證,所以徒1:8講得清楚,耶穌對門徒說,「聖靈降臨在你們身上,你們就可以成為我的見證」.

使徒行傳給我們許多信徒成為耶穌基督復活的見證事蹟,聖經告訴我們,門徒因著聖靈的同在,他們從害怕變成勇敢,不單在人群中宣講釘十字架的基督,高聲說除基督以外別無拯救.他們對迫害他們的祭司和守殿官也毫不害怕,對恐嚇他們的人說:「聽從你們不聽從上帝合理不合理,你們自己考慮吧.我們所看見所聽見的不能不說,順從上帝不順從人是應當的」.聖經記載這群當權者「看見門徒的膽量,又看出他們原是沒有學問的小民就感到很希奇」.門徒領受聖靈後,不單以勇敢來見證基督,他們在困苦患難中,甚至多次被鎖在監裡,不單沒有埋怨,沒有驚惶,相反的,他們在逆境中充滿喜樂平安,同心合意的禱告,不斷讚美上帝.聖經說這種生命見證充滿能力,吸引獄卒歸向基督.

成為上帝的見證,重點不是頭腦,而是我有沒有經歷到聖靈在我生命中的工作?我有沒有將喜樂,平安,在我的生活中活出來?我的生活有沒有活出耶穌基督的生命.耶穌愛那些妓女,稅吏,罪人,痲瘋病人,關懷社會中那些被歧視,被邊沿化,被剝奪權利的,他樂於與被人認為可恥的人為朋友.今日我們的生活,我們的群體也能活出這樣的生命見證嗎?耶穌的國度充滿包容和寬恕,他愛那些仇恨他的人,祝福那些惡待他的人,他的說話充滿祝福,沒有咒詛,我們的生活能為這份寬恕,祝福別人的大愛作見證嗎?上帝的國度教導我們為上帝而活,要積財寶在上,我們的安全和保障是在上帝那裡.我們對錢財態度和運用,能成為上帝不在乎物質豐富的國度作證嗎?當我們認真的將不一樣的信仰活出來,別人就會因此認識我們是屬基督的.今日上帝拯救這個世界的策略,就是透過聖靈去尋找你去成為見證人.

最後一個想同大家分享的宣教的表達重點就是服務,在聖經的傳統,往往指服務的對象是去服待社會中的貧困者.我們要特別留意,路加福音第十章記載耶穌差派70人兩個兩個的出去宣講天國福音後,便隨即提到好撒瑪利亞人的信息,以表達宣講福音和服侍困苦者是不能分開.聖經講基督徒的愛不能只在言語和舌頭上,我們不能對赤身露體缺乏飲食的人說:「平平安安的去吧,願你穿得暖,吃得飽,卻不給他所需要的,若是這樣,你們的信仰是死的.基督徒的愛是去服侍有需要的人.耶穌用好撒馬利亞人的比喻去批評那些在信仰上表面看似很屬靈的人,這些按外貌被認為是應該有更多愛心的人,卻表現出對有需要的人冷漠不仁,視而不見,置身事外.反而那個被猶太人看為不是人的外邦撒馬利亞人,他不單看到這個困苦的人,還動了慈心,以行動包zhak 這個人的傷口,更作出長期的承擔,將他送到客店,請人照顧這困苦者,並願意支付一切費用.怪不得馬太福音25章耶穌講大審判時,那些被上帝稱讚,「你這蒙我父賜福的人,可以入來承受創世以來為你們預備的筵席」,不是你是否每個星期都出席崇拜,每次崇拜都奉獻,在教會中盡力參與事奉,參與詩班獻唱,帶領團契小組,這一切都是重要,但耶穌說:「被上主稱讚的人是那些看到別人餓時給他們吃,看見別人喝的時候,給他水喝.看到別人赤身露體時給他衣服穿.看到別人有病去照顧他,看到別人下在監裡去探望他的人.」宣教就是愛心的服待,尤其是服待那些困苦者.耶穌離世前特別提醒我們,他來到這個世界不是要受人的服待,乃是要服待人,於是他作為一位老師,立刻拿起一條手巾跪在地下為門徒,為不配的人去洗腳.並挑戰我們,你去照樣做吧!

我深願在座每位弟兄姊妹都能活出一份宣教的生命,就是努力學習真理,把握每一個機會宣講福音.追求經歷耶穌的愛和聖靈的改變,成為真實的見證,將上帝的仁愛,公義,喜樂和平安活出來,將眼晴看遠一些,有一份柔軟仁慈的心,敏感別人和社區需要,成為一個服侍困苦者的人.啊們!



\newpage

\section{終極勝利}
\label{sec:340}
\textbf{邢福增教授 - 終極勝利}
\newline
\newline
\hyperref[sec:339]{< < < PREV SERMON < < <}
~
\hyperref[sec:toc]{[返目錄]}
~
\hyperref[sec:341]{> > > NEXT SERMON > > >}


提摩太後書 4:6-8
\newline
\begin{longtable}{cl}
\hline
\hline
章節 & 經文 (和合本修訂版)\\
\hline
4:6 & \begin{tabularx}{0.7\textwidth}{X} 至於我，我已經被澆獻，離世的時候到了。 \end{tabularx} \\ \\ \relax
4:7 & \begin{tabularx}{0.7\textwidth}{X} 那美好的仗我已經打過了，當跑的路我已經跑盡了，該信的道我已經守住了。 \end{tabularx} \\ \\ \relax
4:8 & \begin{tabularx}{0.7\textwidth}{X} 從此以後，有公義的冠冕為我存留，就是按著公義審判的主到了那日要賜給我的；不但賜給我，也賜給凡愛慕他顯現的人。 \end{tabularx} \\ \\
[1ex]
\hline
\hline
\end{longtable}
講題:終極勝利　The Last Race

經文:提摩太後書4:6-8

講員:邢福增教授

場合:香港中文大學崇基學院禮拜堂

日期:2016年10月23日

我有一個習慣,是喜歡逛墳場.每當站在不同的墓碑前,讀著碑上亡者的資料,或是墓誌銘,都會引起不同的思想.各位弟兄姊妹,您有沒有想過,當您離開世界之後,您的墓碑上除了姓名,籍貫及生卒年月日等基本的資料外,會不會有墓誌銘呢?您期望別人站在您的墓碑前,讀到刻在上面的墓誌銘時,會怎樣認識您?以下是我特別從網上找到的幾篇墓誌銘:

十九世紀法國詩人生前自撰的墓誌銘:「米蘭人亨利．貝爾(Marie-Henri Beyle,1783-1842)安眠於此,他曾經活過,寫過,愛過.」

俄國大作家赫爾岑(Aleksander Herzen,1812-1870)的墓誌銘,總結了自己坎坷的命運與非凡成就:「他的母親路易莎．哈格和他的幼子柯立亞,乘船遇難淹死在海；他的夫人娜塔利雅(Natalia Zakharyina)患結核症逝世；他的十七歲女兒麗莎自殺死去,他的一對三歲的雙生兒子患白喉死亡.而他只活了五十八歲,但是苦難不能把一個人白白毀掉.他留下三十卷文集,留下許多至今像火一樣燃燒的文章,它們今天還鼓舞這人們前進.」

英國當代文豪蕭伯納(George Bernard Shaw,1856-1950)一直活到九十四歲,他的墓誌銘很風趣:「我早就知道無論我活多久,這種事一定會發生的.」

曾獲諾貝爾文學獎的美國二十世紀大作家海明威(Ernest Miller Hemingway,1899-1961)的墓誌銘十分簡短:「恕我不起來!」

還記得在三十年前參加青年團契一次夏令營,導師叫我們寫下自己的墓誌銘,坦白說,對一個二十歲的青年人來說,根本談不上對自己的一生作出怎樣的期望或總結；當年在紙上寫下的一切,很快便忘記得一乾二淨.記得我大兒子在童年時,跟我說在長大後的志願是要開餐廳,巴士司機,我並沒有認真看待.因為我相信一年後,他便會有更多新的志願.現在他已二十歲,仍然在尋索其人生路與志向.人生就是漫長的旅程與競賽,隨著自己的成長,在不同的階段,我們都有不同的體會.

保羅在《提摩太後書》第四章對自己一生所作的總結:「至於我,我已經被澆獻,離世的時候到了.那美好的仗我已經打過了,當跑的路我已經跑盡了,該信的道我已經守住了.從此以後,有公義的冠冕為我存留,就是按著公義審判的主到了那日要賜給我的；不但賜給我,也賜給凡愛慕他顯現的人.」(四6-8)如果保羅的墳墓上刻有墓誌銘的話,很大可能就是這段經文.每一次當我讀到這段自述時,內心都會有著複雜的感受.不曉得到我們生命快要終結的時候,是否也可以自信地作出同樣的宣告?──我已經打過美好的仗?跑盡當跑的路?守住所信的道嗎?我相信,如果有人在離世的時候,可以這樣肯定自己的一生,那就真是死而無憾了!

使徒保羅在他晚年的時候,語重心長地寫信給提摩太.當時保羅第二次被囚在羅馬,這次的處境比上一次更險峻.也許保羅亦得知自己的時日無多,在《提摩太後書》---一封被視為新約聖經中保羅最後寫的書信中,保羅總結了自己的生命.如果大家沒有忘記的話,保羅在《腓立比書》第一章裡,也同樣表達了希望早日離世,與主同在的心願.那是他第一次被囚在羅馬的時候.他雖然知道離世是好得無比的,但他卻為了教會的需要,而仍願意活著──「我為你們肉身活著更加要緊」(腓一24).在世的責任要保羅繼續活在世上.到他寫提摩太後書時,他大抵已經無所牽掛了.但是,他在心裡仍惦記著他的愛徒提摩太,一位年青的教會領袖.他願意在這封信中,向提摩太作最後的叮嚀.

保羅特別囑咐提摩太,「務要傳道,無論得時不得時總要專心」(提後四2).在這裡,「得時不得時」主要是指外在的客觀環境.保羅形容提摩太要面對一個拒絕福音的世代.人們將「厭煩健全的教導」,「掩耳不聽真理」.不僅如此,他們更會「隨心所欲地增添好些教師」,「偏向無稽的傳說」(四3-4).這裡所指的「教師」及「無稽的傳說」,大抵就是指社會上流行的各種異教.其實,保羅在此並不是預告一些提摩太不曾面對的處境,而是他切實的經歷.保羅傳道的生涯,時常遭遇被拒絕,被譏諷,被輕渺,他肯定每一位傳道者亦必如此,無一倖免.最重要的是,外在的困難會否動搖傳道者的身分與使命?又會否使信徒膽怯起來,失去了見證信仰的勇氣?保羅的提醒是:基督徒必須謹慎,忍受因著持守真道而來的苦難,並在這世代中作見證,我們如何從上帝領受真道,便要按著所領受的,責備,警戒及勸勉眾人,「做傳福音的工作,盡你的職分」(一5).

各位弟兄姊妹,今天我們身處的時代,同樣是一個容易令我們膽怯的時代,活在多元價值的時代,不少人質疑絕對真理的存在.基督徒再宣稱上帝是獨一真理,耶穌基督才是道路與生命,在這個世代中便顯然變得落伍,甚至是追不上潮流了.近年間,我們常常看見有評論指基督徒在表達其信仰立場時是死守著疆化的教義,與時代脫節.如何能夠緊貼時代,卻又不違背福音信仰?我在多元社會中,是否有傳講福音的勇氣?

保羅除了再三囑咐提摩太要堅守傳道職分外,再從兩方面來作出勸勉.第一,保羅用自己的生命來激勵提摩太.提摩太,不管外在環境如何,福音必須繼續傳揚下去.每一位門徒都有責任堅守己職,就好像把自己的生命向上主「澆奠」一樣,正如保羅在《羅馬書》所言,我們的身體,就是獻給上帝的「活祭」.

值得留意的是,保羅在總結自己的生命時,並不是從外在的成就來衡量的.他沒有以自己的能力與業績來自誇－－看,保羅建立了多少所教會……反倒是突顯自己如何完成託付.他用打仗,賽跑及守道來形容自己,正好反映出這種精神.真的,提摩太並不如保羅,我們也不能與使徒保羅相比.單從所成就的果效來比較,沒有人可以與保羅配比,我們實在有太多的虧欠.但保羅卻藉著自己的生命激勵我們,上主是否賜下公義的冠冕,最重要的,是看我們有否忠於所託──盡力那美好的仗,跑盡當跑的路,守住所信的道.

記得有一年我到美國費城西敏寺神學院開會,會議結束後,有機會逛逛費城市中心,除了著名的國家獨立公園外,亦參觀了國家自由博物館(National Liberty Museum).這所博物館給我留下了深刻的印象,因為館內除了介紹一些對人類歷史(特別是民主自由)作出重要貢獻的所謂偉人外,亦有一個以平民百姓為主角的展覽館.我看到三十一位先天及後天殘障者,如何克服種種挑戰的故事.離開博物館前,有一個展板給參觀者留言,我發現令大部分留言者心靈受到感動的,不是那些舉世知名的偉人,而是那些無名英雄.正如一段展覽的文字所言:他們樹立了最尚高的勇氣,意志及堅強的榜樣.縱然有著嚴重的缺陷,生命仍然可以與別不同!

第二,保羅除了用自己的生命來激勵提摩太外,亦強調基督徒必須立定在信仰的根基之上.他在第四章開始,便指自己的囑咐,是「我在上帝面前,並在將來審判活人死人的基督耶穌面前,憑著他的顯現和他的國度鄭重地勸戒你」(四1).這句經文,充分體現了基督教信仰的終末取向.基督耶穌業已在十字架上勝過死亡,現在祂正座在父上帝的右邊,將來定必再臨,審判活人死人,那時亦是上帝的國度實現的時候.

為甚麼保羅要以終末論來作為他的立論根基?因為基督耶穌的得勝與再來,正正成為基督徒盼望的源頭,並賜給我們面對及忍受苦難與逆境的勇氣.保羅之所以能堅守一切,豈是單靠他個人的能力嗎?當然不是!難道這不是終末的盼望賦予他的勇氣嗎?正是同一根基,讓歷世歷代的信徒,因著仰望主再來,而在當下不如理想的處境中,仍能守著真道,為主作見證.對活在此時此地的我們來言,是否因著末後的盼望而得著更大的勇氣呢?

早前巴西奧運會閉幕,令我想起一位基督徒的生命.1923年,英國田徑選手李愛銳(Eric Liddell)代表國家參加巴黎奧運會,他是100米金牌的大熱門.不過,由於比賽安排在禮拜日舉行,他堅持要參加教會的主日崇拜,故決定放棄比賽.許多人勸他,他只說:「我不在星期日出賽」.後來,另一位英國選手將400米的參選資格讓給他,這並不是李愛銳擅長的項目,但他努力苦練,最後不僅奪得400米賽冠軍,並以47秒6的成績打破世界紀錄.這項紀錄保持了35年,直到1959年才被打破.他以跑步來榮耀上帝,說:「我跑步的時候感覺到主的喜悅.」他的事蹟,後來改編成為電影《烈火戰車》(Chariots of Fire),獲得奧斯卡最佳影片.兩年後,他奉獻自己來到中國天津傳教,因為他的父親也是一位在中國的傳教士.臨行時,當地報紙為他送上祝福說:「他為中國的另一場賽跑,現已開始,心意堅定,筆直前衝,如跑奧運,倘若結果如何,此刻不能輕易得知,我們判定他會跑完,因速度屬於他」.後來有記者問他:「你真的很高興將生命奉獻在現在的事業上嗎?你難道不留戀那閃爍的鎂光燈,蜂擁的人潮,激動的觀眾,大聲的歡呼,珍貴的慶功酒嗎?」他回答:「當然,有時人會很自然地想到這一切,但我更喜歡我在中國所從事的工作,因為這方面的價值遠超過你所說的一切.你知道,這份榮耀是一個永不朽壞的冠冕.」

1941年,面對日本侵華,李將妻女送離中國,自己繼續留在河北的農村服事.後來他被日軍關進山東濰坊集中營.營內環境惡劣,生活艱苦,但李愛銳仍堅守崗位.1945年2月,李愛銳因病死在集中營內,臨終一天仍去探探其他病人,終年43歲.他的下半生事蹟,前陣子又被改編成另一齣電影《終極勝利》.

李愛銳在愛丁堡大學畢業時,引用了一句名言來勉勵眾人:「不論結果是一敗塗地,或是光榮勝利,一個人若盡力而為,那就是他的榮耀了」.後來他在集中營最後一段日子裡,曾經寫過一首詩,表達了他對人生的感慨:「每個人從出生到死亡,雖然都象是站在同一條跑道上,但每個人所做的事又是不同的,因此,生命的意義也便有所不同.」

終極勝利就是他所指的永不朽壞的冠晃,這是基督耶穌所應許賜下的,每一位基督徒要有信心在面對各種的疑惑與考驗.當保羅說:「至於我,我已經被澆獻,離世的時候到了.那美好的仗我已經打過了,當跑的路我已經跑盡了,該信的道我已經守住了.從此以後,有公義的冠冕為我存留,就是按著公義審判的主到了那日要賜給我的；不但賜給我,也賜給凡愛慕他顯現的人」的時候,是他對這分終極的勝利有信心.

我想起另一位人士為自己定的墓誌銘－－「長眠於此的這個中國人,曾做了他應該做的事,說了他自己應該說的話」.他的名字是劉賓雁.他在四十年代加入中國共產黨,在1957年因批評黨的錯誤而被打成右派,開除黨籍.1979年獲平反,但到1987年再次因敢言而被開除黨籍.1989年後,流亡海外.劉賓雁在1985年出版的《第二種忠誠》的報告文學,批評黨內不敢講真話.他形容那些對黨「老實聽話,從無異議」的人為「第一種忠誠」.但他卻指出,中國需要的,是「第二種忠誠」,即是敢於向黨提出不同意見的人,他們不受歡迎,「還要付出從自由,幸福直至生命這樣昂貴的代價」.這位一生身體力行的「第二種忠誠」的奉行者.2005年,劉賓雁在美國病逝,終年80歲.2010年,他的骨灰終於獲準返回中國北京,但他生前為自己所預備的墓誌銘:「長眠於此的這個中國人,曾做了他應該做的事,說了他自己應該說的話」,最終卻不能夠刻在墓碑上……不過,「劉賓雅」這名字,正好展現了一位中國知識分子對國家的「第二種忠誠」.雖然,他的生命因此付上代價,最後客死異鄉,即使骨灰獲準回到祖國,但卻仍無法按其遺願在墓碑上寫了墓誌銘.但是,我們仍盼望,我們仍相信,終有一天,－－「長眠於此的這個中國人,曾做了他應該做的事,說了他自己應該說的話」可以出現在他的墓碑上.因為,終極的勝利,一定不會是黑暗,也不會是罪惡的權勢,這是基督信仰的盼望與相信.

今天,在香港連骨灰龕位都難求,墳地更成為奢侈品.我曾跟家人說,安息後只須將骨灰灑出,就足夠了.不過,如果我仍可以為自己撰寫墓誌銘的話,我是多麼願意借用保羅的名言:「那美好的仗我已經打過了,當跑的路我已經跑盡了,該信的道我已經守住了」.或者,我會再附加一句:「邢福增,一個充滿軟弱的生命,藉信望愛而活」.

各位弟兄姊妹,您的墓誌銘的內容又是甚麼呢?這是你活著生命的總結.同時,如何活著,又在乎你所期望的終極勝利是甚麼?



\newpage

\section{有今生有來世}
\label{sec:341}
\textbf{鄧瑞強博士 - 有今生,有來世}
\newline
\newline
\hyperref[sec:340]{< < < PREV SERMON < < <}
~
\hyperref[sec:toc]{[返目錄]}
~
\hyperref[sec:342]{> > > NEXT SERMON > > >}


路加福音 20:27-28
\newline
\begin{longtable}{cl}
\hline
\hline
章節 & 經文 (和合本修訂版)\\
\hline
20:27 & \begin{tabularx}{0.7\textwidth}{X} 有些撒都該人來見耶穌。他們說沒有復活這回事，於是問耶穌： \end{tabularx} \\ \\ \relax
20:28 & \begin{tabularx}{0.7\textwidth}{X} 「老師，摩西為我們寫下這話：『某人的哥哥若死了，有妻無子，他該娶哥哥的妻子，為哥哥生子立後。』 \end{tabularx} \\ \\
[1ex]
\hline
\hline
\end{longtable}
講題:「有今生,有來世」　Life Now and in the World to Come

經文:路加福音20:27-28

講員:鄧瑞強博士

場合:香港中文大學崇基學院禮拜堂

日期:2016年11月6日

各位弟兄姊妹,早安.

讓我先講一句廢話,我們只熟悉我們熟悉的世界,對陌生的世界我們很陌生.讓我用另一句廢話去解釋一下剛才那句話,我們只能想像我們能想像的,對於不能想像的,我們是不能想像的.這些說話,講了等於無講.20世紀哲學大師維根斯坦(Wittgenstein)說,我們能講我們能講的,對於不能講的,就不能講了.這句聽起來似廢話的句子,卻是20世紀哲學思想裡的一大智慧.

地球人對外星人,有很多描述,有很多故事.最有趣味的,莫於若干年前,電視播放「解剖外星人」的片段.全港巿民圍觀電視,大小報章作出各種評論.現在回想那外星人的模樣,都會從心裡笑出來.那外星人的模樣,除了肚子比常人凸一點以外,外觀幾乎與人類一樣.不單這個外星人,其他關於外星人的描述,大多都是有眼有口,有手有腳,具有人類特徵.我們對外星人的想像,很難脫離「人類」這個熟悉的想像框架.我們只熟悉我們熟悉的世界,對陌生的世界我們其實真的很陌生.

我們的信仰世界,對我們而言,其實也是一個陌生的世界.但我們很不習慣陌生.剛剛這句話,又是廢話.我們當然很不習慣陌生,若習慣了,就不是陌生了.我們很不習慣信仰的世界,我們試圖用我們習慣的方式去構想這信仰世界.有好心的基督徒為了鼓勵信徒多作事奉,會如此說,「你多作事奉,在天堂裡,你便有更多冠冕.想想,你的冠冕又多又大,多令人羨慕.」這真是將人間的想像,帶進對天堂的想像裡.在人間,有競爭,有勝敗,有高下,如此,勝利者的冠冕才令人羨慕.難道,在天堂裡,還需要這冠冕去向其他在天堂裡的人炫耀嗎?還需要這冠冕去肯定自己的價值嗎?還需要這冠冕去吸引別人的目光嗎?我們以人間的邏輯,去想像天堂的邏輯.若天堂真是這樣,還是天堂嗎?

今日的講道經文,是:路加福音20:27－38.撒都該人以人間的想像,去扣問「非人間」的世界.耶穌如何回應他們呢?

路20:27　撒都該人常說沒有復活的事.有幾個來問耶穌說:

路20:28　「夫子!摩西為我們寫著說:『人若有妻無子就死了,他兄弟當娶他的妻,為哥哥生子立後.』

路20:29　有弟兄七人,第一個娶了妻,沒有孩子死了；

路20:30　第二個,第三個也娶過她；

路20:31　那七個人都娶過她,沒有留下孩子就死了.

路20:32　後來婦人也死了.

路20:33　這樣,當復活的時候,她是哪一個的妻子呢?因為他們七個人都娶過她.」

路20:34　耶穌說:「這世界的人(或譯:這時代的人)有娶有嫁；

路20:35　惟有算為配得那世界(或譯:那時代)與從死裡復活的人,也不娶也不嫁；

路20:36　因為他們不能再死,和天使一樣(或譯相似)；既是復活的人,就為神的兒子.

路20:37　至於死人復活,摩西在荊棘篇上,稱主是亞伯拉罕的神,以撒的神,雅各的神,就指示明白了.

路20:38　神原不是死人的神,乃是活人的神；因為在他那裡〔那裡:或譯看來〕,人都是活的.」

在這故事裡,提到撒都該人.撒都該人,是當時建制派的人物,十分現世的,也十分現實的.「現世」到一個地步,他們根本否認「復活」這回事,根本否認在這現世以外還有一個將來的世界.對他們來說,有今生,無來世.一切,都在今生.

如此現世的人,一般而言,都只花時間去經營現世,不太像那些敬虔的法利賽人,由朝到晚都去和耶穌辯論信仰.這次,有幾個撒都該人,要來和耶穌辯論,要去否定耶穌所講的來世,否定復活的可能性.於是,一種極現世的思維方式,和耶穌那種「極奇異」的思維方式相遇.

撒都該人問耶穌.

路20:28　「夫子!摩西為我們寫著說:『人若有妻無子就死了,他兄弟當娶他的妻,為哥哥生子立後.』

路20:29　有弟兄七人,第一個娶了妻,沒有孩子死了；

路20:30　第二個,第三個也娶過她；

路20:31　那七個人都娶過她,沒有留下孩子就死了.

路20:32　後來婦人也死了.

路20:33　這樣,當復活的時候,她是哪一個的妻子呢?因為他們七個人都娶過她.」

這番說話建基於申命記(申25: 5-6)的一條條例.有人娶了妻,死時無後,則他的兄弟要娶這女子,後來若生了兒子,則長子要歸死去的兄弟名下.這是要保存宗族的每條支脈,也為了保存各支脈的土地產業,這些土地產業可以繼續養活該支脈的老弱傷殘.

撒都該人的推理,都頗有趣.既然摩西律法有此婚姻條例,我們都當遵守.設若有一女子,嫁一個,死一個.若她先後嫁了七個男人,則復活時,她是誰的妻子呢?或者說,哪個男人可以得到這個女子呢?七個男子同時擁有一個女子是荒謬的,有歪倫常的,不可接受的.若復活事件令這不可接受的光景出現,則復活事件是不可接受的,不可信的.這麼古怪的問題,撒都該人竟然想得出,可以想像,他們問完這問題後,自己都會笑.

若撒都該人來到今天,他大概會問,若有一個人死去了,他的肝捐給了另外一個人.若這個人死時,又將這肝再捐給另外一人.如是者,七個人都死去了,他們都曾用過同一個肝.復活的時候,那個肝是誰的呢?

撒都該人真有趣,這種問題都想得出來.

我們看看耶穌如何回應.

路20:34　耶穌說:「這世界的人(或譯:這時代的人)有娶有嫁；

路20:35　惟有算為配得那世界(或譯:那時代)與從死裡復活的人,也不娶也不嫁.」

耶穌開宗明義說,不要用這個世界的想法,去言說那個世界；不要用現時代的情景,去構想將來的時代的情景.耶穌呼籲,井底的蛙,不要以井底的光景去論說井外的世界；夏天的蟲,不要以夏天的見識,去構想冰天雪地的世界；住在地球的人類,不要用「人類」的框架去描繪外星人.我們見的世界,是男婚女嫁的世界.但復活的世界,是一種我們想像不到的世界.我們的思想框架,是有娶有嫁的.但耶穌說的新世界,是不娶也不嫁的.我們世界的人之常情,在那個世界,不再是人之常情了.在「不可言說」卻又要「言說」的情況下,耶穌簡單地否定我們的習以為常.在那個世界,人是「不娶也不嫁」的.

撒都該人假定,這個世界有婚姻,那個世界也有婚姻.這個世界需要傳宗接代,那個世界也需要傳宗接代.這個世界需要由婚姻去保障土地擁有權,那個世界也需要由婚姻去保障土地擁有權.這個世界的女人角色是生育,那個世界的女人角色也是生育.這個世界由男人去擁有女人,那個世界同樣由男人去擁有女人.撒都該人以為,這個世界如何,那個世界也如何.撒都該人只熟悉他們熟悉的世界,對陌生的世界他們很陌生.

撒都該人的想法錯了,錯得很利害,錯在以現世的想法去設想將來的世界.耶穌指出,在那個世界裡,那個復活的世界裡,女人與男人的關係,不再由婚姻去界定.這種說法,在耶穌時代,意味著女人不再是「生仔機器」,不再由男人去定義她的存在意義,她也不再是男人的擁有物.男人和女人,在一個我們現在想像不到的世界裡,自由而豐盛地活著.再不受制於生物性的慾望,再不禁錮在刻板的社會制度裡,再沒有權力的比併.耶穌以其生命,以其死亡,以其復活,指向那個我們想像不到的世界.

耶穌繼續說:

路20:36　「因為他們不能再死,和天使一樣(或譯相似)；既是復活的人,就為神的兒子.」

我們是必死的人,面對人的必死性,不同的哲學家,宗教家提出不同的解救之道.希臘早期的哲學家柏拉圖指出,在必死的人的裡面,有一不朽的部分,即:不朽的靈魂.變易的塵世與不朽的靈魂相對立.變易的塵世由於會朽壞,是不好的,不真的,虛幻的,進而否定塵世生活的價值.不朽的靈魂是好的,是真的,不要受塵世影響.到人死的時候,靈魂從肉體的禁錮中釋放出來,進入永恆.

柏拉圖的想法,與耶穌的復活的想法,毫無關連.從復活觀念的角度看,柏拉圖的想法,也只是撒都該人的想法的一種變種.柏拉圖否定變易的塵世,然後以人間的邏輯,構想一相反於變易的「永恆」.然後,設定人性本身有這「永恆性」,即「靈魂不朽性」.最後,以人本身這「不朽」資源去保證自身的「永恆」.若撒都該人的思想特質,就是以本身的資源去構想「另一個世界」,則柏拉圖的做法,也是如此.

耶穌講論復活,復活不是「靈魂不朽」.「靈魂不朽」的想法,會否定這變易世界的意義.復活並不否定這個變易世界的意義,而是要指出,這個變易世界不是全部,這個世界有一我們不能構想的將來.塵世的生活是有意義的,這生活向我們打開來世之門.人是會死的,我們不需用「靈魂不朽」去否定這死亡的意義,也不需用自己的「靈魂不朽」去保證自己的不朽.我們死了,但神要以祂無比的恩典,將我們復活過來.復活,來自「全然他者」(wholly other)的神.若我們用自己的「靈魂不朽」這資源去保證自己的不朽,則這不朽世界仍是我們能構想的,它只是我們的塵世的反面.但若復活不源自我們自身的資源,而是來自神的恩典,則復活對我們而言,是全然陌生的,是來自神的全新生命.

耶穌說,復活的人,不會死,與天使相似.天使是什麼東西,如何生活,仍是不能構想的.只知,他們具有與塵世萬物不一樣的「不死性」.耶穌稱復活的人,是神的兒女.這生命實在是如何,現在仍不清楚,只知神的兒女不單是身分的改變,也是生命形態的改變.神的兒女與神有一不可再分離的關係,神的兒女的「永遠活著」,是由神的恩典去保證的.

最後,耶穌說:

路20:37　「至於死人復活,摩西在荊棘篇上,稱主是亞伯拉罕的神,以撒的神,雅各的神,就指示明白了.

路20:38　神原不是死人的神,乃是活人的神；因為在他那裡〔那裡:或譯看來〕,人都是活的.」

摩西在火燒的荊棘中與神相遇.神自我介紹,祂是「亞伯拉罕的神,以撒的神,雅各的神」.在這裡,耶穌的邏輯是,若神是活人的神,若祂是亞伯拉罕,以撒,雅各的神,則亞伯拉罕,以撒,雅各必然是活人.他們活在哪裡?他們活在神那裡.復活,不是「翻生」,不是「某個靈魂的不朽」,而是永遠活在神那裡.至於這是什麼樣的東西,我們很陌生.我們只能想像我們能想像的,對於不能想像的,我們是不能想像的.

撒都該人和耶穌的這番辯論,對我們有何意義?

一,我們不要像撒都該人那樣,以自己熟悉的世界,去構想神恩典下的陌生世界.不要用資本主義的邏輯,去構想信仰.不要用你自己的想法,去論說神.

二,讓耶穌講的復活新世界,激勵我們努力過信仰生活.不要以為這世界的一切就是全部.在復活的視野下,有今生,也有來世.

三,不要將「不朽」寄托在自己生命裡,「復活」只能源自神的恩典.要走出自我,要離開自義,要信賴神,要仰望神,要在神那裡重新校正自己的人生.

但願:榮耀歸於聖父,聖子,聖靈.阿們.

(後話:本人將早期的崇基講章結集為書,書名《敢於跟隨主》,基道出版,有興趣者可購買.)



\newpage

\section{末世凶兆如何自處}
\label{sec:342}
\textbf{許開明牧師 - 末世凶兆,如何自處?}
\newline
\newline
\hyperref[sec:341]{< < < PREV SERMON < < <}
~
\hyperref[sec:toc]{[返目錄]}
~
\hyperref[sec:343]{> > > NEXT SERMON > > >}


路加福音 21:5-19
\newline
\begin{longtable}{cl}
\hline
\hline
章節 & 經文 (和合本修訂版)\\
\hline
21:5 & \begin{tabularx}{0.7\textwidth}{X} 有人談論聖殿是用美石和供物裝飾的，耶穌就說： \end{tabularx} \\ \\ \relax
21:6 & \begin{tabularx}{0.7\textwidth}{X} 「你們所看見的這一切，日子將到，沒有一塊石頭會留在另一塊石頭上而不被拆毀的。」 \end{tabularx} \\ \\ \relax
21:7 & \begin{tabularx}{0.7\textwidth}{X} 他們問他：「老師，甚麼時候有這些事呢？這些事將臨到的時候有甚麼預兆呢？」 \end{tabularx} \\ \\ \relax
21:8 & \begin{tabularx}{0.7\textwidth}{X} 耶穌說：「你們要謹慎，不要受迷惑，因為將有好些人冒我的名來，說『我是基督』，又說『時候近了』，你們不要跟從他們！ \end{tabularx} \\ \\ \relax
21:9 & \begin{tabularx}{0.7\textwidth}{X} 當你們聽見打仗和動亂的事，不要驚惶；因為這些事必須先發生，但終結不會立刻就到。」 \end{tabularx} \\ \\ \relax
21:10 & \begin{tabularx}{0.7\textwidth}{X} 於是耶穌對他們說：「民要攻打民，國要攻打國， \end{tabularx} \\ \\ \relax
21:11 & \begin{tabularx}{0.7\textwidth}{X} 將有大地震，多處必有饑荒、瘟疫，又有可怕的異象和大神蹟從天上顯現。 \end{tabularx} \\ \\ \relax
21:12 & \begin{tabularx}{0.7\textwidth}{X} 但這一切的事以前，有人要下手拿你們，迫害你們，把你們交給會堂，並且關在監裡，又為我名的緣故拉你們到君王和統治者面前。 \end{tabularx} \\ \\ \relax
21:13 & \begin{tabularx}{0.7\textwidth}{X} 但這些事終必成為你們作見證的機會。 \end{tabularx} \\ \\ \relax
21:14 & \begin{tabularx}{0.7\textwidth}{X} 所以，你們要立定心意，不要預先考慮怎樣申辯； \end{tabularx} \\ \\ \relax
21:15 & \begin{tabularx}{0.7\textwidth}{X} 因為我必賜你們口才和智慧，是你們一切敵人所敵不住、駁不倒的。 \end{tabularx} \\ \\ \relax
21:16 & \begin{tabularx}{0.7\textwidth}{X} 連你們的父母、兄弟、親族、朋友也要把你們交給官府；你們中間也將有被他們害死的。 \end{tabularx} \\ \\ \relax
21:17 & \begin{tabularx}{0.7\textwidth}{X} 你們要為我的名被眾人憎恨。 \end{tabularx} \\ \\ \relax
21:18 & \begin{tabularx}{0.7\textwidth}{X} 然而，你們連一根頭髮也不會損失。 \end{tabularx} \\ \\ \relax
21:19 & \begin{tabularx}{0.7\textwidth}{X} 你們憑著堅忍，就必保全性命。」 \end{tabularx} \\ \\
[1ex]
\hline
\hline
\end{longtable}
講題:末世凶兆,如何自處?　How to Keep Your Cool When Facing Vicious Signs of End Times

經文:路加福音21:5-19

講員:許開明牧師

場合:香港中文大學崇基學院禮拜堂

日期:2016年11月13日

祈禱:上主,我們的救贖主:願我口中的言語,弟兄姊妹的心懷意念,都蒙上主的悅納.阿們.

引言

末世凶兆:這四個字,相信引起許多人的注意；

如何自處:也是許多人的關注!

1過去的一星期最受注意的國際事件是:特朗普爆冷勝出!

許多人可能質異美國人的選舉智慧,也質疑許多專家的預測如此失準!

一位敢言到狂言的人,有可能會引起種族矛盾,甚至性別間歧視的人,又欠缺從政經驗的人做總統?

一位有多段婚姻,私生活混亂的花花公子,美國基督教會竟有不少表示支持他當總統,因為他說過「支持一夫一 制」,他說就信!

2香港的立法會宣誓風波,愈演愈亂,立法會癱瘓,各政黨開會對立,人大要釋法,原本二位議員資格被質疑,現時演變成多位宣誓的議員的資格出現問題!

3國際上的恐怖主義於911後興起,美國的反恐戰爭,不單沒改善,被壓止,反而導致伊斯蘭國興起.一個原先信奉阿拉,講求和平的宗教,何以會演變成殺害大量無辜,傷害其他回殺的部族,製造國際恐怖的組織?

4過往在世界各地,不少不同政見的人,會人間蒸間,近年有馬來西亞成架飛機也會人間蒸發,不知所蹤!

說下去,還可以有很多,這些是否也是末世凶?

不少宗教典籍或流傳千古的傳說都有「末日」之說,

甚至有預言指2012年為世界末日,加上近年特大天災頻頻發生,2008年有汶川大地震,今年的日本311九級大地震等,聖靈塗炭,令末日預兆之說更囂塵上,恐懼日深.

近年香港電台:由資深傳媒工作者李燦榮及張可頤主持的《大凶兆》,請來學者,專家,一同以理性角度分析一系列發生於世界各地的所謂「異象」,深入剖析,讓觀眾對種種現象有更深入的瞭解和認識,消除不必要的憂慮與恐懼.

一連十一集節目中,探討的課題包括「末日密碼」,「世界大地震」,「動物離奇死亡事故」,「第三次世界大戰」,「UFO入侵」,「大瘟疫」,「超級火山大爆發」等等.

末世凶兆

在路21章,耶穌提出末世凶兆,

讓信徒有所提防!

他說出危臉人物,危難的事情,這都是不好聽,使人厭惡的,令人怕的!

但這些都是事實,耶穌不能忍瞞!

1危險人物:假基督,迷惑人,使人跟從

耶穌說:“因為將來有好些人冒我的名來,說,我是基督,”即指那些自稱是基督的人；主這話表明‘假基督’的出現,乃是聖地遭受浩劫以前的第一個豫兆.猶太歷史家約瑟夫曾記載:在主死而復活之後,至耶路撒冷被毀以前的四十年內,有好幾個假基督出現,並且糾集了一班黨羽扈從.

注意:假冒基督的現象,二千年來不曾中斷,且越演越烈,因為這也是末世必要有的情形(參可十三21-22).

“又說,時候近了；你們不要跟從他們,”‘時候’指聖地遭難之時；假基督必以末世將到為其中心資訊,來吸引人們跟從自己.

我們現在先來看 20世紀最後 20 年間,發生的與死亡,毀滅有關的集體自殺事件,都與跟隨假領袖有關:

1978 年 11 月美國印第安那州「人民教堂派」, 在金瓊恩的領導下吞毒藥集體自殺,共死 913 人,其中有 276 人為孩童.

1985 年菲律賓岷答那峨島 60 名阿塔族族人依其宗教領袖馬加亞農命令,吞農藥集體自殺.

1986 年日本本州西部和歌山市海濱有七具屍體, 遺書上說他們是「友誼與真相教會」,在宗教領袖去世後,要追隨宗教領袖而去.

1987 年南韓漢城出現三十二名信徒屍體, 他們跟隨一自認是女神的宗教領袖.他們吞的是不會致命的毒藥,但喉嚨全遭人切斷.

1993 年 2 月,美國德州「大衛教派」與警方對恃 50 餘天後,會眾引火自焚,但後來在屍身上多發現彈孔.宗教領袖克瑞許在跟警察槍戰中死亡.

1994 年十月, 在瑞士發現有 48 具屍體,屬加拿大的太陽神殿教,宗教領袖朱瑞的屍體在加拿大魁北克發現.

1995 年 3 月,日本奧姆真理教在地下鐵車站施放「沙林」毒氣,造成 5500 人死傷.

1996 年美國蒙塔那州東部出現「自由人」, 集住農莊,是宗教團體,有領導者,抗拒憲法與納稅,後經美審慎處理,未曾導致大衛教派事件重演.

1997 年 4 月一群資訊科學人士「天堂之門」集體自殺等候被接升天.

這些自殺事件,絕不只是民間落後信仰與低階層人士參與製造出來的事件,裡面也有屬於高科技,高知識階層的人士.很明顯的,這些集體自殺事件隨著時間迫近 1999 終結,而越來越密集.他們都相信團體中的「救世主」並徹底委身追隨；他們相信救世主所說而且全都是宗教事件.

2危險事情:

A完美的聖殿被毀,遭破壞,不好的事情出現!

聖殿被毀,代表所信上帝不能保護!

聖殿聯緊的中心動搖,步向亡國.

B可怕的人禍:民要攻打民,國要攻打國

C可怕的天災:地震,飢荒,瘟疫

D可怕的異象:

E可怕的迫害:

面對凶兆,如何自處:

1謹慎

不受假基督迷惑,不跟從假基督

【路廿一8】“耶穌說:‘你們要謹慎,不要受迷惑；因為將來有好些人冒我的名來,說:“我是基督.”又說:“時候近了.”你們不要跟從他們.”

﹝原文字義﹞

“謹慎”小心觀察,看得仔細,

“迷惑”引入迷途,欺哄.

2常存忍耐:相信上主保護,不要驚惶

【路廿一18】「然而你們連一根頭髮,也必不損壞.」

【路廿一19】「你們常存忍耐,就必保全靈魂(或作必得生命).」

3注意耶穌在本段經文所引伸的四個意識

四個問題很普遍的問題,卻是很有意義的問題,我們經常被問:

1你驚唔驚?或你擔唔擔心?關心你的人常問!

2你痛唔痛?或辛唔辛苦?醫生見病人常問!

3你做完未?幾時Deadline?你的老板或老師常問!

4你有無希望?

可能涉及1學生,學業?2中年,事業?3老年,健康?

第一個問題:涉及「危機意識」

第二個問題:涉及「苦難意識」

第三個問題:涉及「終未意識」

第四個問題:涉及「希望意識」

這四個意識是本段福音經課的重點,也是耶穌要跟隨他的人,不容忽略的重點!亦是猶太人自古到今非常重視的觀念.

在希羅時代,一般人只有這三個意識「危機,苦難,終末」；

然而,在這個意識後,猶太教與基督教是更進一步提出了第四重要的意義“希望的意識”.

一個具有以上四個意識的人,會成為智慧的人!

事實上,從民族歷史來看,猶太人在各方面都顯出他們的智慧,無論在學術上,經商上,宗教上,處世為人上,均有傑出的表現.

入們具有這四個意識,就不怕末世的凶兆或未來的各種挑戰和衝擊!

猶太人有獨特的時間觀

1.現今時間:當下的,可能是平穩安定的日子!需要有「危機感!」

【路廿一5】“有人談論聖殿,是用美石和供物妝飾的.”

﹝背景批註﹞“談論聖殿,”他們所談論的聖殿,指所羅巴伯所重建的聖殿(參拉三8)被毀之後,約在主前二十年,大希律為討好猶太人,再次重建聖殿,當主耶穌開始事奉神時,已經費了四十六年(參約二20),惟仍未完工.

根據猶太歷史家約瑟夫的記述,整個聖殿是在主後六十四年完工的.這殿建築宏偉,堂皇壯麗,建殿的大理石,有的甚至長達十公尺以上,另有巨型金飾,真是金碧輝煌,光輝耀目.

1.1這輝煌的聖殿給予門徒安全感:這殿反映了上帝的臨在!

以色列人的智慧建設如此偉大的聖殿!他們正處於泰平盛世的好日子.

【路廿一6】“耶穌就說:‘論到你們所看見的這一切,將來日子到了,在這裡沒有一塊石頭留在石頭上,不被拆毀了.’”

﹝文意批註﹞“你們所看見的這一切,”門徒所看的是聖殿的外表.

“將來日子到了,”‘將來’是指在主後七十年所要應驗的事.

“在這裡沒有一塊石頭留在石頭上,不被拆毀了,”

這句話在主後七十年按字義完全應驗了,當時羅馬人在提多率領下,徹底毀滅了耶路撒冷及聖殿建築物,甚至石頭也被撬開,以取得當聖殿被燒毀時從屋頂上熔流其間的金箔.也有人說,當時尚存留有一些石頭未被移動.到了朱理安皇帝的朝代,因他企圖重建聖殿,故把那些存留的頭也移去(使主的預言全部應驗),但他的建殿計畫半途而廢.

1.2耶穌要門徒 「居安思危」有心理的準備!

危險事情出現:

完美的聖殿被毀,遭破壞,不好的事情出現!

聖殿被毀,代表所信上帝不會保護他們!

2.末世時間:這世界要到達終末,世界充滿苦難,天災人禍!

需要有「苦難意識」與「終未意識」!

可怕的人禍:

2.1可怕的戰爭:民要攻打民,國要攻打國

【路廿一9】“你們聽見打仗和擾亂的事,不要驚惶；因為這些事必須先有,只是末期不能立時就到.’”

【路廿一10】“當時耶穌對他們說:‘民要攻打民,國要攻打國；”

2.2可怕的天災:地震,飢荒,瘟疫

2.3可怕的異象:

【路廿一11】“地要大大震動,多處必有饑荒瘟疫；又有可怕的異象,和大神跡,從天上顯現.”

2.4可怕的迫害:

【路廿一12】“但這一切的事以先,人要下手拿住你們,逼迫你們,把你們交給會堂,並且收在監裡,又為我的名拉你們到君王諸侯面前.”

3.未來時間:上帝或基督要來作王,經審判後,世界更更新,新天新地出現!需要有「希望意識」!

在苦難中,仍有希望和恩典

3.1不要驚惶:

3.2作見證:上主予機會

【路廿一12】“但這一切的事以先,人要下手拿住你們,逼迫你們,把你們交給會堂,並且收在監裡,又為我的名拉你們到君王諸侯面前.”

【路廿一13】“但這些事終必為你們的見證.”

“終必為你們的見證,”指上述的情形乃是為主作‘見證’的好機會,顯明所信的是怎樣的一位主.

3.3不要預先籌算:上主賜口才,經歷上主的帶領!

【路廿一14】「所以你們當立定心意,不要豫先思想怎樣分訴；」

【路廿一15】「因為我必賜你們口才智慧,是你們一切敵人所敵不住,駁不倒的.」

3.4常存忍耐:相信上主保護

【路廿一15】「因為我必賜你們口才智慧,是你們一切敵人所敵不住,駁不倒的.」

【路廿一17】「你們要為我的名,被眾人恨惡.」

【路廿一18】「然而你們連一根頭髮,也必不損壞.」

【路廿一19】「你們常存忍耐,就必保全靈魂(或作必得生命).」

﹝原文直譯﹞「要在你們的恆忍裡得著你們的魂生命.」

﹝原文字義﹞「忍耐」恆忍,歡喜著忍受；「保全」得到,買得；「靈魂」魂生命,氣息.

3當擁有四大重要的意識

四大重要的意識:

【】一隻老鼠跌落大米缸?

看見米缸那麼多米,竟直是牠的天堂,三餐無缺,牠就不想跳出!

牠日吃夜吃,米缸的米愈吃愈少!

有一天,牠突然發現米缸的米,給給自己光光了!

牠愈吃愈肥,缸底離缸口又太高,想跳都跳不起來了!

結果?

缺的問題是:欠缺危機感!

心理學家告訴我們:

人受兩股感覺操傱:

1安全感

2不安感

這是對立的感覺,如何平衡?

在安全感中,要有危機感!

否則人會陷入危機中而不自覺!

在危機感中,不安中要達到安全感!否則人在長期不安中,焦慮中,人都巔!

當人從危機感進到安全感後,仍需保持” 危機感!

因此這是一種生命的循環!

【】你試過一條難過的馬路:

車子,單車,行人可左右都有,前後都有!

你的危機感一定要強,才能保命,或保住不被撞傷!

人需要安全感,

> 也需要危機感!

1危機意識

一位智者說:人生最大的危機就是沒有危機感

危機感同時也是一個人進取心的源泉,是一個人成長發展的重要動力.

一個人失去了危機感,就會變得安於現狀.裹足不前,等待他的只有被淘汰的命運.--《只有危機感強烈的人才能生存》作者:王非庶.

人感覺到有事物威脅到自身,並為此緊張,這種緊張是正向的!

“錢包危機感 ”, 婚姻危機感,健康危機感,死亡危機感,

失敗危機感,退休危機感等.

危機意識

孟子所說:“生於憂患,死於安樂.”

是古來的名訓；意識於憂患的存在,因而預想困境的來臨或挑戰,提早做備戰的策略及行動,將所有的策略及反應預存在腦海裡,並動員適當的腎上腺素以便縮短反射時間,則「凡事豫則立」,什麼挑戰皆不足懼矣.這在人生是如此,財務是如此,運動是如此,當然健康也是如此.

要時時保持這種危機意識並非渾然天成,必須勉力維持,但習慣也就成自然；而且,危機意識過猶不及都不甚好.

缺少危機意識,很容易粉飾太平,爭功諉過,甚至麻木不仁.

危機感是一種心理狀態,其存在不一定是事實的逆境和困境.聰明的個人,聰明的企業家,聰明的政府,都善於在逆境下勇敢面對危機感,又善於在順境中保持憂患意識,使自己能夠堅持不懈地努力.所謂居安思危,未雨綢繆,有備無患,就是這個道理.

微軟的比爾•蓋茨總是感到危機感的緊迫存在:“微軟離破產永遠只有18個月.”

海爾的張瑞敏總是感覺:“每天的心情都是如履薄冰,如臨深淵.”

聯想的柳傳志總是認為:“你一打盹,對手的機會就來了.”

百度的李彥宏經常強調:“別看我們現在是第一,如果你30天停止工作,這個公司就完了.”

哈佛商學院教授理查•帕斯卡爾所說的那句名言:

“21世紀,沒有危機感是最大的危機.”

2苦難意識

自古到今哲學家都問:「存在有何意義」

很多時乃在痛苦中人也發問這問題.

我們感到折骨之痛苦,會問︰為甚麼存在要受苦?

痛苦有何意義?

痛苦地存在有何意義?

面對痛苦:

1無奈地叫苦連天,怨天尤人,憤世疾俗!

2自憐自虐:以苦易苦

人們總傾向於借助自憐,或自我譴責,甚至自虐等去逃避當下的痛苦,企圖藉由一種更大的痛苦而終止目前無論由失去愛或傷害人而導致的痛苦,即是將自己陷於痛苦的氛圍中.

正如給蚊叮了一口,傳統的做法就是拍打叮口,使它產生另一種痛,蓋過一小點範圍的痛,因此自虐的行為表面看來是把自己陷於痛苦的深淵中,實質上仍是擺脫痛苦的手法.縱使它不至帶來歡樂或欣喜,起碼是抽離了特定的內容,使人暫時擱置困境.弔詭的是,這種以苦易苦的方式,其實依舊離不開追求舒泰的生物本能.

然而,自虐者的潛意識卻不知道,這種臨時的舒泰效果將導致更深的痛苦,這是第二重的弔詭︰拍打叮口只能短暫地轉移痛的感覺,之後蚊叮的痛癢重現,再加上拍打而來的痛楚,造成雙重的苦痛.

這樣的結果又似乎與自虐者意識的目標 ( 要自己痛苦 ) 相吻合.於是這出乎意料的痛苦在效果上滿足了自虐者自覺的自虐心理.而真正想逃避痛苦的意圖卻給埋藏了.

3轉化苦痛為祝福

智者:努力尋找苦對生命的意義.

苦難是一種磨練,也是一種修煉,所以苦難使人剛強!

把吃苦當吃補!

3“末世及終結的意識”

終結意識:世界會終結,生命會終結,任何物件都會有終結,因著會終結:人產生了兩個對生命的態度:

一是懂得珍惜所擁有的人和物!珍惜我的愛人,家人,友朋

二是不要太執著所失去的人和物!我失去郵票簿,失去的一件衫

末世論:

末日就是基督的再來(parousia,林前十五23；帖前四15),是表明人的歷史來到終局(故言末日),基督的國度(時間與空間)已然展開.希望在明天!

4「希望意識」

猶太人及初期信徒相信:大災難後,上帝或基督要來,將會有新天新地；苦盡甘來,前境充滿希望.

【】三隻青蛙 J067

三隻青蛙

一同跌落鮮裝奶桶中,各有不同的反應.

第一隻: 「這是上帝的旨意!甘願順服!」

第二隻: 「這隻木桶太深了!我實在深有辦法走出去阿!」

第三隻: 「真不幸我跌落在這裡,我可以怎樣脫離困境?

我一定要出去,我不能死在這裡,於是不斷游來游去!」

結果:第一及第二隻都逃不出來了!

第三隻青蛙踫到硬硬的東西,因為鮮奶變成奶油,青蛙站在奶油上,跳出來了!

總言之

面對末世凶兆   ,願你謹慎常有這四大意識?

危機意識,苦難意識,終末意識,希望意識!

在面對末世凶兆?

1我驚唔驚?危機意識!好

2我痛唔痛?苦難意識!好

3我還有多少時間?終末意識!好

4我有無希望?記住上帝是你的希望!希望意識!好



\newpage

\section{化干戈為玉帛}
\label{sec:343}
\textbf{林豪恩先生 - 化干戈為玉帛}
\newline
\newline
\hyperref[sec:342]{< < < PREV SERMON < < <}
~
\hyperref[sec:toc]{[返目錄]}
~
\hyperref[sec:344]{> > > NEXT SERMON > > >}


以賽亞書 2:1-5
\newline
\begin{longtable}{cl}
\hline
\hline
章節 & 經文 (和合本修訂版)\\
\hline
2:1 & \begin{tabularx}{0.7\textwidth}{X} 亞摩斯的兒子以賽亞所見，有關猶大和耶路撒冷的事。 \end{tabularx} \\ \\ \relax
2:2 & \begin{tabularx}{0.7\textwidth}{X} 末後的日子，耶和華殿的山必堅立，超乎諸山，高舉過於萬嶺；萬國都要流歸這山。 \end{tabularx} \\ \\ \relax
2:3 & \begin{tabularx}{0.7\textwidth}{X} 必有許多民族前往，說：「來吧，我們登耶和華的山，到雅各神的殿。他必將他的道教導我們，我們也要行他的路。」因為教誨必出於錫安，耶和華的言語必出於耶路撒冷。 \end{tabularx} \\ \\ \relax
2:4 & \begin{tabularx}{0.7\textwidth}{X} 他必在萬國中施行審判，為許多民族斷定是非。他們要將刀打成犁頭，把槍打成鐮刀；這國不舉刀攻擊那國，他們也不再學習戰事。 \end{tabularx} \\ \\ \relax
2:5 & \begin{tabularx}{0.7\textwidth}{X} 雅各家啊，來吧！讓我們在耶和華的光明中行走。 \end{tabularx} \\ \\
[1ex]
\hline
\hline
\end{longtable}
詩篇 122:1-9
\newline
\begin{longtable}{cl}
\hline
\hline
章節 & 經文 (和合本修訂版)\\
\hline
122:1 & \begin{tabularx}{0.7\textwidth}{X} 我喜樂，因人對我說：「我們到耶和華的殿去。」 \end{tabularx} \\ \\ \relax
122:2 & \begin{tabularx}{0.7\textwidth}{X} 耶路撒冷啊，我們的腳站在你門內。 \end{tabularx} \\ \\ \relax
122:3 & \begin{tabularx}{0.7\textwidth}{X} 耶路撒冷被建造，如同連結整齊的一座城。 \end{tabularx} \\ \\ \relax
122:4 & \begin{tabularx}{0.7\textwidth}{X} 眾支派就是耶和華的支派，上那裡去，按以色列的法度頌揚耶和華的名。 \end{tabularx} \\ \\ \relax
122:5 & \begin{tabularx}{0.7\textwidth}{X} 他們在那裡設立審判的寶座，就是大衛家的寶座。 \end{tabularx} \\ \\ \relax
122:6 & \begin{tabularx}{0.7\textwidth}{X} 你們要為耶路撒冷求平安：「願愛你的人興旺！ \end{tabularx} \\ \\ \relax
122:7 & \begin{tabularx}{0.7\textwidth}{X} 願你城中有平安！願你宮內得平靜！」 \end{tabularx} \\ \\ \relax
122:8 & \begin{tabularx}{0.7\textwidth}{X} 為我弟兄和同伴的緣故，我要說：「願你平安！」 \end{tabularx} \\ \\ \relax
122:9 & \begin{tabularx}{0.7\textwidth}{X} 為耶和華－我們神殿的緣故，我要為你求福！ \end{tabularx} \\ \\
[1ex]
\hline
\hline
\end{longtable}
羅馬書 13:11-14
\newline
\begin{longtable}{cl}
\hline
\hline
章節 & 經文 (和合本修訂版)\\
\hline
13:11 & \begin{tabularx}{0.7\textwidth}{X} 還有，你們要知道，現在正是該從睡夢中醒來的時候了；因為我們得救，現在比初信的時候更近了。 \end{tabularx} \\ \\ \relax
13:12 & \begin{tabularx}{0.7\textwidth}{X} 黑夜已深，白晝將近。所以我們該除去暗昧的行為，帶上光明的兵器。 \end{tabularx} \\ \\ \relax
13:13 & \begin{tabularx}{0.7\textwidth}{X} 行事為人要端正，好像在白晝行走。不可荒宴醉酒；不可好色淫蕩；不可紛爭嫉妒。 \end{tabularx} \\ \\ \relax
13:14 & \begin{tabularx}{0.7\textwidth}{X} 總要披戴主耶穌基督，不要只顧滿足肉體，去放縱私慾。 \end{tabularx} \\ \\
[1ex]
\hline
\hline
\end{longtable}
馬太福音 24:36-44
\newline
\begin{longtable}{cl}
\hline
\hline
章節 & 經文 (和合本修訂版)\\
\hline
24:36 & \begin{tabularx}{0.7\textwidth}{X} 「但那日子，那時辰，沒有人知道，連天上的天使也不知道，子也不知道，惟有父知道。 \end{tabularx} \\ \\ \relax
24:37 & \begin{tabularx}{0.7\textwidth}{X} 挪亞的日子怎樣，人子來臨也要怎樣。 \end{tabularx} \\ \\ \relax
24:38 & \begin{tabularx}{0.7\textwidth}{X} 在洪水以前的那些日子，人照常吃喝嫁娶，直到挪亞進方舟的那日， \end{tabularx} \\ \\ \relax
24:39 & \begin{tabularx}{0.7\textwidth}{X} 不知不覺洪水來了，把他們全都沖去。人子來臨也要這樣。 \end{tabularx} \\ \\ \relax
24:40 & \begin{tabularx}{0.7\textwidth}{X} 那時，兩個人在田裡，一個被接去，一個被撇下。 \end{tabularx} \\ \\ \relax
24:41 & \begin{tabularx}{0.7\textwidth}{X} 兩個女人推磨，一個被接去，一個被撇下。 \end{tabularx} \\ \\ \relax
24:42 & \begin{tabularx}{0.7\textwidth}{X} 所以，你們要警醒，因為不知道你們的主哪一天來到。 \end{tabularx} \\ \\ \relax
24:43 & \begin{tabularx}{0.7\textwidth}{X} 你們要知道，一家的主人若知道晚上甚麼時候有賊來，就必警醒，不讓賊挖穿房屋。 \end{tabularx} \\ \\ \relax
24:44 & \begin{tabularx}{0.7\textwidth}{X} 所以，你們也要預備，因為在你們想不到的時候，人子就來了。」 \end{tabularx} \\ \\
[1ex]
\hline
\hline
\end{longtable}
講題:化干戈為玉帛　Turn swords and spears into shovels and hoes

經文:以賽亞書2:1-5;詩篇122:1-9;羅馬書13:11-14;馬太福音24:36-44

講員:林豪恩先生

場合:香港中文大學崇基學院禮拜堂

日期:2016年11月27日

一.引言

你有沒有到過耶路撒冷呢?你有沒有打算去耶路撒冷呢?今天讀的舊約經課,以賽亞先知看見遠象:有一日,人們從世界各地,各國各族向著耶路撒冷前來.耶路撒冷有甚麼特別,吸引人們從四方八面慕名而來呢?

我沒有去過耶路撒冷,只能從文字中得知一些關於這處地方的資料.原來,耶路撒冷只是一處細小的地方,2006年城區面積為126平方公里(西貢區約136平方公里),人口約724,000人(新界東約803,000).人們遠赴耶路撒冷,所為何事呢?原來,遠赴耶路撒冷的人潮,從前也出現過.當時未有火車和汽車,更沒有飛機,但在那些年間,竟然有合共約200多萬人長逃跋涉,從歐洲各地出發往耶路撒冷去,所為何事呢?歷史有稱這行動為十字軍東征.

從古至今,耶路撒冷的確吸引著很多人從世界各地前往.世界想在耶路撒冷得到甚麼呢?耶路撒冷又能夠給予世界甚麼呢?當我們向耶路撒冷問這些問題的時候,也難免聯想到教會.教會在世界各地建立,世界想在教會得到甚麼呢?教會又能夠送給世界甚麼呢?

二.經文及信息

1．耶路撒冷的自我想像

(以賽亞書)2:1亞摩斯的兒子以賽亞得默示,論到猶大和耶路撒冷.

2:2末後的日子,耶和華殿的山必堅立,超乎諸山,高舉過於萬嶺；萬民都要流歸這山.

是甚麼令耶路撒冷「超乎諸山」,「高舉過於萬嶺」呢?是甚麼吸引「萬民流歸這山」呢?如果把這些問題向當時的以色列人發問,他們可能想起大衛及所羅門王朝當時那種四境來朝的盛況.按列王記上記載,當時來耶路撒冷的人有兩種目的,一種是來進貢:

4:21所羅門統治諸國,從大河到非利士地,直到埃及的邊界.所羅門在世的日子,這些國都向他進貢,服事他.

另一種與之相關的是來觀摩所羅門王朝的盛世:

10:4示巴女王看見所羅門一切的智慧,和他所建造的宮殿,

10:5席上的食物,坐著的群臣,侍立的僕人,他們的服裝,和他的司酒長,以及他在耶和華殿裡所獻的燔祭,就詫異得神不守舍.

列王記上的作者也記錄了下列的資料,讓讀者知道,所羅門的盛世有強大的軍力為後盾:

10:26所羅門聚集戰車騎兵；他有一千四百輛戰車,一萬二千名騎兵,安置在屯車城,在耶路撒冷的王那裡.

同時,也以源源不絕的進貢來支撐:

10:25他們各帶貢物,就是銀器,金器,衣服,兵器,香料,馬,騾子,每年都有一定的數量.

當時的人自然而然地把所羅門的盛世聯想到宗教,示巴女王也以為所羅門王的盛世就是耶和華的榮耀.

被富國強軍吸引眼球,古代的人如是,現在的人也如是.可惜的是,不少人以強盛為上帝,同時以上帝來達至強盛的手段.不過,正如強盛的所羅門王朝轉瞬即逝,以金碧輝煌來吸引人的教會也不能維持多久.

「幾年前,美國水晶宮大教堂破產事件引起社會關注.陷於財政危機與權利糾紛的水晶大教堂讓人對教會產生不佳的印象,甚至動搖教會在一些人心中的信任.但我們也不要忘記,想當年水晶宮大教堂如日中天之際,有多少人把富麗堂皇的水晶宮大教堂當作是上帝的榮耀?在水晶宮大教堂發展處於巔峰之際,同樣引來很多羨慕的眼光,很多牧師夢想有朝一日亦可以把自己的教會建成如水晶宮大教堂般雄偉壯麗.(楊彰興:反思水晶大教堂破產事件:建堂?植堂?殖堂?)

選Donald Trump,美國人可能期望他帶領國家再次強大,帶領人民飛黃騰達.然而,世界向耶路撒冷所要的,是強大之道嗎?是飛黃騰達之路嗎?世界需要甚麼呢?耶路撒冷可以給世界甚麼呢?

以賽亞先知看見,有一天,列國來耶路撒冷,不再是進貢,不再是觀摩耶路撒冷的奢華盛世.那麼,他們來尋找甚麼呢?

2．上帝對耶路撒冷的期許

(以賽亞書)2:3必有許多國的民前往,說:來吧,我們登耶和華的山,奔雅各上帝的殿.主必將他的道教訓我們；我們也要行他的路.因為訓誨必出於錫安；耶和華的言語必出於耶路撒冷.

2:4他必在列國中施行審判,為許多國民斷定是非.他們要將刀打成犁頭,把槍打成鐮刀.這國不舉刀攻擊那國；他們也不再學習戰事.

有一天,世界會發現,無論科技如何進步,武器如何厲害,始終解決不了從古至今的問題,那就是國與國之間的殺戮；人們會發現,就算有能力潛入深海,有辦法飛上太空,始終處理不到從古至今的難題,就是人與人之間的鬥爭.世界尋找一處地方,能夠「在列國中施行審判,為許多國民斷定是非」；人們盼望一個和平的國度.

香港有一隊令人懷念的樂隊,名叫Beyond,他們有一首受歡迎的歌名叫 AMANI,以非洲某地的語言發出呼喚:

AMANI (和平)NAKUPENDA(愛)

NAKUPENDA WE WE(我們愛你)

為甚麼發出如此的呼喚呢?

本來:

祂(和平)主宰世上一切

祂(和平)的歌唱出愛  祂(和平)的真理遍佈這地球

然而:

祂(和平)怎可一去不返

祂(和平)可否會感到烽煙掩蓋天空與未來

因此:

無助與冰凍的眼睛

流淚看天際帶悲憤

是控訴戰爭 到最後傷痛是兒童

為甚麼會如此呢?

權利與擁有的鬥爭

愚昧與偏見的爭鬥

若這裡戰爭　到最後怎會是和平

因此,歌者說:

我向世界呼叫

AMANI NAKUPENDA

NAKUPENDA WE WE

有一天,當「他們也不再學習戰事」的時候,當人們反而想學習「將刀打成犁頭,把槍打成鐮刀」的時候,人們可以往哪裡去得到指導呢?往製造軍火的企業求助嗎?往窮兵黷武的政府求助嗎?向好勇鬥狠的政客請教嗎?因為看見這幅人們尋求和平的遠象,先知向耶和華的選民發出呼籲:

(以賽亞書)2:5雅各家啊,來吧!我們在耶和華的光明中行走.

這不是對列國的呼籲,而是對選民的呼籲.或者說,這不是對世界的呼籲,而是對教會的呼籲.當世界黑暗的時候,先知呼籲選民在耶和華的光明中行走,不是獨善其身,而是做好準備,當世界意識到自己的需要時,能夠尋到光明.

當世界紛亂的時候,教會又可以怎樣呢?在歷史上,教會曾經是衝突及紛亂的製造者.然而,教會也可以如先知所說,協助人「將刀打成犁頭,把槍打成鐮刀」,止息戰亂,協助和解.

在南非還實行種族隔離政策的時候,衝突越來越劇烈,內戰隨時可能發生.幸運地,南非避過了有些非洲國家內戰和種族屠殺的悲劇.在這過程中,曾經獲得諾貝爾和平獎的南非大主教杜圖(Tutu)扮演了和平之子的角色.杜圖雖然極力反對白人政府的種族隔離政策,但他同時呼籲黑人不要對白人報復,他說:「在南非有些人犯下難以置信的暴行,我以最強烈的形容詞來描述這些行為:『禽獸』,『魔鬼』.但是,魔鬼行為卻不會令犯罪者變成魔鬼,一個人決不會完全失去人性,這種人性是每個人被創造時按照神的形象賦予的.」他四處奔走,努力連繫各方力量,致力消弭武裝衝突的危機,以和平的方法解決種族隔離問題,建構公平的社會制度.後來,種族隔離政策結束,杜圖從大主教職位退下來,成立「杜圖和平基金會」,繼續推動和平,搭建和解的橋樑,展開為南非止痛療傷的旅程.杜圖認為世界大戰之後那種以審判和懲處為主調的善後方式並未能帶來和平,他主張以和解為經歷衝突的社會療傷,他希望讓南非向世界宣示饒恕的力量.

3．安全與武裝的掙扎

以賽亞先知呼籲:「2:5雅各家啊,來吧!我們在耶和華的光明中行走.」

因為「2:3必有許多國的民前往……2:4他必在列國中施行審判,為許多國民斷定是非.他們要將刀打成犁頭,把槍打成鐮刀.這國不舉刀攻擊那國；他們也不再學習戰事.」

雅各家為甚麼做不到呢?上帝的選民在這事件上遇到甚麼困難呢?今天詩篇經課上有兩句,或許讓我們得到一點啟發:

122:6你們要為耶路撒冷求平安!耶路撒冷啊,愛你的人必然興旺!

122:7願你城中平安!願你宮內興旺!

第七節的「平安」即是和平  (peace),他們透過甚麼達到呢?就是建造城牆(wall城中)；興旺也可譯作安全(security),如何得到呢?就是透過建造堡壘(宮內towers).雅各家遇到的難題,可能是安全感的挑戰,是自我保護的需要.人們一向理解城牆與堡壘是一種最低限度的自我保護,三十年前有些中國知識份子卻有另一種理解:

當時,中央電視台播放過一套歷史片集,名叫《河殤》.在該片的描述下,萬里長城不是自古以來的輝煌成就,而是失敗和退縮的象徵,代表著封閉,保守,無能,怯弱,是巨大的悲劇紀念碑.

安全和平安是人們最嚮往的狀態之一,矛盾的是,為了得到安全和平安,人們只想到建城牆,築堡壘,甚至進行軍備競賽,耶和華的選民也不能幸免.一次又一次,他們築起城牆,；可惜一次又一次,他們的城牆都不能保護他們.作為上帝的選民,保護從何以來呢?詩人曾說:「有人靠車,有人靠馬,我們卻靠耶和華的名」(詩20:7).靠耶和華的名,就能夠生存下去嗎?還是要先下手為強,制服了敵人,才得安枕呢?這是耶華的選民自古以來的掙扎,也不知道這種掙扎到甚麼時候才能完結.這是否也是我們的掙扎呢?

書信經課(羅馬書)的作者又發出呼籲:

13:11你們曉得,現今就是該趁早睡醒的時候；因為我們得救,現今比初信的時候更近了.

13:12黑夜已深,白晝將近.我們就當脫去暗昧的行為,帶上光明的兵器.

使徒保羅呼籲信徒,是時候醒過來了.從甚麼狀況醒過來呢?

13:13行事為人要端正,好像行在白晝.不可荒宴醉酒；不可好色邪蕩；不可爭競嫉妒.

13:14總要披戴主耶穌基督,不要為肉體安排,去放縱私慾.

原來,保羅呼籲信徒從「暗昧的行為」醒過來,「暗昧的行為」就是「荒宴醉酒」,「好色邪蕩」,「爭競嫉妒」,「為肉體安排,去放縱私慾」.醒過來做甚麼呢?剛剛還說把兵器變為生產工具,現在保羅竟然吩咐信徒帶上兵器.這些兵器是甚麼呢?就是「行事為人要端正,好像行在白晝」,「總要披戴主耶穌基督」.看來,保羅吩咐信徒帶上的兵器,不是用來攻擊別人,甚至也不是用來防衛別人攻擊,而是用來修養自己的.正如有人說,革命,不是革別人的命,而是革自己的命.

使徒呼籲信徒帶上「光明的兵器」,呼應著先知呼籲耶和華的選民「在耶和華的光明中行走.」他們都盼望選民能夠示範給世界看如何和平地處理分歧,如何不動刀槍化解衝突；盼望信徒展範給世界看如何以耶穌基督的心為心,取代隱藏於內心的好戰和攻擊,化解自己爭競和嫉妒的怒火,從私慾的陷溺中拯救自己出來.

三.總結

或者有人嘲笑說,「化干戈為玉帛」這些狀況只在童話世界裡才有,那麼,就讓我們以一個童話故事作為結束:

經過森林,經過峭壁,再跨過這個沙漠,城堡就在望了.王子筋疲力盡,迷迷糊糊摔下馬來,昏了過去.醒來時,發現自己躺在帳幕裡,有一位女子一邊餵他喝藥,一邊說:「你是王子,是為了城堡裡的公主而來的吧?」

在休養期間,王子跟著牧羊女子上路,逐水草而居.幾個月後,王子恢復健康.牧羊少女問:「為甚麼你一定要去那個城堡呢?就是因為你是王子,你就要打敗巨龍,把城堡裡美麗的公主救出來娶為妻子嗎?」

王子回答:「我根本不知道城堡裡面的公主是甚麼樣子.」

牧羊女子問:「那是為了甚麼呢?」

王子答:「為了證明自己.有一次,父皇帶著我們去打獵.他拉弓一射,就把遠處一頭大黑熊射倒了.哥哥們為父皇歡呼,我卻高興不起來,因為我看到中箭的黑熊一邊流著血,一邊掙扎.哥哥們發現我退後,就笑我是膽小鬼.接著,父皇把弓箭交我,對我說:「告訴他們你不是膽小鬼,那裡有兩頭小熊,你把其中一隻殺了給大家看.」我拉開弓箭,瞄準其中一頭小熊,卻看到大黑熊流淚的雙眼.我知道我絕對不能在這時候心軟,最後,我丟下了弓箭,一邊哭,一邊跑掉了.從此以後,皇宮裡所有人都更瞧不起我,連父皇也不再與我說話.後來,我聽聞關於公主被囚禁的事情,於是決定去試試,我要告訴大家我已經不再是那個懦弱的小孩子了.」

王子說完之後,兩個人都沉默了.牧羊少女突然對王子說:「我覺得你很勇敢.欺負弱小的不叫勇敢,真正的勇敢是,即使自己受傷,也要保護弱小.你才是你兄弟中最勇敢的一個.」

過了一些日子,牧羊少女對王子說:「前面就是你要去的遠山,城堡就在山上,我只能陪你走到這裡.但是,你一定要離開我嗎?只因為我不是美麗的公主嗎?」

王子心裡好難過,為了證明自己,是使怎麼不願意,他都必需撇下牧羊少女,繼續前行.

王子剛離開,就聽到牧羊少女的慘叫.王子策馬回頭,找到淹淹一息的牧羊少女.牧羊少女用盡最後力氣說:「城堡裡的巨龍來警告你了.王子,你不要去好不好.」

王子心中燃起熊熊的怒火,大叫道:「我要殺了那龍!我要殺了那龍替你報仇!」

王子跳上馬背,直奔城堡.撞開城堡大門,躍過護城河,看到了那條龍--那條令她失去牧羊少女,讓他夜夜惡夢的巨龍.王子奮力殺龍,最後巨龍摔了下來,在地上痛苦翻滾.正當王子準備殺死巨龍的時候,卻發現巨龍淚水汪汪的雙眼望著城堡頂端.王子抬頭一看,那裡有三條小龍,無助地望著這邊的悲劇.

王子好像明白了甚麼.接著,他收起寶劍,牽過馬兒,調頭就要走了.

「你不打算殺我嗎?你不是為了殺死我,拯救公主而來的嗎?」氣若游絲的巨龍竟然開口說.

王子搖搖頭說:「本來是這樣,但現在那都不重要了.」

巨龍對王子說:「我們原本居住在山邊的湖泊裡,可是你們破壞了我們的家園,我們一路躲,好不容易找到這個人跡罕至的地方,你們又為了同樣愚蠢的傳言而趕盡殺絕.你知道嗎?湖泊裡根本沒有黃金,城堡裡也沒有公主.」

王子說不出話來.他以為自己會仇恨眼前這條帶走牧羊少女的巨龍,但他沒有,他在回想為了證明自己,一路上砍掉了多少樹木,毀壞了多少動物的家園,他失去了牧羊少女,也失去了自己.

王子拿出藥來替巨龍治療,就好像當天牧羊少女照顧他一樣.待巨龍的傷勢康復,他才放心離開.

後來,王子成為了受人愛戴的君王,在他的管治下,沙漠化了的山腳終於重新湧出泉水,為牧羊少女長眼的地方,灌溉出一望無際的芬芳.(陳彥廷:遠山)

故事說完了,講道也結束了,但願經文帶給我們的思考延續下去.



\newpage

\section{我們等待的那個新世界}
\label{sec:344}
\textbf{鄧瑞強博士 - 我們等待的那個新世界}
\newline
\newline
\hyperref[sec:343]{< < < PREV SERMON < < <}
~
\hyperref[sec:toc]{[返目錄]}
~
\hyperref[sec:345]{> > > NEXT SERMON > > >}


以賽亞書 11:1-10
\newline
\begin{longtable}{cl}
\hline
\hline
章節 & 經文 (和合本修訂版)\\
\hline
11:1 & \begin{tabularx}{0.7\textwidth}{X} 從耶西的殘幹必長出嫩枝，他的根所抽的枝子必結果實。 \end{tabularx} \\ \\ \relax
11:2 & \begin{tabularx}{0.7\textwidth}{X} 耶和華的靈必住在他身上，就是智慧和聰明的靈，謀略和能力的靈，知識和敬畏耶和華的靈。 \end{tabularx} \\ \\ \relax
11:3 & \begin{tabularx}{0.7\textwidth}{X} 他必以敬畏耶和華為樂；行審判不憑眼見，斷是非也不憑耳聞； \end{tabularx} \\ \\ \relax
11:4 & \begin{tabularx}{0.7\textwidth}{X} 卻要以公義審判貧寒人，以正直判斷地上的困苦人，以口中的棍擊打全地，以嘴裡的氣殺戮惡人。 \end{tabularx} \\ \\ \relax
11:5 & \begin{tabularx}{0.7\textwidth}{X} 公義必當他的腰帶，信實必作他脅下的帶子。 \end{tabularx} \\ \\ \relax
11:6 & \begin{tabularx}{0.7\textwidth}{X} 野狼必與小綿羊同住，豹子與小山羊同臥；少壯獅子、牛犢和肥畜同群；孩童要牽引牠們。 \end{tabularx} \\ \\ \relax
11:7 & \begin{tabularx}{0.7\textwidth}{X} 牛必與熊同食，牛犢與小熊同臥；獅子與牛一樣吃草。 \end{tabularx} \\ \\ \relax
11:8 & \begin{tabularx}{0.7\textwidth}{X} 吃奶的嬰孩在虺蛇的洞口玩耍，斷奶的幼兒必按手在毒蛇的穴上。 \end{tabularx} \\ \\ \relax
11:9 & \begin{tabularx}{0.7\textwidth}{X} 在我聖山各處，牠們都不傷人，不害物；因為認識耶和華的知識要遍滿全地，好像水充滿海洋一般。 \end{tabularx} \\ \\ \relax
11:10 & \begin{tabularx}{0.7\textwidth}{X} 到那日，耶西的根立作萬民的大旗；列國的人必尋求他，他安歇之所大有榮耀。 \end{tabularx} \\ \\
[1ex]
\hline
\hline
\end{longtable}
講題:我們等待的那個新世界　The New World We are Waiting For

經文:以賽亞書11:1-10

講員:鄧瑞強博士

場合:香港中文大學崇基學院禮拜堂

日期:2016年12月4日

各位弟兄姊妹,早安.

由於今日有洗禮,我的講道會短一點.

今日是「將臨期」第二主日.「將臨期」有雙重意義,一方面,是等待「聖誕節」的來臨,這幫助我們回顧耶穌昔日在伯利恆的降生.另一方面,是期待耶穌的「再來」,這幫助我們對那個有別於現世的「新天新地」保持希望.耶穌第一次降臨時,是卑微的人,生在卑微的馬槽裡.耶穌第二次降臨時,將帶著神聖的榮耀,將萬物帶向終末的圓滿.耶穌第一次降臨時,像一粒種子,埋在塵世,毫不起眼,卻在悠長歲月中默默成長.耶穌第二次降臨時,就是這種子長成巨樹,結出纍纍的生命果實的時刻.如今,我們活在這巨樹一點一滴的成長過程中,或許,我們會窒礙它的成長；或許,我們只是無所事事的旁觀者；或許,我們是努力的灌溉人.

想到耶穌的再來,總讓我們期待那個由耶穌帶來的新世界.想到那個有一天會來臨的新世界,就看到我們現今這個世界的不濟.人性依舊醜陋,公義依舊不彰,和平依舊遙遙無期.

對於終末的新世界,古代的先知曾經作出預言.讓我們看看先知以賽亞的預言:以賽亞書11:1－10.

賽11:1　從耶西的本必發一條；從他更新的枝子必結果實.

賽11:2　耶和華的靈必住在他身上,就是使他有智慧和聰明的靈,謀略和能力的靈,知識和敬畏耶和華的靈.

賽11:3　他必以敬畏耶和華為樂；行審判不憑眼見,斷是非也不憑耳聞；

賽11:4　卻要以公義審判貧窮人,以正直判斷世上的謙卑人,以口中的杖擊打世界,以嘴裡的氣殺戮惡人.

賽11:5　公義必當他的腰帶；信實必當他脅下的帶子.

賽11:6　豺狼必與綿羊羔同居,豹子與山羊羔同臥；少壯獅子與牛犢並肥畜同群；小孩子要牽引牠們.

賽11:7　牛必與熊同食；牛犢必與小熊同臥；獅子必吃草,與牛一樣.

賽11:8　吃奶的孩子必玩耍在虺蛇的洞口；斷奶的嬰兒必按手在毒蛇的穴上.

賽11:9　在我聖山的遍處,這一切都不傷人,不害物；因為認識耶和華的知識要充滿遍地,好像水充滿洋海一般.

賽11:10　到那日,耶西的根立作萬民的大旗；外邦人必尋求他,他安息之所大有榮耀.

先知講出一個遠象,一個我們要努力建造的世界前程.先知的話可分為兩部分,頭部分講公義的遠象,第二部分講和平的遠象.

先講公義的遠象.

賽11:1　從耶西的本必發一條；從他更新的枝子必結果實.

耶西,是大衛的父親,他的後裔開展了一脈帝系.終歸,他們失敗了.這些帝王沒有將神的公義帶進人間,沒有造成永久的和平.從耶西這傾覆的殘幹中,卻生出一條柔弱的枝條.這弱不禁風的枝條,將改變這個世界.

我們知道,這枝條就是主耶穌,他作為耶西的苗裔,他的生命默默地改造著這個世界.耶穌第一次來世時,他展開他的公義工程.他再來時,就是公義圓滿實現之時.

卑微的耶穌,以及耶穌卑微的跟從者,他們如何在世上成就公義呢?

賽11:2　耶和華的靈必住在他身上,就是使他有智慧和聰明的靈,謀略和能力的靈,知識和敬畏耶和華的靈.

要成就公義,必須依靠神的靈.有神的靈在,才能擺脫人性的自私和卑污.神的靈讓人有「智慧和聰明」,這是能辨明正義與不義,能按正義規劃出前景,這是正義之「智」.神的靈同時讓人有「謀略和能力」,這就是有才幹能力實現正義而美善的計劃,這是正義之「勇」.有「智」而無「勇」,有遠象而無執行力,流於「吹水」.有「勇」而無「智」,有實力而無方向,則流於「亂來」.有「智」有「勇」,還要加上「仁心」.神的靈也讓人有「知識和敬畏耶和華的心」,這是領導者的靈性內涵,這是正義之「仁」.「敬畏耶和華的心」,不斷校正「智」與「勇」的方向,也不斷將力量貫注於「智」與「勇」之中.

賽11:3　他必以敬畏耶和華為樂；行審判不憑眼見,斷是非也不憑耳聞；

賽11:4　卻要以公義審判貧窮人,以正直判斷世上的謙卑人,以口中的杖擊打世界,以嘴裡的氣殺戮惡人.

敬畏神的王者,他們審判時,不應單憑表面的觀察,不應道聽塗說,更不應以虛無飄渺的理由去屈枉正直.為窮人伸張實質正義,在貧富懸殊的社會,這要多加留意.為「謙卑人」(即:被社會踐踏的人)制訂實質有益的措施,在金權世界裡,管治者要如何管束自己,才能做到這點?正義的王者,他口中的真理,會指正這世界的邪惡.他正義的判決,要剪除惡人.在正義的世界,窮人與弱者會被建立,惡人會得到應有的懲罰.

這根本就是香港人夢寐以求的世界.惡人被消滅,被踐踏的人被扶起,窮人不被欺凌.暫時,我們與這世界,仍有一段距離.展望將來,也不知何年何日,世界能變成這樣.我們相信,耶穌來過這世界,他自己就是正義的種子,種在人間,抽出幼弱的苗,默默地成長.我們相信,有朝一日,耶穌會再來,世界終會實現圓滿正義的.讓我們不要灰心,懷著信心,堅持盼望,作我們當作的事.終有一天,正義世界是會來臨的.

正義與和平,總是攜手並進的.沒有正義,則無法談論和平.同樣,沒有和平,則無法構想正義.跟著,我們看看先知的和平遠象.

賽11:6　豺狼必與綿羊羔同居,豹子與山羊羔同臥；少壯獅子與牛犢並肥畜同群；小孩子要牽引牠們.

賽11:7　牛必與熊同食；牛犢必與小熊同臥；獅子必吃草,與牛一樣.

賽11:8　吃奶的孩子必玩耍在虺蛇的洞口；斷奶的嬰兒必按手在毒蛇的穴上.

這是一個不可想像的和平世界.先知的和平想像,不再停留在人類社會的層面,而擴大到整個生態的層面.不單是人與人,更是獸與獸之間的和平.不單涉及家畜,更涉及猛獸.這是全面的和平.人與人之間的和平,必須擴大到人與萬物之間的和平.一個人若能愛護小小的生命,他對他人的接納程度也會高一點.豐子愷有一幅出名的漫畫叫《螞蟻搬家》,畫面上,一隊螞蟻在前進,孩童在螞蟻行進的路線上放著矮凳,造成隧道效果,讓螞蟻在下面安然前進.對一隻蟻也能如此,則容易留空間給他人,造成和平.

豺狼與無助的小綿羊和平共處,這在香港是不能想像的.獅子與牛一起,牛又和熊一起,牠們竟能和平共處.吃奶的孩子,是最無知無助的,他們和最毒的毒蛇玩耍,這是匪夷所思的.人世間一切弱肉強食和暴力對抗的場面都消失了,最強大的不再欺凌最弱少的,各種生命一起玩耍,一起享受生命的樂趣,這是我們夢寐以求的和平.

在先知的和平想像裡,獅子都變成吃草動物,這是難以想像的.若獅子都能變成素食者,則人類變成素食者也不是不可能的.人能素食,不單善待動物,也能促進糧食的正義,更能保護生態.我現時不是素食者,但心裡向著這個方向前進.人能愛動物,是善待弱者.善待弱者,是和平的開端.

要創造和平,需要全方位的努力,打破兩極對立競爭的思考方式.做豺狼的,不一定要吃掉綿羊,兩者是可以同居的.贏在起跑線上的,也可以停下來,和跑得慢的,一起享受一起跑步的樂趣.各人多一點遊戲的意識,所謂遊戲,就是放下功利的計算,而享受一種純粹的生命情趣.

要創造和平,你必須首先堅持正義.有了正義,你應當努力實現和平.

正義與和平,是我們的夢想.我們相信基督,他藉十字架的愛,已將正義與和平的種子埋下.我們是等待基督再來的人,我們知道,這種子正發芽生長,茁壯成長.我們要努力,在不義的世界堅持正義,在憎恨的地方播種仁愛.

先知的遠象是:

賽11:10　到那日,耶西的根立作萬民的大旗；外邦人必尋求他,他安息之所大有榮耀.

總有一天,耶西的根苗(即:耶穌)會成為萬國跟隨的王旗.

但願:榮耀歸於聖父,聖子,聖靈.阿們.

(後話:本人將早期的崇基講章結集為書,書名《敢於跟隨主》,基道出版,有興趣者可購買.)



\newpage

\section{自覺我是誰}
\label{sec:345}
\textbf{許開明牧師 - 自覺:「我是誰?」}
\newline
\newline
\hyperref[sec:344]{< < < PREV SERMON < < <}
~
\hyperref[sec:toc]{[返目錄]}
~
\hyperref[sec:346]{> > > NEXT SERMON > > >}


講題:自覺:「我是誰?」　Self- Awareness : “ Who Am I ?”

講員:許開明牧師

場合:香港中文大學崇基學院禮拜堂

日期:2016年12月11日

引言

今日的題目是《自覺:『我是誰』》,人需要自覺到底我是誰.在基督論中,這是一個很重要的討論:耶穌是一出生,就自覺是神的兒子,是基督嗎?還是他慢慢地自覺,並發現這身份?按常理,一個人在年幼時,他是不會預知自己是做特首或皇帝,所以他是長大後,才逐漸有這種自覺.

傳統神學家強調耶穌的神性,認為他生而知之；新派神學家則強調耶穌的人性,認為耶穌是逐漸醒覺,到30歲才知,受洗時發現自己是「神子」.

答案就要問問耶穌!但無論如何,耶穌是自覺能力強的人,人貴乎自覺,自動自覺,自覺,自省,自悟,自動.你的自覺能力如何?耶穌自覺自己的能力,自覺是神之子／基督／夫子,自覺自己的使命,耶穌的自覺,使他成為了不起的人和神.

現代的術語:自覺意識

自覺,指自己有所認識而主動去做,自發,自動,自己感覺到,自己有所察覺等,這是一種意識.自覺即內在自我發現,外在創新的自我解放意識,是人類在自然進化中通過內外矛盾關係發展而來的基本屬性,是人的基本人格,亦是人一切實踐行為的本質規律,表現為對於人自我存在的必然維持,發展,人類自覺本質的維護與發展是自由的真實實現.

人的內在生命本身具有意識本體 (invisible)與物質本體 (visible)的矛盾,構成人的內在本質,其基本屬性就是有意識的維護,發展自我本體.在與外在物質世界的矛盾過程中發展自我本體.自覺意識亦是創造自我的基本規律.

在領導學上,一個領導者沒必要「無所不知」,但必須要懂得「自覺」,誠實瞭解自己的優缺點,並且承擔錯誤.當主管主動檢討自己,也能為團隊帶來強大的力量.作為一個人,特別是好的領導者,你不能沒有自覺.管理專家利普曼(Victor Lipman)在《Forbes》的專欄特別提到:「最有效能的管理者,是對於自己的優缺點具有現實的評估,包括他們對別人的影響,以及需要加強的不足.」 誠實瞭解自己擅長什麼,承擔錯誤,坦承自己沒有答案的「自覺能力」,在高度競爭的職場,看似與為了避免別人懷疑我們的能力,而必須時常表現出無所不知,無所不能的武裝背道而馳.

在團隊人際方面,有所自覺,可以贏得別人的信任,增加你的領導力.反之,試圖隱藏缺失,只會更加凸顯,甚至創造出缺乏誠信,剛愎的形象.當你承認你必須學習,你正為團隊做出良好示範,對組織的好處更大.更重要的是,建立團隊時,瞭解各方的優勢和弱勢,讓你在召募人才和分派任務上有更好的安排,對決策更有盤算.

自覺包括三個層面,第一是自覺,即認識和瞭解自己,我是誰?第二是自省,即改進自己,我的表現如何?第三是自悟,即思想如何彌補自己的不足,我應該如何?

你有時間嗎?如果經常答「我沒有時間」,簡單意思是:「我唔得閒,我無空間留給自己.」嚴重意思是:「我無法掌握我的時間,我對生命無能,我像死人.」大學生常陷入這情況:「我沒有空間留給自己,我已失去了自己,我沒有自己了,要忙讀書,活動,拍拖,交功課,找外快等.」

聖經中,耶穌談到不同的自覺

第一,身份的自覺,耶穌自覺彌賽亞,神的兒子.

第二,行為的自覺,在馬太福音十一章5節:「就是瞎子看見,瘸子行走,長大痲瘋的潔淨,聾子聽見,死人復活,窮人有福音傳給他們.」以上所說的都是耶穌自覺身份後所行的事.

第三,關係的自覺,耶穌從我是誰,引伸到別人是誰.耶穌若是基督,約翰便是基督的先鋒.人要常自覺我與他,我與你,我與人,我與物,我與神的關係.

自知:我是誰?你是誰?

在馬太福音十六章13至20節,耶穌問門徒以下兩個問題:

1「別人說我是誰?」,或可解作別人期望我是誰?門徒答道:屬於施洗約翰 (最末,最近的先知),以利亞(以色列王國的先知),耶利米(被擄時期的先知),或先知裡的一個.

2你們說我是誰?或者期望:我是誰?門徒答:「你是基督,上帝的兒子!」

那麼你是誰呢?我會答:我是我!你要問別人眼中的你是誰?你會否只活在別人眼中呢?別人眼中,大學畢業的我適合去教書；家人眼中大學畢業的我要去工作,千萬不要做傳道人,因為不能賺錢吃飯；自己眼中,大學畢業的我適合繼續去讀書；耶穌眼中,大學畢業的我適合去讀神學,做牧師.結果,我順從主,在神學院讀書.

我們要當心矛盾與掙扎,若太過顧及他人眼中的我,會失去自我；若太過自我,會成為難相處的人.我們該如何平衝?

各位,我們處於一個被支配的世界裡,我們會失去自主,自由,自我.

《我是誰》故事

古時候,有一位公差奉命押解犯人到案,犯人是和尚.不甘心淪為階下囚的和尚,一直在尋找逃跑的時機.他盡力與公差拉關係,百般討好,做出一副恭順合作的樣子.漸漸地,公差的戒備心鬆懈了,甚至晚上住宿時還與和尚同桌吃飯.

一天晚上,兩人投宿一家客棧,因押解的目的地馬上就到了,公差心裡非常高興,就與和尚開懷暢飲起來.和尚見有機可乘,內心狂喜不已,但仍不動聲色與公差划拳飲酒.酒過數十巡,不勝酒力的公差醉得一塌胡塗癱在床上,和尚乘機從公差身上模出鑰匙,打開了自己手上的鐐銬,再把鐐銬銬在公差的身上.仍難消心中忿恨,又找來一把鋒利的刀子,將公差的頭髮剃光,趁著夜色逃之夭夭.

第二天,公差醒來,看不到和尚,心慌了起來,不覺用手摸了摸腦袋,卻摸著一個光頭,心裡頓時鬆了口氣:「原來和尚在這裡!我是和尚!」接著他又檢查了隨身的衣物,盤纏,一切都原封不動.他又愣了半晌,自言自語:「和尚在,衣物,盤纏也都在,那麼,我呢?我到哪兒去了?」

這是一個博君一笑的荒謬笑話.可憐的公差連自己與和尚都分不清,自然難逃丟飯碗的厄運了.「我是誰?好人,壞人?善人,惡人?」這個看似簡單卻又令人無法明確回答的問題,相信你時常會遇到.我們可在現實生活中,發現有許多人像公差一般,不明瞭「自己」到底為何?換句話說,就是渾渾噩噩的活著,為了生活而生活,從來沒有仔細想:

1自己到底是什麼?

2有什麼長處?

3有什麼缺點?

你不了解自己的個性嗎?所以也就常出現許多尷尬的事情.比如人云亦云,沒有自己的主張和見解,安於被他人操縱.又如,盲目的追趕潮流,崇拜偶像,穿衣,說話,舉止都力圖模仿偶像,這些人,都陷入了失去自我的迷思中而不自知.

話又說回來,如果真讓你回答「我是誰」這個問題,相信你也不一定能很準確地回答.其實,這個問題牽涉你能否正確認識「自己」的存在.我們生活在一個紛擾複雜的社會,各式各樣的干擾模糊了我們的視聽,所以有「認知」和「實際」上的差異.

1內在認知:思想和意識.

潛意識中與人一較長短的競爭心,所以我們對自己內在真實的「我」的認知是既模糊又不正確的.

2外在認知:外觀,外在表現

對自己表面的認知,許多人都是很清楚的.如自己的身高,體重,美醜,結交了多少朋友,喜歡什麼,討厭什麼等.

更深層的認知,許多人可能就不太清楚了,如自己的能力,氣質,性格,優缺點,對人對事的態度等.

全面且正確地認識自己,是身處當今競爭激烈的社會中的你必須做一項功課.一個對自己的能力,特質都不清楚的人,是無法集中火力專注前進的,所以,你必得先認識自己,找到一個真實的「我」.這樣,當人們還在迷惘地面對「我是誰」這個問題時,你可以非常沉著而冷靜地回答:「我是我!」

我被支配?

我的時間被支配?因為有無數的Deadlines 控制著你.

我的工作或學業被支配?這樣會失去真正值得大的事,工作最大,讀書最大,賺錢最大,唱歌最大!每天早上8時45分上班,下午5時30分下班,每天,每週,每年有一定的模式,這些固定的模式,框框,使我們像四塊半裡的死人一樣.

我的感覺被支配?快感,快樂感,快快感到甜,咸,味,「咸到死,甜到痺,辣到飛起!」

我的情感,喜好被支配?電影電視要你很輕易喊,笑,憤怒,仇恨,這些是除韓星,男女神,小白臉,小鮮肉外,韓劇所用的技倆之一.

我的思想被支配?廣告,中央,社論,名人論,中環派,面環派,中間派,民主派,假民主派,保皇黨等,使我們無自由,無自主,無自我.

我是誰?我是機器,不斷工作；我是 傀儡,不斷被人控制；我是怪獸,失去人性和傷人；我是人家的偶像,要滿足人家的期望；我是大英雄,救世主,我不會軟弱,我是強壯的；我是被動的,我不是我的主人；我迷失,我不是我,我沒有我自己!每個人都要自覺:問自己:我是誰?

我有價值嗎?

人的自我價值,馬克思主義認為人的價值,就是指人對自己,他人乃至社會需要的滿足；人的價值包含兩個方面,其一是社會價值,其二是人的自我價值.

自我價值,在人之初是通過父母的接納,肯定,承認,讚美,表揚,鼓勵等方式逐漸建立起來的,其核心是自尊.當人的自我價值感很強的時候,人會表現出自我完善的欲望,表現出向上向善的本性；當自我價值為零的時候,人會啟動自我毀滅程式.人類進化出自我提升,保持較高自我價值的需要,較高的自我價值讓我們有勇氣去從事各項事務,去競爭,從而為自己獲得機會.

耶穌以具體的行動,顯示「我是有價值?」

「就是瞎子看見,瘸子行走,長大痲瘋的潔淨,聾子聽見,死人復活,窮人有福音傳給他們.」(太十一5)

我的限制／障礙

自覺高度有三種:

1現實高度,即現實的障礙,例如我真的跳不過10呎高欄.

2心理高度,即心理障礙,其實我跳得過4呎的,但我怕,所以不行.

3靈性高度,即靈性的障礙,超越的上帝高不可攀,不會理會我

《鯊魚和很多小魚》的實驗

A實驗人員將鯊魚和很多小魚放在一個大魚缸裡,在鯊魚和小魚之間,隔著一塊玻璃,鯊魚看到小魚後,出於本能地向小魚遊去 ,想吃小魚,可被玻璃給擋住了,鯊魚又發起進攻,如此幾次失敗後,鯊魚縮到了角落裡.

B後來實驗人員將玻璃悄悄地拿掉,小魚紛紛向鯊魚遊去,奇怪的事情發生了,鯊魚不僅沒吃小魚,反而往後退縮.此時,鯊魚並非喪失了吃小魚的能力,而是因為剛才的錯誤經驗,以為自己吃不了小魚,而頻頻退縮. 小魚也不怕鯊魚.

從心理學角度說,自我設限就是在自己的心裡面默認了一個「高度」,這個「心理高度」常常暗示自己:這麼多困難,我不可能做到的,也無法做到,成功機會幾乎是零.想成功那是不可能的.「心理高度」是人無法取得成就的重要原因之一.它是一塊巨石,頑石,在人生及事業成長道路上,阻礙著人阻礙著人們前進.

約翰自己在獄中,找兩位門徒問耶穌是基督嗎?靈性高度就是當人祈禱不獲應允,或自覺身世遭遇不好時,就產聖靈性的障礙,即無法超越的上帝高不可攀,不會理會我.

各位,你對此怎麼看呢?那麼,又應該怎樣來突破自我設限呢? 每個人確實有自己不願意揭露的弱點,有的人知道了卻要逃避,而作為一個真正想成功的朋友,你必須很勇敢的去面對自己的弱點,來改變弱點,在此過程中,你的內心感受是非常痛苦的.但是,只要你有信念,有成功欲望,就能堅持,就能夠挺住,那任何困難都將不是困難.

產生自我設限的原因

多次的失敗讓自己看不到成功的曙光,因此覺得自己真的是不行,完全喪失了鬥志.或是別人的失敗讓自己望而卻步,對自己很沒有信心,總覺得自己本來就不如人,「他們都沒有做到,自己又怎麼會做到呢?」這都是自卑心理在做怪.

每個人都會在不知不覺中給自己設限,當自己總是和別人比的時候,總是覺得自己不如人家,由於自卑心理作怪,比較容易出現自我設限,因此最好是少和別人去比較.要比較,就要看到別人的優點,來補充自己的缺點,將你身邊較為成功的人作為超越目標.還有要多和以前的自己比較.能夠做到比以前的自己有進步就好了.有句話不是這樣說嘛:每天進步一點點!

我在找什麼?

我在找什麼?追求什麼?

莫言現年60歲,是2012年諾貝爾文學獎得主,為中國人首開記錄第一人.

他曾經解釋過關於「我」這個字 .

有一天「我」 字丟了一撇,成了「找」字,為找回那一撇,「我」 問了很多人,那一撇代表什麼?商人說是金錢；政客說是權力；明星說是名氣；軍人說是榮譽；學生說是分數......最後:「生活」告訴「我」,那一撇是「健康和快樂」,沒有健康和快樂,甚麼都是浮雲一片,這是一種 「悟」

李怡:活出真我 (2016年11月25日崇基學院週會)

到了我這個年紀80,許多人都會回顧一生.一位長期在安寧病房照顧臨終病人的澳洲護士,在她寫的一本書中提到,絕大多數臨終病人說到人生第一憾事,就是後悔沒有勇氣忠於自己去追求夢想,最終發現自己一生活在別人的期望裡.

臨終人最常說的話是:「真希望自己當年有勇氣去做自己真正想做的事,而不是選擇大家期望自己去做的事.」

你們還遠未到回顧一生的時刻.但臨終病人的遺憾是他們年輕時的錯過.要日後無憾,正是要趁年輕時勇敢做想做的事.為甚麼沒有勇氣追夢?

最主要是內心的聲音未能壓倒恐懼.

愛因斯坦說:「把恐懼擺一邊吧!何不聽聽內心深處的吶喊?」因此,時刻聽從內心的吶喊,去勇敢追尋,才不會臨終悔疚枉度此生.

一首令你自覺的好詩

本文翻譯自俄羅斯的一首小詩:

「短」 或

「人生很短,到頭來都是一樣」.

當人太執著一切時,看完後不禁自問,得到什麼啟示?

一天很短,

短得來不及擁抱清晨,就已經手握黃昏!

一年很短,

短得來不及細品初春殷紅竇綠,就要打點素裹秋霜!

一生很短,

短得來不及享用美好年華,就已經身處遲暮!

人生總是經過得太快,領悟得太晚,所以我們要學會珍惜,珍惜人生路上的親情,友情,同事情,同學情 ,朋友情與戰友情.因為一旦擦身而過,也許永不邂逅!

20歲以後,「異鄉」與「家鄉」是一樣的,因為去到哪都能適應；

30歲以後,「晚上」與「白天」是一樣的,經常幾天沒睡,日夜顛倒；

40歲以後,「學歷高」與「學歷低」是一樣,因為學歷低的,反而賺錢賺得多；

50歲以後,「美」與「醜」是一樣,因為再美的,此時皺紋,黑斑,老態都出現了；

60歲以後,「官大」跟「官小」是一樣,因為他們都已退休在家,無官職,階級都一樣；

70歲以後,「房子大」跟「房子小」是一樣,因為老人的關節都退化走不動,所需的空間只一丁點；

80歲以後,「錢多」跟「錢少」是一樣,因為再花也花不了多少錢；

90歲以後,「男人」跟「女人」是一樣,因為已沒有人還能做那種事；100歲以後,「躺著」跟「起來」是一樣,因為即使能起來,也不知要做什麼；離開世界後,住的也是一樣,就是4塊半放在棺木裡或2乘2呎放骨灰位龕!

你我的一生就是如此,差異不大.看開,看懂,看透了,人生不過是爾爾.因此請千萬珍惜您已擁有的,並找回您所失去的,好好把握你餘下的.

祈禱

上主,教導我們常能自覺,自省,自悟.自覺我是誰?自覺我的價值?

自覺我的限制?自覺我的追求?但自覺太執著時,也懂得看開,看透.從自覺中, 得到更新,得到成長,得到有意義的人生!阿們.



\newpage

\section{聖誕默想}
\label{sec:346}
\textbf{許立中先生 - 聖誕默想}
\newline
\newline
\hyperref[sec:345]{< < < PREV SERMON < < <}
~
\hyperref[sec:toc]{[返目錄]}
~
\hyperref[sec:347]{> > > NEXT SERMON > > >}


講題:聖誕默想　A Christmas Reflection

講員:許立中先生

場合:香港中文大學崇基學院禮拜堂

日期:2016年12月18日

聖誕是一個很特別的日子:平日你跟人「講耶穌」,未必有人願意理睬；但到了聖誕,卻似乎人人都興高采烈.以前一到聖誕,很多辦公室的百頁簾就會掛滿聖誕卡；流行用電郵WhatsApp送電子賀卡之後,就很少再看見這個景象.前幾天收到一位歐洲朋友寄來的聖誕卡,裡面分享了一些生活的近況,還夾了一張她女兒的近照,挑起我一種久違了的感覺.

聖誕是個「普世歡騰」的日子；聖誕是個不知為何那麼高興的日子.當然小學生都知道聖誕是記念耶穌誕生.但大多數人高興,肯定不是為了這個原因.這幾年不少人埋怨聖誕天氣太暖,經濟不好,沒有節日氣氛!對於很多人來說,他們是追求一種感覺,一種氣氛.我們今天所認識的聖誕,基本上是商業主導.

當然,聖誕亦有不少商業以外的活動:譬如報佳音,派飯,捐贈寒衣等,都是很有意義的活動.但倘若聖誕節不僅是「星光熠熠耀保良」,「慈善星輝仁濟夜」或「歡樂滿東華」,那麼我們就必須認真地思考聖誕到底是一個怎樣的日子.

事實上,聖誕的確是一個思考,默想的季節.路2:19記載,在天使「報佳音」之後,「馬利亞卻將這一切的事存在心裡,反覆思想.」那麼聖誕對基督徒來說,到底有甚麼意義?

「上帝與我們同在」

在今天的新,舊約經課,我們讀到:

「…主自己要給你們一個兆頭,必有童女懷孕生子,給他起名叫以馬內利.」(賽7:14)

「…這一切的事成就,是要應驗主藉先知所說的話,說:必有童女懷孕生子,人要稱他的名為以馬內利.(以馬內利翻出來就是『上帝與我們同在』.)(太1:23)

「上帝與我們同在」,不單透過我們熟悉的大自然,不僅在「冥冥中」,亦不單只「當我觀看你指頭所造的天,並你所陳設的月亮星宿,便說:人算甚麼,你竟顧念他!世人算甚麼,你竟眷顧他!」「上帝與我們同在」,不單是令康德感到敬畏的遼闊星空和人裡面的道德良心.叫人感到不知所措的,是上帝選擇在肉身與我們同在!

香港回歸將近20年,似乎並沒有使我們對將來更有把握.面對目前的政治環境,不要說計劃人生,就算是前面幾年該怎麼走下去,我們也不肯定:買不買樓,該甚麼時候買?值不值讓它綁整生?抑或儲錢移民更化算?我們對將來充滿疑惑,懷疑自己是否有足夠的條件去掌握自己的生命.

而上帝的道,就是冒險進入一個這樣的現實世界,出生於一個無產階級的家庭,在一個失去主權的殖民地長大,在一個不知道將來的地方去呼召人委身生命.

他在一個馬房出世,是因為約瑟和馬利亞付不起錢.客店沒有地方?你試試亮出你的美國運通白金卡!而馬槽亦並非一個我們想像那麼浪漫的地方,起碼不是我們在聖誕卡看到的.他們的處境和遭遇,只不過是當時一般窮人的待遇,毫不起眼,也沒有人會覺得怎麼特別.

從我們的角度來看,這樣低下和平凡的出生,似乎比較難聯想到甚麼偉大的事物.事實上日後當部分聽他的人知道他的出身之後,都看不起他:「加利利還有甚麼好貨色?」其實窮也不是問題,不少偉大的人物都有卑微的出身,但耶穌的出生是那麼普通,平凡,不起眼,就好像沒發生過甚麼事那樣.

「出世」這事實,本身就令人不安:出世,從來沒有徵詢過當事人的同意.更多時候我們聽到的,是意外懷孕；不少人的出生,甚至是出乎父母意料之外.避孕,亦不過是這幾十年的事.出世,代表沒有選擇!

但在我們這個時代,選擇幾乎被視為不可褫奪的天賦人權:自出娘胎,我們就習慣了有選擇.學校不好,你可以轉校；選錯學科,你可以轉系；待遇欠佳,你可以轉工；有更理想的對象,你可以「換畫」；忍受不了配偶,你可以離婚；有條件的話,你可以移民；活得不耐煩,你可以一死了之.聽說有些地方,甚至可以讓孩子選擇自己的「性別」.

我們忘記了,偶發性和無選擇性是現實世界的特徵.我們生命的起點,在一個很大的程度上決定了我們的命運,而這起點是完全不受我們的操縱!翁倩玉唱「阿信的故事」:「命運是對手,永不低頭!」但事實是,我們只能在既定的環境和條件下,作出我們有限的選擇.

學校不好,你可以轉校,但並不保證出路一定比不轉校好；選錯學科,你可以轉系,卻不代那就是你真正的方向；你當然可以轉工,但無論去到哪個位,都只是你頂頭上司的「阿四」.你換幾多次畫,才發現茫茫人海,所有的選擇其實都不過是偶遇?離了婚,才發現忍受不了的其實是你自己?移民或許會改善環境,卻犧牲了你一直所建立的.你灑脫地自我了斷,世界卻從來不在意你曾經存在.而無論你選擇甚麼性別,都無辦法擺脫感情的煩惱!

這就是「出世」的意思,而未被徵詢,無須同意,是做人一個基本的事實:Aleppo.道成肉身的上帝,並沒有選擇免疫於「出世」以及隨之而來的那一大串限制.

我們可以想像這位嬰孩是如何榮耀,寶貴,但他的出生,無論如何並不算得光彩.值得注意的,是透過跟我們一樣,經歷出生的過程,他是完全接納和認同生命一個基本的事實,就是有開始,有終結；有出生,有死亡.透過他的出生,他是真正背負了我們的命運.

甚麼是真理?

但單單「耶穌是我知心友」,並不能解開我們生命的困惑.單單感到被了解,亦不能解決我們對存在的焦慮.我們要問,這個成為肉身的道,這位與我們同在的上帝,究竟是那一位?為何祂在世上的代言人,往往帶來更多失望和困惑?

舊約聖經第一卷的第一章第一節:「起初,上帝創造天地.」

新約聖經第一卷的第一章第一節:「太初有道,道與上帝同在.道,就是上帝.」(約1:1)

上帝以祂的聖言創造萬有,萬有亦彰顯上帝的永能和神性,叫人無可推諉:「諸天述說上帝的榮耀；穹蒼傳揚祂的手段.這日到那日發出言語；這夜到那夜傳出知識.」(詩19:1-2)弔詭的是,縱然「它的量帶通遍天下,它的言語傳到地極」(4a),上帝本身卻是沉默的:「無言無語,也無聲音可聽.」(3)這無聲之言,不可道之道,留給人無限詮釋的空間.

作為「正統信仰」的捍衛者,教會當然期望上帝說得愈精確愈好,免生爭議.但從「生有涯,知無涯」的角度來說,信仰卻不可能是一套密不透風的理論系統.事實上就算有這麼一個周全的系統,恐怕亦非人所能處理.或許就如Jack Nicholson在“A Few Good Men”裡面對Tom Cruise的一句對白:“You want the truth? You cannot handle the truth!”

事實上嚴格來說,就是在舊約時代,上帝都是透過具體的歷史事件去向人表明祂的心意；往後衍生的書卷,經典,無論是先知書抑或歷史書,都只不過是對已發生的「啟示事件」的記錄和詮釋.到了新約時代,歷史更明確地見證上帝的決定:我們並不需要一套更審慎周密的經典誡律；我們需要的,是一位可以跟從的主.

這正是聖誕的信息:「大哉敬虔的奧祕,無人不以為然,就是上帝在肉身中顯現.」(提前3:16)事實上「從來沒有人看見上帝,只有在父懷裡的獨生子將他表明出來.」(約1:18)

只不過就算道最後成為肉身,坦白說,也無法完全解決人對生命的困惑.畢竟生命是不確定的,而每個人都必須,並且只能活出自己的生命.基督為我們開啟了一條新的道路,但怎樣去行,卻沒有人可以代勞.

道成了肉身,永恆者進入時空之中,是要叫在時空底下的人,可以在肉身中體現永恆.用福音書的說法,就是「成為上帝的兒女」(約1:12).起碼道理上是這樣說.

問題是,為甚麼沒有宗教信仰的人,總覺得這與他們無關?為甚麼他們總認為,信仰是屬於那些感情豐富,意志薄弱和頭腦簡單的人?宗教,狹義地來說,確實不是每一個人都需要；人生,才是所有人必須認真面對的場景.

不少曾經參加教會而最後決定離開的人,都苦澀地察覺到自己並不屬於這個群體.為甚麼?這個群體,有時確實令人感到左右為難.他們喜歡教堂的超脫和寧靜,讓他們在其中找到喘息的空間；可是基督教裡面附帶那麼多的道理規條,有時確實不容易理解.

他們甚至樂於接受十誡,八福,愛人如己這些道理,卻無法嚥下其他許多的東西.而最後 ─ 作為一個必須對自己誠實的人─面對人們捧著這個百孔千瘡的世界要求解釋,而上帝面對各樣的控訴卻最終保持沈默；當對愛心和公義的信念瀕臨破產,而主禱文裡面慈愛的天父可能只不過是個一廂情願的夢想…他們發覺自己不可能是一個基督徒:「我怎麼可能是基督徒?或許我只能同意,那些被賦予信心,能夠接納這一切的人是幸運的(the bliss of ignorance).可是我卻沒有被賦予這樣的信心.這與我無關.」

在信仰分享方面,教會似乎缺乏能夠跟一般人溝通的語言.我們往往從信仰的綱要開始,挑戰人「去到耶穌那裡」,重新檢視生命的優次,將生命交付在主的手中.但該怎樣信?信些甚麼?耶穌在一個怎樣的意義上成為人類的盼望與挑戰?我們卻往往支吾以對.我們假設一個人信了以後,一切問題就自然迎刃而解.但事實上不少人信了半輩子,除了一些耳熟能詳的陳腔濫調,卻始終弄不清楚自己到底信些甚麼.另一邊廂,教會卻以為他們只是「貪愛世界」,「不肯放下自己」.老實說,在這方面,教會一直是文不對題的.

該怎樣做好一個人?

在尋索生命意義的過程中,一個人可以從事很多社會,政治,宗教和關乎個人理想的活動,希望透過這些努力和活動,找到一些做人的端倪.但正如主耶穌曾經向那位少年的官指出,生命中真正重要和絕對不可少的,其實「只有一件」(參可10:21；路10:42).對於這「不可少的一件」到底是指些甚麼,不同的人可能會有不同的理解；但這些不同的理解,卻總是帶著一種「非如此不可」的催逼,意思是做了,就不枉此生；沒有做,就幾乎只是白走一趟.因此簡單來說,信仰要處理的,並不是怎樣做好一個信徒,而是該怎樣做好一個人.

這確實是一個非常重要的課題.甚至可以說,它大概是人們必須考慮信仰的唯一原因:生命中甚麼是絕對的,必須的,不可少的?甚麼只是次要的,有固然好沒有亦無妨?

生命核心價值的具體內容,或許會因應一個人的時空處境而改變 ─ 有時可能是國家興亡,民族尊嚴,有時可能是個人理想,甚至是兒女私情 ─ 但無論如何,沒有改變的,是那絕對肯定的決心.無論是印度人,中國人,英國人抑或丹麥人,好像我們熟悉的祈克果:「重要的是明白我自己,知道上帝要我做的,到底是甚麼；重要的是去找出於我真實的真理,找出一個我可以為之而生並為之而死的信念.」

對於生命的核心價值,或者說,在發現生命的絕對性的過程中,我們往往並不是從最終極的價值開始.在我們實際的經驗中,一個人往往並不是要先找到正確的信仰,他才能夠開始生活,開始做人；相反,只有當一個人認真地生活,認真地做人,他才會覺悟自己需要信仰.

事實上如果一個人不認真地生活,那麼他根本連自己要甚麼也不知道,信仰對於他亦沒有一個真實的意義.這或許便解釋了教會傳教工作失敗的原因.我們往往是過早將答案塞給那些根本未察覺到問題在那裡的人!認真地生活,認真地做人,往往是一個人找到信仰的先決條件!

在宗教上,我們無須論證上帝的存在,我們只需要指出,在人與現實相遇的不同範疇之中－認知方面,倫理方面,美學或社會政治方面－我們不難感受到內在於這些範疇的道德要求(moral imperative).我的意思是,在人生不同的範疇,無論是國家民族,知識探究,個人理想,都隱藏著一種不可否定,難以言說的絕對性,向我們發出要求和催逼.事實上若非這絕對性,我們就完全不能理解為甚麼有人會投身革命,人道救援或無國界醫生,或者好像劉曉波那樣,為了自由的理想而犧牲個人的自由；又或者單純為科學的研究或家族的責任而耗盡自己一生.

換句話說,這不可少的,並不是一個抽象的絕對,而必須透過具體的事物或活動表達出來.難以解釋清楚的,是它可以透過具體事物體現,卻並非具體事物本身.即是說,一個西藏僧人可以為了民族自由而引火自焚,另一個人亦可以 「捨己身叫人焚燒」,卻跟愛完全無關.

在主耶穌跟那少年官的相遇中,我們看見一切宗教上具體的要求,那個少年官誠實地供認他已經做足,並且從小就已經遵行:都已經做足了,還欠些甚麼呢?耶穌一針見血地指出問題並不在於他所做足的其他事,而在於他沒有做的一件事.有了它,其他多少也就沒甚麼相干；沒有它,做得再多也是白做.

每個人的具體關懷可以不一樣,那道德迫切性卻是一致的.它不是從我們有限的經驗推論出來,也不能被我們有限的腦袋所掌握.是這奧秘經驗叫人不期然地問:「為甚麼會有而不是沒有?」或者當人面對一個極端的處境,他會不期然問；「到底生命為何?」

對於那些人生基本的問題,除了那位降生為人,在一切事上跟我們一樣的耶穌之外,我們再找不到其他真正具體的答案.他的人生 - 也就是我們的人生 - 是有意義和方向的,因為上帝永恆的真道,無始無終的奧秘,親自降生成為平凡的人；透過信靠,平凡的人得以在生命中找到一個立足點.透過基督的降生,我們知道上帝並不是一個沉默的宇宙規律,而是來到我們的地方,回應我們面對的現實.

我想那些無神論者,虛無主義者,不可知論者有一件事是對的,就是從歷史的進程來看,他們看不見任何上帝的蹤跡,亦難以證明這個世界是被一「更高意志」所治理.但聖誕卻是從一個完全不同的觀點,去看我們跟世界的關係.我們再不是從宇宙的規律去理解渺小的浮生,而是從馬槽的嬰孩去看這世界的奧秘.

當我們看見耶穌的降生和他在地上的生活,就知道縱然有些事情或處境是我們無法解釋或理解,我們仍然可以相信,天父依然掌管.而在祂的掌管下,我們生活,動作,存留.

二千年前,明亮的晨星引領東方的智者長途跋涉朝見聖嬰,可是這嬰孩將要成就的事,他們卻一無所知；天使天軍向曠野的牧人宣告聖者的來臨,但對於外面仍然沈睡的世界,這夜卻是無聲.

二千年後,普天同慶救主降臨的規模遠超過往歷史；可是無言的上帝,卻仍然只是無聲地向願意聆聽的困乏心靈說話.Ω



\newpage

\section{葉公好龍話聖誕}
\label{sec:347}
\textbf{林豪恩先生 - 葉公好龍話聖誕}
\newline
\newline
\hyperref[sec:346]{< < < PREV SERMON < < <}
~
\hyperref[sec:toc]{[返目錄]}
~
\hyperref[sec:348]{> > > NEXT SERMON > > >}


以賽亞書 52:7-10
\newline
\begin{longtable}{cl}
\hline
\hline
章節 & 經文 (和合本修訂版)\\
\hline
52:7 & \begin{tabularx}{0.7\textwidth}{X} 在山上報佳音，傳平安，報好信息，傳揚救恩，那人的腳蹤何等佳美啊！他對錫安說：「你的神作王了！」 \end{tabularx} \\ \\ \relax
52:8 & \begin{tabularx}{0.7\textwidth}{X} 聽啊，你守望之人的聲音，他們揚聲一同歡唱；因為他們必親眼看見耶和華返回錫安。 \end{tabularx} \\ \\ \relax
52:9 & \begin{tabularx}{0.7\textwidth}{X} 耶路撒冷的廢墟啊，要出聲一同歡唱；因為耶和華安慰了他的百姓，救贖了耶路撒冷。 \end{tabularx} \\ \\ \relax
52:10 & \begin{tabularx}{0.7\textwidth}{X} 耶和華在萬國眼前露出聖臂，地的四極都要看見我們神的救恩。 \end{tabularx} \\ \\
[1ex]
\hline
\hline
\end{longtable}
希伯來書 1:1-4
\newline
\begin{longtable}{cl}
\hline
\hline
章節 & 經文 (和合本修訂版)\\
\hline
1:1 & \begin{tabularx}{0.7\textwidth}{X} 古時候，神藉著眾先知多次多方向列祖說話， \end{tabularx} \\ \\ \relax
1:2 & \begin{tabularx}{0.7\textwidth}{X} 末世，藉著他兒子向我們說話，又立他為承受萬有的，也藉著他創造宇宙。 \end{tabularx} \\ \\ \relax
1:3 & \begin{tabularx}{0.7\textwidth}{X} 他是神榮耀的光輝，是神本體的真像，常用他大能的命令托住萬有。他洗淨了人的罪，就坐在高天至大者的右邊。 \end{tabularx} \\ \\ \relax
1:4 & \begin{tabularx}{0.7\textwidth}{X} 他所承受的名比天使的名更尊貴，所以他遠比天使崇高。 \end{tabularx} \\ \\
[1ex]
\hline
\hline
\end{longtable}
約翰福音 1:1-14
\newline
\begin{longtable}{cl}
\hline
\hline
章節 & 經文 (和合本修訂版)\\
\hline
1:1 & \begin{tabularx}{0.7\textwidth}{X} 太初有道，道與神同在，道就是神。 \end{tabularx} \\ \\ \relax
1:2 & \begin{tabularx}{0.7\textwidth}{X} 這道太初與神同在。 \end{tabularx} \\ \\ \relax
1:3 & \begin{tabularx}{0.7\textwidth}{X} 萬物都是藉著他造的，沒有一樣不是藉著他造的。凡被造的， \end{tabularx} \\ \\ \relax
1:4 & \begin{tabularx}{0.7\textwidth}{X} 在他裡面有生命，這生命就是人的光。 \end{tabularx} \\ \\ \relax
1:5 & \begin{tabularx}{0.7\textwidth}{X} 光照在黑暗裡，黑暗卻沒有勝過光。 \end{tabularx} \\ \\ \relax
1:6 & \begin{tabularx}{0.7\textwidth}{X} 有一個人，是從神那裡差來的，名叫約翰。 \end{tabularx} \\ \\ \relax
1:7 & \begin{tabularx}{0.7\textwidth}{X} 這人來是為了作見證，是為那光作見證，要使眾人藉著他而信。 \end{tabularx} \\ \\ \relax
1:8 & \begin{tabularx}{0.7\textwidth}{X} 他不是那光，而是要為那光作見證。 \end{tabularx} \\ \\ \relax
1:9 & \begin{tabularx}{0.7\textwidth}{X} 那光是真光，來到世上，照亮所有的人。 \end{tabularx} \\ \\ \relax
1:10 & \begin{tabularx}{0.7\textwidth}{X} 他在世界，世界是藉著他造的，世界卻不認識他。 \end{tabularx} \\ \\ \relax
1:11 & \begin{tabularx}{0.7\textwidth}{X} 他來到自己的地方，自己的人並不接納他。 \end{tabularx} \\ \\ \relax
1:12 & \begin{tabularx}{0.7\textwidth}{X} 凡接納他的，就是信他名的人，他就賜他們權柄作神的兒女。 \end{tabularx} \\ \\ \relax
1:13 & \begin{tabularx}{0.7\textwidth}{X} 這些人不是從血生的，不是從情慾生的，也不是從人的意願生的，而是從神生的。 \end{tabularx} \\ \\ \relax
1:14 & \begin{tabularx}{0.7\textwidth}{X} 道成了肉身，住在我們中間，充充滿滿地有恩典有真理，我們也見過他的榮光，正是父獨一兒子的榮光。 \end{tabularx} \\ \\
[1ex]
\hline
\hline
\end{longtable}
講題:葉公好龍話聖誕　To Celebrate or not to Celebrate

經文:以賽亞書52:7-10;詩篇98;希伯來書1:1-4;約翰福音1:1-14

講員:林豪恩先生

場合:香港中文大學崇基學院禮拜堂

日期:2016年12月25日

一．引言

今天是聖誕節,在講聖誕故事之前,先講一個中國成語故事.

春秋時期,楚國有一個叫葉公的人,他常對別人說:「我特別喜歡龍,龍多麼神氣,多麼吉祥啊!」於是當他家裝修房子的時候,工匠們就幫他在房樑上,柱子上,門窗上,牆壁上到處都雕刻上龍,家裡就像龍宮一樣.就連葉公自己的衣服上也繡上了栩栩如生的龍.

葉公喜歡龍的消息傳到了天宮中真龍的耳朵里,真龍想:「沒想到人間還有一個這樣喜歡我的人呢!我得下去看看他.」有一天,龍從天上降下來,來到了葉公的家裡.龍把大大的頭伸進葉公家的窗戶,長長的尾巴拖在地上.葉公聽到有聲音,就走出臥室來看,這一看可不得了,一隻真龍正在那裡瞪著自己,葉公頓時嚇得臉色蒼白,渾身發抖,大叫一聲逃走了.(\url{https://read01.com/nEdz7n.html})

今天的幾段經課串連起來,彷彿述說著一個類似葉公好龍的聖誕故事.

二．釋經及應用

與葉公相差不遠的年代,在一處距離中國很遠的地方,有一位先知名叫以賽亞,他寫了下列一段文字,表達了當時那地方的人的渴望,讓我們聽聽他們望甚麼?

52:7在山上報佳音,傳平安,報好信息,傳揚救恩,那人的腳蹤何等佳美啊!他對錫安說:「你的上帝作王了!」

52:8聽啊,你守望之人的聲音,他們揚聲一同歡唱；因為他們必親眼看見耶和華返回錫安.

52:9耶路撒冷的廢墟啊,要出聲一同歡唱；因為耶和華安慰了他的百姓,救贖了耶路撒冷.

52:10耶和華在萬國眼前露出聖臂,地的四極都要看見我們上帝的救恩.

當時的人渴望甚麼呢?豈不是第八節所說的「耶和華返回錫安」嗎?他們期望耶和華返回鍚安做甚麼呢?豈不就是第七節所說「你的上帝作王了」嗎?當時的人喜歡上帝來臨到達甚麼程度呢?第七節及第八節都寫著「揚聲一同歡唱」.他們愛屋及烏,連那些傳報「耶和華上帝回來作王」這消息的人都因此而受歡迎.

又有詩人寫了一首詩歌,表達類似的渴望.有人說這是王帝登基時候用的詩歌；又有人認為這是寫於被擄時期表達希望回歸的詩歌.無論是那種背景,詩歌所表達的都是與上述以賽亞先知所寫的類似,就是歡迎耶和華來臨,歡迎他來作王.

第一節上半節發出讚美耶和華的呼籲:

98:1你們要向耶和華唱新歌!

第四至六節承接這呼籲,並仔細描述如何歡迎耶和華上帝:

98:4全地都要向耶和華歡呼,要揚聲,歡唱,歌頌!

98:5用琴歌頌耶和華,用琴和詩歌的聲音歌頌他!

98:6用號筒和角聲,在大君王耶和華面前歡呼!

被包含著的第一節下半節至第三節,提供了讚美耶和華的理由:

因為他行過奇妙的事,他的右手和聖臂施行救恩.

98:2耶和華顯明了他的救恩,在列國眼前顯出公義；

98:3記念他對以色列家的慈愛和信實.地的四極都看見我們上帝的救恩.

最後,詩人呼籲天地萬物都來讚美耶和華,歡迎祂來世界主持公義.

98:7願海和其中所充滿的澎湃,願世界和住在其間的發聲.

98:8願大水拍掌,願諸山在耶和華面前一同歡呼；

98:9因為他來要審判全地.他要按公義審判世界,按公正審判萬民.

這位在先知和詩人的詩句中被盼望的,受歡迎的,確實是與眾不同.他受歡迎,被追捧,是實至名歸的.希伯來書的作者對他的簡介其中幾句是這樣說的:

1:1古時候,上帝藉著眾先知多次多方向列祖說話,

1:2末世,藉著他兒子向我們說話,又立他為承受萬有的,也藉著他創造宇宙.

1:3他是上帝榮耀的光輝,是上帝本體的真像,常用他大能的命令托住萬有.他洗淨了人的罪,就坐在高天至大者的右邊.

1:4他所承受的名比天使的名更尊貴,所以他遠比天使崇高.

曾經有電影用「萬世巨星Jesus Christ Superstar」來形容這位聖誕節的主角.事實上,按他的履歷,的確受之無愧.希伯來書的作者指出,按身分來說,他是「上帝的兒子」,「是上帝的榮耀的光輝」,「上帝本體真像」；按工作經驗來說,他「創造宇宙」,並且「用他大能的命令托住萬有」；按資產來說,他是萬有的繼承者.如果要找對象來與他比較,幫助人們明白他的超越,作者選了天使,並且指出他「的名比天使的名更尊貴」,「所以他遠比天使崇高」.這位被描述「坐在高天至大者的右邊」的「萬世巨星」,原來和每一個人都有關係的,這也是作者希望讀者留意的,就是第三節所說「他洗淨了人的罪」.

有一天,這位萬眾歡迎,千古期待的終於來臨了.約翰對於那神聖的事件如此描述:

1:1太初有道,道與上帝同在,道就是上帝.

1:2這道太初與上帝同在.

1:3萬物都是藉著他造的,沒有一樣不是藉著他造的.凡被造的,

1:4在他裡面有生命,這生命就是人的光.

1:5光照在黑暗裡,黑暗卻沒有勝過光.

1:9那光是真光,來到世上,照亮所有的人.

1:10他在世界,世界是藉著他造的,世界卻不認識他.

1:11他來到自己的地方,自己的人並不接納他.

1:12凡接納他的,就是信他名的人,他就賜他們權柄作上帝的兒女.

1:13這些人不是從血生的,不是從情慾生的,也不是從人的意願生的,而是從上帝生的.

1:14道成了肉身,住在我們中間,充充滿滿地有恩典有真理,我們也見過他的榮光,正是父獨一兒子的榮光.

關於這位被期待者,約翰對他的介紹,基本上與先知,詩人及書信作者一致,一同見證他的神聖及超越,並且補充說他「充充滿滿地有恩典有真理」.然而,最奇怪的是,當這位被期待的聖者果然降臨人間的時候,卻沒有了在先知書及詩篇中的歡呼聲及讚美聲,取而代之的是第十節所說的「他在世界,世界卻不認識他」,第十一節說:「他來到自己的地方,自己的人並不接納他」.

這種被稱為「反高潮」的矛盾狀況,在人類社會中不乏例子.例如古代的成語故事中有葉公好龍；在現代的經濟事件中有幾年前廣被「炒作」的滬港通,開通前萬眾期待,甚至被「講股佬」預測每日額度在開市首個小時就會用盡.當正式開通後,交易量並不多,每日額度到現在為止,在絕大多數日子都是用不盡的.因此,有人戲言這列「港滬股市直通車」,其實是「幽靈列車」.

在個人的生活經驗中,也有類似的故事,甚至可能這就是我們自己的故事.有沒有聽聞有人自小就期望或被期望入大學?選擇幼稚園,為了入大學；選擇小學,為了入大學；選擇中學,更為了入大學.期望多年,等候多年的事情終於來臨.入讀大學了,原來卻不是那麼滿足和興奮,反而感到若有所失,心裡充滿疑惑.從小到大的期待終於來到了,不能不讀下去,但讀下去又不知所謂何事.一直期待的事情,出現後竟然成為「雞肋」,食之無味,棄之可惜.

錢鐘書有名句如此說:「圍在城里的人想逃出來,城外的人想衝進去,對婚姻也罷,職業也罷,人生的愿望大都如此.」網絡上有一詞彙名為「雞肋婚姻」,意思不需要解釋,讀者都能領會.令人感到迷糊的是,當幻想多年,期待以久的婚姻臨到的時候,何以卻變為「雞肋」?

今天是聖誕節,我們並不是講婚姻的故事,也不是講讀大學的故事,而是講聖誕的故事.會不會在人們心中,聖誕的主角耶穌基督,也變成了「雞肋」呢?甚至是人們對待祂,如葉公對待那條來訪的龍一樣呢?

我們當如何看待這位聖誕的主角耶穌基督呢?今年是宗教改革五百年的記念年,讓我們聽聽發起宗教改革的馬丁路德的說法,他在A brief introduction on what to look for and expect in the Gospels一文中如此說:

「你應該從兩個向度去認識耶穌基督的訓言,行動及他所承受的苦難.第一,這是一個展示給你的榜樣,讓你跟隨及倣效.正如彼得前書2:21所說:「你們蒙召原是為此；因基督也是為你們受過苦,給你們留下榜樣,叫你們跟隨他的腳縱行.」因此,當你知道耶穌如何祈禱,禁食,幫助人,以愛待人的時候,你也應該如此行.

然而,這只是福音的一小部份,或者甚至不能稱為福音.若只是這樣,耶穌基督對我們的幫助並沒有優於其他聖人,他的生命只仍是他自己的而對我們的生命亳無貢獻.簡而言之,單單以榜樣來理解耶穌基督,只製造偽君子,而不能令人成為基督徒.你應該以另一個更高的層次來認識耶穌基督,雖然這個層次甚少被宣講,但卻是最源遠流長及最美好的.福音的最根本就是,在你認識耶穌基督為榜樣之前,你承認及接受耶穌基督是一份禮物,一份上帝賜給你,也是屬於你的禮物.意思是,當你耳聞或目睹耶穌基督所成就的工作及作承受的苦難,你可以亳無疑問地接受這些工作及苦難都是為你及屬於你的,甚至是等於你所做及受的.看,這就是福音,上帝無比的慈愛並非先知,使徒或天使能夠充份展現的,也不是任何心靈能夠想像的.被上帝對我們如此熱切的愛所承托,天下的心靈和良知得到歡樂,安全和滿足.因此,這稱為福音,德文說這是喜樂,美好,欣慰的信息.」(頁119-120)

三．總結

約翰福音1:12凡接納他的,就是信他名的人,他就賜他們權柄作上帝的兒女.

在結束的時候,讓我們聆聽另一個故事,或者能夠給予我們一點啟發如何接待聖誕節的主角耶穌基督.

就在那一個我永遠不能忘懷的夜晚,我邀請他進入我的心.

我帶主耶穌參觀的第一個地方是書房--思想的圖書室!基督跟我入內,環視架上的書本,桌上的雜誌和牆上的圖畫.我轉身對他說:「主啊,我知道這需要更新變化,你可以幫助我使它歸正嗎?」「一定」他說:「我樂意幫忙.」在以後的歲月中,我發現當我的思想集中在基督身上時,會令我的邪情私慾全然消退.這樣,他幫助我得以抑制那些不良的思想.

我們從書房來到飯廳,我在此耗費了不少精力來滿足自己的胃口和欲望.我對主說:「這裡很寬敞舒服,我肯定你會喜歡我們的飯菜.」於是他跟我坐在飯桌旁,我看見他沒有下筷,便不解地問:「主啊,你不吃嗎?有什麼問題嗎?」「我有食物吃,是你不知道的,我的食物就是遵行那差我來者的旨意.」他又定睛看著我說:「你若想要真正叫你飽足的食物,就好好求天父的旨意,不是你自己的喜好,你自己的欲望,你自己的滿足；你要求天父的旨意,討他的喜悅,那食物就能叫你飽足.」

我們接著走進客廳,這地方十分恬靜舒服,我喜歡這裡.基督看來也很喜歡:「這地方確實討人喜悅,就讓我們常常來吧,我們可以一起團契相交.」初信的我感到喜出望外,我想:「有什麼事比與基督相處溝通更值得嚮往的呢?」他約定我:「我每天必定在這裡等你,你來與我相會.」這樣,早晨復早晨,我來到客廳,然後我們一起細讀聖經,他向我闡明當中的真理.不知從甚麼時候開始,我愈來愈多地失約,幾天不見也如常.有一天,我匆匆忙忙地趕著外出.經過客廳,瞥見主正坐在那裡.我感到好好意思:「他是客人,我請了他來,他要作我家之主,我卻這樣怠慢他.」他說:「我答應了你,我每天在這裡等著與你相會,不是因為甚麼,只因為我愛你,渴望與你團契相交.」

又有一天,我發現主站在大門前等我,神色有異.他對我說:「屋裡有一股怪味,發自一些朽壞腐死的東西.」當他開聲時,我已心知肚明.是的,在樓梯頂端的走廊有一間小屋.在門和鎖的背後,我藏了一兩件個人的私隱,不讓任何人知道.基督指著那間小屋「就在裡面!有些敗死的東西!」我不禁惱怒起來--我已經給他書房,飯廳,客廳,工作間,娛樂室和其他房間,連這一間儲物室也要干預嗎?他說:「我知道你沒有力量處置那些朽壞腐死的東西,請把鑰匙交給我,授權我去清理那地方吧.」我戰戰兢兢地把鑰匙交給他.主開門入內,把一切腐朽發臭的東西都丟棄淨盡.然後再把地方清潔一番,塗上油漆,整個房間頓時變得煥然一新.噢!能把那些腐朽發臭的東西從生命中除掉,真是何等釋放!

後來我有一個想法:「我一直為基督努力保持著心懷的清潔,清理一個又一個的房間.但惱人的是清理一個房間之後,不久另一間又污穢了.為保持清潔的心和順服的生命,真弄得我筋疲力盡,無力支撐!」因此,我斗膽問主:「主阿,有沒有可能,你來接管整間房子,替我清理它們,就像你清理那間小屋一樣?」主說:「一定!我來所要做的就是這樣.你不可能單用自己的力量做個得勝的基督徒,那是不可能的.讓我幫你一起做吧,那樣才成!」

主居我心,呈現出另一片不一樣的天空!(節錄自「主居我心」一文)



\newpage

\section{出死入生的兄弟}
\label{sec:348}
\textbf{鄧瑞強博士 - 出死入生的兄弟}
\newline
\newline
\hyperref[sec:347]{< < < PREV SERMON < < <}
~
\hyperref[sec:toc]{[返目錄]}
~
\hyperref[sec:349]{> > > NEXT SERMON > > >}


希伯來書 2:10-18
\newline
\begin{longtable}{cl}
\hline
\hline
章節 & 經文 (和合本修訂版)\\
\hline
2:10 & \begin{tabularx}{0.7\textwidth}{X} 原來那為萬物所屬、為萬物所本的，為要領許多兒子進入榮耀，使救他們的元帥因受苦難而得以完全，本是合宜的。 \end{tabularx} \\ \\ \relax
2:11 & \begin{tabularx}{0.7\textwidth}{X} 因那使人成聖的，和那些得以成聖的，都是出於一。為這緣故，他稱他們為弟兄也不以為恥， \end{tabularx} \\ \\ \relax
2:12 & \begin{tabularx}{0.7\textwidth}{X} 說：「我要將你的名傳給我的弟兄，在會眾中我要頌揚你。」 \end{tabularx} \\ \\ \relax
2:13 & \begin{tabularx}{0.7\textwidth}{X} 他又說：「我要依賴他。」他又說：「看哪！我與神所給我的兒女都在這裡。」 \end{tabularx} \\ \\ \relax
2:14 & \begin{tabularx}{0.7\textwidth}{X} 既然兒女同有血肉之軀，他也照樣親自成了血肉之軀，為能藉著死敗壞那掌管死權的，就是魔鬼， \end{tabularx} \\ \\ \relax
2:15 & \begin{tabularx}{0.7\textwidth}{X} 並要釋放那些一生因怕死而作奴隸的人。 \end{tabularx} \\ \\ \relax
2:16 & \begin{tabularx}{0.7\textwidth}{X} 誠然，他並沒有幫助天使，而是幫助了亞伯拉罕的後裔。 \end{tabularx} \\ \\ \relax
2:17 & \begin{tabularx}{0.7\textwidth}{X} 所以，他凡事應當與他的弟兄相同，為要在神的事上成為慈悲忠信的大祭司，為百姓的罪獻上贖罪祭。 \end{tabularx} \\ \\ \relax
2:18 & \begin{tabularx}{0.7\textwidth}{X} 既然他自己被試探而受苦，他能幫助被試探的人。 \end{tabularx} \\ \\
[1ex]
\hline
\hline
\end{longtable}
講題:出死入生的兄弟　Jesus, Our Brother

經文:希伯來書2章10-18

講員:鄧瑞強博士

場合:香港中文大學崇基學院禮拜堂

日期:2017年1月1日

各位弟兄姊妹,早安.

新的一年的開始,祝願大家有一新的心靈.舊年,有成有敗,有得有失,想想,成功的事,會使你走歪路嗎,抑或能讓你更敬畏神呢?失敗的事,會叫你心懷怨憤嗎,抑或能教你更謙厚待人呢?事情在神手裡,生命在神手裡,讓我們常常感恩.將過去一年,以感恩的心交給神.今年,我們以感恩的心開始.

今日的講道經文,是今日的書信經課:希伯來書2:10－18.

來2:10　原來那為萬物所屬,為萬物所本的,要領許多的兒子進榮耀裡去,使救他們的元帥,因受苦難得以完全,本是合宜的.

來2:11　因那使人成聖的和那些得以成聖的,都是出於一.所以,他稱他們為弟兄也不以為恥,

來2:12　說:我要將你的名傳與我的弟兄,在會中我要頌揚你；

來2:13　又說:我要倚賴他；又說:看哪,我與神所給我的兒女.

來2:14　兒女既同有血肉之體,他也照樣親自成了血肉之體,特要藉著死敗壞那掌死權的,就是魔鬼,

來2:15　並要釋放那些一生因怕死而為奴僕的人.

來2:16　他並不救拔天使,乃是救拔亞伯拉罕的後裔.

來2:17　所以,他凡事該與他的弟兄相同,為要在神的事上成為慈悲忠信的大祭司,為百姓的罪獻上挽回祭.

來2:18　他自己既然被試探而受苦,就能搭救被試探的人.

這段短短的經文,有點似一篇很短的神學論文,講解:我們有什麼問題,我們活在一個什麼樣的世界裡,耶穌是誰,他做了什麼,為何要這樣做,這樣做帶來什麼效果等等.

(1)耶穌是我們的兄弟

來2:10　原來那為萬物所屬,為萬物所本的,要領許多的兒子進榮耀裡去,使救他們的元帥,因受苦難得以完全,本是合宜的.

經文交代一個宏大的世界觀.「萬物是有所屬的」,或說「萬物的存在是有理由的」.跟著說,「萬物是有所本的」,或說「萬物的存在有一根源」.原來,萬物不是空空而來,不是空空而去.來,從一根源而來,萬物從神而來.去,也不是漫無目的的,生命存在有一目的.萬物來自神,也被神賦予目的.

神對人類的心意,是要領他的兒女(即我們)到榮耀裡去.此世總不圓滿,美善總有缺失,生命總有遺憾,但有朝一日,在神手中,一切會圓滿.誰引領我們進入圓滿裡呢?是神差派來的「元帥」,或說「領頭的人」,「開路的人」.當然,我們知道,這是指耶穌.耶穌做什麼?他「因受苦難得以完全」,或者說,他藉著他經歷的一切苦難,由生到死的種種艱難,他成了一個完全的「領頭的人」.他完全地完成神的使命,完全地展示神的性情.他已經排除萬難,引我們到榮耀裡去的路已打通.經文說,耶穌所做的,「是合宜的」.所謂「合宜」,就是耶穌將神的性情真實地流露出來,耶穌活現了神對人類的心意.

來2:11　因那使人成聖的和那些得以成聖的,都是出於一.所以,他稱他們為弟兄也不以為恥.

「那使人成聖的」指耶穌,「那些得以成聖的」指我們這些信徒.所謂成聖,即我們都歸屬了神.耶穌與我們,都屬於神,都出於一神,都是神的兒女.在這意義下,耶穌稱我們為兄弟,並不以為恥.若我們只是一班偷生怕死的人,毫無道義的人,而耶穌卻稱我們為兄弟而不以為恥,這真是天下間最大的情義.

將兩節經文合起來看,神是愛我們的,我們的人生是有意義的.神的心表明在耶穌生命中,他經歷種種艱難,為要打通到達榮耀的路.耶穌不單是「領頭的人」,更是我們的兄弟.

故此,請不要說人生無意義,人生是有意義的.

請不要說沒有明天,是有榮耀的明天的.

請不要說孤寂難奈,你是有義無反顧的兄弟的.

跟著,第12及13節引經據典去論證耶穌是我們的兄弟,這些不細說了.

(2)耶穌如何做兄弟

來2:14　兒女既同有血肉之體,他也照樣親自成了血肉之體,特要藉著死敗壞那掌死權的,就是魔鬼,

來2:15　並要釋放那些一生因怕死而為奴僕的人.

這裡,經文展示了人類的終極問題.

「一生因怕死而為奴僕」.肉身的死,固然是死.心靈的死,也是死.聖經說,罪的代價是死,故死也是罪的後果.香港人站在「量重磅」上一量,會說「死啦,又肥了」.「肥了」,對現代人來說,也是一種死.或者,照一照鏡,「死啦,又多了一條皺紋」.「皺紋」比死更可怕.多少人因怕肥怕皺紋而成為奴隸.現代神學用「虛無」(nothingness)這個字去展示死亡的內容.死亡是從「一切的有」瓦解為「絕對的無」.虛無,是一切關係的消失,一切意義的消失,一切存在痕跡的消失.「你」好像從無存在過一樣.

腰纏萬貫,或顯赫功業,一死,就意味這一切再與你無關了.或許,你與朋友愛人擁有刻骨銘心的關係,一死,就意味一切都煙消雲散了.引用楊絳《我們仨》的一句話:「我們三人就此失散了.就這麼輕易地失散了.」

為了抓緊這些被「虛無」驅散的東西,大家都知道,人類花了何等精力,何等心思.為了抓住健康,現代人做了多少檢查,食了多少藥,看了多少資料.為了自保,有些人或許會推身邊的人出去受死.就是因為怕握不住種種東西,我們的手常常緊緊握住,連睡覺的時候,也緊緊握住.怕放開了,就永遠捉不到了.

按這裡的經文說,死,或無意義,或虛無,是終極的敵人.魔鬼只是利用人這個弱點,這個恐懼,來讓我們更加陷入虛無裡.這有點似一個人初學踏單車,見到前面有條燈柱,心中恐懼起來,就因這恐懼,反失去控制的能力,最後便撞在燈柱上.魔鬼的工作,只是強化我們的恐懼,然後癱瘓我們克服恐懼的能力.

在這困局中,我們的兄弟耶穌來了.

耶穌是如何做兄弟的?

在一個人人嚷著「你不代表我」的社會裡,不是你說做兄弟,人家便讓你做兄弟的.再沒有「大台」,你不代表我,也不是我兄弟.耶穌是如何做兄弟的?

來2:14　兒女既同有血肉之體,他也照樣親自成了血肉之體,特要藉著死敗壞那掌死權的,就是魔鬼.

人類的「必死性」,來自我們的血肉之體.這是我們所恐懼的死亡的人性根據.耶穌不是認同我們的政見,不是認同我們的做人方式,耶穌認同我們的恐懼,在形成恐懼的人性條件上認同我們,與我們同一,他也照樣親自成了血肉之體.

然後,他面向用死亡和虛無來威脅我們的魔鬼,就在魔鬼面前,大膽地死去；就在魔鬼面前,向魔鬼宣告,他利用死亡向人作出的威脅終歸是沒有威脅力的.

在《伊索寓言》裡,有個寓言叫＜給貓繫鈴＞(Belling the Cat).話說,老鼠開大會,商討用什麼辦法應付最大的敵人－－貓.有隻小老鼠說:「最好的方法,是在貓頸上繫個鈴,這樣,貓來時,聽到鈴聲響,我們便能及時逃命了.」問題是:誰去繫這個鈴呢?讓我擴充一下這個故事.可以想像,老鼠日夜在貓的陰影下過活.有時夜裡夢見貓,便會驚醒.很多老鼠患上精神病,如:憂慮,驚恐症,甚至出現幻覺,常常見到貓似的.有次大老鼠和小老鼠一起出外覓食時,遇到貓,大老鼠將小老鼠迫向後,貓抓住小老鼠而大老鼠得以逃命.沒有老鼠膽敢為貓繫個鈴.這是超出了老鼠的想像,也超出了老鼠的鼠性限制.若耶穌要進入這個故事,解救老鼠,他就要成為一隻老鼠,成為老鼠兄弟.

問題來了,耶穌為何一定要成為老鼠兄弟呢?他用人的角色進入這故事,或以一隻凶狠的狗的角色進入這故事,不是更好嗎?他以人的角色或狗的角色進入故事,將貓綁住或驅趕那貓,不是更好地解決問題嗎?若問題在貓,則解決了貓,當然解決了問題.但問題其實不在貓,而在老鼠本身,在老鼠對死亡的深層次恐懼.貓只是一種外在因素,強化了老鼠的恐懼.沒有貓,這種恐懼仍在.

故此,若耶穌變成人或狗,只將貓解決掉,則這些老鼠無法改變,永遠只是渴望不會再有外在條件激發恐懼,永遠只能是深藏恐懼卻也無法克服恐懼的老鼠,永遠只是隻「未重生」的老鼠.

耶穌一定要成為老鼠,因為他的重點不在貓,在老鼠.他要在恐懼的老鼠中成為他們的兄弟,分享他們的血肉之體,分享他們的恐懼.然後,他做那隻甘心去為貓繫鈴的老鼠.問題不在如何制服貓,問題在如何成為一隻真正的老鼠,一隻進入「榮耀」的老鼠,一隻能帶領其他老鼠同樣進入「榮耀」生命的「領頭老鼠」.耶穌去為貓繫鈴,這是必死的任務.事實上,他死了.他藉大膽的死亡,告訴貓,你並不值得畏懼.他藉大膽的死亡,告訴死亡,我不會在你面前下跪.他藉大膽的死亡,告訴其他老鼠,不要再卑污地活著,不要再做恐懼下的奴隸,而要榮耀地活著.

來2:14　兒女既同有血肉之體,他也照樣親自成了血肉之體,特要藉著死敗壞那掌死權的,就是魔鬼,

來2:15　並要釋放那些一生因怕死而為奴僕的人.

(3)耶穌這兄弟要完成的人性工程

來2:16　他並不救拔天使,乃是救拔亞伯拉罕的後裔.

理由很明顯,因天使沒有血肉之體,不受死亡的威脅.

來2:17　所以,他凡事該與他的弟兄相同,為要在神的事上成為慈悲忠信的大祭司,為百姓的罪獻上挽回祭.

來2:18　他自己既然被試探而受苦(或譯:因受苦而被試探),就能搭救被試探的人.

耶穌作我們的兄弟,受盡痛苦.在種種痛苦中,他屢受試探,要他放棄.但他頂住這些試探,他要讓他的所有兄弟姊妹們明白,魔鬼,死亡,虛無的威脅,是可以頂住的.並且,只要緊緊地跟住他,他會引領我們出死入生.

一隻老鼠,竟敢在貓的頸上繫鈴,這是對所有老鼠的莫大慈悲,也展示了他對宇宙人生的一種信念,他相信:人生存在是有意義的,不要因威脅而懷疑這意義,而是因這意義而克服威脅.萬物從神而來,領受著神的大恩大愛.人生,就是展示這仁愛,在威脅下展示這仁愛.向罪惡說不,向不義說不.耶穌展示了這大膽的仁愛,他成為我們的「領頭人」,我們的「大佬」,我們的「大祭司」.他為我們的罪,獻上挽回祭.或者說,為我們的貪生怕死,自私自利,他獻上自己的生命,要叫我們看到回復榮耀生命的樣式,要叫我們看到回歸天父的路.

或許,有一天,我們會敢於在貓頸上繫個鈴.

但願:榮耀歸於聖父,聖子,聖靈.阿們.

(後話:本人將早期的崇基講章結集為書,書名《敢於跟隨主》,基道出版,有興趣者可購買.)



\newpage

\section{你看見甚麼視覺與人生的探索}
\label{sec:349}
\textbf{許開明牧師 - 你看見甚麼?-視覺與人生的探索}
\newline
\newline
\hyperref[sec:348]{< < < PREV SERMON < < <}
~
\hyperref[sec:toc]{[返目錄]}
~
\hyperref[sec:350]{> > > NEXT SERMON > > >}


約翰福音 1:19-51
\newline
\begin{longtable}{cl}
\hline
\hline
章節 & 經文 (和合本修訂版)\\
\hline
1:19 & \begin{tabularx}{0.7\textwidth}{X} 這是約翰的見證：猶太人從耶路撒冷差祭司和利未人到約翰那裡去問他：「你是誰？」 \end{tabularx} \\ \\ \relax
1:20 & \begin{tabularx}{0.7\textwidth}{X} 他就承認，並不隱瞞，承認說：「我不是基督。」 \end{tabularx} \\ \\ \relax
1:21 & \begin{tabularx}{0.7\textwidth}{X} 他們又問他：「那麼，你是誰？是以利亞嗎？」他說：「我不是。」「是那位先知嗎？」他回答：「不是。」 \end{tabularx} \\ \\ \relax
1:22 & \begin{tabularx}{0.7\textwidth}{X} 於是他們對他說：「你到底是誰，好讓我們回覆差我們來的人。你說，你自己是誰？」 \end{tabularx} \\ \\ \relax
1:23 & \begin{tabularx}{0.7\textwidth}{X} 他說：「我就是那在曠野呼喊的聲音：修直主的道。」正如以賽亞先知所說的。 \end{tabularx} \\ \\ \relax
1:24 & \begin{tabularx}{0.7\textwidth}{X} 那些人是法利賽人差來的。 \end{tabularx} \\ \\ \relax
1:25 & \begin{tabularx}{0.7\textwidth}{X} 他們就問他：「你既不是基督，不是以利亞，也不是那位先知，那麼，你為甚麼施洗呢？」 \end{tabularx} \\ \\ \relax
1:26 & \begin{tabularx}{0.7\textwidth}{X} 約翰回答：「我是用水施洗，但有一位站在你們中間，是你們不認識的， \end{tabularx} \\ \\ \relax
1:27 & \begin{tabularx}{0.7\textwidth}{X} 就是那在我以後來的，我給他解鞋帶也不配。」 \end{tabularx} \\ \\ \relax
1:28 & \begin{tabularx}{0.7\textwidth}{X} 這些事發生在約旦河東邊的伯大尼，約翰施洗的地方。 \end{tabularx} \\ \\ \relax
1:29 & \begin{tabularx}{0.7\textwidth}{X} 第二天，約翰看見耶穌來到他那裡，就說：「看哪，神的羔羊，除去世人的罪的！ \end{tabularx} \\ \\ \relax
1:30 & \begin{tabularx}{0.7\textwidth}{X} 這就是我曾說『那在我以後來的先於我，因為在我以前，他已經存在』的那一位。 \end{tabularx} \\ \\ \relax
1:31 & \begin{tabularx}{0.7\textwidth}{X} 我先前不認識他，如今我來用水施洗，為要使他顯明給以色列人。」 \end{tabularx} \\ \\ \relax
1:32 & \begin{tabularx}{0.7\textwidth}{X} 約翰又作見證說：「我曾看見聖靈彷彿鴿子從天降下，停留在他的身上。 \end{tabularx} \\ \\ \relax
1:33 & \begin{tabularx}{0.7\textwidth}{X} 我先前不認識他，可是那差我來用水施洗的對我說：『你看見聖靈降下來，停留在誰的身上，誰就是用聖靈施洗的。』 \end{tabularx} \\ \\ \relax
1:34 & \begin{tabularx}{0.7\textwidth}{X} 我看見了，所以作證：這一位是神的兒子。」 \end{tabularx} \\ \\ \relax
1:35 & \begin{tabularx}{0.7\textwidth}{X} 又過了一天，約翰同兩個門徒站在那裡。 \end{tabularx} \\ \\ \relax
1:36 & \begin{tabularx}{0.7\textwidth}{X} 他見耶穌走過，就說：「看哪，神的羔羊！」 \end{tabularx} \\ \\ \relax
1:37 & \begin{tabularx}{0.7\textwidth}{X} 兩個門徒聽見他的話，就跟從了耶穌。 \end{tabularx} \\ \\ \relax
1:38 & \begin{tabularx}{0.7\textwidth}{X} 耶穌轉過身來，看見他們跟著，就問他們說：「你們要甚麼？」他們對他說：「拉比，你在哪裡住？」（「拉比」翻出來就是老師。） \end{tabularx} \\ \\ \relax
1:39 & \begin{tabularx}{0.7\textwidth}{X} 耶穌說：「你們來看。」他們就去看他在哪裡住。這一天他們就跟他同住；那時大約是下午四點鐘。 \end{tabularx} \\ \\ \relax
1:40 & \begin{tabularx}{0.7\textwidth}{X} 聽了約翰的話而跟從耶穌的那兩個人，其中一個是西門‧彼得的弟弟安得烈。 \end{tabularx} \\ \\ \relax
1:41 & \begin{tabularx}{0.7\textwidth}{X} 他先找到自己的哥哥西門，對他說：「我們遇見彌賽亞了。」（「彌賽亞」翻出來就是基督。） \end{tabularx} \\ \\ \relax
1:42 & \begin{tabularx}{0.7\textwidth}{X} 於是安得烈領西門去見耶穌。耶穌看著他，說：「你是約翰的兒子西門，你要稱為磯法。」（「磯法」翻出來就是彼得。） \end{tabularx} \\ \\ \relax
1:43 & \begin{tabularx}{0.7\textwidth}{X} 又過了一天，耶穌想要往加利利去。他找到腓力，就對他說：「來跟從我！」 \end{tabularx} \\ \\ \relax
1:44 & \begin{tabularx}{0.7\textwidth}{X} 這腓力是伯賽大人，是安得烈和彼得的同鄉。 \end{tabularx} \\ \\ \relax
1:45 & \begin{tabularx}{0.7\textwidth}{X} 腓力找到拿但業，對他說：「摩西在律法書上所寫的，和眾先知所記的那一位，我們遇見了，就是約瑟的兒子拿撒勒人耶穌。」 \end{tabularx} \\ \\ \relax
1:46 & \begin{tabularx}{0.7\textwidth}{X} 拿但業對他說：「拿撒勒還能出甚麼好的嗎？」腓力說：「你來看。」 \end{tabularx} \\ \\ \relax
1:47 & \begin{tabularx}{0.7\textwidth}{X} 耶穌看見拿但業向他走來，就論到他說：「看哪，這真是個以色列人！他心裡是沒有詭詐的。」 \end{tabularx} \\ \\ \relax
1:48 & \begin{tabularx}{0.7\textwidth}{X} 拿但業對耶穌說：「你從哪裡認識我的？」耶穌回答他說：「腓力還沒有呼喚你，你在無花果樹底下，我就看見你了。」 \end{tabularx} \\ \\ \relax
1:49 & \begin{tabularx}{0.7\textwidth}{X} 拿但業回答他：「拉比！你是神的兒子，你是以色列的王。」 \end{tabularx} \\ \\ \relax
1:50 & \begin{tabularx}{0.7\textwidth}{X} 耶穌回答他說：「因為我說在無花果樹底下看見你，你就信嗎？你將看見比這些更大的事呢！」 \end{tabularx} \\ \\ \relax
1:51 & \begin{tabularx}{0.7\textwidth}{X} 他又說：「我實實在在地告訴你們，你們將要看見天開了，神的使者在人子身上，上去下來。」 \end{tabularx} \\ \\
[1ex]
\hline
\hline
\end{longtable}
講題:你看見嗎?-視覺與人生的探索　What have you seen? How does your sight affect your search for life?

經文:約翰福音1:19-51

講員:許開明牧師

場合:香港中文大學崇基學院禮拜堂

日期:2017年1月15日

引言

我舉起手中的一枝筆,你看見甚麼呢?你看見的當然是一枝筆,這枝筆是紅色的,筆的另一端還可以用來滑智能電話或平板電腦,仔細地看,還可以看到筆上寫著「Chung Chi Chapel」,也就是崇基禮拜堂的筆,是之前在父親節及母親節所送的紀念品.

你看見甚麼?當中可能我們會感困惑,很多時候你看見的,是否就是這樣?或者這只是你所認為或者理解的呢?你看見的,可能只是事物的一部份,並非全部.

剛拍拖時,女友好像經常想試我的視力是否正常.經常問我:「你看見什麼?」我說:「我沒有見到什麼.」她又說:「木瓜樹上好大個木瓜,你看不見麼?」我說:「對啊,好似有一個木瓜.」她說:「不只是一個木瓜,是十個木瓜啊!」我這次的測試顯然是不及格的,下次拍拖我就要打醒十二分精神.

所以很多時候我們看見,都會忽略了一部分,未必不能一次就看見全部,就如今天的有許多看見,遇見,值得我們深思.這段經文是顯現期的經文,特別講到約翰很想人們看見耶穌.

觀察經文:約一29- 51

第一天:施洗約翰見證自己  (19-29)

在29節之前,由第19節起,就是第一天,講述施洗約翰見證自己要成為主的先鋒,修直主的道.

第二天:約翰看見耶穌:引人看神的羔羊 (29-34)

約翰引導人看見耶穌的時候說了兩項見證,就是聖靈以鴿子的形象停在耶穌身上,以及差他來用水施洗的上帝對他說耶穌要來用聖靈為人施洗,然後約翰說:「我看見了」,也即是「我明白,理解了」的意思.

第三天:引門徒看耶穌(35- 42)

約翰向門徒介紹耶穌為神的羔羊,門徒就跟從了耶穌,耶穌問他們要甚麼,原文的另一個意思是「你們要尋找甚麼」,「尋找」這一詞與「看見」有很大的關連,因為「尋找」的英文是seek,就是see(看見)加上k,而我們找東西的時候,必定會先用眼睛四處去看.門徒想與耶穌同住,耶穌回答說:「你們來看.」,這話亦包含了「你們來了解我(耶穌)的意思.後來安得烈把遇見耶穌的事告訴哥哥西門並領他前去見耶穌,耶穌見了西門,就為他起名叫「磯法」,翻出來就是彼得,石頭,磐石,在那時候,石頭就代表「有用的人」的意思,預告了彼得將來要成為教會的磐石.

第四天:耶穌見腓力和拿但業 (43- 51)

第四天,耶穌遇見了腓力,腓力又帶了拿但業來,並向他介紹耶穌,拿但業聽見耶穌是拿撒勒人,就感不悅,因為當時的拿撒勒人最出名的就是木工,但他們就被羅馬人利用,所造的好木都是用來造釘猶太人的十字架,所以當時的以色列人都很討厭拿撒勒人.但後來拿但業聽見耶穌說在無花果樹底下已看見他,他就深信耶穌是神的兒子,並跟隨了耶穌.

本段經文有許多的寶貝──「看見」,很值得我們留意!

人生看東西有三隻眼

第一就是肉眼,是用來感覺的眼,要求清楚,準確,明,透,肉眼對各職業都重要,但對醫生的眼,飛機師的眼,軍人的眼來說,更為重要.

但在蘇格拉底的時期,Parmenides提出了「存在論」,即存在者就會存在.羅馬又有一位Domocrites,他繼承了Parmenides的想法,更提出了「虛無主義」及「原子論」,即不存在者或者是細小的肉眼看不見的,也存在.所以其實人能看見的東西其實很少,只是全部中的一部分.

第二是心眼,是心靈的眼,或作理解力,感應能力.

第三是慧眼,是用作領略人生智慧的眼,或是獨到的慧眼,即有眼光,就能得良人,好助手；但亦有人會「有眼無珠」,錯信人而受騙,受傷,受驚,受害,這人就是缺了慧眼.

在基督教的信仰中,有第四種眼──靈眼,是屬靈的眼,信心的眼睛,讓我們能超越一般的限制,看見我們的神.經文中,約翰多次說:「看哪,上帝的羔羊」,當中的「看」是「凝視」的意思,要更多了解他.

自古到今,人都想要得到別人的注意,例如我的女兒小時候常常對我說:「看看,我多漂亮!」不久的將來,可能某位大人物會公開或喑中說:「看哪,是中央銜點的特首!」這人會是姓曾?姓鄭?姓林?你要注意,你也要深入了解,這個被抬舉的人,對不對版?這可能是香港人近期的掙扎.

人用肉眼看；用心眼了解；用慧眼處世；用靈眼看信仰.看見,是有不同的層次,不同角度和深度,能幫助我們活得更好.

1看見耶穌後的見証

約翰看見耶穌,就見證耶穌,並引人去認識耶穌.約翰看得見,看見聖靈從天而降,降到耶穌身上,又看見耶穌受聖靈的洗；約翰的看見更是深入,因他引用了「神的羔羊」,包含了基督,上帝的祭物,上帝的僕人,得勝者的意思.約翰看得深的另一點是因為羔羊有最少四大項的貢獻:羊血代表生命,因為逾越節時天使去擊殺埃及人的長子時,見到家門上有羊血,就逾過,不殺以色列人；羊肉又是食物；羊皮可以製成羊毛衣保暖；羊奶更比牛奶健康,因為它含的脂肪比牛奶少.

2看見耶穌後跟從

兩個門徒與耶穌同住,就更深入認識耶穌.

3看見耶穌,並領人見他

他們認識耶穌的好,就領身邊的人去見耶穌,就是安得烈領西門彼得,腓力領拿但業去見和認識耶穌.

4看見耶穌後,發現耶穌早已看見他們,了解他們

5看見耶穌後,確信耶穌是上帝兒子,以色列的王

6看見耶穌後,得著信心的看見,看到更大的事,更大的榮耀

從引見,看見,見証,領見,等見,信心之見,人的生命得到了改變,改進,改善,所以人的眼具生命意義,是靈魂之窗,我們卻甚少善用.

留意人的眼神

在羅馬很出名的哲學家哥達維斯(Godas)說過:「人很留意別人的眼神.」因為別人看見你,表示他知道你的存在,而你看著他,也會醒悟自己的存在,自己的使命,知道他人與我是不同的個體,他人的眼神,表示他人看見我的存在,留意我,重視我.正如門徒看見耶穌,才發現耶穌早已看見他們,認識他們,並可以說出他們的名字,知道他們的存在.一個人被知道,被重視很重要的.

你記得誰的眼神嗎?我想起一位快離世的牧者,他臨終前對我說:「我書房的書,你全部拿去吧.」我很快樂,因為我很喜歡看書.

又有另一個眼神,那是我10歲時做完割禮的手術後醒來,一位護士對我的微笑眼神,因為她的微笑讓我知道我沒有死去,還知道痛.

眼與眼的接觸是現代社會很需要的,因為它能讓人彼此看見,建立關係.耶穌看見西門,認識西門,叫他作畿法,這是關係的建立；耶穌看見拿但業,已認識拿但業,拿但業也認識那穌,彼此欣賞,這也是關係的建立.

然而,弔詭的是,人渴慕人的眼神,但同時我們也怕與人眼神接觸?

七十至八十年代起,我們常對著電視；九十年代起,我們就對著電腦,繼而到手機；2010年後就發現成只對著智能手機,因為一機在手,甚麼電都擁有,我們可以用智能電話看電視,聽電台等.我們看手機多於看人,我們看手機上人的相片多於面對面的看人的眼睛.

我們怕望人,怕望人的眼,特別怕望你的對頭人,你不想遇見的人的眼,又或是怕看人的怒眼,不屑一顧的眼,不友善的,無精打采的眼,還有怕你看他但他不看你.假如今天有人看著你超過六秒,可能反而令你害怕,令你擔憂.女的可能會想:那凝視著我的男士會有不軌企圖嗎?男的可能會想:那凝視我的人想向我「埋手」?事實上,人看見或尋找他或它,想了解他或它,人可能便想佔有,控制他,因為人的眼與人的欲望有關.

耶穌的深層看見

耶穌的看見是與別不同的:

1看見約翰的門徒時,耶穌看見他們想要尋找救主,所以就叫他們現在來住看他,親自認識他.

2看見安得烈引彼得時,耶穌改西門名作彼,因他看見彼得的未來

3看見腓力時,耶穌看見他的需要,要腓力來跟從他,改變腓力的一生.

4看見拿但業時,耶穌不單看見他的外表和行為是真以色列人,也看見他的過去,當拿但業看見耶穌時,耶穌引發他信心的眼睛:看見天開了路.而更要的是,沒有人能像耶穌一樣,看見人的過去,現在及未來.

我看見一位朋友突然瘦了30磅,發現了他原來在減肥,減肥的原因是發現身體有毛病,不可過重,再看下去,原來身體出現毛病是由於工作壓力等問題.一個看見,會引發另一個看見,使我們看得更深,更高,更多!在研究上也是,從一個現象,可以看見發現另一個現象.當牛頓看見蘋果跌下,他看見有重量的東西會下跌,然後他看見了地心吸力.我們能否有深度的看見?肉眼看不清,不深,不準；心眼也一樣,會有偏見,而現代人的盲點是看得快,想得快,講得快.我們看人不準,只見眼前的人,或是他所處理的某件事,或是他所說的某句話,他的某篇文章.我們看人常不全面,易有偏見.所以別人看你,批評你,也不全面,對於那些不盡不實的批評,你又何必如此介懷呢.

結語

你們看見什麼?

一位大學老師,拿了張白色大卡紙,然後在中間點了黑點,問大學生:「你們看見什麼?」沒有一位大學生肯回答他的問題.等了一陣子,還是沒有人回答!這位老師,換了張黑色卡紙,取了枝白色筆,又點了一點,又問同學同樣問題,又是沒有人理會他.有位同學忍不住說:「老師,我們不是『小學雞』呀!這麼低智商的問題,有損我們的身份.」老師回答說:「那更加要答,答不到就更有損你作為大學生的身份.」同學說:「好,那我就答.黑色卡紙上,我見到白色的一點.

白色卡紙上,我見到黑色的一點.是人都知道的,老師!」老師問:「各位,對不對?」又無人理會他.

這故事當中有何含意?在黑色卡紙上,見到白色的一點,這是樂觀的人,因為白點代表光明；但在白色卡紙上,見到黑色的一點,這就是悲觀的人,因為四同的都是光的,你只看到黑色的一點.

你看見甚麼,就是甚麼?不是,你看見的,可能不是你認為的,因為你看見的,只是事物的一部份,不是全部；你所看見,可能不是是真相,不是全部,或只是別人想你看到的部份,就例如西九文物館一事.

主阿,我要能看見

這曾經是我們的禱告(可10章51):巴底買他心眼看到耶穌,就喜樂和跳舞,耶穌醫好他肉眼,他看見了世界,他看見了主,就跟從了主耶穌.作為基督徒的我們,也會有巴底買的經驗:當巴底買失明時,他是用心眼,信心的眼睛來看耶穌；當巴底買開眼時,他跟從耶穌,看到耶穌奇妙的作為.我們也需先用信心的眼來看耶穌,然後就能用肉眼看到耶穌的手奇妙的作為.

年近歲晚,朋友送來很多名言,這則房子的啟示很有趣:

「Someone asked himself how to handle life?

有人問自己:『怎樣面對我的生命?』

The flat gave him all the answers:

竟然從房間中得到所有答案:

Roof said: Aim high.

屋頂說:『要高瞻遠矚.』

Fan said: Be cool.

電扇說:『要保持冷靜.』

Clock said: Value time.

時鐘說:『要珍惜光陰.』

Calendar said: Be up to date.

日曆說:『要與時俱進.重大

Wallet said: Save now for the future.

錢包說:『要積穀防飢.』

Mirror said: Always observe yourself.

鏡子說:『要自我反省.』

Lamp said: Light up other’s life.

桌燈說:『要燃亮別人.』

Wall said: Share others’ load.

牆壁說:『要分憂解勞.』

Bed said: Take a rest and dream.

床說 : 『要有休息和夢想.』

Window said: Expand your vision.

窗戶說:『要拓寬視野.』

Floor said: Always be down to earth.

地板說:『要腳踏實地.』

Stairs said: Watch each step you take.

樓梯說:『要步步為營.』

The most inspiring one:

最鼓舞人的一則是:

Toilet bowl said: When it’s time to let go, just let it go…

馬桶說:『要懂得放下.』」

祈禱

祈求上主,賜我們明亮的眼睛,凡事看得清,看得準,看得透；賜我們心眼,懂得思考分析,明辨是非；也賜我們慧眼,為人處事充滿智慧；更賜我們屬靈的眼睛,信心的眼睛,看到上主的帶領和心意!阿們.



\newpage

\section{坐在黑暗裡的百姓看見了大光}
\label{sec:350}
\textbf{李均熊博士 - 坐在黑暗裡的百姓看見了大光……}
\newline
\newline
\hyperref[sec:349]{< < < PREV SERMON < < <}
~
\hyperref[sec:toc]{[返目錄]}
~
\hyperref[sec:351]{> > > NEXT SERMON > > >}


馬太福音 4:12-23
\newline
\begin{longtable}{cl}
\hline
\hline
章節 & 經文 (和合本修訂版)\\
\hline
4:12 & \begin{tabularx}{0.7\textwidth}{X} 耶穌聽見約翰下了監，就退到加利利去； \end{tabularx} \\ \\ \relax
4:13 & \begin{tabularx}{0.7\textwidth}{X} 後來離開拿撒勒，往迦百農去，住在那裡。那地方靠海，在西布倫和拿弗他利地區。 \end{tabularx} \\ \\ \relax
4:14 & \begin{tabularx}{0.7\textwidth}{X} 這是要應驗以賽亞先知所說的話： \end{tabularx} \\ \\ \relax
4:15 & \begin{tabularx}{0.7\textwidth}{X} 「西布倫，拿弗他利，沿海的路，約旦河的東邊，外邦人的加利利— \end{tabularx} \\ \\ \relax
4:16 & \begin{tabularx}{0.7\textwidth}{X} 那坐在黑暗裡的百姓看見了大光；坐在死蔭之地的人有光照耀他們。」 \end{tabularx} \\ \\ \relax
4:17 & \begin{tabularx}{0.7\textwidth}{X} 從那時候，耶穌開始宣講，說：「你們要悔改！因為天國近了。」 \end{tabularx} \\ \\ \relax
4:18 & \begin{tabularx}{0.7\textwidth}{X} 耶穌沿著加利利海邊行走，看見兩兄弟，就是那叫彼得的西門和他弟弟安得烈，正往海裡撒網；他們本是打魚的。 \end{tabularx} \\ \\ \relax
4:19 & \begin{tabularx}{0.7\textwidth}{X} 耶穌對他們說：「來跟從我，我要叫你們得人如得魚一樣。」 \end{tabularx} \\ \\ \relax
4:20 & \begin{tabularx}{0.7\textwidth}{X} 他們立刻捨了網，跟從他。 \end{tabularx} \\ \\ \relax
4:21 & \begin{tabularx}{0.7\textwidth}{X} 耶穌從那裡往前走，看見另外兩兄弟，就是西庇太的兒子雅各和他弟弟約翰，同他們的父親西庇太在船上補網，耶穌就呼召他們。 \end{tabularx} \\ \\ \relax
4:22 & \begin{tabularx}{0.7\textwidth}{X} 他們立刻捨了船，辭別父親，跟從了耶穌。 \end{tabularx} \\ \\ \relax
4:23 & \begin{tabularx}{0.7\textwidth}{X} 耶穌走遍加利利，在各會堂裡教導人，宣講天國的福音，醫治百姓各樣的疾病。 \end{tabularx} \\ \\
[1ex]
\hline
\hline
\end{longtable}
以賽亞書 9:1-4
\newline
\begin{longtable}{cl}
\hline
\hline
章節 & 經文 (和合本修訂版)\\
\hline
9:1 & \begin{tabularx}{0.7\textwidth}{X} 但那受過痛苦的必不再見幽暗。從前神使西布倫地和拿弗他利地被藐視，末後卻使這沿海的路，約旦河東，外邦人居住的加利利地得榮耀。 \end{tabularx} \\ \\ \relax
9:2 & \begin{tabularx}{0.7\textwidth}{X} 在黑暗中行走的百姓看見了大光；住在死蔭之地的人有光照耀他們。 \end{tabularx} \\ \\ \relax
9:3 & \begin{tabularx}{0.7\textwidth}{X} 你使這國民眾多，使他們喜樂大增；他們在你面前歡喜，好像收割時的歡喜，又像人分戰利品那樣的快樂。 \end{tabularx} \\ \\ \relax
9:4 & \begin{tabularx}{0.7\textwidth}{X} 因為他們所負的重軛和肩頭上的杖，並欺壓者的棍，你都已經折斷，如同在米甸的日子一般。 \end{tabularx} \\ \\
[1ex]
\hline
\hline
\end{longtable}
講題:坐在黑暗裡的百姓看見了大光……　The people who sat in the darkness have seen a great light……

經文:馬太福音4:12-23;以賽亞書9:1-4

講員:李均熊博士

場合:香港中文大學崇基學院禮拜堂

日期:2017年1月22日

引言:黑暗與大光的一般的理解

「西布倫,拿弗他利,沿海的路,約旦河的東邊,

外邦人的加利利──

那坐在黑暗裡的百姓看見了大光；

坐在死蔭之地的人有光照耀他們.」(太四16)

今天是顯現期的第三主日.今日經課的福音書引用了以賽亞書九章1-4節的內容.我在神學院是任教舊約的科目,所以在這裡我希望和大家思考這節經文對我們今天的意義.

若然基督徒又或在講壇談到「黑暗中見到大光」,我們很容易便會聯想起天使與魔鬼,善良與邪惡；就算沒有這樣二元化,黑暗中見大光也會在傳統福音派的意識裡代表著罪人在罪惡黑暗中看見耶穌的真光,脫離罪惡黑暗而能重生.其實這種理解比較接近約翰福音的開始,但馬太福音的主要讀者,卻未必會有這樣的聯想.

帝國政治下的黑暗

以賽亞書九章的經文原本是這樣的:「在黑暗中行走的百姓看見了大光；在死蔭之地的人有光照耀他們.」馬太在引用以賽亞書時省略了第一句的動詞,使第一句成為只有一連串地名的子句.這樣馬太便將引文的焦點轉移成為集中那些坐在黑暗百姓的地理位置上.西布倫,拿弗他利和加利利都是屬北方以色列國的境域,對身處南國耶路撒冷的以賽亞先知而言,乃正面臨,甚或已經受亞述帝國侵吞的地區.

在以賽亞書七至八章,經文題到各種結盟,同謀,背叛,其實是指國與國之間的合縱連橫,焦點都是關於亞述對當時地區的威脅和各國的反應；更有可能,第八章所回應的處境,正正是北國被滅,亞述成為統治西布倫,拿弗他利以及加利利地區的帝國,猶大國失卻北國以色列的緩衝而將直接面對亞述的影響.故此到了第九章經文所指的「坐在黑暗的百姓」可以是非常政治性的,而非單單喻意屬靈的狀況.甚至可以照字面解釋,因著亞述迫害被擄被囚的百姓,他們的的確確坐在暗不見天日的地窖囚牢裡所感受到的黑暗和藐視.

所以「外邦人的加利利」不一定指加利利的居民是外邦人,更有可能這裡是指加利利現在是屬於外邦人的地方.耶穌事奉期間加利利地區是屬於監禁甚至殺害施洗約翰的分封王希律安提帕管轄的範圍(十四1-12),那是羅馬的分封王,不全然是猶太人可以自主的地方,加上在馬太成書的時期,學者一般相信是置於主後七十年後,亦即羅馬軍隊燒毀耶路撒冷聖殿的日子之後,所以當時讀者承受政治壓迫,同樣巨大.加上引入這段舊約經文的處境乃施洗約翰下監的消息,也引人思考經文的政治性含意.

馬太福音所描繪關於耶穌的降生與開始事奉的事蹟,基本上是很政治化的.希律殺嬰,合城不安,約瑟避走,約翰下監…路加的聖誕若然是普世歡騰,那麼馬太的聖誕卻是一個憂鬱不安的聖誕.

換言之,當時讀者所關注的,極有可能是政治中的不安和壓迫,而非今日教會常常屬靈化了的光明黑暗；馬太這裡引用的以賽亞書,亦極有可能被理解成在帝國的統治下和政治的黑暗和絕望下,上帝的國要帶來光明的盼望.

這情況可能對今日香港教會來說也許十分貼身.馬太或者他所引用的以賽亞書所提到的黑暗與光明,不是只有屬靈化的意義,也有現世政治的意思.馬太福音如此開始,必然吸引著當時承受政治不安和迫害的信徒的注意力.從此觀之,我們教會的宣講,是否只關注靈魂在罪惡黑暗中得耶穌的拯救,還是可以如馬太福音般對關心政治現實的讀者作出宣講?反過來說,教會的信息又是否吸引在政治現實下面對不安和困境的信徒呢?我們是不是也要反思,耶穌基督的天國的福音是否也能回應今日的政治處境?

但是,馬太所指的光明和上帝國卻又是什麼意思呢?它們的意思是指今世政治的改變嗎?

大光是什麼?──馬太的應用

「光」只有在登山寶訓裡有提及(五14, 16; 六23),說不上是馬太福音的主題,我們也不要太快和約翰的說法相題並論.但馬太引用以賽亞書所說到的加利利地區,與之關聯的迦伯農,卻是馬太福音記載耶穌事蹟的一個很重要的地方.在馬太第九章所謂的「耶穌在自己的城裡」應該也是指迦伯農.在迦伯農,耶穌醫病趕鬼,大行神蹟,分別醫好了羅馬百夫長的僕人(八6-13),西門彼得的岳母(八14-15),他也在自己的城裡使被人用褥子台著的癱子行路(九1-8),使經血不止的病人痊癒(九20-22),使管會堂的主管的女兒死而復生(九18-26),又開盲人的眼(九27-31),趕出啞巴鬼(九32-34)……以至後來約翰的門徒來問耶穌他是否那「將要來的那位」的時候(十一2-6),基本上耶穌的回答:「盲人看見,瘸子行走,痳瘋病人得潔淨,聾子聽見,死人復活,窮人聽到福音…」；就是他在迦伯農和附近所行的一切.我想,這些記載不是更加符合馬太引用的以賽亞書:「那坐在黑暗裡的百姓看見了大光,坐在死蔭之地的人有光照耀他們」嗎?耶穌帶來的光,是為在現實地上面對種種困苦的人民帶來釋放的福音,而非單單屬靈的得救.

天國近了是什麼意思?

從耶穌回答施洗約翰的內容,馬太引用以賽亞書的內容其實要指出耶穌,甚至約翰所講的天國近了,其實就是要應驗先知所說的,上主要來復興以色列,並且上主的王權得顯.但當中一個重要的要求,就是領受的人必須悔改.要上主王權實現人間,驅走現實的黑暗,無論是政治的或屬靈的,均需要心靈的悔改,相信耶穌是神差來的救主彌賽亞.(四17你們要悔改,因為天國近了)

但天國是否如耶穌所說的實現呢?天國在耶穌在生的時候,並不能如期實現,因為上主所要求的悔改,卻並沒有在以色列當中見到.這我們可以從馬太的十一章以後可以見到:哥拉汛,伯賽大是不肯悔改的城,耶穌的家鄉拿撒勒厭棄耶穌,文士法利賽人質疑耶穌的權柄,更認為他是靠鬼王別西卜趕鬼,耶穌所行的神蹟並不能使他們相信,他們要求耶穌顯更多神蹟,目的不是要相信他,其實是要試探他,找把柄去除掉耶穌.

故此在第十三章,馬太轉而談耶穌的比喻,因為用比喻的目的,就是要將天國的奧秘隱藏起來,只能讓相信和跟從的門徒知道,不讓不信抗拒的人知道,「因為百姓的心麻木,耳朵發沉,眼睛閉著」.會不會在馬太福音開始的「天國」是關乎現世的解放,但隨著猶太人的棄絕耶穌,抗拒上主的王權,不相信,不接受,不悔改,之後「天國」便隱藏起來,成為末世性的審判和更新?(參馬太廿四章)但怎樣體會天國現世的實現,又何時天國又只是末世的盼望?耶穌所表現的天國是現世的,但那只能說明天國近了,因為天國另一方面也是末世的,只能透過真心悔放相信才能完全實現.

結語

「天國近了」表達了我們信仰上一個很吊詭的現實.天國已然來臨,卻又未完全實現,仍在途中,仍需等待.面對政治現實,人間的苦情,等待的人如何憂慮心急,我們也只能說「近了」.耶穌來到世界,叫坐在黑暗的加利利人先嘗天國在世上現實的光明解放；但今日,天國卻又是上主再臨末世的審判；天國是既是今世的,政治的,也是末世的,屬靈的.在這之間,我們要如何自處呢?

其實我還在思考,馬太所言的這個天國福音,對我們今天是什麼意思?「天國近了,你們應當悔改」,看來現世的天國實現與末期的天國降臨,關鍵在於我們的悔改.但是否只有人人都悔改,現實坐在黑暗裡的人才有天國的真光?是否只有人人都悔改,今日的政治現實才可以有改變?若然真的人人都悔改,我相信世界的確是可以因而改變!但今天面對每天的景況,我個人就比較悲觀.就如在馬太福音,當人遇到耶穌的釋放和權柄,我們只看到人的反抗,不信和拒絕……又或將我們的眼光放回以賽亞的時代,先知要求的是信靠上主就不動搖,但阿哈斯所代表的人間王權卻只會虛應上主,轉過頭就尋求政治的妥協.政治的現實,政治的壓力迫使我們對現實採取不同的行動.聖經是不是不能回應這等事情?我看不同的聖經內容其實對現世政治是有一定程度的解釋能力的.因為天國是要實現在人間的,正如主禱文說「願你的國降臨,願你的旨意行在地上」；又如啟示錄所展示的,新耶路撒冷是要從天而降在這地上的.但聖經同時叫我們看見,天國是將來的實現,是末期的彰顯；我們還要等候,為什麼要等候?原恩惠不會如彼得後書3章9節所言:不願一人沉淪,而是人人都來悔改?但在這將濟又未濟的間隙,悔改是困難的,悔改是否真實可以達至的呢?對特首說你要悔改,香港就可以得救,那是否有用?

聖經福音的信息不是單一的,乃是多向,動態,處境性的.馬太福音乃甚舊約先知書所呈現的天國同時是地上的,也是末後的,正好幫助我們承載現今複雜的世情,免於過度單純簡化我們信仰的思與行成為只有抽離世界社會的屬靈意義；同時也不能一味的將我們的言行放在只有政治的綱領來考慮量.同時聖經的宣告其實迫使我們不要過度極化,二元化上帝的福音.耶穌已經顯現,真光的確來到世界,天國已然近了；正因為如此,我們更應要思考悔改在今日現實政治世界的真正意義.說到底,在馬太福音的宣講中,耶穌基督的福音已經傳給我們,天國的奧秘也只向相信祂的人顯明,基督也期望我們使萬民成為門徒；我們在世作基督的門徒,我們又真的對天國奧秘有多少掌握?我們的悔改又有多真實?

顯現期上主使我們看見光明.這光明又能否藉我們穿透今日社會政治以至人心中的黑暗?教會又應如何承擔使人悔改以至讓人掌握天國奧秘的責任?今日就讓我們透過馬太的宣講一同思考.



\newpage

\section{何以對蒼天}
\label{sec:351}
\textbf{林豪恩先生 - 何以對蒼天}
\newline
\newline
\hyperref[sec:350]{< < < PREV SERMON < < <}
~
\hyperref[sec:toc]{[返目錄]}
~
\hyperref[sec:352]{> > > NEXT SERMON > > >}


彌迦書 6:1-8
\newline
\begin{longtable}{cl}
\hline
\hline
章節 & 經文 (和合本修訂版)\\
\hline
6:1 & \begin{tabularx}{0.7\textwidth}{X} 當聽耶和華說的話：起來，向山嶺爭辯，使岡陵聽見你的聲音。 \end{tabularx} \\ \\ \relax
6:2 & \begin{tabularx}{0.7\textwidth}{X} 山嶺啊，要聽耶和華的指控！大地永久的根基啊，要聽！因耶和華控告他的百姓，與以色列爭辯。 \end{tabularx} \\ \\ \relax
6:3 & \begin{tabularx}{0.7\textwidth}{X} 「我的百姓啊，我向你做了甚麼呢？我在甚麼事上使你厭煩？你回答我吧！ \end{tabularx} \\ \\ \relax
6:4 & \begin{tabularx}{0.7\textwidth}{X} 我曾將你從埃及地領出來，從為奴之家救贖你，我差遣摩西、亞倫和米利暗在你前面帶領。 \end{tabularx} \\ \\ \relax
6:5 & \begin{tabularx}{0.7\textwidth}{X} 我的百姓啊，當記念從前摩押王巴勒如何籌算，比珥的兒子巴蘭如何回應他，當記念從什亭到吉甲所發生的事，好使你們明白耶和華公義的作為。」 \end{tabularx} \\ \\ \relax
6:6 & \begin{tabularx}{0.7\textwidth}{X} 「我朝見耶和華，在至高神面前跪拜，當獻上甚麼呢？難道獻一歲的牛犢為燔祭來朝見他嗎？ \end{tabularx} \\ \\ \relax
6:7 & \begin{tabularx}{0.7\textwidth}{X} 耶和華豈喜悅千千的公羊，或是萬萬的油河嗎？我豈可為自己的過犯獻我的長子，為自己的罪惡獻我所親生的嗎？」 \end{tabularx} \\ \\ \relax
6:8 & \begin{tabularx}{0.7\textwidth}{X} 世人哪，耶和華已指示你何為善。他向你所要的是甚麼呢？只要你行公義，好憐憫，存謙卑的心與你的神同行。 \end{tabularx} \\ \\
[1ex]
\hline
\hline
\end{longtable}
講題:何以對蒼天　With What Should I Come Before the Lord

經文:彌迦書6章1-8節

講員:林豪恩先生

場合:香港中文大學崇基學院禮拜堂

日期:2017年1月29日

一．引言

今天是年初二,祝各位新年蒙福.農曆新年,又有稱為春節,是華人社會的一個重要節日.相傳春節最早起源於祭天,華夏民族自古就相信,天地哺育眾生,是最高的神.祭祀儀式不能隨便,無論是對象,規模和形式都要講究,不能出錯.例如,祭天,並不是人人都可以做的,只可以由天子來親自主持.有沒有參觀過位於北京的天壇?這便是明清兩朝帝王祭天的場所.現在沒有雖然沒有皇帝,天壇的春節祭天儀式仍然舉行,不過卻以旅遊文化節目形式出現.

天,只有稱為天子的皇帝才可以祭祀,那麼,民間在春節祭祀甚麼呢?傳說地位最高的天,不降臨百姓家中,而是派不同的神執掌不同職能,例如門神,灶神等.有些華人地區的新年祭祀由農曆十二月十六日開始:此為尾牙-拜土地公；十二月二十四日:送神-拜灶神與眾神明；十二月二十九日或三十日:除夕-拜地基主,家神,祖先；正月初一:開春-拜天公；正月初四:迎神-接財神,灶神與眾神明等.

現代城市人還承繼新年拜神的習俗嗎?當看見除夕夜黃大仙祠數以萬計的人爭「上頭炷香」的熱烈狀況,以及年初三往車公廟拜祭的人龍知道答案了.

天若有情,天神對人間有甚麼期望呢?天若有意,天機又是甚麼呢?從另一個角度來看,人們又應如何面對蒼天呢?今天的舊約經課,記載了一段上帝與人關於這個問題的爭辯.為此,他們甚至鬧上法庭.今天,請大家一齊來坐在旁聽席聽審這件案件.

二．釋經及應用

1．開庭

在電視看包青天開庭審訊時,衙役站兩排,先由一人高喊升堂!接著堂下眾人齊聲吶喊威武.驚堂木下--砰的一聲,包公威風凜凜氣勢十足的座下,場景令人印象深刻.

讓我們翻開彌迦書第6章1-2節,看看先知如何描述那次的開庭狀況.

6:1以色列人哪,當聽耶和華的話!要起來向山嶺爭辯,使岡陵聽你的話.

6:2山嶺和地永久的根基啊,要聽耶和華爭辯的話!因為耶和華要與他的百姓爭辯,與以色列爭論.

原來,天地都不是上帝本身,是上帝呼籲來聽審的陪審員.話說回來,到底是甚麼嚴重案件,上帝要呼天喚地來做陪審團,在天地之間開庭審訊呢?

2．耶和華對以色列的指控

6:3我的百姓啊,我向你做了甚麼呢?我在甚麼事上使你厭煩?你可以對我證明.

原來,耶和華上帝要訴訟的對象,竟然是祂的選民以色列.訴訟的內容是甚麼呢?「我的百姓啊,我向你做了甚麼呢?我在甚麼事上使你厭煩?」看來是耶和華上帝的選民以色列對耶和華上帝表示討厭,認為耶和華上帝麻煩.

「你好煩啊!」或者「你可唔可以唔好咁煩?」這些話說,有沒有聽過呢?或者有沒有說過呢?最在甚麼地方聽到呢?甚麼人對甚麼人說呢?

「天氣凍啊,著多件衫啦!」「我唔覺得凍呀,唔好咁煩啦!」

「那麼晚,還不睡覺?」 「你好煩呀!」

「為什麼起床總是不摺被?」 「你好煩呀!」

「考試快到,還不溫習?」「你好煩呀!」

「你好煩呀!」這句說話,很可能我們也聽過,甚至也說過.原來,當時耶和華的選民也向耶和華說:「你好煩呀!」究竟耶和華做了甚麼事令選民感到厭煩呢?讓我們繼續聆聽耶和華的發言.

6:4我曾將你從埃及地領出來,從作奴僕之家救贖你；我也差遣摩西,亞倫,和米利暗在你前面行.

6:5我的百姓啊,你們當追念摩押王巴勒所設的謀和比珥的兒子巴蘭回答他的話,並你們從什亭到吉甲所遇見的事,好使你們知道耶和華公義的作為.

耶和華這段說話的意思是說明,祂沒有做甚麼值得祂的選民厭煩的事情,所做的,就是把以色列人在埃及為奴之地解救出來,並且沿途保護他們,帶領他們進入應許之地.反過來說,若以色列人因為這些事而厭煩耶和華,這是以色列人不懂得感恩圖報,反而埋怨耶和華,即是俗語所說的「良心當狗肺」.

耶和華與祂的選民的爭辯,是否越聽越有熟識的感覺呢?

去年重陽節,佛山大學龍建剛先生寫了一篇名為《我們如此深愛我們的兒女,他們愛我們嗎?》在微信發放,引發社會熱議,當時點擊率日過十萬.

「2016年10月4日晚九點,國慶長假第四天,廣州一所著名大學任教的老友在電話中泣不成聲,斷斷續續講了很久才弄清事情的原委:10月1日清早,他們夫婦倆乘高鐵前往南京,看望在南京大學中文系就讀的女兒.行前為女兒準備了很多她愛吃的零食,比如蛋塔,優酪乳什麼的,打算給女兒一個意外的驚喜.沒想到見面之後卻碰了一鼻子灰:女兒不僅沒有一點驚喜,反而滿肚子怨氣,責怪父母為什麼不經她同意就去南京,對她極不尊重.媽媽說我們想念寶貝女兒了,可女兒說天天微信,還有什麼好想的,你們的感情也太氾濫了.爸爸說媽媽第一次到南京,叫女兒陪去轉轉,女兒回答說和同學約好了,要去蘇州玩,讓爸爸媽媽自己玩,然後就匆忙而去……夫婦倆在南京轉了一天,越想越不是滋味,索性搭高鐵返回廣州.老友在電話中問我:龍兄,你說我錯在哪裡?我不知怎麼回答……(\url{https://kknews.cc/zh-hk/psychology/2jvo9g.html})

那些父母問:到底我做錯甚麼,以至你如此厭煩我?耶和華上帝也問祂的選民類似的問題,祂的選民如何回答呢?各位旁聽嘉賓,讓我們來聽聽被告的答辯:

3．以色列的自辯式回應

6:6我朝見耶和華,在至高上帝面前跪拜,當獻上甚麼呢?豈可獻一歲的牛犢為燔祭嗎?

6:7耶和華豈喜悅千千的公羊,或是萬萬的油河嗎?我豈可為自己的罪過獻我的長子嗎?為心中的罪惡獻我身所生的嗎?

初聽,選民似乎是答非所問.細心嘴嚼,選民的答辯其實是連消帶打.或者,聽眾可以這樣解讀:

我對你還做得不夠嗎?你還要求甚麼呢?你的要求是不是太過分了?我每逢大時大節,初一十五都依足律例獻祭,難道還不夠嗎?你應該知道,出生七日的牛犢就可以用來獻燔祭,難道你要我獻一歲大的小牛嗎?你知道一歲大的比七日大的價錢要貴多少嗎?好,如果你要一歲的小牛才高興的話,我就獻一歲的小牛.你應該知道,獻壇上的燈火不滅就可以了,難道你要我獻一千隻公羊,燒牠們的油好像流水一樣流動,你才滿意嗎?你知道這是皇帝登基時才會這樣做的,難道日日都是皇帝登基嗎?你當想清楚自己的要求是否太過分?會否過分得好像那些外邦人的神一樣,要求他們的信徒殺死自己的兒子作為祭祀.難道你要我們學習住在我們附近的亞瑪力人向他們的神摩洛獻兒童火祭一樣嗎?

耶和華的選民認為自己已經按要求做足了,如果耶和華還嫌不足,就實在太過分了.不知道現代的信徒有沒有類似的感覺呢?我已經恆常返教會崇拜,三餐前祈禱,做足奉獻,如果這樣還不夠,上帝的要求會否有點過分呢?難道除了星期日返崇拜,還要參加星期一晚的神學課程,星期二晚的佈道隊,星期三晚的祈禱會,星期四晚的門徒訓練,星期五晚的團契,以及星期六下午的詩班嗎?除了十一奉獻,難道還要差傳奉獻,擴堂奉獻,辦學奉獻,維修基金奉獻嗎?如果上帝的要求果真是這樣,難怪家中拜祖先的媽媽會說:「阿女,信教還信教,唔好咁迷啊!」

4．先知的回答

聽過控辯雙方的發言,傍聽的你有結論了嗎?或者,你需要時間先想一想,現在,由先知來做結案陳詞:

6:8世人哪,耶和華已指示你何為善.

先知的回應,暗示選民對耶和華的誤會極大.其實耶和華的要求並不是越來越多,更不是越來越苛刻.其實耶和華所期望的,一早已講了,清楚而不繁雜的,正如有一本曾經受歡迎的書,名為“All I really need to know I learned in kindergarten”一樣,人生所應該知道的大原則,例如:友愛,尊重,分享,和平共處等等,我們都在幼稚園裡學過了.先知向以色列人澄清,其實你們上述的質問只是「發悔氣」,並非真正問題,因為你們早已知道答案.

他向你所要的是甚麼呢?只要你行公義,好憐憫存謙卑的心,與你的上帝同行.

來上主面前,並不需要獻上一歲的牛犢為燔祭；也不是千千的公羊和萬萬的油河；更不是為自己的罪過獻長子.何以對蒼天呢?如何面對上主呢?先知重提三方面的要求:第一,是公義.公義配的動詞是「行」,這是行動上的要求 .在舊約聖經中,公義這個字一方面描述上帝的屬性,有譯作公平,公義,公正；同時也與審判相關,包括審判的場景,過程,判決和執行等.因此,公義是社會層面的,是先知談論社會責任的關鍵詞.原來,耶和華期望選民對上主的委身包括對社會的委身.先知在當時重提公義,就是因為在耶和華的選民公義不彰,公平消失.

3:1我說:雅各的首領,以色列家的官長啊,你們要聽!你們不當知道公平嗎?

3:2你們惡善好惡,從人身上剝皮,從人骨頭上剔肉,

3:3吃我民的肉,剝他們的皮,打折他們的骨頭,分成塊子像要下鍋,又像釜中的肉.

3:9雅各家的首領,以色列家的官長啊,當聽我的話!你們厭惡公平,在一切事上屈枉正直；

3:10以人血建立錫安,以罪孽建造耶路撒冷.

3:11首領為賄賂行審判；祭司為雇價施訓誨；先知為銀錢行占卜.

如果我們仔細聽,就會聽得出,先知在這裡所指責的,是首領和官長.理由很簡單,因為首領和官長手握公共權力和公共資源,他們有責任公義地運用這些權力和資源,也有條件濫用.公義要求對社群中的弱者的責任,維護他們的權利,因為社會中的強者自有他們的優勢去保護自己及爭取他們的分,但社會中的弱者卻未必有這樣的能力和資源,若以法律訴訟所涉及的龐大費用就可理解.甚麼人才感受到不公不義的切身之痛呢?很可能不是強勢的人,除非他們感到自己落入弱勢的處境.18/1/2017有關於一位特首參選人的報導如此說:「建制派表示,近日陸續收到來自中聯辦信息,要求支持前政務司長林鄭月娥參選,新民黨田北辰昨形容如「無形之手」操控,令事態走樣.已宣布參選的新民黨主席葉劉淑儀昨說,希望中央給予公平,公正選舉,不用擔心「流選」發生.被問及特首梁振英昨晨稱讚已請辭的政務司長林鄭月娥有擔當,葉太更罕有哽咽,反問「唔通可以話我對香港社會無擔當咩?」」(明報)經歷如此狀況,或許,這位強勢人士會明白,為甚麼公義的考慮應該弱勢者的弱勢出發,而不是從強勢者的高處去看.

第二方面的要求是憐憫,相配的動詞是「好」.公義是外在行為的,憐憫是內心情感的.在舊約聖經中,憐憫也是用來形容上帝的屬性,描述上帝豐盛的, 廣大的,長久的慈愛.這個字在聖經中經常出現,意思包括:善良,慈愛,守約,忠誠,恩慈等.憐憫雖然是內心情感的,但仍然是群體性的,因為這是群體成員互相對待的態度,就正如耶和華對以色列那種「我不會放棄你」的愛.耶和華要求祂的選民以神人之間這種不離不棄的忠誠,延伸至群體之間的關顧與愛護.在彌迦書中,先知也以這個字作結束:

7:18上帝啊,有何神像你,赦免罪孽,饒恕你產業之餘民的罪過,不永遠懷怒,喜愛施恩?

7:19必再憐憫我們,將我們的罪孽踏在腳下,又將我們的一切罪投於深海.

7:20你必按古時起誓應許我們列祖的話,向雅各發誠實,向亞伯拉罕施慈愛.

三.總結

有人問,在這項爭議中,彌迦先知以行公義,好憐憫來回應以色列人關於祭祀的質詢,是否有意思把宗教代約為倫理呢?如果只停留在這理,段章取義,的確容易引起這樣的想法.然而,先知的結案陳詞並不停於止,在行公義,好憐憫的同時,先知接著說:「存謙卑的心,與你的上帝同行」.先知認為,行公義,好憐憫兩項道德倫理要求,背後的動力在於在三項.先知把以色列與上帝的伙伴關係用「與上帝同行」來描述,這是很有意思的.公義與憐憫並不是抽象的思辯,而是選民與上神同行的親身經驗,體現在具體的歷史之中.從這個角度,各位傍聽者可以理解在這場訴訟中,為甚麼耶和華的發言如下:

6:4我曾將你從埃及地領出來,從作奴僕之家救贖你；我也差遣摩西,亞倫,和米利暗在你前面行.

6:5我的百姓啊,你們當追念摩押王巴勒所設的謀和比珥的兒子巴蘭回答他的話,並你們從什亭到吉甲所遇見的事,好使你們知道耶和華公義的作為.

基於各種理由,對上帝的信仰從來都有被異化和形式化的危機,不獨在古代的以色列,也在近代的教會,因此,才有五百年前馬丁路德發起的宗教改革運動.馬可路德批評教會把許多儀節,風俗和傳統攙雜進入了信仰之中,模糊了上帝及其拯救,他起初確立了三個「唯獨」(即「唯獨聖經」,「唯獨恩典」及「唯獨信心」).然而,不要誤會馬可路德只關心個人得救,或者個人修為,他的關注也是社群性的,他在重要文件「基督徒的自由」中如此說:

一.基督徒是全然自由的眾人之主,不受任何人管轄.(這是關於救恩的)

二．基督徒是全然順服的眾人之僕,受任何人管轄.(這是關於社會倫理的)

對跟隨上帝的群體來說,宗教信仰和社會倫理是一體兩面,不能分割的.讓我們以耶穌基督的吩咐來結束:

馬太福音

5:23所以,你在祭壇上獻祭物的時候,若想起有弟兄對你懷恨,

5:24就要把祭物留在壇前,先去跟弟兄和好,然後來獻祭物.



\newpage

\section{活出麗亮人生}
\label{sec:352}
\textbf{許開明牧師 - 活出麗亮人生!}
\newline
\newline
\hyperref[sec:351]{< < < PREV SERMON < < <}
~
\hyperref[sec:toc]{[返目錄]}
~
\hyperref[sec:353]{> > > NEXT SERMON > > >}


講題:活出亮麗的人生　Live a Brilliant Life

講員:許開明牧師

場合:香港中文大學崇基學院禮拜堂

日期:2017年2月5日

引言

「活出亮麗的一生」可能是每一個對人生認真的人都想,探索,追求的.一般人可能認為「活出亮麗的一生」是得到金錢或財富,名利,權力,社會地位,幸福,健康,快樂等；對大學生而言.可能是得到高等的學位,畢業後得到理想的工作；對大學的老師而言,可能是得到教授的名銜或崇高的學術地位.

或者,懂得向成功的人士學習,讀名人傳記,更甚是聽從近代有些專家所講的「生涯規劃」,得到指點就可活出美好的一生,但這種把一生交給專家規劃人生的方式,我是有點存疑的,因為許多人都擁在美好的一生,但不是經過規劃的.在我而言,我的一生就是有上帝的規劃和引導.現代人以為得到一套達致美好人生的觀念,活出精彩的秘笈,就可活出好人生,所以,就有不少怪獸父母出現,要仔女讀名校,跟名師,讀完的後要賺錢,甚至連子女的伴侶都要替他們安排好.

一般人的看法是「如何獲得」或是「如何得到」.耶穌的教導又如何?耶穌對門徒說過,跟隨他的人,就可活出亮麗的人生,這正提出了一個不同的觀念,不是「如何獲得」或是「如何得到」而是「保持不失」 的思想.

在馬太福音5:15 - 20,耶穌認為一個跟隨主的人,他應該已擁有三樣人生的寶貝,如何保持不失?倘這三樣東西保持不失,他就能活出亮麗的生命.這三個寶貝是什麼呢?第一個是「世上之鹽」；第二個是「世上之光」；第三個是「世上之義」.要注意的是,耶穌不是說門徒要成為世界的鹽,而是門徒是世界的鹽；不是要成為世界的光,而是他們是世界的光；不是成為「世上之義」,而是他們是「世上之義」.所以門徒,即跟隨主的人,本身已擁有「世上之鹽」,「世上之光」,「世上之義」,最重要是保持他們,不要失去.

各位,美好的人生就是保持住,不要失去.

一個人最怕是「失去」,保持不失,是人生智慧.就如讀書人不想失分；生意人不想失敗；說話的人不想失言；年青貌美的人不想失去美麗儀容；見工的人不想失望；青年男女不想失戀；年紀大的人不想失健康；生活艱難的人不想失金錢.

可惜,面對人生實況,許多東西隨著時間總會「流失」,美貌會失去,健康會失去,工作的機會與能力會失去,銀行存款會失去,生命年日會流失,最後陪伴著自己的身體也會失去,朋友也會逐漸失去(1983年神學組我的同學,5位中已去3位(30多,40多,50多離去),只餘下2位)使人突然感到心寒.

這些失,導致我們失去亮麗人生.面對這些流失,失去,我們應珍惜一切日子,家人,親友,工作機會(我曾在望覺堂,香港閩南堂,全完堂,崇基禮拜堂工作),同學,同事等,每個人遇上的人,更要盡量不讓它們那麼快失去,要保持健康,保持關係,保持活力.

信息

第一樣:你們是「世上之鹽」,要「不失鹹味」

「你們是世上的鹽；鹽若失了味,怎能叫它再鹹呢?以後無用,不過丟在外面,被人踐踏了.」(太五13)

1鹽有純潔的意思和作用,不失純潔和聖潔

鹽與純潔有關.它白色的光彩,無疑地使這聯想很容易建立.羅馬人說:鹽是萬物中最潔淨的,因它是從最純淨的兩種東西-太陽與海水而來.鹽的確是首先奉獻給上帝的祭物,猶太人在一日的末了獻鹽為祭.因此基督徒若要作世上的鹽,他必須是純潔的榜樣.

2鹽能保持磚熱力,不失熱誠

在巴勒斯坦,普通人家的火爐是建在房子外面,用石塊砌成,底下鋪磚.這種火爐為要保持熱度,就在爐底的磚上放一層厚厚的鹽.經過相當長的時間,鹽全部被燒壞了,就拿起磚來,把鹽拿開,丟棄在爐外.鹽既已失去熱磚的能力,就把它丟棄了.所以鹽就代表人待人處事也要有熱情和熱誠.

3鹽有調味作用  ,不失去與人調和,調好的關係

鹽是調味用的,但是除非它和食物混合,否則一點用也沒有. 鹽調味的作用是鹽最廣大最明顯的特質,沒有鹽的食物就沒有滋味,甚至成了令人討厭的東西.

鹽也有上味的作用.一個有名廚師說:「沒有鹽,就不會有烹調藝術.從廚房把它清除出去,所有吃的情趣會隨之飛逝,它的存在給佳餚以魅力,沒有它,就使珍饈變得乏味,使食客大失所望.」此話一點也不假.因為在廣州,鹽為上味.

有趣的是,主耶穌一千多年前把門徒比方為「世上的鹽」.這說明基督徒在世上的地位,和他所應起的重要作用.

鹽的功能,使不可口變為可口,使淡而無味的人生變為生之樂趣.這須要我們作到以下兩點:一方面先保持有味的鹽；另一方面不要變成失味的鹽,也就是說要作一個與基督徒的名相稱的基督徒.

4鹽有殺菌防腐作用,不失對抗罪惡

古代的世界,鹽是最普通的防腐劑,用以保存東西免致腐壞,排除腐爛與朽壞.保持鹹味不是目的,而是保持了它的新鮮,有魚味,菜味,肉味.鹽保存了食物,使之免於腐朽.基督徒要作世上的鹽,也就必須對生命有防腐性的影響力,保持基督徒的味道.

5鹽象徵智慧,不失去人生的智慧!

猶太的拉比常用鹽來象徵智慧(參西四6),所以希臘語代表失去味道的字實際上又有「變得愚蠢」的意思.(耶穌無疑用的正是亞蘭語裡有這兩個意思的詞 ta-pe-l 一個愚蠢的門徒對世界是沒有影響的,所以基督徒要小心,你會失去鹽味,鹹味,就如你會失去智慧,因此要小心.

嚴格說來,純鹽是不可能失去鹹味的,可是從死海灘上掘出的不純「鹽」,一旦氯化鈉分解了,慢慢也就不鹹了.

保持鹽的鹹味,就是保持人性優良的本質,不失去純潔,不失去熱力,不失去調和及上味作用,不失去防腐作用,不失去智慧.這是多麼美麗的本質.

鹽是上好禮物.世界五大洲許多地方民間都有一種習俗,把鹽當做最崇高的禮品.碰到尊貴的賓客,舉行贈禮儀式的時候,禮物往往就是一小撮潔白的鹽.以蘇聯來說,按照古老民間傳統風俗的記載,也常常讀到以鹽作為珍貴禮品的記述.

人們在現實生活中,無論是工業上,醫藥上,以及日常生活上,鹽都是不可少的東西,那怕是山珍海味,飛潛動植,無論怎樣美好的佳餚,要是缺了一味鹽,那味道雖然並非完全難以下嚥,但是老實說,淡而無味就令人感到食欲不振.

你有想過把自己,作為上好的禮拜獻給天地的主嗎?

第二樣:你們是「世上之光」,不失去光輝

耶穌說我們是鹽和光,然而他警告我們要小心失了鹹味和鎖住了我們及上主的光芒(太五13-20).「是」 世上的光,不是「作」世上的光,而光要使人看見.

持燈者成一路燈

英國作家拉斯金(John Ruskin)與一位弟兄談:「怎樣做個基督徒?」他們坐在家裡,望見窗外遠處,有位點燈的人,手拿著火把,在山坡上點燈.當時天色昏暗,看不見點燈的人本身,但從那些燈火可知他的存在.拉斯金突然有感說:「做一個基督徒就是這樣,你可能看不見他,但從他的好行為,知道他的的存在.因為在他所經過的地方,都點著發燈,都發了光.」

「你們是世上的光；城造在山上,是不能隱藏的.」(太五14)

城造在山上,城在日間本身像光或在晚上發出光,讓尋找城的人看見.「鹽」是重在說到我們的性質如何,而「光」是重在說到我們的行為如何.基督徒的地位如同建造在山上的一座城,是顯明在世人眼前的,所以是不能隱藏的.

「人點燈,不放在斗底下,是放在燈臺上,就照亮一家的人.」(太五15)

人點燈, 放在燈臺上,是要高舉上帝的燈,上帝所發出來的光.以色列人自稱耶路撒冷是「外邦人的光」,有一位著名的拉比被稱為:「以色列的燈」.猶太人所使用的這種表達方式,使我們可以知道耶穌採用它的關鍵.有一件事是猶太人十分清楚的,沒有人能點燃他自己的燈.耶路撒冷的確是外邦人的光,但是上帝點燃了以色列的燈.無論是國家或上帝的僕人所發出來的光,都是借來的.基督徒的光也是如此.耶穌所要求我們的,並不是我們應當發出我們自己的光來.我們必須發出反映祂的光輝.基督徒所發的光,是因為基督他心裡而映照出來的亮光.

「人點燈,不放在斗底下,」十四節是說光「不能隱藏」；這裡是說光「不該隱藏」.「放在燈臺上」意即放在特別顯露的地位上,以照亮周圍.

「放在斗底下,」斗是木造的量器,用以量取糧食,故它象徵生活的掛慮；燈放在斗底下,意思是信徒被生活的問題所扣住,以致遮蔽了亮光,把自己放斗底(憂慮,焦慮中)下.「是放在燈臺上,」燈臺是放置燈的正當器物,它象徵見證主的生活.前節的「城造在山上」是指基督徒站在屬天,高超的地位上,本節的燈「放在燈臺上」(放正當位置),是指基督徒站在正常,合乎體統的地位上.「就照亮一家的人,」前節的「城」是給在城外的人看見的,本節的「燈」是給在家裡的人看見的,這表明基督徒的光,不單能照亮人的「外面」,還能照亮人的「裡面」.

1光輝最主要的是被人看見

巴勒斯坦的房屋十分黑暗,只有一個小小圓形的窗戶,可能直徑只有十八英吋,燈是船形,好像一隻醬油碟子那樣,裡面盛滿油和浮燈蕊.在火柴尚未發明的年代,要重燃一盞燈並非易事.

2光是一種指引.

在任何江河的港口,我們都可以看見有一列燈光在水道上作記號,使船隻得以安全的駛入,即使在城市的街道上,若沒有燈光,我們都知道將會遭遇多大的困難,光可以使路途清晰.

3光時常可以當作警戒主的記號.

當前面有危險的時候,它警告我們停止.能夠使人看得見的光,警戒人的光,引導人的光,就是基督徒所必須作的光.

為上帝發光(太五16)

這裡有兩件最重要的事,一是讓人們看見我們的好行為.在希臘文中有兩個字代表「好」,agathos單單是指一件東西有良好的質地；kalos的意思不單是一件東西好,而且也美麗迷人,引人入勝.這裡所用的「好」字,就是kalos.基督徒的善行,不僅是良善而已,它們必須要吸引人.而二,更要注意的是我們的好行為所引人注意的,並不是我們自己,而是上帝.

耶穌所說的這句,是要完全禁止有人所謂的「誇張性的良善」.穆迪(D. L. Moody)參加一次聚會,有幾位對信仰十分認真的青年人,也來參加聚會.一天晚上,他們舉行了一個通宵的禱告會,第二天早上,他們正要離開的時候,遇見了穆迪.他問他們在作甚麼?他們說明原委之後,接下去說:「穆迪先生,請看我們的臉是多麼光亮.」穆迪很有禮貌的回答說:「摩西並不知道他的臉在發光.」自己感覺到的良善,把注意力集中在本身的良善,這不是基督徒的良善.

第三樣:「世上之義」,不失去律法

「莫想我來要廢掉律法和先知.我來不是要廢掉,乃是要成全.」(太五17)

「廢掉」即拆毀敗壞；「成全」即實現,應驗,填滿,使之完全.「我來不是要廢掉,乃是要成全,」「成全」有兩方面的意思:一是指實現,應驗,根據舊約律法的要求,在主耶穌身上得著完滿的實現,舊約先知的預言,在祂身上也得著完滿的應驗；二是指填滿,使之完全:舊約的律法和先知,只不過是字句和影兒,若沒有主耶穌的到來,就仍有缺陷,還不夠完全,所以祂來填滿律法和先知,使之完全.

「都要成全,」十七節的「成全」和十八節的「成全」意義不同；十七節是指補充律法,使其完全；十八節是指成就律法,使律法中一切的要求都得著實現.

成全又有回到被造的目的,不失去原有目的的意思.換言之,希伯來文中對完全的觀念是功用性的.一件東西倘若完成了它被造的目的,就成為完全.其實這個意思是從這個字引伸出來.teleios是從名詞telos而形成的形容詞.telos意即終極,宗旨,目的,目標.假使一件東西實現了預計的目的,它就是teleios；一個人如果完成了受造及遣入世界的目的,他就是完全的.

我們來作一個簡單的比方,假定在我家裡有一個螺絲釘鬆了,我就出去到鐵器店買了一個螺絲鑽.我發現所買的螺絲鑽大小正適合我使用,不太大也不太小,不太粗也不太滑.

我將螺絲鑽放在螺絲釘承受外力的凹槽,自然也非常地合適,我將螺絲釘扭緊了.按希臘文,尤其新約聖經的意思,螺絲鑽就是teleios,因為它正好達成了我要買它的目的.所以一個人倘使成全了他受造的目的,他就是teleios.

「我告訴你們,你們的義,若不勝於文士和法利賽人的義,斷不能進天國.」(太五20)

「你們的義」是指新約信徒的義行.「文士和法利賽人的義」中的「文士」是當時猶太人中的舊約聖經教師,也就是律法專家,大多數是法利賽人(太三7註解),他們的義是指履行律法規條所得的義,也就是律法上的義(腓三6).本節的意思是說,我們即使是嚴格遵行了舊約律法的每一條例,所得的也不過是「文士和法利賽人的義」,仍然不能進天國.基督徒的義(啟十九8),乃是我們憑著神的生命,履行主耶穌所填滿並提高的律法(這誡命)而有的,所以勝過文士和法利賽人的義.

你們聽見有吩咐古人的話,說,『不可殺人,』又說,『凡殺人的,難免受審判.』只是我告訴你們,凡向弟兄動怒的,難免受審判.(太五21-22)

你們聽見有話說,『不可姦淫.』只是我告訴你們,凡看見婦女就動淫念的,這人心裡已經與他犯姦淫了.(太五27-28)

不失律法或自律來達到正義,這是從人與人間社群的關係去講的.英國經驗主義哲學家強調「存在就是被感知」.休謨(David Hume),提出一套重要觀念「存在的世界是一束的感覺,世界會激發人的情感感知」,他認為美德會帶來快樂的感受,惡德帶來痛苦的感覺.在社群中,不失「正義」,這是「人為的德行」,是美德,帶來「快樂」.筆者更發現這是「社群」間的快樂,「不是獨樂樂,而是眾樂樂」.

耶穌時代的宗教,在嚴謹的正統派猶太人看來,遵守千萬條律法的規則與條例就是事奉上帝了.十誡原是一切律法的基礎和本質,我們可以在其中看到它整個的意義能夠總括在一個字裡面──「尊重」,或者更進一步的「敬畏」.敬畏上帝和上帝的名,敬畏上帝的日子,尊重父母,尊重生命,尊重財產,尊重人格,尊重真理與別人的好名聲,尊重自己,使我們不會被錯誤的慾望所控制──這就是十誡的基本原則.

十誡的基本原則是敬畏上帝,尊重同胞和我們自己.沒有這基本的敬畏與尊重,就不可能有律法,一切的律法都是以此為基礎.耶穌來成全的就是這種敬畏和尊重,祂在實際的生活中指示人敬畏上帝,並尊重所有的人.希臘人說:公義在於將當得之份奉獻給上帝,施捨給人.耶穌來,在生活中指示世人,要將上帝所當受的敬畏獻給祂,並將人應受的尊重給予人.敬畏與尊重並不在於遵守一大堆瑣碎的規則與條例,並不在於祭祀,卻在於憐憫；並不在於律法,卻在於愛心；並不是在禁令中強迫人不做一些甚麼事情,而是在誡命中告訴他們,將他們的生命模塑成主動的,愛人的生命.

敬畏和尊重是十誡的基礎,是永遠也不能廢去的,它們是維護人與上帝以及人與人之間關係的永遠的要素.奧古斯丁(Augustine)說基督徒的生活可以總括在這一句話之內:「知上帝愛你,你愛上帝,作你喜歡作的事.」

當我們發現上帝曾如何愛我們,生命中唯一的願望,就是回報祂的大愛.

敬畏和尊重是十誡的基礎下,最緊要不要「失」,不失鹹味,不失光芒,不失律法/自律/義

活得像個人

漢娜鄂蘭(Hannah Arendt)是一名德籍猶太人,她目睹希特拉的惡行,屠殺猶太人及殺害基督徒.現代人要注意成為「一個人」,人應注意3項:

1勞動(labour):是應付肉體需要,生物性活動,製造生存所需:煮飯,洗衣,清潔.

2工作(work):是製作出工具,建築物,汽車,改善人的生活.

3行動 (action):透過言論進行政治的活動.進行社區的參與,社交的生活.人不只是勞動,工作,還要有行動的生活,才是一個人,不變成是機械人,不懂得思考,尋找意義；這無思想,無思考,會使人幹下惡行,失去人性,引來邪惡.

今天,民眾沒有行動,沒有思考,把一切都交給別人,就產生極權主義,讓獨裁者出現.

今日香港社會,政府及許多議員,未來特首,除了關注如何迎未來,作出創新的改進,改善,改革外,我想耶穌所言:不失鹹味,不失燈台,不失律法和正義也相當重要.以往香港曾經享有的自由,民主,三權分立是否失去了鹹味?我們如何維持?多少基督徒及教會不願,也不敢舉起上帝的燈,沒有在家裡發光,更未能在社會有美好的見證.多少人只關注勞動(生活所需)和工作(表現自我創作),堂會只關注多少人信主,成為Mage Church卻欠缺關注社會與政治的行動,放棄自已的人權,讓獨裁者有機可乘.

結語

一個人最怕是「失去」,漸失去健康,漸失去胃口,漸失去金錢,漸失去耐性,漸失去堅持,漸漸失去生命的一切.然而,無論如何,不要失去我們的人格,名聲,見證.

如果我們保持不失,

1不失鹹味:保持世上的鹽,不失去鹽美好本質；

2不失去燈台:保持世上的光,高舉上帝的燈；

3不失去律法:保持世上的義,敬畏上帝,尊重他人,

我們也不失去生命的美麗和光輝,自然會活出亮麗人生!

祈禱

祈求上主,保守我們不失鹽的美麗本質,不失照亮別人的光輝,不失敬畏上主,尊重他人的誡命,活出亮麗的人生.阿們.



\newpage

\section{君子務本本立而道生}
\label{sec:353}
\textbf{鄧瑞強博士 - 君子務本,本立而道生}
\newline
\newline
\hyperref[sec:352]{< < < PREV SERMON < < <}
~
\hyperref[sec:toc]{[返目錄]}
~
\hyperref[sec:354]{> > > NEXT SERMON > > >}


馬太福音 5:21-26
\newline
\begin{longtable}{cl}
\hline
\hline
章節 & 經文 (和合本修訂版)\\
\hline
5:21 & \begin{tabularx}{0.7\textwidth}{X} 「你們聽過有對古人說：『不可殺人』；『凡殺人的，必須受審判。』 \end{tabularx} \\ \\ \relax
5:22 & \begin{tabularx}{0.7\textwidth}{X} 但是我告訴你們：凡向弟兄動怒的，必須受審判；凡罵弟兄是廢物的，必須受議會的審判；凡罵弟兄是白痴的，必須遭受地獄的火。 \end{tabularx} \\ \\ \relax
5:23 & \begin{tabularx}{0.7\textwidth}{X} 所以，你在祭壇上獻祭物的時候，若想起有弟兄對你懷恨， \end{tabularx} \\ \\ \relax
5:24 & \begin{tabularx}{0.7\textwidth}{X} 就要把祭物留在壇前，先去跟弟兄和好，然後來獻祭物。 \end{tabularx} \\ \\ \relax
5:25 & \begin{tabularx}{0.7\textwidth}{X} 你同告你的冤家還在路上，就要趕快與他講和，免得他把你送交給法官，法官交給警衛，你就下在監裡了。 \end{tabularx} \\ \\ \relax
5:26 & \begin{tabularx}{0.7\textwidth}{X} 我實在告訴你，就是有一個大文錢還沒有還清，你也絕不能從那裡出來。」 \end{tabularx} \\ \\
[1ex]
\hline
\hline
\end{longtable}
講題:「君子務本,本立而道生」　Radical Ethics

經文:馬太福音5:21-26

講員:鄧瑞強博士

場合:香港中文大學崇基學院禮拜堂

日期:2017年2月12日

各位弟兄姊妹,早安.

先講一講今日講道的題目:「君子務本,本立而道生」.這句話出自《論語》第一章＜學而篇＞的第二句,是孔子的學生「有子」的話.「君子務本」,君子致力於更為根本的事情.只是表面上做些功夫,不足為君子.君子要為自己的行為找一更深的基礎.行為背後,是否有更深的義理?義理背後,是否有更深的人性基礎?人性背後,是否隱藏著天理?一步一步,尋索人生的底蘊.「君子務本」,就是說,一個努力止於至善的人,不會人云亦云,不會人做我做,他必然會為自己的行動,找一終極基礎.「本立而道生」,「根本」找到了,終極基礎找到了,按此而活,便能將大道澄現出來.「君子務本,本立而道生」這句話若套用在基督徒的生命上,則可說是:基督徒的生命,不能隨波逐流,人云亦云,必須將生命堅立在神聖的基礎上.如此,生命的踐行,就是將真理的大道顯明出來.

練武的人說:「練拳不練功,到老一場空」,意謂單單練習外在的套路架式,忽略功力的內涵修為,則只會耍出花拳繡腿,中看不中用.有的只是外表的虛招,缺乏厚實的生命勁度.只練拳而不練功的人,到了老時,氣虛力弱,拳也打不出來了.練拳又練功的人,是將拳意收攝入整體生命裡,不再是用手用腳去打拳,而是身心氣神總體地打拳.年紀越大,功力越深.

我們是如何作基督徒的?

只有外在的花拳繡腿,抑或有深厚的靈性修為?遇到困難,便氣虛力弱；抑或遇到困難,更顯得剛勁厚實?只滿足於宗教上的枝節,抑竭力追求靈性的根本?信了一生,是到老一場空,抑越發將生命的豐盛活出來?

今日的講道經文,是馬太福音裡登山寶訓中的一段經文.馬太福音的登山寶訓,是講述基督徒生命應有的特質.這些生命特質,不是死守一些道德規條,而是透露天國的氣息.我們看看主耶穌如何展示天國子民的生命.

講道經文:馬太福音5:21－26

太5:21　「你們聽見有吩咐古人的話,說:『不可殺人』；又說:『凡殺人的難免受審判.』

太5:22　只是我告訴你們:凡〔有古卷在凡字下加:無緣無故地〕向弟兄動怒的,難免受審斷；凡罵弟兄是拉加的,難免公會的審斷；凡罵弟兄是魔利的,難免地獄的火.

太5:23　所以,你在祭壇上獻禮物的時候,若想起弟兄向你懷怨,

太5:24　就把禮物留在壇前,先去同弟兄和好,然後來獻禮物.

太5:25　你同告你的對頭還在路上,就趕緊與他和息,恐怕他把你送給審判官,審判官交付衙役,你就下在監裡了.

太5:26　我實在告訴你,若有一文錢沒有還清,你斷不能從那裡出來.」

這段經文自然地分成三個段落.

(1)引述古人的傳統

太5:21　「你們聽見有吩咐古人的話,說:『不可殺人』；又說:『凡殺人的難免受審判.』」

對這簡單的經文,我們有四點觀察.

一,我們活在古人的傳統裡.無論你生在何方,你總活在傳統裡.對這些傳統,我們可能不聞不問.對傳統的要求,我們可能習以為常,照做可也.

二,傳統告訴當時的人,不可殺人.「不可殺人」是十誡中的一誡,猶太人自少便習慣了.這禁令是負面的,只告訴我們「不可」作什麼,沒有正面的要求.

三,「凡殺人的難免受審判」,告訴我們,違反傳統的禁令是有後遺症的,你要為此付上代價.

四,第四個觀察點是在中文翻譯中看不到的.希臘文的動詞,有時態(tense,如:現在,過去),有語態(voice,如:主動,被動),還有語氣(mood,如:直說,命令),這全都以動詞的語尾去表達.這裡的「不可殺人」,意義是命令式的,但希臘文的用字,在語氣上卻是「直說語氣」(indicative mood).這似乎暗示,在耶穌眼中,這禁令是重要的,但仍不是命令的根本.

(2)殺人的根本原因

太5:22　只是我告訴你們:凡〔有古卷在凡字下加:無緣無故地〕向弟兄動怒的,難免受審斷；凡罵弟兄是拉加〔即:廢柴〕的,難免公會的審斷；凡罵弟兄是魔利〔即:蠢材〕的,難免地獄的火.

對這簡單的經文,我們有五點觀察.

一,主耶穌沒有直接討論「殺人」這主題,他亦沒有按傳統人云亦云.他討論殺人背後的根源:怒氣.

調查凶殺案,通常要找出凶手背後的動機.沒有動機而殺人,是很少有的.不同動機背後,會否有共同元素?或許,怒氣是隱伏在不同動機背後的共同元素.殺人的「本」,在於怒氣.

二,主耶穌的重點不在行為,而在人的生命美德.如何使人「不殺人」,可以有很多方法.嚴刑峻法,斬頭示眾,天羅地網式的監控,都可威嚇人或阻止人不去殺人.主耶穌的重點不在這裡.主耶穌的重點,在人的質素.好樹會結好果子,壞樹會結壞果子.主耶穌要我們正視,我們到底是一棵什麼樣的樹.我們的本性,是流露美善,抑或流露忿恨.

三,忿怒表現為辱罵:你是「廢柴」,你是「死蠢」,你「唔死都無用」.罵人,是羞辱他,低貶他,漠視他的尊嚴.忿怒下的謾罵,是可以沒有底線的.當你認為別人「大逆不道」,「死不足惜」時,殺他一個痛快,便正常不過了.

四,忿怒的後果,越來越嚴重.最初是「受審斷」,跟著是「受公會的審斷」,最後是「難免地獄的火」.從後果來看,是越來越嚴重的.看來,若不好好處理,怒氣有一種變本加厲的傾向.像進入一向下的漩渦,越墮越深.

五,這句經文本身也不是一個禁令.主耶穌沒有說「不要動怒」.若說了,我們也不可能做到.主耶穌像醫生,不是說「不可以病」,而試探究病根,指出病因.讓我們看到,若不好好處理這病根,則後果嚴重.

(3)根治忿怒

太5:23　所以,你在祭壇上獻禮物的時候,若想起弟兄向你懷怨,

太5:24　就把禮物留在壇前,先去同弟兄和好,然後來獻禮物.

太5:25　你同告你的對頭還在路上,就趕緊與他和息,恐怕他把你送給審判官,審判官交付衙役,你就下在監裡了.

太5:26　我實在告訴你,若有一文錢沒有還清,你斷不能從那裡出來.」

殺人的根,在動怒.要克服怒氣,必須修正整個生命人格.這和中醫的想法相似,不是醫病,而是醫人.要修正人格,重點不在避免什麼,不可什麼,而在主動做些什麼,建立什麼.這段經文教我們如何建立正面的生命.

對這段經文,我們有四點觀察.

一,留意23及24兩節.你獻祭時,或崇拜時,若想起有人不滿你,你應先主動與他和解,才來獻祭.經文告訴我們,信仰與我們對待他人的方式或我們展示人格的方式,是息息相關的.禮拜日的崇拜和禮拜一至六的生活,是息息相關的.不要將信仰與平時的生活分開.

二,要主動復和.主動到一個地步,不是我嬲你才復和,而是你嬲我時,我已採取主動行動去復和.不殺人,是表面的.更深是處理怒氣.處理自己的怒氣,很好,但不夠主動,更主動的,是處理別人的不滿,將怒氣的根消除.做到這樣,需要的是一創造復和的人格,不惜多走一里的器量.

三,這在中文翻譯中又是看不到的.在這一大番話中,第一次出現「命令語氣」(imperative mood)的動詞,是出現在第24節裡,並且一出便出現四次.在第24節裡,就把禮物「留」(命令語氣)在壇前,先「去」(命令語氣)同弟兄「和好」(命令語氣),然後來「獻」(命令語氣)禮物.主耶穌不命令我們不去做什麼,這太負面.不許這樣,不許那樣,惹人不滿.主耶穌命令人活出正面的人格,去主動復和.我們不是死守律法者,我們是創造生機者.在隱伏怒氣的環境裡,在殺機四伏的關係裡,是你的責任,去創造復和.是你的責任,以和平之力去取代戰爭.是你的責任,去建造愛的關係.這不是與信仰無關的,這就是信仰的核心.

四,最後兩節經文:25及26兩節.話題一轉,視我們為有債未還的人.要把握機會,和追債的人達成協議.若有一文錢沒有還清,則不能離開監牢.當然,這段經文可照字面解,即:若虧欠人,則要及早償還.這是復和之道,止息忿怒.但是,也可將這段經文,視為神對我們的虧欠的追討.我們有責任創造復和,我們有責任活出信仰的人格,我們有責任展示主所講的復和的世界是一個真實的世界.我們對這一切,負有責任.但我們有虧欠,處處不盡責.主要追討我們.

你可能會問:我們能擔起這責任嗎?我們真能活出這生命人格嗎?我們真能展示復和的力量嗎?

要回答這問題,我們要思想得大一點.

登山寶訓的經文,是置放在主耶穌所講的天國之下的.主耶穌已開展天國.天國,是我們不習慣的新世界,是一個看來不可能的世界.在這裡,「虛心的人有福了!因為天國是他們的.」在心中不戀棧財富而關懷貧窮的人,是有福的,天國是他們的.「哀慟的人有福了!因為他們必得安慰.」心中思念神的事情,而為這世界的崩壞而哀傷的人,是有福的,因為他們必得安慰.「使人和睦的人有福了!因為他們必稱為神的兒子.」努力實現和平及建設圓滿生命的人,是有福的,因為他們必稱為神的兒子.「為義受逼迫的人有福了!因為天國是他們的.」為天國的正義而受苦的人,是有福的,因為天國是他們的.

為何主耶穌講得那麼「實牙實齒」,因為天國已然在祂身上開展出來.如此,我們真能在充滿怒氣的人間創造復和嗎?答案是肯定的.若主耶穌已用大愛奠定天國的復和基礎,我們怎能說復和的努力是軟弱無力而不會有成果的呢?

(4)結語

不要單單滿足於禮拜日「返教會」,要在日常生活中展示你是屬神的人.

在可能引發怒氣的地方,主動一點,創造復和.不要視這為「蝕底」,視之為有分參與天國的福氣.

要克服凶殺,必先克服忿怒.要克服忿怒,必需具備創造復和的人格.要有這人格,必須看到天國的真實,必須甘心跟從主耶穌,必須倚靠恩典而活.

希望你是一個甘心跟從主耶穌的人.

但願:榮耀歸於聖父,聖子,聖靈.阿們.



\newpage

\section{說之間}
\label{sec:354}
\textbf{梁元生院長 - 說「之間」}
\newline
\newline
\hyperref[sec:353]{< < < PREV SERMON < < <}
~
\hyperref[sec:toc]{[返目錄]}
~
\hyperref[sec:355]{> > > NEXT SERMON > > >}


利未記 19:1-2,9-18
\newline
\begin{longtable}{cl}
\hline
\hline
章節 & 經文 (和合本修訂版)\\
\hline
19:1 & \begin{tabularx}{0.7\textwidth}{X} 耶和華吩咐摩西說： \end{tabularx} \\ \\ \relax
19:2 & \begin{tabularx}{0.7\textwidth}{X} 「你要吩咐以色列全會眾，對他們說：你們要成為聖，因為我耶和華－你們的神是神聖的。 \end{tabularx} \\ \\ \relax
19:3 & \begin{tabularx}{0.7\textwidth}{X} 你們各人都當孝敬父母，也要守我的安息日。我是耶和華－你們的神。 \end{tabularx} \\ \\ \relax
19:4 & \begin{tabularx}{0.7\textwidth}{X} 你們不可轉向虛無的神明，也不可為自己鑄造神像。我是耶和華－你們的神。 \end{tabularx} \\ \\ \relax
19:5 & \begin{tabularx}{0.7\textwidth}{X} 「你們宰殺祭牲獻平安祭給耶和華的時候，要獻得使你們可蒙悅納。 \end{tabularx} \\ \\ \relax
19:6 & \begin{tabularx}{0.7\textwidth}{X} 這祭物要在獻的當天或第二天吃；若有剩到第三天的，就要用火焚燒。 \end{tabularx} \\ \\ \relax
19:7 & \begin{tabularx}{0.7\textwidth}{X} 第三天若再吃，這祭物是不潔淨的，必不蒙悅納。 \end{tabularx} \\ \\ \relax
19:8 & \begin{tabularx}{0.7\textwidth}{X} 吃的人必擔當自己的罪孽，因為他褻瀆了耶和華的聖物，這人必從百姓中剪除。 \end{tabularx} \\ \\ \relax
19:9 & \begin{tabularx}{0.7\textwidth}{X} 「你們在自己的地收割莊稼時，不可割盡田的角落，也不可拾取莊稼所掉落的。 \end{tabularx} \\ \\ \relax
19:10 & \begin{tabularx}{0.7\textwidth}{X} 不可摘盡葡萄園的葡萄，也不可拾取葡萄園中掉落的葡萄，要把它們留給窮人和寄居的。我是耶和華－你們的神。 \end{tabularx} \\ \\ \relax
19:11 & \begin{tabularx}{0.7\textwidth}{X} 「你們不可偷盜，不可欺騙，也不可彼此說謊。 \end{tabularx} \\ \\ \relax
19:12 & \begin{tabularx}{0.7\textwidth}{X} 不可指著我的名起假誓，褻瀆你神的名。我是耶和華。 \end{tabularx} \\ \\ \relax
19:13 & \begin{tabularx}{0.7\textwidth}{X} 「不可欺壓你的鄰舍，也不可偷盜。雇工的工錢不可在你那裡過夜，留到早晨。 \end{tabularx} \\ \\ \relax
19:14 & \begin{tabularx}{0.7\textwidth}{X} 不可咒罵聾子，也不可將絆腳石放在盲人面前。你要敬畏你的神。我是耶和華。 \end{tabularx} \\ \\ \relax
19:15 & \begin{tabularx}{0.7\textwidth}{X} 「你們審判的時候，不可不公正；不可偏護貧窮人，也不可看重有權勢人的臉，總要公平審判你的鄰舍。 \end{tabularx} \\ \\ \relax
19:16 & \begin{tabularx}{0.7\textwidth}{X} 不可在百姓中到處搬弄是非，不可陷害鄰舍的性命。我是耶和華。 \end{tabularx} \\ \\ \relax
19:17 & \begin{tabularx}{0.7\textwidth}{X} 「不可心裡恨你的弟兄；要指摘你的鄰舍，免得因他承擔罪過。 \end{tabularx} \\ \\ \relax
19:18 & \begin{tabularx}{0.7\textwidth}{X} 不可報仇，也不可埋怨你本國的子民。你要愛鄰如己。我是耶和華。 \end{tabularx} \\ \\
[1ex]
\hline
\hline
\end{longtable}
詩篇 119:33-40
\newline
\begin{longtable}{cl}
\hline
\hline
章節 & 經文 (和合本修訂版)\\
\hline
119:33 & \begin{tabularx}{0.7\textwidth}{X} 耶和華啊，求你將你的律例指教我，我必遵守到底！ \end{tabularx} \\ \\ \relax
119:34 & \begin{tabularx}{0.7\textwidth}{X} 求你賜我悟性，我就遵守你的律法，且要一心遵守。 \end{tabularx} \\ \\ \relax
119:35 & \begin{tabularx}{0.7\textwidth}{X} 求你叫我遵行你的命令，因為這是我所喜愛的。 \end{tabularx} \\ \\ \relax
119:36 & \begin{tabularx}{0.7\textwidth}{X} 求你使我的心趨向你的法度，不趨向不義之財。 \end{tabularx} \\ \\ \relax
119:37 & \begin{tabularx}{0.7\textwidth}{X} 求你叫我轉眼不看虛假，使我活在你的道路中。 \end{tabularx} \\ \\ \relax
119:38 & \begin{tabularx}{0.7\textwidth}{X} 求你向敬畏你的僕人堅守你的話！ \end{tabularx} \\ \\ \relax
119:39 & \begin{tabularx}{0.7\textwidth}{X} 求你使我所懼怕的羞辱遠離我，因你的典章本為美。 \end{tabularx} \\ \\ \relax
119:40 & \begin{tabularx}{0.7\textwidth}{X} 看哪，我切慕你的訓詞，求你因你的公義賜我生命！ \end{tabularx} \\ \\
[1ex]
\hline
\hline
\end{longtable}
哥林多前書 3:10-11,16-23
\newline
\begin{longtable}{cl}
\hline
\hline
章節 & 經文 (和合本修訂版)\\
\hline
3:10 & \begin{tabularx}{0.7\textwidth}{X} 我照神所給我的恩典，好像一個聰明的工頭，立好了根基，別人在上面建造；只是各人要謹慎怎樣在上面建造。 \end{tabularx} \\ \\ \relax
3:11 & \begin{tabularx}{0.7\textwidth}{X} 因為，那已經立好的根基就是耶穌基督，此外沒有人能立別的根基。 \end{tabularx} \\ \\ \relax
3:12 & \begin{tabularx}{0.7\textwidth}{X} 若有人用金銀、寶石，草木、禾秸，在這根基上建造， \end{tabularx} \\ \\ \relax
3:13 & \begin{tabularx}{0.7\textwidth}{X} 各人的工程必將顯露，因為那日子要將它顯明，有火把它暴露出來，這火要試煉各人的工程怎樣。 \end{tabularx} \\ \\ \relax
3:14 & \begin{tabularx}{0.7\textwidth}{X} 人在那根基上所建造的工程若能保得住，他將要得賞賜。 \end{tabularx} \\ \\ \relax
3:15 & \begin{tabularx}{0.7\textwidth}{X} 人的工程若被燒了，他將損失，雖然他自己將得救，卻要像從火裡經過一樣。 \end{tabularx} \\ \\ \relax
3:16 & \begin{tabularx}{0.7\textwidth}{X} 難道不知你們是神的殿，神的靈住在你們裡面嗎？ \end{tabularx} \\ \\ \relax
3:17 & \begin{tabularx}{0.7\textwidth}{X} 若有人毀壞神的殿，神一定要毀滅那人；因為神的殿是神聖的，這殿就是你們。 \end{tabularx} \\ \\ \relax
3:18 & \begin{tabularx}{0.7\textwidth}{X} 誰都不可自欺。你們中間若有人自以為在今世有智慧，倒不如變為愚拙，好成為有智慧的。 \end{tabularx} \\ \\ \relax
3:19 & \begin{tabularx}{0.7\textwidth}{X} 因為這世界的智慧在神看來是愚拙的。如經上記著：「主使有智慧的人中了自己的詭計；」 \end{tabularx} \\ \\ \relax
3:20 & \begin{tabularx}{0.7\textwidth}{X} 又說：「主知道智慧人的意念，因為它們是虛妄的。」 \end{tabularx} \\ \\ \relax
3:21 & \begin{tabularx}{0.7\textwidth}{X} 所以，無論誰都不可誇耀人；因為萬有都是你們的， \end{tabularx} \\ \\ \relax
3:22 & \begin{tabularx}{0.7\textwidth}{X} 或保羅，或亞波羅，或磯法，或世界，或生，或死，或現今的事，或將來的事，全是你們的， \end{tabularx} \\ \\ \relax
3:23 & \begin{tabularx}{0.7\textwidth}{X} 而你們是屬基督的，基督是屬神的。 \end{tabularx} \\ \\
[1ex]
\hline
\hline
\end{longtable}
馬太福音 5:38-48
\newline
\begin{longtable}{cl}
\hline
\hline
章節 & 經文 (和合本修訂版)\\
\hline
5:38 & \begin{tabularx}{0.7\textwidth}{X} 「你們聽過有話說：『以眼還眼，以牙還牙。』 \end{tabularx} \\ \\ \relax
5:39 & \begin{tabularx}{0.7\textwidth}{X} 但是我告訴你們：不要與惡人作對。有人打你的右臉，連另一邊也轉過去由他打。 \end{tabularx} \\ \\ \relax
5:40 & \begin{tabularx}{0.7\textwidth}{X} 有人想要告你，要拿你的裡衣，連外衣也由他拿去。 \end{tabularx} \\ \\ \relax
5:41 & \begin{tabularx}{0.7\textwidth}{X} 有人強迫你走一里路，你就跟他走二里。 \end{tabularx} \\ \\ \relax
5:42 & \begin{tabularx}{0.7\textwidth}{X} 有求你的，就給他；有向你借貸的，不可推辭。」 \end{tabularx} \\ \\ \relax
5:43 & \begin{tabularx}{0.7\textwidth}{X} 「你們聽過有話說：『要愛你的鄰舍，恨你的仇敵。』 \end{tabularx} \\ \\ \relax
5:44 & \begin{tabularx}{0.7\textwidth}{X} 但是我告訴你們：要愛你們的仇敵，為那迫害你們的禱告。 \end{tabularx} \\ \\ \relax
5:45 & \begin{tabularx}{0.7\textwidth}{X} 這樣，你們就可以作天父的兒女了。因為他叫太陽照好人，也照壞人；降雨給義人，也給不義的人。 \end{tabularx} \\ \\ \relax
5:46 & \begin{tabularx}{0.7\textwidth}{X} 你們若只愛那愛你們的人，有甚麼賞賜呢？就是稅吏不也是這樣做嗎？ \end{tabularx} \\ \\ \relax
5:47 & \begin{tabularx}{0.7\textwidth}{X} 你們若只請你弟兄的安，有甚麼比別人強呢？就是外邦人不也是這樣做嗎？ \end{tabularx} \\ \\ \relax
5:48 & \begin{tabularx}{0.7\textwidth}{X} 所以，你們要完全，如同你們的天父是完全的。」 \end{tabularx} \\ \\
[1ex]
\hline
\hline
\end{longtable}
講題:說「之間」　Reflections on In-betweenness

經文:利未記19:1-2,9-18;詩篇119:33-40;哥林多前書3:10-11,16-23;馬太福音5:38-48

講員:梁元生院長

場合:香港中文大學崇基學院禮拜堂

日期:2017年2月19日

前言:無間,中間,之間

首先,我想說幾句感恩的話.五十年前,我第一次站上這個講臺,作為進入崇基學院的新生,那時我只有十八歲；今天,我六十七歲,已是近望七之年,行將退休,能夠仍然有機會站在這個講臺上, 和大家分享主道及共同崇拜,實在是主的大莫大恩典.故此在此先獻上感謝.

剛才讀經的弟兄姊妹給我們讀出的幾段經文,主要談及舊約時期猶太人和新約時期的基督徒對傳統律法的看法和態度.本來,我對舊約中猶太人的典章律例, 總覺得繁文縟節,不敢多碰；但經過幾個星期的思考,發現一點:就是在這些經文中可以帶出幾個與今日基督徒相關的問題.那就是:我們如何對待文化傳統?我們如何了解歷史和過去,以及我們如何看待今日和明天?這些重大的問題,似乎都從這幾段經文隱隱地透視出來,提供我們作深入的思考和想象.今天讓我們就嘗試以這幾段談律法的經文為基礎,從文化與歷史,傳統與現代,法律與恩典之間的關係,思考一下今日基督徒的處境,身份和選擇等種種問題.進入正題之前,讓我先提出三個關鍵詞,引導下面的討論:它們是:中間,之間,和無間.

讓我先說無間.

無間與中間:從傳統到現代

無間即一體.毫無縫隙,渾然一體的一體.我們指的是一個封閉的自我完整的文化體系或價值系統.當我們生活在一個同質的文化之中,完全浸濡在同一文化,同一的價值系統裡面,而我們的生活和視野,都在一個渾然無間的整體之中,自成一宇宙, 看不見外面的不同的新世界, 或者是境外之天.這就是我說的一體.一般來說,傳統的中國和猶太就是這樣的世界.我們今日所讀的經課,把猶太人傳統說的很清楚.如果我們來看整部利未記,就能更了解猶太人傳統的細節,他們生活得一舉一動,一出一入,都有非常清晰的法則和律例, 這和舊時的中國人的社會,有點像異曲同工,生活都是非常規範化和禮儀化的,待人處事,皆有嚴格的槼條和律法,整個天地合整個文化,皆井然有序,構成一個強大而完整的體系,一切都有規可循,有理可依,無容置啄.這是古時中國和猶太的傳統.在這樣的文化氛圍下,沒有太多的選擇和空間,可以提供你去選擇,即是沒有中間,沒有他途,沒有其他的alternative, 這種情況,是為無間.

那麼,甚麼是中間?大體而言, 在傳統中國和傳統猶太社會之中,走在中間的人不多.因為大部分的人,生活在一個較為嚴謹的,封閉的文化和價值系統裡,思想行為會較為保守, 規行矩步,社會交際和待人接物,都遵守律法的規範和經典的限制.在儒家社會而言,不是兩個不同的選擇. 所追求的不是我者和他者的區分,因為他者不容置疑的受到排拒和隔離,不是一種選擇；在儒家傳統之內, 希望做到的是恰倒好處,是了解「過和不及」的程度差異.所謂中庸之道,非選擇和方向的不同,只是處理是拿捏力度和準繩有別而已.這大概就是上個主日鄧瑞強博士所說的「生於傳統之中」(born into tradition)的意思吧!

而我們今日所處的現代社會,則有明顯的不同.因為現代社會是一個較為開放的,多變的,和邊界不明顯或不清楚的(甚至是無邊界的).雖然,我們之中不乏追隨傳統的人,也不缺激進創新的人, 但更多人可能覺得像是永遠的被夾在中間:生活在傳統與現代,中國與西方,激進與保守的夾縫中間.(在今時今日的香港,也有許多人覺得自己生活在泛民與建制,黃色和藍色的中間.)而中間之路並不寬敞,是條狹窄之路,甚至是條窄縫之路,予人很大的壓迫感.活在中間的感覺是怎樣的?我在許多年前寫過一首小詩,叫做「中間道」, 後來崇基學院音樂系的盧厚敏博士把它譜成曲子,成了一首可以唱出來的歌.以下是「中間道」的歌詞:

我從中間走來,走過崎嶇和顛簸,沒有左邊和右邊的摻扶,只有夾縫中的壓迫.

我從中間走來,走過荊棘和尖石；沒有向左或向右去逢源,只承受冷漠的臉色.

我從中間走來,走過了無限須臾,雙方只叫我不斷的等待,拖延中還夾雜揶揄.

我從中間走來,走過黑暗的渦漩,只因有多种顏色的襯托,竟顯得格外的軒轅.

生活在兩種龐大力量的中間,生活在兩個不同的文化系統的中間,你會有怎樣的感覺?是否感到兩邊拉扯的壓力?是否感到生活在狹窄的夾縫中間?心靈上,精神上,和肩頭上是否感到沉重的負荷,甚至叫你透不過氣?抑或你會感覺良好,感到可以有更多的選擇,甚至可以左右逢源?

中間, the middle way,永遠是兩個不同領域的分界,亦永遠只會是兩個系統的邊沿,而邊沿總是狹小的,充滿張力的,甚至受人歧視和受人排斥的空間. 中國傳統的描述,會叫這些邊沿人做「蠻,夷,戎,狄」, 也會叫他做「怪」,做「賤」,做「鄙」, 做「邪」, 以示主流與末流,核心和邊沿的區別.但我們今日時代不同了,近百年來我們走向現代,學習西方.漸漸地整個社會文化的主流價值觀有了一個嶄新的改變,相對形成了一個以西方文化和現代倫理構成的體系, 而在這個體系的邊沿人,可能是跟不上改變步伐的傳統文化人,也可能是價值觀不相同的一批異見者,在現代的新世界裡,他們也受到排擠和壓迫,同樣是被邊沿化.中間道永遠是狹窄的,是在夾縫裡生活的人.我想說:在中國人的社會裡,基督徒是整個社會和國家從傳統走向現代,東方逐漸接受乃至擁抱西方過程中的產物,中國的基督徒,往往是「中間道」的行者,既受傳統的衛道之士排斥和批評,也不願完全接受西方的洗禮.

總的而言,中間即是兩邊受壓,也是兩方的邊沿, 所處的空間非常狹小,生命就如墮入永恆的張力,被拉緊,被排斥,被壓迫.夾在中間是苦事,兩頭不是最難為.令我想起很多年前一則航空公司的廣告:你要旅行舒適,乘搭本航機,你就永遠不會夾在中間!

既然中間道有其限制和苦處,今天我希望和大家分享的是:走出「中間」之苦,化成「之間」之樂；或曰:走出中間道,做個之間人!中間只有一線,夾在兩個龐大勢力和龐大系統之間,難以舒展,而之間是中間的擴大和拓闊,把原來是中間的夾縫向兩邊伸展,變成寬闊的道路,化壓力為舒泰,變緊張作祥和,本來感到的是緊張和苦惱,現在感到的是自由和愉悅.此即為「之間」的境況.

以下讓我們來從今日的經文看三個「之間」,從中吸取教訓,如何做個「之間人」.

作為「之間人」的基督徒

如何做個「之間人」?「之間」是關係,是互動.人不能傲然孤立,離群獨處.他的生活,甚至說他的生存,依賴著與不同的「他者」所建立的關係.現在讓我們根據今天所讀的經文,思考下列三個「之間」的問題:

是我者與他者之間,第二,是小群與大我之間,第三,是律法與恩典之間.

一．園裡園外:我者與他者之間

讓我們來看利未記十九章9節.那裡提到種田的人的兩個比喻.一是在田裡種禾或者種麥,收成時期都長出穗子來,田主或農夫遂可以收割.另外是在栽種的人在葡萄園裡摘葡萄,滿載而歸.這些舊約所講的比喻,讓我們想到新約時耶穌講過的栽種及葡萄園的故事. 無論是田主或園主, 他們都是這田地和果園的主人,割禾及採摘葡萄都是完全合情合法的. 這裡,我想請大家注意兩種人:第一種人,就是「我者」.從現代的角度看,是指基督徒來說的, 亦即是我等,咱們.基督徒在上帝國裡享有兒子的身份,亦即是田主和園主,他有收割禾穗的權利,有採摘葡萄的權利.我們基督徒,有上帝兒女的名份,擁有神國的基業.第二種人,是「他者」,即並不擁有田地業權的人,以及葡萄園外的人.他們仍然在神的國度之外漂流,在葡萄園外經過和窺探,但沒有進入園中,更遑論採摘葡萄樹的果子.故此,這些人的身份和上帝的兒女有別,他們並不能承受地土和基業,也不能再葡萄園中任意採摘果子,因為他們是外人,是他者.利未記這裡就清楚說明二者的分別.

二者除了身份的不同,我想在此更強調的是二者之間的關係.利未記十九章9至10節那裡說:「在你們的地收割莊稼,不可割盡田角,也不可拾取所遺落的；不可摘盡葡萄園的果子,也不可拾取葡萄園所掉的果子,要留給窮人和寄居的.」「我者」應該如何對待「他者」?之間應如何相待?我與他之間的關係,根據這段聖經,應該不是傳統中國所說的本和末,正與邪的對立和排斥的關係.基督徒雖然得到上帝兒子的名分,但要明白那是白白得來的救恩,就像猶太人之成為選民,是上帝的恩典,而非靠他們的義行.對於還沒有進入神國,未成為基督徒的外人,我們應存著怎樣的態度?這裡說:收割時不要割盡田角,也不要拾取遺下的穗子.在採摘葡萄時,同樣不要摘盡果子,也不要拾取遺下的.意思是:身份雖有不同,但之間不需完全對立,不應割盡,不要摘盡.割盡是盡攬所得,摘盡是毫不留情.這種留有餘地的溫情,體恤弱者的態度,才是基督徒的應有的品格.

二．身份認同:小群與大我之間

如果說利未記得田主與路人,園裡與園外的比喻是指涉基督徒和教外人之間的關係, 那麼,哥林多前書三章那一段經文所提到的,卻是基督徒與另一些基督徒的關係.大家同是教內人,同都稱呼對方為弟兄姊妹的, 但根據林前這段經文來看,在基督徒之中,也分成不少的群體,各有所屬,各有依附,成為多個小群體.有的說:我是屬保羅的,有的說:我是屬亞波羅的,也有的說:我是屬磯法的.如果我們熟悉這些名字,就知道他們都是初期教會的領袖, 都是值得後世景仰的基督教偉人.在人際關係複雜的社會,尋找歸屬感是十分普遍的事.人都喜歡集結成群,組成小圈子,和跟隨有魅力的領袖.俗語稱這種行為叫做“埋堆”.使身有所屬,心有所依.所以同在基督教內,有的說:我是屬保羅的,有的說我是屬亞波羅的,有的說我是屬磯法的.可能我們這批基督徒之中,也會有些人說:我是屬崇基的,我是屬伍牧師的,我是屬許牧師的.所謂「屬」,即歸屬感和認同感.但如果小群的意識膨脹的太大,造成對立和爭競,對於大我有負面的影響,那就不是值得誇耀和鼓勵的事.因此保羅在這裡特別地提醒信徒:(三章22節)「或保羅,或亞波羅,或磯法,或世界,或生,或死,或現今的事,或將來的事,全是你們的,並且你們是屬基督的,基督又是屬神的.」他的意思是:組成小群和小圈子並不要緊,但是要明白一點:無論在甚麼時候,無論在哪种形態,最重要的是屬於基督,屬於上帝.基督徒絕不能只看見所屬的小群,而忽略了長遠的異象及終極的關懷.

三．廢掉/成全:律法與恩典之間

最後,讓我們來看馬太福音五章38-48節那裡關於律法的經文.所謂舊約的律法精神,這裡以「以牙還牙,以眼還眼」作為代表. 當然,我們從利未記,詩篇和其他經典上,可以更加了解猶太人傳統的法規和律例.不過,這些到了新約耶穌基督時期,對律法傳統就有很不同的看法和解釋.不要以惡報惡,不再以牙還牙.反之,“有人打你的右臉,連左臉也轉過來讓他打；有人拿你的裡衣,連外衣也由他拿去；有人強逼你走一里路,你就同他走二里.”這是耶穌基督在新約帶來的價值觀的轉變,以德報怨,以愛化恨. 因為我們活著,不再在律法之下,而是在恩典之下.然而,這並不是說我們在恩典時代生活,就不需要遵守律法了.我們生活在現代社會中,就不需要理會傳統了.我想這段經文原意不是這樣.如果我們往前面看,早些時候,在馬太五章17-18節那裡,耶穌曾經說過:「莫想我來是要廢掉律法和先知.我來不是要廢掉,乃是要成全.就是天地都廢去了,律法的一點一劃也不能廢去,都要成全.」律法和先知,一嚴守律例和槼條,一依賴神跡和法力,是猶太人的兩大傳統,以摩西和以利亞為代表,在新約時代,這兩個傳統的精神一點都沒有廢去, 甚至到了後來的基督教,仍然保存著次兩個大傳統.因此,傳統沒有因為時代轉變就消失或廢掉,它的精神仍然處處可見,在生活的不同方面體現出來,即傳統與現代之間並不對立, 互為因果,互相緊扣,這叫做成全.

結語

今天的信息,可能講得有點太抽象,不易捉摸.所以,容許我做個簡單的歸納如下:做個「之間」人,明白你與別人相交的關係.根據今天所讀的聖經,我們作為基督徒,應如何自處,如何待人,以及如何對待傳統和歷史,如何對待上帝和天國?綜合而言,對他者,我們應盡量採取包容,待以恩情,留有餘地；對同道,雖不排除「埋堆」,結小圈子,但要盡量減少差疑和爭競,不忘大我和遠象,共同為天國大業努力奮鬥；對傳統存著溫情與敬意,對仇敵也懷著寬容和愛心.這樣,我們就能夠走出中間狹窄的縫隙,卸下壓力,帶著愉悅,繼續歡然奔走那連接天人之間的十架道路.



\newpage

\section{靈性高峰}
\label{sec:355}
\textbf{戴浩輝牧師 - 靈性高峰}
\newline
\newline
\hyperref[sec:354]{< < < PREV SERMON < < <}
~
\hyperref[sec:toc]{[返目錄]}
~
\hyperref[sec:356]{> > > NEXT SERMON > > >}


馬太福音 17:1-9
\newline
\begin{longtable}{cl}
\hline
\hline
章節 & 經文 (和合本修訂版)\\
\hline
17:1 & \begin{tabularx}{0.7\textwidth}{X} 過了六天，耶穌帶著彼得、雅各和雅各的弟弟約翰，領他們悄悄地上了高山。 \end{tabularx} \\ \\ \relax
17:2 & \begin{tabularx}{0.7\textwidth}{X} 他在他們面前變了形像，他的臉明亮如太陽，衣裳潔白如光。 \end{tabularx} \\ \\ \relax
17:3 & \begin{tabularx}{0.7\textwidth}{X} 忽然，有摩西和以利亞向他們顯現，與耶穌說話。 \end{tabularx} \\ \\ \relax
17:4 & \begin{tabularx}{0.7\textwidth}{X} 彼得回應，對耶穌說：「主啊，我們在這裡真好！你若願意，我就在這裡搭三座棚，一座為你，一座為摩西 ，一座為以利亞。」 \end{tabularx} \\ \\ \relax
17:5 & \begin{tabularx}{0.7\textwidth}{X} 說話之間，忽然有一朵明亮的雲彩遮蓋他們，又有聲音從雲彩裡出來，說：「這是我的愛子，我所喜愛的。你們要聽從他！」 \end{tabularx} \\ \\ \relax
17:6 & \begin{tabularx}{0.7\textwidth}{X} 門徒聽見，就俯伏在地，極其害怕。 \end{tabularx} \\ \\ \relax
17:7 & \begin{tabularx}{0.7\textwidth}{X} 耶穌進前來，拍拍他們，說：「起來，不要害怕！」 \end{tabularx} \\ \\ \relax
17:8 & \begin{tabularx}{0.7\textwidth}{X} 他們舉目，不見一人，只見耶穌獨自一人。 \end{tabularx} \\ \\ \relax
17:9 & \begin{tabularx}{0.7\textwidth}{X} 下山的時候，耶穌囑咐他們說：「人子還沒有從死人中復活，你們不要把所看到的告訴人。」 \end{tabularx} \\ \\
[1ex]
\hline
\hline
\end{longtable}
出埃及記 24:12-18
\newline
\begin{longtable}{cl}
\hline
\hline
章節 & 經文 (和合本修訂版)\\
\hline
24:12 & \begin{tabularx}{0.7\textwidth}{X} 耶和華對摩西說：「你上山到我這裡來，就在那裡，我要將石版，就是我所寫的律法和誡命賜給你，使你可以教導他們。」 \end{tabularx} \\ \\ \relax
24:13 & \begin{tabularx}{0.7\textwidth}{X} 摩西和他的助手約書亞站起來；摩西上了神的山。 \end{tabularx} \\ \\ \relax
24:14 & \begin{tabularx}{0.7\textwidth}{X} 摩西對長老們說：「你們在這裡等我們，直到我們再回到你們這裡。看哪，亞倫和戶珥與你們同在。誰有訴訟，可以去找他們。」 \end{tabularx} \\ \\ \relax
24:15 & \begin{tabularx}{0.7\textwidth}{X} 摩西上山，有雲彩把山遮蓋。 \end{tabularx} \\ \\ \relax
24:16 & \begin{tabularx}{0.7\textwidth}{X} 耶和華的榮耀駐在西奈山，雲彩遮蓋了山六天，第七天他從雲中呼叫摩西。 \end{tabularx} \\ \\ \relax
24:17 & \begin{tabularx}{0.7\textwidth}{X} 耶和華的榮耀在山頂上，在以色列人眼前，形狀如吞噬的火。 \end{tabularx} \\ \\ \relax
24:18 & \begin{tabularx}{0.7\textwidth}{X} 摩西進入雲中，登上了山。摩西在山上四十晝夜。 \end{tabularx} \\ \\
[1ex]
\hline
\hline
\end{longtable}
講題:靈性高峰　Spiritual Summit

經文:馬太福音17:1-9;出埃及記24:12-18;詩篇2;彼德後書1:16-21

講員:戴浩輝牧師

場合:香港中文大學崇基學院禮拜堂

日期:2017年2月26日

引言

對於馬太福音的作者而言,高山是十分重要的.當耶穌受試探要拜魔鬼,換來世界的權柄時,那是在山上.對路加而言,耶穌在平原教訓人時,馬太就是有山上寶訓.耶穌到山上禱告,今天我們有登山變像.

為何在山上發生的事對馬太而言更重要?馬太的讀者是猶太基督徒,他們對祖宗的信仰傳統十分重視,而山正是神人相遇的地方.首先,舊約開始,西乃山或稱為何烈山是耶和華上帝的住處.當以色列的敵人在山上不能戰勝以色列時,他們就建議要與以色列人在平原爭戰；因為他們相信以色列的神是個山神.山變為一個重要的象徵,代表了神聖和與神同在.最後,上帝在錫安山上統治世界(賽二3；彌四3).因此,登山變像是有其傳統的思想在內,因為摩西和以利亞都在耶和華的山與上帝會面.從這一點我們可以看到為何耶穌在山上變像又有摩西和以利亞.摩西和以利亞就如昔日見上帝那樣,在這高山上覲見聖子上帝,如昔日父神在他們面前一樣.

杜甫的望岳最後兩句「會當凌絕頂,一覽眾山小.」是中國人對山有一定的愛慕之情,因為很多宗教也都以山作為其修行的地方,很多觀與寺都在山上,少林寺就在嵩山.現代人也十分喜歡用山代表很重要的東西,所以常以高峰作為表達重要性,例如:巴黎氣候高峰會議.

我是住在信義宗神學院的,在道風山山上.在過去二十多年我教的神學生畢業時都說:「是我們是時候下山了.」這與我們中國小說中古代的武俠學武藝那樣都是在山上學的,而最後當學完了,他們便下山鋤強扶弱,而我的學生就下山為主作工.

今天顯現期最後一個主日,思想耶穌基督登山變像的事件,讓我們從中學習:「靈性高峰」為了甚麼?是像杜甫那樣的愛泰山的氣勢嗎?今天是顯現期最後一個主日,在教會年中必定是登山變像,對於我們而言,這山有何意義?我希望藉登靈性高峰,參與上帝奧祕和得耶穌之愛作為我們思想的分題.

一  登靈性高峰

基督徒常常為到自己的靈性生活擔憂是一件好事,我生活在道風山也常遇到信主的朋友上山退修,為要走主的道路,退修為要更加學像耶穌基督,退修為要尋求上帝在我們身上美好的旨意.主耶穌登山很可能也有同樣的目的.今天的經文指出,耶穌帶著門徒到一個高山,東方正教認為是在加利利的他泊山,但學者都認為耶穌時代的他泊山有羅馬的駐軍,不大可能在山上留連,又因為他們在黑門山腳的該撒利亞腓立比,大多數學者都認為是黑門山.登山禱告常是我們主耶穌得力的來源,這一次他帶同了門徒中最親密的三個門徒,從耶穌的安排和描述的緊湊[第2節便指出變像]來看,耶穌正在讓門徒登靈性高峰.如果我們從耶穌事工的角度來看,耶穌從開始他的事工,到他登山變像,可以說是耶穌前半段的事工的總結；而當我們從第9節看,登山變像又可以說是後半段事工的開始,是十字架苦路的開始.沒有登山變像的榮耀,作為苦路的起點,我們很難確信苦路和死亡最終是榮耀的復活.沒有登山變像的榮耀,我們很難有力量面對苦路的挑戰.從十七9我們可以肯定,登山變像是上帝對耶穌榮耀復活,升天和再來的一個預示,就好像電影的預告那樣.他帶門徒經歷這預告,並不僅是要門徒記載這事,他要門徒與他登靈性的高峰,為要門徒也與他同走十架的道路.登靈性高峰不是目標,而是途徑,目標是一生走十字架的道路.

各位弟兄姊妹,為要有力量天天背起十架跟隨主,我們是否也與耶穌一起?登靈性高峰?基督徒的見證力量在哪裡?不是美輪美奐的聖堂,又或是神聖的朝聖地點,而是我們每天的生活的場所:我們的家中,學校和工作場所.我們的靈性也不是我們的目標,而是支持我們愛人的力量.

二  參與神的奧祕

耶穌在山上是「臉面明亮如日頭,衣裳潔白如光」,早期教會認為耶穌的變像才是其真正的本像；門徒不僅看見耶穌的真正本相,他們更看見舊約最偉大的先知摩西和以利亞,他們能夠看見上帝的奧祕.在這裡提及的摩西和以利亞,表達了這一刻是歷史的集中點.舊約律法的代表摩西和在瑪拉基書預言要再來的先知以利亞與耶穌一起,是要確立耶穌是救恩歷史的結果,那就是上帝創造人類,但人類犯罪,上帝藉摩西的律法使人類遵行,但人類始終沒有遵行,於是上帝繼續藉先知向他們傳講回轉的信息,正如以利亞所作的；但人的罪惡太深,不能自拔.最後,上帝就差遣他的兒子耶穌基督到這個世界,他本有榮耀的身體,但取了肉身成為人,為要擔當人的罪,死在十字架上,用他的血贖買我們這些不遵守上帝律法的人.門徒得參與上帝奧祕,也成為日後教會的福音內容.不過這並非必然的,我們可以從彼得的回應這個變像的情境時,我們就知道,門徒對榮耀的觀景的嚮往.彼得說:「我們在這裡真好,你若願意,我就在這裡搭三座棚,一座為你,一座為摩西,一座為以利亞.」或許我們會問,他為可要這樣講?首先,棚是住棚節用的,是表示豐收的筵席,棚也有終末的意思,就是彌賽亞在榮耀裡與以色列共享榮耀的筵席.彼得的建議如果是這樣,他心中所想的,就是那一刻已是彌賽亞耶穌永恆榮耀的開始.彼得究竟希望甚麼?從他的獻策他希望此情此景可以繼續下去.是否希望不用下山而享受永恆的榮耀?他希望把握這個時機在手,把本來是上帝的事,在他的主動下,掌握於手中.得參與上帝奧祕的門徒以為永恆的榮耀已經開始,換言之,不用藉十架苦路得來的榮耀已臨到了.人都是不喜歡苦難,人只喜歡榮耀.但上帝所定的旨意是藉基督的十字架,最終人類得著永恆的榮耀.然而,我們只喜愛榮耀,而不接受達到榮耀的苦路.

彼得的希望很快就得到回應,從天上來的聲音,雲中的聲音強調了舊約的三處經文:上帝的兒子(詩二7),上帝所喜悅的(受苦的僕人[賽四十二1]),你們要聽他(申十八15).這幾節是繼太三17後,再次肯定耶穌的身分人.而另一方面,也按太十六13-20彼得宣認耶穌是基督,是永生上帝的兒子後,從天上而來的,正式在登山變像時,得到肯定,有天上的蓋章.

當我們參與上帝的奧祕時,有條不能不走的十架之路放在我們面前；不藉十字苦路而得著的榮耀不是耶穌的道路,所以天上的聲音告訴我們:「你們要聽他!」今天,靈性的高峰也在這一句:我們要聽耶穌的.他說:「要愛你們的仇敵.」我們是否聽他?他說:「你願意人怎樣待你,你便要怎樣待人.」你聽他嗎?這樣,你真的在靈性高峰,參與上帝的奧祕.為著與耶穌同步,我們只應聽從耶穌,正如天上的聲音所說的.為了要有好的傳道人,每個教會都有人事部來考核加入的同工.大概要考的就是神學的知識,但有一間教會的人事部有一位牧者的考核方式三十年都是一樣.他問參與投考的:「你望到窗外那個人嗎?」投考的回答說「我望到.」牧師再問:「你認識他嗎?」投考的說:「不認識.」牧師說:「好,我問你,請你用你所學的神學來形容這個人吧!」這一個問題,你們能回答嗎?牧師分享說:「過去三十年投考的大概有兩個答案:1. 那人是個罪人,需要耶穌基督的救贖.2. 那人是否知道都好,他是上帝兒女,上帝在耶穌基督裡保守他愛他.」牧師說:「兩個答案都有神學,但以我的經驗,回答第二個答案的人通常都會是更好的牧師.」因為他看人是從上帝創造與終極的目標看的.

三  領受耶穌恩典

在舊約裡,不潔的人不能見上帝,見的人必死；摩西在上帝的山,民眾只能在山下等候；因為走近也可以成為殺身之禍.對於門徒來說,他們得見上帝在基督身上的榮耀,又希望維持這榮耀美境,胡亂的提議,可以有很可怕的後果的.第6節指出,門徒極其害怕,不無道理.但他們是在耶穌基督恩典之中,得耶穌之愛.

7-9節表達了上帝的另一面,就是聖子的仁愛,耶穌摸他們,從聽覺的感觀是可怕的,但一對熟悉的手使他們得安慰,這對手曾摸過西門的岳母,使發燒的身體痊癒；使長大痲瘋的得潔淨；使瞎眼的得看見.他們跟著就是得聽安慰的聲音:「不要害怕.」

各位弟兄姊妹,有時當我們遭遇痛苦艱難,我們感觀覺得十分可怕,自己也沒有其他辦法,我們可以怎樣做呢?門徒的唯一方法就是抱頭發抖,我們能作的也差不多.在這個時刻,我深信我們只能依靠那為我們走十架苦路的基督,我們祈求他的醫治的手摸我們,我們也深願能聽到他的聲音對我們說:「不要害怕!」

無論我們登靈性高峰,又或參與上帝的奧祕,我們能走人生的路程的,最重要是得著耶穌的愛.抱頭發抖的門徒最後舉目一看,「只見耶穌」,這句就是我們一生的要訣.我們所走的路無論如何的困難艱辛,只要我們能仰望耶穌,眼只見耶穌一人,我們必能勝過一切的困難和艱辛.

結語

攀登靈性高峰,參與神的奧祕和領愛耶穌恩典是今天登山變像給我們的信息,在此我們學會與耶穌交往,我們甚或參與上帝的奧祕和榮耀,但這一切都要在我們每天的生活中,背起我們的十架跟隨耶穌,最終,我們所初嘗榮耀,將會為我們存留.



\newpage

\section{正邪交鋒愈戰愈勇風起雲湧}
\label{sec:356}
\textbf{許開明牧師 - 正邪交鋒:愈戰愈勇,風起雲湧}
\newline
\newline
\hyperref[sec:355]{< < < PREV SERMON < < <}
~
\hyperref[sec:toc]{[返目錄]}
~
\hyperref[sec:357]{> > > NEXT SERMON > > >}


馬可福音 1:12-13
\newline
\begin{longtable}{cl}
\hline
\hline
章節 & 經文 (和合本修訂版)\\
\hline
1:12 & \begin{tabularx}{0.7\textwidth}{X} 聖靈立刻把耶穌催促到曠野裡去。 \end{tabularx} \\ \\ \relax
1:13 & \begin{tabularx}{0.7\textwidth}{X} 他在曠野四十天，受撒但的試探，並與野獸同在一起，且有天使來伺候他。 \end{tabularx} \\ \\
[1ex]
\hline
\hline
\end{longtable}
路加福音 4:1-13
\newline
\begin{longtable}{cl}
\hline
\hline
章節 & 經文 (和合本修訂版)\\
\hline
4:1 & \begin{tabularx}{0.7\textwidth}{X} 耶穌滿有聖靈，從約旦河回來，聖靈把他引到曠野， \end{tabularx} \\ \\ \relax
4:2 & \begin{tabularx}{0.7\textwidth}{X} 四十天受魔鬼的試探。在那些日子，他沒有吃甚麼，日子滿了，他餓了。 \end{tabularx} \\ \\ \relax
4:3 & \begin{tabularx}{0.7\textwidth}{X} 魔鬼對他說：「你若是神的兒子，叫這塊石頭變成食物吧。」 \end{tabularx} \\ \\ \relax
4:4 & \begin{tabularx}{0.7\textwidth}{X} 耶穌回答說：「經上記著：『人活著，不是單靠食物。』」 \end{tabularx} \\ \\ \relax
4:5 & \begin{tabularx}{0.7\textwidth}{X} 魔鬼又領他上了高山，霎時間把天下萬國都指給他看， \end{tabularx} \\ \\ \relax
4:6 & \begin{tabularx}{0.7\textwidth}{X} 對他說：「這一切權柄和榮華我都要給你，因為這原是交給我的，我願意給誰就給誰。 \end{tabularx} \\ \\ \relax
4:7 & \begin{tabularx}{0.7\textwidth}{X} 你若在我面前下拜，這一切都歸你。」 \end{tabularx} \\ \\ \relax
4:8 & \begin{tabularx}{0.7\textwidth}{X} 耶穌回答他說：「經上記著：『要拜主—你的神，惟獨事奉他。』」 \end{tabularx} \\ \\ \relax
4:9 & \begin{tabularx}{0.7\textwidth}{X} 魔鬼又領他到耶路撒冷去，叫他站在聖殿頂上，對他說：「你若是神的兒子，從這裡跳下去！ \end{tabularx} \\ \\ \relax
4:10 & \begin{tabularx}{0.7\textwidth}{X} 因為經上記著：『主要為你命令他的使者保護你； \end{tabularx} \\ \\ \relax
4:11 & \begin{tabularx}{0.7\textwidth}{X} 他們要用手托住你，免得你的腳碰在石頭上。』」 \end{tabularx} \\ \\ \relax
4:12 & \begin{tabularx}{0.7\textwidth}{X} 耶穌回答他說：「經上說：『不可試探主—你的神。』」 \end{tabularx} \\ \\ \relax
4:13 & \begin{tabularx}{0.7\textwidth}{X} 魔鬼用完了各樣的試探，就離開耶穌，再等時機。 \end{tabularx} \\ \\
[1ex]
\hline
\hline
\end{longtable}
講題:正邪交鋒:愈戰愈勇,風起雲湧　The Battle between Good and Evil: It Will Become More Fierce and Turbulent!

講員:許開明牧師

經文:馬可福音1:12-13;路加福音4:1-13

場合:香港中文大學崇基學院禮拜堂

日期:2017年2月26日

祈禱

上主,在教會年曆大齋期,我們默想主在曠野40日禁食禱告受試探,願我們從中得著生命的自省,勝過試探.阿們.

經課:耶穌受試探(可1．12－13；路4．1－13)

1當時,耶穌被聖靈引到曠野,受魔鬼的試探.

2他禁食四十晝夜,後來就餓了.

3那試探者進前來對他說:「你若是上帝的兒子,叫這些石頭變成食物吧.」

4耶穌卻回答說:「經上記著:『人活著,不是單靠食物,乃是靠上帝口裡所出的一切話.』」

5魔鬼就帶他進了聖城,叫他站在聖殿頂上,

6對他說:「你若是上帝的兒子,就跳下去!因為經上記著:『主要為你命令他的使者,用手托住你,免得你的腳碰在石頭上.』」

7耶穌對他說:「經上又記著:『不可試探主-你的上帝.』」

8魔鬼又帶他上了一座很高的山,將世上的萬國和萬國的榮華都指給他看,

9對他說:「你若俯伏拜我,我就把這一切賜給你.」

10耶穌說:「撒旦,退去!因為經上記著:『要拜主-你的上帝,惟獨事奉他.』」

11於是,魔鬼離開了耶穌,立刻有天使來伺候他.

引言

讀耶穌受試探一事的時候,你有沒有正邪交鋒的感覺?高手過招,愈戰愈勇,風起雲湧,驚天動地!正就是指耶穌基督,神之子,人子(看得見的上帝),上帝國的王；而邪就有那試探人者,魔鬼,撒旦(敵對者),世界之王.

耶穌受試探的地點從曠野,到殿頂,到高山,步步為營；而試探的內容更是一個比一個難於招架.我們要留意這些試探不是一次就完,而是在耶穌一生中,不斷出現,不斷地正邪交鋒,高手過招,愈戰愈勇,風起雲湧(只是試探前奏).還有,這個試探不只是耶穌面對的,也是所有跟隨耶穌的基督徒(小基督)都面對的,更重要的是這些試探是每個人一生都會面對的.

這篇是大學的講章,筆者今早會嘗試從神學,哲學及心理學的角度來探討這段經文.

上世紀90年代,筆者在旺角牧會時,由於色情場所在區內非常嚴重,我們旺角區一班教牧同工聯合起來推動區內的反色情運動.期間最痛心的是發現在警察團契熱心事奉的弟兄,在警署負責掃黃的隊竟然自己也中招,最後更身敗名裂.打擊罪惡的人也會被罪惡所擊,所以我們務必小心.

試探與試練

「試探」,「試煉」,「試驗」,這三個詞語的意思有甚麼分別呢?「試探」(temptation)與「試煉」(trial)這兩個詞的意思大同小異,相同之處在於都有「試驗」(test)的意思,而「試驗」的目的在於證實,所以人若經過「試驗」被證實是失敗,對他而言就是「試探」；反之,人若經過「試驗」被證實是勝過,對他而言就是「試煉」.

聖經觀念

在新約雅各書中,我們也讀到:「人被試探,不可說:『我是被神試探；』因為神不能被惡試探,祂也不試探人.」(雅1:13)可見神不會「試探」人,但神會「試煉」人.

「雅各接著又說:『但各人被試探,乃是被自己的私慾牽引誘惑的.』」(雅1:14)這「私慾」很容易受到魔鬼的誘惑,而魔鬼是專門試探人的.(太4:3)

從舊約約伯記中可以得知:神要「試驗」約伯,以證實他是個完全正直,敬畏神,遠離惡事的人.(伯1:8)就把約伯交在撒旦手中.(伯1:12)撒旦的手段是「無所不用其極」的,他的目的就是叫人遠離神.

「試探」的目的在於「叫人跌倒」；「試煉」的目的在於「叫人爬起」.試探是出自撒旦,要你跌入他的網羅,然後傷害,毀滅你,使你遠離上帝；試煉則是出自上帝,要你經過其樣試煉而變得剛強,使你更接近上帝,甚至更象上帝,而經過試探而不跌倒,這試探可能變成試煉.我們要知道道成肉身的耶穌,不是一出生,一出現就是超人,天下無敵,而是像人一樣受各種的試探和試煉,所以他能明白我們人的遭遇和軟弱.

三個典型的人生試探

第一個:追求感官享受的試探 (不從)

它的性質是要人追求滿足身體的需要.

3那試探者進前來對他說:「你若是上帝的兒子,叫這些石頭變成食物吧.」

4耶穌卻回答說:「經上記著:『人活著,不是單靠食物,乃是靠上帝口裡所出的一切話.』」

在曠野 ,無東西可吃,無水,無人基本的需要,甚至是無助.那試探者的招數就是試探耶穌的能力,要耶穌若是上帝的兒子,就把石頭變成食物,那試探者目的就想耶穌運用自己的能力,來滿足自己的食慾,甚至是傳教上,用食來吸引群眾,不用去釘十架那麼辛苦.耶穌卻回答說:「經上記著:『人活著,不是單靠食物,乃是靠上帝口裡所出的一切話.』」帶出了聽從上帝的話及道比食更重要.

許多人被食困著,被色困著,被煙酒困著,被錢財困著,這些試探相當厲害,最易使許多人跌倒.

第二個:追求名聲的試探(不跳)

它的性質是滿足心靈 (虛榮心)的需要,而地點是在殿頂上,讓人人皆看到,一跳就可成名.

5魔鬼就帶他進了聖城,叫他站在聖殿頂上,

6對他說:「你若是上帝的兒子,就跳下去!因為經上記著:『主要為你命令他的使者,用手托住你,免得你的腳碰在石頭上.』」

7耶穌對他說:「經上又記著:『不可試探主-你的上帝.』」

那試探者就是試探耶穌若是上帝的兒子,就跳下去,因為上帝會為耶穌做準備,使耶穌跳下去也不會死,會得到保護的.那試探者目的就時想耶穌運用自己的能力,使他跳下去不死,從而引起人注意,聲名大噪,更可藉此吸引多人因耶穌的神蹟來跟從他,那麼耶穌就不用上十架那麼辛苦.

耶穌只是對他說:不可試探神.因為無論神或耶穌怎樣做,魔鬼都達到他的目的.從高處跳下來,如果耶穌死了,就是上帝食言,耶穌無能；要是沒有摔死,上帝就好像被迫托住耶穌,上帝被迫失去自由,是耶穌聽了魔鬼的話去迫上帝,叫祂失去自由.加上,從聖殿跳下來一定轟動,跳不死更轟動,更受注目,這神蹟顯示耶穌了不起,上帝為兒子工作.更重要的是這是試探上帝,把上帝堪入測試中,而試探上帝是罪,是不信上帝的表現,是魔鬼想上帝陷入你的測試中,看上帝能否照你的心願.那是你已有個心願,你想上帝為你達成,就是想限制上帝的自由,也是你用上經文的說話,想藉此限制上帝的自由,最嚴重的情況就是製造內部矛盾.

我年輕時剛當上堂主任的時候,有一位同事走來說:「牧師,我們給你這麼難的事,就是想看你如何處理,試探你對他的反應!」我聽了覺得很不開心,你開心嗎?你願意被試探嗎?因為我是上司,下屬卻對你說:「上司,我想試下你的反應,我早就猜到你的反應,我很聰明嗎?」他聰明嗎?又聰明又笨,笨的地方就是他試探上司是想控制上司嗎?上司是他的試驗品嗎?

多少人過不了名利,聲譽那關的試探.

第三個:追求榮華富貴的試探  (不跪)

8魔鬼又帶他上了一座很高的山,將世上的萬國和萬國的榮華都指給他看,

9對他說:「你若俯伏拜我,我就把這一切賜給你.」

10耶穌說:「撒旦,退去!因為經上記著:『要拜主-你的上帝,惟獨事奉他.』」

它的性質是滿足人肉體,心靈,一切的需要.魔鬼選擇的地點是高山,因為那是只有耶穌與魔鬼能及,而其他人不能到的地方,讓耶穌可以靜靜做事,跟魔鬼來個秘密交易.魔鬼自稱他是世界之主擁有萬國和萬國的榮華,拜他就可得榮華富貴,感官享受和名利,滿足肉體和心靈的需要.魔鬼試探的目的是要使耶穌向他下拜,使耶穌的一生只追求名華富貴,並使跟隨耶穌的人也追求此屬世的東西,不能進天國,那麼耶穌也不能完成上帝的使命,會放棄上十字架.耶穌卻是這樣回應:「撒旦,退去! (命令)因為經上記著:『要拜主-你的上帝.』」耶穌知道為榮華而拜撒旦,只會受撒旦操控,

所以耶穌要撒旦走,不准他在身旁!

馬丁路德說過:「我不能阻止鳥在頭頂飛,但可以阻止鳥在我頭上築巢!」所以許多人在試探中跌倒,是因為容許罪惡築巢.

試探就是在試權力與能力的運用,有權不要用盡,更不要亂用,要合適地用權,因為絕對的權力,有可能帶來絕對的腐化.

第一個試探:有關使用自己的權力

把石頭變餅,能力是用在聽從那試探者嗎?能力用在滿足個人的食慾嗎?耶穌有行神蹟的權力,可使石頭變餅,改變不同性質的東西,這是是超越理性和違規的,因那是為了什麼?為了滿足人食慾?為了滿足別人的食慾?為了吸引人來歸信?無疑地,這是容易快速,無痛的方法,這也是現代人追求的方法,但這種能力使用是有問題的.人活著是單靠食物嗎?

第二個試探: 有關借用上帝的能力 (未經上帝允許)

從殿頂跳下,這能力的使用是為搏出名,以借用上帝的能力來使自己成名,並使上帝陷入他的試驗中,這可是更錯.

第三個試探:有關向惡勢力妥協,得其力獲得名華富貴

跪拜魔鬼,以得擁有榮華富貴的能力,這天是天大的錯.自古到今,許多人做的事就是借用魔鬼能力,得到榮華富貴.

從哲學家康德看人的感知能力

康德說:「物自身是不能見的,你只從你理解的物的現像去掌握,你看到許多是物的現象,我想你看到的部份.」大概意思就是人從自己的經驗和感觀去認識世界.

1感性:人的感觀,接受外在事物的刺激；

2悟性:從接觸外在物,得到對物件的悟性；

3理性:進行推理及邏輯思考.

同樣地,耶穌所受的試探也是這三方面.

第一:感性受試探

食是人重要的感覺,是人最易跌倒的部分.耶穌40日不吃不喝,在大饑餓中,最易落在試探中,所以試探者就建議耶穌把石頭變餅.現實中也有很多不同例子,例如我餓的時候我會發惡,易做錯事；我的女兒肚餓時會哭,易騷擾別人；人類始祖亞當,夏娃都勝不過感性的試探.耶穌引出人活著不單是靠感性,滿足感官食慾,還需要「上帝的道」,上帝的道是需要悟性,才可領悟上帝的意思.

第二:悟性受試探

那試探人的,進一步令耶穌的悟性受試探,他引用經文,引出耶穌若是上帝的兒子,從聖殿頂向下跳,根據經文,天使會用手托住耶穌的腳,免他受傷.悟性的偏向就是要成名,引人注意.耶穌的回應卻是:「不可試探主你的上帝.」

第三重是:理性的試探

魔鬼帶他耶穌到高山,對他說:「你若拜我,我就把這一切都賜給你.」理性的偏向是如何獲得榮華富貴,例如去拍馬屁,用錢,用威嚇等各種手段,這就是向撒旦下跪的意思,但耶穌卻用靈性:「單要事奉主」,來擊退試探者.

從始祖偷食禁果

那蛇用了不同的招數來試探亞當和夏娃.(創3:1-6)

一:悟性及理性的試探

1耶和華神所造的,惟有蛇比田野一切的活物更狡猾.蛇對女人說:「神豈是真說不許你們吃園中所有樹上的果子嗎?」

2  女人對蛇說:「園中樹上的果子,我們可以吃,

3惟有園當中那棵樹上的果子,神曾說:『你們不可吃,也不可摸,免得你們死.』」

4  蛇對女人說:「你們不一定死；

5因為神知道,你們吃的日子眼睛就明亮了,你們便如神能知道善惡.

試探者即是蛇,就是使人悟性及理性偏向蛇,使人懷疑上帝所說的話,結果就是女人被絆倒了.

二:感性的試探

6 於是女人見那棵樹的果子好作食物,也悅人的眼目,且是可喜愛的,能使人有智慧,就摘下果子來吃,又給她的丈夫(唯一至愛),她的丈夫也吃了.

這個感性的試探就是至愛的人給的不可不吃；而在理性上,她吃完果子也沒有立即死去,所以也照樣吃了.試探者就是使人感性偏向那果子,先誘惑,迷感人心.

我們要小心,我們的感性,悟性,理性均被試探,都向罪惡,向那惡者傾斜,這就遭糕了.

那惡者用三種方式來試探人使人跌倒:

1利用人感性的偏向,只求感官滿足,以食為先；

2利用人悟性的偏向,偏向罪性；

3利用人理性偏向,偏向罪性.

伊斯蘭國錯誤的思想就是替天行道,以崇拜恐怖的手段,建立一個強大的伊斯蘭帝國,為要粉碎英美帝國,這些都是從悟性,感性的被引走向偏差,趨向感性上的偏差.

我們要注意:悟性,感性的被引走向偏差的方和,趨向感性的偏差不獨是:

1對事情的；

2對人:希特拉所為,他當時得到許多人的支持

3對神:回教徒說阿拉支持他們殺人,且可上天堂.

面對試探耶穌意志堅定,始祖意志簿弱

意志是人自覺地確定目的,並支配行動,克服困難,實現目的的心理過程.意,是指心理活動的一種狀態；志則指對目的方向的堅信,堅持.意志合起來就是對實現目的有方向,有信念地堅持的一種心理活動(含潛意識中的心理活動).

意志的建立有賴三方面:知(即認知),情(即情感),意(即意志),這是是人類心理活動的三種基本形式,人們認識世界的主觀意識過程,通常可以分為三個階段:

1認知階段:目的在於解決「是什麼」或「什麼事實」的問題.人只有首先瞭解事物的外在特性(或外部聯繫)和內在規律(內在本質),即首先瞭解事物「是什麼東西」,才能對它進行其它方面的深入瞭解；

2 評價階段:目的在於解決「有何用」或「有什麼價值」的問題.人只有瞭解在瞭解事物「是什麼東西」以及「對我有何價值」,才能知道如何對它採取正確的處理措施；

3 意志(或決策)階段:目的在於解決「怎麼辦」」或「實施什麼行為」的問題.就是針對事物的品質特性以及每一品質特性對於人的價值,人將選擇一個最合適的行為,以便能夠充分有效地利用事物的價值特性；

沒有認知的評價就是一種盲目的評價,沒有認知的意志就是一種盲目的意志；沒有評價的認知就是一種麻木的認知,沒有評價的意志就是一種無聊的意志；沒有意志的認知就是一種空洞的認知,沒有意志的評價就是一種空洞的評價.

耶穌所受的試探就能體現這三個認個,評價,意志的階段.

第一個試探

1認知事物: 四十日禁食,肚餓,那試探者叫他把石頭變餅

2價值取向:人活著不是單靠食物,乃靠上主的話

3 意志堅持:不做,不把石頭變餅,注重上主的話

第二個試探

1認知事物:那試探者叫他從殿頂跳,不會受傷,因上帝保護

2價值取向:如此試探上帝,無意義

3 意志堅持:不跳,反說:「不可試探主上帝」

第三個試探

1認知事物:魔鬼叫耶穌跪拜得富貴,可得榮華富貴

2價值取向:耶穌追求的不是地上的國,地上的榮華富貴

3 意志堅持:不拜,只拜上帝,並把鬼趕走,而不是耶穌避開鬼

1知己知彼:耶穌清楚知道自己是誰,是上帝的兒子；也知道魔鬼是誰,是撒旦,抵擋者；

2知道魔鬼的詭計；見上述.

3知道正解:耶穌從聖經中得知正解,沒有落入魔鬼的圈套；

4知道要意志堅定,魔鬼要動搖人心.不停搖動人的心,人的手和腳!

因著耶穌以上的行動,魔鬼離開了耶穌,並有天使來伺候他.即使耶穌繼續受試探,仍有天使(上帝的使者)同在.

結束祈禱

上主:

我們總會面對人生不同的試探,祈求上主,使我們在試探中,不致跌倒,好像耶穌基督會還能站起,變得更加堅強!阿們!



\newpage

\section{尋找的就尋見}
\label{sec:357}
\textbf{林豪恩先生 - 尋找的,就尋見}
\newline
\newline
\hyperref[sec:356]{< < < PREV SERMON < < <}
~
\hyperref[sec:toc]{[返目錄]}
~
\hyperref[sec:358]{> > > NEXT SERMON > > >}


講題:尋找的,就尋見　The One Who Seeks Finds

經文:撒母耳記上9

講員:林豪恩先生

場合:香港中文大學崇基學院禮拜堂

日期:2017年3月12日

一．引言

有人說,人生,就是無數尋尋覓覓的故事:牙牙學語的時候,就去尋找學校；踏入畢業年,就開始尋找工作；對不少香港人來說,最重大的任務是買樓,即是尋找一處安居之所.本來,有些尋找是不需要自己操勞的,但由於供應緊張,逼得有些人也不得不在有生之年為自己尋找長眠之地.

在眾多尋覓中,最引人入勝,讓人盪氣迴腸的,可能是尋找伴侶的故事.當人進入青春期,睪固酮與雌激素會在體內變得活躍.這些荷爾蒙創造體驗愛情的欲望,使人們開始尋找伴侶.

YouTube2013年度香港熱門影片首位,點擊超過170萬次的,是一段改篇自一篇尋找伴侶的網絡潮文,文章名為《男人一生,只為尋覓一個肯同自己挨麥記ge 女人》.這篇潮文如此說:

「是咁的,身邊D 朋友成日問我有咩擇偶條件,巴巴閉閉咁分享佢地ge 心得.話要有劉心悠ge 面孔,法拉ge 笑容,雷凱欣ge 身材,何家淑女ge 後台,Angela Monkeyge 整容醫生.我多數就一句車埋去:「咁佢會唔會陪你去食麥記先?」

個宅男肥仔拎住個巨無霸,臉容扭曲到好似畢加索D 名畫咁.聽到依個問題,佢已經從幻想中解放出黎.唔想講都要坦白講,一個男人去到廿幾歲,同個女仔,「去邊食飯」依個問題仲難解拆過「南北韓危機」.有時銀包就得個幾廿蚊,衰D 講句入星巴克買杯Latte扮下高尚,已經要行路返屋企.

試過同一位已成歷史塵埃ge 女朋友走入一間茶記.喂唔差架喎其實!一個客飯配咕嚕肉,唔包餐飲都78蚊,絕對唔cheap!只不過人地就唔係咁諗啦,又話唔乾淨,又話個位逼,又話電視機好嘈.我想即刻飛左佢.但事實上佢即日已經飛左我.

我曾經以為女人都是這樣,我曾經想過做「一休和尚」算數.估唔到有生之年仲搵到一個主動同我食麥記ge 女仔.

一坐低佢就問我食咩,就自己行左出去買. 我心諗「出事!」係麥記兵法之中,只要你戰略得宜,一餐飯既價格要超越California Pizza絕對唔係難事.一個散嗌ge 扭扭粉,掟個靚彎去隔離「麥café」整杯藍梅特飲,加件芝士餅.驀然回首order多客薯條走鹽,加個雞翼併盤,臨尾飄移外賣一個開心樂園餐,話要換禮物比細佬.成壇野已經一百幾廿蚊,仲要未買你個份.

只係說時遲,那時快,女朋友已經買jor 返黎,佢最後要既只係廿一蚊既麥樂雞餐,無加大,無加三蚊變檸檬茶.我見到經眼濕濕,想話比返錢佢,佢仲借D 意話遲下食好ye 先比我請返.我好滾動,流下男兒淚.個一刻,我終於明白,咩叫感動.只有女朋友係你唔太惦ge 情況下,主動拉你去食麥記,仲要顧住你面子,唔令你難堪,掛住個笑樣拖住你隻手仔食薯條,個一下先叫感動.

劉心悠始終會老,法拉都有黑面ge 時候,雷凱欣會爆炸,何家淑女你又唔係陳X聰,Angela Monkey都開始同人撞樣啦!男人一生,尋覓到一個肯同自己挨麥記ge 女人,三個字.娶得過.(香港網絡大典)

尋尋覓覓,過去的日子,你努力尋找甚麼呢?此時此刻,你又在尋找甚麼呢?有沒有想過,你在尋找的時候,你也在被尋找呢?

二．經文與信息

撒母耳記為掃羅拉開序幕的故事,也是一個尋找和被尋見的故事.讓我們一同閱讀這個故事,也一同閱讀自己.

第一幕:尋驢

9:3掃羅的父親基士丟失了幾匹母驢,他就吩咐兒子掃羅說:「起來,帶一個僕人去尋找驢子.」

眾裡尋他千百度,「掃羅走過以法蓮山區,又過沙利沙地,都沒有找著.他們走過沙琳地,驢不在那裡,又走過便雅憫地,也沒有找到.」(9:4-5)

驀然回首,那人也不在燈火闌珊處,「9:5到了蘇弗地,掃羅對跟隨他的僕人說:「我們不如回去……」

是的,何必執著呢?尋不著,又如何?失去了,又怎樣? 窮途似是末路,不如放棄吧!勉強下去,可能有代價的,怎知道尋到的不會比因此而失去的更重要呢呢?掃羅對跟隨他的僕人說:「我們不如回去,恐怕我父親不為驢掛心,反為我們擔憂.」

去或留,堅持還是放棄,不是容易的決定.有時,我們遇到挫折,覺得累了,就會想放棄.但當放棄的一刻,想到過往付出不少努力,已走了一段路而不繼續下去,又不禁問:「說不定還可以再多堅持一會」,「我努力了這麼久,真的要放棄嗎?」有些人行動上放棄了,卻並不見得內心可以放開,心有不甘,情緒糾結.這些狀況,或多或少我們都經歷過.

當我們為堅持還是放棄而痛苦的時候,為去或留計不清數的時候,有沒有想過仰望上主,尋找上天的意思呢?尋找,有時不在於尋回失去的東西,也不在於尋得想要的東西,而在於在尋找的過程中更明白自己,更明白自己的需要,甚至在尋找的過程中讓自己被尋見.

第二幕:尋人

9:6僕人對他說:「看哪,這城裡有一位神人,受人敬重,凡他所說的全都應驗.現在讓我們到他那裡去,或者他能指示我們當走的路.」

果真要去向神人求問嗎?我配得上主的眷顧嗎?我準備好去求見神人嗎?

9:7掃羅對僕人說:「我們若去,有甚麼可以送那人呢?我們囊中的食物都吃盡了,也沒有禮物可以送那神人,我們還有甚麼沒有?」在尋找上主的路上,有些時候,我們感到自己不配,感到自己尚未準備好,有一些顧慮令我們裹足不前.

9:8僕人回答掃羅說:「我手裡有銀子一舍客勒的四分之一,可以送那神人,請他指示我們當走的路.」當我們猶豫的時候,身邊的人顯得很重要,往往是他們把我們帶到上主面前,重要的是我們不拒絕他們的協助及同行.在尋找上主的路上,曾經有人協助過你嗎?你願意接受別人的協助嗎?

9:10掃羅對僕人說:「你說的是,我們可以去.」於是他們往神人所住的城裡去了.

仰望上主,彷如開了一扇窗,看見另一種可能性.到底是誰在尋找誰呢?掃羅在尋找失去的驢的同時,原來有人在尋找他；在未見目標之前,自己卻先被尋見.

第三幕:被尋找

9:11他們上坡要進城,就遇見幾個少年女子出來打水,問她們說:「先見在這裡沒有?」12女子回答說:「在這裡,他在你們前面.快去吧!他今日正到城裡,因為今日百姓要在邱壇獻祭.13在他還沒有上邱壇吃祭物之先,你們一進城必遇見他；因他未到,百姓不能吃,必等他先祝祭,然後請的客才吃.現在你們上去,這時候必遇見他.」14二人就上去；將進城的時候,撒母耳正迎著他們來,要上邱壇去.

在尋尋覓覓的過程中,時機尤關重要.戲劇中緣慳一面的遺憾,往往令觀眾感到無限婉惜.掃羅來得正合時機,他與神人即將碰面了.且慢,這到底這是掃羅尋找神人的故事,還是神人尋找掃羅的故事呢?讓我們回想一下,你與上主相遇,是你尋找上主的故事,還是上主尋找你的故事呢?

9: 15掃羅未到的前一日,耶和華已經指示撒母耳說:16「明日這時候,我必使一個人從便雅憫地到你這裡來,你要膏他作我民以色列的君.他必救我民脫離非利士人的手；因我民的哀聲上達於我,我就眷顧他們.」17撒母耳看見掃羅的時候,耶和華對他說:「看哪,這就是我對你所說的人,他必治理我的百姓.」

魯米,是一位13世紀的伊斯蘭教蘇菲派神秘主義詩人,他的詩歌中有一名句:What you seek is seeking you(你尋找你正尋找你).事實上,聖經中充滿上帝主動來尋找人的事件,是上帝尋找人,而不是人尋找上帝.由於是上帝尋找人而非人尋找上帝,上帝與人之間的鴻溝,對人而言是無法踰越的,但對上帝卻從未構成障礙,「我是否配得」,「我是否有靈智慧根」,「我是否準備好」等因素都不重要.

第四幕:尋見自己

掃羅被撒母耳尋見同時,也是掃羅尋見自己的時刻.

9:21「我不是以色列支派中最小的便雅憫人嗎?我的家族不是便雅憫支派中最小的家族嗎?你為何對我說這樣的話呢?」我是誰呢?我當為何而活呢?或者,這也是所有真誠地活著者的思考.

掃羅是誰呢?作者在故事開始已經作出介紹:

9:1有一個便雅憫人,名叫基士,是便雅憫人亞斐亞的玄孫,比歌拉的曾孫,洗羅的孫子,亞別的兒子,是個大能的勇士.9:2他有一個兒子,名叫掃羅,又健壯,又俊美,在以色列人中沒有一個能比他的；身體比眾民高過一頭.

在世界的標準中,掃羅出類拔萃,鶴立雞群,然而,在上主的面前,那又算得甚麼呢?「我不是以色列支派中最小的便雅憫人嗎?我的家族不是便雅憫支派中最小的家族嗎?」就在那一次尋覓的過程中,掃羅尋得了自己.其實,我是誰呢?

10:5此後你到上帝的山,在那裡有非利士人的防兵.你到了城的時候,必遇見一班先知從邱壇下來,前面有鼓瑟的,擊鼓的,吹笛的,彈琴的,他們都受感說話.

10:6耶和華的靈必大大感動你,你就與他們一同受感說話；你要變為新人.

10:9掃羅轉身離開撒母耳,上帝就改變他,賜給他另一顆心.

尋獲上主所賜的一顆心,變為一個新的人,這更是意外收獲,極大驚喜.然後呢?接下來的人生該如何過呢?為甚麼而活呢?

9:15掃羅未到的前一日,耶和華已經指示撒母耳說:

9:16「明日這時候,我必使一個人從便雅憫地到你這裡來,你要膏他作我民以色列的君.他必救我民脫離非利士人的手；因我民的哀聲上達於我,我就眷顧他們.」

潘霍華說:「我們的上帝是為他人而存活的,因而被釘在十字架上」,「惟有教會在為他人而存在時,才顯出她真正的本色.」

掃羅尋找到自己了!出生於顯赫家族,又健壯,又俊美,又高大,所為何事呢?豈不是為了回應世間的哀聲?豈不是為了對抗侵略者?保護被欺壓的同胞嗎?

9:20「至於你前三日所丟失的幾匹母驢,你心裡不必掛慮.」有時,我們不甘心,誓要尋回我們所失去的,反而錯過了更加寶貴的.在尋尋覓覓之中,尋找得到固然高興,但過程中的發現,不也是珍貴嗎?

三.總結

有人尋尋覓覓,嚐得冷冷清清,落得悽悽慘慘戚戚.守著窗兒,獨自怎生得黑.為何不眺窗遠望,發現另外一種可能性?在彼岸,或許有新的自我發現,或許被賦予新的心靈,或者是全新的境界.

驀然回首,那人卻在燈火闌珊處.那人,正是掃羅自己!是按上主的美意所創造的自己.活出上主所創造的那個自己,正是尋覓之中最大的收穫.



\newpage

\section{因信稱義的因恩欣}
\label{sec:358}
\textbf{鄭漢文博士 - 因信稱義的因,恩,欣}
\newline
\newline
\hyperref[sec:357]{< < < PREV SERMON < < <}
~
\hyperref[sec:toc]{[返目錄]}
~
\hyperref[sec:359]{> > > NEXT SERMON > > >}


羅馬書 5:1-11
\newline
\begin{longtable}{cl}
\hline
\hline
章節 & 經文 (和合本修訂版)\\
\hline
5:1 & \begin{tabularx}{0.7\textwidth}{X} 所以，我們既因信稱義，就藉著我們的主耶穌基督得以與神和好。 \end{tabularx} \\ \\ \relax
5:2 & \begin{tabularx}{0.7\textwidth}{X} 我們又藉著他，因信得以進入現在所站立的這恩典中，並且歡歡喜喜盼望神的榮耀。 \end{tabularx} \\ \\ \relax
5:3 & \begin{tabularx}{0.7\textwidth}{X} 不但如此，就是在患難中也是歡歡喜喜的，因為知道患難生忍耐， \end{tabularx} \\ \\ \relax
5:4 & \begin{tabularx}{0.7\textwidth}{X} 忍耐生老練，老練生盼望， \end{tabularx} \\ \\ \relax
5:5 & \begin{tabularx}{0.7\textwidth}{X} 盼望不至於落空，因為神的愛，已藉著所賜給我們的聖靈，澆灌在我們心裡。 \end{tabularx} \\ \\ \relax
5:6 & \begin{tabularx}{0.7\textwidth}{X} 我們還軟弱的時候，基督就在特定的時刻為不敬虔之人死。 \end{tabularx} \\ \\ \relax
5:7 & \begin{tabularx}{0.7\textwidth}{X} 為義人死，是少有的；為仁人死，或者有敢做的。 \end{tabularx} \\ \\ \relax
5:8 & \begin{tabularx}{0.7\textwidth}{X} 惟有基督在我們還作罪人的時候為我們死，神的愛就在此向我們顯明了。 \end{tabularx} \\ \\ \relax
5:9 & \begin{tabularx}{0.7\textwidth}{X} 現在我們既靠著他的血稱義，就更要藉著他得救，免受神的憤怒。 \end{tabularx} \\ \\ \relax
5:10 & \begin{tabularx}{0.7\textwidth}{X} 因為我們作仇敵的時候，尚且藉著神兒子的死得以與神和好，既已和好，就更要因他的生得救了。 \end{tabularx} \\ \\ \relax
5:11 & \begin{tabularx}{0.7\textwidth}{X} 不但如此，我們既藉著我們的主耶穌基督得以與神和好，也就藉著他以神為樂。 \end{tabularx} \\ \\
[1ex]
\hline
\hline
\end{longtable}
講題:因信稱義的因,恩,欣　Rejoice in Grace Justified by Faith

經文:羅馬書5章1-11節

講員:鄭漢文博士

場合:香港中文大學崇基學院禮拜堂

日期:2017年3月19日

1.「多得而樂」的故事

弟兄姊妹,平安.

今天,我又再回鄉探親一樣,返回崇基禮拜堂這個信仰大家庭.

今天都共聚在此,盼望聖道的臨在.

讓我先思考一個「多得而樂」故事

從前,有一個小朋友,他好快樂,天天歡欣地過日子.

平時食物其實好平常,但他食甚麼都咁開心.

有人問他「點解你咁歡欣?」(Why are you so joyful?)

他說:「因為有得食,慢慢地咀嚼,其實D食物都幾好食.」

家裡面有的玩具其實不多,但這個小朋友,玩得好開心.

有人問他「點解你咁歡欣?」

他說:「因為有得玩,以前玩具不多時,都會玩完又玩,其實D玩具都幾好玩.」

有機會時,這個小朋友會親親大自然,覺得一花一樹都好美麗,小鳥歌唱時他會細聽,就感到好好聽,蝴蝶飛舞在花間,他會細看,就感到好好看,總感到好歡欣.

有人問他「點解你咁歡欣?」

他說:「因為這些都不是我應得的,也就是我多得的,所以我會感到是我在得蒙恩待,所以我會歡欣喜樂.」(I rejoice in grace because I am blessed every day; I rejoice in grace because I am not deserved but blessed with all such blessings every day)

今天,讓我們思考我們在因信耶穌基督在上主面前得蒙恩待可以稱義而歡欣喜樂的信仰.

2.　因信稱義的論述:羅馬書5:1-11解經為五部份的結構

今日經課是《羅馬書》第五章第1至11節.

這段經文其實可以借助三個關鍵字(key terms)來分解當中的道,三個關鍵字就是我們今日講道的題目「因信稱義的因,恩,欣」(Rejoice in Grace Justified by Faith):

第一個關鍵字是「因」(cause/because):是一個邏輯字(logical term),一個邏輯連接詞(logical connective),陳述思考論理的前因後果得以連接起來.

第二個關鍵字是「恩」(grace):是一個名詞(noun),描述人的存有狀態(state of being)以及這個存有狀態的性質.

第三個關鍵字是「欣」(rejoice/joyfulness):是一個動態詞(verb/adverbial),表述人的活動狀態(active state),或一種心靈動態(mental act).

讓我先由第一個關鍵字「因」入手.

先請大家數一數,和合修訂本的十一節經文,一共出現多少個「因」字?

明顯用上「因」字外的有6處,還有潛在的推理字眼,包括:「藉著」,「靠著」,共有8處,代表羅馬書作者保羅在陳述著一個關於「因信稱義」頗複雜的神學說明.

有解經家認為,這段經文其實有一個平行結構,就是5章1至8節與5章9至11節是平行的.即是第1至3節的開首與9至11節的回首是對應地說明「因信稱義」,然後3至5節說明「歡欣何在」,6至8節說明「恩典何在」.

我們可以把這段經文分為五部份.分別是:

一．5:1-2:第1至2節開首地說明「因信稱義」,然後轉入「不但如此,…….」

二．5:3-5:「不但如此,…….」3至5節說明「歡欣何在」

三．5:6-8:6至8節說明「恩典何在」

四．5:9-10:第9至11節是回首地說明「因信稱義」,然後轉入「不但如此,…….」

五．5:11:「不但如此,…….」

請大家聽我讀一次:

部份    和合修訂本   因信稱義的福

第一部份

「因信稱義」

1所以,我們既因信稱義,就藉著我們的主耶穌基督得以與神和好.

2我們又藉著他,因信得以進入現在所站立的這恩典中,並且歡歡喜喜盼望神的榮耀.

第二部份

「歡欣何在」

3不但如此,就是在患難中也是歡歡喜喜的,因為知道患難生忍耐,

4忍耐生老練,老練生盼望,

5盼望不至於落空,因為神的愛,已藉著所賜給我們的聖靈,澆灌在我們心裡.

第三部份

「恩典何在」

6我們還軟弱的時候,基督就在特定的時刻為不敬虔之人死.

7為義人死,是少有的；為仁人死,或者有敢做的.

8惟有基督在我們還作罪人的時候為我們死,神的愛就在此向我們顯明了.

第四部份

「因信稱義」

9現在我們既靠著他的血稱義,就更要藉著他得救,9c免受神的憤怒.

10因為我們作仇敵的時候,尚且藉著神兒子的死得以與神和好, 既已和好,就更要因他的生得救了.

第五部份

「因信稱義」

11不但如此,我們既藉著我們的主耶穌基督得以與神和好,也就藉著他以神為樂.

3.　因信稱義中「因」的陳述:羅馬書5:1-3與9-11平行論述的結構

讓我們先把經文的這個平行結構的開首與回首讀出來,即是第1至3節與9至11節是對應地說明「因信稱義」.

我想請大家由第1節讀起讀到第3節,但讀法是請大家每讀一段話,每見到一個標點符號就停下來,等我讀出第9至11節相應的平行部份,首尾相和:

-------------------------------------------------------------------------------------------------------

|         會眾讀出5: 1-8                   |                講員讀出5: 9-11                           |

--   ----------------------------------   ---   --------------------------------------------------   --

|                                          |                                                          |

| 1所以,                                  | 9現在                                                    |

| (A)    1a我們既因信稱義,                | (a)  9a我們既靠著他的血稱義,                            |

|                                          |                                                          |

| (B1) 1b就藉著我們的主耶穌基督            | (b1) 9b就更要藉著他                                      |

|                                          |                                                          |

| (B2) 1c得以與神和好.                    | (b2) 9c得救,9c免受神的憤怒.                            |

|                                          |                                                          |

|                                          | 10因為                                                   |

|                                          |                                                          |

| (B3) [解釋B2]                            | (b3) [解釋b2]                                            |

|                                          |                                                          |

| 2我們又藉著他,因信                      | 10a我們作仇敵的時候,尚且藉著神兒子的死得以與神和好,    |

|                                          |                                                          |

| 2a得以進入現在所站立的這恩典中,         | 10a既已和好,就更要因他的生得救了.                      |

|                                          |                                                          |

| 2b並且歡歡喜喜盼望神的榮耀.             |                                                          |

|                                          |                                                          |

--   ----------------------------------   ---   --------------------------------------------------   --

| (C) 3不但如此,                          | (c ) 11不但如此,                                        |

--   ----------------------------------   ---   --------------------------------------------------   --

|                                          |                              |                           |

| 就是在患難中也是歡歡喜喜的,             | 2b我們又藉著他,因信……       | 11b也就藉著他以神為樂.   |

| 因為知道患難生忍耐,                     |                              |                           |

|                                          |   並且歡歡喜喜盼望神的榮耀   |                           |

| 4忍耐生老練,老練生盼望,                |                              |                           |

|                                          |   (10b; 11b)                 |                           |

| 5盼望不至於落空,因為神的愛,            |                              |                           |

| 已藉著所賜給我們的聖靈,                 |                              |                           |

| 澆灌在我們心裡.                         |                              |                           |

|                                          |                              |                           |

--   ----------------------------------   ---   --------------------------------------------------   --

|                                          |                              |                           |

| 6我們還軟弱的時候,基督就在特定的        | 2a我們又藉著他,因信得以     | 11a我們既藉著我們的主     |

|  時刻為不敬虔之人死.                    |                              |                           |

|                                          |   進入現在所站立的這恩典中, |    耶穌基督得以與神和好, |

| 7為義人死,是少有的；                    |                              |                           |

|  為仁人死,或者有敢做的.                |   (10a; 11a)                 |                           |

|                                          |                              |                           |

| 8惟有基督在我們還作罪人的時候            |                              |                           |

|  為我們死,神的愛就在此向我們顯明了.    |                              |                           |

|                                          |                              |                           |

--   ----------------------------------   ---   --------------------------------------------------   --

第3節和第11 節都是用「不但如此」作為此句與上述句子的連接,是平行的引伸,引伸出另外兩個主線:3至5節之說明「歡欣何在」,以及6至8節之說明「恩典何在」.

先說神恩,再談歡欣.

4.　因信稱義中的「恩」的描述:羅馬書5:6-8累進結構來類比神

四十年前我在高中信耶穌的時候,就深受羅馬書第5章第6至8節所打動,每次讀出都想落淚.

不過當時我們讀的,背誦的都是和合本.而和合本在第5章第6節是用一個「因」字來開頭,請讓我讀出:

5:6因我們還軟弱的時候,基督就按所定的日期為罪人死.

5:7為義人死,是少有的；為仁人死, 或者有敢做的.

5:8惟有基督在我們還作罪人的時候為我們死,神的愛就在此向我們顯明了.

第5章第6至8節經文用了一個「累進結構」來類比「神對人的大愛」.

這裡有三種人,值不值得為他們而死?:

一)仁人

二)義人

三)罪人

這三類人,值不值得為他們而死?暫時不答.

讓我們先類比現今社會有以下三種人,值不值得為他們而死?

i)富有人

ii)有權人

iii)弱勢人

「值不值得為富有人而死?」

我們天天為老闆打工,為他謀財,操勞致死,為的是分享一些勞動而得的財富,賴以維生,養活家人.我們沒有說,為富者一定不仁.如果富者是仁人,富有者可能好憐憫,你會為為富的仁人而死嗎?

「值不值得為有權人而死?」

當然有人會甘願為當權者賣命,為的是想分有權力.我們平民百姓無法分有合理權能(例如:因為無真普選),被有權有勢的人所欺凌,平民百姓只希望在掌權者下避禍──「山高皇帝遠」,為的是分有一點「帝力於我何有哉」的太平日子.我沒有說,當權者一定不義.如果當權者是義人,當權者可能行公義,你會為當權的義人而死嗎?

「值不值得為弱勢人而死?」

我們每天都可能被城市中弱勢人士的受苦受難所打動,也會因為有些人願意送出「一杯涼水」而動容.至於自己,送一些食物到食物銀行,或送暖到街頭露宿者也會有的,但為弱勢人士而送命的是極其少有.

回到羅馬書第5章提到的這三類人,值不值得為他們而死?:

一是仁人:仁者愛也,有仁心的人,就是他好憐憫的心,是我們所景仰的.

二是義人:義者宜也,有義行的人,就是他行公義的善行非常合宜,是我們值得效法的.

三是罪人:罪者惡也,得罪神的人,就是他的心有惡不容於人與神,他的行為十分不敬虔.

這三種人,值不值得為他們而死?

和合本在第5章第6節以一個「因」字來開頭,說出基督對我們人性軟弱的體諒.

請讓我再讀出一次,在我讀的時候,請你想一想:

是因為你的仁心?基督為你而死──

是因為你的義行?基督為你而死──

還是因為你的罪性……:

5:6因我們還軟弱的時候,基督就按所定的日期為罪人死.

5:7為義人死,是少有的；為仁人死, 或者有敢做的.

5:8惟有基督在我們還作罪人的時候為我們死,神的愛就在此向我們顯明了.

因信稱義中的「欣」的表述:羅馬書5:3-5以神為樂

現在,讓我們再看看《羅馬書》第五章如何談及因信稱義而蒙恩所得到的歡欣.

請大家從第1節讀到第5節,然後以第11節作總結,讀出相關的六段經文,你們不要停,我會重疊讀出每1, 2, 5, 11節的最後一句,讓我試一試效果.

《羅馬書》第五章

1所以,我們既因信稱義,就藉著我們的主耶穌基督得以與神和好.

2我們又藉著他,因信得以進入現在所站立的這恩典中,並且歡歡喜喜盼望神的榮耀.

3不但如此,就是在患難中也是歡歡喜喜的,因為知道患難生忍耐,

4忍耐生老練,老練生盼望,

5盼望不至於落空,因為神的愛,已藉著所賜給我們的聖靈,澆灌在我們心裡.

11不但如此,我們既藉著我們的主耶穌基督得以與神和好,也就藉著他以神為樂.

耶穌說,讓小孩子到我這裡來,因為天國是他們的.

讓我們返回開首的故事.

從前,有一個基督信徒,他好快樂,天天歡欣地過日子.

平時食物其實好平常,但他食甚麼都咁開心.有人問他「點解咁歡欣」(Why are you so joyful?),他說:「因為有得食,慢慢地咀嚼,其實D食物都幾好食.」

他歡欣,因為他欣賞到食物的內在價值(take pleasure in eating it itself),就是好好地食.

他歡欣,因為他欣賞到有食物可食的恩待,是大地有產出,存著感恩的心,好好地食.

從前,有一個基督信徒,他好快樂,天天歡欣地過日子.

家裡面其實工具不多,但他用得好開心.有人問他「點解咁歡欣」,他說:「因為有得用,以前用具不多時,都會用完又用,其實D工都幾好用.」

他歡欣,因為他欣賞到工具的實用價值(take pleasure in usefulness in itself),雖然生活有時依賴有形的工具,有得使用時就好好地用.

他歡欣,因為他欣賞到人生所得夠就夠,可以盡用,知足常樂,有得用時就好好地用.

從前,有一個基督信徒,他好快樂,天天歡欣地過日子.

有機會時,他會親親大自然,覺得一花一樹都好美麗,小鳥歌唱時他會細聽,就感到好好聽,蝴蝶飛舞在花間,他會細看,就感到好好看,

有人問他「點解你咁歡欣」,他說:「因為這些都不是我應得的,也就是我多得的,所有我會感到是我在得蒙恩待,所以我會歡欣喜樂.」

如果有人問你:「點解基督徒可以歡欣?」明明人生多苦困.

你是否可以這樣回答:「因為這些都不是我應得的,也就是我多得的,所有我會感到是我在得蒙恩待,所以我會歡欣喜樂.」(I rejoice in grace because I am blessed every day; I rejoice in grace because I am not deserved but blessed with all such blessings every day)

如何可能?

今天的道,就是說明我們基督徒的生命,都是因信稱義而經歷著:我們都是蒙恩的罪人,雖然不配得卻常存於心靈中的歡欣.

回到開首的「因」,如果「因信稱義」是因為我們自己的「信心」,即是由我們不信,未信到皈信,入信,好像完全是我們自身的努力,我們就很容易因為我們自身的心理強弱,信仰經歷,靈性高低而有所動搖,因而我之「被稱為義」好容易會出現動盪不安.

今天這段經課告訴我們,所有的「因」都是藉著我們的主耶穌基督,為我們還作罪人之時為我們死而讓我們得救,所有的「因」都是靠著我們的上主神的信實,因祂信守信約,就是祂願意放下責怪我們罪行的憤怒,甘願與我們和好.

我雖然是個罪人,因主耶穌基督救我,因神會守約的信約,我才白白得到蒙恩而被視為義,好能可以因著信而歡欣過我在世的日子,無論多苦多難.

人生在世,經受苦難.

面對這個香港這個時代,我們可以如何?

三年前6月我在這樣講了一堂道,是關乎民主公義的結構性罪惡(structural sin concerning democracy and justice),也是關於不讓中環價值佔領我們的心靈.

今天我又在香港一個追求民主公義的關鍵時機下講道,下個主日(3月26日)我作為一位選委要為香港投下一票選出新一任特首.

其實,香港的民主公義有沒有進步過?

人類的歷史,其實有沒有進步過?

今天是宗教改革五百周年.五百年前,馬丁路得改革教會,再一次提出「因信稱義」的信仰解讀.路德和其他改革者藉「五個唯獨」(Five Solas)去陳述上帝的真理,其中包括今天我們提及的兩個:「唯獨信心」(Sola Fide),「唯獨恩典」(Sola Gratia).

「唯獨信心」就是上帝藉耶穌的生和死所成就的愛的禮物,是要人藉相信去接受的.希伯來書十一章1-2節說:

「信就是所望之事的實底,是未見之事的確據.古人在這信上得了美好的證據.」(來十一1-2)

羅馬書四章3節又說:

「……亞伯拉罕信上帝,這就算為他的義.」(羅四3,參創十五16)

沒有條件,你只要信,就可得著耶穌為全人類所預備的.得救的信心像救恩一樣,也是上帝所賜的.

正如以弗所書二章8-9節所說:

「你們得救是本乎恩,也因著信.這並不是出於自己,乃是上帝所賜的；也不是出於行為,免得有人自誇.」(弗二8-9)

我們都是蒙恩的罪人.因著上主的信實,我們得蒙他賜的恩典,可以與神和好,活在歡欣的生活中.

讓我們一起低頭禱告:

「親愛的上主啊,

1我們既因信稱義,就藉著我們的主耶穌基督得以與神你和好.

2我們又藉著他,因信得以進入現在所站立的這恩典中,並且歡歡喜喜盼望你的榮耀.

3不但如此,就是在患難中也是歡歡喜喜的,因為知道患難生忍耐,

4忍耐生老練,老練生盼望,

5盼望不至於落空,因為神你的大愛,已藉著所賜給我們的聖靈,澆灌在我們心裡.

11不但如此,我們既藉著我們的主耶穌基督得以與神你和好,也就藉著他以神你為樂.」

但願榮耀都歸給聖父,聖子,聖靈,永在的三一神.阿們.

(鄭漢文博士現職香港中文大學教育行政與政策學系專業顧問,

同時在崇基學院神學院兼讀神道學學士課程.

曾於2012年1月至6月及2012年9月至2013年4月於崇基禮拜堂作校牧實習.)



\newpage

\section{我對你錯窒住盲堵的陰霾}
\label{sec:359}
\textbf{何崇謙牧師 - 我對你錯:「窒住盲堵」的陰霾}
\newline
\newline
\hyperref[sec:358]{< < < PREV SERMON < < <}
~
\hyperref[sec:toc]{[返目錄]}
~
\hyperref[sec:360]{> > > NEXT SERMON > > >}


創世記 2:1-17
\newline
\begin{longtable}{cl}
\hline
\hline
章節 & 經文 (和合本修訂版)\\
\hline
2:1 & \begin{tabularx}{0.7\textwidth}{X} 天和地，以及萬象都完成了。 \end{tabularx} \\ \\ \relax
2:2 & \begin{tabularx}{0.7\textwidth}{X} 到第七日，神已經完成了造物之工，就在第七日安息了，歇了他所做一切的工。 \end{tabularx} \\ \\ \relax
2:3 & \begin{tabularx}{0.7\textwidth}{X} 神賜福給第七日，將它分別為聖，因為在這日，神安息了，歇了他所做一切創造的工。 \end{tabularx} \\ \\ \relax
2:4 & \begin{tabularx}{0.7\textwidth}{X} 這就是天地創造的來歷。在耶和華神造地和天的時候， \end{tabularx} \\ \\ \relax
2:5 & \begin{tabularx}{0.7\textwidth}{X} 地上還沒有田野的草木，田間的菜蔬還沒有長出來，因為耶和華神還沒有降雨在地上，也沒有人耕種土地。 \end{tabularx} \\ \\ \relax
2:6 & \begin{tabularx}{0.7\textwidth}{X} 但是，有霧氣從地上騰，滋潤整個土地的表面。 \end{tabularx} \\ \\ \relax
2:7 & \begin{tabularx}{0.7\textwidth}{X} 耶和華神用地上的塵土造人，將生命之氣吹進他的鼻孔，這人就成了有靈的活人。 \end{tabularx} \\ \\ \relax
2:8 & \begin{tabularx}{0.7\textwidth}{X} 耶和華神在東方的伊甸栽了一個園子，把所造的人安置在那裡。 \end{tabularx} \\ \\ \relax
2:9 & \begin{tabularx}{0.7\textwidth}{X} 耶和華神使各樣的樹從土地裡長出來，可以悅人的眼目，好作食物。園子當中有生命樹和知善惡的樹。 \end{tabularx} \\ \\ \relax
2:10 & \begin{tabularx}{0.7\textwidth}{X} 有一條河從伊甸流出來，滋潤那園子，從那裡分成四個源頭： \end{tabularx} \\ \\ \relax
2:11 & \begin{tabularx}{0.7\textwidth}{X} 第一條名叫比遜，它環繞哈腓拉全地，在那裡有金子。 \end{tabularx} \\ \\ \relax
2:12 & \begin{tabularx}{0.7\textwidth}{X} 那地的金子很好，在那裡也有珍珠和紅瑪瑙。 \end{tabularx} \\ \\ \relax
2:13 & \begin{tabularx}{0.7\textwidth}{X} 第二條河名叫基訓，它環繞古實全地。 \end{tabularx} \\ \\ \relax
2:14 & \begin{tabularx}{0.7\textwidth}{X} 第三條河名叫底格里斯，它流到亞述的東邊。第四條河就是幼發拉底。 \end{tabularx} \\ \\ \relax
2:15 & \begin{tabularx}{0.7\textwidth}{X} 耶和華神把那人安置在伊甸園，讓他耕耘看管。 \end{tabularx} \\ \\ \relax
2:16 & \begin{tabularx}{0.7\textwidth}{X} 耶和華神吩咐那人說：「園中各樣樹上所出的，你可以隨意吃， \end{tabularx} \\ \\ \relax
2:17 & \begin{tabularx}{0.7\textwidth}{X} 只是知善惡的樹所出的，你不可吃，因為你吃它的日子必定死！」 \end{tabularx} \\ \\
[1ex]
\hline
\hline
\end{longtable}
講題:我對你錯:「窒住盲堵」的陰霾　Judgment in the Haze

經文:創世記2章1-17節

講員:何崇謙牧師

場合:香港中文大學崇基學院禮拜堂

日期:2017年3月26日

為甚麼神好像害怕人能「分辨善惡」,這不是有點違反常情嗎?難道「叫人能分辨善惡」,不是神造人的其中一個目的嗎?若捨棄「分別善惡」的進路,基督徒的人生又如何可以活出神的美善?

創世記記載始祖犯罪,是由於他們吃了「分辯善惡樹」上的果子；起先是夏娃受了魔鬼的誘惑吃,然後又給了她的丈夫亞當吃,聖經說,人便因此而墮落了,罪也因此進入了世界.然而,試想想,人類道德的起始點,不是懂得分辨善惡,並且擇善棄惡嗎?為甚麼懂得「分別善惡」,竟成了人墮落的因由?亞當夏娃的罪在何處?而這裡「分別善惡」,又意思何在?

原來,人類自有史以來所建立的社會,也流行著一種文化,這種文化可以稱之為「分別善惡樹」(或叫做「知善惡」)的文化.即或是在信仰的群體中也不例外,同樣會受到這種文化的影響.我們或許可以說,這正是罪帶來的影響.那麼,「分別善惡樹」的文化是甚麼意思?又有甚麼特色?

在這裡,「分別善惡」當中所指的「善」,是指神所看為好的「善」,「惡」也是指神以為「不好」或「惡」的事物.所以「分別善惡」,其實是指以神的標準來衡量一切.這是神獨有的本性,是神百分之百,自主自決權的屬性.人吃了「分別善惡樹」的果子,意味著他希望能像神一般,有創造主所獨有「分開／決斷」的主權──就如將光與暗分開,將空間以下和以上的水分開,將水與地分開,將物種的類別分開,將男女的性別分開等──這反映出人極盼望能像神一樣,或扮演神所扮演的角色(play God).

「善惡」,本來是一個法律上的用語,是指法官經過考慮,斷定甚麼是對,甚麼是錯.這就如神看見祂所創造的甚是「美好」,又判斷那人獨居「不好」,這就是知道「善」和「惡」的意思.然而,當人以受造者的身份,來承當這「分別善惡」的角色,就變成了一種獨斷,背離了以神作為最高的判決者(Judge),而只按自己的意志做決定,妄顧作為受造者的限制,也忽略了神和祂的啟示在人生中的地位.這樣的人生,造成了一種只以自己的心思,意志為準則,自我欺騙,自我滿足的心心態和行為.而表現出來,就是一種「我對你錯」的判別標準.

所以魔鬼說:「因為神知道,你們吃的日子眼睛就明亮了,你們便如神能知道善惡.」(創三5)而他們吃了果子後,神也確實說:「那人已經與我們相似,能知道善惡.」(創三22)但這並不是一條可喜的路,反而是一條通向死亡的道路!

因為人離開了神,就不再行在生命的路上,而「我對你錯」的文化,帶來的就是一種論斷(Judgmental)的文化.在這裡,我們也嘗試作出一個音意雙符的譯法,那是:「窒住盲堵」的文化.

1.「窒住盲堵」(Judgmental音意雙符的翻譯)的文化特色就是獨斷武斷

亞當夏娃開拓了這個「我對你錯」的文化,便叫後來的人類世世代代也陷在其中,人與人之間,人與神之間,都是我對你錯──男人說,是女人給我吃,而且是你造的女人的錯,因此是神你錯；女人說,是你造的蛇給我吃,所以又是神你錯……,這樣的心態,堵塞了神的命令和心意.聖經告訴我們,當人類社會落在「分別善惡」,「我對你錯的分割關係」時,其自然的結果就是害怕和厭棄神,也害怕和厭棄人,深刻破壞了這兩種關係；此外,亦害怕承認和面對自己的錯,由此而只採取逃避的態度,並以自欺來繞過自己的惡,掩蓋自己的錯,而將別人放在惡和錯的那一邊,形成一種看任何人和事,都以「判人於死地」的態度；甚至與神的關係,也隱藏著一種「我善你惡」的指控──為何要聽你的話,你的話一定對嗎?人懷疑神而不認識自己,這就是Judgmental,是「對人,對神,對己,採取窒住盲堵的行徑」,盲目地堵塞了其他人表白的機會,也讓人和神被拒於千里之外.

這論斷的作風和行為,必然產生兩種對人和對神的表現:Antagonistic,對抗性(音意雙符的翻譯是:「壓根兒來樹敵」)；Repulsive:易拋棄.

2.「硬著tie 控來樹敵」(Antagonistic音意雙符的翻譯)的文化特色是逢權必反

人類的第一個家庭,以及人後來的發展,都將神給我們的恩賜和恩典用在不適當的地方,特別是人與人之間同情共感,心靈相通的能力,互相守望及對此關係的防衛能力,勇士的戰鬥能力,防範外人破壞的能力等,都沒有用在對付撒旦方面,反而用於對付「自己人」.另外,又例如對美麗和知識的敏銳觸覺和渴望,也成了為自己的情慾服務以及獨斷自大的行為.

在家庭如此,在教會,機構等群體生活中也如此.當一個群體活在這種論斷的文化中,人的感情會只收藏不外露,整個群體的氣氛就會變十分冷漠,互不關心,失去彼此間的信任,也失去活力,只會彼此壓制；更有甚者是會出現一些違背福音真理的特質──律法主義,自以為是(因為沒有親密關係的互相勸誡,個人主義抬頭,並唯我獨尊),種族主義,階級主義,山頭主義,虛偽自恃等──造成另一種更差的文化,就是Antagonistic的文化──「壓根兒來樹敵」的文化.

3.「易拋棄」(Repulsive音意雙符的翻譯)的文化特色是遇困就御責

人犯罪後,對神為人所定的位份生出排斥和拋棄的心態,亦不曉得幫助別人成就其位份,甚至對最親近的人也予以排斥.本來神的心意是:「那人獨居不好,我要為他造一個配偶幫助他.」(創二18)後來是「你必戀慕你丈夫,你丈夫必管轄你.」(創三16)「戀慕」原文有「慫恿去犯罪」之意(urge, impulse, 創世記四章7節 說:「你若行得好,豈不蒙悅納,你若行得不好,罪就伏在門前,他必戀慕你,你卻要制伏他.」)家庭中,女人想要居於領導地位,而不是處在伙伴的關係上；男人亦不再視對方為自己骨中的骨般的關係來相處,卻要管轄女人.人類的文化發展始於家庭,家庭中的夫婦關係若落在「易拋棄」的境況中,整個社會,世界也就難免不是如此了!

其實基督徒生命的目標,不只是為主作大事,更應遵照主的吩咐,學習彼此相愛,彼此欣賞,彼此鼓勵,以致我們的生命更趨成熟,能有內在的轉化與更新,這才是真正的事奉；並叫人因此認出我們是主的門徒.我們信主後,和其他信徒走在一起時,有否彼此排斥?本來互為主內弟兄姊妹,有否趨於內向,將人可以接觸到自己的路徑封閉起來,以免受到操控或傷害?或是我們在豎立一座堅固的碉堡,去抗拒別人的愛,認為自己不值得被愛,別人亦不可愛?

生命樹的文化

和「窒住盲堵」的文化相比,生命樹的文化就很不一樣.吃生命樹的果子原是神的命令,是順從神的表現.一個群體若活在順從神的文化中,看到生命的來源是創造主自己,祂掌管我們的人生,是我們生命的主,一切好處和美善都源於祂；我們的好處和善亦遠遠不能與神相比,那麼我們一切的目光,就不再以「我對你錯,我善你惡」的分野作為立足點,可以漸漸地發展出「生命樹的文化」.

神所造的果子「悅人眼目」,吸引人想以它為食物,原本並不是罪,但若美感不能幫助人看見神的能力,恩典和愛,單單追求一己的滿足,這就是罪；而當人沉醉於「悅人眼目」的讚美,或落入「悅人眼目」的情慾,那也是罪.而當我們把信仰衰減為一把量尺,去量度別人的不足,也會落入法利賽人的窠臼,而讓恩典輕輕地溜走了.



\newpage

\section{我們還剩下甚麼}
\label{sec:360}
\textbf{林豪恩先生 - 我們還剩下甚麼}
\newline
\newline
\hyperref[sec:359]{< < < PREV SERMON < < <}
~
\hyperref[sec:toc]{[返目錄]}
~
\hyperref[sec:361]{> > > NEXT SERMON > > >}


以西結書 37:1-14
\newline
\begin{longtable}{cl}
\hline
\hline
章節 & 經文 (和合本修訂版)\\
\hline
37:1 & \begin{tabularx}{0.7\textwidth}{X} 耶和華的手按在我身上。耶和華藉著他的靈帶我出去，把我放在平原中，平原遍滿骸骨。 \end{tabularx} \\ \\ \relax
37:2 & \begin{tabularx}{0.7\textwidth}{X} 他使我從骸骨的四圍經過，看哪，平原上面的骸骨甚多，看哪，極其枯乾。 \end{tabularx} \\ \\ \relax
37:3 & \begin{tabularx}{0.7\textwidth}{X} 他對我說：「人子啊，這些骸骨能活過來嗎？」我說：「主耶和華啊，你是知道的。」 \end{tabularx} \\ \\ \relax
37:4 & \begin{tabularx}{0.7\textwidth}{X} 他又對我說：「你要向這些骸骨說預言，對它們說：枯乾的骸骨啊，要聽耶和華的話。 \end{tabularx} \\ \\ \relax
37:5 & \begin{tabularx}{0.7\textwidth}{X} 主耶和華對這些骸骨如此說：『看哪，我必使氣息進入你們裡面，你們就要活過來。 \end{tabularx} \\ \\ \relax
37:6 & \begin{tabularx}{0.7\textwidth}{X} 我要給你們加上筋，長出肉，又給你們包上皮，使氣息進入你們裡面，你們就要活過來；你們就知道我是耶和華。』」 \end{tabularx} \\ \\ \relax
37:7 & \begin{tabularx}{0.7\textwidth}{X} 於是，我遵命說預言。正說預言的時候，有響聲，看哪，有地震；骨與骨彼此接連。 \end{tabularx} \\ \\ \relax
37:8 & \begin{tabularx}{0.7\textwidth}{X} 我觀看，看哪，骸骨上面有筋，長了肉，又包上皮，只是裡面還沒有氣息。 \end{tabularx} \\ \\ \relax
37:9 & \begin{tabularx}{0.7\textwidth}{X} 耶和華對我說：「人子啊，你要說預言，向風說預言。你要說，耶和華如此說：氣息啊，要從四方而來，吹在這些被殺的人身上，使他們活過來。」 \end{tabularx} \\ \\ \relax
37:10 & \begin{tabularx}{0.7\textwidth}{X} 於是我遵命說預言，氣息就進入骸骨，骸骨就活過來，並且用腳站起來，成為極大的軍隊。 \end{tabularx} \\ \\ \relax
37:11 & \begin{tabularx}{0.7\textwidth}{X} 他對我說：「人子啊，這些骸骨就是以色列全家。他們說：『看哪，我們的骨頭枯乾了，我們的指望失去了，我們滅絕淨盡了！』 \end{tabularx} \\ \\ \relax
37:12 & \begin{tabularx}{0.7\textwidth}{X} 所以你要說預言，對他們說，主耶和華如此說：我的子民，看哪，我要打開你們的墳墓，把你們帶出墳墓，領你們進入以色列地。 \end{tabularx} \\ \\ \relax
37:13 & \begin{tabularx}{0.7\textwidth}{X} 我的子民哪，我打開你們的墳墓，把你們帶出墳墓時，你們就知道我是耶和華。 \end{tabularx} \\ \\ \relax
37:14 & \begin{tabularx}{0.7\textwidth}{X} 我必將我的靈放在你們裡面，你們就要活過來。我把你們安置在本地，你們就知道我---耶和華說了這話，就必成就。這是耶和華說的。」 \end{tabularx} \\ \\
[1ex]
\hline
\hline
\end{longtable}
詩篇 130:1-8
\newline
\begin{longtable}{cl}
\hline
\hline
章節 & 經文 (和合本修訂版)\\
\hline
130:1 & \begin{tabularx}{0.7\textwidth}{X} 耶和華啊，我從深處求告你！ \end{tabularx} \\ \\ \relax
130:2 & \begin{tabularx}{0.7\textwidth}{X} 主啊，求你聽我的聲音！求你側耳聽我懇求的聲音！ \end{tabularx} \\ \\ \relax
130:3 & \begin{tabularx}{0.7\textwidth}{X} 耶和華啊，你若究察罪孽，主啊，誰能站得住呢？ \end{tabularx} \\ \\ \relax
130:4 & \begin{tabularx}{0.7\textwidth}{X} 但在你有赦免之恩，要叫人敬畏你。 \end{tabularx} \\ \\ \relax
130:5 & \begin{tabularx}{0.7\textwidth}{X} 我等候耶和華，我的心等候；我也仰望他的話。 \end{tabularx} \\ \\ \relax
130:6 & \begin{tabularx}{0.7\textwidth}{X} 我的心等候主，勝於守夜的等候天亮，勝於守夜的等候天亮。 \end{tabularx} \\ \\ \relax
130:7 & \begin{tabularx}{0.7\textwidth}{X} 以色列啊，你當仰望耶和華，因耶和華有慈愛，有豐盛的救恩。 \end{tabularx} \\ \\ \relax
130:8 & \begin{tabularx}{0.7\textwidth}{X} 他必救贖以色列脫離一切的罪孽。 \end{tabularx} \\ \\
[1ex]
\hline
\hline
\end{longtable}
羅馬書 8:6-11
\newline
\begin{longtable}{cl}
\hline
\hline
章節 & 經文 (和合本修訂版)\\
\hline
8:6 & \begin{tabularx}{0.7\textwidth}{X} 體貼肉體就是死；體貼聖靈就是生命和平安。 \end{tabularx} \\ \\ \relax
8:7 & \begin{tabularx}{0.7\textwidth}{X} 因為體貼肉體就是與神為敵，對神的律法不順服，事實上也無法順服。 \end{tabularx} \\ \\ \relax
8:8 & \begin{tabularx}{0.7\textwidth}{X} 屬肉體的人無法使神喜悅。 \end{tabularx} \\ \\ \relax
8:9 & \begin{tabularx}{0.7\textwidth}{X} 如果神的靈住在你們裡面，你們就不屬肉體，而是屬聖靈了。人若沒有基督的靈，就不是屬基督的。 \end{tabularx} \\ \\ \relax
8:10 & \begin{tabularx}{0.7\textwidth}{X} 基督若在你們裡面，身體就因罪而死，靈卻因義而活。 \end{tabularx} \\ \\ \relax
8:11 & \begin{tabularx}{0.7\textwidth}{X} 然而，使耶穌從死人中復活的神的靈若住在你們裡面，那使基督從死人中復活的，也必藉著住在你們裡面的聖靈使你們必死的身體又活過來。 \end{tabularx} \\ \\
[1ex]
\hline
\hline
\end{longtable}
約翰福音 11:1-45
\newline
\begin{longtable}{cl}
\hline
\hline
章節 & 經文 (和合本修訂版)\\
\hline
11:1 & \begin{tabularx}{0.7\textwidth}{X} 有一個患病的人，名叫拉撒路，住在伯大尼，就是馬利亞和她姐姐馬大的村莊。 \end{tabularx} \\ \\ \relax
11:2 & \begin{tabularx}{0.7\textwidth}{X} 這馬利亞就是那用香膏抹主，又用頭髮擦他腳的；患病的拉撒路是她的弟弟。 \end{tabularx} \\ \\ \relax
11:3 & \begin{tabularx}{0.7\textwidth}{X} 姊妹兩個就打發人去見耶穌，說：「主啊，你所愛的人病了。」 \end{tabularx} \\ \\ \relax
11:4 & \begin{tabularx}{0.7\textwidth}{X} 耶穌聽見後卻說：「這病不至於死，而是為了神的榮耀，為要使神的兒子藉此得榮耀。」 \end{tabularx} \\ \\ \relax
11:5 & \begin{tabularx}{0.7\textwidth}{X} 耶穌素來愛馬大和她妹妹，以及拉撒路。 \end{tabularx} \\ \\ \relax
11:6 & \begin{tabularx}{0.7\textwidth}{X} 他聽見拉撒路病了，仍在原地住了兩天， \end{tabularx} \\ \\ \relax
11:7 & \begin{tabularx}{0.7\textwidth}{X} 然後對門徒說：「我們再到猶太去吧！」 \end{tabularx} \\ \\ \relax
11:8 & \begin{tabularx}{0.7\textwidth}{X} 門徒對他說：「拉比，猶太人近來要拿石頭打你，你還再到那裡去嗎？」 \end{tabularx} \\ \\ \relax
11:9 & \begin{tabularx}{0.7\textwidth}{X} 耶穌回答：「白天不是有十二小時嗎？人若在白天行走，就不致跌倒，因為他看見這世上的光。 \end{tabularx} \\ \\ \relax
11:10 & \begin{tabularx}{0.7\textwidth}{X} 人若在黑夜行走，就會跌倒，因為他沒有光。」 \end{tabularx} \\ \\ \relax
11:11 & \begin{tabularx}{0.7\textwidth}{X} 耶穌說了這些話，隨後對他們說：「我們的朋友拉撒路睡了，我去叫醒他。」 \end{tabularx} \\ \\ \relax
11:12 & \begin{tabularx}{0.7\textwidth}{X} 門徒就說：「主啊，他若睡了，就會好的。」 \end{tabularx} \\ \\ \relax
11:13 & \begin{tabularx}{0.7\textwidth}{X} 耶穌說這話是指拉撒路死了，他們卻以為他是指通常的睡眠。 \end{tabularx} \\ \\ \relax
11:14 & \begin{tabularx}{0.7\textwidth}{X} 於是耶穌就明白地告訴他們：「拉撒路死了。 \end{tabularx} \\ \\ \relax
11:15 & \begin{tabularx}{0.7\textwidth}{X} 為了你們的緣故，我不在那裡反而歡喜，為要使你們信。現在我們到他那裡去吧。」 \end{tabularx} \\ \\ \relax
11:16 & \begin{tabularx}{0.7\textwidth}{X} 於是那稱為低土馬的多馬對其他的門徒說：「我們也去和他同死吧！」 \end{tabularx} \\ \\ \relax
11:17 & \begin{tabularx}{0.7\textwidth}{X} 耶穌到了，知道拉撒路在墳墓裡已經四天了。 \end{tabularx} \\ \\ \relax
11:18 & \begin{tabularx}{0.7\textwidth}{X} 伯大尼離耶路撒冷不遠，約有六里路。 \end{tabularx} \\ \\ \relax
11:19 & \begin{tabularx}{0.7\textwidth}{X} 有好些猶太人來看馬大和馬利亞，要為她們弟弟的緣故安慰她們。 \end{tabularx} \\ \\ \relax
11:20 & \begin{tabularx}{0.7\textwidth}{X} 馬大聽見耶穌來了，就出去迎接他；馬利亞卻仍然坐在家裡。 \end{tabularx} \\ \\ \relax
11:21 & \begin{tabularx}{0.7\textwidth}{X} 馬大對耶穌說：「主啊，你若早在這裡，我弟弟就不會死了。 \end{tabularx} \\ \\ \relax
11:22 & \begin{tabularx}{0.7\textwidth}{X} 我也知道，即使現在，你無論向神求甚麼，神也必賜給你。」 \end{tabularx} \\ \\ \relax
11:23 & \begin{tabularx}{0.7\textwidth}{X} 耶穌對她說：「你弟弟會復活的。」 \end{tabularx} \\ \\ \relax
11:24 & \begin{tabularx}{0.7\textwidth}{X} 馬大對他說：「我知道在末日復活的時候，他會復活。」 \end{tabularx} \\ \\ \relax
11:25 & \begin{tabularx}{0.7\textwidth}{X} 耶穌對她說：「復活在我，生命也在我。信我的人雖然死了，也必復活。 \end{tabularx} \\ \\ \relax
11:26 & \begin{tabularx}{0.7\textwidth}{X} 凡活著信我的人必永遠不死。你信這話嗎？」 \end{tabularx} \\ \\ \relax
11:27 & \begin{tabularx}{0.7\textwidth}{X} 馬大對他說：「主啊，是的。我信你是基督，是神的兒子，就是那要臨到世界的。」 \end{tabularx} \\ \\ \relax
11:28 & \begin{tabularx}{0.7\textwidth}{X} 馬大說了這話就回去，叫她妹妹馬利亞，私下說：「老師來了，他在叫你。」 \end{tabularx} \\ \\ \relax
11:29 & \begin{tabularx}{0.7\textwidth}{X} 馬利亞聽見了，急忙起來，到耶穌那裡去。 \end{tabularx} \\ \\ \relax
11:30 & \begin{tabularx}{0.7\textwidth}{X} 那時，耶穌還沒有進村子，仍在馬大迎接他的地方。 \end{tabularx} \\ \\ \relax
11:31 & \begin{tabularx}{0.7\textwidth}{X} 那些同馬利亞在家裡安慰她的猶太人，見她急忙起來，出去，就跟著她，以為她要往墳墓那裡去哭。 \end{tabularx} \\ \\ \relax
11:32 & \begin{tabularx}{0.7\textwidth}{X} 馬利亞到了耶穌那裡，看見他，就俯伏在他腳前，對他說：「主啊，你若早在這裡，我弟弟就不會死了。」 \end{tabularx} \\ \\ \relax
11:33 & \begin{tabularx}{0.7\textwidth}{X} 耶穌看見她哭，並看見與她同來的猶太人也哭，就心裡悲嘆，又甚憂愁， \end{tabularx} \\ \\ \relax
11:34 & \begin{tabularx}{0.7\textwidth}{X} 就說：「你們把他安放在哪裡？」他們對他說：「主啊，請你來看。」 \end{tabularx} \\ \\ \relax
11:35 & \begin{tabularx}{0.7\textwidth}{X} 耶穌哭了。 \end{tabularx} \\ \\ \relax
11:36 & \begin{tabularx}{0.7\textwidth}{X} 猶太人就說：「你看，他多麼愛他！」 \end{tabularx} \\ \\ \relax
11:37 & \begin{tabularx}{0.7\textwidth}{X} 其中有人說：「他既然開了盲人的眼睛，難道不能叫這人不死嗎？」 \end{tabularx} \\ \\ \relax
11:38 & \begin{tabularx}{0.7\textwidth}{X} 耶穌又心裡悲嘆，來到墳墓前。那墳墓是個穴，有一塊石頭擋著。 \end{tabularx} \\ \\ \relax
11:39 & \begin{tabularx}{0.7\textwidth}{X} 耶穌說：「把石頭挪開！」那死者的姐姐馬大對他說：「主啊，他現在必定臭了，因為他已經死了四天了。」 \end{tabularx} \\ \\ \relax
11:40 & \begin{tabularx}{0.7\textwidth}{X} 耶穌對她說：「我不是對你說過，你若信就必看見神的榮耀嗎？」 \end{tabularx} \\ \\ \relax
11:41 & \begin{tabularx}{0.7\textwidth}{X} 於是他們把石頭挪開。耶穌舉目望天，說：「父啊，我感謝你，因為你已經聽了我。 \end{tabularx} \\ \\ \relax
11:42 & \begin{tabularx}{0.7\textwidth}{X} 我知道你常常聽我，但我說這話是為了周圍站著的眾人，要使他們信是你差了我來的。」 \end{tabularx} \\ \\ \relax
11:43 & \begin{tabularx}{0.7\textwidth}{X} 說了這些話，他大聲呼叫說：「拉撒路，出來！」 \end{tabularx} \\ \\ \relax
11:44 & \begin{tabularx}{0.7\textwidth}{X} 那死了的人就出來了，手腳都裹著布，臉上包著頭巾。耶穌對他們說：「解開他，讓他走！」 \end{tabularx} \\ \\ \relax
11:45 & \begin{tabularx}{0.7\textwidth}{X} 於是來看馬利亞的猶太人中，有很多人見了耶穌所做的事，就信了他。 \end{tabularx} \\ \\
[1ex]
\hline
\hline
\end{longtable}
題目:我們還剩下甚麼　What Remains of Us

經文:以西結書37:1-14;詩篇130:1-8;羅馬書8:6-11;約翰福音11:1-45

講員:林豪恩先生

場合:香港中文大學崇基學院禮拜堂

日期:2017年4月2日

一,引言

2016年11月25日,李怡先生以「我們還剩下甚麼」為題在崇基學院週會演講,他引用意見把香港分為兩個時期,其一是「努力興建」的時期,其二是「盡情破壞」的時期,他認為這描述概括道出了香港過去和現在的狀況.「我們所要面對的人生,不僅是畢業後薪資少到難以自己獨立生活,這情形在我那個時代一些低學歷的人也是常遇到的.但在「努力興建」的時代,再困難,再窮苦,社會和我們自己的人生都在「努力興建」,因此總有希望.而到了「盡情破壞」的時代……許多事無法確定,甚至是非也難確定.或者你有你的是非,我有我的.而且各執己見,各自固守自己的天地.不像「努力興建」的時代,大家都在公平的法律和規則之下,不同意見即使敵對意見也可以交流,在共同遵守的規則下各顯神通.現在是世事難料,規則常改,龍門會搬,意見分歧不是通過交流而取得進展,而是對立和撕裂.政見不同,影響到家庭和親友的關係,這是過去的時代不會有的……但是,人生需要面對,時代不容選擇.不是你想生活在怎樣的時代,而是時代選擇了你……人生有兩樣東西,一是困難,另一是困擾.困難是盡所有努力,包括找人合作,包括試驗一百次一千次,總有解決的希望；困擾是看不見解決希望,坐困愁城的煩惱.「努力興建」的時代是困難較多而困擾較少,法律與規則穩定,人們只要努力進取,克服困難,就會有回報或至少有希望,沒有許多掙扎和困擾.現在的世道迷離,是真假莫測,是非難辨,困擾多過困難,迫使我們不可以安居樂業過日子,我們不能不掙扎,不能不思考香港社會與個人前途的去向或去留.從生命價值來說,很難說前者的平庸好還是後者的掙扎好.但是,不論我們對這世道是悲觀還是樂觀,都必須積極面對.」

二,釋經及應用

覆巢之下,剩下哀和仇

今天,我們讀了四段經課.若我們以相同的問題問以西結先知:閣下生活的時代,你們還剩下甚麼呢?以西結書37章開首如此回答:

37:1耶和華的靈降在我身上.耶和華藉他的靈帶我出去,將我放在平原中；這平原遍滿骸骨.

37:2他使我從骸骨的四圍經過,誰知在平原的骸骨甚多,而且極其枯乾.

對以西結年代的猶大來說,他們還剩下甚麼呢?先知看見的,是遍滿已經枯乾的骸骨的荒野.

公元前 621 年,猶大王約西亞發動改革的第二年,有一位男嬰出生於耶路撒冷一個祭司家庭,取名為以西結.以西結成長於約西亞改革年代,親身見證國力和靈性的中興時期.可惜,約西亞在位第三十一年,被埃及法老尼哥打敗,重傷而死.隨著約西亞駕崩,宗教改革中斷.包括以西結在內的眾祭司們難免感到失望和挫折.

埃及法老尼哥把猶大變為藩屬國,廢除及擄走猶大百姓所立的約哈斯王,以約雅敬取而代之.根據耶和米先知的記載,這位被埃及法老所立的王,背叛耶和華,實行宗教混合主義,又嚴酷地壓制批判他的先知,甚至派人將逃命至埃及的先知烏利亞捉回耶路撒冷處死.耶利米幸得前朝重臣的兒子亞希甘保護才得保命.這一切,年青的以西結都看在眼內,對他產生的震撼及掙扎可想而知.

埃及打敗猶大之後十五年,北面的巴比倫國力強大得足以搶奪埃及的藩屬國.猶大夾在南北兩強之間,有時想臣服,有時又想背叛.約雅敬王臣服巴比倫三年之後終於背叛,巴比倫大軍壓境,約雅敬王死於戰亂中,他的兒子約雅斤繼位.約雅斤在位三個月,向巴比倫投降.尼布甲尼撒的軍隊洗劫聖殿和皇宮,又將王室,官員,名門望族,工藝技術人員都擄到巴比倫去,祭司以西結就是當時被擄的其中一位.

在被擄之地生活的猶大人,他們的心情是怎樣的呢?我們可以透過詩篇137這首詩歌感受一下:

1我們在巴比倫河邊,坐在那裡,追想錫安,就哭了.

2在一排柳樹中,我們掛上我們的豎琴.

3擄掠我們的在那裡要我們唱歌；搶奪我們的要我們為他們作樂:「給我們唱一首錫安的歌吧!」

4我們怎能在外邦之土唱耶和華的歌呢?

5耶路撒冷啊,我若忘記你,寧願我的右手枯萎；

6我若不記得你,不看你過於我最喜樂的,寧願我的舌頭貼於上膛!

7耶路撒冷攻破的日子,以東人說:「拆毀!拆毀!直拆到根基!」耶和華啊,求你記得!

8將要被滅的巴比倫哪,用你待我們的惡行報復你的,那人有福了.

9抓起你的嬰孩摔在磐石上的,那人有福了.

此情此景,他們還剩下甚麼呢?可能,只剩下兩樣:就是哀和仇.哀是哀自身:「我們在巴比倫河邊,坐在那裡,追想錫安,就哭了.」在巴比倫迦巴魯河邊被擄的猶太人的眼淚,南唐的李後主或者能夠產生共鳴,他在《破陣子》一詞中發出如此哀嘆:

四十年來家國, 三千里地山河.

鳳閣龍樓連霄漢, 玉樹瓊枝作煙蘿. 幾曾識干戈?

一旦歸為臣虜, 沈腰潘鬢消磨.

最是倉皇辭廟日,教坊猶奏別離歌. 垂淚對宮娥.

他們還剩下甚麼呢?除了哀,還有仇.仇不是憂愁的愁,而是仇恨的仇.他們仇恨的對象,是那些在他們痛苦是落井下石的人:「7耶路撒冷攻破的日子,以東人說:「拆毀!拆毀!直拆到根基!」耶和華啊,求你記得!」以及那些令他們落入如此田地的人:「8將要被滅的巴比倫哪,用你待我們的惡行報復你的,那人有福了.9抓起你的嬰孩摔在磐石上的,那人有福了.」

此時此地,我們還剩下甚麼呢?我們的城市,會否也彌漫著哀和仇的情緒呢?有人為香港哀悼,認為香港已經不是從前的香港,已經不是他們熟識的地方,甚至感到香港已死.與此同時,你感受到這個城市的仇怨嗎?不同意見人在網絡上及傳媒上互扣帽子,以陰謀論及咀咒的言語互相攻擊；不同立場的人在社會事件中鬥「晒馬」,甚麼使用肢體暴力攻擊異己.

豈有完卵,剩下的憂和悲

覆巢之下無完卵,如果哀和仇的情緒主導社會,生活於此的個體又剩下甚麼呢?詩人說:

130:1耶和華啊,我從深處向你求告!

深處指身體或靈性處於極深的痛苦,極度危機的狀態.陷入「深處」,令人有墮進黑暗,感到混亂,被黑暗勢力控制,或者被死亡的陰霾籠罩著.詩人沒有具體說明是甚麼境況,不過處身於外在環境惡劣,或者內心混亂的黑暗中掙扎,或多或少,相信每一個人都曾經有過這種「深處」的經歷!

福音書經課(約翰福音11章)具體地展示了其中一種「深處」的痛苦.在這「深處」,他們還剩下甚麼呢?

11:1有一個患病的人,名叫拉撒路,住在伯大尼,就是馬利亞和她姊姊馬大的村莊.

11:2這馬利亞就是那用香膏抹主,又用頭髮擦他腳的；患病的拉撒路是她的兄弟.

11:3她姊妹兩個就打發人去見耶穌,說:「主啊,你所愛的人病了.」

有人說,小病是福.是的,患過大病,或者長期病患,或者正在危疾當中的人,更加同意小病是福,因為他們及他們的家人非常明白疾病如何把他們拖入深淵,令他們感到自己所剩無幾.一位鄰近城市臨床期的醫學院學生有以下的實習分享:

「後來安排的穿刺切片,顯示的結果為惡性腫瘤,傾向認為是胃癌復發轉移到腹腔……最壞的結果發生,陳奶奶和家人的心情盪到谷底……她隨即問,能不能開刀……外科醫師看了電腦斷層影像後,認為無法以外科手術介入.原本希望能從外科醫師口中聽到一條明路的她,難掩落寞地低下頭,輕聲向醫師說聲謝謝,便陷入沉默.

稍後她收到住院期間醫療費用的通知單,她看了一眼,不禁大嘆一口氣……兒子為了照顧陳奶奶,最近把大樓管理總幹事的工作辭去了.因為生了這場病,要讓子女,媳婦輪流請假來醫院照顧她,家中經濟也不甚寬裕,擔心再這樣病下去會拖累家人,內心很是煎熬.我提及醫院有提供住院費用的補助,若有需要可以提出申請,但她只是搖搖頭沒再多說什麼.

後來,病情惡化……不安與絕望瞬間填滿了整間病房,眾人沉默不語,家屬內心的憂慮與不捨全表現在臉上,而任何醫療上的建議,似乎都無法改變結果.」

疾病不單把患者及家人拖進「深處」,也為醫生帶來了沉重的無力感:

「過去在醫學院我學到各式疾病的診斷與治療方法,但面對陳奶奶每況愈下的病情,雖然知道病因但也束手無策,我才親身體認到即使醫學發展到今日仍有其極限.而家屬更要面對他們摯愛的母親來日無多的現實,一股重重的無力感打在我們身上.」(鄭昶傑:「醫學之極限--實務上的煎熬與期待」).

讓我們回到約翰福音的故事,拉撒路病了,還剩下甚麼呢?

11:20馬大聽見耶穌來了,就出去迎接他；馬利亞卻仍然坐在家裡.

11:21馬大對耶穌說:「主啊,你若早在這裡,我兄弟必不死.

11:28馬大說了這話,就回去暗暗地叫她妹子馬利亞,說:「夫子來了,叫你.」

11:29馬利亞聽見了,就急忙起來,到耶穌那裡去.

11:30那時,耶穌還沒有進村子,仍在馬大迎接他的地方.

11:31那些同馬利亞在家裡安慰她的猶太人,見她急忙起來出去,就跟著她,以為她要往墳墓那裡去哭.

11:32馬利亞到了耶穌那裡,看見他,就俯伏在他腳前,說:「主啊,你若早在這裡,我兄弟必不死.」

11:33耶穌看見她哭,並看見與她同來的猶太人也哭,就心裡悲歎,又甚憂愁,

拉撒路最終沒有從病中康復,他離去了,剩下悲歎,剩下憂愁,剩下那流不盡的眼淚.眼淚,代表不能言說的感情,超越言詞的情懷.李清照在 《武陵春》一詞訴說:

風住塵香花已盡,日晚倦梳頭.物是人非事事休,欲語淚先流.

覆巢之下的新巢

以西結看見的,不是戰爭,因此已經沒有反敗為勝的機會；他看見的,是戰敗多年之後已枯乾的骸骨.在乾旱疲乏無力之地,還能開花嗎?在遍滿骸骨的荒野,還剩下甚麼呢?

37:3(耶和華對以西結說):「人子啊,這些骸骨能復活嗎?」

請留意,這不是以西結問的,而是耶和華向以西結提出的問題.在人而言,這是不能想像的狀況,因此也是多餘的問題.然而,耶和華卻以如此夢幻的問題來激發先知的想像.

37:4「你向這些骸骨發預言說:枯乾的骸骨啊,要聽耶和華的話.

37:5主耶和華對這些骸骨如此說:『我必使氣息進入你們裡面,你們就要活了.

37:6我必給你們加上筋,使你們長肉,又將皮遮蔽你們,使氣息進入你們裡面,你們就要活了；你們便知道我是耶和華.』」

37:9主對我說:「人子啊,你要發預言,向風發預言,說主耶和華如此說:氣息啊,要從四方而來,吹在這些被殺的人身上,使他們活了.」

被擄之民以為他們甚麼也沒有了,他們說『我們的骨頭枯乾了,我們的指望失去了,我們滅絕淨盡了.』(37:11)以西結卻看見,他們還有耶和華上帝,這位耶和華上帝,能夠『我的民哪,我必開你們的墳墓,使你們從墳墓中出來,領你們進入以色列地.』(37:12)也能夠「我必將我的靈放在你們裡面,你們就要活了.我將你們安置在本地,你們就知道我－耶和華如此說,也如此成就了.這是耶和華說的.』」(37:14)

破卵的重生

陷入深淵的詩人,還剩下甚麼呢?他還有一位呼求的對象,一位信靠的對象,一位讓他能夠等待的對象:

130:2主啊,求你聽我的聲音!願你側耳聽我懇求的聲音!

130:3主－耶和華啊,你若究察罪孽,誰能站得住呢?

130:4但在你有赦免之恩,要叫人敬畏你.

130:5我等候耶和華,我的心等候；我也仰望他的話.

130:6我的心等候主,勝於守夜的,等候天亮,勝於守夜的,等候天亮.

拉撒路敵不過死亡,他去世了,剩下甚麼呢?

11:33耶穌看見她哭,並看見與她同來的猶太人也哭,就心裡悲歎,又甚憂愁,

11:34便說:「你們把他安放在哪裡?」他們回答說:「請主來看.」

11:35耶穌哭了.

原來,剩下了耶穌的眼淚.那位創造天地萬物的上帝,也是悲天憫人的上帝,祂看見人間的疾苦,分擔著人們的悲歎和憂愁,與人一同灑淚,讓人感受到,也剩下了愛:

11:36猶太人就說:「你看他愛這人是何等懇切.」

拉撒路敵不過死亡,他去世了,剩下甚麼呢?

11:39耶穌說:「你們把石頭挪開.」那死人的姊姊馬大對他說:「主啊,他現在必是臭了,因為他死了已經四天了.」

11:40耶穌說:「我不是對你說過,你若信,就必看見上帝的榮耀嗎?」

11:41他們就把石頭挪開.耶穌舉目望天,說:「父啊,我感謝你,因為你已經聽我.

11:42我也知道你常聽我,但我說這話是為周圍站著的眾人,叫他們信是你差了我來.」

11:43說了這話,就大聲呼叫說:「拉撒路出來!」

11:44那死人就出來了,手腳裹著布,臉上包著手巾.耶穌對他們說:「解開,叫他走!」

11:45那些來看馬利亞的猶太人見了耶穌所做的事,就多有信他的；

拉撒路敵不過死亡,他去世了,剩下甚麼呢?剩下一位能夠戰勝死亡的耶穌,讓所有不能戰勝死亡的人,可以透過信靠這位耶穌而敢於面對死亡,不再懼怕死亡的威嚇.

總結

齊克果認為,死亡並非「致死的疾病」,反而絕望才是.我們今天誦讀的書信經課(羅馬書),保羅提醒我們,縱使有些事情不能逆轉,但卻不只是剩下哀傷和仇怨,也不致絕望,因為在上帝,還有更多的可能性,能夠支撐我們各種的生存狀況:

8:10基督若在你們心裡,身體就因罪而死,心靈卻因義而活.

8:11然而,叫耶穌從死裡復活者的靈若住在你們心裡,那叫基督耶穌從死裡復活的,也必藉著住在你們心裡的聖靈,使你們必死的身體又活過來.



\newpage

\section{住在主裡的重要性}
\label{sec:361}
\textbf{陳天祥博士 - 住在主裡的重要性}
\newline
\newline
\hyperref[sec:360]{< < < PREV SERMON < < <}
~
\hyperref[sec:toc]{[返目錄]}
~
\hyperref[sec:362]{> > > NEXT SERMON > > >}


講題:住在主裡的重要　The Importance of Abiding in Christ

講員:陳天祥博士

場合:香港中文大學崇基學院禮拜堂

日期:2017年4月9日

原委

耶穌來是要叫人得生命,並且得的更豐盛.我們成為基督徒,週日來聚會,所謂何事?是否要得生命,並且得的更豐盛呢?豐盛的屬靈生命應該是我們的目標,今天模成主的樣式,將來進入神的榮耀裡.聖經有許多經文都與這目標有關,其中有約翰15章.四福音在聖經中佔有很重要的位置,而福音書中記載耶穌親口的說話尤為重要.在四福音中,耶穌只有兩段較長篇的講論；太5-7記載耶穌在開始傳道時的講論,是編成的；而約翰13-17是主在離別前對門徒一次過長篇的講論,裝備門徒,應是非常的重要.15章在中間,住在主裡是十五章的主題,卻與13-17有許多共通點,我們在本章裡可以看見13章到底的愛的延續,及彼此相愛的命令；也可見14章的被差遣作主工；在本章後半可見16章有關門徒要面對恨他們的世界；同時,17章的差遣與合一,也躍現在本章中.因此,住在主裡是裝備門徒對自我對主及面對世界的一個重要的連接點.在今天的講論裡,我們會思想三個問題:在主裡是甚麼意思?為何必須住在主裡?以及如何住在主裡?

在主裡是甚麼意思?在這段經文裡,我們來看其中三個意思.

聖經在這裡是用一個比喻來說明,住在主裡如同枝子連在葡萄樹上,葡萄樹是主最後的一週內聖經提及的第三種樹,第一是主騎驢進耶路撒冷城時群眾在路舖上的樹枝及拿著葉的棕樹,它象徵著「義人的興旺」,是凱旋與享受成功歡樂的記號,就如詩128對耶和華的稱謝.第二是像徵以色列國運的無花果樹,因生長不正常而被耶穌認定會枯死,果然就枯死了.第三是葡萄樹,更是象徵以色列選民,是神所揀選和栽培的族類(詩80:8你從埃及挪出一棵葡萄樹,趕出外邦人,把這樹栽上) ,可惜變了野葡萄樹,有名無實,令神失望以致被毀.現另有葡萄樹,合神心意的,重新展現神的國度,因此耶穌自稱是真的葡萄樹.枝子在葡萄樹上就是在這新的國度裡有生命的交流,原文學者 Kenneth Wuest把「住在主裡」譯為maintain a living communion (保持一個生命的交流).枝子沒有了由根部吸收的水份來支持生命,就會枯死.保持生命的交流不單是指參加聚會,或是成為教會的成員.若在我們身上能活出主的樣式,如以賽亞所描述的謙虛的僕人及腓立比書所說主謙卑的心,或結出加拉太書所說的聖靈果子的特點,這些才是生命交流的結徵.即使是初信不太明顯,生命的交流也會構成轉化的動力,信徒會不斷在改變自己,學效主的樣式.

這生命交流的另一個特色是「常在」meno, 是留在,保持,住在之意,英文譯為abide,就是居住的意思,呂挀中譯本,較新的新漢語譯本,及環球聖經譯本都譯作住在；因此我們的題目是住在主裡,可以說是人按著神的恩典積極參與,把救恩俱體的在地上實踐出來.神在創世前在基督裡已揀選了我們,得蒙救贖也是在基督裡,當神看我們的時候,看不見我們的罪惡與虧欠,見到的只是基督的聖潔與公義,基督成了我們的聖潔公義和救贖.這是在主裡,而住在基督裡就是在我們生命中活出這個救恩的事實.「在主裡」是人因著神的恩得生命,「住在主裡」就是人要靠神的恩活出這生命；你若不是先在主裡就不能住在主裡,但不要只在主裡面而不住在主裡.請問你是否在主裡?請問你是否住在主裡?在主裡卻不住在主裡,就像有人有屋但走去訓街,也像有人在餐桌前不吃東西而捱餓,非常愚蠢.

第三個特點:經文說 (:4):「你們要常在我裡面,我也常在你們裡面.」不單是我們住在主裡,主也住在我們裡面.這是雙方面的互動,我們在主裡面,主也在我們裡面,這種互動幾乎可以說是神與人關係的常理,就如神是我們的產業,我們也是神的產業；保羅說他要抓著神所以抓著他的,一方面是神抓著保羅,而另方面是保羅抓著神,互相抓著；人親近神的時候,神就親近我們；我們在地上認他,他就在父面前認我們.神在等待我們行動,以致他更能與我們聯合.由於是互動的關係,若一方面獻上自己,不等於變為無有,反而變得更豐富.在約15章裡,互動的結果是人結出果子.『果子』在聖經中包括有 (1)生命的果子──與神生命相稱的表現 (太3:8)；(2),生活上的果子──透過聖靈活出各樣的美德 (加5:22-23)；(3)工作上的果子──指領人歸主 (羅1:13)；及(4)嘴唇上的果子 (來13:15).

為何必須住在主裡

任何動機都是要取得正面所需而避開負面的結果,住在主裡有正面的果效及逃避負面的結果.

2-1取得正面的果效:經文中提到三方面的果效:有屬靈生命及事奉工作的果效,以至多結果子(:8,16)；有加深與主的關係的果效,以至能住在主的愛中 (:9),而且由本是僕人關係提升至朋友(:14) ；這些都應是基督徒的人生目的,是基督徒所響往的.不但這樣,在主裡更有個人的果效:有主的喜樂(:11).這是甚麼的喜樂呢?主的喜樂應與一般的不同,就如主的食物或滿足與門徒的不同,主的食物是遵行神的旨意；主的平安與世人的不同,使在患難中不至憂愁膽怯,主的喜樂也應與一般的喜樂不同 .「喜樂」原文hara, 從字義來沒有甚麼不同,但按上文所列的屬靈生命果效,這種喜樂是建立在人生屬靈目的的實現.我以為與正向心理學所論的快樂 (Eudaimonia) 相似,這不是一般的歡笑 (happiness),而是做了正確的事而來的暢快.住在主裡,有屬靈生命及事奉的果效,在主的愛中,促成基督徒人生的目的,當然就心中暢快了.

更重要的是這個正面的果效能使自我超越.文中兩次說「不能」,:4說枝子不在葡萄樹上就不能結果子,是特殊的例子.而:5說門徒離了主,就不能作甚麼,是泛指一切的都不能作.因此,住在主裡就由不能至能,這就是心理學的超越自我 (transcedencce).它是指若一個人找到了一個尊崇的人,一個道理,或一個團體,於是把自己獻上,為他/它而活,就能活出超越自我的新意義.這個觀點,與15章的中心思想相似.我們為主而活,主就作我們的資源,就能超越自我.

2-2 避開負面的結果:有甚麼負面的結果呢?開始的時候,只是不能結果子,可能你以為沒有甚麼了不起,就如靈性沒有增長,心志沒有更新,性格沒有得更像主,沒有增加對聖經的明白及實踐,沒有更敏感及順服聖靈的引導,沒有對私慾加強克制的能力,沒有增加對神及對人的愛,也沒有盡力傳福音…等等,我們或許還可以稱為基督徒.但約15:2發出警告:凡屬我不結果子的,神就剪去.:6更說,人若不常在裡面,就像枝子在外面枯乾,人拾起來,扔在火裡燒了.其結果是變為無用而被火焚燒.這段經文在解釋上會有爭議,我們先來注意幾個字,第一個字是「剪去」,原文是airo, 英文是to lift, to take (拿走),在聖經中出現101次,最多的譯法是拿,共23次. 原來不是剪去,而是拿走,明顯已經部分與樹分離了.最令人關注的,是到底被火燒的枝子是否代表已得救的人,是否一次得救就永遠得救.經文指這些人是屬主的「凡屬我不結果子…」,這「屬」字似乎是指沒有活出神生命的『失敗信徒』,若這樣,曾經是信徒的人也會因離棄信仰而滅亡.但也有人強調人是否有生命只能從有沒有果子去分辨,若這樣,這些是指並沒有得救掛名的『假信徒』,真正信的,一定會有果子,他們一次得救就永遠得救.但我以為是否永遠得救並不是本章涉及的問題,比喻只應取其重點.其意思是指即使真信徒,斷絕與主相通,就不能有果子,被神看為廢物,也被人看為廢物.不過,希伯來書清楚說明,那些蒙了光照,嘗過天恩滋味,又於聖靈有分,並嘗過上帝善道的滋味,覺悟來世權能的人(他們應是得救的),若離棄真道,也會滅亡.無論如何,如果你發覺自己是一個不結果子的信徒,應趕快住在主裡,避免負面的結局出現.

如何住在主裡

要配合主的工作,人要委身.住在主裡的可能性是從神開始,而實踐性從人開始.:7『願意』thelo, determine 也可譯想要及立志,在原文的意思是指『意志主動的去要』.然而單靠人就不能作甚麼,要住在主裡,但如何住在主裡呢?大家一定很想知道.

方法只有一個,就是藉著道.在約15章裡,我們看見三重藉著道而住在主裡的功效: 道使人乾淨 (:3),道在人心裡(:7),道是要持守的誡命,就是彼此相愛(:10,12).

道使乾淨-:3「現在你們因我講給你們的道已經乾淨了」

主的話使門徒乾淨,這個乾淨與神修理乾淨不同.修理乾淨原文是kathairo, 英文是prune 其實只有修剪的意思,把不能生產的除掉.使門徒乾淨,這個乾淨的字卻與13:10 「你們是乾淨的,但都不是乾淨的」的「乾淨」是同字,原文是katharos, 英文是 clean, innocent (無辜). 耶穌指除猶大外,門徒都是乾淨的.因此這乾淨可比喻為得救的,是屬於主的,跟隨主是因主有永生之道而不是要取錢囊所存的.這是要住在主裡的首要條件；也就是說,人必須先在主裡,就是接受了耶穌做救主,才是乾淨的,才能住在主裡,有生命的交流.

道在心裡- :7「你們若住在我裡面,我的話也住在你們裡面」

我以為這是住在主裡面的最重要條件.道,就是神的說話,存在你的那一個範疇呢?是在外的儀文,是外在的規條,是教會的信經,還是已內化成你個人的價值觀與人生標準?Kenneth Wuest 把我的話也住在你們裡面譯作 「my word is at home in you」 就是我的話在你裡面安適,無拘束；這就是說,聖經的觀點,標準,價值觀都成了門徒的觀點,標準及價值觀；這就是心理學所說的內化.內化一般來說,有四個層面:外在不能接受的,外在可接受的,是自己認同的,成了自己人生的原則.由第一層到第四層,把在外的轉化成內在的,這就是內化.拿足球的規例來做例子,打茅波的足球員是在第一層,用盡方法去違背規例；有些球員雖然不喜歡這些規例,但在球場上也盡量遵守,這是第二層；有些球員在球場上也主張大家都要用這些規例,這是第三層；如果球員在沒有規例的場合中,仍會主動持守這些規例,這就是第四層了.把聖經內化成為自己處事的原則,單是聽道比較難做到；我的經驗是首先要有強烈的動機要在日常生活中遵行神的旨意,並在研經中把神的話深深地烙印在心內,這不是靈修式的看聖經,而是仔細研究聖經,這樣在日常生活中各樣的選擇,就容易從所研讀的聖經中冒出神的旨意,跟著就是要看人願否要遵行神的旨意而作出一些犧牲,若願意的話,聖經的觀點及價值觀,就漸漸內化成為門徒的觀點及價值觀了.不過,首先要有明白及遵行神旨意的心志,並有研經的生活,如果不能自己研經,就要把握機會經常參加研經班了.若你凡事都按內心醒覺的聖經行事,神的話在你內心無拘束地運作,使你凡事明白神的心意,你也會按神的旨意禱告,這樣,一句驚人的應許就出現了:(請大家同讀15:7,8節).你有沒有動機及毅力,讓這句聖經應許成就在你的身上,使你能超越你自己,在靈性的範疇能做到你現今不能做的呢?

道作為遵守的誡命 :10 「你們若遵守我的誡命,就住在我的愛裡」

7-8節道在人心裡似乎已達高嶧,為何在第10節再說遵守命令,像是衰退/迴歸呢?這當然不是衰退/迴歸,仔細看就發現9-17節是7-8節的重複,只是從不同的角度來重複.7節所指的「我的話」包括很廣,9-17特別以彼此相愛為中心,這是主離去後門徒面對逼害的環境非常重要的功課.同時住在主裡也換上住在主的愛裡.為甚麼呢?這與彼此相愛有關係呢?我以為神賜恩給人必須透過一些地上的機制,以色列人過紅海走乾地也因有強勁的東風,我們平常能謝飯是因地上有農業畜牧工業貸弊及服務業的機制.主愛我們,但若地上沒有一個愛的機制,主的愛也無法達到.彼此相愛就是地上的機制,讓神的愛能達到有需要的人.當然,彼此相愛也促進了門徒的成長或神的事工,在愛中互相勉勵,在愛中互相代禱,在神的工作上互相支援,在物質經濟上互相幫補,這樣,教會在世上才會興旺,果子就常存了(:16-17).

最近我收到一位弟兄的短訊,報告英國有500間教會關閉,同時有423所清真寺建成.我想:如果教會繼續分裂,繼續地盤主義,繼續反智,神的愛,智慧與能力就無法臨到教會,信徒也就不能作甚麼,被取代是預料中事.只盼望這段報告能喚醒地上的教會:彼此相愛.

結論: 約15章住在主裡的經文給予了我們甚麼呢?是一個

應許:是生命相交的應許,是靈性提昇的應許,是神人互動的應許,是超越自我的應許

警告:是被取去的警告,是被看為無用的警告,是被焚燒的警告

指引:乾淨的指引,道在心中的指引,遵守誡命彼此相愛的指引



\newpage

\section{復活的震撼生命存在於死亡之中}
\label{sec:362}
\textbf{許開明牧師 - 復活的震撼:生命存在於死亡之中!}
\newline
\newline
\hyperref[sec:361]{< < < PREV SERMON < < <}
~
\hyperref[sec:toc]{[返目錄]}
~
\hyperref[sec:363]{> > > NEXT SERMON > > >}


約翰福音 20:1-20
\newline
\begin{longtable}{cl}
\hline
\hline
章節 & 經文 (和合本修訂版)\\
\hline
20:1 & \begin{tabularx}{0.7\textwidth}{X} 七日的第一日清早，天還黑的時候，抹大拉的馬利亞來到墳墓，看見石頭已從墳墓挪開了， \end{tabularx} \\ \\ \relax
20:2 & \begin{tabularx}{0.7\textwidth}{X} 就跑來見西門‧彼得和耶穌所愛的那個門徒，對他們說：「有人從墳墓裡把主移走了，我們不知道他們把他放在哪裡。」 \end{tabularx} \\ \\ \relax
20:3 & \begin{tabularx}{0.7\textwidth}{X} 彼得和那門徒就出來，往墳墓去。 \end{tabularx} \\ \\ \relax
20:4 & \begin{tabularx}{0.7\textwidth}{X} 兩個人同跑，那門徒比彼得跑得快，先到了墳墓， \end{tabularx} \\ \\ \relax
20:5 & \begin{tabularx}{0.7\textwidth}{X} 低頭往裡看，看見細麻布還放在那裡，只是沒有進去。 \end{tabularx} \\ \\ \relax
20:6 & \begin{tabularx}{0.7\textwidth}{X} 西門‧彼得隨後也到了，進了墳墓，看見細麻布放在那裡， \end{tabularx} \\ \\ \relax
20:7 & \begin{tabularx}{0.7\textwidth}{X} 又看見耶穌的裹頭巾沒有和細麻布放在一起，是另在一處捲著。 \end{tabularx} \\ \\ \relax
20:8 & \begin{tabularx}{0.7\textwidth}{X} 然後先到墳墓的那門徒也進去，他看見就信了。 \end{tabularx} \\ \\ \relax
20:9 & \begin{tabularx}{0.7\textwidth}{X} 他們還不明白聖經所說耶穌必須從死人中復活的意思。 \end{tabularx} \\ \\ \relax
20:10 & \begin{tabularx}{0.7\textwidth}{X} 於是兩個門徒回自己的住處去了。 \end{tabularx} \\ \\ \relax
20:11 & \begin{tabularx}{0.7\textwidth}{X} 馬利亞卻站在墳墓外面哭。她哭的時候，低頭往墳墓裡看， \end{tabularx} \\ \\ \relax
20:12 & \begin{tabularx}{0.7\textwidth}{X} 看見兩個天使穿著白衣，在安放耶穌身體的地方坐著，一個在頭，一個在腳。 \end{tabularx} \\ \\ \relax
20:13 & \begin{tabularx}{0.7\textwidth}{X} 天使對她說：「婦人，你為甚麼哭？」她對他們說：「因為有人把我主移走了，我不知道他們把他放在哪裡。」 \end{tabularx} \\ \\ \relax
20:14 & \begin{tabularx}{0.7\textwidth}{X} 說了這些話，她轉過身來，看見耶穌站在那裡，卻不知道他是耶穌。 \end{tabularx} \\ \\ \relax
20:15 & \begin{tabularx}{0.7\textwidth}{X} 耶穌問她：「婦人，你為甚麼哭？你找誰？」馬利亞以為他是看園子的，就對他說：「先生，若是你把他移了去，請告訴我，你把他放在哪裡，我去把他移回來。」 \end{tabularx} \\ \\ \relax
20:16 & \begin{tabularx}{0.7\textwidth}{X} 耶穌對她說：「馬利亞。」馬利亞轉過身來，用希伯來話對他說：「拉波尼！」（「拉波尼」就是老師的意思。） \end{tabularx} \\ \\ \relax
20:17 & \begin{tabularx}{0.7\textwidth}{X} 耶穌對她說：「不要拉住我，因為我還沒有升上去見我的父。你到我弟兄那裡去告訴他們，我要升上去見我的父，也是你們的父，見我的神，也是你們的神。」 \end{tabularx} \\ \\ \relax
20:18 & \begin{tabularx}{0.7\textwidth}{X} 抹大拉的馬利亞就向門徒報信：「我已經看見了主。」她又把主對她說的話告訴他們。 \end{tabularx} \\ \\ \relax
20:19 & \begin{tabularx}{0.7\textwidth}{X} 那日（就是七日的第一日）晚上，門徒因怕猶太人，所在的地方門都關了。耶穌來，站在當中，對他們說：「願你們平安！」 \end{tabularx} \\ \\ \relax
20:20 & \begin{tabularx}{0.7\textwidth}{X} 說了這話，他把手和肋旁給他們看。門徒一看見主就喜樂了。 \end{tabularx} \\ \\
[1ex]
\hline
\hline
\end{longtable}
講題:復活的震撼:生命存在於死亡之中　The Impact of Resurrection: Life Exists Amid Death

經文:約翰福音20:1-20

講員:許開明牧師

場合:香港中文大學崇基學院禮拜堂

日期:2017年4月16日

耶穌不見了

1七日的第一日清早,天還黑的時候(看不見時候),抹大拉的馬利亞來到墳墓,看見石頭已從墳墓挪開了,

2就跑來見西門‧彼得和耶穌所愛的那個門徒,對他們說:「有人從墳墓裡把主移走了,我們不知道他們把他放在哪裡.」

3彼得和那門徒就出來,往墳墓去.

4兩個人同跑,那門徒比彼得跑得快,先到了墳墓,

5低頭往裡看,看見細麻布還放在那裡,只是沒有進去.

6西門‧彼得隨後也到了,進了墳墓,看見細麻布放在那裡,

7又看見耶穌的裹頭巾沒有和細麻布放在一起,是另在一處捲著.

8然後先到墳墓的那門徒也進去,他看見就信了.

9他們還不明白聖經所說耶穌必須從死人中復活的意思.

10於是兩個門徒回自己的住處去了.

向抹大拉的馬利亞顯現

馬利亞第二次的探索:感性的流著淚!

11馬利亞卻站在墳墓外面哭.她哭的時候,低頭往墳墓裡看,看不見主的身體

12看見兩個天使穿著白衣 [若無其事],在安放耶穌身體的地方坐著,一個在頭,一個在腳.

13天使對她說:「婦人,你為甚麼哭?」她對他們說:「因為有人把我主移走了,我不知道他們把他放在哪裡.」

14說了這些話,她轉過身來,看見耶穌站在那裡,卻不知道他是耶穌 [視若無睹].

15耶穌問她:「婦人,你為甚麼哭?你找誰?」馬利亞以為他是看園子的,就對他說:「先生,若是你把他移了去,請告訴我,你把他放在哪裡,我去把他移回來.」

馬利亞看不見:

16耶穌對她說:「馬利亞.」馬利亞轉過身來,用希伯來話對他說:「拉波尼!」(「拉波尼」就是老師的意思.)

17耶穌對她說:「不要拉住我,因為我還沒有升上去見我的父.你到我弟兄那裡去告訴他們,我要升上去見我的父,也是你們的父,見我的神,也是你們的神.」

18抹大拉的馬利亞就向門徒報信:「我已經看見了主.」她又把主對她說的話告訴他們.

引言

今日是2017年的復活節,願各位復活節快樂!復活節快樂,因為耶穌已快樂,耶穌復活帶來希望,勝利,生命.

自古到今,人們都在問:

1耶穌為何能夠復活?

2耶穌復活是真的嗎?有何証據?

3耶穌復活有何意義?

一,為何耶穌能夠復活?

這是很難解答的神學問題,簡單來說:因為生命在死亡之中；因為耶穌是上帝兒子,是彌賽亞,是基督具有不朽的生命；因為耶穌說:「我就是生命,我是復活(從死人中復活)!」祂的生命是永活,死亡不能吞噬復活的生命,正如黑暗不能吞噬光明.

今日教會經課,也試圖回答有關的問題?

二,耶穌復活是真的嗎?有何証據?

首先,耶穌為何第三天復活?因為耶穌的預言,他按照預言復活,顯示復活這事掌握上主的手中,對嗎?

a猶太人信念:一般人相信死者的靈魂,前三天在墳墓裡飛來飛去,但後來屍體腐爛使人認不出時,靈魂就離開他了.

b第三天是新一週的開始,之前是安息日,走遠路是違犯律法的.

所以最快來墳墓是七日的頭一日.七日的頭一日代表新的開始,新的創造,具神學意義.

在聖經的世界裡,聖經的作者為何確信耶穌復活是真的呢?

因為當中的三個象徵人物:1馬利亞(感性的象徵),2彼得(理性的象徵),3 那門徒(靈性信仰的象徵)

一 彼得:理性的象徵

根據希羅時代的觀念,男人是理性的,男孩也是有理性的；女人是無理性,或理性較弱的,所以有理性的男人,要管轄少理性的女人,保羅就在這背景下講「妻子當順服丈夫」.

上帝是理性的,是智者,是全知,全能,全在的；但根據柏拉圖的上帝,上帝是看不見的,是無肉身的,無情感的,因為上帝有血有肉有情的話,上帝便會被牽制.

1七日的第一日清早,天還黑的時候(看不見時候),抹大拉的馬利亞來到墳墓,看見( blepo:僅見,很快望一望)石頭已從墳墓挪開了,

2就跑來見西門‧彼得和耶穌所愛的那個門徒,對他們說:「有人從墳墓裡把主移走了,我們不知道他們把他放在哪裡.」

3彼得和那門徒就出來,往墳墓去.

4兩個人同跑,那門徒比彼得跑得快,先到了墳墓,

5低頭往裡看,看見細麻布還放在那裡,只是沒有進去.

6西門‧彼得隨後也到了,進了墳墓,看見 (Theoreo刻意的細心看,看完又看) 細麻布放在那裡,

7又看見 (Theoreo 刻意的細心看,看完又看) 耶穌的裹頭巾沒有和細麻布放在一起,是另在一處捲著.

空墳墓

石洞的入口處,則用餅狀的大石擋住,大石底部嵌放在一條岩床所鑿的軌槽內,須要身強力壯的人才能把大石滾開.誰挪開?要八個大男人才成,要說服八人去做這事很難.憎恨耶穌的人的所為?耶穌都死了,有仇都報了吧.門徒的所為?無可能,因為個個怕得要命,連耶穌都唔認了.挪開大石的目的是什麼?想不到.但偷屍是違法的,會被罰的,更是不尊重先人及其家人的.

裹屍布

細麻布還放在那裡

根據猶太人殯殮的習俗,是將屍身以長條細麻布加香料纏裹.所以令那門徒心裡希奇的事,可能是主的身體不在,但當初包裹的形狀還在,並不混亂.若是被人偷去,必將細麻布和身體一起挪走,或是把細麻布解開,則其形狀一定會弄得很亂.

裹頭巾是另在一處捲著

通常在鑿出來的墳墓裡放置屍體時,是頭部朝裡,腳部朝外(入口處),因此,人從墓口必先見到包裹屍身的細麻布,然後往裡面才會看到裹頭巾.

耶穌身體不見了,裹屍細麻布和裹頭巾卻保持人形,這顯示耶穌復活時,身體離開了那些布,因此裹屍細麻布和裹頭巾原封不動,這要証實耶穌不是被偷的,是自己復活的.因為偷死屍的人,最簡單是把死屍連裹屍布一起搬走；因為偷死屍的人,若只想偷耶穌身體,他不可能把裹屍布拆得如此整齊；因為包扎完好,看似無人動過.

作為有理性的人,以上兩點支持著耶穌身體不見了,不會是被偷.那唯一的可能性就是:耶穌生前所預言,他第三日復活了!這兩項證據,再加上門徒看見復活主的見證成為主復活的歷史證據.

二 那門徒:信仰的悟性象徵

3「彼得和那門徒就出來,往墳墓那裡去.」

4「兩個人同跑,那門徒比彼得跑的更快,先到了墳墓；」

「那門徒比彼得跑的更快,」顯然約翰比彼得年輕.

5「低頭往裡看,就見細麻布還放在那裡；只是沒有去.」

「看」刻意細心的看(表示觀看者希望看到一些重要的東西).

8「先到墳墓的那門徒也進去,看見 (eido 一看就明,明白,領悟真相,心智上理解) 就信了.」

那門徒的看,是心明,是悟性,相信耶穌復活,反而彼得仍是不明白.那門徒有信仰的悟性,這種悟性在信仰上深具意義.當信徒面對人生的種種人和事,好或壞,開心或不開心,成功或失望,平順或曲折,總會探索對我人生的意義和上帝的旨意,從而得到醒覺或悟性.人從理解,分析中,往往能得到悟性.可惜,時下急速的生活節奏,繁忙的工作,眾多的資訊,令人連思想的時間都無,悟性的引發,是要停一停,想一想才會出現.

三 馬利亞:感性的象徵

希伯來人向來重感性,人與外界的接觸透過肉體中眼,耳,口,鼻,引發內在的情感.上帝是智者,是有理性的,但上帝也是慈愛和信實的神,是有感性和感情的神,上帝差遣獨生子耶穌來世,使道成肉與人一起經歷悲歡離合.

馬利亞在本段經文所見是個感情豐富的人:

1七日的第一日清早,天還黑的時候(看不見時候),抹大拉的馬利亞來到墳墓,看見石頭已從墳墓挪開了,

2就跑來見西門‧彼得和耶穌所愛的那個門徒,對他們說:「有人從墳墓裡把主移走了,我們不知道他們把他放在哪裡.」

馬利亞哭過不停,飲泣的哭,她獨自一人,無助時就哭(11節),進入墳墓內,見到天使也哭(12-13節),轉過身來,見到一位男士,不知是耶穌,她又哭!(14-17節)

她傷痛飲泣的眼,加上與心中的成見(耶穌屍體被偷),令她看見細麻布,仍毫不理解；她看見兩天使,仍毫無醒覺；她看見耶穌,仍毫不知情,還誤認耶穌為園丁.感性的眼淚會使人失去理性和悟性,使人見不到真相,見不到天使同在,甚至見不到耶穌就在眼前.

然而,這約翰福音的作者是否否定馬利亞的飲泣?當然不是.

這些眼淚表達了愛,她深深愛耶穌,因為耶穌愛她,曾趕出她身上七隻鬼,使她得平安；因為她愛耶穌由耶穌被釘到死,埋葬入墳墓,她是最遲走的人,她也是安息日後,冒著被捉拿的危險,第一位來探耶穌的人.

淚水有益

淚水有益,每次眨眼睛時出現,它浸潤著我們的眼球.

在情感性眼淚中,當我們經受強烈的感情衝擊時,人體就會用淚水將多餘的化學物質「沖走」.

淚水能舒發悲傷,痛苦的情緒,無奈中亦有出路.

這就是我們在勸慰別人時,會告訴他們說:「盡情哭吧」.

不少男女哭完後,就充滿力量,面對明天.

眼淚是恩典

馬利亞的哭泣

引起天使和夫子耶穌的注意,對她說:「婦人,你為甚麼哭?」,這是出於同情的問候,因而拉近她與天使的關係.

眼淚把我們的交流層次帶到了語句和音節無法企及的高度.

沒有眼淚,人類無法傾心地,盡情地交流.

肉眼見不到主,卻在心中感受愛

耶穌看到馬利亞的眼淚,聽到她的飲泣,耶穌感受到馬利亞的愛,因此特別向她顯現.同樣,在我們的哀傷中,在我們的無助中,在我們看來愈毫無希望的處境中,你要相信復活的主,主見到,主聽到,當你有這份信心時,你會好像馬利亞說:「我已經看見主了!」

愛與關懷

門徒:理性代表,欠缺感性與感情

看見空墳,門徒知這馬利亞非常傷心,對馬利亞卻不發一言；看見馬利亞流淚,仍然不發一言.因為理性的人注重對真相的探索,重對問題的思考而忽略了對人的關懷,智商高的人,可以變得得冷酷無情,自以為是,看不起那些智商低,反應慢的人.在今日講求科技智能,智慧型的社會,追捧智能,強調精英主義,使自覺智商低的人,不敢發言,不敢為自己爭取利益.

復活主的愛

當馬利亞再入墳墓時,她先與天使相遇,再與復活主相遇:

復活主問:「婦人,你為甚麼哭?」但復活主不是說:「不要哭.」因為這是一種壓抑,於是無補；復活主不是說:「不准哭!」因為這是一種命令,使人更想哭；復活主不是說:「你哭了.」因為這是一種描述；而「婦人,你為甚麼哭?」是一種的關心,關心馬利亞哭的原因,然後找出路.

然後復活主問:「你找誰呢?」復活主知馬利亞來到墳墓的目的和需要 .耶穌知道馬利亞找誰,但馬利亞看不見耶穌,她仍在找死去的耶穌(耶穌的屍體)耶穌接著問:「你找誰?」其實想提醒她:「我就是你想找的人」嗎?馬利亞卻把耶穌當園丁,一個完全陌生的人,耶穌多難受呢!但耶穌沒有罵她,知道她肉眼和心眼都見不到主,所以耶穌地呼喚她的名字:「馬利亞!」馬利亞轉過身來,用希伯來話對他說:「拉波尼!」(「拉波尼」就是我的老師 )馬利亞轉過身來,就醒過來了,身轉,心也轉,所以叫了聲:「拉波尼!」馬利亞這隻羊聽到牧人的聲音,牧人也早已聽到馬利亞的聲音,知道馬利亞的心聲.這是多麼親切的相聚,馬利亞萬想不到想找的只是耶穌屍體,竟找到活著的耶穌身體；耶穌親切的「馬利亞」換來一句親切的回應「拉波尼」,這親切的稱呼,拉近了大家的關係.

希伯來人肯定物理世界外有另一個世界,信仰的世界是超越性和神秘性,與上帝建立關係強調耳朵部份的聽覺,著重心靈的眼睛,主觀的,神不需要有限的神像. 你試過在安靜中,想像上主呼喚你的名字嗎?在我面對人生的衝激時,我常想像主稱呼我的名字.我知道復活主,知道我,明白我,我便重新得力了.

不要摸我!

看見復活主,馬利亞隨即想摸和捉住耶穌,耶穌卻說「不要摸我!」意含:不要黏住我,不要抓住我不放,不要害怕(另有譯本),因為這是位復活的主,是神,這害怕是出於對復活主的「敬畏」,對一位至高的主的景仰.

馬利亞為何?

1要摸耶穌- 為要是道這是真耶穌嗎?

2要抓住耶穌-因為不想祂走

3要黏住耶穌-因為想親近主

靈修或默觀,不是停留在享受學耶穌一起的甜密,而是「去」,去作見證.馬利亞開心到捉住耶穌,想擁有耶穌,留住復活主.耶穌卻說:「不要摸我；因我還沒有升上去見我的父；你往我弟兄那裡去,告訴他們說:『我要升上去,見我的父,也是你們的父；見我的神, 也是你們的神.』

大家都有使命:復活主,要升上去見父；馬利亞,要見門徒,告知主復活的好消息.

馬利亞的見證:我已經看見主了

18抹大拉的馬利亞就去告訴門徒說:「我已經看見主了.」多麼興奮的話! 她又將主對她說的這話告訴他們.

這是多種的看見:肉眼看見,心眼看見,靈眼看見,並把主的存在心裡.

馬利亞說:「我已經看見了主!」這位復活主是怎樣的呢?

1祂是從死裡復活的主；

2祂是看見人眼淚的主；

3祂是認識她喚她名字的主；

4祂是可被親近 的主；

5祂是會差遣她去門徒傳好消息的主.

三,為何復活對我們有意義

從現代的哲學和心理學去看這段經文,近代哲學家康德認為人從自己的經驗和感觀去認識世界.

感性:人的感觀,接受外在事物的刺激.例如雪櫃很冷.

感情豐富的馬利亞,她見過主,耶穌跟她說話,感到主的愛,她親身看見復活主.

悟性:從接觸外在物,得到對物件的悟性.感悟可雪藏肉類,不變壞.

那門徒看見細麻布及裹頭巾,從信仰層面看,他就明白:不會是有人偷去主的身體,而是主自己復活.

理性:進行推理及邏輯思維.理性的探討.如何好好用雪作冷藏?

彼得見證他曾經看見 (詳細考察)

從理性層面:a空墳,b細麻布及c 裹頭巾,

從1感受雪很凍- 2感悟雪藏肉類不壞- 3 理性,思考如何用雪來冷藏?

4 心性:意志去完成,意志是人自覺地確定目的,並支配行動,克服困難,實現目的的心理過程.意,是心理活動的一種狀態；志,是對目的方向的堅信,堅持.意志,即對實現目的有方向,有信念地堅持的一種心理活動(含潛意識中的心理活動).

意志的建立

意志的建立有賴三方面:知(即認知),情(即情感,明白意義,產生動力),意(即意志)是人類心理活動.

人們認識世界的主觀意識過程,通常可以分為三個階段:

1認知階段

目的在於解決「是什麼」或「甚麼事實」的問題.人只有首先瞭解事物的外在特性(或外部聯繫)和內在規律(內在本質),即首先瞭解事物「是什麼東西」,才能對它進行其它方面的深入瞭解.

馬利亞認知:耶穌復活是事實!

2評價階段

目的在於解決「有何用」 或「有什麼價值」的問題.人只有瞭解在瞭解事物「是什麼東西」以及「對我有何價值」,才能知道如何對它採取正確的處理措施.

馬利亞評價:耶穌復活是好消息,耶穌是神,祂的預言對現,祂與我們同在.

3意志(或決策)階段

目的在於解決「怎麼辦」 或「實施什麼行為」的問題.就是針對事物的品質特性以及每一品質特性對於人的價值,人將選擇一個最合適的行為,以便能充分有效地利用事物的價值特性.

馬利亞立定意志聽復活主的吩咐,告訴門徒,眾人這好消息,有志者事竟成,復活的故事傳二千多年了.

各位,作為大學的知識份子,從看見感應現象,領悟現像,理解所悟的現象,相信所理解現象,從所建立的信念,有所發現或發明.

自古到今,許多信徒因相信一位復活的主,得到復活多信息,生活得到突破,更新,我是,你又如何?

復活的震撼,因石頭被上主所滾開,復活主從死人中復活因為生命在死亡之中,那就顯得有意義.耶穌說:「我就是生命,我是復活.(生命在我,死亡也在我的掌管!)」祂的生命是永活,死亡不能吞噬復活的生命,正如黑暗不能吞光明,生命在死亡之中,生命力在死亡之中.

生命力在主日

1 七日的頭一日,代表新的日子,新的開始；

2 清早,代表光明勝黑暗,良善勝罪惡,希望勝絕望；

3 石頭挪開,代表把一切障礙都除掉,死亡的石頭被移去；

4 身體復活,生命力量勝過死亡力量,無可能的變有可能；

5 相信復活主的人,從悲傷轉為喜樂,從消沉變積極；

因此,基督教每個主日都在紀念主的復活.

馬利亞原本在悲傷,情緒低落,眼睛模糊,對前境充滿悲觀中,與復活主相遇後,生命力在她像快要死亡中,積極起來了,變得有動力去作見證,傳出去,有使命；有活力；有能力,使死人復活,使悲傷人復活,使失去衝勁的人復活.

恐怖主義不獨屬於伊斯蘭國,美國及俄國大國也有,使人失去動力,使人死亡,死亡存在於活人當中,這正是撒旦,黑暗之子的工作.

面對今日香港,復活神學能否給我們一些亮光?

一國兩制,看來「唔制」的人愈來愈多,因為愈來愈傾斜於一制；五十年不變,由逐漸變到急變,還會有劇變嗎?民主,自由,法治,高度自治空間在縮小,基礎亦像在動搖；許多人都看得不樂觀,不安心,不存希望；許多人看得灰,大學生也是.

我是六十年代從內地到港人士,抵港不久遇溫黛巨風,那時在十號風球下,還跟媽媽到工廠返工,狂風大雨被嚇到面青,接著香港有暴動,戒嚴,制水,居住在沒有電梯的唐樓九樓的我,要到樓下一邊唱著歌一邊擔水回家,之後還經歷六四和九七回歸等.

雖然如此,我個人的感覺香港是愈來愈好,由沒有屋住,到住唐樓,公屋到現在住宿舍；從無書讀,到完成多個學位,社會提供了不少機會.

復活神學,告訴我們,死也有希望,光明定勝黑暗,生命存在於死亡之中,只要我們仍有生命,仍存希望,仍不失理想,仍堅持香港的核心價值,明天仍會更好.

基督徒相信復活主,他是歷史的主,全地的主,生命的主,他在掌管一切.

基督教信仰涉及感性,悟性,理性,德性,靈性及意志部份,使人生命得到充份的發揮；因祂活著,我能面對明.Because He live, I can face tomorrow!



\newpage

\section{既然蒙召}
\label{sec:363}
\textbf{朱建忠先生 - 既然蒙召}
\newline
\newline
\hyperref[sec:362]{< < < PREV SERMON < < <}
~
\hyperref[sec:toc]{[返目錄]}
~
\hyperref[sec:364]{> > > NEXT SERMON > > >}


以弗所書 4:1
\newline
\begin{longtable}{cl}
\hline
\hline
章節 & 經文 (和合本修訂版)\\
\hline
4:1 & \begin{tabularx}{0.7\textwidth}{X} 我為主作囚徒的勸你們，既然蒙召，行事為人就要與你們所蒙的呼召相稱。 \end{tabularx} \\ \\
[1ex]
\hline
\hline
\end{longtable}
講題:既然蒙召　The Calling of The Lord

講員:朱建忠先生

經文:以弗所書4:1

場合:香港中文大學崇基學院禮拜堂

日期:2017年4月23日

讀經:弗4:1 我為主被囚的勸你們,既然蒙召,行事為人就當與蒙召的恩相稱．

引言:

弟兄姊妹當我們讀這一處經文的時候,我們知道保羅當時已經身處監獄之中.但是他卻不因福音被囚而恥,也不因身體拘禁而悲.保羅在林後 4:7-10“我們有這寶貝放在瓦器裡,要顯明這莫大的能力是出於　神,不是出於我們.我們四面受敵,卻不被困住；心裡作難,卻不至失望；遭逼迫,卻不被丟棄；打倒了,卻不至死亡.身上常帶著耶穌的死,使耶穌的生也顯明在我們身上.”使徒保羅掛念的不是他自身安危或個人榮辱,乃是主的教會,他所牧養的群羊,在我們剛才所讀經文弗4:1他勸勉我們:“既然蒙召,行事為人就當與蒙召的恩相稱.”

那麼怎樣的行事為人才是與蒙召的恩相稱呢?聚會準時?詩歌唱的好?工作表現好?…實際上使徒保羅講到既然蒙召,蒙召後的行事為人就當與蒙召的恩相稱,所強調的乃是我們的所是!即我們是怎麼樣的人.這是一種根本性的,實質性的辨認.我們常講做學生要有學生的樣子,老師要有老師的樣子,法官要有法官的樣子,既是如此.你做一個基督徒就要有基督徒的樣子.我們知道聖經裡記載說“門徒被稱為基督徒是從安提阿開始的”(徒11:26),基督徒的原意就是小基督的意思,就是指可以讓世人在你身上看出耶穌基督來.所以,弟兄姊妹,當我們成為蒙召之人後,我們的行事為人就當與蒙召之恩相稱.

今天我們就要趁著上午證道的時間,從三個方面來思想蒙召之恩的意義,盼望我們借著經上的話語,能夠重拾蒙召真諦；反思事奉心志,活出一個與蒙召之恩相稱的生活.

蒙召何人?

我們先看一節聖經,這是當時復活的主差派保羅去做外邦的使徒時候所說的: “我差你到他們那裡去,要叫他們的眼睛得開,從黑暗中歸向光明,從撒旦權下歸向　神；又因信我,得蒙赦罪,和一切成聖的人同得基業.’(徒 26:18)反思我們歸主之前乃是:瞎眼的,在黑暗中的,撒旦權下的.”

一個人眼瞎,看不見,結果就是找不到出路.06年桑美颱風後蒼南之行的經歷:晚上拼車和人一起打的從靈溪到馬站,盤山公路,到了山頂時,遇上大霧,能見度只有一米,前面一片茫然,真不知車輪是在山路上還是在空中.幸好有車燈探路,路標引導,不然無法前行.我們屬靈之路亦是如此,聖經說我們的心眼被魔鬼弄瞎了,看不清出路在哪裡?唯一的出路乃是,被耶穌的真光照亮.約翰福音第一章告訴我們那光是真光,要照亮一切生在世上的人.約 9:39:耶穌又說:“我為審判到這世上來,叫不能看見的可以看見；能看見的反瞎了眼.”

第二個方面聖經告訴我們,我們原本是行在黑暗之中的.約3:19     光來到世間,世人因自己的行為是惡的,不愛光,倒愛黑暗,定他們的罪就是在此.聽一位同工講到北京一教會的信徒,他們對於天堂與世俗之樂只能選擇其一的憂愁.而初期教會:脫離卑賤的事,不效法這個世界.

第三個方面,這是名份的歸屬,根本性的辨識.弗2:2－3說:那時,你們在其中行事為人,隨從今世的風俗,順服空中掌權者的首領,就是現今在悖逆之子心中運行的【邪】靈.我們從前也都在他們中間,放縱肉體的私欲,隨著肉體和心中所喜好的去行,本為可怒之子,和別人一樣.――第4,5節:然而神既有豐富的憐憫,因他愛我們的大愛,當我們(還)死在過犯中的時候,便叫我們與基督一同活過來.

小結:蒙召何人?蒙召乃是一個蒙恩的罪人.

以賽亞蒙召:禍哉,我滅亡了!(賽6:1－5:當烏西雅王崩的那年,我見主坐在高高的寶座上．他的衣裳垂下,遮滿聖殿. 其上有撒拉弗侍立．各有六個翅膀．用兩個翅膀遮臉,兩個翅膀遮腳,兩個翅膀飛翔．彼此呼喊說,聖哉,聖哉,聖哉,萬軍之耶和華．他的榮光充滿全地.因呼喊者的聲音,門檻的根基震動,殿充滿了煙雲.那時我說,禍哉,我滅亡了．因為我是嘴唇不潔的人,又住在嘴唇不潔的民中．又因我眼見大君王萬軍之耶和華.)

――彼得:「主啊!離開我,我是個罪人!」(路5:8)(Luk 5:1－11  耶穌站在革尼撒勒湖邊,他見有兩隻船灣在湖邊；有一隻船是西門的,耶穌就上去,請他把船撐開,稍微離岸,就坐下,從船上教訓眾人.講完了,對西門說:「把船開到水深之處,下網打魚.」西門說:「夫子,我們整夜勞力,並沒有打著什麼.但依從你的話,我就下網.」他們下了網,就圈住許多魚,網險些裂開,甚至船要沉下去.西門彼得看見,就俯伏在耶穌膝前,說:「主啊!離開我,我是個罪人!」 耶穌對西門說:「不要怕!從今以後,你要得人了.」他們把兩隻船攏了岸,就撇下所有的,跟從了耶穌.)

著名聖詩《奇異恩典》作者,約翰牛頓.(詩歌及作者生平簡介:奇異恩典,何等甘甜,我罪已得赦免；前我失喪,今被尋回,瞎眼今得看見.他在七歲的時候母親去世,幼年隨父親航海,染上水手放蕩惡習,吃喝嫖賭,無所不為.後來利用海運做販賣奴隸生意,想不到因鬧出事來,自己反在非洲淪為奴隸,備受痛苦,幸虧父親知道,派人把他救回.在他乘船途中,忽遇大風浪,船幾乎沉沒,他知道自己一生離棄　神,於是認罪悔改,得以平安回國.重生以後,生命有了奇妙改變,獻身終生為主所用.)82歲退休時他說:很多事,我都已經忘記了,但有兩件事,我無法忘記:一,我是個大罪人,二,耶穌是我的救主.

親愛的弟兄姊妹,耶穌說,沒病的不需要醫生,有病的才需要醫生.讓我們大膽承認自己的罪,就像保羅一樣,“在罪人中我是罪魁,然而我蒙了憐憫,...是因耶穌基督要在我這罪魁身上顯明他一切的忍耐,給後來信他得永生的人作榜樣.”(提前1:15,16)蒙召何人:是蒙恩的罪人.

蒙召何由?

為什麼像你我這樣的罪人竟能進入神榮耀,聖潔之國度?有一個故事說,彼得職位管理天國之門,問及每個欲進入之人,憑據為何?...是三代遺傳的基督徒?賺錢三套豪華住宅?擁有三個學位?...

林前 1:26弟兄們哪,可見你們蒙召的,按著肉體有智慧的不多,有能力的不多,有尊貴的也不多.林前 1:27神卻揀選了世上愚拙的,叫有智慧的羞愧；又揀選了世上軟弱的,叫那強壯的羞愧.林前 1:28神也揀選了世上卑賤的,被人厭惡的,以及那無有的,為要廢掉那有的；林前 1:29使一切有血氣的,在　神面前一個也不能自誇.林前 1:30但你們得在基督耶穌裡是本乎　神,　神又使他(耶穌基督)成為我們的智慧,公義,聖潔,救贖.

弟兄姊妹,論到我們的主,聖經上說:(賽 53:2－8)他在耶和華面前生長如嫩芽,象根出於幹地.他無佳形美容,我們看見他的時候,也無美貌使我們羡慕他.他被藐視,被人厭棄,多受痛苦,常經憂患.他被藐視,好象被人掩面不看的一樣,我們也不尊重他.他誠然擔當我們的憂患,背負我們的痛苦；我們卻以為他受責罰,被　神擊打苦待了.哪知他為我們的過犯受害,為我們的罪孽壓傷；因他受的刑罰我們得平安,因他受的鞭傷我們得醫治.

蒙召何由,完全是耶穌捨命流血的恩典與愛把我們撒旦的權下買贖過來.撒下9:1－9“大衛問說,掃羅家還有剩下的人沒有,我要因約拿單的緣故向他施恩. 掃羅家有一個僕人,名叫洗巴．洗巴對王說,還有約拿單的一個兒子,是瘸腿的,叫米非波設.於是大衛王打發人去,從羅底巴亞米利的兒子瑪姬家裡召了他來.掃羅的孫子約拿單的兒子米非波設來見大衛,伏地叩拜．大衛說,米非波設．米非波設說,僕人在此.大衛說,你不要懼怕,我必因你父親約拿單的緣故施恩與你,將你祖父掃羅的一切田地都歸還你,你也可以常與我同席吃飯.”　聖經記載說“米非波設又叩拜說:僕人算甚麼,不過如死狗一般,竟蒙王這樣眷顧. ”弟兄姊妹,實際上我們與主的差別豈不是更比米非波設與大衛的差別更大嗎?然而我們竟蒙了何等大的恩典:當我們還是罪人的時候,耶穌基督為我們釘在十字架上,神的愛就在此顯明了.

蒙召何為?

非洲一宣教士野獸中救土著人的經歷,後來這位土著人帶著一家人並家當來做僕人...思考:我們對主也要將主權交托,因我們也是在魔鬼獅子口中被耶穌基督救出來的.

信主而活:

加 2:20我已經與基督同釘十字架,現在活著的,不再是我,乃是基督在我裡面活著；並且我如今在肉身活著,是因信　神的兒子而活,他是愛我,為我舍己.什麼是保羅所講的:“因信神的兒子而活?”乃是告訴我們若不是因信神的兒子――這一位愛我們,為我們舍己的主耶穌基督而活,那活著就沒有真正的意義.你如果真的認同你的老我已經與基督一同釘在十字架上,那現在活著的不是耶穌基督賜給你的新的生命又是什麼?當耶穌基督在我們裡面活著的時候,他就給我們一個新的生命,這個新生命它會追求聖潔,願意親近神,愛慕天家.保羅說:“我情願離世與主聯合,那是好的無比.”這一句話不僅僅是這樣說說而已,乃是由信仰熱忱而發.這不僅在保羅身上我們可以看到,在初代教會當中,我們可以看到很多這樣的例子.伊格那修(67年－110年)他是初期教會深受教會愛戴的牧長,當他被抓時,很多基督徒非常難過,也有些信徒想要去搭救他,因為當時教會被逼迫,被抓的信徒往往會被投入獅子坑中供羅馬貴族娛樂.但伊格那修在從安提阿被押往羅馬的途中,寫信給羅馬的教會卻這樣說:“..我是上帝的麥子,那麼就讓我被野獸的牙齒磨成粉,以便成為基督的餅.就讓那群野獸出來吧,好讓它們成為我的墳墓,吞沒我的身體.當世界再也見不到我的身體的時候,我就真正成了基督的門徒.替我向基督禱告,使我成為祭品..”.這一幫殉道者們,世界所不配有的基督徒,雖然死了,但因著信仍然說話,他們在告訴我們:我們在這世界上不過是客旅,不過是寄居的,這世界並非是我們的家,我們真正的家,一個更美的家鄉,乃是在天上.

活著為主:

保羅說:“我情願離世與主聯合,那是好的無比.然而我活著,更是于你們有益.”基督徒不怕死,但並不是說是要去死.主為我們釘在十字架上而死,乃是叫我們要為他而活.相比初期教會而言,在這個時代,我們更要有為主而活的實踐.據說,北美清教徒時期,牧師去探訪信徒最先問的不是他(她)的身體情況或家庭,而是要理問威斯敏斯德小要理問答,第1問:人生的首要目的是什麼? 　　標準答案是:人生的首要目的就是榮耀神,永遠以他為樂.林前十31:所以你們或吃或喝,無論作什麼,都要為榮耀神而行.以西結先知:“結33:31    他們來到你這裡如同民來聚會,坐在你面前仿佛是我的民.他們聽你的話卻不去行；因為他們的口多顯愛情,心卻追隨財利.”弟兄姊妹,盼望你不是那“仿佛是神的子民”實際卻不是的人,你乃是愛慕衪的話語,為主而活,行事為人都是為榮耀主,永遠以他為樂的民.

結語:弟兄姊妹,保羅說:“既然蒙召,行事為人就當與蒙召的恩相稱.”過一個蒙召相稱的生活,最重要的是你身份的識別,歸於撒旦?亦或上帝?你是屬於上帝的兒女,而且是蒙了恩召的,這恩召是我們的主耶穌基督用寶血將你買贖回來的,那麼就拿出你的行為來,作光作鹽,傳揚基督的福音,來見證你一個基督徒應有的行事為人.



\newpage

\section{凡有血氣盡都如草}
\label{sec:364}
\textbf{伍渭文牧師 - 凡有血氣,盡都如草}
\newline
\newline
\hyperref[sec:363]{< < < PREV SERMON < < <}
~
\hyperref[sec:toc]{[返目錄]}
~
\hyperref[sec:365]{> > > NEXT SERMON > > >}


使徒行傳 2:14,36-41
\newline
\begin{longtable}{cl}
\hline
\hline
章節 & 經文 (和合本修訂版)\\
\hline
2:14 & \begin{tabularx}{0.7\textwidth}{X} 彼得和十一個使徒站起來，他就高聲向眾人說：「猶太人和所有住在耶路撒冷的人哪，這件事你們要知道，要側耳聽我的話。 \end{tabularx} \\ \\ \relax
2:15 & \begin{tabularx}{0.7\textwidth}{X} 這些人並不像你們所想的喝醉了，因為現在才早晨九點鐘。 \end{tabularx} \\ \\ \relax
2:16 & \begin{tabularx}{0.7\textwidth}{X} 這正是藉著先知約珥所說的： \end{tabularx} \\ \\ \relax
2:17 & \begin{tabularx}{0.7\textwidth}{X} 『神說：在末後的日子，我要將我的靈澆灌凡血肉之軀的。你們的兒女要說預言；你們的少年要見異象；你們的老人要做異夢。 \end{tabularx} \\ \\ \relax
2:18 & \begin{tabularx}{0.7\textwidth}{X} 在那些日子，我要把我的靈澆灌，甚至給我的僕人和婢女，他們要說預言。 \end{tabularx} \\ \\ \relax
2:19 & \begin{tabularx}{0.7\textwidth}{X} 在天上，我要顯出奇事，在地下，我要顯出神蹟，有血，有火，有煙霧。 \end{tabularx} \\ \\ \relax
2:20 & \begin{tabularx}{0.7\textwidth}{X} 太陽要變為黑暗，月亮要變為血，這都在主大而光榮的日子未到以前。 \end{tabularx} \\ \\ \relax
2:21 & \begin{tabularx}{0.7\textwidth}{X} 那時，凡求告主名的都必得救。』 \end{tabularx} \\ \\ \relax
2:22 & \begin{tabularx}{0.7\textwidth}{X} 「以色列人哪，你們要聽我這些話：拿撒勒人耶穌就是神以異能、奇事、神蹟向你們證明出來的人，這些事是神藉著他在你們中間施行，正如你們自己知道的。 \end{tabularx} \\ \\ \relax
2:23 & \begin{tabularx}{0.7\textwidth}{X} 他既按著神確定的旨意和預知被交與人，你們就藉著不法之人的手把他釘在十字架上，殺了。 \end{tabularx} \\ \\ \relax
2:24 & \begin{tabularx}{0.7\textwidth}{X} 神卻將死的痛苦解除，使他復活了，因為他原不能被死拘禁。 \end{tabularx} \\ \\ \relax
2:25 & \begin{tabularx}{0.7\textwidth}{X} 大衛指著他說：『我看見主常在我眼前，他在我右邊，使我不至於動搖。 \end{tabularx} \\ \\ \relax
2:26 & \begin{tabularx}{0.7\textwidth}{X} 所以我心裡歡喜，我的舌頭快樂，而且我的肉身要安居在指望中。 \end{tabularx} \\ \\ \relax
2:27 & \begin{tabularx}{0.7\textwidth}{X} 因你必不將我的靈魂撇在陰間，也不讓你的聖者見朽壞。 \end{tabularx} \\ \\ \relax
2:28 & \begin{tabularx}{0.7\textwidth}{X} 你已將生命的道路指示我，必使我在你面前充滿快樂。』 \end{tabularx} \\ \\ \relax
2:29 & \begin{tabularx}{0.7\textwidth}{X} 「諸位弟兄，先祖大衛的事，我可以坦然地對你們說：他死了，也埋葬了，而且他的墳墓直到今日還在我們這裡。 \end{tabularx} \\ \\ \relax
2:30 & \begin{tabularx}{0.7\textwidth}{X} 既然大衛是先知，他知道神曾向他起誓，要從他的後裔中立一位坐在他的寶座上。 \end{tabularx} \\ \\ \relax
2:31 & \begin{tabularx}{0.7\textwidth}{X} 他預先看見了，就講論基督的復活，說：『他不被撇在陰間；他的肉身也不見朽壞。』 \end{tabularx} \\ \\ \relax
2:32 & \begin{tabularx}{0.7\textwidth}{X} 這耶穌，神已經使他復活了，我們都是這事的見證人。 \end{tabularx} \\ \\ \relax
2:33 & \begin{tabularx}{0.7\textwidth}{X} 他既被高舉在神的右邊，又從父受了所應許的聖靈，就把你們所看見所聽見的，澆灌下來。 \end{tabularx} \\ \\ \relax
2:34 & \begin{tabularx}{0.7\textwidth}{X} 大衛並沒有升到天上，但他自己說：『主對我主說：你坐在我的右邊， \end{tabularx} \\ \\ \relax
2:35 & \begin{tabularx}{0.7\textwidth}{X} 等我使你的仇敵作你的腳凳。』 \end{tabularx} \\ \\ \relax
2:36 & \begin{tabularx}{0.7\textwidth}{X} 故此，以色列全家當確實知道，你們釘在十字架上的這位耶穌，神已經立他為主，為基督了。」 \end{tabularx} \\ \\ \relax
2:37 & \begin{tabularx}{0.7\textwidth}{X} 眾人聽見這話，覺得扎心，就對彼得和其餘的使徒說：「諸位弟兄，我們該怎樣做呢？」 \end{tabularx} \\ \\ \relax
2:38 & \begin{tabularx}{0.7\textwidth}{X} 彼得對他們說：「你們各人要悔改，奉耶穌基督的名受洗，使你們的罪得赦免，就會領受所賜的聖靈。 \end{tabularx} \\ \\ \relax
2:39 & \begin{tabularx}{0.7\textwidth}{X} 因為這應許是給你們和你們的兒女，並一切在遠方的人，就是給所有主—我們的神所召來的人。」 \end{tabularx} \\ \\ \relax
2:40 & \begin{tabularx}{0.7\textwidth}{X} 彼得還用更多別的話作見證，勸勉他們說：「你們當救自己脫離這彎曲的世代。」 \end{tabularx} \\ \\ \relax
2:41 & \begin{tabularx}{0.7\textwidth}{X} 於是領受他話的人，都受了洗；那一天，門徒約添了三千人。 \end{tabularx} \\ \\
[1ex]
\hline
\hline
\end{longtable}
詩篇 116:1-4,12-19
\newline
\begin{longtable}{cl}
\hline
\hline
章節 & 經文 (和合本修訂版)\\
\hline
116:1 & \begin{tabularx}{0.7\textwidth}{X} 我愛耶和華，因為他聽了我的聲音和我的懇求。 \end{tabularx} \\ \\ \relax
116:2 & \begin{tabularx}{0.7\textwidth}{X} 他既向我側耳，我一生要求告他。 \end{tabularx} \\ \\ \relax
116:3 & \begin{tabularx}{0.7\textwidth}{X} 死亡的繩索勒住我，陰間的痛苦抓住我，我遭遇患難愁苦。 \end{tabularx} \\ \\ \relax
116:4 & \begin{tabularx}{0.7\textwidth}{X} 那時，我求告耶和華的名：「耶和華啊，求你救我！」 \end{tabularx} \\ \\ \relax
116:5 & \begin{tabularx}{0.7\textwidth}{X} 耶和華有恩惠，有公義，我們的神有憐憫。 \end{tabularx} \\ \\ \relax
116:6 & \begin{tabularx}{0.7\textwidth}{X} 耶和華保護愚蒙的人；我落到卑微的地步，他救了我。 \end{tabularx} \\ \\ \relax
116:7 & \begin{tabularx}{0.7\textwidth}{X} 我的心哪！你要復歸安寧，因為耶和華用厚恩待你。 \end{tabularx} \\ \\ \relax
116:8 & \begin{tabularx}{0.7\textwidth}{X} 主啊，你救我的命脫離死亡，使我的眼不再流淚，使我的腳不致跌倒。 \end{tabularx} \\ \\ \relax
116:9 & \begin{tabularx}{0.7\textwidth}{X} 我行在耶和華面前，走在活人之地。 \end{tabularx} \\ \\ \relax
116:10 & \begin{tabularx}{0.7\textwidth}{X} 我信，儘管我說：「我受了極大的困苦。」 \end{tabularx} \\ \\ \relax
116:11 & \begin{tabularx}{0.7\textwidth}{X} 我曾驚惶地說：「人都是說謊的！」 \end{tabularx} \\ \\ \relax
116:12 & \begin{tabularx}{0.7\textwidth}{X} 耶和華向我賞賜一切厚恩，我拿甚麼來報答他呢？ \end{tabularx} \\ \\ \relax
116:13 & \begin{tabularx}{0.7\textwidth}{X} 我要舉起救恩的杯，稱揚耶和華的名。 \end{tabularx} \\ \\ \relax
116:14 & \begin{tabularx}{0.7\textwidth}{X} 我要在他的全體百姓面前向耶和華還我所許的願。 \end{tabularx} \\ \\ \relax
116:15 & \begin{tabularx}{0.7\textwidth}{X} 在耶和華眼中，聖民之死極為寶貴。 \end{tabularx} \\ \\ \relax
116:16 & \begin{tabularx}{0.7\textwidth}{X} 耶和華啊，哦，我是你的僕人；我是你的僕人，是你使女的兒子。你已經解開我的捆索。 \end{tabularx} \\ \\ \relax
116:17 & \begin{tabularx}{0.7\textwidth}{X} 我要以感謝為祭獻給你，又要求告耶和華的名。 \end{tabularx} \\ \\ \relax
116:18 & \begin{tabularx}{0.7\textwidth}{X} 我要在耶路撒冷當中，在耶和華殿的院內，在他的全體百姓面前，向耶和華還我所許的願。哈利路亞！ \end{tabularx} \\ \\ \relax
116:19 & \begin{tabularx}{0.7\textwidth}{X} 見上節 \end{tabularx} \\ \\
[1ex]
\hline
\hline
\end{longtable}
彼得前書 1:17-23
\newline
\begin{longtable}{cl}
\hline
\hline
章節 & 經文 (和合本修訂版)\\
\hline
1:17 & \begin{tabularx}{0.7\textwidth}{X} 既然你們稱那不偏待人、按各人行為審判人的主為父，就當存敬畏的心，度你們在世寄居的日子。 \end{tabularx} \\ \\ \relax
1:18 & \begin{tabularx}{0.7\textwidth}{X} 你們知道，你們得以從你們祖先傳下來虛妄的行為中救贖出來，不是靠著會朽壞的金銀等物， \end{tabularx} \\ \\ \relax
1:19 & \begin{tabularx}{0.7\textwidth}{X} 而是憑著基督的寶血，如同無瑕疵、無玷污的羔羊的血。 \end{tabularx} \\ \\ \relax
1:20 & \begin{tabularx}{0.7\textwidth}{X} 基督是神在創世以前所預知，而在這末世才為你們顯現的。 \end{tabularx} \\ \\ \relax
1:21 & \begin{tabularx}{0.7\textwidth}{X} 你們也因著他而信那使他從死人中復活、又給他榮耀的神，好讓你們的信心和盼望都在於神。 \end{tabularx} \\ \\ \relax
1:22 & \begin{tabularx}{0.7\textwidth}{X} 既然你們因順從真理而潔淨了自己的心靈，能真誠愛弟兄，就該以清潔的心彼此切實相愛。 \end{tabularx} \\ \\ \relax
1:23 & \begin{tabularx}{0.7\textwidth}{X} 你們蒙了重生，不是由於會朽壞的種子，而是由於不會朽壞的種子，是藉著神永活常存的道。 \end{tabularx} \\ \\
[1ex]
\hline
\hline
\end{longtable}
路加福音 24:13-35
\newline
\begin{longtable}{cl}
\hline
\hline
章節 & 經文 (和合本修訂版)\\
\hline
24:13 & \begin{tabularx}{0.7\textwidth}{X} 同一天，門徒中有兩個人往一個村子去；這村子名叫以馬忤斯，離耶路撒冷約有二十五里。 \end{tabularx} \\ \\ \relax
24:14 & \begin{tabularx}{0.7\textwidth}{X} 他們彼此談論所發生的這一切事。 \end{tabularx} \\ \\ \relax
24:15 & \begin{tabularx}{0.7\textwidth}{X} 正交談議論的時候，耶穌親自走近他們，和他們同行， \end{tabularx} \\ \\ \relax
24:16 & \begin{tabularx}{0.7\textwidth}{X} 可是他們的眼睛模糊了，沒認出他。 \end{tabularx} \\ \\ \relax
24:17 & \begin{tabularx}{0.7\textwidth}{X} 耶穌對他們說：「你們一邊走一邊談，彼此談論的是甚麼事呢？」他們就站住，臉上帶著愁容。 \end{tabularx} \\ \\ \relax
24:18 & \begin{tabularx}{0.7\textwidth}{X} 兩人中有一個名叫革流巴的回答：「你是在耶路撒冷的旅客中，惟一還不知道這幾天在那裡發生了甚麼事的人嗎？」 \end{tabularx} \\ \\ \relax
24:19 & \begin{tabularx}{0.7\textwidth}{X} 耶穌對他們說：「甚麼事呢？」他們對他說：「就是拿撒勒人耶穌的事。他是個先知，在神和眾百姓面前，說話行事都大有能力。 \end{tabularx} \\ \\ \relax
24:20 & \begin{tabularx}{0.7\textwidth}{X} 祭司長們和我們的官長竟把他解去，定了死罪，釘在十字架上。 \end{tabularx} \\ \\ \relax
24:21 & \begin{tabularx}{0.7\textwidth}{X} 但我們素來所盼望要救贖以色列民的就是他。不但如此，這些事發生到現在已經三天了。 \end{tabularx} \\ \\ \relax
24:22 & \begin{tabularx}{0.7\textwidth}{X} 還有，我們中間的幾個婦女使我們驚奇：她們清早去了墳墓， \end{tabularx} \\ \\ \relax
24:23 & \begin{tabularx}{0.7\textwidth}{X} 不見他的身體，就回來告訴我們，說她們看見了天使顯現，說他活了。 \end{tabularx} \\ \\ \relax
24:24 & \begin{tabularx}{0.7\textwidth}{X} 又有我們的幾個人往墳墓那裡去，所發現的正如婦女們所說的，只是沒有看見他。」 \end{tabularx} \\ \\ \relax
24:25 & \begin{tabularx}{0.7\textwidth}{X} 耶穌對他們說：「無知的人哪，先知所說的一切話，你們的心信得太遲鈍了。 \end{tabularx} \\ \\ \relax
24:26 & \begin{tabularx}{0.7\textwidth}{X} 基督不是必須受這些苦難，然後進入他的榮耀嗎？」 \end{tabularx} \\ \\ \relax
24:27 & \begin{tabularx}{0.7\textwidth}{X} 於是，他從摩西和眾先知起，凡經上所指著自己的話都給他們作了解釋。 \end{tabularx} \\ \\ \relax
24:28 & \begin{tabularx}{0.7\textwidth}{X} 他們走近所要去的村子，耶穌好像還要往前走， \end{tabularx} \\ \\ \relax
24:29 & \begin{tabularx}{0.7\textwidth}{X} 他們卻強留他說：「時候晚了，天快黑了，請你同我們住下吧。」耶穌就進去，要同他們住下。 \end{tabularx} \\ \\ \relax
24:30 & \begin{tabularx}{0.7\textwidth}{X} 坐下來和他們用餐的時候，耶穌拿起餅來，祝福了，擘開，遞給他們。 \end{tabularx} \\ \\ \relax
24:31 & \begin{tabularx}{0.7\textwidth}{X} 他們的眼睛開了，這才認出他來。耶穌卻從他們眼前消失了。 \end{tabularx} \\ \\ \relax
24:32 & \begin{tabularx}{0.7\textwidth}{X} 他們彼此說：「在路上他和我們說話，給我們講解聖經的時候，我們的心在我們裡面豈不是火熱的嗎？」 \end{tabularx} \\ \\ \relax
24:33 & \begin{tabularx}{0.7\textwidth}{X} 於是他們立刻起身，回耶路撒冷去，看見十一個使徒和與他們正在一起的人聚集在一處， \end{tabularx} \\ \\ \relax
24:34 & \begin{tabularx}{0.7\textwidth}{X} 說：「主果然復活了，已經顯現給西門看了。」 \end{tabularx} \\ \\ \relax
24:35 & \begin{tabularx}{0.7\textwidth}{X} 於是，兩個人把路上所遇到，和耶穌擘餅的時候怎麼被他們認出來的事，都述說了一遍。 \end{tabularx} \\ \\
[1ex]
\hline
\hline
\end{longtable}
講題:凡有血氣,盡都如草　All Flesh Is Like Grass

經課:使徒行傳2:14,36-41;詩篇116:1-4,12-19;彼得前書1:17-23;路加福音24:13-35

講員:伍渭文牧師

場合:香港中文大學崇基學院禮拜堂

日期:2017年4月30日(復活期第三主日)

今天的題目來自書信經課彼前一24「凡有血氣的,盡都如草」.這句話表面看來消極.讓我們縱情歡樂吧:世事如黃粱一夢,轉眼成空,何必太執著認真,不如縱情享樂,因為我們快要死了.也不是懼怕老去:有些人抗拒短暫,拒絕老化自然過程,把最絢爛的一刻強留凝固.十五世紀義大利弗羅倫斯座堂經常講道的道明會修士沙方那路勒(Girolamo Savonarola提到一故事.他每次來教會講道前都注意到一位老婦人在聖母像前默默的瞻仰聖像,有一天他與教會一位年長的神甫談及這事,推崇老婦的虔誠.年長神甫說:「她不是虔誠,乃戀慕自已的年青肖像.我們為聖母作像時,雕塑家就請了這婦人作模特兒,當年她容貌真的清純漂亮.」

血氣的短暫和朽壞:

艾略特(T. S. Eliot)就看透血氣之軀的短暫, 他歸信基督三年後1930年發表的〈聖灰日〉(Ash Wednesday) 其中一句說:「真相只能在某一時間,某一地方,呈現一次,我欣悅接納事物的本相.美麗的容貌,動人的聲音,過去就過去了,我不能盼望逆轉時間」(Truth only appears once “What is actual only for one time and only for one place. I rejoice that things are as they are. I renounce the blessed face and renounce the voice because I cannot hope to turn again) .他看透生死輪轉,讓真相自然呈現,不強留,不哀嘆,萬物皆有其時.真相是不能重複的,只能出現一次,一次之後就過去了.傳道書的作者,用詩章道出血氣的朽壞:「銀鍊折斷,金罐破裂,瓶子在泉旁損壞,水輪在井口破爛,塵土仍歸於地,靈仍歸於賜靈的上帝.」(傳十一6, 7)

人的苦惱,是看到血氣生命的短暫的同時,心中卻對永續不朽有響往.我們希望有意義和價值的東西,包括自已的生命,可以永續不朽.所以亞伯蘭得蒙應許成為大國之父,但妻子過了生育期還沒有兒子,就非常焦慮.有一晚上,踏出帳棚望向星空,天上的星星在漆黑中看似就在頭頂,其實相距很遠很遠,很長很長的時間停在那裡,絲毫不動,沒有改變,但自己的人生已過半了.突然間,他感到人生的短暫,無後的焦慮,他渴想不朽永續.

傳記作者Walter Isaacson與喬布斯(Steve Jobs)作了超過四十次訪談,最後一次在他家中的後花園進行,是陽光普照的下午,當時喬布斯身體很不舒服,喬布斯把話題轉到宗教和死亡.當問及他是否相信死後生命(afterlife),是否相信有上帝時,他說:「五十五十.這是大奧秘,人永遠不知道,但我希望死不是生命的終結,我但願累積的智慧不會隨死亡歸於無有,可以永續.但他跟著說:但可能它真的像開關制,click一下,你就離開了；停頓一秒後他笑說:所以我們蘋果產品沒有開關制的.」

惟有主道永存:

彼得前書的作者是積極的,重要在下半句:「凡有血氣的,盡都如草,他的榮美都像草上的花.草必枯乾,花必凋謝；惟有主的道是永存的.所傳給你們的福音就是這道.」(彼前一24, 25) 今日是復活期第三主日,經課就像衛星定位,讓我們思考確立經文的共同概念.若我們細看,四段經課共有的一個主題-基督的死和復活帶給我們不朽的生命.

書信經課,彼前經課一17-23,提到不能朽壞的生命,像永存的種子在我們心中成長,是延續前面第一章3-12節的祝福中所說的不朽生命.這祝福把整卷書的內容握要點出來了,新約書信在開始時都有這祝福格式,因為所有書信都在崇拜的處境中誦讀.「願頌讚歸與我們主耶穌基督的父上帝!他曾照自已的大憐憫,藉耶穌基督從死裡復活,重生了我們,叫我們有活潑的盼望,可以得著不能朽壞,不能玷污,不能衰殘,為你們存留在天上的基業.」(彼前一3, 4)

首經課,使二37「眾人-聽見這話,覺得扎心」.甚麼話?「你們就藉無法之人的手,把他釘在十字架上,殺了.上帝卻將死的痛苦解釋了,叫他復活,因為他原不能被死拘禁.」(使二23, 24)；「因你必不將我的靈魂撇在陰間,也不叫你的聖者見朽壞.」(使二27,引詩十六10)

次經課,詩一一六「死亡的繩索纏繞我；陰間的痛苦抓住我(3節)；……在耶和華眼中,看聖民之死極其寶貴.(15節)……你已經解開我的綁索.(16節)」

福音經課,路二十四13-35, 復活的基督在以瑪忤斯路上,藉講解聖道使兩位門徒的信心活過來.

1 不朽的生命聖化血氣的生命

慢慢成長的種子

基督復活期共五十天,跟著是五旬節,五旬節結束了教會年的上半年,以基督生平救贖事件為主軸的上半年,同時開始了教會年的下半年.上半年關於基督,下半年關於基督徒.下半年又稱常年期,聖帷是代表成長的綠色,直到教會年結束的基督君王主日,基督復臨審判世界,終未來臨.

我們稱下半年為聖靈降臨後,依次第一,二,三……主日.聖靈降臨聖帷是紅色,那天有大響聲,有火焰,有人說起不懂的語言,是一個聽得見,看得見的異常景象.但之後的主日經課,是細涓長流,像種子埋在地下,常常澆灌,拔除野草,修剪枝葉,悉心培育的成長期.成聖的路是漫長的,潛移默化.

復活的生命,在我們短暫朽壞的血氣生命中,播下了永存的種子,這種子不斷成長,伴隨這盡都如草的生命,使其尊貴,得著聖化.血氣的生物界,千姿百態.上帝永存的道,沒有使血氣的牛變成馬,乃變成更好的牛,更好的馬,永存的道轉化了我們血氣的生命.

聖化血氣的生命:

我很喜歡「賜我眾山的力量」的禱文,這禱文把不同時間,不同季節,不同形態的山脈,都賦予屬靈的意義:

看到雪山,想到:讓我光明磊落.願我如雪峰皚皚,臨日之昭,親星之耀,晝夜無間.

看到火山,想到:讓我嫉惡如仇.給我兇猛堅固如火山,內心熔岩滾動,含焰待放.

峻峭的山, 想到:讓我止於至善.讓我攀天淩急風,崢嶸拔地起,抖擻凡塵瑣事,遠離窪谷低地,心無牽掛.

雪山是冰冷的,轉化為光明磊落；火山暴烈,轉化成嫉惡如仇；峻峭的山高不可攀,變作追求止於至善.復活的生命在我們心中起了默默的革命.

賜我眾山的力量:[i]

偉大的上主,賜我眾山的力量.

讓我內心純潔.

賜我力量,清明如練,

像經過夏天雨水洗滌後的翠綠山頭.

讓我寧靜安穩.

靜脈照晚霞,日累西山下,

在黑夜的臂彎中,給我心靈安恬.

讓我止於至善.

讓我攀天淩急風,崢嶸拔地起,

抖擻凡塵瑣事,遠離窪谷低地,心無牽掛.

讓我光明磊落.

願我如雪峰皚皚,

臨日之昭,親星之耀,晝夜無間.

讓我嫉惡如仇.

給我兇猛堅固如火山,

內心熔岩滾動,含焰待放.

偉大的上主,賜我眾山的力量.阿們.

脫去與穿上: 個人倫理與社會倫理

血氣生命轉化的過程中,人也有應盡的責任-作順命的兒女.彼得說:「所以要約束你們的心,謹慎自守……你們既作順命的兒女,就不要效法從前蒙昧無知的時候那放縱私欲的樣子.」(彼前一13, 14) 如何約束血氣本能的放縱呢?要學習真理:「你們既因順從真理,潔淨了自已的心」(彼前一22)；「就要愛慕那純淨的靈奶,像才生的嬰孩愛慕奶一樣,叫你們因此漸長,以致得救.」(彼前二2) 不單學習真理,更要實踐愛心:「當從心裡彼此切實相愛」(彼前一22)

在大學談基督信仰,一般談公義,仁愛,共融的社會倫理,行公義,好憐憫；較少談個人倫理,甚至避談個人倫理,覺得是私隱.但彼得在這裡題到約束你們的心,不要放縱私欲.「你們在外邦人人中,應當行為端正,叫那些毀謗你們是作惡的,因著你們的好行為……歸榮耀給上帝.」(二15) 社會倫理應建基於個人倫理,先改變自已,再推而廣之,改變世界,所以我們稱自已是見證人,見證福音轉化生命的能力.

彼得鼓勵正受逼迫的信徒,以好行為回應逼迫他們的惡人；從自已做起,改變社會.「凡事要存敬畏的心順服主人；不但順服那善良溫和的,就是那些乖僻的也要順服.」(彼前二18) 政治改革沒有吋進,我們都感到氣餒,但仍可以透過個人的好行為,使社會得到革新.與其詛咒黑暗,不如自己發光.

2  不朽的生命,承托血氣生命的脆弱

次經課,詩一一六「死亡的繩索纏繞我；陰間的痛苦抓住我(3節)；……在耶和華眼中,看聖民之死極其寶貴.(15節)……你已經解開我的綁索.(16節)」

復活沒有取消墳墓,但帶走墳墓裡的幽暗和絕望.我們與上主的關系,連死亡都不能間斷,這就是永生,這就是我們的福音.

在希臘的文化中,用朽壞來形容生物界,用不朽形容靈魂的世界.人死後肉體埋在黃土下漸漸朽壞,但靈魂是永恆不滅的.人死後,靈魂脫離肉體的禁固,歸回所屬的永恆世界,就是目測至高的星河,每一顆河漢的星都是離世的人的靈魂,那裡永恆不變,寧靜安穩,不像人世間滾滾紅塵,變幻無常,.

但希伯來聖經的永生,是以和上帝的關系來定義.上帝是亞伯拉罕的上帝,雅各的上帝,以撒的上帝.列祖死了,但他們和上帝的關系沒有死.這不是過去歷史的描述,乃現況的宣告.因上帝是永恆不朽,連於上帝,我們也永恆不朽,這就是永生.永生不是指靈魂不滅,乃指和上主不分離,這永生就在今世接受福音時開始了.「你們蒙了重生,不是由於能壞的種子,乃是由於不能壞的種子,是藉著上帝活潑常存的道.因為凡有血氣的,盡都如草；他的榮美都像草上的花,草必枯乾,花必凋謝.惟有主的道是永存的.所傳給你們的福音就是這道.」(彼前一23- 25)

這位愛我們的上主,已使基督從死裡復活.他不叫基督的靈魂撇在陰間,也不叫相信基督的人的靈魂撇在陰間.(使二27) 不錯,我們會離開陽世,下到陰間,但靈魂不會被撇在陰間.死亡的繩索纏繞我；陰間的痛苦抓住我,但上主已經解開我的綁索.死亡的繩索確能纏繞我們,沒有人免去死亡的繩索,但因基督的復活這索帶已被解開了.這不朽的生命,承托著血氣生命的脆弱.上主也珍惜我們面對死亡時的哀傷,正如詩人所說:「求你把我的眼淚裝在你的皮袋裡.」(詩五十六8)

鄧瑞強博士師母前一段時突然去世,我們同感悲痛.鄧博士寫了一篇感人至深的短文《燕玲.基督.救贖》:「教會的朋友來了,便想起和燕玲的往事,淚就爆出來.家人來了,便想起燕玲和他們的情誼,淚就爆出來.朋友來了,便想起燕玲和他們的交往,淚就爆出來.……一想起她,便又哭了.我發覺,在心底深處,存有一種莫名的力量,很平靜的,像大海一樣廣闊的,像大海一樣深邃的.這種無邊的平靜,像巨大的母親,讓我躺下去,想哭多少便哭多少,想哭多久便哭多久.它無言無語,它祇是在那裡.它看來沒有安慰我,也沒有囑咐我不要傷感,祇是讓我哭,甚至鼓勵我哭,也明白我哭.」

上主看聖民之死,極其寶貴；我流淚,他用皮袋裝.不朽的種子,承托著脆弱的生命.

不朽的生命,有治愈傷痛記憶的能力

前明報總篇輯劉進圖弟兄,2014年二月二十六日被刀狠下毒手,在醫院裡長期臥床,之後坐輪椅,拿柺仗,戴腳套,像小孩子一樣重新學走路,之後做了三年物理治療,至今仍未能跑未能跳.這是用鮮血寫成的記憶,如何忘記呢?怎能饒恕傷害自已的兇手呢?2017年三月七日他在I　Ｑｕｅｓｔ寫了一篇名為《刺痛的記憶》文章,講述得到治愈的經過.

作者引用克羅地亞裔神學家沃弗(Miroslav Volf)所寫的《記憶的力量》一書的建議去治療傷痛的記憶.沃弗由於娶了美國人做太太,又是基督教神學家,在前南斯拉夫的共產政權眼中是典型的間諜,顛覆分子,持續多個月對他實施監控,竊聽,刺探,陷害,然後是無窮無盡的盤問,威脅他坦白從寬抗拒從嚴,軍事法庭隨時判他八年監禁,他感到恐懼,無力,絕望.

沃弗在書中提到醫治受傷的記憶的三股力量:其一是捉緊基督徒乃蒙神憐愛的兒女這個新身分,以此作為安身立命的依據,拒絕讓創傷進佔生命舞台的中心位置；其二是捉緊基督復活令信徒生命不斷湧現新的可能,讓我們的未來不再被過去的經歷主宰,得以重燃盼望；其三是捉緊基督的十字架,將創痛與仇敵的回憶,置於基督十字架的神聖記憶之下,一起經歷饒恕與轉化.劉進圖弟兄以本身的經驗,印證這三點確有效治愈傷痛的記憶.

其實記憶的力量的身份,盼望,與饒恕正是彼得前書所強調的.蒙召的兒女的身份:「作順命的兒女」(彼前一14),「那召你們的」(彼前一5).彼前二9:「惟有你們是被揀選的族類,是有君尊的祭司,是聖潔的國度,是屬上帝的子民,要叫你們宣揚那召你們出黑暗入奇妙光明者的美德.」復活帶來的盼望:「他曾照自已的大憐憫,藉耶穌基督從死裡復活,叫我們有活潑的盼望.」(彼前一3)「你們要將一切的憂慮卸給上帝,因為他顧念你們.」(彼前五7)十字架的救贖和饒恕:「知道你們得贖……是憑基督的寶血.」(彼前一18)而基督為我們受苦,立下榜樣,好叫我們跟隨他的腳蹤.

3 復活帶來不朽的生命,使我們窺探到世界的終末

不朽確是和時間相關的概念-直到永遠,永無窮盡.但聖經的不朽還有目的,終未的概念.不朽是不改變,但祇有上帝是不改變,他的旨意不改變.他預定的藍圖,是最原始的,也是最終未的.「基督在創世以前是預先被上帝知道的,卻在這末世才為你們顯現.」(彼前一20)「你們因信蒙上帝能力保守的人,必能得著所預備,到末世要顯現的救恩.」(彼前一5)「預先」,是上帝永恆的旨意,「這未世」指耶穌第一次來臨,「到末世」指耶穌第二次來臨,也是終未,終未是所有事情的目的揭曉時刻.我們擁有這不朽生命的人,在世上是客旅,行經時間,完成上主交付給我們的任務,這任務是上帝在我們生命中的計劃.

十字架與復活把預先,今世,未世(這未世及到未世)三點凝固在一起.世界的歷史,人類的命運,創造的目的,救贖的目的都在此中求.在啟示錄第五章,天使大聲宣傳說:「有誰配得展開那書卷,揭開那七印呢?」(啟五2)祇有被殺的羔羊,配得及能夠打開書卷.祇能透過十字架,才能明白一切,或說,才能忍受一切.因為上主的旨意隱密難明,但上主的愛,在基督裡顯而易見.因此,因為十字架,我們可以忍受一切.

加爾文認為上主永恆旨意至高無尚,深受其思想影響的英國聖詩之父以撒窩特(Isaac Watt)所寫的When I Survey the Wondrous Cross仰望十架,就把十字架懸掛在宇宙中間,以十字架為坐標,量度整個宇宙,除了十字架,我們沒法明白宇宙存在的目的.英文歌名用survey,量度,中譯「仰望」有虔敬之意味,但失了加爾文強調上帝永恆旨意的原意了.我們不是無緣無故生在世上,是上帝預先的計劃讓我們生於此時此地,這點令我們驚懼.

我們在世是客旅,行經時間,也是上帝的計劃和使命.不錯,我們因血氣之身而會衰敗,最後死亡,正如聖灰日塗灰禮所說:「你要緊記,你來自塵土,仍要歸回塵土,要遠離罪惡,跟隨基督.」但我們懷有不朽的種子,與上帝連結,這關系連死亡也不能奪去,這關系是以基督的寶血所保證.

在這客旅途中,我們也許會遇到種種橫逆,血氣的生命是脆弱的,但上帝不朽的生命承托著我們.也許我們也遭遇很多傷痛的事情,帶來刺痛的回憶,揮之不去,很難饒恕,但復活帶給我們盼望,十字架叫我們學曉饒恕,讓我們把刺痛的記憶,怨憤與基督同釘十字架,也與基督一同活過來.

各位,可能我們生命的書卷不是因為無知被封起來,而是因為我們太了解,不敢面對而封起來.我們總有一些生命的片段覺得幽暗,羞愧,罪疚,不敢面對,把它封起來.今天基督已為你的罪死在十字架,又復活了,基督願意赦免,給我新生.撒旦的控訴,羅疚的拘禁,像死亡的成繩索,纏繞我們,但上主已把綁索解開了.

「死亡的繩索纏繞我；陰間的痛苦抓住我 ,但你已經解開我的綁索.」阿們.

[i] Give me the strength

Of verdant hills

Washed clean by summer rain;

Of purple hills

At peace when weary Day

Sinks quietly to rest

In Night’s cool arms;

Of rugged, wind-whipped hills

That lift their heads

Above the petty, lowland, valley things,

And shake their shoulders free

Of bonds that hold

Them close to earth;

Of snow-capped hills

Sun-kissed by day, by night

Companioned by the stars;

Of grim volcanoes

Pregnant with the fires

Of molten fury!

Grant me strength,

Great God,

Like that of hills!

by Esther Popel (1896-1658), 非裔美國詩人,教育家；伍渭文牧師/譯



\newpage

\section{追追趕趕高高低低}
\label{sec:365}
\textbf{林豪恩先生 - 追追趕趕高高低低}
\newline
\newline
\hyperref[sec:364]{< < < PREV SERMON < < <}
~
\hyperref[sec:toc]{[返目錄]}
~
\hyperref[sec:366]{> > > NEXT SERMON > > >}


詩篇 23:1-6
\newline
\begin{longtable}{cl}
\hline
\hline
章節 & 經文 (和合本修訂版)\\
\hline
23:1 & \begin{tabularx}{0.7\textwidth}{X} 耶和華是我的牧者，我必不致缺乏。 \end{tabularx} \\ \\ \relax
23:2 & \begin{tabularx}{0.7\textwidth}{X} 他使我躺臥在青草地上，領我在可安歇的水邊。 \end{tabularx} \\ \\ \relax
23:3 & \begin{tabularx}{0.7\textwidth}{X} 他使我的靈魂甦醒，為自己的名引導我走義路。 \end{tabularx} \\ \\ \relax
23:4 & \begin{tabularx}{0.7\textwidth}{X} 我雖然行過死蔭的幽谷，也不怕遭害，因為你與我同在；你的杖、你的竿，都安慰我。 \end{tabularx} \\ \\ \relax
23:5 & \begin{tabularx}{0.7\textwidth}{X} 在我敵人面前，你為我擺設筵席；你用油膏了我的頭，使我的福杯滿溢。 \end{tabularx} \\ \\ \relax
23:6 & \begin{tabularx}{0.7\textwidth}{X} 我一生一世必有恩惠慈愛隨著我；我且要住在耶和華的殿中，直到永遠。 \end{tabularx} \\ \\
[1ex]
\hline
\hline
\end{longtable}
講題:追追趕趕高高低低 　Catch My Breath

經文:詩篇23:1-6

講員:林豪恩先生

場合:香港中文大學崇基學院禮拜堂

日期:2017年5月7日

一.引言

追追趕趕高高低低這幾個字,有沒有令你想到陳奕迅主唱的「幸福摩天輪」這首曾經紅極一時的粵語流行曲呢?如果你想起這首歌,這個講道題目可能會給你浪漫的感覺,因為陳奕迅這首歌成功地訴說愛的甜美:

追追趕趕　高高低低 深呼吸然後與你執手相隨 甜蜜中不再畏高 可這樣跟你蕩來蕩去　無畏無懼 追追趕趕　高高低低 驚險的程度叫畏高者昏迷 憑甚麼不怕跌低 多僥倖跟你共同面對　時間流逝 東歪西倒　忽高忽低 心驚與膽戰去建立這親厚關係 沿途就算意外脫軌 多得你　陪我搖曳

追追趕趕高高低低其實不一定是浪漫的,這樣的狀況可以是英文題目所示catch my breath,叫人喘不過氣來.到底是浪漫,還是喘不過氣來呢?是甚麼決定如此截然不同的感覺呢?剛才大家讀詩篇二十三篇的時候,是否也感到溫馨浪漫呢?除了溫馨浪漫,有沒有感到詩歌之中的追追趕趕高高低低那種令人喘不過氣來的催逼呢?讓我們再來讀一次詩篇第二十三篇.

二.經文

23:1耶和華是我的牧者,我必不致缺乏.

23:2他使我躺臥在青草地上,領我在可安歇的水邊.

23:3他使我的靈魂甦醒,為自己的名引導我走義路.

23:4我雖然行過死蔭的幽谷,也不怕遭害,因為你與我同在；你的杖,你的竿,都安慰我.

23:5在我敵人面前,你為我擺設筵席；你用油膏了我的頭,使我的福杯滿溢.

23:6我一生一世必有恩惠慈愛隨著我；我且要住在耶和華的殿中,直到永遠.

三.釋經與應用

1．追與趕-搵食與生存

23:2他使我躺臥在青草地上,領我在可安歇的水邊.2 You have bedded me down in lush meadows, you find me quiet pools to drink from.

對人來說,青草地和溪水旁是休息的好地方.為口奔馳,營營役役之後,若能夠入住於青草地和溪水旁的渡假村休息,是何等鬆弛自在的事情.但對羊來說,剛好相反,青草地和溪水旁不是牠們休息的地方,而是「搵食」的地方,因為草是他們的食物,水是他們的飲品,草和水是牠們生存的必需品.牧羊人帶領他們去青草地和溪水旁,原來是帶他們去「開工」.

對不起,破壞了大家對這首詩的浪漫想像.聽到「開工」,「搵食」,人們的心情就不同了,可能馬上出現「唔使急,最緊要快」催逼感覺,就是那是追追趕趕,永無窮盡的狀態.區瑞強的經典名曲「那天再重聚」裡有幾句這樣說:

在匆匆趕又匆匆追　永遠不說後退

讓身邊生活苦相迫　競爭令人累

朋友　不要淚垂　人生本多失意

競爭　開始了　追不到也是要追

就算是羊,上帝是牧羊人,牠們仍然要為生存而「搵食」.作為人,上帝把我們帶到世界來,我們豈能不為生存而「搵食」嗎?保羅曾說:「若有人不肯做工,就不可吃飯.」(帖後三:10).對羊來說,能存只需要有草有水；但對人來說,生存的需要和考驗就複雜得多.

許冠傑曾唱:「為兩餐乜都肯制前世」.現在,不少人面對的挑戰是能否找到養活自己的工作?對生活在大城市的人來說,「吃」可以豐儉由人,「住」才是困難所在.就算有工作,能否在工作崗位上「生存下去」,也是一項挑戰.工作間的「刀光劍影」,不知令多少人寢食不安,惶惶不可終日.一位教育界前輩如此說:

「公務員事務局在報章刊登招聘廣告,聘請700個助理文書職位,只需中學會考及格.但據人事顧問估計,將會有一大群大學生畢業生湧往申請.大學生爭取中學學歷要求的文書工作,近年成為「常態」.在職場的廝殺中,這類成功者應是喜悅多於失落,而以中學學歷求職應是愈趨困難了.就以筆者較熟悉的教育招聘做例子吧.近年衍生出合約教師這一族群,這其實是一批教育的臨時工,一群隨時可以辭退的教師.合約教師愈積愈多,多到在社交網頁形成一個群組.最近在當中看到一些訴苦文字:「D 位去左邊?就係去左早出世幾年,能力一般,自我感覺良好既人度.呢的老XX好多都好廢,淨係識貪住份人工,書又唔識教,野又唔識,職務又推比我地做……」.(鄭楚雄:香港要有到位的青年「搵食」政策17/8/2016)

求生,一點不浪漫.外間有人以為在大學得到教席就可以安寢無憂,年青的教授們可以告訴你,在所謂(up or out)的挑戰之下,壓力不足為外人道.助理教授要在6年內,若不能晉升為副教授,就「無得留低」.正如區瑞強的歌所唱:「競爭,開始了,追不到也是要追」.

2．趕與追-「搵食je ,犯法呀」

23:3他使我的靈魂甦醒,為自己的名引導我走義路.

有人可能想像退修的場景,靈魂得到滋養和復興,聽到上帝的啟示和引導,踏上人生更有意義的道路.Eugene Peterson把這句譯作:3 True to your word, you let me catch my breath and send me in the right direction.原來,這不是退修的場景,而是誤入迷途,走錯路,心慌意亂,喘不過氣來的危急時刻.

走錯路,與上一句提及的「搵食」很有關係.黃子華棟篤笑有一名句,廣泛流傳:「搵食je ,犯法呀」.「搵食」這壓倒性的理由,讓人理直氣壯地做出違反道德的事,遊走在法律邊緣.他們甚麼會說:「忠忠直直,終需乞食」.不過,其實就算犯法,也不一定「搵到食」.本屆香港金像獎頒獎典禮大贏家《樹大招風》的其中一位導演歐文傑接受訪問是談及他所知道的葉繼歡:

「當時,不少報道將葉繼歡冠以大賊之名,說他設計了一條完美路線,十數分鐘就劫了五間金舖,帶走巨額贓金,但現實上大賊不易做,打完劫,銷贓後贓金多數只剩兩成,加上水腳,彈藥費,再分給兄弟,其實所餘無幾.我們以為大單劫案打完下半世就可歎世界,但其實一單搵很少錢,葉繼歡回到大陸甚至要去做睇場.」(明報12/2/2016《樹大招風》3導演拍3大賊 回望大茶飯年代)

佛家說:苦海無邊,回頭是岸.在人生的旅途上,處處充滿了誘惑,陷阱,如果不懂得「回頭」,終將陷入萬劫不復之境.「搵食je ,犯法呀?」金錢名利,高官厚祿,名譽地位,可知令多少英雄好漢漸漸走上不歸路?

3．高與低

23:4我雖然行過死蔭的幽谷,也不怕遭害,因為你與我同在；你的杖,你的竿,都安慰我.4 Even when the way goes through Death Valley, I’m not afraid when you walk at my side. Your trusty shepherd’s crook makes me feel secure.

雖說「回頭是岸」,然而,「打個電話就搵到幾十萬,甚至幾百萬人民幣」,回頭並不是容易的事呢!沒有人會選擇「偏行己路」,因為人人都以為自己所行的是正路；沒有人會選擇跳入死蔭幽谷,因為人人都認為自己是為勢所逼.就正如沉淪毒海的人士,每個人都有「故事」:「十歲小孩說自己只因不知天高地厚；青春少艾說自己得不到家庭溫暖,追尋早熟歸宿；大好前程專業人 士說抵受不住工作壓力；少婦說只是為了在婆媳關係夾縫中尋找出路；衣食無憂中產中年說不堪事業挫敗.一切生命轉折錯落,是藉口還是理由,正是當局者迷.(香港電台:「禁毒檔案」

詩人的人生路上有死蔭的幽谷,我們的人生路豈會沒有高高低低呢?我們的「死蔭的幽谷」很可能不是毒品,但會否是我們的學業,我們的工作,或是我們的人際關係呢?

4．低與高

23:5在我敵人面前,你為我擺設筵席；你用油膏了我的頭,使我的福杯滿溢.5 You serve me a six-course dinner right in front of my enemies. You revive my drooping head; my cup brims with blessing.

有些人宣稱「我沒有敵人」,不就等於沒有人找你麻煩.在競爭激烈的社會,在「不要輸在起跑線」的威脅之下,有些人自小就被灌輸「視所有人為敵人」的意識.有些人認為,為了自己的生存,逼不得以犧牲別人；有些群體認為,為了自己群體生存,逼不得以犧牲其他群體；有些國家認為,為了我國的生存,逼不得以犧牲他國--不知多少戰爭以保家衛國為號召,以生存危機為去動員.進入敵對關係,等於進入「你贏我輸」或「我贏你輸」的不能共存狀態.人生路上,有多少人是長勝將軍呢?有多少人未曾輸過呢?然而,總沒有人願意輸的,因為輸的感覺太難受了,輸的狀態太痛苦了.MSG把「你用油膏了我的頭」譯作「You revive my drooping head」.詩人輸了,一副垂頭喪氣的樣子,等待被敵人凌辱或處決.這些受凌辱,待處決的威嚇,不一定在戰場才出現,有些人可能在學校經驗過,可能在工作間經驗過.

四.總結-驀然回首,被恩惠慈愛追趕一生

最後一節,詩人說:

23:6我一生一世必有恩惠慈愛隨著我；我且要住在耶和華的殿中,直到永遠.6 Your beauty and love chase after me every day of my life. I’m back home in the house of God for the rest of my life.

在爭扎求存的路途上,怎可能看見「有恩惠慈愛隨著我」呢?Steve Jobs對在史丹福大學畢業禮上曾說:「你沒辦法預見這些點滴如何聯繫,唯有透過回顧,可以看出彼此關聯.所以你必須相信,無論如何,這些點滴會在未來互相連結,有些東西你必須相信,像你的直覺,天命,人生,因果,諸如此類種種.這樣的想法讓我永遠不沮喪灰心,也的確塑造了我人生中所有的不同.」

讓我們也回望一下,在追追趕趕或被被追追趕趕的求生路上,除了「死期」,除了老闆,除了敵人,還有其他在追趕我們呢?詩人是如何發現原來追趕來他的是上帝的恩惠和慈愛呢?

23:2他使我躺臥在青草地上,領我在可安歇的水邊.2 You have bedded me down in lush meadows, you find me quiet pools to drink from.

要爭要搶,往往因為感到「缺乏」的威脅.如果感到足夠,就不需要爭和搶了.創造主把人類放在地球,這片大地能夠出足夠的糧食養活人類嗎?聯合國世界糧食計劃指出:單從數量上說,全球的糧食足夠世界上所有人口所需.樂施會認為,飢餓源於貧窮,貧窮源於不公平.不公平,並非由「搵食」這需要造成,而是由不良的「搵食」方法造成.詩句描述,「躺臥在青草地上」暗示青草茂盛,飽食之餘尚可躺臥；「安歇」是用來形容水,平靜的湖水讓羊群容易飲用.詩人的暗示是,創造主的供應是充裕的.

23:3他使我的靈魂甦醒,為自己的名引導我走義路.3 True to your word, you let me catch my breath and send me in the right direction.

生存的憂慮有時令人陷入不擇手段的引誘；另一方面,勝利的感覺有時令人充昏頭腦.當詩人走離了跑道,上主的恩惠與慈愛透過他的話語展示,當頭棒喝,叫他臨崖勒馬,重回正軌,免得他墮落萬丈深淵.

23:4我雖然行過死蔭的幽谷,也不怕遭害,因為你與我同在；你的杖,你的竿,都安慰我.4 Even when the way goes through Death Valley, I’m not afraid when you walk at my side. Your trusty shepherd’s crook makes me feel secure.

陷入絕境之中,上主的恩惠和慈愛仍追來嗎?原來,上帝是那位「黑暗也不能遮蔽我使你不見,黑夜卻如白晝發亮,黑暗和光明在你看都是一樣」(詩139:12)的上帝.縱使詩人失足跌進死蔭幽谷,他發現上帝也在那裡,並且透過管教和保謢展示祂的恩惠及慈愛,這正是「不怕遭害」的基礎.

23:5在我敵人面前,你為我擺設筵席；你用油膏了我的頭,使我的福杯滿溢.5 You serve me a six-course dinner right in front of my enemies. You revive my drooping head; my cup brims with blessing.

成王敗寇,準備接受凌辱之際,那隻接觸詩人的手,原來是上帝的,把他垂下的頭抬起,挽救他的人格尊嚴,支撐他站起來重新做人.或許,在敵人眼中,甚至在自己眼中,我們失敗了；但詩人告訴我們,上帝是那位「壓傷的蘆葦,祂不折斷；將殘的火把,祂不吹滅 」(以賽亞書42:3)的上帝.詩人回望一生,發現上帝如此待他,因此說出這首詩的第一句:

23:1耶和華是我的牧者,我必不致缺乏.1 God, my shepherd! I don’t need a thing.

幸福摩天輪的主角在「追追趕趕　高高低低」,「驚險的程度叫畏高者昏迷」,「東歪西倒　忽高忽低」的狀況下,能夠「不再畏高」,「不怕跌低」,不怕「意外脫軌」,因為「多僥倖跟你共同面對」.詩篇二十三篇的作者提醒我們,是甚麼讓我們在人生路上安身立命呢?豈不是因為有那位以祂的恩惠與慈愛追著我們不放的上主嗎?



\newpage

\section{時日不多了}
\label{sec:366}
\textbf{許開明牧師 - 時日不多了!}
\newline
\newline
\hyperref[sec:365]{< < < PREV SERMON < < <}
~
\hyperref[sec:toc]{[返目錄]}
~
\hyperref[sec:367]{> > > NEXT SERMON > > >}


約翰福音 14:1-14
\newline
\begin{longtable}{cl}
\hline
\hline
章節 & 經文 (和合本修訂版)\\
\hline
14:1 & \begin{tabularx}{0.7\textwidth}{X} 「你們心裡不要憂愁；你們信神，也當信我。 \end{tabularx} \\ \\ \relax
14:2 & \begin{tabularx}{0.7\textwidth}{X} 在我父的家裡有許多住處；若是沒有，我就早已告訴你們了。我去原是為你們預備地方去。 \end{tabularx} \\ \\ \relax
14:3 & \begin{tabularx}{0.7\textwidth}{X} 我若去為你們預備了地方，就必再來接你們到我那裡去，我在哪裡，叫你們也在哪裡。 \end{tabularx} \\ \\ \relax
14:4 & \begin{tabularx}{0.7\textwidth}{X} 我往哪裡去，你們知道那條路。」 \end{tabularx} \\ \\ \relax
14:5 & \begin{tabularx}{0.7\textwidth}{X} 多馬對他說：「主啊，我們不知道你去哪裡，怎麼能知道那條路呢？」 \end{tabularx} \\ \\ \relax
14:6 & \begin{tabularx}{0.7\textwidth}{X} 耶穌對他說：「我就是道路、真理、生命；若不藉著我，沒有人能到父那裡去。 \end{tabularx} \\ \\ \relax
14:7 & \begin{tabularx}{0.7\textwidth}{X} 既然你們認識了我，也會認識我的父。從今以後，你們就認識他，並且已經看見他了。」 \end{tabularx} \\ \\ \relax
14:8 & \begin{tabularx}{0.7\textwidth}{X} 腓力對他說：「主啊，將父顯給我們看，我們就知足了。」 \end{tabularx} \\ \\ \relax
14:9 & \begin{tabularx}{0.7\textwidth}{X} 耶穌對他說：「腓力，我與你們在一起這麼久了，你還不認識我嗎？看見我的就是看見了父，你怎麼還說『將父顯給我們看』呢？ \end{tabularx} \\ \\ \relax
14:10 & \begin{tabularx}{0.7\textwidth}{X} 我在父裡面，父在我裡面，你不信嗎？我對你們所說的話不是憑著自己說的，而是住在我裡面的父在做他的工作。 \end{tabularx} \\ \\ \relax
14:11 & \begin{tabularx}{0.7\textwidth}{X} 你們要信我，我在父裡面，父在我裡面；即使不信，也要因我所做的工作信我。 \end{tabularx} \\ \\ \relax
14:12 & \begin{tabularx}{0.7\textwidth}{X} 我實實在在地告訴你們，我所做的工作，信我的人也要做，並且要做得比這些更大，因為我到父那裡去。 \end{tabularx} \\ \\ \relax
14:13 & \begin{tabularx}{0.7\textwidth}{X} 你們奉我的名無論求甚麼，我必成全，為了使父因兒子得榮耀。 \end{tabularx} \\ \\ \relax
14:14 & \begin{tabularx}{0.7\textwidth}{X} 你們若奉我的名向我求甚麼，我必成全。」 \end{tabularx} \\ \\
[1ex]
\hline
\hline
\end{longtable}
講題:時日不多了　Time is Running Out

經文:約翰福音14:1-14

講員:許開明牧師

日期:2017年5月14日

各位,今早的經文提及耶穌快要上十字架,快要離開門徒,他在世的時日不多了.在這時刻,耶穌說了什麼?在說這之前,讓我先分享一個故事.

一位病人看完醫生後就打電話給太太,而且不停地哭.太太問:「發生何事?」兄弟說:「醫生說我要服藥,每日一粒,不按時服就會死.」太太答:「那就每天吃藥吧!」兄弟說:「但醫生只給我4粒藥,我只剩四天命.」

如果你是那個人,生命只剩四天,你會如何過日子?時日無多了,用外國人常說的話就是「Time is Running Out!」時間在流失,在消耗中.

到殯儀館的時候,假如我不用當主持,我就愛坐著看靈堂的燭光,燭光會燃燒又會熄滅,同時會發出三樣東西:

1光:帶來照明；

2熱:帶來溫暖；

3煙霧:帶來二氧化炭,使空氣混濁.

人的生命也同樣,燃燒的時候會為身方邊的人帶來好的東西和不好的東西.

我們都會遭遇下列的景況:在世,時日無多了；在家,時日無多了；在學,時日無多了；在職,時日無多了；在位,時日無多了；而我任校牧,時日無多了；個人的健康,快樂,青春,美貌或不美貌,時日無多了；人與人的關係,時日無多了；我們在一起,時日無多了；還有很多很多「時日無多」的事物,這是非常值得我們正視的問題.

時日無多是人生的實況,為甚麼呢?為何不可長生不老?為何不能長期在位?為何不可長期擁有?「時日無多」,為人帶來痛苦,憂慮,不安,無奈,失望,失落等；也因著「時日無多」,人才會珍惜時日,珍惜現在,不只是時間,還有與家人,朋友,同事的關係和自己的身心.

人生是無常與短暫的,而永恆只是人的願望.自古到今,我想永恆永遠活著；永遠持久；永遠不變,年年十八.地球不停轉,行星不停運行,河水不斷向下流,生命不斷波動.初次遇見的人,再見他時已跟初次不同了,你也改變了,成為兩個不同或不完全相同的人了!

1979年,我加入了崇基學院.那時的我,無論樣貌,肥瘦,學識,思想與1983年畢業的我不同,與我2013年9月到任校牧更為不同!今日你與許牧師的關係,到7月後也會改變了,因那時的我不再是這裡的校牧了.

時日不多,也會不同,也會改變,但自古到今,人總想抓住所擁有,所珍惜的,不想順應改變.

小時候的我很愛玩「tum tum 轉」,由慢到快,越轉越快時就想它停下來,甚至想它逆方向轉,要把快速轉動的「tum tum 轉」停下來,就要用雙倍的力,要逆轉就更要用多倍和更大的力,人更會因此而扭傷.

轉得太快會頭暈,更會因失去重心而跌倒；轉得很快時想要衝上去,就易生意外,非常危險.一位同學知道我轉得太快時會頭暈,所以有一次他就利用我這弱點,轉到我暈,使的大衛受驚.

時間是什麼?存在或不存在?看見或不看見?看不見時間,但能感受到存在和長短；看得見時間,卻是因為日出,日落,月圓,月缺映著時間段落.人類又把看不見的變成看得見,時針,時鐘,時分秒的發現或製作是人類偉大貢獻.其實時間是看得見的,是因同在時空中,是佛教中所說的緣份；是基督教中所說的上主的安排.

時間與人生關係緊密,時間有三大層次:昨天,今天,明天,量度時間有秒,分,時,日,月,年,使人有明確時間意識,好使人可以善用和把握時間.

昨天過去了就不能回頭；今天要好好把握；明天要有智慧的,負責任和慎重的計劃.昨天可能會後悔,自責,埋怨；今天可能會浪費,無所事事；明天會焦慮無目標.

時間有量,也有質,快樂與痛苦；滿足與虧欠；自私與慷慨.時間屬於誰?我的時間 ,你的時間,他的時間,是大家的時間:我,你,他和上帝的時間.你想擁有他人的時間,他人又想擁有你的時間,有些人的時間永遠被人或物佔有,被人的佔有例子如老闆,子女,父母,朋友,情人,配偶等；被物質佔有的例子有遊戲機,電腦,工作,毒品,睡眠,正負面情緒,痛苦,憂鬱等.誰能主動操控自己的時間,誰就是自由快樂的人.

耶穌惜別講詞

耶穌是道路,真理,生命

耶穌安慰門徒:「你們心裡不要憂愁」

因為:a在我父的家裡有許多住；

b我若去為你們預備了地方,必再來接你們到我那裡去   c必再來接你們到我那裡去,我往哪裡去,你們知道那條路

多馬說:「他不知耶穌往那裡去? 」

耶穌說:我就是道路,真理,生命；我是帶路者

腓力對他說:「請帶我們去見主,讓我們見到主!」

腓力對他說:「主啊,將父顯給我們看,我們就知足了.」

耶穌對他說:「腓力,你見到我,就是見主了,我就是父!」

耶穌告知腓力兩件事和使命:

使命1:我所做的工作,信我的人也要做,並且要做得比這些更大

使命2 :你們奉我的名無論求什麼,我必成全

耶穌的安慰

安慰,就是加給力量的意思.日子不多的是耶穌,受安慰的應是耶穌,但耶穌反而去安慰人；叫不要憂愁的應是耶穌,耶穌卻叫門徒不用擔憂；去預備地方給自己的應是耶穌,反而是耶穌為門徒預備地方.門徒會被恐懼,憂慮所充滿,活在痛苦中.

你有想過「如何可以殺人不見血」嗎?第二次世界大戰期間,在德國法西斯的集中營裡,納粹份子曾經用過一種殺人不見血的恐怖手段來殺人.首先,他們把犯人關進一間漆黑的小屋子裡,把犯人的手臂從牆上的一個小孔中拉出來,固定在一個架子上.然後,他們會告訴犯人將會放乾他的血,講完便立刻用刀背在犯人的手腕靜脈處來回劃動幾下.再在旁邊以水模彷血往下滴的聲音,令犯人以為自己的血正在不斷下滴,就這樣滴,滴,滴,犯人最終衰竭而死.事實上,犯人的手腕完好無損,納粹份子不用動刀動槍把犯人殺死.只因自己的「恐懼」,那怕死的恐懼,要了他的命.

富蘭克林曾經有句名言「人間最可怕的東西,就是『怕』!」還有,那誤信自己 「時日無多」的恐懼,那種活在恐懼中,生不如死,不想 「活下去」的感覺.

耶穌的安慰是:

1不要憂愁,不被憂慮所佔有,當信上帝應許；

2應許住處,當信耶穌預備地方；

3應許回來,應許接去；

4我往那裹,你們也在那裡；

5我會放下地址和地圖,讓你知道我去那裡,怎樣去；

6叫他們認識耶穌是誰?是他們要找的；

耶穌告知門徒,宣告:「我是 ego eimi (神的自稱)」

1道路:方向,方法,方針,方寸,路向,路標,路牌；

2真理:有價值,有意義,有道路,可依據,依靠；

3 生命:存在的,有生命力,有活力,有延伸.

如何使時間具有價值?

1道路:長短,計算方法

你還有多少時間:可計,也有不可計?

2真理:真實的,有價值的,有意義的,可信靠的,可依賴的

你如何善用,把握你的時間,使你擁有的時間具有意義?你是領導人,你如何善用下屬的時間?

3生命:延續存在與生命力,活力,動力

生命可孕育生命,可激發生命,影響生命.

耶穌關注人多方面的需要

1安慰就是加上力量和心靈需要；相信神就是靈性需要；

2為人預備地方就是為人實際生活需要；

3指引人的生命,即人生智慧,理性需要

4給予托付,使命,使人意志激發,耶穌激發人一 個人全人的需要

如何使時間的質量俱佳?

在資本主義的影響下,人們大多接受社會學家Max Weber的觀點:注重時間的經濟效益,利之所在,人要盡量善用時間:

1工作量化:「朝九晚五」的固定工作模式；

2時間量化:年,月,日,時,分,秒；

3生活量化:創立每天24小時制度；

4盡用一小時:例如崇拜1小時,開會1小時,休息1小時,車程1小時.

舉一個例子:桌上有一個瓶子,可以放入1高爾夫球,2小石,3細沙,4咖啡；就如人們如何分配時間,時間內可以有多項不同的項目混在一起,使人失去特色和焦點,以至人生定位；彼此間無聯繫關係,物與物分割,使人物之間變得沒有關係,變得冷漠；所有東西迫在一起,完全沒有空間,無法呼吸,使人失去自由.

人需要從這時間陷阱跳出來.

1注重專一才能有成效,正如日光照射放大鏡,聚焦才可生火；

2注重關係:人與人之間,人與物,人與環境；

3注重時間取捨,不接太多事情,讓自己得自由,做想做的事.

一般人都認為死亡是虛無,恐欋,悲傷的,但海德格卻抱有積極的看法,就是死亡是早已有決定的,由出生那天,我們的日子有限,步往死亡,所以他會認真地生命,這就是「此在」,即認真地活,意識當下的每刻.

1死亡存有:有死亡意識,因為時間有限,時間無多,人才不會混噩過活；

2在世存有:意識到活在世上與各種人事發生關係,並為他們有所煩憂.

由此推理,時間的運用,注重關係建立:同事關係,親友關係,寵物關係,政府關係.道成肉身神學,上主進入人類的時空,歷史也是著重關係的建立.父母親節的臨近,大家都懂得怎樣做吧?

當我知道將要退休的時候,我在校牧室時日無多了.在學院刊登招聘廣告後,許多堂會及機構都有找我,結果我應邀到中華基督教會堂會服事.

我受到「時日無多」這概念影響,我還有五年,最多十年,應如何發揮我金色年華的生命?我還是強壯,可以多作主工.

1堂會:維繫及培訓在職同工,發展教堂及宣教；

2神學院:培訓教牧及神學生；

3寫作:我有三本學術書,三本靈修書要出版；

4教堂重建:我有三間堂建堂的經驗；

5區會:完成百周年的中華基督教會的神學計劃書.

時間比金錢貴重嗎?

當美國人說:「時間是金錢」,但猶太人十分反對,金錢只可儲蓄並生息,而時間卻不停腳步,而且一去不返.時間就是生命,時間就是一生,時間就是一切.作者李昂斯曾說:「昨天是一張作廢的支票,明天是一張期票,而今天則是你唯一擁有的現金,所以你應當聰明地把握. 」作早李大釗也說過:「昨天喚不回來,明天還不確實,你能確有把握的就是今天.最珍貴的是今天,最容易失去的也是今天. 」

活在當下

印度有個著名的故事,講述人生遇見三樣東西:老虎,毒蛇,水果.

在上的有老虎；在下的有毒蛇；在中間及面前的是水果,假如不吃蘋果就只有等老師吃死；還是享受蘋果後,然後給老虎吃?當你時日無多了,你會怎樣過活呢?

印度人又有另一智慧反思,就是當你面對著一條準備向你進攻的響尾蛇,你會怎樣呢?你應謹慎行事,不惹怒這蛇,才可保護自己,找機會逃走,因為你最重要的是擁有生命時間.

一篇令人回味的散文:

時日不多了!身邊的人要相互珍惜,因為每個人的時間會越來越少,最終還是要互相分離.不要爭執,不要鬥氣,好好說話,相互理解.好好珍惜對你好的人,失去了,怎樣也找不回來；離開了,互聯網也聯繫不上.生命,每個人只有一次,或長或短；生活,每個人都在繼續,或悲或喜；人生,每個人都在旅途,或起或跌.人無完美,事無完滿.有些小人,你不須計較,計教會煩；有些瑣事,你不必在意,在意會疲累.世界很大,人心很細,怎可能沒有遇小人?紅塵很深,人世浮華,怎可能沒有心煩事?想得淺一點,活著就該隨遇而安；看得淡一點,頭上就有一大片藍天.

誰都會害怕失去,因此是拼了命的珍惜,後來卻還是什麼都留不住.

人到老年,還是順其自然吧,無可奈何也好,心態淡然也罷,認認真真地活在當下,自然發覺生活也並不是那麼難捱.用心數數,我們的生命已經過半,還有多少日子讓我們浪費?讓我們消遣?讓我們感嘆?好好的活著,簡簡單單地過好老年的每一天.

生命有限,時日不多了!你會怎樣過活?



\newpage

\section{當耶路撒冷遇上雅典}
\label{sec:367}
\textbf{許立中先生 - 當耶路撒冷遇上雅典}
\newline
\newline
\hyperref[sec:366]{< < < PREV SERMON < < <}
~
\hyperref[sec:toc]{[返目錄]}
~
\hyperref[sec:368]{> > > NEXT SERMON > > >}


使徒行傳 17:22-31
\newline
\begin{longtable}{cl}
\hline
\hline
章節 & 經文 (和合本修訂版)\\
\hline
17:22 & \begin{tabularx}{0.7\textwidth}{X} 保羅站在亞略巴古當中，說：「諸位雅典人！我看你們凡事很敬畏鬼神。 \end{tabularx} \\ \\ \relax
17:23 & \begin{tabularx}{0.7\textwidth}{X} 我到處走走的時候，仔細觀察你們所敬拜的，發現一座壇，上面寫著『獻給未識之神明』。你們所不認識而敬拜的，我現在向你們宣告： \end{tabularx} \\ \\ \relax
17:24 & \begin{tabularx}{0.7\textwidth}{X} 他是創造宇宙和其中萬物的神；他既是天地的主，就不住在人手所造的殿宇裡， \end{tabularx} \\ \\ \relax
17:25 & \begin{tabularx}{0.7\textwidth}{X} 也不用人手去服侍，好像缺少甚麼似的；自己倒將生命、氣息、萬物賜給萬人。 \end{tabularx} \\ \\ \relax
17:26 & \begin{tabularx}{0.7\textwidth}{X} 他從一人造出萬族，居住在全地面上，並且預先定準他們的年限和所住的疆界， \end{tabularx} \\ \\ \relax
17:27 & \begin{tabularx}{0.7\textwidth}{X} 為要使他們尋求神，或者可以揣摩而找到他，其實他離我們各人不遠。 \end{tabularx} \\ \\ \relax
17:28 & \begin{tabularx}{0.7\textwidth}{X} 我們生活、行動、存在都在於他。就如你們的詩人也有人說：『我們也是他所生的。』 \end{tabularx} \\ \\ \relax
17:29 & \begin{tabularx}{0.7\textwidth}{X} 既然我們是神所生的，就不應該以為神的神性像人用手藝和心思所雕刻的金、銀、石像一般。 \end{tabularx} \\ \\ \relax
17:30 & \begin{tabularx}{0.7\textwidth}{X} 世人蒙昧無知的時候，神並不追究，如今卻吩咐各處的人都要悔改。 \end{tabularx} \\ \\ \relax
17:31 & \begin{tabularx}{0.7\textwidth}{X} 因為他已經定了日子，要藉著他所設立的人按公義審判天下，並且使他從死人中復活，給萬人作可信的憑據。」 \end{tabularx} \\ \\
[1ex]
\hline
\hline
\end{longtable}
講題:當耶路撒冷遇上雅典　Jerusalem and Athens

經文:使徒行傳17:22-31

講員:許立中先生

場合:香港中文大學崇基學院禮拜堂

日期:2017年5月21日

當耶路撒冷遇上雅典 (徒17:22-31)

在今天的新約經課,我們比較罕有地看見宗教跟文化的正面交鋒.使徒保羅嘗試跨越宗教,文化,向當時西方世界無論在思想或政制都是最先進的城邦,解釋自己信仰,示範了跨文化信仰對談的可能.

「正統文化」與普及文化

當時希臘大概已經過了她最光輝的日子.真正的哲學大師,都已經成為歷史；留下來的,根據《使徒行傳》的記載,就只有那些拾人牙慧,「只將新聞說說聽聽」的「吹水友」.

事實上,如果我們從使徒行傳十七章開始看,就知道雅典其實並不包括在保羅的「宣教計劃」之內.他們在當時馬其頓的首府帖撒羅尼迦逗留了一段不短的時間,起碼足夠讓他們在當地建立了教會.但保羅有關耶穌的講論,亦挑起了一些誓死捍衛一神信仰的猶太教徒的不滿.

聖經記載那些猶太人招聚市井匪類,聳動群眾反對他們,將收留他們的家庭拉到官府,控告他們違背命令,在凱撒之外另立新王(因基督徒稱耶穌為「主」‘Kyrios’,而那亦是羅馬屬土公民對凱撒的稱呼).保羅和他的拍檔西拉,唯有連夜離開帖撒羅尼迦,去到另一個城市Berea.

但是到了Berea,他們又慣常地去到當地的猶太會堂,透過日日查考聖經,要證明耶穌是基督,就是那位舊約聖經應許將要降臨的彌賽亞.使徒行傳記載,「他們中間多有相信的,又有希臘尊貴(思想開明)的婦女,男人也不少.」

但帖撒羅尼迦的猶太人知道保羅又在Berea傳道,就跟著去,聳動群眾跟他們作對.當地的信徒唯有將保羅送走,留下他的兩個拍檔西拉和提摩太.

送走保羅的人後來輾轉將他帶到雅典.保羅就請他們回去把西拉和提摩太也帶來.

出乎意料之外,希臘的首都雅典,並不是一般人以為那麼高品味文化,「談笑有鴻儒,往來無白丁」的文化之都.當然那裡仍然有依壁鳩魯學派和斯多亞學派的學士,但保羅在雅典市面看見的,並不是滿街都是咖啡廳,「文化沙龍」,而是「滿城都是偶像」.

表面看來,雅典作為一個哲學之都,跟她「滿城都是偶像」似乎有點格格不入.不過一個城市其實並不是那麼純粹的.譬如我們常說「香港人」怎樣怎樣,deserve 些甚麼；但事實上香港是有很多人,而這些人是可以相差很遠的.可以有吳靄儀,也可以有周浩鼎.

香港是一個相當現代化甚至是「知識型」的都市.在這彈丸之地,就有二十間可頒授學位的高等院校,包括八間由大學教育資助委員會(UGC)用公帑資助的院校；我們每年在會議展覽中心舉辦最大型的國際書展,並且每次都並非門庭冷落,門檻羅雀,而是大打蛇餅,將場館擠得水泄不通.我們說這足證香港並非「文化沙漠」.

有趣的是書展的主辦機構,是貿易發展局而不是文化或教育局.我們民間最暢銷的作者,是蔡伯勵,李居明,蘇民峰,麥玲玲.我並非看不起勘輿術數,而是想指出,科學技術,理性文明是絕對可以跟占卦算命,民間宗教共存!

「證據」背後的假設

今日講道的題目「當耶路撒冷遇上雅典」,或許會使你聯想到「北京遇上西雅圖」.無論如何,那是代表了兩種文化的相遇和交鋒.其實當我讀到今日的經課,我是聯想到一位Westminster Theological Seminary的教授Cornelius Van Til.

在Van Til 75歲的時候,學術界的朋友為他編寫了一本紀念文集(Festschrift), 書名就叫做“Jerusalem and Athens”.

Van Til 對基督教神學的最大貢獻,是他提出了在哲學或神學上的所謂「前設論」(‘Presuppositionalism’).他指出任何觀點和論據背後,自覺或不自覺地,其實都有一套想當然的,taken for granted的前設.而在那些被認為是不證自明的,無需討論的,想當然的前設框架底下,任何陳述都可以找到它們的合法理據.

譬如,倘若在你的理論架構已經假設沒有上帝,那麼無論你怎樣推論,都不可能得出上帝存在的結論.因為你已經前設性地否定了這個可能性.當然,如過你已經相信上帝,那麼「自從造天地以來,上帝的永能和神性是明明可知,雖是眼不能見,但藉著所造之物就可以曉得,叫人無可推委」.因此,一切的討論或辯論,倘若不觸及那最根本,底層的假設,就只不過是隔靴搔癢,beat around the bush!

根據《使徒行傳》的記載,雅典當時仍然流行的兩個哲學派別是伊壁鳩魯學派和斯多亞派.

Epicureanism: 一般稱之為「享樂主義」,卻非我們今日所謂「今朝有酒今朝醉」的膚淺及時行樂主義；它不但追求肉體的快樂,更強調精神的快樂,譬如交朋友和欣賞藝術這些較高層次的樂趣.他們甚至相信自我慾望必須節制；節制慾望,才可平和心境,而和平的心境可以幫助我們忍受痛苦.

Epicurus 相信Democritus 的原子論,但他並不認同原子的運動受各種自然法則的支配.他認為人死後,靈魂原子離肉體而去,四處飛散,因此人死後並沒有生命.他說:死亡跟我們沒有關係,因為只要我們一日仍然在生,死亡就不會來臨；而當死亡來臨的時候,我們亦已經不再存在.因此死亡並沒有甚麼可怕的.

Stoicism: 創始人Zeno,主張泛神主義物質一元論,強調神,自然與人為一體,「神」是宇宙的靈魂和智慧,其理性滲透整個宇宙.個體小我必須依照自然的法則而生活,融合於整個大自然.斯多亞派似乎比較合統治者的口味,亞歷山大大帝以後,許多羅馬王帝都自稱是斯多亞派.尼祿皇的mentor Seneca,以及二世紀羅馬皇帝Marcus Aurelius,都是斯多亞派的代表學者.

可以看見,任何哲學去到最底層,往往都是無法證實或證偽的推想或假設.因此現代人比較趨向科學的法則,因為科學比較具體和可以精確推算.不過,科學可以推算的往往只是具體的物質世界；要從科學理論推論出道德價值,甚至應該怎樣做人,恐怕之間還有難以踰越的鴻溝.事實上在哲學的層面,科學的範式/代模(paradigm)亦不斷經歷改變.

譬如學院早前辦了一個「何謂生命?」的跨學科對談,講者是物理學,生物化學及分子遺傳學的專家.他幾堂主要的講論,我都有參加.他在演講中,對生命的種種作出了很多客觀的陳述,有宏觀的和微觀的,有機和無機的互動,有序及無序狀態的並存.那次演講,令我記起他早前在神學樓提到Francis Schaffer的一本書,叫做“How should we then live?”

我的問題正是,當一位科學家從他的實驗室回到家中,他先前對生命的研究,觀察和理解,會怎樣影響和指導他的具體生活?抑或一個科學家其實亦只能擘裂於兩個沒有必然關連的世界(學術世界與現實世界),雖然這兩個世界是同樣地真實.

跨文化的信仰對談

回到我們今日的經文.那些伊壁鳩魯和斯多亞派的學士,因保羅傳講耶穌與復活的道,就把他帶到亞略巴古.

Areopagus:是一個地方也是一個議會.在希臘文,pagus 指大塊的岩石；Areo 可能來自山下有一座供奉Areios 的神廟.後來,羅馬人以羅馬戰神Mars 將它命名為“Mars Hill”,因此阿略巴古又叫做「戰神山」.

阿略巴古在古典時期是雅典刑事和民事案件的上訴法院.到了新約時代,大概仍然保留某種議事的功能.雅典人將保羅帶到阿略巴古,大概就是要評斷一下他所發表的道理.

1644年,英國詩人John Milton向英國國會發表了一篇捍衛出版自由的小冊子,就用了“Areopagitica”作為標題: “A speech of Mr. John Milton for the liberty of unlicensed printing tothe Parliament of England”意思是按著阿略巴古的古希臘審差傳統,也沒有英國國會當時提出要求事先審查的規矩.

事實上今天我們看電視轉播的英國國會辯論,議員可以隨意起來發言,其他議員即時和應,非常活潑和熱鬧,似乎仍保留了早年阿略巴古的希臘式議會辯論!

透過《使徒行傳》的經文,我們幾乎可以感受到保羅站在阿略巴古發言時的認真和莊重:「保羅站在亞略巴古當中,說:眾位雅典人哪,我看你們凡事很敬畏鬼神.我遊行的時候,觀看你門所敬拜的,遇見一座壇,上面寫著『未識之神』.你們所不認識而敬拜的,我現在告訴你們.」

基督教信仰認為人比萬物有更特殊的價值和尊嚴,是因為人按著上帝形象被造；費爾巴克(Ludwig Feuerbach)則倒過來認為人往往擬人化地去想像上帝.因此不同的文化傳統,往往有不同的宗教想像,對上帝也就有不同的認知.

雅典作為一個當代的文化之都,自然有相當強的多元共容能力.他們不但尊敬鬼神,還恐怕因為自己的狹隘而有所遺漏.於是他們索性修一座壇,並註明是「未識之神」.

對於雅典人的滿天神佛,保羅並沒有一開始就予以否定,聖經只是說他「心裡著急」.對於他們的宗教多元主義,保羅只是客觀地表示,他觀察到他們「凡事很敬畏鬼神」,希望他們繼續以一個開放的態度,去考慮他跟著下來要說的.

他說:「創造宇宙和其中萬物的上帝,既是天地的主,就不住人手所造的殿,也不用人手服事,好像缺少甚麼；自己倒將生命,氣息,萬物,賜給萬人.」依壁鳩魯學派或許對於上帝有不同的看法,但是對於神居然要用人手服事這荒謬的觀點,卻是完全可以同意.

當然,說到「他從一本造出萬族的人,住在全地上,並且預先定準他們的年限和所住的疆界」,這似乎是他們比較陌生的聖經觀點.但相比於他們的原子論或者泛神論,這亦不是那麼難以嚥下.他跟著說:「其實他離我們各人不遠；我們生活,動作,存留,都在乎他.就如你們作詩的,有人說:『我們也是他所生的.』」

有學識文化的雅典人,都會知道「我們生活,動作,存留,都在乎他」(“For in Thee we live and move and have our being”)是出於希臘詩人Epimenides.其實更貼切的翻譯,是「我們生活,行動和存有,都是在你裡面」.而這觀點,是斯多亞派的哲學完全能夠擁抱的.

保羅是用他們自己既有的文化傳統去跟他們對質,跟著才順水推舟地指出「我們既是上帝所生的,就不當以為上帝的神性像人用手藝,心思所雕刻的金,銀,石.」保羅在這裡不單是以他們能夠理解的觀點,更是用屬於他們自己的文化理解,去處理他們偶像敬拜的問題.

當然,當保羅提到審判,復活的時候,又再惹來他們的譏誚.但他整篇的講論,卻完美地示範了跨文化信仰對談的可能.而事實上亦有幾個人都因此信了主,包括阿略巴古的官Dionysius 和幾個思想比較開明的希臘人.(17:34)

上帝無宗派,信仰無疆界

由保羅引用希臘詩人,讓我想到其實中國亦有非常豐富的文化傳統.如果我們相信上帝無宗派,信仰無疆界的話,我們又能不能用這些文化的產業,去進行跨文化的信仰對話?

我想起李白的《春夜宴桃李園序》:「夫天地者,萬物之逆旅；光陰者,百代之過客.而浮生若夢,為歡幾何?」

在這裡,李白論到「天地」,論到「光陰」,也就是現代人所說的「空間」與「時間」兩個維度.他指出,「空間」不過是萬物暫時棲息的旅舍,而「時間」則是穿越百代的過客；處身飄浮不定,猶如夢境的人生,歡樂的日子又能有多少?

從一個跨學科對話的角度,我們可以說李白在這裡是將空間與時間,重新安置於人的「生活世界」,為抽象的物理概念賦予血肉；或者說,他讓我們重新感受生活未被科學切割得支離破碎之前的面貌.這回應了我早前提出的問題:對生命的研究,觀察和理解,怎樣影響和指導我們的具體生活?

你不能說李白對時間和空間的描述「科學」抑或「不科學」,畢竟我們今天所理解的「科學」,在李白的時代尚未出現；但對於世代以來欣賞詩詞的人,李白這句詩卻絕對是一個真實到肉的描述,有「深淵與深淵響應」的共鳴.

同樣,被視為出自摩西手筆的《詩篇》九十篇,亦讓人對生命有深刻的體會和想像:「主啊,你世世代代作我們的居所.諸山未曾生出,地與世界你未曾造成,從亙古到永遠,你是上帝!」For in Thee we live and move and having our being!

一直以來,宗教,文化和科學都好像水火不相容.一些人更指出宗教與科學之間存在著基本方法論上的分歧:譬如「創造」是基於信念與權威,而「演化」則基於科學與觀察.但會不會創造與演化在本質上其實屬於兩個不同的範疇,而沒有實質的矛盾?

我的意思是,在萬物演化過程中留下的任何痕跡,都可以成為日後研究,分析和推論的具體數據；那當然是屬於科學的範疇.可是,譬如,在宇宙開始的那一瞬間,並沒有一個人在那裡進行觀察,記錄和數據分析.換句話說,科學當然不能以它所沒有的資料和數據,去論斷時空存在以前的事情.那已是metaphysics 而不是physics 了!

當聖經的作者寫「起初,上帝創造天地」的時候,世界早已進入常態運作；聖經作者當然不可能是創造事件的目擊證人.而信徒亦不可能視之為上帝創造的證據.但科學卻確實沒有比宗教更高的權威,去為宇宙的來由作最終的發言.

又有一次參加一個書院辦的生活午餐講座.他們邀請了一位腦神經科教授,分享一個「活化大腦」的講題.開講之後,才知道講者是講論禪修對大腦和身體功能的好處.我當時的感覺是,當一位學者意外皈依一種宗教/哲學/生活方式,他總能夠將自己的學識專長帶到這興趣裡去,為這意外開發的新領域,賦予一件體面的學術外衣.

有趣的是,在我的席上剛好有一位資深的哲學系教授,講座後甫即忍不住指出,就算我們對嘉賓講者引用的科學數據不予爭議,但怎樣去運用(manipulate)那些數據,卻大有商榷的餘地.只是很多時這些「見證者」往往挾著斐然的學術成就,讓人產生一種錯覺,以為「某學科支持甚至證實某活動的客觀科學性」.更重要的是,你可以從腦神經科學的角度去考慮禪修,但對於佛教徒來說,禪修卻完全是另一個層次的事情.

公平點說,有時某些基督徒動輒就說「科學證明禱告大有功效」,或者「統計顯示基督徒更為成功」,我誠實地對於這類宣稱感到極端不安.因為信仰確實不是這個層次的事情.

有一次,愛恩斯坦被問到他對信仰的體會.他說「催生真正藝術和科學的基本情感,就是對神秘的體驗.要是誰對這樣的情感不再感到好奇,不再感到驚訝,或是根本毫無感知,那他確實無異於行屍走肉,就像一截被掐滅的蠟燭.感覺到萬物背後那心靈無法掌握的存在,它們的幽邃與美麗只會間接觸及我們 － 這就是信仰的涵意.在這個意義上,也只有在這個意義上,我是個具有虔誠信仰的人.」

倘若信仰是關乎「萬物背後那心靈無法掌握的存在」,那麼當你聽到科學或宗教「證明」或「推翻」對方一些甚麼的時候,你幾乎可以肯定至少其中一方是被扭曲了.

我總覺得,一個謙卑誠實的不可知論者,比一個無所不知的基督徒更有指望.

保羅在雅典戰神山舌戰群雄,並不是要表現他的學識淵博,而是因為他看見雅典人的迷失和無知而感到「心裡著急」.他所說的,並不是一堆經過高溫消毒,不經人手,包裝妥當的「信仰懶人包」；他不介意跟不同宗教,不同文化,不同信仰的人對話,因為所說的,只不過是他知道並相信的.而這就是我在今天的經文所學到的功課.Ω



\newpage

\section{尾場的扣問}
\label{sec:368}
\textbf{關瑞文教授  - 尾場的扣問}
\newline
\newline
\hyperref[sec:367]{< < < PREV SERMON < < <}
~
\hyperref[sec:toc]{[返目錄]}
~
\hyperref[sec:369]{> > > NEXT SERMON > > >}


講題:尾場的扣問　Disciples’ Last Question for Jesus

講員:關瑞文教授

場合:香港中文大學崇基學院禮拜堂

日期:2017年5月28日

弟兄姊妹,按教會年曆,今天是復活期第七主日(或升天日後主日),是大公教會聚焦在復活事件的最後一個主日.下主日已經是聖靈降臨主日,是思想聖靈降臨以至初期教會運動正式開展的節期.今主日,就居於主復活和聖靈降臨之間.

復活,是基督徒盼望之所在.保羅說,“17基督若沒有復活,你們的信便是徒然 … 我們若靠基督只在今生有指望,就算比眾人更可憐… 32 若死人不復活,我們就吃吃喝喝吧!因為明天要死了” (林前15).盼望,自有基督信仰以來,其重要性都是重中之重.然而,對生活在今天香港的我們,盼望,就比重要更加重要了.

在過去的十年八載裡,香港經歷了前所未有的變化.不少人覺得,香港的核心價值,以致香港所以成為香港賴以的素質和制度甚至是文化等等,正受到方方面面的威脅,甚至被破壞.年輕人看不見光明的前路,不少成年人也害怕一覺醒來發現香港再不是他們認識的香港.於是,一些人嘗試起來盡點努力來保存他們心目中的獅子山下.可是,無論柔和的舉措,激烈的行動,似乎都屢戰屢敗.情況就好像甩繩馬騮般失控,每下愈況,人在當中越來越無奈,無能感以幾何級數滋長.

無能感是消磨盼望的最佳燃料.心理學的研究指出,充滿盼望的人,內裡是有一股很強的能動感的 (sense of agency).即是說,充滿盼望的人,會自覺“有能為力”的.無能感正是能動感的剋星!

其實,我們不需要一定提升到社會政治層面來談及無能感以致失掉盼望的實況.就是在我們個別的人生,如果我們都是屢戰屢敗,無奈無能,前路未卜,所渴望的久久不能成事,我們還哪裡來生之動力呢?如果我們都覺得現實是充滿未知之數,而人在當中又是屢戰屢敗,總感到無能為力,我們還憑什麼生而有望呢?

就著盼望這個信點,在復活期最後一個主日,我想與大家一起看看,當耶穌升天之際,門徒最後問的一個問題.

這個時刻,是基督道成肉身在世生活三十多年最後的一幕.就在這個尾場,門徒來了一個扣問:第六節:「主啊,你復興以色列國就在這時候嗎?」我們要猜一猜,當耶穌聽到這個發問時,內心會有何種感覺?開心?傷心?認同?無奈?氣憤?今天早上,我邀請大家懷著最大的同理心,跳進耶穌當時的內心,去感受他所感受的.我們會從耶穌的內心感受,去把握基督信仰所賜下的盼望,反思一下我們憑什麼不會丟失盼望.

盼望的保證到底是什麼?或許,有些人會答,是聖靈所賜的能力支撐著我們的盼望,而這能力使我們“有能為力”,勝利凱旋.因為耶穌明明這樣回答門徒的:“8但聖靈降臨在你們身上,你們就必得著能力…”.是聖靈所賜的能力,確保我們不會丟失的盼望嗎?是,也不是.我們必須進深一步理解這段經文.讓我們仔細看看.

在這尾場的一幕,門徒問:「主啊,你復興以色列國就在這時候嗎?」這一問,肯定讓耶穌心裡一沉,沮喪萬分.正如John Stott所言,這個問題的遣詞用字,充分洩露了門徒完全對上帝心意的無知.John Stott說:

·         “The verb restore shows that they were expecting a political and territorial kingdom”: “復興”這動詞,顯示他們以為上帝的國是在地上劃土為疆的政權.

·         “the noun Israel that they were expecting a national kingdom” : “以色列國”這個名詞,顯示他們以為上帝的國是一個國家.

·         “and the adverbial clause at this time that they were expecting its immediate establishment”:“在這時候”這個副詞子句,顯示他們以為上帝的國是現在就完整降臨.

耶穌要離開他們了,他們還是冥頑不靈,愚昧頑固,仍然是一如當初,以為耶穌是軍事式的復國主義彌賽亞.V3-5說:3他受害之後,用許多的憑據將自己活活地顯給使徒看,四十天之久向他們顯現,講說上帝國的事.4耶穌和他們聚集的時候,囑咐他們說:「不要離開耶路撒冷,要等候父所應許的,就是你們聽見我說過的.5約翰是用水施洗,但不多幾日,你們要受聖靈的洗.」耶穌復活後四十天之久,向他們顯現,講說上帝國的事,又預告他們要受聖靈的洗.講了四十天課,門徒竟然以為受聖靈的洗,就是上帝大能復興以色列國的時候!激心吧.如果你是當老師的,這種激心,相信你會很明白.情況就好像你上課時,高高昂昂興興奮奮把你的研究心得教導學生.學生也對你的研究很感興趣,非常專心,時而點頭,時而叫好.可是,到了Q and A了-他們所提出的問題,才讓你發現,原來他們沒有聽懂.那個時刻,你一定感到很沮喪.

不過,門徒如此渴望耶穌復興以色列國,是蠻有原因的.每個猶太人,從小時候開始,每天三次以〈十八頌禱〉Eighteen Benedictions虔誠禱告.Eighteen Benedictions中的第十四條,正正是求主復興大衛的國.猶太人從小受國族文化薰陶,以復興以色列國為盼望的核心.也許,當天早上他們才剛剛念誦〈十八頌禱〉的第十四條,禱告“求主復興大衛的國”,聽到耶穌要離開他們,還以為這個軍事彌賽亞現在就要復興以色列國.雖然門徒是遲鈍無知,但似乎無可厚非.

不過,耶穌可能不是單單慨歎他們的遲鈍無知.耶穌很可能替他們擔心.耶穌心知,復興以色列國,真的不在這時.復活,並沒有帶來復國,並沒有為猶太人帶來他們幾個世紀以來引頸以待的期望.耶穌知道,就是當這批尾場扣問的門徒通通死掉,這代完全過去,以色列國還不會復興,甚至耶穌也還是未曾再來.如果他們繼續把盼望建基在他們的復國夢上,他們早晚喪掉盼望,意志消沉,信心盡失.於是,耶穌就向他們明白解釋,上最後一課:7耶穌對他們說:「父憑著自己的權柄所定的時候,日期,不是你們可以知道的.8但聖靈降臨在你們身上,你們就必得著能力,並要在耶路撒冷,猶太全地,和撒馬利亞,直到地極,作我的見證.」

7耶穌對他們說:「父憑著自己的權柄所定的時候,日期,不是你們可以知道的-在此,耶穌要警告他們,不要把盼望建基在夢想的實現上.你們如果要生而有望,必須轉移視線!

8但聖靈降臨在你們身上,你們就必得著能力,並要在耶路撒冷,猶太全地,和撒馬利亞,直到地極,作我的見證.」-見證,在初期教會,有為主殉道之意.耶穌向他們明言,盼望不單只不是得著什麼夢想的實現,盼望是當你失去所有,連生命也失去,一無所有,還是可以因為聖靈在你們身上,你們就必得著能力.這顯然不是節節勝利的凱旋者的盼望力量,而是被釘十架被世人以為一無所有的人的盼望力量.

夢想落空,一無所有的人,還有盼望力量的嗎?在說笑嗎?耶穌在唱高調嗎?怎麼可能?你試著對一個患了重病渴望痊癒的人說,“你明天就死”,他不是會絕望嗎?怎麼反而有盼望?你試著對一個家暴受害者說,“你的暴力丈夫是不會改變的”,她不是會絕望嗎?怎麼反而有盼望?你試著對一個社會改革者預言,“有生之年,民主不單只不會實現,極權還會有恃無恐,變本加厲”,他不是會絕望嗎?怎麼反而有盼望?耶穌在說笑嗎?在唱不吃人間煙火的高調嗎?然而,這批門徒,在往後的日子裡,確實又真的到死也見不到復國,大部分甚至是死在強權刀下,死得一無所有.可是,他們一生卻又真的活得異常有能力,並且到處燃點盼望.這怎麼可能的呢?我們稍後再談.因為我們還要從經文裡多探討一點.

是哪一點呢?聽到門徒的尾場扣問,除了上面所提到的激心與擔心之外,耶穌的心情會是如何的呢?應該是傷心.

這時刻,是耶穌升天之際.要離開他們了,耶穌希望聽到他們問什麼呢?是「主啊,你復興以色列國就在這時候嗎?」不會吧!門徒的尾場扣問,其實等如在問,“主啊,你現在離開了嗎?是不是現在要把我們想要的給我們呢?”弟兄姊妹,試代入想一想.設若你是一位父親,因工作原因需要到海外生活一段不知有多長的時間.當你的摯愛兒子送你到機場入閘門前時,你期待他對你說什麼?大概是:“爸爸,你什麼時候回來?我會很想你.你那邊是怎樣的?會很冷的嗎?我捨不得你離開!”然後抱著你哭.不過,如果兒子說:“爸爸,你現在就離開了嗎?不久前你說過會買一套樓房給我,現在買了沒有?”你會怎樣反應?你會很傷心.你心裡會泛起一個問號:“到底這段父子關係是怎麼一回事?我與他的關係,真的只是‘給與取’的交易關係嗎?Exchange relationship? 在他心中,我真的只是一種工具性的存在嗎?噢,天啊!”

如果我是耶穌,聽到門徒的尾場扣問, “主啊,你現在離開了嗎?是不是現在要把我們想要的給我們呢?”我會立時想起抹大拉的馬利亞,尤其是主復活開場的那一幕.七天的第一天,抹大拉的馬利亞哭了.哭什麼?當天大清早,抹大拉的馬利亞來找我的屍首.她深知我已經死了,死了的人再沒有可能給她什麼.雖然我在生時滿有大能,也曾替她趕鬼,使他痊癒,她百分百瞭解,事到如今,我已死了,我的屍首再沒有能力為她做什麼.但她竟然肯冒著被羅馬政府逮捕之險來找我-是找我的屍首.我的屍首不見了.兩個門徒也跑來了,冷靜地仔細的研究了空墳一番,確定了抹大拉的馬利亞沒有胡言,就離開了,唯有抹大拉的馬利亞不願離開,還待在墳墓外面哭.遇到人了,就像瘋婦一樣焦急地的問:“先生,若是你把他移了去,請告訴我,你把他放在哪裡,我便去取他.”見到我復活了,她就馬上撲過來擁抱著我!這才是真正不離不棄的關係,這才是真的關係,這才是deep and authentic relationship.這個復活開場的哭,對比起那個復活尾場的扣問,真是高下立見.就是這樣,按初期教會的傳統說法,歷史上第一個傳復活盼望福音的人,並不是彼得,也不是雅各或者是約翰,而是抹大拉的馬利亞,並且抹大拉的馬利亞被稱為第一個使徒(the first apostle),在教會傳統裡又被譽為使徒中的使徒(apostle of the apostles).

這個復活開場的哭,對比起那個復活尾場的扣問,真是高下立見.前者是不離不棄的真關係,是deep and authentic relationship, 後者是“你給我取”的交易關係, exchange relationship.

盼望的源頭當然就是復活,但只有與上帝的 deep and authentic relationship才可以確保我們在一無所有時,在美夢不圓時,在柳暗花不明時,仍然不會丟失盼望.剛才我們問,“這批門徒,在往後的日子裡,到死也見不到復國,大部分甚至是死在強權刀下,死得一無所有.可是,他們一生卻又真的活得異常有能力,並且到處燃點盼望.這怎麼可能的呢?”因為,十天之後,聖靈降臨在他們身上,他們就得著能力.得著能力?什麼能力?節節得勝的能力?沒有,他們之後被拉,被鎖,被殺,始終敵不過強權?得著美夢成真的能力?沒有,他們至死也看不到十八頌禱裡所求的以色列國復興.弟兄姊妹,按上面所說,這種能力,其實是與三一上帝發生了最deep and authentic 的關係而生的能力,是聖靈內住於信徒而來的能力.

(羅8)“14因為凡被上帝的靈引導的,都是上帝的兒子.15你們所受的,不是奴僕的心,仍舊害怕；所受的,乃是兒子的心,因此我們呼叫:「阿爸!父!」16聖靈與我們的心同證我們是上帝的兒女；17既是兒女,便是後嗣,就是上帝的後嗣,和基督同作後嗣.如果我們和他一同受苦,也必和他一同得榮耀.18我想,現在的苦楚若比起將來要顯於我們的榮耀就不足介意了.”

天父遙遙在天,道成肉身的基督住在人群中間,而聖靈就在我們每一個人裡面,以致上帝活在我裡面,我活在上帝裡面.聖靈在我們裡面,每時每刻,提醒我們,尤其是當事事不如意時,美夢不能圓時,提醒我們,基督已經復活,讓復活的盼望,永遠不會在我們裡面丟失.有了聖靈在我們心裡,與三一神建立了deep and authentic relationship, 我們的盼望從此不假外求,不再連繫於外在的成敗得失,不再連繫於所渴望的心願何時才實現.就是失敗也不會失望；就是失去也不會久久哀傷；就是夢想不圓也仍是有動力,能力,繼續不問成敗地簡簡單單的去做應該做的事情,活應該活的生命,知其不可為而為之.這點非常重要,因為基督信仰的盼望,並不是基於打了勝仗帶來復國而有的盼望.基督信仰的盼望,倒是十字架上死得看似一無所有而生的復活盼望.

盼望當然是在盼望某些東西將會實現.盼望感的維持卻是基於 不離不棄充滿著聯繫感的deep and authentic relationship 的.心理學的研究用 generalized hope vs. particularized hope, 或者 definite hope vs. indefinite hope 來理解這個現象.因為時間關係,我們無法詳談這方面.我只是想指出,基督信仰所傳的盼望,並不是癡人說夢,而是可以有經驗科學基礎的.

盼望感的維持是基於不離不棄充滿著聯繫感的deep and authentic relationship ,聽下去好像有點抽象,因此,在結束時,我想用一個我曾牧養的個案,做一個類比.

當年牧養一個單親家庭.媽媽只有三十來歲,爸爸就在兒子出生後不久因患癌症去世.夫婦兩人原是青梅竹馬,愛得非常甜蜜.兒子只有兩歲,經常生病,媽媽常常為兒子,為生計奔波.然而,雖然生活非常艱難,不如意事又真的十常八九,媽媽卻生活得很有動力,生命充滿盼望.有一次,我問她,怎麼在艱難的日子仍然充滿盼望.她拿出手機,讓我看這個圖片:

圖片上的字:“If there ever comes a day when we can’t be together … Keep me in your heart, I’ll stay there forever.” 原來,這是她丈夫生前送給她的.她說:“他活在我的心裡,常常提醒我,無論如何,鼓起勇氣做人,他在這裡,不必氣餒.”

弟兄姊妹,8但聖靈降臨在你們身上,你們就必得著能力,並要在耶路撒冷,猶太全地,和撒馬利亞,直到地極,作我的見證.阿們.



\newpage

\section{願你們平安}
\label{sec:369}
\textbf{許開明牧師 - 願你們平安!}
\newline
\newline
\hyperref[sec:368]{< < < PREV SERMON < < <}
~
\hyperref[sec:toc]{[返目錄]}
~
\hyperref[sec:370]{> > > NEXT SERMON > > >}


約翰福音 20:19-29
\newline
\begin{longtable}{cl}
\hline
\hline
章節 & 經文 (和合本修訂版)\\
\hline
20:19 & \begin{tabularx}{0.7\textwidth}{X} 那日（就是七日的第一日）晚上，門徒因怕猶太人，所在的地方門都關了。耶穌來，站在當中，對他們說：「願你們平安！」 \end{tabularx} \\ \\ \relax
20:20 & \begin{tabularx}{0.7\textwidth}{X} 說了這話，他把手和肋旁給他們看。門徒一看見主就喜樂了。 \end{tabularx} \\ \\ \relax
20:21 & \begin{tabularx}{0.7\textwidth}{X} 於是耶穌又對他們說：「願你們平安！父怎樣差遣了我，我也照樣差遣你們。」 \end{tabularx} \\ \\ \relax
20:22 & \begin{tabularx}{0.7\textwidth}{X} 說了這話，他向他們吹一口氣，說：「領受聖靈吧！ \end{tabularx} \\ \\ \relax
20:23 & \begin{tabularx}{0.7\textwidth}{X} 你們赦免誰的罪，誰的罪就得赦免；你們不赦免誰的罪，誰的罪就不得赦免。」 \end{tabularx} \\ \\ \relax
20:24 & \begin{tabularx}{0.7\textwidth}{X} 那十二使徒中，有個叫低土馬的多馬，耶穌來的時候，他沒有和他們在一起。 \end{tabularx} \\ \\ \relax
20:25 & \begin{tabularx}{0.7\textwidth}{X} 其他的門徒就對他說：「我們已經看見主了。」多馬卻對他們說：「除非我看見他手上的釘痕，用我的指頭探入那釘痕，用我的手探入他的肋旁，我絕不信。」 \end{tabularx} \\ \\ \relax
20:26 & \begin{tabularx}{0.7\textwidth}{X} 過了八日，門徒又在屋裡，多馬也和他們在一起。門都關了，耶穌來，站在當中，說：「願你們平安！」 \end{tabularx} \\ \\ \relax
20:27 & \begin{tabularx}{0.7\textwidth}{X} 然後他對多馬說：「把你的指頭伸到這裡來，看看我的手；把你的手伸過來，探入我的肋旁。不要疑惑，總要信！」 \end{tabularx} \\ \\ \relax
20:28 & \begin{tabularx}{0.7\textwidth}{X} 多馬回答，對他說：「我的主！我的神！」 \end{tabularx} \\ \\ \relax
20:29 & \begin{tabularx}{0.7\textwidth}{X} 耶穌對他說：「你因為看見了我才信嗎？那沒有看見卻信的有福了。」 \end{tabularx} \\ \\
[1ex]
\hline
\hline
\end{longtable}
講題:願你們平安!　Peace be with you!

講員:許開明牧師

經文:約翰福音20:19-29

場合:香港中文大學崇基學院禮拜堂

日期:2017年6月4日

各位,上月我的講題是「時日無多了!」而今日的講題就是:「願你們平安!」

但時日無多,不是越來越令人不安嗎?對,就是因為越來越不安,所以要說:「願你們平安!」.

若大家留意今天的經課,約翰福音二十章中十九至二十九節中出現了三次「願你們平安!」.三次的出現,顯示出這句說話是多麼的重要.讓我們先看第十九節,第一個「願你們平安!」:「那日(就是七日的第一日)晚上,門徒因怕猶太人,所在的地方門都關了.耶穌來,站在當中,對他們說:『願你們平安!』」耶穌這樣說是因為他們在懼怕當中,主特別向他們說的.

接著,在二十一節:「願你們平安!父怎樣差遣了我,我也照樣差遣你們.」這是耶穌對門徒說.願他們有平安,接受上帝的差遣,為上主工作.

第三個「願你們平安!」,是在第二十六節:「過了八日,門徒又在屋裡,多馬也和他們在一起.門都關了,耶穌來,站在當中,說:『願你們平安!』」.大家留意,這段經文記載有一位門徒名叫多馬.在一星期前耶穌向門徒顯示的時候,多馬當時不在場.耶穌特意向他顯現,向他說:「願你們平安!」,因為多馬尚未得著這祝福.另一原因就是多馬曾經懷疑耶穌是否真的復活.門徒相信,唯獨他不相信.但今天他的主在他面前,多馬可能看到耶穌會十分震驚,因為他曾經十分懷疑,他一定十分害怕,耶穌會否質問他?但耶穌沒有,只對他說:「願你們平安!」.

有一地方是需要留意的,整後經文是記載在復活節的晚上,復活節的清晨,因耶穌的復活,是得勝的,光明的,喜樂的.

但復活節的晚上又如何? 一般對黑夜的理解,代表著危險,恐懼不安﹔月黑風高,令人不安.在外的人,在黑夜走路,可能遇上盜賊,野獸 ﹔乘船可能遇上風暴,水怪的危險, 晚上在屋內的人,可能遇上小偷.晚上在戰場,更害怕敵人進攻,犯事的怕被揖拿等.

經文記著當時「 門都關上」,顯示門徒在不安當中.怕人,怕猶太人,怕兵丁……  在懼怕中他們,那一刻最需要就是平安.

其實人類最需要的就是平安: 貧窮的人,因無錢不安﹔富有的人,因太多錢而不安﹔ 有名有利有勢的人,因無健康和時間不安﹔有時間的人,因無機會而不安﹔ 我們總會活在不安之中!

近代心理學家發現,人基本有三種的懼怕,你我經常都踫到:

第一類是怕聽到突然的巨響,延伸的意思是怕突如其來的災難,風災,雨災,火災,空難,海難,地震.

第二類是怕自己跌倒 ,意思是怕失敗,怕出錯,怕站不住,怕不合格,怕失去,失去健康,失去錢,失去生命…

第三類是怕被人遺棄 ,怕失戀,怕被父母,子女遺棄,被出賣,被人陷害,被人取笑,不被重視.

以上三種懼怕 可以是單一發生,但更亦有可能是互有關連的和相通的: 即是第一類懼怕,會引發第二類的懼怕,再引發第三類的懼怕. 使懼怕加劇,令人無法勝過,會抵受不住. 所以懼怕是人生最大的敵人.它足以快速毀滅人,也可以慢慢把人折磨至死.

門徒把門關上,因怕猶太人會突然來捉拿他們, 這是屬第一類(怕巨響的懼怕).門徒怕一旦被捉拿,便會被監禁,失自由；被打,失去健康；被殺失去生命,這是屬第二類(怕跌倒,怕失去的懼怕).門徒怕被捉後,被自已同胞猶太人,甚至家人親友的遺棄. 這明顯是屬第三類(怕被遺棄).

復活主來到,並站在他們當中,讓人人都見到他,並說: 「願你們平安」,這實在意義深長!

耶穌越過牆壁,越過大門,來到他們中間說:「願你們平安」. 到底耶穌所說的平安,是怎樣的平安? 「願你們平安!」

Shalom,就是「平安」的意思,是猶太人見面的互相祝福語. 今天基督徒見面時,也以「平安」作為彼此的問候與祝福； 「平安」的希伯來文是Shalom,在舊約出現有200多次,是「和平」的意思.但這字有三方面積極的含義: 第一,和平－不希望遇到困難或沒有戰爭； 第二是指向身體的健康,福利和安全； 第三指更指引人在凡事上達到最高,至善的地步.

而「平安」的希臘文是Eirene,這個字在新約聖經中,大概出現過90多次.「平安」在27卷新約書卷中都有出現過,因此有人說:「新約是一本談及平安的書.」

英文Peace這個字在中文有很豐富的意思,有三個中文翻譯:

第一,平安－特別指到內心的平靜,就是遇到任何困難都以安寧的態度去面對；

第二,和睦－指出在人際關係上,人與人之間友善的相處,彼此禮讓；

第三,和平－就是人群與人群(社會 ),國與國之間能和平共處,彼此之間沒有衝突,沒有任何爭鬥.

英文把它翻為peace.英文的翻譯(Peace be with you)更貼近聖經的原文, 含有 「平安跟隨 」 之意!

希臘文的Eirene並不是像中國人心中所想的順順利利,沒有困難,沒有煩惱,對於猶太人來說這是沒有可能,人的人生一定會面對不安.它真正的意思,是指向一種完美的幸福,一種令人導向至善的意思. 經歷不安以及經歷不安之後,而依然有平安,能夠勝過一切的風暴.因為平安在心中.

「平安」在聖經另外一個解釋是「連接在一起.」 .指與上帝有完美結連的時候,人就會有平安.很多時候信徒信主多年後,依然在不同的困難,風暴當中沒有平安,可能是因為我們與上主的連接出現了問題.當我們過分緊張,懼怕的時候,會給撒旦一個機會,令我們會做錯事,說錯話.而這個連結的靠著對上帝的信心.

人們一般經歷二種平安的來源. 一種是靠著從外來的環境,力量得著平安.一般人是求平安 是向外的,是靠外力的,但這種外來的,卻難以掌握.

而另一種是從內在的平安:內在的,「自己主場」.從上帝而來出人意外的平安:自己也無能,唯有信靠主大能 .事實上,平安是向內求的.是心態,是在乎自己,不一定向外求.外在環境,有許多人或事會令你不安:不要靠環境,不要靠人.學會敏銳有些什麼事或人最容易把人燒起來? 或可能使你恐懼的是蛇,老鼠,惡夢.

可求賜平安! 使不安心境變平安! 聖經上的平安是完全的,完整的,和諧,好的狀態(well being),不安就是感到不完整,不完全,不和諧,不好的狀態,不安全感!

耶穌在這裡的意思是:「願上帝將各樣的美善賜給你,願你們因此得到喜樂!」 所以當門徒看見復活的主,又得到祝福,得到大大的喜樂! 這「平安」是一種禮物,豐富的恩典. 這「平安」是一種好的願望. 這「平安」是一種力量和能力.

復活主是平安之主,有能力賜平安的主! 在門徒在懼怕之中的時候,耶穌穿牆過壁,突然在他們中間.當中的含義是:復活主的大能； 不用叩門,不用開門,免得他們受驚! 但如今這復活主,復活了! 在門徒眼中,復活主是相當奇妙的,也是奧妙難明的: 祂有身體,卻不被身體限制,仍然能穿牆過壁! 超越肉身,超越牆壁, 超越距離, 超越時間,超越的上帝.不被限制,不被阻攔! 復活主能越過關了的門,越過厚厚的牆壁,進到屋子裡! 甚至進入不信的心,進入懼怕的心, 進入不安的心, 進入懷疑人的心,改變人心,堅固他們的信心.

面對信仰一般人有三種反應,第一是聽而立即相信 ,一種單純的信心.另一種信心是要看見後才能相信.不能用耳朵代替眼睛,這是一般現代人的心態.第三種是要觸摸到才會相信. 這亦是一種從懷疑而進入信.多馬就是這一種.多疑的多馬,一位對復活有懷疑,敢於從理性的角度挑戰復活的門徒.那一刻,復活主竟然出現了!

多馬是否因為懷疑而沒有福呢?耶穌沒有這樣說.多馬從懷疑而進入確信,我覺得都是有福,而且更為有福.他對耶穌的認識是透過他的懷疑而有更深入的了解,對這個信仰是認真的尋索,耶穌同樣滿足他.因此他見到復活主,便宣告:「我的主,我的神」. 日後更把福音傳遍開去(相傳遠及東方的印度,中國),使人得著平安!

還有一個奇妙的地方, 就是耶穌要讓門徒看見,他是經歷過「不安」的主.聖經記載:「說了這話,就把手和肋旁,指給他們看；門徒看見主,就喜樂了.」(約翰福音20:20)耶穌就把手和肋旁,指給他們看,手的釘痕,肋旁的傷痕: 是從死復活的主,是擁有身體的主,不是人的錯覺,不是幻覺. 「看見」cido:不止是看到,是看清楚,看明白.他們就喜樂了! 這復活主也是受過試探,受過苦,受過痛,受過死的主,曾經被人唾棄,被釘在十字架,經歷門徒的離棄.

因此他能明白我們,知道我們的傷和痛的主.他能夠體會我們所經歷的不安,更加明白,關心我們,賜予我們平安.所以基督教強調的平安,不是空中樓閣,而是經歷「不安」過後,而進入「平安」.甚至在「不安」之中,還能保持「平安」的心面對.處於一種不安中的平安,懼怕中的平安,悲傷中的平安,活在不安中的平安.

耶穌明白他們的需要,不使他們受驚,在平安的環境中,把平安賜他們.然而,領受了「平安」,便要勇敢的,受差遣到不安的環境中,傳揚平安的福音.」基督教強調的平安是「在不安中的平安,經歷過不安後的平安」 聖經中的平安: 活在不安中的平安 願你們平安.門徒受差遣去傳道,主動去面對未來不安的境況.

一幅題為“平安”名畫, 給人第一印象是不平安:危岩峭壁,海水波濤凶湧,天空狂風暴雨. 但細心觀看:危岩峭壁中,有一棵樹,樹枝被狂風吹倒,樹上有雀巢,當中有「貓頭鷹」閉眼睡覺,非常安祥!

願你們平安:這種平安,真正意義是什麼呢? 在不安中,站起立祝福說:「願你們平安!」在恐欋中,你看見那位復活主,在暴風雨中的復活主!

海德格(Martin Heidegger,1889年9月26日－1976年5月26日) 說過 一般人看死亡都是虛無,恐欋,悲傷. 但人有積極的看法:如果我們能認清死亡的本相,知道自己生命有限,由出生那天已經步向死亡,我們就會認真地活著,意識當下的每刻,這就是「此在」 死亡的存有都在我們的心裡.有死亡意識,就明白時間有限,時間無多,人才不會混噩過活,使人生活得有秩序,心安理得.

在世存有是意識到活在世上與各種人事發生關係,並為他們有所煩憂.由此推理,時間的運用,要注重關係的建立:同事關係,親友關係,寵物關係,政府關係.

道成肉身神學,上主進人人類的時空,歷史也是著重關係的建立.

復活主差聖靈,賜平安

耶穌知道自己快升天,離開門徒,他差派聖靈! (說了這話,就向他們吹一口氣,說:你們受聖靈!)(約二十:22),聖靈與他們同在,讓他們能經歷到聖靈所賜的平安.

聖靈(Holy Spirit/ Holy Ghost),靈原文是“氣” 氣,令你想起什麼?

看不到, 與生命攸關(空氣,「無氣就無命!」),與力量有關(有氣有力 ,生氣,破壞力量,正氣:做正事的力量 ,勇氣:不怕困難的力量 聖靈是什麼?) 氣就是「氣息」 ,「(耶穌)說了這話,就向他們吹一口氣,說:你們受聖靈!」(約二十:22). 耶穌為何要給門徒吹氣?有象徵意義嗎?有何神學意義?

在創造的故事中,神將生氣吹在亞當鼻孔裡,亞當就成為有靈的活人. 為何神要給亞當吹氣,神說:「我們要照著我們的形像,按著我們的樣式造人,使他們管理海裡的魚,空中的鳥,地上的牲畜,和全地,並地上所爬的一切昆蟲」(創一:26).有靈的活人更加是指出 神要人有與祂相似的形像,因此吹一口氣在亞當的鼻孔裡, 就是要人與神有相類似的形像. 神要人能夠管理治理全地(能力,權柄) ,生養眾多(生育繁衍 ,神賜給他有智慧可以為各種動物命名(智慧).

但是,因亞當夏娃犯罪(失去),亞當夏娃人就失去屬靈的能力,失去人的能力,智慧,德性. 耶和華說:「人既屬乎血氣,我的靈就永遠不住在他裡面；然而他的日子還可到一百二十年」(創六:3)這一節經文是要表明人屬乎血氣的話,神的靈就無法繼續住在人的裡面.

神讓以西結看見遍滿骸骨的平原,「耶和華的靈(注:原文作“手”)降在我身上,耶和華藉他的靈帶我出去,將我放在平原中,這平原遍滿骸骨.」(結三十七:1)「於是,我遵命說預言.正說預言的時候,不料,有響聲,有地震；骨與骨互相聯絡.我觀看,見骸骨上有筋,也長了肉,又有皮遮蔽其上,只是還沒有氣息」(結三十七:7,8). 骸骨可以復活嗎? 可以! 氣息 rûach(風,靈),有神的氣息就有新的活力.當以西結向風發預言,氣息就進入骸骨,骸骨就活了,並且站起來,成為極大的軍隊(結三十七10).以西結發言:有神的氣息,使枯骨生出筋和骨,合起來成為強大的軍隊.

耶穌給門徒吹一口氣,門徒受聖靈,建立教會,掌握天國的匙,赦免誰的罪,誰的罪得赦.人受靈氣,重新再造,勝過懼怕,負起責任和使命,勇往直前.門徒受聖靈,仍然會面對恐懼不安,只是能夠勝過恐懼,在恐懼中前進:聖靈的同在,帶來平安平安是一種定力.

一位兄弟向雜技員學行銅線.先學行麻繩,且離地一尺,學了一年,都是如此!第二年,才行鋼線.為什麼?因要克服恐懼,人才能有定力,勝過懼怕.聖靈的加力,得著氣,得著氣力聖靈的推動力,勇敢面對懼怕!我們得著復活主的平安,領受聖靈作和平之子

中古時的依納爵修士在靈修中,他發現一個有平安的人或有安全感的人,要常常體會我們是上帝所創造的.有下列三個心態:

我是好的  I’m good. (一種完全的,健全的感覺)

我是被愛的  I’m loved.(一種被愛,被重視,被眷顧的感覺)

我是安穩的  I’m safe.(一種和諧的,和平的感覺)

經歷神所賜的平安,這是恩典.結出平安的果子來,這也是恩典,是生命的彰顯.

但基督徒是不應停留在單單領受平安地步,他們的生命應結出果子,也要成為和平之子.是的,每一位基督徒應該是一個和平之子,我們蒙召作神的兒女,更要成為一位傳揚和平的使者.

英文有兩個字,Peacekeeper 和Peacemaker,兩者有所不同.

Peacekeeper是維持或保障自己和別人有平安的人.

作丈夫的你,無平安,你太太,子女也無平安!

作爸,媽的你,你無平安,使整個家庭也無平安!

一個心中有意平安的人,這種人是什樣的人?

最後願以法蘭西的禱文作為結束和互勉:

(聖法蘭西的禱文)

上主,使我作你和平之子

在憎恨之處播下你的愛

在傷痕之處播下你的寬恕

在懷疑之處播下你的信心

在絕望之處播下你的盼望

在幽暗之處播下你的光明

在憂傷之處播下你的歡愉

哦,主啊使我少為自己求

少求受安慰,但求安慰人

少求被瞭解,但求瞭解人

少求愛,但求全心付出愛

在捨去時,我們便有所得

在赦免時,我們便蒙赦免

迎接死亡時,我們便進入永生

Lord, make me an instrument of your peace.

Where there is hatred, let me sow love.

Where there is injury, pardon.

Where there is doubt, faith.

Where there is despair, hope.

Where there is darkness, light.

Where there is sadness, joy.

O Divine Master,

grant that I may not so much seek to be consoled, as to console;

to be understood, as to understand;

to be loved, as to love.

For it is in giving that we receive.

It is in pardoning that we are pardoned,

and it is in dying that we are born to Eternal Life.

Amen.

當有幸我們有一天遇見耶穌,耶穌會對你說什麼呢?

「願你平安!」

「願你成為和平之子,成為Peacekeeper 與Peacemaker」.



\newpage

\section{上帝愛世界}
\label{sec:370}
\textbf{蒲錦昌牧師 - 上帝愛世界}
\newline
\newline
\hyperref[sec:369]{< < < PREV SERMON < < <}
~
\hyperref[sec:toc]{[返目錄]}
~
\hyperref[sec:371]{> > > NEXT SERMON > > >}


創世記 1:1-2:4
\newline
\begin{longtable}{cl}
\hline
\hline
章節 & 經文 (和合本修訂版)\\
\hline
1:1 & \begin{tabularx}{0.7\textwidth}{X} 起初，神創造天地。 \end{tabularx} \\ \\ \relax
1:2 & \begin{tabularx}{0.7\textwidth}{X} 地是空虛混沌，深淵上面一片黑暗；神的靈運行在水面上。 \end{tabularx} \\ \\ \relax
1:3 & \begin{tabularx}{0.7\textwidth}{X} 神說：「要有光」，就有了光。 \end{tabularx} \\ \\ \relax
1:4 & \begin{tabularx}{0.7\textwidth}{X} 神看光是好的，於是神就把光和暗分開。 \end{tabularx} \\ \\ \relax
1:5 & \begin{tabularx}{0.7\textwidth}{X} 神稱光為「晝」，稱暗為「夜」。有晚上，有早晨，這是第一日。 \end{tabularx} \\ \\ \relax
1:6 & \begin{tabularx}{0.7\textwidth}{X} 神說：「眾水之間要有穹蒼，把水和水分開。」 \end{tabularx} \\ \\ \relax
1:7 & \begin{tabularx}{0.7\textwidth}{X} 神就造了穹蒼，把穹蒼以下的水和穹蒼以上的水分開。事就這樣成了。 \end{tabularx} \\ \\ \relax
1:8 & \begin{tabularx}{0.7\textwidth}{X} 神稱穹蒼為「天」。有晚上，有早晨，這是第二日。 \end{tabularx} \\ \\ \relax
1:9 & \begin{tabularx}{0.7\textwidth}{X} 神說：「天下面的水要聚集在一處，使乾地露出來。」事就這樣成了。 \end{tabularx} \\ \\ \relax
1:10 & \begin{tabularx}{0.7\textwidth}{X} 神稱乾地為「地」，稱聚集在一起的水為「海」。神看為好的。 \end{tabularx} \\ \\ \relax
1:11 & \begin{tabularx}{0.7\textwidth}{X} 神說：「地要長出植物，就是含種子的五穀菜蔬，和會結果子、果子裡有種子的樹，在地上各從其類。」事就這樣成了。 \end{tabularx} \\ \\ \relax
1:12 & \begin{tabularx}{0.7\textwidth}{X} 於是地長出了植物：含種子的五穀菜蔬，各從其類；會結果子、果子裡有種子的樹，各從其類。神看為好的。 \end{tabularx} \\ \\ \relax
1:13 & \begin{tabularx}{0.7\textwidth}{X} 有晚上，有早晨，這是第三日。 \end{tabularx} \\ \\ \relax
1:14 & \begin{tabularx}{0.7\textwidth}{X} 神說：「天上要有光體來分晝夜，讓它們作記號，定季節、日子、年份， \end{tabularx} \\ \\ \relax
1:15 & \begin{tabularx}{0.7\textwidth}{X} 它們要在天空發光，照在地上。」事就這樣成了。 \end{tabularx} \\ \\ \relax
1:16 & \begin{tabularx}{0.7\textwidth}{X} 於是神造了兩個大光體，大的管晝，小的管夜，又造了星辰。 \end{tabularx} \\ \\ \relax
1:17 & \begin{tabularx}{0.7\textwidth}{X} 神把這些光體擺列在天空，照在地上， \end{tabularx} \\ \\ \relax
1:18 & \begin{tabularx}{0.7\textwidth}{X} 管理晝夜，分別光暗。神看為好的。 \end{tabularx} \\ \\ \relax
1:19 & \begin{tabularx}{0.7\textwidth}{X} 有晚上，有早晨，這是第四日。 \end{tabularx} \\ \\ \relax
1:20 & \begin{tabularx}{0.7\textwidth}{X} 神說：「水要滋生眾多有生命之物；要有鳥飛在地面以上，天空之中。」 \end{tabularx} \\ \\ \relax
1:21 & \begin{tabularx}{0.7\textwidth}{X} 神就創造了大魚和在水裡滋生的各樣活動的生物，各從其類，以及各樣有翅膀的鳥，各從其類。神看為好的。 \end{tabularx} \\ \\ \relax
1:22 & \begin{tabularx}{0.7\textwidth}{X} 神就賜福給這一切，說：「要繁殖增多，充滿在海的水裡；飛鳥也要在地上增多。」 \end{tabularx} \\ \\ \relax
1:23 & \begin{tabularx}{0.7\textwidth}{X} 有晚上，有早晨，這是第五日。 \end{tabularx} \\ \\ \relax
1:24 & \begin{tabularx}{0.7\textwidth}{X} 神說：「地要生出有生命之物，各從其類，就是牲畜、爬行動物、地上的走獸，各從其類。」事就這樣成了。 \end{tabularx} \\ \\ \relax
1:25 & \begin{tabularx}{0.7\textwidth}{X} 於是神造了地上的走獸，各從其類；牲畜，各從其類；和地上一切的爬行動物，各從其類。神看為好的。 \end{tabularx} \\ \\ \relax
1:26 & \begin{tabularx}{0.7\textwidth}{X} 神說：「我們要照著我們的形像，按著我們的樣式造人，使他們管理海裡的魚、天空的鳥、地上的牲畜和全地，以及地上爬的一切爬行動物。」 \end{tabularx} \\ \\ \relax
1:27 & \begin{tabularx}{0.7\textwidth}{X} 神就照著他的形像創造人，照著神的形像創造他們；他創造了他們，有男有女。 \end{tabularx} \\ \\ \relax
1:28 & \begin{tabularx}{0.7\textwidth}{X} 神賜福給他們，神對他們說：「要生養眾多，遍滿這地，治理它；要管理海裡的魚、天空的鳥和地上各樣活動的生物。」 \end{tabularx} \\ \\ \relax
1:29 & \begin{tabularx}{0.7\textwidth}{X} 神說：「看哪，我把全地一切含種子的五穀菜蔬和一切會結果子、果子裡有種子的樹，都賜給你們；這些都可作食物。 \end{tabularx} \\ \\ \relax
1:30 & \begin{tabularx}{0.7\textwidth}{X} 至於地上一切的走獸、天空一切的飛鳥，並一切在地上爬行的，有生命的動物，我把綠色植物賜給牠們作食物。」事就這樣成了。 \end{tabularx} \\ \\ \relax
1:31 & \begin{tabularx}{0.7\textwidth}{X} 神看一切所造的，看哪，都非常好。有晚上，有早晨，這是第六日。 \end{tabularx} \\ \\ \relax
2:1 & \begin{tabularx}{0.7\textwidth}{X} 天和地，以及萬象都完成了。 \end{tabularx} \\ \\ \relax
2:2 & \begin{tabularx}{0.7\textwidth}{X} 到第七日，神已經完成了造物之工，就在第七日安息了，歇了他所做一切的工。 \end{tabularx} \\ \\ \relax
2:3 & \begin{tabularx}{0.7\textwidth}{X} 神賜福給第七日，將它分別為聖，因為在這日，神安息了，歇了他所做一切創造的工。 \end{tabularx} \\ \\ \relax
2:4 & \begin{tabularx}{0.7\textwidth}{X} 這就是天地創造的來歷。在耶和華神造地和天的時候， \end{tabularx} \\ \\
[1ex]
\hline
\hline
\end{longtable}
詩篇 8:1-9
\newline
\begin{longtable}{cl}
\hline
\hline
章節 & 經文 (和合本修訂版)\\
\hline
8:1 & \begin{tabularx}{0.7\textwidth}{X} 耶和華－我們的主啊，你的名在全地何其美！你將你的榮耀彰顯於天。 \end{tabularx} \\ \\ \relax
8:2 & \begin{tabularx}{0.7\textwidth}{X} 你因敵人的緣故，從孩童和吃奶的口中建立了能力，使仇敵和報仇的閉口無言。 \end{tabularx} \\ \\ \relax
8:3 & \begin{tabularx}{0.7\textwidth}{X} 我觀看你手指所造的天，並你所陳設的月亮星宿。 \end{tabularx} \\ \\ \relax
8:4 & \begin{tabularx}{0.7\textwidth}{X} 人算甚麼，你竟顧念他！世人算甚麼，你竟眷顧他！ \end{tabularx} \\ \\ \relax
8:5 & \begin{tabularx}{0.7\textwidth}{X} 你使他比神微小一點，賜他榮耀尊貴為冠冕。 \end{tabularx} \\ \\ \relax
8:6 & \begin{tabularx}{0.7\textwidth}{X} 你派他管理你手所造的，使萬物， \end{tabularx} \\ \\ \relax
8:7 & \begin{tabularx}{0.7\textwidth}{X} 就是一切的牛羊、田野的牲畜、空中的鳥、海裡的魚，凡游在水裡的， \end{tabularx} \\ \\ \relax
8:8 & \begin{tabularx}{0.7\textwidth}{X} 都服在他的腳下。 \end{tabularx} \\ \\ \relax
8:9 & \begin{tabularx}{0.7\textwidth}{X} 耶和華－我們的主啊，你的名在全地何其美！ \end{tabularx} \\ \\
[1ex]
\hline
\hline
\end{longtable}
哥林多後書 13:11-13
\newline
\begin{longtable}{cl}
\hline
\hline
章節 & 經文 (和合本修訂版)\\
\hline
13:11 & \begin{tabularx}{0.7\textwidth}{X} 末了，弟兄們，願你們喜樂。要追求完全；要接受鼓勵；要同心合意；要彼此和睦。如此，慈愛和平的神必與你們同在。 \end{tabularx} \\ \\ \relax
13:12 & \begin{tabularx}{0.7\textwidth}{X} 你們要用聖潔的吻彼此問安。 \end{tabularx} \\ \\ \relax
13:13 & \begin{tabularx}{0.7\textwidth}{X} 眾聖徒都向你們問安。 \end{tabularx} \\ \\
[1ex]
\hline
\hline
\end{longtable}
馬太福音 28:16-20
\newline
\begin{longtable}{cl}
\hline
\hline
章節 & 經文 (和合本修訂版)\\
\hline
28:16 & \begin{tabularx}{0.7\textwidth}{X} 十一個門徒往加利利去，到了耶穌指定他們去的山上。 \end{tabularx} \\ \\ \relax
28:17 & \begin{tabularx}{0.7\textwidth}{X} 他們見了耶穌就拜他，然而還有人疑惑。 \end{tabularx} \\ \\ \relax
28:18 & \begin{tabularx}{0.7\textwidth}{X} 耶穌進前來，對他們說：「天上地下所有的權柄都賜給我了。 \end{tabularx} \\ \\ \relax
28:19 & \begin{tabularx}{0.7\textwidth}{X} 所以，你們要去，使萬民作我的門徒，奉父、子、聖靈的名給他們施洗， \end{tabularx} \\ \\ \relax
28:20 & \begin{tabularx}{0.7\textwidth}{X} 凡我所吩咐你們的，都教導他們遵守。看哪，我天天與你們同在，直到世代的終結。」 \end{tabularx} \\ \\
[1ex]
\hline
\hline
\end{longtable}
講題:上帝愛世界　For God So Loved the World

經文:創世記1:1-2:4;詩篇8:1-9;哥林多後書13:11-13;馬太福音28:16-20

講員:蒲錦昌牧師

日期:2017年6月11日

一,引言:

各位弟兄姊妹,平安.

六月五日是世界環保日,今天經課所讀的創世記1:1-2:4上及詩8:1-9均與上主的創造有關. 因此,我選了「上帝愛世界」這個主題,並且選了聖經中的不同部分,包括舊約,新約(福音書,書信)來和大家分享.

二,主題:

人類中心主義

自從十七世紀啟蒙時代以來,人類對自己信心十足,並認為可以透過自己的智慧和創造力,藉科學和科技改造世界,成為理想的,豐盛出產的世界, 不用再忍受天氣的改變,猛獸的襲擊,飢渴的痛苦等,以滿足自己的需要.基督教雖然因為相信上主而沒有這種人類的狂妄,但卻受時代人類中心主義思想的影響,甚至在閱讀聖經時也戴上「人類中心主義」的有色眼鏡來看,結果認為上帝創造,基督救贖,聖靈更新和再造的信仰都是以人為目的.例如,我們會解釋創世記第一至二章上帝創造的故事為預備了日月星晨,飛禽走獸,花草樹木,鳥獸蟲魚才創造人類,讓人第六天出現,是為人類準備一切.換而言之,人類才是上主最大的創造目的.墮落的故事和耶穌基督的救贖則只關乎人類,福音是針對人的需要.是嗎?

有拉比解釋創世記中人是最後被造,因為上主希望人謙卑一些.的確,人類若沒有太陽的照耀,清新的空氣,乾淨的食水,菜蔬,果實和各種生物作食物,人類根本不能生存.我們能夠驕傲嗎?我們應該驕傲嗎?

無疑,人類在上主創造中佔有一個非常重要的位置,且具有上主的形象,在萬物中有與眾不同的特質.但是,若以此認為人類是上主創造的最大和唯 一的目的,那麼,就可能誤讀了經文原本不一定有的含義.

讓我們重新閱讀幾段熟悉的經文,看看上主在創造,救贖,更新/再造的心意.

上主的愛,榮耀與創造的美麗(創造)

上主是宇宙萬物的創造者,祂愛顧一切,並且在祂所創造的宇宙中反映了祂的榮美,我們和萬物的存在乃是要榮耀祂.

約3:16:「上帝愛世人,甚至將祂的獨生子賜給他們,叫一切信祂的,不至滅亡,反得永生.」這節經文,我們非常熟悉,但請再細心看看,中文翻譯上有錯誤,希臘文原文是Kosmos,就英文Cosmos的來源,中文「宇宙」的意思.換而言之,這節經文應譯為「上帝愛宇宙／世界,甚至將祂的獨生子賜給他們,叫一切信他的,不至滅亡,反得永生.」英文聖經便是For God so loved the world,聖經寫的明明是宇宙和整個世界,我們卻只看到「世人」.神的關心遠超我們的,祂關懷整個受造界,而耶穌基督的拯救  工作,也無疑有宇宙性的意涵.

詩8:1「耶和華──我們的主啊,你的名在全地何其美!你將你的榮耀彰顯於天.」詩篇有大量的經文,描述上主創造的美麗和所反映上主的榮耀.只要我們細心察看便了解,正如詩8:3-4:「我觀看你指頭所造的天,並你所陳設的月亮星宿,便說:『人算什麼,你竟顧念他!世人算什麼,你竟眷顧他!』」看見上主創造的天之高,月亮星宿在黑暗上空中閃亮,人便驚覺自己的渺小.

我們很多人都喜歡旅行,認識周圍的世界.就算沒時間旅行的,也會看電視感受一下湛藍的天際,各種各樣的地貌,名山大川的美麗,各樣生物及植物的奧妙.可惜,我們許多人平日均營營役役,在城市生活中,舉頭是在灰濛濛一片看不到天空和太陽,晚上也看不到星空,走在街上遇不上人類以外的其他生物,在家中除了撲滅蚊子殺小強以外沒有和什麼生物建立友善的關係,甚至怕大樹危害自己的安全而要伐除.城市不斷擴張,其他生物和植物不斷減少,被邊緣化.我們征服了世界,但同時生活變得孤獨和枯燥乏味,感受不到上主創造的偉大,看不見萬物的豐盛和生命力,也不再為上主所造的一切感到好奇和讚歎上主的奇妙作為!

上主拯救的作為(犯罪,墮落與拯救)

上主的創造不單是為了人類,上主的拯救工作同樣不只是救贖人類於敗壞之中.

在創9:8-17,大家非常熟悉的挪亞方舟,當上主因人的罪惡而要以洪水毀滅大地,卻讓挪亞拯救了他一家和各種生物及植物,並且以彩虹為記,與人類及一切活物立約,不再以洪水毀滅他們.創9:8-11說:「上帝曉諭挪亞和他的兒子說:『我與你們和你們的後裔立約,並與你們這裡的一切活物──就是飛鳥,牲畜,走獸,凡從方舟裡出來的活物──立約.我與你們立約,凡有血肉的,不再被洪水滅絕,也不再有洪水毀壞地了.』」

神的拯救不是只有人,也包括各樣受造物,祂立約的同樣也不只是人類,也包括大地及一切活物.上主的創造是多樣而豐盛,但人的罪惡卻敗壞了自己,也破壞了大地的美善.

先知何西阿在何4:1-3說:「以色列人哪,你們當聽耶和華的話.耶和華與這地的居民爭辯,因這地無誠實,無善良,無人認識上帝.但起假誓,不踐前言,殺害,偷盜,姦淫,行強暴,殺人流血,接連不斷.因此,這地悲哀,其上的民,田野的獸,空中的鳥必都衰微,海中的魚也必消     滅.」原來人的行為與罪惡不單影響自己,也禍及大地,海洋和其上的一切生物.

《環境與基督信仰的關係》,《關懷受造界》等書的作者加爾文‧德威特(Calvin B.DeWitt)在一篇題為「大地環境對基督教福音派的挑戰」的文章中,提到受造界的七種退化現象:

太陽與地球之間能量交換的改變:全球暖化加劇,並且保護地球的臭氧層受到破壞.

土地退化: 土壤侵蝕,鹽化作用與沙漠化不僅造成土地退化,也剝奪了許多生物的棲息地與作物耕地.

森林濫伐:原始森林以每年大約十萬平方公里──約等同於冰島的面積──速率不斷流失.而過度開發利用的結果,亦造成土地以與上述相同的速率不斷退化.

物種滅絕:平均每一天都有三種動植物從地表消失.

水質退化: 包含地下水,湖泊,河川,海洋的污染.

廢棄物的製造與全球毒化:大氣層與洋流中充斥了人類丟到空氣與水中的廢棄物.

人類與文化的墮落:人類社群長久以來與受造物相互依存的可持續生活,因人類與文化的墮落而面臨威脅與衰亡的命運,並導致食物與植物栽培長久以來所形成的多樣性大量消失. ( 羅伯特‧貝瑞編:《關顧受造世界──基督教對環境的關懷與行動》新北市:聖經資源中心,二零一五年五月初版,頁89至110)

因此,受造界一直哀嘆,等待上主拯救的作為.正如羅8:19-22說:「受造之物切望等候上帝的眾子顯出來.因為受造之物服在虛空之下,不是自己願意,乃是因那叫他如此的.但受造之物仍然指望脫離敗壞的轄制,得享上帝兒女自由的榮耀.我們知道,一切受造之物一同歎息,勞苦,直到如今.」

我們是否有為大地的破壞,物種的滅絕感到羞愧嗎?我們覺得人類所造的虧缺了上帝的榮耀嗎?我們也聽到大地和一切受造之物的哀嘆嗎?

4. 終末的遠象與人的角色(終末論)

西1:15-19說:「愛子是那不能看見的上帝的像,是首生的,在一切被造的以先.因為萬有都是靠他造的,無論是天上的,地上的,能看見的,不能看見的,或是有位的,主治的,執政的,掌權的;一概都是藉著他造的,又是為他造的.他在萬有之先;萬有也靠他而立.他也是教會全體之首.他是元始,是從死裡首先復生的,使他可以在凡事上居首位.因為父喜歡叫一切的豐盛在他裡面居住.」

可見,基督耶穌拯救的作為,人類終末的遠象,是基督耶穌在萬物上(天上的,地上的,能看見的,不能看見的)居首位,因為這也是原來上主的心意,基督參與了創造(「萬物都是靠他造的」),最終要萬物榮耀三一真神,以祂為首(「又是為他造的」).

如果人與上主同工,承擔起管家的角色,好好的管理世界,就會出現不同的一番景象.正如詩8:6-8所言:

「你派他管理你手所造的,使萬物,就是一切的牛羊,田野的獸,空中的鳥,海裡的魚,凡經行海道的,都服在他的腳下.耶和華──我們的主啊!你的名在全地何其美!」

有一次,我往大陸公幹,看見北京首都機場與聯合國環境規劃署聯合主辦的「七十億人的夢想:一個星球,關愛型消費」的環境宣傳活動,倡導環保意識和理念,展出的海報有如下內容:

全世界目前有70億人口,到2050年,世界人口將達到90億.這將為地球資源及生態系統的承載能力造成前所未有的壓力.

如果目前的消費和生產模式保持不變,隨著人口的持續增長,到2030年人類將需要兩個地球的資源才得以滿足需求.

我們可以轉變消費模式,在不影響生活質量的基礎上,大幅降低能源,材料的消耗強度.

世界淡水資源並非想像之多:世界的淡水資源(可飲用)只佔地球水資源總量的3\%,其中5\%被凍結在南極,北極和冰川.因此,人類只能依靠剩下的0.5\%來滿足淡水需求.

水資源看起來免費,其實是昂貴的自然資本.世界上的淡水資源並非我們想像之多.請不要浪費每一滴水.

人類污染水源的速度遠遠超過自然界河流和湖泊自我淨化的速度.

每年被浪費的糧食高達十三億噸.與此同時,世界上有逾八億人遭受飢餓之苦,另外八億人被診斷為營養不良.

損失和浪費的糧食的生產佔用了14億公頃的土地,接近全球農業土地的30\%.

今晚剩下的烤雞可用來做明天的三明治.麵包丁可用來做油炸麵包塊.要求餐館打包剩餘的食物,如果你不馬上食用請立即冷藏這些食物.很少有人將剩飯菜從餐館打包回家,但請不要不好意思這樣做!

創新綠色包裝每年可減少垃圾14億件,不必要的包裝與裝飾會帶來更多浪費.食物過度包裝問題突出,且使用著泡沫塑料等不可持續材料.在一些地方,人們正在重新採用更加環保的包裝物料.

如果們希望別人懂得讚歎上主的創造,感覺到「主的名在全地何其美!」那麼,我們便應保護大地,愛護其上一切的生命,並相信上主能藉萬事萬物顯出祂的榮耀,叫人認識上主的美善和能力.過簡樸的生活,不以消費和人自己的榮耀為人生最大的目的,親近大自然,減少破壞大自然,愛惜大地及其上的一切生命,善用資源,以感恩的心享受神所厚賜的一切,均是好管家的表現.

三,結語:

總括而言,上主是創造者,救贖者及更新者,祂藉具有祂形象的人類與祂同工,管理大地,彰顯祂的榮美.願我們也承擔起管家的角色,向上主負 責,使大地再現生機,人類有更美好的未來.阿們!



\newpage

\section{因信稱義}
\label{sec:371}
\textbf{邢福增院長 - 因信稱義}
\newline
\newline
\hyperref[sec:370]{< < < PREV SERMON < < <}
~
\hyperref[sec:toc]{[返目錄]}
~
\hyperref[sec:372]{> > > NEXT SERMON > > >}


羅馬書 5:1-8
\newline
\begin{longtable}{cl}
\hline
\hline
章節 & 經文 (和合本修訂版)\\
\hline
5:1 & \begin{tabularx}{0.7\textwidth}{X} 所以，我們既因信稱義，就藉著我們的主耶穌基督得以與神和好。 \end{tabularx} \\ \\ \relax
5:2 & \begin{tabularx}{0.7\textwidth}{X} 我們又藉著他，因信得以進入現在所站立的這恩典中，並且歡歡喜喜盼望神的榮耀。 \end{tabularx} \\ \\ \relax
5:3 & \begin{tabularx}{0.7\textwidth}{X} 不但如此，就是在患難中也是歡歡喜喜的，因為知道患難生忍耐， \end{tabularx} \\ \\ \relax
5:4 & \begin{tabularx}{0.7\textwidth}{X} 忍耐生老練，老練生盼望， \end{tabularx} \\ \\ \relax
5:5 & \begin{tabularx}{0.7\textwidth}{X} 盼望不至於落空，因為神的愛，已藉著所賜給我們的聖靈，澆灌在我們心裡。 \end{tabularx} \\ \\ \relax
5:6 & \begin{tabularx}{0.7\textwidth}{X} 我們還軟弱的時候，基督就在特定的時刻為不敬虔之人死。 \end{tabularx} \\ \\ \relax
5:7 & \begin{tabularx}{0.7\textwidth}{X} 為義人死，是少有的；為仁人死，或者有敢做的。 \end{tabularx} \\ \\ \relax
5:8 & \begin{tabularx}{0.7\textwidth}{X} 惟有基督在我們還作罪人的時候為我們死，神的愛就在此向我們顯明了。 \end{tabularx} \\ \\
[1ex]
\hline
\hline
\end{longtable}
講題:因信稱義　Justification by Faith

經課:羅馬書5:1-8

講員:邢福增院長

場合:香港中文大學崇基學院禮拜堂

日期:2017年6月18日

難為正邪

在我讀小學的年代,粵語片的橋段,總是邪不能勝正.例如黃飛鴻系列,雖然壞人(如石堅)好像得逞,但最後黃飛鴻一定可以儆惡懲奸的.不過,近年不少電影或電視創作,對所謂正邪或善惡的描繪,大多採模糊手法.也許大家都看過港產片《無間道》,電影描寫警福音派到黑社會的臥底梁朝偉,以及黑社會派入警隊的臥底劉德華兩個身分混亂的男人的故事.劉德華最後千方百計想作一名真的警察,於是把所有知悉他黑社會身分的人都殺死,而一直希望擺脫臥底生涯的梁朝偉,最終也無法如願,甚至被殺害.電影的結局顛覆了邪不能勝正的規律,在中國大陸上映時,這個結局不獲通過,最後換了一個梁朝偉得到平反,劉德華惡有惡報的結局.作為觀眾,不知你喜愛那一個結果?

對,人世間存在著善與惡的力量,一方面,我們相信正義最終必能得勝,但另方面,在過程中,我們又同時感受到事情不一定是那麼黑白分明,難為正邪定分界,正中又邪,邪中又有正,又確是現實的寫照.基督教信仰對人性,也有這樣兩重的理解.一方面,人有上帝的形象,上帝把管家的責任交託我們,治理大地；另方面,人又選擇遠離上帝,不能妥善地作好管家的職責.人類的自我中心,成為歷史上各種紛爭的根源,為其他人,甚至我們身處的世界與大自然,帶來許多的禍害.人就是這樣的複雜,中國偉大的思想家之一的孟子,他有一句名言,是「人之所以異於禽獸者幾希也」,即是說,人與禽獸之間的差別,其實是很細微的.但就是這些細微分別,確立了人之所以為人的本性.孟子的意思是,如果人不努力,就很容易變成禽獸.難怪佛教說:一念天堂,一念地獄.就是一念之間,我們所作的決定,便可以有天淵之別.

藉信赦罪,因信稱義

基督教信仰把這種情況,總結為「罪」.罪是關係性的,指人與創造主之間的關係破裂了；罪又是實存性的,指存在人心裡的罪性,以及我們身處的罪惡的環境.擁有權力與地位的人,如果不懂制約,確是十分可怕的.同時,即使沒有權力的人,也受到罪性的約束,而對別人帶來傷害.

基督教信仰對於罪惡問題的回應,首先是提出重建人與上帝間斷絕的關係,這種復和的關係是通過耶穌基督的死來完成的.當我們接受主耶穌作救主,便代表我們可以重建與上帝間的關係.不過,這只能說解決了剛才我們所指的關係性的罪,事實上,接受了主耶穌的人仍然要面對實存性的罪,因為我們仍然活在這個罪惡的世界,我們的內心仍會受到罪性的誘惑,甚至被罪所轄制.那麼,到底福音對我們有甚麼意思?

剛才我們讀的《羅馬書》經課,對我們思考這問題,可以帶來一些啟迪.《羅馬書》的其中一個中心思想是藉信得救.保羅在第一章中已經陳明了「罪」的真實性.他又說:「罪的工價乃是死」(六23).然後,他便帶出了「因信稱義」的思想.「所以,我們既因信稱義,就藉著我們的主耶穌基督得以與上帝和好」(五1)保羅在第五章開始,劈頭便指出人類的出路端在「與上帝和好」,這是針對人與上帝間因罪惡而帶來的關係破裂.罪人如何與上帝和好呢?答案就是「藉著我們的主耶穌基督」.他在第四章結尾時說:「耶穌被出賣,是為我們的過犯；他復活,是為使我們稱義.」(四25),正好將「因信稱義」的思想點出.基督被釘死在十架上,是「為我們的過犯」,當他第三天復活,勝過死亡的權柄,就是為那些「信」他的人帶來新生命的可能,「使我們稱義」.

靈性掙扎

2017年是馬丁路德宗教改革五百年紀念的大日子.他在1517年,在威登堡(Wittenberg)的城堡教堂門上,釘上「九十五條議文」,表達了對當時的羅馬教會的不滿及改革的訴求,掀起了改教運動的序幕.後人在整理馬丁路德這位改革家的思想時,一定會提及「因信稱義」這教導.他原來在天主教的修道院學習,對於自己的表現,他說:「我是一個很優秀的修道士,我嚴格的遵守修道會的一切規則,我嚴謹到一個地步,如果有任何一個修道士,因為修行的緣故而得以進入天堂,我敢說那就是我了,所有認識我的修道院的弟兄們都可以為我做見證.」看,馬丁路德是多麼的自信,如果有信徒跟你這樣說,那你很大可能會認為,「弟兄,你太自大了」.

不過,馬丁路德還是經歷到罪的真實與人性的軟弱與掙扎.他每一次都是早上一起來就開始認罪禱告,禱告到睡著了,再起來認罪.但當他誠實面對自己時,才發現自己原來沒法確信能夠擁有救恩.他說:「我從前雖然是一個很虔誠的修道士,我的生活無可指摘,我卻常常認為自己是一個大罪人,我在上帝面前坐立不安,我的良心不住的自我控告,我無法相信我所做的一切補罪的功德,可以取悅於上帝,討上帝的喜悅.事實上,我對上帝很生氣,我恨這一位上帝.我恨這一位審判人的公義之神,可憐的罪人,不但在原罪的咒詛之下永無翻身之日,又被舊約的律法壓的透不過氣來.如今,上帝又藉著福音讓我愁上加愁.」信仰的掙扎令他陷入黑暗與深淵之中……

直到他在1518年再讀《羅馬書》.「因為上帝的義正在這福音上顯明出來；這義是本於信,以至於信.如經上所記:『義人必因信得生』.」(一17)他說:「我被自己鞭策得幾乎發瘋,我晝夜不斷的思想保羅在羅馬書這裡的意思,最後我終於澈悟了.義人得生是因著信.就是因著上帝的恩典我們被稱為義.我們是完全被動的領受恩典,我就好像突然重生.」當路德在講這句話的時候,他已經在神學院教書了,他已經拿到神學博士了.路德說:「天國之門就立刻大開,我終於進入樂園了.」這是他釘上九十五條的第二年.

站立在恩典之中

「所以,我們既因信稱義,就藉著我們的主耶穌基督得以與上帝和好」.保羅再一次確認,藉著主耶穌基督,因信被稱為義,罪人得以與上帝和好.當我們聽到「因信稱義」時,往往將之視作艱深的神學教義.對,這是一個很重要的神學理念,但同時,這也是基督徒生命可以經歷到的福氣.

保羅說,「與上帝和好」的結果,就是「我們又藉著他,因信得以進入現在所站立的這恩典中,並且歡歡喜喜盼望上帝的榮耀」(五2)「進入」是當下的經驗,保羅用「站在」在「恩典中」來形容這經驗.聖經學者指.「站在」是指在壓力下的站穩,也就說,這恩典需要站穩,牢固與上主維續著密切的關係.這樣,我們就可以「歡歡喜喜盼望上帝的榮耀」.「盼望」也是建基於對上主的信心,最終罪人可以得著上帝的榮耀.

患難中的盼望

當保羅說「站在」是一種壓力下的站穩,也意味著基督徒的生命會面對不同的考驗.既然如此,又如何能歡歡喜喜地盼望?所以,保羅說:「在患難中也是歡歡喜喜的,因為知道患難生忍耐,忍耐生老練,老練生盼望」(五3-4).「患難」是外在的環境,很可能指向因信仰而承受的逼迫.面對逼迫,保羅的教導是「也是歡歡喜喜的」,似乎很荒謬.難道這是一種教人脫離現實的喜樂嗎?不是,保羅說:「因為知道患難生忍耐」,「忍耐」是「堅忍」,是一種剛毅不屈的精神.在患難中,我們更要「生」(培養)一種堅忍鬥志.不是自怨自艾,怨天尤人,甚至絕望,而是「堅忍」.堅忍的基礎不是自己,而是對上主的信心.「忍耐生老練」,「老練」的一種身經百戰的氣質,這種生命的氣質是在堅忍之中生成的.「老練生盼望」,意思是基督徒的堅忍與毅力所建立的品格,是我們盼望的唯一根基.保羅用一種鎖鏈式的結構來提醒信徒:只有盼望才能面對患難.

人生在世,我們總會經歷不同的患難,並且真實地考驗我們的信仰.不管是政治壓力的逼迫,或是生命裡無法解釋,又可能擺脫的苦.有人經歷患難,整個人徹底被打跨了,也有人卻在能夠展現出信仰無比的力量,在患難中仍有盼望.為甚麼會這樣?

抉擇作怎樣的人

記得五年前的安息年,在內地檔案館查資料,看到一份令我極心寒的檔案!1966年4月文革前夕,兩位青年基督教人士(一是牧師,一是傳道),向黨打報告,匯報了教會內十位青年信徒的「反動」思想及行為!我不敢想像這些被報告的青年人,在未幾爆發的文革中,因這些報告會有甚麼遭遇!我後來把這感受貼在面書,收到不少回應.簡言之:(1)「他們是罪人」.這當然沒有錯.但卻令我想起彼得；(2)「我們都是罪人」,這更沒有錯.但這會否變成五十步笑百步?其實,每個罪人同時也是「被罪者」.我不會說他們沒有選擇,或又者將一切歸咎於社會的錯,但當他們選擇這樣作時,只能簡單地約化成個人的罪嗎?個人與社會／制度間的複雜關係,其實遠超過我們所想像.無論如何,那是個人性被嚴重扭曲的年代!換了是自己,又會如何?跟他們一樣?或是其他?坦白說,我不知道!(彼得不是曾對主說:「我就是必須和你同死,也絕不會不認你.」(太廿六35)嗎?)

記得有一篇文章,給我留下深刻的印象:「在最困難的時候做的選擇,才決定了這個人是什麼樣的人.」作者說:「平常在喝咖啡,在聊天,舒服的躺在床上時,都不會決定你是怎樣的人,當好不容易你的人生混了這麼久了,終於達到一個關鍵時刻,壓力很大,非常兩難,那個時候,你做了什麼決定就會證明你是怎樣的一個人,你會壯烈成仁,為國犧牲,或者變成漢奸,就在這一刻的抉擇.」原來在困難中的選擇,不僅能讓我們認識自己,也決定了自己是誰.

上主的愛

保羅又說:「盼望不至於落空,因為上帝的愛,已藉著所賜給我們的聖靈,澆灌在我們心裡.」原來,讓我們得以持守堅毅,抓緊盼望的,是因為「上帝的愛」.上帝的愛,一方面藉著聖靈,「澆灌在我們心裡」.另方面,是顯明在基督的受苦與犧牲之上.

保羅充分體會到人的罪性.是的,與上帝和好,處理了關係性的罪,人與上帝間因罪而斷絕的關係得以復和,但基督徒仍要真實面對實存性的罪.正如剛才提及馬丁路德的經歷.那麼,1518年後馬丁路德,是否就不再經歷罪的掙扎?不,馬丁路德承認自己仍常受試探.有一次他的信心受到考驗,自己的軟弱和污穢像電影一般,一幕一幕浮現眼前,他憤然提筆在書桌上寫著:「我已受洗」,意思是:他已經與基督一同埋葬,又與祂一同復活；只要生命建立在這個基礎上,就無人可以指控我,包括我自己在內.

「我們還軟弱的時候,基督就在特定的時刻為不敬虔之人死.為義人死,是少有的；為仁人死,或者有敢做的.惟有基督在我們還作罪人的時候為我們死,上帝的愛就在此向我們顯明了.」(五6-8)

我的小兒子很喜愛小寵物,多年前表示想飼養倉鼠.於是我們用港幣20元,買了兩隻回家.後來,有一隻生病了,知道獸醫收費昂貴,但不忍心就此放棄,結果帶去給獸醫.看了兩次,合共費用是千元.當時我心想,一千元,可以買一百隻回家了.如果這隻倉鼠「懂性」,知道我們為它所付出的,跟它本身的價值間的差異時,會否很感動呢?當我這樣想時,馬上說起基督耶穌在十架上為世人所作的,豈不比我為這倉鼠所作的更多嗎?正如保羅說:「為義人死,是少有的；為仁人死,或者有敢做的.惟有基督在我們還作罪人的時候為我們死,上帝的愛就在此向我們顯明了」.面對上主這捨己的愛,我們能不感動嗎?又可以如何回應?

稱義與成聖之路

「因信稱義」不是「因信稱義」,代表我們已經成為「義人」.我們只是因著基督的關係,得以被上帝「稱」為義人.我們都是不配的,唯有恩典才能把我們從「罪」中拯救出來,成為蒙恩的罪人.

今年是我大學畢業30年,也是受洗30年.翻查教會的洗禮禮文,原來領洗時牧師曾問我四個問題:(1)現在我問你願意真心悔罪改過,除去一切邪惡麼?(2)你是否願意誠心接受耶穌基督為救主?(3)你願意終身順服上帝的旨意,遵從上帝的誡命麼?(4)我們聖教會的道理都記在聖經上,你願意勤讀聖經,追求真理麼?只要我們誠實面對自己,便會發覺自己距離目標仍然很遠.捫心自問,內心的邪念罪性仍是揚之不去,善惡兩個律恆常在內心交戰.我們是否真的可以終身順服上帝的旨意,遵從上帝的誡命?從現實人性的角度來看,想起自己所說的「我願意」,難免感到慚愧萬分.

俄國著名基督徒作家托爾斯泰曾說:「我們意識到自己不可能達到理想的完美.我們無法看見自己靠近這完美的程度；我們能夠看見的就只是自己偏離的程度」,確是實況.那麼,基督徒的出路在那裡?答案就是恩典.唯有靠賴基督赦罪恩典,我們才能面對自己,在「稱義」後再次在「成聖」路上奔跑.雖然我們仍然活在這個罪惡的世代,仍會受到罪惡的轄制,仍然會有軟弱的時候,傷害了別人,也讓自己受到傷害.但當我們認清自己的身分,就是上帝的子女,基督的門徒時,我們的生命已經被主的寶血救贖時,我們可以依靠上主的力量,去克勝罪惡,去面對自我.我們可以帶著這樣的盼望,我們可以持守這樣的信心,因為這是創造天地的主,救贖我們的基督,以及保惠師聖靈對我們的應許及幫助.

今天是父親節,盼望每位父親都能夠從天上的父親那裡,學習這份「為父之愛」.



\newpage

\section{首屆地球峰會}
\label{sec:372}
\textbf{王福義長老 - 首屆地球峰會}
\newline
\newline
\hyperref[sec:371]{< < < PREV SERMON < < <}
~
\hyperref[sec:toc]{[返目錄]}
~
\hyperref[sec:373]{> > > NEXT SERMON > > >}


創世記 9:8-17
\newline
\begin{longtable}{cl}
\hline
\hline
章節 & 經文 (和合本修訂版)\\
\hline
9:8 & \begin{tabularx}{0.7\textwidth}{X} 神對挪亞和同他一起的兒子說： \end{tabularx} \\ \\ \relax
9:9 & \begin{tabularx}{0.7\textwidth}{X} 「看哪，我要與你們和你們後裔立我的約， \end{tabularx} \\ \\ \relax
9:10 & \begin{tabularx}{0.7\textwidth}{X} 包括和你們一起所有的生物，就是飛鳥、牲畜、地上一切的走獸，凡從方舟裡出來地上一切的生物。 \end{tabularx} \\ \\ \relax
9:11 & \begin{tabularx}{0.7\textwidth}{X} 我與你們立我的約：凡有血肉的，不再被洪水滅絕，也不再有洪水毀壞這地了。」 \end{tabularx} \\ \\ \relax
9:12 & \begin{tabularx}{0.7\textwidth}{X} 神說：「這是我與你們，以及和你們一起的一切生物所立之永約的記號，直到萬代： \end{tabularx} \\ \\ \relax
9:13 & \begin{tabularx}{0.7\textwidth}{X} 我把彩虹放在雲中，這就是我與地立約的記號了。 \end{tabularx} \\ \\ \relax
9:14 & \begin{tabularx}{0.7\textwidth}{X} 我使雲遮地的時候，會有彩虹出現在雲中， \end{tabularx} \\ \\ \relax
9:15 & \begin{tabularx}{0.7\textwidth}{X} 我就記念我與你們，以及各樣有血肉的生物所立的約：不再有洪水氾濫去毀滅一切有血肉的了。 \end{tabularx} \\ \\ \relax
9:16 & \begin{tabularx}{0.7\textwidth}{X} 彩虹出現在雲中，我看見了，就要記念神與地上一切有血肉的生物所立的永約。」 \end{tabularx} \\ \\ \relax
9:17 & \begin{tabularx}{0.7\textwidth}{X} 神對挪亞說：「這就是我與地上一切有血肉的立約的記號。」 \end{tabularx} \\ \\
[1ex]
\hline
\hline
\end{longtable}
講題:首屆地球峰會　The First Earth Summit

經文:創世記9:8-17

講員:王福義長老

場合:香港中文大學崇基學院禮拜堂

日期:2017年6月25日

一,引言 世界環境日

聯合國大會1972年通過成立聯合國環境署UNEP並在斯德哥爾摩 (Stockholm) 召開人類環境會議Conference on Human Environment,決定6月5日為「世界環境日」.其目的是鼓勵全球民眾環境保護意識並採取環保行動.

由1974年開始世界各地都舉辦活動.

香港的教會最近幾年才有限地響應有關活動和講道,但不少教會還不知道有環境日,也少關心和人與環境的關係,這不只是與物質世界的關係,而是與上帝的關係.

二,地球峰會

1972年後20年來聯合國發現地球生物,生態,生境遭到極大的人為破壞,情況持續惡化；於1992年再召開更大型的會議.由聯合國環境與發展會議 (UNDEP), (是UNEP,UNDP聯合的大會 ) ,在6月3-14在巴西里約熱內盧Rio Summit召開世界地球峰會.172個國家出席,其中有116位國家元首.如果全球有196的國家的話,共有88\%出席.共有2,400非政府組織參加(NGO) .

會議的原因,是囚人類的命運受到環境破壞而面臨險境,包括:

1. 地球暖化─大量使用化石燃料,交通氣體排放,空氣污染

2. 生物絕種,因生境遭破壞

3. 森林的過度砍伐和

4. 過度開發,都市化和污染日增.

會議的成果:建立兩個國際公約 Convention ( 約Covenant)

生物多樣性公約 (Convention of Biological Diversity CBD)

聯合國氣候變化框架公約 (UN Framework Convention on climate Change UNFCCC)

這兩個協約都影響中國和香港

後者發展成為京都議定書Kyoto Protocol (1997).192國家簽署,2005年生效,限制(二氧化碳) Green House Gas GHG的排放.引至 (2015) 巴黎氣候協議

這次峰會的特點是:

1. 由世界組織召開

2. 由各國領袖出席, 有關專家,有權人士

3. 是面對危機的緊急時刻

4. 建立「公約」大家遵守

5. 有標誌,記號以防忘記

其實世界上第一屆的地球峰會不是1992年的巴西Rio Summit而是遠古之年,由上帝親自召開的地球峰會.

三,首屆地球峰會 ( 經文 創9:8-17 )

A. 召集人

由上帝,創造地球與天地的主親自召開 (創9:8 )

祂要與與會者立約

對象包括

(a) 你們( 挪亞一家 )人類代表( 9:9 )

(b) 你們的後裔 ─下一代,未出世的將來的人( 9:9 )

(c) 這裡的一切活物 ─方舟中保留下來的 ( 9:10 ) ；

飛鳥,牲畜,走獸 ( 9:10 ) ；；各樣有血肉的活物 (9:11-12; 15) .

(d) 與地立約( 9:13 ) .

B. 對象 - 出席者

人類,生物,各樣活物,大地,1992年地球峰會

1. 沒有神參與

2. 沒有生物參與

3. 沒有大地參與

4. 沒有將來的人參與

5. 只有政府,NGO和有能力出席的人參與

C. 「約」的內容和性質 ( 9:11 )

這段經文內出現七次 「約」字,顯示重要!

(a) 單方面的由上帝承諾,保證；

(b) 不再以洪水毀壞地；

(c) 人不須承擔任何責任和工作；

(d) 是普世的,關乎全人類和地球；

(e) 是永久性的( 9:12；16 ) 沒有時限的；

(f) 是有記號的( 9:13-14；16 ) ；

(g) 上帝會記念的,放在心中,不會吃言.

這是聖經中第一個上帝與人所立的約；是造物者與被造者之間的約；

是至高無上的,是首個也是關乎普世的.

在舊約還有三個約,分別是上帝與亞伯拉罕,摩西,大衛所訂的約

但性質不同:

(a) 亞伯拉罕( 創17:1-14 )成為選民,以割禮為記號

(b) 摩西之約 ( 出31:16-17 ) 以安息日為記號

(c) 與大衛之約( 撒下:1-17 ) 應允國位堅立

(d) 新約,耶穌設立聖餐( 路12:20 ) ,酒和餅,這杯是用我血所立的新約.

約是承諾,有約束力,有履行的責任和義務.

上帝採取主動促成這次首屆地球峰會,意義重大.

上帝承諾

(a)凡有血肉的不再被洪水滅絕；

(b.)也不再有洪水毀壞地( 9:11 ) .

問題:  上帝不用洪水,會不會用其他方法?

以後所多瑪和蛾摩拉,被天火所焚是局部的,但其理由相同,因為人犯罪.但是上帝不毀滅,人會不會自行毀滅?情況極為可能.

D.約的記號

記號是憑據,是見證,是備忘錄.

虹,天虹,彩虹,天然的現象,已存在,彩虹 英文(Rainbow 雨弓) ,

上帝創造之工,完成,清理之工完畢,休息了,把弓掛在天上為彩虹

是與大地立約的記號!

虹的特色→(a) 雨後,風暴後,陰沉雲層更顯爛爛；

(b) 天然,美麗；

(c) 易見─橫在天空；

(d) 沒有地域之限.

記號的重要性 ( 9:15-16 )

15 我(上帝)便記念；

16 看見就要記念；

這是永遠的約,備忘錄,提醒,記起；

上帝是可信的,決不食言；

上帝辦事公開,透明,信實,公義,

對人沒有要求；

結果,上帝沒再用洪水毀減世界,

但人類用不同手段,以發展為名摧毀了上帝所創造所愛護的地球.

四,首屆地球峰會召開的背景

首屆峰會是在洪水退去後召開,人和生物出了方舟,挪亞築壇獻祭,

結束了一年零10天的洪水淹地 (由2月17日(7:11)至次年二月27日(8:14)),挪亞由600歲到601歲.

耶和華說出心聲:“我不再因人的緣故咒詛地,(人從小時心裡懷著惡念),也不再按著我才行的,滅各種活物了.”(8:21)

大地恢復自然的秩序:地還存留的時候,稼穡,寒暑,冬夏,晝夜,就永不停息了(8:22) .

為甚麼耶和華要用洪水毀滅大地和生物呢?

因為人的罪

6:5  耶和華看見人在地上罪惡很大,終日所思想的盡都是惡

6:11 世界在神面前敗壞,地上滿了強暴 corrupted full of violence

6:12 神觀看世界,見是敗壞了；凡有血氣的人在地上都敗壞了行為

6:13 因為地上滿了他們的強暴.

上帝創造了人地和其中的生物,祂關心,愛護,並非創造完成後不理,而是時常查看,視察.見到情況令祂失望:

(1) 地上罪惡很大；

(2) 世界在神面前敗壞；與神的創造原意相反；

(神是有期望的,有公義的,容不下罪惡 ) ；

(3)人的思想敗壞；

(4)終日所思盡都是惡 - 不停止的,習以為常的,不是偶然的；

(5)有血氣的人在地上敗壞了行為；不但思想還有行動；

(6)地上滿了強暴( 6:11；13 ) ；侵害和欺壓別人和生物.

如該隱殺亞伯 － 無法無天,堆泥,斬樹,破壞生態,強橫無理,喜歡就去做.如所多瑪,蛾摩拉城邑的人,他們一味的行淫,隨從逆性情慾.滿眼是淫色,止不住犯罪.( 彼後2:14 )

上帝的反應:

(a)後悔造人在地上,心中憂傷 ( 6:6 ) ；因為我造他們後悔了 ( 6:7 )

後悔,失望,不及期望,遠離原意,有負所托.

(b)說要將所造的人和走獸,並昆蟲,以及空中的飛鳥,都從地上除滅.

(c)因為地上的強暴,凡有血氣的人；他們的盡頭已經來到我面前

(d)我要把他們和地一併毀滅 ( 6:13 ) .

世界敗壞,充滿暴力,惡人橫行,有違聖意；令造物者上帝失望,後悔,本以為可以遍滿地面,由人代為管理,雖知敗壞無遺.達不到創造的目的,只好拉倒重來,重新洗牌,以洪水清洗一切生物與人類.

人犯罪為何旁及大地,生物?

(a) 始祖犯罪吃了禁果,結果地長出荊棘和疾藜,地必為你的緣故受咒詛,你必終身勞苦 ( 創3:17 ) ；

(b)該隱殺亞伯後,地開了口,現在你必從這地受咒詛,你種地,地不再給你效力,你必流離飄蕩在地上 ( 創4:11-12 ) ；

(c)  所多瑪和蛾摩拉城中人犯罪,耶和華降火,“把那些城和全平原,並城裡所有的居民,連地上生長的都毀滅了.( 創19:25 ) ；

(d) 人犯罪連累其他生物被毀

“我們知道一切受造之物,一同歎息,勞苦,直到如今”( 羅8:22 )

昔日如是,如今也是,現在的情況與洪水前差不多.

我們強暴,貪心,自私,偽善,欺壓,不法,不但對人,也對生物生境,

濫捕魚,破壞珊瑚,毀濕地,砍樹林,殺大象,斬犀牛,捕捉鯊魚,鯨魚

砍香樹,羅漢松…………….不一而足.

是否也會遭報呢?上帝昔日看不過眼,人先敗壞上帝才「毀滅」,

現在有自救之道嗎?

五,救贖之道

希望來自挪亞

「惟有挪亞在耶和華眼前蒙恩」( 創6:8 )

他有何特質可以蒙恩?他的特點是 ( 6:9 )

1.是個義人  不是惡人  (6:5)；不行強暴 ( 6:11) ；恩慈待人待物,合乎 法律和道德標準；

2.在當時的世代是個完全人

健全,遠離罪惡,獨特:

a.與後來的亞伯拉罕一樣 ( 創17:1 )

“你當在我面前作完全人”全心全意的奉獻和信靠；

b.約伯也有同樣特點:“完全正直,敬畏神,遠離惡事”( 伯1:1 )

c. 耶穌說“你們要完全,像你們的天父一樣”( 太5:48)

3.與神同行─如以諾

與神同行300年 ( 創5:22 ) ,明白神的思想和看法,因為挪亞的蒙恩,  蒙神指示建造方舟.與他立約,你同你的妻,與兒子,兒婦,都要入方舟( 6:18 )

4. 順服─凡神所吩咐的,他都照行了 ( 6:22 ) ,

當地上有血肉,有氣息的活物,無一不死之時 ( 6:17 ) ,神藉義人挪亞施  施施行拯救

施行拯救,因在這世代中,我見你在我面前是義人 ( 7:1 )

在當日世界敗壞時有挪亞,在今日世界同樣敗壞時,有誰?

挪亞尊守神吩咐可為生物留種 ( 7:3 ) ,但人類 － 特別是我們中國人,把留下來的都吃了,殺了,賣了,滅了,把生境都發展了!

六 結論  反思

環境保護的基礎

神身為創造者,對環境有關心,查看,關注,顧念,期望,有正義

公義的存在；

人的罪行是環境破壞的主因；

不義的事影響大地和環境；

神因人的罪而毀滅大地；

也因一個義人而施行拯救；

也因愛護而與人,地,生物立約；

祂記念大地,人,生物,但人的表現令神失望,生物正步向滅種.

今天我們可以做甚麼?

看看世界真像,是否惡行,強暴,橫行；

了解人如何破壞世界生態；

在保護方面有不少可行的具體方法,主要的是效法挪亞,成為義人

在上帝眼中的完全人,與神同行─了解神的心意；順服─強過成功；

人與大地生物同是受迼者.

人被神委託為管理者,把大地,生物交付我們的手 (創9:2 ) ；因此我們有加倍的責任做好管家的職份.



\newpage

\section{忠誠的異見者}
\label{sec:373}
\textbf{林豪恩先生 - 忠誠的異見者}
\newline
\newline
\hyperref[sec:372]{< < < PREV SERMON < < <}
~
\hyperref[sec:toc]{[返目錄]}
~
\hyperref[sec:374]{> > > NEXT SERMON > > >}


耶利米書 28:5-9
\newline
\begin{longtable}{cl}
\hline
\hline
章節 & 經文 (和合本修訂版)\\
\hline
28:5 & \begin{tabularx}{0.7\textwidth}{X} 耶利米先知當著祭司和站在耶和華殿裡眾百姓的面，對哈拿尼雅先知說： \end{tabularx} \\ \\ \relax
28:6 & \begin{tabularx}{0.7\textwidth}{X} 「阿們！願耶和華如此行，願耶和華實現你所預言的話，將耶和華殿中的器皿和所有被擄去的人從巴比倫帶回此地。 \end{tabularx} \\ \\ \relax
28:7 & \begin{tabularx}{0.7\textwidth}{X} 然而我在你和眾百姓耳中所要說的話，你應當聽。 \end{tabularx} \\ \\ \relax
28:8 & \begin{tabularx}{0.7\textwidth}{X} 從古以來，在你我以前的眾先知，向多國和大邦說預言，論到戰爭、災禍、瘟疫的事。 \end{tabularx} \\ \\ \relax
28:9 & \begin{tabularx}{0.7\textwidth}{X} 至於那預言平安的先知，到先知的話應驗的時候，人就知道他真是耶和華所差來的。」 \end{tabularx} \\ \\
[1ex]
\hline
\hline
\end{longtable}
講題:忠誠的異見者　The Loyal Dissident

經文:耶利米書28:5-9

場合:香港中文大學崇基學院禮拜堂

日期:2017年7月2日

一.引言

此時此地,「忠誠的異見者」這題目,或者令人產生眾多聯想.本講章的源起,是今天的福音經課馬太福音10:41其中一句:「人因為先知的名接待先知,必得先知所得的賞賜」,舊約經課提示我們從耶利米先知的例子來解讀「因為先知的名接待先知」,以及「得先知所得的賞賜」的意思.

人們以為先知是古代的特殊人物,然而,現在也有人自稱先知.教會更新運動在今天下午舉辦講座,宣傳稿如此說:「是次講座將介紹一些以「先知」,「使徒」為名,活躍於各地華人社區的問題人物,他們到處騙取信徒及堂會的信任,濫權瀆職,斂財騙色,導致各種傷害的情況出現」.

先知是怎樣的人物呢?耶利米可算是其中一位先知的典範,讓我們再細讀今天的舊約經課.

二．經文

28:5先知耶利米當著祭司和站在耶和華殿裡的眾民對先知哈拿尼雅說:

28:6「阿們!願耶和華如此行,願耶和華成就你所預言的話,將耶和華殿中的器皿和一切被擄去的人從巴比倫帶回此地.

28:7然而我向你和眾民耳中所要說的話,你應當聽.

28:8從古以來,在你我以前的先知,向多國和大邦說預言,論到爭戰,災禍,瘟疫的事.

28:9先知預言的平安,到話語成就的時候,人便知道他真是耶和華所差來的.」

三.釋經與應用

1．政治不正確的先知

經文記載的事件發生在猶大王國的晚年,國家處於內憂外患,風雨飄搖之中,朝野上下人人自危,充滿焦慮.如何評估當下的局勢?應該採取甚麼行動?如何面對未來呢?對於這些逼切的問題,朝野內外意見分歧而且對立.經文中提及兩位有資格在朝廷發言的先知,他們應該是重要的意見領袖,分別是耶利米,以及哈拿尼雅.從短短幾節的敘述可見,他們對於當下局勢及回應行動各持不同的見解,甚至針鋒相對.到底他們的見解如何不同呢?他們在爭論甚麼呢?讓我們稍閱前文:

28:1當年,就是猶大王西底家登基第四年五月,基遍人押朔的兒子,先知哈拿尼雅,在耶和華的殿中當著祭司和眾民對我說:

28:2「萬軍之耶和華－以色列的上帝如此說:我已經折斷巴比倫王的軛.

28:3二年之內,我要將巴比倫王尼布甲尼撒從這地掠到巴比倫的器皿,就是耶和華殿中的一切器皿都帶回此地.

哈拿尼雅先知對前景的預測非常樂觀,他認為巴比倫兩年之內就會全方位崩潰,用當下網絡語言,即是「巴比倫爆炸」.因此,他認為不用臣服巴比倫,鼓動猶大乘機爭取獨立,脫離巴比倫的控制.

耶利米則沒有那麼樂觀,他認為當下的巴比倫不是衰敗,而是如日中天.因此,耶利米勸諫猶大王不要以卵擊石,免得聖靈塗炭:

27:12我就照這一切的話對猶大王西底家說:「要把你們的頸項放在巴比倫王的軛下,服事他和他的百姓,便得存活.

27:13你和你的百姓為何要因刀劍,饑荒,瘟疫死亡,正如耶和華論到不服事巴比倫王的那國說的話呢?

其實,當時以色列及猶大的大部領土已經失去,只餘耶路撒冷一城,苟延殘喘.亞述,巴比倫,埃及三個強國爭取世界霸權.巴比倫首先攻破亞述,後來在迦基米施戰役打敗埃及,搶奪了兩強的緩衝地.埃及退後防線,但伺機而動.兩強爭持,如猶大一樣處於狹縫中的小國尋找「漁人得利」的機會.二十七章記載,這些小國曾經為此召開跨國會議,試圖聯合起來爭取脫離巴比倫,尋求獨立.如此看來,哈拿尼雅的意見是有市場的.不過,耶利米對於那些小國的聯合謀略仍然「撥冷水」:

27:2「耶和華對我如此說:你做繩索與軛,加在自己的頸項上,

27:3藉那些來到耶路撒冷見猶大王西底家的使臣之手,把繩索與軛送到以東王,摩押王,亞捫王,泰爾王,西頓王那裡,

27:4且囑咐使臣,傳與他們的主人說,萬軍之耶和華－以色列的上帝如此說:

27:5我用大能和伸出來的膀臂,創造大地和地上的人民,牲畜.我看給誰相宜,就把地給誰.

27:6現在我將這些地都交給我僕人巴比倫王尼布甲尼撒的手,我也將田野的走獸給他使用.

27:7列國都必服事他和他的兒孫,直到他本國遭報的日期來到.

局勢將會如何?是耶利米預測正確,還是哈拿尼雅猜中,當下無人得知.兩位先知都奉耶和華之命發言,耶利米未得到明確啟示之前,也不能指責哈拿尼雅為假先知.他能夠作出的回應,就是我們今天經課誦讀的那幾句,大意是:我也希望你所宣告的平安果真實現.不過,你不要忘記,先知的傳統是提出警告,勸人悔改.

正如各小國的期望,哈拿尼雅對巴比倫兩年內崩潰的預測是受歡迎的,因為這預測迎合了猶大得到解脫的渴望,挽救猶大人的民族尊嚴,鼓舞猶大政權對抗外敵.相對來說,耶利米的言論是不受歡迎的,他勸告猶大臣服巴比倫,放棄爭取獨立,消滅猶大人的盼望,傷害猶大人的民族感情.因此,無論在朝在野,耶利米都遭受非議,甚至被逼害.哈拿尼雅也對他不客氣,拆掉他藉以宣講信息的道具,等於當眾侮辱他:

28:10於是,先知哈拿尼雅將先知耶利米頸項上的軛取下來,折斷了.

28:11哈拿尼雅又當著眾民說:「耶和華如此說:二年之內我必照樣從列國人的頸項上折斷巴比倫王尼布甲尼撒的軛.」於是先知耶利米就走了.

由於耶利米的信息政治不正確,朝野上下都想消滅他的聲音,甚至把他除之而後快:

20:1祭司音麥的兒子巴施戶珥作耶和華殿的總管,聽見耶利米預言這些事,

20:2他就打先知耶利米,用耶和華殿裡便雅憫高門內的枷,將他枷在那裡.

26:11祭司,先知對首領和眾民說:「這人是該死的；因為他說預言攻擊這城,正如你們親耳所聽見的.」

先知,令人想到以利亞的英雄事蹟,又令人想到以利沙的飄逸及超然；耶利米是怎樣的先知呢?何以他涉足國家事務,泥足深陷,以至捲入權力的鬥爭之中而不能抽身呢?作為先知,何以總是發表政治不正確的言論而令自己成為「過街老鼠」呢?何以預言針對自己的國家和人民呢?難道他不愛自己的國家嗎?難道他不關心自己的同胞嗎?

2．為國為民的先知

國家在危急存亡之秋,舉國上下都盼望平安的信息,耶利米為甚麼不配合一下國情呢?為甚麼要發出危言聳聽的預告呢?讓我們再檢視耶利米的言論,嘗試明白一下他的心情:

當巴比倫仍未大軍壓境之前,猶大國仍有足夠時間改革自強的日子,耶利米勸告國王及朝臣:

22:3耶和華如此說:你們要施行公平和公義,拯救被搶奪的脫離欺壓人的手,不可虧負寄居的和孤兒寡婦,不可以強暴待他們,在這地方也不可流無辜人的血.

22:4你們若認真行這事,就必有坐大衛寶座的君王和他的臣僕百姓,或坐車或騎馬,從這城的各門進入.

22:5你們若不聽這些話,耶和華說:我指著自己起誓,這城必變為荒場.』

正如耶利米說,「從古以來,先知向多國和大邦說預言,論到爭戰,災禍,瘟疫的事」,目的並不是顯示自己的預知能力,而是發出警告,勸諫君王及人民悔改.先知預告災禍的目的,就是希望災禍不要發生.

在另一個場合,耶利米對首領們和民眾說:

26:12「耶和華差遣我預言,攻擊這殿和這城,說你們所聽見的這一切話.

26:13現在要改正你們的行動作為,聽從耶和華－你們上帝的話,他就必後悔,不將所說的災禍降與你們.

26:14至於我,我在你們手中,你們眼看何為善,何為正,就那樣待我吧!

細聽耶利米這些言論,怎能說他不愛國呢?他提出警告,呼籲悔改,正是為了國家的前途,為了大衛王室的延續.耶利米對國家的關懷,置於個人的福祉之上,因此他說:「至於我,我在你們手中,你們眼看何為善,何為正,就那樣待我吧!」(26:14)不過,殺了先知,就等於先知所警告的事情不會發生嗎?

耶利米呼籲改革,日復一年,月復一月,無人理會,終於,巴比倫兵臨城下.主戰派決定戰鬥到底,耶利米呢?他竟然提議棄戰投降!一位心繫國家安危的先知,竟然勸告國王向敵人投降,他有顧及自己的形象及個人安危嗎?他在想甚麼呢?

27:12我就照這一切的話對猶大王西底家說:「要把你們的頸項放在巴比倫王的軛下,服事他和他的百姓,便得存活.

27:13你和你的百姓為何要因刀劍,饑荒,瘟疫死亡,正如耶和華論到不服事巴比倫王的那國說的話呢?

先知著緊的,的確不是自己的形象及性命,而是百姓的性命,是耶路撒冷的安危.棄戰投降,不會是將領們的選項,也不是王室及重臣的期盼,因為他們可以預見自己將會喪失權力,甚至被羞辱.先知卻把人民的性命置於政權的榮辱之上,因此,他力揀國王:「你和你的百姓為何要因刀劍,饑荒,瘟疫死亡呢?」(27:13)他冒死抗衡主戰派,公開地對說:「你們不可聽那先知對你們所說的預言.他們說:『耶和華殿中的器皿快要從巴比倫帶回來』；其實他們向你們說假預言.不可聽從他們,只管服事巴比倫王便得存活.這城何致變為荒場呢?」(27:16-17)

「百姓為何要因刀劍,饑荒,瘟疫死亡」呢?「這城何致變為荒場呢?」這是耶利米所著緊的.耶利米拒絕玉石俱焚,就可以指責他不愛國愛民嗎?然而,在內憂外患的焦躁之下,所有選項都偏向極端,還會有人理會耶利米的呼籲嗎?社會還能容得下耶利米的聲音嗎?那怕這呼籲是來自他們所敬拜的耶和華,那信息是來自掌管天地的上主!

3．忠誠的異見者的痛苦

相對於朝臣及宗教領袖的主流意見,耶利米是一位異見者,但他對國家及對人民的忠誠無容置疑.讀者可以想像,如此一位忠誠的異見者,必定是有苦自己知.

20:7耶和華啊,你曾勸導我,我也聽了你的勸導.你比我有力量,且勝了我.我終日成為笑話,人人都戲弄我.

20:8我每逢講論的時候,就發出哀聲,我喊叫說:有強暴和毀滅!因為耶和華的話終日成了我的凌辱,譏刺.

他甚至希望自己不存在於這個世界:

20:14願我生的那日受咒詛；願我母親產我的那日不蒙福!

20:18我為何出胎見勞碌愁苦,使我的年日因羞愧消滅呢?

然而,耶利米確實是一位先知,無論環境如何,無論自己遭遇甚麼對待,無論自己喜歡不喜歡,他仍是忠誠於上主,執行作為一位先知傳述上主對世界的說話的任務:

20:9我若說:我不再提耶和華,也不再奉他的名講論,我便心裡覺得似乎有燒著的火閉塞在我骨中,我就含忍不住,不能自禁.

在現實上,耶利米有機會改變他的命運.當巴比倫控制耶路撒冷後,尼布甲尼撒王認為耶利米有功,願意任其所求給予榮賞:

39:11巴比倫王尼布甲尼撒提到耶利米,囑咐護衛長尼布撒拉旦說:

39:12「你領他去,好好地看待他,切不可害他；他對你怎麼說,你就向他怎麼行.」(卅九12)

尼布甲尼撒王誤會了,耶利米並不是賣國求榮之輩,他不但拒絕封賞,更以對猶大的態度,對如日方中的巴比倫提出警戒:

51:41示沙克何竟被攻取?天下所稱讚的何竟被佔據?巴比倫在列國中何竟變為荒場?

51:44我必刑罰巴比倫的彼勒,使他吐出所吞的.萬民必不再流歸他那裡；巴比倫的城牆也必坍塌了.

耶利米作為耶和華的先知,不討好猶大王,不懼怕巴比倫王；他沒有為自己的政治前途打算,他只是忠誠地傳講耶和華吩咐他的說話.

四．總結

耶利米通常被稱作「流淚的先知」,他的淚為誰流呢?

9:1但願我的頭為水,我的眼為淚的泉源,我好為我百姓中被殺的人晝夜哭泣.

先知何以流淚呢?因為他是一位忠誠的異見者,他對上主忠誠,對人民忠誠,對國家忠誠.在某種歷史時空之下,真正的先知,原來是忠誠的異見者.

有人說,教會或多或少也承繼著先知的角色.耶利米的見證,但願能夠給此時此地的教會一點反省和啟迪.



\newpage

\section{惜別講章與神同行祝福人}
\label{sec:374}
\textbf{許開明牧師 - 惜別講章:與神同行祝福人}
\newline
\newline
\hyperref[sec:373]{< < < PREV SERMON < < <}
~
\hyperref[sec:toc]{[返目錄]}
~
\hyperref[sec:375]{> > > NEXT SERMON > > >}


創世記 12:1-20
\newline
\begin{longtable}{cl}
\hline
\hline
章節 & 經文 (和合本修訂版)\\
\hline
12:1 & \begin{tabularx}{0.7\textwidth}{X} 耶和華對亞伯蘭說：「你要離開本地、本族、父家，往我所要指示你的地去。 \end{tabularx} \\ \\ \relax
12:2 & \begin{tabularx}{0.7\textwidth}{X} 我必使你成為大國，我必賜福給你，使你的名為大；你要使別人得福。 \end{tabularx} \\ \\ \relax
12:3 & \begin{tabularx}{0.7\textwidth}{X} 為你祝福的，我必賜福給他；詛咒你的，我必詛咒他。地上的萬族都必因你得福。」 \end{tabularx} \\ \\ \relax
12:4 & \begin{tabularx}{0.7\textwidth}{X} 亞伯蘭就遵照耶和華的吩咐去了；羅得也和他同去。亞伯蘭離開哈蘭的時候年七十五歲。 \end{tabularx} \\ \\ \relax
12:5 & \begin{tabularx}{0.7\textwidth}{X} 亞伯蘭帶著他妻子撒萊和姪兒羅得，以及他們在哈蘭積蓄的財物、獲得的人口，往迦南地去。他們就來到了迦南地。 \end{tabularx} \\ \\ \relax
12:6 & \begin{tabularx}{0.7\textwidth}{X} 亞伯蘭經過那地，直到示劍地方，摩利橡樹那裡；當時迦南人住在那地。 \end{tabularx} \\ \\ \relax
12:7 & \begin{tabularx}{0.7\textwidth}{X} 耶和華向亞伯蘭顯現，說：「我要把這地賜給你的後裔。」亞伯蘭就在那裡為向他顯現的耶和華築了一座壇。 \end{tabularx} \\ \\ \relax
12:8 & \begin{tabularx}{0.7\textwidth}{X} 從那裡他又遷到伯特利東邊的山，支搭帳棚；西邊是伯特利，東邊是艾。他在那裡又為耶和華築了一座壇，求告耶和華的名。 \end{tabularx} \\ \\ \relax
12:9 & \begin{tabularx}{0.7\textwidth}{X} 後來亞伯蘭漸漸遷往尼革夫去。 \end{tabularx} \\ \\ \relax
12:10 & \begin{tabularx}{0.7\textwidth}{X} 那地遭遇饑荒。亞伯蘭因那地的饑荒嚴重，就下到埃及，要在那裡寄居。 \end{tabularx} \\ \\ \relax
12:11 & \begin{tabularx}{0.7\textwidth}{X} 將近埃及，他對妻子撒萊說：「看哪，我知道你是美貌的女人。 \end{tabularx} \\ \\ \relax
12:12 & \begin{tabularx}{0.7\textwidth}{X} 埃及人看見你會說：『這是他的妻子』，他們就會殺我，卻讓你活著。 \end{tabularx} \\ \\ \relax
12:13 & \begin{tabularx}{0.7\textwidth}{X} 所以，請你說你是我的妹妹，使我可以因你得平安，我的性命也因你存活。」 \end{tabularx} \\ \\ \relax
12:14 & \begin{tabularx}{0.7\textwidth}{X} 亞伯蘭到達埃及時，埃及人看見那女人極其美貌。 \end{tabularx} \\ \\ \relax
12:15 & \begin{tabularx}{0.7\textwidth}{X} 法老的臣僕看見了她，就在法老面前稱讚她。那女人就被帶進法老的宮中。 \end{tabularx} \\ \\ \relax
12:16 & \begin{tabularx}{0.7\textwidth}{X} 法老就因她厚待亞伯蘭，給了亞伯蘭許多牛、羊、公驢、奴僕、婢女、母驢、駱駝。 \end{tabularx} \\ \\ \relax
12:17 & \begin{tabularx}{0.7\textwidth}{X} 耶和華因亞伯蘭妻子撒萊的緣故，降大災擊打法老和他的全家。 \end{tabularx} \\ \\ \relax
12:18 & \begin{tabularx}{0.7\textwidth}{X} 法老召了亞伯蘭來，說：「你向我做的是甚麼事呢？為甚麼沒有告訴我她是你的妻子？ \end{tabularx} \\ \\ \relax
12:19 & \begin{tabularx}{0.7\textwidth}{X} 為甚麼說『她是我的妹妹』，以致我把她接來要作我的妻子呢？現在 ，看哪，你的妻子在這裡，帶她走吧！」 \end{tabularx} \\ \\ \relax
12:20 & \begin{tabularx}{0.7\textwidth}{X} 於是法老吩咐人把亞伯蘭和他妻子，以及他一切所有的都送走了。 \end{tabularx} \\ \\
[1ex]
\hline
\hline
\end{longtable}
講題:惜別講章:與神同行祝福人! Farewell Message: Make My Life Be a Blessing!

經文:創世記12:1-20

講員:許開明牧師

場合:香港中文大學崇基學院禮拜堂

日期:2017年7月9日

一 引言

今早是一篇惜別的講章,因為是我擔任校牧以來,最後的一篇講章.臨別賜言,有特別的意義.

由於7月31日離任,8月1日就上任,負責的堂會有四間超過五百人的堂會及一間佈道所,四或五個主日需要在有關堂會宣道.何時可以回來崇基禮拜堂,真是難說了.

今早的題目是「與神同行祝福人」,經文是創12:1-20節.描述亞伯拉罕如何與神同行並祝福別人!這句話也是我對自己的期望.所以我唱出來也講出來,希望在大家心裡留下深刻的印象.

中國人喜歡問一個問題:「好唔好命? 」

「我好唔好命?」,「你好唔好命?」,「他好唔好命?」.

人期望好命,追求好命,使自己「有一條好命」.但怎樣才算是「好命?」這可謂見仁見智.

但有一點是大部份人認同的:「好命(或命水)」都含有一種運氣.對於一般人的理解,「好命」往往不只是人的努力,還有上天所賜.「福」字的部首是「神」,中國人很早便體會到「福由神降」的道理.

到底「好命」是全然靠人努力或全然都是天賜?如果全然是天賜,人還需努力嗎? 還是「好命」是天賜與人力兩者皆有?自古到今,這些問題都是不易解答的.

另外,到底好命,是單指自己的生命,還是也指使人活得好的生命?

今日我們嘗試從亞伯拉罕的生命去看,「與神同行祝福人」.明顯地,我們可以從聖經中看到他的生命是「好命」的.他的好命不單指上帝對他生命的祝福,而且因為他亦能夠祝福別人.亞伯蘭的意思為「尊貴之父」,他是被上帝所祝福的,亦擁有祝福人的生命.

三種生命存在的模式

一,生命使人受傷或受苦

從古到今,許多人自私自利,把快樂建立在別人的痛苦上.例如最近誘騙未成年男女做壞事,偷運象牙圖利.

不少人的生命不但不能為人帶來祝福,反而為人帶來負累,阻礙,甚至災難.好像昔日始祖夏娃吃完禁果後,也把禁果給亞當吃,夫婦一同犯罪,把罪帶到全世界.結果他們被上帝逐出伊甸園；又如掃羅妒忌大衛不停追殺他,把痛苦帶給大衛.

二,生命叫人得福

中國聖賢,儒家提出了「修身,齊家,治國,平天下」,含有貢獻自己,祝福別人的思想.

從聖經的記載,我們可以發現上帝不單祝福亞伯蘭他自己,而且亦透過他成為別人祝福,祝福他的家族,家人,以及整個以色列的民族.他的生命是「因你得福」及「叫人得福」的生命:

「你要離開本地,本族,父家,往我所要指示你的地去.我必叫你成為大國.我必賜福給你,叫你的名為大；你也要叫別人得福.為你祝福的,我必賜福與他；那咒詛你的,我必咒詛他.地上的萬族都要因你得福.」(創十二1-3).

三,生命的存在可有可無

這可能是最多人的生存模式.這個世界有他或沒有他都是一樣,是可有可無的.

馬克思有一個觀念,認為人本應有自我價值和社會價值.人能夠貢獻社會.但今天有些人活著既無自我價值,也無社會價值.

亞伯蘭(尊貴之父)的生命叫人得福,到底是一個怎麼樣的「福」呢?

首先,是離開吾珥,一個拜月神之地,得以認識真神,並單單敬拜衪.

另外,是離開迦南地,饑餓之地,進入埃及取得食物.然後離開埃及危險之地,生命得以保全還得財富.

最後進入迦南美地,結束游牧民族生活,進入定居,建立民族,成為上帝的子民!

成功與價值?

某記者訪問猶太裔科學家愛因斯坦說︰「你認為自己是成功的人嗎?」

當時,愛因斯坦回答說︰「我從來都不想做個成功的人,只想做個有價值的人.」

記者問︰「為甚麼?」

愛因斯坦回答︰

「因為成功是別人給你的,有價值卻是你去給人的.」

各位,你認為你是位有價值的人嗎?或者是可有可無的人.

我相信:一個有價值的人,他的生命能祝福別人!

亞伯蘭怎樣與神同行?

一,聽從神

與神同行,就要聽從神的吩咐.上主要亞伯蘭離開本地(一個原本熟識和安全的地方),本族(一個原本富庶的地方),父家(一個他已擁有權力的地方),去一個他不認識陌生地方.雖然如此,但亞伯蘭仍願意放下一切聽從上帝的呼召,「亞伯蘭就照著耶和華的吩咐去了.」(創十二4)那顯示是亞伯蘭對上帝的聽從.

甚麼是聽從?我們需要信任神和冒險的精神.

從人的成長表現來看,孩童期較易聽從.到了少年期便開始提問:為何要聽從?而青年期多不會隨便的聽從,特別是對權威的聽從.到了中年人或年長的人要聽從,就更加困難!

亞伯拉罕接受上帝的吩咐時,已75歲了,如此年長要聽從,確實是不容易的事!

倘若亞伯蘭不聽從,以色列民族得不到祝福,救恩便得不到實現,影響深遠.

聽從是蒙福的條件.許多人一生的成功或蒙福都是因為聽從.亞伯蘭如是.

馬利亞聽從上主使者,接受聖靈感孕生子的使命,結果萬代稱她為有福.摩西照上主的吩咐去見法老取得准許帶以色列人出埃及,把以色列人組聖民族,成為民族的祝福.

自古到今,人的悖逆,不肯聽從上帝,而失去福氣,陷入痛苦中!始祖犯下的罪:不聽從,吃禁果!

介

我的聽從

預科畢業後,有機會入大學,我想修讀中文系,上主卻叫我讀宗教,我聽從上主.

我沒有得到教會支持,卻得到大學的獎學金/助學金完成四年的教育.

大學畢業後,有機會教書或牧會.父親希望我教書,但我在靈修中,上主感動我牧會.

牧會二十年後區會鼓勵我進修研究學位:本來我想到美國,上主卻叫我去英國,五年完成兩個學位.

完成後返港,想全職教學.上主卻叫我不要放下牧會,並引領我兼職到神學院任教,所以自學成回港,我在五間神學院兼教,同時牧會及作研究.

當中的過程是很忙,但卻很開心,因為生命能成為別人的祝福.

二,信靠神

與神同行,要信靠神.

信靠什麼?就是肯定神的信實,相信祂的應許會實現.

亞伯蘭每到一處,就會築壇,藉此與上主接觸, 跟他說話,使亞伯蘭行在上主的路上,不致偏離真道,親近上主,表明依靠上主.

三,跟隨神

「亞伯蘭經過那地,到了示劍地方摩利橡樹那裡.那時迦南人住在那地.耶和華向亞伯蘭顯現,說,我要把這地賜給你的後裔.亞伯蘭就在那裡為向他顯現的耶和華築了一座壇.

從那裡他又遷到伯特利東邊的山,支搭帳棚.西邊是伯特利,東邊是艾.他在那裡又為耶和華築了一座壇,求告耶和華的名.」

「亞伯蘭就在那裡為他顯現的耶和華築了一座壇」(創12:7下).昔日以色列人築壇,為要求告耶和華的名,親近上帝,敬拜上帝,並紀念及感謝上帝的作為.亞伯蘭體會到主恩的豐盛:神帶領他平安到示劍,伯特利和希伯倫,於是亞伯蘭便築壇告訴別人,為要常紀念神的作為.亞伯蘭體會到神的祝福,也願成為別人的祝福.

4 亞伯蘭得到什麼「福」?

一,得蒙福的應許

1.賜福的應許

「我必叫你成為大國；我必賜福給你,叫你的名為大,你也要叫別人得福.為你祝福的,我必賜福與他；那咒詛你的,我必咒詛他；地上的萬族都要因你得福.」(創十二2-3)

2.得到示劍地

「亞伯蘭將他妻子撒萊,和姪兒羅得,連他們在哈蘭所積蓄的財物,所得的人口,都帶往迦南地去；他們就到了迦南地.」「亞伯蘭經過那地,到了示劍地方摩利橡樹那裡；那時迦南人住在那地.」「耶和華向亞伯蘭顯現,說,我要把這地賜給你的後裔.亞伯蘭就在那裡為向他顯現的耶和華築了一座壇.」(創十二5-7)

3.得到豐盛離開埃及回到迦南地

「亞伯蘭帶著他的妻子與羅得,並一切所有的,都從埃及上迦南地去.亞伯蘭的金,銀,牲畜極多.」(創十三1-2)

去的時候是饑荒,充滿憂慮,走時卻是豐盛,滿有平安.

4.不是迦南地使亞伯蘭得福,而是亞伯蘭帶財物,食物到迦南地.

這蒙應許給亞伯蘭希望,願景,和動力向前!

我大半生得上主的祝福:

1少時,上主聽禱告,讓我於1963年就抵港,避開文革的批鬥.

2抵港後,生活艱難,卻得到數位牧師的關顧,衣食住均可以在教堂,有美好的學習環境.

3在教會認識太太,三個仔女都在牧師宿舍出生和居住.

上主的祝福和恩典數不盡,我存感恩的心事奉!

二,得可見及不可見的福

亞伯蘭得到財富,牛羊,駱駝,僕婢,土地 (可見)；得到太太:得到尊嚴；得到耶和華保護:得到平安 (不可見的).

三,得變相的祝福

我們所信的是一位奇妙的神,祂容許人犯錯,容許人經歷患難,然後深刻地體會祂的祝福.

上帝祝福亞伯蘭後,到了迦南地後,好像沒有甚麼真正的福給亞伯蘭,反而亞伯蘭和全族人經歷飢荒.當時迦南地有大飢荒,而埃及因有尼羅河滋潤的緣故甚少有飢荒.

亞伯蘭因此帶著妻子離開迦南下到埃及.

將近埃及時,亞伯蘭因妻子撒萊容貌俊美,皮膚白晢,極為美麗,他害怕因妻子的美貌而帶來殺身之禍,故此他把妻子認作妹子,以便得以存活.果然撒萊被帶到埃及王宮作法老的妃子,這對撒萊是大禍,而亞伯蘭亦忍受失去他最愛的人的苦痛.

人生常會面對「失去」,失去其最愛的人或物!

例如歌唱家失去聲音,演奏家傷了手指,愛踢足球的人傷了腳等.上帝有時會容許我們陷在失去的苦況中,以便我們能更多的經歷祂,讓我們的生命得到成長.

信心之父也軟弱

亞伯蘭沒有信心面對挑戰,信心之父亦會因懼怕而軟弱犯錯.

信心之父也是人,也有人性的軟弱.

當時他面對人生兩大能題:

一是飢荒,缺乏食物,這不獨是他自己而且是全族人的難題；

二是強敵法老,古代最強的君主,偶一不慎,他便有殺身之禍,甚至滅族大能.面對這兩大能關,亞伯蘭也因此信心軟弱而說謊,告訴法老撒萊只是他的妹妹.

上帝看來也體諒亞伯蘭的軟弱,讓他在患難中體會上帝祝福的奇妙.但這並不表示信實的上帝並認同亞伯蘭的所為.

神藉法老的口責備亞伯蘭(見創12:18-20).

這兩節話中,法老責問了亞伯蘭三個問題:

你向我作的是甚麼事呢?(你指責,以下犯上!)

為甚麼沒有告訴我他是你的妻子?(你忍瞞事實)

為甚麼說他是你的妹子? (你講大話!)

以致我把她娶來要作我的妻子 (你導致他人犯錯)

亞伯蘭都沒有答覆.他是無詞以對.自己理虧,不敢說甚麼.

亞伯蘭犯下三條死罪!

在禍患中顯祝福

無論如何,上帝在禍患中顯出祂對亞伯蘭的祝福.

亞伯蘭因撒萊的緣故得到法老厚待眾多財物包括牛,羊,駱駝,公驢,母驢和僕婢,這樣解決飢荒問題.

耶和華又因撒萊的緣故,降大災給法老和他的全家,因她對上主的信心和依靠,老公都保不到她,惟一的是上主,使法老讓亞伯蘭可以帶著妻子並在埃及所得的一切財物回迦南,亞伯蘭因此得以保存性命並解除失去妻子之痛苦.

上帝要祝福我們,不一定讓我們凡事一帆風順,有時在苦難中讓我們能更深體會祂的祝福.同時,藉著亞伯拉穻的事跡,讓我們看見苦難落在上帝手中,可成為變相的祝福.其實禍福乃在上帝手中,祂是人生際遇的主宰.

那是否表示:人可以放縱自己去犯罪而得福?當然不是!

亞伯蘭的處境:

認撒萊為妻,可能即被殺!甚至全部人被殺或作奴隸!

認撒萊為妹,可能不死,全家不殺,還可以得財富!

當然,被識穿也是死路一條!使識穿是因為上主使穿:降大災予法老!

耶和華又因撒萊的緣故,降大災給法老和他的全家.

四,得出人意外的福

亞伯蘭在埃及,說謊被揭發後,竟然可以得回生命,得回美麗的太太,還得到財富.為什麼法老對亞伯蘭那麼仁慈?亞伯蘭一個小猶長竟然怕他,或竟然法老知錯悔改?

絕對不是!

法老只是怕亞伯蘭的耶和華!法老是受埃及眾神保護的王,竟然眾神都保不住他的龍體受攻擊!聰明法老知道:他不放亞伯蘭,亞伯蘭的神就不放過他了.在這人為的錯誤,不獨是使亞伯蘭,「蘭嫂」轉危為安,得回生命,至愛,避過饑荒,得到法老大批的財富,更顯示了耶和華勝過一切埃及的神!這是這出人意外的祝福!

五,與神同行,活出祝福人的生命

到底祝福別人生命這觀點,是亞伯蘭自發的思想,還是神的思想?

從整段經文來看,不是亞伯蘭的,而是上主的思想!

上主要亞伯蘭去祝福他人!

上主:祝福的源頭；

亞伯蘭要與神同行,得到神的祝褔,人效法神,去祝福別人!

與神同行祝福人

與神同行,祝褔別人

認識祝福是甚麼?

「福」在聖經原文意思異常豐富,其中較為重要的有四個︰

第一個是指「好關係」,「和諧的關係」,是與人,與神,與自己和諧,含有平安的意思.第二是指「開心」,幸福(Happiness) ,喜樂的感覺.第三是「好的狀態」(Well Being),有福氣或平安.

第四是指「好的生命」,這「好」是有滋長力量,有生命力,感染力,因而使人也得福.

所以,祝福重要的涵義是使人得平安,並且有一種驅策的力量,使蒙福的人去祝福別人.

希臘人「祝福」這個字eulogia,eu是好,Logos是話,說好話就是祝福,當你對別人說感謝的話,讚美的話,肯定的話,你就是在祝福別人.

因蒙神賜福而祝福別人

祝福就好像從上帝那裡「下載」福氣到我們所祝福的人當中,所以我們在家庭中需要彼此祝福.

孩子祝福父母；父母祝福孩子；兄弟姊妹彼此祝福.

祝福我們的家,祝福我們的教會,祝福我們的公司,祝福我們的國家.

祝福的人會先得祝福的好處.當我們祝福我們的家庭,自己的教會,我們會先覺得以自己的家為榮,以自己的教會為榮,以成為一個基督徒為榮.

雖然我們還沒看到改變,但在我們心裡就已經感受到改變,漸漸的我們也會看到我們祝福的對象也在改變.

「惜別分享」:

雖然擔任崇基校牧時間只有四年,但這四年在我來說,是難忘而且是有意義的.許多值得感恩之處.

八十年代初,在崇基讀書,經常來往神學院和教堂,對當年的校牧蘇以葆牧師及師母,和靄可親,經常關顧學生,我也是他所關心的對象.當年還是神學生對校牧工作已嚮往.

2012年9月校牧招聘一出,當時我是全完堂的主任,在堂會牧會已30年,亦剛好55歲,很想轉環境,於是與當年我會的主席及當年的蘇牧師,今天的蘇主教聯絡,他們十分支持我的申請,願意成為我的諮詢人.

願生命成為別人的祝福是我的心願,這幾年的時間,我相信上帝都充份給我機會:

在學生工作方面:

主持週會,每次都特別為不同主題寫禱文,揀經文,尋找“每週靜思”的中西名言,讓自己大得啟發.不少同學亦向我反映在禱文中得到幫助.

擔任明華堂舍監,還我有許多機會,跟學生在一起,吃同學的早餐,晚間的糖水,去宵夜,打波.明華在我在任期間拿過三屆傅杯拔河冠軍,每次我一到就勝,同學說是牧師的祝福,還有牧仔的參與.

在高危時期,關心及轉介一些高危學生給輔導員,並與大學各方的協力下阻止某同學在Esther Lee 頂樓自殺的計劃.

與輔導處同事一起關心和幫助幾位不交功課,不交學費學生,三位是住明華的學生,有國際生及校的地生.

關心犯事的學生,使其改過及成長,後來信主洗禮,成為校牧室的好幫手Student Helper.

關心一些同學遇到邪靈的干擾,或與老師,同學相處有困難.

關心一些信仰或宗教的問題,做宗教的專題研究希望取資料.由於我讀書比一般人長,大學至博士讀的時間都超過20年,資料不少.

在教職員及教堂之友方面:

大部份崇基的委員會會議,聚會及餐會,經常有機會出席,因此與大部份同事及校友都熟落,漸漸可以多點關心他們.

透過教堂的代禱紙,關心他們的需要,為他們代禱或電話的問候.

當中有同事的家人或家庭關係需要代禱.患病需要個別關顧的也不少.感謝許師母,是我的好幫手,特別關顧教堂裡的姊妹.

主日方面:

大學講章,用心的準備,自己最得幫助,然後成為別人的幫助.

讀書會:得到許多思想交流和學習機會.

普通話團契:看到團友熱切認真的追求,每年我辦的講座最多人參加是普團；還有一個個入,一對對的出來.

豐盛團契:因應說粵語的弟兄姊妹團契需要而成立,年前愈來愈豐盛,在信仰上的探討愈來愈深入.

祈禱小組:為中大人和事祈禱,為有需要的肢體祈禱,我們每次均有靈修和祈禱,看見上帝不少奇妙的作為.

基督教事工方面,做得好開心:

有26項的基督教文化事工,林書豪的成功的再思,與醫生及牙醫團契合辦的國際學生醫療特會,還有許多大大小小的音樂會,平安夜聖景燭光歌頌崇拜,Midday Oasis 音樂會等,教會大學尋根團等,感謝上主恩典,校牧室同事的同心協力,使各項活動均榮神益人.

著作方面:

在任期間應崇基65週年校慶編輯委員會邀請,寫了篇《崇基歷任校牧簡介》,讓我對過去15位校牧有所認識,並對校牧的工作有更多的認識.

在任期間應上海復旦大學宗教研究中心的邀請出席基督教高級論壇,完成了《中國教會大學社會服務的在社會的作用》論文.

應華人大學校牧協會邀請寫作一篇專文《從大學校園學生事工探討過去的承傳,現今的挑戰,未來的承擔》,寫此文時,以為自己登陸後,還有幾年做的,所以提出未來的承擔!

應宗教與中國社會研究中心的邀請完成我博士論文的中文版:《中國社會變遷中的趙紫宸人格基督論》.八月將出版.又我十年前英文博士論文 A Study of T. C. Chao’s Christology,也蒙瑞士Peter Lang Academic Publisher出版 ,成為榮休最佳禮物,也是在任期間達成.

個人方面

這四年我們一家五口也住得很開心,牧仔是崇基學生很投入明華堂宿舍,免交通時間,所以可以比其他同學睡多一點,玩多一點,讀多一點,所以他很開心,成績也不錯.師母也很關心兄姊,參加詩班,師母唱歌方面愈唱愈好,我就受罪,專在我準備講道時唱,令我無法專心寫講稿!

感謝主,讓我在校牧室短短四年有美好的事奉,期盼能成為人的祝福

與神同行,祝福人的實踐

1.與神同行,數算日子

「求你指教我們怎樣數算自己的日子,好叫我們得著智慧.」(詩90:12)

「加,減,乘,除」

我們應怎樣數算我們的日子呢?

唐崇榮牧師以「加,減,乘,除」來提醒我們,十分有意思︰

「加號」,就是人早上起床,為神給我新的一天而感恩；

「減號」,就是人晚上睡前,為我人生在世日子又少了一天而珍惜；

「乘號」,就是人盡量善用我的人生做更多有意義的事,不浪費光陰；

「除號」,就是人把人生的時間作均衡的分配,在工作上,家庭上,事奉上,個人靈性上等.

基督徒要常常注意這四個符號對我們人生的啟示.

與神同行,數算主恩

有首詩:「主的恩惠樣樣都要數,主的恩典都要記清楚,

必能叫你驚訝立時樂歡呼.」

認識一位九十多歲的弟兄,他經常說:「每天醒來,起來,行得,走得,食得；主祝福我,我也要祝福別人!」

3 與神同行,祝福他人

「祝福您!」

數年前,偶然在電話錄音中聽到這句話:「祝福您!」十分深刻和回味.縱然是找不到想接觸的人,只聽到致歉的錄音,有點兒失望,但傳來的這結束語,令我不快的情緒一掃而空.

最近,經常與一位敬愛的長者交通,每次在電話談話畢,他總會加上一句結束語:「祝福您!」我也即時回應說:「也祝福您!」令人欣慰.

互相祝福

多年前,我在望覺堂幼稚園畢業禮訓勉,聚會後祝福小朋友.當我舉起雙手時,他們也舉起雙手,我講聖保羅祝福文,他們也講.期間校長十分緊張,連聲向當時擔任校監的我道歉,我回應說:「不要緊!難得我們不可以互相祝福麼!」時隔多年,當時大家互相祝福的情景,印象還非常深刻,我在想,基督徒的生命豈不是要互相祝福的嗎?我們先領受上帝的賜福,然後去祝福別人.

互相祝福,很有意思.

「祝福」擁有一種奇妙的成全能力

聖經裡記述了許多祝福成真的故事,包括:

波阿斯祝福外邦女子路得,結果路得成為他妻子,也成為大衛的祖父母；

撒母耳祝福大衛,結果大衛成為以色列國最偉大的王；

撒迦利亞祝福兒子施洗約翰,結果約翰成為基督的先鋒；

天使迦百列祝福馬利亞,結果萬代稱馬利亞為有福；

耶穌祝福彼得成為教會磐石,結果彼得成為滿有能力的使徒,領多人信主；

復活主祝福門徒,結果門徒由懦弱變剛強,到處為主作見證.

願我們學習,多祝福別人,相信你合神心意的祝福會變成真實.

三,結語

願你的生命也成為別人的祝福

馬利亞的生命曾成為別人的祝福；

主耶穌的生命也成為人的祝福；

亞伯拉罕的生命亦曾成為別人的祝福.

弟兄姊妹們,上帝也願意你的生命成為別人的祝福.

期望你我也能夠:

1數算日子

2數算祝福(恩典)

3去祝福他人

有位同工:「牧師,Give me Blessings」

有個家庭,探完訪:「牧師,請祝福我們一家吧!」

我以祝福人為樂,願大家能彼此祝福.

記得,宣道前,我問「怎樣才算好命?」

現在,我可以清楚回答你: 「一位自覺蒙福又能祝福別人的人」是好命的人.

祈禱:

父上帝,我們的生命是你所賜,你所創造,你所救贖和你所養育的,生命是上主賜福的,願我們認識自己蒙福的生命,懂得去祝福別人,並能彼此祝福,互相扶持! 奉主名求,阿們.

(惜別講章,講於崇基教堂主日學崇拜2017年,7月9日)

獻唱Anthem :  I’ll Walk With God

Word: Paul Francis Webster

Music: Nicholas Brodszky

I’ll walk with God from this day on,

His helping hand I’ll lean upon,

This is my prayer, my humble plea,

May the Lord be ever with me.

There is no death, though eyes grow dim,

There is no fear when I’m near to Him.

I’ll lean on Him forever And He’ll forsake me never.

He will not fail me as long as my faith is strong,

Whatever road I may walk along;

I’ll walk with God, I’ll take His hand,

I’ll talk with God, He’ll understand;

I’ll pray to Him, each day to Him

And He’ll hear the words that I say.

His hand will guide my throne and rod,

And I’ll never walk alone while I walk with God!



\newpage

\section{在泥淖中起舞}
\label{sec:375}
\textbf{鄧瑞強博士 - 在泥淖中起舞}
\newline
\newline
\hyperref[sec:374]{< < < PREV SERMON < < <}
~
\hyperref[sec:toc]{[返目錄]}
~
\hyperref[sec:376]{> > > NEXT SERMON > > >}


羅馬書 8:26-39
\newline
\begin{longtable}{cl}
\hline
\hline
章節 & 經文 (和合本修訂版)\\
\hline
8:26 & \begin{tabularx}{0.7\textwidth}{X} 同樣，我們的軟弱有聖靈幫助。我們本不知道當怎樣禱告，但是聖靈親自用無可言喻的嘆息替我們祈求。 \end{tabularx} \\ \\ \relax
8:27 & \begin{tabularx}{0.7\textwidth}{X} 那鑒察人心的知道聖靈所體貼的，因為聖靈照著神的旨意替聖徒祈求。 \end{tabularx} \\ \\ \relax
8:28 & \begin{tabularx}{0.7\textwidth}{X} 我們知道，萬事都互相效力，叫愛神的人得益處，就是按他旨意被召的人。 \end{tabularx} \\ \\ \relax
8:29 & \begin{tabularx}{0.7\textwidth}{X} 因為他所預知的人，他也預定他們效法他兒子的榜樣，使他兒子在許多弟兄中作長子。 \end{tabularx} \\ \\ \relax
8:30 & \begin{tabularx}{0.7\textwidth}{X} 他所預定的人，他又召他們來；所召來的人，他又稱他們為義；所稱為義的人，他又叫他們得榮耀。 \end{tabularx} \\ \\ \relax
8:31 & \begin{tabularx}{0.7\textwidth}{X} 既是這樣，我們對這些事還要怎麼說呢？神若幫助我們，誰能抵擋我們呢？ \end{tabularx} \\ \\ \relax
8:32 & \begin{tabularx}{0.7\textwidth}{X} 神既不顧惜自己的兒子，為我們眾人捨了他，豈不也把萬物和他一同白白地賜給我們嗎？ \end{tabularx} \\ \\ \relax
8:33 & \begin{tabularx}{0.7\textwidth}{X} 誰能控告神所揀選的人呢？有神稱他們為義了。 \end{tabularx} \\ \\ \relax
8:34 & \begin{tabularx}{0.7\textwidth}{X} 誰能定他們的罪呢？有基督耶穌 已經死了，而且復活了，現今在神的右邊，也替我們祈求。 \end{tabularx} \\ \\ \relax
8:35 & \begin{tabularx}{0.7\textwidth}{X} 誰能使我們與基督的愛隔絕呢？難道是患難嗎？是困苦嗎？是迫害嗎？是飢餓嗎？是赤身露體嗎？是危險嗎？是刀劍嗎？ \end{tabularx} \\ \\ \relax
8:36 & \begin{tabularx}{0.7\textwidth}{X} 如經上所記：「我們為你的緣故終日被殺；人看我們如將宰的羊。」 \end{tabularx} \\ \\ \relax
8:37 & \begin{tabularx}{0.7\textwidth}{X} 然而，靠著愛我們的主，在這一切的事上，我們已經得勝有餘了。 \end{tabularx} \\ \\ \relax
8:38 & \begin{tabularx}{0.7\textwidth}{X} 因為我深信，無論是死，是活，是天使，是掌權的，是有權能的 ，是現在的事，是將來的事， \end{tabularx} \\ \\ \relax
8:39 & \begin{tabularx}{0.7\textwidth}{X} 是高處的，是深處的，是別的受造之物，都不能使我們與神的愛隔絕，這愛是在我們的主基督耶穌裡的。 \end{tabularx} \\ \\
[1ex]
\hline
\hline
\end{longtable}
講題:在泥淖中起舞　Singing in the Dark

經文:羅馬書8:26-39

講員:鄧瑞強博士

場合:香港中文大學崇基學院禮拜堂

日期:2017年7月30日

各位弟兄姊妹,早安.願你們平安.

從上次我在這裡講道,到今天我在這裡講道,已有好一段日子了.在這段時間裡,你們的世界,或許,沒有什麼變化.我的世界,卻已不能再一樣了.大家知道,在這段時間裡,我太太忽然的離開了這個世界,回到天父的懷裡去了.你們的生活,或許,都依舊是這樣.早上起來,吃早餐,早餐不是Ａ餐,便是Ｂ餐.之後,開始工作,工作的情景日日大致如是.之後,回家,與家人吃一頓飯.之後,睡覺,發的夢,也都是差不多的.而我,整個生活變了.早上起來,三十年來,太太煮的早餐,不再有了.之後,開始工作,但工作時某些感覺不再一樣.不再寄望出糧的日子去吃大餐,也不再想著如何安排假期.之後,回家了,空洞無人.之後,睡覺,再沒有夢了.過去的回憶,都變得空白；而將來的計劃,都再沒有內容.

我好似活在這個世界,但又好似不是活在這個世界.心理學講,這是一種「derealization」現象.一個經歷創傷的人,會拒絕現實,會從現實世界中抽離出來,進入一個若即若離,不真不實的世界裡.我明白心理學有這種說法.就我而論,我不是逃離這個現實世界,而是從一個意義世界,因著一件死亡事件,進入到另一個意義世界裡.不太恰當地借用一小節經文來說,便是「舊事已過,都變成新的了」.舊的世界,已漸漸遠去.新的世界,正慢慢形成.但這過程並不乾脆俐落,新的世界正在開展,但舊的人,舊的事,舊的情感,總會無端而來,拉我回到舊日的世界,忍受從其中滲透出來的點點哀愁.

信耶穌,從一個真實的意義來說,就是放下一個舊的世界,進入一個全新的世界裡.保羅在《羅馬書》第六章講到,信耶穌的真義,可以藉洗禮表達的.一個人受洗,浸入水中,代表一種埋葬,與基督同葬,與基督同死.然後,從水中上來,是新生命的開始,與基督一同復活過來.舊的一切已過去了,從水中上來,睜開眼,見到的,就是一個全新的世界,一個奇異的新世界.在這個奇異的新世界裡,不用逃避當下的現實,卻也不完全擁抱現實,因看到現實背後隱伏著神創造的新天新地.在這個奇異的新世界裡,我們肯定看得見的這個世界是真的,也同時肯定有一個看不見的世界,那世界也是真的.在這個奇異的新世界裡,我們看見神成為人,看見至尊貴的神為卑污的人而死,看見不可思議的死而復活事件.一個人,因著信仰,踏入了這個奇異的新世界裡.

進入信仰,就進入一個生命轉化的過程,由「舊的世界」過渡到「新的世界」裡去.但新的生命與舊的生命,在拉鋸中.舊的世界遠去了,但其頑固的力量仍在.進入了新的世界,卻又未能完全體會一種全新的生活的意義.想像一下,你剛剛移民到非洲一條無電力供應的村落裡,晚上,由於過去的習慣,你可能苦悶無比,因再無電視可看了,但更關鍵的,是你還未學懂欣賞滿天的繁星.

新生命與舊生命,有一種張力.張力底下,有一種苦悶.如何是好?

我們看看今天的講道經文:羅馬書8:26－39

羅8:26　況且我們的軟弱有聖靈幫助,我們本不曉得當怎樣禱告,只是聖靈親自用說不出來的歎息替我們禱告.

羅8:27　鑒察人心的,曉得勝靈的意思,因為聖靈照著神的旨意替聖徒祈求.

羅8:28　我們曉得萬事都互相效力,叫愛神的人得益處,就是按他旨意被召的人.

羅8:29　因為他預先所知道的人,就預先定下效法他兒子的模樣,使他兒子在許多弟兄中作長子.

羅8:30　預先所定下的人又召他們來；所召來的人又稱他們為義；所稱為義的人又叫他們得榮耀.

羅8:31　既是這樣,還有甚麼說的呢?神若幫助我們,誰能敵擋我們呢?

羅8:32　神既不愛惜自己的兒子,為我們眾人捨了,豈不也把萬物和他一同白白地賜給我們嗎?

羅8:33　誰能控告神所揀選的人呢?有神稱他們為義了.

羅8:34　誰能定他們的罪呢?有基督耶穌已經死了,而且從死裡復活,現今在神的右邊,也替我們祈求.

羅8:35　誰能使我們與基督的愛隔絕呢?難道是患難嗎?是困苦嗎?是逼迫嗎?是飢餓嗎?是赤身露體嗎?是危險嗎?是刀劍嗎?

羅8:36　如經上所記:我們為你的緣故終日被殺；人看我們如將宰的羊.

羅8:37　然而,靠著愛我們的主,在這一切的事上已經得勝有餘了.

羅8:38　因為我深信無論是死,是生,是天使,是掌權的,是有能的,是現在的事,是將來的事,

羅8:39　是高處的,是低處的,是別的受造之物,都不能叫我們與神的愛隔絕；這愛是在我們的主基督耶穌裡的.

(1)「舊世界」,「舊生活」之下的惆悵歎息

這些經文的作者保羅完全明白,在基督裡,人可得到新生命,永恆的生命.但現實是,我們更多體會到的,是死亡之下的無能為力.我們不單單盼望,也實在已踏入神創造,救贖,聖化的新世界裡,但同時,我們仍活在罪污,醜陋,必死的舊生命中.在信仰裡,舊我與新我,現實與盼望,總有一張力.

羅8:26　況且我們的軟弱有聖靈幫助,我們本不曉得當怎樣禱告,只是聖靈親自用說不出來的歎息替我們禱告.

羅8:27　鑒察人心的,曉得勝靈的意思,因為聖靈照著神的旨意替聖徒祈求.

保羅在這裡,講到舊我的軟弱,講到一種不知應如何祈禱的生命困境.孔子說:「獲罪於天,無所禱也.」得罪了上天,怎樣祈禱也沒有用了.保羅在《羅馬書》曾說:「我真是苦啊!誰能救我脫離這取死的身體呢?」(羅7:24)所謂「取死」,其實指一種取消生命的力量.我們的生命本身,有一種取消生命的頑固力量存在.善意興起時,惡念油然而生.努力去愛時,卻無法驅除自私的念頭.心靈思想天堂時,腳卻不由自主地走向地獄.保羅說:「立志為善由得我,只是行出來由不得我.」(羅7:18)又說:「我覺得有個律,就是我願意為善的時候,便有惡與我同在.」(羅7:21)看來,人生一切的善,都被大惡消滅.生命的火光,被死亡掩蓋.在這種無可如何的光景裡,保羅說:「不曉得當怎樣禱告」.「無所禱也!」

死亡,代表著此生終極的虛無力量,能將生命裡的種種美善一下子摧毀.在自然生命的領域,死亡是邊界,我們無力越過這邊界.在這邊界上,我們很多時也不知該如何祈禱.當我太太在死亡的邊界上時,我也不知該如何祈禱.千頭萬緒,只感到平時站得穩的大地在崩塌.而當大地在崩塌時,你可以如何祈禱呢?

舊世界在崩塌,新生命卻在開展.保羅說:「聖靈親自用說不出來的歎息替我們禱告」.人無所禱,聖靈,上主之靈,卻在人不能言語的時候以其無聲的禱告托著人.聖靈的禱告,是無言的歎息.祂歎息,是因為我們在歎息.

保羅說:「鑒察人心的,曉得勝靈的意思,因為聖靈照著神的旨意替聖徒祈求.」天父明白聖靈的無言歎息,「因為聖靈照著神的旨意替聖徒祈求」.這些無言歎息,本來就是天父的旨意,也是天父內心的無言歎息.天父將人在世的歎息,變成祂自己的歎息.看見人類在罪惡,醜陋,死亡中哀鳴,天父也有感於衷,天人同歎.天父將人生的歎息,放在祂心中.正因這份天心,痛苦的人才有出路.

(2)天心在人間的顯露

天父明白必死的人的悲苦,祂的心默默歎息,祂抹去人類眼淚的手在行動.

羅8:28　我們曉得萬事都互相效力,叫愛神的人得益處,就是按他旨意被召的人.

信耶穌的人,明白上天的心意,他見到一個全新的世界.他看到,這個世界的表面背後,有神的仁慈的手在工作.萬事萬物,在神手中,都滲透著神的善意.願意與天心連合的人,他能體會這份美意,從中得到生命的好處.

羅8:29　因為他預先所知道的人,就預先定下效法他兒子的模樣,使他兒子在許多弟兄中作長子.

羅8:30　預先所定下的人又召他們來；所召來的人又稱他們為義；所稱為義的人又叫他們得榮耀.

這裡提到天父「預先知道」,「預先定下」,重點不在神「命定」一切,重點在神的「主動性」.上天主動施予恩典,像雨一樣,沛然降臨,人要做的,是用心去盛載.天父渴望我們,「效法祂兒子的模樣」.上天的恩雨,就是天父的兒子耶穌.耶穌在地上,卻活出天上的生命.他以無盡的愛,接觸每一個在人間哀歎的人.他進入人類的大限,藉此顯露復活的光明.耶穌在十架上,站在生死的邊界上,他第一句話,是求天父寬恕釘他上十架上的人.他沒有敵人,也希望因著他寬恕的愛,敵人能由「舊的生命」進入「新的生命」中.信徒,是分享耶穌生命的人,在人間的大限中,體會愛,展現愛,進入那超越生死的神聖生命裡.保羅稱這生命的光景為「得榮耀」,生命顯出了光彩.

如何能驅散生命的黑暗,展現光彩?

保羅問了五個修辭性的問題,讓我們體會神聖的新生命,如何正在我們生命中展開.

羅8:31　既是這樣,還有甚麼說的呢?神若幫助我們,誰能敵擋我們呢?

什麼東西會消滅我們生命的光明呢?是心中邪惡,恨意,被傷害後留下的空洞,制度化的暴力,死亡的力量…….但若果神站在你的一邊,你能否看到,情況會大大不同.

羅8:32　神既不愛惜自己的兒子,為我們眾人捨了,豈不也把萬物和他一同白白地賜給我們嗎?

你的生命,會欠缺什麼嗎?若果神連自己的兒子(也即神自己)也送給你,你還會欠缺什麼嗎?

羅8:33　誰能控告神所揀選的人呢?有神稱他們為義了.

人間有種種控訴.你說你沒有敵人,敵人依然會控訴你.但若果神為你辯護,神稱你為「義」(即無罪),你能否體會一種塵世以外的「新生命」呢?

羅8:34　誰能定他們的罪呢?有基督耶穌已經死了,而且從死裡復活,現今在神的右邊,也替我們祈求.

很多東西定我們的罪,外在的政權,我們自己的理想,我們自己的良心,都會定我們的罪.但若果神的兒子為你的罪犧牲了,你會否對新生命的可能,有新的領悟?

羅8:35　誰能使我們與基督的愛隔絕呢?難道是患難嗎?是困苦嗎?是逼迫嗎?是飢餓嗎?是赤身露體嗎?是危險嗎?是刀劍嗎?

羅8:36　如經上所記:我們為你的緣故終日被殺；人看我們如將宰的羊.

羅8:37　然而,靠著愛我們的主,在這一切的事上已經得勝有餘了.

這是最後一個問題.人間的種種危難,生命的種種缺乏,死亡的威脅,這一切,能使我們與基督的愛隔絕嗎?這一切威脅都是真的,但若果耶穌在這種種威脅下仍能活出那種充滿光輝的生命,那麼,我們是否也應對那種在基督生命裡展示的新生命有信心?

保羅不是無視現實,藉著這五個問題,保羅展示塵世生命多方多面的艱難與壓力.然而,他感到聖靈的力量,他在耶穌身上看到一種新生命的可能,他進而體會天父的無邊仁愛.他從舊的生命世界,過渡到新的生命世界.

(3)新生命的宣告

羅8:38　因為我深信無論是死,是生,是天使,是掌權的,是有能的,是現在的事,是將來的事,

羅8:39　是高處的,是低處的,是別的受造之物,都不能叫我們與神的愛隔絕；這愛是在我們的主基督耶穌裡的.

這裡提到很多組合,對比.「生與死」,是一組.「天使與掌權的」,應指「天使與邪靈」,是一組.「現在與將來」,是跨越時間,是一組.「高處與低處」,是跨越空間,是一組.「高處」可能指星宿之力,「低處」可能指陰間之力.此外,「有能的」,可能指生命背後的種種神祕力量.「別的受造之物」,可能指種種意想不到的事情.保羅宣布,這一切,都不能叫我們與神的愛隔絕.

我們是必死的人,在大限的界線內,試圖用種種方式去越過這界線.我曾經站在這界線旁,我知道,人間的一切,在這裡,都顯得無能為力.像跌入泥淖裡,你無法起舞.這時,能撐起整個人的,是聽到聖靈在呢喃,是看到耶穌那復活的光明,是體會天父那份大愛.唯這大愛,人能在泥淖中起舞.願你們也找到這讓人起舞的力量.

但願:榮耀歸於聖父,聖子,聖靈.阿們.



\newpage

\section{上帝的逆道與人間的常道}
\label{sec:376}
\textbf{謝任生牧師 - 上帝的逆道與人間的常道}
\newline
\newline
\hyperref[sec:375]{< < < PREV SERMON < < <}
~
\hyperref[sec:toc]{[返目錄]}
~
\hyperref[sec:377]{> > > NEXT SERMON > > >}


路加福音 1:28-55
\newline
\begin{longtable}{cl}
\hline
\hline
章節 & 經文 (和合本修訂版)\\
\hline
1:28 & \begin{tabularx}{0.7\textwidth}{X} 天使進去，對她說：「蒙大恩的女子，你好，主和你同在！」 \end{tabularx} \\ \\ \relax
1:29 & \begin{tabularx}{0.7\textwidth}{X} 馬利亞因這話就很驚慌，又反覆思考這樣問候是甚麼意思。 \end{tabularx} \\ \\ \relax
1:30 & \begin{tabularx}{0.7\textwidth}{X} 天使對她說：「馬利亞，不要怕，你在神面前已經蒙恩了。 \end{tabularx} \\ \\ \relax
1:31 & \begin{tabularx}{0.7\textwidth}{X} 你要懷孕生子，要給他起名叫耶穌。 \end{tabularx} \\ \\ \relax
1:32 & \begin{tabularx}{0.7\textwidth}{X} 他將要為大，稱為至高者的兒子；主神要把他祖先大衛的王位給他。 \end{tabularx} \\ \\ \relax
1:33 & \begin{tabularx}{0.7\textwidth}{X} 他要作雅各家的王，直到永遠；他的國沒有窮盡。」 \end{tabularx} \\ \\ \relax
1:34 & \begin{tabularx}{0.7\textwidth}{X} 馬利亞對天使說：「我沒有出嫁，怎麼會有這事呢？」 \end{tabularx} \\ \\ \relax
1:35 & \begin{tabularx}{0.7\textwidth}{X} 天使回答她說：「聖靈要臨到你身上；至高者的能力要庇蔭你，因此，那要出生的聖者要稱為神的兒子。 \end{tabularx} \\ \\ \relax
1:36 & \begin{tabularx}{0.7\textwidth}{X} 況且，你的親戚伊利莎白，就是那素來稱為不生育的，在年老的時候也懷了男胎，現在懷孕六個月了。 \end{tabularx} \\ \\ \relax
1:37 & \begin{tabularx}{0.7\textwidth}{X} 因為，出於神的話，沒有一句不帶能力的。」 \end{tabularx} \\ \\ \relax
1:38 & \begin{tabularx}{0.7\textwidth}{X} 馬利亞說：「我是主的使女，願意照你的話實現在我身上。」於是天使離開她去了。 \end{tabularx} \\ \\ \relax
1:39 & \begin{tabularx}{0.7\textwidth}{X} 在那些日子，馬利亞起身，急忙前往山區，來到猶大的一座城， \end{tabularx} \\ \\ \relax
1:40 & \begin{tabularx}{0.7\textwidth}{X} 進了撒迦利亞的家，向伊利莎白問安。 \end{tabularx} \\ \\ \relax
1:41 & \begin{tabularx}{0.7\textwidth}{X} 伊利莎白一聽到馬利亞問安，所懷的胎就在腹裡跳動。伊利莎白被聖靈充滿， \end{tabularx} \\ \\ \relax
1:42 & \begin{tabularx}{0.7\textwidth}{X} 高聲喊著說：「你在婦女中是有福的！你所懷的胎也是有福的！ \end{tabularx} \\ \\ \relax
1:43 & \begin{tabularx}{0.7\textwidth}{X} 我主的母親到我這裡來，為何這事臨到我呢？ \end{tabularx} \\ \\ \relax
1:44 & \begin{tabularx}{0.7\textwidth}{X} 因為你問安的聲音一入我耳，我腹裡的胎就歡喜跳動。 \end{tabularx} \\ \\ \relax
1:45 & \begin{tabularx}{0.7\textwidth}{X} 這相信的女子是有福的！因為主對她所說的話都要應驗。」 \end{tabularx} \\ \\ \relax
1:46 & \begin{tabularx}{0.7\textwidth}{X} 馬利亞說：「我心尊主為大； \end{tabularx} \\ \\ \relax
1:47 & \begin{tabularx}{0.7\textwidth}{X} 我靈以神我的救主為樂； \end{tabularx} \\ \\ \relax
1:48 & \begin{tabularx}{0.7\textwidth}{X} 因為他顧念他使女的卑微；從今以後，萬代要稱我有福。 \end{tabularx} \\ \\ \relax
1:49 & \begin{tabularx}{0.7\textwidth}{X} 因為那有權能的為我做了大事；他的名是聖的。 \end{tabularx} \\ \\ \relax
1:50 & \begin{tabularx}{0.7\textwidth}{X} 他憐憫敬畏他的人，直到世世代代。 \end{tabularx} \\ \\ \relax
1:51 & \begin{tabularx}{0.7\textwidth}{X} 他用膀臂施展大能；他趕散心裡妄想的狂傲人。 \end{tabularx} \\ \\ \relax
1:52 & \begin{tabularx}{0.7\textwidth}{X} 他叫有權柄的失位，叫卑賤的升高。 \end{tabularx} \\ \\ \relax
1:53 & \begin{tabularx}{0.7\textwidth}{X} 他叫飢餓的飽餐美食，叫富足的空手回去。 \end{tabularx} \\ \\ \relax
1:54 & \begin{tabularx}{0.7\textwidth}{X} 他扶助了他的僕人以色列，不忘記施憐憫， \end{tabularx} \\ \\ \relax
1:55 & \begin{tabularx}{0.7\textwidth}{X} 正如他對我們的列祖說過，『憐憫亞伯拉罕和他的後裔，直到永遠。』」 \end{tabularx} \\ \\
[1ex]
\hline
\hline
\end{longtable}
講題:上帝的逆道與人間的常道　God’s subversion to the World

講員:謝任生牧師

經文:路加福音1:28-55

場合:香港中文大學崇基學院禮拜堂

日期:2017年8月6日

版權所有,請勿翻印

引   言

相信我們都會同意 : 聖經是上帝的話語,我們應當明白聖經所言說的上帝之道.可是:

我們對聖經的理解,很多時是以人自己當下的文化處境去解釋昔日經文的內容,文化與社會處境,

我們更會比較少尋求理解他人的理解來理解經文的原意.

我的意思是說,我們對剛才所誦讀的聖經經文,應當怎樣去感悟?感悟不是單單看為理解,而更同時是努力去感悟他人的感悟.現在就讓我們從馬利亞身上去感悟她本人的感悟.就是怎樣從人間常道的轉向上帝的逆道的感悟.

一/     馬利亞的三次人間常道與上帝的逆道

第一件上帝的逆道是天使突然的出現

在這段經文中,馬利亞的三次遇上三次的人間常道與上帝逆道的相遇.就以今天早上我所選用的經文來說,我們對馬利亞的認識很少能以理解馬利亞的理解來感悟這段經文文意.很多時我們多會以羅漫蒂克的處境來美化馬利亞.形容馬利亞的溫柔,順服和美麗.這當然是.只是我們每年的聖誕節的佈置,都是很羅漫蒂克化的.這種心態就是一種人間的常道,以為耶穌基督的降生一定是人間的美境,即是讀者以自己的處境去美化原來的處境.可是,我們不會明白約瑟,馬利亞和耶穌都是著日的難民,是基層人又是貧窮人.這處境是值得我們關注的.例如

1:28 天使進去,對他說,蒙大恩的女子,我問你安,主和你同在了.

1:29 馬利亞因這話就很驚慌,又反復思想這樣問安是甚麼意思.

馬利亞的驚慌不是因為心裡有愧的驚慌,按原來的意思,是指到人因敬畏上帝而表現的驚慌.另外,在馬利亞因天使的話與她突然感到很驚慌,反復思想之間存在有一段短暫時間的空間,當中可能發生出一個問題.就是馬利亞的反應是一種由人間的常道,人間的理性認識,轉向上帝的逆道,即超越人間理性認識的反應.因為天使向馬利亞的顯現是一種逆道,馬利亞的驚慌是人間的一種常道.這是說,人間的常道是指到人以為是真理的角度來理解上帝超越人理性認識的逆道.

按當時以色列人的信念,每個母親都期望自己所生的兒子就是那位以色列民族所期盼的彌賽亞.馬利亞也必然.但是,在歷史的長河中和猶太人女子千萬人中,馬利亞又怎會想到以色列民族所期盼的彌賽亞,竟由她而生?這不竟是人間的逆道而非人間的常道.

第二件上帝的逆道就是馬利亞要未婚生子

1:30 天使對他說,馬利亞,不要怕．你 在神面前已經蒙恩了.

1:31 你要懷孕生子,可以給他起名叫耶穌.

1:32 他要為大,稱為至高者的兒子．主神要把他祖大衛的位給他 ,

1:33 他要作雅各家的王,直到永遠．他的國也沒有窮盡.

1:34 馬利亞對天使說,我沒有出嫁,怎麼有這事呢.

我們必須明白 : 今天的我們奉子成婚沒有甚麼大不了的事.可是,對昔日的猶太人來說,未婚生子是件大惡.可以用石頭擲死.馬利亞是位敬虔的童女,她絕對不能接受這樣的事發生在自己身上,因為這不是人間的常道.

第三次上帝的逆道是要馬利亞的兒子去顛覆羅馬帝國的政權

1:46 馬利亞說,我心尊主為大,

1:47 我靈以 神我的救主為樂.

1:48 因為他顧念他使女的卑微．從今以 後,萬代要稱我有福.

1:49 那有權能的為我成就了大事．他的名為聖.

1:50 他憐憫敬畏他的人,直到世世代代.

1:51 他用膀臂施展大能．那狂傲的人,正心裡妄想,就被他趕散了.

1:52 他叫有權柄的失位,叫卑賤的升高．

1:53 叫飢餓的得飽美食,叫富足的空手回去.

1:54 他扶助了他的僕人以色列,

1:55 為要記念亞伯拉罕和他的後裔,施憐憫,直到永遠,正如從前對我們列祖所說的話.

我們必須明白馬利亞當時的歷史背景 :  當時猶太人正是處於羅馬盛世(Glory of The Roman Empire )的時期,不但軍力強盛,更是太平盛勢.而另一方面,也是猶太人蠢蠢欲動要推翻羅馬政權,例如在約新約和舊約之間,猶太人有被稱為四百年的黑暗時期的馬加比戰爭.但是在這一段的經文中,天使告訴馬利亞,她生下來的兒子耶穌要

用膀臂施展大能．那狂傲的人,正心裡妄想,就被他趕散了,

叫有權柄的失位,叫卑賤的升高,

叫飢餓的得飽美食,叫富足的空手回去.

只是對馬利亞一個16-17歲的童女來說,這真是不可能又不應該的逆道.但是,馬利亞順服上帝的逆道,就是上帝把不可能又不應該成為可能又應該,成就了上帝對世人救贖的終極計劃.這要從人間中常道中的上帝稱謂開始.上帝,甚麼是上帝?或正確的說,「上帝」這名字又是甚麼意思?

二/          上帝,請問尊姓大名?

我們當如何去認識這位常行逆道的上帝?老子說 : 「道可道,非恆道(非常道).名可名,非恆名(非常名)」.上帝是不可知,是有知又是無知的.就是說 : 上帝是逆道的上帝,是指到上帝的作為是人間不能看透的,而且往往是非人的理性可以明白的,更有時是會逆人間之道而行.人間的常道是指到人常以為自己的理性和信念作為真理的指向,並常常堅持指之為真理,但卻不知道這並不一定就是上帝的逆道,而只是人間的常道.既是人間能言說出來的所謂真理,那不是上帝的逆道.只是我們多會以人間的常道來看上帝的逆道,所以構成在基督教派間出現不少爭議,甚至是把人對上帝的認識視為獨家總代理.林語堂把非常道譯為 the Absolut Tao, 就以人對上帝的稱號來看,我們稱之為上帝的上帝,究竟「上帝」這兩個字是指到上帝是甚麼?還是指到上帝是誰?如果是指到上帝是甚麼,那麼我們就無權去問這個問題,因為上帝不是人可以描述出來的.如果是指到上帝是誰,人更是非要小心不可,因為人間所描述的這位上帝,那只是人間的常道立論的一個稱號,絕不能代替上帝本身.人當怎樣稱號上帝本身的名字?這是一個逆道的問題.人怎能詢問上帝的知字?事實上,上帝是無名又有名,是存有又不存有,因為上帝本就是超越存有與不存有的.這就是逆道的展示.

從中國古文化來說,「帝」這個字不是一個專有名詞(Proper Noun),而是一個普通名詞(Common Noun).在中國古文化中,我們有上帝,有中帝,也有下帝.中國詩篇說,「帝即上帝也,順帝之則,順天之道也.」查實中國古代對上帝有相當多的認識,只是到了今天,當我們說「上帝」這個名字時已忘記了中國文化中對上帝的認識.按 R. H. Mathews 在 An Examination of the Terms Used for Translating the Word  ‘God’ into Chinese 一書指出,「上帝」一詞是由「上」與「帝」二字合成.「上」原是由「二」而來,是指及無上的一位(the exalted position of one who is above all others),而「帝」一字原是出自「朿」,是指人束柴獻祭的形象.可惜「上帝」此字的語音(phonetic)己失傳.R. H. Mathews 指出,「這個被讀為『帝』音的音符是不重要的,因為「上帝」二字的古音經已失傳.[1]

如此,我們當怎樣去呼喚這位萬靈萬物的創生者?這也是希伯來人的問題,這就是上帝逆道的基礎.按南懷瑾的解說,常道就是指到人間常規之道.[2]  可是,人間所認為的可道的道不竟不是永恆的,即不是上帝的逆道.因為人不是永恆的,而更是有限的.神學家奧特警告我們說,

給上帝下定義是不可能的.「定-義」意味「限定」,而上帝是無法限定的.定義需要一個更高的類概念.但在上帝之上不存在任何更高的概念.….. 但由於上帝是無法‘定義的,我們永遠不會成功地說出上帝是「甚麼」.上帝是「誰」,我們只能以這種方式表達,即我們解釋我們對上帝或我們以為對上帝有何看法.[3]

另外,神學家庫薩(Cusa) 指出,

我知道,凡是我所知道的都不是上帝,凡是我概括的都不與上帝相似, 毋寧說上帝超越了這些東西.…..上帝不是無,因為無本身還有[無]這個名稱.…..沒有人能夠把握上帝之大,他始終是不可言說的.…..上帝既不是無,不是不存在,也不是既存在又不存在,而是存在與不存在的一切本原的源泉和起源.…..不存在並沒有一個不存在的本原,而是有一個存在的本原.因為,不存在為了存在而需要一個本原.所以,有一個不存在的本原,因為,沒有這一本原,就沒有不存在自身.…..上帝無限地超越於任何被我們理解和稱道為真理的東西之先.…..當我們說這就是他的名稱時,說的並不對,但說的也不錯.因為,說這就是他的名稱,也並沒有錯.我們說的也不是既對又錯,因為,它的單純性超越於一切可稱道的和不可稱道的事物之先. [4]

那是說,上帝雖然是可以被人解釋,但卻不就是上帝本身,因為上帝是不可被人命名的.必須指出的是 : 人所解釋的上帝只不過是有關於上帝的事(about God),而非上帝本身(God Be-Self).我們今天所論說的上帝,只不過是人間常道中所論述的上帝.必須指出的是 : 我們人間一切的言語都不足夠去論述上帝.如此,我們就必須以上帝的逆道來感悟上帝,而非人間的理性來分析上帝.這是修道院的修道意義,即修正人間理性認識上帝的一條道路,以心性來面向上帝.

神學家Karl Barth曾經說過:「我們應該明確兩點:我們的應該和我們的不能夠,並由此敬仰上帝.」這句話正是叫我們把神學與信仰連結起來,並且叫我們可以放心的去談論上帝,因為「敬仰」是宗教的本質,也是神學的基礎.所以,人只對上帝有信仰的知識而沒有對上帝敬畏的智慧,這不是認識上帝.嚴格來說,是對上帝的一種褻瀆,因為我們不是先有信仰才有敬仰,乃是先有敬仰才去對上帝有信仰.例如我們的悔改信主,起因往往不是因為我們在理性上對信仰有所灼見真知,而是因為我們知罪悔改歸向上帝,從而敬仰上帝,信仰上帝.這是十分重要的分別.

三/          保羅以上帝的逆道來認信上帝

那麼,我們當怎樣以心性來認識上帝?這是保羅的逆道,而非用他的知識學問來認識上帝.即是人應當以耶穌基督為心的悟道.耶穌基督對父上帝的認識不是在乎系統神學的理性,而是在乎祂對父上帝的認信與關係.這也是保羅的認信,即凡事都在耶穌基督裡.如果我們稍為留意保羅所寫的新約書信,就必定會發現他常常說在口邊的一句話,就是「在基督裡,在耶穌基督裡,或是在基督耶穌裡」.為何保羅要如此說話?因為他知道他所論述的上帝不是人間言語所能言說得出的上帝.人只有透過道成人身的耶穌基督,曾經在人間存活過的耶穌基督,才能認識和認信上帝.宋泉盛曾經說過 :

我們往往像裁縫師一樣剪裁上帝,使他適合我們有限的想法.我們所裁製的上帝肖像往往強烈地反映著我們文化的偏見和宗教經驗.作為人,我們還能怎樣?是的,我們無法避免這樣做,可是我們必須謹記心中,上帝自身跟我們認知的上帝不是等同的.一旦我們忘記這一點,(事實上我們多麼容易忘記啊!)一旦我們把兩者等同起來,那麼我們就犯了保羅所斥責的偶像崇拜.…從這個觀點來看,我們豈不當好好反省我們從傳統信仰所認識的上帝到底像不像上帝自身? (宋泉盛,1994, P.64)

那麼,我們當怎樣認識上帝?Marcus J. Borg  引自Stephen Mitchell   The  God We  never  Know  論到十誡時引介一個故事 :

人人都知道在西乃山下所發生的故事,但是卻很少人知道在山上發生了甚麼.其實,這是在所難免的,因為我們的眼睛不能看見自己的眼睛,(正如)方程式(是)不能自我證明(的).不過,我們有些聖者曾告訴我們在西乃山上發生過怎麼樣的事 〔他們為了能有所言說,所以都上了山〕.

拉比利末說,「上帝給摩西兩項選擇,一是接受十誡石塊,一是面對面見上帝.摩西知道面對面見上帝會死,就像人面對鏡子而無影像可見的可怕.」

拉比以斯拉說,「摩西真的有接受了十誡,但只是他取了第一誡,因為其他九誡都在沉默中消失了,像所有的顏色在光譜中變成白色的一樣.」

拉比加馬列說,「摩西只取了第一誡中的第一句: 我是不可命名的(I am the Unnamable).」

拉比以哈蘭說,「摩西在山上所見的,就是他在荊棘中所見的一樣.他只得到一個訊息:我是.」

拉比撒母耳說,「不是,那不可命名者只向摩西聽耳語了一個字:我!」

拉比約西說,「在聖者自稱的〔我〕是anokhi : aleph-nun-kaph-yod.所以,摩西所接受的只是這個〔我〕字的首個字母罷了.」

「但是,aleph 是個啞音呀!」

拉比約西說,「正是如此.」那是說,上帝是不可被道言的.

上帝呀!你真係冇得頂呀!

上帝是不可被道言的.所以,神學家Martin Buber 說,有關耶和華一名的最滿意解釋應當是猶太人西閃族對耶和華一名的稱呼”Yah,Yahu, or Yahuvah,”.這句話大概的意思是「祂,這位,就是這位,呀,祂呀!」馬丁布伯認為,這種基本的聲音明顯是西閃族部落的語言,含意著對神一種不可命名的而又神祕,興奮的呼喚.「aleph 是個啞音呀!」這句話正道出了上帝的名字,就是人不能言說出來的.在我年青時我參加過一間教會的青年團契.當中曾發生一件非常不幸的事,就是一個赤貧青年弟兄之死.……  這位弟兄有次被邀請在團契中領禱.只是他剛信主不久,不懂得一些人如水流雲般說出很屬靈的文詞禱語.但又不好意思推卻.他只好厚著頭皮站起來祈禱.只是他一開口卻引來全場大笑.他怎樣禱告的呢?他說 : 上帝呀!上帝,你真係冇得頂呀!這是一次逆道的呼喚上帝!

別問我是誰

讓我們細心想一想,這位弟兄的禱告豈不正是比不少神學家的或教會資深領袖更為真誠,更為準確,更為有意義嗎?所以,上帝的名字就必須是一個啞音,所以猶太人稱上帝為耶和華,耶和華的意思就是「我就是我」(I AM that I AM ).意思是說,我是不能被命名的.我就是我!上帝就是那位真正的別問我是誰.這就是上帝.因為祂是一位逆道的上帝.不能被局限在人間的言詞和理性之中.但是,我們又可以稱謂上帝為上帝,只要我們透過耶穌基督才能正面地認識上帝.換言之,基督教之稱為基督教,不是因為基督教是一種宗教,而是因為耶穌基督的道成人身的原故.沒有耶穌基督,我們就不能直接和具體的認識這位上帝.這就是為何我們禱告後通常要奉耶穌基督之名而求,祈求耶穌基督臨在我們中間,體驗上帝的臨格!前幾天有位姊妹問我 : 假如我當下站在上帝面前,我會怎樣對上帝說出第一句話.我想了一想.我說 : 我不知道怎樣對上帝說出第一句話,但我知道我會以我全身的力量緊擁著上帝.因為能擁抱上帝,一切的話都是多餘的了.

四/          淺描與深描上帝

我們可以知道很多聖經的知識,但就不一定會篤信踐行去尋找上帝與我們個人切身的關係.人若沒有與耶穌基督的生命連繫,我們對上帝的認識只能是認識有關上帝的上帝,是上帝透過人的追尋所投射下來的影子而己.人仍然會認識有關上帝的事情,但卻不是上帝本身.今天不少基督徒時會感到靈性低潮,甚至放棄信仰,就是不但我們對上帝的認識不足,更是對上帝認信的不足.導致這問題當然原恩很多,但其中兩個主因就是我們缺乏如何體會上帝聖言與我們個人的 我‧你 關係.在我們的信仰中,我們會常常把上帝設定為一個他.意思是說,我們把上帝看為一個思考對象,甚至當我們在當需要上帝,向上帝禱告時,也多是把上帝看為有求必應的一尊信仰的對象.然而,我們基本上並未全然的認識到上帝是誰,縱使我們會稱這個「他」為「你」來禱告.但心態上仍然是一個他的上帝.

另一個主因是我們多只在自己的表面思維去讀聖經,即只是以淺描(thin description)而非深描來讀聖經.即是以字面來讀聖經而非深入他人的處境來讀聖經.那是說,讀聖經多是只會從字的平面來讀,是一種人條直線的思考(linear thinking,較少以橫向思維(lateral thinking)及深描(deep description)來讀.

橫向思維是指到我們看一件事物或是一種處境,是從相關方面去思考,而深描思考是指到我們必須要找出事情的原處境來觀察,是去理解他人的理解.[5] 而去與上帝建立我‧你 關係.人若只會從字的平面來讀聖經往往會造成不少基督徒靈性生活感到枯橾,未能進入與上帝活潑常新的上帝逆道的關係.這是為何有不少信徒信主越久,對信仰的認真性,對聖經的渴慕性,對日常生活中的屬靈觸覺性,越會漸漸感到遲鈍,乏力,或是失去興趣.這樣現象可以稱之為神性維度(divine dimension)的缺席.路加福音這一章的經文是不能單單只從字面的常道來理解,而是需要從橫向思維,深描和神性維度來理解.那是說,我們必須把自己的感情去投入馬利亞的處境中.孫寶玲牧師在《逆道》一書中說 :

真正凝視歷史 - 至少是聖經歷史 - 的讀者,應能體會歷史是由連串的「不可能」和「不應該」所組成.這許多的「不可能」和「不應該」,說明生命不是自然的因果邏輯可以完全定義和規限,更不是任何意識形態可以說了就算的.[6]

同樣,在我們人生中總不免有不少的「不可能」和「不應該」.特別是當我們凝視那「不可能」和「不應該」的扭曲苦痛臉容,質疑那自以為是的「次序」和結論,驀然聽見生命的主宰對我們說 :「為甚麼不可能?怎樣不應該?」,我們就會驚覺 : 上帝在我們每個人的生命中,從來沒有停止過逆轉的作為 :原來,歷史並非由客觀中立抽離的記錄,而是要表達和顯現信念,憧憬與價值的盛載.這就是希伯來書十一章所述及的信心問題.馬利亞之所以被上帝所揀選,就是她能放下人間的常道,去以信心順服上帝的逆道.

盧雲神父不只是一位靈修學學者,他同時也是一位神學家.他能夠把神學和靈性結合一起,並且體會到神學的終極意義是甚麼.他說 :

「神學」(theology) 的本意是「在禱告中與上帝聯合」.現在神學成為了一門學術性的科目,跟其他科目看齊,並且神學家時常覺得難以禱告.未來的基督徒領袖需要重拾神學教導中的默觀禱告層面,好叫自己所講的每一句話,每一個勸導,每一個策略都是出自一顆與上帝有親密關係的心,這是非常重要的.….. 他們的領袖觀必須紮根在道成肉身的道-耶穌-建立恆久親密的關係上.他們需要在耶穌身上尋找自己所講的說話,勸告和引導的源頭.… 要使基督徒領導真的有果效,基督徒領袖必須要從一個道德者轉為默觀者.… 我們最大的試棎莫過於認為權力就是傳福音的有效工具.…這些人每每宣稱自己是貧窮和無權勢的耶穌的跟隨者.…但基督徒領袖的道路並非向上爬而是向下爬,直達十字架為止. [7]

我相信盧雲神父並非要為難從事神學和獻身唸神學的人.只是他看出今日神學的一個危機,就是只問人間的常道,忽視上帝的逆道.盧雲進一步解釋說,… 未來的基督徒領袖必須是神學家,認識上帝的心懷.在許多似乎是隨機發生的事情上,他們要藉著禱告,研究,仔細分析等訓練,來表明上帝救贖大工的神性事件.思想未來的基督徒領袖,使我深信我們需要神學上的領導.要實現這個構想,必定要從神學院和修道院開始. [8]

可是,一個認真與上帝同行的人,是會對日常生活的大小事情或人物,甚至是自己,都會看成滿有意義與靈性修為的好機會.杜理安說過 :基督徒內在生命的第一步,是(要)在日常生活中將耶穌基督的十字架成為一個實實在在的經驗.

我們的生存,及其一切困難和悲哀的元素,在被釘十字架的耶穌基督身上,都成為有意義的並有作用的.[9]

以上引文的意思指出,我們可以在自己的日常生活中,透過上帝的創生造物,以及人與人之關係中所存有上帝的神性維度,從上帝創生論中轉向耶穌基督的救贖論.這應是一個認真信仰的基督徒的信仰與生活.這就是為何我們每當禱告完結時要奉耶穌基督的知而求.這句話的原意就是要請耶穌基督在我們當中作證,為中保.因為人若不藉著耶穌基督,沒有人能到父那裡去.

不是總結的總結

請容許我對講道與聽道的一點感受作為總結.今天,我們是否只以人間的常道來聽道和崇拜?還是,我們更懂得以上帝的逆道來聽道和崇拜?我的意思是說 : 我們不但要以一個理性的腦袋來聽道和崇拜,而同時必須要以一個敬虔的心靈來與上帝相遇.同樣,我們身為教牧的人,更當是如此.這也可以沙士比亞的一個問題 :  To Be or Not To Be 的問題.我們的生命也當是如此! 我的老師崇基神學院的戴智民博士曾經教導我們,一小時的講道要用十小時的準講.我不能青出於藍,所以我這篇講章接近用了100小時的準備.不過,要說明一點,我還是不很滿意呢!多謝大家海諒.

[1]  見「R. H. Mathews , An Examination of the Terms Used for Translating the Word  ‘God’ into Chinese

[2] 南懷瑾,《老子他說》.\url{http://sns.91ddcc.com/t/71879}.取閱日期 11-07-2017.

[3]  奧特,1990,P.30-31.

[4]  Cusa,1994,頁14-20.

[5]

[6] 孫寶玲著,《逆道》.(香港 : 德慧文化圖書有限公司,2016年二版),頁7.

[7]  盧雲.《奉耶穌的名》,頁32.

[8]  盧雲.《奉耶穌的名》,頁65-66.

[9] 杜理安著,周天和譯.《現代人與靈性生活》(香港:基督教文藝出版社,1986年三版),頁1.



\newpage

\section{哀莫大於心死}
\label{sec:377}
\textbf{林豪恩先生 - 哀莫大於心死}
\newline
\newline
\hyperref[sec:376]{< < < PREV SERMON < < <}
~
\hyperref[sec:toc]{[返目錄]}
~
\hyperref[sec:378]{> > > NEXT SERMON > > >}


列王記上 19:9-18
\newline
\begin{longtable}{cl}
\hline
\hline
章節 & 經文 (和合本修訂版)\\
\hline
19:9 & \begin{tabularx}{0.7\textwidth}{X} 他在那裡進了一個洞，在洞中過夜。看哪，耶和華的話臨到他，說：「以利亞，你在這裡做甚麼？」 \end{tabularx} \\ \\ \relax
19:10 & \begin{tabularx}{0.7\textwidth}{X} 他說：「我為耶和華－萬軍之神大發熱心，因為以色列人背棄了你的約，毀壞了你的壇，用刀殺了你的先知，只剩下我一人，他們還要追殺我。」 \end{tabularx} \\ \\ \relax
19:11 & \begin{tabularx}{0.7\textwidth}{X} 耶和華說：「你出來站在山上，在耶和華面前。」看哪，耶和華從那裡經過。在耶和華面前有烈風大作，山崩石裂，耶和華卻不在風中；風後有地震，耶和華也不在其中； \end{tabularx} \\ \\ \relax
19:12 & \begin{tabularx}{0.7\textwidth}{X} 地震後有火，耶和華也不在火中；火以後，有輕微細小的聲音。 \end{tabularx} \\ \\ \relax
19:13 & \begin{tabularx}{0.7\textwidth}{X} 以利亞聽見，就用外衣蒙臉，出來站在洞口。聽啊，有聲音向他說：「以利亞，你在這裡做甚麼？」 \end{tabularx} \\ \\ \relax
19:14 & \begin{tabularx}{0.7\textwidth}{X} 他說：「我為耶和華－萬軍之神大發熱心，因為以色列人背棄了你的約，毀壞了你的壇，用刀殺了你的先知，只剩下我一人，他們還要追殺我。」 \end{tabularx} \\ \\ \relax
19:15 & \begin{tabularx}{0.7\textwidth}{X} 耶和華對他說：「去吧，從原路回去，往大馬士革的曠野去。到了那裡，你要膏哈薛作亞蘭王， \end{tabularx} \\ \\ \relax
19:16 & \begin{tabularx}{0.7\textwidth}{X} 又膏寧示的孫子耶戶作以色列王，並膏亞伯‧米何拉人沙法的兒子以利沙作先知接續你。 \end{tabularx} \\ \\ \relax
19:17 & \begin{tabularx}{0.7\textwidth}{X} 將來逃過哈薛之刀的，必被耶戶所殺；逃過耶戶之刀的，必被以利沙所殺。 \end{tabularx} \\ \\ \relax
19:18 & \begin{tabularx}{0.7\textwidth}{X} 但我在以色列中留下七千人，是未曾向巴力屈膝，未曾親吻巴力的。」 \end{tabularx} \\ \\
[1ex]
\hline
\hline
\end{longtable}
講題:哀莫大於心死　No Grief Greater than Utter Despair

經文:列王記上19章9-18節

講員:林豪恩先生

場合:香港中文大學崇基學院禮拜堂

日期:2017年8月13日

一．引言

今天的舊約經課,又是我們熟識的以利亞先知的故事.

近代其中一種常見的解讀,認為以利亞患了憂鬱症,其中一篇講章如此說:

「聖經有段經文,所敘述的光景,能用來了解抑鬱症的成因和症狀,並且,其中涵蓋了得醫治的過程,方法,指引.見列王記上十九章1-8節:(即是今天經課的上文)

「亞哈將以利亞一切所行的和他用刀殺眾先知的事,都告訴耶洗別.耶洗別就差遣人去見以利亞,告訴他說:「明日約在這時候,我若不使你的性命像那些人的性命一樣,願神明重重地降罰於我.」以利亞見這光景,就起來逃命,到了猶大的別是巴,將僕人留在那裡.自己在曠野走了一日的路程,來到一棵羅騰樹下；就坐在那裡求死說:「耶和華啊,罷了!求你取我的性命,因為我不勝於我的列祖.」他就躺在羅騰樹下睡著了.有一個天使拍他說:「起來吃吧!」他觀看,見頭旁有一瓶水與炭火燒的餅,他就吃了喝了,仍然躺下.耶和華的使者第二次來拍他說:『起來吃吧!因為你當走的路甚遠.』他就起來吃了喝了,仗著這飲食的力,走了四十晝夜,來到 神的山,就是何烈山.」

講者指出,「此時,以利亞罹患了重度的憂鬱症,失去人生的盼望,渴求死去.」(「以利亞的憂鬱症」: 趙鏞基牧師)

這種解讀的應用就是:從以利亞的經歷看醫治抑鬱症的方法:

1．有天使來安撫以利亞,表示抑鬱症患者需要別人安撫；

2．天使供給以利亞水和餅,表示抑鬱症患者需要飲食；

3．以利亞行路四十天,表示抑鬱症患者需要運動；

4．以利亞面見上帝,表示抑鬱症患者需要靈修.

如果抑鬱症患者從安撫,飲食,運動,靈修這四項得到幫助,也是好事.不過,這四項很可能不只是抑鬱症患者的需要,也應該是所有人的需要.

以抑鬱症來解讀以利亞的故事也好像說得通.不過,以精神病患來解釋某人的行為和表現,也可以是危機:其一是把現象個人化,歸咎於當事人的不濟,暗示當事人自作自受.按此歸因,解決方法就是要求當事人自我改善,奮發自強；第二,更可怕的是把精神病患標籤運用作為對付異己的方法,民間的例子有:「他是瘋子,不用理會他說的話!」國家層面的例子有:「蘇聯境內沒有不贊同蘇聯共產黨的人,沒有不贊同蘇聯共產主義制度的人,而只有精神病患者.」記得是誰說的嗎?說這話的人是前蘇共領袖赫魯曉夫,所以,他將那些敢於對蘇共表示不滿的異議人士送進精神病院.

以利亞是精神病／或者我們現在稱為情緒病患者嗎?他的故事只提供了似是而非的「克服抑鬱症四步曲」嗎?讓我們再細閱以利亞的故事,尋找其他解讀的角度,及對此時此地的讀者的啟迪.

二．經文

19:9他在那裡進了一個洞,在洞中過夜.看哪,耶和華的話臨到他,說:「以利亞,你在這裡做甚麼?」

19:10他說:「我為耶和華－萬軍之上帝大發熱心,因為以色列人背棄了你的約,毀壞了你的壇,用刀殺了你的先知,只剩下我一人,他們還要追殺我.」

19:11耶和華說:「你出來站在山上,在耶和華面前.」看哪,耶和華從那裡經過.在耶和華面前有烈風大作,山崩石裂,耶和華卻不在風中；風後有地震,耶和華也不在其中；

19:12地震後有火,耶和華也不在火中；火以後,有輕微細小的聲音.

19:13以利亞聽見,就用外衣蒙臉,出來站在洞口.聽啊,有聲音向他說:「以利亞,你在這裡做甚麼?」

19:14他說:「我為耶和華－萬軍之上帝大發熱心,因為以色列人背棄了你的約,毀壞了你的壇,用刀殺了你的先知,只剩下我一人,他們還要追殺我.」

19:15耶和華對他說:「去吧,從原路回去,往大馬士革的曠野去.到了那裡,你要膏哈薛作亞蘭王,

19:16又膏寧示的孫子耶戶作以色列王,並膏亞伯‧米何拉人沙法的兒子以利沙作先知接續你.

19:17將來逃過哈薛之刀的,必被耶戶所殺；逃過耶戶之刀的,必被以利沙所殺.

19:18但我在以色列中留下七千人,是未曾向巴力屈膝,未曾親吻巴力的.」

三.釋經與應用

1．以利亞是和君王打交道的先知

以利亞是一位來無影去無縱的先知.按列王記敘述,他彷彿從天而降,臨到當時的以色列王亞哈的面前,對他說:「我指著所事奉永生耶和華－以色列的上帝起誓,這幾年我若不禱告,必不降露,不下雨.」(王上17:1)請留意兩件事,第一,以利亞在舞台上的第一位對手,不是泛泛之輩,而是以色列王亞哈；第二,他對亞哈說的話,並不是個人私事,而是公共事務,下雨與否,在當時的農業社會,關乎響整個國家的福祉.可以這樣說,以利亞先知出場的舞台,是政治公共舞台.

17章餘下的內容(2-24節)記載了關於以利亞所行的神蹟,例如他在撒勒法的寡婦家中,把只夠一天的食物變成「長食長有」,又使寡婦的兒子死而復生.這些神奇事件,是要說明以利亞是先知.這先知的身份,在17章結束時透過寡婦的口說了出來:「現在我知道你是神人,耶和華藉你口所說的話是真的.」

18章回到17章1節的場景,即是先知以利亞與以色列王亞哈的對手戲.其中最戲劇性的是以利亞要求亞哈安排他和巴力先知比試,結果他一人戰勝450個巴力先知.讓我們重新思考這場比試,以利亞的對手是誰呢?比武場上,以利亞的對手是巴力先知,然而,按作者的敘述,巴力先知只是丑角而已,難道讀者會假設耶和華的先知會被巴力先知打敗嗎?因此,這場比試的結果早已在預料之中.那麼,在這舞台上,誰才是以利亞的真正對手呢?

2．以利亞是在政治上被挫敗的先知

以利亞挑戰巴力先知這場戲是做給誰看的呢?他擊敗巴力先知所為何事呢?請聽以利亞的說話:

18:21以利亞前來對眾民說:「你們心持兩意要到幾時呢?若耶和華是上帝,就當順從耶和華；若巴力是上帝,就當順從巴力.」眾民一言不答.

當以利亞以耶和華的先知的身分戰勝巴力先知之後,在場的眾民改變了態度:

18:39眾民看見了,就俯伏在地,說:「耶和華是上帝!耶和華是上帝!」

眾民的態度改變了,顯示以利亞不單嬴了巴力先知,還嬴了民意.不幸的是,羸了巴力先知,似乎並不能改變目標人物；民意,對這位目標人物沒有多大影響.這目標人物是誰呢?呼之欲出,就是以色列王亞哈,他才是以利亞的真正對手,是以利亞要改變的對象.比試結束之後,劇情又回到以利亞和亞哈的對手戲.以利亞為甚麼想改變亞哈呢?因為他是以色列王,以利亞盼望透過改變他來改變國家的權力結構.到底以利亞盼望改變甚麼呢?

16:30暗利的兒子亞哈行耶和華眼中看為惡的事,比他以前的列王更甚,

16:31犯了尼八的兒子耶羅波安所犯的罪；他還以為輕,又娶了西頓王謁巴力的女兒耶洗別為妻,去事奉敬拜巴力,

16:32在撒馬利亞建造巴力的廟,在廟裡為巴力築壇.

以色列的巴力崇拜是誰帶來的,是誰在背後支撐呢?就是王后耶洗別.如果以色列王亞哈能夠限制王后耶洗別的影響力和權力,那麼,以色列才能夠清除巴力崇拜,而回歸崇拜耶和華.

讓我們再問,以利亞打敗巴力先知,嬴了民意,他能夠令以色列王亞哈去限制王后耶洗別的權力,改變以色列的權力結構嗎?

19:1亞哈將以利亞一切所行的和他用刀殺眾先知的事都告訴耶洗別.

亞哈的確把以利亞戰勝巴力先知的事情告訴王后耶洗別,不過,結果卻與以利亞的期望相反:

19:2耶洗別就差遣人去見以利亞,告訴他說:「明日約在這時候,我若不使你的性命像那些人的性命一樣,願神明重重地降罰與我.」

19章3節敘述以利亞的反應:「以利亞見這光景就起來逃命……」.以利亞看見甚麼光景呢?就是19章1-2節所描述的狀況,以色列王亞哈看見以利亞打敗巴力先知,回到王宮,不是削弱王后耶洗洗別的權力,而是向耶洗別報告；耶洗別的權力不但沒有減少,從她在第2節的說話可見,她仍然大權在握,如果不是擁有調軍遣將的權力,也應該擁有指揮秘密警察的權力.

以利亞辛苦奔波,呼風喚雨,英雄式地打敗巴力先知,所作的一切,最後都是徒然.以利亞作為耶和華的先知,已經展示正義與邪惡的分別,呼喚國家改邪歸正,棄惡從善,然而,國家的政治狀況絲亳沒有改變,甚至比他的想像更邪惡；以利亞作為耶和華的先知,已經展示先知的能力,然而,國家的權力結構如巖石一般堅硬,動也不動,甚至比他的想像更頑強.他要說的,都已叫喊了,要做的,都已經奔波了,如果要檢討有甚麼果效,就只能證明自己果然是脆弱的雞蛋,而亞哈的朝廷是堅固的高牆,自己反而落得被追殺的劣境.這堅固的高牆,不為正義所改變,不為民意所改變,不為能力所改變,對以利亞來說,除了挫敗,除了失落,除了消沉,除了絕望,除了求去,還可以有甚麼合適的心情去面對呢?

此情此境,那麼遠,卻又這麼近；似陌生,卻又感共鳴；個別的,卻又是集體的.

3．上帝的回應

3．1．上帝有能力,上帝卻不是能力

19:11耶和華說:「你出來站在山上,在耶和華面前.」看哪,耶和華從那裡經過.在耶和華面前有烈風大作,山崩石裂,耶和華卻不在風中；風後有地震,耶和華也不在其中；

19:12地震後有火,耶和華也不在火中；火以後,有輕微細小的聲音.

以力服人,是世界的遊戲規則.從以利亞單槍匹馬挑戰四百五十位巴力先知的比武,讓人猜想以利亞其實也是這種遊戲規則的信徒.所以當色列的朝廷不為所動,以利亞深陷力有不逮的挫敗之中.以利亞可能陷於是自己作為先知力有不逮,還是差遺自己的上帝力有不逮的苦思.耶和華向以利亞展示推倒高牆,搖動大地的能力；然而,耶和華也向以利亞啟示,上帝有能力,但上帝不是能力本身,上帝的遊戲規則,不是以力服人.力量可以製造山崩石裂的嘈音,可以發出烈風火焰氣勢,嘈音和氣勢,都只是吸引注意力的戲碼,信息,卻在嘈音和氣勢之外,在沉寂之中的「輕微細小的聲音」.作為先知,以利亞應該事奉甚麼呢?以利亞應該以甚麼來衡量事物呢?是「力量The Power」,還是「信息The Word」?約翰福音啟示,上帝不是力量,上帝是道(the Word / Message),這是耶和華希望祂的先知以利亞領悟的,也很可能是上主希望所有跟隨祂的人領悟的.

3．2．邪惡雖頑固,上帝卻仍在

19:15耶和華對他說:「去吧,從原路回去,往大馬士革的曠野去.到了那裡,你要膏哈薛作亞蘭王,

19:16又膏寧示的孫子耶戶作以色列王,並膏亞伯‧米何拉人沙法的兒子以利沙作先知接續你.

19:17將來逃過哈薛之刀的,必被耶戶所殺；逃過耶戶之刀的,必被以利沙所殺.

上帝彷彿對以利亞說,世界不按你的議程行事,不回應你的努力,不為你的熱心所感動,失望,沮喪,挫敗等感受是可以理解的,然而,不就等於黑暗已經戰勝光明,邪惡已經取代正義,巴力已經殺死耶和華.耶和華向以利亞預告君王的更替,政局的變遷,並不代表以利亞的期望會實現,只是告訴他知道,上帝仍在,祂仍掌管.作為耶和華的先知,上帝期望他放開自己的議程,接受上帝的安排,相信上帝的世界比他的世界要大.縱使經歷挫敗和沮喪,感到失望和無能為力,上帝不容祂的先知在生理上或心理上求去,甚至不讓他隱姓埋名,反而差派他回去那個政治舞台,重投「熱廚房」.不過,再登舞台的以利亞,已經不再是從前的以利亞,而是被耶和華更新的以利亞.

3．3．成功不必我在,功成其中有我

19:18但我在以色列中留下七千人,是未曾向巴力屈膝,未曾親吻巴力的.」

以利亞多次說:「我為耶和華－萬軍之上帝大發熱心,因為以色列人背棄了你的約,毀壞了你的壇,用刀殺了你的先知,只剩下我一人,他們還要追殺我.」以利亞以為自己是唯一的正義代表,耶和華的獨家代理,因此難免有捨我其誰的情懷,也難免感到自己的失敗,就代表著大勢已去.耶和華告訴祂的先知,一個人倒下,不代表正義消失；一個人被打敗,不代表世界末日,在他那個時空,還有七千個同樣地堅持著的同路人,只是他未必知道而已.七千,是一個代表數,意即在上帝來說,尚有足夠的人在這路上繼續奔跑.

四.總結

哀莫大於心死,可能是有些積極參與公共事務的朋友走過這幾年累積的心境,又或者是關心社會時事的朋友被最近幾個月發生的事情衝擊而引發的反應；也可能是在家庭關係或工作間遇到「好心沒好報」或「好心著雷劈」的悲憤.有些時候,「天氣的當不似預期」,別人也不接受我們的好意,令人感到沮喪.耶和華教導以利亞,這不是放棄的時候,而是在上帝面對謙卑自省,重整信仰,以更新的心態重新出發的機會.



\newpage

\section{不要奉旨}
\label{sec:378}
\textbf{鄧瑞強博士 - 不要「奉旨」}
\newline
\newline
\hyperref[sec:377]{< < < PREV SERMON < < <}
~
\hyperref[sec:toc]{[返目錄]}
~
\hyperref[sec:379]{> > > NEXT SERMON > > >}


馬太福音 15:21-28
\newline
\begin{longtable}{cl}
\hline
\hline
章節 & 經文 (和合本修訂版)\\
\hline
15:21 & \begin{tabularx}{0.7\textwidth}{X} 耶穌離開那裡，退到推羅、西頓境內。 \end{tabularx} \\ \\ \relax
15:22 & \begin{tabularx}{0.7\textwidth}{X} 有一個迦南婦人從那地方出來，喊著說：「主啊，大衛之子，可憐我！我女兒被鬼纏得很苦。」 \end{tabularx} \\ \\ \relax
15:23 & \begin{tabularx}{0.7\textwidth}{X} 耶穌卻一言不答。門徒進前來，求他說：「這婦人在我們後頭喊叫，請打發她走吧。」 \end{tabularx} \\ \\ \relax
15:24 & \begin{tabularx}{0.7\textwidth}{X} 耶穌回答：「我奉差遣只到以色列家迷失的羊那裡去。」 \end{tabularx} \\ \\ \relax
15:25 & \begin{tabularx}{0.7\textwidth}{X} 那婦人來拜他，說：「主啊，幫幫我！」 \end{tabularx} \\ \\ \relax
15:26 & \begin{tabularx}{0.7\textwidth}{X} 他回答：「拿孩子的餅丟給小狗吃是不妥的。」 \end{tabularx} \\ \\ \relax
15:27 & \begin{tabularx}{0.7\textwidth}{X} 婦人說：「主啊，不錯，可是小狗也吃牠主人桌上掉下來的碎屑。」 \end{tabularx} \\ \\ \relax
15:28 & \begin{tabularx}{0.7\textwidth}{X} 於是耶穌回答她說：「婦人，你的信心很大！照你所要的成全你吧。」從那時起，她的女兒就好了。 \end{tabularx} \\ \\
[1ex]
\hline
\hline
\end{longtable}
講題:不要「奉旨」　Nothing is Taken for Granted

經文:馬太福音15:21-28

講員:鄧瑞強博士

場合:香港中文大學崇基學院禮拜堂

日期:2017年8月20日

各位弟兄姊妹,早安.願你們平安.

早幾日,三名參與社會運動的學生領袖被判入獄,香港社會的撕裂無可避免地擴大.同情他們的群眾,與反對他們的群眾,看來沒有什麼對話的空間.這些學生領袖追求的,是一個還未出現的社會,是一個在盼望中的世界,是一個看來遙遠但仍可實現的夢想.反對他們的,持守著現實的原則.看來手中的麵包,比不知如何才能實現的公義更加真實.或許,眼前的現實,真是我們能得到的最後的東西了.保羅說:「若死人不復活,我們就吃吃喝喝吧!因為明天要死了.」(林前15:32)他的意思是說,若現實就是最後的真實,沒有那種全新的生命的可能,則算了吧,像動物一樣活著吧,反正明天就像動物一樣死去.但是,如果有新的生命的可能呢?如果明天,學生能自由及友愛地學習,大人能有尊嚴地工作和生活,老人能不用憂慮地安享晚年,如果這一切是可能的呢?如果人類自古以來的天國夢真有實現的可能,則我們是否需要為這個美夢而奮鬥?

人,總是要發夢的.現實,總是會叫人的夢破滅的.平凡的人,在夢多次破滅後,便滿足於吃吃喝喝,不再發夢了.有些非凡的人,看見現實不斷地破滅夢想,便努力不懈地,以夢想之力去改造現實.當然,這是漫長的人類工程.

黃之鋒判刑前3小時,他母親發出一公開的帖文,有一段是這樣寫的:「…法庭的判決,估計我們會感到十分失望.之鋒已履行社會服務令,律政司卻鍥而不捨,誓要以言入罪(「重奪公民廣場」好暴力?),扼殺年輕人的滿腔熱誠和理想,對社會的願景和抱負.為何香港墮落如斯,如此對待這一代的孩子?……」

可以想像,這位母親,應曾向神祈禱,求神免去她兒子牢獄的劫難.神沒有回應.神很沉默.其實,神很多時都很沉默.就算是耶穌的祈禱,在十架前夕的祈禱,神都是很沉默的.

神奉旨聽我們的祈禱嗎?沒有奉旨的,神有時很沉默.

社會的撕裂能否修補?年輕人的渴望能否實現?人類的夢能否成真?這一切都是未知之數.沒有什麼是「奉旨」的.在這個幸福不是必然的世界,我們能做什麼?抑或,我們什麼也不能做?

我們看看今日的講道經文:馬太福音十五21－28

這裡講到一位母親,她遇到莫大的困境,對她而言,幸福不是奉旨的.

我們看看她的反應.

太15:21　耶穌離開那裡,退到泰爾,西頓的境內去.

太15:22　有一個迦南婦人,從那地方出來,喊著說:「主啊,大衛的子孫,可憐我!我女兒被鬼附得甚苦.」

太15:23　耶穌卻一言不答.門徒進前來,求他說:「這婦人在我們後頭喊叫,請打發她走吧.」

太15:24　耶穌說:「我奉差遣不過是到以色列家迷失的羊那裡去.」

太15:25　那婦人來拜他,說:「主啊,幫助我!」

太15:26　他回答說:「不好拿兒女的餅丟給狗吃.」

太15:27　婦人說:「主啊,不錯；但是狗也吃牠主人桌子上掉下來的碎渣兒.」

太15:28　耶穌說:「婦人,你的信心是大的!照你所要的,給你成全了吧.」從那時候,她女兒就好了.

馬太福音為這故事,設定一個場景.耶穌往泰爾,西頓的邊境方向去.泰爾,西頓,在以色列的界線外.在耶穌眼中,這兩個城巿算是最敗壞的了.在馬太福音裡,耶穌有個講法:「哥拉汛哪,你有禍了!伯賽大啊,你有禍了!因為在你們中間所行的異能,若行在泰爾,西頓,他們早已披麻蒙灰悔改了.但我告訴你們,當審判的日子,泰爾,西頓所受的,比你們還容易受呢!」(太11:21-22)這即是說,你若不悔改,泰爾,西頓都好過你了.言下之意,泰爾,西頓算是最差的了.按照下文的講法,也按照馬太一貫的神學,耶穌只在以色列地區傳道,所以,耶穌並沒有越過界線,並沒有踏足泰爾,西頓.耶穌在邊界上走.

「有一個迦南婦人,從那地方出來」.很奇怪,馬太用「迦南」這個字眼.同一個故事,在馬可福音,只稱這女人是「希臘人,屬敘利腓尼基族」(可7:26).「迦南人」,在舊約,是以色列人的世仇.用到「迦南」這字眼,就會引發厭惡,仇恨,抗拒的情緒.耶穌在邊界行走,引來了一個「迦南婦人」.一個邊界外的人,闖入了耶穌的界線.

她有苦情,她有所求,她無計可思.

太15:22　有一個迦南婦人,從那地方出來,喊著說:「主啊,大衛的子孫,可憐我!我女兒被鬼附得甚苦.」

女兒被鬼所附,有一強大外力宰制著她的生命,她毫無自由,像在牢獄中,生命被折磨,苦不堪言.母親愛女情深,只能哀求她能找到的唯一希望.

她求耶穌可憐.她稱呼耶穌為「主」.在一個以奴隸為主要生產力的世界,稱呼一個人為「主」,就是成為附屬於他的「奴隸」.為了她的女兒,這個女人「低聲下氣」,懇求耶穌憐憫.她還稱耶穌為「大衛的子孫」,這是以色列人的民族用語,不應是迦南人說的.看來,這女人勇於越過界限,進入以色列人的文化裡.

若你是耶穌,若你有醫治大能,遇到這樣的一個女人,她哀求你,你有何反應?

我想,我們大部分人,會即時應允她的呼求.

經文說:「耶穌卻一言不答.」

她遇到耶穌的沉默.她遇到神的沉默.

如此卑微的一位母親,為著女兒無邊的苦楚來哀求,但耶穌卻一言不答.耶穌不是奉旨要回應我們的哀求的.憐憫人的神,不是奉旨要做點什麼的.

上年(2016年),導演馬田史高西斯(Martin Scorsese)的電影《沉默》(Silence)上畫.講到豐臣秀吉年代,日本禁止天主教的事蹟.很多教徒殉道,神父被迫放棄信仰.神父懇切祈禱,但神沉默.在深深的苦難裡,神父祈禱,他說:「我害怕.你沉默的重量令人畏懼.我祈禱,但我很迷惘.我是否只是向著虛無祈禱?虛無,因為你根本就不存在.」

我們常常祈禱,為自己,為身邊人,為社會,為世界,但常常經驗神的沉默.神在嗎?

這個迦南女人,忍受著耶穌的沉默.

我想,連耶穌的門徒,也忍受不了耶穌的沉默.

太15:23　……門徒進前來,求他說:「這婦人在我們後頭喊叫,請打發她走吧.」

看來,門徒厭煩了這個女人,他們求耶穌做點東西,打發她走.

耶穌的反應有點出乎意外.

太15:24　耶穌說:「我奉差遣不過是到以色列家迷失的羊那裡去.」

耶穌拒絕門徒的懇求.美國總統特朗普說:「America First」.在這裡,耶穌似是說:「Israel Only」.不單是「以色列優先」,更是「唯獨以色列」.耶穌是說,他只幫某些人,他的工作是有特定範圍的,這個女人卻在我的工作範圍以外.

門徒好好歹歹,算是為那女人求情.但就算是門徒的代求,也不能奉旨得到耶穌的應允.耶穌不單沉默,也正面拒絕.

若你求人幫忙,人家卻說:「這在我的能力之外,幫不上忙了.」一般情況下,你大概便放棄了.

這女人沒有放棄.

太15:25　那婦人來拜他,說:「主啊,幫助我!」

那女人「拜」耶穌,她應該是跪在耶穌面前了.耶穌變本加厲地拒絕她,她卻變本加厲地高舉耶穌的神權.她繼續稱呼耶穌為主,繼續求耶穌的幫助.

當上帝沉默,毫無反應時,這女人沒有沉默.當上帝看來拒絕她的所求時,當形勢不利她時,這女人沒有放棄.

當社會在瓦解,當強大的力量扭曲一切,當信徒也埋怨神為何不聽祈禱的時候,有些人,像這女人一樣,竟然仍在敬拜神,仍在呼求神的幫助.他們不顧現實的逆境,他們單單仰望神.

耶穌再進一步的回應,違反了我們對耶穌的一般想像.

太15:26　他回答說:「不好拿兒女的餅丟給狗吃.」

耶穌稱那女人為「狗」.不少學者認為耶穌的講法不合時宜,認為他「大男人主義」,民族主義色彩也太濃.有些學者為了減低耶穌這講法的冒犯性,認為耶穌這裡講的,是「小狗」,是「puppy」,是家裡的寵物.這講法有一定道理,因為聖經原文的用字,真是「小狗」的意思.不過,按照這故事的情節發展,耶穌是有意的,一步一步地加強他對這女人的拒絕的.簡言之,這用字應有羞辱成分.

但我們不必專注於追問:耶穌為何羞辱人,因為這裡的文學體裁,重點明顯不在於耶穌對這女人的羞辱,而在於以看來羞辱的方法去讓這女人達到生命的轉化.耶穌的羞辱,及那女人機智的回答,很似我們禪宗的公案.在公案裡,徒弟問師父問題,師父用棒便打.若你問,那師父為什麼打人,這不是羞辱求學者嗎?若你這樣問,便很難明白這種求道與覺悟的文學體裁了.

我再舉個例子,就是張良的故事.

話說張良,他後來成了漢初的謀臣.有日,在橋上散步,遇到一老翁.老翁故意將鞋掉落橋下,然後叫張良把鞋拾回來.張良先是不悅,但後來想到老人家畢竟年紀大,便幫他將鞋拾回來.誰知那老翁將腳一伸,叫張良替他穿鞋.張良壓住怒氣,才能替老翁把鞋穿好.然後,老翁覺得這年輕人孺子可教,便將絕世的兵法傳授給他.

讀到這些故事,若你只著眼於追問為何那老翁羞辱張良,則便錯失這種故事的用意了.這些只是權宜的機遇,讓人在屈辱中成長.

說回那迦南女人.耶穌說她是狗,不配得到兒女的餅.那女人也回答得妙.

太15:27　婦人說:「主啊,不錯；但是狗也吃牠主人桌子上掉下來的碎渣兒.」

人家說她是狗,她沒有一巴掌打過來.她承認生命的卑微,也仰賴主的恩典.這母親說,我不是要吃主人兒女的餅,我只是吃掉下來的碎渣兒.她沒有要求過於自己應享的權利,她只渴求恩典中的滿足.耶穌拒絕走過界線去接觸她,但她這個回答,是將自己看成是家中的寵物,已成為家中的一員,已走過了耶穌設定的界線而成為耶穌家中的一份子.

看看這個女人,這個母親.耶穌沉默,上帝沉默,天不助她,她沒有「天時」.耶穌拒絕她,她在耶穌的工作範圍外,她沒有「地利」.耶穌貶抑她,羞辱她,她沒有「人和」.她什麼也沒有,沒有任何外力的幫助,孤零零地面對生命的橫逆,面對著身邊人的痛苦.叫天不應,叫地不聞.她有的,是一赤子之心,是溫柔的謙虛,是對所求的事的鍥而不捨,是絕望時的無比希望,是當不斷被拒絕時仍堅持不被拒絕,是當神看來不仁不義時仍堅持神的至仁至義.唯有這種人,當歷史的前景顯得灰暗,而人人逃命時,他能屹然不倒.唯有這種人,當社會崩壞時,他能堅持真理.唯有這種人,當黑暗籠罩大地時,仍然能堅持發光.

最後,我們看看耶穌對她的評語.

太15:28　耶穌說:「婦人,你的信心是大的!照你所要的,給你成全了吧.」從那時候,她女兒就好了.

耶穌讚她說:「你的信心是大的!」這句經文,原文沒有動詞,並且字的次序是:「偉大你的信心」.英文的翻譯,一般是 “Great thy faith”.耶穌稱讚她「偉大」.在馬太福音,這個字從來沒有用來形容一個人的生命特質,單單用於這個女人,「偉大」.

人間的一切幸福,都沒有「奉旨」的,唯那些配稱「偉大」的人物,他們以無比的信心,毅力,謙和,站在天地間.神因著他們的「偉大」,將幸福賜予人間.

希望世上有多幾個這樣的人.希望你是其中的一份子.

但願:榮耀歸於聖父,聖子,聖靈.阿們.



\newpage

\section{毋忘初心}
\label{sec:379}
\textbf{鄧瑞強博士 - 毋忘初心}
\newline
\newline
\hyperref[sec:378]{< < < PREV SERMON < < <}
~
\hyperref[sec:toc]{[返目錄]}
~
\hyperref[sec:380]{> > > NEXT SERMON > > >}


啟示錄 2:1-7
\newline
\begin{longtable}{cl}
\hline
\hline
章節 & 經文 (和合本修訂版)\\
\hline
2:1 & \begin{tabularx}{0.7\textwidth}{X} 「你要寫信給以弗所教會的使者，說：『那右手拿著七顆星，在七個金燈臺中間行走的這樣說： \end{tabularx} \\ \\ \relax
2:2 & \begin{tabularx}{0.7\textwidth}{X} 我知道你的行為、勞碌、忍耐，也知道你不容忍惡人。你也曾察驗那自稱為使徒卻不是使徒的，看出他們是假的。 \end{tabularx} \\ \\ \relax
2:3 & \begin{tabularx}{0.7\textwidth}{X} 你能忍耐，曾為我的名勞苦而不困倦。 \end{tabularx} \\ \\ \relax
2:4 & \begin{tabularx}{0.7\textwidth}{X} 然而，有一件事我要責備你，就是你把起初的愛心拋棄了。 \end{tabularx} \\ \\ \relax
2:5 & \begin{tabularx}{0.7\textwidth}{X} 所以你要回想你是從哪裡墜落的，並且要悔改，做起初所做的工作。你若不悔改，我要到你那裡去，把你的燈臺從原處挪去。 \end{tabularx} \\ \\ \relax
2:6 & \begin{tabularx}{0.7\textwidth}{X} 然而你還有一件可取的事，就是你恨惡尼哥拉派的行為，這種行為也是我所恨惡的。 \end{tabularx} \\ \\ \relax
2:7 & \begin{tabularx}{0.7\textwidth}{X} 凡有耳朵的都應當聽聖靈向眾教會所說的話。得勝的，我必將神樂園中生命樹的果子賜給他吃。』」 \end{tabularx} \\ \\
[1ex]
\hline
\hline
\end{longtable}
講題:毋忘初心　Leave not Your First Love

經文:啟示錄2:1-7

講員:鄧瑞強博士

場合:香港中文大學崇基學院禮拜堂

日期:2017年9月10日

各位弟兄姊妹,早安.願你們平安.

少許題外話.我們教會崇拜每週的讀經,是根據一固定的「經課表」(Lectionary)來編排的,這「經課表」是信義宗,長老宗,聖公宗,天主教等共同採用的.即是說,你若今天參加信義宗教會的崇拜,他們的讀經和我們的讀經,是一樣的.這「經課表」三年一循環,即是說,三年前的今個禮拜,讀的經文和今日讀的經文,是完全一樣的.我在這裡講道,已接近十二年了,即是說,我見證著同一「經課表」差不多完成它在這裡的第四次循環了.在講道的事奉上,我通常的做法,是根據「經課表」編排的經文來講道的.並不是我自己揀選一段經文,講我想講的東西；而是經文揀選我,讓我去傳遞它本身的信息.而講過的經文,我會盡量避開,而盡量選用我未講過的經文.經過「經課表」四次循環之後,很多經文我已講過了.故此,我想用一段時間,講講「經課表」以外的經文.「經課表」很少《啟示錄》的經文,從今天起,我會以《啟示錄》的經文作一系列的講道.當然,去到聖誕節等大節期時,我仍會採用「經課表」的經文的.

今日的講道經文:啟示錄2:1－7

啟2:1　「你要寫信給以弗所教會的使者,說:『那右手拿著七星,在七個金燈臺中間行走的,說:

啟2:2　我知道你的行為,勞碌,忍耐,也知道你不能容隱惡人.你也曾試驗那自稱為使徒卻不是使徒的,看出他們是假的來.

啟2:3　你也能忍耐,曾為我的名勞苦,並不乏倦.

啟2:4　然而有一件事我要責備你,就是你把起初的愛心離棄了.

啟2:5　所以,應當回想你是從哪裡墜落的,並要悔改,行起初所行的事.你若不悔改,我就臨到你那裡,把你的燈臺從原處挪去.

啟2:6　然而你還有一件可取的事,就是你恨惡尼哥拉一黨人的行為,這也是我所恨惡的.

啟2:7　聖靈向眾教會所說的話,凡有耳的,就應當聽!得勝的,我必將神樂園中生命樹的果子賜給他吃.』」

啟示錄,是作者約翰寫給亞細亞七間教會的書信.這七間教會,現在位於土耳其,當時是羅馬帝國的疆土.羅馬擁有強大的軍事力量,又藉著君王崇拜去統一帝國的意識形態,強大而統一,成就了一羅馬盛世.耶穌就是死在這盛世裡.沒有軍事強權,耶穌不至於死.沒有君王崇拜的意識形態,也不會封殺耶穌那種天國王權.在那密不透風的,強大而統一的世界裡,容不下新鮮的空氣.

以弗所,羅馬帝國東面的繁華城巿.位於地中海東邊,是一個港口城巿.得益於地中海的貿易,以弗所的商業很發達.若你現在去以弗所旅遊,通常會參觀古城遺蹟,那裡有一圖書館遺址.這Celsus圖書館是古代世界三大圖書館之一.另一個可能會去的旅遊點,是亞底米女神殿(Temple of Artemis),這神殿被喻為古代世界七大建築奇蹟之一.以弗所,在文化,宗教,經濟上,都是出色的.

在軍力強大的帝國裡,在繁華的異教城巿裡,教會的存在很重要.教會有責任指出暴力的經濟之外,有和平的互愛共享；在皇帝與偶像的敬拜外,有超越的上帝.教會要在這個看似和平,實質暴力；看似繁華,實質空洞的世界裡,講論仁愛的天國.可以想像,當一班黑社會大佬拿著刀劍去「劈友」時,你擋在前面「講福音」,是驚險萬分的.當時的教會,正是如此,忍受著沉重的十架.

啟2:1　你要寫信給以弗所教會的使者,說:「那右手拿著七星,在七個金燈臺中間行走的.」

這是主基督向以弗所教會顯明的身分.前一節聖經(啟1:20)已解釋:「七星就是七個教會的使者(或牧者),七燈臺就是七個教會」.「七」,一個象徵「圓滿」的數字.主基督的右手拿著(或者譯:堅守著)所有教會的牧者,他也在所有教會中間行走.教會在亂世,可能覺得「危危乎」,無所依靠.主基督說,我堅守著你.政治的壓力是大的,金錢的力量是大的,但都不是教會的主.唯獨基督,掌握著教會的命運.主在教會中行走,他看到教會裡發生的一切事情.有時,我們都會擔心教會的前景,但不用擔心,主基督持守著,看顧著.

啟2:2　我知道你的行為,勞碌,忍耐,也知道你不能容隱惡人.你也曾試驗那自稱為使徒卻不是使徒的,看出他們是假的來.

啟2:3　你也能忍耐,曾為我的名勞苦,並不乏倦.……

啟2:6　然而你還有一件可取的事,就是你恨惡尼哥拉一黨人的行為,這也是我所恨惡的.

主基督知道教會的所作所為,主基督知道一切.有時,單單知道,已很有安慰力量.太太打理了一天的家務,丈夫向她說:「我知道你的勞苦」,這已很有安慰力量.我們生命的一切,主基督知道.

主基督在多方面都讚賞以弗所教會.主基督讚他們有美好的行為表現,讚他們做善事時是不怕勞碌,讚他們能堅忍地面對困難,讚他們不容忍敵擋真理的惡人,讚他們能辨別出假的教會領袖來,讚他們實踐真理永不言倦.最後,也讚他們能斷然拒絕「尼哥拉」一黨人的行為,這也是主所恨惡的.我們並不能確定「尼哥拉」一黨人的行為,到底指什麼,只能推斷,他們是教會內的某種派系,他們表面有信仰,但內裡卻以世俗的想法和行為滲入信仰裡.他們作出很多妥協,在信仰裡容許很多「反信仰」的元素.以弗所教會堅決反對他們.

在以弗所這種大都會教會,經濟上充滿花花綠綠的銀紙的誘惑,政治上面對崇拜統治者的高壓,社會上彌漫著崇拜偶像的氛圍,處處充斥著利益交換和各種妥協,要做到以弗所教會所做到的,其實是極艱難的.

想想,我們能否堅忍地行善?

想想,我們是否有辨別能力,知道哪位領袖講的,能引導你進入永生,哪位領袖講的,會引導你走向滅亡?

想想,我們能否意識到,在某些地方,我們已向世俗的價值觀,作了過多的妥協?

以弗所教會做到這一切,其實已很好.若遇到一間像以弗所教會這樣的教會,我們應該很慶幸.

然而,然而,主基督卻責備他們.

啟2:4　然而有一件事我要責備你,就是你把起初的愛心離棄了.

「責備你」,原文直譯是「反對你」.主基督說:「我反對你,我站在你的對立面.」主基督對教會,對我們,又愛又恨.他用右手扶持著我們,卻又直接反對我們.他讚賞我們的美善,卻對我們的缺失,直斥其非.面對神,不能單單訴之溫情,不能單單說,我的軟弱,我的失腳,主必憐憫.主總有憐憫的一面,但針對我們的錯,主卻要明確地反對我們.不要沉醉在錯誤與罪惡必蒙憐憫的溫馨美夢裡,神要反對你,就反對你；要放棄你,就放棄你,你並不是宇宙的中心.

以弗所教會,錯在什麼地方呢?

他們「把起初的愛心離棄了」.中文和合本這樣的翻譯,將經文的重心放在信徒原初的愛心上.問題的重心,變成是個人德行的問題.有沒有愛心,變成問題所在.但觀乎整本《啟示錄》,個人的愛心德行,從來都不是《啟示錄》的信息重心.整本《啟示錄》,討論的,是神至高的王權,及信徒對這王權徹底的忠誠.如此,斷不會在《啟示錄》的開首位置,主基督特意責備教會信徒沒有個人的愛心.

「起初的愛心」,英文聖經一般譯作「First Love」.照字面來看,「First」就是「起初」；「Love」就是「愛心」.事實上,當我們讀到英文「First Love」時,想到的,不是「起初的愛心」,而是「初戀」.「初戀」者,是純潔的心全然傾注在愛的對象上.以弗所教會的問題,不是忘記了原先的愛心,而是遺忘了初戀,遺忘了他們選擇做基督徒的時候,他們是全心全意地將生命交給神.以弗所教會的問題,不在於重拾愛心.要有溫馨的愛,浪漫的愛,施予的愛,其實沒有什麼太大困難.以弗所教會的問題,不在這裡.他們的問題,是忘記了那改變他們生命的原初的愛的傾注.他們忘記了生命的焦點所在,忘記了賦予意義的委身,忘記了安身立命的根基.他們的問題,是遺忘了初心.主基督責備他們,是希望他們「毋忘初心」.「毋忘初心」,就是今日講道的題目.若由我去翻譯啟示錄2:4的經文,我會譯:「然而有一件事我要責備你,就是你把初心離棄了.」

原來,一間教會,縱使有種種佈道事工,種種栽培事工,種種團契生活,種種善行；縱使能教導真理,為真理辯護,斥退謊言；縱使能潔身自愛,有美好的品格；縱使有這一切,然而,若遺失初心,一切就沒有了.

一位丈夫,他在公司勤奮地工作,賺很多錢,買一間大屋來住,送很多禮物給太太,帶太太環遊世界,也沒有在外拈花惹草,縱使有這一切,但若他已遺忘了與太太之間原初的愛,若他已遺忘初心,這一切就都失去意義了.

或許,我們會說,這個男人已經很難得的了.或許,我們會說,就算只有個軀殼,都好過什麼也沒有.或許,我們會說,只有看開一點,這已算是幸福的了.在世人眼中,的確如是.但主基督反對這一切,他要求信徒,毋忘初心.這份「初心」,是意義的根基.他要求信徒,信仰的外表背後,有信仰的實質.大家在這裡崇拜上帝,平時的生活就要離棄各種偶像.只有上帝是你生命的主人,除此以外,任何東西都不能主宰你的生命.主基督要求的,就是這份初心.

啟2:5　所以,應當回想你是從哪裡墜落的,並要悔改,行起初所行的事.你若不悔改,我就臨到你那裡,把你的燈臺從原處挪去.

主基督的要求,是絕對的.若已遺忘了初心,遺忘了將生命全然交給神,就要悔改.若不悔改,永恆的生命便失去了.

啟2:7　聖靈向眾教會所說的話,凡有耳的,就應當聽!得勝的,我必將神樂園中生命樹的果子賜給他吃.

這是主基督向以弗所教會最後的叮嚀,主基督對這教會充滿期盼.主基督像是說,你行的,你會得勝的,你會明白生命的永恆意義的,你會展示生命的光明的,你會進入那永恆的神聖國度的.

在我們的社會,可以預見,政治的影響力會越來越大,政治背後有種種權力和金錢的計量,越來越多的力量影響法治,法治原先對權力的制約慢慢變成權力對法治的制約,言說真理的先知會受苦,假領袖會混入教會,信仰會在壓力下作出種種妥協,教堂的十字架會被拆掉,或許,有一天,這裡也不再有崇拜,就算如此,請各位弟兄姊妹相信,有一位神,在高天之上.就算如此,請各位弟兄姊妹相信,基督堅守著教會,他看見一切.就算如此,請各位弟兄姊妹,仍將生命傾注於上帝那裡,在那裡領受永恆的生命.就算如此,請各位弟兄姊妹,毋忘初心.

但願:榮耀歸於聖父,聖子,聖靈.阿們.



\newpage

\section{吃甚麼成了甚麼}
\label{sec:380}
\textbf{伍渭文牧師 - 吃甚麼成了甚麼}
\newline
\newline
\hyperref[sec:379]{< < < PREV SERMON < < <}
~
\hyperref[sec:toc]{[返目錄]}
~
\hyperref[sec:381]{> > > NEXT SERMON > > >}


哥林多前書 11:23-24
\newline
\begin{longtable}{cl}
\hline
\hline
章節 & 經文 (和合本修訂版)\\
\hline
11:23 & \begin{tabularx}{0.7\textwidth}{X} 我當日傳給你們的是從主所領受的。主耶穌被出賣的那一夜，拿起餅來， \end{tabularx} \\ \\ \relax
11:24 & \begin{tabularx}{0.7\textwidth}{X} 祝謝了，就擘開，說：「這是我的身體，為你們捨的；你們要如此行，為的是記念我。」 \end{tabularx} \\ \\
[1ex]
\hline
\hline
\end{longtable}
講題:吃甚麼成了甚麼　We Are What We Eat

經文:哥林多前書11:23-24

講員:伍渭文牧師

場合:香港中文大學崇基學院禮拜堂

日期:2017年9月3日

2010年十一月香港電視台藝人羅君左(又稱「斌仔」)因糖尿病併發症離世.他令我印象深刻,因為早年不改變飲食習慣已截腿坐輪椅,但乃不改變飲食習慣,最後拼發症離開,終年五十一歲. 香港營養師協會會長在其報章專欄說: 「瘋狂吃喝與慢性自殺無異.」

我們吃甚麼成了甚麼We are what we eat,屬靈的食糧也如是.前一些日子,看到一視頻:一位被人發現時像啞巴一樣不說話,滿身是病,被疹斷難以生存下來的孩童,母親是一個妓女,不知誰是爸爸.最初聽的時候,覺得這人真可憐,但後來發現這人就是他自己,這人竟然成為一位講道家.因為得到基甸送經會的一本聖經,命運坎坷的他研讀並相信而且順服聖經,變成另一個人.是的,我們吃甚麼成了甚麼.

其實吃是貫穿聖經的主題:聖經開始,始祖因為吃了分別善惡樹的禁果而墮落.分別善惡表示獨立自主,吃禁果是要離開上主,拒絕和背叛造物主,就如一棵樹告訴樹枝,你不能脫離樹幹,否則會枯乾死亡.聖經結束,新天新地出現,有羔羊的筵席.約翰福音七個神蹟中,有著名的五餅二魚吃飽五千人,耶穌為神蹟作出解釋:他就是生命的糧,天上賜下來的嗎哪.(約六35)

讀經:林前十一23, 24「我當日傳給你們的,原是從主領受的,就是主耶穌被賣的那一夜,拿起餅來,(自己)祝謝了,就擘開,說:(他者)這是我的身體,為你們捨的.你們應當如此行,為的是(上主)記念我.」「(教會)你們每逢吃這餅,喝這杯,是表明主的死,直等到他來.」(林前十一26)

這段經文說明:1 聖餐是教會的標記-身份証；2 聖餐是到父神的憑據-通行證,提起基督的工作,在聖靈裡,到父面前.記念的主體是父上主,父神因基督的工作,記念我們.3 因為基督的救贖,我們心存感戴,效法基督,對世人說:這是我的身體,為你們所捨的；世界說:這是你的身體,為我所享用.4 因為基督的救贖,我們心存感戴,過一個祝謝感戴的人生.

1 教會的身份証:

聖餐是基督信仰的核心.在正典確立之先[i],信經出現之先[ii],是信徒身証的認證.剛才讀出的哥林多前成書很早,約公元53年,但在此之前已有「從主領受的」傳統,傳到保羅那裡.「我們吃甚麼成了甚麼」也是聖餐的意義.教會要領洗後才能領聖餐是很早的傳統,洗禮是入門禮儀(rite of initiation),聖餐是持續禮儀(rite of continuation)；過紅海離開埃及為奴之地一次就夠了,進入迦南在曠野飄流需要每天從天而來的嗎哪.過紅海是一次過的聖洗,嗎哪是持續的聖餐.

基督三重身體.在神學上,基督的身體有三個不同的表述: (1) 道成肉身的身體(Incarnate body),一如信經所說:「童女懷孕,在本丟彼拉多手下遇難,被釘十字架…….」(2) 教會是基督的身體(Ecclesial body),亦稱奧體(Mystical body),世人不明白,以為教會不過是一般的社會組織,但聖經說,我們連於元首基督,互為肢體.「我們不拘是猶太人,是希利尼人,是為奴的,是自主的,都從一位聖靈受洗,成了一個身體,飲於一位聖靈.」(林前十二13) (3) 聖餐的餅酒是基督的聖體(Eucharistic body) 耶穌「拿起餅來,祝謝了,就擘開,說:這是我的身體,為你們捨的.」.基督三重身體:肉體,奧體,聖體.基督升天後,坐在父神的右邊,將來復臨,但不是缺席世界,缺席教會.對世界來說,教會是他的身體,世人若看見信徒彼此相愛,就認出基督在信徒中間；對教會來說,若有兩三個人奉基督的名聚會,基督就在他們中間.基督藉聖道和聖餐(聖體),與教會同在.基督藉教會(奧體),與世界同在.

破壞奧體的合一就混亂聖餐.我們知道哥林多前書十一章是對聖餐最詳盡的教導,其實十一章的聖餐,和十二章的恩賜教導,及十三章論愛是一氣呵成的.聖餐(十一章) 是例子,說明教會同是基督的身體,互為肢體,以恩賜彼此服侍(十二章),最大的恩賜是愛(十三章).保羅強調,我們共領同一個餅,分領同一的杯,就不能分門別類.領聖餐,得著永生的食物,生命連於基督,成為基督的身體,聖餐是信徒的身份證明.他們領聖餐表明生命連於基督,同屬一個身體,互為肢體,但在聖餐前的愛筵,表現卻不像同屬一個身體.

「你們聚會的時候彼此分門別類……算不得吃主的晚餐；因為吃的時候,各人先吃自己的飯,甚至這個饑餓,那個酒醉.你們要吃喝,難道沒有家嗎?還是藐視神的教會,叫那沒有的羞愧呢?」(林前十一 20-22)

「所以無論何人,不按理喫主的餅,喝主的杯,就是干犯主的身主的血了.人應當自己省察,然後喫這餅,喝這杯.因為人喫喝, 若不分辨是主的身體,就是喫喝自己的罪了.」(林前11:22-19)

他們甚麼時候吃喝不分辨是主的身體呢?是愛筵的時刻,不是聖禮的時刻.當時信徒肯定了解流行的神秘宗教(mystery religion),重視靈體臨在的時刻,他們的問題是在吃愛筵時出了問題,吃愛筵表示相愛合一,但他們分黨分派,「吃自己的飯」,不重視奧體,到聖餐認真吃聖體也無濟於事了；奧體,聖體同是基督的身體. 我們吃甚麼成了甚麼；或說我們如何吃就成了如何的人,We are what we eat; we are how we eat.聖餐表徵信徒合而為一,同屬基督的身體,整個聖餐禮儀,每個信徒都積極參與,不是旁觀者；他們落場打波,不是坐著看波.

每個信徒都參與呈獻.早期教會每次聚會,信徒都拿去餅酒作為呈獻[iii],然後襄禮的執事拿來(部份)給主禮的主教,主教在聯同眾長老按手或覆手在餅酒之上,然後拿起餅來,祝謝祈禱,然後擘開,遞給會眾領受.普世教會都按這四重禮儀行動:獻呈(酒餅)Offertory-祝謝Prayer-擘開Fraction-分領Communication,進行聖餐禮儀.東方士每拿主教與西方羅馬主教共祝聖餐,東西方地域,文化雖然不同,因禮儀相似.多年前我在香港一間教會主持聖餐,教會怕我程序跟慣常不同,做成混亂,清楚用文字記述如何一步一步進行聖餐.

信徒拿去餅酒呈獻,也代表把工作的成果(labour)獻呈,工作(work)和崇拜(worship)結合,餅酒是當時的食物,拿去的餅酒祇有少部份用作聖餐,其餘分派給缺乏的人,聖餐的餅酒是食物銀行.

這四重禮儀在中古時期斷裂了,信徒到教堂不需帶餅酒呈獻,也無須分領餅和杯,改教前信徒根本不重視領餐,祇看餅酒成聖體聖血神蹟.首(呈獻)尾(分領)沒有了,祇看一個認為是「神蹟」的誕生,聖餐集中在二,三的環節,聖職人員如何使餅酒祝聖成為基督的身體.

祝聖變成方程式,無論誰人,要是按一定的儀文,基督就臨在餅酒之中,我們稱這為因功生效(ex opere operato, out of the work it works).但甚麼是祝聖(consecration)?祝聖就是分別為聖,呈獻給上主使用.祝聖甚麼時候開始?不是聖禮人員拿起餅來祝謝那一刻,乃是信徒離開家庭,拿去餅酒到教堂那一刻；甚至是把種子撒在地裡澆水灌溉那一刻,因為期待有一天小麥長大,割下磨成麵粉,烘成餅,拿去教會呈獻成為基督的身體.

若把祝聖放在神職人員唸著禱文擘餅的一刻,祝聖禱文就像方程式,唸罷餅酒變作基督的身體,就像採用母乳配方粉(formula)哺嬰,但配方粉永遠不能替代母乳.雖然跨國奶粉商聲稱奶粉加入了神奇的添加劑PHD, phospholipid,形聖所有細胞膜的重要元素,包括腦細胞,父母就想當然吃了會令孩子聰明,因為PHD使人聯想到博士.母乳形成有一過程:母親的飲食,運動,作息,心情等等；哺乳時又有身體親密的互動,還可以把抗體傳給嬰兒,這是配方粉不能給予的.

祝聖貫徹四重禮儀行動-呈獻-禱告-擘開-分領,不是禱告祝聖餅酒那一刻.所以今天聖餐我恢復餅酒呈獻一禮儀,聖餐進行之始,由襄禮端餅杯進場,這是信徒的呈獻.這行動幫助我們唸祝聖禱文時更真實:「在基督祭司的職分裡,我們與他聯合在一起,向你獻上這祭禮.」聖餐不單是教會身體的認證,信徒的身份證,也是到父神面前的通行證.我們看看「你們如此行,為的是記念我」是甚麼意思.

2 為的是記念我:

記念anamnesis,來自舊約Le-zikkaron, 是聖經禮儀常用語.記起(memory)是悟性追憶,想起一樁不在場的事情或人物.但記念可翻作memorial是聖經禮儀用詞,跟向上主禱告獻祭有關.使10:4 彼得跟哥尼流說:「你的禱告和你的賙濟已達到上主面前.」哥尼流也印證:「前四天這個時候,我在家中守著申初的禱告,忽然有一個人,穿著光明的衣裳,站在我面前,說:哥尼流,你的禱告,已蒙垂聽,你的賙濟,達到神面前已蒙紀念了.」使10: 30,31.誰是記念的主體呢?是上主.

再看先知以賽亞的話:賽43:25「惟有我為自己的緣故塗抹你的過犯,我也不記念你的過犯,我也不記念你的罪惡.」以西結書18:22「他所犯的一切罪過都不被記念,因所行的義,他必存活.」因為上主的拯救,成了上主的子民,不被記念罪過.尼希米就握住這子民的身份求上主記念.尼5: 19「我的神阿,求你記念我為這百姓所行的一切,施恩與我.」.

人不是記念的主體,上主才是記念的主體；是父上主因耶穌為我們所作的一切記念我們.耶穌叫我們記念他在父面前為我們所作的一切,因為惟有藉基督,在聖靈裡,才能到父面前.我們到父面前,求他記念耶穌為我們所做的一切.

記念anamnesin:也許這字前置片語幫助我更明白多一點記念anamnesin的意思.Ana是提起,耶路撒泠在傳說是耶穌被埋的墳墓其上所蓋的教堂,西方教會(羅馬天主教)稱主墓堂,而東方教會(正教)稱復活堂(anastasis),復活一字在希臘文站起來.Anamnesin記念重點不是回到過去的各各他,乃把各各他提到現在,求上主記念基督為我們所做的一切,直等到他來不單是副詞,不斷記念,直到基督再來；更多是記念的內容.

直等到他來:保羅說:「你們每逢吃這餅,喝這杯,是表明主的死,直等到他來.」(林前十一26) 表明,是信仰的宣告,其內容包括他的生,死,埋葬,復活,升天,坐在父的右邊,再來.我們提起整個基督救贖事件,求父神記念.我們也與基督聯合,以餅酒為祭,把自已獻給上主,這是我們的感恩祭(eucharist).

記念重點不在我們主觀的記念,不在我們如何敬虔,如何反思；我們對上主的愛是非常有限的,容易枯竭,也容易放棄.        記念的主體是父,因為基督的原故,上主接納我們.基督的愛像湧流不住的泉源,用之不竭.藉基督,在聖靈裡,到父面前.當作為配偶的,做了不該做的事,尋求饒恕時,提起結婚戒指,不是回到過去結婚帶上婚戒的時刻,乃把過去的事件,提起,放在當下:那天我們彼此承諾過:「無論疾病健康,富厚貧窮,我必守助你,我必愛護你,直至終身.」然而,人間的承諾也會因人性的脆弱,無法履行；但基督的救恩,大有能力.我們記念基督的工作,求父神記念,基督是我們到上主面前的通行証.所以聖餐重點不在追思基督的死,認罪自省,乃感恩祝謝,是一個歡樂的筵席.

這記念是歡愉的.聖餐(Holy Communion)又稱主餐(The Lord’s Supper)和感恩祭(Eucharist).聖餐一詞強調契通,跟上主和信徒彼此的契通；主餐強調基督是我們的屬靈糧食,生命的糧；感恩祭強調呈獻,就如我們剛才提及,在合一聖餐禮文所說:「在基督的職份裡我們聯合在一起,向你獻上這祭禮.」感恩的概念來自祝謝,拿起餅來,祝謝了就擘開.這祝謝也是信徒對基督的救贖工作的回應.其實崇拜基本的觀念是回應,感謝的回應.相比之下,世界的記念碑記念館是悲傷的,也隱含是控訴.

紀念碑是追思也是控訴.不少國家民族都立有紀念碑和紀念館,追思過去難忘的事件:日本原炸紀念館,南京大屠殺紀念館,紐約九一一紀念館,猶太人在各地都有大屠殺(holocaust)紀念館.這些記念不單對受害者表達哀思,也隱含對加害者的控訴.但基督在士架上,把所有的苦難承受過來,他吸納了仇恨,仇恨就止住了.

記念,anamnesin提起,在夜深人靜時,我們的良心也提起一些記憶,一些講過的說話,一些做過的事；我們感到不安,懊悔,愧疚.撒旦最拿手就是控訴,他被稱為控訴者(accuser),常常以負面的事控訴我們?(啟12: 10)但發生了的事已經發生了,我們不能逆轉時間重新來過,像電影中的超人,把地球逆轉,用盡自己的能力,愈來愈快反方向轉動,回到過去意外尚未發生的時間,把死去的愛人拯救過來.我們不能回到過去,抹去後悔做過的事情.但可以來到到聖餐桌前,提到基督為我們所做的一切.拿起餅來,主說:「這是我的身體,為你而捨的」,拿起杯來,主說:「這是赦罪的血,為你流出」.我們用聖餐的恩典提起,來抗衡撒旦控訴的記提起.提起(記念)基督的工作,求父記念；提起(記念)基督的工作,放下各樣的重擔.「讓我們放下各樣的重擔,脫去容易纏累我們的罪,存心忍耐,奔那擺在我們前頭的路程；徙仰望為我們信心創始成終的耶穌.」來十二1, 2.

這是我的身體:

基督如何在聖餐中,臨在親近我們呢?我們在這裡沒有時間討論各種臨在的教義,但臨在是真實的.「因為人吃喝,若不分辨是主的身體,就是吃喝自己的罪了.因此,在你們中間有好些刊-軟弱的,與患病的,死的也不少.」(林前十一30) 但從路加福音以馬忤斯路上的故事我們得到耶穌臨在的啟發.

以馬忤斯路上的臨在.馬丁路德及其跟隨者提到聖餐基督的臨在用了三個很簡單的前置詞,不多解釋:in, under, with在餅酒之中,餅酒之下,伴隨餅酒.我想這三個前置詞在路二十四以馬忤斯路上的故事,也可適用解釋基督的臨在:在神話語中(in),隱藏在同路人中(under),及在伴隨有需要的人(with).兩位門徒何時發現耶穌?在接待陌生人一起用餐時.耶穌豈不是說:「因為我餓了,你們給我吃,渴了,你們給我喝.我作客旅,你們留我住.」(太二十五35)那裡有苦難,那裡有受苦的基督；作在小子的身上,就是作在基督的身上.他們也記起耶穌解聖經時,心裡火熱.基督沒有離開我們,以聖道陪伴我們.聖體包含兩個元素:道和物質,聖洗物質是水,聖餐物質是餅酒.聖道是聽得見,聖禮是看得見,但都有上主的話語.還有那位從耶路撒泠下來的香客,一路陪伴聆聽,祇是他們眼睛迷糊看不出來,是上主使他們眼睛迷糊,看不清楚；上主隱藏在同路人中伴隨我們.

他們聽聖經時,祇是心裡火熱,直到以善行接行陌生人,使饑餓的得飽足,無家者得寄宿,看到「時候晚了,日頭已經平西了」,就發出邀請,「請你同我們住下吧.」就在這時候,基督顯現.聖餐的精神:「這是我的身體,為你們捨的.」(This is my body given for you) 當下的消費主義,剛剛相反,「這是你的身體,為我所享用.」(This is your body taken by me)因為記念基督的救恩,因為祇能憑著信靠基督,在聖靈裡,才能到父面前.到神面前,別沒他法,惟獨信心.但心存感戴,到人面前,就要顯出愛心,作為對上主恩典的回應.如何使我們的信心活潑常存生生不息呢?就是操練愛心,faith active in love.愛心使工作變得輕省愉快,以愛為動力,在愛中沒有懼怕,工作是感戴上主救贖的回應,也是聖餐祝謝的意思.

祝謝了

最後,我以一個小故事說明祝謝,感恩祝謝.1982年我到芝加哥參加一個教牧研討會,講者是一著名神學院的宗教育系主任.他說這講座令他承受很大的壓力,因為聽眾都是很有經驗的牧師.起初他想盡辦法要預備一篇精彩的講演,不想失禮,也不想聽眾失望,其實不想失禮多一點.但壓力隨著講演的日子愈近我大.他惟有祈禱舒解壓力,在禱告中,他覺得應轉移焦點,不看自己的表現(performance),看為服侍的機會.幫助牧師,更好幫助信徒,頓時壓力輕了很多.祝謝是感恩,感謝上主給機會服侍,祝謝是以他者為目標,感戴上主在基督裡的恩典,服侍他者.

「我當日傳給你們的,原是從主領受的,就是主耶穌被賣的那一夜,拿起餅來,(自己)祝謝了,就擘開,說:(他者)這是我的身體,為你們捨的.你們應當如此行,為的是(上主)記念我.」

[i]正典(量度的尺)確立見亞他拿修(Athanasius) c.367 Festal Epistle新約27書卷.

[ii]信經形成前: 公元390安波羅修(Ambrose) Epistle 42, 5提到信經.

[iii] 希坡律陀 (d. 235)使徒傳統 Hippolytus Apostolic Tradition (三世紀) (\url{http://www.orthodox.cn/patristics/hippolytusapostrad_tw.htm2007}年2月22日初期中譯來自香港及東南亞都主教教區的台北聖三堂)

(按立主教後的聖餐)「然後執事必須獻上祭品(呈獻Offertory),而他必須把手放在祭品上,祝謝了,由議會的全體長老陪同,(序經對答Preface, dialogue)說:3願主與你們同在.而所有人要回答:也與你的心靈同在.主教說:讓我們舉心向上.眾人回答:我們已舉心向主.主教說:讓我們為主致上感謝.眾人回答:這是適當而且正確的.然後主教繼續說(祝謝Eucharist Prayer):4我們在此向你獻上感恩,藉由你所喜愛的兒子,你將祂先前的時光a賜給我們,作為救主,並作為傳遞你意旨的使者.5祂與你的聖言密不可分,藉由祂你成就一切,在祂裡面你大大的歡喜,6你從一位處女的子宮裡把祂從天上派遣下來給我們,在她裡面懷胎時使有了血肉之軀,成為你的兒子,從聖靈與處女所生.7 也是祂成就了你的意旨,並為你得著一群神聖的人們,祂伸出手來接受苦難,為了就是要解放那些因為信你而受苦的人.8當祂被遣b往自願的受難時,就是為了要 消滅死亡,並毀壞魔鬼的鎖鏈,踐踏地獄,把義人帶到光裡面,釋放受束縛的人,並展現祂的復活.9取來餅,並向你致上感謝,說:「拿去吃吧,因為這是我的身 體,為你們擘開.」同樣地拿起杯,說(立餐話語words of institution):「這是我為你們流的血.10你們當這樣行,為的是紀念我.」11因此,記得祂的死亡與復活時,我們(信徒拿來,執事拿去,[長老聯同]主教拿起[餅酒])向你獻上餅與杯, 向你致上感謝,是你讓我們配得站在你面前,並以你的司祭身份服事你.12(聖靈臨在餅酒和教會,使聖體,奧體合一)我們祈求你派遣聖靈,來到你神聖教會的祭品當中.在大家的聚會裡,將聖靈的充滿給 予所有參與你奧秘的人,堅定他們在真理中的信心,13那麼我們便讚美你,榮耀你,藉由你的兒子耶穌基督,你得以彰顯榮耀,聖父與聖子,連同聖靈,在你神聖 的教會裡,從現在直到永遠,世世無盡,阿民.」

(復活節領洗後的聖餐)20. 10已領過浸的人在此不需攜帶聖餐禮的聖器前來,因為那是聖餐禮才使用的.在同一時刻,讓每個人攜帶祭品前來是很適當的.為常常會有人帶祭品來,而這是不能拒絕的.



\newpage

\section{淚眼泯恩仇}
\label{sec:381}
\textbf{林豪恩先生 - 淚眼泯恩仇}
\newline
\newline
\hyperref[sec:380]{< < < PREV SERMON < < <}
~
\hyperref[sec:toc]{[返目錄]}
~
\hyperref[sec:382]{> > > NEXT SERMON > > >}


創世記 50:15-21
\newline
\begin{longtable}{cl}
\hline
\hline
章節 & 經文 (和合本修訂版)\\
\hline
50:15 & \begin{tabularx}{0.7\textwidth}{X} 約瑟的哥哥們見父親死了，就說：「也許約瑟仍然懷恨我們，會照我們從前待他一切的惡，重重報復我們。」 \end{tabularx} \\ \\ \relax
50:16 & \begin{tabularx}{0.7\textwidth}{X} 他們就傳口信給約瑟說：「你父親未死之前曾吩咐說： \end{tabularx} \\ \\ \relax
50:17 & \begin{tabularx}{0.7\textwidth}{X} 『你們要對約瑟這樣說：從前你哥哥們惡待你，你要饒恕他們的過犯和罪惡。』現在求你饒恕你父親的神之僕人們的過犯。」他們對約瑟說了這話，約瑟就哭了。 \end{tabularx} \\ \\ \relax
50:18 & \begin{tabularx}{0.7\textwidth}{X} 他的哥哥們又來俯伏在他面前，說：「看哪，我們是你的奴隸。」 \end{tabularx} \\ \\ \relax
50:19 & \begin{tabularx}{0.7\textwidth}{X} 約瑟對他們說：「不要怕，我豈能代替神呢？ \end{tabularx} \\ \\ \relax
50:20 & \begin{tabularx}{0.7\textwidth}{X} 從前你們的意思是要害我，但神的意思原是好的，要使許多百姓得以存活，成就今日的光景。 \end{tabularx} \\ \\ \relax
50:21 & \begin{tabularx}{0.7\textwidth}{X} 現在你們不要害怕，我必養活你們和你們的孩子。」於是約瑟安慰他們，講了使他們安心的話。 \end{tabularx} \\ \\
[1ex]
\hline
\hline
\end{longtable}
羅馬書 14:1-12
\newline
\begin{longtable}{cl}
\hline
\hline
章節 & 經文 (和合本修訂版)\\
\hline
14:1 & \begin{tabularx}{0.7\textwidth}{X} 信心軟弱的，你們要接納，不同的意見，不要爭論。 \end{tabularx} \\ \\ \relax
14:2 & \begin{tabularx}{0.7\textwidth}{X} 有人信甚麼都可吃；但那軟弱的，只吃蔬菜。 \end{tabularx} \\ \\ \relax
14:3 & \begin{tabularx}{0.7\textwidth}{X} 吃的人不可輕看不吃的人；不吃的人也不可評斷吃的人，因為神已經接納他了。 \end{tabularx} \\ \\ \relax
14:4 & \begin{tabularx}{0.7\textwidth}{X} 你是誰，竟評斷別人的僕人呢？他或站立或跌倒，自有他的主人在，而且他也必會站立，因為主能使他站穩。 \end{tabularx} \\ \\ \relax
14:5 & \begin{tabularx}{0.7\textwidth}{X} 有人看這日比那日強；有人看日日都是一樣。只是各人要在自己的心意上堅定。 \end{tabularx} \\ \\ \relax
14:6 & \begin{tabularx}{0.7\textwidth}{X} 守日子的人是為主守的。吃的人是為主吃的，因他感謝神；不吃的人是為主不吃的，他也感謝神。 \end{tabularx} \\ \\ \relax
14:7 & \begin{tabularx}{0.7\textwidth}{X} 我們沒有一個人為自己而活，也沒有一個人為自己而死。 \end{tabularx} \\ \\ \relax
14:8 & \begin{tabularx}{0.7\textwidth}{X} 我們若活，是為主而活；我們若死，是為主而死。所以，我們或死或活總是主的人。 \end{tabularx} \\ \\ \relax
14:9 & \begin{tabularx}{0.7\textwidth}{X} 為此，基督死了，又活了，為要作死人和活人的主。 \end{tabularx} \\ \\ \relax
14:10 & \begin{tabularx}{0.7\textwidth}{X} 可是你，你為甚麼評斷弟兄呢？你又為甚麼輕看弟兄呢？因我們都要站在神的審判臺前。 \end{tabularx} \\ \\ \relax
14:11 & \begin{tabularx}{0.7\textwidth}{X} 經上寫著：「主說，我指著我的永生起誓：萬膝必向我跪拜；萬口必稱頌神。」 \end{tabularx} \\ \\ \relax
14:12 & \begin{tabularx}{0.7\textwidth}{X} 這樣看來，我們各人一定要把自己的事在神面前交代。 \end{tabularx} \\ \\
[1ex]
\hline
\hline
\end{longtable}
馬太福音 18:21-35
\newline
\begin{longtable}{cl}
\hline
\hline
章節 & 經文 (和合本修訂版)\\
\hline
18:21 & \begin{tabularx}{0.7\textwidth}{X} 那時，彼得進前來，對耶穌說：「主啊，我弟兄得罪我，我當饒恕他幾次呢？到七次夠嗎？」 \end{tabularx} \\ \\ \relax
18:22 & \begin{tabularx}{0.7\textwidth}{X} 耶穌說：「我告訴你，不是到七次，而是到七十個七次。 \end{tabularx} \\ \\ \relax
18:23 & \begin{tabularx}{0.7\textwidth}{X} 因為天國好像一個王要和他僕人算賬。 \end{tabularx} \\ \\ \relax
18:24 & \begin{tabularx}{0.7\textwidth}{X} 他開始算的時候，有人帶了一個欠一萬他連得的僕人來。 \end{tabularx} \\ \\ \relax
18:25 & \begin{tabularx}{0.7\textwidth}{X} 因為他沒有甚麼償還之物，主人下令把他和他妻子兒女，以及一切所有的都賣了來償還。 \end{tabularx} \\ \\ \relax
18:26 & \begin{tabularx}{0.7\textwidth}{X} 那僕人就俯伏向他叩頭，說：『寬容我吧，我都會還你的。』 \end{tabularx} \\ \\ \relax
18:27 & \begin{tabularx}{0.7\textwidth}{X} 那僕人的主人就動了慈心，把他釋放了，並且免了他的債。 \end{tabularx} \\ \\ \relax
18:28 & \begin{tabularx}{0.7\textwidth}{X} 那僕人出來，遇見一個欠他一百個銀幣的同伴，就揪著他，扼住他的喉嚨，說：『把你所欠的還我！』 \end{tabularx} \\ \\ \relax
18:29 & \begin{tabularx}{0.7\textwidth}{X} 他的同伴就俯伏央求他，說：『寬容我吧，我會還你的。』 \end{tabularx} \\ \\ \relax
18:30 & \begin{tabularx}{0.7\textwidth}{X} 他不肯，卻把他下在監裡，直到他還了所欠的債。 \end{tabularx} \\ \\ \relax
18:31 & \begin{tabularx}{0.7\textwidth}{X} 同伴們看見他所做的事就很悲憤，把這一切的事都告訴了主人。 \end{tabularx} \\ \\ \relax
18:32 & \begin{tabularx}{0.7\textwidth}{X} 於是主人叫了他來，對他說：『你這惡奴才！你央求我，我就把你所欠的都免了； \end{tabularx} \\ \\ \relax
18:33 & \begin{tabularx}{0.7\textwidth}{X} 你不應該憐憫你的同伴，像我憐憫你嗎？』 \end{tabularx} \\ \\ \relax
18:34 & \begin{tabularx}{0.7\textwidth}{X} 主人就大怒，把他交給司刑的，直到他還清了所欠的債。 \end{tabularx} \\ \\ \relax
18:35 & \begin{tabularx}{0.7\textwidth}{X} 你們各人若不從心裡饒恕你的弟兄，我天父也要這樣待你們。」 \end{tabularx} \\ \\
[1ex]
\hline
\hline
\end{longtable}
講題:淚眼泯恩仇　Grudges washed away in tear

經文:創世記50:15-21;詩篇103;羅馬書14:1-12;馬太福音18:21-35

講員:林豪恩先生

場合:香港中文大學崇基學院禮拜堂

日期:2017年9月17日

一．引言

金庸先生的武俠小說膾炙人口,有些拍成了電視和電影.不知道你對金庸武俠小說認識多少呢?

一位中年人問年輕人:「你看過金庸的小說嗎?」

年輕人說:「沒有,只有看過電視劇.」

中年人說:「那你知道金庸寫的14部小說的書名的第一個字,串起來會成為一首詩嗎?『飛雪連天射白鹿,笑書神俠倚碧鴛』」

年輕人說:「不知道,但是我有看羅琳(哈利波特作者)的小說,你知道這七本小說書名的第一個字串起來是什麼嗎?」

中年人:不知道.

年輕人:她寫的七本小說書名的第一個字串起來是:『哈哈哈哈哈哈哈』!

有人喜歡金庸的武俠小說,有人喜歡羅琳的哈利波特.不知道你有沒有發現這些小說在劇情上有甚麼共同之處呢?沒錯,就是「報仇」.原來,在小說世界裡,總少不了恩怨情仇,例如:

《射鵰英雄傳》中,郭靖前半段一心報父仇,誓要殺掉迫害自己父親的幕後主使完顏洪烈.而在後半段,由於師傅江南五怪死於桃花島,他要報的由父仇變為師仇,報仇這一主線也是連接整部小說的線索.《神鵰俠侶》中,楊過為父仇所困,要尋找殺父仇人.裘千尺居深潭多年也是為報仇雪恨；還有李莫愁雖名為莫愁,但卻愁怨極深,因為內心飽受情傷而充滿著情仇.

恩怨情仇出現於武俠小說,也在現實世界上演.基督教裡面也有恩怨情仇嗎?當然有,所以中國神學研究院院長李思敬博士那本舊約聖經導讀,就名叫「恩怨情仇論舊約」.

今天我們誦讀的舊約經課,可算是恩怨情仇的其中一幕.

二．經文

50:15約瑟的哥哥們見父親死了,就說:「或者約瑟懷恨我們,照著我們從前待他一切的惡足足地報復我們.」

50:16他們就打發人去見約瑟,說:「你父親未死以先吩咐說:

50:17『你們要對約瑟這樣說:從前你哥哥們惡待你,求你饒恕他們的過犯和罪惡.』如今求你饒恕你父親上帝之僕人的過犯.」他們對約瑟說這話,約瑟就哭了.

50:18他的哥哥們又來俯伏在他面前,說:「我們是你的僕人.」

50:19約瑟對他們說:「不要害怕,我豈能代替上帝呢?

50:20從前你們的意思是要害我,但上帝的意思原是好的,要保全許多人的性命,成就今日的光景.

50:21現在你們不要害怕,我必養活你們和你們的婦人孩子.」於是約瑟用親愛的話安慰他們.

三．釋經與應用

古人說:冰封三尺,非一日之寒.此時此地的人說:深層次矛盾.如何處理恩怨情仇,如何能夠和解?的確是非常複雜的事情.畢竟,沒有人能夠逼令加害者誠心尋求寬恕,也沒有人能夠勉強受害者真心給予寬恕.約瑟與兄弟們的故事,只是世上無數恩怨情仇的其中一個例子,而且是家族裡的事情,未必能夠普及至其他層面,因此,我們應用時必需謹慎.

使徒保羅曾經教導信徒「不可含怒到日落」,然而,解決恩怨情仇,就正如保羅另一句話:「立志行善由得我,只是行出來由不得我」；除了意願之外,時機也是非常重要的因素.再讀約瑟與兄弟的故事時,希望我們能從中尋找一點啟迪.

3．1．約瑟與兄弟們淚眼泯恩仇的過程

第一:舊秩序失效,新秩序待建的轉捩點

50:15約瑟的哥哥們見父親死了,就說:「或者約瑟懷恨我們,照著我們從前待他一切的惡足足地報復我們.」

仇怨不一定會顯露,仇恨也不一定能夠報復,當有大家都尊重的父老在維持秩序的狀況下,大家都「俾面」父老,雙方就會在這種秩序之下和平共處.不過,這種表面的和諧未必是真正的和平,可能雙方之間的恩怨只是被壓了下去,沒有爆發出來而已.因此,當那位大家都「俾面」的父老消失,維繫秩序的力量也消失,埋藏了的恩怨情仇就有可能以各種形式浮現出來.有時,傳媒報導某些富豪家庭在父親過身後爆發家人之間的衝逡,或者也是這種類似劇情.

「約瑟的哥哥們見父親死了」,即是說原來那種穩定的兄弟,父親,約瑟這三角關係瓦解了,變成兄弟與約瑟之間的二者關係,即是說,兄弟與約瑟將要直接面對他們之間埋藏了的恩怨情仇.家庭研究學者發現,當兩個人的關係充滿張力時,雙方都重視的第三者的出現有助穩定關係.例如夫妻不和,但雙方在「為了子女有一個家」的原因之下,繼續維持關係.但有一天,當子女離家,父親,母親,子女這相對穩定的三角關係變成了丈夫和妻子直接面對的兩人關係,如果夫妻兩人本來的矛盾沒有處理好,又不能在新狀況下重新調整,危機就會出現.這就是為甚麼婚姻危機的另一個高峰期竟然是在子女獨立之時.

約瑟和兄弟們,在沒有父親這第三者維持秩序的狀況下,張力立即湧現.按經文的敘述,卻不是雙方都感覺到這種張力,有反應的,是兄弟們,而不是約瑟.對著新的狀況,兄弟們似乎感到驚慌.他們感到驚慌是有背景的,其一是他們曾經傷害約瑟,其二是現在約瑟有能力向他們報仇.於是,他們擔心:「或者約瑟懷恨我們,照著我們從前待他一切的惡足足地報復我們.」

這種舊秩序失效,新秩序待建的轉捩點非常重要.這個非常時期提供了非常的機會給雙方重新檢視彼此的關係,也讓雙方能夠以全新的態度調整當下的關係.舊秩序失效,新秩序待建,是機,也是危.危者,可以陷入叢林法則,或者如武俠小說般怨怨相報的情節:從前兄弟加害約瑟,現在約瑟向兄弟報仇；若干年後,兄弟的後人強盛起來,去找約瑟報仇；又若干年後,約瑟的後人奮發圖強,去找兄弟的後人報仇.惡性循環,沒完沒了.

約瑟與兄弟的故事會陷入冤冤相報的惡性循環嗎?他們如何回應積壓多年的恩怨情仇呢?

第二:犯錯者認錯及請求饒恕

50:16他們就打發人去見約瑟,說:「你父親未死以先吩咐說:

50:17『你們要對約瑟這樣說:從前你哥哥們惡待你,求你饒恕他們的過犯和罪惡.』如今求你饒恕你父親上帝之僕人的過犯.」

約瑟和兄弟們的怨仇能夠走出冤冤相報的惡情循環的第二步,是兄弟們向約瑟承認「從前你哥哥們惡待你」,承認所做的是「過犯和罪惡」.向受害者認錯,其實是不容易的,有些犯錯者會倒果為因說:「其實我們當年所做的,都是時勢所需.如果當年不是這樣做,就不會有今天的成就!」還有很多不肯承認錯誤的語言藝術:明明是軍隊入侵鄰國,卻說成軍隊「進出」鄰國.所以,德國在戰後真誠地向歐洲的受害國度歉是難得的表現.1970年12月7日,時任西德總理來到華沙猶太人死難者紀念碑下.他向紀念碑獻上花圈後,肅穆垂首,突然雙腿下跪,並祈禱說:「上帝饒恕我們吧,願苦難的靈魂得到安息.」總理以此舉向二戰中無辜被納粹黨殺害的人表示沉痛哀悼,並虔誠地為納粹時代的德國認罪.

第三:受害者感到安慰

50:17他們對約瑟說這話,約瑟就哭了.

勃蘭特在波蘭猶太人紀念碑前下跪謝罪,淡化了飽受納粹蹂躪的波蘭人民沉積在心底里的憤怒,他們為勃蘭特的舉動感動得熱淚盈眶,為真正的和解產生了極為深遠的影響.

兄弟們來承認對自己的惡行,請求饒恕,約瑟必然百感交雜.可能想到自己多年來所受的辛酸；可能傷感兄弟們一直生活在惶恐之中,被遭報復的擔心所驚嚇；可能無奈兄弟們原來一直視自己對他們的寬愛厚待為虛偽……萬千感受湧上心頭,此時此刻,不能言語,都變為哭聲；無言以對,都溶於淚水之中.

第四:犯錯者願意承擔後果

50:18他的哥哥們又來俯伏在他面前,說:「我們是你的僕人.」

對犯錯者來說,口頭上的認錯只是和解的第一步,以行動來承擔後果是口頭道歉的兌現.戰後的德國以國家的名義賠償給受害者,並且直至現在為止,仍然禁止任何納粹德國的意識及行為.今年八月,兩名中國遊客在德國旅遊時模仿納粹德軍的敬禮,結果被拘捕,由此可見德國與納粹時期所犯的錯誤一刀兩斷的決心.相比之下,日本處理二次大戰的罪行的態度,卻令被傷害的鄰國一直放不下這心結.

約瑟的兄弟們不單單在口頭上認錯,而且願意按他們那個年代的方式承擔可能的後果.約瑟如何回應呢?

第五:受害者給予寬恕

50:19約瑟對他們說:「不要害怕,我豈能代替上帝呢?

50:20從前你們的意思是要害我,但上帝的意思原是好的,要保全許多人的性命,成就今日的光景.

50:21現在你們不要害怕,我必養活你們和你們的婦人孩子.」於是約瑟用親愛的話安慰他們.

約瑟如此超然的回應,並非因為他瀟灑.其實,約瑟也是凡人,也有作為人的七情六慾.只是,約瑟早已處理了他與兄弟間的恩怨情仇.你還記得約瑟曾經在甚麼時候掙扎著如何面對兄弟們嗎?你還記得約瑟曾經在甚麼時候為此而痛哭流淚嗎?沒錯,就是當他的父家遇到餓荒,兄弟們來埃及求助的時候,只是兄弟們未必能體會約瑟的心情,也未必能信任約瑟當時對他們的寬恕而已.此時此刻,約瑟沒有因兄弟們的不信任再次請求寬恕而憤怒,反而再次向兄弟們發出寬恕及和好的保證,讓犯錯求恕者安心.

3．2．和解的三角關係

約瑟在這裡提出的理解,其實之前已經對兄弟們說過一次,這理解對於處理恩怨情仇是非常重要的,讓我們再以剛才提及的三角關係來思考約瑟提出的理解.

我們發現,兄弟們視父親為他們與約瑟之間的第三者,是高於兄弟們與瑟兩者的權威,他們依靠這個權威來維持秩序,並且依靠這個權威來保證瑟不會向他們報復.正因如此,這個兄弟們擔心而向約瑟請求饒恕的片段,是在父親去世後才出現的.

且看兄弟們憑甚麼向約瑟請求饒恕?

50:16他們就打發人去見約瑟,說:「你父親未死以先吩咐說:50:17『你們要對約瑟這樣說:從前你哥哥們惡待你,求你饒恕他們的過犯和罪惡.』如今求你饒恕你父親上帝之僕人的過犯.」

從這段說話可見,兄弟們沒有轉換原來那種三角關係的框架,仍然以父親的權威作為力量去處理他們與約瑟的恩怨情仇,以父親的影響力去請求約瑟寬待.人間的權威在維持秩序上的確有作用,就如被有效及正確執行的法律能夠阻止仇殺和私刑.不過,法律不一定能夠清洗仇怨,保證復和.君不見有些爭產案在法庭上判決了,但家人之間的恩怨並沒有隨著判決而化解,反而從此成為陌路人.約瑟的兄弟們離不開這種依靠人間的第三權威來處理衝突的思維,因此,他們所期望的最好結果,仍然是充滿權力差異的主僕關係,而不是平等的兄弟關係.且再聽他們如何說:

50:18他的哥哥們又來俯伏在他面前,說:「我們是你的僕人.」

約瑟有甚麼不同框架去理解他與兄弟們的恩怨情仇呢?原來也離不開三角關係.不過,約瑟所依據的,不是父親這個人間的第三力量.那麼,他所依據的第三者是誰呢?所產生的結果又有甚麼不同呢?

50:19約瑟對他們說:「不要害怕,我豈能代替上帝呢?

50:20從前你們的意思是要害我,但上帝的意思原是好的,要保全許多人的性命,成就今日的光景.

對約瑟來說,處理他與兄弟間的恩怨情仇的第三者,不是父親,也不是法老,而是上帝.約瑟邀請上帝介入他與兄弟的恩怨情仇之中,因此,除了自己的感受和利益,除了人間當時的行為標準,除了他的個人可運用的權力,他多了一個角度去回顧事情,多了一種理解去疏理經歷,那就是上帝的意思.兄弟的三角關係是自己-父親-約瑟,因此,父親去世,這三角關係就動搖；約瑟的三角關係是兄弟-上帝-自己,縱使父親去世,也沒有動搖.約瑟接受上帝成為第三者,因此和解所產生出來的不是權力差異的主僕關係,而是兄弟間的互助關係；不是有你死我活的鬥爭社會,而是共同存活的生存狀態:

50:21現在你們不要害怕,我必養活你們和你們的婦人孩子.」於是約瑟用親愛的話安慰他們.

四．總結

可能有人會質疑:現實怎會如此完美?沒錯,現實是不完美的,因此,這世界不可能沒有恩怨情仇,正如人們常說「沒有無緣無故的恨」.正因如此,這個世界需要尋找回應恩怨情仇的方法.雙方都認定自己有冤要報,以牙還牙,以眼還眼,是其中一種回應,也是常見的回應.不過,這種回應的結果可能是大家都沒有牙齒,也變成失明.如何可能超越這種回應方式呢?如果沒有一種超越雙方存亡的存在,如果沒有一種高於雙方利益的理想,如果沒有一種維護雙方尊嚴的保障,如果沒有一種治療雙方傷痛的力量,誰能夠冒如此大的風險去化解怨仇,創造復和呢?

對給予寬恕的約瑟來說,這位超越的第三者是耶和華上帝；對那位尋求寬恕的德國總理來說,那位超越的第三者也很可能是上帝,他跪地道歉時如此說:「上帝,饒恕我們吧!」上帝介入,不是黑白不分地每人打五十大板,也不是以威逼利誘製造虛假的和諧.他邀請我們在他面前謙卑自省,回想他如何對待我們,以他的角度重新檢視纏繞我們的恩怨情仇.

今天我們讀的書信經課中保羅如此回應說(羅馬書):

14:10你這個人,為甚麼論斷弟兄呢?又為甚麼輕看弟兄呢?因我們都要站在上帝的臺前.

如果用約瑟的說話來總結保羅對信徒之間的爭執的指導,就是:「我豈能代替上帝呢?」在你我之上,還有上帝呢!詩篇經課建議我們撫心自問,我們自己是怎樣的人?上帝又是如何對待我們呢?

103:10他沒有按我們的罪過待我們,也沒有照我們的罪孽報應我們.

103:11天離地何等的高,他的慈愛向敬畏他的人也是何等的大!

103:12東離西有多遠,他叫我們的過犯離我們也有多遠!

我們雖然是「罪人」,但上帝還是看我們為「人」,不是以我們的「罪」來代替了我們這個「人」.在我們的眼中,我們還看見得罪我們的人,或者與我們不同意見的人是一個「人」嗎?今天的福音經課耶穌以另一個角度演繹相同的教導:

18:23天國好像一個王要和他僕人算帳.

18:24才算的時候,有人帶了一個欠一千萬銀子的來.

18:25因為他沒有甚麼償還之物,主人吩咐把他和他妻子兒女,並一切所有的都賣了償還.

18:26那僕人就俯伏拜他,說:『主啊,寬容我,將來我都要還清.』

18:27那僕人的主人就動了慈心,把他釋放了,並且免了他的債.

18:28「那僕人出來,遇見他的一個同伴欠他十兩銀子,便揪著他,掐住他的喉嚨,說:『你把所欠的還我!』

18:29他的同伴就俯伏央求他,說:『寬容我吧,將來我必還清.』

18:30他不肯,竟去把他下在監裡,等他還了所欠的債.

18:31眾同伴看見他所做的事就甚憂愁,去把這事都告訴了主人.

18:32於是主人叫了他來,對他說:『你這惡奴才!你央求我,我就把你所欠的都免了,

18:33你不應當憐恤你的同伴,像我憐恤你嗎?』

18:34主人就大怒,把他交給掌刑的,等他還清了所欠的債.

18:35你們各人若不從心裡饒恕你的弟兄,我天父也要這樣待你們了.」

釋放與監禁有甚麼分別呢?曼德拉曾經說過:「當我走出囚室,經過通往自由的監獄大門時,我知道自己若不能把悲傷與怨恨留在身後,那麼我其實仍在獄中.」

願上主幫助我們.



\newpage

\section{信仰的重軛}
\label{sec:382}
\textbf{許立中先生 - 信仰的重軛}
\newline
\newline
\hyperref[sec:381]{< < < PREV SERMON < < <}
~
\hyperref[sec:toc]{[返目錄]}
~
\hyperref[sec:383]{> > > NEXT SERMON > > >}


馬太福音 11:16-19,25-30
\newline
\begin{longtable}{cl}
\hline
\hline
章節 & 經文 (和合本修訂版)\\
\hline
11:16 & \begin{tabularx}{0.7\textwidth}{X} 「我該用甚麼來比這世代呢？這正像孩童坐在街市上向同伴呼喊： \end{tabularx} \\ \\ \relax
11:17 & \begin{tabularx}{0.7\textwidth}{X} 『我們為你們吹笛，你們不跳舞；我們唱哀歌，你們不捶胸。』 \end{tabularx} \\ \\ \relax
11:18 & \begin{tabularx}{0.7\textwidth}{X} 約翰來了，既不吃也不喝，人們就說他是被鬼附的； \end{tabularx} \\ \\ \relax
11:19 & \begin{tabularx}{0.7\textwidth}{X} 人子來了，也吃也喝，他們又說這人貪食好酒，是稅吏和罪人的朋友。而智慧是由它的果子來證實的。」 \end{tabularx} \\ \\ \relax
11:20 & \begin{tabularx}{0.7\textwidth}{X} 那時，耶穌在一些城行了許多異能。因為城裡的人不肯悔改，他就責備那些城說： \end{tabularx} \\ \\ \relax
11:21 & \begin{tabularx}{0.7\textwidth}{X} 「哥拉汛哪，你有禍了！伯賽大啊，你有禍了！因為在你們中間所行的異能若行在推羅、西頓，他們早已披麻蒙灰悔改了。 \end{tabularx} \\ \\ \relax
11:22 & \begin{tabularx}{0.7\textwidth}{X} 但我告訴你們，在審判的日子，推羅和西頓所受的，比你們還容易受呢！ \end{tabularx} \\ \\ \relax
11:23 & \begin{tabularx}{0.7\textwidth}{X} 迦百農啊，你以為要被舉到天上嗎？你要被推下陰間！因為在你那裡所行的異能，若行在所多瑪，它還可以存留到今日。 \end{tabularx} \\ \\ \relax
11:24 & \begin{tabularx}{0.7\textwidth}{X} 但我告訴你們，在審判的日子，所多瑪地方所受的，比你們還容易受呢！」 \end{tabularx} \\ \\ \relax
11:25 & \begin{tabularx}{0.7\textwidth}{X} 那時，耶穌說：「父啊，天地的主，我感謝你！因為你把這些事向聰明智慧的人隱藏起來，而向嬰孩啟示出來。 \end{tabularx} \\ \\ \relax
11:26 & \begin{tabularx}{0.7\textwidth}{X} 父啊，是的，因為你的美意本是如此。 \end{tabularx} \\ \\ \relax
11:27 & \begin{tabularx}{0.7\textwidth}{X} 一切都是我父交給我的；除了父，沒有人知道子；除了子和子所願意啟示的人，沒有人知道父。 \end{tabularx} \\ \\ \relax
11:28 & \begin{tabularx}{0.7\textwidth}{X} 凡勞苦擔重擔的人都到我這裡來，我要使你們得安息。 \end{tabularx} \\ \\ \relax
11:29 & \begin{tabularx}{0.7\textwidth}{X} 我心裡柔和謙卑，你們當負我的軛，向我學習；這樣，你們的心靈就必得安息。 \end{tabularx} \\ \\ \relax
11:30 & \begin{tabularx}{0.7\textwidth}{X} 因為我的軛是容易的，我的擔子是輕省的。」 \end{tabularx} \\ \\
[1ex]
\hline
\hline
\end{longtable}
講題:信仰的重軛　The Unbearable burden of Faith

經文:馬太福音11:16-19,25-30

講員:許立中先生

場合:香港中文大學崇基學院禮拜堂

日期:2017年9月24日

「我可用甚麼比這世代呢?好像孩童坐在街市上招呼同伴,說:我們向你們吹笛,你們不跳舞；我們向你們舉哀,你們不捶胸.約翰來了,也不吃也不喝,人就說他是被鬼附著的；人子來了,也吃也喝,人又說他是貪食好酒的人,是稅吏和罪人的朋友.但智慧之子總以智慧為是.」

「那時,耶穌說:父啊,天地的主,我感謝你!因為你將這些事向聰明通達人就藏起來,向嬰孩就顯出來.父啊,是的,因為你的美意本是如此.一切所有的,都是我父交付我的；除了父,沒有人知道子；除了子和子所願意指示的,沒有人知道父.凡勞苦擔重擔的人可以到我這裡來,我就使你們得安息.我心裡柔和謙卑,你們當負我的軛,學我的樣式；這樣,你們心裡就必得享安息.因為我的軛是容易的,我的擔子是輕省的.」

「凡勞苦擔重擔的人可以到我這裡來,我就使你們得安息.」這節經文往往被誤讀為對罪人,未信者的福音呼召.但倘若我們將它放回其前文後理去理解,就會知道耶穌在那裡所講的,其實是作為一個有信仰的人的難處.

信仰其中一個難處,就是為了利便傳播,教會往往傾向將它說得太過簡單.你明白也好,不明白也不要緊,據說你只需要憑著單純的信心,接受那堆基要的教義,就可以成為一個「重生得救」的基督徒.

這確實有點匪夷所思,也脫離我們的生活經驗.我們不明白,明明是超級理性化的教義陳述,怎麼一下子成為「改變生命」的口訣?如果我們對信仰有甚麼不明白,往往亦得不到理性上的認真對待,反而被視為「不夠信心」.

傳統上,教會習慣單向地從個人,靈性的角度去理解一切問題,因此亦往往只能得出一些近乎陳腔濫調的結論.譬如我們傾向將一切問題歸因到個人跟上帝之間的關係:我們之所以遇到困擾我們的家庭,工作或者感情問題,歸根究底,都是因為我們跟上帝的關係出現問題；而要改善跟上帝的關係,當然就是多讀聖經和多祈禱.

我說它近乎陳腔濫調,是因為它幾乎是一隻「百搭」的牌,可以解決任何信仰上的問題.而我們知道,一個可以解答所有問題的答案,實際上就是一個解決不了任何問題的答案.

有時,我們確實將太多的寄望放在「讀經祈禱」上.徒八:26-40提到腓利在耶路撒冷往迦薩的路上,遇到一個埃提阿伯的太監,在車上讀先知以賽亞的書.腓利問他:你所念的,你明白嗎?太監說:沒有人指教我,怎能明白呢?於是腓利上車,並沒有批評他「不夠信心」,而是向他講解當中的意思.最後,也不用上洗禮班,腓利就在就近的水源為他施洗.

「讀經」不是一個答案,因為我們的問題正正在於不明白聖經；根本不明白,再讀一百次也未必會明白.有時甚至愈勤力去讀,誤解和成見反而愈牢固.錯誤和偏見,是需要被指正而不是重複reinforce.因此,朝朝上教堂祈禱,並不保證一個人不會行差踏錯；而有些人明明是剛愎自用,卻誠實地以為自己「每一步都遵從自己的信仰」,行在上主的心意中.

過分強調個人,內心,有時會令我們忽略一些明顯而同樣重要的其他因素,譬如社會當時的文化氛圍,所屬行業階層的次文化,信仰群體的屬靈氣質,甚至是個人性格上的弱點和盲點,都影響著信徒對信仰和生活的理解.而信徒卻往往將一切歸因於信仰,以為自己一切的言行思想,都服膺在上主的掌管之下.這也解釋了為甚麼兩個同樣認真和敬虔的基督徒,可以站在完全對立的立場上寸步不讓.因為無論是藍營或綠營,黃絲或藍絲,他們都覺得自己是站在真理的一方.

其實作基督徒也實在不容易.

我認識一位在基督教機構負責行政工作數十年的姊妹,早幾年才退休.有一次聚會,她跟朋友討論一些問題,其實也不是甚麼大不了的問題,但她卻激動得滿面通紅,青筋暴現.平伏下來之後,朋友問她為甚麼會這樣?她說自己多年來在工作上要跟不同的人開會,很多時遇見一些很「離譜」的人和處事的方式,但礙於自己是基督徒,不可以發火甚至出聲,於是就經常「頂住度氣」.久而久之,有時在一些討論中,就很容易不能自制地產生強烈的情緒反應.

有時做基督徒難,是因為當你告訴人你是基督徒的時候,他們就會用另一種眼光來看你.而你的一舉一動,一言一行,都不能稍有差池.更要命的,是有時連我們自己也是這樣看自己,要求自己.

還有,假如你有一個很難相處的親人,或者一個很難頂的同事,那麼雖然你們之間的關係不好不一定是你的問題,但周圍的人多數會說:「即係,你地基督徒又咪係咁!」又如果你的下屬疏忽職守,你提醒責備他幾句,他轉過頭就向其他同事數落你,如何沒有愛心,簡直不像一個基督徒!又如果你從事公職,有時要執行一些自己也未必完全理解或者認同的行動或政策,情況就更複雜.所以,有時我也體諒到有些基督徒不輕易在人面前承認自己基督徒的身分,以免令原本已經難搞的問題,變得更加複雜.

成為基督徒之後,有時我們確實要應付不少奇怪的要求或指控:「你們基督徒不是應該這樣的嗎?」,「怎麼你們基督徒都可以那樣!?」更可笑的是,那些向基督徒提出質詢的人,對信仰其實毫無認識.他們只是想當然地幻想基督徒應該怎樣.但為免「失見證」,很多時我們就任由他們批評指點,教我們應該如何做一個基督徒.而我要說,這樣做基督徒確實是很辛苦的,同時亦繼續強化一般人對信仰的誤解.

但更辛苦的,有時反而是來自信仰群體的壓力.

「我可用甚麼比這世代呢?好像孩童坐在街市上,招呼同伴,說:我們向你們吹笛,你們不跳舞.我們向你們舉哀,你們不捶胸.約翰來了,也不喫,也不喝,人就說他是被鬼附著的.人子來,也喫,也喝,人又說他是貪食好酒的,是稅吏和罪人的朋友」(太十一:16-19)

基督徒時常強調「生活見證」,做每一件事都怕會絆倒人.我們最常引用的經文是林前十:23:「凡事都可行,但不都有益處.凡事都可行,但不都造就人.」

推到一個極端,我們甚麼事都不敢做,甚麼事都要看人面色；恐怕一下不為意就會絆倒人.到最後,聽見人吹笛,我們不敢跳舞,恐怕有人會認為不合聖徒體統；有人舉哀,我們不敢搥胸,恐怕被人誤會我們失去盼望.再直接點說,既然凡事都有絆倒人的可能,那在實際上,我們說「凡事都可行」其實不是認真的.

而因為我們不斷為了體貼人的軟弱而否定自己的想法,我們跟自己真正的感受便愈來愈脫節.到一個地步我們不再知道自己喜歡甚麼,不喜歡甚麼,想做甚麼,不想做甚麼.簡單來說,我們不再知道自己是誰,而弟兄姊妹亦不知我們裡面實在是個怎樣的人.我們是迷失甚至是消失在教會的「大我」裡面,成為一個斯斯文文,客客氣氣而面目模糊的「基督徒」.有性格的人在教會裡並不受歡迎.事實上,我想就連主耶穌亦難以符合那些所謂「基督徒的標準」!

這段經文告訴我們一個事實:「約翰來了,也不喫,也不喝,人就說他是被鬼附著的.人子來,也喫,也喝,人又說他是貪食好酒的人,是稅吏和罪人的朋友.」

跟一些有道德問題的人交往,也是一個問題.好幾年前,一位牧師跟一些有黑社會背景的人交往,最後帶了他們信主還替他們施洗.有些輿論就說他的做法「具爭議性」,而教會雖未至於割蓆,反應亦相當謹慎,不敢正面肯定他的努力.換句話說,無論我們做甚麼,不做甚麼,都會有人因我們而跌倒!那我們做還是不做呢?主耶穌的態度是:「凡不因我跌倒的,就有福了!」(太十一:6)

我並不是要將耶穌的說話和保羅的教訓針鋒相對；我要指出的是,就算是最正確的教導,都會有矯枉過正的可能.對於做或者不做,我們都必須誠實地在上帝面前作出一個屬於自己的決定.因為你做,可以是出於勇氣,也可能是出於血氣；你不做,可能是因為謙讓,亦可能是因為懦弱!可以說,沒有一段經文可以擔保我們萬無一失.

加五:1:「基督釋放了我們,叫我們得以自由,所以要站立得穩,不要再被奴僕的軛挾制.」而保羅在這裡所指的「奴隸的軛」,是來自信仰的群體,就是加拉太教會裡面的保守勢力!這也是保羅寫信給加拉太教會的原因.

在信仰的實踐裡,基督徒有很多「是就是…不過…」.這堆理論跟實踐的差異,將我們的信仰閹割成「講就天下無敵,做就有心無力」的空言.

而教會裡面躁動的青年,聽出講台的宣講不過是煞有介事的虛張聲勢,所謂‘Quote part of the Scripture and you can make the Bible say anything’,並無實質的屬靈權柄和真理基礎,也就明目張膽地挑戰教會的權威,反正你懂得引用聖經我亦懂得挑戰你的詮釋.或許就像人大釋法那樣,宗教改革運動所高舉的Sola Scriptura,就逐漸消失於信徒的互相攻擊之中.

信仰的理論,返得三兩年教會,大概都已經可以掌握得七七八八.我們迴避談論的,是這些理論確實的聖經根據和神學基礎.譬如,基督徒是不是不可以吸煙喝酒?傳福音是不是基督徒人生最終極的意義和目的?這些都是我們從來不敢討論的價值標準,甚至視之為信仰的底線.

當然,吸煙可以致癌是健康常識,但如果你在旺角或銅鑼灣上班,那致癌的機會亦相當高.如果你喜歡吃東坡肉或豬大腸,那亦可以相當致命.我的意思是,我們可以以健康的理由勸人戒煙或少吃豬大腸,但不需要將它提升到一個道德甚至信仰的層次.至於酒,如果你去過幾次基督教的婚禮,你就會知道主耶穌在世上行的第一件神蹟,就是將水變酒!我尊重不吸菸不喝酒作為揀女婿或擇偶條件,但這確實不是成為基督徒的信仰門檻.

信仰的問題是可以討論的.一個問題是否涉及信仰,或者怎樣涉及信仰,亦可以透過神學反省來處理.

說到這裡,我突然明白當日的宗教領袖為何對耶穌恨之入骨.他們心想:我們辛辛苦苦建立一個忠心順服,聽教聽話的信仰群體,你就來說行為完全不重要,人得救是靠賴恩典.你這樣說,我們還可以怎樣要求他們循規蹈矩?我完全明白牧會者的擔心.

我想當信徒未有判斷能力,你可以給他們一些簡單的正確答案；但他們開始長大成人,進入社會,你不能再控制他們接觸到些甚麼,亦無法保護他們不受到世界的引誘和影響.這個時候,他們需要的,再不是一些正確的答案,而是如何找到一些他們能夠真心信服的答案.而在這方面,我只能說教會實在有太多未付的帳單.

至於傳福音是不是基督徒人生最終極的意義和目的?我早年曾經寫過一篇文探討這個問題,叫做「大使命有多大?」企圖從一個常識的角度,亦不違反聖經解釋的原則,去為這個問題重新定位.無論如何,作為基督徒,我們能夠做的,亦只能回歸主耶穌有關的教訓.

我喜歡主耶穌說:你們是「世上的鹽」,「世上的光」.只要鹽沒有失了味,光是放在燈臺上而不是放在斗底下,我們的存在本身就是基督的見證,而不是我們刻意去做些甚麼.

這當然跟我們個人和集體的佈道工作沒有衝突.事實上,有具體生命見證的支撐,口述的福音才不致於變得理論化和抽象.真實的生命,是一切佈道事工的基礎和憑據.

相反來說,假如我們的生命貧瘠和蒼白,那麼縱然我們能說萬人的方言,並天使的話語,將天上的雀仔都引了下來,我們仍然只不過是鳴的鑼,響的鈸；沒有一個湧流的生命,傳福音就像莎士比亞在Macbeth中對生命的慨歎:It is a tale told by an idiot, full of sound and fury, Signifying nothing!

返回勞苦擔重擔的問題.將基督徒壓到透不過氣的,很可能不是外在的壓力,而是我們自己不切實際的理想和執著.

記得參加過一個聚會,領唱突然要求會眾安靜,作以下的「反省」:「…我們說的每一句話,做的每一件事,是否都為了主?是否都合乎主的心意?」

每一句話,每一件事?答案也不用想吧,怎可能「是」呢?除了道成肉身的那一位,又有誰可以一言一行都合乎天父的旨意呢?就是主耶穌在世的時候,亦經歷過不少掙扎,需要獨自退到曠野祈禱,處理那些矛盾.他甚至曾經祈求天父「倘若可行,求你叫這杯離開我…」

有時過度的自省,對靈命未必有好處.就以這個例子來說,在上述的「反省」之後,領唱還帶領各人「立約」:「讓我們立志,今生所說的每一句話,所做的每一件事,都全然為主.」很崇高的想法和說法,卻明顯不可能辦到.信誓旦旦地說出一個明知道不可能實踐的誓言,對靈命有何益處?

宗教的陳腔濫調,如同拙劣的廣告,以為將「事實」說成百份之一百五十,經過五十個百分點的損耗,就能得到百份百的效果.我不反對自省,也不反對立約,正如婚約中的誓言:「無論是富貴或貧窮,疾病或健康,都承諾愛對方,照顧於對方」,這至少是可以實踐；但若將承諾改為:「我願意每時每刻都想著你,每個思想每個行為都不忘記愛你」,那就不單只是虛情假意,簡直就是自欺欺人了.

主耶穌說:「是,就說是；不是,就說不是；若再多說,就是出於那惡者.」(太五:37)烹飪,調味是必須的；然而過度加鹽添醋,整碟菜都要搞砸.

傳道書五:2:「你在上帝面前不可冒失開口,也不可心急發言；因為上帝在天上,你在地下,所以你的言語要寡少.」講得太多,很容易就會變成花言巧語.讓我們小心宗教的術語,盡量戒除一切自己不是真正相信的花言巧語,無論聽起來如何「屬靈」.因為不經大腦地說,水過鴨背地聽.久而久之,這類說話也就變成了毫無實質意義的空言.

這裡牽涉兩個重要的生命教育的問題:

宗教術語令人心靈麻木；完美/理想主義滋生宗教偽善

被政府以罪犯身分拘禁,兩個月前在獄中病逝的諾貝爾和平獎得主劉曉波在他一本著作中指出:每個時代,每個民族,每個人都有自己的理想.理想之於人類的推動作用,主要不在於理想本身的價值,不管這理想是多麼美好；亦不在於它總有一天能夠實現,不管人們多麼相信；理想的意義和價值,乃在於它能夠為我們提供一個可以努力並且值得努力的目標,擺脫現實的束縛.

用劉曉波的講法:理想與現實之間的關係,永遠是對立的,衝突的,表現出精神與肉體之間,無限與有限之間,永恆與短暫之間,企求超越與無法超越之間的矛盾.理想的核心價值,並不在理想本身,而在於理想對現實的參照價值,激勵人類不斷向自身的現實發出質問.

用一個大家比較熟悉的例子:主禱文說「願你的國降臨,願你的旨意行在地上,如同行在天上.」這是一個祝願和理想,也是一個指導我們應該怎樣在世上生活的盼望；但是在主再來之前,這理想從來未曾亦永遠不會成為現實.如果你因此以為教會就是上帝的國度,人間天堂,那你很快會理想幻滅.

劉曉波認為,如果人類僅僅陶醉於理想之中,將理想本身看為現實人生的絕對價值,那麼理想便會失去它的激勵作用,而變成了麻醉劑.它使人不敢或不願意正視現實的艱難和不足,對百孔千瘡的人生採取逃避的態度,沉溺於永遠無法實現的烏托邦之中,每日過著一種不真實和自欺欺人的生活.事實上這也是馬克思對宗教最主要的批評.

而這個情況曾經在理想主義的共產中國出現過.中共建國的初期,人民都生活在浪漫而不真實的憧憬當中,喊一些偉大到不得了但無可能實踐的口號；它推翻了宗教的救世主,轉而寄望一位世俗的救世主,憧憬著人間的天國就在前面的轉角.

劉曉波更以中國儒家那種「威武不能屈,貧賤不能移,富貴不能淫」的人格理想作為舉例,指出如果這種人格理想是一種參照,是我們努力的目標,那麼它會發出奪目的光輝.但是,中國人的愚昧,在於把它作為現實人格的絕對標準,去要求人必須做這樣的人,而人自身的弱點,又使人不可能進入此境界.因此,這種人格理想便造就成批的偽君子,表面上是這種人格,背後卻無惡不作；而比較誠實的人,就因為看見在現實中根本達不到這個理想,變得憤世嫉俗.而這兩類人,我們在教會都見得太多了.

單就這段論述來說,劉曉波是準確而深切的表達了保羅在羅馬書特別是第七章對自身處境的掙扎.

反觀今日教會的主流教導,卻正正是標榜一種「合乎上帝標準」(即完美無暇)的理想人格,把它看為現實人格的絕對標準,去要求基督徒達到(「好就好,不過做不到」),就像主耶穌說:「把難擔的重擔,捆起來擱在人的肩上」.教會若不及早醒覺,堵住這個破口,恐怕我們一切佈道的努力,至終只是應驗主耶穌的責備:「你們走遍洋海陸地,勾引一個人入教；既入了教,卻使他作地獄之子,比你們還加倍」!

25:「那時,耶穌說:父啊,天地的主,我感謝你,因為你將這些事,向聰明通達的人,就藏起來,向嬰孩,就顯出來；父啊,是的,因為你的美意本是如此.」

我們實在需要重新反省,「效法基督」到底是甚麼意思.聖經說主耶穌道成肉身來到世界,我們卻往往將成為肉身的基督削骨去肉,推回上天堂,將它變成一個完美的抽象符號.然後就說「是就是,不過我卻不是耶穌啊!」

聖經說「耶穌的智慧和身量,並上帝和人喜愛他的心,都一齊增長.」而既說「增長」,就必然還有「增長的空間」(room for improvement).絕對的完美,是抽象和靜止的,用佛家的說法就是「涅盤」；道成肉身的完美,卻是有成長的動力和進步的空間.保羅在林後十三:9說:「我們所求的,就是你們作完全人」；在腓三:16他卻補充說:「然而我們到了甚麼地步,就當照著甚麼地步行.」

有時我們不知道如何去教導下一代,是因為我們本身亦感到迷惘.主耶穌說他的軛是容易的,他的擔子是輕省的,因為相對於其他人加諸我們身上的重軛,無論是社會的,學校的,道德的,宗教的,我們只需要向那位獨一的真神,愛我們的天父交代和負責.而主耶穌已經為我們立下榜樣,以致我們可以「心裡柔和謙卑」地面對!

賽三十:15「主耶和華以色列的聖者曾如此說:你們得救在乎歸回安息.你們得力在乎平靜安穩.你們竟自不肯.你們卻說:不然,我們要騎馬奔走.所以你們必然奔走.又說:我們要騎飛快的牲口.所以追趕你們的,也必飛快.」

用今日的話來說:「你們得救在乎返回我的懷抱.你們得力在乎唔好搞咁多野.你們竟然不肯,說,我們要努力做到最好!所以你地必然做到筋疲力盡!你們又說:我們要盡一切努力去達到上帝的標準!所以你地必然做到趴響度,而你們的最好,有排都未到你們自己想像的目標,更不要說上帝的標準!」

在這個背景之下,我們重新理解主耶穌的提醒和安慰:「凡勞苦擔重擔的人,可以到我這裡來,我就使你們得安息.我心裡柔和謙卑,你們當負我的軛,學我的樣式,這樣,你們心裡就必得享安息.因為我的軛是容易的,我的擔子是輕省的.」(太十一:28-30)



\newpage

\section{基督無能為力嗎}
\label{sec:383}
\textbf{白德培牧師 - 基督無能為力嗎?}
\newline
\newline
\hyperref[sec:382]{< < < PREV SERMON < < <}
~
\hyperref[sec:toc]{[返目錄]}
~
\hyperref[sec:384]{> > > NEXT SERMON > > >}


腓立比書 2:1-13
\newline
\begin{longtable}{cl}
\hline
\hline
章節 & 經文 (和合本修訂版)\\
\hline
2:1 & \begin{tabularx}{0.7\textwidth}{X} 所以，在基督裡若有任何勸勉，若有任何愛心的安慰，若有任何聖靈的團契，若有任何慈悲憐憫， \end{tabularx} \\ \\ \relax
2:2 & \begin{tabularx}{0.7\textwidth}{X} 你們就要意志相同，愛心相同，有一致的心思，一致的想法，使我的喜樂得以滿足。 \end{tabularx} \\ \\ \relax
2:3 & \begin{tabularx}{0.7\textwidth}{X} 凡事不可自私自利，不可貪圖虛榮；只要心存謙卑，各人看別人比自己強。 \end{tabularx} \\ \\ \relax
2:4 & \begin{tabularx}{0.7\textwidth}{X} 各人不要單顧自己的事，也要顧別人的事。 \end{tabularx} \\ \\ \relax
2:5 & \begin{tabularx}{0.7\textwidth}{X} 你們當以基督耶穌的心為心： \end{tabularx} \\ \\ \relax
2:6 & \begin{tabularx}{0.7\textwidth}{X} 他本有神的形像，卻不堅持自己與神同等； \end{tabularx} \\ \\ \relax
2:7 & \begin{tabularx}{0.7\textwidth}{X} 反倒虛己，取了奴僕的形像，成為人的樣式；既有人的樣子， \end{tabularx} \\ \\ \relax
2:8 & \begin{tabularx}{0.7\textwidth}{X} 就謙卑自己，存心順服，以至於死，且死在十字架上。 \end{tabularx} \\ \\ \relax
2:9 & \begin{tabularx}{0.7\textwidth}{X} 所以神把他升為至高，又賜給他超乎萬名之上的名， \end{tabularx} \\ \\ \relax
2:10 & \begin{tabularx}{0.7\textwidth}{X} 使一切在天上的、地上的和地底下的，因耶穌的名，眾膝都要跪下， \end{tabularx} \\ \\ \relax
2:11 & \begin{tabularx}{0.7\textwidth}{X} 眾口都要宣認：耶穌基督是主，歸榮耀給父神。 \end{tabularx} \\ \\ \relax
2:12 & \begin{tabularx}{0.7\textwidth}{X} 我親愛的，這樣看來，你們向來是順服的，不但我在你們那裡，就是我現在不在你們那裡的時候更是順服的，就當恐懼戰兢完成你們自己得救的事； \end{tabularx} \\ \\ \relax
2:13 & \begin{tabularx}{0.7\textwidth}{X} 因為是神在你們心裡運行，使你們又立志又實行，為要成就他的美意。 \end{tabularx} \\ \\
[1ex]
\hline
\hline
\end{longtable}
講題:基督無能為力嗎?　A powerless Christ?

經文:腓立比書2:1-13

講員:白德培牧師

場合:香港中文大學崇基學院禮拜堂

日期:2017年10月1日

Phil 2:1-13

If then there is any encouragement in Christ, any consolation from love, any sharing in the Spirit, any compassion and sympathy, make my joy complete: be of the same mind, having the same love, being in full accord and of one mind. Do nothing from selfish ambition or conceit, but in humility regard others as better than yourselves. Let each of you look not to your own interests, but to the interests of others. Let the same mind be in you that was in Christ Jesus,

who, though he was in the form of God,

did not regard equality with God as something to be exploited,

but emptied  himself,

taking the form of a slave, being born in human likeness.

And being found in human form,

he humbled himself and became obedient to the point of death - even death on a cross.

Therefore God also highly exalted him

and gave him the name that is above every name,

so that at the name of Jesus

every knee should bend, in heaven and on earth and under the earth,

and every tongue should confess

that Jesus Christ is Lord, to the glory of God the Father.

Therefore, my beloved, just as you have always obeyed me, not only in my presence, but much more now in my absence, work out your own salvation with fear and trembling; for it is God who is at work in you, enabling you both to will and to work for his good pleasure.

1.      As an introduction: A visit to Cluny and Taizé

A few weeks ago, in August, I was working with a production team from Shanghai on a documentary about the 500 year anniversary of the Reformation. We visited various places of the Reformation in Germany, Switzerland, and France. Our final stop was in Cluny and Taizé in France. These places didn’t play a role in the Reformation; they don’t actually belong to the Reformation history. But they remind us that Reformation should not be confined solely to something that had happened 500 years ago. Reformation has happened before and happened since then.

The Cluny movement began in the 10th century. It was a spiritual and monastic movement that wanted to liberate the church from interference by worldly powers and from control by political leaders. It was so successful that within one hundred years it turned the relationship between state and church, between political rulers and spiritual leaders upside down. While previously political leaders tended to dominate ecclesial affairs, the church now began to dominate and control political leaders. One point of constant conflict was the appointment of bishops. When a great king of Germany, Henry IV, insisted on appointing bishops, the pope excommunicated him and the king had to do penance before the pope. He had to go to the castle of Canossa, high in the mountains of Italy, and wait 3 days in the coldness of winter, until the pope finally decided to readmit him to the church. The famous penance of Canossa in 1077 was the peak of the Cluny movement.

The church built in Cluny was for several centuries, from the 12th to the 16th century, the largest church in the world. It was built as a place of constant worship and glorification of God.

Nowadays we can only see the ruins of this formerly huge church.

What happened?

The church turned into an instrument of power and was more and more perceived as an expression of a corrupt church. What was originally an expression of a will to glorify God, turned into self-confidence and eventually into pride. A church that is powerful in the world is ultimately subject to the same corruption as the world. In the French Revolution at the end of the 18th century people stood up against the church and destroyed it. First it was turned into a store room, later into a quarry where people could find stones to build their houses.

At the entry of the museum of Cluny, which gives access to the ruins of this formerly so enormous church there is a very profound sentence - maybe a form of late repentance for the attempted splendor of the church. It says: “The monk is the one who wants to tell to the world a word (message) that is so enormous that he is obliged to remain silent.”

Indeed, it is simplicity, fragility, or silence that best reflect God’s glory.

From Cluny it is only 10km to Taizé. Taizé was founded by a Christian brother from Switzerland during the Second World War. It is a place of worship in unity across the divisions within Christianity. It is a place of reconciliation.

Cluny and Taizé stand for two radical opposites within the history of the church. From Cluny to Taizé we cross from one side of the church to the radical other. Here we encounter the church in its most simple form. Taizé stands for humility and fragility. In the entry pavilion there is a picture of the founder, Brother Roger. It is a picture taken as he is about to embrace a woman after the worship. It turns out that this woman is the person who is about to kill him. She was a mentally disturbed woman who pulls out a knife and kills him. It turns out that Brother Roger embraces his own murderer.

This is really “the same mind that was in Christ Jesus”.

Cluny and Taizé - power and humility: This is also what the Bible text of today is about - the power and humility, powerlessness of Christ.

2.      The Philippian hymn

Today’s reading is one of the most famous texts in the Bible where Paul teaches about how to understand the ministry of Jesus Christ. The text is like a creed for the early church. Paul tells about the meaning of Christ and the journey that Christ took.

·         Christ did not strive for equality with God.

·         He emptied himself and took on the form of a slave and was born as human.

·         He gave up on his divine character. Paul calls this process a process of self-emptying - in Greek it is called kenosis.

·         He humbled himself and became obedient to the point of death - even death on the cross.

Paul reminds us that it is only because of going through suffering and humiliation that Christ gains victory. And this journey of Christ is presented as a model for each of our own spiritual journey. He exhorts us “Let the same mind be in you that was in Christ Jesus”.

What does it mean for us as a church to follow Christ in this process of self-emptying?

Maybe no other text has so much put humility into the center of Christian virtues. Humility is more than simply being polite - it is a genuine attitude of knowing that one is in no way better than what the world regards as a failed person. This is really going at the heart of our value system. What does it mean in our community life? And what does it mean in our personal life?

3.      Humility as a corporate virtue

I find that the vision of Christ going through a process of self-emptying actually stands quite in contrast to what we aim at in our church life. Most of us will do our best to have the church grow - to proclaim the gospel and to invite people to join - that is indeed part of what we are called for. Yet there is some ambiguity to this - because it ultimately means that we naturally move towards a large and more powerful church. Large churches are admired for the successful business model they apply or for the dedicated offering of its members, for the variety of mission programs offered. So-called mega-churches have arisen in many modern cities - particularly famous are those in South Korea. But we also have some smaller mega-churches in Hong Kong. I do in no way want to belittle these churches. They grew out of dedication and commitment and of a natural tendency to grow. Doesn’t it make a difference whether many people believe? Isn’t it more uplifting and encouraging worshiping within a large group of committed fellow Christian brothers and sisters? But large churches are easily misleading us into believing that they are somehow better churches. They are neither better nor worse than a small or a poor church with little outreach and with possibly more internal problems. It’s only that large churches may more easily create the illusion that they are strong.

Yet, it seems to me that large churches stand in some contrast to the vision that Paul offers in his letter to the Philippians. The model he offers is not based on strength and power but on humility, brokenness and suffering. It is only through these that God lifts us up - as he lifted Christ up. Christ-likeness can much more be found in a community that reflects the cross. Maybe this is one of the reasons why I am so attracted to the ministry in prison because very obviously it is a community that reflects the up-and-down journey that Phil 2 expresses.

4.      Actuality of humility in our personal faith

Now, what does kenosis - self-emptying - in our personal life mean?

Of course, first it means to show humility in your encounter with other people. It is indeed worth to think about what humility means in our relationships. Do we really treat our staff - if we are in a senior position - as equal? Do we treat the many workers at the low end of the pay scale as equal and with respect? Or, how do we arrange our time? Does our time arrangement reflect equal respect to everyone? How do we organize our social relations? Are they designed only to bring us advantage?

It is surely worthwhile to think about such questions, but it is equally necessary to remember that humility is actually a universal value. It is not just Christians who propagate humility. And I think all of us know - at least in theory - enough about what humility is. And all of us need to be constantly reminded not just to propagate humility but to actually practice it. So, today’s Bible text shall be such a reminder.

What is, however, more interesting, is how the idea of humility actually challenges our faith - this is what I want to consider now.

Do we really believe in a God who revealed himself in the brokenness of Christ? Do we not again and again try to find refuge in a powerful god?

There is one theologian who put particular emphasis on Christ’s process of self-emptying: Dietrich Bonhoeffer. He is one of the most famous theologians of the 20th centuries - famous not only for his writings but also for the way his life reflected his faith, and even more famous for his martyrdom. He was a professor of theology and during Hitler’s government he understood that it was his duty as a Christian to stand up against Hitler. He was involved in a plan to overthrow him and was arrested and spent 2 years in prison. Only 31 days before the Second World War ended he was killed. Throughout his time in prison he showed confidence - not confidence in release, but confidence in God’s comforting presence.

During his time in prison, he wrote the famous many letters to his outside friends. These letters are now published as Letters and Papers from Prison - one of the most important theological books of the 20th century. For Bonhoeffer, following Christ in Christ’s process of self-emptying meant first to turn towards the world and becoming worldly. That doesn’t mean to accept the logic and values of the world but to get involved in the struggles of the world. Often theologian like to remain on a theoretical level. He understood that the belief in Christ’s incarnation led Christians on a path that called them to follow Christ into the messiness of the world.

But even more importantly, it meant for him to give up the belief in God as a powerful God. He criticized a kind of belief that looked at God’s intervention like a child looking for the help of his father: whenever we can do it on our own, we don’t rely on God. We forget God. We do it ourselves. And as soon as we fail or are in troubles, we cry to God for help. Bonhoeffer regarded such a faith as wrong.

i.            It is wrong because it is a childish faith. And as we grow up, we cry less and can do it on our own. We don’t always need to cry for help to God. It is like being attacked and crying for the big brother helping us. As a result of our growing up, we gradually push God out of the world.

ii.            It is also wrong because the more we are able to do things ourselves and explain things without reference to God, the less we care about God. Look for instance at medicine: in old times, people didn’t know the causes for many illnesses. They interpreted them in a religious framework, regarded them possibly as a punishment from God. Accordingly, they called on God to heal them. As medical science grew, the hypothesis of God was not necessary anymore because an illness could be explained as caused by some germs. And God disappeared. In this way, God gradually became irrelevant, until he only had some remaining small area of credibility at the margins of our life.

iii.            And thirdly, it is wrong because it is not biblical. The Philippian hymn, Phil 2:7, reminds us that Jesus gave up his divine powers as he entered this world. The cross means the revelation of God’s powerlessness in this world. It is a radical break with a faith that holds on a dominion or power pattern; that regards God simply as the more powerful friend who would protect me from my enemies and return their bullying me.

For Bonhoeffer, this kind of God was pushed out of the world on the cross. Bonhoeffer maintained faith in God in the midst of our life, not just at the margins of our life. The cross stands for God’s radical acceptance of his own powerlessness in this world. Bonhoeffer gave up on this because he realized that Christian faith is something else than believing in a God who turns out or is expected to be more powerful than all others. He didn’t give up on his faith. But he understood the mystery of Christ’s self-emptying.

Indeed, Christ does not help us through his strength. Being a Christian is not about reviving the idea of a God who is powerful in this world but rather to accept the powerlessness of God in this world. It is the belief that God joined us in suffering in the same way as God suffers and necessarily does so in the godlessness of the world. To be a Christian is to find strength in that God also only could suffer in this world and cannot prove his divinity by powerfully overcoming the world’s godlessness.

Or as Bonhoeffer said: “The God who is with us is the God who forsakes us (Mark 15:34). The God who lets us live in the world without the working hypothesis of God is the God before whom we stand continually. Before God and with God we live without God. God lets himself be pushed out of the world on to the cross. He is weak and powerless in the world, and that is precisely the way, the only way, in which he is with us and helps us. (Bonhoeffer, Letter from prison, 16 July 1944)

I think all of us may experience such presence of God in moments of counseling or joint prayer - when we help someone not through our power and our money but when we help someone simply by being with another person - equally powerless to change a situation. I often experience such helpless togetherness when in prison - I don’t have a key to take them out - and I cannot change their sentences - but it is in insisting in fellowship where there is only hopelessness that we jointly experience the presence of Christ.

·         In moments when I felt overwhelmed and unable to help

·         In moments when I only share with someone else who is suffering

·         Possibly even joining in tears

It is in such moments that we give up our fantasies of power and experience the closeness of God’s mystery. This is also how it is with all of us: When we help someone through our money, through our expertise, through our contacts - of course this will make us feel good. But this has nothing to do with Christian faith. Christian faith - and Christian ministry - insists  where there is no benefit for yourself.

5.      Conclusion

Today’s text is a constant reminder that our faith leads us in a journey that follows the journey of Jesus Christ. It is a journey that leads through the cross - and only through the cross - to new life. It does not lead us into victory but first into suffering and weakness - but in suffering and weakness we discover depth, significance, and new life - that is the resurrection.

It is for this reason that Christians are drawn to those suffering - not just to help, but to be with, to suffer jointly. It is when we give up on our fantasies of power that we experience new depth. It may be in our visiting ministry or in our fellowship or in our worships - it is at the rare points where we are standing naked before God - powerless, you and me - that we experience the transcendence of God - another kind of victory.



\newpage

\section{富裕的貧窮}
\label{sec:384}
\textbf{鄧瑞強博士 - 富裕的貧窮}
\newline
\newline
\hyperref[sec:383]{< < < PREV SERMON < < <}
~
\hyperref[sec:toc]{[返目錄]}
~
\hyperref[sec:385]{> > > NEXT SERMON > > >}


啟示錄 3:14-22
\newline
\begin{longtable}{cl}
\hline
\hline
章節 & 經文 (和合本修訂版)\\
\hline
3:14 & \begin{tabularx}{0.7\textwidth}{X} 「你要寫信給老底嘉教會的使者，說：『那位阿們、誠信真實的見證者、神創造的根源這樣說： \end{tabularx} \\ \\ \relax
3:15 & \begin{tabularx}{0.7\textwidth}{X} 我知道你的行為，你也不冷也不熱；我巴不得你或冷或熱。 \end{tabularx} \\ \\ \relax
3:16 & \begin{tabularx}{0.7\textwidth}{X} 既然你如溫水，也不冷也不熱，我要從我口中把你吐出去。 \end{tabularx} \\ \\ \relax
3:17 & \begin{tabularx}{0.7\textwidth}{X} 你說：我是富足的，已經發了財，一樣都不缺，卻不知道你是困苦、可憐、貧窮、瞎眼、赤身的。 \end{tabularx} \\ \\ \relax
3:18 & \begin{tabularx}{0.7\textwidth}{X} 我勸你向我買從火中鍛鍊出來的金子，使你富足；又買白衣穿上，使你赤身的羞恥不露出來；又買眼藥抹你的眼睛，使你能看見。 \end{tabularx} \\ \\ \relax
3:19 & \begin{tabularx}{0.7\textwidth}{X} 凡我所疼愛的，我就責備管教。所以，你要發熱心，也要悔改。 \end{tabularx} \\ \\ \relax
3:20 & \begin{tabularx}{0.7\textwidth}{X} 看哪，我站在門外叩門，若有聽見我聲音而開門的，我要進到他那裡去，我與他，他與我一起吃飯。 \end{tabularx} \\ \\ \relax
3:21 & \begin{tabularx}{0.7\textwidth}{X} 得勝的，我要賜他在我寶座上與我同坐，就如我得了勝，在我父的寶座上與他同坐一般。 \end{tabularx} \\ \\ \relax
3:22 & \begin{tabularx}{0.7\textwidth}{X} 凡有耳朵的都應當聽聖靈向眾教會所說的話。』」 \end{tabularx} \\ \\
[1ex]
\hline
\hline
\end{longtable}
講題:富裕的貧窮　The Poverty Richness

經文:啟示錄3:14-22

講員:鄧瑞強博士

場合:香港中文大學崇基學院禮拜堂

日期:2017年10月8日

各位弟兄姊妹,早安.願你們平安.

剛剛過了中秋節,願各人生命圓滿.

窗,對一居所而言,很重要.沒有窗的房間,不能親身體會春夏秋冬的變化,不能感受日與夜的交替,也不能自由地呼吸新鮮的空氣.我們都喜歡窗,需要窗.四面圍牆的生活,我們會跼促不安,不開心,鬱悶.窗,是空的,其重要性正正在於它是空的.填滿了,就沒有窗了.窗,空無一物,卻讓我們接觸這個世界,體會時間的流逝,呼吸新鮮的空氣.問題是,在這個強調物質的世界裡,我們心靈的窗,常常被物質堵塞了.沒有窗,四面只是圍牆,我們活得很不開心.主耶穌來,提示我們,打開心窗.

今日講道經文:啟示錄3:14－22

啟3:14　「你要寫信給老底嘉教會的使者,說:『那為阿們的,為誠信真實見證的,在神創造萬物之上為元首的,說:

啟3:15　我知道你的行為,你也不冷也不熱；我巴不得你或冷或熱.

啟3:16　你既如溫水,也不冷也不熱,所以我必從我口中把你吐出去.

啟3:17　你說:我是富足,已經發了財,一樣都不缺；卻不知道你是那困苦,可憐,貧窮,瞎眼,赤身的.

啟3:18　我勸你向我買火煉的金子,叫你富足；又買白衣穿上,叫你赤身的羞恥不露出來；又買眼藥擦你的眼睛,使你能看見.

啟3:19　凡我所疼愛的,我就責備管教他；所以你要發熱心,也要悔改.

啟3:20　看哪,我站在門外叩門,若有聽見我聲音就開門的,我要進到他那裡去,我與他,他與我一同坐席.

啟3:21　得勝的,我要賜他在我寶座上與我同坐,就如我得了勝,在我父的寶座上與他同坐一般.

啟3:22　聖靈向眾教會所說的話,凡有耳的,就應當聽!』」

(1)主耶穌的自稱(啟3:14)

啟3:14　你要寫信給老底嘉教會的使者,說:「那為阿們的,為誠信真實見證的,在神創造萬物之上為元首的……」

主耶穌被稱為「阿們的,為誠信真實見證的,在神創造萬物之上為元首的」.

用了三重稱呼,來描繪主耶穌.

「阿們」,即確確實實的,穩穩陣陣的,可信賴的.

「為誠信真實見證的」,即主耶穌是可靠的,講真話的,見證真理的.

這樣描繪主耶穌,用香港話來說,即:主耶穌是「堅」的,不是「流」的.這是要凸顯出,老底嘉的信徒,很「流」.

「在神創造萬物之上為元首的」.「元首」,也可譯作「始源」,「基礎」.主耶穌是萬有的創造者,是萬有的根基.

倚靠萬有,而不倚靠萬有的始源,根基,創造者,是短視,是盲目,是愚蠢.偏偏,短視,盲目,愚蠢的信徒,多的是.

(2)失去生命的「功能」(啟3:15-16)

啟3:15　我知道你的行為,你也不冷也不熱；我巴不得你或冷或熱.

啟3:16　你既如溫水,也不冷也不熱,所以我必從我口中把你吐出去.(嘔吐之意)

在老底嘉附近,有溫泉,水是熱的,可治病.歌羅西城也在附近,有冷泉,水是冷的,清涼好飲.偏偏,老底嘉的水,不是溫泉,卻有溫泉的餘溫,有溫泉水的臭味,入口令人作嘔.熱水或冷水,都對人有益.但如溫水,不冷也不熱,不能飲,也不能治病,只能被人吐出來,令人作嘔.

不冷不熱,與熱心與否無關,而關乎生命的內在質素能否產生有意義的效果.若生命的質素太差,越熱心,造成的禍害越大.一隻骯髒的蒼蠅,越熱心去吻人的面龐,越引人反感.溫泉的熱水,或冷泉的涼水,或冷或熱,其存在都是有意義的.這裡講的溫水,其存在全無意義.一個人,若只是行屍走肉,對上天的仁愛毫無領會,對人間的美善全不在乎,對他人的福祉全不關心,每天想著的,只是吃喝玩樂,他的生命,就是溫水,他活過與不曾活過,毫無分別,他不會記念任何人,任何人也不會記念他.主耶穌說,他不喜歡這種生命,這種生命令他「作嘔」.

一個人全無內在質素,是令人作嘔的.舉例說,若一個律師講的話,是在藐視法律的正義,是在利用法律去打壓異己,是在顯示自己是法律的主人,他就失去律師的內在質素,他的存在對社會產生不到意義.面對這種「溫水」一樣的人,主耶穌只能作嘔.

現在,出現問題的,是教會.教會之為教會,是站在地上,卻望向上天.教會像祭司,與痛苦的人站在一起,為他們的苦祈求上天；也與基督站在一起,向痛苦的人施以援手.教會像先知,指正錯誤,言說真理.但若教會裡各人已失去尋求上帝的心志,已對實踐真正的價值毫無興趣,已不知「敬拜神」與「不敬拜神」有何分別,則這教會可能已不成教會,存在與否,毫無分別.面對這教會,主耶穌只能作嘔.

教會搞到如此,是因為價值的誤判.

(3)價值的誤判(啟3:17-18)

啟3:17　你說:我是富足,已經發了財,一樣都不缺；卻不知道你是那困苦,可憐,貧窮,瞎眼,赤身的.

啟3:18　我勸你向我買火煉的金子,叫你富足；又買白衣穿上,叫你赤身的羞恥不露出來；又買眼藥擦你的眼睛,使你能看見.

問題的核心,是老底嘉教會太富裕,而造成一種錯覺,以為自己的生命很豐盛.他們表面的富足,使他們忘記了內裡的空洞.生命的歡愉,令他們不再羨慕靈魂的救贖.

老底嘉教會說:「我是富足,已經發了財,一樣都不缺」.

「我是富足」的,開宗明義,他們感到生命充溢著各種美好的東西,物質供應是豐富的,享樂是多姿多采的,滿足是唾手可得的.

「已經發了財」,指他們手上擁有的物質和財富的數量之大,足以讓他們感到生命的絕對穩妥.

「一樣都不缺」,這是他們的心理狀態.我們熟悉的詩篇23篇說:「耶和華是我的牧者,我必不致缺乏」,我一無所缺.老底嘉教會卻說,財富在我手上,我一無所缺.

他們擁有太多物質,以致封死了心靈的窗.他們的生活太豐富了,以致他們察覺不到內裡的空洞.他們實在太富有,以致他們低估了自己的貧窮.

有一個這樣的故事:(參:\url{http://paulocoelhoblog.com/2015/09/04/the-fisherman-and-the-businessman/})

在巴西,海邊有一條村莊.有一商人在海邊坐著.他看到一漁夫划著一小艇回來,船上有不少大魚.商人問:「捕這些魚,花了多少時間?」漁夫答:「只一會兒.」商人很奇怪,問:「為何不花多一點時間,多捕一些魚?」漁夫答:「夠養活家人了.」商人又問:「那你整天餘下的時間,做什麼?」漁夫答:「與兒女玩耍,和太太一起睡睡午覺,黃昏時和朋友飲杯酒,彈吉他,唱歌,跳舞.」

商人說:「我可助你成為一成功人士.聽我說,由現在開始,你盡量花時間在海中,捕盡量多的魚.賺多一些錢,之後,買更大的船,捕更多的魚.再之後,成立公司,生產罐頭.再之後,成立更大的公司,管理生產與銷售系統.再之後,上巿.那時,你就會成為有錢人.」

漁夫問:「那之後,怎麼樣?」

商人答:「那之後,你可以在海邊買間屋,閒來釣釣魚,花時間和小孩子玩,和太太睡睡午覺,黃昏時和朋友飲杯酒,彈吉他,唱歌,跳舞.」

漁夫問:「這不就是我現在的生活嗎?」

富有的人有一種追求富有的思考方式,有時,這思考方式讓人離開了真正的生活.太多的物質,會堵塞心靈的窗,讓人在四面圍牆裡,看不到真正的世界,也呼吸不到新鮮的空氣.

面對這些自覺豐足而一無所缺的人,主耶穌說:「你們實在是困苦,可憐,貧窮,瞎眼,赤身的.」明顯,老底嘉教會,對自己的光景,嚴重誤判.

老底嘉城,有三大置富元素.它是當時的銀行中心,類似現在的瑞士,富人將錢藏在這裡.它是眼藥的生產中心.它的黑羊毛,也極有名,這些黑羊毛有藍紫色的光澤.

老底嘉城,有錢,有眼藥,有名貴羊毛,但主耶穌卻說,你們是「貧窮,瞎眼,赤身的」,你們實在是「困苦,可憐」的.他們擁抱的價值,在主眼中,並不重要.

主耶穌說:「我勸你向我買火煉的金子,叫你富足；又買白衣穿上,叫你赤身的羞恥不露出來；又買眼藥擦你的眼睛,使你能看見.」老底嘉教會的三大強項,正正就是他們的三大盲點.他們將生命的豐盛和穩妥,押注在這三大強項上,就在這裡,他們注目世上的事物,而遠離了創造萬物的神.

香港社會是非常富有的,財富本身,不是問題所在,問題在於我們是否只是注目財富,而忽略內在生命的靈性與德行.孔子說:「不義而富且貴,於我如浮雲.」(《論語》〈述而〉)取之無道,用之不義,則生命仍是空洞無物.孟子說:「人之有道也,飽食,煖衣,逸居而無教,則近於禽獸.」(《孟子》〈滕文公上〉)人的生命,應有高貴的內涵,飽食,煖衣,逸居,這都是好的,但卻「無教」,沒有靈魂,「則近於禽獸」.

(4)主的叮嚀(啟3:19-22)

啟3:19　凡我所疼愛的,我就責備管教他；所以你要發熱心,也要悔改.

啟3:20　看哪,我站在門外叩門,若有聽見我聲音就開門的,我要進到他那裡去,我與他,他與我一同坐席.

啟3:21　得勝的,我要賜他在我寶座上與我同坐,就如我得了勝,在我父的寶座上與他同坐一般.

主耶穌愛他的教會,他嚴厲的責備,是出於愛,出於對我們的期望.

他站的位置,是門外,他等待信徒接待他.主耶穌應為此哀嘆,他只是站在門外,而未能進到信徒的生命中.他的信徒只關心世上的豐足,忘記了他.

他期望我們,認清永恆價值之所在,以致終能分享神永恆的榮美.

(5)結語

老底嘉教會,其實是很多教會及信徒追求的目標,富裕,自信,滿足.從現實而言,他們真的「一樣都不缺」.他們缺乏的,是生命內涵,是價值質素,是靈性品格.看來,他們過著純物質的生活,卻對主耶穌毫無內在關係.這種生命,徒有信仰的外表,卻令主耶穌作嘔.他們活著,卻不會在世上留下任何意義.他們應是世上的光,但卻是晦暗不明.他們誤將物質的豐盛,看作是生命的豐盛.他們對什麼是真正價值,完全誤判.他們自以為富足,但主耶穌說,你們貧窮不已.

每個人都追求富足與豐盛,但願我們追求的,是有永恆意義的富足與豐盛.

願:榮耀歸於聖父,聖子,聖靈.阿們.



\newpage

\section{我施平安又降災禍}
\label{sec:385}
\textbf{關瑞文教授 - 我施平安,又降災禍}
\newline
\newline
\hyperref[sec:384]{< < < PREV SERMON < < <}
~
\hyperref[sec:toc]{[返目錄]}
~
\hyperref[sec:386]{> > > NEXT SERMON > > >}


以賽亞書 45:1-7
\newline
\begin{longtable}{cl}
\hline
\hline
章節 & 經文 (和合本修訂版)\\
\hline
45:1 & \begin{tabularx}{0.7\textwidth}{X} 耶和華對所膏的居魯士如此說，他的右手我曾攙扶，使列國降服在他面前，列王的腰帶我曾鬆開，使城門在他面前敞開，不得關閉： \end{tabularx} \\ \\ \relax
45:2 & \begin{tabularx}{0.7\textwidth}{X} 「我要在你前面行，修平崎嶇之地。我必打破銅門，砍斷鐵閂。 \end{tabularx} \\ \\ \relax
45:3 & \begin{tabularx}{0.7\textwidth}{X} 我要將暗中的寶物和隱藏的財富賜給你，使你知道提名召你的就是我－耶和華，以色列的神。 \end{tabularx} \\ \\ \relax
45:4 & \begin{tabularx}{0.7\textwidth}{X} 因我的僕人雅各，我所揀選的以色列，我提名召你；你雖不認識我，我也加給你名號。 \end{tabularx} \\ \\ \relax
45:5 & \begin{tabularx}{0.7\textwidth}{X} 我是耶和華，再沒有別的了；除了我以外再沒有神。你雖不認識我，我必給你束腰。 \end{tabularx} \\ \\ \relax
45:6 & \begin{tabularx}{0.7\textwidth}{X} 從日出之地到日落之處使人都知道除我以外，沒有別的。我是耶和華，再沒有別的了。 \end{tabularx} \\ \\ \relax
45:7 & \begin{tabularx}{0.7\textwidth}{X} 我造光，又造暗；施平安，又降災禍；做成這一切的是我－耶和華。 \end{tabularx} \\ \\
[1ex]
\hline
\hline
\end{longtable}
講題:我施平安,又降災禍　I Make Weal and Create Woe

講員:關瑞文教授

經文:以賽亞書 45章1-7節

場合:香港中文大學崇基學院禮拜堂

日期:2017年10月22日

(1)我-耶和華所膏的塞魯士；我攙扶他的右手,使列國降伏在他面前.我也要放鬆列王的腰帶,使城門在他面前敞開,不得關閉.我對他如此說:(2)我必在你前面行,修平崎嶇之地.我必打破銅門,砍斷鐵閂.(3)我要將暗中的寶物和隱密的財寶賜給你,使你知道提名召你的,就是我-耶和華,以色列的神.(4)因我僕人雅各,我所揀選以色列的緣故,我就提名召你；你雖不認識我,我也加給你名號.(5)我是耶和華,在我以外並沒有別神；除了我以外再沒有神.你雖不認識我,我必給你束腰.(6)從日出之地到日落之處使人都知道除了我以外,沒有別神.我是耶和華；在我以外並沒有別神.(7)我造光,又造暗；我施平安,又降災禍；造作這一切的是我-耶和華.

與大家分享一首我自少就非常喜歡的詩歌.每年,我在某課堂上,都與學生分享這首詩歌.有深度的詩歌往往能使我們更多明白聖經的真理.因為,在每首詩歌背後的屬靈經驗,常常讓聖經真理的現代意義,表達得淋漓盡致.讓我們看看蘊含在一首詩歌裡的屬靈經驗,可如何啟發我們的人生.詩歌的名稱是「我心靈得安寧」:

(第一節歌詞)

有時享平安,如江河平又穩,有時遇悲傷似浪滾；

不論何環境,我已蒙主引領,我心靈得安寧,得安寧.

記得小時侯,在少年團契學會這首詩歌,一唱鐘情.當時,覺得詩歌的旋律非常平靜,還以為詩歌背後是一段非常甜蜜的故事.不久,偶然經過一間書室,隨便拿來一本介紹詩歌的書,才知道這首詩背後,是一個非常傷感卻又動人的故事.

原來,詩人名叫Spafford,生於19世紀,是個正直而勤奮的商人,與妻育有五個兒女,家庭挺幸福,生意也興隆,一家人都是基督徒.1873年,他們計劃了一次郵輪旅遊,一家人要享受輕鬆愉快的假期.然而,在登船之前不久,因為Mr. Spafford 生意事務纏身,最終沒有與妻兒同行.他只好承諾遲些與他們在下船的地方匯合.誰不知,天意弄人,原來登船那天,是他們一家人最後團聚的日子.怎會呢?郵輪出發後四天,在大西洋某處遇上了致命的海難.幸好,Mrs. Spafford 獲救且被送往醫院.她醒來時,傷痛欲絕,因為所有兒女已命喪黃泉.不久,仍在家裡的 Mr.Spafford,收到 Mrs.Spafford遠方傳來一紙只有兩隻字的電報: “SAVED, ALONE” (獲救了,但只剩我一個人).晴天霹靂,禍福無常!Mr. Spafford 隨即訂了船票,要趕赴妻子身旁.Mr. Spafford悲痛不已,在船上回憶以往一家人的生活片段,他只能哭著問上帝,為何會如此天意弄人.可是,問了一次,又再問一次,上帝都依然默不作聲,沒有給詩人什麼回應.四天后,船開到海難之處,當天卻是風平浪靜.Spafford 一個人站在微風與海浪聲中,望著吞噬他兒女的浩瀚大海,整個人被日落的天空包圍著,海面金光閃閃.很奇妙,他心裡從未如此平靜過.他頓悟了-無論如何,上帝是掌管歷史,掌管人生的.擦乾眼淚,他趕快回了房間,寫出了這首平靜而有力的詩歌.

有時享平安,如江河平又穩,有時遇悲傷似浪滾；

不論何環境,我已蒙主引領,我心靈得安寧,得安寧.

這是怎麼樣的屬靈經驗呢?當好人遇上壞事,是怎麼一回事呢?災難是誰作弄的呢?讓我們翻開賽四十五1－7.

原來,賽四十四24－28及賽四十五1－8有密切關係.這兩段經文是一首詩的兩個部分.頭一部分,即四十四章24－28,是個引子,用來介紹在四十五章1節出現的塞魯士.塞魯士是誰呢?等下我們會談到,現在讓我們看看這首詩的寫作背景.

這段經文是何時寫成的呢?一般學者相信,它是寫成於西元前五百多年之間.對於以色列南國猶大人來說,這段時間是他們民族命運的關鍵時刻.原來以色列國的國運一路都不大亨通,以色列北面的土地早在西元前七百多年已被亞述帝國吞併,其中的以色列人,一早已經因與異族通婚而被同化.因此,所謂以色列,一半人一早已經在歷史中消失,剩下來弱小南面的猶大國,也一早成為亞述的傀儡國.但因為猶大國有一定的利用價值,於是,就好像香港起初回歸中國時一樣,亞述一直都容許猶大國享有高度自治.但好景不常,西元前六百幾年的時候,巴勒斯坦鄰近興起了梟勇善戰但出名殘酷的巴比倫帝國.在西元前612年,亞述帝國陷落,作為亞述的傀儡國的猶大國,也在西元前586年淪陷,落入巴比倫之手.巴比倫一如以往,兇殘成聖,將整個猶大地在一夜間變成荒涼.大家可以閉上眼睛,試想一想當時慘痛的情況:很多人在一夜間痛失家園,弧兒寡婦,哭啼處處,隨時在街上遇見的面容潰爛的屍體,可能有一具是昨天還活生生與你一起吃晚飯的摯愛良朋或者家人.辛辛苦苦經營幾百年才得的經濟成就,完全毀於一旦,各人人心惶惶,本來繁榮安定的生活,在這次政權易手當中,完全成為往事只能如煙.這次政權易手所帶來的創傷,不單是在肉體上,也不單是在心靈上,也在靈性上.為何?因為一如以往的作風,巴比倫攻陷耶路撒冷後,就將聖殿移為平地.眼巴巴看著一向象徵著「平安-shalom」的宏偉聖殿被移為平地,猶大人除了感到無能為力之外,更重要的是他們心底裡發問:「我們一向所信奉所敬拜的耶和華,竟然連自己的居所也保不住了!作為他的子民,剩下來的只有絕望.」他們這樣想,對他們來說,是很有道理的.原來當時一般人有兩個信念:第一,耶和華是實實在在居住於聖殿之中,因此,聖殿就是耶和華所在之處；第二,當時的人認為,每個民族都是由他們所信奉的神明來保護的.兩個國家在地上打仗,背後就是兩個民族神在天上戰爭,地上哪個國家戰敗,就是天上哪個民族神戰敗.所以,很自然,眼見聖殿被毀,就只有一個結論,就是萬軍的耶和華原來不及別人的神.事實上,根據歷史記載,當時就有很多猶大人改而信奉了巴比倫的神.不肯「轉軑」的,就始終不明白,為何猶大國會落得如此收場.

巴比倫亦一如既往,攻陷了猶大國之後,就將城中精英分子全部擄走,帶回巴比倫國參與建設.這群精英,當中不乏不肯「轉軑」的耶和華敬拜者,以賽亞書四十五章的作者,應該就是其中一人.他們離鄉別井,天天思念故鄉,思念從前的輝煌,雖然生活尚可,但眼見暗黑的歷史巨輪無情的向前滾,晦暗的時光來勢洶洶,而人力又始終不能逆轉時勢,心裡忍不住晝夜不停地發問,為何萬軍的耶和華竟然是戰敗者.他們的痛苦,心裡的鬱結,信仰的動搖,看來旁人無法明白.

就是在這樣痛苦的煎熬下,以賽亞與一同被擄的人一日一日的捱下去,盼望耶和華終有一日如以往般施行拯救,在以色列人民中興起戰士,舉兵復國.可是,日復日,年復年,始終見不到耶和華的拯救,一些人在絕望中老死,一些人則堅持不住,放棄信仰,改投信奉巴比倫神,背祖忘宗.唯獨是以賽亞一夥人,仍然堅守信念,毋忘初衷,深信領他們出埃的神,必然再次施行拯救,重新創造以色列民.

一年又一年,被擄三十幾年後,終於出現了一線曙光.三十幾年不是一段短的時間,你可以想一想,如果你現在是三十來歲,你應該是年青力壯的.三十幾年後,你已是一個六,七十歲兩鬢斑白的人了.被擄三十幾年後,東方興起了一個名叫波斯的國家,當時在位的君王是塞魯士.塞魯士一向享有仁君的聲譽,以往被其攻佔的國家,都享有宗教自由.當塞魯士開始要攻打巴比倫時,揚言若果戰勝,就讓被擄的猶大人歸回原地,重建家園,重建聖殿.以賽亞書四十四章24至四十五章8就是在見到這個曙光時寫成的.

從四十五章1－8節看,經過三十多年的慘痛經驗,猶大人對上帝產生了至少兩個新的認識:

(1)上帝的作為奇妙難測,不隨人意-讓我們看看第一節:「我-耶和華所膏的塞魯士；我攙扶他的右手,使列國降伏在他面前.我也要放鬆列王的腰帶,使城門在他面前敞開,不得關閉.」當中的「所膏的」,在希伯來原文有彌賽亞之意,而彌賽亞即是上帝的受膏者,又即是君王.這是一個很突破的新瞭解,一路以來忠於耶和華的猶大人都以為,耶和華必定在以色列民中興起君王拯救他們的,而他們也天天為此祈求.誰不知耶和華並不一定要按人的本子辦事,神立意在異教世界中興起以色列民族的救星.明白了上帝的安排,以賽亞也欣然接受,並且執筆寫了這首詩,破天荒地在詩的前半部份,即四十四章28歌頌異族波斯王塞魯士.以賽亞的經歷,叫他明白了一點,就是上帝之所以是真上帝,在於他並不按人的常規辦事.他所行之事,是人不可測透,在他面前,我們只好欣然接受由他所導演的戲.上帝之所以為真神,在於他並不聽你吩咐,按你的本子辦事.

(2)第二個新發現,在第七節中反映出來:「我造光,又造暗；我施平安,又降災禍；造作這一切的是我-耶和華.」一般人認為,光明,幸福,平安是來自神,而黑暗,災禍,痛苦是來自撒旦的.但是經歷過喪家亡國,現在又看見復國曙光,以賽亞一夥人終於明白到,獨一真神是造光,又造暗,施平安又降災禍的上帝.明白到這包括順人意的幸福與叫人悲痛的絕望一切都是出於耶和華,他們才發現,被極權殘忍的權力者統治也是出於耶和華的,被拯救於水深火熱的當然也都是耶和華.「我造光,又造暗；我施平安,又降災禍；造作這一切的是我-耶和華.」經過三十年的苦痛經歷,以賽亞明白到,烏雲本身絕不是荒謬的,在上帝的掌管下,連黑暗本身都是有意義的.

從上分析,我們瞭解到,以色列人的亡國經驗,叫他們明白了兩件重要事情:(1)上帝之所以為上帝,在於他並不按人的本子辦事.(2)在上帝的掌管下,連黑暗本身都是有意思的.

弟兄姊妹,上帝的教訓是歷久彌新的,這兩點,對於我們今天又有甚麼意義呢?

(1)    上帝之所以為上帝,在於他並不按人的本子辦事-你祈禱求神幫助,是求他幫助,還是教他如何幫助你?有很多禱告沒有應驗,是因為上帝沒有幫助你,還是因為上帝沒有按照你的意思幫助你,以致他幫了你時你也不察覺.上帝是不是一定要善有善報,惡有惡報呢?為什麼我們總要認為上帝一定要照人的常規辦事.是不是我們一生愛主,就不會遇到太多挫折或疑難雜症?為什麼我們認為上帝一定要遵守我們想當然的遊戲規則?是不是當人世間公義不彰,大道不行,我們就要絕望地呼喊「上帝啊 上帝啊 你為甚麼離棄我」?還是我們相信「我造光,又造暗；我施平安,又降災禍；造作這一切的是我-耶和華」,以致我們在落難而深感無能為力時,依然能夠有氣有力地各守本分,等待第八節所言的發生:「諸天哪,自上而滴,穹蒼降下公義；地面開裂,產出救恩,使公義一同發生；這都是我-耶和華所造的.」弟兄姊妹,上帝之所以是真神,正正在於他並不接受你的吩咐行事.

(2)在上帝的掌管下,連黑暗本身都是有意思的-你是否能夠相信暗黑時光本身是有意義的呢?(失敗,失戀,無力…)人生不如意事十常八九,若每當陷於傷心失意時,我們都只懂怨天尤人,自怨自憐,慨嘆撒旦攻擊,你是否已經忘記,在上帝的掌管下,烏雲本身是有意義的,並且,太陽本來就是長存於烏雲之上呢?每個危機都可以是轉機,只要我們認定烏雲本身也是上帝的使者就是了.

結論

上帝並不按人的本子辦事,與及烏雲本身在上帝裡是有意義的這兩個道理,大抵就是「我心寧得安寧」這首詩的作者最終所體會到的.無論我們在何景況,願意我們能分享到詩人及以賽亞先知的信仰體會,以致我們有足夠的勇氣面對人生幽暗的一面,阿們.



\newpage

\section{不斷改革的教會}
\label{sec:386}
\textbf{賴品超教授 - 不斷改革的教會}
\newline
\newline
\hyperref[sec:385]{< < < PREV SERMON < < <}
~
\hyperref[sec:toc]{[返目錄]}
~
\hyperref[sec:387]{> > > NEXT SERMON > > >}


啟示錄 3:1-6
\newline
\begin{longtable}{cl}
\hline
\hline
章節 & 經文 (和合本修訂版)\\
\hline
3:1 & \begin{tabularx}{0.7\textwidth}{X} 「你要寫信給撒狄教會的使者，說：『那有神的七靈和七顆星的這樣說：我知道你的行為，就是名義上你是活的，實際上你是死的。 \end{tabularx} \\ \\ \relax
3:2 & \begin{tabularx}{0.7\textwidth}{X} 你要警醒，堅固那些剩下、快要死的，因為我發現你的行為，在我神面前沒有一樣是完全的。 \end{tabularx} \\ \\ \relax
3:3 & \begin{tabularx}{0.7\textwidth}{X} 所以，要記得你所領受和聽見的；要遵守，並要悔改。你若不警醒，我必如賊一樣來到；我幾時來到你那裡，你絕不會知道。 \end{tabularx} \\ \\ \relax
3:4 & \begin{tabularx}{0.7\textwidth}{X} 然而，在撒狄你還有幾位是未曾污穢自己衣服的，他們會穿白衣與我同行，因為他們是配穿的。 \end{tabularx} \\ \\ \relax
3:5 & \begin{tabularx}{0.7\textwidth}{X} 得勝的必這樣穿白衣，我也不從生命冊上塗去他的名；我要在我父面前，和我父的眾使者面前，宣認他的名。 \end{tabularx} \\ \\ \relax
3:6 & \begin{tabularx}{0.7\textwidth}{X} 凡有耳朵的都應當聽聖靈向眾教會所說的話。』」 \end{tabularx} \\ \\
[1ex]
\hline
\hline
\end{longtable}
講題:不斷改革的教會　Continuing Reformation of the Church

經文:啟示錄3章1-6節

講員:賴品超教授

場合:香港中文大學崇基學院禮拜堂

日期:2017年10月29日

後天(2017年10月31日)是記念宗教改革500週年的大日子.在過去的一年,香港教會及神學界,有很多不同的活動去慶祝,記念和反思.本人也出席了一些相關學術會議,當中經常探討的問題是,500年前的宗教改革,對當代華人教會有何意義.今天想藉啟示錄3章 1-6節的一段經文,分享一下個人的一些想法.

翻查網上資料,經文提及的撒狄(Sardis),即今天的土耳其的Sart市.撒狄是歷史名城,是古代的呂底亞(Lydia)國的首都,由於地理上有天然屏障的保護,十分安全,富庶；雖然曾先後被波斯大軍及亞歷山大帝攻陷,一直仍頗為繁榮,唯在公元17年的地震後,元氣大傷,無法恢復昔日的興盛.可以說,在啟示錄成書的時代,撒狄作為一個城市,雖然仍是赫赫有名,但恐怕是名大於實.也許,撒狄的教會也有類似的情況.

正如啟示錄的其他六封書信一樣,給撒狄教會的信也是先介紹發信者,就是那有上帝的七靈和七星的(3:1),而這是一種修辭,指的是那掌控物質與心靈的大能者,而他是完全知道教會的情況,而信內對撒狄教會的判斷主要有三點.

第一,「我知道你的行為,按名你是活,其實是死的」(3:1).這裡所說的「名」,並非指 「撒狄」這個名字本身的意義(有人猜測「撒狄」意指「剩餘的,喜樂王子,逃出者」等),而是指名聲,名氣(reputation).正如新譯本也是譯作「名聲」.換言之,一般人聽到是撒狄的教會,也許會假返設這是一間有活力的教會,然而在上帝眼中,其實是死的,是沒有生命的,名存實亡.

第二,「因我見你的行為,在我神面前,沒有一樣是完全的」(3:2).與第一個判語相似的是,這也是一句當頭棒喝,都是指出問題的嚴重性.相對來說,第一句判語重點在指出,上帝的審判不是憑表面上有某種名聲,而是看內在的生命；而第二句則指出,施行審判的是上帝,不是按人的相對標準分高下,而是分辨是否達到上帝所要求的完全；而在上帝面前,撒狄教會正是沒有一樣是完全的.有時信徒很容易自覺很不錯,例如常到教會,奉獻,食飯祈禱,每日讀經,常傳福音,在別人眼中是虔誠信徒,在教牧眼中的絕世好教友,那又是否足夠?是否也行公義,好憐憫,存謙卑的心與上帝同行?作出這些行為,是否出於愛?在上帝面前是否完全?

第三,「然而在撒狄你還有幾名是未曾污穢自己的,他們要穿白衣與我同行,因為他們是配得過的.」(3: 4).相對於前二者屬整體概括性的負面判斷,第三句是比較針對個別的少數的情況.這不僅反映審判者的全知,也反映了審判的公允,而非以偏蓋全.當中的審判,不是先定下名額(quota)或百分比(percentage),然後在每間教會中找相對較好的,給予代表聖潔的白衣；而是看他們是否「配得」.換言之,衪的審判不是用norm referencing,不是透過教會與教會之間的比較,或信徒與信徒之間,找出相對較好的,那些在平均水準以上,那些在以下；而是採用criteria referencing,就是一種絕對的標準來衡量信徒是否完全,而按此標準來衡量.

由這些判語,或許我們會得出一個想像,就是雖然仍有少數很好的信徒,但整體來說,這是一間很不濟,死氣沉沉的教會,甚至快要消失.然而,我想也可能有另一種情況,就是它看起來仍是很活躍,甚至有很好的名聲,只是實際上與死無異.有一本書,名叫Excellence without a Soul: How a Great University Forget Education ,書名直譯為:「沒有靈魂的卓越:一所偉大的大學如何忘記教育」.作者Harry R. Lewis是哈佛大學的教授,曾經擔任Harvard College(即本科生所屬的學院)的院長,而書名所講的「一所偉大的大學」所指的正是哈佛大學.哈佛大學是舉世知名的大學,在世界大學排名榜經常位列前茅,但在作者看來,這所公認為卓越或偉大的大學卻是沒有靈魂的,因為它已經忘記了教育的使命.哈佛的情況雖有一定的複雜性及特別的原因,但對於在大學體制內工作多年的人來說,相類似的現象卻絕不陌生,並且也不難理解.在不同的世界大學排名中,往往只注意可以量化的指標,例如師生比例,出版的數目及被引用的次數,國際化的程度(即多少師生並非本國人),甚至畢業生的收入等；至於教學的質素,由於難以量化和比較,根本不會計算.當世界眾多大學不斷追求提升排名,不少大學教授在自願,被迫或半推半就的情況下,往往以研究及出版(包括申請研究經費)為首要,甚至是唯一重要的任務時,教育,尤其本科生的教育,往往備受忽視.當一所大學失落自己在教育上的異象與使命(vision and mission)時,只顧同行之間的互相比拼,很容易會隨波逐流,雖然在某些方面十分活躍(包括:研究,出版,籌款),並在排名榜上力爭上游,甚至獲得同行的讚許,有好的名聲,但其實是沒有靈魂的,就有如電影中的喪屍一樣.

如果大學可以如此失去靈魂,教會也可以.當教會領袖只著眼於可見的增長(例如教堂面積,聚會人數,奉獻金額,物業數量,分堂數目等),追求多,快,好,省的教會增長,這很容易獲得好的名聲,甚至成為羨慕,仿效的對象.問題是,會友靈性如何?生活中的見證又如何?是否徹底按主的道來生活?我們需要謹記,教會要面對的審判,並非來自「顧客」(教友)或「同行」(與其他教會比較),而是上帝；上帝的審判不會「拉curve」,而所要求的,不是平均水準以上(above average),也不是相對於別人的卓越,而是絕對的完全(perfect).

由這種視角看,可以肯定絕大多數的教會都不是完全的.那怎麼辦?做些甚麼?對此,經文用了五個字詞(中文是十字真言): 儆醒,堅固,回想,遵守,悔改.

1．儆醒:這字詞出現二次,正面講「你要儆醒」(3:2),反面說「若不儆醒,我必臨到你那裡如同賊一樣,我幾時臨到,你也決不能知道.」(3:3) 儆醒就是保持覺醒(state alert),尤其是意識到審判即將在無法預測的時候出現.如果不想用「賊」來比喻主耶穌的來臨,或許可以想像為突擊測驗；因此,我們不能先怠慢或說“hea”一段時間,然後在預計的時間才開始準備；而是要馬上及早作好準備,並且是恆常地作好準備,隨時可以進行測驗／考試.用武俠小說的講法,最安全的地方也是最危險的地方,因為以為安全,很容易失去警覺性,反而更危險,最大危險在不知危險.

2．堅固:「堅固那剩下將要衰微的」(3:2).當醒覺過來後,很容易覺今是而昨非,想馬上徹底改變現況,並展開鴻圖大計.但聖經在這裡的教導是十分務實,就是先「堅固那剩下將要衰微的」；正如對於有生命危險的病人,最先可能是急救止血,死不去之後,才談如何健身甚或跑馬拉松.有些很好的靈性上的習慣或傳統,一旦失去／死掉,很難重新恢復,最好還是在死去之前盡力保存.

3．回想:「回想你是怎樣領受,怎樣聽見的」(3:3) .當大多數人以某種“普世”的標準來衡量和排名,為免隨波逐流,以別人的眼光來代替上帝的審判,甚至成為喪屍,「勿忘初心」非常重要.不時回想最初所領受的異象與使命,不一定是一成不變照搬數百,甚至過千年所定立的一切,但若並以此對照今天的所行的路,最少有助批判地反思並重新導向.

4．遵守:在回想之後,「又要遵守」(3:3).在回想原領受的福音,誡命,異象與使命等,當然是予以遵守,否則就如照完鏡之後,不予梳洗整理,便如沒照一樣.然而,遵守是完全的遵守,並且持之以恆.過去一年,改了生活習慣,晚上跑步,減了大約二十磅,遇見醫學院一位教授,他的忠告是,最大的挑戰不是短時間內透過大量運動來「操 fit」並減肥,這其實不難,最難是之後能否長期維持下去.我們遵守主道,不應是三分鐘熱度,而是需要長期並全面的持守,直到見主的日子.

5．悔改:「又要遵守,並要悔改」(3:3).次序上這有點怪.如果是嚴格遵守,就不用悔改.有些佛經是以「皆大歡喜,信受奉行」作為結結束,就有如一些童話故事的快樂結尾(happy ending),自始公主和王子很快樂的生活下去.為何聖經在這是講完遵守再加「並要悔改」?最近與一些久未見的別的學院同事「飯聚」(不是「犯罪」),很健康地談起做運動.一位商學院的教授說他曾找不同的專業教練,一對一的學習不同的球類技巧,雖然收費貴,但物有所值,因為可以很快學到很多很專業的技巧.一位理工科的教授則表示,比較喜歡用自己的方法慢慢探索如何改善技巧,這才是樂趣之所在.不知兩者的不同取向是否反映學術背景上的差異,但商學院的教授則補充了很重要的一點,就是教練最重要的功能,不是教正確的姿態及技巧,而是從訪旁糾正,因為我們學了新的,正確的方法後,很容易在不自覺的情況下,回到舊的,錯的,但早已習慣的方法,因此需要教練不斷提點,直至我們能習慣成自然地使用正確的方法.同樣地,在重新開始遵守主之後,我們的實踐也會不自覺地,慢慢地出現走樣,變形,打折扣,因此是需要經常的悔改.

後天是記念宗教改革500週年的大日子,因為500年前的10月31日馬丁路德將他的95條論綱釘在威登堡大教堂的門上,提出辯論,繼而引起一連串的辯論和改革,不單在神學上有很多的辯論,包括因信稱義,聖經,聖餐等,在教會的體制以至政教關係也出現了翻天覆地的改變.當然,實際情況十分複雜,但教會一般是以此,就是1517年10月31日的事件,作為十六世紀一場席捲整個歐洲,而影響力甚至蔓延全球的宗教改革運動的序幕!然而,馬丁路德在95條論綱主要討論的,根本沒有觸及因信稱義,聖經權威等問題,他所特別關注的是悔改的問題,就是信徒用金錢買赦罪券,對真實的悔改有何意義,還是無助於達至真實的悔改.

經過五百年,不少基督新教的教會,仍十分珍惜宗教改革的傳統,尤其所帶來在教義及教制上的更新,甚至為宗教改革所帶來的在文化與歷史上的轉化而自豪.然而,我們不可忘記,十六世紀的宗教改革運動,也有它的種種局限,不足,甚至是陰暗面,包括引起教會的分裂,對異教人士的敵視等.因此,我在近期的一個學術會議上,刻意的從生態神學的視角,批判地反思十六世紀的宗教改革的傳統,指出當中最有價值的遺產,也許不是唯獨信心,唯獨聖經等教義,而是那種不斷悔改,不斷改革的精神.

也許基督新教的教會,會認為我們是已經悔改了,已經成功改革了,已經撥亂反正的教會,由於十六世紀的宗教改革,我們有了純正的福音,正統信仰等.然而,留意啟示錄的這段經文,當中強調的是「行為」,「我知道你的行為」,「因我見你的行為」.不錯,按著名字或名聲,我們的教會算是已經改革過的,但我們是否也活出不斷悔改,不斷改革的生命?

願那創立並又不斷引導教會的三一上帝,繼續啟迪和引領教會,讓我們可以不斷悔改,不斷改革,以至能更像上帝那樣完全,榮歸上主.



\newpage

\section{從衝突到共融}
\label{sec:387}
\textbf{蒲錦昌牧師 - 從衝突到共融}
\newline
\newline
\hyperref[sec:386]{< < < PREV SERMON < < <}
~
\hyperref[sec:toc]{[返目錄]}
~
\hyperref[sec:388]{> > > NEXT SERMON > > >}


哥林多前書 10:17
\newline
\begin{longtable}{cl}
\hline
\hline
章節 & 經文 (和合本修訂版)\\
\hline
10:17 & \begin{tabularx}{0.7\textwidth}{X} 因為餅只是一個，我們雖然人多，仍是一體，我們同享一個餅。 \end{tabularx} \\ \\
[1ex]
\hline
\hline
\end{longtable}
以弗所書 4:1-6
\newline
\begin{longtable}{cl}
\hline
\hline
章節 & 經文 (和合本修訂版)\\
\hline
4:1 & \begin{tabularx}{0.7\textwidth}{X} 我為主作囚徒的勸你們，既然蒙召，行事為人就要與你們所蒙的呼召相稱。 \end{tabularx} \\ \\ \relax
4:2 & \begin{tabularx}{0.7\textwidth}{X} 凡事要謙虛、溫柔、忍耐，用愛心互相寬容， \end{tabularx} \\ \\ \relax
4:3 & \begin{tabularx}{0.7\textwidth}{X} 以和平彼此聯繫，竭力保持聖靈所賜的合一。 \end{tabularx} \\ \\ \relax
4:4 & \begin{tabularx}{0.7\textwidth}{X} 身體只有一個，聖靈只有一位，正如你們蒙召，是為同有一個指望而蒙召， \end{tabularx} \\ \\ \relax
4:5 & \begin{tabularx}{0.7\textwidth}{X} 一主，一信，一洗， \end{tabularx} \\ \\ \relax
4:6 & \begin{tabularx}{0.7\textwidth}{X} 一神－就是萬人之父，超越萬有之上，貫通萬有，在萬有之中。 \end{tabularx} \\ \\
[1ex]
\hline
\hline
\end{longtable}
講題:從衝突到共融　From Conflict to Communion

經文:哥林多前書10:17;以弗所書4:1-6

講員:蒲錦昌牧師

場合:香港中文大學崇基學院禮拜堂

日期:2017年11月5日

一,引言:

各位弟兄姊妹,平安!

今年是教會改革500周年.對普世教會而言,這是一個值得高興的日子,但也是一個叫我們傷心的日子.我們高興,因為基督教由此而誕生,代表了教會不斷改革更新的力量；我們又傷心,因為這是教會自東西方分裂為天主教(公教)與東正教(正教)之後,又一次大分裂.究竟分裂的教會能否重新復合呢?衝突能否走向共融呢?

天主教會和信義宗教會(世界信義宗聯會)早幾年前已開始思考應如何紀念今年這個特別的日子,最後共同撰寫了《從衝突到共融--2017信義宗--天主教共同紀念改革運動》這個歷史性文件,打破了前四百多年衝突,分裂的狀態,一起以行動來表達合一的決心和共融的方向.

香港基督教協進會和天主教會也趁這個機會舉辦合一之旅,走到德國威登堡,奧斯堡,天特/特倫托,梵蒂岡這些歷史現場,重新認識改革運動與探索教會合一的未來.我們又舉行教牧講座和神學院之間的研討會,希望互相認識,對話和了解,重新審視過去的歷史,用正面的態度對對方,治療教會昔日分裂的創傷,一起攜手走合一的道路.

今天,我們讀了兩段聖經,分別是林前1:10-17,弗4:1-6,一段講到分裂的危機,另一段講到合一,共融的基礎.

二,主題

1.從合一的角度出發

第一段經文是弗4:1-6節,保羅勸勉當時以弗所教會持守合一的情操:「謙虛,溫柔,忍耐,用愛互相寬容,以和平彼此聯繫」,這個合一不是憑我們的努力可以達至的,而是「聖靈所賜的合一」.但是,我們也要努力與上主的工作配合,「竭力保持」.

教會合一的基礎在於教會只是一個,我們所信仰的上主也只有一位.用保羅的話講:「身體(教會)只有一個,聖靈只有一位,正如你們蒙召,是為同有一個指望而蒙召,一主,一信,一洗,一上帝--就是萬人之父,超越萬有之上,貫通萬有,在萬有之中.」(4:4-6)不同的教會傳統和宗派,其實都有這些共通點.

怎樣由分裂的危機和處境,走向合一和共融呢?

《從衝突到共融》這份文件,作出了教會合一的五點建議,我在下面選出其中一部份和大家分享.

《從衝突到共融》239節:

「第一項要務:儘管差異更容易被看出和體驗到,但是為了強化大家共同之處,天主教與信義宗雙方應該時時都以合一的角度,而不是從分歧的觀點出發.」

我　　們有很多共通點,都相信同一位上主,相信同一位救主--耶穌基督,接受同一位聖靈／聖神的感化.我們都是基督的跟從者,都受過洗禮,成為基督徒.我們主日都返教會,敬拜上主,閱讀聖經,領受上主的話,進行聖餐禮,領受上主的恩典,接受神職人員的領導和牧養.

我們就好像同一個家庭成長的兄弟姊妹,由於血緣相同,有共同成長的經驗和歷史,雖然我們自己感覺到不少差異,但是,別人往往一眼看出我們是一家人,是兄弟姊妹,樣貌相似,言行舉止也相近.

當我們穿著神職人員的服飾和戴著硬領在街上走,別人一定不會誤會我們是佛教的僧人或道教的道長,因為穿的不是袈裟,道袍,反而,有時候有人會叫我神父,正如各位神父可能有人叫你牧師一樣.

如果能夠從相通而非相異的角度看事物,我們可能更懂得欣賞和我們不同的教會傳統和宗派.

例如,天主教有聖人,好像聖伯多祿,聖保祿,聖雅各伯,聖若望這些聖經人物,也有聖方濟,聖依納爵,聖佳蘭,大德蘭,小德蘭等.其實,基督教也有各自的屬靈領袖,傳奇人物或偉人,像信義宗的馬丁‧路德,墨蘭頓,聖公宗的湯瑪斯‧克蘭墨大主教(Thomas Cranmer),長老宗的加爾文,循道衛理宗的約翰‧衛斯理,中華基督教會的誠靜怡,高伯蘭.我們都承繼他們的傳統,傳述他們的事蹟,研讀他們的著作.

不少基督新教的宗派沒有七聖禮的觀念,但是,仍然會為青少年人進行堅信禮,接受聖職時舉行按立禮,教友結婚時為其在教會裡舉行婚禮,在教友臨終時到醫院探望,為他們禱告,陪伴他們走最後一段人生路,沒有告解禮但會聽教友做錯事後的懺悔和為他們宣赦.

只要我們多從相同的角度看,原本以為分歧的可能不一定是分歧.

2. 打破分歧,創造合一

第二段經文是林前1:10-17,是保羅寫給哥林多教會書信的開首,一開始便點出了整卷書的主題--教會分裂的危機和處境,與及合一的迫切性.

當時,哥林多教會有非常多的內部問題,(呈現在整卷書信中),其中最嚴重的是他們以保羅,亞波羅,彼得(磯法)甚至耶穌基督之名而分成不同黨派.保羅勸告他們:「你們中間不可分裂,只要一心一意彼此團結.」(1:10,《和修本》)哥林多教會面對分裂的危機,保羅也坦白承認教會這個問題.

他以基督信仰勸勉他們,耶穌基督十架的智慧--在一般人看為愚拙的智慧,來勝過世間的愚昧--在一般人眼中是智慧.正如這段經文所講:「上帝揀選了世上愚拙的,為了使有智慧的羞愧.上帝也揀選了世上卑賤的,被人厭惡的,以及一無所有的,為要廢掉那樣樣都有的,是凡血肉之軀的,在上帝面前,一個也不能自誇.」(1:27-29)謙卑是教會合一最重要的德行,而自我中心和驕傲是合一最大的障礙.

《從衝突到共融》的241節:

「第三項要務:天主教與信義宗雙方應該繼續致力於尋求可見的合一,共同計劃具體的步驟,並朝著這目標不斷努力.」

因為歷史的緣故,我們無疑承繼了前人所出現的分歧,這是不容否認的.因此,如果我們要合一,就要努力打破這些分歧,並主動創造有利的條件,以表達教會的合一.

第一個最應該努力的是聖經的翻譯.由於香港和華人教會承繼了教會分裂的遺產,我們在聖經翻譯上也各行各路,現在天主教會用《思高譯本》,基督教較為通行的則有《和合本》,《和合本修訂版》.大家在人名,地名,專有名詞等一大堆其實並無教義分歧的東西上也變成不同,這根本是不必要的,也是華人教會特有的問題.我們應該努力,使這些分歧縮小甚至消失.

過去,曾經有人倡議天主教和基督教學者共同翻譯聖經,我們應該支持和繼續這方面的努力,直至有具體的成果.希望有一天,我們只有一本共同的譯本.

第二個可以努力的方向是我們的日常用語,有些和聖經有關,有些不是.天主教叫「天主」,基督教叫「上帝」,「神」,可否多用「上主」,甚至有一天分歧的叫法消失呢?現在,連社會上一般人也把我們看為兩個宗教,因此,香港有「六宗教領袖座談會」,即天主教,基督教,佛教,道教,伊斯蘭教,孔教.大陸有「五大宗教」,即天主教,基督教,佛教,道教,伊斯蘭教.這是很大的誤會.我建議香港的六宗教領袖座談會改名,就叫「宗教領袖座談會」,不要再把天主教和基督教視為兩個宗教來看待.一般人不了解我們,以為天主教是信天主,而基督教是信基督耶穌的.這是多麼大的一個笑話!

由於天主教,基督教均有一個「教」字,因此很容易被人誤會與佛教,道教,伊斯蘭教同樣是一個教.我建議,正如佛教中有禪宗,密宗,天台宗,法相宗等,而仍然是同一個教,天主教可以叫「大公宗」(Catholic Church),東正教叫「東正宗」(Orthodox Church),基督教/基督新教叫「改革宗」(Reformed Churches),那麼,便不會被別人認為是兩個或三個宗教了.至於「基督教」應該對應英語的“Christianity”而泛指基督信仰的不同派別,正等如「基督徒」對應英語的“Christian”一樣.

使徒保羅／保祿勸告哥林多／哥林多教會不要以基督之名把教會分裂了:「我的意思是你們各人說:『我是屬保羅的』,『我是屬阿波羅的』,『我是屬璣法的』,『我是屬基督的』.基督是分開的嗎?保羅為你們釘了十字架嗎?你們是奉保羅的名受了洗嗎?」(林前1:12-13)當「天主」,「基督」只成為一部份教會傳統的名稱,都是不合宜的.

我們是否能夠移風易俗,把這些分歧改變過來呢?希望有一天,我們能夠做到.不再有「天主」,「上帝」,「神」的差異；不再有「聖神」,「聖靈」的分別；不再有「伯多祿」和「彼得」,「保祿」和「保羅」看似兩個人其實是一個人的誤會.至於「神父」和「牧師」,「主教」和「監督」,「彌撒」和「崇拜」,「信友」和「教友」,「天主教徒」和「基督教徒」,「感恩祭」和「聖餐」,「鐸職」和「牧職」,「因信稱義」和「因信稱義」等有翻譯又有神學理解不同的地方,希望透過誠懇的對話和一起的努力,也可以慢慢的減少或消除.

由於過去所形成的分歧和大家關係的對立,特別在基督教這方面,我們要扭轉昔日敵對的態度和思想,的確要花一些功夫,尤其在信徒的層面.不少基督教徒均可能讀過吳主光所寫《基督教與天主教的分別》這一類小冊子,不但只從分歧的角度看大家的關係,而且沒有向天主教方面求證而單方面的批評對方,以至製造了不少誤解,這是對天主教十分不公平和有礙教會創傷的癒合的.我的第三個建議,是由天主教和基督教的學者,教會領袖一同聯合出版小冊子和書籍,介紹教會的歷史和信仰,讓信徒和非信徒認識我們是在合一中的教會,而不是兩個敵對的宗教.

正如現在多了歷史學者寫全球史,希望可以扭轉只從西方人的角度看世界的錯誤,也有中,日,韓的歷史學者希望走出三個國家各自描述過去歷史時的盲點,而共同撰寫東北亞近代史,同樣,天主教和基督教也可以攜手重新演繹我們的歷史.

在日常教會生活中,基督新教因有這麼多的宗派,各自成一體系,結果自然發展出不同的禮儀.每個主日的崇拜中,基督徒採用的禮儀,穿的禮袍(或不穿),唱的詩歌,念的禱文,讀的聖經版本,均可能不同,這也構成了宗派之間的隔膜.試想想,如果某一個宗派的會友走到另外一個宗派的教會崇拜,發覺唱的詩歌,進行的禮儀和他／她的教會是很接近甚至相同的,他／她肯定會產生一種熟悉而親近的感覺.

過去,中國教會不同宗派合力完成的《普天崇拜》和《普天頌讚》正是這方面努力的成果,使不同宗派的教會有共同的詩歌與禮儀可以採用.天主教,正教會和基督新教,也大可以在崇拜禮儀,詩歌等方面多交流,以至形成互相學習,共同承繼的遺產.

協進會合一與教會關係委員會和天主教基督徒合一委員會兩年多前共同成立神學對談小組,就匯聚了一群不同宗派的神學教育工作者,聖經學者和牧者,完成了尼西亞信經的合一譯本(試用版),並正收集各宗派的意見以進行修訂,期望日後可以在不同宗派中一同採用,日後大家在誦念尼西亞信經時就不會出現不同文本的問題了.如果這個工作成功的話,日後希望主禱文等在崇拜中常用的禱文,也有合一譯本.

因此,我的第四個具體建議,是以主日崇拜禮儀為核心,創作更多大家可以共用的聖詩,並大力推廣在主日崇拜中運用；共譯禱文及祝文,在不同主日採用；研究聖餐禮的禮文,尋找共同之處(其實聖公宗,循道衛理宗和中華基督教會的聖餐禮文有不少共通之處),而不是各自進行禮儀更新.那麼,日後教友不論到那個教會崇拜,也會感到熟悉,舒服,溫暖和體會教會的合一共融.

3. 一同宣揚福音和服務世人

保羅在腓4:8勸勉腓立比教會的弟兄姊妹說:「弟兄們,我還有未盡的話:凡是真實的,可敬的,公義的,清潔的,可愛的,有美名的,若有甚麼德行,若有甚麼稱讚,這些事你們都要思念.」基督徒要思想好的事情,不單如此,還要把這些好事情做出來.正如保羅在腓4:9跟著說:「你們在我身上所學習的,所領受的,所聽見的,所看見的,這些事你們都要去行,賜平安的上帝就必與你們同在.」

在路9:49-50記載,約翰曾問耶穌:「夫子,我們看見一個人奉你的名趕鬼,我們就禁止他,因為他不與我們一同跟從你.」耶穌說:「不要禁止他；因為不敵擋你們的,就是幫助你們的.」

只要是有益於他人的事,就可以和其他人一起做,甚至不一定要對方有相同的身份認同.

我們作為基督的跟從者,更應該和同是基督徒的一起合作,見証上主.

《從衝突到共融》243節:

「第五項要務:天主教與信義宗雙方應該一同向世界宣揚福音和服務世人,以見證上主的憐憫.」

教會經常都舉辦宣揚福音的活動,有些規模較小,有些規模較大,基督教不同宗派就曾多次聯合起來舉辦佈道聚會.我聽過有人建議,天主教和基督教一起舉辦福傳聚會.如果這個建議得到實現,將是歷史性的,因為從來未試過.如果怕工程太大,可以由有興趣的地區教會做起.那麼,別人便不會視我們為兩個敵對的宗教,而是一個教會.這對福音／福傳工作更有說服力!

香港各教會在社會服務,教育,醫療等方面均有不少參與,在社會公義,勞工,環保等課題上,亦有所關心,若不同宗派能夠聯合起來,一起見證,力量自然更大,效果當然會更加明顯.

當今世界正面對嚴重的生態危機,很多教會和宗派都注意到這一個問題,並從信仰角度作出回應.教宗方濟各在二零一五年度發表了《願你受讚頌》通諭,表達出天主教會對生態議題的關懷,正教會的君士坦丁堡宗主教巴爾多祿茂一世亦很重視環保,有綠色宗主教(Green Patriarch)之稱.今年,兩大教會傳統就定9月1日(正教會紀念天主創造宇宙日)至10月4日(聖方濟瞻禮)為「受造界時期」,推廣環保.據我所知,聖公宗,信義宗,改革宗,循道衛理宗等教會的大會,也關注這方面的議題.我在七月份也參加了一個洛桑運動有關環保的會議,顯示出福音派教會也有相同的關注.這次會議從香港去的有20多人,都是對環保有熱情的弟兄姊妹,牧者和神學院老師.可見,環保是不同宗派及教會傳統都關注的議題.我們一起努力推動環保也很自然,因為香港是我們的家.一起做環保工作不單只可以產生更大的社會效果,而且也可以促進教會合一.

三, 結語:

教會未來是走向合一或更分裂,走向共融或更分歧呢?這要看我們有多大的信心,愛心和對未來的盼望.願上主合一的靈在我們內心動工,叫我們更謙卑,更包容,更體會上主的心腸,以至我們更有決心和勇氣.願上主使用我們,打破昔日所製造出來的分裂,創傷與仇恨,一同聯手,彼此合作,在世界裡見證上主福音的大能!阿們!



\newpage

\section{真相寶座敬拜}
\label{sec:388}
\textbf{鄧瑞強博士 - 真相,寶座,敬拜}
\newline
\newline
\hyperref[sec:387]{< < < PREV SERMON < < <}
~
\hyperref[sec:toc]{[返目錄]}
~
\hyperref[sec:389]{> > > NEXT SERMON > > >}


啟示錄 4:1-11
\newline
\begin{longtable}{cl}
\hline
\hline
章節 & 經文 (和合本修訂版)\\
\hline
4:1 & \begin{tabularx}{0.7\textwidth}{X} 這些事以後，我觀看，看見天上有一道門開著。我頭一次聽見的那好像吹號的聲音對我說：「你上這裡來，我要把此後必須發生的事指示你。」 \end{tabularx} \\ \\ \relax
4:2 & \begin{tabularx}{0.7\textwidth}{X} 我立刻被聖靈感動，見有一個寶座安置在天上，有一位坐在寶座上。 \end{tabularx} \\ \\ \relax
4:3 & \begin{tabularx}{0.7\textwidth}{X} 那坐著的，看來好像碧玉和紅寶石；又有彩虹圍著寶座，光彩好像綠寶石。 \end{tabularx} \\ \\ \relax
4:4 & \begin{tabularx}{0.7\textwidth}{X} 寶座的周圍又有二十四個座位，上面坐著二十四位長老，身穿白衣，頭上戴著金冠冕。 \end{tabularx} \\ \\ \relax
4:5 & \begin{tabularx}{0.7\textwidth}{X} 有閃電、聲音、雷轟從寶座中發出。在寶座前點著七支火炬，就是神的七靈。 \end{tabularx} \\ \\ \relax
4:6 & \begin{tabularx}{0.7\textwidth}{X} 寶座前有一個如同水晶的玻璃海。寶座的周圍，四邊有四個活物，遍體前後都長滿了眼睛。 \end{tabularx} \\ \\ \relax
4:7 & \begin{tabularx}{0.7\textwidth}{X} 第一個活物像獅子，第二個像牛犢，第三個的臉像人臉，第四個像飛鷹。 \end{tabularx} \\ \\ \relax
4:8 & \begin{tabularx}{0.7\textwidth}{X} 四個活物各有六個翅膀，遍體內外都長滿了眼睛。他們晝夜不住地說：「聖哉！聖哉！聖哉！主—全能的神；昔在、今在、以後永在！」 \end{tabularx} \\ \\ \relax
4:9 & \begin{tabularx}{0.7\textwidth}{X} 每逢四活物將榮耀、尊貴、感謝歸給那坐在寶座上、活到永永遠遠者的時候， \end{tabularx} \\ \\ \relax
4:10 & \begin{tabularx}{0.7\textwidth}{X} 二十四位長老就俯伏敬拜坐在寶座上活到永永遠遠的那一位，又把他們的冠冕放在寶座前，說： \end{tabularx} \\ \\ \relax
4:11 & \begin{tabularx}{0.7\textwidth}{X} 「我們的主，我們的神，你配得榮耀、尊貴、權柄，因為你創造了萬物，萬物因你的旨意被創造而存在。」 \end{tabularx} \\ \\
[1ex]
\hline
\hline
\end{longtable}
講題:真相,寶座,敬拜　Reality, Throne and Worship

經文:啟示錄4:1-11

講員:鄧瑞強博士

場合:香港中文大學崇基學院禮拜堂

日期:2017年11月12日

各位弟兄姊妹,早安.願你們平安.

有人說,當喬布斯(Steve Jobs)死後,蘋果公司最大的損失,是失去「願景」,「遠象」,vision.喬布斯走的每一步,可以說,由「願景」引導著,他同時以「願景」引導著他人.他的產品,不是滿足顧客的需要,而是創造顧客的需要.他不是要在已有的產品上力臻完善,而是不斷實現人類未曾想過的夢想.簡單說,他不是活在這個現實的世界上,而是活在一個想像的世界裡.

周星馳在電影《少林足球》裡的一句對白,差不多已變成香港的諺語,周星馳說:「做人如果無夢想,同條鹹魚有什麼分別」.

夢想,遠象,願景,都考驗人的想像力.

有沒有試過,在想像力的邊界上,想像那邊界外的世界?

今日,讓我們試試,看看那邊界外的景象.

今日的講道經文:啟示錄4:1－11

啟4:1　此後,我觀看,見天上有門開了.我初次聽見好像吹號的聲音,對我說:「你上到這裡來,我要將以後必成的事指示你.」

啟4:2　我立刻被聖靈感動,見有一個寶座安置在天上,又有一位坐在寶座上.

啟4:3　看那坐著的,好像碧玉和紅寶石；又有虹圍著寶座,好像綠寶石.

啟4:4　寶座的周圍又有二十四個座位；其上坐著二十四位長老,身穿白衣,頭上戴著金冠冕.

啟4:5　有閃電,聲音,雷轟從寶座中發出；又有七盞火燈在寶座前點著；這七燈就是神的七靈.

啟4:6　寶座前好像一個玻璃海,如同水晶.寶座中和寶座周圍有四個活物,前後遍體都滿了眼睛.

啟4:7　第一個活物像獅子,第二個像牛犢,第三個臉面像人,第四個像飛鷹.

啟4:8　四活物各有六個翅膀,遍體內外都滿了眼睛.他們晝夜不住的說:聖哉!聖哉!聖哉!主神是昔在,今在,以後永在的全能者.

啟4:9　每逢四活物將榮耀,尊貴,感謝歸給那坐在寶座上,活到永永遠遠者的時候,

啟4:10　那二十四位長老就俯伏在坐寶座的面前敬拜那活到永永遠遠的,又把他們的冠冕放在寶座前,說:

啟4:11　我們的主,我們的神,你是配得榮耀,尊貴,權柄的；因為你創造了萬物,並且萬物因你的旨意被創造而有的.

(1)「在空間的高處」,「在時間的盡處」看到的真相

啟4:1　此後,我觀看,見天上有門開了.我初次聽見好像吹號的聲音,對我說:「你上到這裡來,我要將以後必成的事指示你.」

天上的世界,是世間界限以外的世界,在空間的至高處.

「以後必成的事」,我理解為在歷史終點時發生的事情,在時間的盡頭.

啟示錄的作者,在神的恩典裡,看到天門大開,並得以越過天門,進到天界,看到人類遠象的極致,領悟「空間至高處」及「時間的盡處」的「異象」,vision.

夢想,遠象,願景,都是看到「現實的這樣」和「時間的此刻」以外的未出現的,卻是更真的世界.信仰,是夢想,遠象,願景的極致,讓人去到人的界限處,接觸神明.讓平凡的人,與我們無法想像的神,有一交往.以致我們這些必死的人,體會死亡以外的永恆.

沒有了遠象,沒有了信仰,人怎能在如此沉重的人生中活下去.

在2015年,香港的中學文憑試,中文作文題有一條這樣的題目:「夢想看似不切實際,其實很有意義」或「夢想看似很有意義,其實不切實際」.其實,這題目無論前半部抑或後半部,其實都假定夢想不切實際,然後問,不切實際的東西,是否可能有意義.這問題,問得很好,也問得很道地,問得很香港.言下之意,是假定不設實際的東西,不可能有意義.若假定不設實際的東西有意義,這問題便問不出來.若有人認為不設實際的東西有意義,便有責任講出理由.當然,我們知道,夢想之有意義,正正在於其不切實際.若「實際」真的那麼重要,就不會有喬布斯,不會有大家手拿著的智能電話.若「實際」真的那麼重要,就不會有那位高喊「我有一個夢」的馬丁路德金(Martin Luther King Jr.),他那不設實際的夢改變了一個國家.若「實際」真的那麼重要,則人人只能做一條鹹魚.

早一陣子,有位小學生,問香港的教育局局長:「局長,你有什麼夢想嗎?」局長答得很誠實:「細個時候有過很多夢想,很多都是不切實際…..到大個為生活所折磨……現在成為教育局長,很想努力為教育做一點事情.」(參:\url{http://thestandnews.com/society/}\%E7\%95\%B6\%E5\%B0\%8F\%E5\%AD\%A9\%E5\%AD\%90\%E5\%95\%8F\%E6\%A5\%8A\%E6\%BD\%A4\%E9\%9B\%84\%E6\%9C\%89\%E4\%BB\%80\%E9\%BA\%BC\%E5\%A4\%A2\%E6\%83\%B3/)這個答案很香港.但若喬布斯也是這樣想,則我們可能仍在使用原始的手提電話.可以想像,香港的教育,未來幾年,不必有什麼想像了.

眼界太低,目光短淺,肯定是問題.中國人說:「螳螂捕蟬,黃雀在後」.這隻螳螂,專注於眼前的小利,不曾留心生命的危在旦夕.若牠的視野能從高一點的角度看,若牠的視野能從後一點的時間回望,便能看到自身的危險,而作明智的抉擇.可惜,很多人都像這螳螂,追逐當下利益,放棄了生命.

啟示錄,讓我們從空間的最高點,時間的最盡處,看宇宙人生的真相.或許,看到真相後,我們會選擇得好一點.

(2)至高處的寶座

至高處的異象,是什麼?

啟4:2　我立刻被聖靈感動(在靈中),見有一個寶座安置在天上,又有一位坐在寶座上.

啟4:3　看那坐著的,好像碧玉和紅寶石；又有虹圍著寶座,好像綠寶石.

啟4:4　寶座的周圍又有二十四個座位；其上坐著二十四位長老,身穿白衣,頭上戴著金冠冕.

啟4:5　有閃電,聲音,雷轟從寶座中發出；又有七盞火燈在寶座前點著；這七燈就是神的七靈.

啟4:6　寶座前好像一個玻璃海,如同水晶.寶座中和寶座周圍有四個活物,前後遍體都滿了眼睛.

啟4:7　第一個活物像獅子,第二個像牛犢,第三個臉面像人,第四個像飛鷹.

啟4:8　四活物各有六個翅膀,遍體內外都滿了眼睛.……

宇宙人生的終極真相是:在至高處,神坐著為王.

神沒有形像,只能以代表聖潔,尊貴,榮耀的寶石來形容.有「虹」環繞神的寶座,舊約以「虹」為神人立約的標記.寶座周圍,有二十四位長老,代表新舊約所有聖徒.有閃電,聲音,雷轟從寶座中發出,代表神的威榮及審判.燈台代表教會.燈台上的燈,是教會發出的光明,這光明不源自教會自身,源自「靈」,聖靈.聖靈在教會裡,教會才能發出光明.寶座前有「玻璃海」,按照西亞的文明,海裡有代表邪惡的怪獸,「玻璃海」正正顯示,這海怪已被消滅了,海不再顛簸,平靜得像「水晶」一樣.寶座前有神物,獅子代表百獸之王,牛贖代表家畜之王,鷹代表飛鳥之王,人代表萬物之靈.這些神物,可說是萬物的代表,他們都圍在神的旁邊,事奉神.這四種動物,其實同時是埃及和西亞文化的神像常用的形像,如獅身人面像,同時用了獅子及人的形像.亞述的神像,是牛身,鷹的翅膀,人面的三合一.這四種動物出現在神面前,也可代表任何人間的神明,也不真是神明,只能降服在真神之下.神物滿身是眼,可代表神對萬物的明察秋毫.

這景象的細節,可能很複雜,但其核心信息,卻是很簡單.萬事萬物,無論多神聖,多利害,多吸引,都不是神,都不是人追求的終極對象,都不是人將命運交付其上的主宰.萬事萬物,都只是被造之物,降服在獨一的創造主之下.只有一位神,坐在至高的寶座上.若拜其他東西,都只能是偶像崇拜.

莊子的《逍遙遊》講到有隻大鵬鳥,「有鳥焉,其名為鵬,背若太山,翼若垂天之雲【背像泰山,翼像雲彩】,摶扶搖羊角而上者九萬里【乘著旋風而直上九萬里】,絕雲氣,負青天【超越雲層,背負青天】,然後圖南,且適南冥也【去南海,意指去到海的極處】.」地上的小麻鳥,卻譏笑這隻鵬鳥,說:「牠這樣做,有什麼意義呢?你看我,飛起來不高,飛一會便落下來.在這草叢裡飛來飛去,多快活.」莊子的評語是:「此小大之辯也.」這就是「知道什麼是偉大」和「眼光短淺如豆」的分別了.

若你在心靈的至高點看到神,若你在生命的終點看到神,則你對現在的生活,可能有不同的評價.見到至高至遠,便知何謂低俗短視.

(3)敬拜那配得敬拜的

啟4:8　……他們(四活物)晝夜不住的說:聖哉!聖哉!聖哉!主神是昔在,今在,以後永在的全能者.

啟4:9　每逢四活物將榮耀,尊貴,感謝歸給那坐在寶座上,活到永永遠遠者的時候,

啟4:10　那二十四位長老就俯伏在坐寶座的面前敬拜那活到永永遠遠的,又把他們的冠冕放在寶座前,說:

啟4:11　我們的主,我們的神,你是配得榮耀,尊貴,權柄的；因為你創造了萬物,並且萬物因你的旨意被創造而有的.

天界所有的生命,都向獨一真神發出頌讚.

頌讚誰,敬拜誰,就是以他為最有價值者,生命的全盤意義,都在乎他.你敬拜誰,就是將生命的全部意義,押在他身上.

我們總要敬拜一點什麼,我們總要找些東西給生命賦予意義,我們忍受不了意義的失落.

在獨裁的世界,如北韓,獨裁者索性扮演神的角色,要人民敬拜他們,由他們去決定人民的意義.

香港是一個金錢世界,money talks.金錢決定大學生在大學裡讀什麼科,金錢決定一個人做什麼工,金錢也決定一個人追求一種什麼樣的生活.我們再沒有空間去想像,再沒有閒暇去發夢,再沒有勇氣去反抗「凡事向錢看」的生活方式.金錢,是現代生活的上帝.

我們的社會,訓令我們滿足於現狀.可以炒樓,可以炒股,可以盡情賺錢,可以盡情揮霍.只要不越過界線,只要向社會樹立的神像下拜,只要不挑戰這神像背後的虛無,則你可安然生活,像鹹魚一樣安然.請不要在靈魂的高處,講論生命的永恆價值.請不要站在歷史的前端,渴望社會的體制變得更好.請不要站得高,請不要看得遠,現在已有的一切,已經很好了.

啟示錄的遠象,卻站在靈魂的最高處,時間的最前端,看到那終極永恆的上帝.然後,看穿一切偶像崇拜的虛假,看穿我們平時所崇拜所追求的東西的虛無,看穿我們平時賴以為生的意義結構的不堪一擊.然後,我們明白到,真正的價值在何處,真正值得我們崇拜的東西是什麼,就是那位掌管著歷史,並創造萬物的上帝.

但願,我們的靈性視野高一點,遠一點,像啟示錄的作者一樣,在至高之處,在至極之時,能看到上帝,並敬拜他.

但願:榮耀歸於聖父,聖子,聖靈.阿們.



\newpage

\section{財富才幹與在地若天}
\label{sec:389}
\textbf{鄭漢文博士 - 財富,才幹與在地若天}
\newline
\newline
\hyperref[sec:388]{< < < PREV SERMON < < <}
~
\hyperref[sec:toc]{[返目錄]}
~
\hyperref[sec:390]{> > > NEXT SERMON > > >}


馬太福音 25:14-30
\newline
\begin{longtable}{cl}
\hline
\hline
章節 & 經文 (和合本修訂版)\\
\hline
25:14 & \begin{tabularx}{0.7\textwidth}{X} 「天國又好比一個人要出外遠行，就叫了僕人來，把他的家業交給他們。 \end{tabularx} \\ \\ \relax
25:15 & \begin{tabularx}{0.7\textwidth}{X} 他按著各人的才幹，給他們銀子：一個給了五千，一個給了二千，一個給了一千，就出外遠行去了。 \end{tabularx} \\ \\ \relax
25:16 & \begin{tabularx}{0.7\textwidth}{X} 那領五千的立刻拿去做買賣，另外賺了五千。 \end{tabularx} \\ \\ \relax
25:17 & \begin{tabularx}{0.7\textwidth}{X} 那領二千的也照樣另賺了二千。 \end{tabularx} \\ \\ \relax
25:18 & \begin{tabularx}{0.7\textwidth}{X} 但那領一千的去掘開地，把主人的銀子埋藏了。 \end{tabularx} \\ \\ \relax
25:19 & \begin{tabularx}{0.7\textwidth}{X} 過了許久，那些僕人的主人來了，和他們算賬。 \end{tabularx} \\ \\ \relax
25:20 & \begin{tabularx}{0.7\textwidth}{X} 那領五千的又帶著另外的五千來，說：『主啊，你交給我五千。請看，我又賺了五千。』 \end{tabularx} \\ \\ \relax
25:21 & \begin{tabularx}{0.7\textwidth}{X} 主人說：『好，你這又善良又忠心的僕人，你在少許的事上忠心，我要派你管理許多的事，進來享受你主人的快樂吧！』 \end{tabularx} \\ \\ \relax
25:22 & \begin{tabularx}{0.7\textwidth}{X} 那領二千的也進前來，說：『主啊，你交給我二千。請看，我又賺了二千。』 \end{tabularx} \\ \\ \relax
25:23 & \begin{tabularx}{0.7\textwidth}{X} 主人說：『好，你這又善良又忠心的僕人，你在少許的事上忠心，我要派你管理許多的事，進來享受你主人的快樂吧！』 \end{tabularx} \\ \\ \relax
25:24 & \begin{tabularx}{0.7\textwidth}{X} 那領一千的也進前來，說：『主啊，我知道你，你是個嚴厲的人：沒有種的地方也要收割，沒有播的地方也要收穫， \end{tabularx} \\ \\ \relax
25:25 & \begin{tabularx}{0.7\textwidth}{X} 我就害怕，去把你的一千銀子埋藏在地裡。請看，你的銀子在這裡。』 \end{tabularx} \\ \\ \relax
25:26 & \begin{tabularx}{0.7\textwidth}{X} 他的主人回答他說：『你這又惡又懶的僕人，你既知道我沒有種的地方也要收割，沒有播的地方也要收穫， \end{tabularx} \\ \\ \relax
25:27 & \begin{tabularx}{0.7\textwidth}{X} 就該把我的銀子放給兌換銀錢的人，到我來的時候可以連本帶利收回。 \end{tabularx} \\ \\ \relax
25:28 & \begin{tabularx}{0.7\textwidth}{X} 把他這一千奪過來，給那有一萬的。 \end{tabularx} \\ \\ \relax
25:29 & \begin{tabularx}{0.7\textwidth}{X} 因為凡有的，還要加給他，叫他有餘；沒有的，連他所有的也要奪過來。 \end{tabularx} \\ \\ \relax
25:30 & \begin{tabularx}{0.7\textwidth}{X} 把這無用的僕人丟在外面黑暗裡，在那裡他要哀哭切齒了。』」 \end{tabularx} \\ \\
[1ex]
\hline
\hline
\end{longtable}
講題:財富,才幹與在地若天　Wealth, Talent and the World as Heaven

經文:馬太福音25章14-30節

講員:鄭漢文博士

場合:香港中文大學崇基學院禮拜堂

日期:2017年11月19日

1.「按才幹接受託付的比喻」的主流解讀

弟兄姊妹,平安.

今日,我又再回鄉探親一樣,返回崇基禮拜堂這個信仰大家庭.

今日,我們共聚在此,盼望聖道的臨在.

今日,我們重新閱讀大家已經耳熟能詳的故事,是記載在馬太福音25章14至30節.我們用三把聲線來讀,我讀一句,弟兄讀一句,然後姊妹讀一句,之後再循環.

25:14「天國又好比一個人要往外國去,就叫了僕人來,把他的家業交給他們,

25:15按著各人的才幹給他們銀子:一個給了五千,一個給了二千,一個給了一千,就往外國去了.

25:16那領五千的隨即拿去做買賣,另外賺了五千.

25:17那領二千的也照樣另賺了二千.

25:18但那領一千的去掘開地,把主人的銀子埋藏了.

25:19過了許久,那些僕人的主人來了,和他們算帳.

25:20那領五千銀子的又帶著那另外的五千來,說:『主啊,你交給我五千銀子.請看,我又賺了五千.』

25:21主人說:『好,你這又良善又忠心的僕人,你在不多的事上有忠心,我要把許多事派你管理；可以進來享受你主人的快樂.』

25:22那領二千的也來,說:『主啊,你交給我二千銀子.請看,我又賺了二千.』

25:23主人說:『好,你這又良善又忠心的僕人,你在不多的事上有忠心,我要把許多事派你管理；可以進來享受你主人的快樂.』

25:24那領一千的也來,說:『主啊,我知道你是忍心的人,沒有種的地方要收割,沒有散的地方要聚斂,

25:25我就害怕,去把你的一千銀子埋藏在地裡.請看,你的原銀子在這裡.』

25:26主人回答說:『你這又惡又懶的僕人,你既知道我沒有種的地方要收割,沒有散的地方要聚斂,

25:27就當把我的銀子放給兌換銀錢的人,到我來的時候,可以連本帶利收回.

25:28奪過他這一千來,給那有一萬的.

25:29因為凡有的,還要加給他,叫他有餘；沒有的,連他所有的也要奪過來.

25:30把這無用的僕人丟在外面黑暗裡；在那裡必要哀哭切齒了.』」

現在讓我們先按主流解經的方式想一想.我們現在看的和合本聖經加了一個標題是:

[按才受託的比喻；按才幹接受託付的比喻 ]

有一份教材,叫〈按才幹受責任,比喻中的兩種僕人〉,是這樣解讀馬太這個比喻:

「比喻中,主人是耶穌的自喻,僕人指的是門徒,銀子代表恩賜.

按才幹受責任的比喻,呈現兩種截然不同的僕人型態:良善忠心的僕人與又惡又懶的僕人.

良善忠心的僕人還具備幾個特點:

1.　認知僕人的角色責任:今日,神將家業(教會)託付給我們管理,我們責無旁貸,因為我們是僕人,這是信徒應有的體認.

2.　善用主人給予的恩賜:從這裡我們也可以發現一項屬靈原則:越加使用恩賜,恩賜就越加增長.因此,善用恩賜,不但個人恩賜可以倍增,神的家業也得到壯盛,凡是神的僕人,當以此為念.

3.　彼此分工不爭競比較:神給每個人的恩賜不同(林前七7),乃是為了因應教會各樣的事工,我們憑著所領受的盡心盡力去做,就能得到主的喜悅與賞賜.

至於又惡又懶的僕人:

1.　輕看自己的才幹與恩賜:如此怠惰,終於惹動主人的怒氣.這提醒我們,不要輕忽所得的恩賜(提前四14),當在小事上忠心(路十六10),才能成就更多更大的善工.

2.　對主人認識不足且偏差:當知,我們行事為人要對得起主,凡事蒙祂喜悅,在一切善事上結果子,漸漸地多知道神(西一10).屬神的人,千萬不要因為無知而得罪神,成為祂眼中的惡僕.

總言之,善用恩賜,不但個人恩賜可以倍增,神的家業也得到壯盛.」

今日講道講完.

2.「按才幹接受託付的比喻」的另類解讀

大家是否滿足?要不要想下去,這段經文,可否有其他解說?

我從本科到博士讀了十四年哲學.哲學要求的思考方法,包括:「回到基本,另類構想.」(Back to basic; alternative imagination)

主流想法以外,可否有另類解讀?以上主流解讀有三個前設:

一)主流是讚賞前二位僕人(五千,二千)的做法,低眨第三位(一千).

二)主流把主人看成是上主.

三)主流把天國成員,類比成在地國上對主人忠心的僕人.

在甚麼觀點下,可以平反一下一千僕人的評價?例如,他是不公義的抗爭者.

如果主人不是上主,而是指地上的資本家或老闆或當權者?會出現甚麼新觀點?

耶穌講的是天國的比喻,究竟從中可以讀出我們可以如何在地若天地生活?

若回到基本,基本在那裡?

大家或許知道,我過去二十多年在中大教育學院是發展價值教育的.價值教育就是返回學習的根本意義,我把數拾年思考的道理化為兩句說話:「心動力變,力用事成；事成在力,力用在心.」(Heart determines Powers realizing Tasks)

現在,讓我們從價值教育的角度看馬太福音25章14至30節,想一想這段經文,究竟在教育下一代甚麼的價值觀?

我們分為三部曲的解法.

對這段經文的解讀:呈現甚麼價值意識?宣揚甚麼教育意義?引發甚麼信念實踐?

三種解法:

是從創富,投資與回報來解(善用資本,需要在事功上在世間能夠產生分別)

是從才幹,擅長與能力來解(善用資本,需要返回到心的質素之所養成)

是從天國,地國與在地若天來解(善用資本,需要返回到在地存活的天國子民這身分)

3.「三高人士」的解讀:高資本＋高投資回報率＋高估值

這個比喻隱含對「三高人士」的高度評價,三高指的是「高資本＋高投資回報率＋高估值」.

重視成事的人間世,唯有完成事功的人才會被稱讚的.

增加資本(capital),提升投資回報率(return rate),贏取最佳評估(good will),是現代資本主義的邏輯.

香港人一定明白創造財富的道理.

若有五百萬,如何可以在最短時間來創造一倍?

如果放貸,年利率15\%以複息計,五年就可以倍增.暫時不談高利貸是否惡行.

如果放在樓市(房地產),也可以.暫時不談買賣樓房是否惡行.

但這段經文的出現,還未有資本主義的制度?

耶穌點解會用資本主義的思維去作教導?更加不會鼓勵大家投資.

這類思維正正在體現了第一種人生價值觀:價值就是願望達成(Desire-Satisfaction)；這是一種高效達標的人生觀.

成功人士其實指的是在這個世界甚有事功(worldly successful),廣東話叫:「你好叻!」

重回根本再作另類地解讀:這個比喻會否是耶穌對人世間的真實描述.

我們的世界充斥著這樣的主人,好像資本家,老闆,我們是僱員,但做到好似奴隸.

五千,二千和一千的分法,在這個故事裡,好像在討論高層,中層和低層的三層工人.

這個主人在低層工人眼中是:「主啊,我知道你是忍心的人,沒有種的地方要收割,沒有散的地方要聚斂」,這個主人在低層工人心中是:「我就害怕」.

經文沒有交代其他兩位中高層工人如何看主人,可能一樣,但只是用業績表現(performance)來克服恐懼.

老闆的資本交給我們僱工去創(造財)富,但財富不屬於我們的.做得好,只會讓自己要做更多.三個工人,其中一位在低層的被栽員(連遣散費都沒有),只餘中高層二人,但其中一位要「做埋佢個份」.(下次投資回報由五千加到一萬再變成一萬一,是否可能持續?)

那麼,我們願不願意做這種「三高人士」?

又或者問,我們其實已經在自覺或不自覺做了這種「三高人士」,跟我們的信仰有何關係?在天國運動中我的人生要繼續如此嗎?

4.「三德人士」的解讀:良善＋忠心＋可託付

這段經文,有一個有趣的地方,就是在21和23節,主人對五千和二千兩個僕人的評價都是一樣的,讓我再讀一次,但我們今次需要各位作一次想像,代入兩種僕人的心境:

會眾坐在左邊的是五千的高層人士(向他們讀出):

25:21主人說:『好,你這又良善又忠心的僕人,你在不多的事上有忠心,我要把許多事派你管理；可以進來享受你主人的快樂.』

會眾坐在右邊的是二千的中層人士(向他們讀出):

25:23主人說:『好,你這又良善又忠心的僕人,你在不多的事上有忠心,我要把許多事派你管理；可以進來享受你主人的快樂.』

讀得好!「你好乖!」

高層工人和中層工人的考績(appraisal)都寫同一段文字,在一個機構會否容許出現嗎?這個主人會否有點太也懶惰?而且會否有惡意?

無論如何,這裡有三點值得思考:

一)事功上高中層二人都在創富上翻了一翻(double),即五千多賺五千,二千多賺二千,創富回報率(return rate)是一樣的,所以二人的評價都一樣.

二)不論事功上成就多少,二人的評價都在德性方面(virtues),可以說二人都是「三德人士」,具備「良善＋忠心＋可託付」的優良品德(good qualities).

三)主人的回應,是把成事的因素返回人心(heart).事成在力,力用在心.事成與否,才幹能力(talents)是否願意用得適當,解釋項,變項是在人心.

所以,這段文字包含我們如何品評有優秀品德的(virtuous)人士,同時表示人也該如此.

第二個角度看,「三德人士」,固然地國也需要這樣的子民,天國運動需要嗎?

讓我們逐一看看「三德人士」的三種美德:

1) 「忠心」(faithful),比較容易明白,就是忠於主人的吩咐,廣東話叫:「你好乖!」

2) 「可託付」(trust-worthy),所託若是創富,也看到了.但若所託是善用才幹?

3) 「良善」(good)?究竟如何看到五千和二千這二人是良善的?

若這段經文是比喻天國,那麼天國需要怎樣的人?

這裡體現了第二種人生價值觀:價值就是品質實現(Quality-Actualization)；這是另一種展潛盡性的人生觀(potential-actualizing life).

若成功人士其實指的是在這個世界甚有事功(worldly successful),那麼成德人士指的,是活在這個界上成就在自身的品格德性(morally praiseworthy).

若根據這天國的比喻,我們如何投影到天國運動去?在地若天,就是養成這類美好的品德,好能「心動力變,力用事成.」

5.財富,才幹與在地若天

立功和立德,都有相應的評價用語.

我1995年返母校教書,第一段時期的研究是「品辭研究」,研究人如何品評人物?所用的品辭有何邏輯?

例如,在香港,家長如何稱讚兒女?老師如何稱讚學生?老闆如何稱讚下屬?

我發現常用的品辭都會用上「又叻又乖」.

其實,掌權者都希望子民「又叻又乖」.

我的解說是:

在經濟上要叻,有才幹可以創富,成為富戶,最好成為富甲一方的首富；

在政治上要乖,心會順應權力持有者而運用才幹,成為順民.

但是這是一個天國的比喻,耶穌講這個比喻一定不會叫我們在地國上要成為首富,成為順民.

所以,我們要在故事的結果來看看第三個角度:就是有份與無份(sense of belonging).

我在價值教育方面的主張,就是價值教育無法不從身分構成或身分轉化(Values education as identity formation / transformation)的角度入手.

三位僕人的人生境況是:

25:21／25:23主人說:『好,你這又良善又忠心的僕人,你在不多的事上有忠心,我要把許多事派你管理；可以進來享受你主人的快樂.』

25:30把這無用的僕人丟在外面黑暗裡；在那裡必要哀哭切齒了.』」

最終是前二者有份留在主人領域中與第三位是無份留下(「無得留低」).

這裡體現了第三種人生價值觀:價值就是身分履行(Identity-Enactment)；這是又一種盡義履分的人生觀.

我們可以叫這種人為「稱義人士」,指的是在所屬的世界有份,用百般的心與力把當中的美好體現出來.

天國子民,就是一種身分.第三種人生價值觀來解說在地若天,我們在地國上學習成為天國子民的身分轉化.

一千僕人有何值得欣賞的地方?

比喻沒有說明五千和二千僕人如何在有限時間內倍增財富,主人也不管當中的是非對錯.

有限時間內謀取暴利其實可以是作了惡行.主人也不管當中的是非對錯.(可能他們只是跟隨主人的慣常做法.)

但一千僕人卻沒有仿效同行,只作「保本」的動作,不作惡行謀取暴利,不為世間價值觀所動.所以,代價就是被世間所放棄.

作主門徒,可能會被被世間所放棄.

一千原銀奉還,主人的批評是(看26-27節):

25:26「主人回答說:『你這又惡又懶的僕人,你既知道我沒有種的地方要收割,沒有散的地方要聚斂,

25:27就當把我的銀子放給兌換銀錢的人,到我來的時候,可以連本帶利收回.」

連本帶利,其實是最低度的創富法.為何不做?

這裡有一個懸謎:究竟是一千僕人在才幹上是不懂做,還是他在心意上不願做?是力的問題,還是心的問題?

作為天國子民,我們都正是活在一個無間道,在眼見不到的天國與苟要存活的地國之間,要出世而入世,要在資本主義可以活得下時又要體現出更大的世界.

在地若天的信念實踐是怎樣的一回事?在這個要求資本善用的世代(the age of best use of capitals),我們如何活出一個為世所用的人生?

讓我們一起祈禱,求主啟導.



\newpage

\section{我是誰}
\label{sec:390}
\textbf{梁元生院長 - 「我是誰?」}
\newline
\newline
\hyperref[sec:389]{< < < PREV SERMON < < <}
~
\hyperref[sec:toc]{[返目錄]}
~
\hyperref[sec:391]{> > > NEXT SERMON > > >}


講題:“我是誰?--信仰的考問”　Who Am I?

經文:馬太福音16

講員:梁元生院長

場合:香港中文大學崇基學院禮拜堂

日期:2017年11月26日

一．前面的話

各位弟兄姊妹,各位會友:

今天的講章是我為兩個月前的崇拜講道而準備的.根據的經文是既定的聖經經課(Lectionary)所列的馬太福音十六章.但是那天打風,不好風球高掛,沒有崇拜,所以林豪恩校牧問我可否在十一月底用原來的講章和大家分享.我說:上天既然禁止我在八月講道,是否我的講章有問題.因此我就仔細地把原來的講稿多看幾次,做了一些修改.最大的改變是:原來的預算,是以“我是誰?”這個問題帶出兩方面的思考:一種是哲學的思考,另一種是信仰上的思考.前者是人對自己的叩問和身份的尋索,甚麼是存在的意義, 甚麼是活著的本質?換言之,“我是誰?”是從人出發,所問的是一個存在的問題,也可以說是一個哲學的問題.“我”究竟是什麼,為何而存在?活著的真正意義在哪裡?“我是誰?”這個問題的背後,是尋求認識自我(self)和尋索身份(identity)的嘗試,希望藉此找到人生存在的意義(the meaning of existence).我們問這個問題,背後是想要對“活著”和“存在”的意義作出嚴肅和認真的反思,其最終目的是要認識自己.由於人不是一個孤立的存在,在他追求認識自己的過程中,他不免要向周遭的環境看,包括往前面看及往後面看,因為人活著並不孤立,而是活在眾多的人際關係裡和歷史的長河中,他看到自己在混雜於簇擁的人群中,他看到自己在歷史的長河中漂流,前面有滾滾水流,而後面則有接踵而至的來的巨浪洪濤.連綿不斷,生生不息,故哲人嘆曰:逝者如斯乎,不舍晝夜.因此之故,我們不但會問:“我是誰?”也會問:“我從那裡來?我往那裡去?”,意圖使人明白其存在不是一個孤立的個體,而他的存在實際上是在一個連續不斷,生生不息的處境之中,因而其存的意義在及其身份的建構,也因為處於這個處境而產生了特別的意義.

我當然不奢望在今天的半個小時中去解決這些人生中永恆的哲學難題.今天要討論的是第二個方面,即是從“我是誰?”這個問題去思考信仰和檢視我們自己的信仰.由於“我是誰?”是耶穌向門徒發出的問題,而非人想自己發出的問題,所以今天的題目應該加上一個副標題,那就是“信仰的考問”.

二．耶穌的問題:“我是誰?”

上個主日鄭漢文教授用一套反傳統的方法,帶領我們從新演繹主人和分派銀子的比喻,今天我想回到傳統的解經方式,從經文的時,地,人的背景入手, 去瞭解這段聖經故事.

首先,讓我們回顧一下耶穌提出這個問題是在甚麼時候和怎樣的背景下提出的?“我是誰?”是耶穌自己提出來的問題,請大家留意祂是在怎樣的背景下把這個問題提出來的,祂問這個問題時在甚麼地方? 背後又有甚麼原因?

讓我們再翻看馬太福音十六章13-14節:

聖經說:“耶穌到了該撒利亞腓立比的境內,就問門徒說:人說我人子是誰?他們

說:有人說是施洗的約翰,有人說是是以利亞,又有人說是耶利米或是先知裡的一

位. 耶穌說:你們說我是誰?西門彼得回答說:你是基督,是永生神的兒子.”

耶穌在馬太十六章這裡問的問題:“我是誰?”相信大家都耳熟能詳,不知讀過多少遍,不知聽牧師在講臺上講過多少遍.尤其是我們對於西門彼得的回答(算是標準答案吧)都非常熟悉:耶穌問:我是誰?彼得回答說:你是基督,是永生神的兒子.

既然答案已經非常清楚,那麼,我們還為什麼要重新思考這個問題的意義在哪裡?正如我在開頭時說過:我是誰?這個問題,可以連帶出“我從那裡來,我往那裡去?”的相關問題.那麼,讓我們就從這個相關的問題出發,當耶穌問“我是誰?”這個問題時,當時的時間和環境是怎樣的?

我們在甚麼時候會問自己:“我是誰?”很可能是對人生感到迷惘,對生命的價值和意義有所懷疑的時候.耶穌為什麼在這個時候提出“我是誰?的問題呢?是他感到人生的困惑嗎?還是他的門徒到此時陷入迷惘和失去方向?讓我們先弄清楚當時的背景因素.聖經很清楚的告訴我們,耶穌問這個問題的時候,處身在加利利海以北該撒利亞腓立比的地方.那時,門徒已經跟隨他又一段日子,聽過他在山上講道,教訓人群,也看過他治病趕鬼,叫破腳的人行走,叫瞎眼的人看見,叫死了的人復活；有隨著他坐船渡過加利利海,平靜風和浪,也看過他跳下海中,踏波而行.不少人都親耳聽過他的宣講,也有不少人親眼看過他的神跡. 換言之,已經有很多人知道耶穌是個異人,奇人,是個不平凡的人.故此,當耶穌問:人說我是誰?的時候,這幾個答案就出來了:施洗約翰,以利亞,耶利米,或先知裡的一位.這些答案的出現, 是人們對耶穌有了一些認識,因為看到耶穌所做的工作,便在猶太人的傳統裡尋找可以比擬的人物.人們首先想到的,是耶穌是以利亞.以利亞是猶太人傳說中的神人,滿有大能的先知.他有治病趕鬼的能力,也能叫死人復活.兩個多月之前,林豪恩校牧在這裡, 用舊約列王記上十六和十七章比較詳細地介紹過以利亞這位來去無蹤,與阿哈王對抗的先知.大有能力,能叫天火焚祭,能叫天降大雨,能叫油用不盡,能叫死人復活.以利亞在猶太人的宗教傳統中,是個能夠治病趕鬼,顯示奇跡的神人,這是猶太宗教最核心的部分.以利亞代表著非凡的力量,能夠在這個理性的凡間充分發揮超越的神聖的力量.他們認為耶穌的所做所為,與以利亞相近.故此以耶穌為以利亞的化身.也許我們當中也有人作如是想.我們冀求一種超凡的神力,介入我們的生命,無論是醫治疾病也好,是扭轉乾坤也好,我們希望耶穌像以利亞一樣,用神跡改變世界.

也有人以為耶穌是耶利米,這也是從猶太人的歷史和文化傳統中去尋找依據的. 耶利米是在以色列人被擄時期的先知, 他和他們一同受苦,一同流淚,對被擄為奴的猶太人抱著悲憫之情,使他們在困苦中不至失去希望. 把耶穌看作耶利米的人,大概感受到他的同情,愛心和憐憫.我們相信耶穌,是否也是因為感受到他的仁慈和愛心呢? 至於說他是施洗約翰的人,一方面可能知道耶穌和他的表兄的親密關係,另一方面可能以他們的生活態度和言行教誨相比較,宿荒山,行曠野,喊著說:天國近了,你們要悔改.其生活言行都有相同之處,故此認為他就是約翰.認為耶穌是施洗約翰的人,也許會是因為約翰敢於反抗當代的統治者希律王,他的抗議和不屈,最後賠上他的性命.在某些人看來,他是個勇敢犧牲的革命志士.這些答案都有一些根據,不是天外飛仙,胡亂猜測.這是我們一般人常常會做的事:用傳統的成說,或生活中相近的事物作比附,作為我們認識一個人或一件事的基礎.

耶穌問:我是誰?有人說是以利亞,有人說是耶利米,有人說是施洗約翰.這些都無關重要,甚至是totally irrelevant. 這裡的“有人”,實際上是“他者”,是其他的人.這並不是耶穌詢問的對象.甚麼是信仰?甚麼才是信仰的叩問?信仰不是他人的答案,不是“有人說我是誰.”而是你的真誠回答.耶穌在這裡所要強調的是:“你們說我是誰?”他問的是“你們”, 而不是他者.他要的答案,也不是在傳統內憑想像的臆測,不是只靠人云亦云. 耶穌要求的是“你”的答案.耶穌問:“你們說:我是誰?”這也是今天要問大家的問題:“你們說:我是誰?”“你們說,耶穌是誰?”

我給大家半分鐘的靜默時間, 你在心內去誠實作答.不用發聲,更加不用喧嘩.對於你個人來說,耶穌是誰?是文化傳統中的能人,神人如以利亞?是用愛心和憐憫撫慰你內心傷口的人,就像耶利米?抑或代表著勇於抗爭,挑戰權威的鬥士,如施洗約翰?許多人或許都不會直接回答,因為他們不用開口,大家都借用了彼得的答案.當耶穌問:“你們說我是誰” 時 ,彼得就立刻上前,沖口而出說:“你是基督,是永生神的兒子.”這個答案被耶穌所肯定和稱讚,因此也就成了我們的“標準答案.”

如果耶穌向我們問:你們說我是誰?我們很多人都應該會學彼得這樣回答:主啊,你是基督,是永生神的兒子.

今天讓我們來檢驗一下這個“標準答案”.

三．門徒的答案:“你是基督!”

彼得說:“你是基督,是永生神的兒子.”這個答案吧耶穌的身份說的正確,和說得清楚, 連耶穌也稱讚.然而,從後來一些事的發生,叫我們對彼得提出的答案也產生了疑問. 因為這次耶穌問完不久之後,他和門徒離開了該撒利亞,轉向南行.來到一處高山,就帶了彼得等幾個門徒,暗暗地上了高山,在他們面前變像,衣服放光,他的旁邊也出現了兩個猶太歷史中的兩個偉人,一個就是剛才說過的以利亞,另一個是摩西.這就顯示了以利亞和摩西在猶太人歷史,文化和宗教上的不容或缺的重要地位.而耶穌是基督,甚麼是基督?就是兩個重要傳統中見的核心人物,是注重神跡和法力的傳統,以及注重典章和法律的傳統之間的連繫人物,是調和兩個傳統的中道者.最後,兩者俱不見,只餘一耶穌.可見認識基督是最為重要的,不可取代的.以利亞和摩西所代表的猶太人的傳統,或者說神跡奇事及典章制度所代表的基督教元素,都很重要,但最終皆煙消雲散,不是信仰的核心和本體.故聖經說:門徒四看空無一人,只有耶穌.

彼得的確認識耶穌是基督,是永生神的兒子.但他所理解的基督和耶穌說的基督仍有頗大的差異.彼得所說的基督是“王者基督”, 是榮耀之王, 是將要降臨的彌賽亞.他之來到世間,恍如君臨天下,使凡是跟隨他的也覺得與有榮焉.而且,與他同在,除了得榮耀,也要得著無憂與安樂.故此,當他看到耶穌與身邊的兩位偉人顯現時,他即作出呼籲:由他建造三座帳棚,給他們三人居住.與主同住帳棚的意思,就是沒有憂愁, 不再思念山下的事,留在山上的帳棚中,與主基督共用歡愉快樂.彼得雖然對基督的身份有所認識,但他對基督使命的理解,顯然和耶穌所想的又很大的不同.換言之,他的認識“基督”,止於“基督身份”(Identity)的層面,而對“基督命運”(destiny)的認識,尚未能有真正的瞭解.換言之,他對基督的身份有正確的認知,但對基督的使命還是不清楚和不明白.究竟基督,這位永生神的兒子,要帶領他們往那裡去呢?

我們從“我是誰?”這句話作為引子,就希望帶出“身份”(identity)

和“命運”(destiny)的分別.這也是耶穌在問完“我是誰”之後希望門徒思考的問題.

門徒跟隨耶穌,聽他的教訓,看他的治病趕鬼的能力,希望對這位夫子認識更多更深.他們跟隨他走遍加利利海周邊地區,來到該撒利亞腓立比,亦即他們走過的最北之地.此時此地,耶穌問他們:“我是誰?”為要門徒對信仰有正確的認識.這事之後,耶穌就帶著他們轉道向南,朝耶路撒冷進發.當中經過迦伯農,撒瑪利亞等地,最終來到耶路撒冷.轉道往南,來到耶路撒冷,耶穌很清楚這是他的命運:到了耶路撒冷, 他將會受到捉拿,審問,拷打,最後被釘十字架.耶穌知道自己的悲劇命運,就從轉到向南那一刻開始.然而為了成就救贖大業,既知其苦,仍然上路.故此一路以來,他多次向門徒講到“基督”不是要作王,而是要受苦的資訊.“我往那裡去?”“我要回到天國,為你們預備地方.” 門徒如彼德之輩在認識到耶穌即基督之後,都感到身為彌賽亞的隨從而自傲,但卻不知道悲劇性的命運即將臨近.他們誤會以為耶穌將以王者的身份駕臨聖城,甚至認為他會成為榮耀的新王,開邦立國,快將創立一個新的以色列國.“我往那裡去?”耶穌和門徒對這個問題的答案有很大的差異,有很不同的憧憬.

你對耶穌有怎樣的認識和理解?如果你好像彼得,知道所信的是誰,知道他是基督,是永生神的兒子.那麼,你的信仰根基是穩固的和實在的, 你對耶穌的身份有著清楚地瞭解.但是,這種信仰的理解會對你的人生,你的命運,會有怎樣的影響?你的人生和命運,和你的信仰是否是兩個不相關的問題?

四．結語:往南走

我們在人生的道路上,是否經歷一片蒼茫,不辨方向；抑或時常都遇著迷惘,失去信心?我們不斷地尋找可以使我們重拾信心,重新上路的引路人.也許我們當中有些人會遇到過某位高僧,某位道長,或者某位先知.但在座的大多數的我們,遇見了基督耶穌.當我們的人生之路到了該撒利亞腓立比的時候,我們遇見基督.亦即是說,當我們的人生走到最北的地方,走入寒冬的季節,我們遇見基督.他問:我是誰?這個問題叫我們認識他,追尋他,去明白他的身份,認識他是永生神的兒子.這是我們找到信仰的開始.

有了信仰,我們的人生將會經歷一個大轉移,或者說是大轉彎,大轉向.(a great turnaround) 以前是往北走的,是越走越寒冷的人生, 現在有了耶穌,明白他是基督之後,人生就有一大逆轉,轉向南行,奔向聖城,奔向耶路撒冷.這是我們人生中最重要的抉擇,也是生命中最為重要的信仰的轉向(the spiritual turn). 雖然信仰的真知和信仰帶來的命運,仍然是糾纏不清,是叫人不易明白的一回事.

然而,自此之後,你的信仰將會在你的生命中發揮著積極的作用和影響.今天,就暫時和大家思考到這裡.願主繼續祝福各位!使各位的人生經歷屬靈的大轉向, 從此一路向南,把命運交托耶穌!



\newpage

\section{無人不以為然的奧秘}
\label{sec:391}
\textbf{伍渭文牧師 - 無人不以為然的奧秘}
\newline
\newline
\hyperref[sec:390]{< < < PREV SERMON < < <}
~
\hyperref[sec:toc]{[返目錄]}
~
\hyperref[sec:392]{> > > NEXT SERMON > > >}


提摩太前書 3:16
\newline
\begin{longtable}{cl}
\hline
\hline
章節 & 經文 (和合本修訂版)\\
\hline
3:16 & \begin{tabularx}{0.7\textwidth}{X} 敬虔的奧祕是公認為偉大的：神在肉身顯現，被聖靈稱義，被天使看見，被傳於外邦，被世人信服，被接在榮耀裡。 \end{tabularx} \\ \\
[1ex]
\hline
\hline
\end{longtable}
哥林多前書 11:23-29
\newline
\begin{longtable}{cl}
\hline
\hline
章節 & 經文 (和合本修訂版)\\
\hline
11:23 & \begin{tabularx}{0.7\textwidth}{X} 我當日傳給你們的是從主所領受的。主耶穌被出賣的那一夜，拿起餅來， \end{tabularx} \\ \\ \relax
11:24 & \begin{tabularx}{0.7\textwidth}{X} 祝謝了，就擘開，說：「這是我的身體，為你們捨的；你們要如此行，為的是記念我。」 \end{tabularx} \\ \\ \relax
11:25 & \begin{tabularx}{0.7\textwidth}{X} 飯後，他也照樣拿起杯來，說：「這杯是用我的血所立的新約；你們每逢喝的時候，要如此行，來記念我。」 \end{tabularx} \\ \\ \relax
11:26 & \begin{tabularx}{0.7\textwidth}{X} 你們每逢吃這餅，喝這杯，是宣告主的死，直到他來。 \end{tabularx} \\ \\ \relax
11:27 & \begin{tabularx}{0.7\textwidth}{X} 所以，任何不按規矩吃了主的餅，喝了主的杯，就是干犯主的身體和主的血了。 \end{tabularx} \\ \\ \relax
11:28 & \begin{tabularx}{0.7\textwidth}{X} 人應該省察自己，然後吃這餅，喝這杯。 \end{tabularx} \\ \\ \relax
11:29 & \begin{tabularx}{0.7\textwidth}{X} 因為人吃喝，若不分辨是主的身體，他的吃喝就是定自己的罪了。 \end{tabularx} \\ \\
[1ex]
\hline
\hline
\end{longtable}
講題:無人不以為然的奧秘    Great Is the Mystery of Godliness

經文:提摩太前書3:16;哥林多前書11:23-29

講員:伍渭文牧師

場合:香港中文大學崇基學院禮拜堂

日期:2017年12月3日

引言:將臨期是終末的開始 (Beginning of the End)

在第二次大戰1942年,英軍在北非阿拉曼戰役(El Alamein)首次打敗沙漠之狐隆美爾(Johnannes Eugen Rommel)統領的德軍.之前德軍未輸過,之後英軍未敗過.這役之後英軍士氣大振,說戰爭完了,我們得勝了.首相邱吉爾(Winston Churchill)發表演說:「現在不是結束,也不是結束的開始,也許或且可以說:是開始的結束(End of the beginning).」Now this is not the end. It is not even the beginning of the end. but it is, perhaps, the end of the beginning.

有一場戰役,決定性的屬靈戰役,各各他之役,基督藉著死敗壞那掌死權的,就是魔鬼,帶來終末的開始Beginning of the end.撒旦雖然現在還裝腔作勢,若無其事,但基督在十架其實已一役竟其功,徹底打敗撒旦,我們現在祇是執埋手尾,等候基督再來,更新萬有.

各各他一役給我們看窺探到終末的世界:罪得赦免,身體復活；羔羊與獅子共處,社會一體共融,不分高瑞低端；沒有眼淚,沒有疼痛的新天新地.這是奧秘,因這一切前所未聞,想也未想過,竟然發生了.阿拉曼一役是開始的結束End of the beginning, 但各各他一役是終末的開始Beginning of the end,因我們看到終末的景象了.

這引發虔敬的奧秘感,攝住了早期教會的敬拜者,他們在崇拜中頌唱:「大哉!敬虔的奧秘,無人不以為然,就是.神在肉身顯現.被聖靈稱義,被天使看見,被傳於外邦,被世人信服,被接在榮耀裡.」(提前三16)

將臨期也是終末的開始Beginning of the end.將臨期不單預備迎接聖誕,也預備迎接基督的再來.今天所讀的經課其實說到終末的審判,這審判的主題,剛過去的主日是教會年最後的主日,稱為基督君王主日,福音經課太二十五章,論及分別羔羊與山羊的終末審判.

但到終末的路漫漫兮修遠,我們不知基督何時再來.個人信仰的經驗也如是,信耶穌是新生命的開始,舊生命的結束,但舊生命仍頑抗,我們舊人仍然有「盈餘的邪惡」(prevalent filth and evil).新舊生命的糾纏,靈欲相爭,至死方休.將臨期是自省的節期,聖帷紫色,紫色是王者尊貴的顏色,迎接光明的來臨,就要抹去生命中的幽暗,儆醒等候.

在等候基督復臨的漫漫征途上,耶穌設立了聖餐,保護我們,對付撒旦的攻擊.聖餐亦稱為奧秘；我們透過聖餐禮儀,學習明白基督福音的奧秘.

1. 聖餐是奧秘

奧秘mysterion是希臘文,拉丁文繙作聖禮sacramentum.在東方教會,耶路撒冷的大主教西銳爾(Cyril of Jerusalem, 313-386 A.D.)的義理講解 (Catechetical Lectures) ,在講論十八:論道,一聖教會,肉身復活並得永生之後,跟著五篇短短的講論 (十九至二十三篇),統稱為奧秘 講解 (Mystagogics Lectures).五篇依次論奧秘(mystery),論聖洗(holy baptism),論抹油(chrism) ,論基督的身體,論聖禮及領餐,均在復活日領洗後的一周講解.[i]

在領洗前夕,總結過去義理講論的第十八講論說:「但現在神聖的逾越節臨近了,你們在基督裡蒙愛的,將被重生的洗光照,你們將會再被教導,若上主許可,甚麼是必須的；當召喚你們時,應以何等的虔敬和態度回應聖洗奧秘的施行.你們應以何等敬畏和態度,從聖洗走向上主的聖壇(聖餐),享用它們屬靈和天上的奧秘.」[ii]

當代天主教也依照古教會的傳統稱聖餐為奧秘.「感恩祭禮儀:基督徒受洗歸入基督的死和復活的逾越奧秘(Paschal Mystery),得著罪蒙赦免的豐盛救恩.每次的彌撒(崇拜)我們慶賀這奧秘,記念基督的愛及感謝讚美上主.」

這奧秘就是上主的愛:藏於三一,成於歷史,證於禮儀上主的愛.

2.藏於三一(hidden in Trinity),成於歷史(accomplished in history),證於禮儀(witnessed by liturgy)的奧秘

沒有人曉得藏於上主本體父,子,靈的情格相融相愛,直到基督降世為世人贖罪,把父神的愛表明出來.在大祭司的禱告中,耶穌更以父與子的合一和相愛,作為信徒內部團契的楷模.「我在他們裡面,你在我裡面,使他們完完全全合而為一,像我們合而為一.……也知道你愛他們如同愛我」(約十七23).對外的宣教,也以父差遣子的模式為典範:「你怎樣差我到世上,我也照樣差他們到世上.」(約十七18)

這奧秘直等到基督再來時,才圖窮匕現,完全明白.在基督再來之前,我們用崇拜的禮儀來見證這行經時間,顯在空間,成於歷史的奧秘.「你們每逢吃這餅,喝這杯,是表明主的死,直等到他來.」(林前十一26)

奧秘顯明上主的愛.魯益斯(C.S.Lewis)說,道成肉身的奧秘是最大的神蹟.最大的神蹟不是復活,是道成肉身.道若能成為肉身,為甚麼不能水變酒呢?為甚麼不能使瞎眼看見,跛腿行路,死人復活呢?相比之下,這些神蹟輕而易舉哩.基督成為肉身,就像一個人,深深呼吸,閉著氣,潛入深海至深之處,尋找失去的寶貝.愈潛入海底,就遠離陽光,愈來愈黑暗,愈來愈寒冷,最後到達了那幽黑冰凍的海底,把淤泥撥開,找到埋在下面的一顆明珠,然後緊緊握著這明珠,慢慢升上海面,高舉這顆明珠,使它原來的光澤重現世間.我們就是被埋在深海至深之處的明珠,罪沾污掩埋了我們原來的光澤.基督成為肉身,把我們尋回.

3.如何吃就是如何的人 (We are how we eat)-如何吃是回應上主的愛

聖餐教導我們與教會,與上主,與世界,與自己的關系.(1)聖餐是我們的身份證,吃聖禮,成為基督的奧體.(2)聖餐是我們的通行證,提起基督的工作,求父記念我們；藉基督,在聖靈裡,到父面前.這兩點是上次講道時的重點,是上主給我們的愛.今日的講道,我們思考如何回應上主的愛:(3)與世界的關係.聖餐提醒我們回應上主的愛,效法基督對世人說:這是我們身體,為你們捨的.(4) 教我們過一個感恩呈獻的生命.基督徒是最順從的僕人.

上次講道的題目是「吃甚麼就成了甚麼」(We are what we eat).但成為甚麼,又不單單繫於我們知道吃甚麼,乃在於我們如何吃(We are how we eat).吃甚麼是知識,如何吃是實踐,信心與行為不能分割,篤信力行.真正的虔誠,是看顧患難中的孤兒寡婦,並且不沾染世俗.

哥林多教會混亂聖餐,是在聖餐進行之前的愛筵,是如何吃的問題.「你們聚會的時候彼此分門別類……算不得吃主的晚餐；因為吃的時候,各人先吃自己的飯,甚至這個饑餓,那個酒醉.你們要吃喝,難道沒有家嗎?還是藐視神的教會,叫那沒有的羞愧呢?」(林前十一 20-22) 請注意,他們的罪是叫那沒有的羞愧,用最流行的話就是叫教會內的「低端人口」羞愧.他們吃愛筵的表現,不像基督的身體.如何吃是愛的問題,他們沒有愛弟兄.

審判的關鍵是篤信力行的愛.在剛過去的基督君王主日提到基督再來的審判,福音經課選自太二十五31-46.審判的日子,不是按我們所宣稱的信,乃以行為實踐出來的信.「因為我餓了,你們給我吃,渴了,你們給我喝；我作客旅,你們留住我…….」(太二十五35)

基督在我們願意分享時顯現.我們從以馬忤斯的故事,看到基督在聖餐的臨在.基督在何時向兩位門徒顯現呢?他們聽完耶穌解釋聖經後,以行動實踐信仰,接待一位陌生人共用晚膳,分享食物,並請他留宿.(路二十四29)兩位門徒所做的,正是耶穌在聖餐中所作的:「這是我的身體,為你們所捨的」.這是我的食物,與你共享；這是我的居所,請你留宿.就在此刻們看到了基督.

4.基督徒是最順從的僕人(Christian is the most dutiful servant)

今年是馬丁路德改教運動五百周年紀念,路德著作豐碩,單是美國英文版(1957-1986)已五十五冊,德國威瑪(Weimar, 1883-2009)德文,拉丁文版共一百二十一卷.路德認為最重要也是自己至喜愛的有兩本著作,1529年的教導孩童義理的小問答 (Small Catechism) ,及1520年的論基督徒的自由Freedom of the Christian Man.兩本都是薄薄的小書.論基督徒的自由,行文溫婉平和.

這書開始不久,有兩句話廣為人引用的話,不少它路德宗教會把它刻在牆壁:「基督徒是最自由的君王,不受任何牢制；基督徒是最順從的僕人,服侍每一個人」(A Christian man is the most free lord of all, and subject to none; a Christian man is the most dutiful servant of all, and subject to every one).最自由的君王,因為連罪都不能制牢,甚麼能制牢我們呢?我們不受罪的制牢,因為信,憑信得到罪的赦免.但到人面前,就要顯出愛心,作為對上主的愛的回應.而信心在愛中,才能得以活潑常存 (faith active in love).憑信,獲得基督奧體的身份證；憑信,得到往父那裡的通行證.但回應上主的恩典,我們到人面前就要顯出愛心,效法基督對世人說:「這是我們身體,為你們捨的」(This is my body given for you).

消費主義剛剛相反:「這是你的身體,為我所享用.」(This is your body taken by me).旅游公司為了利潤,讓司機超時因工作,結果產生交通意外.工廠為了省錢沒做足安全設施,引至工人嚴重工傷.色情文化產業,物化身體,當作商品售買.但聖餐把藏於三一上主的愛與共融,藉基督道成肉身顯在歷史,再藉崇拜禮儀重演在眼前.「這是我的身體,為你們而捨的」,顛覆了世界的價值觀.

因為聖餐,我們重視城市中的「低端人口」.將臨期的燭光,照亮我們的眼睛,看清楚在至高之處,猶如從外太空看地球,地面沒有高瑞,低端,一切都是不同顏色,不同圖案的平面.在至高之處,榮耀歸於上主,在至低的地上,憐恤歸於被邊緣化的異鄉人.我們求滿有憐恤的上主,將這些人舉自塵土,使他們信上富足.基督臥馬槽,成為木匠之子；叫我們曉得,簡樸可以充實,卑微亦能顯揚尊貴.人看來軟弱的聖嬰,他朝終成大能的救主.

5.以感恩為祭的人生

我們聖餐禮文有一句:「在基督的祭司的職份裡,我們聯合在一起,向你獻上這祭禮」.我們呈獻這感恩祭,直到終未基督再來.我們以感謝為祭,過每天的生活.我們行經時間,把生命所結的果子,呈獻給未日顯現的基督；這點令我我想起保羅.

保羅信了基督之後,看萬事如糞土,為要得著基督,在基督再來時呈獻禮物給他.「我們的盼望和喜樂並所誇的冠冕是甚麼呢?豈不是我們主耶穌再來的時候,你們在他面前站立得住嗎?因為你們是我們的榮耀,我們的喜樂.」(帖前二19)甚麼是保羅在基督再來的時呈獻的禮物呢?

基督再來原文是parousia.這字形容君王來到他的屬土時,當地的子民將呈獻一頂冠冕給王,作為順服王權的象徵.每個省市都用最好的金銀珠寶,競爭成為最佳冠冕,引以自豪誇耀的冠冕.保羅說:我所誇耀的就是你們.你們是福音奧秘所結的果子,我像母親一樣乳養你們成長,像父親一樣凡對你們有益的,毫不避諱告訴你們.

上次講道結束時,我提到一故事.多年前我在芝加哥參加教牧研討會,講者是一位神學院宗教教育系主任.他坦言跟我們分享,預備這場講演,感到極大壓力,睡不著.後來轉移了講演的動機:從關注自己表現,轉到建立聽眾.表現(performance) 是「我在你們面前站立得住」；保羅建立信徒的生命,關心「你們在他(基督)面前站立得住」.

為喬布斯(Steve Jobs)寫傳記的作者華萊.以撒遜(Walter Isaacson)與喬布斯作了超過四十次訪談,最後一次在他家中的後花園進行,是陽光普照的下午,當時喬布斯身體很不舒服,喬布斯把話題轉到宗教和死亡.當問及他是否相信死後生命(afterlife),是否相信有上帝時,他說:「五十五十.這是大奧秘,人永遠不知道,但我希望死不是生命的終結,我但願累積的智慧不會隨死亡歸於無有,可以永續.但他跟著說:但可能它真的像開關制,click一下,你就離開了；停頓一秒後他笑說:所以我們蘋果產品沒有開關制的.」

傳記作者記述了喬布斯六歲的時候,他的同年鄰居曉得他是收養回來的便告訴他:「你父母遺棄你,不想要你.」立時小喬布斯跑回家問過究竟,養父母對他說:「你是特別的,我們選上你.」這事令喬布斯感到自己與眾不同.我在想,會不會喬布斯一生為了證明自己與別不同,他追求完美,追求卓越,甚至諸多挑剔,辱罵同事平庸不達標,因為他要證明自己是與別不同的,以表現來肯定自己的價值.

保羅因為基督的愛,使他感到自己的價值.表現(performance)會失手的,表現令我們一生誠隍誠恐.我們的良心會久不久提起,我們過去做過,令我們感到懊悔,羞愧的事情,撒旦更藉此控訴我們.但耶穌叫我們記念他為我們所作的,記念就是提起anamnesin,給我赦罪的確據,不再誠隍誠恐.在基督裡,我們已不被定罪,是最自由的君王.此生我們一無牽掛,專心一志以感恩為祭,以愛心作眾人的僕役,回應上主的恩典.

結語:萬王之王的來臨時,我們要獻上冠冕,這是我們給基督的禮物,因為他首先把生命獻上,作禮物賜給我們,使我得著永生,不被罪的控訴,成為最自由的君王,不受任何牢制.我們以感恩為祭,甘心為僕人,回應基督的愛,服侍有需要的人.

[i] 「The Catechetical Lectures of Cyril of Jerusalem」, Archbishop of Jerusalem, vol. VII, 《The Nicene and Post-Nicene Fathers》, 6-157；「奧秘,講解」,144-157.)

[ii] “But now the holy day of the Passover is at hand, and ye, beloved in Christ, are to be enlightened by Laver of regeneration. Ye shall therefore again be taught what is requisite, if God so will; with how great devotion and order you must enter in when summoned, for what purpose each of the holy mysteries of Baptism is performed and with what reverence and order you must go from Baptism to the Holy Altar of God, and enjoy its spiritual and heavenly mysteries.” (Ibid, 142)



\newpage

\section{花褪殘紅青杏小}
\label{sec:392}
\textbf{鄧瑞強博士 - 花褪殘紅青杏小}
\newline
\newline
\hyperref[sec:391]{< < < PREV SERMON < < <}
~
\hyperref[sec:toc]{[返目錄]}
~
\hyperref[sec:393]{> > > NEXT SERMON > > >}


以賽亞書 40:1-11
\newline
\begin{longtable}{cl}
\hline
\hline
章節 & 經文 (和合本修訂版)\\
\hline
40:1 & \begin{tabularx}{0.7\textwidth}{X} 你們的神說：「要安慰，安慰我的百姓。 \end{tabularx} \\ \\ \relax
40:2 & \begin{tabularx}{0.7\textwidth}{X} 要對耶路撒冷說安慰的話，向它宣告，它的戰爭已結束，它的罪孽已赦免；它為自己一切的罪，已從耶和華手中加倍受罰。」 \end{tabularx} \\ \\ \relax
40:3 & \begin{tabularx}{0.7\textwidth}{X} 有聲音呼喊著：「要在曠野為耶和華預備道路，在沙漠為我們的神修直大道。 \end{tabularx} \\ \\ \relax
40:4 & \begin{tabularx}{0.7\textwidth}{X} 一切山窪都要填滿，大小山岡都要削平；陡峭的要變為平坦，崎嶇的必成為平原。 \end{tabularx} \\ \\ \relax
40:5 & \begin{tabularx}{0.7\textwidth}{X} 耶和華的榮耀必然顯現，凡有血肉之軀的都一同看見，因為這是耶和華親口說的。」 \end{tabularx} \\ \\ \relax
40:6 & \begin{tabularx}{0.7\textwidth}{X} 有聲音說：「你喊叫吧！」我說：「我喊叫甚麼呢？」凡有血肉之軀的盡都如草，他的一切榮美像野地的花。 \end{tabularx} \\ \\ \relax
40:7 & \begin{tabularx}{0.7\textwidth}{X} 耶和華吹一口氣，草就枯乾，花也凋謝。百姓誠然是草； \end{tabularx} \\ \\ \relax
40:8 & \begin{tabularx}{0.7\textwidth}{X} 草必枯乾，花必凋謝，惟有我們神的話永遠立定。 \end{tabularx} \\ \\ \relax
40:9 & \begin{tabularx}{0.7\textwidth}{X} 報好信息的錫安哪，要登高山；報好信息的耶路撒冷啊，要極力揚聲。揚聲不要懼怕，對猶大的城鎮說：「看哪，你們的神！」 \end{tabularx} \\ \\ \relax
40:10 & \begin{tabularx}{0.7\textwidth}{X} 看哪，主耶和華必以大能臨到，他的膀臂必為他掌權；看哪，他的賞賜在他那裡，他的報應在他面前。 \end{tabularx} \\ \\ \relax
40:11 & \begin{tabularx}{0.7\textwidth}{X} 他要像牧人牧養自己的羊群，用膀臂聚集羔羊，抱在胸懷，慢慢引導那乳養小羊的。 \end{tabularx} \\ \\
[1ex]
\hline
\hline
\end{longtable}
講題:花褪殘紅青杏小　No matter how hopeless, in God we hope

經文:以賽亞書40:1-11

講員:鄧瑞強博士

場合:香港中文大學崇基學院禮拜堂

日期:2017年12月10日

各位弟兄姊妹,早安.願你們平安.

「花褪殘紅青杏小.燕子飛時,綠水人家繞.枝上柳綿吹又少,天涯何處無芳草!」

這是蘇東坡的《蝶戀花》,在絕望中,滲透著希望.

「花褪殘紅」,花已殘,色已褪,凋零而絕望,蘇東坡卻留意到「青杏小」.小小的青杏,生命的種子,卑微而不起眼,卻在凋零中準備展示它的生命.

「燕子飛時,綠水人家繞」,蘇東坡看到生命在飛揚,看到他所關心的人文世界.

「枝上柳綿吹又少,天涯何處無芳草」,枝上的柳絮,越吹越少,生命的美麗像是遠去了.然而,畫面一轉,卻看到種子散落大地,處處展現生機,天涯何處無芳草.

學者估計,這首詞寫於蘇東坡晚年.他六十歲時,被貶去廣東的惠州,這是帝國的邊陲地帶,落後多瘴,九死一生.他一生都不太如意,被一貶再貶,在遙遠他鄉,滿懷失落時,他寫下這首詞,表達絕處中的希望.

耶路撒冷,在主前587或586年陷落,巴比倫大軍消滅了猶大國.聖殿及皇宮被燒,聖殿的財寶被掠奪,城牆被毀,人民被殺.未被殺的,被擄去巴比倫作奴隸.如此弱小的一個國家,衰敗凋零,「花褪殘紅」.看來,他們將會湮沒在歷史的流逝中,看不到「小小的青杏」了.

真的如此嗎?我們看看今日的講道經文:以賽亞書40:1－11.

以賽亞書,分為三個部分.第1至39章,是第一部分,一般而言,講述神對人間罪惡的審判,其結尾部分,預告神對罪惡的猶大國的審判.第40至55章,是第二部分,論述被擄的猶太人將要歸回故鄉.第56至66章,是第三部分,論述猶太人歸回後的新生活.

大家應留意到,今日的講道經文,正處於以賽亞書第一部分及第二部分的交接位置,這是黑夜深處及黎明將至的交界,是無告的絕望與新生的希望的轉折點.這個交接點橫跨的時間,是漫漫長夜.以賽亞書第39章指向主前587或586年猶大的亡國事件.第40章講論的,是主前539或538年的被擄回歸.兩件事件,在時間上有一斷裂,中間約有50年的差距.半個世紀了,時間長嗎?對受苦的人而言,一天已是無盡歲月,何況是半個世紀!

想像一下,一群在異鄉做奴隸的人,遙想故國,望眼欲穿.一年又一年,回家無望.一個又一個,客死異邦.在那古老的信仰裡,因著盼望耶和華的解救所燃點起來的希望之火,在這五十年的歲月裡,也漸漸暗淡了.

就在「花褪殘紅」的盡頭,在絕望深處,先知發出希望的信息,他們言說「小小的青杏」.

這群絕望的人,聽到什麼希望的信息呢?

賽40:1　你們的神說:你們要安慰,安慰我的百姓.

賽40:2　要對耶路撒冷說安慰的話,又向她宣告說,她爭戰的日子已滿了；她的罪孽赦免了；她為自己的一切罪,從耶和華手中加倍受罰.

賽40:3　有人聲喊著說:在曠野預備耶和華的路,在沙漠地修平我們神的道.

賽40:4　一切山窪都要填滿,大小山岡都要削平；高高低低的要改為平坦,崎崎嶇嶇的必成為平原.

賽40:5　耶和華的榮耀必然顯現,凡有血氣的必一同看見,因為這是耶和華親口說的.

賽40:6　有人聲說:你喊叫吧!有一個說:我喊叫什麼呢?說:凡有血氣的盡都如草；他的美容都像野地的花.

賽40:7　草必枯乾,花必凋殘,因為耶和華的氣吹在其上；百姓誠然是草.

賽40:8　草必枯乾,花必凋殘,惟有我們神的話必永遠立定.

賽40:9　報好信息給錫安的啊,你要登高山；報好信息給耶路撒冷的啊,你要極力揚聲.揚聲不要懼怕,對猶大的城邑說:看哪,你們的神!

賽40:10　主耶和華必像大能者臨到；他的膀臂必為他掌權.他的賞賜在他那裡；他的報應在他面前.

賽40:11　他必像牧人牧養自己的羊群,用膀臂聚集羊羔抱在懷中,慢慢引導那乳養小羊的.

這段經文,分別有四個人發言,我們聽聽他們給受苦的人什麼希望.

(1)賽40:1－2:雙倍的痛苦,雙倍的安慰

賽40:1　你們的神說:你們要安慰,安慰我的百姓.

賽40:2　要對耶路撒冷說安慰的話,又向她宣告說,她爭戰的日子已滿了；她的罪孽赦免了；她為自己的一切罪,從耶和華手中加倍受罰.

第1節經文,原文是這樣開頭的:「安慰,安慰……」.

在那苦不堪言的漫漫長夜裡,聽到的,是破空而來的「安慰」,「安慰」.這種感覺,有點似:你困在冰天雪地裡,忽然,有人帶著火把走過來,向著你說,終於找到你了.

「安慰」(nacham)這個字,在希伯來文裡,和「後悔」是同一個字.「安慰」和「後悔」都涉及一種改變,舊的東西不要了,要有新的開始.神定意安慰,神定意要有新的開始.猶大國太弱少,他們自己不可能創造新的開始.新的開始,只能由神而來.「安慰,安慰」,連說兩次,神定意讓痛苦的人得到希望.

第2節:「要對耶路撒冷說安慰的話」,《呂振中譯本》直譯,譯作「要和耶路撒冷談心」.耶路撒冷當時是一荒廢之城,被遺棄,遭蹂躪.經文這裡將耶路撒冷視為一女性,一受蹂躪的女性,神卻珍之重之,與她談心,與她相戀.神不理會世間人對她的偏見與遺忘,他排除萬難,愛他們到底.

「又向她宣告說,她爭戰的日子已滿了；她的罪孽赦免了；她為自己的一切罪,從耶和華手中加倍受罰.」這裡解釋猶大亡國的原因.一國之興,在於敬畏神.一國之亡,在於遠離神.猶太人又亡國又被擄,苦難是雙重的,故此神要「安慰,安慰」她,加倍的重建她.先知宣告,一切苦難過去了.

(2)賽40:3－5:新的出埃及事件

賽40:3　有人聲喊著說:在曠野預備耶和華的路,在沙漠地修平我們神的道.

賽40:4　一切山窪都要填滿,大小山岡都要削平；高高低低的要改為平坦,崎崎嶇嶇的必成為平原.

賽40:5　耶和華的榮耀必然顯現；凡有血氣的必一同看見；因為這是耶和華親口說的.

另一人發聲了.他講到一條路,一條回歸故鄉的路,一條重獲身分的路,一條尋回生命意義的路.犯罪而被罰的人,罪與罰都過去了,面前,是一條康莊大道.

這是一條在曠野,在沙漠中開出來的路.明顯,這是「出埃及」的主題的再現.曠野,沙漠,都是死亡之地,卻要在這死亡之地中開出生命的大道.

新機,是神賜予的；開闢這條大道,卻是我們的責任.真正的希望,是從神而來的；預備自己迎接新生,卻是我們的本分.要努力排除生命的沙石和雜質,「一切山窪都要填滿,大小山岡都要削平；高高低低的要改為平坦,崎崎嶇嶇的必成為平原」.在黑暗中的人,有責任要保持光明的心靈.在痛苦裡的人,要努力保持希望的願景.抑鬱,作為一種病,令抑鬱的人無法擺脫心靈的低沉.絕望,作為一種心靈的病,使絕望的人無法期盼光明的大道.在這裡,宣告希望的使者,呼喚我們盼望光明,盼望自由,盼望新生.明天,不是由我們的絕望去定義的；明天,掌握在神的手中.明天將要出現的光明,是絕望者和痛苦者想像不到的.第5節:「耶和華的榮耀(或說:神聖的光明)必然顯現,凡有血氣的必一同看見,因為這是耶和華親口說的.」明天,你會見到光明的.這不會被你的絕望所阻礙,也不會被世間的局勢所中止,因為,這是掌管明天的耶和華親口說的.

(3)賽40:6－8:生命短暫,惟神的話永存

賽40:6　有人聲說:你喊叫吧!有一個說:我喊叫什麼呢?說:凡有血氣的盡都如草；他的美容都像野地的花.

賽40:7　草必枯乾,花必凋殘,因為耶和華的氣吹在其上；百姓誠然是草.

賽40:8　草必枯乾,花必凋殘,惟有我們神的話必永遠立定.

另一人發聲了.「有人聲說:你喊叫吧!有一個說:我喊叫什麼呢?」

在最悲苦的時刻,人會呼天搶地.「你喊叫吧!」為生命的悲苦呼喊吧!

另一人說:「我喊叫什麼呢?」人生如此,苦不堪言,可以喊叫什麼呢?

《和合本》的翻譯,將這問題之後的經文視為回答者喊叫的內容,加了個「說」字.在我看來,之後的經文,是解釋我無言以對的理由.在人生的大苦裡,我無話可說,因為:「凡有血氣的盡都如草,他的美容都像野地的花」.生命短暫如斯,有什麼好喊叫呢?樂也好,苦也好,片刻便過去了.人在有限中,一切也無能為力,可以追求什麼呢?可以喊叫什麼呢?有限的人,看來只能在苦中默然消逝.

「草必枯乾,花必凋殘,因為耶和華的氣吹在其上；百姓誠然是草.」人生短暫,這是事實.「草必枯乾,花必凋殘,惟有我們神的話必永遠立定.」詩人再說一次,人生的確短暫,這是千真萬確的事實.有限的人,的確可以指著自己的有限,然後說,人生的努力是徒然的,人生的盼望是空洞的,人生的光明是騙人的.然而,經文說,在人的生死流逝背後,有長存的東西,支持著有限的人.萬物一去不返,「惟有我們神的話必永遠立定」.受苦的人若仍有盼望的話,這盼望不能來自我們自身的有限性,而是來自神那創造光明的說話.

面對「花褪殘紅」,蘇東坡卻道「青杏小」.和蘇東坡不一樣,面對「草必枯乾,花必凋殘」,先知卻說,大自然的生死循環本身,並沒有給予我們超越苦難的力量.永存者,惟神的道.

(4)賽40:9－11:看哪,你們的神,他來了!

賽40:9　報好信息給錫安的啊,你要登高山；報好信息給耶路撒冷的啊,你要極力揚聲.揚聲不要懼怕,對猶大的城邑說:看哪,你們的神!

賽40:10　主耶和華必像大能者臨到；他的膀臂必為他掌權.他的賞賜在他那裡；他的報應在他面前.

賽40:11　他必像牧人牧養自己的羊群,用膀臂聚集羊羔抱在懷中,慢慢引導那乳養小羊的.

最後一個人發言,他透過一個傳喜訊者,通報盼望的實現.盼望實現的關鍵,是神的臨在.「看哪,你們的神!」

我們記得,40章第1節是說:「安慰,安慰我的百姓,你們的神說」.現在,看啊,你們的神,他不只是說,他已然臨在,顯出他的光明,成就他的安慰.

我們的神,如何成就他的安慰呢?

先知講出兩方面.一,「主耶和華必像大能者臨到,他的膀臂必為他掌權.他的賞賜在他那裡,他的報應在他面前.」神以大能戰士的身分出現.他將世間一切暴力掃清,將種種混亂變成秩序,主持人間正義.被擄的,得到釋放.痛苦的,得到安慰.受欺壓的,神為你主持正義.不用絕望,神會臨在.

二,神的臨在像牧人,「他必像牧人牧養自己的羊群,用膀臂聚集羊羔抱在懷中,慢慢引導那乳養小羊的.」

神的膀臂不單消滅暴力,這膀臂同時也用來保護羊群.

這日,就是「將臨期」期待的那一日.這日,是神安慰受苦者的日子,是被奴役的人出埃及的日子,是有限的人見到神永恆光明的日子.我們的信仰,期待著這一日.

(5)結語

蘇東坡那首《蝶戀花》,還有下闕.

「牆裡秋千牆外道.牆外行人,牆裡佳人笑.笑漸不聞聲漸悄,多情卻被無情惱.」

蘇東坡雖然豁達,但他的真正想望,仍與他隔著一道牆.他的理想,在牆內,佳人在蕩秋千,在笑.但自己卻在牆外,在圓滿的幸福之外.「笑漸不聞聲漸悄」,在有限的生命裡,漸行漸遠,幸福的想望越來越渺茫了.「多情卻被無情惱」,有夢想的人,被這無夢的世界拒絕了.留下的,是無邊苦惱.我們以為,蘇東坡能在大自然的「青杏小」裡,看到他自己的夢想的實現.現實卻是,他在有限的人生裡,找不到令自己不被無情所惱的永恆者.

先知卻說,有一位神,能主動打破禁錮生命的牆,拒絕絕望,以其永存的說話,在不可能有光明的地方,創造光明.神以其不受限制的創造力,從空無中創造萬有,保證人在漫漫長夜裡,不致有絕望的理由.

希望各位弟兄姊妹,在這「將臨期」,仰望來臨中的基督,找到希望之光.

但願:榮耀歸於聖父,聖子,聖靈.阿們.



\newpage

\section{燕子飛時綠水人家繞}
\label{sec:393}
\textbf{林豪恩先生 - 燕子飛時　綠水人家繞}
\newline
\newline
\hyperref[sec:392]{< < < PREV SERMON < < <}
~
\hyperref[sec:toc]{[返目錄]}
~
\hyperref[sec:394]{> > > NEXT SERMON > > >}


約翰福音 1:6-8,19-28
\newline
\begin{longtable}{cl}
\hline
\hline
章節 & 經文 (和合本修訂版)\\
\hline
1:6 & \begin{tabularx}{0.7\textwidth}{X} 有一個人，是從神那裡差來的，名叫約翰。 \end{tabularx} \\ \\ \relax
1:7 & \begin{tabularx}{0.7\textwidth}{X} 這人來是為了作見證，是為那光作見證，要使眾人藉著他而信。 \end{tabularx} \\ \\ \relax
1:8 & \begin{tabularx}{0.7\textwidth}{X} 他不是那光，而是要為那光作見證。 \end{tabularx} \\ \\ \relax
1:9 & \begin{tabularx}{0.7\textwidth}{X} 那光是真光，來到世上，照亮所有的人。 \end{tabularx} \\ \\ \relax
1:10 & \begin{tabularx}{0.7\textwidth}{X} 他在世界，世界是藉著他造的，世界卻不認識他。 \end{tabularx} \\ \\ \relax
1:11 & \begin{tabularx}{0.7\textwidth}{X} 他來到自己的地方，自己的人並不接納他。 \end{tabularx} \\ \\ \relax
1:12 & \begin{tabularx}{0.7\textwidth}{X} 凡接納他的，就是信他名的人，他就賜他們權柄作神的兒女。 \end{tabularx} \\ \\ \relax
1:13 & \begin{tabularx}{0.7\textwidth}{X} 這些人不是從血生的，不是從情慾生的，也不是從人的意願生的，而是從神生的。 \end{tabularx} \\ \\ \relax
1:14 & \begin{tabularx}{0.7\textwidth}{X} 道成了肉身，住在我們中間，充充滿滿地有恩典有真理，我們也見過他的榮光，正是父獨一兒子的榮光。 \end{tabularx} \\ \\ \relax
1:15 & \begin{tabularx}{0.7\textwidth}{X} 約翰為他作見證，喊著說：「這就是我曾說：『那在我以後來的先於我，因為在我以前，他已經存在。』」 \end{tabularx} \\ \\ \relax
1:16 & \begin{tabularx}{0.7\textwidth}{X} 從他的豐富裡，我們都領受了恩典，而且恩上加恩。 \end{tabularx} \\ \\ \relax
1:17 & \begin{tabularx}{0.7\textwidth}{X} 律法是藉著摩西頒佈的；恩典和真理卻是由耶穌基督來的。 \end{tabularx} \\ \\ \relax
1:18 & \begin{tabularx}{0.7\textwidth}{X} 從來沒有人見過神，只有在父懷裡獨一的兒子將他表明出來。 \end{tabularx} \\ \\ \relax
1:19 & \begin{tabularx}{0.7\textwidth}{X} 這是約翰的見證：猶太人從耶路撒冷差祭司和利未人到約翰那裡去問他：「你是誰？」 \end{tabularx} \\ \\ \relax
1:20 & \begin{tabularx}{0.7\textwidth}{X} 他就承認，並不隱瞞，承認說：「我不是基督。」 \end{tabularx} \\ \\ \relax
1:21 & \begin{tabularx}{0.7\textwidth}{X} 他們又問他：「那麼，你是誰？是以利亞嗎？」他說：「我不是。」「是那位先知嗎？」他回答：「不是。」 \end{tabularx} \\ \\ \relax
1:22 & \begin{tabularx}{0.7\textwidth}{X} 於是他們對他說：「你到底是誰，好讓我們回覆差我們來的人。你說，你自己是誰？」 \end{tabularx} \\ \\ \relax
1:23 & \begin{tabularx}{0.7\textwidth}{X} 他說：「我就是那在曠野呼喊的聲音：修直主的道。」正如以賽亞先知所說的。 \end{tabularx} \\ \\ \relax
1:24 & \begin{tabularx}{0.7\textwidth}{X} 那些人是法利賽人差來的。 \end{tabularx} \\ \\ \relax
1:25 & \begin{tabularx}{0.7\textwidth}{X} 他們就問他：「你既不是基督，不是以利亞，也不是那位先知，那麼，你為甚麼施洗呢？」 \end{tabularx} \\ \\ \relax
1:26 & \begin{tabularx}{0.7\textwidth}{X} 約翰回答：「我是用水施洗，但有一位站在你們中間，是你們不認識的， \end{tabularx} \\ \\ \relax
1:27 & \begin{tabularx}{0.7\textwidth}{X} 就是那在我以後來的，我給他解鞋帶也不配。」 \end{tabularx} \\ \\ \relax
1:28 & \begin{tabularx}{0.7\textwidth}{X} 這些事發生在約旦河東邊的伯大尼，約翰施洗的地方。 \end{tabularx} \\ \\
[1ex]
\hline
\hline
\end{longtable}
講題:燕子飛時　綠水人家繞　A Witness to Testify to the Light

經文:約翰福音1:6-8,19-28

講員:林豪恩先生

場合:香港中文大學崇基學院禮拜堂

日期:2017年12月17日

一　引言

今天的講道題目「燕子飛時,綠水人家繞」,明顯是拾鄧瑞強博士的牙慧.相信在座不少人都好像我一樣,很喜歡聽鄧博士講道,能夠從中拾到一點碎屑,已經感到飽足.因此,跟隨鄧博士,最多只是暴露自己的不足,但一定不會跟錯,也不會遇上「跟車太貼」的風險.

讓我們想像一下,「燕子飛時,綠水人家繞」是怎樣的圖畫呢?燕子在天空飛舞,清澈的河流圍繞著村落人家,這種自然景象告訴人們,冬末春初乍暖還寒的日子漸漸過去,溫暖的初夏正在來臨.燕子不是溫暖的初夏,綠水也不是溫暖的初夏,但燕子在寒冬之後再次飛舞,綠水在乾涸的水道再次流動,是初夏來臨的「見證」,讓人們知道日子會越來越和暖.現在是將臨期,提示人們耶穌基督的來臨,人們又如何得知耶穌基督來臨呢?為耶穌基督的來臨而作的見證是怎樣的呢?讓我們再讀今天的福音經課,看看施洗約翰的例子能帶給我們甚麼啟發.

二　經文

1:6有一個人,是從上帝那裡差來的,名叫約翰.

1:7這人來,為要作見證,就是為光作見證,叫眾人因他可以信.

1:8他不是那光,乃是要為光作見證.

1:19約翰所作的見證記在下面:猶太人從耶路撒冷差祭司和利未人到約翰那裡,問他說:「你是誰?」

1:20他就明說,並不隱瞞,明說:「我不是基督.」

1:21他們又問他說:「這樣,你是誰呢?是以利亞嗎?」他說:「我不是.」「是那先知嗎?」他回答說:「不是.」

1:22於是他們說:「你到底是誰,叫我們好回覆差我們來的人.你自己說,你是誰?」

1:23他說:「我就是那在曠野有人聲喊著說:『修直主的道路』,正如先知以賽亞所說的.」

1:24那些人是法利賽人差來的；

1:25他們就問他說:「你既不是基督,不是以利亞,也不是那先知,為甚麼施洗呢?」

1:26約翰回答說:「我是用水施洗,但有一位站在你們中間,是你們不認識的,

1:27就是那在我以後來的,我給他解鞋帶也不配.」

1:28這是在約旦河外伯大尼,約翰施洗的地方作的見證.

三　釋經與應用

這段經文的主角明顯是施洗約翰,但他卻說自己不是主角；他在舞台上光芒四射,但他澄清自己不是那光；他站在歷史的轉捩點,但堅持自己不是那位.約翰的表現,在那種「爭做主角,爭曝光,有位必上」的世界中,令人感到大惑不解.事實上,當時的猶太人也感到困惑,經文指出「猶太人從耶路撒冷差祭司和利未人到約翰那裡,問他說:「你是誰?」」「你是誰?」這個問題對於耶穌基督的見證者來說太重要了,這是第一個問題,也是一生都要問的問題.「你是誰?」這個問題,有時並不容易正面回答.有人問我:「你是牧師嗎?」我說:「不是.」又問:「你是宣教師嗎?」我說:「我沒有接受過宣教師的按立.」又問:「你是傳道人嗎?」我回答:「在大學裡沒有傳道人這個職位.」我知道如此的回答會令人感到迷惑,不過,事實如此.經文記載當施洗約翰被祭司和利未人問:「你是誰?」的時候,他的回答及作者的敘述也是有點令習慣以身份來定義人的人感到迷惑.

1．「不是甚麼」相對於「是甚麼」

1:19約翰所作的見證記在下面:猶太人從耶路撒冷差祭司和利未人到約翰那裡,問他說:「你是誰?」

1:20他就明說,並不隱瞞,明說:「我不是基督.」

1:21他們又問他說:「這樣,你是誰呢?是以利亞嗎?」他說:「我不是.」「是那先知嗎?」他回答說:「不是.」

約翰作為一位見證者,指出自己「不是甚麼」可能比介紹自己「是甚麼」更有作見證的意義.因為,承認自己不是那光,才不遮蓋那光,不遮蓋那光,才能讓人看見真光.承認自己不是基督,才能引導人尋得基督.承認自己不是以利亞,也不是那先知,才不會令人誤會自己是「復興古國盛世」的歷史繼承人,才不會誤導人沉醉於追求「民族復興」的夢,反而令人意識到新時代的來臨,以開放的心懷去迎接「眼未曾看見,耳未曾聽過」的新世界.

然而,人們總是希望自己是「somebody」而不願意接受自己是「nobody」,教會和信徒也未必能夠免疫.在古時,有一男子名叫約伯,他遭逢厄運.他有三位正義朋友來安慰他,結果卻為了維護對上帝的信念,又或者證明自己是「somebody」,在人生重大議題上必須發言,而與約伯展開了好幾個回合的唇槍舌戰.最後,上帝發聲了:「誰用無知的言語使我的旨意暗昧不明?」歷世歷代,會否都有不少以為自己在維護上帝的見證者,反而令人看不見上帝呢?

教會,本來是在地上作為耶穌基督的見證,會否也變成令人只見教會而不見耶穌呢?牧者及信徒,本來是指引人往耶穌基督那裡去,會否也變成只見自己而不見耶穌呢?

如何讓人看見基督呢?最底線的是承認自己不是基督.相反來說,最能夠取消基督的,就是明示或暗示自己就是基督.有人會如此狂妄嗎?事實上是有的,而且也能夠吸引不少人.有人做過資料搜集,展示近百年來有1100人自稱基督.較近期的例子有這個:

《時代雜誌》在2001年11月刊登的一篇名為《耶穌回來了,她是中國人》的文章.一位河南鄭州的鄧姓女子,她自稱為二次降臨的女基督,名叫「閃電」.

上述例子是明言,並不隱藏.下列是暗示,需要自行領會.

南華早報十一月中報導,中國江西一地方政府的《微信》公眾號發佈的消息顯示,政府推動「信教向信黨轉化」,在基督教信徒眾多的江西上饒市餘干縣等地,清除基督徒家中的宗教字畫,取而代之的是張貼國家主席習近平的畫像.

以下不是明言,也不是暗示,但也可以產生令人不想看見上帝的效果.一位在友校任教的老師寫下了他的感慨:

「很多人自稱信徒,是為了多拿一片葉子來遮掩他們的邪惡本質.可能正因如此,林鄭話自己是天主教徒,梁美芬話自己是基督徒,周浩鼎又話自己係基督徒,就連面對鏡頭粗口爛舌,撩事鬥非,喊打鹹殺的何君堯選舉時都報稱係基督徒……講俾人聽佢地信乜ye 宗教,就算完全唔返教會,冇乜幾何入廟拜神,都可以把宗教信仰變成亞當夏娃身上那片葉,可以遮羞蔽醜,一個人的空疏無知,心裡的歪念及醜惡動機,以為都可以遮蓋起來……幾年前,有朋友講,好擔心將來上到天堂,會見番晒呢一班咁ge 人.」(鍾劍華:這樣的信徒信甚麼?)

施洗約翰作為耶穌基督的見證,他的自處是:「我不是甚麼」,而不是宣揚自己「我是甚麼」.當日約翰以不掩蓋基督,來見證基督,會否值得今日的教會和信徒參考呢?

2．見證之處:建制力量,圍爐取暖,曠野之聲

說說回頭,施洗約翰並非完全不回答「你是誰?」這個問題,只是他沒有把自己套入當時社會上任何一個身份來回答,而是以「我做甚麼」的行動描述回答:

1:23他說:「我就是那在曠野有人聲喊著說:『修直主的道路』,正如先知以賽亞所說的.」

身份與行動可以是兩回事.約翰為耶穌基督作見證,希望人們留意他的行動,而不執著於審查他的身份.此時此地,人們卻非常重視身份.例如,在信徒群體中,標示自己是社會成功人士；在政治活動中,宣揚自己是信徒.

其實,雪亮的觀眾眼睛是不被身份瞞騙的,是信徒也好,是牧師也好,是主教也好,真正產生見證果效的是行事為人,而不是身份.約翰是誰呢?這不重要.重要的是他所作的:『修直主的道路』.約翰如何「修直主的道路」,幫助人預備接受耶穌基督的來臨呢?這個問題同時與約翰選擇在甚麼場景作見證相關,我們一併討論.

在甚麼地方迎接耶穌基督來臨呢?其實約翰有三項選擇:

第一,在當時的宗教建制之中

為甚麼約翰有這項選擇呢?大家還記得約翰的父母是誰嗎?路加福音如此記載:

1:5在希律作猶太王的時候,亞比雅班裡有一個祭司,名叫撒迦利亞；他妻子是亞倫的後代,名叫伊利莎白.

1:8撒迦利亞按班次在上帝面前執行祭司的職務,

1:9照祭司的規矩抽籤,進到主的殿裡燒香.

原來,約翰的父親撒迦利亞是祭司,屬於亞比雅班次；而約翰的母親伊利莎白,是第一位大祭司亞倫的後代.因此,約翰其實根正苗紅,進入當時的宗教建制供職,是順理成章的選擇.如果,約翰留在宗教建制之內去「修直主的道路」,會否令整個祭司及文士系統接納耶穌就是他們期盼的彌賽亞呢?會否因此而帶領當時整個猶太人社會迎接耶穌的來臨呢?這種想法,正如某些人盼望,如果國家領導人變成了基督徒,就能夠帶領全國人民信耶穌一樣.聖經沒有解釋約翰為甚麼不選擇在聖殿做祭司,但稍後我們討論他的行動時,或者能夠找到一點端倪.

第二,參加昆蘭社團,清心等候彌賽亞

當時的猶太社會主要分成三種教派,撒都該人(Sadducees),法利賽人(Pharisees)和隱士派人(Essenes),又稱愛色尼人.大家對撒都該和法利賽人比較熟識,因為聖經提及他們.至於隱士派猶太人,由於對建制宗教不滿而選擇住在離開城市的郊外,住在一起,凡物共有,相傳他們守獨身,不蓄奴,親操農務,穿白衣,強調生活上的聖潔以等候彌賽亞來臨.

施洗約翰的形象雖然獨持,他也盼望彌賽亞的來臨,但卻沒有加入這個遺世獨立的群體.或者,他不參加這種獨善其身的群體,就是因為他要為彌賽亞的來臨作見證.

第三,作為曠野之聲 (a voice in the wilderness)

那麼,約翰要作的見證是甚麼呢?『修直主的道路』又是甚麼意思呢?讓我們看看福音書的記載:

馬太福音

3:2「你們要悔改!因為天國近了.」

馬可福音

1:4照這話,施洗約翰來到曠野,宣講悔改的洗禮,使罪得赦.

1:5猶太全地和全耶路撒冷的人都出去,到約翰那裡,承認他們的罪,在約旦河裡受他的洗.

路加福音

3:3他就走遍約旦河一帶地方,宣講悔改的洗禮,使罪得赦.

3:7約翰對那出來要受他洗的眾人說:「毒蛇的孽種啊,誰指示你們逃避那將要來的憤怒呢?

3:8你們要結出果子來,和悔改的心相稱.不要自己心裡說:『我們有亞伯拉罕為祖宗.』我告訴你們,上帝能從這些石頭中給亞伯拉罕興起子孫來.

3:9現在斧子已經放在樹根上,凡不結好果子的樹就砍下來,丟在火裡.」

3:10眾人問他:「這樣,我們該做甚麼呢?」

3:11約翰回答:「有兩件衣裳的,就分給那沒有的；有食物的,也該這樣做.」

3:12也有稅吏來要受洗,對他說:「老師,我們該做甚麼呢?」

3:13約翰對他們說:「除了規定的數目,不要多收.」

3:14也有士兵問他說:「我們該做甚麼呢?」約翰說:「不要勒索任何人,也不要敲詐人；自己有糧餉就該知足.」

原來,施洗約翰修直主的道路的行動,就是呼籲人們悔改,迎接天國的來臨,而且他呼籲的悔改是非常針對性及行動性的.聖經形象化地描述約翰站在河中施洗,用本地話來說,是一個「落水」的形象.約翰這種站在流水中「落水」的見證,相對於昆蘭社團隱世修行式的生活,哪一種更能幫助人們迎見天國呢?約翰站在曠野,相對於站在祭司制度中,哪一處更能夠指責包括建制內的罪惡,呼籲包括建制內的人們悔改呢?

其實,約翰選擇「落水」,選擇發出「曠野之聲」是充滿危機的.相對來說,若選擇作為祭司,有整個宗教建制的保護；選擇加入昆蘭社團,遠離世事,也相對安全.然而,約翰站在河流中作見證,就難免亳無保護地受到流水的衝擊,為此,他要付出沉重的付價:(馬太福音)

14:3原來,希律為他兄弟腓力的妻子希羅底的緣故,把約翰抓住綁了,關進監獄,

14:4因為約翰曾對他說:「你佔有這婦人是不合法的.」

14:5希律就想要殺他,可是怕民眾,因為他們認為約翰是先知.

14:6到了希律的生日,希羅底的女兒在眾人面前跳舞,使希律歡喜,

14:7於是希律發誓許諾隨她所求的給她.

14:8女兒被母親指使,就說:「請把施洗約翰的頭放在盤子裡,拿來給我.」

14:9王就憂愁,然而因他所發的誓,又因同席的人,就下令給她；

14:10於是打發人去,在監獄裡斬了約翰,

14:11把頭放在盤子裡,拿來給那女孩,她拿去給她母親.

代價如此具大的見證值得嗎?那要看由誰來評價.耶穌曾如此評論施洗約翰:

11:7他們一走,耶穌就對眾人談到約翰,說:「你們從前到曠野去,是要看甚麼呢?看風吹動的蘆葦嗎?

11:8你們出去到底是要看甚麼?看穿細軟衣服的人嗎?那穿細軟衣服的人是在王宮裡.

11:9你們出去究竟是要看甚麼?是先知嗎?是的,我告訴你們,他比先知大多了.

11:10這個人就是經上所說的:『看哪,我要差遣我的使者在你面前,他要在你前面為你預備道路.』

11:11我實在告訴你們,凡女子所生的,沒有一個比施洗約翰大.

四　總結

「燕子飛時,綠水人家繞」,燕子不是太陽,但燕子在寒冬後再次飛翔,帶給人溫暖的盼望；流水也不是太陽,但流水在乾涸的河道再次湧現,帶給人新生的盼望.事實上,燕子總有離去的日子,流水也總會有乾涸的時候,重要的是在他們出現的日子,曾經帶給人們溫暖及新生的盼望.約翰福音三章30節記下了施洗約翰的名言,值得所有為耶穌基督作見證的教會及信徒參考,他說:「祂必興旺,我必衰微.」作為見證者的教會及信徒,可能以為我們的興旺就是見證,以及我們的衰微就是失見證.我們的成敗得失,上主一切都看在眼內；我們的心思意念,上主也無不知曉.因此,就算花褪殘紅,也無需傷感,有否記得另有詩說:「落紅不是無情物,化作春泥更護花」呢?見證者的心,上主是知道的；見證者的情,是潤澤人間的.



\newpage

\section{可道與不可道}
\label{sec:394}
\textbf{許立中先生 - 可道與不可道}
\newline
\newline
\hyperref[sec:393]{< < < PREV SERMON < < <}
~
\hyperref[sec:toc]{[返目錄]}
~
\hyperref[sec:395]{> > > NEXT SERMON > > >}


提摩太前書 3:16
\newline
\begin{longtable}{cl}
\hline
\hline
章節 & 經文 (和合本修訂版)\\
\hline
3:16 & \begin{tabularx}{0.7\textwidth}{X} 敬虔的奧祕是公認為偉大的：神在肉身顯現，被聖靈稱義，被天使看見，被傳於外邦，被世人信服，被接在榮耀裡。 \end{tabularx} \\ \\
[1ex]
\hline
\hline
\end{longtable}
講題:可道與不可道　When the Metaphysical became Tangible

經文:提摩太前書3章16節

講員:許立中先生

場合:香港中文大學崇基學院禮拜堂

日期:2017年12月24日

「大哉敬虔的奧秘,無人不以為然；就是上帝在肉身顯現,被聖靈稱義,被天使看見,被傳於外邦,被世人信服,被接在榮耀裡.」(提前3:16)

傳統上,聖誕是一個比較感性的節日.就是對非信徒來說,聖誕的詩歌和音樂,燈飾和氣氛,很容易就讓人產生愉快和cosy的感覺,希望跟家人或朋友一起分享生命的美好.

但正如聖誕對一般人來說就只不過是一種感覺,一種氣氛,聖誕過後,自然就生活依舊,工作如常.年紀還小,或許還會期待下個聖誕的來臨,就像那緊接下來的農曆新年一樣；年事稍長,聖誕就只不過是一個長一點的假期而已.

事實上也不僅是聖誕,現代社會似乎將所有傳統節日壓平: Happy Holiday!無論是冬至,聖誕,新年,中秋,幾乎就連清明我們的祝賀語都是XX快樂!除了農曆新年我們終於可以「恭喜發財」,其他節日都變得沒有內容,沒有性格!

當然,對基督徒來說,聖誕除了是一個「大節」,它更是記念救主的降生,就是基督教之所以是基督教的關鍵事件.我接著下來要分享的,在內容上其實並沒有甚麼「聖誕氣氛」,但在主題上,卻希望指出聖誕確實不僅是一種氣氛.

上個月崇基週會委員會邀請了演藝界的黃秋生先生到學生週會分享.雖然主持的老師並沒有要求或引導黃先生分享他的信仰觀點,但可能是禮堂的氣氛,週會的程序,包括讀經和詩班的獻唱,令他很自然地分享了一點他自己對信仰的看法.他說他並沒有明確的宗教信仰,勉強可以說有點泛神論的傾向；但聽他具體的分享,似乎更接近我們所謂「舉頭三尺有神明」那種感覺,而非像泛神論那樣直接將宇宙等同於上帝本身.他說倘若一定要歸邊,則他比較接近道家思想.

印象中黃先生比較受歡迎,或許也跟他對生活,做人比較有反省有關,以致相對於其他演藝工作者較為「貼近生活」和樂於表達.對宇宙人生有比較宏觀的思考,多少影響一個人對工作,對生活,對做人的基本態度.正如舊約《箴言書》指出:「一個人怎樣思量,他就是怎樣的人.」(23:7)

2011年十二月,龍應台被邀請在港大醫學院畢業典禮上演講,她的演講一貫地充滿洞見和個人反省.她說:「一個人所抵抗的以及所堅持的,匯成一個總體,就叫做『信仰』.但是信仰,依靠的不是隆重的大聲宣告；信仰深藏在日常生活的細節裡,信仰流露在舉手投足之間最尋常最微小的決定裡.」

在那裡,龍應台當然不是具體討論到宗教信仰,而是一個人內心的信念.所以問題反而是,我們的信仰有沒有機會沉澱到我們的內心,融入我們的生活?而這個信念,「依靠的不是隆重的大聲宣告；信仰深藏在日常生活的細節裡,信仰流露在舉手投足之間最尋常最微小的決定裡.」

我甚至認為,一個人是一個怎樣的人,更多是決定於他想些甚麼,而不是他信些甚麼!這當然跟信徒往往在信仰和思想上分家,近乎精神分裂地生活在兩個毫不相干的世界有關.我的意思是,一個人可以信得馬虎求其,或者信得嚴肅認真但跟現實生活脫節；但他真正的想法,卻是具體地孕育出他每個行為和每個抉擇.因此一個人想些甚麼,確實比他信些甚麼更重要!而在華人知識分子當中,根據我有限的觀察,恐怕是接近老莊多於奧古斯丁或多瑪阿奎拿!

老莊思想給人的印象是比較豁達,自由,不拘一格.《道德經》第25章:「人法地,地法天,天法道,道法自然.」即天,地,人都互相依循道的法則,而道則表彰於自然的規律和現象.但其實在25章的開頭,老子對道有較仔細的描述:「有物渾成,先天地生,寂兮寥兮,獨立不改,周行而不怠,可以為天下母,吾不知其名,字之曰道,強為之名曰大.」

用白話來說:「冥冥中有些甚麼,先於我們所認識的宇宙而存在.在孤寂的渾沌中,它獨立自主,不受任何外在因素所影響；循環不息,而沒有任何疲態,或許它就是眾生之母.我不知該叫它甚麼,姑且稱之為道；找不到更適合的形容,或可叫它做大.」

這樣的描述,很容易令我們聯想到新約聖經《約翰福音》第一章:「太初有道,道與上帝同在,道就是上帝.這道太初與上帝同在,萬物是藉著他造的,凡被造的,沒有一樣不是藉著他造的.生命在他裡頭,這生命就是人的光.」

當然,無論是老子對道的體會,抑或希臘哲學λογος的概念,都是指向一種非人化的宇宙規律或法則.是我們熟悉文天祥的《正氣歌》,進一步將這客觀的宇宙規律賦予道德的內涵:「天地有正氣,雜然賦流形,下則為河嶽,上則為日星,於人曰浩然,沛乎塞蒼冥.」

換句話說,文天祥認為這道不僅是物質世界的規律法則,更加表現為天地間的「正氣」,人世間的道德法則,為道注入一點人性化的味道.是《約翰福音》大膽地將道跟上帝拉上關係,不單只是「道與上帝同在」,更進一步宣稱「道就是上帝」!這樣,《約翰福音》實質上是指出道並不是一種獨立,抽象的存在；道之所以為道,是因為上帝是上帝.

弗4:4b－6:「…正如你們蒙召,同有一個指望,一主,一信,一洗,一上帝,就是眾人的父,超乎眾人之上,貫乎眾人之中,也住在眾人之內.」

而正因為道這超越性,涵蓋性和內蘊性,以致老子認為真正的道是不可表達的.因為當你要表達他的時候,你其實必須將他放在你既有的理念框架之中,而當你嘗試這樣做的時候,你是無可避免地扭曲他,因為沒有任何理念框架,能夠盛載令這理念框架本身成為可能的原因和條件.舉一個比較形象化的例子,一個訓練有素的舉重選手,往往能夠輕易舉起超過自己身體的重量,但他又可以用甚麼作為它的著力點,讓他能夠舉起自己?

使徒保羅在《羅馬書》第九章用另一個類比表達同樣的意思,就是yiu 匠跟瓷器的關係:一件瓷器不能挑戰造它的yiu 匠:「你為甚麼這樣造我?」因為瓷器之所以成為瓷器,完全因為yiu 匠這樣將它塑造出來；那麼它又依據甚麼去挑戰yiu 匠的做法?故此「道可道非常道」.真正的道確實是不能表達的.而這正是聖誕的悖論(paradox)!

小時候看報紙的四格漫畫,久不久就會看到以「無題」為題的漫畫.無題,有時是因為漫畫的信息較為含蓄,只能意會難以言傳,於是索性讓讀者自已揣摸漫畫要表達的主旨.

因此雖說是「無題」,其實還不是沒有主題,只不過意思比較難用三言兩語表達得清楚.正如《道德經》第一句就指出「道可道非常道」,最後卻用了八十一章的篇幅,去表達那個不可道之道.這也是一種悖論.

同樣,人生在世其實也不能夠沒有主題,雖然那主題往往是隱晦和很難說得清楚.有了主題作為參照的框架,我們才知道如何去自我定位,作出恰當的回應.有了主題,我們才知道該如何抉擇,如何安排生活的優次.

昔日上帝將「野地各樣走獸,和空中各樣飛鳥,都帶到那人面前看他叫甚麼；那人怎樣叫各樣的活物,就是他的名字.」定名,可以說是人理解世界,探究生命的最原始和基本的嘗試,也是一切知識學問的基本任務.

《創世記》第一章:「起初上帝創造天地.地是空虛混沌,淵面黑暗…」老子在《道德經》第一章:「無名,天地之始；有名,萬物之母.」定名,可以說是所有意義問題的核心.無名,是天地的起始；有名,才產生了歷史.有了名字,有了定義,事件才串連出一個有意義的故事.

幾年前學書法的時候,曾臨摹唐代詩人王昌齡的《出塞》,感覺相當超時空:「秦時明月漢時關,萬里長征人未還,但使龍城飛將在,不教胡馬渡陰山」.叫我感到震撼的是第一句:漢代邊關守將,抬頭望見的其實仍是秦朝的那個月亮!無論是在秦朝抑或漢代,甚至是在今日,我們抬頭望見的,還是那個月亮,那片天空.而歷史的事件,跟這片空間彷彿毫無關係.我的意思是,今晚的夜空跟1937年的南京沒有甚麼分別.但生活在不同的時空,面對的命運卻是差天共地.

生命的虛幻性,乃在於它的randomness,它的偶然性,它的「無題」 ─ 就是那種甚麼都好像是,卻又甚麼都好像不是的曖昧性.生命中很多重要的抉擇,最初其實都是一些很偶然的事件.譬如我們的事業,我們的配偶.昨天參加一個婚禮,新郎在致詞的時候多謝一位朋友,因為如果不是那次在停車場碰到他,他就不會認識新娘；不認識新娘,自然就不會跟她做朋友；沒有跟她做朋友,也就不會跟她拍拖；沒有跟她拍拖,當然就不會跟她結婚.

又如果我們不是在那次公開試失手,就不會退而求其次選擇主修另一科；沒有選擇那一科,畢業後也就不會入那一行；沒有入那一行,自然就沒有今日在那行業的成就,或者根本不會在這個行業度過大半生.

當然,我們每一個選擇都是認真和嚴肅的,但這些那麼認真的決定,追溯到底卻可能只是一個偶然.如果是這樣,人生會不會就只是一連串漫無目的的機緣巧合,而所有的抉擇和努力都變得毫無意義?

是那太初原有的生命之道,為這個看似「無題」的世界,重新賦予主題:「道成了肉身,住在我們中間,充充滿滿的有恩典有真理.」有恩典,因為並不在乎我們如何努力去為虛幻的生命強加意義；有真理,是因為祂具體地彰顯出上帝創造的原意.故此,基督徒的任務,其實是將意義,主題重新灌注入這個被遺忘了主題的世界.

是在這背景底下,我們去理解聖誕的意義:就是那位無限的成為有限,那「不可道之道」成為可道的「非常道」:甚至可以說那位絕對的妥協他的絕對性,進入相對的世界.這對猶太人來說是褻瀆,對希臘人來說是contradiction in terms.

猶太教和基督教都屬於「一神信仰」,相信一位創天造地的上帝.它們之間的分別,在於基督信仰相信這位造物主曾經以人的身分來到這個世界,而猶太教則認為永恆的上主不可能以人的身分出現.畢竟在無限與有限,永恆和短暫,絕對與相對,上帝和人之間,是不可能有任何拖泥帶水的灰色地帶.正如一些神學家說,上帝是The Totally Other.

法利賽人出身的保羅,當然明白和肯定這個道理.他折衝的解釋,是一方面肯定上帝的獨一性,卻指出這位獨一的上帝曾經有過不平凡的舉動:「只有一位上帝,在上帝和人中間只有一位中保,乃是降世為人的基督耶穌；他捨自己作萬人的贖價,到了時候,這事必證明出來.」(提前2:5-6)是那位成為肉身的道,不可思議地成為人的一分子,以致可以將人帶到那位獨一不死,住在人不能靠近的光中的上帝,重新檢視自己,檢視人生.

道成肉身,或者說那位無限者成為有限,就意味著上帝必須在某個程度上放棄他的絕對性,以致他跟有限的人有溝通的語言和可能.「你們當以基督的心為心.他本有上帝的形象,不以自己與上帝同等為理所當然的,反倒虛己,取了奴僕的形象,成為人的樣式；既有人的樣子,就自己卑微,存心順服,以至於死,且死在十字架上.」(腓2:5-8)

因此,「上帝在肉身顯現」,並不是君臨天下的排場.他甚至不像釋迦牟尼那樣生為王子,然後拋棄榮華富貴,印證他所覺悟的道.更令人難堪的是,他甚至比我們更平凡:「他無佳形美容,我們看見他的時候,也無美貌使我們羨慕他.」(賽53:2)作為成為肉身的真理,他並沒有將他的絕對性強加於人身上,因為他知道,真理在複雜和相對的人間,往往表現為不同的面相.

為此,他甚至似乎不大必要地經歷了作為人的童年.福音書記載耶穌的童年不多,但也用了不少篇幅講述他的降生,甚至詳細列出他的家譜,為的就是要確立他並不是一般人想像的觀音托世或菩薩顯靈,而是切切實實地成為一個人.

耶穌是實實在在成為一個無助的嬰孩,以致馬利亞和約瑟必須帶他遠赴他方,以逃避希律王的追殺.不錯,他是降卑到一個地步甚至無法保護自己!

之後,福音書就只是簡單記載「孩子漸漸長大,強健起來,充滿智慧,又有上帝的恩在他身上.」(路2:40)直到他十二歲的時候,馬利亞和約瑟帶他上耶路撒冷守逾越節.節期完畢,他們準備返家鄉的時候,馬利亞和約瑟以為耶穌跟同鄉小孩混在一起,所以沒有理會.「走了一天的路程」,才發現不見耶穌,於是唯有截返耶路撒冷,恐怕耶穌是因為人太多而走失了.

回到耶路撒冷,又過了三天,他們驚訝地發現耶穌竟然還留在聖殿,跟文士和教師討論經卷的問題.這是一幅令人感到錯愕的圖畫:一個沒有受過經學訓練的孩子,竟然跟一班除了聖典之外就幾乎沒有其他興趣的學究,討論經卷的問題!

想像一個十二歲的小朋友流落京城四,五天,而所有遇到他成年人,並不關心替他找回家人,而是興致勃勃地跟他討論信仰問題,這是不是有點奇怪?

更值得注意的是,那位本來「不可道」的道,並沒有以權威的姿態去跟那班學者專家對質.聖經只是記載他「一面聽,一面問」,並沒有嘗試解釋或回應任何問題；奇怪的是,「凡聽見他的,都希奇他的聰明,和他的應對.」

一個十二歲的孩子,有十二歲孩子的表達方式.而那位成為肉身的永恆的道,並沒有超越人性,以神童的姿態強勢地自我呈現.他反而完全安於一切作為人的定限,恰如其分地表達自己.畢竟所有的孩子都喜歡問問題!而當他這樣做的時候,真理反而是毫不刻意地呈現為聽者心中的希奇和信服.

我記起《以賽亞書》42:「看哪,我的僕人─我所扶持所揀選,心裡所喜悅的!我已將我的靈賜給他；他必將公理傳給外邦.他不喧嚷,不揚聲,也不使街上聽見他的聲音.…他憑真實將公理傳開.」在少年的耶穌身上,我們認識到真理是不用大聲疾呼,甚至無須急於提供答案.

在這裡,我們亦觀察到一個相當一致的應對模式,就是成為肉身的道,並不是因為詞窮,或不知道如何應對,而用問題作為應付問題的援衝.而是他知道,在一個複雜和相對的處境,任何問題背後都可能有不同的動機,並且可以衍生出無數後續的問題.

因此在大多數的情況下,耶穌都拒絕為一個簡單的問題提供一個直接的答案.

譬如,有人認為耶穌德高望重,希望他能夠替他和他的兄長分家產.耶穌並沒有沾沾自喜,反說:「你這個人!誰立我作你們斷事的官,給你們分家業呢?」(路12:13-14)然後若有所指地對其他人說:「你們要謹慎自守,免去一切的貪心,因為人的生命不在乎家道豐富.」

當然分家產只是一些雞毛蒜皮的私人問題,不值得我們過度引伸,卻提醒我們很多時那些來請教意見的人,其實並不是要得到一個公正的答案,而是希望確認自己的想法.一旦得不到自己想要的意見,他自然會覺得你不公道.可以說,最公正的公正行或「調解員」,都無法令爭論雙方完全滿意.就算是法庭的頒令(court order),其實亦只能夠終止爭論,而不能說服任何一方.這類問題,根本不可能有一個令所有人信服的絕對答案.

另一次,有人帶一個捉姦在床的女人來見耶穌,告訴他按著摩西的律法,應該將她用石頭打死,然後問他應該如何處置?他們是將一件已經有明確案例的道德判斷,擺在耶穌面前,要求一個答案,一個說法.而無論他們有何居心,他們想要的,明顯不是一個絕對的答案,真正的道理.

又再有一次,有人拿著一個羅馬的錢幣問耶穌:作為上帝的子民,納稅給凱撒應該不應該?一個簡單的問題,要求一個直接的答案.耶穌拒絕提供一個直接的答案,反問他們,錢幣上的像是誰?他們說,是凱撒.耶穌這才回答他們:凱撒的物當歸給凱撒,上帝的物當歸給上帝.

在一個複雜和相對的處境,任何問題背後都可能有不同的動機,並且可以衍生出無數後續的問題.耶穌當代的人,往往就是拿著這些相對,模稜兩可或者點答都錯的問題,要求耶穌提供一個絕對的答案.那麼他們真正想要的,又是不是「真理」或事實的真相?

那位成為肉身的道給我們的啟示,就是真理往往透過問題,而不是透過答案而呈現.絕對的真理,只能夠existentially存在於問者的心中,透過良知的自我質問,而呈現出來.因此,對於他們提出的道德殘局,耶穌永遠拒絕提供一個客觀地絕對的答案.或者不是沒有,而是人在罪惡和自義無法接受.

這並不是將真理縮約到一個純粹主觀的判斷,而是正如祈克果曾經指出:「發現那些所謂的客觀真理,搞通所有的哲學系統,有需要的話指出每一個體系的內部矛盾又有甚麼用?倘若我能夠建構一套社會的理論,將所有的細節整合成一個完滿的體系…卻並不是我實在生活於其中的世界,那又有甚麼益處呢?倘若我能夠解釋基督教的道理,卻對我自己和我的生命沒有深一層的意義,那又有甚麼意思?倘若真理冰冷和赤裸地站在我面前,完全不在乎我是否認識她,叫我感到不寒而慄而不是恭謹信靠,那又於我何益?」

當然我們最為熟悉的是這一段:「我真正缺乏的,是不肯定自己應該做些甚麼,而不是自己應該知道些甚麼…重要的是要明白我自己,找出上帝真正要我做的,究竟是甚麼；重要的是要找出一個於我真實的真理,找出一個我可以為之而生並為之而死的信念.」

在彼拉多審問耶穌的那個晚上,耶穌明確指出,「我的國不屬這世界；我的國若屬這世界,我的臣僕必要爭戰,使我不至於被交給猶太人.只是我的國不屬這世界.」彼拉多就問他說:「這樣,你是王嗎?」耶穌回答說:「你說我是王.我為此而生,也為此來到世間,特意為真理作見證.凡屬真理的人,就聽我的話.」彼拉多說:「真理是甚麼呢?」(約18:33-38)

真理是甚麼呢?這似乎不僅是個哲學問題.成為肉身的道,指出自己來到世間,是「特意為真理作見證」.彼拉多或許不是一個好官,但我們心中認同的問題:真理是甚麼呢?而耶穌又為甚麼答得那麼迂迴?其實,耶穌的迂迴是一貫的,因為「道可道非常道」,他只是拒絕落在人的語言圈套中.

很可惜,對於真理是甚麼,甚至連信徒亦往往不求甚解,只將真理等同於所屬宗派所認信的教義.你想多問兩句,他們就說那是關乎信心,難怪教外人往往對我們所宣講的「真理」不當一回事!

其實不少現代人也是在問著同樣的問題:真理到底是甚麼?甚至到底有沒有真理?對於在學術圈子工作的人來說,學術研究當然就是發現真理的最佳途徑.它標榜客觀公正,反對盲從,迷信,是這個世代最受歡迎和信任的答案.

對於一些人來說,學術的真理就是一切的真理.它給我們看見自然界如何運作,人類的歷史如何演進,甚至是人本身的思想或心理背後操縱的因素和原則.這些知識道理帶給我們莫大的喜樂和滿足.甚至基督徒亦樂於認同知識帶來的喜悅: “All truth is God’s truth, wherever it is found.”

不過很多時就算學術研究跟生活意義之間亦沒有一種必然的關係.譬如一個研究生命基因的遺傳因子工程專家,未必能夠告訴一個考慮自殺的人為何值得生存下去.

輔導心理學特意將真相(truth)跟人們對它的知覺(perception)分開來處理:我們對事物真相的掌握,往往只是一個知覺的表象,而這個表象的形成,又無可避免地滲入了我們本身的種種主觀條件,例如我們的背景與訓練,創傷與情結,感受與聯想.這也是真理問題難以簡單回應的主要原因:真相可能只有一個,卻可以有多向度的認知和理解.

真理的知識論(The Epistemology of Truth):

知識的主觀性:一切的知識(包括對真理的認知)都必然牽涉「知者」,因此一切有關真理的認知,亦必然同時牽涉我們作為人的一切限制.絕對真理落到我們的手,無可避免地變成相對.

歷史作為一個進程,就意味著我們所知道的一切,將來必然會被更廣闊的視野,更高層次的理論所超越和取代.

文字的有限性:縱然一件事實幾乎可以從無限的角度去詮釋和理解,人的語言卻只能夠從一個或頂多幾個角度去呈現和表達.

由此發展出來的,是事物的多面性與相對性.有些人視之為對真理的挑戰.但從積極的角度來看,事物的多面性與相對性,反而強化了它的層次與立體感,也讓人有不斷發掘與進深的空間.因此當我們說「真理不是鐵板一塊」,並沒有絲毫妥協真理的絕對性,反而是提醒我們在真理面前永遠要存著謙卑恭謹的態度,明白「吾生也有涯,而知也無涯.」

絕對的真理確實是不可道的；而可道的真理亦不可能是絕對的.《希伯來書》形容上帝是 「烈火」(12:29)；保羅在提前6:16形容上帝是「那獨一不死,住在人不能靠近的光裡,是人未曾看見,也是不能看見的」.

今年的聖誕,就讓我們從一個不同的角度,去重新思考聖誕的意義:「大哉敬虔的奧秘,無人不以為然,就是上帝在肉身顯現!」那位創造宇宙萬物的主宰,那位不可道的道,那位真理的本身,從「人不能靠近的光」降卑來到世間,讓我們透過他,可以從一個嶄新的角度去思考自己的人生!



\newpage

\section{以試煉為喜樂}
\label{sec:395}
\textbf{朱建忠博士 - 以試煉為喜樂}
\newline
\newline
\hyperref[sec:394]{< < < PREV SERMON < < <}
~
\hyperref[sec:toc]{[返目錄]}
~
\hyperref[sec:396]{> > > NEXT SERMON > > >}


雅各書 1:1-12
\newline
\begin{longtable}{cl}
\hline
\hline
章節 & 經文 (和合本修訂版)\\
\hline
1:1 & \begin{tabularx}{0.7\textwidth}{X} 神和主耶穌基督的僕人雅各問候散居在各處的十二個支派的人。 \end{tabularx} \\ \\ \relax
1:2 & \begin{tabularx}{0.7\textwidth}{X} 我的弟兄們，你們遭受各種試煉時，都要認為是大喜樂， \end{tabularx} \\ \\ \relax
1:3 & \begin{tabularx}{0.7\textwidth}{X} 因為知道你們的信心經過考驗，就生忍耐。 \end{tabularx} \\ \\ \relax
1:4 & \begin{tabularx}{0.7\textwidth}{X} 但要讓忍耐發揮完全的功用，使你們能又完全又完整，一無所缺。 \end{tabularx} \\ \\ \relax
1:5 & \begin{tabularx}{0.7\textwidth}{X} 你們中間若有缺少智慧的，該求那厚賜與眾人又不斥責人的神，神必賜給他。 \end{tabularx} \\ \\ \relax
1:6 & \begin{tabularx}{0.7\textwidth}{X} 只要憑著信心求，一點也不疑惑；因為那疑惑的人，就像海中的波浪被風吹動翻騰。 \end{tabularx} \\ \\ \relax
1:7 & \begin{tabularx}{0.7\textwidth}{X} 這樣的人不要想從主那裡得到甚麼。 \end{tabularx} \\ \\ \relax
1:8 & \begin{tabularx}{0.7\textwidth}{X} 三心二意的人，在他一切所行的路上都搖擺不定。 \end{tabularx} \\ \\ \relax
1:9 & \begin{tabularx}{0.7\textwidth}{X} 卑微的弟兄要因高升而誇耀， \end{tabularx} \\ \\ \relax
1:10 & \begin{tabularx}{0.7\textwidth}{X} 富足的卻要因被降卑而誇耀，因為富足的人要消逝，如同草上的花一樣。 \end{tabularx} \\ \\ \relax
1:11 & \begin{tabularx}{0.7\textwidth}{X} 太陽出來，熱風颳起，草就枯乾，花也凋謝，它美麗的樣子就消失了；那富足的人在他一生的奔波中也要這樣衰殘。 \end{tabularx} \\ \\ \relax
1:12 & \begin{tabularx}{0.7\textwidth}{X} 忍受試煉的人有福了，因為他經過考驗以後必得生命的冠冕，這是主應許給愛他之人的。 \end{tabularx} \\ \\
[1ex]
\hline
\hline
\end{longtable}
講題:以試煉為喜樂　Joy Through Trials

經文:雅各書1章1-12節

講員:朱建忠博士

場合:香港中文大學崇基學院禮拜堂

日期:2017年12月31日

經文:雅各書 1 章:1 做神和主耶穌基督僕人的雅各,請散住十二個支派之人的安!2 我的弟兄們,你們落在百般試煉中,都要以為大喜樂,3 因為知道,你們的信心經過試驗就生忍耐.4 但忍耐也當成功,使你們成全,完備,毫無缺欠. 5 你們中間若有缺少智慧的,應當求那厚賜於眾人也不斥責人的神,主就必賜給他. 6只要憑著信心求,一點不疑惑；因為那疑惑的人就像海中的波浪,被風吹動翻騰. 7 這樣的人不要想從主那裡得什麼, 8心懷二意的人在他一切所行的路上都沒有定見.9 卑微的弟兄升高,就該喜樂； 10 富足的降卑也該如此,因為他必要過去,如同草上的花一樣.11太陽出來,熱風刮起,草就枯乾,花也凋謝,美容就消沒了；那富足的人在他所行的事上,也要這樣衰殘.12 忍受試探的人是有福的!因為他經過試驗以後,必得生命的冠冕,這是主應許給那些愛他之人的.

引言:當我們聽到“試煉”時會聯想到什麼呢?...今天的經文上寫道:你們落在百般試煉中,都要以為大喜樂.這怎麼可能?我們的經驗好像不是這樣?要麼就是聖經有問題或是我們的經驗並不代表生命的真相.我想作神的兒女都不會懷疑聖經的真實性,所以是我們的經驗:在試煉中的悲愁,煎熬並非是我們生命的本質,意即,並非我們經歷試煉的時候,必然的會是悲愁,煎熬,而是可以視之為喜樂的.

一. 經試煉邁向完全 V1-4

1 做神和主耶穌基督僕人的雅各,請散住十二個支派之人的安!2 我的弟兄們,你們落在百般試煉中,都要以為大喜樂,3 因為知道,你們的信心經過試驗就生忍耐. 4但忍耐也當成功,使你們成全,完備,毫無缺欠.

聖經在這裡告訴我們可以經試煉為喜樂的原因是:3 信心經過試驗就生忍耐. 4 但忍耐也當成功,使你們成全,完備,毫無缺欠.第四節中的也當成功,使.. 成全完備,毫無缺欠,其實是同一個意思.也就是,信心經過試煉可以讓我們達到完全.聖經很強調信徒的生命要達到完全.太五:48 耶穌說,“你們要完全,像你們的天父完全一樣.” 保羅說,“這不是說,我已經得著了,已經完全了.我乃是竭力追求...我只有一件事,就是忘記背後,努力面前,向著標竿直跑... 所以我們中間凡是完全人,總要存這樣的心.”(腓 3:12-15) 希伯來書的作者也勉勵他的讀者:“存心忍耐,奔那擺在我們前頭的路程.”(來 12:1) 這些聖經經句都告訴我們,神的兒女不是信了就好了,結束了,而是一個開始,神的兒女需要邁向完全.

讓我們回到雅各書的背景來瞭解一下,為何《雅各書》如此強調信心要經過試煉,並且在試煉中邁向完全.我們來看這卷書的受書人是:散住十二個支派之人,意思就是散住在十二支派的猶太基督徒.而寫作的時間當時正是發生災荒的時候(一說是西元四十多年,另一說六十多年正是羅馬鎮壓猶太人的時候),作為猶太基督徒更是受到多重的逼迫.對我們來說,當試煉來的時候,有時我們可能會抗拒,但力不能勝；也有時,我們可能會灰心,在壓力下投降,甚至懷疑神是否顧念我們,並且因種種困難而不滿發怨言.但聖經在這裡提醒我們,可以積極的來看待試煉,將試煉看成是生命操練的機會,視之為喜樂.我們要明白,“這試煉臨到我,是神容許的.祂對我有美好的心意.雖然我現在不明白,但我會嘗試尋求.他希望神的心意在我生命中成就.” 這正是雅各提議的,你們落在百般試煉中,都要以為大喜樂.

舉例說明:

二. 求神賜智慧

5 你們中間若有缺少智慧的,應當求那厚賜於眾人也不斥責人的神,主就必賜給他. 6 只要憑著信心求,一點不疑惑；因為那疑惑的人就像海中的波浪,被風吹動翻騰. 7 這樣的人不要想從主那裡得什麼, 8 心懷二意的人在他一切所行的路上都沒有定見.

但當我們走在這條繼續完全的完全之旅上時,有時我們也會困惑,後退,就如雅各講的出現“缺乏”,這正是我們生命的實況.所以,我們需要禱告祈求. “5 你們中間若有缺少智慧的,應當求那厚賜於眾人也不斥責人的神,主就必賜給他.”

為什麼是求智慧而不是其它?如醫治?富足?智慧與試煉,問題,苦難有什麼關係呢?當我們遇到問題與苦難時,我們不是經常尋求解決方案嗎?如生病求醫治.但仔細思想,你就會明白,這正是我們為什麼有時會感覺並不需要神的錯覺.如果我們只是聚焦在問題以及答案,我們怎麼會還需要神呢?當我們遇到苦難時,若只是向神要答案,而忽視了尋求神本身,這不正是我們生命缺乏的更深 層次原因嗎?

不是說求知識,尋方法,要答案不對,而是一心尋求這些容易使我們對生命產生錯覺.一心只求解答我的問題,解決我的需要,容易將問題看得比生命更大, 以至於錯過與神的相遇. “5 你們中間若有缺少智慧的,應當求那厚賜於眾人也 不斥責人的神,主就必賜給他.” 雅各說的智慧,並非是指給出一個簡單的答 案,而是引導我們認識神的慷慨和美善,進而可以更加的認識神及親近神.這種 4 智慧不是要按照需要者提供所想要的,真正來說,是使之產生信心,盼望與愛.

(進一步說明“信”,“望”,“愛”):

三.  以智慧看人生跌宕

9 卑微的弟兄升高,就該喜樂； 10 富足的降卑也該如此,因為他必要過去,如同草上的花一樣. 11 太陽出來,熱風刮起,草就枯乾,花也凋謝,美容就消沒了；那富足的人在他所行的事上,也要這樣衰殘.

卑微的弟兄升高,就該喜樂,我們也為此感謝神；但除非一個人緊緊跟從神,走在完全的路上,不然高升可能是跌倒的第一步.所以讓人真正喜樂的不是高升,因為它像地上的花草一樣,必要過去.生命比升職更大,而神比生命更大.同樣,當富足的降卑,也可以喜樂.因為降卑如同升高同樣會過去.生命比降卑更值得 看重.

我們的經驗自然是追求高升恐怕降卑,積累財富躲避貧窮,但這些都如同地上的花草一樣轉瞬即逝.一旦讓這些如同花草的東西主導我們的生命,上主的子 民就會離開那令他們完全的道路,也忘記了信,望,愛.這樣的話,生活只是憑 添技巧,知識與本錢,而沒有生命的智慧. 舉例說明:

結語:親愛的弟兄姊妹,讓我們再次確信聖經所應許的:“忍受試探的人是有福 的!因為他經過試驗以後,必得生命的冠冕,這是主應許給那些愛他之人的.” 現在當你再聽到試煉時,又會想到些什麼呢?



\newpage

\section{誰配揭開歷史之謎}
\label{sec:396}
\textbf{鄧瑞強博士 - 誰配揭開歷史之謎?}
\newline
\newline
\hyperref[sec:395]{< < < PREV SERMON < < <}
~
\hyperref[sec:toc]{[返目錄]}
~
\hyperref[sec:397]{> > > NEXT SERMON > > >}


啟示錄 5:1-14
\newline
\begin{longtable}{cl}
\hline
\hline
章節 & 經文 (和合本修訂版)\\
\hline
5:1 & \begin{tabularx}{0.7\textwidth}{X} 我看見坐在寶座那位的右手中有書卷，正反面都寫著字，用七個印密封著。 \end{tabularx} \\ \\ \relax
5:2 & \begin{tabularx}{0.7\textwidth}{X} 我又看見一位大力的天使大聲宣告說：「有誰配展開那書卷，揭開那七個印呢？」 \end{tabularx} \\ \\ \relax
5:3 & \begin{tabularx}{0.7\textwidth}{X} 在天上、地上、地底下，沒有人能展開、能閱覽那書卷。 \end{tabularx} \\ \\ \relax
5:4 & \begin{tabularx}{0.7\textwidth}{X} 因為沒有人配展開、閱覽那書卷，我就大哭。 \end{tabularx} \\ \\ \relax
5:5 & \begin{tabularx}{0.7\textwidth}{X} 長老中有一位對我說：「不要哭。看哪，猶大支派中的獅子，大衛的根，他已得勝，能展開那書卷，揭開那七個印。」 \end{tabularx} \\ \\ \relax
5:6 & \begin{tabularx}{0.7\textwidth}{X} 我又看見寶座和四個活物，以及長老之中有羔羊站著，像是被殺的，有七個角七隻眼睛，就是神的七靈，奉差遣往普天下去的。 \end{tabularx} \\ \\ \relax
5:7 & \begin{tabularx}{0.7\textwidth}{X} 這羔羊前來，從坐在寶座上那位的右手中拿了書卷。 \end{tabularx} \\ \\ \relax
5:8 & \begin{tabularx}{0.7\textwidth}{X} 他一拿了書卷，四活物和二十四位長老就俯伏在羔羊面前，各拿著琴和盛滿了香的金爐；這香就是眾聖徒的祈禱。 \end{tabularx} \\ \\ \relax
5:9 & \begin{tabularx}{0.7\textwidth}{X} 他們唱新歌，說：「你配拿書卷，配揭開它的七印；因為你曾被殺，用自己的血從各支派、各語言、各民族、各邦國中買了人來，使他們歸於神， \end{tabularx} \\ \\ \relax
5:10 & \begin{tabularx}{0.7\textwidth}{X} 又使他們成為國民和祭司，歸於我們的神；他們將在地上執掌王權。」 \end{tabularx} \\ \\ \relax
5:11 & \begin{tabularx}{0.7\textwidth}{X} 我又觀看，我聽見寶座和活物及長老的周圍有許多天使的聲音；他們的數目有千千萬萬， \end{tabularx} \\ \\ \relax
5:12 & \begin{tabularx}{0.7\textwidth}{X} 大聲說：「被殺的羔羊配得權能、豐富、智慧、力量、尊貴、榮耀、頌讚。」 \end{tabularx} \\ \\ \relax
5:13 & \begin{tabularx}{0.7\textwidth}{X} 我又聽見在天上、地上、地底下、滄海裡和天地間一切所有被造之物，都說：「願頌讚、尊貴、榮耀、權勢，都歸給坐在寶座上的那位和羔羊，直到永永遠遠！」 \end{tabularx} \\ \\ \relax
5:14 & \begin{tabularx}{0.7\textwidth}{X} 四活物就說：「阿們！」眾長老也俯伏敬拜。 \end{tabularx} \\ \\
[1ex]
\hline
\hline
\end{longtable}
主題:誰配揭開歷史之謎?　Who is Worthy to Open the Mystery of History?

經文:啟示錄5:1-14

講員:鄧瑞強博士

場合:香港中文大學崇基學院禮拜堂

日期:2018年1月7日

各位弟兄姊妹,早安.在這新的一年,祝願各人能更深體會神拯救的大能,並能更深明白生命的意義.

舊的一年過去,新的一年剛剛開始,我們或會懷緬一下過去的事情,檢討一下去年的成敗,思考一下今年的計劃.面對過去,我們會追憶.面向將來,我們會期盼.我們是「歷史動物」,有特殊的歷史意識,對時間的流逝,有微妙的感覺.我們會討論歷史的成敗,會「數風流人物」.我們會想像美好的明天,會叫人「還看今朝」.就算是個人的生活歷史,我們都會下評語,「過去一年過得真好」,或者,「過去一年真失敗」.似乎,我們直覺地知道,時間的流逝不可能是空白的,無內容的,無意義的.時間的流逝,總有某種含意,某種目的.問題只是,歷史有何意義,時間的流逝引領我們走向何方?

啟示錄是一卷探討歷史意義的書卷,它討論歷史的終局.在歷史終局的角度裡,看看人間的帝王將相,英雄豪傑誰能經得起歷史的浪花淘盡.是什麼人物,在日子消逝後,留下深刻的足跡,留下不會磨滅的名字?啟示錄又將人的視野從地上帶到天上,讓人能脫離一下塵世的混濁,無止盡的貪婪,利慾薰心的盲目,在明淨的藍天之上,細看短視而墮落的人為了一塊骨頭而爭得你死我活.啟示錄教我們在時間的終點看現在,使我們有歷史的智慧；也教我們在空間的高點看當下,使我們能看清楚醜陋與高貴.

今日,讓我們透過一段啟示錄的經文,在空間的極處,在人生的高點,在塵世之上的天界,去思考歷史的意義.

今日的講道經文:啟示錄5:1－14

啟5:1　我看見坐寶座的右手中有書卷,裡外都寫著字,用七印封嚴了.

啟5:2　我又看見一位大力的天使大聲宣傳說:「有誰配展開那書卷,揭開那七印呢?」

啟5:3　在天上,地上,地底下,沒有能展開,能觀看那書卷的.

啟5:4　因為沒有配展開,配觀看那書卷的,我就大哭.

啟5:5　長老中有一位對我說:「不要哭!看哪,猶大支派中的獅子,大衛的根,他已得勝,能以展開那書卷,揭開那七印.」

啟5:6　我又看見寶座與四活物,並長老之中有羔羊站立,像是被殺過的,有七角七眼,就是神的七靈,奉差遣往普天下去的.

啟5:7　這羔羊前來,從坐寶座的右手裡拿了書卷.

啟5:8　他既拿了書卷,四活物和二十四位長老就俯伏在羔羊面前,各拿著琴和盛滿了香的金爐；這香就是眾聖徒的祈禱.

啟5:9　他們唱新歌,說:你配拿書卷,配揭開七印；因為你曾被殺,用自己的血從各族,各方,各民,各國中買了人來,叫他們歸於神,

啟5:10　又叫他們成為國民,作祭司歸於神,在地上執掌王權.

啟5:11　我又看見且聽見,寶座與活物並長老的周圍有許多天使的聲音；他們的數目有千千萬萬,

啟5:12　大聲說:曾被殺的羔羊是配得權柄,豐富,智慧,能力,尊貴,榮耀,頌讚的.

啟5:13　我又聽見在天上,地上,地底下,滄海裡,和天地間一切所有被造之物,都說:但願頌讚,尊貴,榮耀,權勢都歸給坐寶座的和羔羊,直到永永遠遠!

啟5:14　四活物就說:「阿們!」眾長老也俯伏敬拜.

(1)誰「配」展開那書卷?

在高天之上,在價值至高之處,在永恆的領域,有一寶座.在寶座上,坐著創造歷史並將歷史帶到終點的神.他創造歷史,創造人類,創造萬有,並讓人類在大地上肩負著歷史的責任,以致神的美意在歷史的過程中顯明出來.歷史不是漫無目的的,我們在大地上的生活也不是蒼白空洞的.歷史滲透著神的美意,但我們卻將歷史書寫成血腥的戰爭史.征服身邊的異族,征服身邊的土地,征服身邊的文化,然後,造成一個統一的大帝國；然後,說這便是歷史的任務；然後,將欲望的無窮滿足看成是歷史的終點.

啟5:1　我看見坐寶座的右手中有書卷,裡外都寫著字,用七印封嚴了.

在天上,坐寶座的神,右手中有書卷.這書卷,就是神的旨意,就是神的歷史計劃,就是神要我們知道的歷史的意義.有學者指出,「坐寶座的右手中有書卷」應理解為「坐寶座的右邊有書卷」,而這「寶座」也不是一單座位的座椅,而是一類似現代「沙發」的「躺椅」,這躺椅可容多於一人咒詛.歷史的秘密,就放在坐寶座者的右邊.這書卷被七個火蠟印章封住了.

一個重要問題:誰配揭開這歷史之謎?

啟5:2　我又看見一位大力的天使大聲宣傳說:「有誰配展開那書卷,揭開那七印呢?」

這裡最關鍵的用字,就這一個「配」字.誰「配」?誰「配」坐在至高者的右邊,揭示歷史的真義?誰「配」?

自古以來,所有英雄故事,都是先要人克服無盡劫難,完成一超凡任務,然後才「配」成為一英雄.所謂「配」,就是在考驗中合格.就如《西遊記》的故事,唐僧使徒要走過多少異域,克服多少來自上天下地的妖魔鬼怪,才能取得西經.話說,唐僧使徒四人經歷完八十次劫難,已取得西經,理應已「配」成為天界人物.但原來「配」與「不配」,自有一種天理在,西經到手,也未算得是「配」.原來,圓滿的考驗要求一完美的考驗數目,九九八十一才算圓滿.八十難還差一難,他們必須通過八十一次考驗,才算是通過了完整的考驗,然後才「配」成為天界人物.菩薩最後要刻意製造多一苦難來,讓他們多受一劫.考驗完滿了,成功了,他們就名正言順地「配」成為天界人物.

在啟示錄裡,那天使問:「誰能在人生中完美地勝過人性的種種軟弱?誰能經得起神的考驗?誰配向世人揭示神在歷史中的大計?誰配?」

啟5:3　在天上,地上,地底下,沒有能展開,能觀看那書卷的.

啟5:4　因為沒有配展開,配觀看那書卷的,我就大哭.

啟示錄表達出一種悲哀,一種無以名狀的悲哀,一種你想聆聽真理卻發現竟然無人言說真理的悲涼,一種你想明白歷史的意義卻發現歷史失落了所有意義的蒼茫,一種你想活出生命的真實內容卻發覺所有人都只是行屍走肉的空洞.一種大哀!一種天地間的大哀!啟示錄的作者在荒涼的大地上「大哭」.

(2)被殺的羔羊,他「配」展開那書卷

在茫茫然的大哀中,有從天而來的安慰.有人歷盡人間的磨難與考驗,並且得勝,這個人「配」打開那藏著歷史奧秘的書卷.

啟5:5　長老中有一位對我說:「不要哭!看哪,猶大支派中的獅子,大衛的根,他已得勝,能以展開那書卷,揭開那七印.」

啟5:6　我又看見寶座與四活物,並長老之中有羔羊站立,像是被殺過的,有七角七眼,就是神的七靈,奉差遣往普天下去的.

啟5:7　這羔羊前來,從坐寶座的右手裡拿了書卷.

這人是誰?他是「猶大支派中的獅子」,這個講法出自創世記49章,在歷史的始點,神已預備好適當人選.「大衛的根」,這個意念出自先知以賽亞,預告新的王者的來臨.在神的計劃裡,在某時某刻,大地自有聖人出.他來,以其特有的生命,克服世間一切考驗,展示歷史的真義.

他來了,他得勝了,他配展開那裡卷.他是誰?他是「被殺的羔羊」.剛剛才說,他是「獅子」,是「王者」；現在卻說,他是「羔羊」,是「被殺的」.「獅子」與「羔羊」,是兩種不能相容的元素,像「生命」和「死亡」,不能混為一談.偏偏,以「羔羊」的生命活出「獅子」的王道,以「捨己」的「死亡」去活出「不朽」的「生命」,以全然卑微去顯出生命的偉大,原來這卻是破解歷史之謎的關鏈.生命的光榮與永恆,在於卑微地承受屈辱的十架犧牲,在於為了愛世人而將自己的命運置之度外,在於為了向歷史負責任而在不可言說真理的世代堅持言說真理.這是甘心被殺的羔羊,卻顯出獅子般的無比偉大.是這樣的耶穌,他「配」揭開神的心,呈示人行走在歷史上的真義.

這羔羊有七角七眼,「角」代表「勇氣」與「力量」,「七角」就是大無畏的「勇氣」與不能屈的「力量」.關於「眼」,經文說,「七眼」就是「神的七靈」,「神的靈」是湧現生命的泉源,也是煥發生命的力量.耶穌,作為羔羊,以無比的勇氣,背起了歷史的十架,犧牲了自己,卻喚醒人類人性深處潛藏的生命種子,也啟導歷史的新方向.這生命,配打開那書卷,昭示歷史的意義.

如今,這位揭示歷史生命的意義的耶穌,坐在天父右邊,引導著歷史走向終點.

(3)三回合的讚美

強者的暴力,人性的貪婪,享樂的沉溺,短視的目標,低下的品格,利益的計算,朝生暮死的生活態度,都會令歷史失去方向.耶穌,在世俗的嘈吵中堅持靜心聆聽天父的聲音,在自私的世界毅然為他人而犧牲,在人人眼目向下時勇於望向上天.他點出了生命的意義,指明歷史的方向,也以被殺羔羊的生命,引導歷史走向圓滿的終點.這樣的人物,如此的神聖,配受眾生的讚美.

經文描繪出對耶穌作出的三個回合的讚美.

啟5:8　他既拿了書卷,四活物和二十四位長老就俯伏在羔羊面前,各拿著琴和盛滿了香的金爐；這香就是眾聖徒的祈禱.

啟5:9　他們唱新歌,說:你配拿書卷,配揭開七印；因為你曾被殺,用自己的血從各族,各方,各民,各國中買了人來,叫他們歸於神,

啟5:10　又叫他們成為國民,作祭司歸於神,在地上執掌王權.

這讚美講出耶穌如何改變歷史.他使人歸向神,成為一特殊群體,作為溝通天地的祭司.

啟5:11　我又看見且聽見,寶座與活物並長老的周圍有許多天使的聲音；他們的數目有千千萬萬,

啟5:12　大聲說:曾被殺的羔羊是配得權柄,豐富,智慧,能力,尊貴,榮耀,頌讚的.

這讚美將人間一切美好而尊貴的品質歸給耶穌.

啟5:13　我又聽見在天上,地上,地底下,滄海裡,和天地間一切所有被造之物,都說:但願頌讚,尊貴,榮耀,權勢都歸給坐寶座的和羔羊,直到永永遠遠!

啟5:14　四活物就說:「阿們!」眾長老也俯伏敬拜.

耶穌坐在至高者的右邊,接受全宇宙萬物的敬拜.這是歷史的終點,萬物一致地面向神,高舉神,歸向神.

我們總會敬拜某些東西.我們敬拜的對象,就是我們心之所向,也是我們生活之所依.這左右我們的生活,也決定我們在歷史中的尋索.你敬拜物質,則歷史的意義便在乎物質的豐富,自己的物資不夠,便去搶別人的,搶得越多,搶得越成功,便被看為是歷史的進步.啟示錄講的敬拜,是敬拜那位在愛中捨己的耶穌.實踐他的生命,才會活出生命的意義,才會揭露歷史的真義.敬拜的對象錯了,生活便迷失,歷史便黯淡.敬拜的對象對了,生活便豐盛,歷史便光明.

教會的生活,敬拜神是不可少的.教會要向世界表明,世人當敬拜誰,生命當何所依,歷史當何所往.我知道我們禮拜堂的詩班,會有些變動,但希望詩班裡各弟兄姊妹明白,敬拜之為敬拜,是我們將心靈的最高貴獻禮呈奉給那最值得尊崇的神.我們需要有詩班,代表全會眾向神作出這種獻禮.

這種獻禮,有如耶穌的羔羊獻禮,將生命專一的獻給神.這種獻禮,無關乎環境的變化,無關乎處境的轉換,無關乎人事的更替.用啟示錄的語言講,正正就在那看似無神的世界,正正就在那看來混亂的局勢,正正就在那歷史不明的時刻,更需要忠勇地敬拜神.在不能敬拜的時刻,在不想敬拜神的時刻,在不願意敬拜神的時刻,仍堅持敬拜,這就是對神的尊榮的無限肯定.

神高高在上,也站在歷史的終點看著我們.我們會像耶穌一樣捨己地生活嗎?我們會像耶穌一樣,在愛的生活中,揭示歷史的意義嗎?我們會否因為看到神在歷史中的美意,而全心敬拜神呢?

但願:榮耀歸於聖父,聖子,聖靈.阿們.



\newpage

\section{喜事起事}
\label{sec:397}
\textbf{林豪恩先生 - 喜事起事}
\newline
\newline
\hyperref[sec:396]{< < < PREV SERMON < < <}
~
\hyperref[sec:toc]{[返目錄]}
~
\hyperref[sec:398]{> > > NEXT SERMON > > >}


馬太福音 22:1-14
\newline
\begin{longtable}{cl}
\hline
\hline
章節 & 經文 (和合本修訂版)\\
\hline
22:1 & \begin{tabularx}{0.7\textwidth}{X} 耶穌又用比喻對他們說： \end{tabularx} \\ \\ \relax
22:2 & \begin{tabularx}{0.7\textwidth}{X} 「天國好比一個王為他兒子擺設娶親的宴席。 \end{tabularx} \\ \\ \relax
22:3 & \begin{tabularx}{0.7\textwidth}{X} 他打發僕人去，請那些被邀的人來赴宴，他們卻不肯來。 \end{tabularx} \\ \\ \relax
22:4 & \begin{tabularx}{0.7\textwidth}{X} 王又打發別的僕人，說：『你們去告訴那被邀的人，我的宴席已經預備好了，牛和肥畜已經宰了，各樣都齊備，請你們來赴宴。』 \end{tabularx} \\ \\ \relax
22:5 & \begin{tabularx}{0.7\textwidth}{X} 那些人不理就走了，一個到自己田裡去，一個做買賣去。 \end{tabularx} \\ \\ \relax
22:6 & \begin{tabularx}{0.7\textwidth}{X} 其餘的抓住僕人，凌辱他們，把他們殺了。 \end{tabularx} \\ \\ \relax
22:7 & \begin{tabularx}{0.7\textwidth}{X} 王就大怒，發兵除滅那些兇手，燒燬他們的城。 \end{tabularx} \\ \\ \relax
22:8 & \begin{tabularx}{0.7\textwidth}{X} 於是王對僕人說：『喜宴已經齊備，只是所邀的人不配。 \end{tabularx} \\ \\ \relax
22:9 & \begin{tabularx}{0.7\textwidth}{X} 所以你們要往岔路口上去，凡遇見的，都邀來赴宴。』 \end{tabularx} \\ \\ \relax
22:10 & \begin{tabularx}{0.7\textwidth}{X} 那些僕人就出去，到大路上，凡遇見的，不論善惡都招聚了來，宴席上就坐滿了客人。 \end{tabularx} \\ \\ \relax
22:11 & \begin{tabularx}{0.7\textwidth}{X} 王進來見賓客，看到那裡有一個沒有穿禮服的， \end{tabularx} \\ \\ \relax
22:12 & \begin{tabularx}{0.7\textwidth}{X} 就對他說：『朋友，你到這裡來怎麼不穿禮服呢？』那人無言可答。 \end{tabularx} \\ \\ \relax
22:13 & \begin{tabularx}{0.7\textwidth}{X} 於是王對侍從說：『捆起他的手腳，把他扔在外邊的黑暗裡；在那裡他要哀哭切齒了。』 \end{tabularx} \\ \\ \relax
22:14 & \begin{tabularx}{0.7\textwidth}{X} 因為被召的人多，選上的人少。」 \end{tabularx} \\ \\
[1ex]
\hline
\hline
\end{longtable}
講題:喜事起事  Conflicts in a Wedding Banquet

經文:馬太福音22章1-14節

講員:林豪恩先生

場合:崇基學院禮拜堂主日崇拜

日期:2018年1月21日

一. 引言

有一首粵語詩歌,歌名叫「回家」,其中幾句如此說:

步伐是否感覺疲倦了

現實令你感到無奈

只懂每日每夜忙碌裡

兜兜轉轉似沒了沒完

但願讓你知道神是愛

在路上每一秒同在

可否接受這份純真愛

一生歡欣快樂精彩

回家將傷痛放下

看透俗世謊話

迷失漆黑中祂總把你念掛

回家即使新世代

全然皆冰冷

在這屋裡有主的愛別再怕

這是受歡迎的佈道會詩歌,因為詩歌唱出了不少現代城市人的心聲.人在江湖,身不由己,競爭處處,經常處於作戰狀態,難免感到冰冷和疲乏,令人盼望在紛擾的世界中,找到一處淨土.尤其是那些不喜歡衝突,喜歡和諧的朋友,到教會來,就是希望避開塵世的是非,追求一刻的平靜.這樣的期望,其實都可以理解,也可說是卑微的.通常,「入教未深」的朋友,都能夠如願以償.然而,有些「入教已深」的朋友,卻感到迷惑,因為他們漸漸發現教會似是另一個「江湖」,信仰群體似是另一個「是非之地」.他們漸漸懷疑,信仰其實是否他們要尋找的淨土?教會是否只是紛擾世界的縮影?再聽到今天的福音經課中這個天國的比喻,不但沒有釋疑,甚至帶來更重的迷茫與失落感覺.

二. 經文

22:1耶穌又用比喻對他們說:

22:2「天國好比一個王為他兒子擺設娶親的筵席,

22:3就打發僕人去,請那些被召的人來赴席,他們卻不肯來.

22:4王又打發別的僕人,說:『你們告訴那被召的人,我的筵席已經預備好了,牛和肥畜已經宰了,各樣都齊備,請你們來赴席.』

22:5那些人不理就走了；一個到自己田裡去；一個做買賣去；

22:6其餘的拿住僕人,凌辱他們,把他們殺了.

22:7王就大怒,發兵除滅那些兇手,燒毀他們的城.

22:8於是對僕人說:『喜筵已經齊備,只是所召的人不配.

22:9所以你們要往岔路口上去,凡遇見的,都召來赴席.』

22:10那些僕人就出去,到大路上,凡遇見的,不論善惡都召聚了來,筵席上就坐滿了客.

22:11王進來觀看賓客,見那裡有一個沒有穿禮服的,

22:12就對他說:『朋友,你到這裡來怎麼不穿禮服呢?』那人無言可答.

22:13於是王對使喚的人說:『捆起他的手腳來,把他丟在外邊的黑暗裡；在那裡必要哀哭切齒了.』

22:14因為被召的人多,選上的人少.」

三. 釋經與應用

耶穌所說的這個故事,是用來比喻天國的.在人們的想像中,天國應該是安全的遊樂場,是和平,友愛,舒暢的樂土.可惜,這個天國故事卻充滿張力,衝突處處.

熱鬧歡愉對孤單落寞

22:2「天國好比一個王為他兒子擺設娶親的筵席

耶穌為天國的故事所定下的場景,的確是熱鬧歡愉的,因為這是一個喜慶場合,是一個娶親的筵席.我常常在婚禮訓勉首先恭賀雙方父母,因為結婚不單是兩個人的事情,更是兩個家庭的事情.看見兒女成家立室,父母會感到自己完成人生的一項重要任務,因而感到安慰.平民百姓的婚筵尚且是喜氣洋洋,更何況耶穌比喻中的是國王的婚筵呢?或者我們可以想像當年英國皇儲查理士與戴安娜舉行婚禮時舉國歡騰的熱鬧.這個比喻的奇怪之處在於故事預設的是熱鬧歡愉的場景,但故事的結束卻是孤單落寞.

22:13於是王對使喚的人說:『捆起他的手腳來,把他丟在外邊的黑暗裡；在那裡必要哀哭切齒了.』

熱鬧歡愉的婚筵,與一個人在黑暗裡哀哭切齒,是兩個非常不配合的畫面,是反差非常大的對比.這樣充滿矛盾與張力的故事,竟然用來比喻天國.

和好結合對兵戎相見

22:2「天國好比一個王為他兒子擺設娶親的筵席

娶親婚筵,是結合,而不是分裂的場景.正如上段所述,婚姻是兩個人的結合,也是兩個家庭的結合.在古代社會,甚至是兩個群族,或者兩個國家的結合.因此,沒有必要問西漢的昭君有多愛匈奴的單于,昭君出塞,被稱為和番的行動,希望化解當時漢人與匈奴之間的對抗和衝突.耶穌所說的天國比喻的開始是和好結合的場景,但故事的發展卻是分裂和對抗.

22:6其餘的拿住僕人,凌辱他們,把他們殺了.

22:7王就大怒,發兵除滅那些兇手,燒毀他們的城.

和好結合的婚筵,與對抗殺戮的場面,又是兩個非常不配合的畫面,又是反差非常大的狀況.如此不協調的場面,竟然也用來比喻天國.

認真主人對不認真客人

22:4王又打發別的僕人,說:『你們告訴那被召的人,我的筵席已經預備好了,牛和肥畜已經宰了,各樣都齊備,請你們來赴席.』

在當時的猶太人社會,請人赴宴有兩個步驟:主人先打發人發出初步的邀約,屆時再差派僕人通知客人赴宴.這是一種特別有禮的做法.故事提到王差派僕人邀請各人赴宴,顯示第一次的邀請已經發出,並已被接受.現在,筵席要開始了,美酒佳餚都準備好了,王就派僕人再往邀請,以表達王給予客人的重視和榮譽.故事吊詭之處,就是如此認真的婚筵主人,如此把榮譽給予客人,客人的反應卻非常不認真,甚至敷衍.

22:5那些人不理就走了；一個到自己田裡去；一個做買賣去；

認真的主人,遇上不認真的客人,又是一種落差大的對比.原來這種狀況也與天國有關.

好客主人對好戰客人

22:4王又打發別的僕人,說:『你們告訴那被召的人,我的筵席已經預備好了,牛和肥畜已經宰了,各樣都齊備,請你們來赴席.』

22:6其餘的拿住僕人,凌辱他們,把他們殺了.

好客的主人遇上敷衍的客人,已經是掃興及可以令主人失預算的事情.更令人困惑的,是客人的激烈的回應.相傳古代戰爭之中,有所謂「兩國交兵,不斬來使」的共識,何況故事中的僕人只是來提醒被邀請的客人往赴婚筵,何以發生「拿住僕人,凌辱他們,把他們殺了」的殘酷事件呢?在王辦喜事的時候,客人乘機起事,超出了赴會與否的考慮,難免令人懷疑這批客人另有目的.

名單客人對隨機客人

22:8於是對僕人說:『喜筵已經齊備,只是所召的人不配.

22:9所以你們要往岔路口上去,凡遇見的,都召來赴席.』

22:10那些僕人就出去,到大路上,凡遇見的,不論善惡都召聚了來,筵席上就坐滿了客.

如果那些「拿住僕人,凌辱他們,把他們殺了」的客人果真藉喜事而起事,國王平亂後,仍然按計劃舉行婚筵,他要與眾同樂,廣宴天下的心意沒有因亂事而取消.只是由原來的有名有姓的名單客人,變成了「岔路口上去,凡遇見的,都召來赴席」的隨機客人.如果故事結束在這裡,仍可算是風波後的平安,仍然有喜慶收場的氣氛,因為20節那句說:「筵席上就坐滿了客」.路加福音敘述類似的故事時,就結束在這裡,但馬太福音卻繼續反高潮的劇情,拒絕大團圓結局.似乎又有人在喜事中起事……

穿禮服客人對不穿禮服客人

22:11王進來觀看賓客,見那裡有一個沒有穿禮服的,

22:12就對他說:『朋友,你到這裡來怎麼不穿禮服呢?』那人無言可答.

22:13於是王對使喚的人說:『捆起他的手腳來,把他丟在外邊的黑暗裡；在那裡必要哀哭切齒了.』

其望大團圓結局的讀者,會感到很惋惜,只因為有一位賓客沒有穿禮服,就破壞了大團圓結局的完美.這位沒穿禮服的客人事實上也很奇特,在當時的猶太人社會,赴宴一般都會穿上禮服,或者比較清潔白淨的衣服,如果自己沒有,就向鄰居借.借禮服赴會並不是甚麼奇怪的事情,在幾十年前的香港仍然可見,穿得端莊是對主人及筵席的尊重.在比喻中,主人曾經詢問,那位客人如果有難處仍可以解釋,然而他卻無言可答.

有人說,那些客人是在大街上臨時被邀請回來的,怎可能有禮服呢?這問題也有道理,然而,如果所有客人都穿了禮服,只有一個人沒有穿,卻可以假設縱使客人是臨時邀請來的,穿禮服仍是可能的.如何可能呢?有聖經學者提出,禮服是由舉辦婚筵的王提供的,所以所有賓客其實都有禮服可穿,只是那位人兄沒穿而己.

四.  總結

耶穌的比喻講完了,你聽得明白嗎?當時有人聽得明白:

21:45祭司長和法利賽人聽見他的比喻,就看出他是指著他們說的.

21:46他們想要捉拿他,只是怕眾人,因為眾人以他為先知.

接著在22:1寫道:「耶穌又用比喻對他們說」,就是剛才我們讀過的那個故事.其實,這是連續三個關於天國比喻的最後一個,第一個記載在28-32節,題為「兩個兒子的比喻」；第二個記載在33-44節,名為「惡園戶的比喻」.耶穌講天國的比喻,目的何在呢?施洗約翰在馬太福音3章2節曾高喊:「天國近了,你們應當悔改!」；耶穌也在第4章17節發出同樣的呼籲:「天國近了,你們應當悔改!」.聰明的祭司長和法利賽人聽得出耶穌的比喻是指著他們說的,但卻不是走上悔改的路,而是:

22:15當時,法利賽人出去商議,怎樣就著耶穌的話陷害他……

故事內充滿張力,其實故事外也充滿張力,耶穌生活在充滿張力的真實世界.在現實世界中,我們不明白為甚麼總會有衝突,為甚麼總有人「搞事」；在天國的比喻中,讀者也不明白為甚麼會有如此奇怪的客人,不但不與喜樂的人同喜樂,反而要在別人辦喜事的時候起事；也不明白為甚麼有個別人士如此怪異,當所有人都穿著禮服去參加婚筵的時候,他偏偏不穿.有時候,我們在想,如果這個世界沒有這些奇怪人士,世界就會和平得多；如果社會中沒有這些好勇鬥狠的人士,社會就能夠繁榮穩定；如果我們的群體中沒有這些破壞者,我們的生活就會舒服得多.誰是在喜事中起事的人呢?

祭師長和法利賽人對號入座,認為自己是比喻中第一批被邀請但拒絕赴會的客人.他們的理解獲得保羅的肯定,使徒行傳記載保羅和巴拿巴如此對猶太人說:

13:46「上帝的道先講給你們原是應當的；只因你們棄絕這道,斷定自己不配得永生,我們就轉向外邦人去.」

今天,當我們坐在這裡聽這個天國比喻的時候,我們又認為自己是甚麼角色呢?相信沒有人自大得認為自己是那位擺設筵席的王吧!或者,我們被教導自己一定是在大街上被邀請而又有穿禮服的第二批客人.讓我們對號入座想像一下:我們坐好了,等待開席,就議論原本被邀請而拒絕赴宴的第一批客人如何反叛,如何粗暴,如何罪有應得,如何不配享受筵席.當筵席要開始時,我們看見有一個沒有穿禮服的人被抽了出來,我們又竊竊私語,當我們都穿上禮服的時候,他竟然沒穿,怪責他破壞喜慶氣氛.

是的,故事中那批答應了邀請而又沒有赴會的客人,失落了享受筵席的機會；那批恩將仇報,凌辱及殺害來使的惡人,也自食惡果；那位沒有穿著禮服的朋友,也被處理了.然而,我們真的如此肯定,自己的位置坐得那麼穩定,那麼有資格去批評不肯赴會的人,怪責沒有穿禮服的人破壞婚筵的歡愉喜慶嗎?有沒有可能--只試探討可能性,我們其實是那些答應過但又拒絕赴會的其中一位呢?會不會我們其實是沒有穿禮服的那位呢?如果,萬一,有這個可能性,會不會我們有份破壞婚筵的歡愉喜慶呢?即是說,會不會帶來張力的其實是我們?衝突發生,我們也有「功勞」呢?

有一位教授最近在網上發表一篇文章(「這樣的信徒信什麼?」鍾劍華),有幾句這樣說:

「很多人返教會,是為了多拿一片葉子來遮掩他們的邪惡本質……所謂宗教信仰,只係抽耶穌水,只係借黃大仙過橋,只係想把其他教友拉落水,只係意圖以宗教來包裝佢地ge 惡行,只係以信徒這個身份來點妝佢地ge 醜陋面貌.只要有好處,毒蛇蠍子也會引用聖經.反正冇損失,何妨偶然向黃大仙上柱香.妄用神之名,過去幾年在香港社會真的是屢見不鮮.而且越來越核突,越來越露骨,越來越不知所謂……」

如此的反省也很破壞氣氛,也可能令人感到張力.不過,如果與感恩歡慶的態度相比,以謙恭自省的態度來閱讀耶穌這個天國的比喻也未必太過離題.從馬太福音敘述這個比喻看來,沒有「同歡唱齊慶賀」的目的,反而充滿審判的味道.面對著審判,有些人的回應可能是「抵佢死!」,或者「殺無赦」；然而,最合適的回應還是聽從施洗約翰及耶穌的呼籲:「天國近了,你們應當悔改!」

最近,聽到一些豪言壯語如此說:「如果佢呢Ｄ咁既人都上到天堂,我寧願落地獄!」其實大家又不用擔心天堂收人太濫,耶穌這個天國的比喻,就是要解釋:「被召的人多,選上的人少」這狀況.但願我們不就是我們所批評的那些人,或者不會變成我們所議論的那些人就好了.



\newpage

\section{高端信徒與低端信徒}
\label{sec:398}
\textbf{林豪恩先生 - 高端信徒與低端信徒}
\newline
\newline
\hyperref[sec:397]{< < < PREV SERMON < < <}
~
\hyperref[sec:toc]{[返目錄]}
~
\hyperref[sec:399]{> > > NEXT SERMON > > >}


哥林多前書 8:1-13
\newline
\begin{longtable}{cl}
\hline
\hline
章節 & 經文 (和合本修訂版)\\
\hline
8:1 & \begin{tabularx}{0.7\textwidth}{X} 關於祭過偶像的食物，我們曉得「我們都有知識」，但知識使人自高自大，惟有愛心能造就人。 \end{tabularx} \\ \\ \relax
8:2 & \begin{tabularx}{0.7\textwidth}{X} 若有人自以為知道甚麼，他其實仍不知道他所應當知道的。 \end{tabularx} \\ \\ \relax
8:3 & \begin{tabularx}{0.7\textwidth}{X} 若有人愛神，他就是神所認識的人了。 \end{tabularx} \\ \\ \relax
8:4 & \begin{tabularx}{0.7\textwidth}{X} 關於吃祭過偶像的食物，我們知道「偶像在世上算不得甚麼」；也知道「神只有一位，沒有別的」。 \end{tabularx} \\ \\ \relax
8:5 & \begin{tabularx}{0.7\textwidth}{X} 雖然在天上或地上有許多所謂的神明，就如他們中間有許多的神明，許多的主， \end{tabularx} \\ \\ \relax
8:6 & \begin{tabularx}{0.7\textwidth}{X} 但是我們只有一位神，就是父，萬物都出於他，我們也歸於他；並只有一位主，就是耶穌基督，萬物都是藉著他而有，我們也是藉著他而有。 \end{tabularx} \\ \\ \relax
8:7 & \begin{tabularx}{0.7\textwidth}{X} 可是，不是人人都有這知識。有人到現在因拜慣了偶像，仍以為所吃的是祭過偶像的食物；既然他們的良心軟弱，也就污穢了。 \end{tabularx} \\ \\ \relax
8:8 & \begin{tabularx}{0.7\textwidth}{X} 其實，食物不能使我們更接近神，因為我們不吃也無損，吃也無益。 \end{tabularx} \\ \\ \relax
8:9 & \begin{tabularx}{0.7\textwidth}{X} 可是，你們要謹慎，免得你們這自由竟成了軟弱人的絆腳石。 \end{tabularx} \\ \\ \relax
8:10 & \begin{tabularx}{0.7\textwidth}{X} 若有人見你這有知識的在偶像的廟裡坐席，而這人的良心是軟弱的，他豈不放膽去吃那祭過偶像的食物嗎？ \end{tabularx} \\ \\ \relax
8:11 & \begin{tabularx}{0.7\textwidth}{X} 因此，基督為他死的那軟弱弟兄，也就因你的知識沉淪了。 \end{tabularx} \\ \\ \relax
8:12 & \begin{tabularx}{0.7\textwidth}{X} 你們這樣得罪弟兄，傷了他們軟弱的良心，就是得罪基督。 \end{tabularx} \\ \\ \relax
8:13 & \begin{tabularx}{0.7\textwidth}{X} 所以，食物若使我的弟兄跌倒，我就永遠不吃肉，免得使我的弟兄跌倒了。 \end{tabularx} \\ \\
[1ex]
\hline
\hline
\end{longtable}
題目:高端與低端信徒   High-end and Low-end Believers

經文:哥林多前書8章1-13節

講員:林豪恩先生

場合:崇基學院禮拜堂主日崇拜

日期:2018年1月28日

一．引言

今天誦讀的書信經課是哥林多前書,哥林多這個城市原來有些地方可以與香港相提並論.

第一,哥林多享有地理優勢,它獨佔從雅典來的東,西陸路與西北,東南的兩灣之間的交匯處.地理優勢有甚麼好處?香港人是明白的.香港旅遊發展局如此介紹:「香港位處中國東南端的珠江口岸,面向南中國海,憑藉其優越的地理位置,以及優良的維多利亞港,得以發展成為一個國際貿易港,轉口港,目前更是亞洲的金融及服務中心.」地理優勢正是當年哥林多成為富庶及繁榮的商業城市的原因之一.

第二,哥林多是一個新城.古哥林多城在主前146年被羅馬大軍摧毀.一個世紀之後,該撒猶流(Julius Caesar)將軍下令重建,成為富有羅馬文化的「新」都市.為重建哥林多,許多羅馬退伍軍人被安置於此地,也有從各地來的自由人(即被釋放的奴隸)與工匠來到此地.一處新發展的地方有甚麼優勢呢?香港人會記得,當市場尚未被瓜分的曾經何時,機會處處,白手興家的故事隨處可聽,人們四面而來尋求迅速致富的機會.在一個相對自由,無傳統文化包袱的地方,人人也有致富的機會,其中不乏女性, 保羅提及的革來氏與非比,很可能都是在哥林多一帶經商致富的女性商人.

以下兩項與香港不同,香港有人說體育界和宗教界沒有經濟貢獻,在哥林多,體育界和宗教界也是經濟的重要來源之一.

大家知道有四年一度的奧林匹克運動會,當年,在哥林多則有二年一度的伊蘇米亞運動會,是古希臘三大泛希臘的運動會之一.哥林多人懂得把運動會逐步擴大,包括對羅馬皇帝的尊崇,詩歌朗誦與音樂比賽,駕馭馬車比賽等,有些項目容許女性參與.伊蘇米亞運動會吸引大量的遊客,為當地創造了大筆的收益.

宗教活動也有助哥林多的經濟.重建後的哥林多雖是羅馬文化的城市,但是傳統希臘的神明,源自埃及的神明,羅馬的國家與帝王崇拜等,繼續香火頂盛.香港也是一處中西神明林立的地方,不過,哥林多人卻懂得把宗教發展為興盛的產業.

商旅運輸,運動盛會,宗教朝聖,大量流動人口,就會帶動旅遊相關行業.當時的住宿問題,多半用活動的帳棚來解決.因此,保羅與他的同工在哥林多製造帳棚是能夠維生的.

保羅在哥林多逗留了18個月,在當地建立了一群基督徒會眾.在一個多種族,多宗教,多元文化,傳統道德束縛相對薄弱,處處都是賺錢機會又充滿競爭,有些人講究享受和放縱肉體,而有些人拘泥律法的傳統,這個基督徒群會遇到不少信仰實踐上的挑戰.香港信徒處身的經濟文化社會狀況,有點似哥林多信徒的處境,在這種狀況下出現的信仰問題,在耶路撒冷以猶太人及猶太文化為主體的地方未必會出現.今天書信經課提及的事件,是在這多元文化處境之下信仰生活的其中一項挑戰.保羅如何回應這困難呢?他的回應又給我們甚麼啟發呢?

二．經文

8:1論到祭偶像之物,我們曉得我們都有知識.但知識是叫人自高自大,惟有愛心能造就人.

8:2若有人以為自己知道甚麼,按他所當知道的,他仍是不知道.

8:3若有人愛神,這人乃是神所知道的.

8:4論到吃祭偶像之物,我們知道偶像在世上算不得甚麼,也知道神只有一位,再沒有別的神.

8:5雖有稱為神的,或在天,或在地,就如那許多的神,許多的主；

8:6然而我們只有一位神,就是父－萬物都本於他；我們也歸於他－並有一位主,就是耶穌基督－萬物都是藉著他有的；我們也是藉著他有的.

8:7但人不都有這等知識.有人到如今因拜慣了偶像,就以為所吃的是祭偶像之物.他們的良心既然軟弱,也就污穢了.

8:8其實食物不能叫神看中我們,因為我們不吃也無損,吃也無益.

8:9只是你們要謹慎,恐怕你們這自由竟成了那軟弱人的絆腳石.

8:10若有人見你這有知識的,在偶像的廟裡坐席,這人的良心若是軟弱,豈不放膽去吃那祭偶像之物嗎?

8:11因此,基督為他死的那軟弱弟兄,也就因你的知識沉淪了.

8:12你們這樣得罪弟兄們,傷了他們軟弱的良心,就是得罪基督.

8:13所以,食物若叫我弟兄跌倒,我就永遠不吃肉,免得叫我弟兄跌倒了.

三．釋經與應用

經文提及的問題是:基督徒可不可以吃祭過偶像的食物呢?這個問題除了出現在哥林多教會之外,也會出現於以其他宗教信仰為主流的地方,例如華人教會.按保羅的描述,信徒群體中對這問題有兩種相反的意思,一批信徒認為可以吃,另一批認為不可以吃.至於哪一批是高端信徒,哪一批是低端信徒,則見仁見智了.當時,也有兩種信仰風格同時在教會中激蕩,其一是希臘尚智文化,另一種是尚靈的熱情.從保羅引述的觀點可見,贊成吃祭過偶像的食物的那批信徒,很可能傾向尚智文化,他們以知識及理性指出吃祭過偶像的食物是沒有問題的.

1．可能自覺高端的尚智派信徒

8:1論到祭偶像之物,我們曉得我們都有知識.

8:4論到吃祭偶像之物,我們知道偶像在世上算不得甚麼,也知道神只有一位,再沒有別的神.

8:5雖有稱為神的,或在天,或在地,就如那許多的神,許多的主；

8:6然而我們只有一位神,就是父－萬物都本於他；我們也歸於他－並有一位主,就是耶穌基督－萬物都是藉著他有的；我們也是藉著他有的.

此派信徒的推理是:

神只有一位.同意吧!既然只有一位,其他被稱為神的,都不真的是神.也同意吧!既然偶像都是假的,就算食物用來祭過偶像,也不會改變食物本身的性質.既然祭過偶像的食物與沒有祭過偶像的食物沒有分別,為甚麼不能吃呢?

保羅如何看這派信徒的理解呢?他是同意的.保羅在字裡行間,把自己加進了這種理解,例如:「8:1論到祭偶像之物,我們曉得我們都有知識.」；「8:4我們知道偶像在世上算不得甚麼」.

既然保羅同意尚智派信徒對吃祭過偶像的食物的理解,決定不是已經有了嗎?以保羅的身份,應該可以一言九鼎,一槌定音了吧!值得留意的是,但保羅卻不是如此處理信徒群體中的分歧.

2．可能被評為低端的尚靈派信徒

保羅坦誠地表示自己在道理上認同尚智派信徒對吃祭過偶像的食物的理解,然而,他卻留意到另一批信徒對這事情的觀感:

8:7但人不都有這等知識.有人到如今因拜慣了偶像,就以為所吃的是祭偶像之物.他們的良心既然軟弱,也就污穢了.

8:10若有人見你這有知識的,在偶像的廟裡坐席,這人的良心若是軟弱,豈不放膽去吃那祭偶像之物嗎?

從字裡行間可見,保羅這段說話仍然是向著贊成吃祭過偶像的食物的信徒說的.

第一,你們有知識去明白「祭過偶像的食物與沒有祭過偶像的食物沒有分別」這道理,但有些人卻沒有這樣的知識去明白你們所明白的.

第二,這些人的經驗塑造了他們的世界,他們成為基督徒之前,拜慣了偶像,無論真假,他們與偶像之間總有一些體驗.經驗是有延續性的,並不容易好像電腦硬盤一樣,可以完全清洗,然後重新載上資料.就算頭腦上知道「祭過偶像的食物與沒有祭過偶像的食物沒有分別」,看著那些祭過偶像的食物,豈能不想起從前拜偶像的狀況嗎?心裡豈能沒有感受嗎?

第三,如果這些被稱為「沒有知識」的信徒,看見「有知識」的信徒如此做,就傾向跟隨「有知識」的信徒進入廟宇吃祭過偶像的食物.如此,他們的行動與他們的經驗產生矛盾,引起內心的不安,也動搖他們的信仰.

在當時的哥林多教會,有些人認為吃祭過偶像的食物是沒有問題的,有些人則認為不應該吃.當信徒群體對某事情在信仰的立場上有相反的意思,如何處理呢?哥林多教會出現這情況,絕對不是獨特的.香港教會在社會中作為少數群體,面對著傳統習俗,或者潮流玩意,也經常會問「基督徒可不可以這樣做」,也經常出現對事情持不同,甚至相反的信仰立場!

最近,一位姓何的前問責局長因為在美國被起訴而成為新聞人物,令人想起他在任時,基督徒群體曾經為他的行動而爭論,題目是:「基督徒可否進廟求籤?」(不幸地當時他求得的是下籤)我們可以想像,以尚智派的觀點來看,進廟求籤算不得甚麼；以尚靈派的觀點來看,基督徒當然不應該參與拜祭偶像的活動.幸好,當年對這事件的爭論並沒有令信徒不和.

對生活行為的不同信仰立場還有很多,例如:「基督徒可不可以打麻雀?」「基督徒可不可以參加萬聖節派對?」「基督徒可不可以穿短褲及小背心參加崇拜?」等「基督徒可以不可以乜乜物物」的爭論.有趣的是,雙方往往都能夠舉出聖經經文和神學論述支持自己的立場.也是幸好,這些不同立場並沒有在信徒群體中做成廣泛而深遠的衝突和分裂.

不過,幾年前因為政改而引發的連串事件,卻令信徒群體出現劇烈的爭論,甚至互相攻擊,幾年後的今天,不知餘波消散了沒有.原來,政治問題對信仰群體的震撼力往往比傳統習俗問題和潮流玩意問題更嚴重.

3．保羅的回應

當年哥林多教會也出現了不少生活與信仰的難題,保羅在哥林多書信中回應了不少.關於可否吃拜祭過偶像的食物的爭論,保羅如此回應:

3．1．知識與愛心(right and kind)

8:1論到祭偶像之物,我們曉得我們都有知識.但知識是叫人自高自大,惟有愛心能造就人.

8:2若有人以為自己知道甚麼,按他所當知道的,他仍是不知道.

8:3若有人愛神,這人乃是神所知道的.

人們傾向把生活所有部份都變成一場辯論比賽,目的是要證明自己對,對方錯,因此,容易造成誓不兩立的對抗局面.保羅引導哥林多信徒從爭論對錯的層次,轉向愛心的層次.這是不容易的,特別是對於自覺有知識人.保羅提醒,當我們以為自己知道很多的時候,其實可能不知的更多.所以,何必一定要證明自己對而對方錯呢?就算有人上知天文下知地理,在上帝看來也算不得甚麼.對上帝來說,甚麼是算數的呢?保羅指出是愛.回應哥林多信徒之間諸多的分歧和紛爭,保羅在13章指出,現今常存的有信有望有愛,其中最大的是愛.保羅並非不評理,但他認為基督徒之間和平共處的鑰匙是愛.在辯論上打敗對手,未必能帶來和平,正如宣傳廣告說:「贏jor 場交,輸jor 個家,值得咩?」保羅對哥林多教會的這項教導,可否成為我們處理信徒之間的矛盾的參考呢?

2.神看與人看

8:8其實食物不能叫神看中我們,因為我們不吃也無損,吃也無益.

8:9只是你們要謹慎,恐怕你們這自由竟成了那軟弱人的絆腳石.

對於爭辯雙方來說,往往為自己的立場引經據典,聲稱是嚴肅的神學或信仰問題,甚至提升至「上帝怎樣看」的層次,因此得出「上帝喜悅這樣做」或者「上帝不喜悅這樣做」的判斷,然後把這判斷用以指責別人.保羅建議爭持不下的哥林多信徒,不如先放下食物,食物其實算不得甚麼,不用把食物提升到那麼高的層次,不用把食物與上帝輕易相連起來.放下手上的食物,就會看見人,看見你的弟兄,你的姐妹,再思想,難道手上的食物比眼前的人更重要嗎?難道維護一知半解的立場比保護眼前的弟兄姊妹更值得嗎?有些爭論,可能並不如我們想像中那麼是非分明的信仰問題,也未必是上帝如何看的問題,只是人與人之間如何相處的問題.回歸人與人之間的相處,或者有另一個角度看事情.

例如有人說:「你穿短褲小背心去崇拜,是不尊重上帝的表現.試想想,你若有機會去見特首,你也會穿得整齊得體,你朝見上帝比見特首穿得更差,豈不是尊重特首比尊重上帝更多嗎?」話好像說得有道理,然而,讓我們退一步想想,上帝真的介意你穿甚至去朝見祂嗎?難道你赤條條洗澡時不也在朝見上帝嗎?真正介意你穿甚麼上教堂的,未必是上帝,而是與你一齊崇拜的弟兄姊妹.不如把問題回歸人與人之間的不同,協商群體能夠接受的共識,而不用凡事「擺上帝上台」.

3．3．基督珍惜,我們不可以輕視

8:11因此,基督為他死的那軟弱弟兄,也就因你的知識沉淪了.

8:12你們這樣得罪弟兄們,傷了他們軟弱的良心,就是得罪基督.

8:13所以,食物若叫我弟兄跌倒,我就永遠不吃肉,免得叫我弟兄跌倒了.

東晉名臣王導曾說:「我不殺伯仁,伯仁卻因我而死」.保羅的底線是,「免得叫我弟兄跌倒」.理由何在呢?是不是我喜歡這弟兄呢?未必是.是不是弟兄與我同一陣線呢?未必是.是不是弟兄對教會有貢獻呢?未必是.保羅提出的理由,與我們個人無關的,甚至與教會的利益和關的,只因為「基督為他死」.「基督為他死」,就是弟兄的價值；「基督為他死」,他就是基督的人,保羅指出,我們傷害基督的人,就等同傷害基督.既然我們從來沒有意思要傷害基督,就避免傷害「基督為他死」的人.

四．總結

贊成吃拜祭過偶像的食物的信徒可能認為自己明白事理,與時並進；可能批評反對者迷信迂腐,落伍離地.在這裡,保羅呼籲有知識的人體諒及接納沒有相同知識的人,保護他們.其實,反對吃拜祭過偶像的食物的信徒,也可能認為自己信仰虔誠,嚴守主道；批評反對者不屬靈,不順服.在互相批評之下,大家也在製造壓力給對方,令教會生活感到辛苦.在第十章,保羅轉過來勸募自覺屬靈的信徒不用對拜祭過偶像的食物太過敏感,基督徒也要尋找能夠生存在多元宗教及文化的社會之道.在哥林多信徒群體之間激蕩著的兩種信仰風格有可能協調而成為信仰資源嗎?保羅自己就是一個好例子,他飽讀經書,明白事理,在猶太和希臘文化之間遊走,同時也有豐富的個人靈性經驗,尚智和尚靈,在他的生命中相輔相承.

高端與低端,是一個建構出來的陷阱,踩進去,損人害己,隨時泥足深陷.保羅建議我們互相視對方為基督所珍惜的人,努力拆除這個陷阱.



\newpage

\section{不能承受生命的速}
\label{sec:399}
\textbf{伍渭文牧師 - 不能承受生命的速}
\newline
\newline
\hyperref[sec:398]{< < < PREV SERMON < < <}
~
\hyperref[sec:toc]{[返目錄]}
~
\hyperref[sec:400]{> > > NEXT SERMON > > >}


約拿書 3:1-5,10
\newline
\begin{longtable}{cl}
\hline
\hline
章節 & 經文 (和合本修訂版)\\
\hline
3:1 & \begin{tabularx}{0.7\textwidth}{X} 耶和華的話第二次臨到約拿，說： \end{tabularx} \\ \\ \relax
3:2 & \begin{tabularx}{0.7\textwidth}{X} 「起來，到尼尼微大城去，把我告訴你的信息向其中的居民宣告。」 \end{tabularx} \\ \\ \relax
3:3 & \begin{tabularx}{0.7\textwidth}{X} 約拿就照耶和華的話起來，到尼尼微去。尼尼微是一座極大的城，約有三天的路程。 \end{tabularx} \\ \\ \relax
3:4 & \begin{tabularx}{0.7\textwidth}{X} 約拿進城，走了一天，宣告說：「再過四十天，尼尼微要傾覆了！」 \end{tabularx} \\ \\ \relax
3:5 & \begin{tabularx}{0.7\textwidth}{X} 尼尼微人就信服神，宣告禁食，從最大的到最小的都穿上麻衣。 \end{tabularx} \\ \\ \relax
3:6 & \begin{tabularx}{0.7\textwidth}{X} 這消息傳到尼尼微王那裡，他就從寶座起來，脫下朝服，披上麻布，坐在灰中。 \end{tabularx} \\ \\ \relax
3:7 & \begin{tabularx}{0.7\textwidth}{X} 他叫人通告尼尼微全城，說：「王和大臣有令，人、畜、牛、羊都不可嘗任何東西，不可吃，也不可喝水。 \end{tabularx} \\ \\ \relax
3:8 & \begin{tabularx}{0.7\textwidth}{X} 人與牲畜都要披上麻布，切切求告神。各人要回轉離開惡道，離棄自己掌中的殘暴。 \end{tabularx} \\ \\ \relax
3:9 & \begin{tabularx}{0.7\textwidth}{X} 誰知道神也許會回心轉意，不發烈怒，使我們不致滅亡。」 \end{tabularx} \\ \\ \relax
3:10 & \begin{tabularx}{0.7\textwidth}{X} 神察看他們的行為，見他們離開惡道，神就改變心意，原先所說要降與他們的災難，他不降了。 \end{tabularx} \\ \\
[1ex]
\hline
\hline
\end{longtable}
詩篇 62:5-12
\newline
\begin{longtable}{cl}
\hline
\hline
章節 & 經文 (和合本修訂版)\\
\hline
62:5 & \begin{tabularx}{0.7\textwidth}{X} 我的心哪，你當默默無聲，專等候神，因為我的盼望是從他而來。 \end{tabularx} \\ \\ \relax
62:6 & \begin{tabularx}{0.7\textwidth}{X} 惟獨他是我的磐石，我的拯救；他是我的庇護所，我必不動搖。 \end{tabularx} \\ \\ \relax
62:7 & \begin{tabularx}{0.7\textwidth}{X} 我的拯救、我的榮耀都在於神；我力量的磐石、我的避難所都在於神。 \end{tabularx} \\ \\ \relax
62:8 & \begin{tabularx}{0.7\textwidth}{X} 百姓啊，要時時倚靠他，在他面前傾心吐意；神是我們的避難所。（細拉） \end{tabularx} \\ \\ \relax
62:9 & \begin{tabularx}{0.7\textwidth}{X} 人真是虛空，人真是虛假；放在天平裡就必浮起，他們一共比空氣還輕。 \end{tabularx} \\ \\ \relax
62:10 & \begin{tabularx}{0.7\textwidth}{X} 不要仗勢欺人，也不要因搶奪而驕傲；若財寶加增，不要放在心上。 \end{tabularx} \\ \\ \relax
62:11 & \begin{tabularx}{0.7\textwidth}{X} 神說了一次、兩次，我都聽見了，就是能力屬乎神。 \end{tabularx} \\ \\ \relax
62:12 & \begin{tabularx}{0.7\textwidth}{X} 主啊，慈愛也是屬乎你，因為你照著各人所做的報應他。 \end{tabularx} \\ \\
[1ex]
\hline
\hline
\end{longtable}
哥林多前書 7:29-31
\newline
\begin{longtable}{cl}
\hline
\hline
章節 & 經文 (和合本修訂版)\\
\hline
7:29 & \begin{tabularx}{0.7\textwidth}{X} 弟兄們，我是說：時候不多了。從此以後，那有妻子的，要像沒有一樣； \end{tabularx} \\ \\ \relax
7:30 & \begin{tabularx}{0.7\textwidth}{X} 哀哭的，不像在哀哭；快樂的，不像在快樂；購買的，像一無所得； \end{tabularx} \\ \\ \relax
7:31 & \begin{tabularx}{0.7\textwidth}{X} 享受這世界的，不像在享受這世界；因為這世界的局面將要過去了。 \end{tabularx} \\ \\
[1ex]
\hline
\hline
\end{longtable}
馬可福音 1:14-20
\newline
\begin{longtable}{cl}
\hline
\hline
章節 & 經文 (和合本修訂版)\\
\hline
1:14 & \begin{tabularx}{0.7\textwidth}{X} 約翰下監以後，耶穌來到加利利，宣講神的福音， \end{tabularx} \\ \\ \relax
1:15 & \begin{tabularx}{0.7\textwidth}{X} 說：「日期滿了，神的國近了。你們要悔改，信福音！」 \end{tabularx} \\ \\ \relax
1:16 & \begin{tabularx}{0.7\textwidth}{X} 耶穌沿著加利利的海邊走，看見西門和西門的弟弟安得烈在海上撒網；他們本是打魚的。 \end{tabularx} \\ \\ \relax
1:17 & \begin{tabularx}{0.7\textwidth}{X} 耶穌對他們說：「來跟從我，我要叫你們得人如得魚一樣。」 \end{tabularx} \\ \\ \relax
1:18 & \begin{tabularx}{0.7\textwidth}{X} 他們立刻捨了網，跟從他。 \end{tabularx} \\ \\ \relax
1:19 & \begin{tabularx}{0.7\textwidth}{X} 耶穌稍往前走，又見西庇太的兒子雅各和他弟弟約翰在船上補網。 \end{tabularx} \\ \\ \relax
1:20 & \begin{tabularx}{0.7\textwidth}{X} 耶穌隨即呼召他們，他們就把父親西庇太和雇工留在船上，跟從了耶穌。 \end{tabularx} \\ \\
[1ex]
\hline
\hline
\end{longtable}
講題:不能承受生命的速-這世界的樣子將要過去  The Unbearable Transience of Life-The Form of the World Is Passing Away

經課:約拿書3:1-5,10;詩篇62:5-12;哥林多前書7:29-31;馬可福音1:14-20

講員:伍渭文牧師

場合:香港中文大學崇基學院禮拜堂

日期:2018年2月4日

我教中學時有一位同事是短跑好手,1967年代表香港參加日本東京世界大學運動會.比賽期間認識很多世界各地的選手,歡樂滿滿,畢生難忘.到了閉幕禮結束時,全場燈光熄滅,頓時感到失落；這一剎那,他感到不能承受生命的速,促發他追求基督信仰,信主後非常熱心.這個聖誕假期,跟很早移民北美一位弟兄聚舊.2008年世界金融風暴,他為了保住下屬的工作,預早宣布完成手頭工程,兩年後便退休.港幣二十五億的醫院工程終於完成了,他是總負責的工程師,等待入伙紙的一天也來到了,公司預備了一個盛大的慶祝會,但他突然情緒低落沒有參加,良久凝固在辦公室的坐椅上.這一刻,要和工作了三十三年的公司告別了,他感到不能承受生命的速.醫生後來跟他說,他患上抑鬱症.

因為有不能承受生命的速,人就說人一世物一世:「我們就吃吃喝喝吧!因為明天要死了.」(林前十五32).人生就像短短的機場鐵路,在香港站上車,那裡很有氣派,名店林立,高聳的國際金融大廈(International Financial Centre)就在那裡.第二個九龍站,圓方那裡也熱鬧,不單有美食名店,將來更接連西九龍高鐵站,這是地下最深也是最昴貴的建築.下一站是青衣,青衣城也有美食名店,但感覺不一樣了,因為下一站是香港國際機場,終站到了.

1 Imminent End將臨的終末 (林前七29-31)

世界短暫,生命永恆.在今天的書信經課,保羅也說時間短速:「弟兄們,我對你們說,時候減少了……這世界的樣子將要過去了.」(林前七29, 31)為何時候減少了?因為基督很快復臨,為上主工作的時間不多了.有些信徒甚至放下工作,在耶穌的升天的橄欖山等候他再來,保羅在帖撒羅尼迦後書三10說:「若有人不肯做工,就不可吃飯.」但請注意,保羅沒有說生命短速,死亡來到,甚麼都沒有了.保羅乃指出世界的短速,世界可以給我們的也是短速.但基督來到世界,為我們的罪釘死在十字架,第三天復活,帶給我們永恆的盼望.從永恆的視角透視,地上的生命何其短暫,有不能承受的速.

對保羅來說,機場站不是終站,他還要上飛機,飛向一個更美的家鄉,他看到更終末的目的地,迎見為他捨命的基督.從機場鐵路的起點到終點,全長35.3公里,但香港機場到洛杉磯,11,640公里.若人生不是從香港站到機場站,到了機場還要飛往洛杉磯,那麼在香港站,九龍站,青衣站的美食和購物並不是那麼重要了.到了機場終站,也不用絕望,因為我們將騰空而起,奔向萬頃窮碧,要到達永久的家鄉.

剛才我提及的工程師弟兄,分享他第一次建造橋的經驗,對人生很有啟發性.他請教有建橋經驗的合作公司的工程師,從完成的橋這端,開始說到起頭.這是比較難也需要好好思考的解釋,他想知道如何能最有效益來到橋的這一端.聯合國第二任秘書(1953-61)長韓馬紹(Dag Hammarskjöld)說得好,若我們不認識死亡,就不曉得如何生活.他說:「我們不需要尋求死亡,死亡會找上門的,但我們要找到往死亡的路,使死亡顯得有意義」 (Do not seek death. Death will find you. But seek the road which makes death a fulfillment) .他又說:「命運本身不能逃避,也不需響往,它不是沒有原委的隱密,因為它說明這世界和人類歷史的軌跡,都有意義.」 (Destiny is something not be to desired and not to be avoided, a mystery not contrary to reason, for it implies that the world, and the course of human history, have meaning.)

生命的催逼感.從上主永恆的國度看世界和地上的生命,世界顯得短暫,生命顯得急速.耶穌比喻天國像一個人,發現地裡埋了寶貝,就變買所有來買下這塊地.天國也像有人發現最好的珍珠,就變賣所有來買這顆珍珠.不能承受生命的速催生了生命的迫切感,召命感.「弟兄們,我對你們說,時候減少了……這世界的樣子將要過去了.」從終未的國度來看,很多以前認為極其重要的事情如婚姻家庭,都變得沒有那麼重要了,人可以結婚,也可以不結婚.保羅因為神國的出現,看事情的視角拉闊了:「時候減少了,從此之後,那有妻子的,要像沒有妻子；哀哭的,要像不哀哭；快樂的,要像不快樂；置買的,要像無有所得；用世物的,要像不用世物,因為這世界的樣子將要過去了.」(林前七29-31)

這世界的樣子將要過去(The form of the world is passing away).保羅說「這世界的樣子將要過去了」是甚麼意思呢?世界的樣子不是上主所造,眼目所見,寄居其上的物質世界.上主所造的世界原是美好,在基督再來時將更新成為新天新地.世界的樣子(the form of the world)指當時控制著人思想的價值觀,人生觀.

(1) 男尊女卑, 一夫多妻,妻子是丈夫的財富,用以傳宗接代；但保羅說夫婦是伴侶(partner).「男子當各有自己的妻子,女子也當有自己的丈夫.丈夫當用合宜之分待妻子,妻子待丈夫也要如此.妻子沒有權柄主張自己的身子,乃在丈夫；丈夫也沒有權柄主張自己的身子,乃在妻子.夫妻不可彼此虧負,除非兩相情願,暫時分房,為了專心禱告方可；以後仍要同房,免得撒旦趁著你們情不自禁引誘你們.」(林前七2-5)

(2) 靈魂肉體對立,肉體是靈魂的囚牢.因為肉體是短暫而且朽壞,所以禁欲,專注永恆的靈魂.非當弔詭,專注養育靈魂,輕忽克制肉體,所以縱欲.他們以為信了基督,就要離開不信的配偶,免得沾污自己的身體.保羅說不可離開未信的配偶,恩典可治愈自然(grace heals nature).「倘若某弟兄有不信的妻子,妻子也情願和他同住,他就不要離棄妻子.妻子有不信的丈夫,丈夫也情願和她同住,她就不要離棄丈夫.因為不信的丈夫就因為妻子成了聖潔,並且不信的妻子因著丈夫成為聖潔.不然,你們的兒女就不潔淨,但如今他們是聖潔的了.」(林前七)

(3) 工作聖俗有別.羅馬帝國的異教信仰是全覆蓋的,市場買到的的肉都拜過偶像.保羅說,偶像算不得甚麼,這些祭祀沒有甚麼魔力影響肉食,但為了愛心不絆倒他人就不吃罷.從事修路築橋的工人也會參與祭祀異教的守護神,我們禮儀(liturgy)一字希臘文是子民(laity)的工作(ergon),作為帝國的子民,要為國家修路築橋,而每個城邦都有守護神祇,所以開工時進行祭祀儀式.售賣肉食或修路築橋的工作,有些信徒不能接受,認為是背棄信仰.保羅說,偶像算不得甚麼,信仰釋放了我們,信仰免除我們被偶像所轄制.

(4) 身份貴賤分明.當時的犬儒學派cynics說:「真實的人,從未生而為奴,也許位降(戰敗被擄)為奴；不真實的人,從未真正自由,始終為奴.」(“True man can never be a slave in nature although he may be in status and that a false man can never be a free man in reality but is always a slave.”  William Barclay, The New Daily Study Bible. The Letter to the Corinthians, 2002, p. 76) 但保羅說:「各人蒙召的時候是甚麼身份,仍要守住這身份.你是作奴隸蒙召的嗎?不要因此憂慮.若能以自由,就求自由更好.因為作奴僕蒙召於主的,就是主所釋放的人；作自由之人蒙召的,就是基督的奴僕.你們是重價買來的,不要作人的奴僕.弟兄們,你們各人蒙召的時候是甚麼身份,仍要守住這身份.」(林前七20-24) 信仰給人心靈的自由,尊貴的身份.信仰不是狂熱,信仰對現實保持清醒(common sense), 若能夠自由,就爭取自由；但信仰看到現實表象看不到的屬靈真相.我們真正的波士是上主,我們是基督的僕人,因此所有工作是神聖尊貴.

(5)國族彼此征戰掠奪.羅馬和平Pax Romana是美麗的口號,真正是鐵蹄殖民,掠奪資源.啟示錄六章四隻馬,首為白馬,騎士手拿著弓代表征戰, 跟著的紅馬騎士手裡拿著刀,奪去太平,是內戰.今天舊約經課約拿書三章提到猶太先知約拿到宿敵亞述國的首都傳道,使他們免被上主審判,是和好共融,雖然約拿開始時很不願意.亞述人出名殘暴,曾滅北國以色列,雜其婚姻,亂其宗教,斷其傳統,使戰敗的以色列人成為後來猶太所厭惡的撒馬利亞人.

小結:這世界的樣子將要過去,因為聖靈像淨化一切的風,靜靜吹襲這世界.基督埋在地下那三天,不動聲色,像春雨一樣,隨風潛入夜,潤物細無聲,恩典已滲透這世界,恩典治愈自然.我們要回應上主的召喚,在不同的崗位,作個得人如魚的漁夫.

2 Immediate Call即時的呼召(可一14-20).

今天福音書可一14-20耶穌呼召門徒:「約翰下監以後,耶穌來到加利利,宣傳神國的福音,說:”日期滿了,神的國近了!你們當悔改,信福音.”耶穌順著加利利海邊走,看見西門和西門的兄弟安得烈在海邊撒網,他們本是打魚的.耶穌對他們說:”來跟從我!我要叫你們得人如得魚一樣.”他們就立刻捨了網,跟從了他.耶穌稍往前行,又見到西庇太的兒子雅各和雅各的兄弟約翰在船上補網.耶穌隨即招呼他們,他們就把父親西庇太和僱工留在船上,跟從耶穌去了.」

耶穌呼召,門徒就立即撇下一切跟隨耶穌.但門徒並沒有放棄擁有的船,也沒有停止搖櫓的技巧.撇下一切,是告別過去,過去是要得魚,現在用來得人.耶穌講道時,很多人圍在海邊,門徒把船開出,離岸不遠來宣道.耶穌要獨處禱告,門徒把他送到對岸.崇基神學院有一位國內從事媒體工作的神學生,她丈夫是中國好聲音編導,太太唸完神學,夫婦製作記錄片,介紹宣教士在華的工作,去年趁改教五百年紀念,製作了馬丁路德生平記錄片,水平很高,也被國內電視台播出.

每個職業都有它的目標,為奴隸的目標是贖身得自由,作漁夫的目標是得魚,每個行業都有它們要得的魚,但保羅教導我們要得人如魚.因為看到終末很快來到,我們把追求魚的目標成為得人的目標.前校長沈祖堯教授曾公開分享,經2003年沙 士疫症,看到生命的脆弱,以後不是看病,乃醫治病人.前律政司長黃仁龍在崇基周會分享,信仰幫助他在一堆堆文件和數字中,看到人的臉龐；他不單處理事,更是公正地幫助人.

所有職業都是得人如魚,因為職業是召命,職業英文vocation一字,源自拉丁文呼召calling,不同的身份或崗位是見証信仰的地方,也是對人顯出愛的地方.所有職業都有愛的元素,養妻活兒是對家庭的愛,促進經濟繁榮使人人就業也是愛,因為失業是痛苦的事.為奴的當然希望得到自由,這是常識理性,但得不到自由也不用氣餒,等候機會,因為我們已被基督釋放,像路德所說:「基督徒是最自由的君王,不受任何的牢制.基督徒也是最順服的僕人,甘心服侍眾人.」路德說,所有工作都是召命,除了一項,就是修道士,因為修院遠離人群.

上主在我們人生不同的季節有不同的身份,這些身份有些是與生俱來的,有些是我們選擇的.我是我父,母親的兒子,我哥哥的弟弟,妹妹弟弟的兄長,我又是我兒子的父親,孫兒的爺爺.但我的工作崗位,是自己選擇的身份.林前七20:「各人蒙召時是甚麼身份,仍是守住這身份.」身份是服侍人的崗位,見証信仰的地方.

這世界的樣子將要過去,因為恩典已把世界改變了,這世界的樣子要過去,不是地震海嘯災難式的突然終止,乃被基督的仁愛與誠信真實改變了,而教會就是改變世界的推手,被召在不同崗位,守住不同的身份,得人如魚.

3 崩壞面前的穩妥Immovable Fortress (詩六十二5-12).

這世界的樣子將要過去,我們的身體也要漸漸過去,慢慢朽壞,歸回塵土.在微信群看到一信息,題目「宮崎駿77歲生日,我的時間不多了」:「1月5號今天,是宮崎駿77歲生日.他早就鬚髮皆白,戴個老花鏡,手握鉛筆,在一堆畫稿面前一坐就是一天,幾十年如此.只是最近,這個老人總是喃喃地念叨一句話:“我的時間不多了.”這句話,讓人聽了心裡突然一疼.原來那個製造童話的人,原來那個小心翼翼保護了我們天真的人,原來那個用力溫暖了全世界的人,也是會老去的啊.歲月還是沒有饒過你.」

歲月沒有饒過我的是外在的肉體,但內心的自由不能奪去.保羅稍後給哥林多教寫了另一封信:「所以,我們不喪聸.外體雖然毀壞,內心卻一天新似一天.我們這至暫至輕的苦楚,要為我們成就極重無比永遠的榮耀.原來我們不是願念所見的,乃是願念所不見的；因為所見的是暫時的,所不見的是永遠的.」(林後五16-18)看不見的比看得見的更重要,也是顯現期的提醒.

今天是顯現期第五主日.顯現期(Epiphany)的其中一個表記是東方賢士看見巨星來朝拜,他們首先找到王宮的希律王,但新生王不是表象所說必居王宮的,表象的顯現可以有欺哄性.賢士立刻調較新方向,像衛星導航重新調較(recalculating),最後找到臥在馬槽的新生王.之後,他們很有智慧,沒有回到希律王宮,改道而行,避過兇險.

在人生的征途上,我們遇到很多表面看來(Epiphany)不明白的橫逆,信心受到試驗.有時行經死蔭的幽谷,信仰失去焦點,這時候我們需調較方向,察驗上主的路徑.這段時間,我認識的朋友中,不少身體軟弱,先有突然崩壞的震驚,跟著是長期潰敗的磨難,生命像將倒下來的牆.詩篇經課詩六十二5-12,給我們面對崩壞的穩妥；上主是磐石和高台臺,我必不動搖.前面3,4節提及大家攻擊一人,把他毀壞,如同毀壞歪斜的牆,將倒的壁.但全詩重複三次「惟獨他是我的磐石,我的高台,我必不動搖」(2,6,7)

上主是高台,是崩壞面前的穩妥,對香港人也是適時的安慰.2018年香港電台首次城市論壇,回顧2017及前瞻2018政經大事,結束時主持問每位出席嘉賓新一年的展望,一位說:「我不能希望明天會更好,祇希望明天不會更壞,壞的時候也不要太快.」我留意到社交媒體出現不少香港陳年舊事的視頻,「東方之珠」,「獅子山下」MTV大受歡迎；不少人嘆息已往的制度漸漸潰敗,香港變得跟以前不一樣了,香港的底氣漸漸消失,我用「底氣」這詞大家應明白了.我想,若教會此刻能活出信仰,守住身份,香港就不會那麼快崩潰.

結論:有人問愛因斯坦,你最想問的一個問題是甚麼?你最想知道的是甚麼?他說:「我想知道這浩瀚的宇宙,是否對人類友善.」相對整個宇宙,人類太渺少了,生命也太脆弱了,若是宇宙對人類不友善,禍哉!我們滅亡了.但感謝上主,宇宙之主,已降生為人,為我們的罪被釘在十字架上,第三天復活.我們不能承受生命的速,因為時候少了,這世界的樣子要過去.但生命不是隨死亡而結束,機場鐵路站不是我們的終站,我們還要一飛沖天,遨遊那萬頃窮碧,到達那永久的家鄉.這世界和其上的都要過去,讓我們守住身份,得人如魚,這是上主給我們的召命,也是存到永遠的福音果子.



\newpage

\section{預你一份}
\label{sec:400}
\textbf{鄧瑞強博士 - 預你一份}
\newline
\newline
\hyperref[sec:399]{< < < PREV SERMON < < <}
~
\hyperref[sec:toc]{[返目錄]}
~
\hyperref[sec:401]{> > > NEXT SERMON > > >}


啟示錄 6:1-17
\newline
\begin{longtable}{cl}
\hline
\hline
章節 & 經文 (和合本修訂版)\\
\hline
6:1 & \begin{tabularx}{0.7\textwidth}{X} 我看見羔羊揭開七個印中第一個印的時候，聽見四活物中的一個活物，聲音如雷，說：「你來！」 \end{tabularx} \\ \\ \relax
6:2 & \begin{tabularx}{0.7\textwidth}{X} 我就觀看，看見一匹白馬，騎在馬上的拿著弓，並有冠冕賜給他。他出來征服，勝而又勝。 \end{tabularx} \\ \\ \relax
6:3 & \begin{tabularx}{0.7\textwidth}{X} 羔羊揭開第二個印的時候，我聽見第二個活物說：「你來！」 \end{tabularx} \\ \\ \relax
6:4 & \begin{tabularx}{0.7\textwidth}{X} 就另有一匹馬出來，是紅色的；有權柄賜給了那騎馬的，要從地上奪去太平，使人彼此相殺；他又接受了一把大刀。 \end{tabularx} \\ \\ \relax
6:5 & \begin{tabularx}{0.7\textwidth}{X} 羔羊揭開第三個印的時候，我聽見第三個活物說：「你來！」我就觀看，看見一匹黑馬；騎在馬上的，手裡拿著天平。 \end{tabularx} \\ \\ \relax
6:6 & \begin{tabularx}{0.7\textwidth}{X} 我聽見在四個活物中似乎有聲音說：「一個銀幣買一升麥子，一個銀幣買三升大麥；油和酒不可糟蹋。」 \end{tabularx} \\ \\ \relax
6:7 & \begin{tabularx}{0.7\textwidth}{X} 羔羊揭開第四個印的時候，我聽見第四個活物說：「你來！」 \end{tabularx} \\ \\ \relax
6:8 & \begin{tabularx}{0.7\textwidth}{X} 我就觀看，看見一匹灰色馬；騎在馬上的，名字叫作「死」，陰間也隨著他；有權柄賜給他們，可以用刀劍、饑荒、瘟疫、野獸，殺害地上四分之一的人。 \end{tabularx} \\ \\ \relax
6:9 & \begin{tabularx}{0.7\textwidth}{X} 羔羊揭開第五個印的時候，我看見在祭壇底下有曾為神的道，並為作見證而被殺的人的靈魂， \end{tabularx} \\ \\ \relax
6:10 & \begin{tabularx}{0.7\textwidth}{X} 大聲喊著說：「神聖真實的主宰啊，你不審判住在地上的人，為我們所流的血伸冤，要到幾時呢？」 \end{tabularx} \\ \\ \relax
6:11 & \begin{tabularx}{0.7\textwidth}{X} 於是有白袍賜給他們各人；又有話吩咐他們還要歇息片刻，等到與他們同作僕人的，和他們的弟兄，像他們一樣被殺的人的數目湊足的時候。 \end{tabularx} \\ \\ \relax
6:12 & \begin{tabularx}{0.7\textwidth}{X} 羔羊揭開第六個印的時候，我看見地大震動，太陽變黑像粗麻布，整個月亮變紅像血， \end{tabularx} \\ \\ \relax
6:13 & \begin{tabularx}{0.7\textwidth}{X} 天上的星辰墜落在地上，如同無花果樹被大風搖動，落下未熟的果子一樣。 \end{tabularx} \\ \\ \relax
6:14 & \begin{tabularx}{0.7\textwidth}{X} 天就裂開，好像書卷被捲起來；山嶺海島都被移動離開原位。 \end{tabularx} \\ \\ \relax
6:15 & \begin{tabularx}{0.7\textwidth}{X} 地上的君王、臣宰、將軍、富戶、壯士，和一切為奴的、自主的，都藏在山洞和巖石穴裡， \end{tabularx} \\ \\ \relax
6:16 & \begin{tabularx}{0.7\textwidth}{X} 向山和巖石說：「倒在我們身上吧！把我們藏起來，躲避坐寶座者的臉面和羔羊的憤怒； \end{tabularx} \\ \\ \relax
6:17 & \begin{tabularx}{0.7\textwidth}{X} 因為他們遭憤怒的大日子到了，誰能站得住呢？」 \end{tabularx} \\ \\
[1ex]
\hline
\hline
\end{longtable}
講題:預你一份　Count You In

經文:啟示錄6章1-17節

講員:鄧瑞強博士

場合:香港中文大學崇基學院禮拜堂

日期:2018年2月11日

各位弟兄姊妹,早安.農曆新年快到,預祝各人新年快樂.

上個月,《原子科學家公報》(Bulletin of the Atomic Scientists)將「末日鐘」(Doomsday Clock)撥快半分鐘,距離人類滅亡的午夜,如今只差兩分鐘.當然,這是一種象徵講法,卻又實實在在的提示我們,核軍備競賽及環境災難造成的危機,近在眼前.

啟示錄本身就是一「末日鐘」,展示危機,展示政治危機,經濟危機,人性沉淪的危機,因著罪惡深不可測的力量,這些危機已將人類推向生死存亡之秋.啟示錄的「末日鐘」,已撥到午夜前的一秒鐘.真理的大廈,美善的國度,神聖的領域,看來都在傾覆的前夕.就在這午夜前的一秒,啟示錄言說盼望.

今日的講道經文:啟示錄6:1－17

啟6:1　我看見羔羊揭開七印中第一印的時候,就聽見四活物中的一個活物,聲音如雷,說:「你來!」

啟6:2　我就觀看,見有一匹白馬；騎在馬上的,拿著弓,並有冠冕賜給他.他便出來,勝了又要勝.

啟6:3　揭開第二印的時候,我聽見第二個活物說:「你來.」

啟6:4　就另有一匹馬出來,是紅的,有權柄給了那騎馬的,可以從地上奪去太平,使人彼此相殺；又有一把大刀賜給他.

啟6:5　揭開第三印的時候,我聽見第三個活物說:「你來!」我就觀看,見有一匹黑馬；騎在馬上的,手裡拿著天平.

啟6:6　我聽見在四活物中似乎有聲音說:「一錢銀子買一升麥子,一錢銀子買三升大麥；油和酒不可糟蹋.」

啟6:7　揭開第四印的時候,我聽見第四個活物說:「你來!」

啟6:8　我就觀看,見有一匹灰色馬；騎在馬上的,名字叫作死,陰府也隨著他；有權柄賜給他們,可以用刀劍,饑荒,瘟疫〔或譯:死亡〕,野獸,殺害地上四分之一的人.

啟6:9　揭開第五印的時候,我看見在祭壇底下,有為神的道,並作見證被殺之人的靈魂,

啟6:10　大聲喊著說:「聖潔真實的主啊,你不審判住在地上的人,給我們伸流血的冤,要等到幾時呢?」

啟6:11　於是有白衣賜給他們各人；又有話對他們說,還要安息片時,等著一同作僕人的和他們的弟兄也像他們被殺,滿足了數目.

啟6:12　揭開第六印的時候,我又看見地大震動,日頭變黑像毛布,滿月變紅像血,

啟6:13　天上的星晨墜落於地,如同無花果樹被大風搖動,落下未熟的果子一樣.

啟6:14　天就挪移,好像書卷被捲起來；山嶺海島都被挪移離開本位.

啟6:15　地上的君王,臣宰,將軍,富戶,壯士,和一切為奴的,自主的,都藏在山洞和巖石穴裡,

啟6:16　向山和巖石說:「倒在我們身上吧!把我們藏起來,躲避坐寶座者的面目和羔羊的忿怒；

啟6:17　因為他們忿怒的大日到了,誰能站得住呢?」

今日的講道,接續我之前的講道.經文這裡講到的情景,仍是天上的情景.天父的右手,拿著被七印封住的書卷,這書卷藏著神對歷史的終極心意.當「印」被逐一揭開時,歷史的意義便一步一步顯明了.誰配揭開這些印?是「被殺的羔羊」耶穌.一個順從天父之命,在人間活出大愛的人,他配揭開歷史之謎.

啟6:1　我看見羔羊揭開七印中第一印的時候,就聽見四活物中的一個活物,聲音如雷,說:「你來!」

啟6:2　我就觀看,見有一匹白馬；騎在馬上的,拿著弓,並有冠冕賜給他.他便出來,勝了又要勝.

第一印揭開了,有騎白馬的戰士出場.過去的釋經,通常指這是戰爭的力量.現代的釋經,有指這是基督的出場.在啟示錄19章,再出現騎白馬者,那明顯是指基督.這裡提及「雷」,「雷」在啟示錄,是指神的臨在.再者,這裡提及「冠冕」,在啟示錄,有兩個字可譯作「冠冕」.一是這裡的「冠冕」,希臘文是「stephanos」,是指比賽得勝者可戴的「桂冠」.耶穌受難時戴的「荊棘冠冕」,那「冠冕」也是用這個希臘文.另一個用字是「diadēma」,是強權者顯示權力的冠冕.這裡,騎白馬者不是戴這種權力冠冕,而是戴比賽得勝者的「桂冠」.綜合以上觀點來看,經文這裡,騎白馬者是指得勝的基督.揭開第一印,揭開歷史的謎,帶領著歷史前進的,在歷史的複雜比賽中終歸得勝的,是基督.

第二印揭開了(6: 3-4),來了騎紅馬的騎士.明顯,這代表戰爭的殺戮.

第三印揭開了(6: 5-6),來了騎黑馬的騎士,他手裡拿著天平.有聲音說:「一錢銀子買一升麥子,一錢銀子買三升大麥；油和酒不可糟蹋」.明顯,這裡提及經濟的危機.「一錢銀子」,是當時人一日的工資.「一升」糧食的分量,是當時人一日食用的分量.自己賺錢自己食,可買小麥.賺錢養活一家,便要買較粗糙的大麥了.人的生活,僅能糊口.「油和酒不可糟蹋」,一般中文聖經的翻譯,都譯作「不可糟蹋」,但一般英文聖經的翻譯,都譯作「do not damage/harm」,這翻譯準確一點.重點不在於資源匱乏,故不要蹧蹋.重點在不要因大麥和小麥價格高昂而只種大麥小麥,而不再種橄欖取油和不種葡萄取酒.有些地方,最適宜種橄欖；有些地方,最適宜種葡萄,這是最佳的本土經濟,也能維持經濟的多樣化.取消多樣化,經濟單一,謀求最佳利潤,最後經濟會陷入危機.

第四印揭開了(6: 7-8),來了匹「死灰色」的馬,坐在其上的,是「死亡」.這代表政治的殘暴,經濟的困局,邪惡的力量結合起來所帶來的悲慘局面,死人無數.

表面上,這與香港的情況相去甚遠.事實上,的確是相去甚遠,我們也不應渲染「末日意識」.啟示錄要講的,是直指引發一切危機背後的核心元素.啟示錄讓我們看見,政治的危機,核心根源是暴力的無限量使用.「強姦」是暴力的原型(archetype).有力者,強暴者,否定對方的意願,視對方為物件,一件沒有人性尊嚴的物件,一件可被任意蹂躪的物件.這物件被粗暴使用,只為滿足自己卑劣的慾望.在暴力的世界裡,人性被物化,被工具化,用來滿足強暴者的目的,在暴力裡,任何人都可以被暴力ＤＱ,任何人都可以被暴力封口,然後,任何人都只能滿足於被餵飼的,豬一樣的生活.

啟示錄也讓我們看見,經濟的危機,源自人的生活僅能糊口,及對經濟多樣化的摧毀.我不知道香港現時的經濟結構算不算是「地產霸權」,但看來小商舖的營利會被高昂的租金掠奪,小巿民的收入會被高昂的住屋租金所蠶食.年輕人,看見賺的錢,都花在住屋上,當然很焦躁.僅能糊口的生活,使生活死氣沉沉.

政治的壓力,和經濟的壓力,帶來「心」的死亡.哀莫大於心死.心死了,生活便是覓食與享樂.這是崩壞的前夕.

人總不想如此生活.在人心底裡,總渴望真理,仁愛,神聖的意義.我們總渴望一種政治,是「為人民服務」的.我們總渴望一種經濟,是人人共享的.我們總渴望一個世界,在那裡,人是面向神聖而活的.在歷史的巨輪前面,渴想這些東西的人,像擋車的螳螂,被壓在無情的巨輪下.

啟6:9　揭開第五印的時候,我看見在祭壇底下,有為神的道,並作見證被殺之人的靈魂,

啟6:10　大聲喊著說:「聖潔真實的主啊,你不審判住在地上的人,給我們伸流血的冤,要等到幾時呢?」

啟6:11　於是有白衣賜給他們各人；又有話對他們說,還要安息片時,等著一同作僕人的和他們的弟兄也像他們被殺,滿足了數目.

有學者說,這三節經文,是啟示錄的「鑰節」.這裡提及殉道士,為真理而死的人,他們為人類背負起歷史的十架,為高貴的價值不惜犧牲自己.他們在神那裡,在天上,在祭壇下,他們哀求,他們祈禱.這是啟示錄唯一的哀求祈禱.他們代表為真理而死的一切人,向神發出一個天問.「聖潔真實的主啊,你不審判住在地上的人,給我們伸流血的冤,要等到幾時呢?」特別留意,這裡稱神為「聖潔真實的」(holy and true),凸顯這世界的「邪惡不潔」與「虛謊造假」.殉道者問:這麼多人在暴力下受苦受死,神為何不主持正義?世界如此虛偽,不講道理,無視真相,天國何時降臨?

神的回應有點令人驚訝:「還要安息片時,等著一同作僕人的和他們的弟兄也像他們被殺,滿足了數目.」

天國未臨,因為為了堅持真理而不惜犧牲性命的人,數目不夠.或者說,當講真話的人不夠時,講真話的國度怎會出現?當堅持神聖價值的人不足時,神聖的國度怎會降臨?

不要問神為何不主持正義,要問問自己,你是否活得正義.不要問天國為何不早日降臨,要問問自己,你是否為天國的降臨盡上了自己的責任.在歷史的終點前夕,你活了一個怎麼樣的人生?

這又回到信仰裡最基本的問題,我們為什麼而生,為什麼而死?何謂生,何謂死?為真理而死的人,是找死嗎,抑或是找到永恆的生?在這歷史時刻,我們是無關宏旨,抑總有我們卑微的角色?在不潔面前,如何堅持聖潔?在虛假前面,如何堅持真理?在這歷史的瞬間的背後,誰在統合歷史而對歷史中每個當負責任的人作出審問?我們在這刻,當思考什麼,預備什麼,作什麼?

或許,有一天,大話會變成時尚,虛假會成為主流.在那個時刻,言說真理便需要殉道的精神了.論到「謊言」,時事評論家李怡引用前蘇聯諾貝爾獎得主索忍尼辛(Solzhenitsyn)的一段話:「我們知道他們在說謊,他們也知道自己是說謊,他們也知道我們知道他們在說謊,我們也知道他們知道我們知道他們說謊,但是他們依然在說謊」.然後,李怡評論說:「上引的一段話,可以說是他(索忍尼辛)對執政者說謊的經典論述,富文學性,既超現實又極真實.政治說謊者的確如他所說,明知是謊言並且人人知道是謊言,仍然要把謊話不停說下去.不是要人相信,自己也不相信,純粹出於掌權的需要.」(引自:\url{http://app1.rthk.hk/elearning/1minreading/index.php?director=1 and page=151})

暴力需要謊言去支撐,謊言令暴力變得可信.這是歷史崩壞前的光景,啟示錄帶我們去到這真理將亡之際,然後,讓我們好好想想,生活的真正意義.

要完成一幅拼圖,每塊小塊都不能少.神的歷史圖畫的完成,也需要有足夠的小塊去完成.這圖畫有很多空格,有很多遺漏,小塊的數目不夠,完成不了.你就是其中一塊小塊,神預了你的一份.你願意成為神的歷史圖畫裡的其中一塊嗎?

最後,第六印揭開了,是歷史終點的前夕,神主持正義的時刻到了.不正義與虛謊者,面對神的審判.人能配合神的旨意,成為真理的守護者,而不淪為神的正義追討的對象,多好!

但願:榮耀歸於聖父,聖子,聖靈.阿們.



\newpage

\section{悔改靈洗信福音}
\label{sec:401}
\textbf{鄭漢文博士 - 悔改,靈洗,信福音}
\newline
\newline
\hyperref[sec:400]{< < < PREV SERMON < < <}
~
\hyperref[sec:toc]{[返目錄]}
~
\hyperref[sec:402]{> > > NEXT SERMON > > >}


馬可福音 1:9-15
\newline
\begin{longtable}{cl}
\hline
\hline
章節 & 經文 (和合本修訂版)\\
\hline
1:9 & \begin{tabularx}{0.7\textwidth}{X} 那時，耶穌從加利利的拿撒勒來，在約旦河裡受了約翰的洗。 \end{tabularx} \\ \\ \relax
1:10 & \begin{tabularx}{0.7\textwidth}{X} 他從水裡一上來，就看見天裂開了，聖靈彷彿鴿子降在他身上。 \end{tabularx} \\ \\ \relax
1:11 & \begin{tabularx}{0.7\textwidth}{X} 又有聲音從天上來，說：「你是我的愛子，我喜愛你。」 \end{tabularx} \\ \\ \relax
1:12 & \begin{tabularx}{0.7\textwidth}{X} 聖靈立刻把耶穌催促到曠野裡去。 \end{tabularx} \\ \\ \relax
1:13 & \begin{tabularx}{0.7\textwidth}{X} 他在曠野四十天，受撒但的試探，並與野獸同在一起，且有天使來伺候他。 \end{tabularx} \\ \\ \relax
1:14 & \begin{tabularx}{0.7\textwidth}{X} 約翰下監以後，耶穌來到加利利，宣講神的福音， \end{tabularx} \\ \\ \relax
1:15 & \begin{tabularx}{0.7\textwidth}{X} 說：「日期滿了，神的國近了。你們要悔改，信福音！」 \end{tabularx} \\ \\
[1ex]
\hline
\hline
\end{longtable}
講題:悔改,靈洗,信福音　Repent, Be Spirited and Believe

經文:馬可福音1:9-15

講員:鄭漢文博士

場合:香港中文大學崇基學院禮拜堂

日期:2018年2月18日(年初三)

「耶穌出道」的敘事(可1:9-15)

弟兄姊妹,今天是年初三,首先跟大家拜年,祝大家出入平安,靈力漸長.

今日是年初三,中國人叫赤口,即係容易因言語不合而生爭吵.

去年2017年11月19日我在此講道時解讀馬太福音25章14-30節,講及「財富,才幹與在地若天」(Wealth, Talents and the World as Heaven),之後被評為「顛覆性解經」.

所以,我今日講道會好小心,對今日的經課,我採取一種好傳統派的解經.單看題目:「悔改,靈洗,信福音」,今日之前我已經被人評論:「哇,好傳統.仲有咩好講?」

今日,我們重新閱讀「耶穌出道」的敘事,是記載在馬可福音1章9至15節.請聽我讀出:

9那時,耶穌從加利利的拿撒勒來,在約旦河裡受了約翰的洗.

10他從水裡一上來,就看見天裂開了,聖靈仿佛鴿子降在他身上.

11又有聲音從天上來,說:“你是我的愛子,我喜愛你.”

12聖靈立刻把耶穌催促到曠野裡去.

13他在曠野四十天,受撒旦的試探,並與野獸同在一起,且有天使來伺候他.

14約翰下監以後,耶穌來到加利利,宣講神的福音,

15說:“日期滿了,神的國近了.你們要悔改,信福音!”

短短七節,可以分為三個重點,在其他福音書都有記載:

馬可1章第9至11節講:「耶穌受洗」.(太3:13-17；路3:21-22；約1:19-34)

馬可1章第12至13節講:「耶穌受試探」.(太4:1-11；路4:1-13)

馬可1章第14至15節講:「耶穌開始傳道」.(太4:12-17；路4:14-15)

我們今天用「四福音平行經文合參」的方法,即係「以經解經」,去思考這三個重點.

新約聖約研究有一個「雙源假說」(Two-Source Hypothesis),認為馬可福音是最早完成的福音書,另加一個失傳的口述記錄Q來源,此二者是為馬太和路加二福音書所參考.

也有研究發現,馬可福音是馬可記錄使徒彼得口述寫成的,所以有叫做:彼得福音.

今天不作誰先誰後的學術爭論.但可以看出,這短短七節的經文看來是一個最簡單的記述,然後其他三卷福音把相關的細節擴展開來,我們可以借助四卷福音書對讀出今日經課的意義.

(所以,我們會翻查四段經文,請大家準備定用四個書籤把四段經文標示出來.)

從「耶穌受洗」讀出「靈洗」的第一重意義:確認身分(約1:19-35參讀)

首先是今日經課馬可福音第9至11節所講的「耶穌受洗事件」,其他三卷福音書都有記載.

路加福音在第3章21-22節,用兩節記載跟馬可福音差不多一樣的文字.

馬太福音在第3章13-17節,用五節記載,比馬可福音多了以下一段對話(讓我讀出):

13當下,耶穌從加利利來到約旦河,見了約翰,要受他的洗.

14約翰想要攔住他,說:「我當受你的洗,你反倒上我這裡來嗎?」

15耶穌回答說:「你暫且許我,因為我們理當這樣盡諸般的義.」於是約翰許了他.

約翰與耶穌這段對話,究竟二人在拉扯甚麼,爭議甚麼?

耶穌請約翰為他施洗,提供的是約旦河的水洗.

約翰想要阻止耶穌,說:“我應該受你的洗,你怎麼到我這裡來呢?”

耶穌可以提供甚麼的洗?是比約翰的水洗更應領受?

約翰與耶穌二人是甚麼關係?

究竟約翰是「邊一位」?

約翰福音在第1章19-35節,有較詳盡的解說.有標題命之為:「施洗約翰的見證」,用的是約翰的視點,去敘述耶穌受洗的前因後果以及其中過程.

讓我們對讀讀出約翰福音在第1章19-35節這十多節經文,我讀單節,會眾請讀雙節:

施洗約翰的見證(約翰福音1章19-35節；同時可見於:太3:1-12；可1:1-8；路3:1-18)

19約翰所作的見證記在下面:猶太人從耶路撒冷差祭司和利未人到約翰那裡,問他說:「你是誰?」

20他就明說,並不隱瞞,明說:「我不是基督.」

21他們又問他說:「這樣,你是誰呢?是以利亞嗎?」他說:「我不是.」「是那先知嗎?」他回答說:「不是.」

22於是他們說:「你到底是誰,叫我們好回覆差我們來的人.你自己說,你是誰?」

23他說:「我就是那在曠野有人聲喊著說:『修直主的道路』,正如先知以賽亞所說的.」

24那些人是法利賽人差來的；

25他們就問他說:「你既不是基督,不是以利亞,也不是那先知,為甚麼施洗呢?」

26約翰回答說:「我是用水施洗,但有一位站在你們中間,是你們不認識的,

27就是那在我以後來的,我給他解鞋帶也不配.」

28這是在約旦河外伯大尼,約翰施洗的地方作的見證.

29次日,約翰看見耶穌來到他那裡,就說:「看哪,神的羔羊,除去世人罪孽的!

30這就是我曾說:『有一位在我以後來,反成了在我以前的,因他本來在我以前.』

31我先前不認識他,如今我來用水施洗,為要叫他顯明給以色列人.」

32約翰又作見證說:「我曾看見聖靈,彷彿鴿子從天降下,住在他的身上.

33我先前不認識他,只是那差我來用水施洗的,對我說:『你看見聖靈降下來,住在誰的身上,誰就是用聖靈施洗的.』

34我看見了,就證明這是神的兒子.」

以上用的是約翰的視點,去敘述耶穌受洗的事件,是約翰對耶穌受洗的見證.這個見證的重點何在?

回到今日經課馬可福音第1章9節前的(馬可福音)第7-8節「施洗約翰的見證」:

7他傳道說:「有一位在我以後來的,能力比我更大,我就是彎腰給他解鞋帶也是不配的.

8我是用水給你們施洗,他卻要用聖靈給你們施洗.」

按馬可福音第1章4節,約翰所提供的只是水洗,是「悔改的洗禮,使罪得赦」,按這施洗實踐者的第一身見證,受他水洗的耶穌其實比他能力大,廣東話講法是「佢仲勁」,因為耶穌是用聖靈給你們施洗(baptizing in the Spirit),看約翰福音在第1章32-33節:

32約翰又作見證說:「我曾看見聖靈,彷彿鴿子從天降下,住在他的身上.

33我先前不認識他,只是那差我來用水施洗的,對我說:『你看見聖靈降下來,住在誰的身上,誰就是用聖靈施洗的.』

耶穌在接受施洗約翰水洗的同時,接受了靈洗(baptized by the Spirit; baptism of the Spirit),在靈洗(be Spirited)的同時,耶穌得到從上而來的身分確定,由聖父親自認定為神的兒子.

作為神的兒子,耶穌可以提供怎樣的靈洗?

從「耶穌受試探」讀出「靈洗」的第二重意義:超越試探(太4:1-11參讀)

聖靈的洗禮包括接受試探.馬可福音第1章12至13節記載耶穌在靈洗後,

12聖靈立刻把耶穌催促到曠野裡去.

13他在曠野四十天,受撒旦的試探,並與野獸同在一起,且有天使來伺候他.

馬可福音用了非常精簡的方式來描述這段四十天的經歷,包括了「時,空,物,性」:

時間(Time):四十天

空間(Space):曠野

人物(Character):耶穌以外,有撒旦,有野獸,有天使,還有聖靈.

本性(Nature):有人性,獸性,魔性,還有神性.

馬太福音第4章1至11節以及路加福音第4章1至13節都對耶穌受試探這事件作了較詳盡的描述.

但是,我今天不會進入這段經文的討論,我只會讀出馬太福音的版本,請會眾在聽我讀出的同時,借助耶穌的經歷,默想自己曾經經歷近似的試探,在試探中如何經歷從聖靈而來的心靈洗淨(my spirit being purified through the Spirit).馬太福音4章1至11節這樣記載:

1當時,耶穌被聖靈引到曠野,受魔鬼的試探.

2他禁食四十晝夜,後來就餓了.

3那試探人的進前來,對他說:「你若是神的兒子,可以吩咐這些石頭變成食物.」

4耶穌卻回答說:「經上記著說:人活著,不是單靠食物,乃是靠神口裡所出的一切話.」

5魔鬼就帶他進了聖城,叫他站在殿頂上,

6對他說:「你若是神的兒子,可以跳下去,因為經上記著說:主要為你吩咐他的使者用手托著你,免得你的腳碰在石頭上.」

7耶穌對他說:「經上又記著說:『不可試探主－你的神.』」

8魔鬼又帶他上了一座最高的山,將世上的萬國與萬國的榮華都指給他看,

9對他說:「你若俯伏拜我,我就把這一切都賜給你.」

10耶穌說:「撒旦,退去吧!因為經上記著說:當拜主－你的神,單要事奉他.」

11於是,魔鬼離了耶穌,有天使來伺候他.

如何戰勝試探,就是靈洗的必要部份,我們可以如何超越試探?

經文告訴我們:三次試探,耶穌都是引述神的話語,頂住魔性之攻心,找到超越之出路.

中國神學研究院前院長余達心牧師(Rev. Carver Yu)在他的神學導論《聆聽:神學言說的開端》中曾經對神學作一解釋:「神學就是聆聽神」(theology as listening to God).

靈修又是甚麼?靈修就是不受世間雜音所困而得以聆聽神,聆聽神向心靈深處說話?

從聖經,我們可以讀到歷代信徒聆聽神聲音而寫下的話.

聖靈的洗禮,不單單是在某時某地發生的一項事件(event),而是在晝夜思想神的話下的生命操練,往往是在經歷種種試煉以及戰勝種種試探中靈力不斷成長(growth).

從「耶穌傳道」讀出「靈洗」的第三重意義:神國靠近(路3:23-38參讀)

現在進入今日經課的第三個重點.

馬可福音第1章14-15節這樣敘述耶穌受試探後耶穌開始傳道:

14約翰下監以後,耶穌來到加利利,宣講神的福音,

15說:“日期滿了,神的國近了.你們要悔改,信福音!”

洗禮之後,其實就是要接力,前行者交下了的棒,後來者要接棒.

施洗約翰下監的記載我不多講了.之後被害死的事件也不提了.

施洗者已退下,領洗者接棒,耶穌開始傳道.

究竟耶穌所傳的道與施洗約翰所傳的道又何不同?

按馬可福音第1章第4節所記載:

4施洗約翰來到曠野,宣講悔改的洗禮,使罪得赦.

耶穌呢?按馬可福音第1章第15節所記載:

14……耶穌來到加利利,宣講神的福音,

15說:“日期滿了,神的國近了.你們要悔改,信福音!”

耶穌所傳的道與施洗約翰所傳的道有所異同?

二者所同在於:「你們要悔改.」悔改的行動就是水洗；悔改的後果就是罪得赦免.

二者不同在於:約翰為悔改者施行悔改的洗禮,使罪得赦.只是水洗,但還未到靈洗.

施洗約翰見證,悔改後還要等待聖靈的洗禮,而施行靈洗的就是神子耶穌.

耶穌所傳的道,就是要人信福音.

那麼耶穌所講的福音指甚麼?跟神的國近了有何關係?

耶穌宣講的福音,就是他的身分:「耶穌是神的兒子.」(Jesus is the son of God.)

信福音就是相信耶穌是神的兒子.因為道成肉身,進入世界,天國就近了.

可1:10,太3:16,路3:2,約3:29-34,都記載耶穌領受聖靈的洗禮.

四卷福音書都記載耶穌在洗禮後,天上有聲音說:「這是我的愛子.」

寫路加福音的作者好像感到還不夠,在第3章用了23至38節通過耶穌的家譜,上塑了七十六代,論證耶穌是神的兒子:

23耶穌開始傳道,年紀約有三十歲.依人看來,他是約瑟的兒子,約瑟是希裡的兒子,

……[ 前後句式:「甲是乙的兒子」共講了76次；最後一句……]

38以挪士是塞特的兒子,塞特是亞當的兒子,亞當是神的兒子.

耶穌宣講的福音,就是他的身分:「(我)耶穌是神的兒子.」

因為信他的都成為神的兒女,因為信他的也同時領受靈洗成為天國子民,好讓神的國更靠近.

悔改,靈洗,信福音(可1:8,太3:11,路3:16,約1:19-34參讀)

今天,我們嘗試借助聖經神學(Biblical Theology)的方法,用「四福音平行合參」,解讀今天馬可福音短短七節的經課.

可1:10,太3:16,路3:2,約3:29-34,都記載耶穌領受聖靈的洗禮.

可1:8,太3:11,路3:16,約1:19-34,都強調耶穌施行聖靈的洗禮.

按聖經神學的取向,耶穌受洗這事件,已用上了三一神的語言.首先已得見我們所信仰三位一體神的呈現,其次往後的發展再呈現聖父,聖子與聖靈的互動.

今天的經課讓我們看到,我們的信仰有三個重要元素.

首先,悔改(repent)是最基本的.

悔改,認了罪,罪得赦免.悔改就是向世界回轉朝向神.

然而單單悔改,卻沒有信福音,也沒有靈洗,也只是一種道德行為(moral conduct).

其次,悔改之餘,也信福音(believe in the Gospel).

信福音就是跟隨神道成肉身進入世界.悔改信福音,受洗加入教會,成為天國運動一份子.

可是,悔改加上信福音,卻沒有靈洗,也只算是一種宗教行為(religious behavior).

今天我們留意到信仰有第三個重要元素,就是靈洗(be Spirited).

在悔改與信福音的同時,領受靈洗,就是我們的信仰.你有沒有經歷過靈洗?

教會傳統,由施洗約翰開始施行悔罪的洗禮,但也只是水禮.

教會傳統,水禮這個聖禮(sacrament)也是公開認信,加入教會的一種儀式.這個聖禮是當盡諸般的義的一種.

水洗以外,你如何經歷靈洗(baptized by the Spirit; baptism of the Spirit)?

如果悔改是向世界回轉朝向神,信福音是跟隨神道成肉身進入世界,靈洗就是一個悔改到福音這兩端之間的連續體(a continuum between two poles),時時刻刻活在其中.

但願我們明白我們的信仰,是:「悔改信福音,無時不靈洗.」

但願我們活出我們的信仰:在悔改中確認我們都是聖父神的子女這身分後,借著聖靈的同在,在世不斷學習聆聽神並依靠神的話語,超越試煉與戰勝試探,效法聖子耶穌基督,活出讓神的國靠近人世間的見證.

讓我們一起祈禱,求主啟導.

聖靈啊,

求你洗淨我們的生命,讓我們可以活出我們的信仰.

求你教我們悔改,在悔改中確認我們都是聖父神的子女這身分.

聖靈啊,

求你常充滿我們的生命,教我們在世不斷學習聆聽神並依靠神的話語,超越試煉與戰勝試探.

求你教我們效法聖子耶穌基督,活出讓神的國靠近人世間的見證.

但願榮耀都歸給聖父,聖子,聖靈,永在的三一神.阿們.



\newpage

\section{捨與得}
\label{sec:402}
\textbf{邢福增院長 - 捨與得}
\newline
\newline
\hyperref[sec:401]{< < < PREV SERMON < < <}
~
\hyperref[sec:toc]{[返目錄]}
~
\hyperref[sec:403]{> > > NEXT SERMON > > >}


馬可福音 8:31-38
\newline
\begin{longtable}{cl}
\hline
\hline
章節 & 經文 (和合本修訂版)\\
\hline
8:31 & \begin{tabularx}{0.7\textwidth}{X} 從此，他教導他們說：「人子必須受許多的苦，被長老、祭司長和文士棄絕，並且被殺，三天後復活。」 \end{tabularx} \\ \\ \relax
8:32 & \begin{tabularx}{0.7\textwidth}{X} 耶穌明白地說了這話，彼得就拉著他，責備他。 \end{tabularx} \\ \\ \relax
8:33 & \begin{tabularx}{0.7\textwidth}{X} 耶穌轉過來看著門徒，斥責彼得說：「撒但，退到我後邊去！因為你不體會神的心意，而是體會人的意思。」 \end{tabularx} \\ \\ \relax
8:34 & \begin{tabularx}{0.7\textwidth}{X} 於是他叫眾人和門徒來，對他們說：「若有人要跟從我，就當捨己，背起自己的十字架來跟從我。 \end{tabularx} \\ \\ \relax
8:35 & \begin{tabularx}{0.7\textwidth}{X} 因為凡要救自己生命的，必喪失生命；凡為我和福音喪失生命的，必救自己的生命。 \end{tabularx} \\ \\ \relax
8:36 & \begin{tabularx}{0.7\textwidth}{X} 人就是賺得全世界，賠上自己的生命，有甚麼益處呢？ \end{tabularx} \\ \\ \relax
8:37 & \begin{tabularx}{0.7\textwidth}{X} 人還能拿甚麼換生命呢？ \end{tabularx} \\ \\ \relax
8:38 & \begin{tabularx}{0.7\textwidth}{X} 凡在這淫亂罪惡的世代，把我和我的道當作可恥的，人子在他父的榮耀裡與聖天使一同來臨的時候，也要把那人當作可恥的。」 \end{tabularx} \\ \\
[1ex]
\hline
\hline
\end{longtable}
創世記 17:1-7,15-16
\newline
\begin{longtable}{cl}
\hline
\hline
章節 & 經文 (和合本修訂版)\\
\hline
17:1 & \begin{tabularx}{0.7\textwidth}{X} 亞伯蘭九十九歲時，耶和華向他顯現，對他說：「我是全能的神。你當在我面前行走，作完全的人， \end{tabularx} \\ \\ \relax
17:2 & \begin{tabularx}{0.7\textwidth}{X} 我要與你立約，使你的後裔極其繁多。」 \end{tabularx} \\ \\ \relax
17:3 & \begin{tabularx}{0.7\textwidth}{X} 亞伯蘭臉伏於地；神又對他說： \end{tabularx} \\ \\ \relax
17:4 & \begin{tabularx}{0.7\textwidth}{X} 「看哪，這就是我與你立的約，你要成為多國的父。 \end{tabularx} \\ \\ \relax
17:5 & \begin{tabularx}{0.7\textwidth}{X} 從今以後，你的名字不再叫亞伯蘭，要叫亞伯拉罕，因為我已經立你作多國之父。 \end{tabularx} \\ \\ \relax
17:6 & \begin{tabularx}{0.7\textwidth}{X} 我必使你生養極其繁多；國度要從你而立，君王要從你而出。 \end{tabularx} \\ \\ \relax
17:7 & \begin{tabularx}{0.7\textwidth}{X} 我要與你，以及你世世代代的後裔堅立我的約，成為永遠的約，是要作你和你後裔的神。 \end{tabularx} \\ \\ \relax
17:8 & \begin{tabularx}{0.7\textwidth}{X} 我要把你現在寄居的地，就是迦南全地，賜給你和你的後裔永遠為業；我也必作他們的神。」 \end{tabularx} \\ \\ \relax
17:9 & \begin{tabularx}{0.7\textwidth}{X} 神又對亞伯拉罕說：「你和你的後裔一定要世世代代遵守我的約。 \end{tabularx} \\ \\ \relax
17:10 & \begin{tabularx}{0.7\textwidth}{X} 這就是我與你，以及你的後裔所立的約，是你們所當遵守的，你們所有的男子都要受割禮。 \end{tabularx} \\ \\ \relax
17:11 & \begin{tabularx}{0.7\textwidth}{X} 你們要割去肉體的包皮，這是我與你們立約的記號。 \end{tabularx} \\ \\ \relax
17:12 & \begin{tabularx}{0.7\textwidth}{X} 你們世世代代的男子，無論是在家裡生的，或是用銀子從外人買來而不是你後裔生的，都要在生下來的第八日受割禮。 \end{tabularx} \\ \\ \relax
17:13 & \begin{tabularx}{0.7\textwidth}{X} 你家裡生的和你用銀子買的，都必須受割禮。這樣，我的約就在你們肉體上成為永遠的約。 \end{tabularx} \\ \\ \relax
17:14 & \begin{tabularx}{0.7\textwidth}{X} 不受割禮的男子都必從民中剪除，因他違背了我的約。」 \end{tabularx} \\ \\ \relax
17:15 & \begin{tabularx}{0.7\textwidth}{X} 神又對亞伯拉罕說：「至於你的妻子撒萊，不可再叫她撒萊，她的名要叫撒拉。 \end{tabularx} \\ \\ \relax
17:16 & \begin{tabularx}{0.7\textwidth}{X} 我必賜福給她，也要從她賜一個兒子給你。我必賜福給撒拉，她要興起多國；必有百姓的君王從她而出。」 \end{tabularx} \\ \\
[1ex]
\hline
\hline
\end{longtable}
羅馬書 4:13-25
\newline
\begin{longtable}{cl}
\hline
\hline
章節 & 經文 (和合本修訂版)\\
\hline
4:13 & \begin{tabularx}{0.7\textwidth}{X} 因為神給亞伯拉罕和他後裔承受世界的應許不是藉著律法，而是藉著信而得的義。 \end{tabularx} \\ \\ \relax
4:14 & \begin{tabularx}{0.7\textwidth}{X} 若是屬於律法的人才是後嗣，信就落空了，應許也就失效了。 \end{tabularx} \\ \\ \relax
4:15 & \begin{tabularx}{0.7\textwidth}{X} 因為律法是惹動憤怒的，哪裡沒有律法，哪裡就沒有過犯。 \end{tabularx} \\ \\ \relax
4:16 & \begin{tabularx}{0.7\textwidth}{X} 所以，人作後嗣是出於信，因此就屬乎恩，以致應許保證歸給所有的後裔，不但歸給那屬於律法的，也歸給那效法亞伯拉罕之信的人。 \end{tabularx} \\ \\ \relax
4:17 & \begin{tabularx}{0.7\textwidth}{X} 亞伯拉罕所信的是那叫死人復活、使無變為有的神，在這位神面前亞伯拉罕成為我們眾人的父，如經上所記：「我已經立你作多國之父。」 \end{tabularx} \\ \\ \relax
4:18 & \begin{tabularx}{0.7\textwidth}{X} 他在沒有盼望的時候，仍存著盼望來相信，就得以作多國之父，正如先前所說：「你的後裔將要如此。」 \end{tabularx} \\ \\ \relax
4:19 & \begin{tabularx}{0.7\textwidth}{X} 他將近百歲的時候，雖然想到自己的身體如同已死，撒拉也不可能生育，他的信心還是不軟弱， \end{tabularx} \\ \\ \relax
4:20 & \begin{tabularx}{0.7\textwidth}{X} 仍仰望神的應許，總沒有因不信而起疑惑，反倒因信而剛強，將榮耀歸給神， \end{tabularx} \\ \\ \relax
4:21 & \begin{tabularx}{0.7\textwidth}{X} 且滿心相信神所應許的必能成就。 \end{tabularx} \\ \\ \relax
4:22 & \begin{tabularx}{0.7\textwidth}{X} 所以這也就算他為義。 \end{tabularx} \\ \\ \relax
4:23 & \begin{tabularx}{0.7\textwidth}{X} 「算他為義」這句話不是單為他寫的， \end{tabularx} \\ \\ \relax
4:24 & \begin{tabularx}{0.7\textwidth}{X} 也是為我們將來得算為義的人寫的，就是為我們這些信神使我們的主耶穌從死人中復活的人寫的。 \end{tabularx} \\ \\ \relax
4:25 & \begin{tabularx}{0.7\textwidth}{X} 耶穌被出賣，是為我們的過犯；他復活，是為使我們稱義。 \end{tabularx} \\ \\
[1ex]
\hline
\hline
\end{longtable}
講題:捨與得　Losing to Gain

經課:馬可福音8:31-38;創世記17:1-7,15-16;羅馬書4:13-25

講員:邢福增院長

場合:香港中文大學崇基學院禮拜堂

日期:2018年2月25日

得了甚麼?失了甚麼?

1992年諾貝爾經濟學獎得主貝克(Gary Becker),擅長運用經濟學理論來解釋人的行為:人生很多抉擇,事實上都是基於「得」,「失」的經濟性考慮,例如為甚麼選擇這份工作,為甚麼結婚,離婚,諸如此類……一切的抉擇都是出於在權衡「得」,「失」之後,深信是對自己有利才作出的.有一次當他發表演說後,有人問他:為甚麼一些人(例如傳教士),會作出那麼大的自我犧牲?顯然,這些人的選擇不僅僅是為了一己的「得」,「失」了.貝格在思想一會後,回答說:其實,這種人仍是在「得」與「失」之間權衡,問題是他們所期望得到的,與其他人不一樣而已.他們得到了別人的讚揚,尊重……成為一個公認為偉大的人物(如德蘭修女),這些都不是金錢可以換取的.歸根究柢,基督徒也是受「得」「失」的標準定律影響.

貝格的說法令我想起徐小鳳一首多年前的名曲,是1978年電視長劇「大亨」的主題曲.這是一齣四十年前極受歡迎的電視劇,那時我在讀初中,一直在追看.顧名思義,「大亨」就是描寫男主角鄭少秋為了爭取成功而不擇手段,最終成為社會的名流大亨的故事.主題曲的歌詞有有意思,說:

他也在找,我也在找,找到名利幾多；他拚命追,我拚命追,追到又如何?

追到甚麼?找到甚麼?收到又幾多?得了甚麼?失了甚麼?可有認真算過?

何必呢?何必呢?可知一切他朝都會身外過.得的多?還失的多?升得高的終於都會低墮.何必呢?何必呢?拋開一切束縛身心韁鎖.且向心內,仔細追尋,找那安然ge 我.

人生在世,到底在尋找甚麼?為了追求名與利,又要付上甚麼代價?就算真的得到名與利,失去的又是甚麼?這又是否值得呢?

基督是誰?

今主日是大齋節第二主日,福音書的經課是馬可八章,也是一個與得失有關的教導.當時耶穌剛醫治了伯賽大的盲人,在往該撤利亞．腓立比的路上,問門徒:「人們說我是誰?」門徒將聽到的不同說法說出來(施洗約翰,以利亞,先知的一位……),接著耶穌再追問:「你們說我是誰?」當時,彼得說:「你是基督」(八27-30).

彼得似乎將標準答案講了出來,但其實他並不知道「基督」的真正意思.馬可經課一開始說:「從此,他教導他們……」(31a),這裡,「從此」的原文,是「開始」,即耶穌正式跟門徒講述關於「基督」的身分:「人子必須受許多的苦,被長老,祭司長和文士棄絕,並且被殺,三天後復活.」(31b)這是一位「受苦」的彌賽亞,「必須」是指人子為實現道成肉身,降世為人的目的,必須經歷被棄絕與被殺.雖然耶穌提及「三天後復活」,但門徒根本完全聽不明白.經文描述「耶穌明白地說了這話」,「明白」有譯本譯作「講得清清楚楚」,就是毫不隱晦地向門徒表白了「基督」的身分與使命.讀到這裡時,我想起校牧室署理校牧Walter林豪恩,當他被崇基學院方永平院長召見時,也是「明白地說」「講得清清楚楚」:現在你「必須」署任校牧,對你而言,這會「受許多的苦」,在我們找到新校牧後,你便可以「復活」.當時,我相信Walter對「找到校牧後復活」沒有多大的感覺,而只是聽到「必須」「受苦」……

我不知道Walter的太太知道他答應方院長後,會否像彼得對耶穌一樣:「就拉著他,責備他」(據馬太福音十六章22節,彼得說 :「主啊,千萬不可如此!這事絕不可臨到你身上.」).彼得的反應很直接,「責備」反映是很肉緊的流露,類似用粵語很強語氣地說:「你咁樣點得架」!彼得是愛護自己的老師,不想耶穌受苦嗎?也許是,但相信他更大的考量是自己的前途:我們捨下一切來跟從你,從沒有想過你會被棄絕及被殺的.那麼,我們「點算好」?我們是否跟錯大佬?「咁點得架」!

作門徒真好?

到馬可第九章,彼得在山上見到改變形象的耶穌與摩西及以利亞在一起時,是何等的興奮(這是兩周前的經課呢).他說:「拉比,我們在這裡真好!我們來搭三座棚,一座為你,一座為摩西,一座為以利亞.」(九5)「我們在這裡真好」跟責備耶穌「千萬不可如此」(太十六22)正是出於同樣的心態.彼得及門徒跟從耶穌,也許確被耶穌的教導吸引,但他們也在計算得失,我捨下所有來跟從你,到底又會得到甚麼?跟從耶穌會得到甚麼?如今,耶穌的決定令他們前功盡廢,換了是你,不反對才怪呢.

弟兄姊妹,跟從耶穌會「得」到甚麼呢?我們會否被創世記十七章的經課吸引著:耶和華對亞伯蘭說:「你要成為多國的父」(4)「我必使你生養極其繁多；國度要從你而立,君王要從你而出」(6)「要賜一個兒子給你」(15)跟從耶穌,作主門徒,就是因為這是能夠給我許多好處的:土地代表財富,多國的父代表地位,在年老時仍能生兒子代表神蹟……這些,如果信主後都會得到,多好?

這正是耶穌當日要校正門徒的心態,耶穌備責彼得為「撒旦」,並不是說彼得就是撒旦,而是說當彼得「不體會上帝的心意」,只「體會人的心意」時,這種想法無疑是撒旦的意念.耶穌甘於忍受人世間的苦難,正是體會上帝心意的表現.因為,他清楚知道自己的目的與使命,正如保羅在羅馬書經課中說:「耶穌被出賣,是為我們的過犯,他復活,是為使我們稱義.」(羅四25)即使受苦,但人子為了救贖世人的罪,使罪人在上帝面前得稱為義,是必須要受苦的.受苦是「捨」己,卻使罪人「得」著新的生命.福音正是充滿著這種吊詭的張力:為了別人(他者)而捨棄自己,正是得著生命.

捨己,背十架,跟從主

接著,耶穌進一步教導門徒關於「捨棄」與「得著」的吊詭性.他指出,作門徒就是「跟從」基督,必須要先「捨己」,再「背起自己的十字架」(34),「跟從」,「捨己」與「背起自己的十字架」是連貫在一起的行動.甚麼是「捨己」?是佛家所說的「破我執」,免去人世間所有的「執著」嗎?捨己當然要放下自己某些執著,但耶穌所言,應有比「破我執」更大的意思.因為,除了「捨己」,還要「背起自己的十字架」來「跟從」耶穌.背十字架顯然就是預表了耶穌要走的路,因此,跟從耶穌的門徒,也要像基督一樣,背起自己的十字架.那麼,耶穌是指每一個跟從他的人都要走上殉道的路嗎?注意,耶穌說的是「自己的十字架」,每人都有自己的十字架,這是獨特而不一樣的十字架,卻是為了「體會上帝的意思」而背上的.接下來,在35至38節,耶穌進一步作說明這種生命的「吊詭性」:

第一,「因為凡要救自己生命的,必喪失生命；凡為我和福音喪失生命的,必救自己的生命.」(35)從體會人的意思來「救」自己,其實是「喪失」自己.相反,體會上帝的意思(「為我和福音」)好像換來「喪失」自己,卻吊詭地是「救」了自己.

第二,「人就是賺得全世界,賠上自己的生命,有甚麼益處呢?」(36)體會人的意思,無疑可以「賺得」全世界,但代價卻是「賠上」自己.

第三,「人還能拿甚麼換生命呢?」(37)耶穌所指的「生命」,不是「性命」,而是「自我」.體會人的意思的「自我」,其實是一個扭曲了的「自我」.一個真正完整的「自我」,必須是按著上帝的心意而活的生命.

第四,「凡在這淫亂罪惡的世代,把我和我的道當作可恥的,人子在他父的榮耀裡與聖天使一同來臨的時候,也要把那人當作可恥的.」(38)耶穌對門徒說的再清楚不過,跟從耶穌,作主門徒,是會被這個「淫亂罪惡的世代」視作「可恥」的.因為,這個世代希望每一個人都只為「自己」而活,為了自己,可以「雙重標準」,可以「特事特辦」,可以「以權謀私」,可以「不擇手段」,可以「推諉他人」,可以「卑躬屈膝」,可以「直鹿為馬」,可以「濫用公權」,可以「識時務」「看風駛」……總之,「人不為己,天誅地滅」.作主門徒的意思,卻是要「捨棄」這種自我中心的標準與價值,改為按照上帝國的價值與方向行事為人.背起十字架既是各人付上不同的代價,但同時也是代表門徒群體一起彰顯出的上帝國度.

我是誰?

原來,耶穌藉著問門徒「我是誰?」,將「耶穌是誰?」的問題,轉化成為「門徒是誰?」跟從耶穌走的道路,是拒絕認同這個與「世代」的價值與標準,因為有一個更根本的價值──「為了我和我的道」──就是上帝的國度.耶穌清楚說明,背十字架一定與自己的意思相違,但吊詭地,這個「捨棄」卻是真正的「得著」自己.

我們讀創世記十七章的經課,從人的意思,只讀到財富,地位與神蹟,卻忽略了耶和華向亞伯蘭顯現時說:「我是全能的上帝.你當在我面前行走,作完全的人,我要與你立約,使你的後裔極其繁多.」(1-2)耶和華要與亞伯蘭建立「約」的關係,與上主立約,是要在祂面前行走,「作完全的人」.「世世代代的後裔要堅立我的約,成為永遠的約,是要作你和你後裔的上帝」(7)「一定要世世代代遵守我的約」(9).跟從基督,作主門徒,就是說我們的生命與自我,需要放置在與主的關係之中.我們的生命,已經被基督贖回.正如羅馬書的經課說:「耶穌被出賣,是為我們的過犯,他復活,是為使我們稱義.」(羅四25)

是的,要「作完全的人」是不容易的,完全不是自以為義,而是代表因著上主的道,使人格變得正直與完整(integrity),這個「世代」一定會把「上主的道」及其實行者「當作可恥的」.對此,耶穌卻從終末的角度給予門徒盼望與應許:有一日,作主門徒,會見到「人子在他父的榮耀裡與聖天使一同來臨的時候」,那時,被顛倒的是非會得到公正的審判,並且,人子會將那些把「上主的道」當作可恥的人「當作可恥的」.那麼,這一天,將要等多久呢?

想像一下,當亞伯拉罕的後代,在亡國被擄的日子讀創世記十七章時的光景:這時,以色列國已經國破家亡,君王與百姓被擄,沒有土地,沒有地位,也沒有神蹟.但創世記的編修者卻在提醒被擄群體,這位全能的上帝與他們之間,仍有永遠的約.這段經文,正是給予被擄群體面對黑暗的盼望與應許.沒有人知道黑暗的日子有多長,上主的子民與百姓還要被視作「可恥」有多久,唯一可以持守的,只有仰望上主的盼望.

在沒有盼望的時候,仍存著盼望來相信

最後,跟大家分享一位103歲的中國老牧師的生命.他剛於2月17日凌晨安息主懷.他的名字是范愛侍,是1957年被打為中央級別的五位基督教界右派之一.他在回憶錄中,記下這段一生難忘的痛苦經歷:當他知道自己被同道攻擊,劃為右派後.在1957年11月6日晚上,上到會議場地的天台,「跳下去吧!了此無告的人生……!」就在此時,他想起自己的蒙召與事奉,「今壯志未酬,捲入了政治運動,背上了滔天罪名,有口難分.功尚未成,我即去乎?悲痛,迷惘,無助,無告……之情,直無法形容.我猶如一葉扁舟,被狂風追逐於汪洋大海之上,迷失了方向,不知所歸.又想到愛妻和年齡1至13歲的5個孩子們.我怎忍心棄他們而去呢?」「『你不可跳!』上來了另一個意思.此時,人生天平的指針:在生與死之間劇烈地擺動.任何微小的砝碼加在任何一邊,都將導致決定性的後果.」「我忽然想到神子耶穌受撒旦試探的一幕:撒旦曾誘使他『從殿頂上跳下去』時,耶穌洞悉其詭計:斥其『退去』!我想,我該效法救主的腳步,懸崖勒馬.從而避免了悲劇的發生.但為此,我付上了沉重的代價:無情的挨批,痛苦的改造,屈辱的人生,艱難的歲月……」

文化大革命結束後,范牧師在1978年5月獲平反,換言之,他戴了右派的帽合共22年(即8030天)……他後來說:「如果1957年11月6日晚,北京新僑飯店平頂上,我預知未來改造之路是如此艱巨而漫長,說不定我人生天平的指針,已向另一邊傾斜!」真的,如果范愛侍牧師預知他需要忍受22年,他大概不能熬過這8030天.但他卻是在靠著上主,一天又一天地熬過.他真的喪失了許多,賠上了許多,捨棄了許多,但吊詭地,在上主的國度裡,他卻救了自己,賺了,及得著更多.

大齋節期讓我們再思生命.在上帝國度裡,我們要學習捨棄的是甚麼?為了他者,甚麼是我們的十字架?我們會面對各式各樣的試探,我們的回應與選擇,也是在主面前回答:「我是誰?」這問題.在某些地方,仍有人要為主殉道.但是,這也可以是為了家庭兒女而奉獻自己生命的母親,為照顧摯親而擺上自己的人,與邊緣弱勢者的同行,因為改變社會而選擇公民抗命的被囚青年人或流亡者,甚至是為了改善權益而被甘冒被資方懲處的勞方……看似「捨棄」卻是為了「他者」而擺上.今天,我們所持守的「主和主的道」又是甚麼?為此,捨了甚麼?得了甚麼?

聽說Walter在辦公室門上貼了一張「校牧快來」的揮春.願使徒保羅的提醒也臨在我們:「在沒有盼望的時候,仍存著盼望來相信」(羅四18 ).我們真的不知道終末的榮耀何時實現,因此,更要仰賴上主的恩典來渡過每一天.



\newpage

\section{耶穌都有火}
\label{sec:403}
\textbf{許開明牧師 - 耶穌都有「火」!}
\newline
\newline
\hyperref[sec:402]{< < < PREV SERMON < < <}
~
\hyperref[sec:toc]{[返目錄]}
~
\hyperref[sec:404]{> > > NEXT SERMON > > >}


約翰福音 2:13-22
\newline
\begin{longtable}{cl}
\hline
\hline
章節 & 經文 (和合本修訂版)\\
\hline
2:13 & \begin{tabularx}{0.7\textwidth}{X} 猶太人的逾越節近了，耶穌上耶路撒冷去。 \end{tabularx} \\ \\ \relax
2:14 & \begin{tabularx}{0.7\textwidth}{X} 他看見聖殿裡有賣牛羊和鴿子的，還有兌換銀錢的人坐著， \end{tabularx} \\ \\ \relax
2:15 & \begin{tabularx}{0.7\textwidth}{X} 耶穌就拿繩子做成鞭子，把所有的，包括牛羊都趕出聖殿，倒出兌換銀錢之人的銀錢，推翻他們的桌子， \end{tabularx} \\ \\ \relax
2:16 & \begin{tabularx}{0.7\textwidth}{X} 又對賣鴿子的說：「把這些東西拿走！不要把我父的殿當作買賣的地方。」 \end{tabularx} \\ \\ \relax
2:17 & \begin{tabularx}{0.7\textwidth}{X} 他的門徒就想起經上記著：「我為你的殿心裡焦急，如同火燒。」 \end{tabularx} \\ \\ \relax
2:18 & \begin{tabularx}{0.7\textwidth}{X} 因此猶太人問他：「你能顯甚麼神蹟給我們看，表明你可以做這些事呢？」 \end{tabularx} \\ \\ \relax
2:19 & \begin{tabularx}{0.7\textwidth}{X} 耶穌回答他們說：「你們拆毀這殿，我三日內要把它重建。」 \end{tabularx} \\ \\ \relax
2:20 & \begin{tabularx}{0.7\textwidth}{X} 猶太人說：「這殿造了四十六年，你三日內就能重建嗎？」 \end{tabularx} \\ \\ \relax
2:21 & \begin{tabularx}{0.7\textwidth}{X} 但耶穌所說的殿是指他的身體。 \end{tabularx} \\ \\ \relax
2:22 & \begin{tabularx}{0.7\textwidth}{X} 所以他從死人中復活以後，門徒想起他曾說過這事，就信了聖經和耶穌所說的話。 \end{tabularx} \\ \\
[1ex]
\hline
\hline
\end{longtable}
講題:耶穌都有「火」!　Jesus Had a “Temper”!

經課:約翰福音2章13-22節

講員:許開明牧師

場合:香港中文大學崇基學院禮拜堂

日期:2018年3月4日

引言:

今早很高興有機會返來崇基學院禮拜堂,時間過得很快,離開大家已有七個月了!

感謝上主,我於上年七月三十一日榮休後,八月一日旋即應邀成為中華基督教會合一堂主任牧師,讓我於退休之年,還是強壯,仍可以繼續事奉!合一堂主任牧師跟一般堂會牧師不同,這職要督導幾間座堂,還有社會服務事工,差傳事工,教育事工,教會發展事工等.實在給我有許多發揮的機會.

今天的經課,約2:13一26

【約二13】「猶太人的逾越節近了,耶穌就上耶路撒冷去.」

【約二14】「看見殿裡有賣牛羊鴿子的,並有兌換銀錢的人,

坐在那裡.」

【約二15】「耶穌就拿繩子作成鞭子,把牛羊都趕出殿去；

倒出兌換銀錢之人的銀錢,推翻他們的桌子.」

【約二16】「又對賣鴿子的說:『把這些東西拿去；

不要將我父的殿,當作買賣的地方.』」

本文提到在耶路撒冷的聖殿,耶穌在聖殿發火,這火看來非常猛烈,使人震驚!我們不期然會問:點解耶穌發咁大火!耶穌平日給人的形像不是謙卑柔和,頗有教養的謙謙君子,也是受人尊重的夫子(老師),何以今次舉止如此粗魯?

中國人所謂:「佛都有火」,這四字讓我想到「耶穌都有火! 」

“佛都有火”是指一個好脾氣的人,都忍不住氣,要發火,表示有人的所作所為,令到好脾氣人也難以忍受!正所謂「佛都有火」,「耶穌都有火」 , 那麼我們都會有火, 人人都有火!事實上, 發火是我們常見的情緒!

發火好嗎?

我們可能會遇見兩種極端的人:

一種人,無乜火氣；另一種人,太有火氣,好易火爆.無乜火氣的人,會被人話沒有火氣,過份和善!這種人易被欺負!太有火氣,常常火爆,易得罪人,更容易在失控下,幹出大錯特錯的事情.無論如何,兩種極端,都各有弊端!

【憤怒乘客,憤怒車長,引致嚴重車禍】

2018年2月10日在大埔公路上,發生嚴重車禍,一架872線巴士在高速下翻側, 釀成19人死亡,65人傷, 實在慘不忍睹.目擊者言,車長因不忿乘客斥責而刻意高速駕駛, 這宗完全是人為釀成的車禍.由乘客的火,觸動車長的火,車長怒火中燒, 高速轉彎, 導致巴士失控,釀成嚴重的傷亡!怒氣, 發火, 自古到今, 不知造成多少意外,多宗傷亡, 每個人都要小心!

憤怒是不是罪?

憤怒是否不合宜的?是否是罪?

憤怒的本身是罪,或者不當表達憤怒才是罪?

這些觀念均會影響我們處理憤怒情緒.

從基督徒的觀點來看,雷． 布威克 (    Ray Burwick   ) 認為自私的怒氣是罪,不自私[義怒]的怒氣則可以是正當的.學者 黎曦庭(Tim LaHaye),腓立普斯 (Bob Phillips)也認同這個看法.柯蓋瑞(Gary R. Collins)認為怒氣是正常的,不一定有罪,若然處理不當,卻往往導致犯罪.Augsburger 則區分憤怒情緒本身以及如何處理憤怒情緒,認為前者並無對錯,而後者才有道德上的對錯.

怒氣的種類:

新舊聖經「怒氣」的涵義

舊約中「怒氣」, 在希伯來文是“鼻子或鼻孔” (怒氣是噴著怒火的鼻子).所以,希伯來文的諺語,長鼻子表示「不輕易發怒」;「短鼻子」, 則表示「容易發怒」.

新約希臘文的「怒氣」, 分開兩類: 一是「怒氣」,乃為「自然的情緒」;二是「忿怒」,乃為「強烈的,持久的情緒」.

聖經不同類型的怒氣:

1義怒:是由於看見別人行為不公正,不公義,引起憤怒,並有份強烈的矯正的願望.這是好的怒氣!值得推崇的義怒!

2憤怒或怒氣:發怒的情緒,常伴有報復的欲望.憤怒的人有氣憤的感覺,希望宣洩出來.這是會帶來破壞,使人受傷!倘不能抒洩和受控時,怒氣更會逐漸升級!

3暴怒或忿怒:強烈的怒氣,往往會使人失去理智,不能自控,造成對人莫大的傷害或毀壞.不加以制止時,情況會非常嚴重!

自私的生氣:能粉碎溝通,離間感情,並毀掉無數人的健康和幸福.可悲的是,人們總是為生氣找藉口而不對它負責.在不同程度上,大家都在與怒氣爭戰.感謝上主,聖經包含了以屬靈方式應對生氣和克服怒氣的原則.

不自私的生氣:是義怒:這種生氣是聖經允許的,常被稱為「公義的憤慨」. 上帝也是發怒的神(詩篇7:11,馬可福音3:5).新約中用了兩個發怒希臘詞,上帝的發怒含有「激情和活力」的；人的發怒:含有「不安,熾熱」的.根據聖經,上帝賜給人怒氣這個能量, 推動我們解決問題.請注意「義憤」,不是為了自己,而是為保護他人或維護正確的原則.

然而,當生氣有自私的動機(雅各書1:20),人便想報復!當這怒氣不消時(以弗所書4:26-27),生氣就變成了罪.由生氣產生的能量攻擊的不是面對的問題,而是眼前的人,這就導致衝突和傷亡.

當我們掌握了一些聖經上對憤怒的陳述後,現在回到今天的經課:

耶穌為何在聖殿發怒?!

筆者發現有三大原因:

1聖殿成為不潔,不公平買賣的地方

當中有疑似「特許經營店」,藉此發聖殿財.馬太福音21章12至17節,記載了在聖殿中兌換銀錢和販賣鴿子的事情；參考《約翰福音》2章13至16節所描述的,連牛和羊等牲畜都在聖殿中販賣.其實,兌換銀錢是不公道,不公平和不公義的交易.當時朝聖的人,因旅途之遙遠,他們無法帶同牲畜一同上路前往聖殿,為著獻祭,他們即時在聖殿購買.如此,便衍生「發聖殿財」的機會.

以今天的語言重新描述這件事,就是有一小撮人,不知怎樣可得到特許經營權,在聖殿中販賣用作獻祭的牲畜.即使前往敬拜的人,自行攜帶牲畜也好,在聖殿中工作的祭司是有權決定那自備的牲畜可否作為牲祭之用.他們會對自備的牲畜左挑右惕,指出所帶來的牲畜有瑕疵、不合乎標準、不合資格獻上；以致朝拜者必須向那些「特許經營店」購買,收費會較高昂.

至於兌換銀錢,可能「匯率」甚高,「食水」甚深.對於到聖殿朝拜的人來說,這些「特許經營店」或「特許找換店」是對欺壓朝拜者.但我希望大家留意,這些欺壓和欺詐的行為,是在聖殿任職的神職人員,進行「教商」勾結而生.

當中可惡之處,是教會人士與商人勾結活動,以「聖殿敬拜」之名進行.因此,聖殿本稱為萬國禱告的殿,卻被這些神職人員作為不公義貿易和謀利之場所!無怪耶穌對他們說:『經上記著說:「我的殿必稱為禱告的殿.你們倒使他成為賊窩了!」』」(太 21:12-13) 從這方面瞭解當時的境況,更可理解為何主耶穌毫不留情地稱聖殿做「賊窩」；全因為有人以「聖殿敬拜」的名,卻用不公義的手法榨取敬拜者的金錢、甚至是他們的敬拜情感.

2 聖殿成為不寧靜,變成不能敬拜的場所

耶穌看見(約 2:14): (1)“看見殿裡有賣牛,羊,鴿子的,並有兌換銀錢的人,坐在那裡”(約 2:14).賣牛,羊,鴿子,兌換銀錢,本來都是為了方便外地來的,或遠方來的敬拜者.解經家發現,過去商人們的攤位是設在汲淪溪對岸的,現在,那些賣牛,羊,鴿子的,並有兌換銀錢的人,把買賣搬到殿裡來了.這裡“殿”,主要是指聖殿的外院,也稱為外邦人的院,是外邦人唯一能進入禱告的地方,現在猶太人把它變為吵雜的市場.

3 還有更嚴重的事是猶太人侵犯外邦人可以敬拜,禱告神地方

這表示猶太人根本不想外邦人敬拜上主,相反卻要在神聖的地方搬賣,奪取他們金錢.這樣的吵雜,和污穢, 嚴重影響到所有來聖殿敬拜神的人,影響到所有人在敬拜時對神的敬虔,以至破壞了整個敬拜的氣氛,破壞了人的心藉著敬拜可以歸向神,也即失去了整個敬拜的目的和意義.因此耶穌「又對賣鴿子的說:『把這些東西拿去；不要將我父的殿,當作買賣的地方.』」【約二16】

耶穌「憤怒」的特質

舊約

聖經中也說到了聖潔的憤怒.這是一種宗教上的義憤(Righteous indignation)是指對罪惡無法容忍,聖潔心靈表現出的極度不滿.

上帝的憤怒就包含了這些因素,上帝期望人應該成為良善的,然而他卻犯罪,因此惹動上帝怒氣.

「因 這 地 的 人 離 棄 了 耶 和 華 他 們 列 祖 的 神 , 領 他 們 出 埃 及 地 的 時 候 與 他 們 所 立 的 約 ,去 事 奉 敬 拜 素 不 認 識 的 別 神 , 是 耶 和 華 所 未 曾 給 他 們 安 排 的 .」同樣是基於這種因素,當摩西看到金牛犢和以色列人拜偶像時,他的憤怒在西奈山上爆發,導致兩塊法版在地上摔得粉碎.信徒需要理解合宜的和非合宜的兩種憤怒,並要確定自己的發怒是符合當時情況,就像上帝發怒時一樣.

新約

在新約中,提及耶穌的義怒:耶穌的發怒代表上帝的怒氣,還有群眾對聖殿當權者的憤怒.耶穌有發怒,是沒有罪的(來7:26-28),因此合宜地表達怒氣並不是犯罪.人發怒的本質不是犯罪,但怒氣處置不當卻常帶來犯罪.然而,怒氣無論是自私地為己而發或公義地為人而發,最重要是學習如何適當處置怒氣.人們合宜地動怒有時候也稱之為「上帝的怒氣」(the wrath of God).

耶穌的「憤怒」,是社會性的,是整體社會都感到「憤怒」的感覺.當我們的尊嚴、安全、或社會所持重要的價值被冒犯時,以致我們有「大眾憤怒」的感覺.更可悲的是,若果有朝一日,大家對於尊嚴被冒犯而不再感到憤怒的話,可能我們已經被馴化了.又或者,我們成為旁觀者,而感到無能為力,這個時候便更危險了.在主耶穌潔淨聖殿的事件中,我們學習因為面對不公義而產生的憤怒,可以作為一股正能量.我們面對不公義,不會變得麻木.我們可以學習主耶穌,將憤怒轉化為行動.

耶穌的憤怒來自其火熱的心

耶穌的門徒就想起[經上]記著說: 「我為你的殿心裡焦急,如同火燒」(約 2:17). 當時的門徒,可能會非常害怕和吃驚,因為他們是第一次看見,耶穌會發如此大的火, 而且是公開挑戰猶太人的當權者.經文這裡的「想起」,原意:記得,表示這給他們留下了 極深刻的印象,他們將耶穌的發火和舊約聯繫起來, 「因我為你的殿心裡焦急,如同火燒」 (詩 69:9 上).這裡「如同火燒」,原意:火熱激昂; 「焦急」原意:吞吃; 意思是「為神的家那種火熱激昂,都將要把我給吞吃了」.

耶穌憤怒的結果

1 引起注意

「耶穌進了　神的殿,趕出殿裡一切作買賣的人,推倒兌換銀錢之人的桌子和賣鴿子之人的凳子.對他們說:『經上記著說:「我的殿必稱為禱告的殿.你們倒使他成為賊窩了!」』」(太 21:12-13)

「耶穌憤怒」,引起當權者的注意!這憤怒轉化出來的行動.主耶穌就是因為面對父的聖殿被拖累成為賊窩,這些人對神不敬而使主耶穌怒氣爆發.

這段經文的記載,給我學習的一個重點,是將憤怒化為可以抗議的、表達不滿的、表達給對方知道不可再過份侵害社會等能力,對準那些行不公義的人發出抗議；即是要將憤怒轉化為有針對性的行動.

這種「社會性的憤怒」的出現,是提醒我們社會上有不公平、損害人的尊嚴等等事情正在發生.若果在憤怒產生的同時,只有無奈感,就置身於「高牆與雞蛋」的處境中；高牆,代表當權者,雞蛋,指一般平民,無權無勢的人!情況更壞的是,這種「社會性的憤怒」長期被抑壓下去,沒有有效的途徑給我們表達的話,始早社會被一種巨大的無力感所癱瘓.

「高牆與雞蛋」的力量懸殊、表達的空間愈來愈小、要針對的對象可能是高不可攀的權貴和制度,成效可謂是一個謎.回頭望主耶穌在聖殿所做的,給了我們很好的提醒.這一次的驅趕可能是唯一的一次,因為主耶穌當時入聖殿後,便入耶路撒冷城受死.

那怕只是一次,主耶穌潔淨聖殿的行動成功了.祂這次行動是凸顯了悠來已久的事情荒謬之處,就是在上帝的殿中,在神的居所裡,有人行如此不公義的買賣、並經年累月的發生而無人反抗.就算主耶穌未能除滅這種不公義的事情,也很有力地表達面對不公義而感到之憤怒.耶穌的憤怒和行動提醒我們,要將憤怒轉化為一種力量去揭露不公義的事情.

2.引起關注

猶太人,當權者的關注:因此猶太人問祂說:“你既作這些事,還顯什麼神跡給我們看呢?”(約2:18).主耶穌當眾向猶太人的當權者發出挑戰,那些當權者不敢責備,和制止祂,因為他們 心裡面是理虧的.但是,他們不服氣,他們要挑戰耶穌的權柄,他們要求主耶穌行一個 神跡來證明祂潔淨聖殿的權柄.

義怒卻要於適當的時候才發出來!

面對不公不義的社會,我們的義怒,有時要小心的,謹慎的表達；甚至被迫敢怒而不敢言.當你知敢言可能傷及家人和朋友.懷著義怒,在適當的時候,適當的場景,適當的對象下,耶穌才發火!

有時,太過急,不適當的時候發火,反被不仁不義的人,把你消滅,不但燒不到對方,反而使你引火自焚!或被惡人抓到機會把你消滅!

如此觀點,耶穌不是也中招嗎?祂對猶太人來說!此舉是自投羅網,甚至引火自焚,祂得罪了撒都該人,大祭司! 不是走進死路嗎! 耶穌走上耶路撒冷就要準備早已準備成為被殺的羔羊!但耶穌被殺,死了,祂第三天可以復活.衪是神兒子,事實上衪是神,你卻不是!耶穌能從死裡復活,反敗為勝!你也能嗎!

3.引起反思

耶穌回答說:“你們拆毀這殿,我三日內要再建立起來”(約 2:19). 耶穌用“你們拆毀”,“我再建立”來對比.「你們拆毀這殿」的意思:你們污穢了聖殿,破壞了敬拜,你們這是在拆毀這殿;「我三日內,必要再建立起來」的意思: 被拆毀的殿,被破壞的神人關係,只有我來才能被重建.注意這裡的“殿”字,原意: 指聖殿裡的“聖所”,呂振中譯成“殿堂”, 可以理解成是最神聖的神人相遇的地方.

4.引出真知

猶太人便說:“這殿是四十六年才造成的,你三日內就再建立起來嗎?”(約 2: 20).

1猶太人完全不明白

「這殿是四十六年才造成的」,可見猶太人完全不明白耶穌的意思,他們思想的殿,還是指外面的建築物.當時的聖殿,是主前二十年左右,大希律王開始建造的, 至當日已經造了四十六年了,耶穌三日內就可以再建起來? 他們當然不會相信.猶太人要 看的神跡就是肉眼所能見的神跡,他們關心的就是神跡給他們帶來的好處,他們不懂得神跡只是一個“記號”,神跡是要引導人去認識行神跡的神,是要引導人的心去歸 向神.

2耶穌講出神蹟:三日建殿的神蹟

「但耶穌這話是以祂的身體為殿」(約 2:21).

這裡的“殿”字,指聖所.「耶穌這話是以祂的身體為殿」,意思是,建築物的聖殿,將以耶穌基督自己來取代.這也是在預言祂的死而復活.猶太人要看神跡,即“記號”;耶穌的意思,“記 號”最根本的指向,就是要他們去仰望十字架;主耶穌告訴我們,祂的義怒不僅僅是為了潔淨,更重要的是為了重建.祂來,就是要除去舊的,建立新的,這才是祂死而 復活真正的意義.

3主耶穌的身體才是殿!這實在是個極重要的啟示和教導

這“殿”字,1從外院, 2到神的家,3到聖所,4最後到主耶穌的身體;什麼是主耶穌的身體? 那麼, 既然耶穌的身體才是殿,才是“聖所”,表示對我們新約信徒來講,只有在耶穌基督裡,我 們才能和神有真正的相交.

耶穌成功地處理憤怒

1)耶穌的憤怒有適當的理由.也就是說,祂有正當的理由.耶穌的憤怒不是出於自私的緣故,並未有引起爭論 . 在義怒中,耶穌推倒買賣的?已指出憤怒原因是:「把父的殿變成買賣的地方!」

2)耶穌的憤怒對中的焦點.祂只是針對那些有罪的行為和 不公不義.

4)耶穌的憤怒有良好的自我控制.耶穌沒有情緒失控,即使處在憤怒中.殿中的祭司長並不喜歡祂潔淨聖殿(路加福音19:47), 祂控制了自己的情緒,而不是被情緒所控制.

5)耶穌的憤怒有合適的處理.祂沒有讓自己的忿恨轉化為怨恨或忿怒,沒有發大會在好的時機處理處理得當.

6)耶穌的憤怒有很好的結果.引人注意,引人關注,引人深思,    引出新知等!使門徒及跟隨者得學習.

耶穌的憤怒出於對神的敬虔,是不可避免的.正如祂的情緒,用來檢查神的話語.因此,耶穌的回應常用來完成上主的旨意.

從耶穌身上學習處理怒氣

心理學者教導我們處理情緒, 可以歸納出三個重要觀念,:

一.接受怒氣是一種情緒.

二.情緒需要紓解與控制.

三.和解或重建是最好的處理辦法.

耶穌三方面都做到了!

認識及處理憤怒階段

尋找火頭(火苗)

牧會有段時間,我好易發火,我要尋找火頭,找到了就消滅火苗!消防員臨近火場,最首要做的一件事,是要迅速找出火的源頭,火苗在那裡,並盡快樸滅,任由火苗橫行,會帶來嚴重後果.倘若我們經常發火,你需要清楚你的火苗在那裡?心理學家認為,怒火的源頭是:當人的某種需要未被滿足,或某種期望落空時,產生一種強熱的情緒,以表達不滿的情緒.怒氣是一種正常的情緒反應,可以產生強大的力量!然而這力量可以是正向的,也可能是負向的.

學者富勒曼隆格(Robert E. Freeman-Longo)和柯藍 (Murray Cullen)認為「憤怒爆發循環」是:

a「假裝正常階段」,

b「漸增憤怒階段」,

c「憤怒爆發階段」,

d「後悔或防衛機轉階段」四階段進行循環.

a假裝正常時憤怒是隱藏的,雖然表面沒問題但其實憤怒仍然存在,碰到任何刺激,就會爆發出來；

b漸增憤怒時,憤怒透過對刺激事件的認知與解讀而漸漸累積,並可能進行一些成癮行為,開始幻想或計劃發洩憤怒；

c憤怒爆發的階段則可能以怒罵,暴力,自虐等方式表現出來；

d後悔或防衛機轉階段則會後悔並啟動自我防衛機轉,否認自己有憤怒的問題,並可能產生「假悔悟 的情形去掩飾過錯.」

富勒曼隆格和柯藍建議應該將「憤怒爆發循環」調整為「憤怒管理循環」,包括:

a「正常階段」,

b「漸增憤怒及運用阻斷階段」,

c「憤怒管理階段」,

d「問題解決階段」.

猶太人問:你既作這事,還顯什麼神蹟(上帝的作為)給我們看?【你做這些事,叫我地點信你是神子?】

深入思想耶穌的回答:「你們所作所為正在拆毀上帝這聖殿,敬拜上帝的地方；雖然你們拆毀,不過我三日內會重建這上帝的殿,建立起信徒聚集的教會來,使它成為敬拜上帝的地方!這就是神蹟(上帝的作為)」

a正常階段不須假裝,而是承認自己有處置憤怒的問題,並且避開刺激情境；耶穌發出憤怒,沒有逃避,但祂卻避開刺激的情境,面對當權者,不提推倒買賣者一事,免引起火爆場面；對方也不提, 因當權者理虧!

b漸增憤怒階段會運用新的認知和解讀去阻斷憤怒的累積；

耶穌間接而溫和地指出他們所為正在拆毀聖殿,其實是好大罪呀!藉此抒發心中的不滿!

c憤怒管理階段會用有建設性,不具破壞性的方法去面對引起憤怒的問題；耶穌主要用言語指出在聖殿買賣之錯,針對事,非針對人!耶穌用繩子不用鞭子驅逐買賣者!不作破壞性的行為!

猶太人反問:「這殿是用了46年才建成,你三日內怎可以建成呢?」耶穌沒有回答,讓他們自己想!因回答只會起爭論!

其實耶穌在心中已答了他們「猶太人如果你信上帝,這便是我的神蹟(上帝在我身上的作為)了」.

在此,寫聖經的約翰補充說:耶穌說這話,是以祂的身體為聖殿,從死裡復活的基督,使為人敬拜的對像,那裡有信復活主的人聚集,那裡就有教會,那裡就成為聖殿,敬拜基督的所在了!

明顯地耶穌已使門徒明白其中真理了!

d問題解決階段,則可學習如何用更好的方式處置憤怒.

耶穌的方式是:引起注意,引起關注,引人反思,引出新知等!你也可以試下!

當我們憤怒時,往往難於自控和自制,或者常有一些不合理的焦點.我們做不到以上原則的一點或者多個點.這就是人的憤怒,雅各教導我們應該知道的是“你們各人要快快地聽,慢慢地說,慢慢地動怒.因為人的怒氣,並不成就神的義”(雅各書1:19-20).耶穌沒有表現出人的怒氣,而是上帝的聖怒.這聖怒已達到它的目的

[一位亨特Hunter弟兄的一段話]

如果你放任怒氣,就會變成一頭怪獸,使深愛你的人恐懼,甚至到傷害.消滅這頭怪獸的最好方法,就是你願意住在基督裡,讓基督的平安,和平,寬恕充滿你心!

當人,當事情,當環境達不到我們的期望時,我們會發怒!

當別人發怒時也會觸發我們心裡的怒火；當不公不義的事情發生在自己身上時,更會引發怒火!無論如何,記住,當你長久懷怒時,你正在懲罰自己,你的敵人或魔鬼知道時,他一定暗暗地發笑!

[開明心屝]

活過半百才發現,一個善於處理你怒氣的人, 是個幸福的人. 並會將幸福帶給他周圍的人! 相反, 任由怒氣瀑發的人, 終日生活在自己建造的火爐中, 生活在痛苦裡!

結束祈禱

感謝上主,賜人情緒,特是憤怒的情緒來抒發情感.祈求上主叫我們敢於發義怒,藉此維護正義,保護弱勢；亦能控制自己的怒氣,慢慢的動怒,智慧地抒發,也不長久懷怒,建立榮神益人的生命.奉主名求,阿們.



\newpage

\section{數不盡的精兵}
\label{sec:404}
\textbf{鄧瑞強博士 - 數不盡的精兵}
\newline
\newline
\hyperref[sec:403]{< < < PREV SERMON < < <}
~
\hyperref[sec:toc]{[返目錄]}
~
\hyperref[sec:405]{> > > NEXT SERMON > > >}


啟示錄 7:1-17
\newline
\begin{longtable}{cl}
\hline
\hline
章節 & 經文 (和合本修訂版)\\
\hline
7:1 & \begin{tabularx}{0.7\textwidth}{X} 此後，我看見四位天使站在地的四角，執掌地上四方的風，使風不吹在地上、海上和各種樹上。 \end{tabularx} \\ \\ \relax
7:2 & \begin{tabularx}{0.7\textwidth}{X} 我又看見另有一位天使從日出之地上來，拿著永生神的印。他向那得到權柄能傷害地和海的四位天使大聲喊著， \end{tabularx} \\ \\ \relax
7:3 & \begin{tabularx}{0.7\textwidth}{X} 說：「你們不可傷害地、海和樹林，等我們在我們神眾僕人的額上蓋了印。」 \end{tabularx} \\ \\ \relax
7:4 & \begin{tabularx}{0.7\textwidth}{X} 我聽見以色列人各支派中受印的數目有十四萬四千； \end{tabularx} \\ \\ \relax
7:5 & \begin{tabularx}{0.7\textwidth}{X} 猶大支派中受印的有一萬二千；呂便支派中有一萬二千；迦得支派中有一萬二千； \end{tabularx} \\ \\ \relax
7:6 & \begin{tabularx}{0.7\textwidth}{X} 亞設支派中有一萬二千；拿弗他利支派中有一萬二千；瑪拿西支派中有一萬二千； \end{tabularx} \\ \\ \relax
7:7 & \begin{tabularx}{0.7\textwidth}{X} 西緬支派中有一萬二千；利未支派中有一萬二千；以薩迦支派中有一萬二千； \end{tabularx} \\ \\ \relax
7:8 & \begin{tabularx}{0.7\textwidth}{X} 西布倫支派中有一萬二千；約瑟支派中有一萬二千；便雅憫支派中受印的有一萬二千。 \end{tabularx} \\ \\ \relax
7:9 & \begin{tabularx}{0.7\textwidth}{X} 此後，我觀看，看見有許多人，沒有人能計算，是從各邦國、各支派、各民族、各語言來的，站在寶座和羔羊面前，身穿白衣，手拿棕樹枝， \end{tabularx} \\ \\ \relax
7:10 & \begin{tabularx}{0.7\textwidth}{X} 大聲喊著說：「願救恩歸於坐在寶座上我們的神，也歸於羔羊！」 \end{tabularx} \\ \\ \relax
7:11 & \begin{tabularx}{0.7\textwidth}{X} 眾天使都站在寶座和眾長老，以及四個活物的周圍，俯伏在寶座前，敬拜神， \end{tabularx} \\ \\ \relax
7:12 & \begin{tabularx}{0.7\textwidth}{X} 說：「阿們！頌讚、榮耀、智慧、感謝、尊貴、權能、力量都歸於我們的神，直到永永遠遠。阿們！」 \end{tabularx} \\ \\ \relax
7:13 & \begin{tabularx}{0.7\textwidth}{X} 長老中有一位回應我說：「這些穿白衣的是誰？是從哪裡來的？」 \end{tabularx} \\ \\ \relax
7:14 & \begin{tabularx}{0.7\textwidth}{X} 我對他說：「我主啊，你是知道的。」他向我說：「這些人是從大患難中出來的，他們曾用羔羊的血把衣裳洗得潔白。 \end{tabularx} \\ \\ \relax
7:15 & \begin{tabularx}{0.7\textwidth}{X} 所以，他們在神寶座前，晝夜在他殿中事奉他；那坐在寶座上的要用帳幕覆庇他們。 \end{tabularx} \\ \\ \relax
7:16 & \begin{tabularx}{0.7\textwidth}{X} 他們不再飢，不再渴；太陽必不傷害他們，任何炎熱也不傷害他們， \end{tabularx} \\ \\ \relax
7:17 & \begin{tabularx}{0.7\textwidth}{X} 因為寶座中的羔羊必牧養他們，領他們到生命水的泉源；神必擦去他們一切的眼淚。」 \end{tabularx} \\ \\
[1ex]
\hline
\hline
\end{longtable}
講題:數不盡的精兵　Onward Christian Soldiers

經文:啟示錄7章1-17節

講員:鄧瑞強博士

場合:香港中文大學崇基學院禮拜堂

日期:2018年3月11日

各位弟兄姊妹,早安.

今是是預苦期(或「大齋期」,Lent)第四主日,是預備我們的心,面向基督的十架苦難.基督的十架苦難,源自人類的罪惡.基督甘受這苦難,是為了罪惡的人類.在這段期間,我們應多反省:我們生命的有限,我們罪惡的深不可測,我們對神的反抗,我們在破碎中對救贖的渴望.「預苦期」提示我們思考基督沉重的十架,也鼓勵我們在十架那裡找到不再哭泣的盼望.

今日的講道,講到在苦難世界中信徒的身分,及信徒在這苦難世界中可以有的盼望.

今日的講道經文:啟示錄7:1－17

啟7:1　此後,我看見四位天使站在地的四角,執掌地上四方的風,叫風不吹在地上,海上,和樹上.

啟7:2　我又看見另有一位天使,從日出之地上來,拿著永生神的印.他就向那得著權柄能傷害地和海的四位天使大聲喊著說:

啟7:3　「地與海並樹木,你們不可傷害,等我們印了我們神眾僕人的額.」

啟7:4　我聽見以色列人各支派中受印的數目有十四萬四千.

啟7:5　猶大支派中受印的有一萬二千；呂便支派中有一萬二千；迦得支派中有一萬二千；

啟7:6　亞設支派中有一萬二千；拿弗他利支派中有一萬二千；嗎哪西支派中有一萬二千；

啟7:7　西緬支派中有一萬二千；利未支派中有一萬二千；以薩迦支派中有一萬二千；

啟7:8　西布倫支派中有一萬二千；約瑟支派中有一萬二千；便雅憫支派中受印的有一萬二千.

啟7:9　此後,我觀看,見有許多的人,沒有人能數過來,是從各國,各族,各民,各方來的,站在寶座和羔羊面前,身穿白衣,手拿棕樹枝,

啟7:10　大聲喊著說:「願救恩歸與坐在寶座上我們的神,也歸與羔羊!」

啟7:11　眾天使都站在寶座和眾長老並四活物的周圍,在寶座前,面伏於地,敬拜神,

啟7:12　說:「阿們!頌讚,榮耀,智慧,感謝,尊貴,權柄,大力都歸與我們的神,直到永永遠遠.阿們!」

啟7:13　長老中有一位問我說:「這些穿白衣的是誰?是從哪裡來的?」

啟7:14　我對他說:「我主,你知道.」他向我說:「這些人是從大患難中出來的,曾用羔羊的血把衣裳洗白淨了.

啟7:15　所以,他們在神寶座前,晝夜在他殿中事奉他.坐寶座的要用帳幕覆庇他們.

啟7:16　他們不再飢,不再渴；日頭和炎熱也必不傷害他們.

啟7:17　因為寶座中的羔羊必牧養他們,領他們到生命水的泉源；神也必擦去他們一切的眼淚.」

啟示錄的第五章,講到天上有一藏著神的歷史大計的書卷,那書卷被七個印封著.第六章講到,那七個封印被逐一揭開,揭露出因著人類罪惡的緣故,大地受審判的景象.在第六章的結尾,問了一個問題:「神忿怒的大日到了,誰能站得住呢?」第七章就是要回答這個問題.

第七章設定的場景,仍在天上.第1至2節講到,大地在暴風雨的前夕,一切看來在傾覆的邊緣.當審判臨到大地的時候,誰能站得住呢?有些人,他們能站在風浪之上；這些人,得到神特別的保護.

啟7:3　「地與海並樹木,你們不可傷害,等我們印了我們神眾僕人的額.」

「印」,是主人的印.受了誰的印,生命就歸屬誰.生命總有所屬.按啟示錄的講法,你不屬於神,便屬於撒旦.你不屬於真理,便屬於虛謊.你不屬於正義,便屬於邪惡.在正邪之間,你不能模棱兩可.人總有所屬,問題只是屬於誰.世上有一些人,他們的靈魂刻著神的標記,他們屬於神,也受到神的看顧.

有時,在街上,會看見「仁愛傳教女修會」(Missionaries of Charity)的修女,她們穿著同德蘭修女一樣的,藍間白底的修女袍.見到她們,我都會肅然起敬.她們跟隨德蘭修女,在神面前立下誓言,除了一般修女的誓言外,她們更立誓,要全心服侍貧窮人中的至貧窮者.她們在溝渠中,抱起垂死的人,潔淨他們,照料他們.如此服侍他人的人,他們的生命,刻著神的印記.

啟7:4　我聽見以色列人各支派中受印的數目有十四萬四千.

14萬4千,當然不是一個實數,是一個象徵數目,代表所有真實的信徒.舊約,以色列民有12支派.新約,耶穌有12門徒.12乘12就是144.啟示錄喜歡以14萬4千這數目來代表神的全部子民.

一般而言,在聖經裡,數點人數,便是數點軍人,準備作戰.啟示錄的第8章講到號角,是作戰的號角.顯然,經文這裡,是讓信徒明白自己「軍人」的身分.信徒,基督徒,應是精兵,為真理而戰,為正義而戰,為真神而戰.但在基督教信仰群體裡,在一般教會裡,現在要多找幾個精兵,都不容易了.很多人返教會,是等上天堂.見不到他們在工作中,是真心造福他人.見不到他們在生活裡,跟一個迷信而自私的人有何分別.不少信徒的祈禱,是求神祝福他買的股票升值.神看來只是一個祝福人不勞而獲的神仙.不少信徒只是求神祝福他這樣那樣,看不到他求神賜他力量去祝福其他人.一個人只看到自己,只著緊自己的幸福,只求神祝福自己,怎可能是精兵?印度的甘地,有次這樣說:「我不抗拒基督,我愛基督.只是太多基督徒根本不似基督.若基督徒真按聖經所載的基督的教訓而活,全印度今天就都是基督徒了.」(“Oh, I don’t reject Christ. I love Christ. It’s just that so many of you Christians are so unlike Christ. If Christians would really live according to the teachings of Christ, as found in the Bible, all of India would be Christian today.”)孔子說:「一日克己復禮,天下歸仁焉.」對我們來說,或許,可以這樣講:「一日作主精兵,天下歸神焉.」沒有精兵的軍隊,必敗.沒有精兵的基督教,怎可能叫人相信.

啟7:9　此後,我觀看,見有許多的人,沒有人能數過來,是從各國,各族,各民,各方來的,站在寶座和羔羊面前,身穿白衣,手拿棕樹枝,

這是用另一角度去描寫14萬4千精兵,用另一角度去描寫神的子民,他們無窮無盡,來自五湖四海,四面八方.然後,斗轉星移,時間轉換,現在已是大戰完畢的光景.剛才,還是天地崩壞之際,正邪大戰一觸即發之時.轉眼間,已是戰事已完,高唱凱歌的日子.時間的轉換,是叫我們從另一角度看信徒的戰爭.

看看這群人,他們已卸下戰袍,身穿白衣,放下武器,手拿棕樹枝.

地點沒有變,這群人仍在天上,站在寶座和羔羊前.作為神的精兵,他們總是能站在生命的高處,低頭看人間.人間有善惡,有正邪,有真理與虛假之戰.他們不能置之度外,不能獨善其身.事實上,信徒活在世上,正正就是要在凶險的世道中行正義,講真話,顯大愛.他們要以戰士的勇毅投身在天國的建設中.但同時,在信仰裡,他們的心靈與主在一起,能從靈性的高度看塵世；在信仰裡,他們的視野能看到永恆,能從時間的盡處看當下.在信仰中,他們看到上帝至終的勝利.因著基督復活所顯示的勝利,他們在人間的艱難奮鬥中已預嘗天上的勝利.神保證他子民的勝利,他的子民在這種保證中奮勇地為真理作戰.

他們身穿白衣,這是勝利者的衣裳,純潔,正義,光明,這也是真信徒生命應有的氣質.人生的作戰,不是以「厚黑」取勝.要得勝,靠的是光明的人格.光明的人格本身不保證得勝,得勝,全因為神.他們手拿棕樹枝,棕樹枝代表勝利,凱旋,和平.主耶穌在世的最後一個禮拜,進入耶路撒冷時,人群也拿棕樹枝歡迎他,他像王一樣進入京城,後來,卻帶著荊棘王冠上十架,以十架的愛完成善惡之戰.若將這裡的棕樹枝連繫耶穌十架故事的棕樹枝,則這裡的棕樹枝代表著十架中的勝利.勝利,由苦難達成.生命,由死亡開展.

啟7:10　大聲喊著說:「願救恩歸與坐在寶座上我們的神,也歸與羔羊!」

天上的聖徒唱凱歌,他們明白,勝利,不是由於自己的功勞.他們能站在天上,能有分於凱旋的行列,全因為神的保護,神的引領,神的救恩.在信仰裡,人能得勝,能有所作為,能造福他人,毫無可誇之處.能得勝,全因為神.

啟7:11　眾天使都站在寶座和眾長老並四活物的周圍,在寶座前,面伏於地,敬拜神,

啟7:12　說:「阿們!頌讚,榮耀,智慧,感謝,尊貴,權柄,大力都歸與我們的神,直到永永遠遠.阿們!」

天使,天上的長老,神面前的神獸,都一起敬拜神.七重的讚詞,完全的讚美,全歸乎神.只有神,配受讚美.在這裡,沒有半點讚美,落入凡間.有時,我們會對老闆說:「老闆英明.」有時,我們會對老婆說:「老婆英明.」有時,我們會對人間的帝王說:「皇上英明.」但只要你說某人英明,將某人奉為主,則很易你自貶為奴隸,或他貶你為奴隸,在你生命裡打上他的印記,然後,任他魚肉.人間的罪惡,根源於此,就是有人自命為神,暴力地消滅他人的尊嚴,將他人玩弄於指掌之中.啟示錄斷然拒絕將任何人神化.

最後一個段落,交代天上這群精兵經歷了什麼,以什麼方式得勝,他們的前景如何等等問題.

啟7:13　長老中有一位問我說:「這些穿白衣的是誰?是從哪裡來的?」

啟7:14　我對他說:「我主,你知道.」他向我說:「這些人是從大患難中出來的,曾用羔羊的血把衣裳洗白淨了.

這群人是誰?他們「是從大患難中出來的」.一路走來,一步一十架.

這大患難是什麼?

按啟示錄的講法,患難源自人類自命為神,與撒旦握手,成為撒旦的爪牙,造成人間政治,經濟,社會上的欺壓,使人活得不似人.

天上的聖徒,在這光景中,如何為主作戰?

他們「曾用羔羊的血把衣裳洗白淨了」.這句何解?

他們是什麼也不用做,被動地等待主耶穌解救嗎?他們信了主,被動地由主耶穌的寶血洗淨,然後等上天堂嗎?看來不是.

由於啟示錄第六章講到「殉道士」,也講到天國未臨,全因甘心為神的真理殉道的精兵數目不夠,故此,這裡講,他們「曾用羔羊的血把衣裳洗白淨了」,就只能表示,這群人主動以「羔羊耶穌」流血的方式,將生命的崇高意義純淨地活出來.在患難中,在顛倒是非的世界,在將他人「整死」的制度裡,這群精兵,像主耶穌一樣,堅守天父的旨意,堅持對弱者的愛,堅定地將人從「吃人的制度」中救出來.為此,他們遭遇主耶穌的遭遇,在真理的十架上殉道.

就是這些人,真理的忠勇戰士,他們的生命在地上卻如同在天上,穿著白色的義袍,搖著勝利的棕樹枝,唱著凱旋的詩歌.

他們的前景是什麼?

啟7:15　所以〔因此〕,他們在神寶座前,晝夜在他殿中事奉他.坐寶座的要用帳幕覆庇他們.

啟7:16　他們不再飢,不再渴；日頭和炎熱也必不傷害他們.

啟7:17　因為寶座中的羔羊必牧養他們,領他們到生命水的泉源；神也必擦去他們一切的眼淚.」

他們的前景,是站在永恆的神面前.「坐寶座的要用帳幕覆庇他們」,從舊約的「會幕」,聖殿,到新約裡基督以自己的身體為殿,「帳幕」都指向神親自的臨在.神以自身的神聖,作為真信徒永恆的居所.人總要找安身立命的居所,有人在成就裡找,有人在學問裡找,神的子民卻在神聖的臨在裡,找到生命安頓的地方.

啟示錄繼續論述真信徒的至福.他們「不再飢,不再渴」,生命的基本需要,不再成為重擔,不再能逼令他們向不義低頭.「日頭和炎熱也必不傷害他們」,邪惡的嚴苛不再能虐待他們.他們找到「生命的活水」,「神也擦去他們一切的眼淚」,因為在天父寶座前的羔羊必牧養他們.

平平白白,啟示錄講出生命的真相.人生如戰場,善惡相遇,醜陋與優美對峙,卑污與高尚交戰.神呼籲我們,作精兵,不作逃兵,為生命的神聖而戰.這是一種靈性的戰鬥,非暴力,不用卑污的手段,而以生命的高潔,以純淨的真理,以無私的大愛去完成,像釘十架的基督一樣.同樣,也像復活的基督一樣,展望生命的永恆.這一切,全在乎我們信賴的上帝,也在乎我們對上帝的信賴.是他,親自牧養我們.是他,擦去我們一切的眼淚.

但願:頌讚,榮耀,智慧,感謝,尊貴,權柄,大力都歸與我們的神,直到永永遠遠.阿們!



\newpage

\section{棕枝主日的驢工}
\label{sec:405}
\textbf{李均熊博士 - 棕枝主日的「驢」工}
\newline
\newline
\hyperref[sec:404]{< < < PREV SERMON < < <}
~
\hyperref[sec:toc]{[返目錄]}
~
\hyperref[sec:406]{> > > NEXT SERMON > > >}


講題:棕枝主日的「驢」工　”Donkey” Work in Palms Sunday

講員:李均熊博士

場合:香港中文大學崇基學院禮拜堂

日期:2018年3月25日

1.) 引言:棕枝主日的「棕枝」

今天是教會傳統的棕枝主日,經課是關於耶穌在受死復活前的一週進入耶路撒冷的事蹟,為耶穌基督的受苦打開序幕.但今年經課年選讀的經文是出自馬可福音,當中並無記載任何的棕樹枝,何解今日會稱為「棕枝主日」?在四福音中,馬太,馬可以及路加福音均沒有提及棕樹枝,唯獨約翰福音記載在耶路撒冷預備過逾越節的人拿著棕樹枝出城迎接耶穌.反而馬可,或者馬太或路加,用上大部份篇幅提到耶穌吩咐門徒到對面村子尋找驢駒的事情,約翰只一句「耶穌找到了一匹驢駒」,便輕輕帶過.

稱呼這日作「棕枝主日」明顯不是基於對觀福音對耶穌預備進入耶穌路冷的情節,而是出於約翰福音對於迎接耶穌的群眾的描寫而有.對於後來教會要「慶祝」耶穌進入耶路撒冷,就很容易找來一些有像徵意義的道具去協助.試想想我們今日要找一隻驢進入禮拜堂會是如何困難!

約翰福音所記載關於耶穌進耶路撒冷的記載,所強調的是耶穌應驗舊約撒迦利亞先知九章所預言的耶路撒冷的王要坐著驢駒,和平的進入京城.約翰福音的焦點是上主之民迎接和平之君.但回到今天我們的經課,馬可福音卻沒有任何關於君王的說法,甚至也沒有引用撒迦利亞書去指出耶穌騎驢的象徵意義.在馬可十一章1-11節中,佔了一半以上的篇幅,是關於那隻驢是怎樣找回來的.今天就讓我們再思一下在棕枝主日中馬可這段記載的意義為何,又對我們基督徒的生命有甚麼啟發.

2.) 馬可的「驢」工

細看馬可福音,耶穌進入耶路撒冷的描寫其實十分反高潮,其他人看來歡天喜地,但「耶穌到了耶路撒冷,進入聖殿,看了周圍的一切.天色已晚,他就和十二使徒出城,往伯大尼去.」就是這樣看了看,就走了.什麼潔淨聖殿呢?舌戰律法師呢?那些精彩的「情節」都不在眾人注目耶穌進城的當天發生的.反而於入城前為著找一隻給耶穌騎進城內的驢駒,卻是馬可福音所關注的事情:

1耶穌和門徒快到耶路撒冷,來到伯法其和伯大尼,在橄欖山那裡.耶穌打發兩個門徒,2對他們說:「你們往對面村子裡去,一進去的時候會看見一匹驢駒拴在那裡,是從來沒有人騎過的,把牠解開,牽來.3若有人對你們說:『為甚麼做這事?』你們就說:『主要用牠,但會立刻把牠牽回到這裡來.』」4他們去了,看見一匹驢駒拴在門外街道上,就把牠解開.5在那裡站著的人,有幾個說:「你們解開驢駒做甚麼?」6門徒照著耶穌的話說,那些人就任憑他們牽去了.7他們把驢駒牽到耶穌那裡,把自己的衣服搭在上面,耶穌就騎上.(可十一1-7)

由指派門徒到實行吩咐,馬可福音的記載十分詳細,但當中很多記載其實對我們來說都十分費解的.究竟是耶穌有超自然的能力使事情順利,還是其實祂在更早之時已經預備妥當?究竟驢駒的主人是誰,他又是否和耶穌一起,以至「主要用牠」應譯作「驢的主人要用牠」呢?說到底,為什麼這樣詳盡記載「那裡站著的人」的反應?整段記載究竟有可重要?

正因為費解,有解經家認為馬可福音作者也不知到耶穌派門徒找驢駒的這段記載意義為何,只因為不願流失這傳統而記下.馬太省略了實行的細節,並和約翰福音一樣,把整件事連上撒迦利亞的預言,騎驢成為和平之君的記號.(馬太的改動更大,一隻驢變成兩隻驢).路加則和馬可相仿,把指派門徒和實行的細節詳細寫下,可能路加也只是純粹忠於史料「按著次序」的筆法.

但只要我們詳讀整卷馬可福音,我們就會發現馬可福音記載充滿著反諷(ironic)的手法.照字面解可能我們會一頭霧水,但比對馬可福音之前和之後的記述手法,會否作者其實迫使我們再思人子作為彌賽亞的意思是什麼?以敘事分析研究福音書稱著的學者Alan Culpepper,07年出版的馬可福音注釋書的一段說話正好讓我們玩味馬可福音的反諷語氣:

Nevertheless, the scene remains puzzling, unless Mark intended it to introduce the ironic incongruities that follow. Scripture is fulfilled in the course of everyday, inconsequential acts and by disciples who do not understand what they are doing. The Son of God enters the Holy City on a colt, acclaimed by a crowd, but by the end of the week another crowd will call for his death. The temple appears to be full of activity, but it is like a barren fig tree. The Messiah will be anointed, but he will be anointed not by the high priest but by a woman. When he is acknowledged as Son of God, it is again not the high priest but a Gentile, the Roman centurion, who voices the confession. Throughout, Mark’s account of Jesus’ death is filled with pathos and irony. Just as the Easter message will be announced by women who say nothing because they are afraid, so the account of Jesus’ last days in Jerusalem is introduced by a report of the mundane activities of disciples and bystanders who did not know the meaning of what they were doing. (R. Alan Culpepper, Smyth  and  Helwys Bible Commentary: Mark, p. 366f)

(大意:此事的確令人費解,除非馬可刻意諷剌以後事情的不協調:古經的應驗發生在日常而又毫不起眼的行為中,又透過完全不明其行動意思的門徒實行出來；人子坐著驢駒入聖城,被群眾簇擁,但一週後卻被另一班群眾高呼釘死祂；聖殿熱熱鬧鬧,但卻像枯槁的無花果樹；受膏者彌賽亞果然受膏,但卻非由大祭司進行,而是透過一個女人；耶穌的確被稱為神的兒子,不過也不是由大祭司宣告,而是透過外邦羅馬百夫長的告白.馬可描寫耶穌的死充滿悲愴和反諷!就如復活的信息將要由一班因為害怕而甚麼也不告訴人的婦女去宣講,耶穌在耶路撒冷的最後日子也要由門徒刻板的舉動以及一些不明就裡站著的局外人來引入.)

3.)Donkey work--承受苦悶又刻板的信仰生活

對於傳統認為耶穌是應驗撒迦利亞預言,以和平之君的身份進入聖城,這其實是後來的人反思而得的結論,就如約翰福音明明的寫道:「門徒當初不明白這些事,等到耶穌得了榮耀後才想起這些話是指他寫的……」(約十一16)在耶穌進城的當日,一如馬可福音的記載,大家其實仍然不明就裡,只是按吩咐而行.但反過來想,若然耶穌真是以君臨天下的姿態入城,在過節時期特派的羅馬駐防軍可能於耶穌一入城就將祂拘捕了.耶穌怎可以「出師未捷身先死」?不過,當時有接近十萬人來到耶路撒冷過節,一少群人在城門外的起哄未必會引來因節期而增派的羅馬軍的關注.這也解釋了馬可福音的細節,耶穌騎著的是一隻從別處村莊借來,又未經人騎過,亦即未受載人訓練的驢駒.這些令人費解的預備工夫,其實有可能是在應驗先知說法的同時,刻意隱藏祂作為彌賽亞的身份.約翰福音的注解就是,耶穌的時候未到,一切也不能說明.這也是學者認為馬可福音一貫保持的彌賽亞秘密.

我們可能在棕枝主日十分願意慶祝耶穌作為和平之君,榮耀的進入耶路撒冷,但我們又有沒有想過,耶穌的不被明白的那份孤單,面對門徒和群眾的一知半解的困惱,面對自己將要面對的荒謬世情的那種無奈?馬可福音要讓我們看見的,並不是一個大能和操控一切的大羅神仙,而是那甘於平凡,甚至甘於被誤解,那真正的人之子.為要成就救贖的工作,神子耶穌願意隱藏於平凡一般的背後.

作為耶穌的門徒,我們又應如何跟從這樣真正謙和的主耶穌呢?在耶穌來到耶路撒冷之前,馬可福音記載了門徒在將近進入聖城時對耶穌的要求:

35西庇太的兒子雅各和約翰進前來,對耶穌說:「老師,我們無論求你甚麼,願你為我們做.」36耶穌對他們說:「要我為你們做甚麼?」37他們對他說:「在你的榮耀裡,請賜我們一個坐在你右邊,一個坐在你左邊.」38耶穌對他們說:「你們不知道所求的是甚麼.我所喝的杯,你們能喝嗎?我所受的洗,你們能受嗎?」39他們對他說:「我們能.」耶穌對他們說:「我所喝的杯,你們要喝；我所受的洗,你們也要受.40可是坐在我的左右,不是我可以賜的,而是為誰預備就賜給誰.」41其餘十個門徒聽見,就對雅各和約翰很生氣.42耶穌叫了他們來,對他們說:「你們知道,外邦人有君王作主治理他們,有大臣操權管轄他們.43但是在你們中間,不可這樣.你們中間誰願為大,就要作你們的用人；44在你們中間誰願為首,就要作眾人的僕人.45因為人子來,並不是要受人的服事,乃是要服事人,並且要捨命作多人的贖價.」(可十35-45)

在我為這篇講道搜集資料時,我就在網上看到一位解經者有這樣的想像:耶穌吩咐的那兩個對對面村尋找驢駒的門徒,就是西庇太的兩個兒子,雅各和約翰.耶穌將要進入榮耀裡(當然這裡是指祂的受苦受死),這兩位門徒要在天國裡做我的左右承相嗎?給我去牽著驢來吧!但雅各和約翰在路上可能就會覺得,找牲口這些低下功夫為什麼找我們兩個去做?驢出名難服待,更可況未被人騎過的驢?我們本身是漁夫嘛!!這麼唐突要在別家村口牽走別人的驢駒,會否被人看為偷驢賊?!那時有人追問我們,耶穌的說法是否真的可以蒙混過去?到時我們應如何應對??

4.) 結語

信仰基督並不會天天精彩的:崇拜也未必每次令人興奮；靈修讀經祈禱也不一定有所得著；在教會服待也必須要承受很多Donkey Work (刻板平庸甚至乎沉悶卻又必須的工作).你會怎樣廁身信仰的這些沉悶刻板的時候?

馬可提醒我們的主耶穌並沒有佳形美容使我們羨慕,服侍祂的工作也不一定精彩偉大.在進入受難週的今天,或者就讓我們多些思想這位經常隱藏的彌賽亞,如何在刻板,平凡,甚至顯得有些荒謬的日常中,與我們同行,我們又可以如何服侍衪.



\newpage

\section{祂在我們以先}
\label{sec:406}
\textbf{戴浩輝牧師 - 祂在我們以先}
\newline
\newline
\hyperref[sec:405]{< < < PREV SERMON < < <}
~
\hyperref[sec:toc]{[返目錄]}
~
\hyperref[sec:407]{> > > NEXT SERMON > > >}


講題:祂在我們以先　He is ahead of us

講員:戴浩輝牧師

場合:香港中文大學崇基學院禮拜堂

日期:2018年4月1日

引言

這個多星期我的whatsapp群組每天都在報告和要求代禱,因為在神學院界有一位教授的兒子面臨死亡,我們都在懇切為他禱告.記得幾年前的同一個時間,我就送別了岳父返回天家.他的一生十分坎坷,他一生都信祖先庇佑,有時又相信人死如燈滅,有時人死後又會做鬼.幾年前糖尿病嚴重,又中風,不能走動,糖了上腳可能要折肢,於是他說,我死都不要折肢,因為我已失明,如果還要折肢,將來做鬼又盲又跛,做鬼都被人欺侮.人死後是否做鬼?死亡的另一面是怎樣?今天是復活節,是紀念耶穌基督從死裡復活,在這復活日中,讓我們藉上帝道使我們活在復活的生命中,而不是活在人世所延續的恐懼中.死亡是人生的一部分；聖經說:「人人都有一死,死後且有審判.」在死亡的陰霾下,我們又如何活得更豐盛,更精彩和有能力?在個問題的癥結就在於我們如何看人生,我們從死亡看抑或從永恆的生命看.今天我們歡慶主耶穌的復活,今天我們的生命可以不從死亡的角度看人生,因為我們的主已經復活了,成為我們這些必死的人將來復活的盼望.讓我們從可十六1-8中看上帝如何在我們以先為我們成就大事.

一) 在我們以先清除障礙

今天的福音說,七日的頭一日,婦女們買了為死人安葬用的香膏,膏耶穌的身體.但她們有一個困難,「誰給我們把石頭從墓門滾開呢?」這些婦女,並不因為這個難題而不來,難題是在她們前頭；或許,如果換了其他人,可能就因著這難題留在家中.但因著她們對主耶穌的愛,她們還是盡力的去做,去想辦法.因著愛,那大石不能阻止她們不來.

最能夠表示愛的能力是主耶穌為我們死在十字架上,祂對我們的愛是沒法比較的.各位,主耶穌愛你,祂為你死在十字架上,他又為你復活了,祂為你勝過死亡,祂要將永恆的生命賜給你,你相信嗎?祂愛你,你相信嗎?

我們愛一個人,不因為他已離開我們而失去這愛,我們的感情和記憶還是繼續對那人的情誼.從復活的事件中,我們可以活得更豐盛,因為我們有主耶穌的愛,不僅是石頭不能使屬主的人不行動,甚至是死亡,也不能使這愛的力量減低.各位弟兄姊妹,我們有沒有在生活中感到無力?連返教會的心力都不逮,主愛我們,祂的愛使我們繼續努力以赴.我們人生中當然有不少大石小石,有的阻礙我們進步,有的是在我們心理中做成障礙.但我們可以肯定的:主愛我們,祂在我們以先已把它一一的挪去,我們只要相信和仰望,我們必能得著復活的主所賜的能力.

二) 在我們以先得勝死亡

最近英國有一位偉人去世,一位只靠電腦與外界溝通的霍金教授,他熱愛生,相比起另一位英國全身嚴重癱瘓的男子尼克林森,要求法院准許醫生合法結束他的生命,有強烈的對比.2005年,當時58歲的尼克林森因中風患上了閉鎖綜合症.這種疾病使他雖然神智清醒,卻全身癱瘓.他一直尋求法庭保護任何願意幫助他結束生命的醫生.尼克林森形容自己目前的生活單調,痛苦,沒有尊嚴和無法忍受.尼克林森的妻子說,他爭取合法結束他的生命,獲得廣大公眾支持.

當人無法有意識的經驗生命,或與其他人有意義的關聯時,死亡臨到其身上可能是件好事.但對大多數人來說,死亡是我們最後的敵人.保羅說:「儘末了所毀滅的仇敵就是死」(林前十五26).當人到中年,自己愈來愈近死亡,漸漸想到不要買太多書,因為我可能無法把它們讀完.我也可能無法再見我想見的人,可能無法再去我想去的地方,可能無法做我想做的事.死亡是最後的敵人,特別可從我們目送至親進入棺材時感到.我們內心呼喊著可否多一年,一個月,一日,一小時,甚至一刻?可以給我對他說「我愛你」?可以收回自己做錯的事或說錯的話?但事實告訴我們:「不可能」.在幾年前,曾在一個科學資訊電視節目中看到,人可能長生不老,原因是科學家正在研究,人類細胞在分裂時其DNA的密碼鏈會失去一些末梢部分,最後這鏈因為過短會停止分裂.科學家認為,如果能把控制末梢的酉每增強,密碼鏈就不會斷,人就能有不斷的新陳代謝,長生不老.但最近,科學家又有新的發現和理論.在世界中生物細胞分為線狀和圓環兩樣,有單細胞是圓環形的,也是長生不老的,但壞處是沒有改良的能力,沒有進步的能力.而世界上大多數生物則屬於線狀,壞處是會有一天斷到不能再斷,即是有死亡的一日,好處就是可以隨環境改進,人類的生命就是這樣的發展.科學家斷言,就算人類可以控制酉每,使密碼鏈不斷,這細胞也遲早變為癌細胞,人還是遲早死.

但感謝我們的上帝,靠著耶穌從死裡復活,把我們接到生命之源的上帝那裡.在我們以先主復活了,保羅說:基督已經從死裡復活成為睡了之人初熟的果子(林前十五20).我將會肉身死亡,但我仍相信,我要在末後復活,因為基督耶穌已經復活了 !

三) 在我們以先成就旨意

復活的主能帶給我們能力,聖經說他比門徒先往加利利去.為何要到加利利?耶穌是在加利利長大的,祂想返鄉下?加利利一向是外邦人的地方和象徵,那麼與耶穌和門徒約定在加利利又有何關係?在十四27-28這樣說:「耶穌對他們說:你們都要跌倒了,因為經上記著說:我要擊打牧人,羊就分散了.但我復活以後,要在你們以先往加利利去.」門徒的跌倒與耶穌往加利利去有何關係呢?

a. 從十五37-39馬可表達了主的死代表世界的拯救,而加利利既是外邦事工的開始,復活的主就在門徒以先往那裡去了.換言之,加利利與耶穌基督的約定就是大使命傳福音的開始.誰需要福音?就是所有未相信的外邦人,未聽聞耶穌基督的人.

b. 在十四27預言門徒跌倒與耶穌復活後在加利利又有何關係?復活的主在門徒以先已到加利利去,主在那裡使門徒重新得力.他們最終能從極其害怕的情況變為不怕死的見證者.

c. 復活的主今天在我們以先已成就祂在我們身上的旨意,你相信嗎?正如今天我們在這復活主日,我們相信主已復活了嗎?我父家全家都是基督徒,縱然我們幾兄弟的信仰都有出入,我有一個弟弟長期受佛教的影響,在他看來,他死後是否復活並不重要,他只看重在這人生中,他能否活得快樂和有意義.他很重親情,我們母親一生要我們信耶穌,我們就聽從了,而這一位弟弟也聽從,每主日都返教會,在教會做執事,服事人,他所看重的是今生能有意義的生活.在他心中來生並不重要,但這不能阻礙他在基督裡服侍其他人.我們到死的那一日也可能懷疑,我身體已火化了,我還可以復活嗎?我只能相信主耶穌基督的話,他宣告說:「復活在我,生命也在我,信我的人雖然死了也必復活.」(約十一25)甚至到這一刻,我們大家都可能與馬可福音九章25節一個被鬼附身的孩子的父親的呼喊:「我信,求主幫助我的不信 !」縱然今天我們還懷疑,但我們仍能藉祂所賜的聖靈愛人,服事有需要的人；我們甚至不知道,我們能做到這些,但最終是主耶穌在我們以先成就祂在我們身上的旨意.

回到我的岳父,我師母照顧他的生活,也照顧他的靈性,她駕車接送看病,她與她分享信仰,正如她向她妹妹傳福音,最後,我就在家裡為我的岳父施行洗禮,並由家附近的教會幫助他的信仰上的成長 .當他的生命指數下跌到一個不能接受的地步時,師母在他耳邊講述他將要看見,將要行走,將要與主耶穌一起時,生命指數特別高升,還未信的兒子也感到力量的大.最後人終有一死,但藉死亡,他與主一起,因為主耶穌已在他以先得勝死亡,在他以先為他準備永生的家,為他和將來我們一家相聚.

結語

感謝上帝,祂藉基督凡事在我們以先為我們開路,去除障礙,成就旨意,賜予我們永生.



\newpage

\section{天上靜然無聲半小時}
\label{sec:407}
\textbf{鄧瑞強博士 - 天上靜然無聲半小時}
\newline
\newline
\hyperref[sec:406]{< < < PREV SERMON < < <}
~
\hyperref[sec:toc]{[返目錄]}
~
\hyperref[sec:408]{> > > NEXT SERMON > > >}


啟示錄 8:1-6
\newline
\begin{longtable}{cl}
\hline
\hline
章節 & 經文 (和合本修訂版)\\
\hline
8:1 & \begin{tabularx}{0.7\textwidth}{X} 羔羊揭開第七個印的時候，天上寂靜約有半小時。 \end{tabularx} \\ \\ \relax
8:2 & \begin{tabularx}{0.7\textwidth}{X} 我看見那站在神面前的七位天使，有七枝號賜給他們。 \end{tabularx} \\ \\ \relax
8:3 & \begin{tabularx}{0.7\textwidth}{X} 另有一位天使拿著金香爐來，站在祭壇旁邊；有許多香賜給他，要和眾聖徒的祈禱一同獻在寶座前的金壇上。 \end{tabularx} \\ \\ \relax
8:4 & \begin{tabularx}{0.7\textwidth}{X} 那香的煙和眾聖徒的祈禱從天使的手中一同升到神面前。 \end{tabularx} \\ \\ \relax
8:5 & \begin{tabularx}{0.7\textwidth}{X} 天使拿著香爐，盛滿了壇上的火，倒在地上；就有雷轟、響聲、閃電、地震。 \end{tabularx} \\ \\ \relax
8:6 & \begin{tabularx}{0.7\textwidth}{X} 拿著七枝號筒的七位天使預備好要吹號。 \end{tabularx} \\ \\
[1ex]
\hline
\hline
\end{longtable}
題目:天上靜然無聲半小時 Silence in Heaven for about Half an Hour

經文:啟示錄8章1-6節

講員:鄧瑞強博士

場合:香港中文大學崇基學院禮拜堂

日期:2018年4月8日

各位弟兄姊妹,早安.

今日講道的經文:啟示錄8:1-6

啟8:1 羔羊揭開第七印的時候,天上寂靜約有二刻(即半個鐘).

啟8:2 我看見那站在神面前的七位天使,有七枝號賜給他們.

啟8:3 另有一位天使,拿著金香爐來,站在祭壇旁邊.有許多香賜給他,要和眾聖徒的祈禱一同獻在寶座前的金壇上.

啟8:4 那香的煙和眾聖徒的祈禱從天使的手中一同升到神面前.

啟8:5 天使拿著香爐,盛滿了壇上的火,倒在地上；隨有雷轟,大聲,閃電,地震.

啟8:6 拿著七枝號的七位天使就預備要吹.

1957年,上映了一齣電影,是瑞典導演英瑪褒曼(Ingmar Bergman)拍的《第七封印》(The Seventh Seal),這電影現在已成為經典.電影的開場白,便是剛才讀的經文的第一句:「羔羊揭開第七印的時候,天上寂靜約有半個鐘」.這句話點出電影的一個主線,當死亡的恐怖包圍整個世界時,神沉默得令人吃驚.

電影的場景,設定在十四世紀的瑞典.當時黑死病蔓延,四周都是死人.隱伏在各種事物背後的,看來不是上帝,而是死神.主角是一個剛剛參加完十字軍的騎士,他看來是一個甘願為上帝戰鬥至死的人.但他看來覺得這場仗無聊至極,他與他的僕人,在回家路上.在一個清晨,在海邊,他看見死神.其實,死神一直在他身邊,現在,死神要帶他走了.騎士說,他其實隨時準備死,但只是希望在死前,做件有意義的事.於是,騎士建議,與死神下棋.只要棋局未完,他就仍能活在世上.死神答應了.於是,整齣電影常常見到棋局.人生如一盤有終結的棋局,主角藉棋局爭取時間,在大地漫遊,尋找生命的意義.

由於死亡籠罩大地,各人以不同方式去驅散死亡的威逼.有群眾列隊,用自虐的方式去贖罪,用鞭鞭打自己,用膝頭跪著行,用荊棘刺穿額頭.有人熱衷捉女巫,抓住無辜婦女,指控她交鬼,以致帶來瘟疫,這是找弱者作代罪羔羊.有人在死人中找財物,瘟疫是他發達的機會.有神父以死亡威嚇人,藉此抬高自己的神聖地位.形形式式,林林總總,各人在懸崖邊拼命抓住支撐身體的一根草.

主角卻不斷問,生命的意義何在?神存在嗎?一切到頭來是否只是空無?他說:信仰是一種折磨,就像你愛一個人,但這個人在黑暗中,永不露面；無論你叫他叫得多大聲,他都沒有回應.

最後,騎士遇上一個小家庭,丈夫叫Jof(或許,這暗示約瑟),太太叫Mia(或許,這暗示馬利亞),還有一個嬰孩.他們是小人物,巡迴做小丑表演.他們唱著歌,吃著野草莓,談論著兒子將來會變成怎樣,看來,只有他們,在死亡張狂的世界,談論著明天的希望.我想,騎士在他們身上,找到生命的平靜.

有一次,騎士又和死神下棋.騎士意識到,死神可能盯上Jof一家,於是,他故意碰跌棋子,轉移了死神的視線,讓Jof一家及時離開.這時,死神下了一步關鍵的棋,騎士輸了.死神對騎士說,你的時候到了,你能做到件有意義的事嗎?騎士說,做到了.

當上天的審判籠罩大地,死亡的力量在周圍肆虐時,我們還有希望嗎?

這是今天講道經文要回答的問題.

啟示錄8:1-6的經文,呈現「三明治」結構.我們逐層來看.

(1)「三明治」結構的最外圍:第1-2節及第6節,審判快要開始

啟8:1 羔羊揭開第七印的時候,天上寂靜約有二刻(即半個鐘).

啟8:2 我看見那站在神面前的七位天使,有七枝號賜給他們.

啟8:6 拿著七枝號的七位天使就預備要吹.

這三節經文是啟示錄8:1-6的「三明治」結構的最外圍,前後都講到「七位天使」及「七枝號」.我們知道,羔羊耶穌握著的書卷,是天父對歷史及人生的終極心意,這書卷被七個印封住.現在,第七個印打開了,卻引來了「七枝號」.這時,天上靜然無聲半小時.

天上為何寂靜了?

或許,天上的詩班歌聲太大,聽不到人間受苦者的祈禱.或許,靜一靜,讓所有人去預期,去注視世事下一步的發展.或許,是一種禮儀性的,獻祭前的安靜,讓人注視神.

為何是半個鐘頭?

我想,這與啟示錄的時間象徵有關.啟示錄常常講「三年半」,這是「七年」的一半.「7」若是圓滿的時間,則「3.5」便是事情去到一半.事情在發展中,但終局已被神安排好了.「半個鐘頭」,我想,也有這種含意.人間越來越險惡,堅持真理的信徒越來越苦.神設定的終局雖然遙遙在望,但艱苦的信仰旅程只行到一半.「一半」,意味著事情在進展中,終局可期,但仍需信徒努力.

打開了七印,我們得知神的心意.神的心意,是「七號」.「七號」指向神對罪惡更深的審判,神不容認罪惡.「號」,也是爭戰的信號.在罪惡受審判的日子,罪惡的力量會反撲,人間會暗無天日.在這時,信徒要更加心志堅定,努力奮戰.

(2)「三明治」結構的內圍:第3節上半節及第5節,審判臨到大地了

啟8:3 另有一位天使,拿著金香爐來,站在祭壇旁邊.……

啟8:5 天使拿著香爐,盛滿了壇上的火,倒在地上；隨有雷轟,大聲,閃電,地震.

「另有一位天使」,顯示這是另一段落的開始.這段落有「香爐」,「祭壇」.

「祭壇」,是獻祭的地方.啟示錄6:9說:「我看見在祭壇底下,有為神的道,並為作〔真理〕見證被殺之人的靈魂.」敬拜上帝,明白真理的人,在這祭壇上,獻上了生命.古往今來,多少堅持真理的人,都在祭壇上,獻上了自己.早幾日(4月4日),是美國民權領袖馬丁路德金(Martin Luther King Jr.)遇刺50周年的紀念日子.他在那個黑人不可以有夢想的年代,高聲說:「我有一個夢」.在他的夢裡,他看見一個平等友愛的世界.然後,一粒子彈便要了他的命.他倒在生命的祭壇上.

經文這裡提到天使拿著「香爐」,這應像是現在天主教或聖公會崇拜時,事奉人員手裡拿的「香爐」.這「香爐」被幾條長鏈吊著,裡面點著乳香.事奉人員搖著「香爐」,發出乳香的芬芳,代表著神聖的臨在.「香爐」原是優雅的物件,然而,啟示錄8:5說:「天使拿著香爐,盛滿了壇上的火,倒在地上」.「香爐」成了審判的器具.這裡有三樣物件,「香爐」關乎真敬拜,「壇上的火」代表獻祭,「倒在地上」指向審判.這三種東西合在一起,意義是什麼?或許,意思是:人間最可貴的東西,正是真敬拜,神的真正臨在,人人不做「王」,不自稱「神」,這是「香爐」所代表的.然而,人間的混亂與邪惡,正是人人要做王,人人要做神,人人都拒絕真神.假話消滅真話,劣幣驅逐良幣,惡人逼害善人,這是「壇上的火」所代表的.神在天上看著這一切,不能無動於衷,他要主持正義,他要審判罪惡,他要創造新天新地,這是「倒在地上」所代表的.簡言之,以「香爐」代表的真神,以「獻在壇上」的真信徒為寶貴,以「倒在地上」的審判,去懲治邪惡,伸張正義.

當「壇上的火」倒在地上時,「隨有雷轟,大聲,閃電,地震」.這是神臨在的記號,也是「終末」的記號.神要將歷史的終局帶來.

簡單來說,神在人間,他看見不義,他看見「有勢有力」者對「無權無勢」者的欺壓,他看似沉默不言,毫無反應,事實卻是,神要為一切討回公道,回復正義,帶來萬物的終局.他等待時機,他也希望人學曉信靠與忍耐,更重要的,他要信徒配合他的行動.

如何配合?這便是今日講道經文的「三明治」結構的中心部分.

(3)「三明治」結構的中心:第3節下半節及第4節,神真正的心意

啟8:3 ……有許多香賜給他,要和眾聖徒的祈禱一同獻在寶座前的金壇上.

啟8:4 那香的煙和眾聖徒的祈禱從天使的手中一同升到神面前.

這裡的重點是:香,聖徒的祈禱.

啟示錄5:8說過:「這香就是眾聖徒的祈禱」.聖經有視「祈禱」如「香」的傳統(如:詩141:2).但這裡的經文說:「有許多香賜給他(天使)」,明顯,這香由神所賜.看來,這裡的「香」不是聖徒的祈禱.經文說,這「香」要和聖徒的祈禱,一同升到神面前.看來,這「香」代表著神可見的臨在,代表著一種上升的力量,代表著承托生命的恩典.信徒如何知道自己的祈禱已到達神面前?啊,有從神而來的「香」,護送著信徒的祈禱到達神面前.是神主動施予的恩典,保證我們在人生的哀求,能上達天庭.神會聽祈禱嗎?神承托我們生命的心意,保證他會聽我們的祈禱.我們如此軟弱無力的祈禱,有效嗎?由神而來的「香」,會伴隨我們的禱告,達到神的耳中.

最後一個要處理的問題是:這些聖徒的祈禱是什麼?

我們在啟示錄中去找答案.

其中一個可能是啟示錄6:10的祈禱,那裡是殉道士的哀求:「聖潔真實的主啊,你不審判住在地上的人,給我們伸流血的冤,要等到幾時呢?」若這就是聖徒祈禱的內容,那麼,神要施行的審判,便是對這祈禱的回答.

不過,我想,聖徒的祈禱應就是第7章眾聖徒的呼喊,啟示錄7:10說:「願救恩歸與坐在寶座上我們的神,也歸與羔羊!」這不是哀求,是認信,是敬拜,是完全的信賴.若這理解是對的話,則在黑暗,死亡,審判重重包圍下,信徒生命的中心,仍是平靜安穩.

如此,經文的用意便清楚了.今日的講道經文,是「三明治」結構.外圍是重重的審判,但中心卻是信徒生命向世界展示的救恩,平靜,安穩.我們知道,啟示錄的審判結構,又分三重,就是:七印,七號,七碗.七號在中間,重重包圍著七號的,是七印和七碗的審判.今日的講道經文,是七號,外圍是更多的審判,但在這所有審判的核心,就是聖徒的祈禱.神其實不愛審判,神愛人向他祈禱.神不愛罪人死亡,他愛罪人悔改得生命.中心不是死亡,而是生命.

我開頭時,講到《第七封印》那電影.那位騎士,他生活的世界,被黑死病重重包圍,身邊每個人都沾染著死亡的氣息,而他也常與死神博奕.他看見的東西,盡是空虛混沌,毫無意義.最後,只希望在輸給死神之前,做件有意義的事情.最後,他做了,他故意碰跌棋子,轉移了死神的視線,讓一家善良的人安然離去.這位騎士,以這一簡單善舉,安頓了自己的心.死神帶他走之前,他作了一個死前祈禱:「憐憫我們,因我們又渺小,又驚恐,又無知.」

在這看來意義不彰的世界,神的聲音隱約難辨,在死亡面前,或許,我們也只能如那騎士那樣祈禱.但是,啟示錄卻教導我們,在黑暗,死亡,審判重重包圍下,請像聖徒一樣祈禱,相信神的救贖,相信神聽祈禱,相信神能主持正義.從而,在生命的核心,保持盼望,信賴,安穩.

但願:榮耀歸於聖父,聖子,聖靈.阿們.



\newpage

\section{向南走向北走}
\label{sec:408}
\textbf{林豪恩先生 - 向南走,向北走}
\newline
\newline
\hyperref[sec:407]{< < < PREV SERMON < < <}
~
\hyperref[sec:toc]{[返目錄]}
~
\hyperref[sec:409]{> > > NEXT SERMON > > >}


使徒行傳 8:26-40
\newline
\begin{longtable}{cl}
\hline
\hline
章節 & 經文 (和合本修訂版)\\
\hline
8:26 & \begin{tabularx}{0.7\textwidth}{X} 有主的一個使者對腓利說：「起來！向南走，往那從耶路撒冷下迦薩的路上去。」那路是曠野。 \end{tabularx} \\ \\ \relax
8:27 & \begin{tabularx}{0.7\textwidth}{X} 腓利就起身去了。不料，有一個埃塞俄比亞人，是個有大權的太監，在埃塞俄比亞女王甘大基的手下總管銀庫，他上耶路撒冷去禮拜。 \end{tabularx} \\ \\ \relax
8:28 & \begin{tabularx}{0.7\textwidth}{X} 回程中，他坐在車上，正念著以賽亞先知的書， \end{tabularx} \\ \\ \relax
8:29 & \begin{tabularx}{0.7\textwidth}{X} 聖靈對腓利說：「你去！靠近那車走。」 \end{tabularx} \\ \\ \relax
8:30 & \begin{tabularx}{0.7\textwidth}{X} 腓利就跑到太監那裡，聽見他正在念以賽亞先知的書，就說：「你明白你所念的嗎？」 \end{tabularx} \\ \\ \relax
8:31 & \begin{tabularx}{0.7\textwidth}{X} 他說：「沒有人指教我，怎能明白呢？」於是他請腓利上車，與他同坐。 \end{tabularx} \\ \\ \relax
8:32 & \begin{tabularx}{0.7\textwidth}{X} 他所念的那段經文是這樣：「他像羊被牽去宰殺，又像羔羊在剪毛的人手下無聲，他也是這樣不開口。 \end{tabularx} \\ \\ \relax
8:33 & \begin{tabularx}{0.7\textwidth}{X} 他卑微的時候，得不到公義的審判，誰能述說他的身世？因為他的生命從地上被奪去。」 \end{tabularx} \\ \\ \relax
8:34 & \begin{tabularx}{0.7\textwidth}{X} 太監回答腓利說：「請問，先知說這話是指誰，是指自己，還是指別人呢？」 \end{tabularx} \\ \\ \relax
8:35 & \begin{tabularx}{0.7\textwidth}{X} 腓利就開口，從這段經文開始，對他傳講耶穌的福音。 \end{tabularx} \\ \\ \relax
8:36 & \begin{tabularx}{0.7\textwidth}{X} 二人正沿路往前走，到了有水的地方，太監說：「看哪！這裡有水，有甚麼能阻止我受洗呢？」 \end{tabularx} \\ \\ \relax
8:37 & \begin{tabularx}{0.7\textwidth}{X} 腓利說：『你若是一心相信，就可以。』他回答：『我信耶穌基督是神的兒子。』 \end{tabularx} \\ \\ \relax
8:38 & \begin{tabularx}{0.7\textwidth}{X} 於是他吩咐把車停下來，腓利和太監二人一同下到水裡，腓利就給他施洗。 \end{tabularx} \\ \\ \relax
8:39 & \begin{tabularx}{0.7\textwidth}{X} 他們從水裡上來，主的靈把腓利提了去，太監再也看不見他了，就歡歡喜喜地上路。 \end{tabularx} \\ \\ \relax
8:40 & \begin{tabularx}{0.7\textwidth}{X} 後來有人在亞鎖都遇見腓利；他走遍那地方，在各城宣揚福音，一直到凱撒利亞。 \end{tabularx} \\ \\
[1ex]
\hline
\hline
\end{longtable}
講題:向南走,向北走　Go South, Go North

經文:使徒行傳8章26-40節

講員:林豪恩先生

場合:香港中文大學崇基學院禮拜堂

日期:2018年4月29日

你或你的父母是從北南下而來到這裡的嗎?或者,你會曾經向遠方的北走,走到北美或北歐,然後再向南回流到這裡呢?向南走,向北走,很可能在你的人生中留下很多難忘的故事,有些故事可能改變你的人生方向.在你向南走,向北走的故事中,可能與其他人的向北走,向南走的故事交錯相遇,令故事更精彩而深刻.今天的經課--使徒行傳8章26-40節敘述了一個向南走向北走的故事,另一個向北走向南走的故事,以及兩者相遇而產生出來的第三個故事.

經文

8:26有主的一個使者對腓利說:「起來!向南走,往那從耶路撒冷下迦薩的路上去.」那路是曠野.

8:27腓利就起身去了.不料,有一個埃塞俄比亞人,是個有大權的太監,在埃塞俄比亞女王甘大基的手下總管銀庫,他上耶路撒冷去禮拜.

8:28回程中,他坐在車上,正念著以賽亞先知的書,

8:29聖靈對腓利說:「你去!靠近那車走.」

8:30腓利就跑到太監那裡,聽見他正在念以賽亞先知的書,就說:「你明白你所念的嗎?」

8:31他說:「沒有人指教我,怎能明白呢?」於是他請腓利上車,與他同坐.

8:32他所念的那段經文是這樣:「他像羊被牽去宰殺,又像羔羊在剪毛的人手下無聲,他也是這樣不開口.

8:33他卑微的時候,得不到公義的審判,誰能述說他的身世?因為他的生命從地上被奪去.」

8:34太監回答腓利說:「請問,先知說這話是指誰,是指自己,還是指別人呢?」

8:35腓利就開口,從這段經文開始,對他傳講耶穌的福音.

8:36二人正沿路往前走,到了有水的地方,太監說:「看哪!這裡有水,有甚麼能阻止我受洗呢?」

8:38於是他吩咐把車停下來,腓利和太監二人一同下到水裡,腓利就給他施洗.

8:39他們從水裡上來,主的靈把腓利提了去,太監再也看不見他了,就歡歡喜喜地上路.

8:40後來有人在亞鎖都遇見腓利；他走遍那地方,在各城宣揚福音,一直到凱撒利亞.

三．釋經與應用

1．腓利:「向南走,向北走」

向南走

8:26有主的一個使者對腓利說:「起來!向南走,往那從耶路撒冷下迦薩的路上去.」那路是曠野.8:27腓利就起身去了.

向南走,是腓利的故事,也是不少香港人的故事.

「他17歲考入金城銀行,獲派在集團內的天津航運公司服務,時為1928年.由於才幹出眾,不久獲華北航運界推舉為天津航運業公會副會長.1940年成立民營的中國航運信託公司,掛上英國及巴拿馬旗幟,航行中國沿海和遠東口岸.大陸1949年易主後,他將基地遷往香港,成立金山輪船國際有限公司.隨香港,台灣,日本及東南亞地區的經濟於五,六十年代蓬勃發展,進出口貿易大增,他的海上王國同步急速擴張.1969年成立東方海外,並在1973年上市.(蘋果日報2017年07月09日)這是香港首任特首董建華的父親董浩雲的故事.

「1956年5月中旬,他連夜騎馬趕到泰來縣火車站,乘坐齊齊哈爾往大連的旅客列車到達大連,再乘船往上海……他只好繼續南逃,歷盡艱辛,多次靠吃老鼠,螞蟻,棉花充飢,走了三個月的路到達廣州,路上還發揮其刻製印章的技術瞞過關防人員,才成功逃離大陸,後經澳門,於1957年7月成功偷渡到香港.

這是被喻為「香港四大才子」之一倪匡的故事.

1945年,香港人口50萬,1950年220萬,1980年510萬,差不多以每十年100萬的速度高速增長.這些增長的人口,其中不少是如上述「向南走」的故事.這些並不是遙遠的故事,可能就是我們的祖輩,父輩,叔伯兄弟,甚或是我們自己的故事.

向南走是不是單程路呢?南方是不是終點呢?

8:39主的靈把腓利提了去……8:40後來有人在亞鎖都遇見腓利；他走遍那地方,在各城宣揚福音,一直到凱撒利亞.

腓利向南走,後來向北走；曾經何時,人們向南走來香港,後來又回頭向北走.

香港政府統計處網頁有以下一段文字:「不少港人都經常因公務而須往返於中港兩地之間.然而,究竟有多少香港人曾在內地工作呢?他們在內地又是從事甚麼工作的呢?根據政府統計處在二零零三年一月至三月(15年前)期間進行的一項統計調查顯示,約有238 200名香港居民在統計前十二個月內曾往中國內地工作……當中,男性佔78.1\%.(可想而之產生了甚麼社會及家庭問題)此外,與全港就業人口比較,他們的教育程度亦普遍較高,其中有35.2\%更具專上教育程度.2003年的十年之後的狀況又有甚麼不同呢?2012年的電影「春嬌與志明」告訴大家,北上的已經不只是男性,還有年輕女性；北上之處已經不只是珠三角和長三角,已經抵達北京.

北上,不只是春嬌與志明的故事,可能是我們的家人,朋友,甚至是自己的故事.

向南走,向北走,有多少是自主的選擇?有多少是勢所難免呢?腓利的旅程似乎是不由自主的,故事開始時指出8:26「有主的一個使者對腓利說:「起來!向南走,往那從耶路撒冷下迦薩的路上去.」那路是曠野.」腓利向北走,更是「身不由己」:「8:39主的靈把腓利提了去……8:40後來有人在亞鎖都遇見腓利」.縱使可能是「身不由己」,然而,不等於沒有意義,腓利向南走向北走的旅程是有意義的.我們這些向南走向北走的人,我們的旅程又有甚麼意義呢?

2．銀庫總管:「向北走向南走」

腓利「向南走向北走」的故事包含著另一個「向北走向南走」的故事.與腓利的故事不同,這個向北走向南走的行程應該是自願的.

向北走

8:27不料,有一個埃塞俄比亞人,是個有大權的太監,在埃塞俄比亞女王甘大基的手下總管銀庫,他上耶路撒冷去禮拜.

向南走

8:28回程中,他坐在車上

8:31他說:「沒有人指教我,怎能明白呢?」於是他請腓利上車,與他同坐.

8:39他們從水裡上來,主的靈把腓利提了去,太監再也看不見他了,就歡歡喜喜地上路.

埃塞俄比亞女王甘大基的銀庫總管北上,所為何事呢?據經文描述,他向北走,是往耶路撒冷去禮拜.長途跋涉往耶路撒冷去禮拜,我們可以想像他是一個有所追求的人,北上去尋找一些在他處身的南方所欠缺的東西.他北上了,又南下回去,他找到了嗎?他滿足了嗎?在宗教的處境來看,人們往宗教場所去禮拜,多多少少懷著敬虔的心,盼望得到一些甚麼.禮拜後,離開時,得到了嗎?心靈滿足了嗎?還是如埃塞俄比亞女王甘大基的銀庫總管懷著「沒有人指教我,怎能明白呢」這種淡淡的無奈而離去呢?

回看這城曾經走過的路,最初製造業北上,後來影藝娛樂事業北上,銀行金融會計界北上,餐飲服務業北上,醫療界也北上,政治人物也紛紛北上.向北走,很可能為了尋找一些東西,希望得到一些東西.有些人得到財富,有些人得到權力,然而,等於找到人生的安穩,心靈的滿足了嗎?還是如故事中那位有財富又有權力的埃塞俄比亞女王甘大基的銀庫總管一樣,想做的事情已經做了,但卻感到仍欠缺甚麼似的,若有所失地迷茫而回呢?

在人生的旅程中,難免有類似的無奈和迷茫.不過,也不一定是這樣,就算是這樣,也可以只是過程,並不是結局.按經文在39節描述,他最後是「歡歡喜喜地上路」.故事中的銀庫總管遇到甚麼人,發生了甚麼事,在旅程中從「沒有人指教我,怎能明白呢?」的迷茫,轉變為「歡歡喜喜地上路」的輕省呢?這可能是他向北走向南北的故事,與腓利向南走向北走的故事相遇而併發出來的新故事.

3．「向北走向南走」,「向南走向北走」兩個故事的相遇

8:28回程中,他坐在車上,正念著以賽亞先知的書,

8:29聖靈對腓利說:「你去!靠近那車走.」

8:30腓利就跑到太監那裡,聽見他正在念以賽亞先知的書,就說:「你明白你所念的嗎?」

8:31他說:「沒有人指教我,怎能明白呢?」於是他請腓利上車,與他同坐.

銀庫總管是認真的人,希望盡自己的努力,撥開雲霧,為自己的狀況尋找出路.閱讀聖經,代表著透過知識認識信仰；力求明白,代表著透過理性接觸信仰.然而,知識無窮無盡,理性有所不及,始終,好像還有欠缺.所欠缺的是甚麼呢?或者,就如祈克果所說的「信心的跳躍」.這是一份投入,好像游泳者不只是站在池邊學習,而是投入水中暢泳.信靠上帝的人必然對上帝有某程度的認識,同時,也必然是讓自己投進上帝的懷中,讓上帝懷抱著自己.

8:36二人正沿路往前走,到了有水的地方,太監說:「看哪!這裡有水,有甚麼能阻止我受洗呢?」

8:38於是他吩咐把車停下來,腓利和太監二人一同下到水裡,腓利就給他施洗.

在岸上觀察,認識,理解的銀庫總管,終於跳入水裡了.

黎志添教授:《宗教研究與詮釋學》一書中的總結說:「我們認為研究者無論如何同情地投入宗教信仰者的經驗世界,也不能忘記他始終是研究者而不是宗教信仰者.研究者的旁觀者身分更是清楚的表現在研究者的研究作品成果上.當研究者作為宗教經驗的旁觀者時,即使他投入宗教信仰者本身的經驗與意識,他的研究成果仍然是帶著研究者作為旁觀者的眼光與價值去看宗教.」

8:39他們從水裡上來,主的靈把腓利提了去,太監再也看不見他了,就歡歡喜喜地上路.

縱使腓利不見了,銀庫總管卻「歡歡喜喜地上路」,因為這次相遇,並沒有隨著分開而煙消雲散；這次相遇,留下了讓銀庫總管「歡歡喜喜地上路」的來源,填補了銀庫總管縱使擁有財富和權力仍若有所欠的洞穴.那是甚麼呢?在向南走向北走的故事中,甚麼是中心呢?

4．向南走向北走的故事中心

8:32他所念的那段經文是這樣:「他像羊被牽去宰殺,又像羔羊在剪毛的人手下無聲,他也是這樣不開口.

8:33他卑微的時候,得不到公義的審判,誰能述說他的身世?因為他的生命從地上被奪去.」

8:34太監回答腓利說:「請問,先知說這話是指誰,是指自己,還是指別人呢?」

8:35腓利就開口,從這段經文開始,對他傳講耶穌的福音.

在寫作結構上,這是故事的中心；對基督信仰者而言,這也是信仰的中心.

有一位從大陸往北美從事西方法學研究的學者名叫劉同蘇,他有如此的體會:「文化基督徒」往往把信仰視作一種理念形態,卻不知信仰就是生命,而生命就是「活著」.」他說:「我以前就是一個「文化基督徒」……念很多神學的書……在理性的範圍內思考和基督的關係……但是,觀念上的改變並不可以叫我全人發生改變.真正全人的改變是怎麼樣的呢?除非有活生生的生命進入,也以活生生的生命去接受那個活生生的生命,那才是真正轉變的關鍵.」(劉同蘇:從「文化基督徒」到有生命的基督徒)

三．總結

曾經,向南走來到這個小島的人,以為只是稍避亂局,結果一年又一年,結婚生子,漸漸變為落地生根.又曾經,以為已經落地生根的人,又連根拔起,向地球另一半更遠的南方,或者更遠的北方走去.或者,從北而向南來的,仍然會繼續向南走,或者終會向北走,「向南北向北走」的故事沒有終結.這個故事提示我們,在「向南走向北走」的旅程中,甚麼才是我們的終極關懷呢?甚麼才能令我們安身立命呢?帶著甚麼「向南走向北走」與其他「向南走向北走」的人們相遇才有意義呢?

人們說,要強大,要富有,才能解決問題,才能滿足需要.

埃塞俄比亞女王甘大基的手下總管銀庫應該是有財富及權力的人,但不要以為有財富及權力就能解決問題,就滿足人的需要.故事用「有大權的太監」及「在埃塞俄比亞女王甘大基的手下總管銀庫」來描述的人的欠缺,正是由寂寂無名,無權無物的小人物所解決和滿足的.

對一般人來說,並不是要追求強大和富有,只是盼望能保住自由及自主,擔心受到限制及壓迫.將來如何,沒有人能說得定.

然而,不要以為只有自由自主者才有能力去解決問題,才能滿足人的需要,故事中那位按自己意願向北走向南走的銀庫總管的欠缺,正是被那位不由自主地「向南走向北走」的腓利所解決和滿足.若我們再追查腓利的過去,他正是在耶路撒冷被逼害而要四處走動的逃難者呢!

找住信仰的中心,成為人生的中心,無論向南走,或是向北走,都不失去方向；無論是自主的旅程,還是被動的旅程；無論得到甚麼,失去甚麼,都不失去人生的美善.讓我們以幾句記載在希伯來書13:14勉勵初期的信徒的話互相提醒:

「在這裡,我們本沒有永存的城,而是在尋找那將要來的城……16只是你們不可忘記行善和分享,因為這樣的祭物是上帝所喜悅的.」

這是難能可貴的.願上主恩待我們.



\newpage

\section{流淌不住的水像神聖的洗禮}
\label{sec:409}
\textbf{伍渭文牧師 - 流淌不住的水,像神聖的洗禮}
\newline
\newline
\hyperref[sec:408]{< < < PREV SERMON < < <}
~
\hyperref[sec:toc]{[返目錄]}
~
\hyperref[sec:410]{> > > NEXT SERMON > > >}


使徒行傳 10:44-48
\newline
\begin{longtable}{cl}
\hline
\hline
章節 & 經文 (和合本修訂版)\\
\hline
10:44 & \begin{tabularx}{0.7\textwidth}{X} 彼得還在說這些話的時候，聖靈降在一切聽道的人身上。 \end{tabularx} \\ \\ \relax
10:45 & \begin{tabularx}{0.7\textwidth}{X} 那些奉割禮的信徒和彼得同來，見聖靈的恩賜也澆在外邦人身上，就都驚奇； \end{tabularx} \\ \\ \relax
10:46 & \begin{tabularx}{0.7\textwidth}{X} 因聽見他們說方言，稱讚神為大。於是彼得回答： \end{tabularx} \\ \\ \relax
10:47 & \begin{tabularx}{0.7\textwidth}{X} 「這些人既受了聖靈，跟我們一樣，誰能阻止用水給他們施洗呢？」 \end{tabularx} \\ \\ \relax
10:48 & \begin{tabularx}{0.7\textwidth}{X} 他就吩咐奉耶穌基督的名給他們施洗。於是他們請彼得住了幾天。 \end{tabularx} \\ \\
[1ex]
\hline
\hline
\end{longtable}
詩篇 98:1-9
\newline
\begin{longtable}{cl}
\hline
\hline
章節 & 經文 (和合本修訂版)\\
\hline
98:1 & \begin{tabularx}{0.7\textwidth}{X} 你們要向耶和華唱新歌！因為他行過奇妙的事，他的右手和聖臂施行救恩。 \end{tabularx} \\ \\ \relax
98:2 & \begin{tabularx}{0.7\textwidth}{X} 耶和華顯明了他的救恩，在列國眼前顯出公義； \end{tabularx} \\ \\ \relax
98:3 & \begin{tabularx}{0.7\textwidth}{X} 記念他對以色列家的慈愛和信實。地的四極都看見我們神的救恩。 \end{tabularx} \\ \\ \relax
98:4 & \begin{tabularx}{0.7\textwidth}{X} 全地都要向耶和華歡呼，要揚聲，歡唱，歌頌！ \end{tabularx} \\ \\ \relax
98:5 & \begin{tabularx}{0.7\textwidth}{X} 用琴歌頌耶和華，用琴和詩歌的聲音歌頌他！ \end{tabularx} \\ \\ \relax
98:6 & \begin{tabularx}{0.7\textwidth}{X} 用號筒和角聲，在大君王耶和華面前歡呼！ \end{tabularx} \\ \\ \relax
98:7 & \begin{tabularx}{0.7\textwidth}{X} 願海和其中所充滿的澎湃，願世界和住在其間的發聲。 \end{tabularx} \\ \\ \relax
98:8 & \begin{tabularx}{0.7\textwidth}{X} 願大水拍掌，願諸山在耶和華面前一同歡呼； \end{tabularx} \\ \\ \relax
98:9 & \begin{tabularx}{0.7\textwidth}{X} 因為他來要審判全地。他要按公義審判世界，按公正審判萬民。 \end{tabularx} \\ \\
[1ex]
\hline
\hline
\end{longtable}
約翰一書 5:1-6
\newline
\begin{longtable}{cl}
\hline
\hline
章節 & 經文 (和合本修訂版)\\
\hline
5:1 & \begin{tabularx}{0.7\textwidth}{X} 凡信耶穌是基督的，都是從神生的；凡愛生他之神的，也必愛從神生的。 \end{tabularx} \\ \\ \relax
5:2 & \begin{tabularx}{0.7\textwidth}{X} 我們愛神，又實行他的命令，由此就知道我們愛神的兒女了。 \end{tabularx} \\ \\ \relax
5:3 & \begin{tabularx}{0.7\textwidth}{X} 我們遵守神的命令，這就是愛他了，而且他的命令並不是難守的。 \end{tabularx} \\ \\ \relax
5:4 & \begin{tabularx}{0.7\textwidth}{X} 因為凡從神生的就勝過世界；使我們勝過世界的就是我們的信心。 \end{tabularx} \\ \\ \relax
5:5 & \begin{tabularx}{0.7\textwidth}{X} 勝過世界的是誰呢？不就是那信耶穌是神兒子的嗎？ \end{tabularx} \\ \\ \relax
5:6 & \begin{tabularx}{0.7\textwidth}{X} 這藉著水和血而來的，就是耶穌基督，不是單用水，而是用水又用血，並且有聖靈作見證，因為聖靈就是真理。 \end{tabularx} \\ \\
[1ex]
\hline
\hline
\end{longtable}
約翰福音 15:9-17
\newline
\begin{longtable}{cl}
\hline
\hline
章節 & 經文 (和合本修訂版)\\
\hline
15:9 & \begin{tabularx}{0.7\textwidth}{X} 我愛你們，正如父愛我一樣；你們要常在我的愛裡。 \end{tabularx} \\ \\ \relax
15:10 & \begin{tabularx}{0.7\textwidth}{X} 你們若遵守我的命令，就會常在我的愛裡，正如我遵守了我父的命令，常在他的愛裡。 \end{tabularx} \\ \\ \relax
15:11 & \begin{tabularx}{0.7\textwidth}{X} 「我已對你們說了這些事，是要讓我的喜樂存在你們心裡，並讓你們的喜樂得以滿足。 \end{tabularx} \\ \\ \relax
15:12 & \begin{tabularx}{0.7\textwidth}{X} 你們要彼此相愛，像我愛你們一樣，這是我的命令。 \end{tabularx} \\ \\ \relax
15:13 & \begin{tabularx}{0.7\textwidth}{X} 人為朋友捨命，人的愛心沒有比這個更大的了。 \end{tabularx} \\ \\ \relax
15:14 & \begin{tabularx}{0.7\textwidth}{X} 你們若遵行我所命令的，就是我的朋友。 \end{tabularx} \\ \\ \relax
15:15 & \begin{tabularx}{0.7\textwidth}{X} 以後我不再稱你們為僕人，因為僕人不知道主人所做的事；但我稱你們為朋友，因為我從我父所聽見的一切都已經讓你們知道了。 \end{tabularx} \\ \\ \relax
15:16 & \begin{tabularx}{0.7\textwidth}{X} 不是你們揀選了我，而是我揀選了你們，並且派你們去結果子，讓你們的果子得以長存，好使你們奉我的名，無論向父求甚麼，他會賜給你們。 \end{tabularx} \\ \\ \relax
15:17 & \begin{tabularx}{0.7\textwidth}{X} 我這樣命令你們，是要你們彼此相愛。」 \end{tabularx} \\ \\
[1ex]
\hline
\hline
\end{longtable}
講題:流淌不住的水,像神聖的洗禮　The Flowing Water, Like the Holy Baptism

經課:使徒行傳10:44-48;詩篇98:1-9;約翰一書5:1-6;約翰福音15:9-17

講道經文:馬可福音1章1-9節

講員:伍渭文牧師

場合:香港中文大學崇基學院禮拜堂

日期:2018年5月6日(復活期第六主日)

引言:今天是復活期第六主日,復活節為第一主日.教會年分兩部分,上半年有關基督,下半年有關基督徒,連結兩部分應是關鍵的一天,那天不是復活日,是五旬節.要是基督復活日是教會年的高潮,五十天的復活期便成了反高潮.耶穌復活後升天前四十天對門徒顯現,教導他們有關天國的道理,因為耶穌救贖事件的高潮不在他的死和復活,乃在聖靈降臨,教會的誕生.將臨期的燭光(candlelight),到顯現期成為耀光(brilliance),到被釘十架從正午到申初成為黝黑,到復活日變成正午日照(Midday sun),在聖靈降臨日成為火焰(flame),落在門徒的身上,教會開始了,這才是基督救贖事件的高潮.耶穌復活後用四十天教導門徒,就是這天國的奧秘.

在十架的基督,像母親孤單去經歷生產之苦痛,但母親並不孤單,而且充滿喜樂,因為她知道這生產之苦帶來生命的誕生.基督被棄絕,帶來了神人的修和,和友愛群體的誕生:「我要將你的名傳與我的弟兄,在會中我要讚美你.」(詩二十二 22) 希伯來書把這句話用在教會身上,但在前面加了:「原來那為萬物所屬,為萬物所本的,要領許多的兒子進榮耀裡去……所以,他稱為他們為弟兄也不以為恥,說:我要將你的名傳與我的弟兄,在會中我要頌揚你.」(來二10-12)希伯來書的作者著跟說:「兒女既同有血肉之體,他也照樣親自成了血肉之體,特要藉著死敗壞那掌死權的,就是魔鬼,並要釋放一生因怕死而為奴僕的人.」(來二14, 15)

天國的奧秘就是在十架孤單被棄絕的基督,要帶來愛的群體,這也是復活期一連五個主日,讀出約壹書的主題-彼此相愛.今天福音書約十五9-17,把彼此相愛的命令,連繫到三一上主之間的愛:「我愛你們,正如父愛我一樣,你們要常在我的愛裡.你們若遵守我的命令,就常在我的愛裡；正如我遵守了我父的命令,常在他的愛裡.」(約十五9,10)耶穌到約但河接受施洗約翰的洗,顯出他對父的愛和順服；在約但河的洗禮,耶穌順服被父差遣作為彌賽亞的命運,為世人受死贖罪.

1約旦河流淌的水,把我們的罪洗滌潔淨

今天題目:流淌不住的水,像神聖的洗禮,來自看過一首英文詩.早期教會有教規十二使徒遺訓(Didache, The Teaching of Twelve Disciples)說到洗禮首選是要到流動的水中,因耶穌在約但河領洗,約但河的水流淌不住.耶穌經過近三十年的沉寂,在約旦河接受洗禮後,便以彌賽亞的身份出來公開傳道,趕鬼,醫病的工作,最後被釘在十架,為世人贖罪.

(1)約旦河洗禮不是按立耶穌為基督.耶穌洗禮時,有什麼異常景象呢?“他從水裡一上來就看見天裂開了,聖靈彷彿鴿子降在他身上.又有聲音從天上來說,你是我的愛子,我喜悅你.”(可一10,11)這是聖經唯一同時間出現父,子,靈.父做見證:“你是我的愛子,我喜悅你.”這句話是彌賽亞的名號:“看哪,我的僕人,我所扶持,所揀選,心裡所喜悅的,我已將我的靈賜給他,他必將公理傳給外邦.”(賽42:1)耶穌從水裡上來時,聖靈像鴿子降在耶穌身上.

根據這段經文,早期異端就提出嗣子論:穌這人蒙上帝揀選,在約旦河受施洗約翰悔改的一刻,被聖靈充滿成為彌賽亞.在此之前,耶穌是一個敬畏上主的人,在洗禮後被上主選立為基督,天人合一,洗禮是按立為基督的禮儀.這是異端,因為不相信基督道成肉身,賢士們對聖嬰獻禮敬拜就是褻瀆上帝,怎可以當人為上帝呢?

(2)這水是洗滌人的罪.經文說到聖靈彷彿像鴿子降在他身上.聖經在此之前,何處曾說到從水裡上來,有鴿子降在其上呢?是當洪水漫過大地四十天后,水止住了,大地露出來,大地從水裡上來,鴿子從方舟放出去,嘴裡含著新擰下來的橄欖葉子.(創8:6-12)耶穌站到約旦河,不是被按立為彌賽亞,他生下來就是彌賽亞.他到約旦河受施洗約翰悔改的洗禮,是表明他要為人擔當罪孽.父說:“你是我的愛子,我喜悅你.”是指向耶穌的命運.同樣的話,在顯現節最後的福音經課出現,當耶穌登山變象時,有聲音從雲彩裡出來說:“這是我的愛子,你們要聽他的.”(可9:7)聽什麼呢?“人子必須受許多的苦,被長老,祭司,文士棄絕,並且被殺,過三天復活.”(可8:31)

耶穌踏進約旦河接受施洗約翰的洗,當流淌的水漫過其身時,就是為世人接受審判.若我是電影導演,這時候會用分割鏡頭,把流淌的河水變作洶湧澎湃,卷為巨浪,匯成洪水,引入南亞海嘯的鏡頭,把整個大地淹沒,這是上主洪水的審判.然後把一隻方舟拋起,在水裡劇烈激蕩,但方舟安然無恙.然後經過四十天洪水漸漸退去.露出旱地,有鴿子從方舟放出去,回來的時候,嘴裡含著橄欖葉,這時候天上出現彩虹,審判成為過去,一切有了新的開始.這時候,鏡頭一轉,耶穌從水裡上來,那鴿子降在耶穌身上.然後播出音樂:彩虹下的約定:

我 空虛的心靈

終於不再流淚

期待著雨後 繽紛的彩虹

訴說你我的約定

我 不安的腳步

終於可以停歇

主你已為我 擺設了生命的盛宴

與你有約 是永恆的約

彩虹為證 千古不變

我要高歌為 生命喜悅

萬物歌頌你的慈愛

大地訴說你的恩典

所以,彼得前書三20-21就以洪水漫過方舟比喻信徒的洗禮:“借著水得救不多,只有八個人,這水所表明的洗禮,借著基督的復活,也拯救你們.”約旦河流淌的水,帶給我們歡樂,因我們被潔淨.在顯現日耶穌洗禮的禱文中,提及耶穌在約旦河潔淨他的新婦教會.”今天教會與他天上的配偶結合了,因為基督在約旦河洗去了他的罪惡；博士們帶著禮物迅速前來,奔赴王家的婚宴；變水為酒,使賓客歡心,哈利路亞.”

顯現節是歡樂的節日,我們的聖帷是白色或金色,是聖誕期的高潮和終結.屈梭多模(金口約翰John Chrysostom)主教說:”今天基督進入了約旦河,以洗滌世人的罪過,今天上主的聲音,響徹水面,你是我的愛子,我喜悅你.今天,聖靈借鴿子降臨在水面上,報告人類永遠滅亡的災禍已停止.」”(每日頌禱[Liturgia Horarum]一,頁308.)

今天詩篇經課九十八1-9篇也洋溢這歡樂:「你們要向耶和華唱新歌!因為他行過奇妙的事,他的右手和聖臂施行救恩.2耶和華顯明瞭他的救恩,在列國眼前顯出公義；3記念他對以色列家的慈愛和信實.地的四極都看見我們上帝的救恩.4全地都要向耶和華歡呼,要揚聲,歡唱,歌頌!5用琴歌頌耶和華,用琴和詩歌的聲音歌頌他!6用號筒和角聲,在大君王耶和華面前歡呼!7願海和其中所充滿的澎湃,願世界和住在其間的發聲.8願大水拍掌,願諸山在耶和華面前一同歡呼；9因為他來要審判全地.他要按公義審判世界,按公正審判萬民.」審判的洪水,因救贖成拍掌的大水.

看到流水,不再單單慨嘆:「逝者如斯,不舍晝夜；我們不能把腳放進同樣流動的水中兩次.」我們明白現今我們仍然伏在律法之下,因著罪我們的日子短少；早晨發芽生長,晚上割下.看到流淌的水,我想起靈魂被洗滌潔淨.我們的肉體雖然被歲月磨蝕,但裡面的新人卻日漸壯大,因為在約旦河中我們已經與基督一起埋葬,也與他一起復活.

不少聖詩以約旦河喻為死亡的寒河,因為要跨過約旦河,才能進入迦南美地.今天崇拜最後一首荒原行旅歌(普天頌讚,1978版,456首) 第三節:”到時足踐約旦河濱,求壯我膽免心驚,滅死威權,搓死鋒芒,使我安抵迦南境.”馬丁路德1519年講章預備死亡說:”水禮是上主可見的記號,人必須像列祖雅各行經約但河時,以堅強的信心緊緊握著.這杖引導我們,定睛直視死蔭的幽暗,罪惡和地獄.……這記號是教恩的幫助,不會撒謊欺哄我.上主已說了,上主不會欺騙.”水禮是雅各歸家的柺杖,幫他渡過約但河,回到列祖應許之地.我們有一天都要歸家,渡過死亡的冷河,水禮是天路客的杖,支撐著我們行走天路.我們不害怕約旦河的寒流,因為基督在約旦河的洗禮流,已改變了我們的命運.

2流淌不住的水,成為別人的祝福

流淌的水是不會停留的.加利利海和死海其實都是兩個湖,加利利湖匯入的水也匯出,魚類眾多,生氣渤渤；死海納而不放,一條魚也沒有,因為海水不是流淌不住.在約旦河基督已潔淨我們,我們也要成為別人的祝福.我們怎樣領受,也要怎樣付出.西元325年尼西亞會議通過一法典.在船上,若有慕道者生命危殆,而沒有主教或長老神職人員,受了洗的信徒,只要不是多妻,可以為他施洗,上岸後報備教會.特土良也說明這是古訓:”如何接受,也可以如何付出”.在宗教改革運動初期,馬丁路德強調信徒藉洗禮已受按立成為君尊的祭司,承擔各樣事奉.1520年致德國諸侯書,就以此鼓勵德國諸侯,起來維護聖經福音信仰的純正,投進改教運動.

我喜愛一輯CD系列“流淌的歌聲”,這系列以通俗與美聖潔合,舊歌翻唱,期望唱片“像小河淌水,它有著清晨露珠般的清澈,有著夜色般的朦朧,聽著它-讓它進入心境如夢的境界.”就讓我們的生命像小河淌水,滋潤人的生命.在基督受洗的事上,我們可以從其身上學到兩樣東西,成為別人的祝福.

3 流淌不住,蓄勢後發.流淌不住是行動,行動可以先發制人,也可以後發制人.耶穌三十歲才領洗,叫我們學習打好根基,蓄勢待發.三十年是不短的光陰,三十年韜光養誨,用三十年預備作三年多的工,誰也按捺不住.耶穌的榜樣,教我們專注當下的責任,不要祇顧即時的滿足.耶穌三十年才公開工作,三年多的工作,大部分專注訓練十二位門徒,到了人生最後一星期,才騎進耶路撒冷,公開與敵人衝突.在三十歲受洗之前,他在十二歲上聖殿時,已經表明以父的事為念,但他安靜在家中,度過寂寂無名的日子,因為時候還沒有到.有人說他要擔起約瑟離世後的責任.聖經沒有交待穌三十歲前的事蹟,除了十二歲上聖殿,聖經祇說了一句:「耶穌的智慧和身量,並神和人喜愛他的心都一起增長.」

莊家熟了,要趕快收割.但用多一點時間把刀磨好,收割的時候用少了時間,收獲更多.剛去世的葛培理有一次在史陶德(John Stott)牧師主持的教牧研討會做嘉賓致辭,公開承認讀書不夠,若可以再來,會花90\%時間讀書,10\%時間工作.今天你會否因為責任有被困的感覺?等候忍耐,接受磨煉,預備承擔更重要的責任.船停在港彎,是整固帆杆,預備糧草,等待風起雲湧時,我們揚帆遠去.

4 願意給人服事,肯定別人的服事

耶穌在約旦河接受約翰的洗禮,是盡諸般的義.他肯定約翰的事奉,不想太早出來.當看到別人被上主使用時,你會為他們鼓掌嗎?我欣賞終身學習,孜孜不倦的鄭德音老牧師.他雖身為香港中國播道總會監督,但仍常常出席公開的教牧研討會,拿著筆記本,與我們晚輩一起學習.你可能覺得鄭牧師自己念書不多,以勤補拙.但相信大家對崇基學院前院長,香港中文大學前文學院院長沈宣仁教授的學問不會懷疑,深深佩服.我跟沈教授見面不多,但幾次單獨見面談起各種話題,他總是蠻有興趣,認真聆聽,有時表示讚賞,不是為你做結論指點你.

跟沈宣仁前院長談話有一種平等的感覺,也有被建立的感覺.他不會以學養來擊倒(overpower)你,反而他專心聆聽,誠懇提問,展示他推崇的蘇格拉底式的叩問求真,但使人感到受尊重,被充權(empower).耶穌在施洗約翰手下受洗,肯定是施洗約翰事奉的高峰,得到極大的鼓舞.他深受感動,喊出:「他必興旺,我必衰微」；耶穌謙卑,令施洗約翰更加謙卑更謙卑.

我們要學習服事人,也要學習被人服事.但接受人的服事比較困難,1980年前我們放下香港牧會工作,參與母校師資培訓計劃,到美國深造神學,接受一些教會,甚至我們不認識的教會經濟上的支持.推薦我加入計劃的恩師張子華牧師 (已安息主懷)鼓勵我說:「我們需要恩典,才能施予；我們也需要恩典,才能接受施予.」至今我還牢記這句話.所以聖經教導我們要彼此服事,互相代求.

結論:1912年春天,五十歲的哲學家桑塔亞納George Santayana)母親過世,承繼了一筆遺產.有一天他在哈佛大學授課時,一隻知更鳥(Robin)飛來站在窗檻上, 他看了這鳥片刻,有所感悟,回過頭來對學生說:“我與陽春有約”.接著宣佈下課,然後向學校辭職,退隱著書.耶穌三十歲來到約但領洗,當他從水裡上來時,聖靈彷彿像一隻鳥(鴿子)降在他身上,他看到這鴿子,知道歷史與他有約,他走上十字架；被釘,埋葬,第三天復活,四十日升天,再過十天,「五旬節到了,門徒都聚集在一處.忽然,從天上有戀聲下來,好像一一陣大風吹過,充滿了他們所坐的居子；又有舌頭如火焰顯出來,分分開落在他們各人的頭上.」(使1-3)

這大響與火焰,告訴門徒,更新萬有的陽春-上主的國降臨了,他們與這陽春有約,這是彩虹下的約定,我們藉基督承繼了永恆的產業(弗一11),而領受的聖靈,就是產業的憑證(弗一14).



\newpage

\section{重述福音}
\label{sec:410}
\textbf{鄧瑞強博士 - 重述福音}
\newline
\newline
\hyperref[sec:409]{< < < PREV SERMON < < <}
~
\hyperref[sec:toc]{[返目錄]}
~
\hyperref[sec:411]{> > > NEXT SERMON > > >}


啟示錄 11:1-19
\newline
\begin{longtable}{cl}
\hline
\hline
章節 & 經文 (和合本修訂版)\\
\hline
11:1 & \begin{tabularx}{0.7\textwidth}{X} 有一根蘆葦，像丈量的杖，賜給我；且有話說：「起來！將神的殿和祭壇，以及在殿中禮拜的人，都量一量。 \end{tabularx} \\ \\ \relax
11:2 & \begin{tabularx}{0.7\textwidth}{X} 只是殿外的院子不用量，因為這是要給外邦人的；他們將踐踏聖城四十二個月。 \end{tabularx} \\ \\ \relax
11:3 & \begin{tabularx}{0.7\textwidth}{X} 「我要賜權柄給我那兩個見證人，穿著粗麻衣說預言一千二百六十天。」 \end{tabularx} \\ \\ \relax
11:4 & \begin{tabularx}{0.7\textwidth}{X} 他們就是那站在世界之主面前的兩棵橄欖樹和兩個燈臺。 \end{tabularx} \\ \\ \relax
11:5 & \begin{tabularx}{0.7\textwidth}{X} 若有人想要害他們，就有火從他們口中噴出來，燒滅仇敵；凡想要害他們的都必須這樣被殺。 \end{tabularx} \\ \\ \relax
11:6 & \begin{tabularx}{0.7\textwidth}{X} 這二人有權柄關閉天空，使他們說預言的日子不下雨；又有權柄使水變為血，並且能隨時隨意用各樣的災害擊打大地。 \end{tabularx} \\ \\ \relax
11:7 & \begin{tabularx}{0.7\textwidth}{X} 他們作完見證的時候，那從無底坑裡上來的獸要跟他們交戰，並且得勝，把他們殺了。 \end{tabularx} \\ \\ \relax
11:8 & \begin{tabularx}{0.7\textwidth}{X} 他們的屍首將倒在大城的街道上；這城按著靈意叫所多瑪，又叫埃及，就是他們的主釘十字架的地方。 \end{tabularx} \\ \\ \relax
11:9 & \begin{tabularx}{0.7\textwidth}{X} 從各民族、支派、語言、邦國中有人觀看他們的屍首三天半，又不許人把屍首安放在墳墓裡。 \end{tabularx} \\ \\ \relax
11:10 & \begin{tabularx}{0.7\textwidth}{X} 住在地上的人會因他們而歡喜快樂，互相饋送禮物，因為這兩位先知曾使住在地上的人受痛苦。 \end{tabularx} \\ \\ \relax
11:11 & \begin{tabularx}{0.7\textwidth}{X} 過了這三天半，有生命的氣息從神那裡進入他們裡面，他們就站起來；看見他們的人都大大懼怕。 \end{tabularx} \\ \\ \relax
11:12 & \begin{tabularx}{0.7\textwidth}{X} 兩位先知聽見有大聲音從天上對他們說：「上這裡來。」他們就駕著雲上了天，他們的仇敵也看見了。 \end{tabularx} \\ \\ \relax
11:13 & \begin{tabularx}{0.7\textwidth}{X} 正在那時候，地大震動，城倒塌了十分之一；因地震而死的有七千人，其餘的都恐懼，歸榮耀給天上的神。 \end{tabularx} \\ \\ \relax
11:14 & \begin{tabularx}{0.7\textwidth}{X} 第二樣災禍過去了；看哪，第三樣災禍快到了。 \end{tabularx} \\ \\ \relax
11:15 & \begin{tabularx}{0.7\textwidth}{X} 第七位天使吹號，天上就有大聲音說：「世上的國已成了我們的主和他所立的基督的國了。他要作王直到永永遠遠！」 \end{tabularx} \\ \\ \relax
11:16 & \begin{tabularx}{0.7\textwidth}{X} 在神面前，坐在自己座位上的二十四位長老都俯伏在地上敬拜神， \end{tabularx} \\ \\ \relax
11:17 & \begin{tabularx}{0.7\textwidth}{X} 說：「今在昔在的主—全能的神啊，我們感謝你！因你執掌大權作王了。 \end{tabularx} \\ \\ \relax
11:18 & \begin{tabularx}{0.7\textwidth}{X} 外邦發怒，你的憤怒臨到了。審判死人的時候也到了；你的僕人眾先知、眾聖徒及敬畏你名的人，連大帶小得賞賜的時候到了；你毀滅那些毀滅大地者的時候也到了。」 \end{tabularx} \\ \\ \relax
11:19 & \begin{tabularx}{0.7\textwidth}{X} 於是，神天上的聖所開了，在他聖所中，他的約櫃出現了；隨後有閃電、響聲、雷轟、地震、大冰雹。 \end{tabularx} \\ \\
[1ex]
\hline
\hline
\end{longtable}
講題:重述福音　Retelling the Gospel

經文:啟示錄11章1-19節

講員:鄧瑞強博士

場合:香港中文大學崇基學院禮拜堂

日期:2018年5月13日(復活期第七主日)

各位弟兄姊妹,早安.

今日是母親節,祝願各位母親,身心康健,蒙神祝福,能教養兒女,明白上主的真理,一生走正直的道路.

近來,有個問題很困擾我,這個問題是:「我的母語是什麼?」是的,我的母語到底是什麼?有人說,這個問題太無聊,不用回答.本來,真是的,這問題實在太無聊,真的不用回答.但是,當有人說,我的母語不是我的真正母語時,這個問題就變得複雜了.

見到一隻動物,我說「鹿」,但有人說,這是「馬」.這明明是「鹿」,我從少到大,都叫這動物為「鹿」,怎會是「馬」呢?有人說,你以為你的母語稱這為「鹿」,但這只是你對母語的一種錯誤認識,這種動物在你的真正母語中,應稱為「馬」.這時,問題就複雜了.我的母語,到底是什麼?

母語,應是胎兒在母胎中已聽慣的語言,也是嬰孩來到這個世界第一種學習的語言.每種語言,都用一種獨特的方式去理解這個世界,都用一種獨特的方式去建構一個意義世界.這個不難理解,舉例說,德語的每個名詞,都分性別,分陽性,陰性,中性.你一用德語,世界上的每種事物,便呈現性別.中文不是這樣的.如此,用德語看的世界,和用中文看的世界,當然是不一樣的世界.在這意義上,一種語言,開出一個世界；多用一種語言,就多活一個世界.母語所描繪的世界,所建構的世界,就是每個人最基本的生活世界.

除了中文,英文這些語言之外,哲學家擴大了「語言」的範圍,按哲學家維根斯坦(Wittgenstein)的思路,生活裡的不同領域差不多都自成一獨立的語言系統.飲紅酒時,有人說,這酒很「乾身」(dry).若你是初哥,你便需要學習一新的語言系統,一杯液體,竟會是「乾」的.若你去打籃球,又是另一語言系統,在球場上,對方會「拆你屋」(alley-oop).基督教信仰,其實也自成一獨特的「語言系統」,有自己的用語,有自己的文法,有自己看到的世界.譬如說,主禱文第一句,「我們在天上的父」,這句話作為宗教語言表達的含意,與其中文文字表面所講的意思,有很大差異.「天上的父」的「天上」,不是我們肉眼見的「天上」,不是物理意義上的「天上」,它意指另一種世界.「天上的父」的「父」,不是男性,沒有太太,亦沒有鬍鬚,這與我們日常講的「父」不是同一個意義.「我們在天上的父」,這句話描繪出一個與平常世界不一樣的世界,在這信仰世界裡,我們看每個人都是同一個「父」之下互有關連的兄弟姊妹.你身邊的人,不是和你毫無關連的人.「我們在天上的父」,這句簡單的話,述說一個新世界,也引導我們活在這個奇異的世界裡.

由於「信仰本身」就是一種語言系統,涉及不同的文法與用語,也描繪著不一樣的世界,故此,縱使我們的信仰是用中文表達的,我們也要將之作一種「翻譯」.成功的翻譯,是將我們從這個世界,移到另一個世界裡.

今日的講道經文,充滿宗教象徵,用基督教特有的語言,揭露人間的真相.看看我能否成功翻譯.

今日的講道經文:啟示錄11:1－19

啟11:1　有一根葦子賜給我,當作量度的杖；且有話說:「起來!將神的殿和祭壇,並在殿中禮拜的人都量一量.

啟11:2　只是殿外的院子要留下不用量,因為這是給了外邦人的；他們要踐踏聖城四十二個月.

啟11:3　我要使我那兩個見證人,穿著毛衣,傳道一千二百六十天.」

啟11:4　他們就是那兩棵橄欖樹,兩個燈臺,立在世界之主面前的.

啟11:5　若有人想要害他們,就有火從他們口中出來,燒滅仇敵.凡想要害他們的都必這樣被殺.

啟11:6　這二人有權柄,在他們傳道的日子叫天閉塞不下雨；又有權柄叫水變為血,並且能隨時隨意用各樣的災殃攻擊世界.

啟11:7　他們作完見證的時候,那從無底坑裡上來的獸必與他們交戰,並且得勝,把他們殺了.

啟11:8　他們的屍首就倒在大城裡的街上；這城按著靈意叫所多瑪,又叫埃及,就是他們的主釘十字架之處.

啟11:9　從各民,各族,各方,各國中,有人觀看他們的屍首三天半,又不許把屍首放在墳墓裡.

啟11:10　住在地上的人就為他們歡喜快樂,互相餽送禮物,因這兩位先知曾叫住在地上的人受痛苦.

啟11:11　過了這三天半,有生氣從神那裡進入他們裡面,他們就站起來；看見他們的人甚是害怕.

啟11:12　兩位先知聽見有大聲音從天上來,對他們說:「上到這裡來.」他們就駕著雲上了天,他們的仇敵也看見了.

啟11:13　正在那時候,地大震動,城就倒塌了十分之一,因地震而死的有七千人；其餘的都恐懼,歸榮耀給天上的神.

啟11:14　第二樣災禍過去,第三樣災禍快到了.

啟11:15　第七位天使吹號,天上就有大聲音說:世上的國成了我主和主基督的國；他要作王,直到永永遠遠.

啟11:16　在神面前,坐在自己位上的二十四位長老,就面伏於地,敬拜神,

啟11:17　說:昔在,今在的主神─全能者啊,我們感謝你!因你執掌大權作王了.

啟11:18　外邦發怒,你的忿怒也臨到了；審判死人的時候也到了.你的僕人眾先知和眾聖徒,凡敬畏你名的人,連大帶小得賞賜的時候也到了.你敗壞那些敗壞世界之人的時候也就到了.

啟11:19　當時,神天上的殿開了,在他殿中現出他的約櫃.隨後有閃電,聲音,雷轟,地震,大雹.

(1)這段經文的位置

我們記得,神永恆的心意,記載在一「書卷」裡,這書卷被七個封印封著.耶穌逐個印打開,揭開第七印時,天上寂靜了半個鐘.書卷打開了,其內容是七枝號.

第一號至第六號,明顯是神對罪惡世界的最後審判.第六號的結尾,再出現那書卷,這是啟示錄的第十章的內容.這書卷的內容,就是神對世界的終極心意.今日的講道經文,就是這書卷的核心內容,也就是神的終極心意.

(2)信徒的特質一:被神保護

啟11:1　有一根葦子賜給我,當作量度的杖；且有話說:「起來!將神的殿和祭壇,並在殿中禮拜的人都量一量.

啟11:2　只是殿外的院子要留下不用量,因為這是給了外邦人的；他們要踐踏聖城四十二個月.

「量度的杖」,可以是王者的「權杖」,也可以是牧羊人的「牧杖」.兩者合璧,即王者像牧人,他的杖,他的竿,都安慰我們.這裡的「殿」與「壇」,都指向敬拜神的人.敬拜神的人,在殿中敬拜,以自己的生命為活祭獻在壇上.神要以王者的牧杖,保護這群信徒.

至於殿外的人,即不敬拜神的人,神任讓他們在自己的罪惡中滅亡.他們要踐踏真理與正義四十二個月.四十二個月,即三年半.「七」是圓滿的數字,「七年」是神設定的圓滿期限,「三年半」即事情在進展中,未到終點.「三年半」的意義,有點似「易經」的最後一卦「未濟卦」.罪惡在人間,仍是猖獗；神的正義未達,卻遙遙可望.

(3)信徒的特質二:被召作真理的見證人

啟11:3　我要使我那兩個見證人,穿著毛衣,傳道一千二百六十天.」

啟11:4　他們就是那兩棵橄欖樹,兩個燈臺,立在世界之主面前的.

啟11:5　若有人想要害他們,就有火從他們口中出來,燒滅仇敵.凡想要害他們的都必這樣被殺.

啟11:6　這二人有權柄,在他們傳道的日子叫天閉塞不下雨；又有權柄叫水變為血,並且能隨時隨意用各樣的災殃攻擊世界.

「兩個見證人」,代表整體信徒.猶太人傳統,要有兩個見證人,其見證方為有效.「一千二百六十天」,以1年為360日計,即「三年半」.在天國未圓滿實現的這段時間裡,信徒要做的,教會要做的,就是為真理作證.

「兩棵橄欖樹」,這是典故,出自舊約撒迦利亞書4:3,與「燈台」及「聖靈」有關.「燈台」在啟示錄裡,就是指教會.教會要倚靠「聖靈」之力,在歷史中,見證神的真理.

叫火從天上降下來,又叫天閉塞不下雨,這是舊約以利亞先知的事蹟.叫水變為血,是摩西的事蹟.以利亞和摩西,是耶穌登山變像時,在山上顯現的兩個神聖人物.他們和耶穌,相討救世的事,相討一個新的出埃及故事.每個信徒,在人間,在歷史裡,都要像以利亞和摩西,宣講神的救世事蹟.

(4)信徒的特質三:為見證真理而受逼害

啟11:7　他們作完見證的時候,那從無底坑裡上來的獸必與他們交戰,並且得勝,把他們殺了.

啟11:8　他們的屍首就倒在大城裡的街上；這城按著靈意叫所多瑪,又叫埃及,就是他們的主釘十字架之處.

啟11:9　從各民,各族,各方,各國中,有人觀看他們的屍首三天半,又不許把屍首放在墳墓裡.

啟11:10　住在地上的人就為他們歡喜快樂,互相餽送禮物,因這兩位先知曾叫住在地上的人受痛苦.

信徒面對的,是深淵一般的邪惡力量,這力量像怪獸一樣,要撕裂信徒.信徒受死的地方,是「大城」,即當時的羅馬,這是集政治,經濟,軍事強權於一身的巨大怪獸.人間的大城,按靈意來說,是罪惡滔天的所多瑪,是奴役人的埃及.主耶穌因這罪惡而死,也為釋放被奴役的人而死,故此經文說,這「就是他們的主釘十字架之處」.信徒面對的,就是反抗神的最極端邪惡.信徒被殺害,暴屍「三天半」,當然,這又是指等待天國實現所需的時間,邪惡當道,神的正義在來臨.邪惡的世界,會為信徒的死而高興.黑暗之為黑暗,就是完全容忍不到半點光明.

你會問:為何神不早早介入?為何神容讓邪惡在人間肆虐「三年半」?

答案是:這是神行事的方式,這是基督教這種「語言系統」獨有的文法,就是讓信徒在艱難的歲月裡能明白自己的「天命」,那也是主耶穌領受的「天命」,以真理的言說者的身分,去對抗人間的謊話與邪惡.這就是羔羊耶穌得勝罪惡的方式,以犧牲的受苦去勝過殘暴.這是人打你右臉,連左臉也給他打的一種文法.

十字架,受苦,受死,就是結局嗎?

不是的,「三年半」完結後,就是終極的天地.

(5)信徒的特質四:信徒最終會被神肯定

啟11:11　過了這三天半,有生氣從神那裡進入他們裡面,他們就站起來；看見他們的人甚是害怕.

啟11:12　兩位先知聽見有大聲音從天上來,對他們說:「上到這裡來.」他們就駕著雲上了天,他們的仇敵也看見了.

耶穌死後三天復活,為真理受苦受死的信徒,死後「三天半」復活.不是「三天」而是「三天半」,很明顯,是就著啟示錄的時間象徵來說的.

時候到了,神的生氣吹來了,他們像基督一樣,復活升天,去到永恆的神那裡.

(6)信徒的特質五:信徒生命的「基督化」

我們發現,信徒面對的,經歷的,正是基督面對的,經歷的.在啟示錄的描繪下,每個信徒像是「重活基督」的生命一次.活在人間,我們就是基督的代表.

我們像基督一樣,宣講真理,宣講天國的即將來臨.我們像基督一樣,要與邪惡的力量對峙,要為真理受苦.我們像基督一樣,在神面前,在時間的終點,被神肯定,被神賜予永恆的生命.

(7)信徒事奉的結果:人的悔改

啟11:13　正在那時候,地大震動,城就倒塌了十分之一,因地震而死的有七千人；其餘的都恐懼,歸榮耀給天上的神.

留意信徒努力事奉後帶來的效果.

阿摩司書5:3講到,在神的審判裡,是消滅十分之九,只留下十分之一.在這裡,因為信徒的努力,現在多人悔改,神只消滅十分之一,留下十分之九.

列王記上19:18講到,在那邪惡世界,神的忠信子民,只有七千人,其餘都是敗壞的.在這裡,因為信徒的努力,現在多人悔改,神消滅的不信者,只有七千人,其餘都是悔改的人.

(8)神的心意的揭露

經文來到這裡,天上的書卷打開了,神的心意也揭露了.

神不願意全然毀滅罪人,他的心意是呼召信徒,呼召他們像基督一樣生活,堅強地言說真理,剛毅地面對苦難,忍耐地等待神的正義.

信徒這樣生活,不是無意義的.信徒這樣生活,能使罪人有悔改的可能.

神這樣的計劃,就是「福音」.啟示錄用一種電影畫面般的語言,為我們「重述福音」.就因為神有這樣的計劃,世上有這樣的一批信徒,世界不至黑暗無光.

你和我,像微小的螢火蟲,在漆黑的大地,斷斷續續的,發出柔弱的光.就因這一點點的光,使黑暗中的人,感到絲絲溫暖,抱著希望.

(9)神應許的終局

啟11:15　第七位天使吹號,天上就有大聲音說:世上的國成了我主和主基督的國；他要作王,直到永永遠遠.

啟11:16　在神面前,坐在自己位上的二十四位長老,就面伏於地,敬拜神,

啟11:17　說:昔在,今在的主神─全能者啊,我們感謝你!因你執掌大權作王了.

啟11:18　外邦發怒,你的忿怒也臨到了；審判死人的時候也到了.你的僕人眾先知和眾聖徒,凡敬畏你名的人,連大帶小得賞賜的時候也到了.你敗壞那些敗壞世界之人的時候也就到了.

啟11:19　當時,神天上的殿開了,在他殿中現出他的約櫃.隨後有閃電,聲音,雷轟,地震,大雹.

這是最榮耀的時刻.神要作王了,天國降臨了,邪惡被消滅了.

(10)結語

基督教信仰,是一種很特殊的語言系統,有自己的語彙,有自己的文法,有自己要描述的一個奇特的世界.我希望,這種語言,會慢慢變成我們的真正母語.

但願:榮耀歸於聖父,聖子,聖靈.阿們.



\newpage

\section{保羅身旁的三個夥伴}
\label{sec:411}
\textbf{戴國璋牧師 - 保羅身旁的三個夥伴}
\newline
\newline
\hyperref[sec:410]{< < < PREV SERMON < < <}
~
\hyperref[sec:toc]{[返目錄]}
~
\hyperref[sec:412]{> > > NEXT SERMON > > >}


提摩太後書 4:10-14
\newline
\begin{longtable}{cl}
\hline
\hline
章節 & 經文 (和合本修訂版)\\
\hline
4:10 & \begin{tabularx}{0.7\textwidth}{X} 因為底馬貪愛現今的世界，已經離棄我，往帖撒羅尼迦去了；革勒士往加拉太去；提多往撻馬太去； \end{tabularx} \\ \\ \relax
4:11 & \begin{tabularx}{0.7\textwidth}{X} 只有路加在我這裡。你來的時候把馬可帶來，因為他在服事上於我有益。 \end{tabularx} \\ \\ \relax
4:12 & \begin{tabularx}{0.7\textwidth}{X} 我已經打發推基古往以弗所去。 \end{tabularx} \\ \\ \relax
4:13 & \begin{tabularx}{0.7\textwidth}{X} 我在特羅亞留給加布的那件外衣，你來的時候要帶來，那些書也帶來，特別是那幾卷羊皮的書。 \end{tabularx} \\ \\ \relax
4:14 & \begin{tabularx}{0.7\textwidth}{X} 銅匠亞歷山大多方害我；主必照他所行的報應他。 \end{tabularx} \\ \\
[1ex]
\hline
\hline
\end{longtable}
講題:保羅身旁的三個夥伴　The three fellow workers of Paul

經文:歌羅西書4:10-14;提摩太後書4:9-11

講員:戴國璋牧師

場合:香港中文大學崇基學院禮拜堂

日期:2018年5月27日(聖靈降臨後第一主日-三一主日)

歌羅西書及提摩太後有一些共通點；它們都寫於羅馬,且都寫在保羅遭囚禁的境況中.但兩者相隔了約七年.前者是保羅因上訴該撒,第一次到羅馬城,在自己租用房子遭軟禁的兩年間寫就,時約 AD57.後者是在上訴該撒事得直後獲釋,再經一輪的傳道生活後重返羅馬,適值該撒尼錄對基督教進行迫害,保羅作為當時基督教的首領人物而遭殺害.後書就寫於遭殺害殉教之前的囚禁中間,時約 AD64.兩者比較下,保羅的心境就有著很大的變化.我們所讀的是兩書問候話,當中提及一些共同的人物.今天我們思考這三個人物對保羅的重要.

1. 忠心的陪伴 - 路加醫生

路加是保羅一生傳道工作上的幫手,是保羅生命中忠心的朋友,這是眾所周知的.有關路加的出生,信主的經過,甚至如何成為保羅傳道的拍檔,我們都缺乏資料.學者們只能從有限的了解中嘗試重構這個人物.路加是個學識淵博的人物,他的作品 (福音及行傳)文筆優美,結構嚴謹,考據詳實.在在說明他的學養深厚,路加是個典型的拉丁名字,意為「光；帶來光的人；帶來知識的人」.自身相信是個希獵人身相信是個希獵人.有學者猜想路有學者猜想路加曾經是一個奴隸,後來獲得自由身.在當時,有許多學識淵博的者,擁有專業技能,- 如醫生 -的希臘人,因戰爭而成為奴隸,又因主人的賞識而獲自由的.路加可能是其中之一.他在認識保羅之前,或許已經是基督徒.他是在特羅亞(亞西亞的一個海港)遇上保羅和他的團隊.當時約為AD48.此後,無論保羅在馬其頓,以弗所或返回耶路撒冷,遭囚於該撒利亞,到達羅馬,開始他西班牙的工作,再回到羅馬並遭難,作為醫生的路家都陪伴保羅旁邊.

在寫後書的時間,不同的夥伴因不同的原由離他而去.「獨有 」路加在保羅的身旁,成為他重要的支持.一個這樣忠心的夥伴,一生難求.

2. 於人有益的同工 - 馬可傳道

馬可跟保羅認識比路加更長久.馬可是一個拉丁名字,意思是 「鬥士；尚武之人 」.這在羅馬人中是個很受歡迎的名字 (羅馬有位英雄叫 Mark Anton,生活在耶穌 稍前 ) .但我們知道這位「馬可」,原叫「約翰 」是一位猶太人,是巴拿巴的表弟,或許也是來自居比路(塞浦路斯).這位馬可的約翰曾經一度使保羅心煩,在保羅與巴拿巴拍檔的第一次傳道中,中途離隊.此事觸發保羅與巴拿巴的矛盾.在第二次傳道開始 ,兩人分道揚鑣.從馬可乃巴拿巴表弟一事看,或許他比較年輕.年輕人總叫年長的人擔心行事輕浮而不踏實.這也許是保羅這個力求認真的人容不下馬可的原因.若這是聖經中對馬可事蹟唯一的記述,這或許是他的一生蓋棺定論.但感謝神,不是.在保羅後來的傳道歷程中,馬可成為一個可靠的同工.在他最末後的描述,他甚至稱馬可為「於我有益」.馬可不是一個「妨礙」,是「益處」.雖則我們不知是怎樣的益處.

馬可果然是叫人有益處的一個同工,不但對保羅而言,對偉大使徒彼得亦然.他稱馬可為「我兒子」.這叫我們想起保羅與他心愛的門生提摩太.早期教會流傳著這樣一個說法:「馬可是彼得的傳道幫手,記述彼得傳道時講述的耶穌事蹟,並輯錄成福音書」.若沒有馬可,我們也就沒有馬可福音.

3. 無法有始有終的信徒 - 底馬先生

底馬(Demas)是個希臘名字,原意是受「歡迎的人」,是保羅傳道上的一位同工.保羅稱他為所親愛的,和路加一樣.卻在最危難的關頭,離棄同伴往繁華璀璨的帖撒羅尼迦.一個一直都走得這麼好的信徒,怎麼會離棄信仰成為叫主傷心的人呢 - 值得深思.



\newpage

\section{事奉耶和華}
\label{sec:412}
\textbf{吳碧珊牧師 - 事奉耶和華}
\newline
\newline
\hyperref[sec:411]{< < < PREV SERMON < < <}
~
\hyperref[sec:toc]{[返目錄]}
~
\hyperref[sec:413]{> > > NEXT SERMON > > >}


撒母耳記上 3:1-10
\newline
\begin{longtable}{cl}
\hline
\hline
章節 & 經文 (和合本修訂版)\\
\hline
3:1 & \begin{tabularx}{0.7\textwidth}{X} 那孩子撒母耳在以利面前事奉耶和華。在那些日子，耶和華的言語稀少，不常有異象。 \end{tabularx} \\ \\ \relax
3:2 & \begin{tabularx}{0.7\textwidth}{X} 那時，以利在自己的地方睡覺；他眼目開始昏花，不能看見。 \end{tabularx} \\ \\ \relax
3:3 & \begin{tabularx}{0.7\textwidth}{X} 神的燈還沒有熄滅，撒母耳睡在耶和華的殿內，神的約櫃就在那裡。 \end{tabularx} \\ \\ \relax
3:4 & \begin{tabularx}{0.7\textwidth}{X} 耶和華呼喚撒母耳，撒母耳說：「我在這裡！」 \end{tabularx} \\ \\ \relax
3:5 & \begin{tabularx}{0.7\textwidth}{X} 他跑到以利那裡，說：「你叫我嗎？我在這裡。」以利說：「我沒有叫你，回去睡吧。」他就回去睡了。 \end{tabularx} \\ \\ \relax
3:6 & \begin{tabularx}{0.7\textwidth}{X} 耶和華又呼喚撒母耳。撒母耳起來，到以利那裡，說：「你叫我嗎？我在這裡。」以利說：「我兒，我沒有叫你，回去睡吧。」 \end{tabularx} \\ \\ \relax
3:7 & \begin{tabularx}{0.7\textwidth}{X} 那時撒母耳還未認識耶和華，耶和華的話也未曾向他啟示。 \end{tabularx} \\ \\ \relax
3:8 & \begin{tabularx}{0.7\textwidth}{X} 耶和華第三次再呼喚撒母耳。撒母耳起來，到以利那裡，說：「你叫我嗎？我在這裡。」以利才明白是耶和華呼喚這小孩。 \end{tabularx} \\ \\ \relax
3:9 & \begin{tabularx}{0.7\textwidth}{X} 以利對撒母耳說：「你回去睡吧。他若再叫你，你就說：『耶和華啊，請說，僕人敬聽！』」撒母耳就回去，仍睡在原處。 \end{tabularx} \\ \\ \relax
3:10 & \begin{tabularx}{0.7\textwidth}{X} 耶和華來站著，像前幾次呼喚：「撒母耳！撒母耳！」撒母耳說：「請說，僕人敬聽！」 \end{tabularx} \\ \\
[1ex]
\hline
\hline
\end{longtable}
講題:事奉耶和華　Serve the Lord

講員:吳碧珊牧師

經文:撒母耳記上3:1-10

場合:香港中文大學崇基學院禮拜堂

日期:2018年6月3日

撒上3:1-10 耶和華向撒母耳顯現

3:1那孩子撒母耳在以利面前事奉耶和華.在那些日子,耶和華的言語稀少,不常有異象.

3:2那時,以利在自己的地方睡覺；他眼目開始昏花,不能看見.

3:3上帝的燈還沒有熄滅,撒母耳睡在耶和華的殿內,上帝的約櫃就在那裡.

3:4耶和華呼喚撒母耳,撒母耳說:「我在這裡!」

3:5他跑到以利那裡,說:「你叫我嗎?我在這裡.」以利說:「我沒有叫你,回去睡吧.」他就回去睡了.

3:6耶和華又呼喚撒母耳.撒母耳起來,到以利那裡,說:「你叫我嗎?我在這裡.」以利說:「我兒,我沒有叫你,回去睡吧.」

3:7那時撒母耳還未認識耶和華,耶和華的話也未曾向他啟示.

3:8耶和華第三次再呼喚撒母耳.撒母耳起來,到以利那裡,說:「你叫我嗎?我在這裡.」以利才明白是耶和華呼喚這小孩.

3:9以利對撒母耳說:「你回去睡吧.他若再叫你,你就說:『耶和華啊,請說,僕人敬聽!』」撒母耳就回去,仍睡在原處.

3:10耶和華來站著,像前幾次呼喚:「撒母耳!撒母耳!」撒母耳說:「請說,僕人敬聽!」

詩139:1-6, 13-18 耶和華的全知和眷顧

139:1耶和華啊,你已經鑒察我,認識我.

139:2我坐下,我起來,你都曉得；你從遠處知道我的意念.

139:3我行路,我躺臥,你都細察；你也深知我一切所行的.

139:4耶和華啊,我舌頭上的話,你沒有一句不知道的.

139:5你前後環繞我,按手在我身上.

139:6這樣的知識奇妙,是我不能測的；至高,是我不能及的.

139:13我的肺腑是你所造的,我在母腹中,你已編織我.

139:14我要稱謝你,因我受造奇妙可畏,你的作為奇妙,這是我心深知道的.

139:15我在暗中受造,在地的深處被塑造；那時,我的形體並不向你隱藏.

139:16我未成形的體質,你的眼早已看見了；你所定的日子,我尚未度一日,都在你的冊子寫上了.

139:17上帝啊,你的意念向我何等寶貴!其數何等眾多!

139:18我若數點,比海沙更多；我睡醒的時候,仍和你同在.

林後4:5-12;

4:5我們不是傳自己,而是傳耶穌基督為主,並且自己因耶穌作你們的僕人.

4:6那吩咐光從黑暗裡照出來的上帝已經照在我們心裡,使我們知道上帝榮耀的光顯在耶穌基督的臉上.

4:7我們有這寶貝放在瓦器裡,為要顯明這莫大的能力是出於上帝,不是出於我們.

4:8我們處處受困,卻不被捆住；內心困擾,卻沒有絕望；

4:9遭受迫害,卻不被撇棄；擊倒在地,卻不致滅亡.

4:10我們身上常帶著耶穌的死,使耶穌的生也在我們身上顯明.

4:11因為我們這活著的人常為耶穌被置於死地,使耶穌的生命在我們這必死的人身上顯明出來.

4:12這樣看來,死是在我們身上運作,生卻在你們身上運作.

可2:23-3:6安息日摘麥穗,治好萎縮了的手

2:23有一個安息日,耶穌從麥田經過.他的門徒走路的時候,摘起麥穗來.

2:24法利賽人對耶穌說:「看哪!他們為甚麼做安息日不合法的事呢?」

2:25耶穌對他們說:「大衛和跟從他的人飢餓需要食物時所做的事,你們沒有念過嗎?

2:26他在亞比亞他作大祭司的時候,怎麼進了上帝的居所,吃了供餅,又給跟從他的人吃呢?這餅除了祭司以外,人都不可以吃.」

2:27他又對他們說:「安息日是為人設立的,人不是為安息日設立的.

2:28所以,人子也是安息日的主.」

3:1耶穌又進了會堂,在那裡有一個人,他的一隻手萎縮了.

3:2眾人為了要控告耶穌,就窺探他會不會在安息日醫治那人.

3:3耶穌對那手萎縮了的人說:「起來站在當中!」

3:4他又問眾人:「在安息日行善行惡,救命害命,哪樣是合法的呢?」他們都不作聲.

3:5耶穌怒目環視他們,因他們的心剛硬而憂傷,就對那人說:「伸出手來!」他把手一伸,手就復原了.

3:6法利賽人出去,立刻同希律一黨的人商議怎樣除掉耶穌.

引言:

今天香港教會太安舒,特別昔日在殖民地政府管治下,確實給予教會很多傳福音的機會與空間,我們容易習以為常,認為是必然的.但其實今天仍有不少地方教會的發展是受到一定情度的限制,如中國,印度等等.

我們香港教會在不知不覺間被寵壞了,被世俗化,我們被市場經濟效益思想主導著.以事工成效衡量教會的發展,慢慢失去了那本來的召命,教會的存在乃是為事奉耶和華.

查理狄更斯《雙城記》開場的第一段.「這是最好的時代,這是最壞的時代；這是智慧的年代,這是愚蠢的年代；這是信仰的時代,這是懷疑的時代；這是光明的時代,這是黑暗的時代…」

活在這時代的你及我,如何描述這所處的境況!如何理解此刻我們事奉耶和華的意思?今天想與大家一起分享,互相勉勵,再思想:事奉耶和華的意思與這身份帶來的挑戰!於舊約聖經中,事奉耶和華這句子出現了40次.

撒上3:1-10

撒母耳自一斷奶時,已被母親哈拿交予先知以利,好使撒母耳能終身於聖殿服事耶和華.

現刻以利年紀已經很老了,眼目昏花,於晚間耶和華呼喚撒母耳；撒母耳以為是以利找他,直至第三次,以利才明白,就提醒撒母耳,是耶告華上帝親自呼喚他.若再聽到聲音,要回答說:「耶和華啊,請說,僕人敬聽!」

於是撒母耳就回去自己原先睡覺的地方.忽然像前三次呼喚說:「撒母耳啊!撒母耳啊!」撒母耳即時回答說:「請說,僕人敬聽!」」

經文3:7「那時撒母耳還未認識耶和華,也未得耶和華的默示.」

未認識是指撒母耳尚無直接與耶和華上帝接觸的經驗,未有與耶和華上帝直接對話.耶和華也未曾向他有任何直接的指示,啟示.

以利教導撒母耳再聽到聲音,回應的說:「僕人敬聽」.意思是:要知道自己乃是奴僕卑微的身份,願意靜聽,隨時準備遵依吩咐,及決意順服於這位主的吩咐.

在經文中我們發現；

是上帝親自去,主動的呼召,揀選撒母耳,去事奉祂.

縱然撒母耳當時對上帝的接觸,認識有限,上帝揀選人去事奉祂,不在他的才能,資歷,位份

人未能明白為何上帝偏偏選中他,不是人預計

正如詩139:1-6, 13-18

所描述；這位耶和華在詩人仍於母腹中,在未成形體的時候,他已被耶和華寫上在冊子裡.是人不能測的；這至高者,是人不能及的.無論人坐下,起來,行路或躺臥,祂都曉得,都細察.

詩人一再強調；告訴我們,這位我們所相信的耶和華是滿有權柄和能力,祂是知悉我們每一個人的特性及限制,祂是主動邀請人去回應祂的揀選.

祂是等待人的回應；撒母耳在聆聽到這細微聲音!即時的答:「請說,僕人敬聽!」

上帝在邀請人去事奉祂的時候,也會給予一個空間,一個選擇.我們所見撒母耳也樂意,謙卑的去回應,順服這被使用的機會.

林後4:5-12;

聖經林後4:5-12;同時也在提醒我們,回應上帝的揀選去事奉祂的時候,自己是要清楚這角色.

事奉者的目標是指向耶穌基督,是為福音的見證者,是參與上帝拯救的工作,也是和上帝同工,不是標榜自己的才智能力,而是要謙卑地順服.保羅提到他和與他同工的伙伴不過是福音的「瓦器」.

「瓦器」,這在較早的社會是非常普遍的家庭用品,它是用泥土曬乾後送入窯裡燒硬,蓋在房子上面以防下雨或日曬.非常普遍的用品.

保羅用「瓦器」來形容事奉者的角色,乃是要表達人是不配,但卻是上帝的憐憫和恩典,揀選了他,使他有這樣的福份.因此,「瓦器」裡面貴重的「物品」並不屬於他,而是上帝的.

雖然「瓦器」是那樣脆弱的物品,但卻沒有被「壓碎」；是；在事奉的路上遇到許多困難,保羅描述「我們處處受困,卻不被捆住；內心困擾,卻沒有絕望；遭受迫害,卻不被撇棄；擊倒在地,卻不致滅亡.」但那事奉的心志是沒有被擊倒.

保羅提醒事奉者,作基督的見證人；並不是很「平安」的一件事,是會有「受困」,「內心困擾」,「甚而迫害」,「被擊倒在地」.在事奉者的身上常帶著耶穌的死,以致能將耶穌的生也在他們身上顯明.

4:11保羅描述

「和合本修訂版」

因為我們這活著的人常為耶穌被置於死地,使耶穌的生命在我們這必死的人身上顯明出來.

「思高聖經版本」

因為我們這些人雖然活著﹐卻永遠為了耶穌的緣故被送交於死﹐好叫耶穌的活也在我們這能死的肉身上顯明出來.

「當代中文譯本」

我們這些活著的人,已經預備好了隨時為耶穌的緣故,不惜被人置諸死地；這樣好使耶穌的生命,在我們這終有一死的肉身上表現出來.

「現代中文譯本聖經修訂版」

我們的一生,常常為了耶穌的緣故冒著死亡的危險,為要使他的生命能夠在我們這必朽的身上顯明出來.

保羅要那些願意跟從,那事奉主基督的人,要準備,要知道,為了成全,成就耶穌基督的吩咐,這信仰,這身份是會危險,是要冒險的,是會被人置諸死地.

耶穌也是如此的以順服的生命去回應耶和華上帝對祂的揀選作為事奉耶和華的人,主耶穌選擇了這甘願冒險的路,堅定冒著死亡的來回應這的召命.

可2:23-3:6

耶穌門徒在安息日摘麥穗及耶穌治好萎縮了的手事蹟中,經文記載著；耶穌基督與門徒在安息日經過麥田,門徒肚子「餓了」,就伸手去摘麥穗吃.問題並不是出在摘麥穗,因為若是肚子餓了,經過別人的麥田去摘麥穗,這是摩西法律所允許的行為.

問題是在摘麥穗下來時的動作,是「用手搓著吃」,這才是關鍵問題.因為是「搓」,需要雙手用力,這就違反了安息日規定,也是這些法利賽人所認定的「做工」之動作.就是違反安息日的誡命.耶穌的門徒這樣做,就是大逆不道.他們為此特地來找耶穌,向他提出抗議.

耶穌提起昔日大衛被掃羅追殺的時候,他跑到挪伯城祭司亞希米勒那裡；有關事蹟在撒上21-22章也有記述.那時大衛很肚餓,正在逃亡.他就欺騙亞希米勒,說他是因公出差,肚子非常餓,希望能吃點東西.亞希米勒因為沒有準備其它食物,就同意大衛拿祭壇上的祭品來吃.

依照摩西法律,這種食物只有祭司才可以吃.就因為這樣,這件事被掃羅知道了,因恨惡大衛,這筆帳算在亞希米勒的身上,掃羅藉此大屠殺該城所有的祭司,而亞希米勒的家族全都遭到屠殺.

現在耶穌以這件殘酷的大屠殺事件,來反問這些宗教領袖法利賽人.摩西法律,設有安息日,乃是為了幫助人專心去敬拜上帝,而不是用了成為借口去阻塞人,甚而成為藉口,計謀去傷害人..

耶穌提出了一個中心問題:安息日,律法,規條,以往運作模式,不能,不可以成為箝制人生命的金科玉律.

可3:1-6

也是繼續發生於安息日；耶穌是尊重律法的人,在安息日和眾人一起去猶太人會堂參加聚會.這刻祂入到會堂.祂看到一個手枯萎了的人；.

有一本新發現的「希伯來福音書」,裡面就記載著:

「有一位泥水匠,他靠著手工供作養活家庭.因為他的右手枯萎了,知道耶穌在會堂,就去懇求耶穌醫治他的手,以免他淪為乞丐.」

這些宗教領袖們向對耶穌基督的行徑相當不滿,處心積慮地想要找機會對耶穌基督下手.他們再想要用摩西的法律來陷害耶穌基督.

耶穌請這個手枯萎的病人站在他們面前,然後就提出了這個非常基本的生命問題:摩西的法律到底是要救人,還是要害人?如果法律是為了要害人,那樣的法律對嗎?若是要救人,那麼即使是安息日,難道就不可以救人嗎?會堂裡所有的人都「一聲不響」!

這些宗教領袖們沒有想到耶穌基督竟然會提這樣的反問,使他們無言以對.

第5節作者在這裡用「耶穌怒目環視左右」這句話,來形容耶穌基督心中的忿怒.原因是看見他們心腸剛硬.

耶穌的憤怒,是祂看到這個病人的需要,他的枯乾的手,是他的一家人生活已經受到影響,為何卻還被人利用來試探耶穌基督.

我們可以看出耶穌對人的真實的需要與價值,與法利賽人所詮釋的法律有很大的出入.祂甘願承受這行動帶來的一切危險.

若是耶穌真實被這些人去定罪,他們可以馬上就將耶穌基督逮捕移送猶太人最高法院去審問,而違背安息日的條例是可以判處死刑的.現在這些人不敢,只是自己也知道所提的理據很薄弱.

當耶穌基督用這問題來問這些宗教領袖之後,可以想像得到這些法利賽人反應.定會非常憤怒,對耶穌的仇恨就越深,更會找緊其他機會,將祂除去.

耶穌是冒著生命的危險,去忠於自己被選召作事奉耶和華者的身份.

總結:

今天幾段經文對我們有何提醒呢!

我們每一個人都在那創天造地萬軍耶和華手中,祂是那至高者,我們未出娘胎,祂已經鑒察我們,認識我們.人的知識與能力是未可明暸的.

是祂主動去愛人,去揀選人事奉祂,而不是人有努力,或自己可以做任何事,得到這被揀選作祂的事奉者這身份.

「事奉」這個詞,原文就是成為奴隸的意思,事奉耶和華有如成為耶和華上帝的奴隸,任憑衪用.耶和華要求衪的子民,事奉祂的人要忠心,誠心誠意地事奉衪.不能帶有任何其他動機去事奉.

我們都如粗陋的「瓦器」,但卻是上帝的憐憫和恩典,以這「瓦器」去盛載裡面貴重的「物品」,那就是上帝.

耶和華是我們的主人,是我們事奉的對象.事奉耶和華的重點,不是為祂做甚麼；事奉者就是要全心全意的以自己整個人的生命,每天的生活去活出耶和華上帝的心意.

事奉必須倚靠上帝得以成就,重點不是事奉者的技巧,事奉者的能力,而是上帝的行動.

再者,經文提醒我們,事奉見證耶和華並不是很「平安」的一件事,在事奉者的身上是要帶著耶穌的死,因是要抗衡罪的權勢,會因為此緣故被人置諸死地；是會冒著被攻擊,死亡的危險,以便能將耶穌的生也在他們身上顯明.

人可以作出回應,選擇是否願意,去事奉耶和華的.這抉擇乃是一個生命方向的選擇,要知道在抗衡世俗的主流意識與價值時,是要抗衡那些欺壓人的強權,古舊被人簡化的傳統或處事模式,是會有危險及艱辛的.但這條路卻又是會豐盈,充滿祝福.

這選擇是要求生命的改變,不是行為上的改變,是整個態度和方向,是質素,是心的改變.整個人生都是為要耶和華這主人的心意得到彰顯,得到榮耀.是要將耶穌的生命藉此在自己身上顯明.是以一個「瓦器」去盛載「至寶」.

這個選擇是挑戰人；是否願意放下自己的喜好；以及自身努力而得來的安全感,甚而願意為放下自我的安危,去為軟弱,受壓迫,剛硬的人重新尋見上帝.

今天我們香港在這高舉金錢,成效利益為取向的社會,就如大家都在談論,大灣區的經濟發展,香港的獨特經濟優勢會被減弱.我們現正積極鼓勵開創香港的創新科技,否則難以保持香港的競爭能力.

我們看見,今天年青人少談自己人生的理想；反之是問；大學那個學位和高級文憑好,邊樣搵工易D?邊樣吃香D?

又見那些消費主義,享樂文化充斥著各階層；如各方媒體都鼓勵借貸,借錢去完成他們的夢想,創立事業這個夢想,不用再受老板氣,又或可實現一個華禮夢想婚禮等等.

作為一個被主揀選的人,以著耶和華上帝的心意而去活的人；是否能分辨這些意識形態,文化價值取向而say no呢!為理想；不計薪酬福利的微弱,而選擇返工時間長,福利少的自願團體工作,能有勇氣以仍用小米手機而自豪,這個暑假不去歐洲旅行改去長洲,而感到滿足呢!

另方面在這回歸二十多年以來整個社會的深層次矛盾,充滿危機,不同意見,對立與撕裂.「繁榮穩定」,「馬照跑,舞照跳」,「高度自治」,「港人治港」論據下；作為事奉耶和華的人,如何就社會不公義,不平等而奮鬥,同時珍惜及尊重不同意見的表達!有勇氣去冒險,輕看被標籤,被指控,被限制的危險,而堅持活出自己的信念呢!

一個事奉耶和華的人,是祂主動邀請,也要自己作出抉擇,是冒險的抉擇,也是如此才能明白及活出這位永恆的主的吩咐,成為祂的見證人.



\newpage

\section{屋頂上的提琴手}
\label{sec:413}
\textbf{鄧瑞強博士 - 屋頂上的提琴手}
\newline
\newline
\hyperref[sec:412]{< < < PREV SERMON < < <}
~
\hyperref[sec:toc]{[返目錄]}
~
\hyperref[sec:414]{> > > NEXT SERMON > > >}


啟示錄 12:1-17
\newline
\begin{longtable}{cl}
\hline
\hline
章節 & 經文 (和合本修訂版)\\
\hline
12:1 & \begin{tabularx}{0.7\textwidth}{X} 天上出現了一個大兆頭：有一個婦人身披太陽，腳踏月亮，頭戴十二顆星的冠冕； \end{tabularx} \\ \\ \relax
12:2 & \begin{tabularx}{0.7\textwidth}{X} 她懷了孕，在生產的陣痛中疼痛地喊叫。 \end{tabularx} \\ \\ \relax
12:3 & \begin{tabularx}{0.7\textwidth}{X} 天上又出現了另一個兆頭：有一條大紅龍，有七個頭十個角；七個頭上戴著七個冠冕。 \end{tabularx} \\ \\ \relax
12:4 & \begin{tabularx}{0.7\textwidth}{X} 牠的尾巴拖拉著天上星辰的三分之一，把它們摔在地上。然後龍站在那將要生產的婦人面前，等她生產後要吞吃她的孩子。 \end{tabularx} \\ \\ \relax
12:5 & \begin{tabularx}{0.7\textwidth}{X} 婦人生了一個男孩子，就是將來要用鐵杖管轄萬國的；她的孩子被提到神和他寶座那裡去。 \end{tabularx} \\ \\ \relax
12:6 & \begin{tabularx}{0.7\textwidth}{X} 婦人就逃到曠野，在那裡有神給她預備的地方，使她在那裡被供養一千二百六十天。 \end{tabularx} \\ \\ \relax
12:7 & \begin{tabularx}{0.7\textwidth}{X} 天上發生了爭戰。米迦勒同他的使者與龍作戰，龍同牠的使者也起來應戰， \end{tabularx} \\ \\ \relax
12:8 & \begin{tabularx}{0.7\textwidth}{X} 牠們都打敗了，天上再也沒有牠們的地方。 \end{tabularx} \\ \\ \relax
12:9 & \begin{tabularx}{0.7\textwidth}{X} 大龍就是那古蛇，名叫魔鬼，又叫撒但，是迷惑普天下的；牠被摔在地上，牠的使者也一同被摔下去。 \end{tabularx} \\ \\ \relax
12:10 & \begin{tabularx}{0.7\textwidth}{X} 我聽見在天上有大聲音說：「我神的救恩、能力、國度，和他所立的基督的權柄現在都來到了。因為那個在我們神面前、晝夜控告我們弟兄的，已經被摔下去了。 \end{tabularx} \\ \\ \relax
12:11 & \begin{tabularx}{0.7\textwidth}{X} 弟兄勝過那條龍是因羔羊的血，和因自己所見證的道。雖然至於死，他們也不惜自己的性命。 \end{tabularx} \\ \\ \relax
12:12 & \begin{tabularx}{0.7\textwidth}{X} 所以，諸天和住在其中的，你們都快樂吧！只是地和海有禍了！因為魔鬼知道自己的時候不多，就氣憤憤地下到你們那裡去了。」 \end{tabularx} \\ \\ \relax
12:13 & \begin{tabularx}{0.7\textwidth}{X} 龍見自己被摔在地上，就迫害那生男孩子的婦人。 \end{tabularx} \\ \\ \relax
12:14 & \begin{tabularx}{0.7\textwidth}{X} 於是有大鷹的兩個翅膀賜給婦人，讓她能飛到曠野，到自己的地方，躲避那蛇。她在那裡受供養一載二載半載。 \end{tabularx} \\ \\ \relax
12:15 & \begin{tabularx}{0.7\textwidth}{X} 蛇在婦人背後，從口中噴出水來，像河一樣，要將婦人沖走。 \end{tabularx} \\ \\ \relax
12:16 & \begin{tabularx}{0.7\textwidth}{X} 地卻幫助了婦人，開口吞了從龍口噴出來的水。 \end{tabularx} \\ \\ \relax
12:17 & \begin{tabularx}{0.7\textwidth}{X} 於是龍向婦人發怒，去與她其餘的兒女作戰，就是與那些遵守神命令、為耶穌作見證的。那時龍站在海邊沙灘上。 \end{tabularx} \\ \\
[1ex]
\hline
\hline
\end{longtable}
講題:屋頂上的提琴手　Fiddler on the Roof

經文:啟示錄12章1-17節

講員:鄧瑞強博士

場合:香港中文大學崇基學院禮拜堂

日期:2018年6月10日

各位弟兄姊妹,早安.

屋頂上的提琴手(Fiddler on the Roof),是百老匯的一齣音樂劇,後來被拍成電影(在香港,片名譯作《錦繡良緣》).故事的背景,是1905年的俄國,還是沙皇統治的年代,那裡有一條村,住著一群猶太人.在那裡,他們安樂地生活.忽然間,俄國政府驅逐他們離開.於是,全村猶太人各散東西.有人去歐洲,有人去巴勒斯坦,有人去美國.原先,他們有一個家；忽然,他們便無家可歸.原先,他們在熟悉的環境,過著愉快的生活；忽然,一切變得陌生,要在艱難的路上前行.原先,他們在和諧的生活中,經驗上帝的臨在；忽然,他們要在不知前路的旅途中,思考上帝的真實.

主角說:「我們都是屋頂上的提琴手.」

為什麼呢?

一個人為什麼要站在屋頂上呢?

因為我們無論活在這裡或那裡,這都不是我們永恆的家園.我們活在資本主義的世界裡,但我們知道,「資本」抑或背後的「主義」,都不能叫我們安身立命.或許,我們只想簡單地活在一個小社區裡,但小社區很多時被巨大的力量蠶食或吞噬.看來,我們很難在大地上找到一個穩固的立足點.馬克思在《共產黨宣言》中的一句話,或者很能描述我們當下的光景:「一切穩固的東西都煙消雲散」.由於無地立足,於是,要站在「屋頂上」.「屋頂」,是無處立足時勉強的立足處,無處安頓時的臨時安頓處,是飄浮在大海中唯一可以寄托的一塊浮木.

為什麼在「屋頂上」的是一個「提琴手」?

因為人無論多麼艱難,多麼漂泊,人總想追求意義,言說希望,拉出一首美妙的樂章.越是無家可歸,越是以希望為家.越是淒涼,越需要動人的音符.在無處立足的「屋頂上」,心靈深處的「提琴手」,要奏出托著生命的旋律.

一個「提琴手」,如何能在「屋頂上」演奏,而不會失足跌下來?

主角說,保持平衡的關鍵,是「傳統」(tradition).所謂「傳統」,就是指,一種有「意義秩序」的生活.這種有「意義秩序」的生活,背後是「深厚的信仰」.「深厚的信仰」背後,是一位創造生命,救贖生命,並賦予生命以全面意義的上帝.是上帝,支撐著我們,讓我們這些無家可歸的人,這些站在「屋頂上」的人,能像「提琴手」,奏出生命的讚歌,而不至於失平衡而掉下來.我們總需要一個「傳統」,一個生命的支撐點,去托著我們飄搖的人生.

今日講道的經文的內容,其實言說著同樣的故事,有一種強大的力量,不斷驅逐我們,令我們無家可歸,令我們走上「屋頂」.同時,有另一種力量,幫助我們平衡,使我們不至在「屋頂上」掉下來,並且使我們能像「提琴手」,奏出生命之歌.這種力量,就是我們的「傳統」,即:我們的「信仰」.在這「信仰傳統」背後,便是那創造生命,救贖生命,並賦予生命以全面意義的上帝.

今日的講道經文:啟示錄12:1－17

(第一段)

啟12:1　天上現出大異象來:有一個婦人身披日頭,腳踏月亮,頭戴十二星的冠冕.

啟12:2　她懷了孕,在生產的艱難中疼痛呼叫.

啟12:3　天上又現出異象來:有一條大紅龍,七頭十角；七頭上戴著七個冠冕.

啟12:4　牠的尾巴拖拉著天上星辰的三分之一,摔在地上.龍就站在那將要生產的婦人面前,等她生產之後,要吞吃她的孩子.

啟12:5　婦人生了一個男孩子,是將來要用鐵杖轄管〔轄管:原文是牧〕萬國的；她的孩子被提到神寶座那裡去了.

啟12:6　婦人就逃到曠野,在那裡有神給她預備的地方,使她被養活一千二百六十天.

(第二段)

啟12:7　在天上就有了爭戰.米迦勒同他的使者與龍爭戰,龍也同牠的使者去爭戰,

啟12:8　並沒有得勝,天上再沒有牠們的地方.

啟12:9　大龍就是那古蛇,名叫魔鬼,又叫撒旦,是迷惑普天下的.牠被摔在地上,牠的使者也一同被摔下去.

啟12:10　我聽見在天上有大聲音說:「我神的救恩,能力,國度,並他基督的權柄,現在都來到了!因為那在我們神面前晝夜控告我們弟兄的,已經被摔下去了.

啟12:11　弟兄勝過牠,是因羔羊的血和自己所見證的道.他們雖至於死,也不愛惜性命.

啟12:12　所以,諸天和住在其中的,你們都快樂吧!只是地與海有禍了!因為魔鬼知道自己的時候不多,就氣忿忿地下到你們那裡去了.」

(第三段)

啟12:13　龍見自己被摔在地上,就逼迫那生男孩子的婦人.

啟12:14　於是有大鷹的兩個翅膀賜給婦人,叫她能飛到曠野,到自己的地方,躲避那蛇；她在那裡被養活一載二載半載.

啟12:15　蛇就在婦人身後,從口中吐出水來,像河一樣,要將婦人沖去.

啟12:16　地卻幫助婦人,開口吞了從龍口吐出來的水〔原文是河〕.

啟12:17　龍向婦人發怒,去與她其餘的兒女爭戰,這兒女就是那守神誡命,為耶穌作見證的.那時龍就站在海邊的沙上.

(1)第一段(啟12: 1-6):婦人,紅龍,曠野

啟12:1　天上現出大異象來:有一個婦人身披日頭,腳踏月亮,頭戴十二星的冠冕.

啟12:2　她懷了孕,在生產的艱難中疼痛呼叫.

一個異象,所謂「異象」,即人間事情背後的,深層次的真實內容.

這裡有一個婦人,她懷的孩子,明顯是耶穌.這女人包含了夏娃,馬利亞,教會的三重形象.創世記說,夏娃的後裔與蛇(或龍)的後裔彼此為敵,彼此為仇.人間的義者,必與人間的不義者,彼此為敵,這是人間無可避免的情景.在新約的福音書裡,馬利亞是耶穌的母親.馬利亞的兒子,與邪惡的力量對抗,並要克服人間一切的不義.至於教會,教會是耶穌的至愛,以啟示錄的象徵來說,教會是基督的「新娘」.教會是一女性形象,她的生命深處,懷著基督的生命.

這婦人「身披日頭,腳踏月亮,頭戴十二星的冠冕」,「十二」是神的子民的數目,這婦人,代表教會,散發著日與月的光明.「她懷了孕,在生產的艱難中疼痛呼叫」.新約常用「懷孕」來代表「有盼望的艱難」.「懷孕」是艱難的,是漫長的,但同時是有盼望的,只要孩子出世,就會歡天喜地,忘記過去的艱難.這就是神子民的光景,生命是光明的,但無可避免地與黑暗對抗.這一定是艱難的,但這艱難被盼望所覆蓋著.

啟12:3　天上又現出異象來:有一條大紅龍,七頭十角；七頭上戴著七個冠冕.

啟12:4　牠的尾巴拖拉著天上星辰的三分之一,摔在地上.龍就站在那將要生產的婦人面前,等她生產之後,要吞吃她的孩子.

「大紅龍」:「紅」,代表血腥暴力；「龍」,是近東常見的神獸.在德國的柏林,有一間「佩加蒙博物館」(Pergamon Museum),或叫「別迦摩博物館」,當中有兩項大型的重要展品.一是別迦摩的祭壇,這可能就是啟示錄說的別迦摩城的「撒旦的座位」.二是古巴比倫的「伊什塔爾城門」(Ishtar Gate),這城門上很多神獸,其中一隻便是「龍」.在古亞述或古巴比倫的文明裡,「龍」有時代表異教神明,有時代表皇帝.啟示錄在這裡,以「龍」代表魔鬼,牠是邪惡的根源.

這「龍」有「七頭十角」.「七頭」代表什麼?啟示錄的作者為免我們亂猜,自行清楚說明.啟17:9說:「那七頭就是女人所坐的七座山」,這裡提及的女人,是另一個女人,啟示錄稱之為「大淫婦」.「女人所坐的七座山」意指什麼?原來,是指當時羅馬錢幣上的圖案,羅馬女神坐在七座山之上.羅馬,自古以來,被稱為「七山之城」.「龍」的「七頭」,其實指向「羅馬」,這是當時最強大,最倚仗軍力的帝國.「角」,意指「強力」.「十角」,就是指「十足的暴力」.這隻「七頭十角的大紅龍」,集合古亞述,古巴比倫,古羅馬的暴力帝國的性質,代表著人間自古以來所有的強大暴力.這強大暴力,掠奪財富,奴役人民,摧毀人性.這隻怪獸,要吞滅懷孕婦人要生的孩子,這孩子就是基督,是人類救贖的希望.

啟12:5　婦人生了一個男孩子,是將來要用鐵杖轄管〔轄管:原文是牧〕萬國的；她的孩子被提到神寶座那裡去了.

啟12:6　婦人就逃到曠野,在那裡有神給她預備的地方,使她被養活一千二百六十天.

婦人要生的「孩子」,是牧養萬國的.他像牧人,引導軟弱無助的羊,找到生命的水源,找到安歇的地方.龍要吞吃他,神卻保護他.

婦人,即教會,即神的子民,被人間的邪惡和暴力所迫,要逃到「曠野」.「曠野」,是遠離文明的地方,是什麼都沒有而只有神的地方,是無立足處的人的「屋頂」.一個人能在「屋頂」上拉「提琴」嗎?啟示錄說,能.在「龍」的逼迫之下,有神的救贖,有神預備的地方,有神定好的日子時間.

(2)第二段(啟12: 7-12):十架的勝利

啟12:7　在天上就有了爭戰.米迦勒同他的使者與龍爭戰,龍也同牠的使者去爭戰,

啟12:8　並沒有得勝,天上再沒有牠們的地方.

啟12:9　大龍就是那古蛇,名叫魔鬼,又叫撒旦,是迷惑普天下的.牠被摔在地上,牠的使者也一同被摔下去.

這是「天啟文學」(apocalyptic literature)的特色,人間的善惡爭鬥,背後有更深的靈性根源.具體的邪惡背後,有靈性上的邪惡.天使米迦勒站在真理的神的那一方,去與邪惡的魔鬼爭戰.這「大龍」是誰?就是「迷惑普天下」的邪惡力量.「偶像崇拜」,「暴力威脅」,「虛言假語」,「金錢美色」,都可以迷惑人.啟示錄藉這異象,叫我們看清楚,我們生活背後的種種迷惑力量.這些迷惑力量,叫我們走向「虛假」,陶醉在「虛假」裡,也被「虛假」麻醉.

邪惡的力量是真實的,但啟示錄並不承認它們是終極的.天使米迦勒與大龍戰鬥,並且勝過大龍.大龍的邪惡力量,不是終極的.

啟12:10　我聽見在天上有大聲音說:「我神的救恩,能力,國度,並他基督的權柄,現在都來到了!因為那在我們神面前晝夜控告我們弟兄的,已經被摔下去了.

啟12:11　弟兄勝過牠,是因羔羊的血和自己所見證的道.他們雖至於死,也不愛惜性命.

啟12:12　所以,諸天和住在其中的,你們都快樂吧!只是地與海有禍了!因為魔鬼知道自己的時候不多,就氣忿忿地下到你們那裡去了.」

這段經文是本章的核心.邪惡的力量有強大的威嚇性,卻沒有終極性.信徒在這場真假善惡之戰中,扮演著關鍵的角色.信耶穌,請不要常常想著上天堂.信耶穌,是向世界展示對抗虛假的真理,是在「屋頂」上向世人展示仍然可以演奏「提琴」的理由,是在這場必死的戰鬥中展示出生命的遠象.

世上總有一些人,能為了世上的正義與美善,視死如歸.他們的生命,告訴我們,在可見的一切背後,有某種永恆而神聖的東西在.耶穌,作為軟弱的羔羊,以十架的愛,以十架的死,展示一種連死亡也不能克制的生命力,展示一種無邊的邪惡也不能取消的善,展示一種比暴力更強大的柔弱.耶穌勝過邪惡,是透過承擔苦難,是透過苦難中滲透出來的仁愛,是透過絕對的和平.

他的死亡,對抗著塵世一切「成者為王,敗者為寇」的價值.他以這種反抗,去對抗「大紅龍」.他的死,展示了上帝的生命力.這不是失敗,而是一種勝利,一種分享上帝的永恆仁愛的勝利.

信仰的深刻之處,是明白耶穌這種死亡中的勝利.信徒的尊榮,是分享基督的生命,在卑微中展示偉大,在捨棄中顯露豐盛,在受苦中活出榮耀.「十字架」,就是我們賴以存活的「傳統」(tradition).就是這「傳統」,使我們能在「屋頂上」,不至掉下來,反能奏出動人的旋律.

這是我們的盼望,當我們分享基督的生命,活出基督的仁愛時,人間的魔性會被克服,世間的邪惡會被消滅,這就是天國的臨現,這也是啟示錄終極的盼望.

(3)第三段(啟12: 13-17):婦人,紅龍,曠野

啟12:13　龍見自己被摔在地上,就逼迫那生男孩子的婦人.

啟12:14　於是有大鷹的兩個翅膀賜給婦人,叫她能飛到曠野,到自己的地方,躲避那蛇；她在那裡被養活一載二載半載.

再出現第一段經文的情節,龍逼害那婦人.在無處立足時,神給婦人一對翅膀.這象徵對應舊約以賽亞先知所說:「但那等候耶和華的必從新得力,他們必如鷹展翅上騰.」(賽40: 31)婦人飛到哪裡去?到曠野去.在那裡,等待三年半,等待天國的終極來臨.有逼迫,有救贖,有等待.在逼迫中,走上「屋頂」.在救贖中,拿出「提琴」,奏出生命之歌.

啟12:15　蛇就在婦人身後,從口中吐出水來,像河一樣,要將婦人沖去.

啟12:16　地卻幫助婦人,開口吞了從龍口吐出來的水〔原文是河〕.

啟12:17　龍向婦人發怒,去與她其餘的兒女爭戰,這兒女就是那守神誡命,為耶穌作見證的.那時龍就站在海邊的沙上.

這裡重現聖經裡一個重要的救贖主題,就是「出埃及」的主題.水,代表混亂與邪惡,龍要用水去消滅神的子民,正如出埃及時,神的子民面對的是無法越過的紅海.當時,神叫紅海分開,神的子民在海中走過旱地.經文如今,重複這主題,神叫我們在水的混亂中走過旱地.

龍與神的子民的戰鬥,會持續一段時間.信徒能經歷的,是神的解救.在這段時間,信徒要做的,是持守神的誡命,為耶穌作證,即是說,懷著孕婦般的盼望,在不義的社會中,更堅持正義；在謊話的世界,更堅持真相；在暴力的人間,更堅持仁愛.然後,我們便能分享基督復活的勝利.

(4)結語

啟示錄第12章,教我們什麼?

教我們看清楚世界的本相,世界背後隱伏著邪惡,這是迷惑的力量,引導我們沉醉於暴力及虛謊的成果,遠離真理與仁愛.信徒要看清楚邪惡的真面目,要反對之.生命裡的善惡之爭,不能免.信徒看來被邪惡逼得無處立足,但神的救贖在.「出埃及」的故事,永遠在發生著.而我們,有當參與的角色,就是分享基督的生命,讓神透過我們,開展「出埃及」的故事.是我們,讓當今的「出埃及」故事,成為可能.我們是那「屋頂上的提琴手」,在那容易掉下來的「屋頂」,奏出救贖的凱歌.

但願:榮耀歸於聖父,聖子,聖靈.阿們.



\newpage

\section{香柏與芥菜}
\label{sec:414}
\textbf{林豪恩先生 - 香柏與芥菜}
\newline
\newline
\hyperref[sec:413]{< < < PREV SERMON < < <}
~
\hyperref[sec:toc]{[返目錄]}
~
\hyperref[sec:415]{> > > NEXT SERMON > > >}


講題:香柏與芥菜　A Cedar and a Mustard

講員:林豪恩先生

經文:以西結書17:22-24,31:2-11,37:5-10

場合:香港中文大學崇基學院禮拜堂

日期:2018年6月17日

一. 以色列的香柏夢

好像每個地方都有夢想.

例如中國最近常說中國夢,根據百度百科:「中國夢,是中國共產黨第十八次全國代表大會召開以來,習近平總書記所提出的重要指導思想和重要執政理念,正式提出於2012年11月29日.習總書記把「中國夢」定義為「實現中華民族偉大復興,就是中華民族近代以來最偉大夢想」,並且表示這個夢「一定能實現」.「中國夢」的核心目標也可以概括為「兩個一百年」的目標,也就是:到2021年中國共產黨成立100週年和2049年中華人民共和國成立100週年時,逐步並最終順利實現中華民族的偉大復興,具體表現是國家富強,民族振興,人民幸福,實現途徑是走中國特色的社會主義道路,堅持中國特色社會主義理論體系,弘揚民族精神,凝聚中國力量,實施手段是政治,經濟,文化,社會,生態文明五位一體建設.」

美國夢呢?維基百科說:「是一種相信只要經過努力不懈的奮鬥便能在美國獲致更好生活的信仰,亦即人們必須通過自己的工作勤奮,勇氣,創意,和決心邁向富裕,而非依賴於特定的社會階級和他人的援助.許多歐洲移民都是抱持著美國夢的理想前往美國的.」

如果古代的以色列有一個以色列夢,這夢會是怎樣的呢?我們可以透過今天的舊約經課以西結書去窺探這個以色列夢:

17:22主耶和華如此說:「我要從香柏樹高高的樹梢摘取並栽上,從頂端的嫩枝中折下一嫩枝,栽於極高的山上,

17:23栽在以色列高處的山上.它就生枝,結果,成為高大的香柏樹,各類飛禽中的鳥都來宿在其下,宿在枝子的蔭下.

17:24田野的樹木因此就知道是我─耶和華使高樹矮小,使矮樹高大,使綠樹枯乾,使枯樹發旺.我─耶和華說了這話,就必成就.」

我們暫且稱它為「香柏夢」.

聖經百科全書說:「香柏木」學名「雪松」是古代世界常見的一種松屬常青樹.除了幾個例外的情況外,聖經一般提到的「香柏樹」都是指著名的「利巴嫩的香柏樹」(黎巴嫩雪松).雪松是一種高貴的樹木,是以色列人所認識的樹木中最高和最大的一種.其生長速度頗快,可長達120呎高,樹幹直徑可達8呎.黎巴嫩的雪松為人鍾愛,不單因為它生命力強,壯麗高雅,壽命極長,也因為它氣味芳香和木質堅硬.由於如此受歡迎,黎巴嫩雪松在1998年被列入《世界自然保護聯盟》瀕危物種紅色名錄.在文化層面,雪松是用來象徵莊嚴,有力,高貴,雄壯,巍峨和宏偉的.在以西結書十七章3節,22節至24節,先知用這些高貴的樹中之王來象徵世間的力量,權能和威榮.

三十一章2至18節

31:2「人子啊,你要對埃及王法老和他的軍隊說:論到你的強盛,誰能與你相比呢?

31:3看哪,亞述是黎巴嫩的香柏樹,枝條榮美,蔭密如林,極其高大,樹頂高聳入雲.

31:4眾水使它生長,深水使它長高；所栽之地有江河環繞,汊出的水道流至田野的樹木.

31:5所以它高大超過田野的樹木；生長時因水源豐沛,枝子繁多,枝條增長.

31:6空中所有的飛鳥在枝子上搭窩,野地所有的走獸在枝條下生子,所有的大國也在它的蔭下居住.

31:7它樹大枝長,極為榮美,因它的根在眾水之旁.

31:8上帝園中的香柏樹不能遮蔽它；松樹不及它的枝子,楓樹不及它的枝條,上帝園中的樹都沒有它榮美.

31:9我使它枝條繁多,極為榮美；在上帝的園中,伊甸所有的樹都嫉妒它.」

二. 香柏夢竟然變成芥菜夢

香柏樹令人想到王權的尊貴和國家的強盛,今天誦讀的詩篇更以香柏樹來比喻個人的昌盛:

92:12義人要興旺如棕樹,生長如黎巴嫩的香柏樹.

92:13他們栽於耶和華的殿中,發旺在我們上帝的院裡.

92:14他們髮白的時候仍結果子,而且鮮美多汁,

92:15好顯明耶和華是正直的；他是我的磐石,在他毫無不義.

自從大衛所羅門王朝以來,古代的以色列人都在想像香柏夢,追求香柏夢,耶穌道成內身來到他們當中,宣講的卻不是香柏夢,竟然是一個芥菜夢:

4:30耶穌又說:「我們可用甚麼來比擬上帝的國呢?可用甚麼比喻來說明呢?

4:31它像一粒芥菜種,種在地裡的時候,雖比地上所有的種子都小,

4:32但種下去以後,它長起來,比各樣的菜都大,又長出大枝,以致天上的飛鳥可以在它的蔭下築巢.」

在耶穌所講的芥菜夢裡面,仍然有「國」這個概念,仍然有「由小變大」這元素,因此,可能仍讓人覺得「芥菜夢」與「香柏夢」一脈相承.不過,若面對現實,仔細觀察,芥菜和香柏實在相差太離了.

2.1. 就算不是菜,也只是灌木,而不是樹

香柏是高大的喬木,長於高山上；芥菜呢?聽在廣東人的耳中,只會想到「芥菜鹹蛋肉片湯」的芥菜,是一種種於菜田的草本植物而已.有人說,廣東人誤會了,這不是煮湯的芥菜,而是用來做「芥末醬」的芥末樹.在以色列地方,這些芥末樹可長成一棵高過人的樹,甚至到八尺,十尺.不過,無論八尺還是十尺,芥末樹始終是生長於一般地方的灌木,而不是長於高山上的高大的喬木.芥菜長大後,雖然比它的種子要大了很多倍,但芥菜仍然是芥菜,它沒有變成香柏,也沒有任何部分似香柏.

2.2. 普通植物,亳不突出

教會所追求的是「香柏夢」還是「芥菜夢」呢?教會在社會的形象是香柏樹還是芥末樹呢?在我們的社會中,提到教會,人們想到甚麼呢?可能是聖誕節和復活節長假期,可能是莊嚴宏偉的教堂,可能是教會名校,教會醫院,或者基督教社會服務機構等.基督教似乎是社會主流文化之一,傳媒每年報導大主教發表的聖誕及新年文告,教會也成為政府在教育,醫療,社會服務的重要夥伴.教會如此的「高大」,未必是因為她本身是「香柏樹」,而是在英治下眾多政治及社會條件所形成的特殊狀況.這些條件會隨著政治及社會環境改變而改變,當條件消失的時候,教會就會發現原來自己本來沒有那麼高大,沒有那麼主流,也未必有實力能夠支撐得起那麼龐大的身軀.畢竟,基督徒只是人口的少數,在新的管治者眼中,教會或許只是被要求與新時代相適應的邊緣群體.

懷著「香柏夢」的,不容易接受「芥菜夢」.或者,「香柏夢」的確令人感覺良好,這份飄飄然的感覺,可以令人失去自知之明,也可以令人不自量力.以西結先知曾說:

31:10所以主耶和華如此說:「因它高大,樹頂高聳入雲,心高氣傲,

31:11我要把它交給列國中強人的手裡,他們必定按它的罪惡懲治它.我已經驅逐它.

以賽亞先知也說:

9:8主向雅各家發出言語,主的話臨到以色列家.

9:9眾百姓,就是以法蓮和撒瑪利亞的居民,都將知道；他們憑驕傲自大的心說:

9:10「磚塊掉落了,我們要鑿石頭重建；桑樹砍了,我們要改種香柏樹.」

曾經何時,在某種歷史條件之下,教會以為自己是香柏樹:莊嚴,有力,高貴,雄壯,巍峨和宏偉.九七前,聖公會主教在政府的禮賓排名極高,九七後,基督教領袖在特區政府的禮賓排名下調,與其他宗教領袖相若.不過,大型基督教活動卻喜歡邀請政府高層出席.某年,某大型公開祈禱會,邀請當時民望低迷的特首董建華出席,為他祈禱和表達支持.事件當時在教內外引起頗大反響.在新時代之下,其實教會如履薄冰.如果,有朝一日,政府撤去對教會辦教育,醫療,社會服務等的資源,教會還有能力去維持這些服務嗎?如果教會沒有能力維持這些服務,恐怕會連寄居在這些服務單位的聚會地方都失去.這並不是匪夷所思,在一河之隔的大地,這是常態.

唯願教會及眾信徒心不驕,氣不傲,知道自己只是一棵小芥菜而己.

2.3. 芥菜燒不盡,靈風吹又生

聖經百科全書指出,「芥菜屬於十字花科,當地有白芥菜和黑芥菜兩種,這兩種芥菜都是一年生草本植物.在耶穌的時代,芥菜的種子是一般人所知道的種子中最小的.芥菜專供栽培生產辛辣的種子和可食的葉片,二千年來一直是以色列重要經濟作物.黑芥菜種子氣味強,適於做芥末醬,而白芥菜種子的氣味較淡,為泡菜,沙拉醬的防腐劑.兩者的種子均可磨成粉加入麵糰中或製油.」這介紹告訴我們,芥菜是常見及平凡的植物,或者由於它的常見及平凡,它無處不在,生命力頑強,就算是炎熱的地方,乾燥的土地,都能夠生長；同時,也由於芥菜的種子細小,能夠輕易散播,就算被收集,也容易有「漏網之魚」.正如中文所說:「野草燒不盡,春風吹又生」.

百度網頁上有這樣一段關於中國教會的文字:「作為中國五大宗教之一的基督教,於「文革」開始不久,即受到嚴厲的攻擊:教堂被佔領,聖經遭焚燒,聚會被禁止.無論是「三自會」的愛國進步領袖或一般的教牧人員和信徒,多被視為「牛鬼蛇神」,「專制物件」,遭受殘酷的批鬥和不同程度的逼迫.有的被活活摧殘至死；有的經不起長期折磨而精神錯亂或斷然結束個人的生命；也有的為求自保而公開否認其信仰,甚或積極陷害,檢舉和批判教會中的牧師,信徒,以表示對黨的忠心.然而,不少基督徒在淩辱困苦中仍堅守他們的信仰,就是在最艱難的關頭也沒有停止聚會或放棄傳福音的機會.故此,有形的教會雖被取消,無形的教會卻在中國每一個角落悄然生長.」

野草燒而不盡,等待春風吹來.春風是甚麼呢?聖經經常以「風」來比喻「聖靈」:

以西結先知在另一著名篇章如此描述:

37:5主耶和華對這些骸骨如此說:『看哪,我必使氣息進入你們裡面,你們就要活過來.

37:6我要給你們加上筋,長出肉,又給你們包上皮,使氣息進入你們裡面,你們就要活過來；你們就知道我是耶和華.』」

37:7於是,我遵命說預言.正說預言的時候,有響聲,看哪,有地震；骨與骨彼此接連.

37:8我觀看,看哪,骸骨上面有筋,長了肉,又包上皮,只是裡面還沒有氣息.

37:9耶和華對我說:「人子啊,你要說預言,向風說預言.你要說,耶和華如此說:氣息啊,要從四方而來,吹在這些被殺的人身上,使他們活過來.」

37:10於是我遵命說預言,氣息就進入骸骨,骸骨就活過來,並且用腳站起來,成為極大的軍隊.

耶穌在約翰福音中也曾如此描述聖靈:

3:8風隨著意思吹,你聽見風的聲音,卻不知道是從哪裡來,往哪裡去；凡從聖靈生的也是如此.

風,能夠把細小的芥菜種子輕易吹起,飄落於不同的地方；風,也為飄落於不同地方的種子帶來生機,發芽生長.

三. 總結:香柏與芥菜都能蔭鳥於下

香柏樹長於高山上,高聳入雲,高貴,有力,雄壯,巍峨和宏偉；芥菜樹是平凡的,常見的植物,既不高大,也不華麗,沒有令人景仰之處,甚至也沒有吸引人注目的地方.縱使兩者如此不同,讀者有沒有留意,在聖經作者的筆下,它們有一處地方是相同的?

舊約經課以西結書說:

17:22主耶和華如此說:「我要從香柏樹高高的樹梢摘取並栽上,從頂端的嫩枝中折下一嫩枝,栽於極高的山上,

17:23栽在以色列高處的山上.它就生枝,結果,成為高大的香柏樹,各類飛禽中的鳥都來宿在其下,宿在枝子的蔭下.

福音經課馬可福音說:

4:30耶穌又說:「我們可用甚麼來比擬上帝的國呢?可用甚麼比喻來說明呢?

4:31它像一粒芥菜種,種在地裡的時候,雖比地上所有的種子都小,

4:32但種下去以後,它長起來,比各樣的菜都大,又長出大枝,以致天上的飛鳥可以在它的蔭下築巢.」

無論是高大的香柏樹,還是矮小的芥菜樹,它們都能夠惠及蒼生,而惠及蒼生的狀況竟然都是一樣的,就是蔭護飛鳥,讓飛鳥宿於其下.這會否是上主在古時對祂的選民的期望?會否也是耶穌所帶來的上帝的國的特徵呢?

在某些歷史時空之下,教會可能長得像香柏樹一樣高大,在那些時候,教會會否變得心高氣傲,忘記了蔭疪有需要的人呢?又在某些歷史時空之下,教會可能只是平平無奇的芥菜樹,教會會否顧影自憐,完全沒有想像過耶穌所說的比喻,芥菜樹也能為有需要的人提供護蔭呢?

五月下旬,我們舉辦了一次往印度的宗教文化考察及服務團,十八位學生說服了他們的家人讓他們去冒一次險.印度是一個高度宗教化的地方,有資料說十多億人口中約80\%是印度教徒,約14\%的伊斯蘭教徒,基督教(包括天主教徒)只有2\%.由數字可見,基督徒是社會中的極小數.這個小數群體,常常受到壓力.2008年,在印度東部發生宗教滅絕行動,激進印度教徒瘋狂攻擊當地天主教徒,有人被燒死,有修女被強姦,多達50000人被逼逃亡.當時,梵蒂岡嚴厲譴責印度教徒的暴行.

這棵搖搖欲墜的小芥菜,卻沒有放棄護蔭在印度社會中被歧視和壓迫的賤民,他們稱為Dalits,是untouchable的意思.達利人基金會(Dalit Foundation)指出:平均每一個小時就有兩位達利人遭受到暴力相待；每一天就有三位達利婦女被強暴,兩位達利人被殺害,兩戶達利人家被燒毀.有誰願意接觸不可接觸的呢(touch the untouchable)?人們聽過德蘭修女及她的修會.除此之外,還有不少宣教團體進入Dalits群體,據教會估計,賤民佔印度基督徒六成之多.

我們拜訪了為Dalits爭取人權的天主教機構,也拜訪了一個細小的基督教團體.這團體在貧窮地區提供為兒童及婦女提供服務,我們就是在他們主辦的一所學校作義務教學.

如果上帝的選民長得像香柏樹一樣高大,他們應該如以西結先知所說,讓「各類飛禽中的鳥都來宿在其下,宿在枝子的蔭下」；如果只是芥菜樹,也可以如耶穌基督所說的,「天上的飛鳥可以在它的蔭下築巢」.



\newpage

\section{返璞歸真}
\label{sec:415}
\textbf{王家輝牧師 - 返璞歸真}
\newline
\newline
\hyperref[sec:414]{< < < PREV SERMON < < <}
~
\hyperref[sec:toc]{[返目錄]}
~
\hyperref[sec:416]{> > > NEXT SERMON > > >}


哥林多後書 8:7-15
\newline
\begin{longtable}{cl}
\hline
\hline
章節 & 經文 (和合本修訂版)\\
\hline
8:7 & \begin{tabularx}{0.7\textwidth}{X} 既然你們在信心、口才、知識、萬分的熱忱，以及我們對你們的愛心上，都勝人一等，那麼，當在這慈善的事上也要勝人一等。 \end{tabularx} \\ \\ \relax
8:8 & \begin{tabularx}{0.7\textwidth}{X} 我說這話，並不是命令你們，而是藉著別人的熱忱來考驗你們愛心的真誠。 \end{tabularx} \\ \\ \relax
8:9 & \begin{tabularx}{0.7\textwidth}{X} 你們知道我們主耶穌基督的恩典：他本是富足，卻為你們成了貧窮，好使你們因他的貧窮而成為富足。 \end{tabularx} \\ \\ \relax
8:10 & \begin{tabularx}{0.7\textwidth}{X} 我在這事上把我的意見告訴你們，是對你們有益，因為你們開始辦這事，而且起此心意已經有一年了。 \end{tabularx} \\ \\ \relax
8:11 & \begin{tabularx}{0.7\textwidth}{X} 如今就當辦成這事，既然有願做的心，也當照你們所有的去辦成。 \end{tabularx} \\ \\ \relax
8:12 & \begin{tabularx}{0.7\textwidth}{X} 因為人只要有願做的心，必照他所有的蒙悅納，並不是照他所沒有的。 \end{tabularx} \\ \\ \relax
8:13 & \begin{tabularx}{0.7\textwidth}{X} 我不是要別人輕鬆，你們受累，而是要均勻： \end{tabularx} \\ \\ \relax
8:14 & \begin{tabularx}{0.7\textwidth}{X} 就是要你們現在的富餘補他們的不足，使他們的富餘將來也可以補你們的不足，這就均勻了。 \end{tabularx} \\ \\ \relax
8:15 & \begin{tabularx}{0.7\textwidth}{X} 如經上所記：多收的沒有餘，少收的也沒有缺。 \end{tabularx} \\ \\
[1ex]
\hline
\hline
\end{longtable}
講題:返璞歸真　Back to Basics

經文:哥林多後書8:7-15

講員:王家輝牧師(香港中文大學崇基學院神學院教牧事工部主任兼校牧)

場合:香港中文大學崇基學院禮拜堂

日期:2018年7月1日

1弟兄們,我把上帝賜給馬其頓眾教會的恩告訴你們,2就是他們在患難中受大試煉的時候,仍有滿足的快樂,在極窮之間還格外顯出他們樂捐的厚恩.3我可以證明,他們是按著力量,而且也過了力量,自己甘心樂意地捐助,4再三地求我們,准他們在這供給聖徒的恩情上有分；5並且他們所做的,不但照我們所想望的,更照上帝的旨意先把自己獻給主,又歸附了我們.6因此我勸提多,既然在你們中間開辦這慈惠的事,就當辦成了.7你們既然在信心,口才,知識,熱心,和待我們的愛心上,都格外顯出滿足來,就當在這慈惠的事上也格外顯出滿足來.8我說這話,不是吩咐你們,乃是藉著別人的熱心試驗你們愛心的實在.9你們知道我們主耶穌基督的恩典:他本來富足,卻為你們成了貧窮,叫你們因他的貧窮,可以成為富足.10我在這事上把我的意見告訴你們,是與你們有益；因為你們下手辦這事,而且起此心意,已經有一年了,11如今就當辦成這事.既有願做的心,也當照你們所有的去辦成.12因為人若有願做的心,必蒙悅納,乃是照他所有的,並不是照他所無的.13我原不是要別人輕省,你們受累,14乃要均平,就是要你們的富餘,現在可以補他們的不足,使他們的富餘,將來也可以補你們的不足,這就均平了.15如經上所記:多收的也沒有餘；少收的也沒有缺.

引言

今天是「香港特別行政區成立紀念日」,有些教會根據香港基督教協進會的建議,把這個主日定為「香港主日」.期望這個主日的崇拜信息跟香港有直接關係.此外,在今日亦會有不同團體舉行「慶祝香港主權回歸祖國」的活動,而另一邊,也會有遊行示威,顯示在香港主權回歸中國後,政府的施政表現未符期望,同時也包括各項政治,經濟,民生問題的訴求.

有基督徒會參加慶回歸活動,有基督徒會去遊行示威,估計有更多基督徒什麼也不會參與.無論你的選擇如何,既然我們身處這個城市,我們需要對這個城市有些想像,因為在這裡發生的一切對我們有直接或間接的影響.

若然基督徒要對一個地方作出想像,那就當然要從聖經開始.在今個主日的經課之中,哥林多後書八章7-15節正好給我們為香港作想像帶來一點提示,而「返璞歸真」就是我在閱讀這段經文後所得的反思.為了讓大家對這段經文所講述的情況有更準確的了解,故此我們需要從第1節開始講.

馬其頓教會的榜樣

一般人會在什麼情況下,才願意跟別人分享自己所有?大概是在自己有剩餘,有能力的時候才會做吧!假若自己也處於短缺的情況,那又怎會有能力實踐分享呢?所謂「自顧不暇」,自己的困難也沒有能力解決,何來資源或能力助人解困呢?可是馬其頓眾教會卻有不一樣的示範.保羅希望哥林多教會可以向他們借鏡.

保羅從提多的口中知道哥林多教會是願意慷慨分享的教會,然而若跟馬其頓眾教會相比,保羅認為馬其頓眾教會更是值得效法,堪作哥林多教會學習的榜樣.

第1節中提到「馬其頓眾教會」,即是不指一間教會.當中包括那些教會呢?在新約時代,腓立比,帖撒羅尼迦和庇哩亞等地均在馬其頓省內,因此馬其頓眾教會也該包括這些教會在內.當時這些地方的教會正面對很大的困難(參腓1:29-30；帖前1:6；2:14；帖後1:4等).他們身處外邦人的社群,故此當地不信主的外邦人排斥他們.同時間猶太人(因為拒絕承認基督是救主)亦攻擊他們.然而這些信徒卻是堅守信仰,並且滿有喜樂.

第2節形容他們的狀況:「在極窮之間」.他們可能因受逼迫以致生活艱苦.又或是因著羅馬政府苛徵重稅,以致整個地方根本不會富裕.進而教會的經濟就更是雪上加霜了.縱然如此,馬其頓眾教會並沒有只關注處理自己的難題.他們在缺乏中仍堅持慷慨分享.參腓4:10-18,他們多次支持保羅在佈道旅程中的需要.他們所作出的奉獻支持,在保羅眼中似乎已超出他們可負擔的.

馬其頓眾教會的慷慨分享,絕非易事,他們這樣作並非為顯出自己的實力,而是照著上帝給他們的信心而擺上.馬其頓教會所作的可說是超越常理.他們不是因自己有餘而幫助有需要的人,他們同樣在艱難和缺乏之中,但卻沒有忘記有些地方比他們更困難,更需要援手.

自己也在難關中,但仍然願意幫助人,結果會是怎樣呢?這使我想起在列王記上十七章8-16記載一個外邦寡婦怎樣幫助先知以利亞.當時是旱災,因為以利亞曾向亞哈王說,若他幾年不禱告,就不會降霜,不會降雨,結果真的如以利亞所言一樣.沒有水,沒有糧的日子當然十分困難,就在這個時候,耶和華上帝吩咐以利亞到外邦人的地方(就是西頓的撒勒法).耶和華在那裡預備了一位寡婦幫助以利亞.因為她和兒子只有少許的麵粉和油,頂多可以做到充飢一餐的餅.至於下一餐如何,她們也不能確定,根本上無法接濟以利亞.可是那位寡婦只是憑藉愛心和信心,結果麵粉和油一直沒有缺乏,他們安渡難關.當然這是神蹟,否則我們也想不到有其他的可能.

香港的境況

在我們今日的生活中,我們是否仍有這麼大的愛心和信心呢?我估在今時今日的香港,這是極不容易做的到的事情.因為香港人活在巨大的經濟壓力之下.有剛發表的研究指出,香港奪得「全球外派員工生活成本最高的城市」.去年的榜首是非洲國家安哥拉的首都羅安達.第二名是東京,新加坡,上海和北京則分別排第4,第7及第9.

在這樣的情況下,我們視其他人為競爭有限資源的對手.或者,我們會想盡辦法在對方身上獲取最多的利益.「賺到盡」就是我們這個城市的精神!「賺到盡」跟「慷慨分享」和「在缺乏中仍願意分享」絕對是兩碼子的事.香港的住屋問題正好反映了「賺到盡」的情況.

在葵涌有一幢逾40年樓齡的私人樓宇,一個近300呎的單位「一tong 四」,每個tong 房實用面積不足30呎,月租卻要2,000元,即平均呎租60元.

另有一些tong 房設計成像「太空艙」(實質是一個床位).例如有一個是位於德輔道西,鄰近港鐵西營盤站,業主以800多萬元購入約700呎,三房一廳的單位.如果整個單位出租,每月租金大概是40,000元.可是買家把它tong 成總共有12個「太空艙」,艙內床褥約有2米乘1米,每月以約4,800元租出.換言之,若全部可以租出,業主每月可以有57,600元的租金收入.

越住越貴,亦越住越細是今日許多香港人的經驗.在擠迫的空間,工時又長的情況下,代價是什麼呢?當然是犧牲了健康,精神壓力大,人際關係亦出現問題.我估這些也是大家不會感到陌生的情況.大家為自己的生活已經是感到筋疲力竭,還有什麼餘力可以去支援別人呢?

保羅給哥林多教會的挑戰

哥林多教會因位處哥林多城,故此比起馬其頓省內各個城市,哥林多是相對富裕的.他們在優越的環境下發展,能夠不忘他人的需要,實在難得,因此保羅也對他們的分享行動予以肯定.只是他們有了計劃,還未完成,故此保羅提醒他們不要停留在只有感動的層面,而是認真把這一切化為實際行動.

保羅不但以馬其頓眾教會的榜樣鼓勵哥林多教會,他更以耶穌基督的生命來提醒哥林多教會要效法主.主耶穌放棄天上榮華,降生為人,且卑微至死(見腓2:6-8),成就救恩,為要使哥林多人和世上所有人因信稱義,承受天上的產業.因此凡屬主的人應受耶穌捨身大愛的感動,以祂的犧牲為榜樣,慷慨分享自己所有一切.這個提醒是十分重要,而且先後也不能倒轉.

我們能夠跟別人分享自己所有的,全因為上主已給我們很大的祝福,充滿恩典.既是如此,我們要懂得感恩.感恩的方法就是把自己得到的恩典,也成為別人的祝福,這是最好的感恩表現.

正如我們當中有不少是老師.我們從老師身上得到許多教導,學有所成.我們怎樣報答老師呢?我估並非要送什麼禮物給老師或請他一起進餐.當然這樣做,老師也會歡喜,但老師更大的快樂是來自,你能夠學以至用,貢獻社會,造福他人.這比起任何禮物更加珍貴.

保羅在13至15節用這樣的說話勉勵哥林多教會「13我原不是要別人輕省,你們受累,14乃要均平,就是要你們的富餘,現在可以補他們的不足,使他們的富餘,將來也可以補你們的不足,這就均平了.15如經上所記:多收的也沒有餘；少收的也沒有缺.」

「多收的也沒有餘；少收的也沒有缺」.是引自出埃及記16章18節；以色列人在曠野時,各按所需,蒙上帝供應嗎哪,沒有人可以得到更多,亦沒有人會缺乏.這是表明上帝的心意,讓信徒互相補足,富裕的賙濟貧窮的,而貧窮的亦不用擔憂生活.

這是極美麗的圖畫,可是因為人的貪婪和自私,要實現這樣的理想是極其困難的.教會應當在世界成帶來示範作用.

為這城市的想像

既然「多收的也沒有餘；少收的也沒有缺」是重要的方向,那麼這也是我們為這城市所作想像的參考.是否以共產主義的制度套進香港就可以實現「多收的也沒有餘；少收的也沒有缺」呢?我估就是今日奉行共產主義的國家也知道不可能.

我想集中讓教會想想,我們可以怎樣在此時此地實現這句說話.當然由教會辦理的慈惠事業可以發揮到一定作用,例如扶貧工作,或處理跨代貧窮的問題.這些工作教會向來也做得不錯.然而我估還有其他的可能性.

現時居住的空間越來越有限,有不少小孩子在狹窄的環境成長,連做功課,溫習的空間也缺乏.教會有相對比較大的空間,而除了舉行聚會時需要較多空間外,其餘的時間,也是可以釋放一些空間給貧窮家庭的小孩子.

還有教會中有不少人可以擔任導師指導這些小孩子,因為他們的家長,可能也是工作時間非常長,在教育水平方面可能也有限,難以幫助小孩子在學習上的需要.教會的空間和人力正好可以幫助這些有需要的人.

此外,我在崇基學院神學院中認識一些兼讀的同學,他們各自有自己的專業和工作,更有些是提早退休的.他們有豐富的人生經驗,學識,經濟基礎和廣闊的人際網絡.他們有部份是在商界多年的.他們集合起來以「使命營商」為理念,開設社會企業,把賺得來的金錢,部份投放在慈善用途,而另外亦僱用弱勢群體(如新移民,少數族裔,精神病康復者,更生人士等).這是非常有意思,也是在實踐聖經所說的「多收的也沒有餘；少收的也沒有缺」.

這正正是一個城市能夠「返璞歸真」的最佳演繹.求主恩助和加我們力量,使眾人同得恩惠與祝福.阿們.



\newpage

\section{耶穌論施捨}
\label{sec:416}
\textbf{戴國璋牧師 - 耶穌論施捨}
\newline
\newline
\hyperref[sec:415]{< < < PREV SERMON < < <}
~
\hyperref[sec:toc]{[返目錄]}
~
\hyperref[sec:417]{> > > NEXT SERMON > > >}


馬太福音 6:1-4
\newline
\begin{longtable}{cl}
\hline
\hline
章節 & 經文 (和合本修訂版)\\
\hline
6:1 & \begin{tabularx}{0.7\textwidth}{X} 「你們要謹慎，不可故意在人面前表現虔誠，叫他們看見，若是這樣，就不能得你們天父的賞賜了。 \end{tabularx} \\ \\ \relax
6:2 & \begin{tabularx}{0.7\textwidth}{X} 「所以，你施捨的時候，不可叫人在你前面吹號，像那假冒為善的人在會堂裡和街道上所做的，故意要得人的稱讚。我實在告訴你們，他們已經得了他們的賞賜。 \end{tabularx} \\ \\ \relax
6:3 & \begin{tabularx}{0.7\textwidth}{X} 你施捨的時候，不要讓左手知道右手所做的， \end{tabularx} \\ \\ \relax
6:4 & \begin{tabularx}{0.7\textwidth}{X} 好使你隱祕地施捨；你父在隱祕中察看，必然賞賜你。」 \end{tabularx} \\ \\
[1ex]
\hline
\hline
\end{longtable}
講題:耶穌論施捨　Jesus’ Words on Giving Alms

經文:馬太福音6章1-4節

講員:戴國璋牧師

場合:香港中文大學崇基學院禮拜堂

日期:2018年7月8日

在登山寶訓的中段,耶穌提及三件「善事 」:「施捨」,「禱告」與「禁食」.「善事」,即「善美的事」,原來就是「義」這個字.指的是猶太教傳統所認為關乎虔敬的三事.不久前,耶穌曾說過:「你們的義,若不勝過文士和法利賽人的義,斷不能進天國.」對這三項獨特的宗教行為,耶穌都有新的見解.今天我們將焦點放在第一項:施捨.

一,施捨的真義

按中國人的慣用,施捨是指把財物給行乞的人.這是一個狹義的用法.這詞本指佛教的善信對行腳僧人的布施.佛教的僧侶(小乘)修行時以嗟食行乞為生.傳說佛祖喬達摩修行時也是隨從此傳統風俗.在猶太教中,那是指較廣義的周濟貧窮人.耶穌用的吹號隱喻,正反映了這一點 .

「吹號」,其實是一個相關語.在猶太的俗語中作為對周濟貧窮人的捐輸事.在耶穌時代,猶太教已發展到十分重視社會的關懷:周濟窮人.並視之為宗教行為上的 美德.這項美德遭到耶穌的肯定,他的許多教訓都跟周濟窮人有關.事實上,早期的基督教會從她姊妹猶太教那裡繼承了這項美善事業,且發揚光大.這種可誇,我們能否持續下去呢 ?今天,香港教會競相建立超級教會,建做美輪美奐的教堂,其豪奢之勢叫人咋舌.建立的許多輔助事業,如醫療,教育等事,則只顧服務大客戶,盡其力量吸金.這些事都不是我們前輩所樂為的.

二,施捨的真情

耶穌說人施捨之時勿在前收吹號,這不是藉以炫耀的東西.相反施捨時要叫左手不知右手所作.這句話一直帶來解釋的困難.有認為耶穌不過用誇張的修辭表達不造作的重要.也有人認為這是一句猶太的慣用語,意思是要人不在心中多想念自己的善.今天,許多人會為自己做過的好事不斷吹詡,不斷提醒,甚至在未曾出過力的事上爭功 .這樣地解釋也未嘗不可.但有沒有一個解釋可以合宜耶穌說話那字面的意義呢?耶穌的話相信指行善事要出乎自然, 自然而然地,不存造作,那就是最好的.在行善時,若是下意識地去行,而不是意識地,那就是左手不知右手所行.

出於意識而行善助人,其實值得表揚.基督徒許多善德的鍛煉吩咐,都是這樣.

保羅說:凡事要察驗,善美的要持守,惡事要禁戒.若無意識又怎斷善惡?但能下意識地行善,那是出於為善的生命本質,值得慶賀.

三,施捨的真回報

人在做,天在看,這是一句近年常聽見的說話.其實後面還跟著一句話:天會報應.這句今天常聽的說話主要取負面意義.但耶穌對善事的解說,主要取正面的意義.耶穌對三項善事都以此作結:「你父在暗中察看,必然報答你.」何時報答,未有明言,卻必然報答.聖經就對義人之義的回報,有多年後回報,也有在永生裡回報,但必有回報.要到永生方回報,有廣東人說:「風水佬呃人」的嫌疑.但我們許多前輩就是這樣地相信.

耶穌不單叫我們莫藉善行取人的榮耀,更進一步,他教導要把施捨為作在不能回報你的人身上.因為人既不能回報你,鐵定要由天回報,豈不更好.



\newpage

\section{從迦密山回望世界}
\label{sec:417}
\textbf{梁元生院長 - 從迦密山回望世界}
\newline
\newline
\hyperref[sec:416]{< < < PREV SERMON < < <}
~
\hyperref[sec:toc]{[返目錄]}
~
\hyperref[sec:418]{> > > NEXT SERMON > > >}


列王記上 18:17-22
\newline
\begin{longtable}{cl}
\hline
\hline
章節 & 經文 (和合本修訂版)\\
\hline
18:17 & \begin{tabularx}{0.7\textwidth}{X} 亞哈見了以利亞，就說：「真的是你嗎？你這使以色列遭殃的人！」 \end{tabularx} \\ \\ \relax
18:18 & \begin{tabularx}{0.7\textwidth}{X} 以利亞說：「使以色列遭殃的不是我，而是你和你的父家，因為你們離棄耶和華的誡命，去隨從巴力。 \end{tabularx} \\ \\ \relax
18:19 & \begin{tabularx}{0.7\textwidth}{X} 現在你要派人去召集以色列眾人，以及耶洗別所供養的四百五十個巴力的先知和四百個亞舍拉的先知，叫他們都上迦密山到我這裡來。」 \end{tabularx} \\ \\ \relax
18:20 & \begin{tabularx}{0.7\textwidth}{X} 亞哈就派人到以色列眾人那裡，召集先知上迦密山。 \end{tabularx} \\ \\ \relax
18:21 & \begin{tabularx}{0.7\textwidth}{X} 以利亞近前來對眾百姓說：「你們心持二意要到幾時呢？如果耶和華是神，就當順從耶和華；如果是巴力，就當順從巴力。」百姓一言不答。 \end{tabularx} \\ \\ \relax
18:22 & \begin{tabularx}{0.7\textwidth}{X} 以利亞對百姓說：「作耶和華先知的只剩下我一個；巴力的先知卻有四百五十人。 \end{tabularx} \\ \\
[1ex]
\hline
\hline
\end{longtable}
馬可福音 6:14-29
\newline
\begin{longtable}{cl}
\hline
\hline
章節 & 經文 (和合本修訂版)\\
\hline
6:14 & \begin{tabularx}{0.7\textwidth}{X} 耶穌的名聲傳開了，希律王也聽見。有人說：「施洗的約翰從死人中復活了，因此才有這些異能在他裡面運行。」 \end{tabularx} \\ \\ \relax
6:15 & \begin{tabularx}{0.7\textwidth}{X} 但別人說：「他是以利亞。」又有人說：「是先知，正如先知中的一位。」 \end{tabularx} \\ \\ \relax
6:16 & \begin{tabularx}{0.7\textwidth}{X} 希律聽見卻說：「是我所斬的約翰，他復活了。」 \end{tabularx} \\ \\ \relax
6:17 & \begin{tabularx}{0.7\textwidth}{X} 原來，希律為他兄弟腓力的妻子希羅底的緣故，派人去抓了約翰，把他綁了在監獄裡，因為希律已經娶了那婦人。 \end{tabularx} \\ \\ \relax
6:18 & \begin{tabularx}{0.7\textwidth}{X} 約翰曾對希律說：「你佔有你兄弟的妻子是不合法的。」 \end{tabularx} \\ \\ \relax
6:19 & \begin{tabularx}{0.7\textwidth}{X} 於是希羅底懷恨他，想要殺他，只是不能。 \end{tabularx} \\ \\ \relax
6:20 & \begin{tabularx}{0.7\textwidth}{X} 因為希律怕約翰，知道他是義人，是聖人，所以就保護他，雖然聽了他的講論十分困惑，仍然樂意聽他。 \end{tabularx} \\ \\ \relax
6:21 & \begin{tabularx}{0.7\textwidth}{X} 有一天，恰巧是希律的生日，希律擺設宴席，請了大臣、千夫長和加利利的領袖。 \end{tabularx} \\ \\ \relax
6:22 & \begin{tabularx}{0.7\textwidth}{X} 他的女兒希羅底進來跳舞，使希律和同席的人都很高興。王就對女孩說：「無論你要甚麼，向我求，我都會給你」； \end{tabularx} \\ \\ \relax
6:23 & \begin{tabularx}{0.7\textwidth}{X} 又對她多次起誓說：「無論你向我求甚麼，就是我國家的一半，我也會給你。」 \end{tabularx} \\ \\ \relax
6:24 & \begin{tabularx}{0.7\textwidth}{X} 她就出去對她母親說：「我該求甚麼呢？」她母親說：「施洗約翰的頭。」 \end{tabularx} \\ \\ \relax
6:25 & \begin{tabularx}{0.7\textwidth}{X} 她就急忙進去見王，求他說：「我願王立刻把施洗約翰的頭放在盤子裡給我。」 \end{tabularx} \\ \\ \relax
6:26 & \begin{tabularx}{0.7\textwidth}{X} 王就很憂愁，然而因他所發的誓，又因同席的人，不願食言， \end{tabularx} \\ \\ \relax
6:27 & \begin{tabularx}{0.7\textwidth}{X} 就立刻派一個衛兵，吩咐拿約翰的頭來。衛兵就去，在監獄裡斬了約翰， \end{tabularx} \\ \\ \relax
6:28 & \begin{tabularx}{0.7\textwidth}{X} 把頭放在盤子裡，拿來給那女孩，她就給她母親。 \end{tabularx} \\ \\ \relax
6:29 & \begin{tabularx}{0.7\textwidth}{X} 約翰的門徒聽到了，就來把他的屍體領去，放在墳墓裡。 \end{tabularx} \\ \\
[1ex]
\hline
\hline
\end{longtable}
講題:從迦密山回望世界　The Legacy of Mount Carmel

經文:列王記上18:17-22;馬可福音6:14-29

講員:梁元生院長(香港中文大學文學院院長；崇基學院前任院長)

場合:香港中文大學崇基學院禮拜堂

日期:2018年7月15日

一．前言:噢,是“以利亞”!

今天早上我們讀的經文,取自新約馬可福音六章14-29節.在那裡我們看到一幕令人驚悸和震撼,以及充滿戲劇性的場面,.首先,從經文的記述,我們看到一個正在舉行的華麗而熱鬧的宮廷餐舞會.希律王在宮殿中招待賓客,衣香繽影,酒池肉林,來賓有羅馬將領和本地士官, 正是冠蓋雲集,高朋滿座,場面極其熱鬧.席過片時,忽然閒大家屏息靜氣,定睛看著一位美少女登場,她婀娜多姿,隨歌起舞,艷驚四座.主人希律王也全然陶醉,舞蹈表演之後和少女傾談幾句,然後,宴會繼續舉行.但故事還沒有結束,高潮卻在這個時候出現了!一個王的手下官員由宮廷的側門走進來,手裡托著一個大銀盤,像上菜時的樣子,銀盤上面罩著蓋子.大家再次靜默,鴉雀無聲.就在此刻,希律王下令打開蓋子,嘩啦啦,老天爺,這次可不得了!原來銀盤上面載著的並不是什麼山珍海錯或美味佳肴,而是一顆人頭,一枚血淋林的首級.頸上鮮血未乾,是剛砍下來的人頭.眾人目瞪口呆,那是誰?但有人不期然地叫了出來:噢,是以利亞!

我們熟悉這段聖經的人都知道:被殺的人並非以利亞,而是施洗約翰.他因為不滿希律王娶弟婦希羅底為后,得罪了這個婦人,暗暗懷恨於心,跳舞的少女是希羅底的女兒,因跳舞而得著王的歡心,答應為她做任何一件事.於是希羅底就吩咐女兒向王要施洗約翰的頭.這就是這段經文和這個故事的背景.這是個極具戲劇性和令人震撼的故事,但不是我今天要講的主題.我今天要說的是另一個也極具戲劇性及令人震撼的故事.那是發生在施洗約翰被殺之前將近一千年的故事.那故事的主角也出現在馬可福音六章的這段論述之中.

我想要大家注意的,就是這段關於希律王宴會的經文之前,群眾提到的一個人:這個個人就是以利亞.他的名字不斷地重覆地出現在以色列人的歷史上.每當以色列的群眾看到神跡,或者令人震撼的事,包括銀盤子裡的血人頭,總會不假思索的叫出「以利亞」這個名字.為什麼是以利亞?而不是別的人呢?我們要知道:以利亞雖然是猶太人歷史上一位著名的先知,但他畢竟是一個生活於九百多年前舊約時代的人!為甚麼以色列人遇到大事時總記起他,呼喚著他的名字?

明顯可見,以利亞是以色列人的集體記憶中一個極為重要的核心人物, 這個記憶經歷千年而不斷,不止毫不褪色,反而隨時浮現,左右著以色列人的思想,牽動著他們的情緒,影響著他們的生活和行為.在各種不同的場合裡,他是猶太人首先想起的人；例如,在耶穌問門徒「我是誰?」時,他們認為耶穌是甚麼人呢,他們就立刻想到以利亞和摩西；在耶穌登山變像的時候,門徒見到的也是以利亞和摩西.由此可見,在門徒的心目中,在猶太人的集體記憶裡,在絕大多數人的想象中,以利亞佔據著非常重要的地位.我們會問:一個舊約時期的人,一個在過去歷史上出現過的人,憑著甚麼功績或做過什麼事情,叫猶太人長久地記住他及懷念他?換言之,怎樣才可以在民族歷史上留下印記?更重要的是,怎樣才可以在信仰傳統上留下一席位?

今天就讓我們從聖經上看一看這個猶太人歷史和集體記憶中最為重要的人物-以利亞,檢視一下他這個人和他的世界.

二．以利亞和他的時代

如果各位記得,署理校牧林豪恩先生在前一個月主日崇拜的講道中已經介紹過以利亞這位舊約的先知.所以,今天在此我只簡單講述他的生平事跡.願意了解更多關於以利亞生平既其所行的神跡奇事,可以閱讀列王記上一書.簡言之,以利亞是舊約時代的先知,生於主前(公元前)第九世紀.在他以前,猶太人曾經有過一段光輝歲月,在大衛和所羅門作王的時候,各族臣服,百幫來朝,堪稱盛世.但在所羅門王後期,國勢逐漸由盛轉衰,其後眾子爭權,國土分為南北兩國,北國叫以色列,南國則叫猶大.兩國都沒有堅守祖訓,以耶和華為信仰的對象,反之, 他們沾染了外邦各族拜偶像之風,叫上帝厭惡.先知不斷告誡,但百姓仍行惡如儀.以利亞所處時代,就是眾民離棄上帝,事奉偶像,爭權奪利,社會分化,貧富懸殊的時代.舊約列王記上記載了許多關於以利亞的事跡,今天我們只拿出兩樁事情來討論和反省,第一件事是以利亞在迦密山和巴力先知的爭戰,第二件事是以利亞如何幫助陷於飢餓與死亡邊沿的孤兒寡婦.

首先,以利亞的時代是個信仰迷失,心靈迷惘的時代.當前的世界,當前的中國,當前的香港,同樣是秩序失衡,信念迷失,偶像充斥的社會.迦密山的爭戰代表著僅存的信徒(以利亞)以信心和原則對抗物欲橫流的洪潮和偶像泛濫的社會.以利亞代表著滾滾濁流中的一股清泉,是萬眾隨波逐流時的唯一堅持.當時,亞哈在位為王,皇后是耶洗別,信奉偶像巴力；她對猶太人諸多逼害,要所有的人都背棄耶和華,改信巴力.故此,許多猶太先知均被殺害,而拜巴力的人數則大大增加.對以利亞來說,他所信奉的耶和華上帝已經沒有許多的追隨者,他所面對的幾乎全是信仰迷失的群眾,異教風行,偶像處處.同行的人越來越少,堅守原則的人越來越罕有.信仰迷失的時代,也是自覺孤獨的時代.既有信仰覆亡的危機意識,又有孤單獨行的滄桑感覺.全世界都好像離開了耶和華,倒向巴力.列王記上第十八章的描寫,就是他孤身一人與巴力的數百先知在迦密山比拼的畫面.耶和華的先知只剩下以利亞一人,而巴力的先知則有八百五十人.這是一個以勢欺人,以眾凌寡的場面,也是那個時代的真實寫照.八百五十巴力先知,對一個人的以利亞.這是不公平的角力,不公平的競爭.每一個巴力先知,後面可以代表著不同的偶像,除了異教,還可以是錢財,權力,房產,成績,體面,虛榮感等等所造成的偶像.在這場不公平的競爭中,不要說亞哈這一邊的先知人數眾多,他還有數以千計的軍兵為後盾,軍隊,警察,威權和武力,皆是亞哈的支持者.換言之,以利亞所面對的,是宗教與皇權聯合的力量.這種巨大的力量,要借著競爭和比拼之名,要把耶和華僅存的力量和寂寞的聲音吞噬.這是以利亞時的社會:一個以勢欺人,以眾凌寡的社會.或許我們會說:這是古代時期的君權專制社會,現在已經一去不復了.真的是這樣嗎?

讀經至此,使我有無限的感慨.因為我們當今所處的社會,也是一個以眾凌寡,以勢欺人的社會.在位而有權柄的以權謀私,鼓動民意的實則上是以眾凌寡,而立例和民意,大家都以競爭作為比拼的準則.在大學裡也是如此,甚麼都以“競爭”為標準(Competition-based)去衡量表現,學問,研究和做人.教學要看學期末學生對教師的評分,而這評分有要拿來和其他老師,甚至不同院系教師來做比較；做研究也如此,你的文章有多少人引述,刊登你文章的期刊的學術分量有多重?大學排名,甚至你的科系在世界上排名第幾,都是今日教育界普遍應用的指標.在小學,甚至在幼兒園,也不能逃脫競爭及其造成的壓力.父母千方百計要在競爭中取得優勢,只有靠著各種辦法去搬動起跑線,讓孩子得著先機.這個競爭和比拼的文化,在以利亞時代如此,至今仍是一樣.以競爭為名,鬥個你死我活.

我們回顧以利亞的時代,看迦密山上的大比拼和大角力,看到一場驚天動地的鬥爭:包括行雷閃電,天火焚祭,巴力先知被擊殺,以及後來的天降大雨,水淹大地.這是一場驚天地,泣鬼神的鬥爭,而以以利亞的勝利終場.今天我們也許沒有機會看到一幅驚天動地的比拼場面,然而,隱伏在平和與微笑下的殺機和壓力,絕對不少於三千年前的迦密山時代.千年的比拼與角力,其實並未止息.

以利亞在迦密山與巴力的四百五十位先知對壘,凸顯了當時社會注重競賽和比拼的行為.而這場比拼(confrontation)與競爭(competition)背後的精神, 正正是現代社會的標準和現代教育的特徵.換言之,今日的社會同樣以競爭和比拼作為衡量成功的標尺,在競爭之中,同樣也會動員大眾,以眾凌寡,以強凌弱.而不同的是,在古代和現代的競爭中,有很不同的結局. 在舊約歷史中的這場競爭中,以利亞祈求上帝,降下天火,燒盡燔祭,以一人之力,打敗巴力數百先知和眾多的軍兵.這個迦密山的故事藉著以利亞的性格講出做人的三種特質:一是信心--就是對上帝的忠誠和依賴；二是勇氣--以一敵百,堅持原則和不隨波逐流的勇氣；三是能力--擁有一種從上帝而來的超凡力量.從以利亞身上我們能夠學會的三種功課是:有信心,有原則,和有能力.但是,在今天的競爭中,比較的是財富,是資源,是民意,是權力,而非向上帝的信心,道德的原則,和非凡的能力.

以利亞的時代,除了激烈的競爭與比拼,也是一個貧富懸殊和極之不公平的社會.以利亞從迦密山下來,所看到的不是錦衣美食,而是絕望和貧窮.聖經記載著他和窮寡婦的故事.這個寡婦死了丈夫,兒子也得病頻臨於死.家徒四壁,無食物以充飢,無生油以製餅, 無立錐睡覺之地.這時以利亞竟然來求她製餅療饑.她把家裡剩下最後一點麵酵也拿出來,那知越做越有；她也把瓶子裡最後的一點油拿出來,那知油越倒越有,流之不盡.這和迦密山的比拼一樣,也是個以少變多的神跡.我想藉這個“倒不盡的油”的故事來問一下各位,世界上真有這種以少變多,越用越有的事情嗎?你有一百塊錢,分給人十塊後,你還剩下多少?九十塊,對了.再分給人二十,你還剩下多少?物質,金錢,都不會源源無盡,終會有完結的時刻.你身上有甚麼東西是拿出去卻用不完的呢?甚麼才是以利亞的油呢?甚麼東西是倒出來卻流不盡的呢?我告訴你,你身上也有以利亞的油,是流不盡的,那就是愛心和同情心--愛護弱者,憐憫貧窮的人和關愛有需要的家庭和小眾.你付出的愛心和憐憫,是不會用盡的,反而是越用越多,無窮無盡的.這個神跡,後來在他的徒弟以利沙身上又再出現,救了另一個貧窮家庭和孤兒寡婦.(見列王記下)原來,愛心,憐憫,施與和救助等迦密山的賦予,是善行,也是神跡,但都是可以重覆實施的(duplicable).縱然是三千年前以利亞時代的神跡,在二十一世紀的今天,這種愛心的施與,無私的救濟,靠神而有的善心和神跡,仍然是會繼續不斷地湧現,為寡婦,為孤兒,為堅守信仰的人,為敬畏上帝的人,提供一曲“屋頂上的鳴奏曲”,借用數周前鄧瑞強博士的講題和電影錦繡良緣的主題曲.

三．從迦密山回望世界

今天我們探討和觀察以利亞的時代和他所處的世界:一個以眾凌寡,以勢欺人的世界,一個社會絕望,貧病交迫的世界.九百年後,我們聽到有人再次呼叫以利亞時,不再是亞哈王當政,不再是耶洗別為后,卻是羅馬帝國時代,是希律和希羅底勢力控制的時代,但仍然是權力專制,社會上充滿失望和絕望的時代,正因如此,才會有血淋淋的施洗約翰人頭放在銀盤上的情形出現；如是者又過了二千年,直到今天, 我們以為再也不會有耶洗別和希羅底了吧?!錯了,我們環顧世界,信奉巴力的到處都有,以鮮血饅頭置於餐盤上的還是大有人在.從九世紀到一世紀,又從一世紀到二十一世紀,歷史經歷了千年的變化.但人性和這個世界改變了多少?貧窮和絕望還是這個世界和這個社會中的主要旋律,我們仍然在冀求一個以少勝多,以弱制強的神跡,我們仍然在渴望油瓶裡有倒不盡的油,和用不完的麵酵.從迦密山回望這個世界,我們不期然地會用以利沙的話說:願那加在以利亞身上的靈,加倍的在我身上!



\newpage

\section{不足夠}
\label{sec:418}
\textbf{林豪恩先生 - 不足夠}
\newline
\newline
\hyperref[sec:417]{< < < PREV SERMON < < <}
~
\hyperref[sec:toc]{[返目錄]}
~
\hyperref[sec:419]{> > > NEXT SERMON > > >}


列王記下 4:42-44
\newline
\begin{longtable}{cl}
\hline
\hline
章節 & 經文 (和合本修訂版)\\
\hline
4:42 & \begin{tabularx}{0.7\textwidth}{X} 有一個人從巴力‧沙利沙來，帶著初熟果子的食物、二十個大麥做的餅和新麥穗，裝在袋子裡送給神人。神人說：「把這些給大家吃。」 \end{tabularx} \\ \\ \relax
4:43 & \begin{tabularx}{0.7\textwidth}{X} 僕人說：「這些豈可擺在一百人面前呢？」以利沙說：「你只管給大家吃吧！因為耶和華如此說，他們必吃了，還有剩下的。」 \end{tabularx} \\ \\ \relax
4:44 & \begin{tabularx}{0.7\textwidth}{X} 僕人就擺在他們面前，他們吃了，還有剩下，正如耶和華所說的。 \end{tabularx} \\ \\
[1ex]
\hline
\hline
\end{longtable}
詩篇 145:10-18
\newline
\begin{longtable}{cl}
\hline
\hline
章節 & 經文 (和合本修訂版)\\
\hline
145:10 & \begin{tabularx}{0.7\textwidth}{X} 耶和華啊，你一切所造的都要稱謝你，你的聖民也要稱頌你。 \end{tabularx} \\ \\ \relax
145:11 & \begin{tabularx}{0.7\textwidth}{X} 他們要傳講你國度的榮耀，談論你的大能， \end{tabularx} \\ \\ \relax
145:12 & \begin{tabularx}{0.7\textwidth}{X} 好讓世人知道你大能的作為和你國度威嚴的榮耀。 \end{tabularx} \\ \\ \relax
145:13 & \begin{tabularx}{0.7\textwidth}{X} 你的國是永遠的國！你執掌的權柄存到萬代！耶和華一切的話信實可靠，他一切的作為都有慈愛。 \end{tabularx} \\ \\ \relax
145:14 & \begin{tabularx}{0.7\textwidth}{X} 耶和華扶起所有跌倒的，扶起所有被壓下的。 \end{tabularx} \\ \\ \relax
145:15 & \begin{tabularx}{0.7\textwidth}{X} 萬有的眼目都仰望你，你按時給他們食物。 \end{tabularx} \\ \\ \relax
145:16 & \begin{tabularx}{0.7\textwidth}{X} 你張手，使一切有生命的都隨願飽足。 \end{tabularx} \\ \\ \relax
145:17 & \begin{tabularx}{0.7\textwidth}{X} 耶和華一切所行的，無不公義，一切所做的，都有慈愛。 \end{tabularx} \\ \\ \relax
145:18 & \begin{tabularx}{0.7\textwidth}{X} 耶和華臨近凡求告他的，臨近所有誠心求告他的人。 \end{tabularx} \\ \\
[1ex]
\hline
\hline
\end{longtable}
以弗所書 3:14-21
\newline
\begin{longtable}{cl}
\hline
\hline
章節 & 經文 (和合本修訂版)\\
\hline
3:14 & \begin{tabularx}{0.7\textwidth}{X} 因此，我在父面前屈膝— \end{tabularx} \\ \\ \relax
3:15 & \begin{tabularx}{0.7\textwidth}{X} 天上地上的各家都是從他得名的－ \end{tabularx} \\ \\ \relax
3:16 & \begin{tabularx}{0.7\textwidth}{X} 為要他按著他豐盛的榮耀，藉著他的靈，使你們內心的力量剛強起來； \end{tabularx} \\ \\ \relax
3:17 & \begin{tabularx}{0.7\textwidth}{X} 又要他使基督因著你們的信住在你們心裡，使你們既在愛中生根立基， \end{tabularx} \\ \\ \relax
3:18 & \begin{tabularx}{0.7\textwidth}{X} 能夠和眾聖徒一同明白基督的愛是何等的長、闊、高、深， \end{tabularx} \\ \\ \relax
3:19 & \begin{tabularx}{0.7\textwidth}{X} 並知道這愛是超過人的知識所能測度的，為要使你們充滿神一切的豐盛。 \end{tabularx} \\ \\ \relax
3:20 & \begin{tabularx}{0.7\textwidth}{X} 神能照著運行在我們心裡的大能充充足足地成就一切，超過我們所求所想的。 \end{tabularx} \\ \\ \relax
3:21 & \begin{tabularx}{0.7\textwidth}{X} 願他在教會中，並在基督耶穌裡，得著榮耀，直到世世代代，永永遠遠。阿們！ \end{tabularx} \\ \\
[1ex]
\hline
\hline
\end{longtable}
約翰福音 6:1-15
\newline
\begin{longtable}{cl}
\hline
\hline
章節 & 經文 (和合本修訂版)\\
\hline
6:1 & \begin{tabularx}{0.7\textwidth}{X} 這些事以後，耶穌渡過加利利海，就是提比哩亞海。 \end{tabularx} \\ \\ \relax
6:2 & \begin{tabularx}{0.7\textwidth}{X} 有一大群人因為看見他在病人身上所行的神蹟，就跟隨他。 \end{tabularx} \\ \\ \relax
6:3 & \begin{tabularx}{0.7\textwidth}{X} 耶穌上了山，和門徒一同坐在那裡。 \end{tabularx} \\ \\ \relax
6:4 & \begin{tabularx}{0.7\textwidth}{X} 那時猶太人的逾越節近了。 \end{tabularx} \\ \\ \relax
6:5 & \begin{tabularx}{0.7\textwidth}{X} 耶穌舉目看見一大群人來，就對腓力說：「我們到哪裡去買餅給這些人吃呢？」 \end{tabularx} \\ \\ \relax
6:6 & \begin{tabularx}{0.7\textwidth}{X} 他說這話是要考驗腓力，他自己原知道要怎樣做。 \end{tabularx} \\ \\ \relax
6:7 & \begin{tabularx}{0.7\textwidth}{X} 腓力回答他：「就是兩百個銀幣的餅也不夠給他們每人吃一點點。」 \end{tabularx} \\ \\ \relax
6:8 & \begin{tabularx}{0.7\textwidth}{X} 有一個門徒，就是西門‧彼得的弟弟安得烈，對耶穌說： \end{tabularx} \\ \\ \relax
6:9 & \begin{tabularx}{0.7\textwidth}{X} 「這裡有一個孩子，帶著五個大麥餅和兩條魚，但是分給這麼多人還算甚麼呢？」 \end{tabularx} \\ \\ \relax
6:10 & \begin{tabularx}{0.7\textwidth}{X} 耶穌說：「你們叫大家坐下。」那地方的草多，人們就坐下，男人的數目約有五千。 \end{tabularx} \\ \\ \relax
6:11 & \begin{tabularx}{0.7\textwidth}{X} 耶穌拿起餅來，祝謝了，就分給坐著的人，也同樣分了魚，都照他們所要的來分。 \end{tabularx} \\ \\ \relax
6:12 & \begin{tabularx}{0.7\textwidth}{X} 他們吃飽後，耶穌對門徒說：「把剩下的碎屑收拾起來，免得糟蹋了。」 \end{tabularx} \\ \\ \relax
6:13 & \begin{tabularx}{0.7\textwidth}{X} 他們就把那五個大麥餅的碎屑，就是大家吃剩的，收拾起來，裝滿了十二個籃子。 \end{tabularx} \\ \\ \relax
6:14 & \begin{tabularx}{0.7\textwidth}{X} 人們看見耶穌所行的神蹟，就說：「這真是那要到世上來的先知！」 \end{tabularx} \\ \\ \relax
6:15 & \begin{tabularx}{0.7\textwidth}{X} 耶穌知道他們要來強迫他作王，就獨自又退到山上去了。 \end{tabularx} \\ \\
[1ex]
\hline
\hline
\end{longtable}
講題:不足夠 　Not Enough

經課:列王記下4:42-44;詩篇145:10-18;以弗所書3:14-21;約翰福音6:1-15

講員:林豪恩先生

場合:香港中文大學崇基學院禮拜堂

日期:2018年7月29日

一．引言

有一個問題,直接問自己可能會感到焦慮,所以轉為代其他人回答:「你認為香港即將退休的普羅大眾儲夠錢過退休後的幾十年沒有呢?」

有一間資產管理公司(施羅德)於2017年12月發表一項調查顯示,本港95\%已退休投資者後悔當初未為退休積累更多資金,調查亦發現本港61\%受訪未退休投資者擔憂退休後的收入不足以提供舒適生活.(文匯報2017/12/13)

2014年5月,有記者訪問在中大教財務學的老師李兆波先生,「請問幾多錢至夠退休呢?」李先生答:「以60歲退休後仍要生活25年,如要求中產生活,退休前至少需1,500萬元儲蓄.」(晴報2014/5/8)

無論是1500萬,還是500萬,對一般「打工仔」來說都是高不可攀的「天文數字」.香港人這種不足夠的狀況和感覺,比環球及亞洲區高出約30\%.所以,香港人常常感到沉重的壓力,以及揮之不去的焦慮是有原因的；香港人那麼著急「搵錢」,也是有背景的.

社會的氣氛,以及我們個人的感覺,都常常警惕著我們「不足夠」的狀況,例如:退休金不足夠,醫療保險不足夠,房屋供應不足夠,年輕人競爭力不足夠......

「不足夠」這種狀況和感覺,並非現代產物,古已有之.今日誦讀的舊約和福音經課,都提及「不足夠」這現象.這些古時「不足夠」的故事,對今時今日的我們有甚麼啟發呢?

二．內容

1．兩個不足夠的故事

第一個「不足夠」的故事記載在列王記下4:42-44:

故事的大背景在4:38「以利沙回到吉甲,那地正有饑荒.」

飢荒,是指一個國家或地區在某一個時期內沒有足夠的糧食供應,導致局部或全域性的人們逃難,或因飢餓導致的佔人口比例較大的非正常死亡的現象.令人困惑的是,現代科技比幾千年前進步得多,生產力也比幾千年前強大得多,但飢荒仍然出現.非謀利組織ACAPS的「人道概況:2018重大危機分析」指出,埃塞俄比亞將與尼日利亞東北部,索馬里,南蘇丹和也門等地區,面臨飢荒風險.

4:42有一個人從巴力‧沙利沙來,帶著初熟果子的食物,二十個大麥做的餅和新麥穗,裝在袋子裡送給神人.神人說:「把這些給大家吃.」

4:43僕人說:「這些豈可擺在一百人面前呢?」

在不足夠的大背境中,小量的供應,只是杯水車薪,感覺是無濟於事,不能解決問題,也不能解除不足夠的焦慮.

第二個「不足夠」的故事記載在約翰福音6:1-15:

6:5耶穌舉目看見一大群人來,就對腓力說:「我們到哪裡去買餅給這些人吃呢?」

6:7腓力回答他:「就是兩百個銀幣的餅也不夠給他們每人吃一點點.」

其實,二百銀幣並非小數目,大概是當時一個工人八個月的薪金.不過,腓力指出用這款項購買到的餅也不夠給他們每人吃一點點,可想而之現場人數多,需要的巨大,更何況他們上了山,哪裡可以買到那麼多餅呢?

6:8有一個門徒,就是西門‧彼得的弟弟安得烈,對耶穌說:

6:9「這裡有一個孩子,帶著五個大麥餅和兩條魚,但是分給這麼多人還算甚麼呢?」

安得烈說得有道理,「就是兩百個銀幣的餅也不夠給他們每人吃一點點」,五個大麥餅和兩條魚,「分給這麼多人還算甚麼呢」?

不足夠是事實.對著不足夠的狀況,有些回應是「不作為」,因為他們認為杯水車薪,做了也解決不了問題:好像以利沙故事中的僕人:「這些豈可擺在一百人面前呢?」又例如五餅二魚故事中的腓力:「就是兩百個銀幣的餅也不夠給他們每人吃一點點.」及安得烈:「但是分給這麼多人還算甚麼呢?」

幾年前,有一個香港人名叫余偉業,創辦了社企「要有光」,邀請業主以低於市值的租金把空置單位租給他們,他們再分租給有逼切居住需要的市民.當時,也有回應說:「這樣的計劃能夠提供到多少個單位呢?對於解決香港的居住問題有甚麼作用呢?」

「不足夠」的感覺的確會令人感到無能為力,做了也不能解決問題.這種無力感令人無奈地「不作為」.然而,有需要的人仍然存在,如何是好呢?其他三卷福音書都記載了門徒的回應:「這地方偏僻,而且時候已經晚了,請叫眾人散去,他們好進村子,自己買些食物.」(馬太福音14:15)相隔幾千年,門徒的想法和此時此地我們香港人的想法很接近:「自己的需要自己滿足,自己的問題自己解決」,「有錢的住豪宅,幸運的住公屋,無錢的住tong 房,走頭無路的去訓街,各人都吃自己種的果」,「如果沒有條件在香港生活,就不要留在這裡.」物競天擇,適者生存,香港,只適合有足夠競爭力的人存在.在這種目光之下,香港沒有窮人,只有懶人.

以利亞卻仍然看見人們的需要:4:43他對僕人說:「你只管給大家吃吧!」耶穌也並不以為人們聽了他講的道就等於吃飽了,他仍然肯定人們對食物的需要:三卷福音書都記載了他對門徒的吩咐:「你們給他們吃吧!」(馬太福音14:16)面對不足夠,有些回應是否定需求,有些回應是壓抑需求,不過,有些基本的需求不能被壓抑,也不應被壓抑.《世界人權宣言》第25條第一部份如此寫著:「人人有權享受為維持他本人和家屬的健康和福利所需的生活水準,包括食物,衣著,住房,醫療和必要的社會服務；在遭到失業,疾病,殘廢,守寡,衰老或在其他不能控制的情況下喪失謀生能力時,有權享受保障.」

2．兩個足夠的故事

作為讀者,大家已經「劇透」了這兩個「不足夠」的故事,知道兩個故事的結局都變成「足夠」.

第一個足夠的故事是如何發生的呢?

4:42有一個人從巴力‧沙利沙來,帶著初熟果子的食物,二十個大麥做的餅和新麥穗,裝在袋子裡送給神人.神人說:「把這些給大家吃.」

因為有一個人在不足夠的狀況之下,仍記得為地上的出產,為所得的收成而感恩,按著他們的信仰傳統獻祭.

另一個足夠的故事是如何發生的呢?

6:9「這裡有一個孩子,帶著五個大麥餅和兩條魚,但是分給這麼多人還算甚麼呢?」

6:10耶穌說:「你們叫大家坐下.」那地方的草多,人們就坐下,男人的數目約有五千.

6:11耶穌拿起餅來,祝謝了,就分給坐著的人,也同樣分了魚,都照他們所要的來分.

因為有一個孩子,不知道他是否太過天真,把自己的五個餅和兩條魚拿出來；還是好像在小學課室一樣,被旁邊的同學作弄,把他的五個餅和兩條魚拿出來.

如果是兒童主日學,故事講到這裡就可以總結了.老師問:「各位同學,這兩個故事教訓我們甚麼呢?」同學們都會回答標準答案:「拿出來與人分享.」老師就會回應:「答對了!上主甚麼都有,在祂沒有不足夠,祂能夠滿足我們的需要,所以我們不用憂慮,只管把自己有的拿出來與人分享.」沒錯,上主是豐足的,今天的經課中詩篇及書信也如此回應:

詩篇145:10-18

詩人以雛鳥等待母鳥餵哺,比喻上主供應的足夠:「145:15萬有的眼目都仰望你,你按時給他們食物.145:16你張手,使一切有生命的都隨願飽足.」

以弗所書3:14-21

3:18-19能夠和眾聖徒一同明白基督的愛是何等的長,闊,高,深,並知道這愛是超過人的知識所能測度的,為要使你們充滿上帝一切的豐盛.

3:20上帝能照著運行在我們心裡的大能充充足足地成就一切,超過我們所求所想的.

3．把不足夠變為足夠

我們沒有懷疑上帝是豐足的,但不足夠也是我們的真實經驗和切身感受,兩者之間好像存在著巨大洪溝,難以跨越.

好像那位從巴力‧沙利沙來的虔誠信徒一樣,把初熟果子拿來獻祭,或者如那位被傳頌的小孩子拿五個大麥餅和兩條魚出來,就能夠解決我們所面對的不足夠的問題了嗎?現實地想想,可能大家心裡都有疑惑.

去年特首選舉前,其中一位參選人向一名行乞老婦捐贈500元,引起廣泛輿論.當事人重申自己出於善心及憐憫才給錢婆婆,呼籲人們不應用來做政治炒作.當事人對此事的感觀和市民的感觀可能存在相當的距離,當事人視之為個人行為,市民視之為處事者的公共行為.市民心裡有感受:捐贈500元給乞丐算是甚麼?我們當中也有很多人有需要,政府如何回應我們的需要呢?

相信社會上不少人都如這位特首參選人一樣有某程度的善心及憐憫,在某程度上都願意與人分享自己所擁有的.然而,在付出的時候,難免心裡會問:我所交出的會否所託非人呢?他們如何運用我提供的資源呢?我如此做真的能夠對事情有幫助嗎?

今年的財政預算案是否「派錢」又成為激烈的討論,有不少人認為香港社會既然在醫療,福利,教育,房屋等各公共領域都有強大的需要,為甚麼仍要「派錢」呢?為甚麼不投放更多資源在有需要的地方呢?從前香港人嘲笑鄰近的政府「Hea得只懂派錢」,今時今日,香港人也群起要求政府派錢.為甚麼呢?有些人不是否定社會的需要,而是不信任政府能夠把資源運用得好,擔心自己所託非人.

這兩個不足夠變為足夠的故事,其實不一定是常態,當中有各種條件的配合.讀者常常留意到的,是有人願意奉獻自己僅有的,但不常留意的是接收奉獻者是誰及如何處理.如果以利沙收到那人獻上初熟的果子,自己一口吃完,或者藏為己用,這故事的結局就不同了；如果耶穌收到五餅二魚,剛足夠自己的晚餐,或者與十二位門徒每人吃一點,故事的結局也就不同了.會有這樣的事情發生嗎?當然有,而且多的是.在某些地方,官員的財產比他的薪酬多千萬倍,如此龐大的財產從那裡來的呢?會不會當中不少本來是公共資源呢?在一些發生天災的地方,捐款有多少真正用於救災及重建呢?又有多少落入處事者的手中呢?

願意付出的人,希望把所付出的交給以利沙和耶穌這些值得信任的人,今天,如以利沙和耶穌一樣值得信任的人在哪裡呢?如果沒有這樣的人或組織或政府或運作系統,為了應付不足夠而增加的供應,也未必能夠幫助有需要的人.例如金融危機時,各國政府大印銀紙,銀紙去了有錢人的的戶口中,還是貧窮者的手中呢?

有餘的人,或者僅有的人,或者並非不願意分享,他們可能希望把資源交給如以利沙和耶穌這些能夠運用他們所交出的去解決大家的問題,滿足大家的需要的人.今天,如以利沙和耶穌這些能夠運用資源去滿足人們需要的人或組織或政府或運作系統在哪裡呢?

現在仍要等待神蹟嗎?讓我們試試以古代人的眼睛來看看,現代世界處處都是神蹟.運用農業技術而能夠大量生產糧食是奇蹟,工廠運用機械大量製造產品是奇蹟,飛機能夠把人及貨物短時間送到地球上每一個角落也是奇蹟.現今是一個生產力非常強大的世界,奇怪的是這個世界仍有不少人受飢荒之苦；現今是一個科技非常先進的世界,奇怪的是這個世界仍有嬰兒飲營養不足的假奶粉及注射不合格的疫苗；現在是一個物質非常豐富的世界,奇怪的是在這個地界不少人仍處於不足夠的焦慮中.以食物為例,現在我們的世界可能並非處於不足夠的狀況,問題是我們如何處理食物.

2018年7月6日是一個令人屏息以待的日子,不少人在關心中美貿易戰會否開打.運載美國大豆到中國的散裝貨輪「飛馬峰號」於美國西雅圖時間6月18日晚上7時許開出,預計將在本港時間7月6日下午1時抵達.該艘設計時速12.5海里的貨船,以「極速」的14海里橫過黃海,直奔大連.飛馬峰號最後在當天下午5時07分抵達大連港外海,最終趕不及在加徵25\%關稅前到達.網上消息指,賣家因買方取消訂單,已將大豆傾倒入大海.這批大豆有多少呢?是8萬噸,估計價值1.5億元人民幣.(東網2018/7/8)這不是個別事件,根據英國時尚精品Burberry去年的報告,這家奢華精品燒毀價值高達2860萬英鎊的庫存衣飾及香水,並解說此舉在時裝業相當常見.(法新社2018/7/20)

四．總結

讓我們想像一下,今天,足夠與不足夠,會否未必再是產量的問題,而可能是價值和操作的問題呢?又讓我們想像一下,人類堅持個體獨自解決不足夠的焦慮,會否令大家繼續陷於不足夠的焦慮之中呢?如果群體一齊面對不足夠的焦慮,會否有可能大家都足夠,而不用焦慮呢?有沒有類似的例子呢?

全球富豪排行榜由2000年開始至今,首十位沒有一位是北歐人；目前世界上最強大的跨國企業也不在北歐,但在眾多幸福調查中,北歐國家總是名列前茅.例如在丹麥,稅收很高,高收入人群的稅率可能高達51.5\%.但是,收上來的稅錢通過一系列社會計劃重新投入社會,例如免費大學教育,免費醫療,優厚的產假待遇和失業救濟等.哥本哈根智庫「幸福研究院」的負責人Meik Wiking說:「我們不是在交稅,我們是在投資社會,我們是在購買生活質量.」(2017年 3月 28日BBC中文網)

Silvia是一位土生土長的香港女子,嫁了去芬蘭,無車,無樓,無穩定工作,婚禮簡單地進行.她說:「在芬蘭,福利保障做得很好,即使失業,都不會餓死和無屋住.既然如此,錢就沒那麼重要,快樂與錢也沒有那麼相關.」(2017年12月03蘋果日報)

如何把不足夠的狀況改變為足夠呢?或者,所需要的不再是壓榨大地再加倍出產,不是勞役人類日以繼夜地工作,而是建立一個共享的社群.如何能夠建立共享的社群呢?人對其他人的信任不可少,每個人的價值都被重視,人與人之間的尊重不可少,人們都有機會去參與及建構一個能被信任的公共權力.信任,生命價值,平等,尊重,公平參與等這些價值,基督教傳統中有不少資源可供參考,有待我們去發掘及實踐.



\newpage

\section{作上帝蒙慈愛的兒女}
\label{sec:419}
\textbf{戴國璋牧師 - 作上帝蒙慈愛的兒女}
\newline
\newline
\hyperref[sec:418]{< < < PREV SERMON < < <}
~
\hyperref[sec:toc]{[返目錄]}
~
\hyperref[sec:420]{> > > NEXT SERMON > > >}


以弗所書 5:1-7
\newline
\begin{longtable}{cl}
\hline
\hline
章節 & 經文 (和合本修訂版)\\
\hline
5:1 & \begin{tabularx}{0.7\textwidth}{X} 所以，作為蒙慈愛的兒女，你們該效法神。 \end{tabularx} \\ \\ \relax
5:2 & \begin{tabularx}{0.7\textwidth}{X} 要憑愛心行事，正如基督愛我們，為我們捨了自己，當作馨香的供物和祭物獻給神。 \end{tabularx} \\ \\ \relax
5:3 & \begin{tabularx}{0.7\textwidth}{X} 至於淫亂和一切污穢，或是貪婪，在你們中間連提都不可，這才合乎聖徒的體統。 \end{tabularx} \\ \\ \relax
5:4 & \begin{tabularx}{0.7\textwidth}{X} 淫詞、妄語和粗俗的俏皮話都不合宜；總要說感謝的話。 \end{tabularx} \\ \\ \relax
5:5 & \begin{tabularx}{0.7\textwidth}{X} 要確實知道，無論是淫亂的，是污穢的，是貪心的（貪心的就是拜偶像的），在基督和神的國裡都得不到基業。 \end{tabularx} \\ \\ \relax
5:6 & \begin{tabularx}{0.7\textwidth}{X} 不要被人虛浮的話欺騙了，因這些事，神的憤怒必臨到那些悖逆的人。 \end{tabularx} \\ \\ \relax
5:7 & \begin{tabularx}{0.7\textwidth}{X} 所以，不要與他們同夥。 \end{tabularx} \\ \\
[1ex]
\hline
\hline
\end{longtable}
講題:作上帝蒙慈愛的兒女　Be Imitators of God as Behaved Children

經文:以弗所書5章1-7節

講員:戴國璋牧師

場合:香港中文大學崇基學院禮拜堂

日期:2018年8月5日

主題名稱來自選讀經文的首節, 通俗表達為: 「做上帝愛ge 仔女」.

在以弗所書, 保羅的倫理教導開始於第四章.也就是說, 從第四章至第六章是要教導我們做個怎樣的基督徒.為了易於明了, 保羅連串地運用了幾個譬如:「新人與舊人」,「父親與兒女」及「光明和幽暗」.

今天,我們會將焦點集中在第二對.這雖然是譬如, 卻是基於一個神學的事實:「我們果然地是上帝的兒女」.

按照保羅的信念, 聖靈引導我們呼叫上帝作「阿爸」(父啊)!這就是我們兒女身份的明証.我們確信自己是天父上帝的兒女, 也就要活得像個天父的兒女.這就是保羅在此處囑咐的根據所在.

兒女酷肖父母,天經地義

感謝神, 在我們的教會生活中, 上帝常賜新生命給一些家庭.我們最常聽見的一句恭維話:「好可愛,長得像極爸爸／媽媽!」等.這類話其實有點「多餘」.孩子是父母的愛情結晶. 當然長得像自己的父母親.兒女酷肖父母是天經地義之事.兒女不肖父母,反而是離經背道的.

兒女跟父母酷肖,更成為上帝啟示的手段.當始祖亞當被造時,聖經說:「他是按上帝的形像樣式而造他們,所以亞當被稱為上帝的兒子」.當先祖塞特出生時,聖經的用語是:「亞當活到一百三十歲,生了一個兒子,形像樣式和自己相似…. .」(創15:3)暗示往後的人遺傳了始祖的軟弱.當救主耶穌來到人世間,以後從死而復生,顯明是上帝兒子時,聖經就說他是「上帝的像」,指他酷肖上帝,能顯露上帝的真面目.這樣,當我們被稱為「上帝兒女」時,酷肖上帝就是必然的了.所以保羅在此教導說:「你們當酷肖你們的父上帝,成為蒙他喜愛的兒女.

兒女學效父母,自然而然

兒女肖似自己父母,故然根由於遺傳了父母的基因,更因兒女跟父母的親近,耳濡目染,會不經意地把父母的神情,性格,價值觀甚至習慣的動作模仿了而不自知.許多時候,人的評語:「某人像極他父／母親」,指的更多地是指這個方面.所以才有「虎父無犬子」,「龍生龍,鳳生鳳」等的血統論理論.

與其說是血統遺傳達至這種相似,不如說是耳濡目染作成的結果.

保羅說:「你們該效法上帝,好像蒙慈愛的兒女一樣」正是這個意思.

不經意地習染是正常,刻意模仿學效更是應當.

在「登山寶訓」中,耶穌的一個總結:「你們要完全,好像你們的天父完全」正是這個意思.

兒女取悅父母,順理成章

作為上帝兒女,酷肖上帝,天經地義；要蒙上帝寵愛,學效父親的嘉言美行,自然而然.然後,保羅勸眾人要學效耶穌憑愛心而行事.他是怎樣地憑愛心行事呢?乃是「愛我們,為我們捨了自己,當作馨香的供物和祭物獻與上帝」(2節).請注意,耶穌愛我們而行的捨己行為,被描述成馨香的供物祭物,是用作取悅上帝的禮物.耶穌成為你我的一個榜樣:他怎樣取悅父上帝,我們也要怎樣地取悅父上帝.主耶穌是怎樣取悅上帝的呢?用他自己的話:「我常作父所喜悅的事.」(約8:29).

你有否想過你應行事去取悅你所愛的家人,而不是去刺激你的家人,叫他們難受,不安,憤怒甚至失去盼望.

我們都是父上帝的兒女,上帝是我們的父親.這種關係,沒甚事物能破壞.準備好學父親的聖潔,慈憐,以成聖生活取悅他嗎?



\newpage

\section{使徒保羅的使命人生}
\label{sec:420}
\textbf{周宇牧師 - 使徒保羅的使命人生 (講道以普通話進行)}
\newline
\newline
\hyperref[sec:419]{< < < PREV SERMON < < <}
~
\hyperref[sec:toc]{[返目錄]}
~
\hyperref[sec:421]{> > > NEXT SERMON > > >}


腓立比書 3:7-14
\newline
\begin{longtable}{cl}
\hline
\hline
章節 & 經文 (和合本修訂版)\\
\hline
3:7 & \begin{tabularx}{0.7\textwidth}{X} 只是我先前以為對我是有益的，我現在因基督的緣故而當作是有損的。 \end{tabularx} \\ \\ \relax
3:8 & \begin{tabularx}{0.7\textwidth}{X} 不但如此，我已把萬事當作是有損的，因我以認識我主基督耶穌為至寶。我為他已經丟棄萬事，看作糞土，為要贏得基督， \end{tabularx} \\ \\ \relax
3:9 & \begin{tabularx}{0.7\textwidth}{X} 並且得以在他裡面，不是有自己因律法而得的義，而是有信基督的義，就是基於信，從神而來的義， \end{tabularx} \\ \\ \relax
3:10 & \begin{tabularx}{0.7\textwidth}{X} 使我認識基督，知道他復活的大能，並且知道和他一同受苦，效法他的死， \end{tabularx} \\ \\ \relax
3:11 & \begin{tabularx}{0.7\textwidth}{X} 或許我也得以從死人中復活。 \end{tabularx} \\ \\ \relax
3:12 & \begin{tabularx}{0.7\textwidth}{X} 這不是說我已經得著了，已經完全了；而是竭力追求，或許可以得著基督耶穌所要我得著的。 \end{tabularx} \\ \\ \relax
3:13 & \begin{tabularx}{0.7\textwidth}{X} 弟兄們，我不是以為自己已經得著了；我只有一件事，就是忘記背後，努力面前的， \end{tabularx} \\ \\ \relax
3:14 & \begin{tabularx}{0.7\textwidth}{X} 向著標竿直跑，要得神在基督耶穌裡從上面召我來得的獎賞。 \end{tabularx} \\ \\
[1ex]
\hline
\hline
\end{longtable}
講題:使徒保羅的使命人生 (講道以普通話進行)

Apostle Paul: A Life on Mission (conducted in Putonghua)

經文:腓立比書3章7-14節

講員:周宇牧師

場合:香港中文大學崇基學院禮拜堂

日期:2018年8月12日

保羅的個人成長背景:

保羅在腓立比書3:5-6對自己的個人成長背景有一個簡明扼要的介紹:第八天受割禮,是以色列族便雅憫支派的人,是希伯來人所生的希伯來人.就律法來說,是法利賽人,就熱心來說, 是逼迫教會的,就律法的上義來說, 是無可指摘的.

從這樣一個自我介紹來看, 保羅似乎是在形容自己的人生上半場幾乎是完美無缺-almost perfect!

但緊接著保羅話鋒一轉卻說,“只是我先前以為與我有益的,我現在因基督都當作有損的.不但如此,我也將萬事當作有損的,因我以認識我主基督耶穌為至寶.我為他已經丟棄萬事,看作糞土,為要得著基督.”腓3:7-8

讓我們來重溫一下究竟發生了什麼事,導致了保羅生命及職業生涯(career path)的巨變.

保羅的皈信及呼召:

保羅原名掃羅,出生成長在羅馬的(太平)盛世(Pax Romana).雖是猶太人,卻享有羅馬公民身份.他生在羅馬帝國基利家省的大數,長在耶路撒冷城,並在嚴謹的律法師／拉比 迦瑪列門下受教.

按照今天的話說,保羅也算得上是出生名門,就讀名校,前程似錦. He seems to have a secured future and on the right track to a successful career path!

保羅從“法學院”畢業之後也有一個很好的實習機會Internship.有一天在耶路撒冷公會court裡幫助證人看衣服,並親眼目睹“叛教者” 司提反被眾人打死-Actually, he approved of (批准／同意)their killing of Stephen-中文和合本翻譯為“掃羅也喜悅他被害”.

毫無疑問,這次事件對保羅的事業發展意義深遠,他開始變本加厲地“殘害教會,進各人的家,拉著男女下在監裡.由於他的出色表現,保羅迅速從一個實習生被委以重任. Acts. 9: 1-2

在Acts 9章一開頭我們看見,他主動請纓,向大祭司要了文書,去傳達給大馬士革的各會堂,抓捕一切信奉這“道”的人.可就在去大馬士革執行任務的路上,他先被耶穌抓到了(single-handedly apprehended by Jesus)!類似當年上帝在燃燒的荊棘叢中呼召摩西那樣地戲劇化,耶穌這次親自在大光中向保羅顯現,並吩咐他到外邦人和君王並以色列人面前宣揚耶穌的名.It seemed this was a mission impossible but it actually happened. 這一系列劇情發展詳細記錄在Acts 9-28章裡面,前不久也被Sony Picutures 拍攝成為一部大片Paul-Apostle of Christ.

保羅的使命人生:異像,使命,策略

今天給大家的只是一些片段snap shots, 重點聚焦在保羅如何領受異像,承接使命,並很有策略地做到使命必達-Mission Possible! Paul was indeed a visionary, missionary , and great strategist-我們可以從基督教頭一百年的迅猛發展得出這樣一個結論.

與千百年前上帝吩咐摩西帶領以色列民出埃及為奴之地不同,那次是一個史詩般的宏大敘事,這次倒更像是潤物細無聲的內在核裂變.就像耶穌教導門徒的那樣,“一粒麥子不落在地裡死了,仍舊是一粒；如是死了,就結出許多子粒來.”

當保羅為了要得著耶穌而把萬事看作糞土丟棄,以認識耶穌基督為至寶時,他實際上是被耶穌基督得著了.用他自己的話說,“我已經與耶穌同釘十字架,現在活著的不在是我,而是基督在我裡面活著；並且我如今在肉身活著,是因信神的兒子而活,他是愛我, 為我捨己.”

如果你是保羅的家人,或是關心他的同事朋友,看到他的這種狀態,會給他一些什麼忠告或專業意見呢?你也許會說,保羅,你是不是瘋了,你有羅馬公民身份,你有良好的專業訓練,看你的CV,你完全有可能成為下一屆的“大律師公會”主席,究竟發生了什麼事,讓你如此瘋癲,去跟隨你曾經極力逼迫的那道-the Way?!

保羅在腓立比書3:9-14節對他新的選擇和人生目標有進一步的解釋(paraphrase):以前我是靠嚴格遵守律法稱義,而現在呢,是因信基督被稱為義.因為惟有耶穌才能滿足所有律法的要求-Jesus has fulfilled all the requirements of the Law! 我現在是在基督裡面,我渴望更深地認識基督,願意和他一同受苦,通過效法他的死,來經歷他復活的大能.

保羅繼續道,“這不是說我已經得著了,已經完全了,我乃是竭力追求,或者可以得著基督耶穌所以得著我的.弟兄們,我不是以為自己已經得著了,我只有一件事,就是忘記背後,努力面前的,向著標竿直跑,要得神在基督耶穌裡從上面召我來得的獎賞”.

保羅在哥林多後書第5章對自己的“為主癲狂”更是一語道破,Yes! I am crazy for God as you all have labeled me, so what?-“我們若果癲狂,是為神；若果謹守,是為你們.原來基督的愛激勵我們.因我們想,一人既替眾人死,眾人就都死了；並且他替眾人死,是叫那些活著的人不再為自己活,乃為替他們死而復活的主活.“

自從保羅與耶穌有了這樣一次的面對面-living Encounter之後,他得著了那從天上來的異像,便從此開始了他的使命人生,義無反顧,無問西東(totally sold out for Christ)!

摯愛耶穌成為保羅一生宣教的動力,保羅的宣教工作也並沒有結束在使徒行傳28:31,而是未完待續.更準確地說,是一切都剛剛開始……

在接下來的三百年間,基督的跟隨者們把福音傳遍了整個羅馬帝國,甚至君士坦丁大帝也皈依基督,並隨後頒布基督教成為國教.這些驚天動地的事情又是如何發生的?

仔細研究一下保羅的宣教策略和模式就會給我們提供一些線索.Now we move from his being to his doing.

保羅的使徒性團隊及早期教會的複合性網絡Complex Network

有聖經學者通過對使徒行傳和保羅書信的詳細研究發現,保羅總共有69位可以叫出名字的同工,這其中有9位是保羅-提摩太使徒類型的領袖(sodality leaders),53位地方教會同工領袖(modality leaders), 有13位婦女領袖,及16位經濟財務支持者(benefactors).

複合型網絡(Complex Network)是一個當代的科技名詞,物理學家Albert-LaszLo Barabasi 在他的鏈接Linked: Everything is Connected to Everything Else and What It Means for Business, Science, and Everyday Life 一書中以使徒保羅為例,認為他是最懂得複合型網絡並把它用到極致的第一人!類似當今世界各大航空公司都有自己的Regional Strategic Hub,保羅在當時的羅馬帝國也建立了宣教的戰略樞紐站,比如Antioch, Ephesus, and Rome,並在當時的31個重要城市建立網絡型的家庭教會,規模從2-3人到50-70人不等.

另外一位宗教社會學家 Rodney Stark 在他的《基督教的興起》 Rise of Christianity一書中也有類似的描述和定量分析,關於這些城鎮教會對當時羅馬社會的深刻影響.

我們可以看出,保羅不但是在心誌上(heart)與耶穌同死同埋葬同復活,他的心智(beautiful mind)也同時被上帝使用.他被聖靈充滿,決心把福音從耶路撒冷,猶太全地,和撒瑪利亞,一直傳揚到地極.他實在是一個把自己全人(heart, mind, soul, and body)獻上當作活祭,為耶穌癲狂,為主燃燒的人!

能做到如此地心靈與心智,情感與理智的整全合一All for Jesus,在今天受過高等教育的信徒中間著實不易.這也許和我們自啟蒙運動以來的日益碎片化的神學教育發展有很大關係.

Dr. Edward Farley of the Vanderbilt University pointed out this issue in his book 神學教育的一致性及碎片化Theologia-The Unity and Fragmentation of Theological Education

“今天的神學已成為一種學有專攻的學術訓練,為要預備專業的神職人員；而在過往,神學被理解成為獲得智慧的靈魂指南,是日常生活中人人都需要的.”

“Today theology has become a mastery of academic disciplines for the purpose of professional ministerial preparation; in the past theology was understood as the orientation of the soul for the purpose of acquiring wisdom, which all men and women need regardless of their station in life. ”

不單是神學教育, 是我們整個教育哲學 (philosophy of education) 似乎都出現了問題.Dr. Harry Lewis, the former Dean of Harvard College,在他的 《失去靈魂的優秀─哈佛如何忘卻其教育宗旨? 》(Excellence Without A Soul: How a Great University Forgot Education) 一書中指出,包括哈佛大學在內的美國大學已丟棄了本科教育的根本-即旨在幫助那些十八九歲的年輕人長大成人,認識他們的身份,探尋生命的真意,當他們二十一二歲從大學畢業時可以成為更好的人.

He sees as the abandonment by American universities, including Harvard, of the fundamental job of undergraduate education… to turn eighteen- and nineteen-year-olds into twenty-one- and twenty-two-year-olds, to help them grow up, to learn who they are, to search for a larger purpose for their lives, and to leave college as better human beings.

哈佛大學如何忘記了Education for Life這一教育宗旨的確值得我們反思!頗有意思的是,在1642 哈佛大學章程裡有這樣的一款宣稱:「基督是一切真正知識學問的基礎」

Christ, in whom are hidden all the treasures of wisdom and knowledge-Colossians 2:3 “…一切智慧知識都在他(基督)裡面藏著.“

我想起詩篇11:3,“根基若毀壞,義人還能做什麼呢?”“When the foundations are being destroyed, what can the righteous do?”

也讓我想起哥林多前書3章10-11的說話,“我照神所給我的恩,好像一個聰明的工頭,立好了根基,有別人在上面建造,只是個人要謹慎怎樣在上面建造.因為那已經立好的根基就是耶穌基督,此外沒有人能立別的根基. ”

今天我很榮幸地站在這裡,再一次被提醒崇基學院的辦學使命 - 旨在繼承國內十三間基督教大學之辦學使命；本“崇奉基督”之精神,命名“崇基學院”. “崇基”亦兼含“崇高基礎”之意.

謝謝Walter校牧的邀請及給我用我的母語普通話證道的機會!我也想藉此機會分享一下我在大學期間信主的見證.

我是在國內出生長大,雖然外婆和母親都是基督徒,小的時候也跟著她們一起在家裡聚會；但自從踏入小學校門,就開始接受系統的無神論教育,讀到高中時甚至當上了學校的團支部書記,我真的是要準備好做共產主義的接班人了!

但就在考取大學的那一年發生了89學潮,我帶著極大的困惑開始了大學生活.我讀的是電子工程專業,但對哲學發生了極大的興趣-I was searching for a larger purpose for life after my belief in the so-called Communism.當時盛行的存在主義哲學確實給我的理性思考提供了一個新的框架,但遠遠不能給我的心靈帶來任何的滿足和盼望,我陷入極大的空虛苦悶之中.

這時,一位也是名叫保羅的美國留學生開始引起我的注意.他的陽光,喜樂,和樂於助人讓我想起了我們曾經的好榜樣“雷鋒”.但當我問他為什麼能做到這一點時,他的答案使我大失所望,也大吃一驚-他說是因為耶穌!

耶穌?!就是那位我外婆常常講的,又醫治了我媽媽的那位釘在十字架上的耶穌?

How could that be? 我開始對聖經有興趣了,為的是要看看這位耶穌究竟是誰?在短短的兩個月內,我和保羅的聖經學習使我興奮不已.特別是在約翰福音里耶穌的說話,“你們必曉得真理, 真理必叫你們得以自由”.之後耶穌又說到,“我就是道路,真理,生命,若不藉著我,沒有人能到父那裡去.”

這些震撼性的宣稱彷彿一束大光射入我的心靈深處!在91年的春天,我接受耶穌基督做我個人生命的救主.我在耶穌基督裡重新回到天父的懷抱,也在剛剛成立的校園團契裡成長和學習服事.

我今天站在崇基教堂的講壇上,心中充滿了感恩! 27年過去了,我經歷了耶穌所應許的豐盛生命.我感恩一直以來為我禱告的母親,我感恩給我把福音講明的保羅,我感恩在我屬靈成長旅程中的一個個恩師,與我一同承受生命之恩的太太,在職場和宣教工場上的同路人. . .

談到宣教,我越來越確信當時使徒保羅的策略和模式今日依然行之有效.我們在使徒行傳中和保羅的書信中發現他不但有Passion for Christ,並滿有智慧和策略.在保羅的宣教旅程當中,他大有能力地宣揚福音 Kerygma,也孜孜不倦地教導(didache)建立門徒和教會,並設立／委認地方教會領袖(長老和執事).保羅被解往羅馬之前的三次宣教行程／策略可以概括為這樣一個循環:Evangelize strategic cities, Establish and strengthen local churches, and Entrust to faithful men.

而他的簡單且有效的傳播媒介就是他的書信,他的信使就是他嘔心瀝血培養出來的門徒.這些門徒／信使在各地教會傳閱的豈只是保羅所寫的一封封書信,他們身上carry的豈不更是使徒保羅生命的DNA-耶穌基督!難怪他們當時被戲稱為Little-Christ-Christian.

在保羅的13封書信當中,我們也可以看出與之對應的信息.比方說,他的早期書信(eg 加拉太書和羅馬書)側重於對福音要義Kerygma的闡明,中期書信(以弗所和歌羅西書)則是側重於教會的建立,晚期書信(提摩太前後書)側重於勸勉領袖.

就在保羅離世前給他用福音所生的兒子的一封信中,他再次勸勉提摩太:“我兒啊,你要在耶穌基督的恩典上剛強起來.你在許多見證人面前聽見我所教訓的,也要交託那忠心能教導別人的人.”

就是這種代代相傳,生生不息的繁衍倍增organic multiplication,基督教會在接下來的三百年間重塑了羅馬帝國日益滑坡,分崩離析的道德面貌moral fabric.從宏觀層面來看,以Judeo-Christian為主導的世界觀,價值觀,和人生觀在過往的兩千年裡有力地支撐和維繫著人類文明的健康發展.勿容置疑,中國近代的高等教育和醫療衛生發展都是基督教文明的受惠者.

在今日高度全球化,數字化,多元化的背景下,做為基督徒,我們面臨著更嚴峻,更複雜的挑戰和機會.How Should then We Live? 保羅的使命人生又可以給我們提供那些屬靈的營養呢?

使徒保羅已毫不含糊地給我們發出邀請:“你們該效法我,像我效法基督一樣林前11:1 “Follow my example, as I follow the example of Christ.”

Yes, we can live and follow the Way of Christ and His Apostles!

當年的保羅處在一個表面盛世,內裡卻動盪不安的Pax Romana.他個人的身份認同,文化認同,及信仰體系在這種大時代背景下正醞釀著一場風暴.他沒有機會在“孔子學院”學習如何“靜心誠意,修身齊家,治國平天下”,但自從與道成肉身的耶穌基督在大光中相遇之後,領受了從天上來的異像和人生使命,便毅然決然地和過去的成就地位熱愛說bye-bye.

也期盼各位,無論您是在生命的哪一個階段,都可以像保羅那樣選擇以耶穌基督為至寶,把生命的主權交給上帝.因為他為你所預備的是你眼睛未曾看見,耳朵未曾聽見,人心也未曾想到的美好福音,豐盛生命!

只要您願意來就近耶穌,他正在叩響你的心門,等待你的邀請進入到你的生命當中…..

“若有人在基督裡,他就是新造的人,舊事已過,都變成新的了.”

Are you ready for a fresh start? A new beginning?

林後5:18-6:3. “(這)一切都是出於神,他藉著基督使我們與他和好,又將勸人與他和好的職分賜給我們.這就是神在基督裡叫世人與自己和好,不將他們的過犯歸到他們身上,並且將這和好的道理託付了我們.”

所以,我們做基督的使者,就好像神借我們勸你們一般,我們替基督求你們與神和好.神使那無罪的替我們成為罪,好叫我們在他裡面成為神的義.

我們與神同工的,也勸你們不可徒受他的恩典.因為他說:“在悅納的時候,我應允了你；在拯救的日子,我搭救了你.”看哪,現在正是悅納的時候!現在正是拯救的日子!

Let’s pray……

我 們 在 天 上 的 父 : 願 人 都 尊 你 的 名 為 聖 .

願 你 的 國 降 臨 ； 願 你 的 旨 意 行 在 地 上 , 如 同 行 在 天 上 .

我 們 日 用 的 飲 食 , 今 日 賜 給 我 們 .

免 我 們 的 債 , 如 同 我 們 免 了 人 的 債 .

不 叫 我 們 遇 見 試 探 ； 救 我 們 脫 離 兇 惡 .因 為 國 度 , 權 柄 , 榮 耀 , 全 是 你 的 , 直 到 永 遠 .阿 們!



\newpage

\section{生活}
\label{sec:421}
\textbf{林豪恩先生 - 生.活}
\newline
\newline
\hyperref[sec:420]{< < < PREV SERMON < < <}
~
\hyperref[sec:toc]{[返目錄]}
~
\hyperref[sec:422]{> > > NEXT SERMON > > >}


箴言 9:1-6
\newline
\begin{longtable}{cl}
\hline
\hline
章節 & 經文 (和合本修訂版)\\
\hline
9:1 & \begin{tabularx}{0.7\textwidth}{X} 智慧建造房屋，鑿成七根柱子， \end{tabularx} \\ \\ \relax
9:2 & \begin{tabularx}{0.7\textwidth}{X} 宰殺牲畜，調好美酒，又擺設筵席， \end{tabularx} \\ \\ \relax
9:3 & \begin{tabularx}{0.7\textwidth}{X} 派遣女僕出去，自己在城中至高處呼喚： \end{tabularx} \\ \\ \relax
9:4 & \begin{tabularx}{0.7\textwidth}{X} 「誰是愚蒙的人，讓他轉到這裡來！」又對那無知的人說： \end{tabularx} \\ \\ \relax
9:5 & \begin{tabularx}{0.7\textwidth}{X} 「你們來，吃我的餅，喝我調的酒。 \end{tabularx} \\ \\ \relax
9:6 & \begin{tabularx}{0.7\textwidth}{X} 你們要離棄愚蒙，就得存活，並要走明智的道路。」 \end{tabularx} \\ \\
[1ex]
\hline
\hline
\end{longtable}
詩篇 34:9-14
\newline
\begin{longtable}{cl}
\hline
\hline
章節 & 經文 (和合本修訂版)\\
\hline
34:9 & \begin{tabularx}{0.7\textwidth}{X} 耶和華的聖民哪，你們當敬畏他，因敬畏他的一無所缺。 \end{tabularx} \\ \\ \relax
34:10 & \begin{tabularx}{0.7\textwidth}{X} 少壯獅子尚且缺食忍餓，但尋求耶和華的甚麼好處都不缺。 \end{tabularx} \\ \\ \relax
34:11 & \begin{tabularx}{0.7\textwidth}{X} 孩子們哪，來聽我！我要將敬畏耶和華的道教導你們。 \end{tabularx} \\ \\ \relax
34:12 & \begin{tabularx}{0.7\textwidth}{X} 有誰喜愛生命，愛慕長壽，得享美福？ \end{tabularx} \\ \\ \relax
34:13 & \begin{tabularx}{0.7\textwidth}{X} 你要禁止舌頭不出惡言，嘴唇不說詭詐的話。 \end{tabularx} \\ \\ \relax
34:14 & \begin{tabularx}{0.7\textwidth}{X} 要棄惡行善，尋求和睦，一心追求。 \end{tabularx} \\ \\
[1ex]
\hline
\hline
\end{longtable}
以弗所書 5:15-20
\newline
\begin{longtable}{cl}
\hline
\hline
章節 & 經文 (和合本修訂版)\\
\hline
5:15 & \begin{tabularx}{0.7\textwidth}{X} 你們要謹慎行事，不要像無知的人，要像智慧的人。 \end{tabularx} \\ \\ \relax
5:16 & \begin{tabularx}{0.7\textwidth}{X} 要把握時機，因為現今的世代邪惡。 \end{tabularx} \\ \\ \relax
5:17 & \begin{tabularx}{0.7\textwidth}{X} 不要作糊塗人，要明白主的旨意如何。 \end{tabularx} \\ \\ \relax
5:18 & \begin{tabularx}{0.7\textwidth}{X} 不要醉酒，酒能使人放蕩；要被聖靈充滿。 \end{tabularx} \\ \\ \relax
5:19 & \begin{tabularx}{0.7\textwidth}{X} 要用詩篇、讚美詩、靈歌彼此對說，口唱心和地讚美主。 \end{tabularx} \\ \\ \relax
5:20 & \begin{tabularx}{0.7\textwidth}{X} 凡事要奉我們主耶穌基督的名常常感謝父神。 \end{tabularx} \\ \\
[1ex]
\hline
\hline
\end{longtable}
約翰福音 6:51-58
\newline
\begin{longtable}{cl}
\hline
\hline
章節 & 經文 (和合本修訂版)\\
\hline
6:51 & \begin{tabularx}{0.7\textwidth}{X} 我就是從天上降下來生命的糧；人若吃這糧，必永遠活著。我為世人的生命所賜下的糧就是我的肉。」 \end{tabularx} \\ \\ \relax
6:52 & \begin{tabularx}{0.7\textwidth}{X} 因此，猶太人彼此爭論說：「這個人怎能把他的肉給我們吃呢？」 \end{tabularx} \\ \\ \relax
6:53 & \begin{tabularx}{0.7\textwidth}{X} 耶穌對他們說：「我實實在在地告訴你們，你們若不吃人子的肉，不喝人子的血，在你們裡面就沒有生命。 \end{tabularx} \\ \\ \relax
6:54 & \begin{tabularx}{0.7\textwidth}{X} 吃我肉、喝我血的人就有永生，並且在末日我要使他復活。 \end{tabularx} \\ \\ \relax
6:55 & \begin{tabularx}{0.7\textwidth}{X} 我的肉是真正可吃的；我的血是真正可喝的。 \end{tabularx} \\ \\ \relax
6:56 & \begin{tabularx}{0.7\textwidth}{X} 吃我肉、喝我血的人常在我裡面，我也常在他裡面。 \end{tabularx} \\ \\ \relax
6:57 & \begin{tabularx}{0.7\textwidth}{X} 永生的父怎樣差我來，我又怎樣因父活著，照樣，吃我肉的人也要因我活著。 \end{tabularx} \\ \\ \relax
6:58 & \begin{tabularx}{0.7\textwidth}{X} 這是從天上降下來的糧，不像你們的祖宗吃過嗎哪還是死了；吃這糧的人將永遠活著。」 \end{tabularx} \\ \\
[1ex]
\hline
\hline
\end{longtable}
講題:生.活　Live a Life that Matters

經文:箴言9:1-6;詩篇34:9-14;以弗所書5:15-20;約翰福音6:51-58

講員:林豪恩先生

場合:香港中文大學崇基學院禮拜堂

日期:2018年8月19日

一．引言

有一個問題,如果你能夠回答,答應必定是「是」,而不會是「否」,這問題是:「你現在仍然活著嗎?」然而,如果問題改為:「你活著的狀況如何呢?」就不容易回答,而且答案也多樣化.

今天的講道題目是「生.活」,英文是Live a Life that Matters,因為今天的四段經課都提到Live和Life,稍後再閱讀經文時詳談.生與活好像分不開,然而,有些人卻體會到活著不等於生活,他們說:「只有活著,沒有生活」.以下的自白不是香港人的故事,大家不用對號入座.

「每天我下班回到家裡後,第一件事是開燈,第二件事是開音響.光和聲音都是好東西,它們永遠可以一下子灌滿房間.每月交完房租,感覺人生淒涼慘澹.在公司得不到重視,開始懷疑自己的能力.每次接到電話,爸媽都勸你「賺不到錢就回家吧」.你半夜的時候,被客戶挨千刀的微信吵醒.你站在商場裡,算計著這個月可以隨意支配的錢還剩多少.你根本不介意出差,畢竟所到之處感覺都不是家.人生迷茫,不知道自己要什麼,找不到滿意的工作,無論多努力,付出多少,仍然徘徊在生存線上.這一切讀起來都那麼辛酸而真實.只有活著,沒有生活；人如機器,活著乏味.(\url{http://www.rsdyy.com/archives/view-10962-1.html})

或者不少人也有同感,人生好像一部啟動了的機器,不知被甚麼推動著,但卻不能不繼續運作下去.這位朋友感慨地說,不知道自己為甚麼而活,我們又知道自己為甚麼而活嗎?知道自己被甚麼推動著嗎?《標竿人生》一書列出五種支配人生的力量:第一種是罪咎；第二種是忿恨；第三種是恐懼；第四種是物質享受；第五種是別人的期望.我們的人生會否也被這些力量的其中一種駕馭呢?

讓我們再閱讀今天的幾段經課,尋找對這些駕馭我們人生的力量的回應,以及啟發我們如何可以有生活地活著.

二. 內容

2.1. 箴言9章:以智慧生活擺脫被罪咎駕馭

9:1智慧建造房屋,鑿成七根柱子,

9:2宰殺牲畜,調好美酒,又擺設筵席,

9:3派遣女僕出去,自己在城中至高處呼喚:

9:4「誰是愚蒙的人,讓他轉到這裡來!」又對那無知的人說:

9:5「你們來,吃我的餅,喝我調的酒.

9:6你們要離棄愚蒙,就得存活(and live),並要走明智的道路.」

箴言講智慧,現代人更講智慧,例如政府也曾經提出把香港建設成為「智慧城市」.現代人的智慧和箴言的智慧不同,現代的「智慧城市」是Smart City,箴言的智慧是Wisdom；現代人雖然每人都有Smart Phone,但都不見得每人都有智慧.今天我們讀的箴言提出的不是「智慧城市」,而是「智慧生活」,呼籲以智慧生活擺脫被罪咎駕馭的人生.

箴言九章頭幾節在說怎樣的智慧呢?讓我們一併讀最後的幾節,它們是首尾呼應的.

13愚昧的女子喧嚷,她是愚蒙,一無所知.

14她坐在自己家門口,在城中高處的座位上,

15呼喚過路的,向那些在路上直走的人說:

16「誰是愚蒙的人,讓他轉到這裡來!」又對那無知的人說:

17「偷來的水是甜的,暗藏的餅是美的.」

在世界這個大市集上,智慧與愚昧都擺設筵席,邀請人們跟隨它們.人們不容易分辨,因為它們站的位置是一樣的(在城中高處),它們的呼籲也是一樣的(誰是愚蒙的人,讓他轉到這裡來!).其中一把聲音叫喊說:「你們要離棄愚蒙,就得存活,並要走明智的道路.」說得有道理,不過卻是老生常談；另外一把聲音說:「偷來的水是甜的,暗藏的餅是美的.」聽來有點那個,但卻有吸引力.偷有兩種吸引人的地方,其一是實用性的,例如建築工程中偷工減料,即是節省了開支,增加了利潤.遠處的例子有四川的豆腐渣學校,近事的例子有香港的沙中線工程.偷的另一種吸引力是刺激感,正如明代的《雪濤小說》中那句名言說,「妻不如妾,妾不如婢,婢不如妓,妓不如偷.」

在奇怪的制度下,建築豆腐渣學校的人沒有被法律制裁,不過,午夜夢迴,看見老師和學生被壓死的慘況,聽到死者家人呼天搶地的哭號,他們能夠安睡嗎?曾經有人對韓戰及越戰死亡的美國士兵進行病理解剖,發現他們之中很多人血管硬化,這些士兵大多數是年青人,為甚麼出現老化的現象呢?研究人員認為其中一個原因是殺人令他們內心深處感到罪惡感而影響到身體的狀況.(林仁和:老人心理學)

「偷來的水是甜的,暗藏的餅是美的.」聽來很吸引,很有刺激感,但做了壞事,必然會受到良心的譴責,為了掩飾和粉飾,難免一生都要被自欺欺人的謊言所糾纏.有人以為做到「神不知鬼不覺」是聰明.是的,當人們看見一個新建的鐵路站,怎會想到下面暗藏了甚麼呢?然而,正是這些在地低下暗藏著的,在看不見的地方被偷走的,腐蝕我們的城市.難免令人推想,這些年間,我們的城市還有多少美好的價值被暗藏著偷走呢?智慧呼籲:「你們要離棄愚蒙,就得存活,並要走明智的道路.」有生活的活著,不是偷偷摸摸的,而是光明正大的.如何獲得智慧呢?明智何處尋呢?箴言第九章中間的部份說清楚了:「10敬畏耶和華是智慧的開端,認識至聖者便是聰明.」第10節「聰明」這個字,與第6節「明智」本是同一個字.

2.2. 詩篇34篇:以敬畏上主擺脫憤恨與恐懼的駕馭

34:11孩子們哪,來聽我!我要將敬畏耶和華的道教導你們.

34:12有誰喜愛生命(which of you desire life),愛慕長壽,得享美福?

34:13你要禁止舌頭不出惡言,嘴唇不說詭詐的話.

34:14要棄惡行善,尋求和睦,一心追求.

若非策略性質或表演性質(例如在某些地方的議會辯論中),口出惡言的人很有可能是感到憤恨,把內心的激烈情緒宣之於口.若非攻於心計或圖謀不軌,詭詐的話通常出自恐懼的心,希望能逃避懲罰或責任.

父母對孩子口出惡言,例如:「生舊叉燒好過生你」,其實是發洩恨鐵不成鋼的憤怒；父母對孩子說詭詐的話,例如:「你看,鄰居明仔幾乖,不會像你那麼頑皮.」(其實明仔是否乖也未可知),以為這樣說能夠激勵孩子,其實被自己內心對孩子不中用的恐懼所控制.

以上這些只是小例子,然而,有些人的一生被憤恨駕馭,他們出現,其他人就能感受到他們的「氣場」；有些人一生被恐懼駕馭,他們不但自己恐懼,還會把恐懼散播給接近的人.被憤恨或恐懼駕駛的人,雖然活著,但不會感到生命是美好和值得的.

詩人說,喜愛生命,愛慕長壽,得享美福,負面來說是要「禁止舌頭不出惡言,嘴唇不說詭詐的話」,而真正治標又治本的,是擺脫內心的憤恨和恐懼.詩人又說,「我要將敬畏耶和華的道教導你們」,敬畏耶和華能夠回應我們內心的憤恨和恐懼嗎?

《竭誠為主》的作者章伯斯(Oswald Chambers)曾說:“The remarkable thing about God is that when you fear God, you fear nothing else, whereas if you do not fear God, you fear everything else.”

有一位退休醫生,在七十五歲時分享以下的見證:

「我父親是商人,家境富裕.姊弟六人,最小的是個男孩,我排第四,姊弟們備受寵愛,我被父母忽視.讀書,工作,結婚,生子,退休,匆匆走過了幾十年.回顧七十歲前的我,活在一片黑暗中,尤其在那充滿鬥爭的動亂年代,人人自危,總是膽顫心驚.我好吵架是出了名,能為一點小事大吵大鬧,像機關槍不停地叫,藉此宣洩內心的怨氣.中國有句古話:『江山易改,本性難移.』確實,誰都無法改變我這個死不認錯的老婦人.二○○四年十一月底的一天夜裡,我在憤恨,絕望,痛苦的孤獨中,一口氣嚥下半瓶安眠藥,但卻被救回.十二月,一位素未謀面的老牧師知我的遭遇,特來慰問.他領我讀約翰福音一章12節:『凡接待祂的,就是信祂名的人,祂就賜他們權柄作上帝的兒女.』他問:『妳願意打開心門接受耶穌,承認祂是妳的救主嗎?』我答:『願意.』他便帶領我做了決志禱告.決志禱告中,我痛哭流淚,淚水像湧泉不停地流,我感到上帝把我堅硬冰冷的心溶化了,祂進入我生命中,是那樣的溫暖,撫平了我心靈一切傷痛.我像一個被丟棄,被遺忘,受盡委屈的孩子,來到父母的面前盡情地哭訴.在安靜中我彷彿聽見天父慈愛的召喚:『孩子,回家吧!』我在世間漂泊了七十年,今年終於回家了!回到慈父的懷抱.上帝的愛把我埋藏了數十年之久的怨恨,苦毒,不平通通都溶化掉了.」(程瑤珍:上帝醫治我的憂鬱)

詩篇34篇提示我們,以敬畏上主擺脫憤恨與恐懼的駕馭,可以重新體會生命的溫暖和安穩.

2.3. 以弗所書:以感恩的心擺脫物質享受的駕馭

5:15你們要謹慎行事(be careful then how you live),不要像無知的人,要像智慧的人.

5:16要把握時機,因為現今的世代邪惡.

5:17不要作糊塗人,要明白主的旨意如何.

5:18不要醉酒,酒能使人放蕩；要被聖靈充滿.

5:19要用詩篇,讚美詩,靈歌彼此對說,口唱心和地讚美主.

5:20凡事要奉我們主耶穌基督的名常常感謝父上帝.

現今的世代有多邪惡呢?它把人生的意義壓縮至擁有,擁有愈多就愈快樂,愈重要,也愈安全.它認為能夠放任地享受物質及娛樂,才算有生活.以弗所書卻提醒人們「要謹慎行事」.類似箴言,作者用「無知的人」與「智慧的人」對比,17節呼籲人們「不要作糊塗人」,如何做呢?就是「要明白主的旨意如何」,接著以「不要醉酒,酒能使人放蕩」補充「不要作糊塗人」,以「要被聖靈充滿.要用詩篇,讚美詩,靈歌彼此對說,口唱心和地讚美主.凡事要奉我們主耶穌基督的名常常感謝父上帝」來補充「要明白主的旨意如何」.

作者提出以「常常感謝父上帝」的態度來擺脫物質享受的駕馭.物質支持我們活著,但生活的品質不一定與物質多少成正比.心懷感恩的人,會感到生命比較值得過及有意義.心理學研究發現,感恩傾向較強的人較常體驗正向的情緒,如愉悅,熱情,愛,快樂,樂觀等,培養感恩心讓人較不易受制於極負面的衝動,如嫉妒,不平,貪婪,怨恨.此外,擅於感恩的人較能因應日常生活的壓力,身體較健康,生病也較容易恢復.感恩心將惜福的心提升到最高點--能夠珍惜人生,人際關係及上帝的恩惠.當你生存在這樣的狀態,會比較容易感到活著是有生活,生活是有意思的.

2.4. 約翰福音:以為上主而活擺脫被別人的期望駕馭

6:51我就是從天上降下來生命的糧；人若吃這糧,必永遠活著(live forever).我為世人的生命所賜下的糧就是我的肉.」

6:52因此,猶太人彼此爭論說:「這個人怎能把他的肉給我們吃呢?」

6:53耶穌對他們說:「我實實在在地告訴你們,你們若不吃人子的肉,不喝人子的血,在你們裡面就沒有生命(you have no life in you).

6:54吃我肉,喝我血的人就有永生(have eternal life),並且在末日我要使他復活.

6:55我的肉是真正可吃的；我的血是真正可喝的.

6:56吃我肉,喝我血的人常在我裡面,我也常在他裡面.

6:57永生的父怎樣差我來,我又怎樣因父活著(I live because of the father),照樣,吃我肉的人也要因我活著(whoever eats me live because of me).

6:58這是從天上降下來的糧,不像你們的祖宗吃過嗎哪還是死了；吃這糧的人將永遠活著(live forever).」

約翰福音不但講及活著,更提出「永遠活著(live forever)」.其實,永遠活著,不一定是好消息,特別是對於開首時引述的那位感慨「只有活著,沒有生活」的朋友來說.那位朋友說:「在公司得不到重視,開始懷疑自己的能力.每次接到電話,爸媽都勸你“賺不到錢就回家吧”.」他很可能感覺自己是失敗者.失敗,是達不到公司的期望,得不到公司的重視；失敗,是雖然父母肯收留,但卻感到未能滿足父母盼望自己自立的合理期望；失敗,是連自己對自己的期望都達不到......如果在長期達不到別人及自己的期望地活著,還能說甚麼生活呢?從另一個角度說,如此的生存狀況,等於讓別人的期望操縱自己的人生,往往迷失在別人的期望之中.有人會說:「總不能不理會其他人而生活的.」說得對,一個目中無人或者生活在孤島的人,遲早都會感到活著而沒有生活,連耶穌也說:「一個人不能事奉兩個主.」即是我們的生存是有方向的,有信靠的對象,有交待的對象.我們以不同的東西作為人生信靠和交待的對象,會產生不同的生存狀況及生活質數.

耶穌是如何活著的呢?他的活著,超越了別人的期望,甚至也超越自己的期望,他說:「6:57永生的父怎樣差我來,我又怎樣因父活著(I live because of the father)」他這種因父活著的生命,是一種能夠永遠活著的生命.我們呢?他說:「照樣,吃我肉的人也要因我活著(whoever eats me live because of me)」.因耶穌而活,讓我們活著有方向,有意義,幫助我們超越別人的期望,也脫離自我的困局.

三．總結

有人說,年輕一代越來越得懂得追求生活.求職網站jobsDB公布《2017年十大企業調查報告》,26至33歲的「千禧世代」認為「工作與生活平衡」是最重要,也是三大求職考慮因素之一.「工作與生活平衡」其實並不一定令人感到活著有生活,準時返工放工的人也不一定感到生存得有意思,而長時間投入工作的人也不一定感到生活乏味,今天的幾段經課啟發我們去思考,活得有智慧,存敬畏上主的態度活著,存感恩的心活著,按上主的差遣而活著,如此活著,會感到生命沒有枉過,如此的生命是值得延伸至永恆的.

我看見有些大學畢業生奮力追求有意義的人生,期望自己能為身處的社區及人群做一點好事.有人在土瓜灣做社區建設,期望能夠協助舊區的居民守望相助；有人在太子和深水Po 地區協助婦女自力更生；有人投身環境保育；有人甚至跑往非洲某國參與義務的種族復和工作.日子漸過,年紀漸長,他們都遇到龐大的壓力,因為這個城市並不獎勵這些追求意義而長期做前線工作的人.當其他同學都有車有樓,將來生活有保障的時候,難免令他們質疑自己所行的路.這些同學的處境令我們思考,當社會呼籲年輕人追求有意義的人生,為使命而活,為他人的福祉而活的時候,制度也當與社會對他們的呼籲配合.



\newpage

\section{難為正邪定分界也得裝身作靈戰}
\label{sec:422}
\textbf{鄭漢文博士 - 難為正邪定分界也得裝身作靈戰}
\newline
\newline
\hyperref[sec:421]{< < < PREV SERMON < < <}
~
\hyperref[sec:toc]{[返目錄]}
~
\hyperref[sec:423]{> > > NEXT SERMON > > >}


講題:難為正邪定分界 也得裝身作靈戰

Be Prepared as Spiritual Fighter amidst Unclear

講員:鄭漢文博士

場合:香港中文大學崇基學院禮拜堂

日期:2018年8月26日

1. 難為正邪定分界

各位弟兄姊妹,好久沒有在講壇相見.

今年五月一連三周看書打字寫文,終於完成論文,我終於七月底神道學學士畢業,結束了神攻十年.今年8月1日開始兼職傳道人.今日第一次在新身分下講道.

八月底講道,其實重疊著開學的備課.

我在中大教育學院的工作包括開發價值教育.自2006年開創價值教育文學碩士課程,到了這一年是最後一屆,每屆我都負責任教一個必修科MAVE6000「價值教育的原則,實踐與教學」.在這個科中,其中有一課我一定會教以下的論點:價值語言其實有三種.

第一種是「是非對錯」的語言(language of right and wrong)；

第二種是「好壞利害」的語言(language of good and bad)；

第三種是「善惡正邪」的語言(language of good and evil).

今日的經課,我選了以弗所書6章10-20節來跟大家一同思考.

今日的經課,涉及第三種價值語言,關心人如何面對善惡正邪.正與邪,容易分嗎?

一講及善惡正邪,我就想起《難為正邪定分界》這首歌,課上我一定請同學一起唱.這首歌是香港電視劇集《飛越十八層》主題曲.

但今日,讓我們聽一聽原唱者的心聲:

《難為正邪定分界》(作詞:鄭國江；作曲:顧嘉輝；合唱:葉振棠/麥志誠)

人:　對抗命運但我永不怕捱 過去現在難題迎刃解

人生的彩筆蘸上悲歡愛恨 描畫世上百千態

魔:　控制命運任我巧安排 看似夢幻凡人難盡瞭解

人間的好景給我一朝破壞 榮辱愛恨任分派

人:　努力未願平賣 人性我沒法賤賣

魔:　今天死結應難解

人:　努力興建

魔:　盡情破壞

合:　彼此也在捱

人:　世界腐敗犯法那需領牌(看吧邪力正強大)

人:　法理若在為何強盜滿街(看吧強盜滿街)

人:　人海的衝擊比那蒼海更大(啊啊)

合:　難為正邪定分界

這種人魔大戰,但又難為正邪定分界,是在敘說那個時代?

今天香港,特別應景.

我們能夠逃避這種正邪的人魔大戰嗎?

2. 難以不參戰(6:10-12):人在神魔大戰中有分

今日的經課以弗所書6章10-20節告訴我們,我們是無法逃避這種正邪的人魔大戰的.

讓我們一起閱讀以弗所書6章第10至12節:

6:10	我 還 有 末 了 的 話 , 你 們 要 靠 著 主 , 倚 賴 他 的 大 能 大 力 , 作 剛 強 的 人 .

6:11	要 穿 戴 　 神 所 賜 的 全 副 軍 裝 , 就 能 抵 擋 魔 鬼 的 詭 計 .

6:12	因 我 們 並 不 是 與 屬 血 氣 的 爭 戰 , 乃 是 與 那 些 執 政 的 , 掌 權 的 , 管 轄 這 幽 暗 世 界 的 , 以 及 天 空 屬 靈 氣 的 惡 魔 爭 戰 . 〔 兩 爭 戰 原 文 都 作 摔 跤 〕

人在神魔大戰中,難以中立.不參戰,也會被捲入.

我們在常識上不難經歷是非對錯以及好壞利害,那麼人如何經驗善惡(good and evil)中的邪惡?如何在電影中描繪人所經驗的邪惡(evil being portrayed in cinema)?.

Hollywood電影裡,有那些電影中描繪人如何與惡魔為戰?

其中一個系列大家都會看過:星球大戰Star Wars.

其中一位最深刻的惡魔典型要算Darth Vader,穿上日式武裝,黑頭盔,機械呼吸器下發聲的黑武士.

究竟他如何成長?如何走入魔道?

讓我們思考一個在光明面的年青人如何走向黑暗的魔道.這個故事記載在《星球大戰前傳三之黑帝君臨》(Star Wars: Episode III Revenge of the Sith).

電影中敘述進入青年期的Anakin Skywalker,在絕地武士的學藝上日漸精進,為共和國而每戰每勝立功無數,但絕地武士議會Jedi Council通過他只能當列席成員,而未能成為正式大師master.他心中出現才能未能被確認的憤恨不平rage/anger of de-recognition.

在公共世界中,他充滿難言的不平,憤恨.

在私自共世界中,Anakin在內心充滿不安,恐懼.

Anakin成長時經歷母親被害身亡,憤而殺盡所有害母族群.

他經歷師傅Qui-Gon死於戰鬥,然後師從師兄Obi-Wan Kenobi,但只感到新師傅處處批評他,低估他的能力.

Anakin夜夜睡夢中,預感正在懷孕的秘密婚妻Padmé將會受苦下死亡.

他不願再被死亡所制約,要找到增強力量,超越死亡的方法……

共和議會議長Palpatine被Anakin發現真身是黑帝Sith Darth Sidious.

黑帝這樣誘騙Anakin:

一是絕地武士議會驚怕他原力超過他們；

二是唯有動用原力的黑暗力量dark side of the Force,才能拯救其愛妻免於痛苦死亡；

三是唯有黑帝懂得教他動用原力的黑暗面力量；

所以,黑帝邀請Anakin跟隨他.

Anakin心中嚮往原力的黑暗面力量可以救亡,絕地武士議會議長快要殺死黑帝時,為救活懂得救妻於死亡的議長Sith Darth Sidious,唯有出手:

一是先殺絕地武士議會議長,之後無法回頭；

二是跪向黑帝求收為徒,被命名為Darth Vader,立即執行新拜師傅下的清殺令

三是再殺所有童年絕地武士,此惡行成為Jedi清理門戶追殺他的根據.

從此,Darth Vader成為黑帝帝國的殺手.

Padmé找他時,師傅Obi-Wan跟來,讓Anakin以為愛妻Padmé出賣他而緊握她的頸以至昏迷.

Anakin對愛人說:”Love won’t save you, only my power can.” 愛無法救,只有權能.信權能不再信愛的他,已入魔道.

他對權能這樣想:”I am more powerful than the Chancellor. I can overthrow him. I can rule the galaxy. ” 我快會更有能力,超過黑帝,可以稱帝,雄霸宇宙.宇宙最強的他,成為魔道.

他對師傅Obi-Wan表示:”If you are not with me, you are my enemy. ” 不與我同行即是與我為敵,這就是魔道.

使徒大戰下,Anakin被師傅斬去雙腳,在火湖中被燒雙手,行將快死,師傅離他而去.

此時感到被放棄的Anakin對師傅說:”I hate you.”

對徒弟失望的Obi-Wan說:”You are my brother. I love you.” 但同時痛心地掉頭走了.

相隔不久,黑帝來火湖救Anakin,把他重新改裝成黑武士維達大臣Lord Vader.

Lord Vader醒來第一件事問新主人黑帝: 愛妻如何?Is she safe? Is she alright?

黑帝答: “It seemed in your anger, you killed her.”　他大叫: “No!”

邪惡力量,是借更大力量與權能的渴望而會讓惡由心生.

Jedi 不宜有自己的possession and attachments,該有compassion.Anakin因為私心,失去大公.

這是一個武士失落的故事(The fall of a Jedi).

因愛成恨,渴望力量權能無限擴大,不平於他人對自己的未被確認……都會讓人走入魔道.

3. 首在五防守,腰胸腳手頭(6:13-17):

真理腰帶,公義護胸心鏡,平安福音鞋,信德籐牌,救恩頭盔

究竟與屬靈氣的惡魔爭戰所要穿戴神所賜的全副軍裝是怎樣的?

讓我們一起閱讀以弗所書6章第13至17節:

6:13	所 以 要 拿 起 　 神 所 賜 的 全 副 軍 裝 , 好 在 磨 難 的 日 子 , 抵 擋 仇 敵 , 並 且 成 就 了 一 切 , 還 能 站 立 得 住 .

6:14	所 以 要 站 穩 了 , 用 真 理 當 作 帶 子 束 腰 , 用 公 義 當 作 護 心 鏡 遮 胸 ．

6:15	又 用 平 安 的 福 音 , 當 作 預 備 走 路 的 鞋 穿 在 腳 上 ．

6:16	此 外 又 拿 著 信 德 當 作 籐 牌 , 可 以 滅 盡 那 惡 者 一 切 的 火 箭 ．

6:17a	並 戴 上 救 恩 的 頭 盔

6:17b	拿 著 聖 靈 的 寶 劍 , 就 是 　 神 的 道 ．

許多解經家都指出這副軍裝有六部份,首五部份都是防衛性的(defensive),單單只有第六部份是攻擊性的(offensive).我們先探討這軍裝五部份的性質.

第13節指出裝身的目的是要站立得住(stand strong).

之後4節提及「腰胸腳手頭」軍裝全副共有以下五部份:

首先是腰,用真理當作帶子束腰.

其次是胸,用公義當作護心鏡遮胸.

第三是腳,用平安的福音當作預備走路的鞋穿在腳上.

第四是手,拿著信德當作籐牌,可以滅盡那惡者一切的火箭.

最後是頭,戴上救恩的頭盔.

除了剛才提及,這一身裝備,是要以保命為上的,目的在於免受戰敵所傷.

值得一提的是屬靈氣的惡魔爭戰所要穿戴的其實是五重的屬靈品性(fivefold spiritual qualities),腰胸腳手頭需要培養相應的一真理,二公義,三平安,四信心與五救恩.

我們可以從相反的角度解讀這五重的屬靈品性:

第一是真理的腰帶,好重要,不撐真理,就會造假,用香港話:「無腰骨撐真理.」

第二是公義的遮胸護心鏡,是要護心,不公因為有私心,不義因為無良心,保守你的心勝過萬有.

第三是平安福音鞋,不平,不安就容易受傷.

香港人祝福友人:「出入平安」.出門最緊要平安.

我想起去年暑假和今年暑假兩次出門都受傷.

去年七月,我去了英國朋友家住了十日,有天在屋內二樓下來從樓梯失足滑跌落地,折傷了右腳首二腳指,第二指骨節錯置,同時在離開朋友家時因意見不合也傷了友誼.

今年六月,往高雄出任訪問教授25天,回想是當時在倦乏不堪中動身,心緒不靈中起行,上機落機路上拿了過重的行李,拉傷了右肩,之後運轉不靈,時常陣痛.回來醫治至今依然未康復,影響到手字的書寫.

我的治療師問我:「為何你每一次長途旅程都受重傷?」

「平安」真的好重要.

回到第四是信德盾,人失信,甚麼箭都擋不了.

Hollywood英雄電影類型中,漫威電影Marvel 與 DC 兩個系列一直在爭勁,各有粉絲(fans).

兩大系列各有一位英雄用盾,用籐牌(shield):復仇者聯盟Avengers中的美國隊長Captain American用盾,正義聯盟Justice League中的神奇女俠Wonder Woman也用盾.

在電影中,二者都會用盾作攻守.守時,可以保護全身抵擋攻擊,攻時,可以飛盾傷敵.

表面看來,用盾時似是一個人的進退,實際是為團隊打頭陣,首出擋炮火,擋火箭.

另一套電影《戰狼三百》中,有一位斯巴達(spartan)殘障人因未能舉盾過肩而未被重用,因為羅馬戰陣是每位軍人舉起方盾作一整體式的防守與進攻,必須暨可保護自己又能隨時舉盾過肩保護戰友.

信德作籐牌,是讓信心同時群己兼顧.

在《復仇者聯盟3:無限之戰》(Avengers 3: Infinity War)中,Thanos在取得六粒無限寶石,就可在彈指之間,在偶然揀選下讓一半人口消逝,好讓餘下的一半活得更好.這是他的使命,是正是邪?

宇宙英雄合力阻止他的,各有裝備.在《復仇者聯盟3:無限之戰》最後一役進行群戰驅敵前,Wakanda的國王黑豹Black Panther指著Captain America說:”Get that man a shield”.

神奇女俠Wonder Woman拿起盾作進擊前,她會說:”Shall we?”(讓我們一起作戰.)

信仰在作戰時不是一個人的事.”Shall we?”(讓我們一起作戰.)

信心這屬靈德性(faithfulness as spiritual virtue)從來不單單是一個人在心理上的強弱,而是一群信徒的進退.

中國人有一句話:「窮則獨善其身,達則兼善天下.」

流行著這樣的解釋:「一個人在條件充足時當然可以顧及他人,但一個人在條件不足時唯有只顧自己.」

我認為這個解釋未夠深刻.我是這樣解說的:

一個人在條件不足時首先要顧及自己,不讓惡臨在自己身上,好好保留一己的善性.香港人所謂:發窮惡,就是在條件差時人就變惡了,容讓邪惡佔領了自己的心,失掉最基本的善性.

一個人先要守住自己的善性──不入魔道,在條件充足時就可以顧及他人,讓天下的惡一一被清除.

先養善心,進退間不入魔道,行有餘力,即可以除惡降魔.

全副軍裝,最後是頭,因為斬首,可以致命.頭受少少傷,都可以變成腦殘.元首也者,頭也,領導全身.首領也者,領導全群.

頭盔保護人的大腦,讓人保持清明.在他人看來,頭盔也表明披戴者的身分與戰意.

Iron Man當然全身裝配都好重要,但最重要是他的頭盔,因為他的頭盔不單跟全身裝配一樣堅硬可以護身,而且是人工智能與人的智能的結合,往往在千鈞一髮間在頭盔內思索進退取捨.Iron Man電影中有不少鏡頭放在Tony Stark的面目來倒影頭盔內智能化了的思考.

以弗所書在這裡提的是救恩的頭盔,是甚麼意思呢?

救恩,是神所給與我們的,不是我們自己建立的,所以不再於我們的行為.

救恩,是神早早已經給與我們的,作戰時就是我們的身分,表明為神而戰.

救恩,是神所給的,是敵人無法取走的戰意,在未戰之時已知必勝.

裝備若被拿走了,我們還有甚麼?還是甚麼?

在雷神索爾Thor系列中我們看到裝備的深層意義.

在Thor成長時以為因可用神錘(hammer)就所向無敵,有點不可一世,父王就教訓他,把神錘投落地球,也把Thor投落凡間,直至有一天他配得敬重時才能再運用此神錘.

到得Thor成熟時,(在《雷神索爾3:諸神黃昏》Thor 3: Ragnarok)才發現自己原來有個家姐,返回Asgard奪權,濫殺無辜.電神Thor以為可以用神錘攻勝,誰知一下子就被家姐單手粉碎了個錘.忽然失去裝備,力量也相應失去.

在戰意最低,快要戰敗時,爸爸提醒他是Thor of Thunder而非Thor of Hammer,他才發現,力量源頭不在神錘這個工具,而是在運用神錘工具的心,就是神本身.

Wonder Woman也有相近的經歷,神奇女俠以為拿著「殺神劍」(Godkiller, Sword of Athena)就可一劍殺死誘惑人間戰亂的戰神Ares,誰知戰神一下子就熔化了她以為無堅不摧的神劍.戰神Ares跟她說:”Only god can kill god”.原來,神奇女俠本身才是「殺神劍」.

以弗所書6章第17節明言,唯有可以戰勝惡魔的是聖靈的寶劍.

4. 五四一陣式下四中場助聖靈進攻(6:17b-20):

神道寶劍,儆醒不倦,禱告代求,當盡本分,放膽講福音

全副戰裝最重要的就是神本身,唯有可戰的是聖靈的寶劍,就是神的道.

讓我們一起閱讀以弗所書6章第17b至20節:

6:17b	拿 著 聖 靈 的 寶 劍 , 就 是 　 神 的 道 ．

6:18	靠 著 聖 靈 , 隨 時 多 方 禱 告 祈 求 , 並 要 在 此 儆 醒 不 倦 , 為 眾 聖 徒 祈 求 ,

6:19	也 為 我 祈 求 , 使 我 得 著 口 才 , 能 以 放 膽 , 開 口 講 明 福 音 的 奧 秘 ,

6:20	( 我 為 這 福 音 的 奧 秘 , 作 了 帶 鎖 鍊 的 使 者 ) 並 使 我 照 著 當 盡 的 本 分 , 放 膽 講 論 .

真正與惡魔大戰的是神,我們極其量只是助陣,當然首先要以五部裝身保護自己.

後三節經文可以列出四種助陣的信仰行動:

一是儆醒不倦,

二是祈禱代求,

三是當盡本分,

四是放膽講福音.

但是,這四種信仰行動是助陣的,是依存於對聖靈的信靠.

我是在中四升中五的暑期福音營信主的,九月開學後跟同學返教會,先返少年團契.當年見到一位神智不清的姊妹,常自言自語一些信仰宣告(後來知道是類近護教學的論述).不明所以,後來較年長的告訴我:她認為摩門教是異端邪說,曾經單人匹馬走入敵營,她單刀赴會在人地教會大辯大論三日三夜,之後就痴痴呆呆,自成一台.

這個人物的鮮活故事在我信仰成長時期是非常深刻難忘的,至今依然歷歷在目時刻提醒我,而另一深刻記憶的是當時團契常唱的一首歌的其中一句:「非勢力,非才能,乃靠靈成事.」

儆醒不倦與祈禱代求這兩者是我們天天要學的屬靈操練,然後有當盡的本分就去履行,有機會時就按我們口才恩賜的大小放膽講論福音.

但是,這四式都是助陣的,好像足球比賽的中場,助攻二傳的,我們只有單箭頭,就是聖靈的寶劍,神的道.若要中場傳送做得好,必須時時勤加練習,配合單箭頭的靈動.

我開始認真地準備今天這次講道,是在七月世界盃最高峰的時期,一邊觀戰,一邊備戰.

所以,我們試一試用足球戰術來類比今日的經課.

這段經文排的戰術是五四一陣式.

在前場作指揮的隊長是聖靈,單箭頭.入球取勝的希望在前鋒,中場助攻,後防不失.

中場分為儆醒不倦,祈禱代求,當盡本分,放膽講福音四個位,但全要靠著隊長聖靈作進退.

真正在前場攻球的是聖靈,四位在中場的只是傳球,二傳給聖靈攻堅.

後防分為腰胸腳手頭五位,後防的目的就是要站立得住(stand strong),抵擋攻擊免受損,力保不失,一定不可輕上.

後防若出錯可以錯在那裡?可以想像:五個後防,表面好強,底子裡,點知一個打假波,一個私下落注買波賭自己隊輸,一個穿了鴛鴦鞋跑路不平衡,一個從來都無自信,一個似中堅但用後腦頂頭搥好易擺烏龍.無真理,無公義,無平安,無信德,無恩待,問你死未?

附帶一句:現在的香港,好似一隊後防出錯的球隊.

好在還有四位中場在守望.那麼,香港的靈,還在嗎?

不說了,只有唏噓.

5. 最後,如何裝身作靈戰?

最後,一個人如何裝身作靈戰?

讓我們想一想魔戒Lord of the Rings中 Frodo的故事.

他不過是中土不起眼的Hobbit人,不單身型矮細,智能也普通,但正因對高超沒有渴望,天將降大任於斯人也,他最能拒絕魔戒的誘惑:因為任何人載上魔戒即可上接魔道權能號令天下.Frodo的清心,形成了一個守戒團契The Fellowship of the Ring,集合各式戰士好友,彼此守望,各自裝身,一起打這場靈魔大戰,護送Frodo上路.多次戰役,多有死傷,多番掙扎,多有波折,Frodo終於完成授命,把魔戒送回魔戒原地被毀於火湖中.

一個人如何裝身作靈戰?在神眼中,我們每一個信徒都好寶貴.

我還記得少年團契的團訓,提摩太前書4章12節:

「不 可 叫 人 小 看 你 年 輕 , 總 要 在 言 語 , 行 為 , 愛 心 , 信 心 , 清 潔 上 , 都 作 信 徒 的 榜 樣 .」

一個人如何裝身作靈戰?我在此提出「一二多全一」(One, Two, Some, All, One)的五部曲靈修法.

一是要儆醒不倦,就是敏感於邪與惡的誘惑,多面向正與善的力量,並建立「抗邪力」──對邪與惡的抗拒力(resistance to temptation of evil).人,不論大中小,都好易走入魔道.

我不過是個小兵,那麼就成為邪魔的嘍囉.

我其實是名大將,那更易成為黑帝一人之下萬人之上的黑將軍.

我不高不低,也好容易落入正邪不分的灰色地界,終日陷溺在永劫的無間道之中.

二是要祈禱代求,最好有互相代禱的守望者,二人同行.

一個信徒,當然可以建立多個二人同行的守望配對(Triple P / PPP: prayer partners in pairs).

在家庭,在工作間,在教會,在朋友圈,都最好有不同的配對,生命同行.

三是要入群入組.有一年我們曾經在自己教會推行「人人有群或有組,信仰生活會更好」的運動,就是希望信徒不要落單.就算沒有入團契,沒有參與小組的,也多參加飯團,教會有許多「飯頭」,你只要肯跟.因為我們需要fellowship of the faith.

四是常念教會整體.之前提及,許多解經家都指出這副軍裝有六部份,但首五部份是防衛性的(defensive),單單第六部份是攻擊性的(offensive).

防衛性的軍裝好像是講個人的事.其實不會,因為若然有一個人穿上全副軍裝獨自跑到街上,路人會認為他不是賣藝的就只是cosplay,否則就一定是「痴線的/short short地」(神經病).

軍裝總是集體的呈現.穿上全副軍裝就是整體戒備,每個人做好自己的分內事,整體才會也　就會完好.

許多解經家都認為這段經文是講教會的,教會有責任裝備每位信徒,為的是「助靈作戰」.

五是返回到一,「一二多全一」的最後的一,就是在靈裡合一,時刻一起在靈內作戰.

先要做好自己的靈修,讓我自己成為屬靈人.

難為正邪定分界,也得裝身作靈戰.阿們.

(Amidst unclear darkness, be prepared as spiritual fighter)



\newpage

\section{智慧的傳承}
\label{sec:423}
\textbf{賴品超牧師 - 智慧的傳承}
\newline
\newline
\hyperref[sec:422]{< < < PREV SERMON < < <}
~
\hyperref[sec:toc]{[返目錄]}
~
\hyperref[sec:424]{> > > NEXT SERMON > > >}


箴言 3:1-12
\newline
\begin{longtable}{cl}
\hline
\hline
章節 & 經文 (和合本修訂版)\\
\hline
3:1 & \begin{tabularx}{0.7\textwidth}{X} 我兒啊，不要忘記我的教誨，你的心要謹守我的命令， \end{tabularx} \\ \\ \relax
3:2 & \begin{tabularx}{0.7\textwidth}{X} 因為它們必加給你長久的日子，生命的年數與平安。 \end{tabularx} \\ \\ \relax
3:3 & \begin{tabularx}{0.7\textwidth}{X} 不可使慈愛和誠信離開你，要繫在你頸項上，刻在你心版上。 \end{tabularx} \\ \\ \relax
3:4 & \begin{tabularx}{0.7\textwidth}{X} 這樣，你必在神和世人眼前蒙恩惠，有美好的見識。 \end{tabularx} \\ \\ \relax
3:5 & \begin{tabularx}{0.7\textwidth}{X} 你要專心仰賴耶和華，不可倚靠自己的聰明， \end{tabularx} \\ \\ \relax
3:6 & \begin{tabularx}{0.7\textwidth}{X} 在你一切所行的路上都要認定他，他必使你的道路平直。 \end{tabularx} \\ \\ \relax
3:7 & \begin{tabularx}{0.7\textwidth}{X} 不要自以為有智慧；要敬畏耶和華，遠離惡事。 \end{tabularx} \\ \\ \relax
3:8 & \begin{tabularx}{0.7\textwidth}{X} 這便醫治你的肉體，滋潤你的百骨。 \end{tabularx} \\ \\ \relax
3:9 & \begin{tabularx}{0.7\textwidth}{X} 你要以財物和一切初熟的土產尊崇耶和華， \end{tabularx} \\ \\ \relax
3:10 & \begin{tabularx}{0.7\textwidth}{X} 這樣，你的倉庫必充滿有餘，你的酒池有新酒盈溢。 \end{tabularx} \\ \\ \relax
3:11 & \begin{tabularx}{0.7\textwidth}{X} 我兒啊，不可輕看耶和華的管教，也不可厭煩他的責備， \end{tabularx} \\ \\ \relax
3:12 & \begin{tabularx}{0.7\textwidth}{X} 因為耶和華所愛的，他必責備，正如父親責備所喜愛的兒子。 \end{tabularx} \\ \\
[1ex]
\hline
\hline
\end{longtable}
講題:智慧的傳承　Transmission of Wisdom

經文:箴言3章1-12節

經課:申命記4:1-2, 6-9,詩篇15:1-5,雅各書1:17-27,馬可福音7:1-8, 14-15, 21-23

講員:賴品超牧師

場合:香港中文大學崇基學院禮拜堂

日期:2018年9月2日

今天的經課,給我一點困惑.申命記4:1-2, 6-9,詩篇15:1-5,雅各書1:17-27似乎都強調行道,外在行為的聖潔,甚至要嚴守一切的律法.馬可福音7:1-8, 14-15, 21-23,某些宗教人士批評主耶穌的門徒沒有嚴格按照一種「聖潔」的方法去洗手,而這種聖潔的洗手方法,據說是指:指尖向上,澆水,讓水由指尖流到手腕,重覆共七次.(當然目的主要是宗教上的潔淨,不同於香港醫護界的建議:用潔液搓手,再唱兩次“happy birthday”.)

主耶穌沒有直接的講應否採用這種洗手的方法,而是採用轉移視線的技巧, 提出不必拘泥於外在行為上有關潔淨的禮儀的細節,而是應更多關注內心的潔淨,甚至反客為主,提出對人的遺傳的批判.那麼,主耶穌的回應是否泛指外在行為完全不重要,甚至律法也不重要?似乎也不是.問題的關鍵在於,是否本末倒置,重視人的遺傳而忽視上帝的誠命,是否「拘守」外在的細節而忽視更重要的內心潔淨.

我的困惑是,人的遺傳與上帝的道(或上帝的誠命)是否及／或如何嚴格區分? 「拘守」與否又如何劃分?更具體的問題是:作為牧師及神學教育工作者,我所教導的,不太可能完全局限在聖經／上帝的話語,例如也會有教會宗派的傳統與歷史,不同的神學流派,靈修傳統,不同的釋經方法,甚至一般的治學方法或論文寫作技巧.我教的這些,似乎都是人的遺傳.明天開學,我要問:是否仍要教這些?為何教這些?如何教?

申命記4: 6提及,在外人眼中,以色列人謹守律法其實是一種智慧.這給我另一個思考的角度--智慧.這就是說,可以用「智慧」來概括我們所承接及傳遞的.因此,今天另選一段來自箴言的經文來分享「智慧的傳承」.

在聖經中,箴言是很特別的一卷.

就作者而言,有一部份宣稱是來自所羅門王的箴言,但有些沒有那麼明確；例如按照25:1所說,第25-29章應也來自所羅門,但卻是在後來的猶大王希西家身的人謄錄.更有些似乎並非以色列人,例如第 30章是亞古珥的真言,第31章來自利慕伊勒王的真言(更準確說是來自他的母親,但沒有具體名字).

就文體而言,基本上是由一些看似有智慧的語句組成,沒有明顯的歷史背景或起承轉合的故事,不同的語句之間用語有點重覆,邏輯性不強,有較大的詮釋空間.

就內容而言,有些內容似與狹義上的「宗教」無關,不少是關於日常生活不同方面,包括投資,持家,交友,也有些是對宇宙人生的觀察.

這種以「智慧」的概念為基礎,以「箴言」的形式來表達,並不違反律法與誡命,只是它不以法律與誡命那種比較權威性與刻版的形式表達；它可以涉及生活中的不同領域,但又容讓讀者有想像的空間,也更能落實在一同的處境中.(其實,

箴言的「導言」(和合本似以1: 1-6為導言,但解經家多認為1:1-7)簡單地交代了:作者(1:1)＋寫作動機,功能,適用範圍,如何使用(1:2-6)＋總綱(1:7).據1:1,作者是所羅門,具有王者的尊榮,但卻是以父親的心腸,甚至語調去寫 ,例如1:8, 2:1, 3:1 都是以「我兒」開始,4: 1用「眾子」.有趣的是,在內容上,它卻不是所羅門王的治國秘笈傳給繼位的兒子.箴言並不是教君王或儲君如何運用各種權術計謀,使人民敬畏君王,而是教人敬畏上帝(1:7),使人領受仁義,公平,正直的訓誨(1:2),目的是讓不同人得智慧,包括:愚頑人,少年人,智慧人,聰明人(1:4-5).簡言之,是要使所有人都明白箴言,智慧的言語；而這些箴言,智慧的言語,可能有很多,包括一些不知來源,易使人困惑(例如:「古語」有云:人無夢想與鹹魚何異?--實為現影對白),因此需要指出重點與目的--總綱是敬畏上帝.(正因如此,即使是來自外邦人的智慧也可以吸收,只要能符合敬畏上帝的原則與目的.)

導言之後的部份1: 8-31; 2:1-22, 3: 1-12,基本是在進一步發揮甚麼叫「敬畏上帝」,今特選 3: 1-12 來細讀.當中可再細分數小段:

對誡命,法則

3:1  and  2  --反面地,說不要忘記「我的」法則；正面講,心中要謹守「我的」誡命.不只是大腦中記得(知性),而是心中謹守(認同).留意「我的」,不是我所拒絕或懷疑的,強加在你身上,或許不是甚麼放諸四海而皆準的絕對的定律,也不一定是甚麼獨家的驚人大發現,但這是我自己所「承」接,實踐(信受奉行),總結了我的經驗,發覺這是很好的,現在要「傳」給你(我兒),是為了你的好處,就是「長久的日子,生命的年歲,平安」.

對待別人

3: 3  and  4 -- 待人以 慈愛及誠實,不單刻在心中,更是掛在頸上(如十架般),不可離身,為的是提醒自己,也是要讓人看見.慈愛及誠實都是美德,在上帝面前蒙恩寵,有聰明(3:4)這是可理解,但為何會在世人眼前也是蒙恩寵,有聰明? 我們的電視,電影甚至社會文化,有時會質疑這些美德.慈愛?人不為己天誅地滅,婦人之仁,對敵人仁慈=對自己殘忍,弱肉強食的世界,他們窮是因為懶不值得可憐…..誠實?忠忠直直終需乞食,做人要學識語言「偽術」,懂得欺騙才是聰明…..誠實 ,即可信靠.曾有做生意的弟兄分享,吃虧也要履行合約,為的是信譽,長遠更有利.電影橋段,今天出賣你的朋友,明天你也會被出賣.據我有限的觀察,你對人沒有慈愛,別人也如此待你.慈愛包含寬恕；試想一下,如果有朋友偶然得罪你,你卻絕不寬恕,從此絕交,可能剩下沒多少朋友的.據一些與博奕論(game theory)有關的囚徒困境(Prisoner’s dilemma)的研究,最能建立長遠合作的最佳策略是「寬大的一報還一報」(Generous Tit for Tat),就是先對別人好,別人正面回應,自己也正面回應,如果別人背叛,不一定立即報復,而是偶然一次寬恕.這種策略,正是誠實加上慈愛(內含寬恕).

對神

3: 5-8--專心仰賴／信靠神.反面來說是不可倚靠自己的聰明,而是在「一切」所行的事上認定神,並且是「專一」的倚靠神,不是時靠時不靠,有些事靠有些事不靠.不反對有聰明,智慧,但不要倚賴自己的聰明多過神,甚至自作聰明,為倚靠神的事加上限制:「屬靈之事(教會生活)倚靠神, 私人的私事或公司的公事要靠自己.」.具體而言,要遠離惡事.有時人犯罪正是因為自以為聰明,有智慧,夠定力,可以試一試.不少吸毒的,犯法的也是如此想:我如此聰明不易被扶抓,我意志力堅強不易上癮.最佳的預防:遠離 -- 越遠越好.不是無膽,愚蠢,軟弱,而是因為敬畏神,是智慧的表現；十分有福,因為神必指引前路,不易掉進某些與惡事有關陷阱,最少身體可以健康:「醫治肚臍,滋潤百骨」.

3: 9-10--尊榮神,尤其指對神的奉獻(獻上財物及初熟土產).換言之,是要過一個樂於奉獻,施予,付出的人生,以榮耀上帝；那就可以富有,現在擁有很多(倉房充滿有餘),更是有可持續的富足(新酒盈溢).這似有違常理,奉獻了就是少了.經文不是說,越多奉獻,越多回報,越有錢,也不是說奉獻的人有錢過不奉獻的人.只說,尊榮神的,奉獻給神的人,也可以有豐足,無憂的人生.正如休耕,休漁期, 表面上或短期內是少了,但可以更長遠地擁有,更具可持續性,更少些擔心；相反,惡性爭捕,提早耗盡.施比受是否一定更為有福,很難比較,但施予者肯定有福.曾有學者研究,對照之下,樂於助人,包括捐獻的人,人生更快樂,更長壽,更能抗逆.(牧師,傳道人不可怕醜,講道要定期講奉獻,因這是讓人得福的途徑.)然而,是否如此行就一定平安大吉,順風順水?不一定,人生總有逆境.如何面對?(這也可以算是與上帝的關係的第三方面.)

面對人生境遇

3:11-12-- 重提「我兒」,不要輕看上帝的(對兒子般的(管教.試想,主耶穌雖是上帝的兒子,也是在苦難中學學會順從(希伯來書5:8,12:1-13).神的管教,責備出於愛,視逆境為上帝的管教,不必懷疑神的愛,更不必馬上便怨天尤人,反而更多自省是否自己也出了問題.是因為沒有信實,慈愛,沒有尊榮神?信徒相信受苦是有意義,不是純屬偶然,嘗試找當中的美意,這有助抗逆.信靠神的人生,不保證一帆風順,天色常藍,花香常漫,而是也會遇困難,當中要更信靠神,不一定能克服或跨過困境, 最少有助忍受逆境與苦難.

讀這段經文提醒,教會需要將敬畏神這種智慧傳承下去.如何?

1. 動機／心態:年輕時較難代入以父母的心情教學生或會眾,但我們需要有為父母的「心腸」去教,不一定在形式上以某種父母的口吻去「教訓」兒女,更不是「虎媽式」嚴控兒女以成就我的旨意,讓父母得榮耀,而是成全兒女的生命,為他們的好處.你的「關顧」是出於真誠的愛,還是只是工作／職業上的需要,年輕人對此也可以很敏感.

2. 方式:正如這段經文,不是權威式的說:不用問理由,你只管照著行；而是要講「理」,既講要怎樣做,也講有甚麼益處:因為會有長壽,平安(3:2),可以在上帝和世人眼前蒙恩寵,有聰明(3:4),健康(3:8),富有(3:10)；並且是以經驗分享的形式,讓人自行反思領悟.這或許是當代人,尤其年輕一代比較容易接受的方式.

3. 內容:專業知識(knowledge)或技能(skill),這些有可能很快會過時,不再適用；而人生的態度(attitude)與價值(value),可以是一生受用,可能傳承更久.可惜的是,我們的學校,包括大學教育,太注重知識或技能,因容易量化考核.那麼,教會的培育又怎樣?最重要的是(較易量度的)聖經知識與事奉技巧?還是(較難量度的)敬畏上帝,慈愛,誠實?

4. 影響:教育可以成為獨裁統治的工具,是「進化了的獨裁者」(來自《獨裁者的進化》一書)甚或「完美的獨裁者」(來自《完美的獨裁》一書)所重視.有如中國的科舉制度,令讀書人常常支持統治階層,缺少對獨裁帝制的批判性反思.香港過去也是積極吸納社會精英進入政府的行政體系,公務員即使如何「hea 做」也不會反(米飯班主)政府,而教會學校過去在製造優秀的公務員上頗為成功.問題是,這是我們所最想見到的嗎?我們傳承的智慧,是讓人敬畏神,信靠神,最少不是自以為是,自以為聰明.獨裁者卻往往是最自以為是,自以為聰明,覺得其他人,一般人蠢,民主必大亂,最好由我這最聰明的人去獨裁管治,但這卻很可能是災難的開始.中國法家的韓非子(如:〈難勢〉)或馬基雅維利(Machiavelli, 1469-1527)的《君王論》(1532),教君主如何有效施行獨裁,鞏固王權.所羅門所要傳承的智慧,不是教人如何獨裁統治,而是教大多數的人如何敬畏上帝,以誠實及慈愛待人,使人領受仁義,公平,正直的訓誨,但願這種智慧的傳承能夠造福社會.



\newpage

\section{來跟從我}
\label{sec:424}
\textbf{劉國偉先生 - 來跟從我}
\newline
\newline
\hyperref[sec:423]{< < < PREV SERMON < < <}
~
\hyperref[sec:toc]{[返目錄]}
~
\hyperref[sec:425]{> > > NEXT SERMON > > >}


約翰福音 11:1-16
\newline
\begin{longtable}{cl}
\hline
\hline
章節 & 經文 (和合本修訂版)\\
\hline
11:1 & \begin{tabularx}{0.7\textwidth}{X} 有一個患病的人，名叫拉撒路，住在伯大尼，就是馬利亞和她姐姐馬大的村莊。 \end{tabularx} \\ \\ \relax
11:2 & \begin{tabularx}{0.7\textwidth}{X} 這馬利亞就是那用香膏抹主，又用頭髮擦他腳的；患病的拉撒路是她的弟弟。 \end{tabularx} \\ \\ \relax
11:3 & \begin{tabularx}{0.7\textwidth}{X} 姊妹兩個就打發人去見耶穌，說：「主啊，你所愛的人病了。」 \end{tabularx} \\ \\ \relax
11:4 & \begin{tabularx}{0.7\textwidth}{X} 耶穌聽見後卻說：「這病不至於死，而是為了神的榮耀，為要使神的兒子藉此得榮耀。」 \end{tabularx} \\ \\ \relax
11:5 & \begin{tabularx}{0.7\textwidth}{X} 耶穌素來愛馬大和她妹妹，以及拉撒路。 \end{tabularx} \\ \\ \relax
11:6 & \begin{tabularx}{0.7\textwidth}{X} 他聽見拉撒路病了，仍在原地住了兩天， \end{tabularx} \\ \\ \relax
11:7 & \begin{tabularx}{0.7\textwidth}{X} 然後對門徒說：「我們再到猶太去吧！」 \end{tabularx} \\ \\ \relax
11:8 & \begin{tabularx}{0.7\textwidth}{X} 門徒對他說：「拉比，猶太人近來要拿石頭打你，你還再到那裡去嗎？」 \end{tabularx} \\ \\ \relax
11:9 & \begin{tabularx}{0.7\textwidth}{X} 耶穌回答：「白天不是有十二小時嗎？人若在白天行走，就不致跌倒，因為他看見這世上的光。 \end{tabularx} \\ \\ \relax
11:10 & \begin{tabularx}{0.7\textwidth}{X} 人若在黑夜行走，就會跌倒，因為他沒有光。」 \end{tabularx} \\ \\ \relax
11:11 & \begin{tabularx}{0.7\textwidth}{X} 耶穌說了這些話，隨後對他們說：「我們的朋友拉撒路睡了，我去叫醒他。」 \end{tabularx} \\ \\ \relax
11:12 & \begin{tabularx}{0.7\textwidth}{X} 門徒就說：「主啊，他若睡了，就會好的。」 \end{tabularx} \\ \\ \relax
11:13 & \begin{tabularx}{0.7\textwidth}{X} 耶穌說這話是指拉撒路死了，他們卻以為他是指通常的睡眠。 \end{tabularx} \\ \\ \relax
11:14 & \begin{tabularx}{0.7\textwidth}{X} 於是耶穌就明白地告訴他們：「拉撒路死了。 \end{tabularx} \\ \\ \relax
11:15 & \begin{tabularx}{0.7\textwidth}{X} 為了你們的緣故，我不在那裡反而歡喜，為要使你們信。現在我們到他那裡去吧。」 \end{tabularx} \\ \\ \relax
11:16 & \begin{tabularx}{0.7\textwidth}{X} 於是那稱為低土馬的多馬對其他的門徒說：「我們也去和他同死吧！」 \end{tabularx} \\ \\
[1ex]
\hline
\hline
\end{longtable}
講題:來跟從我　Come! Follow Me

經文:約翰福音11章1-16節

講員:劉國偉先生

場合:香港中文大學崇基學院禮拜堂

日期:2018年9月9日

事件的大背景:耶穌的傳道歲月正走向尾聲

事件的直接背景:耶穌的非常好友拉撒路患重病,他的姊妹馬利亞和馬大就打發人去見耶穌

作者告訴我們,耶穌素來愛馬大和她妹子並拉撒路；作為主所愛的人,對主有所期望,屬人之常情.

耶穌在所居之地仍住了兩天,祂所掌握的時間,有時超越我們的理解! 到底拉撒路的家人,特別是同寢共食的門徒,於那兩天在想些什麼?門徒對於耶穌的行徑有什麼想法?耶穌用兩天時間讓門徒的想法發酵?對於耶穌的做法,門徒抱有疑惑,也屬人之常情.

門徒提醒耶穌:「拉比,猶太人近來要拿石頭打你,你還往那裡去嗎?」我們都傾向避開風險,實屬人之常情.

耶穌回答說:「白日不是有十二小時嗎?人在白日走路,就不致跌倒,因為看見這世上的光. 若在黑夜走路,就必跌倒,因為他沒有光.」 …耶穌回答的風馬牛不相及?祂叫門徒要把握時機,選擇光明,跟隨光走,就不會跌倒.耶穌的作為和教導,有時叫我們感到難以明瞭,這豈非信仰旅程中,並不罕有的經驗?

耶穌說「拉撒路睡了,我去叫醒他.」門徒回應「主啊,他若睡了,就必好了.」…門徒顯然不明白主的心意…我們又是否太多的先入為主,太多自己既定的想法,佔據了我們的心思?

門徒不想去冒險實屬人之常情…何必要冒險?「穩穩陣陣唔好咩?」但耶穌呼召我們,並非叫我們去過安逸無風無浪的人生…「好叫你們相信」…信仰之於生活,就正正是這句話所包含.

多馬的回應「我們也去和他同死吧.」,一句情緒話,甚至是發悔氣!…人與事不如己願,生活遇上風險,不想去經歷自己不喜歡的環境際遇,根本不接受眼前的種種......多馬的名句:我非看見他手上的釘痕,用指頭探入那釘痕,又用手探入他的肋旁,我總不信.

到底「來跟從我」是什麼意思?是我們期待人生從心所願,心想事成,或者至少一分耕耘一分收穫,卻可能沒有想過,原來我們期望耶穌跟隨著我們的心願走,賜福和成就我們的想法和心願,卻沒有認真背起十字架來緊緊跟從祂!

要怎樣跟從主?…Obedience before Understanding順服擺在明白以先!



\newpage

\section{異鄉人}
\label{sec:425}
\textbf{許立中先生 - 異鄉人}
\newline
\newline
\hyperref[sec:424]{< < < PREV SERMON < < <}
~
\hyperref[sec:toc]{[返目錄]}
~
\hyperref[sec:426]{> > > NEXT SERMON > > >}


彼得前書 4:1-6
\newline
\begin{longtable}{cl}
\hline
\hline
章節 & 經文 (和合本修訂版)\\
\hline
4:1 & \begin{tabularx}{0.7\textwidth}{X} 既然基督在肉身受苦，你們也該將這樣的心志作為兵器，因為在肉身受過苦的已經與罪斷絕了， \end{tabularx} \\ \\ \relax
4:2 & \begin{tabularx}{0.7\textwidth}{X} 使你們從今以後不再隨從人的情慾，只順從神的旨意，在世度餘下的光陰。 \end{tabularx} \\ \\ \relax
4:3 & \begin{tabularx}{0.7\textwidth}{X} 因為你們從前隨從外邦人的心意，生活在淫蕩、情慾、醉酒、荒宴、狂飲和可憎的偶像崇拜中，時候已經夠了。 \end{tabularx} \\ \\ \relax
4:4 & \begin{tabularx}{0.7\textwidth}{X} 在這些事上，他們見你們不與他們同奔放蕩無度的路就以為怪，毀謗你們。 \end{tabularx} \\ \\ \relax
4:5 & \begin{tabularx}{0.7\textwidth}{X} 他們必須在那位將要審判活人死人的主面前交賬。 \end{tabularx} \\ \\ \relax
4:6 & \begin{tabularx}{0.7\textwidth}{X} 為此，死人也曾有福音傳給他們，要使他們的肉體按著人受審判，他們的靈卻靠神活著。 \end{tabularx} \\ \\
[1ex]
\hline
\hline
\end{longtable}
講題:異鄉人 　The Alien

經文:彼得前書4章1-6節

講員:許立中先生

場合:香港中文大學崇基學院禮拜堂

日期:2018年9月30日

「基督既在肉身受苦,你們也當將這樣的心志作為兵器；因為在肉身受過苦的,就已經與罪斷絕了；你們存這樣的心,從今以後,就可以不從人的情慾,只從上帝的旨意,在世度餘下的光陰.因為往日隨從外邦人的心意,行邪淫,惡慾,醉酒,荒宴,群飲,並可惡拜偶像的事,時候已經夠了；他們在這些事上,見你們不與他們同奔那放蕩無度的路,就以為怪,毀謗你們；他們必在那將要審判活人死人的主面前交賬.為此,就是死人也曾有福音傳給他們,要叫他們的肉體按著人受審判,他們的靈性卻靠上帝活著.」

「受苦」是《彼得書信》一個經常出現的主題.作者在這裡進一步闡釋信徒在世上是客旅,是寄居者的主題(彼前2:11),具體地解釋了信仰對信徒生活態度應有的指導作用.

但將受苦作為一種「心志」,忍受苦難作為一種生活態度,有時或許會令人感到煩厭,甚至懷疑會不會是某些基督徒時常被人詬病的「失敗主義」,用來掩飾甚至合理化自己的失敗和軟弱.特別是對於那些在社會經濟起飛之後出生,從來沒有真正受過苦的人,更是感到難以理解.

這有點像五,六十年代成長的一代,經過艱難困苦的日子,住木屋,板間房,徙置區,用腐乳鹹菜sung 白粥；他們的父母節衣縮食供兒女讀書；沒有能力供所有孩子上學,就犧牲一個兩大的借身分證去做工廠,讓弟妹繼續讀書.這些孩子依循著教育提供的階梯,在社會上拾級而上,最後買屋,買車,請菲傭,成為新中產階層.

矛盾的是,他們自己艱苦奮進,為下一代提供了舒適的成長環境,為他們遮風擋雨,卻在不覺間同時剝奪了他們受磨練的機會,削弱了他們應付問題的能力.他們小時候所沒有的一切,他們都毫不吝嗇地買給他們的孩子.然後他們問,我們不是已經為他們提供了最理想的成長環境,怎麼他們反而欠缺學習動機,不懂得跟別人相處,情緒智商偏低,好像樣樣都不如我們那一代?歸根究底,大概是因為這些孩子根本不知道個「苦」字怎麼寫.

我們的邏輯是:如果透過更充分的準備,更豐厚的物質條件,可以避免一些無謂的痛苦,走少一些冤枉路,那不是更加划算和理想?

《彼得前書》的作者卻似乎在這裡指出,基督徒受苦是必須的,甚至應當以「在肉身受苦」作為與罪斷絕的手段,彷彿如果沒有在肉身上吃過點苦頭,所謂離罪成聖就不過是紙上談兵,得個講字.

從某個角度來看,「受苦」似乎並不是一個‘Option’,而是任何生物可以絀壯成長的必要條件.我們花錢送我們的子女去Outward Bound,讓他們暫時離開受保護的環境,攀山涉水,直接面對不徇情面的大自然,就是這個意思.只可惜一,兩個禮拜的課程之後,卻一曝十寒,無以為繼.

對於「受苦」的意義,我們可以從不同的角度去理解:哲學的,醫學的,信仰的,人道的.... 可是基督徒對「受苦」的理解,特別是對「基督為我們受苦」的理解,卻往往流於表面.

我們很多時會聽見人描述主耶穌在十字架上的痛苦,甚至好像Mel Gibson在多年前拍的「受難曲」那樣,繪形繪聲地形容那個荊棘冠冕,怎樣插入耶穌的頭皮；那條加料的皮鞭,怎樣重複抽打在主的身上,直至皮開肉裂.這背後的想法,當然是為了激發我們的感恩之心和悲憫之情,想到主既然這樣為我們受苦,我們也當將生命獻給主.這是我們經常思考,也是我們習以為常的角度.

偶然從這方面去默想主,當然不能說不對.但如果我們就只有這個向度去理解「基督為我們受苦」,恐怕我們的情緒,很快便會因為長期超載而疲勞甚至麻木,以致需要不斷增加劑量,才能產生同樣的反應.更嚴重的,是這種想法很容易令我們誤解基督受苦,受死的意義,以一種廉價的情緒反應,去取代福音的核心信息.

我們發現,恆常地運用我們的想象力,去幻想基督受苦的景象,不但讓我們感到吃力；不斷重複想象所造成的回報損耗(diminishing returns),亦令到基督受苦的景象日漸失去激勵的效用.情況就好像我們在一些特別的日子思念一些特別的人,感受會來得集中和深刻；要你天天這樣做,時時刻刻這樣做,效果可能就適得其反了.

事實上,論到刑罰的殘忍程度,中國人發明的刑罰,譬如炮烙或五馬分屍,和種種用來逼供的酷刑,就比釘十字架更加殘酷.從已經公開的資料顯示,現代人用來虐待戰俘和拷問罪犯的酷刑,也都讓人觸目驚心.

所以,基督受苦的意義,並不在于它的殘酷性.否則我們就要好像一些地方的信徒那樣,在一些宗教節日,抬著很重的十字架巡遊,用帶刺的皮鞭抽自己,甚至刀將自己割得遍體鱗傷.這些都是我們在一些新聞特輯見過的.

論到這一類宗教苦行,保羅在西2:20-23指出,「你們若是與基督同死,脫離了世上的小學,為什麼仍像在世俗中活著,服從那『不可拿,不可嘗,不可摸』等類的規條?…這些規條使人徒有智慧之名,用私意崇拜,自表謙卑,苦待己身,其實在克制肉體的情慾上,是毫無功效.」

但倘若基督受苦的意義並不在于肉體所受的痛苦,那麼它又怎樣成為我們與罪斷絕的兵器呢?

彼得前書 3:18「因基督也曾一次為罪受苦,就是義的代替不義的,為要引我們到上帝面前.按著肉體說,他被治死；按著靈性說,他復活了.」

彼得在這裡指出,基督受苦的意義,並不在于祂在肉體所受的痛苦.基督受苦的意義,是「義的代替不義的,為要引我們到上帝面前」.

這裡牽涉一個很重要的教義,是往往被我們掛在嘴邊,但卻是知其然而不知其所以然:知道是這樣,卻不知道為甚麼是這樣:「因信稱義」.是否真正認識這教義,對於我們能不能「從今以後,可以不從人的情慾,只從上帝的旨意,在世度餘下的光陰」息息相關!

當年馬丁路德看羅馬書的時候,重新發現「因信稱義」的道理,結果掀起劃時代的宗教改革運動.可惜正如牛津學者Alister McGrath曾經指出,「因信稱義」這教義是一條神學上的恐龍,曾經在它的時代顯赫一時,可是作為一個迫切的神學課題,卻早已經絕了種.它變成一件博物館裡的古董,被學者觀賞和研究,但跟現代的信仰生活已經拉不上關係.

這是我非常同意的.事實上這也是我作為一個信徒多年來的問題:倘若「因信稱義」真是那麼重要,它又為何那麼抽象難明?當然,在理論上,這教義實在簡單得難以置信:倘若你相信上帝,你的罪就立時歸到那位代替你死而復活的基督身上.

複雜點用路德的「交換理論」:基督將他的「義」賜與相信他的人,而信徒的「罪」就轉移到耶穌基督身上.這在理論上是非常簡單淺白的.問題是面對著這樣的講法,你要就信,要就是不信,並沒有太多可以理解或討論的餘地和空間.

而這問題並不是無關痛癢的.

我們常常在教會聽到「得救確據」的問題:「我如何肯定基督其實也為我死?」,而所謂「基督為我死」的確實意思又是甚麼?我怎知道自己真的「得救」?當然,一般的答覆,往往是說「這正是我們需要行使信心的地方!」只是當我們孤注一擲地「信」了之後,那種患得患失和一廂情願的感覺,卻總是揮之不去.而「因信稱義」這個事實,除了影響我日後所謂的永恆歸宿,在實存意義上又跟我每日的生活行止有甚麼關係?

這是我們很難不去問的問題.

要解決這個信仰困局,我相信我們首先需要克服一個重要的文化語言的問題:據我了解,對於兩岸三地的中國人來說,「稱義」這個詞,根本不存在於我們的日常語言之中.

除了在教會聽道之外,你幾時聽過有人用「稱義」這兩個字?我們絕對不應低估這個問題的嚴重性,因為不存在於我們語言系統的,大概亦很難在我們的意義架構佔一個位置,或對我們的日常生活,抉擇產生甚麼影響.這也直接解釋了我們的信仰有時難於「落地」的原因.

「稱義」(Justification) 在中文的詞彙中,其實並不是一個最好或者唯一的翻譯.困難卻在於似乎沒有一個簡單,現成的現代詞彙,可以完全包含這個詞原本的概念和意思.

按照字典的解釋,「稱義」(justify)這個詞,一般是指「證明清白,合理或正確」的意思.很明顯這個詞彙在神學甚至一般用途上的「法理」味道,跟這個字典的解釋有很強的相關性.而我進一步觀察到,所謂生命意義的追尋,在一個頗大的程度上,其實就是企圖去找出以致證明自己的生命,自己的存在,是清白,正確和合理.做人,最緊要是頂天立地,仰俯無愧.

這至少在心理學的層面,對人類行為有相當強的解釋能力.事實上田立克(Paul Tillich)就曾經嘗試將「稱義」(being justified)這個概念現代化為「被接納」(being accepted),去探討「因信稱義」在心理治療上的涵義.

而如果我們停下來誠實地自我省察一下,可能我們也會驚訝地發現,我們的生活是否過得充實,在一個頗大的程度上,乃決定於我們「自我解釋」的能力.尼采的名言:“He who has a why to live for can bear with almost any how.” (「一個有目的,理由而生存的人,幾乎可以忍受任何處境.」)因此,「稱義」或者「證明為正確或合理」的教義,事實上是構成了人類道德價值一個重要的核心.

透過重新解讀「因信稱義」,我們可以真正肯定生命的意義和價值:基督的受苦,弔詭地顯示生存可以有的勇氣,他說「我父作事直到如今,我也作事」；基督的受苦,展示和證明一個人是可以堅持遵行天父的旨意,甚至去到一個死的地步.他的復活,表明上帝認可一個這樣完全投靠祂的生命.透過耶穌基督的死和復活,我們知道自己不必為眼前的利害得失而折腰,亦無須為了要證明自己而惶惶不可終日.

因此,從「因信稱義」的角度來看,基督曾一次為罪受苦,叫我們可以與生命的創造主重新和好；透過他,我們認識那位差他來的.

回到「稱義」的問題.面對自己的一生,一個人怎樣可以得稱為義?或者說,一個人怎樣去justify自己的一生是有價值,有意義呢?並不是靠他作任何事可以證明.

用一個簡單的例子:一個角色或者一件事怎樣可以變得有意義?並不是那個角色或那件事本身可以做些甚麼令自己變得有意義；你只要將那個角色放回那原本盛載它的故事,將那件事放回它原本所屬的情節便是了.

當一個角色抽離了盛載他的故事,那個角色,就變成一個無所依賴的游魂；當一個橋段離開了整個故事的context,那個橋段,本身無論如何精彩,就只是一堆支離破碎的事件:有gag,沒有劇情.

一個角色的意義,是完全被它所屬的故事所決定；一個故事的意義,是完全取決於作者對故事的構思和鋪排.換句話說,不是我們做甚麼去證明我們的價值,而是我們純粹憑信歸回我們原來的本位.

沒有作者的觀點作為定位,一個故事就無法發展下去.情況就好像一個人站在銅鑼灣或旺角的街頭,看見整條街人來人往,人山人海,卻不知自己要往哪裡去,或者這個外面的世界跟他有甚麼關係.那種感覺,就像卡繆(Albert Camus)筆下的「異鄉人」,肉身還是生活在這個世界,心靈卻好像飄浮在一個不屬自己的異域,不知自己到底是誰,想做些甚麼,感覺生命並不真實.

這種感覺絕不抽象,對於上個世紀的中國人感受尤其深刻.由抗日戰爭中國人被日本人屠殺,到國共內戰中國人互相廝殺；由所謂「抗美援朝」派百多萬散兵游勇去北朝鮮做砲灰,到1954年毛澤東告訴當時印度總理尼克魯他不怕原子彈,「中國那麼多人,炸不完.炸死一千萬兩千萬算不得甚麼」；由三反五反,反右運動,大躍進,文化大革命到台灣的二二八事件,軍事戒嚴,白色恐怖；由中國人終於站起來,到鍾祖康話「來生不做中國人」.我們這一代戰後出生的中國人,恐怕仍然難以擺脫我們上一代的那種失落感.

面對這種強烈的不踏實的疏離感,人們必須以感官的刺激,剎那的光輝,物質的成就,去掩蓋生命的虛空,存在的焦慮.就好像保羅描繪那些對生命沒有盼望的人那樣:「讓我們吃吃喝喝罷,因為明天要死了!」有時,消極負面的精神狀態,亦可能產生表面進取,裡面空洞的繁華景象.

從這方面來看,生活在這個荒謬世界而有上帝形象的,其實都是精神,靈性上的異鄉人.這個世界並不能讓他們得到真正需要的.正如彼得在這裡指出,他們必須透過「行邪淫,惡慾,醉酒,荒宴,群飲,並可惡拜偶像的事」,去壓抑和麻痺心靈的痛苦和空虛.他們是陷墮在罪的惡性循環之中；他們犯罪也同時被罪壓逼,以致不能自制地繼續犯罪.

而相對於這些迷失的罪人,那些透過道成肉身的基督,重新確定自己身分的人,在這個以失落和犯罪為常態的社會,也是一班有時不知所措的異鄉人.意思是他們當然不再需要隨從往日「外邦人的心意,行邪淫,惡慾,醉酒,荒宴,群飲,並可惡拜偶像的事…」；但昔日的同道中人,正正因為缺乏這個視野和醒覺,所以當他們看見這班曾經熟悉的老友竟然「不與他們同奔那放蕩無度的路,就以為怪,甚至毀謗他們」,以為他們只不過是在裝胸作勢,自鳴清高.這班人的清醒,凸顯了他們的放縱和迷失.這是他們難以忍受的.

因此,面對這個格格不入的世界,信徒有時亦會感到迷失.希伯來書的作者強調,基督徒在世界不過是客旅,是寄居者,是異鄉客,是局外人.彼前1:1開宗明義就說:耶穌基督的使徒彼得,寫信給那分散在本都,加拉太,加帕多家,亞西亞,庇推尼寄居的…」這世界並不是我們的家鄉!

這裡還有一個有待處理的問題:「為此,就是死人也曾有福音傳給他們,要叫他們的肉體按著人受審判,他們的靈性卻靠上帝活著.」從字面去理解,這段經文明顯跟聖經其他部分格格不入.

“For this is the reason the gospel was preached even to those who are now dead, so that they might be judged according to human standards in regard to the body, but live according to God in regard to the spirit.”(v.6, NIV)

「因為這個緣故,福音曾經傳給那些甚至已經不再人世的,以致他們的肉身縱然按著人間的標準受到審判,他們的靈性卻是向上帝而活著.」

這裡的意思是:在肉身方面,這些已經不在人世的,確實不能逃避承擔他們在蒙昧無知的時候所種下的惡果,得到他們應得的審判；但在靈性方面,他們卻因為覺悟上帝的權能,看見他們的生命故事,在上帝手中有不同的可能.

這確實也是我們的經驗.無論我們在甚麼階段信主,已經在我們生命中發生的,就已經發生,我們不能夠改變.這些已經發生的事,很多時令我們感到痛苦和後悔.我們寧願它們從來沒有發生.但已經發生的,沒有人能夠改變.我們只有承擔這些事所帶來的後果.在這個意義上,可以說我們的肉身,是按著人間的標準受到審判.但從一個屬靈的角度來看,我們的過去,一切的成敗得失,透過永生的盼望,卻賦予了新的意義.

彼得這個說法,叫我想到對於一個生活了七,八十年的人,無論他們曾經生活得如何糊塗混噩,我們都不應假設他們對生命一無所知.並且如果你告訴他們,認識基督之前的生命都是空白一片,完全沒有意義,那對他們來說是惡耗而不是福音:往者已矣,覆水難收；是好是歹,他們的生命已過了一大截,難道要他們再鑽入母腹重新做人?這是尼哥底母在約翰福音曾經問耶穌的問題.

聖經在這裡給我們的亮光是:福音是告訴他們,透過那位宇宙創造主成為肉身來到這個世界,我們知道生命是確實的,並不是一場夢,無論是美夢還是惡夢.而生命既然不是一場夢,他們曾經成功或失敗的努力,似乎是白費或沒有白費的心血,都被記念；他們所有發自心靈的懊悔,都被聽到；所有的眼淚,都被上帝知道.

透過為世人罪孽捨身的基督,他們那邁向夕陽的生命可以得到肯定；他們可以用一個新的眼光,去重新閱讀自已的生命故事,並且在當中看見上帝那雙無形的恩手 ─ 就在他們還是愚昧無知的時候 ─ 就一直在拖帶著他們.在他們面前的,並不是「等收工」的風燭殘年,而是一幅接近完成的圖畫,一個邁向結局的故事.這,對他們來說才是真正的好消息,真正的福音.

或許,這才是「他們的肉體按著人受審判,他們的靈性卻靠上帝活著」的意思.

亦唯有面對這個事實,我們才可以坦然無懼地面對命運的不確定性,知道真正的問題,並不在於生命的不確定,而在於我們根本不知道怎去面對充滿可能性的生命；亦只有當我們的生命有一個重心,我們可以面對我們自己的愚昧和缺失,和我們曾經錯失種種的機會,知道每一個錯誤的選擇,背後都有一個新的可能!這是一個跟我們實際生命息息相關的真正福音!

可惜的是,這個天大的喜訊,落到我們的手,卻往往變成了一個不痛不癢,跟真實生命毫不相關的宗教課題.

讓我們再一次被提醒:「基督既在肉身受苦,你們也當將這樣的心志作為兵器；因為在肉身受過苦的,就已經與罪斷絕了；你們存這樣的心,從今以後,就可以不從人的情慾,只從上帝的旨意,在世度餘下的光陰.」



\newpage

\section{給我一個笑臉}
\label{sec:426}
\textbf{吳炳華牧師 - 給我一個笑臉}
\newline
\newline
\hyperref[sec:425]{< < < PREV SERMON < < <}
~
\hyperref[sec:toc]{[返目錄]}
~
\hyperref[sec:427]{> > > NEXT SERMON > > >}


講題:給我一個笑臉  　Give Me a Smiling Face

經文:詩篇42,43

講員:吳炳華牧師

場合:香港中文大學崇基學院禮拜堂

日期:2018年10月7日

引言:在人與人的接觸或相處裡面,笑臉往往使人有一種親切,歡迎,仁慈的表示.給人有一份鬆弛,安舒,釋放,開心,溫暖,得鼓舞,激勵的感受.

聖經翻譯者用這種擬人的筆法來強調詩人心中的渴求,詩人此刻最想見是神的臉,其實,就是神自己,是神的臨在.

(因為從來沒有人見過神的面)

什麼原因?

為什麼詩人咁想見到神?

在進入內容之先,先掌握這兩篇詩的結構

大部份的聖經學者都同意42,43原屬一篇,分開最有可能是因禮儀程序上的需要.而且明顯的理由,是它們有一個很清的分段分隔每一段都是由一句幾乎是完全相同的句子,像詩歌的副歌一樣42:5,11,43:5

這三段也表達出詩人在困苦中所經歷的三個階段.每一個階段都有它的重點和主題.

讓我們回到剛才的問題:詩人為何咁想見到神?

I  枯乾   第一階段(圖畫)   <你的神在那裡?>

詩人首先就用一隻乾渴至幾乎喪命的鹿,亂咁跑來跑去要急於尋找水源,來喻作自己心底裡對神的深深渴慕之情.因為他的心田好像一幅龜裂的地土.他的心靈十分枯乾.

從詩人所發出的問題中,可能比到一些線索,點解詩人咁切慕上帝.

有學者猜測他是被放逐,要離鄉別井,遠離家園.離開了他的群體,遠離神的聖殿,又或者是因身患重病而無辦法去到耶和撒冷敬拜神.

「朝見上帝」就是指到聖殿敬拜神

昔日的敬拜(聖殿)與今日,我相信有很大的分別,

昔日:聖殿乃代表著神的居所,神的臨在,

在那裡:上帝與人相遇

上帝與人說話

上帝救贖以色列人的事件被記念

以色列人獻呈他們的禱告,並得著上帝的應許

以色列人因上帝應允他們的禱告而獻上感恩和讚美的祭.

聖殿就是他們的生命與上帝接觸的重心.

但此刻對詩人來講,回到聖殿去敬拜神似乎是一件無可能發生的事.

他已經有一段頗長的日子經歷不到與神親密之接觸.神似乎離他很遠很遠.

當他想起昔日他帶領會眾一同在聖殿守節,那興奮熱鬧的場面,他就更難過.

這種情況連敵人也察覺得到,因而譏笑他“你的神在哪裡?”

這句充滿敵意又深具破壞信心的話.言外之意是:指詩人所信的神若不是無情地毫不關心他的死活,就是根本沒有神的存在.

這真是令詩人感到痛苦難堪而又無言以對.

就惟有以眼淚當飲食苦不堪言.

詩人面對著:國破家亡,離鄉別井,受人欺負,寄人籬下,不能與自己的百姓在聖殿敬拜神,對他來講實在是有太多失去(lost)

但相信最大的損失/失去就是因找不到上帝而感受到那份乾渴.感受不到衪的臨在.似乎神離他很遠.這也是今日很多基督徒的經驗.今日,我們生活在現代化和科技進步的時代往往與神的方式有很大的差距今天我們都講求效率和方便.按一下按扭->交功課,找資料,買東西

但原來神並不是一副機器任人控制.神完全不受人的規範.

甚麼事會令人感到心靈的枯乾?

忙碌的工作(壓力)

趕急的生活節奏

平淡單調的生活

身體疾病

失去(朋友,親人)

動盪不安的環境

若我問大家的心境如何?

有否感到枯乾? 知否原因在那裡?你又怎樣面對這情況呢?

亞摩司書8:11 人饑餓非因無餅,人乾渴非因無水,乃因不聽耶和華的話

又或者今日我們不會問神“你在哪裡?”多的是神問我們“你在哪裡?”

II 被淹沒  (V6-10)   <我的神在哪裡?>

現在我們要看第二幅圖畫  (V6-10)

在這裡詩人啤我們一幅很不同的圖畫,就是他落在深水裡,都是同水有關的,但與第一幅形成一個好強烈的對比.

第一幅:詩人把自己的心田繪畫為龜裂,枯乾(無水)

第二幅:卻描述自己身陷大水裡

瀑布深淵和波浪都比喻所受的痛苦.

詩人的內心好像被厚厚的烏雲掩蓋著,他的心也像跌進「深陷的淤泥中」,越掙扎也就越深沉下去一樣.他不單只覺得神已忘記了他,甚至神己丟棄jor 他.

詩人的哀哭仍然是環繞著這一句責問的話:「你的神在哪裡?」但今次不是由別人發出,乃是來自自己內心的響聲:「我的神在哪裡呢?」.

例:唐慕華 Marva  Dawn她是一位神學家,教育家,作家

2009年她寫了一本書“在軟弱中享安好”

她在書中提到2008年 那時,她剩下一隻腳和一隻眼,幾個月之後

更確診患有癌症對她來說是終極的致命一擊,她形容自己好像被洪濤過她身般

她覺得很難相信神,當時,若不是神為她預備了很多關心她的親人,她真是站不住了.<這個見證,可能在講的時候會有增減>

III  活水的泉源   (43:1-4)

首先,我們會留意到一個很明顯的轉變,就是詩人從回憶和自述中進到向神發出呼求.從哀怨神似乎是隱藏了的聲中變為請求沉默的神作為他的審判官.

詩人是直接向神說話,他求神為他伸冤.他求神作出裁判,來伸明敵人的譏笑和辱罵是錯.上帝並沒有離開他,忘記他,放棄他.

從他的祈求中,他對神慢慢恢復信心,他相信上帝是永活的,他相信神的亮光和真理,會引導他,終有一天,他會再有機會與眾信徒在聖殿敬拜神.

這一段似乎詩人要帶出他已恢復對那永活的神的信心,有如找到活水的泉源一般充滿著盼望,並且要向敵人和自己宣告<我的神在這裡>

轉淚點在那裡呢?->就在每一段的副歌

詩人向自己的心說話,鼓勵自己提醒自己,不需要成為情緒的無助受害者,

不要讓自己的的感覺主宰你.

引導自己:

要仰望那永活的神

仰望=是一種耐性及充滿希望的等待,等待神的行動

告訴自己,讚美的日子終於會來臨

雖然,什麼時候來臨,那是照神的時間,不是你自己的時間.

結束的例子:當我知道大哥離世的消息後,我有一整個星期的時間,午飯後在房間祈禱,向神傾訴我的苦情最後,我似乎聽到祂對我說:「我明白」

這一句話對於我來講就是“神的笑臉”.



\newpage

\section{為乜你可以去到幾盡}
\label{sec:427}
\textbf{鄭漢文博士 - 為「乜」你可以去到「幾盡」?}
\newline
\newline
\hyperref[sec:426]{< < < PREV SERMON < < <}
~
\hyperref[sec:toc]{[返目錄]}
~
\hyperref[sec:428]{> > > NEXT SERMON > > >}


馬可福音 10:17-22
\newline
\begin{longtable}{cl}
\hline
\hline
章節 & 經文 (和合本修訂版)\\
\hline
10:17 & \begin{tabularx}{0.7\textwidth}{X} 耶穌剛上路的時候，有一個人跑來，跪在他面前，問他：「善良的老師，我該做甚麼事才能承受永生？」 \end{tabularx} \\ \\ \relax
10:18 & \begin{tabularx}{0.7\textwidth}{X} 耶穌對他說：「你為甚麼稱我是善良的？除了神一位之外，再沒有善良的。 \end{tabularx} \\ \\ \relax
10:19 & \begin{tabularx}{0.7\textwidth}{X} 誡命你是知道的：『不可殺人；不可姦淫；不可偷盜；不可作假見證；不可虧負人；當孝敬父母。』」 \end{tabularx} \\ \\ \relax
10:20 & \begin{tabularx}{0.7\textwidth}{X} 他對耶穌說：「老師，這一切我從小都遵守了。」 \end{tabularx} \\ \\ \relax
10:21 & \begin{tabularx}{0.7\textwidth}{X} 耶穌看著他，就愛他，對他說：「你還缺少一件：去變賣你所有的，分給窮人，就必有財寶在天上；然後來跟從我。」 \end{tabularx} \\ \\ \relax
10:22 & \begin{tabularx}{0.7\textwidth}{X} 他聽見這話，臉就變了色，憂憂愁愁地走了，因為他的產業很多。 \end{tabularx} \\ \\
[1ex]
\hline
\hline
\end{longtable}
講題:為「乜」你可以去到「幾盡」? 　For What, How Far Would You Go?

講員:鄭漢文博士

經文:馬可福音10:17-22

場合:香港中文大學崇基學院禮拜堂

日期:2018年10月14日

1.為「乜」你可以去到「幾盡」?

各位弟兄姊妹,今日我在這個講壇相見後,好大可能半年後才再見.所以,今日我會為「講道」可以去到「都幾盡」.

今日的經課係取自馬可福音十章17至31節.

在座若返教會一段時間,都會對「財主問永生」的這個事件.所以,當我準備這次講道時,我就在想可以有甚麼新意可講?有甚麼新方法去講解這段經文?

這個星期,剛好為崇基學生提供未圓湖畔系列活動,其中一項叫「念師恩」(Remembering Teachers),就是與這一代大學生分享我如何想念老師教育我的恩情.

我讀崇基本科時除了對我好好的沈宣仁院長和教我中國哲學的勞思光教授,還有何秀煌老師,他是我碩士指導老師,因為他們仨我去了牛津研究語言哲學,心靈哲學,宗教哲學和政治哲學.

何秀煌老師在他學問的後期,提出一個叫「空心詞」的思考方法.我為今天講道命題時,運用了他的方法.

起題時,係用了香港粵語來起題.

其實係借用自《狂舞派》這套電影,主角阿花狂熱於跳舞,但在比賽前傷了右腳,不良於行,只能坐輪椅.在她放棄比賽時,最有機會和她那隊比拼取勝的對手反而追問她:〈為「跳舞」你可以去到「幾盡」?〉

今天的講題係:〈為「乜」你可以去到「幾盡」?〉(For what, how far would you go?)

這個講題有兩個空心詞,每個空心詞都可以填一個詞進去.

〈為「乜」〉的「乜」是空心詞,你會填些甚麼?例如:永生,自尊,實任,上位(權位)

〈為「乜」〉英文是for what,其實是指〈為何why〉.

講題後半部問:〈你可以去到「幾盡」?〉

〈去到「幾盡」?〉是How的問題,how far, to what extent?

香港新高中的通識教育科/(Liberal Studies) LS考卷經常這樣出題:「多大程度上同意或不同意以上命題或立場(to what extent would you agree with such proposition or position…)」

有人可以〈為「謀財」而可以去到「害命咁盡」〉.

有人可以〈為「上位」(爭奪權位)而可以去到「官商鄉黑鉤結咁盡」〉.

當年我考入了中大後,我見證過不少可以去得好盡的師兄師姐.

其中一位師兄,為了「捍衛人權」他可以去到好盡──「連自己的大學學位都放低」.

之後,成長時,不斷見到好多值得我景仰的人物,都為一些重大價值而去得好盡.

最近,香港出現不少年青人,為求土地公義,為政治公義,放下自由而被入獄.

近日,我又見到一位為了「香港的民主」的中大同事,他可以去到「連大學教職都放低,提早退休」咁盡.

你呢,為了「香港好」,你可以去到「幾盡」?

你呢,你一生為了「乜ye 」?你可以去到「都幾盡」──拼命至死?

(For what in your life, how farthest would you go?)

各位,為了「永生」,你可以去到「幾盡」?

今日馬可的經文,正是一個年青人面對耶穌提的這個問題.

這段經文可以分為三個段落,可以看為都是有人問耶穌一個問題,按耶穌教學法是因心施教,應機作答,答法各有不同但可累進.

首先是起問的財主,然後是門徒,最後是彼得.

以下我們按經文的三段段落,一同思考〈為「乜」你可以去到「幾盡」?〉

2.為「永生」你可以去到「幾盡」?

讓我們再讀今日經課的第一個段落,馬可福章十章17至22節:

17耶穌出來行路的時候,有一個人跑來,跪在他面前,問他說:「良善的夫子,我當做什麼事才可以承受永生?」

18耶穌對他說:「你為什麼稱我是良善的?除了神一位之外,再沒有良善的.

19誡命你是曉得的:不可殺人；不可姦淫；不可偷盜；不可作假見證；不可虧負人；當孝敬父母.」

20他對耶穌說:「夫子,這一切我從小都遵守了.」

21耶穌看著他,就愛他,對他說:「你還缺少一件:去變賣你所有的,分給窮人,就必有財寶在天上；你還要來跟從我.」

22他聽見這話,臉上就變了色,憂憂愁愁地走了,因為他的產業很多.

我們做老師的時常要回答學生提問的問題.回答是否得當,就要留意學生的問法.

耶穌發現這個求永生之道的問道者,其實問錯問題,因為問的人認為「永生是要做事才可以得到的.」所以,他希望叩問「永生是要做咩ye 事才可以得到?」(Doing what deeds to get eternal life?)永生好像可以靠功德行善事.

他的問法是:「我當做什麼事才可以承受永生?」問的人是追問「當做何事」(a deed to be done)

所以耶穌順著這個問法而答,讓他得知錯問了永生.

首先,耶穌問他有沒有跟誡命做.

然後列出十誡的後六誡關於人的善事.

這個求道人答:「這一切我從小都遵守了.」(All of these I have already been doing.)

我們把這些所有善行之事(All Already)稱為A.

究竟A是否得永生的充分條件?(Is A sufficient condition for inheriting eternal life?)即是單靠A就可得永生.

如果不是,還有那個「當做何事」(D for deed)是必要的條件?(What D is necessary as a condition for inheriting eternal life?),做A+D就足以取得永生,缺D就不得.

耶穌指出有一個「當做何事」D構成必要條件,但原來D還有三個層次:

可10: 21「你還缺少一件:

Level 1:D1去變賣你所有的,

Level 2:D2分給窮人,就必有財寶在天上；

Level 3:D3你還要來跟從我.」

英文是:Doing one threefold action in order to get something called ‘eternal life’

廣東俗語叫:「一雞三味」.其實,「變賣所有,分給窮人,跟從耶穌」真的好難為.

你只要你自己代入想像下,你就是求永生之道的求道者,耶穌這樣叫你:

Level 1:D1去變賣你所有的──如果你現在賣哂你D樓,你也許物業好多.現在樓市係咪時候套現?我要否留一間給自己安居?

Level 2:D2分給窮人,就必有財寶在天上──係咪在街上給乞丐HKD500就可以?究竟誰是窮人?我是否看到在職貧窮?扶貧我做到幾多?

Level 3:D3你還要來跟從我──跟你,大佬,你係「乜水」?有無方向?有無路線圖?

耶穌答完永生之道後,故事這樣結局:

22他聽見這話,臉上就變了色,憂憂愁愁地走了,因為他的產業很多.

因為求永生之道的問道者的問法是:「我當做什麼事才可以承受永生?」問的人是追問「當做何事」(a deed to be done).

用了Doing mode of question: a matter of give and take - giving all in this world for taking eternal life (by the act of mankind through good deeds)

所以耶穌順著這個問法而答,讓他得知錯問了永生.

為了「永生」,這個求道者會否可以去到咁盡:「變賣所有,分給窮人,跟從耶穌」?

他連第一關都做不到:因為他的產業很多.

為了「永生」,你可以去到「幾盡」?

3.為「自力」你可以去到「幾盡」?

讓我們轉入今日經課的第二個段落,馬可福章十章23至27節:

23耶穌周圍一看,對門徒說:「有錢財的人進神的國是何等地難哪!」

24門徒希奇他的話.耶穌又對他們說:「小子,倚靠錢財的人進神的國是何等地難哪!

25駱駝穿過針的眼,比財主進神的國還容易呢.」

26門徒就分外希奇,對他說:「這樣誰能得救呢?」

27耶穌看著他們,說:「在人是不能,在神卻不然,因為神凡事都能.」

當時按希伯來人信仰,猶太人傳統,認為蒙神祝福的人就會亨通富貴,貧窮不順的人一定是蒙神咀咒的.

這種思路也好似我們今天流行的成功神學:認為蒙神祝福的人就會成功,失敗的人一定是蒙神咀咒的.

所以,門徒好驚訝,連這麼財運亨通的官家少年都入不到永生之門,上帝係咪自相矛盾?又祝福他們以致富貴起來,但又因財富阻撓他們進天國?

所以,門徒又擔心,我咁窮,憑咩入到天國?經文記載,門徒就分外希奇:「這樣誰能得救呢?」

神的救法若是以一種條件式的思維(conditional thinking),即是人必須履行一些條件C才能得救,那麼不具備這些條件C就不能得救,那麼因為這些條件C的必要的,C就會令一些人不能得救.由無這些條件C到具備了C,這些條件C就是充分條件.

剛才我們探討過:「變賣所有,分給窮人,跟從耶穌」構成這個產業很多的求道者得救入天國的必要的條件C.

這個問題可以倒轉問:「這樣誰不能得救呢?」

四十多年前自信耶穌以來,我一直有這樣的偏見,有三種人好難得救:

一是財富多的人

二是權力強的人

三是才智高的人

Wealth could be posed as hindrance to salvation (I am rich enough to do anything; there is nothing I cannot do)

我懷疑我是自小仇富,可能因為自己貧窮出身,所以我偏向以為財主不能入天國.今日這段經課,正好像給了我們一個聖經基礎.真的大快人心.

富甲天下,有錢可以試得鬼推磨.財富卻成為富人入天國的負累.

耶穌明明說:「倚靠錢財的人進神的國是何等地難哪!駱駝穿過針的眼,比財主進神的國還容易呢.」

唯有放下今生的財富,才得著永生,才積聚財富於天國.

Status (positional power) could be posed as hindrance to salvation (I am powerful enough to do anything; there is nothing I cannot do)

路加福音十章22節說這位求道人是一個官,是有政治地位的.地位,包括政治權位,是一種在位置而有的力量.例如公權力,是對公權所涉的人民福祉為依歸.一下子把超過一半以上的儲備填海起人工島,多少市民受惠?誰最得益?用的是每位香港市民約一百萬的儲備.

權力是競爭性價值(competitive good),越少人佔有就越大權能,在帝國時代是一人為王加上皇親國戚,現時在香港社會是被選上位者的一小群而非民選首長──學者稱為群帶資本主義(crony capitalism).用新馬克思主義(neo-Marxian)的觀點講,是因為權位可以壟斷生產關係以及生產資料而產生霸權(hegemony).我不能講下去了,因為在座有不少知識分子會認為我的思維未夠水平.

Intelligence (intellectual excellence) could be posed as hindrance to salvation (I am intelligent enough to do anything; there is nothing I cannot do)

我的一生除了升學直到大學,之後工作,香港中文大學是我唯一的僱主.

我其實是校園的動物,生物,但我在校園見到好多怪物,死物.

我的經驗告訴我,向知識分子傳福音,難過向財主傳福音,因為知識分子認為自己的才智足以改善世界,拯救世界,包括自己.

大學的教授,追求卓越,好容易會自大,傲慢得不會認為基督教福音值得一聽.

基督教的福音過不了知識分子的智力檢測,對於知識分子,福音只是智能不足,心靈軟弱等弱者的心理需要.

馬太福音十九章20節說這位求道人是一個少年人,幼受庭訓,好有文化德行,從小他就被教導以律法教育成長,所以,他非常自豪地回答耶穌:「這一切我從小都遵守了.」

按今日經課的思路,智能才華成了知識分子的財富,恃才傲物,不放下其實好難進入天國.

總而言之,按剛才題及這三種財/才/材三才觀:有錢財的財,有運用的材,有智能的才,就可以改變世界,今生我已經可以「自力」,為所欲為.何必要有神的信仰?更加不需要「得救」.

4.為「放下」你可以去到「幾盡」?

讓我們進入今日經課的第三個段落,最後一段,馬可福章十章28至31節:

28彼得就對他說:「看哪,我們已經撇下所有的跟從你了.」

29耶穌說:「我實在告訴你們,人為我和福音撇下房屋,或是弟兄,姊妹,父母,兒女,田地,

30沒有不在今世得百倍的,就是房屋,弟兄,姊妹,母親,兒女,田地,並且要受逼迫,在來世必得永生.

31然而,有許多在前的,將要在後,在後的,將要在前.」

第三段落由作為門徒之首的彼得的感嘆,因大家已經好擔心.彼得想耶穌回應他們的實況,說:「看哪,我們已經撇下所有的跟從你了.」

用另一種語言講,彼得覺得我們已經由having something轉成having nothing,由doing變成following Jesus / being Christ-followers.

我考入香港中文大學時第一份讀書報告的功課是書寫Eric Fromm的一本書The art of loving.Eric Fromm 有另一本好出名的書叫To have or to be,描繪人的存在有兩種狀態,一種是擁有的狀態(having mode of existing),另一種是存有的狀態(being mode of existing).這段經文,就是讓我們由前者進入後者.

耶穌在此是否在指出信仰的路上的一個原理?

29耶穌說:「我實在告訴你們,人為我和福音撇下……

31然而,有許多在前的,將要在後,在後的,將要在前.」

這個原理,可以用一個比喻來說明.

甲乙丙丁四個比賽跑步一百公尺.這個比賽在起步安排上是這樣的:

P1:按財富比例有先走距離的優惠,

P2:但同時要背起財富比例相應的重量.

設想甲乙丙丁的財富比例是10:5:2:1,富甲一方的甲就背起十公斤負擔的背包跑,但起點有1 m的優惠；貧窮白丁就背起一公斤負擔的背包跑,起點只有0.1 m的優惠,乙丙各有0.5 m及0.2 m的起步優惠,但分別要背起五,二公斤負擔的背包跑.

設想大家的跑速都是一樣地跑一百公尺要12秒 (12 s / 100 m = 0.12 s/m),因應負重而打相應的折扣比例,即是負重累事而減速,比例是:10:5:2:1.

請問,四人誰最快到終點?以下是計算:

甲的計算: (12 s / 100 m) x 110\% x (100 m - 1 m) = 12/100 x 110\% x 99 m = 13.068 s

甲跑完一百米,要13.068秒.

乙的計算: (12 s / 100 m) x 105\% x (100 m - 0.5 m) = 12/100 x 105\% x 99.5 m = 12.537 s

乙跑完一百米,要12.537 秒.

丙的計算: (12 s / 100 m) x 102\% x (100 m - 0.2 m) = 12/100 x 102\% x 99.2 m = 12.142 s

丙跑完一百米,要12.142秒.

丁的計算:(12 s / 100 m) x 101\% x (100 m - 0.1 m) = 12/100 x 101 \% x 99.9 m = 12.108 s

丁跑完一百米,要12.108秒.

甲雖然贏在起跑線上,但跑得最後才到.換言之,因擁有最多而佔有的起點優惠,也因擁有最多而負累更多,負累抵銷了優惠,在前的反而在後.

相反,丁雖然輸在起跑線上,可是,因擁有最少因而負累更少,負累抵銷了優惠,在後的反而在前.

你在世擁有甚麼,會否成為你進天國,得永生的負累?

你在世擁有甚麼財富,權位與才智,這些條件會否成為你進天國,得永生的負累?

我們如何不把今生的擁有成為進天國,得永生的負累?反而,成為神恩待你而可以讓你:「變賣所擁有,分給窮人,成為耶穌跟從者.」

耶穌所講的「有許多在前的,將要在後,在後的,將要在前」這個信仰原理,可以從苦難入手看.最早受患難的,累積堅忍,反而成為最能捱到最後的人.一直無苦無困的人,後來一出現苦難,就會最早離場.

當你受苦受困時,最早跟你同行的,是否真的可以跟你捱到最後?

如果無法避免困苦,你希望在那一年齡?許多家長寧願孩子早點捱苦,反而可捱得起.

把擁有的存在模式轉化為存有的存在模式,也許可以進入神的時間,以為最早有苦難的就是最後可以捱到苦難.

5.為「跟從神」甚麼都可以去到「幾盡」?

當一個人當自己作無有(having nothing),就可以為「跟從神」甚麼可以去到「幾盡」,連生命都可以交上.

但這生命的交上,並不是找死,雖然死也不怕.

古言有云:民不畏死,奈何以死懼之.

即是說:對於一個不怕死的人,用死威嚇他,係威嚇唔到佢.

生命的交上,是倒空自己,視自己一無所有,有的都是神的恩待,為的是在地若天地行神的道.換言之,若有多少的財富,權能與才智,為的都是成為祝福他人的條件,好讓大家跟從神.

耶穌是一個例子,本是神的兒子,反倒虛己,變成人的樣式,住有一無所有的被邊緣化的人群中間.請問,他有何怕?

既然天父會叫他三天復活,就可以為「跟從神」可以去到「幾盡」:被釘於十字架上為我們代死.

反倒虛已(emptying oneself)是一種存有的境界,意思是我的生命不再是我的,乃是回天家前能在今生活多少則多少.

你對這樣的信仰境界領會了多少?

阿們.



\newpage

\section{立於亂世之中}
\label{sec:428}
\textbf{鄧瑞強博士 - 立於亂世之中}
\newline
\newline
\hyperref[sec:427]{< < < PREV SERMON < < <}
~
\hyperref[sec:toc]{[返目錄]}
~
\hyperref[sec:429]{> > > NEXT SERMON > > >}


啟示錄 13:1-18
\newline
\begin{longtable}{cl}
\hline
\hline
章節 & 經文 (和合本修訂版)\\
\hline
13:1 & \begin{tabularx}{0.7\textwidth}{X} 我又看見一隻獸從海裡上來，有十個角七個頭；在十個角上戴著十個冠冕，七個頭上有褻瀆的名號。 \end{tabularx} \\ \\ \relax
13:2 & \begin{tabularx}{0.7\textwidth}{X} 我所看見的獸，形狀像豹，腳像熊的腳，口像獅子的口。那條龍將自己的能力、座位和大權柄都給了牠。 \end{tabularx} \\ \\ \relax
13:3 & \begin{tabularx}{0.7\textwidth}{X} 我看見獸的七個頭中，有一個似乎受了致命傷，那傷卻醫好了。全地的人都很驚訝，跟從了那隻獸。 \end{tabularx} \\ \\ \relax
13:4 & \begin{tabularx}{0.7\textwidth}{X} 他們都拜那條龍，因為牠把自己的權柄給了獸；又拜那隻獸，說：「誰能比這隻獸，誰能與牠交戰呢？」 \end{tabularx} \\ \\ \relax
13:5 & \begin{tabularx}{0.7\textwidth}{X} 龍又賜給那隻獸說誇大褻瀆話的口，又賜給牠權柄可以任意行事四十二個月。 \end{tabularx} \\ \\ \relax
13:6 & \begin{tabularx}{0.7\textwidth}{X} 那獸就開口向神說褻瀆的話，褻瀆神的名和他的帳幕，就是那些住在天上的。 \end{tabularx} \\ \\ \relax
13:7 & \begin{tabularx}{0.7\textwidth}{X} 牠又被准許與聖徒作戰，並且得勝，也賜給牠權柄，可以制伏各支派、各民族、各語言、各邦國。 \end{tabularx} \\ \\ \relax
13:8 & \begin{tabularx}{0.7\textwidth}{X} 凡住在地上、名字從創世以來沒有記在被殺羔羊的生命冊上的人都要拜牠。 \end{tabularx} \\ \\ \relax
13:9 & \begin{tabularx}{0.7\textwidth}{X} 凡有耳朵的都聽吧！ \end{tabularx} \\ \\ \relax
13:10 & \begin{tabularx}{0.7\textwidth}{X} 該被擄掠的，必被擄掠；該被刀殺的，必被刀殺。在此，聖徒要有耐心和信心。 \end{tabularx} \\ \\ \relax
13:11 & \begin{tabularx}{0.7\textwidth}{X} 我又看見另一隻獸從地裡上來。牠有兩個角如同羔羊，說話好像龍。 \end{tabularx} \\ \\ \relax
13:12 & \begin{tabularx}{0.7\textwidth}{X} 牠在第一隻獸面前施行第一隻獸所有的權柄，並且使地和住在地上的人拜那致命傷被醫好了的第一隻獸。 \end{tabularx} \\ \\ \relax
13:13 & \begin{tabularx}{0.7\textwidth}{X} 這隻獸又行大奇事，甚至在人面前使火從天降在地上。 \end{tabularx} \\ \\ \relax
13:14 & \begin{tabularx}{0.7\textwidth}{X} 牠得了權柄在第一隻獸面前能行奇事，迷惑住在地上的人，告訴他們要為那受過刀傷還活著的獸造個像。 \end{tabularx} \\ \\ \relax
13:15 & \begin{tabularx}{0.7\textwidth}{X} 又有權柄賜給牠，讓那隻獸的像有生氣，並且能說話，又使所有不拜獸像的人都被殺害。 \end{tabularx} \\ \\ \relax
13:16 & \begin{tabularx}{0.7\textwidth}{X} 牠又使眾人，無論大小、貧富，自主的、為奴的，都在右手上，或是在額上，打一個印記； \end{tabularx} \\ \\ \relax
13:17 & \begin{tabularx}{0.7\textwidth}{X} 這樣，除了那有印記，有獸的名或有獸名數字的，都不得買或賣。 \end{tabularx} \\ \\ \relax
13:18 & \begin{tabularx}{0.7\textwidth}{X} 在此，要有智慧：讓有悟性的人解開獸的數目吧，因為這是一個人的數字，那數字是六百六十六。 \end{tabularx} \\ \\
[1ex]
\hline
\hline
\end{longtable}
講題:立於亂世之中 　Walk on Troubled Waters

經文:啟示錄13章1-18節

講員:鄧瑞強博士

場合:香港中文大學崇基學院禮拜堂

日期:2018年10月21日

各位弟兄姊妹,早安.

小時候,很喜歡看「鹹蛋超人」,超人用「十字死光」打怪獸,保衛地球.當超人打死怪獸時,小朋友們都一起歡呼.怪獸從何而來?原來,怪獸出沒,很多時是因為人類攪擾了大自然.譬如,有隻怪獸Gubira,就因為人類在海底進行建設而吵醒了牠.日本的怪獸之王哥斯拉,牠因核能而出現,因牠以核能為養份.看來,怪獸與人的科技息息相關.人強化自己力量之時,也帶來獸性.在《蜘蛛俠》電影系列裡,反派角色通常就是人與科技結合而成的「怪獸」.譬如,在《蜘蛛俠2》裡,反角「八爪魚博士」(Dr. Octopus)就是人與四隻有智能的巨型機械臂合成的「怪物」.人假借科技的力量,將自己變成「怪獸」.最恐怖的怪獸,就是人直接變成獸,「狼人」是其中一種.明明是人,卻因著某些力量的牽引,變成了狼.沒有人性,只有獸性.

哈佛大學的學者王德威寫了《歷史與怪獸》(The Monster that is History),探討「現代性」與「怪獸性」.看來,科技越進步,控制人的能力越大,獸性越強.

因著科技之助,現代社會,可能已變成一隻怪獸.不是外面來了一隻怪獸,而是整個社會變成了一隻怪獸,而我們在怪獸之內,奉獻自己的生命,去強化這隻怪獸的力量.因我們在怪獸內,是怪獸的一部分,我們可能以為,根本沒有怪獸存在.

二次大戰時期,納粹德國將猶太人輸送入焚化爐中.焚化爐,是吞噬人的現代機器,是食人的大怪獸.但當大屠殺的主要負責人艾希曼(Adolf Eichmann)在耶路撒冷受審時,他毫不覺得自己有罪,他說,我只是一個執行命令的人而已.他成為怪獸的一部分,也就看不見怪獸.

怪獸無處不在,不斷變身.人類與怪獸的搏鬥,歷久常新.歷史就是人類打怪獸的記錄.

今日,我們看一段啟示錄的經文,這裡記載著人與怪獸相遇的故事.

今日的講道經文:啟示錄13:1－18

我將這經文分成兩段,先看第一段.

(第一段)

啟13:1　我又看見一個獸從海中上來,有十角七頭,在十角上戴著十個冠冕,七頭上有褻瀆的名號.

啟13:2　我所看見的獸,形狀像豹,腳像熊的腳,口像獅子的口.那龍將自己的能力,座位,和大權柄都給了牠.

啟13:3　我看見獸的七頭中,有一個似乎受了死傷,那死傷卻醫好了.全地的人都希奇跟從那獸,

啟13:4　又拜那龍──因為牠將自己的權柄給了獸,也拜獸,說:「誰能比這獸,誰能與牠交戰呢?」

啟13:5　又賜給牠說誇大褻瀆話的口,又有權柄賜給牠,可以任意而行四十二個月.

啟13:6　獸就開口向神說褻瀆的話,褻瀆神的名並他的帳幕,以及那些住在天上的.

啟13:7　又任憑牠與聖徒爭戰,並且得勝；也把權柄賜給牠,制伏各族,各民,各方,各國.

啟13:8　凡住在地上,名字從創世以來沒有記在被殺之羔羊生命冊上的人,都要拜牠.

啟13:9　凡有耳的,就應當聽!

啟13:10　擄掠人的,必被擄掠；用刀殺人的,必被刀殺.聖徒的忍耐和信心就是在此.

這裡提到一隻「海獸」,牠背後的力量根源是「龍」.「龍」是邪惡的終極力量,這邪惡的終極力量需要透過一隻具體的怪獸表現出來.譬如:「龍」是「暴力」這種邪惡的根源,但若沒有暴力的人或暴力的政權,「龍」的「暴力」仍是不可見的.「龍」是「驕傲」的根源,而那些目空一切的驕傲者,將「龍」的這種本性活現出來.在啟示錄裡,邪惡的「龍」透過「海獸」將其邪惡活現出來.今日,任何吃人的制度,泯滅人性的思想,都是怪獸,都來自邪惡的「龍」.

這裡的「海獸」,其意念來自舊約.在舊約,有兩隻海獸.一隻叫「拉哈伯」(Rahab)(參:伯9:13, 26:12; 詩89:10; 賽30:7, 51:9)；一隻叫「利維坦」(Leviathan)(參:伯41:4, 詩74:14, 詩104:26, 賽27:1.《和合本》譯作「鱷魚」；《現代中文譯本》譯作「海怪」,「海獸」).在以賽亞書51:9-10裡,將「拉哈伯」等同「埃及」(也參:賽30:7).埃及,是奴役人的帝國,被看成是一頭怪獸.當日,以色列人被解放,不用再作奴隸.離開埃及時,他們過紅海.那時,紅海分開,先知以賽亞將「紅海分開」視為是「斬開了」拉哈伯這隻海獸.人性的解放,生命的自由,全賴和這隻怪獸對抗.

至於「利維坦」,後來成了英國一本政治哲學書的書名.作者霍布斯(Thomas Hobbes)認為,權力集中的威權國家,力量最強大,最能保護人民.當然,為了統一意志,集中力量,這國家的權力是不容人民質疑和挑戰的.如此,國家成了一隻「大海獸」利維坦.到了今日,我們看到,這種權力集中的怪獸,是吃人的.

科技的力量,人心的驕傲,引來一隻又一隻的怪獸,神的救贖,就是克服一隻又一隻的怪獸.

回到啟13:1-5的經文.啟示錄這隻「海獸」,「十角七頭」.第12章講到,「龍」是「七頭十角」的.可以看到,「海獸」其實是「龍」的翻版.「七頭」,按啟示錄17:9的解說,是指「七山之城羅馬」.羅馬,是奴役人的新埃及.「角」是「武力」,海獸具備十足的武力.「形狀像豹,腳像熊的腳,口像獅子的口」,都指向它的暴力.舊約的海獸,指向奴役人的埃及.啟示錄的海獸,指向以暴力欺壓人的羅馬.都是邪惡的化身,是奴役人的力量,是扭曲人的生命的暴力.

這隻怪獸,甚至有「起死回生」的表面力量,有逆境求全的能力.暴力,能清除反對者,異見者,帶來表面的和諧,有一時的繁榮,但這只是表面的安好.

啟13:5　又賜給牠說誇大褻瀆話的口,又有權柄賜給牠,可以任意而行四十二個月.

啟13:6　獸就開口向神說褻瀆的話,褻瀆神的名並他的帳幕,以及那些住在天上的.

短短兩節經文,出現三次「褻瀆」.「褻瀆」,就是看不起神明,對神明不敬.膽敢如此行的人,定是自視為神.暴力的政權,容易有造神運動.一個人變了神,他做的,便變得合法；他的意志,便變成絕對.

啟13:8　凡住在地上,名字從創世以來沒有記在被殺之羔羊生命冊上的人,都要拜牠.

當政治變成宗教,統治者變成神明,社會生活便變得謹慎恐懼,只怕得罪了神明一樣的統治者.這神明不能錯,若錯了,便用謊話掩飾.如此,弄虛作假,變成生活的常態.前捷克總統哈維爾(Vaclav Havel)寫過一篇長文《無權力者的權力》(The Power of the Powerless),描述在捷克極權時期的生活.他講到一個蔬果店的經理,在櫥窗上貼上一幅標語:「全世界的工人們,聯合起來!」他這樣做目的何在?

哈維爾這樣寫:「我敢斷言,大多數商店經理們對於櫥窗上標語的意義是從來不會過問的,他們也不會用那樣的標語來表達自己真實的意見.」「但要是不貼,就會有麻煩.上面會指責他不按規定佈置櫥窗,有人甚至會控告他叛亂.為了要過得下去,他非照章辦事不可.這樣的小事千千萬萬,做了才能保障能過如他們所說的『與社會和諧』的相對平靜的生活.」「這標語確實是一個符號,並且他本身包含了一個說不出口來但卻十分明確的信息.從字面上說起來,可以這樣表達:『我,水果商某某,住在這兒,並且清楚地知道自己應該做什麼,我是個靠得住的人,無可挑剔,我很聽話,所以我有權利過平安日子.』」「因此之故,他們必須活在謊言之中.人們不必去接受謊言；他們只要去接受與謊言為伍,活在謊言中的生活,這就夠了.就是這樣,人們確認了這個體制,完善了這個體制,製造了這個體制；人們就是這個體制.」(哈維爾:《無權力者的權力》(左岸文化事業,2003),頁77,78,82.)

被好好教育的人民,作出無意義的虛假行為,為的是自保.然而,這種自保,便將自己變成這種怪獸體制的一部分.

在這暴力,虛假,統治者神化自己的怪獸體制裡,基督徒因著對真神的敬拜,會反對將人神化,反對虛假,反對暴力.

啟13:9　凡有耳的,就應當聽!

啟13:10　擄掠人的,必被擄掠；用刀殺人的,必被刀殺.聖徒的忍耐和信心就是在此.

這經文講出基督徒的命運,要被抓去坐監的,真的會坐監；要被逼害的,真的會被逼害.這時,我們會問,神在哪裡?神在堅持真理的信徒的受苦中,神在認信真理的信徒的見證中,神在為真理而受苦的十架上.

我再引用哈維爾的話,這個曾因「顛覆國家」罪名而坐牢的人,他說:「在謊言中生活的人,只要不遇到在真實中生活的人,那些揭露虛偽的觀點,就不會存在.只要一有不同的選擇,表象和在謊言中生活的本質和所有一切就會受到威脅.同時,這個不同選擇所佔據的空間究竟有多大,完全不重要,因為這個不同選擇的力量並不在其物理的屬性,而是在其揭露體制支柱和不穩定基礎所投射出的光明.」(哈維爾:《無權力者的權力》,頁93.字眼有少許改動.)

微少的火光,可以驅散黑暗.啟示錄說:「聖徒的忍耐和信心就是在此.」

(第二段)

啟13:11　我又看見另有一個獸從地中上來有兩角如同羊羔,說話好像龍.

啟13:12　牠在頭一個獸面前,施行頭一個獸所有的權柄,並且叫地和住在地上的人拜那死傷醫好的頭一個獸.

啟13:13　又行大奇事,甚至在人面前,叫火從天降在地上.

啟13:14　牠恩賜給牠權柄在獸面前能行奇事,就迷惑住在地上的人,說:「要給那受刀傷還活著的獸作個像.」

啟13:15　又有權柄賜給牠,叫獸像有生氣,並且能說話,又叫所有不拜獸像的人都被殺害.

啟13:16　牠又叫眾人,無論大小,貧富,自主的,為奴的,都在右手上或在額上受一個印記.

啟13:17　除了那受印記,有了獸名或有獸名數目的,都不得作買賣.

啟13:18　在這裡有智慧:凡有聰明的,可以算計獸的數目；因為這是人的數目,它的數目是六百六十六.

這裡有隻「地獸」,有兩角如同羊羔.早一陣子,在香港歷史博物館有一個展覽,主題叫＜奢華世代:從亞述到亞歷山大＞,其中展出一個刻有亞歷山大大帝肖像的錢幣,那個亞歷山大大帝的頭,便長有兩隻羊角.原來,頭上有羊角,是羅馬皇帝當時的形象,代表力量.

這隻「地獸」的主要工作,是推行極端的君王崇拜,靠的是「迷惑」人的謊言.對不願向統治者低頭的人,「地獸」可以封鎖其經濟生活,甚至殺害他.這是全面的強權,政治,經濟,社會生活,全受影響.「地獸」可以說是統治者的打手.拍統治者馬屁的人,他們鑑貌辨色,為了討好主人,他們會用上極端手段.

「地獸」的數目是「人的數目」,牠的數目是666.希伯來文及希臘文,字母都有數值.希伯來人名或希臘人名,都可以根據其名字的字母,算出一個數值來.「666」到底是指誰?最多人同意的講法是指當時殺害基督徒的「尼祿」皇帝.後來,歷史上所有的大壞蛋,都有人認為就是「666」.不過,我想,「666」所指的,是人的本身,這是「人的數目」.任何有限的人,忘記了自身的有限性,自封為神,都是「怪獸」.

啟示錄第13章,講到一怪獸世界.有限的人,藐視神明,自視為神,橫行無忌,令世界充滿獸性.怪獸是凶殘的,食人的.一般人,只得向怪獸低頭,在虛情假意的世界,虛情假意地活著.在敬拜獸的時候,將自己變成獸.正正是在這種處境中,一個人敬拜誰,將生命交給誰,便成為最關鍵的問題.敬拜神,不是禮拜日在教堂裡坐坐,然後和會友吃餐飯.敬拜神,是將生命交給那又公義又慈愛的神.當怪獸橫行時,在死亡威迫下,一個人仍拒絕敬拜怪獸,仍堅持敬拜獨一的真神,這就是真信仰.你願意持守這種真信仰嗎?

曾受苦的基督應許說:「得勝的,我要賜他在我寶座上與我同坐,就如我得了勝,在我父的寶座上與他同坐一般.」(啟3:21)

願有一天,各人能在基督的寶座上,與基督同坐.

但願:榮耀歸於聖父,聖子,聖靈.阿們.



\newpage

\section{復興之道}
\label{sec:429}
\textbf{林豪恩先生 - 復興之道}
\newline
\newline
\hyperref[sec:428]{< < < PREV SERMON < < <}
~
\hyperref[sec:toc]{[返目錄]}
~
\hyperref[sec:430]{> > > NEXT SERMON > > >}


耶利米書 31:7-9
\newline
\begin{longtable}{cl}
\hline
\hline
章節 & 經文 (和合本修訂版)\\
\hline
31:7 & \begin{tabularx}{0.7\textwidth}{X} 耶和華如此說：「你們當為雅各歡樂歌唱，為萬國中為首的歡呼。當傳揚，頌讚說：『耶和華啊，求你拯救你的百姓，拯救以色列的餘民。』 \end{tabularx} \\ \\ \relax
31:8 & \begin{tabularx}{0.7\textwidth}{X} 看哪，我必將他們從北方之地領來，從地極召集而來；同他們來的有盲人、瘸子、孕婦、產婦；他們必成群結隊回到這裡。 \end{tabularx} \\ \\ \relax
31:9 & \begin{tabularx}{0.7\textwidth}{X} 他們要哭泣而來。我要照他們懇求的引導他們，使他們在河水旁行走正直的路，他們在其上必不致絆跌；因為我是以色列的父，以法蓮是我的長子。 \end{tabularx} \\ \\
[1ex]
\hline
\hline
\end{longtable}
講題:復興之道  　Ideas of Restoration

講員:林豪恩先生

經文:耶利米書31:7-9

場合:香港中文大學崇基學院禮拜堂

日期:2018年10月28日

一. 引言

復興這兩個字,我們有時在基督教界聽到,大概是呼籲信徒從「不冷不熱」的狀況,靠著聖靈的能力,重新火熱.具體表現是積極投入事奉,廣傳福音.按字典的解釋,「復興」一般的意思是:「衰落後再興盛起來」.剛巧,近年不少國家都在談論復興.

大約十年前,中央電視台製作了一部介紹中國近代史的政論片,名為《復興之路》,透過回首過去,展望未來,探討中華民族的偉大復興如何實現.在地球的另外一方,特朗普以「令美國再次強大」的口號勝出了總統選舉,在東西之間,土耳其重新推動鄂圖曼精神,期望復興輝煌歷史.

剛才大家聆聽經課時,有沒有聽到復興的聲音呢?今天的講題是Ideas of Restoration,有沒有看見耶利米書描繪以色列人被擄回歸的景象?有沒有聽到詩篇頌唱重建的歌聲?有沒有想像到希伯來書介紹那更新了的祭司職任?最後,有沒有留意到那位福音書記載中恢復視力的朋友?這四段經課對我們思考復興有甚麼啟發呢?帶給我們怎樣的復興想像呢?

二. 經文

耶利米書

31:7耶和華如此說:「你們當為雅各歡樂歌唱,為萬國中為首的歡呼.當傳揚,頌讚說:『耶和華啊,求你拯救你的百姓,拯救以色列的餘民.』

31:8看哪,我必將他們從北方之地領來,從地極召集而來；同他們來的有盲人,瘸子,孕婦,產婦；他們必成群結隊回到這裡.

31:9他們要哭泣而來.我要照他們懇求的引導他們,使他們在河水旁行走正直的路,他們在其上必不致絆跌；因為我是以色列的父,以法蓮是我的長子.

1.1. 社會制度的復興:社會應該包括「盲人,瘸子,孕婦,產婦」等弱勢群體

亡國被擄的,有一天竟然能夠踏上回歸故土重建家園的道路,他們歡樂歌唱是可以明白的,他們熱淚盈眶也是可以理解的.相對特別的,卻是第八節所提及在回歸人群中的幾類人:「同他們來的有盲人,瘸子,孕婦,產婦」.這現象啟發我們思考復興的第一個層面,是關於人,或者說,是關於群體,用現代的詞彙說,是社會制度.

在回歸的群體中,有「盲人,瘸子,孕婦,產婦」這些人從何而來呢?應該不是被擄往巴比倫的人,因為沒有戰勝國會愚蠢地把戰敗國的弱勢群體帶到自己國家去照顧.戰勝國通常擄走政治上的重要人物,以及有生產力的壯丁.列王記下24:14節記述巴比倫王尼布甲尼撒攻破耶路撒冷後的行動,符合我們的假設:「又將耶路撒冷的眾民和眾首領,並所有大能的勇士,共一萬人,連一切木匠,鐵匠都擄了去；除了國中極貧窮的人以外,沒有剩下的.」耶利米書提及的「盲人,瘸子,孕婦,產婦」,可能就是那些連尼布甲尼撒也不願帶走的「國中極貧窮的人」,或者是不良於行的人,或者是那些在戰爭中或逃亡中受傷的人.無論來歷如何,這些都是需要被照顧的人,而不是在回歸重建的重要時期能夠在國防上和生產上貢獻的人.上帝的選擇很特別,祂竟然說:「31:8看哪,我必將他們從北方之地領來,從地極召集而來；同他們來的有盲人,瘸子,孕婦,產婦；他們必成群結隊回到這裡.」即是說,無論他們是強勢或弱勢,是有經濟能力還是需要被照顧,是有防衛能力還是需要被保護,他們都屬於同一個回歸群體,屬於同一個回歸的隊伍,都是重建中的社會的一份子.

從耶利米書的敘述可見,上帝期望我們的社會應該是包容性(inclusive)的,而不是排他性(exclusive)的.然而,在追求復興或強大的過程中,有些想法是盡量清除拖累社會的成員,保留生產力強盛的精英.大家一定知道第二次世界大戰時期,納粹德國屠殺猶太人,但有否聽過他們連自己德國人也屠殺呢?1939 年 10 月,納粹開始屠殺那些被認為「不配擁有生命」的德國人.起初,他們鼓勵醫院無視這些病人,讓病人在饑餓和疾病的折磨下死去.後來,成批的「顧問」來到醫院,決定誰應當死亡.那些病人被送到各個屠殺中心,死於致命的注射或毒氣.在整個二時期間,有人估計有多達200,000 身體或精神有缺陷的德國人被屠殺.

或者有人說,這是非常時期的非常手段.然而,為了強大而清除弱者,為了進步而趕走落後者,為了經濟增長而把人以GDP的貢獻分等級,這些價值豈不也是一樣危險嗎?

追求復興的時候,我們不能忘記,復興最終的目標是人民,而「盲人,瘸子,孕婦,產婦」這些需要幫助的人也應該是人民中的一份子.有人說:「一個社會的文明程度取決於它對弱者的態度」.

詩篇

126:1當耶和華使錫安被擄的人歸回的時候,我們好像做夢的人.

126:2那時,我們滿口喜笑,滿舌歡呼；那時,列國中就有人說:「耶和華為他們行了大事!」

126:3耶和華果然為我們行了大事,我們就歡喜.

126:4耶和華啊,求你使我們這些被擄的人歸回,好像尼革夫的河水復流.

126:5流淚撒種的,必歡呼收割!

126:6那帶種流淚出去的,必歡呼地帶禾捆回來!

2.1. 經濟運作的復興:付出有回報,辛苦有收穫,是一個可期待的社會

呂振中譯本把「被擄歸回」翻譯為「恢復故業」,幾個常用的英文譯本都接近這個意思.輝煌的故業復興,過往的幸福重拾,當然是值得歡喜；失而復得,如乾涸了的河道復流,難怪詩人說:「我們好像做夢一樣的人」.正如上一段所說,為這種狀況歡慶並不是特別的事情.最後的兩句,卻引導我們思考復興的第二個層面,關於勞動與回報,用現代的語言,是經濟制度.

126:5流淚撒種的,必歡呼收割!

126:6那帶種流淚出去的,必歡呼地帶禾捆回來!

這兩節是平行的,第六節具體地呼應第五節,表達相同的意思,就是勞動者會得到勞動的成果,辛苦付出會有正面回報.如果用許冠傑的歌來說,就是「出左半斤力,想話lor 返七八量」.今時今日,當各國都把復興作為口號的時候,有沒有考慮經濟上的公平性呢?有沒有留意勞動者是否得到勞動的成果呢?

剛才提過,特朗普以「讓美國再次強大」為口號勝出了總統選舉,他這個「再次強大」的口號與東方領導人「復興」的口號不謀而合,遙相呼應.要再次強大,即是說曾經衰弱,那是甚麼時候?甚麼狀況呢?容易令人想到的,是2008年的金融海嘯及引發的經濟衰退.Now新聞最近有一篇關於美國的報導:

【Now新聞台9/10/2018】「金融海嘯後,經濟逐漸遠離寒冬,理應越來越多人可以安居樂業,但現實並非如此,無家可歸的人反而與日俱增.專家把問題歸咎於財富分配不均,低技術工人被淘汰,房屋供應緊絀導致租金上漲,政府減少房屋及醫療資助,令貧窮戶百上加斤.洛杉磯無家可歸者人數五年內增加了四成九.市長直言,「無家可歸」已成為當今世代最大的道德與人道危機.」

領導人吹捧自己已經令國家再次強大,並且傳言連任的競選口號是「繼續強大」.國家強大了,人民幸福嗎?資產價格上升了,勞動者得到勞動的回報嗎?國家經濟指數亮麗了,人民的生活得到改善嗎?

詩人並非期望不勞而獲,只是提出重建一個社會的時候,經濟制度上最基本的要求:「126:5流淚撒種的,必歡呼收割!126:6那帶種流淚出去的,必歡呼地帶禾捆回來!」我們每天工作,掃地的,售貨的,駕駛的,辦公室工作的……我們的付出,所得的回報能夠讓我們好好地生活嗎?我們一生人在工作,所得的回報能夠讓我們放心退休嗎?如果能夠,這才是真正的復興,真正的強大.

希伯來書

7:23一方面,那些成為祭司的數目本來多,是因為受死亡限制不能長久留住.

7:24另一方面,這位既是永遠留住的,他具有不可更換的祭司職任.

7:25所以,凡靠著他進到上帝面前的人,他都能拯救到底,因為他長遠活著為他們祈求.

7:26這樣一位聖潔,無邪惡,無玷污,遠離罪人,高過諸天的大祭司,對我們是最合適的；

7:27他不像那些大祭司,每日必須先為自己的罪,後為百姓的罪獻祭,因為他只一次將自己獻上就把這事成全了.

7:28律法所立的大祭司本是有弱點的人,但在律法以後,上帝以起誓的話立了兒子為大祭司,成為完全,直到永遠.

3.1. 祭司職分的復興:包括安身立命,信仰及心靈世界

希伯來書宣布,在新的時代,不但不會取消祭司制度,祭司制度更會被更新.我們剛才讀的那段經文,主要講述更新了的祭司制度下的大祭司.作者在26節指出這位大祭司「對我們是最合適的」,在個人特質上,他「聖潔,無邪惡,無玷污,遠離罪人,高過諸天」(26),在執行任務上,他「只一次將自己獻上就把這事成全了」(27),在果效上,「凡靠著他進到上帝面前的人,他都能拯救到底,因為他長遠活著為他們祈求」(25).

在新時代的復興中,信仰到底有甚麼位置呢?有人曾說:「人民有信仰,民族有希望,國家有力量.」即是說,信仰是復興不可或缺的元素,與希望和力量有密切關係.重要的是,該「信仰」帶領人民往朝見誰呢?指示人民往哪裡去呢?從古至今,中國人都不單單追求物質的進步,更不停地尋索「安身立命」之道.能夠讓人安身立命的信仰,必需能夠給予人-種深度的安全感,跨過高山深谷,穿越生死也感到平安；也必需能夠給予人一種深度的滿足感,在患難中有忍耐,在忍耐中有盼望,因盼望帶來喜樂.中國人也知道這種平安和喜樂不能單單從人而來,因此而懂得仰望上蒼.

中國人也有以上天為價值源頭的道德精神傳統,文天祥的「正氣歌」說:「天地有正氣,雜然賦流形.下則為河岳,上則為日星.於人曰浩然,沛乎塞蒼冥.」高師寧是中國宗教社會學學者,她對企業家群體的商業道德做過調查,發現有信仰的人和不信仰的人不一樣.信基督教或天主教的企業家,不能說就盡善盡美,但相比那些不信教的企業家而言,他們對職工會更好,對客戶會更講誠信,還有很多人熱衷於慈善公益事業.

希伯來書的作者強調,新的時代,是祭司制度復興的年代,這位新的大祭司貫通天地,帶領人直接進到上帝面前,引導人通往安身立命之處,支持人好好地生活.

馬可福音

10:46他們到了耶利哥.耶穌同門徒並許多人離開耶利哥的時候,有一個討飯的盲人,是底買的兒子巴底買,坐在路旁.

10:47他聽見是拿撒勒的耶穌,就喊了起來,說:「大衛之子耶穌啊,可憐我吧!」

10:48有許多人責備他,不許他作聲,他卻越發喊著:「大衛之子啊,可憐我吧!」

10:49耶穌就站住,說:「叫他過來.」他們就叫那盲人,對他說:「放心,起來!他在叫你啦.」

10:50盲人就丟下衣服,跳起來,走到耶穌那裡.

10:51耶穌回答他說:「你要我為你做甚麼?」盲人對他說:「拉波尼,我要能看見.」

10:52耶穌對他說:「你去吧!你的信救了你.」盲人立刻看得見,就在路上跟隨耶穌.

4.1. 視力的復興:恢復視力,恢復視野,看見甚麼,跟隨甚麼

最後一段經課,馬可福音10章46-52節,講述一個盲人恢復視力的故事.耶穌問:「你要我為你做甚麼?」盲人對他說:「我要能看見.」這位朋友終於恢復視力,他看見甚麼呢?看見又如何呢?視力正常的朋友,你又看見甚麼呢?看見又如何呢?有時,看見甚麼,未必由視力決定,而是由視野決定.

經文沒有明言那位盲人恢復視力後看見甚麼,但卻記述了他的行動:「就在路上跟隨耶穌」.我們合理地假設,他看見了耶穌,因此選擇跟隨耶穌.其實,對一位恢復視力的盲人來說,這個五光十色的世界實在有太多東西以前看不到現在可以看了,有太多行動以前做不到現在可以做了,但他似乎在耶穌身上看見了道路,真理,生命,因此就選擇了跟隨道路,真理,生命.

當人們追求復興的時候,他們看見甚麼呢?當人們追求強大的時候,他們看見甚麼呢?這不是視力的問題,而是視野的問題.

視野,有時會被阻擋,令人們看見特定的事物,而被限制去看其他事物.

何光滬是中國宗教研究學者,他指出:「無數的中國人提起傳教士,還只說那都是「鴉片販子」,「外國間諜」,「文化侵略」等等,這不是讓人覺得這個民族驚人地無知,或驚人地健忘嗎?但我認為,這絕非中國人民的過錯,而是那些故意誤導者的罪過.」他接著說:「越來越多不抱偏見的學者,都認識到基督教對中國現代化確實是很大的積極的促進力量.從歷史上看,辛亥革命以前,基督教就對中國社會做出了很大貢獻.例如在鴉片盛行期間,大部分在華傳教士,特別是英國的傳教士,早就大聲疾呼反鴉片,寫文章,畫漫畫,辦戒煙館,勸阻中國人吸鴉片.一些英國傳教士發現在中國呼籲沒有用,就跑回英國去呼籲,又寫文章,又辦雜誌,抨擊鴉片貿易違反道德違反信仰,還成立反鴉片組織,聯絡各界人士,一次次呼籲簽名,向國會請願,要求立法取消鴉片貿易.又例如,辛亥革命前後,天主教和新教傳教士已經在中國各辦了幾千所小學中學,還有中國最早的十幾所大學.又例如,傳教士還興辦了當時中國幾乎所有的醫院,為中國引入了現代醫療和醫學教育機構.」

馬可福音描述施洗約翰說:1:2看哪,我要差遣我的使者在你前面,預備道路.1:3在曠野有人聲喊著說:預備主的道,修直他的路.

今天,修直那條道路,可能不是用鋤頭,也不用推土機,而是撥開雲霧,洗刷天空,恢復視野.當人們恢復視野的時候,但願都好像故事中那位重獲視力的朋友一樣:「立刻看得見,就在路上跟隨耶穌.」

三. 總結

今天的四段經課,啟發我們思考restoration的幾個向度,重建包容性的社會,重建公平的經濟,重建安身立命的信仰,最後而且不能輕視的,是恢復視野,看見包容性的社會令人民守望相助,公平的經濟令人民安居樂業,安身立命的信仰令人民生活得平安及喜樂.



\newpage

\section{神要擦去他們一切的眼淚}
\label{sec:430}
\textbf{伍渭文牧師 - 神要擦去他們一切的眼淚}
\newline
\newline
\hyperref[sec:429]{< < < PREV SERMON < < <}
~
\hyperref[sec:toc]{[返目錄]}
~
\hyperref[sec:431]{> > > NEXT SERMON > > >}


以賽亞書 25:6-9
\newline
\begin{longtable}{cl}
\hline
\hline
章節 & 經文 (和合本修訂版)\\
\hline
25:6 & \begin{tabularx}{0.7\textwidth}{X} 在這山上，萬軍之耶和華必為萬民擺設宴席，有肥甘與美酒，就是滿有骨髓的肥甘與精釀的美酒。 \end{tabularx} \\ \\ \relax
25:7 & \begin{tabularx}{0.7\textwidth}{X} 在這山上，他必吞滅那纏裹萬民的面紗和那遮蓋列國的遮蔽物。 \end{tabularx} \\ \\ \relax
25:8 & \begin{tabularx}{0.7\textwidth}{X} 他已吞滅死亡直到永遠。主耶和華必擦乾各人臉上的眼淚，在全地除去他百姓的羞辱；這是耶和華說的。 \end{tabularx} \\ \\ \relax
25:9 & \begin{tabularx}{0.7\textwidth}{X} 到那日，人必說：「看哪，這是我們的神，我們向來等候他，他必拯救我們。這是耶和華，我們向來等候他，我們必因他的救恩歡喜快樂。」 \end{tabularx} \\ \\
[1ex]
\hline
\hline
\end{longtable}
詩篇 24:1-10
\newline
\begin{longtable}{cl}
\hline
\hline
章節 & 經文 (和合本修訂版)\\
\hline
24:1 & \begin{tabularx}{0.7\textwidth}{X} 地和其中所充滿的，世界和住在其中的，都屬耶和華。 \end{tabularx} \\ \\ \relax
24:2 & \begin{tabularx}{0.7\textwidth}{X} 他把地建立在海上，安定在江河之上。 \end{tabularx} \\ \\ \relax
24:3 & \begin{tabularx}{0.7\textwidth}{X} 誰能登耶和華的山？誰能站在他的聖所？ \end{tabularx} \\ \\ \relax
24:4 & \begin{tabularx}{0.7\textwidth}{X} 就是手潔心清，意念不向虛妄，起誓不懷詭詐的人。 \end{tabularx} \\ \\ \relax
24:5 & \begin{tabularx}{0.7\textwidth}{X} 他必蒙耶和華賜福，又蒙救他的神使他成義。 \end{tabularx} \\ \\ \relax
24:6 & \begin{tabularx}{0.7\textwidth}{X} 這是尋求耶和華的族類，是尋求你面的雅各。（細拉） \end{tabularx} \\ \\ \relax
24:7 & \begin{tabularx}{0.7\textwidth}{X} 眾城門哪，要抬起頭來！永久的門戶啊，你們要被舉起！榮耀的王將要進來！ \end{tabularx} \\ \\ \relax
24:8 & \begin{tabularx}{0.7\textwidth}{X} 這榮耀的王是誰呢？就是有力有能的耶和華，在戰場上大有能力的耶和華！ \end{tabularx} \\ \\ \relax
24:9 & \begin{tabularx}{0.7\textwidth}{X} 眾城門哪，要抬起頭來！永久的門戶啊，你們要高舉！榮耀的王將要進來！ \end{tabularx} \\ \\ \relax
24:10 & \begin{tabularx}{0.7\textwidth}{X} 這榮耀的王是誰呢？萬軍之耶和華是榮耀的王！（細拉） \end{tabularx} \\ \\
[1ex]
\hline
\hline
\end{longtable}
啟示錄 21:1-6
\newline
\begin{longtable}{cl}
\hline
\hline
章節 & 經文 (和合本修訂版)\\
\hline
21:1 & \begin{tabularx}{0.7\textwidth}{X} 我又看見一個新天新地，因為先前的天和先前的地已經過去了，海也不再有了。 \end{tabularx} \\ \\ \relax
21:2 & \begin{tabularx}{0.7\textwidth}{X} 我又看見聖城，新耶路撒冷由神那裡，從天而降，預備好了，就如新娘打扮整齊，等候丈夫。 \end{tabularx} \\ \\ \relax
21:3 & \begin{tabularx}{0.7\textwidth}{X} 我聽見有大聲音從寶座出來，說：「看哪，神的帳幕在人間！他要和他們同住，他們要作他的子民。神要親自與他們同在。 \end{tabularx} \\ \\ \relax
21:4 & \begin{tabularx}{0.7\textwidth}{X} 神要擦去他們一切的眼淚；不再有死亡，也不再有悲哀、哭號、痛苦，因為先前的事都過去了。」 \end{tabularx} \\ \\ \relax
21:5 & \begin{tabularx}{0.7\textwidth}{X} 那位坐在寶座上的說：「看哪，我把一切都更新了！」他又說：「你要寫下來，因為這些話是可信靠的，是真實的。」 \end{tabularx} \\ \\ \relax
21:6 & \begin{tabularx}{0.7\textwidth}{X} 他又對我說：「成了！我是阿拉法，我是俄梅戛；我是開始，我是終結。我要把生命的泉水白白賜給那口渴的人喝。 \end{tabularx} \\ \\
[1ex]
\hline
\hline
\end{longtable}
約翰福音 11:32-44
\newline
\begin{longtable}{cl}
\hline
\hline
章節 & 經文 (和合本修訂版)\\
\hline
11:32 & \begin{tabularx}{0.7\textwidth}{X} 馬利亞到了耶穌那裡，看見他，就俯伏在他腳前，對他說：「主啊，你若早在這裡，我弟弟就不會死了。」 \end{tabularx} \\ \\ \relax
11:33 & \begin{tabularx}{0.7\textwidth}{X} 耶穌看見她哭，並看見與她同來的猶太人也哭，就心裡悲嘆，又甚憂愁， \end{tabularx} \\ \\ \relax
11:34 & \begin{tabularx}{0.7\textwidth}{X} 就說：「你們把他安放在哪裡？」他們對他說：「主啊，請你來看。」 \end{tabularx} \\ \\ \relax
11:35 & \begin{tabularx}{0.7\textwidth}{X} 耶穌哭了。 \end{tabularx} \\ \\ \relax
11:36 & \begin{tabularx}{0.7\textwidth}{X} 猶太人就說：「你看，他多麼愛他！」 \end{tabularx} \\ \\ \relax
11:37 & \begin{tabularx}{0.7\textwidth}{X} 其中有人說：「他既然開了盲人的眼睛，難道不能叫這人不死嗎？」 \end{tabularx} \\ \\ \relax
11:38 & \begin{tabularx}{0.7\textwidth}{X} 耶穌又心裡悲嘆，來到墳墓前。那墳墓是個穴，有一塊石頭擋著。 \end{tabularx} \\ \\ \relax
11:39 & \begin{tabularx}{0.7\textwidth}{X} 耶穌說：「把石頭挪開！」那死者的姐姐馬大對他說：「主啊，他現在必定臭了，因為他已經死了四天了。」 \end{tabularx} \\ \\ \relax
11:40 & \begin{tabularx}{0.7\textwidth}{X} 耶穌對她說：「我不是對你說過，你若信就必看見神的榮耀嗎？」 \end{tabularx} \\ \\ \relax
11:41 & \begin{tabularx}{0.7\textwidth}{X} 於是他們把石頭挪開。耶穌舉目望天，說：「父啊，我感謝你，因為你已經聽了我。 \end{tabularx} \\ \\ \relax
11:42 & \begin{tabularx}{0.7\textwidth}{X} 我知道你常常聽我，但我說這話是為了周圍站著的眾人，要使他們信是你差了我來的。」 \end{tabularx} \\ \\ \relax
11:43 & \begin{tabularx}{0.7\textwidth}{X} 說了這些話，他大聲呼叫說：「拉撒路，出來！」 \end{tabularx} \\ \\ \relax
11:44 & \begin{tabularx}{0.7\textwidth}{X} 那死了的人就出來了，手腳都裹著布，臉上包著頭巾。耶穌對他們說：「解開他，讓他走！」 \end{tabularx} \\ \\
[1ex]
\hline
\hline
\end{longtable}
講題:神要擦去他們一切的眼淚 　God Will Wipe Every Tear from Their Eyes

經課:以賽亞書25:6-9;詩篇24:1-10;啟示錄21:1-6;約翰福音11:32-44

講員:伍渭文牧師

場合:香港中文大學崇基學院禮拜堂

日期:2018年11月4日

引言:葛培理安息了.在安息禮拜舉行前一段時間不斷回顧他的生平和講道,其中有一段提到他如何面對自己的離世.他說:「有一天,有人會說葛培理死了,離開我們了.其實他沒有死,只是改了一個地址.」地址最後會寫上「新界」(New Territories).

葛培理說得對,他沒有死-如果人死如燈,他活在另一個世界-新界.聖經說,上主是亞伯拉罕的神,以撒的神,雅各的神,這不是歷史的描述,是信仰的表白.這些列祖已不在人世,但乃然活著,活在永生神的安息中.甚麼是永生?永生就是連死亡也不能把我們與上主的關系隔斷的關系.這關系在我們相信耶穌,接受他做我們的救主,生命的主就開始了.

教會每年十一月一日為諸聖日,十一月第一主日為諸聖主日,用來紀念離世信徒.俗世社會十月三十一日為萬聖節,其實是諸聖日的前夕All Hallows Eve.羅馬的教會在主後835 年,將原本於五月中旬舉行的歷代無名聖徒記念日,改往每年十一月一日舉行,並稱這日為「諸聖日」(All Saints’ Day),用基督信仰聖化(Christianize) 異教的鬼節,移風易俗.異教的萬聖節,原來用來驅逐厲鬼或籠絡遊魂,帶來合境平安.教會的萬聖節,用來思念聖徒的嘉言懿行和信德.我們相信上主是掌管萬有的,無懼幽暗的世界,我們在主裡與離世聖徒連成一體,同屬一個大公教會,每主日我們認信:我信聖而公的教會,我信聖徒相通.在地的聖徒稱作戰的教會(church militant),離世的聖徒稱凱旋的教會(church triumphant).離世信徒並沒有死去,祇是遷到新界New Territories.

讓我們藉今天諸聖日經課,認識聖經對離世信徒的盼望.賽二十五6-9；詩二十四; 啟二十一1-6a; 約十一32-44.

1. 基督有智慧-帕子終揭開.

舊約賽二十五6

(1)復興萬有,揭開帕子:今天舊約經課和書信經課展視了終未的新天新地,自由沒壓制的新天新地；死亡被吞滅,眼淚被抹乾.相對賽二十四章的荒原世界:「荒涼的城拆毀了」(v.10),「在街上因酒有悲嘆的聲音.」(v.11)新天新地提到是上主的聖山-聖殿雄立其上的錫安山.「他又必在這山上,萬軍之耶和華必為萬民用肥甘設擺筵席,用陳酒和滿髓的肥甘,並澄清的陳酒,設擺筵席.他又必在這山上除滅遮蔽萬國蒙臉的帕子.他已經吞滅死亡直到永遠.主耶和華必擦去各人臉上的眼淚.」(賽二十五6-8)啟示錄也有這對比,「巴比倫大城傾倒了!傾倒了!」(啟十八2),「聖城新耶和撒冷由上主那裡從天而降.」(啟二十一2)

個人.人生常出現與死亡相關的帕子,像忽罹頑疾,或離世前受疾病的折磨.我弟弟最近問我,為何母親信了耶穌,每主日回教會參加崇拜,常讀經祈禱,她最不想發生的竟然發生在她身上.家母晚年中風,臥床兩年多,靠胃管餵食,不能說話.她很怕要人照顧的.我說我不知道,但我知道她現在安息在上主懷中,等候身體復活.

詩人.詩篇二十二篇的作者亦有同樣的帕子蒙臉,看不透為何他病得生死懸一線.「我如水被倒出來；我的骨頭都脫了節；我的心在我裡面如臘鎔化.我的精力枯乾,如同瓦片；我的舌頭貼在我牙床上.你將我安置在死地的塵土中.」(詩二十二14, 15)這病是否是上主離棄他的證據呢?所以一開始時詩人呼喊:「我的神,我的神!為何離棄我?」

我們很容易聯想到,上主使我們面臨死亡,是一種懲罰,死亡確是可怕的災難.若我們看舊約經課的另外選項,次經所羅門智訓第三章就說得更清楚了.「然而義人受上帝保護永遠不遭磨難.認為義人死亡以及死亡會給他們帶來可怕的災難.這乃是一個愚蠢的錯誤.他們離開我們.但這並非一場災難.實際上.義人是在平安裡.表面上看來他們似乎遭受了懲罰.然而他們懷有永生的希望.他們的遭遇.比起他們將要得到的祝福來.那是微不足道的.上帝考驗他們.如同爐火煉金.發現他們配得上與他同在.」(三1-6) (But the souls of the righteous are in the hand of God, and there shall no torment touch them.2 In the sight of the unwise they seemed to die: and their departure is taken for misery,3 And their going from us to be utter destruction: but they are in peace.4 For though they be punished in the sight of men, yet is their hope full of immortality.5 And having been a little chastised, they shall be greatly rewarded: for God proved them, and found them worthy for himself.6 As gold in the furnace hath he tried them, and received them as a burnt offering.)

(2) 我是始,我是終,都成了.書信經課啟二十一1-6a,說到新天新地,「神要擦去他們一切的眼淚；不再有死亡,也不再有悲哀,哭號,疼痛,因為以前的事都過去了.」(啟二十一4 )「都成了!我是阿拉法,我是俄梅嘎；我是始,我是終.」(啟二十一 6a)

新天新地更是計劃的完成.我是始,表示任何事都在上主的計劃中,沒有天意弄人,也沒有盲目的命運,雖然我們不明白.我是終,表示這計劃將按上主的藍圖成就.都成了,是隱密大計的完成,到那天就真相大白,帕子將被揭開.能打開封起來的書卷,祇有稱自己為我是始我是終的被殺羔羊,他會揭開死亡的帕子.對沒有信仰的人,死亡不單猜不透,更是極大的荒謬.

死亡的荒謬.網上看到一篇王丹寫的文章,他引用卡謬「薛西弗斯的神話」一書說到死亡的荒謬.「人生最大的荒謬,就是關於生命本身:我們不願看到自己老去,我們不願生命終結.……生命的每一步都帶領著我們走向衰老,最終走向死亡.……卡繆在《薛西弗斯的神話》這本書中給我們描述的世界的荒謬圖景:眾神懲罰(希臘神話)薛西弗斯,命他不停地推著一塊巨石上山,到了山頂,巨石又以自身的重量滾落下來,如此循環往復.一切努力看上去都是徒勞,這是神話為世人展現的荒謬,如此清晰,如此殘酷.」\url{https://tw.news.yahoo.com/le-mythe-sisyphe-012910293.html}).

但聖經對死亡的解釋並不含糊,事出有因.「罪從一人入了世界,死又是從罪來的,於是死就臨到眾人,因為眾人就犯了罪.」(羅五12) 「按著定命,人人都有一死,死後且有審判.」(來九27) 但有一個人勝過罪惡,他沒有犯罪,死亡不能拘禁他,他從死裡復活.因為這一個人的義,眾人可以脫離罪的咒詛；凡相信他的,也不被死亡的拘禁,死了也身體復活,疾病得到終極的醫治.我們對很多疾病的成因都不曉得,但最關心的是疾病得醫治,但病醫好了最終也要離開世界.

基督在世上三重事工:傳道,醫病,趕鬼；層層相扣.傳悔改使罪得赦的道,與神和好,回歸創造秩序.醫病,恢復肉身的秩序.趕鬼,恢復屬靈的秩序.這只最初的秩序,也是最終的秩序.這秩序已由那位宣稱「我是阿拉法,我是俄梅嘎；我是始,我是終」那一位恢復過來了.因為他,惟有他是「手潔心清,不向虛妄,起誓不懷詭詐的人.」我們看看今天的詩篇經課詩篇二十四篇.

2. 基督有大愛-伴隨走過死蔭的幽谷

詩二十四

(1) 誰是手潔心清?上主在山上必為萬民用肥甘設擺筵席,他又必在這山上除滅遮蔽萬國蒙臉的帕子,擦去各人臉上的眼淚.但誰能登耶和華的聖山呢?誰能站在他的聖所?詩篇二十四篇清楚指出,「就是手潔心清,不向虛妄,起誓不懷詭詐的人.」(vv. 3, 4)

詩篇二十四篇傳統被視為約櫃重返聖殿的詩篇.以色列人因為拜偶像,上主讓約櫃被非利士人擄去.撒下六章記載大衛挑選三萬人要把約櫃運回耶路撒冷,自己率眾在約櫃前作樂跳舞.但中途牛失前蹄,負責運送的烏撒就伸手扶住約櫃.「神向烏撒發怒,因這錯誤擊殺他,他就死在約櫃旁.」(撒下六7)因為民數記四15明言,搬運聖所一切聖物,先要妥善用海狗皮包裹.「不可摸聖物,免得他們死亡.」明白這背景,就使我們更瞭解3, 4節:「誰人登耶和華的山?誰能站在他的聖所?就是手潔心清,不向虛妄,起誓不懷詭詐的人.」烏撒用手觸摸約櫃而遭擊殺,大衛驚怕,不肯把約櫃運進大衛的城,約櫃就停放在迦得人(非利士人)俄別以東(拜以東偶像)的家中三個月,期間上主大大賜福俄別以東全家,大衛就歡歡喜喜把約櫃運進大衛城了.上主祝福本是拜偶像的非利士人,因他願意接立約櫃在家中,我們也相信因這事情,俄別以東離開拜偶像,追求以色列人獨一真神.

約櫃是整個聖殿的核心晶片,沒有約櫃,聖殿不成聖殿.聖殿是上主與人相遇的地方,但人憑甚麼能親近上主呢?誰能登耶和華的山,誰能站在他的聖所?每年贖罪日,大祭司拿著祭牲的血,身繫一條繩子,恐防不測得以把屍首拖出來,誰人都不敢進入至聖所.至聖所之稱為至聖所,因為安置了約櫃.約櫃頂有天使基路帕守護,大祭司把血灑蓋在約櫃頂四角為人贖罪,所以約櫃頂稱施恩座,施恩座原文是遮蓋的意思.約櫃安放了兩塊立約的法版,這是第二法版,第一法版在西乃山因為拜金牛犢被摩西摔碎了.約櫃安放了第二法版表示上主的憐憫,是上主的憐憫,人才能夠親近上主.因為基督成了被殺的羔羊,因著他寶血的遮蓋,我們才能坦然無懼到父面前.「我們祇管坦然無懼地到施恩的寶座前,為要得憐恤,蒙恩惠,作隨時的幫助.」(來四16)

詩二十四篇是慶賀約櫃重返聖殿的詩篇.第一,二節讚美上主之後,出現禮儀對答:朝聖者問:「誰能登耶和華的聖山?誰能站在他的聖所?」(v. 3) 然後祭司答:「就是手潔心清,不向虛妄,起誓不懷詭詐的人.」(.4) 我們可以想像,這班要登山進入聖殿的朝聖者,中間有抬著約櫃的人,他們因為約櫃引帶,轉被叩問為凱旋的歡呼:「眾城門哪,你們要抬起頭來!永久的門戶,你們要被舉起!那榮耀的王將要進來」(v.7)

祇有一位是「手潔心清,不向虛妄,起誓不懷詭詐」,但他已成為贖罪祭,他的血已塗在約櫃的頂成為施恩座.施恩座要進入耶路撒冷至高之錫安山的金頂聖殿；他被殺的羔羊,要能登上至高上主的寶座,坐在其右.因被殺的羔羊,已成為掌管萬有的主.「榮耀的王是誰呢?萬軍之耶和華,他是榮耀的王!」(v. 19) 所以,這詩篇也是基督升天日讀出的詩篇,因著與基督連結,我們與基督不單能昂首進入聖所敬拜上主,也與基督升到天上,坐在父上主的右邊.「他又叫我們與耶穌一同復活,一同坐在天上.」(弗二6)

基督以被殺羔羊的身份,嚐過死味,然後復活,成為掌管萬有的主,登上聖山,我們經歷死亡的信徒,跟他聯合,來到父面前,安息在父的懷中.詩二十二篇描述因頑疾折磨生死懸一線的作者,發出哀歌:「我的神,我的神!為甚麼離棄我?」這也是基督在十架上的哀歌,他也親身經歷「我如水被倒出來；我的骨頭都脫了節；我的心在我裡面如臘鎔化.我的精力枯乾,如同瓦片.」(vv.14, 15)基督經歷死亡,所以他能陪伴我們行過死蔭的幽谷,與我們同在,在臨終之時,陪伴我們.

查理衛斯理一生寫了七千多聖詩,1788年臨終口授最後的詩「在年邁虛弱時」:

在年邁虛弱時,

誰能拯救我這無助的蟲呢?

耶穌,我唯一的盼望,

是我衰敗身心的力量；

容我捉住你展顏一笑,墮進永恆!

(In age and feebleness extreme,

Who shall a helpless worm redeem?

Jesus, my only hope Thou art,

Strength of my failing flesh and heart:

O could I catch one smile from Thee,

And drop into eternity!)

人老離世,如花經盛而衰,最後脫離花托,戛然掉落,化作污泥.詩人原文用Drop into eternity, 墮進永恆.血脈已盡,魂歸上主,自然不過.但詩人以捉住(耶穌)展顏一笑,才安然墮進永恆.臨終是時間的盡頭,詩人用想像力把時間轉換成空間:捉住(耶穌)展顏一笑,墮進永恆.

3. 基督有大能:叫我們起來.

「我的朋友拉撒路睡了,我去叫醒他」.「拉撒路出來!」(約十一11, 43)

經課:約十一32-44.拉撒路離世,耶穌不單心裡難過,哭了,v.33指出他還心裡悲嘆 (新譯本: 「就心裡激動,難過起來」；NIV : he was deeply moved in spirit and troubled) .耶穌心裡難過,憂傷引至哭泣,我們明白,但耶穌為何心裡激動哩!這字有發怒,憤慨的意思(anger, indignation) ,他為何憤慨呢?我們看看三十八節.當耶穌來到拉撒路的墳前,這字再次出現,此時耶穌不再哭了,不再難過了,但面對墳墓,他再度憤慨,發怒.耶穌悲傷,因為看到死亡帶來人的苦痛,死亡奪去拉撒路,死亡使馬大,馬利亞失去兄弟拉撒路而悲痛；他激動憤慨,因為看到死亡對人的宰制.

無論是才智聰穎的,愚昧無知的,富有的,貧窮的,都要在死亡面前俯首稱臣.有人嬰孩夭折,有人中年棄世,有人壽終正寢.耶穌來到墳前,就像一個摔角手,直視對手,要擊倒對手,要把對手-死亡-拿下來,所以在墳墓前,耶穌心裡激動憤慨.直視死亡,怒不可遏.不久他要進耶路撒冷,被釘十字架,藉著死敗壞那掌死權的,要釋放一生因怕死而為奴的人.

我們稱信徒離世是睡了.「我的朋友拉撒路睡了,我去叫醒他」(約十一11).「論到睡了的人,我們不願意弟兄們不知道,恐怕你們憂傷,像那些沒有指望的人一樣.我們若信耶穌死而復活了,那已經在耶穌裡睡了的人,神也必將他們與耶穌一同帶來.」(帖前四13, 14)安睡的人,在基督再來時都被叫醒,他們現在安睡的墓地都被視為神聖的空間.

今天下午有機會在一個公開講座,與天主教禮儀學者羅國輝神父對談「神聖空間,標記」.在預備過程中,查閱公教法典大全對神聖空間的定義,除了崇拜的教堂,還包括信徒的墓地.前一段時間,新聞提到幾位中國大陸遊客,被瑞典警方,從酒店送到墳塋.其實不是那麼恐怖,是把他們送到地鐵站附近,那裡有一著名教堂,教堂一如在歐洲或美國,附有墓地.信徒生前到教堂崇拜,死後也希望葬在教堂的墓地,有些大教堂,更把著名信眾和國家的聖王侯,偉人葬在教堂之內.西敏寺大教堂差不多每年公演莎士比亞名劇,因劇中帝侯就葬在教堂之內.

公教法典(Code of Canon Law) 信徒的墓地也是神聖空間.神聖空間指用於崇拜及埋葬信徒的空間,藉禮儀書為以上用途完成了呈獻或祝聖.(\#1205) (Sacred places are those which have been designated for divine worship or for the burial of the faithful through a dedication or blessing which the liturgical books prescribe for this purpose.)

法典把崇拜上主和埋葬信徒並列,因為離世的信徒安息在上主懷中,安葬他們的墓園也成為神聖了.

另一個原因,應和基督信仰的身體復活有關.神聖的墓穴:把永恆不朽的上帝和短暫朽壞的世界聯系在一起,在耶穌時期的泛希臘羅馬世界,是嶄新的信仰,產生前所未有的宗教現象-聖徒的崇敬,甚至敬拜(The cult of saints).「從亙古到永遠,你是上帝」是當時耳熟能詳的觀念,人們認為祇有神祇才是永恆不朽,凡人終歸塵土,人死後靈魂便脫離朽壞的皮囊,升到星河(Milky Way) .星河是古代人可目測到宇宙的至高處,那裡繁星顆顆,千萬年就停在穹蒼之頂,永恆不動絲毫,那裡平靜如鏡,不沾凡塵.

我們肉體在月亮之下流轉變幻,最後朽爛成一堆黃土；肉體的歸宿是地上,靈魂的歸宿在天上,各有所屬,不能逾越.肉體是靈魂的牢獄,死亡使靈魂得釋放,祇有羅馬皇帝死後連同肉體羽化升天(apotheosis) ,證明他們有神性,是神祇地上的代表,強化其神權為外衣的極權統治,但皇帝死後連同肉體升天也是一廂情願的想法.如果沒有復活的盼望,祇有像哥林多的人說:「我們就吃吃喝喝吧!因為明天死了」.(林前十五32)

若河漢是吾家,我們連同復活不朽的身體,也居住其間；那躺在地上聖徒的墳墓,就成為天堂人間的交匯處,這是前所未有的宗教現象,因為泛希臘羅馬文化,貶抑朽壞的肉身,不相信身體復活；認為短暫朽壞之軀,不可承受永恆之福. 但早期教會對聖徒的墳墓極其崇敬,他們的墓地成了朝聖者參觀祈禱的地方；他們的行狀,也成為信眾的榜樣,他們朽壞的身體,已連上不朽的上帝,分享上主永恆的生命.

結論:他有全備的智慧.死亡,連同隨伴死亡的諸般苦難和病疾是帕子,也許此刻我們不明白,但終有一天帕子除掉,眼淚被抹乾.基督說:我是始,我是終.沒有天意弄人的荒謬,祇有隱密的計劃.

他有捨命的愛.因為那位手潔心清,起誓不懷虛妄的基督,甚愛我們,為我們成為被殺的羔羊,帶同我們能登上聖山,升上至高者的右邊.因著愛,他嚐過死味,在我們臨終時,陪同我們行過死蔭的幽谷,不被死亡所傷害.

他有叫醒我們的能力.而且他視我們是朋友,死亡是安睡在上主懷中,到那天,他會叫醒我們.



\newpage

\section{音樂起了}
\label{sec:431}
\textbf{鄧瑞強博士 - 音樂起了}
\newline
\newline
\hyperref[sec:430]{< < < PREV SERMON < < <}
~
\hyperref[sec:toc]{[返目錄]}
~
\hyperref[sec:432]{> > > NEXT SERMON > > >}


啟示錄 14:1-13
\newline
\begin{longtable}{cl}
\hline
\hline
章節 & 經文 (和合本修訂版)\\
\hline
14:1 & \begin{tabularx}{0.7\textwidth}{X} 我又觀看，看見羔羊站在錫安山，和他在一起的有十四萬四千人，都有他的名和他父親的名寫在額上。 \end{tabularx} \\ \\ \relax
14:2 & \begin{tabularx}{0.7\textwidth}{X} 我聽見從天上有聲音，像眾水的聲音和大雷的聲音，我所聽見的聲音好像琴師所彈的琴聲。 \end{tabularx} \\ \\ \relax
14:3 & \begin{tabularx}{0.7\textwidth}{X} 他們在寶座前，和在四活物及眾長老前唱新歌，除了從地上買來的那十四萬四千人以外，沒有人能學這歌。 \end{tabularx} \\ \\ \relax
14:4 & \begin{tabularx}{0.7\textwidth}{X} 這些人未曾沾染婦女，他們原是童身。羔羊無論往哪裡去，他們都跟隨他。他們是從人間買來的，作為初熟的果子歸給神和羔羊。 \end{tabularx} \\ \\ \relax
14:5 & \begin{tabularx}{0.7\textwidth}{X} 在他們口中找不出謊言，他們是沒有瑕疵的。 \end{tabularx} \\ \\ \relax
14:6 & \begin{tabularx}{0.7\textwidth}{X} 我又看見另一位天使在空中飛翔，有永遠的福音要傳給住在地上的人，就是各邦國、各支派、各語言、各民族。 \end{tabularx} \\ \\ \relax
14:7 & \begin{tabularx}{0.7\textwidth}{X} 他大聲說：「要敬畏神，把榮耀歸給他，因為他施行審判的時候已經到了。要敬拜那創造天、地、海和水源的主。」 \end{tabularx} \\ \\ \relax
14:8 & \begin{tabularx}{0.7\textwidth}{X} 另有第二位天使接著說：「傾覆了！那曾叫列國喝淫亂、烈怒之酒的大巴比倫傾覆了！」 \end{tabularx} \\ \\ \relax
14:9 & \begin{tabularx}{0.7\textwidth}{X} 另有第三位天使接著他們，大聲說：「若有人拜那隻獸和獸像，在額上或在手上受了印記， \end{tabularx} \\ \\ \relax
14:10 & \begin{tabularx}{0.7\textwidth}{X} 他也必喝神烈怒的酒；這酒是斟在神憤怒的杯中的純酒。他要在聖天使和羔羊面前，在火與硫磺之中受痛苦。 \end{tabularx} \\ \\ \relax
14:11 & \begin{tabularx}{0.7\textwidth}{X} 使他們受痛苦的煙往上冒，直到永永遠遠。那些拜獸和獸像，受了牠名字的印記的人，晝夜不得安寧。」 \end{tabularx} \\ \\ \relax
14:12 & \begin{tabularx}{0.7\textwidth}{X} 在此，遵守神命令和堅信耶穌真道的聖徒要有耐心。 \end{tabularx} \\ \\ \relax
14:13 & \begin{tabularx}{0.7\textwidth}{X} 我聽見從天上有聲音說：「你要寫下：從今以後，在主裡死去的人有福了。」聖靈說：「是的，他們要從自己的勞苦中得安息，因為工作的成果永隨著他們。」 \end{tabularx} \\ \\
[1ex]
\hline
\hline
\end{longtable}
講題:音樂起了 　The Music Starts

經課:啟示錄14:1-13

講員:鄧瑞強博士

場合:香港中文大學崇基學院禮拜堂

日期:2018年11月11日

各位弟兄姊妹,早安.

上個月(2018年10月14日),天主教的教宗方濟各(Pope Francis)在聖伯多祿廣場宣布七位人士列為聖人.我們的信仰傳統沒有聖人,但有些信徒的確在信仰的路上留下深刻的軌跡,細味一下他們的生命,對我們的信仰人生,有很好的提示.這七位聖人中,有一位叫羅梅洛(Oscar Romero).他是誰?他是中美洲國家薩爾瓦多(El Salvador)人,是首都聖薩爾瓦多(San Salvador)教區的總主教(Archbishop).

他一生致力反對貧窮,社會不義,政治暗殺.教宗方濟各讚揚他為了親近窮人和自己的人民,不顧生死.(\url{https://www.theguardian.com/world/2018/oct/14/pope-francis-canonises-oscar-romero-and-pope-paul-vi})羅梅洛親近窮人,窮人是被遺忘的人.他親近自己的人民,因為他的同胞活在反人民的政權下.他走向受苦者,因為他認為一個被釘十架的上帝是一個受苦的上帝.受苦的上帝,親近受苦的人.

他努力尋求正義,但效果緩慢.他說:「天國不單在我們的能力以外,更在我們的視野之外.我們一生,只能完成神的偉大工程裡的一丁點兒.我們做的,無一樣是完成的.這是說,天國總在我們能力之外.沒有句子能完全表達所當說的.沒有祈禱能完全表達我們的信仰.沒有認罪能帶來完美.沒有牧者的探訪能帶來生命的整全.沒有活動節目能完成教會的使命.沒有任何目標十全十美.這就是我們的本來面目.我們只是播種,這種子有一天或會成長.我們澆水,盼望將來.我們建造地基,為了將來的工程.…這一切都不圓滿,但這是開始,向前踏上一步,這造成一種機遇,讓神的恩典能進入其中,然後帶來圓滿.我們或許永遠看不到最後成果,但這正是建築師和工人的分別.我們只是工人,不是建築師.我們是牧者,不是救世主.」(\url{https://apprenticeshiptojesus.wordpress.com/2007/09/12/oscar-romero-quotes/})

羅梅洛總主教是一個懷著希望的人,他在真理的國度上緩步前進.

有一齣電影,以《羅梅洛》(Romero (1989))為名,講述他的一生.在結尾部分,是一個彌撒的場合.羅梅洛講他人生中最後的一篇道.他說:「我對軍人作一特別呼籲.兄弟們,你們每一個都是我們的一分子.我們是同一樣的人民.你們殺死的農民是你們的兄弟姊妹.當你們聽到有人告訴你去殺一個人時,請想想神的話,『不可殺人』.沒有軍人可被強迫去執行一個違反上帝誡命的命令.以上帝之名,也以我們當中已受很多苦的人民之名,這些人民的哀求上達天庭,我懇求你們,我乞求你們,我命令你們,停止一切鎮壓.」

道講完了,當他舉起餅和杯祝聖聖餐時,一個槍手向他開了一槍,他的血就混和了聖餐的血,他的身軀倒在聖餐桌旁.按天主教神學的理解,聖餐是一個獻祭,聖餐桌是祭壇,羅梅洛名副其實地死在祭壇上.

在現實裡,羅梅洛總主教在聖餐桌前講了一段禱文,他說:「願那為人類被殺的〔基督〕身體和被獻上的〔基督〕寶血養育我們,好叫我們也奉獻自己的身體和血,承受痛苦,像基督一樣,不是為自己,而是為了將和平與正義的豐收帶給我們的人民.」(\url{https://apprenticeshiptojesus.wordpress.com/2007/09/12/oscar-romero-quotes/})這段禱文讀完後,子彈就穿過他的身體.

他死前幾日,曾對記者說:「你去告訴人們,若他們成功地殺死我,我會寬恕及祝福做這事的人.我希望他們知道,這是浪費他們的時間.一個主教死了,但神的教會,就是神的子民,是永遠不會消亡的.」(\url{https://www.uscatholic.org/culture/social-justice/2009/02/oscar-romero-bishop-poor})

這是一個「人」的故事,一個「真正的人」的故事.他真實地活在歷史的某一點,活在某個反人性的國家裡.他沒有變成一個失去人性的人,他活出人性,他爭取窮人的人性,他呼喚人間去聆聽神的呼喚.一個可以充滿神聖的世界,為何變成地獄?他生活的世界,認為他是一個不應存在的人.但在神的眼中,他是一個聖者.

今日講道的經文,描繪著同樣的張力.一個信徒在怪獸橫行的世界裡,是活在生與死之間.怪獸要消滅他,神卻將他放在永恆的國度裡.

講道經文:啟示錄14:1－13

啟14:1　我又觀看,見羔羊站在錫安山,同他又有十四萬四千人,都有他的名和他父的名寫在額上.

啟14:2　我聽見從天上有聲音,像眾水的聲音和大雷的聲音,並且我所聽見的好像彈琴的所彈的琴聲.

啟14:3　他們在寶座前,並在四活物和眾長老前唱歌,彷彿是新歌；除了從地上買來的那十四萬四千人以外,沒有人能學這歌.

啟14:4　這些人未曾沾染婦女,他們原是童身.羔羊無論往哪裡去,他們都跟隨他.他們是從人間買來的,作初熟的果子歸與神和羔羊.

啟14:5　在他們口中察不出謊言來；他們是沒有瑕疵的.

啟14:6　我又看見另有一位天使飛在空中,有永遠的福音要傳給住在地上的人,就是各國,各族,各方,各民.

啟14:7　他大聲說:「應當敬畏神,將榮耀歸給他!因他施行審判的時候已經到了.應當敬拜那創造天地海和眾水泉源的.」

啟14:8　又有第二位天使接著說:「叫萬民喝邪淫,大怒之酒的巴比倫大城傾倒了!傾倒了!」

啟14:9　又有第三位天使接著他們,大聲說:「若有人拜獸和獸像,在額上或在手上受了印記,

啟14:10　這人也必喝神大怒的酒；此酒斟在神忿怒的杯中純一不雜.他要在聖天使和羔羊面前,在火與硫磺之中受痛苦.

啟14:11　他受痛苦的煙往上冒,直到永永遠遠.那些拜獸和獸像,受牠名之印記的,晝夜不得安寧.」

啟14:12　聖徒的忍耐就在此；他們是守神誡命和耶穌真道的.

啟14:13　我聽見從天上有聲音說:「你要寫下:從今以後,在主裡面而死的人有福了!聖靈說:「是的,他們息了自己的勞苦,作工的果效也隨著他們.」

我將這經文分為五個段落.

(1)啟14:1-5:音樂起了

啟14:1　我又觀看,見羔羊站在錫安山,同他又有十四萬四千人,都有他的名和他父的名寫在額上.

世界是紛亂的,但在神給人的遠象裡,看到靈性世界的平靜祥和.在神聖的景象中,有「羔羊」.「羔羊」是被人宰殺的,但卻又復活過來的聖者.他站在錫安山上.這是王者的山,是得勝者的山.同他一起的,有十四萬四千忠誠的信徒.十四萬四千,代表著神的子民的總數.他們的額上,有基督和天父名號.這代表著他們的生命屬於神而不屬於這個世界.這個世界有希望,是因為有神,也因為這個世界總有活在世上卻又反抗世上的不義的人,他們不屬乎這個世界,他們屬乎神.

啟14:2　我聽見從天上有聲音,像眾水的聲音和大雷的聲音,並且我所聽見的好像彈琴的所彈的琴聲.

啟14:3　他們在寶座前,並在四活物和眾長老前唱歌,彷彿是新歌；除了從地上買來的那十四萬四千人以外,沒有人能學這歌.

音樂起了,這是天籟,這是世人不懂的生命之歌.

美妙的歌聲,總是誘人的.在希臘神話裡,有些像美人魚的女妖Siren,在岸邊唱歌,吸引水手聆聽,而陶醉失魂的水手,便會將船撞向礁石,令船沉沒.這或許就是世界的歌聲,令人意亂情迷,最後失去生命.啟示錄所講的天籟,都會令人像基督那樣放下生命.但這歌聲裡,講的是人性,講的是神聖,講的是一個可以想像卻又不可企及的天國.羅梅洛總主教聽過這歌聲,他追隨這歌聲,最後死在基督的祭壇上.

啟14:4　這些人未曾沾染婦女,他們原是童身.羔羊無論往哪裡去,他們都跟隨他.他們是從人間買來的,作初熟的果子歸與神和羔羊.

啟14:5　在他們口中察不出謊言來；他們是沒有瑕疵的.

這裡表達信徒的聖潔特質,他們緊緊地跟著羔羊耶穌.「在他們口中察不出謊言來」,這表達他們與真理為伍.在這虛假得明目張膽的時代,堅持真理的人,是難能可貴的.啟示錄講得明白,像基督一樣堅持真理的人,他們的生命也會像基督一樣,在死裡顯出永恆的生命.

(2)啟14:6-7:第一位天使的宣告

啟14:6　我又看見另有一位天使飛在空中,有永遠的福音要傳給住在地上的人,就是各國,各族,各方,各民.

啟14:7　他大聲說:「應當敬畏神,將榮耀歸給他!因他施行審判的時候已經到了.應當敬拜那創造天地海和眾水泉源的.」

這是第一位天使,他講論福音.福音是什麼?在這裡,福音是敬畏神,將榮耀歸給神.人總有敬畏的人或東西,人總會將榮耀歸給某些人或某些東西,但請不要敬畏一個錯誤的對象,也不要將榮耀歸給不該有榮耀的東西.請不要欣賞以虛謊上位的人,也不要戀慕強權下的虛名.在不正義的世界,有來自不正義的威脅.在威脅下就範,是人之常情.在這裡,這位天使鼓勵我們,不要就範,不要停止講真話,不要對正義的國度失去期盼.

人能在混亂的世界站得住腳,是因為他知道誰創造天地,誰掌管萬有.體會天地間有一浩然的正氣,這需要一點心靈的感應.這天使說,這天地的主宰要裁斷善惡,而這審判的時刻已到了.是非善惡,並不是混然不分的.人應明辯是非,在是非不分的時刻,毅然走在正路上.

(3)啟14:8:第二位天使的宣告

啟14:8　又有第二位天使接著說:「叫萬民喝邪淫,大怒之酒的巴比倫大城傾倒了!傾倒了!」

第二位天使宣告,那將人貶為奴隸的暴力國度巴比倫傾倒了.巴比倫,在啟示錄的時代,就是指羅馬帝國,一個軍力強大,巧取豪奪的帝國.天使說,這帝國的本質是「邪淫」的.在聖經傳統裡,「邪淫」的重點不在性的混亂,而在於拜偶像.「拜偶像」的神學定義很簡單,將「被造之物」看成是「神」,就是「拜偶像」.羅馬帝國敬拜的,是君王,是強大的軍力,是奢華的生活.他們以此為他們的「神」,他們敬拜君王,愛好征服,羨慕奢華.這些偶像無處不在,今日,這些偶像依舊誘惑我們去依仗權勢,掠奪窮人的財富,以過奢華的生活.

天使說,這邪淫的酒,是大怒的酒.或者這樣翻譯更簡明:「這邪淫的酒,引來神的大怒.」神不會對不義坐視不理.天使直接宣告,這不正義的體制和不正義的生活方式,傾倒了.若有人維繫這人間的不義體制,渴望過那不義的生活,將會和「巴比倫」一同滅亡.

(4)啟14:9-11:第三位天使的宣告

啟14:9　又有第三位天使接著他們,大聲說:「若有人拜獸和獸像,在額上或在手上受了印記,

啟14:10　這人也必喝神大怒的酒；此酒斟在神忿怒的杯中純一不雜.他要在聖天使和羔羊面前,在火與硫磺之中受痛苦.

啟14:11　他受痛苦的煙往上冒,直到永永遠遠.那些拜獸和獸像,受牠名之印記的,晝夜不得安寧.」

第三位天使提醒我們,千萬不要成為怪獸的一部分.

第9節及第11節,提及「拜獸和獸像」,這和第一位天使呼籲我們「敬拜真神」相對應.

中間的第10節,提及「神大怒的酒」,「神忿怒的杯」,這和第二位天使提及的「邪淫,大怒之酒」相對應.

明顯,這段經文提示我們,認清「怪獸世界」的虛假,認清虛假的生活世界的即將滅亡,呼籲我們千萬不要成為怪獸的一部分,否則,就會承受生命的永遠喪亡.這需要人十分清醒地活,要對世界的不義元素有清晰的了解,要對不義有明確的拒絕.

(5)啟14:12-13:結語

啟14:12　聖徒的忍耐就在此；他們是守神誡命和耶穌真道的.

啟14:13　我聽見從天上有聲音說:「你要寫下:從今以後,在主裡面而死的人有福了!聖靈說:「是的,他們息了自己的勞苦,作工的果效也隨著他們.」

在虛假的世界,嚮往真理,要一步一步地走過去,需要莫大的「忍耐」.

神學家莫特曼(Moltmann)在其《盼望神學》(Theology of Hope)裡這樣說,人有時過份狂妄,強求所希望的東西即時實現.人有時又太早絕望,認為所希望的東西永遠不會實現.其實,兩者都不是希望.真正的希望,具有一種「旅途中」的性質.在旅途中,向著目標前進,總需要「忍耐」.(參:Jürgen Moltmann: Theology of Hope (SCM Press, 1967), p.23.)

記得羅梅洛總主教的話:「天國不單在我們的能力以外,更在我們的視野之外.我們一生,只能完成神的偉大工程裡的一丁點兒.我們做的,無一樣是完成的.……這就是我們的本來面目.我們只是播種,這種子有一天或會成長.我們澆水,盼望將來.……這一切都不圓滿,但這是開始,向前踏上一步,這造成一種機遇,讓神的恩典能進入其中,然後帶來圓滿.」

天上的聲音說:「在主裡面而死的人有福了!」羅梅洛死在祭壇上,他會說:這種死,是一種蒙福的死.

聖靈說:「是的,他們息了自己的勞苦,作工的果效也隨著他們.」這是在安息禮拜裡經常會誦讀的經文.真的,一個人死了,不就是死了.他一生的惡,會永隨他.他一生的善,也會永隨他.我深深希望,在虛謊的世界,我們永遠堅持真理,永遠堅持正義,永遠堅持良善.這一切一切,會永遠伴隨我們.

但願:榮耀歸於聖父,聖子,聖靈.阿們.



\newpage

\section{主愛我們到底}
\label{sec:432}
\textbf{劉國偉先生 - 主愛我們到底}
\newline
\newline
\hyperref[sec:431]{< < < PREV SERMON < < <}
~
\hyperref[sec:toc]{[返目錄]}
~
\hyperref[sec:433]{> > > NEXT SERMON > > >}


約翰福音 13:1-15
\newline
\begin{longtable}{cl}
\hline
\hline
章節 & 經文 (和合本修訂版)\\
\hline
13:1 & \begin{tabularx}{0.7\textwidth}{X} 逾越節以前，耶穌知道自己離世歸父的時候到了。他一向愛世間屬自己的人，就愛他們到底。 \end{tabularx} \\ \\ \relax
13:2 & \begin{tabularx}{0.7\textwidth}{X} 晚餐的時候，魔鬼已把出賣耶穌的意思放在加略人西門的兒子猶大心裡。 \end{tabularx} \\ \\ \relax
13:3 & \begin{tabularx}{0.7\textwidth}{X} 耶穌知道父已把萬有交在他手裡，且知道自己是從神出來的，又要回到神那裡去， \end{tabularx} \\ \\ \relax
13:4 & \begin{tabularx}{0.7\textwidth}{X} 就離席站起來，脫了衣服，拿一條手巾束腰， \end{tabularx} \\ \\ \relax
13:5 & \begin{tabularx}{0.7\textwidth}{X} 隨後把水倒在盆裡，開始洗門徒的腳，並用束腰的手巾擦乾。 \end{tabularx} \\ \\ \relax
13:6 & \begin{tabularx}{0.7\textwidth}{X} 到了西門‧彼得跟前，彼得對他說：「主啊，你洗我的腳嗎？」 \end{tabularx} \\ \\ \relax
13:7 & \begin{tabularx}{0.7\textwidth}{X} 耶穌回答他說：「我所做的，你現在不知道，但以後會明白。」 \end{tabularx} \\ \\ \relax
13:8 & \begin{tabularx}{0.7\textwidth}{X} 彼得對他說：「你絕對不可以洗我的腳！」耶穌回答他：「我若不洗你，你就與我無份了。」 \end{tabularx} \\ \\ \relax
13:9 & \begin{tabularx}{0.7\textwidth}{X} 西門‧彼得對他說：「主啊，不僅是我的腳，連手和頭也要洗！」 \end{tabularx} \\ \\ \relax
13:10 & \begin{tabularx}{0.7\textwidth}{X} 耶穌對他說：「凡洗過澡的人不需要再洗，只要把腳一洗，全身就乾淨了。你們是乾淨的，然而不都是乾淨的。」 \end{tabularx} \\ \\ \relax
13:11 & \begin{tabularx}{0.7\textwidth}{X} 耶穌已知道要出賣他的是誰，因此說「你們不都是乾淨的」。 \end{tabularx} \\ \\ \relax
13:12 & \begin{tabularx}{0.7\textwidth}{X} 耶穌洗完了他們的腳，就穿上衣服，又坐下，對他們說：「我為你們所做的，你們明白嗎？ \end{tabularx} \\ \\ \relax
13:13 & \begin{tabularx}{0.7\textwidth}{X} 你們稱呼我老師，稱呼我主，你們說的不錯，我本來就是。 \end{tabularx} \\ \\ \relax
13:14 & \begin{tabularx}{0.7\textwidth}{X} 我是你們的主，你們的老師，尚且洗你們的腳，你們也應當彼此洗腳。 \end{tabularx} \\ \\ \relax
13:15 & \begin{tabularx}{0.7\textwidth}{X} 我給你們作了榜樣，為要你們照著我為你們所做的去做。 \end{tabularx} \\ \\
[1ex]
\hline
\hline
\end{longtable}
講題:主愛我們到底 　The Full Extent of His Love

經課:約翰福音13章1-15節

講員:劉國偉先生

場合:香港中文大學崇基學院禮拜堂

日期:2018年11月18日

你可以去到幾盡?

2013年港產電影《狂舞派》上映,引發一鼓熱潮.故事講述在豆腐店長大的少女阿花,自小熱愛跳舞,她參加了大學舞團 BombA,卻可惜因為內部不和,阿花憤然離隊；後來當BombA隊長和隊員消失鬥志,阿花決定重返舞團,準備為舞團參加舞林大賽,但卻遇上意外而導致嚴重腿傷,無法出賽,正當坐在輪椅上的阿花感到夢想破滅之際,主要對手的隊長出現在她面前,問了她一句:「So how far are you willing to go for dance ? (為了跳舞,你可以去到幾盡?)」

當年引起的熱潮,主要不是跳舞,而是這電影的金句: 為了跳舞,你可以去到幾盡?過去5年,我在好幾次基督徒聚集的場合,聽過有人套用這電影金句: 為了耶穌,為了信仰,你可以去到幾盡?這確實是一個不容易回答,或者我們試圖逃避的問題.不過,在我們考慮如何回答這個問題之前,我今天與大家一起思考到底耶穌為了愛我們,去到幾盡.

約翰福音由第12章開始,進入了耶穌在世最後一個星期的敘事.第13章更開展了耶穌被捉拿前與門徒相聚的一個重要片段,故事情節充滿張力,直至全卷書的結尾.作者使徒約翰以十分獨特的引子,把他心目中的耶穌如此描繪出來:

「逾越節以前,耶穌知道自己離世歸父的時候到了.他既然愛世間屬自己的人,就愛他們到底.」(約翰福音13章1節) 用我們的地道說話來演繹,大概就是「他既然愛世間屬自己的人,就去到盡.」

愛他們到底

「逾越節以前,耶穌知道自己離世歸父的時候到了.他既然愛世間屬自己的人,就愛他們到底. 」(約翰福音13章1節)

一般相信,約翰福音的成書時間在主後85至90年左右,那時,約翰已經是一個長者,同輩的使徒都已經被主接去,留在他的記憶中的往事,成為他十分重要的生命伴侶.他深情的一句「就愛他們到底」,憶述半世紀以前的片段,那必定是彌足珍貴的記憶.

「就愛他們到底」這句話開啟了往後接連五章經文,就是在那個吃晚餐的樓房內的敘事.因為「愛他們到底」,耶穌給予門徒心理準備,預告自己被出賣及彼得三次不認主；因為「愛到底」,他吩咐門徒要彼此相愛；「愛到底」,耶穌勸導門徒不要難過,並鼓舞他們及告訴他們,祂就是道路真理生命,並且賜下保惠師,又留下平安給他們；耶穌更加告訴門徒,祂已經勝過世界；當然,因為「愛到底」,耶穌作了一個深情的禱告,並且以這句話來結束:「使你所愛我的愛在他們裡面,我也在他們裡面.」

耶穌為門徒洗腳

「耶穌就離席站起來脫了衣服,拿一條手巾束腰,隨後把水倒在盆裡,就洗門徒的腳,並用自己所束的手巾擦乾.」(約翰福音13章4-5節)

使徒約翰,一個以「耶穌所愛的那門徒」來形容自己的使徒,採用了異常細膩的筆觸,去描述這件陳年往事────這是已經存在於他的腦海中,超過半個世紀的片段.年老的約翰接連用了至少八個動詞去描繪耶穌的動作:離席,站起來,脫衣服,拿手巾,束腰,倒水,洗腳,擦乾,這種寫作手法,在約翰福音中並不多見,按整本聖經文本而言也甚為罕見,正因如此,愈發突顯這兩節經文的獨特之處.

我們知道在巴靳斯坦的古老風俗中,為賓客洗腳是奴僕的職責,是一種卑微的服侍,作為拉比的耶穌,竟然擔當這個卑微的角色,是相當使人訝異的,至少在約翰的回憶中,整件事情真的仿如昨日,依然歷歷在目,他記得每一個細節,耶穌每一個動作都在他心中留下烙印,過了半個世紀,他依然難以忘懷,並要逐筆逐句的,把他的深刻印象繪畫出來.

在人生旅程中,我們都相信耶穌與我們同行,但卻不一定察覺耶穌在我們的生活中所留下的痕跡.正如我們都認同耶穌愛我們,卻未必可以在貼地的生活中,擁有「耶穌愛我們到底」的經驗.有些人會為巿井生涯中的點點滴滴,甚至看似平凡的片段和故事,深深感謝耶穌的同在,細味耶穌的足印；但也有些人,把生活中的許多許多場景,視為理所當然,彷彿一切在自己的掌握之中,遺漏了耶穌為我們所做的一切,祂的恩典,祂的同行,不曾是生命中難忘的畫面.

你洗我的腳麼?

彼得對祂說:「主啊!你洗我的腳麼?」耶穌回答說:「我所作的,你如今不知道,後來必明白.」彼得說:「你永不可洗我的腳.」耶穌說:「我若不洗你,你就與我無份了.」(約十三6-8 )

西門彼得對主的發問,表達了他心中的疑惑:尊貴的拉比怎可洗我的腳!正如我們一樣,彼得受既有的思維所局限,在這個氣氛異於尋常的樓房裡,彼得無法頓悟耶穌為門徒洗腳的意義.耶穌並沒有因此不洗彼得的腳,祂耐心地向彼得說:「你如今不知道,後來必明白.」表達出主對彼得的體諒,接納他所面臨的局限,理解他無法參透即將走上十字架的主的言行舉措.

看人間略影,總常有「如今不知道」的局面,年老的使徒約翰回首前事,細緻地記下了彼得和耶穌的對話,也許走過半個世紀的人生路,約翰會充份體味「後來必明白」的生活智慧.耶穌這句話除了包含體諒,接納,更告訴門徒一種生活智慧:人生之中有些事情總要經過生活的熬煉,歲月的洗滌,並智慧的滋長,方能明白箇中含意,很顯然的,彼得未能立時掌握,就像我們在生活中常有的經驗:我們會很有「效率」地,把以為成熟的「結論」表達出來:「你永不可洗我的腳.」

從彼得的疑惑到他的回應,一而再顯示彼得無法消化當前的處景:耶穌面向十字架的刑場,門徒卻活在爭論誰為大的不愉快經歷之中；三年的同枕共食生活,終於來到尾聲…….在面向太多的不明朗前景,緊抓著舊有的思維方式,以一種相對「安全」的思路去理解眼前的景象,似乎是一種自然的選擇,但也注定沒有可能去解讀一個全新的畫像,而仍然呆站於「你永不可洗我的腳」的思維之中.

耶穌並沒有給彼得長篇大論的解說.只輕輕的說:「我若不洗你,你就與我無份了.」耶穌帶出了一個核心的信息,祂所作的一切,都要叫人與祂「有份」,有份於祂的生命,祂的呼召,祂的奉獻,祂的虛己……讓耶穌服侍我們,就是讓我們有份於祂,這是何等奇特的觀念,在我們人間庸俗的宗教情操中,從來只有凡人向神明膜拜朝奉,然後得一點恩惠,怎會有道成肉身的主宰,降下凡塵為人洗腳,最後把生命也捨棄?

人生路上,市井之中,不明白不理解的事情和際遇不知凡幾,活在舊思維中的彼得和我們,也常會對生命作出「九唔搭八」的回應,若果我們要經歷心意更新,似乎先要放下舊有的牢囿,也許這是獻上活祭────生命中一種死亡的經歷的更深一層含義.

祂在尋找孤單的人

耶穌說:「……你們是乾淨的,然而不都是乾淨的.」耶穌原知道要賣他的是誰,所以說,你們不都是乾淨的.(約十三10-11)

按照路加的記載,在最後晚餐的一夜,門徒起了爭論,他們中間那一個可算為大,(路廿二24)這已不是頭一趙,據馬太和馬可的記載,在耶穌最後一週行程,進入耶路撒冷之前,西庇太兄弟(也包括他們的母親,參太廿20及可十35)也曾引起一次類似的爭論.

再想一想這個晚上接著要發生的事情:耶穌對門徒的深入談話,分離的禱告,客西馬尼園的經歷,被賣,捉拿,審訊……這確是一個沈重的夜晚.

你可會想像當耶穌說:「你們是乾淨的,然而不都是乾淨的.」祂是在一種怎樣的心境和情懷之中?三載歲月,門徒在這關鍵時刻,竟爭論誰為大；在日夕共處的團隊中,竟有人出賣自己的老師!那時,耶穌正在為門徒洗腳,祂可能蹲在彼得的腳旁,正要服事他,就在與彼得對話的過程中,說出了這一句當時沒有人明白的話.

耶穌面對這個沈重的夜晚,眼看著即將發生的一切,就在替門徒洗腳的時候,祂說出門徒當中有人心懷不乾淨的念頭,總給人幾分感慨,幾分心痛的味道.但這種經歷並沒有人明白,那種孤單是何等的深!這一夜,看來有親密的門徒相伴,但當下明白祂的恐怕沒有一個,直至夜深,在客西馬尼園的掙扎,仍是在一種「有人在身傍,卻沒有共鳴」的孤單境遇之中.

人海之中,我們都遇過孤單的時刻,有人因此選擇走在大隊之中,以避過這種孤單的滋味,也有人孤身上路.無論如何,當我們發覺自己身處同伴之中,卻沒有人可以明白自己內心的說話,那種並非孤獨,但卻孤單的心境,惹人苦惱,也覺沈重.

孤單的主,在道成肉身住在人間的旅程中,要尋找孤單的人,告訴他們生命並不孤單.耶利米哀歌寫道:「耶和華啊,求你紀念我們所遭遇的事……我們是無父的孤兒……」(哀五1,3)約翰記下了耶穌在當晚說的一句話,恰好是對哀歌的回應:「我不撇下你們為孤兒,我必到你們這裡來.」(約十四18)在人類的歷史長廊中,是耶穌為我們走在孤單的道路上,卻因為祂的同在,我們可以面對生命中孤單的時刻.

你們也當彼此洗腳

我是你們的主,你們的夫子,尚且洗你們的腳,你們也當彼此洗腳.我給你們作了榜樣,叫你們照著我向你們所做的去做.(約十三14-15)

年老的約翰如此寫道:「祂既然愛世間屬自己的人,就愛他們到底.」(約13:1下),這是約翰要告訴我們的故事背景,他體味到夫子為門徒洗腳,表達出「門徒與主有份」(參約13:8下)那種生命連繫,也突顯出「你們當彼此洗腳」的教導(約13:14下),「彼此洗腳」帶有彼此有份及互相為僕的強烈意象,相信對使徒約翰而言,這些都是深藏心底又記憶猶新的教導.

在這一切畫面和聲音的背後,是主同行的愛,是道成了肉身,充充滿滿有恩典的住在人間.在人生旅程中,耶穌何時為我們拿起手巾束腰,把水倒在盆裡,為每一個不配的我們洗腳,再用溫柔的手,以束腰的手巾把我們的雙腳擦乾?

是祂叫我有份於祂的恩情,以奴僕的角色和捨己的行動,呼召我們去跟隨祂.

New International Version如此翻譯約翰福音13章1 節: Having loved his own who were in the world, he now showed them the full extent of his love.

約翰告訴我們:就是耶穌愛他們到底的愛情,是一種「去到盡」的愛(Full extent of his love)

耶穌愛到底,去到盡?……13章至17章記載了在樓房內的事情,由第18章開始,耶穌經歷被捕,審問,交押,判刑,被釘,受死,埋葬,復活與顯現,實在好盡好盡!

北美著名基督徒作家及神學家唐曼華(Marva Dawn)寫道……

Offering our bodies has all sorts of practical implications…. We cannot be Christians merely intellectually; rather, we must respond to God’s love by loving him with our words and attitudes, our emotions and actions.

面對耶穌愛我們到底的愛情,一種「去到盡」的愛(Full extent of his love),我們實在需要學習在生活之中,在市井的每一個角落裡,細心察看耶穌在我們生命中留下的痕跡,祂叫我們與祂有份,祂等待我們更明白祂,在孤單旅程中與人同行,在我們沒有彼此洗腳的時候,He now shows us the full extent of his love.



\newpage

\section{轉化心靈的教育}
\label{sec:433}
\textbf{吳梓明教授 - 轉化心靈的教育}
\newline
\newline
\hyperref[sec:432]{< < < PREV SERMON < < <}
~
\hyperref[sec:toc]{[返目錄]}
~
\hyperref[sec:434]{> > > NEXT SERMON > > >}


約翰一書 4:7-8,12,16
\newline
\begin{longtable}{cl}
\hline
\hline
章節 & 經文 (和合本修訂版)\\
\hline
4:7 & \begin{tabularx}{0.7\textwidth}{X} 親愛的，我們要彼此相愛，因為愛是從神來的。凡有愛的都是由神而生，並且認識神。 \end{tabularx} \\ \\ \relax
4:8 & \begin{tabularx}{0.7\textwidth}{X} 沒有愛的就不認識神，因為神就是愛。 \end{tabularx} \\ \\ \relax
4:9 & \begin{tabularx}{0.7\textwidth}{X} 神差他獨一的兒子到世上來，使我們藉著他得生命；由此，神對我們的愛就顯明了。 \end{tabularx} \\ \\ \relax
4:10 & \begin{tabularx}{0.7\textwidth}{X} 不是我們愛神，而是神愛我們，差他的兒子為我們的罪作了贖罪祭；這就是愛。 \end{tabularx} \\ \\ \relax
4:11 & \begin{tabularx}{0.7\textwidth}{X} 親愛的，既然神這樣愛我們，我們也要彼此相愛。 \end{tabularx} \\ \\ \relax
4:12 & \begin{tabularx}{0.7\textwidth}{X} 從來沒有人見過神，我們若彼此相愛，神就住在我們裡面，他的愛在我們裡面得以完滿了。 \end{tabularx} \\ \\ \relax
4:13 & \begin{tabularx}{0.7\textwidth}{X} 因為神將他的靈賜給我們，由此我們知道我們是住在他裡面，而他也住在我們裡面。 \end{tabularx} \\ \\ \relax
4:14 & \begin{tabularx}{0.7\textwidth}{X} 父差子作世人的救主，這是我們所看見並且作見證的。 \end{tabularx} \\ \\ \relax
4:15 & \begin{tabularx}{0.7\textwidth}{X} 凡宣認耶穌為神兒子的，神就住在他裡面，而他也住在神裡面。 \end{tabularx} \\ \\ \relax
4:16 & \begin{tabularx}{0.7\textwidth}{X} 我們知道並且深信神是愛我們的。神就是愛，住在愛裡面的就是住在神裡面；神也住在他裡面。 \end{tabularx} \\ \\
[1ex]
\hline
\hline
\end{longtable}
講題:轉化心靈的教育 　Beyond Education of the Mind

講員:吳梓明教授

經文:約翰一書4:7-8,12,16

場合:香港中文大學崇基學院禮拜堂

日期:2018年11月25日

請先聽我讀出一段經文:『親愛的弟兄(姊妹)啊! 我們應當彼此相愛,因為愛是從上帝而來的; 凡有愛心的,都是由上帝而生,並且認識上帝… 因為上帝就是愛.… 我們若彼此相愛,上帝就住在我們裡面,愛他的心在我們裡面得以完全了… 上帝就是愛,住在愛裡面的,就是住在上帝裡面,上帝也住在他裡面』(約翰壹書 4:7-8,12,16).

各位主內的弟兄姊妹,及各位崇基友好平安. 很多謝林豪恩弟兄的邀請,讓我今日有機會與大家- 崇基學院禮拜堂的兄弟姊妹們一同聚集,一同敬拜上帝. 回到崇基,我也自然地想及教育的問題,因為崇基給我們的,是不一樣的教育,它是一種超越純理性的教育(beyond the education of the mind),更是一種關懷及轉化心靈的教育. 所以我今日揀選這個題目,與大家一同思考『轉化心靈的教育』.

大家有否留意到?原來我們一生是離不開教育的,出世的時候我們在家中接受家庭教育,長大了被送到學校接受學校教育；出來工作後,仍是不斷地跟上司學習做事的方法,體驗做人的道理,現在我們也許還需要教導我們的下屬學識做事的方法和做人的道理；甚至若果我們是為人父母的話,我們一生都要肩負教育子女的責任.就算你說已經放棄了這個責任,你仍然是以自己的生命影響子女和其他人的生命,這就是教育的工作.當我們從事教育工作的時候,又有否留意到我們身邊的人原來不只是一塊木頭(無生命,無腦的東西),或是一隻動物(如愚蠢的豬,身邊的貓狗)而已,事實上,在接受我們教育或影響的委實是一群有活潑心靈的人.所以,教育也不僅是學習一些技能,一些做事的方法,甚或是做人的道理,它更也是一種『轉化心靈,以生命影響生命』,是一種極富有意義的活動.

九月期間,我曾回到崇基學院,與崇基校友們分享<我所敬仰的凌道揚博士>. 凌道揚(1888-1993)博士是我們崇基學院第二任的校長(1955-1960),他也是民國時期在中國提倡『農林救國』,推廣農林科學知識教育,並倡議『中國植樹節』的農林學家.凌道揚不單是從事農林科學教育,也是從事基督教教育的工作,當他被委任為崇基學院第二任校長的時候,他更提倡了『以靈育代替美育』一種新的教育取向.一般人講『五育』,是指『德,智,體,群,美』,基督教教育傳入中國之後,蔡元培先生(北大的校長)曾經建議『以美育代宗教教育』(他不是反對宗教教育,只是建議擴大美育的範疇去吸納宗教教育,因為宗教教育也是追求『真,善,美』的教育),但凌道揚博士做崇基校長的時候卻是提出了另一個創新中國教育文化的反建議,即是『以靈育代替美育』. 他認為:『崇基(學院)的宗旨是要注重德,靈,智,體,群五育』[1] .所以說:崇基的教育是『不一樣的五育』,它是『以靈育代替美育的基督教教育』,因為基督教教育除了是尋真善美的教育外,更也是一種『轉化心靈的教育』.

我一直來也關心宗教與教育的關係,我發現原來基督教教育就是一種幫助人轉化生命的教育(基督教是勸人悔改,幫助人提升生命的素質,結出聖靈的果子; 保羅說:『聖靈所結的果子,就是仁愛,喜樂,和平,忍耐,恩慈,良善,信實,溫柔和節制』(加拉太書 5:22-23)；也是尋求真善美,信望愛等高質素的生命教育),也即是一種『轉化心靈』的教育,我希望今天能夠與大家分享我近年的一些發現.

三年前,我獲邀請協助校牧室帶團往國內考察有關中國教會大學的歷史,因為它也是與崇基學院的基督教教育有關的,這些歷史不僅是教會大學的歷史,更也是中國基督教舊約的歷史.我們時常問:『基督教教育是甚麼?』,究竟過去的200多年來,基督(新)教傳教士將怎麼樣的基督教帶來了中國,是甚麼樣的基督教教育?今日仍然能夠在中國人心中留下深刻印象,對中國社會仍具深遠的影響力的,又是怎麼樣的基督教教育呢?

2016年我們曾造訪南京師範大學,遇到一位金陵女子大學的校友- 甘克超女士. 當我們遇到她的時候,她已經是91歲了. 她是1948年畢業,至今剛巧已70年了；她曾親身見證新中國的建立,經歷過文化大革命的滄桑歲月,雖然現在已年過九十,但她仍沒有忘記基督教教育對她的影響,她仍然能夠勇敢地對我們說:是『金女大改變了我的生命,使我成為了今日的我』. 這真是一個十分感人的見證!她也不斷地提醒我們,金女大的校訓是:『厚生』,是指人的生命不僅是要活下去,更是要活得精彩,活得豐盛的生命；她解釋說,『厚生』(abundant life)是『豐盛的生命』,就是能夠幫助自己和別人,實踐基督服務的精神,以生命影響生命.

2017年我們到了武漢華中師範大學訪問一位退休的校長- 章開沅教授.他是我們崇基學院的老朋友,是崇基學院1994-1995年度黃林秀蓮訪問學人. 我是1989年已經認識章教授的,他當時是華中師範大學的校長,他邀請我參加了5月31日到6月3日在華中師範大學舉行的一個學術研討會,所以6月4日我還在武漢休息,6月5日才返回香港的.章開沅教授原本是在南京金陵大學讀書,後來是在華中師範大學(前身是華中大學)教書及當校長,兩所均是與教會大學有關. 更特別的是當他擔任校長一職時,他卻獲得一個奇特的挑戰(是上天給他的一個呼召) [2],開始推動有關中國教會大學史的研究；他本來是辛亥革命史的專家,後來更因為他從事教會大學史的研究,現在也成為了南京大屠殺歷史研究的專家,專門研究抗日戰爭期間位於南京的兩所教會大學(金陵大學和金陵女子大學)及當中的傳教士們如何救護中國難民及見證南京大屠殺的歷史事實,亦間接確認了傳教士當時是如何奮不顧身地救護中國的人民. 在2017年的分享中,章教授(已經是91歲的老人家)總結他在教會大學工作的經驗時,還確鑿地說:『基督教就是愛的宗教』,教會學校讓他學會了『教育的靈魂就是愛,沒有愛的教育就是沒有靈魂的教育』；而他印象中的基督教教育就是一種表達基督的大愛和服務精神的教育,是透過基督的愛而轉化中國人民生命的教育.

大家有否到過南京師範大學? 我們也參觀過南師大的教育大樓,在教育樓前面亦豎立了另一個銅像 - 是陶行知博士的銅像. 陶行知是誰?大家有聽過這個名字嗎?陶行知是近代中國一位偉大的教育家,在新中國成立之初,國家主席毛澤東與周恩來總理均尊稱他為『中國人民教育家』. 所以,新中國的教育學者們都是十分敬重陶行知博士,尊崇他為近代中國一位偉大的教育家. 陶行知主張甚麼?大家或許會聽過他的一句名言:『捧著一顆心來,不帶半根草走』. 他就是主張:『愛的教育』和『愛滿天下』的,有否聽過『四粒糖』的故事嗎?『有一次,陶行知遇到一個學生,發現他曾經用泥塊砸自己班上的同學,陶校長立時制止他,並命令他放學後立刻到校長室去.放學後,陶行知來到校長室時,那位學生早已等著校長,準備挨聽校長的訓話了.可是陶行知卻笑著掏出一顆糖果送給他,對他說:“這是獎給你的,因為你準時來到校長室,而我卻遲到了.”那學生接過糖果後,校長隨即又掏出第二顆糖果放到他的手裡,說:“這也是獎勵你的,因為我不准你出手打人時,你就立即住手了,這亦表明你很尊重我,我應該是獎勵你的.”那學生很驚訝地望著校長.這時候,校長又掏出第三顆糖果塞到學生手裡,說:“我調查過了,你用泥塊砸那些同學,是因為他們欺負其他女生；你砸他們是說明你是一個很正直又善良的學生,你更且是有跟壞人作鬥爭的勇氣,我也是應該獎勵你的!”那些學生十分感動,他即時流著眼淚並帶著後悔的態度說:“校長,我知錯了,我砸的不是壞人,而是我的同學……”陶行知校長終於滿意地對著他,隨即又掏出第四顆糖果給那位學生,說:“為了你能夠正確地認識自己的錯誤,我再獎給你一塊糖果；好了,我沒有更多的糖果了,我們的談話也可以結束吧!』.

奇怪嗎? 為甚麼一個中國教育家會講『愛的教育』?  我們會否質疑: 這絕對不會是從中國的教育,哲學,文化傳統出來的.『愛的教育』是從何而來?你們會這樣問嗎?(大約是20年前,我的一位博士生就是提出這樣的問題,終於完成了他的博士論文,他的博士論文就是那麼簡單,是要證明今日中國人- 國內教育界學者們所尊崇的『中國人民教育家 - 陶行知』原來是一位基督徒教育家,他所倡議的『愛的教育』其實就是『基督教教育』) [3]. 陶行知原來是在南京金陵大學(一所教會大學)讀書,在大學期間信教,1914年他有機會往美國深造,本想讀政治,為孫中山先生建立的新中國服務；但1915年夏天,他在美國參加了YMCA青年會為中國留學生舉辦的夏令會,並決志獻身以教育救國. 回國後,他與蔡元培,胡適等推動中國的教育改革,但他獨倡平民教育和愛的教育.在1927年,陶行知創辦曉莊學校,是一所實驗師範學校. 曉莊師範學校並不是一所基督教學校,它卻是注入了基督教信仰的一個重要的元素:『愛的教育』. 陶行知這樣說:『曉莊是從愛裡產生出來的.沒有愛便沒有曉莊.… 愛之所在即曉莊所在.』.[4]  『曉莊是從愛裡產生出來的』,這樣的思想是從那裡來的?對中國傳統文化教育來說,它不會是來自中國的哲學或是甚麼的主義思想吧；『愛的教育』,『愛滿天下』就是一種創新的文化(教育)思維. 大家記得我開始的時候所讀的經文嗎?這就是陶行知最熟識的聖經教導:『因為愛是從上帝而來的!凡有愛心的,都是由上帝而生,並且認識上帝. 沒有愛心的,就不認識上帝; 因為上帝就是愛』.[5] 對陶行知來說,基督教教育就是愛的教育,因為只有從上帝而來的愛,才能改變人的生命,才能轉化人的心靈.而陶行知相信這就是中國人所需要的宗教教育,也是最容易被中國人所接納的一種基督教教育.

傳教士帶來中國的基督教就是愛的宗教,這份愛不僅是教育方面表現出來,成為愛的教育,也可以在醫療及社會服務方面,見證著上帝的愛,讓中國人感受到上帝的愛,生命也被改變過來. 今年6月,我們曾經訪問四川大學,尋找崇基學院副校長啟真道(Leslie Kilborn,1895-1972)博士的根,因為啟博士是在四川成都樂山出世,家族三代都是在中國從事醫療教育的工作,啟博士擔任過成都華西協合大學醫學院的院長.他的父親啟爾德(Omar Kilborn)在1894年在成都開設了第一所西醫醫院,透過醫療服務見證上帝的大愛,愛及每一個身,心,靈有需要的中國人.這所醫院就是今日我們所認識的成都市第二人民醫院.

有否想過? 我們現在是努力地為主作工,100年後會是怎樣呢?會否仍有人記得我們的工作嗎?華西傳教士是在差不多130年前來中國從事醫療教育和宣教的工作,他們雖然在1950年代被趕出中國,但他們的工作卻並沒有白費. 雖然是經歷過文革,天災,人禍等,但中國人民仍是懂得飲水思源,懂得感恩的!(前幾天我們有否慶祝美國人的感恩節?原來中國人也是懂得感恩的!) 他們並沒有忘記傳教士百多年前所帶來的基督大愛的精神的.

大家還記得另一位英國的傳教士李提摩太(Timothy Richard, 1845-1919)嗎?他曾建議必須用庚子賠款為中國人興建大學,提升中國人民的教育水平；他在1902年發動興建山西大學堂,到了2002年,當山西大學慶祝創校100周年時,他們仍然紀念李提摩太的功績,心存感恩. 大概是七年前(2011年),一位四川大學畢業的中國學者李雅詩(Ms. Aisi Li) 剛在英國完成她的博士論文- 是專門研究李提摩太與山西大學的.她在2011年約同一位加拿大的教授 (Ruth Hayhoe 許美基督教授,她是近代中國高等教育史的專家,著有China’s Universities, 1895-1995),另一位是美國的學者 (Eunice Johson,她曾在山西省教授英語,後來回美國完成她的博士論文,題目是:Timothy Richard’s Vision: Education and Reform in China, 1880-1910),他們三人一同到英國倫敦參拜這位英國傳教士李提摩太的墓. 奇妙嗎?這四位不同國籍的人仕,他們有著不同的信念和關懷,也是活在不同的年代中,卻是因為見證基督的大愛,使他們連在一起. 這位中國的學者當時還滿懷感觸地寫下一段感想,說:

『我希望有一天,當華人學生想及中國高等教育之時,他們不會忘記傳教士在現代中國及其大學教育中所曾扮演的重要角色.我亦衷心盼望: 這些傳教士在他們的工作上所表現的熱忱,堅毅,及信心,能夠激勵年青的華人學生,有同樣的信心,排除萬難,追尋實現他們的夢想』.[6]

剛才提及啟真道博士一家三代在成都從事的醫療教育和服務的工作,也同樣地被歌頌和記念,是因為他們所宣揚的是基督的大愛,實踐愛的教育,他們更是身體力行,活現轉化心靈的教育,改變了許多(千千萬萬) 中國人的生命.

我上個月還有機會在杭州參加了一個很特別的學術研討會,是由浙江省兩會及浙江省神學院所主辦,為紀念中國20世紀一位傑出的神學家,他也是一位愛國愛教的神學家- 趙紫宸教授(1888年生) 誕生130周年的學術研討會. 大家有否聽過趙紫宸這個名字?他本來是蘇州東吳大學一位傑出的教授,何明華會督很欣賞他的才華,1941年時安排他來香港,在主教府為他舉行一個三級跳的儀式- (1) 是從循道會轉為聖公會的教友; (2) 是按立他為聖公會的會吏；(3) 是在同一個儀式中按立他為牧師,所以稱之為三級跳. 由此亦可以見到趙紫宸真是一個不平凡的人,他後來在燕京大學教書,擔任燕京大學宗教學院院長.他在1980年在北京離世,香港聖公會卻為他在聖約翰座堂舉行一個追思禮拜. 後來,他的家鄉浙江省胡州德清縣新市鎮又為他建設一個公園,名為『紫宸園』(在中國大陸從來沒有一位華人牧師,華人神學家是受到國家市政府這樣的尊崇和優待的,趙紫宸是唯一的一位).

趙紫宸在新中國建立之時,曾主動地站在前線,呼籲教會必須進行改革. 1950年初,中華全國基督教協進會在北京成立一個專責籌備委員會時,他獲選擔任改革組組長,負責草擬中國教會改革的新方案. 在短短的半年時間內,他撰寫了三篇文章,分別是: 《中國基督教教會改革的途徑》 (3月10日完成), 《用愛心建立團契》 (4月5日完成), 及《審判下教會的革新》 (7月26日完成). 這三篇文章,也不僅是表達了趙紫宸對中國教會改革的看法,更也是表現出他在回應大時代轉變時所需要的一種『基督教中國化』的創新神學思維.尤其是第二篇的題目是:《以愛心建立團契》,趙紫宸認為教會就是一個用愛心建立的團契；我們要建立教會,就必須是用愛心建立教會的團契. 事實上,我們不論是在教會,學校,或是社會服務機構工作,我們所有的基督教精神就是愛的精神,表現基督的大愛就是我們的見證. 這也該是我們崇基學院必須見證和傳承下去的基督教的大愛和服務的精神吧!

還記得陶行知怎樣說嗎?陶行知是身體力行地告訴我們: 最容易為中國人所接受的基督教教育就是『愛的教育』. 陶行知創辦『曉莊學校』的時候,他說:『曉莊是從愛裡產生出來的.沒有愛便沒有曉莊.[7]  因為聖經的教導就是:『愛是從上帝而來的! 凡有愛心的,都是由上帝而生,並且認識上帝. 沒有愛心的,就不認識上帝; 因為上帝就是愛』.[8]

祈求上帝祝福我們每一位,無論我們是在任何一個角落,或是任何的工作崗位中,我們都能夠認識清楚,基督教就是一個愛的宗教,在教會內,我們要以愛心建立團契；在教育的場所中,我們要實踐基督教教育,我們所從事的教育工作就是一種『以愛心轉化人的心靈』的工作,願上主祝福使用大家!  誠心所願!

[1]參凌道揚:《我們 - 崇基的使命》.1955年2月4日就職典禮上致詞.《崇基校刊》第7期,1955年10月出版,頁3-4.

[2]可參閱吳梓明:《30年來中國教會大學史研究的反思》.香港: 香港中文大學崇基學院宗教與中國社會研究中心,2017年出版,頁1.

[3]可參閱何榮漢:《陶行知- 一位基督徒教育家的再發現》.香港: 香港中文大學研究院宗教與神學學部,宗教研究哲學博士論文,2002年7月.

[4]參陶:<曉莊三歲敬告同志書>,刊於《鄉村教師》,第7期,1930年3月15日.轉載於文明國編:《陶行知自述》,合肥:安徽文藝出版社,2013年,頁13.

[5]《新約聖經》《約翰壹書》4:7-8 .

[6]參 Aisi Li, “Afterword”, in Eunice V. Johnson, Timothy Richard’s Vision: Education and Reform in China, 1880-1910. (edited by Carol Lee Hamrin), Eugene, Oregon: Pickwick Publishing Co., 2014, p. 131.

[7]參陶:<曉莊三歲敬告同志書>,刊於《鄉村教師》,第7期,1930年3月15日.轉載於文明國編:《陶行知自述》,合肥:安徽文藝出版社,2013年,頁13.

[8]《新約聖經》《約翰壹書》4:7-8 .



\newpage

\section{無力的時代有力的信仰}
\label{sec:434}
\textbf{袁天佑牧師 - 無力的時代,有力的信仰}
\newline
\newline
\hyperref[sec:433]{< < < PREV SERMON < < <}
~
\hyperref[sec:toc]{[返目錄]}
~
\hyperref[sec:435]{> > > NEXT SERMON > > >}


耶利米書 33:14-16
\newline
\begin{longtable}{cl}
\hline
\hline
章節 & 經文 (和合本修訂版)\\
\hline
33:14 & \begin{tabularx}{0.7\textwidth}{X} 「看哪，日子將到，我應許以色列家和猶大家的恩言必然實現。這是耶和華說的。 \end{tabularx} \\ \\ \relax
33:15 & \begin{tabularx}{0.7\textwidth}{X} 在那些日子、那時候，我必使大衛公義的苗裔長起來；他必在地上施行公平和公義。 \end{tabularx} \\ \\ \relax
33:16 & \begin{tabularx}{0.7\textwidth}{X} 在那些日子，猶大必得救，耶路撒冷必安然居住，他的名必稱為『耶和華－我們的義』。 \end{tabularx} \\ \\
[1ex]
\hline
\hline
\end{longtable}
帖撒羅尼迦前書 3:9-13
\newline
\begin{longtable}{cl}
\hline
\hline
章節 & 經文 (和合本修訂版)\\
\hline
3:9 & \begin{tabularx}{0.7\textwidth}{X} 我們在神面前，因著你們滿有喜樂。為這一切喜樂，我們能用怎樣的感謝為你們報答神呢？ \end{tabularx} \\ \\ \relax
3:10 & \begin{tabularx}{0.7\textwidth}{X} 我們晝夜切切祈求要見你們的面，來補足你們信心的不足。 \end{tabularx} \\ \\ \relax
3:11 & \begin{tabularx}{0.7\textwidth}{X} 願我們的父神自己和我們的主耶穌，為我們開路到你們那裡去。 \end{tabularx} \\ \\ \relax
3:12 & \begin{tabularx}{0.7\textwidth}{X} 又願主使你們彼此相愛的心，和愛眾人的心，都能增長，充足，如同我們愛你們一樣， \end{tabularx} \\ \\ \relax
3:13 & \begin{tabularx}{0.7\textwidth}{X} 好堅固你們的心，使你們在我們的主耶穌同他眾聖徒來臨的時候，在我們父神面前，成為聖潔，無可指責。阿們！ \end{tabularx} \\ \\
[1ex]
\hline
\hline
\end{longtable}
路加福音 21:25-36
\newline
\begin{longtable}{cl}
\hline
\hline
章節 & 經文 (和合本修訂版)\\
\hline
21:25 & \begin{tabularx}{0.7\textwidth}{X} 「日月星辰要顯出預兆，地上的邦國也有困苦，因海中波浪的響聲而惶惶不安。 \end{tabularx} \\ \\ \relax
21:26 & \begin{tabularx}{0.7\textwidth}{X} 人想到那要臨到世界的事，就都嚇得魂不附體，因為天上的萬象都要震動。 \end{tabularx} \\ \\ \relax
21:27 & \begin{tabularx}{0.7\textwidth}{X} 那時，他們要看見人子帶著能力和大榮耀駕雲來臨。 \end{tabularx} \\ \\ \relax
21:28 & \begin{tabularx}{0.7\textwidth}{X} 一有這些事，你們就當挺身昂首，因為你們得救贖的日子近了。」 \end{tabularx} \\ \\ \relax
21:29 & \begin{tabularx}{0.7\textwidth}{X} 耶穌對他們講了一個比喻說：「你們看無花果樹和各樣的樹， \end{tabularx} \\ \\ \relax
21:30 & \begin{tabularx}{0.7\textwidth}{X} 樹葉一長出來，你們看了自然就知道夏天近了。 \end{tabularx} \\ \\ \relax
21:31 & \begin{tabularx}{0.7\textwidth}{X} 同樣，當你們看見這些事發生，就知道神的國近了。 \end{tabularx} \\ \\ \relax
21:32 & \begin{tabularx}{0.7\textwidth}{X} 我實在告訴你們，這世代還沒有過去，一切都要發生。 \end{tabularx} \\ \\ \relax
21:33 & \begin{tabularx}{0.7\textwidth}{X} 天地要廢去，我的話卻絕不廢去。」 \end{tabularx} \\ \\ \relax
21:34 & \begin{tabularx}{0.7\textwidth}{X} 「你們要謹慎，免得被貪食、醉酒和今生的憂慮壓住你們的心，那日子就忽然臨到你們， \end{tabularx} \\ \\ \relax
21:35 & \begin{tabularx}{0.7\textwidth}{X} 如同羅網一樣，因為那日子要臨到所有居住在地面上的人。 \end{tabularx} \\ \\ \relax
21:36 & \begin{tabularx}{0.7\textwidth}{X} 你們要時時警醒，常常祈求，使你們能逃避這一切要來的事，得以站立在人子面前。」 \end{tabularx} \\ \\
[1ex]
\hline
\hline
\end{longtable}
講題:無力的時代,有力的信仰　Time of Powerless, but Faith is Our Power

經課:耶利米書33:14-16;帖撒羅尼迦前書3:9-13;路加福音21:25-36

講員:袁天佑牧師

場合:香港中文大學崇基學院禮拜堂

日期:2018年12月2日

今天的講題,是採用我剛剛出版的一本新書的書名.其實我沒有任何意願想寫書,只是在退休後,在臉書上分享一些講章.有出版社覺得很有意思,便將之輯錄成書出版.去年出版了《走進時代的信仰》,今年出版了《無力的時代,有力的信仰》這書.書名不是我自己想出來的,這是出版社社長的心思.我想,或許不是書中的講章內容有甚麼特別,只是書名改得很有意思,所以銷售情況都很理想.吸引讀者相信特別是「無力的時代」這標題.或許這真是不少香港人的感受.政治的無力感,極權的臨近,在前面的好像只有黑暗.有幾次有不同的教會邀請我講及在這政局下有甚麼出路,有甚麼曙光,我都覺得有點困難去講.時代真的無力?

狄更斯(Charles Dickens)在他的名著《雙城記》中說了一句經典名句:「那是最好的時代,也是最壞的時代......」.甚麼時候才是最好的時代?甚麼時候是壞的時代?甚麼時候是無力的時代?甚麼時候是有力的時代?甚麼時候是黑暗的時代?甚麼時候是光明的時代?我在書的序言中指出,好抑或壞,黑暗抑或光明,有力或無力,基本上不是在甚麼時間,或外在的環境,而是在於自己.在好的,在有力的,或是光明的時刻中,其實正是最壞,最無力和最黑暗的時代,因為在這些時刻中,人反為放鬆自己,沉醉於安逸之中,好的,有力的或是光明的日子,反變成壞,無力和黑暗.在壞的,在無力和黑暗之中,覺得是壞,無力和黑暗,不思進取,當然會變得更壞,更無力和黑暗.但如果在壞,無力和黑暗的時刻,反思過去的失敗,有進取的心,這反為會是好的,有力的和光明的時刻.

上星期,九龍西進行立法會補選.大家都知道泛民敗選.這未必不是一件好事.失敗教我們更多反思原因,成功反會令我們驕傲而不自覺.

中文大學社會學系的陳健民教授在他「最後的一堂課」中便指出,「在黑暗的時候,才會看見星星.」不過,重要點是:我們願否成為那黑暗中的星星呢?

今主日的經課,都是同樣的指向這方面,在壞,在無力和黑暗的日子中,要存著盼望和信心.

耶利米書三十三章14~16節

耶利米先知,生活在主前第七世紀末至第六世紀左右,這也是猶大國末期的時代.猶大國先後受鄰近多個國家所控制,包括亞述,埃及和巴比倫.國王的無能,宗教信仰的敗壞,這是猶大國最黑暗的時期.不單只猶大國處於黑暗之中,耶利米的一生,都是黑暗,被稱為「流淚的先知」.他向猶大的君王和百姓宣佈猶大國將亡,上帝的懲罰會臨到他們當中.當然這些說話不能被君王和百姓所接納,耶利米亦因此經歷人們的反對,鞭打,監禁等逼迫(參看耶十一18~21).

當你處在這困境時,你會如何?會感到黑暗,灰心嗎?

但耶利米並沒有因此而感到黑暗或灰心.耶利米書二十九章記載了耶利米寫信給被擄至巴比倫的百姓,「你們要建造房屋,住在其中；要開墾田園,吃園中所出產的；要娶妻生兒養女,為你們的兒子娶妻,使你們的女兒嫁人,生兒養火,你們要在那裡生養眾多,不可減少.我使們們被擄到的那城,你們要那城求平安,為那城向耶和華祈求,因為那城得平安,你們也隨著得平安.」(5~7節)耶利米這番話的意思,即是叫他們「如常生活」.而耶利米更指出,這是耶和華要對百姓所說的.

耶利米書三十二記載耶利米被囚於猶大的王宮之中,但他仍用錢買地,預備建造房屋,田地和種葡萄.在耶利米書三十三章14~18節,耶利米更宣告:「看哪,日子將到,我應許以色列家和猶大家的恩言必然實現.這是耶和華說的.在那日子,那時候,我必使大衛公義的苗裔長起來；他必在地上施行公平和公義.在那些日子,猶大必得救,耶路撒冷必安然居住,他的名必稱為『耶和華--我們的義』.」

在黑暗和看來沒有盼望的日子,耶利米仍存希望,仍看見光明.

帖撒羅尼迦前書三章9~13節

帖撒羅尼迦教會是保羅在第二次傳道旅程所建立的教會.按照《使徒行傳》的記載,保羅在離開腓立比後,便前往帖撒羅尼迦,在那裡連續三個安息日在會堂裡講解有關耶穌復活的事.有人相信耶穌,包括希臘人和一些尊貴的婦女,但猶太人因嫉妒,聚集一些市并流氓,煽動群眾,並指保羅等人是違背凱撒的,把他們趕離帖撒羅尼迦(徒十七1~9).

後來保羅來到庇哩亞,傳道亦有效果,但帖撒羅迦城的猶太人繼續往這地,煽動群眾.保羅只有離開,前往雅典(徒十七10~15).保羅雖然身在雅典,但心仍掛念著帖撒羅尼亞教會,派提摩太去探望.提摩太後來回報,帶來好的消息,知道教會在患難和逼迫中,不單領受真道,還把真道傳往別處(帖前一6~8).不過,保羅亦從提摩太知悉,教會中有人誤解真理,認為已死的人,不能得著主再來的祝福(帖前四13~14),於是在哥林多寫信給帖撒羅尼迦教會.這也是保羅最早的一封書信.

明顯的,不論是保羅,或是帖撒羅尼迦教會,他們都是處於困難之中,但保羅並沒有失望,他仍存著盼望,希望能再見到教會眾人,在困難和逼迫中,期望教會信眾能在信心,愛心和盼望上,能有所長進.保羅說:「我們晝夜切切祈求要見你們的面,來補足你們信心的不足.願我們的父上帝自己和我們的主耶穌,為我們開路到你們那裡去.又願主使你們彼此相愛的心,和愛眾人的人,都能增長,充足,如同我們愛你們一樣,好堅固你們的心,使你們在我們的主耶穌同他眾聖徒來臨的時候,在我們父上帝面前成為聖潔,無可指責.」(四10~13)

路加福音二十一章25~36節

路加福音二十一章,大部份都是記載耶穌論及耶路撒冷和聖殿被毀,甚至有人認為是論及將來末日來臨的預言,於是去揣測末日來臨的日子.但在馬太和馬可的記述中,都有提及耶穌這樣說:「那日子,那時辰,沒有人知道,連天上的天使也不知道,子也不知道,惟有父知道.」(太二十四36；可十三32)既然是無法知曉,揣測是沒有意思的.耶穌要指出的,我們已生活在末世的日子之中.不單是聖殿和耶路撒冷被毀,我們也是生活在戰爭,大自然的災害,宗教信仰受到壓制的日子當中,昔日如是,今天也是這樣.

看來這是壞的日子,是黑暗的日子,耶穌並沒有教導我們去揣測末日的日子,反而有兩個提醒.

第一,要相信那日子必定來臨:「那時,他們要看見人子帶著能力和大榮耀來臨.一有這些事,你們就當挺身昂首,因為你們得救贖的日子近了.(二十一27~28)

耶穌預言聖殿將會拆毀,但他也曾說過:「你們拆毀這殿,我三日內要把它重建.」(約二19)耶穌所指的是他死後三天復活的事,這事已成為真實,所以主再來的日子也必定來臨.我們要有這信心.

第二,耶穌提醒門徒,生活在當下,要警醒.他說:「你們要謹慎,免得被貪食,醉酒和今生的憂慮壓住你們的心,那日子就忽然臨到你們,如同羅網一樣,因為那日子要臨到所有居住在地面上的人.你們要時時警醒,常常祈求,使你們能逃避這一切要來的事,得以站立在人子面前.」(路二十一34~36)

要警醒謹守的意思,明顯就是在這末世的日子中,能持守真理.

我開始講道時指這兩年,出版社為我編輯了兩本書.第一本書出版時,社長為書改了一個書名,《走在時代前的信仰》.但我覺得我所講的道理,沒有甚麼新意,沒有帶人走向前,只是在這時代中,堅守真道,所以我改了書名,《走進時代的信仰》.信仰是不變的,但怎樣在這時代中去堅守而已.

第二本書,《無力的時代,有力的信仰》.我聽聞有人希望能拿來在團契聚會中討論,但後來他這樣說:「書的內容說得很對,但好像沒有新的題材,沒有對現在政局說了些甚麼.」因為只是聽聞,所以我也沒法回應,但心中想,我們的信仰有甚麼新的題材呢?

雨傘運動後,指我是「雨傘牧師」,亦以為我很關心政治.但其實我不大懂政治,所以沒法為香港政治前途作出任何建議.我退休後,仍是牧師,我所關心的,就是信徒,無論時代變得如何,是好是壞,是有力或無力,是光明抑黑暗,信徒皆能持守信仰,學習怎樣在這時代中實踐信仰.我也盼望香港市民,雖然看見極權的來臨,可以有點兒灰心失望,但不要絕望,仍要持守初衷.不少人在雨傘後說:「要在社區中深耕細作.」這是甚麼意思?其實就是在社區中實踐「行公義,好憐憫」.基督徒更會加上「與你的上帝同行」(彌六8)我相信這是在現今的時代,其實也是在任何時代中,要做的事.

今主日是「將臨期第一主日」.大家都很清楚:在將臨期,我們要預備紀念和慶祝主耶穌第一次的已臨,另一方面,也要預備迎接主耶穌第二次再臨.不過,很多時候,教會在這節期間,多是預備和安排慶祝聖誕的活動,較少去呼籲信眾,也安排活動,讓信眾預備身心,去迎接主耶穌第二次的來臨.

將臨期,是一個處於「既濟和未濟」(already but not yet)的張力的節期.一方面,我們仍生活在耶穌已第一次降臨的日子中,他來帶來光明,但另一方面,我們仍生活在耶穌第二次來臨的盼望中,在這段時間,人仍生活在罪惡和黑暗籠罩的日子中.在這末世的日子,其實是好與壞,有力與無力,黑暗與光明同時存在的日子.在這充滿著張力的日子中,人最重要不要放棄好(美善)的事,光明的事,人仍要作個有力的人.而基督徒更靠著我們的信仰,在這世代中作有力的見證.

陳健民在他的「最後的一堂課」中指出:「我有信仰,但沒有宗教.」這句話令我想起一位猶太藉無神論的哲學家布洛霍(Ernst Bloch)所說的一句話:「凡是有盼望的地方,就有宗教；但有宗教的地方,郤不一定有盼望.」(Where there is hope, there is religion.  Where there is religion, there is not always hope.)

將臨期是提醒我們,在壞,黑暗和無力的時代中,彰顯盼望.這才是宗教.我們要在黑暗中燃點光明,成為光照耀世人.我們的信念是:「光照在黑暗裡,黑暗郤沒有勝過光.」(約一5)



\newpage

\section{傳而有信}
\label{sec:435}
\textbf{林豪恩先生 - 傳而有信}
\newline
\newline
\hyperref[sec:434]{< < < PREV SERMON < < <}
~
\hyperref[sec:toc]{[返目錄]}
~
\hyperref[sec:436]{> > > NEXT SERMON > > >}


講題:傳而有信 　The Messenger and the Message

講員:林豪恩先生

經文:瑪拉基書3

場合:香港中文大學崇基學院禮拜堂

日期:2018年12月9日

引言

三代經課,四年重覆一次.今天的經課,似曾相識.上網搜查,發現原來四年前,即是2015年的將臨期第二主日,剛巧是我講道.網上資料繼續顯示,再四年前,即是2011年的將臨期第二主日,又是我講道.這項發現,帶來了遠慮和近憂,遠慮是:原來在這講台上所講的每一句說話,都被公諸於世,並記錄在案,不知道會被如何使用,也不知道會帶來甚麼後果?近憂是:經課都已經講過兩次了,還有甚麼好講呢?在苦惱之中,突然想起一位前輩講過的話,大意是說,作者推出一本新書,通常只有百分之二十的內容是新的,百分之八十都是舊的.如果有更多新內容,就會留待出版第二本書的時候使用.現在,給大家一項挑戰,以下所講的內容,哪些是重述兩篇舊講章的百分之八十?哪些是新加上的百份之二十?

言歸正傳,其實,對於一位傳道者來說,沒有道可傳,是相當焦慮的.因為我們知道,對於聽道者來說,沒有道可聽,也是十分沮喪的.楊牧谷牧師在二十世紀結束前,發表「信息衰竭乃謝幕的徵兆」這短文,在教會界一石激起千重浪.文章有幾句如此說:「一個機構,一本雜誌,一個傳道人,它/他就是一個攜帶及傳遞信息的媒體,假如日久曠廢,信息流失,它/他就再沒有存在的需要,要退下來便成了自然又合理的結局.給它/他換件新裝是無濟於事的,就如說:為機構蓋一座大樓,為傳道人換上一套Armani西服,為雜誌由印刷改用電子網絡出版,可以延續它的生命嗎?……基本上,牧師是沒有生存保障的,他唯一的保障就是上帝委任他傳的道,失道?他就失了存在的價值,任何解釋都是多餘.」難怪有牧師閱讀這篇短文後,用「聞者驚心」來形容他的反應.

將臨期第二主日經課介紹的,正正是傳信息的人.這些傳信息的人,舊約經課瑪拉基書稱為「使者」(messenger),路加福音中撒迦利亞的詩篇中稱為「先知」(prophet),福音書經課路福音稱為「曠野的呼聲」(the voice in the wilderness),書信經課腓立比書中稱為耶穌基督的僕人(servants of Christ Jesus),他們的共同地方,用楊牧谷牧師的說法,都是「攜帶及傳遞信息的媒體」.將臨期到了,是「報佳音」的季節了,有甚麼佳音向此時此地的人們宣呢?將臨期的經課所報的又是甚麼佳音呢?讓我們重遇這些傳信息者及他們所傳的信息,從中學習及參考.

二．內容

1．瑪拉基書:瑪拉基的預告

3:1萬軍之耶和華說:「看哪,我要差遣我的使者在我前面預備道路.你們所尋求的主必忽然來到他的殿；立約的使者,就是你們所仰慕的,看哪,快要來到.」

3:2他來的日子,誰能當得起呢?他顯現的時候,誰能立得住呢?因為他如煉金匠的火,如漂洗者的鹼.

3:3他必坐下如煉淨銀子的人,必潔淨利未人,熬煉他們像金銀一樣；他們就憑公義獻供物給耶和華.

3:4那時,猶大和耶路撒冷所獻的供物必蒙耶和華悅納,彷彿古時之日,上古之年.

這段經文中有三位使者.第一節上半部說:「我要差遣我的使者在我前面預備道路」,這句說話在路加福音7:27節被引述來指出,這位使者就是施洗約翰；第一節下半部說:「立約的使者,就是你們所仰慕的,看哪,快要來到.」這第二位使者,相信就是他們盼望的彌賽亞.那麼,第三位使者在那裡呢?就是正在傳達這信息的先知瑪拉基.瑪拉基這名字的意思就是「神的使者」.瑪拉基這位使者在這裡向當時的以色列人民介紹另外兩位使者,有甚麼意義呢?或者,我們需要稍稍閱讀上文.

2:17你們用言語煩瑣耶和華,你們還說:「我們在何事上煩瑣他呢?」因為你們說:「凡行惡的,耶和華眼看為善,並且他喜悅他們」；或說:「公義的上帝在哪裡呢?」

「凡行惡的,耶和華眼看為善,並且他喜悅他們」,「公義的上帝在哪裡呢?」這些聲音,可能表達了人們對當時局勢的埋怨.這是一個怎樣的時代呢?這是一個怎樣的社會呢?中國人有話說:「殺人放火金腰帶,修橋補路無屍骸 」.善不見得有善報,惡不見得有惡報.當時的以色列人民質問,如果有公義的上帝,祂在哪裡呢?他們甚至已經放棄了上帝是公義的,指責上帝正邪不分,視惡為善,賞惡罰善.古代以色列作為一個宗教社會,不可能宣布「上帝已死」,但如此的上帝,怪不得在瑪拉基書中描述了那麼多以色列人,包括祭司「藐視」耶和華上帝的狀況.他們為甚麼如此無奈呢?他們為甚麼如此憤怒呢?他們經歷了甚麼?他們看見了甚麼呢?

3:14你們說:『事奉上帝是徒然的,遵守上帝所吩咐的,在萬軍之耶和華面前苦苦齋戒,有甚麼益處呢?

3:15如今我們稱狂傲的人為有福,並且行惡的人得建立；他們雖然試探上帝,卻得脫離災難.』」

此時此地,有聲音說:關心社會還有用嗎?社會還肯讓我參與嗎?「和平理性非暴力」的方法用過,有用嗎?「勇武抗爭」的方法用過,有用嗎?票投過了,有用嗎?選舉舉行過了,有用嗎?議會繼續開會,有用嗎?有人說,君不見某陣營派甚麼人出來參選都能夠勝出?又有人說,君不見某陣營一盆散沙,互相攻擊,成事不足敗事有餘嗎?你還能期望甚麼呢?你還堅持甚麼呢?誰是同路人,誰是對手都不容易分得清,陰謀論鋪天蓋地,你還能信任誰呢?你還能相信甚麼呢?

如果我們深切感受自身的處境,或者我們能夠明白多一些當時的以色列人,特別是在瑪拉基書被指責的祭司群體,為何如此輕視上帝.或者,這是殘酷的現實帶來的無助,正義不彰做成的「心淡」.

面對如此的社會狀況,瑪拉基這位傳道者還有甚麼道可傳呢?他有靈丹妙藥嗎?沒有.他能撥亂反正嗎?不能.但他呼籲身處混亂狀況的祭司和人民不要放棄,因為上帝並沒有放棄.如何知道上帝沒有放棄呢?這就是瑪拉基要介紹兩位將要來的使者的原因.上帝已經差遣了兩位使者,一位來預備道路,另一位「你們所仰慕的」,「快要來到」.他「如煉金匠的火,如漂洗者的鹼.他必坐下如煉淨銀子的人,必潔淨利未人,熬煉他們像金銀一樣；他們就憑公義獻供物給耶和華」.真金不怕紅爐火,公義仍是我們所信仰的耶和華上帝所肯定和喜悅的.

今年的將臨期,有甚麼佳音可以在此時此地傳播呢?願瑪拉基的呼籲激勵我們,不要被眼前的亂象消滅信念,因為上帝繼續差派使者,並且快將來到.在亂世之中,信念仍有用嗎?或者,以下這個故事能夠給我們一點啟發.

James Stockdale是越戰時期被俘虜的美國海軍上將,他曾反復承受極端殘酷的拷問長達七年之久,卻沒有對俘虜他的人絲毫讓步.James Stockdale被關押8年後放了出來,有人問他:「你為什麼能熬過這艱難的8年?」他說:「因為我有一個信念,相信自己一定能出來,能夠再見到我的妻子和孩子,這個信念一直支撐著我,使我生存了下來.」那人又問:「那你的同伴中最快死去的又是哪些人呢?」他回答說:「是那些太樂觀的人.」這不是很矛盾嗎?為什麼那些樂觀的人會死得很快呢?他說:「他們總想著耶誕節可以被放出去了吧?耶誕節沒被放出去；就想復活節可以被放出去,復活節沒被放出去；就想著感恩節,而後又是耶誕節,結果一個失望接著一個失望,他們逐漸喪失了信心,再加上生存環境的惡劣,於是,他們鬱鬱而終.」他總結說:「對長遠我有一個很強的信念,相信自己一定能夠活著出去,一定能再見到我的妻子和小孩；但是我又正視現實的殘酷.」不放棄終極得勝的信念,同時正視不理想的現實,或者正是此時此地需要的佳音.

2．路加福音:撒迦利亞的詩歌

瑪拉基先知的信息,在歷史長河中的某一刻,得到回響.有另一位傳信息者-祭司撒迦利亞,以下列這首詩歌回應:

1:68「主-以色列的上帝是應當稱頌的!因他眷顧他的百姓,為他們施行救贖,

1:69在他僕人大衛家中,為我們興起了拯救的角,

1:70正如主藉著古時候聖先知的口所說的,

1:71『他拯救我們脫離仇敵,脫離一切恨我們之人的手.

1:72他向我們列祖施憐憫,記得他的聖約,

1:73就是他對我們祖宗亞伯拉罕所起的誓,

1:74-75叫我們既從仇敵手中被救出來,就可以終身在他面前,無所懼怕地用聖潔和公義事奉他.

懷著堅強信念的人,知道上帝沒有視而不見,也沒有放棄.在被擄回歸,面對內憂外患的混亂局勢中,有瑪拉基這位神的使者呼籲人民不要被眼前的亂象消滅信念,不要因為感到世道險惡而同流合污；在身處被羅馬帝國殖民統治的壓抑時期,仍然有祭司撒迦利亞唱出安慰與盼望的詩歌.現在,我們受到仇敵壓迫,但終有一天,我們會得到釋放,在一個自由及不用懼怕的狀況下,活出聖潔,活出公義.

1:76孩子啊,你要稱為至高者的先知；因為你要走在主的前面,為他預備道路,

1:77叫他的百姓因罪得赦,認識救恩；

1:78因我們上帝憐憫的心腸,叫清晨的日光從高天臨到我們,

1:79要照亮坐在黑暗中死蔭裡的人,把我們的腳引到和平的路上.』」

上次我用「父親的詩句」來講撒迦利亞這些漂亮的詩句,感覺溫馨浪漫.今天,我們把這首詩放在當時的歷史處境中,感覺完全不一樣.如果仍然要用「美」這個字來形容這幾句詩,可能是「淒美」.讓我們想一想,差不多所有為信念而投身公共事務,推動社會正義的前線朋友,最擔心,也最內疚的是連累家人.如果要在不同政見,不同陣營的朋友當中找一個共識,「擔心連累家人」可能是最有共嗚的.祭司撒迦利亞,竟然「擺自己的兒子上檯」,為那位「將要來的」預備道路,為社會延續信念和盼望,參與為「坐在黑暗中死蔭裡的人」帶來光明,為面對「仇敵」與處身「仇恨」的人「的腳引到和平的路上」.如果瑪拉基是那位被時代選上的先知,撒迦利亞是那位被時代選上的祭司,約翰就是那位被時代選上的孩子.其實,他們都是被上帝選上的使者,把信念及盼望一代一代地延續下去.

3．路加福音:曠野的聲音

信念和盼望,是這麼近,又是那麼遠.瑪拉基的名字已經載入史冊,撒迦利亞也從舞台上退隱,世界似乎沒有甚麼大改變,架構仍然是如此穩固.讓我們看看路加如何介紹當時的局勢:

3:1凱撒提庇留在位第十五年,本丟‧彼拉多作猶太總督,希律作加利利分封的王,他兄弟腓力作以土利亞和特拉可尼地區分封的王,呂撒聶作亞比利尼分封的王,

3:2亞那和該亞法作大祭司.

南就猶大,中部撒馬利亞,北部加利利,全部都屬於羅馬帝國的版圖.羅馬皇帝透過政治分贓及直接派遣,有代理人嚴密地看守著那片土地及其上的人民.在被看管得密不透風的環境中,卻有一位孩子漸漸長大.看來,瑪拉基的信息快將得到印證,撒迦利亞的盼望漸近實現.

那時,撒迦利亞的兒子約翰在曠野裡,上帝的話臨到他.

3:3他就走遍約旦河一帶地方,宣講悔改的洗禮,使罪得赦.

3:4正如以賽亞先知書上所記的話:「在曠野有聲音呼喊著:預備主的道,修直他的路!

3:5一切山窪都要填滿；大小山岡都要削平!彎彎曲曲的地方要改為筆直；高高低低的道路要改為平坦!

3:6凡有血肉之軀的,都要看見上帝的救恩!」

朝代更替,然而歷史並沒有終斷；帝國頑強,但是盼望沒有消失.瑪拉基在生時沒有見到自己預告的那位「預備道路」的使者出現,以賽亞在生時也尚未聽到自己筆下的那「曠野的聲音」,然而,在某年某月某日,那位「預備道路」的使者果真出現了,那「曠野的聲音」響徹雲霄.

縱使世界未如理想,上帝沒有放棄,差遣了要來的,終必會來.當下,我們應如何生活呢?使徒保羅對腓立比信徒的忠告值得我們參考:

4．腓立比書:使徒的忠告

1:1基督耶穌的僕人保羅和提摩太寫信給住腓立比,在基督耶穌裡的眾聖徒,以及諸位監督和執事.

1:2願恩惠,平安從我們的父上帝和主耶穌基督歸給你們!

1:3我每逢想念你們,就感謝我的上帝,

1:4每逢為你們眾人祈求的時候,總是歡歡喜喜地祈求,

1:5因為從第一天直到如今,你們都同心合意興旺福音.

1:6我深信,那在你們心裡動了美好工作的,到了耶穌基督的日子必完成這工作.

1:7我為你們眾人有這樣的想法原是應當的,因為你們常在我心裡；無論我是在捆鎖中,在辯明並證實福音的時候,你們都與我一同蒙恩.

1:8我以基督耶穌的心腸切切想念你們眾人,這是上帝可以為我作證的.

1:9我所禱告的就是:要你們的愛心,在知識和各樣見識上,不斷增長,

1:10使你們能分辨是非,在基督的日子作真誠無可指責的人,

1:11更靠著耶穌基督結滿仁義的果子,歸榮耀稱讚給上帝.

如果世界上充滿是非不分的聲音,使徒保羅的囑咐信徒仍要「分辨是非」；如果社會中滿佈不是你死就是我亡的競爭,使徒保羅仍然為信徒祈禱「要你們的愛心,在知識和各樣見識上,不斷增長」；如果處身於「劣幣驅逐良幣」的環境,使徒保羅仍然鼓勵信徒「結滿仁義的果子」,「作真誠無可指責的人」.有用嗎?有的,因為使徒保羅深信「你們心裡動了美好工作」,「到了耶穌基督的日子必完成這工作」(V6).

三．總結

「又到聖誕,又到聖誕」,我們有甚麼佳音向社會分享呢?1989年黃霑先生等音樂人「搞笑過癮玩改詞」產生了一系列聖誕改編歌,改篇自Santa Claus Is Coming To Town的「又到聖誕」當中的末世情懷,或者與當下不遑多讓.其中幾句歌詞如此說:

又到聖誕　又到聖誕

下個聖誕有沒今年甘歎

鄧小平is coming to town

趁又到聖誕　人就要歎喇

第處聖誕有沒香港甘歎

阿鉅子誠都怕且就yang

八番次共你聚聚　唱歌搭住晚飯

九七先至念大獲飯　系要今年玩多晚

將臨期其實也充滿末世情懷,正如使徒保羅所指出,有「耶穌基督的日子」(1:6),「基督的日子」(1:10).不過,將臨期的末世情懷是充滿盼望的,因為「到了耶穌基督的日子必完成這工作」；將臨期的信息是積極的,因為當我們「靠著耶穌基督結滿仁義的果子」,就有「歸榮耀稱讚給上帝」的回報.如果我們就是傳信息的人,那麼我們與我們所傳的信息是分不開的.願上主幫助我們,也願我們互相勉勵.



\newpage

\section{我們多麼需要星星}
\label{sec:436}
\textbf{鄧瑞強博士 - 我們多麼需要星星}
\newline
\newline
\hyperref[sec:435]{< < < PREV SERMON < < <}
~
\hyperref[sec:toc]{[返目錄]}
~
\hyperref[sec:437]{> > > NEXT SERMON > > >}


馬太福音 2:2-12
\newline
\begin{longtable}{cl}
\hline
\hline
章節 & 經文 (和合本修訂版)\\
\hline
2:2 & \begin{tabularx}{0.7\textwidth}{X} 「那生下來作猶太人之王的在哪裡？我們在東方看見他的星，特來拜他。」 \end{tabularx} \\ \\ \relax
2:3 & \begin{tabularx}{0.7\textwidth}{X} 希律王聽見了，就心裡不安；耶路撒冷全城的人也都不安。 \end{tabularx} \\ \\ \relax
2:4 & \begin{tabularx}{0.7\textwidth}{X} 他就召集了祭司長和民間的文士，問他們：「基督該生在哪裡？」 \end{tabularx} \\ \\ \relax
2:5 & \begin{tabularx}{0.7\textwidth}{X} 他們說：「在猶太的伯利恆。因為有先知記著： \end{tabularx} \\ \\ \relax
2:6 & \begin{tabularx}{0.7\textwidth}{X} 『猶大地的伯利恆啊，你在猶大諸城中並不是最小的；因為將來有一位統治者要從你那裡出來，牧養我以色列民。』」 \end{tabularx} \\ \\ \relax
2:7 & \begin{tabularx}{0.7\textwidth}{X} 於是，希律暗地裡召了博學之士來，查問那星是甚麼時候出現的， \end{tabularx} \\ \\ \relax
2:8 & \begin{tabularx}{0.7\textwidth}{X} 就派他們往伯利恆去，說：「你們去仔細尋訪那小孩子，找到了就來報信，我也好去拜他。」 \end{tabularx} \\ \\ \relax
2:9 & \begin{tabularx}{0.7\textwidth}{X} 他們聽了王的話就去了。忽然，在東方所看到的那顆星在前面引領他們，一直行到小孩子所在地方的上方就停住了。 \end{tabularx} \\ \\ \relax
2:10 & \begin{tabularx}{0.7\textwidth}{X} 他們看見那星，就非常歡喜； \end{tabularx} \\ \\ \relax
2:11 & \begin{tabularx}{0.7\textwidth}{X} 進了房子，看見小孩子和他母親馬利亞，就俯伏拜那小孩子，揭開寶盒，拿出黃金、乳香、沒藥，作為禮物獻給他。 \end{tabularx} \\ \\ \relax
2:12 & \begin{tabularx}{0.7\textwidth}{X} 因為在夢中得到主的指示，不要回去見希律，他們就從別的路回自己的家鄉去了。 \end{tabularx} \\ \\
[1ex]
\hline
\hline
\end{longtable}
講題:我們多麼需要星星       Oh, the Star is There

經文:馬太福音2:2-12

講員:鄧瑞強博士

場合:香港中文大學崇基學院禮拜堂

日期:2018年12月24日

在這平安夜,願各位朋友,各位弟兄姊妹平安.

若果前幾年,香港打仗,今天下午才簽和約,那麼,今晚我們在這裡過平安夜,互祝平安,就會別有一番滋味.戰爭越激烈,平安越寶貴.若果我們體會過人間的爭鬥,親嚐過生命的幽暗,經驗過靈魂的絕望,我們就會珍惜神所賜的平安.

耶穌出生的那個晚上,是漆黑的歷史長夜的一個盡頭.當舊約最後一位先知死去之後,差不多有四百年,才等到耶穌的出生.在這漫長的四百年間,人間再沒有先知,再聽不到神的聲音,再看不到神指示的前路.想想,四百年,人聽不到神的聲音,多麼可怕.這四百年,是靈魂的黑夜.這四百年,像是無邊無際的黑洞,連點出來的火,都是黑暗的,不發光的,死寂的.好漫長的黑暗,好廣大的黑暗,好深沉的黑暗.在黑暗中的人,盼望著星光.

經過上世紀的兩次世界大戰,經過無數的革命,無數的解放戰爭,人類已多少明白,要合作,相愛,光明才會來臨.二戰後組成的聯合國,其憲章的「序言」,是這樣寫的:「欲免後世再遭今代人類兩度身歷慘不堪言之戰禍,重申基本人權,人格尊嚴與價值,以及男女與大小各國平等權利之信念,創造適當環境,俾克維持正義,尊重由條約與國際法其他淵源而起之義務,久而弗懈,促成大自由中之社會進步及較善之民生.(為達此目的)力行容恕,彼此以善鄰之道,和睦相處.集中力量,以維持國際和平及安全.」

然而,在今天,這個世界看來又忘記了痛苦換來的教訓.在國際領域,很多國家強調自己優先,高舉狹隘的民族主義,不再寬容互助,而是以自己為中心,自私自利,種下仇恨的種子.在香港,政制的發展停滯,政府行事欠缺公道,法治的根基一步一步地鬆動,很多人,特別是年輕人,感到前景暗淡.

當國與國只顧自己,當社會裡的年輕人憂慮不安,當社會的運作只以強權來操控,天地又變得黑暗了.不知黑暗的日子多長,只知人類製造出來的黑暗真是可怕.

當最後一個先知死了的時候,以色列人只能在漫漫長夜中,苦候神的指引.我們或許與他們一樣,在黑暗中等待.四百年了,他們望向上蒼,看不見星星,只看見包圍著星星的無邊黑暗.四百年了,或許有人說,放棄吧!無望了!的確,四百年,已足夠長,去消磨人的意志與盼望.在夜空裡,還有沒有星星的呢?

黑夜已深,長夜漫漫,我們多麼需要星星.那怕是一點點的星光,已足夠燃點起希望.夜這麼深,還有星星嗎?

就是在這種時候,有人說:在最黑暗的環境,才會見到星星.

真有星星嗎?真有星星嗎?真有星星嗎?

太2:1　當希律王的時候,耶穌生在猶太的伯利恆.有幾個博士從東方來到耶路撒冷,說:

太2:2　「那生下來作猶太人之王的在哪裡?我們在東方看見他的星,特來拜他.」

四百年的黑暗,這麼悠長的黑暗,人連盼望和光明都忘記了.就在這晚,守望人類命運的智者,他們看見了「星」.這「星」,像是梵高(Vincent van Gogh)名畫《星夜》裡的星,畫面是陰沉的,前景的樹像是黑色的一團火,叫人不安.那時,梵高本人身處在精神病院裡.他自身的無望,畫面的無望,他感到的人類存在的無望,卻依然壓不住那幅畫上面那些明亮的星星所散發出來的光明和希望.那些星星,又大,又圓,又明亮,又有動感,讓人看後,忘記了畫面的憂鬱,而領悟到星星背後那個光芒萬象的神聖世界.智者們,在黑夜裡,看見了星.

剛才讀的經文,展示出一幅充滿張力的畫面.在一個王面前,那幾個智者,要尋索另一個王.希律王,是出名殘暴的.為了防止別人推翻他,會大開殺戒.為此,他殺了一個太太及三個兒子.在這希律王面前,智者們沒有向他下拜,而要去拜另一個王,這事之風險之大,猶如在獨裁者面前尋索民主.但是,因為智者們看到星星,看到從神而來的希望,看到渺少的人類想像不到的前景,他們敢於尋索那看來不可能尋到的,他們敢於去夢那不可能的夢想,他們在巨大的黑暗面前言說光明.

無可避免,希律王與「新生的王」會正面對峙,這是黑暗與光明的相遇,是暴力與和平的對決,是邪惡與神聖的交戰.

太2:3　希律王聽見了,就心裡不安；耶路撒冷合城的人也都不安.

必死的人,面對生命的主,必然是會「不安」的.從歷史的長河看,不義的政權與暴君,總會有天煙消雲散.只有正義長存,只有真理長存,只有愛長存.

太2:4　他就召齊了祭司長和民間的文士,問他們說:「基督當生在何處?」

太2:5　他們回答說:「在猶太的伯利恆.因為有先知記著,說:

太2:6　猶大地的伯利恆啊,你在猶大諸城中並不是最小的；因為將來有一位君王要從你那裡出來,牧養我以色列民.」

太2:7　當下,希律暗暗地召了博士來,細問那星是甚麼時候出現的,

太2:8　就差他們往伯利恆去,說:「你們去仔細尋訪那小孩子,尋到了,就來報信,我也好去拜他.」

希律王召見他的智囊,向他們問話,搜集情報,欺騙智者來作為他的探子,當然,其潛在目的只有一個,就是消滅那位初生的耶穌,消滅敵人於萌芽狀態.而耶穌,只是一嬰孩,一個手無寸鐵的嬰孩,一位完全無意使用暴力的嬰孩,但作為從神而來的光明,希望,救贖,已叫暴力的希律王亂了章法.這是人類靈魂的最終抉擇,去敬拜這位代表「和平與愛」的嬰孩,抑或以暴力去消滅他?

在這平安夜,總有人選擇拒絕平安.

然而,希律王的計謀,是不會得逞的.來自上帝的星星,再次出現,引導著智者.

太2:9　他們聽見王的話就去了.在東方所看見的那星忽然在他們前頭行,直行到小孩子的地方,就在上頭停住了.

太2:10　他們看見那星,就大大地歡喜.

這個世界,不是任由希律王主宰的.這個世界,有神默默地引導著.縱然黑夜漫漫,前路渺茫,但神仍在,祂眷顧人間,從祂而來的星光,仍照著大地,仍引導著高貴的靈魂,仍指示真理的大道,這就是平安夜最重要的信息.

太2:11　進了房子,看見小孩子和他母親馬利亞,就俯伏拜那小孩子,揭開寶盒,拿黃金,乳香,沒藥為禮物獻給他.

太2:12　博士因為在夢中被主指示不要回去見希律,就從別的路回本地去了.

真正的智者,是那些放棄世上的權勢,而看到「謙和」與「卑微」的偉大的人.在嬰孩耶穌那裡,他們體會到上天對世人的眷顧,他們看到黑暗中的希望,他們領悟到神對卑微者的救助.他們的心降服了.在那無比光明的嬰孩面前,他們俯伏下拜.在神聖的星光下,他們放下了自己.他們拿黃金,乳香,沒藥獻給嬰孩耶穌.黃金,王者之物,是獻給真正的王者的.乳香,是獻祭的材料,是獻給永恆的祭司用的.耶穌將永遠打通人走向神的道路.沒藥,是膏抹死者的用品,是獻給受死的耶穌膏抹屍體用的.耶穌要背負起蒼生的苦難,為人類的罪惡獻上自己為贖罪祭.在黑暗的世界,他燃點了自己.他的生命,定義了「平安」.

平安夜講的平安,不是輕音樂之下內心的浪漫平靜,也不是暴力滅聲後人人不敢作聲的死寂,而是明白天父的心,敢於在黑暗中相信光明,在光明中勇於去愛,在愛中散播希望.這種「平安」,就是聖經講的「Shalom」,就是生命的圓滿.

在這平安夜,願你見到星星,願你找到希望,願你心靈平安.

祝福你們.

但願:榮耀歸於聖父,聖子,聖靈.阿們.



\newpage

\section{人之子}
\label{sec:437}
\textbf{許立中先生 - 人之子}
\newline
\newline
\hyperref[sec:436]{< < < PREV SERMON < < <}
~
\hyperref[sec:toc]{[返目錄]}
~
\hyperref[sec:438]{> > > NEXT SERMON > > >}


路加福音 2:41-52
\newline
\begin{longtable}{cl}
\hline
\hline
章節 & 經文 (和合本修訂版)\\
\hline
2:41 & \begin{tabularx}{0.7\textwidth}{X} 每年逾越節，他父母都上耶路撒冷去。 \end{tabularx} \\ \\ \relax
2:42 & \begin{tabularx}{0.7\textwidth}{X} 當他十二歲的時候，他們按著過節的規矩上去。 \end{tabularx} \\ \\ \relax
2:43 & \begin{tabularx}{0.7\textwidth}{X} 守滿了節期，他們回去，孩童耶穌仍舊在耶路撒冷。他的父母並不知道， \end{tabularx} \\ \\ \relax
2:44 & \begin{tabularx}{0.7\textwidth}{X} 以為他在同行的人中間，走了一天的路程才在親屬和熟悉的人中找他， \end{tabularx} \\ \\ \relax
2:45 & \begin{tabularx}{0.7\textwidth}{X} 既找不著，就回耶路撒冷去找他。 \end{tabularx} \\ \\ \relax
2:46 & \begin{tabularx}{0.7\textwidth}{X} 過了三天，他們發現他在聖殿裡，坐在教師中間，一面聽，一面問。 \end{tabularx} \\ \\ \relax
2:47 & \begin{tabularx}{0.7\textwidth}{X} 凡聽見他的人都對他的聰明和應對感到驚奇。 \end{tabularx} \\ \\ \relax
2:48 & \begin{tabularx}{0.7\textwidth}{X} 他父母看見就很驚奇。他母親對他說：「我兒啊，為甚麼對我們這樣做呢？看哪，你父親和我很焦急，到處找你！」 \end{tabularx} \\ \\ \relax
2:49 & \begin{tabularx}{0.7\textwidth}{X} 耶穌對他們說：「為甚麼找我呢？難道你們不知道我應當在我父的家裡嗎？」 \end{tabularx} \\ \\ \relax
2:50 & \begin{tabularx}{0.7\textwidth}{X} 他所說的這話，他們不明白。 \end{tabularx} \\ \\ \relax
2:51 & \begin{tabularx}{0.7\textwidth}{X} 他就同他們下去，回到拿撒勒，並且順從他們。他母親把這一切的事都存在心裡。 \end{tabularx} \\ \\ \relax
2:52 & \begin{tabularx}{0.7\textwidth}{X} 耶穌的智慧和身量，並神和人喜愛他的心，都一齊增長。 \end{tabularx} \\ \\
[1ex]
\hline
\hline
\end{longtable}
講題:人之子 　The Son of Man

講員:許立中先生

經文:路加福音2:41-52

場合:香港中文大學崇基學院禮拜堂

日期:2018年12月30日

「每年到逾越節,他父母就上耶路撒冷去.當他十二歲的時候,他們按著節期的規矩上去.守滿了節期,他們回去,孩童耶穌仍舊在耶路撒冷.他的父母並不知道,以為他在同行的人中間,走了一天的路程,就在親屬和熟識的人中找他,既找不著,就回耶路撒冷去找他.過了三天,就遇見他在殿裡,坐在教師中間,一面聽,一面問.凡聽見他的,都希奇他的聰明和他的對應.他父母看見就很希奇.他母親對他說；『我兒!為甚麼向我們這樣行呢?看哪,你父親和我傷心來找你!』耶穌說:『為甚麼找我呢?豈不知我應當以我父的事為念嗎?』他所說的這話,他們不明白.他就同他們下去,回到拿撒勒,並且順從他們.他母親把這一切的事都存在心裡.耶穌的智慧和身量,並上帝和人喜愛他的心,都一齊增長.」(路2:41-52)

今天要跟大家思想的,是本週的福音經課,就是我們剛才讀過的路加福音2:41-52.這裡是一段其他福音書沒有記載,關於耶穌童年發生的一件事.不知道你有沒有留意到,這段童年往事特別之處,是在耶穌出生到他三十歲出來傳道之間,聖經唯一有記載關於他的事蹟.雖說「道成了肉身,住在我們中間」,但作為一個人,我們對於耶穌確實是認識甚少.

事實上每次當我在教會背誦《使徒信經》的時候,都會不期然感到有點不解.我的不解,乃在於《信經》中有關耶穌在世間的描述,竟然是幾乎完全越過；「我信我主耶穌基督,上帝獨生的子,因聖靈感孕,由童貞女馬利亞所生」,跟著下一句就已經是「在本丟彼拉多手下受難,被釘於十字架,受死,埋葬,降在陰間…」.

《使徒信經》在他的生和他的死之間,竟然連一句簡單的描述也沒有,彷彿他在世上的生活毫不重要,甚至不值一提.這讓人一個印象,就是他生存唯一的目的或意義,就是為了死.或許你也聽過這樣的論調,就是耶穌基督降生的意義,是在於他的死,而不是在於他的生.

我們暫且不去討論這個講法背後教義上的原因.但這樣的鋪排,確實很困難讓我們跟這位主建立一種比較親近的關係:你如何跟一個你根本不認識的人建立關係?而在華人教會,特別是福音派信仰傳統,與主建立一個親密的「個人關係」佔有一個相當重要的位置.只不過起碼在《信經》裡面的基督,確實令人感到有點面目模糊,難以接近.

在我們喜愛的聖詩中,有不少是描繪信徒跟主的關係,譬如“I come to the garden alone, while the dew is still on the roses, and the voice I hear falling on my ear, the Son of God discloses. And he walks with me and he talks with me, telling me I am his own.  And the joy we share as we tarry there, none other has ever known.” (“In the Garden”)在中文詩歌方面,如果你有點年資,或許都會記得那首冬令會必唱,蘇佐揚牧師的《你愛主嗎?》當然還有那首「主若是玫瑰一朵,我就是綠葉一片,主與我心心相印,緊緊地以愛結連.」然而,到底又是甚麼在支撐著這樣浪漫,淒美的關係?

當然,對於編寫《信經》的人來說,能夠囊括信仰的核心教義,可能是最首要的考慮,以致當這個群體不斷增長,人數不斷加添的時候,信徒能夠很快,很精確地記住幾個信仰的基本要義,就如一個國家的憲章,一個地方的基本法,一個社團的規章那樣.教會,有《使徒信經》.

但正如一部完善的國家憲法,並不能激勵人愛國；一本顧慮周全的基本法,亦不能阻止人們將它不斷重新解釋,到最後甚至完全脫離它原本的精神；一套社團的規章,同樣不能夠阻止人只是為了它的會所設施而加入成為會員.這樣,我們亦難以期望最正統基要的《使徒信經》,能夠對信徒群體起到甚麼關鍵的作用.它就只是信仰群體的一個共同認信.

無論如何,公元325年在安提阿舉行了第一次大公教會會議,處理當時不同信仰群體出現的信仰混亂.其中最重要的,當然就是基督信仰最核心的問題:耶穌到底是誰? 作為上帝的獨生子,他如何成為一個人?而明明是從母腹而生,他又在一個怎樣的意義上與上帝「同質」(homo-ousios)?而這一切,聖經似乎並沒有一個斬釘截鐵的現成答案.

我們打開新約聖經,保羅跟其他使徒所寫的書信,主要是對「基督事件」的註釋,引申.就算是四卷福音書,亦大部分是耶穌出來傳道之後的言行和講論,而甚少對他出來傳道之前的敘述和描繪.我們對耶穌的認識和印象,不少其實是從主日學,宗教畫,和一些聖經以外的媒介而獲得.

最近在互聯網上讀到一篇相當有趣的文章,題目為“Forensic Scientists Recreate the Face of Jesus - This is How he ‘Really’ Looked like”.

那篇文章指出,世代以來,人們很多時將他想像為一個長髮留鬚的中年漢,或許還加上憂鬱的眼神和有點蒼白的膚色.這個形象根深蒂固於人們的想像,是因為當歐洲人四出拓張他們的領土,征服新大陸的時候,都會帶同他們的神明；而那神明的形象,是按著他們最美好的想像建構出來,就好像上一代查爾登希士頓飾演的摩西那麼深入民心.

因此,我們對耶穌都有一些相當牢固,先入為主的想像,譬如從一些聖畫:左手抱著羊,右手拿著牧羊的杖；或者頭有光環,心打開了,攤開雙手,對眼充滿溫柔憐憫；我記得小時候,家裡就有一幅在客西馬尼園祈禱的耶穌.另外,當然還有那些月曆牌和荷里活電影:印象比較深刻的,是那齣由奧莉花荷西飾演馬利亞的耶穌傳,片中飾演童年耶穌的是個碧眼金髮,可愛到不得了的小男孩.

而剛才提到那篇文章推斷出來的耶穌,即根據當時的歷史,地理,文化,生活習慣,考古等重新塑造出來的耶穌,看起來竟然跟我們平常想像的施洗約翰那種形態差不多,是一個比較粗獷硬朗,黑黑實實的男人.用一個現代的想像,他就算不是查理士布朗臣,至少也該是羅素高爾.

事實上這個形象,確實比較接近長年從事木匠工作的人大概會有的外貌.事實上由耶穌出生到他三十歲出來傳道,從聖經的記載(或缺乏記載)推斷,他基本上是跟一個常人無異.以致當他開始出來傳道,福音書記載人們見到的時候,心裡會說:「那不過是個從拿撒勒出來的鄉巴佬!」而並非我們想像那麼英偉和文質彬彬.正如《以賽亞書》所預言,他並沒有「佳形美容」,以致我們看見的時候可以羨慕他.

當世界的潮流開始講indigenization(本土化,去殖民地化)的時候,我們亦開始從一些現代藝術看見非洲裔,黑皮膚的耶穌畫像.初看有點奇怪,但一直以來歐洲版的耶穌,並不見得比這些本土化的耶穌想像更接近事實的真相.

這一切「本土化」的努力,只能理解為人們渴望尋找一位他們能夠真正認同的對象,就是除了那位在高天之上坐著為王的上主,他們需要一位能夠體恤他們的軟弱和困局的同行者.而事實上這也是「道成肉身」一個重要的意義.

種種證據顯示,信徒並沒有困難接受耶穌是「上帝獨生的子」,卻似乎難以真心接受他同時是和我們一樣性情的人,以致每當在生命中遇到甚麼難題,信徒總是抱怨「我知道信仰對我有甚麼要求,可我卻不是耶穌啊!」

簡單來說,在他們的心目中,要實踐信仰,就要好像耶穌那樣有些非凡本事!久而久之,每次提到「效法基督」,其實就是提醒我們根本不是基督!那些所謂「基督的榜樣」或「基督徒的標準」,根本是強人所難的崇高理想:好是好,但是做不到.這是信徒在信仰實踐上經常遇到的困局.

然而反過來看,在福音書記載中,耶穌卻更多是以「人子」作為自稱:「人子來,為要尋找拯救失喪的人.」路19:10；「狐狸有洞,天空的飛鳥有窩,人子卻沒有枕頭的地方.」太8:20；「因為人子來,並不是要受人的服事,乃是要服事人,並且要捨命作多人的贖價.」可10:45

那麼,「人子」這個稱呼又是甚麼意思呢?這個名詞在舊約《以西結書》2:1及《但以理書》7:13-14曾經用過,分別泛指作為人的處境,以及那位將會在末日審判眾生的審判者.我想到一個我們可能比較容易明白的例子.

拖延了好幾年,政府終於決定預約拘捕當年發起「讓愛與和平佔領中環運動」的所謂「佔中三子」,控告他們涉嫌煽惑和引發後來歷時兩個多月的「雨傘運動」.對於這件案,不同的人按著他們不同的政治取向會有不同的判斷.不過我留意到有傳媒形容他們為「香港之子」,意思不單只是指他們土生土長香港人的身分,更暗示他們為了維護香港而盡心竭力.我想以此概念去推斷「人子」的意思,大概亦不會相距太遠.就是耶穌生而為人,也站在人的身分去維護人的利益和福祉.而耶穌作為一個人,福音書的記載是毫不含糊的.

《路加福音》:「當那些日子,凱撒奧古士督有旨意下來,要天下人民都報名上冊.這是居里扭作敘利亞巡撫的時候,頭一次行報名上冊的事.眾人各歸各城,報名上冊.」(2:1-3)

這個簡單的事實陳述,傳遞了一個重要的信息:耶穌是在一個特定的時間,出生於一個特定的地方.那時,羅馬帝國的開國君主Gaius Octavius權傾一時,被稱為凱撒奧古士督(他死後,元老院決定將他列入「神」的行列,並且將八月稱為「奧古士督」月,這也是歐洲稱八月份為August的來由).至於當時誰是敘利亞的總督,哪個分封王管轄猶太地區,都有確鑿的歷史記載.

而耶穌就是在這樣一個時候,出生於一個毫不起眼的窮鄉僻壤.他出生於伯利恆,彌迦書5:2形容為「在猶太諸城中是最小的」.他出生的時候,周圍就只有一些看羊人和他們的牲畜.正如福音書記載,當時個個都回鄉做戶籍登記,以致他的父母走遍全城,也找不到一個可以落腳的客棧,最後唯一有瓦遮頭的地方,就是一個馬棚.

就在那個邊緣的鄉鎮,那位本來不屬於這個世界的「道」成為了歷史,成為了這個塵世的一部分,成為一個軟弱無助的蒼生.但反過來說,透過他的來到,歷史本身,當然也包括我們的歷史,可以透過「道」成為有意義的存在.

關於基督的降生,讓我們不要忘記一個簡單的事實,就是好像所有嬰孩的一樣,耶穌是經過十月懷胎,產前的陣痛,出生於羅馬帝國管治下一個卑微的殖民地.如果我們忘記這個基本的事實,耶穌的出生就會變成一個抽象的符號,詩意的願景.而離開這個歷史的座標,聖誕事件就會失去它原本的意義.

就在今年的聖誕,我聽到不少基督徒對時局感到悲觀無奈.他們認為聖誕應該是充滿盼望的日子,但他們覺得在可見的未來,都看不到有甚麼出路.那慶祝聖誕還有甚麼意義?

或許這正是潘霍華在他的《獄中書簡》的提醒:信仰確實並不是一種理想主義的浪漫樂觀,相信只要有信心,願望就必定能夠達成,並且最好明天就發生.潘霍華在獄中發現到,「一個人必須完全過著今世的生活才能學習到信心…」而所謂「完全過著今世的生活」,就是「負起生命的一切責任和困難,成功和失敗,一切經驗和無可奈何之事.就在這樣的生命中,我們才把自己無條件地放置在上帝的手裡,參與祂在世上的苦難,與基督在客西馬尼園一同警醒.」而這種態度,對今天的基督徒來說是完全陌生的.今日的信徒,基本上是生活在一個奇異的幻想世界,以為做對了程序,問心無愧,願望就應該達成；期望稍有落差,就會立時不知所措,不知道該怎樣面對.

會不會在我們的信仰世界,確實一直都缺乏了一個「現世」的層面,就是「負起生命的一切責任和困難,成功和失敗,一切經驗和無可奈何之事.…將自己無條件地放置在上帝的手裡」,接受自己其實看不得多遠,改變不了多少?又會不會在耶穌的生平中,我們只是看見「道」,而看不見盛載著「道」那個沈重的「肉身」,其實亦受著現實世界種種的規限和制肘?

我當然不是企圖以「肉身」的具體性去暗示「道」不夠實在,而只是想指出,我們往往側重飄逸的靈性而令肉身變得異常沈重.今天的福音經課,正好提醒我們那位太初原有的生命之道,並不是好像我們的高官探訪公屋那樣,坐下飲啖茶,閒話家常幾句就收工；又或者好像有一次,某機構找幾個有錢人去住幾日tong 房,紆尊降貴地體驗一下基層市民的生活,然後表示多麼明白基層市民的處境和需要.但耶穌基督卻是實實在在成為一個好像你和我的人!

回到我們開始時的經文,在耶穌出生到他三十歲出來傳道之間,福音書就只是簡單記載「孩子漸漸長大,強健起來,充滿智慧,又有上帝的恩在他身上.」(路2:40)直到他十二歲,馬利亞和約瑟就帶他上耶路撒冷守逾越節.

按著猶太人的習俗,猶太的少年是在十三歲行成人禮,成為一個「律法之子」,正式成為猶太會堂的成員,承擔起割禮背後的一切義務和責任.因此,耶穌當時是差不多到了要行成人禮的年齡.他跟隨父母上耶路撒冷守逾越節,很可能是他行成人禮前的一次預習之旅；福音書沒有明確記載,但那有可能是他第一次隨父母上耶路撒冷.

經文記載,節期完畢,他們準備返家鄉的時候,馬利亞和約瑟以為耶穌跟同鄉小孩混在一起,所以沒有理會.況且耶穌未出過城,量他也不敢不跟著大隊.因此,「走了一天的路程」,他們才發現不見耶穌,於是唯有截返耶路撒冷,恐怕耶穌是因為人太多而走失了.

回到耶路撒冷,又過了三天,他們驚訝地發現耶穌竟然還留在聖殿,跟文士和教法師討論信仰問題.聖經沒有記載這幾天耶穌是怎麼過,吃些甚麼,住在哪裡?總之他就是氣定神閒地,在聖殿裡跟大人傾談關於聖經和信仰的問題.

這是一幅令人感到錯愕的圖畫:一個沒有受過經學訓練的孩子,竟然跟一班除了聖典之外就幾乎沒有其他興趣的學究,討論經典的問題!

想像一個十二歲的小朋友流落京城四,五天,而所有遇到他成年人,並不關心替他找回家人,而是興致勃勃地跟他討論信仰問題,這是不是有點奇怪?

更值得注意的是,那位本來「不可道」的道,並沒有以權威的姿態去跟那班學者專家對質.聖經只是記載他「一面聽,一面問」；奇怪的是,「凡聽見他的,都希奇他的聰明,和他的應對.」這不禁令人想起Socratic dialogues,用問題的方式,抽絲剝繭地引出問題的核心.而一個傳統猶太教的小孩,竟然用最經典的希臘哲學方式去探究真理.這不能不令人感到驚奇!

一個十二歲的孩童,有十二歲孩童的表達方式.而那位成為肉身的道,並沒有越過人性的定限,以神童的姿態自我呈現,技驚四座,叫人拍案叫絕.他反而完全安於人世間一切歷史,社會,文化,外在和內在的定限,恰如其分地表達自己.

讓我們不要忘記,他所做的其實也沒甚麼了不起:畢竟所有小孩子都喜歡問問題!而當他這樣做的時候,真理就很自然地透過點中問題的核心要害,令聽者由衷感到希奇和信服.

真理並不需要像聖旨那樣被宣告出來,而領受真理的人亦無須下跪接旨.更多時候,真理反倒是透過我們對真實問題的醒悟而呈現出來.也許我們都有經驗,我們是否聽得進一些道理,往往是決定於所說的是否也是我們心中的問題.在少年的耶穌身上,我們認識到真理是不需要用「大聲公」的.

在約瑟,馬利亞找到耶穌之後,他們向耶穌表達他們的擔憂.馬利亞對他說；『我兒!為甚麼向我們這樣行呢?看哪,你父親和我傷心來找你!』他們的反應其實也是人之常情.畢竟當時耶穌只是一個十二歲的孩童.反而耶穌當時的答覆,從任何角度來看都又點不近人情:『為甚麼找我呢?豈不知我應當以我父的事為念嗎?』因此,聖經說「他所說的這話,他們不明白」,其實完全可以理解.那個答覆,亦顯示他對自己的身分有所醒覺.

聖經跟著記載:「他就同他們下去,回到拿撒勒,並且順從他們.他母親把這一切的事都存在心裡.」在隨後的十八年,聖經就再沒有任何其他記載,只是簡單地說:「耶穌的智慧和身量,並上帝和人喜愛他的心,都一齊增長.」而從他對自己的身分有所醒覺,到他正式出來傳道這十八年,耶穌大概就是跟隨約瑟學習木匠的手藝,成為一個木匠,老老實實地做一個人.

在西方民主社會,基督教是一個相當體面的中產群體,有時甚至會不自覺地對相對落後的地方,制度和人表現得不理解甚至不耐煩.我們雖然不需要為自己的背景,出身和社會地位感到歉疚,但讓我們不要忘記,上帝的兒子,曾經出生於一個不顯眼的社群,在當時來說是一個已經失去「國家」身分的民族.他選擇出生的地方,是沒有出過任何顯赫人物的窮鄉加利利；他在一個基層的社區渡過他人生大部分的時間,向一班未見過世面的鄉下人宣揚一個仁愛和公義的國度.他甚至不像保羅那樣,有機會去到京城羅馬向有識之士申辯.他在生的日子,確實是藉藉無名.

基督信仰是一個歷史的信仰.上帝在耶穌基督身上揭示自己；而透過他,上帝同時也在人類歷史中揭示自己,特別是大部分無名的群眾.而這部分的歷史,往往是被正統歷史書所遺忘.對正統歷史來說,這些人只不過是一個毫無意義的統計數字.因此其實也只有以此作為起點,上帝信仰方才變得可能,意思是祂並不是體面的知識分子的專利.

我們必須在自己具體生命的處境中學習信心:不單只在各樣壓制和逼迫當中,也在我們周邊出現的掙扎和盼望之中；我們不僅反對違反人性的各種形色的獨裁主義,也挑戰不同旗幟下剝削人民的需要和夢想的民主政體.我們不將政治的訴求等同於信仰的理想,將政治的理想等同於上主的國度；明白到人間的制度,永遠不能絕對地克服人性的墮落.

曾經成為肉身的主,並沒有被黑暗所吞噬,或者被自己百姓的拒絕所壓倒.他的亮光能夠勝過一切的幽暗.我們必須進入自己當下的歷史,以周邊受壓迫的大眾所展示的生存意志,去滋養我們對生命的盼望.我不能說得比潘霍華更好；他說:「這就是信仰,就是悔改,也就是一個人所以成為人和基督徒的意義.這樣,當我們以生活於今世而參與上帝的苦難,成功怎能教我們傲慢,失敗又怎能使我們走入迷途呢?」



\newpage

\section{當東方智慧遇上基督}
\label{sec:438}
\textbf{白德培牧師 - 當東方智慧遇上基督}
\newline
\newline
\hyperref[sec:437]{< < < PREV SERMON < < <}
~
\hyperref[sec:toc]{[返目錄]}
~
\hyperref[sec:439]{> > > NEXT SERMON > > >}


馬太福音 2:1-12
\newline
\begin{longtable}{cl}
\hline
\hline
章節 & 經文 (和合本修訂版)\\
\hline
2:1 & \begin{tabularx}{0.7\textwidth}{X} 在希律作王的時候，耶穌生在猶太的伯利恆。有幾個博學之士從東方來到耶路撒冷，說： \end{tabularx} \\ \\ \relax
2:2 & \begin{tabularx}{0.7\textwidth}{X} 「那生下來作猶太人之王的在哪裡？我們在東方看見他的星，特來拜他。」 \end{tabularx} \\ \\ \relax
2:3 & \begin{tabularx}{0.7\textwidth}{X} 希律王聽見了，就心裡不安；耶路撒冷全城的人也都不安。 \end{tabularx} \\ \\ \relax
2:4 & \begin{tabularx}{0.7\textwidth}{X} 他就召集了祭司長和民間的文士，問他們：「基督該生在哪裡？」 \end{tabularx} \\ \\ \relax
2:5 & \begin{tabularx}{0.7\textwidth}{X} 他們說：「在猶太的伯利恆。因為有先知記著： \end{tabularx} \\ \\ \relax
2:6 & \begin{tabularx}{0.7\textwidth}{X} 『猶大地的伯利恆啊，你在猶大諸城中並不是最小的；因為將來有一位統治者要從你那裡出來，牧養我以色列民。』」 \end{tabularx} \\ \\ \relax
2:7 & \begin{tabularx}{0.7\textwidth}{X} 於是，希律暗地裡召了博學之士來，查問那星是甚麼時候出現的， \end{tabularx} \\ \\ \relax
2:8 & \begin{tabularx}{0.7\textwidth}{X} 就派他們往伯利恆去，說：「你們去仔細尋訪那小孩子，找到了就來報信，我也好去拜他。」 \end{tabularx} \\ \\ \relax
2:9 & \begin{tabularx}{0.7\textwidth}{X} 他們聽了王的話就去了。忽然，在東方所看到的那顆星在前面引領他們，一直行到小孩子所在地方的上方就停住了。 \end{tabularx} \\ \\ \relax
2:10 & \begin{tabularx}{0.7\textwidth}{X} 他們看見那星，就非常歡喜； \end{tabularx} \\ \\ \relax
2:11 & \begin{tabularx}{0.7\textwidth}{X} 進了房子，看見小孩子和他母親馬利亞，就俯伏拜那小孩子，揭開寶盒，拿出黃金、乳香、沒藥，作為禮物獻給他。 \end{tabularx} \\ \\ \relax
2:12 & \begin{tabularx}{0.7\textwidth}{X} 因為在夢中得到主的指示，不要回去見希律，他們就從別的路回自己的家鄉去了。 \end{tabularx} \\ \\
[1ex]
\hline
\hline
\end{longtable}
講題:當東方智慧遇上基督

Eastern Wisdom Meeting the Christ

經文:馬太福音2:1-12

講員:白德培牧師

場合:香港中文大學崇基學院禮拜堂

日期:2019年1月6日

The Magi Visit the Messiah

1 After Jesus was born in Bethlehem in Judea, during the time of King Herod, Magi from the east came to Jerusalem 2 and asked, “Where is the one who has been born king of the Jews? We saw his star when it rose and have come to worship him.” 3 When King Herod heard this he was disturbed, and all Jerusalem with him. 4 When he had called together all the people’s chief priests and teachers of the law, he asked them where the Messiah was to be born. 5 “In Bethlehem in Judea,” they replied, “for this is what the prophet has written:

6 “‘But you, Bethlehem, in the land of Judah, are by no means least among the rulers of Judah; for out of you will come a ruler who will shepherd my people Israel.’”

7 Then Herod called the Magi secretly and found out from them the exact time the star had appeared. 8 He sent them to Bethlehem and said, “Go and search carefully for the child. As soon as you find him, report to me, so that I too may go and worship him.” 9 After they had heard the king, they went on their way, and the star they had seen when it rose went ahead of them until it stopped over the place where the child was. 10 When they saw the star, they were overjoyed. 11 On coming to the house, they saw the child with his mother Mary, and they bowed down and worshiped him. Then they opened their treasures and presented him with gifts of gold, frankincense and myrrh. 12 And having been warned in a dream not to go back to Herod, they returned to their country by another route.

1. Intro

The story of the wise men from the East coming to seek the new born baby is one of the beautiful legends surrounding the birth of Jesus. It is with this story that we celebrate epiphany. It is a great story that perfectly fits into nativity scenes and into children’s performances of the Christmas story.

Yet behind this simply nice story there is a deeper story happening. It is the story of people from the East searching God. The question thus is:

a) who are these men?

b) why did they seek Jesus?

c) And what does this visit mean for us?

2. Who are the wise men?

It is commonly assumed that there were 3 wise men - but actually we do not know how many came. People commonly think of three wise men because three gifts are mentioned. Interestingly, orthodox Christians believe that there were 12 magi.

Maybe more interesting is who they are? The Bible introduces them as ‘magoi’ - ‘magi’ - meaning people who have special divine knowledge the special skill of religious divination and of astrology. They may have been religious advisers to the rulers in Persia. Many scholars suggest that they were Zoroastrians. Zoroastrianism was the religion of the Persian empire. Its founder was Zarathustra or Zoroaster in the 6th century before Christ. It had many parallels to Christianity but was also an important competitor to Christianity. It similarly claimed their leader being born by a virgin. Like Jesus Zarathustra started his ministry at the age of 30. And it similarly expected an eschatological savior like the Jews’ expectation of a Messiah.

But why did the wise men from the East come to worship Jesus?

3. Why did they come to seek Jesus?

Maybe the first question to ask is why Matthew included this story. Matthew is the only gospel telling this story. When Matthew introduces this story he may have had several reasons to do so. And these reasons may partly give us an idea why the wise men from the East have come to seek Jesus.

a) Receptivity of people from outside of the Jewish establishment

A first reason is that Matthew wanted to show that it is not the Jewish establishment that first discovers the Christ but people at the margins of society, people from outside Judaism such as the magi. We could call it a social reason. Indeed, I believe it is an on-going reminder for us that people from outside the establishment may be those who are better able to discover Christ in new contexts. Or in other words: God first reveals himself not to the social establishment but to the people at the margins of society. People from outside the establishment are more receptive to the Christian faith because establishment people are not interested in religious change. They are usually more interested in extending their own privileges, maintaining their power, and guaranteeing a stable society.

Matthew thus introduces the people from outside, Gentiles, as those who first come to worship the Christ. The magi stand in opposition to Herod who is not only the official king of the Jews but also a devout Jew who tried to show off his devotion by establishing strong relations with Jewish religious leaders.

b) Christ as better Zoroaster

But there may be a second reason why Matthew introduces this story: the story may suggest that the wise men from the East were dissatisfied with the teaching of their traditional faith and were looking for truth beyond their own tradition. We could call it a spiritual reason. Or in other words: Zoroastrians are seeking the Christ because they discover in Him the better Zoroaster. They are moved to seek beyond their own religious tradition. Again, this is a common pattern in people seeking Christ: they seek in him a revelation that stands in continuity to their traditional teaching but changes it in crucial areas.

c) Christ as king

Finally, there may be a third reason why Matthew introduces this story: the baby in the manger stands in starkest contrast to the power of the king Herod. We could call it a political reason. They worship him and bring him gifts that are typically given to a king. By doing so, they elevate him into the status of a king. Meanwhile, the king is left in Jerusalem, in his palace, seeking for a way to kill any attempt that undermines his rule. Remember, at the time of Jesus, two major empires faced each other: the Roman Empire and the Persian Empire. Israel laid at the fringes of the Roman Empire, not too far from the Persian Empire. The visit by assumedly high-ranking Persian religious figures coming to the Roman Empire and elevating Jesus to the rank of a king was a political challenge to both the Roman and the Persian Empire. The wise men from the East implicitly expressed that the true king is beyond the empire.

The story has tremendous actuality if we think of the attempts of present-day rulers to eradicate any threat to their rule. As we celebrate the coming of Christ, we declare him the only true king and power and thus challenge all worldly claims for kingship.

For Matthew, the worship of Christ begins with people from outside the social establishment, from outside the mainstream of society, from outsiders.

Ultimately, we do not know what drove the magi, but what we should remember is that it is their own religious tradition, their astrology and divination that moved them to seek Jesus. The motivation to seek Jesus was, we could say, intrinsic to their own tradition. It was not brought to them from the Christian tradition, but it came to them through their own religious sources.

4. Asians seeking Christ

Brothers and sisters, what moves me most in this story is still something else: 2000 years ago wise men, scholars from the East, left their home to seek Christ. This movement still goes on today. Nowadays many people in the East leave their tradition to seek truth in Christ. All of you know how in the past 40 years a tremendous growth of Christianity in China has happened. Historically speaking, this has been a unique and unprecedented growth movement. We could say that people in the East, people in China, increasingly seek the Christ. Among them are ordinary people, peasants, workers, elderly people. And among them there are also scholars

5. Scholars from the East seeking Christ today

A good number of scholars have turned to the Christian faith because they found a deep flaw in their own tradition. It is not that traditional faith is all bad; rather, it means that people discover inherent weaknesses. It may be the Confucian focus on hierarchical social relations and on obedience, which is what really allowed dictatorship to grow. Or it may be traditional religion of ancestor worship that too narrowly focuses on familial well-being and lacks a concept of love beyond the family.

Some of these scholars believe that the political model of democracy needs a different spiritual fundament - namely the Christian faith.

The question for these scholars as well as for the ordinary seekers from the East is: How do we find the Christ today? I believe that in this regard the beautiful Bible story still tells us an important lesson

6. How to find the Christ today?

The wise men from the East first go to Jerusalem, but soon learn from religious advisers in Jerusalem that the king is to be found in Bethlehem. I think there are two points we should note:

First, they do engage with local wisdom to find out more. They do not ignore local mainstream knowledge.

But then, and that is the second point, they are willing to move on from the capital of Jerusalem to the poor village around 9 km outside of Jerusalem. I think this is still an important lesson we should learn: Christ is not found in the capital, in the beautiful palaces, and in the centers of power, but at the margins - maybe in the slums, maybe in centers for refugees.

In other words: First, learn the knowledge of the center. But then, move beyond the center.

In fact, Jerusalem and Bethlehem are so near each other. Yet they are worlds apart. It was so in the past. And it is so today. Today, they even belong to different countries - Jerusalem is claimed by Israel and also by Palestine. But Jerusalem is more and more occupied by Israel. Bethlehem is within Palestine.

Nowadays, many Christians from HK make pilgrimages to the so-called holy land - which is so unholy and so full of war and conflict. The sad thing is that most of them go to Israel, but when going to Bethlehem they are led by guides from Israel and see Palestine only through the perspective of Israeli guides. This reinforces a common negative perception of Palestinians as only Muslim and as a place of terrorists.

As we are seeking Christ today, we should not go to the beautiful centers of power but we should find Christ among the poor. We need to learn to see reality through their eyes.

Yet, there is still another important lesson in this story:

7. Seeking Christ based on Eastern wisdom

The wise men from the East do not leave their own religious tradition. The story of the wise men from the East is not a conversion story where the wise men put down their own faith and adopt a new faith. No, it is a story that Zoroastrians seek Christ and worship him always while maintaining their own faith.

In the past term I had the opportunity to teach a course on the history of Christianity in Asia. It was the first time that I taught this course. Previously, we had courses on the history of Western Christianity and a course on the history of Christianity in China and a course on history of Christianity in HK, but never one on the history of Christianity in Asia. I wanted to fill this gap because there is a rich history of Christianity in Asia. It was an opportunity to address some basic issues of myself biographically, being born in the West and living in the East. And it was an opportunity to address issues of Christianity more generally, particularly regarding the common perception of Christianity as a Western religion. The course reminded me and my students that Christianity has always been an Asian religion and had simultaneously to the development in the West also a significant historical development in Asia.

However, one of the challenges that I encountered during the course was that much of Christianity in Asia had completely left their own tradition. They adopted a Western Christ instead of seeking Christ from within their own religious tradition.

Or in other words: they did not go to Bethlehem but they went to Athens, Rome, Wittenberg - or nowadays maybe Texas or Toronto. And, after going there, they did not return to their own context. The story of the wise men is an important reminder that we should seek the Christ in Bethlehem, which is part of Asia - not in Europe or North America.

I find that this great step of seeking Christ from within the Asian tradition has not yet happened.

Has traditional Christianity properly engaged with Asian wisdom?

And do present-day seekers of Christ really use their own religious tradition to worship Christ?

Are we not too often trying to substitute traditional religion?

One theologian who expressed it most beautifully is Kosuke Koyama (小山 晃佑 ) who was a missionary in Northern Thailand. He tried to find expression of Christ in a way that would not replace the wide-spread Buddhist faith. He maintained appreciation for the Thai people’s Buddhist belief and talked about the cool Buddha and the hot Jesus.

8. Conclusion

Brothers and sisters - we started with a beautiful and famous Christmas story. And we end up being challenged.

The story challenges traditional Christianity, reminding us that it is the Gentiles, people from outside the mainstream, from outside the establishment structures, who first recognize the Christ and worship him like a king!

It is a story that challenges Eastern wisdom, reminding us that the Christ is to be sought not in the centers of power but at the margins of the empire. How much are we ready to seek Christ not in the centers of power but among the marginalized? How much are we ready to go to Bethlehem and not to Jerusalem?

Finally, it is challenges us, Asian Christians: How much are we ready to go to Bethlehem instead of Rome and Wittenberg and Toronto? How much are we ready to find Christ from within our own tradition?



\newpage

\section{死亡眼淚與生命}
\label{sec:439}
\textbf{劉國偉先生 - 死亡,眼淚與生命}
\newline
\newline
\hyperref[sec:438]{< < < PREV SERMON < < <}
~
\hyperref[sec:toc]{[返目錄]}
~
\hyperref[sec:440]{> > > NEXT SERMON > > >}


約翰福音 11:17-44
\newline
\begin{longtable}{cl}
\hline
\hline
章節 & 經文 (和合本修訂版)\\
\hline
11:17 & \begin{tabularx}{0.7\textwidth}{X} 耶穌到了，知道拉撒路在墳墓裡已經四天了。 \end{tabularx} \\ \\ \relax
11:18 & \begin{tabularx}{0.7\textwidth}{X} 伯大尼離耶路撒冷不遠，約有六里路。 \end{tabularx} \\ \\ \relax
11:19 & \begin{tabularx}{0.7\textwidth}{X} 有好些猶太人來看馬大和馬利亞，要為她們弟弟的緣故安慰她們。 \end{tabularx} \\ \\ \relax
11:20 & \begin{tabularx}{0.7\textwidth}{X} 馬大聽見耶穌來了，就出去迎接他；馬利亞卻仍然坐在家裡。 \end{tabularx} \\ \\ \relax
11:21 & \begin{tabularx}{0.7\textwidth}{X} 馬大對耶穌說：「主啊，你若早在這裡，我弟弟就不會死了。 \end{tabularx} \\ \\ \relax
11:22 & \begin{tabularx}{0.7\textwidth}{X} 我也知道，即使現在，你無論向神求甚麼，神也必賜給你。」 \end{tabularx} \\ \\ \relax
11:23 & \begin{tabularx}{0.7\textwidth}{X} 耶穌對她說：「你弟弟會復活的。」 \end{tabularx} \\ \\ \relax
11:24 & \begin{tabularx}{0.7\textwidth}{X} 馬大對他說：「我知道在末日復活的時候，他會復活。」 \end{tabularx} \\ \\ \relax
11:25 & \begin{tabularx}{0.7\textwidth}{X} 耶穌對她說：「復活在我，生命也在我。信我的人雖然死了，也必復活。 \end{tabularx} \\ \\ \relax
11:26 & \begin{tabularx}{0.7\textwidth}{X} 凡活著信我的人必永遠不死。你信這話嗎？」 \end{tabularx} \\ \\ \relax
11:27 & \begin{tabularx}{0.7\textwidth}{X} 馬大對他說：「主啊，是的。我信你是基督，是神的兒子，就是那要臨到世界的。」 \end{tabularx} \\ \\ \relax
11:28 & \begin{tabularx}{0.7\textwidth}{X} 馬大說了這話就回去，叫她妹妹馬利亞，私下說：「老師來了，他在叫你。」 \end{tabularx} \\ \\ \relax
11:29 & \begin{tabularx}{0.7\textwidth}{X} 馬利亞聽見了，急忙起來，到耶穌那裡去。 \end{tabularx} \\ \\ \relax
11:30 & \begin{tabularx}{0.7\textwidth}{X} 那時，耶穌還沒有進村子，仍在馬大迎接他的地方。 \end{tabularx} \\ \\ \relax
11:31 & \begin{tabularx}{0.7\textwidth}{X} 那些同馬利亞在家裡安慰她的猶太人，見她急忙起來，出去，就跟著她，以為她要往墳墓那裡去哭。 \end{tabularx} \\ \\ \relax
11:32 & \begin{tabularx}{0.7\textwidth}{X} 馬利亞到了耶穌那裡，看見他，就俯伏在他腳前，對他說：「主啊，你若早在這裡，我弟弟就不會死了。」 \end{tabularx} \\ \\ \relax
11:33 & \begin{tabularx}{0.7\textwidth}{X} 耶穌看見她哭，並看見與她同來的猶太人也哭，就心裡悲嘆，又甚憂愁， \end{tabularx} \\ \\ \relax
11:34 & \begin{tabularx}{0.7\textwidth}{X} 就說：「你們把他安放在哪裡？」他們對他說：「主啊，請你來看。」 \end{tabularx} \\ \\ \relax
11:35 & \begin{tabularx}{0.7\textwidth}{X} 耶穌哭了。 \end{tabularx} \\ \\ \relax
11:36 & \begin{tabularx}{0.7\textwidth}{X} 猶太人就說：「你看，他多麼愛他！」 \end{tabularx} \\ \\ \relax
11:37 & \begin{tabularx}{0.7\textwidth}{X} 其中有人說：「他既然開了盲人的眼睛，難道不能叫這人不死嗎？」 \end{tabularx} \\ \\ \relax
11:38 & \begin{tabularx}{0.7\textwidth}{X} 耶穌又心裡悲嘆，來到墳墓前。那墳墓是個穴，有一塊石頭擋著。 \end{tabularx} \\ \\ \relax
11:39 & \begin{tabularx}{0.7\textwidth}{X} 耶穌說：「把石頭挪開！」那死者的姐姐馬大對他說：「主啊，他現在必定臭了，因為他已經死了四天了。」 \end{tabularx} \\ \\ \relax
11:40 & \begin{tabularx}{0.7\textwidth}{X} 耶穌對她說：「我不是對你說過，你若信就必看見神的榮耀嗎？」 \end{tabularx} \\ \\ \relax
11:41 & \begin{tabularx}{0.7\textwidth}{X} 於是他們把石頭挪開。耶穌舉目望天，說：「父啊，我感謝你，因為你已經聽了我。 \end{tabularx} \\ \\ \relax
11:42 & \begin{tabularx}{0.7\textwidth}{X} 我知道你常常聽我，但我說這話是為了周圍站著的眾人，要使他們信是你差了我來的。」 \end{tabularx} \\ \\ \relax
11:43 & \begin{tabularx}{0.7\textwidth}{X} 說了這些話，他大聲呼叫說：「拉撒路，出來！」 \end{tabularx} \\ \\ \relax
11:44 & \begin{tabularx}{0.7\textwidth}{X} 那死了的人就出來了，手腳都裹著布，臉上包著頭巾。耶穌對他們說：「解開他，讓他走！」 \end{tabularx} \\ \\
[1ex]
\hline
\hline
\end{longtable}
講題:死亡,眼淚與生命 Death, Tear and Life

經文:約翰福音11:17-44

講員:劉國偉先生

場合:香港中文大學崇基學院禮拜堂

日期:2019年1月13日

引子

死亡與我們的人生經驗,甚至人類歷史糾纏.不同性格,信念,又或者不同的宗教和文化背景的人,自然對死亡有不同的詮釋和觀點.

使徒約翰在垂暮之年寫下耶穌福音紀事的回憶錄,記載了拉撒路復活的獨家故事,在約翰福音第11章,用上了與死亡相關的辭彙15次,牽引著讀者從一次死而復活的神蹟中,認識耶穌更深,也認識自己和人性更多.

突破創辦人蘇恩佩女士在她的著作「死亡,別狂傲」中,引用英國十六世紀著名詩人約翰多恩(John Donne)的作品:

死亡,別狂傲,縱或有人稱你

聲勢駭人,然而並非如此；

那些你自信可以推翻的人

是不滅的,可憐的死亡,你未能殺我.

那是恩佩姊妹對死亡的吶喊,對於一個活在死亡邊緣二十年的人,是對死亡一個透徹的認識.

死亡

約翰把整件事情仔細地記錄下來:馬大和馬利亞兩姊妹特意打發人,去告知耶穌有關拉撒路病重的消息,耶穌聽見後,卻仍在原處逗留兩天,作為這個家庭的一位摯友,耶穌沒有立即動身,確是叫人費解.兩天之後,耶穌告訴門徒要再往猶太地,往訪馬大一家,耶穌的門徒難免有點莫名其妙,要不兩天前就馬上啟程,拉撒路既已死去,又為何冒著人身安全,再踏足猶太地,面對被猶太人拿著石頭追打的際遇!

在約翰的筆下,他記述了最少三種對死亡的回應.

啟程之際,多馬怨氣地說:「我們也和拉撒路同死罷!」(11:16)表達了死亡是終結,不要白白的犧牲,「咁樣死值得咩?」

「有好些猶太人來看馬大和馬利亞,要為她們的兄弟安慰她們.」(11:19)死亡是一回十分傷痛的事,需要別人安慰.

「主啊!你若早在這裡,我兄弟必不死.」(11:21)死亡是一個不可逆轉的結局,而死去與未死,有時只是一剎之間.

耶穌等了兩天,走了兩天路程,拉撒路已經辭世4天.猶太民間相傳,人死了4天之後,靈魂就離開身體,換言之,拉撒路決不可能復活再生.

馬大知道耶穌正要進村莊,便趕忙出去迎接,碰面的第一句話,相信是她埋藏心底,纏繞好幾個日夜的由衷之言:「主啊!你若早在這裡,我兄弟必不死.」這個小片段剛過去,馬大的妹妹馬利亞也知道這位遲來的夫子已到了村口,她同樣急忙起來要馬上與耶穌見面,身旁的人也馬上跟在後面,還以為她要去哭墳.

當馬利亞遇見耶穌,劈頭第一句話竟與姊姊馬大所說的完全一樣:「主啊!你若早在這裡,我兄弟必不死.」,看來,這已是她們姊妹二人的一致結論,當中也許沒有怨恨,但那種失望,不解,感慨,和對事情無法挽回的嗟嘆,都全部在她們的臉容,眼淚和身體語言中表露無遺.

眼淚

使徒約翰記下了福音書中耶穌哭泣的一個片段,那肯定是深入魂魄的一個記憶.耶穌與門徒身處在一種張力之中,走了兩天路程,方才抵達伯大尼,就遇上了馬大和馬利亞相同的說話,還有,耶穌與馬大的對話,似有若無.

馬大對耶穌說:「主啊,你若早在這裡,我兄弟必不死.就是現在,我也知道,你無論向上帝求甚麼,上帝也必賜給你.」耶穌說:「你兄弟必然復活.」馬大說:「我知道在末日復活的時候,他必復活.」耶穌對她說:「復活在我,生命也在我.信我的人雖然死了,也必復活；凡活著信我的人必永遠不死.你信這話嗎?」馬大說:「主啊,是的,我信你是基督,是上帝的兒子,就是那要臨到世界的.」 (11:21-27)

馬大對教義的掌握,是清晰及準確的,卻正如我們一樣,往往被自己的經驗所局限,以為上帝會做甚麼,或不會做甚麼! 我們愈有生活經驗,或者經常作決策的人,愈容易陷入這種定見和定局之中,沒有一個足夠開放的心靈,甚至不願意開放自己的心靈,去碰觸信仰旅程中的新的體驗.

就在這個時空背景之下,使徒約翰記載了不同的人所流下的眼淚.

至親的眼淚

那些同馬利亞在家裡安慰她的猶太人,見她急忙起來出去,就跟著她,以為她要往墳墓那裡去哭.馬利亞到了耶穌那裡,看見他,就俯伏在他腳前,說:「主啊,你若早在這裡,我兄弟必不死.」耶穌看見她哭,(11:31-33上)

馬利亞為拉撒路的死去而哭泣,是完全可以理解.猶太人以為她要去哭墳,也是正常的估算.

親友的眼淚

耶穌看見她哭,並看見與她同來的猶太人也哭(11:33)

與哀哭的人同哭,乃人之常情……

無論是至親的眼淚,或者親友的眼淚,都是基於死亡是一個不可逆轉的結局,悲從中來.耶穌呢? 祂明明知道頃刻之間,拉撒路就要復活,而且祂就是復活,就是生命,何以祂也哭起來??

3.「耶穌哭了.」(約11:35)

耶穌哭了:祂以生命與人間的情愫糾結

此刻,馬利亞俯伏在耶穌腳前哭泣,陪著馬利亞前來的友好同樣帶著淚眼,大抵都認同她們姊妹倆所下的結論 : 夫子來晚了!.在淚眼汪汪之間,一片愁雲慘霧,耶穌再沒有解說甚麼,在眾人心中,一切已成定論,耶穌那份無人明白的苦惱,正加倍凝結.

使徒約翰憶起耶穌那一刻心裡悲嘆,又甚憂愁.天地的主理當胸有成竹,知道頃刻之間眼淚將變成歡笑,一切的嗟嘆將化成喜悅,然而,祂卻與人間的情愛相纏,他難過地問:「你們把他安放在那裡?」(難道祂不知道嗎?),就在大夥兒動身前往墳墓之際,耶穌哭了.

門徒的怨氣和不解,馬大和馬利亞的失望和感嘆,親友同哀的眼淚,身處人間那種被死亡所捆綁和別離時的傷痛,在哀哭的人當中,耶穌悲從中來,祂體味到凡塵的苦,在前因後果的交纏之間,祂像凡人一樣流下了眼淚.

生命

身處現場的猶太人開始有所議論,約翰記下了兩段說話:

猶太人就說:「你看他愛這人是何等懇切.」 (11:36)

其中有人說:「他既然開了瞎子的眼睛,豈不能叫這人不死嗎?」 (11:37)

信仰有時就是如此直接和簡單! 你看,神愛你這個人是何等懇切! 祂既然開了救恩的道路給你,豈不能叫你從死亡般的困境走出來嗎?

與我們同行這條生命之道的主,是一位以生命與人間的情愫糾結的主!哭泣的主不單讓我們認識他真是一個人,祂更呈現出神與人同行的赤裸裸的事實,祂親自進入凡人的生離死別的苦痛裡,深情地在人間活過.祂哭起來――一個會因為人神之間的糾葛而哭泣的主,加倍叫我認識祂對人間的愛有多深.

死亡與生命之間

在拉撒路的死亡與復活之間,是耶穌的眼淚,因為上帝同行的愛,叫我們可以面對死亡,活出榮耀上帝的生命.正如耶穌說:「我不是對你說過,你若信,就必看見上帝的榮耀嗎?」(11:40)

是耶穌,不是別人,大聲呼叫說:「拉撒路出來!」那死人就出來了,手腳裹著布,臉上包著手巾.(11:43-44上) 「拉撒路」就是「神是我的幫助」的意思.在這個愈來愈多人自我設限,自我審查的社會,多少活在害怕得失權勢的擔憂甚至恐懼之中,我們可會像使徒約翰般銘記耶穌就是復活,就是生命,祂以深情的愛與我們從死亡走往生命,確信及經歷「神是我的幫助」.

哥林多前書15章55-56節 : 「死啊!你得勝的權勢在哪裡? 死啊!你的毒鉤在哪裡? 死的毒鉤就是罪」(林前15:55-56上)

面對死亡,面對罪的權勢,我們要記住這一個畫面 : 耶穌對他們說:「解開,叫他走!」



\newpage

\section{來唱歌吧}
\label{sec:440}
\textbf{鄧瑞強博士 - 來唱歌吧}
\newline
\newline
\hyperref[sec:439]{< < < PREV SERMON < < <}
~
\hyperref[sec:toc]{[返目錄]}
~
\hyperref[sec:441]{> > > NEXT SERMON > > >}


啟示錄 15:1-4
\newline
\begin{longtable}{cl}
\hline
\hline
章節 & 經文 (和合本修訂版)\\
\hline
15:1 & \begin{tabularx}{0.7\textwidth}{X} 我看見在天上有另一兆頭，大而且奇，就是七位天使掌管末了的七種災難，因為神的烈怒在這七種災難中發盡了。 \end{tabularx} \\ \\ \relax
15:2 & \begin{tabularx}{0.7\textwidth}{X} 我看見彷彿有攙雜火的玻璃海；又看見那些勝了那獸和獸像，以及牠名字的數字的人，都站在玻璃海上，拿著神的豎琴。 \end{tabularx} \\ \\ \relax
15:3 & \begin{tabularx}{0.7\textwidth}{X} 他們唱神僕人摩西的歌和羔羊的歌，說：「主—全能的神啊，你的作為又偉大又奇妙！萬國之王啊，你的道路又公義又真實！ \end{tabularx} \\ \\ \relax
15:4 & \begin{tabularx}{0.7\textwidth}{X} 主啊，誰敢不敬畏你，不把榮耀歸於你的名？因為只有你是神聖的。萬民都要來，在你面前敬拜，因你公義的作為已經彰顯了。」 \end{tabularx} \\ \\
[1ex]
\hline
\hline
\end{longtable}
講題:來唱歌吧  Come Sing a Song

經文:啟示錄15章1-4節

講員:鄧瑞強博士

場合:香港中文大學崇基學院禮拜堂

日期:2019年1月20日

各位弟兄姊妹,早安.

不知你洗澡時會不會唱歌?很多人洗澡時唱歌.聽說,因為這是人最放鬆的時刻；還有,在熱水下,血液循環加速,人較興奮；再加上獨自一人,不用理會別人的批評；而且,在密室中,有點回音效果；加上少許水氣,矇矇矓矓,有點舞台效果,於是,人就會唱歌,將內心的情緒舒發出來.

日本有位詩人谷川俊太郎(Shuntaro Tanikawa),寫了首詩《我歌唱的理由》,講他為何唱歌.這首詩如此說:

我歌唱／是因為一隻小貓崽／被雨澆透後死去／一隻小貓崽

我歌唱／是因為一棵山毛櫸／根糜爛掉枯死／一棵山毛櫸

我歌唱／是因為一個孩子／瞠目結舌,顫驚呆立／一個孩子

我歌唱／是因為一個單身漢／蹲下來背過身子往別出看／一個單身漢

我歌唱／是因為一滴淚／滿腹委屈和焦躁不安／一滴清淚

(田原譯,登於:\url{https://read01.com/DGe5POA.html}\#.XDnwuNSF5XQ)

想想這首詩的畫面:一隻貓仔,被雨濕透,死去.一棵山手櫸,理應很茁壯,根卻壞掉,枯死.一個孩子,在驚恐中,不知所措.一個單身漢,卑屈地蹲下來,卻不知望向何方.一滴淚,在委屈與不安下,悲苦地流下.

這就不是洗澡時隨意唱一首歌,而是有感於生命的柔弱與無助.心中很想向這些至柔弱者伸出援手,但又知自己的手伸不過去,在想作出幫助而又作不到幫助的痛苦裡,詩人不得不唱出心中的鬱悶,不得不唱出悲憫的歌.這就是詩人唱歌的理由,他唱出他對蒼生的悲憫.

美國有一隊唱福音民歌的二人組合,叫All Sons and Daughters,唱了一首歌,歌名叫 “Reason to Sing”(歌唱的理由).詩歌大致說:世事像碎片,我感受不到神.我恐懼無助,恨惡一切.在這時候,I need a reason to sing.多麼需要知道神仍然以手托著這個世界,而這就是歌者仍能唱歌的理由.

一個人唱歌,可能是洗澡時在花灑下感到舒暢,可能是體會到萬物如此無助而產生的悲憫,又或者,是因為在世界崩壞的時候,看到神仍掌管著萬有.

今日講道的經文,寫在邪惡當道之時,那時怪獸橫行,謊言成了社會運作的常態.這時,作者仍能找到歌唱的理由,因為他看到活著的上帝,仍在人間主持正義,扭轉乾坤.今日,讓我們看看他歌唱的理由.

講道經文:啟示錄15:1－4

啟15:1　我又看見在天上有異象,大而且奇,就是七位天使掌管末了的七災,因為神的大怒在這七災中發盡了.

啟15:2　我看見彷彿有玻璃海,其中有火攙雜.又看見那些勝了獸和獸的像並牠名字數目的人,都站在玻璃海上,拿著神的琴,

啟15:3　唱神僕人摩西的歌,和羔羊的歌,說:主神─全能者啊,你的作為大哉!奇哉!萬世之王啊,你的道途義哉!誠哉!

啟15:4　主啊,誰敢不敬畏你,不將榮耀歸與你的名呢?因為獨有你是聖的.萬民都要來在你面前敬拜,因你公義的作為已經顯出來了.

啟15:1　我又看見在天上有異象,大而且奇,就是七位天使掌管末了的七災,因為神的大怒在這七災中發盡了.

世事紛陳,擾擾攘攘,作者不從人間的角度看,他仰望上天,他從天上看.他看見天上有異象,大而且奇.「大」,因涉及整個世界及歷史的命運.「奇」,是因為意想不到.看看世界,邪惡當道,還以為這世界會在邪惡中沉淪,還以為堅持真理的人會死得不明不白,但原來事實不是這樣的.神在人間,正道在人間,公義在人間.神對邪惡震怒,發出祂的審判.按舊約「出埃及」事件的記述,神審判法老前,先有「十災」.神作最後審判前,先有人間的災劫.經文這裡提到「七災」,是審判的前奏.神讓人在災劫中反省,悔改.但邪惡之謂邪惡,總是在任何災劫中都不會承認自己的錯,造成災劫的總是別人的錯.當人人都指控魔鬼是邪惡時,惟魔鬼會走出來,真誠地,大聲說,我不是邪惡的,你們才是邪惡的.魔鬼殺一個人時,都是正氣凜然的.

面對人間的邪惡,啟示錄是從天上的角度去俯瞰,同時以歷史的終局的觀點去回望,邪惡好像橫行霸道,但永遠不是歷史的主人,也永遠不是主宰歷史的神明.只要我們的視野「有從天上看」的高度和「有從終點看」的遠大,我們便能與啟示錄的作者一樣,看到邪惡的最後崩壞.

啟15:2　我看見彷彿有玻璃海,其中有火攙雜.又看見那些勝了獸和獸的像並牠名字數目的人,都站在玻璃海上,拿著神的琴,

啟15:3　唱神僕人摩西的歌,和羔羊的歌………

在舊約,當以色列人擺脫法老,不再當奴隸時,他們逃走,走到紅海邊,前無去路,後面有法老的追兵.在這個絕境裡,「大而且奇」的事就出現了.摩西舉起杖,紅海就分開.按猶太人的一個傳統說,紅海分開時,海水凝結成冰,整個海面變成玻璃海.過了紅海後,在海邊,他們唱起讚頌上帝救贖的歌.(參:\url{http://www.schechter.edu/visual-midrash-on-the-parting-of-the-red-sea-parashah-beshalach/})

啟示錄用了相關的意象,當神終極審判邪惡時,神的子民也站在玻璃海上,唱著救贖的歌.「海」,在近東文明,代表混亂與邪惡.海裡有怪獸,令人不安.如今,海被凝結成玻璃,不再顛簸,也變得透明,怪獸無法隱藏.海被克服了,邪惡被克服了.多麼希望我們的政治世界也變成一個玻璃海,一切運作變得透明,也無怪獸隱藏在裡面.玻璃海裡有「火」,因為神的審判在其中.神對子民的救贖,伴隨著對邪惡的審判.

站在海邊的,是神的子民,他們是「勝了獸和獸的像並牠名字數目的人」.神的子民勝過三樣東西:獸,獸像,獸的數目.「獸」是邪惡的根源.邪惡要發揮力量,需要邪惡的媒介或體制,這就是「獸的像」.「獸像」是偶像化的崇拜體制,這不單單是宗教體制,也同時是政治,社會,經濟的體制.舉例說,當我們以為金錢比一切都重要時,金錢就不只是經濟的東西,而是控制人生活的偶像神明.偶像之為偶像,就在於它的「假」.偶像崇拜,就是將「假」當成「真」.人將「假」當成「真」的時候,就會將「真人性」淪落為「假人性」,將人性「物化」,將自己淪為邪惡的工具.邪惡的怪獸要控制人,透過「獸像」,透過將人「物化」的體制,使人失去人性.人選擇進入「獸像」的體制敬拜「獸」,生命就被打上「獸的數目」,生命屬於「獸」,成為獸的一份子.

我想起魯迅寫的《狂人日記》,小說的主人翁是一個「狂人」,他的日記這樣寫:「我翻開歷史一查,這歷史沒有年代,歪歪斜斜的每頁上都寫著『仁義道德』幾個字.我橫豎睡不著,仔細看了半夜,才從字縫裡看出字來,滿本都寫著兩個字是『吃人』!」這個狂人認為他身邊的人都是密謀要吃他的.他這樣寫:「我認識他們是一夥,都是吃人的人.可是也曉得他們心思很不一樣,一種是以為從來如此,應該吃的；一種是知道不該吃的,可是仍然要吃.」這是一個「人吃人」的世界.這日記的最後兩行字是:「沒有吃過人的孩子,或者還有?救救孩子……」

魯迅筆下的世界,是一個「怪獸化」的世界,「人吃人」的世界,非人化的世界.每個人的生命,都是怪獸體系的一份子.他們吃人的時候,口中是說著「仁義道德」的.任何為了某種理想,某個體制的利益而漠視他人生命的,都是「怪獸」體系.在其中,人順從了「獸像」,打上「獸的數目」,成了「獸」的一份子.魯迅的悲憫是,看看還有沒有未吃過人的孩子,看看還有沒有仍保持人性的人,若有的話,要救救他們.這個世界,期待神的真子民.神的真子民,是那些「勝了獸和獸的像並牠名字數目的人」,這些人能救救孩子,教他們不去吃人,教他們像耶穌一樣,以仁愛及真理拒絕暴力與謊言.

神的子民,站在玻璃海上,拿著神的琴,唱神僕人摩西的歌,和羔羊的歌.他們唱歌,總有唱歌的理由.在艱難的歲月裡,他們為何能唱歌呢?

啟15:3　唱神僕人摩西的歌,和羔羊的歌,說:主神─全能者啊,你的作為大哉!奇哉!萬世之王啊,你的道途義哉!誠哉!

啟15:4　主啊,誰敢不敬畏你,不將榮耀歸與你的名呢?因為獨有你是聖的.萬民都要來在你面前敬拜,因你公義的作為已經顯出來了.

神被稱為「主,神,全能者」.主,是僕人效忠的對象.神,是被造物敬拜的對象.全能者,是戰無不勝的戰士,是被欺壓的弱者倚靠的對象.

在這位神面前,我們為何能唱歌呢?

回看過去,這位神為我們成就了大事.經文說,祂的作為「大」哉「奇」哉,「大而且奇」,重複第一節的講法.祂的道途「義」哉「誠」哉,義哉,正義也,賞善罰惡；誠哉,誠實也,堅持真理.神「大而且奇」的作為,是讓正義明明白白地彰顯出來.作為全能者,祂保證擊敗邪惡,主持正義.在耶穌的十架受死裡,神顯明邪惡不一定能迫使人變成邪惡的人,耶穌以仁愛與寬恕向邪惡說「不」,祂的復活顯明這種生命能進入永恆.暴力與謊言,最多能欺騙人十年百年,惟正義與真理永存不朽.

就現在來說,神希望我們明白「獨有祂是聖的」.整卷啟示錄,就以這句「獨有你是聖的」為中心.魔鬼在人間,製造很多偶像,誘惑我們也迫使我們去敬拜它們.金錢,利益,權勢,名譽,享樂,種種東西,都以「仁義道德」的名義,吸引我們向其下拜.在我們毫不察覺的情況下,使我們成為「吃人的人」,成為怪獸的一份子.啟示錄在這群魔亂舞的世代,大聲說:「獨有神是聖的」.

展望將來,神希望萬民來到祂面前敬拜祂,這是教會的使命.人生在世,不要為自己而活.為自己而活,很容易跌入世間種種偶像崇拜裡,成為「吃人的人」.只要你為自己而活,常常想著自己的好處,偶像就能乘虛而入,以虛假的成功換取你的靈魂.在太初,蛇能成功引誘亞當夏娃,就是讓他們以為跟從牠的做法會帶來自身的好處.教會在世的使命,是效法基督,惟獨順服天父的旨意,因為深深知道惟獨天父是聖的.魯迅呼籲:「沒有吃過人的孩子,或者還有?救救孩子……」教會效法基督,就能救救孩子.在如此多偶像的世界,人真的很易失去人性,變成吃人的怪獸.神的子民真的要加把勁,堅持真理,活得像人,實踐天父的旨意,渴望天國的來臨.

我希望,我們每個人都能加入這個詩班,站在玻璃海上,唱讚頌神的詩歌.親愛的弟兄姊妹,來,唱首歌,唱首羔羊的歌.當有一天,有人問我們,我們歌唱的理由何在?我們可以答:

我歌唱／是因為一隻小貓崽／被雨澆透後／耶穌溫暖著牠

我歌唱／是因為一棵山毛櫸／根看來爛掉了／卻抽出新芽

我歌唱／是因為一個孩子／瞠目結舌,顫驚呆立／耶穌卻呼喚他

我歌唱／是因為一個單身漢／禁閉在痲瘋的孤寂中／耶穌摸他

我歌唱／是因為一滴淚／無法理解的／從天父眼中滴落世界的一滴眼淚

(改寫自:谷川俊太郎的《我歌唱的理由》)

但願:榮耀歸於聖父,聖子,聖靈.阿們.



\newpage

\section{講定唔講做得到先好講}
\label{sec:441}
\textbf{林豪恩先生 - 講定唔講?做得到先好講}
\newline
\newline
\hyperref[sec:440]{< < < PREV SERMON < < <}
~
\hyperref[sec:toc]{[返目錄]}
~
\hyperref[sec:442]{> > > NEXT SERMON > > >}


講題:講定唔講?做得到先好講 Over-Promise, Under-Deliver

講員:林豪恩先生

場合:香港中文大學崇基學院禮拜堂

日期:2019年1月27日

一．引言

「講定唔講」?參加選舉的候選人一定講,而且通常「有咁大講到咁大」,力求爭取選票.選舉時開出的「支票」有多少能夠兌現呢?大家心中有數.有民間調查顯示,高達83\%的美國選民覺得總統未能兌現諾言.有位鄰近地區的總統幕僚發表文章為被批抨「走數」的總統辯解,題目為《競選承諾不一定都要兌現》,文中反問:「世界上有那位當選人競選承諾都兌現了的,不能全兌現是常態.」當選了,總不能沒有政績,於是市面上出現了各式各樣有趣的橫額,例如旺角街頭曾出現「成功爭取綠燈延長兩秒 保障街坊過路安全」的橫額,又有「成功爭取增設汽水機」,「成功爭取免費數碼電視頻道直播4場世界盃賽事」的橫額.

有些管理人員則選擇相反的方向,「唔講咁大」,以免讓受眾產生過高的期望,當結果公佈後受眾會發現比原來想像的好,受眾自然感到喜出望外.例如薪酬調整,很多時公司會預先發放經營困難的消息,讓員工對加薪幅度不要抱太高期望,但當公佈的加薪幅度比預期高的時候,員工都會笑逐顏開.

教會呢?有些教會認為一定要make noise,否則怎能為上帝作見證呢?怎能突顯正邪的分別呢?在長洲太平清醮活動期間,每年都舉辦飄色巡遊,廣為傳媒報導.某年,某教會在太平清醮節期組織會眾巡遊,身穿黃色 T恤,手持木製盾牌及長槍,高喊「耶穌在長洲」及「耶穌得勝」等口號,似乎想與太平清醮活動分庭抗禮,雙方一度發生衝突.

事件中的教會可能認為,教會有責任傳揚上帝,為上帝作見證.如此的認知應該是沒錯的,只是這樣的行動是否傳揚上帝,是否正在為上帝作見證呢?今天我們誦讀的詩篇,啟發我們思考以另外一種傳揚上帝,為上帝作見證的境界.這是一種怎樣的境界呢?

二．內容

1．無言的頌讚

19:1諸天述說上帝的榮耀；穹蒼傳揚他的手段.

19:2這日到那日發出言語；這夜到那夜傳出知識.

第二節「這日到那日發出言語」的是誰呢?是「諸天」；「這夜到那夜傳出知識」是誰呢?是「穹蒼」.第一及二句是平行句,因此,「諸天」與「穹蒼」同義,而「上帝的榮耀」也與「他的手段」(和修譯:他手的作為)也同義.「這日到那日」與「這夜到那夜」合理來,即是日連夜,二十四小時不停.不停做甚麼呢?不停地「發出言語」及「傳出知識」.試想像,二十四小時每日每夜不停地講,真是嘈吵得不得了,煩厭得不得了.這就是現代社會,二十四小時新聞,二十四小時廣播,無論是在家中,在街上,在商場中,在交通工具上……叫人避無可避.可惜的是,有些地方的宗教活動及宣傳也唯恐落後,有樣學樣,二十四小時唱詩歌敬拜,馬拉松奮興會等,各式各樣聚會,馬不停蹄.

19:3無言無語,也無聲音可聽.

幸好,「諸天述說上帝的榮耀」是「無言無語」的；「穹蒼傳揚他的手段」是「無聲音可聽」的.那麼,諸天如何每日不停地發出言語呢?穹蒼如何每夜不停地傳出知識呢?它們發出甚麼言語呢?他們傳出甚麼知識呢?原來,諸天的出現,穹蒼的存在,默默地按著自然定律日月推移,靜靜地按著所定軌跡四時運轉,這就是述說上帝榮耀的言語,就是傳揚他手的作為的知識.

有人譏諷基督徒,「講就天下無敵,做就有心無力」.有時,「講多無謂」,事情做好了,就已經是好見證.創世記記載了上帝創造人類的心意:

1:27上帝就照著他的形像創造人,照著上帝的形像創造他們；他創造了他們,有男有女.

1:28上帝賜福給他們,上帝對他們說:「要生養眾多,遍滿這地,治理它；要管理海裡的魚,天空的鳥和地上各樣活動的生物.」

如果我們能夠協助人類活出上帝的形像,並且好好管理大地,上帝的榮耀就被彰顯了.

曾經有電視廣告說:「月黑風高殺蟲夜,黑旋風忍者為民除害,唔使用劍嘎」.就算在月黑風高的危險處境中,也不一定要用劍的,「無言無語也無聲音可聽」的穹蒼,只要照常運行,就已經在傳揚上帝.

陳健民教授的最後一課的最後一句說話是:「在最黑的環境,我們才見到星星」.見到星星又如何呢?他說:「希望大家不要放棄」.星星繼續肯「開夜班」,就已經帶給人們希望.使徒保羅說:「如今常存的有信,有望,有愛」(林前13:13),盼望正是通往上帝的其中一條主要道路.

讓我們岔開一下話題,最近讀到一篇網上節文,也談及「講定唔講」這話題.文章名為「好媽媽就是三個字」,第一個字 是「靜」:

「很多媽媽眼尖,一眼能看出孩子三個問題,而且立刻指出,還反覆強化,甚至當著更多人的面說孩子就是哪裡不好.母親看到自己孩子的優點和缺點的時候都不要動,不要慌亂,為什麼?因為孩子需要母親用良好的情緒去孕育和保護,只有母親看到孩子的優,缺點都不露聲色的時候,才能賦予孩子內心的成長動力.母親情緒不穩定,整個家庭就跟著翻騰,孩子是活在母親這條船裡的,一點點動靜都會影響到孩子,破壞孩子的內在穩定狀態,孩子單單平息自己的內心緊張情緒就要消耗很多心理能量,還能留下多少心理能量用在學習,用於自我成長呢?」

「唔講」,也不等於放棄；「唔講」,不等於死寂.「無言無語」,其實可以是很有活力,很有作為,讓我們繼續讀詩篇19篇:

19:4它的量帶通遍天下,它的言語傳到地極.上帝在其間為太陽安設帳幕；

19:5太陽如同新郎出洞房,又如勇士歡然奔路.

19:6它從天這邊出來,繞到天那邊,沒有一物被隱藏不得它的熱氣.

「無言無語」,但「它的言語」卻能夠「傳到地極」；「無聲音可聽」,但「它的量帶」卻能夠「通遍天下」.它光亮熱誠,動力十足,「如同新郎出洞房」,「又如勇士歡然奔路」.它如何「述說上帝的榮耀」呢?它如何「傳揚他的手段」呢?它不用擴音器,而是用它的忠誠和穩定,每日「從天這邊出來,繞到天那邊」；又用的熱氣,把溫暖放射到無一個角落.它的量帶就是它的言語,它的溫暖所到之處,就是它的言語所及之地.

剛才提及「好媽媽就是三個字」這篇短文的另一個字是「慈」:

「在孩子面前指手劃腳,不是扯著嗓門叫喊就是繃著臉嘮叨的母親,會讓孩子慢慢感覺這個世界沒有任何一個地方是舒適的,是有意思的,孩子會失去心中的方向,而母親也將沒有任何引導能力,變得更加抓狂.一個母親真正的力量,是她的慈祥,讓人感到大海般的包容接納,這就是孩子成長的寬闊的天地.」

教會和基督徒也可以參考,未必需要與世界「鬥大聲」,現任天主教教宗勉勵教會及信徒:「慈愛就是見證」.

2．祂講而不是我講

諸天及穹蒼不用言語,人又如何呢?人當然需要言語,正如電影「大話西遊」的其中一段經典對白說:「悟空,你想要呀?你想要就要講啦!做咩唔出聲je ?你望住我做乜je ?雖然你咁有誠意噉望住我,我係好高興,你都要講架.」「講定唔講?」詩篇19篇的作者認為,話總是要講,「唔講點知呢?」重要的是,誰講?講甚麼?

19:7耶和華的律法全備,能甦醒人心；耶和華的法度確定,能使愚人有智慧.

19:8耶和華的訓詞正直,能快活人的心；耶和華的命令清潔,能明亮人的眼目.

19:9耶和華的道理潔淨,存到永遠；耶和華的典章真實,全然公義－

19:10都比金子可羨慕,且比極多的精金可羨慕；比蜜甘甜,且比蜂房下滴的蜜甘甜.

詩人指出,耶和華所講的是:全備,確定,正直,清潔,潔淨和真實.詩人又指出,耶和華所講的能夠:甦醒人心,使愚人有智慧,快活人的心,明亮人的眼目.誰所講的可以「存到永遠」呢?,誰所講的是「全然公義」呢?誰所講的「比金子可羨慕」,「且比蜂房下滴的蜜甘甜」呢?這不可能是人間的言語,一定是從上而來的說話.上主發言了,人應該如何回應呢?

19:11況且你的僕人因此受警戒,守著這些便有大賞.

19:12誰能知道自己的錯失呢?願你赦免我隱而未現的過錯.

19:13求你攔阻僕人不犯任意妄為的罪,不容這罪轄制我,我便完全,免犯大罪.

19:14耶和華－我的磐石,我的救贖主啊,願我口中的言語,心裡的意念在你面前蒙悅納.

正如以利對撒母耳說:「你仍去睡吧；若再呼喚你,你就說:耶和華啊,請說,僕人敬聽!」(撒上3:9)「僕人敬聽」,應該是人面對上主說話的最適合回應.詩人指出,耶和華的說話對我們有兩層重要幫助,其一是過去的,幫助我們省察自己的「錯失」,尋求上主寬恕；其二是將來的,「攔阻僕人不犯任意妄為的罪」.在上主的話語面前,我們只能謙卑地說:「願我口中的言語,心裡的意念在你面前蒙悅納」.

三．總結:「靜默有時,言語有時」

處身於現代社會中,四方八面都是聲音,畫面,文字等各式各樣的資訊,有時,我們都禁不住言語.陳奕迅其中一首流行曲「浮誇」把這種狀況表達淋漓盡致:

有人問我 我就會講 但是無人來

我期待到無奈 有話要講 得不到裝載

我的心情猶像樽蓋 等被揭開 咀巴卻在養青苔

人潮內 愈文靜 愈變得不受理睬 自己要搞出意外

像突然地高歌 任何地方也像開四面台

著最閃的衫 扮十分感慨 有人來拍照要記住插袋

你當我是浮誇吧 誇張只因我很怕

似木頭 似石頭的話 得到注意嗎

其實怕被忘記 至放大來演吧

很不安 怎去優雅

世上還讚頌沉默嗎

不夠爆炸 怎麼有話題 讓我誇 做大娛樂家

為甚麼要浮誇呢?因為怕.怕什麼呢,怕無人理解,被人嘲笑,努力得不到認同.這不是想像,而是經驗:

那年十八 母校舞會 站著如嘍囉

那時候 我含淚發誓各位必須看到我

在世間 平凡又普通的路太多 屋村你住哪一座

情愛中 工作中 受過的忽視太多 自尊已飽經跌墮

為了不被忽視,為了被看見,為了脫離普通,為了自尊,歌者認為有多浮誇就要多浮誇:

你叫我做浮誇吧 加幾聲噓聲也不怕

我在場 有悶場的話 表演你看嗎 夠歇斯底里嗎

以眼淚淋花吧 一心只想你驚訝 我舊時似未存在嗎

加重注碼 青筋也現形 話我知 現在存在嗎

當我們在K房聲嘶力竭地唱著這首歌的時候,多多少少感受到自己心裡也有這些被忽視的懼怕和被重視的追求.可惜的是,就算青筋暴現,很多時未必能夠驅散懼怕；大喊大叫,通常都不會帶來優雅.傳道書說:「靜默有時,言語有時」(3章7節).現代社會以為只有言語能展示力量,然而,在靈性的境界,靜默的力量卻不可小觀.

2018年6月23日,泰國清萊府的少年足球隊「野豬隊」(Moo Pa)12名少年和教練,進入睡美人洞探險,未料當天卻適逢泰國雨季,一行人因暴雨灌入洞內,逃到高處避難,被困於洞中.為尋找失蹤的13人,在兩週的時間裡,來自世界各地的數百名志願者投入救援,其中不乏軍警人員,潛水專家,醫療人員,工程公司,救難團隊等,來到救援工作中協助或提供援助.7月10日,受困第18天,全員救出.

被困在洞穴內這18天,這些十多歲的男孩子及25歲的教練是怎樣渡過的呢?在被困洞穴的第八日,有拯救專家接受訪問時說:「如果他們仍然生存,目前面對的最大考驗是心理問題,因為他們至今都不知道會否獲救.」他們如何面對驚慌,焦慮,不安,絕望的情緒呢?團隊之間出現互相指責,埋怨,甚至衝突嗎?

第9天,兩名英國潛水員在洞穴內發現這13人還活著.他們沒有哭泣,只是靜靜地坐在黑暗中打坐冥想.一個受困男孩的母親透過視像看到洞穴內的情況,告訴美聯社說:「看看他們有多麼平靜,坐在那裡等待.沒人在哭或什麼的.這是令人驚訝的.」參與救援的潛水人員卡拉迪奇(Ivan Karadzic)對BBC表示:「我實在不明白這些孩子怎麼能這麼冷靜,他們受困洞穴內狹小空間長達兩周,而且媽媽不在身邊.不可思議的堅強孩子,簡直令人難以置信.」

講述這個事件,並不是要推薦打坐冥想,只是希望大家留意靜默的力量.在靈性的世界,浮誇是無補於事的,靜默,反而能帶來當下所需要的安定和存在感.

泰國清萊府的少年足球隊事件不禁令我想到,如果這事件發生在基督教會的少年團契,我們的十多歲男孩在洞穴內如何自處呢?我們的二十五歲導師能夠帶領少年人做甚麼事情呢?或者,基督教群體需要重尋失落多時的靈性操練,而不單單是宗教知識傳遞.



\newpage

\section{拆毀與建立}
\label{sec:442}
\textbf{戴浩輝牧師 - 拆毀與建立}
\newline
\newline
\hyperref[sec:441]{< < < PREV SERMON < < <}
~
\hyperref[sec:toc]{[返目錄]}
~
\hyperref[sec:443]{> > > NEXT SERMON > > >}


路加福音 4:21-30
\newline
\begin{longtable}{cl}
\hline
\hline
章節 & 經文 (和合本修訂版)\\
\hline
4:21 & \begin{tabularx}{0.7\textwidth}{X} 耶穌對他們說：「你們聽見的這段經文，今天已經應驗了。」 \end{tabularx} \\ \\ \relax
4:22 & \begin{tabularx}{0.7\textwidth}{X} 眾人都稱讚他，並對他口中所出的恩言感到驚訝；他們說：「這不是約瑟的兒子嗎？」 \end{tabularx} \\ \\ \relax
4:23 & \begin{tabularx}{0.7\textwidth}{X} 耶穌對他們說：「你們一定會用這俗語向我說：『醫生，你醫治自己吧！我們聽見你在迦百農所做的事，也該在你自己的家鄉做吧。』」 \end{tabularx} \\ \\ \relax
4:24 & \begin{tabularx}{0.7\textwidth}{X} 他又說：「我實在告訴你們，沒有先知在自己家鄉被人接納的。 \end{tabularx} \\ \\ \relax
4:25 & \begin{tabularx}{0.7\textwidth}{X} 我對你們說實話，在以利亞的時候，天閉塞了三年六個月，遍地有大饑荒，那時，以色列中有許多寡婦， \end{tabularx} \\ \\ \relax
4:26 & \begin{tabularx}{0.7\textwidth}{X} 以利亞並沒有奉差往她們中任何一個人那裡去，只奉差往西頓的撒勒法一個寡婦那裡去。 \end{tabularx} \\ \\ \relax
4:27 & \begin{tabularx}{0.7\textwidth}{X} 在以利沙先知的時候，以色列中有許多痲瘋病人，但除了敘利亞的乃縵，沒有一個得潔淨的。」 \end{tabularx} \\ \\ \relax
4:28 & \begin{tabularx}{0.7\textwidth}{X} 會堂裡的人聽見這些話，都怒氣填胸， \end{tabularx} \\ \\ \relax
4:29 & \begin{tabularx}{0.7\textwidth}{X} 就起來趕他出城。他們的城造在山上；他們帶他到山崖，要把他推下去。 \end{tabularx} \\ \\ \relax
4:30 & \begin{tabularx}{0.7\textwidth}{X} 他卻從他們中間穿過去，走了。 \end{tabularx} \\ \\
[1ex]
\hline
\hline
\end{longtable}
講題:拆毀與建立  To Tear Down and to Build

講員:戴浩輝牧師

經文:路加福音4:21-30

場合:香港中文大學崇基學院禮拜堂

日期:2019年2月3日

引言

凡是挑戰我們所珍惜的假象,又或把我們從我們的安穩地帶踢出來的人,我們不歡迎甚至鄙視他們.但真理卻是:當耶穌要準備拯救我們,他先要對付我們的假象,或所有與基督的靈相反的偏見.比利‧桑德(Billy Sunday)可謂是葛培理之前一代的佈道家,是一位職業的壘球員,他在某一個城傳福音,在他一次講道中,他批評當地的僱主給予工人的工作條件很差；在聚會完結後,有幾位當地的老板藉當地的牧師給他一些忠告:比利,不要理會工人的問題,只要集中拯救靈魂,不要涉及政治問題；因為你這樣做等同「反掃貓毛」.比利‧桑德在回應的紙條寫著:如果我是「反掃貓毛」的話,叫貓兒轉過身來吧!

今天的經文描述,耶穌在自己長大的城拿撒勒傳道,同樣面對反對,耶穌的話就好像比利‧桑德的一樣,要破除人很多的假象.讓我們從這段經文中去探索,在我們當中是否也有若干假象要被拆毀,好讓耶穌基督在我們生命中,重新建立合乎上帝心意的生命.

一. 拆毀假象,建立真理

路四21-30是承接上主日的經課,耶穌在自己家鄉宣講,第21節是這宣講的撮要,而第22節是各人的反應,「眾人都稱讚他,並希奇他口中所的恩言；」但為何到第28節指出「他們都怒氣滿胸?」究竟「這不是約瑟的兒子」是甚麼意思?為何耶穌會以一句俗語回答他們?

如果我們把22和23節聯在一起來思想,我們便得出一個結果；「這不是約瑟的兒子」和「行在你自己家鄉裡」正是耶穌要拆毀他們假象的焦點.你既然這樣了得,怎麼「肥水流往別人田」?這裡並沒有否定耶穌能力和祂所宣講的道,但這裡要反映的,是耶穌沒有像在迦百農所行的好處在「自己人」的地方也行出來.

很多時我們都可能有拿撒勒人的假象,我們請同工工作,當然是要服事弟兄姊妹,某某教師或牧師很久也沒有問候我,好像並不關心我,自然很可能也不會關心其他弟兄姊妹,那麼他就是偷懶,失職.我也曾作為監督,當時也偶爾聽到投訴同工的罪狀.很多時投訴者假設教牧應當如何,要如何為他們謀幸福.而最多投訴的不是返教會不久的初信者,而是信了很久的所謂老臣子.當教牧並不按他們的程序辦事,他們便認為教牧不濟.教會中有不少人有很多這樣的假象,如果不被拆毀,很難建立起教會的真相.

還記得一九八九年在德國慕尼黑的學生宿舍生活,有一位曾榮輝牧師由加拿大信義會來到我那裡,說要在南德國宣教牧養；當時,他講了一句話,給我很大的挑戰,我來不是要牧養你,而是要你去牧養人.當時我是學生,但他說明他沒有時間在慕尼黑,因為他的基地在二百公里外的史圖加特.那樣,我就作了義務牧師,成立了慕尼黑堂,牧養該堂,直到我返港工作.我要拆毀我自己的假象,有牧師到來,就是牧養我這個留德的學生.教會內很多衝突是因為我們堅持我們有權,而沒有理會,我們的權是否站得住腳.現今世界的衝突同樣就是強權而不需要真理,這種強權式所建立的一切,最終要被拆毀,因為都是假象.

二 拆毀定型,建立信心

耶穌還未為到自己家鄉作工,鄉里便認定他要做這做那,他們為耶穌假設了很多,他們似乎各人對耶穌的認識有不同的焦點,而耶穌又不能滿足他們.在今天,我相信每一位弟兄姊妹對禮拜堂和教牧都有不同的期望,但這不一定能滿足的.我們每一個人都可能活在其他人對自己的要求下,甚至我們有時的不濟被視為我們的宿命,我們被「睇死」,被定型!

可口可樂行政總裁加奧(Donald R. Keough)在退休之前向艾瑪利(Emory)大學的畢業生演講,面對著一群將要面對艱難和批判的世界的青年男女,他說:我有一個朋友是建築師,他曾說:我可以拿

一座由最好的建築師所建的最新的樓,然後你給我有不同距離鏡頭的相機,我只要對焦在幾個不完美點上,我便可以把這座樓拍成將要倒下來的危樓.在這個社會,有些人也用生活的相機對準我們,如果你容許他們用他們的「相機」對焦你的人生,我們因著這些「照片」便會感到失望,為我們的生命而感到害怕,不知所措.所以我們要警覺,有人會這樣向你對焦,看你的短處和不足,他們向耶穌這樣做,也會向你這樣做.因此我們靠著聖靈,從主耶穌身上學習,拆毀人家對自己的定型,建立信心.很多人都聽過林肯總統的故事,他曾經參加總統的競選七次都落選,最後才被選上當總統；他可謂抗逆力極高,永不言敗的人.荷蘭畫家凡高一生從繪畫所賺取的,換算港幣只得680元.一百年之後,他一幅加撒特醫生(Dr. Gachet)便值六億六千萬港幣.同樣,兩千多年前,在中東地方一位傳道人被自己家鄉的人和當時宗教領袖所排斥,但祂永不言敗,把福音藉祂的生命傳開.如果林肯和凡高讓反對和逆境掌管他們的生命,使他們不能釋放地講出或做出他們認為對的事,我們

這個世界將會十分貧乏.同樣,如果拿撒勒人耶穌讓反對和逆境掌管他們的生命,使他們不能釋放福音的真義,我們這個世界還是沉淪,沒有得救的途徑.今天不論教牧,我們崇基禮拜堂的弟兄姊妹都同樣藉聖靈的能力,拆毀定型或宿命,建立信心.

三 拆毀界限,建立仁愛

耶穌繼續教訓他們,藉以利亞和以利沙兩位北國以色列先知的事蹟來表明祂的事工.我相信大家對這兩個故事都不陌生,在以利亞的故事中,以利亞要逃避飢荒和朝廷的迫害,最後到了西頓寡婦的家中,寡婦的慷慨使他為這個家帶來平安.而以利沙的故事敘述當時最強大的敘利亞,有一位將軍,有權有勢,但長了大痲瘋；然而,他相信一位擄獲作婢女的以色列小女孩的話,來到以利沙的面前尋

求醫治.耶穌把這兩個故事說出一方面是要拆毀拿撒勒同鄉對上帝國臨在的錯誤想像,就好像以利亞和以利沙時代的以色列那樣.他們生活在虛假的安全感中,以為他們是上帝的子民,無論怎樣,上帝只會幫助他們,上帝只愛他們.很明顯地,耶穌的鄉里甚至比對迦百農,都抱怨耶穌在那裡所作的,而不做在他們那裡.耶穌的宣講一方面不讓自己的鄉里操控(domesticated),另一方面,就是要把他們這些界限拆毀.同樣,今天我們當中可能還有這些界限,例如:教會只為會友而存在的話,這樣的教會就是一所俱樂部；相反的,教會是為所有的人而設的話,教會就成為救生站.這救生站縱然有人來到而沒有相信耶穌,仍得到教會關懷幫助；這樣,便真的是耶穌的教會.當然,在香港也有教會以關懷服侍為工具,傳福音才是目的.我雖是一個福音派基督徒,但我仍相信基督徒服侍見證是有其內在的價值的.當然這並不表達我們排除以口傳揚福音,因為惟有藉服侍每一個有需要的人,這教會或基督徒才稱得上是屬基督的,因為主就是愛.耶穌以這兩位先知所接觸的對象給我們一個天國的對象.外邦的窮寡婦,在以色列人的眼中可謂卑微得無可比的了.敘利亞的大將乃縵,可謂權位都有.上帝都愛他們,這兩個故事是要建立上帝對世人大愛的信息,上帝愛每一個人,不因為他們的身分地位和性別.讓我們靠著聖靈所賜的能力,活出這樣的愛,拆毀我們心靈的界限,實踐仁愛.

結語

耶穌來到自己的地方,自己的人卻不真認識祂,因為他們有太多的假象和界限；耶穌也不因他們對自己批評而成為自己的宿命和定型；祂把這一切都拆毀.祂來是要傳揚福音真理,因為祂是真理.祂來是為要建立我們每人的信心和愛心,信靠祂的人是有福的,阿們.



\newpage

\section{你們甚麼都不懂我們懂得甚麼}
\label{sec:443}
\textbf{劉國偉先生 - 「你們甚麼都不懂!」……我們懂得甚麼?}
\newline
\newline
\hyperref[sec:442]{< < < PREV SERMON < < <}
~
\hyperref[sec:toc]{[返目錄]}
~
\hyperref[sec:444]{> > > NEXT SERMON > > >}


講題:「你們甚麼都不懂!」……我們懂得甚麼? “You Know Nothing At All!”……What We Do Know About?

經文:約翰福音11:45-57,12:1-11

講員:劉國偉先生

場合:香港中文大學崇基學院禮拜堂

日期:2019年2月10日

引子

麥哥登Gordon James MacDonald, 美國著名牧者,作家,神學工作者,亦是美國前總統克林頓陷入靈性及政治危機的時候,其中一位主要成員,為這位前總統進行輔導與心靈重建的牧者.

1984年,麥哥登一本十分成功的著作Ordering Your Private World面世 (中譯「心意更新」在香港出版),使他攀上高峰,成為最暢銷書作者.1987年,一封寄往出版社的匿名信揭發了麥哥登的婚外情,麥哥登勇敢地承認他的過犯,他羞辱了主的名,傷害了家庭,孩子,破壞了牧職和教會對他的信任,毀掉了自己的人格,使他如日中天的聲望和事業,在一夜之間完全拆毀.

麥哥登牧師走上生命重整的道路,而他另一本著作Rebuilding Your Broken World面世(中譯「從破碎到重建」在香港出版),他以過來人的身份,談論到人為甚麼會失德.麥哥登指出人無論在情慾,道德,金錢,權力層面失德,其實都指向一個淵深的自我(deeper person),失德之罪是漸次累積造成的,特別當權力愈來愈大的時候,能否駕馭內心深處黑暗的自我,就成為特別是領袖人物的嚴峻的挑戰.

麥哥登在”Ordering Your Private World”談論到被慾望驅使的人A Driven Person : They can become so pre-occupied with success and achievement that they have little time to stop and ask if their inner person is keeping pace with the outer process. Usually it is not, and there is an increasing gap, a breakdown in integrity. People like this often become progressively deceitful; and they not only deceive others, they deceive themselves…they lie to themselves about motives; values and morals are compromised. Shortcuts to success becomes a way of life.

他們為了取得成就,幾乎沒有時間停下來尋問他們的內心,是否跟外在的行為同步.通常不是,並且存在越來越大的差距,直至誠信的崩潰.像這樣的人,往往會逐漸慣於欺騙; 他們不僅欺騙別人,還欺騙自己……他們對自己的動機撒謊; 在價值和道德上妥協.以捷徑取得成功,就成為一種生活方式.

麥哥登牧師曾經活在欺騙別人也自我欺騙的黑暗中,對於那些巧言辭令,去掩藏內心的真正動機的人,有深刻的第一身經歷；在福音書的記載中,也出現了這一種被慾望驅使,為了掩藏內心的真正動機的人,他們更搬弄出大是大非的理由.

玩耍人命的政客

今天講道的經文選取了約翰福音11章45節到12章11節.整個大段落可以分作三個小段落 : 11章45至57節記述了猶太公會商議要殺害耶穌的前因後果; 12章1至8節穿插了馬利亞在伯大尼的筵席,以香膏膏耶穌的事跡; 然後在12章9至11節記述了祭司長連拉撒路也要殺死的構想.

「那些來看馬利亞的猶太人見了耶穌所做的事,就多有信他的；但其中也有去見法利賽人的,將耶穌所做的事告訴他們.」(約11:45-46)

當拉撒路復活的消息傳至法利賽人,再傳至祭司長,就引起了整個猶太公會的高度關注,而公會的討論議題確實提升到相當嚴重的層次:

「這人行好些神蹟,我們怎麼辦呢?若這樣由著他,人人都要信他,羅馬人也要來奪我們的地土和我們的百姓.」 (約11:47-48)

拉撒路從死裡復活,引發更多猶太人相信耶穌,無疑是整件事件的轉捩點.對大祭司及公會而言,與其說是他們要與耶穌在宗教地位爭一日之長短,不如明明白白的說是一次權力捍衛,他們對愈來愈多猶太人投向耶穌,表示不安,更進一步擔心耶穌的影響力再蔓延開去,則他們在民眾中間的影響力便相對日漸萎縮,而且,萬一真的釀成亂局,便大大催化羅馬政權介入巴勒斯坦事務,倘真如此,他們所擁有的一切權力,地位,物質享受便可能化為烏有.

在那時候,祭司都由撒都該人出任.撒都該人享受著社會所賦予的上等階層的身份,他們有財有勢,政治上親羅馬政權,在猶太人中享有相當政治影響力,這一切既得利益,叫他們有強烈的維持現狀的動機,美其名是穩定壓倒一切,沒有動亂便能保障猶太人免於羅馬鐵蹄的洗禮,其實骨子裡卻是一次精密的政治計算.

耶穌遇害那年的大祭司名叫該亞法,他從公元18年至36年擔任此職,他口所出的那句名言:「你們不知道甚麼. 獨不想一個人替百姓死,免得通國滅亡,就是你們的益處.」現代中文譯本如此翻譯:『你們甚麼都不懂!讓一個人替全民死,免得整個民族被消滅.難道看不出這對你們是一件合算的事嗎?』(約11:49-50)

該亞法這句話充份暴露他的政客的本相,以玩耍人命來展現他的「英明領導」,用耶穌一命換取整個民族的利益,既是精打細算,更加「表明」他以民眾的集體利益為決策的基礎.然而,身為大祭司,要判斷一個人竟不是基於他所作的事是否正直,要把一個人鏟除也不是基於對耶和華一神信仰的捍衛,他的目的只是要維持政治現狀,這已成為他的最高信仰原則.如此看來,該亞法這句話實在把公會中人商議要殺耶穌的動機,美化得太過離譜.實在的計算是只要耶穌一死,便除去心腹大患,大祭司和公會的宗教領導地位可保,穩定的政局可以避免羅馬政權的干涉,也就讓猶太貴胄得以續享繁榮安定.於是,「從那日起,他們就商議要殺耶穌.」(約11:53) 猶太人的領袖不是單單討論一下事情,更加落實計劃如何殺害耶穌.

『你們甚麼都不懂!』

「有許多猶太人知道耶穌在那裡,就來了,不但是為耶穌的緣故,也是要看他從死裡所復活的拉撒路.但祭司長商議連拉撒路也要殺了；因有好些猶太人為拉撒路的緣故,回去信了耶穌.」(約12:9-11)

耶穌固然要除掉,但從墳墓中走出來的拉撒路,也是個叫他們頭痛的人物.既然已向耶穌動了殺機,要把拉撒路也一併除掉,並非一個叫人感到意外的想法!只是,這個構想竟然出於宗教領袖,他們不單要耶穌消失於人世,更立心要毀滅拉撒路這活生生的神蹟,確實顯得狠毒非常的殺人兼毀滅證據的惡行!何竟拉撒路再次與死亡走近,竟因為他曾經真的死過,他的死而復生的獨特經歷,竟成為他再次面向死亡的催化劑!

很多做領袖的都過不了覺得跟從者威脅自己的考驗,要不把他鏟除,或者設法打擊他,讓他不可能成為自己的競爭對手；能謙虛躬身自省,尋求更新變化之道的,畢竟少數.祭司長及公會計劃連拉撒路也要殺掉,正是因為許多猶太人為拉撒路的緣故離開他們(約12:10,參現代中文譯本),於是擺弄政治手段,愚民也愚己,把耶穌與拉撒路一起根除,竭盡所能地阻止離心他投的猶太人去相信耶穌.

歷史中浮過不少掌權高官,政客謀士,以政治手腕,權術賣弄,假意談公義,口口聲聲為人民之福祉,卻到頭來只為自己及自己的利益集團效命；在人間市井之中,也屢見以政治交易尋求個人利益,以此作為生活最高目標之市井中人.好些人將玩弄政治成為生命信仰所在,那只會背向真理而且愈走愈遠.該亞法定睛於政治利益,排斥異己,玩耍人命,褻瀆了他所擁有的神聖位份,這成為一切在上位者,在公共空間,職場,教會和團體中擁有行政權力的人,一個沉重的提醒.

該亞法說 :『你們甚麼都不懂!』正告訴人們他並非普通人,他比別人懂得更多,這一種自視高人一等的感覺,並以捷徑來達成自己的慾望,那個深淵的自我,催使他成為一個被慾念驅使的人,A Driven Person.他不僅欺騙別人,還欺騙自己……他對自己的動機撒謊; 在價值和道德上妥協.然後,他認為別人『甚麼都不懂!』

我們懂得甚麼

穿插在兩小段落之間的是馬利亞膏耶穌的片段.

由於時候未到,耶穌到靠近曠野的一個小城暫避風頭 :

所以,耶穌不再顯然行在猶太人中間,就離開那裡往靠近曠野的地方去,到了一座城,名叫以法蓮,就在那裡和門徒同住. (約11:54)

逾越節近了,猶太人在討論耶穌會否來耶路撒冷過節,公會的人也佈置線眼,作好捉拿耶穌的準備 :

猶太人的逾越節近了,有許多人從鄉下上耶路撒冷去,要在節前潔淨自己.他們就尋找耶穌,站在殿裡彼此說:「你們的意思如何,他不來過節嗎?」那時,祭司長和法利賽人早已吩咐說,若有人知道耶穌在哪裡,就要報明,好去拿他.(約11:55-57)

就在這暴風雨來臨的前夕,耶穌再訪伯大尼這引起是非爭議的地方,耶穌和拉撒路兩個爭議人物一同坐席 :

逾越節前六日,耶穌來到伯大尼,就是他叫拉撒路從死裡復活之處.有人在那裡給耶穌預備筵席；馬大伺候,拉撒路也在那同耶穌坐席的人中.(約12:1-2)

筵席之間,馬利亞拿著一斤極貴的真哪噠香膏,抹耶穌的腳,又用自己頭髮去擦,老約翰清楚記得哪噠香膏使屋裡滿了膏的香氣.

馬利亞就拿著一斤極貴的真哪噠香膏,抹耶穌的腳,又用自己頭髮去擦,屋裡就滿了膏的香氣.(約12:3)

哪噠香膏油是由一種黎巴嫩的植物Nardostachys提煉而成.和合本翻譯的「真」原文包含有珍貴及純度的含義.馬利亞所獻上的哪噠香膏油,表達了她對耶穌那份委身的愛,那是真摯的,珍貴的,高純度的,及充滿香氣的!

按當時的情景,耶穌很可能身處一間古羅馬的躺臥餐廳,飲宴的賓客以半倚的坐姿進食,雙腳就平放在長椅上,服侍的人就在賓客的中間奉上食物和飲品,馬利亞就在外圍膏耶穌的腳.以當時的文化,女人的長頭髮是她的榮耀(林前11:15),馬利亞以自己的頭髮,謙卑地去擦耶穌的腳,相對於膏抹由頭部開始的習俗,馬利亞亦表達了自己的不配,那是一種很深情的表達,傾盡最真摯的最珍貴的愛,呈獻給耶穌.

猶大並不同意,他認為拿一個普通工人一年的工資去膏一雙腳十分浪費……我們再一次看見一個巧言辭令的人,以一個堂而皇之的道理,去掩藏內心的真正動機和慾念,老約翰作為目擊者,他毫不客氣指證猶大口是心非.

有一個門徒,就是那將要賣耶穌的加略人猶大,說:「這香膏為甚麼不賣三十兩銀子賙濟窮人呢?」他說這話,並不是掛念窮人,乃因他是個賊,又帶著錢囊,常取其中所存的.(約12:4-6)

馬利亞深知耶穌迎向死亡,早就為耶穌的死預備了這珍貴的哪噠香膏,正如耶穌說「只是你們不常有我.」 馬利亞就把握了這個珍貴時刻.

耶穌說:「由她吧!她是為我安葬之日存留的.因為常有窮人和你們同在,只是你們不常有我.」 (約12:7-8)

親愛的弟兄姊妹,若果今天有一位該亞法,位高權重,有權有勢,他可能站於精英階層而自我感覺良好,認為自己有不一樣的視野和觀點,並且有不平凡的判斷力,然後向我們說『你們甚麼都不懂!』,你會如何回應?

從文學佈局的角度來看,老約翰把馬利亞膏抹耶穌的片段,穿插在玩耍人命的兩段敘事之間,那可能是他寫作過程的自然流暢的回憶,但這個交錯編排一起的安排,卻產生了強烈的對比效果.我認該亞法在故事裡是一個被慾望驅使的人,A Driven Person,而馬利亞卻是一個被愛所呼喚的人,A Called Person.

我相信馬利亞膏耶穌的事跡,是回應該亞法『你們甚麼都不懂!』的評語的最有力回應,那是關乎耶穌以捨命來成就救恩,關乎馬利亞所呈獻的委身的愛,是真摯的,珍貴的,高純度的,及叫四週的人感受到滿室香氣的愛!

弟兄姊妹,也許我們走過江湖,多少懂得或經歷過權力遊戲,被慾望驅使的人和事也不斷在眼前上演,但我們跟從耶穌,叫我們真正懂得甚麼? 那是基督為我們捨己,然後我們學習如何真誠地愛祂.當世人為自己的利益,攪盡腦汁作出種種計算,我們為耶穌獻上的真摯的,珍貴的,高純度的,及充滿香氣的愛,在我們的生活當中,又是怎麼一回事?

親愛的弟兄姊妹,在你的人生旅程中,若果你真的很想真摯地愛主耶穌,你可曾把握著生命的各種機遇,去回應耶穌基督那份捨己的愛,就像馬利亞為主安葬之日存留了一瓶真哪噠香膏? 我希望你可以坦誠地告訴自己,對你而言,到底你那一瓶極其貴重的哪噠香膏油是什麼?



\newpage

\section{不義而富且貴於我如浮雲}
\label{sec:444}
\textbf{鄧瑞強博士 - 不義而富且貴,於我如浮雲}
\newline
\newline
\hyperref[sec:443]{< < < PREV SERMON < < <}
~
\hyperref[sec:toc]{[返目錄]}
~
\hyperref[sec:445]{> > > NEXT SERMON > > >}


啟示錄 18:1-20
\newline
\begin{longtable}{cl}
\hline
\hline
章節 & 經文 (和合本修訂版)\\
\hline
18:1 & \begin{tabularx}{0.7\textwidth}{X} 此後，我看見另一位有大權柄的天使從天降下，地由於他的榮耀而發光。 \end{tabularx} \\ \\ \relax
18:2 & \begin{tabularx}{0.7\textwidth}{X} 他以強而有力的聲音喊著說：「傾覆了！大巴比倫傾覆了！她成了鬼魔的住處，各樣污穢之靈的巢穴，各樣污穢之鳥的窩，各樣污穢可憎之獸的出沒處。 \end{tabularx} \\ \\ \relax
18:3 & \begin{tabularx}{0.7\textwidth}{X} 因為列國都喝了她淫亂大怒的酒；地上的君王和她行淫；地上的商人因她極度奢華而發了財。」 \end{tabularx} \\ \\ \relax
18:4 & \begin{tabularx}{0.7\textwidth}{X} 我又聽見另一個聲音從天上說：「我的民哪，從那城出來吧！免得和她在罪上有份，受她所受的災殃； \end{tabularx} \\ \\ \relax
18:5 & \begin{tabularx}{0.7\textwidth}{X} 因她的罪惡滔天，神已經記得她的不義。 \end{tabularx} \\ \\ \relax
18:6 & \begin{tabularx}{0.7\textwidth}{X} 她怎樣待人，也要怎樣待她，按她所行的加倍地報應她；用她調酒的杯加倍調給她喝。 \end{tabularx} \\ \\ \relax
18:7 & \begin{tabularx}{0.7\textwidth}{X} 她怎樣榮耀自己，怎樣奢華，也要使她照樣痛苦悲哀。因她心裡說：『我坐了皇后的位，並不是寡婦，絕不至於悲哀。』 \end{tabularx} \\ \\ \relax
18:8 & \begin{tabularx}{0.7\textwidth}{X} 所以在一天之內，她的災殃要一齊來到，就是死亡、悲哀、饑荒。她將被火燒盡，因為審判她的主神大有能力。」 \end{tabularx} \\ \\ \relax
18:9 & \begin{tabularx}{0.7\textwidth}{X} 地上的君王，與她行淫、一同奢華的，看見燒她的煙，就必為她哭泣哀號； \end{tabularx} \\ \\ \relax
18:10 & \begin{tabularx}{0.7\textwidth}{X} 因怕她的痛苦，就遠遠地站著，說：「禍哉，禍哉，這大城！堅固的巴比倫城啊！一時之間，你的審判要來到了。」 \end{tabularx} \\ \\ \relax
18:11 & \begin{tabularx}{0.7\textwidth}{X} 地上的商人也都為她哭泣悲哀，因為沒有人再買他們的貨物了； \end{tabularx} \\ \\ \relax
18:12 & \begin{tabularx}{0.7\textwidth}{X} 這貨物就是金、銀、寶石、珍珠、細麻布、絲綢、紫色和朱紅色衣料、各樣香木、各樣象牙的器皿、各樣極寶貴的木頭和銅、鐵、大理石的器皿， \end{tabularx} \\ \\ \relax
18:13 & \begin{tabularx}{0.7\textwidth}{X} 和肉桂、豆蔻、香料、香膏、乳香、酒、油、細麵、麥子、牛、羊、馬、馬車，以及奴隸、人口。 \end{tabularx} \\ \\ \relax
18:14 & \begin{tabularx}{0.7\textwidth}{X} 「你所貪愛的果子離開了你；你一切的珍饈美味和華美的物件都從你那裡毀滅，絕對見不到了。」 \end{tabularx} \\ \\ \relax
18:15 & \begin{tabularx}{0.7\textwidth}{X} 販賣這些貨物、藉著她發財的商人，因怕她的痛苦，就遠遠地站著哭泣悲哀， \end{tabularx} \\ \\ \relax
18:16 & \begin{tabularx}{0.7\textwidth}{X} 說：「禍哉，禍哉，這大城！她穿著細麻、紫色、朱紅色的衣服，用金子、寶石、珍珠為妝飾。 \end{tabularx} \\ \\ \relax
18:17 & \begin{tabularx}{0.7\textwidth}{X} 一時之間，這麼多的財富就歸於無有了。」所有的船長和到處航海的，水手以及所有靠海為業的，都遠遠地站著， \end{tabularx} \\ \\ \relax
18:18 & \begin{tabularx}{0.7\textwidth}{X} 看見燒她的煙，就喊著說：「有哪一個城能跟這大城比呢？」 \end{tabularx} \\ \\ \relax
18:19 & \begin{tabularx}{0.7\textwidth}{X} 於是他們把灰塵撒在頭上，哭泣悲哀地喊著說：「禍哉，禍哉，這大城！凡有船在海中的，都因她的珍寶成了富足。她在一時之間就成為荒蕪。 \end{tabularx} \\ \\ \relax
18:20 & \begin{tabularx}{0.7\textwidth}{X} 天哪，眾聖徒、眾使徒、眾先知啊！你們都要因她歡喜，因為神已經在她身上為你們伸了冤。」 \end{tabularx} \\ \\
[1ex]
\hline
\hline
\end{longtable}
講題:不義而富且貴,於我如浮雲 Fallen is Babylon the Great!

經文:啟示錄18章1-20節

講員:鄧瑞強博士

場合:香港中文大學崇基學院禮拜堂

日期:2019年2月17日

各位弟兄姊妹,祝願你們身心平安.

在農曆新年期間,我們最常聽到的話是「恭喜發財」.要發財,其實是很專業的,業餘玩票,很難發財.要發財,你的生命目標要對準金錢,每刻盤算的是金錢,發的夢是金錢,交的朋友是金錢,緊張的是金錢,開心是因為金錢,憂心是因為金錢,夠膽以生命為籌碼去撈取金錢,如此專業,才能有機會發財.簡單而言,真的要發財,你要賣命給金錢.你要得到金錢,必先讓金錢得到你.正如一個人若要做科學家,他必須先讓科學的真與美吸引他,充滿他,擁有他.

一旦你盯上金錢,金錢也會盯上你.隔不多久,你便要看看股票巿場的各種指數.之後,你的情緒指數便與股票巿場的指數同步,巿場的指數起,你的情緒便起；它跌,你便跌.生命成了金錢世界裡的一隻棋子.

啟示錄以天上的觀點俯覽塵世,當然會看到我們這個金錢世界,啟示錄揭露這金錢世界的本相.啟示錄又以歷史終點的角度點評人間,它道出金錢世界的虛幻.

講道經文:啟示錄18:1－20(我採用邊讀邊解的方式)

啟18:1　此後,我看見另有一位大權柄的天使從天降下,地就因他的榮耀發光.

【從天而來的天使,代表著從天而來的觀點,他的觀點照亮不明朗的世情.】

啟18:2　他大聲喊著說:巴比倫大城傾倒了!傾倒了!成了鬼魔的住處和各樣污穢之靈的巢穴,並各樣污穢可憎之雀鳥的巢穴.

【「巴比倫大城」當然指「羅馬」,它代表金錢世界的極致,以暴力為支撐,以奴役人為手段,它最終傾倒了,成了鬼魔和鳥獸居住的荒涼之地.】

啟18:3　因為列國都被她邪淫大怒的酒傾倒了.地上的君王與她行淫；地上的客商因她奢華太過就發了財.

【金錢世界有其吸引力.「邪淫大怒的酒」令人傾倒.所謂「邪淫」,即極盡不道德.所謂「大怒」,即極高漲的情緒.「邪淫大怒」者,即不道德的狂歡,列國為這傾倒,人人都趨之若鶩.所謂「地上的君王」,即社會的管治階層,他們也以這種生活為目標.極盡狂歡的生活,需要很多奢侈品,售買奢侈品的商人當然高興.】

啟18:4　我又聽見從天上有聲音說:我的民哪,你們要從那城出來,免得與她一同有罪,受她所受的災殃；

啟18:5　因她的罪惡滔天；她的不義,神已經想起來了.

【經文對著我們說的,「你們…」.我們要有所行動,要離開,要分別為聖.信徒要明確拒絕敗壞的生活模式.】

啟18:6　她怎樣待人,也要怎樣待她,按她所行的加倍地報應她；用她調酒的杯加倍地調給她喝.

啟18:7　她怎樣榮耀自己,怎樣奢華,也當叫她照樣痛苦悲哀,因她心裡說:我坐了皇后的位,並不是寡婦,決不至於悲哀.

啟18:8　所以在一天之內,她的災殃要一齊來到,就是死亡,悲哀,饑荒.她又要被火燒盡了,因為審判她的主神大有能力.

【「巴比倫大城」,金錢世界,奢華生活,到了滅亡的時候了.】

啟18:9　地上的君王,素來與她行淫,一同奢華的,看見燒她的煙,就必為她哭泣哀號.

啟18:10　因怕她的痛苦,就遠遠地站著說:哀哉!哀哉!巴比倫大城,堅固的城啊,一時之間你的形罰就來到了.

【「地上的君王」,社會的管治階層,見到代表世上極盡奢華的「巴比倫大城」傾覆,發出哀嘆.他們渴望保持距離,但若他們不悔改,這也將是他們的命運.】

啟18:11　地上的客商也都為她哭泣悲哀,因為沒有人再買他們的貨物了；

【倚靠「巴比倫大城」的奢華生活來發達的商人,也發出哀鳴.】

啟18:12　這貨物就是金,銀【這來自歐洲的西班牙】,寶石,珍珠【這可能來自亞洲的印度】,細麻布,紫色料【這可能來自小亞細亞】,綢子(絲)【這可能來自中國】,朱紅色料,各樣香木,各樣象牙的器皿【這來自非洲】,各樣寶貴的木頭,和銅,鐵,漢白玉的器皿,

啟18:13　並肉桂,荳蔻,香料,香膏,乳香【這來自阿拉伯】,酒,油,細麵,麥子,牛,羊,車,馬,和奴僕,人口【這來自各方各面】.

【這個貨物清單,有些值得留意的地方.

一,這些貨物來自五湖四海.若以羅馬為中心,這些貨物來自東方的中國,印度,阿拉伯,小亞細亞；來自西方的西班牙；來自南方的非洲等,在當時,要維持這種貿易,需要極強大的軍力.

二,要過奢華生活,很難靠本土經濟.要享受各種奇珍異品,需要頻繁而花費的國際物流.

三,特別留意的是,這清單最後的貨物,是「人」.這裡,聖經希臘原文用了兩個字來描述:sōma和psychē,中文和合本聖經翻譯作「奴僕,人口」,這翻譯有點累贅,奴僕其實即人口.英文聖經KJV翻譯作「slaves and souls of men」,即奴隸與人的生命.NKJV翻譯作「bodies and souls of men」,這其實是希臘原文的字面翻譯,身體及人的生命.無論如何翻譯,都有點怪.Sōma這個用字,直譯是「身體」,這是對奴隸的最確切描述,奴隸不是人,他們是物質,是貨物,只是一種軀殼.

啟示錄反抗這種對人的貶視,反抗這種對人的商品化,反抗這種金錢世界的貿易本相,啟示錄採用聖經對人的稱謂,稱之為psychē,這用字出現在創世記2:7的希臘文譯本裡,神造人之後,吹口氣在他鼻孔裡,人就成了一living psychē,活著的「有生氣者」.人有神的氣息,具特別的尊嚴,不只是純粹的軀殼.羅馬人稱奴隸為「sōma軀殼」,啟示錄偏偏要稱他們為「psychē具有神的氣息的有生氣者」.故此,這貨物清單的結尾,用一種意譯的方式來表達的話,則是:最後的貨物是只有軀殼的奴隸,不,不,他們其實是人,是神聖的人.

四,貨物清單以奴隸作結,其實也點出金錢世界的運作本相.金錢世界很容易將人分為兩級,用馬克思的說話來講,即是「有產階級」和「無產階級」.有產階級過的,可以是極盡奢華的生活.這一切之所以可能,很多時是以無產階級的苦難為代價的.窮苦的人,活得像奴隸一樣.

五,美國國務院發表《2018年販運人口報告》(Trafficking in Persons Report 2018),香港被列入「第二級別的觀察名單」,這是什麼含意呢?按這報告的說明,這是指:香港沒有全面達到《人口販運受害者保護法》的最低標準,但正在努力中,與此同時,受害人卻正在增加,而政府未能提供證據顯示與上年相比加強了打擊人口販賣的努力,或作出相關承諾.(參:\url{https://www.state.gov/documents/organization/285921.pdf})

在香港,人像奴隸被賣,主要有兩類:一是家庭傭工的問題,2016年的一份報告指出,1/6的家庭傭工受到剝削.有人一日做17小時,忍受語言,性,身體的暴力等.二是將女性誘騙到香港後強迫賣淫.

(參:\url{https://www.state.gov/documents/organization/282801.pdf},頁215-217.)】

六,啟示錄的受信者,是七個教會,他們的地理位置,正處於羅馬的商業大道上.其中三個城巿:以弗所,推雅推喇,撒狄,都有奴隸巿場.啟示錄這裡講的貨品,是受信教會日日見到的買賣.這是要受信人思考,我們是無意識地融入社會既有的經濟體制中,抑或有所取捨?

七,啟示錄常常講人生命中的「印記」,這正是奴隸的標誌.人總是某個主人的奴隸,我們不是打上「666」的印記,就是打上「神」的印記.不是邪惡的奴隸,就是神的奴隸.你屬於誰,要好好選擇.

啟18:14　巴比倫哪,你所貪愛的果子離開了你；你一切的珍饈美味,和華美的物件也從你中間毀滅,決不能再見了.

啟18:15　販賣這些貨物,藉著她發了財的客商,因怕她的痛苦,就遠遠的站著哭泣悲哀,說:

啟18:16　哀哉!哀哉!這大城啊,素常穿著細麻,紫色,朱紅色的衣服,又用金子,寶石,和珍珠為妝飾.

啟18:17上　一時之間,這麼大的富厚就歸於無有了.

【不單社會的管治階層發出哀鳴,與金錢世界打交道的商人都哀嘆.「巴比倫大城」傾覆了,他們發財的機會也沒有了.】

啟18:17下　凡船主和坐船往各處去的,並眾水手,連所有靠海為業的,都遠遠的站著,

啟18:18　看見燒她的煙,就喊著說:有何城能比這大城呢?」

啟18:19　他們又把塵土撒在頭上,哭泣悲哀,喊著說:哀哉!哀哉!這大城啊.凡有船在海中的,都因她的珍寶成了富足!她在一時之間就成了荒場!

【連繫統治者與商人的物流業,也哀嘆,發財的機會也沒有了.】

啟18:20　天哪,眾聖徒,眾使徒,眾先知啊,你們都要因她歡喜,因為神已經在她身上伸了你們的冤.

【當「大巴比倫」傾覆時,人人哀鳴,天哪,眾聖徒,眾使徒,眾先知等卻是歡喜快樂.「大巴比倫」,代表強大的壓迫制度,代表將人奴隸化的體制,代表漠視人性的生活方式,它的傾覆,是人類根深蒂固的渴望.人類渴望不再做奴隸,人類渴望真正的解放,人類渴望人性化的生活.當這刻實現時,才是人類真正歡喜的時刻.我們渴望這天的來臨.】

長長的一段經文,是哀歌,是為人類在金錢世界的墮落而流下的眼淚,是為強權者在奢華生活中對他人的不擇手段而歎息.

長長的一段經文,是先知的抗議,他抗議世人向金錢下拜,他抗議世人不懂尋求神的正義,他抗議這種奴役人的經濟制度,他抗議人在種種物質的享受中甘願沉溺,他抗議社會的管治者不顧民情,他抗議商人巧取豪奪,他抗議國際之間只顧做買賣而不顧人民疾苦.

長長的一段經文,是受苦者的控訴,他們控訴這世界無天理,他們控訴人間的貧富懸殊得違反常情,他們控訴管治者只顧奢華生活,他們控訴世人只顧做生意,他們控訴這敗壞的制度竟無一人出來為這光景負責.

長長的一段經文,是神慈聲的呼喚,祂呼喚追求真理的人,離開這慾望的色聲犬馬；祂呼喚心靈夢想自由的人,遠離這讓人奴役的生活模式；祂呼喚渴想永恆的人,逃離這種墮落的奢華；祂呼喚世人,做一個真正的人.

長長的一段經文,是神嚴厲的判語.祂宣判「大巴比倫城」的滅亡,祂宣判這無視道德,殺害忠良,泯滅人性的強權制度的滅亡.祂宣告正義,祂宣告救贖,祂宣告所有人類的正直呼聲被應允.

我們呢?是無視滅亡的警告而繼續墮落的生活,抑遠離墮落而敬虔度日?

但願:榮耀歸於聖父,聖子,聖靈.阿們.



\newpage

\section{饒恕的中途站}
\label{sec:445}
\textbf{區祥江博士 - 饒恕的中途站}
\newline
\newline
\hyperref[sec:444]{< < < PREV SERMON < < <}
~
\hyperref[sec:toc]{[返目錄]}
~
\hyperref[sec:446]{> > > NEXT SERMON > > >}


創世記 45:1-15
\newline
\begin{longtable}{cl}
\hline
\hline
章節 & 經文 (和合本修訂版)\\
\hline
45:1 & \begin{tabularx}{0.7\textwidth}{X} 約瑟在所有侍立在他旁邊的人面前情不自禁，就喊叫說：「每一個人都離開我，出去吧！」約瑟和兄弟相認的時候沒有一人站在他那裡。 \end{tabularx} \\ \\ \relax
45:2 & \begin{tabularx}{0.7\textwidth}{X} 他放聲大哭，埃及人聽見了，法老家中的人也聽見了。 \end{tabularx} \\ \\ \relax
45:3 & \begin{tabularx}{0.7\textwidth}{X} 約瑟對他兄弟們說：「我就是約瑟。我的父親還在嗎？」他兄弟們不敢回答他，因為他們在他面前都很驚惶。 \end{tabularx} \\ \\ \relax
45:4 & \begin{tabularx}{0.7\textwidth}{X} 約瑟又對他兄弟們說：「靠近我一點。」他們就近前來。他說：「我是被你們賣到埃及的兄弟約瑟。 \end{tabularx} \\ \\ \relax
45:5 & \begin{tabularx}{0.7\textwidth}{X} 現在，不要因為把我賣到這裡而憂傷，對自己生氣，因為神差我在你們以先來，為要保全性命。 \end{tabularx} \\ \\ \relax
45:6 & \begin{tabularx}{0.7\textwidth}{X} 現在這地的饑荒已經二年了，還有五年不能耕種，沒有收成。 \end{tabularx} \\ \\ \relax
45:7 & \begin{tabularx}{0.7\textwidth}{X} 神差我在你們以先來，為要給你們在世上存留餘種，大施拯救，保全你們的性命。 \end{tabularx} \\ \\ \relax
45:8 & \begin{tabularx}{0.7\textwidth}{X} 這樣看來，差我到這裡來的不是你們，而是神。他又使我如同法老之父，作他全家之主，和埃及全地掌權的人。 \end{tabularx} \\ \\ \relax
45:9 & \begin{tabularx}{0.7\textwidth}{X} 你們要趕緊上到我父親那裡，對他說：『你兒子約瑟這樣說：神已立我作全埃及之主，請你下到我這裡來，不要耽擱。 \end{tabularx} \\ \\ \relax
45:10 & \begin{tabularx}{0.7\textwidth}{X} 你和你的兒子孫子，羊群牛群，以及一切所有的，都可以住在歌珊地，與我相近。 \end{tabularx} \\ \\ \relax
45:11 & \begin{tabularx}{0.7\textwidth}{X} 我要在那裡奉養你，因為還有五年的饑荒，免得你和你的家屬，以及一切所有的，都陷入窮困中。』 \end{tabularx} \\ \\ \relax
45:12 & \begin{tabularx}{0.7\textwidth}{X} 看哪，你們的眼睛和我弟弟便雅憫的眼睛都看見，是我親口對你們說話。 \end{tabularx} \\ \\ \relax
45:13 & \begin{tabularx}{0.7\textwidth}{X} 你們要把我在埃及一切的尊榮和你們所有看見的事情都告訴我父親，也要趕緊請我父親下到這裡來。」 \end{tabularx} \\ \\ \relax
45:14 & \begin{tabularx}{0.7\textwidth}{X} 於是約瑟伏在他弟弟便雅憫的頸項上哭，便雅憫也在他的頸項上哭。 \end{tabularx} \\ \\ \relax
45:15 & \begin{tabularx}{0.7\textwidth}{X} 他又親眾兄弟，伏著他們哭。過後，他的兄弟就和他說話。 \end{tabularx} \\ \\
[1ex]
\hline
\hline
\end{longtable}
講題:饒恕的中途站 Middle Station of Forgiveness

講員:區祥江博士

經文:創世記45:1-15

場合:香港中文大學崇基學院禮拜堂

日期:2019年2月24日

饒恕從來都不是一個簡單,容易的過程,當中會經過不同的階段,或歷程.我們可以說,創傷愈大,饒恕需要處理的時間就愈長.當然,傷害與被傷害,饒恕與被饒恕各自都有自己不同的心路歷程,這也令饒恕在這人世間所呈現的不同的面貌和不同的結果,變得多變和不容易找到一條放之四海皆準的定律.

剛巧今天的經訓就來到創世記 45:1-15節的一段約瑟饒恕兄弟的故事,故事描述一段讚人熱淚的相認,饒恕,甚至給人一個能達至復和的寄盼的片斷.但,這只是饒恕的一個中途站.對約瑟來說,能來到這一步已經絕對不容易.我們若看完創世記50 章的結局,我們就知道.原來饒恕與被饒恕的一方仍有很大的距離.這個感人的場面只是饒恕的一個中途站.

我們先來一個簡單的歷史回顧.

約瑟是一個傳奇人物,他的故事打從創世記37章出場到創世記50章結束,都有他的足跡. 他的故事被改編為音樂劇,Joseph and the Amazing Technicolor Dreamcoat 由 Tim Rice 填詞, Andrew Lloyd Webber 作曲.

他的故事情節十分精彩,有不少重覆出現的主題,以衣服為例,父親的彩衣本來是給大兒子的,因為約瑟是雅各深愛的妻子所生的第一個兒子,雖然在兄弟中他是排行第十一的,但彩衣卻落在他身上,這就惹了其他兄弟的嫉妒,更動了殺機,雖然,最後能免了一死,他的彩衣卻染上血蹟,用作他死亡的證據.第二件出現的衣服是成為波提乏家中的管家之後,雖然得到主人的信任,但卻受到波提乏妻子的引誘,逃走時留下一件外衣,卻成為指控他的證據.當然被鎖入監就要穿上囚衣,最後一件是法老賞識他而賜給他一件皇室的外袍,這是他得到尊榮的見證.所以,衣服的更替也標致著一個人的際遇.Tim Rice 在音樂劇中說得好:「Give me my colored coat My amazing colored coat 」,約瑟可能在眾多的衣服中,仍然最愛的是父親給他的彩衣.

其實,約瑟的彩衣是整個家庭傷痛與饒恕的起始點.家家有本難念的經,不容易定事情的對錯.有些家庭的互動是隔代留下來的影響.正如上面提過,約瑟的父親雅各特別寵愛約瑟,在17歲的時候,雅各送給他一件彩衣,在當時彩衣應該是送給長子的,雅各這樣的行徑釀成十個兄長對約瑟的妒忌.其實,當再推上一代,雅各的母親利伯加就是因為偏愛雅各,以致釀成日後以掃追殺雅各的地步.由此可見,約瑟被賣,錯可以追溯到利伯加和以撒身上.當我們去看家庭的恩怨,誰是誰非的時候,要把問題拉闊來看.

約瑟饒恕的第一個站:從神的眼光看傷害的意義

寬恕是個人在神面前,以神的眼中檢視過去,而看到神在傷害背後的意義.寬恕的第一步是有勇氣去檢視我們家庭所發生的事,傷害我們的人可能曾被人傷害.約瑟的哥哥們也許正因為雅各偏愛約瑟而受到傷害,每個人背後都有他的故事.寬恕並不是易事,約瑟被兄弟出賣,要寬恕他們確實不容易.

由被賣到埃及,被主母控告,被囚,被酒政忘恩負義,後來才被記起,為法老解夢,最終成為法老的宰相.當中的起跌數次,直至到最終才回望,看到這些苦難背後有神的心意,他才放下怨恨.

正如今天的經文中,約瑟多次指出他看到神是幕後的白手,祂是事情背後的看守人,叫傷害的事化成扭轉民族命運的契機.

45章

5節:現在,不要因為把我賣到這裡自憂自恨.這是神差我在你們以先來,為要保全生命.

7節:神差我在你們以先來,

9節:神使我作全埃及的主,

另一個端倪是在聖經創世記41章51-52節,,當時約瑟差不多 37 歲,因為聖經記載他30歲開始侍候法老,荒年未到以前,約瑟給他的兩個兒子起名,大兒子起名叫嗎哪西,意思是「使之忘了」,而次子的名字為以法蓮,意思是「使之昌盛」.由約瑟被賣,被陷害,被遺忘,到被法老重用,回望這20年,他開始明白今天落入這個處境是為了什麼.如他替兒子改名字便可見他已經寬恕了他的哥哥們.當我們用神的眼光去看待事情,便能夠帶給我們對傷害有一個新的看法,而這看法能幫助我們去寬恕.我們要學會寬恕,應先退一步先在神面前安靜,思想神為什麼容讓事情發生,當中有著神的恩典和計劃. 當然,我們也不應強迫身邊受傷的人一定要看到傷害背後神的意思,因為有不少些時候是未必找到答案,苦難本身是不容易給予一個合理的答案,我們才有約伯記對苦難的尋問的記錄.

另一個中途站是:傷痛的眼淚

不過,縱然自己私底下處理了自己的怨憤,放下對家人的憎恨和憤怒,在面對面處理這些家人的恩怨的時候,總會牽動我們不少的眼淚.約瑟在與兄弟修補關係的過程中,流了多少次的淚.

有時是要避開去哭(創42:23-24 )

「23他們不知道約瑟聽得出來,因為在他們中間用通事傳話. 24約瑟轉身退去,哭了一場」

今天經文也提到他放聲大哭的場面:創 45:1-2

1 約瑟在左右站著的人面前情不自禁,吩咐一聲說:「人都要離開我出去!」約瑟和弟兄們相認的時候並沒有一人站在他面前. 2他就放聲大哭；埃及人和法老家中的人都聽見了.

也有忍不住相擁而哭的(45:14-15).那些感人的場面給我們很大的盼望和安慰,因為如此大的被出賣,被陷害,令自己落到如此艱難的處境,人世間最大的傷痛莫過於是親人給自己的傷,最終,約瑟都能原諒了他們.相擁的淚是最大的明證.

饒恕是一條雙行路:我們也要分辨傷害人的是否有悔意

42-44章其實是約瑟測試兄弟是否有悔意的. 第一次買糧時,他要哥哥們帶便雅憫到埃及時,他偷聽到他們如何理解當前的苦難:「21他們彼此說:「我們在兄弟身上實在有罪.他哀求我們的時候,我們見他心裡的愁苦,卻不肯聽,所以這場苦難臨到我們身上.」 22呂便說:「我豈不是對你們說過,不可傷害那孩子嗎?只是你們不肯聽,所以流他血的罪向我們追討.」」(42:21-22)

接著約瑟設計用占卜用的杯子留住便雅憫.猶大頂罪那段說話最感人不過:「30我父親的命與這童子的命相連.如今我回到你僕人－我父親那裡,若沒有童子與我們同在, 31我們的父親見沒有童子,他就必死.這便是我們使你僕人－我們的父親白髮蒼蒼,悲悲慘慘地下陰間去了. 32因為僕人曾向我父親為這童子作保,說:『我若不帶他回來交給父親,我便在父親面前永遠擔罪.』 33現在求你容僕人住下,替這童子作我主的奴僕,叫童子和他哥哥們一同上去. 34若童子不和我同去,我怎能上去見我父親呢?恐怕我看見災禍臨到我父親身上.」」(44:30-34)

這頂罪的做法,預表猶大家所出的耶穌為我們頂罪的典範.也賺取了約瑟對他們真心悔罪的信心.這是饒恕的重要的一個站.

真饒恕是一個艱辛旅程,我們來看當中是雙方都要有行動和心意的轉化的.

犯錯/罪者懊悔,悔改,在罪疚中得釋放

受傷者得醫治,離開傷害,不捆在報復循環的詛咒中,願意施恩接納,饒恕

iii. 最終,冒險進入修補的關係中,破碎的關係得復和

被饒恕者那邊廂的心路歷程

不過,寬恕跟復和原來是兩件可以分開來看的事 .復和是雙方間的事,有時候寬恕未必帶來復和,這是兩個過程,寬恕是我們單方面可以做到,但復和卻需要別人也同時願意才可,而在未能完全復和之前,我們應以恩典待人.就像約瑟的10個哥哥,雖然約瑟寬恕了他們,但他們仍然放不低,依然懼怕約瑟.

在今天的經文中,我們不難看到那些兄弟的恐懼.

45:3 「約瑟對他弟兄們說:「我是約瑟.我的父親還在嗎?」他弟兄不能回答,因為在他面前都驚惶」

從經文中,我們也看到一個權力懸殊的情況,那被害人約瑟今天強盛起來,傷害者變成脆弱的一方,這令接受饒恕一方有一些不確定和恐懼.而其實約瑟在言談之間那種權力的氣勢是十分強盛的.

45:8 他又使我如法老的父,作他全家的主,並埃及全地的宰相.

45:13  你們也要將我在埃及一切的榮耀和你們所看見的事都告訴我父親,

那親疏有別的次序是明顯的:「父親,便雅憫,其他兄弟」

言者無心,聽者有意.這種權力懸殊帶來的恐懼其實一直未散,他們擔心約瑟不是真心饒恕他們.

我們可以從雅各去世後兄弟的恐慌和部署要討好約瑟可見:

「約瑟的哥哥們見父親死了,就說:「或者約瑟懷恨我們,照著我們從前待他一切的惡足足地報復我們.」  他們就打發人去見約瑟,說:「你父親未死以先吩咐說: 『你們要對約瑟這樣說:從前你哥哥們惡待你,求你饒恕他們的過犯和罪惡.』如今求你饒恕你父親神之僕人的過犯.」他們對約瑟說這話,約瑟就哭了.他的哥哥們又來俯伏在他面前,說:「我們是你的僕人.」」(創50:15-18)

到這時,約瑟仍然以眼淚作回應,因為他為兄弟信不過他已經寬恕了他們而傷心,原來他寬恕了,不代表兄弟能釋懷,所以接著,「約瑟對他們說:「不要害怕,我豈能代替神呢? 從前你們的意思是要害我,但神的意思原是好的,要保全許多人的性命,成就今日的光景現在你們不要害怕,我必養活你們和你們的婦人孩子.」於是約瑟用親愛的話安慰他們.」(創50:19-21)

約瑟的回答反影他是一個有氣量,善解人意,約瑟決定恩待他的哥哥們,願意供養他們的妻兒.而且約瑟更用親愛的話去安慰他們.

從約瑟寬恕的故事,我們看到他能夠寬恕其實有不同的站,當中神的角色是不可或缺的.

第一,約瑟相信神是最終的審判者,他會覺得若自己去報復是代替了神作審判者的角色.第二,他相信神能轉化人不良的動機為美善的結局,他沒有否定他們真的立心不良,真的做了傷害他的事,不過,他卻而回望式的角度來理解傷害背後,原來有神的美意在當中.當我們寬恕別人的時候,我們用神的眼光去看事情,我們有空間去看神讓這些事情發生,是有祂特定的計劃,第三,更進一步,有神的眼光看傷害,看傷害自己的人.原來這些傷害自己的人都有他們受傷的地方,例如雅各的偏愛和約瑟不懂收斂自己帶動兄弟仇恨的兩個夢,都傷了兄弟的感情.這部份聖經沒有詳細記載.最後,能否從饒恕到復和原來也牽涉到權力的懸殊,被饒恕者對饒恕者的信心問題.創世記 50 章故事結尾其實也沒有寫下兄弟們對約瑟那眼淚,那安慰的話的反應,故事結尾的留白,似乎在邀請我們也當作出自己的決定呢!



\newpage

\section{幽黑愈下瞥見榮光}
\label{sec:446}
\textbf{伍渭文牧師 - 幽黑愈下,瞥見榮光}
\newline
\newline
\hyperref[sec:445]{< < < PREV SERMON < < <}
~
\hyperref[sec:toc]{[返目錄]}
~
\hyperref[sec:447]{> > > NEXT SERMON > > >}


路加福音 9:28-31
\newline
\begin{longtable}{cl}
\hline
\hline
章節 & 經文 (和合本修訂版)\\
\hline
9:28 & \begin{tabularx}{0.7\textwidth}{X} 說了這些話以後約有八天，耶穌帶著彼得、約翰、雅各上山去禱告。 \end{tabularx} \\ \\ \relax
9:29 & \begin{tabularx}{0.7\textwidth}{X} 正禱告的時候，他的面貌改變了，衣服潔白放光。 \end{tabularx} \\ \\ \relax
9:30 & \begin{tabularx}{0.7\textwidth}{X} 忽然有摩西和以利亞兩個人同耶穌說話； \end{tabularx} \\ \\ \relax
9:31 & \begin{tabularx}{0.7\textwidth}{X} 他們在榮光裡顯現，談論耶穌去世的事，就是他在耶路撒冷將要完成的事。 \end{tabularx} \\ \\
[1ex]
\hline
\hline
\end{longtable}
講題:幽黑愈下,瞥見榮光 A Glimpse of Glory in the Dimming World

講員:伍渭文牧師

經文:路加福音9:28-31

場合:香港中文大學崇基學院禮拜堂

日期:2019年3月3日

高清電視出現,對一些電視藝員和新聞主播來說是夢魘,臉龐風光,一覽無遺.手機的美照p圖又帶來安慰.有次朋友把合照發給我,我驚奇為何這樣?不像我,這是把我易容,transfiguration.但耶穌登山變像不是易容,是顯容,他是上主的榮光,為了成為世界的光,易容取了人的形像,成為木匠之子.今天是Transfiguration Sunday,譯作顯容主日較為正確.跟著的周三是聖灰日,開始為期四十天的預苦期了.聖灰日也是顯容日,聖灰日的塗灰禮主禮把去年棕枝主日的棕枝燒成灰燼,塗在額上:「你來自塵土,仍要歸回塵土,要遠離罪惡,跟隨基督.」登山變像顯出耶穌容貌,塗灰禮顯出我們的容貌.

耶穌是世界的光,我們可以用光的不同形態來表述教會年基督的生平:將臨期就是燭光(Candlelight).顯現期是耀光(Brilliance),他使水變酒,瞎眼得見光明,瘸腿行走；他的權能和慈愛,如光照耀.預苦期聖周這光曾在正午到申初三小時,變成黑暗(Darkness),但第三天復活成為不能仰視的正午日照(Mid-day Sun)；再過五十天,聖靈降臨,這光成為熊熊火焰.到了世界的終末,這光成為榮耀(Glory),與眾天使從天而降.今天福音書經課,給我們窺探到他的榮耀,瞥見榮光.

1 上主榮光,除去黑暗(路九28-31)

「耶穌帶著彼得,約翰,雅各上山去禱告.正禱告的時候,他的面貌就改變了,衣服潔白放光.忽然有摩西,以利亞兩個人同耶穌說話；他們在榮光裡顯現,談論耶穌去世的事,就是他在耶路撒冷將要成的事.」(路九28-31)

耶穌來到世界,直視世界的黑暗,和黑暗周旋,甚至被釘在十架上,撲滅幽暗的權勢.在十架的清影下,我們才能明白馬槽的意義.教會用四個主日的將臨期迎接聖誕,但聖誕期祇得十二天.聖誕翌日二十六日,教會年不是Boxing Day (拆禮物日),是聖使提反日,紀念第一位殉道者；二十七日聖約翰日,唯一沒有殉道的老使徒,他是信仰堅貞的代表,晚年因信仰被放逐拔摩海島.史提反是紅色的見證人,約翰是白色的見證人,他們都是殉道士(見證人).二十八日是嬰孩被殺日(Holy Innocents Feast),教會紀念死於希律王刀下的兩歲孩童,視他們為信仰而犧牲,他們的靈魂被上主接收.這天或跟著的主日,我們在崇拜中紀念不能保護自己的孩童:被虐待的兒童,被販賣的雛妓,非法的童工,國內的留守兒童.聖誕後的主日稱聖家主日(Holy Family Sunday),福音經課讀出約瑟與馬利亞攜同幼童耶穌逃往埃及,作政治難民.

這些聖誕後緊隨的節日和經課,顯示幽黑愈下,罪惡囂張.黑暗要消滅光,光要驅走黑暗.耶穌榮耀的生命,在十架上一役竟其功,撲滅了黑暗的權勢.

到了一月六日,聖誕期結束,顯現期開始,耶穌在約但受洗,顯現期的結束是耶穌登山顯容；一前一後,都顯明耶穌彌賽亞的身份,他要為人贖罪.洗禮和顯容,都有父從天上發聲.領洗時,父從天上發聲:「這是我的愛子,我所喜悅的.」(可一10)我所喜悅就是彌賽亞的名號.山上顯容:「有聲音從雲彩裡出來,說:這是我的兒子,我所揀選的,你們要聽他.」(路九35)聽他甚麼呢?就是前面所說:「人子必須受許多的苦,被長老,祭司長,和文士棄絕,並且被殺,第三日復活.」(路九22)摩西,以利亞在榮光裡顯現,同耶穌說話,談論耶穌去世的事,就是他在耶路撒冷將要成的事.(路九28-31)

耶穌在山上顯容,顯出原本的榮耀,但為了親近我們,耶穌易容成為世上的光,且死在十架上,引領我們進入光明.

2 上主榮光,顯出憐憫  (出三十四29,30)

出三十四29,30「摩西手裡拿著兩塊法版下西奈山的時候,不知道自己的面皮因耶和華和他說話就發了光.亞倫和以色列眾人看見摩西的面皮發光就怕挨近他.」

第二法版.首先,這兩塊法版是第二法版,能夠有這二法版是上主的憐憫.原來的第一法版已被摩西打碎了,因為百姓按捺不住,見摩西久久還未下山,造了金牛犢代替上主,摩西聽到跳舞的喧嘩聲音,大大震怒,知道闖出禍了,就把法版砸碎.他求上主赦免百姓的罪,甚至願意把自己的名字從生命冊塗抹.

上主的背.跟著摩西希望上主顯出他的榮耀,好叫他壯膽,帶領這班難搞的百姓.「摩西說,求你顯出你的榮耀給我看.耶和華說:「我要顯我一切的恩慈,在你面前經宣告我的名.我要恩待誰就恩待誰；憐憫誰就憐憫誰；又說:你不能看見我的面,因為人見我我的面不能存活.耶和華說:看哪,在我這裡有地方,你要站在磐石穴中,用我的手遮掩你,等我過去,然後我要將我的手收回,你就得見我的背,卻不得見我的面.」(出三十三18-23)為何摩西求看榮耀,上主顯出他的憐憫呢?上主以他的憐憫為他的榮耀.人不能直視上主的臉,祇能看到上主的背,這背就是耶穌,道成肉身的耶穌,耶穌代表上主的憐憫.

以憐憫為榮耀.當耶穌被掛在十架,從正午至申初三個小時,最後喊叫:成了!此期間,白晝轉為玄黑,天昏地暗,磐石崩裂,聖殿分隔神人幔子從上到下裂開,就是最榮耀的一刻,上主以他的憐憫為他的榮耀.為何白晝變成黑暗?當然這事驚天動地,宇宙搖盪,令天象大變.

另一個解釋是,因為我們的罪白晝轉為黑暗.2014年香港展出意大利巴洛克善用光影的藝術大師卡拉瓦喬(Caravaggio)名畫「以馬忤斯的晚餐」,主辦單位為這名畫買了上億港元的保險.主題:光影Light and Shadow.同場亦有幾位本港藝術家展出以光影為主題的作品.一位藝術家在一潔白的熒幕上打出十誡的誡命:不可殺人…….一條幼幼長長的字條在潔白的熒幕上轉動,速度愈來愈快,熒幕愈來愈暗,瞬間成為玄黑.這作品很有啟發性,在十架上的基督,此刻背負著世人的罪過,你和我的罪過也在其上,是我們的罪過,使正午轉為黑暗.

在人幽暗的時刻,陪伴他們.上主以他的憐憫,為他的榮耀.我們甚麼時候,叫人看到上主的榮光呢?當人落在幽黑中,當人看不到任何的光明,我們代表基督,陪伴他們,他們就瞥見榮光.

在我牧養數十年的經驗中,有兩件難忘的事,使我感到傳道職事的榮耀.有一次趕到醫院,一進病房,一位年輕太太就抱著我哭不成聲,她三十歲的丈夫剛剛離開.我不能說些甚麼,就陪著她,她不斷飲泣.這位失去丈夫的姊妹,稍後她就轉了教會,我們也轉了協助另一間教會,失去聯絡,但師母跟我常常提起她.多年後,偶然再見這位姊妹,我們感到很興奮,曉得她有了新的家庭,生活安定.

另一件是家庭探訪的經驗.一位在小學作校工的弟兄,工作時間長,家庭責任重,不多來聚會.他太太早年離世,育有兩個讀初小的頑皮男孩.每天工作回家,身體疲倦,還要忙著弄晚餐給兩個小傢伙,父兼母職.他問我,為何生活如此艱難?我祇陪伴他,聽他吐苦水,最後為他禱告就離開.他家在麗瑤公共屋村,在歸家途中,小巴下山,涼風拂臉,暮色蒼茫,心裡悲涼,眼眶濕潤.同時也感到很有榮耀,因為能代表基督,陪伴安慰困苦的人.盧雲神父說:「我們可否願意不逃避痛苦,無能為力時不故作忙碌,站在哀傷的人旁邊一起面對死亡?」

3 瞥見榮光,榮耀職事 (林後三12；路九37-40；詩九十九6,7)

林後三7-12「那用字刻在石頭上屬死的的職事尚且有榮光,甚至以色列人因摩西臉上的榮光,不能定睛看他的臉；這榮光原是漸漸退去的,何況那屬靈的職事豈不是更有榮光嗎?……我們既有這樣的盼望,就大瞻講說,不像摩西將帕子蒙在臉上.」

(路九37-40)「第二天,他們下了山,就有許多的人迎見耶穌.其中有一人喊叫說:「夫子!求你看顧我的兒子,因為他是我的獨生子.他被鬼抓住就忽然喊叫；鬼又叫他抽瘋,口中流沫,並且重重地傷害他,難以離開他.我求過你的門徒,把鬼趕出去,他們卻是不能.」

榮耀職事.耶穌是世界的光,驅走心中的黑暗絕望,使我們得著光明,也得著福音的職事.這福音職事使我們直視幽暗勢力造成扭曲的臉容.耶穌給彼得,雅各,約翰在山上看見他榮耀的臉相,也在山腳人世間看到被黑暗勢力扭曲的臉相,直視這些,不要逃避.但要曉得,我們不能靠自己與黑暗力量對峙,祇能靠耶穌,因他是上主的榮光.

當我們在香港街道上,看到老婆婆很吃力推著大堆紙皮上斜路時；當我們看見在國內山區,十歲的女孩背著三歲的弟弟在背上,在燒柴弄晚餐時,我們的心都破碎了,想哭了.這些年長和年幼的正需要幫助的人,社會竟然遺棄了他們,不顧他們.當我看見一班人享用豪宴,大大桌子中央堆著鮮花,內圍是滿布花瓣的水池,一位少女站在水中,為客人端菜,她臉露惶恐,因盛菜的器皿很重很熱.上載這照片是一位作家,她說:「願這位女孩無恙吧.」我看見這圖片時,不是心酸,是憤怒.

顯容主日,在山上瞥見榮光,目的為了山下被罪扭曲,被罪壓制的臉相.世界的光,目就是為了驅走幽暗,使人看見光明.但我們緊記,憑自己我我們不能勝過這幽暗的力量.「他被鬼抓住就忽然喊叫；鬼又叫他抽瘋,口中流沫,並且重重地傷害他,難以離開他.我求過你的門徒,把鬼趕出去,他們卻是不能.」(路九40 ) 尼采說得對:「與怪獸搏鬥的人要謹防自己因此而變成怪獸.如果你長時間的盯著深淵,深淵也會盯著你.」所以雅各書提醒我們:「在上主我們的父面前,那清潔沒有沾污的虔敬,就是看顧在患難中的孤兒寡婦,並且保守自己不沾染世俗.」

摩西,亞倫,撒母耳.今天詩篇經課九十九篇,列出三位為人代求的祭司先知(priest-prophets),他們是履行榮耀職事的好榜樣.「在他的祭司中有摩西和亞倫；在求告他的名的人中有撒母耳.他們求告耶和華,他就應允他們.他在雲柱中對他們說話他們遵守他的法度和他所賜給他們的的律例.」(詩九十九6,7)

他們是滿有憐恤的代求者.摩西在西乃山為百姓代求,甚至願意把自己的名字從生命冊抹去(出三十二至三十四).亞倫為挑戰摩西和自己的敵人代禱.利未人中可拉党為首糾眾挑戰摩西和亞倫的領導地位,他們煽動百姓叛變,上主降下瘟疫.亞倫快快拿著香爐走到百姓當中禱告,止住瘟疫.(民十六44-48).撒母耳為以色列人面對非利士人的攻擊代求.「願你不住地為我們呼求耶和華-我們的上主,求我們脫離非利士人的手.」(撒上七8)

他們聆聽律法並遵守律法.「他(上主)在雲柱中對他們說話他們遵守他的法度和他所賜給他們的律例.」(詩九十九7)他們為人代禱並教導律法,他們是祭司先知.他們的力量來自親近上主,聆聽上主.這是保守自己直視深淵而不墮進深淵,保守自己與怪獸搏鬥而不變成怪獸的秘訣.

耶穌在世除了趕鬼醫病,還有傳揚天國的福音,福音就是上主的話語.基督信仰不能被約化成基督教人道主義,我們有信息.凡前來就光的,心中就明亮；凡悔改的,罪就得著赦免.上主的話語就是腳前的燈,路上的光.要直視黑暗,首先要被上主的話語光照.上主的聖道,保守我們看顧患難中的孤兒寡婦時,自己不沾染世俗.

結論:基督登山變像不是易容,是顯容,他本是滿有榮光的上主,為了我們易容成為人,目的是親近我們,他是世界的光,被黑暗的世界拒絕並釘死在十架上,從午正到申,白晝因為我們的罪轉成黑暗.耶穌在十架的犧牲,顯出了上主的憐憫,這憐憫也是上主所誇的榮耀.我們已蒙這榮耀所光照,就讓我們走進幽黑愈下的世界,顯出上主的憐憫,讓被罪壓傷的人,瞥見榮光.

下周三的聖灰日,是預苦期的開始,聖灰日也是我們的顯容日:「你來自塵土,仍要歸回塵土,要遠離罪惡,跟隨基督.」讓我們歸回塵土前多作主工,走進世人的苦難中,讓他們瞥見上主榮光.



\newpage

\section{耶穌為猶大洗腳}
\label{sec:447}
\textbf{劉國偉先生 - 耶穌為猶大洗腳}
\newline
\newline
\hyperref[sec:446]{< < < PREV SERMON < < <}
~
\hyperref[sec:toc]{[返目錄]}
~
\hyperref[sec:448]{> > > NEXT SERMON > > >}


約翰福音 13:1-2,12-30
\newline
\begin{longtable}{cl}
\hline
\hline
章節 & 經文 (和合本修訂版)\\
\hline
13:1 & \begin{tabularx}{0.7\textwidth}{X} 逾越節以前，耶穌知道自己離世歸父的時候到了。他一向愛世間屬自己的人，就愛他們到底。 \end{tabularx} \\ \\ \relax
13:2 & \begin{tabularx}{0.7\textwidth}{X} 晚餐的時候，魔鬼已把出賣耶穌的意思放在加略人西門的兒子猶大心裡。 \end{tabularx} \\ \\ \relax
13:3 & \begin{tabularx}{0.7\textwidth}{X} 耶穌知道父已把萬有交在他手裡，且知道自己是從神出來的，又要回到神那裡去， \end{tabularx} \\ \\ \relax
13:4 & \begin{tabularx}{0.7\textwidth}{X} 就離席站起來，脫了衣服，拿一條手巾束腰， \end{tabularx} \\ \\ \relax
13:5 & \begin{tabularx}{0.7\textwidth}{X} 隨後把水倒在盆裡，開始洗門徒的腳，並用束腰的手巾擦乾。 \end{tabularx} \\ \\ \relax
13:6 & \begin{tabularx}{0.7\textwidth}{X} 到了西門‧彼得跟前，彼得對他說：「主啊，你洗我的腳嗎？」 \end{tabularx} \\ \\ \relax
13:7 & \begin{tabularx}{0.7\textwidth}{X} 耶穌回答他說：「我所做的，你現在不知道，但以後會明白。」 \end{tabularx} \\ \\ \relax
13:8 & \begin{tabularx}{0.7\textwidth}{X} 彼得對他說：「你絕對不可以洗我的腳！」耶穌回答他：「我若不洗你，你就與我無份了。」 \end{tabularx} \\ \\ \relax
13:9 & \begin{tabularx}{0.7\textwidth}{X} 西門‧彼得對他說：「主啊，不僅是我的腳，連手和頭也要洗！」 \end{tabularx} \\ \\ \relax
13:10 & \begin{tabularx}{0.7\textwidth}{X} 耶穌對他說：「凡洗過澡的人不需要再洗，只要把腳一洗，全身就乾淨了。你們是乾淨的，然而不都是乾淨的。」 \end{tabularx} \\ \\ \relax
13:11 & \begin{tabularx}{0.7\textwidth}{X} 耶穌已知道要出賣他的是誰，因此說「你們不都是乾淨的」。 \end{tabularx} \\ \\ \relax
13:12 & \begin{tabularx}{0.7\textwidth}{X} 耶穌洗完了他們的腳，就穿上衣服，又坐下，對他們說：「我為你們所做的，你們明白嗎？ \end{tabularx} \\ \\ \relax
13:13 & \begin{tabularx}{0.7\textwidth}{X} 你們稱呼我老師，稱呼我主，你們說的不錯，我本來就是。 \end{tabularx} \\ \\ \relax
13:14 & \begin{tabularx}{0.7\textwidth}{X} 我是你們的主，你們的老師，尚且洗你們的腳，你們也應當彼此洗腳。 \end{tabularx} \\ \\ \relax
13:15 & \begin{tabularx}{0.7\textwidth}{X} 我給你們作了榜樣，為要你們照著我為你們所做的去做。 \end{tabularx} \\ \\ \relax
13:16 & \begin{tabularx}{0.7\textwidth}{X} 我實實在在地告訴你們，僕人不大於主人；奉差的人也不大於差他的人。 \end{tabularx} \\ \\ \relax
13:17 & \begin{tabularx}{0.7\textwidth}{X} 你們既知道這些事，若是去實行就有福了。 \end{tabularx} \\ \\ \relax
13:18 & \begin{tabularx}{0.7\textwidth}{X} 我不是指著你們眾人說的，我知道我所揀選的是誰；但是要應驗經上的話：『吃我飯的人用腳踢我。』 \end{tabularx} \\ \\ \relax
13:19 & \begin{tabularx}{0.7\textwidth}{X} 事情還沒有發生，我現在先告訴你們，讓你們到事情發生的時候好信我就是那位。 \end{tabularx} \\ \\ \relax
13:20 & \begin{tabularx}{0.7\textwidth}{X} 我實實在在地告訴你們，接納我所差遣的就是接納我；接納我的就是接納差遣我的那位。」 \end{tabularx} \\ \\ \relax
13:21 & \begin{tabularx}{0.7\textwidth}{X} 耶穌說了這些話，心裡憂愁，於是明確地說：「我實實在在地告訴你們，你們中間有一個人要出賣我。」 \end{tabularx} \\ \\ \relax
13:22 & \begin{tabularx}{0.7\textwidth}{X} 門徒彼此相看，猜不出他說的是誰。 \end{tabularx} \\ \\ \relax
13:23 & \begin{tabularx}{0.7\textwidth}{X} 門徒中有一個人，是耶穌所愛的，側身挨近耶穌的胸懷。 \end{tabularx} \\ \\ \relax
13:24 & \begin{tabularx}{0.7\textwidth}{X} 西門‧彼得就對這個人示意，要問耶穌是指著誰說的。 \end{tabularx} \\ \\ \relax
13:25 & \begin{tabularx}{0.7\textwidth}{X} 於是那人緊靠著耶穌的胸膛，問他：「主啊，是誰呢？」 \end{tabularx} \\ \\ \relax
13:26 & \begin{tabularx}{0.7\textwidth}{X} 耶穌回答：「我蘸一點餅給誰，就是誰。」耶穌就蘸了一點餅，遞給加略人西門的兒子猶大。 \end{tabularx} \\ \\ \relax
13:27 & \begin{tabularx}{0.7\textwidth}{X} 他接了那餅以後，撒但就進入他的心。於是耶穌對他說：「你要做的，快做吧！」 \end{tabularx} \\ \\ \relax
13:28 & \begin{tabularx}{0.7\textwidth}{X} 同席的人沒有一個知道耶穌為甚麼對他說這話。 \end{tabularx} \\ \\ \relax
13:29 & \begin{tabularx}{0.7\textwidth}{X} 有人因猶大管錢囊，以為耶穌是對他說「你去買我們過節所需要的東西」，或是叫他拿些甚麼給窮人。 \end{tabularx} \\ \\ \relax
13:30 & \begin{tabularx}{0.7\textwidth}{X} 猶大受了那點餅以後立刻出去。那時候是夜間了。 \end{tabularx} \\ \\
[1ex]
\hline
\hline
\end{longtable}
講題:耶穌為猶大洗腳 Jesus Washed the Feet of Judas

經文:約翰福音13:1-2,12-30

講員:劉國偉先生

場合:香港中文大學崇基學院禮拜堂

日期:2019年3月10日

引子

捷克共和國前總統哈維爾,在1978年發表了一篇對後世影響甚深的文章「無權力者的力量」The Power of the Powerless；那時,正值哈維爾於1977年向當時的執政者,統一工人黨的總書記胡薩克寫一封公開信,剛好入獄一年之後.

在公開信中,哈維爾直接指出當時社會制度下的人性的腐敗和制度本身的衰敗,哈維爾深信人性的彰顯和人權的釋放才是當時社會的唯一出路.哈維爾其後被傳訊,在1977年10月被判以「危害共和國利益」之罪名,判處十四個月有期徒刑；然後,他在1978年10月仍在獄中的時候,發表了「無權力者的力量」The Power of the Powerless；接著,在1979年,哈維爾被判以「顛覆共和國」之罪名,再被判處有期徒刑四年半

在「無權力者的力量」這篇文章,哈維爾寫下了一段文字特別吸引我 :「正是平凡人如我們,未必有勇氣站到最前,不過大部分人有能力鼓起小勇氣,於平凡生活中不順從,不妥協,不隨著大水流漂浮,每一個日常崗位,都是堅持的陣地.」

約翰福音第13章記載耶穌為門徒洗腳.面對地方權力龐大的猶太公會,那些尊貴的領袖,擁有宗教和地方權力,有財富和聲望,也有一群追求安定生活的百姓群眾,更有來自羅馬的祝福；耶穌呢? 一個沒有實質權力的拉比,在讀者眼中,他與猶太公會博奕的過程中,顯得勝算無望,當祭司長要策劃殺害耶穌之時,耶穌能做的,只可以吃飯和洗腳?

猶大的心思

「逾越節以前,耶穌知道自己離世歸父的時候到了.他既然愛世間屬自己的人,就愛他們到底. 2 吃晚飯的時候,魔鬼已將賣耶穌的意思放在西門的兒子加略人猶大心裡.」(約13:1-2)

從少年時期信主開始,閱讀約翰福音很多次,一直到二零零一年九月,當我再次走進約翰福音第十三章,講及耶穌在為門徒洗腳的情景,我實在被眼前的影像弄得發呆,久久不能離開.那是一次在完全沒有心理準備的情況下,被邀請往一個機構的同工計劃營中主領聚會,我在匆忙之間,決定以耶穌為門徒洗腳的動作為默想和分享的重點,讓我意外地展開這段耶穌為門徒洗腳的旅程.那一次,我深深進入使徒約翰這位老人家的超越半世紀的仔細回憶,既驚訝主耶穌那份深厚精緻的愛,亦訝異於約翰那刻入靈魂骨髓的記憶.

再過兩年,我為一份基督教刊物構思一系列的文章,這段經文又再浮現腦際,於是我再次深入研讀這段經文.我沈思於耶穌為門徒洗腳的每個動作,我更驚訝於耶穌為猶大洗腳!也許我的心靈真的過於狹小,我不懂得如何體會,為出賣自己的人洗腳,是一種怎樣的情操,兩個星期過去,我始終無法下筆,也從沒試過讀聖經讀得這般思潮起伏.

然後,在一次聖餐崇拜中,我再次想起了使徒約翰的話:「祂就愛他們到底」,NIV譯本這樣寫道:he now showed them the full extent of his love(約13:2),一想到耶穌為賣祂的人洗腳,我的眼淚便不能自禁.

十幾年來,我都在思想猶大的心思.約翰寫道:「魔鬼已將賣耶穌的意思,放在西門的兒子加略人猶大心裡.」按其他三卷福音書(太26:14-16；可14:10-11；路22:3-6)更強調在晚餐之前,特別按照馬太和馬可的記載,在伯大尼西門家裡耶穌被澆以極品香膏之後,如路加所言「撒旦入了猶大的心.」(路22:3)從符類福音得知,猶大在賣主這事上相當主動,包括接觸祭司長和守殿官(祭司長唯一可以調動的武裝部隊就是聖殿的衛警),並與他們商量具體的行動方案,顯然猶大對整件事情有所評估及計算,絕非一時衝動.

有人認為猶大賣主的真正動機是一宗歷史懸案,還是使徒約翰最直接,他以敘事者的身份直截了當的告訴讀者猶大是個賊(約12:6),猶大批評馬利亞浪費了可以賣得三十兩的香膏,接著,猶大就去主動聯絡祭司長要出賣耶穌.

你們也當彼此洗腳

耶穌洗完了他們的腳,就穿上衣服,又坐下,對他們說:「我向你們所做的,你們明白嗎? 你們稱呼我夫子,稱呼我主,你們說的不錯,我本來是.我是你們的主,你們的夫子,尚且洗你們的腳,你們也當彼此洗腳.我給你們作了榜樣,叫你們照著我向你們所做的去做.(約13:12-15)

猶大接觸祭司長一天之後,在這個特別的夜晚,耶穌為門徒洗腳.除了西門彼得與耶穌有一番對話之外,看來其他門徒都默默的接受耶穌的服事,耶穌逐個逐個的洗,再擦乾,當然包括猶大.

聖經沒有特別記載耶穌為猶大洗腳的情景,使徒約翰沒有落墨的主要原因,很可能是當時根本沒有任何特別的事情發生,就像其他十個沒有拒絕耶穌洗腳的門徒一樣,耶穌用他的雙手觸及猶大的腳,用水洗淨,再拿手巾擦乾.我無法探究當中的張力:耶穌為賣祂的人洗腳,猶大正接受他快要出賣的老師的服事,耶穌的雙手觸摸猶大的雙腳,對耶穌和猶大來說,那是一種甚麼滋味?!

按約翰的回憶和判斷,當耶穌為門徒洗腳的時候,猶大已經決定把自己賣給魔鬼,要作一個背叛師傳,出賣老師的門生,但耶穌卻洗他的腳,與他同台吃飯!別說是出賣我們的人,即使是辱罵,抵毀,和冤枉我們的人,我們尚且不容易原諒他們,即或可以,也未必有足夠動力去主動服事他們,即或也可以,在我們知道對方要出賣我們之前,甚至快要把我們置諸死地,我們是否仍有充足的愛去觸摸他那一雙污穢和疲乏的腳?  耶穌竟然那樣做過.

就在這刻,耶穌這位看似無權力的拉比,甚至是即將被出賣的拉比,向門徒說 :「我是你們的主,你們的夫子,尚且洗你們的腳,你們也當彼此洗腳.」(約13:14) 耶穌吩咐門徒要彼此洗腳,這是耶穌當晚向門徒發出的第一個命令,祂說:「我給你們作了榜樣,叫你們照著我向你們所作的去作.」(約13:15)

你怎樣理解「你們也當彼此洗腳」和「叫你們照著我向你們所作的去作」這兩句話? 耶穌一再要求門徒仿傚祂,以祂作為榜樣,這可不是鬧著玩的,可曾想過甚麼是「以祂作為榜樣」?基督的榜樣透過祂設立的聖餐清楚地透析出來,那是一種捨棄自己的愛,包含了我們實踐所達至的最深程度,就是一種為對方捨命的行為,可以是生死攸關的事,可真不是鬧著玩的!

耶穌這位看似無權力的拉比,甚至看來連自己的命也保不住,祂問門徒,也問猶大:「我向你們所作的,你們明白麼?」(約13:12)耶穌選擇了服侍和捨己的道路,我相信猶大和祭司長當然明白服侍和捨己的含義,但在他們的世界中,這些生命素質是不會被認真看待的.所以耶穌說:「現在要應驗經上的話,說:『同我吃飯的人用腳踢我.』」(約13:18b)(這應驗了詩篇41篇9節) 我們難以相信猶大不明白耶穌說話的意思,然後,耶穌就把話說得更明顯了.

那時候是夜間了

耶穌說了這話,心裡憂愁,就明說:「我實實在在地告訴你們,你們中間有一個人要賣我了.」(約13:21) 按當時的坐姿, 大家都以左手托頭,半倚在長椅上,頭則靠近台邊,若果主所愛的門徒(約翰)可以靠在耶穌的胸懷去問祂:「主啊! 是誰呢? 」那自然是坐在耶穌的右邊.按馬太福音的記載,當耶穌說你們中間有一個人要賣我了,門徒就一個一個問祂 : 主啊!是我麼?  馬太刻意記載了猶大與耶穌的對話――賣耶穌的猶大問他說:「拉比,是我嗎?」耶穌說:「你說的是.」(太26:25)按約翰福音的描述,耶穌對猶大的對答,在當時並沒有引起門徒有什麼反應,加上猶大的座位,大概是耶穌手所能及之處,才自然地蘸了一點餅,遞給猶大(約13:26) ,猶大極可能就坐在耶穌的左側,那是與主人相近的座位,表達了一個尊貴及親密的位份!而且,按猶太人風俗,主人蘸餅是主人十足友好的表示.奈何,這麼近卻那麼遠,猶大與耶穌所想所作有很大分歧,或者也可以說,耶穌洗猶大的腳,猶大不為所感,他有他自己的議程,耶穌洗完了約翰的腳,老約翰將那一幕牢牢記住；耶穌為猶大洗腳,但猶大卻仍然按照他的議程行事.

並非很多人願意被別人改變,也實在有很多人想要改變身邊的人.猶大不認同耶穌所為,並按自己的議程去強行改變耶穌的議程,意圖利用耶穌來達成自己的目標.信仰生命裡其中一處吊詭的地方,就是不願意被別人改變的我們,要接受耶穌去改變我們的生命,不是我們去改變耶穌,反過來是我們必須被耶穌所改變.人類所面臨的悲劇,不單在於我們忘記了造物主,很多時是因為我們自以為比上帝更了解這個世界,和自己的處境,並取代上帝的位置,自訂議程,而沒有讓上帝真正改變我們.

猶大並沒有及時終止賣耶穌的計劃,他落入一種週密部署,明知故犯的處境當中,面對耶穌洗腳的服侍,並多番提點,猶大仍然裝作若無其事,使徒約翰說是魔鬼控制了猶大的心(約13:2,現代中文譯本),或者,我們可以換一個角度,其實是猶大出賣了自己,讓邪惡控制自己的生命.

耶穌對猶大說 :「你所作的,快去作罷!」(約13:27)耶穌坦然面對,被一個同宿共膳三年的門徒出賣,祂沒有強行阻止猶大,更沒有意圖躲避眼前這殘忍的結果,於是,猶大受了那點餅,立刻就出去. (約13:30) 約翰語帶雙關的提醒讀者「那時候是夜間了.」猶大這一次離開之後,就走進無盡的黑夜裡,自以為按自己計劃行事,靠攏祭司長這權力核心,卻把自己推往絕望的深淵.

我一直不敢把猶大說得太壞.豈不是嗎?我們在人生旅途徘徊,何曾沒有試過出賣自己,不忠於所信,行事不符聖道,甚至出賣別人?我們豈沒有試過忘記上帝,或者訂下明知不符聖經教導的議程?正是我們這種罪性,出賣了耶穌,也正因如此,耶穌要為著那些出賣他的人,走往各各他.

總結

我準備這篇講章的時候,哈維爾的「無權力者的力量」The Power of the Powerless很自然就出現在我的腦海.耶穌――一個沒有實質權力的拉比,以無權者的力量,為門徒立下洗腳的榜樣,接受弟子的出賣,選擇被當權者以假見證送上十字架,然後成就救恩；而猶大卻選擇與謀害自己師傅的當權者合作,試圖擁有因為靠攏權力而得到的好處.耶穌與猶大走的是完全相反的道路.

今天這個社會,有很多權力交易,意圖靠攏權力去獲取好處,即使要說謊話,違背真理,甚至違反常理,仍然趨之若鶩.親愛的弟兄姊妹,就在這預苦期的第一個主日,我們看見聖壇的聖幃上的釘和荊棘編成的冠冕,我們就加倍的思念受苦的基督.祂問跟從祂的人「我向你們所作的,你們明白麼?」,你如何作答? 祂吩咐跟從祂的人「我給你們作了榜樣,叫你們照著我向你們所作的去作.」,你又如何回應? 弟兄姊妹,答案是十分清楚的 : 「耶穌為猶大洗腳」,然後,祂吩咐門徒「你們也當彼此洗腳」.願我們在這紫色的聖幃面前,預表著這位走在受苦和捨己道路的尊貴主,成為我們所景仰,所效法和效忠的對象；或者如此,我們會更樂意成為一個為別人洗腳的人.



\newpage

\section{上帝的家在人間了}
\label{sec:448}
\textbf{鄧瑞強博士 - 上帝的家在人間了}
\newline
\newline
\hyperref[sec:447]{< < < PREV SERMON < < <}
~
\hyperref[sec:toc]{[返目錄]}
~
\hyperref[sec:449]{> > > NEXT SERMON > > >}


啟示錄 21:1-27
\newline
\begin{longtable}{cl}
\hline
\hline
章節 & 經文 (和合本修訂版)\\
\hline
21:1 & \begin{tabularx}{0.7\textwidth}{X} 我又看見一個新天新地，因為先前的天和先前的地已經過去了，海也不再有了。 \end{tabularx} \\ \\ \relax
21:2 & \begin{tabularx}{0.7\textwidth}{X} 我又看見聖城，新耶路撒冷由神那裡，從天而降，預備好了，就如新娘打扮整齊，等候丈夫。 \end{tabularx} \\ \\ \relax
21:3 & \begin{tabularx}{0.7\textwidth}{X} 我聽見有大聲音從寶座出來，說：「看哪，神的帳幕在人間！他要和他們同住，他們要作他的子民。神要親自與他們同在。 \end{tabularx} \\ \\ \relax
21:4 & \begin{tabularx}{0.7\textwidth}{X} 神要擦去他們一切的眼淚；不再有死亡，也不再有悲哀、哭號、痛苦，因為先前的事都過去了。」 \end{tabularx} \\ \\ \relax
21:5 & \begin{tabularx}{0.7\textwidth}{X} 那位坐在寶座上的說：「看哪，我把一切都更新了！」他又說：「你要寫下來，因為這些話是可信靠的，是真實的。」 \end{tabularx} \\ \\ \relax
21:6 & \begin{tabularx}{0.7\textwidth}{X} 他又對我說：「成了！我是阿拉法，我是俄梅戛；我是開始，我是終結。我要把生命的泉水白白賜給那口渴的人喝。 \end{tabularx} \\ \\ \relax
21:7 & \begin{tabularx}{0.7\textwidth}{X} 得勝的要承受這些為業；我要作他的神，他要作我的兒子。 \end{tabularx} \\ \\ \relax
21:8 & \begin{tabularx}{0.7\textwidth}{X} 至於膽怯的、不信的、可憎的、殺人的、淫亂的、行邪術的、拜偶像的和一切說謊話的人，他們將在燒著硫磺的火湖裡有份；這是第二次的死。」 \end{tabularx} \\ \\ \relax
21:9 & \begin{tabularx}{0.7\textwidth}{X} 拿著七個金碗、盛滿末後七種災禍的七位天使中，有一位來對我說：「你來，我要給你看新娘，就是羔羊的妻子。」 \end{tabularx} \\ \\ \relax
21:10 & \begin{tabularx}{0.7\textwidth}{X} 我在聖靈感動下，天使帶我到一座高大的山，給我看由神那裡、從天而降的聖城耶路撒冷， \end{tabularx} \\ \\ \relax
21:11 & \begin{tabularx}{0.7\textwidth}{X} 這城有神的榮耀，它光輝如同極貴的寶石，好像碧玉，明如水晶。 \end{tabularx} \\ \\ \relax
21:12 & \begin{tabularx}{0.7\textwidth}{X} 它有高大的牆，有十二個門，門上有十二位天使，門上又寫著以色列人十二個支派的名字。 \end{tabularx} \\ \\ \relax
21:13 & \begin{tabularx}{0.7\textwidth}{X} 東邊有三個門，北邊有三個門，南邊有三個門，西邊有三個門。 \end{tabularx} \\ \\ \relax
21:14 & \begin{tabularx}{0.7\textwidth}{X} 城牆有十二個根基，根基上有羔羊十二使徒的名字。 \end{tabularx} \\ \\ \relax
21:15 & \begin{tabularx}{0.7\textwidth}{X} 那對我說話的天使拿著金的蘆葦當尺，要量那城、城門和城牆。 \end{tabularx} \\ \\ \relax
21:16 & \begin{tabularx}{0.7\textwidth}{X} 城是四方的，長寬一樣。天使用蘆葦量那城，共有一萬二千斯他迪，長、寬、高都是一樣。 \end{tabularx} \\ \\ \relax
21:17 & \begin{tabularx}{0.7\textwidth}{X} 他又量了城牆，按著人的尺寸，就是天使的尺寸，共有一百四十四肘。 \end{tabularx} \\ \\ \relax
21:18 & \begin{tabularx}{0.7\textwidth}{X} 牆是碧玉造的；城是純金的，如同明淨的玻璃。 \end{tabularx} \\ \\ \relax
21:19 & \begin{tabularx}{0.7\textwidth}{X} 城牆的根基是用各樣寶石修飾的：第一個根基是碧玉，第二是藍寶石，第三是綠瑪瑙，第四是綠寶石， \end{tabularx} \\ \\ \relax
21:20 & \begin{tabularx}{0.7\textwidth}{X} 第五是紅瑪瑙，第六是紅寶石，第七是黃璧璽，第八是水蒼玉，第九是紅璧璽，第十是翡翠，第十一是紫瑪瑙，第十二是紫晶。 \end{tabularx} \\ \\ \relax
21:21 & \begin{tabularx}{0.7\textwidth}{X} 十二個門是十二顆珍珠；每一個門是一顆珍珠造的。城內的街道是純金的，好像透明的玻璃。 \end{tabularx} \\ \\ \relax
21:22 & \begin{tabularx}{0.7\textwidth}{X} 我沒有看見城內有殿，因主—全能者神和羔羊就是城的殿。 \end{tabularx} \\ \\ \relax
21:23 & \begin{tabularx}{0.7\textwidth}{X} 那城內不用日月光照，因為有神的榮耀光照，又有羔羊為城的燈。 \end{tabularx} \\ \\ \relax
21:24 & \begin{tabularx}{0.7\textwidth}{X} 列國要藉著城的光行走；地上的君王要把自己的榮耀帶給那城。 \end{tabularx} \\ \\ \relax
21:25 & \begin{tabularx}{0.7\textwidth}{X} 城門白晝總不關閉，在那裡沒有黑夜。 \end{tabularx} \\ \\ \relax
21:26 & \begin{tabularx}{0.7\textwidth}{X} 人要將列國的榮耀尊貴帶給那城。 \end{tabularx} \\ \\ \relax
21:27 & \begin{tabularx}{0.7\textwidth}{X} 凡不潔淨的，和那行可憎與虛謊之事的人，都不得進那城，只有名字寫在羔羊生命冊上的才得進去。 \end{tabularx} \\ \\
[1ex]
\hline
\hline
\end{longtable}
講題:上帝的家在人間了 The Dwelling of God is with Men

經文:啟示錄21章1-27節

講員:鄧瑞強博士

場合:香港中文大學崇基學院禮拜堂

日期:2019年3月17日

各位弟兄姊妹,早安.

講了啟示錄很久了,是時候結束了.

啟示錄這封長長的信,是如何結束的呢?

啟示錄若是一封情信,它就是訴說一對戀人,分隔兩地,互相牽掛的故事.那女子是一即將出嫁的新娘,但新郎一時之間未能到來娶她.女子在窗前候盼,生活孤苦,被人欺負,憑藉的,只是從新郎那裡送來的叮嚀,鼓勵.新郎說,他在途中.新娘說,願你快來.苦等的新娘,最難過的是相思.這相思如何了斷,正是新郎的到來.這情信,以戀戀相依作結.

啟示錄若是一封君主與將軍之間的戰地書信,它就是訴說這場戰役的慘況.將軍在守孤城,這城是最重要的關口,守護蒼生,猶如真理守護著虛偽的人間.敵軍重重包圍著,危機四伏.君主知道險情,來信加以鼓勵,並說救兵已發,即日可達.這困局如何開脫,正是天兵的到達.君主來信告之,御駕親臨.

啟示錄若是學者討論歷史世情的學術書信,它就是陳述歷史中的種種磨難,訴說帝國對世人的欺凌,討論朝代興亡的大道.這討論如何作結,正是對「正義」的盼望.這書信,以看見證義在人間彰顯作為圓滿結局.

今日,我們看看,啟示錄這封長長的信,是如何結束的?

今日的講道經文是:啟示錄21: 1-27.

由於經文很長,我不先讀一次,只能邊讀邊解.

啟21:1　我又看見一個新天新地；因為先前的天地已經過去了,海也不再有了.

這裡講到「新天新地」,一個全新的世界.這「新天新地」的主要特徵,是不再有「海」.海,代表動盪,也代表爭鬥不絕的政治,這一切,不再有了.在啟示錄裡,「海」裡有怪獸.「海」是怪獸的溫床,「海」是邪惡出現的條件,現在不再有了.在新約,提到「海」,通常是指「地中海」,這是羅馬帝國的勢力範圍.羅馬帝國,是當時奴役人的暴力集團,現在沒有了.

再沒有奴役,再沒有暴力,再沒有邪惡,一個新天新地.這是人類未曾見過的,故謂之曰「新」.

啟21:2　我又看見聖城新耶路撒冷由神那裡從天而降,預備好了,就如新婦妝飾整齊,等候丈夫.

在這新天新地裡,有一從天而來的新城巿,新的聖城,新耶路撒冷.她像什麼?她像新娘,等候來娶她的丈夫.在啟示錄19:7裡,早已指出耶穌就是「新郎」,而忠信的信徒就是「新娘」.「聖城」指向「信徒」.我們等一會,再回到這意象.

啟21:3　我聽見有大聲音從寶座出來說:「看哪,神的帳幕在人間.他要與人同住,他們要作他的子民.神要親自與他們同在,作他們的神.

啟21:4　神要擦去他們一切的眼淚(這經文的直譯:神要擦去他們眼睛流出來的每一滴眼淚)；不再有死亡,也不再有悲哀,哭號,疼痛,因為以前的事都過去了.」

這新城耶路撒冷,就是神在人間的帳幕,神在人間的住處.新城耶路撒冷,是神與人永遠同在的居所.我們的神,喜歡走向人間.耶穌降生人間,正好說明這事實.馬太福音說,耶穌被稱為「以馬內利」,這詞的意思是「神與人同在」.在啟示錄,在新城耶路撒冷裡,「神與人同在」得到終極實現.神在,我們流的每一滴眼淚,都被擦乾了.死亡,悲哀,哭號,疼痛,都過去了.

啟21:5　坐寶座的說:「看哪,我將一切都更新了!」又說:「你要寫上；因這些話是可信的,是真實的.」

啟21:6　他又對我說:「都成了!我是阿拉法(alpha,希臘文的第一個字母),我是俄梅戛(omega,希臘文的最後一個字母)；我是初,我是終.我要將生命泉的水白白賜給那口渴的人喝.

啟21:7　得勝的,必承受這些為業:我要作他的神,他要作我的兒子.

啟21:8　惟有膽怯的,不信的,可憎的,殺人的,淫亂的,行邪術的,拜偶像的,和一切說謊話的,他們的分就在燒著硫磺的火湖裡；這是第二次的死.」

神是創始成終的主,他將歷史帶到終點.在終點,是對人生的判定.得勝的,活出真理的,他們要作神永恆的兒女.邪惡的,敵擋真理的,他們遭受永遠的消滅.不要說,天無眼.神看著人間,他判定是非.

啟21:9　拿著七個金碗,盛滿末後七災的七位天使中,有一位來對我說:「你到這裡來,我要將新婦,就是羔羊的妻,指給你看.」

啟21:10　我被聖靈感動,天使就帶我到一座高大的山,將那由神那裡,從天而降的聖城耶路撒冷指示我.

「羔羊」就是基督,「羔羊的妻」就是忠信的信徒.天使要向作者指出「羔羊的妻」,最後,指出的,是「新城耶路撒冷」.再一次,經文指出,在歷史終點裡要顯明的,是:信徒的生命,就是在人類盼望中的「新聖城」.

啟21:11　城中有神的榮耀；城的光輝如同極貴的寶石,好像碧玉,明如水晶.

啟21:12　有高大的牆,有十二個門,門上有十二位天使；門上又寫著以色列十二個支派的名字.

啟21:13　東邊有三門,北邊有三門,南邊有三門,西邊有三門.

啟21:14　城牆有十二根基,根基上有羔羊十二使徒的名字.

「新城耶路撒冷」,城門上寫的,是神在舊約裡12支派的名字,城的根基上寫的,是神在新約裡基督的12門徒的名字.又一次,經文指出:「新城耶路撒冷」就是神的子民.

啟21:15　對我說話的,拿著金葦子當尺,要量那城和城門城牆.

啟21:16　城是四方的,長寬一樣.天使用葦子量那城,共有四千里(更好的翻譯是:12000 stadion,「斯塔德」,重點在12000.用現在的講法,這是2220公里.由香港到北京的距離,約為2000公里.12000 stadion 約是香港到北京的距離再加百分之10),長,寬,高都是一樣(若很形象的想像,這是一個巨大的正立方體城巿)；

啟21:17　又量了城牆,按著人的尺寸,就是天使的尺寸,共有一百四十四肘(城樓的高度是65公尺,大約是香港一般住宅樓宇20層左右).

啟21:18　牆是碧玉造的；城是精金的,如同明淨的玻璃.

啟21:19　城牆的根基是用各樣寶石修飾的:第一根基是碧玉；第二是藍寶石；第三是綠瑪瑙；第四是綠寶石；

啟21:20　第五是紅瑪瑙；第六是紅寶石；第七是黃璧璽；第八是水蒼玉；第九是紅璧璽；第十是翡翠；第十一是紫瑪瑙；第十二是紫晶.

啟21:21　十二個門是十二顆珍珠,每門是一顆珍珠.城內的街道是精金,好像明透的玻璃.

請不要問:為何精金可以似透明的玻璃.這是神學書寫,「精金」代表王者的尊貴,得勝的信徒與基督一同作王,讓大地回復正義.在啟示錄第15章,講到忠信的信徒站在「玻璃海上」,唱勝利之歌.故此,他們站的地,就像「明透的玻璃」.

在這段經文裡,「新城耶路撒冷」最令人深刻的地方,有兩點.第一點,是它與12000和144相關.很明顯,這兩組數字,和12有關,這就是神子民相關的數字.我們記得,在啟示錄,神子民的數目,就是14萬4千,就是144.再一次,經文指出,這裡講的「新城耶路撒冷」,就是指信徒本身.「新城耶路撒冷」從天而降,這顯明忠信的信徒,他們的生命是屬天而不屬地的.

第二點,「新城耶路撒冷」是一正立方體.在聖經裡,唯一一個正立方體的建築物,就是聖殿裡最神聖的地方:「至聖所」(參:王上6: 20).很明顯,啟示錄想講,「新城耶路撒冷」就是新的「至聖所」.「至聖所」正是神臨在的地方.

來到這裡,啟示錄的用意已很明顯,它想指出,在新天新地裡,核心是「聖城新耶路撒冷」,這聖城其實不是「空間意義」下的城,這聖城其實就是信徒本身.這就是神住在人間的帳幕,我們的生命本身就是神臨在的「至聖所」.

啟21:22　我未見城內有殿,因主神─全能者和羔羊為城的殿.

啟21:23　那城內又不用日月光照；因有神的榮耀光照,又有羔羊為城的燈.

啟21:24　列國要在城的光裡行走；地上的君王必將自己的榮耀歸與那城.

啟21:25　城門白晝總不關閉,在那裡原沒有黑夜.

啟21:26　人必將列國的榮耀,尊貴歸與那城.

啟21:27　凡不潔淨的,並那行可憎與虛謊之事的,總不得進那城；只有名字寫在羔羊生命冊上的才得進去.

在「新耶路撒冷」裡,沒有聖殿,神親自為「殿」.再一次,以一種「生命存在」來作為一建築物.神為「殿」,神完全臨在人間,萬分神聖,光芒萬象.

忠信的信徒,是「至聖所」,神住在「至聖所」裡,即:神住在信徒生命中.同時,神是「聖殿」,信徒是「至聖所」,「至聖所」位於「聖殿」裡,信徒位於神裡面,即:信徒住在神的生命中.如此,神住在信徒生命中,信徒住在神的生命中.這就是約翰福音17: 21裡,耶穌向天父發出的禱告,他說:「……父,你在我裡面,我在你裡面,使他們也在我們裡面,叫世人可以信你差了我來.」這就是神與人的生命的終極連合,神在我們裡面,我們在神裡面.

至此,啟示錄想說什麼?啟示錄想表達,人間悠長的歷史,正是要表明,神多麼渴望住在人的生命中,透過信徒,實現永恆的正義與和平.我們在歷史中的使命,是努力住在神裡面,努力以神為家,好叫我們的生命,成為聖潔的生命,成為與世人不一樣的生命,爆發出神聖的光明.這光明像流星,劃過漆黑的夜空,讓人仍能想望未來.在歷史的漫漫長夜裡,信徒的生命點燃著光明.在謊話的世界,堅持真理.在暴力的世界,堅持和平.在絕望的世界,堅持希望.在世上,信徒做的,就是這種堅持,那怕多麼渺少.就是這種堅持,世人看見何謂神,何謂福音,何謂希望.啟示錄好想我們好好活著,以致世人在我們生命中,看見神.

讓我對一整系列的啟示錄講道,作一總結.

啟示錄是一封信,寫給七個教會,讓當中的信徒讀.這群信徒可能只是想簡簡單單信耶穌,享受一下愛的友誼,在人間多一點支援.但是,只要對信仰認真,就發現這信仰要求著不一樣的生命.在人間,有怪獸一樣的暴政,逼使你助紂為虐.有以「大淫婦」為象徵的奢華生活,這生活形態吸引著人,叫人沉醉在肉慾的歡愉裡.有各式各樣的偶像,吸引人走向虛假的人生.這時,一個信徒應該如何?不同流合污,將失去很多利益,甚或有殺身之禍.神在嗎?神看見嗎?神對此無能為力嗎?

啟示錄幫助信徒,仰視上蒼,看到在那裡,神高高地坐著為王.看到天父差聖子耶穌,下凡救世.看到耶穌的方式,是在種種反對真理的情景裡,堅持真理.在死亡的威嚇下,不因死亡而放棄.看到耶穌以十字架的方式,去克服死亡和邪惡的威迫.耶穌以天父為家,活出了不一樣的人生.他呼喚我們,像他一樣,以天父為家,好叫我們活的一生,能讓人看到,天父活在我們生命中.生命的沉淪與得勝,黑暗與光明,我們都看見了.讓我們好好選擇人生.

但願:榮耀歸於聖父,聖子,聖靈.阿們.



\newpage

\section{五十步笑百步}
\label{sec:449}
\textbf{林豪恩先生 - 五十步笑百步}
\newline
\newline
\hyperref[sec:448]{< < < PREV SERMON < < <}
~
\hyperref[sec:toc]{[返目錄]}
~
\hyperref[sec:450]{> > > NEXT SERMON > > >}


路加福音 13:1-9
\newline
\begin{longtable}{cl}
\hline
\hline
章節 & 經文 (和合本修訂版)\\
\hline
13:1 & \begin{tabularx}{0.7\textwidth}{X} 正當那時，有些在場的人把彼拉多使加利利人的血攙雜在他們祭物中的事，告訴耶穌。 \end{tabularx} \\ \\ \relax
13:2 & \begin{tabularx}{0.7\textwidth}{X} 耶穌對他們說：「你們以為這些加利利人比其他的加利利人更有罪，所以受這害嗎？ \end{tabularx} \\ \\ \relax
13:3 & \begin{tabularx}{0.7\textwidth}{X} 我告訴你們，不是的！你們若不悔改，都同樣要滅亡！ \end{tabularx} \\ \\ \relax
13:4 & \begin{tabularx}{0.7\textwidth}{X} 從前西羅亞樓倒塌，壓死了十八個人，你們以為那些人比一切住在耶路撒冷的人更有罪嗎？ \end{tabularx} \\ \\ \relax
13:5 & \begin{tabularx}{0.7\textwidth}{X} 我告訴你們，不是的！你們若不悔改，都照樣要滅亡！」 \end{tabularx} \\ \\ \relax
13:6 & \begin{tabularx}{0.7\textwidth}{X} 於是，耶穌用比喻說：「有一個人在葡萄園裡栽了一棵無花果樹。他前來在樹上找果子，卻找不到， \end{tabularx} \\ \\ \relax
13:7 & \begin{tabularx}{0.7\textwidth}{X} 就對園丁說：『看哪，我這三年來到這棵無花果樹前找果子，竟找不到。把它砍了吧，何必白佔土地呢？』 \end{tabularx} \\ \\ \relax
13:8 & \begin{tabularx}{0.7\textwidth}{X} 園丁回答：『主啊，今年且留著，等我在樹周圍掘開土，加上肥料， \end{tabularx} \\ \\ \relax
13:9 & \begin{tabularx}{0.7\textwidth}{X} 以後若結果子便罷，不然再把它砍了。』」 \end{tabularx} \\ \\
[1ex]
\hline
\hline
\end{longtable}
講題:五十步笑百步 The Pot Calling The Kettle Black

講員:林豪恩先生

經文:路加福音13:1-9

場合:香港中文大學崇基學院禮拜堂

日期:2019年3月24日

引言

三月六日星期三(聖灰日)開始了四十天的預苦期,提醒我們藉著真誠的懺悔與默想,準備迎向耶穌基督的受死與復活.今天是預苦期第三主日,預苦期的第十九日,即是差不多過了一半,我們的靈性有否跟隨主耶穌走近十字架的腳步前行呢?預苦期第三主日的福音經課選自路加福音13:1-9節,讓我們再閱讀經文,細味耶穌這一步的足跡.

二．釋經與應用

1．他們是「抵死」的

13:1正當那時,有人將彼拉多使加利利人的血攙雜在他們祭物中的事告訴耶穌.

「正當那時」是甚麼場景呢?按12章第1節描述,當時有成千上萬的人聚集,人多擠壓,甚麼彼此踐踏.耶穌開講,所講的內容充滿審判的味道,呼籲聽眾自我省察.例如:

12:3你們在暗中所說的,將要在明處被人聽見；在內室附耳所說的,將要在房上被人宣揚.

12:20上帝卻對(田產豐盛的財主)說:『無知的人哪,今夜必要你的靈魂；你所預備的要歸誰呢?』

12:49我來要把火丟在地上,倘若已經著起來,不也是我所願意的嗎?

12:56假冒為善的人哪,你們知道分辨天地的氣色,怎麼不知道分辨這時候呢?」

12:57「你們又為何不自己審量甚麼是合理的呢?

正當耶穌呼籲聽眾自省的那時,就有人把一些加利利人在獻祭或者往獻祭的途中被彼拉多殺害的事告訴耶穌.這位聽眾在耶穌與群眾對話的時候提出這事件,所謂何事呢?單純的理解,是回應耶穌所講關於審判的主題,為這主題提供例子.

13:2耶穌說:「你們以為這些加利利人比眾加利利人更有罪,所以受這害嗎?

耶穌承接著聽眾提出的例子,透過反問揭示這位聽眾的想法:耶穌,我同意你所說,有審判的.你看,那麼多加利利人都平安無事,只有那幾個人被殺,一定是那幾個人做了甚麼事而招來惡果.

「咁多人無事,只有你有事,一定係你有Ｄ唔妥」,這也是我們共同的想法.這種想法能夠令我們事情合理化,維持心理平衡及邏輯一致.讓我們舉一個沒有那麼沉重,甚至有點無聊的例子:如果像我這種高大及英俊都無分的男人,娶到一位漂亮女子為妻,「花生友」很可能會說:「他一定很有錢!」否則於理不合,「花生友」會感到理性上難以理解,心裡很不舒服.這種狀況,類似心理學稱為認知失調的狀況.

認知失調理論指出,當人們在經歷兩種截然對立的「認知」(也就是想法,信念,意見或觀察)時,所產生的一種心理上的緊張狀態.為了緩解或移除這種緊張狀態及其所帶來的不適,人們必須做出某種內在調解,使原本相互牴觸的認知,重新歸為一統.使自己的認知恢復協調一致的過程非常人性化,也就是說,凡身為人皆有此一需求,因為沒有人可以承擔認知長期失衡所造成的痛苦.

這種要求認知協調的負面例子,除了那位聽眾提出的「彼拉多使加利利人的血攙雜在他們祭物中的事」之外,耶穌順著他們的思路,還提出了另外一項:

13:4從前西羅亞樓倒塌了,壓死十八個人；你們以為那些人比一切住在耶路撒冷的人更有罪嗎?

那麼多人住在耶路撒冷,只是這十八個人被塌樓壓死,為甚麼呢?按著聽眾的邏輯,就是因為「那些人比一切住在耶路撒冷的人更有罪」.沒有說出來的潛台詞是:「他們抵死!」

雖然要求認知協調是人的自然需要,但是,這種再基本不過的內在人性需求卻大有可能導致各式各樣泯滅人性的災難性後果,因為這樣做等於把不幸合理化,甚至有可能把邪惡行動正常化.

2004年發生南亞大海嘯,造成巨大人命及財物損失.有些人總要為災難找個道德原因.當時基督徒圈子之中流傳著一種言論,認為這場世紀大海嘯是上帝藉此審判「拜偶像」的國家(印尼有全球最多的回教徒,泰國和斯里蘭卡人民多信奉佛教,印度則是佛教發源地,而當地人民多信奉印度教).不知道這種言論是否追求認知協調的心理反應,還是藉此「抽水」護教,證明自己信奉的宗教蒙福,而印證拜偶像的人必如偶像一樣毀滅.無論是甚麼動機,面對聖靈塗炭的災難發出如此的言論,都是令人心寒的.

更邪惡的,是為對手的不幸製造污名化的標籤,打擊對手.有人猜想,那些被彼拉多殺死的加利利人,其實是革命份子,他們不服彼拉多的暴政,因此被彼拉多除掉.又有人猜想,那些在西羅亞樓倒塌中被壓死的是親建制人士,他們在耶路撒冷京城協助管治者壓逼人民,所以遭天遣.無論事情是否屬實,或者證據多少,發出如此言論的動機及效果,都是要製造「他們抵死」的定調.

「他們抵死」這定調,成為了多少種族屠殺的口號?也成為了多少政治運動中批鬥對手或異見者的理由?

2．我們不比他們好多少

耶穌如何回應人們為他關於審判的信息所提供的例子呢?

13:3我告訴你們,不是的!

13:5我告訴你們,不是的!

無論發出這些「怪罪受害者」的言論人著甚麼動機,耶穌勇敢為受難者澄澈:這些加利利人被殺,不是因為他們比眾加利利人更有罪；那些西羅亞樓倒塌壓死的十八個人,不是因為他們比其他住在耶路撒冷的人更有罪.

接著,耶穌勇敢地說出可能冒犯聽眾的事實:

13:3 你們若不悔改,都要如此滅亡!

13:5 你們若不悔改,都要如此滅亡!

耶穌向群眾宣布審判的信息,原不是要人用來互相審判,而是呼籲人自省及把握機會悔改.不過,人總是喜歡替別人聽道,而不是為自己聽道.

在網上看到有位牧師說:「有的時候我講完道,在門口送弟兄姊妹.有人說『啊,牧師,你今天講得好感動!我跟你講,我覺得好可惜,我先生沒來.這個道就應該是我先生聽的.』」

事實是:「我們若不悔改,都要如此滅亡!」事實是,我們不比他們好得多少,甚至是不比他們好.如果亡者能開口,他們可能會告誡生者:今日吾軀歸故土,他朝君體也相同.其實,生者也應該自覺:今日君體歸故土,他朝吾軀也相同.在滅亡的終極來看,我們和他們並無兩樣.在悔改的呼籲之下,我們和他們都是一樣的.

3．你知道你被恩待嗎?

在基督教信仰中,審判是重要主題,悔改也是重要主題,然而,如果只有審判及悔改,信仰就顯得機械化,也變得冷冰冰.朋友們,基督教信仰對你來說是不是機械化及冷冰冰的呢?其實,耶穌在這個事件中的話還未說完,他接下來說了一個比喻:

13:6於是用比喻說:「一個人有一棵無花果樹栽在葡萄園裡.

這個人有點奇怪,何以見得呢?那個年代,無花果樹比葡萄樹粗生,人們通常粗放栽種無花果樹；而葡萄樹需要較好土壤生長,也需要多加照料修剪.這個人用葡萄園來栽種無花果,顯示出他對該棵無花果特別栽培的用意.用另一個角度來看,這棵無花果樹得到的栽培和照顧,本不是它配得的,可以說這是恩典.

他來到樹前找果子,卻找不著.

13:7就對管園的說:『看哪,我這三年來到這無花果樹前找果子,竟找不著.把它砍了吧,何必白佔地土呢!』

一般來說,無花果每年都應該結果子,三年都沒有收成,真是白佔了地土.那棵無花果樹,對於自己被恩待,無動於衷.「把它砍了吧」,這是審判的警告.

13:8管園的說:『 主啊,今年且留著,等我周圍掘開土,加上糞；

13:9以後若結果子便罷,不然再把它砍了.』」

這白佔地土而不結果子,本應被砍掉的無花果樹,還有人出來為它求情,寬容時限,再給機會,而且願意花上更多心思及資源去栽培,這明顯又是恩典.

不知道當時的聽眾們有否領悟到,他們就是那些被寬容的人,他們就是那些活在恩典之中的人.他們仍然存留,沒有隨被彼拉多殺死的那幾位加利利人而去,沒有隨那十八位在西羅亞樓倒塌中壓死的而去,不是因為他們比離去者美好,而只是被恩典承托著而已.

審判,悔改,恩典,都是基督教信仰的重要主題,其實也是人生的重要主題.在基督教信仰中,審判是上主的位置,悔改的是自我要求,恩典是待人之道；不過,在日常生活中,我們常常把次序混亂,變成:我們坐在審判的位置,要求別人悔改,並且從沒有感到自己被恩待,也不會恩待別人.

誰會意識到自己活在恩典之中呢?可能是類似剛才講笑時提及的那位既不高大,又不英俊,無樓又無車,一事無成的男人,竟然有一女子願意與他長相廝守.最近從朋友分享的訊息看到一篇四格漫畫,內容是這樣的:第一格,甲男子向乙男子說:「我覺得我太太很愛我.」第二格,乙男子回應說:「你如此難相處,你老婆怎會愛你!」第三格,甲男子說:「正因如此,我才覺得她愛我.」第四格,甲男子解釋說:「愛是恆久忍耐,她肯忍我那麼耐,證明她很愛我.」意識到自己的不足或過錯,自知自己不比別人好得多少,可能會比較容易感受到自己活在恩典之中.

現實上,很多人都不容易感到自己活在恩典之中,其中一個可能是大家都不察覺自己的不足或過錯,或者不願意接受自己在事件中的參與及責任.有一本書名叫「錯不在我?」作者以深入淺出的文筆總結了過往認知失調研究的相關學術發現,同時更進一步列舉現實生活中的各種事例加以印證,小至夫妻間的相處,大至國與國之間的殺伐.

例如,一般人都會自我感覺良好,至多只有程度上的差別而已.所以,我們或多或少都會覺得自己聰明,可愛,善良,理性與值得尊敬.但當某人的行為牴觸了正面的自我概念或自我意象時(例如,我竟然打了我的老婆,怎麼可能?只有爛咖才會打女人啊!),就出現了認知失調的狀況.

為了解決認知失調,你可以有兩種選擇:一種是設身處地去同理老婆的感受,徹底反省,幡然悔悟,承認自己不是個好人,儘管一時難堪,卻能夠達到心靈上的自我解放,也將大有機會與老婆重修舊好,但這顯然不是一條好走的路(我怎麼可能不是一個好人?).

第二條路,則是將維持正面的自我意象(我是個高尚的人)列為最高指導原則,從而邁入一種自我辯護的迴圈:我這麼高尚,不可能打女人,更不會打老婆,但現在老婆真的被我打了,唯一的理由就是:她該被打,因為她不是一個好女人(她怎麼可以對我出言不遜?都是她先開始的!).

作者指出,「如果錯誤發生了,記憶會提醒我那是別人造成的」.於是動粗的老公開始重組他的記憶,老婆過往的溫柔可人全都遭到刪除,僅保留老婆各種負面的特質,配合自身固有的盲點與某種程度的自我暗示,恰好足以建構出一個全身佈滿污點的枕邊人形象.長此以往,施加在老婆身上的暴行只會愈演愈烈,因為關於她該被打的證據愈來愈多.(翁志遠:書評《錯不在我?》)

作者繼續引經據典地告訴讀者,自我辯護的動機是如何強大到足以讓人突然「發現」實際上從未存在的過去,並對之深信不移；讓司法體系源源不絕的產生冤案；以及讓殘忍無道的虐待戰俘事件與種族大屠殺永遠不會在人類歷史中消失.

這本書的簡要核心訊息呼應耶穌的教導,審判讓我們知罪,悔改是我們應有的回應,正如中國古人說「人非聖賢,孰能無過!過而能改,善莫大焉.」耶穌的教導比這本書多了重要一點,就是恩典,有機會悔改,也是恩典.我們是活在恩典之中,不是恐懼之中；我們在恩典中悔改,而不是在恐懼中改過.體昧恩典,讓我們產生追求美善的動力:有人珍惜我們,我們怎能不珍惜自己呢?有人對我們好,我們怎能自暴自棄呢?

總結

最後,讓我們以本主日的書信經課,哥林多前書第十章,在預苦期的懺悔與默想中警戒自己:

10:1弟兄們,我不願意你們不曉得,我們的祖宗從前都在雲下,都從海中經過,

10:2都在雲裡,海裡受洗歸了摩西；

10:3並且都吃了一樣的靈食,

10:4也都喝了一樣的靈水.所喝的,是出於隨著他們的靈磐石；那磐石就是基督.

10:5但他們中間多半是上帝不喜歡的人,所以在曠野倒斃.

10:6這些事都是我們的鑑戒,叫我們不要貪戀惡事,像他們那樣貪戀的；

10:7也不要拜偶像,像他們有人拜的.如經上所記:「百姓坐下吃喝,起來玩耍.」

10:8我們也不要行姦淫,像他們有人行的,一天就倒斃了二萬三千人；

10:9也不要試探主,像他們有人試探的,就被蛇所滅.

10:10你們也不要發怨言,像他們有發怨言的,就被滅命的所滅.

10:11他們遭遇這些事都要作為鑑戒,並且寫在經上,正是警戒我們這末世的人.

同時,也靠賴上主的恩典與保守繼續前行:

10:12所以,自己以為站得穩的,須要謹慎,免得跌倒.

10:13你們所遇見的試探,無非是人所能受的.上帝是信實的,必不叫你們受試探過於所能受的；在受試探的時候,總要給你們開一條出路,叫你們能忍受得住.



\newpage

\section{以何心為心為何所用}
\label{sec:450}
\textbf{鄭漢文博士 - 以何心為心?為何所用?}
\newline
\newline
\hyperref[sec:449]{< < < PREV SERMON < < <}
~
\hyperref[sec:toc]{[返目錄]}
~
\hyperref[sec:451]{> > > NEXT SERMON > > >}


講題:以何心為心?為何所用? Be Used by Whom Like the Heart of Whose

講員:鄭漢文博士

場合:香港中文大學崇基學院禮拜堂

日期:2019年4月14日

全心得力講好道

大家平安!我今日是今年初成為崇基榮譽助理校牧之後第一次講道,有些緊張.希望你們在禱告中紀念我的講道服事.

我為何會講道?故事從2007年十月講起,我參加了自己教會辦的以馬忤斯之旅,得到靈性復興後,立志以講道服事,希望能夠有天講道,講學和講笑三合一.其實2007年之前,已經有教會邀請我在教育主日講道,我當時認為我取得教育哲學博士,所以被邀請.我感到自己借了世間的榮譽,借了與神學無關的博士學位來講道,我理應謙卑接受神學訓練.所以,我2008年9月開始兼讀神學研究文學碩士,三年後完成此課程後的2011年9月開始兼讀神道學學士(同時開始在崇基禮拜堂實習,講道及生活),一讀七年,去年才完成.我是在崇基完成學士後,在中大完成碩士,然後往英國讀完博士之後十四年,再在中大讀完碩士,然後在崇基讀完學士,剛好一個循環(full circle),回到起點.

我今天在此告解(confession),我之前是「無牌講道」(preaching without license)的.但是今日正式「有牌講道」(licensed to preach),主若願意我會「有排講道」(life-long to preach).

我叫這十年的兼讀神學做「神學攻讀共十年」,我打算寫一本書叫《神攻十年》,分享在崇基讀神學的要訣心得.

在這「神攻十年」其實不是必修講道學,但是我選修了三科講道學.所以,我的講道是師出名門的,我是師承自李炳光,吳振智,伍渭文三位牧師,可是,他們的講道法完全是三種不同的講道法.

以前我講道,通常我會只選用三種講道法ABC其中一種加上我自己的變奏,已經夠用.今天,是我出任榮譽助理校牧的第一次開工大吉,我打算三者皆用上,並且加上我自己的變奏.

讓我先把我從他們所得的領受,化為口訣與大家分享,現在公告天下,然後你可以鑑賞一二:

1) 李炳光牧師講道法是:融入旁徵博引,聽道後留神一句.

2) 吳振智牧師講道法是:透入人間罪性,言必歸基督十架.

3) 伍渭文牧師講道法是:出入三代經課,串讀出更新生活.

好啦,我們開始,希望我能全心得力講好道.

「以何心為心?為何所用?」的三種可能生命型態

今天我講道題目是兩個問題,分別各有一個空心詞:以「何」心為心?為「何」所用?

我發現這兩個問題的「何」字可以三種可能性.

第一種可能答案,可能生命型態是:(L1)「 以私心為心,萬有為我所用.」

(all things are used by my self)

原來人,生來就是自我中心(egocentric),從自己的眼睛出發看世間事物(seeing things from one’s own eyes).換言之,你是無法不從你的角度看身邊的事物.

大家試一試,現在做一個實驗:看在你左邊或在右邊的朋友,你看到他甚麼?你只是看到你所看的一面,只是片面,你是無法同時看到他的另一面,你看到的他是無法全面.

就算是你的伴侶,你以為日子久了就會看得全面,剛好相反,越相處你會越看不透,有時莫名其妙.

自我中心不一定就是自私自利,在哲學術語上講,是好容易從「存有論上的自我中心」(ontological egocentricism)陷入「價值論上的自我主義」 (axiological egoism).

甚至有人性論認為:人生下來就有私心,所有身邊的人與物都是為我所用.

俗語說:他這樣做,有何居心?

在社會科學解釋人的行為時,首先會用自利來解釋,即是說人之所以「做A而不做B」,是因為「A比B對他更有利」.

所以,教育的功能之一就是道基督教育,我在中大教育學院教了二十五年,每年都會開教道基督教育這門課.道基督教育,就是讓人學會怎樣面對自己的私心.

特別在無人知識自己的進退取捨時,如何面對一己的利慾之心?如何不是在心底裡不是在利用他人?中國古代傳統叫「君子慎獨」.

不自私的人有兩種,一種是不利用人以滿足自己,另一種是懷有公心.

有些行為,真的無法用自己得益來解釋的.他這樣做,有何居心?

有些行為,真的對自己不利也會做的,是因為懷有公心,作用於更大的世界.

大家都知我剛剛教完「行傳六講」的主日學,我們選用「時,地,人,事」四個角度拆解28章的《使徒行傳》,其中「人物志」是一條線索,當中我們選了三位使徒來思考:

在3月3日,10日,17日的第二至四講,我們探討了《使徒行傳》中的第二個「三一」(Trinity),就是三位重要人物所構成的一體,就是彼得,司提反,保羅三位聖靈同行的「學神」,看看他們的生命如何和為甚麼去到盡.成為一個簡單的人物志:

人物志第一回,我們第二講講彼得,探討「聖道讓你的生命不一樣?」

人物志第二回,我們第三講講司提反,探討「你的生命為甚麼去得盡?」

人物志第三回,我們第四講講保羅,探討「你如何由掃羅變成保羅?」

《使徒行傳》是講甚麼的?是講:彼得,司提反,保羅等早期教會人物,如何因為耶穌基督而人生得以改變,成為早期教會的信徒領袖,成為我們後世信徒的榜樣.

我們提到他們三個信徒的生命見證了三種「去到盡」的生命境界:

(一)從彼得,我們體會到第一種去到盡的生命,就是讓聖靈充滿,從而生命可以去到盡高的水平,讓聖道因我們的言語而臨在,讓其他生命也可以去到盡,造就教會.

(二)從司提反,我們體會到第二種去到盡的生命,就是讓聖靈充滿,從而生命可以去到盡頭的水平,讓聖道因我們的見證而臨在,成為信群的領袖,結果可能要殉道,為主而死.

(三)從保羅,我們體會到第三種去到盡的生命,就是讓聖靈充滿,從而生命可以去到盡轉的水平,讓我們生命任何時間都可以為主而作180度的全轉向──完全的大逆轉,從而為神踏上不同使命的人生路,並且在種種限制中依然自由地為神作見證.

保羅為何放棄他在猶太宗教群體的尊貴地位,為何願意成為羅馬階下囚?

有些行為,真的無法用自己得益來解釋的.許多人都不明白:

明明本來是可以安樂讀書的學生,為何要犯法被判有罪而甚至入獄?

明明本來是可以安樂教書的教授,為何要犯法被判有罪而甚至入獄?

明明本來是可以安樂退休的牧師,為何要犯法被判有罪而甚至入獄?

明明本來是可以安樂議政的議員,為何要犯法被判有罪而甚至入獄?

所以,可以出現第二種人生型態,就是:(L2)「 以公心為心,萬有經我而有用.」

(all things are used through my self but for the public good)

根據這種觀點看,人可以出於公心,所作行為是為了一個更美好的公共世界.

教育不就是要培養為世所用的人才?

人類的文明,發展出不少觀念,可以用作描繪出這個人值得無私地付出以有用於產生更美好的公共世界.

例如英文Justice這個理念,在漢語就有三個不同的中文對應詞:

一是公義,是社會對任何人都有適當的對待方式,不因其地位而有特權.

二是公正,是社會制度對所有人都有重視同等權益的程序安排,不因其地位而有差異.

三是公道,是社會資源分配較為合情合理,不會遍向權貴階層.

四是正義,……夠了夠了,否則這次講道就變成社會哲學(Social Philosophy)入門課.

有些無私行為,真的無法用自己得益來解釋的.

唯一解釋是:他們在追求正義,民主,法治等等,是讓社會變得更美好,在這個更美好了的社會下生活的所有人,都同時受惠得益.

但是,為何是因為他們無私行動的人要入獄,反而是我們受惠得益?

想一想,耶穌基督也是無私地被釘死在十字架上,反而是我們受惠得益?

有些人解釋說:因為他們追求榮譽,希望自己歷史留名.

的而且確,人的追求除了利益外,是會追求榮譽.

你看,我也是個榮譽的助理校牧,不過這樣冠名,其實指我無受薪,我只是義工一名.

但想深一層,人是可以自義(self righteous)的,自以為義.

在我少年成長時,在座堂教堂的高處,寫著「榮歸上主」(Glory to the Lord),我是時常對著這個橫牌祈禱的.念念不忘,必有迴響.讓我明白,人原來可以追求事事「榮歸自己」(Glory to oneself).

年日漸長,開始明白,原來有些人可以假借以公心為名,實際上「萬有經我而用,榮歸自己(為我榮譽之名所用)」.從此,每次進退取捨,我都反問一下自己,是否「榮歸自己」?

因為世間有「榮歸自己」的人,更體會「榮歸上主」的生命可能型態,正是我們基督徒所追求的境界!

由此引出第三種生命型態,就是:(L3)「以基督的心為心,萬有為主何用」

(為我上主之名所用)(all things are used by the Lord through my heart like the heart of Christ)

甚麼叫以基督的心為心?

基督事件的前中後:三段式基督頌

上述三種生命型態,引出今日所要同大家思考的第一段經課,是取自腓立比書2章5至11節.是關於「基督事件」的一篇三段式的基督頌(Hymn on Christ, Christ Event).

我是16歲開始在教會成長,每次聽到這段經文,我都會好感動,莫明地哭泣起來.

請大家打開聖經,今天請聽我讀出:

5 你們當以基督耶穌的心為心:

6 他本有上帝的形像,

不以自己與上帝同等為強奪的；

7 反倒虛己,

取了奴僕的形像,

成為人的樣式；

8 既有人的樣子,就自己卑微,

存心順服,以至於死,

且死在十字架上.

9 所以,上帝將他升為至高,

又賜給他那超乎萬名之上的名,

10 叫一切在天上的,地上的,和地底下的,

因耶穌的名無不屈膝,

11 無不口稱「耶穌基督為主」,

使榮耀歸與父上帝.

腓立比書這段文字好美麗,用了敘事體,分前中後三段,借用哲學術語,分解如下:

第一段5至6節是講基督的本然,起始上「他本有上帝的形像,與上帝可以同等.」(alpha: same as God)

第二段7至8節是講基督的願然,中段說他甘願「反倒虛己,取了奴僕的形像,成為人的樣式；既有人的樣子,就自己卑微,存心順服,以至於死,且死在十字架上.」(humble Himself by taking an human form)

第三段9至11節是講基督的終極實然,最後是「 上帝將他升為至高,又賜給他那超乎萬名之上的名,因耶穌的名無不屈膝,無不口稱「耶穌基督為主」,但最重要是使榮耀歸與父上帝.」 (omega: Glory to God)

甚麼叫「以基督的心為心」?(let my heart like the heart of Christ)

本然上,就是本來我們的生活可以好安樂,享受難得的福樂.

願然上,我們甘願放下自己的安樂而去助人,甚至為令得世界更美好而付出,以至於入獄.

然而到最後,終極實然是「因耶穌的名無不屈膝, 無不口稱『耶穌基督為主』,使榮耀歸與父上帝.」

這樣地「以基督的心為心」,你願意嗎?

不用擔心,今日另一經課詩篇118篇第1節和第29節,第一節和最後一節都重申本為善的上主會用我們得其所的:

1 你們要稱謝耶和華,因他本為善；他的慈愛永遠長存!

……

29 你們要稱謝耶和華,因他本為善；他的慈愛永遠長存!

耶穌基督極謙卑時也極榮耀,因為主要用他

耶穌在世時,無時無處不示範可以如何「以上主的心為心而為主所用」.

今日經課另一段取自《路加福音》19章28至40節.今天是棕枝主日,這段經文,是記載耶穌要進入耶路撒冷.

跟從他的信群以及在旁觀的人群以為他要做些大事,一些留名的大事,「威水事」(glorious act).

但對耶穌來說,他是要去送死,明知故犯地上路.中文語叫:「從容就義」.

這正是耶穌基督極謙卑也是極榮耀的事件,何解?因為「主要用他」.

這段經文有一個關鍵句子,落在第31節和34節,就是「主要用他」.

一頭從未被人騎過的驢駒,從他的主人手取過來用,是耶穌指明要用他,因為主要用他.

這裡出現三個「主」字.

首先,是在世的主人.我們好像驢駒,在世做奴如駒,被僱主驅策所用,終日生活無望.一方面又怕無用而被遺棄(made redundant),又怕有用但被用到死.

但是有一天,主耶穌出現成為勝過死亡的基督,我們得以解放,從世間主子的奴役,轉為跟隨耶穌作為我們生命的主.所以,當我們信耶穌是基督,是我們的救主,我們還要跟隨耶穌為主.合而稱之:「我信主耶穌基督」.(這是《使徒信經》的第二段,我正向校牧申請,下一次開主日學試探討《使徒信經》.)

然而,進一步,主耶穌說,「為主所用」有一重更高的意思,就是「為上主所用」,耶穌也是這樣地「為上主所用」而成為基督.上主就是天父上帝,一切榮耀都是歸給上帝.

以基督的心為心,效法基督,為上主所用.

如果這是一條苦路,你願意踏上嗎?

用言語扶助疲乏的人

耶穌成為基督,成為上主用作拯救世人的僕人,必先受苦,才能救世人於苦海.

這個「受苦僕人」(suffering servant)的意象,早在以賽亞書已經出現.

聖經學者認為,以賽亞書寫有四首僕人之歌,預表基督作為主耶和華的僕人,為主所用時的心.

講道的最後,我想引用以賽亞書50章4-9a節作結.

今日經課第50章的是第三首僕人之歌.

請聽我讀出,不過當我讀出時,除了耶穌基督外,請你也可想像,紀念一位現今正為主所用的僕人(包括我在內),可能是在香港的,也可能是在內地的:

(新標點和合本)

4 主耶和華賜我受教者的舌頭,

使我知道怎樣用言語扶助疲乏的人.

主每早晨提醒,

提醒我的耳朵,

使我能聽,像受教者一樣.

5主耶和華開啟我的耳朵,

我並未違背,也未退後.

6人打我的背,我任他打；

人拔我兩頰的鬍鬚,我由他拔；

人侮辱我,向我吐唾沫,我並不掩面.

7主耶和華必幫助我,

所以我不抱愧.

我硬著臉面好像堅石,

也知道我必不致蒙羞.

8稱我為義的與我相近；

誰與我爭論,

讓我們來對質；

誰與我作對,

讓他近前來吧!

9看哪,主耶和華必幫助我,

誰能定我有罪呢?

以上的頌讀,希望能扶助到為主行公義但疲乏的人.

至於我,我過去四十多年,都是按第50章第4節這一句話作為我志向,四十多年來這句話都隱含著上主給我的恩賜,讓我可以借講學,講笑與講道用來服侍人.

讓我與大家分享我四十多年來的領受,但我比較喜歡和合修訂本的中譯,請聽我讀出:

4 主耶和華賜我受教者的舌頭,

使我知道怎樣用言語扶助疲乏的人.

主每天早晨喚醒,喚醒我的耳朵,

使我能聽,像受教者一樣.

但願我的講道,跟我的講學與講笑,都是因言語的運用而扶助疲乏的人.

如果剛才我的言語能令得你可以在清晨補足你的睡眠,也算是對你疲乏的一次扶助.

在結束前,我想分享,這個星期在準備這一篇講章時,除了希望得自剛才提及三位牧師的真傳,還浮現了第四位牧師的講道法,所領受的是一種心法.

我是好想效法他,他是這樣說:

「我只是一個敲鐘者,希望發出警號,讓人們知道不幸和災難正在發生,期望喚醒人們的良知,共挽狂瀾……如果我仍有氣力,必繼續在教會敲鐘,在世上敲鐘,在人心敲鐘」.

阿們.



\newpage

\section{為甚麼在死人中找活人呢}
\label{sec:451}
\textbf{鄧瑞強博士 - 為甚麼在死人中找活人呢?}
\newline
\newline
\hyperref[sec:450]{< < < PREV SERMON < < <}
~
\hyperref[sec:toc]{[返目錄]}
~
\hyperref[sec:452]{> > > NEXT SERMON > > >}


路加福音 24:1-12
\newline
\begin{longtable}{cl}
\hline
\hline
章節 & 經文 (和合本修訂版)\\
\hline
24:1 & \begin{tabularx}{0.7\textwidth}{X} 七日的第一日，黎明的時候，那些婦女帶著所預備的香料來到墳墓那裡， \end{tabularx} \\ \\ \relax
24:2 & \begin{tabularx}{0.7\textwidth}{X} 發現石頭已經從墳墓滾開了， \end{tabularx} \\ \\ \relax
24:3 & \begin{tabularx}{0.7\textwidth}{X} 她們就進去，只是不見主耶穌的身體。 \end{tabularx} \\ \\ \relax
24:4 & \begin{tabularx}{0.7\textwidth}{X} 正在為這事困惑的時候，忽然有兩個人站在旁邊，衣服放光。 \end{tabularx} \\ \\ \relax
24:5 & \begin{tabularx}{0.7\textwidth}{X} 婦女們非常害怕，就俯伏在地上。那兩個人對她們說：「為甚麼在死人中找活人呢？ \end{tabularx} \\ \\ \relax
24:6 & \begin{tabularx}{0.7\textwidth}{X} 他不在這裡，已經復活了。要記得他還在加利利的時候怎樣告訴你們的， \end{tabularx} \\ \\ \relax
24:7 & \begin{tabularx}{0.7\textwidth}{X} 他說：『人子必須被交在罪人手裡，釘在十字架上，第三天復活。』」 \end{tabularx} \\ \\ \relax
24:8 & \begin{tabularx}{0.7\textwidth}{X} 她們就想起耶穌的話來。 \end{tabularx} \\ \\ \relax
24:9 & \begin{tabularx}{0.7\textwidth}{X} 於是她們從墳墓那裡回去，把這一切事告訴十一個使徒和其餘的人。 \end{tabularx} \\ \\ \relax
24:10 & \begin{tabularx}{0.7\textwidth}{X} 把這些事告訴使徒的有抹大拉的馬利亞、約亞拿，和雅各的母親馬利亞，還有跟她們在一起的婦女。 \end{tabularx} \\ \\ \relax
24:11 & \begin{tabularx}{0.7\textwidth}{X} 她們這些話，使徒以為是胡言，就不相信。 \end{tabularx} \\ \\ \relax
24:12 & \begin{tabularx}{0.7\textwidth}{X} 彼得起來，跑到墳墓前，俯身往裡看，只見細麻布，就回去了，因所發生的事而心裡驚訝。 \end{tabularx} \\ \\
[1ex]
\hline
\hline
\end{longtable}
講題:為甚麼在死人中找活人呢? Why Seek Ye the Living among the Dead?

講員:鄧瑞強博士

經文:路加福音24:1-12

場合:香港中文大學崇基學院禮拜堂

日期:2019年4月21日

各位弟兄姊妹,早安.

今天是復活節,我們應開心雀躍.漫長的冬夜已過,乍見溫暖的春光.罪惡的暴虐止住了,仁愛的花朵到處綻放.死亡不再肆無忌憚,生命的種子開始發芽.這是開心的日子.

古教父大馬士革的約翰(St. John of Damascus)如此說:

「現在,萬物充滿光明,無論是在天地之間,抑或深在大地之下.蒼生唱頌基督的復活.藉這復活,萬物得以堅立.」

(St. John of Damascus, “Now all things have been filled with light, both heaven and earth and those beneath the earth; so let all creation sing Christ’s rising, by which it is established.”)(\url{https://prabook.com/web/john.of_damascus/3733347})

古教父屈梭多模(St. John Chrysostom)的復活節講道如此說:

「請同來參與信仰的筵席,請同來領受美善的豐盈.請不要哀悼你的貧乏,因為天國已向你顯明.請不要再為罪行難過,因為從墳墓而來的是寬恕.請不要恐懼死亡,因為救主的死已讓我們飛翔.……死亡啊,你的毒鉤在哪裡?陰間啊,你得勝的權勢在哪裡?基督已復活,陰間傾覆了.基督已復活,邪魔失敗了.基督已復活,天使在歡欣.基督已復活,生命在掌權.基督已復活,墳墓清空了.」(John Chrysostom: “Let all partake of the feast of faith. Let all receive the riches of goodness. Let no one lament his poverty, for the universal kingdom has been revealed. Let no one mourn his transgressions, for pardon has dawned from the grave. Let no one fear death, for the Saviour’s death has set us free…… O death, where is thy sting? O hades, where is thy victory? Christ is risen, and you are overthrown! Christ is risen, and the demons are fallen! Christ is risen, and the angels rejoice! Christ is risen, and life reigns! Christ is risen, and not one dead remains in a tomb!”)

(\url{https://orthodoxwiki.org/Paschal_Homily})

另一古教父拿先斯的貴格利(St. Gregory of Nazianzus)如此說:

「昨日,我與他同釘十架；今日,我與他一同得榮耀.昨日,我與他同死；今日,我與他同活.昨日,我與他同埋葬；今日,我與他同復活.我們要向這位為我們受苦及復活的耶穌作奉獻.……讓我們奉獻自己,對神而言,這是最寶貴的,最合宜的.我們按神的形像受造,我們就向神奉獻這形像.我們要肯定〔這形像賦予給〕我們的尊嚴,我們要榮耀我們的形像的本源,我們要領悟這奧秘的大能,我們要明白基督為了什麼而死.」

(St. Gregory of Nazianzus, “Yesterday I was crucified with Him; today I am glorified with Him; yesterday I died with Him; today I am quickened with Him; yesterday I was buried with Him; today I rise with Him. But let us offer to Him Who suffered and rose again for us…… Let us offer ourselves, the possession most precious to God, and most fitting; let us give back to the Image what is made after the Image. Let us recognize our Dignity; let us honour our Archetype; let us know the power of the Mystery, and for what Christ died.”)

(\url{http://mb-soft.com/believe/txuc/nazian26.htm})

從基督的復活裡,古教父看到天地因此而改變,人類生命因此而改變,人生方向也因此而改變.復活是宇宙間的頭等大事.無怪乎保羅說:「若基督沒有復活,我們所傳的便是枉然,你們所信的也是枉然.」(林前15:14)

今日的講道,我們看一個復活故事:路加福音24: 1-12

路24:1　七日的第一日,黎明的時候,那些婦女帶著所預備的香料來到墳墓前,

路24:2　看見石頭已經從墳墓滾開了,

路24:3　她們就進去,只是不見主耶穌的身體.

路24:4　正在猜疑之間,忽然有兩個人站在旁邊,衣服放光.

路24:5　婦女們驚怕,將臉伏地.那兩個人就對她們說:「為什麼在死人中找活人呢?

路24:6　他不在這裡,已經復活了.當記念他還在加利利的時候怎樣告訴你們,

路24:7　說:『人子必須被交在罪人手裡,釘在十字架上,第三日復活.』」

路24:8　她們就想起耶穌的話來,

路24:9　便從墳墓那裡回去,把這一切事告訴十一個使徒和其餘的人.

路24:10　那告訴使徒的就是抹大拉的馬利亞和約亞拿,並雅各的母親馬利亞,還有與她們在一處的婦女.

路24:11　她們這些話,使徒以為是胡言,就不相信.

路24:12　彼得起來,跑到墳墓前,低頭往裡看,見細麻布獨在一處,就回去了,心裡希奇所成的事.

(1)場景:墳墓

這個故事的場景是墳墓.

墳墓當然指向「死亡」.哲學家海德格(Martin Heidegger)指出,人性的特點是向前展望,想像未來,這就是人的「超越性」.所謂「超越性」,就是人不單單活在當下,更力圖超出當下,走向將來.然後,這種超越的企圖,遇到不可超越的界限,就是死亡.人可以肆無忌憚地設想明天,想像未來的生命藍圖,但無論你有多大遠見,有多大膽色,你的藍圖就是不能規劃死後的世界,就是不能越過死亡的界線.若果你的想像力像是一望遠鏡,在鏡頭內,能望到的極限,就是你的死亡.你不能了無止境地望,你的張望,在一界限內終止.若果你的生命像大海,廣闊無邊,但這無邊卻終止在地平線上.死亡就是這地平線,我們無法走到地平線的彼岸.當然,你可以向前行,地平線像是往後退,但無論你去到那裡,你的極限,就是你的地平線.我們可以希望,可以規劃,可以預期,但關乎將來,最能確定的,是總會來臨的死亡.死亡是我們生命的地平線,我們的界限.

墳墓,就是我們的界限,也是我們走向將來時,最終必然來臨的事件.

海德格是一個謙卑的哲人,守護著人類的界限.讓我們每個人盡情去想像未來吧,但我們最能確切把握的,就是終有一天走入墳墓.我想,人類是唯一能去希望的動物,也是唯一能因為「希望」被「界限」終止而產生「絕望」的動物.

「墳墓」是絕望之地.絕望者,希望被終止之處.神學家田立克(Paul Tillich)指出,人類的「墳墓」,人類的絕望情景,有三種類型.第一種類型,是物質性的(physical),就是我們身體的「死亡」.在死亡面前,我們無能為力,我們絕望.第二種絕望類型,是道德性的(moral),是指一種道德體驗,就是無論如何,你總不能盡上人的天職,常在「上天譴責」的絕望下.第三種絕望類型,是精神性的(spiritual),是指一種存在感受上的空洞,感受不到存在的意義,每日生活在空虛之下.每天,張開眼,第一個問題是:我為什麼繼續活下去呢?

現代醫學昌明,令我們稍稍遠離身體死亡的絕望.現代生活沒有了無處不在的道德規範,令我們不那麼感到道德上被譴責的絕望.現代人感受最深的,是精神上的空洞.百無聊賴,精神蒼白.抑鬱與精神病,變成致命的疾病.

我們面對的,是人的界限,是人的絕望之處,是墳墓.

(2)絕望中的絕望:「在死人中找活人」

婦女們在墳墓裡,要找死去的耶穌.她們當然找不到,因為耶穌復活了.若她們在墳墓裡要找復活的耶穌,也找不到,因為復活的耶穌也不在墳墓裡.

她們的困局,一種無可奈何的困局,是她們走不出人的界限,走不出大海的地平線,走不出墳墓.這是她們的限制,也是我們的限制.

天使對婦女們說:「為什麼在死人中找活人呢?」

這真是一個好問題.「為什麼在死人中找活人呢?」為什麼在死地找生機呢?為什麼在不可能有希望的地方找希望呢?

小時候,很喜歡走入水族店裡,看那些水族箱裡的魚.有種小魚不值錢,叫「山坑魚」.我想,「山坑魚」應活在「坑」裡.我無錢買魚,我想,或許,我能在「坑」裡捉到一條.我當時的世界,只是我住的徙置區的範圍.在這範圍內,叫得做「坑」的,就只有「坑渠」.我有一段頗長時間,經過任何「坑渠」的時候,就注目搜索,看看有沒有「山坑魚」.找了好久,都沒有找到,我其實是很失望的.這其實是天地間的一樁淒涼事件.一個小孩子,他有一卑微的夢想.他只是渴望,能找到一條不值錢的「山坑魚」,然後,滿足於擁有一條屬於自己的,不值錢的「山坑魚」.他沒有辦法走出他的世界,一個狹小的世界.他能看到的,不是山間的清泉,而是人類傾倒廢水的「坑渠」.他的無知,是他不知這是「坑渠」,而不是「山坑」.他的悲哀,是他只活在這個只有「坑渠」的世界裡.他的淒涼,是一日復一日的,在不可能有「山坑魚」的死坑渠裡,尋找「山坑魚」.

這是一個小孩子的悲哀,也常常是人類的悲哀.不知何解,我們常常試圖在沒有出路的地方找出路,沒有生命的地方找生命,在死人中找活人.

我們的生活世界的界限是墳墓,這是我們的界限所在.我們的一切努力,都走不出這個界限,走不出墳墓.人類故事的盡頭,神的故事開始了.

(3)「墳墓」外的力量

路24:6　天使說:「他不在這裡,已經復活了.」

「已經復活了」,這裡的「復活」是被動式,是「被復活」.有一個主動者,在墳墓之外,在死亡之外,在人類的界限之外,他主動介入了,他復活了耶穌.

這真是福音.在我們的界限以外,有一主動者,他能給我們在我們的界限以內找不到的幫助.若果我們全部人都跌入了深淵,而深淵裡的資源都不足以讓我們爬出去的話,則我們就需要深淵外的人幫我們,他放一條繩子下來,我們才能得救.耶穌「被復活」,這事件告訴我們,在死亡的深淵之外,在死亡的勢力範圍之外,有一主動者,他介入了,改變了死亡的困境.

又或者說,我們都像困在礦坑裡的人,出路被封住了,活著的世界被限制在一漆黑而混濁的空間裡,我們敲打著通往外面世界的金屬管道.我們敲打著,靜候著.若果只聽到自己的敲打聲,沒有回應的響聲,則這礦坑將就是我們的墳墓.若果聽到回應的響聲,知道外面有人,則能重燃希望.耶穌「被復活」,就是從墳墓外,礦坑外傳來的響聲,告訴我們,外面有人,他們在行動,救贖已開始.

耶穌「被復活」,這事件讓我們看見,墳墓外有一主動者,他不單在行動,且已經成功救回一個人.耶穌,作為第一個「被復活」的人,向我們顯明,解開死亡的禁錮,開脫道德的重負,填滿精神的空白,這是有可能的.耶穌,作為第一個「被復活」的人,向我們顯明,解救人脫望絕望的深淵,這救援工程是能成功的.這救援的工程仍在進行中,下一個被救出來的,就是你.

2018年6月底,泰國有一支少年足球隊,連教練在內,總共13人,走入一洞穴探險.忽然,洪水來了,他們被困在洞穴裡.憑他們自己的力量,是無法走出山洞了.外在的救援人員來了,但由於山洞內的環境複雜,天氣的變化無定,少年人的虛弱狀況,各種因素加起來,都未能確定能否救得到他們.當第一個少年人被救出來時,所有人都看到光明.今日,是紀念第一個人從死亡的深淵中被救出來的日子,是紀念耶穌復活的日子,讓我們看到光明.

但願:榮耀歸於聖父,聖子,聖靈.阿們.



\newpage

\section{人子得榮耀的時候到了}
\label{sec:452}
\textbf{劉國偉先生 - 人子得榮耀的時候到了}
\newline
\newline
\hyperref[sec:451]{< < < PREV SERMON < < <}
~
\hyperref[sec:toc]{[返目錄]}
~
\hyperref[sec:453]{> > > NEXT SERMON > > >}


約翰福音 12:12-50
\newline
\begin{longtable}{cl}
\hline
\hline
章節 & 經文 (和合本修訂版)\\
\hline
12:12 & \begin{tabularx}{0.7\textwidth}{X} 第二天，有一大群上來過節的人聽見耶穌要來耶路撒冷， \end{tabularx} \\ \\ \relax
12:13 & \begin{tabularx}{0.7\textwidth}{X} 就拿著棕樹枝出去迎接他，喊著：「和散那，以色列的王！奉主名來的是應當稱頌的！」 \end{tabularx} \\ \\ \relax
12:14 & \begin{tabularx}{0.7\textwidth}{X} 耶穌找到了一匹驢駒，就騎上，如經上所記： \end{tabularx} \\ \\ \relax
12:15 & \begin{tabularx}{0.7\textwidth}{X} 「錫安的兒女啊，不要懼怕！看哪，你的王來了；他騎在驢駒上。」 \end{tabularx} \\ \\ \relax
12:16 & \begin{tabularx}{0.7\textwidth}{X} 門徒當初不明白這些事，等到耶穌得了榮耀後才想起這些話是指他寫的，並且人們果然對他做了這些事。 \end{tabularx} \\ \\ \relax
12:17 & \begin{tabularx}{0.7\textwidth}{X} 當耶穌呼喚拉撒路，使他從死人中復活出墳墓的時候，同耶穌在那裡的眾人就作見證。 \end{tabularx} \\ \\ \relax
12:18 & \begin{tabularx}{0.7\textwidth}{X} 眾人因聽見耶穌行了這神蹟，就去迎接他。 \end{tabularx} \\ \\ \relax
12:19 & \begin{tabularx}{0.7\textwidth}{X} 法利賽人彼此說：「你們看，你們一事無成，世人都隨著他去了。」 \end{tabularx} \\ \\ \relax
12:20 & \begin{tabularx}{0.7\textwidth}{X} 那時，上來過節禮拜的人中，有幾個希臘人。 \end{tabularx} \\ \\ \relax
12:21 & \begin{tabularx}{0.7\textwidth}{X} 他們來見加利利的伯賽大人腓力，請求他說：「先生，我們想見耶穌。」 \end{tabularx} \\ \\ \relax
12:22 & \begin{tabularx}{0.7\textwidth}{X} 腓力去告訴安得烈，然後安得烈同腓力去告訴耶穌。 \end{tabularx} \\ \\ \relax
12:23 & \begin{tabularx}{0.7\textwidth}{X} 耶穌回答他們說：「人子得榮耀的時候到了。 \end{tabularx} \\ \\ \relax
12:24 & \begin{tabularx}{0.7\textwidth}{X} 我實實在在地告訴你們，一粒麥子不落在地裡死了，仍舊是一粒；若是死了，就結出許多子粒來。 \end{tabularx} \\ \\ \relax
12:25 & \begin{tabularx}{0.7\textwidth}{X} 愛惜自己性命的，就喪失性命；那恨惡自己在這世上的性命的，要保全性命到永生。 \end{tabularx} \\ \\ \relax
12:26 & \begin{tabularx}{0.7\textwidth}{X} 若有人服事我，就當跟從我；我在哪裡，服事我的人也要在哪裡；若有人服事我，我父必尊重他。」 \end{tabularx} \\ \\ \relax
12:27 & \begin{tabularx}{0.7\textwidth}{X} 「我現在心裡憂愁，我說甚麼才好呢？說『父啊，救我脫離這時候』嗎？但我正是為這時候來的。 \end{tabularx} \\ \\ \relax
12:28 & \begin{tabularx}{0.7\textwidth}{X} 父啊，願你榮耀你的名！」於是有聲音從天上來，說：「我已經榮耀了我的名，還要再榮耀。」 \end{tabularx} \\ \\ \relax
12:29 & \begin{tabularx}{0.7\textwidth}{X} 站在旁邊的眾人聽見，就說：「打雷了。」另有的說：「有天使對他說話。」 \end{tabularx} \\ \\ \relax
12:30 & \begin{tabularx}{0.7\textwidth}{X} 耶穌回答說：「這聲音不是為我，而是為你們來的。 \end{tabularx} \\ \\ \relax
12:31 & \begin{tabularx}{0.7\textwidth}{X} 現在正是這世界受審判的時候；現在這世界的統治者要被趕出去。 \end{tabularx} \\ \\ \relax
12:32 & \begin{tabularx}{0.7\textwidth}{X} 我從地上被舉起來的時候，我要吸引萬人來歸我。」 \end{tabularx} \\ \\ \relax
12:33 & \begin{tabularx}{0.7\textwidth}{X} 耶穌這話是指自己將要怎樣死說的。 \end{tabularx} \\ \\ \relax
12:34 & \begin{tabularx}{0.7\textwidth}{X} 眾人就回答他：「我們聽見律法書上說，基督是永存的；你怎麼說，人子必須被舉起來呢？這人子是誰呢？」 \end{tabularx} \\ \\ \relax
12:35 & \begin{tabularx}{0.7\textwidth}{X} 耶穌對他們說：「光在你們中間為時不多了，應該趁著有光的時候行走，免得黑暗臨到你們；那在黑暗裡行走的，不知道往何處去。 \end{tabularx} \\ \\ \relax
12:36 & \begin{tabularx}{0.7\textwidth}{X} 你們趁著有光，要信從這光，使你們成為光明之子。」 \end{tabularx} \\ \\ \relax
12:37 & \begin{tabularx}{0.7\textwidth}{X} 他雖然在他們面前行了許多神蹟，他們還是不信他。 \end{tabularx} \\ \\ \relax
12:38 & \begin{tabularx}{0.7\textwidth}{X} 這是要應驗以賽亞先知所說的話：「主啊，我們所傳的有誰信呢？主的膀臂向誰顯露呢？」 \end{tabularx} \\ \\ \relax
12:39 & \begin{tabularx}{0.7\textwidth}{X} 他們所以不能信，因為以賽亞又說： \end{tabularx} \\ \\ \relax
12:40 & \begin{tabularx}{0.7\textwidth}{X} 「主使他們瞎了眼，使他們硬了心，免得他們眼睛看見，他們心裡明白，回轉過來，我會醫治他們。」 \end{tabularx} \\ \\ \relax
12:41 & \begin{tabularx}{0.7\textwidth}{X} 以賽亞因看見了他的榮耀，就說了關於他的這話。 \end{tabularx} \\ \\ \relax
12:42 & \begin{tabularx}{0.7\textwidth}{X} 雖然如此，官長中卻有好些信他的，只因法利賽人的緣故不敢承認，恐怕被趕出會堂。 \end{tabularx} \\ \\ \relax
12:43 & \begin{tabularx}{0.7\textwidth}{X} 這是因他們愛人給的尊榮過於愛神給的尊榮。 \end{tabularx} \\ \\ \relax
12:44 & \begin{tabularx}{0.7\textwidth}{X} 耶穌喊著說：「信我的人不是信我，而是信差我來的那位。 \end{tabularx} \\ \\ \relax
12:45 & \begin{tabularx}{0.7\textwidth}{X} 看見我的，就是看見差我來的那位。 \end{tabularx} \\ \\ \relax
12:46 & \begin{tabularx}{0.7\textwidth}{X} 我就是來到世上的光，使凡信我的不住在黑暗裡。 \end{tabularx} \\ \\ \relax
12:47 & \begin{tabularx}{0.7\textwidth}{X} 若有人聽見我的話而不遵守，我不審判他，因為我來不是要審判世人，而是要拯救世人。 \end{tabularx} \\ \\ \relax
12:48 & \begin{tabularx}{0.7\textwidth}{X} 棄絕我、不領受我話的人自有審判他的；我所講的道在末日要審判他。 \end{tabularx} \\ \\ \relax
12:49 & \begin{tabularx}{0.7\textwidth}{X} 因為我沒有憑著自己講，而是差我來的父已經給我命令，叫我說甚麼，講甚麼。 \end{tabularx} \\ \\ \relax
12:50 & \begin{tabularx}{0.7\textwidth}{X} 我也知道他的命令就是永生。所以，我講的正是照著父所告訴我的，我就這麼講了。」 \end{tabularx} \\ \\
[1ex]
\hline
\hline
\end{longtable}
講題:人子得榮耀的時候到了 The Hour has Come for the Son of Man to be Glorified

經文:約翰福音12章12-50節

講員:劉國偉先生

場合:香港中文大學崇基學院禮拜堂

日期:2019年4月28日

引子

孩提年代,在就讀的幼稚園安排下,我參演了直屬小學的聖誕劇,在偌大的小學禮堂的舞台扮演牧羊人.我行在前頭,身旁有兩位同學也扮演牧羊人,還有全身穿上羊毛,扮演小羊的三位同學在旁邊跪行.老師指導我們三位牧羊人,當看見天使及主的榮光時,所要表達的情緒,先是懼怕,當聽到天使說「在至高之處榮耀歸與神,在地上平安歸與祂所喜悅的人.」就轉為開心.我認為那是相當忠於原著的.

我記得老師教導我們用瑟縮和後退表達恐懼,開心則以臉上的笑容和舉目望上天花來表達.問題來了,當有同學問甚麼是榮耀,老師大概沒料到有此一問,只是含糊地說 : 就是很光很光.回想起來,老師是以舞台的燈光效果來解釋,事實上要對幾歲孩童詮釋甚麼是榮耀,也頗具挑戰性!

在使徒約翰的筆下,耶穌又怎樣詮釋祂自己的榮耀?當我們回望棕枝主日那天,耶穌騎著驢駒進入耶路撒冷,他向耶路撒冷的百姓說,人子得榮耀的時候到了,到底耶穌在說甚麼? 我們又明白多少?這讓我想起扮演牧羊人的那一幕,我按照老師的指示,臉上帶著笑容和舉目望上天花,但其實我並不知道榮耀是甚麼意思.

今天講道的經文選取了約翰福音12章12節到50節.約翰在整個大段落中提及了榮耀共8次,很自然成為文本的其中一個焦點.讓我們順著經文的脈絡,嘗試去思想人子得榮耀,是什麼意思,對我們的人生,又有什麼意義.

榮耀在那裡?

因為過逾越節的緣故,耶路撒冷變得人山人海,更因為拉撒路復活的神蹟,就越法吸引群眾去觀看耶穌,所以當耶穌進入耶路撒冷之時引起連番哄動,是完全可以理解的,約翰繼續以他的深刻記憶和細膩筆觸,記下好幾個細節: 第二天,有許多上來過節的人聽見耶穌將到耶路撒冷,就拿著棕樹枝出去迎接他,喊著說:和散那!奉主名來的以色列王是應當稱頌的!(約12:12-13) 「和撒那」是現在拯救的意思,這詩句來自詩篇118篇25及26節,情緒高漲的百姓以迎接君王的氣勢,去歡迎耶穌進城,這是百姓第一次公開呼叫耶穌是以色列王.

耶穌以甚麼回應這些百姓呢? 祂騎著一隻驢駒! 按舊約記載的好幾個片段,讓我們知道驢駒並非卑劣之象徵,就好像士師睚珥的三十個兒子(士10:4),王室重臣亞希多弗(撒下17:23),掃羅王的孫子米非波設(撒下19:26),都以尊貴的身份騎於驢駒身上.所以,耶穌並非選擇了一隻卑微的動物.

另一方面,耶穌並非一位騎著馬的君王,即將要從事與戰爭有關的行動,祂騎著驢駒就正如先知以賽亞所描述的和平之君,確實恰如其分.耶穌的門徒事後就醒起先知撒迦利亞的預言(亞9:9),約翰如此撮錄了下來: 錫安的民哪,不要懼怕!你的王騎著驢駒來了.(約12:15)

這個震撼的場面,就連法利賽人都彼此說:「看哪,你們是徒勞無益,世人都隨從他去了.」(約12:19),耶穌的聲勢一時無兩,看來百姓正要擁立祂為王,但這是耶穌所理解的榮耀嗎? 我們可以肯定說不,因為作者如此記述: 這些事門徒起先不明白,等到耶穌得了榮耀以後,才想起這話是指著他寫的(約12:16),也就是說,作者認為耶穌進入耶路撒冷的那一刻,並沒有叫祂得著榮耀.

然後,約翰加插了有幾個希臘人期望見耶穌的回憶.這不尋常,所以作者把它記載下來,這代表了耶穌的名聲已經引起了外邦人的注意,祂的影響力即將從巴勒斯坦向外延展.以色列王的身份代表尊貴,統治和權力,外邦人的關注,象徵了知名度,名譽和成就,但這是榮耀之所在嗎?

人子得到的榮耀

耶穌說: 人子得榮耀的時候到了.(約12:23) 祂所指的含義是甚麼?

祂說: 我實實在在地告訴你們,一粒麥子不落在地裡死了,仍舊是一粒,若是死了,就結出許多子粒來.(約12:24) 耶穌的榮耀是先死後生,死了就結出更多生命的果實來,這個道理我們不會陌生,我們所相信的主,就是如此把自己的生命捨棄,來換取我們的生命! 我常常想,我們在人海中營營役役去經營生存之道,我們會否只顧著求生存,忘卻基督徒其實是要過一個捨己的人生? 特別當我們面臨求生與捨己的抉擇時刻,我們如何面對?

耶穌接著說: 愛惜自己生命的,就失喪生命；在這世上恨惡自己生命的,就要保守生命到永生.(約12:25) 若果用比較地道的語言去演繹愛惜生命,我會選用惜(音:錫)身,人為什麼要愛惜自己的生命,做人做得好惜身? 第一個原因可能是本能反應:怕死,怕不安全,怕損失,怕辛苦,第二個原因可能是自私,不想離開安舒區,不願意為別人犧牲自己的利益,不想自己成為蝕底的一個,不想減少自己的收入或利潤,不想降低自己的物質享受……耶穌的榮耀的道路卻是一條不愛惜自己生命的道路,祂不惜身的程度,甚至把自己的生命都捨棄.今日,有人為了理想,為了自己生活的城市,把自己的自由擺上,由本來可以好安穩的生活,走進監獄的牢房,你有沒有想過這可能是耶穌的榮耀的道路?

然後,耶穌說: 若有人服事我,就當跟從我；我在哪裡,服事我的人也要在那裡；若有人服事我,我父必尊重他.(約12:26) 原來榮耀包含得到天父的尊重,這些被天父尊重的人是誰? 就是那些服事神的人,而且,真心服事主的人,是那些緊緊跟從耶穌的人,耶穌在那裡,真正服事主的人也在那裡.那麼,耶穌在那裡? 耶穌已經告訴我們,祂走在那一條不愛惜自己生命的道路,祂把生命埋在地裡死了,祂不惜身,以捨己來換取別人生命的好處.所以,我們事奉主的人,應該緊緊跟從耶穌,祂在哪裡,服事祂的人也要在那裡.

我原是為這時候來的

耶穌所得的榮耀,就是天父的榮耀,所以耶穌說: 父啊,願你榮耀你的名!」當時就有聲音從天上來,說:「我已經榮耀了我的名,還要再榮耀.」(約12:28) 原來我們不是為了為自己去取得榮耀,卻是為了榮耀上帝,正如主禱文的開首: 我們在天上的父,願人都尊你的名為聖.就在耶穌迎向死亡,心裡憂愁之際,祂說: 我原是為這時候來的,那是耶穌為了榮耀上帝,甘心成為那一粒落在地裡死去的麥子,這一粒麥子,將要吸引萬人歸祂,就結出許多子粒來.

這次對百姓的講論,作者如此總結: 耶穌說了這話,就離開他們隱藏了.他雖然在他們面前行了許多神蹟,他們還是不信他.(約12:36-37),這條捨己的道路實在不容易走,即使有許多神蹟,即使耶穌甘願捨棄自己的生命,不等如就此一帆風順,所以,先知以賽亞感嘆說: 主啊,我們所傳的有誰信呢?以賽亞有因為這感嘆而停止傳講信息麼? 捨己的道路本來就包括了捨棄對成就的追求.

作者亦從另一個角度,表達這種困難: 雖然如此,官長中卻有好些信他的,只因法利賽人的緣故,就不承認,恐怕被趕出會堂.這是因他們愛人的榮耀過於愛上帝的榮耀. (約12:42-42) 使徒約翰在這裡以人的榮耀和上帝的榮耀作對比,身為官長,當然擁有名譽,地位和權力,被趕出會堂就等同失去了身份,領袖的身份也就保不住,甚至影響身家財富,這些愛人的榮耀的人難免要向自己的良知說不,甚至埋葬自己的信仰,不願意說真話,不想尋真理,久而久之,把自己說的謊話信以為真.

這就是耶穌面對的時代處境,又何嘗不是我們的處境,二千年過去,我們人類仍然活在人的榮耀和上帝的榮耀的張力之中,但我們信耶穌的人在人的榮耀和上帝的榮耀之間,應該義無反顧地選擇上帝的榮耀.

父的命令就是永生

作者在44節至50節加插了耶穌另一段公開講話,這段公開講話的結束部份是這樣的: 因為我沒有憑著自己講,惟有差我來的父已經給我命令,叫我說甚麼,講甚麼.我也知道他的命令就是永生.故此,我所講的話正是照著父對我所說的.」(約12:49-50) ,他的命令就是永生這句話的翻譯,和合本較為貼近原文,而現代中文譯本比較讓中文世界的讀者易於掌握: 他的命令會帶來永恆的生命.

看不透現實的人,就會覺得永生很遙遠,難以理解和觸摸,於是,被現實生活種種的追求,包括一個理想世界的追尋,困在現世有限的時空裡,無暇亦無願去走近永生.這種情況,固然在未信主的朋友當中絕不罕見,但服事主跟從主的人又如何? 我們面對現實的掙扎,也可能處於兩難之間,甚至一方面努力服事主,卻又不想緊緊跟從主.

我就想起了我兒時扮演牧羊人那一幕,我在舞台上交了功課,可能在觀眾的眼中,也算入戲的了,但我其實沒有真正掌握甚麼是榮耀歸與神.今天,我算是明白了多一些,求主幫助我,也幫助大家,繼續活在基督的榮耀裡,愛上帝的榮耀遠遠多過人的榮耀.



\newpage

\section{雞叫以先}
\label{sec:453}
\textbf{劉國偉先生 - 雞叫以先}
\newline
\newline
\hyperref[sec:452]{< < < PREV SERMON < < <}
~
\hyperref[sec:toc]{[返目錄]}
~
\hyperref[sec:454]{> > > NEXT SERMON > > >}


約翰福音 12:12-50
\newline
\begin{longtable}{cl}
\hline
\hline
章節 & 經文 (和合本修訂版)\\
\hline
12:12 & \begin{tabularx}{0.7\textwidth}{X} 第二天，有一大群上來過節的人聽見耶穌要來耶路撒冷， \end{tabularx} \\ \\ \relax
12:13 & \begin{tabularx}{0.7\textwidth}{X} 就拿著棕樹枝出去迎接他，喊著：「和散那，以色列的王！奉主名來的是應當稱頌的！」 \end{tabularx} \\ \\ \relax
12:14 & \begin{tabularx}{0.7\textwidth}{X} 耶穌找到了一匹驢駒，就騎上，如經上所記： \end{tabularx} \\ \\ \relax
12:15 & \begin{tabularx}{0.7\textwidth}{X} 「錫安的兒女啊，不要懼怕！看哪，你的王來了；他騎在驢駒上。」 \end{tabularx} \\ \\ \relax
12:16 & \begin{tabularx}{0.7\textwidth}{X} 門徒當初不明白這些事，等到耶穌得了榮耀後才想起這些話是指他寫的，並且人們果然對他做了這些事。 \end{tabularx} \\ \\ \relax
12:17 & \begin{tabularx}{0.7\textwidth}{X} 當耶穌呼喚拉撒路，使他從死人中復活出墳墓的時候，同耶穌在那裡的眾人就作見證。 \end{tabularx} \\ \\ \relax
12:18 & \begin{tabularx}{0.7\textwidth}{X} 眾人因聽見耶穌行了這神蹟，就去迎接他。 \end{tabularx} \\ \\ \relax
12:19 & \begin{tabularx}{0.7\textwidth}{X} 法利賽人彼此說：「你們看，你們一事無成，世人都隨著他去了。」 \end{tabularx} \\ \\ \relax
12:20 & \begin{tabularx}{0.7\textwidth}{X} 那時，上來過節禮拜的人中，有幾個希臘人。 \end{tabularx} \\ \\ \relax
12:21 & \begin{tabularx}{0.7\textwidth}{X} 他們來見加利利的伯賽大人腓力，請求他說：「先生，我們想見耶穌。」 \end{tabularx} \\ \\ \relax
12:22 & \begin{tabularx}{0.7\textwidth}{X} 腓力去告訴安得烈，然後安得烈同腓力去告訴耶穌。 \end{tabularx} \\ \\ \relax
12:23 & \begin{tabularx}{0.7\textwidth}{X} 耶穌回答他們說：「人子得榮耀的時候到了。 \end{tabularx} \\ \\ \relax
12:24 & \begin{tabularx}{0.7\textwidth}{X} 我實實在在地告訴你們，一粒麥子不落在地裡死了，仍舊是一粒；若是死了，就結出許多子粒來。 \end{tabularx} \\ \\ \relax
12:25 & \begin{tabularx}{0.7\textwidth}{X} 愛惜自己性命的，就喪失性命；那恨惡自己在這世上的性命的，要保全性命到永生。 \end{tabularx} \\ \\ \relax
12:26 & \begin{tabularx}{0.7\textwidth}{X} 若有人服事我，就當跟從我；我在哪裡，服事我的人也要在哪裡；若有人服事我，我父必尊重他。」 \end{tabularx} \\ \\ \relax
12:27 & \begin{tabularx}{0.7\textwidth}{X} 「我現在心裡憂愁，我說甚麼才好呢？說『父啊，救我脫離這時候』嗎？但我正是為這時候來的。 \end{tabularx} \\ \\ \relax
12:28 & \begin{tabularx}{0.7\textwidth}{X} 父啊，願你榮耀你的名！」於是有聲音從天上來，說：「我已經榮耀了我的名，還要再榮耀。」 \end{tabularx} \\ \\ \relax
12:29 & \begin{tabularx}{0.7\textwidth}{X} 站在旁邊的眾人聽見，就說：「打雷了。」另有的說：「有天使對他說話。」 \end{tabularx} \\ \\ \relax
12:30 & \begin{tabularx}{0.7\textwidth}{X} 耶穌回答說：「這聲音不是為我，而是為你們來的。 \end{tabularx} \\ \\ \relax
12:31 & \begin{tabularx}{0.7\textwidth}{X} 現在正是這世界受審判的時候；現在這世界的統治者要被趕出去。 \end{tabularx} \\ \\ \relax
12:32 & \begin{tabularx}{0.7\textwidth}{X} 我從地上被舉起來的時候，我要吸引萬人來歸我。」 \end{tabularx} \\ \\ \relax
12:33 & \begin{tabularx}{0.7\textwidth}{X} 耶穌這話是指自己將要怎樣死說的。 \end{tabularx} \\ \\ \relax
12:34 & \begin{tabularx}{0.7\textwidth}{X} 眾人就回答他：「我們聽見律法書上說，基督是永存的；你怎麼說，人子必須被舉起來呢？這人子是誰呢？」 \end{tabularx} \\ \\ \relax
12:35 & \begin{tabularx}{0.7\textwidth}{X} 耶穌對他們說：「光在你們中間為時不多了，應該趁著有光的時候行走，免得黑暗臨到你們；那在黑暗裡行走的，不知道往何處去。 \end{tabularx} \\ \\ \relax
12:36 & \begin{tabularx}{0.7\textwidth}{X} 你們趁著有光，要信從這光，使你們成為光明之子。」 \end{tabularx} \\ \\ \relax
12:37 & \begin{tabularx}{0.7\textwidth}{X} 他雖然在他們面前行了許多神蹟，他們還是不信他。 \end{tabularx} \\ \\ \relax
12:38 & \begin{tabularx}{0.7\textwidth}{X} 這是要應驗以賽亞先知所說的話：「主啊，我們所傳的有誰信呢？主的膀臂向誰顯露呢？」 \end{tabularx} \\ \\ \relax
12:39 & \begin{tabularx}{0.7\textwidth}{X} 他們所以不能信，因為以賽亞又說： \end{tabularx} \\ \\ \relax
12:40 & \begin{tabularx}{0.7\textwidth}{X} 「主使他們瞎了眼，使他們硬了心，免得他們眼睛看見，他們心裡明白，回轉過來，我會醫治他們。」 \end{tabularx} \\ \\ \relax
12:41 & \begin{tabularx}{0.7\textwidth}{X} 以賽亞因看見了他的榮耀，就說了關於他的這話。 \end{tabularx} \\ \\ \relax
12:42 & \begin{tabularx}{0.7\textwidth}{X} 雖然如此，官長中卻有好些信他的，只因法利賽人的緣故不敢承認，恐怕被趕出會堂。 \end{tabularx} \\ \\ \relax
12:43 & \begin{tabularx}{0.7\textwidth}{X} 這是因他們愛人給的尊榮過於愛神給的尊榮。 \end{tabularx} \\ \\ \relax
12:44 & \begin{tabularx}{0.7\textwidth}{X} 耶穌喊著說：「信我的人不是信我，而是信差我來的那位。 \end{tabularx} \\ \\ \relax
12:45 & \begin{tabularx}{0.7\textwidth}{X} 看見我的，就是看見差我來的那位。 \end{tabularx} \\ \\ \relax
12:46 & \begin{tabularx}{0.7\textwidth}{X} 我就是來到世上的光，使凡信我的不住在黑暗裡。 \end{tabularx} \\ \\ \relax
12:47 & \begin{tabularx}{0.7\textwidth}{X} 若有人聽見我的話而不遵守，我不審判他，因為我來不是要審判世人，而是要拯救世人。 \end{tabularx} \\ \\ \relax
12:48 & \begin{tabularx}{0.7\textwidth}{X} 棄絕我、不領受我話的人自有審判他的；我所講的道在末日要審判他。 \end{tabularx} \\ \\ \relax
12:49 & \begin{tabularx}{0.7\textwidth}{X} 因為我沒有憑著自己講，而是差我來的父已經給我命令，叫我說甚麼，講甚麼。 \end{tabularx} \\ \\ \relax
12:50 & \begin{tabularx}{0.7\textwidth}{X} 我也知道他的命令就是永生。所以，我講的正是照著父所告訴我的，我就這麼講了。」 \end{tabularx} \\ \\
[1ex]
\hline
\hline
\end{longtable}
講題:雞叫以先 Before the Rooster Crows

經文:約翰福音12章12-50節

講員:劉國偉先生

場合:香港中文大學崇基學院禮拜堂

日期:2019年5月12日

引子

二十二年前,當我還在現代詩歌運動的前線,相對活躍地參與作曲,填詞和監製的工作,我創作了一首詩歌《夜盡時》,收錄在一張鐳射光碟《有一隻羊》之中.《夜盡時》這首詩歌描寫彼得三次不認主的片段,歌詞是這樣的 :

1. 誰這麼否認? 誰這麼躲開? 誰曾確說道:縱死也不怕,絕不會否認?

2. 在冷冰深夜,伴火堆火光,誰曾確說道,再三的宣告:我不是與他一伙?

副歌:誰人何苦要苦苦追問? 情懷仍在才偷偷查聽,難明難堪卻偏偏否認,雞啼喚我醒,痛哭夜盡時!

我創作這首詩歌的時候,特別想念到彼得當夜的複雜情緒和心懷.不是嗎? 耶穌為門徒洗腳,耶穌告訴有人要賣祂,然後在客西馬尼園被捉拿,情急之下,彼得嘗試抵抗,削去對方差役的一隻耳朵,接著耶穌被解往亞那的住處,再來已是大祭司的院子,一連串的事情的發生,都是過去3年從來沒有遇見過,是如此突如其來,而且就在瞬間發生.

於是,我試圖代入彼得的思緒,以旁述的身份,為彼得反問在大祭司院子查問他是否與耶穌一起的人 : 誰人何苦要苦苦追問? 情懷仍在才偷偷查聽!  正是彼得對耶穌有情,才在這個危急關頭,冒險尾隨耶穌,好打聽事情的發展,只是想不到被人「點相」,情急之下匆匆否認,但卻繼續被人苦苦追問!

今天講道的經文並不是有關彼得否認主耶穌的片段,卻是事件發生之前幾個小時,耶穌向剩下的門徒的談話.剩下,是因為猶大已經出去了,那時候是夜間了.

在13章31節到17章26節這個大段落中,記載了耶穌自承是道路真理生命,耶穌應許賜保惠師,講論葡萄樹與枝子的比喻,安慰門徒,還有分離前的禱告.這些都是主線分明,有系統的主題論述,而這個大段落的第一小段,約翰福音13章31節到38節這8節經文,驟眼看來,顯得有點零碎,也有點格格不入,使徒約翰何以如此鋪排?

仔細看真一些,這8節經文,其實蘊含著4個人生的根本問題,大部份活在世上的人,包括你和我,在人生旅程中都或多或少要面對和回答,而且這幾個問隨著我們不同的人生階段,會不斷來回拜訪我們.若果雞叫是這個生死關頭的夜晚的一個警號,又或者是把人敲醒的鐘,那麼我就覺得在這鐘聲敲出之前,看看我們自己,或者這個社會,要回答什麼問題,就更顯得十分有意義,有價值.

人生終極追求

在這個看似零碎的段落,作者記下耶穌所說的第一段說話:「如今人子得了榮耀,上帝在人子身上也得了榮耀.上帝要因自己榮耀人子,並且要快快地榮耀他.」(約13:31-32) 我們在兩星期前的講道,曾經分享過人的榮耀與上帝的榮耀的對比,耶穌選擇了上帝的榮耀,那是一粒麥子落在地裡死了,就結出許多子粒來.(約12:24) 耶穌的榮耀是先死後生,祂選擇不愛惜自己的生命,或者唔惜(音:錫)身到一個地步,把自己的生命都捨棄,祂更教導門徒服事主的人,是那些緊緊跟從耶穌的人,耶穌在那裡,真正服事主的人也在那裡.耶穌在那裡? 耶穌就走在那一條不愛惜自己生命的道路,祂把生命埋在地裡死了,祂不惜身,以捨己來換取別人生命的好處.

耶穌在這個充滿張力的夜晚,提醒跟從祂的人,若果人生走到盡頭,我們可會為自己的一生下一個結論 : 如今我得了榮耀,上帝在我身上也得了榮耀.以得榮耀以及上帝在我身上得榮耀作為人生的終極追求,固然是十分崇高的理想,卻並非離天萬丈,無法可及的境界,我們要做的,就是甘心成為那一粒落在地裡死去的麥子,走捨己的道路,不要做貪愛人的榮耀的人,不要向自己的良知說不,不要埋葬自己的信仰,要說真話,追尋真理,不會久而久之,把自己說的謊話信以為真.告訴自己,就由生活中一些微少的事情開始,做一個得著上帝的榮耀,而上帝在我身上也得到榮耀的人.

我會簡化第一個根本問題是人生的終極追求.

生命方向與時機

第二個根本問題是生命方向與時機.定下了人生目標,卻走錯了方向,或者一而再錯失時機,那是人生一種很大的遺憾,也差不多等同空談理想和信仰.

約翰記載耶穌的第二個小段的說話是: 「小子們,我還有不多的時候與你們同在；後來你們要找我,但我所去的地方你們不能到.這話我曾對猶太人說過,如今也照樣對你們說.」(約13:33) 從約翰福音記載的續後發展來解讀,這一節經文的語意十分清晰,就是耶穌正式告訴門徒祂即將與他們分離,並且去一處他們現在不能去的地方.其實耶穌先後在約翰福音7章34節和8章21節對猶太人說過類似的話,加上從上文下理,我們可以確切理解耶穌所要去的地方是指著回到父那裡.

耶穌這段說話對我們有什麼啟發和意義?其實我們的人生要往那裡去,什麼時候去,又豈真的可以由自己徹底操控嗎?耶穌提醒門徒要珍惜與主一起的時光,這叫我立時就想起,耶穌在逾越節前六日,在伯大尼的筵席中,提醒門徒「你們不常有我」(約12:8下).弟兄姊妹,什麼時刻是我們生命中應當特別珍惜?有那些生命時機是我們必須牢牢掌握的?什麼人生方向是我們必須把持的?當耶穌說我所去的地方你們不能到,可會想過我們要如何作好準備,最終走進上帝的榮耀裡?

彼此切實相愛

接著,耶穌說 : 「我賜給你們一條新命令,乃是叫你們彼此相愛；我怎樣愛你們,你們也要怎樣相愛.你們若有彼此相愛的心,眾人因此就認出你們是我的門徒了. 」(約13:34-35)在生死危急關頭,彼此相愛其實是很大的考驗! 耶穌所要求門徒的,不單是一種掛在咀邊的普普通通的愛,而是一種捨己忘生的愛,因為耶穌說我怎樣愛你們,你們也要怎樣相愛,As I have loved you,在黑暗掌權的深夜,在警世的鐘聲敲響之際,我們在求取生活的安穩,甚至求取生存的空間之時,可以像耶穌那樣去愛嗎?

對於老約翰來說,他大概不會猶豫,因為他如此寫道 :主為我們捨命,我們從此就知道何為愛；我們也當為弟兄捨命.(約壹3:16) 所以,使徒約翰提醒我們 : 「 小子們哪,我們相愛,不要只在言語和舌頭上,總要在行為和誠實上.」 .(約壹3:18) 親愛的弟兄姊妹,很多時,主的家裡的人最難愛! 願我們彼此勉勵,珍惜弟兄姊妹的情誼,互相欣賞各人的恩賜,接納彼此的服侍,靠主恩典,主怎樣愛我們,我們也要怎樣相愛,眾人因此就認出我們是主的門徒了.

這是我想提及的人生第三個根本問題 : 如何彼此切實相愛?

跟從與犧牲

這段經文引發我思想的人生第4個根本問題是,如何去面對跟從和犧牲的考驗.

西門‧彼得問耶穌說:「主往哪裡去?」耶穌回答說:「我所去的地方,你現在不能跟我去,後來卻要跟我去.」彼得說:「主啊,我為甚麼現在不能跟你去?我願意為你捨命!」(約13:36-37)作者如實記載彼得的豪情壯語,與接著在夜半發生的事情,這段對話竟成為反諷的情節.

彼得並非對耶穌所說的話的含意完全懵然不知,「有人要賣我了」,「你們不都是乾淨的」,難道他體察不到眼前的氣氛和情緒嗎? 正因為彼得覺得情勢不好,就鏗鏘有力地說: 我願意為你捨命!I will lay down my life for you! 我形容這是亂勢豪情,有熱情卻未必有足夠的深度,當有如火的試驗臨到,很可能就要馬上退縮了,上帝因我們得著榮耀,不應該是剎那間的情節,而是一生持守的信仰,憑著信以致於信,熬過漫長的黑夜,直到夜盡黎明!

持守信仰不是憑著一份衝動,因為那是整個人生的歷練,耶穌對彼得說「後來卻要跟我去」(約13:36下),我繼續走信仰的道路,可能因為人生種種,包括這個愈來愈不理想的社會及整個時代給我們的考驗,信仰的路愈走愈難愈窄,但卻看見恩典愈多!就在這個情勢不好的黑夜,我們又可會問自己,我們願意為耶穌擺上什麼? 借用彼得的說話:What we will lay down for Jesus?

雞叫以先

讓我們再重溫《夜盡時》這首詩歌的副歌歌詞: 誰人何苦要苦苦追問? 情懷仍在才偷偷查聽,難明難堪卻偏偏否認,雞啼喚我醒,痛哭夜盡時!

生命總難免有經不起考驗的時刻,因為我們也不過是凡人,有時我們會身處難明難堪卻偏偏跌倒的境況,就正如保羅在羅馬書說的: 我所願意的善,我反不做；我所不願意的惡,我倒去做.(羅7:19) 但耶穌清楚知道彼得,祂明言雞叫以先,你要三次不認我.耶穌必定沒有輕看雞啼痛哭的彼得,並且更在提比哩亞海再次委以重任予彼得.

弟兄姊妹,雞叫以先,我們思考的四個人生的根本問題: 人生終極追求,生命方向與時機,彼此切實相愛 和跟從與犧牲的考驗,用地道一點的話去描述大概是: 人生為乜?  要去邊度? 要點樣愛? 犧牲乜ye ?

但願敲鐘的聲音響起時,我們可以有紮紮實實的答案,即或我們偶然經不起考驗,耶穌清楚我們,祂必定不輕看我們在雞叫時候的覺醒,也會盛載我們懊悔的眼淚.



\newpage

\section{在黑暗與榮耀之間}
\label{sec:454}
\textbf{鄧瑞強博士 - 在黑暗與榮耀之間}
\newline
\newline
\hyperref[sec:453]{< < < PREV SERMON < < <}
~
\hyperref[sec:toc]{[返目錄]}
~
\hyperref[sec:455]{> > > NEXT SERMON > > >}


約翰福音 13:31-35
\newline
\begin{longtable}{cl}
\hline
\hline
章節 & 經文 (和合本修訂版)\\
\hline
13:31 & \begin{tabularx}{0.7\textwidth}{X} 猶大出去後，耶穌說：「如今人子得了榮耀，神在人子身上也得了榮耀。 \end{tabularx} \\ \\ \relax
13:32 & \begin{tabularx}{0.7\textwidth}{X} 如果神因人子得了榮耀，神也要因自己榮耀人子，並且要立刻榮耀他。 \end{tabularx} \\ \\ \relax
13:33 & \begin{tabularx}{0.7\textwidth}{X} 孩子們！我與你們同在的時候不多了；你們會找我，但我所去的地方，你們不能去。這話我曾對猶太人說過，現在也照樣對你們說。 \end{tabularx} \\ \\ \relax
13:34 & \begin{tabularx}{0.7\textwidth}{X} 我賜給你們一條新命令，乃是叫你們彼此相愛；我怎樣愛你們，你們也要怎樣彼此相愛。 \end{tabularx} \\ \\ \relax
13:35 & \begin{tabularx}{0.7\textwidth}{X} 你們若彼此相愛，眾人因此就認出你們是我的門徒了。」 \end{tabularx} \\ \\
[1ex]
\hline
\hline
\end{longtable}
講題:在黑暗與榮耀之間 Between Darkness and Glory

經文:約翰福音13章31-35節

講員:鄧瑞強博士

場合:香港中文大學崇基學院禮拜堂

日期:2019年5月19日

各位弟兄姊妹,早安.

在過去的某一天,忽然間,你進入這個世界.在將來的某一天,忽然間,你離開這個世界.我們意識到自己在一旅程中,從一端走到另一端.兩端都是神秘的,不可測的,只有中間這段旅程,我們可以見到沿途的風光.或許,有一天,你會坐下來,想一想,我當如何走這段旅程.

有人想,這旅程太艱苦,於是,坐下來,不再走了.人可以「物化」自己,將自己視為「物」,不再有感情,不再有期望.像石頭一樣,立在路旁.天地如何變化,與我無關.不談過去,也沒有將來.他可能成為宅男,只活在電腦世界裡.他可能用藥物麻醉自己,只活在當下的迷糊裡.他可能是一個消費狂,只活在物質提供的當下快感裡.

有人想,這旅程太艱苦,連做一塊石頭也不想,就這樣離去了.

有人在途中,不斷掠奪資源,不斷推倒旁人,不斷爭先恐後.我們在旅途中最怕遇上這種人,他們拖著的行李箱,很多時會滾過我們的腳掌.然後,忽然,有一天,他們的旅途中止了,面對他們最恐懼的神秘.

有人在途中,想著將來的神秘,想像這神秘就是天堂,想著天堂裡人與人之間的微笑,然後,決定以那種天堂的微笑度過人間的旅途.他們仍在旅途中,仍未到達彼岸,卻將彼岸的微笑帶來這人間,讓人間感受一下天堂,讓在艱苦旅途中同行的人,得到扶持,得到鼓勵,得到希望.我想,信仰的覺悟正是如此.從生命到死亡的過程中,努力去愛,讓在旅途中的所有人,找到前進的動力.

耶穌的一生,就是這樣.

今日講道的經文,是今日的福音書經課:約翰福音13:31－35

我們連著13:30來讀,這經文提供一背景資料.

約13:30　猶大受了那點餅,立刻就出去.那時候是夜間了.

約13:31　他既出去,耶穌就說:「如今人子得了榮耀,神在人子身上也得了榮耀.

約13:32　神要因自己榮耀人子,並且要快快地榮耀他.

約13:33　小子們,我還有不多的時候與你們同在；後來你們要找我,但我所去的地方你們不能到.這話我曾對猶太人說過,如今也照樣對你們說.

約13:34　我賜給你們一條新命令,乃是叫你們彼此相愛；我怎樣愛你們,你們也要怎樣相愛.

約13:35　你們若有彼此相愛的心,眾人因此就認出你們是我的門徒了.」

這段經文的開首:

約13:30　猶大受了那點餅,立刻就出去.那時候是夜間了.

在最後晚餐,出賣耶穌的猶大,離開耶穌,離開同伴,進入黑暗.

猶大,一路上與耶穌在一起,一路上看到耶穌對他人的愛,一路上聽耶穌言說那天國,但很可惜,他有自己的人生規劃,他竟然看不到耶穌對所有在旅途中的人的意義,他選擇排斥身邊的人,選擇出賣身邊的耶穌,選擇行自己的路.聖經說,他進入了黑暗.對猶大而言,他看不到終點的光明,看不到終點的光明對生命的意義.他猛力地反抗,嘲諷看起來軟弱的仁愛,拒絕耶穌講的希望,試圖運用自己的「強力意志」闖開一條出路,終歸將自己推向人性的黑暗深處.

當黑暗漫延時,經文即時轉向光明.

約13:31　他既出去,耶穌就說:「如今人子得了榮耀,神在人子身上也得了榮耀.

約13:32　神要因自己榮耀人子,並且要快快地(或譯:立刻)榮耀他.

兩節經文,在耶穌的說話裡,所有動詞都是「榮耀」.四個「榮耀」的動詞,像四盞射燈,照著走向死亡的耶穌.第31節的兩個「榮耀」是「被動語態」的,人子(耶穌)「被榮耀」,而天父因著耶穌也「被榮耀」.耶穌「被榮耀」,當然,是被天父榮耀.天父的光明,在耶穌走向十架時,充滿耶穌.天父的光明,尋找顯露出來的罅隙,十字架,正是讓這光明展露出來的罅隙.

從另一方面說,當耶穌走向十架時,天父「被榮耀」,天父的光明在十架中被萬物認識.天父的光明,在耶穌幽暗的十架裡顯露.耶穌的十架,看來黑暗,卻將天父的光明顯露出來.天父的心在人間,在愛他人而承受痛苦的人身上顯明.為了愛而背負十架的心,就是天心.天心在,天父的光明便顯露,天父便「被榮耀」.神聖的光明,不在功名那裡,不在財富那裡,而在愛的十架裡.

第32節的兩個「榮耀」是「主動語態」的,天父主動的「榮耀」耶穌,並且要即時的「榮耀」耶穌.我們會問:「天父為何要主動的『榮耀』耶穌?」經文答:天父「因自己榮耀人子」,這經文直譯是「神在自己內榮耀人子」.這意思是:天父在自己的本性內(因著自己的本性),找到榮耀人子的理由.天父的本性,就是為愛人類而承受愛的苦難.耶穌的十架行動,反映出天父的本性.天父主動的「榮耀」耶穌,因為耶穌反映天父自身的光明.

經文又提及,天父要「立刻榮耀」耶穌.我們會問:為何要「立刻」?為何不等到十架的黑暗過去,才談論「榮耀」與「光明」?理由是:十架本身,就是天父彰顯光明的時刻.天父並不棄絕十架,然後等「雨過天青」後才顯露光明.人甘願承擔十架,本身就是「雨過天青」,本身就是光明.

在人生的旅途上,我們尋找光明,尋找榮耀,尋找上帝,尋找旅途中的意義,尋找面對死亡的方式.這個人生旅程,耶穌走過一次了.他承擔了人生的十架,展示對他人的毫無保留的愛,他就以這看來幽暗的十架,讓我們看到光明,看到榮耀,看到上帝,看到人生旅途的意義,看到面向死亡的最好的活著方式.

托爾斯泰有篇短篇小說,叫＜人靠什麼而活＞,故事是這樣的.

在某條村落,住著一個鞋匠的家庭,他們很貧窮.在某個寒冷的日子,鞋匠出去收取工作後的工錢,卻一無所獲.回家時,路過教堂,看見一個赤身的人.他在想,這人可能是壞人,若我走過去,說不定他會搶我的東西.於是,他走開了.但他受到良心的譴責,最後,他還是走過去.在那裡,他看到一個赤身的年輕人,在寒冷中瑟縮.鞋匠將自己的衣服給他穿.想著,讓他在外頭,肯定會凍死.於是,便帶他回家.

鞋匠回到家裡,叫妻子拿點吃的東西來.妻子頓時大發雷霆,一家人自己吃,食物也不夠,丈夫還多帶一個陌生人回來.鞋匠只管對妻子說,難道你不愛上帝嗎?聽到這裡,妻子便心軟,拿出僅有的食物來,讓他們吃掉.這時,陌生人笑了.

陌生人住了下來,鞋匠教他做鞋.陌生人學得很快,甚至比鞋匠做得更好.很多人來找他們做鞋,鞋匠的生活,因而也好轉了.

有一天,一個有錢人,帶著貴重的皮革,來到鞋匠這裡,要訂做一對皮靴.他要求這對靴必須能夠穿上一年,不能走樣,不能斷線.否則,要拉他們坐監的.當然,能做到的話,人工是很高的.接不接這單生意,鞋匠望望陌生人.陌生人點頭,於是,鞋匠接下這單生意.這時,陌生人又笑了.

基本的工序由陌生人來做,最後,陌生人沒有做出一雙皮靴,而做出一雙拖鞋來.鞋匠責備他說,你這是怎麼搞的,你把我害慘了,人家要求我們做皮靴,你卻做了拖鞋來,我怎賠得起這貴重的皮子呢!

當鞋匠責罵時,門外有客人來.客人原來是那有錢人家裡的僕人.僕人說,我家老爺用不著靴子了,他已經去世啦.我家太太派我來這兒,著你們用那皮子做一雙拖鞋便可以了.這時,陌生人便將那已做好的拖鞋交給僕人,讓他去了.

時光飛逝,陌生人已在鞋匠家裡過了六年.有一天,來了一個婦人,她帶著兩個女孩,婦人希望鞋匠為這兩個女孩做鞋.鞋匠招呼著她們,卻發現陌生人定睛的看著那兩個女孩,好像他認識她們似的.鞋匠妻子和那婦人聊天,問道:「你是她們的媽媽嗎?」婦人說:「不是.我和她們非親非故,我只是收養了她們.」話題一打開,那婦人便將往事一一道來.事情是這樣的:六年前,兩個女孩的父母,一周內先後去世.她們的父親是一個樵夫,砍樹時,一棵樹倒在他身上,他就死了.過幾天,她們的母親生下她們,也死了.婦人是她們的鄰居,看見這兩個可鄰的女孩,便收養了她們.這時,陌生人又笑了.

婦人的故事說完了,訂鞋的事交代清楚了,便走了.這時,陌生人身上忽然發出光芒.鞋匠夫婦很奇怪,便問他:「你是什麼人?」陌生人說:我其實是死亡天使,上帝差我來凡間,要帶走人的靈魂.可是,我有一次,沒有聽從上帝的話.上帝派我去取一個婦人的靈魂.我下到凡間,見到一個女人,她剛剛生下一對女嬰.她知道我是去取她的靈魂的.她求我說:「天使啊,我的丈夫死了,我也無親無故,沒有人會照顧這對孤女,請不要帶走我的靈魂,容我養大她們.」

天使看到她的慘況,便放過她們.天使回到上帝那裡,對上帝說,我不能取走這女人的靈魂.

上帝對天使說:你去,取走那女人的靈魂.但由於你違命,你暫時不得重返天界,直到你在凡間,找到以下三個問題的答案為止.一,「人性裡有什麼?」二,「人無法知道什麼?」三,「人靠什麼而活?」

天使再次來到凡間,帶走那女人的靈魂,卻無法再回天界了.他無處可去,便在教堂那裡藏身.飢寒交迫之際,竟然有人肯將自己的衣服給他穿,帶他回家.當鞋匠的妻子將僅有的食物給他吃時,他領悟到第一個問題的答案.「人性裡有什麼?」人性裡有仁愛.於是,天使笑了.

天使說,當那有錢人來做鞋時,我看見我的同伴在他後面.這有錢人來做鞋,盤算著一年後的事,卻不知道,死亡天使當天就要帶走他的靈魂.於是,我領悟到第二個問題的答案.「人無法知道什麼?」人無法知道自己真正需要什麼.於是,我又笑了.

天使繼續說,當那婦人帶著那兩個女孩來的時候,我一眼認出她們.六年前,我以為,沒有了女孩的母親,她們便活不成.於是,我試圖放過她們的母親.我想不到,有人會乳養她們,視她們為女兒.我在她們的養母身上,看到活生生的上帝.於是,我又笑了.我終於明白第三個問題的答案.「人靠什麼而活?」人不是為自己而活,人靠愛而活.人類能夠活下去,靠人與人之間的愛.在愛中,上帝露面,因為神就是愛.

說罷,天使重獲翅膀,衝天而去.

我們回到經文.

約13:33　「小子們,我還有不多的時候與你們同在；後來你們要找我,但我所去的地方你們不能到.這話我曾對猶太人說過,如今也照樣對你們說.」

明顯,耶穌要回到天界,回到天父的身邊.他回到天父身邊的途徑,是藉著十字架,藉著承擔責任,藉著愛.對外人(猶太人)來說,這是匪夷所思的.就算對耶穌的跟隨者來說,這也不見得容易明白.若果基督徒明白,並且努力實踐,我想,這個世界就會比現在的情況好得多了.

約13:34　「我賜給你們一條新命令,乃是叫你們彼此相愛；我怎樣愛你們,你們也要怎樣相愛.

約13:35　你們若有彼此相愛的心,眾人因此就認出你們是我的門徒了.」

十字架,就是愛.愛,就是十字架.十字架,徹底的愛,犧牲的愛.現代人習慣了消費,很容易將愛「浪漫化」,將愛變成「消費品」,將愛看成是一樣「令自己快樂」的東西.與身邊的人在一起,不再快樂了,便理解成「再沒有愛了」.

我們習慣將「愛」看成是一種可以「享受」的東西,卻很少將「愛」看成是我們承擔責任,付上辛勞而創造出來的成果.德蘭修女對愛的定義,講得又簡單,又好.她說:「Love until it hurts.」「愛,直至受傷」.真愛,總會令你受傷.真愛的典範,就是十字架.為了愛,完全放下自己,承受責任,忍受痛苦.德蘭修女繼續說:「I have found the paradox, that if you love until it hurts, there can be no more hurt, only more love.」若我們愛,直至受傷時,我們會發現,這根本不算是受傷,而是更多愛.

耶穌說:「我怎樣愛你們,你們也要怎樣相愛.」若耶穌以十架的愛去愛我們,我們也當以十架的愛彼此相愛.

猶大,他走向黑暗.耶穌,他活出光明.我們,處在黑暗與榮耀之間.我們可以伴隨猶大,走向黑暗.也可以學效耶穌,活出光明.走向黑暗,是容易的.活出光明,是艱難的.我希望,我們都能活出光明.

但願:榮耀歸於聖父,聖子,聖靈.阿們.

(附記:本人在崇基禮拜堂講道的講章已結集成《荷花池畔春草綠》講章集,有興趣的朋友可購買.試讀本可參閱:\url{http://www.news.ccc.cuhk.edu.hk/do/wp-content/uploads/2019/04/CCC-SK-Tang-Sermon-Text-Sample.pdf}.訂購表:\url{http://www.news.ccc.cuhk.edu.hk/do/wp-content/uploads/2019/04/CCC-SK-Tang-Sermon-Order-Form.pdf})



\newpage

\section{何處是吾家}
\label{sec:455}
\textbf{林豪恩先生 - 何處是吾家?}
\newline
\newline
\hyperref[sec:454]{< < < PREV SERMON < < <}
~
\hyperref[sec:toc]{[返目錄]}
~
\hyperref[sec:456]{> > > NEXT SERMON > > >}


使徒行傳 16:6-15
\newline
\begin{longtable}{cl}
\hline
\hline
章節 & 經文 (和合本修訂版)\\
\hline
16:6 & \begin{tabularx}{0.7\textwidth}{X} 因為聖靈禁止他們在亞細亞講道，他們就經過弗呂家、加拉太一帶地方。 \end{tabularx} \\ \\ \relax
16:7 & \begin{tabularx}{0.7\textwidth}{X} 到了每西亞的邊界，他們想要往庇推尼去，耶穌的靈卻不許。 \end{tabularx} \\ \\ \relax
16:8 & \begin{tabularx}{0.7\textwidth}{X} 他們就越過每西亞，下特羅亞去。 \end{tabularx} \\ \\ \relax
16:9 & \begin{tabularx}{0.7\textwidth}{X} 夜間，有異象向保羅顯現。有一個馬其頓人站著求他說：「請你過來，到馬其頓來幫助我們！」 \end{tabularx} \\ \\ \relax
16:10 & \begin{tabularx}{0.7\textwidth}{X} 保羅既看見這異象，我們就立即設法往馬其頓去，認為神呼召我們傳福音給那裡的人。 \end{tabularx} \\ \\ \relax
16:11 & \begin{tabularx}{0.7\textwidth}{X} 我們從特羅亞開船，直行駛到撒摩特喇，第二天到了尼亞坡里； \end{tabularx} \\ \\ \relax
16:12 & \begin{tabularx}{0.7\textwidth}{X} 從那裡來到腓立比，就是馬其頓這一帶的一個重要城市，也是羅馬的駐防城。我們在這城裡住了幾天。 \end{tabularx} \\ \\ \relax
16:13 & \begin{tabularx}{0.7\textwidth}{X} 在安息日，我們出城門，到了河邊，知道那裡有一個禱告的地方，我們就坐下來對那些聚會的婦女講道。 \end{tabularx} \\ \\ \relax
16:14 & \begin{tabularx}{0.7\textwidth}{X} 有一個賣紫色布的婦人，名叫呂底亞，是推雅推喇城的人，素來敬拜神。她在聽著，主就開導她的心，使她留心聽保羅所講的話。 \end{tabularx} \\ \\ \relax
16:15 & \begin{tabularx}{0.7\textwidth}{X} 她和她一家都領了洗，就求我們說：「你們若以為我是真心信主的，請到我家裡來住。」於是她堅決請我們留下。 \end{tabularx} \\ \\
[1ex]
\hline
\hline
\end{longtable}
講題:何處是吾家 Where is Home

講員:林豪恩先生

經文:使徒行傳16:6-15

場合:香港中文大學崇基學院神學院容啟東校長紀念樓G/F新聖堂

日期:2019年5月26日

一．引言

大家知否每年一度「全球罕見的人口流動」是甚麼盛事呢?這就是春運.2018年春運大軍是29.7億人次,相當於非洲,歐洲,美洲,大洋洲的總人口搬一次家.

原來,春運是與家有關.有網絡作者如此描述春運:「甚麼是春運呢?就是有工作的地方沒有家；有家的地方沒有工作.他鄉容不下靈魂,故鄉安不了肉身,肉體和靈魂互相絞撕,便有了春運.」春運,顯現了中國人千百年來的觀念:家是與鄉土分不開的.

當下居住在香港這個城市的大多數人,想到家,似乎不一定想到鄉土.香港人對家的想像又是怎樣的呢?

有地產公司的廣告說:「給你一個五星級的家」,另一個地產公司的廣告說:「以心建家」.對香港人來說,家會否與地產公司分不開呢?不久之前,某本地電視台製作實況節目《有樓萬事足》,內容圍繞香港人的「上樓夢」.其中一位參演市民說:「有樓有高潮」,「無樓唔好白撞」,刺激了不少香港人的神經,引起廣泛討論.事實上,無論是認同者,批評者,嘲諷者,大部分香港人都向著「買樓」的目標前進,不少買了樓的人為了供樓而工作,而尚未買樓的人為著買不到樓而焦慮.幾年前有一篇潮文名為「女朋友帶你返屋企食飯」,被改編成YouTube短片,片中「伯母」一角,活現了對居住於不同種類居所的未來女婿的不同反應,點擊率一時無兩.

其實,香港人把家與買樓聯想在一起,並不是新鮮事物.由1970年代開始,大學生圈子已經有言論以「四仔主義」概括人生四大目標,包括「屋仔」,「車仔」,「老婆仔」與「生番個仔」.「四仔主義」也可以借用來描述香港人對家的想像:一個家應該要有「屋仔」,「車仔」,「老婆仔或老公仔」,「生埋個BB仔」,那就是典型的家了.

今天的四段經課,雖然並非直接談論家,也不是教導家人之間的相處之道,但卻可以引發我們對家的想像.這些想像,並不是要反對我們以上談及社會普遍對家的想法,而是嘗試擴充我們的想像,思考其他的狀況,以涵概更多的可能性,例如沒有「屋仔」,或者沒有「老婆仔／老公仔」,或者沒有「ＢＢ仔」的生活狀況.

二．內容

1．以使命為家

第一段經文取自使徒行傳16章6-15節:

16:6因為聖靈禁止他們在亞細亞講道,他們就經過弗呂家,加拉太一帶地方.

16:7到了每西亞的邊界,他們想要往庇推尼去,耶穌的靈卻不許.

16:8他們就越過每西亞,下特羅亞去.

16:9夜間,有異象向保羅顯現.有一個馬其頓人站著求他說:「請你過來,到馬其頓來幫助我們!」

16:10保羅既看見這異象,我們就立即設法往馬其頓去,認為上帝呼召我們傳福音給那裡的人.

16:11我們從特羅亞開船,直行駛到撒摩特喇,第二天到了尼亞坡里；

16:12從那裡來到腓立比,就是馬其頓這一帶的一個重要城市,也是羅馬的駐防城.我們在這城裡住了幾天.

16:13在安息日,我們出城門,到了河邊,知道那裡有一個禱告的地方,我們就坐下來對那些聚會的婦女講道.

16:14有一個賣紫色布的婦人,名叫呂底亞,是推雅推喇城的人,素來敬拜上帝.她在聽著,主就開導她的心,使她留心聽保羅所講的話.

16:15她和她一家都領了洗,就求我們說:「你們若以為我是真心信主的,請到我家裡來住.」於是她堅決請我們留下.

這個段落記述了保羅及同工第二次宣教旅程的其中一些行程,他本來計劃往亞細亞探訪第一次行程所建立的教會,但卻體會到聖靈禁止他們在亞細亞講道,他們打算北上庇推尼,但又體會到聖靈的阻止.保羅本來沒有計劃跨過愛琴海遠赴歐洲,但卻遇到馬其頓異象,於是越過大海,遠赴彼岸.在短短的幾節經文裡面,作者記錄了保羅經過的地方名稱,分別有:亞細亞,弗呂家,加拉太,每西亞,庇推尼,特羅亞,馬其頓,撒摩特喇,尼亞坡里,腓立比.透過閱讀這些陌生的地方名稱,讀者不難感到到保羅四處流徙,無處為家的漂泊.家在哪裡呢?最後一節,終於有「家」出現了:16:15「她和她一家都領了洗,就求我們說:「你們若以為我是真心信主的,請到我家裡來住.」於是她堅決請我們留下.」

此刻,腓立比呂底亞的家收留了保羅,不過,對保羅來說,這只是「借來的地方,借來的時間」.「借來的地方,借來的時間」這句熟識的說話,是從何而來的呢?區家麟考究過,1968年,長駐東亞與香港的澳洲籍記者Richard Hughes,曾以此為書名書寫香港.不過,此語並不是Richard Hughes原創,乃源自一位旅居香港,叫Tom Wu的上海商人.當年逃難到香港的上海資本家,亂離中,洞察香港特質,感觸一句,能夠流傳久遠,成為香江名句.這處「借來的地方,借來的時間」收留過好些漂泊流徙的人,成為了不少無家者的家.或者,有人認為這是一個「漏洞」,然而,對不少人來說,這是a home away from home.

保羅說自己是「屬於以色列族,便雅憫支派,是希伯來人所生的希伯來人」,然而,他的家卻不在以色列境內,而在現今土耳其境內的大數城,這不是他的選擇,而是他的先祖的選擇.當保羅認識耶穌是基督之後,他不再留在大數,卻以安提阿為基地,並經常到當時的亞細亞及希臘地區宣教.他的行程,他所到之處,有些是他的計劃,有些是回應呼喚,其實有不少是被逼離開,甚至是逃亡而走的.

當下生活在這個城市的人,彷彿是自己的選擇,但當中有多少是歷史的偶然呢?就算在本地出生,這也可能是父輩或祖父輩的決定.曾經何時,我們以為自己找到了a home away from home,然而,我們彷如驚弓之鳥,風吹草動,我們感到此地不可留,又要拔營起行.

或者,「安定繁榮」是一種期盼,從歷史來看,「在途上」也並非罕有的事情.保羅的一生大部份時間都是「在途上」.「在途上」彷彿無家,也可以有很多收留他的家；「在逃上」自己沒有家,也可以協助其他人建立家.

中文大學前歷史系教授劉義章博士寫了一本關於調景嶺的書,其中有一段如此說:「隨著1949年政局劇變……基督教和天主教的外國宣教士被逼離開中國……大都暫時棲身在九龍尖沙咀漆咸道兵房中所謂「自由避風港」……當時難民絡繹不絕抵達香港,等待救援.瑞典籍宣教士戴瑞蘭與美籍宣教士西門英才,包美達等人,帶著一大袋藥品去摩星嶺探望患病難民,贈醫施藥.信義差會挪威籍顧永榮及其同工則用他們僅有的一美元,在摩星嶺創辦了一所「難民兒童義務學校」,讓兒童上學讀書.當難民被遷入調景嶺,戴瑞蘭等人亳不猶豫與他們同往,住進簡陋的「Ａ字棚」.隨後他迅即向香港社會局調景嶺營辦公處申請並獲准為難為提供免費醫療服務.1950-1955年成為教會在調景嶺發展的黃金時期.在小區成立的眾教會全由輾轉來自大陸多個省市的傳教士所建立.教會和西教士以耶穌基督的大愛關顧這些孤苦無依的難民,成功把調景嶺居民聚合在一起,使他們努力及勇敢地生活下去.」《孤島扁舟:見證大時代的調景嶺》的作者劉義章博士成長於調景嶺,早年就讀調景嶺小學及中學.另一位作者計超,1949年生於上海,1961年由滬來港定居,翌年經靈實肺病療養院司務道護士長介紹,入住由戴大衛先生創辦的調景嶺學生輔助社馬可紀念之家(即現在香港學生輔助會的前身),並就讀於區內信義小學和慕德中學.

一般來說,人們認為家很應該有「屋仔」,「車仔」,「老婆仔或老公仔」,「生埋個BB仔」,然而,家也可以不由「四仔」來定義,對如保羅來說,家是由使命來定義.無論是自己選擇停留在那裡,還是被逼停留在那裡,保羅都視自己為福音的使者,為耶穌基督作見證.今天,我們居住在這個城市,明天,我們可能繼續在這裡,也可能在上海,台灣,英國,加拿大,美國,澳洲或新西蘭,如果繼續有使命,我們就有機會有家.

2．約翰福音14章23-29節:此心安處是吾家

14:23耶穌回答他說:「凡愛我的人就會遵守我的道,我父也會愛他,並且我們要到他那裡去,與他同住.

14:24不愛我的人就不遵守我的道.你們所聽見的道不是我的,而是差我來之父的.

14:25「我還與你們在一起的時候,已對你們說了這些事.

14:26但保惠師,就是父因我的名所要差來的聖靈,他要把一切的事教導你們,並且要使你們想起我對你們所說的一切話.

14:27我留下平安給你們,我把我的平安賜給你們.我所賜給你們的,不像世人所賜的.你們心裡不要憂愁,也不要膽怯.

14:28你們聽見我對你們說過,我去了還要回到你們這裡來.你們若愛我,就會因我到父那裡去而喜樂,因為父比我大.

14:29現在事情還沒有發生,我預先告訴你們,使你們在事情發生的時候會信.

約翰福音14章23-39節的家在哪裡呢?就是在23節耶穌的應許裡面:「凡愛我的人就會遵守我的道,我父也會愛他,並且我們要到他那裡去,與他同住.」英文NRSV譯作we will come to them and make our home with them.14:29耶穌向門徒提示,有事情將會發生,會令人感到憂愁,感到膽怯.常言道:「人生不如意事十常八九」,在個人際遇而言,誰能保證天色常藍,花香常漫呢?在金融經濟而言,誰能長期「托住個市」使其永遠向上呢?在全球或地邊政治而言,那有長治久安,永遠和平呢?雖然這些都是我們的期盼,也是我們努力的目標,然而,在我們的意識之內,或者預計之外,總是「有事情將會發生的」.

幾年前有評論指責香港首富李嘉誠由「愛國」變成「不愛國」；由在內地「投資」變成「撤資」,質疑他究竟以何處為家?李嘉誠引用詩詞回應明志:「此心安處是吾家」.李氐家族的投資包括「貨櫃碼頭,能源,零售,水廠,電網,地產,火車,飛機租賃等業務,遍及52個國家」,當然不願意被捆綁在某一處地方,因此,回答何處為家時,好一句「此心安處是吾家」.問題是,世間何處能安心呢?耶穌對門徒說:「14:27我留下平安給你們,我把我的平安賜給你們.我所賜給你們的,不像世人所賜的.你們心裡不要憂愁,也不要膽怯.」世間無處能安心,然而,對某些人來說世間處處能安心,能安心之處,就是耶穌所應許的家:「凡愛我的人就會遵守我的道,我父也會愛他,並且我們要到他那裡去,與他同住.」有聖父與聖子同住的家,讓我們有平安的心靈迎向不安的世界,並在不安的世界締造平安.

3．這世界非我家／這世界是我家

詩篇67:1-7節:

67:1願上帝憐憫我們,賜福給我們,使他的臉向我們發光,(細拉)

67:2好讓全地得知你的道路,萬國得知你的救恩.

67:3上帝啊,願萬民稱謝你!願萬民都稱謝你!

67:4願萬族都快樂歡呼；因為你必按公正審判萬民,引導地上的萬族.(細拉)

67:5上帝啊,願萬民稱謝你!願萬民都稱謝你!

67:6地已經出了土產,上帝,我們的上帝,要賜福給我們.

67:7上帝要賜福給我們,地的四極都要敬畏他!

基督徒有時會說,「這世界非我家」,正如詩歌所唱:

這世界非我家,我停留如客旅,

我積財寶在天,時刻仰望我主,

天門為我大開,天使呼召迎迓,

故我不再貪愛這世界為我家.

我將永遠居住天上榮美福地,

千萬得贖聖徒一同歡呼勝利,

他們讚美歌聲,在天各方回應,

故我不在貪愛這世界為我家.

詩篇67篇的作者說:「67:6地已經出了土產,上帝,我們的上帝,要賜福給我們」,他卻似乎暗示,「榮美福地」不只是在天上,而是上帝所賜福的大地.他又說:「67:2好讓世界得知你的道路,萬國得知你的救恩」,似乎「千萬聖徒,選美歌聲」也不只是在天上,而是上帝的憐憫及公義所及之處.無論世界非我家還是世界是我家,「萬國」與「世界」帶給香港人的聯想,的確是與家有關的,因為那兩稱為「萬國」與「世界」的兩座建築物,它們的英文名是funeral home,而且是眾人最後的家.

4．啟示錄22:10,22-22:5:永遠的家

21:10我在聖靈感動下,天使帶我到一座高大的山,給我看由上帝那裡,從天而降的聖城耶路撒冷,

21:22我沒有看見城內有殿,因主-全能者上帝和羔羊就是城的殿.

21:23那城內不用日月光照,因為有上帝的榮耀光照,又有羔羊為城的燈.

21:24列國要藉著城的光行走；地上的君王要把自己的榮耀帶給那城.

21:25城門白晝總不關閉,在那裡沒有黑夜.

21:26人要將列國的榮耀尊貴帶給那城.

21:27凡不潔淨的,和那行可憎與虛謊之事的人,都不得進那城,只有名字寫在羔羊生命冊上的才得進去.

22:1天使又讓我看一道生命水的河,明亮如水晶,從上帝和羔羊的寶座流出來,

22:2經過城內街道的中央；在河的兩邊有生命樹,結十二樣 的果子,每月都結果子；樹上的葉子可作醫治萬民之用.

22:3以後不再有任何詛咒.在城裡將有上帝和羔羊的寶座.他的僕人都要事奉他,

22:4也要見他的面.他的名字將寫在他們的額上.

22:5不再有黑夜；他們也不需要燈光或日光,因為主上帝要光照他們.他們要作王,直到永永遠遠.

啟示錄告訴我們,funeral home其實不是我們最後的家,正如耶穌所應許,我們永遠家是上主與我們同住之處,是上主的憐憫及公義所及之處,是上主的光明照遍之處.

三．總結

如果,我們不單單以「四仔」來定義家,那麼,家就不一定要有「四仔」.有使命的地方也可以是家,有上主同在的地方也可以是家.就個人來說,雖然我們以「四仔」為家,然而,在公共的層面,卻有不少呼籲要求我們以那些地方為家,例如公司建議我們以公司為家,我們居住的城市勸誘我們以這城市為家,政權要求我們以國為家,如此種種,其中一個目的是要求我們對公司,或組織,或城市,或國家忠誠.如果,實踐上主的使命之處就是家,上主同在之處才是家,那麼,提示我們應當把我們的忠誠投放何處,托付於誰?



\newpage

\section{忽然從天上有響聲下來}
\label{sec:456}
\textbf{林豪恩先生 - 忽然從天上有響聲下來}
\newline
\newline
\hyperref[sec:455]{< < < PREV SERMON < < <}
~
\hyperref[sec:toc]{[返目錄]}
~
\hyperref[sec:457]{> > > NEXT SERMON > > >}


使徒行傳 2:1-21
\newline
\begin{longtable}{cl}
\hline
\hline
章節 & 經文 (和合本修訂版)\\
\hline
2:1 & \begin{tabularx}{0.7\textwidth}{X} 五旬節那日到了，他們全都聚集在一起。 \end{tabularx} \\ \\ \relax
2:2 & \begin{tabularx}{0.7\textwidth}{X} 忽然，有響聲從天上下來，好像一陣大風吹過，充滿了他們所坐的整座屋子； \end{tabularx} \\ \\ \relax
2:3 & \begin{tabularx}{0.7\textwidth}{X} 又有舌頭如火焰向他們顯現，分開落在他們每個人身上。 \end{tabularx} \\ \\ \relax
2:4 & \begin{tabularx}{0.7\textwidth}{X} 他們都被聖靈充滿，就按著聖靈所賜的口才說起別國的話來。 \end{tabularx} \\ \\ \relax
2:5 & \begin{tabularx}{0.7\textwidth}{X} 那時，有從天下各國來的虔誠的猶太人，住在耶路撒冷。 \end{tabularx} \\ \\ \relax
2:6 & \begin{tabularx}{0.7\textwidth}{X} 這聲音一響，許多人都來聚集，各人因為聽見門徒用他們各自的鄉談說話，就甚納悶， \end{tabularx} \\ \\ \relax
2:7 & \begin{tabularx}{0.7\textwidth}{X} 都詫異驚奇說：「看哪，這些說話的不都是加利利人嗎？ \end{tabularx} \\ \\ \relax
2:8 & \begin{tabularx}{0.7\textwidth}{X} 我們每個人怎麼聽見他們說我們生來所用的鄉談呢？ \end{tabularx} \\ \\ \relax
2:9 & \begin{tabularx}{0.7\textwidth}{X} 我們帕提亞人、瑪代人、以攔人，和住在美索不達米亞、猶太、加帕多家、本都、亞細亞、 \end{tabularx} \\ \\ \relax
2:10 & \begin{tabularx}{0.7\textwidth}{X} 弗呂家、旁非利亞、埃及的人，並靠近古利奈的利比亞一帶地方的人，僑居的羅馬人， \end{tabularx} \\ \\ \relax
2:11 & \begin{tabularx}{0.7\textwidth}{X} 包括猶太人和皈依猶太教的人， 克里特人和阿拉伯人，都聽見他們用我們的鄉談講論神的大作為。」 \end{tabularx} \\ \\ \relax
2:12 & \begin{tabularx}{0.7\textwidth}{X} 眾人就都驚奇困惑，彼此說：「這是甚麼意思呢？」 \end{tabularx} \\ \\ \relax
2:13 & \begin{tabularx}{0.7\textwidth}{X} 還有人譏誚，說：「他們是灌滿了新酒吧！」 \end{tabularx} \\ \\ \relax
2:14 & \begin{tabularx}{0.7\textwidth}{X} 彼得和十一個使徒站起來，他就高聲向眾人說：「猶太人和所有住在耶路撒冷的人哪，這件事你們要知道，要側耳聽我的話。 \end{tabularx} \\ \\ \relax
2:15 & \begin{tabularx}{0.7\textwidth}{X} 這些人並不像你們所想的喝醉了，因為現在才早晨九點鐘。 \end{tabularx} \\ \\ \relax
2:16 & \begin{tabularx}{0.7\textwidth}{X} 這正是藉著先知約珥所說的： \end{tabularx} \\ \\ \relax
2:17 & \begin{tabularx}{0.7\textwidth}{X} 『神說：在末後的日子，我要將我的靈澆灌凡血肉之軀的。你們的兒女要說預言；你們的少年要見異象；你們的老人要做異夢。 \end{tabularx} \\ \\ \relax
2:18 & \begin{tabularx}{0.7\textwidth}{X} 在那些日子，我要把我的靈澆灌，甚至給我的僕人和婢女，他們要說預言。 \end{tabularx} \\ \\ \relax
2:19 & \begin{tabularx}{0.7\textwidth}{X} 在天上，我要顯出奇事，在地下，我要顯出神蹟，有血，有火，有煙霧。 \end{tabularx} \\ \\ \relax
2:20 & \begin{tabularx}{0.7\textwidth}{X} 太陽要變為黑暗，月亮要變為血，這都在主大而光榮的日子未到以前。 \end{tabularx} \\ \\ \relax
2:21 & \begin{tabularx}{0.7\textwidth}{X} 那時，凡求告主名的都必得救。』 \end{tabularx} \\ \\
[1ex]
\hline
\hline
\end{longtable}
講題:忽然從天上有響聲下來 Suddenly There Came a Sound from Heaven

經文:使徒行傳2:1-21

講員:林豪恩先生

場合:香港中文大學崇基學院神學院容啟東校長紀念樓G/F新聖堂

日期:2019年6月9日

一．引言

四年前,2015年的聖靈降臨日,主日崇拜也是在這裡舉行.當時,我短暫離開了校牧室,被邀請回來講道,經文剛好也是使徒行傳2:1-21節,當時的講題是「聖靈降臨的新世界」.這個新世界一方面是一個來自東南西北四方人士能夠彼此溝通,互相聽得明白的新世界,另一方面是一個所有年齡群體,不但是成人,連兒童及長者都可以有夢想和參與的新世界.四年過去了,東方和西方似乎更聽不明白對方,或者更不願意聆聽對方；南面和北面的溝通不見得更順暢,往往只是權力的表述和較量.或者,有人認為這世界本來如此,世界的遊戲規則本來如此,人人都繼續埋首投進遊戲中,每天都仍然有人中獎,每年都仍然有人進身百萬富翁行列.人人都盼望有賭未為輸,若然未有回報,只是時辰未臨到.在人人都營營役役,返工放工的日子,其實有幾聲巨響出現過.這些響聲有否敲醒人們呢?人們又如何回應呢?聖靈降臨的當天,也有一聲巨響,那一聲巨響出現後,世界從此不再一樣.今天,我們用另一個角度重述這個聖靈降臨的故事.

二．經文

2:1五旬節到了,門徒都聚集在一處.

2:2忽然,從天上有響聲下來,好像一陣大風吹過,充滿了他們所坐的屋子,

2:3又有舌頭如火焰顯現出來,分開落在他們各人頭上.

2:4他們就都被聖靈充滿,按著聖靈所賜的口才說起別國的話來.

2:5那時,有虔誠的猶太人從天下各國來,住在耶路撒冷.

2:6這聲音一響,眾人都來聚集,各人聽見門徒用眾人的鄉談說話,就甚納悶；

2:7都驚訝希奇說:「看哪,這說話的不都是加利利人嗎?

2:8我們各人怎麼聽見他們說我們生來所用的鄉談呢?

2:9我們帕提亞人,米底亞人,以攔人,和住在美索不達米亞,猶太,加帕多家,本都,亞細亞,

2:10弗呂家,旁非利亞,埃及的人,並靠近古利奈的利比亞一帶地方的人,從羅馬來的客旅中,或是猶太人,或是進猶太教的人,

2:11克里特和阿拉伯人,都聽見他們用我們的鄉談,講說上帝的大作為.」

2:12眾人就都驚訝猜疑,彼此說:「這是甚麼意思呢?」

2:13還有人譏誚說:「他們無非是新酒灌滿了.」

2:14彼得和十一個使徒站起,高聲說:「猶太人和一切住在耶路撒冷的人哪,這件事你們當知道,也當側耳聽我的話.

2:15你們想這些人是醉了；其實不是醉了,因為時候剛到巳初.

2:16這正是先知約珥所說的:

2:17上帝說:在末後的日子,我要將我的靈澆灌凡有血氣的.你們的兒女要說預言；你們的少年人要見異象；老年人要做異夢.

2:18在那些日子,我要將我的靈澆灌我的僕人和使女,他們就要說預言.

2:19在天上,我要顯出奇事；在地下, 我要顯出神蹟；有血,有火,有煙霧.

2:20日頭要變為黑暗,月亮要變為血；這都在主大而明顯的日子未到以前.

2:21到那時候,凡求告主名的,就必得救.

三．內容

1．朝聖之日常

按律法的規定,我們作為猶太人中的男人每年應該三次去耶路撒冷過節.不過我們並不是本地猶太人,而是生活在外邦地方的猶太外僑,去耶路撒冷是不容易的.那一年,我和幾位叔伯兄弟從亞細亞內陸出發,經過以弗所,下到港口米利都,在那裡上了一艘開往帕大喇的船.第二日,我們在帕大喇轉往一艘開往腓利基的船,在船上碰見有些從希臘,甚至從地極之處的羅馬而來的同胞,大家都不約而同往耶路撒冷過節.船到達了泰爾,卸了貨,又上了一些旅客,最後來到凱撒利亞.我們幾位叔伯兄弟,加上本來在船上的人,以及後來上船的,共幾十人,從凱撒利亞一同上耶路撒冷.我們當中,有些已經在外地生活了好幾代,對祖宗的語言已經有點陌生,然而,大家連群結隊上耶路撒冷,也好有個照應.

我們沒有走沿海的路,聽說那裡從前是非利士人的地方.我們從凱撒利亞進入內陸,經過古時的以法連和便雅憫地,從山地進入耶路撒冷.在路上,我們遇見從不少本地人,有些從加利利而來,有些從低加坡里而來,聽他們說,耶路撒冷最近發生了一些事,有一個人名叫耶穌,被羅馬巡撫判了死刑.過後,又平靜了.其實,名叫耶穌的人多的是,就正如你們中國人叫國偉,瑞強一樣.在我們那個年代,羅馬統治者要維持帝國穩定,把任何反抗都消滅於萌芽狀態,以十字架死刑來侍候威脅社會安寧的人,不要說在耶路撒冷這個敏感的地方,在帝國內其他地方也不陌生.有些事情,發生了就發生了,當作茶餘飯後話題就算了,不用太認真,免招麻煩.畢竟,我們是來朝聖的,來盡我們作為猶太人的宗教義務.喜歡也好,不喜歡也好,就算是行禮如儀,有些事情總是要做的.我們在耶路撒冷找個地方棲身,第二天,一齊去聖殿獻祭.聖殿是我們的目的地,獻祭是我們上耶路撒冷要做的事情,也順道聽聽聖殿神聖的大祭司的教導,接受祝福,完成任務便回家.

你們的時代有一位女子叫鄭素素,她在她的書內寫著:「每天返工放工,OT,然後去飲酒,應酬,第二天逼地鐵上班,車上人人都帶著同一個垂死的表情.」聽說這就是你們這個時代的生活的寫照.這就是日常,日常累積起來的就是人生,被生了下來,就要繼續生存下去,千萬不要學你們城市中某位男子,他留下遺書表示每天返工放工的營役令他感覺生活枯燥而尋死.白頭人送黑頭人是傷心欲絕的.那位叫鄭素素的女子,她克服每天返工放工的枯燥的方法是培養跑步的習慣,並且著書立說教人跑步.

2．忽然從天上有響聲下來

將近聖殿的時候,突然聽到巨大響聲.我們以為是打雷聲,但天氣卻沒有變化,而且這個季節不應有如此巨大的打雷聲.這響聲太特別了,大家都東張西望,看看是不是某處有特別事情發生了.我們這些在外地來的人,雖然不是第一次到耶路撒冷,但也需要特別小心,免得節外生枝.就在聖殿附近,有一群人聚會,有些人好像在演講,另好些人在圍觀.演講的人無論是外表和口音,都像我們在路上遇到的加利利人,奇怪的是,我卻聽得明白他們的說話,甚至親切得好像我們在亞細亞說的話一樣.其他地方來的人,無論是從東邊帕提亞來的人,北邊本都來的人,南邊埃及來的人,或是西邊從羅馬來的人,無論是是猶太人,還是進猶太教的人,大家都說聽得明白那些人所說的話.他們都有我的感覺,就是那些人說的話親切得好像我們在外地生活所說的話一樣容易聽得到.

我們不知道有甚麼事情發生在這些人身上,也不知道他們經歷了甚麼事情,只聽到他們也在談論上帝,好像說那個我們在路上聽人說過被釘死在十字架上的耶穌復活了,又說那個復活的耶穌就是上帝應許的拯救者彌賽亞.今次上耶路撒冷看見的情況真有趣,平常我們是要進入聖殿聆聽教導的,今次信息竟然在聖殿外宣講；平時是聖殿的神職人員教導的,今次宣講的竟然是名不經傳的普通人；平時我們主要是來獻祭的,現在聽了他們這樣說,到底我們還要不要去獻祭呢?說笑而已,來了耶路撒冷不獻祭,如何回去向父老交待呢!不過,今次真的和以前不同了,是甚麼產生不同呢?難道就是剛才那巨大的響聲?

在過去幾年間,聽聞在你們的城市中也出現過幾次巨大的響聲.有一次巨響帶著催淚彈的味道,連續響了79日.又有一次巨響帶著魚彈的味道,夾雜著磚頭的飛舞.今天是六月九日,可能又是一次幾十萬人咆哮的巨響.這些響聲,人們聽到嗎?人們被喚醒過來嗎?有人認真對待這些巨響嗎?對個人而言,你的身體曾經發出響聲嗎?你的情緒曾經發出響聲嗎?你的關係曾經發出響聲嗎?你被喚醒過來嗎?你如何回應這些響聲呢?

3．有些人質疑響聲

人群中有人開始回應.群眾交頭接耳,有人驚訝「這是甚麼意思呢?」有人譏誚說:「他們無非是新酒灌滿了.」那些演講說耶穌就是基督的人,聽到後就反駁說:

2:15你們想這些人是醉了；其實不是醉了,因為時候剛到巳初.

2:16這正是先知約珥所說的:

2:17上帝說:在末後的日子,我要將我的靈澆灌凡有血氣的.你們的兒女要說預言；你們的少年人要見異象；老年人要做異夢.

2:18在那些日子,我要將我的靈澆灌我的僕人和使女,他們就要說預言.

2:19在天上,我要顯出奇事；在地下, 我要顯出神蹟；有血,有火,有煙霧.

2:20日頭要變為黑暗,月亮要變為血；這都在主大而明顯的日子未到以前.

2:21到那時候,凡求告主名的,就必得救.

群眾中有些人覺得他們說得有道理,有些人仍然質疑.質疑的人並不是聽不到響聲,只是認為世事紛亂,怎能分真假呢?正如你們的潮語所說:「認真就輸了」.

三十年前,在北方的某地,曾經有巨大的響聲發出,不但傳到來這處南方邊陲,更響徹整個地球.隨著時日過去,其實也有不少人質疑自己當年是否耳嗚,聽錯了聲音.是耳嗚還是果真有如此大的巨響呢?耳嗚也好,打雷也好,既然也說不清楚,不如放下過去,面向將來,努力面前吧.有些事情,不用太認真,反正認真也沒有甚麼好處.

4．這響聲出現後,就不再一樣了

想不到,有些人卻認真起來.這響聲彷彿敲醒了他們,令他們感到扎心.他們問那些傳講耶穌的人:「我們當怎樣行?」他們竟然聽從那些人的教導,悔改,受洗,並且立志「脫離這彎曲的世代」(2:40).聽說他們組成了一個新的群體,聚在一起,凡物公用,一齊生活,一齊吃飯.我和叔伯兄弟對他們的狀況一方面感到好奇,一方面也感到驚嘆:我們都不是小孩子,而是飽經生活磨鍊的成人,怎會輕易承認自己有悔改的需要呢?怎會輕易改變自己的想法呢?怎會輕易放棄「有競爭才有進步」的遊戲而投入另外一種互通有無,互相分享的生活模式呢?或者,這響聲對他們來說實在太巨大了,不但敲醒了他們,更敲破了世界灌輸給他們的信念,敲碎他們本來追隨的遊行規則.被敲醒了的人,不能再裝睡,正如保羅所說:「若有人在基督裡,他就是新造的人,舊事已過,都變成新的了.」(林後5:17)

在我們那個年代,這些認真的人是有代價的,越認真,代價就越大.我聽說,當時那些傳揚耶穌就是基督的帶頭人,沒有多少個能壽終正寢的.希望在你們的年代,認真的人會有較好的際遇.

5．有些人消滅響聲

有些人投入,有些人猶豫,這是可以理解的.然而,奇怪的是有些人的反應卻非常大.有一次,當那位名叫彼得和他的同伴約翰正向百姓說話的時候,祭司們,守殿官和撒都該人來了,他們下手拿住他們,並把他們押在拘留所.聽說是因為他們接受不了門徒傳揚耶穌復活的事,又接受不了有死人復活這樣的事.(4:1-3)那些門徒努力說明那次響聲是上帝的作為,上帝親臨他們身上,如果不是,他們不會有如此的改變,也不會有如此的勇氣和能力去建立悔改和相愛的群體.他們反問祭司們和守殿官,難道你們沒有聽到當天那巨響嗎?聖殿附近所有人到聽到,你們也在聖殿內外,怎會聽不到呢?祭司們和守殿官彼此相望,他們知道自己是聽到你,但他們決定否認有這響聲.最後,他們決定連承認聽到這響聲的人都要把他們除掉.他們不但要消滅曾經有響聲的事實,也要把堅持聽到響聲的人滅聲.

如果在你們二千多年後仍有這樣的事情,也不用感到驚訝,只是證明人類的基因遺傳得相當完整,或者證明人類進化的速度其實是相當緩慢的.但有一項測試你們可以留意,如果二千年後的你們仍然聽到當年那些人所宣講的耶穌是基督的信息,就證明滅聲行動是失敗的.所以,雖然有人要消滅那響聲,甚至除掉那些堅持說聽過響聲的人,然而,真的假不了,假的也真不了.響聲出現過就是出現過,不會被個人意願所塗抹的.

6．有些人假裝沒有聲過響聲

踏在回程的路上,我們仍在談論在耶路撒冷發生的事情.那些被巨響喚醒了的人後來怎樣呢?會被那些要否認那巨響的人消滅嗎?還是他們最後互相諒解,和平共存呢?在回程的人群當中,有些人曾經參加過那個宣稱耶穌就是基督的新群體,並且立志把這好消息帶回去居住地,在當地建立這些新群體.有些人的確這樣做了.不過,出乎意料之外,後來有位使徒叫保羅周遊列國傳道時,他們其中一些人卻攻擊保羅,指斥保羅褻瀆上帝及律法.我本來以為被喚醒了的人不能再裝睡,但原來被喚醒過,轉過頭來可以再沉睡.究其原因,可以是「覺今是而昨非」,可以是「棄明投暗／棄暗投明」,人在江湖,原因可以是多得數不清.你明的!

可怕的不但是要消滅響聲的力量,更是假裝沒有聽過響聲的力量.當年雷曼爆破了,很多人被喚醒,但過後又重投那個製造泡沫的遊戲,醉生夢死.槍聲響起,很多人被喚醒,但時不久之後,又重投那個權力遊戲,美化槍聲為爆竹聲.身體健康的警號響起,及時救回,不久之後又重投煙花之地,酒肉之中.大能不死,「執返條命」,不一定會有後福的.【明報專訊】美國20年前科倫拜恩校園槍擊案生還者之一尤班克斯上在18/5/2019在家中逝世,終年37歲,死因初步沒可疑但有待確認.他的家人表示,他最終在濫藥問題上「打敗仗」.

四．總結

今天的講章如果結束在上一個段落,會較為正面及有激勵性.因為二千年過去了,我們仍然聽到當年那些人所宣講的耶穌是基督的信息,證明滅聲行動是失敗的.然而,本講章卻畫蛇添足地加了最後一個段落,揭示有人會聽到響聲後會假裝沒有聽過,這是反高潮,如果帶來挫敗及無奈感,實在抱歉.不過,這才是我們的真相,這才是我們的現實.正因我們是如此軟弱,如此易變,我們時刻都需要從上主而來的恩典,就是那忽然後天上來的響聲,再次敲醒我們,把我們在世界的遊戲中玩得太投入之中拯救出來.



\newpage

\section{真理在途上}
\label{sec:457}
\textbf{鄧瑞強博士 - 真理在途上}
\newline
\newline
\hyperref[sec:456]{< < < PREV SERMON < < <}
~
\hyperref[sec:toc]{[返目錄]}
~
\hyperref[sec:458]{> > > NEXT SERMON > > >}


約翰福音 16:12-15
\newline
\begin{longtable}{cl}
\hline
\hline
章節 & 經文 (和合本修訂版)\\
\hline
16:12 & \begin{tabularx}{0.7\textwidth}{X} 「我還有好些事要告訴你們，但你們現在擔當不了。 \end{tabularx} \\ \\ \relax
16:13 & \begin{tabularx}{0.7\textwidth}{X} 但真理的靈來的時候，他要引導你們進入一切真理。因為他不是憑著自己說的，而是把他所聽見的都說出來，並且要把將要來的事向你們傳達。 \end{tabularx} \\ \\ \relax
16:14 & \begin{tabularx}{0.7\textwidth}{X} 他要榮耀我，因為他要把從我領受的向你們傳達。 \end{tabularx} \\ \\ \relax
16:15 & \begin{tabularx}{0.7\textwidth}{X} 凡父所有的都是我的，所以我說，他要把從我領受的向你們傳達。」 \end{tabularx} \\ \\
[1ex]
\hline
\hline
\end{longtable}
講題:真理在途上 Truth is on the way

經文:約翰福音16:12-15

講員:鄧瑞強博士

場合:香港中文大學崇基學院神學院容啟東校長紀念樓G/F新聖堂

日期:2019年6月16日

各位弟兄姊妹,早安.

過去一個禮拜,香港發生了很多事情,社會充滿張力,好像去到爆發的臨界點似的,令人擔心,令人心痛,令人心傷.但願天父憐憫這個城巿!

耶穌生活的那個世界,在羅馬帝國的統治之下.羅馬帝國是和平的,表面的和平背後,卻是殘酷的暴力.耶穌行走在羅馬的土地上,有12個門徒常常圍著他.他的門徒中,有奮銳黨人,這是勇武抗爭的人；有稅吏,這是在建制中「搵食」的人；有出賣耶穌的猶大,他可能是理想主義者,也可能是機會主義者；但更多是漁夫,平民百姓.無論他們是誰,他們在耶穌身上,看到一種獨特的生命,這生命創造著一個新的世界,這新的世界展現著神的臨在.

耶穌是誰?他憑什麼創造一個新的世界?這是一個什麼樣的世界?

就著這些問題,羅馬派來的總督彼拉多曾經與耶穌有一段對話.彼拉多問耶穌:「你是猶太人的王嗎?」耶穌的回答是:「我的國不屬這世界.」彼拉多再問:「這樣,你是王嗎?」耶穌回答:「你說我是王.我為此而生,也為此來到世間,特為給真理作見證.凡屬真理的人就聽我的話.」彼拉多問:「真理是什麼呢?」

耶穌是誰?

他是王,但不是人間權力那種意義上的帝王,他的王權不是屬於這個世界的,他的國度也不是政治意義上的國度,但他的確是創造一個新世界.

他憑什麼創造一個新的世界?

憑真理.耶穌的國度,是真理的國度.他以真理創造一個真理的世界.真理揭露虛假,對抗虛假,反對虛假.真理揭露「最低的武力」本是武力,真理揭露「苦口婆心的偽善」本是偽善,真理揭露「無道理的說詞」本是無道理.耶穌手中沒有武器,他只言說真理.他沒有什麼權力,他只見證真理.他亦沒有什麼政治意圖,他只堅持真理.

耶穌在言說真理中,創造一個真理的世界.真理的世界,是一個什麼樣的世界?

真理的世界,沒有地圖上的領域,沒有可量度的疆界.在真理被顯揚的地方,在虛假被揭露的地方,就是真理的世界.正如,光明顯露,黑暗被驅散的地方,就是光明的世界.耶穌說,「凡屬真理的人就聽我的話」,耶穌的真理世界,就是由堅持真理的人所組成的群體.這個群體的存在,彰顯了神的臨在.

彼拉多問:「真理是什麼呢?」

問得真好,真理是什麼呢?耶穌當時沒有回答這個問題.事實上,又怎能用三言兩語回答這問題呢?最怕是,耶穌答了幾句話,然後,後來的人便執著這幾句話,將之當成是「絕對真理」.手握「絕對真理」的人最可怕,因他會消滅任何敵擋他的真理的人.耶穌沒有回答,讓真理在沉默中開顯.

是時候看看我們今天的講道經文了,是今天崇拜的福音書經課:約翰福16: 12-15,透過這經文,我們領悟一下耶穌所講的真理.

約16:12　「我還有好些事要告訴你們,但你們現在擔當不了〔或譯:不能領會〕.

約16:13　只等真理的聖靈來了,他要引導你們明白〔原文是進入〕一切的真理；因為他不是憑自己說的,乃是把他所聽見的都說出來,並要把將來的事告訴你們.

約16:14　他要榮耀我,因為他要將受於我的告訴你們.

約16:15　凡父所有的,都是我的；所以我說,他要將受於我的告訴你們.」

這段經文是放在耶穌和門徒吃最後晚餐的場景下說的,耶穌知道自己的命運,為真理而背負的十架正等待著他.

耶穌說:「我還有好些事要告訴你們,但你們現在擔當不了〔或譯:不能領會〕.」

事情一言難盡,說盡了也無補於事,因有限的人在有限的歷史處境裡未能有領會事情真相的能力.耶穌說,事情還在開展,事情的真相會進一步揭露,要多等一會兒.耶穌當下面對十架的苦難,他當然曾向門徒言說十架,但門徒不能領悟,好好地一位善良的老師為何會受死.當耶穌一面講十架時,他們一面討論誰為大.門徒不能明白耶穌受苦的真實意義.耶穌死後,會復活,他也曾向門徒預告復活,但預告了又如何,當耶穌真的上十架時,他們在絕望中四散.耶穌復活後,門徒的人生和使命才正式開始.但當下的門徒,怎能理解他們在世界歷史的進程中將扮演著重要的角色.他們將有份於帶領歷史走向神聖的終點,就是神臨在人間,為萬有的主.他們怎能理解?怎能領悟?

事實上,耶穌講了很多了,講他背負人間的十架,講死後復活,講新的生命的出現,講新的世界的誕生,講他的再來,講他對萬物的審判,講神最終臨在人間,耶穌講了很多很多了,我們需要在歷史的進程中去領悟.

耶穌死後,會復活升天.升天後,會差派聖靈來接續他的真理工程.聖靈來,在歷史的過程中,幫助我們一步一步明白事情的真相.

約16:13　「只等真理的聖靈來了,他要引導(will guide)你們明白〔原文是進入〕一切的真理；因為他不是憑自己說的(will not speak on his own authority),乃是把他所聽見的都說出來(whatever he will hear he will speak),並要把將來的事告訴(will declare)你們.

約16:14　他要榮耀(will glorify)我,因為他要將受(will receive)於我的告訴(will declare)你們.」

短短兩句經文,用了八個「未來時態」的動詞.為了凸顯這些「未來時態」的動詞,我將這兩節經文誇張地重譯如下:「等到真理的聖靈來了,他將引導你們進入一切的真理,因為他將要說的,不是憑自己的；他將要說的,是他將要聽到的,並將向你們宣告要來的事.他將要榮耀我,因為他將領受屬於我的東西,並將宣告給你們.」

兩句經文,全部指向將來.耶穌不在人間,但他的靈在人間.耶穌不見了,但他的靈仍在.歷史會不斷發展,耶穌的靈(或聖靈)會伴隨著信徒,在歷史的過程中點明真理.這兩節經文,講了幾件事.

一,在暗晦的歷史過程中,有沒有真理的呢?有.

二,真理在哪裡?在耶穌的生命與說話裡.真理的聖靈,不言說自己,只言說從耶穌那裡聽到的,只言說從耶穌那裡領悟的.歷史中發生的事件無窮無盡,真理的聖靈就在不同的事件中讓人認出基督,也讓人認出敵擋基督的人與事.

三,人能說自己掌握了真理嗎?看來不能.在這經文裡,是聖靈,在歷史的過程中,解說耶穌的生命和說話,並將之向人顯明.活著的基督,在歷史的過程中,提示聖靈,他的生命和說過的話當如何理解.聖靈指向基督,基督引導聖靈.聖靈在歷史事件中點出基督的生命和說話,基督的心意在歷史事件中向聖靈顯明.在基督和聖靈的互動中,人是被動的,人不能說自己已握著不能變的真理.脫離了耶穌對聖靈的指引,脫離了聖靈對耶穌的揭露,人獨立地說自己握著真理,是很危險的.歷史在流動,要在流動的過程中,看出哪件事件讓基督的生命呈現出來,需要聖靈不斷的指引.有限的人,只能在聖靈中不斷摸索,去辨認基督在當下的呈現.我們必須謙虛,永遠地探索,真理總是在開顯的途上.

最怕的是,人不在真理前謙虛,反而是狂妄地認為自己掌握著真理.一個母親,若認為自己想的便是真理,那麼,對她有異議的孩子,必定是錯.若孩子反對她,則必不能讓步,否則便是縱容了孩子的錯.但,有沒有可能,是孩子對,而這位母親錯?聖靈能在誰人身上,讓耶穌的生命和說話活現世上?

約16:15　「凡父所有的,都是我的；所以我說,他要將受於我的告訴你們.」

在這裡,耶穌說,我的生命和說話,就是天父的生命和說話.聖靈將我的生命和說話告訴你們,其實就是顯明天父的心意.天地間,有一天父,或者說,有一天心,或者說,有一正氣.這是一有情的生命,這有情的生命充滿天地.耶穌將這充滿天地的正氣活現出來.就像空氣中充滿著電波,而唯有校正了頻道的接受器將這電波的內容顯露出來.耶穌是天父生命完美的接收器,他顯露出天父的有情生命,讓無形的天心在人的生命中顯現出來.天父的心,正是真理所在.

天父的心,不是客觀的命題,而是與人交往的生命.這交往,落實在耶穌的生命中.耶穌升天後,這交往延續在聖靈的工作上.人生漫漫長路,聖靈與我們同行,隨處點化,讓我們看到何人何事正揭露耶穌的生命和說話,何人何事正埋沒耶穌的生命和說話.在聖靈中,真理向我們揭露,耶穌的生命和說話向我們揭露,天父的心向我們揭露.這種揭露,不是叫我們明白一種道理,而是呼喚我們進入真理,活像基督,彰顯天心.我們是天心的流露,像基督一樣.我們渴望著有一天,天心大露,正氣漫延,神臨在人間.

過去一個禮拜,香港發生了很多事情,年輕人問了很多問題.為香港悲哀的人,問悲哀的意義?用心建設香港的人,問應如何思考自己的身分?在困局中的人,問盼望何來?渴望天國的人,問天國何在?願意回應神的人,問神的本性為何?找不到出路的人,問出路何在?受傷或被控的人,問神的看顧在哪裡?除了唱詩外,基督徒可以做什麼?耶穌的一生,對我們有何啟發?

問得真好,問的人有血有肉有靈魂.其實,這已不是問,而是生命的吶喊.

面對這一切問號,我不懂回答,我只是看到,天父的心在這一切吶喊中揭露自己.我只是看到,聖靈在人群中運行.聖靈會讓不同的人在不同的事件中看到耶穌,然後,各人在耶穌中明白真相,自己領悟出自己問題的答案.悲哀的人,會看到耶穌在獅子山頭,為這城巿悲哀.用心建設香港的人,會看到耶穌和他們一起,開展天國事業.在困局中的人,會與耶穌一同從墳墓裡復活,體會從神而來的新生命.渴望天國的人,看到人與人之間顯露了耶穌的愛,也知道天國就在這愛中.願意回應神的人,就看到耶穌正正因為要回應神,走在十架的路上.找不到出路的人,會在耶穌的十架中,看到人類的救贖.受傷或被控的人,看到自己與耶穌同釘十架,有份於耶穌的救贖工程.除了唱詩外,基督徒看到耶穌與所有受苦的人手牽手,一起前行.耶穌的一生,對我們有何啟發?耶穌的一生,教曉我們看到天父的真實,看到耶穌自己將天父的心實現出來的真實,看到耶穌所開展的天國的真實.

這一切是真的嗎?

在歷史的過程中,聖靈會引導我們領悟這一切都是真的.

但願:榮耀歸於聖父,聖子,聖靈.阿們.

(附記:本人在崇基禮拜堂講道的講章已結集成《荷花池畔春草綠》講章集,有興趣的朋友可購買.試讀本可參閱:\url{http://www.news.ccc.cuhk.edu.hk/do/wp-content/uploads/2019/04/CCC-SK-Tang-Sermon-Text-Sample.pdf}.訂購表:\url{http://www.news.ccc.cuhk.edu.hk/do/wp-content/uploads/2019/04/CCC-SK-Tang-Sermon-Order-Form.pdf})



\newpage

\section{你們找誰}
\label{sec:458}
\textbf{劉國偉先生 - 你們找誰?}
\newline
\newline
\hyperref[sec:457]{< < < PREV SERMON < < <}
~
\hyperref[sec:toc]{[返目錄]}
~
\hyperref[sec:459]{> > > NEXT SERMON > > >}


約翰福音 18:1-14
\newline
\begin{longtable}{cl}
\hline
\hline
章節 & 經文 (和合本修訂版)\\
\hline
18:1 & \begin{tabularx}{0.7\textwidth}{X} 耶穌說了這些話，就同門徒出去，過了汲淪溪。在那裡有一個園子，他和門徒進去了。 \end{tabularx} \\ \\ \relax
18:2 & \begin{tabularx}{0.7\textwidth}{X} 出賣耶穌的猶大也知道那地方，因為耶穌和門徒屢次在那裡聚集。 \end{tabularx} \\ \\ \relax
18:3 & \begin{tabularx}{0.7\textwidth}{X} 猶大領了一隊兵，以及祭司長和法利賽人的聖殿警衛，拿著燈籠、火把和兵器來到園裡。 \end{tabularx} \\ \\ \relax
18:4 & \begin{tabularx}{0.7\textwidth}{X} 耶穌知道將要臨到自己的一切事，就出來對他們說：「你們找誰？」 \end{tabularx} \\ \\ \relax
18:5 & \begin{tabularx}{0.7\textwidth}{X} 他們回答他：「拿撒勒人耶穌。」耶穌對他們說：「我就是。」出賣他的猶大也同他們站在一起。 \end{tabularx} \\ \\ \relax
18:6 & \begin{tabularx}{0.7\textwidth}{X} 耶穌一對他們說「我就是」，他們就退後，倒在地上。 \end{tabularx} \\ \\ \relax
18:7 & \begin{tabularx}{0.7\textwidth}{X} 他又問他們：「你們找誰？」他們說：「拿撒勒人耶穌。」 \end{tabularx} \\ \\ \relax
18:8 & \begin{tabularx}{0.7\textwidth}{X} 耶穌回答：「我已經告訴你們，我就是。你們若找的是我，就讓這些人走吧。」 \end{tabularx} \\ \\ \relax
18:9 & \begin{tabularx}{0.7\textwidth}{X} 這要應驗耶穌說過的話：「你所賜給我的人，我一個也不失落。」 \end{tabularx} \\ \\ \relax
18:10 & \begin{tabularx}{0.7\textwidth}{X} 西門‧彼得帶著一把刀，就拔出來，把大祭司的僕人砍了一刀，削掉了他的右耳，那僕人名叫馬勒古。 \end{tabularx} \\ \\ \relax
18:11 & \begin{tabularx}{0.7\textwidth}{X} 於是耶穌對彼得說：「收刀入鞘吧！我父給我的杯，我豈可不喝呢？」 \end{tabularx} \\ \\ \relax
18:12 & \begin{tabularx}{0.7\textwidth}{X} 那隊兵、千夫長和猶太人的警衛拿住耶穌，把他捆綁了， \end{tabularx} \\ \\ \relax
18:13 & \begin{tabularx}{0.7\textwidth}{X} 先帶到亞那面前，因為他是那年的大祭司該亞法的岳父。 \end{tabularx} \\ \\ \relax
18:14 & \begin{tabularx}{0.7\textwidth}{X} 這該亞法就是從前向猶太人忠告說「一個人替百姓死是有利的」那個人。 \end{tabularx} \\ \\
[1ex]
\hline
\hline
\end{longtable}
講題:你們找誰? Who is it you want?

經文:約翰福音18章1-14節

講員:劉國偉先生

場合:香港中文大學崇基學院禮拜堂

日期:2019年6月23日

引言

香港島銅鑼灣鬧市一棟大廈的天台,於2003年架起「耶穌是主」四個霓虹燈大字,到如今已經16年.這個巨型招牌,日間紅色,在晚間亮燈之後就更加耀眼,在九龍對岸都可以清楚看到.

這個霓虹招牌,往往觸發我去思想,我們期盼或者相信耶穌是我們身處的這個城市的主,但也許前題是,我們作為基督徒的,有否以耶穌為我們每日每夜的生活的主,又或者問:耶穌對我們來說到底是誰?

信仰旅程從來都是考驗.2003年豎起銅鑼灣「耶穌是主」招牌之時,正是香港近二十年來的一個困難時期,包括醫療,政治,經濟,都面臨很大的考驗.今年6月,香港經歷了異常的困難,我們聽到許多指控:掌權者出賣港人；示威人士是暴徒及暴動；警方過份使用武力；還有很多很多不同的評論和指控.

二千多年前的一個夜晚,從橄欖山腳的客西馬尼園開始,天地之主耶穌經歷了出賣,捉拿,捆綁,押解,盤問,侮辱,誣告,還有很多很多.我們就從客西馬尼園的一個小片段,去思想耶穌問那些捉拿祂的人的一個問題:「你們找誰?」,或者對我們所身處的處境,可以帶來一些提醒與感動.

捉拿耶穌的人

約翰繼續用他細膩的筆觸,仔細描寫耶穌在客西馬尼園被捉拿的過程.

「耶穌說了這話,就同門徒出去,過了汲淪溪.在那裡有一個園子,他和門徒進去了.賣耶穌的猶大也知道那地方,因為耶穌和門徒屢次上那裡去聚集.猶大領了一隊兵,和祭司長並法利賽人的差役,拿著燈籠,火把,兵器,就來到園裡. 」(約18:1-3)

當晚,猶太教領袖安排了大量兵丁,差役去捉拿耶穌,真有這個必要嗎?到底他們如此部署,擔心什麼?想達到什麼目的?

約翰筆下所形容的一隊兵,希臘原文是Speira,是羅馬軍隊的一種編制,可以表示一個600人的軍隊,也可以表示擁有240個騎兵和760個步兵的1000人的軍團.無論是600人還是1000人,如臨大敵來捉拿一個手無寸鐵的拉比,實在叫人難以想像!約翰更告訴我們,除了羅馬兵丁,也包含了大祭司和法利賽人的差役,那應該就是聖殿的差役,也就是猶太教領袖所藉以維持聖殿和猶太律法的治安力量.

大祭司該亞法說過:你們不知道什麼!獨不想一個人替百姓死,免得通國滅亡,這就是你們的益處(約12:50) 為了維持他們心目中的安定和繁榮,自從拉撒路復活以後,他們就密謀商議要殺死耶穌,所以祭司長和法利賽人早已通報他們的線眼,若有人知道耶穌在哪裡,就要報明,好去捉拿耶穌.

猶太教領袖為要萬無一失,竟然動員了如此不合比例的武裝力量,務求在這一個夜晚,以極快速及極有效的方式,short and sharp,去捉拿耶穌,大概預期耶穌和祂的門徒會反抗及逃跑罷.

你們找誰?

猶太教領袖對付耶穌的手段和動機是明顯不過的,他們只不過裝胸作勢,然後要跑過場,把耶穌帶去前大祭司亞那的院子進行盤問,然後押往現任大祭司該亞法的院子,做一些程序,找一些証人,虛構一些明眼人都知道不真實的理據,又企圖找法律的支點,本著他們的位份和權力,來定性耶穌所作的事,務求速戰速決,以防計劃生變.試想想,星期四晚的深夜在客西馬尼園捉拿耶穌,然後在星期五的凌晨時分,押解耶穌到亞那和該亞法的院子進行盤問和審訊,天亮之後,再召開猶太議會審問耶穌,再經過彼拉多,希律王,彼拉多的幾番折騰,星期五下午三時左右,耶穌就死在十字架上.速度及效度,short and sharp,成為掌權者的完美計策.

兵丁和差役沒有想過耶穌會是第一個站出來的人吧!耶穌面對著拿燈籠,火把,兵器,並穿上全副軍服的數以百計甚至上千的羅馬兵丁,在黑夜的火光之中,祂勇敢地站前來問:你們找誰?當捉拿的人回答說「找拿撒勒人耶穌」,耶穌說「我就是」.「我就是」這個回答充滿了權柄,其實是約翰福音的一個寫作特色.

在約翰福音總共記載了耶穌七句「我就是」的宣告.在6章35節耶穌說「我就是生命的糧」；在8章12節耶穌說「我是世界的光」；在9章5節再一次重申祂是「世上的光」；第10章耶穌說「他是羊的門及好牧人」；第5個「我就是」是拉撒路復活事件中耶穌說「復活是我,生命是我；信我的人,雖然死了也必復活.」在14章6節耶穌說「我就是道路真理生命」；15章第1節記載了耶穌說我是「真葡萄樹」.

約翰記載了七次耶穌的自承「我就是」之後,在18章這段經文記載了另一次最震撼的回覆「我就是」,這個回覆涵蓋著18章4節約翰的描述:「耶穌知道將要臨到自己的一切事」.「我就是」除了是一個簡單的回覆,讓別人辨認自己,對身處客西馬尼園的耶穌來說,其實更加是一種承擔,耶穌已經準備了被拘捕,受審,受折磨,並且死於十字架上.

耶穌一說「我就是」,捉拿耶穌的人就戲劇性地退後,跌倒了,這個細節刻劃出耶穌站出來的勇敢,那是汗如血滴的儆醒禱告之後的勇氣,絕非不理智的匹夫之勇,卻把他們嚇壞了,一時間不知應對.當他們仍未回過神來,耶穌再問「你們找誰?」他們奉上相同的答案是要找拿撒勒人耶穌,耶穌再回答「我就是」,祂心繫門徒,就請兵丁放過祂的門徒.

跟隨耶穌的門徒,這刻大概知道耶穌準備束手就擒,跟從夫子三年了,在一股更新猶太教信仰的旅程中,經歷非凡,幾天之前剛剛經歷了光榮進入耶路撒冷,被百姓前呼後擁,怎麼在這個深夜,耶穌好像投降一樣!於是,不甘於如此局面的西門,就用刀砍掉了大祭司僕人馬勒古的右耳.我們當然明白彼得的行為於事無補,耶穌就提醒彼得:「父所給我的那杯,我豈可不喝呢?」按照路加的記載,耶穌治好了馬勒古的耳朵,而彼得就收刀入鞘,眼看耶穌被人帶走.

誰出賣了誰

耶穌的勇氣和承擔,並清晰的自承「我就是」,在這段經文中,與約翰筆下的猶大形成了文學上的人物對比.在約翰筆下,對猶大的特寫共有三次:

猶大知道耶穌這個深夜在那裡,是他提供地點情資,與猶太教領袖合作,進行捉拿行動.(約18:2)

猶大領頭,帶著數以百計的武裝去捉拿耶穌(約18:3),路加更告訴我們猶大以親吻作暗號出賣耶穌.(路22:48)

約翰更特別提醒讀者,猶大不只提供情資,領路,他其實在捉拿的過程中一直在現場(約18:5)

約翰要我們注意什麼?…

猶大出賣耶穌十分賣力,交足貨,參與了一次有組織,有預謀的拘捕行動；

猶大的自私(是個賊)與虛偽(門徒當中沒有人知道他要作什麼,他裝作若無其事),對比耶穌迎向十字架的勇氣,構成一個很大的角色對比和落差.

作為讀這段聖經的人,不禁會問這怎麼可能發生?猶大以耶穌來換取自己的利益,即使夫子為他洗腳,一再規勸和警告,三載恩情,竟可以無恩如此(disgrace)?

對我們來說,出賣耶穌的現實含義是什麼?耶穌教導我們天天背起十字架跟隨祂；叫我們服侍弟兄中一個最小的；教導我們好撒瑪利亞人的故事；教導門徒一粒麥子死了落在土裡的比喻……親愛的弟兄姊妹,若果我們沒有行道,甚至背道而馳,又何嘗不是對信仰的某一個層次的出賣?

猶大出賣了耶穌,作出了一個無可挽回的選擇,在客西馬尼園的夜晚,他竟然帶著想要殺害耶穌的人,去捉拿他跟隨了三年的夫子耶穌,為要成就自己的慾望,最終把自己推向生命的絕路.

約翰告訴我們,那一隊兵和千夫長並猶太人的差役,把耶穌捆綁起來,送往大祭司該亞法的岳父那裡,亦即是前任大祭司亞那的院子,作者再次提醒我們這個該亞法就是曾經說過「一個人代替人民死,這是有益的」那一個人.

付錢給猶大的猶太教領袖,其實嚴重違反了他們自己制定的司法程序規則,包括公會不得在安息日或節期的前夕,就重大案件召開會議；不得在審判當日就裁定死刑；已經不要說他們所呈上的是假證據和假見證,還有在公會審問期間,耶穌所受到的羞辱和暴力對待.

與其說猶大出賣耶穌,我更加認為猶大出賣了自己,大祭司和猶太公會領袖又何嘗不是出賣了公會的神聖和使命,並他們自己的靈魂?

我們找誰

若果耶穌今天問你問我「你們找誰?」我們會怎樣回答他?我們想找到的那位耶穌到底是一個怎樣的耶穌?我們是要找到按我們自己心願而塑造出來的耶穌,還是那位道成肉身,透過聖經向我們啟示祂自己的耶穌?

我們找誰?你所跟隨的主,是否生命的糧?世界的光?是羊的門及好牧人?是復活和生命?是道路真理生命?是真葡萄樹?是那位定意將自己成為一粒落在地裡的麥子,死了,就長出許多子粒的主?

弟兄姊妹,在這個困難的境況與時刻,讓我們仰望耶穌,確認祂是我們的主,在客西馬尼園的黑暗的晚上,前景模糊不清之際,原來最重要最優先不是要選定跟隨或者不跟隨猶太議會,支持或不支持大祭司該亞法,最要緊的是要辨識耶穌對我是誰,然後我要怎樣跟隨祂.

銅鑼灣耶穌是主的霓虹招牌,今夜將會仍然耀眼,但願這四個字,可以照亮我們往後的道路! 客西馬尼園的剎那間,門徒四散所帶著的迷惘與無力感,讓客西馬尼園看似是通向死亡的絕路,誰知其實是復活的開端.



\newpage

\section{陶造新世代}
\label{sec:459}
\textbf{朱光華牧師 - 陶造新世代}
\newline
\newline
\hyperref[sec:458]{< < < PREV SERMON < < <}
~
\hyperref[sec:toc]{[返目錄]}
~
\hyperref[sec:460]{> > > NEXT SERMON > > >}


詩篇 78:1-
\newline
\begin{longtable}{cl}
\hline
\hline
章節 & 經文 (和合本修訂版)\\
\hline
78:1 & \begin{tabularx}{0.7\textwidth}{X} 我的子民哪，要側耳聽我的訓誨，豎起耳朵聽我口中的言語。 \end{tabularx} \\ \\
[1ex]
\hline
\hline
\end{longtable}
講題:陶造新世代 Nurture the New Breed

經文:詩篇78:1-

講員:朱光華牧師

場合:香港中文大學崇基學院禮拜堂

日期:2019年6月30日

引言:

亞薩的訓誨詩

祖宗子孫,承先啟後

一脈相承,代代相傳

確定目的,改善手法

實踐方向

充實這一代,教導下一代

教會與家庭,並肩齊作戰

這幾個星期香港經歷一連串翻天覆地的事件,遊行,示威,激盪的人群,混亂的場面,即使我們不是身在其中,單看新聞,已足夠令我們心情起伏不定,有失望,有憤怒,有無奈,有疲倦,甚至恐懼.

鼓勵大家要留意自己的情緒,如果很困擾,腦海經常浮現一些畫面或者事件,揮之不去,建議大家搵老師搵校牧傾下,或者有專業輔導的專家幫下手,甚至要暫時中斷接收相關的消息,暫停唔睇新聞唔上網,讓我們的思想安靜沉澱一下.

多謝校牧的邀請和安排,讓我能夠離開崇基將近40年,可以重新踏足崇基禮拜堂,在這裡崇拜,甚至講道.

今日我想藉著詩篇七十八篇,跟大家分享一個題目,「陶造新世代」,78篇很長,共72節,我們只是集中1-8節.

第1節開始,作者這樣叮囑勉勵讀者,詩七十八:1,「我的民哪,要側耳聽我的訓誨,豎起耳朵聽我口中的言語」,

英文翻譯是Give ear, O my people.誰可以稱「我的民哪」O my people?這種口吻彷彿是皇帝向自己的子民,頒佈聖旨.

如果我們讀這詩篇是領聖旨,我們應該以什麼態度來聽呢?經文清楚指示我們要點樣聽.要側耳專心聽,要豎起耳朵留心聽,要必恭必敬,因為聖旨只是講一次,不會重覆,而且聽完聖旨之後,要領命領旨之後,要照著執行,我們恐怕聽漏了一句半句,甚至要做筆記.

嘩,讀這首詩要那麼隆重,那麼認真,讀者自然會產生期望,接下來的信息非常重要.

詩七十八:2 - 6,「2我要開口說比喻,我要解開古時的謎語,3是我們所聽見,所知道的,也是我們的祖宗告訴我們的.4我們不將這些事向他們的子孫隱瞞,要將耶和華的美德和他的能力,並他奇妙的作為,述說給後代聽.5因為,他在雅各中立法度,在以色列中設律法；是他吩咐我們祖宗要傳給子孫的,6使將要生的後代子孫可以曉得；他們也要起來告訴他們的子孫,」

我們留意到,這段經文有兩組重覆的詞句:第一組是祖宗,祖先,前輩,第二組是子孫,後代.這兩組的詞點出了兩個方向,一是承先,一是啟後.

承先,是我們要回顧上一代,回憶我們的祖先,回顧我們過往歷史,要追溯要重尋我們的根,要思想他們的嘉言懿行.

這是承先,還有另一個向度.啟後,是向下傳遞信息,教導教訓我們的後代,啟發啟導子孫,開導引導後輩,甚至訓導誘導下一代.

這種承先啟後,是一脈相承,代代相傳,值得留意,究竟經文列出幾多代?

詩七十八:5 - 6,「5他在雅各中立法度,在以色列中設律法；是他吩咐我們祖宗要傳給子孫的,6使將要生的後代子孫可以曉得；他們也要起來告訴他們的子孫,」

他,就是耶和華上帝,第一代

我們祖宗﹝our forefathers﹞,第二代

傳給子孫,就是我們讀者這一代,祖宗傳給我們這第三代

子孫,將要生的後代子孫﹝the children yet to be born ﹞,第四代

他們﹝剛才所提及的第四代﹞也要告訴他們的子孫,第五代

如果用一個圖表來表示,就很清楚,

第一代

耶和華上帝

第二代

我們的祖宗

第三代

我們

第四代

我們的子孫

第五代

我們子孫的子孫

祖宗子孫,承先啟後,

一脈相承,代代相傳,

美國芝加哥三一神學院,有一位好出名的新約學者,Don Carson卡遜,他已經退休,他的釋經講道,很受歡迎,曾經被邀請來香港的港九研經培靈會的大會講員,有一次他接受訪問,問及他教養兒女的心得,

「你認為你教養兒女是否成功?」

「雖然我有用心教養我的兒女,雖然我的兒女都接受我的教導,我也見到他們有按聖經真理行事,但我仍未知是否成功!」

卡遜能夠咁樣做,照計已很難得.

「那麼你何時才能確定你教養兒女是否成功?」

「要等,要等到我的兒女結婚,等他們成家立室,等他們有自己的兒女,而且能夠將我教導他們的聖經,教導他們的兒女,到那時,我才知道我的教養兒女成功.」

是的,一代傳一代,

陶造新世代,祖宗子孫,承先啟後,

陶造新世代,一脈相承,代代相傳,

陶造新世代,經文確定教導的目標和手法,

詩七十八:6 - 7,「6使(so that)將要生的後代子孫可以曉得；……7好叫(so that)他們仰望神,不忘記神的作為,惟要守他的命令.」

經文提及三個動詞,教導他們之後,希望達致的目標,他們能夠產生三個動作,

曉得﹝know﹞,

仰望﹝trust﹞,

守命令﹝keep the commandment﹞,

這三個動作有一個進程,先曉得,才能夠仰望,才會守命令.

我看過一齣電影,是劉德華與張學友的對手戲,劉德華責備張學友的一段對白,

講你又唔聽

聽你又唔明

明你又唔做

做你又做錯

錯你又唔認

認你又唔改

我們當中做老師的,一定非常明白,有教唔等於有聽,有聽唔等於明白,明白唔等於會做,有做又唔等於唔會做錯.

如果用詩78篇的教訓,要使他們曉得,使他們明白,理解,融會貫通,使他們仰望,使他們信靠,使他們同意擁抱這些命令,使他們喜歡這些命令,使他們甘心情願的守命令,使他們能夠化成行為,成為實踐,成為習慣.

有一次,有一個媽媽向我這個牧師投訴他的兒子:「我俾我個仔激死,佢好頑皮,又唔讀書,又唔肯返教會,重顧住打機,教極都唔聽!」

「你怎樣教他?」

「我話佢囉!我話佢好多次架啦!硬係唔聽.」

「話佢」,即是責備他,罵他,囉唆他,斥責他.這媽媽以為責備他罵他,甚至罰他,就等於教曉他.其實我們知道,責備他罵他不等於就係教育了他.於是我問那位媽媽,

「那麼你多次責備他罵他,責備他罵他這麼多年,是否奏效?有沒有效果?」這其實是明知故問的問題.

「當然不奏效!」

「既然不奏效,為什麼你仍選擇堅持繼續用這不奏效的方法教他?」

有些方法,是不能帶來培育陶造的效果,栽培下一代,要思考培育他們的方式,詩七十八篇,提及好幾種不同的教導方式.

「1我的民哪,你們要留心聽我的訓誨,側耳聽我口中的話.2我要開口說比喻；我要說出古時的謎語,……4要將耶和華的美德和他的能力,並他奇妙的作為,述說給後代聽.5他在雅各中立法度,在以色列中設律法；」

這裡提及幾種方式:

1節,訓誨,我口中的話,直接正面的教導教訓,

2節,比喻,謎語,猜謎語,用懸疑的手法,用比喻,透過有趣的學習方式,活動教學,遊戲中學習,

4節,耶和華的美德和他的能力,並他奇妙的作為,上帝在以色列百姓中的作為,包括歷史,傳記,英雄事蹟,悲劇故事(例如參孫),出人意表的事件,上帝幽默手法,上帝如何憐憫人,如何與世人同行,

5節,立法度,設律法,也就今我們所講的set boundary,劃界線,設立範圍,

其實這林林總總的培育方法,新約也有提及,提後三:16,「16聖經都是神所默示的,於教訓,督責,使人歸正,教導人學義都是有益的,」

聖經是有益的,不過特別注明在四方面有益的,在教訓,在督責,在使人歸正,在教導人學義這四方面,都是有益的有用有好處的,都係useful,profitable,beneficial,

教訓:teaching - directing you where to go告訴你何去何從,教導你那條路是正路

督責:rebuking - telling you where you have gone wrong老老實實的指出,你是如何走錯路走歪路

使人歸正:correcting - showing you the right path  and  how to get back 在你面前展示一條回轉歸正的路,並且幫助你怎樣糾正重返正途

教導人學義:training - keeping you in the right direction, 如果你沒有走歪路,如果你站在正確的路徑上,聖經會訓練你如何保持在這正確的方向,持守聖潔

其實詩篇七十八篇不單教導觀念,道理,這首詩更示範demonstrate如何教導,不單是命題propositions,而且還要傳遞歷史.詩78篇,從第9節開始,記載耶和華神拯救以色列人的歷史,從出埃及,過紅海,到曠野,有食物有食水,到應許之地,一件一件的事蹟.其實就是講故事,好像用話劇的方式,一幕一幕鋪陳出來,演譯出來.

中國人講故事的格式,有起承轉合.西方人講故事的格式不同,可以用5個T來歸納:有Tale故事,tension張力,矛盾,角色衝突,twist峰迴路轉的轉折位,teaching這故事教訓我們,講故事高手,能夠繪影繪聲的將故事的tone氣氛都講到出來

那麼多不同的培育方式,如何決定用什麼方法教導?那就要審時度勢,要隨機應變,要留意栽培的對象處於什麼狀態,他有什麼需要.

例如這幾個星期,外面山雨欲來,風雨飄搖,如果你的兒子經過一整天遊行示威回來,這個時候似乎不應講太多道理,更唔不應說「la la la ,話jor 你架啦,人多的地方唔好去,叫你唔好去搞事,而家焦頭爛額」,

這個時候,對方需要的是愛,是關懷,是鼓勵,是慰問,是拍一拍膊頭的支持,是體諒的眼神.

正如路十五,浪子故事的父親,浪子在外面遇到一連串的逆境,麻煩,困難,挫敗,現今回來,爸爸遠遠見到,跑去,擁抱,連連親咀,吩咐下人,tong 雞殺鴨,慶祝一翻,3個open:Open door, open arm, open heart,

Open door,我們家的大門永遠為你打開,

Open arm,用「廣闊的胸襟,強勁的臂彎」,擁抱他接納他歡迎他,

這些動作,代表開放的心靈,open heart.

經文透過了這三個要點,陶造新世代:

祖宗子孫,承先啟後

一脈相承,代代相傳

確定目的,改善手法

究竟這段經文,對我們有什麼可作借鏡反思呢?應該點樣實踐?我想同大家分享兩點.

第一,充實這一代,教導下一代,

如果要教導下一代,我們必須先充實這一代,我們自己有沒有學聖經?若果我們自己有沒有查經,沒有上主日學,我們憑什麼理由道理,去說服兒女要上主日學?我們自己唔學習聖經,卻勉強他們地要學習聖經,他們會覺得,我們成年人只是講一套做一套,他們會輕看我們的上帝,鄙視我們的聖經.

而且我們可以強迫多久,年輕時還可以強迫,高中,大學,他們獨立,他們有他們自己的生活,我們還可以勉強嗎?

我們在座的父親母親會問:牧師,我何嘗不知道要勤讀聖經,但是年紀大,讀極都唔入腦,怎辦?

在農村,有一個年長的信徒去找牧師,「牧師,我讀聖經,讀完就忘記,人家讀聖經,過目不忘,我讀聖經過目即忘.怎辦?」

「我現在給你一個筲箕,你去河邊,用筲箕裝滿一桶水回來給我.」

年長的信徒一頭霧水,不過很順服,照做,過了一個星期,這個信徒來見牧師,

「牧師,這個筲箕漏窿,不能盛水.」

「讓我看看那個筲箕.……前一個星期,我給你的筲箕,很骯髒,很多污漬,但經過你這個星期盛水,是的,不能盛到水,但筲箕變得很乾淨.

同樣,如果你每日都讀聖經,你可能轉眼就忘記,但你的生命得潔淨,你的心靈得洗滌.那就夠了!」

弟兄姊妹,我們作為上一代,我們有沒有學?有沒有領受?我們有沒有讓聖經潔淨洗滌我們的生命?我們是否享受學習聖經?

除了充實這一代,教導下一代,要建立下一代,

第二,教會家庭,雙管齊下,

教會家庭,要並肩作戰,我們做父母,都會將孩子帶到教會,希望教會能夠幫助他們的兒女成長,有好的道基督教育,有質素的品格教育,傳遞他們正確的價值觀,甚至能夠有一班好朋友,建立友誼友情.這是好的,但不足夠!

我跟大家計一條數.一個星期168小時,我們的兒女在教會多久?極其量只是3小時,4小時,他們在家有多少小時?最少都廿小時,三十小時.請問那個地方比較容易改變他們,教育他們,影響他們?是在教會的主日學老師?抑或是在家中的父母?

不是說參加主日學不重要,但沒有家庭教育,主日學的蜻蜓點水式的教導,可以給予什麼?有什麼長遠果效呢?

所以教導下一代,除了教會,還需要家庭父母,父母教導下一代,我們需要四隻手,

要動手,不要只是動口,不能只是講,這不可,那不可,不可玩水,不可爬高,不可快跑.

有一個小故事.父親問3歲的孩子,貓貓怎麼叫:喵喵喵；狗狗怎麼叫?汪汪汪；媽媽怎麼叫?no no no.

只是造訴孩子NO是不足夠的,要實踐,要有行動,不是單講做人要有信用,而是要父母親身實踐,有誠信的行動,我們要實踐誠信.

要攜手,不可以只是媽媽教導孩子,爸爸躺在沙發看電視,打電腦.也不可以只是父親負責訓誨,而母親只負責換片洗衣煮飯做菜,要父母同心同步同行,教導培養兒女,是父親母親的共同責任,甚至不可以父母親二人全心全意工作賺錢,由工人姐姐承擔教導,由公公婆婆,爺爺麻麻教導,

要放手,當孩子漸長,要讓小孩子嘗試,甚至讓他們碰下釘,撞下板,闖小禍,以致他們能夠從錯誤中學習.

我見過很多年青人,他一生都未試過失敗,因為一直以來,但凡出問題,他們的父母就挺身而出,為他們排難解紛,結果剝奪了他從失敗從錯誤中學習的機會.

要舉手,我的女兒三十多歲,結婚,身在南半球,我即使關心她,要培育她,也不可能,我只能為她舉手祈禱,為女婿祈禱,為他們的婚姻,為他們的家庭,為他們的專業工作,為他們的教會事奉祈禱.

要為孩子祈禱,為下一代仰望神,人世間很多事,不是我們做父母可以控制,有一日我們的孩子要自己出去闖,我們只能及早將兒女交託上帝,讓神塑造佢地.

充實這一代,教導下一代

教會和家庭,並肩齊作戰

有一次,一位女士帶著一個六個月大的嬰孩來到教會聚會,要求牧師為這嬰孩舉行嬰兒洗禮,牧師同意,但有一個條件,嬰孩的父母必須表示願意用真理培育這孩子.

但原來這孩子是私生子,孩子的父親已經不知去向.這母親感到很為難,因為只有她單身一個人.但牧師鼓勵她,雖然沒有父親,照樣在所定的洗禮的日子,帶來孩子來教會,為嬰孩作洗禮.

到了孩童洗禮當日,牧師為這母親手中抱著的嬰孩作洗禮,牧師先向這媽媽問一個問題:

誰願意將這孩子獻給神?

誰願意在這孩子有需要,困難的時候扶育他,

誰願意用聖經真理教養這孩子?

誰願意帶領這孩子歸向神?

誰願意不斷為這孩子的成長祈禱?

當這個單親媽媽正要開口說「我願意」的時候,全體會眾站起來,大聲說「我們願意」.

西人有一句諺語,It takes a Whole Village to Raise a Child我們需要整條村的人,一齊培養一個孩子.

弟兄姊妹,我們需要整個教會的人,uncle auntie,哥哥姐姐,甚至公公婆婆,一齊黎培養一個孩子.即使我們單身,即使我們兒女成長離家,但我們總可以參予,為主陶造下一代.



\newpage

\section{你為甚麼問我呢}
\label{sec:460}
\textbf{劉國偉先生 - 你為甚麼問我呢?}
\newline
\newline
\hyperref[sec:459]{< < < PREV SERMON < < <}
~
\hyperref[sec:toc]{[返目錄]}
~
\hyperref[sec:461]{> > > NEXT SERMON > > >}


約翰福音 18:15-27
\newline
\begin{longtable}{cl}
\hline
\hline
章節 & 經文 (和合本修訂版)\\
\hline
18:15 & \begin{tabularx}{0.7\textwidth}{X} 西門‧彼得跟著耶穌，另一個門徒也跟著；那門徒是大祭司所認識的，他就同耶穌進了大祭司的院子。 \end{tabularx} \\ \\ \relax
18:16 & \begin{tabularx}{0.7\textwidth}{X} 彼得卻站在門外。大祭司所認識的那個門徒出來，對看門的使女說了一聲，就領彼得進去。 \end{tabularx} \\ \\ \relax
18:17 & \begin{tabularx}{0.7\textwidth}{X} 那看門的使女對彼得說：「你不也是這人的門徒嗎？」他說：「我不是。」 \end{tabularx} \\ \\ \relax
18:18 & \begin{tabularx}{0.7\textwidth}{X} 僕人和警衛因為天冷生了炭火，站在那裡取暖；彼得也同他們站著取暖。 \end{tabularx} \\ \\ \relax
18:19 & \begin{tabularx}{0.7\textwidth}{X} 於是，大祭司盤問耶穌有關他的門徒和他教導的事。 \end{tabularx} \\ \\ \relax
18:20 & \begin{tabularx}{0.7\textwidth}{X} 耶穌回答他：「我一向都是公開地對世人講話，我常在會堂和聖殿裡，就是猶太人聚集的地方教導人，我私下並沒有講甚麼。 \end{tabularx} \\ \\ \relax
18:21 & \begin{tabularx}{0.7\textwidth}{X} 你為甚麼問我呢？去問那些聽過我講話的人，我所說的，他們都知道。」 \end{tabularx} \\ \\ \relax
18:22 & \begin{tabularx}{0.7\textwidth}{X} 耶穌說了這些話，旁邊站著的一個警衛打了他一耳光，說：「你這樣回答大祭司嗎？」 \end{tabularx} \\ \\ \relax
18:23 & \begin{tabularx}{0.7\textwidth}{X} 耶穌回答他：「假如我說的不對，你指證不對的地方；假如我說的對，你為甚麼打我呢？」 \end{tabularx} \\ \\ \relax
18:24 & \begin{tabularx}{0.7\textwidth}{X} 於是亞那把耶穌綁著押解到大祭司該亞法那裡。 \end{tabularx} \\ \\ \relax
18:25 & \begin{tabularx}{0.7\textwidth}{X} 西門‧彼得正站著取暖，有人對他說：「你不也是他的門徒嗎？」彼得不承認，說：「我不是。」 \end{tabularx} \\ \\ \relax
18:26 & \begin{tabularx}{0.7\textwidth}{X} 大祭司的一個僕人，是被彼得削掉耳朵那人的親屬，說：「我不是看見你同他在園子裡嗎？」 \end{tabularx} \\ \\ \relax
18:27 & \begin{tabularx}{0.7\textwidth}{X} 彼得又不承認，立刻雞就叫了。 \end{tabularx} \\ \\
[1ex]
\hline
\hline
\end{longtable}
講題:你為甚麼問我呢? Why question me?

經文:約翰福音18章15-27節

講員:劉國偉先生

場合:香港中文大學崇基學院禮拜堂

日期:2019年7月7日

引言

我很喜歡的北美神學家,教育家和作家唐曼華博士Dr Marva Dawn在她的一本著作”I’m lonely Lord-How Long?” 中譯「我在孤單等待」中,有兩段信息,給我很好的提醒,也與我今天所分享的講道內容,有所共鳴.

第一段信息是她引用詩篇25篇15至18節:「我的眼目時常仰望主,因為祂必將我的腳從網裡拉出來.求你轉向我,憐恤我,因為我是孤獨困苦.我心裡的愁苦甚多,求你救我脫離我的禍患.求你看顧我的困苦,我的艱難,赦免我一切的罪.」

唐曼華特別提到「轉向」這個詞彙,並說這個詞彙深深震撼著她.「轉向」這個希伯來詞語強調的是「關注」或「注視」,考慮到我們所有的錯誤和叛逆,我們有什麼權利期望神會轉過來注視自己?然而,詩人呼叫「求你轉向我,憐恤我」這個請求提醒了我們:祂總是俯就我們,注視及憐恤我們.

第二段信息,是唐曼華引用了詩篇56篇3至4節「我懼怕的時候要倚靠你.我倚靠神,我要讚美祂的話.我倚靠神,必不懼怕.血氣之輩能把我怎麼樣呢?」

「勇敢並不表示你不害怕.」唐曼華寫道:「恐懼和孤單互相增強……大衛承認他也有感到害怕的時候,正正當我們害怕之際,就是倚靠的時候……在倚靠的時候,我們讚美神的話語,因為在神的話語裡面,向我們顯明了祂的信實……神的話語亦關乎神與我們立約的關係.」

唐曼華藉著大衛的詩篇告訴我們,即使勇敢的人也會恐懼,不過出現恐懼的時候,正就是我們學習倚靠神的上好時機.

兩個交錯一起的故事板塊

約翰福音18章15至27節這個段落,包含了兩個不同的故事板塊在交錯進行:耶穌被前任及現任大祭司盤問,當中交叉記載彼得三次不認主的過程,作者採用這個寫作手法而產生出來的文學效果,帶給讀者就兩個敘事人物的觀照,是用文學詮釋來研讀這段經文時,不能遺漏的觀察.

我們嘗試走進經文,細心閱讀.段落開始的時候,作者交代西門彼得跟著耶穌,並且因為另一個門徒的幫忙,得以進去大祭司的院子.這裡所指的大祭司,其實是已經遭羅馬撤去大祭司職位的亞那,亞那在主後六至十四年擔任大祭司,其後被羅馬總督革職,但由於在猶太傳統中,大祭司是終身職事,而且亞那在祭司家族中仍然相當有影響力,現任大祭司該亞法,就是亞那的女婿,所以作者亦隨從百姓,稱亞那為大祭司.

18章16節告訴我們,那位可以出入大祭司院子的門徒,和看門的使女說了一聲,就領了彼得進去；然而,這引起了使女的注意:「你不也是這人的門徒麼?」彼得馬上回應:「我不是.」(約18:17) 提問的使女其實也沒有很大把握,去確認彼得就是耶穌的門徒,於是,彼得也只是簡單的回覆,這個危機看來要過去了.

其實彼得大可以在這時離去,不過他對耶穌的情懷仍在,而且仍然深深關切夫子的安危,他還是留下來.誠然彼得是勇敢的,但「勇敢並不表示你不害怕」,也因為害怕,彼得第一次否認了和耶穌的關係,不過彼得雖然害怕,但他仍是勇敢的,所以他留下.「僕人和差役,因為天冷,就生了炭火,站在那裡烤火；彼得也同他們站著烤火.」(約18:18)

耶穌的勇氣

由19節開始,鏡頭就轉向被盤問的耶穌.亞那其實有充分的動機想要參與這場審問.亞那雖然已經不再是羅馬總督眼中的大祭司,但他卻仍然是大祭司家族的幕後領導人,尤其是有關與聖殿獻祭相關的生意上,在任的大祭司不方便處理,就最好交由亞那管理；在聖殿的外邦人院,有專門售賣祭牲的商販攤檔,以及兌換錢銀的商人,因為這些生意是聖殿範圍近乎壟斷的業務,那些懷著敬虔的心去聖殿獻祭的百姓,就要接受非常不合理的價格,被這個祭司集團及他們指定的供應商嚴重地剝削,而亞那和他的同黨就發大財.所以主耶穌潔淨聖殿的所作所為,其實是衝著這群管理耶路撒冷聖殿的掌權者和既得利益者,而耶穌的行為,當然損害了祭司集團的利益,但更影響深遠的,是耶穌削弱了祭司集團在外邦人院壟斷獻祭生意的正當性.

如此,我們不難理解亞那要參與審問這個攪事的耶穌,也更容易理解,其實整個過程就是要確定耶穌罪名成立,只是最初連罪名都未想好,就匆匆上馬.所以,光明磊落的耶穌回答大祭司的時候說:「我從來是明明的對世人說話……我在暗地裡,並沒有說什麼.」(約18:20) 耶穌正要告訴亞那,其實祂仍然未聽到自己的罪狀是什麼.

耶穌在亞那面前受審,對猶太公會的司法程序來說,其實是一個很大的諷刺.首先,亞那其實並非羅馬授權的,具有司法實權的仍然在位的大祭司；再者,按照猶太公會的審訊原則,應該根據證人的供詞,和各項證據,先確立疑犯的罪狀,不可以在提出控告罪狀之前,就在公會直接盤問被告,企圖從被告的口中,抓著把柄,然後入罪.所以耶穌反問亞那:「你為什麼問我呢?你可以問那聽見的人.」(約18:21) 耶穌的反問,其實要當眾提醒亞那,以及對意圖陷害祂的人提出異議,指出他們其實沒有按照猶太公會正規的司法程序來行事.「你為什麼問我呢?」這個反問也反映了耶穌在權力之前,在看來毫無勝算,理應感到無力的時候,祂對信念及使命那份堅定與堅持.

耶穌勇敢地站在謀害祂的人當中,祂坦承自己所說的一切話,可以經得起查驗.然而,耶穌的勇氣和正氣,在差役的眼中顯然是一種冒犯權威的錯誤行為,所以,他當眾用手掌打了耶穌.一個仍然未有罪狀,當然仍然未定罪的人,卻在權力和制度暴力下,被預期要與權力一方充分合作；如今耶穌不單不合作,更訴諸程序公義,換來的結果,就是耶穌被身邊的差役打祂的臉 struck him in the face,然後說:「你這樣回答大祭司嗎?」 差役濫權,打耶穌的原因是因為耶穌說話的態度,和對權力的不敬.

耶穌回答的時候,仍然堅持猶太公會的司法程序:「我若說的不是,你可以指證那不是；我若說的是,你為甚麼打我呢?」短短幾節經文,出現了三個問題:

你為甚麼問我呢?

你這樣回答大祭司嗎?

3. 我若說的是,你為甚麼打我呢?

耶穌自審問開始直至結束,都堅持提出公平正義的問題,誰知在耶穌兩個問題中間,卻穿插了一個可能不少人會認同,卻其實是反映權力的荒謬的一條問題,大祭司有崇高的身份和地位,公會更加是一處神聖莊嚴的地方,豈可容許耶穌在這個場景如此挑戰亞那.

耶穌沒有否認自己傳道生涯所說的話,更沒有否認這些說話背後所蘊含祂的身份和信仰價值,即使面對非法的審問,及濫權的對待,耶穌仍然牢牢持守信仰的價值.亞那在違反正義的審問過程中,抓不著耶穌的把柄,「就把耶穌解到大祭司該亞法那裡,仍是捆著解去的. 」(約18:24) 約翰刻意提醒讀者,耶穌仍是被捆綁著,被押解去該亞法那裡.一個人連控告的罪狀都未有,卻一直由客西馬尼園被綑綁押解上庭,在那個寒冷的夜晚,份外顯得冰冷.

彼得的恐懼

故事的另一個板塊,是彼得仍然勇敢卻又害怕地,站在一個看到耶穌的位置烤火,他對拉比的關心是無容置疑的.然而,在艱難的時刻,信仰的考驗往往隨時又反覆出現!此時,又有一個人問彼得:「你不也是他的門徒嗎?」彼得繼續不承認,說:「我不是.」(約18:6)

這是彼得第二次在查問的人面前跟耶穌割席,彼得先後兩次被人查問是否耶穌的門徒,而耶穌門徒這個身份帶來的潛在風險,因為對手是大祭司和整個猶太公會,其實是可以大得難以估量的.面對死亡,能夠有豁出去的勇氣,其實很不容易,而彼得求生存的本能反應,叫他霎時間否認耶穌,也背棄了自己.

耶穌回應大祭司的問題:「你為甚麼問我呢?」這刻,彷彿成為襯托彼得內心掙扎的一個窗口,彼得心裡會否暗地裡埋怨這些好管諸事之人「為甚麼問我呢?」然而,我相信彼得連番招架,根本未及反思,在情緒混亂之際,第三問來得非常堅實.

「有大祭司的一個僕人,是彼得削掉耳朵那人的親屬,說:「我不是看見你同他在園子裡嗎?」(約18:26) 圍著火盤的人只有彼得是耶穌的門徒,彼得在這個處境底下,其實非常孤單.唐曼華說恐懼和孤單互相增強,當時彼得的恐懼應該不斷攀升,於是彼得就第三次不承認耶穌.立時,雞就叫了.

耶穌轉眼在我身上

親愛的弟兄姊妹,在我們每天每夜的現實生活中,孤單和恐懼時常突襲我們,在艱難的時刻,更可能如影隨形.因為恐懼,我們有時無奈逼使自己去違背信念和原則,為了生存,即使勇敢和熱誠如彼得,亦不幸接二連三次否認主,當我們遇上有若彼得的處境時,我們會怎樣反應?

耶穌的勇氣對照彼得的恐懼,耶穌的堅定相比彼得的否認,我相信作者要我們思想,耶穌其實比孤單的彼得更加孤單.不過,耶穌必然確信天父注視祂,在這經歷受苦的時間,亦成為耶穌倚靠天父的緊密時刻.

當我在對觀福音查考彼得三次不認主的敘事時,我找到解開兩個故事板塊交錯一起的鑰匙,路加醫生記載了雞叫之後,一個很重要的片段:「主轉過身來,看彼得.彼得便想起主對他所說的話……他就出去痛哭.」(路22:61-62) 唐曼華提到「轉向」這個詞彙,深深震撼著她,她問讀者 : 我們有什麼權利期望神會轉過來注視自己?弟兄姊妹,就在主耶穌經歷受苦之際,祂轉眼注視彼得,因為祂憐恤正在孤單,恐懼和無力感當中的彼得,亦是這位復活之主,在提比哩亞湖邊再一次確立對彼得的呼召,堅立彼得去勇敢面對未來,迎向作主門徒的更大挑戰.

弟兄姊妹,我再次讀出兩段詩篇.

詩篇25篇15至18節:「我的眼目時常仰望主,因為祂必將我的腳從網裡拉出來.求你轉向我,憐恤我,因為我是孤獨困苦.我心裡的愁苦甚多,求你救我脫離我的禍患.求你看顧我的困苦,我的艱難,赦免我一切的罪.」

詩篇56篇3至4節「我懼怕的時候要倚靠你.我倚靠神,我要讚美祂的話.我倚靠神,必不懼怕.血氣之輩能把我怎麼樣呢?」

主與我們同在,也與香港同在.



\newpage

\section{智者出口聖者出走仁者出手}
\label{sec:461}
\textbf{林豪恩先生 - 智者出口,聖者出走,仁者出手}
\newline
\newline
\hyperref[sec:460]{< < < PREV SERMON < < <}
~
\hyperref[sec:toc]{[返目錄]}
~
\hyperref[sec:462]{> > > NEXT SERMON > > >}


路加福音 10:25-37
\newline
\begin{longtable}{cl}
\hline
\hline
章節 & 經文 (和合本修訂版)\\
\hline
10:25 & \begin{tabularx}{0.7\textwidth}{X} 有一個律法師起來試探耶穌，說：「老師！我該做甚麼才可以承受永生？」 \end{tabularx} \\ \\ \relax
10:26 & \begin{tabularx}{0.7\textwidth}{X} 耶穌對他說：「律法上寫的是甚麼？你是怎樣念的呢？」 \end{tabularx} \\ \\ \relax
10:27 & \begin{tabularx}{0.7\textwidth}{X} 他回答說：「你要盡心、盡性、盡力、盡意愛主—你的神，又要愛鄰如己。」 \end{tabularx} \\ \\ \relax
10:28 & \begin{tabularx}{0.7\textwidth}{X} 耶穌對他說：「你回答得正確，你這樣做就會得永生。」 \end{tabularx} \\ \\ \relax
10:29 & \begin{tabularx}{0.7\textwidth}{X} 那人要證明自己有理，就對耶穌說：「誰是我的鄰舍呢？」 \end{tabularx} \\ \\ \relax
10:30 & \begin{tabularx}{0.7\textwidth}{X} 耶穌回答：「有一個人從耶路撒冷下耶利哥去，落在強盜手中。他們剝去他的衣裳，把他打個半死，丟下他走了。 \end{tabularx} \\ \\ \relax
10:31 & \begin{tabularx}{0.7\textwidth}{X} 偶然有一個祭司從那條路下來，看見他就從另一邊過去了。 \end{tabularx} \\ \\ \relax
10:32 & \begin{tabularx}{0.7\textwidth}{X} 又有一個利未人來到那裡，看見他，也照樣從另一邊過去了。 \end{tabularx} \\ \\ \relax
10:33 & \begin{tabularx}{0.7\textwidth}{X} 可是，有一個撒瑪利亞人路過那裡，看見他就動了慈心， \end{tabularx} \\ \\ \relax
10:34 & \begin{tabularx}{0.7\textwidth}{X} 上前用油和酒倒在他的傷處，包裹好了，扶他騎上自己的牲口，帶他到旅店裡去，照應他。 \end{tabularx} \\ \\ \relax
10:35 & \begin{tabularx}{0.7\textwidth}{X} 第二天，他拿出兩個銀幣來，交給店主，說：『請你照應他，額外的費用，我回來時會還你。』 \end{tabularx} \\ \\ \relax
10:36 & \begin{tabularx}{0.7\textwidth}{X} 你想，這三個人哪一個是落在強盜手中那人的鄰舍呢？」 \end{tabularx} \\ \\ \relax
10:37 & \begin{tabularx}{0.7\textwidth}{X} 他說：「是憐憫他的。」耶穌對他說：「你去，照樣做吧！」 \end{tabularx} \\ \\
[1ex]
\hline
\hline
\end{longtable}
講題:智者出口,聖者出走,仁者出手 The Wise Talk, the Sacred Pass, the Good Do

經文:路加福音10:25-37

講員:林豪恩先生

場合:香港中文大學崇基學院禮拜堂

日期:2019年7月14日(聖靈降臨後第五主日)

引言

今天的福音經課講述的「好撒瑪利亞人」故事,大家都耳熟能詳.當大家讀這個故事的時候,有沒有感到這故事的諷刺和張力呢?這個故事很可能一次過刺激了三種人,其中兩種是「明批」,另外一種是「暗串」.何以見得呢?因為這三種人在故事中都被安排為負面角色.第一種可能感到不愉快的是律法師,故事中的角色是「起來試探耶穌」；第二種可能感到不是味兒的是祭司和利未人,他們在故事中的形象是「見死不救」.第三種呢?故事好像沒有提及第三種負面角色,不過,卻把正面的角色安排給撒瑪利亞人,而不是猶太人.讀者可以從這段經文的上文,路加福音9章51節看見當時的猶太人與撒瑪利亞人之間的群族關係:

9:51耶穌被接上升的日子將到,他就定意向耶路撒冷去,

9:52便打發使者在他前頭走.他們到了撒馬利亞的一個村莊,要為他預備.

9:53那裡的人不接待他,因他面向耶路撒冷去.

9:54他的門徒雅各,約翰看見了,就說:「主啊,你要我們吩咐火從天上降下來燒滅他們,像以利亞所做的嗎?」

在這種不友好的群族關係之下,故事把正面的角色安排給撒瑪利亞人而不是猶太人,可能會令猶太人不滿.如此講述這個「救急扶危」的故事,如果用今天的公關角度來看,真是公關災難.律法師和祭司群體都是當時的宗教領袖,而猶太人正是耶穌和門徒的自己人,開罪他們,真是自討苦吃.如此這般,耶穌仍要如此講述這個故事,門徒仍要記錄這個故事,對我們有甚麼啟發呢?

二．經文

10:25有一個律法師起來試探耶穌,說:「老師!我該做甚麼才可以承受永生?」

10:26耶穌對他說:「律法上寫的是甚麼?你是怎樣念的呢?」

10:27他回答說:「你要盡心,盡性,盡力,盡意愛主-你的上帝,又要愛鄰如己.」

10:28耶穌對他說:「你回答得正確,你這樣做就會得永生.」

10:29那人要證明自己有理,就對耶穌說:「誰是我的鄰舍呢?」

10:30耶穌回答:「有一個人從耶路撒冷下耶利哥去,落在強盜手中.他們剝去他的衣裳,把他打個半死,丟下他走了.

10:31偶然有一個祭司從那條路下來,看見他就從另一邊過去了.

10:32又有一個利未人來到那裡,看見他,也照樣從另一邊過去了.

10:33可是,有一個撒瑪利亞人路過那裡,看見他就動了慈心,

10:34上前用油和酒倒在他的傷處,包裹好了,扶他騎上自己的牲口,帶他到旅店裡去,照應他.

10:35第二天,他拿出兩個銀幣來,交給店主,說:『請你照應他,額外的費用,我回來時會還你.』

10:36你想,這三個人哪一個是落在強盜手中那人的鄰舍呢?」

10:37他說:「是憐憫他的.」耶穌對他說:「你去,照樣做吧!」

三．釋經與應用

1．智者出口

10:25有一個律法師起來試探耶穌,說:「老師!我該做甚麼才可以承受永生?」

10:26耶穌對他說:「律法上寫的是甚麼?你是怎樣念的呢?」

10:27他回答說:「你要盡心,盡性,盡力,盡意愛主-你的上帝,又要愛鄰如己.」

10:28耶穌對他說:「你回答得正確,你這樣做就會得永生.」

10:29那人要證明自己有理,就對耶穌說:「誰是我的鄰舍呢?」

在當時的宗教社會,律法師一方面是解釋經文的專家,同時也是解釋法律的權威,因此,若要問故事中誰是智者?非他們莫屬了.律法師的問題:「老師!我該做甚麼才可以承受永生?」是十分正當的,耶穌沒有理由不回答.不過,這又是一個不懷好意的問題,因為作者指出問這個問題的動機並不是為了獲得答案,而是「試探耶穌」.無論是真情或是假意,耶穌真誠地回答:

10:26耶穌對他說:「律法上寫的是甚麼?你是怎樣念的呢?」

10:27他回答說:「你要盡心,盡性,盡力,盡意愛主-你的上帝,又要愛鄰如己.」

10:28耶穌對他說:「你回答得正確,你這樣做就會得永生.」

可能因為耶穌這樣的示範,使徒彼得寫彼得前書時也如此教導後輩:

3:14不要怕人的威嚇,也不要驚慌；

3:15只要心裡尊主基督為聖.有人問你們心中盼望的緣由,就要常作準備,以溫柔,敬畏的心回答各人.

這位律法師問:「老師!我該做甚麼才可以承受永生?」,其實他是沒有打算按著耶穌的建議去做的,他來,只是為了找耶穌打「口水戰」.撇開「試探」這種如此兇險的狀況,在生活上,你有沒有見過別人來問問題,其實並不是想聽答案,只是想講自己的意見?有沒有試過有人來尋找建議,其實並不打算聽建議,只是尋找支持的聲音?有沒有試過有人來請教工作如何做,其實並不打算做,只是打算說服你不用他做?

媽:阿仔,今次考試成績咁差,你要想辦法改善啦!

仔:阿媽,今年Ｄ老師教得咁差,有乜辦法成績唔差呀!

媽:如果老師真係教得唔好,你更加要用多Ｄ時間溫書啦!

仔:阿媽,隔離屋個明仔夠用好多時間溫書啦,又唔見佢成績好好?

媽:阿仔,阿媽讀得書少,唔識教你,不如俾錢你去補習啦!

仔:阿媽,唔好sai 錢啦,你睇下,幾多人去補習,又唔見佢地個個成績都咁好!(反駁所有建議,即是甚麼也不用做!)

古語有云:

春天不是讀書天,夏日炎炎正好眠.過得秋來冬又到,收拾書本好過年.

(結論就是:任何時候都不適合讀書!)

政府制定政策前,喜歡請學者做研究,提供數據,分析及建議.不知道在座的老師們有沒有接過類似的委託呢?幾年前,政府委託一位長期關注退休保障制度的學者,帶領其團隊研究退休保障方案,隨後向政府力薦「全民退休保障」,結果卻遭受冷待,更被當時的政務司司長批評方案「不切實際」,指斥其研究團隊沒有「認真學術研究應有的態度」.該教授回應「作為研究者感到憤怒」,更批評政府向公眾推出的諮詢文件「不盡不實」,「只當我們是傻瓜」.教授經此一役,感到政府只說不做,向公眾表示:「今屆政府,百分之一百不會實施退休保障改革,大家不用寄以厚望.」

從耶穌與律法師的答問來看,律法師是知道答案的,而且耶穌也肯定他所知的:

10:28耶穌對他說:「你回答得正確,你這樣做就會得永生.」

10:29那人要證明自己有理,就對耶穌說:「誰是我的鄰舍呢?」

律法師當然要證明自己有理,如果反駁成功,就代表不用按著耶穌的建議去做,畢竟,律法師是沒有打算切實地去做的.用說話去逃避行動,或者用說話去推卸責任,在現實上也屢見不鮮,或者,不少人也有如成為律法師一樣的智者潛質.智者可能知識比人豐富,所知道的也正確,而且能言善辯,不過,在這個故事中卻是負面的角色,並不是耶穌建議讀者學習的榜樣.

2．聖者出走

10:29那人要證明自己有理,就對耶穌說:「誰是我的鄰舍呢?」

10:30耶穌回答:「有一個人從耶路撒冷下耶利哥去,落在強盜手中.他們剝去他的衣裳,把他打個半死,丟下他走了.

10:31偶然有一個祭司從那條路下來,看見他就從另一邊過去了.

10:32又有一個利未人來到那裡,看見他,也照樣從另一邊過去了.

如果律法師是智者,聖者一定是祭司和利未人.祭司和利未人「看見他就從另一邊過去了」,常見的有三種解釋.第一種是說祭司及利未人可能擔心這位受害者不是律法所界定的「鄰舍」,而現場狀況又不容許他們確認對方是否受割禮的人,因此還是留給其他人協助為妙.第二種解釋:律法禁止祭司觸摸屍體,如果摸了,就是不潔淨.祭司和利未人或者是為了潔淨的緣故而對路旁的不幸者避之則吉.第三種解釋認為,祭司和利未人擔心劫匪還埋伏在附近.若停下來幫助受害者,自己也可能成為劫匪的目標.無論我們用甚麼理由來為祭司和利未人「聖者出走」的表現解釋,故事對他們的評價卻是呼之欲出的.從另一個角度看,無論祭司和利未人如何合理化自己的表現,維持聖者的身分而自我感覺良好,其他人對他們的觀感難免是「見死不救」.人們對於祭司和利未人這種「離地的神聖」是有感受的,因此,當耶穌走遍各城各鄉傳道,切實關心及回應人們身心靈的需要時,人們擁戴耶穌而令祭司和利未人妒忌.今時今日,人們的選擇比二千年前更多,也更自由,人們更會用腳去投票.

在西方世界早已有信徒離開教會的現在.David Kinnaman於《論盡基督徒》(UnChristian,中文版二○一○年)剖析教會文化六項因素導致十八歲至三十歲青年信徒離家出走:教會過度保護,信仰經驗膚淺,對科學欠缺理性,對同性戀過度論斷,信仰過度排斥性,對疑惑者極不友善.這些不同因素交織,使有理想的青年信徒感到難以在堂會內健康成長.(胡志偉:再思青年信徒出走教會現象)

在亞洲,基督教人口比例最多的韓國,教會正經歷一個大問題,就是整個世代的信徒出走. 彭書睿分析,信二代,信三代對韓國教會內部領導層習染官僚氣感到厭倦,加上韓國社會競爭激烈,工作,教育,世俗主義都帶來教會牧養很大的衝擊.(基督日報:韓國宣教盛勢下降 牧者:信二代出走接棒無援)

香港呢?在一份報章(經濟日報)讀到一篇名為「10個關於信徒離開教會的原因」的文章,其中幾項如下:

當你唔熟聖經,唔理解神學,佢叫你多花時間鑽研.

當你花時間鑽研,佢同你講知識叫人自高自大,要有愛心,惟有愛心能造就人.

當你研究聖經原文,佢同話你浪費時間,原文唔重要,要了解當中精義,不能單憑字面.

佢一邊怪責信徒唔關心教會,當有信徒關心教會發展,問點解要成十億擴堂時,佢地又怪責人太多事太多聲音.

當教會領袖犯錯,佢地話唔好論斷人,唔好掉石頭,唔好因為犯錯而抹殺佢ge 教導,普通信徒犯錯,有一系列ge 教會紀律!

佢地話教會唔應該講政治,唔准提雨傘運動種種,但佢地又樂意請基督徒高官及議員講見證.

聖者為了維持自己的神聖,一定找到理由「看見他就從另一邊過去了」.不過,觀眾的眼睛是雪亮的,他們會看見耶穌並不讚賞這樣的「神聖」；人們是有獨立思考的,自然會用腳去做選擇.

3．仁者出手

10:33可是,有一個撒瑪利亞人路過那裡,看見他就動了慈心,

10:34上前用油和酒倒在他的傷處,包裹好了,扶他騎上自己的牲口,帶他到旅店裡去,照應他.

10:35第二天,他拿出兩個銀幣來,交給店主,說:『請你照應他,額外的費用,我回來時會還你.』

這個故事被稱為「好撒瑪利亞人」,這位撒瑪利亞人「好」在哪裡呢?答案也是呼之欲出,他的「好」,是因為他路見受傷的人而採取行動救助.耶穌講這故事時對這位救助者的行動描繪得非常仔細:上前,用油和酒倒在他的傷處,包裹,扶他騎上自己的牲口,帶他到旅店裡去,照應他,拿出兩個銀幣交給店主,吩咐『請你照應他』,保證『額外的費用,我回來時會還你.』救助步驟如此具體,應該是鼓勵讀者仿傚.「好撒瑪利亞人」的「好」,是因為他「出手救助」,然而,「出手救助」這好行為,是出於一棵「好」心.耶穌是如此描述他的:「有一個撒瑪利亞人路過那裡,看見他就動了慈心」.明代詩人曹學佺的著名對聯說:「仗義每多屠狗輩,負心多是讀書人」.「好撒瑪利亞人,可以不是如智者的律法師,可以不是如聖者的神職人員,甚至也可以不是自覺正統的猶太人.

此時此地,誰是「落在強盜手中／被剝去衣裳／被打個半死／被丟下」的受害者等著救助呢?相信大家想到的各有不同.我想到有一種工人,他們返工隨時「搵條命去搏」,而且「搏」完之後也可能沒有甚麼保障.

你有沒有使用過港珠澳大橋呢?你知道這宏偉的建築是如何成就的嗎?有工人組織指出,香港段工程在通車前已有11名工友喪生,內地段工程亦有9人死去,合共20名工人因工程而喪命,受傷人數粗略估計超過600人.

工傷之後如何呢?有一位處理工傷的大律師寫下她的所見所感:(《看不見的傷痛》:二度傷害)

「工傷工友追討賠償的過程,往往是更深的二次傷害.整個制度的設計,都在假定傷者的主觀感受是不可靠的……總是對傷者的不信任……你說你行動不便?代表顧主的保險公司派私家偵探日夜跟蹤,說拍攝到你看似健步如飛,談笑風生的樣子……作為律師,見過太多的案件,本來可能只是小傷,但往往因為訴訟過程中的各種壓力和困擾,發展出比原本的「身體傷害」嚴重得多的精神問題及家庭問題.就算最後打贏官司,精神傷害,經濟壓力. 個人尊嚴,以年計的精神折磨. 與親人關係的破碎,又豈是金錢能補償得了的呢?甚麼時候,我們的制度才能給工傷工友應有的尊重和尊嚴,不會再對他們一沉百踩呢?」

今年五月,有一位名叫陳錦康的朋友病逝.他不是政要,也不是明星,社會大眾多不知道他是誰,但知道他做過甚麼事情的人都為他的離去而惋惜.他24歲加入工業傷亡權益會,直至病逝時,服務了36年.每逢有工業意外發生,他總會出現在意外現場,醫院甚至殮房,幫助及安撫徬徨無助的家屬,協助他們處理後事,援助基金等事宜,甚至事後與傷亡者的公司力爭賠償等.記者問他是什麼驅使他一直做落去?他說了個故事,「很多年前有一單沉箱意外,他到醫院找不到家屬,知道家屬住在粉嶺的臨時房屋,他漏夜去拜訪,門打開,屋內只有一盞微弱的燈,媽媽抱住兩個年幼女兒,正在抱頭痛哭,傾完一輪,知道呢家人全副身家原來只剩得八蚊.」

四．總結

10:36(耶穌問)你想,這三個人哪一個是落在強盜手中那人的鄰舍呢?」

如果你就是那位「落在強盜手中那人」,你希望期遇到誰呢?律法師?祭司和利末人?還是那位撒瑪利亞人呢?如果我們不幸成為受害者,我們期望受到甚麼對待呢?在故事中的律法師是知道的

10:37他說:「是憐憫他的.」

接著耶穌對他說:「你去,照樣做吧!」

從古至今,社會都高抬智者,人人都景仰聖者,耶穌所說的故事提示人們,鄰舍需要的是仁者.其實,智者和聖者也可以是仁者,如何變成仁者呢?不需要學位,也毋須修練,看見受苦者,好像那位撒瑪利亞人一樣「動了慈心」就是了.



\newpage

\section{不可少的只有一件}
\label{sec:462}
\textbf{許立中先生 - 不可少的只有一件}
\newline
\newline
\hyperref[sec:461]{< < < PREV SERMON < < <}
~
\hyperref[sec:toc]{[返目錄]}
~
\hyperref[sec:463]{> > > NEXT SERMON > > >}


路加福音 10:38-42
\newline
\begin{longtable}{cl}
\hline
\hline
章節 & 經文 (和合本修訂版)\\
\hline
10:38 & \begin{tabularx}{0.7\textwidth}{X} 他們繼續前行，耶穌進了一個村莊。有一個女人，名叫馬大，接他到自己家裡。 \end{tabularx} \\ \\ \relax
10:39 & \begin{tabularx}{0.7\textwidth}{X} 她有一個妹妹，名叫馬利亞，在主的腳前坐著聽他的道。 \end{tabularx} \\ \\ \relax
10:40 & \begin{tabularx}{0.7\textwidth}{X} 馬大伺候的事多，心裡忙亂，進前來，說：「主啊，我的妹妹留下我一個人伺候，你不在意嗎？請吩咐她來幫助我。」 \end{tabularx} \\ \\ \relax
10:41 & \begin{tabularx}{0.7\textwidth}{X} 主回答說：「馬大，馬大，你為許多的事操心煩惱， \end{tabularx} \\ \\ \relax
10:42 & \begin{tabularx}{0.7\textwidth}{X} 但是不可少的只有一件。馬利亞已經選擇了那上好的福分，是沒有人能從她奪去的。」 \end{tabularx} \\ \\
[1ex]
\hline
\hline
\end{longtable}
講題:不可少的只有一件 The One and The Many

講員:許立中先生

經文:路加福音10:38-42

場合:香港中文大學崇基學院禮拜堂

日期:2019年7月21日

「他們走路的時候,耶穌進了一個村莊,有一個女人名叫馬大,接他到自己家裡.他有一個妹子名叫馬利亞,在耶穌腳前坐著聽他的道.馬大伺候的事多,心裡忙亂,就進前來說,主阿,我的妹子留下我一個人伺候,你不在意麼?請吩咐他來幫助我.耶穌回答說,馬大,馬大,你為許多的事,思慮煩擾；但是不可少的只有一件；馬利亞已經選擇那上好的福分,是不能奪去的.」

(路10:38-42)

在今天的福音經課裡,我們讀到一個十分熟悉的記載.這件事並沒有出現在其他的福音書.經文記載,耶穌和門徒到處宣揚天國的道理.有一次,他們經過一個村莊,作者並沒有清楚說明在哪裡,但從其他的記載,我們知道馬大和馬利亞是住在伯大尼,離開耶路撒冷大概只有兩英里.經文給我們的印象,馬大和馬利亞是兩姊妹,馬大是家姐,是耶穌在當地的一個追隨者.她邀請耶穌到她的家裡作客.馬大的妹妹馬利亞把握這個難得的機會,坐在耶穌的腳前,聽她講道.反而馬大衛了歡迎這位稀客,因此辦得比平時鋪張,一個人在廚房忙得團團轉.她對馬利亞只是坐在那裡陪著耶穌,感到非常不滿.

最後馬大終於忍不住,向耶穌投訴:「主阿,我的妹子留下我一個人伺候,你不在意麼?請吩咐她來幫助我.」語氣中明顯有怪責的意思.不過耶穌的回答反倒是相當溫柔的.他說:「馬大,馬大」,有點像父親對兒子說:「仔阿,仔」那樣:「馬大,馬大,你為許多的事,思慮煩擾；但是不可少的只有一件；馬利亞已經選擇那上好的福分,是不能奪去的.」

有另外一些文本說:「馬大,馬大,你為許多的事,思慮煩擾；但真正需要的其實不多,或許只有一件」.但無論我們選擇哪一個文本,耶穌似乎是說,馬大可能是捉錯用神,為了太多事務性的東西忙得團團轉,卻沒有把握到事情的重點,有點徒勞無功的意思.

但是至於「那不可少的一件」到底是甚麼,耶穌當時卻沒有說明白.他只是說「馬利亞已經選擇那上好的福分,是不能奪去的」.而很多時我們想當然的聯想,「那不可少的一件」,一定是關乎一個人的屬靈狀況,或者是他的宗教責任.很多人很快就跳到經文的應用:事奉上帝,最重要的並不是工作,而是態度!

我們很肯定地指出,你是一個怎樣的人,比你做些甚麼更加重要.我們甚至回到《創世記》第四章,指出耶和華首先是看中亞伯,然後悅納他的供物；首先是看不中該隱,然後才是拒絕他的供物.我們說:「你是一個怎樣的人,比你做甚麼更重要.」

我當然不能說這個想法不對.事實上這類說話,你甚麼時候說都是對的.問題是,這是不是這段經文,這個故事的核心重點?

事實上基督徒對這個故事是熟悉到沒有甚麼可以再說的了.真是要找些新的點子的話,於是有人開始想到要為馬大「平反」:他們說其實馬大之所以搞咁多ye ,都只不過是出於他對耶穌的尊崇.比作是你,如果有一個大人物要到你家裡吃飯,難道你就不會斬多幾味,或者親自下廚為這位貴賓預備幾個特別的菜色?難道你就會落街買幾個飯盒招呼貴賓嗎?肯定是不會的.

還有,馬利亞之所以能夠坐在耶穌腳前聽道,豈不是全賴有馬大在廚房忙碌?如果兩姊妹都安靜地坐在耶穌腳前聽道,那麼晚餐又由誰來預備?事實上沒有馬大的忙碌,就不可能有馬利亞的安靜聽道.

所以有時我們看見一些堂會一個週末開成五,六堂聚會,都只能夠替背後的工作人員感到辛苦,而不敢批評他們太過忙碌,失去安靜敬拜的意義.事實上我對上一次講道,就是要在兩日之內,在同一間堂會將同一篇講章重複三次.至少我是並不喜歡這樣的經驗.並且「量變質變」,以前你是說參加第幾「堂」聚會,現在的信徒是說參加哪一「場」聚會,就好像你去看電影,聽演唱會那樣.說遠了一點.

無可否認,這些都是一些很有趣並且很現實的角度,亦似乎在某些方面呼應著這段經文.但我們仍然要問,那其實馬大的問題,到底在那裡?

在這裡,我想打岔提一提另一段似乎不大相關的記載,希望最後可以帶出它們的關連性.(太19；可10；路18)有一次一個人無端走來攔著耶穌,問他一個看來十分簡單的問題:「夫子,我當作甚麼事,才可以承受永生?」

從《馬太福音》的記載,我們知道他是一個「少年人」,而從《路加福音》,我們知道他是一個官.加起來,我們就知道他是一個少年的官,當然在不同的文化風俗,對「少年」會有不同的界定.而做得官的,恐怕亦不會太年輕.

無論如何,他問了耶穌一個簡單的問題,看來亦期望得到一個簡單的答案.於是耶穌也真的給他一個簡單答案:「誡命你是曉得的:不可姦淫,不可殺人,不可偷盜,不可作假見證,當孝敬父母,如此類推.」那個人似乎又覺得這答案太過簡單,他說:「喂,這一切我從小都遵守了.」事實上我想他是有點奇怪,耶穌竟然會提出這幾乎是主日學程度的答案.

我想說的是:一個自覺「從小就遵守」一切宗教義務和規條的人,竟然還會問,他還需要作甚麼事,才可以承受永生.換句話說,他基本上是盡了一切宗教上的要求,然後他還覺得若有所失.耶穌聽了,就說,「你還缺少一件」.而我覺得,耶穌在第十八章說,那位少年官所缺少的一件,跟他在這裡對馬大所說的「那不可少的只有一件」,似乎有相當重要的關連.

在那少年官來說,那缺少的一件,就是「要變賣你一切所有的,分給窮人,你還要來跟從我」.非常明確,毫不含糊.但對馬大,耶穌卻沒有具體說明那不可少的一件是甚麼.他只是說,「馬利亞已經選擇那上好的福分,是不能奪去的.」

從另一個角度來看,馬大所做的,其實也不算不好,只是並非「上好」；就好像那位少年的官,他從小就遵守所有的誡命並不是不好,至少他是循規蹈矩,行得正,企得正.這裡的重點是,他問這個問題,是反映連他自己也知道,他根本還未曾觸及問題真正的核心.

當然,生活在二十一世紀的香港,我們可能不會好像那少年官那麼急於知道「我當作甚麼才可以承受永生?」事實上不少人甚至不齒於以宗教作為手段,去逃避面對生命種種難處和不確定性.對於知識分子來說－或者也包括不少的基督徒－「未知生,焉知死」仍然是他們的生活態度.並不是因為他們不相信永生,而是因為他們相信此生仍然是一個很獨特的經驗,仍然有很多責任他們需要去承擔.

因此拉闊來看,很多人反而會問:在一個存在著無限可能的世界,一個人可以怎樣找到一條比較明確,屬於自己的道路?我們可能會記得Robert Frost的名句“The Road Not Taken”:面對兩條分岔路,到底應該選擇哪條?Robert Frost選那條少人行的.那是一個拒絕隨流的選擇,亦是一條有探索空間的道路.另一方面,在政治上,在高牆與雞蛋之間,村上春樹說他永遠會選擇雞蛋.一個相對地簡單的道德抉擇.但在這個多元化的社會,有時人多未必就錯,人少亦未必正確.

譬如社會上大多數人都是守法的,他們等巴士會排隊,而守法或者排隊並不會因為多數人都這樣做而成為錯誤；相反來說,政治領袖或者恐怖份子永遠只是少數,他們亦不會因為人少而變得正確.所以如果純粹以參與的人數作為判斷的標準,有時亦未必可靠.

那我們到底可以怎樣在這個豐富多元(diversified)的世界找到統一性(unity),以致我們的生命可以有一個比較確定的座標(frame of reference),去理解周圍發生的事情,從而作出至少「過到人,過到自己」的抉擇?

我們的問題是,萬事萬物在這個繽紛多元的世界,彼此之間到底有沒有任何關連?如果沒有,生命就只是一連串毫無關係的偶發事件,而我們亦不可能從這些偶發事件梳理出任何頭緒.

你這樣做,或者你不這樣做；你走這邊,或你走那邊；你維護建制,或者你反對建制,在本質上其實並無分別.一切所謂的「意義」,都只不過是你一廂情願的想法；而你亦休想以這套想法去說服其他人,因為其他人也有他們自己的想法.

但如果事物之間的確有關連,那又會是一個怎樣的關連?而這種關連或統一性,又會否抹煞生命豐富的多元性?我的意思是,如果生命不是偶然,這繽紛多元的世界,其實都是指向一個目標,一個實體,那麼在這個overriding的統一性,規律,原則底下,又如何能不至抹煞事物本身的個別性?這樣說好像有點抽象.

譬如基督徒相信這個世界是上帝創造,人最大的福祉,就是按著上帝的心意生活.信徒相信,這是構成生活意義的統一性.但是在實踐上,個別的人或群體,在這個多元化的世界,又如何體現這統攝性的大原則呢?

在這二十多年來,基督徒慢慢發覺自己無法抽身於這個光怪陸離的社會,對於風起雲湧的政治形勢,亦難以置身事外.他們一直為了教會在世界,在社會的角色爭辯不休.最後大家似乎終於找到一個可以接納的共識,就是《彌迦書》六章八節:「世人啊,耶和華已指示你何為善；他向你們所要的是甚麼呢?只要你行公義,好憐憫,謙卑地與你們的上帝同行.」

原則是有了.但當你發現不同的宗派,堂會,在實踐這原則時自由演繹,各自發揮,甚至到一個彼此排斥,互相否定的地步,就表明這個統一性其實已經失去駕馭分歧的能力.當一個所謂原則發展到南轅北轍,面目全非,就好像當年八三一人大釋法那樣,就表示所謂原則與實踐,統一與多元之間,已經失去了本質上的關係.他們所聲稱的,跟他們實際所做的,可以說是毫不相干.

這裡似乎牽涉一個哲學問題,就是「一」與「多」(The One and the Many)的問題:那位少年官自小就遵守一切誡命,滿足了很多宗教上的要求；但耶穌說他「還缺少一件」,並且可能是最關鍵的一件.馬大忙忙碌碌,施展渾身解數,要預備最好的去招呼耶穌,可是耶穌卻對她說,「你弄那麼多幹嘛?那不可少的,其實只有一件」.

「一」是核心性的,往往也是抽象的,「多」是具體的呈現.耶穌所質疑的,是他們所呈現的,還是不是那原本的初衷?「一」與「多」,The One and the Many,在哲學上並不是對立的.相反,它們之間其實是密切地互相依賴的.「一」的統攝性,是體現於具體事物的多元性.譬如「狗」這個物種,就統攝了唐狗,番狗,日本狗,狼狗,臘腸,獅子狗這一切被稱為「狗」的生物.

可以說,如果沒有這些不同品種,體型,國籍,毛色,大小的狗,我們就不知道狗是一種怎樣的生物.我們是透過一隻隻,一類類不同的狗,去理解那被統稱為「狗」的物種.這是「一」和「多」的關係.你大概不知道,非洲的鬣狗(hyaena),在動物分類上是更接近貓科.

你當然會問,這個關於「一」和「多」的哲學問題,跟我們今天思考的經文有甚麼關係?確實是大有關係.

在宗教上,一個人可以很熟悉一些宗教儀式的細節和流程,譬如殯儀館的堂倌；但他們對道教或基督教背後的哲學或信仰,可能完全沒有興趣.我們亦見過不少在中國傳統宗教或基督教家庭長大的孩子,他們對於宗教的習俗和實踐可能瞭如指掌,但對信仰本身卻不甚了了.可以說,他們好像了解很多,但實際掌握多少就很難說了.

那位少年的官,自小就遵守那些世代流傳下來的誡命,他以為自己已經清楚每一條誡命的內容和要求.但有學者指出,他其實根本不明白誡命,因此才會以為自己從小就遵行這些誡命.

那些真正明白誡命的,譬如使徒保羅,就體會到「立志行善由得我,行出來卻由不得我」；而誡命的精義,就是耶穌所說的:「你們聽過吩咐古人的話:不可殺人,但我告訴你們,凡無緣無故向弟兄發怒的,都要受審判；你們聽過這樣的話:不可姦淫,但我告訴你們,凡看見婦女就動淫念的,他在心裡已經犯了姦淫.」這些就是誡命的精神,律法的要義.那你真的以為有人可以從小就遵守這一切的誡命嗎?

那位少年官的問題,正是他了解到每條誡命的具體細節,卻不明白這些誡命到底指向一個怎樣的目的和實體.就好像人們看見hyaena像條狗,就以為它是狗.

他以為自己從小就遵守這一切的誡命,卻仍然對誡命背後的總綱一無所知,以致他去求教耶穌:「夫子,我該作甚麼才可以承受永生?」他努力遵守眾多的誡命或許不能算錯,只不過在他的「多」的背後,並沒有「一」的支撐,因此到最後只是徒勞無功.

而馬大的問題,並不在於她想盡辦法去討好耶穌,表達對耶穌的著緊和敬意,乃在於她被自己想出來的辦法和方案,掩蓋了這些辦法和方案原本要達成的目的,就是讓這位遠道而來的稀客,感到賓至如歸.

情況就有點像一個人邀請朋友回家吃飯.朋友一入門口,就請他入客廳坐下,欣賞那部72寸5K大電視,然後入廚房預備精彩的菜色.晚飯間,他向朋友介紹每道特別的菜色,講解那支法國紅酒的年份和背景,順便解釋產地的土壤,水質和氣候,怎樣影響那支酒的味道和結構.講解完了,也食得七七八八,於是他再拿出那隻珍藏的藍光影碟,招呼朋友在客廳欣賞,自己才走到廚房幫忙太太善後.一切打點完,也就差不多是送客的時候.結果忙一大輪,他甚至忘記問候那位朋友剛入院的太太.

我們都珍惜生命,明白時間的寶貴.我們愛惜光陰,活在當下,努力過一個充實的,不枉此生的生活.但有時當我們驀然回首－往往是當我們面對生命階段的一些轉折位,或者當身體發出一些警號,或者我們一向視為理所當然的狀況或條件突然發生變化－我們回望過去,重新審視,思考自己走過的路.我們發覺自己一直的執著與堅持,或許未必好像自己一直以為那麼理所當然.

在這充滿偶發性,有時好像雜亂無章的現實世界,我們一直所秉持的原則,其實未必能夠統攝,駕馭生命的複雜和多元性.我們很多時只不過是方便地將那些難以理解,不能歸類的,擱置在一旁,或暫時掃在地毯底.到最後,這一切被我們掃埋一邊的片段,時刻,或會回頭向我們提出質問.

就像那位一直理直氣壯,以為自己俯仰無愧的少年官那樣,我們走到生命的主面前,向他求問:「我還要做甚麼,才可以不枉此生?」又或者好像馬大那樣,忙忙碌碌,充充實實地過了大半生,然後我們問:「我為你勞碌了大半生,你不在意嗎?請吩咐我的子女,朋友,那個我服務大半生的機構,團體或者教會,回報我這些年來作出的貢獻和犧牲.」然後我們聽見主耶穌說:「你還不明白嗎?我並沒有要求你作出你覺得委曲的貢獻和犧牲.你只需要如實地活出自己的生命,做你真正覺得正確和舒坦的事情,並且在當中享受生命.而這個權利,是沒有人能夠剝奪你的.」

這個,我覺得就是主今日透過這段經課,給我們的提醒.



\newpage

\section{真理是甚麼呢}
\label{sec:463}
\textbf{劉國偉先生 - 真理是甚麼呢?}
\newline
\newline
\hyperref[sec:462]{< < < PREV SERMON < < <}
~
\hyperref[sec:toc]{[返目錄]}
~
\hyperref[sec:464]{> > > NEXT SERMON > > >}


約翰福音 18:28-38
\newline
\begin{longtable}{cl}
\hline
\hline
章節 & 經文 (和合本修訂版)\\
\hline
18:28 & \begin{tabularx}{0.7\textwidth}{X} 他們把耶穌從該亞法那裡押解到總督府。那時是清早。他們自己卻不進總督府，恐怕染了污穢，不能吃逾越節的宴席。 \end{tabularx} \\ \\ \relax
18:29 & \begin{tabularx}{0.7\textwidth}{X} 於是彼拉多出來，到他們那裡，說：「你們告這人是為甚麼事呢？」 \end{tabularx} \\ \\ \relax
18:30 & \begin{tabularx}{0.7\textwidth}{X} 他們回答他說：「這人若不作惡，我們就不會把他交給你了。」 \end{tabularx} \\ \\ \relax
18:31 & \begin{tabularx}{0.7\textwidth}{X} 彼拉多對他們說：「你們自己帶他去，按著你們的律法問他吧！」猶太人說：「我們沒有殺人的權柄。」 \end{tabularx} \\ \\ \relax
18:32 & \begin{tabularx}{0.7\textwidth}{X} 這是要應驗耶穌所說，指自己將要怎樣死的話。 \end{tabularx} \\ \\ \relax
18:33 & \begin{tabularx}{0.7\textwidth}{X} 於是彼拉多又進了總督府，叫耶穌來，對他說：「你是猶太人的王嗎？」 \end{tabularx} \\ \\ \relax
18:34 & \begin{tabularx}{0.7\textwidth}{X} 耶穌回答：「這話是你說的，還是別人論到我時對你說的呢？」 \end{tabularx} \\ \\ \relax
18:35 & \begin{tabularx}{0.7\textwidth}{X} 彼拉多回答：「難道我是猶太人嗎？你的同胞和祭司長把你交給我。你做了甚麼事呢？」 \end{tabularx} \\ \\ \relax
18:36 & \begin{tabularx}{0.7\textwidth}{X} 耶穌回答：「我的國不屬於這世界；我的國若屬於這世界，我的部下就會為我戰鬥，使我不至於被交給猶太人。只是我的國不屬於這世界。」 \end{tabularx} \\ \\ \relax
18:37 & \begin{tabularx}{0.7\textwidth}{X} 於是彼拉多對他說：「那麼，你是王了？」耶穌回答：「是你說我是王。我為此而生，也為此來到世界，為了給真理作見證。凡屬真理的人都聽我的話。」 \end{tabularx} \\ \\ \relax
18:38 & \begin{tabularx}{0.7\textwidth}{X} 彼拉多對他說：「真理是甚麼呢？」 \end{tabularx} \\ \\
[1ex]
\hline
\hline
\end{longtable}
講題:真理是甚麼呢? What is Truth?

經文:約翰福音18:28-38

講員:劉國偉先生

場合:香港中文大學崇基學院禮拜堂

日期:2019年7月28日

引言

猶太裔著名的政治及思想家漢娜•鄂蘭(Hannah Arendt)在1963年出版了《艾希曼在耶路撒冷--關於艾希曼審判的報告》,提出了其後國際著名的“平庸之惡”,或者更應該譯作”邪惡之平庸”這個概念(banality of evil).

1961年4月11日,以色列政府在耶路撒冷進行一次審訊,是有關二次世界大戰時,一位納粹德國軍官阿道夫·艾希曼(Adolf Eichmann)的審訊,審訊持續到同年5月31日,結果艾希曼被判處絞刑.

當時,漢娜•鄂蘭以《紐約客》特約撰稿人的身份報導這次審訊.然後於1963年出版了上述那本書,《艾希曼在耶路撒冷--關於艾希曼審判的報告》,並提出了著名的“平庸之惡”或者”邪惡之平庸”這個概念.

阿道夫·艾希曼(Adolf Eichmann),是納粹德國的高級軍官,亦是猶太人大屠殺的“最終方案”的主要負責官員.1942年艾希曼出席萬湖會議,被任命負責執行屠殺猶太人的最終方案,更獲晉升為中校.

二次世界大戰之後,艾希曼曾被美軍俘虜,但後來成功逃脫,並經過漫長的逃亡生涯,最後於1961年,被以色列的情報部門在阿根廷逮捕.耶路撒冷以人道罪名等十五條罪名起訴艾希曼,結果罪成,並在1962年6月1日處死.

漢娜•鄂蘭記錄了艾希曼的自辯,反覆強調他自己只不過是齒輪系統中的一環,作出了傳動的作用.艾希曼相信自己所做的事,都是當時國家法律所允許的；他是一名軍人,只在服從和執行上級的命令.

漢娜•鄂蘭甚至形容艾希曼並不陰險,完全不像一個劊子手,而且彬彬有禮坐在犯人欄內.漢娜•鄂蘭認為極權統治者有極端之惡,這是第一種邪惡；而第二種邪惡是管治的參與者和被統治者的平庸之惡,那就是對於明顯不過的邪惡行為不予抵制,及不加限制,甚至因制度和機制,而甘心直接參與其中,這就是平庸之惡.鄂蘭的結論是,很多邪惡事情的發生,只因有平庸的人不假思索地執行惡令.

我想起使徒保羅在哥林多前書愛的篇章所說:愛是不喜歡不義,只喜歡真理；真正的愛包含辨識不義,並且不會喜歡不義,更努力追尋真理,因為只喜歡真理.若果有人在不義的世代,意想置身事外,單單追求個人的安穩逸樂,對出現於生活四週的邪惡的事不聞不問,嘗試以平庸作為一個選項,那並沒有認真地活出愛,甚至有可能助長邪惡的滋長.

猶太人與彼拉多

今日的講道經文,約翰福音18章28至38節上這個大段落,又可以分成兩個小段.第一段由28節至32節,另一段由33節至38節上.

第一個小段的焦點,在於猶太人與彼拉多之間的互動.首先,使徒約翰告訴我們,該亞法與猶太人十分謹慎地遵從他們的猶太教法規.18章28節有這個記載:「他們自己卻不進衙門,恐怕染了污穢,不能吃逾越節的筵席.」熟悉舊約的人,都知道猶太人極少進到外邦人的住處,因為這樣等同沾染污穢.當時是星期五的上午,到那天黃昏就是逾越節了,若果來不及沐浴更新,帶著污穢就不能吃逾越節的筵席.也許這種謹慎,是值得尊重甚至敬佩的,不過,當彼拉多說:「你們自己帶他去,按著你們的律法審問他吧.」猶太人卻說:「我們沒有殺人的權柄.」(約18:31)就反映出該亞法與猶太人一方面煞有介事地遵守猶太法規,卻另一方面竟然誣害耶穌,甚至動了殺機,並打算藉著彼拉多的手置耶穌於死地.流無辜人的血當然嚴重地違反律法精神,我們看見的是該亞法和猶太人遵守法規的虛偽,卻同時進行踐踏律法的殺機和行動.

彼拉多與耶穌

第二個小段記載了彼拉多與耶穌的對話；這段對話包含了彼拉多的5個問號.

你是猶太人的王麼?(約18:33；現中/你是猶太人的王嗎?)

彼拉多進入他的總督府,開始審問耶穌,這就是彼拉多向耶穌發出的第一個問題.耶穌並沒有直接回答他,卻反問彼拉多:「這話是你自己說的,還是別人論我對你說的呢?」耶穌打開一個對話的空間,讓彼拉多有機會想一想到底他對耶穌的印象,是基於猶太人的論述,還是他自己的親身印象.很可惜,彼拉多無意把握這個與耶穌對話的機會,卻向耶穌提出一個反問.

我豈是猶太人呢?(約18:35；現中/你以為我是猶太人嗎?)

彼拉多這個反問,表達了他完全無意去捲入耶穌是猶太人的王這個爭論,他一直想置身事外,明哲保身,不想招惹麻煩.他向耶穌表明,全是因為猶太人和祭司長把耶穌交給他,才會在他的總督府進行這次審問.於是,彼拉多就簡單直接的問耶穌.

你作了甚麼事呢?(約18:35；現中/你做了什麼事呢?)

耶穌的回答,也可能叫彼拉多摸不著頭腦,他說:「我的國不屬這世界,我的國若屬這世界,我的臣僕必要爭戰,使我不至於被交給猶太人；只是我的國不屬這世界.」(約18:36)不過,這個回覆至少讓彼拉多知道,耶穌在政治的層面,對羅馬帝國並沒有任何實質的威脅.猶太人意圖捏造耶穌的罪名,就是以「猶太人的王」作為罪證,指控耶穌背叛凱撒,從而構成因政治罪名,而引致死刑.

這樣,你是王麼?(約18:37；現中/那麼,你是王了?)

彼拉多再次查問耶穌到底他是否認為自己是王；這可以是提供耶穌一個脫罪的機會,只要耶穌說不是,事情也許會有轉機；不過,耶穌面向死亡,卻無意退縮,他明明白白地向彼拉多說:「我為此而生,也為此來到世間.」(約18:37)

真理是甚麼呢?(約18:38；現中/真理是什麼?)

彼拉多意圖不沾鍋.把耶穌推向希律王的計謀失敗了；把耶穌退回給猶太人的計劃,也失敗了；期待耶穌為自己辯護,讓他自己可以簡單輕便的處理這次審訊,但因為耶穌並沒有他所預期的合作態度,叫彼拉多進退不得.

彼拉多提出真理是甚麼?但他真的需要答案嗎?他沒有等待耶穌的回應,就又走到總督府外邊,向猶太人說話.不過,使徒約翰卻巧妙地把這個問題的答案寫在前面.

「我就是真理」

「真理」這個詞彙是約翰福音一個特定選用並且經常採用的詞彙.約翰採用有關真理這個詞彙的時候,較多採用名詞的方式,而且都帶有獨特的宗教含意.

約翰福音是以希臘文寫作,希臘文「真理」alētheia這個字,含有真實和誠實的含意,他的反義詞就是虛有其表和謊言,主要是屬於認知的範疇,透過知識去分辨.

對應希伯來文「真理」èmet這個詞彙,亦具有堅固,信實的含意,不過它不單屬於認知的範疇,更屬於道德的範疇.

使徒約翰成功地為希臘文alētheia「真理」這個字注入了希伯來的思想養份.在希臘思想哲學中,有關真理的知識具有客觀性,真理是經由認知的過程來吸收；在希伯來的思想中,真理具有位格,與人的內心和良知有十分重要的關連.在希伯來思想中,妨礙一個人獲得真理的,主要並非因為他的知識不足或認知能力稍遜,卻因為一顆剛硬的心,就正如先知以賽亞所感慨:「耶和華啊,你為何使我們走差,離開你的道,使我們心裡剛硬,不敬畏你呢?」(賽63:17),當一個人驕傲和自以為是,他的內心就失去尋求上帝和真理的動力和傾向,縱使他有聰明和知識,卻可能仍然與真理相去甚遠.

當耶穌說「我就是真理」,那其實是一個屬於希伯來人的思想方式,為真理alētheia這個希臘字注入了位格,而這個位格就是聖子耶穌.

在約翰福音中,這個帶有宗教獨特含意的詞彙,主要有三個用法.第一,表達從上帝而來的真理教導,涉及行為實踐.「行真理的必來就光,要顯明他所行的是靠上帝而行.」(約3:21) 或者「你們是出於你們的父魔鬼,你們父的私慾你們偏要行.牠從起初是殺人的,不守真理,因牠心裡沒有真理.」(約8:44)

第二,在上帝裡的真理,真真實實呈現出來,指向上帝的真實.「道成了肉身,住在我們中間,充充滿滿地有恩典有真理.」(約1:14)又例如「律法本是藉著摩西傳的；恩典和真理都是由耶穌基督來的.」(約1:17)

第三,真理就是耶穌基督本體

「我就是道路,真理,生命；若不藉著我,沒有人能到父那裡去.」(約14:6)

至少有兩處經文同時展現了全部三個用法.約翰福音8章32節「你們必曉得真理,真理必叫你們得以自由.」這節經文就是當年燕京大學校訓的主要基礎,燕京大學的校訓是「因真理,得自由,以服務.」

「必曉得真理」既表示走在上帝道路中的人,知道上帝的真實和祂的教導,而這金句也可以同時指向耶穌基督的本體,因為耶穌基督就是「真理」,是耶穌基督叫我們得著真正的自由.

另一處同時展現了全部三個用法的經文記載於約翰福音17章17節:「求你用真理使他們成聖；你的道就是真理.」這節金句表示真理使我們成聖,就是上帝的真理教導,上帝的真實和聖子耶穌,叫我們成聖；下半節更直接指明「你的道就是真理」,這裡所用的「道」這個字,就是太初有道的道,Logos,簡單直接地說就是耶穌基督.

真理與我相遇

猶太人在大祭司該亞法及猶太公會的領導下,以平庸的身份,或默不作聲,或甚至不假思索地直接參與他們殺害耶穌的邪惡行動.

提問「真理是甚麼呢?」的彼拉多,也作出了一個平庸的決定,他為了避免猶太領袖以「不是凱撒的中神」的指控,沒有好好運用他的權柄,最終埋沒良知,參與誣害和殺死耶穌.

真理與我們相遇的其中一個重要意義,是真理叫我們成聖,若果我們認真地看待成聖這回事,那我們一定不會讓自己置身於邪惡的平庸.

成聖sanctify,希臘文hagiazo 即分別出來預備作聖潔的用途,亦含有成為聖潔的意思.換言之,與真理相遇的人,生活變得很不一樣,特別在邪惡的世代中,做一個追尋真理,堅守真理的人,我們才可以去經歷上帝的真實,並且在真理的聖靈引導下,進入真理,也就是緊密地有耶穌基督同在.

陳健民教授在他的最新一輪的網上專欄「獄中書簡」中,剛好提及《平庸之惡》的作者漢娜•鄂蘭及她的丈夫海因裡希•布呂赫 (Heinrich Blücher).陳健民形容他們活在德國納粹時期,最終流亡海外,一直思考公民在亂世的責任.海因裡希說活在那個世代,「悲觀是懦弱,但樂觀卻是愚蠢」.我們亦是活在一個考驗勇氣和智慧的時代.但願我們知道怎樣回答甚麼是真理,靠賴聖靈進入真理,並且因真理得自由,因真理成聖,宣揚上帝的愛,又因為愛的緣故,我們只喜歡真理.



\newpage

\section{作為基督新婦的教會}
\label{sec:464}
\textbf{伍渭文牧師 - 作為基督新婦的教會}
\newline
\newline
\hyperref[sec:463]{< < < PREV SERMON < < <}
~
\hyperref[sec:toc]{[返目錄]}
~
\hyperref[sec:465]{> > > NEXT SERMON > > >}


以弗所書 5:26-27,31-32
\newline
\begin{longtable}{cl}
\hline
\hline
章節 & 經文 (和合本修訂版)\\
\hline
5:26 & \begin{tabularx}{0.7\textwidth}{X} 以水藉著道把教會洗淨，使她成為聖潔， \end{tabularx} \\ \\ \relax
5:27 & \begin{tabularx}{0.7\textwidth}{X} 好獻給自己，作榮耀的教會，毫無玷污、皺紋等類的缺陷，而是聖潔沒有瑕疵的。 \end{tabularx} \\ \\ \relax
5:28 & \begin{tabularx}{0.7\textwidth}{X} 丈夫也應當照樣愛妻子，如同愛自己的身體；愛妻子就是愛自己了。 \end{tabularx} \\ \\ \relax
5:29 & \begin{tabularx}{0.7\textwidth}{X} 從來沒有人恨惡自己的身體，總是保養愛惜，正像基督待教會一樣， \end{tabularx} \\ \\ \relax
5:30 & \begin{tabularx}{0.7\textwidth}{X} 因我們是他身體的肢體。 \end{tabularx} \\ \\ \relax
5:31 & \begin{tabularx}{0.7\textwidth}{X} 「為這個緣故，人要離開父母，與妻子結合，二人成為一體。」 \end{tabularx} \\ \\ \relax
5:32 & \begin{tabularx}{0.7\textwidth}{X} 這是極大的奧祕，而我是指基督和教會說的。 \end{tabularx} \\ \\
[1ex]
\hline
\hline
\end{longtable}
講題:作為基督新婦的教會 The Church As the Bride of Christ

經文:以弗所書5:26-27,31-32

講員:伍渭文牧師

場合:香港中文大學崇基學院神學院容啟東校長紀念樓G/F聖堂

日期:2019年8月4日

讀經:正如基督愛教會,為教會捨己.要用水藉著道把教會洗淨,成為聖潔,可以獻給自己作個榮耀的教會.(弗五26, 27)人要離開父母,與妻；子聯合,二人成為一體.這是極大的奧秘,但我是指著基督和教會說的.(弗五31,32)

有人用「主人與貓」比喻中港關係.他開始飼養小貓時,很多朋友指貓會抓壞傢俬,但他沒有用籠困住小貓,「我發現同佢相處ge 過程中,佢好乖,咁我就比好大ge 空間佢」.撇開比諭本身是否合宜,香港一向以來被主人視為掌上明珠-東方之珠,突然變成寵物,總覺有點不習慣.

聖經用夫婦比諭教會與基督的關系.上主與信徒的關系,不是人與寵物的關系,乃人世間最親密的夫婦平等關系.婚姻禮在羅馬天主教(公教)被視為七個聖事(聖禮)中,唯一不須經由有效的禮儀和有效的主禮者而確立的聖事,如按立聖秩(聖品),須經主教；婚禮成就聖事在新人相方的盟誓,禮儀婚約採用盟約Covenant而不是民事婚姻採用的合約Contract.盟約正確是切約,cut the covenant,立約雙方把祭牲切開分兩面擺設,然後在其中行過,以生命保證約的莊嚴恆久；若然背約,將如祭牲切開.創十五17,18上:「日落天黑,不料有冒煙的爐並燒著的火把從那些肉塊中經過.當那日,耶和華與亞伯蘭立約.」

貫穿整個舊約聖經.約確立上主與選民的關系,有挪亞方舟之約,亞伯拉罕割禮之約,摩西西乃之約.這三個約都是聖洗的預表,而聖洗就是基督用來禮潔淨教會,成為他的新婦.讓我們從三個預表聖洗的約,思考教會作為基督新婦的意思,思考我們與基督的關系.

一 挪亞方舟之約-歡愉的對換

「當時進入方舟,藉著水得救的不多,祇有八個人.這水所表明的洗禮,現在藉著耶穌基督復活也拯救你們；這洗禮本不在乎除掉肉體的污穢,祇求在上主面前有無愧的良心.」(彼前三20,21)

二 亞伯拉罕割禮之約-信心的後裔

「你們在他裡面也得了不是人手所行的割禮,乃是基督使你們脫去肉體情慾的割禮,你們既受洗與他一同埋葬,也就在此與他一同復活.」(西二11,12)

三 摩西西乃之約-邁向應許地

「弟兄們,我不願意你們不匙曉得,我們的祖宗從前都在雲下,都從海中經過,都在雲裡,海裡受洗歸入摩西.」(林前十1,2)

一 婚姻是歡愉的對換

1基督以自己的生命為聘禮迎娶我們.

基督本不須到約但河領悔改的洗禮,因他沒有罪.他是以彌賽亞身份站在約但河,接受象徵洪水審判的水禮.耶穌洗禮時,有什麼異常景象呢?“他從水裡一上來就看見天裂開了,聖靈彷彿鴿子降在他身上.又有聲音從天上來說,你是我的愛子,我喜悅你.”(可一10,11) “你是我的愛子,我喜悅你.”這句話是彌賽亞的名號:“看哪,我的僕人,我所扶持,所揀選,心裡所喜悅的,我已將我的靈賜給他,他必將公理傳給外邦.”(賽42:1)

耶穌從水裡上來時,聖靈像鴿子降在耶穌身上,象徵審判過去,喜樂來臨.聖經在此之前,何處曾說到從水裡上來,有鴿子降在其上呢?是當洪水漫過大地四十天后,水止住了,大地露出來,大地從水裡上來,鴿子從方舟放出去,嘴裡含著新擰下來的橄欖葉子.(創8:6-12)

一月六日紀念耶穌領洗的顯現日,是歡樂的節日,聖帷是白色或金色,是十二天聖誕期的高潮和終結.屈梭多模(金口約翰John Chrysostom)主教說:”今天基督進入了約旦河,以洗滌世人的罪過,今天上主的聲音,響徹水面,你是我的愛子,我喜悅你.今天,聖靈借鴿子降臨在水面上,報告人類永遠滅亡的災禍已停止.」”(每日頌禱[Liturgia Horarum]一,頁308.)

顯現日耶穌洗禮的禱文中,提及耶穌在約旦河潔淨他的新婦教會.「今天教會與他天上的配偶結合了,因為基督在約旦河洗去了他的罪惡；博士們帶著禮物迅速前來,奔赴王家的婚宴；變水為酒,使賓客歡心,哈利路亞.」顯現後第二主日福音經課為約二1-11,迦拿婚宴水變酒.

進入婚姻的人,甘心樂意分享自己所有,和分擔對方所需.在婚禮中新人盟誓:「我以此戒指為徵記,與你結婚,我以我的身心尊敬你,我以我的財物,與你分享.」,用在基督和教的關系中,路德說這是歡愉的對換.基督在約但河領洗說明他為我們的罪接受洪水的審判,他就是那方舟,他被洪水擊打,但保護進入方舟的人,得免於難.他把我們的罪孽擁攬在身上,被釘十架,他以自己的生命作為聘禮,迎娶我們.聖靈提醒我們:「這是我的身體,為你們而捨的.」

2永恆的約.婚姻歡愉的對換,上主沒有按我們應得的來審判我們,乃按我們的需要來饒恕我們,基督就像那被洪水擊打的方舟,讓我們在裡面的得免洪水的審判.挪亞方舟之約以彩虹為証,是永恆之約.「我與你們立約,凡有血肉的,不再被洪水滅絕,也不再有洪水毀壞地了.……我把虹放在彩雲中,這就可作我與地立約的記號了.」(創九11,13)

愛因斯坦曾經被人問:「你對這宇宙最想知道的是甚麼?」他說:「我很想知道這宇宙是否對我們友善?」相對宇宙,地球顯得渺小脆弱,要是宇宙不友善,飛來巨大隕石,引至地震海嘯,洪水災難再現.這約不像人的聯合聲明,隨時會變成歷史文件.

應用:聖經以人間夫妻關系類比基督和教會的關系.「但我是指著基督和教會說的」(弗五32)夫妻之相處,要效法基督如何愛教會,為教會捨己,顯出憐憫,饒恕,恩典,新的開始.憐憫按需要,律法按應得.感謝上主,憐憫是向審判跨勝.

二 亞伯拉罕割禮之約-信心的後裔

「你們在他裡面也得了不是人手所行的割禮,乃是基督使你們脫去肉體情慾的割禮,你們既受洗與他一同埋葬,也就在此與他一同復活.」(西二11,12；創十七)創十七詳細記載亞伯拉罕要為族人施行割禮,嬰孩生下來第八天就行割禮,他自己九十九歲也行了割禮,作為立約的記號.在割禮這一章,特別多談到信心的後代.

1屬靈的傳生.「我是全能的神.你當在我面前作完全人,我就與你立約,使你的後裔極其繁多.」(創十七1)

「我要將你現在寄居的地,就是迦南全地,賜給你和你的後裔永遠為業,我也必作他們的神.」(創十七8)

「你和你的後裔必世世代代遵守我的約.你們所有的男子都要受割禮:這就是我與你並你的後裔所立的約,是你們所當遵守的.」(創十七9)

亞伯拉罕是信心之父,但他都要過去.我們祇能夠親自服事自己那一代的人,就過去了.但用生命誕生生命,養育裁培生命,生命就茁壯成長,生生不息.

婚禮儀節說:婚姻聖禮乃為遵行上帝旨意繁殖人類而設,教養子女敬愛上帝,頌揚主之聖名.

婚姻有傳生(procreation, creation for and on behalf of the Creator為創造主創造生命-造人仔)的召命.婚姻的意義,超越夫婦兩人之間的愛情,參與上主創造生命,承擔傳生的天職,生兒育女.男女的眷愛和生兒育女是婚姻的冠冕和峰頂.男女的眷戀是人的選擇,傳生是上主的主權.從來沒有人被之諮詢過才生下來,祇有那一位,但他答應了.婚姻成立家庭,家庭是教會和社會最基本的單位.婚姻使生命繁衍,文化延續,天國早日降臨.挪亞一家得救,開出萬族,創九9:「我與你們和你們的後裔立約」；亞伯拉罕蒙召成一族,後裔多如天上的星地上的塵；到摩西從西乃山頒下一國的典章,在萬民中作屬上主的子民,君尊的祭司,聖潔的國度.

信徒婚後若蒙上主賜下兒女,要用真理養育他們,成為信仰的後裔.尚未結婚和尚未有兒女的,也藉履行大使命,使萬民作主門徒,這是屬靈的傳生使命.在教會裡,我們可以培育後輩；巴拿巴助初信的保羅一把,保羅當兒子般培養提摩太,行經時間,在他人生命中留下痕跡.

保羅在帖前二19,20說:「我們的盼望和喜樂,並所誇的冠冕是甚麼呢?豈不是我們主耶穌來的時候,你們在他面前站立得住嗎?因為你們就是我們的榮耀,我們的榮耀,我們的喜樂.」我們等候最終的羔羊婚宴,這位將要再來的新郎,也是一位君王.他來時像當時的王來巡視其屬土,每個城邦在王來的時候,用最好的珠寶造成冠冕,呈給王帶上.保羅說,你們就是我要呈給王的冠冕,令我喜樂和榮耀的冠冕.

2改名進入新世界:亞伯蘭改名亞伯拉罕.改名表示新的關系,你的名字,我的姓氏.改名也表示進入嶄新的世界.小孩,兩,三歲時開始為世誡命名,起初爸爸,媽媽,跟著爺爺,麻麻…….進一步講述他們在這新世界的故事,講故事就找到在新世界的身份.爺爺是爸爸的爸爸,我有一個曾弟弟.結婚時,女的改名冠夫姓,結束婚禮步出禮堂:隆重介紹:某某先生夫人.新婚的人也進入嶄新的姻親世界,命名:岳父岳母,老爺奶奶.甚麼姨媽姑姐.信主奉聖父聖子聖靈的名受洗,開始崇拜也奉父,子,靈的名.這是聖堂,看的是聖經,唱的是聖詩.我們也述說自己在這嶄新世界的故事,說明我們的身份.領聖餐與眾肢體分享基督的身體,等候最後審判,身體復活,羔羊的婚筵.

三 摩西西乃之約-邁向應許地

婚禮宣諭:為使夫婦在福樂困苦之中,皆能互相親愛,彼此扶助安慰.(終身朋伴)婚姻是一生的承諾.証婚牧師問新人:「你是否願意與某某結為夫婦,遵行上帝婚姻聖禮之旨意,與他共同生活,無論疾病康健,你是否願意時常敬愛他,安慰他,尊重他,扶助他,與他一生共守,決無異心.」

婚姻和信仰,也是漫長曲折的旅程.摩西西乃之約,叫我們決心跨過紅海,離開舊土,來到充滿煙雲之西乃山,與上主立約,領受律法,邁向應許之地.在人生曠野路上,婚姻賜我們良友共走,行經時間,路遙不知疲.

但以色列人為何要離開埃及?出埃及原因是(1)得著應許之地 ,出三8「我下來是要救他們脫離埃及人的手,領他們到那地,到美好,寬闊,流奶與蜜之地.」(2)來到曠野敬拜上主.上主對摩西說:「對他(法老)說……容我的百姓去,好在曠野事奉我(worship me, NIV).」(出七16,亦參出八1; 九1,13；十3)

決心離開埃及.法老最初希望摩西在埃及地經驗西乃之旅.「你們去,在這地祭祀你們的神吧.」(出八25)法老之後讓步,以為摩西視西乃的敬拜是尋求特別的宗教經驗,之後還要回到埃及地.法老說,男的可以去,但留下女和兒童；男,女,兒童都可以去,但要留下牲口.往西乃路程祇需三天(出八27),來回不要一周時間,西乃敬拜和埃及地不需割切.牲口就是「畜質」,留下回來埃及的牽引；牲口是生產工具,沒有牲口,祇有在曠野成餓殍了(出十24).摩西三度拒絕法老的建議,要把一切都帶走,分毫不留埃及地,因為他們響往列祖應許之地.

婚姻首要是放下行李,很多婚姻產生困難,因為結婚多年,仍然未放下行李,把配偶和其人比較.「人要離開父母,與妻子聯合,二人成為一體.」離開就是放下行李.到應許之地首先要決心過紅海.

但何處是應許之地呢?其實迦南或埃及,流奶與為奴,分別在西乃山的敬拜.往後的歷史證明,失去對上主的敬拜,就失去應許之地,甚至被擄去異邦.耶穌同在,就是天堂.雖然是三天的路程,因為頑梗不順服,也會延宕至三個月.「滿了三個月的那一天,就來到西乃的曠野.」 (出十九1, 2)

西乃並非中途竭息之所,乃安頓心靈應許之地.沒有西乃的敬拜生活,踏足之沃土如浮沙,所見之美地像蜃樓,出埃及是為了來到西乃.事實上,以色列被譽級為經中之經的五經,不在應許之地迦南獲得,是在西乃曠野頒下.

四 羔羊婚筵.

基督用水藉著道(聖洗)潔淨他的新婦教會,這是絢爛旅程的開始,而用聖餐提醒我們這旅程的終結,聖餐是終末羔羊的預嘗.

1終末的婚宴,世界終局是歡樂,充滿盼望.所有事情的發生,都邁向上主的計劃,完成他隱密的旨意,沒有偶然的事情,祇是當時不明白.不錯,我們要像聰明的童女,儆醒等候新郎的來臨,也要敬畏最後的審判,但終末是一個婚宴.最後的審判,不是最後是審判,是最後是歡樂的婚宴.生命樹是終末,分別善惡的樹都要過去.

2終末的婚宴是永遠的婚宴.人為上主而被造,不是為一個男士或女而被造.真正的滿足,不在家庭,不在婚姻,否則我們會失望.是在基督裡的婚姻才能使我們滿足,不是婚姻使我們滿足.因此,婚姻是一項選擇,在教會裡我們也有朋伴,也有屬靈的兒女,靠著上主給予的恩典,獨身並不孤身.

結語:往昔雖可留戀,未來更值期待.過去一星期就在這地方舉行為期六天的聖樂營,250位來自境外的參加者,連同三十位導師共280人坐滿這聖堂,晨禱晚頌,專心思考大會主題:「基督與教會的婚盟」.昨天收到一位營友給我的專題回饋.「為赴那終末的婚宴,我們聚集敬拜仿若一場又一場的彩排,唯願新郎喜悅,因他已以生命做聘禮娶了他的新娘……今日不記得在哪個環節了,當耳邊環繞詩班輪流的讚美歌聲,我眼前出現如電影般的畫面-看見仿佛列車的門打開了,自己被裹挾在赴終末婚宴的人流之中身不由己往前走,遠遠見到一人在人流中逆流而來,奮力向我靠近,欣然發現是我的先生!我們彼此相擁而笑,似乎並不覺得分開許久,也不詫異這樣的相聚,仿佛早就約好了的一般,他那興奮的樣子,就像是在一票難求的新年音樂會上給我占了一個好位置!這突如其來的意像,令我淚水奔湧而出,滑落臉頰,久不能抑!一種從未有過的感受油然而生:往昔雖可留戀 ,未來更值期待.」



\newpage

\section{釘死你們的王}
\label{sec:465}
\textbf{劉國偉先生 - 釘死你們的王?}
\newline
\newline
\hyperref[sec:464]{< < < PREV SERMON < < <}
~
\hyperref[sec:toc]{[返目錄]}
~
\hyperref[sec:466]{> > > NEXT SERMON > > >}


約翰福音 19:1-16
\newline
\begin{longtable}{cl}
\hline
\hline
章節 & 經文 (和合本修訂版)\\
\hline
19:1 & \begin{tabularx}{0.7\textwidth}{X} 於是，彼拉多命令把耶穌帶去鞭打了。 \end{tabularx} \\ \\ \relax
19:2 & \begin{tabularx}{0.7\textwidth}{X} 士兵用荊棘編了冠冕，戴在他頭上，給他穿上紫袍， \end{tabularx} \\ \\ \relax
19:3 & \begin{tabularx}{0.7\textwidth}{X} 又走到他面前，說：「萬歲，猶太人的王！」他們就打他耳光。 \end{tabularx} \\ \\ \relax
19:4 & \begin{tabularx}{0.7\textwidth}{X} 彼拉多又出來對眾人說：「看，我帶他出來見你們，讓你們知道我查不出他有甚麼罪狀。」 \end{tabularx} \\ \\ \relax
19:5 & \begin{tabularx}{0.7\textwidth}{X} 耶穌出來，戴著荊棘冠冕，穿著紫袍。彼拉多對他們說：「看哪，這個人！」 \end{tabularx} \\ \\ \relax
19:6 & \begin{tabularx}{0.7\textwidth}{X} 祭司長和聖殿警衛看見他，就喊著說：「釘十字架！釘十字架！」彼拉多對他們說：「你們自己把他帶去釘十字架吧！我查不出他有甚麼罪狀。」 \end{tabularx} \\ \\ \relax
19:7 & \begin{tabularx}{0.7\textwidth}{X} 猶太人回答他：「我們有律法，按照律法，他是該死的，因為他自以為是神的兒子。」 \end{tabularx} \\ \\ \relax
19:8 & \begin{tabularx}{0.7\textwidth}{X} 彼拉多聽見這話，越發害怕， \end{tabularx} \\ \\ \relax
19:9 & \begin{tabularx}{0.7\textwidth}{X} 又進了總督府，對耶穌說：「你是哪裡來的？」耶穌卻不回答。 \end{tabularx} \\ \\ \relax
19:10 & \begin{tabularx}{0.7\textwidth}{X} 於是彼拉多對他說：「你不對我說話嗎？難道你不知道我有權柄釋放你，也有權柄把你釘十字架嗎？」 \end{tabularx} \\ \\ \relax
19:11 & \begin{tabularx}{0.7\textwidth}{X} 耶穌回答他：「若不是從上頭賜給你的，你就毫無權柄辦我，所以，把我交給你的那人罪更重了。」 \end{tabularx} \\ \\ \relax
19:12 & \begin{tabularx}{0.7\textwidth}{X} 從此，彼拉多想要釋放耶穌，無奈猶太人喊著說：「你若釋放這個人，你就不是凱撒的忠臣。凡自立為王的就是背叛凱撒。」 \end{tabularx} \\ \\ \relax
19:13 & \begin{tabularx}{0.7\textwidth}{X} 彼拉多聽見這些話，就帶耶穌出來，到了一個地方，叫「鋪華石處」，希伯來話叫厄巴大，就在那裡坐堂。 \end{tabularx} \\ \\ \relax
19:14 & \begin{tabularx}{0.7\textwidth}{X} 那日是逾越節的預備日，約在正午。彼拉多對猶太人說：「看哪，你們的王！」 \end{tabularx} \\ \\ \relax
19:15 & \begin{tabularx}{0.7\textwidth}{X} 他們就喊著：「除掉他！除掉他！把他釘十字架！」彼拉多對他們說：「要我把你們的王釘十字架嗎？」祭司長回答：「除了凱撒，我們沒有王。」 \end{tabularx} \\ \\ \relax
19:16 & \begin{tabularx}{0.7\textwidth}{X} 於是彼拉多把耶穌交給他們去釘十字架。 \end{tabularx} \\ \\
[1ex]
\hline
\hline
\end{longtable}
講題:釘死你們的王? Shall I crucify your king?

經文:約翰福音19:1-16

講員:劉國偉先生

場合:香港中文大學崇基學院神學院容啟東校長紀念樓G/F聖堂

日期:2019年8月11日

耶路撒冷東門

五月的時候,與基督徒領袖訓練計劃CLAP的十多位同學,一起到以色列學習,耶路撒冷當然是學習的重點地方.在行程第六天的早上,我們去到橄欖山,遠望對面的聖殿山,當我向同學解說眼前耶路撒冷東面的城牆時,我很自然就提起了耶穌基督當年從耶路撒冷東門進城的景象.

當然,我們現在看見的耶路撒冷城牆,主要是十六世紀土耳其鄂圖曼帝國所建造的,與耶穌基督進入耶路撒冷的日子,相距一千五百多年；考古學家相信,十六世紀建造的城牆的方位,與耶穌基督的年代基本相似,而耶穌年代的東門,就是在我們如今看見的這個東門的底下.

橄欖山與聖殿山之間是一個V型山谷,就是汲淪溪谷,從耶路撒冷東門進城,即是從城牆外的郊區向上攀行,換言之,耶穌與門徒向上行走的時候,是不可能清楚看見聖殿山上偌大的平台廣場的景象,不過,他們應該可以看見城門口人頭湧湧的情景,而且愈走近群眾那裡,就愈發聽到人聲鼎沸,更看見民眾拿著棕樹枝迎接,呼喊著說:「和散那! 奉主名來的以色列王是應當稱頌的!」

「和散那」就是「現在拯救」的意思；民眾稱耶穌是以色列王,這是發生在耶路撒冷東門的一起十分重要的歷史事件.由東門直上,就是通往聖殿的主要通道,那時因為逾越節期將近,耶路撒冷合成都充滿了遊人和客旅,尤其是在聖殿山的平台廣場上.

當天,那些激情地在東門口拿著棕樹枝,呼喊耶穌是以色列王的民眾,也許萬料不到,只不過在短短的時間,耶穌竟然因為「以色列王」的稱號,而被置於死地,彼拉多更命人在耶穌的十字架上方,釘上了一個寫著希伯來,羅馬,希臘三種文字的牌子:「猶太人的王,拿撒勒人耶穌.」

我們今天站在橄欖山遠望的東門,已經被封閉,所以我們只能夠看見一個城門的痕跡,而原來的通道則成為一幅石牆,更不用說,耶穌年代的東門而埋藏在現在我們眼見的城牆底部.這幅石牆,或者這個已經封閉,甚至埋藏了的東門,彷彿把那年那個歷史時刻永久地封閉了,猶太人曾經稱呼耶穌是以色列王的歷史事件,也彷彿被埋葬在已經封閉的記憶中.

我當天與同學分享,耶穌這位天國君王,面向聖殿山,在今天主哭耶京堂的位置,為耶路撒冷哭泣；然後在今天萬國教會相傳是客西馬尼園的地點,汗如血滴禱告,在凌晨時分,等候那些捉拿祂的人.聖殿山和東門成為這位天國君王在聖週/受難週的一個重要焦點:祂幾乎每天經東門往返橄欖山和聖殿山；祂在這神聖的山頭上潔淨聖殿,在聖殿公開講論天國之道,在迎接死亡之前,祂再三展現作為天國君王的勇氣和風采.

這位以色列君王,以眼淚,血汗和生命,選擇了一條叫當日在東門擁戴他的群眾感到失望的一條道路,他迎向十字架,走上各各他.判刑那天,彼拉多問猶太人:「我可以把你們的王釘十字架嗎?」

一位被鞭打和被羞辱的君王

約翰福音19章1-16節可分為四個小段落,分別記載了耶穌被鞭打和侮辱(1-3節),彼拉多意圖釋放耶穌(4-7節),彼拉多與耶穌的對話(8-11節),彼拉多向猶太人妥協處耶穌以死(12-16節)

使徒約翰以記載彼拉多鞭打耶穌,來展開第一個小段落；接著,記述了兵丁用荊棘編作冠冕,戴在耶穌頭上,又給耶穌穿上紫色的袍.這幾位兵丁看來十分享受侮辱耶穌的過程,他們挨近耶穌說:「恭喜猶太人的王啊.」然後就打耶穌.

兵丁對耶穌的嘲諷,其實是承接著約翰福音18章39節彼拉多向猶太人所說的話:「但你們有個規矩,在逾越節要我給你們釋放一個人,你們要我給你們釋放猶太人的王麼?」也承接著彼拉多審問耶穌時的對話,彼拉多問耶穌:「你是王麼?」耶穌回答說:「你說我是王,我為此而生,也為此來到世間,特為給真理作見證.」(約18:37)

耶穌所受到的五項肉體和心靈的傷害,都帶著一些含義:

被鞭打涉及他罪犯的身份,彼拉多意圖以鞭打作為刑罰,來免去耶穌的死罪；羅馬人的皮鞭上面鑲嵌著骨頭或金屬粒,鞭打人的時候就產生了皮開肉裂的效果.

把冠冕戴在頭上,當然就是要嘲諷「耶穌是猶太人的王」這個講法.荊棘刺穿了耶穌的額頭和頭蓋,耶穌自然血流滿面,頭蓋後部的血更隨隨流到嚴重受傷的背部.冠冕本來是榮耀的象徵,然而,荊棘是上帝咒詛土地的表記(創3:17-18),荊棘編成的冠冕就代表了咒詛,痛苦和羞辱.

兵丁給耶穌穿上紫色的袍,紫色的布料是富有人家所使用的材料,紫色的袍就是貴族和皇室所喜歡穿著的衣服.在穿衣服和脫衣服的過程中,受鞭傷的耶穌都會因背部的傷口,帶來極大的皮肉之痛.也許耶穌短短的一生,就從來沒有穿過紫色的衣服!這時,耶穌從頭至腳,從頭上的荊棘冠冕,加上這件袍子,還有臉上的血和汗,成為一個徹底的被嘲弄的形象.

兵丁就出言嘲弄耶穌恭喜他成為猶太人的王.羅馬兵丁在對待囚犯方面,也許並沒有什麼好名聲,不過,他們並沒有痛恨耶穌的理由,卻如此起勁地戲弄耶穌,叫人搖頭嘆息,甚至憤慨,不禁要問,在縱容使用暴力的時候,人性是否往往變得更加醜陋?

然後,兵丁又打耶穌,除了約翰福音記載以手掌直接擊打,符類福音更記載兵丁用蘆葦做成的藤條,打耶穌的頭.你可以想像嗎?以藤條打一個皇帝的臉,其實是對那個皇帝何等大的羞辱?

我們看這個人!

彼拉多以為鞭打和侮辱耶穌,可以使猶太人涉去部份憤恨,換來一個釋放耶穌的下台階.他向猶太人說:「我帶他出來見你們,叫你們知道我查不出他有甚麼罪來.」彼拉多做足功課,他從總督府屋內,帶耶穌出來,猶太人看見的是一個戴著荊棘冠冕,穿著紫袍,滿臉傷痕的曾經被受民眾擁帶,呼喊「和撒那」的那位.彼拉多對他們說:「你們看這個人!」 (約19:4-5)

當下,跟從大祭司的群眾可能一時未能反應過來,約翰告訴我們,是該亞法與聖殿差役首先起哄喊叫說「釘他十字架!釘他十字架!」彼拉多就第二次宣告「我查不出他有甚麼罪來.」 彼拉多的意思,就是按照羅馬法律,耶穌並沒有招致死刑的罪名.然而,要置耶穌於死地的猶太人卻「搬龍門」,指出按照猶太律法,耶穌以自己為神的兒子是該死的.約翰告訴讀者,既然「猶太人的王」這控訴不能入罪,就改為指控「自稱為神的兒子」,反正結論一早就定好了,猶太人毫不介意用什麼藉口來達至這個結論,無論是猶太公會,或者是總督府的審訊,程序正義已經不再重要,最重要的是要達成既定的結論.

弟兄姊妹,若果你身處現場,看見這位被人從頭到腳惡意羞辱的君王的容貌和身影,你又有何感應?你有沒有想過, 二千年前的那一天,主耶穌所經歷的羞辱和不義的對待,和你有什麼關係?聖壇上的十字架,象徵耶穌為我們所承擔的羞辱,咒詛和痛苦,今天,我們又怎樣看這個人?

權柄是誰的

約翰告訴我們,手握權力的彼拉多,心理負擔越來越重,而且由害怕變得越發害怕.手握權力的人,很多時表現出來,都是一副胸有成竹的模樣；約翰卻捕捉了這個鏡頭,按他的細微觀察,彼拉多其實一直都處於一種害怕的心理狀態,當他聽到耶穌自稱是神的兒子這個指控,他可能越發覺得案中有案,稍一不慎,就會做錯決定,影響他自己的政治前途.

彼拉多簡單直接的問耶穌:「你是從哪裡來的?」從彼拉多這個問題,我們知道他對這一宗看來是猶太律法的爭議,或者宗教糾紛,或者出於猶太高層的嫉妒和仇恨動機,他一直想置之不理,把耶穌退回給該亞法和猶太人,如今被祭司長,祭司長的隨從和洶湧的猶太人,逼他行使公權力,他必須弄清楚耶穌的來頭,好算計風險,萬一殺害了耶穌,卻又惹起另一些權貴和百姓的抗議,倒不是他想要得到的結果.

所以,當耶穌不回答,彼拉多就說:「你不對我說話嗎?你豈不知我有權柄釋放你,也有權柄把你釘十字架嗎?」 (約19:10)耶穌回答直接了當:「若不是從上頭賜給你的,你就毫無權柄辦我.所以,把我交給你的那人罪更重了.」 (約19:11) 我們還記得耶穌說過:「沒有人奪我的命去,是我自己捨的.我有權柄捨去,也有權柄取回來.」(約10:8) 彼拉多作為羅馬皇帝賦權的地方官員,深以為自己有權力去作出判刑和政治決定；然而,愛我們的主是超越地上君王和一切掌權的,正因為祂的國度超越現世的國度,祂的生命超越地上掌權的人,或者簡單點說,當基督走在十字架的路上,捨身忘死,地上的權柄怎可能轄制祂,在真理中,得著真正的自由.

彼拉多的問號

從約翰福音第18章第28節開始,彼拉多出現在耶穌受審的過程中,而且還扮演著非常重要的角色.彼此多在整個事件中發出了很多個問號:

你們告這個人是為什麼是呢?(約18:29)

你是猶太人的王麼?(約18:33)

我豈是猶太人呢?(約18:35)

你作了什麼事呢?(約18:35)

這樣,你是王麼?(約18:37)

真理是什麼呢?(約18:38)

你們要我給你們釋放猶太人的王麼?(約18:39)

你是那裡來的?(約19:9)

你不對我說話麼?(約19:10)

你豈不知我有權柄釋放你,也有權柄把你釘十字架麼?(約19:10)

我可以把你們的王釘十字架嗎?(約19:15)

這十一個問題當中,其實都環繞著「耶穌是以色列的王」這個中心,當中總共有四條問題直接涉及耶穌是猶太人的王這個講法.其中最決定性的一個問號,也就是最後一個問號.彼拉多心想要釋放耶穌,無奈抵擋不住猶太人的壓力:「你若釋放這個人,就不是凱撒的中神,凡以自己為王的,就是背叛了凱撒.」

彼拉多但求穩住局面,尤其不想引起羅馬政府對自己的負面評價,當他看見自己先後兩次宣告查不出耶穌有什麼罪來,又意圖以巴拉巴作為最後一著棋子,他仍然無法說服該亞法和猶太人,在他們喊叫「除掉他!除掉他!」的聲浪中,彼拉多就在判刑的鋪華石處,發出了最後一個問號,然後,就將耶穌判以十字架的刑罰.

主耶穌,願你為王!

祭司長該亞法說出了叫一切真誠遵守摩西律法和敬拜上主的猶太教徒心碎的一句話「除了凱撒,我們沒有王.」

按照舊約先知撒母耳的傳統,耶和華上主才是以色列人的王,當年上主對撒母耳說:「百姓向你說的一切話,你只管依從,因為他們不是厭棄你,乃是厭棄我,不要我作他們的王.」(撒上8:10) 如今,竟然是出於大祭司的口,無論在猶太人的社會民情中,或者是宗教感情上,都難以說出「除了凱撒,我們沒有王」這句話.

按照約翰的記載,第一個稱呼耶穌是以色列王的門徒是拿但業,當耶穌與拿但業首次相遇的時候,拿但業說:「拉比,你是上帝的兒子,你是以色列的王!」(約1:49 )來自加利利的城市迦拿的拿但業,符類福音稱呼他為巴多羅買,他是約翰筆下最早出現於約翰福音的門徒之一,但在符類福音,卻沒有記載拿但業任何事蹟.這位被主耶穌讚許是「真以色列人」的門徒,啟導我去思想,「基督君王」帶給我們的信仰含義,並且在生活中的實踐意義.

教會年曆設有「基督君王日」,標誌著教會年曆一個年度的終結.教會年曆由聖誕前的將臨期開始,讓信徒預備迎接道成肉身的救主,然後進入聖誕期和之後的顯現期,基督將祂的神性顯現給世人知道,再進入預苦期後,我們紀念主耶穌十架捨身,復慶祝祂從死裡復活,然後是一段為期二十多個星期的聖靈降臨期,信徒透過研讀基督的事跡和教導,學習如何實踐作主的門徒,直到年曆的最後一個主日,即今年11月24日,我們慶祝主基督為君王,統管萬有.

基督君王日,引發我們再一次將焦點集中在耶穌基督這位和平之君,經歷將臨期,聖誕期,顯現期,預苦期,復活期及聖靈降臨期之後,重申基督的王權,祂不單是我們每天每夜的生活之主,祂更是那位帶著榮耀再次回來的君王,祂要回來審判萬民,因為基督的王權,超過世上一切國度和掌權者.

因為基督君王的信念,使我們在艱難和困苦之中,在人生的風暴裡,在看不見盡頭的黑夜中,我們仍然可以懷著盼望,腳踏實地的遠望永恆的天家.

大祭司不願意亦無法應對耶穌所提出的宗教更新議程,「若這樣由著他,人人都要信他.」(約11:47) 祭司長和公會領袖既然無法處理耶穌提出的問題,就乾脆把提出問題的耶穌除掉,所以當彼拉多問他們說「我可以把你們的王釘十字架麼?」,該亞法就狠狠地回答了那句羞恥的謊話.

弟兄姊妹,我們所跟隨的這位基督君王,祂其實是因為我們生命中的羞恥和過犯,為我們承擔罪價,死在十字架上.當我們在生活中感到迷失,無力,沮喪,生活的大環境充滿壓力和幽暗,甚至覺得自己與上帝的關係疏遠,我們只需要回到主耶穌的十字架那裡,學拿但業一樣確認祂是神的兒子和我們人生的君王,從基督的角度來重新觀看自己的生活和人生,並且靠著真理的聖靈,放下叫我們失去自由的重擔,我們定必重新得力.

今天,在這個聖堂內,我們都看見拿但業,又即是巴多羅買的徽號,徽號中的匕首,記載著拿但業的殉道的一頁歷史.耶穌讚許他是心中沒有詭詐的人,這位真心跟隨耶穌的使徒,始終都走在十字架的路上,與他同時代的,那一群在耶路撒冷東門擁戴耶穌的猶太人,甚至在總督府呼喊「除掉他!釘他十字架」的猶太人,是何等強烈的,鮮明的對比.求主憐憫我們,在艱難的歲月中,我們仍然堅定以基督為我們自己,我們的家庭,我們的教會,和我們的社會的君王.



\newpage

\section{耶穌係呢度}
\label{sec:466}
\textbf{林豪恩先生 - 「耶穌係呢度!」}
\newline
\newline
\hyperref[sec:465]{< < < PREV SERMON < < <}
~
\hyperref[sec:toc]{[返目錄]}
~
\hyperref[sec:467]{> > > NEXT SERMON > > >}


路加福音 12:49-56
\newline
\begin{longtable}{cl}
\hline
\hline
章節 & 經文 (和合本修訂版)\\
\hline
12:49 & \begin{tabularx}{0.7\textwidth}{X} 「我來是要把火丟在地上，假如已經燒起來，不也是我所希望的嗎？ \end{tabularx} \\ \\ \relax
12:50 & \begin{tabularx}{0.7\textwidth}{X} 我有當受的洗還沒有受，在這事完成之前，我是多麼地焦急！ \end{tabularx} \\ \\ \relax
12:51 & \begin{tabularx}{0.7\textwidth}{X} 你們以為我來是要使地上太平嗎？不！我告訴你們，是使人紛爭。 \end{tabularx} \\ \\ \relax
12:52 & \begin{tabularx}{0.7\textwidth}{X} 從今以後，一家五個人將要紛爭，三個和兩個相爭，兩個和三個相爭： \end{tabularx} \\ \\ \relax
12:53 & \begin{tabularx}{0.7\textwidth}{X} 父親和兒子相爭，兒子和父親相爭；母親和女兒相爭，女兒和母親相爭；婆婆和媳婦相爭，媳婦和婆婆相爭。」 \end{tabularx} \\ \\ \relax
12:54 & \begin{tabularx}{0.7\textwidth}{X} 耶穌又對眾人說：「你們看見西邊起了雲彩，就說：『要下大雨了』，果然就有； \end{tabularx} \\ \\ \relax
12:55 & \begin{tabularx}{0.7\textwidth}{X} 起了南風，你們就說：『要燥熱了』，也就有了。 \end{tabularx} \\ \\ \relax
12:56 & \begin{tabularx}{0.7\textwidth}{X} 假冒為善的人哪，你們知道分辨天地的氣象，怎麼不知道分辨這是甚麼時代呢？」 \end{tabularx} \\ \\
[1ex]
\hline
\hline
\end{longtable}
講題:「耶穌係呢度!」 Jesus Is Already Here!

講員:林豪恩先生

經文:路加福音12:49-56

場地:香港中文大學崇基學院神學院容啟東校長紀念樓G/F聖堂

日期:2019年8月18日

引言

今年六月,有市民提出司法覆核,申請書指出,警員向神職人員說「叫你個耶穌落黎見我地」,明顯屬荒謬及不合邏輯,因耶穌屬於「非自然人」,所有教友對此感到冒犯及難受,故要求法庭按《道歉條例》,命令警務處長向公眾道歉,以及宣告涉事警員的行為越權.法官周家明引述案例指出,若要納入司法覆核管轄中,有關決定須影響決策者以外的人,即改變他們司法上的權利或義務,或剝奪其某種利益.司法覆核亦不能解決所有糾紛,只能涵蓋下一級法庭,審裁機構,行使公權的機構,並在其決定屬非法,越權,全不合理或沒有符合法定程序時,才能干預.周官認為,這宗申請超越司法覆核恰當範圍,亦非可作合理爭辯,決定駁回申請.(2019年6月26日明報)

今年七月,有一部日本電影,英文名叫Jesus,在香港上影譯名為《耶穌真係落jor 黎》,相信是配合近期社會事件而欲收宣傳之效.其實,該電影的日文名的意思是「我討厭耶穌」.故事講述小男孩一家搬到鄉郊生活,被送入基督教學校讀書,正感徬徨無助之際,迷你版耶穌卻突然顯現眼前!這個小耶穌會打相撲又會變法術,對男孩可謂有求必應!但此時男孩的朋友卻遭逢意外,令他對神的信任開始動搖.究竟那個小耶穌是真耶穌顯靈,抑或只是主角幻想出來的玩伴,最後都沒有定論.

「叫你個耶穌落黎見我地」!假如耶穌出現,召喚的人能確認眼前所見的就是耶穌嗎?另一方面,眼前沒有一個具體的「耶穌」出現,就可以肯定耶穌不在嗎?人們又如何能夠得知耶穌在還是不在呢?讓我們透過本主日經課的四段經文思考這個問題,希望能獲得一些啟發.

經文

不憑眼見,而是分辨；不是時間,而是時機

路加福音

12:49「我來要把火丟在地上,倘若已經著起來,不也是我所願意的嗎?

12:50我有當受的洗還沒有成就,我是何等地迫切呢?

12:51你們以為我來,是叫地上太平嗎?我告訴你們,不是,乃是叫人紛爭.

12:52從今以後,一家五個人將要紛爭:三個人和兩個人相爭,兩個人和三個人相爭；

12:53父親和兒子相爭,兒子和父親相爭；母親和女兒相爭,女兒和母親相爭；婆婆和媳婦相爭,媳婦和婆婆相爭.」

12:54耶穌又對眾人說:「你們看見西邊起了雲彩,就說:『要下一陣雨』；果然就有.

12:55起了南風,就說:『將要燥熱』；也就有了.

12:56假冒為善的人哪,你們知道分辨天地的氣色,怎麼不知道分辨這時候呢?」

在對號入座之前,請大家留意耶穌這段說話是向誰說的?「12:56假冒為善的人哪,你們知道分辨天地的氣色,怎麼不知道分辨這時候呢?」在福音書被耶穌稱得上「假冒為善的人」通常是文士和法利賽人,就是被耶穌指責扭曲律法的宗教和民間領袖.作為宗教和民間領袖,文士和法利賽人致力保護他們建立的宗教詮釋系統不被挑戰,祭司和撒都該人努力守衛管轄地的穩定,防止亂事出現,免得被羅馬帝國有理由出兵來平亂而失去僅有的自治權.耶穌說:「12:51你們以為我來,是叫地上太平嗎?」在那種的狀態之下,文士和法利賽人,祭司和撒都該人耶穌的出現能夠幫助社會穩定,這是可以理解的.或者,他們對耶穌這種期望,也是歷世歷代當權者對宗教的要求.在這種期望之下,人們自然認為如果「耶穌落左黎既話」,社會一定和諧；如果仍然有紛爭,即是等於「耶穌無落到黎」.當時,耶穌對文士和法利賽人說:「我告訴你們,不是,乃是叫人紛爭.」文士和法利賽人一定認為那個耶穌好有可疑,甚至可以肯定眼前的那個耶穌一定不是上帝的兒子.

耶穌曾對祂的跟隨者說:「我心裡柔和謙卑,你們當負我的軛,學我的樣式」(太十一29),可見耶穌並不是教門徒惹是生非.所以,不少跟隨耶穌的人都願意為和平而努力,甚至願意配合社會創造和諧的環境.有一本書名為《基督教與和諧社會建設》,是2013年中國社會科學院基督教研究中心與北京燕京神學院聯合舉辦的國際合作論壇的講稿結集,探討基督教在構建和諧社會中可能的貢獻.然而,事情往往並非一相情願就能夠成就的,當處境太過極端的時候,「和諧頻道」也被逼發出被認為「不和諧」的聲音.

在那個晚上,耶穌獨自跪下禱告,(路加福音22:42說:「父啊!你若願意,求你將這杯撤去；然而,不是照我的意願,而是要成全你的旨意.」如果可以開懷暢快地喝甜酒,沒有人想喝苦杯.然而,耶穌「12:50我有當受的洗還沒有成就」.耶穌的杯,耶穌的洗,所指的是十字架的道路.耶穌被釘上十字架犧牲,對於認罪悔改的人是拯救,對於繼續行惡的人是審判,12:49所說「我來要把火丟在地上」的「火」就是審判的火.與耶穌同釘十字架的左右兩個人,就展現出既是拯救,又是審判的狀況,也顯示出彼此紛爭的狀況:「23:39同釘的犯人中有一個譏笑他,說:「你不是基督嗎?救救你自己和我們吧!」23:40另一個就應聲責備他,說:「你是一樣受刑的,還不怕上帝嗎?23:41我們是應得的,因為我們是自作自受,但這個人沒有做過一件不對的事.」23:42他對耶穌說:「耶穌啊,你進入你國的時候,求你記念我.」23:43耶穌對他說:「我實在告訴你,今日你要同我在樂園裡了.」」耶穌在哪裡呢?在十字架上,有人否認是耶穌,有人認信了這是耶穌.他們的分別在於一個仍然知錯不改,一個願意謙卑認錯.

施洗約翰曾經呼籲:「你們要悔改,因為天國近了」.耶穌帶來的天國,就是那個來臨中的世代,無可避免地與那個將要過去的世代產生衝突.耶穌說:「12:51你們以為我來,是叫地上太平嗎?我告訴你們,不是,乃是叫人紛爭.」不知道大家有沒有唱過一首叫「萬世戰爭」的詩歌呢?有人笑說這首詩歌適合在婚禮上唱,因為婚禮為「萬世戰爭」揭開序幕.「萬世戰爭」原名叫「The Conflict of the Ages」,the ages是指兩個世代,「萬世戰爭」所講的並不是時間,而是來臨中的世代與將要過去的世代之間的衝突.

12:54耶穌又對眾人說:「你們看見西邊起了雲彩,就說:『要下一陣雨』；果然就有.12:55起了南風,就說:『將要燥熱』；也就有了.12:56假冒為善的人哪,你們知道分辨天地的氣色,怎麼不知道分辨這時候呢?」「這個時候」,原文是kairos,是聖經中兩個與時間有關的詞彙的其中一個,另一個是(chronos).(chronos)是「時序性」的時間,(kairos)是「時機性」的時間.「時序」時間屬於時鐘時間,它是客觀時間,強調先後之別.「時機」時間關乎那特別時刻,要求人回應和行動.它是主觀時間,因為各人有其特別時刻.從古至今都有人預言耶穌再來的日期,有一位叫肯頓畢肖爾,他預告耶穌再臨會落在2018至2028年之間,其理據是因為1948年以色列建國再經過一個聖經世代(70至80年),即是2018至2028年之間.這樣的計算都是以「時序」時間(chronos)來推測耶穌的出現日期,然而,耶穌說:「怎麼不知道分辨這時候呢?」那是關乎特別的時刻,要求人分辨和回應.因此,在同一個「時序」時間,在同一個地方,有人高呼:「叫你個耶穌落黎見我地」；有人高唱「Sing Hallelujah to the Lord」.

不憑身分,而辨內容；不靠安撫,而分善惡

23:23耶和華說:「我豈為近處的上帝呢?不也為遠處的上帝嗎?」

23:24耶和華說:「人豈能在隱密處藏身,使我看不見他呢?」耶和華說:「我豈不充滿天地嗎?

23:25我已聽見那些先知所說的,就是託我名說的假預言,他們說:『我做了夢!我做了夢!』

23:26說假預言的先知,就是預言本心詭詐的先知,他們這樣存心要到幾時呢?

23:27他們各人將所做的夢對鄰舍述說,想要使我的百姓忘記我的名,正如他們列祖因巴力忘記我的名一樣.

23:28得夢的先知可以述說那夢；得我話的人可以誠實講說我的話.糠秕怎能與麥子比較呢?這是耶和華說的.」

23:29耶和華說:「我的話豈不像火,又像能打碎磐石的大錘嗎?」

如果耶利米先知被要求「叫你個耶穌落黎見我地」,他可能會如此回應:23:23耶和華說:「我豈為近處的上帝呢?不也為遠處的上帝嗎?」23:24耶和華說:「人豈能在隱密處藏身,使我看不見他呢?」耶和華說:「我豈不充滿天地嗎?」耶穌在不在呢?問題可能只是我們的眼睛被蒙蔽而看不到他,反過來說,其實他正在看著我們.如果意識到上帝正在看著我們,我們準備好「佢落黎見我地」了嗎?還是謙虛一點,並不是「叫佢落黎見我地」,而是洗心革面去朝見祂呢?

在耶利米的時代,縱使上主「充滿全地」,但人民仍看不見祂,為甚麼呢?因為有些本來應該帶引人去上主面前的領袖,反而蒙蔽人民的眼睛.這些領袖如何蒙蔽人民的眼睛呢?「23:25我已聽見那些先知所說的,就是託我名說的假預言,他們說:『我做了夢!我做了夢!』23:26說假預言的先知,就是預言本心詭詐的先知,他們這樣存心要到幾時呢?23:27他們各人將所做的夢對鄰舍述說,想要使我的百姓忘記我的名,正如他們列祖因巴力忘記我的名一樣.」

這些先知的假信息是怎樣的呢?聽眾如何分辨真假的信息呢?觀其行,聽其言:第一,假先知的行為「23:14我在耶路撒冷的先知中曾見可憎惡的事；他們行姦淫,做事虛妄,又堅固惡人的手,甚至無人回頭離開他的惡.」第二,假先知的宣講:「23:17他們常對藐視我的人說:『耶和華說:你們必享平安』；又對一切按自己頑梗之心而行的人說:『必沒有災禍臨到你們.』」正如耶穌指責祭司和法利賽人「假冒為善」一樣,耶利米年代的先知也是「假冒為善」.先知本應代表善,但他們與惡同行,卻仍然保留善的名號；先知代表上帝發言,呼籲悔改,但他們卻為惡人開脫,為惡人護航,卻仍然自稱先知.耶利米建議人們去分辨,不要只靠身分,而要對內容慎思明辨；不要只信安慰的話,而要分辨善惡.

不憑權力,而憑公義；不靠高貴,而護弱勢

82:1上帝站在有權力者的會中,在諸神中行審判,

82:2說:你們審判不秉公義,徇惡人的情面,要到幾時呢?(細拉)

82:3你們當為貧寒的人和孤兒伸冤；當為困苦和窮乏的人施行公義.

82:4當保護貧寒和窮乏的人,救他們脫離惡人的手.

82:5你們仍不知道,也不明白,在黑暗中走來走去；地的根基都搖動了.

82:6我曾說:你們是神,都是至高者的兒子.

82:7然而,你們要死,與世人一樣,要仆倒,像王子中的一位.

82:8上帝啊,求你起來審判世界,因為你要得萬邦為業.

那些戲言「叫你個耶穌落黎見我地」的人,其實知不知道如果耶穌果真出現在他們面前會是怎樣的一回事?「82:1上帝站在有權力者的會中,在諸神中行審判,82:2說:你們審判不秉公義,徇惡人的情面,要到幾時呢?82:7然而,你們要死,與世人一樣,要仆倒,像王子中的一位.」他們其實應該要求「耶穌唔好咁快落黎見我地」,讓他們有機會「82:3你們當為貧寒的人和孤兒伸冤；當為困苦和窮乏的人施行公義.82:4當保護貧寒和窮乏的人,救他們脫離惡人的手.」耶穌再次「落黎見我地」,是會審判那些欺善怕惡的在地上審判者,保護那些被惡人欺壓的貧寒和窮乏的人.

從世界不配有的人中看見耶穌

最後,希伯來書提示我們,有時人們從偉大人物領略到上主的臨在,然而,有時在被世界唾棄的人也可以是上主的臨在,希伯來書沒有稱這些人是「被世界唾棄的人」,而是稱他們為「世界不配有的人」.無論是偉人,平凡人,或者「世界不配有的人」,他們的共同地方是甚麼呢?讓我們一齊從這段經文中發掘.

11:29他們因著信,過紅海如行乾地；埃及人試著要過去,就被吞滅了.

11:30以色列人因著信,圍繞耶利哥城七日,城牆就倒塌了.

11:31妓女喇合因著信,曾和和平平地接待探子,就不與那些不順從的人一同滅亡.

11:32我又何必再說呢?若要一一細說,基甸,巴拉,參孫,耶弗他,大衛,撒母耳,和眾先知的事,時候就不夠了.

11:33他們因著信,制伏了敵國,行了公義,得了應許,堵了獅子的口,

11:34滅了烈火的猛勢,脫了刀劍的鋒刃；軟弱變為剛強,爭戰顯出勇敢,打退外邦的全軍.

11:35有婦人得自己的死人復活.又有人忍受嚴刑,不肯苟且得釋放,為要得著更美的復活.

11:36又有人忍受戲弄,鞭打,捆鎖,監禁,各等的磨煉,

11:37被石頭打死,被鋸鋸死,受試探,被刀殺,披著綿羊,山羊的皮各處奔跑,受窮乏,患難,苦害,

11:38在曠野,山嶺,山洞,地穴,飄流無定,本是世界不配有的人.

11:39這些人都是因信得了美好的證據,卻仍未得著所應許的；

11:40因為上帝給我們預備了更美的事,叫他們若不與我們同得,就不能完全.

願最後這兩節,成為我們在這些日子的鼓勵.

12:1我們既有這許多的見證人,如同雲彩圍著我們,就當放下各樣的重擔,脫去容易纏累我們的罪,存心忍耐,奔那擺在我們前頭的路程,

12:2仰望為我們信心創始成終的耶穌.他因那擺在前面的喜樂,就輕看羞辱,忍受了十字架的苦難,便坐在上帝寶座的右邊.



\newpage

\section{合理解綁帶來神的榮耀}
\label{sec:467}
\textbf{區祥江博士 - 合理解綁帶來神的榮耀}
\newline
\newline
\hyperref[sec:466]{< < < PREV SERMON < < <}
~
\hyperref[sec:toc]{[返目錄]}
~
\hyperref[sec:468]{> > > NEXT SERMON > > >}


路加福音 4:16-19
\newline
\begin{longtable}{cl}
\hline
\hline
章節 & 經文 (和合本修訂版)\\
\hline
4:16 & \begin{tabularx}{0.7\textwidth}{X} 耶穌來到拿撒勒，就是他長大的地方。在安息日，照他素常的規矩進了會堂，站起來要念聖經。 \end{tabularx} \\ \\ \relax
4:17 & \begin{tabularx}{0.7\textwidth}{X} 有人把以賽亞先知的書交給他，他就打開，找到一處寫著： \end{tabularx} \\ \\ \relax
4:18 & \begin{tabularx}{0.7\textwidth}{X} 「主的靈在我身上，因為他用膏膏我，叫我傳福音給貧窮的人；差遣我宣告：被擄的得釋放，失明的得看見，受壓迫的得自由， \end{tabularx} \\ \\ \relax
4:19 & \begin{tabularx}{0.7\textwidth}{X} 宣告神悅納人的禧年。」 \end{tabularx} \\ \\
[1ex]
\hline
\hline
\end{longtable}
講題:合理解綁帶來神的榮耀 Rightful Untying Brings God’s Glory

講員:區祥江博士

經文:路加福音4:16-19

場合:香港中文大學崇基學院禮拜堂主日崇拜

日期:2019年8月25日

跟據記載,耶穌可能是最後一次入會堂教訓人,因為猶太人的領袖愈來愈對耶穌有敵意.今天經文的故事就是一個敵意挑戰耶穌在安息日工作的描述.我們先重溫耶穌第一次進會堂所說的話,是記載在路加福音 4:16-19:

「16耶穌來到拿撒勒,就是他長大的地方.在安息日,照他平常的規矩進了會堂,站起來要念聖經. 17有人把先知以賽亞的書交給他,他就打開,找到一處寫著說:

18主的靈在我身上,

因為他用膏膏我,

叫我傳福音給貧窮的人；

差遣我報告:

被擄的得釋放,

瞎眼的得看見,

叫那受壓制的得自由,

19報告神悅納人的禧年.

這段經文可說是耶穌出來傳道的使命宣言.」

我們轉回今天我們的經課,這是講述耶穌在安息日治好駝背的女人的故事.是一個得釋放,得醫治和得自由的故事,正正是耶穌使命宣言的具體實現.

我們先對這故事來一些觀察.

時間是安息日,

地點是會堂.

人物有四,包括:

被鬼附著,病了十八年,腰彎得一點直不起來的女人,大家知道路加是醫生,所以他的描述是十分生動的.這女人的病最有可能是英文為ankylosing spondylitis,國際簡稱為AS,是一種長期脊椎關節有發炎症狀的關節炎,通常脊椎連接骨盆的關節會受到影響,偶爾還會侵犯其他部位的關節,如肩膀或臀部,眼睛和腸道問題也可能發生.症狀:背痛,關節僵硬.

耶穌,當然是帶來醫治的主角.

管會堂的領袖,耶穌出來傳道到在這個階段,宗教領袖們一直留意著耶穌的一舉一動,等待一有機會便要抓住他的錯處,

最後是一班會眾,聽管會堂訓話的,亦是最後因耶穌的反搏後歡喜的會眾,為數不詳.

人物的行動:

耶穌:

12耶穌看見,便叫過她來,對她說:「女人,你脫離這病了!」 13於是用兩隻手按著她；

和修本改了:「婦人,你的病好了!」

其實和合本更有意思:「女人,你脫離這病了!」

大家留意到是耶穌看見就憐憫這女人,主動叫她過來,不是女人要求耶穌醫治她的.耶穌用:「脫離」是語帶相關的.英文是:「untie」,與之後解開槽上的牛,驢,解開撒旦的捆綁有關的.

女人:她立刻直起腰來,就歸榮耀與神.

這婦人雖然有這樣的病症,仍經常到會堂去禮拜,這可說是她信心的表現,所以神賜恩給她.可見先顧靈魂的人,往往也得肉體的益處.這女人真心是亞伯拉罕的後裔,因為她一痊愈,就自然的將榮耀歸給神.

管會堂的人

跟著管會堂的便出面干涉,但他卻沒有直接面對耶穌的勇氣,他只是對那些群眾說話,雖然這些話是擺明了衝著耶穌而來的.耶穌為什麼非要在安息日治病不可? 說話時是帶著憤怒的:

14管會堂的因為耶穌在安息日治病,就氣忿忿地對眾人說:「有六日應當做工；那六日之內可以來求醫,在安息日卻不可.」

在他看來,安息日的「儀式」比一切更重要.他說這件事不合「程序」,破壞了安息日的儀式條文.對他來 說,那婦人是否得醫治是無關緊要的.他根本不考慮這婦人得醫治的事,也不管撒旦的邪惡勢力被遏制的意義.他只想到在安息日作任何形式的工都是不可的.

他的說話表面是合理的.而事實上耶穌並不是來要廢除律法,當時猶太人都有不少在安息日限制種類工作如何合法地得到容許的討論.

耶穌卻看到他假冒為善的一面.他們容許自己在安息日解開槽上的牛,驢,牽去飲水.對自己有容許(allowance),對別人卻強加不合理的要求.意思是說,「如果我伸手解開這婦人的捆綁,是在作工,那麼你解開槽上的牛驢的捆綁,牽去飲水,那就不是在工作嗎?」

這裡耶穌用了拉比慣用的辯證方法.From lesser to greater.

人不比動物更加寶貴嗎?屬靈的後裔不比一般人重要嗎?被撒旦捆綁得釋放不比讓牲畜解開捆綁去飲水更重要嗎?

耶穌更說:16何況她本是亞伯拉罕的後裔,被撒旦捆綁了十八年,不該在安息日這天解開她的綁嗎?」

耶穌對何為應該或不當來了一個新的詮釋.

對耶穌來說,人比條文律法重要得多了.

在屬靈層面上,這女人本是亞伯拉罕的後裔,被撒旦捆綁了這十八年了.

或者這管會堂的可以反搏:「這女人都病了18年,為何不可以等多一日,不在安息日去醫好她」

耶穌行動所反映出來的信息是:「沒有藉口延長痛苦和推遲一個人能整全的過活,不是還有一天,而是我們等不及了.」

在網上看到有一篇對這段經文的靈修分享,網頁名稱:跟隨耶穌,每日靈修.很有意思:

「耶穌這番話值得我們深思:我們常常停留在「習以為常,順理成章」的事情上,因為太習慣了,即使漸漸遠離了神的心意也不自知,就像管會堂的根本不覺得自己有那裡不對一樣.追根究底,我們安於待在「舒適區 Comfort Zone」裡,對一項工作駕輕就熟之後,自然不願輕易作出改變.但若不時常察驗,很容易就陷在所謂「傳統和儀文」的框架之中,空有敬虔的外表,卻無敬虔的實意,導致事重於人,形式重於愛心.」

\url{https://www.hoc5.info/letsfollowjesus/main/daily/0820.html}

最後我們看故事的結局:

17耶穌說這話,他的敵人都慚愧了；眾人因他所行一切榮耀的事,就都歡喜了.

我們看到耶穌言之有理,甚至帶來敵人的慚愧.我們也巴不得自我反思,我們有什麼事情是本末倒置,自以為義的,也當被耶穌挑戰和提醒,慚愧可以是一個很好的反應呢,因為知錯才有動力去改.有人不肯認錯道歉,可以帶來很大的後果.

眾人因他所行一切榮耀的事,就都歡喜了.令我想起以賽亞書 58:13-14有關安息日的教道.

和修版比較易讀:

「13你若禁止自己的腳踐踏安息日,

不在我的聖日做自己高興的事,

稱安息日為「可喜樂的」,

稱耶和華的聖日為「可尊重的」,

尊敬這日,

不走自己的道路,

不求自己的喜悅,

也不隨意說話；

14那麼,你就會以耶和華為樂.

耶和華要使你乘駕於地的高處,

又要以你祖先雅各的產業養育你；

這是耶和華親口說的.」

耶穌安息日治好駝背的女人是「可喜樂的」和「可尊重的」事.

也合乎:「不走自己的道路,不求自己的喜悅,也不隨意說話；」的原則.

在網上看到一篇猶太人有關在安息日應否示威遊行的文章.作者 Laurie Green 是 一位在 Washington, D. C. 的拉比和屬靈領袖.她引用了Abraham Heschel 一位稱為 Professor of Mysticism and Ethics 的猶太神學家的例子,他寫了一本很有名的安息日的小書:「Sabbath」主題是:Shabbat as a Sanctuary in Time,The Sabbaths are our great cathedrals, the Jewish equivalent of sacred architecture.

今天我們不是要談他本書,而是說一個這麼看重安息日的猶太神學家,在1965年3月25日,他一起參予馬丁路德金帶領成千上萬的非暴力示威者前往阿拉巴馬州(Alabama)蒙哥馬利市(Montgomery)的國會大廈,經過5天54英里的阿拉巴馬州塞爾瑪(Selma)遊行.當天Heschel 說:「Today we are praying with our feet」「今天我們用我們的腳祈禱」成為佳話.

這篇網上的文章提了幾個問題,可能都是值得我們香港的基督徒面對政治上的動盪,以個人的名義是否在星期日參予一些「和理非非」的遊行示威呢?

文章作者 ( Laurie Green) 說:在我們的困難和民主社會,幾乎每個週末都有抗議活動. 所以我們必須要問自己:

難道我們自己不值得安息日嗎?

我的思想和精神需要一天遠離壓力的世界?

我的身體是否需要在自然界中進行身體休息,而不是另一天遊行,人群和噪音?

我的家人或社區需要我嗎?

對我來說重要的事情與什麼是真正的社會或政治上緊急情況之間是否存在差異?

我對這些問題的回答是,“是的.”因此,我不會在安息日參予每一次集會或抗議活動. 我保存我對安息日的抗議活動的參予於:

a)大型國家集會,或

b)有時限的緊急行動.

安息日是一個快樂的時刻. 我的希望是,當我們不可避免地被呼喚在安息日“追求公義”時,我們是否可以與我們的家庭和社區一起,在屬靈的心境中這樣做,謙卑地為我們上帝賦予的改變世界的能力而歡欣鼓舞. 阿們.



\newpage

\section{就在這裡把祂釘十字架}
\label{sec:468}
\textbf{劉國偉先生 - 就在這裡,把祂釘十字架!}
\newline
\newline
\hyperref[sec:467]{< < < PREV SERMON < < <}
~
\hyperref[sec:toc]{[返目錄]}
~
\hyperref[sec:469]{> > > NEXT SERMON > > >}


約翰福音 19:17-27
\newline
\begin{longtable}{cl}
\hline
\hline
章節 & 經文 (和合本修訂版)\\
\hline
19:17 & \begin{tabularx}{0.7\textwidth}{X} 耶穌背著自己的十字架出來，到了一個地方，名叫「髑髏地」，希伯來話叫各各他。 \end{tabularx} \\ \\ \relax
19:18 & \begin{tabularx}{0.7\textwidth}{X} 他們就在那裡把他釘在十字架上，還有兩個人和他一同被釘，一邊一個，耶穌在中間。 \end{tabularx} \\ \\ \relax
19:19 & \begin{tabularx}{0.7\textwidth}{X} 彼拉多又寫了一個牌子，釘在十字架上，寫的是：「猶太人的王，拿撒勒人耶穌。」 \end{tabularx} \\ \\ \relax
19:20 & \begin{tabularx}{0.7\textwidth}{X} 有許多猶太人念這牌子，因為耶穌被釘十字架的地方靠近城，而且牌子是用希伯來、羅馬、希臘三種文字寫的。 \end{tabularx} \\ \\ \relax
19:21 & \begin{tabularx}{0.7\textwidth}{X} 猶太人的祭司長就對彼拉多說：「不要寫『猶太人的王』，要寫『那人說：我是猶太人的王』。」 \end{tabularx} \\ \\ \relax
19:22 & \begin{tabularx}{0.7\textwidth}{X} 彼拉多回答：「我寫了就寫了。」 \end{tabularx} \\ \\ \relax
19:23 & \begin{tabularx}{0.7\textwidth}{X} 士兵把耶穌釘在十字架上以後，把他的衣服拿來分為四份，每人一份。他們又拿他的內衣，這件內衣沒有縫，是上下一片織成的。 \end{tabularx} \\ \\ \relax
19:24 & \begin{tabularx}{0.7\textwidth}{X} 他們就彼此說：「我們不要撕開，我們抽籤，看是誰的。」這要應驗經上的話說：「他們分了我的外衣，為我的內衣抽籤。」士兵果然做了這些事。 \end{tabularx} \\ \\ \relax
19:25 & \begin{tabularx}{0.7\textwidth}{X} 站在耶穌十字架旁邊的，有他的母親、姨母、革羅罷的妻子馬利亞，和抹大拉的馬利亞。 \end{tabularx} \\ \\ \relax
19:26 & \begin{tabularx}{0.7\textwidth}{X} 耶穌見母親和他所愛的那門徒站在旁邊，就對母親說：「母親，看，你的兒子！」 \end{tabularx} \\ \\ \relax
19:27 & \begin{tabularx}{0.7\textwidth}{X} 又對那門徒說：「看，你的母親！」從那刻起，那門徒就接她到自己家裡去了。 \end{tabularx} \\ \\
[1ex]
\hline
\hline
\end{longtable}
講題:就在這裡,把祂釘十字架! Here, they nailed Him to the Cross

經文:約翰福音19章17-27節

講員:劉國偉先生

場合:香港中文大學崇基學院禮拜堂主日崇拜

日期:2019年9月8日

引言

去年9月,我第一次在我們的禮拜堂講道,開始了過去一年約翰福音釋經講道的服侍.

我當時沒有想過,在我處理經文和預備講章的過程中,不只一次因為藉著經文進入香港的處境,發覺竟然如斯吻合,而深受感動.我相信這是因為「上帝的道是活潑的,是有功效的,比一切兩刃的劍更快,甚至魂與骨,骨節與骨髓,都能刺入剖開,連心中的思念和主意都能辨明.」(來4:12)

希伯來書的作者更進一步說:「並且被造的,沒有一樣在祂面前不顯然的；原來萬物,在那與我們有關係的主眼前,都是赤露敞開的.」(來4:13)  今天,我以「就在這裡,把祂釘十字架! 」為題,去解讀約翰福音19章17至27節,在香港這個十分艱難的時刻,深願主繼續向我們說話.

使徒約翰在這段經文,記載了耶穌釘十字架的四個小片段.他繼續發揮第四福音的特色,就是為早已面世的三卷符類福音,作出了重要的補充,一方面是一些獨家的記載,即使記載與符類福音相同的片段,約翰往往也會有一些獨特的觀點.從這個角度去閱讀約翰福音19章這四個小片段,我們可以從約翰的筆下,重組耶穌在各各他被釘於十字架的過程中,我們得到的一些重要觀察和信息.這四個小片段分別是三個十字架,三種文字,四個兵丁,四位婦女.

三個十字架

他們就把耶穌帶了去.耶穌背著自己的十字架出來,到了一個地方,名叫「髑髏地」,希伯來話叫各各他. 他們就在那裡釘他在十字架上,還有兩個人和他一同釘著,一邊一個,耶穌在中間. (約19:17-18)

約翰與其他三位福音書的作者一樣,都告訴讀者耶穌釘十字架的地方名叫髑髏地,就是在城外一處地方.還有,耶穌是與另外兩個人一同被釘於那裡,馬太和馬可形容那兩人是強盜,路加則稱他們為犯人.四卷福音書,尤其是馬太福音,有較強的寫作傳統,將耶穌事蹟與應驗舊約經文掛勾,以賽亞書53章12節:「因為祂將命傾倒,以致於死,祂也被列在罪犯之中；祂卻擔當多人的罪,又為罪犯代求.」從客觀的事實來說,那釘在耶穌左右兩旁的人既是強盜,犯人,所以,符類福音作者以強盜或犯人來形容耶穌左右兩人,可以說是意料之內,而耶穌本身被判死刑,又與他們釘在同一個地方,也的確是被列在罪犯之中.

不過,約翰卻沒有採用類似的手法去處理這個畫面,他只是輕輕帶過那釘在耶穌左右兩邊的是兩個人.普普通通一個人,大抵不會認為自己是罪犯,更加不會認為自己是強盜；然而,約翰輕輕說他們是兩個人,而不稱他們為強盜或犯人,就讓讀者可以有更寬闊的聯想:原來你和我都是人,在聖潔的上帝眼中,世人都犯了罪,虧缺了神的榮耀,換句話說,若果我們真的是一個誠實面對自己的人,其實我們也是一個在神的律法底下的罪犯,甚至強盜,耶穌卻因為愛的緣故,甘心捨棄自己的生命,與我們本來要受死的人一起,列在罪犯之中.

路加告訴我們,有一個犯人譏誚耶穌,這是在肉體非常痛苦的十字架受刑者,常有的發洩情緒的表現；但竟然另一個應聲責備他,並且向耶穌說:「耶穌啊!你得國降臨的時候,求你記念我.」(路23:42)耶穌的回應,叫一切在主裡面認罪悔改的人,抱有喜樂的盼望:「我實在告訴你,今日你要同我在樂園裡了.」(路23:43)

使徒約翰告訴我們:「凡接待祂的,就是信祂名的人,祂就賜他們權柄,作神的兒女.」(約1:12)  這是被列在罪犯之中的彌賽亞,所給予全人類充充滿滿的恩典和真理.

三種文字

彼拉多又 用牌子寫了一個名號,安在十字架上,寫的是:「猶太人的王,拿撒勒人耶穌.」有許多猶太人念這名號；因為耶穌被釘十字架的地方與城相近,並且是用希伯來,羅馬,希臘三樣文字寫的.猶太人的祭司長就對彼拉多說:「不要寫『猶太人的王』,要寫『他自己說:我是猶太人的王』.」彼拉多說:「我所寫的,我已經寫上了.」 (約19:19-22)

四卷福音書都有記載在這塊釘在十字架上的木牌,三種文字當中,希伯來人的日常用語亞蘭文,代表了以色列人在宗教文化上所信靠的耶和華一神信仰；羅馬人採用的拉丁文,代表了當時領導世界的權力,和發展強大的整個國家政體；希臘文則代表古代世界在哲學思想的領導地位.所以,這塊木牌本來是耶穌的罪狀,卻成為向世人的重大宣示 : 以色列的王.

我們相信彼拉多當然沒有這份心思,他只是採用了當時最流行的三種語文,為要讓本地的猶太人和路過耶路撒冷的外邦人,都可以知道這個荒謬的事情,就是猶太人竟然讓他們的王釘在十字架上,這也是彼拉多對祭司長和猶太人施壓於他的回應.祭司長當然向彼拉多表示異議,要求彼拉多改為「他自己說我是猶太人的王」,意圖把可能引起的尷尬和羞辱移除,並且徹底與耶穌劃清界線.但彼拉多毫不留情的會回答:「我所寫的,

我已經寫上了.」(約19:22)

在耶穌的十字架之下,人類在玩耍權力遊戲,這是何等大的落差和對比；彼拉多對耶穌和猶太人的嘲諷,對猶太人的羞辱和愚弄,卻成就了一個偉大的宣告:「猶太人的王」.使徒約翰告訴我們:「祂到自己的地方來,自己的人倒不接待祂.」(約1:11)正如拿但業說耶穌是以色列的王,以色列的王來到了,以色列人卻不接待祂,並且要把祂送上十字架.耶穌不單是以色列的王,也是全人類的君王,這個訊息,就是在彼拉多和祭司長互相角力之下,以三種文字被公諸於世.

四個兵丁

兵丁既然將耶穌釘在十字架上,就拿他的衣服分為四分,每兵一分；又拿他的裡衣,這件裡衣原來沒有縫兒,是上下一片織成的. 他們就彼此說:「我們不要撕開,只要拈鬮,看誰得著.」這要應驗經上的話說:「他們分了我的外衣, 為我的裡衣拈鬮.」兵丁果然做了這事. (約19:23-24)

四卷福音書都有提及兵丁以抽籤的方式,去分耶穌的衣服,唯獨使徒約翰非常仔細的補充了一個小故事,就是有關耶穌的裡衣.

按照羅馬人的規定,行刑的兵丁以四人作一組看守罪犯,兵丁可以分了犯人的衣服.釘耶穌的兵丁,各自取得一份後,發現耶穌的裡衣相當特別,約翰告訴讀者,這件裡衣原來沒有縫兒,是上下一片織成的.

這並不表示這件裡衣價值連城,但若果撕為四份,就必然大大降低它的價值.約翰告訴我們,兵丁這麼做,正奇妙地應驗了詩篇22篇18節的預言:「他們分我的外衣,為我的女兒拈鬮.」約翰連這一點細節也不輕易放過,他帶著有一點驚訝的語調,說:「兵

丁果然作了這事.」

約翰作出這一個細節的補充,就帶引我們回到舊約聖經的詩篇22至24篇.按釋經傳統理解,這是三部曲的彌賽亞詩篇.22篇預表的是那位為羊群捨命的好牧人,也是新約作者最常引用的詩篇；23篇表述上主就是照料羊群的大牧人；24篇則表明這位大牧人,上主,在榮耀中顯現,是榮耀的王.

耶穌赤裸裸被釘在十字架上,連裡衣也被兵丁抽籤取去,是對天地的主徹徹底底的羞辱.約翰提醒我們耶穌的教導:好牧人為羊捨命,對照貪圖利益也毫無憐憫的兵丁,這幅在各各他的畫像,實在叫人神傷.也許有人會解說,當兵丁也只是一份工作,而且被差派擔任行刑,更是一份厭惡的差事,羅馬容許兵丁分得這些好處,大家也都習以為常.不過,正正是這種習以為常,就叫我們的世界與真理愈拉愈遠.

四位婦女

站在耶穌十字架旁邊的,有他母親與他母親的姊妹,並革羅罷的妻子馬利亞,和抹大拉的馬利亞.耶穌見母親和他所愛的那門徒站在旁邊,就對他母親說:「母親( - 原文是婦人),看,你的兒子!」又對那門徒說:「看,你的母親!」從此,那門徒就接她到自己家裡去了. (約19:25-27)

約翰記載在十字架的旁邊有四位婦女:第一位是耶穌的母親馬利亞；第二位是耶穌母親的姊妹；第三位是革羅罷的妻子馬利亞；第四位就是抹大拉的馬利亞.四卷福音書的記載,雖然略有出入,但是仔細對照之後,大體上是可以互為補充解釋的.約翰再一次扮演了補充資料的角色,他提到耶穌的母親馬利亞,那是符類福音沒有記載的.約翰更記載了耶穌的兩句說話:「母親!看你的兒子!」和「看你的母親!」

這四位婦女,都與耶穌有緊密的關係,那是人間的真情真愛,目睹他們親愛的耶穌一步一步邁向死亡.有血有肉的耶穌,在死亡之際,身體和心靈正在承受極度痛苦的時候,仍然牽掛人間的親情,表現出對母親的深切關懷,就把母親交托給自己所愛的門徒.

這位道成肉身的彌賽亞,深入人間的愛情當中,與人建立有血有肉的關係,他真正在人間活過.

道成了肉身

約翰福音的序言這樣說:「太初有道,道與神同在,道就是神.這道太初與神同在.生命在他裡頭,這生命就是人的光.那光是真光,照亮一切身在世上的人.」來到約翰福音近尾聲的時候,這位道與真光的神,又會是一番怎樣的光景?

各各他是道成肉身的地理上一個總結,耶穌在那裡完結了祂的人生旅程,結束了祂作為「人子」的一段與人類血肉之軀情愛交纏的經歷.

十字架是道成肉身的「住在我們中間」的最高峰的體現,他不單成為血肉之軀,也不單住在我們當中,更為我們捨身十字架,為要解救我們脫離罪惡,和這個必朽壞的身軀.

今天這段經文,除了為符類福音補充了珍貴的觀點角度,和獨家記載的小片段,約翰讓我們看見這位道成肉身的彌賽亞,帶著四重身份去完成祂的人生歷程和十字架旅途:罪犯,君王,好牧人,人子.

罪犯 : 親愛的弟兄姊妹,我們本來就是罪犯,但因為耶穌為我們成為一個罪犯,在十字架上為我們捨命承擔贖價,我們就得到永生的盼望；

君王 : 這位死在十字架的以色列王,就是我們所景仰的基督君王,我們是忠心服侍祂的臣僕,讓祂統領我們的人生；

好牧人 : 耶穌是我們的好牧人,耶穌教導我們:「我的羊聽我的聲音,我也認識他們,他們也跟著我.我又賜給他們永生；他們永不滅亡,誰也不能從我手裡把他們奪去.」(約10:27-28)我們要敏銳辨識主的聲音,也要緊緊跟隨祂；

人子 : 在福音書的記載中,耶穌喜歡以「人子」自稱.祂確然是馬利亞所生的兒子,祂以這個身份進入人間,是道成了肉身的根本,也是祂可以為人類承擔罪債的根本.事實上,耶穌不單在各各他被釘十字架；對我們每一個人的人生旅程來說,人子耶穌都在我們的生命中,為我們釘十字架!

這位被列為罪犯的人子,祂是我們的好牧人和君王,祂為你我走上十字架的旅途.親愛的弟兄姊妹,我們也有屬於我們自己的十字架旅途.你與我本來就是那些接受耶穌呼召,天天背十字架跟隨主的人,我們不應有一個選項是隔天,或隔周,或隔年去背起我們的十字架.

走到生命晚程的使徒約翰向你和我作見證:「道成了肉身,住在我們中間,充充滿滿的有恩典有真理.我們也見過他的榮光,正是父獨生子的榮光.」

當黑暗的權勢籠罩我們,罪惡越發顯得猖獗,我們背著十字架走窄路,也顯得越來越艱難,但我們知道耶穌今天仍然住在我們中間,他充充滿滿的恩典和真理,讓我們在艱難歲月中,仍然看見祂的榮光,那是詩篇24篇所言的榮耀的王的榮光；於是,我們懷著信心和盼望,學習主耶穌住在人群當中,並且學習那位捨命的好牧人,就如詩篇22篇所記載的擺上和犧牲；更因為我們確信十字架是我們的榮耀,我們天天背起十字架,有血有肉的在黑暗的世代裡,去反映基督的榮光,學效詩篇23篇那位大牧人,去保護群羊.

主必與我們同在!



\newpage

\section{他們如此容易便墮落了}
\label{sec:469}
\textbf{鄧瑞強博士 - 他們如此容易便墮落了}
\newline
\newline
\hyperref[sec:468]{< < < PREV SERMON < < <}
~
\hyperref[sec:toc]{[返目錄]}
~
\hyperref[sec:470]{> > > NEXT SERMON > > >}


出埃及記 32:1-14
\newline
\begin{longtable}{cl}
\hline
\hline
章節 & 經文 (和合本修訂版)\\
\hline
32:1 & \begin{tabularx}{0.7\textwidth}{X} 百姓見摩西遲遲不下山，就聚集到亞倫那裡，對他說：「起來！為我們造神明，在我們前面引路，因為領我們出埃及地的那個摩西，我們不知道他遭遇了甚麼事。」 \end{tabularx} \\ \\ \relax
32:2 & \begin{tabularx}{0.7\textwidth}{X} 亞倫對他們說：「你們去摘下你們妻子、兒女耳上的金環，拿來給我。」 \end{tabularx} \\ \\ \relax
32:3 & \begin{tabularx}{0.7\textwidth}{X} 眾百姓就摘下他們耳上的金環，拿來給亞倫。 \end{tabularx} \\ \\ \relax
32:4 & \begin{tabularx}{0.7\textwidth}{X} 亞倫從他們手裡接過來，用模子塑造它，把它鑄成一頭牛犢。他們就說：「以色列啊，這是領你出埃及地的神明！」 \end{tabularx} \\ \\ \relax
32:5 & \begin{tabularx}{0.7\textwidth}{X} 亞倫看見，就在牛犢面前築壇。亞倫宣告說：「明日要向耶和華守節。」 \end{tabularx} \\ \\ \relax
32:6 & \begin{tabularx}{0.7\textwidth}{X} 次日清早，百姓起來獻燔祭和平安祭，就坐下吃喝，起來玩樂。 \end{tabularx} \\ \\ \relax
32:7 & \begin{tabularx}{0.7\textwidth}{X} 耶和華吩咐摩西：「下去吧，因為你從埃及領上來的百姓已經敗壞了。 \end{tabularx} \\ \\ \relax
32:8 & \begin{tabularx}{0.7\textwidth}{X} 他們這麼快偏離了我所吩咐的道，為自己鑄了一頭牛犢，向它跪拜，向它獻祭，說：『以色列啊，這就是領你出埃及地的神明。』」 \end{tabularx} \\ \\ \relax
32:9 & \begin{tabularx}{0.7\textwidth}{X} 耶和華對摩西說：「我看這百姓，看哪，他們真是硬著頸項的百姓。 \end{tabularx} \\ \\ \relax
32:10 & \begin{tabularx}{0.7\textwidth}{X} 現在，你且由著我，我要向他們發烈怒，滅絕他們，但我要使你成為大國。」 \end{tabularx} \\ \\ \relax
32:11 & \begin{tabularx}{0.7\textwidth}{X} 摩西就懇求耶和華－他的神，說：「耶和華啊，你為甚麼向你的百姓發烈怒呢？這百姓是你用大能大力的手從埃及地領出來的！ \end{tabularx} \\ \\ \relax
32:12 & \begin{tabularx}{0.7\textwidth}{X} 為甚麼讓埃及人說：『他領他們出去，是要降災禍給他們，在山中把他們殺了，將他們從地上除滅』呢？求你回心轉意，不發你的烈怒，不降災禍給你的百姓。 \end{tabularx} \\ \\ \relax
32:13 & \begin{tabularx}{0.7\textwidth}{X} 求你記念你的僕人亞伯拉罕、以撒、以色列。你曾向他們指著自己起誓說：『我必使你們的後裔像天上的星那樣多，並且我要將所應許的這全地賜給你們的後裔，讓他們永遠承受為業。』」 \end{tabularx} \\ \\ \relax
32:14 & \begin{tabularx}{0.7\textwidth}{X} 於是耶和華改變心意，不把所說的災禍降給他的百姓。 \end{tabularx} \\ \\
[1ex]
\hline
\hline
\end{longtable}
講題:他們如此容易便墮落了 They corrupted themselves so easily

經文:出埃及記32章1-14節

講員:鄧瑞強博士

場合:香港中文大學崇基學院禮拜堂主日崇拜

日期:2019年9月15日

各位弟兄姊妹,早安.

中國有一種國技,聽說是不傳外人的,叫「變臉」.揮一揮手,轉一轉身,臉就變了.從天使的臉,一轉身,便能變成惡魔的臉,毫無破綻,令人驚奇,這是中國特有的藝術.

香港,一轉身,面目就變了,只是一瞬間,已變得不再一樣.我現在乘地鐵,老老實實,我真的有點驚.小朋友現在說,很怕見到警察.我教書的,有時會用到鐳射筆,我現在也不敢帶鐳射筆在身上.從來都不是這樣的,只是一瞬間,香港就變了,變得如此陌生.

建造一座城巿,建造人心,建造社會的和諧,需要漫長歲月,但一個大話,再加一個大話,又再加一個大話,一個美好的世界便迅速崩壞了.為了煮一窩好湯,煮了好幾個鐘頭,夠火候了,夠美味了,忽然一隻蒼蠅掉了下去,便前功盡廢.真可惜呀,由好變壞,只是一瞬間的事.

要走上山頂,很緩慢,一步一步的上.但要墮入深谷,則很快,只要一失足,便直墮深淵.在《天路歷程》這本故事書裡,講到:在天堂的門口,便有一條通往地獄的路.千辛萬苦才能走到天堂的門口,但只要走錯一步,便會瞬間直達地獄.人的墮落及一個社會的墮落,比想像中快.

今日,我們看看神子民的一個墮落故事,他們的墮落,比想像中快.

我們看看今日的舊約經課,但為了使故事完整一點,我們多看一點.

經文:出埃及記32:1－14

出32:1　百姓見摩西遲延不下山,就大家聚集到亞倫那裡,對他說:「起來!為我們做神像,可以在我們前面引路；因為領我們出埃及地的那個摩西,我們不知道他遭了甚麼事.」

出32:2　亞倫對他們說:「你們去摘下你們妻子,兒女耳上的金環,拿來給我.」

出32:3　百姓就都摘下他們耳上的金環,拿來給亞倫.

出32:4　亞倫從他們手裡接過來,鑄了一隻牛犢,用雕刻的器具做成.他們就說:「以色列啊,這是領你出埃及地的神.」

出32:5　亞倫看見,就在牛犢面前築壇,且宣告說:「明日要向耶和華守節.」

出32:6　次日清早,百姓起來獻燔祭和平安祭,就坐下吃喝,起來玩耍.

出32:7　耶和華吩咐摩西說:「下去吧,因為你的百姓,就是你從埃及地領出來的,已經敗壞了.

出32:8　他們快快偏離了我所吩咐的道,為自己鑄了一隻牛犢,向它下拜獻祭,說:『以色列啊,這就是領你出埃及地的神.』」

出32:9　耶和華對摩西說:「我看這百姓真是硬著頸項的百姓.

出32:10　你且由著我,我要向他們發烈怒,將他們滅絕,使你的後裔成為大國.」

出32:11　摩西便懇求耶和華─他的　神說:「耶和華啊,你為甚麼向你的百姓發烈怒呢?這百姓是你用大力和大能的手從埃及地領出來的.

出32:12　為甚麼使埃及人議論說『他領他們出去,是要降禍與他們,把他們殺在山中,將他們從地上除滅』?求你轉意,不發你的烈怒,後悔,不降禍與你的百姓.

出32:13　求你記念你的僕人亞伯拉罕,以撒,以色列.你曾指著自己起誓說:『我必使你們的後裔像天上的星那樣多,並且我所應許的這全地,必給你們的後裔,他們要永遠承受為業.』」

出32:14　於是耶和華後悔,不把所說的禍降與他的百姓.

(1)恐懼

三個月前,他們還在埃及,是埃及的奴隸,生活苦不堪言,活在社會的底層.他們的哀聲上達天庭,神聽到了.神派摩西去帶領他們出埃及.奴隸終於解放了,浩浩蕩蕩,離開埃及,往夢想中的家園前進.在紅海前,他們親眼見到,紅海分開,他們走過乾地.後有法老的追兵,他們親眼見到,紅海的水復合,淹過他們.從做奴隸,到得自由,歷過天險,追兵被滅,他們都親身經歷過了.神的真實,神的解救,神的同在,他們都明明白白地體會過了,還有什麼需要懷疑的呢?

到了西乃山,神召摩西上山,向他傳授十誡,教他生活之正道,指導他如何建造親近神的會幕.摩西在山上,群眾在山下.

出32:1　百姓見摩西遲延不下山,就大家聚集到亞倫那裡,對他說:「起來!為我們做神像,可以在我們前面引路；因為領我們出埃及地的那個摩西,我們不知道他遭了甚麼事.」

人在陌生的曠野,四圍都潛伏著可怕的敵人,而領導他們的摩西不見蹤影,群眾害怕了.他們希望做個具體的神像,代替摩西,在他們前面引路.

他們信神嗎?他們信,不過,他們更倚靠一個具體的摩西.具體的摩西不見了,他們信神嗎?他們信,不過,他們更渴望一個具體的神像.難道他們不能相信那位不可見的神嗎?他們信,不過,他們總得用自己的方法,將那位神變得可見,可觸摸,可把握.

本來,是一個美麗的故事,由神主導的一個救贖故事,神主動的解放奴隸,神主動的顯示大能,神主動的與人同行.但是,群眾覺得由神主導的故事自己掌控不了,他們見不到神,見到的,是在陌生之地種種出現的危機.他們渴望能自己主導故事的發展,他們渴望造一個可見的神像,能領導他們.簡言之,他們渴望從神手中奪取事情的主導權,讓一切變得可被操控.

他們信神嗎?他們信,但他們不將命運交給神,因這是不可控制的,他們將命運交給自己製造出來的偶像,因這是可控的.在前行的路上,他們覺得由不可見的神帶領,是充滿危機的.反而,他們覺得,由他們造的神像帶領,是很穩當的.

這是一種很微妙的轉變,由「被神帶領」變成「我製造一個帶領我的神」,由「不可被人掌控的神」變成「一個可被人掌控的神」,由「神是主人」變成「我的需要才是主人」,這種微妙的轉變,實質是將「信仰神」變成「神為我服務」.祈禱,再不是將生命交給神,任由神處理,而是神必需應允我的呼聲,我有這樣那樣的需要.這種信徒會祈禱嗎?會,還會相當熱切的祈禱,神要滿足他們很多東西.他們越來越看不到神了,越來越看到自己.

(2)牛犢

出32:4　亞倫從他們手裡接過來,鑄了一隻牛犢,用雕刻的器具做成.他們就說:「以色列啊,這是領你出埃及地的神.」

出32:5　亞倫看見,就在牛犢面前築壇,且宣告說:「明日要向耶和華守節.」

出32:6　次日清早,百姓起來獻燔祭和平安祭,就坐下吃喝,起來玩耍.

他們造了一隻牛犢,來代表不可見的耶和華,他們將神變成可操控的東西.

「牛」,是當時常見的神明形象.不過,「牛」其實也是現代都巿常常膜拜的神明.在各地的金融中心,證券交易所前面,通常見到「牛」的形象.在以色列人的歷史裡,後來的北國以色列國,真的是拜「牛」的.他們在「伯特利」和「但」這兩個地方,立了「牛」的像,代表耶和華,讓人膜拜.這種崇拜還延續了很長歷史.若果《出埃及記》最後編輯的年代是較後期的話,則這裡的故事有批評後來以色列國膜拜金牛的作用.

在當時,由「埃及」到「美索不達米亞」,「牛」是神明常見的形象.一般而言,「牛」的崇拜與「繁殖」或「力量」有關.在這裡,重點應在「力量」而不在「繁殖」,因為他們當時缺乏的,是一個帶領他們的領袖.他們要造像,是因為他們想有個神像「可以在他們前面引路」,而當「牛」造成時,他們說:「以色列啊,這是領你出埃及地的神」.要能「在他們前面引路」,要能「出埃及」,重點都不在「繁殖」,而在「力量」.他們需要明確的力量形像,去對抗沿途的敵人,去應付路上的不測.他們向「牛」像「獻燔祭和平安祭」,之後,「就坐下吃喝,起來玩耍」,這應不是崇拜「繁殖」之神可能有的淫亂行為,而應是預祝戰勝的慶祝.或許,他們覺得這「牛」的形象很威猛,「牛」的角很有力量,能嚇怕敵人,能戰無不勝,能引領他們到達夢想之地.他們因這「牛」的形象而覺得安心,他們獻祭,宴樂,慶祝.

他們的心安,建基於自己造出來的偶像,而不是相信創造萬有的耶和華.表面上,這只是耶和華神的具體形象,實質上,他們相信自己的力量多於相信神的同在.

好好的一個民族,本來帶著神救贖的深刻記號,自由的,往夢想之地前進.但因著心底的恐懼,因感到神的不實在,便走向人間常見的「偶像」,膜拜「牛」,擁抱「牛角」所代表的「武力」,然後,覺得穩當,覺得安心.

三個月前,他們才出了埃及,真想不到,他們這麼快便忘記了神的真實,他們這麼快便不信任神的帶領,他們如此容易便墮落了.

(3)神的仁慈

地上的群眾膜拜「牛」,因為他們認為神看不見他們.神真的看不見他們嗎?

天上的神看見.

出32:7　耶和華吩咐摩西說:「下去吧,因為你的百姓,就是你從埃及地領出來的,已經敗壞了.」

神對人的不信與墮落,作出判決.

出32:9　耶和華對摩西說:「我看這百姓真是硬著頸項的百姓.

出32:10　你且由著我,我要向他們發烈怒,將他們滅絕,使你的後裔成為大國.」

我們看看摩西的回應,他真是一個好的領導.

出32:11　摩西便懇求耶和華─他的神說:「耶和華啊,你為甚麼向你的百姓發烈怒呢?這百姓是你用大力和大能的手從埃及地領出來的.

出32:12　為甚麼使埃及人議論說『他領他們出去,是要降禍與他們,把他們殺在山中,將他們從地上除滅』?求你轉意,不發你的烈怒,後悔,不降禍與你的百姓.

出32:13　求你記念你的僕人亞伯拉罕,以撒,以色列.你曾指著自己起誓說:『我必使你們的後裔像天上的星那樣多,並且我所應許的這全地,必給你們的後裔,他們要永遠承受為業.』」

基本上,摩西講了三個重點.

一,神的「大力和大能」,是用來將人「從埃及地領出來的」.力量不是用來殲滅對方的,是用來救助弱者的.摩西懇求神以力量救贖,而不是殲滅.當然,這也是對地上所有人的提醒,不要膜拜「牛」,不要視「暴力」為膜拜的對象,要善用力量去救助人.誤用「力量」,只會招致神的審判.

二,摩西向神說:「為甚麼使埃及人議論說:你領他們出去,是要降禍與他們,把他們殺在山中,將他們從地上除滅呢?」言下之意,神若擊殺這群墮落的人,會引起別人的嘲笑,「又話救他們,原來是想殺他們」.神要為自己的榮譽而戰.若訴諸神的榮譽,人就有得救的可能.若訴諸人的行為,則死路一條.

三,摩西求神記念昔日與列祖立的盟約.不要看當下人的背信棄義,請看當日神你自己的信誓旦旦.同樣,訴諸人的守約,死路一條.訴諸神的信實,則有得救的可能.

摩西沒有訴諸人的任何優點,只訴諸神的憐憫.事實上,人的恐懼,軟弱,無依無靠,令人常常走向可倚靠,可掌控的暴力,但這只會使一個群體沉淪與滅亡.要相信神,神看來不可見,但事實上,祂常常看顧著我們.

最後,神按祂的仁愛,寬恕了敗壞的人.

出32:14　於是耶和華後悔,不把所說的禍降與他的百姓.

我們的社會都需要有領袖在我們前面引路,希望這領袖不會帶領我們膜拜「牛」,而能帶領我們體會仁愛的力量,帶領我們到達神那應許之地,帶領我們看見那不能看見的神.

但願:榮耀歸於聖父,聖子,聖靈.阿們.

(附記:本人在崇基禮拜堂講道的講章已結集成《荷花池畔春草綠》講章集,有興趣的朋友可到各大基督教書局購買.試讀本可參閱:www.news.ccc.cuhk.edu.hk/do/wp-content/uploads/2019/07/CCC-SK-Tang-Sermon-Text-Sample-2.pdf.)



\newpage

\section{資源位置與社群共融}
\label{sec:470}
\textbf{林豪恩先生 - 資源,位置,與社群共融}
\newline
\newline
\hyperref[sec:469]{< < < PREV SERMON < < <}
~
\hyperref[sec:toc]{[返目錄]}
~
\hyperref[sec:471]{> > > NEXT SERMON > > >}


阿摩司書 8:4-7
\newline
\begin{longtable}{cl}
\hline
\hline
章節 & 經文 (和合本修訂版)\\
\hline
8:4 & \begin{tabularx}{0.7\textwidth}{X} 你們這些踐踏貧窮人、使這地困苦人衰敗的，當聽這話！ \end{tabularx} \\ \\ \relax
8:5 & \begin{tabularx}{0.7\textwidth}{X} 你們說：「初一幾時過去，我們好賣糧；安息日幾時過去，我們好擺開穀物；我們要把伊法變小，把舍客勒變大，以詭詐的天平欺哄人， \end{tabularx} \\ \\ \relax
8:6 & \begin{tabularx}{0.7\textwidth}{X} 用銀子買貧寒人，以一雙鞋換貧窮人，把壞的穀物賣給人。」 \end{tabularx} \\ \\ \relax
8:7 & \begin{tabularx}{0.7\textwidth}{X} 耶和華指著雅各的驕傲起誓說：「他們這一切的行為，我必永遠不忘。 \end{tabularx} \\ \\
[1ex]
\hline
\hline
\end{longtable}
提摩太前書 2:1-7
\newline
\begin{longtable}{cl}
\hline
\hline
章節 & 經文 (和合本修訂版)\\
\hline
2:1 & \begin{tabularx}{0.7\textwidth}{X} 所以，我勸你，首先要為人人祈求、禱告、代求、感謝； \end{tabularx} \\ \\ \relax
2:2 & \begin{tabularx}{0.7\textwidth}{X} 為君王和一切在位的，也要如此，使我們能夠敬虔端正地過平穩寧靜的生活。 \end{tabularx} \\ \\ \relax
2:3 & \begin{tabularx}{0.7\textwidth}{X} 這是好的，在我們的救主神面前可蒙悅納。 \end{tabularx} \\ \\ \relax
2:4 & \begin{tabularx}{0.7\textwidth}{X} 他願意人人得救，並得以認識真理。 \end{tabularx} \\ \\ \relax
2:5 & \begin{tabularx}{0.7\textwidth}{X} 因為只有一位神，在神和人之間也只有一位中保，是成為人的基督耶穌。 \end{tabularx} \\ \\ \relax
2:6 & \begin{tabularx}{0.7\textwidth}{X} 他獻上自己作人人的贖價；在適當的時候這事已經證實了。 \end{tabularx} \\ \\ \relax
2:7 & \begin{tabularx}{0.7\textwidth}{X} 我為此奉派作傳道，作使徒，在信仰和真理上作外邦人的教師。我說的是真話，不是說謊。 \end{tabularx} \\ \\
[1ex]
\hline
\hline
\end{longtable}
路加福音 16:1-13
\newline
\begin{longtable}{cl}
\hline
\hline
章節 & 經文 (和合本修訂版)\\
\hline
16:1 & \begin{tabularx}{0.7\textwidth}{X} 耶穌又對門徒說：「某財主有一個管家，有人向主人告管家浪費他的財物。 \end{tabularx} \\ \\ \relax
16:2 & \begin{tabularx}{0.7\textwidth}{X} 主人叫他來，對他說：『我聽到了，你做的是甚麼事？把你所經管的交代清楚，你不能再作我的管家了。』 \end{tabularx} \\ \\ \relax
16:3 & \begin{tabularx}{0.7\textwidth}{X} 那管家心裡說：『主人辭我，不用我再作管家，我將來做甚麼呢？鋤地嘛，沒有力氣；討飯嘛，怕羞。 \end{tabularx} \\ \\ \relax
16:4 & \begin{tabularx}{0.7\textwidth}{X} 我知道怎麼做，好叫人們在我不作管家之後，接我到他們家裡去。』 \end{tabularx} \\ \\ \relax
16:5 & \begin{tabularx}{0.7\textwidth}{X} 於是他把欠他主人債的，一個一個地叫了來，問頭一個說：『你欠我主人多少？』 \end{tabularx} \\ \\ \relax
16:6 & \begin{tabularx}{0.7\textwidth}{X} 他說：『一百簍油。』管家對他說：『拿你的賬，快坐下，寫五十。』 \end{tabularx} \\ \\ \relax
16:7 & \begin{tabularx}{0.7\textwidth}{X} 他問另一個說：『你欠多少？』他說：『一百石麥子。』管家對他說：『拿你的賬，寫八十。』 \end{tabularx} \\ \\ \relax
16:8 & \begin{tabularx}{0.7\textwidth}{X} 主人就誇獎這不義的管家做事精明，因為今世之子應付自己的世代比光明之子更加精明。 \end{tabularx} \\ \\ \relax
16:9 & \begin{tabularx}{0.7\textwidth}{X} 我又告訴你們，要藉著那不義的錢財結交朋友，到了錢財無用的時候，他們可以接你們到永遠的住處去。 \end{tabularx} \\ \\ \relax
16:10 & \begin{tabularx}{0.7\textwidth}{X} 人在最小的事上忠心，在大事上也忠心；在最小的事上不義，在大事上也不義。 \end{tabularx} \\ \\ \relax
16:11 & \begin{tabularx}{0.7\textwidth}{X} 若是你們在不義的錢財上不忠心，誰還把那真實的錢財託付你們呢？ \end{tabularx} \\ \\ \relax
16:12 & \begin{tabularx}{0.7\textwidth}{X} 如果你們在別人的東西上不忠心，誰還把你們自己的東西給你們呢？ \end{tabularx} \\ \\ \relax
16:13 & \begin{tabularx}{0.7\textwidth}{X} 一個僕人不能服侍兩個主；他不是恨這個愛那個，就是重這個輕那個。你們不能又服侍神，又服侍瑪門。」 \end{tabularx} \\ \\
[1ex]
\hline
\hline
\end{longtable}
講題:資源,位置,與社群共融 Resource, Position, and Social Solidarity

經文:阿摩司書8:4-7;詩篇113;提摩太前書2:1-7;路加福音16:1-13

講員:林豪恩先生

場合:香港中文大學崇基學院禮拜堂主日崇拜

日期:2019年9月22日

弟兄姊妹平安.邀請大家一齊再閱讀及思考今天經課的四段經文.

一.阿摩司書:資源差異造成的社群分裂

8:4你們這些要吞吃窮乏人,使困苦人衰敗的,當聽我的話!

8:5你們說:月朔幾時過去,我們好賣糧；安息日幾時過去,我們好擺開麥子；賣出用小升斗,收銀用大戥子,用詭詐的天平欺哄人,

8:6好用銀子買貧寒人,用一雙鞋換窮乏人,將壞了的麥子賣給人.

8:7耶和華指著雅各的榮耀起誓說:他們的一切行為,我必永遠不忘.

經文中提及當時以色列社會中分開了兩種人,這兩種人是被甚麼分開的呢?阿摩司稱一邊為稱另一群為「窮乏人」,「困苦人」,「貧寒人」,稱另一群為「你們」.經文用第二人稱,因此「你們」應該是阿摩司的目標聽眾.「你們」是誰呢?按阿摩司的描述,是「賣糧」的人.古代以色列的理想社會,是每個家庭有自己的土地,理論上糧食是自給自足的.阿摩司時代出現「賣糧」的人,顯示社會中有些家庭已經失去土地,無法自給自足；也顯示有些家庭已經進入了農業商業化的階段,生產糧食作商業用途.失去土地的原因可能各有不同,晉身成為糧食供應商的也各有條件.在阿摩司時代,社會存在著這兩個不同擁有資源差異的群體,同時,資源的差異又再成為擴大兩者之間的差異的條件.賣糧的與買糧的這兩個群體的關係如何呢?如何互相對待呢?經文描述賣糧的如此對待買糧的:

8:5你們說:月朔幾時過去,我們好賣糧；安息日幾時過去,我們好擺開麥子

第一,賣糧者任何時間也想著如何能夠從買糧者賺取更多.阿摩司時代的賣糧者連被視為神聖的安息日也不耐煩,因為阻礙了他們「做生意」.現代人是能夠理解的,「阻人發達,罪大惡極」這俗語就表達了我們的價值.與阿摩司的時代不同,我們已經除去了所有人為的阻礙,無論是年初一還是年初二,每論是早晚還是午夜,生意仍可以繼續做.不過,令人困擾的是大自然,例如強烈風暴就令停工停市,是「阻人發達」的「惡罪」.因此,傳言有「李氏力場」阻擋風暴來襲,避免因停工,停市而招致經濟損失.

8:5「賣出用小升斗,收銀用大戥子,用詭詐的天平欺哄人／將壞了的麥子賣給人」

第二,賣糧的人總會用各種方法去從買糧的人身上賺取更多.阿摩司時代對顧客的欺騙手法是肉眼可以見的,這些具體可見的欺騙方法,幾千年後仍大行其道.不過,現代科技提供了看不見的欺騙方法,在資訊就是資源的社會,「消息」是賺錢的工具.阿摩司時代賣「壞了的麥子」,我們的時代賣「有毒的奶粉」,其實,都是為了賺錢而不理其他人的死活,由此看來,現代人和古代人的差距並不太遠.

8:6「好用銀子買貧寒人,用一雙鞋換窮乏人」

第三,賣糧的不但賣糧給買糧的人,他們更「買起」買糧的人.聰明的人會知道,人比貨物更有價值,因為人可以不斷生產貨物,因此,買賣貨物演變成買賣人口.在阿摩司時代,為了生存(不單是自己的生存,更是家人的生存),沒錢買糧,就唯有「賣身」.幾千年後,仍然有人「買起」其他人嗎?關注現代奴隸議題的澳洲慈善機構「行走自由基金會」(Walk Free Foundation) 的「2016年全球奴役指數」(Global Slavery Index 2016)就指出,全球有4580萬遭別人以剝削為目的,透過暴力,威脅或其他手段(包括誘拐,欺詐,欺騙,濫用權力或濫用脆弱境況等)被販賣的人.

在我們的城市中,可能很少人需要為糧食而「賣身」,然而,不少人要為另一種生存必需品而「賣身」.有沒有聽過「房奴」或「樓奴」這些詞語,兩者其實都指那些為了支付高額的住房按揭貸款,而使自身社會生活受到極大影響的人.就像古代的奴隸社會,奴隸為了生存而為奴隸主賣命幹活一樣,房奴為了償還住房貸款而賣命工作.而最脆弱的一群,是連晉身房奴都沒有資源,這些人住在租金尺價比豪宅更貴的tong 房,沒有獨立電錶和水錶,業主收取的電費水費比電力公司及水務署所收的更貴.

資源多的如此對待資源少的,與基督教信仰有甚麼關係呢?最後一節經文如此說:

8:7耶和華指著雅各的榮耀起誓說:他們的一切行為,我必永遠不忘.

甚麼事情令耶和華上帝那麼上心呢?甚麼事情令天地的主耿耿於懷呢?在阿摩司時代,就是那些賣糧的人欺壓買糧的人；那麼,在我們的年代呢?

二.詩篇:位置差異造成的社群分裂

113:1你們要讚美耶和華!耶和華的僕人哪,你們要讚美,讚美耶和華的名!

113:2耶和華的名是應當稱頌的,從今時直到永遠!

113:3從日出之地到日落之處,耶和華的名是應當讚美的!

113:4耶和華超乎萬民之上；他的榮耀高過諸天.

113:5誰像耶和華－我們的上帝呢?他坐在至高之處,

113:6自己謙卑,觀看天上地下的事.

113:7他從灰塵裡抬舉貧寒人,從糞堆中提拔窮乏人,

113:8使他們與王子同坐,就是與本國的王子同坐.

113:9他使不能生育的婦人安居家中,為多子的樂母.你們要讚美耶和華!

詩篇113篇主要提及兩種人,一種是上文阿摩司書提及的「貧寒人」,「窮乏人」,另一種是「王子」,「本國的王子」.阿摩司描述「賣糧的人」與「貧寒人」的差異主要在資源方面,詩人提及的「貧寒人」與「王子」的差異主要在位置方面.位置不同,在社會狀況而言是地位和尊卑的差異.詩人描述「貧寒人」的位置與「灰塵」有關,即是無關重要,在社會「無份」；「窮乏人」的位置在「糞堆」中,即是比一般垃圾更令人厭惡.「王子」的位置呢?詩人描述「王子」在座位上,這不是一般的座位,是高尚,尊貴,榮耀的座位.有傳說合理化人的尊卑,例如「婆羅門是原人的嘴,剎帝利是原人的雙臂,吠舍是原人的大腿,首陀羅是原人的腳.至於賤民,則被排除在原人的身體之外.」《吠陀》不過,聖經的教導不是這樣的.創世記如此說:

1:26上帝說:「我們要照著我們的形像,按著我們的樣式造人」.如果每個人都是照著上帝的形像,按著上帝的樣式做的,何以合理化有些人輕於灰塵,卑於糞堆,而另一些人高高在上,尊貴榮華呢?

人被分高低貴賤,是因為人自己制造出來的價值觀念,例如詩113最後一句:「113:9他使不能生育的婦人安居家中,為多子的樂母」反映出當時的人認為婦女多產是有福及有地位,不能生育是被咀咒和被看不起；人被分高低貴賤同時也是因為人自己制造出來的制度,撒母耳先知當年對以色列人要建立王國制度的警告就是一個清楚的事例:

8:11「管轄你們的王必這樣行:他必派你們的兒子為他趕車,跟馬,奔走在車前；

8:12又派他們作千夫長,五十夫長,為他耕種田地,收割莊稼,打造軍器和車上的器械；

8:13必取你們的女兒為他製造香膏,做飯烤餅；

8:14也必取你們最好的田地,葡萄園,橄欖園賜給他的臣僕.

8:15你們的糧食和葡萄園所出的,他必取十分之一給他的太監和臣僕；

8:16又必取你們的僕人婢女,健壯的少年人和你們的驢,供他的差役.

8:17你們的羊群他必取十分之一,你們也必作他的僕人.

基督教信仰又與這些把人分為高低尊卑的狀況有甚麼關係呢?詩人說:

113:7他從灰塵裡抬舉貧寒人,從糞堆中提拔窮乏人,

113:8使他們與王子同坐,就是與本國的王子同坐.

詩113是一首讚美詩,詩人多次呼籲「你們要讚美耶和華!」我們有沒有想過,我們讚美的耶和華是一位怎樣的上帝呢?詩113指出因為甚麼而讚美耶和華呢?這些問題值得我們思考.

三.提摩太前書:各按各職,敬虔,端正,平安無事地度日

2:1我勸你,第一要為萬人懇求,禱告,代求,祝謝；

2:2為君王和一切在位的,也該如此,使我們可以敬虔,端正,平安無事地度日.

2:3這是好的,在上帝我們救主面前可蒙悅納.

2:4他願意萬人得救,明白真道.

2:5因為只有一位上帝,在上帝和人中間,只有一位中保,乃是降世為人的基督耶穌；

2:6他捨自己作萬人的贖價,到了時候,這事必證明出來.

2:7我為此奉派作傳道的,作使徒,作外邦人的師傅,教導他們相信,學習真道.我說的是真話,並不是謊言.

剛才我們引述創世記的經文,從創造的角度發現每個人都按著上帝的形象和樣式被創造,因此,每個人都擁有相同的尊嚴和價值.保羅在提摩太前書帶引我們從救贖的角度去看:

2:4他願意萬人得救,明白真道.

2:5因為只有一位上帝,在上帝和人中間,只有一位中保,乃是降世為人的基督耶穌；

2:6他捨自己作萬人的贖價,到了時候,這事必證明出來.

每一個人都非常寶貴,每一個人都非常有價值,每一個人都重要到一個程度聖子耶穌基督捨自己作他／她的贖價.那麼,人們應該如何相待呢?保羅的教導是:

2:1我勸你,第一要為萬人懇求,禱告,代求,祝謝；

2:2為君王和一切在位的,也該如此

在保羅的時代,正如其他時代,人們在社會中處身於不同的位置,有些是平民百姓,有些是君王和在位的,即是有管治者及被管治者.雖然人們處身不同的位置,但前題是每個人都有價值得勝子耶穌基督為他／她捨棄性命,因此,管治者和被管治者都同樣有價值,同樣值得被重視.保羅特別提出在祈禱的對象中應該包括君王和在位的,讓我們猜想當時的信徒群體未必會為君王和在位的祈禱,其實,這也是可以理解的.基督徒在當時處身於被管治的位置,而且是被邊緣化的弱勢群體,自身不保,還要為那些高高在上的君王和在位者祈禱?那些君王和在位者對基徒有多友善,有多少幫助,以致要為他們祈禱呢?其實可能沒有,以保羅的個人經驗來看,少打兩棒已經很好了.保羅的勸導,不是單單從現實出發,而是從神學出發,每一個人都重要到一個程度聖子耶穌基督捨自己作他／她的贖價.保羅突破自己處身的位置去看其他人,突破個人的遭遇去看當時的君王和在位者,保羅的胸襟實在令人大開眼界.

不過,保羅呼籲信徒為君王和在位的祈禱,並不是景仰他們,也不是奉承他們,而是祈求上主令他們在他們的位置上做好他們的角色,只有他們做好他們的角色,建立良好的社會氣氛和條件,人們才能夠「敬虔,端正,平安無事地度日」.君王和在位者應該創造怎樣的社會環境呢?保羅提出三項:第一,「平安無事」,讓人們可以在一個安全和免於恐懼的狀況下生活；第二,「端正」,讓人們可以生活得有品格,有尊嚴；第三,「敬虔」,讓人們可以自由地按著自己的信仰敬拜上主.

四.路加福首:要藉著那不義的錢財結交朋友,到了錢財無用的時候,他們可以接你們到永存的帳幕裡去

如果保羅在提摩太前書的教導提示我們如何回應人與人之間位置差異的狀況,耶穌在路加福音所講的故事讓我們思想如何回應人與人之間資源差異的狀況.

16:1耶穌又對門徒說:「有一個財主的管家,別人向他主人告他浪費主人的財物.

16:2主人叫他來,對他說:『我聽見你這事怎麼樣呢?把你所經管的交代明白,因你不能再作我的管家.』

16:3那管家心裡說:『主人辭我,不用我再作管家,我將來做甚麼?鋤地呢?無力；討飯呢?怕羞.

16:4我知道怎麼行,好叫人在我不作管家之後,接我到他們家裡去.』

16:5於是把欠他主人債的,一個一個地叫了來,問頭一個說:『你欠我主人多少?』

16:6他說:『一百簍油.』管家說:『拿你的帳,快坐下,寫五十.』

16:7又問一個說:『你欠多少?』他說:『一百石麥子.』管家說:『拿你的帳,寫八十.』

16:8主人就誇獎這不義的管家做事聰明；因為今世之子,在世事之上,較比光明之子更加聰明.

16:9我又告訴你們,要藉著那不義的錢財結交朋友,到了錢財無用的時候,他們可以接你們到永存的帳幕裡去.

16:10人在最小的事上忠心,在大事上也忠心；在最小的事上不義,在大事上也不義.

16:11倘若你們在不義的錢財上不忠心,誰還把那真實的錢財託付你們呢?

16:12倘若你們在別人的東西上不忠心,誰還把你們自己的東西給你們呢?

16:13一個僕人不能事奉兩個主；不是惡這個愛那個,就是重這個輕那個.你們不能又事奉上帝,又事奉瑪門.」

讀者不手用爭論故事中那位管家的行為是否值得仿傚,因為16:8已明言這位是「不義的管家」,不過,那位管家所問的問題,其實每個人也應該問:

16:3那管家心裡說:『主人辭我,不用我再作管家,我將來做甚麼?鋤地呢?無力；討飯呢?怕羞.

而那位管家的找到方法解決「將來做甚麼?往哪裡去?」的問題:

16:4我知道怎麼行,好叫人在我不作管家之後,接我到他們家裡去.』

其他人又怎樣呢?是否也應該思考一下自己的將來呢?如果以為有錢財就有將來,那麼,就要聽聽耶穌藉著這個故事的提醒了:

16:9我又告訴你們,要藉著那不義的錢財結交朋友,到了錢財無用的時候,他們可以接你們到永存的帳幕裡去.

阿摩司先知時代的人,運用他們資源上的優勢去令其他人受苦,令社群分化,耶穌提醒,擁有資源優勢的應該運用資源去建立人與人之間的情誼,創造社群共融.資源運用與資源分配,似乎是一個社會議題,最終卻是一個信仰議題,因為耶穌總結這個故事時說:

16:13一個僕人不能事奉兩個主；不是惡這個愛那個,就是重這個輕那個.你們不能又事奉上帝,又事奉瑪門.」

五.總結

我們身處的制度和文化有時會限制我們的想像,以為一定是這樣的,沒有其他可能性.其實世界很大,社群之間人與人的關係有很多可能性.有沒有想像過社群中人與人可以互相體諒呢?有沒有想像過可以互相照顧呢?有一位香港學生往歐洲交流,認識了一位從法羅群島來的交流生,後來嫁到法羅群島去.她發現在法羅群島,有錢人的屋和普通人的房屋都差不多,原因是有錢人不想自己的屋太豪華而令其他人不好受.人們願意交重稅,令大學生有生活津貼,孕婦有一年的有薪產假,失業和沒有工作能力的人有足夠生活津貼,種種的福利令國民享有平等的生活.那裡沒有陽光與海灘,一年有三分之二的時間都在下雨,即使夏天平均溫度也只有13度.在這個長期陰冷和灰蒙蒙的國度,她卻發現人民的笑容比其他地方來得燦爛.



\newpage

\section{前輩教後輩打美好的仗}
\label{sec:471}
\textbf{鄭漢文博士 - 「前輩」教「後輩」打美好的仗}
\newline
\newline
\hyperref[sec:470]{< < < PREV SERMON < < <}
~
\hyperref[sec:toc]{[返目錄]}
~
\hyperref[sec:472]{> > > NEXT SERMON > > >}


講題:「前輩」教「後輩」打美好的仗 “Oldies” teach “Youngsters” how to fight a good war

三代經課:提摩太前書6:6-19

講員:鄭漢文博士

場合:香港中文大學崇基學院禮拜堂主日崇拜

日期:2019年9月29日

保羅教提摩太:

自己當打的年華已消逝,消化過去的經驗傳給下一代.

大家好,今天是我今年第二次在此講道,今次講完之後,會否還有得講?

自從1995年我開始任教於中大教育學院,二十五年來,每次教課我都感到可能是最後一堂課,教埋就無得再教了.所以,每一堂課,我都會盡心盡力教好書.

同樣,十多年前開始講道,每次都是用同樣的心境去講道,可能講埋這一次就無得講.所以,每一堂道,我都會盡心盡性盡意盡力去講好道.

今日的經課係取自提摩太前書6章6節至19節,讓我們現在一起讀一次,我先讀一節,姊妹接著讀下一節,弟兄接著讀再下一節,然後循環讀完:

6然而,敬虔加上知足的心便是大利了；

7因為我們沒有帶甚麼到世上來,也不能帶甚麼去.

8只要有衣有食,就當知足.

9但那些想要發財的人,就陷在迷惑,落在網羅和許多無知有害的私慾裡,叫人沉在敗壞和滅亡中.

10貪財是萬惡之根.有人貪戀錢財,就被引誘離了真道,用許多愁苦把自己刺透了.

11但你這屬神的人要逃避這些事,追求公義,敬虔,信心,愛心,忍耐,溫柔.

12你要為真道打那美好的仗,持定永生.你為此被召,也在許多見證人面前,已經作了那美好的見證.

13我在叫萬物生活的神面前,並在向本丟‧彼拉多作過那美好見證的基督耶穌面前囑咐你:

14要守這命令,毫不玷污,無可指責,直到我們的主耶穌基督顯現.

15到了日期,那可稱頌,獨有權能的萬王之王,萬主之主,

16就是那獨一不死,住在人不能靠近的光裡,是人未曾看見,也是不能看見的,要將他顯明出來.但願尊貴和永遠的權能都歸給他.阿們!

17你要囑咐那些今世富足的人,不要自高,也不要倚靠無定的錢財；只要倚靠那厚賜百物給我們享受的神.

18又要囑咐他們行善,在好事上富足,甘心施捨,樂意供給人,

19為自己積成美好的根基,預備將來,叫他們持定那真正的生命.

在崇基禮拜堂講道,需要至少三個月前就開始準備,一兩個月前就要提交講道題目.

當六月天開始燒起城市動盪時,我在七月被自己母會循道衛理聯合教會《會訊》編輯邀請,為教會在此時此地的社會參與可以如何回應,寫篇短函.

作為一個年輕人眼中的「廢老」,我不認為可以給「反送中抗爭運動」的「年青運動員」甚麼建議.但編輯叫到,我就講幾句,「arm 聽就參考o下」,不過這段文字刊出時會否已經過時,不得而知.

當編輯邀請我寫千字文時,我正要為9月29日以榮譽助理校牧身分在崇基禮拜堂講道而定題目,我因要寫這篇文章而生發靈感,定了這個教友可能會感到出奇或難受的命題:

《「廢老」教「廢青」打美好的仗》

(Oldies teach youngsters how to fight a good war).

9.29當日經課是提前6章6至19節.我8.13那日想到的講道大綱暫列如下:

保羅教提摩太: 他自己當打的年華已消逝,消化過去的經驗傳給下一代.

夠就夠的知足: 屬靈德性需有當下的節制

非物質主義: 心靈的財富不在貪戀錢財

美好的仗: 成為美好見證為未來設希望

廢老教廢青: 還有未盡的話

九月《會訊》終於出版了,我這篇題為《一個「廢老」給年輕信徒的牧養家書》的千文字沒有被審查抽走或被刪改,得享受了一次充分的言論自由,詳情大家可以上網看看我寫了甚麼,網址在此:

\url{https://www.methodist.org.hk/media/monthlynews/pdf/2019/09/11/iss364.pdf}

(我取了兩份實體印本,現在傳給你們看,傳完請放在出口交還我.)

今日講道的講題是:《「前輩」教「後輩」打美好的仗》

(Oldies teach youngsters how to fight a good war).

前輩能給後輩的是打過美好的仗的經驗,但經驗好易過時,經驗需要借助T來提煉成為智慧,後輩在更新變化中,才會對新處境有所對應,應時地回應.

T究竟是甚麼?

T可以是Theory(理論);

T可以是Theology(神學)；

T可以是Sacred Text(聖經);

T可以是Tradition(傳統).

今日講道,我就試用四個不同的T來解讀保羅這位前輩教提摩太這位後輩打過美好的仗的經驗,對我們今日信徒在世信仰實踐的提示.

第一個T是Theory(理論).

在2006年我所創建的價值教育碩士課程裡,我發展了一個叫VIP的價值教育理論模型(VIP model as theoretical framework for Values Education).

VIP的V指的可以是Values價值,可以是Visions願景或遠象或生命境界,可以是Virtues品性或德性.

VIP的I指的是Identities身分,身為X就有Y的本分,有待履行.

VIP的P指的是Practice(s)實踐,就是多人同行的行動方式.

價值轉化身分,身分強化實踐.

(Values transform identities which empower practices.)

我發現,今日經課提前6章6至19節可以借VIP分為三段來解讀.

2.　甚麼是我們的價值Values,遠象Visions與德性Virtues?

首先,提摩太前書6: 6-11與V相關.讓我們先重看.

首先,這裡描繪了一個價值觀上的進退取捨,涉及財富與生命境界.

按價值教育三分法,可以這樣分出對待財富的三個生命層次或生命境界Visions of Life:

第一種人,財富少過自己所需,繼續增加財富,為的是要脫貧.

第二種人,財富足夠自己所需,努力創富,為的可能是安全感,也可為扶貧.

第三種人,財富多過自己所需,繼續增加財富,為的是甚麼呢?

前輩保羅力勸後輩提摩太,甚麼是要退避三舍的?

第9及10節這樣具體描述第三種人生命境況的惡:

9但那些想要發財的人,就陷在迷惑,落在網羅和許多無知有害的私慾裡,叫人沉在敗壞和滅亡中.

10貪財是萬惡之根.有人貪戀錢財,就被引誘離了真道,用許多愁苦把自己刺透了.

貪,是萬惡之根.貪得無厭,令人討厭.

前輩保羅力勸後輩提摩太,活出「夠就夠」的知足的生命境界.

提摩太前書6: 6-8節勸說:

6然而,敬虔加上知足的心便是大利了；

7因為我們沒有帶甚麼到世上來,也不能帶甚麼去.

8只要有衣有食,就當知足.

要有當下的節制,才可以養成屬靈德性Spiritual Virtues.

用現在香港的流行話說,有研究年青人的學者也這樣說,2014年之後的年輕一代,出現了一種新的價值觀,叫做:「後物質主義」(Post-materialism),就是物質的追求夠就夠,有餘的精力寧願花在追求心靈的富足.

第一個明顯跨代差異(cross-generational gap)就是新生代不做樓奴,認為住的地方不需要買入,不必為做業主而成為打工樓奴,更會批判以土地資源作買賣是惡的根源,提出土地公義的訴求.

在上世紀八,九十年代成長的,有許多努力創富的長輩認為這一代年輕人毫不長進,所以稱他們為「廢青」.這班前輩,不少是老闆或上司的或私下或公開地下判斷:「廢青」打工的態度不夠投入,「抵佢一世無樓」.

相反,這一代年輕人認為這些自認前輩的「廢老」,仍然活在資本主義,物質主義,工具主義的奴役中,無法進入新時代,非常離地,不懂前線在打甚麼的仗.

「Facebook老餅,難以連登.」

我們有神信仰的,怎樣看我們的生命呢?

保羅在6章11節說:

11但你這屬神的人要逃避這些事,追求公義,敬虔,信心,愛心,忍耐,溫柔.

也就是說:心靈的財富不在貪戀錢財,更要在遠離財富生產機制的惡,返回一個十分重要的身分認同,就是屬神的人.

唯有屬神的人,我們才可能養成公義,敬虔,信心,愛心,忍耐,溫柔等屬靈品性.

甚麼是我們的身分Identity?作個見證甚麼的見證人?

「屬神的人」就是身分認同.

為讓大家對身分之從認識到認同,我們可以引入第二個T,就是Theology神學,而一般信徒要認識的最少的神學(minimal theology of common believers),其實就在我們每次崇拜都公開認信而頌讀的《使徒信經》.

《使徒信經》的每一句,都是用「我信」開首.

我徵求了署理校牧同意,十一月開主日學,就是與大家一同研讀《使徒信經》.

所以詳細的分析,屆時再講論.以上是植入式廣告.

當我們以《使徒信經》認信時,就認信了聖父,聖子,聖靈的三一神學.

用哲學邏輯的術語,每次崇拜認信,就在「言辭行為」(speech act)上作了「神學上的委身」(theological commitment),向公共宣認自己相信一個神,祂同時是創造主的聖父,是救贖主的聖子,是生命主的聖靈,是三位一體的.

如是者,當T是三一神信仰時看人生在世的種種進退取捨,就可以有三層反問值得思而後行:

一,我作為聖父上帝天國子民,在地國這樣行,恰當嗎?

二,我作為聖子主耶穌基督門徒,在社會這樣行,適宜嗎?

(WWJD: What would Jesus do?)

三,我作為聖靈感動的生命,在人間這樣行,適當嗎?

以上三問,若有任何不安,找找其他信徒商量下.

其實,早在上世紀八十年代的當代思潮,已經出現身分政治(identity politics)的論述.

許多教會領袖都主張教會遠離政治.

其實信徒作為一種身分,政治就擺脫不了.

當一個人認信為天國子民後,面對地上國度,能真的政治中立嗎?

我攻讀神道學學士時,需要有兩次實習.第一次來到崇基禮拜堂,第二次去了浸會大學校牧室,兩次都被牧師分派到牧養普通話內地生團體.

有一次在浸會大學崇拜講道時,向來港留學的內地生分解天國子民的「本體論上的委身」(ontological commitment),就是當你相信耶穌基督所傳的福音是天國近了,你們悔改,成為天國子民的同時,就是表明在優先次序上天國比任地上的國更重要,要知效忠的是天國.

即是說,在兩國爭戰時,你會站在天國立場,拒絕為地國權勢而傷害天國,傷害天國子民.

一個信徒越接近地國權勢,越有機會遠離天國.

對於自命信徒的地上領袖,我們的提醒是:「你遠離天國了,你還不悔改?」

這就是我們信徒現在在打的仗.你敢言嗎?

對於保羅要教提摩太打美好的仗,他這樣對我們這等後輩說:

12你要為真道打那美好的仗,持定永生.你為此被召,也在許多見證人面前,已經作了那美好的見證.

13我在叫萬物生活的神面前,並在向本丟‧彼拉多作過那美好見證的基督耶穌面前囑咐你:

14要守這命令,毫不玷污,無可指責,直到我們的主耶穌基督顯現.

15到了日期,那可稱頌,獨有權能的萬王之王,萬主之主,

16就是那獨一不死,住在人不能靠近的光裡,是人未曾看見,也是不能看見的,要將他顯明出來.但願尊貴和永遠的權能都歸給他.阿們!

我們今日在世的身分,就是為打那美好的仗而被召,也在許多見證人面前,作了那美好的見證.

我們要效法的,就是作過那美好見證的基督耶穌.

我們這場仗越打得美好,就越把主耶穌基督顯現出來.

美國上世紀八,九十年代的神學,發展出「敘事神學」(Narrative Theology)這一學派,較為重視「故事」(Story)與「敘事」(Narrative)在神學思考的重要性.

這個學派大概再分出有三支來.其中一支流的一位代表人物叫James Wm. Clarendon,他寫了一套三部曲的系統神學,第一部談Ethics倫理,第二部談Doctrine教義,第三部談Witness見證.

按此思路,每一個信徒的生命故事,就是一個見證,見證:

一．如何為真道打那美好的仗,持定了永生.

二．如何在在許多見證人面前,作了你自己要作的見證.

三．你所作的見證是如何呈現美好.

再用哲學分析來澄清身分的意義,認同一個身分,就是意含:立志履行這個身分的相關本分,哲學上叫practical necessitation of duties.

例如:我是香港人,就要加油.

例如:我是中大人,就要守住中大人文傳統.

例如:我是基督徒,就要……

究竟Christian identity有甚麼本分?

認同一個身分,就是意含:立志履行這個身分的相關本分,也就會產生與你具有同一身分的那些走在前的前輩所留下一些囑咐,要求著認同的那些走在後的後輩在實踐上去履行.

所以,讓我轉入VIP的第三部分Practice實踐去思考經課後半部所出現的三個囑咐.

甚麼是我們的實踐Practice?讓世人得見基督的同在

今日經課的後半段出現了三個囑咐.

剛才我們已經用身分意識化為見證來思考了在13至16節的第一個囑咐:

13我在叫萬物生活的神面前,並在向本丟‧彼拉多作過那美好見證的基督耶穌面前囑咐你:

14要守這命令,毫不玷污,無可指責,直到我們的主耶穌基督顯現.

15到了日期,那可稱頌,獨有權能的萬王之王,萬主之主,

16就是那獨一不死,住在人不能靠近的光裡,是人未曾看見,也是不能看見的,要將他顯明出來.但願尊貴和永遠的權能都歸給他.阿們!

第二和第三個囑咐在17至18節:

17你要囑咐那些今世富足的人,不要自高,也不要倚靠無定的錢財；只要倚靠那厚賜百物給我們享受的神.

18又要囑咐他們行善,在好事上富足,甘心施捨,樂意供給人,

這兩個囑咐指向我們在VIP的V部份所提及的三種生命境況,再提一次:

第一種生命境界是:財富少過自己所需,繼續增加財富,為的是脫貧.

第二種生命境界是:財富足夠自己所需,努力創富,為的可能是安全感,也可為扶貧.

第三種生命境界是:財富多過自己所需,繼續增加財富,為的是貪財.

保羅的教導簡單地說,就是拒絕第三種貪財人生,站住第二種知足而感恩的人生,扶助第三種貧困人.

保羅提醒提摩太以後教導後人,今世富足的人不要以為自己「好醒,好掂」,不要持財傲物,忘了最後倚靠的是神.

保羅提醒提摩太以後教導後人,財富足夠自己所需,努力創富,為的是可以扶貧.

所謂創富,是創造心靈上的富足.

後物質主義,正是「廢老」稱這一代年輕人為「廢青」所追求的.

這些都是作為前輩的保羅的經驗之談,若我們加上兩個T來解說,以回應當今香港我們信徒如何回應,我想到了一段公共神學領域常引用的經文(Sacred Text)以及教會形成了的一點傳統(Church Tradition).

當信徒參與公共空間時,常引用到的經文是彌迦書六章8節,這段經文指向四個層次的基督徒實踐(Christian practice):

一是行公義

二是好憐憫

三是存謙卑的心

四是與上主同行

這四層次的基督徒同在的實踐(practice of Christian presence)讓人群得以看見基督同在(Christian presence).

從2014年雨傘運動到2019年反送中抗爭運動,人們如何得見基督同在?

在警民對抗的前線,在記者,醫護,社工,議員站在最前外,牧者也加入,站在張力最前線,作愛與和平之子.

6.12那天,非信徒從Sing Hallelujah to the Lord的歌聲中得見基督同在的和平.基督同在除了信徒進入人群空間外,另一方式是讓人群進入信徒與基督同在的空間(教會).

6.9至今,暴力出現在一區又一區,那區有暴力,那區教會(local churches)就得思考開放與否──讓有心靈休息需要的人群可以進入教會的空間?

再者,若T是傳統,四福音書記戴了耶穌的實踐空間,可以分為四大類:

一在內室:在被邊緣化者家裡坐席

二在室內:在聖殿中證道,在會堂內教導

三在街頭:在路上醫治,在海邊船上平靜風浪

四在山上:宣講八福登山寶訓,施行五餅二魚神蹟

效法基督,今天信徒的基督同在,可以考慮以下四種做法:

一在內室:家中私禱,多互相聆聽

二在室內:開祈禱會,辦研討會,多作心靈休息服侍

三在街頭:合法遊行,野餐,夢遊,唱詩,公禱

四在網上:聯署,做鍵盤戰士,轉發信息,善用社交媒體.(耶穌年代無互聯網,但今日「廢青」有.)

兄弟姊妹爬山,各自努力.

5.　廢老教廢青: 還有未盡的話

但是,我還有未盡的話,「廢老」總係「死唔斷氣」!

我近天國了,如何可以做到「廢老不廢」呢,這是我死前在思考的.

近天國了,還不悔改?

我在中大教了二十五年,二十五年來做老師好多遺憾,在此我要做個告解(confession),我最想寫的一本書是《懺悔錄》.以下是我三個歉疚:

我要向我的前輩道歉,有負師恩,弟子不才,無能把學問做得更好,無法把老師傳給我的好好傳給下一代.

我要向我的同輩道歉,有負同行,未能同心協力做好事,無法站在前線,與你們一同打美好的仗.

我要向我的後輩道歉,有負初心,曾經信託我的學生,我卻對你造成傷害,請你原諒我.你的原諒,可讓我死得無憾,你也不要帶著悲傷怒恨的心地生活下去.

回首往生,曾經對老師未能承命,對同行未能承志,對學生未能承託.如今夜夜難眠,晚禱有淚.

甚麼時候,我可以好像保羅對自己一生的總結,正如提摩太後書四章6至8節這樣的說:

「至於我,我已經被澆獻,離世的時候到了.那美好的仗我已經打過了,當跑的路我已經跑盡了,該信的道我已經守住了.從此以後,有公義的冠冕為我存留,就是按著公義審判的主到了那日要賜給我的；不但賜給我,也賜給凡愛慕他顯現的人.」

今日經課,還有未盡解說的話.

崇基學院神學院院長邢福增教授在2019年9月23日「週一加油站」講道也是用上今日經課,以「那美好的仗,真的能打嗎?」為題分享信息,這是他對今日經課的總論,借之補充我的不足:

「院長指出保羅以「打美好的仗」來勉勵基督徒,人生如打仗,不管你是否喜歡都在打仗.大家在人生不同階段有著不同目標,同時也會面對各種困難與挑戰.能否迎向挑戰,或是在挫敗後再次勇敢上路,都是人生的真實寫照.

那麼,我們如何才能像保羅所言,去打一場美好的仗?院長指出有時人因自我迷失與面對外在艱難環境而帶來的重擔,結果無法堅持.他勉勵大家堅找信仰的價值,不被物質及財富所吸引,並持守召命,莫忘初衷,堅持走下去,打一場美好的仗.

美好的仗,是仰望著那終極的勝利,在各種的疑惑,考驗,絕望,淚水或憂傷之中,都堅持不放棄.」

其實,我好焦慮今日講道我會越界,所以我向他借來他的講章作參考.他好慷慨,借了給我.謹此多謝他.我拿他的開場白作為我的結束語:

「《提摩太前後書》是保羅寫給提摩太的信,是晚年保羅對年輕教會領袖的叮嚀.

在這情辭逼切的書信中,保羅用了「打仗」來借喻基督徒的生命,囑咐提摩太要「打那美好的仗」(提前一18),「為信仰打那美好的仗」(提前六12).

保羅更留下廣為傳誦的自白,為自己生命作總結:「那美好的仗我已經打過了,當跑的路我已經跑盡了,該信的道我已經守住了……」(提後四7)

為何保羅會用「打仗」來作借喻?當然,標準答案是:基督徒的生命是一場屬靈爭戰,所以我們要full gear,戴上全副軍裝(弗六14-18).當我們說屬靈爭戰時,並不代表這只是一場發生在「靈界」的對抗,正如保羅說:這場爭戰的對象包括:「有血有肉的人」,「那些執政的,掌權的,管轄這幽暗世界的」,還有「天空靈界的惡魔」(弗六12).

其實,不管你是否喜歡,人生如戰場,我們都在打仗.在人生不同階段都有目標,同時也會面對各種困難與挑戰.能否迎向挑戰,或是在挫敗後再次勇敢上路,確是人生的真實寫照.

既然基督徒的生命就是打仗,那麼,我們如何才能像保羅所言,去打一場美好的仗?或者,這場仗,真是美好的嗎?我們真的能打嗎?」

我2018年8月28日在崇基禮拜堂曾經以「難為正邪定分界,也得裝身作靈戰」為題講道,去演繹以弗所書6章10-20節所談及這場戰役,今日好像續集,補充如何打這美好的仗.

保羅是我前輩,叮嚀作為學者的我,小心自己所引以為傲的學問.

所以,我以今日經課沒有包括但是接著的第20至21節來結束,稍作修改,構作為一個醒察心靈的祈禱作結:

20 [ 鄭漢文 ] 啊,你要保守所託付你的,躲避世俗的虛談和那敵真道,似是而非的學問.

21已經有人自稱有這學問,就偏離了真道.

願恩惠常與你們同在!

但願榮耀都歸給聖父,聖子,聖靈,永在的三一神.阿們.



\newpage

\section{我們是無用的僕人}
\label{sec:472}
\textbf{袁天佑牧師 - 我們是無用的僕人}
\newline
\newline
\hyperref[sec:471]{< < < PREV SERMON < < <}
~
\hyperref[sec:toc]{[返目錄]}
~
\hyperref[sec:473]{> > > NEXT SERMON > > >}


路加福音 17:5-10
\newline
\begin{longtable}{cl}
\hline
\hline
章節 & 經文 (和合本修訂版)\\
\hline
17:5 & \begin{tabularx}{0.7\textwidth}{X} 使徒對主說：「請加增我們的信心。」 \end{tabularx} \\ \\ \relax
17:6 & \begin{tabularx}{0.7\textwidth}{X} 主說：「你們若有信心像一粒芥菜種，就是對這棵桑樹說：『你要連根拔起，栽在海裡』，它也會聽從你們。」 \end{tabularx} \\ \\ \relax
17:7 & \begin{tabularx}{0.7\textwidth}{X} 「你們當中誰有僕人耕地或是放羊，從田裡回來，就對他說『你快來坐下吃飯』呢？ \end{tabularx} \\ \\ \relax
17:8 & \begin{tabularx}{0.7\textwidth}{X} 他豈不對僕人說『你給我預備晚飯，束上帶子伺候我，等我吃喝完了，你才可以吃喝』嗎？ \end{tabularx} \\ \\ \relax
17:9 & \begin{tabularx}{0.7\textwidth}{X} 僕人照所吩咐的去做，主人還謝謝他嗎？ \end{tabularx} \\ \\ \relax
17:10 & \begin{tabularx}{0.7\textwidth}{X} 這樣，你們做完了一切所吩咐的，要說：『我們是無用的僕人，所做的本是我們該做的。』」 \end{tabularx} \\ \\
[1ex]
\hline
\hline
\end{longtable}
講題:我們是無用的僕人 We are unworthy servants

講員:袁天佑牧師

經文:路加福音17:5-10

場合:香港中文大學崇基學院禮拜堂主日崇拜

日期:2019年10月6日

我給與大家的講題,大家都知道是出自路加福音十七章5~10節,其中的第10節,僕人做完主人吩咐的工作,除了不能向主人邀功外,他只能說:「我們是無用的僕人,所做的本是我們該做的.」

這句說話,是我過去事奉了40年的一種感受,特別是最後的幾年,甚至是在教會中擔任會長的工作,雖有若干權力,但仍感到自己無用,一無是處.

我想起關俊棠神父所寫的一本書,《步入紅塵》,在最後一篇文章「願作人間雨」,他寫下這段他自己的心路歷程:

「在修道院受訓期間,胸懷大志,要為基督贏取全世界.初出道時,仍雄心勃勃,雖然已開始覺得世界有點太大.其後不斷修訂自己的大志；由世界縮小到中國,再由中國回到香港社會......社會到本地教會......教會到神職班......當要把原來的理想大志,在現實下屢屢不得不作修正時,心中難免激盪出許多的難過,不滿和無奈.最後,悟出一條道理:多年來自己被自己的大志所困.於是不再胸懷大志.」

想起讀神學的初期,覺得自己應該是一個很有本事的人.因40年前,大學畢業,甚至拿取了理學士一級榮譽的學歷,去讀神學,在教會中去擔任牧職,在當時實在是少有.但畢業後,回到教會事奉,實在遇到挫敗之多,實在難以計算.

畢業時,在一間大的堂會學習事奉,在一位長輩牧師指導下工作.他沒有吩咐我作甚麼,也沒有叫我提出甚麼大計,甚至連小的計劃也沒有.他只是叫我,「他做甚麼,你便做甚麼」.兩個最深印象,其中一個是叫我去收拾士多房,看看有那些東西有用便保留,沒有便掉去.另一件事,就是去抄寫教友的紀綠.過去沒有電腦,所以需要每幾年更新抄寫過一篇.心中便這樣想,堂堂大學生,又神學畢業,去收拾士多房,抄寫教友資料,這不是幹事都可以做的嗎?是否大材少用呢?

這樣子,就渾渾噩噩的過了兩年,後來調派到觀塘一間細小禮拜堂去事奉.除了堂會工作外,亦要負責一個社區服務中心的工作.我當然樂意接受這工作,認為可以大展拳腳.

當時的觀塘,可以說是中下階層的社區,不少工廠,所以大部份住在區內的都是勞工階層.他們所面對的,是勞工的權益問題,長時間工作也引起家庭的問題.要為他們解決實在是談何容易!例如勞工權益,大家都知道,經過了幾十年的今天,勞資關係仍很惡劣,最低工資,標準工時,公會的權力,仍是沒完沒了的有衝突.最近更因政治原因,在一間航空公司忠心服務的員工都要被解僱.所以回想在這裡的工作,有甚麼成果呢?

後來又有多年的時間,被調往一些大堂會去工作,服務中上階層的會友,又同時有20多年的時間,負責教會的學校教育工作,管理幾十間中,小學和幼稚園.

首先講講學校教育的工作.不知大家覺得學校教育在過去20多年年的發展是怎樣呢?我個人覺得教育的改變,政治化多於教育原則.例如學校法團化,中學的教育語言改革,國民教育,通識教育的爭議等等,究竟提出這些改革是因教育的原則抑或是政治呢?政治實在比教育的原則更重要,所以今天的學校教育,教育界非常的擔心,但似乎是無法對抗.

又談到教會的牧養工作,我常常問:「究竟每星期的講道,能教化信眾多少?真的聽到你講又願意去實踐的又有多少?要來得到教訓的人多,抑或來訓覺的人多?」當然你可以說,因牧師講道太悶,所以才訓覺.

不過,這是事實.不少基督徒信仰上帝,目的多是希望內心得到平安,生活蒙上帝的祝福,最後可以上到天堂.但基督教的信仰核心是否這樣呢?耶穌說:「人子來,並不是要受人的服事,乃是要服事人,並且要捨命作多人的贖價.」(可十45)耶穌又說:「一粒麥子不落在地裡死了,仍舊是一粒；若是死了,就結出許多子粒來.愛惜自己性命的,就喪失性命；那恨惡自己在這世上的性命的,要保全性命到永生.」(約十二24~25)究竟信仰是推動我們去服事人,抑或是生活成功,受人服事呢?我們的信仰是要我們保守和賺取自己的生命,抑或要願意付出,甚至是犧性?

耶穌又說:「健康的人用不著醫生；有病的人才用得著.我不是來召義人悔改,而是召罪人悔改.」(路五32)最近的幾個月,因發生這個反修例運動,在不同社區中教會開放給示威人士,教會中有人也發出這問題:「教會應否開放給他們,因他們可能是暴徒?」但我想,耶穌為甚麼與稅吏,妓女,甚至是激進黨的西門,與他們同席呢?

牧會差不多40年,能讓多少教友明白這些道理?

回顧過去40年,在社會服務,在學校教育,甚至是堂會牧養,能做到多少呢?所以關俊棠神父所講的那番說話,實在令我非常認同.

不過,我今天來不是要為大家帶來負面的感覺.關神父在他所寫的文章中,他有繼續這樣的分享:

「一天,暮春三月,外邊下著毛毛細雨,似有所悟.......能真正滋潤大地,使萬物欣欣向榮的,並非洶湧澎湃,驚心奪目的急流大川,而是這種不起眼,不爭妍,不澎湃,不傾盆,密密綿綿的絲絲細雨.忽然間,我發覺,我雖未能為基督贏得全世界,我能為中國同胞所做的又是那麼小,我連自己教會內許多不理想的情況都束手無策,但至少有一樣事情是我可以做的,那就是:如果每天這世界上,至少有一個人能因為我的緣故而活得快樂一點,有多點做人的信心,那我也不算白活!從那一天起,我把過去的大志統統除下,為自己換上這份小志,幾年下來,漸漸感到倒也不枉此生!

今天我的志向已不再在要求『改變』甚麼,而是在提醒自己要事事『珍惜』,珍惜自己,也珍惜別人.珍惜,就是不糟蹋自己,也不糟蹋別人,更不去隨便糟蹋自己與別人建立的關係.......珍惜,就是視自己和別人都是上天賜給人間的一份厚禮.『改變』是生命自然流露,有它自己的時間表.我們做的只是為『改變』提供更恰當的環境和時機；我相信,『珍惜』是最積極的準備.」

人有大志,要去改變,改變國家,改變社會,改變教會,可能都是一事無成.但一份的小志和珍惜,可能可以為改變作出準備,最低限度,如果有人因我們的鼓勵,能活得快樂,有信心生活,那豈不是更加重要.

路加福音十七章5節,門徒對耶穌說:「請加增我們的信心.」很多時候,我們會以為信心越大,能力越大,能為上帝做的事越多和越偉大.但這完全是錯誤的,這就正如耶穌所講的「僕人的本分」的比喻那樣,不是僕人可以影響主人,只有主人去吩咐僕人去做事.我們的信心也不影響究竟我們能為上主做甚麼.耶穌說:「你們若有信心像一粒芥菜種,就是對這棵桑樹說:『你要連根拔起,栽在海裡』,它也會聽從你.」(十七6)

只要一粒芥菜種那麼小的,便能將一顆高大的桑樹連根拔起.很小的信心,耶穌從來沒有輕看小的信心,因為這細小的信心,是一顆自覺一無是處的信心,但也是一顆能完全擺上的信心,一顆不是要來操控上主的信心,因為當我們自覺有大的信心時,就會覺得上主應給我做更大的事.

回想當我預備奉獻攻讀神學時,我寫了一篇獻身的文章,題目是「沒有荊棘的火」.相信大家都曉得這題目的背景,就是上主呼召摩西時,摩西看見荊棘,有火但沒有燒著.對於上主對摩西的呼召,摩西這樣回應,他說:「我是甚麼人?竟能去見法老,把以色列人從埃及領出來呢?」(出三11)他又對上主說:「主啊,求求你,我並不是一個能言善道的人,以前這樣,就是你對僕人說話以後也是這樣,因為我是拙口笨舌的.」(出四10)

這其實也是我的感受.大學時,我讀的是化學,對於哲學文學歷史等一方面沒大的興趣,另一方面也覺得很難應付.但神學就是哲學文學歷史的研究.信心和上帝的恩典就是這樣.幾年神學的訓練,後來更有機會去進修,拙口笨舌的我也當了牧師40年的日子,沒有感到這些訓練或是自己的弱點,會成為事奉上帝的攔阻.讀化學再諗神學,或許使我有點「神化」.雖然化學的學問已完全忘記了,但那種有邏輯和有系統的訓練,也成為我諗神學的幫助.

還記得我前面講及在擔任牧師的初期去收拾士多房和抄寫教友紀錄嗎?今天回想起來,這都不是沒有意義的,就正如考古學那樣,都是從垃圾堆中去尋找的.過去前人留下的資料,有時我們覺得沒用,放在士多房中,但當你有心慢慢整理時,你會發現有很多寶貴的歷史資料.另外,抄寫教友資料,令我對他們的認識更多和能記得更多,反之,現在用電腦,就沒有那麼深刻了.

所以不要以為一些微不足道的工作是沒有意義的.只是當你覺得自己是有用時,你就會覺得這些是沒用.但當你覺得自己是無用的僕人時,樂意去做一些微不足道的事,很奇妙的,在上帝手中,會結出你意想不到的果效.豈碼在你自己身上,也會顯出奇妙的恩典.

路加福音十七章10節這樣說:「我們是無用的僕人.」但我想引用哥林多後書十二章9節所記載上主所說的話作總結.上主說:「我的恩典是夠你用的,因為我的能力是在人的軟弱上顯得完全.」

這兩節經文,是我一生事奉得到的體會.



\newpage

\section{我已經看見了主}
\label{sec:473}
\textbf{劉國偉先生 - 我已經看見了主}
\newline
\newline
\hyperref[sec:472]{< < < PREV SERMON < < <}
~
\hyperref[sec:toc]{[返目錄]}
~
\hyperref[sec:474]{> > > NEXT SERMON > > >}


約翰福音 20:1-18
\newline
\begin{longtable}{cl}
\hline
\hline
章節 & 經文 (和合本修訂版)\\
\hline
20:1 & \begin{tabularx}{0.7\textwidth}{X} 七日的第一日清早，天還黑的時候，抹大拉的馬利亞來到墳墓，看見石頭已從墳墓挪開了， \end{tabularx} \\ \\ \relax
20:2 & \begin{tabularx}{0.7\textwidth}{X} 就跑來見西門‧彼得和耶穌所愛的那個門徒，對他們說：「有人從墳墓裡把主移走了，我們不知道他們把他放在哪裡。」 \end{tabularx} \\ \\ \relax
20:3 & \begin{tabularx}{0.7\textwidth}{X} 彼得和那門徒就出來，往墳墓去。 \end{tabularx} \\ \\ \relax
20:4 & \begin{tabularx}{0.7\textwidth}{X} 兩個人同跑，那門徒比彼得跑得快，先到了墳墓， \end{tabularx} \\ \\ \relax
20:5 & \begin{tabularx}{0.7\textwidth}{X} 低頭往裡看，看見細麻布還放在那裡，只是沒有進去。 \end{tabularx} \\ \\ \relax
20:6 & \begin{tabularx}{0.7\textwidth}{X} 西門‧彼得隨後也到了，進了墳墓，看見細麻布放在那裡， \end{tabularx} \\ \\ \relax
20:7 & \begin{tabularx}{0.7\textwidth}{X} 又看見耶穌的裹頭巾沒有和細麻布放在一起，是另在一處捲著。 \end{tabularx} \\ \\ \relax
20:8 & \begin{tabularx}{0.7\textwidth}{X} 然後先到墳墓的那門徒也進去，他看見就信了。 \end{tabularx} \\ \\ \relax
20:9 & \begin{tabularx}{0.7\textwidth}{X} 他們還不明白聖經所說耶穌必須從死人中復活的意思。 \end{tabularx} \\ \\ \relax
20:10 & \begin{tabularx}{0.7\textwidth}{X} 於是兩個門徒回自己的住處去了。 \end{tabularx} \\ \\ \relax
20:11 & \begin{tabularx}{0.7\textwidth}{X} 馬利亞卻站在墳墓外面哭。她哭的時候，低頭往墳墓裡看， \end{tabularx} \\ \\ \relax
20:12 & \begin{tabularx}{0.7\textwidth}{X} 看見兩個天使穿著白衣，在安放耶穌身體的地方坐著，一個在頭，一個在腳。 \end{tabularx} \\ \\ \relax
20:13 & \begin{tabularx}{0.7\textwidth}{X} 天使對她說：「婦人，你為甚麼哭？」她對他們說：「因為有人把我主移走了，我不知道他們把他放在哪裡。」 \end{tabularx} \\ \\ \relax
20:14 & \begin{tabularx}{0.7\textwidth}{X} 說了這些話，她轉過身來，看見耶穌站在那裡，卻不知道他是耶穌。 \end{tabularx} \\ \\ \relax
20:15 & \begin{tabularx}{0.7\textwidth}{X} 耶穌問她：「婦人，你為甚麼哭？你找誰？」馬利亞以為他是看園子的，就對他說：「先生，若是你把他移了去，請告訴我，你把他放在哪裡，我去把他移回來。」 \end{tabularx} \\ \\ \relax
20:16 & \begin{tabularx}{0.7\textwidth}{X} 耶穌對她說：「馬利亞。」馬利亞轉過身來，用希伯來話對他說：「拉波尼！」（「拉波尼」就是老師的意思。） \end{tabularx} \\ \\ \relax
20:17 & \begin{tabularx}{0.7\textwidth}{X} 耶穌對她說：「不要拉住我，因為我還沒有升上去見我的父。你到我弟兄那裡去告訴他們，我要升上去見我的父，也是你們的父，見我的神，也是你們的神。」 \end{tabularx} \\ \\ \relax
20:18 & \begin{tabularx}{0.7\textwidth}{X} 抹大拉的馬利亞就向門徒報信：「我已經看見了主。」她又把主對她說的話告訴他們。 \end{tabularx} \\ \\
[1ex]
\hline
\hline
\end{longtable}
講題:我已經看見了主 I have seen the Lord

經文:約翰福音20:1-18

講員:劉國偉先生

場合:香港中文大學崇基學院禮拜堂主日崇拜

日期:2019年10月13日

引言

弟兄姊妹平安!

今天的講題是「我已經看見了主」,這句極有歷史意義的說話,記載於約翰福音20章18節,出自跟隨耶穌基督的一位婦女,抹大拉的馬利亞.她是福音書作者所見證,第一位看見復活的基督的人.

天主教禮儀及聖事部遵循教宗方濟各的意願,於2016年6月10日公布法令,將羅馬禮儀年曆於7月22日慶祝的聖女抹大拉的馬利亞(瑪利亞瑪達肋納)紀念日,提升為慶日,與其他宗徒地位同等.抹大拉的馬利亞的慶日日期不變,仍是7月22日.按天主教的禮儀年曆,慶日(festum)比紀念日(memoria)更為教會所重視.

天主教聖座禮儀聖事部在頒布這道法令時表示,聖女瑪利亞瑪達肋納是復活基督的首位見證人,也是向門徒報告主耶穌復活消息的首位使者.將聖女瑪利亞瑪達肋納的紀念日提升為慶日,強調了教會的福傳工作,就是不分男女,種族,人民,語言,國家,一律予以接納,向他們宣講福音喜訊,陪伴他們的塵世旅程,與他們分享天主救恩的奇妙,闡明這位聖女愛耶穌並被耶穌所愛的重要意義.

瑪利亞瑪達肋納在死亡之地宣布復活的生命,在羅馬禮儀年曆中將聖女瑪利亞瑪達肋納的瞻禮作為慶日,與其他宗徒有同等地位是合情合理的,進一步彰顯了這位教會女性楷模的特殊使命.

提起這位重要的門徒,我就想起本年5月與基督徒領袖訓練計劃的同學,一起在以色列加利利遊歷的片段,我們用一整個上午從拿撒勒出發,徒步踏上耶穌路徑(Jesus Trail),在攝氏40多度高溫之下,穿越半曠野路程,抵達加利利湖西岸,靠近抹大拉的地方.相對而言,今天的抹大拉的社經地位,已難以與第一世紀的抹大拉相比.今天的抹大拉是一個藉著考古遺跡發展旅遊經濟的地方,而馬利亞時期的抹大拉,則是一個優良的漁港,漁業發達,位處南北走向的國際大道,曾經是加利利湖西岸最大之城市.

抹大拉的馬利亞

福音書共有13節經文提及抹大拉的馬利亞,分別是馬太福音3次,馬可福音5次,路加福音兩次及約翰福音3次,亦是較為罕有四卷福音書都有提及的女性.

將四福音合併整理,我們可以按時序,得到以下的敘事結論.

路加福音8章2 節記載: 「還有被惡鬼所附,被疾病所累,已經治好的幾個婦女,內中有稱為抹大拉的馬利亞(曾有七個鬼從她身上趕出來)」,經文接著說: 「…和好些別的婦女,都是用自己的財物供給耶穌和門徒.」自此,抹大拉的馬利亞大概就跟隨耶穌基督的團隊,走遍各城各鄉傳講神國的福音.

抹大拉的馬利亞離開她熟悉的加利利,與其他門徒一起,跟從耶穌南下耶路撒冷及週邊地區,直至遇上基督被拘捕,審訊,判刑,釘十字架.當耶在各各他受刑的時候,馬太如此描寫: 「內中有抹大拉的馬利亞,又有雅各和約西的母親馬利亞,並有西庇太兩個兒子的母親.」(太27:56) 馬可(可15:40)和約翰(約19:25)都有類似的記載.

當大多數門徒因耶穌被捕已經四處躲藏,抹大拉的馬利亞始終身處現場.馬太告訴我們,耶穌安葬的時候,有抹大拉的馬利亞和那個馬利亞在那裡,對著墳墓坐著.(太27:61) 馬可同樣寫道: 抹大拉的馬利亞和約西的母親馬利亞都看見安放他的地方.(可15:47)

七日的頭一日,主基督復活的日子,馬太寫道: 安息日將盡,七日的頭一日,天快亮的時候,抹大拉的馬利亞和那個馬利亞來看墳墓.(太28:1)馬可也有類似的記載: 過了安息日,抹大拉的馬利亞和雅各的母親馬利亞並撒羅米,買了香膏要去膏耶穌的身體.(可16:1)

然後,馬可告訴我們: 在七日的第一日清早,耶穌復活了,就先向抹大拉的馬利亞顯現(耶穌從她身上曾趕出七個鬼).(可16:9)

安息日之後的黎明

約翰告訴我們: 七日的第一日清早,天還黑的時候,抹大拉的馬利亞來到墳墓那裡,看見石頭從墳墓挪開了,就跑來見西門‧彼得和耶穌所愛的那個門徒,對他們說:「有人把主從墳墓裡挪了去,我們不知道放在哪裡.」 (約20:1-2)

一直身處現場的馬利亞,在安息日完結,夜裡四更天的時候,急切來到安葬耶穌的墳墓,她一心想再到現場,並且帶著香膏,為想膏抹所愛的主,幾名婦女可能正思量如何輥開墓門的巨石,誰知她們急切的心情,卻轉為驚恐.(可16:5)

鏡頭轉至彼得和約翰:  彼得和那門徒就出來,往墳墓那裡去.兩個人同跑,那門徒比彼得跑得更快,先到了墳墓,低頭往裡看,就見細麻布還放在那裡,只是沒有進去.西門‧彼得隨後也到了,進墳墓裡去,就看見細麻布還放在那裡,又看見耶穌的裹頭巾沒有和細麻布放在一處,是另在一處捲著.先到墳墓的那門徒也進去,看見就信了.(因為他們還不明白聖經的意思,就是耶穌必要從死裡復活.)於是兩個門徒回自己的住處去了. (約20:3-10)

彼得和約翰親眼目睹這個空墳墓,還有裹屍布與包頭巾安放在原來放置耶穌身體的地方.十分明顯,偷去遺體所呈現的結果,不可能是眼前這個模樣.所以約翰乾脆告訴讀者,他看見就信,耶穌確如祂自己曾經預言過,祂死後會復活過來,然而,這眼前事實太新奇了,一時間仍未明白,需要消化,他們就回家去了.

耶穌向抹大拉的馬利亞顯現

抹大拉的馬利亞並沒有離去,按上下文所呈現,抹大拉的馬利亞很可能是當時唯一留下來的門徒.約翰寫道: 馬利亞卻站在墳墓外面哭.哭的時候,低頭往墳墓裡看,就見兩個天使,穿著白衣,在安放耶穌身體的地方坐著,一個在頭,一個在腳.天使對她說:「婦人,你為甚麼哭?」她說:「因為有人把我主挪了去,我不知道放在哪裡.」(約20:11-13)

馬利亞顯然以為有人挪去耶穌的身體,她的眼淚是我們不難明白的,因為那是雙重的失去: 死了,卻連遺體也不見了.對馬利亞來說,這是一個充滿問號的現場,這三天三夜發生的事,超乎經驗,超乎理解,也超乎想像.那一位滿懷理想,宣揚真理,指斥公會和宗教領袖,又把馬利亞從魔鬼的黑暗權勢中,釋放出來的主耶穌,竟然被殺害,卻又連遺體都被人搬走了.

正當馬利亞這位深愛著耶穌,身體力行地跟隨耶穌,耶穌在那裡,她就設法在那裡,如今卻身陷情緒困境之時,耶穌向她顯現.約翰以動人的筆觸寫下來: 說了這話,就轉過身來,看見耶穌站在那裡,卻不知道是耶穌.耶穌問她說:「婦人,為甚麼哭?你找誰呢?」馬利亞以為是看園的,就對他說:「先生,若是你把他移了去,請告訴我,你把他放在哪裡,我便去取他.」耶穌說:「馬利亞.」馬利亞就轉過來,用希伯來話對他說:「拉波尼!」(拉波尼就是夫子的意思.)耶穌說:「不要摸我,因我還沒有升上去見我的父.你往我弟兄那裡去,告訴他們說,我要升上去見我的父,也是你們的父,見我的上帝,也是你們的上帝.」 (約20:14-17)

弟兄姊妹,若果耶穌問你「你為甚麼哭?你找誰呢?」,此時當下的你,又會怎答祂? 若果我們哭,我們為什麼哭? 我們知道馬利亞因為不見了耶穌而哭,她正在找耶穌,但她要找的是她熟悉的人子耶穌,然而,基督復活了,祂並沒有失蹤,祂已經勝過死亡,滿有復活的大能,開通了人神之間可以重新契合的道路.耶穌也提醒馬利亞不要拉住祂,耶穌鼓勵這位愛慕及渴望與人子耶穌一起的馬利亞起來行動,天亮了,救恩已經完成,從這黑暗的猶太地開始,救恩要傳開去,傳到普天下.

抹大拉的馬利亞就去告訴門徒說:「我已經看見了主.」她又將主對她說的這話告訴他們. (約20:18)

親愛的弟兄姊妹,走人生窄路,特別是夜裡四更天的時候,覺得耶穌不見了,失蹤了,有時難免會問: 上帝你在那裡?  要面對這種裡裡外外的情境,實在並不容易,標準答案很容易就可以說出來,然而,只有真心去愛,真誠地跟從耶穌走進生命現場的人,才可以牢牢抓住應許,看見耶穌始終在我們當中.

我特別喜歡天主教聖座禮儀聖事部稱許抹大拉的馬利亞: 在死亡之地宣布復活的生命.親愛的每一位,但願我們即使在很大很大的困境中,無論是個人的,社會的,普世的,我們總有力量在耶穌的空墳墓前宣告:「我已經看見了主.」

復活的主與我們同在,也與香港同在.



\newpage

\section{渡口}
\label{sec:474}
\textbf{鄧瑞強博士 - 渡口}
\newline
\newline
\hyperref[sec:473]{< < < PREV SERMON < < <}
~
\hyperref[sec:toc]{[返目錄]}
~
\hyperref[sec:475]{> > > NEXT SERMON > > >}


創世記 32:6-33:11
\newline
\begin{longtable}{cl}
\hline
\hline
章節 & 經文 (和合本修訂版)\\
\hline
32:6 & \begin{tabularx}{0.7\textwidth}{X} 使者回到雅各那裡，說：「我們到了你哥哥以掃那裡。他正迎著你來，並且有四百人和他一起。」 \end{tabularx} \\ \\ \relax
32:7 & \begin{tabularx}{0.7\textwidth}{X} 雅各就很懼怕，而且愁煩。他把跟他同行的人和羊群、牛群、駱駝分成兩隊， \end{tabularx} \\ \\ \relax
32:8 & \begin{tabularx}{0.7\textwidth}{X} 說：「以掃若來擊殺其中一隊，剩下的另一隊還可以逃脫。」 \end{tabularx} \\ \\ \relax
32:9 & \begin{tabularx}{0.7\textwidth}{X} 雅各說：「耶和華—我祖父亞伯拉罕的神，我父親以撒的神啊，你曾對我說：『回你本地本族去，我要厚待你。』 \end{tabularx} \\ \\ \relax
32:10 & \begin{tabularx}{0.7\textwidth}{X} 你向僕人所施的一切慈愛和信實，我一點也不配得。我先前只用我的一根杖過這約旦河，如今我卻成了兩隊。 \end{tabularx} \\ \\ \relax
32:11 & \begin{tabularx}{0.7\textwidth}{X} 求你救我脫離我哥哥的手，脫離以掃的手，因為我怕他來殺我，連母親和兒女都不放過。 \end{tabularx} \\ \\ \relax
32:12 & \begin{tabularx}{0.7\textwidth}{X} 你曾說：『我必定厚待你，使你的後裔如同海邊的沙，多得不可勝數。』」 \end{tabularx} \\ \\ \relax
32:13 & \begin{tabularx}{0.7\textwidth}{X} 當夜，雅各在那裡住宿，就從他手中所擁有的拿禮物要送給他哥哥以掃， \end{tabularx} \\ \\ \relax
32:14 & \begin{tabularx}{0.7\textwidth}{X} 就是二百隻母山羊、二十隻公山羊、二百隻母綿羊、二十隻公綿羊、 \end{tabularx} \\ \\ \relax
32:15 & \begin{tabularx}{0.7\textwidth}{X} 三十匹哺乳的母駱駝和牠們的小駱駝、四十頭母牛、十頭公牛、二十匹母驢和十匹公驢。 \end{tabularx} \\ \\ \relax
32:16 & \begin{tabularx}{0.7\textwidth}{X} 他把每種牲畜各分一群，交在僕人手中，對僕人說：「你們要在我的前頭過去，使群和群之間保持一段距離」。 \end{tabularx} \\ \\ \relax
32:17 & \begin{tabularx}{0.7\textwidth}{X} 他又吩咐領頭的人說：「我哥哥以掃遇見你的時候，問你說：『你是誰的人？要往哪裡去？你前面這些是誰的？』 \end{tabularx} \\ \\ \relax
32:18 & \begin{tabularx}{0.7\textwidth}{X} 你就說：『是你僕人雅各的，是送給我主以掃的禮物。看哪，他自己也在我們後面。』」 \end{tabularx} \\ \\ \relax
32:19 & \begin{tabularx}{0.7\textwidth}{X} 他又吩咐第二、第三和所有趕畜群的人說：「你們遇見以掃的時候要照這樣的話對他說， \end{tabularx} \\ \\ \relax
32:20 & \begin{tabularx}{0.7\textwidth}{X} 你們還要說：『看哪，你僕人雅各在我們後面。』」因雅各說：「我藉著在我前面送去的禮物給他面子，然後再見他的面，或許他會寬容我。」 \end{tabularx} \\ \\ \relax
32:21 & \begin{tabularx}{0.7\textwidth}{X} 於是禮物在他前面過去了；那夜，雅各在營中住宿。 \end{tabularx} \\ \\ \relax
32:22 & \begin{tabularx}{0.7\textwidth}{X} 他夜間起來，帶著兩個妻子，兩個婢女和十一個孩子，過了雅博渡口。 \end{tabularx} \\ \\ \relax
32:23 & \begin{tabularx}{0.7\textwidth}{X} 他帶著他們，送他們過河，他所有的一切也都過去， \end{tabularx} \\ \\ \relax
32:24 & \begin{tabularx}{0.7\textwidth}{X} 只剩下雅各一人。有一個人來和他摔跤，直到黎明。 \end{tabularx} \\ \\ \relax
32:25 & \begin{tabularx}{0.7\textwidth}{X} 那人見自己勝不過他，就摸了他的大腿窩一下。雅各的大腿窩就在和那人摔跤的時候扭了。 \end{tabularx} \\ \\ \relax
32:26 & \begin{tabularx}{0.7\textwidth}{X} 那人說：「天快亮了，讓我走吧！」雅各說：「你不給我祝福，我就不讓你走。」 \end{tabularx} \\ \\ \relax
32:27 & \begin{tabularx}{0.7\textwidth}{X} 那人說：「你叫甚麼名字？」他說：「雅各。」 \end{tabularx} \\ \\ \relax
32:28 & \begin{tabularx}{0.7\textwidth}{X} 那人說：「你的名字不要再叫雅各，要叫以色列，因為你與神和人較力，都得勝了。」 \end{tabularx} \\ \\ \relax
32:29 & \begin{tabularx}{0.7\textwidth}{X} 雅各問他說：「請告訴我你的名字。」那人說：「何必問我的名字呢？」於是他在那裡為雅各祝福。 \end{tabularx} \\ \\ \relax
32:30 & \begin{tabularx}{0.7\textwidth}{X} 雅各就給那地方起名叫毗努伊勒 ，說：「我面對面見了神，我的性命仍得保全。」 \end{tabularx} \\ \\ \relax
32:31 & \begin{tabularx}{0.7\textwidth}{X} 太陽剛出來的時候，雅各經過毗努伊勒，他的大腿就瘸了。 \end{tabularx} \\ \\ \relax
32:32 & \begin{tabularx}{0.7\textwidth}{X} 因此，以色列人不吃大腿窩的筋，直到今日，因為那人摸了雅各大腿窩的筋。 \end{tabularx} \\ \\ \relax
33:1 & \begin{tabularx}{0.7\textwidth}{X} 雅各舉目觀看，看哪，以掃來了，有四百人和他一起。雅各就把孩子們分開交給利亞、拉結和兩個婢女。 \end{tabularx} \\ \\ \relax
33:2 & \begin{tabularx}{0.7\textwidth}{X} 他叫兩個婢女和她們的孩子走在前頭，利亞和她的孩子跟在後面，而拉結和約瑟在最後。 \end{tabularx} \\ \\ \relax
33:3 & \begin{tabularx}{0.7\textwidth}{X} 他自己卻走到他們前面，一連七次俯伏在地才挨近他哥哥。 \end{tabularx} \\ \\ \relax
33:4 & \begin{tabularx}{0.7\textwidth}{X} 以掃跑來迎接他，將他抱住，伏在他的頸項上親他，他們都哭了。 \end{tabularx} \\ \\ \relax
33:5 & \begin{tabularx}{0.7\textwidth}{X} 以掃舉目看見婦人和孩子，就說：「這些和你一起的是誰呢？」雅各說：「這些孩子是神施恩給你僕人的。」 \end{tabularx} \\ \\ \relax
33:6 & \begin{tabularx}{0.7\textwidth}{X} 於是兩個婢女和她們的孩子前來下拜， \end{tabularx} \\ \\ \relax
33:7 & \begin{tabularx}{0.7\textwidth}{X} 利亞和她的孩子也前來下拜，隨後約瑟和拉結也前來下拜。 \end{tabularx} \\ \\ \relax
33:8 & \begin{tabularx}{0.7\textwidth}{X} 以掃說：「我所遇見的這些畜群是甚麼意思呢？」雅各說：「是為了要在我主眼前蒙恩。」 \end{tabularx} \\ \\ \relax
33:9 & \begin{tabularx}{0.7\textwidth}{X} 以掃說：「弟弟啊，我的已經夠了，你的你自己留著吧！」 \end{tabularx} \\ \\ \relax
33:10 & \begin{tabularx}{0.7\textwidth}{X} 雅各說：「不，我若在你眼前蒙恩，就請你從我手裡收下這禮物；因為我見了你的面，如同見了神的面，並且你也寬容了我。 \end{tabularx} \\ \\ \relax
33:11 & \begin{tabularx}{0.7\textwidth}{X} 請你收下我帶來給你的禮物，因為神恩待我，使我一切都充足。」雅各再三求他，他才收下。 \end{tabularx} \\ \\
[1ex]
\hline
\hline
\end{longtable}
講題:渡口 The Ford Jabbok

講員:鄧瑞強博士

經文:創世記32:6-33:11

場合:香港中文大學崇基學院禮拜堂主日崇拜

日期:2019年10月20日

各位弟兄姊妹,早安.

十六世紀,有位靈修大師,一般稱他為「十架若望」(John of the Cross),寫了本靈修學的書,書名叫《心靈的黑夜》(Dark Night of the Soul).書的第二部分第五章講到,人越接近光明,眼睛越注目光明,越看不見東西,因為我們的眼睛忍受不到這種光明.面對強烈的光明,我們看不到外在的世界.就此而論,我們像盲了一樣,像在黑暗中摸索一樣.另一方面,人越接近光明,越體會自己內心的黑暗,越體會自己的無知與軟弱,越體會自己的無助,越認識到過去所倚靠的一切東西的不可倚靠,就此而論,他的內心也墮入黑暗,對曾經熟悉的世界不再熟悉.這種外在與內在的黑暗狀態,「十架若望」稱之為「心靈的黑夜」.

只要我們活得足夠久,忍受過足夠的苦難,思想上有足夠的深度,對現實的空虛有足夠的敏感,對真理有深度的渴望,對人生意義有足夠的尋索,或多或少,或遲或早,我們都會體會這種「心靈的黑夜」.

那一刻,你會忽然間覺得,整個熟悉的世界變得不再熟悉,平時追逐的意義忽然失去意義,不知生存下去的意義何在,想做點什麼卻又發現自己的無能與無助.回望過去,很懷疑過去的做法是否正確.想望將來,發覺明天的事情完全在自己的掌控之外.看看當下的自己,很迷茫,很陌生.想抓著一點什麼,卻發覺一切都不可靠.只要你的心靈純潔,奮勇向上,越迷茫的時候,便會越加渴求意義.越空洞,便越嚮往神聖.就在此時,在黑暗中,隱隱約約,恍恍惚惚,會體會神聖的臨在.唯有當一切次一等的意義都放下時,最終極的意義才出現.唯有當生活裡不是最重要的東西放下時,最重要的東西才浮上來.唯有當塵世的一切都顯得不可靠時,才能體會上帝才是最可靠的.但上帝在何方,卻仍晦暗不明.唯有在黑暗中,堅決前行.總有一天,去到生命的轉折點,去到旅程的拐彎處,去到人生的渡口,然後,以上天的恩典,以人的信心跳躍,走過去,便體會到新的世界,新的生命,新的意義.

我們看看今日舊約經課的一個故事,以色列人的先祖雅各,來到他心靈的黑夜,來到命運的轉折點,來到一個決定性的渡口.這一晚,過去見不得人的墮落,未來不能掌控的危機,面對上帝時不能隱藏的罪行,面對手足時無法啟齒的卑污,前塵往事,一一湧上心頭.他如何過渡這個渡口呢?

未讀這個故事,交代少許背景.

以色列民族的始祖亞伯拉罕生以撒,以撒生了一對孖生子,大哥叫以掃,弟弟叫雅各.「以掃」這個名字的意思是「多毛」,生在香港,可能就叫「長毛」.「雅各」這個名字的意思是「抓住,捉住」,他一出生,便抓住哥哥的腳.

雅各的一生,真的是抓得就抓,搶得便搶,騙得就騙.他首先騙了哥哥的長子名分,後來又騙了父親給哥哥的祝福.因著欺騙哥哥,哥哥老羞成怒,動了殺機.於是,雅各便要亡命天涯,遠走他方.在離鄉別井的漫長歲月裡,輪到雅各被人欺騙,受盡欺凌.二十年後,飽經風霜的雅各,想回老家了.問題是:哥哥還記著過去的一切嗎?還有仇恨嗎?還有殺機嗎?他不知道,他起程回家了.

雅各派僕人作先頭部隊,打個招呼,看看光景如何?

創世記32:6-8如此記載:

創32:6　所打發的人回到雅各那裡,說:「我們到了你哥哥以掃那裡,他帶著四百人,正迎著你來.」

創32:7　雅各就甚懼怕,而且愁煩,便把那與他同在的人口和羊群,牛群,駱駝分做兩隊,

創32:8　說:「以掃若來擊殺這一隊,剩下的那一隊還可以逃避.」

這是叫人提心吊膽的,以掃帶著四百人走過來,來者不善.若只是兄弟相見,幹嗎要帶四百人來.

雅各的策略是叫僕人將一批又一批的牲畜送過去,作為禮物,自己走在後頭,看看光景如何?

創世記32:20下-21如此記載:

創32:20下　……因雅各心裡說:「我藉著在我前頭去的禮物解他的恨,然後再見他的面,或者他容納我.」

創32:21　於是禮物先過去了；那夜,雅各在隊中住宿.

我要讀以上兩節經文,因這兩節經文對於理解今日講道的主要經文有重要作用.這兩節短短的經文,在原文中出現了五次「面」(face)這個字.

照原文字面直譯,這兩節經文是這樣的:

創32:20下　……他〔雅各〕說:「我去安撫他的面,以禮物,在我面前送去的.然後,我見他的面,或者,他會抬起我的面.」

創32:21　於是,禮物送過去了,在他面前；而他住宿,那晚,在營中.

用字用得很妙,五個「面」字,他的面,我的面,他的面,我的面,最後是旁觀者的描述:「他〔雅各〕的面」.伏筆是:「以掃的面又將如何?」經文藉著文字的編織,形象地表達兩張面的相遇,到最後,面對面時,將是一幅什麼樣的圖畫?是一張面消滅另一張面?抑或,是一張面擁吻另一張面?這是懸而未決的.在這裡,留下一思想謎團:這將是一種什麼樣的「面對面」?

過了這一晚,明天,就會面對哥哥的面了.那是異常漫長的一晚.

今日講道的主要經文:創世記32:22-32

創32:22　他夜間起來,帶著兩個妻子,兩個使女,並十一個兒子,都過了雅博渡口,

創32:23　先打發他們過河,又打發所有的都過去,

創32:24　只剩下雅各一人.有一個人來和他摔跤,直到黎明.

創32:25　那人見自己勝不過他,就將他的大腿窩摸了一把,雅各的大腿窩正在摔跤的時候就扭了.

創32:26　那人說:「天黎明了,容我去吧!」雅各說:「你不給我祝福,我就不容你去.」

創32:27　那人說:「你名叫什麼?」他說:「我名叫雅各.」

創32:28　那人說:「你的名不要再叫雅各,要叫以色列；因為你與神與人較力,都得了勝.」

創32:29　雅各問他說:「請將你的名告訴我.」那人說:「何必問我的名?」於是在那裡給雅各祝福.

創32:30　雅各便給那地方起名叫毗努伊勒〔就是神之面的意思〕,意思說:「我面對面見了神,我的性命仍得保全.」

創32:31　日頭剛出來的時候,雅各經過毗努伊勒,他的大腿就瘸了.

創32:32　故此,以色列人不吃大腿窩的筋,直到今日,因為那人摸了雅各大腿窩的筋.

那一晚,在這人生的渡口,在這心靈的黑夜,發生了什麼事?

看看這位雅各,在這渡口,在這黑夜,他面臨最深刻的危機.他可能會死,可能會失去多年來用奸詐騙回來的財富,可能會失去多年來辛苦經營的家庭.「有」與「無」,只是片刻之事.這危機從何而來?從他過去一生做人的方式而來,從他過去對待兄弟的手法而來.他一生都在「抓」,抓得多少是多少.面對手足,毫無親情.手足只是欺騙的對象.他不擇手段,要得到他要得的東西.但這得到的一切,能永恆嗎?就因他過去種的「因」,造成今日的危機的「果」.

在這一晚,他應深深懊悔他過去對待親人的方式,他過去抓取世物的方式,他過去活著的方式,或者說,他應全面懷疑他過去的存在,他應懷疑他過去賴以存在的意義結構,他應懷疑他過去對他人,對世界,對自己的看法.熟悉的世界變得陌生,過去一切可靠的東西現在變得不再可靠,一個什麼都抓住的人現在什麼都不能抓住.明天呢?明天不可捉摸,不能想像,無法控制.明天,是一片虛無,在人類的計算之外.當下,是一片黑暗,毫無光明,無邊的黑夜.漫漫長夜,往事湧上心頭,又無聲遠去.際此良夜,能說什麼呢?無比空洞,無比寂寞,無比惆悵.

那一晚,有一場心靈的交戰.

請不要太快將之理解為「神與雅各摔交」或「天使與雅各摔交」.

創32:24　只剩下雅各一人.有一個人來和他摔跤,直到黎明.

經文刻意講,「有一個人」來和雅各摔跤.

創32:25　那人見自己勝不過他…….

上帝不會勝不過人.

創32:26　那人說:「天黎明了,容我去吧!」雅各說:「你不給我祝福,我就不容你去.」

上帝不會害怕天光,只有夢才會在天光中消失.

所以,首先,我們應理解這是一個「人」,這「人」在夢境一樣的背景中出現.雅各和這人糾纏.雅各看來比他頑強,但這人又明顯能傷害雅各.

創32:25　那人見自己勝不過他,就將他的大腿窩摸了一把,雅各的大腿窩正在摔跤的時候就扭了.

這人讓雅各從今以後瘸腿了.瘸腿,就不能快,不能搶了.香港人說:「執輸行頭,慘過敗家」,從今以後,雅各將不能和人搶了.這一晚的糾纏,將雅各的做人方式徹底地改變了.一個靠搶的人,從今以後,只能謙和地行走了.雅各要和這人糾纏下去,便得付起這個代價.雅各要在這人身上得到什麼呢?他要的,是這人的祝福.這人最後祝福了雅各.

創32:30　雅各便給那地方起名叫毗努伊勒〔就是神之面的意思〕,意思說:「我面對面見了神,我的性命仍得保全.」

在這節經文裡,出現「面對面」這字眼.之前,我們講過,那個晚上,在雅各心中,有一個大問題.「他的面,我的面,他的面,我的面」,然後,我的面,再然後,將會是「他的面」,雅各與哥哥面對面時,將會如何?「面對面」這主題,在這裡再出現.在這心靈交戰中,雅各發現,「面對面」時,他的性命竟然得到保全.

當然,我知道,經文這裡是寫「我面對面見了神」,不是見到以掃.亦因經文這樣寫,一般人理解雅各這晚的經歷,是神人交戰.要理解這經文,我們要多看幾節經文.

創世記33:3-4, 8-11:

創33:3　他自己在他們前頭過去,一連七次俯伏在地才就近他哥哥.

創33:4　以掃跑來迎接他,將他抱住,又摟著他的頸項,與他親嘴,兩個人就哭了.……

創33:8　以掃說:「我所遇見的這些畜群是什麼意思呢?」雅各說:「是要在我主面前蒙恩的.」

創33:9　以掃說:「兄弟啊,我的已經夠了,你的仍歸你吧!」

創33:10　雅各說:「不然,我若在你眼前蒙恩,就求你從我手裡收下這禮物；因為我見了你的面,如同見了神的面,並且你容納了我.

創33:11　求你收下我帶來給你的禮物；因為神恩待我,使我充足.」雅各再三地求他,他才收下了.

看看雅各的身體恣勢,面對哥哥的態度.

創33:3　他自己在他們前頭過去,一連七次俯伏在地才就近他哥哥.

就是這個態度,這種示弱的態度,這種瘸腿的態度,在他與哥哥的生命糾纏中,他贏了他的哥哥,他贏回親情,他贏回自己的生命.

聖經中最動人的一幕.

創33:4　以掃跑來迎接他,將他抱住,又摟著他的頸項,與他親嘴,兩個人就哭了.

可能出現的悲劇,竟然變成喜劇.有點超現實,有點不可思議.問題在於:在那黑夜,在那危機前夕,在那生命的渡口,雅各作出了什麼改變?只要有真正的生命銳變,仇恨就能化解,黑夜就會消失,敵人擁抱的畫面就會出現.

不單雅各的身體姿勢有改變,他面對哥哥時的角色扮演也有所改變.

創33:8　以掃說:「我所遇見的這些畜群是什麼意思呢?」雅各說:「是要在我主面前蒙恩的.」

創33:9　以掃說:「兄弟啊,我的已經夠了,你的仍歸你吧!」

創33:10　雅各說:「不然,我若在你眼前蒙恩,就求你從我手裡收下這禮物；因為我見了你的面,如同見了神的面,並且你容納了我.

創33:11　求你收下我帶來給你的禮物；因為神恩待我,使我充足.」雅各再三地求他,他才收下了.

以前,雅各搶奪父親給哥哥的祝福；現在,他主動向哥哥獻上禮物.掠奪者,變成施予者.哥哥不要這些禮物,雅各也要「再三地求他」收下.雅各再與哥哥糾纏,他要求哥哥接納他的禮物,接納他這個人,祝福他.

特別留意,在這裡,雅各說:「我見了你的面,如同見了神的面,並且你容納了我」.這經文似曾相識,面對面,見到神,生命卻仍存留.在雅各心中,哥哥的面,如同神的面,同樣的難於理解,難於捉摸.面對面時,可能是殺身之禍,也可能是愛的復原.若果,在雅各眼中,他哥哥的面,如同神的面,則我們可以理解,那個晚上,雅各是與夢魘中的哥哥糾纏.他哥哥能傷害他,但他堅持要的,是哥哥祝福他.如何能得到哥哥的祝福?就是生命的全然改變.以前是欺負哥哥的,現在是在哥哥面前俯伏在地.以前是搶的,現在是獻上禮物的.這一改變,決定於那一晚的心靈黑夜,從靈修而言,這同時是人在神面前的生命抉擇,抉擇是否願意生命全然更新.在這意義上,那一晚的糾纏,也是人面對神的一種經驗.面對面的,面對哥哥,同時,也是面對面的,面對神.面對恨自己的人,如何面對,如何化解,這是一種深刻的靈性體會,也是面對面的見到神的經驗.

那一晚,雅各徹底地變了,覺今是而昨非.

創32:28　那人說:「你的名不要再叫雅各,要叫以色列；因為你與神與人較力,都得了勝.」

那人給他改了一個新的名字,以後不叫「雅各」了,不叫「抓住」了,叫「以色列」,這名字的意義是「上帝勝利」.從今以後,雅各能在神面前,在人面前,都得勝,不再是憑自己的力去爭,而是讓上帝得勝.當一個人心中有上帝,讓上帝勝過自己時,他才算是真正勝利的人.以色列,即「上帝得勝」,這名字真好,我們需要有更多擁有這名字的人.

在香港,有很多要爭天下的人,有很多年年考第一的人,他們都是早期的雅各,雅各的名字就是「抓住,爭取」的意思.這個雅各,欺凌了哥哥,搞亂了一個家,造成無法修補的仇恨.但我相信,我相信,這無數的雅各,總有一晚,夜闌人靜,會想起前塵往事,會有心靈交戰.這是心靈的黑夜,這是生命的拐彎處,這是生命超越的渡口.想想,只要瘸下腿來,只要堅持得到對方的祝福,生命就會煥然一新.這同時是一個面對神的渡口.只要大膽放下自己,只要大膽說「願神得勝」,就可以不叫「雅各」,叫「以色列」,這是神子民的名號.我渴望,我渴望,見到的,不是仇殺,而是仇敵的擁抱,是仇敵的親吻,是新天新地.

願上帝得勝.

願:榮耀歸於聖父,聖子,聖靈.阿們.

(附記:本人在崇基禮拜堂講道的講章已結集成《荷花池畔春草綠》講章集,有興趣的朋友可到各大基督教書局購買.試讀本可參閱:\url{http://cc.cuhk.edu.hk/dev/images/content/fund-raising-campaigns/university-sermons-by-dr-sk-tang/CCC-SK-Tang-Sermon-Text-Sample-2.pdf}.較早期的崇基講章,已結集成《敢於跟隨主》,也可到各大基督教書局購買.)



\newpage

\section{這麼近那麼遠}
\label{sec:475}
\textbf{林豪恩先生 - 這麼近那麼遠}
\newline
\newline
\hyperref[sec:474]{< < < PREV SERMON < < <}
~
\hyperref[sec:toc]{[返目錄]}
~
\hyperref[sec:476]{> > > NEXT SERMON > > >}


耶利米書 14:4-22
\newline
\begin{longtable}{cl}
\hline
\hline
章節 & 經文 (和合本修訂版)\\
\hline
14:4 & \begin{tabularx}{0.7\textwidth}{X} 因為無雨降在地上，土地就乾裂，農夫為此蒙羞抱頭。 \end{tabularx} \\ \\ \relax
14:5 & \begin{tabularx}{0.7\textwidth}{X} 田野的母鹿因為無草也撇棄才生的小鹿。 \end{tabularx} \\ \\ \relax
14:6 & \begin{tabularx}{0.7\textwidth}{X} 野驢站在光禿的高地喘氣，好像野狗；牠們的眼目因無草而失明。」 \end{tabularx} \\ \\ \relax
14:7 & \begin{tabularx}{0.7\textwidth}{X} 耶和華啊，雖然我們的罪孽控告我們，求你為你名的緣故行動吧！我們本是多次背道，得罪了你。 \end{tabularx} \\ \\ \relax
14:8 & \begin{tabularx}{0.7\textwidth}{X} 以色列所盼望，在患難時作他救主的啊，你在這地為何像寄居的，又如旅行的只住一夜呢？ \end{tabularx} \\ \\ \relax
14:9 & \begin{tabularx}{0.7\textwidth}{X} 你為何像受驚嚇的人，像不能救人的勇士呢？耶和華啊，你在我們中間，我們是稱為你名下的人，求你不要離開我們。 \end{tabularx} \\ \\ \relax
14:10 & \begin{tabularx}{0.7\textwidth}{X} 耶和華論到這百姓如此說：「這百姓喜愛遊蕩，不約束自己的腳步，所以耶和華不悅納他們。現今他要記起他們的罪孽，懲罰他們的罪惡。」 \end{tabularx} \\ \\ \relax
14:11 & \begin{tabularx}{0.7\textwidth}{X} 耶和華又對我說：「不要為這百姓求福。 \end{tabularx} \\ \\ \relax
14:12 & \begin{tabularx}{0.7\textwidth}{X} 他們禁食的時候，我不聽他們的呼求；他們獻燔祭和素祭，我也不悅納。我卻要用刀劍、饑荒、瘟疫滅絕他們。」 \end{tabularx} \\ \\ \relax
14:13 & \begin{tabularx}{0.7\textwidth}{X} 我就說：「唉！主耶和華，看哪，那些先知常對他們說：『你們必不見刀劍，也不遭饑荒；耶和華要在這地方賞賜你們真正的平安。』」 \end{tabularx} \\ \\ \relax
14:14 & \begin{tabularx}{0.7\textwidth}{X} 耶和華對我說：「那些先知託我的名說假預言，我並未差遣他們，沒有吩咐他們，也沒有對他們說話；他們向你們預言的是虛假的異象、占卜、虛無，以及心中的詭詐。 \end{tabularx} \\ \\ \relax
14:15 & \begin{tabularx}{0.7\textwidth}{X} 所以耶和華如此說：『論到託我名說預言的那些先知，我並未差遣他們；他們說這地不會有刀劍、饑荒，其實那些先知自己必被刀劍、饑荒滅絕。 \end{tabularx} \\ \\ \relax
14:16 & \begin{tabularx}{0.7\textwidth}{X} 聽他們說預言的百姓必因饑荒、刀劍被扔在耶路撒冷的街道上，無人埋葬。他們連妻子帶兒女，都是如此。我必將他們的惡倒在他們身上。』」 \end{tabularx} \\ \\ \relax
14:17 & \begin{tabularx}{0.7\textwidth}{X} 你要向他們說這些話：願我眼淚汪汪，晝夜不息，因為少女---我百姓受了重大的打擊，傷口極其嚴重。 \end{tabularx} \\ \\ \relax
14:18 & \begin{tabularx}{0.7\textwidth}{X} 我若出到田間，看哪，有被刀殺的；我若進入城內，看哪，有因饑荒患病的；先知和祭司也在各地往來經商，不知如何是好。 \end{tabularx} \\ \\ \relax
14:19 & \begin{tabularx}{0.7\textwidth}{X} 你全然棄絕猶大嗎？你的心厭惡錫安嗎？你為何擊打我們，使我們無法得醫治呢？我們指望平安，卻得不著福氣；指望痊癒，看哪，受了驚惶。 \end{tabularx} \\ \\ \relax
14:20 & \begin{tabularx}{0.7\textwidth}{X} 耶和華啊，我們承認自己的罪惡和我們祖先的罪孽，因我們得罪了你。 \end{tabularx} \\ \\ \relax
14:21 & \begin{tabularx}{0.7\textwidth}{X} 求你為你名的緣故，不厭惡，不輕視你榮耀的寶座。求你記念，不要違背你與我們所立的約。 \end{tabularx} \\ \\ \relax
14:22 & \begin{tabularx}{0.7\textwidth}{X} 外邦虛無的神明中有能降雨的嗎？天能自降甘霖嗎？耶和華－我們的神啊，不是你嗎？我們要等候你，因為這一切都是你所造的。 \end{tabularx} \\ \\
[1ex]
\hline
\hline
\end{longtable}
講題:這麼近那麼遠 So Close Yet So Far

講員:林豪恩先生

經文:耶利米書14:4-22

場合:香港中文大學崇基學院禮拜堂主日崇拜

日期:2019年10月27日

一.引言

遠近如果是數學題,可量度,可計算,比較容易處理；然而,若遠近是關係問題,尤其是情侶之間的關係,就考起很多人了.近代最偉大的物理學家之一的霍金,能夠計算宇宙星體間的距離,但也表示掌握不到情感的距離.我們這些凡夫俗子更加容易在遠近之間迷失.有投入愛河的男孩子來訴苦說:「我在她身邊的時候,她抱怨我黐身；我去跟朋友玩的時候,她投訴我不理會她.我真是苦啊!誰能救我脫離這取死的遠近呢?」我聽著他們訴苦,心中暗道:「少年,你太年輕了.這個遠近的迷宮,只要你進入,終生也未必能夠拆解.不要以為公主和王子快快樂樂地住在一起,從此就解決了遠近的難題,若你不方便回家問你的爸爸媽媽,張學友這首歌都可以告知一二.

「這麼近那麼遠」　作詞:黃偉文    作曲:張學友

一天一天

日日夜夜 面對面

既相處 也同眠

一點一點

逐漸逐漸 便發現

縱相對 卻無言

靜靜默默

望著熟悉的背面

一彎身影

原來離我多麼的遠

像天涯 那一端

沒法行 前一吋

我 留著你在身邊

心 仍然很遠

我想伸手拉近點

竟觸不到那邊

就欠一點點

但這一點點

最近幾個月,也令不少人驚覺有些平時一起生活的家人,一起出遊的朋友,一起崇拜的弟兄姊妹,大家的距離原來是那麼遠的；同時又發現,有些平日不常交往的人,竟然是那麼近的.人際關係和社會事件,也令不少信徒與上帝的關係陷入掙扎之中.雖說人與上帝的距離只是一個禱告那麼近,但禱告了幾個月,也看不見上帝回應,難免質疑上帝是否遙遠得聽不到祈禱的聲音,還是自己不夠屬靈而上帝不聽禱告,難免令人懷疑自己與上帝關係到底是近還是遠.

今天崇拜,我們不是講愛情專題,也不是研討社會時事,而是透過經課的四段經文思考我們與上帝的遠近.

二.釋經及應用

1.耶利米書:遠近不取決於宗教行為,而取決於行為善惡

14:7耶和華啊,雖然我們的罪孽控告我們,求你為你名的緣故行動吧!我們本是多次背道,得罪了你.

14:8以色列所盼望,在患難時作他救主的啊,你在這地為何像寄居的,又如旅行的只住一夜呢?

14:9你為何像受驚嚇的人,像不能救人的勇士呢?耶和華啊,你在我們中間,我們是稱為你名下的人,求你不要離開我們.

耶利米時代的以色列投訴耶和華離開他們,不與他們同在.他們為甚麼認為耶和華離開他們呢?上文指出當時以色列陷入困境:

14:4因為無雨降在地上,土地就乾裂,農夫為此蒙羞抱頭.

14:5田野的母鹿因為無草也撇棄才生的小鹿.

14:6野驢站在光禿的高地喘氣,好像野狗；牠們的眼目因無草而失明.」

旱災不是一朝一夕的事情,一日復一日,仍然看不見雨水降下,他們開始質疑耶和華視而不見,或者見死不救:

14:8以色列所盼望,在患難時作他救主的啊,你在這地為何像寄居的,又如旅行的只住一夜呢?

14:9你為何像受驚嚇的人,像不能救人的勇士呢?

他們最後說:「求你不要離開我們」.以色列覺得與耶和華的距離很遠,他們的理解是耶和華離他們而去.在逆境長時間沒有改善的情況下,他們的感受及想法,此時此地的我們是可以理解的.上帝是否已經離我們而去?上帝是否離我們太遠,因此對我們的困境鞭長莫及呢?

耶和華如何回應以色列的抱怨和質疑呢?

14:10耶和華論到這百姓如此說:「這百姓喜愛遊蕩,不約束自己的腳步,所以耶和華不悅納他們.現今他要記起他們的罪孽,懲罰他們的罪惡.」

以色列以為耶和華離他們而去,耶和華指出其實是百姓離祂而去.當時的以色列應該會感到大惑不解,甚至覺得冤枉.為甚麼呢?因為下文提到,以色列是做足宗教行為的:

14:12他們禁食的時候,我不聽他們的呼求；他們獻燔祭和素祭,我也不悅納.

既然禁食,燔祭和素祭等都做足,以色列在何事上離耶和華而去呢?人傾向以宗教行為來量度與上帝的離近,例如返教會就等於虔誠,不過,先知一再強調:

「耶和華喜悅燔祭和平安祭,豈如喜悅人聽從他的話呢?」(撒上 15:22)

14:10節下半節指出當時以色列的真正問題:「現今他要記起他們的罪孽,懲罰他們的罪惡.」罪惡不單是個人層面,而是整個以色列的.然而,社會的罪惡未必容易被發現的,尤其是那些在遊戲中勝出的人,更不覺得有甚麼問題,有時,連先知也協助粉飾太平:

14:13「唉!主耶和華,看哪,那些先知常對他們說:『你們必不見刀劍,也不遭饑荒；耶和華要在這地方賞賜你們真正的平安.』」

物必先腐後而蟲生,耶利米指出當時以色列的困境只是內裡腐敗的外顯而已.

14:19你全然棄絕猶大嗎?你的心厭惡錫安嗎?你為何擊打我們,使我們無法得醫治呢?我們指望平安,卻得不著福氣；指望痊癒,看哪,受了驚惶.

14:20耶和華啊,我們承認自己的罪惡和我們祖先的罪孽,因我們得罪了你.

14:21求你為你名的緣故,不厭惡,不輕視你榮耀的寶座.求你記念,不要違背你與我們所立的約.

14:22外邦虛無的神明中有能降雨的嗎?天能自降甘霖嗎?耶和華－我們的上帝啊,不是你嗎?我們要等候你,因為這一切都是你所造的.

14:19「我們指望平安,卻得不著福氣；指望痊癒,看哪,受了驚惶.」這句說話在此時此地的人讀來,或有同感.面對那樣的處境,耶利米的兩項呼籲值得我們參考:

一,14:20「耶和華啊,我們承認自己的罪惡和我們祖先的罪孽,因我們得罪了你.」這點大多數人都明白也同意,如果要拉近距離,修補關係,「承認自己的罪惡和我們祖先的罪孽」可能是必須的第一步.

二,14:22「外邦虛無的神明中有能降雨的嗎?天能自降甘霖嗎?耶和華－我們的上帝啊,不是你嗎?我們要等候你,因為這一切都是你所造的.」這是要求我們確認所信的是誰,我們的主是誰.從耶利米書的記載中,可以看到那個時代的以色列經常掙扎投靠北面的強國亞述和巴比倫,還是南面的強國埃及.其實投靠北方強國或是南方強國意義是一樣的,就是相信強權和武力.耶和華希望他們投靠祂,這才是平安的基礎.投靠耶和華是甚麼意思呢?是超越眼前的強權和武力,相信人是上主創造的,有價值及尊嚴,應以平等和友愛相處；相信上帝討厭罪惡和喜愛公義,並以此建立人際社會.

耶利米書這個片段指出,與上帝的離或近,不是由宗教行為決定,而是與行為的善與惡更有關係.

2.詩篇及福音書:遠近不取決於地理狀態,而取決於心靈狀態

詩篇

84:1萬軍之耶和華啊,你的居所何等可愛!

84:2我羨慕渴想耶和華的院宇,我的內心,我的肉體向永生上帝歡呼.

84:3萬軍之耶和華－我的王,我的上帝啊,在你祭壇那裡,麻雀為自己找到了家,燕子為自己找著菢雛之窩.

84:4如此住在你殿中的有福了!他們不斷地讚美你.

84:5靠你有力量,心中嚮往錫安大道的,這人有福了!

84:6他們經過「流淚谷」,叫這谷變為泉源之地；且有秋雨之福蓋滿了全谷.

84:7他們行走,力上加力,各人到錫安朝見上帝.

84:8萬軍之耶和華上帝啊,求你聽我的禱告!雅各的上帝啊,求你側耳而聽!

84:9上帝啊,我們的盾牌,求你觀看,求你垂顧你受膏者的面!

84:10在你的院宇一日,勝似千日；寧可在我上帝的殿中看門,不願住在惡人的帳棚裡.

84:11因為耶和華上帝是太陽,是盾牌,耶和華要賜下恩惠和榮耀.他未嘗留下福氣不給那些行動正直的人.

84:12萬軍之耶和華啊,倚靠你的人有福了!

這是一首朝聖詩,朝聖者認為,錫安是耶和華的山,所以「各人到錫安朝見上帝」；建在錫安山上的聖殿是耶和華的居所,難怪「在你的院宇一日,勝似千日」.人們都相信地理上的遠近與關係的遠近有關,所以長途跋涉往聖地朝聖,為的是希望更接近神明.

現實往往是多樣性的,在聖殿可能更接近上帝,也可能更遠離上帝.讓我們看看經課的福音書記載那個在聖殿的片段:

路加福音

18:9耶穌向那些自以為義而藐視別人的人講了這比喻:

18:10「有兩個人上聖殿去禱告,一個是法利賽人,一個是稅吏.

18:11法利賽人獨自站著,自言自語地禱告說:『上帝啊,我感謝你,我不像別人勒索,不義,姦淫,也不像這個稅吏.

18:12我每週禁食兩次,凡我所得的都獻上十分之一.』

18:13那稅吏遠遠地站著,連舉目望天也不敢,只捶著胸,說:『上帝啊,開恩可憐我這個罪人!』

處身神聖之地,不等於心靈就會聖化；站在聖者面前,不等於就會謙卑下來.在聖殿,有些人也難以脫去「藐視別人」的高傲；在聖者面前,有些人也掩不住「自以為義」的飄飄然.說到底,聖地只是一處地理位置,聖殿也只是一棟建築物,與上帝遠或近,未必取決於地理狀態,卻在於心靈狀態.

有些狀態是很吊詭的,自以為近得可以拍肩膀,原來伸手摸著的是空氣；自以為遠得遙不可及的,那人卻在燈火闌珊處.與上帝的或遠或近,如何判斷呢?是遠是近,誰來判斷呢?

18:14我告訴你們,這人回家去比那人倒算為義了.因為凡自高的,必降為卑；自甘卑微的,必升為高.」

這人是誰呢?就是那位「遠遠地站著,連舉目望天也不敢」又被旁邊那位自義者「抽水」的稅史,他沒有成績表交給上帝,只能夠「捶著胸,說:『上帝啊,開恩可憐我這個罪人!』」

或者,我們也不需要太過計較自己與上帝的遠或近.當我們堆砌證據,展示自己與上帝接近的時候,可能就是遠離的時候.主日崇拜程序的第一部分「預備禮儀」,會否暗示我們「遠遠地站著,連舉目望天也不敢」呢?第二部分「省察心靈」,是否引領我們「捶著胸,說:『上帝啊,開恩可憐我這個罪人!』」呢?我們有甚麼成就值得向上帝展示嗎?我們有甚麼偉業向上帝邀功嗎?沒有,我們只是等候上帝說:「必要赦免你們的罪,洗淨你們一切的不義」,等候帷幕打開,等候聖靈光照.

3.提摩太後書:遠近不取決於人的認同,而取決於跑完當跑的路

4:6至於我,我已經被澆獻,離世的時候到了.

4:7那美好的仗我已經打過了,當跑的路我已經跑盡了,該信的道我已經守住了.

4:8從此以後,有公義的冠冕為我存留,就是按著公義審判的主到了那日要賜給我的；不但賜給我,也賜給凡愛慕他顯現的人.

4:9你要趕緊到我這裡來.

4:10因為底馬貪愛現今的世界,已經離棄我,往帖撒羅尼迦去了；革勒士往加拉太去；提多往撻馬太去；

4:11只有路加在我這裡.你來的時候把馬可帶來,因為他在服事上於我有益.

4:12我已經打發推基古往以弗所去.

4:13我在特羅亞留給加布的那件外衣,你來的時候要帶來,那些書也帶來,特別是那幾卷羊皮的書.

4:14銅匠亞歷山大多方害我；主必照他所行的報應他.

4:15你也要防備他,因為他極力抗拒我們的話.

4:16我初次上訴時,沒有人前來幫助,竟都離棄了我,但願這罪不歸在他們身上.

4:17惟有主站在我身邊,加給我力量,使我能把福音完整地傳開,讓所有的外邦人都聽見；我也從獅子口裡被救出來.

4:18主必救我脫離一切的兇惡,也必救我進他的天國.願榮耀歸給他,直到永永遠遠.阿們!

馬丁路德金在1968年2月6日一次演講中說:「每當有事情發生的時候,怯懦者會問:這樣安全嗎?投機者會問:這樣權宜嗎?虛榮者會問:這會受歡迎嗎?但是良心會問,這是對的嗎?總有一天,一個人必須採取一種既不安全,也不權宜,也不受歡迎的立場,但他必須採取這種立場,因為良心告訴他這是正確的.」(“Cowardice asks the question, ‘Is it safe?’ Expediency asks the question, ‘Is it politic?’ Vanity asks the question, ‘Is it popular?’ But, conscience asks the question, ‘Is it right?’ And there comes a time when one must take a position that is neither safe, nor politic, nor popular, but one must take it because one’s conscience tells one that it is right”)

使徒保羅在他的人生路上,遇上「4:10底馬貪愛現今的世界,已經離棄我」,「4:14銅匠亞歷山大多方害我」,「4:16我初次上訴時,沒有人前來幫助,竟都離棄了我」,或許,他們都覺得保羅的行動冒險,太不明智,也具爭議.然而,有人離去不等於上帝也離去,保羅深信「4:17惟有主站在我身邊,加給我力量,使我能把福音完整地傳開」.保羅不需要問:「我與上帝是近還是遠?」正如互相信任的夫妻不用問:「你是否還愛我?」如果能夠肯定自己所打的是「那美好的仗」,所跑的是「當跑的路」,所守的是「該信的道」,那麼,我們就可以有保羅那分安心,而不需要花心思去捕捉那種飄忽不定的與上帝遠或近的感覺.



\newpage

\section{今天救恩到了這家}
\label{sec:476}
\textbf{戴浩輝牧師 - 今天救恩到了這家}
\newline
\newline
\hyperref[sec:475]{< < < PREV SERMON < < <}
~
\hyperref[sec:toc]{[返目錄]}
~
\hyperref[sec:477]{> > > NEXT SERMON > > >}


路加福音 19:1-10
\newline
\begin{longtable}{cl}
\hline
\hline
章節 & 經文 (和合本修訂版)\\
\hline
19:1 & \begin{tabularx}{0.7\textwidth}{X} 耶穌進了耶利哥，要從那裡經過。 \end{tabularx} \\ \\ \relax
19:2 & \begin{tabularx}{0.7\textwidth}{X} 有一個人名叫撒該，作稅吏長，是個財主。 \end{tabularx} \\ \\ \relax
19:3 & \begin{tabularx}{0.7\textwidth}{X} 他要看看耶穌是怎樣的人，只因人多，他的身材又矮，所以看不見。 \end{tabularx} \\ \\ \relax
19:4 & \begin{tabularx}{0.7\textwidth}{X} 於是他跑到前頭，爬上桑樹，要看耶穌，因為耶穌要從那裡經過。 \end{tabularx} \\ \\ \relax
19:5 & \begin{tabularx}{0.7\textwidth}{X} 耶穌到了那裡，抬頭一看，對他說：「撒該，快下來！今天我必須住在你家裡。」 \end{tabularx} \\ \\ \relax
19:6 & \begin{tabularx}{0.7\textwidth}{X} 他就急忙下來，歡歡喜喜地接待耶穌。 \end{tabularx} \\ \\ \relax
19:7 & \begin{tabularx}{0.7\textwidth}{X} 眾人看見，都私下議論說：「他竟然到罪人家裡去住宿。」 \end{tabularx} \\ \\ \relax
19:8 & \begin{tabularx}{0.7\textwidth}{X} 撒該站著對主說：「主啊，我把所有的一半給窮人；我若勒索了誰，就還他四倍。」 \end{tabularx} \\ \\ \relax
19:9 & \begin{tabularx}{0.7\textwidth}{X} 耶穌對他說：「今天救恩到了這家，因為他也是亞伯拉罕的子孫。 \end{tabularx} \\ \\ \relax
19:10 & \begin{tabularx}{0.7\textwidth}{X} 人子來是要尋找和拯救失喪的人。」 \end{tabularx} \\ \\
[1ex]
\hline
\hline
\end{longtable}
講題:今天救恩到了這家 Today salvation has come to this house

講員:戴浩輝牧師

經文:路加福音19:1-10

場合:香港中文大學崇基學院禮拜堂主日崇拜

日期:2019年11月3日

引言

福音是藉甚麼可以到自己的家裡呢?「福音到我家」這個題目是從路十九9耶穌對撒該所說的「今天救恩到了這家」而來的,因此,我也很想藉撒該這個很可愛的事蹟與大家分享這個題目.

路十九1-10這段經文的確是十分吸引的一段經文,我還記得約在二十多年前,我帶領兒童詩歌班唱過名為「撒該」的一首歌是這樣的:「從前有一個人名叫撒該,他是一個矮小的人,有一天他定意爬上桑樹要看耶穌究竟是誰；耶穌從桑樹下經過,抬頭看見他在那裡,『今天我必住在你家,今天我必住在你家.』」耶穌到撒該的家那就是我們題目所說福音到他的家,因為聖經所載的就是耶穌基督的福音,耶穌基督在那裡,那裡就有祂所帶來的福音.然而,耶穌如何到撒該的家?換言之,福音如何到達撒該的家?撒該作了些甚麼使耶穌的福音臨到他的家?這個問題的答案可謂十分關鍵,因為這問題的答案可以給我們一個示範作用,讓我們同樣的去做,使耶穌的福音臨到我們的家.

一 渴慕見主

在這個事蹟中,福音臨到的第一個條件是撒該渴慕見耶穌.聖經沒有記述他為何「要看看耶穌是怎樣的人」,從聖經的記述中我們只可以想像,他是稅吏長,他是個財主,但他並不開心滿足,為何?因為稅吏在巴勒斯坦猶太人的地方並不是為本地政府工作,而是為敵人即羅馬政府工作,稅吏被視為賣國賊,而撒該是稅吏長,即賣國頭子；可想而知,他可能在眾稅吏中得到人家的奉承,但也知道這些稅吏各懷鬼胎,沒有真心.在自己的社群當中是一個如何不受歡迎的人,我們是這樣的人嗎?在今天可能有些公僕都面臨著撒該所面對的問題.聖經描述他是個矮小的人,當群眾都圍著耶穌的時候,如果他在群眾當中受尊重的話,他根本無需爬上桑樹上看耶穌.但他想見主面心切,他不惜在其下屬面前放下尊嚴,爬上樹去,為要見主一面.他這種渴慕見主的心,便成為福音到他家的第一個條件.各位弟兄姊妹,福音如何能到我的家?豈不是我們先有一個要認識主,追求主的心嗎?渴慕尋求主是撒該的心,故事的結尾卻提醒我們,耶穌來是要尋找我們.這裡第3節中的「看看」(尋求要看)ἐζήτει (ezetei)和第10節的「尋找」ζητῆσαι (zetesai)原文是同一個字,作者在這裡要提醒我們,我們渴慕主的時候,主耶穌一定使我們的渴慕得到滿足.

渴慕見主實在的意思就是天天的與主有交往,我們每天的讀經和禱告就是表明我們渴慕見主.當今天我們面對香港這個局面,我們更多要尋求見主,使我們得著福音,得著安慰.

二 接待救主

撒該渴慕見主的心得到回應,主耶穌叫喚他下來,並要到他家裡住宿；撒該的反應是怎樣的呢?第6節記載:「他就急忙下來,歡歡喜喜地接待耶穌.」福音如何到我們的家呢?這個故事提醒我們,正如撒該一樣,福音是藉著我們接待主而臨到我們的.想到這裡我認為我們每一位都希望我們全家歸主,但很多時我們又成為我們全家歸主的絆腳石,因為我們的言行很多時都成為家人對信仰的誤會和拒絕.我們都是蒙恩的罪人,在一天的生活中,我們的自我很可能使我們的基督徒身份受家人的攻擊,我們都不希望這樣,但單靠我們自己,我們只有壞見證.如果耶穌不到我們家裡來,我們家裡沒有福音.究竟耶穌怎樣來?我們如何接待救主到我們家?這個接待耶穌到我們家的第一步應該就是我們常禱告邀請主聖靈在我們家裡工作,充滿我們,使我們行事為人學像耶穌,以至人看見我們,便像看到耶穌一樣.

第二,是我們邀請基督的身體到我們家裡,那就是基督徒肢體到我們家中,讓我們家人都認識我們基督徒朋友,讓他們把耶穌的愛活現在我們家中.當然,這不是要放棄我們在家中的好榜樣,而是在第三者的不受家人偏見所影響下,讓基督徒的愛臨在我們的家中.讓我們的弟兄姊妹也成為我們家人的弟兄姊妹,這讓基督的福音在我們家中便有接觸點,就有可以往下扎根,往上結果的可能了.撒該是歡歡喜喜地迎接耶穌,讓我們也是歡歡喜喜地歡迎我們主內的弟兄姊妹到我們家裡,這樣也不知不覺的把福音接到我們自己的家中.

三 決心悔改

救恩到我家的第三個條件是決心悔改,這個故事有一個特別的地方,群眾兩次的把耶穌與撒該分開,一次是肉身上的,一次卻是心靈上的.所謂肉身是指群眾好像沒有給予撒該這個矮小的人機會見耶穌,他們築成人牆,把撒該的視線阻擋了.第二次是在第7節,他們私下議論,群眾用言論來阻止撒該親近耶穌.這兩個事件的主要原因可謂出於撒該平日如何對待群眾.一個財主如果是正當得財,又過著行律法的生活,一定受人家的歡迎的,撒該是個財主,但他不受歡迎.他是稅吏長,是戕害同胞的人,他又怎能得著同胞的接納?然而,這些阻隔卻可以被打破,那就是撒該的真心悔改.撒該要把一半所有的分給窮人,凡被他訛詐過的,他都還他四倍.所謂悔改,就是從從前的路回轉過來,因著與上帝在耶穌裡和好的緣故,也和上帝的造化一同復和.按撒該而言,他真正的悔改了,他首先把阻擋他與耶穌的障礙清除,那就是那些群眾；藉著他表達愛他們的心:分一半給窮人,訛詐了誰便還他四倍.假若他家財有十元,先給了五元分給了窮人,如果他真的訛詐了某人,那怕只是一元,他就要付出四元；他的財產只剩下一元.如果訛詐的是兩元,他還要欠債.撒該決心悔改給我們甚麼提醒?我認為最主要是好像撒該一樣,先清除阻隔我們與耶穌關係的問題；而這些問題很多時都不是他人的問題,而是自我的問題.家人不肯信主很多時都不是家人的問題,而是我們使他們不肯相信我們的上帝是那位又真又活的神；因為在我們的生命裡,並沒有表現這位又真又活的神.基督徒與其他人的分別在於他們能反省悔改,這是與我們罪性生命的機制相反的,我們罪性的機制都認為自己常常是對的,沒有錯誤的；但這是不真實的.當人說我們錯時,我們自我防衛的機制開動,我們拒絕與那些說我們錯的人來往,正所謂「忠言逆耳」,我們甚至會攻擊那些勸導我們的人.在今天,香港的問題豈不是我所眼見的就是真理嗎?於是我們都只能認為對方是錯的.於是,我們最後再也沒有人勸導我們,而我們便一錯再錯了.各位弟兄姊妹,福音到我家的第三個條件就是我們天天反省悔改,勝過昨日的我,而不是鬥爭勝過我們的家人或朋友.撒該得著耶穌,便如保羅一樣,將萬事看作糞土；我們得著耶穌,是否也將我們的所謂尊嚴,看法,不肯饒恕的心放下,以耶穌為至寶呢 (腓三8-9)?

或許,當我們思考到此時,我們可能仍看不到福音在哪兒?福音就是我們渴慕見主,接待救主和決心悔改嗎?這不是福音,都是我的行為!不錯,我們迎接救恩到我們的家是帶有行動的,而這救恩─耶穌基督的福音就是祂對撒該的宣告:「人子來是要尋找和拯救失喪的人.」

結語

今天撒該的故事提醒我們,救恩如何能在我們的家裡?我們的家人如何能得救?豈不是我們都有一個渴慕見主的心?這個心的實底就是我們天天藉靈修親近主.第二,福音到我們的家是藉接待主到我們家,實際的方法就是讓基督徒肢體參與我們家人得救恩的計劃,他們到我們家裡成為我們家人的朋友,他們的愛心激勵我們家人接受福音.最後,我們藉耶穌基督愛我們這些罪人的力量,天天反省悔改,特別使我們和家人和好,使和好的福音真正臨到我們家中.主耶穌今天向你我宣告:「人子來是要尋找和拯救失喪的人.」阿們!



\newpage

\section{那沒有看見就信的有福了}
\label{sec:477}
\textbf{劉國偉先生 - 那沒有看見就信的有福了}
\newline
\newline
\hyperref[sec:476]{< < < PREV SERMON < < <}
~
\hyperref[sec:toc]{[返目錄]}
~
\hyperref[sec:478]{> > > NEXT SERMON > > >}


約翰福音 20:19-29
\newline
\begin{longtable}{cl}
\hline
\hline
章節 & 經文 (和合本修訂版)\\
\hline
20:19 & \begin{tabularx}{0.7\textwidth}{X} 那日（就是七日的第一日）晚上，門徒因怕猶太人，所在的地方門都關了。耶穌來，站在當中，對他們說：「願你們平安！」 \end{tabularx} \\ \\ \relax
20:20 & \begin{tabularx}{0.7\textwidth}{X} 說了這話，他把手和肋旁給他們看。門徒一看見主就喜樂了。 \end{tabularx} \\ \\ \relax
20:21 & \begin{tabularx}{0.7\textwidth}{X} 於是耶穌又對他們說：「願你們平安！父怎樣差遣了我，我也照樣差遣你們。」 \end{tabularx} \\ \\ \relax
20:22 & \begin{tabularx}{0.7\textwidth}{X} 說了這話，他向他們吹一口氣，說：「領受聖靈吧！ \end{tabularx} \\ \\ \relax
20:23 & \begin{tabularx}{0.7\textwidth}{X} 你們赦免誰的罪，誰的罪就得赦免；你們不赦免誰的罪，誰的罪就不得赦免。」 \end{tabularx} \\ \\ \relax
20:24 & \begin{tabularx}{0.7\textwidth}{X} 那十二使徒中，有個叫低土馬的多馬，耶穌來的時候，他沒有和他們在一起。 \end{tabularx} \\ \\ \relax
20:25 & \begin{tabularx}{0.7\textwidth}{X} 其他的門徒就對他說：「我們已經看見主了。」多馬卻對他們說：「除非我看見他手上的釘痕，用我的指頭探入那釘痕，用我的手探入他的肋旁，我絕不信。」 \end{tabularx} \\ \\ \relax
20:26 & \begin{tabularx}{0.7\textwidth}{X} 過了八日，門徒又在屋裡，多馬也和他們在一起。門都關了，耶穌來，站在當中，說：「願你們平安！」 \end{tabularx} \\ \\ \relax
20:27 & \begin{tabularx}{0.7\textwidth}{X} 然後他對多馬說：「把你的指頭伸到這裡來，看看我的手；把你的手伸過來，探入我的肋旁。不要疑惑，總要信！」 \end{tabularx} \\ \\ \relax
20:28 & \begin{tabularx}{0.7\textwidth}{X} 多馬回答，對他說：「我的主！我的神！」 \end{tabularx} \\ \\ \relax
20:29 & \begin{tabularx}{0.7\textwidth}{X} 耶穌對他說：「你因為看見了我才信嗎？那沒有看見卻信的有福了。」 \end{tabularx} \\ \\
[1ex]
\hline
\hline
\end{longtable}
講題:那沒有看見就信的有福了 Blessed are those who believe without seeing

經文:約翰福音20章19-29節

講員:劉國偉先生

場合:香港中文大學崇基學院禮拜堂主日崇拜

日期:2019年11月10日

引言

今年8月20日晚,我與太太一起,參加了由福戲網絡主辦,在循道衛理聯合教會香港堂舉行的一個講座:「歷史的潘霍華,現代的潘霍華」(Dietrich Bonhoeffer),當晚的講員潘錦華導演分享「由《馬丁路德》到《潘霍華的情書》:表演藝術與我何干」,另一位講員是我們禮拜堂的榮譽助理校牧鄧瑞強博士,瑞強兄分享的題目是「思想與行動之間:潘霍華呈現的教會與信徒的社會責任」,當晚座無虛席,連地庫的副堂也坐滿了,多少反映出在香港處於水深火熱的情景下,弟兄姊妹想透過潘霍華的生命故事和他的神學思想,一種由歷史走到當下的熱切的期盼.潘導演和瑞強兄的分享都很精彩,會後,勾起了我很多思緒.

我在1979年初次聽聞潘霍華,緣於在中文大學圖書館閱讀一篇有關潘霍華神學思想的研究論文,然後我開始閱讀潘霍華的著作,首先是1974年道聲出版社出版的「追隨基督」中譯本,當時已經出版了5年.當中有蘇恩佩姊妹撰寫的「代序」.恩佩姐在文章的第一句,引用潘霍華的名句「當基督呼召一個人,祂是召他來為祂死.」雖然當時已經不是第一次閱讀這兩句句子,卻再一次震撼著我的心房,我不期然地想,那是一種怎樣的信心?潘霍華的好友,英國哲齊斯特主教(Bishop of Chichester)貝爾(Bishop Bell)如此演繹潘霍華這名句:「當基督呼召一個人,祂是召他來為祂死:死的形式有多種,不過作基督門徒的真諦,可說完全包括在這句話裡面了.」

然後今年9月21日,我與家人一起觀賞了在文化中心上演的話劇「潘霍華的情書」(Bonhoeffer’s Love Letters),一個熟悉的故事,對於我而言,潘霍華並非一個陌生人,當晚觀賞這齣話劇,有很多場景和段落,都深受感動,也重新刺激我去問自己:潘霍華所見證所活出的是怎樣的信心?

於是,我決定再次翻開「獄中書簡」(Letters and Papers from Prison),赫然看見當年的一些閱讀筆記和在書本劃下的記號和痕跡.我現在讀出我曾經標記的一個小段落:「一個人必須完全過著今世的生活,才能學習到信心.這就是我所說的現世之意義,即負起生命的一切責任與困難,成功與失敗,一切經驗與無可奈何之事.就在這樣的生命中,我們才把自己無條件地放置在上帝手裡,參與祂在世上的苦難,與基督在客西馬尼園一同警醒.這就是信仰……這樣,當我們生活於今世而參與上帝的苦難,成功怎能叫我們傲慢,失敗又怎能使我們走入迷途呢?」

信心:在恐懼中倚靠主賜的平安

今天講道的題目是「那沒有看見就信的有福了」,是從死裡復活的主,向門徒顯現時所說的一句話,承接我剛才所問的,我們所見證所活出的是怎樣的信心?

使徒約翰記載了三次主耶穌向門徒顯現.約翰所說的門徒,明顯是指向耶穌所特別揀選的十一個門徒,因為猶大自殺身亡了,所以當時就少了一個.

第一次顯現記載在約翰福音20章19至25節；第二次顯現記載在同一章聖經的26至29節；第三次記載在第21章1至23節.我們今天主要集中思想,約翰筆下所描述主耶穌向門徒第一及第二次顯現的兩個小段落.

約翰寫道:「那日(就是七日的第一日)晚上,門徒所在的地方,因怕猶太人,門都關了.耶穌來,站在當中,對他們說:「願你們平安!」說了這話,就把手和肋旁指給他們看.門徒看見主,就喜樂了.(約21:19-20)

約翰所形容的情景,是星期日的晚上,門徒仍然害怕猶太人,於是把門關了,原文「關」這個動詞,其實是「鎖上」的意思.很明顯的,抹大拉的馬利亞向門徒見證耶穌基督已經復活,彼得和約翰親眼看見空墳墓的環境證據,再加上前往以馬忤斯路上的兩個門徒,見證曾經與復活的主相遇,在這非常特別的一天,再三聽到耶穌已經復活的消息,卻仍然未能消除他們心中的恐懼.弟兄姊妹,黑夜是可怕的,往往埋藏了很多黑暗的人和事,甚至危機,我們也經歷了許多漫長的黑夜,有時候真的想把自己關鎖在一個安全的房子裡,以圖與外間的黑夜隔開.

門徒是有福的,在他們經過從星期四深夜到星期五凌晨,耶穌被捉拿開始,到星期日從早到晚十多個小時所聽見既是奇妙,又一時間難以明白和接受的見證,當他們仍然活在恐懼之中,而耶穌就在這情景底下,向他們顯然.

Shalom「願你們平安」在此時此刻,不再是一句日常慣用的問候語,聽在這一群感到羞愧,不解,恐懼和疑惑當中的門徒的耳中,這一句問候如今象徵神與人和好,將門徒從對死亡和黑暗的恐懼的光景,帶進基督復活大能的一句滿有恩典的說話.

約翰告訴我們,主耶穌說了「願你們平安」這句話之後,就把他的傷痕指給門徒看.此時此刻,基督的傷痕,或者說是曾經死過的印記,成為恩典的記號,最重要的,是門徒看見復活的主,實實在在的體驗了復活的大能,就喜樂了.

第一句「願你們平安」是耶穌基督邀請門徒進入祂的復活的生命,第二句「願你們平安」,就帶著使命和差遣.

在星期四的晚餐聚會,主耶穌跟門徒分離前的禱告,說:「你怎樣差我到世上,我也照樣差他們到世上.」(約17:18)如今復活的主,在完成了救贖大功之後,再一次提出父差遣祂,祂要照樣差遣門徒.我就想起了潘霍華的名句:「當基督呼召一個人,祂是召他來為祂死.」接受耶穌基督的差遣,就包含了在生命旅途中,要為信仰的緣故,冒著大大小小的風險,或者直接點說,基督要求我們背起十字走天路,本身就是一個冒險的旅程 A risk taking journey.

在這個冒險的旅程中,從神而來的平安,就成為非常重要的保障.這份平安,不在乎識見和閱歷,也不在於意志和訓練,乃是聖靈的工作.所以,主耶穌說了平安與差遣的話,就向門徒吹一口氣,說:「你們受聖靈.」

21至23節的經文如此說:「耶穌又對他們說:「願你們平安!父怎樣差遣了我,我也照樣差遣你們.」說了這話,就向他們吹一口氣,說:「你們受聖靈!你們赦免誰的罪,誰的罪就赦免了；你們留下誰的罪,誰的罪就留下了.」 (約20:21-23)

正如馬太記載了主耶穌對教會的教導:「我實在告訴你們,凡你們在地上所捆綁的,在天上也要捆綁；凡你們在地上所釋放的,在天上也要釋放.」(太18:18)勉勵門徒和教會,我們從主所接受的差遣,是生死攸關的任務,我們要傳揚的是叫人出死入生的基督救恩,凡接待主的,就是信主名的人,主就賜他們權柄,作神的兒女.(約1:12)

那個星期日的夜晚,是黑暗與光明交錯的時空,也是恐懼和喜樂相纏的晚上,我想像不到主耶穌顯現之後又離去,到底會是一種怎樣的感覺!不過我確信,因為主耶穌的顯現,還有聖靈的臨在,門徒的生命就起了徹底的質變,他們從新體悟信心的道路,不在於無需冒險,或者沒有危險,卻原來是當我們進入基督的復活大能,和靠著聖靈的能力,我們所得著的從主而來的平安,就幫助我們勇敢地接受基督的差遣.

我總不信

使徒約翰告訴我們,當夜,多馬不在那房子,也就是說,他仍然未曾親眼見過復活的基督.

「那十二個門徒中,有稱為低土馬的多馬；耶穌來的時候,他沒有和他們同在.那些門徒就對他說:「我們已經看見主了.」多馬卻說:「我非看見他手上的釘痕,用指頭探入那釘痕,又用手探入他的肋旁,我總不信.」(約20:24-25)

多馬鏗鏘有力的一句話「我總不信」,表達了他仍然活在一個講求合理推斷,要求穩穩陣陣的思維方式之中.作為福音書的讀者,我們千萬不要以為,多馬是唯一的一個表示不信,其實在馬可福音也有類似的描述:「在七日的第一日清早,耶穌復活了,就先向抹大拉的馬利亞顯現；耶穌從他身上曾趕出七個鬼.她去告訴那向來跟隨耶穌的人；那時他們正哀慟哭泣.他們聽見耶穌活了,被馬利亞看見,卻是不信.這事以後,門徒中間有兩個人,往鄉下去；走路的時候,耶穌變了形象向他們顯現,他們就去告訴其餘的門徒；其餘的門徒,也是不信.」(可16:9-13)

弟兄姊妹,信心旅程並非一蹴即就,生命的轉化和成長需要時間.不要說多馬沒有經歷耶穌基督第一次向門徒顯現的每一個細節,因而暫時失去一個信心跳躍騰飛的機會,就是其他的門徒在基督顯現過了八日之後,他們仍然謹慎地在房子裡把門繼續關鎖,這實在與五旬節聖靈降臨之後,門徒的集體信仰面貌有很大的落差.

經文說:「過了八日,門徒又在屋裡,多馬也和他們同在,門都關了.耶穌來,站在當中說:「願你們平安!」就對多馬說:「伸過你的指頭來,摸我的手；伸出你的手來,探入我的肋旁.」多馬說:「我的主!我的上帝!」 (約20:26-28)

耶穌當然理解,門徒仍然要活在一個關鎖了門的房子的生命狀態,耶穌這個復活的身體穿過那關鎖了的大門,再次以「願你們平安」接納門徒當下的生命境況,顯現在門徒當中,充滿了對門徒的體恤和憐愛.

走信心的道路,本來就應該全然信靠,事事講求證據,兆頭,印證,本身就是一種「我信,但我信不足」的屬靈狀況.親愛的弟兄姊妹,這是否我們人生旅程的真實寫照?特別當我們身處漫長的黑夜,包括一些臨到自己的意想不到的傷害,無可奈何的際遇,還有不合理,無恩情的人和事,在等候之間遲遲沒有答案,未見出路,更不知何時見黎明,這些時候,信心就成為非常真實的考驗.

主耶穌並非單單體恤和憐愛多馬,祂對多馬說:「你因看見了我才信；那沒有看見就信的有福了.」 我相信主耶穌所說的這句話,不單只是向多馬說,也向眾多門徒說,並且今天仍然在我們當中如此說.是的,「不要疑惑,總要信!」是信德的基礎,疑惑本來就抵銷信心,不過疑惑之後,能夠靠著信心走下去,我們又會為人生旅程跨進了信心的一大步,主耶穌總會在那個漫長的黑夜,還有關鎖了門的房子裡,向我們說「願你們平安」,給我們信心和勇氣,走進祂復活的大能裡,也同樣賜我們平安,差遣我們進入這個世界,憑信心走上冒險的旅程.

憑著信

潘霍華,1906年2月4日生於德國,是一位將生命獻在殉道者祭壇上,自始至終忠心追隨基督的人,他確實被公認為一位出色的神學家,不過我更喜歡形容潘霍華是一個忠心追隨基督的門徒.潘霍華被納粹秘密警察捉拿,人生的最後幾個星期被關押在一個大監牢裡,他沒有看見希特拉和納粹德國的結局,在1945年4月9日曙光初露的時刻,被納粹黨送上了刑場,數天之後,盟軍宣告勝利.

潘霍華是憑著信心,在世走過39年之後,離開這個世界的,當他從囚倉被押送往刑場之際,他與囚友說:「這就是終局,但在我卻是生命的開始.」.

結束今天講道的時候,我選讀「獄中書簡」記載潘霍華在1945年元旦所寫的禱文:

有了善的力量同在與導引我,安慰與鼓勵,遠超過一切的恐懼.

舊歲還在慢慢地苦楚著我的心靈,沉痛冗長的日子仍然在拖延著.父啊,求你容許你所磨練的靈魂,得到你所允許的救治.

倘若我們得喝乾這悲慘的杯,直到飲盡痛苦的屑滓,在你的命令下,我們決不躊躇,願以感激的心接受,你恩手之所施賜.

讓燭光發射輝煌的祝賀,這豈不就是黑暗中從你來的亮光,要把我們帶往那渴望中的相聚.

現在夜更深沉,求容許我們聽見你兒女的聲音,從圍繞著我們的黑暗世界,發出那普世讚頌的歌聲!

一切善的力量來幫助我們,我們將勇敢地面對將來,不論晨昏,上帝必恩待我們,每逢新歲開始,我們這信念更為堅定.

弟兄姊妹,「那沒有看見就信的有福了」,因為我們活著是憑著信心,不是憑著眼見,We live by faith, not by sight.  願主加添我們信心,願榮耀歸給三一真神!



\newpage

\section{從相分之痛到相聚之盼}
\label{sec:478}
\textbf{林豪恩先生 - 從相分之痛到相聚之盼}
\newline
\newline
\hyperref[sec:477]{< < < PREV SERMON < < <}
~
\hyperref[sec:toc]{[返目錄]}
~
\hyperref[sec:479]{> > > NEXT SERMON > > >}


耶利米書 23:3-6
\newline
\begin{longtable}{cl}
\hline
\hline
章節 & 經文 (和合本修訂版)\\
\hline
23:3 & \begin{tabularx}{0.7\textwidth}{X} 我要從我趕他們到的各國召集我羊群中剩餘的，領他們歸回本處；他們必生養眾多。 \end{tabularx} \\ \\ \relax
23:4 & \begin{tabularx}{0.7\textwidth}{X} 我必設立牧人照管他們，牧養他們。他們不再懼怕，不再驚惶，沒有一個失喪的。這是耶和華說的。 \end{tabularx} \\ \\ \relax
23:5 & \begin{tabularx}{0.7\textwidth}{X} 「看哪，日子將到，我要為大衛興起公義的苗裔； 他必掌王權，行事有智慧，在地上施行公平和公義。這是耶和華說的。 \end{tabularx} \\ \\ \relax
23:6 & \begin{tabularx}{0.7\textwidth}{X} 在他的日子，猶大必得救，以色列也安然居住。他的名必稱為『耶和華－我們的義』。 \end{tabularx} \\ \\
[1ex]
\hline
\hline
\end{longtable}
講題:從相分之痛到相聚之盼 Suffering and Eschatological Reconciliation

講員:林豪恩先生

經文:耶利米書23:3-6

場合:香港中文大學崇基學院禮拜堂主日崇拜

日期:2019年11月24日

一.引言:離散的經驗

「從相分之痛到相聚之盼」這講題不是為最近的情況而定的, 因為我們兩個月前就已經要提交講道題目.經文也不是為目前的狀況而挑選的,因為我們按照經課講道,這經課是普世教會共用,已經編排好三年一個循環的.如果今天的經文及課題配合我們當下的處境,或許是上主對我們的恩典.

大家仍記得這段禱文嗎?

上主啊:

我們在暫停之地,坐在那裡,追想中文大學和崇基校園,就哭了.在一排排障礙物之外,我們關上我們的風琴,我們不能同唱我們的歌.

上主啊:

我們怎能忘記你呢?我們怎能不看你過於我最喜樂的呢?

上主啊:

我們能夠在崇基禮拜堂外唱耶和華的歌嗎?你是全地的主,你聽到我們在任何一處向你發出的歌聲及祈求.

這是我們上一個主日的禱文.因為崇基禮拜堂不能夠如常進行主日崇拜,我們不能共聚一起,不知道會否有人感到分離的無奈及傷感.

這種離散的突然及傷感,在那個星期特別濃烈.十一月中,學期過了一大半,同學們開始適應及投入學習和校園生活,突然之間,來自不同地區和國家的學生,拖著行李,從宿舍步行往大學出口,上山下山,經過重重路障,到校園出口,再步行往有交通工具可到達的位置,經陸路或水路離開大學,再以各種方法離開香港.居住在宿舍的教職員家庭,扶老攜幼,對長者和小童來說,像是攀山越嶺,撤離大學尋找臨時居所.有些同事在網上分享表示有走難的感覺.

說到走難,年長的香港人其實不陌生；分離,更是連續幾代香港人的經驗.七十多年前二次大戰日軍攻佔香港,有些在香港生活的人走難上內地,回鄉逃避戰火.五十年代內地政治運動一波未平一波又起,不少人南下香港,他們當時的計劃只是暫避,等待局勢穩定後回家重聚.六七十年代文化大革命狂潮,能夠走得動的都逃難,不少人家離子散.由於分離在毫無準備之下發生,家人都不知何時再可以再通音訊,更不知何時能夠重聚.

有些離散家庭後來漸漸在香港重聚,打算落地生根之際,卻又要再上路.就在我們讀中學的那些年,常常聽到有同學舉家移民,這可能是我們那一代的共同回憶.這次離開雖然是有計劃,然而,分離所產生的家庭和社會問題,也為不少人帶來困擾,例如「太空人」的孤寂和引誘.

「但願人長久,千里共嬋娟」,古今中外的人們都經歷分離.今天的經課或者能夠與我們產生一點共鳴.上主如何回應人間的分離呢?能夠賜給我們怎樣的安慰和盼望呢?

二. 經文與內容

1.耶利米書:羊群的分離與重聚

23:1耶和華說:「那些殘害,趕散我草場之羊的牧人有禍了!」

23:2耶和華－以色列的上帝斥責那些牧養他百姓的牧人,如此說:「你們趕散我的羊群,並沒有看顧他們；我必討你們這行惡的罪.這是耶和華說的.

23:3我要將我羊群中所餘剩的,從我趕他們到的各國內招聚出來,領他們歸回本圈；他們也必生養眾多.

23:4我必設立照管他們的牧人,牧養他們.他們不再懼怕,不再驚惶,也不缺少一個；這是耶和華說的.」

23:5耶和華說:「日子將到,我要給大衛興起一個公義的苗裔；他必掌王權,行事有智慧,在地上施行公平和公義.

23:6在他的日子,猶大必得救,以色列也安然居住.他的名必稱為『耶和華－我們的義』.」

在耶利米的時代,為甚麼有離散的事情發生呢?人民為甚麼流離失所呢?耶和華上帝指出,是因為管治者失責,因為領袖失職.

「那些殘害,趕散我草場之羊的牧人有禍了!」(23:1)

羊群是誰?「耶和華－以色列的上帝斥責那些牧養他百姓的牧人」(23:2)經文所講的羊群是當時的百姓,不是現在我們說的信徒.牧人又是誰呢?(2:8祭司都不說,耶和華在哪裡呢?傳講律法的都不認識我.官長違背我；先知藉巴力說預言,隨從無益的神.)是當時的宗教,社會及政治領袖,不是現代教會的牧師.因此,耶利米論及的是國家,不是我們現在的教會.國家或地方的領袖失職,可以為人民帶來極大的痛苦.在耶利米時代,甚至家破人亡,四處逃命,或被擄他方.這些悲劇在歷史中並不陌生,在地球上從未停止.

社會崩壞,人民流離失所,為甚麼向領袖問責呢?常言道,權力越大,貢獻可以越大,遺害也可以更大.人民承受他們的錯失帶來的痛苦,無辜又無奈,不知如何是好.這一切,耶和華上帝都看見,並揚言會給他們應有的報應:

「我必討你們這行惡的罪」(23:2)

除此之外,經文中尚有證據證明耶和華上帝把社會動盪,人民流離失所的慘況歸因於管治者,因為經文中提供的答案就是新的管治者:

23:4我必設立照管他們的牧人,牧養他們.他們不再懼怕,不再驚惶,也不缺少一個；這是耶和華說的.」

23:5耶和華說:「日子將到,我要給大衛興起一個公義的苗裔；他必掌王權,行事有智慧,在地上施行公平和公義.

人民如何能夠「不再懼怕,不再驚惶,也不缺少一個」呢?經文指出領袖或管治者「行事有智慧,在地上施行公平和公義」.

2.詩篇:萬國的分離與重聚

46:1上帝是我們的避難所,是我們的力量,是我們在患難中隨時的幫助.

46:2所以,地雖改變,山雖搖動到海心,

46:3其中的水雖匉訇翻騰,山雖因海漲而戰抖,我們也不害怕.(細拉)

46:4有一道河,這河的分汊使上帝的城歡喜；這城就是至高者居住的聖所.

46:5上帝在其中,城必不動搖；到天一亮,上帝必幫助這城.

46:6外邦喧嚷,列國動搖；上帝發聲,地便鎔化.

46:7萬軍之耶和華與我們同在；雅各的上帝是我們的避難所!(細拉)

46:8你們來看耶和華的作為,看他使地怎樣荒涼.

46:9他止息刀兵,直到地極；他折弓,斷槍,把戰車焚燒在火中.

46:10你們要休息,要知道我是上帝!我必在外邦中被尊崇,在遍地上也被尊崇.

46:11萬軍之耶和華與我們同在；雅各的上帝是我們的避難所!

詩篇46篇充滿安慰和激勵,其實,讀者是能夠從字裡行間感受到那動盪的局勢和戰爭的味道,也能夠想像到人民驚惶失措的不安狀況.相傳詩篇的背景是列王記下18:13-19:37或以賽亞書36-37章所記的事.希西家王四十年,亞述王西拿基立派大軍來攻打猶大.希西家因為懼怕,就將殿裡所有的金子銀子都送給亞述.但亞述王後來又派將軍拉伯沙基來罵陣,譏笑他們不夠可騎二千匹馬的人；又勸他們投降,因他們的神耶和華不能救他們,因沒有一國的神能救他們脫離他的手.希西家聽了這話,極其難過,就上聖殿去禱告.本詩就是希西家去禱告時,詩人對王所唱的詩.

除了耶利米書所指出國家內部問題,人民離散的另外一個原因是國與國之間的戰爭.聖經記載的被稱為上帝選民的以色列人經歷過戰爭,歷世歷代普天下不少人也經歷戰爭,第一次及第二次世界大戰雖然已經在歷史書中,但就是現在,戰爭仍然是人民離散,流離失所的其中一個主要原因.

過去幾年十大能民國之首,你知道是來自哪一個國家的人民嗎?是敘利亞.為甚麼呢?相信大家都知道原因的.敘利亞人要逃難不是因為貧窮,而是因為打仗,家園被毀；他們有些可能是在自己的地區被攻打之前收到消息,預先逃走.戰事開始前,敘利亞雖然不算是一個有錢的國家,但也不太窮.世界銀行把敘利亞的國民收入劃分為中下,2010年戰事開始前,敘利亞人均GDP為＄2803美元,以購買力平價計算的人均GDP(PPP)大概在今天的越南和緬甸之間. 集體離開家園,只因為逃避戰火.(港台電視視點31)

內部管治失效,導致政局混亂,民不聊生；國與國之間戰火連年,生命及家園都受到威脅,人民流離失所,骨肉分離,上主看見嗎? 上主知道人們分離的痛苦嗎?上主如何回應呢?

3.路加福音:聖子與聖父的分離與重聚

23:33到了一個地方,名叫「髑髏地」,就在那裡把耶穌釘在十字架上,又釘了兩個犯人:一個在左邊,一個在右邊.

23:34當下耶穌說:「父啊!赦免他們；因為他們所做的,他們不曉得.」兵丁就拈鬮分他的衣服.

23:35百姓站在那裡觀看.官府也嗤笑他,說:「他救了別人；他若是基督,上帝所揀選的,可以救自己吧!」

23:36兵丁也戲弄他,上前拿醋送給他喝,

23:37說:「你若是猶太人的王,可以救自己吧!」

23:38在耶穌以上有一個牌子寫著:「這是猶太人的王.」

23:39那同釘的兩個犯人有一個譏笑他,說:「你不是基督嗎?可以救自己和我們吧!」

23:40那一個就應聲責備他,說:「你既是一樣受刑的,還不怕上帝嗎?

23:41我們是應該的,因我們所受的與我們所做的相稱,但這個人沒有做過一件不好的事.」

23:42就說:「耶穌啊,你得國降臨的時候,求你記念我!」

23:43耶穌對他說:「我實在告訴你,今日你要同我在樂園裡了.」

官府嗤笑得有道理:「他救了別人；他若是基督,上帝所揀選的,可以救自己吧!」何以基督不救自己脫離十字架呢?

兵丁戲弄他似乎也說得通:「你若是猶太人的王,可以救自己吧!」何以王不救自己脫離被戲弄的狀況呢?

聖子耶穌基督的痛苦是被鞭打,皮開肉裂,被釘上十字架,喘一口氣也痛不欲生,然而,這都只是身體的痛.被嘲笑,被戲弄,是尊嚴的痛.與門徒分離,被門徒出賣呢?如果這些都算是外人,那麼,肉身的母親馬利亞看著自己被掛在十字架上,於心何忍呢?你能夠想像被掛在十字架上的耶穌看來下面傷心欲絕的媽媽會是怎樣的心情嗎?打在兒身?痛在父母心,這是為人父母都能夠同感的.然而,最痛若的,是聖子與聖父的分離,如此的分離,連聖子也發出哀號:

馬可福音15章34節:到了下午三點鐘,耶穌大聲喊:「以羅伊,以羅伊,拉馬撒巴各大尼?」意思是:「我的上帝,我的上帝,你為什麼離棄我?」

潘霍華認為,主耶穌在十架上痛苦叫喊,是聖經中最重要的一句話.這是耶穌一生最堅定的時刻.當上帝讓耶穌承擔人類最可怕的痛,最深沉的悲傷,包括死亡時,祂也必須經歷遭上帝遺棄,甚至喪失信心的時刻,因為這就是人生的實況.

人世間歷盡離散之痛,天地之主竟然也參與在這痛苦之中,或者沒有人會明白,沒有人會理解,甚至很多人不接受.是的,正如耶穌在十字架上所說:

「父啊!赦免他們；因為他們所做的,他們不曉得.」(23:34)

他們不曉得,我們曉得嗎?

4.歌羅西書:在十字架上所成就的萬物重聚

1:11照他榮耀的權能,得以在各樣的力上加力,好叫你們凡事歡歡喜喜地忍耐寬容；

1:12又感謝父,叫我們能與眾聖徒在光明中同得基業.

1:13他救了我們脫離黑暗的權勢,把我們遷到他愛子的國裡；

1:14我們在愛子裡得蒙救贖,罪過得以赦免.

1:15愛子是那不能看見之上帝的像,是首生的,在一切被造的以先.

1:16因為萬有都是靠他造的,無論是天上的,地上的；能看見的,不能看見的；或是有位的,主治的,執政的,掌權的；一概都是藉著他造的,又是為他造的.

1:17他在萬有之先；萬有也靠他而立.

1:18他也是教會全體之首.他是元始,是從死裡首先復生的,使他可以在凡事上居首位.

1:19因為父喜歡叫一切的豐盛在他裡面居住.

1:20既然藉著他在十字架上所流的血成就了和平,便藉著他叫萬有－無論是地上的,天上的－都與自己和好了.

在動盪不安的環境,在矛盾衝突尖銳的時刻,在禮崩樂壞的社會,在暴力與仇恨威脅之下,我們心中的遠像是怎樣的呢?

耶利米先知對被趕散的以色列人說:

23:3我要將我羊群中所餘剩的,從我趕他們到的各國內招聚出來,領他們歸回本圈；他們也必生養眾多.

23:4我必設立照管他們的牧人,牧養他們.他們不再懼怕,不再驚惶,也不缺少一個；這是耶和華說的.」

23:5耶和華說:「日子將到,我要給大衛興起一個公義的苗裔；他必掌王權,行事有智慧,在地上施行公平和公義.

23:6在他的日子,猶大必得救,以色列也安然居住.他的名必稱為『耶和華－我們的義』.」

詩人對大軍壓境的猶太王說:

46:9他止息刀兵,直到地極；他折弓,斷槍,把戰車焚燒在火中.

46:10你們要休息,要知道我是上帝!我必在外邦中被尊崇,在遍地上也被尊崇.

46:11萬軍之耶和華與我們同在；雅各的上帝是我們的避難所!

耶穌對十字架上認罪悔改的犯人說:

23:43「我實在告訴你,今日你要同我在樂園裡了.」

縱使現場的人不去體諒聖子與聖父分離的痛苦,也不明白當中的意義；世上始終有人明白,使徒保羅說:

1:20既然藉著他在十字架上所流的血成就了和平,便藉著他叫萬有－無論是地上的,天上的－都與自己和好了.

在我們的前景中,有沒有和平的想像呢?在我們的想像中,有沒有和好的圖畫呢?在我們的圖畫中,有沒有包含眾生呢?先知,詩人,使徒都為我們描繪了那幅遠像,聖子耶穌基督親自為成就這遠象而親歷最深刻的分離之痛,帶領我們向著這遠像而前進.此時此地,烽煙四起,前路矇矓,但願我們心裡的道路仍然清晰,我們心裡的方向仍然堅定.



\newpage

\section{超越的愛}
\label{sec:479}
\textbf{陳衍昌法政牧師 - 超越的愛}
\newline
\newline
\hyperref[sec:478]{< < < PREV SERMON < < <}
~
\hyperref[sec:toc]{[返目錄]}
~
\hyperref[sec:480]{> > > NEXT SERMON > > >}


哥林多前書 13:1-13
\newline
\begin{longtable}{cl}
\hline
\hline
章節 & 經文 (和合本修訂版)\\
\hline
13:1 & \begin{tabularx}{0.7\textwidth}{X} 我若能說人間的方言，甚至天使的語言，卻沒有愛，我就成為鳴的鑼、響的鈸一般。 \end{tabularx} \\ \\ \relax
13:2 & \begin{tabularx}{0.7\textwidth}{X} 我若有先知講道的能力，也明白各樣的奧祕，各樣的知識，而且有齊備的信心，使我能夠移山，卻沒有愛，我就算不了甚麼。 \end{tabularx} \\ \\ \relax
13:3 & \begin{tabularx}{0.7\textwidth}{X} 我若將所有的財產救濟窮人，又犧牲自己的身體讓人誇讚，卻沒有愛，仍然對我無益。 \end{tabularx} \\ \\ \relax
13:4 & \begin{tabularx}{0.7\textwidth}{X} 愛是恆久忍耐；又有恩慈；愛是不嫉妒；愛是不自誇，不張狂， \end{tabularx} \\ \\ \relax
13:5 & \begin{tabularx}{0.7\textwidth}{X} 不做害羞的事，不求自己的益處，不輕易發怒，不計算人的惡， \end{tabularx} \\ \\ \relax
13:6 & \begin{tabularx}{0.7\textwidth}{X} 不喜歡不義，只喜歡真理； \end{tabularx} \\ \\ \relax
13:7 & \begin{tabularx}{0.7\textwidth}{X} 凡事包容，凡事相信，凡事盼望，凡事忍耐。 \end{tabularx} \\ \\ \relax
13:8 & \begin{tabularx}{0.7\textwidth}{X} 愛是永不止息。先知講道之能終必歸於無有；說方言之能終必停止；知識也終必歸於無有。 \end{tabularx} \\ \\ \relax
13:9 & \begin{tabularx}{0.7\textwidth}{X} 我們現在所知道的有限，先知所講的也有限， \end{tabularx} \\ \\ \relax
13:10 & \begin{tabularx}{0.7\textwidth}{X} 等那完全的來到，這有限的必消逝。 \end{tabularx} \\ \\ \relax
13:11 & \begin{tabularx}{0.7\textwidth}{X} 我作孩子的時候，說話像孩子，心思像孩子，意念像孩子；既長大成人，就把孩子的事丟棄了。 \end{tabularx} \\ \\ \relax
13:12 & \begin{tabularx}{0.7\textwidth}{X} 我們現在是對著鏡子觀看，模糊不清；到那時，就要面對面了。我如今所認識的有限，到那時就全認識，如同主認識我一樣。 \end{tabularx} \\ \\ \relax
13:13 & \begin{tabularx}{0.7\textwidth}{X} 如今常存的有信，有望，有愛這三樣，其中最大的是愛。 \end{tabularx} \\ \\
[1ex]
\hline
\hline
\end{longtable}
講題:超越的愛 The Love that transcends

經文:哥林多前書13:1-13

講員:陳衍昌法政牧師

場合:香港中文大學崇基學院禮拜堂主日崇拜

日期:2019年12月1日

一.引: 超越的愛

愛,具有超越,改變,轉化,提升的力量.

約翰21:15-19 ；林前13:1-13 愛篇

愛,是「道成肉身」；多於意念本質,必須付諸行動:

愛是放下,進入處境.本色化,見於「道成肉身」的行為.在聖公會慶祝在華100周年的研討會,本人以Trinity inspires: the witness of a young cathedral 為題分享本色化的信仰意義,把「道成肉身」理解為「拒絕把生命鎖定於在時間和空間的某一點」.唯有這樣,我們的生命才會顯得更豐富,正如主耶穌自己所說:「我來了,是要叫人得生命,並且得的更豐盛.」約10:10下

二.基督的愛超越個人,找到和歸屬家

更豐盛生命的其中一個指標,可以說是超越個人的局限,找到更大的歸屬感.聖保羅在以弗所書有以下的教導:

「這樣,你們不再是外人或客旅,是與聖徒同國,是上帝家裡的人了.」

以弗所書 2:19 RCUV

\url{https://www.bible.com/139/eph.2.19.rcuv}

在2014年,我在英國南部一間曾經事奉的聖公會教會崇拜,親身體驗到聖保羅這些說話的含意.當天的講員是剛剛從對衡大學(University of Durham)退休,大名鼎鼎的新約聖經學者鄧雅各(Professor James Dunn),崇拜後茶會.他知道我認識他一位學生盧龍光牧師,傾談間他告訴我:我是循道會會友,在這裡找到了家.

鄧教授的話肯定了基督的愛,叫我們連繫和歸屬於以祂的愛為根基的家.

因此,愛首先叫人超越自我,進入以愛為基的家.努力學習以愛的行動,放下以往麻木不仁,向有需要的人伸出援手.正如我個人十分喜歡的祝福內容:平安進入世界,勇敢為人,保持善德,勿以惡報惡,鼓勵灰心者,扶持軟弱者,援助困苦者,愛主事主,在聖靈的能力裡歡欣喜樂.

三.愛叫人認識人生階段的無知

「我作孩子的時候,說話像孩子,心思像孩子,意念像孩子；既長大成人,就把孩子的事丟棄了. 我們現在是對著鏡子觀看,模糊不清 ；到那時,就要面對面了.我如今所認識的有限,到那時就全認識,如同主認識我一樣.」

哥林多前書 13:11-12 RCUV

\url{https://www.bible.com/139/1co.13.11-12.rcuv}

聖保羅這些話,提醒我們人類知識的有限, 有幼兒到少年,少年到青年,中年到晚年,都有階段的局限.然而, 最大的謬誤是基於對自我形象的了解.

離世不久的現代天主教教育家保羅費哀里(Paulo Freire)對人性的體會是最為透徹的洞悉:世代以來各國各族之中,不仁不義的世道人心,導致弱者備受欺壓的情況無處不在；基於這些情況,他認為教育的重要任務就是喚醒良知, 在他看來,喚醒良知的教育,正人類與生俱來的召喚(ontological vocation).證諸於發生於香港各區,中大和理大的近日事件,豈非至理名言?

四.愛呼喚人以基督的心為心

「你們當以基督耶穌的心為心: 他本有上帝的形像, 卻不堅持自己與上帝同等 ； 反倒虛己, 取了奴僕的形像, 成為人的樣式； 既有人的樣子, 就謙卑自己, 存心順服,以至於死, 且死在十字架上.」

腓立比書 2:5-8 RCUV

\url{https://www.bible.com/139/php.2.5-8.rcuv}

大理國君主段(智興)王爺 ->一燈大師

瑛姑(劉瑛),嬰兒被裘千仞鐵掌打傷,段王爺拒絕為嬰兒療傷.瑛姑懷恨在心,設計幫助郭晴黃蓉闖過 一燈大師四大弟子:漁,樵,耕,謮闞口,獲得一燈大師,現為南僧(前稱南帝)動慈心用一陽指為黃蓉療傷.瑛姑則算準時間,現身攻擊一燈.一燈原來一早己領到她的計謀,欣然面對他捨王位作僧人前的愛妃瑛姑的復仇.這故事告訴我原來愛能夠叫人徹悟改變.

事實上,新約聖經約翰福音也記載了相同的故事:

「他們吃完了早飯,耶穌對 西門．彼得 說:「 約翰 的兒子 西門 ,你愛我比這些更深嗎?」 彼得 對他說:「主啊,是的,你知道我愛你.」耶穌對他說:「你餵養我的小羊.」 耶穌第二次又對他說:「 約翰 的兒子 西門 ,你愛我嗎?」 彼得 對他說:「主啊,是的,你知道我愛你.」耶穌說:「你牧養我的羊.」 耶穌第三次對他說:「 約翰 的兒子 西門 ,你愛我嗎?」 彼得 因為耶穌第三次對他說「你愛我嗎」,就憂愁,對耶穌說:「主啊,你無所不知,你知道我愛你.」耶穌說:「你餵養我的羊. 我實實在在地告訴你,你年輕的時候,自己束上帶子,隨意往來；但年老的時候,你要伸出手來,別人要把你束上,帶你到不願意去的地方.」

約翰福音21:15-18 RCUV

\url{https://www.bible.com/139/jhn.21.15-18.rcuv}

在這裡,我想到聖保羅寫給哥林多信徒的另一段說話:

「你們知道我們主耶穌基督的恩典:他本是富足,卻為你們成了貧窮,好使你們因他的貧窮而成為富足.」

哥林多前書8:9 RCUV

\url{https://www.bible.com/139/2co.8.9.rcu}

近兩年,我在溫哥華女兒和女婿的見證,叫我對這經文有新的體會.他們為了讓大女兒在數學科趕上水平,把她送到每學費高達加幣40000元的私校上學,Eaten Arrowsmith Magnussen School.

她現在這學校完成了兩年的課程,今年將會是最後一年,然後便會回歸主流學校.她的經驗告訴我們,她缺乏信心的數理科目都明顯進步了.

五.結語:鳴鑼響鈸的警告

我家一對書架上文字很矛盾:Do as I say, don’t do as I do!這是基督徒的大忌.

多年前遊劍橋,發現了一段有趣的諺語:

為什麼學習?Why study?

學習越多,懂得越多；The more I study, the more I know.

懂得越多,忘掉越多；The more I know, the more I forget.

忘掉越多,懂得越少；The more I forget, the less I know.

那麼,為什麼學習?. Then, why study?

事實上,這幾句說話,和愛篇起頭的金句有異曲同工之妙；姑且稱之為:

為什麼為善?

我若能說人間的方言,甚至天使的語言,卻沒有愛,我就成為鳴的鑼,響的鈸一般. 我若有先知講道的能力,也明白各樣的奧祕,各樣的知識,而且有齊備的信心,使我能夠移山,卻沒有愛,我就算不了甚麼. 我若將所有的財產救濟窮人,又犧牲自己的身體讓人誇讚 ,卻沒有愛,仍然對我無益.」

哥林多前書 13:1-3 RCUV

\url{https://www.bible.com/139/1co.13.1-3.rcuv}

談到這裡,在我們又回到了原來的起點:愛是行動.

不錯,愛,特別是基督徒的愛,並不可以停留在抽象的本質和意念之中.

這一點,著名的文學家和神學大師魯益師(C.S. Lewis)在《四種愛》(The Four Loves) 的序言中,用一個比喻說得最清楚,他指出:透過實踐上的努力,讓自己一步一步靠近主.

假設我們的家位在山谷裡,而我們現正走在回家的山道上.中午時分,我們來到一處崖頂,我們的家就位於山崖底下.從空間上來說,我們這時離自己的家非常接近,擲一顆石子下去,甚至可以打得著屋頂.但我們不是攀岩專家,無法就這樣垂直往下爬.要回家,只能繞路而行.在迂迥續

繞道的過程中,我們反而比站在崖頂上的時候離家更遠了.

站在崖頂上,我們確實離家更近,但這種近是沒有意義的,因為我們站在崖頂上愈久,就愈晚才回得了家.

根據創世記的說法,我們與生俱來帶有上帝的形象；不過單靠與上帝形象上的相似,並沒法使我們多接近多一些.只有在行動上的接近,實踐上的努力我們才有望更接近祂.

讓我們同心祈求上主,幫助我們常以鳴鑼響鈸為警誡,努力尋求行動的見證,讓愛的行動改變自己,帶動自己,超越原來的我,每一天,每一刻都認主日深,愛主日親,隨主日近!



\newpage

\section{跟從我吧}
\label{sec:480}
\textbf{劉國偉先生 - 跟從我吧!}
\newline
\newline
\hyperref[sec:479]{< < < PREV SERMON < < <}
~
\hyperref[sec:toc]{[返目錄]}
~
\hyperref[sec:481]{> > > NEXT SERMON > > >}


約翰福音 21:1-19
\newline
\begin{longtable}{cl}
\hline
\hline
章節 & 經文 (和合本修訂版)\\
\hline
21:1 & \begin{tabularx}{0.7\textwidth}{X} 這些事以後，耶穌在提比哩亞海邊又向門徒顯現。他怎樣顯現記在下面。 \end{tabularx} \\ \\ \relax
21:2 & \begin{tabularx}{0.7\textwidth}{X} 西門‧彼得、叫低土馬的多馬、加利利的迦拿人拿但業、西庇太的兩個兒子，和另外兩個門徒，都在一起。 \end{tabularx} \\ \\ \relax
21:3 & \begin{tabularx}{0.7\textwidth}{X} 西門‧彼得對他們說：「我打魚去。」他們對他說：「我們也和你一起去。」他們就出去，上了船；那一夜並沒有打著甚麼。 \end{tabularx} \\ \\ \relax
21:4 & \begin{tabularx}{0.7\textwidth}{X} 天剛亮的時候，耶穌站在岸上，門徒卻不知道他是耶穌。 \end{tabularx} \\ \\ \relax
21:5 & \begin{tabularx}{0.7\textwidth}{X} 耶穌就對他們說：「孩子們！你們有吃的沒有？」他們回答他：「沒有。」 \end{tabularx} \\ \\ \relax
21:6 & \begin{tabularx}{0.7\textwidth}{X} 耶穌對他們說：「你們把網撒在船的右邊，就會得到。」於是他們撒下網去，竟拉不上來了，因為魚很多。 \end{tabularx} \\ \\ \relax
21:7 & \begin{tabularx}{0.7\textwidth}{X} 耶穌所愛的那門徒對彼得說：「是主！」那時西門‧彼得赤著身子，一聽見是主，就束上外衣，跳進海裡。 \end{tabularx} \\ \\ \relax
21:8 & \begin{tabularx}{0.7\textwidth}{X} 其餘的門徒因離岸不遠，約有二百肘，就坐著小船把那網魚拉過來。 \end{tabularx} \\ \\ \relax
21:9 & \begin{tabularx}{0.7\textwidth}{X} 他們上了岸，看見那裡有炭火，上面有魚和餅。 \end{tabularx} \\ \\ \relax
21:10 & \begin{tabularx}{0.7\textwidth}{X} 耶穌對他們說：「把剛才打的魚拿幾條來。」 \end{tabularx} \\ \\ \relax
21:11 & \begin{tabularx}{0.7\textwidth}{X} 西門‧彼得就上船，把網拉到岸上，網裡滿了大魚，共一百五十三條；雖然魚這樣多，網卻沒有破。 \end{tabularx} \\ \\ \relax
21:12 & \begin{tabularx}{0.7\textwidth}{X} 耶穌對他們說：「你們來吃早飯。」門徒中沒有一個敢問他：「你是誰？」因為他們知道他是主。 \end{tabularx} \\ \\ \relax
21:13 & \begin{tabularx}{0.7\textwidth}{X} 耶穌走過來，拿餅給他們，也照樣拿魚給他們。 \end{tabularx} \\ \\ \relax
21:14 & \begin{tabularx}{0.7\textwidth}{X} 耶穌從死人中復活後向門徒顯現，這是第三次。 \end{tabularx} \\ \\ \relax
21:15 & \begin{tabularx}{0.7\textwidth}{X} 他們吃完了早飯，耶穌對西門‧彼得說：「約翰的兒子西門，你愛我比這些更深嗎？」彼得對他說：「主啊，是的，你知道我愛你。」耶穌說：「你餵養我的小羊。」 \end{tabularx} \\ \\ \relax
21:16 & \begin{tabularx}{0.7\textwidth}{X} 耶穌第二次又對他說：「約翰的兒子西門，你愛我嗎？」彼得對他說：「主啊，是的，你知道我愛你。」耶穌說：「你牧養我的羊。」 \end{tabularx} \\ \\ \relax
21:17 & \begin{tabularx}{0.7\textwidth}{X} 耶穌第三次對他說：「約翰的兒子西門，你愛我嗎？」彼得因為耶穌第三次對他說「你愛我嗎」，就憂愁，對耶穌說：「主啊，你無所不知，你知道我愛你。」耶穌說：「你餵養我的羊。 \end{tabularx} \\ \\ \relax
21:18 & \begin{tabularx}{0.7\textwidth}{X} 我實實在在地告訴你，你年輕的時候，自己束上帶子，隨意往來；但年老的時候，你要伸出手來，別人要把你束上，帶你到不願意去的地方。」 \end{tabularx} \\ \\ \relax
21:19 & \begin{tabularx}{0.7\textwidth}{X} 耶穌說這話，是指彼得會怎樣死來榮耀神。說了這話，耶穌對他說：「你跟從我吧！」 \end{tabularx} \\ \\
[1ex]
\hline
\hline
\end{longtable}
講題:跟從我吧! Follow me!

經文:約翰福音21:1-19

講員:劉國偉先生

場合:香港中文大學崇基學院禮拜堂主日崇拜

日期:2019年12月8日

引言

今天所選取的經文,記載在約翰福音最後一章,亦是我自去年9月開始,在我們禮拜堂參與講道服侍以來的第十五篇講章,我們來到約翰福音最後一章,也是我一年多以來有關約翰福音講道的最後一篇.願上主恩待我們,當我們思想這十九節經文的時候,求聖靈感動和引導我們.

貼地的信仰經歷總是刻骨銘心

我從事學生工作接受查考聖經的訓練,加上幾十年來,在本地及境外帶領查經所累積的經驗,其中一個重要的觀察,就是查考聖經的人,很多時都會忽略了文本的細節,尤其是敘事經文.約翰福音21章記載基督在提比哩亞海邊顯現的事蹟,無論是作者使徒約翰下筆,或者是他口傳,由他的門徒或助手執筆,又甚至有些觀點認為這章經文是成書之後,才增訂的編輯之選,但一個最根本的基礎,就是這個事跡,正如福音書好些敘事經文,都是深藏當事人心靈深處的記憶,好些在讀者眼中看來不經意的筆觸和描述,卻原來蘊含叫人深思的空間.

較多學者認為約翰福音的成書日期是公元一世紀末,簡單直接來說,那是事情發生超過最少半個世紀之後的著作,老約翰這位極可能是最後辭世的使徒,在他年邁之時,寫下超逾半世紀的回憶.我認為能夠留在這位飽歷風霜的老使徒心坎裡的,應該是經過歲月清洗之後的刻骨銘心的記憶,是的,我相信貼地的信仰經歷總是刻骨銘心.

所以,我常常告訴自己,查考聖經時應在意一些敘事細節,細味作者留給我們的點點滴滴,因為對置身於現場的作者來說,那些細節可能是非常深刻的記憶,對於讀者而言,若能看透當中的一鱗半爪,我會形容為有若尋寶.

重操故業?

這些事以後,耶穌在提比哩亞海邊又向門徒顯現.他怎樣顯現記在下面:有西門‧彼得和稱為低土馬的多馬,並加利利的迦拿人拿但業,還有西庇太的兩個兒子,又有兩個門徒,都在一處.西門‧彼得對他們說:「我打魚去.」他們說:「我們也和你同去.」(約21:1-3上)

我聽過也閱讀過有信徒批評彼得和六位門徒重操故業,甚至貪愛世界! 他們認為既然彼得和六位門徒,當初回應耶穌的呼召,明明已經捨棄了漁船,捨棄了自己的職業,走上跟從耶穌的道路,參與傳天國福音的職事,就不應該再去捕魚了,這好像為自己的未來,留下了一個空間,方便自己走回頭路.

其實,這短短的兩節半經文,我們用平常心走進門徒當時的處境,就不會走得太遠,得出他們重操故業,甚至貪愛世界的結論.不是嗎? 彼得和團隊不能不顧及眼前的需要,他們打魚的根本動機當然包括了食物的需要,甚至藉此賺取一些生活費.按著這個我相信是較為合理的推論,充其量,彼得與門徒只不過是開展了一種類似自養自傳的做法,一方面為生計籌謀,另一方面,按基督的吩咐,在加利利等候主的吩咐.若果我們帶著同理心去了解門徒的處境,明白到彼得和門徒在等待復活的主的過程中,肚子還是會餓的,而且每天也有必須的生活開支,在團隊方向未明,管錢的那位更已經自盡己命的情況下,剩下來的門徒想想辦法去解決眼前生活所需,按他們的技能去賺取一點收入,若果我們投入門徒這個現實處境,就會對他們接下來的遭遇,有更深的體會.

沒有漁獲的晚上

他們就出去,上了船；那一夜並沒有打著甚麼.天將亮的時候,耶穌站在岸上,門徒卻不知道是耶穌.耶穌就對他們說:「小子!你們有吃的沒有?」他們回答說:「沒有.」耶穌說:「你們把網撒在船的右邊,就必得著.」(約3:3下-6)

約翰福音記載基督在提比哩亞海邊顯現的事蹟,有一個非常清晰的時序: 七位門徒夜裡打魚,經過整夜勞累,即使他們當中包括專業漁夫,卻竟然毫無收穫,天將亮的時候,他們就與耶穌相遇.

人生旅途常常遇上叫人洩氣的時候,有時我們整夜勞碌,卻似乎苦無收獲,即使本來有把握,有經驗的事,卻意外落空,投入的心力和時間,彷彿枉然!說到底,門徒只是為了生活去打漁而已……

很明顯的,「那一夜並沒有打著甚麼」這個結果,相信並非門徒所預期,彼得應當判斷那是一個可以正常捕魚的夜晚,按上下文記載,那個晚上也沒有狂風大雨,專業漁夫從事自己很有把握的工作,卻整夜勞累,毫無收穫.弟兄姊妹,你可有想過耶穌何不讓他們有少許收穫,不至於賠上整個晚上,卻仿似白做一場?又或者,耶穌何不早一點出手,讓門徒早一點收工?何苦要他們體力,意志,忍耐都近乎耗盡的時候才出手相助?

當然,我們作為旁觀者,我們知道耶穌站在岸上,約翰沒有告訴我們,到底耶穌在岸邊站了多久,他記載的重點是「門徒卻不知道是耶穌」.親愛的弟兄姊妹,這會否是我們的寫照?原來耶穌明明在我們當中,我們卻在生活的重重責任底下,追求人生成就的過程裡,在忙碌的生活節奏之中,甚至為生活奔波糊口的疲累裡,卻不察覺耶穌的存在!

親愛的弟兄姊妹,你又有沒有注意到,在那個黎明時份,耶穌對門徒說的第一句話是什麼?那並非「有漁獲嗎?」或者「收穫有多少?」,卻是「小子!你們有吃的沒有?」,你可曾細味過這句問候背後,所蘊藏對門徒的情愛?

原來耶穌在

他們便撒下網去,竟拉不上來了,因為魚甚多.耶穌所愛的那門徒對彼得說:「是主!」那時西門‧彼得赤著身子,一聽見是主,就束上一件外衣,跳在海裡.其餘的門徒離岸不遠,約有二百肘,就在小船上把那網魚拉過來.他們上了岸,就看見那裡有炭火,上面有魚,又有餅. (約21:6下-9)

這實在是十分感人的畫面:因為耶穌的介入,神蹟發生了,撒下去的網,竟然拉不上來了.門徒渡過那看似沒有漁獲的夜晚,卻在黑夜的盡頭,新一天正要來臨的時候,耶穌為門徒預備了滿網的大魚.

這時候,門徒認得是主耶穌,西門‧彼得就一馬當先,然後其餘的門徒陸續上岸,一個更感人的畫面出現了:那裡有炭火,上面有魚,又有餅.這對於整夜勞碌,本來毫無所獲的七位壯男,這是何等貼心的服侍!是主耶穌為門徒預備的早飯,是非常體貼的舉措,那是炭火仍在烤煮的熱食!

耶穌明知故問對門徒所說:「小子!你們有吃的沒有?」,原來向門徒提問以先,耶穌已經預設了這個體貼門徒的安排.信仰的經歷,當然包括捕獲的魚甚多那種澎湃,卻在人生的同一個段落裡,看見主為門徒預備早飯.那位創天造地,捨身十架的主,成就了何等偉大的救恩!就是這位行大而可畏的事的主,明白我們切身的需要,在我們的生活小節中,顯露祂的慈愛.但願門徒和彼得的醒覺:「是主!」,是我們的每天生活的重要發現.

活在現場中

耶穌對他們說:「把剛才打的魚拿幾條來.」西門‧彼得就去( - 或譯:上船),把網拉到岸上.那網滿了大魚,共一百五十三條；魚雖這樣多,網卻沒有破.耶穌說:「你們來吃早飯.」門徒中沒有一個敢問他:「你是誰?」因為知道是主. 耶穌就來拿餅和魚給他們.耶穌從死裡復活以後,向門徒顯現,這是第三次.(約21:10-14)

若果說耶穌的介入,讓門徒經歷這一個偉大的神蹟,叫人觸動；而耶穌為門徒預備早飯所顯露的憐愛,本來就已經叫人十分感動,但更加觸動我心的,是耶穌吩咐門徒:「把你們剛才打的魚,拿幾條來.」(約21:10) 從故事的發展而言,即使省卻了「把剛才打的魚拿幾條來」這個小節,整個記事仍然可以推進下去,反倒這個細節的出現,需要我們停下來想一想.

弟兄姊妹,你有沒有想過,既然耶穌要為門徒準備一頓貼心的早飯,為什麼耶穌不多帶幾條魚,甚至施行神蹟多變幾條魚出來,卻要門徒從滿網的大魚中,拿幾條過來?

敘事者沒有向讀者交代耶穌的動機,但文本卻清楚顯示,無論是捕獲153條大魚的神蹟,又或者要準備早飯這一件愛心美事,耶穌都在事件過程中騰出一些位置,讓門徒有份其中.

弟兄姊妹,原來神聖相遇是如此貼地,耶穌並沒有叫跟隨祂的門徒,單單去做信仰旅途的觀眾,安舒地坐在觀眾席上,欣賞耶穌的表演,然後我們鼓掌讚美,卻其實是耶穌把門徒拉進祂的大大小小的作為當中,要門徒成為一個活在現場的跟從者.

這個從深夜到黎明的場景,門徒經歷體力,意志,忍耐都近乎耗盡的艱辛,成就了隨之而來的刻骨銘深的信仰體會.福音信仰從來都不離地,卻很多時是因為我們將信仰和現場生活割離,然後反過來認為聖經的教導脫離現實.主耶穌就是如此貼地,體貼門徒面對生活的現實,賜他們百多條大魚,讓他們有份參與這個神蹟和這頓早飯.

早飯預備好以後,耶穌對門徒說:「你們來吃早飯.」然後,是耶穌親自拿餅和魚給門徒,對老約翰來說,這些點點滴滴,又成為他另一個深刻記憶.

撫心自問,我們可有注意過生活裡的點點滴滴,可有細味耶穌在我們的生活中留下的痕跡?信仰從來都不離地!我們千萬不要做自己信仰旅途上的旁觀者.

你跟從我吧!

他們吃完了早飯,耶穌對西門‧彼得說:「約翰的兒子西門,你愛我比這些更深嗎?」彼得說:「主啊,是的,你知道我愛你.」耶穌對他說:「你餵養我的小羊.」耶穌第二次又對他說:「約翰的兒子西門,你愛我嗎?」彼得說:「主啊,是的,你知道我愛你.」耶穌說:「你牧養我的羊.」第三次對他說:「約翰的兒子西門,你愛我嗎?」彼得因為耶穌第三次對他說「你愛我嗎」,就憂愁,對耶穌說:「主啊,你是無所不知的；你知道我愛你.」耶穌說:「你餵養我的羊.我實實在在地告訴你,你年少的時候,自己束上帶子,隨意往來；但年老的時候,你要伸出手來,別人要把你束上,帶你到不願意去的地方.」(耶穌說這話是指著彼得要怎樣死,榮耀上帝.)說了這話,就對他說:「你跟從我吧!」(約21:15-19)

彼得經過整夜勞碌,網到153條大魚,裹腹之後,天色已亮之際,耶穌跟彼得作出這次深情對話.耶穌向彼得三次查問「你我愛我比這些更深麼?」,然後叮囑彼得「你餧養我的羊.」

「你愛我比這些更深麼?」是生命中優先次序的考驗,我們也許會問自己,我愛這些愛那些,愛這個愛那個,到底是否愛主勝過愛其他一切.若果我們足夠誠實,要作出這個決定已經很不容易.然而,耶穌託付彼得「你餧養我的羊.」卻更加是一個涉及生與死的召命.

我們不禁想起,在耶穌的教導當中,牧羊人是要怎樣當的?耶穌說:「好牧人為羊捨命.」(約10:11),一個好的牧羊人以小羊的生命為自己的負擔,是生命與生命互相連結.主耶穌給彼得的挑戰,是叫彼得為主而活,為他者而生.

使徒約翰如此寫道:「耶穌說這話是指著彼得要怎樣死,榮耀上帝.說了這話,就對他說:「你跟從我吧!」(約21:19)  我就再次想起潘霍華的名句「當基督呼召一個人,祂是召他來為祂死.」耶穌在這個早上先後兩次對彼得說「你跟從我吧!」,另一次記載於21章22節,是福音書所記載耶穌向門徒所發出的最後的呼召,那是隱含彼得要跟隨至死的事實.感謝神,耶穌基督在歷世歷代繼續呼召門徒,在我們這個時代,照樣呼召你和我.

遊歷以色列有不少地方都叫我印象深刻,而且即使去了好幾次,每次都有新的感動.其中一處地方就是彼得受職堂(Church of St. Peter’s Primacy).彼得受職堂位於迦百農附近的塔加(Tabgha) ,由天主教聖方濟各會於1933年修建,位處加利利湖邊,湖水直接衝近教堂側的一個古代碼頭遺址,這景點盛載著我們今天所講述的事蹟.

本年5月,我第4次踏足那裡,那份現場感覺濃烈,我想起了主耶穌向彼得的呼召「你跟從我吧!」.

我求上帝在整夜勞碌,似乎苦無收獲,即使本來有把握,有經驗的事,卻意外落空,彷彿枉然之際,我定要仔細察看主的作為,深信在最適當的時候,主就介入那漫長的黑夜,讓我看見祂.

我再次提醒自己,是上帝的恩典,讓我可以參與祂的大能奇事,是祂預留我可以參與和服侍的空間,讓我不致於做信仰旅途上的旁觀者.弟兄姊妹,不要叫信仰和現場生活割離,千萬不要讓自己活像一個信仰旅途上的旁觀者.

耶穌基督在歷世歷代繼續呼召門徒,在我們這個時代,照樣呼召你和我.即使

看起來,好像又是一個沒有漁獲的夜晚,但我們知道,耶穌定睛看著我們,不論我們經歷過豐收的喜悅,也曾試過沿地踏步,甚至徒勞無功,但耶穌始終相伴,並呼召我們跟隨祂.



\newpage

\section{那將要來的是你嗎}
\label{sec:481}
\textbf{鄧瑞強博士 - 那將要來的是你嗎?}
\newline
\newline
\hyperref[sec:480]{< < < PREV SERMON < < <}
~
\hyperref[sec:toc]{[返目錄]}
~
\hyperref[sec:482]{> > > NEXT SERMON > > >}


馬太福音 11:2-11
\newline
\begin{longtable}{cl}
\hline
\hline
章節 & 經文 (和合本修訂版)\\
\hline
11:2 & \begin{tabularx}{0.7\textwidth}{X} 約翰在監獄裡聽見基督所做的事，就派他的門徒去， \end{tabularx} \\ \\ \relax
11:3 & \begin{tabularx}{0.7\textwidth}{X} 問耶穌：「將要來的那位就是你嗎？還是我們要等候另一位呢？」 \end{tabularx} \\ \\ \relax
11:4 & \begin{tabularx}{0.7\textwidth}{X} 耶穌回答他們：「你們去，把所聽見、所看見的告訴約翰： \end{tabularx} \\ \\ \relax
11:5 & \begin{tabularx}{0.7\textwidth}{X} 就是盲人看見，瘸子行走，痲瘋病人得潔淨，聾子聽見，死人復活，窮人聽到福音。 \end{tabularx} \\ \\ \relax
11:6 & \begin{tabularx}{0.7\textwidth}{X} 凡不因我跌倒的有福了！」 \end{tabularx} \\ \\ \relax
11:7 & \begin{tabularx}{0.7\textwidth}{X} 他們一走，耶穌就對眾人談到約翰，說：「你們從前到曠野去，是要看甚麼呢？看風吹動的蘆葦嗎？ \end{tabularx} \\ \\ \relax
11:8 & \begin{tabularx}{0.7\textwidth}{X} 你們出去到底是要看甚麼？看穿細軟衣服的人嗎？那穿細軟衣服的人是在王宮裡。 \end{tabularx} \\ \\ \relax
11:9 & \begin{tabularx}{0.7\textwidth}{X} 你們出去究竟是要看甚麼？是先知嗎？是的，我告訴你們，他比先知大多了。 \end{tabularx} \\ \\ \relax
11:10 & \begin{tabularx}{0.7\textwidth}{X} 這個人就是經上所說的：『看哪，我要差遣我的使者在你面前，他要在你前面為你預備道路。』 \end{tabularx} \\ \\ \relax
11:11 & \begin{tabularx}{0.7\textwidth}{X} 我實在告訴你們，凡女子所生的，沒有一個比施洗約翰大；但在天國裡，最小的比他還大。 \end{tabularx} \\ \\
[1ex]
\hline
\hline
\end{longtable}
講題:那將要來的是你嗎? Are you the one who is to come?

經文:馬太福音11:2-11

講員:鄧瑞強博士

場合:香港中文大學崇基學院禮拜堂主日崇拜

日期:2019年12月15日

各位弟兄姊妹,早安.

若你身染頑疾,臥病在床,在公立醫院裡,等著做一個生死攸關的重大手術.若你又知道,在這醫院裡,只有兩個醫生會做這手術.一個剛剛才做過一個同樣的手術,而病人死了.一個是資深醫生,做過無數次同樣的手術,而沒有一個病人死亡.你只知道這些資料,你也認得這兩個醫生,你只是不知道誰會負責你的手術.而剛好,那位資深醫生就站在你的病床邊.這時,你很可能會問的一個問題是:「那將要來的是你嗎?」這真是一個生死攸關的問題.

小時候,住在「屋村」,偶爾,不知什麼原因,一個樓層的小朋友會和另一個樓層的小朋友衝突起來,然後發生一些不太嚴重的毆鬥.有一次,衝突又起,我落了單,被另一個樓層的小朋友追著,遠處,見到有另一群小朋友迎過來,其中一個像是我的哥哥,我心想:「那將要來的是我的哥哥嗎?」雖然未能確定,但也只能走過去,因這是僅有的逃命機會.走得近了,才發現,這群小朋友也是另一個樓層的.最後,被圍困著,孤立無援,四面楚歌,被人打了一身.任何一個被圍困的人,若有陌生人接近,都會想這個問題:「那將要來的是你嗎?」這個「你」,當然就是幫助我們,解救我們的人.

已經是二十一世紀了,人類應該很文明的了,但人圍捕人的故事,仍經常發生.剛剛這個禮拜,緬甸的昂山素姬便要出席海牙聯合國國際法院的聆訊,為羅興亞人受的苦作出辯解.羅興亞人的故事,又是一群人被另一群人圍捕的故事.他們孤立無援,四面楚歌,連昂山素姬也承認軍方對他們或有使用過分武力.對於被圍捕的人而言,對於羅興亞人而言,他們期待的,是解救者.他們心中的問題,常常是:「那將要來的是你嗎?」那將要來的,是能解救他們脫離困局的解救者嗎?

香港,在一片迷惘裡,人與事都會急劇變化.或許,很多人會離開他們的位置,很多人會冒起來.不同的群體,等待著不同的真命天子.在亂世,我們都希望真正的解救者.用中國哲人張橫渠的講法,我們渴望有人能「為天地立心,為生民立命,為往聖繼絕學,為萬世開太平」.我們等待一些人,能實現上天仁愛的心,能尊重每個人的生命尊嚴,能守住支撐著人類文明的普世價值,能對歷史的將來負起真切的責任.關心香港的人,心中都會問:「那將要來的,是這樣的一個人嗎?」

經文:馬太福音11:2－11

太11:2　約翰在監裡聽見基督所做的事,就打發兩個門徒去,

太11:3　問他說:「那將要來的是你嗎?還是我們等候別人呢?」

太11:4　耶穌回答說:「你們去,把所聽見,所看見的事告訴約翰.

太11:5　就是瞎子看見,瘸子行走,長大痲瘋的潔淨,聾子聽見,死人復活,窮人有福音傳給他們.

太11:6　凡不因我跌倒的就有福了!」

太11:7　他們走的時候,耶穌就對眾人講論約翰說:「你們從前出到曠野是要看甚麼呢?要看風吹動的蘆葦嗎?

太11:8　你們出去到底是要看甚麼?要看穿細軟衣服的人嗎?那穿細軟衣服的人是在王宮裡.

太11:9　你們出去究竟是為甚麼?是要看先知嗎?我告訴你們,是的,他比先知大多了.

太11:10　經上記著說:『我要差遣我的使者在你前面預備道路』,所說的就是這個人.

太11:11　我實在告訴你們,凡婦人所生的,沒有一個興起來大過施洗約翰的；然而天國裡最小的比他還大.

這段經文,很自然地分成三個段落.

(1)約翰的問

太11:2　約翰在監裡聽見基督所做的事,就打發兩個門徒去,

太11:3　問他說:「那將要來的是你嗎?還是我們等候別人呢?」

這裡的約翰,是「約洗約翰」,他在約旦河裡為人施洗.他是一個正直的人,不畏強權.他直指統治者不道德,指斥統治者淫亂.當時的統治者是小希律,奪取了哥哥的妻子.約洗約翰直斥他的不是,直斥他毫無規矩,無法無天.於是,被統治者捉拿,下獄,最後也死在獄中,又是一個為了正義而死在獄中的人.

這樣的一個人,在獄中,當然有很多想法,很多期待.一個正義的人,縱然在獄中,也很留心監獄外社會正義的發展.他聽到耶穌的事情,看來,耶穌也推動著某種社會正義的發展.約洗約翰很著緊,耶穌這個人,是否是真命天子?看來就是他了,他能接續自己的信仰工程,宣揚天國的臨近,呼喚人悔改,在地上實現正義與和平.若真是他,則自己可安心了.若自己有甚麼不測,也有人承接自己的志向了.正義的人,在獄中,最關心的,不是自己在獄中,而是在監獄外的正義.這些人是真正自由的人.在不正義的世界裡,全然漠視世界的不正義,這些人才是在監獄裡.他們的心,不可能是自由的.在監獄裡,約洗約翰打發自己的門徒,去問耶穌:「那將要來的是你嗎?還是我們等候別人呢?」

(2)耶穌的答

太11:4　耶穌回答說:「你們去,把所聽見,所看見的事告訴約翰.

太11:5　就是瞎子看見,瘸子行走,長大痲瘋的潔淨,聾子聽見,死人復活,窮人有福音傳給他們.

太11:6　凡不因我跌倒的就有福了!」

耶穌的回答,是講論他自己所帶來的一個新世界,一個前所未有的世界.因著耶穌的出現,完全絕望的景況出現新機.因著耶穌,瞎子能看見.耶穌受難前夕,最後一個神蹟,就是醫治瞎子,讓瞎眼的能夠看見.先醫瞎子,後上十架,這故事編排有一重大含意,就是說,人若因某些因素令自己變成盲目,對發生在眼前的真相視而不見,則他不可能領悟十字架的意義.盲目,有很多種,肉眼的盲,是一種；良心的盲,是一種；靈性的盲,又是另一種.良心盲了,則在他眼前,發生一件滅絕天良的事,他也看不見,視為平常.靈性盲了,則正在步向滅亡,仍以為自己步向天堂.有時,人盲到一個地步,全世界都看得見的事,他卻視而不見.耶穌告訴施洗約翰,他會讓人看見真相,讓事情水落石出.

因著耶穌,瘸子能行走.瘸子是無能為力的人.約翰福音有一個醫治瘸子的經典故事(約5:1-9).耶路撒冷有一個水池,叫「畢士大」池.相傳,池水動的時候,是天使攪動池水,有病的人,誰先下水,誰就得醫治.耶穌來到這水池,見有一個瘸了三十八年的人.問他說:「你要痊愈嗎?」請留意這人的回答,他說:「先生,水動的時候,沒有人把我放在池子裡；我正去的時候,就有別人比我先下去.」多麼絕望的回答.完全沒法自救,也得不到別人的救助.在生活裡,完全無助.耶穌卻叫他起來行走,他就起來,走了.在我們社會,有很多叫天不應,叫地不聞的無助者,耶穌聽到他們的呼聲.耶穌告訴施洗約翰,他要讓無助的人得到救援,讓他們有力地活下去.

耶穌叫長大痲瘋的,得到潔淨.長大痲瘋的人,在當時,是untouchable的,是被隔離的,是被遺忘的.耶穌刻意的接觸他們,讓他們感覺到從人而來的溫暖.讓這些被非人化的人重拾人性尊嚴.就算在香港,都會有人被非人化地對待,被呼以污名,被凌辱,被虐待.耶穌告訴施洗約翰,他要讓這些人重拾人性尊嚴.

因著耶穌,聾子能聽見.有些聾子,是生出來的,這是悲劇.有些聾子,是裝出來的,這製造悲劇.你講甚麼,你問甚麼,你要求甚麼,他總是聽不見.耶穌卻讓聾子聽得見他人的聲音,重新恢復人間的溝通.耶穌告訴施洗約翰,他重建人間的對話,讓各人聽得明對方的心聲.

因著耶穌,死人能復活.死,是絕望事件,是終結性的,再無生機可言.耶穌,卻要在絕望中帶來生機.上個月,保安局局長說,今年6月至9月的自殺個案為256宗,比去年同期多34宗.發現屍體,送院前或送院時死亡個案為2537宗,比去年同期多311宗.就這些數字的表面來看,香港近來是多了些不太正常的死亡.這些死亡事件背後,特別是自殺事件,看來是一個又一個的絕望故事.自殺的人,是絕望的.耶穌告訴施洗約翰,他叫死人復活,他要取消絕望故事的絕望,讓人間看見前所未有的希望.

最後,耶穌說,有美好的信息傳給貧窮的人.通常,好的信息是沒有窮人的份的；他們有的,只是悲慘的生活.耶穌告訴施洗約翰,唉聲嘆氣的窮人,終於聽到美好的信息.

大家想想,因著耶穌,這世界發生了甚麼變化.一種強大的生命力,取消社會上的盲目,瘸子一樣的無助,長大痲瘋一樣的被非人化對待,拒絕溝通的耳聾,死亡的絕望,貧窮的淒慘.這種強大的仁愛力量,由耶穌那裡發源,藉聖靈進入我們的生命中.我們活在世界,正是要將這種仁愛力量繼續發揚出去,讓大地的黑暗,被光明驅散,讓天國的輪廓漸漸地顯明出來.

耶穌講完以上一番話,便說:「凡不因我跌倒的就有福了!」(太11:6)為何會因他的話跌倒?是因為他的話與這世界的話相差太遠,太不可信.這世界的人不相信窮人是有福的.這世界的人不相信死人可以復活.這世界的人不相信絕望的深淵還有希望.耶穌卻說,我要在死亡的絕望裡帶來復活的新希望.相信耶穌的話,需要堅持希望,需要因希望而忍耐,需要為希望而受苦.耶穌希望我們做得到,而不是在絕望中放棄.

耶穌以上述這番話回答施洗約翰,聽到這番話,施洗約翰應看到在耶穌身上天國已經開展.他不用再等別人了.這世界要等的,就是耶穌.

(3)耶穌的問與答

太11:7　他們走的時候,耶穌就對眾人講論約翰說:「你們從前出到曠野是要看甚麼呢?要看風吹動的蘆葦嗎?

太11:8　你們出去到底是要看甚麼?要看穿細軟衣服的人嗎?那穿細軟衣服的人是在王宮裡.

太11:9　你們出去究竟是為甚麼?是要看先知嗎?我告訴你們,是的,他比先知大多了.

太11:10　經上記著說:『我要差遣我的使者在你前面預備道路』,所說的就是這個人.

太11:11　我實在告訴你們,凡婦人所生的,沒有一個興起來大過施洗約翰的；然而天國裡最小的比他還大.

如今,輪到耶穌去評論這位在監獄裡仍掛心社會正義,仍掛心天國降臨的施洗約翰.這位施洗約翰,在位處曠野的約旦河為人施洗.他是誰,他為何能吸引這麼多人去受洗呢?

耶穌問:那些去受洗的人,他們去曠野,是看甚麼呢?要看風吹動的蘆葦嗎?看大自然的風景,不會吧!看像蘆葦一樣的普通人物,不會吧!要看穿細軟衣服的人嗎?穿得這般高貴的人,不會在曠野,他們在王宮裡.是要看先知嗎?這就對了.你們到曠野,要是聽先知在曠野的呼聲.你們要尋找,那展示高貴的道德,無偽的真理,純潔的神聖的先知.對了,你們尋找這些.你們渴求先知,施洗約翰的確是先知,並且是最大的先知,因為他為耶穌開路.耶穌是誰?是神在人間.他以死和復活,建立了天國.無論施洗約翰如何偉大,畢竟,只是耶穌建立的天國之前的人物.他無法體會,天國的真實展現,他無法體會世界因著耶穌而有的改變.今日,我們活在耶穌所開展的天國中,在耶穌的復活裡,我們體會神在死亡深處,在絕望深淵,依然掌權,依然能展示我們無法想像的生命,依然能綻放復活的光明.耶穌在評論施洗約翰時,他啟示了自己的真正身分及事工.

今日,香港仍在混亂中,仍有很多事情不明不白,前景仍然暗晦不明.我們會問:要等待誰來收拾殘局?我們會問耶穌:那將要來的是你嗎?

耶穌做了甚麼?他是我們要等的那個人嗎?

耶穌說:我叫瞎子看見,瘸子行走,長大痲瘋的潔淨,聾子聽見,死人復活,窮人有福音傳給他們.

他所做的,是徹底地改變了人性及世界.

我們不等他,還可以等誰呢?

作為基督徒,我們應該已體會過耶穌的大愛所造成的影響力,或許,與其問要等誰,不如問我們靠這位耶穌能多做一點甚麼.邪惡無法永存,天國總會實現.

但願:榮耀歸於聖父,聖子,聖靈.阿們.

(附記:本人在崇基禮拜堂講道的講章已結集成《荷花池畔春草綠》講章集,有興趣的朋友可到各大基督教書局購買.試讀本可參閱:\url{http://cc.cuhk.edu.hk/dev/images/content/fund-raising-campaigns/university-sermons-by-dr-sk-tang/CCC-SK-Tang-Sermon-Text-Sample-2.pdf}.較早期的崇基講章,已結集成《敢於跟隨主》一書,也可到各大基督教書局購買.)



\newpage

\section{等候神}
\label{sec:482}
\textbf{許立中先生 - 等候神}
\newline
\newline
\hyperref[sec:481]{< < < PREV SERMON < < <}
~
\hyperref[sec:toc]{[返目錄]}
~
\hyperref[sec:483]{> > > NEXT SERMON > > >}


以賽亞書 40:27-31
\newline
\begin{longtable}{cl}
\hline
\hline
章節 & 經文 (和合本修訂版)\\
\hline
40:27 & \begin{tabularx}{0.7\textwidth}{X} 雅各啊，你為何說，以色列啊，你為何言，「我的道路向耶和華隱藏，我的冤屈神並不查問」？ \end{tabularx} \\ \\ \relax
40:28 & \begin{tabularx}{0.7\textwidth}{X} 你豈不曾知道嗎？你豈未曾聽見嗎？永在的神耶和華，創造地極的主，他不疲乏，也不困倦；他的智慧無法測度。 \end{tabularx} \\ \\ \relax
40:29 & \begin{tabularx}{0.7\textwidth}{X} 疲乏的，他賜能力；軟弱的，他加力量。 \end{tabularx} \\ \\ \relax
40:30 & \begin{tabularx}{0.7\textwidth}{X} 就是年輕人也要疲乏困倦，強壯的也必全然跌倒。 \end{tabularx} \\ \\ \relax
40:31 & \begin{tabularx}{0.7\textwidth}{X} 但那等候耶和華的必重新得力。他們必如鷹展翅上騰；他們奔跑卻不困倦，行走卻不疲乏。 \end{tabularx} \\ \\
[1ex]
\hline
\hline
\end{longtable}
講題:等候神 Waiting on God

經文:以賽亞書40章27-31節

講員:許立中先生

場合:香港中文大學崇基學院禮拜堂主日崇拜

日期:2019年12月22日

「雅各阿,你為何說,我的道路向耶和華隱藏?以色列阿,你為何言,我的冤屈神並不查問?你豈不曾知道麼?你豈不曾聽見麼?永在的神耶和華,創造地極的主,並不疲乏,也不困倦；他的智慧無法測度.疲乏的,他賜能力；軟弱的,他加力量；就是少年人也要疲乏困倦,強壯的也必全然跌倒.但那等候耶和華的,必從新得力,他們必如鷹展翅上騰,他們奔跑卻不困倦,行走卻不疲乏.」(賽40:27-31)

如果不想跟時代太過脫節,這段日子的講壇事奉實在並不容易.我想除了直接呼喚上主出來,收拾這個爛攤之外,恐怕無論是誰說甚麼,都不會被所有人接受.我甚至會說,就算上主今日親自臨鑑,恐怕也無法得到所有人的信服.這並不是我憑空臆測.二千多年之前,上主成為肉身降生人世,卻是到處受到挑戰和招惹爭議 － 特別是在聖殿和會堂.無論如何,教會並不能抽身於人世.俄國作家杜斯妥也夫斯基在他的經典小說 “The Brothers Karamazov”,就曾經透過主角的想像,安排當時的大主教,也就是德高望重的宗教裁決庭總長,跟耶穌見面.其實他們最初的相遇,大主教是以「擾亂公眾秩序」捉了當時正在行神蹟的耶穌.然後就像當年尼哥底母那樣,有一晚大主教微服去監牢見耶穌.他們之間的談話,或者說,是大主教鏗鏘有力的獨白,充分反映了信仰和宗教之間的矛盾和張力.

大主教指出,耶穌渴望的,是得到他的子民發自內心對上主愛的回應；但他的信眾,其實不過是誠惶誠恐地想得到教會明確的訓示和祝福.耶穌指出透過真理得到的自由,對他們只是不必要的負累.因為自由背後是責任,而他們覺得生活已經太沈重.他們不需要更多的責任；自由對他們沒啥用處.

最後裁決總長勸耶穌離開,並且不要再回來,因為只要他一聲令下,那些曾經追隨他左右的群眾,將會毫不猶疑地將處決異端的炭火堆在他腳下.因此,成為肉身的基督,其實並不如我們想像那麼受歡迎.慶祝聖誕,或者盼望上主來臨,很可能只是不明白聖誕意義的人的浪漫想像.

有人說「黃藍是政見,黑白是良知」.這本來是很精闢到位的提醒.但事實是良知並不能抽離具體的社會,政治處境而運作.就算是同一個地域,不同的社會階層,文化修養和人生閱歷,都可以孕育出相當不同的道德價值.而道德價值的分歧,就解釋了良知其實並沒有一個絕對的準則.箴21:2「人所行的,在自己眼中都看為正；惟有耶和華衡量人心.」

不同的人基於他們的道德良知,可以作出迥異的道德判斷；而他們是完全可以無愧於他們的上主.因此我們看見信徒憑著他們的良知,在分歧的政見底下,做成不少教會的撕裂.原本只是政見上的分歧,最後竟上升到真理的層次,帶來信仰的裁決:你傾向某種政見,你就是離經叛道!「憑著他們的果子,就可以認出他們來.荊棘上豈能摘葡萄呢?蒺藜裡豈能摘無花果呢?」(太7:16) 曾認錯宮粉羊蹄甲是洋紫荊!

其實在舊約時代,信仰跟政治確實是密不可分的；大部分先知書的提醒和警告,都跟當時的社會,政治形勢有關.先知以賽亞,就是在以色列民族的歷史危機中,成為信仰和道德良知最重要的發言人.

有猶太傳統推測以賽亞為王族之後,因此應該有點「建制」背景.我們只知道他在主前740年耶路撒冷聖殿中見異象,蒙召成為先知,經歷南國猶大四代君王,在朝廷盡忠職守40年.當時正值國家危急存亡之秋,北國以色列在主前722年亡於亞述；而南國猶大,亦只能俯首稱臣,年年進貢,以苟延殘喘.

學者一般將《以賽亞書》分成兩部分.前半部(1-39章)大概是主前701年,耶和華將亞述大軍消滅於耶路撒冷城外之後不久完成(37:36-37).而因為當時的以賽亞似乎得到上主的「明撐」,所以也可以說是他先知事業的最高峰.而後半部(40-66章)則相信是在他半退隱後才繼續完成.

我們今天要思考的經文,恰好就是《以賽亞書》後半部的頭一卷.整章書卷可說是相當磅bok,充滿氣勢.在聖誕前後,我們經常會聽到韓德爾的神曲《彌賽亞》,開始的時候,就是用了《以賽亞書》四十章第一節:「你們的神說:你們要安慰,安慰我的百姓.」

這裡所描繪的上主雖然充滿大能,卻是極其溫柔:「主耶和華必像大能者臨到…(但他亦)像牧人牧養自己的羊群,用膀臂聚集羊羔抱在懷中,慢慢引導那乳養小羊的.」(40:10,11)在這裡,可以說是充滿了一般用來描繪女性的那種溫柔.

跟著,以賽亞幾乎是用了《創世記》甚至是《約伯記》的筆觸和想像,去描繪這位充滿溫柔慈愛的上主,形成非常強烈的對比:「誰曾用手心量諸水,用手虎口量蒼天,用升斗盛大地的塵土,用秤稱山嶺,用天平平岡陵呢?誰曾測度耶和華的心,或作他的謀士指教他呢?他與誰相議,誰教導他,誰將公平的路指示他,又將知識教訓他,將通達的道指教他呢?

「看哪,萬民都像水桶的一滴,又算如天平上的微塵；他舉起眾海島,好像極微之物.黎巴嫩的樹林不夠當柴燒,其中的走獸也不夠作燔祭.萬民在他面前好像虛無,被他看為不及虛無,乃為虛空.你們究竟將誰比神,用甚麼形象與神相比呢?…你們豈不知道麼?你們豈不曾聽見麼?從起初豈沒有人告訴你們麼?自從立地的根基,你們豈沒有明白麼?神坐在地球大圈之上,地上的居民好像蝗蟲.他鋪張穹蒼如幔子,展開諸天如可住的帳棚.」

在這一大堆的鋪陳之後,終於來到我們今日要思想的經文:「雅各阿,你為何說,我的道路向耶和華隱藏?以色列阿,你為何言,我的冤屈神並不查問?」簡單來說,他們是埋怨上主為何不理睬他們.這令我們想起,當日摩西帶領以色列人出埃及,他們亦經常向上主埋怨這樣那樣,甚至發晦氣懷念起昔日在埃及的種種好處.

其實也許雅各和以色列是對的:在上主面前「好像虛無,被他看為不及虛無,乃為虛空」的萬民,在那位大能者眼中又算得甚麼?雅各選擇一條怎樣的路,以色列受到甚麼冤屈,這位大能者又怎麼會「上心」?

我們豈不是說:每個人走哪條路都是自己揀的.早知今日,何必當初?我們都要為自己所選的路負責.怎能夠說贏了就自己袋袋平安,輸了就要上主包底,「唔算數,再來過」?

至於冤屈,有誰不是覺得自己受冤屈,被虧負?我們之所以出手,只不過是因為他們先動手；我們之所以做十五,完全因為對方先做初一.我們每一個行動,都有充分的解釋；我們走的每一步,都有充足的理據.這樣,我們又怎能說上主不看我們,不理我們的冤屈呢?

當年有兩兄弟,覺得自己在家產分配上被虧負.他們各持己見,認為道理在自己一邊.爭持不下,他們決定去找一位沒有利益衝突的第三方做公正.他們最後同意耶穌比較公正持平.但耶穌卻是斷然拒絕這個一般只有德高望重的人才會得到的邀請.他說「誰立我作你們之間斷事的官?」這些半斤是否足八兩的婆仔數,你怎好意思去找上主做裁決?

事實上這類糾紛,正如我在開頭的時候指出,恐怕就是上主親自出面,最後都會被認為不公正,有私心.畢竟人所做的,在自己眼中都看為正.

最近有位朋友在他的面書上貼了一段引述英國哲學家羅素說話: “The whole problem with the world is that fools and fanatics are always so certain of themselves, and the wiser people so full of doubts.”(Bertrand Russell)這當然是秀才遇著兵的無奈.羅素慨嘆,那些嘈喧巴閉,聲大夾惡的,都是被人擺佈的愚民；真正的智者,反而充滿疑惑.但在這個時候貼這些話,當然是意有所指.

我在他的貼文後面加了一句留言: “So true. But then it would not be easy for anyone to be certain about anything anymore.”我不錯是同意那句話.但他這樣說,還有誰敢再說些甚麼?因為只有蠢人才會對自己充滿信心,智者永遠懷疑自己的判斷；為了避免被誤會是蠢人,還有誰好意思說些中肯的話?今天這個社會,就是落在一個說也不是,不說也不是的局面.

無論如何,在這裡「雅各說:我的道路上主並不關心!以色列說:我的冤屈上主並不理睬!」我們反而感到有點理解,因為我們大概都領教過人自義的本性,並且完全盲目於自己的盲點.或許我們自己,在不同程度上也是這樣的人.大多數人想找的「公道」,其實主要是對自己有利的意見.

面對上主並不體察民情的投訴,以賽亞的回覆是:「你豈不曾知道麼?你豈不曾聽見麼?永在的神耶和華,創造地極的主,並不疲乏,也不困倦；他的智慧無法測度.」驟耳聽來,以賽亞好像有點牛頭不搭馬嘴.

基本上,以賽亞要指出的是,上主並不像他們想的那樣.「永在的神耶和華,創造地極的主,並不疲乏,也不困倦」；意思是沒有甚麼對祂來說是太複雜,太難或者太煩.「他的智慧無法測度」.

相比之下,我們就比較煩惱.譬如以色列和巴勒斯坦,誰才是耶路撒冷的真正主人?中,美各懷鬼胎,這是政治現實,但夾在中間,又可以怎樣自處?又或者面對政治衝突,應該據理力爭抑或息事寧人?在不同的時候,你可能會有不同的答案.但叫人頭痛的是,無論是甚麼問題,也無論你站在哪一邊,都總會有人說你不對,而他們又總是似乎言之成理!

我認識好幾位上世紀七十年代中期移民海外的香港留學生,理工科畢業.八十年代初內地開始改革開放,他們差不多是第一批入內地參與建設的海外菁英,得到的待遇和條件遠超過他們在外國所得.雖然曾經在外國生活,但作為改革開放的既得利益者,他們對內地的印象非常積極正面,最後甚至決定在內地置業,過他們優悠的退休生活.我又認識另外一些這個圈子以外,跟內地關係非常密切的朋友.他們對內地的認識亦相當深入,卻對現時的社會和時局有截然不同的觀感和看法.

那你問到底誰的認識更加真實?我認為這兩班朋友其實都同樣真實.但是你大概沒有可能說服任何一邊同意另外一邊的觀點.他們對內地的認識,都是第一身,並非道聽塗說.分別是他們身處於非常不同的圈子和環境.可以說如果人與人之間不可能達成共識,那是因為每個人都有其無可爭議的親身經歷.你不可能跟一個人的親身經歷爭論:「咪玩啦,我響裡面住jor 二十幾年!」

先知以賽亞指出,「永在的神耶和華,創造地極的主,並不疲乏,也不困倦；他的智慧無法測度.」在這裡,我們首先留意到的是:相對於我們的有限和短暫,上主是「永在」的.我們只不過目擊歷史一個很細小的片段,然後企圖從這片段作最大限度的引伸.

但就算我們讀埋希伯來大學歷史系教授Harari最新的全球暢銷書《人類簡史》(“Sapiens: A Brief History of Mankind”),再向前展望他的《未來簡史》(“Homo Deus: A Brief History of Tomorrow”),或者再精讀他寫給二十一世紀的《二十一課》,我們都只是好像管中窺豹,對人類的過去有個大概的印象,對將來有個大致上的展望和憧憬；並且這個印象,只是一個有個人立場的詮釋.畢竟沒有人真是能夠站在時間的起點,或者明日之後,去準確地描述人類文明的故事.

作者當然比我們有更淵博的學識,但跟我們一樣,他亦無法超越「不識廬山真面目,只緣身在此山中」的侷限.一個受過訓練的人,往往可以輕易舉起比自己體重更重的重量,但他卻永遠無法在原地舉起自己.就算湯恩比(Arnold Toynbee)再世,亦無法擺脫作為「當局者」的主觀困局.而相對於人的智慧,「永在的神,創造地極的主,他的智慧無法測度」.

面對這樣一位上主,不同的人有不同的態度.有些人認為,人實在太過渺小,不足以左右歷史的進程,成就任何真正有價值的事情.因此他們只能夠旁觀世態,靜觀上主的作為.他們成為歷史的旁觀者,見證人,卻始終站在生命的外圍,裹足不前.

另外一些人認為,有這無可匹敵的強大後盾,他們可以挑戰任何的歷史任務.他們憑著簡單的信心勇往直前,有時興高采烈,有時焦頭爛額.他們認為這樣才不枉此生.

但當風浪太過猛烈,黑暗勢力太過強大,他們也會感到灰心氣餒.畢竟「就是少年人也要疲乏困倦,強壯的也必全然跌倒.」在這些時候,他們多麼希望上主會說「疲乏的,可以休息；軟弱的,可以退場,讓我來收拾這個爛局」!

但上主卻有不同的對策.他說:「疲乏的,他賜能力；軟弱的,他加力量」.自己的問題,我們總不能袖手旁觀.如果每個太難的問題都有人代勞,我們就永遠學不會去處理自己的問題.上主應允加力,關鍵是我們要「等候」:「但那等候耶和華的,必從新得力,他們必如鷹展翅上騰,他們奔跑卻不困倦,行走卻不疲乏.」但這「等候」,卻似乎不是我們想的那麼簡單.

香港這個城市有不少特色,其中最突出的,大概是講求速度效率:千祈不要叫我等!

當年「地鐵」的廣告:「話咁快就到!」

11:「梗有一間係左近!」

香港電梯的開關速度,據聞是全世界最快!

擁有無線電話的人均比例,亦屬全球之冠!

你十分鐘內沒有覆WhatsApp,立刻會受到對方的批評.

曾經打電話找一位朋友,「飛線」到他的「秘書台」:「這裡是黃先生的秘書台,有甚麼可以幫到你?」沒有想到不是黃先生接聽,我稍為猶疑了兩秒,對方即說:「那麼我收線了,拜拜!」

當我們習慣了速度,效率,到一個地步它控制我們生活大小事項,要求絕對的尊重,那麼它可能成為一個問題.

事實上「等候」在聖經中亦是一個重要的主題.主耶穌說:「惟有那忍耐到底的,必然得救.」究竟忍耐或者「等待」背後包含了些甚麼?有何積極意義?抑或只是一種消極的無可奈何?你或許會想起Samuel Beckett的《等待果陀》.

「等」或「忍」是一種非常微妙的心理狀態.一方面你等,表明你仍未得到你所等的；但你之所以願意等,卻又預設了一個被你所等,並且最終要被你等到的(或許是一個人,或許是一個機會).在等的過程,你並沒有得到你所等的；但你所等待的,卻影響甚至支配著那等待中的你.如果你所等的在預期的時間未出現,那麼之後的一分一刻都會是煎熬:會不會出了甚麼意外?等錯了地方?記錯了時間?會不會是惡作劇?若等,再等多久?若不等,就在你不等的時候出現怎辦?就在這一切不肯定的張力中,你重新發現自己跟被等待那位的關係,是你在慣常的接觸中不會察覺的.

在這個意義上,或許人與上主之間的關係,最首要的就是:等.等的背後表明:人並不擁有上主.人永遠是看不清,摸不透,知不了,掌握不到上主.我們不能將上主放在我們的褲袋,成為我們隨時的方便.

希伯來書:「上主乃住在人不能靠近的光中.」信仰一個最基本的事實是:上帝是超乎我們能夠理解的.

我們即時的反應可能是:有沒有搞錯?不理解怎麼信?不理解當然可以信!你明不明白一架飛機的內部構造,和令它能夠起飛的aerodynamics? 至少我第一次坐飛機的時候,是連最皮毛的物理常識也沒有.你理不理解一部電腦的內部結構和運作原理?直到今天,我還沒有動機買本 “computer for dummies” 來看,但這篇講章,卻是我坐在電腦前面拼湊出來的.如果要我們明白了才做,理解了才用,那麼恐怕我們早上連門口也出不了.

我當然不是要在這裡解剖信心.我只不過是想回應剛才的問題:超乎我們可以理解的,我們當然可以信,並且是只能夠相信 － 正因為我們不理解!這可能也是Anselm說:“I believe in order that I can understand” 的意思.

「但那等候耶和華的,必從新得力,他們必如鷹展翅上騰,他們奔跑卻不困倦,行走卻不疲乏.」

鷹能夠高飛,當然跟它龐大的身軀,寬廣的翅膀有關；但如果你看過一些自然教育的圖書,你就會知道,鷹能夠在高空盤旋,其實並非靠它強而有力的雙翼,而是靠那比它更強大的氣流；去到某一個位,它只要張開雙翼,氣流就自然會將它承托.因此飛得很高的鷹,其實反而是很悠閒的；就像在空中翱翔的滑翔降傘,靠的並不是本身有甚麼強大的動力,而只不過是將自己放置在看不見卻明確感受到的氣流當中,被它所承托.這樣上騰高飛,又怎會疲乏困倦?

不錯,無論是鷹抑或滑翔降傘,當中確實也涉及一些刻意的「操作」:或者是鷹調節它身體不同的部分,或者是人控制綁在他身上的降傘.但這些「操作」卻不是他們能夠在空中飛翔的原因.他們只不過是將自己置身於比他強大的氣流,感受它的流動,跟它合作；靠著它的力量,毫不費力,輕描淡寫地在空中活動.

這對基督徒來說有很重要的引伸:我們能夠高飛,並且重新得力,靠的也不是本身的勇氣和毅力,堅忍和決心.這些都有它們的位置,卻不是最關鍵的因素.最關鍵的,是我們將自己全然投靠在那位「永在的神」,「創造地極的主」裡面,感應祂在世間的運行,在當中尋找屬於自己的位置.

然後我們體會到,「等候神」最基本的意思其實就是「讓神為神」(let God be God),拒絕去假設自己站在上主的位置,作出超過自己能夠理解和承擔的舉動和決定.而正如剛才所說,這並不代表我們無所作為.我們乃是隨著氣流將我們帶到的高度,在那裡做我們知道,我們明白的事情.

然後我們發現,「他們奔跑卻不困倦,行走卻不疲乏」,原來並不是修辭學上的poetic expression,而是基督徒生活實況的描寫!

「風隨著意思吹,你聽見風的響聲,卻不曉得從哪裡來,往哪裡去；凡從聖靈生的,也是如此」(約3:8)

「但那等候耶和華的,必從新得力,他們必如鷹展翅上騰,他們奔跑卻不困倦,行走卻不疲乏.」Ω



\newpage

\section{那就是記號了}
\label{sec:483}
\textbf{劉國偉先生 - 那就是記號了}
\newline
\newline
\hyperref[sec:482]{< < < PREV SERMON < < <}
~
\hyperref[sec:toc]{[返目錄]}
~
\hyperref[sec:484]{> > > NEXT SERMON > > >}


路加福音 2:8-14
\newline
\begin{longtable}{cl}
\hline
\hline
章節 & 經文 (和合本修訂版)\\
\hline
2:8 & \begin{tabularx}{0.7\textwidth}{X} 在伯利恆的野外有牧羊人，夜間值班看守羊群。 \end{tabularx} \\ \\ \relax
2:9 & \begin{tabularx}{0.7\textwidth}{X} 有主的一個使者站在他們旁邊，主的榮光四面照著他們，牧羊人就很懼怕。 \end{tabularx} \\ \\ \relax
2:10 & \begin{tabularx}{0.7\textwidth}{X} 那天使對他們說：「不要懼怕！看哪！因為我報給你們大喜的信息，是關乎萬民的： \end{tabularx} \\ \\ \relax
2:11 & \begin{tabularx}{0.7\textwidth}{X} 因今天在大衛的城裡，為你們生了救主，就是主基督。 \end{tabularx} \\ \\ \relax
2:12 & \begin{tabularx}{0.7\textwidth}{X} 你們要看見一個嬰孩，包著布，臥在馬槽裡，那就是給你們的記號。」 \end{tabularx} \\ \\ \relax
2:13 & \begin{tabularx}{0.7\textwidth}{X} 忽然，有一大隊天兵同那天使讚美神說： \end{tabularx} \\ \\ \relax
2:14 & \begin{tabularx}{0.7\textwidth}{X} 「在至高之處榮耀歸與神！在地上平安歸與他所喜悅的人！」 \end{tabularx} \\ \\
[1ex]
\hline
\hline
\end{longtable}
講題:那就是記號了 This is the Sign

經文:路加福音2章8-14節

講員:劉國偉先生

場合:香港中文大學崇基學院禮拜堂平安夜燭光崇拜

日期:2019年12月24日

引言

今天講道的題目「那就是記號了」源出於路加福音第2章第12節.記號的英文直譯sign含有跡象,徵兆的意思,有時亦會譯作神蹟.這節經文「你們要看見一個嬰孩,包著布,臥在馬槽裡,那就是記號了.」所選用的「記號」這個字的希臘原文是σημεὶον(simeìon),可以譯作記號,表徵,神蹟.

在我們的生活現場裡,其實充滿了各式各樣的記號,標誌和表徵,從日常生活,商業範疇,政治活動,以至宗教生活,都存在很多記號和標誌.

此外,很多人也熱衷於尋求徵兆,兆頭；占卜,觀星,相命,透過一些徵兆去求問前程.亦有科學家和未來學家,藉科學觀察和研究一些物質現象,去歸納和推斷,從而解構科學,天文,又或者為人類預測未來.

原來在我們的生活經驗當中,記號,標誌,表徵從來都是生活中的重要部份,甚至影響我們對人,事,物的聯想,預測和判斷.

耶穌基督的降生所帶給我們的,是什麼記號? 那記號盛載著什麼含義?在那個黑暗的晚上,並與我們所身處的紛亂世代,這記號又有什麼意義?

畜棚裡的馬槽

「約瑟也從加利利的拿撒勒城上猶太去,到了大衛的城,名叫伯利恆,因他本是大衛一族一家的人,要和他所聘之妻馬利亞一同報名上冊.那時馬利亞的身孕已經重了.他們在那裡的時候,馬利亞的產期到了,就生了頭胎的兒子,用布包起來,放在馬槽裡,因為客店裡沒有地方.」(路2:4-7)

作者路加醫生刻意強調約瑟是「大衛一族一家的人」這身份,提醒讀者耶穌出自王族家譜.耶穌的王族背景跟接著下來發生的事,形成了強烈的對比:這一位生於王族的嬰孩,竟然降生在畜棚之內,盛載這個細小軀體的竟然只是馬槽.

因為人口普查的緣故,百姓各歸各鄉,細小的伯利恆承載不起忽然湧來的回鄉旅客,客店再沒有房間,店主就把畜棚臨時開放,供給很有需要,或者部份窮苦的旅客住宿.

父母總會為剛剛誕生的嬰孩,找著足夠的保護,馬槽是餵飼牲口,放置飼料的容器.很明顯的,馬槽是當時可以找到的相對安好的空間,就成為一張嬰兒床,那就是救主以血肉之軀來到世間的第一個躺下歇息的地方.

天上的大君王,躺身於卑微的馬槽,看似是時勢及環境所逼,也許對約瑟而言是無可奈何的選擇,但掌管歷史的主,難道被世情時局所擺佈,身不由己?斷乎不是!從天上的使者告訴牧羊人的說話中,我們清楚知道,馬槽就是救主降生的記號的一個組成部份,是計劃的一部份.讓我們一起再誦讀路加福音第2章8至12節.

在伯利恆之野地裡有牧羊的人,夜間按著更次看守羊群.有主的使者站在他們旁邊,主的榮光四面照著他們；牧羊的人就甚懼怕.那天使對他們說:「不要懼怕!我報給你們大喜的信息,是關乎萬民的；因今天在大衛的城裡,為你們生了救主,就是主基督.你們要看見一個嬰孩,包著布,臥在馬槽裡,那就是記號了.」

這個記號出現在伯利恆

伯利恆位處耶路撒冷以南大約12公里的一個山上,海拔約790公尺,自古以來,直至今日,仍然是一個天然的牧羊場地.

主前八世紀的彌迦先知,曾經預言以色列一位君王來自伯利恆這個小城:

「伯利恆,以法他啊,你在猶大諸城中為小,將來必有一位從你那裡出來,在以色列中為我作掌權的；他的根源從亙古,從太初就有.」(彌迦書5: 2)

這個被天使稱為大衛的城的伯利恆,從主前十世紀建立的大衛王朝,延綿至耶穌降生,經過一千年的等待,伯利恆再次引起世人的注視.這個城雖然小,卻被猶太人理解為王族的城,彌迦先知更告訴我們,因為救主的緣故,這個小城連接著亙古,連接著太初,顛覆了我們對什麼是大,什麼是小的概念.

我們要怎樣衡量一個城市的地位? 是她的經濟生產總值? 她的政治地位? 她的人口及教育水平? 她的文化和歷史厚度? 她的軍警戰鬥力量?  先知彌迦告訴我們,伯利恆之歷史地位在於她出現了一位為上主掌權的掌權者,路加福音的作者則告訴我們,這位掌權者降生的記號出現在伯利恆.弟兄姊妹,也許其中一個簡單直接的角度去認識一個城市,在於她出現了什麼偉大的人,或者那裡的人具備什麼生命素質.

請你記住,上主揀選了伯利恆去盛載救主降生的記號.

天使把這個記號告訴牧羊人

伯利恆之野地,也就是牧羊人放羊的地方,是屬於半沙漠的乾燥地帶,這個超過千年的牧場,被視為伯利恆的標誌.這片曠野地,與城內相比,有明顯的分別,這是牧羊人工作,居住和生活的地方.

牧羊人不是城市人,在猶太人的生活習性中,牧羊工作被視作卑微和低下階層的工作,然而,他們除了對經濟生產有所貢獻,他們所牧養的羊是獻祭的主要材料.

先知彌迦接著寫道:「他必起來,倚靠耶和華的大能,並耶和華他神之名的威嚴,牧養他的羊群.」(彌迦書5:4上)先知奇妙的預言告訴我們,這位掌權者,同時是一位牧羊人,伯利恆正好成為一個十分相襯的背景.

牧羊人的身份地位當然不及君王,權貴,然而上帝不單沒有忽略他們,反倒揀選他們,成為天使宣佈耶穌誕生的好消息的對象.

試想,希律王,及可能一些上層社會的貴胄,豈不從智者口中,得悉「那生下來作猶太人之王」的訊息?他們的反應與牧羊人成了鮮明的對比:牧羊人急忙進城,要看上主成就的事,希律王卻選擇把伯利恆一帶,兩歲以下的男孩都殺掉! 用心去聆聽天上的聲音的人,不可能是那些自命不凡,以我為先,以為人定勝天,喜歡別人歌功頌德的人.希律王聽見「那生下來作猶太人之王」,就認定是對他的威脅,以極殘忍的暴力,奪去了無數嬰孩的性命,也傷透無名無姓無面孔的老百姓的心；牧羊人卻用心傾聽,他們甚至聽到天兵天使的美妙的讚美之聲.

牧羊人與希律的對比,顛覆了我們對歷史中無名無姓無面孔的老百姓的想像.上帝就是揀選了城外曠野的牧羊人,他們只不過是社會基層的小人物,並沒有什麼重要性的一群,但他們被告知救主降生的記號,更親身聽見天軍天使那美妙的聲音說:「在至高之處榮耀歸與神!在地上平安歸與祂所喜悅的人!」(路2:14)  弟兄姊妹,你認為誰人心中有真正的平安? 是希律王嗎? 還是得知救主降生記號的牧羊人?

那就是記號了

到底那個記號是怎樣的?這個記號有3個重要元素:一個嬰孩,包著布,臥在馬槽.

馬槽象徵卑微,按人來說,那不可能是君王誕生躺臥的地方,然而,這是上主的選擇.道成肉身的主,捨棄天上榮華,降生在人間一處卑微的地方,這是基督君王之道,這顯然與攀附權貴,為維護權力,地位和既得力益而不擇手段的價值觀格格不入,但基督的救恩的道路卻是以馬槽作為起點,然後以十字架為終點.弟兄姊妹,救恩之路就是自甘卑微和捨己.

嬰孩耶穌被布包著,從文學的佈局來說,作者路加醫生藉著兩種不同用途的布,作出福音書的首尾呼應.路加福音23章53至54節如此記載:「這人去見彼拉多,求耶穌的身體,就取下來用細麻布裹好,安放在石頭鑿成的墳墓裡.」

包著嬰孩身體的布,代表一個新生命的誕生,而裹屍布則象徵生命的終結.從生到死,都是上主的選擇,生的時候降生在馬棚,死的時候卻葬於財主的墳墓.弟兄姊妹,我們今天又要做什麼樣的選擇?我們要如何生?若果必要的話,我們會為什麼而死?

最後,這記號最重要的,當然是那個嬰孩.你可有想過,救主要來這個世間拯救世人,竟不是雷霆萬軍,帶著天使天軍駕臨,呼風喚雨,氣勢磅礡,與邪惡決一死戰,卻選擇以嬰孩的樣式,從小做起?

大概人類歷史可以說是一部戰爭歷史,除了是兵戎相見,子彈炸彈橫飛的真實戰爭,即使在相對太平的日子,我們很多時都眼見或者甚至身處人世間種種權力漩渦之中,尤有甚者,我們甚至會形容為人生戰場.

在權慾橫流的世界,有人弄權為樂,有人有權用盡,又有人以權謀私,將信任,盼望和仁愛磨掉；把弄權力的人,往往以扭曲的言詞,偽裝的良善和粗暴的手段行事,將真誠,正直,可靠,謙卑,溫柔等美德淘空.

救主以嬰孩的弱小身軀降臨人間,無權無勢,而且被有權殺人的希律王討命,誰知這個嬰孩,就是那位成全律法,捨身救贖,自己稱自己為道路,真理,生命的主.

我很喜歡今天晚上其中一段詩歌旁白,在結束講道的時候,我再一次讀給大家聽:

這是一個充滿奇妙神蹟的季節；上帝正在對世人說出祂的愛.讓我們聚集一起,凝視那滿有恩典記號的馬槽,看見天國君王選擇了慈愛的搖籃,而不是權力的寶座.耶穌,我們的救世主誕生了.

弟兄姊妹,救主降生的記號......出現在伯利恆,聽進野地的牧羊人的心裡,這個記號看似卑微,看來是如此脆弱,卻帶著上帝救贖世人的奧秘,是主耶穌向罪惡和死亡跨勝的人生旅程的起點.從伯利恆到各各他,從馬槽到十字架,主耶穌向我們顯現的那份深厚的愛,若果我們真心聆聽,虛心接受,必然看見主的榮耀,得享主所賜的平安,因為我們要成為一個上帝所喜悅的人.

願榮耀歸與三一真神!又確信平安歸與祂所喜悅的人!



\newpage

\section{寒夜數星星}
\label{sec:484}
\textbf{林豪恩先生 - 寒夜數星星}
\newline
\newline
\hyperref[sec:483]{< < < PREV SERMON < < <}
~
\hyperref[sec:toc]{[返目錄]}
~
\hyperref[sec:485]{> > > NEXT SERMON > > >}


講題:寒夜數星星 Recount Gracious Deeds in Distress

講員:林豪恩先生

場合:香港中文大學崇基學院禮拜堂主日崇拜

日期:2019年12月29日

一.引言

「寒夜數星星」這題目,令人想起關正傑的舊歌--〈星〉:

閉起雙眼睛　心中感覺清靜

再張開眼睛　怕觀望前程

夜冷風更清　這一片荒野地

沿途是歧路　我方向未能明

啊　不見朗月　導我迷途　只有星

啊　荒野路　伴我獨行　是流螢

縱步獨行　沿途寂靜　似只有呼吸聲

緩步前往　決意走崎嶇山徑

今天是2019年12月29日,一年將要完結.這一年跟過去很多年都不一樣,發生了很多事情,是很多人從沒想過會在這個城市出現的.回想一幕幕煙火連天的畫面,就算閉起雙眼睛,心中未必感覺清靜；再張開眼睛,的確怕觀望前程.「夜冷風更清」正是不少人的慨嘆；繁華的鬧市,彷彿變成了「荒野地」.「沿途是岐路,方向未能明」也是很多人的徬徨.黑夜裡,甚麼也看不見,好像甚麼也沒有了.歌詞說,不是甚麼也沒有,而是景物不同了.縱使不見朗月,但有星星導我迷途；縱使在荒野路,但有流螢相伴；縱使獨行寂靜,但有呼吸聲相隨.2019年,我們失去的確實是很多,甚至有人感到甚麼也沒有了.在這一年完結的時候,或者,我們可以用另一種角度回望,未必是甚麼也沒有了,可能是景物不同了.事實上,寒夜並不是空無,也不是黑暗,因為天上仍有或多或少,或暗或明的星星.讓我們與聖經作者們一齊嘗試在寒夜裡數星星,看看還有沒有值得我們珍惜的事物.

二.經文及內容

以賽亞書63:7-9

63:7我要照耶和華一切所賜給我們的,並他憑憐憫與豐盛的慈愛所賜給以色列家的大恩,述說他的慈愛和美德.

63:8他說:「他們誠然是我的百姓,未行虛假的子民.」這樣,他就作了他們的救主.

63:9他們在一切苦難當中,他也同受苦難,並且他面前的使者拯救他們(七十士譯本是「沒有使者或天使,只有主親自拯救他們」).他以慈愛和憐憫救贖他們,在古時的日子時常抱他們,背他們.

先知以賽亞在甚麼狀況之下呼籲猶大人述說上帝的慈愛和美德呢?上文63:1-6節論及以東.提到以東,當時的猶大人會想到甚麼呢?以東人如何對待猶大人呢?整卷俄巴底亞書(其實只有一章)就是談論以東,俄巴底亞先知如此說:

11當陌生人擄掠雅各的財物,當外邦人進入他的城門,為耶路撒冷抽籤分取財物的日子,你竟站在一旁,像與他們同夥.

12你兄弟遭難的日子,你不該瞪著眼看；猶大人被滅的日子,你不該幸災樂禍；他們遭難的日子,你不該說狂傲的話.

13我子民遭災的日子,你不該進他們的城門；他們遭災的日子,你不該瞪著眼看他們受苦；他們遭災的日子,你不該伸手搶他們的財物.

14他們遭難的日子,你不該站在岔路口剪除他們逃脫的人,你不該交出他們的倖存者.

以賽亞書63:1-6節耶和華上帝審判以東,側面肯定了猶大艱難和痛苦的遭遇.甚麼都沒有了,還述說甚麼恩典和慈愛呢?以賽亞先知提醒猶大人,處境的確改變了,失去的的確很多,然而,並不等於一無所有,還有三項是他們仍然擁有並且應該珍惜的,這些是他們寒夜中的星星,是黑夜中的盼望.

第一,以賽亞先知代表上帝向猶大人宣告,「他們誠然是我的百姓」,這是身份的肯定.人最迷亂的時候,往往是失去身份,不知自己是誰的時候.不知道自己是誰,不是沒有名字,而是所遭遇的令我們自我懷疑.例如被家人誤解或拒絕時,那種委屈感和孤獨感令人懷疑自己是否仍被視為家庭的一份子；被上司批評或者失去工作時,懷疑自己是否很無用；表遠了意見又採取了行動,對方仍然固我,甚至被視為敵人時,不禁令人懷疑自己是否仍是這個社會中的持份者.這些懷疑未必是空穴來風,而往往是經歷某些對待或遭遇而產生的.在過去幾個月的刺激與衝突之中,激盪起人們反思自己都底是誰,追求甚麼?不單單是動物性地生存,生命只是搵食養活自己,直至死亡.有人開始發現自己的良知,發現自己的正義感,對著某些事情原來是會憤怒的,有反應的.有人更深刻地問自己是誰?與誰走在一起?與誰認同?

當時猶大人的遭遇令他們懷疑自己,懷疑自己的信仰,懷疑自己的群體,因為他們彷彿成為了孤兒,無處為家.耶和華上帝在猶大人遭遇艱難,感到迷失的時候向他們重申,「他們誠然是我的百姓」.在過去幾個月中,或者我們也懷疑自己,對自己的身份感到迷茫,也感到這裡花果漂零,甚至感到無處為家的淒涼以及焦躁.我們歸屬哪邊呢?哪邊視我們為自己人呢?或者,我們都期望聽到掌管天地的上帝對我們說:「你們誠然是我的百姓」.屬於天地的主,就是天國的子民；天國的子民,按天國價值觀生活,也以天國價值觀審視現實.知道自己屬於誰以及站立於甚麼價值觀之上,是寒夜中的一顆小星星.

第二,以寒亞先知說,耶和華上帝「他們在一切苦難當中,他也同受苦難」,這是同在與同感.人們常常期望上帝點石成金,動一動指頭就令人化涕為笑,因此,與人同受苦難的上帝是難以難解,難以接受,也不受歡迎.然而,馬丁路德卻認為這正正是上帝愛的表現.就是在聖嬰誕生的事件而不是釘十字架事件,受苦的上帝已經顯現淋漓盡致.聖誕故事吸引之處,是充滿著神蹟奇,令人驚訝,例如童女竟然懷孕,看見天使往返飛翔,並且開口說話.但是,在路德看來,最大的神蹟是上帝竟然道成肉身,如嬰孩一般從婦女的肚腹中懷孕出生.最大的神蹟是全能的上帝竟然變成一個需要被照顧,需要長大的嬰兒.人們期待的上帝是超越的上帝,是一位坐在天上的上帝,但受苦的上帝卻是一位降下凡塵,與我們同在的上帝.這受苦的上帝對我們來說,有甚麼意義?受苦的上帝是表示在最艱難,最痛苦,最無助的地方,上帝都在那裡,創造天地的主宰也經歷了被人遺棄,死於十字架上這等痛苦,所以,在你受苦的時刻,不要以為只有你自己最孤單,上帝也同受苦難.此時此地,或許不少香港人感受到社會撕裂的痛,在一切困苦的人中間,上帝也在那裡同受苦難.

第三,這位與我們同苦難的上帝,卻同時又是拯救我們的上帝.以賽亞先知對當時的猶大人說:「並且他面前的使者拯救他們(七十士譯本是「沒有使者或天使,只有主親自拯救他們」)他以慈愛和憐憫救贖他們,在古時的日子時常抱他們,背他們」這事如何可能呢?讓我們看看希伯來書的作者對這事如何體會及如何告訴我們.

希伯來書2:10-18

2:10原來那為萬物所屬,為萬物所本的,為要領許多兒子進入榮耀,使救他們的元帥因受苦難而得以完全,本是合宜的.

2:11因那使人成聖的,和那些得以成聖的,都是出於一.為這緣故,他稱他們為弟兄也不以為恥,

2:12說:「我要將你的名傳給我的弟兄,在會眾中我要頌揚你.」

2:13他又說:「我要依賴他.」他又說:「看哪!我與上帝所給我的兒女都在這裡.」

2:14既然兒女同有血肉之軀,他也照樣親自成了血肉之軀,為能藉著死敗壞那掌管死權的,就是魔鬼,

2:15並要釋放那些一生因怕死而作奴隸的人.

2:16誠然,他並沒有幫助天使,而是幫助了亞伯拉罕的後裔.

2:17所以,他凡事應當與他的弟兄相同,為要在上帝的事上成為慈悲忠信的大祭司,為百姓的罪獻上贖罪祭.

2:18既然他自己被試探而受苦,他能幫助被試探的人.

潘霍華在被囚在處死之前說:「只有受苦的上帝才能夠幫助我們.」基督的受苦是真實的,祂雖然死而復活,但祂所受的苦改變了祂,我們可以在復活的基督身上見到,因為祂身上永留傷痕.神學家莫特曼指出,耶穌基督手上的釘痕仍存留在復活的身體上,永不消失,成為苦難的印記.因此,拯救與苦難並不對立,也不取代,拯救與苦難是共存的.莫特曼提醒基督徒,盼望若抽離不平的苦難,則必然不真實並失去解放的力量.只有以苦難為基礎的盼望才更具體實在,更具深刻的向度.基督的苦難並不只是祂個人的苦難,祂表達上帝親自參與我們的苦難,並將我們的痛苦攬在祂自己身上.

希伯來書說2:11「他稱他們為弟兄也不以為恥」,又說2:14「既然兒女同有血肉之軀,他也照樣親自成了血肉之軀」.那位「受苦的人子」成了我們中間的一位,以致於無數無名,被鞭打的,遭人遺棄的人都成了祂的兄弟姊妹.復活的基督就是被釘十架的基督,2:14-15「為能藉著死敗壞那掌管死權的,就是魔鬼,並要釋放那些一生因怕死而作奴隸的人」,先有受苦而後才有得勝.

莫特曼繼續說,「與上帝平安意味著與世界不平安,應許的未來如同尖銳的棒子,無情地刺進每一刻尚未實現的現在.」莫特曼的意思是,對上帝的盼望進入了現在,使人無法容忍現實,因為上帝國不是目前所見這個樣子,因此與上帝平安就是與世界不安.在迎向平安的盼望當中,過程卻有許多掙扎的不平安,也就是在不平安當中,繼續因盼望而得到平安.然而,只要上帝國與世界仍有不同,就會有兩者間的衝突,如果不去面對,可能是自我欺騙；如果勇於去面對,一定會帶來很多痛苦,然而在痛苦當中才能經歷更深層的平安.痛苦中的平安,也是另一顆寒夜裡的小星星.

以賽亞說63:9「並且他面前的使者拯救他們(七十士譯本是「沒有使者或天使,只有主親自拯救他們」)」希伯來書說:2:15「並要釋放那些一生因怕死而作奴隸的人」,不同時代都有寒夜,不同時代都有黑暗得只剩下星星的狀況,因此,上帝的拯救行動從沒停止,一直進行釋放奴隸的行動.以賽亞用出埃及經歷鼓勵猶大人:63:9「他以慈愛和憐憫救贖他們,在古時的日子時常抱他們,背他們」,馬太福音也以出埃及的意象來描述耶穌基督降世及使命.

馬太福音2:13-23

2:13他們走後,忽然主的使者在約瑟夢中向他顯現,說:「起來!帶著小孩子和他母親逃往埃及,住在那裡,等我的指示；因為希律要搜尋那小孩子來殺害他.」

2:14約瑟就起來,連夜帶著小孩子和他母親往埃及去,

2:15住在那裡,直到希律死了.這是要應驗主藉先知所說的話:「我從埃及召我的兒子出來.」

2:16希律見自己被博學之士愚弄,極其憤怒,差人將伯利恆城裡和四境所有的男孩,根據他向博學之士仔細查問到的時間,凡兩歲以內的,都殺盡了.

2:17這就應驗了耶利米先知所說的話:

2:18「在拉瑪聽見號咷大哭的聲音,是拉結哭她兒女；她不肯受安慰,因為他們都不在了.」

2:19希律死了以後,在埃及,忽然主的使者在約瑟夢中向他顯現,

2:20說:「起來,帶著小孩子和他母親回以色列地去!因為要殺害這小孩子的人已經死了.」

2:21約瑟就起來,帶著小孩子和他母親進入以色列地去.

2:22但是他因聽見亞基老繼承他父親希律作了猶太王,怕到那裡去；又在夢中得到主的指示,就往加利利境內去了.

2:23他們到了一座城,名叫拿撒勒,就住在那裡.這是要應驗先知所說的話:「他將稱為拿撒勒人.」

馬太福音記載「希律要搜尋那小孩子來殺害destroy(摧毀)他」,與希伯來書2:14「既然兒女同有血肉之軀,他也照樣親自成了血肉之軀,為能藉著死敗壞(摧毀)那掌管死權的,就是魔鬼」這句遙遙相應,希律摧毀嬰孩耶穌不成,耶穌反而「藉著死敗壞(摧毀)那掌管死權的,就是魔鬼」.

關於嬰孩耶穌這短短的敘述,馬太福音多次指出是應驗先知的預告:

2:15這是要應驗主藉先知所說的話:「我從埃及召我的兒子出來.」

2:17這就應驗了耶利米先知所說的話:2:18「在拉瑪聽見號咷大哭的聲音,是拉結哭她兒女；她不肯受安慰,因為他們都不在了.」

2:23這是要應驗先知所說的話:「他將稱為拿撒勒人.」

聖子道成肉身,人們以為嬰孩耶穌裂天而降,馬太福音從另一個角度指出並不是突如其來,而是一脈相承.承繼甚麼呢?就是先知的預言,即是上帝參與世間,施行拯救,釋放奴隸的傳統.這釋放傳統的代表事件是以色列人出埃及事件,貫穿著被擄回歸,接下去對彌賽亞的盼望.

馬太福音指出聖誕故事是出埃及的延續及更新,也是被擄回歸的延伸及更新,是上主打敗邪惡力量,使受壓迫者得釋放的故事,是人能夠自由地與天地一齊敬拜上主的故事.(出埃及記5:1)「摩西和亞倫去對法老說:「耶和華－以色列的上帝這樣說:『放我的百姓走,好讓他們在曠野向我守節.』」

『放我的百姓走,好讓他們在曠野向我守節.』這是出埃及故事的主題,是被擄回歸故事的主題,是聖誕故事的主題,會不會也是我們當下的城市故事的主題呢?故事仍未結束,劇情仍在發展,但結局卻已經寫下,就是萬國萬民最終會被釋放,在新天新地與上主相會.想到這裡,我們可以參與詩人的讚美.

詩篇148:1-14

詩人首先呼籲天上的讚美耶和華:

148:1哈利路亞!你們要從天上讚美耶和華,在高處讚美他!

148:2他的眾使者啊,要讚美他!他的諸軍啊,都要讚美他!

148:3太陽月亮啊,要讚美他!放光的星宿啊,都要讚美他!

148:4天上的天和天上的水啊,你們都要讚美他!

148:5願這些都讚美耶和華的名!因他一吩咐就都造成.

148:6他將這些設定,直到永永遠遠；他訂了律例,不能廢去.

接著,詩人呼籲地上的也讚美耶和華:

148:7你們哪,都當讚美耶和華:地上一切所有的,大魚和深洋,

148:8火和冰雹,雪和霧氣,成就他命令的狂風,

148:9大山和小山,結果子的樹木和一切香柏樹,

148:10野獸和一切牲畜,昆蟲和飛鳥,

148:11世上的君王和萬民,領袖和世上所有的審判官,

148:12少年和少女,老人和孩童,

148:13願這些都讚美耶和華的名!因為獨有他的名被尊崇,他的榮耀在天地之上.

最後,詩人提出選民更加要讚美耶和華:

148:14他高舉自己百姓的角,使他的聖民以色列人,就是與他相近的百姓得榮耀.哈利路亞!(或譯「他一切聖民以色列人, 就是與他相近的百姓,都讚美他」)

三．總結:

引言提及關正傑「星」這首歌,仍有最後再段,讓我們作為此篇講章的結束.

踏過荊棘　苦中找到安靜

踏過荒郊　我雙腳是泥濘

滿天星光　我不怕風正勁

滿心是期望　過黑暗是黎明

啊　星也燦爛　伴我夜行　給我影

啊　星光引路　風之語　輕輕聽

帶著熱情　我要找理想　理想是和平

尋夢而去　那怕走崎嶇險徑

明日誰步過　這星也帶領



\newpage

\section{在基督裡的多重福氣}
\label{sec:485}
\textbf{蘇成溢牧師 - 在基督裡的多重福氣}
\newline
\newline
\hyperref[sec:484]{< < < PREV SERMON < < <}
~
\hyperref[sec:toc]{[返目錄]}
~
\hyperref[sec:486]{> > > NEXT SERMON > > >}


講題:在基督裡的多重福氣 The Spiritual Blessings in Christ

講員:蘇成溢牧師

經文:女人弗所書1:1-14

場合:香港中文大學崇基學院禮拜堂主日崇拜

日期:2020年1月5日

引言

願恩惠平安從我們的父上帝和主耶穌基督歸與你們.這一句問安語,來自以弗所書1章2節,是保羅書信常用的.對生活在公元一世紀的基督徒來說,他們大多是一無所有的人；但對他們來說,最重要的並非是財富與權位,乃是恩惠與平安.我相信,經歷過2019年各種衝擊與挑戰的香港基督徒,包括今天在這裡的你與我,也會認同恩惠與平安的重要性.首先,我認為恩惠能使我們知道,無論順逆,只要全心仰賴上帝,總有出路.我辦公的地方是太子道西,很接近港鐵太子站,距離旺角警署.過去數個月,不知有多少個下午,黃昏或晚上,在太子道西與彌敦道交界旺角警署旁,發生示威,集會,堵路,激烈衝突,封站等.然而,就算因為晚上有會議,最終我也能平安回家,這是因為旺角太子區的其他街道,例如花墟一帶的洗衣街,界限街等,仍有路可走.為此我堅信一件事,就是那些為上帝做事的人,縱使遇上困境,仍有出路.另一方面,平安教我們要實踐與人和睦,為所居之城求平安.當然,平安並非簡單地指人與人之間和平相處,平安無事,乃是以平等,公義,真理,彼此守望,互相關心,和諧共處等來實現.也即是說,得著平安之福的基督徒,同時領受平安使命,為全地得享合乎上帝國標準的完滿平安而努力.

在進入講解今天的經課前,我想提有關此段經文的兩點特別之處:一是一氣呵成,二為在基督裡的重要.第一方面:以弗所書1: 3-14整段內容的文法,句式和用語十分特別.3-14節是連成一句,一氣呵成的一整段.即從3節起至14節結束是一句話.這表明在保羅的心中,實在有很豐富的一段話要講.當中的內容雖涉及不同的範疇,卻環環緊扣,全是關乎福音信仰,當中包含了上帝的奧妙旨意,生命如何成為聖潔,宇宙萬物的終極走向等.這段話可以視為保羅闡釋上帝救贖計劃的總論,對何謂基督福音的概要介紹.第二方面:按中文繙譯,十二節經文裡共有六次提到 “在/藉基督裡” ,分別出現在3, 5, 9, 10, 11, 12節.可見在保羅的思想中,在基督裡何等重要.

在以下的經課講解,我會以在基督裡定為最重要的思想核心 (弗1: 3-14) ,並分為兩個部份與大家一同思考:1) 在基督裡究竟如何  2) 在基督裡指甚麼

第一:在基督裡究竟如何

i. 首先,我們一起看每一節所談到的內容3節 得各樣屬靈福氣 (如下) ；4節 揀選我們,成為聖潔無瑕,滿有愛心；5-7 節  得兒子名份 (白白得恩典:蒙救贖,赦免過犯) ；8-9節 以諸般智慧使我們知道上主旨意的奧秘；10節  使一切同歸於一 (完全聯合,總結一起) ；11-12節 得著基業；13-14節 因福音而相信,受聖靈印記/基業憑據.總合來說共有九樣 (揀選,聖潔,愛心,兒子名份,上主旨意,同歸於一,基業,信心,聖靈印記) . 從以上所數算的屬靈福氣,其實包含個人的和整體兩方面的.個人的屬靈福氣有揀選,聖潔,愛心,兒子名份,基業,聖靈印記等；整體的屬靈福氣有恩典,救贖,福音,同歸於一等.今天世人所追求的福,與保羅所列舉的,截然不同.對很多人來說,收入高,居住環境好,生活有保障而無憂,就是福氣；亦有人認為,毋須看重金錢和物質享受,但要自由自在:做喜歡做的事,去喜歡去的地方,吃喜歡吃 的菜式,買喜歡買的東西等,才算是有福氣.各位弟兄姊妹,你又如何衡量自己是否個有福氣的人.我看過一個關於在一個小鎮內一名基督徒護士的故事.一名醫科生在實習時,認識當地小鎮的一名護士.二十年來只有她一人默默工作,服務病人.對這位醫科生來說,這名護士所充滿的忍耐,剛毅和歡愉,令了驚訝.她為了急症,即使要半夜趕來,亦毫無倦意,但薪酬卻過於微薄.有一深夜,經過了一整天艱辛工作,這名醫科生要為了替她抱不平,向她說:” 你為甚麼不叫他們付多點給你? 上帝也知道你是值得的!” 她回答說:” 如果上帝也知道我是值得的,我夫復何求.”原來,她不是為人工作,更非為酬報,乃是為上帝工作.

ii. 另一方面,保羅在長長的經文分享中,他清晰及不嫌其煩地指出:從個人到群體甚至整個宇宙,若要得福氣 ,基督就是那唯一的中心.衪為那獨一的途徑,基礎,開始與終結.生命和萬有的完滿和成全,全在乎基督.換言之,基督是上帝計劃和救世福音的關鍵者；反過來說,人若沒有得著基督,恩典,揀選,救贖,成聖,兒子名份和福氣,均不可能.十天前大家以不同方式慶祝聖誕節,然而,除了聖誕日崇拜或聖誕子夜彌撒外,我發現,在香港及世界很多地方的聖誕節,是沒有耶穌的.聖誕樹,聖誕燈飾,聖誕禮物,聖誕食品,聖誕老人,聖誕餐會,所謂的聖誕歌曲,那一樣是與這位使人得屬靈福氣的耶穌有關的呢?求主幫助我們,只專注於稱為基督的耶穌身上,因為祂是唯一的主角

第二:在基督裡指甚麼

i. 首先,我認為在基督裡是指生命與基督結連,即我的生命中有基督,也確知自己在基督生命之中 (約15: 1-8) .就如自己來自甚麼家庭,生於何地,活在何處,生命自然流露出相融和的特質,習慣和生活方式.來自四川與來自廣東的人口味當然不同.我出生於澳門及在當地受中小學教育,之後來香港,在崇基讀神學,成為牧者,至今逾四十年.因此,我十分了解澳門及香港兩地的社會及文化特質,如何塑造澳門人及香港人的特質和生活方式.各位基督徒,請問,你會怎樣描述,自己因生命中有基督,有些甚麼特質?

ii. 另一方面,一個在基督裡的人,不單擁有基督所賜的身份,更得著豐盛恩典.也即是說,因為在基督裡,人被賦予了一個獨特身份,此身份能使人完全享受主為他所預備一切的豐盛恩典.一次我到牙買加出席會議,會議場地是在一間渡假式酒店.那酒店原來是那種住,食,消閒,活動全包的,入住者唯一需要做的,就是配帶酒店提供的住客手帶以確認身份.當你帶著手帶進入了酒店,你所具備的身份,便可完全地享受那裡所提供的,預備的,精心策劃的一切.感謝主,祂已為我們預備了身份和豐盛恩典,來吧,讓我們一起因信成為在基督裡的人.

透過以下兩個提問,我邀請你與我為今天的經課作反省

第一,今天作為基督徒的你與我,會如何描述自己與基督結連?在我的生命中,究竟有那些基督的特質?

第二,你經歷過在基督裡的豐盛恩典嗎? 你怎樣數算所受的恩呢? 聖詩【數算主恩】的副歌歌詞說:Count your blessings name them one by one, Count your blessings see what God has done, Count your blessings name them one by one, Count your many blessings see what God has done. 主的恩典樣樣都要數,主的恩典都要記清楚；主的恩典,樣樣都要數,必能叫你驚訝立時樂歡呼.



\newpage

\section{神啊求你不要耽延}
\label{sec:486}
\textbf{鄧瑞強博士 - 神啊,求你不要耽延}
\newline
\newline
\hyperref[sec:485]{< < < PREV SERMON < < <}
~
\hyperref[sec:toc]{[返目錄]}
~
\hyperref[sec:487]{> > > NEXT SERMON > > >}


講題:神啊,求你不要耽延 Do Not Delay, O My God

講員:鄧瑞強博士

經文:詩篇40

場合:香港中文大學崇基學院禮拜堂主日崇拜

日期:2020年1月19日

各位弟兄姊妹,早安.

有隊愛爾蘭樂隊,叫做U2,1983年出版了一張CD專輯,這專輯名為「War」(戰爭).這是一張控訴戰爭的CD.這CD的第一首歌,叫「Sunday Bloody Sunday」,講述1972年英軍在北愛地區開槍射殺手無寸鐵的遊行人士的事件.這事件後來就稱為「Bloody Sunday」.這首歌的第一句是:「I can’t believe the news today」(我不敢相信今天的新聞).不敢相信,是因為一切看來都匪夷所思.看起來不真實的景象,看起來不可能發生的事件,卻是真的,就發生在身邊.這CD的最後一首歌,歌名叫「40」.歌詞很簡單,是這樣的:「我曾耐性等候上主；他垂聽我的呼求.他從禍坑裡,從淤泥中,把我拉上來.他使我的腳立在磐石上,使我腳步穩當.許多人必看見而懼怕.我要唱一首新歌,一首新歌.但還要等多久呢?How long, how long, how long?」全首歌,迴盪著這問號,How long,還要等多久呢?還要等多久,戰爭才會過去,人能愉快地唱一首新歌呢?

U2以「40」這首歌來結束「War」這CD專輯,這首歌留下一大問號,還要等多久?這首壓尾歌的歌名叫「40」,就是指「詩篇40篇」.這是今天的詩篇經課,也是我今天講道的經文.

經文:詩篇40篇

詩40:1　我曾耐性等候耶和華；他垂聽我的呼求.

詩40:2　他從禍坑裡,從淤泥中,把我拉上來,使我的腳立在磐石上,使我腳步穩當.

詩40:3　他使我口唱新歌,就是讚美我們神的話.許多人必看見而懼怕,並要倚靠耶和華.

詩40:4　那倚靠耶和華,不理會狂傲和偏向虛假之輩的,這人便為有福!

詩40:5　耶和華─我的神啊,你所行的奇事,並你向我們所懷的意念甚多,不能向你陳明；若要陳明,其事不可勝數.

詩40:6　祭物和禮物,你不喜悅；你已經開通我的耳朵.燔祭和贖罪祭非你所要.

詩40:7　那時我說:看哪,我來了!我的事在經卷上已經記載了.

詩40:8　我的神啊,我樂意照你的旨意行；你的律法在我心裡.

詩40:9　我在大會中宣傳公義的佳音；我必不止住我的嘴唇.耶和華啊,這是你所知道的.

詩40:10　我未曾把你的公義藏在心裡；我已陳明你的信實和你的救恩；我在大會中未曾隱瞞你的慈愛和誠實.

詩40:11　耶和華啊,求你不要向我止住你的慈悲!願你的慈愛和誠實常常保佑我!

詩40:12　因有無數的禍患圍困我,我的罪孽追上了我,使我不能昂首；這罪孽比我的頭髮還多,我就心寒膽戰.

詩40:13　耶和華啊,求你開恩搭救我!耶和華啊,求你速速幫助我!

詩40:14　願那些尋找我,要滅我命的,一同抱愧蒙羞!願那些喜悅我受害的,退後受辱!

詩40:15　願那些對我說阿哈,阿哈的,因羞愧而敗亡!

詩40:16　願一切尋求你的,因你高興歡喜!願那些喜愛你救恩的,常說:當尊耶和華為大!

詩40:17　但我是困苦窮乏的,主仍顧念我；你是幫助我的,搭救我的.神啊,求你不要耽延!

我將這首詩分成三個段落.

(1)感恩:1-6節

詩40:1　我曾耐性等候耶和華；他垂聽我的呼求.

詩人曾經遇到艱難,曾經祈禱,也曾耐心等候,而最終神聽了他的祈禱.這是感恩詩的典型開場白.「他垂聽我的呼求」,這句經文的原文的字面意思,是神彎了腰,側了身,伸耳朵過來,來聽我的呼求.神不遙遠,神俯就卑微,很親近.

神做了什麼呢?我感恩什麼呢?

詩40:2　他從禍坑裡,從淤泥中,把我拉上來,使我的腳立在磐石上,使我腳步穩當.

「禍坑」,「淤泥」,是困住人的,是陷阱,是毀滅生命的,是置人於死地的,人在其中動彈不得,不能自救,只能等死.詩人曾經身處一莫大的困局中.

神將他拉上來,將他放在磐石上.磐石是穩固的.「磐石」這字眼,有懸崖的含意.通常,堡壘就築在這些懸崖上,很安全,很穩當.

淤泥與磐石,深坑與高臺,真是天淵之別.神的作為,讓人離開地底的淤泥,而立於生命的高崗,這就是救贖的含意.為這徹底的改變,詩人感恩不盡.

詩40:3　他使我口唱新歌,就是讚美我們神的話.許多人必看見而懼怕,並要倚靠耶和華.

神讓詩人能口唱新歌,這是讚美神的詩歌,這也是勝利之歌.香港人近來都在淤泥中,舉步維艱,但總有一天,香港人能一起唱一首歌,一首新歌.「許多人必看見而懼怕」,所謂「懼怕」,即fear,即「敬畏」.人看見神不可思議的救贖,生起敬畏的心.我想起「出埃及」的故事,在紅海邊,神不可思議地將海水分開,做奴隸的以色列人,走過紅海,成為自由人.在紅海那邊,他們剛剛見過神奇妙的作為,心生敬畏,油然地,便唱起救贖之歌.

人在必死的苦境中,經歷到極大的解救,對生命必然有所悟.第4至5節道出這種生命之悟.

詩40:4　那倚靠耶和華,不理會狂傲和偏向虛假之輩的,這人便為有福!

詩40:5　耶和華─我的神啊,你所行的奇事,並你向我們所懷的意念甚多,不能向你陳明；若要陳明,其事不可勝數.

經歷過救贖,詩人領悟到:「那倚靠耶和華,不理會狂傲和偏向虛假之輩的,這人便為有福!」有福的人生,原來就是這樣的.

農曆新年快到了,我們又會聽到很多祝福的說話,如:恭喜發財,心想事成…,在詩人眼中,這些東西無福可言.(原來,我們在新年裡,彼此祝賀一些無福的東西.)什麼是福?詩人說:「那倚靠耶和華,不理會狂傲和偏向虛假之輩的,這人便為有福!」

「倚靠耶和華」這經文中的「倚靠」,是名詞,有「避難所」這含意.所謂「倚靠耶和華」,就是以耶和華為生命終極的蔭庇,以神為自己靈魂的家園.什麼事也向神祈禱,不一定要倚靠神.有人會為自己能找到一份不用做而又人工好的工作而祈禱,這不叫做「倚靠神」.能安於一份工作而努力在工作中祝福人,這才叫做「倚靠神」.「倚靠神」,是以神為家,以神的旨意為生命的依歸.

「不理會狂傲和偏向虛假之輩」,「狂傲」的希伯來文是rahab(讀音:rahav),這個字解「狂傲」,但同時也是神話中「海怪」的名字.海怪,很多時是指政治裡深不可測的毀滅力量.若這裡的「狂傲」真是指「海怪」,則「虛假之輩」就是指「假神」,「人造的偶像」,「權力的偶像」.

什麼是「有福」?那倚靠耶和華,不理會狂傲和偏向虛假之輩的,這人便為有福!這些人以神為家,拒絕虛假,拒絕謊言,拒絕偶像,拒絕毀滅他人的權力,也不向邪惡的力量低頭.這些人才是有福的.

詩40:5　耶和華─我的神啊,你所行的奇事,並你向我們所懷的意念甚多,不能向你陳明(《新譯本》:「沒有人可以和你相比」)；若要陳明,其事不可勝數.

這是心靈的一種境界,體會到生命的無比豐盛.從神而來的,數之不盡的奇事,充滿心靈.眼目所見,無一不是神救贖的奇妙.路上所遇,無一不是出於神的意念.詩人在此,對大劫後的生命奇蹟充滿驚歎,猶如在森林大火後,見到小草裂地而出,蚱蜢在灰燼中蹦蹦跳.數不盡的,就是這些從神而來的,創發生命的奇事.

一個感恩的人,看到的,是一個充滿恩典的世界.

(2)天命:6-12節

隨著感恩而來的,是體會到自身的天命.人懂得感恩,便體會到天心.感應到天心,便體會到天心在我們心內的呼喚,因而體會到上蒼對我們的要求,從而促使我們努力以有限之人生去實現天父無限的仁慈.

6-12節講:在感恩中,天心下達人間.

詩40:6　祭物和禮物,你不喜悅；你已經開通我的耳朵.燔祭和贖罪祭非你所要.

這是一簡單的,三文治結構的句子.頭和尾,都是講述神不喜悅的東西,神不喜悅無靈魂的奉獻.經文中間是餡,是核心信息,神喜悅的,是開通的耳朵.這是神所開通的,希望我們甘心聆聽,聆聽神的吩咐,聆聽天心,領受天命.

詩40:7　那時我說:看哪,我來了!我的事在經卷上已經記載了.

當我們感激上蒼時,我們體會到,冥冥之中,上天為我們譜寫了角色與使命.來到人間,好像就是為了背起放在這裡的十字架.詩人看到這天命,他說:「好,讓我來吧」.當天地崩壞之時,他準備好,以雙手撐起天地.

我想起2017年那齣電影《Darkest Hour》(《黑暗對峙》)裡的邱吉爾.德國粹納逐步侵佔歐洲的不同國家,英國的戰時內閣力勸邱吉爾簽和約,因為他們認為打不過納粹.邱吉爾卻拍抬而起,堅持作戰,因這不單關乎英國的安危,也關乎整個歐洲的命運,更關乎人類文明的前程.在最黑暗的時刻,他舉起燭光.

詩40:8　我的神啊,我樂意照你的旨意行；你的律法在我心裡.

在這個只顧自己的社會裡,他看不到自己,他只看到神.在這個再不講良心的時代,他不計較得失,他心中有神的法度.在這個目無法紀,藐視上帝的世界,他堅信上帝,堅持按神的律法而行.這個人,放下了生死,超越於得失,心無罣礙,只有上帝.這種生命,純然無雜,像清澈的一池春水,反映出明亮的上天.

這種人生,流露出來的,不是自己,而是上天的心.上天的心是什麼?是正義,仁愛,真理.

詩40:9　我在大會中宣傳公義的佳音；我必不止住我的嘴唇.耶和華啊,這是你所知道的.

詩40:10　我未曾把你的公義藏在心裡(按:隱藏在心中而不宣示出來)；我已陳明你的信實和你的救恩；我在大會中未曾隱瞞你的慈愛和誠實(即:真理).

所謂「不止住」,「未曾隱藏」,「未曾隱瞞」,即是說,「天心」如如地在這生命中流露出來,無法止住,不會隱藏.這真是「天人合一」的境界.他向眾人言說的,沒有私意,而只是上天的公義,仁愛,真理.在不仁不義,充滿謊話的世界,堅持仁義真理,要有「雖千萬人吾往矣」的大勇.

這樣的一個踐行天命的人,卻不自視偉大,而是時刻仰賴上天的恩典,深知自己的卑污.

詩40:11　耶和華啊,求你不要向我止住你的慈悲!願你的慈愛和誠實常常保佑我!

詩40:12　因有無數的禍患圍困我,我的罪孽追上了我,使我不能昂首；這罪孽比我的頭髮還多,我就心寒膽戰.

需要的,是神的恩典,是神的保佑,毫不以為自己有什麼了不起.自身的真面目,是深不可測的罪污.數不盡的,是自己的罪孽.

前一大段,論到感恩,結尾的第5節,講述神數不盡的恩情.

今段,論到領受天命,結尾的第12節,講述的,是人自申訴不盡的罪污.

這是基督教對人的深刻理解.就算是一個流露天心的聖人,仍自覺自己只是一個罪人.人不可能全然神聖,因為唯有神是「聖」的.但在中國文化裡,一個聖人,大抵真的會認為自己是神聖的.基督教始終對人性深處的邪惡,有深刻的理解.

(3)祈求:13-17節

詩40: 13-17與詩70篇基本上是一樣的,這段落是可以從這首詩中分離出來的.這解釋了我為何以13-17節作為一段,而不是按《和合本》那樣將12-17節作為一段.

詩40: 13-17是祈求,這是有點怪的,因為在感恩詩裡,很少祈求的.能夠感恩,是因為危難己過.若要祈求,就是因為正面對危難.到底危難過了沒有呢?

過了,便感恩；未過,便祈求.但這詩開頭是感恩,結尾是祈求,將兩者合而為一,這只能是因為危難過了,卻又未完全消失.這是一種already but not yet(已然未然)的狀態.危機過了,故有得勝的把握.但危機仍未完全消失,故仍要奮鬥.這有點似:公司已聘用你,你可以感恩；但試用期未過,故仍要努力.

基督徒的人生,就在這probation期間內.基督的救贖已完成,這是值得感恩的.但我們的天命仍有待完成,這是我們需時刻警醒祈禱的.這首詩「感恩」的部分,我們看過了,「天命」的部分,我們也看過了.在這probation期內,我們當祈求什麼呢?

詩40:13　耶和華啊,求你開恩搭救我!耶和華啊,求你速速幫助我!

詩40:14　願那些尋找我,要滅我命的,一同抱愧蒙羞!願那些喜悅我受害的,退後受辱!

詩40:15　願那些對我說阿哈,阿哈的,因羞愧而敗亡!

詩40:16　願一切尋求你的,因你高興歡喜!願那些喜愛你救恩的,常說:當尊耶和華為大!

看來,詩人的性命危在旦夕,有人要害他.在這緊迫時刻,他將生命寄托於神.

在這紛亂社會裡,在這人吃人的世界裡,在這只能仰賴上帝的人生裡,詩人最後的領悟是:

詩40:17　但我是困苦窮乏的,主仍顧念我；你是幫助我的,搭救我的.神啊,求你不要耽延!

他體會到自身的卑微,在茫茫苦海中飄浮,只有一樣東西支撐著他,叫他不放棄,就是那顧念人,搭救人的神.

我們都體會過,天父的恩典.我們都領略到,生命的天職.但在生活的重壓下,很多人絕望,有人甚至放棄生命,但願在重壓下,我們仍相信上帝,向神祈禱.只要你願意祈禱,就能保持希望.我們相信,神聽人祈禱.在這裡,在香港,有無數人在祈禱.神啊,求你不要耽延!

但願:榮耀歸於聖父,聖子,聖靈.阿們.

(附記:本人在崇基禮拜堂的講章已結集成《荷花池畔春草綠》講道集,試讀本可參閱:\url{http://cc.cuhk.edu.hk/dev/images/content/fund-raising-campaigns/university-sermons-by-dr-sk-tang/CCC-SK-Tang-Sermon-Text-Sample-2.pdf}.較早期的崇基講章,也已結集成《敢於跟隨主》一書,有興趣的朋友可到各大基督教書局購買.)



\newpage

\section{不可分黨一心相合}
\label{sec:487}
\textbf{鄭漢文博士 - 不可分黨,一心相合.}
\newline
\newline
\hyperref[sec:486]{< < < PREV SERMON < < <}
~
\hyperref[sec:toc]{[返目錄]}
~
\hyperref[sec:488]{> > > NEXT SERMON > > >}


哥林多前書 1:10-18
\newline
\begin{longtable}{cl}
\hline
\hline
章節 & 經文 (和合本修訂版)\\
\hline
1:10 & \begin{tabularx}{0.7\textwidth}{X} 弟兄們，我藉我們主耶穌基督的名勸你們說話要一致。你們中間不可分裂，只要一心一意彼此團結。 \end{tabularx} \\ \\ \relax
1:11 & \begin{tabularx}{0.7\textwidth}{X} 我的弟兄們，革來氏家裡的人曾對我提起你們，說你們中間有紛爭。 \end{tabularx} \\ \\ \relax
1:12 & \begin{tabularx}{0.7\textwidth}{X} 我的意思是，你們各人說：「我是屬保羅的」；「我是屬亞波羅的」；「我是屬磯法的」；「我是屬基督的。」 \end{tabularx} \\ \\ \relax
1:13 & \begin{tabularx}{0.7\textwidth}{X} 基督是分裂的嗎？保羅為你們釘了十字架嗎？你們是奉保羅的名受了洗嗎？ \end{tabularx} \\ \\ \relax
1:14 & \begin{tabularx}{0.7\textwidth}{X} 我感謝神，除了基利司布和該猶以外，我沒有給你們中的任何一個人施洗， \end{tabularx} \\ \\ \relax
1:15 & \begin{tabularx}{0.7\textwidth}{X} 免得有人說你們是奉我的名受洗的。 \end{tabularx} \\ \\ \relax
1:16 & \begin{tabularx}{0.7\textwidth}{X} 我曾為司提法那家施過洗；此外我已記不清有沒有給別人施過洗。 \end{tabularx} \\ \\ \relax
1:17 & \begin{tabularx}{0.7\textwidth}{X} 因為基督差遣我不是為施洗，而是為傳福音；並不是用智慧的言論，免得基督的十字架落了空。 \end{tabularx} \\ \\ \relax
1:18 & \begin{tabularx}{0.7\textwidth}{X} 因為十字架的道理，在那滅亡的人是愚拙，在我們得救的人卻是神的大能。 \end{tabularx} \\ \\
[1ex]
\hline
\hline
\end{longtable}
講題:不可分黨,一心相合. Unity in Diversity

講員:鄭漢文博士

經文:哥林多前書1:10-18

場合:香港中文大學崇基學院禮拜堂主日崇拜

日期:2020年1月26日

在教會中有紛爭

大家好,先和大家拜年,新年快樂.祝大家身心靈都平安.

今天是我今年第一次在此講道,今年我被校牧安排了每單數月份最後的星期日講道.下次是三月底,五月底,七月底…….

新一年,我給自己講道一點新意,新挑戰.

以前在四段經課,我多數先選福音書來講道.

今年,我希望先選書信來講,而且我會有意在「邊緣位置」多講一些.甚麼叫「邊緣位置」(marginal space)?你出現聽o下就會知?

我和許立中弟兄曾經合作出版過一本書叫《邊緣人間》,你可以找來看以此作為線索…….

今日的經課係取自哥林多前書1章10節至18節,讓我再讀一次:

10弟兄們,我藉我們主耶穌基督的名勸你們都說一樣的話.你們中間也不可分黨,只要一心一意,彼此相合.

11因為革來氏家裡的人曾對我提起弟兄們來,說你們中間有紛爭.

12我的意思就是你們各人說:「我是屬保羅的」；「我是屬亞波羅的」；「我是屬磯法的」；「我是屬基督的」.

13基督是分開的嗎?保羅為你們釘了十字架嗎?你們是奉保羅的名受了洗嗎?

14我感謝神,除了基利司布並該猶以外,我沒有給你們一個人施洗,

15免得有人說,你們是奉我的名受洗.

16我也給司提法那家施過洗,此外給別人施洗沒有,我卻記不清.

17基督差遣我,原不是為施洗,乃是為傳福音,並不用智慧的言語,免得基督的十字架落了空.

18因為十字架的道理,在那滅亡的人為愚拙；在我們得救的人,卻為神的大能.

這段經文,聖經以「在教會中有紛爭」作為主題.我們把這段經文分為三段來思考:

首先是10至12節是保羅對分爭的描述；

跟著13至16節是保羅對分爭的質問；

最後17至18節是保羅對分爭的解答.

2.　不可分黨

先看10至12節,保羅如何對分爭描述?

10弟兄們,我藉我們主耶穌基督的名勸你們都說一樣的話.你們中間也不可分黨,只要一心一意,彼此相合.

11因為革來氏家裡的人曾對我提起弟兄們來,說你們中間有紛爭.

12我的意思就是你們各人說:「我是屬保羅的」；「我是屬亞波羅的」；「我是屬磯法的」；「我是屬基督的」.

A Church Divided Over Leaders

10 I appeal to you, brothers and sisters, in the name of our Lord Jesus Christ, that all of you agree with one another in what you say and that there be no divisions among you, but that you be perfectly united in mind and thought.

11 My brothers and sisters, some from Chloe’s household have informed me that there are quarrels among you.

12 What I mean is this: One of you says, “I follow Paul”; another, “I follow Apollos”; another, “I follow Cephas”; still another, “I follow Christ.”

這些經文好像在呈現,有份創立哥林多教會的保羅透露,有人向他打小報告,對哥林多教會「篤灰」的是革來氏家裡的人(Chloe’s household),教會前線出現群而分黨,出現了至少四派的教友,他們以源出不同的屬靈領袖而分群分黨,造成教會紛爭.用香港人的話翻譯,就係在教會裡教友「各自跟了不同的大佬」.

哥林多教會的分黨分派的情景,在今日的香港教會是否似曾相識呢?

所以,有些信徒就會認為:「聖經其實不贊成政治中的政黨(political parties).」

這個「所以」,是一種對紛爭厭惡的信仰反應,情有可原.但在新約聖經的時代,還未有「政黨政治」(Party Politics)這種政治體制.第10節的「不可分黨」跟政黨無關.

我大學時副修政治學,「政黨政治」(Party Politics)正正是有效地代表社會裡不同的政治取態的公民,而全世界發展出來的都偏向「兩黨或多黨輪替管治」的民主體制.英國有保守黨和工黨,美國有共和黨和民主黨,台灣有國民黨和民進黨.至於中國,…….大家明白,不用宣之於口了.

所以,或贊成或反對「政黨政治」(Party Politics)其實都沒有聖經基礎.但是,有沒有信仰基礎?

教會向來教導要合一,所以,向來反對差異.我不同意,想平反「不違反合一的差異」(diversity congruent with unity).

我初入大學階段時,已經信了耶穌五,六年.

當時不明白為何教會史上出現不同的大分裂(schism),分出天主教(Catholic),東正教(Orthodox),基督新教(Protestant)?為何會有不同的宗派(denominations),包括我所屬的循道宗(Methodist)如何從聖公宗(Anglican)分出來,還有浸信會,崇真會,宣道會,中華基督教會?

我當時哲學主修的同班同學中有一位我非常敬重的信仰前輩,大學時已經在教會作牧養工作.我便向他請教.

他這樣解釋給我知:「因為神愛世人,讓有不同差異需要的信徒可以適應,可以按宗派傳統而有信群可以歸屬(diversity accommodates particularity so that each of us can have one’s own faith community to belong to).」

三十年來,我一直跟著這位信仰前輩的理解,去解讀教會中出現的分殊差異(differentiated diversity).

你們在座的,其實都認識這位信仰前輩,他就是在這裡多年的榮譽助理校牧鄧瑞強博士.

所以,我會公開承認:「我是跟鄧瑞強博士的.」(I follow Dr. Tang.)

你呢?

回顧你的信仰成長,你對那位信仰前輩或屬靈領袖較為靠近呢?

例如,在崇基禮拜堂,你會否在以下選項中選擇其一:

「我是屬鄧瑞強博士的.」(I follow Dr. Tang.)

「我是屬Simon校牧的.」(I follow Simon.)

「我是屬Walter的.」(I follow Walter.)

「我是屬Roger的.」(I follow Roger.)

我在這裡,不是鼓勵大家分群分黨,在這個教會搞紛爭.

我認為一個信徒的信仰成長,對一些信仰前輩或屬靈領袖較為靠近,是人之常情.

我認為一群信徒的信仰成長,對不同的信仰前輩或屬靈領袖各有不同程度的較為靠近,是人之常情.

所以,我們值得反思是,「怎樣的差異才是良好的,健全的(what make diversity a good thing)?」

我提出三點讓大家考慮,可以叫做「分殊差異的倫理」(Ethics of Diversity / Difference):

差異是倒映個別人的特點或特殊需要(diversity is attentive to particulars).

差異呈現群體的豐富(diversity contributes to richness of the community).

差異最終是為了更好的合一(diversity converges onto a better unity).

請容許我「埋身地」演繹以上三點,就是借用大家在崇基禮拜堂的經驗來說明.

自我四十三年信主後的教會經驗,我現在在崇基禮拜堂正在經歷著以上差異的三重精彩!

我自己在這裡作為其中一位聚會的教友,我也會同時說:

「我是屬鄧瑞強博士的.」(I follow Dr. Tang.)

「我是屬Simon校牧的.」(I follow Simon.)

「我是屬Walter的.」(I follow Walter.)

「我是屬Roger的.」(I follow Roger.)(我自我分裂,還是自戀?)

其他三位的講道我是不停地追蹤的,我自己講道時也會參詳他們的精彩.

我相信包括我在內,四位的講道各自各精彩.再加不時邀請來的外來講者,「五者講道,缺一不可.」

我作為牧養團隊其中一份子,從未如此見過這般美好的團隊(one of the best teams I ever join).

我其實非常捨不得現在的光景,但也明白好景不常,能多敘一刻就多一刻.(every moment is so precious)因此,大家好好珍惜當前的美好.

所以,快些報讀主日學,多享受不同型態的信仰學習.(植入式廣告!)

所以,事實上出現了的差異,不一定構成合一的障礙.

差異可以是合一的豐富.例如,大家在這裡可以同時說:

「我是跟得到Simon校牧的」；

「我是跟得到Walter的」；

「我是跟得到鄧瑞強博士的」；

「我是跟得到Roger的」.

在這裡,你發現了差異的美麗.

基督是分開的嗎?

可是,當時哥林多教會沒有這樣差異的美麗,反而構成合一的障礙.

第13至16節我們讀到了,保羅發了義怒,責罵哥林多教會他們:

13基督是分開的嗎?保羅為你們釘了十字架嗎?你們是奉保羅的名受了洗嗎?

14我感謝神,除了基利司布並該猶以外,我沒有給你們一個人施洗,

15免得有人說,你們是奉我的名受洗.

16我也給司提法那家施過洗,此外給別人施洗沒有,我卻記不清.

13 Is Christ divided? Was Paul crucified for you? Were you baptized in the name of Paul?

14 I thank God that I did not baptize any of you except Crispus and Gaius,

15 so no one can say that you were baptized in my name.

16 (Yes, I also baptized the household of Stephanas; beyond that, I don’t remember if I baptized anyone else.)

保羅連珠發了三個質問:

13基督是分開的嗎?保羅為你們釘了十字架嗎?你們是奉保羅的名受了洗嗎?

因此,是甚麼因素令得差異變成不良而有惡的,是因為違背了以下三個原則:

教會作為基督的身體,不可分開.

洗禮作為教會聖禮之一,不可構成合一的障礙.

信仰成長中出現的屬靈領袖,只是過渡意義,沒有終極意義.

按基督新教的聖禮觀,只有兩個聖禮:一是洗禮,一生只會做一次；二是聖餐禮,一生重複地做--只要你曾經經歷過前者.(According to the Protestant tradition of Christianity, there are only two sacrements: baptism and holy communion).

洗禮作為兩大聖禮之一,是發生在某時某地,在某人即某牧者手下進行的.好像出生一樣的有生日,例如:出生在某醫院.

受洗是歸入基督,成為基督身體的一部分.一經受洗就可以領受聖餐,聖餐禮才是基督身體的永恆重現.

那裡洗禮?誰施行洗禮?這些差異,不可構成合一的障礙.Baptism is just a milestone, never a final home.

所以,聖餐叫做holy communion,是在主裡合一的存活.(being one in Christ)

正如洗禮是進入聖餐合一的入口(baptism being entrance to holy communion),信仰成長中出現的屬靈領袖,只是過渡意義,終極意義在指向為我們死在十字架上的基督.

所以,若在我領的小組中有組員領洗,我只是做了渡人邁向基督.我的差異,造就了合一.

免得終極意義的基督的十字架落空

今日經課的最後兩節,就是指出任何教會活動的終極意義.有給這兩節一個主題,是:「基督是神的能力和智慧」.

讓我再讀出第17至18節:

17基督差遣我,原不是為施洗,乃是為傳福音,並不用智慧的言語,免得基督的十字架落了空.

18因為十字架的道理,在那滅亡的人為愚拙；在我們得救的人,卻為神的大能.

17 For Christ did not send me to baptize, but to preach the gospel-not with wisdom and eloquence, lest the cross of Christ be emptied of its power.

Christ Crucified Is God’s Power and Wisdom

18 For the message of the cross is foolishness to those who are perishing, but to us who are being saved it is the power of God.

我在本科讀中國哲學的時候,老師曾介紹宋明儒學的「理一分殊」(Unity in Diversity)的觀點.簡單說,上天有個「理」,但人間能夠體現到這個「同一的理」卻是「分殊」的個體性,因為個別的人有「心」可以「心同此理」,但活出這個「理」來可以是「分殊差異」的人生.

不同的施洗,甚至不同的傳道,也只是「分殊差異」的個體性.若我們的心能同此理,不停住在這些過渡意義的「分殊差異」,而是指向終極意義的合一,就是基督的十字架,而這道理是表明「基督是神的能力和智慧」.

5.　 一心相合,在於基督的大愛

快完結了.再返今日的經課哥林多前書1章第10節,保羅勸說的指向,讓我再讀一次:

10弟兄們,我藉我們主耶穌基督的名勸你們都說一樣的話.你們中間也不可分黨,只要一心一意,彼此相合.

一心相合是指向基督十字架所呈現的愛.一心相合,需要基督的信心,愛心與望心.

哥林多前書這卷書最多人背誦的要算第13章.

哥林多前書第13章被稱為「愛篇」(love chapter),常在婚禮中被用上.愛篇中所呈現的看似二人相愛在境界上可企及的高峰.

若明白哥林多前書其實是保羅得知哥林多教會出現甚多的紛爭,所以又勸又責弟兄姊妹當中的分裂以及諸種惡行,因而寫出第13章「勸愛篇」(love-persuasion chapter).

結束時,請大家打開哥林多前書第13章,讓我們一同站立誦讀這段愛篇,作為回應的禱文:

「1我若能說萬人的方言,並天使的話語,卻沒有愛,我就成了鳴的鑼,響的鈸一般.

2我若有先知講道之能,也明白各樣的奧祕,各樣的知識,而且有全備的信,叫我能夠移山,卻沒有愛,我就算不得甚麼.

3我若將所有的賙濟窮人,又捨己身叫人焚燒,卻沒有愛,仍然與我無益.

4愛是恆久忍耐,又有恩慈；愛是不嫉妒；愛是不自誇,不張狂,

5不做害羞的事,不求自己的益處,不輕易發怒,不計算人的惡,

6不喜歡不義,只喜歡真理；

7凡事包容,凡事相信,凡事盼望,凡事忍耐.

8愛是永不止息.先知講道之能終必歸於無有；說方言之能終必停止；知識也終必歸於無有.

9我們現在所知道的有限,先知所講的也有限,10等那完全的來到,這有限的必歸於無有了.

11我作孩子的時候,話語像孩子,心思像孩子,意念像孩子,既成了人,就把孩子的事丟棄了.

12我們如今彷彿對著鏡子觀看,模糊不清,到那時就要面對面了.我如今所知道的有限,到那時就全知道,如同主知道我一樣.

13如今常存的有信,有望,有愛這三樣,其中最大的是愛.」

願恩惠常與你們同在!

但願榮耀都歸給聖父,聖子,聖靈,永在的三一神.阿們.



\newpage

\section{聖所門口的要理問答}
\label{sec:488}
\textbf{光華牧師 - 聖所門口的要理問答}
\newline
\newline
\hyperref[sec:487]{< < < PREV SERMON < < <}
~
\hyperref[sec:toc]{[返目錄]}
~
\hyperref[sec:489]{> > > NEXT SERMON > > >}


詩篇 15:1-5
\newline
\begin{longtable}{cl}
\hline
\hline
章節 & 經文 (和合本修訂版)\\
\hline
15:1 & \begin{tabularx}{0.7\textwidth}{X} 耶和華啊，誰能寄居你的帳幕？誰能居住你的聖山？ \end{tabularx} \\ \\ \relax
15:2 & \begin{tabularx}{0.7\textwidth}{X} 就是行為正直、做事公義、心裡說實話的人。 \end{tabularx} \\ \\ \relax
15:3 & \begin{tabularx}{0.7\textwidth}{X} 他不以舌頭讒害人，不惡待朋友，也不隨夥毀謗鄰舍。 \end{tabularx} \\ \\ \relax
15:4 & \begin{tabularx}{0.7\textwidth}{X} 他眼中藐視匪類，卻尊重那敬畏耶和華的人。他發了誓，雖然自己吃虧也不更改。 \end{tabularx} \\ \\ \relax
15:5 & \begin{tabularx}{0.7\textwidth}{X} 他不放債取利，不受賄賂以害無辜。做這些事的人必永不動搖。 \end{tabularx} \\ \\
[1ex]
\hline
\hline
\end{longtable}
講題:聖所門口的要理問答 Entrance catechism

經文:詩篇15:1-5

講員:朱光華牧師

場合:香港中文大學崇基學院禮拜堂主日崇拜

日期:2020年2月2日

每個主日早上我們都如常來到崇基禮拜堂崇拜,

大家踏入這禮堂時,迎接我們的是一張張笑臉,

表示對我們的歡迎,期待,樂意,

我們習慣了每週來參加崇拜,

但是這一兩個星期,當我們進入禮拜堂門口,

我們要問自己幾個問題:

我是否戴了口罩?

我是否用酒精搓手液清潔雙手?

我是否去過武漢或者內地?

我有否發燒徵狀?

問過這些問題,通過這些問題,才可以進入禮拜堂.

所以,在這新年時節,我祝願大家身體健康,百毒不侵,藥到病除,

想起兩位中國歷史人物的名字,

一位姓霍名「去病」,一位姓辛名「棄疾」,

我也祝願大家「去病」和「棄疾」.

不過,今天早上跟大家講道的題目是「聖所門口的要理問答」,

不是問剛才那幾個問題,

我們要問自己更基本的問題:

我來到禮拜堂,來到聖所,

我是否有資格來朝見上帝?

我是否有足夠條件來敬拜上帝?

我是否達到進入聖所的標準?

我有什麼因素,足以令我被排抗於聖所門外?

請大家打開聖經,詩篇15篇,

詩篇15篇的作者大衛,

在第一節問了一個這麼樣的問題,

「1耶和華啊,誰能寄居你的帳幕?誰能居住你的聖山?」

Lord, who may dwell in your sacred tent? Who may live on your holy mountain?

這首詩擘頭第一句問了這樣的問題,

照理大衛這類有質素的信徒,那麼屬靈的信徒,

他大可不必問這些問題,

他不必懷疑他會不能通過這些問題.

這首詩不單問了一個問題,

在接下來第2節開始,更提供了這問題的答案,

-------------------------------------------------------------------------

| 1 節,問題   |  1 耶和華啊,誰能寄居你的帳幕?誰能居住你的聖山?      |

|              |                                                        |

|-----------------------------------------------------------------------|

| 2-5 節,答案 |  2 就是行為正直,做事公義,心裡說實話的人.            |

|              |                                                        |

|              |  3 他不以舌頭讒害人,不惡待朋友,也不隨夥毀謗鄰舍.    |

|              |                                                        |

|              |  4 他眼中藐視匪類,卻尊重那敬畏耶和華的人.            |

|              |                                                        |

|              |    他發了誓,雖然自己吃虧也不更改.                    |

|              |                                                        |

|              |  5 他不放債取利,不受賄賂以害無辜.                    |

|              |                                                        |

|-----------------------------------------------------------------------|

| 5 節,確據   |    做這些事的人必永不動搖.                            |

|              |                                                        |

-------------------------------------------------------------------------

做這些事的人必永不動搖.

原來這種問題答案的格式,不是詩篇15篇獨有,

原來這種問題答案的格式,不是詩篇15篇獨有,

在另一首詩篇,詩篇廿四篇,也有相同的格式,

同樣是出自大衛手筆的詩篇,

詩廿四:3,「3誰能登耶和華的山?誰能站在他的聖所?」

詩廿四:4,提供了問題的答案,

「4就是手潔心清,意念不向虛妄,起誓不懷詭詐的人.」

我們讀者不禁要問:這些問題,是否只是禮儀上的要求?是否只是循例詢問?

似乎不是!

根據歷代志下廿三:19,「19耶何耶大又設立守衛把守耶和華殿的各門,無論因何事而不潔淨的人,都不准進去.」

這節經文告訴我們兩件事:

第一,原來當時祭司耶何耶大,

安排了保安守著聖殿門口,

這些保安要把關,要篩查進聖殿的人,

若有需要,他們會禁止人入聖殿,

第二,他們用什麼尺度來量度和衡量誰可以進入?

若果不潔淨的人,都不准進入.

這裡所謂不潔淨,不是指他們沒有沐浴,沒有洗面洗手之類,

而是指他們的道德行為是否妥當,是否符合律法要求.

所以,有聖經學者告訴我們,

這類格式的詩歌,就好像是「聖所門口的要理問答」,

就好像是《韋斯敏斯德要理問答》或者是《海德堡要理問答》,

詩十五:1,「1耶和華啊,誰能寄居你的帳幕?誰能居住你的聖山?」

上半句和下半句的意思是重覆的,

誰能有聖所的居留權,

誰能在神的家定居下來,settle下來,

詩十五:2,先提出兩個一般性generic的條件,

「2就是行為正直,做事公義……的人.」

行為正直,就是那些行事無可指摘的人,無瑕疵,

這是借用利未記,摩西律法教導以色列百姓獻祭,

所用的祭物是要完美的,沒有殘疾,沒有defect的祭物,

進入聖所的第一個條件,

必須行為正直純全,無可指摘,無瑕疵,blameless的人.

經文提到第二個條件是做事公義的人,

今天這個世代,我們很重視爭取公義,

我們常自覺我們常是站在公義的一邊,我們要仗義執言,

為了公義我們甘願以雞蛋的姿態,對撼高牆,

這勇氣值得嘉許,但問題是:什麼是公義?

我覺得對就對,我覺得有道理就是公義,

我覺得不義就是不義,

講到底,誰是義與不義的最終決定者?

我!

我的判斷,我的理性,我的愛惡.

我以前有一段日子在澳洲生活,

澳洲天冷的季節都很寒冷,

所以洗澡時,總是凍到「雪雪」聲,

於是很多人的家中的浴室,都安裝一盞熱燈,

一盞熱燈,有四個火數250火的燈膽,

四個加起來就是1000火,

所以洗澡時,開著熱燈,雖然剝了衣服,

一點都不覺得凍,而是熱「笠笠」,

從此,寒冬下洗澡,成為了一大樂事.

不過這個裝置,有一個缺點,

就是在1000火的強光下,

浴室裡面的污漬,頭髮,微塵,全部現形,

但這些污漬微塵,平日普通的燈膽下,70火,100火,發現不到,看不見,

但是在強光下,無所遁形.

弟兄姊妹,平日我們覺得自己幾好,

是一個不錯的基督徒,

我們固定參加教會聚會,參予事奉,恆常奉獻,

週末回到教會,跟弟兄姊妹問候幾句,彼此關心,

三兩個小時後,我們回家,繼續我們如常的生活,

我們覺得我們是很不錯的基督徒,自我感覺良好,

但在與上帝面對面,在神的道就如手術室的手術燈的強光照射下,

不同角度的照射,

我們平日留意唔到的污點,雜質,罪惡,

都現形,都暴露出來.

我想強調的:

不是我們自己決定什麼善,什麼是惡,

什麼義,什麼是不義,

量度義和善的標準,不是人,而是上帝.

我們很喜歡的聖經經文,

彌迦書六:8,「8世人哪,耶和華已指示你何為善.他向你所要的是甚麼呢?只要你行公義,好憐憫,存謙卑的心與你的神同行.」

很多時,我們將注意力放在其中的兩小句「行公義,好憐憫」,

但這番話,是由誰發出?是誰給的指示?

「耶和華已指示你何為善.他向你所要的是甚麼呢?」

He has told you, O man, what is good; and what does the Lord require of you.

發出指示的,是永恆主,

確立善與義的標準,同樣是這位永恆上帝,

而且,「行公義,好憐憫」後面,還加上「存謙卑的心與你的神同行」,

與上帝同行,與上帝齊上齊落,與上帝的步伐一致,與上帝的方向相同.

公義,美善,正直,無瑕疵,

不單是我們外在的行為,

不單是見得到visible的表現,

更是我們見不到invisible的行事動機,見不到的念頭,見不到的想法,見不到的目的,

有一次,在美國發生了一件這樣的事,

有一個信主很多年的姊妹,是一位基督徒社工,

經朋友介紹,認識一個妓女,多來往,多交談,慢慢熟了,

很有負擔向這妓女傳福音,談人生,講信仰,

有一次,教會舉辦佈道會,這位姊妹邀請她參加佈道會,

希望她能信主,結果這妓女被這社工的熱誠感動,

一生人第一次踏足教會,參加佈道會聽耶穌,

佈道會完jor 之後,這社工姊妹問這妓女,

「點啊,聽完道,覺得點?對我們教會牧師講道,覺得他講得好唔好?」

「你們的聚會幾好啊,教會的人幾nice,你們的牧師,我識佢,佢成日來幫襯我的.」

公義和美善,不單是見得到的行為表現,

更是見不到的生命質素,品格誠信,

就好像是建築物,

不單是建築物的見得到的上層建築Superstructure,

更加是建築物的見不到的底層地基Substructure,

「看不見的」比「看得見的」重要

第2節,先從整體,用一個generic的角度,講行為做事,要完美之外,

第2節的下半句開始,用不同的角度,勾劃這個義人的樣式,

第2節的下半和第3節,講人的口,人的舌頭,

「2就是行為正直,做事公義,心裡說實話的人.3他不以舌頭讒害人,不惡待朋友,也不隨夥毀謗鄰舍.4他眼中藐視匪類,卻尊重那敬畏耶和華的人.他發了誓,雖然自己吃虧也不更改.5他不放債取利,不受賄賂以害無辜.」

人的舌頭有正反兩面的指標,

正面,要說實話,說出心中的真實話,君子坦蕩蕩,快樂老實人,

在大學做研究,用正當的手法取得數據,不可用虛假數據,不用誤導的數據,更加不抄襲,

反面,不以舌頭讒害人,

舌頭是可以有很大殺傷力,

羅三:10,「10沒有義人,連一個也沒有」

接下來的經文13-14節,勾劃不義的人的形像,

用了四句說話,表達人的舌頭的殺傷性,

羅三:13 - 14,「13他們的喉嚨是敞開的墳墓；他們的舌頭玩弄詭詐.他們的嘴唇裡有毒蛇的毒液,14滿口是咒罵苦毒.」

等同墳墓,詭詐,毒液,苦毒,

詩十五:3節的下半,「3他不以舌頭讒害人,不惡待朋友,也不隨夥毀謗鄰舍.」

不傷害不中傷朋友,不詆譭鄰舍,

不講別人壞話,閒話,反話,

這裡特別提出不隨夥,

意思是我們平日一個人時,不會講人閒話壞話,

不過當我們一班人一齊,茶餘飯後,開始談及其他人時,我們傾向會講人壞話閒話,

特別是我們當中,很多是身居要職,我們的說話很有份量,我們的說話,就容易對人產生傷害,

講得好,對人很大幫助,

講得不好,對人的傷害同樣很大,

特別我是牧師,有人會問我對某人某事的意見,

我就要小心,一句無心快語,甚至一個慨歎,一個眼神,就可能殺人於無形,

例如,未經fact-check,我轉發出去的信息,對有關的人同樣很大傷害.

而且毀謗,不一定要說話,

有時沉默,一個簡單的面部表情,也可以毀謗,

當一班人在批評另一個人,

他們的批評不合乎事實的,

我身在其中,我聽了,不表示意見,不出聲反對或者澄清,就已經對當事人帶來傷害.

詩十五:4,「4他眼中藐視匪類,卻尊重那敬畏耶和華的人.他發了誓,雖然自己吃虧也不更改.」

詩人甚至勾劃這個可以進入聖所的人的眼神,

他眼中藐視匪類,

「尊重那敬畏耶和華的人」,

不是尊重耶和華,而是尊重那敬畏耶和華的人,

心中響往,重視,佩服這些敬畏耶和華的人,

珍惜value敬畏耶和華的人,

「他發了誓,雖然自己吃虧也不更改」,

應承jor ,即使吃虧也不更改,

即使帶來不方便也不改,

這是指我們的信用,

人家對我們的信任,是經年累月的建立積存起來,

但一次失信,可以一夜之間完全摧毀.

詩十五:5,「5他不放債取利,不受賄賂以害無辜.」

不會剝削窮人,不會待薄他們,

我們在家中如何對待工人姐姐?

當他們是自己人,一家人,抑或使到佢地盡?

星期日,做完屋企ye ,然後才俾佢出街放假?

兒女對工人姐姐放肆無禮貌,呼喝藐視,我們縱容兒女的行為?

工人姐姐,做jor 4年半,隨便找個理由辭退他們,

如果唔係,做滿5 年就要給他們長期服務金?

我們沒有時間每一句都詳盡的討論,

原來這首詩共列出十項條件,

5項正面的,5項反面的.

--------------------------------------------------------------

|         | 正面條件             |        | 反面條件         |

|         |                      |        |                  |

| 2       | 行為正直             | 3      | 不以舌頭讒害人   |

|         |                      |        |                  |

|         | 做事公義             |        | 不惡待朋友       |

|         |                      |        |                  |

|         | 心裡說實話           |        | 不隨夥毀謗鄰舍   |

|         |                      |        |                  |

| 4       | 眼中藐視匪類         | 5      | 不放債取利       |

|         |                      |        |                  |

|         | 尊重那敬畏耶和華的人 |        | 不受賄賂以害無辜 |

|         |                      |        |                  |

--------------------------------------------------------------

若果這首詩是踏入禮拜堂前的要理問答,

我們要答對了,進入神的帳幕,其實不難,

只要我們背熟,只要我們有好的記憶力,

只要我們識字,我們就可以過關.

但要理問答的目的不是但求過關,

而是我們心口如一的作答,出自真誠的作答,

這10項條件,可以作為我們自我檢視的條件.

不單答案正確,而且行事正確,

不單是正統的神學orthodoxy,

而且是正確的行為orthopraxy,

正統的教義和正直的操守The right doctrine and the right conduct,

聖經不單給我們消息資訊information,更要求我們生命演進transformation.

要理問答給我們不單是一個可以進去的地方place,

更是一條可以引往那個地方的通道path.

到了新約,我們明白,這是耶穌基督的救恩,

希伯來書十:19,「19弟兄們,我們靠著耶穌的血得以坦然進入至聖所,」

舊約要回答要理問答,要檢視生命,才可以進入聖所,

到了新約,我們可以進入至聖所,是因為耶穌釘十字架,

祂的受苦,祂的死,祂的代贖,

而且是「坦然」進入至聖所,

也就是堂堂正正,挺胸昂首,大搖大擺,很神氣,很有自信的,進入至聖所,

完全是因為耶穌,

耶穌說:我就是道路,真理,生命,若不藉著我,沒有人能到父那裡去.

唯有在基督裡的人,與上主有新關係的人,

所以,我們不單藉著耶穌進入聖所,

我們在進入聖所之後,要仰望基督,要尋找恩主,要靠近耶穌,要緊貼上主的胸懷.

稍後我們要一齊領受聖餐,就是要用聖餐餅和聖餐杯,

來記念主耶穌對我們的恩情,

來堅固我們對主耶穌的信心,

來確認我們對主耶穌的忠心.

願聖父聖主聖靈,三位一體的神賜福大家.



\newpage

\section{願榮光歸上帝}
\label{sec:489}
\textbf{劉國偉先生 - 願榮光歸上帝}
\newline
\newline
\hyperref[sec:488]{< < < PREV SERMON < < <}
~
\hyperref[sec:toc]{[返目錄]}
~
\hyperref[sec:490]{> > > NEXT SERMON > > >}


講題:願榮光歸上帝 Glory to God

講員:劉國偉先生

場合:香港中文大學崇基學院禮拜堂主日崇拜

日期:2020年2月9日

引言

今天的講題是「願榮光歸上帝」.擬定這個題目的時候,讓我聯想起美國著名的作詞人芬尼·克羅斯貝Fanny Crosby所創作的一首詩歌To God Be the Glory,華人教會比較多選唱的中譯版本,是宣道出版社青年聖歌的「榮耀歸於真神」:

榮耀歸於真神,

祂成就大事,

為愛世人甚至賜下獨生子,

獻上祂生命,

為人贖罪受害,

永生門已大開,

人人可進來.

(To God be the glory, great things He hath done, So loved He the world that He gave us His son, Who yielded His life, an atonement for sin, And opened the Life Gate that all may go in.)

這一首歌詞所提及的是大家耳熟能詳的救恩要道,Fanny Crosby所填寫的詩詞,充滿力量,簡單有力地告訴人,歸榮耀予上帝的根本原因.不過,可能我們信主時日長久,有時就難免淡忙了初信主時,那份心靈悸動和興奮,對於Fanny Crosby所寫的「榮耀歸於真神,祂成就大事」,就可能感覺有點淡然.

今天,對我們來說,在生活中所經歷的事情,甚麼才算是大事?對香港人來說,從去年6月開始,不斷演化的社會運動,並最近的新型冠狀病毒的威脅,當然是大事當中的大事,而且我更會形容這是又大又難的事.但願當我們經歷這些又大又難的事的時候,我們心中牢記著上帝為我們所成就,耶穌基督的救恩這件大事.這樣,我們就可以懷著信心,常存盼望,即使在最艱難的歲月中,仍然以歸榮耀給上帝,作為我們的人生導向.

耶和華的榮光

今天的經課以賽亞書58章所論述的焦點,是審判與救贖,先知作為守望者,受上帝託付,「要大聲喊叫,不可止息；揚起聲來,好像吹角.向我百姓說明他們的過犯；向雅各家說明他們的罪惡. 」(賽58:1)

主前八世紀,正直南北兩國,面臨亡國的挑戰.在此危急關頭,看起來,選民好像天天尋求上帝,樂意明白上帝的道,好像行義的國民,不會離棄上帝的典章,向上帝求問公義的判語,喜悅親近上帝,他們在宗教的追求上,甚至認真到進行禁食.然而,面對國力衰弱,強敵當前,北國亡於亞述之事實(主前722年),看似大能臨頭之際,主前八世紀的猶太人不禁要問上帝:我們禁食,你為何不看見呢? 我們刻苦己心,你為何不理會呢?

上帝的回應是:「 看哪,你們禁食的日子仍求利益, 勒逼人為你們做苦工.你們禁食,卻互相爭競,以凶惡的拳頭打人.你們今日禁食,不得使你們的聲音聽聞於上.這樣禁食豈是我所揀選,使人刻苦己心的日子嗎?豈是叫人垂頭像葦子,用麻布和爐灰鋪在他以下嗎?你這可稱為禁食,為耶和華所悅納的日子嗎? 」(賽58:2-5)用我們今日的語言,就是有信仰的外貌,但在神的眼中,卻並沒有信仰的實質.

那麼,信仰的實質是什麼呢?上帝指向信仰的實踐:「我所揀選的禁食不是要鬆開凶惡的繩,解下軛上的索,使被欺壓的得自由, 折斷一切的軛嗎?不是要把你的餅分給飢餓的人,將飄流的窮人接到你家中, 見赤身的給他衣服遮體,顧恤自己的骨肉而不掩藏嗎?」(賽58:6-7)

弟兄姊妹,真實的敬拜是將生命放在生活的祭壇上,唯有如此,我們才可以發出生命的光輝,並且得到上帝的醫治:「你的光就必發現如早晨的光；你所得的醫治要速速發明.」(賽58:8上)然後,「耶和華的榮光必作你的後盾.」 (賽58:8下)

這段經文給我最大的安慰,是這幾句話:「那時你求告,耶和華必應允；你呼求,祂必說:我在這裡.」上帝願意親自向我們說「我在這裡」；神不離開我們,比甚麼都重要!當我們感到前景暗昧不明,卻仍然堅持信仰的實踐,耶和華的榮光是可以作我們的後盾；在人間的黑暗歲月中,有神的同在,我們的光就可以如早晨的光一樣,驅走黑暗的長夜.

在黑暗中有光

我們今天誦讀的詩篇112篇,其實蘊含著與以賽亞書58章相同的意象:在黑暗中有光.

詩篇112篇是一篇離合體詩篇,即是22行詩句的每一行的第一個字,是按照希伯來文22個字母,順序寫成,可以用作公共崇拜時所採用的詩章,是用心用意的創作.

詩人以讚美的詩句開始,鼓勵會眾要讚美耶和華!宣告敬畏耶和華,甚喜愛祂命令的, 這人便為有福!詩人進一步詮釋,敬畏上帝,喜愛祂的命令的人,為何有福.一個真心以讚美及敬拜上帝,作為信仰核心的人,在每天生活裡就要活得正直,他的公義存到永遠,即是在他離開人世之後,仍然有功效,這樣正直的人是有福的. 詩人特別強調「正直人在黑暗中,有光向他發現；他有恩惠,有憐憫,有公義. 」(詩112:4)現代中文譯本如此翻譯:「光為著義人在黑暗中照耀,為著慈愛,憐憫,正直的人照耀.」

這個讚美及敬拜上帝的正直人,詩人形容「他必不怕兇惡的信息；他心堅定,倚靠耶和華. 他心確定,總不懼怕,直到他看見敵人遭報.」(詩112:7-8) 「敵人」亦可以按第10節詮釋為「惡人」；而「遭報」可譯作「衰敗」；所以,「直到他看見敵人遭報」可以譯作「直到他看見惡人的衰敗」.

弟兄姊妹,也許兇惡的信息時常環繞我們,但我們的心堅定,我們確信黑暗之中有光,照耀著我們,因為我們活得有恩惠,有憐憫,有公義.我們也確信惡人不會長久.我想起了湖北武漢的李文亮醫生,他本著誠實,正直說真話,頂著被公權力嚴厲警告和訓誡的壓力,發出了病毒到臨的警告信息,他也在這場搶救生命的戰爭中,獻上了自己的生命.我相信,正如詩人說:「李文亮醫生的仁義存到永遠. 他的角必被高舉,大有榮耀.」(詩122:9下)

基督的榮耀

談到歸榮耀予上帝,保羅在哥林多前書第2章的教導,也提出了一個很好的論點,去啟導我們的思考.

保羅一再強調,要倚靠聖靈的大能去證實所傳的福音,並且以上帝的大能作為信仰的基礎:

「我說的話,講的道,不是用智慧委婉的言語,乃是用聖靈和大能的明證,叫你們的信不在乎人的智慧,只在乎上帝的大能.」(林前2:4-5)

保羅所用「大能」(dunamis)這個希臘字,在新約著作中出現了很多次,散見於十多卷著作中,核心的含義是能力,權柄,剛強,可以涉及異能,奇事,大多數都指向上帝.

正如保羅所言,華麗的詞藻,高深的學問,委婉動聽的智慧之言,都不足以證實所傳的福音,靠的是聖靈的大能,信仰亦自然以上帝的大能做基礎.「大能」這個字正要表明給我們看,人類那個不能自救的真相；是上帝把我們帶進祂的榮耀裡,而並非我們可以做什麼,或者獻上什麼,可以換取上帝的救贖；所以保羅說:「上帝在萬世以前,預定使我們得榮耀」(林前2:7下) 這與詩篇112篇的作者是一致的:「光為著義人在黑暗中照耀」,都是上帝的大能,讓我們可以活在祂的榮耀中.

當這個世界有權有位的人(林前2:6；2:8)自以為有智慧,有能力,甚至耶路撒冷的猶太精英,將提出問題的耶穌釘在十字架上,顯然他們認為自己決定英明,而這些人是難以領悟和體會,耶穌基督就是榮耀的主.保羅怎樣去形容這位榮耀之主,去為我們預備那榮耀的永恆所在?

正是眼睛沒有見過,

耳朵沒有聽過,

也從來沒有人想到的!(林前2:9)

弟兄姊妹,我們所渴慕的榮耀,並非這個世界的追求,他們也不會明白歸榮耀予上帝的人生取向.

光照人前

親眼見過這位榮耀之主的使徒約翰,以榮光來描述祂:「我們也見過祂的榮光,正是父獨生子的榮光.」(約1:14下)「光」這個意象在福音書中被普遍採用,相關的字辭出現超過50次.我們今日的福音經課,記載了主耶穌的教導,也與光有關:「你們是世上的光.城造在山上是不能隱藏的.人點燈,不放在斗底下,是放在燈臺上,就照亮一家的人.你們的光也當這樣照在人前,叫他們看見你們的好行為,便將榮耀歸給你們在天上的父.」 (太5:14-16)

主耶穌強調門徒要成為世上的光,說出這個比喻的背景,是一個充斥著文士和法利賽人的律法主義的信仰世界,與主前八世紀的以賽亞時代的猶太人的世界,竟然如此相近:有信仰的外貌,甚至是一個嚴緊的外貌,但在耶穌的眼中,卻並沒有信仰的實質.故此,耶穌說明得十分清楚:把光照在人前,以好行為將榮耀歸給在天上的父.

無論是以賽亞所言:公義必在你前面行, 耶和華的榮光必作你的後盾；又或者詩人在讚美敬拜時發出的詩句:正直人在黑暗中,有光向他發現；他有恩惠,有憐憫,有公義；這些經文都在黑暗長夜裡,向走在窄路中的天路客,發出激盪心靈的邀請；我們知道這完全在乎上帝的大能,我們就順著聖靈的能力,把光照在人前,叫活在這個紛亂的世界中的人,看見我們的好行為,便將榮耀歸給我們在天上的父.

願榮光歸上帝

Fanny Crosby在1820年,生於美國紐約東南面一個小鎮,因為家貧和誤診,六個星期大的時候,因感冒引起眼炎,最後完全失明.在她一歲的時候,父親早逝,由母親和祖母撫養長大.

這位自六星期大,就活在黑暗無光的現實世界中的作詞家,一生的作品數以千計,被譽為史上最多產的教會詩歌填詞人.

Fanny於1915年辭世,為這個世界留下了無數的美好的詩詞,還有一個榮耀上帝的生命見證.聖靈使用Fanny Crosby的詩詞,和她的生命見證,叫很多人蒙福,包括基督徒和很多尋求信仰的人.

在她的墓碑上,刻著她所寫的另一首名作的幾句詩詞:

有福的確據 Blessed Assurance

Blessed assurance, Jesus is mine!

Oh, what a foretaste of glory divine!

有福的確據,耶穌屬我!

噢!我今預嚐,主榮耀恩典!

我羨慕Fanny Crosby在黑暗中,有光向她發現；她的光就如早晨的光；她的光照在人前；倚靠上帝的大能,活出了榮耀主的人生.

願榮光歸予聖父,聖子,聖靈.



\newpage

\section{這樣的因這樣的果}
\label{sec:490}
\textbf{鄧瑞強博士 - 這樣的因,這樣的果}
\newline
\newline
\hyperref[sec:489]{< < < PREV SERMON < < <}
~
\hyperref[sec:toc]{[返目錄]}
~
\hyperref[sec:491]{> > > NEXT SERMON > > >}


申命記 30:15-20
\newline
\begin{longtable}{cl}
\hline
\hline
章節 & 經文 (和合本修訂版)\\
\hline
30:15 & \begin{tabularx}{0.7\textwidth}{X} 「看，我今日將生死禍福擺在你面前。 \end{tabularx} \\ \\ \relax
30:16 & \begin{tabularx}{0.7\textwidth}{X} 我今日所吩咐你的，就是要愛耶和華－你的神，遵行他的道，謹守他的誡命、律例、典章，使你可以存活，增多，而且耶和華－你的神必在你所要進去得為業的地上賜福給你。 \end{tabularx} \\ \\ \relax
30:17 & \begin{tabularx}{0.7\textwidth}{X} 倘若你的心偏離，不肯聽從，卻被引誘去敬拜別神，事奉它們， \end{tabularx} \\ \\ \relax
30:18 & \begin{tabularx}{0.7\textwidth}{X} 我今日向你們申明，你們必定滅亡；在你過約旦河進去得為業的地上，你的日子必不長久。 \end{tabularx} \\ \\ \relax
30:19 & \begin{tabularx}{0.7\textwidth}{X} 我今日呼天喚地向你作見證：我已經將生與死，祝福與詛咒，擺在你面前。所以你要揀選生命，好使你和你的後裔都得存活。 \end{tabularx} \\ \\ \relax
30:20 & \begin{tabularx}{0.7\textwidth}{X} 要愛耶和華－你的神，聽從他的話，緊緊跟隨他，因為他是你的生命，必使你的日子得以長久，可以在耶和華向你列祖亞伯拉罕、以撒、雅各起誓要給他們的地上居住。」 \end{tabularx} \\ \\
[1ex]
\hline
\hline
\end{longtable}
講題:這樣的因,這樣的果 What You Choose is What You Become

經文:申命記30:15-20

講員:鄧瑞強博士

場合:香港中文大學崇基學院禮拜堂主日崇拜

日期:2020年2月16日

各位弟兄姊妹,早安.在這疫症蔓延的日子,願各人平安.在平常的日子,我們很容易覺得「平安」是理所當然的,但在疫症蔓延的日子,我們會明白「平安」原來已是從天而來的恩典,從神而來的寶貴禮物.

疫症,容易觸發人對生死的感悟,容易讓人體會生死的無常,容易讓人明白人的有限和無助.哲學家海德格(Heidegger)分析人意識到自己會死這個事實,他指出,人在平常生活中,會營營役役,會迷失在群眾中,會遺忘存在的意義；但是,當人面對自己的死亡時,他會猛然醒覺,這是他自己要親身面對的經驗,這是一種「無人可代替你」的經驗,無論你身邊有多少人,你總要孤獨地親歷自己的死亡,這時,一個人的「真正自己」,他的「個我存在」,便會從人群中凸現出來,是「你」自己,不是「他人」,去經歷你自己的死亡,無人可代替你去經歷你這個最切身的經驗.這時,唯有這時,人在天地間,意識到自己的獨特性,意識到自己要為自己如何面對死亡負上全責,同時,意識到自己要為自己整個一生的存在意義負起無可推卸的責任.在生死的大限前,人深刻的意識到,自己要為一生的意義,一生的禍福,負起全責.

今日的講道經文,是今日的舊約經課:申命記30:15-20.這經文的核心內容,就是:生命的禍福,我們自己要負起全責.我再讀這段經文時,請留意有兩個字詞常常出現,一是「生命」或與「生命」有關的字詞,二是「今日」這字詞.

申30:15　「看哪,我今日將生與福,死與禍,陳明在你面前.

申30:16　〔原文有「今日」的字眼〕吩咐你愛耶和華─你的神,遵行他的道,謹守他的誡命,律例,典章,使你可以存活,人數增多,耶和華─你神就必在你所要進去得為業的地上賜福與你.

申30:17　倘若你心裡偏離,不肯聽從,卻被勾引去敬拜事奉別神,

申30:18　我今日明明告訴你們,你們必要滅亡；在你過約旦河,進去得為業的地上,你的日子必不長久.

申30:19　我今日呼天喚地向你作見證；我將生死禍福陳明在你面前,所以你要揀選生命,使你和你的後裔都得存活；

申30:20　且愛耶和華─你的神,聽從他的話,專靠他；因為他是你的生命,你的日子長久也在乎他.這樣,你就可以在耶和華向你列祖亞伯拉罕,以撒,雅各起誓應許所賜的地上居住.」

「生命」或與「生命」有關的字詞出現六次,「今日」這字詞也出現六次,短短六節經文,平均每節都出現一次,好像向我們說,你「今日」就在決定生與死的關頭,你要選擇「生命」.

這段經文有雙重的歷史背景,第一重歷史背景是文學上的(textual),是經文文字編織出來的歷史背景,在這歷史背景裡,這段經文是摩西的臨別贈言.摩西帶領群眾離開埃及後,在曠野流浪了四十年,就在進入迦南地前夕,但同時也是摩西即將離世的前夕,他說了剛才經文這番話.那時,摩西站在摩押平原,前面是約旦河.過了約旦河,就是耶利哥.踏入耶利哥,就進入迦南地,也就是神應許之地.摩西遙望那地,他遙望到將來,他站在歷史的緊要關頭.過了這河,過了這關,歷史將會改變.但他自己死期將至,進不了迦南.歷史的將來,是不確定的.就在這歷史轉折的關頭,摩西說,明天,是福是禍,是生是死,不是由天注定,是由人今天的抉擇鑄成的.

剛才講到,這段經文有雙重的歷史背景,第二重的歷史背景,牽涉到申命記最後編輯成書的年代.大家都知道,申命記的結尾,提到摩西的死,故申命記不可能全由摩西寫成.摩西留下的思想遺產,由後人加以編輯修訂.我們現在相信,申命記的最後編纂,發生在公元前六世紀,這是被擄巴比倫的以色列人被釋放而可以回歸故土的前夕.這是另一個歷史時空.申命記的編纂者,是透過這編輯過程,讓經文對另一個時代的人說話.若大家細心留意申命記29: 25-28,都可看到編纂者留下的文字,那裡明示他們離開摩西的年代已有一段歷史了.長話短說,申命記的經文,對應著另一個歷史背景,它向著即將離開巴比倫而回歸故土的以色列人說話.這些人已住在巴比倫70年,整整一代人,他們根本不知道故土是什麼模樣,他們對回歸後前景如何毫無頭緒.要他們回歸,是面對不確定的將來.就在這歷史轉折的關頭,經文的編纂者,透過摩西的口,說:明天,是福是禍,是生是死,不是由天注定,是由人今天的抉擇鑄成的.

在歷史的轉折點,面對不測的未來,人應如何自處?人應如何抉擇?

我們聽聽摩西的話.

申30:15　「看哪,我今日將生與福,死與禍,陳明在你面前.

今日,當下,生死禍福,就在我們面前.我們像昔日的以色列人一樣,站在歷史的轉折點上,前路未卜.猶如行在脆弱的薄冰上,如何行,腳應踏在何處,一步一驚心.走對了,就能渡到彼岸；走錯了,冰碎了,便墮入冰海,失去生命.如何抉擇,關乎生死.

申30:16　〔原文有「今日」的字眼〕吩咐你愛耶和華─你的神,遵行他的道,謹守他的誡命,律例,典章,使你可以存活,人數增多,耶和華─你神就必在你所要進去得為業的地上賜福與你.

今日,當下,摩西說,人一生的禍福,就在當下決定.不要想,今天錯就錯吧,讓明天才去補救.不要這樣想.讓每一刻的當下,都是我們面對神的永恆時刻.我們當下的想法,就判定我們將來的生死禍福.

摩西在這裡,吩咐我們三樣東西.一,是「愛耶和華－你的神」；二,是「遵行他的道」；三,是「謹守他的誡命,律例,典章」.有了這三樣,就能走在生命的路上.

一,愛神,就是一生純一地渴慕神,心無罣礙,不受誘惑.這涉及「情之所至」.二,「遵行他的道」,重點是「行」,walk in his ways,不行世人的路,唯獨行神的路,或者說,行耶穌走過的路.這涉及「行之所向」.三,「謹守他的誡命,律例,典章」,重點是「守」,keep his commandments,看神的誡命為至寶,守在心靈的夾萬裡.這涉及「意之所止」.

站在歷史的轉折點上,前路未卜,生死未可預料.人的種種謀算不能保證你走在生命路上.唯有以上這三樣東西,才能讓你一生平安.

近來,生活裡,常常面對抉擇,雖然算不上是什麼生死抉擇,但或可以有嚴重後果.如:有人說,某某地方有口罩賣.我應不應該去排隊?有人說,某某網上商店可以預訂口罩,是不用預先付錢的.嘩,這麼著數,不用付任何代價.我應不應該上網訂購?有人說,這裡有米.有人說,那裡有廁紙.我應怎樣做呢?我沒有排隊去買口罩,也沒有上網訂購,我想,我還有少許口罩,我盡量省著用,應該有更多人比我有更迫切的需要,希望他們買到吧.但是,但是,若有一天,我也用完我僅有的口罩,而又買不到,那時,我會後悔嗎?我不知道.有人說,現在不去買廁紙,有一天,你就可能要用上Ａ4紙.這或許是真的.那時,我會後悔嗎?我不知道.不過,我知道,我知道,我應該愛神,也應該努力愛人如己.我知道應該讓機會給更有需要的人.我知道,我應努力行在耶穌走過的路上,他沒有為明天憂慮,我要努力學他.我知道,我應努力「謹守神的誡命,律例,典章」,這些才是我應努力「守護」的,而不是努力去「守護」,去「keep」住那些廁紙.我相信摩西所講的,這才是「生」路,帶來「生命」的路.

在生死關頭,人會恐慌,人會搶奪,人會失去靈性,獸性大發,人會為保住自己而不擇手段,這是死路.在這時候,更要明白,在生死關頭,是我們面對神的永恆時刻.我們的抉擇,有永恆的重量.

申30:17　倘若你心裡偏離,不肯聽從,卻被勾引去敬拜事奉別神,

申30:18　我今日明明告訴你們,你們必要滅亡；在你過約旦河,進去得為業的地上,你的日子必不長久.

摩西說,當人心中不再純一地愛神,而去事奉別的神時,他的生命便走向滅亡了.當人認為某種世物是很重要的,重要到一個人可以遺忘信仰生命的品格時,這個人便是視這東西為「終極」的,實際上便是視之為「神」了.摩西說,以別神代替耶和華的,他便走在滅亡的路上.有些人搶廁紙,一搶便搶十條廁紙,這一搶,連人應有的「格」也丟失了.廁紙成了他的「終極關懷」,成了他的「神」,他事奉了這神,離棄了耶和華.在人生的路上,他選擇了死亡的路.

申30:19　我今日呼天喚地向你作見證；我將生死禍福陳明在你面前,所以你要揀選生命,使你和你的後裔都得存活；

申30:20　且愛耶和華─你的神,聽從他的話,專靠他；因為他是你的生命,你的日子長久也在乎他.這樣,你就可以在耶和華向你列祖亞伯拉罕,以撒,雅各起誓應許所賜的地上居住.」

第20節有個字眼值得留意,「專靠」他.「專靠」,指生命的連合,緊緊的依靠.聖經第一次出現這字眼,就是創世記講到神創造男與女的時候.創世記2: 24說:「人要離開父母,與妻子連合(這就是「專靠」那個字),二人成為一體.」

摩西說,不單愛神,不單行在神的路上,還要「專靠」神,與神的生命連合,用新約的話來說,就是「活像基督」,「活出基督」,「活著就是基督」.這樣,我們就能得著生命.

1962年,美蘇的對峙,爆發了古巴的核戰危機.這時美蘇核彈的數量與威力,都遠遠超過二戰時的原子彈.美蘇雙方,為了克制對方,都在對方領土的邊緣,部署核導彈.美國在土耳其部署了瞄準蘇聯的核導彈.蘇聯便嘗試在美國旁邊的古巴同樣部署核導彈.裝載核導彈設備的蘇聯艦艇開往古巴,美國便在古巴沿海封鎖古巴,阻止蘇聯艦艇前進.雙方對峙,劍拔弩張.這時,圍封古巴海岸線的美國艦艇發現到海底有蘇聯的潛水艇,於是,投下沒有摧毀性的深水炸彈去迫使它浮上來.美國人不知,這潛水艇載有核彈.由於這潛艇在海底已很久,也沒法收到莫斯科的信號,遭遇這些深水炸彈,潛艇上的軍官以為美蘇已開戰.按先前訂好的戰略,潛艇有權發射核彈.那邊廂,美國的偵察機在古巴上空被擊落,美國認為蘇聯不打算談判,戰爭的按扭一觸即發.戰爭一發動,這就不可能是一般意義的戰爭,而是核彈橫飛,連地球也受到嚴重摧毀的戰爭.這是歷史轉折的重要關頭.在這關頭作抉擇的人,他們作的抉擇,都關乎千萬人的生死.潛艇上的蘇聯軍官,最終沒有按下發射核彈的按扭.美國的軍官,同樣沒有按下戰爭的按扭.甘迺迪叫停了戰爭,赫魯曉夫也叫停了戰爭.人類最後避過一劫.

這個故事我講了這麼長時間,是要說明歷史的將來,生命的將來,是禍是福,是生是死,很在乎我們當下如何抉擇.前景未明,前途未卜,將來是開放的.擺在我們的前路,容讓我們去選擇,有生路,有死路.生路與死路,關鍵是我們到底是否愛神?到底是否在愛神的心態下小心行每一步?到底是否相信神給信靠他的人的生命應許?今天,是考驗我們信仰的關鍵時刻.

請再聽摩西的話:

申30:19　我今日呼天喚地向你作見證；我將生死禍福陳明在你面前,所以你要揀選生命,使你和你的後裔都得存活；

申30:20　且愛耶和華─你的神,聽從他的話,專靠他；因為他是你的生命,你的日子長久也在乎他.」

但願:榮耀歸於聖父,聖子,聖靈.阿們.

(附記:本人在崇基禮拜堂的講章已結集成《荷花池畔春草綠》講道集,試讀本可參閱:\url{http://cc.cuhk.edu.hk/dev/images/content/fund-raising-campaigns/university-sermons-by-dr-sk-tang/CCC-SK-Tang-Sermon-Text-Sample-2.pdf}.較早期的崇基講章,也已結集成《敢於跟隨主》一書.有興趣的朋友可到各大基督教書局購買這兩書.)



\newpage

\section{遇見神}
\label{sec:491}
\textbf{林豪恩先生 - 遇見神}
\newline
\newline
\hyperref[sec:490]{< < < PREV SERMON < < <}
~
\hyperref[sec:toc]{[返目錄]}
~
\hyperref[sec:492]{> > > NEXT SERMON > > >}


馬太福音 17:1-9
\newline
\begin{longtable}{cl}
\hline
\hline
章節 & 經文 (和合本修訂版)\\
\hline
17:1 & \begin{tabularx}{0.7\textwidth}{X} 過了六天，耶穌帶著彼得、雅各和雅各的弟弟約翰，領他們悄悄地上了高山。 \end{tabularx} \\ \\ \relax
17:2 & \begin{tabularx}{0.7\textwidth}{X} 他在他們面前變了形像，他的臉明亮如太陽，衣裳潔白如光。 \end{tabularx} \\ \\ \relax
17:3 & \begin{tabularx}{0.7\textwidth}{X} 忽然，有摩西和以利亞向他們顯現，與耶穌說話。 \end{tabularx} \\ \\ \relax
17:4 & \begin{tabularx}{0.7\textwidth}{X} 彼得回應，對耶穌說：「主啊，我們在這裡真好！你若願意，我就在這裡搭三座棚，一座為你，一座為摩西 ，一座為以利亞。」 \end{tabularx} \\ \\ \relax
17:5 & \begin{tabularx}{0.7\textwidth}{X} 說話之間，忽然有一朵明亮的雲彩遮蓋他們，又有聲音從雲彩裡出來，說：「這是我的愛子，我所喜愛的。你們要聽從他！」 \end{tabularx} \\ \\ \relax
17:6 & \begin{tabularx}{0.7\textwidth}{X} 門徒聽見，就俯伏在地，極其害怕。 \end{tabularx} \\ \\ \relax
17:7 & \begin{tabularx}{0.7\textwidth}{X} 耶穌進前來，拍拍他們，說：「起來，不要害怕！」 \end{tabularx} \\ \\ \relax
17:8 & \begin{tabularx}{0.7\textwidth}{X} 他們舉目，不見一人，只見耶穌獨自一人。 \end{tabularx} \\ \\ \relax
17:9 & \begin{tabularx}{0.7\textwidth}{X} 下山的時候，耶穌囑咐他們說：「人子還沒有從死人中復活，你們不要把所看到的告訴人。」 \end{tabularx} \\ \\
[1ex]
\hline
\hline
\end{longtable}
講題:遇見神 Encounter with God

經文:馬太福音17:1-9

講員:林豪恩先生

場合:香港中文大學崇基學院禮拜堂主日崇拜

日期:2020年2月23日

一. 引言

今天是「登山變像日」,當年門徒就在那天見到耶穌神聖的形象.你有沒有遇見神的經驗呢?或者退後一步問,你想不想遇見神呢?遇見神會是怎樣的經驗呢?張敬軒曾經唱過一首歌,歌名就是「遇見神」,描述的經驗是這樣的:

看天邊飄過雲海,告訴世界幻變常在,

從來人在月缺下盼月圓.

天天總有上演,告訴你我動人故事,

動人情節留下伏線,細心飾演!

如何人心粉碎像微塵,無言留淚,滿身抖震,

如何遇見神,被擁抱,再得起身!

明明平安彷彿天邊遠,流連遊盪,暗失方寸.

然而遇見神,被指引,再生精采心內滲.

誰無遐想一生精采佳美?

由誰來導,演出好戲?

人能遇見神,是福氣,最終高飛.

從來難數高低多少次,何時成就?那天失意?

前途幻變時,路雖遠,有這福氣心內暖.

歌詞提到遇見神是福氣,這福氣的具體表現有三重:第一,「如何人心粉碎像微塵,無言留淚,滿身抖震,如何遇見神,被擁抱,再得起身」；第二,「明明平安彷彿天邊遠,流連遊盪,暗失方寸.然而遇見神,被指引,再生精采心內滲」；第三,「人能遇見神,是福氣,最終高飛」.總括而言就是:「再得起身」,「再生精采」,「最終高飛」.今天的福音書經課,記載了門徒遇見神的經驗,讓我們透過這片段,思考遇見神可以是是一種怎樣的狀況,以及帶給遇見神的人怎樣的影響.

二. 經文及解釋

17:1過了六天,耶穌帶著彼得,雅各,和雅各的兄弟約翰,暗暗地上了高山,

17:2就在他們面前變了形像,臉面明亮如日頭,衣裳潔白如光.

17:3忽然,有摩西,以利亞向他們顯現,同耶穌說話.

17:4彼得對耶穌說:「主啊,我們在這裡真好!你若願意,我就在這裡搭三座棚,一座為你,一座為摩西,一座為以利亞.」

17:5說話之間,忽然有一朵光明的雲彩遮蓋他們,且有聲音從雲彩裡出來,說:「這是我的愛子,我所喜悅的.你們要聽他!」

17:6門徒聽見,就俯伏在地,極其害怕.

17:7耶穌進前來,摸他們,說:「起來,不要害怕!」

17:8他們舉目不見一人,只見耶穌在那裡.

17:9下山的時候,耶穌吩咐他們說:「人子還沒有從死裡復活,你們不要將所看見的告訴人.」

供奉神明

門徒見到耶穌的神聖形象的第一個反應就是:

17:4彼得對耶穌說:「主啊,我們在這裡真好!你若願意,我就在這裡搭三座棚,一座為你,一座為摩西,一座為以利亞.」

對於彼得的回應,有人對他提出兩項指責,第一是「上得到上卻不下了山」,二是「貶低耶穌與摩西和以利亞同等」.在這裡,我想替彼得說幾句平凡的話.

第一,關於「上得到上卻不下了山」的指責,意思即是他們只顧享受山上與主同在的美好,而逃避山下現實世界的需要.其實,門徒可能不是要留在山上不想下山,因為彼得建議的,只是搭三座棚,而不是建三座房子,三座棚並不是長期生活的好地方.此外,門徒建議搭三座棚,可能只是沿襲列祖的傳統,舊約記載他們曾經在遇見神的位置留下痕跡,例如亞伯拉罕和雅各都這樣做過.

第二,關於「貶低耶穌與摩西和以利亞同等」,指責者認為,憑彼得的口供,證據確鑿,因為他曾說:「我就在這裡搭三座棚,一座為你,一座為摩西,一座為以利亞」.不過,從我們的民間信仰的角度來看,又未必可以肯定彼得把耶穌視為與摩西及以利亞同等的.

我們身處的沙田區,有一間出名的廟宇,每年農曆年初二,新界鄉事領袖都會去那廟宇為香港求籤.這廟宇主殿中央放置了巨大的車公神像,而在主殿內兩邊盡頭,則分別是福祿宮(財神殿)及元辰殿.其實不少廟宇也有類似的設計,除了主神殿之外,旁邊還有陪伴或守護的神殿.由此看來,主次和輕重還是有分的.在疑點歸於被告的原則之下,請陪審團裁定彼得的罪名不成立.

各位陪審員,多謝你們以我們民間信仰的觀念來為彼得平凡,如此,彼得被釋放了,不過,困難卻留下了給我們.如果用我們民間信仰的觀念去理解,為耶穌,摩西,以利亞搭三座棚,其實是供奉神明的概念.我們撫心自問,我會否把上主視為供奉的對象呢?把上主視為供奉的對象即是甚麼意思呢?網絡上很多資訊指導人們如何供奉神明,其中一篇名為:「為甚麼要奉茶?神明,祖先絕對要這三杯」如此說:「這三杯茶水要勤於更換,且絕對不能空杯,因這三杯茶水被視為神明與祖先所賜的「財」,更是由神明與祖先庇佑的「財源」,因此不可不供茶水,也需勤勞更換,保持茶水乾淨,絕不乾涸,如此,家運,財運,事業運才能財運滾滾.」我們可以反駁,我們沒有向耶穌像或者向十字架奉茶,更沒有上香,怎能說我們如民間信仰一樣供奉神呢?重點不在有沒有奉茶或上香,而是把神視為「家運,財運,事業運」的交換對象.我們提供給神的交易物品不是茶水或香燭,而可能是守禮拜和守禮節；尤有甚者,連禮拜和禮節都不守,兩手空空而來,向神進行感情勒索:「若你是愛我的,你應該保守我身體健康,家庭幸福,工作順利,事業亨通,心想事成……」.讓我們問一問自己,我們對上主的態度與民間信仰供奉神明相差多遠呢?又讓我們從上主的那一邊想想,上主會否期望我們如民間信仰供奉神明一樣對待祂呢?

害怕／敬畏神

讓我們繼續看看彼得提出這建議之後有甚麼事情發生.

17:5說話之間,忽然有一朵光明的雲彩遮蓋他們,且有聲音從雲彩裡出來,說:「這是我的愛子,我所喜悅的.你們要聽他!」

雲彩及聲音出現之後,門徒的反應值得我們深思:

17:6門徒聽見,就俯伏在地,極其害怕.

遇見神怎麼會感到害怕呢?害伯神這種感覺和反應是否已經離我們很遠了?

巴刻在《認識神》第十六章有一段如此說:

「現代人已慣於把神的恩慈和他的嚴厲分開,今天在一般人當中,已拒絕所有神忿怒和審判的內容,取消神的嚴厲,只做留神的仁慈,把神當作天上聖誕老人.很多人堅信,神是無限的忍耐和仁慈,也不會想到:神對我的態度,會受到我是否照他所說的去做所影響.任何人若嘗試叫別人在神面前存敬畏的心,在他的話語中戰兢,立刻會被認為老土得無可救藥.」

「害怕」這個字一字兩義,一方面是「驚恐,懼怕某人或某事」,同時另一方面是「尊敬,有很深的敬意」,其實,聖經多次用「害怕」來描述人們對耶穌的回應,例如馬大福音9:1-8記述耶穌醫好一個癱子後人們的回應:

9:8眾人看見都驚奇(害怕),就歸榮耀與神,因為他將這樣的權柄賜給人.

就是在末後的日子,啟示錄15:4也如此說:

「主啊,誰敢不敬畏(害怕)你,不將榮耀歸與你的名呢?因為獨有你是聖的.萬民都要來在你面前敬拜,因你公義的作為已經顯出來了.」

有人說,害怕／敬畏神是古代人無知的表現,現代人已經擺脫了過去的愚昧,慷慨激昂地宣告:「從來就沒有什麼救世主」.不過,當現代社會把神驅趕之後,理想世界就出現了嗎?還是進入一個「遇神殺神,遇佛殺佛」更可怕的世界呢

大約十年前,有一位風水師成為城中熱話,因為他稱自己是龔如心千億遺產的唯一受益人而展開漫長的官司.案件峰迴路轉,他中途被控偽造及行使虛假文書罪,罪成判刑12年.後來有傳風水師信了耶穌,有署名韓世光在報紙以「心中有鬼好過目中無神」為題發表評論:

「惡貫滿盈的陳振聰到底是真心誠意信主懺悔,還是像臨急抱佛腳那樣尋求另一種超自然力量來庇佑自己?當然不排除有人誤把主耶穌當作黃大仙,不過有一點是值得玩味的,正如法官所言,也許此君心靈上有所改變,因此成為基督徒之後已經沒有再上庭作假口供.

姑勿論如何,壞人尋求宗教慰藉是一件好事,如果壞人相信為非作歹會受到上天懲罰,傷天害理之時起碼沒有那麼兇狠……

共產主義者採用的是無神論,目中無神也就無所敬畏,而無畏無懼的後果往往是無法無天.方言如今的神州大地,無官不貪,無商不奸……

所以我說,心中有鬼好過目中無神,敬神怕鬼的人通常比較膽小而善良,絕對不會販賣坑渠油毒害同胞.」東方日報2013年7月12日 (五) 韓世光:「心中有鬼好過目中無神」

在C. S. Lewis(魯益師)的名著《納尼亞傳奇》中,人們曾經為獅王亞斯蘭爭論,有人說,亞斯蘭只不過是一隻如貓一般的動物,只會喵喵叫幾聲；有人說亞斯蘭是兇猛的野獸,有人卻說亞斯蘭是可以親近的獅王.當亞斯蘭出現的時候,輕視祂如小貓一隻的人嚇破了膽,原來祂果真是兇猛的獅子.但是,兇猛的獅子怎可能又是可親近的呢?讓我們回到耶穌登山變象的故事給我們的的啟發:

17:7耶穌進前來,摸他們,說:「起來,不要害怕!」

原來,這位我們需要害怕的上帝也是解除我們害怕的上帝；這位令人俯伏在地的上帝也是進前來扶起我們的上帝；這位令人驚懼的上帝也是安慰的上帝.

三. 應用

最後,請我們思考門徒見到登山變象的耶穌的經歷,如何應用在我們的信仰與生活之中.第一個層面是關於我們基督徒的,第二個層面嘗試延伸至人類整體.

1．基督徒

首先,讓我們思考第一個層面.基督徒即是跟隨耶穌基督的人,雖然也是跟隨耶穌基督,但由於生存在不同時空,應該沒有如福音書中的門徒般遇見過耶穌,也沒有保羅往大馬色路上遇見耶穌的經歷.不過,我們也不需要因此而妄自菲薄,因為耶穌曾經說過:(約翰福音20:29)「你因看見了我才信；那沒有看見就信的有福了.」我們可以把每一次參加崇拜,都是一次與神相遇的機會.現在,讓我們發揮想像,把耶穌登山變象的故事重述為我們參加崇拜的故事.

17:1過了六天

主日崇拜在星期日舉行,我們剛好過了六天又參加下一次主日崇拜.

耶穌帶著彼得,雅各,和雅各的兄弟約翰,暗暗地上了高山,17:2就在他們面前變了形像,臉面明亮如日頭,衣裳潔白如光.

我們去參加主日崇拜,同時也可以理解為上主呼喚我們去崇拜祂；常這道,人尋找神,其實是神尋找人.崇拜開始的詩歌,帶領我們思想及進入神的榮耀.

17:6門徒聽見,就俯伏在地,極其害怕.

在神的榮耀,威嚴和聖潔之下,崇拜者「俯伏在地,極其害怕」是合適也合理的反應.在崇拜中,這是省察心靈認罪求赦的環節.問一問自己,面對聖潔公義的上帝,心情是怎樣的?心裡在想著甚麼呢?我想起耶穌在路加福音18:10-13節講過的故事:

「有兩個人上殿裡去禱告:一個是法利賽人,一個是稅吏.法利賽人站著,自言自語地禱告說:『神啊,我感謝你,我不像別人勒索,不義,姦淫,也不像這個稅吏.我一個禮拜禁食兩次,凡我所得的都捐上十分之一.』那稅吏遠遠地站著,連舉目望天也不敢,只捶著胸說:『神啊,開恩可憐我這個罪人!』

威嚴榮耀,聖潔公義的上帝呼喚我們省察自己的時候,難道我們仍然可以看見別人眼中的刺,而不想到自己眼中有樑木嗎?難道我們可以仍然堅持自己沒有錯,所有都是對方的錯嗎?難道我們可以不為自己的惡念和不見得光的行為而羞愧嗎?

17:7耶穌進前來,摸他們,說:「起來,不要害怕!」

這是崇拜中我們認罪後上主的宣赦,接納及復和,聖靈再一次充滿我們,賜給我們平安及力量更新心靈,繼續前行.

17:5說話之間,忽然有一朵光明的雲彩遮蓋他們,且有聲音從雲彩裡出來,說:「這是我的愛子,我所喜悅的.你們要聽他!」

接著我們領受聖道,透過經課及講解聆聽「聽他」.

17:9下山的時候,耶穌吩咐他們說:「人子還沒有從死裡復活,你們不要將所看見的告訴人.」

參加完崇拜離開禮拜堂,我們可以「將所看見的告訴人」,因為人子已經從死裡得活了.這是差遣,也是使命.教會的其中一款崇拜結束差遣文是這樣的:

「你們平安入世,勇敢為人,保持善德,勿以惡報惡；鼓舞灰心的人,扶持軟弱的人,援助困苦的人,尊敬眾人,愛主事主,在聖靈的能力中歡欣喜樂.」

2．社會大眾

現實上,社會中不是所有人都願意成為基督徒,因此,他們未必擁有上述的機會及經歷.不過,上帝是創造主,是全人類的上帝,門徒見到登山變像的耶穌的經歷,如何可以成為對人類整體有意義呢?門徒在這個經歷中體會到對神聖的敬畏,我們可否延伸至對上主所創造的世界呢?對上主創造的世界還著敬畏之情有可能是跨越宗教信仰的嗎?有一位朋友如此分享大自然令他產生的敬畏感,不知道你有沒有類似的經驗呢?

「當我在新疆塔克拉瑪幹沙漠邊,看到茫茫沙海中頑強挺立的活著的,死去的胡楊樹時；當我在青藏高原上通往珠峰大本營的公路上,望見雪山延綿,矗立在我周圍時；當我第一次看到海潮怒吼著湧向岸邊時；當我在飛機上看到大地伸展的褶皺時；當我仰望夜空,看到無數顆星在浩瀚宇宙中閃耀時……對大自然的敬畏,是悟道的時刻.」(網上文章,作者:落葉滿山谷)

對上主創造的世界心存敬畏有何重要呢?前幾天剛在社交群組收到一張圖片,下面有這幾句文字:

「一直以來,都是人類把動物關進籠子.今年春節,動物成功地把十幾億人關進籠子.」

這段日子,不但是出現在中國的病毒,還有在澳洲的山火,在東非出發的蝗蟲,這一切,都與人類如何對待大自然有關.

對上主的敬畏,延伸至對上主創造的世界的敬畏,就不會視大自然為剝削的對象,而隨意破壞.



\newpage

\section{塗抹的手才能寫出真實}
\label{sec:492}
\textbf{伍渭文牧師 - 塗抹的手,才能寫出真實}
\newline
\newline
\hyperref[sec:491]{< < < PREV SERMON < < <}
~
\hyperref[sec:toc]{[返目錄]}
~
\hyperref[sec:493]{> > > NEXT SERMON > > >}


講題:塗抹的手,才能寫出真實 Only the Hand That Erases Can Write True Thing

講員:伍渭文牧師

經文:詩篇32

場合:香港中文大學崇基學院禮拜堂主日崇拜

日期:2020年3月1日

引言:每一個人都會為一些過去的說話和行為,感到懊悔和愧疚.有時夜欄人靜,午夜夢迴,又想起這些事情.為何自己這麼衝動?為何自己這麼愚蠢?甚至為何自己這麼邪惡?你希望從未講過這些話,做過這些事情,你希望可以把這些塗抹.或者此刻,你就活在這些愚蠢和邪惡的行為中,因這些帶給你肉體的歡悅,但之後又感到愧疚.你渴望可以把過去抹掉,改寫自己的生命.若是這樣,今天的信息可以幫助你.畢德生(Eugene Peterson)的譯本《信息》(Message)翻譯詩三十二1譯得很傳神:「你真幸運,你一定開心死了.你可以洗底,從新開始.」(Count yourself lucky, how happy must be-you get a fresh start, your slate’s wiped out )

今日講道題目「塗抹的手,才能寫出真實」,來自中古著名修士愛克哈特大師(Meister Eckhard).塗抹的手不是文過飾非,乃渴想改寫自己的生命.悔改不單洗底,且有新的眼光回眸過去,展望將來.古教父奧古斯丁(A.D. 354-430)悔改後,寫了《懺悔錄》,以嶄新的眼光回眸過去,檢視自己的生命.他說最重要的知識是甚麼?就是「認識自己是一個罪人,且藉悔改認罪,得到上主的赦免.」這也是詩篇三十二篇的信息.奧古斯丁很喜歡這詩篇,他在病榻的牆壁寫上這詩篇,常常誦唸默想.面對肉身的苦痛,死亡的威嚇,不憂不懼.「為此,凡虔誠人都當趁你可尋找時候禱告你,大水泛溢的時候,必不能到他那裡.你是我藏身之處,你必保佑我脫離苦難,以得救的樂歌四面環繞我.」(vv. 6, 7)

大齋期或預苦期第一主日選讀詩篇三十二篇,因為這詩篇點出整個節期的主題:誠實悔改認罪,接受上主的拯救.剛過去的星期三(Ash Wednesday)聖灰日,啟動了為期四十天的大齋期.主禮把去年棕枝主日揮動的棕枝燒為灰燼,塗在領灰的額上說:你來自塵土,也必歸回塵土,要離棄罪惡,跟隨基督.大齋期是基督受難的預備期,我們進入心靈的曠野,效法耶穌在福音經課(太四1-11)四十天祈禱和禁食；並賙濟有需要的人.

預苦期也是準備領洗者,進入學習義理的最後階段.他們將於復活日前一晚(Easter Vigil) 領洗,聖洗禮儀周末晚上開始,直到翌日復活清晨結束,復活日崇拜是聖洗崇拜禮儀的末段.預苦期也是被逐或自我放逐的信徒回歸教會的日子.預苦期以聖灰日悔罪開始,復活日歡騰中結束.因著基督的復活,我們罪得赦免,得救的樂歌四面環繞我們.

詩篇的結構屬交叉配列A,B,B’A’,訓誨與感恩文體相間. A訓誨文體在首尾(vv. 1,2; 9,10), 以第三身論述律法和智慧.B感恩文體在中間,以第一身分享個人經驗(vv. 3-5；6-8).訓誨開始正面點出主旨:誠實認罪必蒙赦免,得著上主的慈愛:「得赦免其過,遮蓋其罪的,這人是有福的,凡心裡沒有詭詐,耶和華不算為有罪的,這人是有福的.」(vv. 1,2) 結束反面提醒:不認罪惹來管教,倚靠上主必有慈愛環繞他.(訓誨者勸戒)「你不可像那無知的騾馬,必用嚼環轡頭勒住它,不然就不能順服.惡人必多受苦楚；惟獨倚靠耶和華的,必有慈愛環繞他.」 (vv. 9,10)

首尾的訓誨,包裹中間第一身闡述的感恩經驗,先是反面不認罪帶來的苦痛:「我閉口不認罪的時候,因終日唉哼而骨頭枯干.黑夜白日,你的手在我身上沉重；我的精液耗盡,加同夏天的干旱.我向你陳明我的罪,不隱瞞我的惡,我要向耶和華承認我的過犯,你就赦免我的罪惡.」(vv. 3,4,5).跟著正面分享認罪後的歡樂,有上主指示當行的路,危難有上主的保護,呼喚人向上主禱告:「為此,凡虔誠人都當趁你可尋找時候禱告你,大水泛溢的時候,必不能到他那裡.你是我藏身之處,你必保佑我脫離苦難,以得救的樂歌四面環繞我.『[上主說:]我要教導你,指示當行的路,我要定睛在你身上勸戒你.』」(vv. 6,7,8)

1 塗抹的手,坦承過犯

1.1認識罪.在認罪之先,讓我們先認識罪.第一,二節用了三個字來描述罪.「得赦免其過,遮蓋其罪的,這人是有福的,凡心裡沒有詭詐,耶和華不算為有罪的,這人是有福的.」(vv. 1,2) 這裡用三個不同的希伯來字描述罪.和合本:其過(NIV transgressions),其罪(眾數sins), 有罪(單數Sin,邪惡 ).罪是過犯,得罪上主；罪是罪行,傷害弟兄；罪是邪惡,扭曲人性.

a其過(Transgression).其過指人踰越界限,不接受本體的限制,所以聖經稱罪人為悖逆之子.以弗所書二2「你們死在過犯之中……隨從今世的風俗,順服空中掌權者的首領,就是現今在悖逆之子心中運行的邪靈.」罪人稱為悖逆之子,過犯(Transgressions)就是反叛上主,撕裂長久忠誠與愛的關系.創世記不是早上創造,晚上墮落.分辨善惡的果子是界限,有一些關系是有界限來保護的,像夫婦的關系,越過界限就是背叛.能分辨善惡表示獨立自主,吃禁果表示他不需要上主,懷疑上主的愛.「神豈是真說……蛇對女人說,你們不一定死.」(創三1, 4)

但人不順服上主,以為從此得著自由了,但以弗所書說人祇是轉向順從撒旦.「隨從今世的風俗,順服空中掌權者的首領.」哥德劇作的主角浮士德就最好的例子.浮士德為了個人的卓越成就,與撒旦交易,出賣自己的靈魂,像國家衛健生委一早跑到武漢,拿到新型冠狀病毒最新資料,趕緊在國際學刊發表,而沒有關心公共衛生,促使有關單位公佈疫情,盡快防疫.有些人趁機提高物價,甚至倒賣救援物資斂財.瘟疫暴露了人性之醜惡.

浮士德與魔鬼交易,出賣自己的靈魂,並沒有得到真正的滿足.他是一個靈魂永不安息的人,永遠在追逐中.從學習哲學轉向法律,再轉向醫學和神學,最後轉向神秘的玄學.浮士德是不斷的追尋者,但同時得不到滿足,似乎掌握著,但卻溜走了的宿命中.撒旦最後會回收取本全和利息的,結果浮士德死得很恐怖.

b其罪(Sins 罪行).第二個字「其罪」指具體的罪行,像十誡的誡命,社會的刑法,罪是偏離正途,傷害鄰舍,甚至最信任自己的朋友.當人離開上主,失去以上主為向心,人彼此恐懼.有上主為向心,人不怕分散全地,但來到無垠無際的示拿平原,他們要聚在一起,建造巴別塔,統一言語,消除異見,因為沒有上主為向心,人懼怕人；該隱因嫉妒殺兄弟亞伯.居莎士比亞名劇「凱撒大帝」著名的對白:「布魯圖,你也參與同謀!」凱撒在羅馬元老院的台階被刺殺,其中致命的一刀,來自他最信任的朋友布魯圖.我認識一位弟兄,很早期在內地做開玩具廠,員工二千多.但最後被基督徒合夥人出賣排斥,經濟一落千丈,後來上主恩待,轉向做文具,蒙上主祝福.也有特殊學校校長悉心栽培的接班人,待自己退休接手後變臉.這跟該隱殺亞伯,猶出賣耶穌是一脈相承的.罪使人彼此相害.不愛上主時,也很難愛弟兄.

C罪(Sin, 邪惡).第三個字「罪」有邪惡,扭曲之意,罪使人自欺.「凡心裡沒有詭詐」(v.2)從上下文的認罪語境,詭詐指自欺,失去誠實的能力.人若長期活在罪中,日子久遠,習慣說慌,連自己也欺騙.

詩篇第一篇是善與惡的比較,詩三十二是悔改與不悔改的比較,兩者開始都用「福哉斯人」(NIV Blessed is the man/ he who).有福的人(義者):「不從惡人的計謀,不站罪人的道路,不坐褻慢人的坐位.」(詩一1)罪像一滑坡,每況愈下.試探的開始,是要解決一具體問題,惡人的確有計謀,協助你走出困境.但像毒品一樣,罪中之樂慢慢會上癮成為習慣.站罪人的道路,習以為常,行在其上.再進一步,成為自己的價值觀,並把扭曲的道理教導別人,坐穩褻慢人的座位了,以前文士是坐著教導門生的.耶穌說:文士和法利賽人坐在摩西的位上,凡他們所吩咐你們的,你們都要謹守遵行.」(太二十三2, 3)

長期活在自欺,表裡不一致,就會影響心理和生理,產生情緒生理(psychosomatic)病.「我閉口不認罪的時候,因終日唉哼而骨頭枯干.黑夜白日,你的手在我身上沉重；我的精液耗盡,加同夏天的干旱.」(vv. 3, 4 ) 靈性的嚴重扭曲,甚至拜撒旦.我以前任教的一間神學院,一位老師是美國的專業心理輔導員,她輔導的對象中,有撒旦教信徒.這些人有些行為非常扭曲:淫邪的性幻想,非常暴力傾向,響往自殺.罪使人變得詭詐,扭曲人性,輕藐神聖的事,甚至褻瀆上主.

總的來說,罪是悖逆上主,傷害他者,扭曲自己.所以在大齋期我們對上主禱告,對自己禁食,對他者賙濟.增強對罪的免疫力,不讓罪蠶食我們的靈魂,扭曲我們的人性,隨從今世的風俗(弗二2),飄離信仰.魯益斯警告我們,絕大部分信徒沒有經過天人交戰,深刻反思後而離開信仰,他們是不知不覺的飄離(drifted away)信仰,這才要命.

要逃避試探,知己知彼,守護自己的罩門.防疫期間,我們帶上口罩,勤洗手,因為口與手容易被病毒入侵的罩門.同樣,每個人都有容易被試探攻陷的罩門.保羅以前的好同工底馬,因為貪愛現今的世界,離保羅而往帖撒羅尼迦去了.(提後四10)帖撒羅尼迦這大都會的魅力,不是浪得虛名.猶大因為貪愛三十兩子舉報耶穌,因為戀上金錢,魔鬼有機可乘.「魔鬼己將賣耶穌的意思放在猶太心裡」(約十三2)；「撒旦入了猶太的心」(路二十二3).掃羅的權力欲使他不能忍受百姓說:「掃羅殺死千千,大衛殺死萬萬.」(撒上十八8),大衛情感豐富,作詩彈琴,敗於感情欲念的牽引.小心罩門,做好防護,逃避試探.

1.2認識悔改,莫待悔無能.「凡虔誠人都當趁你可尋找的時候禱告你.」(v. 6) 趁可尋找的時候是時機,在武漢肺炎的事上,因為隱瞞遲公報,錯過黃金防疫窗口期,令疫情全球擴散,付出昂貴的代價. 同理,悔罪及時,莫待悔無能.不錯,甚麼時候悔改,甚麼時候得到赦免,天父就像浪子故事那位慈父,等候我們歸家.但當我們一次又一次硬心,消滅聖靈的感動時,就失去悔罪的能力,並付上沉重的代價.「你不可像那無知的騾馬,必用嚼環曫頭勒住它.」(v. 9) 當被聖靈觸動,不要消滅這感動,趕快悔改,用塗抹的手,扇動死灰復燃(Glow),扇出火焰(Flame)來.

2塗抹的手,甦醒靈魂

「我要教導你,指示你當行的路；我要定睛在你身上勸戒你.」(詩三十二8) 塗抹的手,撥開霧霾,使靈魂甦醒,帶來屬靈的饑渴,渴慕上主的教導,走在當行的路上.年假最後的一天初四,我帶上口罩,外出往沙田街市買菜,大學火車站祇有幾個人的,大學校園更是空無一人,寂靜得令人不安.因為疫臨城下,災難愈來愈近.這是熟悉的城市,又是陌生的城市,我看到以前未見過的境象,進入閾域(Limbo)的臨界中.塗抹的手,把熟悉的環境抹掉,把慣常閣置,擺脫慣常的制約,思想活潑起來,思考平時不會思考的東西.

2.1 生命脆弱的覺醒.瘟疫使我們認識自己的限制,最有醫學知識的醫護人員,面對沒有藥物治療,同時可以沒有病徵傳播病毒,令人惶恐.之前當權者以強悍的武力「止暴制亂」,覺得很有力量.部分勇武示威者也以滲入暴力回敬,但此刻武力,暴力都沒有用了,我們感到生命如斯脆弱.我們也感到生命可以如斯短暫.若生命如斯短暫,欲望又無限,那甚麼是生命最重要的呢?若此刻上主取去我的生命,我會後悔嗎?有甚麼東西我會後悔沒有完成呢?但是那些東西真是值得我後悔嗎?

2.2走在當行的路上.悔改不單洗底,且有新的眼光,明白我們不是無原無故生下來的.「指示你當行的路」表示上主在我們一生有他的計劃,完成他的計劃,才是真實的人生.「人若自潔,脫離卑賤的事,就必作貴重的器皿,成為聖潔,合乎主用,預備行各樣的善事.」(提後二21) 我們要尋求認識當行的路,行在其中.也許塗抹的手,打翻了安舒的巢窩,讓我們有機會展翅高飛,遨翔萬頃窮碧.也許塗抹的手,解開我們繫在碼頭的繩索,讓我們的船離開港灣,揚帆遠去,見識美麗的海島,物阜民豐的大陸.

2.3當盡明顯的本份.瘟疫是屬於上主隱密的事,為何上主容許發生,我們不知道,因為隱密的事屬於上主,但明顯的事,屬於我們和我們的子孫.(申二十九29).在瘟疫期間,明顯的事是分享自己不多的口罩給沒有的人；到超市購物,留一些必需的糧食,廁紙給其他人,不要掃貨積存.報導說,武漢一家一百人的庭教庭教會,發起全國信徒三天為抗疫禁食禱告,之後每天派出穿防護衣的信徒,對基層信徒和居民派口罩送飯盒,被政府認同.

3 塗抹的手,是上主的手

「你們義人應當依靠耶和華歡喜快樂；你們心裡正直的人都當歡呼.」(vv. 11)

1歡樂的主調.詩篇三十二篇勸勉人誠實認罪,得著赦免.但整卷篇詩的主調不是嚴苛的督責,乃是赦罪的喜樂和福氣.第一節:「得赦免其過,遮蓋其罪的,這人是有福的.」認罪悔改,得蒙赦免,不是就脫離一切的苦難.武漢不少信徒也染上冠狀病毒肺炎,甚至因而離世.但苦難不能成為網羅,綑縛他們,因為上主的慈愛像保護罩環繞著他們.「你必保佑我脫離苦難,以得救的樂歌,四面環繞我.」(v.7)「依靠耶和華的,必有慈愛四面環繞他.」(v. 10)

2生與死的自由.瘟疫使我們醒悟人生無常,不少人甚至相信生與死,是不能逆轉的宿命.傳道書三1,2豈不是說:「凡物都有定期,天下萬物都有定時,生有時,死有時.」我們沒有被諮詢過就生下來,也沒有得自己的同意就要離開這世界.生與死是不能逃避,也無能為力的宿命.

今日書信經課羅五12-21,說明人有兩種生命的狀態,在亞當裡和在基督裡；在肉體裡和在聖靈裡；在亞當裡是日光之下的創造論,在基督裡是日光之上的救贖論.日光之下的人生是虛空,人生像是不能逆轉的宿命.但從救贖論看,認罪悔改得著新生命,在基督裡,生死不再是無能為力的宿命；生與死乃生命的自由.我們可以有自由決定自己屬靈生命的誕生,也可以天天向罪死,使屬靈的生命茁壯成長.無論得時不得時,我們都有這自由.

多帕多迦三教父之尼森主教聖國瑞(Gregory of Naziansus)在《訓道篇[傳道書]釋義》說:「當一個人,由於敬畏上主而懷孕,經過心靈的產痛而生產自己的得救,這是適時的生產……我們可以說,我們就是自己的父親.……但是,如保羅所說,如果『基督的形體』沒有在我們內形成,我們便使我們自己成為流產兒,不完善和患病的人.……保羅又說:『主,為了你,我們天天受死.』……這就是適時的死亡,這死亡產生真生命.」(【每日頌禱】一冊: 502, 506)

在基督裡的自由,世界上的橫逆都不能奪去,「大水泛溢的時候,必不能到他那裡.」(v.6).詩篇第一篇提到義人「像一棵樹栽在溪水旁,按時候結果子.」(詩一3)按甚麼時候?按他自己的時候:無論甚麼時候,瘟疫的時候,平安的時候,無論得時不得時,都可以為上主結果子,就像剛才所說的武漢教會.

總結:塗抹的手,是我們的手,坦誠悔改,必蒙赦免.塗抹的手,也是上主的手,他抹掉我們的過去,可以洗底,從新開始.他以慈愛環繞我們,賜我們平安,瘟疫也不能奪去的平安.他藉聖經指引我們當行的路,成就他在我們生命中的計劃.



\newpage

\section{本於信以至於信}
\label{sec:493}
\textbf{劉國偉先生 - 本於信,以至於信.}
\newline
\newline
\hyperref[sec:492]{< < < PREV SERMON < < <}
~
\hyperref[sec:toc]{[返目錄]}
~
\hyperref[sec:494]{> > > NEXT SERMON > > >}


詩篇 121:1-8
\newline
\begin{longtable}{cl}
\hline
\hline
章節 & 經文 (和合本修訂版)\\
\hline
121:1 & \begin{tabularx}{0.7\textwidth}{X} 我要向山舉目，我的幫助從何而來？ \end{tabularx} \\ \\ \relax
121:2 & \begin{tabularx}{0.7\textwidth}{X} 我的幫助從造天地的耶和華而來。 \end{tabularx} \\ \\ \relax
121:3 & \begin{tabularx}{0.7\textwidth}{X} 他不叫你的腳搖動，保護你的必不打盹！ \end{tabularx} \\ \\ \relax
121:4 & \begin{tabularx}{0.7\textwidth}{X} 保護以色列的必不打盹，也不睡覺。 \end{tabularx} \\ \\ \relax
121:5 & \begin{tabularx}{0.7\textwidth}{X} 保護你的是耶和華，耶和華在你右邊蔭庇你。 \end{tabularx} \\ \\ \relax
121:6 & \begin{tabularx}{0.7\textwidth}{X} 白日，太陽必不傷你；夜間，月亮也不害你。 \end{tabularx} \\ \\ \relax
121:7 & \begin{tabularx}{0.7\textwidth}{X} 耶和華要保護你，免受一切的災害，他要保護你的性命。 \end{tabularx} \\ \\ \relax
121:8 & \begin{tabularx}{0.7\textwidth}{X} 你出你入，耶和華要保護你，從今時直到永遠。 \end{tabularx} \\ \\
[1ex]
\hline
\hline
\end{longtable}
講題:本於信,以至於信. By faith from first to last.

講員:劉國偉先生

經文:詩篇121:1-8

場合:香港中文大學崇基學院禮拜堂主日崇拜

日期:2020年3月8日

引子

馬丁路德(1483-1546),16世紀的宗教改革家,在1520年出版的他的著作〈基督徒的自由〉(On the Freedom of a Christian)中如此寫道:「許多人以為基督徒的信,乃是一件容易之事,他們其所以如此行,是因他們沒有信的經驗,也從來沒有嘗到信的大能.因為除非一個人在某一個時候,處於試煉的壓迫之下,嘗過信所給予的勇氣,他就無法將信好好描寫出來,也無法對凡論信所正確描寫的加以領會.」很明顯的,馬丁路德以「因信稱義」作為教義的核心,他更進一步把這核心教義,引伸至信仰人生的不同層面.

如此說來,「因信稱義」不單是一個人怎樣得救的課題,這個課題當然非常重要,然而,「因信稱義」其實涵蓋整個信仰人生,是基督徒在世生活的重要基石.

我們將鏡頭轉向第一世紀.在羅馬書第1章16至17節,我們看見使徒保羅如此說道:「我不以福音為恥,這福音本是神的大能,要救一切相信的,先是猶太人,後是希利尼人.因為神的義,正在這福音上顯明出來.這義本於信,以至於信,如經上所記:「義人必因信得生.」」

我們可以說這兩節經文,是羅馬書的提綱挈領的經文,道出了保羅在羅馬書接下來要討論的神學課題的核心.「本於信,以至於信.」是今天宣道的講題,正是取材於此經節.使徒保羅在此經節所表述的教導,是要告訴大家,福音顯明了上帝怎樣使人稱義:由信心開始,用信心結束.新國際譯本NIV如此翻譯:a righteousness that is by faith from first to last；這信心是對所信的上帝的委身,是生命的投靠,也是理性和意志的接納；所以,一煞那的感動,片刻的立志,可能只是建立真正信心的觸發點,與使徒保羅及馬丁路德所談及的,信心有如道路,開了個頭,還是要真心真意走下去,有明顯的落差.

我的幫助從何而來

我們今天所選讀的詩篇121篇,是一首《上行詩》.由120篇到134篇結集一起的上行詩,明顯在以色列人的敬拜生活中,有它的實際用途.到底上行詩如何具體應用,至今仍未有一致的看法,但不論是那些登上聖殿山,去敬拜神的朝聖者,在旅途上的詩集,又或者是祭司在耶路撒冷聖殿,登上十五級階梯時所誦唱；也有可能兩者皆是,所以《上行詩》又稱《登階詩》.那些到耶路撒冷朝聖和到聖殿敬拜的人,其實與我和你一樣,從衣食住行,柴米油鹽,閒言閒語,是非曲直的現實生活之中,走進朝見神聖上主的旅途上.

詩篇121篇的作者,將生活中的一個貼地的問題,帶進這個向山舉目的朝聖之旅:我的幫助從何而來?用廣東人的慣常講法,可以轉換為「邊個係我靠山?」回應的答案來得直接清晰:我的幫助從造天地的耶和華而來.詩人對上主的這一份信心,乃建基於對祂的認識；從第3至8節,詩人總共採用「保護」這個詞彙6次,反覆誦念他如何相信上主的保護.

詩人告訴大家,我們可以在登山的旅途中,步行穩定,或者站立得穩,因為保護我們的上主,不會像巴力那樣或是睡覺,或是遠行(王上18:27),上主必不打盹,也不睡覺,祂不會疲倦,也決不置身事外；(v.3-4)

這位上主叫白日熾熱的太陽,夜間那象徵不可知的幽暗力量的月亮,都不會傷害信靠祂的人,祂就是那位在我們右邊蔭庇我們,叫我們即使身陷戰陣之中,也得到祂特別保護的主；(v.5-6)

在這篇詩結束的時候,詩人告訴大家,即使災害會臨到我們,也許我們每天出入會面對危險,但他確信大家會得到上主的保護,而且是終極的永遠的保護.(v.7-8)

信靠就是順服

今天我們誦讀的舊約經課,是有關「信心之父」亞伯拉罕的一個小片段.簡短的幾節經文,展現在亞伯蘭面前的,是一次人生的冒險之旅.「離開本地,本族,父家」幾乎是一次徹底的從零開始,面對一個陌生的未來,當然伴隨不可預知的風險.

經文以最簡潔的方式,向我們表達了亞伯蘭對上主的信靠:「亞伯蘭就照著耶和華的吩咐去了.」(創12:4a)沒有爭辯,沒有抗拒,就順服前行.一個人能夠對上主的呼召與差遣,如此順服,其實是他對上主絕對的信靠,那是一份密切的人神關係.真實的信靠就帶來真誠的順服,原來人對神的信心,不在乎理性的辨識,也不在於意志的選擇,而是生命關係,深厚的人神關係,就有牢固的信心,就帶來真誠的順服.

叫一切信祂的……

至於今天福音經課提及的尼哥德慕,使徒約翰告訴我們,他是一個法利賽人,而且是猶太人的官(約3:1),按約翰福音原文用字archon,其實就是猶太公會的一員；耶穌也確認他的教師的身份(v.10).簡要地說,尼哥德慕是猶太社會的精英,有地位,權力,學問,而且他有渴慕真理的心.但是,「人若不重生,就不能見上帝的國.」(v.3)這等艱澀的道理和信仰,又怎可能倚靠一個人的見識,學養,或者權力和地位去掌握和內化?

耶穌對尼哥德慕說:「上帝愛世人,甚至將祂的獨生子賜給他們,叫一切信他的,不致滅亡,反得永生.」(v.16)這節耳熟能詳的金句,再一次提醒我們,我們來到神面,是憑著那單純的信心,完全不是因為我們有什麼可以自恃和誇口的,而且這條憑信心踏上的救恩之路,也就是基督為我們捨身,而成為那條通往上帝那裡的道路,也完全是因為上帝的愛,叫我們可以得到永生.

弟兄姊妹,我們都是得贖之民,我們除了記掛死後的永生,又要如何面對今生的使命?我就想起馬丁路德的話:「許多人以為基督徒的信,乃是一件容易之事,他們其所以如此行,是因他們沒有信的經驗,也從來沒有嘗到信的大能.因為除非一個人在某一個時候,處於試煉的壓迫之下,嘗過信所給予的勇氣,他就無法將信好好描寫出來,也無法對凡論信所正確描寫的加以領會.」但願當我們身處風暴之中的時候,我們發現在我們心裡的,是亞伯拉罕那份對上主完全順服的信心,而不是地位,權力,學問充斥著我們的心,並且視之為我們的靠山.

本於信,以至於信……

使徒保羅在羅馬書4章3節引用創世記:「亞伯拉罕信上帝,這就算為他的義.」(創15:6)再次演繹「因信稱義」的核心,是在於恩典:「所以人得為後嗣是本乎信,因此就屬乎恩,叫應許定然歸給一切後裔；不但歸給那屬乎律法的,也歸給那效法亞伯拉罕之信的.」(羅4:16)

人不能自誇,不單叫我們在主面前謙卑,信靠,其實也提醒我們,當面臨人生的考驗與難關時,認知自己的有限,就更應當憑著信,仰望那位死而復活的主,因為我們憑信心而活.

因著信,我們不以福音為恥(羅1:16),並且好像使徒保羅一樣,情願盡上我們的力量,將福音傳開.(羅1:15)也因著信,我由信心開始,走在信心的道路上,成就一個信靠,順服上帝的人生旅程.

弟兄姊妹,我們可以怎麼做?詩篇的作者啟導我們,去細味和記下上帝每一次的幫助.回想因為上帝的保護,你不曾搖動的腳,渡過了白日與黑夜的傷害,跨過了每天出出入入的危害和風險.

又或者,我們可以參考馬丁路德的觀點:「…不能坐享安逸…應該藉著禁食,儆醒,勤勞,以及各種合理的訓練來留意訓練身體,好叫身體順服內心的人與信.」

「從信就流出在主裡面的愛與喜樂來,從愛又流出快樂,甘願,與自由的心來,樂意服事鄰舍,不計感恩與忘恩,毀與譽,得與失.」

By faith from first to last!

願榮耀歸給三一真神!



\newpage

\section{上帝是不是在我們中間呢}
\label{sec:494}
\textbf{鄧瑞強博士 - 上帝是不是在我們中間呢?}
\newline
\newline
\hyperref[sec:493]{< < < PREV SERMON < < <}
~
\hyperref[sec:toc]{[返目錄]}
~
\hyperref[sec:495]{> > > NEXT SERMON > > >}


出埃及記 17:1-7
\newline
\begin{longtable}{cl}
\hline
\hline
章節 & 經文 (和合本修訂版)\\
\hline
17:1 & \begin{tabularx}{0.7\textwidth}{X} 以色列全會眾遵照耶和華的吩咐，從汛的曠野一段一段地往前行。他們在利非訂安營，但百姓沒有水喝。 \end{tabularx} \\ \\ \relax
17:2 & \begin{tabularx}{0.7\textwidth}{X} 百姓就與摩西爭鬧，說：「給我們水喝吧！」摩西對他們說：「你們為甚麼與我爭鬧呢？你們為甚麼試探耶和華呢？」 \end{tabularx} \\ \\ \relax
17:3 & \begin{tabularx}{0.7\textwidth}{X} 百姓在那裡口渴要喝水，就向摩西發怨言，說：「你為甚麼把我們從埃及領出來，使我們和我們的兒女，以及牲畜都渴死呢？」 \end{tabularx} \\ \\ \relax
17:4 & \begin{tabularx}{0.7\textwidth}{X} 摩西就呼求耶和華說：「我要怎樣對待這百姓呢？他們差一點就要拿石頭打死我了。」 \end{tabularx} \\ \\ \relax
17:5 & \begin{tabularx}{0.7\textwidth}{X} 耶和華對摩西說：「你帶著以色列的幾個長老，走在百姓前面，手裡拿著你先前擊打尼羅河的杖，去吧！ \end{tabularx} \\ \\ \relax
17:6 & \begin{tabularx}{0.7\textwidth}{X} 看哪，我要在何烈的磐石那裡，站在你面前。你要擊打磐石，水就會從磐石流出來，給百姓喝。」摩西就在以色列的長老眼前這樣做了。 \end{tabularx} \\ \\ \relax
17:7 & \begin{tabularx}{0.7\textwidth}{X} 他給那地方起名叫瑪撒，又叫米利巴，因為以色列人在那裡爭鬧，並且試探耶和華，說：「耶和華是否在我們中間呢？」 \end{tabularx} \\ \\
[1ex]
\hline
\hline
\end{longtable}
講題:上帝是不是在我們中間呢? Is the Lord among us or not?

經文:出埃及記17:1-7

講員:鄧瑞強博士

場合:香港中文大學崇基學院禮拜堂主日崇拜

日期:2020年3月15日

各位弟兄姊妹,早安,祝願各人平安.疫情仍在,祈求天父繼續保守著我們.

2020年奧斯卡最佳電影的得獎作品是南韓的《上流寄生族》(Parasite).內容大致上說:一個家庭,因著種種社會原因,活下來,就是一個活在社會底層的貧窮家庭.從畫面看,他們真是活在底層的.原本,他們活在街道地面之下的地下室；後來,他們與另一個貧窮家庭爭奪一富有人家隱閉的地下室.最殘酷的爭奪,發生在貧窮人之間,他們爭奪一個難得的寄生位置.這個貧窮家庭,父母兄妹,其實各自有種種謀生技能,只要給他們機會,他們是發揮得很好的.但社會制度就是劃下了一道無形界線,一道很難越過的界線,一道窮人與富人之間的界線.他們無法越過這界線,於是,他們始終就只能是社會裡最底層的人,靠寄生於富有人家而活下去.那位爸爸講了一句令人深思的說話,他說:「如果我有錢,我都是個好人」.後來,這位爸爸因為抵受不了那有錢人嫌惡他身上的氣味,一種窮人洗不掉的寒酸味,他殺了那個有錢人.之後,他只能隱藏在那富有人家家裡一個不為人知的地下室裡,更無法擺脫他的存在的特有氣味,只能成為更深依附上流人家的寄生蟲.他的兒子,知道他藏在那地下室裡.這兒子的渴望,就是成為有錢人,買下這間豪宅.我們知道,這兒子的夢,是不會實現的.

就著這劇情,我們思考三個問題:

一,一種依附他人的生存狀態,是否不能打破?

二,是否真的要有錢,才能做好人?面對貧窮所帶來的逆境,是否真的不能有道德上上進的可能?是否一定要回到寄生於他人的生存狀態?

三,對整個劇情的一個形而上學的追問,一個終極之問:這種殘缺的生存狀態,是只能永恆地輪迴下去的嗎?寄生於他人…努力活出自我…寄生於他人…又再努力活出自我…結果仍是要寄生於他人,這是無法擺脫的命運嗎?

請大家記著這三個問題,我們等一會會回答它們.

我們看一個聖經故事,一班寄生上流,寄生他人的人尋索生命自由的故事.這故事就是今日崇拜讀經裡的舊約經課:出埃及記17:1－7

出17:1　以色列全會眾都遵耶和華的吩咐,按著站口從汛的曠野往前行,在利非訂安營.百姓沒有水喝,

出17:2　所以與摩西爭鬧,說:「給我們水喝吧!」摩西對他們說:「你們為甚麼與我爭鬧?為甚麼試探耶和華呢?」

出17:3　百姓在那裡甚渴,要喝水,就向摩西發怨言,說:「你為甚麼將我們從埃及領出來,使我們和我們的兒女並牲畜都渴死呢?」

出17:4　摩西就呼求耶和華說:「我向這百姓怎樣行呢?他們幾乎要拿石頭打死我.」

出17:5　耶和華對摩西說:「你手裡拿著你先前擊打河水的杖,帶領以色列的幾個長老,從百姓面前走過去.

出17:6　我必在何烈的磐石那裡,站在你面前.你要擊打磐石,從磐石裡必有水流出來,使百姓可以喝.」摩西就在以色列的長老眼前這樣行了.

出17:7　他給那地方起名叫瑪撒〔就是試探的意思〕,又叫米利巴〔就是爭鬧的意思〕；因以色列人爭鬧,又因他們試探耶和華,說:「耶和華是在我們中間不是?」

(1)一種依附他人的生存狀態,是否不能打破?

出17:1　以色列全會眾都遵耶和華的吩咐,按著站口從汛的曠野往前行,在利非訂安營.百姓沒有水喝.

這句經文,一個小故事的開場白,有一個歷史背景.《出埃及記》第14章,記載以色列人離開埃及,過紅海.個多月後,便來到這節經文講的「汛的曠野」,平沙漠漠,大地蒼茫,然後,他們在「利非訂」這個地方安營.這個小故事,背後是一個大故事.原來,是一群寄生於他人的奴隸的出走故事.再講遠一點,這群人為何在埃及做了別人的奴隸?原來,是他們的祖先無飯吃,他們到了埃及,寄居別人的土地,在那裡活了下來.後來,他們越來越多人,威脅到主人家的安全,主人家便將他們打壓成奴隸.他們寄生在別人家裡,成了最底層的人.

回到我們剛才問的三個問題,第一個問題是:一種依附他人的生存狀態,是否不能打破?

在《上流寄生族》那齣戲裡,我們看到那家庭努力打破寄生一族的宿命.哥哥努力親近富有人家的千金小姐,他常常想:是否真的不能融入富有階級的生活裡?那位爸爸,嘗試用其他洗衣的方法,去除那窮人的寒酸味,去除那低下階層的特有標籤.他們是努力的,但在電影裡,他們的努力並不成功.

聖經的故事,是一救贖故事.那群在埃及裡作奴隸的人,他們是能打破奴隸的宿命的.其實,奴隸本身是很難走出奴隸的宿命的,正如一個患上「強迫症」(Obsessive-Compulsive Disorder)的人是很難自己走出強迫行為的宿命的.早一陣子,有一電視節目探討「強迫症」.你很難叫一個洗手強迫症的人停止洗手的.你對他說,你是自由的,你有自由不去洗手的.他其實是做不到的,當他停止洗手時,他不是覺得自由,而是覺得污穢,他只能透過洗手,去清除那污穢.

奴隸的結構是很複雜的.多層次的奴隸結構,令一個奴隸逃脫不出奴隸的命運.你可能說:「我不是奴隸,我有自由買這樣東西,買那樣東西,我可以自由地享受這些東西給我的快樂.」我們真的如此自由嗎?當購物「淘寶化」的時候,你其實不是自由地選這個牌子,買那件東西的,你會看別人購買後打的分數,你到頭來不是自由地選那件東西,而是別人替你選的.當然,你可能知道,這些所謂「別人的意見」也可以是製造出來的.你以為自由地買一件東西,你卻不知有多少力量已操控著你這行為.再深一層,你以為能購物便能帶給你快樂,買到那層見到海豚跳上水的樓盤便算是幸福人,這種想法本身,便已是消費社會塑造出來的消費人格.我們可稱之為「意識形態」的枷鎖,它框住我們的思想範圍及理想方向.我們買甚麼?並不那麼自由.我們為何要買?也並不那麼自由.奴隸的枷鎖,層層疊疊,奴隸自己是很難掙脫的.

《上流寄生族》的那個家庭,會同意說,我們掙不脫這奴隸枷鎖.聖經卻言說「出埃及」的故事.奴隸自己的確掙不脫這枷鎖,但神能夠.神差派他的使者,行一些前人未見的事,指導他們走向一前所未見之地.這一片新天新地,是奴隸未曾有過的渴望.如今,這群奴隸向著那新天新地前進.

(2)是否一定要回到寄生於他人的生存狀態?

出17:1　以色列全會眾都遵耶和華的吩咐,按著站口從汛的曠野往前行,在利非訂安營.百姓沒有水喝,

出17:2　所以與摩西爭鬧,說:「給我們水喝吧!」摩西對他們說:「你們為甚麼與我爭鬧?為甚麼試探耶和華呢?」

出17:3　百姓在那裡甚渴,要喝水,就向摩西發怨言,說:「你為甚麼將我們從埃及領出來,使我們和我們的兒女並牲畜都渴死呢?」

整幅圖畫,有神在看顧著的,是神帶領他們去到「利非訂」的.他們沒有水渴,這是神知道的.我們沒有口罩,這也是神知道的.在人生的旅途上,遇到困難,是常態,也是神考驗我們的時刻.

在《上流寄生族》裡,那位爸爸說:「如果我有錢,我都是個好人.」言下之意,如果無錢,偷呃拐騙也在所不惜了.事實上,他們在戲中,是用奸計逼走原先的司機及女管家,而令自己上位的.沒有錢,談不上道德的上進.到了片尾,他們走向金錢路的最大能阻又出現了,就是那窮人抹不走的氣味,那是底層人的存在的永恆標籤,那是貧富之間闖不過的界線.遇著這難阻,那爸爸便沒有甚麼道德可言了,他殺了那有錢人,於是,他又回到地下室,回到寄生蟲的生命.

在《出埃及記》這個故事裡,群眾在奔往自由的旅程中,遇到他們無法解決的難題:他們沒有水渴.這是生死攸關的問題.他們向摩西發怨言,說:「你為甚麼將我們從埃及領出來,使我們和我們的兒女並牲畜都渴死呢?」言下之意,做一個奴隸豈不比現在的光景更好.事實上,這群人是試圖返回埃及的.

遇到困難,人有甘做奴隸的傾向.社會心理學家佛洛姆(Erich Fromm)寫了一本書,書名叫《逃避自由》(Escape from Freedom).書中講到,人獨立而獲得自由後,他會發現,獨立意味著孤獨,自由意味著要自己面對問題,自己要背負運用自由後的代價.這責任太重了,這後果太大了,人在孤單與無力的情況,唯有靠向提供保護的權威.有奶便是娘,給我飲食的便是老闆,給我安全感的便是主人.於是,越自由的地方,在危機下,越容易產生極權.

這群以色列人,在無水的情況下,他們眷戀被充沛的尼羅河水澆灌的豐饒埃及.他們不做自己面對困難的自由人,他們寧做寄生於埃及的奴隸.做一隻被有錢人養的狗,好過做一個面對困苦的自由人.

(3)一個形而上學的追問,一個終極之問:寄生於他人,是無法擺脫的命運嗎?

面對人生的悲劇,一個形而上學的追問,一個終極之問:寄生於有權有錢的人,是窮人無法擺脫的命運嗎?

在《上流寄生族》裡,這的確是一無法擺脫的命運.住在地下室的人,無論如何向上爬,終歸又返回地下室.那個兒子想發達,想買起那間屋,好讓藏在屋裡地下室的爸爸能走出來.但畢竟那不是現實,那只是夢,且是一個不可能的夢.這是無法走出來的輪迴,這是無法改變的命定,這是沒有出路的地下室.

出17:4　摩西就呼求耶和華說:「我向這百姓怎樣行呢?他們幾乎要拿石頭打死我.」

出17:5　耶和華對摩西說:「你手裡拿著你先前擊打河水的杖,帶領以色列的幾個長老,從百姓面前走過去.

出17:6　我必在何烈的磐石那裡,站在你面前.你要擊打磐石,從磐石裡必有水流出來,使百姓可以喝.」摩西就在以色列的長老眼前這樣行了.

出17:7　他給那地方起名叫瑪撒〔就是試探的意思〕,又叫米利巴〔就是爭鬧的意思〕；因以色列人爭鬧,又因他們試探耶和華,說:「耶和華是在我們中間不是?」

聖經這個小故事,其終極之問是:「耶和華是在我們中間不是?」在這種看來死路一條的生活裡,到底有神沒有?若有,他在不在?這真是苦不堪言的人的天問.而這故事,就是回答這終極之問,神在我們中間,神在絕望的人中間.

聖經這個小故事,是要開我們的眼睛,讓我們看到我們原先看不到的出路,讓我們看到我們原先看不到的神.那群以色列人,面對艱難時,他們寧願回去做奴隸,他們認命,他們要回到地下室去.當我們看不到出路時,我們的抉擇,也可能是這樣.聖經的故事,是要開我們的眼睛,讓我們看到,我們面對的,不是無法擺脫的命運,不是永恆的奴隸輪迴,而是慈愛的上帝.當我們無水喝的時候,這位上帝賜我們水喝.不是帶我們到水泉那裡去,而是在根本無可能有水的磐石那裡,叫水流出來.我們眼見的,就只是乾地,只是絕境,只是硬石.但神希望我們看到,這不是甚麼乾地,這是神帶領我們去的地方.這不是甚麼絕境,這是我們難得的,經驗神的慈愛的地方.這裡有的,更不是甚麼硬石,這是活水的源頭.神呼喚我們,不要回到地下室,不要再做奴隸,不要再做寄生蟲.要向前行,要望向應許之地,要望到那顧念蒼生的主.

從現實的困局,我們觸碰到人性內的奴性邏輯.奴性邏輯的背後,我們追問終極的形而上學之問.到底,這是否只是一地下室,抑或,神就在我們中間?

我希望大家看到:神真的是在我們中間.

但願:榮耀歸於聖父,聖子,聖靈.阿們.

(附記:本人在崇基禮拜堂的講章已結集成《荷花池畔春草綠》講道集,試讀本可參閱:\url{http://cc.cuhk.edu.hk/dev/images/content/fund-raising-campaigns/university-sermons-by-dr-sk-tang/CCC-SK-Tang-Sermon-Text-Sample-2.pdf}.較早期的崇基講章,也已結集成《敢於跟隨主》一書.有興趣的朋友可到各大基督教書局購買這兩書.)



\newpage

\section{失明的才能看見}
\label{sec:495}
\textbf{關瑞文教授 - 失明的才能看見}
\newline
\newline
\hyperref[sec:494]{< < < PREV SERMON < < <}
~
\hyperref[sec:toc]{[返目錄]}
~
\hyperref[sec:496]{> > > NEXT SERMON > > >}


講題:失明的才能看見 Only the Blind can See

講員:關瑞文教授

經文:約翰福音9

場合:香港中文大學崇基學院禮拜堂主日崇拜

日期:2020年3月22日

弟兄姊妹,早晨.今天三代經課,讀到約翰福音第九章,對我們現在處身的國際和本地局勢,實在太有啟迪了.我們要以充滿感恩的心去讀.

如何形容目前的國際和本地局勢呢?如果我用四個字來形容,相信會有不少人共鳴.這四個字,就是「嚴重撕裂」.在國際方面,中美之間的摩擦,一波未平一波又起.

這邊廂要發中國夢,那邊廂要 Keep America Great.先有已經持續一段時間的中美貿戰.最近還因為鬥嘴到底武漢病毒源頭是武漢還是美國,弄得美國國務卿龐培歐與中國中央政治局委員兼外事工作委員會辦公室主任楊潔篪,在電話中罵戰.近月來,解放軍戰機在台海中線附近多次實施夜航,而美國海軍則剛剛於早幾天以「羅斯福」號為主的航空母艦在南海海域展開聯合訓練作為挑釁還擊.大家似乎都要透過挑起對外矛盾而增強自己在內部的認受性和團結.還有,就是俄羅斯,沙特以致美國之間的石油大戰,弄得環球股市大跌,在武漢病毒已經出現「全球大流行」的同時,雪上加霜.國際間的撕裂已經到了分分鐘擦槍走火的地步.香港本地的撕裂,也不遑多讓-紅黃藍白黑!

在因各種政治等等因素而至撕裂紛爭橫行的世代裡,教會應該怎樣面對,信徒應該如何認清上帝心意?約翰福音第九章給了我們很重要的教導.雖然約翰福音第九章沒有直接讓我們知道上帝在今時今日的心意,但它卻提醒我們,如何不會在紛亂中,跌入各種不同的陷阱.而這些陷阱,經常使人即使當上帝的心意是一清二楚時,都總是視而不見的.我們將要從約翰福音第九章看見一個悖論-就是「失明的才能看見」上主的心意.

請翻開約翰福音第九章.

當天發生的事,精彩逼人.話說於某安息日,耶穌與門徒在街上巧遇一個生來瞎眼的乞丐,門徒問耶穌是誰犯罪以至瞎子生來受苦.這一問顯然是發自一種道德倫理推論:(1)前設-生來受苦乃源自罪；(2)事實-瞎子生來受苦；(3)結論-一定有人犯罪；問題只是誰犯罪?耶穌不以門徒所問為對,略略解釋後就使瞎子重見光明.路人甲乙丙見之希奇,尋根究底卻不知醫者耶穌何許人,便求問宗教權威法利賽人.法利賽人則不問情由,攻擊耶穌,欲加之罪.判詞是「這個人不是從上帝來的,因為他不守安息日.」不過卻又有人反問:「一個罪人怎能行這樣的神蹟呢?」那他們就唯有一口咬定,這個所謂瞎子根本從來都不是瞎眼,他只是個騙子,與耶穌同夥騙人,意思是要指控耶穌所謂行神蹟,全都是騙局.於是他們找來瞎子的父母質詢.父母因怕得罪法利賽人而被趕出會堂,便支支吾吾,推卸責任.法利賽人只好再召見瞎子,嚴加拷問,要逼瞎子承認耶穌是罪人.然而,瞎子真的「唔識死」,因不能否認鐵一般的被醫治經驗,大膽頂撞法利賽人,說:「他〔耶穌〕是個罪人不是,我不知道；有一件事我知道,從前我是眼瞎的,如今能看見了.」然後繼續為耶穌辯護.結果,法利賽人也判定瞎子是罪人-v34他們回答說:「你全然生在罪孽中,還要教訓我們嗎?」於是把他趕出去了.後來,被趕走的瞎子在路上巧遇耶穌,經過一番的對話,他終於看見,原來醫好他的耶穌,是主,就拜他.瞎子最初以為自己看見的耶穌只是一個人 (v11),後來在 v17說耶穌是先知,到 v33進一步發覺耶穌是「從上帝來的」,v35-36發現耶穌是「上帝的兒子」,最後認出耶穌是主.他由瞎眼,後來慢慢迷迷糊糊看見,最後清清楚楚看到.耶穌則提出一個悖論作結語:「我為審判到這世上來,叫不能看見的,可以看見；能看見的,反瞎了眼.」

「失明的才能看見」上主!

弟兄姊妹,你是否有留意到,為何「不能看見的,可以看見；能看見的,反瞎了眼」呢?如果要深入了解,我們必須要知道約翰福音成書的大背景,以致約翰福音聚焦所關注的問題.

其實,約翰福音是為信徒身處在撕裂的大環境而寫的.所謂撕裂,是隨著當時政治社會等等的撕裂,為猶太教內部帶來了莫大的撕裂,以致當時的基督徒生活在惶恐與疑慮當中.

大部分聖經學者都認為,約翰福音是在公元70 年後寫成.很大可能是寫於80-95年間.公元70年,是猶太教的重大歷史轉折點.這些年間,正是猶太人反羅馬殖民戰爭進行得如火如荼之際,歷史稱為第一次猶太大起義；同時也是羅馬帝國內部權鬥,帝皇寶座爭奪戰白熱化的時候.

公元 69 年,維斯帕先成功奪權稱帝,成為羅馬弗拉維王朝的開朝皇帝.公元70年,維斯帕先的兒子揮軍殺入耶路撒冷,大致平定了猶太起義,並且焚毀了聖殿,大批猶太人被俘為奴.

聖殿被毀之後,猶太教受到重創.之前的猶太教有四大主流派別:法利賽人,撒都該人,艾賽尼派,奮銳黨.而基督徒,在當時也是猶太教的一個新興派別.聖殿被毀之後,撒都該人,艾賽尼派,奮銳黨都消失了, 就只剩下法利賽人和基督徒.約翰福音成書的時候,正是法利賽人和基督徒在爭奪猶太教正統話語權的時候.當然,基督徒勢孤力弱,而法利賽人就與羅馬新政權有千絲萬縷的關係,甚至政教勾結來管治猶大地,因此基督徒群體是節節敗退的.當時,法利賽人已經決定,要把所有認耶穌是基督的人,趕出猶太人的會堂.對猶太教而言,被趕出會堂,就等如被宣判死刑,因為會堂是猶太教的禮拜與舉行贖罪儀式的地方.被趕出會堂,就是被逐出教,與耶和華的救恩無份

被趕出會堂,希臘文是 aposynagōgos. 約翰福音,就是環繞著基督徒因aposynagōgos而產生的惶恐與疑慮而寫的,要讓他們看見福音的真相而安心.所以我們會在約翰福音9:22; 12:42; 16:2多次讀到aposynagōgos,而第九章就很完整地表達了aposynagōgos ,「被趕出會堂」的處境.

在這個社會以致教內撕裂裡面,何解會原本「不能看見的,可以看見；能看見的,反瞎了眼」?為什麼「失明的才能看見」上主的心意呢?

第九章裡,除了耶穌之外,有五批人:1. 法利賽人,2. 門徒,3. 路人甲乙丙,4. 瞎子的父母,5. 瞎子.前四者是「能看見的,反瞎了眼」,只有瞎子是失明的反而看見.為甚麼呢?

第一批,法利賽人.法利賽人是「得勢當權派」.他們的理想是在危機年代裡保全自己的宗教命脈,甚至委曲求存,與焚毀聖殿的羅馬政權合作.運用自己所得的權力和話語權,不問真相,不問情由,總之要把帶來不穩定的基督徒清除,如 v22 說,他們「商議定了,若有認耶穌是基督的,要把他趕出會堂.」法利賽人「能看見的,反瞎了眼」,因為他們太接近權力了,也從權力得到太多好處了.

第二批,第九章事件中的門徒.他們是被法利賽人欺壓的弱勢無權份子.為什麼連耶穌的門徒在這事情上都是「能看見的,反瞎了眼」,看不懂在他們面前所發生的事呢?從他們向耶穌的發問,可以得知.v1-2耶穌過去的時候,看見一個人生來是瞎眼的.門徒問耶穌說:「拉比,這人生來是瞎眼的,是誰犯了罪?是這人呢?是他父母呢?」他們把受苦與犯罪拉上了因果關係.受苦的人,就是因為道德有問題,因干犯上帝的義怒被懲罰.其實,受苦與犯罪是否有直接因果關係,在當日的拉比傳統裡,其實沒有共識的.但門徒似乎想當然地要站在道德高地理解人世間的曲折,也期望耶穌對這生來瞎眼的人作出道德裁判.結果,耶穌一盤冷水倒在他們的頭上,告訴他們事情殊不是如他們所想象的.門徒是被法利賽人欺壓的弱勢無權份子,甚至被界定為異端.被人欺壓以致成為弱勢無權的一方.作為面對高牆的雞蛋,很易會誤墮「站在道德高地」的陷阱,被逼「窮得只剩下道德」的身份,成為「道德高牆派」.可是,這也一樣會讓人「能看見的,反瞎了眼」,看不懂上帝的心意,就如當日的門徒一樣.

第三批,路人甲乙丙.他們是「盲從權力派」,是愚蠢的群眾.v13是「他們把從前瞎眼的人帶到法利賽人那裡」的.他們遇有不明的事情,第一想起的就是求問於有話語權的人.他們盲目相信有權的人,成為盲目的群眾.誰大聲誰對,誰有權誰高.「盲從權力派」,自然「能看見的,反瞎了眼」.

第四批,瞎子的父母.他們是「遠離紛爭自保派」.法利賽人要逼問瞎子父母.父母因畏強權,便支支吾吾,推卸責任,就是連自己的親生骨肉,都可以放棄不理.v20父母回答說:「他是我們的兒子,生來就瞎眼,這是我們知道的.21至於他如今怎麼能看見,我們卻不知道；是誰開了他的眼睛,我們也不知道.他已經成了人,你們問他吧,他自己必能說.」福音書作者說,v22他父母說這話,是怕猶太人；因為猶太人已經商議定了,若有認耶穌是基督的,要把他趕出會堂.他們怕事自保,面對紛爭,半句公道說話都不敢說,寧願犧牲親人,「各人自掃門前雪,莫管他人瓦上霜」.「遠離紛爭自保派」,當然「能看見的,反瞎了眼」.

就只有瞎子,是「失明的才能看見」,甚至到最後認出耶穌是主.怎麼會這樣的呢?他生來瞎眼,在當時文化裡是不祥人,因此父母也棄之於街頭,他自少是個乞兒仔.身無一物,也沒有人教過他什麼對錯道德.所以他不會如法利賽人得勢當權,也不懂如門徒般站在道德高牆看世界.他也不知權威為何物,因此不會如路人甲乙丙盲從權力.他也本來無一物,前途一片空白,所以不需要如父母一樣計算個人利益安危.所以他的膽子很大,竟然公然頂撞法利賽人,為真理化身的耶穌說項:瞎子回答法利賽人,v25說:「他是個罪人不是,我不知道；有一件事我知道,從前我是眼瞎的,如今能看見了.」… v27「我方才告訴你們,你們不聽,為甚麼又要聽呢?莫非你們也要作他的門徒嗎?」… v30「他開了我的眼睛,你們竟不知道他從哪裡來,這真是奇怪!31我們知道上帝不聽罪人,惟有敬奉上帝,遵行他旨意的,上帝才聽他.32從創世以來,未曾聽見有人把生來是瞎子的眼睛開了.33這人若不是從上帝來的,甚麼也不能做.」哇!真的是「唔識死」「赤子之心」「天真得令人敬佩」的公道說話!結果,真的被法利賽人 v34「把他趕出去了.」不過,不怕,他就是懷著一夥天真的心,在路上遇見耶穌,並且認出耶穌是主的真相,真人示範了「失明的才能看見」上主的悖論.

弟兄姊妹,我們都是活在撕裂紛爭的世代裡,也許真的不易看清方向,但是我們至少可以避免迷路:

1. 我們要避免成為「得勢當權派」.不要靠近權力,不要想從權力乞求好處,即使是為保存信仰的緣故；

2. 如果我們在撕裂紛爭裡是被壓受屈的弱勢,要保持清醒,免得衰化為「道德高牆派」,以為自己是審判人的主；

3. 我們也不要成為愚蠢的群眾,盲從權力,權威；

4. 我們也要求主給我們勇氣,不要淪為「遠離紛爭自保派」,「各人自掃門前雪,莫管他人瓦上霜」,連半句公道說話都不敢說.

5. 在撕裂紛爭中,求主賜下一顆天真而懂得分辨是非的心,讓我們能夠在路上遇見基督耶穌我們的主.

阿們.



\newpage

\section{體因罪而死心靈因義而活}
\label{sec:496}
\textbf{鄭漢文博士 - 身體因罪而死,心靈因義而活}
\newline
\newline
\hyperref[sec:495]{< < < PREV SERMON < < <}
~
\hyperref[sec:toc]{[返目錄]}
~
\hyperref[sec:497]{> > > NEXT SERMON > > >}


羅馬書 8:6-11
\newline
\begin{longtable}{cl}
\hline
\hline
章節 & 經文 (和合本修訂版)\\
\hline
8:6 & \begin{tabularx}{0.7\textwidth}{X} 體貼肉體就是死；體貼聖靈就是生命和平安。 \end{tabularx} \\ \\ \relax
8:7 & \begin{tabularx}{0.7\textwidth}{X} 因為體貼肉體就是與神為敵，對神的律法不順服，事實上也無法順服。 \end{tabularx} \\ \\ \relax
8:8 & \begin{tabularx}{0.7\textwidth}{X} 屬肉體的人無法使神喜悅。 \end{tabularx} \\ \\ \relax
8:9 & \begin{tabularx}{0.7\textwidth}{X} 如果神的靈住在你們裡面，你們就不屬肉體，而是屬聖靈了。人若沒有基督的靈，就不是屬基督的。 \end{tabularx} \\ \\ \relax
8:10 & \begin{tabularx}{0.7\textwidth}{X} 基督若在你們裡面，身體就因罪而死，靈卻因義而活。 \end{tabularx} \\ \\ \relax
8:11 & \begin{tabularx}{0.7\textwidth}{X} 然而，使耶穌從死人中復活的神的靈若住在你們裡面，那使基督從死人中復活的，也必藉著住在你們裡面的聖靈使你們必死的身體又活過來。 \end{tabularx} \\ \\
[1ex]
\hline
\hline
\end{longtable}
講題:身體因罪而死,心靈因義而活 Sinful Body Dies, Righteous Spirit Stays Alive

講員:鄭漢文博士

經文:羅馬書8:6-11

場合:香港中文大學崇基學院禮拜堂主日崇拜

日期:2020年3月29日

1. 死路,生路,待尋道

各位觀看今日崇拜直播或網錄的弟兄姊妹或者朋友,我首先向你問安: “Stay safe and be healthy”.「事事小心,保持健康」,成了我近三個月給學生友朋的問安.

也先向大家做一個告解式的致歉(a confession, an apology).

這是我人生第一次做網上直播講道,面對沒有平時的會眾,好像我們這個學期要在網上教學一樣沒有了學生的面容,我未必好適應,所以我今天的講法,與平日有些不同.

今日的講題,其實在兩三個月前已經早定.

其實,當時我不知道我將會講甚麼,在三代經課中我選了書信經課,即係羅馬書8章6至11節.讀到了第10節,就從原句抽出文字來作講題,這樣引用經文作主線,是比較安全的做法.

其實,當時一月下旬我已想到了幾個都幾出位的講題,有些構想較為語不驚人死不休,現在發現好在我最後沒有用上.

其實,當時疫情正開始令人擔心,無法想像到今天無法聚會,變成網上直播.講道的言說對象變成不是一個個具體在前的會眾,變成身在疫情,可能在家求平安的香港市民.難得大家在面對死路,生路的當下依然尋道.

講道需要回應時代.所以,今天我用一種類比思維,將經文與經驗互相參讀會讀出甚麼道來.

最近我們都常常被困在家,我作為中大教師教書時除了Zoom來也Zoom去,有閒時都上社交媒體.

最近有不少疫情相關的片段令人發笑,令人發愁.也有些讓人發怒,讓人心省.

因為這數天才終定講章的撰寫,讓我可以在今日用上世人面對疫情幾個片段(包括視聽影音以及文字故事)來映照今日經課的思考.多謝片段中的他們,給我靈感,取材借來思考活在聖靈的生命境界.

面對從武漢開始廣傳至全球的新冠狀病毒,大概可以歸結抗疫的基本功,讓我們借今天直播或網錄再來呼籲一下,包括這五項:

(1)勤洗手wash hands,

(2)戴口罩wear mask to protect yourself and protect others,

(3)留在家stay home for saving lives,

(4)擴群交距離extend social distancing

(5)增強抗疫力strengthen immunity

抗疫五功,缺一不可.

近日看到了好多世人面對疫情的文字故事以及影音影像,都在這五項工夫上作了有趣的二次創作.

例如,有人仿傚Bee Gees唱Stayin’ Inside …… 記得留在家.

又例如,有人仿傚Beatles唱I gotta wash my hands …… 記得勤洗手.

而教洗手的影片也非常多.(洗手七式20秒.)

如果我現在用Zoom教課的話,我就一定會播給學生看,一來可以佔點時間,這是「養生教學法」的精要,二來可以開心o下,在不安的世代留下一點生命氣息,「活靈活現一下」.

如果可以給今天講章一個副題,我會命名為「帶罪之身難呼吸,活靈活現五生態」.

好,回到今日剛才讀過的經課:羅馬書8章6至11 節.

我們會分三個段落來思考這六節經文.

首先是第一段(第6節)關於生命與和平的.

跟著是一連三節(第7-9節)關於上帝的.

最後是最後兩節(第10-11節)關於基督與復活的.

而貫串六節三段落的是聖靈,是展現活在聖靈中的生命所活現的五大特性,我叫「活靈活現的靈命五態,缺一不可」.

2. 活靈活現生命,活靈活現平安(羅8: 5-6)

首先兩項活現,是活在聖靈中的生命,會活現生命,會活現平安.

生命原來是一場無妄之災.

原來生命可以掌握的實在好少,那麼,最重要的是求生命求平安.

那麼,甚麼才可以保生命與保平安?

回到今日經課第一句,我加上前一節作為聯讀,以下是第5至6節:

5因為,隨從肉體的人體貼肉體的事；隨從聖靈的人體貼聖靈的事.

6體貼肉體的,就是死；體貼聖靈的,乃是生命,平安.

中文版本有一個觀念是消失了,但在英文是好明顯的.我用NIV版讀一次:

5 Those who live according to the flesh have their minds set on what the flesh desires; but those who live in accordance with the Spirit have their minds set on what the Spirit desires.

6 The mind governed by the flesh is death, but the mind governed by the Spirit is life and peace.

Mind這個字眼,在英文版是好關鍵的.

保羅在這段話,有三點值得思考.

一．人有體與靈(flesh-spirit)的二分.後來發展出身心二元論(Mind-Body Dualism).

二．這是死路與生路的選擇(it concerns a matter of life and death).如何出死入生?

三．在二者之間,人有一個mind可以定向:隨從甚麼?體貼甚麼?叫做mind-set.我把mind譯成意志,意志是可以定向的,中國人叫立志.捨棄一個生命境況進而取另一種生命境界.

最近我們常困在家,大家上社交媒體.其中有不少片段令人發笑,令人發愁.

其中一段是一個兩三年歲的BB女叫Layla,媽媽告訴她平時常去的食店都關了,常叫的外賣也沒法送了,唯一可以食的是媽咪煮的食物.這段母女對話的影像,循環地有BB女說話媽媽說話BB女大喊,……我忍不住放上了我的Facebook,寫了一段評述:

我的習慣是每日三評,每天選三段資訊作評論.(已經十年了,我邀請通識教育科和價值教育的學生跟著追讀原資訊.)

我的習慣是中文資訊中文評,英文資訊英文評.

這次我是用英文寫的.我放之在講章上,你可以下載看,或到我Facebook看:

“So poor this girl can’t eat at eating places outside home, nor order any deliveries, but eat only food cooked by mom, for this her mom apologized.

Multiple sadness ……

Girl can’t eat at any eating places nor delivered food.

Girl can but eat unpreferred food cooked by mom.

Mom needs to feed this daughter who disliked her cooking result.

Mom needs to apologize.

This scene was filmed by (her dad?)

This film is circulated on social media (for this I apologized).

(This is) Best advertisement for names of eating places mentioned by this crying girl.”

所寫的意思大概係:

「呢個女仔好慘情,無法出街食平時食店,又無得叫外賣,唯一可食是老媽子手藝.為此媽咪不斷道歉.

這裡出現多重哀傷……

女孩無法出街食平時食店；

女孩無得叫外賣；

女孩唯一可食的是老媽子手藝(但好明顯唔想食)；

媽咪要餵食這個女兒,但女兒不喜歡食老媽子手藝；

為此媽咪不斷道歉；

是誰拍下這段影像?(是老爸/老豆嗎?)

是誰把這影像放上社交媒體?(為此我也道歉!)

片段中大哭大喊BB女叫出不少食店名字,可算是最佳廣告.」

我放上面書後,收到不少likes及轉載.

一兩日後,在我的post下收到一個comment,是一位男士R標記一位女士L,寫道:“This is us.”(「係我地.」)我讀到的意思係:老公對老婆說「係我地.」

我立即覆返句安慰兼致歉的回應.

正所謂:出街食又死,在家食又死.出街食係冒生命危險被感染而死,在家食係悶死.

我們死在肉身的食慾.

3. 活靈活現上帝(羅8: 7-9)

我們都是「帶罪之身」,我們的罪性由肉身開始.

甚麼叫「帶罪之身」?讓我作個類比,我們都是「潛在帶菌者」.

其實這不準確,因為這次從武漢開始廣傳至全球的是新冠狀病毒,而不是細菌.

所以,需要改口說,我們都是「潛在帶病毒者」.

病毒的傳染性非常不清楚,正如罪性的傳染性.

「帶病毒之身」,就當隔離,唔好累街坊.

你其實不知自己係咪「隱形帶病毒者」,所以最好戴口罩.

有手帶的人,可以自律嗎?人有罪性係會違規,人人有天性嚮往自由,需要被紀律.

「帶病毒之身」就是帶罪之身.

帶罪立功就自我隔離,道理咁簡單,點解依然會走來走去?甚至五百萬人一起出走?

因為人的罪性.關於罪的傳染性,保羅說出了三個基本原則:

一是若我們的心志是體貼肉體的就多數中招.

二是若我們的心志是妄視更高的法度.

三是若我們的心志是拒絕上帝的靈住入人心.

先說第一個原則:單單食,就可以是罪的根源.

幾年前因為身體健康要改變飲食,所以查閱了一些飲食哲學書.其中一句話感到震撼:”You are what you eat.” 聞說病源起於武漢人食野味.

世界正常時,好難搵食.現在病毒橫行時,仲難搵食.

食,真係好重要.

有一段新聞(我修輯如下):

(source: \url{https://skypost.ulifestyle.com.hk/article/2598782/}【無視隔離令】戴監察手帶少女不戴口罩出街　鄰座食客揭發窮追300米到酒店)

「一名男網民在3月23日深夜在Facebook群組上載一段影片,指當晚約10時在沙田某日式餐廳用膳期間,家人發現鄰座一名戴鴨舌帽的少女右手戴有監察手帶,當時該名少女指因正在食飯,故未有佩戴口罩,身旁有一位年長男士陪同.

該名網民指原本打算避開對方了事,惟擔心該名少女會繼續肆無忌憚地在社區遊走,遂決定即時報警,並上前直指對方不是,該對男女隨後結帳離開餐廳,最後返回約300米外的一間酒店.

新冠肺炎疫情持續肆虐,政府規定入境人士須接受14日強制家居檢疫,惟23日再有人被揭懷疑違禁外出.警方已在XX酒店尋回該名13歲少女,已送往YY的檢疫中心檢疫.」

民以食為天?體貼一下肉體……

也許,回流留學生在機上不去洗手間又不吃任何食物呆坐了十數小時,下機入住酒店時立即去食餐好.這也許是接應少女的那位男士(爸爸?)非常體貼.但體貼肉體變成不義行為.

追拍她二人的男士不過是在食一餐飯,到後來才發現鄰座食客戴有監察手帶,自己可能被感染.大叔大姐,食飯je ,洗唔洗lor 人命!真係無妄之災.所以要追究下去.

我看到的是悲傷:這個少女如何消化這段經驗?

因為體貼肉體,而忘記了法與義.

忘記法與義的有另一個案例,是一位執法女警.

明報(2020年3月28日星期六)報導:

「昨日(26日)確診感染新型冠狀病毒的22歲女警,早前於明愛醫院抽血等候期間沒佩戴口罩並咳嗽,導致兩名抽血員被列為密切接觸者須隔離檢疫,其中一人因身體不適到急症室求醫.警方在facebook發帖回應稱,該警員防疫意識不足,就事件對醫護人員及服務構成影響,並引起社會不安深感抱歉,已即時聯絡她了解,並要求她必須遵照醫護人員的指示,保護他人避免受到感染.警方又稱,會提醒全體警務人員積極配合醫護人員的防疫工作,避免同類事件再發生.」

(source: \url{https://news.mingpao.com/ins/}\%E6\%B8\%AF\%E8\%81\%9E/article/20200327/

s00001/1585309715037/確診女警醫院沒戴口罩-警-引起社會不安深感抱歉)

對這事件,大家都好驚訝!難以相信.如何反應呢?

「公共醫療醫生協會會長馬仲儀表示,女警拒戴口罩不可饒恕,她說自醫院啟動「緊急應變級別」後,偶有長者因不知疫情情況下沒戴口罩,「但他們好乖,勸一下就會戴」,女警不是因為年紀大或語言不通而無法溝通,年輕的她應了解疫情.馬仲儀說,兩名抽血員要隔離是完全可避免,「相信所有醫護(對此事)都感到不開心」.」

明報(2020年3月28日星期六)

「香港護士協會主席,衛生服務界立法會議員李國麟批評女警拒戴口罩是「缺德」,她身為執法者應「知道什麼是紀律」.」

明報(2020年3月28日星期六)

醫管局員工陣線如何回應?

【強烈要求醫院管理局譴責病人不負責任,嚴重危害員工健康】

(source: \url{https://www.facebook.com/107732697364538/posts/154856235985517/})

「醫管局員工陣線(下稱本會)對以上事件表示極度不滿,更對醫管局不願譴責病人不負責任之行為感到非常遺憾.於記者會上,何婉霞稱只能盡量勸喻病人佩戴口罩,更呼籲病人與醫護人員合作.然而本會認為事件非常嚴重,尤其兩名醫護已經需要接受隔離檢疫,直接導致前線抗疫人手減少；若前線醫護不幸染病,身為僱主的醫管局更加責無旁貸.另一方面,醫管局只以口頭勸喻方式建議病人佩戴口罩,無疑妄顧該病人對其他無辜病人傳播病毒之風險.本會要求醫管局公開譴責有關病人不負責任之行為,並就本會以下三項建議作出回應……

武漢肺炎疫情由年初發展至今,政府堅拒封關.一月時,政府和醫管局皆異口同聲相信求診人士並不會隱瞞旅遊史,引致一個又一個的輸入個案,全因為一個「信」字.今天,病人拒絕配戴口罩讓醫護人員曝露於危險當中,醫管局居然聲稱只能盡量勸喻,背後又是一個「信」字.面對不守規則的病人危害員工安全,醫管局的角色是否應如斯渺小?」

對於武漢肺炎確診女警不戴口罩之累人累物,人的惡性,真的莫名其妙?連公共衛生都不服?

一位神學教授在Facebook這樣解說:

「受感染是身體有問題,

去醫院唔肯帶口罩是精神有問題,

勸喻後仍不肯帶是靈性有問題,

一個精神及靈性都有問題的人,是配槍執法的人.」

這就讓我們明白今天經課的第二段經文(第7至9節):

7原來體貼肉體的,就是與上帝為仇；因為不服上帝的律法,也是不能服,

8而且屬肉體的人不能得上帝的喜歡.

9如果上帝的靈住在你們心裡,你們就不屬肉體,乃屬聖靈了.人若沒有基督的靈,就不是屬基督的.

一個「潛在帶病毒者」被確診為「帶病毒之身」,因為惡性,便成為「帶罪之身」.

保羅說出了三個基本原理來明白人之行惡,是人的罪性深植人心:

一是若我們的心志是體貼肉體的就多數中招.

二是若我們的心志是妄視更高的法度.

三是若我們的心志是拒絕上帝的靈住入人心.

附帶一提,之前主日學教過《使徒信經》最重要的三句認信:

我信上帝

我信我主耶穌基督

我信聖靈

這段羅馬書8章7至9節三節經文又再呈現我們所信的是一位三一神.(三一神論是保羅神學常現的特徵.)

4. 活靈活現基督,活靈活現復活(羅8: 10-11)

我們如何選擇出死入生?

我們如何選擇出生入死?

「帶病毒之身」有兩種反應.在同一個地方,出現兩個故事.

34歲女護士Daniela Trezzi在意大利北部倫巴第(Lombardy)一間醫院的加護病房負責前線抗疫工作,該地區為全國疫情最嚴重的地區.她出現新冠肺炎病徵後接受檢測,3月10日證實確診,其後因非常擔心自己會傳染病毒給其他人而自殺了斷生命.意大利國家護士聯會對事件表示痛心及遺憾,並指「每位護士都是為了拯救生命做好事而選擇這個專業,但都會有不好的情況發生」.

(source: \url{https://travel.ulifestyle.com.hk/news/detail/31328/}意大利34歲護士確診肺炎-擔心傳染他人選擇自殺了斷/1/)

這又何苦呢?她不過是「帶病毒之身」,何必就義?

另有一個故事,也是就義的生命.

題目:意大利倫巴第疫區醫生的見證

翻譯:Debby

時間:2020.3.22

證人收集人:Gianni Giardinelli

Témoignage de Julian Urban 38ans · Docteur en Lombardie:

(Source: \url{http://www.lavoieduciel-garabandal.fr/archives/2020/03/22/38123116.html})

座落意大利城市倫巴第這所醫院的醫生寫的原文是法文,以下是Debby的漢譯.

朱利安·厄本(Julian Urban)的證詞38歲倫巴第醫院的醫生 :

「我從來沒有想過在,如同最黑暗的噩夢中經歷了在我醫院裡這三周里發生的一切.噩夢如江河在流淌,河水越來越大.最開始的時間,來了幾個病人 ,接著是幾十個,然後是數百個,現在我們已經不再是醫生了,我們如同流水線上的分揀員,我們要決定誰應該留下治療,誰應該被送回家等死.儘管他們勞作​​一生給意大利政府繳了稅.

直到兩週之前,我和我的同事都是無神論者.這是正常的,因為我們是醫生,我們相信科學,而且科學排斥上帝的存在.

我總是嘲笑我那去教堂的父母.

九天前,一位75歲的牧師來到醫院.他是一個好人,有嚴重的呼吸困難,但總有一本聖經伴隨他左右.他給我們留下了很深刻的印象,他忙著在醫院裡給臨終前病人 讀聖經,並緊緊地抓住他們的手.

當我們有時間聽他讀聖經的時候,我們都已經是疲倦,沮喪,心理和身體不堪重負.

現在,我們必須承認,作為人,我們已經達到了極限,我們不能做的事情越來越多,每天的死亡人數越來越多.

我們已經筋疲力盡,我們有兩個同事死了,而另一些同事已被感染.

我們意識到,在人無能為力的時候,我們需要上帝,當我們只有幾分鐘的空閒時間,我們便開始向禱告呼求他的幫助. 我們同事們彼此交談,說我們不能相信之前我們還是堅定的無神論者,但是現在我們每天都在尋求上帝賜給我們平安,請求上帝賜予我們力量幫助我們堅持下去,以便我們能夠照顧病人.

昨天,這位75歲的牧師去世了,儘管在3週內有120多人死亡,我們都精疲力竭,崩潰.這位牧師,儘管他的健康狀況,和我們每天的困難處境下,他成功地把上帝而來的平安帶給我們,是我們曾經不再希望能夠找到的.

牧羊人去了他的主耶穌那裡了,如果情況繼續這樣下去,我們很快也會跟隨他而去.

我已經六天沒有回家了,我想不起來我上次吃飯的時間,我意識到自己在這個地球上毫無用處,我想竭盡全力幫助別人.當我被同胞的痛苦和死亡所包圍時,我因著我能夠重新回到上帝的懷抱而感到平安幸福.」

一個生命的就義,活現了基督精神,讓另一個人得見上帝.從此,看破生死.

我們如何選擇出死入生?

誰又為我們出生入死?在香港,為我們出生入死的會不會是從一月開始要求封關的醫護人員?

政府不聽封關,唯有在家自封.

最近看到,呼籲遠離群集留在家,醫護人員都手執紙牌,寫著:

“We stay in work for you.  You stay at home for us."

「我們為你們留守在工作崗位,你們為我們留在家.」

醫護持牌人包括我們中大前任校長沈祖堯教授.真証傳播2011年訪問當時還在位的校長,youtube(\url{https://www.youtube.com/watch?v=_lqtYcBW4Ok})題目名為:「虛心的人有福了──從沙士英雄到中大校長」,內中提到:

「甚麼驅使這位沙士英雄在退休之年,重新披上戰甲,去打一場更花心力,更耗精神,更加困難的硬仗?原來事情背後,竟有上帝的呼召在其中,叫他成為聖經中的約書亞,決心把上帝的子民帶到應許的美地!」

沙士好難不令我們想起另一位中大人:謝婉雯醫生.

今天身體復活,我們寧願在2003年沙士死去的謝婉雯醫生可以復活過來,再與我們一起抗疫.

10基督若在你們心裡,身體就因罪而死,心靈卻因義而活.

11然而,叫耶穌從死裡復活者的靈若住在你們心裡,那叫基督耶穌從死裡復活的,也必藉著住在你們心裡的聖靈,使你們必死的身體又活過來.

大家還記得我在較早前在主日學教《使徒信經》時有三句版及八句版兩種解法,提到第三句「我信聖靈」的認信,究竟跟最後五句有何關係?

Inclusive reading 會認為最後五句是從屬於第三句「我信聖靈」.

使徒保羅在今日經課這段經文肯定聖靈住入我心裡就讓我們同時可以肯定:

我信罪得赦免.

我信身體復活.

我信永生.

罪歸於取死的身體,心靈才是自由之源.

身體就因罪而死,正如生命因體貼肉身的欲望而敗壞.

心靈卻因義而活,正如生命因基督為我們代死而得以稱義,因而罪得赦免,並且身體復活.

基督取了耶穌的身體,人這一個罪身,但他全然無罪,卻因我們的罪,代死在十字架上.

復活節快到,還可以再聚集在一起崇拜嗎?

也許,下一次不用網上直播,而是同在天家.因為我信永生.

5. 五大活現,缺一不可(羅8: 14-17)

有聖靈注入的生命,活靈活現靈命五態:

活靈活現生命,

活靈活現平安,

活靈活現上帝,

活靈活現基督,

活靈活現復活.

五大活現,缺一不可.

聖靈好像呼吸清新機一樣,把可能走入我們身體傷害呼吸道的病毒可以一一清除,這個病毒叫做罪性,而我們因聖靈而得以做出合乎義的進退.

按孟子理解,義者,宜也.甚麼才是合宜?

上述這五大活現又建基於心靈因應以下的五個合宜,活出義人該走的道路.

在開始時我們提到五項抗疫基本功,以下是五項屬靈抗疫基本功:

(1)勤洗手,洗走貼在肉體上的欲望；

(2)戴口罩,讓隱性的罪性不會交差感染；

(3)留在家,留在主的天家；

(4)擴群交距離,跟社會制度的惡保持距離；

(5)增強抗疫力,讓聖靈充滿生命,讓惡不能臨在身上.

今天講道有一思路,我們從身體所在,一直追到身分所在.

所以用個Mind,立志的成果就是身分的確立,我們都是蒙恩的罪人,卻同時成為了上帝的兒女.

聖靈讓人心有所屬,歸屬上帝.

聖靈恩賜讓人重生得力,能行奇事.

聖靈恩賜是關於未來的,讓身體得生命.

關於聖靈,我跟著5月29日會用使徒行傳2章1至21節以及在7月26日會用羅馬書8章26至39節再多演繹.後會有期,未來十架下除口罩相見.

容我借羅馬書第8章14至17節作為今日講道結束的祈禱:

上帝啊,

你曾承諾,因為凡被上帝的靈引導的,都是上帝的兒子.

我們所受的,不是奴僕的心,仍舊害怕；所受的,乃是兒子的心,因此我們呼叫:「阿爸!父!」

聖靈與我們的心同證我們是上帝你的兒女；既是兒女,便是後嗣,就是上帝你的後嗣,和基督同作後嗣.

如果我們和他一同受苦,也必和他一同得榮耀.

上帝啊,

我們期待你的榮光臨在香港,讓我們活靈活現,戰勝罪惡,戰勝疫情.

阿們.



\newpage

\section{基督是我們的生命}
\label{sec:497}
\textbf{劉國偉先生 - 基督是我們的生命}
\newline
\newline
\hyperref[sec:496]{< < < PREV SERMON < < <}
~
\hyperref[sec:toc]{[返目錄]}
~
\hyperref[sec:498]{> > > NEXT SERMON > > >}


講題:基督是我們的生命 He is our Life

講員:劉國偉先生

場合:香港中文大學崇基學院禮拜堂主日崇拜

日期:2020年4月12日

引言

弟兄姊妹平安!在復活節的早晨,願我們一起思想基督的復活,對我們的信仰有什麼意義,我們如何經歷「基督是我們的生命」這寶貴信息.

兩個重要的問題

今天選讀的福音經課,記載抹大拉的馬利亞在七日的頭一日的清早,來到埋葬耶穌的墳墓,按馬可記載,抹大拉的馬利亞與同行的婦女,買了香膏要去膏耶穌的身體(可16:1) .然而,她們看見封閉墳墓的石頭被挪開了,抹大拉的馬利亞看見的是一個空墳墓,所以她急忙跑去見西門彼得和約翰(耶穌所愛的那門徒),並對他們說「有人把主從墳墓裡挪了去,我們不知道放在哪裡.」(約20:2)

約翰見證他親眼看見那空墳墓,並且墓室裡的細麻布和裹頭巾不尋常地整齊的放在裡面,他就相信耶穌所說,有關祂自己死後復活的預言,約翰回憶彼得與他先後進入那空墳墓,然後就返回他們的住處,仍然堅持留在現場的抹大拉的馬利亞,卻站在墳墓外面哭泣.(v.11)

當復活的耶穌顯現,馬利亞還以為祂是看園的人,她仍在哭泣,正如我去年10月,在這裡分享同一段聖經信息的時候,曾經說過,馬利亞的眼淚是不難明白的,因為那是雙重的損失,耶穌死去了,如今連遺體也不見了.對馬利亞來說,這是一個充滿問號的現場,這三天三夜發生的事,超乎經驗,超乎理解,也超乎想像.那一位滿懷理想,宣揚真理,又把馬利亞從魔鬼的黑暗權勢中,釋放出來的主耶穌,竟然被殺害,卻又連遺體都被人搬走了.

正當這位深愛著耶穌,身體力行地跟隨耶穌,就是耶穌在那裡,她就設法在那裡的馬利亞,在淚眼之中,耶穌問她說:「婦人,為甚麼哭?你找誰呢?」(v.15)弟兄姊妹,你有否注意到,這是復活的主,在祂從死裡復活過來以後的第一句說話,或者說得更準確一點,那是祂向第一位看見祂復活的人,所提出的問題,若果我們看得闊一點,我更願意嘗試把這個重要的問題,看成是耶穌向人類的提問.弟兄姊妹,若果耶穌問你「你為甚麼哭?你找誰呢?」,此時當下的你,又會怎回答這位復活主? 若果我們哭,我們為什麼哭?若果我們尋找什麼人,我們最想尋著的是誰?又或者倒過來問,當我們失去什麼就會牽動我們的眼淚?誰人是我們最不願意失去? 這其實是人生兩個很重要的問題:「最不情願失去的是什麼?最想找得到的是什麼?」

兩個問題的答案

我們看看今天選讀的詩篇,嘗試在其中尋找這兩個問題的答案.

詩篇118篇是一篇感恩詩,內容聚焦於上主的拯救,全詩可分作三大段落,我們今天所誦讀的第1及2節,屬於整篇詩的引言部份,感謝上帝的慈愛永遠長存；然後,我們誦讀了14至18節,屬於第二個段落,19至24節屬於第三個段落.

當詩人為上帝的大能與拯救表達出一種歡呼的情緒時(詩118:14-16),他清楚表明,因上主的拯救,他「必不致死」(v.17),而且上主「未曾將他交於死亡」(v.18).詩人告訴我們,他把生命看得極其寶貴,很明顯的,他不願意因恨他的人對付他,而奪去他的生命.回到我剛才的提問:「最不情願失去的是什麼?」詩人的答案很可能是「生命」.是的,在去年的社會衝突中,我們為到好幾位以死明志的運動參與者,還有社會仍在努力追問的死亡案例,感到心靈陣陣抽搐,因為我們都看重人的生命；這幾個月持續在全球擴散的新冠狀病毒,全球死亡人數累計已超過十萬,我們感到難以接受,心靈陣陣哀傷,因為我們都看重人的生命.

詩篇118篇19節如此寫道:「給我敞開義門,我要進去,稱謝耶和華.」這就展開了敬拜者進入聖殿的步伐.詩人說這是義人要進去的門,是上主親自成為他的拯救,(v.20-21) 所以,在上主的聖殿裡,詩人最想找到的是神自己.詩人說:「匠人所棄的石頭,已成了房角的頭塊石頭.」(v.22) 從新約的記載來對照,是主耶穌親自引用了這節詩句,並套用在自己身上；(太21:42) 使徒彼得面對猶太長官和權貴之時,亦引用這詩句去見證基督,(徒4:11)  使徒保羅亦同樣引用聖殿的房角石來指向基督(弗2:20-21) ；所以,從詩篇118篇延展至新約,對一個追隨基督的人來說,若果要回答「最想找得到的是什麼?」,我會選擇耶穌基督.詩人所寫,「因為你已經應允我,又成為了我的拯救.」(v.21) 正預表了基督親自成就救贖,那是被人棄絕和殺害的主,卻在我們眼中看為至寶.正如使徒保羅的名句:「我也將萬事當作有損的,因我以認識我主基督耶穌為至寶.我為他已經丟棄萬事,看作糞土,為要得著基督.」(腓3:8)

弟兄姊妹,在人生旅途中,我們時常經歷很多誘惑和考驗,有時為了生存,有時為了在慾望之中追追逐逐,我們若果真誠回答「最想找得到的是什麼?」,你的答案會與保羅所說的一樣嗎?

基督是我們的生命

事實上,使徒保羅在歌羅西書論及基督的復活時,他直接說「基督是我們的生命」,(西3:4),這就讓我們更深入了解何以保羅說:「我以認識我主基督耶穌為至寶」.在人生的晚程大部份時間失去人身自由,甚至身陷牢獄的使徒保羅提醒我們,我們若真與基督一同復活,就當求上面的事,(v. 1) 引導我們去思想,一個學習去經歷復活大能的基督徒,需要弄清楚,甚麼才是我們的人生追求.

所謂思念上面的事,不要思念地上的事,並非指向一個割裂的信仰,只談論形而上的宗教理念,不食人間煙火,保羅所指涉的地上的事是「要治死你們在地上的肢體,就如淫亂,污穢,邪情,惡慾,和貪婪(貪婪就與拜偶像一樣)」(西3:5),他又吩咐信徒當棄絕一切的惱恨,忿怒,惡毒( - 或譯:陰毒),毀謗,並口中污穢的言語.(西3:8) 用聖經的簡單扼要講法,那就是邪情私慾.(加5:24) 那麼上面的事又指向什麼? 按保羅的教導,上面有基督坐在上帝的右邊,那是上帝榮耀的寶座,是祂掌權的國度,我們追隨基督的人,就是要活出主禱文「願你的旨意行在地上,如同行在天上」的生命,遵行上帝的旨意,反映上帝的榮耀,並見證祂的權能和國度.

「基督是我們的生命」反映了我們與基督一同復活的生命經歷,向罪我們死了,靠著基督我們活過來,成為新造的人:「若有人在基督裡,他就是新造的人,舊事已過,都變成新的了.」(5:17) 所以,我十分羨慕使徒保羅的自白「 因我活著就是基督」(腓1:21上)這是我一生在主面前努力追求的生命標竿,那是一種與基督同復活的穩行在生命高處的經歷,我這兩年回到母校擔任校牧職事,特別讓我有更深的體會和經歷.

在死亡之地宣佈復活的生命

當使徒彼得衝破心裡障礙,順從聖靈的引導,從約帕前去該撒利亞,向意大利營的百夫長哥尼流傳道之時,他說:「上帝怎樣以聖靈和能力膏拿撒勒人耶穌,這都是你們知道的.他周流四方,行善事,醫好凡被魔鬼壓制的人,因為上帝與他同在.他在猶太人之地,並耶路撒冷所行的一切事,有我們作見證.他們竟把他掛在木頭上殺了.第三日,上帝叫他復活,顯現出來.」(徒10:38-40) 我們都明白福音信息並非單單在於耶穌基督三年的傳道和教導,並治病,趕鬼,也並非止於祂捨身十架,而必須包含祂的復活,彼得的論述亦正如使徒保羅所言:「若基督沒有復活,我們所傳的便是枉然,你們所信的也是枉然.」(林前15:14) 若果基督對我們的愛只停於祂的死亡,沒有從死裡復活,向罪惡與死亡的權勢跨勝,則救贖計劃並沒完成,我們所信的就是枉然.

事實上,見證基督復活是使徒行傳記載的其中一個重點,在使徒行傳中相關於「復活」這個字彙共出現了25次,新約書信更出現了69次,但「復活」不單純是用來討論的教義或神學觀點,這是我們信仰生命所追求,以我們自己本來必死的生命,去經歷基督復活的大能,在每天每夜的生活裡,讓基督的大能在我們的生命中彰顯.

從香港去年的社會事件,到如今全球嚴峻的疫情,我們經歷過衝突,打壓,撕裂,仇恨,隔離,甚至死亡, 在看似幽暗的黎明前夕,身處現場見證基督復活的抹大拉的馬利亞,「在死亡之地宣佈復活的生命」(天主教禮儀及聖事部用語),她告訴門徒說:「我已經看見了主.」

(約20:18)  弟兄姊妹,但願我們都是基督復活的見證人,在死亡之地宣稱「我已經看見了主」.弟兄姊妹,我們不情願失去生命,我們最想得到的就是耶穌基督,就在這復活節的早晨,願我們立志要更深刻去經歷復活的真實,特別在這疫情當中,重新審視自己的生活和人生追求,思念上面的事,不思念地上的事.

復活的主與我們同在,與香港同在,也特別與普世活在艱難,困苦的人同在.願榮耀歸給三一真神.



\newpage

\section{困局中的突破}
\label{sec:498}
\textbf{鄧瑞強博士 - 困局中的突破}
\newline
\newline
\hyperref[sec:497]{< < < PREV SERMON < < <}
~
\hyperref[sec:toc]{[返目錄]}
~
\hyperref[sec:499]{> > > NEXT SERMON > > >}


約翰福音 20:19-31
\newline
\begin{longtable}{cl}
\hline
\hline
章節 & 經文 (和合本修訂版)\\
\hline
20:19 & \begin{tabularx}{0.7\textwidth}{X} 那日（就是七日的第一日）晚上，門徒因怕猶太人，所在的地方門都關了。耶穌來，站在當中，對他們說：「願你們平安！」 \end{tabularx} \\ \\ \relax
20:20 & \begin{tabularx}{0.7\textwidth}{X} 說了這話，他把手和肋旁給他們看。門徒一看見主就喜樂了。 \end{tabularx} \\ \\ \relax
20:21 & \begin{tabularx}{0.7\textwidth}{X} 於是耶穌又對他們說：「願你們平安！父怎樣差遣了我，我也照樣差遣你們。」 \end{tabularx} \\ \\ \relax
20:22 & \begin{tabularx}{0.7\textwidth}{X} 說了這話，他向他們吹一口氣，說：「領受聖靈吧！ \end{tabularx} \\ \\ \relax
20:23 & \begin{tabularx}{0.7\textwidth}{X} 你們赦免誰的罪，誰的罪就得赦免；你們不赦免誰的罪，誰的罪就不得赦免。」 \end{tabularx} \\ \\ \relax
20:24 & \begin{tabularx}{0.7\textwidth}{X} 那十二使徒中，有個叫低土馬的多馬，耶穌來的時候，他沒有和他們在一起。 \end{tabularx} \\ \\ \relax
20:25 & \begin{tabularx}{0.7\textwidth}{X} 其他的門徒就對他說：「我們已經看見主了。」多馬卻對他們說：「除非我看見他手上的釘痕，用我的指頭探入那釘痕，用我的手探入他的肋旁，我絕不信。」 \end{tabularx} \\ \\ \relax
20:26 & \begin{tabularx}{0.7\textwidth}{X} 過了八日，門徒又在屋裡，多馬也和他們在一起。門都關了，耶穌來，站在當中，說：「願你們平安！」 \end{tabularx} \\ \\ \relax
20:27 & \begin{tabularx}{0.7\textwidth}{X} 然後他對多馬說：「把你的指頭伸到這裡來，看看我的手；把你的手伸過來，探入我的肋旁。不要疑惑，總要信！」 \end{tabularx} \\ \\ \relax
20:28 & \begin{tabularx}{0.7\textwidth}{X} 多馬回答，對他說：「我的主！我的神！」 \end{tabularx} \\ \\ \relax
20:29 & \begin{tabularx}{0.7\textwidth}{X} 耶穌對他說：「你因為看見了我才信嗎？那沒有看見卻信的有福了。」 \end{tabularx} \\ \\ \relax
20:30 & \begin{tabularx}{0.7\textwidth}{X} 耶穌在他門徒面前另外行了許多神蹟，沒有記錄在這書上。 \end{tabularx} \\ \\ \relax
20:31 & \begin{tabularx}{0.7\textwidth}{X} 但記載這些事是要使你們信耶穌是基督，是神的兒子，並且使你們信他，好因著他的名得生命。 \end{tabularx} \\ \\
[1ex]
\hline
\hline
\end{longtable}
講題:困局中的突破 Newness in a Closed System

經文:約翰福音20:19-31

講員:鄧瑞強博士

場合:香港中文大學崇基學院禮拜堂主日崇拜

日期:2020年4月19日

各位弟兄姊妹,早安.

在這疫症蔓延的日子,祝願各人身體健康,心靈平安.

在聖經的用語裡,「平安」是指一個人身心靈的健全,一個人整體的圓滿.這種「平安」,只能從神而來.耶穌的救贖工程,就是進入人類至深的恐懼裡,進入死亡裡,廢掉死亡的權勢,然後復活,將「平安」賜給我們.

近來,疫症瀰漫世界,我們整日的思慮,都被疫症佔據.有一晚,我忽發奇想,想像「新冠狀病毒」來探訪我,它用它的角度,來訴說它與我們人類相遇的故事.

「新冠狀病毒」這樣說:

【人啊,人,你們好,我們老朋友又相見了.「我的家族」與「你們人類的家族」世世代代都相識相知.你們人類改變了我們的演化軌跡,我們也改變了你們的歷史命途.

它繼續說:

我和人類相識已久,但不知何解,你們看來很怕我,但我其實平平無奇,我只是一個頑童,走到你們這些人類面前,忽然掀走你們的面具,讓你們看到自己本來的面目,那空洞而脆弱的生命,只此而已.

你們怕我,說我帶來死亡,其實,我沒有帶來死亡,死亡只是你們自身的本相.你們不是怕我,你們是怕面對自己,你們不肯承認自己的虛無.

人啊,人,你們應該多謝我,我讓你們那不斷加快的腳步停下來.你們停下來後,才有空間和時間去重估一切價值,重新認識事情的優先次序,重新釐清價值的高下,重新體會親情的可貴,重新學習慢慢地品嚐一頓飯.人啊,人,是我,讓你們明白安息日的重大意義.你們違反神的安息日很久了.

人啊,人,你們真應該多謝我.是我,讓你們明白,人最有貢獻的作為,就是停下來不作為,讓萬物有空間成就自己.人與人之間,最重要的是與人有一社交距離,孤獨地面對自己.你們的哲人祈克果曾經說過,孤獨地面對自己,才能活出生命的意義.人能真誠地面對自己,才能認識自己不是自己命運的主人,不是世界的主人,不是歷史的主人,而是多麼需要天地間的大恩典去承托的脆弱生命.

人啊,人,你們應該多謝我,我讓你們看到一個全新的世界.與我相遇之後,你們更有可能活出一全然不同的人生.】

以上說的,都是「新冠狀病毒」留下的話.對我們來說,這疫情仍是一危機,但其實每一危機,都是一難得的機遇,讓我們停下來,審視一下生活.唯有這樣,新的人生才有可能.

耶穌的死,對跟從他的人而言,是一大危機.上個禮拜是復活節,我們記念耶穌的復活.耶穌的復活,是在危機中透露的機遇,人在耶穌的復活裡,體驗到全新的可能.

今日的講道經文,是今日的福音書經課:約翰福音20: 19-31

透過這段經文,讓我們體會耶穌的復活帶給我們的新機.

約20:19　那日(就是七日的第一日)晚上,門徒所在的地方,因怕猶太人,門都關了.耶穌來,站在當中,對他們說:「願你們平安!」

約20:20　說了這話,就把手和肋旁指給他們看.門徒看見主,就喜樂了.

約20:21　耶穌又對他們說:「願你們平安!父怎樣差遣了我,我也照樣差遣你們.」

約20:22　說了這話,就向他們吹一口氣,說:「你們受聖靈!

約20:23　你們赦免誰的罪,誰的罪就赦免了；你們留下誰的罪,誰的罪就留下了.」

約20:24　那十二個門徒中,有稱為低土馬的多馬；耶穌來的時候,他沒有和他們同在.

約20:25　那些門徒就對他說:「我們已經看見主了.」多馬卻說:「我非看見他手上的釘痕,用指頭探入那釘痕,又用手探入他的肋旁,我總不信.」

約20:26　過了八日,門徒又在屋裡,多馬也和他們同在,門都關了.耶穌來,站在當中說:「願你們平安!」

約20:27　就對多馬說:「伸過你的指頭來,摸〔原文是看〕我的手；伸出你的手來,探入我的肋旁.不要疑惑,總要信.」

約20:28　多馬說:「我的主!我的神!」

約20:29　耶穌對他說:「你因看見了我才信；那沒有看見就信的有福了.」

約20:30　耶穌在門徒面前另外行了許多神蹟,沒有記在這書上.

約20:31　但記這些事要叫你們信耶穌是基督,是神的兒子,並且叫你們信了他,就可以因他的名得生命.

這個故事的開頭,是一時間標記:「那日(就是七日的第一日)晚上」.要明白這些時間標記的意義,我們需要知道約翰福音的神學思考.約翰福音是反映創世記的創造神學的.

約翰福音開頭的第一節經文:「太初有道」,「太初」這字眼,就是創世記開頭的第一節經文「起初,神創造天地」的「起初」.若大家留意約翰福音的第一章及第二章,你們會留意到有一「七日創造」的結構.

創世記第一日講神創造光明.約翰福音1: 14講「道成了肉身,住在我們中間,充充滿滿地有恩典有真理.我們也見過他的榮光,正是父獨生子的榮光.」這裡也講「光明」,不是普通的「光明」,是「真光」,這真光來到這個世界.

創世記第二日講神創造空間,將水分開,稱穹蒼為「天」.約翰福音1:29那裡說:「次日」,即第二日,這一段落也講「水」,講施洗約翰的「水」禮,然後,也講「天」,在耶穌的受洗裡,聖靈彷彿鴿子,從天而降.

創世記第三日講神創造地,讓萬物有可居之地.約翰福音1:35那裡說:「再次日」,即第三日,這一段落講兩個門徒在耶穌那裡找到可居之地.

時間關係,我不能細說了,我直接跳到第七日.創世記的第七日,神停下來,安息了.安息日的重點不在休息.其重點在什麼地方?約翰福音開場白的第七日,將重點表明出來.約翰福音開頭的第七日,正是2:1打後的故事,即「迦拿婚宴」的「水變酒」的故事.這是婚禮.安息日的重點,正是婚禮,是神與人的婚禮,神與人在愛中共在相連.

問題只是,我們能否在耶穌的救贖工程裡,看到神和你,神和大地,神和萬物在愛中的緊密相連?

約翰福音的第一章及第二章的「七日創造」結構,刻意遺漏第五天及第六天.這要留待日後以耶穌的受難情節來補足.我們知道,禮拜日是第一天,第五天即禮拜四,第六天即禮拜五.在受難故事中,第五天(禮拜四)的高潮是耶穌為門徒洗腳,以愛的服事將生命連結起來,然後耶穌講「葡萄樹與枝子」的比喻,透露祂與信徒的生命連繫.這反映創世記第五天的創造,神在這天創造萬物的生命.原來,真正的生命,要連於耶穌.

受難故事的第六天(禮拜五),耶穌被彼拉多審問.當中有一場景,耶穌戴著荊棘冠冕,穿著紫袍,彼拉多對群眾說:「看啊,這個人」.這是天地間一個真正的人的出場.這反映創世記第六天的創造,神在這天,以祂自己的形像造人.具神形像的人,在世間,何處可見?就在耶穌身上可見.

創世記的第七天,神安息了.受難故事的第七天(禮拜六),耶穌死了,祂安息了.這安息的意義,我剛才講了,就是成就神和人的婚盟,一種愛的永恆結連.在耶穌的救贖裡,神和人復和了.然後,是禮拜日,七日的第一日,耶穌復活了,耶穌成就了新的創造,祂開創了一個全新的天地.

回到今日的講道經文.

約20:19　那日(就是七日的第一日)晚上,門徒所在的地方,因怕猶太人,門都關了.

「七日的第一日」,這時間標記讓我們感覺到神的創造好像還未開始.雖然耶穌已復活,但由於信徒還未領悟,他們仍未體會到復活耶穌所創造的新天新地.那個時間仍是「晚上」,像是創世記的開場白,「地是空虛混沌,淵面黑暗」,他們仍未體會「神的靈運行在水面上」.他們「怕猶太人」,「猶太人」在約翰福音裡,扮演著敵擋神的角色.門徒在黑暗中,在敵擋神的勢力下,在恐懼中,因此「門都關了」.

這些日子以來,罕有地,我們全世界的不同的門都關了.有些國家封關了,其他人不得進入.有些城巿封城了,巿民不能進出.公司的門關了,工廠的門關了,商店的門關了,教堂的門關了,我們的家門也關了.因著疫情,我們互相不再探訪.因著疫情,我們不再同台吃飯.因著疫情,我們甚至不再握手.人人都在恐懼中,人人都在自保中,我們活得好像復活主從來沒有復活一樣,我們活得好像耶穌的救贖工程從沒有開展過一樣,我們活得像無望的人一樣.

「門都關了」,但復活的耶穌忽然闖進來.就好像在創世記裡,神說:「要有光」,就有了光.沒有什麼先決理由的,我們不用問:神用什麼創造光明?神是創造的主,祂說有,就有.「門都關了」,我們的生命進入一密閉狀態,我們想像不到任何新的可能性.但忽然,新的可能性出現了,沒有理由的,超乎我們想像的,復活的耶穌忽然闖了進來,打開了一道肉眼看不見的出口.耶穌向我們展示一全新的天地,全新的創造,全新的人生.

約20:19　……耶穌來,站在當中,對他們說:「願你們平安!」

約20:20　說了這話,就把手和肋旁指給他們看.門徒看見主,就喜樂了.

約20:21　耶穌又對他們說:「願你們平安!父怎樣差遣了我,我也照樣差遣你們.」

約20:22　說了這話,就向他們吹一口氣,說:「你們受聖靈!

約20:23　你們赦免誰的罪,誰的罪就赦免了；你們留下誰的罪,誰的罪就留下了.」

在我們的不安裡,耶穌宣告「平安」.然後,耶穌向他們展示一些東西.耶穌不是向他們展示復活的光輝燦爛,耶穌是向他們展示充滿死亡氣息的傷痕.要參透復活,必需正視受難與死亡.這是因愛他人而承受的苦難與死亡.在這十架的愛的深處,就是天父的心.觸碰這顆心,人就看見新天新地.很多基督徒毫不喜樂,毫無平安,心門常關,是因為不甘願實踐這「承受十架艱難」的愛.門徒看見耶穌的受苦記號,他們明白了,他們便喜樂了.

門徒看見了耶穌透過死而復活所創造的新天新地,他們成為新的人類.新人類有新的人生使命.耶穌向他們說:「願你們平安!父怎樣差遣了我,我也照樣差遣你們.」門徒成了延續基督使命的人.如何延續呢?力量何在呢?使命是什麼呢?

約20:22　〔耶穌〕說了這話,就向他們吹一口氣,說:「你們受聖靈!」

約翰福音採用了創世記裡神造人故事的情節,神向人吹了一口氣,耶穌也向門徒吹一口氣,如此,新的人類誕生了,他們秉承著「神的形像」,「神的新形像」,「神的真正的形像」,即「基督的形像」.門徒要像基督一樣活著.他們如何能如此活呢?憑著聖靈.聖靈瀰漫天地,以其創造的生命力,承托著神的子民.神的子民的使命是什麼?

約20:23　「你們赦免誰的罪,誰的罪就赦免了；你們留下誰的罪,誰的罪就留下了.」

很明顯,這是一種「判斷」的使命,像審判官一樣,去判斷誰屬乎真理,誰屬乎光明,誰屬乎愛?這判斷的使命意味著行動的使命,我們要堅持真理,堅持光明,堅持愛.如此,我們便有份於復活耶穌所創造的新天新地,也有份與耶穌一起去拓展這新天新地的領域.

復活的故事,要講的內容,不是復活的耶穌來帶走我們,帶我們到天堂那裡,而是講:復活的耶穌創造了一個新天新地,在這天地裡,死亡被消滅,恐懼被克服,邪惡的勢力被征服,黑暗被驅散,緊閉的門被開啟.而耶穌呼喚我們,像祂一樣,在塵世,好好地活,活出真理,活出光明,活出愛,好叫更多的人領悟這新天新地的福樂.

以上的故事,還有下集.

約20:24　那十二個門徒中,有稱為低土馬的多馬；耶穌來的時候,他沒有和他們同在.

約20:25　那些門徒就對他說:「我們已經看見主了.」多馬卻說:「我非看見他手上的釘痕,用指頭探入那釘痕,又用手探入他的肋旁,我總不信.」

約20:26　過了八日,門徒又在屋裡,多馬也和他們同在,門都關了.耶穌來,站在當中說:「願你們平安!」

約20:27　就對多馬說:「伸過你的指頭來,摸〔原文是看〕我的手；伸出你的手來,探入我的肋旁.不要疑惑,總要信.」

約20:28　多馬說:「我的主!我的神!」

多馬,代表著門徒以外的一個人,門徒向他講述他們看見的一個新世界.他不信,他有自己的一套不信的理由,他的出現令門依然緊緊關上.

各人應該留意到,經文這裡的時間標記是十分特別的,經文說:「過了八日」.我們知道,門徒在「七日的第一日」(即禮拜日)聚集,他們以耶穌的復活日來作為敬拜神的日子.我們預期,經文這裡也會再一次寫「七日的第一日」,或寫「過了七天」,但出乎我們預期,經文寫「過了八日」.這裡,「八」字是重點所在.在創世記,神以七天完成創造.第八天,就是全新的天地了.約翰福音這裡,凸顯第「八」天,就是說,復活的耶穌已重造天地,也重造了具有「神的真正形像」的基督門徒,一個全新的天地已出現了.剩下來的情節,就是門徒在人類歷史裡,效法基督,背負起人類的十架,實現天父的大愛.

在這故事框架裡,出現了多馬.按故事的情節發展,應是門徒以十架的愛,去令多馬相信耶穌.但又再出乎意外,是復活的耶穌親自來,向多馬展示因愛而受苦的記號.或許,約翰這樣描述,正正是要表明一個信仰真理,就是:當信徒向未信者展示十架的愛時,就等於耶穌親自向這人展示十架的愛.

看見這十架的愛,多馬作出最徹底的認信:「我的主!我的神!」他找到生命的主人,他找到敬拜的終極對象.

這些日子,疫情蔓延,各種門都關了.就在這種時刻,讓我們明白耶穌復活對這世界的改造.他造了一個新天新地,一個死亡沒有終極話事權的天地,一個以愛為本的天地,他也造了一群「新造的人」,這群人(即你和我)懷著使命,以真理,以光明,以愛,讓這新天新地成為人人可見的現實.

今日,我們這群看見這真相的信徒,應以更大的勇氣,在這令人不安的,黑暗的,人心緊閉的時刻,活出更大的愛.願天國早日實現,願神的旨意成就在人間.

但願:榮耀歸於聖父,聖子,聖靈.阿們.

(附記:本人在崇基禮拜堂的講章已結集成《荷花池畔春草綠》講道集,試讀本可參閱:\url{http://cc.cuhk.edu.hk/dev/images/content/fund-raising-campaigns/university-sermons-by-dr-sk-tang/CCC-SK-Tang-Sermon-Text-Sample-2.pdf}.較早期的崇基講章,也已結集成《敢於跟隨主》一書.對這兩本書有興趣的朋友,可到各大基督教書局購買.)



\newpage

\section{踏上這無盡旅途}
\label{sec:499}
\textbf{林豪恩先生 - 踏上這無盡旅途}
\newline
\newline
\hyperref[sec:498]{< < < PREV SERMON < < <}
~
\hyperref[sec:toc]{[返目錄]}
~
\hyperref[sec:500]{> > > NEXT SERMON > > >}


使徒行傳 2:14,36-41
\newline
\begin{longtable}{cl}
\hline
\hline
章節 & 經文 (和合本修訂版)\\
\hline
2:14 & \begin{tabularx}{0.7\textwidth}{X} 彼得和十一個使徒站起來，他就高聲向眾人說：「猶太人和所有住在耶路撒冷的人哪，這件事你們要知道，要側耳聽我的話。 \end{tabularx} \\ \\ \relax
2:15 & \begin{tabularx}{0.7\textwidth}{X} 這些人並不像你們所想的喝醉了，因為現在才早晨九點鐘。 \end{tabularx} \\ \\ \relax
2:16 & \begin{tabularx}{0.7\textwidth}{X} 這正是藉著先知約珥所說的： \end{tabularx} \\ \\ \relax
2:17 & \begin{tabularx}{0.7\textwidth}{X} 『神說：在末後的日子，我要將我的靈澆灌凡血肉之軀的。你們的兒女要說預言；你們的少年要見異象；你們的老人要做異夢。 \end{tabularx} \\ \\ \relax
2:18 & \begin{tabularx}{0.7\textwidth}{X} 在那些日子，我要把我的靈澆灌，甚至給我的僕人和婢女，他們要說預言。 \end{tabularx} \\ \\ \relax
2:19 & \begin{tabularx}{0.7\textwidth}{X} 在天上，我要顯出奇事，在地下，我要顯出神蹟，有血，有火，有煙霧。 \end{tabularx} \\ \\ \relax
2:20 & \begin{tabularx}{0.7\textwidth}{X} 太陽要變為黑暗，月亮要變為血，這都在主大而光榮的日子未到以前。 \end{tabularx} \\ \\ \relax
2:21 & \begin{tabularx}{0.7\textwidth}{X} 那時，凡求告主名的都必得救。』 \end{tabularx} \\ \\ \relax
2:22 & \begin{tabularx}{0.7\textwidth}{X} 「以色列人哪，你們要聽我這些話：拿撒勒人耶穌就是神以異能、奇事、神蹟向你們證明出來的人，這些事是神藉著他在你們中間施行，正如你們自己知道的。 \end{tabularx} \\ \\ \relax
2:23 & \begin{tabularx}{0.7\textwidth}{X} 他既按著神確定的旨意和預知被交與人，你們就藉著不法之人的手把他釘在十字架上，殺了。 \end{tabularx} \\ \\ \relax
2:24 & \begin{tabularx}{0.7\textwidth}{X} 神卻將死的痛苦解除，使他復活了，因為他原不能被死拘禁。 \end{tabularx} \\ \\ \relax
2:25 & \begin{tabularx}{0.7\textwidth}{X} 大衛指著他說：『我看見主常在我眼前，他在我右邊，使我不至於動搖。 \end{tabularx} \\ \\ \relax
2:26 & \begin{tabularx}{0.7\textwidth}{X} 所以我心裡歡喜，我的舌頭快樂，而且我的肉身要安居在指望中。 \end{tabularx} \\ \\ \relax
2:27 & \begin{tabularx}{0.7\textwidth}{X} 因你必不將我的靈魂撇在陰間，也不讓你的聖者見朽壞。 \end{tabularx} \\ \\ \relax
2:28 & \begin{tabularx}{0.7\textwidth}{X} 你已將生命的道路指示我，必使我在你面前充滿快樂。』 \end{tabularx} \\ \\ \relax
2:29 & \begin{tabularx}{0.7\textwidth}{X} 「諸位弟兄，先祖大衛的事，我可以坦然地對你們說：他死了，也埋葬了，而且他的墳墓直到今日還在我們這裡。 \end{tabularx} \\ \\ \relax
2:30 & \begin{tabularx}{0.7\textwidth}{X} 既然大衛是先知，他知道神曾向他起誓，要從他的後裔中立一位坐在他的寶座上。 \end{tabularx} \\ \\ \relax
2:31 & \begin{tabularx}{0.7\textwidth}{X} 他預先看見了，就講論基督的復活，說：『他不被撇在陰間；他的肉身也不見朽壞。』 \end{tabularx} \\ \\ \relax
2:32 & \begin{tabularx}{0.7\textwidth}{X} 這耶穌，神已經使他復活了，我們都是這事的見證人。 \end{tabularx} \\ \\ \relax
2:33 & \begin{tabularx}{0.7\textwidth}{X} 他既被高舉在神的右邊，又從父受了所應許的聖靈，就把你們所看見所聽見的，澆灌下來。 \end{tabularx} \\ \\ \relax
2:34 & \begin{tabularx}{0.7\textwidth}{X} 大衛並沒有升到天上，但他自己說：『主對我主說：你坐在我的右邊， \end{tabularx} \\ \\ \relax
2:35 & \begin{tabularx}{0.7\textwidth}{X} 等我使你的仇敵作你的腳凳。』 \end{tabularx} \\ \\ \relax
2:36 & \begin{tabularx}{0.7\textwidth}{X} 故此，以色列全家當確實知道，你們釘在十字架上的這位耶穌，神已經立他為主，為基督了。」 \end{tabularx} \\ \\ \relax
2:37 & \begin{tabularx}{0.7\textwidth}{X} 眾人聽見這話，覺得扎心，就對彼得和其餘的使徒說：「諸位弟兄，我們該怎樣做呢？」 \end{tabularx} \\ \\ \relax
2:38 & \begin{tabularx}{0.7\textwidth}{X} 彼得對他們說：「你們各人要悔改，奉耶穌基督的名受洗，使你們的罪得赦免，就會領受所賜的聖靈。 \end{tabularx} \\ \\ \relax
2:39 & \begin{tabularx}{0.7\textwidth}{X} 因為這應許是給你們和你們的兒女，並一切在遠方的人，就是給所有主—我們的神所召來的人。」 \end{tabularx} \\ \\ \relax
2:40 & \begin{tabularx}{0.7\textwidth}{X} 彼得還用更多別的話作見證，勸勉他們說：「你們當救自己脫離這彎曲的世代。」 \end{tabularx} \\ \\ \relax
2:41 & \begin{tabularx}{0.7\textwidth}{X} 於是領受他話的人，都受了洗；那一天，門徒約添了三千人。 \end{tabularx} \\ \\
[1ex]
\hline
\hline
\end{longtable}
講題:踏上這無盡旅途 Walk onto this Endless Journey

講員:林豪恩先生

經文:使徒行傳2:14,36-41

場合:香港中文大學崇基學院禮拜堂主日崇拜

日期:2020年4月26日

一．引言

「踏上這無盡旅途」,不是用來描述抗疫之路.雖然抗疫之路比很多人預期的漫長,也比人們預期的艱辛很多,但人們仍然相信,疫情總有結束的一天.「踏上這無盡旅途」是Supper Moment的經典作品「無盡」的其中一句歌詞.不知道你有沒有聽過這首歌,對於某一代的人,Supper Moment的「無盡」可以媲美Beyond的「海闊天空」.今天,我們不是講疫情,也不是講流行曲,而是講信仰.「踏上這無盡旅途」可否用來描述你的信仰路呢?你對這旅途有甚麼感覺呢?精彩?苦悶?平淡?渴望繼續向前行?想過離開?迷茫徘徊?讓我們透過今天的四段經課來重溫及探索這無盡旅途的點滴.

二．經文及思考

1．使徒行傳:聖道呼喚起航

2:14彼得和十一個使徒站起,高聲說:「猶太人和一切住在耶路撒冷的人哪,這件事你們當知道,也當側耳聽我的話.

2:36「故此,以色列全家當確實地知道,你們釘在十字架上的這位耶穌,神已經立他為主,為基督了.」

2:37眾人聽見這話,覺得扎心,就對彼得和其餘的使徒說:「弟兄們,我們當怎樣行?」

2:38彼得說:「你們各人要悔改,奉耶穌基督的名受洗,叫你們的罪得赦,就必領受所賜的聖靈；

2:39因為這應許是給你們和你們的兒女,並一切在遠方的人,就是主－我們神所召來的.」

2:40彼得還用許多話作見證,勸勉他們說:「你們當救自己脫離這彎曲的世代.」

2:41於是領受他話的人就受了洗.那一天,門徒約添了三千人,

使徒行傳第二章這個片段描述的是耶穌復活升天後門徒第一次公開傳福音的情況.為甚麼那麼多聽眾呢(2:41)?因為當時是五旬節,除了耶路撒冷本地人之外,還有旅居各地的猶太人僑胞回來過節.門徒傳講給猶太人的福音重點記載在36節:「故此,以色列全家當確實地知道,你們釘在十字架上的這位耶穌,上帝已經立他為主,為基督了.」對猶太人來說福音就是「耶穌是基督」.

他們聽到福音之後有甚麼回應呢? 「2:37眾人聽見這話,覺得扎心」,於是受了洗,加入了門徒的行列.用今天講題的話說,他們「踏上這無盡旅途」.這旅途是順是逆,是平坦還是艱辛,在那一刻,相信他們不知道,或者沒有問這個問題.

讓我們回想自己是如何踏上這旅途的呢?會不會好像當時的猶太人一樣,適逢其會聽聞福音而決志信主?對香港某一代基督徒來說,大球場舉行佈道會可能是集體回憶之一,聽說決志信主的人曾經多得站滿整個大球場草地.

有人認為信仰沒有旅途,因為做基督徒就是為了上天堂,最好等臨死之前,那麼開始就是結束了.相傳君士坦丁大帝就是其中一個例子,他一直不洗禮,免得接受信仰的要求,於臨死之前才領洗.有人等到結束才開始,有人以為開始就是結束,邱吉爾曾說:「現在並不是結束,甚至不是結束的開始,現在可能是開始的結束而已.」

你的信仰旅途起步了嗎?或許你已經觀察了一段日子,聽了不少次福音,何不踏出信心的一步呢?進入這旅途之中,才能親自體會到路上的種種,不再是聽別人說,而是親身經歷,正如經課詩篇作者的自述.對於已經起步的朋友,你的旅途經驗是怎樣的呢?有沒有為有些信徒的熱情而感到驚奇呢?是甚麼令他們在信仰路上有如此熱情呢?讓我們參考詩人的經驗.

2．詩篇:聖跡激發回應

116:1我愛耶和華,因為他聽了我的聲音和我的懇求.

116:2他既向我側耳,我一生要求告他.

116:3死亡的繩索纏繞我；陰間的痛苦抓住我；我遭遇患難愁苦.

116:4那時,我便求告耶和華的名,說:耶和華啊,求你救我的靈魂!

116:12我拿甚麼報答耶和華向我所賜的一切厚恩?

116:13我要舉起救恩的杯,稱揚耶和華的名.

116:14我要在他眾民面前向耶和華還我的願.

116:15在耶和華眼中,看聖民之死極為寶貴.

116:16耶和華啊,我真是你的僕人；我是你的僕人,是你婢女的兒子.你已經解開我的綁索.

116:17我要以感謝為祭獻給你,又要求告耶和華的名.

116:18-19我要在他眾民面前,在耶和華殿的院內,在耶路撒冷當中,向耶和華還我的願.你們要讚美耶和華!

詩人開首就示愛說(1):「我愛耶和華」,為何詩人對上帝如此熱情和坦率呢?在中間詩人又宣告(13):「我要舉起救恩的杯,稱揚耶和華的名」,到最後詩人更表示(19):「我要在他眾民面前,在耶和華殿的院內,在耶路撒冷當中,向耶和華還我的願.」為甚麼詩人如此激烈地回應呢?因為詩人不單單聽到傳道者的宣講,聽道的回應很可能只是相信,親身在生活中經驗到上帝的參與,則產生出激動及回應.詩人在甚麼情況之下經驗到上帝在他的生命中的參與呢?116:3「死亡的繩索纏繞我；陰間的痛苦抓住我；我遭遇患難愁苦」,他可能面對死亡威脅,心情被驚懼愁苦所籠罩,上帝把他從這危難和黑暗之中救拔出來.

在你的信仰旅途之中,有沒有經驗過上帝在你生命中的參與呢?有些人在病痛中經歷上帝的醫治,有些人在人生挫折中經歷上帝的扶持,有些人在人際傷痛中經歷上帝的安慰,有些人在追求公義中經歷上帝的肯定.這些經歷,激發了他們對上帝的熱情,產生出回應的行動.

經典聖詩「奇異恩典」的作者約翰牛頓(John Newton)成長於基督教背景,剛滿20歲,卻成為奴隸販賣者.1748年3月10日,23歲的牛頓從睡夢中驚醒,暴風雨撲向他的船隻.那一天死亡直視他的眼睛,和他面對面交鋒,他情不自禁的呼喊:『上帝阿,憐憫我們!如果我上得了岸,我要過不一樣的生活.』那一夜他逃脫了風暴的死劫.這經歷扭轉了他的人生,驚覺自己罪孽深重,決志重新做人,將自己獻給基督,踏上另一條人生旅途.1787年,牛頓打破沉默,出版了《論非洲奴隸貿易》(Thoughts upon the African Slave Trade),勇敢向社會坦承自己曾參與販賣黑奴.深受這本書啟發,1788年,威伯福斯提出廢奴法案,1807年,牛頓過世的前幾個月,廢奴法案終於被通過.

你可能覺得上帝在約翰牛頓的生命中的參與是奇異的,因此他的詩歌叫奇異恩典.事實上,詮釋可以是個人性的,如果當時約翰牛頓沒有把那次經歷理解為上帝的參與,而只是幸運地大能不死,他的人生旅途可能完全不同了.奇異恩典不單單是約翰牛頓的詩歌,也可以是你和我的詩歌,我們如何詮釋發生在我們生命中的事情呢?我們有沒有在人生的旅途中看見上帝的參與呢?再說一次,詮釋可以是個人性的,未必是上帝沒有介入你的生命,可能只是你沒有讓上帝參與在你對自身經歷的詮釋中.

話說回頭,有沒有發現約翰牛頓在1748年經歷上帝的拯救,而在1787年才公開向社會承認自己在黑奴販賣的參與呢?其實這不奇怪,因為生命需要時間成長,信仰旅途是一步一步向前的；以為走入電話亭就立即變成超人,往往是危機的開始.曾幾何時,有些團體安排信主不久的藝人分享見證,藉他們的名氣吸引觀眾.但這些藝人後來怎麼樣了?如果信仰的旅途是無盡的,就不用爭朝夕,也不用急於一時,慢工出細貨,炆火烹出好味道.詩人展示了信仰的熱情,接下來的彼得前書帶引我們進入信仰的深度行.

3．彼得前書:聖言默化生命

1:17你們既稱那不偏待人,按各人行為審判人的主為父,就當存敬畏的心度你們在世寄居的日子,

1:18知道你們得贖,脫去你們祖宗所傳流虛妄的行為,不是憑著能壞的金銀等物,

1:19乃是憑著基督的寶血,如同無瑕疵,無玷污的羔羊之血.

1:20基督在創世以前是預先被上帝知道的,卻在這末世才為你們顯現.

1:21你們也因著他,信那叫他從死裡復活,又給他榮耀的上帝,叫你們的信心和盼望都在於上帝.

1:22你們既因順從真理,潔淨了自己的心,以致愛弟兄沒有虛假,就當從心裡彼此切實相愛.

1:23你們蒙了重生,不是由於能壞的種子,乃是由於不能壞的種子,是藉著上帝活潑常存的道.

「1:22潔淨了自己的心」,如何能夠呢?「1:22沒有虛假」,如何做到呢?前幾天在社交網絡收到一幅漫畫:老師訓示學生說:「你天天說謊,將來能幹甚麼呢?」學生回答:「當領導!」潔淨心靈,修養品格,不是舉手決志就有,甚至不是某次特別經歷就得到的.相信你會同意,不是自稱基督徒的人就擁有基督徒的品格.俗語有云:「學壞三日,學好三年」,正如詩歌所說:「Take time to be holy」.彼得前書告訴我們,生命的轉化是靠那顆放在我們生命中的種子.隨著種子茁壯成長,我們的生命就有所不同,正如古人說「有諸內必而形諸外」.這種子靠甚麼而生長呢?彼得說是「藉著上帝活潑常存的道」.彼得的提醒,呼應著約翰福音15章所記載耶穌的教導:

15:7你們若常在我裡面,我的話也常在你們裡面,凡你們所願意的,祈求,就給你們成就.

15:8你們多結果子,我父就因此得榮耀,你們也就是我的門徒了.

果子是甚麼呢?使徒保羅在加拉太書如此說:

5:22聖靈的果子就是仁愛,喜樂,和平,忍耐,恩慈,良善,信實,溫柔,節制.

在信仰的旅途上,聖道呼換我們起行,聖跡激發我們回應,還有聖言默化我們的生命.保羅說,聖經一方面「教訓,督責,使人歸正」,另一方面「教導人學義,使人得以完全,預備行各樣的善事」.越行前一步,你越發現旅途的豐富及精彩.不但轉化你的生命,更開你視域,讓你看見奇妙的境界.

4．路加福音:聖禮察亮眼目

24:13正當那日,門徒中有兩個人往一個村子去；這村子名叫以馬忤斯,離耶路撒冷約有二十五里.

24:14他們彼此談論所遇見的這一切事.

24:15正談論相問的時候,耶穌親自就近他們,和他們同行；

24:16只是他們的眼睛迷糊了,不認識他.

24:17耶穌對他們說:「你們走路彼此談論的是甚麼事呢?」他們就站住,臉上帶著愁容.

24:18二人中有一個名叫革流巴的回答說:「你在耶路撒冷作客,還不知道這幾天在那裡所出的事嗎?」

24:19耶穌說:「甚麼事呢?」他們說:「就是拿撒勒人耶穌的事.他是個先知,在神和眾百姓面前,說話行事都有大能.

24:20祭司長和我們的官府竟把他解去,定了死罪,釘在十字架上.

24:21但我們素來所盼望,要贖以色列民的就是他!不但如此,而且這事成就,現在已經三天了.

24:22再者,我們中間有幾個婦女使我們驚奇；她們清早到了墳墓那裡,

24:23不見他的身體,就回來告訴我們,說看見了天使顯現,說他活了.

24:24又有我們的幾個人往墳墓那裡去,所遇見的正如婦女們所說的,只是沒有看見他.」

24:25耶穌對他們說:「無知的人哪,先知所說的一切話,你們的心信得太遲鈍了.

24:26基督這樣受害,又進入他的榮耀,豈不是應當的嗎?」

24:27於是從摩西和眾先知起,凡經上所指著自己的話都給他們講解明白了.

24:28將近他們所去的村子,耶穌好像還要往前行,

24:29他們卻強留他,說:「時候晚了,日頭已經平西了,請你同我們住下吧!」耶穌就進去,要同他們住下.

24:30到了坐席的時候,耶穌拿起餅來,祝謝了,擘開,遞給他們.

24:31他們的眼睛明亮了,這才認出他來.忽然耶穌不見了.

24:32他們彼此說:「在路上,他和我們說話,給我們講解聖經的時候,我們的心豈不是火熱的嗎?」

24:33他們就立時起身,回耶路撒冷去,正遇見十一個使徒和他們的同人聚集在一處,

24:34說:「主果然復活,已經現給西門看了.」

24:35兩個人就把路上所遇見,和擘餅的時候怎麼被他們認出來的事,都述說了一遍.

剛才我們提過,上帝有沒有參與在我們的生命中,視乎我們有沒有發現祂的參與.有些時候,上帝親臨,我們也不一定看見的,正如路加福音記載的門徒一樣:

24:15正談論相問的時候,耶穌親自就近他們,和他們同行；

24:16只是他們的眼睛迷糊了,不認識他.

他們甚麼時候才看見那位與他們談話的就是耶穌呢?

24:30到了坐席的時候,耶穌拿起餅來,祝謝了,擘開,遞給他們.

24:31他們的眼睛明亮了,這才認出他來.

從那兩個門徒告訴其他人的的說話也證實了這要點:

24:35兩個人就把路上所遇見,和擘餅的時候怎麼被他們認出來的事,都述說了一遍.

「拿起餅來,祝謝了,擘開」這個行動,對那兩個門徒有非常特別的記憶.這個行動成為了教會的聖禮.以馬忤斯路上的門徒,在耶穌向他們講解聖經的時候也認不出他,直至耶穌擘餅的時候認出他來,這真是一個奧秘.「聖禮」就是奧秘,sacraments一詞,是從拉丁文 sacramentum 演變而來的,意思是指著基督教一切奇妙與難以了解的事,以此字來代表希臘文的「奧秘」(mystery).在信仰的旅途中,有沒有想過聖禮可以開啟我們的視域,能夠看見奧秘,能夠看見更真實,更美好,更長遠的景象呢?

三．總結:無盡的旅途

C.S. Lewis在《納尼亞傳奇7:最後一戰》以豐富的想像力描繪出這旅途中非常吸引的一幕:

(在一個小小的茅草屋頂的馬廄裡)「他們從門裡向外望到的景像是夠奇怪的了.但,比上述任何景象更奇怪的是,他們在門內向四周打量,竟發現自己置身於溫暖的白晝,蔚藍的天空在他們的頭上,繁花在他們的腳邊,笑意在阿斯蘭的眼睛裡.阿斯蘭迅速轉過身來,蹲得更低,用尾巴甩打自己的身體,然後像金箭似的躥出去了.

「來,朝更深處跑!來,朝更高處跑!」阿斯蘭回過頭來喊道.他們跟著阿斯蘭走去.邁向的地方,大家有熟識的感覺,但又說不出來.

把大家的感受概括起來的,是獨角獸:「我終於到家了!這是我的真正的國土!我屬於這兒.這是我生平一直在尋找的國土,儘管我直到現在才知道它.為什麼我們愛老的納尼亞呢?理由是它有時候看上去有點兒像這個新的納尼亞.」

卻說這時出現一樁非常奇怪的事情.其他的人和獸也都開始奔跑了,使他們自己大吃一驚的是,他們都能趕得上獨角獸:不僅狗兒和人,而且連胖胖的小驢子和矮腿小矮人波金也都趕得上.風猛吹在他們的臉上,彷彿他們是在一輛沒有擋風玻璃的,疾馳如飛的汽車裡,鄉村在他們的身旁飛快地掠過.他們愈跑愈快,好像飛翔一樣,但沒有一個感到熱不可耐,疲倦或喘不過氣來.

「我明白了,」露茜終於沉思地說道,「現在我明白了.這花園就像那馬廄.裡邊遠比外邊大得多.」

「當然啦,夏娃的女兒,」羊怪說道,「你愈是往高處深處走去,一切東西就變得愈大.裡邊比外邊大.」

露茜仔細打量著花園,發覺它確實壓根兒不是一個花園,而是一個完全的世界,有它自己的江湖,森林,海洋和山嶺.但它們並不陌生:她一一都認識它們.

「我明白了,」她說道,「這仍舊是納尼亞,它比馬廄門外的納尼亞更真實更美麗!我明白了,世界中的世界,納尼亞里的納尼亞……」

這是不斷發現豐富,精彩,令人期盼的無盡旅途.



\newpage

\section{更大的事}
\label{sec:500}
\textbf{劉國偉先生 - 更大的事}
\newline
\newline
\hyperref[sec:499]{< < < PREV SERMON < < <}
~
\hyperref[sec:toc]{[返目錄]}
~
\hyperref[sec:501]{> > > NEXT SERMON > > >}


講題:更大的事 Even Greater Things

講員:劉國偉先生

場合:香港中文大學崇基學院禮拜堂主日崇拜

日期:2020年5月10日

引言

各位弟兄姊妹平安.

今天是復活期第五主日,轉眼之間教會年曆即將過半,再過三個星期(5月31日)就是五旬節,亦即聖靈降臨主日,開展教會年曆的下半年.我們可以概括地說,由將臨期到復活期,是有關耶穌基督的降生,職事,受死,復活和升天,構成教會年曆的上半；由聖靈降臨節到基督君王日,是有關聖靈的引導,是下半年.臨近教會年曆的中間,是一個合適的時間,去反思基督為我們犧牲,復活,所帶來的救贖,對於我們每天每夜的生活有什麼意義,同時也反省在聖靈的引導和充滿下,我們又可以為拯救我們的神,在生命的祭壇上,如何以生命活祭獻給主.

今天亦是社會人士為母親慶祝母親節的日子.今年的母親節在全球疫情的困難光景下,顯得額外特別,對身為母親的女士們,或者對於作為子女的,都是一個特別的考驗,但願每個家庭親子相愛,跟從主的我們從聖經中得著亮光,從聖靈得著力量,知所方向,亦知所進退,或順或逆,都為基督而活.

今天我們分享的講題是「更大的事」,我們主要思想的經課,是今天誦讀的三代經文中的福音書,亦即約翰福音14章1至14節.這段經文的背景發生於主耶穌受難之前,是聖週的星期四晚上,當我們在復活期第五主日誦讀這段經文的時候,就更容易意會原來由聖週直至五旬節,門徒都身陷一種困惑,憂愁,甚至恐懼的情緒當中,就是在這情景底下,耶穌對門徒說:「你們心裡不要憂愁」.在這時刻,我們即將迎向聖靈降臨節的時候,再一次思想耶穌基督向門徒所說:「我所做的事,信我的人也要做,並且要做比這更大的事」到底蘊含什麼信仰和生活的意義.

第四福音的文學佈局

從一個較為宏觀的角度,去分析第四福音的文學佈局和結構,我們可以把整卷福音書分成五大部份.

第一部份是「序言」,記載於第1章1至18節,以「道成肉身」作為主題,為整卷福音書定調.

第二部份記載於1章19節至51節,當中主要涉及有關施洗約翰的記載,和主耶穌呼召門徒的記事,當中包括記敘耶穌與拿但業相遇的一個片段.

第三部份由2章1節開始直至12章50節,當中記載了耶穌基督的傳道生涯,包括祂的職事和教導,約翰福音所記述主耶穌在受難前的七個神蹟,就是在這個大段落中.

第四部份由13章1節至20章31節,記載了有關耶穌基督在聖週的事蹟,並且祂的受死,復活和顯現.當中有特別長的篇幅,記述聖週星期四晚上,耶穌與門徒聚餐及相關的談論和祈禱,今天選讀的經文,就是在這個大段落中.

第五部份是整卷書的「跋」,有若「後記」,記載於21章1至25節,是有關主耶穌在提比哩亞海邊向門徒顯現的事蹟.

綜合來說,這五大部份當中的第二部份「耶穌呼召門徒」,第三部份「耶穌的傳道職事」及第四部份「有關聖週及復活的記載」,可以說構成了約翰福音敘事的主要內容.原來在約翰福音14章12節耶穌基督所說「更大的事」這個片語,是約翰福音獨家採用,而除了14章12節之外,分別在1章50節和5章20節都有出現.出現三次的「更大的事」這個片語,剛好散落於約翰福音核心敘事的上述三個部份,即是「耶穌呼召門徒」,「耶穌的傳道職事」及「有關聖週及復活的記載」,帶著前後呼應的文學佈局.這個觀察結果,叫我們理解及接受這是作者的佈局安排,卻很難理解為這是偶然意外,沒有任何寫作動機的寫作結果.

讓我們進入經文之中,了解這三次出現「更大的事」的相關敘事.

更大的事

我聽過解釋約翰福音14章12節「我所做的事,信我的人也要做,並且要做比這更大的事」,主要從字面的意思去處理,其中一種較流行的講法,是耶穌一生都只是在巴勒斯坦遊走,但祂的門徒則把福音傳至四方,在地域範圍來說,門徒比耶穌所作的更大；另一種亦頗常見的傳釋,是門徒所接觸的人數,以至歸主人數,比耶穌所做的在數目上更大.

「更大」在約翰福音採用的希臘原文是meìzona,其實這個字在這段落中很難翻譯成「更大的數量」,雖然meìzona這個字,有時可以用作表達數量更多；約翰福音採用這個字的時候,側重於更好的品質,和更重要的意思,例如耶穌說「我有比約翰更大的見證」(約5:36)；又或者猶太人問耶穌「難道你比我們的祖宗亞伯拉罕還大麼?」(約8:53)若果按這方向思考,meìzona所指向的事,看來不會是門徒比耶穌所做的事,覆蓋地域更大,又或者歸信主的人更多.

第一次出現「更大的事」這片語記載於約翰福音1章50節,我們從47節開始閱讀:

耶穌看見拿但業來,就指著他說:「看哪,這是個真以色列人,他心裡是沒有詭詐的.」拿但業對耶穌說:「你從哪裡知道我呢?」耶穌回答說:「腓力還沒有招呼你,你在無花果樹底下,我就看見你了.」拿但業說:「拉比,你是上帝的兒子,你是以色列的王!」耶穌對他說:「因為我說『在無花果樹底下看見你』,你就信嗎?你將要看見比這更大的事」；又說:「我實實在在地告訴你們,你們將要看見天開了,上帝的使者上去下來在人子身上.」

正如先知撒迦利亞的寫作:「......我必使我僕人大衛的苗裔發出......當那日,你們各人要請鄰舍坐在葡萄樹和無花果樹下.這是萬軍之耶和華說的.」 (亞3:8下及3:10) 猶太人此後聚集在葡萄樹及無花果樹底下,進行默想,研討有關彌賽亞來臨的美景.耶穌當然知道拿但業在無花果樹底下,盼望彌賽亞來臨的虔敬心靈,所以耶穌說出「你在無花果樹底下,我就看見你了」,其實是直接抓住拿但業的內心深處.

耶穌對拿但業說:「你將要看見比這更大的事」,那麼,這「更大的事」是指向什麼?使徒約翰記下了這個答案:「你們將要看見天開了,上帝的使者上去下來在人子身上.」我們知道創世記記載了族長雅各的一個夢,他夢見一個由地通天的梯子,有神的使者在梯子上去下來.(創28:19)與這段舊約敘事比較,耶穌把自己比喻作梯子,正如祂所說:「我就是道路,真理,生命；若不藉著我,沒有人能到父那裡去.」(約14:6)祂就是人神之間的一條通路,彌賽亞帶來的真正美景,就是人類可以透過耶穌,去到父那裡.用我們今天簡單的說法,耶穌成就救恩,我們看見天開了,人神之間重新有生命的連結,這就是耶穌向拿但業說的「更大的事」.

耶穌與拿但業的對話,從文學的角度來說,埋下了兩個非常重要的伏筆,一個是拿但業所說「你是以色列的王」,第二個就是「更大的事」.兩個伏筆,都在約翰福音敘事的第四部份有關耶穌基督在聖週,以至受死,復活的記事中再次出現.

第二次出現「更大的事」這片語記載於約翰福音5章20節,我們從19節讀起.

耶穌對他們說:「我實實在在地告訴你們,子憑著自己不能做甚麼,惟有看見父所做的,子才能做；父所做的事,子也照樣做.父愛子,將自己所做的一切事指給他看,還要將比這更大的事指給他看,叫你們希奇.父怎樣叫死人起來,使他們活,子也照樣隨自己的意思使人活.」(約5:19-21)

耶穌這段說話之前,正正就是他在安息日於耶路撒冷的畢士大池,醫治了一個患病三十八年的人,猶太人知道了,就來逼迫耶穌,因為他們不滿意耶穌干犯了安息日.耶穌回應:「我父作事直到如今,我也作事.」(約4:17)

就是在這個背景底下,耶穌提出了「更大的事」.「更大的事」在這段經文中的含義,非常清晰,那就是「父怎樣叫死人起來,使他們活,子也照樣隨自己的意思使人活.」那是清楚指向耶穌基督叫人出死入生的救恩.

耶穌所作的事工,重點不在於如何遵守安息日,也不在於有沒有神蹟奇事,真正的核心乃是人子叫死人起來,隨自己的意思使人活過來.

沒有比這更大的事

「更大的事」出現於約翰福音記事的第二個大段落.就是耶穌呼召門徒的記事,弟兄姊妹,你有沒有想過,耶穌呼召門徒去跟從祂,那個核心使命就是要門徒親身經歷和見證這更大的事.耶穌就是道路,他成為那通天的梯子,叫人可以與上帝生命相連,與神復和.弟兄姊妹,這是不是你回應召命的根本和初心?還是在你籌謀的人生中,有另外的想法,你最想見證的,並不是耶穌所指的「更大的事」?

「更大的事」也出現於約翰福音的第三個大段落中,那是耶穌基督的職事和教導.耶穌基督的教訓,言論,與及祂所行的神蹟奇事,祂周遊四方,對人施予憐憫,我們不禁深深仰慕,對祂所說所做覺得精彩異常,甚至眼花撩亂,不過,在我們對耶穌的職事和教導深深佩服之餘,有沒有看到整個事工的核心,乃在於耶穌基督叫死人活過來的奇妙恩典?在這個充滿黑暗權勢,甚至嗅聞到死亡氣味的世代,我們會否因為耶穌基督的救恩,而叫你真的感到心裡希奇?

我們回到約翰福音第十四章.作者告訴我們,門徒當時心裡憂愁,完全不能掌握耶穌基督所分享和教導的信息；多馬問耶穌:「我們不知道你往那裡去,怎麼知道那條路呢?」(v.5) 腓力又對耶穌說:「求主將父顯給我們看,我們就知足了.」(v.8)很明顯的,門徒不單憂愁,其實他們感到困惑,在劣勢之中,對未來毫無把握,耶穌就回應腓力說:「我與你們同在這樣長久,你還不認識我嗎?」(v.9上)弟兄姊妹,我把耶穌這個問題,套用在我們身上:「耶穌與我們同在這麼長久,我們還不認識耶穌嗎?」耶穌向門徒說:「你們當信我,我在父裡面,父在我裡面；即或不信,也當因我所做的事信我.」(v.11)弟兄姊妹,但願我們擦亮眼睛,在我們的生命經歷中,好好體察耶穌為我們所作的事,然後,我們憑著信心,準備要承受耶穌給予我們的應許:「我所做的事,信我的人也要做,並且要做比這更大的事,因為我往父那裡去.」(v.12) 基督往父那裡去,就差派聖靈保惠師與門徒同在,五旬節到臨,門徒被聖靈充滿,教會開展了劃時代的一天,他們確實做了比耶穌基督三年傳道職事更大的事,重點不在於地域更大,或者數量更多,而是他們認定自己是被上帝所呼召的人,以基督的救恩為人生更大的事:若人生中真有些什麼大事,那耶穌基督是道路,真理,生命,就必然是比所有大事更大的事,而且沒有比這更大的事.

有人渴望生活平安,身心康泰,家庭和諧,有人追求社會公義,有人努力尋索學問,甚至有人追求神蹟奇事,這些人生大事,都是人生的重要事情,但那一群在五旬節被聖靈充滿的使徒,好像司提反一樣,認定了更大的事,甚至為所信,為基督丟棄萬事.

弟兄姊妹,當生活翻起暗湧,風浪仍然在前面,道路仍然好像看不通透,讓我們以今天誦讀的彼得前書的一節金句,作為彼此的勉勵:「唯有你們是被揀選的族類,是有君尊的祭司,是聖潔的國度,是屬神的子民,要叫你們宣揚那召你們出黑暗入奇妙光明者的美德.」(彼前2:9)彼得勉勵跟隨耶穌的人,應該怎樣回應基督叫人出死入生的救恩,就是約翰福音記敘的「更大的事」:「叫你們宣揚那召你們出黑暗入奇妙光明者的美德.」

弟兄姊妹,五旬節就在前頭,深願我們與聖靈同工,讓靈風吹起,為主作更大的事.願榮耀歸給三一真神.



\newpage

\section{古來聖賢皆寂寞}
\label{sec:501}
\textbf{鄧瑞強博士 - 古來聖賢皆寂寞?}
\newline
\newline
\hyperref[sec:500]{< < < PREV SERMON < < <}
~
\hyperref[sec:toc]{[返目錄]}
~
\hyperref[sec:502]{> > > NEXT SERMON > > >}


約翰福音 14:15-21
\newline
\begin{longtable}{cl}
\hline
\hline
章節 & 經文 (和合本修訂版)\\
\hline
14:15 & \begin{tabularx}{0.7\textwidth}{X} 「你們若愛我，就會遵守我的命令。 \end{tabularx} \\ \\ \relax
14:16 & \begin{tabularx}{0.7\textwidth}{X} 我要求父，父就賜給你們另外一位保惠師，使他永遠與你們同在。 \end{tabularx} \\ \\ \relax
14:17 & \begin{tabularx}{0.7\textwidth}{X} 他就是真理的靈，是世人不能接受的。因為他們既看不見他，也不認識他；你們卻認識他，因他常與你們同在，也要在你們裡面。 \end{tabularx} \\ \\ \relax
14:18 & \begin{tabularx}{0.7\textwidth}{X} 我不會撇下你們為孤兒，我必到你們這裡來。 \end{tabularx} \\ \\ \relax
14:19 & \begin{tabularx}{0.7\textwidth}{X} 再過不久，世人不再看見我，你們卻會看見我，因為我活著，你們也要活著。 \end{tabularx} \\ \\ \relax
14:20 & \begin{tabularx}{0.7\textwidth}{X} 到那日，你們就會知道我在父裡面，你們在我裡面，我也在你們裡面。 \end{tabularx} \\ \\ \relax
14:21 & \begin{tabularx}{0.7\textwidth}{X} 有了我的命令而又遵守的人，就是愛我的；愛我的人，我父要愛他，我也要愛他，並且要親自向他顯現。」 \end{tabularx} \\ \\
[1ex]
\hline
\hline
\end{longtable}
講題:古來聖賢皆寂寞? Are all sages lonely?

經文:約翰福音14:15-21

講員:鄧瑞強博士

場合:香港中文大學崇基學院禮拜堂主日崇拜

日期:2020年5月17日

各位弟兄姊妹,早安,願各人平安.

疫症仍在,我們還不能聚集在一起崇拜,但無論各人身處何方,我們都能在聖靈的感通裡,回到基督的生命中；透過基督,將生命獻給天父.

我們這禮拜堂,位於香港中文大學崇基學院裡.崇基學院的大門口,有一對聯,寫著:「崇高惟博愛　本天地立心　無間東西　溝通學術；基礎在育才　當海山勝境　有懷胞與　陶鑄人群」.

上聯講「心」,講「本體」,講「內在精神」.上聯寫著「本天地立心」,乃溯源天地的本心,開示其內容.天地的本心的內容,便是「博愛」(即第一節的「崇高惟博愛」),而「博愛」的精神,乃一切教育的核心價值.現代教育的墮落,在於將教育變成「投資」,變成「生意」,變成「工具」,失去靈魂的價值.

下聯講「境」,講「作用」,講「外顯功能」.下聯寫著「當海山勝境」,這不單寫崇基學院的校園被山水環抱,也寫「天地之本心」外顯為「海山」一樣的胸襟.唯有如此胸襟,才能「有懷胞與　陶鑄人群」.「有懷胞與」這講法,出自北宋儒學大師張載(張橫渠)的「民吾同胞,物吾與也」,即後來講的「民胞物與」,視一切人為自己的同胞,視一切物為自己的同類(「與」即「同類」,「物與」即「視物為同類」).有天地之本「心」,便有這樣的生命「境」界.教育的功用,是成就一大愛的人.

關於這位中國哲人張載先生,我會多講幾句,以帶入我今日的講道內容.在中國哲學史上,張載的一項思想貢獻,是他的「氣論」.「氣」,天地有正「氣」的「氣」.傳統儒家講「性善」,但「性善」的根源在哪裡?孟子說:「此天之所與我者.」(《告子上》)「性善」的根源在於「天」.問題是:人在這邊,天在那邊,「人」如何明白「天」?張載的「氣論」正是解決這問題.張載說,萬物的本源就是「氣」,氣是創造的生機,氣也是萬物的底蘊.氣有流動,有聚散,這形成人與萬物.人與萬物都是「氣」的外現,就這而論,人與他人,人與萬物,皆是一體,皆是一氣,此即所謂「民胞物與」.人若活出「民胞物與」的大愛,就體現出這天地之正氣,就在人生中展示終極的意義.張載在萬事萬物中看到那創生萬物的「氣」,而渴望在自己的生命中,被這「氣」充滿,好活出創發生命的大愛,展示萬物的一體,而這就是人生在世的意義.

「氣」這個字,在聖經裡,就是「靈」這個字.無論是舊約希伯來文的「ruach」,抑或新約希臘文的「pneuma」,中文聖經譯作「靈」,但其本意都有「氣」這個含意.關於耶穌的降生,馬太福音和路加福音說,是「聖靈」感孕馬利亞懷孕,「靈」即「氣」,是天地的「靈氣」聚焦起來,透過耶穌這個人物,展現其內容.

天地的靈氣的內容是什麼?

在路加福音裡,耶穌出道的第一篇講章,就是引用以賽亞書,說:「主的靈(即:氣)在我身上,因為祂用膏膏我,叫我傳福音給貧窮的人；差遣我報告:被擄的得釋放,瞎眼的得看見,叫那受壓制的得自由,報告神悅納人的禧年.」耶穌然後說:「今天這經應驗在你們耳中了.」(路4: 18-19, 21)

靈氣在耶穌身上,這靈氣的內容就是耶穌日後的行動.天地的靈氣支配著耶穌的行動,這行動就是對弱者和受苦者的大愛.

今日的講道經文,是今日的福音書經課:約翰福音14:15－21.

約翰福音的神學角度,與馬太福音和路加福音的神學角度不同,不在天地的靈氣聚焦在耶穌身上,而在耶穌散發出天地的靈氣.用神學的技術語言來說,馬太和路加是從「聖靈論」去談論「基督論」,約翰是從「基督論」去開展「聖靈論」.現在,我們放下這些神學觀念,看一看福音書的經文.

約14:15　「你們若愛我,就必遵守我的命令.

約14:16　我要求父,父就另外賜給你們一位保惠師〔或譯:訓慰師；下同〕,叫他永遠與你們同在,

約14:17　就是真理的聖靈,乃世人不能接受的；因為不見他,也不認識他.你們卻認識他,因他常與你們同在,也要在你們裡面.

約14:18　我不撇下你們為孤兒,我必到你們這裡來.

約14:19　還有不多的時候,世人不再看見我,你們卻看見我；因為我活著,你們也要活著.

約14:20　到那日,你們就知道我在父裡面,你們在我裡面,我也在你們裡面.

約14:21　有了我的命令又遵守的,這人就是愛我的；愛我的必蒙我父愛他,我也要愛他,並且要向他顯現.」

約14:15　「你們若愛我,就必遵守我的命令.

開宗明義,「信耶穌」,或「跟從耶穌」,其重點是「愛耶穌」,所謂「愛耶穌」,就是渴望與耶穌的生命合一.如何與耶穌的生命合一?就是想耶穌所想的,追求耶穌所追求的,實踐耶穌所實踐的,簡言之,就是活他的生命,就是遵守他的命令,就是讓耶穌的生命透過自己流露出來.

現代人信耶穌信得太容易了,舉一舉手,決一決志,返一返教會,便說:我信了耶穌了.這種生命可能與耶穌的生命南轅北轍.聽聽他如何祈禱便知道了.若他只是按著世俗的渴望求耶穌應承這樣,應承那樣,沒有放下自己的心念而祈求力量去履行神的旨意,則這個人仍然不明白耶穌,與耶穌無關係,離耶穌很遠.信耶穌,就是愛耶穌,就是實踐耶穌的命令,活耶穌的生命.

約14:16　我要求父,父就另外賜給你們一位保惠師〔或譯:訓慰師；下同〕,叫他永遠與你們同在,

約14:17　就是真理的聖靈,乃世人不能接受的；因為不見他,也不認識他.你們卻認識他,因他常與你們同在,也要在你們裡面.

耶穌講這番話的時候,是最後晚餐,是離世的前夕.當他被釘十字架的時候,人類所有的神學問題便會湧出來.有神嗎?神還在嗎?神是仁愛的嗎?有天理嗎?神盲了嗎?神殘廢了嗎?生命有意義嗎?人生還有希望嗎?我們還能渴望什麼嗎?真有解救的可能嗎?邪惡是否就是大地的主宰?歷史的前進是否就只剩下殘暴的痕跡?

在一切疑問的前夕,耶穌安靜的說.

「我要求父」,天父還在,天理還在,世界還未失控,天父仍在掌控.

耶穌求什麼?

他向天父求聖靈,求一浩然的靈氣.

天父會另外賜給你們一位保惠師,叫他永遠與你們同在.這裡講「另一位」保惠師,意思是耶穌本身就是頭一位保惠師.耶穌的生命會以另一形式接續下去.保惠師,希臘文是「paraklētos」,意思是:為你辯護的律師,代求者,安慰者,幫助者.耶穌離開這個世界後,他的生命所散發的「靈氣」,會永留人間,成為人永恆的扶助者.

這「保惠師」,扶助者是誰?

是「真理的聖靈」,是將我們帶入真理裡的「靈氣」.

天父的心,顯明在耶穌的生命中.耶穌的十架行動,完全將天父的心顯明出來.耶穌的十架行動後,世人再看不見他,但他卻留下永恆的靈氣.這「靈氣」當然不只是耶穌的影響力,這「靈氣」有神性的生命力.這靈氣瀰漫大地,幫助人明白耶穌,愛耶穌,活出耶穌的生命,體會耶穌生命中所顯露的天心,讓人進入耶穌所展示的真理裡.

這「靈氣」感化世人的工作,是艱巨無比的.反對耶穌的,也會反對這聖靈.經文說:「〔這〕真理的聖靈,乃世人不能接受的；因為不見他,也不認識他.」看不到耶穌的十架意義的,當然也看不到耶穌所流露的靈氣.愛耶穌的,活出耶穌生命的,「卻認識他〔這靈氣〕,因他常與你們同在,也要在你們裡面」.用心實踐耶穌生命的人,體會這靈氣充溢著自己的生命,支撐著自己.

約14:18　我不撇下你們為孤兒,我必到你們這裡來.

約14:19　還有不多的時候,世人不再看見我,你們卻看見我；因為我活著,你們也要活著.

耶穌不撇下我們為孤兒,所謂「孤兒」,即無父者,無人扶助者,無人保護者.耶穌所散發的天地之靈氣,會永遠地支撐著愛耶穌的人.今次講道的題目是「古來聖賢皆寂寞?」古來聖賢,是否都是無所依靠的寂寞人?耶穌說,聖賢不是寂寞人,他們是被天地的靈氣感召,被這靈氣承托的人.他們永遠不是一個人活著,而是在「靈氣」的陪伴下活著.

耶穌離世後,世人再看不見他.或者說,在人類殘酷的人性表露無違時,在暴力宣示主權時,在絕望橫行大地時,世人真的看不見神,看不見真理,看不見希望.但是,總有人在耶穌留下的靈氣裡,體會到創生萬物的大愛力量.經文說:「你們卻看見我；因為我活著,你們也要活著」.耶穌的靈氣還在,我們在這靈氣中,明白天心,明白基督的可貴,明白大愛的力量,明白人生的意義.

約14:20　到那日,你們就知道我在父裡面,你們在我裡面,我也在你們裡面.

約14:21　有了我的命令又遵守的,這人就是愛我的；愛我的必蒙我父愛他,我也要愛他,並且要向他顯現.」

結尾的21節,呼應開頭的15節,講到愛耶穌的人,就是實踐耶穌命令的人,就是活出耶穌生命的人,這些人,最後會明白天地間的終極奧秘.

這終極奧秘是什麼?

「到那日」,就是當耶穌不在時,我們這些愛耶穌的人,因著實踐耶穌的命令,我們將耶穌留下的靈氣活出來.當我們流露基督的靈氣時,我們在自身的生命中,讓耶穌活現出來.神在,因耶穌活現在我們的生命中,這就是耶穌所說:「你們在我裡面,我也在你們裡面」以及「愛我的必蒙我父愛他,我也要愛他,並且要向他顯現」的意思.

不要問:為何神不在?

神是在的.耶穌的靈氣在吹拂,在感召人,在扶助人,幫助人明白耶穌的大愛.愛耶穌的人,在實踐中,明白耶穌的真理,也讓耶穌活現出來.耶穌的顯現,就是天心的揭露.神是在的,他在信徒的真理實踐中.信徒在實踐耶穌的命令中,將耶穌活現出來.

與其問:在這悲慘的人間,神為何不在?

不如問:我如何能更真切的實踐耶穌的命令?

唯有這真切的實踐,展露出天地間的靈氣的創生之力,讓基督活現出來,並呈現天父對人間的美意.

中國哲人張載體會到那天地間的正氣,他領悟到生命當行的事,他說:「為天地立心,為生民立命,為往聖繼絕學,為萬世開太平」.他要展現天地的本心,為世人找安身立命的根據,顯明聖人的真理,為人類的前景謀幸福.或者,用張載的講法,我們也應努力:「明天父本心,顯聖靈之力,效基督繼絕活,為天國現和平」.

願榮耀歸上主,願人間享和平.

(附記:本人在崇基禮拜堂的講章已結集成《荷花池畔春草綠》講道集,試讀本可參閱:\url{http://cc.cuhk.edu.hk/dev/images/content/fund-raising-campaigns/university-sermons-by-dr-sk-tang/CCC-SK-Tang-Sermon-Text-Sample-2.pdf}.較早期的崇基講章,也已結集成《敢於跟隨主》一書.有興趣的朋友可到各大基督教書局購買這兩書.)



\newpage

\section{作我的見證}
\label{sec:502}
\textbf{豪恩先生 - 作我的見證}
\newline
\newline
\hyperref[sec:501]{< < < PREV SERMON < < <}
~
\hyperref[sec:toc]{[返目錄]}
~
\hyperref[sec:503]{> > > NEXT SERMON > > >}


使徒行傳 1:6-14
\newline
\begin{longtable}{cl}
\hline
\hline
章節 & 經文 (和合本修訂版)\\
\hline
1:6 & \begin{tabularx}{0.7\textwidth}{X} 他們聚集的時候，問耶穌：「主啊，你就要在這時候復興以色列國嗎？」 \end{tabularx} \\ \\ \relax
1:7 & \begin{tabularx}{0.7\textwidth}{X} 耶穌對他們說：「父憑著自己的權柄所定的時候和日期，不是你們可以知道的。 \end{tabularx} \\ \\ \relax
1:8 & \begin{tabularx}{0.7\textwidth}{X} 但聖靈降臨在你們身上，你們就必得著能力，並要在耶路撒冷、猶太全地和撒瑪利亞，直到地極，作我的見證。」 \end{tabularx} \\ \\ \relax
1:9 & \begin{tabularx}{0.7\textwidth}{X} 說了這些話，他們正看的時候，他被接上升，有一朵雲彩從他們眼前把他接去。 \end{tabularx} \\ \\ \relax
1:10 & \begin{tabularx}{0.7\textwidth}{X} 他升上去的時候，他們定睛望天，看哪，有兩個人身穿白衣站在他們旁邊， \end{tabularx} \\ \\ \relax
1:11 & \begin{tabularx}{0.7\textwidth}{X} 說：「加利利人哪，你們為甚麼站著望天呢？這離開你們被接升天的耶穌，你們見他怎樣升上天去，他也要怎樣來臨。」 \end{tabularx} \\ \\ \relax
1:12 & \begin{tabularx}{0.7\textwidth}{X} 有一座山，名叫橄欖山，離耶路撒冷不遠，有安息日可行走的路程。那時，門徒從那裡回耶路撒冷去， \end{tabularx} \\ \\ \relax
1:13 & \begin{tabularx}{0.7\textwidth}{X} 他們一進城，就上了所住的樓房；在那裡有彼得、約翰、雅各、安得烈、腓力、多馬、巴多羅買、馬太、亞勒腓的兒子雅各、激進黨的西門，和雅各的兒子猶大。 \end{tabularx} \\ \\ \relax
1:14 & \begin{tabularx}{0.7\textwidth}{X} 這些人和幾個婦人，包括耶穌的母親馬利亞，和耶穌的兄弟，都同心合意地恆切禱告。 \end{tabularx} \\ \\
[1ex]
\hline
\hline
\end{longtable}
詩篇 68:1-10,32-35
\newline
\begin{longtable}{cl}
\hline
\hline
章節 & 經文 (和合本修訂版)\\
\hline
68:1 & \begin{tabularx}{0.7\textwidth}{X} 願神興起，使他的仇敵四散，使那恨他的人從他面前逃跑。 \end{tabularx} \\ \\ \relax
68:2 & \begin{tabularx}{0.7\textwidth}{X} 你驅逐他們，如煙被吹散；惡人見神的面就消滅，如蠟被火熔化。 \end{tabularx} \\ \\ \relax
68:3 & \begin{tabularx}{0.7\textwidth}{X} 惟有義人必然歡喜，在神面前快樂，他們要在喜樂中歡欣。 \end{tabularx} \\ \\ \relax
68:4 & \begin{tabularx}{0.7\textwidth}{X} 你們當向神唱詩，歌頌他的名；為那駕車經過曠野的修平道路。他的名是耶和華，你們要在他面前歡樂！ \end{tabularx} \\ \\ \relax
68:5 & \begin{tabularx}{0.7\textwidth}{X} 神在他的聖所作孤兒的父，作寡婦的伸冤者。 \end{tabularx} \\ \\ \relax
68:6 & \begin{tabularx}{0.7\textwidth}{X} 神使孤獨的有家，使被囚的出來享福；惟有悖逆的要住在乾旱之地。 \end{tabularx} \\ \\ \relax
68:7 & \begin{tabularx}{0.7\textwidth}{X} 神啊，當你走在百姓前頭，在曠野行進，（細拉） \end{tabularx} \\ \\ \relax
68:8 & \begin{tabularx}{0.7\textwidth}{X} 地見神的面就震動，天也降雨；西奈山見以色列神的面也震動。 \end{tabularx} \\ \\ \relax
68:9 & \begin{tabularx}{0.7\textwidth}{X} 神啊，你降下大雨；你的產業以色列疲乏的時候，你使他堅固。 \end{tabularx} \\ \\ \relax
68:10 & \begin{tabularx}{0.7\textwidth}{X} 你的會眾住在境內；神啊，你在恩惠中為困苦人預備所需的。 \end{tabularx} \\ \\ \relax
68:11 & \begin{tabularx}{0.7\textwidth}{X} 主發命令，傳好信息的婦女成了大群： \end{tabularx} \\ \\ \relax
68:12 & \begin{tabularx}{0.7\textwidth}{X} 「統領大軍的君王逃跑了，逃跑了！」在家等候的婦女也分得了掠物。 \end{tabularx} \\ \\ \relax
68:13 & \begin{tabularx}{0.7\textwidth}{X} 你們躺臥在羊圈，好像鴿子的翅膀鍍銀，翎毛鍍金一般。 \end{tabularx} \\ \\ \relax
68:14 & \begin{tabularx}{0.7\textwidth}{X} 全能者在境內趕散列王的時候，勢如飄雪在撒們。 \end{tabularx} \\ \\ \relax
68:15 & \begin{tabularx}{0.7\textwidth}{X} 巴珊山是極其宏偉的山，巴珊山是多峰多嶺的山。 \end{tabularx} \\ \\ \relax
68:16 & \begin{tabularx}{0.7\textwidth}{X} 你們多峰多嶺的山哪，為何以妒忌的眼光看神所願居住的山？耶和華必住這山，直到永遠！ \end{tabularx} \\ \\ \relax
68:17 & \begin{tabularx}{0.7\textwidth}{X} 神的車輦累萬盈千；主在其中，好像在西奈聖山一樣。 \end{tabularx} \\ \\ \relax
68:18 & \begin{tabularx}{0.7\textwidth}{X} 你已經升上高天，擄掠了俘虜；你在人間，就是在悖逆的人中，受了供獻，使耶和華神可以與他們同住。 \end{tabularx} \\ \\ \relax
68:19 & \begin{tabularx}{0.7\textwidth}{X} 天天背負我們重擔的主，就是拯救我們的神，是應當稱頌的！（細拉） \end{tabularx} \\ \\ \relax
68:20 & \begin{tabularx}{0.7\textwidth}{X} 神是為我們施行拯救的神；人能脫離死亡是在乎主---耶和華。 \end{tabularx} \\ \\ \relax
68:21 & \begin{tabularx}{0.7\textwidth}{X} 但神要打破他仇敵的頭，就是那常犯罪之人的頭顱。 \end{tabularx} \\ \\ \relax
68:22 & \begin{tabularx}{0.7\textwidth}{X} 主說：「我要使百姓從巴珊歸來，使他們從深海轉回， \end{tabularx} \\ \\ \relax
68:23 & \begin{tabularx}{0.7\textwidth}{X} 好叫你打碎仇敵，使你的腳踹在血中，使你狗的舌頭也有份。」 \end{tabularx} \\ \\ \relax
68:24 & \begin{tabularx}{0.7\textwidth}{X} 神啊，你是我的神，我的王；人已經看見你行走，進入聖所。 \end{tabularx} \\ \\ \relax
68:25 & \begin{tabularx}{0.7\textwidth}{X} 歌唱的行在前，作樂的隨在後，都在擊鼓的童女中間： \end{tabularx} \\ \\ \relax
68:26 & \begin{tabularx}{0.7\textwidth}{X} 「從以色列源頭而來的啊，你們當在各會中稱頌神---耶和華！」 \end{tabularx} \\ \\ \relax
68:27 & \begin{tabularx}{0.7\textwidth}{X} 在那裡，有統管他們的小便雅憫，有猶大的領袖和他們的一群人，有西布倫的領袖，有拿弗他利的領袖。 \end{tabularx} \\ \\ \relax
68:28 & \begin{tabularx}{0.7\textwidth}{X} 你的神已賜給你力量；神啊，求你堅固你為我們所成全的事！ \end{tabularx} \\ \\ \relax
68:29 & \begin{tabularx}{0.7\textwidth}{X} 因你耶路撒冷的殿，列王必帶貢物獻給你。 \end{tabularx} \\ \\ \relax
68:30 & \begin{tabularx}{0.7\textwidth}{X} 求你斥責蘆葦中的野獸和公牛群，並萬民中的牛犢。直到他們帶著銀塊來朝貢；神已經趕散好戰的萬民。 \end{tabularx} \\ \\ \relax
68:31 & \begin{tabularx}{0.7\textwidth}{X} 埃及的使臣要出來，古實人要急忙向神伸出手來。 \end{tabularx} \\ \\ \relax
68:32 & \begin{tabularx}{0.7\textwidth}{X} 地上的國度啊，你們要向神歌唱，要歌頌主，（細拉） \end{tabularx} \\ \\ \relax
68:33 & \begin{tabularx}{0.7\textwidth}{X} 就是那駕行在亙古的諸天之上的主！聽啊，他發出聲音，是極大的聲音。 \end{tabularx} \\ \\ \relax
68:34 & \begin{tabularx}{0.7\textwidth}{X} 你們要將能力歸給神；他的威榮在以色列之上，他的能力顯在天上。 \end{tabularx} \\ \\ \relax
68:35 & \begin{tabularx}{0.7\textwidth}{X} 神啊，你從聖所顯為可畏，以色列的神是那將力量權能賜給他百姓的。神是應當稱頌的！ \end{tabularx} \\ \\
[1ex]
\hline
\hline
\end{longtable}
約翰福音 17:1-11
\newline
\begin{longtable}{cl}
\hline
\hline
章節 & 經文 (和合本修訂版)\\
\hline
17:1 & \begin{tabularx}{0.7\textwidth}{X} 耶穌說了這些話，就舉目望天，說：「父啊，時候到了，願你榮耀你的兒子，使兒子也榮耀你； \end{tabularx} \\ \\ \relax
17:2 & \begin{tabularx}{0.7\textwidth}{X} 因為你曾賜給他權柄掌管凡血肉之軀的，使他把永生賜給你所賜給他的人。 \end{tabularx} \\ \\ \relax
17:3 & \begin{tabularx}{0.7\textwidth}{X} 認識你—獨一的真神，並且認識你所差來的耶穌基督，這就是永生。 \end{tabularx} \\ \\ \relax
17:4 & \begin{tabularx}{0.7\textwidth}{X} 我在地上已經榮耀你，你交給我做的工作，我已完成了。 \end{tabularx} \\ \\ \relax
17:5 & \begin{tabularx}{0.7\textwidth}{X} 父啊，現在求你使我在你面前得榮耀，就是在未有世界以前，我同你享有的榮耀。 \end{tabularx} \\ \\ \relax
17:6 & \begin{tabularx}{0.7\textwidth}{X} 「你從世上賜給我的人，我已把你的名顯明給他們。他們本是你的，你把他們賜給我，他們也遵守了你的道。 \end{tabularx} \\ \\ \relax
17:7 & \begin{tabularx}{0.7\textwidth}{X} 現在他們知道，你所賜給我的一切都是從你那裡來的； \end{tabularx} \\ \\ \relax
17:8 & \begin{tabularx}{0.7\textwidth}{X} 因為你所賜給我的話，我已經賜給他們，他們也領受了，又確實知道，我是從你出來的，並且信你差了我來。 \end{tabularx} \\ \\ \relax
17:9 & \begin{tabularx}{0.7\textwidth}{X} 我為他們祈求，不為世人祈求，卻為你所賜給我的人祈求，因他們本是你的。 \end{tabularx} \\ \\ \relax
17:10 & \begin{tabularx}{0.7\textwidth}{X} 凡是我的都是你的，你的也是我的，並且我因他們得了榮耀。 \end{tabularx} \\ \\ \relax
17:11 & \begin{tabularx}{0.7\textwidth}{X} 我到你那裡去；我不再留在世上，他們卻在世上。聖父啊，求你因你的名，就是你所賜給我的名，保守他們，使他們像我們一樣合而為一。 \end{tabularx} \\ \\
[1ex]
\hline
\hline
\end{longtable}
講題:作我的見證 Be My Witnesses

經課:使徒行傳1:6-14;詩篇68:1-10,32-35;彼得前書4:12-14,5:6-11;約翰福音17:1-11

講員:林豪恩先生

場合:香港中文大學崇基學院禮拜堂主日崇拜

日期:2020年5月24日

使徒行傳:作我的見證

1:6他們聚集的時候,問耶穌說:「主啊,你復興以色列國就在這時候嗎?」

1:7耶穌對他們說:「父憑著自己的權柄所定的時候,日期,不是你們可以知道的.

1:8但聖靈降臨在你們身上,你們就必得著能力,並要在耶路撒冷,猶太全地,和撒馬利亞,直到地極,作我的見證.」

1:9說了這話,他們正看的時候,他就被取上升,有一朵雲彩把他接去,便看不見他了.

1:10當他往上去,他們定睛望天的時候,忽然有兩個人身穿白衣,站在旁邊,說:

1:11「加利利人哪,你們為甚麼站著望天呢?這離開你們被接升天的耶穌,你們見他怎樣往天上去,他還要怎樣來.」

1:12有一座山,名叫橄欖山,離耶路撒冷不遠,約有安息日可走的路程.當下,門徒從那裡回耶路撒冷去,

1:13進了城,就上了所住的一間樓房；在那裡有彼得,約翰,雅各,安得烈,腓力,多馬,巴多羅買,馬太,亞勒腓的兒子雅各,奮銳黨的西門,和雅各的兒子猶大.

1:14這些人同著幾個婦人和耶穌的母親馬利亞,並耶穌的弟兄,都同心合意地恆切禱告.

經文記述的是死而復活的耶穌與門徒最後一次見面,這次見面之後,耶穌就駕雲上升了.最後一次會面,門徒問耶穌甚麼問題呢?這個問題對門徒來說應該是非常重要的,也應該反映出門徒的關注.

現在是五月底,大學的考試結束了,意味著畢業班的同學已經完成學業,基督徒同學問的問題也反映他們的處境:「我做甚麼工作是上主的心意呢?」或者:「上主甚麼時候才賜給我一份工作呢?」

我們的問題反映我們的關注,也反映我們的處境.過去一年,成人們祈禱求問甚麼問題呢?會不會重問二十多年前的問題:「是不是要離開香港呢?」「移居是不是上主的心意呢?」

在我們的印象中,耶穌的門徒多是漁夫及一般平凡百姓,加上一些婦女,若我們仔細留意,他們向耶穌提出的問題卻是令人們感到驚訝的:

1:6他們聚集的時候,問耶穌說:「主啊,你復興以色列國就在這時候嗎?」

那是一個怎樣的時代?復國的議題竟然成為漁夫的首要關注.那是一個怎樣的歷史處境?政治議題竟然牽動尋常百姓的心.處身於此時此地的朋友們,你想像得到嗎?我相信你的明白的.當大學生,甚至中學生放下書本而走上街頭,這是一個怎樣的時代呢?當政治爭論走入尋常百姓的飯桌而令家人反目,這是一個怎樣的社會狀況呢?

公元前63年,羅馬帝國征服了巴勒斯坦,以色列就落在羅馬的統治之下.羅馬把猶太,加利利,比利亞和特拉可尼等四個省交給以東人希律王.當時的以色列人就生存在遠有羅馬帝國的權力,近有外邦人希律王的管治之下.

大希律死後,他的三個兒子分別治理他的封地.他們在四福音書中皆稱為希律王:亞基老(4BC-6AD)分得猶太省,安提帕(4BC-39AD)佔有加利利和比利亞,腓力(4BC-34AD)得到特拉可尼.其中亞基老暴虐無道,作王十年後即被猶太人請願,遭羅馬皇帝罷黜後流放,改由羅馬派巡撫管理.

簡述這段歷史,希望我們知道門徒處身的時代,也能夠想像他們那個時代的以色列人的集體經驗和盼望,以至能夠明白他們為甚麼在重要的時刻向耶穌提出復興以色列國的問題.或許,當時的人認為,只要以色列國被復興,他們就不會再被欺壓,他們就會重拾民族自尊及自信,重建傳說中的大衛王朝.站起來,富起來,強起來,自古以來都是國家和民族的盼望,也是擁有國家和民族強烈認同的百姓的盼望,這些情感和盼望,甚至是超越宗教信仰的.

第二次世界大戰時候,德國教會為甚麼會支持希特拉呢?一方面是納粹政權全面滲透,改組和操控整個德國教會,另一方面是是德國人民的配合.他們力求走出第一次世界大戰戰敗的羞辱,恢復德國的國際地位,重新走在昔日「光輝歲月」的道路之上.希特拉及其納粹政權的政治理念,正符合人民的期望,許多基督徒也接受了,視之為恢復民族自豪與富強的最好機會.

各位弟兄姊妹,各位朋友,此時此地,甚麼是你首要的關注呢?甚麼是你追求的光輝和榮耀呢?你正為甚麼而戰呢?你關注的都可能是有基礎的,對你來說都應該是重要的；你為之而戰的,或許都是有理由的.讓我們停一停,聆聽耶穌如何回答門徒的問題,借來思考一下我們的問題和關注.

1:7耶穌對他們說:「父憑著自己的權柄所定的時候,日期,不是你們可以知道的.

耶穌沒有說門徒不應該問這問題,只是重新天父的主權.耶穌的回答,一方面表示天父並沒有抽身於歷史而讓歷史漫無目的地亂撞,另一方面肯定天父是掌管歷史的主,有路線圖及時間表,如傳道者所認信的:「上帝造萬物,各按其時成為美好」.

然而,當時當刻,對於門徒,耶穌卻另有吩咐:

1:8但聖靈降臨在你們身上,你們就必得著能力,並要在耶路撒冷,猶太全地,和撒馬利亞,直到地極,作我的見證.」

當時的門徒所領受的使命,就是作耶穌基督的見證.這使命是超越地域和國族的,無論他們身處耶路撒冷,猶太全地,和撒馬利亞,甚至直到地極,都照樣適用.

今天,畢業班同學可能問:上帝想我做甚麼工作呢?上帝安排我在甚麼地方工作呢?他們得的回應可能是:無論你做甚麼工作,無論你在哪裡工作,「作我的見證」.成人問:我應該離開還是留下呢?回應可能也是:無論你往哪裡去,「作我的見證」.使徒行傳記載了門徒在那個時代以他們的身分在他們的處境如何作上主的見證,其他人在不同的時代以不同的身分又如何作上主的見證呢?讓我們繼續閱讀其他幾段經課以作參考.

二．詩篇:見證上帝的公義及恩惠

68:32世上的列國啊,你們要向上帝歌唱；願你們歌頌主!

68:33歌頌那自古駕行在諸天以上的主!他發出聲音,是極大的聲音.

68:34你們要將能力歸給上帝.他的威榮在以色列之上；他的能力是在穹蒼.

68:35上帝啊,你從聖所顯為可畏；以色列的上帝是那將力量權能賜給他百姓的.上帝是應當稱頌的!

詩人呼籲歌頌上帝,歌唱甚麼呢?呼籲稱頌上帝,稱頌甚麼呢?

68:1願上帝興起,使他的仇敵四散,叫那恨他的人從他面前逃跑.

68:2他們被驅逐,如煙被風吹散；惡人見上帝之面而消滅,如蠟被火鎔化.

68:3惟有義人必然歡喜,在上帝面前高興快樂.

68:4你們當向上帝唱詩,歌頌他的名；為那坐車行過曠野的修平大路.他的名是耶和華,要在他面前歡樂!

詩人呼籲,要見證上帝是那位「義人必然歡喜,在上帝面前高興快樂」,「惡人見上帝之面而消滅,如蠟被火鎔化」的上帝.即是說,正義良善的人得到肯定和感到安心,奸險邪惡的人得到懲罰和感到不安.這彷彿是老生常談,不過,現實卻未必按常理展現.環顧四週,正義良善的人得到肯定嗎?能免於恐懼嗎?奸險邪惡的人得到懲罰嗎?他們會感到不安嗎?信仰上帝的個人及教會,能夠讓義人感到快樂,令惡人感到懼怕嗎?還是反過來欺負好人,懼怕惡人呢?

68:5上帝在他的聖所作孤兒的父,作寡婦的伸冤者.

68:6上帝叫孤獨的有家,使被囚的出來享福；惟有悖逆的住在乾燥之地.

68:7上帝啊,你曾在你百姓前頭出來,在曠野行走.(細拉)

68:8那時,地見上帝的面而震動,天也落雨；西奈山見以色列上帝的面也震動.

68:9上帝啊,你降下大雨；你產業以色列疲乏的時候,你使他堅固.

68:10你的會眾住在其中；上帝啊,你的恩惠是為困苦人預備的.

詩人還呼籲人們頌讚上帝甚麼呢?「作孤兒的父」,「作寡婦的伸冤者」,「叫孤獨的有家」,「使被囚的出來享福」,「以色列疲乏的時候,你使他堅固」,總括而言,詩人呼籲人們見證,上帝是那位「恩惠是為困苦人預備的」上帝,簡而言之,見證上帝是困苦人的上帝.

不過,在人人都追求成功的社會中,見證上帝是困苦人的上帝還有市場嗎?反過來說,可能宣揚上帝是成功人士的上帝會更受歡迎,正因如此,神學也要「成功」,有教會創造了「成功神學」.初期教會不是這樣的,初期教會在羅馬帝國的逼害之下,自己就是困苦的人,但他們仍然努力見證上帝是困苦人的上帝.第二世紀初士每拿主教坡旅甲勉勵信徒:「你們有能力行善,就不要遲延,因為慈惠賙濟能叫自己脫離死亡.」初期教會史家優西比烏記述,三世紀的羅馬教會供養著超過一千五百名寡婦和貧民.早期基督徒對困苦者的關懷和付出,在急難時候更顯珍貴,見證著上帝真是困苦人的上帝.

三．彼得前書:縱使受苦仍要見證

4:13倒要歡喜；因為你們是與基督一同受苦,使你們在他榮耀顯現的時候,也可以歡喜快樂.

4:14你們若為基督的名受辱罵,便是有福的；因為上帝榮耀的靈常住在你們身上.

5:6所以,你們要自卑,服在上帝大能的手下,到了時候,他必叫你們升高.

5:7你們要將一切的憂慮卸給上帝,因為他顧念你們.

5:8務要謹守,警醒.因為你們的仇敵魔鬼,如同吼叫的獅子,遍地遊行,尋找可吞吃的人.

5:9你們要用堅固的信心抵擋他,因為知道你們在世上的眾弟兄也是經歷這樣的苦難.

5:10那賜諸般恩典的上帝曾在基督裡召你們,得享他永遠的榮耀,等你們暫受苦難之後,必要親自成全你們,堅固你們,賜力量給你們.

5:11願權能歸給他,直到永永遠遠.阿們!

4:13倒要歡喜；因為你們是與基督一同受苦,使你們在他榮耀顯現的時候,也可以歡喜快樂.

《彼得前書》是寫給散居在小亞細亞北部一帶的基督徒,4:13說:「你們是與基督一同受苦」,為何他們「受苦」呢?

3:17上帝的旨意若是要你們因行善受苦,這總比因行惡受苦好.

剛才提及的詩篇經課,詩人呼籲見證上帝是那位「義人必然歡喜,在上帝面前高興快樂」,「惡人見上帝之面而消滅,如蠟被火鎔化」的上帝,在某種社會處境中,可以是一項非常大的挑戰.教會或個別信徒如果要作這樣的見證,艱難可以是相當巨大的,代價可能是非常沉重的.

上文提及的二戰時期德國基督徒,也不是鐵板一塊的.當納粹政府通過「雅利安條款」時,有些基督徒已經感到不妥當,因為條款驅逐非雅利安人,特別是猶太人離開所有公職,包括教會牧職.後來事實證明,這是排斥和屠殺猶太人的起始點.其中一位反對者潘霍華,首先被禁止出版,接著被禁止教學.他領導的地下神學院開設兩年後,也被封殺,四處遊走逃避蓋世太保的追蹤.

4:14你們若為基督的名受辱罵,便是有福的；因為上帝榮耀的靈常住在你們身上.

對當時小亞細亞的基督徒來說,又是甚麼令他們「受辱罵」呢?

4:2使你們從今以後不再隨從人的情慾,只順從上帝的旨意,在世度餘下的光陰.

4:3因為你們從前隨從外邦人的心意,生活在淫蕩,情慾,醉酒,荒宴,狂飲和可憎的偶像崇拜中,時候已經夠了.

4:4在這些事上,他們見你們不與他們同奔放蕩無度的路就以為怪,毀謗你們.

當時小亞細亞的基督徒,不是因為站出來做了甚麼轟天動地的見證而受辱罵,只是因為他們不肯站出來隨眾放蕩而己.在某些歷史時空之下,原來不「淫蕩,情慾,醉酒,荒宴,狂飲和拜偶像」,「不與他們同奔放蕩無度的路」,就足以被排斥和被惡待.

彼得明白落在如此處境中的信徒的艱辛和困惑,因此肯定他們,被辱罵,並不是做了甚麼壞事,只是「順從上帝的旨意,在世度餘下的光陰」；受苦,並不是他們做了惡事,只是惡的社會逼害好人.彼得用耶穌的經歷鼓舞他們,說明義人受逼害的情況也在他們的救主身上出現過.彼得勸勉他們,若然繼續為主作見證,必須有預備受苦的心志.

四．約翰福音:完成上主所託付的就是最榮耀的見證

17:1耶穌說了這話,就舉目望天,說:「父啊,時候到了,願你榮耀你的兒子,使兒子也榮耀你；

17:2正如你曾賜給他權柄管理凡有血氣的,叫他將永生賜給你所賜給他的人.

17:3認識你－獨一的真神,並且認識你所差來的耶穌基督,這就是永生.

17:4我在地上已經榮耀你,你所託付我的事,我已成全了.

17:5父啊,現在求你使我同你享榮耀,就是未有世界以先,我同你所有的榮耀.

17:6「你從世上賜給我的人,我已將你的名顯明與他們.他們本是你的,你將他們賜給我,他們也遵守了你的道.

17:7如今他們知道,凡你所賜給我的,都是從你那裡來的；

17:8因為你所賜給我的道,我已經賜給他們,他們也領受了,又確實知道,我是從你出來的,並且信你差了我來.

17:9我為他們祈求,不為世人祈求,卻為你所賜給我的人祈求,因他們本是你的.

17:10凡是我的,都是你的；你的也是我的,並且我因他們得了榮耀.

17:11從今以後,我不在世上,他們卻在世上；我往你那裡去.聖父啊,求你因你所賜給我的名保守他們,叫他們合而為一像我們一樣.

甚麼是最完美,最極致的見證呢?約翰福音1:18:「從來沒有人見過上帝,只有在父懷裡獨一的兒子將他表明出來.」

能夠將上帝表明出來,這應該是最完美,最極致的見證了.耶穌說:17:4「我在地上已經榮耀你」耶穌是如何榮耀了上帝呢?17:4下半節說:「你所託付我的事,我已成全了.」

有人以為,考第一拿獎學金才是美好的見證,這可能只是教育制度的遊戲；有人以為,事業有成才是美好見證,這可能只是市場的標準；有人以為,溫馨家庭生活無憂才是美好見證,這可能只是中產階級的浪漫.或許,我們應該停下來想一想耶穌的見證:「你所託付我的事,我已成全了.」上帝託付我的事是甚麼呢?我有理會過嗎?進展如何呢?有朝一日,如果我們能夠如羅一樣說出下面這段說話,就可以肯定自己已經作了美好的見證,無悔上帝賜給我們在地上的日子了:

提摩太後書4:7那美好的仗我已經打過了,當跑的路我已經跑盡了,該信的道我已經守住了.

今天福音經課是記述耶穌向天父的禱告,其中有一句這樣說:17:10「凡是我的,都是你的；你的也是我的,並且我因他們得了榮耀.」

願耶穌的禱告實現,祂因我們得到榮耀.



\newpage

\section{靈動我心力行奇事言說基督見證上帝}
\label{sec:503}
\textbf{鄭漢文博士 - 靈動我心力行奇事,言說基督見證上帝}
\newline
\newline
\hyperref[sec:502]{< < < PREV SERMON < < <}
~
\hyperref[sec:toc]{[返目錄]}
~
\hyperref[sec:504]{> > > NEXT SERMON > > >}


使徒行傳 2:1-21
\newline
\begin{longtable}{cl}
\hline
\hline
章節 & 經文 (和合本修訂版)\\
\hline
2:1 & \begin{tabularx}{0.7\textwidth}{X} 五旬節那日到了，他們全都聚集在一起。 \end{tabularx} \\ \\ \relax
2:2 & \begin{tabularx}{0.7\textwidth}{X} 忽然，有響聲從天上下來，好像一陣大風吹過，充滿了他們所坐的整座屋子； \end{tabularx} \\ \\ \relax
2:3 & \begin{tabularx}{0.7\textwidth}{X} 又有舌頭如火焰向他們顯現，分開落在他們每個人身上。 \end{tabularx} \\ \\ \relax
2:4 & \begin{tabularx}{0.7\textwidth}{X} 他們都被聖靈充滿，就按著聖靈所賜的口才說起別國的話來。 \end{tabularx} \\ \\ \relax
2:5 & \begin{tabularx}{0.7\textwidth}{X} 那時，有從天下各國來的虔誠的猶太人，住在耶路撒冷。 \end{tabularx} \\ \\ \relax
2:6 & \begin{tabularx}{0.7\textwidth}{X} 這聲音一響，許多人都來聚集，各人因為聽見門徒用他們各自的鄉談說話，就甚納悶， \end{tabularx} \\ \\ \relax
2:7 & \begin{tabularx}{0.7\textwidth}{X} 都詫異驚奇說：「看哪，這些說話的不都是加利利人嗎？ \end{tabularx} \\ \\ \relax
2:8 & \begin{tabularx}{0.7\textwidth}{X} 我們每個人怎麼聽見他們說我們生來所用的鄉談呢？ \end{tabularx} \\ \\ \relax
2:9 & \begin{tabularx}{0.7\textwidth}{X} 我們帕提亞人、瑪代人、以攔人，和住在美索不達米亞、猶太、加帕多家、本都、亞細亞、 \end{tabularx} \\ \\ \relax
2:10 & \begin{tabularx}{0.7\textwidth}{X} 弗呂家、旁非利亞、埃及的人，並靠近古利奈的利比亞一帶地方的人，僑居的羅馬人， \end{tabularx} \\ \\ \relax
2:11 & \begin{tabularx}{0.7\textwidth}{X} 包括猶太人和皈依猶太教的人， 克里特人和阿拉伯人，都聽見他們用我們的鄉談講論神的大作為。」 \end{tabularx} \\ \\ \relax
2:12 & \begin{tabularx}{0.7\textwidth}{X} 眾人就都驚奇困惑，彼此說：「這是甚麼意思呢？」 \end{tabularx} \\ \\ \relax
2:13 & \begin{tabularx}{0.7\textwidth}{X} 還有人譏誚，說：「他們是灌滿了新酒吧！」 \end{tabularx} \\ \\ \relax
2:14 & \begin{tabularx}{0.7\textwidth}{X} 彼得和十一個使徒站起來，他就高聲向眾人說：「猶太人和所有住在耶路撒冷的人哪，這件事你們要知道，要側耳聽我的話。 \end{tabularx} \\ \\ \relax
2:15 & \begin{tabularx}{0.7\textwidth}{X} 這些人並不像你們所想的喝醉了，因為現在才早晨九點鐘。 \end{tabularx} \\ \\ \relax
2:16 & \begin{tabularx}{0.7\textwidth}{X} 這正是藉著先知約珥所說的： \end{tabularx} \\ \\ \relax
2:17 & \begin{tabularx}{0.7\textwidth}{X} 『神說：在末後的日子，我要將我的靈澆灌凡血肉之軀的。你們的兒女要說預言；你們的少年要見異象；你們的老人要做異夢。 \end{tabularx} \\ \\ \relax
2:18 & \begin{tabularx}{0.7\textwidth}{X} 在那些日子，我要把我的靈澆灌，甚至給我的僕人和婢女，他們要說預言。 \end{tabularx} \\ \\ \relax
2:19 & \begin{tabularx}{0.7\textwidth}{X} 在天上，我要顯出奇事，在地下，我要顯出神蹟，有血，有火，有煙霧。 \end{tabularx} \\ \\ \relax
2:20 & \begin{tabularx}{0.7\textwidth}{X} 太陽要變為黑暗，月亮要變為血，這都在主大而光榮的日子未到以前。 \end{tabularx} \\ \\ \relax
2:21 & \begin{tabularx}{0.7\textwidth}{X} 那時，凡求告主名的都必得救。』 \end{tabularx} \\ \\
[1ex]
\hline
\hline
\end{longtable}
講題:靈動我心力行奇事,言說基督見證上帝 Acts of Spirited hearts speaking for Christ as witness of God

經課:使徒行傳2:1-21

講員:鄭漢文博士

場合:香港中文大學崇基學院禮拜堂主日崇拜

日期:2020年5月31日

1. 「他們」第一代信徒如何經歷聖靈,因而動心得力行奇事?

弟兄姊妹們,你們好,大家平安.

今天我能夠面對面講道,真好.這樣的相見,對我來說,日子也不多了.世事多變,世事難料.

今天好大可能會是我在崇基禮拜堂講道的倒數第二次……下一次講道安排了在6月28日,好大可能是我在這裡最後一次講道.

因為我在香港中文大學任教了四份一世紀,七月底正式「被離開」了.

我再說一遍,我並不是退休,更不是榮休,我是工作了十三個合約,不被再續約而要離開.我還在找工作中……

為何最後?因為我出任榮譽助理校牧的委任,是基於我在中大全職工作的身分.所以,七月31日我中大教員身分便完結了.

現在請你們不要問我將來會如何,因為我也在問上帝,我將來會如何?等上帝回答了我,我再告訴你們.

上帝好似對我說,再等一等,六月下旬就會知.所以,我會在6月28日講道時與大家分享我的動向.

今天是聖靈降臨主日,我選講使徒行傳2章1-21節,經課主線是:「聖靈降臨」.我費剎思量想到今天講道的題目:「靈動我心力行奇事,言說基督見證上帝」 ,將整卷《使徒行傳》共28章要說的用這十六個字概括了.

讓我們一同誦讀今日的經課《使徒行傳》2章1-21節這段聖靈降臨的記載,請打開聖經到《使徒行傳》第二章第1節讀到第21節(以下講章會用和修本),一同思考今天的經課,思考「他們」第一代信徒如何經歷聖靈,因而動心得力行奇事?

我們現在開始把這21節經文分開三段來思考,用三個What的提問來考量經文說了甚麼?

Q1.首先的What──發生了甚麼事的問題?What has happened?

Q2.其次的What──甚麼人發生了甚麼變化的問題?What changes have been brought about?

Q3.第三個What──甚麼人因變化而有甚麼反應的問題?What responses to these changes?

Q1.首先的What──發生了甚麼事的問題?What has happened?

1五旬節那日到了,他們全都聚集在一起.

2忽然,有響聲從天上下來,好像一陣大風吹過,充滿了他們所坐的整座屋子；

3又有舌頭如火焰向他們顯現,分開落在他們每個人身上.

《使徒行傳》第二章第1至3節這段文字記錄了五旬節聖靈降臨的現象,這現象是一個事實的屬靈描述(spiritual description of a fact).

耳,聽到甚麼呢?聽到有響聲,在空間上是從天上下來.

如何描述呢?好像一陣大風吹過,是好像,是一個比喻.

眼,看見甚麼呢?又有舌頭如火焰向他們顯現,有如火焰,是另一個比喻.

有何體驗呢?分開落在他們每個人身上.就要追看每個人的生命改變.

Q2.其次的What──甚麼人發生了甚麼變化的問題?What changes have been brought about?

《使徒行傳》第二章第4節記錄了五旬節聖靈降臨的第一個出現的不一樣:

4他們都被聖靈充滿,就按著聖靈所賜的口才說起別國的話來.

NIV英文譯本是這樣記載:

4 All of them were filled with the Holy Spirit and began to speak in other tongues as the Spirit enabled them.

我們值得留意,其實對聖靈降臨的比喻為「舌頭如火焰」(tongues of fire)分開落在他們每個人身上所產生的變化,就是使徒能力上改變,能有「如火焰的舌頭」(tongues to fire),由聖靈恩賜口才,說起別國的話來.

我1995年開始在中大任教,若把二十五年來我所研究的學問,濃縮為一句話,就是:「心動力變,力用事成.」

我是用這個思路,來組構今天的講道主題的上半句,就是「靈動我心,得力行奇事.」

聖靈降臨對人類歷史產生了最大的分別,就是聖靈感動我們的心,因而得著能力,去行出奇事神蹟來.(inspired by the Holy Spirit to be empowered to bring about works of wonder)

《使徒行傳》第二章第4節記錄了五旬節聖靈降臨的第一個出現的不一樣,在場出現的人見證了這樣的不一樣.

Q3.第三個What──甚麼人因變化而有甚麼反應的問題?What responses to these changes?

讓我再看看第5至12節:

5那時,有從天下各國來的虔誠的猶太人,住在耶路撒冷.

6這聲音一響,許多人都來聚集,各人因為聽見門徒用他們各自的鄉談說話,就甚納悶,

7都詫異驚奇說:「看哪,這些說話的不都是加利利人嗎?

8我們每個人怎麼聽見他們說我們生來所用的鄉談呢?

9我們帕提亞人,瑪代人,以攔人,和住在美索不達米亞,猶太,加帕多家,本都,亞細亞,

10弗呂家,旁非利亞,埃及的人,並靠近古利奈的利比亞一帶地方的人,僑居的羅馬人,

11包括猶太人和皈依猶太教的人, 克里特人和阿拉伯人,都聽見他們用我們的鄉談講論上帝的大作為.」

12眾人就都驚奇困惑,彼此說:「這是甚麼意思呢?」

「靈動我心,得力行奇事」的結果是人的心被打開,追問Why?「這是甚麼意思呢?」

講道其實都是講經常的道理,但最難是聽道的人把心打開,所以其實講道最重要是直指人心.【華語】如果我現在用華語把今天的講道講下去,你們當中誰會覺得很開心?請舉手!

【English】What if I use English to continue my preaching, who else would feel delighted? Please raise up your hand!

【French】Voulez-vous parlez le Francaise? Comprenez-vous l’HongKongese?

其實,我是「外語有限公司」.我係沒有這些舌頭,說方言的能力.

但當時的使徒,就因為聖靈如「舌頭如火焰」(tongues of fire)恩賜得能力,說起各種的話語來(tongues to fire),讓聖道直入人心,燒起各人的心火.

但是,這真是一種不可思議的奇事,人未必接受得到,就會用各種解釋方式去「去奇化」.

那時,有虔誠的猶太人從天下各國來,住在耶路撒冷.當時從羅馬來耶路撒冷的客旅中,有些猶太人,有些進猶太教的人,成長於不同的地方,也就講不同的方言,即是鄉下話.

各人聽見這些門徒用眾人的鄉下話說話,就十分不明不白.反問這些人不都是生來都在加利利,未曾到過我們不同的出生地,未曾有機會學習我們的鄉下話,為甚麼會用我們生來就懂的方言說話?

---------------------------------------------------

|      |                                          |

| 2:13 |  還有人譏誚,說:「他們是灌滿了新酒吧!」|

|      |                                          |

---------------------------------------------------

五旬節聖靈降臨後第一個講章彼得十二位使徒一起為聖道作見證.

請續看第14至21節,彼得在五旬節的講道,開首是確認聖靈降臨的預言應驗了(我把彼得講的話在此講一次):

14彼得和十一個使徒站起來,他就高聲向眾人說:「猶太人和所有住在耶路撒冷的人哪,這件事你們要知道,要側耳聽我的話.

15這些人並不像你們所想的喝醉了,因為現在才早晨九點鐘.

16這正是藉著先知約珥所說的:

17『上帝說:

在末後的日子,

我要將我的靈澆灌凡血肉之軀的.

你們的兒女要說預言；

你們的少年要見異象；

你們的老人要做異夢.

18在那些日子,我要把我的靈澆灌,

甚至給我的僕人和婢女,

他們要說預言.

19在天上,我要顯出奇事,

在地下,我要顯出神蹟,

有血,有火,有煙霧.

20太陽要變為黑暗,

月亮要變為血,

這都在主大而光榮的日子未到以前.

21那時,凡求告主名的都必得救.』

2. 四個「三一」解讀《使徒行傳》的時空人事

其實今日講道的經課已經講完.

聽道的你們,滿意嗎?(收貨嗎?)

今日的經課已經講完,但今天要講的道,正式開始.

尋道的入手位,可以問:這段經課有甚麼前因,後果?

先說後果.

五旬節聖靈降臨後,第一個講章是由彼得為聖道作見證,言說基督,見證上帝.

見證之後,群眾如何反應?

按《使徒行傳》第2章37-41節的記載:一次講道,改變了三千人,受洗之後,形成教會.

這是教會的出現,第一個教會的建立,初期教會的故事,教會史的開始.

使徒行傳在整本聖經的位置是獨特的,因為在裡面我們讀到「教會的創世記」(genesis of the Church).

《使徒行傳》第2章37-41節的記載,當時他們得聽聖道而生命有所改變,當時五旬節來自不同地方的人,信而受洗,建立了教會.

五旬節聖靈降臨的後果是教會形成.

那麼,今日2章1-21節這段經課有甚麼前因?

讓我們再打開聖經,先看一看《使徒行傳》全書的關鍵,就在第一章第8節,耶穌基督升天前對門徒說:

-------------------------------------------------------------------------

|      |                                                                |

| 1:8  |  但聖靈降臨在你們身上,你們就必得著能力,並要在耶路撒冷,      |

|      |                                                                |

|      |  猶太全地和撒瑪利亞,直到地極,作我的見證.                    |

|      |                                                                |

-------------------------------------------------------------------------

耶穌升天跟聖靈降臨有甚麼關係?聖靈降臨為要成就甚麼?

耶穌升天,是回到天父那裡,就是從回到原來處.

但他不是就此不理我們,反而他要賜我們一份最大的禮物,就是聖靈.

耶穌升天,是好讓天父的應許得以實現:聖靈降臨.

聖靈降臨為要成就甚麼?就是要讓教會形成.

今日經課是在《使徒行傳》中的第一個「三一」(Trinity)的脈絡之中,就是三件重要事情構成一體,對使徒行傳的一個簡單見解:「耶穌升天,聖靈降臨,教會形成」.

《使徒行傳》可分為三個主要部分,記述耶穌早期的門徒怎樣在聖靈引導下,「在耶路撒冷,猶太全地和撒瑪利亞,直到地極」傳福音,敘述在日益擴大的地區中有關耶穌的福音正在迅速傳開,而且更多的教會得以建立:

一．耶穌升天後在耶路撒冷的活動(1:1－8:3)；

二．向猶太及撒瑪利亞全境發展(8:4－12:25)；

三．更進一步向地中海周圍發展,直到羅馬(13:1－28:31).

我從2011年在崇基學院神學院兼讀神道學學士,2018年完成.七年間需要兩次實習.

2012年9月至2013年4月曾於香港中文大學的崇基禮拜堂作第一次神學生實習,也有服侍內地來的留學生和留在香港生活的內地生.

第二次是從去年2013年9月至今年2014年4月在浸會大學校牧室作神學生實習,專門服侍內地生.這八個月,我被安排在普通話主日崇拜用普通話講道,我用《使徒行傳》二十八章講了九次道.

2014年我自己循道衛理聯合教會當時主恩堂堂主任劉麗詩牧師知道我在專講《使徒行傳》,便邀請我到主恩堂講道,2014年作了行傳四講.2014年六月在北角堂神學主日講了一次《使徒行傳》,2019年六月藍田堂也講了一次.

2019年1月,我被委任為崇基禮拜堂榮譽助理校牧,二,三月也開講了六課主日學,研習《使徒行傳》.

我把去年主日學的重點放入講章中,讓錯過了主日學的弟兄姊妹可以下載一些表列(tables)來自修.(全篇講章共一萬4千字,我講時只會選段來講,你們自己下載.)

從2008年到2018年攻讀神學這十年專心研究《使徒行傳》,我認為可以用四個「三一」(Trinity)去分解《使徒行傳》二十八章,形成「行傳結構」如下:

----------------------------------------------------------------------------

|               |              |              |                            |

| 三項事件:    |   耶穌升天   |   聖靈降臨   |          教會形成          |

|               |              |              |                            |

----------------------------------------------------------------------------

|               |              |              |                            |

| 三位人物:    |     彼得     |    司提反    |            保羅            |

|               |              |              |                            |

----------------------------------------------------------------------------

|               |              |              |                            |

| 三處地方(一):|   耶路撒冷   |   猶太全地   |            地極            |

|               |              |              |                            |

----------------------------------------------------------------------------

|               |              |              |              |      |      |

| 三處地方(二):|              |              |    安提阿    | 雅典 | 羅馬 |

|               |              |              |              |      |      |

----------------------------------------------------------------------------

|               |              |              |              |             |

| 三種情態:    |心動(聖靈充滿)|力變(奇事神蹟)|言說(見證講道)|             |

|               |              |              |              |             |

----------------------------------------------------------------------------

去年2月17日,我們完成了《使徒行傳》中的第一個「三一」(Trinity),就是三件重要事情構成一體,對使徒行傳的一個簡單見解:

首先,第一個重點我們講到耶穌升天,探討「你的生命在等待甚麼?」

其次,第二個重點我們講到聖靈降臨,探討「聖靈讓你的生命不再一樣?」

最後,第三個重點我們講到教會形成,探討「你願意成為教會的一份子嗎?」

-----------------------------------------------------------------------

|    |              |           |          |          |               |

| 使 | 217三項事件:|  耶穌升天 | 聖靈降臨 | 教會形成 | (第1章第8節)|

|    |              |           |          |          |               |

-----------------------------------------------------------------------

使徒行傳是講甚麼的?是講:「耶穌升天,聖靈降臨,教會形成」.

《使徒行傳》全書的關鍵,就在第一章第8節.

之後三周日,分別進行第二講至第四講,我們會探討《使徒行傳》中的第二個「三一」(Trinity),就是三位重要人物所構成的一體,成為《使徒行傳》的一個簡單的人物志:

3月3日,我們會講第二講彼得,探討「聖道讓你的生命不一樣?」(第1至12章)

3月10日,我們會講第三講司提反,探討「你的生命為甚麼去得盡?」(第6至7章)

3月17日,我們會講第四講保羅,探討「你如何由掃羅變成保羅?」(第7至28章)

--------------------------------------------------------------------------

|    |                  |          |          |          |               |

| 使 |17/2三項事件:    | 耶穌升天 | 聖靈降臨 | 教會形成 | (第1章第8節)|

|    |                  |          |          |          |               |

|------------------------------------------------------------------------|

|    |                  |          |          |          |               |

| 徒 | 3/3三位人物之一:|   彼得   |          |          | (第1至12章) |

|    |                  |          |          |          |               |

|    |-------------------------------------------------------------------|

|    |                  |          |          |          |               |

|    |10/3三位人物之二:|          |  司提反  |          | (第6至7章)  |

|    |                  |          |          |          |               |

|    |-------------------------------------------------------------------|

|    |                  |          |          |          |               |

|    |17/3三位人物之三:|          |          |   保羅   | (第7至28章) |

|    |                  |          |          |          |               |

--------------------------------------------------------------------------

使徒行傳是講甚麼的?是講:彼得,司提反,保羅等早期教會人物志,如何因為耶穌基督而人生得以改變,成為早期教會的信徒領袖,成為我們後世信徒的榜樣.

去年3月24日第五講,我們會探討《使徒行傳》中的第三個「三一」(Trinity),就是三個重要地方所構成的一體,成為《使徒行傳》的一個簡單空間圖.

講到《使徒行傳》的第三個「三一」(Trinity):(一)耶路撒冷猶太全地地極；而地極的開展,另開另一個三一:(二)安提阿,雅典,羅馬.空間上背後有一個地理上的理解:

--------------------------------------------------------------------------

|    |                  |          |          |          |               |

| 使 |17/2三項事件:    | 耶穌升天 | 聖靈降臨 | 教會形成 | (第1章第8節)|

|    |                  |          |          |          |               |

|------------------------------------------------------------------------|

|    |                  |          |          |          |               |

| 徒 | 3/3三位人物之一:|   彼得   |          |          | (第1至12章) |

|    |                  |          |          |          |               |

|    |-------------------------------------------------------------------|

|    |                  |          |          |          |               |

|    |10/3三位人物之二:|          |  司提反  |          | (第6至7章)  |

|    |                  |          |          |          |               |

|    |-------------------------------------------------------------------|

|    |                  |          |          |          |               |

|    |17/3三位人物之三:|          |          |   保羅   | (第7至28章) |

|    |                  |          |          |          |               |

|-----------------------------------------------------------------------|

|    |                  |          |          |                          |

| 行 |  三處地方(一):  | 耶路撒冷 | 猶太全地 |          地極            |

|    |                  |          |          |                          |

|    |-------------------------------------------------------------------|

|    |                  |          |          |          |       |       |

|    |  三處地方(二):  |          |          |  安提阿  |  雅典 |  羅馬 |

|    |                  |          |          |          |       |       |

--------------------------------------------------------------------------

《使徒行傳》主要記述耶穌早期的門徒怎樣在聖靈引導下,「在耶路撒冷,猶太全地和撒瑪利亞,直到地極」傳福音(徒1:8).

我們以時空人事來探討《使徒行傳》中的四個「三一」(Trinity):第一講講三事件,第二至四講講三人物,第五講講三空間再生三空間,到第六講講「心動力變行奇事」這「心,力,事」的三時態/情態,心動力變行奇事是指「聖靈充滿,得力行奇事神蹟」,其中之一是「言說基督,見證上帝,用靈語講道」.

2019年3月31日第六講(最後一講),我們會探討《使徒行傳》中的第四個「三一」(Trinity),就是聖靈出現的三種情態所構成的一體,成為《使徒行傳》的一個簡單聖靈圖:

----------------------------------------------------------------------------

|               |              |              |              |             |

| 三種情態/情態:|心動(聖靈充滿)|力變(奇事神蹟)|言說(見證講道)|             |

|               |              |              |              |             |

----------------------------------------------------------------------------

所以有學者認為《使徒行傳》其實是《聖靈行傳》,講章較後會有機會再從聖靈在《使徒行傳》的出現再反思這個觀點.

2019年2月17日至3月31日《使徒行傳》主日學六講的時間表及學習關鍵詞如下:

--------------------------------------------------------------------------

|    |                  |          |          |          |               |

| 使 |17/2三項事件:    | 耶穌升天 | 聖靈降臨 | 教會形成 | (第1章第8節)|

|    |                  |          |          |          |               |

|------------------------------------------------------------------------|

|    |                  |          |          |          |               |

| 徒 | 3/3三位人物之一:|   彼得   |          |          | (第1至12章) |

|    |                  |          |          |          |               |

|    |-------------------------------------------------------------------|

|    |                  |          |          |          |               |

|    |10/3三位人物之二:|          |  司提反  |          | (第6至7章)  |

|    |                  |          |          |          |               |

|    |-------------------------------------------------------------------|

|    |                  |          |          |          |               |

|    |17/3三位人物之三:|          |          |   保羅   | (第7至28章) |

|    |                  |          |          |          |               |

|------------------------------------------------------------------------|

|    |                  |          |          |                          |

| 行 |24/3三處地方(一): | 耶路撒冷 | 猶太全地 |          地極            |

|    |                  |          |          |                          |

|    |-------------------------------------------------------------------|

|    |                  |          |          |          |       |       |

|    |24/3三處地方(二): |          |          |  安提阿  |  雅典 |  羅馬 |

|    |                  |          |          |          |       |       |

|----|-------------------------------------------------------------------|

|    |                  |          |          |          |       |       |

| 傳 |31/3三種時態/情態:|   心動   |   力變   |   言說   |       |       |

|    |                  |(聖靈充滿)|(奇事神蹟)|(見證講道)|       |       |

|    |                  |          |          |          |       |       |

--------------------------------------------------------------------------

總而言之,2019年我們在崇基禮拜堂以主日學方式一同研習《使徒行傳》,我把曾作九講和四講的結構,轉化為六講,取了題目為:「活出聖靈同行的生命」.回應這個主題,相應提出了十個問題:

------------------------------------------------------------------------------------

| 講次   | 日期   | 四類     | 重點     | 講題                        | 經文／鑰段 |

------------------------------------------------------------------------------------

| 第1講1 | 190217 | 第一事   | 耶穌升天 | 你的生命在等待甚麼?        | 徒1:7-9    |

|        |        |          |          |                             |            |

| 第1講2 | 190217 | 第二事   | 聖靈降臨 | 聖靈讓你的生命不再一樣?    | 徒2:1-4    |

|        |        |          |          |                             |            |

| 第1講3 | 190217 | 第三事   | 教會形成 | 你願意成為教會的一分子嗎?  | 徒2:46-47  |

|        |        |          |          |                             |            |

| 第2講  | 190303 | 第一人   | 彼得     | 聖道讓你的生命不再一樣?    | 徒2:38-39  |

|        |        |          |          |                             |            |

| 第3講  | 190310 | 第二人   | 司提反   | 你的生命會為甚麼「去到盡」?| 徒7:55-60  |

|        |        |          |          |                             |            |

| 第4講  | 190317 | 第三人   | 保羅     | 你如何由掃羅變成保羅?      | 徒9:17-20  |

|        |        |          |          |                             |            |

| 第5講1 | 190324 | 第一地   | 安提阿   | 被稱為基督徒有何意義?      | 徒11:22-26 |

|        |        |          |          |                             |            |

| 第5講2 | 190324 | 第二地   | 雅典     | 基督的愛讓群體不再一樣?    | 徒17:16-23 |

|        |        |          |          |                             |            |

| 第5講3 | 190324 | 第三地   | 羅馬     | 你的羅馬在哪里?            | 徒28:23-31 |

|        |        |          |          |                             |            |

| 第6講  | 190331 | 心力事三 | 聖靈行傳 | 如何得見聖靈同行的生命?    | 全書28章   |

|        |        | 時態     |          |                             |            |

------------------------------------------------------------------------------------

希望一同完成從上而來的使命:完成行傳六講,研讀共28章的《使徒行傳》,讓自己的生命回應這十問下,所盼望的是可以一起學習成為:「活出聖靈同行的生命」 .

3. 「我」如何經歷聖靈,因而動心得力行奇事?

至於我自己,是如何經歷聖靈,因而動心得力行奇事?大概可以分為四大段落.

我快離開了,今天我想留下我的見證在這裡.

第一個段落開始自1976年信主.我在中四升中五時參加中學在香港中文大學的暑期福音營,在營裡被神的靈感動而信主的(我的信主過程好奇妙,有機會再分享).之後,我便跟中學同學返循道衛理聯合教會香港港.到了十八歲後的聖誕節我就主動請求受洗了,也就加入了這個宗派的教會.至今四十四年了.

這個循道衛理聯合教會的宗派,成為我教會生活的根.這個根的英文名字有兩重:

(1)叫循道Methodist: 講究「信仰操練和生活法度」；

(2)叫衛理Wesleyan: 追隨John Wesley(會祖約翰衛斯理)的靈性精神和成聖渴慕,對我來說,就是「學做聖靈同行的學神」.

因為香港堂少年時的四年培育,我1980年考入香港中文大學,完成主修哲學的學士和碩士後,留學英國牛津大學作研究生.1993年復活節離開牛津,轉往倫敦大學教育研究所,完成教育哲學博士於1994年底.1995年1月,我返回母校香港中文大學,任教於教育行政與政策學系,至今二十五年.我是工作了十三個合約,不被再續約而要離開.我在香港中文大學任教了四份一世紀,七月正式被離開了.我並不是退休,更不是榮休,因為我在找工作,我相信上主在安排一個地方,讓我未來可有三五年,整合我四十多年開發的學問,回應這個時代.

當我從英國學成歸港後因自1995年1月開始在香港中文大學教書,所以住入了沙田,曾經返過沙田堂,之後在香港堂重建前後返過香港堂.也因生了兩個男孩子,我們從灣仔的香港堂轉會到沙田的馬鞍山堂.一去十多年.馬鞍山堂這個教會成為了我們一家四口的家:兩個孩子在這裡成長,我和內子多年都在此服侍.2008年因為出任青少年部部長,為更好的服侍,啟動了內子和我攻讀神學的回應.

所以,第二個段落是從1995年至2008年.在這第二個階段的十四年中,我研發了一個叫「生命三維敘事法」(SQUID Narrative Method)的工具,並合著出版了一書:《尋求美好人生的自我理解:生命三維敘事法》.我創建這個方法時是解決教育哲學中的一個問題,簡單說是六育教育觀之德智體群美靈的靈育.我發現了一體兩面的品質,身分與願望三維合一的敘事方法,一面應用是整理生命的價值教育,另一應用是可以讓信徒經歷聖靈,以致有靈命的成長.原來四份一世紀以來,上帝給了我一個恩賜,是發明種種助人成長的工具,當中有些是可以讓聖靈動工的工具.

第三個段落的啟動是在2007年秋天,我參加了我自己教會總議會主辦的第48旅以馬忤斯之旅,經歷了一次靈命的更新,立心要讀神學,同時學做學神,並希望以講道為我到不同的教會作服侍.

2008-2011年,我在香港中文大學的崇基學院神學院完成神學研究文學碩士,隨即在2011年兼讀神道學學士,在前年2018年畢業.這份論文研究一位叫Dr.Shen的崇基精神的奠定者,詳情我會在6 月 14 日下午 1 時半「讀•書•人•會」(禮拜堂活動中心舉行的)聚會分享沈宣仁著《三十年來情與理》這本書時再跟大家分享我這位老師的基督教教育神學.

回頭看,我的生命分為四個十四年.

從小學一年級到完成中學,我進修基礎教育,歷經第一個十四年.

從大學一年級到完成博士,我專修哲學的學業,歷經第二個十四年.

從1995年工作至2008年,我專作教育的事業,歷經第三個十四年.

2008年開始了我第四個十四年,就是讀神學的志業,開發轉化生命的學問,就是展開與信仰相關建立生命的企業.

並期望在這十四年內(還有兩年),我能結合「哲學,教學與神學這三學」為「轉化生命的學問」.

希望未來能結合上主給我的三種恩賜為一種服侍:就是「講道,講學和講笑」的三合一事工.

以上是我成為聖靈同行的學神的故事.

以上是我經歷聖靈,因而動心得力行奇事的故事.

我相信到6月28日最後一次講道時,好快進入第四個階段的啟動.

為何最後,因為我出任榮譽助理校牧的委任,是基於我在中大全職工作的身分.所以,七月31日我中大教員身分便完結了.

但我相信上主在安排一個地方,讓我未來可有三,五年,整合我四十多年來開發轉化生命的的學問,以回應這個時代.

4. 學做聖靈同行的「學神」:靈動我心力行奇事,言說基督見證上帝

在我兼讀神學十年內除了講道講《使徒行傳》外(單單今天經課徒2:1-21也講了四次),也寫了多篇「行傳研究」(Acts Studies),其中一篇文章試圖探討這個問題:「對《使徒行傳》作出敘事分析可以讀出怎樣的聖靈觀?」.文章的文題命名為:「《使徒行傳》中聖靈的敘事功能」,是作者試圖以把「聖靈」這個字辭(term)在《使徒行傳》如何出現,對於《使徒行傳》這個敘事體(作為一種文體)所起的作用作出研究,從而呈現讀者可以得出的聖靈觀可以是怎樣的.

我做了這個研究,產生了兩個表,第一個是《使徒行傳──漢文樣色52次聖靈出現心力事版》,第二個是《使徒行傳──九分28章為三段各三節版》(若想索取這兩個附錄,請電郵給我:cu2020hk.rogercheng@gmail.com).

研究過程會保留在上載了的講章,在這裡我只分享結論.

首先,按聖經和修本版,「聖靈」這個雙字詞在《使徒行傳》出現了52次.

其次,究竟有沒有一個歸類法,可以把52項歸為甚麼類?多少類?

1994完成教育哲學博士後,我1995年開始任教於香港中文大學教育學院,漸漸把教育哲學的原理歸結為一句話:「心動力變,力用事成」.自2000年開始把此句轉化為生命教育和敘事教學法的運用,得出「心動力變成故事」的敘事解讀.依此思路,一個作者(者)因為心動而力變,力變變作言與行,以致有所成事.回頭敘述,就是把故事中的作者的心動力變加以說明出來,從而呈現相關的意義.

我的研究是從敘事功能(narrative function)去作歸類,筆者提出一個「言行心者」四分法:

演說(言)inspired speeches,

敘事(行)spirited acts,

充滿(心)filled with the Holy Spirit　及

作者(者)Holy Spirit as an agent,

借之以達致敘事理解(narrative understanding).

如果耶穌是「道成肉身」(incarnation),聖靈對人之影響可以稱為「人成道身」(embodiment),而「人成道身」在敘事上可以四類狀態.

一是言類:聖靈影響人的言說,人的言說影響故事發展.換言之,因為聖靈所動,說出特別的講話來,就是言辭的聖道臨在功能(inspiring),影響事情的發展.共有8次「聖靈」出現在力言spirited speeches,言辭行為見證講道.有學者研究《使徒行傳》內中包含36個演說,其中30個是信徒講的.

二是行類:聖靈影響人的行動,人的行動影響故事發展.換言之,因為聖靈所動,作出特別的行為來,就是「人成道身」的體現功能(embodying),影響事情的發展.共有15次「聖靈」出現在力行spirited acts,行動行出奇事神蹟.

三是心類:聖靈充滿影響人物,人物影響到故事發展.換言之,因為聖靈充滿了當事人,活出特別的狀態來,就是充滿的恩賜臨在功能(empowering),影響事情的發展.共有22次「聖靈」出現在心動filled with the Holy Spirit,聖靈充滿.

但52次聖靈用字的出現,有些無法合乎以上三類的,可以歸為第四類:

四是者類:聖靈自身構成為人物,行者影響到故事發展.換言之,聖靈自己是一個作者,是一個行者,是聖靈的獨自行動功能(initiating),不依人類而影響事情的發展.共有7次「聖靈」出現在行者Holy Spirit as an agent,聖靈有自性,可自主,可行動.

總結而言,《使徒行傳》其實是《聖靈行傳》,從中看到使徒「靈動我心力行奇事,言說基督見證上帝」的傳奇.

單單從《使徒行傳》28章全卷經文本身來閱讀(back to the text itself),究竟《使徒行傳》此文本會呈現聖靈在故事發展的怎樣的面貌(樣式或格局,英文指pattern)?

按大部份的相關參考,《使徒行傳》都被視為是早期教會發展的歷史,而歷史是敘事體的一種(敘事體還可以有個人傳記,小說等).所以,本文重點分析「聖靈」這個字辭在《使徒行傳》敘事過程中的敘事功能,從而單一突顯敘事法的觀點角度(isolating the narrative perspective by suspending the other perspectives).

然而,本文所採取的關於「基督徒讀者」(Christian readership)的一種策略,就是設想:若然讀者是只具備普通教育水平,沒有希臘原文的語言能力,沒有成書歷史脈絡編修方法(或其他方法)的掌握,沒有既定的聖靈論(不論是來自神學的,聖經的或教會的來源).設定讀者是沒有以上角度的能力,甚至設定讀者是沒有聖靈的信念或特定的理解.

筆者以及讀者第一個要面對的問題,就是任何「基督徒讀者」(Christian reader)要面對的文本問題.

若然讀者是只具備普通教育水平,例如活在香港的一個30多歲的職工,在慕道中或剛信耶穌不久,會讀哪個文本?

筆者採納了香港聖經公會的和合修訂本(上帝版).所作的程序是如下的:

首先,在香港聖經公會的網址(\url{http://rcuv.hkbs.org.hk})網上下載選用了的是「和合修訂本上帝版」了全卷28章的漢語經文成為一個WORD檔案,連分題及分段也一併保留.Default為新細明體.

運用WORD的尋找功能,鍵入「聖靈」雙字詞.得出52次的出現.所有52次的出現都標以紅色.

運用WORD的尋找功能,鍵入「靈」單字.得出14多次的出現.所有14次的出現都標以粉紅色(或深紅色).

由第一章第1節開始,續句看下去,把直述句內(及引述句)的文字改為標楷體,而在直述句中的引述句則加以斜體化.

以上程序,加上香港聖經公會對《使徒行傳》簡介作為首頁,就得出〈附錄(一)香港聖經公會《使徒行傳》漢語和合修訂本28章全卷,多顏色經節彙編,包括「聖靈」以紅色標示52次出現〉.

於是,就產生了《使徒行傳──漢文樣色52次聖靈出現心力事版》這個表格的基礎.

以下進一步的論述以此文本為基礎來開發出另外兩個表列.

自1994完成教育哲學博士後,我1995年開始任教於香港中文大學教育學院,漸漸把教育哲學的原理歸結為一句話:「心動力變,力用事成」.自2000年開始把此句轉化為生命教育和敘事教學法的運用,得出「心動力變成故事」的敘事解讀.依此思路,一個作者(者)因為心動而力變,力變變作言與行,以致有所成事.回頭敘述,就是把故事中的作者的心動力變加以說明出來,從而呈現相關的意義.

意義可以細分為四分節,續一討論52次用語如何通過演說(言),敘事(行),充滿(心)及作者(者)這四類來達致敘事理解(narrative understanding).

至於《使徒行傳》如何分段?不同的《使徒行傳》研究專論,不同的學者,會有不同的分段.

以保羅信主前後的生平作分段,只是其中一種分法.例如,以下把《使徒行傳》28章按保羅出現分為前中後三段,因為後段有三次宣教,所以再細為後段為後前,後中,後後段(共五段).然後把52次聖靈出現所涉及的四類再分配.

----------------------------------------------------------------------------------------------

|    28章     | 前中後三段        | 後段再細分為三               |  言  |  行  |  心  |  者  |

|             |                   |                              |      |      |      |      |

| ----------- - ----------------- - ---------------------------- - ---- - ---- - ---- - ---- |

|             |                   |                              |      |      |      |      |

|   1:1-8:36  | 前段:保羅信主前   |                              | 演說 | 行事 | 心動 | 作者 |

|             |                   |                              |      |      |      |      |

| ----------- - ----------------- - ---------------------------- - ---- - ---- - ---- - ---- |

|             |                   |                              |      |      |      |      |

|  9:1-11:18  | 中段:保羅初信成長 |                              | 演說 | 行事 | 心動 | 作者 |

|             |                   |                              |      |      |      |      |

| ----------- - ----------------- - ---------------------------- - ---- - ---- - ---- - ---- |

|             |                   |                              |      |      |      |      |

| 11:19:17:15 | 後段:保羅宣教     | 後段:保羅宣教(一)e.g.Antioch | 演說 | 行事 | 心動 | 作者 |

|             |                   |                              |      |      |      |      |

| ----------- - ----------------- - ---------------------------- - ---- - ---- - ---- - ---- |

|             |                   |                              |      |      |      |      |

| 17:16-18:23 |                   | 後段:保羅宣教(二)e.g.Athens  | 演說 | 行事 | 心動 | 作者 |

|             |                   |                              |      |      |      |      |

| ----------- - ----------------- - ---------------------------- - ---- - ---- - ---- - ---- |

|             |                   |                              |      |      |      |      |

| 18:24-28:31 |                   | 後段:保羅宣教(三)e.g.Rome    | 演說 | 行事 | 心動 | 作者 |

|             |                   |                              |      |      |      |      |

----------------------------------------------------------------------------------------------

但是,這種以保羅生平為中心的分段法,是把《使徒行傳》讀成《使徒保羅傳記》.雖然保羅佔後半篇幅頗多,但在保羅佔重要性之前半已有其他使徒如彼得和司提反等,既名為《使徒行傳》,前半對其他使徒的敘事應有作用.把《使徒行傳》讀成《使徒保羅傳記》對於前段八章在耶路撒冷的重要事件有所忽略,特別是沒有理會第一章1至8節耶穌升天應許聖靈的卷首語(寫作佈局預示),也沒有對五旬節聖靈降臨對信徒的重要影響,也沒有對第1章到第28章是展開「聖靈降臨在你們身上得著能力在耶路撒冷,猶太全地和撒瑪利亞,直到地極,作見證」的空間上的敘事.

然後可以再細分為三分節,續一討論52次用語如何借用敘事段落(narrative conclusion)以及敘事結句(narrative conclusion)這兩個指示,來把《使徒行傳》全卷28章加以分為三大段,各段又可再細分為三小段,共分九個段落.在第3節會論述,這樣的分段法,可以回應第1章到第28章是展開「聖靈降臨在你們身上得著能力在耶路撒冷,猶太全地和撒瑪利亞,直到地極,作見證」的空間上的敘事.

Shepherd (1994)在 The Narrative Function of the Holy Spirit as a Character in Luke-Acts 主張把「聖靈」看成一個故事中的人物或角色(character),因此值得研究聖靈這個用字在《使徒行傳》敘事過程中的敘事功能是如何發揮的.

如果耶穌是「道成肉身」,聖靈對人之影響可以稱為「人成道身」,而「人成道身」在敘事上可以三類狀態.

一是言類:聖靈影響人的言說,人的言說影響故事發展.換言之,因為聖靈所動,說出特別的講話來,影響事情的發展.

二是行類:聖靈影響人的行動,人的行動影響故事發展.換言之,因為聖靈所動,作出特別的行為來,影響事情的發展.

三是心類:聖靈充滿影響人物,人物影響到故事發展.換言之,因為聖靈充滿了當事人,活出特別的狀態來,影響事情的發展.

但52次聖靈用字的出現,有些無法合乎以上三類的,可以歸為第四類:

四是者類:聖靈自身構成為人物,行者影響到故事發展.換言之,聖靈自己是一個作者,是一個行者,不依人類而影響事情的發展.

用human agency這個英文字再表述,就是言行心這三類人都是在co-agency with the Holy Spirit,而第四類則是Holy Spirit alone as divine agent.以下可以用一個架構陳示聖靈這四種敘事功能:

------------------------------------------------------------------------------------------------------------------

|                    |                    |                      |                      |                        |

| 聖靈自身構成為人物 | 聖靈充滿影響人物   | 聖靈影響人的行動     | 聖靈影響人的言說     | 聖靈出現的四種敘事功能 |

| 行者影響到故事發展 | 人物影響到故事發展 | 人的行動影響故事發展 | 人的言說影響故事發展 |                        |

|                    |                    |                      |                      |                        |

| ------------------ - ------------------ - -------------------- - -------------------- - ---------------------- |

|                    |                    |                      |                      |                        |

|                    |                    |                      |   (力言)言辭行為   |      一．言類          |

|                    |                    |                      |   Inspired speeches  |                        |

|                    |                    |                      |                      |                        |

| ------------------ - ------------------ - -------------------- - -------------------- - ---------------------- |

|                    |                    |                      |                      |                        |

|                    |                    |     (力行)行動     |                      |      二．行類          |

|                    |                    |    Spirited acts     |                      |                        |

|                    |                    |                      |                      |                        |

| ------------------ - ------------------ - -------------------- - -------------------- - ---------------------- |

|                    |                    |                      |                      |                        |

|                    |   (心動)充滿     |                      |                      |      三．心類          |

|                    |  Filling with the  |                      |                      |                        |

|                    |    Holy Spirit     |                      |                      |                        |

| ------------------ - ------------------ - -------------------- - -------------------- - ---------------------- |

|                    |                    |                      |                      |                        |

|    (作者)行者    |                    |                      |                      |      四．者類          |

|   Holy Spirit as   |                    |                      |                      |                        |

|      an agent      |                    |                      |                      |                        |

------------------------------------------------------------------------------------------------------------------

究竟在《使徒行傳》中「聖靈」用字52次出現,那些是聖靈對人之影響可以稱為「人成道身」的敘事功能中第一種言類:聖靈影響人的言說,人的言說影響故事發展.換言之,因為聖靈所動,說出特別的講話來,影響事情的發展.

聖靈的敘事功能之一,就是言辭的聖道臨在功能(inspiring).講不講?如何講?從36個演說看出,有聖靈臨在的講話是 inspired by the Spirit,而且講話會inspiring the audience,這樣的speech act有making the turn of history的作用.我把8次出現歸為此類.

究竟在《使徒行傳》中「聖靈」用字52次出現,那些是聖靈對人之影響可以稱為「人成道身」的敘事功能中第二種行類:聖靈影響人的行動,就是「人成道身」的體現功能(embodying),人的行動影響故事發展.換言之,因為聖靈所動,作出特別的行為來(但並非講話),影響事情的發展.我把15次出現歸為此類.

聖靈的敘事功能之三,是充滿的恩賜臨在功能(empowering).人的心被聖靈所充滿,除了說方言,教導及傳道的靈動之言外,還可以趕鬼,治病和預知等行類.充滿的恩臨,是一個活不活的存有分別.活著,就是歷史的一個見證.我把22次出現歸為此類.

聖靈的敘事功能之四,是聖靈的獨自行動功能(initiating).前面三類都涉及human agency as sole agent or co-agent,是人令歷史進展,推進故事.但凡有聖靈字辭出現的情節,就出現co-agency.而還有的是第四類sole agency of the Holy Spirit.聖靈有自性,可自主,可行動.我把7次出現歸為此類.

香港聖經公會在網上(\url{http://rcuv.hkbs.org.hk/RCUV_2/ACT/}Ｓ/)是這樣說明:

「《使徒行傳》的一個重要特徵就是強調聖靈的工作.聖靈的能力在五旬節那天降臨在耶路撒冷信徒身上,並且在本書所敘述的各件事蹟上繼續帶領及堅立教會和教會的領袖們.初期教會的信息可以從本書幾篇重要的講章看得到.書中所記述的許多重大事件顯示這信息怎樣影響了初期的信徒和教會的團契生活.」

第一章1至8節耶穌升天應許聖靈的卷首語,其實是《使徒行傳》的寫作佈局的預示,預示空間上的敘事,就是對第1章到第28章是展開「聖靈降臨在你們身上得著能力在耶路撒冷,猶太全地和撒瑪利亞,直到地極,作見證」的三環擴展.

李斐德(Walter L. Liefeld)在《詮釋使徒行傳》一書中提及按敘事段落的結句作為指示.他一共列出12個例子來.

正如上文提及,按學者如Soards (1994)在The Speeches in Acts中,就列出當中的演說(speeches in direct speech)共有36段.單閱讀本文附錄(一)之《使徒行傳──漢文樣色52次聖靈出現心力事版》,便發現加了開關引號的言說(直述句direct speech及引述句quotation)(以藍色標示)在整個28章的敘事體中構成顯要部份.

而另一位學者Padilla, Osvaldo (2008) 在The Speeches of Outsiders in Acts一書中,也發現《使徒行傳》內中包含許多演說,他更注意到有些演說是信徒講的,另有一些是非信徒所講的,他稱後者為the speeches of outsiders(外人),36段中佔6段.

總結以上論證,再看附錄(一)之《使徒行傳──漢文樣色52次聖靈出現心力事版》,信徒講的演說講話,或多或少是因為在前受聖靈所影響的,可以稱為inspired speeches,在後令得傳道的故事有變化,例如《使徒行傳》作者會敘述:信主的人增多了.所以這些不單是speeches,還是speech acts──言辭行為,就是在敘事體中這些言說其實同時做了一些有後果的行為或行動(performance or action),影響著故事的發展.由此而產另一個附錄(二)《使徒行傳──九分28章為三段各三節版》.

(若想索取以上兩個附錄,請電郵給鄭漢文:cu2020hk.rogercheng@gmail.com)

5. 「你們」這一代信徒又如何經歷聖靈,因而動心得力行奇事?

今日講道是從第二章見到「他們」,反思到「我自己」,現在到「你們」.

「你們」這一代信徒,又如何經歷聖靈,因而動心得力行奇事?

《使徒行傳》的一個重要特徵就是強調聖靈的工作,所以有人稱之為《聖靈行傳》.

《使徒行傳》第二章,就是聖靈的能力在五旬節那天降臨在耶路撒冷信徒身上,以及在本書所敘述的各件事蹟上繼續帶領及堅立教會和教會的領袖們.

五旬節聖靈降臨後,第一個講章是由彼得為聖道作見證,見證之後,群眾如何反應?

使徒行傳第2章37-41節這樣記載:

37眾人聽見這話,覺得扎心,就對彼得和其餘的使徒說:「諸位弟兄,我們該怎樣做呢?」

38彼得對他們說:「你們各人要悔改,奉耶穌基督的名受洗,使你們的罪得赦免,就會領受所賜的聖靈.

39因為這應許是給你們和你們的兒女,並一切在遠方的人,就是給所有主-我們的上帝所召來的人.」

40彼得還用更多別的話作見證,勸勉他們說:「你們當救自己脫離這彎曲的世代.」

41於是領受他話的人,都受了洗；那一天,門徒約添了三千人.

一次講道,改變了三千人,受洗之後,形成教會.

這是教會的出現,第一個教會的建立,初期教會的故事,教會史的開始.

當時他們得聽聖道而生命有所改變,當時五旬節來自不同地方的人,信而受洗,建立了教會.最早期的教會生活是怎樣的?

《使徒行傳》第2章37-41節這樣記載:

42他們都專注於使徒的教導和彼此的團契,擘餅和祈禱.

43眾人都心存敬畏；使徒們又行了許多奇事神蹟.

44信的人都聚在一處,凡物公用,

45又賣了田產和家業,照每一個人所需要的分給他們.

46他們天天同心合意恆切地在聖殿裡敬拜,且在家中擘餅,存著歡喜坦誠的心用飯,

47讚美上帝,得全體百姓的喜愛.主將得救的人天天加給他們.

初期教會的資訊可以從本書幾篇重要的講章看得到.書中所記述的許多重大事件顯示這資訊怎樣影響了初期的信徒和教會的團契生活.

回想今天的我們這一代信徒,如何在教會成長?

一個人加入教會前要:

聽道.

悔改.

認罪.

受洗.

一個人加人教會後的生活:

恆心遵守使徒的教訓,彼此交接,擘餅,祈禱.

信的人都在一處,凡物公用．並且賣了田產家業,照各人所需用的分給各人.

他們天天同心合意恆切的在殿裡,讚美　神.

且在家中擘餅,存著歡喜誠實的心用飯.

一個人加入教會後的影響:

眾人都懼怕.

使徒又行了許多奇事神跡.

得眾民的喜愛.

主將得救的人,天天加給他們.

究竟從五旬節經驗,對「如何經歷聖靈,因而動心得力行奇事」這問題有何啟示?

簡單說,聖靈可以讓我們的生命不再一樣的不一樣是有不同的層次的,研究聖靈論的學者會由不同的分法.我們在這裡,可以從最厲害,最高峰最特殊的改變(五旬節經驗),到最基本最持久,最通常的改變,列為五層如下:

(一)聖靈充滿的最高表現,就是得著不學習而能運用的能力,可以行奇事神蹟(wonders and signs)──例如,五旬節聖靈降臨後十二使徒所作所為(說方言,趕鬼,治病等)

(二)長出屬靈恩賜(spiritual gifts)來,結出聖靈的果子──而最大的是仁愛

(三)常被聖靈感動,我們可以效法基督,活出可以感動他人的生命來──與受苦的人同受苦

(四)聖靈是保惠師,保守我們與聖父和聖子相連,且成為我們禱告的仲介者──靈裡代求

(五)聖靈讓我們活在平安,安息與安靜之中(peace, rest and still)──心靈的安寧與更新

今天,在香港,你心靈安寧嗎?如果,你心有不安,學習回應聖靈吧.

各位,今天是2020年5月31日.

今天發生了甚麼事?今天想起,九個月前發生了甚麼事?

今天中大發生了甚麼事?(下午我會為崇基學院神學院步行籌款而導賞崇基.)

今天的香港會否好像耶路撒冷聚集各種言語的人在此城市?今天的香港發生了甚麼變化?

今天的香港會否像五旬節一樣會發生大變化?就看我們因變化而有甚麼反應?

今天讓我們迎接聖靈,好讓靈動我心,得力行奇事,言說基督,見證上帝.

讓我們一起祈禱:

「創造生命的上主啊,

讓我們活出你創造的奇工,讓我們充滿你的智慧.

為我們釘十字架的耶穌基督啊,讓我們的生命活出超越取死的身體.

求你憐憫我們,讓我們活出不一樣的生命.

保惠我們的聖靈啊,

求你充滿我們,更新我們的心靈,讓我們活出不一樣的生命.

求你與我們同行,好讓我們的生命能夠去到盡,那怕是生生死死的盡頭.

神啊神,今天讓我們迎接聖靈你,好讓我們:靈動我心,力行奇事；言說基督,見證上帝.

阿們.」

願恩惠常與你們同在!

但願榮耀都歸給聖父,聖子,聖靈,永在的三一神.阿們.



\newpage

\section{奉命行事}
\label{sec:504}
\textbf{王家輝牧師 - 奉命行事}
\newline
\newline
\hyperref[sec:503]{< < < PREV SERMON < < <}
~
\hyperref[sec:toc]{[返目錄]}
~
\hyperref[sec:505]{> > > NEXT SERMON > > >}


馬太福音 28:16-20
\newline
\begin{longtable}{cl}
\hline
\hline
章節 & 經文 (和合本修訂版)\\
\hline
28:16 & \begin{tabularx}{0.7\textwidth}{X} 十一個門徒往加利利去，到了耶穌指定他們去的山上。 \end{tabularx} \\ \\ \relax
28:17 & \begin{tabularx}{0.7\textwidth}{X} 他們見了耶穌就拜他，然而還有人疑惑。 \end{tabularx} \\ \\ \relax
28:18 & \begin{tabularx}{0.7\textwidth}{X} 耶穌進前來，對他們說：「天上地下所有的權柄都賜給我了。 \end{tabularx} \\ \\ \relax
28:19 & \begin{tabularx}{0.7\textwidth}{X} 所以，你們要去，使萬民作我的門徒，奉父、子、聖靈的名給他們施洗， \end{tabularx} \\ \\ \relax
28:20 & \begin{tabularx}{0.7\textwidth}{X} 凡我所吩咐你們的，都教導他們遵守。看哪，我天天與你們同在，直到世代的終結。」 \end{tabularx} \\ \\
[1ex]
\hline
\hline
\end{longtable}
講題:奉命行事 To act on orders

經文:馬太福音28:16-20

講員:王家輝牧師

場合:香港中文大學崇基學院禮拜堂主日崇拜

日期:2020年6月7日

16十一個門徒往加利利去,到了耶穌約定的山上.17他們見了耶穌就拜他,然而還有人疑惑.18耶穌進前來,對他們說:「天上地下所有的權柄都賜給我了.19所以,你們要去,使萬民作我的門徒,奉父,子,聖靈的名給他們施洗 .20凡我所吩咐你們的,都教訓他們遵守,我就常與你們同在,直到世界的末了.」

「奉命行事」是忠心的表現,對於講求紀律,服從上級命令的軍隊就更加重要.因為個別士兵只看見眼前的情況,所以未必能夠掌握整個戰局與戰略.倘若只按自己的歡喜或判斷行事,最終非但送上自己性命,更可能會連累其他同袍.然而,在極端的情況下,怎樣執行上級命令卻有很大的討論空間.

在1989年2月,當時分隔德國為東德和西德的柏林圍牆依然存在.有兩位東德青年Chris Gueffroy和Christian Gaudian那時誤以為「射殺任何翻越柏林圍牆的人」的命令已經被撤銷,故此決定翻越柏林牆,希望進入西德.怎料被四位守衛圍牆的東德士兵發現,而開槍阻止.結果Chris Gueffroy被當場擊斃,成為東西德統一合併前最後一位嘗圖翻越圍牆而被槍殺的東德遇難者,而Christian Gaudian則被捕入獄.

在1991年9月(當時東西德已合併大概已有一年)柏林地區法院起訴這四名守衛圍牆的前東德士兵.士兵們的辯護律師Rolf Bossi認為根據東德法律,東德民眾沒有自由離開國家的權利,槍擊非法離境是合法的執勤.而四名士兵亦堅稱自己是執行上級的命令,沒有違法.

主審法官Theodor Seidel在裁定殺射Chris Gueffroy的士兵Ingo Heinrich為有罪時表示,雖然作為士兵,他們是執行東德法律與連串上級命令的最低層,不過「不是一切合法的事情就是對的」(not everything that is legal is right).此外,他還補充「在20世紀末,在代表權力結構殺人時,沒有人有權忽視自己的良心」(At the end of the 20th Century … , no one has the right to ignore his conscience when it comes to killing people on behalf of the power structure).

後來Ingo Heinrich和另一被定罪的同袍決定上訴,而另一法官雖然同樣定Ingo Heinrich有罪,不過卻以不同的法律觀點.法官認為Ingo Heinrich按法律執勤是正確的,但卻表示槍殺手無寸鐵的逃離國家者是不成比例的武力,故Ingo Heinrich仍然被定罪,但後來監禁的判刑則改以緩刑執行.

由此可見「奉命行事」看來直接,卻原來並不簡單!

基督徒顯然需要聽命於基督,也要竭力遵行基督所吩咐,然而在不同處境之中怎樣執行卻又有許多值得深思的地方.今天我們就以幾乎是所有基督徒也會熟悉的馬太福音28章16至20節來再思想基督徒該奉基督之命行什麼事?

山上的約定

在耶穌復活後顯現的尾聲,耶穌要頒布被稱為「大使命」的命令給11位使徒.只剩下11位,當然就是沒有出賣主的猶大在其中了.

第16節記載,門徒到加利利,上了耶穌跟他們所約定的山上.按照馬太福音4:12所載,耶穌是在加利利開始他的傳道工作.現在已是他留在世上最後的日子,而門徒將要延續耶穌的工作,故此加利利是耶穌工作的起點與終點,亦是門徒受命與工作的起點.

究竟耶穌和門徒約定在那座山呢?從這裡的記載,沒有清楚指明,不過這座山似乎是有標誌性的.在馬太福音第五章有稱為「登山寶訓」的耶穌教導,同樣沒有提及山的名稱.在馬太福音十七章記載了耶穌帶著彼得,雅各,約翰上了高山,之後在山上變了形像,跟前面一樣也是沒有提及山的名稱.雖然不知道山的名稱,可是山卻隱含了最重要的教導與最重要的關係在其中,而當耶穌要跟門徒分別之前,就把最重要的命令在這山上頒佈.

上山學藝予人的感覺是非比尋常.那是進入另外一個境界之前,要經過必須的鍛鍊,提升能力,準備心志,在學成之日,就是下山之時.將苦學多年的學以致用,一展所長.到崇基學院神學院讀神學的同學也有這樣的感覺,幾年的神學訓練像是上山學藝:聖經研究,神學思考,研讀教會歷史,並有實踐嘗試,完成之後,下山成為上主所使用的僕人.

門徒在山上領受使命,看著基督升天,之後就是下山開展新的時代,將基督福音傳開.

任何人也可以是使徒

耶穌所呼召的門徒,並不如我們所想像,是社會中的精英份子,能力超卓,膽大包天的人.他們基本上都是十分平凡的人,對耶穌有不同期望,甚至可能是在計算可以從耶穌身上得著什麼好處的人.

這幾年我更有興趣想知道,在耶穌的門徒中,他們的政治光譜是怎樣的?在12門徒中,有管錢並且會「穿櫃桶」(即私下盜用公款)的猶大.有稅吏馬太,又名利未,明顯地跟羅馬政府有密切關係.更有奮銳黨西門,不惜以武力對抗羅馬政府,希望以色列國復興.這群極之不同的人同樣成為耶穌核心的門徒.他們一起相處,張力之大是我們今天完全可以想像到的.

耶穌準備離開他們了,他們在山上遇見耶穌,又敬拜他的,也有懷疑他的(第17節是這樣記載的).如果綜合其他福音書和使徒行傳,我們就可以估到,有些門徒對耶穌是否真的復活一直採半信半疑的態度.到了如今在山上集合的時候,按使徒行傳1章6節所載,門徒以為耶穌準備啟動復興以色列國的運動,顯示他們仍未真正認識耶穌的身份和使命.

原來耶穌所呼召和差遣的也是極為普通的人,或者更直接地說,耶穌呼召和差遣任何人.不論你怎樣看自己,配或不配,有能力還是沒有能力,充滿信心或是信心不足,你也可以被耶穌呼召和差遣,你也在任何人之列!

既然任何人也可以被呼召和差遣,關鍵就不在人,而是在於基督.因為他得著絕對的主權(第18節),那是天上,地下所有的權柄.在主禱文中有這樣的一句:「願你的旨意行在地上,如同行在天上」(馬太福音6章10節),這樣的禱告顯示了天上,地下因著罪的緣故而有隔閡,不然就不用這樣告祈求了.

現在罪的隔閡因基督的犧牲與復活,得以被修補.天上地下再不是分隔,這是回復上帝原初的創造,是一個新的創造經歷,亦解釋了為何今個主日經課會選用創世記1章1至2節的原因.

大使命的三重吩咐

使徒因著基督的權柄接受了差遣.這個差遣是委任,是命令,而這個差遣是包含三部份的.1.要去使萬民成為主的門徒；2.奉父,子,聖靈的名施洗；3.教導人遵守主所吩咐的一切.

要去使萬民成為主的門徒

「你們要去」就是使徒們要主動把福音傳開,而不是等待別人到來找他們.這個做法是有違慣常的情況.從古至今,大部份時間也是學生去訪尋名師,而鮮有老師去找學生.我相信這個情形在大學的研究院是十分普遍的.有學生如果想讀博士,他首先要做的不是寫研究計劃,而是要先確定是否有指導老師願意收他為學生.如果要老師去找學生,大概有兩種可能:其一是那位老師恐怕自己畢生所學無人繼承,所以要主動找尋學生.其二是那位老師不能吸引到學生注意,故此要周圍尋找學生.

香港小學和中學界早一段日子曾出現縮班,殺校危機,原因是適齡學童人數減少,以致學額過剩.為了保著學校的生存,學校要用盡各種方法吸引家長帶子女報讀.其實這個情況亦影響大學,不少學系也會面對報名人數不足,令學系得到撥款和其他的資源減少.為著保持競爭力和在大學生存,不少學系也要想辦法做宣傳,希望增加收生.

使徒不能等待人來認識有關耶穌基督的道理,他們必須出去,甚至乎到自己未曾到過的地方,這才符合「萬民」的要求.不過關鍵乃是在「門徒」這個字.

今天絕大部份的教會和基督徒也十分認真看待傳福音的工作,同時也會在教會的使命宣言提出建立門徒是核心的使命.可是實際上,教會有時會把重點放在數字上,而加入教會的人也忘記自己是要成為門徒.當重點放錯了,會有什麼結果出現呢?當我們看重人數,加入教會的人,可能只成為教會的「會員」(member).會員所著緊的可能就是福利,像是加入什麼購物會,健身中心等,因為有不同的優惠與好處,你才會參加.你要付出的就是你的金錢作消費.難怪有時候我們會以消費者的心態來評價教會,以服務滿意度來量度教會各樣事情,例如:講道是否有水準,設備是否足夠,一起崇拜的人是否友善等等.

「門徒」不是會員,不是消費者,成為門徒需要付出代價.對於使徒而言,他們十分清楚那是怎樣的一回事.面對強大羅馬帝國的權勢,跟隨基督就幾乎是等於把性命也押上.實踐天國的道理是世人難以理解和接受的,不然耶穌也不會被釘死了.

今日當我們要成為門徒,也希望使人成為門徒時,我們就要有心理準備,不會成為受歡迎的人物.你不一定得到物質的好處,更可能因為堅持真實的信仰而面對不同的挑戰.這就是作門徒的實況.

奉父,子,聖靈的名施洗

除了使人成為門徒這個吩咐外,奉父,子,聖靈的名施洗,明顯是早期教會在施洗時在禮儀行動中必有的部份.

受洗是一個公開認信的程序,也是委身與重生的象徵.水是生命來源的象徵,人從水中出來乃像生命新的開始.一個若認信跟隨基督,他進入水中再出來,那就是生活的新開始.這是重要的記號.

不知何鮮,有人會有這樣的錯覺,認為由誰為你施洗是十分關鍵.等級越高,施洗的效果越好.即是年輕的牧師不及資深的牧師,資深牧師不及主任牧師,主任牧師不及主教,而主教亦不及大主教.當然按教制不同,聖職是有等級,但若按基督的吩咐,重點乃在不是奉人的名去施洗,乃是父,子,聖靈的名.因為不是人使你得救,而是神才能拯救你.

教導人遵守主所吩咐的一切

大使命的第三重吩咐是要教導人遵守主所吩咐的一切.這大概就是馬太福音作者寫成這部福音的其中一個關鍵.他所記下一切有關耶穌的教導就是讓使徒把曾經聽過的也教導其他人.不是按他們的感趣,喜好,而是按他們親身所領受過的.

這看來是簡單,但其實一點也不簡單.如果你希望得到更多聽眾或跟隨者,你大概會講一些大部份人喜歡聽的說話.可是在耶穌的教導中,對不少人而言是難以接受的.單單講「要背起十字架」來跟從耶穌(馬太福音16章24節)已經是相當令人難以接受.

無怪乎,今日有不少傳道人也被批評,他們在傳講的是幸福音,而不是真正的福音.幸福音的重點是關乎個人的福祉,可以對周圍莫不關心.我們當不斷反問自己,究竟是教導人遵守主的吩咐,抑或是按著人的喜好來調整了主的吩咐.這是我們需要認真的反思.

基督的終極應許

當我們把焦點都放在大使命的吩咐時,我們容易忽略基督的終極應許.耶穌說:「我就常與你們同在,直到世界的末了.」

這個應許我們懂得背誦,但卻沒有仔細理解這是怎樣的一回事.馬太福音作者刻意提到耶穌說「我就常與你們同在」.這個其實是呼應馬太福音1章23節,耶穌被稱為「以馬內利」.這個名字的意思就是「上帝與我們同在」.

耶穌的降生是顯明「上帝與我們同在」,他縱然是復活後升天,不在世間,但因著他的門徒能遵守他的吩咐,世上每個人也會經驗到上帝仍然與我們同在.

今日我們是否讓人經驗到上帝仍然與我們同在呢?或許我們也在疑問,上帝真的與我們同在嗎?若上帝在,為何有許麼多的苦難,打擊臨到我們?

弟兄姊妹,這是上帝的問題,還是我們的問題呢?當我們沒有盡力去遵守主的吩咐,也未能讓萬民作主的門徒時,恐怕我們真的難以經驗上帝仍然與我們同在.不過主的應許是十分真實的,我們若真的奉基督的命而行,我們必然經驗到他的同在.

我認識不少的弟兄姊妹,他們信主的經歷多數在生命最低沉的時刻.在無助與孤單之中,有基督徒的朋友真誠相伴,雖然事情未必有轉變,但至少感覺有人關心,有人在支持,這是同行同在的力量,使人有力氣可以撐下去.

可以預計,未來的日子也許有很多難關,我們作為主的門徒,奉基督的命令去使人作他的門徒.讓人親身經驗上帝仍然與我們同在.

你是否已領了這個命令,你會奉命行事嗎?求主賜恩加力.阿們.



\newpage

\section{背在鷹的翅膀上}
\label{sec:505}
\textbf{劉國偉先生 - 背在鷹的翅膀上}
\newline
\newline
\hyperref[sec:504]{< < < PREV SERMON < < <}
~
\hyperref[sec:toc]{[返目錄]}
~
\hyperref[sec:506]{> > > NEXT SERMON > > >}


出埃及記 19:2-8
\newline
\begin{longtable}{cl}
\hline
\hline
章節 & 經文 (和合本修訂版)\\
\hline
19:2 & \begin{tabularx}{0.7\textwidth}{X} 他們從利非訂起程，來到西奈的曠野，在那裡的山下安營。 \end{tabularx} \\ \\ \relax
19:3 & \begin{tabularx}{0.7\textwidth}{X} 摩西到神那裡，耶和華從山上呼喚他說：「你要這樣告訴雅各家，對以色列人說： \end{tabularx} \\ \\ \relax
19:4 & \begin{tabularx}{0.7\textwidth}{X} 『我向埃及人所行的事，你們都看見了， 我如鷹將你們背在翅膀上，帶你們來歸我。 \end{tabularx} \\ \\ \relax
19:5 & \begin{tabularx}{0.7\textwidth}{X} 如今你們若真的聽從我的話，遵守我的約，就要在萬民中作屬我的子民，因為全地都是我的。 \end{tabularx} \\ \\ \relax
19:6 & \begin{tabularx}{0.7\textwidth}{X} 你們要歸我作祭司的國度，為神聖的國民。』這些話你要告訴以色列人。」 \end{tabularx} \\ \\ \relax
19:7 & \begin{tabularx}{0.7\textwidth}{X} 摩西去召了百姓中的長老來，將耶和華吩咐他的話當面告訴他們。 \end{tabularx} \\ \\ \relax
19:8 & \begin{tabularx}{0.7\textwidth}{X} 百姓都同聲回答：「凡耶和華所說的，我們一定遵行。」摩西就將百姓的話回覆耶和華。 \end{tabularx} \\ \\
[1ex]
\hline
\hline
\end{longtable}
講題:背在鷹的翅膀上  On Eagles’ Wings

講員:劉國偉先生

經文:出埃及記19:2-8

場合:香港中文大學崇基學院禮拜堂主日崇拜

日期:2020年6月14日

引言

弟兄姊妹平安!在一個三號強風信號的早上,一個低氣壓的生活場景中,我們在這裡有一扇可以透過十字架,遠望上天的窗戶,想像上帝可以陪伴我們,引導我們高飛天際,思念上面的事,所謂上面,就是有基督坐在上帝的右邊,那是上帝榮耀的寶座,是祂掌權的國度,我們追隨基督的人,活在世上的人,就是以自己的生命去遵行上帝的旨意,反映上帝的榮耀,見證祂的權能和國度,活在地上,卻見證上面的榮耀.

今天的講題也是與高飛上天相關連的:「背在鷹的翅膀上」.

剛才我們誦讀的舊約經課,出埃及記19章2至8節,記敘了以色列人出埃及之後大約三個月,來到西乃曠野,在那裡安營,並與神立約.摩西轉述耶和華曉諭祂的子民的說話:「看見我如鷹將你們背在翅膀上,帶來歸我.」(出19:4下)這就是今日講題的出處.

有人稱西乃曠野一帶地方,叫作「鷹之地」,那裡是飛鷹開展牠們寬闊有力的翅膀,在空中盤旋的地方.所以,以色列人離開埃及之後,來到西乃曠野,在山下的平原安營,在那個特別的場景,尤其是整個民族的特別的經歷,從埃及人的勞役中得到自由以後,以色列人在西乃山下,聽見上帝的信息,說祂像鷹一樣展翅,背著以色列人離開埃及,其實是一個充滿現場感的比喻,十分貼切,相信對當代以色列人來說,十分感動.

出埃及之後差不多四十年,摩西再一次用鷹作比喻,來形容上主與祂的子民的關係.「又如鷹攪動巢窩,在雛鷹以上搧展,接取雛鷹,背在兩翼之上.」(申32:11)

聖經多處地方提及過鷹,可以說是以色列人所熟悉的一種動物,用鷹作比喻去論說信仰,對以色列人來說,就是一個很好的把信仰融入生活的例子.

絕對「信得過」的上主

在我們進一步研讀今天的舊約經課前,我選擇從詩篇作者的信仰體悟開始.今天誦讀的詩篇100篇,是一篇稱謝詩.這詩篇是由93篇開始,一連8篇讚美和稱謝神作王作主的一個詩篇結集.

這篇詩可以分作前後兩半.前半部份是第1至3節,作者在這篇詩的開首,即邀請全世界的人向上主歡呼,他鼓勵我們要樂意事奉上主,當來向祂歌唱.和合本翻譯的「當來向祂」其實直譯的意思,就是來到祂所在的地方.現代中文譯本翻譯得直截了當:「到祂面前」.這固然是對朝聖客來到聖殿敬拜的鼓勵,不過,當中的信息和含意已經超越了純粹一個特定的地點,「到祂面前」是一種面對面的關係,是深度的認信,所以第3節如此寫道:「我們是祂做的,也是屬祂的,我們是祂的民,也是他草場的羊.」

如此事奉,敬拜上主,包含著當中深度的認信,正如摩西在出埃及記19章5至6節所表述耶和華曉諭的話:「在萬民中作屬我的子民……你們要歸我作祭司的國度,為聖潔的國民.」要成為來到聖殿敬拜的人一種持之以恆的信仰動力.正是如此認信的人,就會「稱謝進入祂的門；讚美進入祂的院.」(詩100:4上)作者在這篇詩的後半提醒我們,「上主本為善,祂的慈愛存到永遠,祂的信實直到萬代.」(詩100:5)

詩人告訴我們,上帝的信實直到萬代,祂是絕對「信得過」的上主,我們所投靠的是那位本為善,而且慈愛存到永遠的主.所以,我們要敬拜,事奉祂,並且在與上主面對面的關係中,累積我們的屬靈資產和動力,去面對世界的艱難和挑戰.

上帝向我們顯明的愛

信仰的路從來都充滿挑戰,我們明知這是進窄門,走窄路的生命歷程.使徒保羅也是如此鼓勵我們,他提醒我們生活中有患難,我們又如何可以歡歡喜喜在患難之中盼望上帝的榮耀?(羅5:2下)保羅告訴我們一個秘訣:「因為知道患難生忍耐,忍耐生老練,老練生盼望.」(羅5:3下-4)信仰之所以能夠演進到這境地,全是「因為所賜給我們的聖靈,將上帝的愛澆灌在我們心裡.」(羅5:5)

保羅在我們今日所誦讀的新約書信中,再一次提及因信稱義的核心,仍然在於上帝的愛,保羅告訴我們:「因我們還軟弱的時候,基督就按所定的日期為罪人死.為義人死,是少有的；為仁人死,或者有敢做的.唯有基督在我們還作罪人的時候為我們死,上帝的愛就在此向我們顯明了.」(羅5:6-8)

讀到保羅所寫的這幾節金句,實在給我們很大的鼓勵,按我們的本相和本性,不堪如我們,耶穌竟為我們死,就是我們還軟弱的時候,還作罪人的時候,耶穌為我們死,聖靈將這份愛澆灌在我們心裡,我們就可以站在上帝的恩典中,歡歡喜喜盼望上帝的榮耀.正如詩篇100篇的作者所說,神的慈愛存到永遠,就是這份永恆的愛,可以幫助我們渡過人間的困境.

如鷹攪動巢窩

剛才我們說過,以色列人出埃及之後差不多四十年,摩西再一次用鷹作比喻,來形容上主與祂的子民的關係:「又如鷹攪動巢窩,在雛鷹以上搧展,接取雛鷹,背在兩翼之上.」(申32:11)

申命記第32章是摩西呼天喚地,向以色列全會眾的充滿情懷的呼喊.弟兄姊妹,你可有認真想過「耶和華的分,本是祂的百姓；祂的產業,本是雅各.」(申32:9)是什麼意思?自始至終,上帝愛他的子民如同保護眼中的瞳人,所以摩西說:「耶和華遇見他在曠野荒涼野獸吼叫之地,就環繞他,看顧他,保護他如同保護眼中的瞳人.」(申32:10) 常人的想法,所謂環繞,看顧和保護,大概就是我們生活中的平順,安穩,豐裕和幸福；然而,上帝卻如鷹攪動巢窩,讓雛鷹離開安舒區,從懸崖峭壁高山崗嶺的巢窩中,把雛鷹推出去,在不著邊際的天空中下墜,在毫無把握的高速下滑中,學習飛翔.上帝的環繞,看顧和保護,在於我們自覺生活苦無出路,甚至非常危難之際,祂展開兩翅,有力地接取雛形,背在祂強大的兩翼之上.

上帝子民所走過的路絕不簡單,我們可能以為那個攪動巢窩的歲月,就是他們在西乃曠野40年的歷練；又或者我們以為是40年前,摩西與法老角力的時候開始,以色列人在勞役與危難之間,經歷上帝向埃及地所施行的十災；不過,最困難的起點,是有一位不認識約瑟的埃及的新的法老,起來治理埃及,轄制和苦害以色列人,那是他們在尼羅河三角洲東面的歌珊美地,在肥沃和安穩的土地上,繁衍了四個世紀之後的鉅變.

弟兄姊妹,回頭看看香港過去四十年,由上世紀中英談判香港前途開始,香港人在安定繁榮的旗幟底下追求生活,直至近幾年,特別是過去一年,我們所經歷的種種,我們可以肯定的說,我們看來活得不再安舒.

真的如此嗎?上主告訴我們,祂其實如鷹一樣,將我們背在他的翅膀上,並向我們說「帶來歸我」,又吩咐我們要在萬民中作祂的子民.(出19:4-5) 我們所追求的,是一個安舒的巢窩,還是上帝的吩咐和召命?

活出使命

今天所誦讀的福音經課如此說:「耶穌走遍各城各鄉,在會堂裡教訓人,宣講天國的福音,又醫治各樣的病症.」(太9:35)最觸動我的心是這一句:「祂看見許多的人,就憐憫他們；因為他們困苦流離,如同羊沒有牧人一般.」(太9:36) 耶穌形容要收的莊稼多,所以他要打發工人出去收莊稼.弟兄姊妹,你以為這個收莊稼的召命,只是停留於福音書所記載的十二門徒身上麼?

這份主日崇拜程序表上列寫了今天事奉人員的名字,我為到弟兄姊妹的服侍感恩.不過,我相信上帝給我們的召命,並不是停留在我們的禮拜堂的事奉人員名單,還有今天在這個禮拜堂當中座席的大家,以及透過電子平台與我們一同敬拜的弟兄姊妹,你們也是有名有姓的耶穌基督的門徒,耶穌的召命臨到每一個立志跟從祂的門徒.

這位為罪人死的慈愛的主,今天仍然呼召我們每一個跟從他的人「你們白白地得來,也要白白地捨去.」弟兄姊妹,靠著聖靈澆灌在我們心中的愛,也靠著主給我們的權柄,我們去趕逐污鬼,在罪惡掌權的世界,宣講天國的福音,又醫治各樣的病症,憐憫身邊的人,愛我們的鄰舍.

上帝攪動巢窩,除了熬練我們的生命,也是要挑戰和裝備我們,去為祂作生命的見證.願我們乘著母鷹的雙翅高飛天際,也學會向天際翱翔,即使有時乏力,受傷,總相信可以再次伏在上帝的強大的翅膀上,穩行高處.

求上帝賜我們力量,在艱難之中仍然確信可以高飛.



\newpage

\section{生命中不能承受的輕}
\label{sec:506}
\textbf{邢福增教授 - 生命中不能承受的輕}
\newline
\newline
\hyperref[sec:505]{< < < PREV SERMON < < <}
~
\hyperref[sec:toc]{[返目錄]}
~
\hyperref[sec:507]{> > > NEXT SERMON > > >}


耶利米書 20:7-13
\newline
\begin{longtable}{cl}
\hline
\hline
章節 & 經文 (和合本修訂版)\\
\hline
20:7 & \begin{tabularx}{0.7\textwidth}{X} 耶和華啊，你欺哄了我，我也被你欺哄了。你比我強，並且得勝。我終日成為笑柄，人人都戲弄我。 \end{tabularx} \\ \\ \relax
20:8 & \begin{tabularx}{0.7\textwidth}{X} 我每逢講話的時候，就哀嘆，我喊叫：「有暴力和毀滅！」因為耶和華的話終日成了我的凌辱和譏刺。 \end{tabularx} \\ \\ \relax
20:9 & \begin{tabularx}{0.7\textwidth}{X} 我若說：「我不再提耶和華，也不再奉他的名講論」，我心裡便覺得似乎有燒著的火悶在我骨中，我忍受不住，不能自禁。 \end{tabularx} \\ \\ \relax
20:10 & \begin{tabularx}{0.7\textwidth}{X} 我聽見許多的毀謗，四圍都是驚嚇；連我知己朋友都看著我跌倒：「告他吧，我們要告他！或者他被引誘，我們就能勝他，在他身上報仇。」 \end{tabularx} \\ \\ \relax
20:11 & \begin{tabularx}{0.7\textwidth}{X} 然而，耶和華與我同在，好像可怕的勇士。因此，迫害我的都絆跌，不能得勝；他們大大蒙羞，由於行事沒有智慧，必永遠受那不能忘懷的羞辱。 \end{tabularx} \\ \\ \relax
20:12 & \begin{tabularx}{0.7\textwidth}{X} 考驗義人、察看人肺腑心腸的萬軍之耶和華啊，求你使我得見你在他們身上報仇，因我已將我的案件向你稟明了。 \end{tabularx} \\ \\ \relax
20:13 & \begin{tabularx}{0.7\textwidth}{X} 你們要向耶和華唱歌！要讚美耶和華！因他救了窮人的性命脫離惡人的手。 \end{tabularx} \\ \\
[1ex]
\hline
\hline
\end{longtable}
羅馬書 6:1-11
\newline
\begin{longtable}{cl}
\hline
\hline
章節 & 經文 (和合本修訂版)\\
\hline
6:1 & \begin{tabularx}{0.7\textwidth}{X} 這樣，我們要怎麼說呢？我們可以仍在罪中使恩典增多嗎？ \end{tabularx} \\ \\ \relax
6:2 & \begin{tabularx}{0.7\textwidth}{X} 絕對不可！我們向罪死了的人，豈可仍在罪中活著呢？ \end{tabularx} \\ \\ \relax
6:3 & \begin{tabularx}{0.7\textwidth}{X} 難道你們不知道，我們這受洗歸入基督耶穌的人，就是受洗歸入他的死嗎？ \end{tabularx} \\ \\ \relax
6:4 & \begin{tabularx}{0.7\textwidth}{X} 所以，我們藉著洗禮歸入死，和他一同埋葬，是要我們行事為人都有新生的樣子，像基督藉著父的榮耀從死人中復活一樣。 \end{tabularx} \\ \\ \relax
6:5 & \begin{tabularx}{0.7\textwidth}{X} 我們若與他合一，經歷與他一樣的死，也將經歷與他一樣的復活。 \end{tabularx} \\ \\ \relax
6:6 & \begin{tabularx}{0.7\textwidth}{X} 我們知道，我們的舊人和他同釘十字架，使罪身滅絕，叫我們不再作罪的奴隸， \end{tabularx} \\ \\ \relax
6:7 & \begin{tabularx}{0.7\textwidth}{X} 因為已死的人是脫離了罪。 \end{tabularx} \\ \\ \relax
6:8 & \begin{tabularx}{0.7\textwidth}{X} 我們若與基督同死，我們信也必與他同活， \end{tabularx} \\ \\ \relax
6:9 & \begin{tabularx}{0.7\textwidth}{X} 因為知道基督既從死人中復活，就不再死，死也不再作他的主了。 \end{tabularx} \\ \\ \relax
6:10 & \begin{tabularx}{0.7\textwidth}{X} 他死了，是對罪死，只這一次；他活，是對神活著。 \end{tabularx} \\ \\ \relax
6:11 & \begin{tabularx}{0.7\textwidth}{X} 這樣，你們也要看自己對罪是死的，在基督耶穌裡對神卻是活的。 \end{tabularx} \\ \\
[1ex]
\hline
\hline
\end{longtable}
馬太福音 10:24-39
\newline
\begin{longtable}{cl}
\hline
\hline
章節 & 經文 (和合本修訂版)\\
\hline
10:24 & \begin{tabularx}{0.7\textwidth}{X} 「學生不高過老師，僕人不高過主人。 \end{tabularx} \\ \\ \relax
10:25 & \begin{tabularx}{0.7\textwidth}{X} 學生所遭遇的與老師一樣，僕人所遭遇的與主人一樣，也就夠了。既然有人罵一家的主人是『別西卜』，更何況他的家人呢？」 \end{tabularx} \\ \\ \relax
10:26 & \begin{tabularx}{0.7\textwidth}{X} 「所以，不要怕他們，因為掩蓋的事沒有不顯露出來的，隱藏的事也沒有不被知道的。 \end{tabularx} \\ \\ \relax
10:27 & \begin{tabularx}{0.7\textwidth}{X} 我在暗中告訴你們的，你們要在明處說出來；你們耳中所聽的，要在屋頂上宣揚出來。 \end{tabularx} \\ \\ \relax
10:28 & \begin{tabularx}{0.7\textwidth}{X} 那殺人身體但不能滅人靈魂的，不要怕他們；惟有那能在地獄裡毀滅身體和靈魂的，才要怕他。 \end{tabularx} \\ \\ \relax
10:29 & \begin{tabularx}{0.7\textwidth}{X} 兩隻麻雀不是賣一銅錢嗎？你們的父若不許，一隻也不會掉在地上。 \end{tabularx} \\ \\ \relax
10:30 & \begin{tabularx}{0.7\textwidth}{X} 就是你們的頭髮也都數過了。 \end{tabularx} \\ \\ \relax
10:31 & \begin{tabularx}{0.7\textwidth}{X} 所以，不要懼怕，你們比許多的麻雀還貴重！」 \end{tabularx} \\ \\ \relax
10:32 & \begin{tabularx}{0.7\textwidth}{X} 「所以，凡在人面前認我的，我在我天上的父面前也必認他； \end{tabularx} \\ \\ \relax
10:33 & \begin{tabularx}{0.7\textwidth}{X} 凡在人面前不認我的，我在我天上的父面前也必不認他。」 \end{tabularx} \\ \\ \relax
10:34 & \begin{tabularx}{0.7\textwidth}{X} 「你們不要以為我來是帶給地上和平，我來並不是帶來和平，而是刀劍。 \end{tabularx} \\ \\ \relax
10:35 & \begin{tabularx}{0.7\textwidth}{X} 因為我來是要叫『人與父親對立，女兒與母親對立，媳婦與婆婆對立。 \end{tabularx} \\ \\ \relax
10:36 & \begin{tabularx}{0.7\textwidth}{X} 人的仇敵就是自己家裡的人。』 \end{tabularx} \\ \\ \relax
10:37 & \begin{tabularx}{0.7\textwidth}{X} 愛父母勝過愛我的，不配作我的門徒；愛兒女勝過愛我的，不配作我的門徒。 \end{tabularx} \\ \\ \relax
10:38 & \begin{tabularx}{0.7\textwidth}{X} 不背自己的十字架跟從我的，不配作我的門徒。 \end{tabularx} \\ \\ \relax
10:39 & \begin{tabularx}{0.7\textwidth}{X} 得著性命的，要喪失性命；為我喪失性命的，要得著性命。」 \end{tabularx} \\ \\
[1ex]
\hline
\hline
\end{longtable}
講題:生命中不能承受的輕 Unbearable Lightness of Being

經課:耶利米書20:7-13;羅馬書6:1-11;馬太福音10:24-39

講員:邢福增教授

場合:香港中文大學崇基學院禮拜堂主日崇拜

日期:2020年6月21日

生命的選擇

今主日的講題,是來自捷克文學家米蘭昆德拉(Milan Kundera)的小說──《生命中不可承受的輕》(The Unbearable Lightness of Being),早年曾改編為電影「布拉格之戀」.小說是以1968年蘇軍入侵捷克,鎮壓布拉格之春民主運動為背景的愛情故事,帶出值得深思的問題:人生所作的決定,最理想當然是可以把每一個可能的選擇都試過,然後再比較到底那一個選擇是最好的.但現實是,這根本是不存在的.人生只有一次,我們永遠不知道這選擇是否比其他選擇更好或更壞,對不對?

弟兄姊妹,人生難免要作選擇,有些選擇無關痛癢,比如今早在家上網崇拜,還是來到現場(甚至別教會)；或中午到哪食飯……但有些選擇,卻令人感到極其沉重,擔心選錯了,就錯恨難返.選擇背後,我們感受到輕與重的拉扯,這正是貫穿今主日經課的主題.

假使召命原來不像你預期

有時,我們明明滿懷自信地做某個選擇,但慢慢發現,事情卻不像原先般發展.正如一首頗令人有共鳴的流行曲──「假使世界原來不像你預期」,歌詞第一句:「容我割開這肉身,拿下碎的心重嵌」,完全將那份內心的痛表白出來.

《耶利米書》廿章,正是一位先知突然發現,世界,其至是上主的召命原來不像預期時,發自內心掙扎的獨白.其實,這也是向耶和華的呼喊.他認為自己的遭遇,是被耶和華「欺哄」.天主教《思高版聖經》將「欺哄」釋作「誘惑」.耶利米對被召作先知,確有一定期望,與其說這是主觀的想像,毋寧說是耶和和華給他的「合理期望(誘惑)」.但現在他發現「世界原來不像他預期」……先知不得不發出「我也被你欺哄了」的埋怨,深感自己信錯了耶和華!

那麼,耶和華如何呼召耶利米?讓我們看看《耶利米書》一章的記載:在耶路撒冷被擄前,耶和華的話臨到耶利米:「我尚未將你造在母腹中,就已認識你；你未出母胎,我已將你分別為聖,派你作列國的先知.」(一5)「你不要怕他們,因為我與你同在,要拯救你.這是耶和華說的.」(一8)

顯然,耶利米發現自己被騙上了賊船.在二十章開始時,耶利米按耶和華吩咐說預言後,便受到聖殿的總管祭司巴施戶珥毆打,並將他鎖在聖殿北面的便雅憫門.雖然在第二日,耶利米繼續對巴施戶珥宣講耶和華的話(二3-6),但第7節開始,耶利米立即向耶和華表達不滿,「我終日成為笑柄,人人都戲弄我.」先知的職分對耶利米而言,變成了極大的羞辱.

在8至10節,耶利米剖白其內心世界:「我每逢講話的時候,就哀嘆,我喊叫:『有暴力和毀滅!』因為耶和華的話終日成了我的凌辱和譏刺.」(廿8)「我聽見許多的毀謗,四圍都是驚嚇；連我知己朋友都看著我跌倒」(廿10a).他恨不得「不再提耶和華,也不再奉他的名講論」(廿9a),這似乎是令他可以輕鬆的解脫,但這樣他又面對內心極大爭戰:「我心裡便覺得似乎有燒著的火悶在我骨中,我忍受不住,不能自禁」(廿9b).面對外在環境(暴力,毀滅,凌辱,譏刺,毀謗,驚嚇),唯一的選擇就是「劈炮」(辭職),但耶利米內心卻仍有一團火在燃燒,迫使他再次反思召命.

在拉扯中的堅持

在拉扯之中,耶利米如何選擇?「忍受不住,不能自禁」催迫耶利米繼續去發聲.那麼,他如何面對這不似預期的世界?耶利米只能回想耶和華昔日的應許──「同在」.他在11節筆鋒一轉,說「然而,耶和華與我同在,好像可怕的勇士」.他以「可怕的勇士」來形容耶和華,「可怕」是向著那些迫害先知的人的,因著耶和華的同在,耶利米深信,迫害他的都要「蒙羞」.這位耶和華,又是一位「考驗義人,察看人肺腑心腸的萬軍之耶和華」(廿12),因此,耶利米的痛苦及掙扎,耶和華一定知道,最後,他不得不發出讚美(廿13).

表面上看來,經文由不滿,埋怨到信靠及讚美,似乎很戲劇化.對於同樣面對類似拉扯及掙扎的我們,當然不是機械公式化地應用,而需要真實面對其中的張力.當我們再看14至18節時,見到耶利米甚至還在咒詛自己的生日,作生不如死的控訴.所以,耶和華的信實,並沒有馬上將外在威脅消除,先知仍真真實實地面對著暴力,毀滅,凌辱,譏刺,毀謗,驚嚇……但即或如此,耶利米卻選擇繼續承擔召命,也承受伴隨著召命而來的各種挑戰,而讓他可以堅持下去的,是抓緊耶和華的同在.耶利米好像在輕與重之間作出了選擇,但與其說這是選擇,倒不如說是他的堅持.接下來廿一章,耶和華的話再次臨到,耶利米便遵命,宣告耶路撤冷將被毀……

拒絕活在罪中

面對一個不像預期的世界,先知曾想過放棄召命,而比這更徹底的,是改變自己,完全與這個世界認同.

《羅馬書》六章的經課,保羅在處理基督徒的生活,特別是如何面對罪的問題.他在第五章討論了「因信稱義」後,在第六章一開始便問:如果因信稱義代表恩典,那麼,基督徒是否可以繼續活在罪中,藉此證明恩典的偉大?對這問題,保羅並不是說:「感謝主,真係好!」而是斬釘截鐵地回答:「絕對不可」.為何保羅會拋出這問題?真的有基督徒會以犯罪來證明恩典嗎?還是,基督徒因著自己的軟弱,卻又無法擺脫活在罪中的生活,因而需要找個更神聖的理由來為自己開脫?

保羅處理的,同樣是理想與現實的差距問題,如何避免基督徒陷在罪惡之中?如果擺脫不了罪,倒不如接納在罪中活的自己,是否更符合現實?再以上帝為理由為自己站台,更是無懈可擊!這種做法,表面上仍肯定及高舉恩典,但實際上,無疑是向罪降服,成為「罪的奴隸」,完全跟罪惡的現實認同.

保羅指出,基督徒必須認定自己的身分,這是新生活的基礎:基督徒的生命要跟基督連結在一起,關鍵在於,每一位基督徒已經藉洗禮歸入基督的死,再與基督「一同埋葬」,成為「向罪死了的人」(六3-4).如果與基督同死同埋葬是已實現的事實,那麼,接下來便要迎接那仍未實現的,就是與基督一同復活.這是將來的榮耀,藉著盼望來支撐基督徒當下的生活(六5).因著這盼望,基督徒要宣告,「舊我」已經死去,「不再作罪的奴隸」,就是說,生命可以「脫離」罪惡權勢的轄制(六6-7).更重要,基督徒的生命,成為與基督同活的生命,這也是「在基督耶穌裡」對上帝活的生命,就是向上帝負責的生命(六8-11).

我們與惡的距離

早前,台灣有電影劇「我們與惡的距離」,深刻地揭示出人性真實的一面.導演林君陽曾說:「我們都以為自己是好人,但我們很多時候都有惡念而不自知,有些人的惡念展現出來,有些人沒有,而每個人的善念與惡念都主導了現在這個世界的面貌.」我相信,作為基督徒,我們應對此有更深的體會.是的,我們活在一個惡的世界之中,基督徒並不是「善」的化身,而是覺醒到「惡」的真實,在善惡掙扎中,但仍致力追求善的人.保羅後來真的說:「感謝主呀」,但在這句前,他卻說:「立志為善由得我,只是行出來由不得我」(羅七18b),「我真苦啊!誰能救我脫離這必死的身體呢?」(七24),接著才是:「感謝上帝,靠著我們的主耶穌基督就能!」(七25a) 原來,我們與善的距離是多麼的遠,我們的惡的距離又是多麼的近……

那麼,基督徒的出路在那裡?答案就是因信稱義的恩典.唯有靠賴基督赦罪恩典,我們才能面對自己.雖然我們仍然活在這個罪惡的世代,仍會面對罪惡的誘惑,仍然會有軟弱的時候,但當基督徒認清自己的身分,我們的生命已經被主的寶血救贖,是向上帝負責,而不是作罪奴隸的生活；與主同活,在善與惡交戰之中,不斷依靠上主的力量,去克勝罪惡,面對自我,盼望將來的榮耀.

作主門徒,不要怕

是的,面對理想與現實的矛盾,活在一個原來不似預期的世界之中,基督徒要堅持召命,要活出與自己身分相稱的生活,真的很沉重,很辛苦,甚至還要面對外在環境更猛烈的打壓.這是《馬太福音》經課中耶穌向門徒的教導.耶穌在揀選及差遣十二位門徒說:「看哪!我差你們出去,如同羊進入狼群」(十16),「你們要為我的名被眾人憎恨」(十22)

耶穌對即將「送羊入狼口」的門徒說:「所以,不要怕他們」.「他們」是誰?就是那些痛恨耶穌的宗教權威及政治領袖.因為耶穌的教導,既顛覆了猶太教傳統,也被視為威脅現有政權；他們務要將耶穌除掉.但耶穌先後說:「不要怕他們」,「不要懼怕」(十26,28,29)耶穌知道,日後門徒將面對更大規模的逼害與挑戰,門徒需要勇敢地公開為主作見證:「你們要在明處說出來」,「要在屋頂上宣揚出來」,在在需要極大的勇氣(十26-28)；要「在人面前認我」(十32),並且要過一個「配作我的門徒」的生活(十37-38),就是「背自己的十字架」來跟從耶穌的生活(十38).

門徒背十字架,當然不完全是指向被釘在十架上的刑罰,而是以此來象徵各人要為信仰付上代價.「代價」好像予人負面印象,但耶穌卻通過「得著性命的,要喪失性命；為我喪失性命的,要得著性命」(十39)的對比,吊詭地指出,背十架的背後,其實是每位門徒要回答的問題:為了耶穌基督,為了這個信仰,你願意捨棄自己的生命,或是生命中看為貴重的人或事嗎?

身處此時此地的香港,這段經文對我們有何有意義呢?我們會否因公開為信仰作見證而被逼害,甚至捨命?我們內心會否懼怕表達或堅持自己的信仰立場?還是我們根本沒有可以表達或堅持的立場與信念?在我們生命中,耶穌基督是否勝過一切?還是有太多的人與事,佔據了我們的心?「有事鍾無艷,無事夏迎春」?我們是不敢為信仰付代價?還是將某些東西置於更高的位置,將上帝邊緣化?口裡說「為了信仰」,骨子裡卻是另一回事?

是與不,輕與重

德國神學家潘霍華(Dietrich Bonhoeffer)被納粹德軍處決後,他的好友英國哲齊斯特主教(Bishop of Chichester)貝爾(George Bell)曾說:「潘霍華說過:『當基督呼召一個人,祂召他來為祂死.』死的形式有多種,不過作基督門徒的真諦可說是完全包括在這句話裡面了.潘霍華在受刑之前已經殉過許多次的道.」( ‘When Christ calls a man’, says Dietrich Bonhoeffer, ‘he bids him come and die.’ There are different kinds of dying, it is true; but the essence of discipleship is contained in those words. And the marvelous book is a commentary on the cost. Dietrich himself was a martyr many times before he died. )對,他的殉道不僅是他的被捕與行刑,正如蘇恩佩姊妹說:潘的殉道早在他在意志上對基督說「是」,對自己說「不」的時候就開始了!

相信我們都不想成為殉道者,但弟兄姊妹,我們最近是否曾對基督說「是」,對自己說「不」?還是對自己說「是」,對基督說「不」?對基督說「是」,對自己說「不」,是考驗我們,是否願意順服上帝的主權,也就從自我中心改為以上帝為中心的生活.即或這個世界原來(仍然)不似預期,但卻堅持以天國的價值來重新審視自己,以及與其他人的關係.甚麼是我們要堅持的?甚麼是我們的底線?這問題是要求我們去回答的.

米蘭昆德拉在小說中提出「非如此不可!」(It must be)的拷問.他說:「只有必需的東西才是重的,只有重的東西才有價值」(only necessity is heavy, and only what is heavy has value).無疑,這是沉重地負起自己命運的選擇,昔日的先知及門徒,也面對同樣的選擇,還要(能)揹起生命的「重」嗎?潘霍華在《獄中書簡》(Letters and Papers form Prison)說:「我們必須完全過著今生的生活,才能學習到信心.」(only learns to have faith by living in the full this-worldliness of life)他所指的今生,就是「負起生命的一切責任與困難,成功與失敗,一切經驗與無可奈何之事.就在這樣的生命中,我們才把自己無條件地放置在上帝手裡,參與祂在世上的苦難,與基督在客西馬尼園一同警醒」(living fully in the midst of life’s tasks, questions, successes and failures, experiences, and perplexities-then one takes seriously no longer one’s own sufferings but rather he suffering of God in the world. Then one stays awake with Christ in Gethsemane).是的,這種負起真的十分沉重……

昆德拉說:「最沉重的負擔同時也是最激越生命實現的形象.負擔愈沉重,我們的生命就愈貼近地面,生命就愈寫實也愈真實.相反的,完全沒有負擔會讓人的存在變得比空氣還輕,會讓人的存在飛起,遠離地面,遠離人世的存在,變得只是似真非真,一切動作都變得自由自在,卻又無足輕重.那麼,我們該選擇哪一個呢?重?還是輕?」(The heavier the burden, the closer our lives come to the earth, the more real and truthful they become. Conversely, the absolute absence of a burden causes man to be lighter than air, to soar into the heights, take leave of the earth and his earthy being, and become only hale real, his movements as free as they are insignificant. What then shall we choose, Weight or lightness?)米蘭昆德拉的問題很吊詭:選擇不可承受的「輕」?還是有價值的「重」?耶穌的回答是:「得著性命的,要喪失性命；為我喪失性命的,要得著性命」(太十39)

容許我再引用「假使世界原來不像你預期」的歌詞來結束:「當一切完全不像你預期,才明白不可心死,沿路再走幾千公里,要做最好的你.假使世界原來不像你預期,才明白不要放棄,沉住愛反覆的心理,風雨不會沒了期,終於會等到夢寐.全城在變遷,不減你是你.」我們活在一個不似預期的世界之中,也需要在其中作出各種的選擇,求主幫助我們,曉得堅持怎樣的信念與價值,過怎樣的生活,作怎樣的基督徒(人).



\newpage

\section{惡待善待恩待等等接待的信仰意義}
\label{sec:507}
\textbf{鄭漢文博士 - 惡待,善待,恩待等等……接待的信仰意義}
\newline
\newline
\hyperref[sec:506]{< < < PREV SERMON < < <}
~
\hyperref[sec:toc]{[返目錄]}
~
\hyperref[sec:508]{> > > NEXT SERMON > > >}


馬太福音 10:40-42
\newline
\begin{longtable}{cl}
\hline
\hline
章節 & 經文 (和合本修訂版)\\
\hline
10:40 & \begin{tabularx}{0.7\textwidth}{X} 「接納你們的就是接納我；接納我的就是接納差遣我來的那位。 \end{tabularx} \\ \\ \relax
10:41 & \begin{tabularx}{0.7\textwidth}{X} 把先知當作先知接納的，必得先知的賞賜；把義人當作義人接納的，必得義人的賞賜。 \end{tabularx} \\ \\ \relax
10:42 & \begin{tabularx}{0.7\textwidth}{X} 無論誰，只因門徒的名，就算把一杯涼水給這些小子中的一個喝，我實在告訴你們，他一定會得到賞賜。」 \end{tabularx} \\ \\
[1ex]
\hline
\hline
\end{longtable}
講題:惡待,善待,恩待等等……接待的信仰意義 Fideistic Meanings of Receptive Responsibility in Bad, Good and Gracious Treatment

經文:馬太福音10:40-42

講員:鄭漢文博士

場合:香港中文大學崇基學院禮拜堂主日崇拜

日期:2020年6月28日

1. 我原是過客,謝謝大家的恩待,今日反思接待的信仰意義

弟兄姊妹們,你們好,大家平安.

今天我能夠面對面講道,真好.這樣的相見,對我來說,日子無多.世事多變,世事難料.

今天好大可能會是我在崇基禮拜堂最後一次講道,至少是我以榮譽助理校牧身分講道最一次,是我的告別講道(farewell sermon).

會否真的在6月28日的今天告別這講壇?他朝會否再站在這裡再講道?將來的事,將來再說.

我原來是個過客,謝謝大家的恩待,今日有意和大家反思「接待的信仰意義」(Fideistic Meanings of Reception).

今天的關鍵字是接待的「待」,對待的「待」,款待的「待」,等待的「待」.

「待」這個字可以有那些構詞組合的可能呢?可以給與我們信仰經驗的意義得以詮釋解讀.

我們今日經課只有三節經文,是馬太福章10章40-42節,我用和合本版本.我先讀一次:

40「人接待你們就是接待我；接待我就是接待那差我來的.

41人因為先知的名接待先知,必得先知所得的賞賜；人因為義人的名接待義人,必得義人所得的賞賜.

42無論何人,因為門徒的名,只把一杯涼水給這小子裡的一個喝,我實在告訴你們,這人不能不得賞賜.」

我打算用五個脈絡(contexts)來解讀這個文本(text).這五個脈絡構成五個視點,透現我這個惜別講道(farewell sermon):

第一,是馬太福音第十章:這三節經文的上文在說甚麼?以致以這三節作為結語.

第二,是路加福音6章27-38節:當中耶穌提到惡待與善待,得出相待的四種可能型態.

第三,是普世倫理(Universal Ethics)的兩個律:己所不欲勿施於人,以及己所欲施於人.

第四,是「接待」的神學哲學:我提出「接待回應力」(receptive responsibility)作為哲學與神學可資應用在教育的觀念.

第五,是我個人經驗的印證:回顧四十年來情與理,以及分享我未來動向的敘事想像.

若有聖經,可以打開馬太福音,從頭到尾讀一次第十章,由第1至39節的論說,如何抵達今天40-42三節經文作為結語?

馬太福音第十章被解經家稱為「使命差遣篇」(chapter on commission),就是耶穌基督親自揀選十二個門徒,把使命託付給他們(mission being commissioned),差遣他們實現這個使命:

7隨走隨傳,說『天國近了!』

為此他們將會被賦予靈界的能力,就是:

8醫治病人,叫死人復活,叫長大痲瘋的潔淨,把鬼趕出去.

傳道,治病,趕鬼.是三大使命,也是三大天國近了的福音.

門徒(disciples)與使徒(apostles)有何分別呢?

跟隨耶穌認他為基督的都是門徒,基督門徒就是基督跟隨者(Christ-followers),之後被為稱基督徒(Christians).

而使徒只能是耶穌基督親自揀選的門徒,由他親自把使命託付給你,差遣他們你實現這個使命.英文就是Jesus Christ commissions you as his commissioner for a mission.

四福音只記載了耶穌基督親自揀選的門徒成為使徒的只得十二位.後來出了個保羅,他自稱他在大馬色路上耶穌基督親自揀選了他實現傳道給外邦人使萬國萬民信主的使命.

上次5月31日講道〈靈動我心力行奇事,言說基督見證上帝〉,我已概論了整本《使徒行傳》.但《行使徒傳》裡又沒有記載所有十二位使徒的行動,卻又記載了許多有不此清單上的其他使徒的心動力行成奇事.

所以,門徒與使徒其實不是二件事,而是一個連續體,按跟隨主看是門徒,按使命看是使徒.

耶穌基督在使命差遣前預告了使徒會面對的強弱機危,並早給了應對的錦囊.

活出使命的門徒要面對的危機,從第16至39節提到有以下四種:

一是政治危機:會被送交會堂君王；但聖靈會賜話語,會得作見證.

二是家庭危機:仇敵是家人,地上動刀兵；只管一心愛耶穌,過於家人.忍耐到底的必然得救.

三是抹黑攻擊:罵耶穌是鬼王；不要怕,因為隱藏的事沒有不被人知道的.

四是死亡威脅:要殺你身體；不要怕,他們不能殺你靈魂.

耶穌教導手法的總念是在第16節:「我差你們去,如同羊進入狼群；所以你們要靈巧像蛇,馴良像鴿子.」

耶穌教導心法的總念是在第38-39節: 「38不背著他的十字架跟從我的,也不配作我的門徒.39得著生命的,將要失喪生命；為我失喪生命的,將要得著生命.」

原來做耶穌門徒,是要為耶穌賣命的.這個福音從來不是廉價的.

經文到此,忽然轉入賞賜的討論作結,為何如此?

而且40至42這三節轉到另一視點,就是從接待門徒的接待者視點看.

解讀的方向是,耶穌所預告給十二門徒實現使命時,會面對不同的對待.會出現四種的惡待,當然也會遇上善待門徒的.善待你的人是配得賞賜的.他們得到甚麼賞賜?為何如此?

所以,現在我們進入今天講道題目的前半部份,談一談被接待時可能出現三種被對待的方式,就是惡待,善待,恩待.

用價值學的語言,對待人可以出現高,中,低三種待人之法(3 ways of being treated: badly treated, well treated, graciously treated).

對待人的三個型態之從低等,經中等,到高等,就是:惡待,善待,恩待.

2. 相待的四個型態:惡待惡待,善待善待,惡待善待,善待惡待

我曾經在2019年2月24日到自己母會循道衛理聯合教會的天水圍堂講道,是日是社會關懷主日,經課是路加福音6章27-38節,講道題目是:《惡待,善待,恩待中得見罪與愛》.

我自己覺得講章寫得好精彩.我現在決定善待大家,所以把講章吸收在這裡,與大家分享,也因而可以流傳.

從路加福音6:27-38我們可以惡待,善待,恩待之低中高得出相待的四種型態.

分節:

(1) 相待的四種型態；惡待,善待與恩待中得見罪與愛

(2) 善待那善待自己的人:黃金第二律(路6:31正讀)

(3) 惡待那惡待自己的人:黃金第一律(路6:31反讀)

(4) 惡待那善待自己的人:忘恩負義,恩將仇報的罪人(路6:32-34)

(5) 善待那惡待自己的人:愛你仇敵(路6:27-30),像恩待我們罪人的慈悲的父(路6:35-36)

讓我先讀出路加福音6章27-38節這段經文:

論愛仇敵(太5．38－48；7．12)(和修本)

27「可是我告訴你們這些聽的人,要愛你們的仇敵!要善待恨你們的人!

28要祝福詛咒你們的人!要為淩辱你們的人禱告!

29有人打你的臉,連另一邊也由他打.有人拿你的外衣,連內衣也由他拿去.

30凡求你的,就給他；有人拿走你的東西,不要討回來.

31「你們想要人怎樣待你們,你們也要怎樣待人.

32你們若只愛那愛你們的人,有什麼可感謝的呢?就是罪人也愛那愛他們的人.

33你們若善待那善待你們的人,有什麼可感謝的呢?就是罪人也是這樣做.

34你們若借給人,希望從他收回,有什麼可感謝的呢?就是罪人也借給罪人,再如數收回.

35你們倒要愛仇敵,要善待他們,並要借給人不指望償還,你們的賞賜就很多了,你們必作至高者的兒子,因為他恩待那忘恩的和作惡的.

36你們要仁慈,像你們的父是仁慈的.」

彼此相待的四種型態:惡待,善待與恩待中得見罪與愛

關於待人的態度,在今日經文路加福音6章中27-38節找到三個關鍵字:惡待,善待與恩待.我們用這三個關鍵字來思考這段經文.

讓我們先看一看人與人相待有多少可能?

首先,對惡待自己的人,我可以還以善待或惡待他這兩種可能.

同樣,對善待自己的人,我可以還以善待或惡待他這兩種可能.

二乘二等於四.

從待人之有善惡二分,可以產生彼此相待的四種型態:

\texttt{---------------------------------------------------}

\texttt{|                                                 |}

\texttt{| 善待那善待自己的人:  善待 --------------> 善待 |}

\texttt{|                         \          /            |}

\texttt{| 惡待那惡待自己的人:     ----\ /--              |}

\texttt{|                               x                 |}

\texttt{| 惡待那善待自己的人:     ---/  \---\            |}

\texttt{|                         /           \           |}

\texttt{| 善待那惡待自己的人:  惡待 --------------> 惡待 |}

\texttt{|                                                 |}

\texttt{---------------------------------------------------}

讓我分四小節來對解讀這段經文.

從而看看如何在惡待,善待與恩待中得見罪與愛.

從而反思在善待和惡待以外,在甚麼情況下,會出現恩待這另一種態度?

善待那善待自己的人:黃金第二律(路6:31正面讀)

第一種彼此相待的型態,就是善待那善待自己的人.

讓我們先看一看今日經課路加福音第6章第31節:

31你們願意人怎樣待你們,你們也要怎樣待人.

在道德倫理上,耶穌這句話是在表述一個叫做普世倫理的黃金律(Gold Rule of Universal Ethics)之第二律:「己所欲,施於人.」這是從正面的,積極的角度去陳構一個倫理.

即是自己認為好的事,也作在他人身上.自己心裡所意願所想望的,體察其他人都同樣會意願會,想望的.

例如:教會有一百人共食愛筵,添加了三文魚壽司約一百件.你好想食,心裡想若排在尾就慘,自己不希望別人多取到無餘而自己無得食,將心比心,所以自己若排隊排得前也只取一件,等每個人取食物的機會都至少過了一次,見壽司有餘才再多取.

中國古語說,幼吾幼以及人之幼,老吾老以及人之老.推己及人,愛屋及烏.

例如:你好倦時上了公車,其實好想有個位坐下,你想像有人體會你的需要而讓座給你你會好感激.若你見到比你更有需要的人,你體會到他也想有個位坐下,你也許會讓座.我恭喜你,你在善待一個有需要的陌生人.

那麼,為甚麼要有「關愛座」?

因為人不一定能做得到「己所欲,施於人.」所以,需要設立一些制度,建立一些做法,讓大家願不願意都好,都會跟隨一種「讓座給有需要的人」的做法.

但甚麼叫做「有需要」?只讓座給「老弱傷殘」嗎?這些人的需要是明眼看得到的?(盲都睇到啦!)但有些需要,真的是留神都未必睇得到的.

例如:一個學生剛剛考試不及格被老師罵又懶又蠢,又被同學恥笑……

例如:一個少女剛剛得知男朋友是同時與自己的閨密拍拖而鬧分手……

例如:一個年青人剛剛被忽然解僱但家有老人由他供養……

例如:一個媽媽剛剛得知驗血報告,得知絕症死期而想念家中兩個幼孩……

你知道若你落入這樣的人生光景,上了公車,真的好倦,好想坐下,低頭哭泣.這時,讓座給你的人真的是對你最大的恩待.

愛你的鄰舍,留意他的需要,主動助人.

惡待那惡待自己的人:黃金第一律(路6:31反面讀)

其實,剛才只是普世倫理黃金律(Gold Rule of Universal Ethics)的第二律,而黃金律是先有,還有第一律:「己所不欲,勿施於人.」這是從反面的,消極的角度去陳構一個倫理.

即是自己認為不好的事,不要作在他人身上.你不想他人待你以某種令人難受的方式,就不要以同樣方式待人.

就是把路加6章31節這個正向要求倒轉來讀,讀出一個反面的教導來.從積極角度轉為消極角度:由do轉為don’t,由該做轉為不該,轉化後的第一律,踐行起來相對較容易.

其實「己所不欲,勿施於人.」這句話,是來自孔子的《論語》.

人有自然性,就是你既然對我不好,我當然合情合理地對你不好.所以,惡待那惡待自己的人,好像是人之常情.

但想深一層,其實我不想別人對我不好,當我這樣想時,我就不要對別人不好.

例如:你晚上好怕鄰舍好嘈,嘈到你到午夜都無法睡,你就體諒家中各人不會因你晚上不睡還唱歌睇電視而不得安睡.

例如:愛鄰舍,不要把垃圾放在鄰居的地方.

例如:你不想他人侵犯你的私隱,你就不要侵犯他人的私隱.所以,在教會,在團契分享甚麼都需要保守秘密,不要講給不在場,不相關的人.

自由社會的最基本道德是自由的尊重.你是自由的,大家都尊重你的自由.但若你自由的運用是損害他人的自由,你的自由就受到限制,以至於被關入獄.

講起尊重,讓我們談一談對他人身體的尊重.最近爆出好多未有尊重他人身體的案例,有些已經證實,有些有待確證,我不是想談個別事件,而是提醒大家,男女授受不親,男性不宜隨便觸踫女性的身體,女性也小心不要容讓男性落入危困.

除了有意侵犯他人的男性,或有意引人犯罪的女性.許多時人不自覺地落入如此狀況,也許因為忽略「男女太不同」.

女性的己所不欲(接受程度),即被觸踫的程度或模式就是侵犯(例如因為開心事而手握手甚至擁抱),可能是男性之為男性所想像不到的,想像不到女性之所不欲(不能接受).同理,男性的己所不欲(親近程度)(例如因為開心事而手握手甚至擁抱),即被吸引的程度或模式就是邀請親近,可能是女性之為女性所想像不到的,想像不到男性之欲罷不能.

這段話如果可能大家聽不明白我講甚麼,就由他吧.總而言之,我不想別人對我身體有何不尊重,我最好不要不尊重他人的身體.

身體是聖靈的殿,每個人都是守殿人.

3. 待人的三個型態之低中高:惡待,善待,恩待

惡待那善待自己的人:忘恩負義,恩將仇報的罪人(路6:32-34)

其實,做到以上兩個倫理的人,已經在人世間不得了,已經難能可貴.

但是,若我們細看路加6章32-34節,可以見到耶穌並不滿意此二者,因為罪人都會這樣做,都會做到以上二者.讓我一起再讀耶穌的話:

32你們若只愛那愛你們的人,有什麼可感謝的呢?就是罪人也愛那愛他們的人.

33你們若善待那善待你們的人,有什麼可感謝的呢?就是罪人也是這樣做.

34你們若借給人,希望從他收回,有什麼可感謝的呢?就是罪人也借給罪人,再如數收回.

換言之,罪人都會善待那善待自己的人,以及惡待那惡待自己的人.

耶穌看透人性,就算罪人都會這樣做,做到以上二者.So what?　咁有點je ,好威咩!

當然,因為人的罪性,未必做到以上二者.但做到以上二者的,依然可以是個罪人.

因為我們可能在踐行這二律時,依然會惡待那善待自己的人.就是第三種型態.

因為人的罪性,我們甚至忘恩負義,恩將仇報.

上主已經恩待人,但人依然忘恩作惡.

你看耶穌明明是神的兒子是基督,是來打救世人的,卻被當時的人所釘死.

我們明知耶穌為我們而死,且死在十字架上,為的是讓神的恩典臨到我們身上可以讓罪得赦而與神和好,但我們偏偏要偏行己路.

普世二律,是講人的倫理.耶穌要講的,是神的倫理.俗語說:「正一唔係人咁品.」

行此二律,在人眼中,已是好人.

行此二律,在耶穌眼中,在神眼中,依然是罪人.因為罪人都會善待那善待自己的人,以及惡待那惡待自己的人.

善待那惡待自己的人:愛你仇敵(路6:27-30),像恩待我們罪人的慈悲的父(路6:35-36)

在路加福音今日的經課,耶穌給我們要講究的,是一個更高的倫理:「愛你仇敵」.

回顧一下,從待人之有善惡二分,就此可以產生彼此相待的四種型態:

一．善待那惡待自己的人

二．惡待那惡待自己的人

三．惡待那善待自己的人

四．善待那惡待自己的人

我們已經探討過首三種型態,現在讓我們想一想第四種待人的型態,就是善待那惡待自己的人,有另一個關鍵字,叫:恩待.

誰可恩待有罪(性)的人?只有慈悲的父.今日經課最後一節(6章36節)這樣說:

36你們要慈悲,像你們的父慈悲一樣.

恩待是甚麼意思呢?若用善待和惡待來陳構恩待,恩待人就是善待那惡待你的人.就是第四種型態.

有甚麼例子?耶穌好激,他這樣舉例:

27「可是我告訴你們這些聽的人,要愛你們的仇敵!要善待恨你們的人!

28要祝福詛咒你們的人!要為淩辱你們的人禱告!

29有人打你的臉,連另一邊也由他打.有人拿你的外衣,連內衣也由他拿去.

30凡求你的,就給他；有人拿走你的東西,不要討回來.

咁高水平的動作,無專家指導,你地千祈唔好隨便學.

專家是誰?是牧師會吏?是導師老師?是父母家長?是uncle auntie?

都是,都不是.

指導我們的專家是聖父,是聖子,是聖靈.要慈悲好像聖父,要跟從耶穌效法他,要等待聖靈感動我們的生命.

讓有聖父,聖子,聖靈同行的人作我們同行者,少惡待,多善待,學恩待.

今日經課的開首6章27-30節看完,再借結尾6章35-36節作總結參讀:

35你們倒要愛仇敵,要善待他們,並要借給人不指望償還,你們的賞賜就很多了,你們必作至高者的兒子,因為他恩待那忘恩的和作惡的.

36你們要仁慈,像你們的父是仁慈的.」

所以,耶穌給與我們一個更高的倫理,是神的倫理,有三點結論:

你們倒要愛仇敵,要善待他們,Be kind to your enemies.

並要借給人不指望償還,你們的賞賜就很多了.賞賜包括:你們必作至高者的兒子,因為他恩待那忘恩的和作惡的.Be the children of the Most High.

你們要仁慈,像你們的父是仁慈的.Be merciful like our Heavenly Father.

今天講道至此,值得思考以下三個問題:

Q1:對於惡待,善待與恩待,有甚麼經驗之談?可以與你的好友分享一二.

Q2:對於人的倫理與神的倫理,看的分別在那裡?有何體會?

Q3:今日聽道後,有何提醒不再做?有何立志要去做?

4. 接待的賞賜:被接待者身分所載的福分得以恩臨在接待者身上

剛才我們從普世倫理到耶穌倫理得出待人之道,就是少惡待,多善待,學恩待.

現在我嘗試解提出「接待回應力」可用在教學上的神學哲學的觀念(Philosophico-Theological Pedagogical use of the idea of Receptive Responsibility).

接待除了可分善惡外,還可以分厚薄.

厚待方面,歌都有得唱.

陳奕迅Eason Chan在《我不好愛》係咁唱:

你 又美麗 又偉大 又這麼慷慨

我 被厚待 亦相當意外

給你寵愛 這戀愛 結果倒下來

大概只是 我不好愛

我教價值教育碩士課程一定教唱這首歌.

至於待薄,未必是惡待,但肯定不好受.

看看第40節:

40「人接待你們就是接待我；接待我就是接待那差我來的.

首先,接待一個人代表接待差這個人所代表的更高者.

英文係:Ability to respond through receiving x as x representative of G (Christ representing God)

香港話係:「你知唔知係邊個叫我黎?我係代表緊d咩?」

其次,任何人都有許多身分,有重有輕.

以被接待者所有的某一身分的重來接待這個人,就是重視了這個身分的重要性.

反之,不以他身分的重來接待,就是待薄了,那就輕視了這個身分所代表的分量.

這分量可以是神的言--他是神的代言人,也可以是神的義--他是神眼中的義人.

所以第41節說:

41因應先知之身分接待先知就如先知,因應義人之身分接待義人就如義人.

而我們作為門徒,可大可小.無論我多渺小,你都厚待我,在我饑渴時餵養我,你們得到的是神親自答謝你的賞賜.把最渺小的以最高格調的方式去款待的就是最高水平的接待回應力,英文有一個字好能表達:Hospitality.

所以第42節說:

42無論何人,因為門徒的名,只把一杯涼水給這小子裡的一個喝,我實在告訴你們,這人不能不得賞賜.」

我好喜歡Les Miserables 這個故事,電影中文名叫《孤星淚》或《悲慘世界》.

其中一幕是主角Jean Veajean坐了十九年冤獄後,出獄找不到工作,又被惡待.

他去到主教府時被主教收容,又食又住.他夜裡偷走銀器.但被警察捉到,拿回與主教對證.主教說:「好弟兄,你為何走得這麼快,又這麼客氣,我送你金器和銀器,你為何不拿金只拿銀?」

警察放了他,離開後,主教對Jean Veajean說:「我用這銀器買贖了你的靈魂,從此你的生命活出上主創造你的高貴.」

唔得,我好喜歡這一段,要唱一次給你們聽:

“But remember this, my brother

See in this some higher plan

You must use this precious silver

To become an honest man

By the witness of the martyrs

By the Passion and the Blood

God has raised you out of darkness

I have saved your soul for God! ”

當一個生命經歷過被恩待,會被激活出上主創造的高貴.

這個原理用在教學上是好有效的.只要你善待學學,厚待學生,學生所成的一定高過他自己的想像.

「對待回應力」有一種負面力量叫「差等對待」(differentiating treatment),這一面叫側鋒,會造成不小心的傷害.我舉一個例.

這裡有署理校牧和三位助理校牧,個個都有講道.為何講道結集只出版某一位榮譽助理校牧的講道,只出版他一本專集?

難道我的講道不夠好?其他講道都好好喎?

自我在崇基學生時代,即是上世紀八十年代,這禮拜堂的崇拜的講道叫大學講章.當年出版了數本大學講章集,群花共舞,有何不好?為何不把數年來我們四位的講道結集?

「你知唔知大細超其實好難頂!」

我做了爸爸二十多年,明白到為一個兒子做生日會之時,要同時準備一份禮物給另一兒子.

如果你出席好朋友子女的生日會,不單要給生日主角送禮物,還要給配角預備,就算禮物有厚有薄,總要善待當中最小的,最被忽略的.

大家看看這個可容一千人的禮拜堂,你有看到門徒中最小的嗎?其實他可能是先知,是義人,是代表上主的.

二十五年來,我見到每位學生,我都是想像他可以成為最大的可能.你怎知道上主在他生命不會創奇?見證上主的高貴?這是接待回應力的教育應用.

5. 四十年來「被對待」的經歷,以及未來的被恩待

今天我要告別了.我是不配得你們厚愛的.

我是1980年考入香港中文大學,開始在這個禮拜堂生活.愛上一個愛好文學的少女後的數年,在這裡進行婚姻祝福禮.

兩周前6月14日下午,我負責禮拜堂讀•書•人•會6月份聚會.活動介紹是這樣的:

「鄭漢文自1980年考入中大便受教於沈宣仁博士,由必修「大學修學指導」到選修「宗教哲學」,從為學到為教都深受Dr. Shen的影響.2018年為完成崇基學院神學院B.D.神道學學士的論文《沈宣仁博士在崇基學院所實踐的基督教教育神學》而重讀此書,從而演繹出「靈智兼修的品格群體」作為基督教教育神學的核心,同時回顧(1995年至今年快要離開中文大學與崇基學院)任教了四份一世紀的教學生涯.是次活動Roger會選讀《三十年來情與理》部份篇章段句來分享《四十年來恩(待)與(承)教》,借此作念師恩的惜別.」

在崇基,中大成長時領受老師恩待,路上得到不少同行人的善待.

我是工作了十三個合約,不被再續約而要離開.自1995年1月我重返母校任教,我在香港中文大學任教了四份一世紀.過去二十五年,正職都工作得不愉快.但兼任的工作,包括在崇基學院,崇基學院神學院以及崇基禮拜堂都十分愉快.全部都是冠以榮譽甚麼的.能在這講壇被恩待,不配講道,得到不少一同服侍者的善待.

其實我四十年來在修煉的就是耶穌教導馬太福音十章第16節:「我差你們去,如同羊進入狼群；所以你們要靈巧像蛇,馴良像鴿子.」同時兼修亞里士多德的實踐智慧(Aristotle’s practical wisdom / phronesis),在適當時候活現多元的應變能力.

讓我可以與大家分享我的未來,以下是鄭漢文Roger未來動向的公布:

「我係Roger,香港出生,1980年考入香港中文大學,從攻讀學士到完成博士共十四年多,然後1995年返回母校教育學院從事教學工作,工作至今已四份一世紀有多,到了今年七月便告別生活了四十年的中大.

十分蒙恩,得台灣的長榮大學禮聘為副教授,八月開始成為東南亞文化與產業學士學程的教員,同時參與大學通識教育的教學,在行政職務上兼任入學服務處執行秘書,一方面被委託照顧在台南的長榮大學就讀的香港生,另方面專責香港收生事務.

(若有年青人想在本年九月於台南的長榮大學開學,請於七月三日前申請,結果會於八月三日放榜.以後可以直接聯絡我.這是我最後一個植入式廣告.”rogercheng@cuhk.edu.hk“這個中大正規電郵地址在2020年7月31日後失效,不想與Roger失聯者請早日電郵至”cu2020hk.rogercheng@gmail.com” or “rogercheng@mail.cjcu.edu.tw”.)」

謝謝今日為我中譯英的Natalie,我還是要親口用英文講一句:

Today farewell sermon here is like my spiritual press conference, in which I announce my future, a plan God has prepared for me and my family, and a new mission which redeems all my past 40 years of wandering journey.

四十年走過曠野,一直依靠上主嗎哪來生活.

曾被惡待,不必記仇,常被善待,必要報恩.

我有學生在場聽我講道,臨別提醒,你們記住:

曾被惡待,切莫變惡人,若被善待,念恩成善人.

今天再被上主恩待,又被台灣的長榮大學恩待,八月開始,台南除口罩再見.

情意過濃,胡言亂語,今天講道到此,讓我們一起低頭祈禱:

天父啊,寬恕我們忘恩又作惡,讓我們活得好像你一樣仁慈,

學會恩待他人,就算他們忘恩或作惡.

耶穌啊,教我們善待他人,不單善待那善待我們的人,就算他們惡待我們,也能善待他們.

聖靈啊,求你更新我們的生命,讓我們在關懷社會,愛仇敵時成為「罪人得恩待」的見證.

奉主耶穌的名祈禱,阿們.

願恩惠常與你們同在!

但願榮耀都歸給聖父,聖子,聖靈,永在的三一神.阿們.



\newpage

\section{荒謬時代中信仰的實踐}
\label{sec:508}
\textbf{袁天佑牧師 - 荒謬時代中信仰的實踐}
\newline
\newline
\hyperref[sec:507]{< < < PREV SERMON < < <}
~
\hyperref[sec:toc]{[返目錄]}
~
\hyperref[sec:509]{> > > NEXT SERMON > > >}


馬太福音 11:16-19,25-30
\newline
\begin{longtable}{cl}
\hline
\hline
章節 & 經文 (和合本修訂版)\\
\hline
11:16 & \begin{tabularx}{0.7\textwidth}{X} 「我該用甚麼來比這世代呢？這正像孩童坐在街市上向同伴呼喊： \end{tabularx} \\ \\ \relax
11:17 & \begin{tabularx}{0.7\textwidth}{X} 『我們為你們吹笛，你們不跳舞；我們唱哀歌，你們不捶胸。』 \end{tabularx} \\ \\ \relax
11:18 & \begin{tabularx}{0.7\textwidth}{X} 約翰來了，既不吃也不喝，人們就說他是被鬼附的； \end{tabularx} \\ \\ \relax
11:19 & \begin{tabularx}{0.7\textwidth}{X} 人子來了，也吃也喝，他們又說這人貪食好酒，是稅吏和罪人的朋友。而智慧是由它的果子來證實的。」 \end{tabularx} \\ \\ \relax
11:20 & \begin{tabularx}{0.7\textwidth}{X} 那時，耶穌在一些城行了許多異能。因為城裡的人不肯悔改，他就責備那些城說： \end{tabularx} \\ \\ \relax
11:21 & \begin{tabularx}{0.7\textwidth}{X} 「哥拉汛哪，你有禍了！伯賽大啊，你有禍了！因為在你們中間所行的異能若行在推羅、西頓，他們早已披麻蒙灰悔改了。 \end{tabularx} \\ \\ \relax
11:22 & \begin{tabularx}{0.7\textwidth}{X} 但我告訴你們，在審判的日子，推羅和西頓所受的，比你們還容易受呢！ \end{tabularx} \\ \\ \relax
11:23 & \begin{tabularx}{0.7\textwidth}{X} 迦百農啊，你以為要被舉到天上嗎？你要被推下陰間！因為在你那裡所行的異能，若行在所多瑪，它還可以存留到今日。 \end{tabularx} \\ \\ \relax
11:24 & \begin{tabularx}{0.7\textwidth}{X} 但我告訴你們，在審判的日子，所多瑪地方所受的，比你們還容易受呢！」 \end{tabularx} \\ \\ \relax
11:25 & \begin{tabularx}{0.7\textwidth}{X} 那時，耶穌說：「父啊，天地的主，我感謝你！因為你把這些事向聰明智慧的人隱藏起來，而向嬰孩啟示出來。 \end{tabularx} \\ \\ \relax
11:26 & \begin{tabularx}{0.7\textwidth}{X} 父啊，是的，因為你的美意本是如此。 \end{tabularx} \\ \\ \relax
11:27 & \begin{tabularx}{0.7\textwidth}{X} 一切都是我父交給我的；除了父，沒有人知道子；除了子和子所願意啟示的人，沒有人知道父。 \end{tabularx} \\ \\ \relax
11:28 & \begin{tabularx}{0.7\textwidth}{X} 凡勞苦擔重擔的人都到我這裡來，我要使你們得安息。 \end{tabularx} \\ \\ \relax
11:29 & \begin{tabularx}{0.7\textwidth}{X} 我心裡柔和謙卑，你們當負我的軛，向我學習；這樣，你們的心靈就必得安息。 \end{tabularx} \\ \\ \relax
11:30 & \begin{tabularx}{0.7\textwidth}{X} 因為我的軛是容易的，我的擔子是輕省的。」 \end{tabularx} \\ \\
[1ex]
\hline
\hline
\end{longtable}
講題:荒謬時代中信仰的實踐 Practice of Faith in this ridiculous era

經文:馬太福音11:16-19,25-30

講員:袁天佑牧師

場合:香港中文大學崇基學院禮拜堂主日崇拜

日期:2020年7月5日

「沒有最荒謬,只有更荒謬」,是今天不少香港市民對香港的情況常說的一句話.不要說過去一年,只是過去幾個月,市民看到以下的事情,都會感到荒謬.

新冠肺炎,抗疫應該是要快和準.雖然香港抗疫也算做得不錯,但我們看見特區政府往往的回應,往往都給人留下有點兒荒謬的情形.封關,猶疑不決.限聚令,沒有科學根據,甚至是是自相矛盾,例如,可以容讓教會有聚會唱詩,但不讓市民在商場唱歌.最近限聚令的放寬至50人,但是茶樓飲食店就不限,教堂可座滿八成,但室外就限於50人,究竟病毒在室內傳播快抑或是室外快呢?口口聲聲說,沒有政治考慮,但又沒有科學理據,那麼限聚令是根據甚麼原則制定呢?

因著新冠肺炎,學生不能上課,影響最大的是那些參加文憑試的學生.學生已受到疫情影響學習,但竟然一條很簡單的歷史題目,因為「傷害一些人的感情」,沒有討論的空間,可以在學生努力作答後取消.最近教育局頻頻的致函學校,多次提出,如果學生對國歌又或是罷課等事,屢勸不改,可以報警.但教育是否用報警可以達成呢?

當然過去一年,不論你有甚麼政治立場,我們都會看到香港社會的撕裂.警民衝突,大家都認同要成立具公信力的獨立委員會去調查,但特首郤不願意,只是以監警會來找數,而監警會的報告竟然是由警隊一哥提出的.所謂「香港的真相」,被自己找來的外聘專家指這是「新真相」.「監警會」被人看成是「警監會」.

一個存在已30多年的嬉笑怒罵的節目,因為警察的投訴可以被終止.香港遲遲未為23條立法,大家都知道是擔心公民權利被剝奪,但特區政府可以配合中央,將在香港沒任何咨詢的法例強加在香港......這都是不少香港人感到不開心的事,令香港人會說,這是「更荒謬」的日子.

不過,「荒謬」並非這個時代獨有的.其實每個時代都有其更荒謬的事.在耶穌時代,耶穌就這樣描述他時代的荒謬:「我該用甚麼來比這世代呢?這正像孩童坐在街市上向同伴呼喊:『我們為你們吹笛,你們不跳舞；我們唱哀歌,你們不捶胸.』約翰來了,既不吃也不喝,人們就說他是被鬼附的；人子來了,也吃也喝,他們又說這人貪食好酒,是稅吏和罪人的朋友.」(太十一16-19)

所以,「沒有最荒謬,只有更荒謬」,是每個時代都發生的事,只是人怎樣在荒謬的日子中去生活,有信念的人怎樣去實踐自己的信念,面對荒謬而已.

耶穌在指出他時代的荒謬後,他向人發出悔改的呼籲(太十一20-24),並向人提出挑戰,到他那裡,向他學習(十一25-30).經課只選擇了後者,所以在此,我也集中與大家分享後者,我們要學習的功課.

在這短短幾節經文中,耶穌提出了兩方面向我們的挑戰:

一,作個謙卑的人

在第25-27節,耶穌這樣說:「父啊,天地的主,我感謝你!因為你把這些事向聰明智慧的人隱藏起來,而向嬰孩啟示出來.父啊,是的,因為你的美意本是如此.一切都是我父交給我的；除了父,沒有人知道子；除了子和子所願意啟示的人,沒有人知道父.」

耶穌說「這些事」,向某些人隱藏,向某人啟示出來.「這些事」是甚麼事呢?當然是指十一20-24所說有關審判的事(經課中沒有選讀的).但我們可以更廣闊的了解,「這些事」指上帝對世界的救贖計劃,甚至是上帝怎樣處理世界的事.事實上,有哪些事,不是上帝的事呢?

耶穌說:「上帝的事向聰明智慧的人隱藏起來,向嬰孩郤啟示出來.」如果樣樣事情都是上帝的事,但向聰明智慧的人隱藏起來,那麼聰明和有智慧的人都可以說是成為了世上無知的人；反之,沒有聰明才智,無知的嬰孩,甚至是無權無勢的人,可能才是有真智慧的人.

不過,誰才是有真智慧的人?那實在不易指出,但誰是聰明智慧郤是愚昧的人?那就隨時可以找到了!

容許我這樣說,我們的特首,就是有聰明智慧的人.她誇口自己「年年考第一」,表示自己是聰明有智慧的人.當有人提及現在社會的矛盾時,她又說:「我是讀社會學出身,我怎麼會不知道香港的深層次矛盾是甚麼?」不錯,她中學時年年都考第一,在香港大學修讀社會學,但是否真的了解香港的矛盾呢?我不敢說完全不了解,就算理解,她也沒有去處理.她所說的這幾句誇耀自己聰明的說話,反給人拿來嘲笑而已.

一個自認自己是聰明有智慧的人,聖經有這樣一句說話:「若有人自以為知道甚麼,他其實仍不知道他所應當知道的.」(林前八2)

為甚麼自以為聰明的人,反而是無知的人呢?明顯的是由於人的驕傲.特首自己曾這樣說:「官到無求膽自大.」這句說話本來是「人到無求品自高」,一個人到了「無求」的階段,就會好好的培植自己的人格.《論語.為政》中有這樣一句話:「子曰:『吾十有五而志於學,三十而立,四十而不惑,五十而知天命,六十而耳順,七十從心所欲,不逾矩.』」人年紀越大,慢慢會變得謙遜.但自以為聰明的人,就不會謙遜,特首以為自己懂得教育,用「雞籠」來描述教育；不查究說話的來源,便隨便去引用和誤指是孟德拉的名句,去支持自己DQ一條DSE的考題.......結果便引來大家的批評和嘲笑.驕傲帶來無知.

我剛才引用哥林多前書八章2節的經文,指自以為是知道的人,郤是無知的人.在哥林多前書八至十章,保羅與信徒討論一個「吃祭肉」的問題.時間關係,我不在此解釋這問題.吃與不吃祭肉,除了不是因為有知識的問題,便甚麼事都可作之外,保羅更重要要指出的,就是在這幾章經文最後的幾句話.他說:「凡事都使眾人喜歡,不求自己的益處,只求眾人的益處,使他們得救.」(十39)人以為自己有知識,除了令自己驕傲外,就是只顧自己,不放人在眼裡,做成自己更大的無知.

耶穌說:「凡自高的,必降為卑；自甘卑微的,必升為高.」(太二十三12)耶穌基督自己就是信徒的榜樣.在腓立比書二章1-11節,保羅就清楚指出耶穌怎樣從虛己而進入榮耀之中,他亦訓示信徒要向耶穌學習.謙卑不是有沒有知識的事,而是一種態度,能夠「不可自私自利,不可貪圖虛榮；只要心存謙卑,各人看別人比自己強.各人不要單顧自己的事,也要顧別人的事.」(二3-4)

耶穌說:「上帝的事向嬰孩啟示出來.」嬰孩,是一個目不識丁,他只懂倚賴他的父母.漸漸長大,知道父母的愛顧,會加上信心.他對知識的追求,也是一步一步的向前,而不是完全知悉.

我記得差不多50年前,當我讀大學時,有一位老師便這樣對我說:「當你讀大學一年級時,你會覺得有很多浩瀚的知識在你面前要去尋找.第二年時,你會覺得有很多東西你都不知道.讀完第三年畢業時,拿了個學位,Bachelor,你就會發現原來你是一無所知.Bachelor,就是寡佬的意思,沒有老婆也沒有屋,知識都沒有.」

信徒對上帝的事情的了解也是如此.特別是在今天,當我們看見很多荒謬崩壞的事時,我們都想知道,究竟上帝「搞邊科」?他會怎樣收拾這個「爛攤子」?

其實這也是我過去40年牧會中常問的問題.我常問:「跟著會有甚麼事發生,我要怎樣的去回應才對呢?」我都好像沒有答案似的,但只有一樣知道的,就是無論環境如何改變,我仍要努力的去事奉上帝,忠於我的工作,服侍信眾和教會的鄰舍.我亦發現一個很重要的道理,就是所學到更多的知識,不是在大學或讀神學時得到,而是在服侍中學習而得的.我不需要明白上帝完全的計劃,只需知道的就是在今天荒謬時代中我要作的是甚麼.

一位舊約的聖經學者范涵(Terence Fretheim)曾這樣說:「當人更多明白上帝,就更突顯上帝的奧秘......你掌握更多,就更知道自己無知.」自以為聰明的,反成愚拙；謙卑的,反成為智者.

二,放下勞苦重擔,負起輕省的軛

耶穌說:「凡勞苦擔重擔的人都到我這裡來,我要使你們得安息.我心裡柔和謙卑,你們當負我的軛,向我學習；這樣,你們的心靈就必得安息.因為我的軛是容易的,我的擔子是輕省的.」(十一28-30)

能放下,就算要拿起另一擔子,而且是輕省的,那是多麼好呢!

的確,今天不少人信耶穌,都是因為這寶貴的應許.信耶穌,生活不一定富裕,但也生活無憂；不一定富貴榮華,但也可有高官厚祿；能夠參與教會,有事奉,有奉獻,更有平安的生活.那實在容易和輕省.

但耶穌指他的軛是否真的容易,他的擔子真的輕省呢?耶穌自己甚至為他要完成的救贖計劃被釘在十字架.他又告訴我們:「若有人要跟從我,就當拾己,背起自己的十字架來跟從我.」(太十六24)

那麼,拿起耶穌的擔子實在不是輕省和容易!但為甚麼耶穌說,他的擔子是輕省和容易呢?

我看到有三個原因:

第一,是信念的問題:

耶穌描述他的時代,指約翰「既不吃也不喝」,但人說他是被鬼附；他自己「也吃也喝」,人又指他貪食好酒(十一18-19).所以可以說,你怎樣做,都會有人指你的不是.所以重要的,是你自己怎樣看自己所做的事,是否有意義,是否值得.假若覺得有意義,值得的話,別人的閒言閒語,並不為重.不過不少人總是生活在別人的期望下,造成不少壓力重擔.輕省容易的重擔,是因為你覺得值得和有意義的擔子.

第二,是價值觀的問題:

的確,吃喝或不吃喝,怎樣吃喝,吃喝得好與不好,都是人關心的問題.不少人就是活在不斷為糊口,甚至是超越糊口所需要的奔波勞碌,造成人生中的重擔.

我不是說,糊口並不重要,但除了糊口外,有沒有其他超越物質的更高的理想?匈牙利的詩人裴多菲(Petofi Sandor, 1823-1849)一句名言:「生命誠可貴,愛情價更高,若為自由故,兩者皆可拋.」有更高的理想,能對生命對愛情也可放下.我們人生的理想和價值是甚麼呢?

對基督徒而言,我們有甚麼理想和價值觀呢?耶穌說:「不要為那會壞的食物操勞,而要為那存到永生的食物操勞.」(約六27)

大家要特別留意,耶穌同樣用「操勞」來描述對那會壞的食物的追求和為那存到永生的食物的追求,表示兩者都要「操勞」而得.為會壞的食物,當然要操勞,但為甚麼追求永生,也要操勞呢?

當然假若信耶穌只是為生活平安富裕,有永生,天堂有位的話,那實在不必操勞,只要每星期返教會參加崇拜,或更好的每日都看下聖經便可以了.但真正的信仰,相信耶穌,是要背起十字架,要在今生生活中「行公義,好憐憫」(彌六8),這就不易了.特別在荒謬時代中,行公義,更不容易.從歷史和今天所發生的事,我們都看到,行公義,實踐真理的人,多面對打壓和逼迫.耶穌背起他的十字架,也指出跟從他的人,也要背起自己的十字架來跟從他.

但為甚麼耶穌說,他的擔子是輕省和容易的呢?就是因為這樣才有永恆的價值,才是真正永生之道.所以彼得說:「主啊,你有永生之道,我們還跟從誰呢?」(約六68)

第三,擔子輕重與人的驕傲自負或是謙遜有關:

耶穌說:「我心裡柔和謙卑,你們當負我的軛,向我學習；這樣,你們的心靈就必得安息.因為我的軛是容易的,我的擔子是輕省的.」(十29-30)

驕傲自負的人,是無知的人,往往做事愚蠢,造成自己的擔子,也將擔子加在別人身上.很坦白的說,我看見這幾年在香港所發生的事,實在與現今的管治團隊有關.聰明的高官,甚至是IQ160的,對於過去一年中所引起的問題,實在是責無旁貸.

謙遜的人,特別是我在上面所說的,不看自己比別人強,不自私,對別人有關愛的,這人不會自覺自己有能力去承擔自己沒有能力去承擔的事,我深信他不會為自己做勉強的事,也不會將難擔的手擔子加在別人身上.不過,他會因為別人的需要,去承擔一些難擔的擔子.但因為這抱負,他所負起的擔子會變得輕省和容易.

總結:

在結束時,我想引用不久前因疫情閒在家中追看的一齣韓劇《同伊》作為總結.這劇集是在10年前所播出的.不知大家有沒有看過,如果沒有,可以在網上找找看,但因為有60集之多,沒有時間就不要追看了.當然大家也可以想想另一劇集《大長今》,內容也差不多.

韓劇和一些韓國的電影,通常都會對充滿著荒謬的韓國政權作出諷刺和批判.《同伊》或是《大長今》便是這樣的劇集,當中當然是描述在古時朝鮮權貴間爭權奪利的事.

當我重看這些劇集時,我有3點的反思.第一,權力可以將真相壓下.為掩蓋真相,可以繼續以更大的錯去蓋著早前所犯下的小錯.第二,有權的人通常都會爭權奪利,看風使舵,當誰得勢時便倚靠誰.這兩點特別在有權力鬥爭的社會中常出現的事,亦是製造荒謬事情的原因.香港的政治情況是否也是如此呢?

劇集當然不是這麼悲觀.同伊和大長今這些主角,怎樣從賤民變成皇室成員或是非凡人物,這對出身卑微的人,可以是很大的鼓勵.卑微並不重要,存著謙卑的心,不驕不躁,是智慧的所在.另一點,我覺得更重要的,是她們對真相的執著,對公義的追求,對仁愛的實踐,那種鍥而不捨的精神,正是我覺得劇集中最重要學習的事.但要執著真相,堅守公義,彰顯仁愛,並不是一件容易的事,這是難擔的擔子.但那願意不為自己,捨己的精神,會讓這擔子變得輕省和容易.

智慧和承擔,是荒謬時代中最需要的.耶穌說:「凡勞苦擔重擔的人都到我這裡來,我要使你們得安息.我心裡柔和謙卑,你們當負我的軛,向我學習；這樣,你們的心靈就必得安息.因為我的軛是容易的,我的擔子是輕省的.」(十一28-30)這是耶穌向我們的挑戰,我們是否願意接受嗎?



\newpage

\section{道在心田}
\label{sec:509}
\textbf{劉國偉先生 - 道在心田}
\newline
\newline
\hyperref[sec:508]{< < < PREV SERMON < < <}
~
\hyperref[sec:toc]{[返目錄]}
~
\hyperref[sec:510]{> > > NEXT SERMON > > >}


講題:道在心田 The Word Sown on the Soil of Heart

講員:劉國偉先生

經文:馬太福音13

場合:香港中文大學崇基學院禮拜堂主日崇拜

日期:2020年7月12日

引言

弟兄姊妹平安.

我們今天選讀的舊約經課和詩篇,經文內容都是與豐收的景象相關.我對農業的認識非常膚淺,但是在過去多次聖地行的旅程中,有難得的機會去粗略認識以色列的農業科技發展,包括去參觀一些農業科技創新的初創公司.最深刻的一次,是參觀耶路撒冷的希伯來大學,有機會認識聞名世界的滴水灌溉科技,負責介紹的同事那引以自豪的眼神,給我留下非常深刻的印象.由上世紀60年代開始,以色列的水利工程師首次提出滴水灌溉的設想,並成功研製出有關的裝置, 60年之後的今天,以色列已經超過80\%的農耕土地使用滴水灌溉.

作為旅客或訪客,我每一次停留台拉維夫的時候,總會參觀位於當地市中心的迦密市場(Carmel Market).我特別喜歡在攤檔之間鑽動,身處色彩繽紛的各式各樣蔬菜,新鮮水果,乾果和香料之間,實在是一個有形有景的豐收影像,很多遊客都喜歡在擺滿農產品的攤檔面前拍照,那裡也是我購買手信的好去處.

其實以色列位於乾旱地帶,尤其是南部大部份地區都是曠野,水源十分缺乏.當你走過以色列的東西南北,看見相對乾旱的地貌,其實你心中很難構想出詩篇65篇所形容的豐收景象.自古以來,農業耕種都是以色列人在應許之地的長期挑戰.

今天選讀的福音經課,是大家都十分熟悉的「撒種的比喻」,三卷符類福音都有記載這個教導.耶穌採用農耕經驗來講論天國的道理,其實非常應景,也是十分貼地,對當時的社會來說,即使你不是農夫,你也會發覺耶穌所描述的內容不難理解,甚至產生共鳴.

馬太福音的五段講章

馬太福音包含了五大段耶穌的教導,分別是第5至7章的登山寶訓,第10章有關差遣門徒,第13章天國的比喻,第18章門徒的關係,還有第24至25章有關末世的儆醒.按今日之用語,這五篇講論可依次分類為:

(1) 門徒素質

(2) 門徒承擔使命

(3) 門徒的天國觀

(4) 門徒群體的建立

(5) 門徒在危難裡的考驗

撒種的比喻位於五篇講論的中間,其重要性不在於這章經文是研究耶穌比喻的基礎經文,而是這是耶穌從會堂作為教訓與傳道的基地,轉移到海邊的一個轉折位.我們不能說會堂的門已經向耶穌關上,但事實上,耶穌卻經歷猶太宗教領袖的批評和挑戰,然後再發展下去,就是尋找耶穌的把柄,最後決意要除掉祂.當權者意圖將耶穌驅逐離開他們管轄的範圍,認為這就可以終止耶穌所帶來的影響,卻沒有想到,耶穌在湖邊,路旁,山上,樹下,城門口,有若戶外的聖殿,向心靈有所渴慕的群眾,宣講天國的福音,和教訓天國的道理.到頭來,是好些控制猶太教的權貴,要把耶穌排除在會堂以外,甚至將耶穌置於死地,但天國的信息卻在會堂以外的世界,開花結果.

馬太在第11及12章記述了不同的人對耶穌的反應,承接上文,馬太以「當那一天,耶穌從房子出來,坐在海邊.」(太13:1)作為開始,正要交代撒種的比喻,跟第12章所記載,包括與法利賽人的爭論,其實有很大關聯.

馬太告訴我們,「有許多人到祂那裡聚集,祂只得上船坐下,眾人都站在岸上.」(太13:2) 耶穌坐下來,是拉比講論重要教導時的情態,讓我們來看,按照馬太記載,耶穌這第一篇海邊的教導.

四種心田

巴勒斯坦的農田形貌,以較狹長的長方形為多,兩塊田之間是行人走道,日子久了,走道被踏踩得十分堅硬,種子落在這裡,是難以伸入泥土的.

石頭上的淺土地,是巴勒斯坦常見的土質,在石灰岩的上面,薄薄的一片土層,這層泥土大概有三四吋深,即大概10厘米左右,種子落在這淺土層,由於石頭吸收較多熱力,種子發芽很快,可是種子的根往下生長時,卻碰上堅硬的石頭,不容易找到充足的養料和水份,日頭再曬,就因為缺水而枯萎.

撒在荊棘裡的種子,雖然可以發芽生長,卻因為稍事成長之後,給週邊的荊棘擠住,無法攝取足夠的陽光和養份,就不能茁壯成長.

好土有足夠的深度,種子可以深入土裡,吸收水份和養料,沒有荊棘的擠壓,成長就有更好的空間,並且結出三十倍,六十倍和一百倍.

耶穌就解說,這四種不同的田地,是代表什麼.

路旁硬地,就好像閉塞的心靈,無法接受上帝的道,即使聽見了,心裡也不會明白,撒旦就趁著這個機會,把心田上面的道奪去了.一個人的心田怎麼會像硬土一般,成為對天國的道理的閉塞的心靈?崇拜權力的人,喜歡擁有權力和賣弄權力,很難接受叫人捨己的天國度理；對福音信仰存有偏見,認為這世界沒有神,又或者耶穌只不過是一位教師,不可能是神的兒子......這些偏見可以使一個人的心眼盲掉；自傲自恃的人不願意做一個受教的人,又怎會明白一個叫我們將生命主權交給上帝的道理? 我有一位好朋友,是一位非常成功的商界精英,二十五年前已經是某國際超級品牌的第一位華人董事,他那時常說自己的人生哲學就是追求卓越,努力奮鬥,而且自己嚴守不騙人不害人的宗旨,若果上帝因為他不相信主而死後真的要下地獄,他只能控訴這個上帝是不公義的；十五年前,他與太太邀請我跟他二十歲的兒子傾談……

淺土地上種子發芽甚快,初聽道理時歡喜接受,但那心田沒有深度,當為了信仰而接受考驗,遭逢患難,或者受了逼迫,立刻就跌倒了.畢德生如此意譯這種生命狀況:so when the emotions wear off and some difficulty arrives, there is nothing to show for it. (The Message) 信仰若只是建基於一時的熱情,一剎那間的衝動,純粹感覺,則來得快,去得也快.弟兄姊妹,沒有根的信仰,很多時成為我們真誠走窄路的很大的障礙,看來是接受了神的道,卻並沒有天天背起十字架來跟隨主,沒有根的信仰,不單難以成長,更可能會枯萎.

耶穌說撒在荊棘裡的,就是人聽了道,後來有世上的思慮,錢財的迷惑,把道擠住了,不能結實.現代人的生活實在愈來愈擁擠,想方設法要更快獲取更多,也不甘心自己失去曾經擁有的,於是一場得與失的攻防戰,不斷在內心交織.是的,真正的競爭者,從來不在市場,不是你的同行,也不是你身邊跟你比拼誰更快上位的同事,而是那個要擠住你,阻礙你在信仰路上去多走一步的自己.現代中文譯本將「把道擠住了」譯作「窒息了信息的生機」,弟兄姊妹,為何有些人的心田是一塊窒息的土壤?為何人類常常用物質世界的得失追逐,去把心田裡最重要的道種擠住? 我們又是否生活在這種狀態中?

耶穌說撒在好地上的,就是人聽道明白了.整段經文所提及的道,道理,在原文都是同一個字Logos,按約翰福音的信息,道就是耶穌；而「明白」所指涉的意思,不僅是頭腦上的領會,其實更包含了相信和接受之意.所以,「人聽道明白了」,最深層的意思是一個人相信和接受主,並讓自己的心田,成為結出因道種而長出來的果實.弟兄姊妹,在一個艱難的世代,你願意耕耘自己的心田,為主結出生命的果實嗎? 三十倍,六十倍和一百倍.

土質會變好變壞

我們都知道農田的土質會隨著時間和環境轉變,可能好土變成壞土,壞土可能變成好土,請問你的那塊心田,現在是什麼光景?你需要上帝幫助你,尤其是感到窒息,欠缺水份,欠缺養份的時候,求上帝改變我們的土質,滋養我們的心田,又或者我們是一塊好土,就求上帝保守我們,不要讓荊棘生長在我們的心田中.

收成可期

若果你是一個為天國培育生命的人,我會認為你是一個撒種的人.對於撒種的人來說,有時見到道種在路旁,在淺土地上,在荊棘叢中,枯乾窒息,難免會失望.不過,耶穌這個比喻,要告訴我們,總有收成,而且可能是豐收一百倍,弟兄姊妹,即使仍然未見果效,也不必灰心,求上帝給我們智慧和信心,努力撒種,歡呼收割,也許會迎來生命的豐收.

讓我們再一次欣賞詩篇65篇的選段:

你眷顧地,降下透雨,

使地大得肥美.

上帝的河滿了水；

你這樣澆灌了地,

好為人預備五穀.

你澆透地的犁溝,潤平犁脊,

降甘霖,使地軟和；

其中發長的,蒙你賜福.

你以恩典為年歲的冠冕；

你的路徑都滴下脂油,

(將恩典堆在你的頭上,車輪走過都揤起脂油!)

滴在曠野的草場上.

小山以歡樂束腰；

草場以羊群為衣；

谷中也長滿了五穀.

這一切都歡呼歌唱.

弟兄姊妹,我投身學生福音運動及青年牧養工作已經超過四十年,經歷過流淚撒種的滋味,在事奉路途上被主內肢體傷害,還有撒旦的攻擊,但我不曾放棄,仍然耕耘心田,繼續撒種,期盼前面繼續有生命的豐收.

願神的道在我們的心田結實成長,更願我們都是一個為天國撒種的人,並經歷生命的豐收.

願榮耀歸予三一真神.



\newpage

\section{生命中能承受的重}
\label{sec:510}
\textbf{林豪恩先生 - 生命中能承受的重}
\newline
\newline
\hyperref[sec:509]{< < < PREV SERMON < < <}
~
\hyperref[sec:toc]{[返目錄]}
~
\hyperref[sec:511]{> > > NEXT SERMON > > >}


以賽亞書 44:6-8
\newline
\begin{longtable}{cl}
\hline
\hline
章節 & 經文 (和合本修訂版)\\
\hline
44:6 & \begin{tabularx}{0.7\textwidth}{X} 耶和華－以色列的君王，以色列的救贖主－萬軍之耶和華如此說：「我是首先的，也是末後的；除我以外再沒有神。 \end{tabularx} \\ \\ \relax
44:7 & \begin{tabularx}{0.7\textwidth}{X} 自從古時我設立了人，誰能像我宣告，指明，又為自己陳說呢？讓他指明未來的事和必成的事吧！ \end{tabularx} \\ \\ \relax
44:8 & \begin{tabularx}{0.7\textwidth}{X} 你們不要恐懼，也不要害怕。我豈不是從上古就告訴並指示你們了嗎？你們是我的見證人！除我以外，豈有神呢？誠然沒有磐石，就我所知，一個也沒有！」 \end{tabularx} \\ \\
[1ex]
\hline
\hline
\end{longtable}
講題:生命中能承受的重 The Bearable Heaviness of Being

講員:林豪恩先生

經文:以賽亞書44:6-8

場合:香港中文大學崇基學院禮拜堂主日崇拜

日期:2020年7月19日

引言

有網友留言說:「人生在世,每個人的心底,都有一塊沉重的石頭,這塊石頭抬不起,砍不爛,就呆立不動在那裡,跟隨你變老,死去.最沉重的是生命本身嗎?是繁複無盡的工作嗎?是壓得人喘不過氣來的樓宇按揭嗎?是一切都需要花費心力認真經營的生活嗎?」

你還記得,感到最沉重的是甚麼時候嗎?你還記得,險些壓垮你的是甚麼事情嗎?你還感覺到,那種沉重得喘不過氣來的感覺嗎?

今天經課的四段經文,讓我們想到,有些重擔是群體性的,你感到沉重,因為你是群體的一分子,分擔著群體的命運；有些重擔是際遇性的,人把重擔加在你身上時,可能你也把重擔加在他身上；有些重擔是來自個人內在的,你被自己不止息的慾望牽引,或者被強烈的失敗或罪咎感拖累而抬不起頭來；有些重擔是從另一個世界而來的,影響到你,但不為你所影響.

二．如何能夠承受群體性的沉重?(以賽亞書44:6-8)

44:6耶和華－以色列的君,以色列的救贖主－萬軍之耶和華如此說:我是首先的,我是末後的；除我以外再沒有真神.

44:7自從我設立古時的民,誰能像我宣告,並且指明,又為自己陳說呢?讓他將未來的事和必成的事說明.

44:8你們不要恐懼,也不要害怕.我豈不是從上古就說明指示你們嗎?並且你們是我的見證!除我以外,豈有真神嗎?誠然沒有磐石,我不知道一個!

以賽亞書跨越三個時代與場景,分別是第一至三十九章,主要以南國的烏西雅,約坦,亞哈斯與希西家王的時代作背景；第二部分由第四十至五十五章,以被虜回歸的年代作背景；第三部分由五十六至六十六章,以回歸後以斯拉及尼希米的年代作背景,前後跨越約 300 年.如果你是生存在以賽亞書時代的人,視乎你生存在這段歷史中的甚麼時期,你可能經歷過戰亂,或被逼流落他方,或從寄居之地回歸故士從頹垣敗瓦中重建家園.這些重擔,基本上與你個人的功過無關,也不是你憑個人的努力可以改變,只因為你是這段歷史中的以色列群體的一份子.

群體性的重擔可以被動地承受,也可以主動地承擔.古代知識份子往往有一種憂國憂民,以天下萬民之福祉為己任的志向.宋代的范仲淹的「岳陽樓記」展現的正是這種古代知識份子的人格風度:「居廟堂之高則憂其民；處江湖之遠則憂其君.是進亦憂,退亦憂.然則何時而樂耶?其必曰“先天下之憂而憂,後天下之樂而樂”乎?」

讓我們回到以賽亞書,44第1節耶和華對以色列說:

「我的僕人雅各,我所揀選的以色列啊,現在你當聽.

以色列的重擔,不是因為他被遺棄,反而是因為他被揀選.他被揀選來做甚麼呢?

44:8你們不要恐懼,也不要害怕.我豈不是從上古就說明指示你們嗎?並且你們是我的見證!

以色列被揀選,作為耶和華的見證.從舊約聖經的記述來看,以色列整個群體作為耶和華的見證,是他們的榮耀,同時也是他們的重擔.

今天,作為基督徒的我們有這種被揀選的意識嗎?若果有重擔,我們有想過承擔被揀選者的重擔嗎?彼得前書的作者對新約時期的信徒說:「你們是被揀選的一族,是君尊的祭司,是神聖的國度,是屬神的子民,要使你們宣揚那召你們出黑暗入奇妙光明者的美德.」(2:9)

舊約時期的以色列群體,如何承擔得起這重擔呢?

44:2那位造你,使你在母腹中成形,並要幫助你的耶和華如此說:我的僕人雅各,我所揀選的耶書崙哪,不要害怕!

面對重擔感到害怕是人之常情,上主以母親懷孕產子比喻祂與以色列的關係及感情,親切地安慰以色列「不要害怕」.

44:3因為我要把水澆灌乾渴的地方,使水湧流在乾旱之地.我要將我的靈澆灌你的後裔,使我的福臨到你的子孫.

44:4他們要在草叢中生長,如溪水旁的柳樹.

南國被滅之後,以色列深深感受到民族滅絕的威脅,「乾渴」,「乾旱」是比喻「後裔」,「子孫」減少及滅絕的危險.上主向他們保證,如水湧流在旱地,聖靈充滿他們的後裔和子孫,群體會重生和重建.

44:6耶和華－以色列的君,以色列的救贖主－萬軍之耶和華如此說:我是首先的,我是末後的；除我以外再沒有真神.

以色列如何能夠承擔起作為耶和華上主的見證這群體重擔呢?耶和華上主向以色列說明,祂是超越歷史的唯一,在歷史之先,也在歷史之後；同時,耶和華也是參與歷史之中,是以色列的君,也是以色列的救贖主.

三．如何能夠承受人生際遇中的沉重?(詩篇86:11-17)

86:14上帝啊,驕傲的人起來攻擊我,又有一黨強橫的人尋索我的命；他們沒有將你放在眼中.

常言道,多一個朋友總好過多一個敵人,然而,一生人總會遇上對你懷著敵意的人；做父母的祈求兒女一生遇上的都是好人,不幸的是,塵世中總有「衰人」,有時甚至是在意想不到的時候遇上.最近新聞說,「一名70歲婦人去年9月接獲騙徒來電提供低息貸款,其間一名60歲男子以「花言巧語」向她展開攻勢,二人隨後發展成情侶關係,騙徒多次到女事主的住所,以結婚及官司纏身為由向她索取金錢,女事主更變賣家產及向財務公司借款.該名涉案的60歲男騙徒,前後共騙取女事主590萬元後消聲匿跡.」

在人生的際遇中,可能會遇到有人把重擔加給你,令你感到無辜,感到憤怒,甚至令你陷入困境之中.

如果詩人是大衛,他的敵人可能是因為他的角色而來的,他的沉重是因為他的身分而起的,他的際遇,與他的角色和身分密切相關.你的情況類似嗎?如果不是這個角色,或者這個身分,可能你的際遇就會不同,你承受的沉重也會不同.

如何能夠承受這些重擔呢?最好當然是能夠避開,遠離「衰人」,大家可以參考詩人的建議:

86:11耶和華啊,求你將你的道指教我；我要照你的真理行；求你使我專心敬畏你的名!

86:12主－我的上帝啊,我要一心稱讚你；我要榮耀你的名,直到永遠.

如果你「行得正企得正」都遇到「衰人」,又如何可以承受呢?詩人指出,把重擔加給他的人其中一個特徵是:「他們沒有將你放在眼中」,而詩人的回應是:「求你使我專心敬畏你的名」.詩人的做法有甚麼值得參考呢?

試想想那位不幸遇上「衰人」的婆婆,她很可能經常想著這個「衰人」,想著他如何欺騙自己,或者想著如何向這個「衰人」報復,結果就是與那個「衰人」糾纏不休.如何能夠趕走「衰人」那揮之不去的陰影呢?詩人的建議是,第一,「求你使我專心敬畏你的名」,想著上主,取代想著那個「衰人」；第二,「照你的真理行」,無論「衰人」的行徑如何「衰」,詩人警惕自己不要讓「衰人」令自己也變成如他一樣「衰」,反而提醒自己要繼續按上主的真理行事.

86:15主啊,你是有憐憫有恩典的上帝,不輕易發怒,並有豐盛的慈愛和誠實.

86:16求你向我轉臉,憐恤我,將你的力量賜給僕人,救你婢女的兒子.

86:17求你向我顯出恩待我的憑據,叫恨我的人看見便羞愧,因為你－耶和華幫助我,安慰我.

詩人不採取「你做初一我做十五」的回應,或者也沒有能力「以其人之道,還治其人之身」,那麼可以怎樣呢?詩人仰賴上帝的憐憫和恩典,慈愛和誠實.詩人認為最重要的,是上帝對他的恩待.為甚麼呢?因為一方面這就是對他的肯定,另一方面也是折服對手的好方法.

三．如何能夠承受自己內心衝突的沉重?(羅馬書8:12-25)

有時候,我們的沉重與別人無關,而是因為自己內心交戰而來.

8:12弟兄們,這樣看來,我們並不是欠肉體的債去順從肉體活著.

8:13你們若順從肉體活著,必要死；若靠著聖靈治死身體的惡行,必要活著.

這種內心交戰而帶來的沉重,小孩子也能夠體會得到.有一位小學生如此描述他的內心交戰:

「我喜歡玩手機,一玩起來就沉迷.最渴望的是沒有人叫我停止,我想玩多久就玩多久.其實,我也知道時間一分一秒的過去,也感到總應該要停止的.但有一個邪惡的聲音不停跟我說:「唔緊要啦!」我費力的掙扎著,手機就像AA超能膠一樣貼在我的手上,無法脫離.又有一個正義的聲音敲著我的腦袋:「你這樣長時間玩手機會近視呀!你的功課做完了嗎?下週要考試了!」這聲音令我感到自疚,真的玩得太久了,該有點節制了吧!兩種聲音交替著,有時我才不管正義的聲音,玩到被爸爸罵才停止.其實很不開心的,一方面被爸爸罵,另一方面管不住自己的感覺很差.」(劉聿:放下手機)

小孩子都感受到內心交戰的沉重,成人的體會應該更多.有些是無傷大雅的,例如你知道自己不應該吃太多,但你的手和你的口卻不聽你的話；又例如你知道自己應該去做運動,你卻不能指揮你的腳站起來走出去.古代教會曾經列出人類的七宗罪,包括:傲慢,貪婪,色慾,嫉妒,暴食,憤怒及怠惰,這些都是人所共知的,然而也是經常痛苦掙扎的.如何承受得起這是掙扎的沉重呢?使徒保羅認為你要知道自己是誰,你不是那兩股力量的獵物,也不是慾望的載體,而是上帝的兒子,最終是自由的.

8:14因為凡被上帝的靈引導的,都是上帝的兒子.

8:21但受造之物仍然指望脫離敗壞的轄制,得享上帝兒女自由的榮耀.

保羅補充,我們仍然活在肉體之中的有生之年,靈慾交戰難以停止,因此,千萬不要以為自己已經昇華成為聖人.如果是這樣,如何能承受這靈慾交戰的沉重呢?保羅提出,能夠幫助我們的是盼望:

8:24我們得救是在乎盼望；只是所見的盼望不是盼望,誰還盼望他所見的呢?

8:25但我們若盼望那所不見的,就必忍耐等候.

仍然與罪惡掙扎的人不是輸了,只是未完全贏.向罪惡投降的,出賣靈魂的,就在慾望因此而得到滿足的那一刻,其實已經輸了.

四．如何能夠承受善惡混雜的沉重(馬太福音13:24-30, 36-43)

上文曾經提及過,有一種沉重,是你在人生的際遇中別人加給你的.因此,不少人都可能曾經發出如此的嘆息:「為甚麼這個世界有那麼多「衰人」呢?」或者盼望:「如果那些「衰人」在地球上消失,世界就美好了!」

13:24耶穌又設個比喻對他們說:「天國好像人撒好種在田裡,

13:25及至人睡覺的時候,有仇敵來,將稗子撒在麥子裡就走了.

13:26到長苗吐穗的時候,稗子也顯出來.

13:27田主的僕人來告訴他說:『主啊,你不是撒好種在田裡嗎?從哪裡來的稗子呢?』

13:28主人說:『這是仇敵做的.』

13:36當下,耶穌離開眾人,進了房子.他的門徒進前來,說:「請把田間稗子的比喻講給我們聽.」

13:37他回答說:「那撒好種的就是人子；

13:38田地就是世界；好種就是天國之子；稗子就是那惡者之子；

13:39撒稗子的仇敵就是魔鬼；

這個世界只有好人,沒有人就完美了.所以,當發現稗子的時候,那位僕人的反應是可以理解的:

『你要我們去薅出來嗎?』

曾經,有一些人認為另一些人是衰人,不應該存在於他們的社會之中,於是他們驅逐了那些人；曾經,有一些民族認為另一個民族是劣等民族,不應該存在於世界上,於是他們進行種族屠殺.從人類的歷史來看,大家都知道這是悲劇,但仍禁不住問類似的問題:「你要我們去薅出來嗎?」用中國的古語說,就是「替天行道」.不過,縱使人們如何努力清除異己,如何奮勇替天行道,那些被認為不應存在的卻仍然存在,被判為應該滅絕的仍舊生存.在比喻中,主人如此回答僕人的建議:

13:29『不必,恐怕薅稗子,連麥子也拔出來.

13:30容這兩樣一齊長,等著收割.

那麼,等到何時呢?

13:39收割的時候就是世界的末了；收割的人就是天使.

身邊環繞著你認為不應存在的人,社會存在著不是自己陣營的人,往往令人感到煩厭,感到困擾,甚至感到威脅,人們如何承受得起這沉重呢?

13:41人子要差遣使者,把一切叫人跌倒的和作惡的,從他國裡挑出來,

13:42丟在火爐裡；在那裡必要哀哭切齒了.

13:43那時,義人在他們父的國裡,要發出光來,像太陽一樣.有耳可聽的,就應當聽!」

為甚麼對於「替天行道」的做法要三思和謹慎呢?因為「天」自己會行道.有人說,真的假不了,假的真不了；到那日,「叫人跌倒的和作惡的,從他國裡挑出來,丟在火爐裡；在那裡必要哀哭切齒了」,另一方面,「義人在他們父的國裡,要發出光來,像太陽一樣」.如果你認為自己是「麥子」,你的周圍可能有「稗子」,除去「稗子」不是你的任務,而是天使的任務,你的任務是繼續做「麥子」,按著「麥子」的DNA生長.

五．總結

講道開始引述的那位網友問:「最沉重的是生命本身嗎?是繁複無盡的工作嗎?是壓得人喘不過氣來的樓宇按揭嗎?是一切都需要花費心力認真經營的生活嗎?」他回答說:「恐怕都不是.是絕望,是想要獲得自我救贖但卻無能為力時那種無邊無際的絕望.」古時耶和華上主對以色列說:「我是首先的,我是末後的；除我以外再沒有真神.」耶穌向門徒揭示終極的真實說:「人子要差遣使者,把一切叫人跌倒的和作惡的,從他國裡挑出來,丟在火爐裡……那時,義人在他們父的國裡,要發出光來,像太陽一樣.」使徒保羅勉勵信徒說「我們得救是在乎盼望；只是所見的盼望不是盼望,誰還盼望他所見的呢?但我們若盼望那所不見的,就必忍耐等候.」但願我們都像詩人一樣回應:「主啊,你是有憐憫有恩典的上帝,不輕易發怒,並有豐盛的慈愛和誠實.」



\newpage

\section{誰能使我們與基督的愛隔絕呢}
\label{sec:511}
\textbf{劉國偉先生 - 誰能使我們與基督的愛隔絕呢?}
\newline
\newline
\hyperref[sec:510]{< < < PREV SERMON < < <}
~
\hyperref[sec:toc]{[返目錄]}
~
\hyperref[sec:512]{> > > NEXT SERMON > > >}


講題:誰能使我們與基督的愛隔絕呢? Who Shall Separate Us from the Love of Christ?

講員:劉國偉先生

經文:馬太福音13

場合:香港中文大學崇基學院禮拜堂主日崇拜

日期:2020年7月26日

引言

弟兄姊妹平安.

今天是我以校牧的身份,最後一次在崇基禮拜堂講道,也許你以為這是刻意安排,其實不然.話說兩個多月前,收到本來今天負責講道的講員的指示,他未能在今天來我們當中分享,所以我就擔當了這份榮耀的替工,負責今天的講道職事,對我來說,能夠在離任前最後一次崇拜講道,確實是上帝給我的一份額外的恩情.

兩個月前,當我提交講道題目之後,我就想像今天講道的情景會是怎樣的.那時,我們仍然未曾恢復有現場座席,後來卻又恢復了,所以我以為今天的講道,是會有現場會眾參與其中的,但到頭來,我們又恢復實時轉播的安排,今天現場所見的,都是參與司職的親愛的弟兄姊妹,至於大家,就只得隔空相遇.

還記得7月5日的主日崇拜,我第一次向大家交代我即將離開崇基校牧室的消息,當天就有弟兄姊妹提出,與我一起拍照留念,我當然樂意並感到榮幸,那天我還笑說,尚有幾個星期,仍然有機會可以拍照,然而世事難料,今天我們就沒有辦法面對面,肩並肩地,在禮拜堂拍照留念了!

世事確實難料,但我們作為跟隨主的人,清楚知道無論世情變化如何,始終永遠不變的是上帝,上帝就是愛,上帝的愛從亙古到永遠,始終不變,這是我們在變幻不定的世界中所堅定相信的.

天國是我們所珍愛的

我們今天選讀的福音經課,總共提及了五個天國的比喻.

天國好像一粒芥菜種(太13:31),天國又像麵酵(v33),這兩個比喻都在告訴我們,天國運動開展的時候,規模可能很小,甚至乎並不顯眼,然而耶穌基督的國度,就有如芥菜種和麵酵一般,可以成長起來,帶來很大的影響.

天國又好像網撒在海裡(v47),到世界末了的時候,天使要從中把惡人分別出來,丟在火爐裡.所以天國並不單單著眼於數量的擴展,更注重生命和信仰的本質.

天國又好像寶貝藏在地裡(v44),又好像重價的珠子(v45),尋著的人就去變賣一切所有的,為要得著天國.

這些比喻告訴我們,追隨基督的信仰道路,從來都是萬分認真的事情.天國信仰可以有爆炸性的能力,影響既深且遠；因為真正走在天國路上的人,都是那些把天國看為無價珍寶的人,願意歡歡喜喜的變賣一切,為要得著天國,所以追隨基督的人當中,充滿了生命能量,也正因如此,仇敵要在麥子當中撒下稗子的種子(v25),也像那些混在海裡的惡人,攪一些破壞的工作,但最終都被丟在火裡.(v50)

弟兄姊妹,對天國的那份珍愛,就是認定天國是生命中的寶貝,天國又是重價的珠子,或者好像潘霍華(Dietrich Bonhoeffer)所說的那是重價的恩典,所以我們要珍之,重之,愛之,並且願意變賣一切所有的,為要得著基督.親愛的弟兄姊妹,我們作為跟隨主的人,是否都有這份情懷?

上帝浩大不變的愛

今天所選讀的舊約經課,是大家都非常熟悉的敘事片段.我們都讚賞所羅門向上帝求智慧,然而我們很多時都忽略了,加上和合本翻譯帶來的限制,我們很容易忽略,就是所羅門向上帝求智慧,其實乃基於他對上帝的愛的理解和經歷.

列王記的作者告訴我們,所羅門認定神以極大的慈愛對待自己的父親大衛,現代中文譯本如此翻譯列王記上3章6節,所羅門回答:「你一向對你的僕人我父親大衛大施慈愛,因為他對你忠心,誠實,有純真的心.你不斷地給他浩大不變的愛……」上帝因為所羅門求智慧就回答他說:「你既然求這事,不為自己求壽,求富,也不求滅絕你仇敵的性命,但求智慧可以聽訟,我就應允你所求的……」(v10-12a)

弟兄姊妹,一個愛神的人,所追求的並非長壽,也不必苦心求富貴,卻是祈求可以更好的事奉上帝和服侍人.他心裡面想著的,不會是設法除掉對自己的威脅,去除自己的眼中釘,甚至消滅自己的敵人.所羅門確信在上帝那裡,有浩大不變的愛,所以年輕的所羅門選擇了一條上帝所喜悅的道路.

上帝持之以恆的愛

今天選讀的詩篇,來自於特別長篇的詩篇119篇其中一個選段.這個選段表達出作者對上帝的教誨和律法的尊重.當中有一句關鍵的金句,也就是今天宣召的經文119篇132節:「求你轉向我,憐憫我,好像你素常待那些愛你名的人.」

原來我們對上帝話語的渴慕和尊重的那份動力,是源於愛.詩人告訴我們,他渴望遵守上帝的命令,原來在於上帝對我們的垂顧和憐憫.上帝轉向我們,憐憫我們,本來是他素常所作的(v132),也就是說一種持之以恆的愛,因為上帝持之以恆的愛,叫我們渴慕上帝的命令,並祈求上帝藉著祂的話語,使我們的腳步穩當,不為罪孽所轄制.

弟兄姊妹,但願我們渴慕上帝的愛,多於任何其他在這世上可以得著的情愛.

羅馬書的佈局

看過今天所選讀的福音書,舊約,詩篇三段經文,我挑戰大家我們有沒有那份因為珍愛天國,而願意變賣一切的情懷；在人生的追求中,我們又是否看得見其實是上帝浩大不變的愛相伴同行?我們又是否體會到上帝那份持之以恆的愛,成為我們渴慕上帝的原動力?

今天的講題「誰能使我們與基督的愛隔絕呢?」,源於今天的書信經課,我們讀出的經文是位於羅馬書整卷書卷的中間段落,在文本結構中,有著承先啟後的作用.我們先粗略地看看羅馬書的整體佈局.

羅馬書可分作以下六個部份.

引言/ 1:1-17

世人活在罪中需要上帝/ 1:18-3:20

因信稱義/ 3:20-8:39

神的揀選/ 9:11-11:36

信徒生活/ 12:1-15:13

結束與及問候/ 15:14-16:27

我們今天所誦讀的選段第8章26至39節,乃羅馬書中段完結前的一個小總結,正如我先前所說,從文學結構而言,帶有承先啟後的作用.當中又可以分做四個小段落,我們嘗試走進這段經文,讓我們對今天的講題,有更深刻的體會和領受.

這四個小片段分別是:

聖靈替我們禱告 (8:26-27)

萬事都互相效力( 8:28-30)

神若幫助我們(8:31-34)

誰能使我們與基督的愛隔絕呢?(8:35-39)

聖靈替我們禱告

「況且我們的軟弱有聖靈幫助,我們本不曉得當怎樣禱告,只是聖靈親自用說不出來的歎息替我們禱告.鑒察人心的,曉得勝靈的意思,因為聖靈照著神的旨意替聖徒祈求.」( 羅8:26-27)

這裡出現一個連接詞「況且」,就提醒我們這個段落的信息是與前文緊緊相扣的.前文談及「得救是在乎盼望」.原來在盼望,忍耐,等候的過程中,我們總會因為種種際遇,而經歷軟弱的時刻.這時,聖靈就來俯就我們,成為我們生命中最堅定的幫助,幫助原文帶有「手拉手」和「彼此相連」的含意,是聖靈陪伴著我們去面對那些軟弱的時刻,甚至乎我們可能會軟弱到一個地步,以致我們不曉得當怎樣禱告,當我們的禱告言辭缺乏,聖靈就親自用說不出來的歎息,替我們禱告.

弟兄姊妹,你有沒有想過,當我們為自己的軟弱嘆息的時候,聖靈和我們一同嘆息?聖靈把我們的疲乏,厭倦,自困和無力感,陳明在神面前,求上帝體恤和提升我們.弟兄姊妹,這是非常深層次的愛.

萬事都互相效力

「我們曉得萬事都互相效力,叫愛神的人得益處,就是按祂旨意被召的人.因為祂預先所知道的人,就預先定下效法祂兒子的模樣,使祂兒子在許多弟兄中作長子.預先所定下的人又召他們來；所召來的人又稱他們為義；所稱為義的人又叫他們得榮耀.」(羅 8:28-30)

保羅在羅馬書的引言中談及「本於信,以致於信.」(v1:17)就展開了有關信心和救恩的討論；當他闡述世人活在罪中,卻可以「因信稱義」,走在得救的路途上,保羅就以「愛」去總結「因信稱義」這偉大的救恩之道.他告訴我們,不但有聖靈為我們禱告,其實上帝在我們人生中所鋪排的一切事情,沒有一件是例外和偶然的,「萬事」所指涉的就是一切臨到我們身上的事,那是上帝為我們所預備的人生道路.至於萬事互相效力,是否能夠得著益處,關鍵乃在於是否「愛神」.接近的際遇,甚至乎是相同的經歷,不愛上帝的人可能毫無得益,愛神的人卻可以得到豐盛的收穫.我們所得著的最大得益處是什麼?就是可以效法耶穌基督的模樣.

「就預先定下效法祂兒子的模樣」(v8:29)這句子,按原文直譯就是「模塑成為祂兒子的形象」.走在這條得救的信心道路上,我們的生命日漸被更新變化,我們越是愛上帝,就越可以因生命的種種際遇,而得著益處,並且邁向與主越來越相似的方向.

保羅又告訴我們,我們回應上帝呼召的人,不單被上帝稱為義,所稱為義的人上帝又叫他們得榮耀.保羅在這這裡所採用的「稱義」和「得榮耀」的時態都是過去式,按上帝永恆的國度,稱義和得榮耀都已經是既成事實,就是說所有愛神的人,無論生命際遇如何,我們因為得稱為義,也就已經得著榮耀.我們也許在世人眼中際遇平凡,甚至坎坷,但上帝叫愛神的人定必得著益處,都必定得著從上帝而來的榮耀.

弟兄姊妹,當你觀看世情,又或者親身體會,這個顛倒黑白,弄權濫權的世代,你感到灰心嗎?又或者憤憤不平?你可會有一份信念,神定必幫助我們,藉萬事互相效力,叫我們得著益處?

神若幫助我們

「既是這樣,還有甚麼說的呢?神若幫助我們,誰能敵擋我們呢?神既不愛惜自己的兒子,為我們眾人捨了,豈不也把萬物和祂一同白白地賜給我們嗎?誰能控告神所揀選的人呢?有神稱他們為義了.誰能定他們的罪呢?有基督耶穌已經死了,而且從死裡復活,現今在神的右邊,也替我們祈求.」(羅8:31-34)

保羅信心滿滿的說:「神若幫助我們,誰能敵擋我們呢?」(v8:31)我們愛上帝,努力效法基督,並且一天比一天更似基督,因此我們滿有從上帝而來的榮耀,而上帝定要幫助我們如此.試問宇宙天地之間,有誰比我們所相信的神更強大,更有能力,以至阻礙我們繼續在信心的道路上前進呢!

另一方面,我們又何須抱有懷疑,生怕上帝不幫助我們走這條信心的道路,因為我們明明白白的知道,神愛我們到一個地步,甚至連祂自己的獨生子也不愛惜,為我們眾人捨了,祂也必定成就我們這條榮耀的道路.

保羅接連以反問「誰能控告神所揀選的人呢?」,「誰能定他們的罪呢?」追隨基督的人,有神稱我們為義,也有復活的基督,如今在上帝的右邊,替我們祈求,那麼,我們所走的這條信心道路,就可以成為一條得勝,並且得勝有餘的道路.

誰能使我們與基督的愛隔絕呢?

「誰能使我們與基督的愛隔絕呢?難道是患難嗎?是困苦嗎?是逼迫嗎?是飢餓嗎?是赤身露體嗎?是危險嗎?是刀劍嗎?如經上所記:『我們為你的緣故終日被殺；人看我們如將宰的羊.』然而,靠著愛我們的主,在這一切的事上已經得勝有餘了.因為我深信無論是死,是生,是天使,是掌權的,是有能的,是現在的事,是將來的事,是高處的,是低處的,是別的受造之物,都不能叫我們與神的愛隔絕；這愛是在我們的主基督耶穌裡的.」(羅8:35-39)

正如麥子與稗子的比喻,又或者撒網在海裡的比喻,仇敵魔鬼豈會袖手旁觀,對信主的人客客氣氣,不施予任何挑戰或攻擊麼?所以保羅提醒我們,我們會遇上患難,困苦,逼迫,飢餓,赤身露體,危險,刀劍.弟兄姊妹,當有一天我們因為心中所信,而需要面對患難,困苦和逼迫,我們不要感到稀奇；我們甚至可能會面對,有如飢餓那種極端貧窮的狀況,和面對赤身露體般的恥辱.若果面對各種各樣的危險,甚至帶著死亡威嚇的刀劍,那當然不好過,不過這一切能叫我們與基督的愛隔絕嗎?

保羅又說,靠著愛我們的主,在這一切的事上都已經得勝有餘了.「這一切的事」與前文所說的「萬事」前後呼應；保羅並且繼續引伸,表明他內心所深信的道理,就是無論是人間的生與死,又或者來自靈界的勢力,就是天使,掌權的和有能力的；也許時間帶給我們考驗,無論是現在的事,或者將來的事；還有空間帶來的挑戰,無論恩膏處或低處所帶來的困難；也包括我們所不認識和不理解的別的受造之物,都不能叫我們與神的愛隔絕.保羅更加在句末強調「這愛是在我們的主基督耶穌裡的」,是那位道成肉身,就是那位實實在在地愛我們到底的主.

結語

親愛的弟兄姊妹,我們千萬不要對上帝的愛抱有懷疑,也千萬不要因為生命中的際遇,遠離上帝.若果我們愛祂,我們深知道沒有誰,又或者任何人生的際遇,可以叫我們與基督的愛隔絕.兩年前,當我順服上帝的帶領,投身崇基校園的牧養工作,我知道那是一個巨大的挑戰,我原不知道要走的路有多長,我只是忠於上帝的呼召,一天接一天努力的走著,一轉眼就已屆退休之年,當我帶著不捨之情,要離開我所鍾愛的同學,禮拜堂的弟兄姊妹和同工,我清楚知道,前面仍有愛我的主與我同行,呼召我繼續走在事奉的路途上,為天國效命.

還記得兩年前有聘牧委員在面試中誠意問我,轉職到崇基校牧室,在薪金上可能有一個相當幅度的下調,我當時的回應是,我在職場生涯中已經有兩次,因為順服上帝的呼召,如此試過了.弟兄姊妹,我們欣賞所羅門求智慧,我們也可以為事奉上帝和服侍人去求智慧,不求取長壽和尊貴,對我來說,我所祈求的,是上帝保守我走在祂呼召我的道路中；我願意大家為主求智慧,但我更願意大家求一個加倍愛主的心.若果你問我誰能使我們與基督的愛隔絕,我認為最簡單直接的答案,可能就是「自己」.所以保羅嚴厲地說:「若有人不愛主,這人可詛可咒.」(林前16:22) 親愛的弟兄姊妹,在這個非常艱難的大時代裡,愛主是我們的必然選項,當我們越深層的體會上帝對我們的愛,我們也就會愛主更多,然後我們會對上帝的愛有更深的體會,於是我們就會加倍愛主.

親愛的弟兄姊妹,還有五天我便正式卸任校牧室主任的職務,讓我在這裡先行向各位弟兄姊妹告別.能夠在這裡服侍大家,得到眾多弟兄姊妹的愛護,鼓勵和欣賞,實在是上帝給我非常深厚的恩典.弟兄姊妹,但願在這個非常艱難的大時代裡,我們彼此相愛,也愛主更多.

願榮耀歸給三一真神,又願愛神的人得著生命的益處.



\newpage

\section{我們這裡只有五個餅兩條魚}
\label{sec:512}
\textbf{高國雄牧師 - 「我們這裡只有五個餅,兩條魚.」}
\newline
\newline
\hyperref[sec:511]{< < < PREV SERMON < < <}
~
\hyperref[sec:toc]{[返目錄]}
~
\hyperref[sec:513]{> > > NEXT SERMON > > >}


馬太福音 14:13-21
\newline
\begin{longtable}{cl}
\hline
\hline
章節 & 經文 (和合本修訂版)\\
\hline
14:13 & \begin{tabularx}{0.7\textwidth}{X} 耶穌聽到了，就從那裡上船，私下退到荒野的地方去。眾人聽到後，從各城來，步行跟隨他。 \end{tabularx} \\ \\ \relax
14:14 & \begin{tabularx}{0.7\textwidth}{X} 耶穌出來，見有一大群人，就憐憫他們，治好了他們的病人。 \end{tabularx} \\ \\ \relax
14:15 & \begin{tabularx}{0.7\textwidth}{X} 傍晚的時候，門徒進前來，說：「這地方偏僻，而且時候已經晚了，請叫眾人散去，他們好進村子，自己買些食物。」 \end{tabularx} \\ \\ \relax
14:16 & \begin{tabularx}{0.7\textwidth}{X} 耶穌對他們說：「不用他們去，你們給他們吃吧！」 \end{tabularx} \\ \\ \relax
14:17 & \begin{tabularx}{0.7\textwidth}{X} 門徒說：「我們這裡只有五個餅、兩條魚。」 \end{tabularx} \\ \\ \relax
14:18 & \begin{tabularx}{0.7\textwidth}{X} 耶穌說：「拿過來給我。」 \end{tabularx} \\ \\ \relax
14:19 & \begin{tabularx}{0.7\textwidth}{X} 於是他吩咐眾人坐在草地上，就拿著這五個餅和兩條魚，望著天祝福，擘開餅，遞給門徒，門徒又遞給眾人。 \end{tabularx} \\ \\ \relax
14:20 & \begin{tabularx}{0.7\textwidth}{X} 他們都吃，並且吃飽了。門徒把剩下的碎屑收拾起來，裝滿了十二個籃子。 \end{tabularx} \\ \\ \relax
14:21 & \begin{tabularx}{0.7\textwidth}{X} 吃的人中，男的約有五千，還不算婦女和孩子。 \end{tabularx} \\ \\
[1ex]
\hline
\hline
\end{longtable}
講題:「我們這裡只有五個餅,兩條魚.」 “We have nothing here but five loaves and two fish.”

經文:馬太福音14:13-21

講員:高國雄牧師

場合:香港中文大學崇基學院禮拜堂主日崇拜

日期:2020年8月2日

13耶穌聽見了,就上船從那裡獨自退到野地裡去.眾人聽見,就從各城裡步行跟隨他. 14耶穌出來,見有許多的人,就憐憫他們,治好了他們的病人. 15天將晚的時候,門徒進前來,說:「這是野地,時候已經過了,請叫眾人散開,他們好往村子裡去,自己買吃的.」 16耶穌說:「不用他們去,你們給他們吃吧!」 17門徒說:「我們這裡只有五個餅,兩條魚.」 18耶穌說:「拿過來給我.」 19於是吩咐眾人坐在草地上,就拿著這五個餅,兩條魚,望著天祝福,擘開餅,遞給門徒,門徒又遞給眾人. 20他們都吃,並且吃飽了,把剩下的零碎收拾起來,裝滿了十二個籃子. 21吃的人,除了婦女孩子,約有五千.

耶穌用五餅二魚給五千人吃飽,相信可能是其中一個在福音書中最多人認識的神蹟,場面震撼,見證人眾多,四部福音書也有記載(這極少有),讓這事件在教會傳統中也有著其獨特的地位,以色列有記念此神蹟的五餅二魚堂,世界各地有多個以五餅二魚命名的福音機構,更有數之不盡以五餅二角為主題的藝術作品,一個在兒童主日學必定講過的聖經故事,在團契中也必曾經查考過的經文,相信各位弟兄姊妹一定非常熟悉,今天希望透過此我們看似熟悉的,作更多的思考及學習…

經文一開始提到

「耶穌聽見了,就上船從那裡獨自退到野地裡去.」

到底是什麼事情讓耶穌要退到野地呢?當時耶穌的心情又如何呢?

馬太福音14:1-12節記載施洗約翰的門徒來告訴耶穌有關施洗約翰被斬首的死訊,施洗者約翰是耶穌的親戚,也是為耶穌洗禮的先鋒,深信當時耶穌的心一定傷透…

馬可福音6:31他就說:「你們來,同我暗暗地到曠野地方去歇一歇.」這是因為來往的人多,他們連吃飯也沒有工夫.

這或許也同樣是常態,因為每天的忙碌實在也太辛苦了…

辛苦的工作,難過的心情,受壓的環境,與你我今天的生活相似嗎?

我們每天最關心的是甚麼呢?自己,家人及友好的健康及工作?香港疫情的數字?國際對香港及中國的態度?…

耶穌眼中所看到的是什麼呢?

V14耶穌出來,見有許多的人,就憐憫他們,治好了他們的病人.

作為人子,耶穌同樣有身心靈的需要,與我們一樣,但憐憫在經文中出現,這憐憫亦馬上化成行動,接待,講論,治病…直到天晚.

今天的另一段經課,使徒保羅向羅馬的弟兄姊妹分享心事:

羅馬書9:1-3「我在基督裡說真話,並不謊言,有我良心被聖靈感動,給我作見證:我是大有憂愁,心裡時常傷痛； 3為我弟兄,我骨肉之親,就是自己被咒詛,與基督分離,我也願意.」

保羅愛弟兄的表達非常極致,為憐惜所愛的願意放棄一切,甚至與基督分離!…

我們又如何呢?耶穌的教導,保羅的榜樣:要愛就要到底,愛更要有所行動,要改變你的關注點,愛人就是要看到他人的需要比自己優先.

V15天將晚的時候,門徒進前來,說:「這是野地,時候已經過了,請叫眾人散開,他們好往村子裡去,自己買吃的.」

人的需要再次在經文中出現,人餓了就要食物,非常自然正常不過的事(耶穌敵擋魔鬼試探時也只是用「人活著不是單靠食物」,不是單靠,即是清楚表示我們也需要依靠食物),門徒的關心明顯地是出於對人的愛,關懷群體生活的所需,向耶穌所提出的建議也非常合理,理應值得稱讚,表揚,及按此而行,但耶穌卻有完全想像不到的回應:

V16「不用他們去,你們給他們吃吧!」

耶穌這要求是否太過份呢?實在太離地了,到底耶穌清楚現況嗎?

門徒(我們)如何回應呢?

V17 we have nothing here but five loaves and two fish (NRSV)

我們什麼也沒有,只有五個餅,兩條魚.

約翰福音6:7腓力回答說:「就是二十兩銀子的餅,叫他們各人吃一點也是不夠的.」

有學者指二十兩銀子大概等如一個人200天的工資(按今天香港情況計算,可能大概\$150,000),明顯地腓力的回應是用邏輯假設的思考模式而決定的,「就算假如有…也不足夠,沒有辦法,我們幫忙不了…」我們今天是否就是用這思考模式來回應上帝呢?

門徒真的什麼也沒有嗎?我們過去如何回應主的呼召及吩咐呢?

「主啊我在這裡,請差遣…他!因此他比我好,比我屬靈,比我優秀…」

我們總喜歡與人比較,甚至炫耀自己所有擁有的比別人多,但在信仰內卻往往故作「謙卑」,不斷只集中看到自己不足的地方,而不願意以自己所有的恩賜來回應.

耶穌吩咐門徒把所有的分出去,真正的擁有由分享開始.

我們只見到自己的限制,不足?還是能夠先見到他人的需要,而動了憐憫的心呢?

神蹟如何發生的呢?

這是一個非常震撼的神蹟,五個餅,二條魚吃飽了5000個成年男人(如若按比例計算,總人數可能有15000-20000人),剩下零碎裝滿十二個籃子.

學者努力研究到底神蹟如何能發生:

耶穌展示讓食物質變的超自然能力?

如果是這樣,為何會有零碎剩下?是計算失誤嗎?

再者,這神蹟的意義只在於說明耶穌擁有我們所沒有的超自然能力?我們在其中又有何學習的地方呢?

還有其他可能嗎?

有學者嘗試用另一個角度去理解這神蹟:

約翰福音6:8-9有一個門徒,就是西門‧彼得的兄弟安得烈,對耶穌說:「在這裡有一個孩童,帶著五個大麥餅,兩條魚,只是分給這許多人還算甚麼呢?」

萬多人中理應有如這孩童一般帶著一點食物…

難道是耶穌施展個人魅力感動眾人,使眾人願意分享?

這理解可能與我們傳統想像的有分別,首先我們要清楚以上二種對此神蹟如何發生的「猜想」也是合理,並且應該是有可能的,可惜同時我們到現為止仍然是無法知道當日事件是如何發生的(相信要到天家時再問問主了)；我們比較難相信後者這種猜想可能是因為覺得這理解會大大降低這神蹟的震撼,甚至對耶穌的能力有所質疑；其實完全相反,這不但沒有減低「神蹟」的能力感,更沒有削弱耶穌的「神性」,上帝就是那位創造萬物,賜予我們一切的主,改變「物質」,超越「自然」等對耶穌上帝來說當然是輕而易舉的事,但當上帝給予我們自由意志之後,改變人心才是最難的地方,需要人的降服及參與,這才真正是耶穌/信仰要展示的最大力量!

今天另外的經課:

詩篇145:15「萬民都舉目仰望你；你隨時給他們食物.」

以賽亞書55:2-3「你們為何花錢買那不足為食物的?用勞碌得來的買那不使人飽足的呢?你們要留意聽我的話就能吃那美物,得享肥甘,心中喜樂.你們當就近我來；側耳而聽,就必得活.」

所以從來我們就根本不用擔心生活的所需,上帝清楚知道,又會按時賜給我們,問題只在於我們一直只會注視自己的不足,沒有看到上帝要我們去看的；如果你真的深信「萬物都是從主而來」,你就能真正沒有難題地「把從主而來的獻與主」；與其不斷掙扎,努力尋找生活(搵食),你更應該努力尋找(搵)那萬有的主,賜生活真豐富的上帝.

從安得烈與腓力的表達中,清楚表示他們也同樣覺得不足夠,也可能同樣覺得耶穌的要求「無理」,「過份」,「離地」,但腓力只留在假設的地步,而安得烈就把真正謹有的,拿出來分享,可能當他分享的時候也覺沒有大用處,對主也是小信(甚至自覺沒有信心),但只要願意以行動擺上,當上帝使用,一切不再一樣.

無論這神蹟是如何發生,其實一點也不重要…

重點應該在問「神蹟是在甚麼情況下才會發生?」

這才真正與我們生活有關係,耶穌施行神蹟的目的不應只是要展現出祂何等的厲害(這一點我們作為追求信仰,認識主的人應該一早清楚知道並順服),而是讓我們也能經歷,參與,並在其中體會主的心意,更有機會在我們身上展現出同樣的效益…

約翰福音14:12耶穌說:「我實實在在地告訴你們,我所做的事,信我的人也要做,並且要做比這更大的事…」

你信這話嗎?

耶穌向門徒的應許:「在人所不能的事,在上帝凡事都能.」

今天我們自己,這城市,社會,世界需要神蹟嗎?又或者問需要什麼神蹟呢?

自我的成就?一切如己意?事惰由壞變好?解決所有「問題」?

耶穌在五餅二魚這神蹟中展現了上帝的大能,看到別人的需要而放低自己,在有限當中仍然願意獻上所有的…

上帝就會透過這獻上的行動而使身邊每一個人都飽足,並且有餘.

我們今天又如何在這世代為主作更大的事呢?

在限制中,又軟弱的我們如何成就這神蹟呢?

我們會是在不斷假設,思考,而結論就是停在「我無能為力」的一位?(腓力)

是那一位雖然小信,仍然努力找出所謹有的,帶著懷疑的心也盡力回應一下主的?(安得烈)

又或許是願意把自己有的奉獻出來,回應當下的需要?(孩童,群眾?)

還是會願意望著天祝福,把有的擘開,遞出去,又傳給眾人,最終大家都被祝福,而能一同經歷豐富呢?



\newpage

\section{只剩下我一個人}
\label{sec:513}
\textbf{張美珍博士 - 「只剩下我一個人」}
\newline
\newline
\hyperref[sec:512]{< < < PREV SERMON < < <}
~
\hyperref[sec:toc]{[返目錄]}
~
\hyperref[sec:514]{> > > NEXT SERMON > > >}


講題:「只剩下我一個人」 “I am the only one left”

講員:張美珍博士

場合:香港中文大學崇基學院禮拜堂主日崇拜

日期:2020年8月9日

各位弟兄姊妹平安.

沒想過在這樣的情況下再度踏上這個講台. 台下沒有聽眾,但我知有很多透過互聯網絡一起崇拜的會眾.大家都安好嗎?

我對上一次在這裡講道,是2015年7月. 當時的講題是:「上帝,你還在嗎?」之後便封了咪 (原本應承兩位前校牧每年暑假在這裡講一次道, 因為鄧瑞強博士暑期會放假),封咪乃因過去數年 面對人生太多困擾,覺得自己沒資格亦沒足夠的智慧站在這個講台上,與大家分享上帝的道.

今次多謝新校牧高牧師的邀請,再次站在這裡,一方面為表示對他的支持,另一方面是開始想通了,原來當我們對這個世界,對自己認識愈多,人生的困擾也會愈多,不會停止.只要謙卑自己,接受自己的有限及不足,我依然可以與會眾一起透過上帝的道,作出反省,尋求人生的智慧.

在預備今天的講道時,忽地覺得今日的經課熟口熟面,原來在六年前的今天,即2014年8月第2個星期,也是唸同一組經課(不知大家有沒有發覺我們採用的經課,是三年一個循環).當年我以同一段經文(列王記上19章)講了一篇道,題目是:「堅強與軟弱」,思考為何一向堅強的先知以利亞,可以突然變得很軟弱；這情況亦是會偶然出現在我們身上的.

今日雖然採用同一段經文,但再讀的時候卻有另一種體會.今日的講題是:「只剩下我一個人.」大家覺得這題目是否有點悲涼呢?

若我說「只剩下我一隻貓」,或「一隻狗」,大家會否有同樣的感覺?

我可以告訴大家,如果只剩下一隻貓,牠一定會好開心.原因是我家現有五隻貓,經常為了爭寵或霸地盤而打架,我方才發現,原來貓並非群居動物!

如果「只剩下我一隻貓」,應該並非壞事.我第一次養貓時只有一隻,那時牠非常悠然自得,尤其在加拿大時天氣比較凍,屋內有太陽的地方就會找到牠在進行日光浴,很懂得享受「貓」生.

但人類是群居動物,離開人群是無法獨自生活.在創世記講述上帝創造人類之時,上帝說那人獨居不好,於是創造了伴侶,便因此衍生了家庭及社群.

因此,當想到「只剩下我一個人」的時候,難免有點悲涼的感覺.今日所讀的舊約經課,即列王記上19:9-18,這句話出現兩次,都是從一位上帝所重用而又相當傳奇的人物--以利亞--的口中.

以利亞之所以傳奇,乃因他的出現,好像彷彿從天而降(王上17章首次出現他的名字,但背景沒人知曉.與另一位偉大的先知摩西不同,由出生到成長,於出埃及記都有詳細紀錄)；最後更神奇的,就是以利亞沒有經歷死亡,最終是被上帝以旋風接他升天(王下2章).

以利亞出場後的兩章經文(王上17及18章),記載了他如何藉上帝的大能,經歷及施行神蹟,令接濟他的窮寡婦的兒子從死裡復活,並靠著耶和華上帝的名,令天降大雨終止旱災,又令火從天而降.

然而,當他在山上與巴力的先知比試勝利後,以色列王亞哈的妻子耶洗別向他發出了追殺令.以利亞逃命,去了別是巴,在曠野走了一日的路程,在一棵羅騰樹下尋死.有天使供應他水及餅,再走了四十晝夜,來到何烈山(何烈山又名叫西乃山,這兩個名字常常互用；摩西頒十誡就在這個山上.)之後進了一個洞,住在洞中躲起來.

耶和華的話臨到他說:「以利亞啊,你在這裡做什麼?」 他說:「我為耶和華萬軍之神大發熱心.因為以色列人背棄了你的約,毀壞了你的壇,用刀殺了你的先知,只剩下我一個人,他們還要尋索我的命.」他這句話重覆了兩次,反映他是何等的孤單及恐懼,怕死怕得要命.

很悲涼,是嗎?一個為耶和華萬軍之神大發熱心,能夠透過禱告掌控大自然,可以持續三年零六個月不降雨；之後又能呼風喚雨,消除旱災,又能夠令火從天而降的神人,為何竟然因為一個追殺令落荒而逃,孤伶伶的一個人躲起來?

過去這大半年,面對新冠肺炎疫情及生死威脅之時,大家是否也有小小共鳴呢?限聚令雖然不是追殺令,但執法人員每日在街頭巡查,對違例者來說有如追殺令,因為隨時收告票；每日政府發放確診者但未知源頭的數目,以及在社交媒體互傳確診者所到過的大廈,餐廳,街市等等,也的確令我們精神緊張.

只不過是半年的光景,世界變得不一樣.到現在我們仍然無法弄清楚到底發生什麼事.我們從沒想過以往帶我們飛上天空,去體驗世界的航空公司,會瀕臨破產邊緣；以前我放年假都喜歡去郵輪,因為不想搬酒店,但以後是否有這膽量再踏足郵輪都是未知之數.

過往我們視作理所當然的日常生活,就連平日與親人的相聚或飯局,去酒樓飲茶,去街市買sung ,帶小朋友去公園,這些生活中最平常不過的事情,在今時今日竟然變得彌足珍貴,可望而不可即.

方言世界,中美關係的緊張局面,如箭在弦,彷彿世界大戰一觸即發.不論你多大年紀,相信我們都從未經歷過一個如此人心惶惶,不知所措的世界.

最要命的,就是以往若遇到任何困擾,可以見見談得來的朋友吹吹水解憂,或者找社工或牧師談談.現在我們被叮囑要躲起來,因為你所接觸的每一個人,包括你的至親,其至你自己,你都無法肯定是否帶菌者,結果要暫停一切活動,以免被病毒感染.

我想起三十多年前在崇基讀書的時候,有天我們幾位同學約了沈宣仁老師行山,之後他帶著我們去看了一套戲,片名叫The Thing, 中譯「怪形」.那時我是第一次看如此可怖的電影,至今記憶猶新,很多鏡頭至今仍歷歷在目.

故事發生在天寒地凍,人煙稀少的南極.有數個科學家連同一些支援人員,駐守在基地作研究.有一晚,在基地中,一隻死裡逃生的雪橇 (Husky) 突然變成怪物,攻擊基地中其他狗與科學家.原來有一來自外星的變種物體,在基地殺害生物後,能夠完美地變成被害者的容貌,並模仿他們的思考模式和記憶,混入他們當中.

這群科學家並工作人員開始奮力對抗外星生物,但他們也不知道是否已有隊友中招,身邊的隊友說不定會是外星生物假扮的!他們在驚恐之中互相猜疑,繼而變得不信任,彼此展開對抗.最痛苦的,就是自己都不知道自己是否已經給外星生物入侵,在懷疑隊友又懷疑自己之間,有些選擇強制隔離關起隊友,有些則自我隔離而躲起來.這個情景,是否與我們現在的境況非常相似?

當以利亞躲在山洞時,耶和華兩次問他說:「以利亞啊,你在這裡做什麼?」 他說:「我為耶和華萬軍之神大發熱心.因為以色列人背棄了你的約,毀壞了你的壇,用刀殺了你的先知,只剩下我一個人,他們還要尋索我的命.」

「只剩下我一個人」--假若內心充滿恐懼與鬱結,感到孤單無助,的確是頗為悲涼的；不但會令人情緒低落,嚴重者會患上抑鬱,甚至走上絕路.

英國政府在2017年首次設立了一個部門並委任了一位部長(Minister for Loneliness),負責解決社會的孤獨問題. “Nobody should feel alone or be left with no one to turn to. Loneliness is a serious issue that affects people of all ages and backgrounds and it is right that we tackle it head on. Our Strategy sets out a powerful vision for addressing this generational challenge.”

可惜她在兩年後以缺乏資源為由辭了職.否則我們可以借鏡如何解決孤獨問題.

過去數月,尤其過去這兩星期,抗疫及限聚令,有否令你感到孤單?政府官員在近日疫情記者會或抗疫的宣傳片中說,明白市民有抗疫疲勞,呼籲市民繼續忍耐,不要出街；但其實我們今日見市面仍有那麼多人在街上,並非抗疫疲勞那麼簡單,而是不同階層或年齡層的人,有不同的需要,不論是肉身的還是心靈的；並非一句不要出街便能解決.(早前因為食飯這個最簡單及最基本的肉身需要,搞出了一個大頭佛,相信決策制定者或專家顧問 上了寶貴一課.)

你的內心有否曾經產生疑問:我是一個基督徒,一直做好自己本分,平常又樂於助人,但為何辛苦經營的生意有可能倒閉?我工作勤奮盡責,表現出色,但一直賴以為生的工作怎麼有可能會失去?我和丈夫或太太夫妻恩愛,孩子聽話聽教,原本和諧的家庭為何竟然變得家嘈屋閉(因為孩子不能上學,父母在家工作,在狹小的空間裡爭奪資源?)?

我辛辛苦苦考入大學,作為大學畢業生,或是去外國留學生,本應開開心心展開人生的新一頁,但為何畢業至今仍未找到工作?而留學生亦被迫暫停學業?

我們每一個人的內心有否曾經如以利亞般,在心靈的洞穴中,充斥著憂慮,恐懼或者鬱結?

我上學年要教兩堂學院通識課,內容為「法理情的道德判斷」,與同學們探討作出道德抉擇的過程,並人類的思維模式.在備課時看了一本書,名為《快思慢想》,作者是康納曼(Daniel Kahneman),指出大腦原來有兩個系統,並分析他們如何運作.

系統一是自動化的模式,主要憑印象,直覺和感覺做成,毫不費力地運作；但系統二平常則處於很舒適的低費力模式中,只動用一小部分的資源在運作；但若系統一能得到系統二認可,  系統一的印象和直覺會轉化為系統二的信念,而情感的衝動便可化作持久的行動.

簡單來說,當以利亞退縮至一個洞穴,看到只剩下他自己一個人,感到憂心忡忡,這是事實,但這是系統一的運作.耶和華去尋找以利亞時,問他:「你在這裡做什麼?」,間接啟動了以利亞的系統二.當他向上帝連續兩次表達了內心的孤單與恐懼後,上帝叫他離開洞穴,吩咐他回去完成任務,並告訴他尚有七千人,是未曾向巴力屈膝的.

以利亞在洞穴裡,跌入了情緒陷阱,在孤單與惶恐當中,只看到只剩下自己一人.他忘了耶和華的存在,更忽略了這個世界原來還有很多人與他一起並肩作戰.

當自己感到孤單,憂傷,疲累之時,如何能夠從系統一的思維跳出來,轉化至系統二?有時單靠自己力量,未必能夠做得到.以利亞的故事告訴我們,耶和華上帝一直沒有離開以利亞,不過上帝展現的姿態,未必如我們所想像的.

「 耶和華說:你出來站在山上,在我面前.那時耶和華從那裡經過,在他面前有烈風大作,崩山碎石,耶和華卻不在風中；風後地震,耶和華卻不在其中； 地震後有火,耶和華也不在火中；火後有微小的聲音. 」(王上19:11-12)

上帝的出現,有時只是微小的聲音；祂可能透過家人,朋友的一句話,一首歌,一本書,甚或社交媒體的一個分享而出現.有時你對朋友的一句不經意的問候,可能上帝就在其中.

最後,與大家分享一首在過去數月於網絡流傳的詩We Fell Asleep作結.作者是一位歌手,出生於英國的巴基斯坦人.原文是英文版(見附件),對於我們面對全球抗疫的境況有很好的描述及提醒.我對網上的中文版作出了一些修改:

「一覺醒來,世界變了.

突然,迪士尼不再美妙,巴黎不再浪漫,紐約不再璀燦,萬里長城無法抵禦病毒,麥加已成為空城.

擁抱和親吻突然成為武器,不探望父母和朋友反而是愛的關懷.忽地之間,你體會了權力,美貌和金錢原來一文不值,它們在生死關頭都無法提供氧氣讓你呼吸生存.

當人類被關起來,地球持續運行,世界依然美好.她正在給我們傳遞一個訊息:人類並非必要.空氣,土壤,天空和流水,沒有你們依然美好.所以當你們走出籠子的時候,請記得你們只是地球的客人,並非主人.」

一首寫得非常精簡而優美的詩.

地球的主人當然不是我們,是上帝；我們雖然是過客,但別忘記我們亦是上帝所託付的管家.今日我們被迫暫時退入洞穴,日常生活被打斷,計劃被迫中止,但希望大家不要因而跌入情緒陷阱--以為「只剩下我一個人」,被病毒尋索我們的命.

地球持續運行,世界依然美好,上帝未曾離開過我們,而且還有很多任務等待我們去完成；最重要的,就是尚有千千萬萬的人與我們一起並肩作戰.當疫情受控及限聚令放寬之時,希望能在這禮拜堂見到大家,讓我們可以互相支持打氣.

在這艱難時刻,願望上帝賜恩加力,幫助我們一起作上帝忠心的管家,令這個世界因我們的存在而變得更美好! 阿們.

We Fell Asleep

By Haroon Rashid in April 2020

“We fell asleep in one world, and woke up in another.

Suddenly Disney is out of magic,

Paris is no longer romantic,

New York doesn’t stand up anymore,

The Chinese wall is no longer a fortress, and Mecca is empty.

Hugs  and  kisses suddenly become weapons, and not visiting parents  and  friends becomes an act of love.

Suddenly you realize that power, beauty  and  money are worthless, and can’t get you the oxygen you’re fighting for.

The world continues its life and it is beautiful. It only puts humans in cages. I think it’s sending us a message:

“You are not necessary. The air, earth, water and sky without you are fine.

When you come back, remember that you are my guests. Not my masters.”



\newpage

\section{教會這個群體}
\label{sec:514}
\textbf{許立中先生 - 教會這個群體}
\newline
\newline
\hyperref[sec:513]{< < < PREV SERMON < < <}
~
\hyperref[sec:toc]{[返目錄]}
~
\hyperref[sec:515]{> > > NEXT SERMON > > >}


馬可福音 3:13-19
\newline
\begin{longtable}{cl}
\hline
\hline
章節 & 經文 (和合本修訂版)\\
\hline
3:13 & \begin{tabularx}{0.7\textwidth}{X} 耶穌上了山，把自己所要的人召來，他們就來到他那裡。 \end{tabularx} \\ \\ \relax
3:14 & \begin{tabularx}{0.7\textwidth}{X} 於是他設立十二個人，又稱他們為使徒，要他們常和自己同在，也要差他們去傳道， \end{tabularx} \\ \\ \relax
3:15 & \begin{tabularx}{0.7\textwidth}{X} 並給他們權柄趕鬼。 \end{tabularx} \\ \\ \relax
3:16 & \begin{tabularx}{0.7\textwidth}{X} 他設立的十二個人有西門－耶穌又給他起名叫彼得， \end{tabularx} \\ \\ \relax
3:17 & \begin{tabularx}{0.7\textwidth}{X} 還有西庇太的兒子雅各和雅各的弟弟約翰—耶穌又給他們起名叫半尼其，就是雷的兒子— \end{tabularx} \\ \\ \relax
3:18 & \begin{tabularx}{0.7\textwidth}{X} 又有安得烈、腓力、巴多羅買、馬太、多馬、亞勒腓的兒子雅各、達太和激進黨的西門， \end{tabularx} \\ \\ \relax
3:19 & \begin{tabularx}{0.7\textwidth}{X} 還有出賣耶穌的加略人猶大。 \end{tabularx} \\ \\
[1ex]
\hline
\hline
\end{longtable}
講題:教會這個群體 This Community we called the Church

講員:許立中先生

經文:馬可福音3:13-19

場合:香港中文大學崇基學院禮拜堂主日崇拜

日期:2020年8月16日

「耶穌上了山,隨自己的意思叫人來；他們便來到他那裡.他就設立十二個人,要他們常和自己同在,也要差他們去傳道,並給他們權柄趕鬼.這十二個人有西門(耶穌又給他起名叫彼得),還有西庇太的兒子雅各和雅各的兄弟約翰(又給這兩個人起名叫半尼其,就是雷子的意思),又有安德烈,腓力,巴多羅買,馬太,多馬,亞勒腓的兒子雅各,和達太,並奮銳黨的西門,還有賣耶穌的加略人猶大.」(可3:13 - 19)

引發我思考今天這個題目的,是近日社會上的政治風波.當然,教會講壇或許並不是一個適合談論政治的地方,但作為每一個人生活作息的宏觀背景,政治卻確實是無所不在.事實上,從一個歷史的角度,教會跟政治確實大有關連.

初期的教會,受到當時羅馬政府迫害而四散,不少信徒被逼離開自己的家鄉,也同時將福音帶到他們所到之處.可以說,是政治的原因,促使了福音的傳播.後來羅馬皇帝信了福音,將基督信仰定為國教,從此教運興隆,出入將相,廢王立王,舉足輕重.這明顯也是政治.

到了中世紀,政教的勾結,基本上已經使教會腐化到一個難以忍受的地步.馬丁路德本來只是想提出內部改革,但陰差陽錯,當時教廷的固執,結合整個社會的政治氣候,結果是身不由己地從羅馬教會分裂出來,成為「抗羅宗」或更正教.而沒有教廷在信仰上的一錘定音,更正教就更是百花齊放,衍生出各自的流派.

而香港在過去百多年作為英國的殖民地,英國政府對教會的工作基本上是寬容和友善的.外國的差會,無論是天主教或更正教,亦紛紛透過參與扶貧,救災,辦學,建立地方教會,傳播福音,對整個社會也作出了積極的貢獻.而這一切,跟社會當時的政治氣候,都大有關係.

因此在個多月前,當中國政府單方面提出實施《港區國安法》的時候,本地教區一位樞機主教,立即出來安慰他的牧民,國安法的實施,對教會的運作毫無影響的時候,我感到有點意外.當年中國新政權的建立,對於當時的教會就產生了翻天覆地的影響.就算在近年中梵建交的過程,梵蒂岡在主教任命的問題上,也被逼作出了相當重大的妥協.以致有人覺得如果以前宗教是過分滲入政治,現在的問題大概是政治完全操控了宗教,而兩者都有違「政教分離」的原則.這樣,香港最近在政治上的重大改變,又怎可能對教會毫無影響?

這個情況,讓我進一步思想到一個更基本的問題:甚麼是教會?教會到底是一個怎樣的群體?

首先,當一個人決定成為基督徒,他不僅是接納基督作為他生命的主(那是比較個人和內心)；他更是具體地參與一個群體組織:譬如香港中文大學崇基禮拜堂,或者他們最初接觸福音的地方.事實上,他們亦是透過這麼一個群體,去認識福音,並因此做了一個那麼重要的個人決定.

我們必須了解,今天向我們發出邀請和挑戰的真理,都是透過教會這個信仰群體所傳遞,而並非「不假人手」從天而降.而基督教的信息,亦不是發明於今日；它是透過整整二十個世紀的歷史進程,經過種種困難和波折傳遞下來.因此,了解教會的本源,最原初的成員是一群怎樣的人,對於信徒今天該如何自處,就顯得相當重要了.

我的意思是:今天說到基督教,不同的人會有不同的聯想.事實上今日基督教會琳瑯滿目的宗派和傳統,有時甚至連信徒亦感到有點迷惘.你今天加入一間怎樣的堂會,在一個頗大的程度上就決定了你信仰的向度或「屬靈的氣質」.今天當你受洗加入教會的時候,你是受洗進到一個怎樣的群體?

有關教會到底是一個怎樣的群體,我們在教會聽到的,往往是一個教義性的,概念上的,所謂「合乎上帝心意」的教會,譬如說教會是基督的身體；教會是被建造在使徒和先知的根基上,有基督自己為房角石,是聖潔的,榮耀的,是那位充滿萬有者所充滿的.在剛過去的「港九培靈研經大會」中,就有牧者用了十堂的聚會,透過以弗所書,重申這方面的道理.但這一切又如何兌換到現世中,一般人所認識的教會?新約聖經,特別是保羅書信,論到教會的本質,對信徒又有何實質的意義?難道單單就是為了對照出現世教會的虧欠和不足?

論到教會的起源,一般人都是看《使徒行傳》.但今天這段經文,卻讓我們稍稍看見比初期教會更原始的信徒群體雛型(Prototype).

教會是一個怎樣的群體?

似乎在不同的階段,教會都有不同的表現.徒2:43-47:「信的人都在一處,凡物公用.並且賣了田產家業,照各人所需用的分給各人.」 我想馬克思在1867年寫《資本論》的時候,大概亦曾經從這段經文得到靈感.但教會在當時其實仍未是一個有組織的社團,群體.是到了《使徒行傳》第六章,當說希臘話的猶太人投訴希伯來人疏忽照顧他們的寡婦,教會才開始思考比較具體的架構組織,去應付日常運作上實際的需要.而門徒稱為基督徒,更是從安提阿開始(11:26).

反觀在開頭提到馬可福音三章的經文,描述耶穌最初設立的十二個門徒.讓我們不要忘記,這十二位耶穌親自揀選的門徒,就是日後教會最核心的成員.耶穌基督是教會唯一的根基,但「使徒權柄」,無論在教會的承傳或聖經正典的確立上,都是地上教會運作的具體準則.

○耶穌上了山…(v.13)

關於耶穌呼召門徒,其實在第一章已略有所述.在這裡,只不過是較完整地重申十二門徒的設立.

在這裡,我注意到耶穌並沒有「出去」呼召人,反而是離開群眾「上了山」.相對於我們熟悉的「所以你們要出去,使萬民作我的門徒」(太28:20,21),這「上去」似乎有很大的反差.

在主耶穌出來傳道之前,他曾被聖靈催促,獨自退到曠野禁食四十晝夜.(可1:12,13)很明顯,在他出去宣講之前,重要的,是他必須先確定自己有一個信息.也正正是這個緣故,在他設立門徒的時候,他並沒有急於出去找尋,而是先「上了山」.

我們實在不必擔心人會不會信.事實上現在的人,就是想盡辦法也要信點甚麼.可以說只要你夠膽講,就有人夠膽信:無論你要他信的是甚麼.(「雞都教」,「東方閃電教」…)

當我還住在沙田,有一次我去花園城附近食麵,經過天主教聖本篤堂.我看見兩個阿婆經過,其中一個忽然走上石級,合起雙手向門口的保羅銅像拜了幾下.我看了覺得非常有趣.中國人在信仰方面一向很「隨和」:「入屋叫人,入廟拜神」,無論是甚麼,總之「有拜錯,沒放過」.

香港人一向現實,但回歸二十三年,經歷零三年的沙士疫症,接二連三的人大釋法重新解釋《基本法》和「一國兩制」；一四年持續七十九日,遍地開花的社會運動,去年的反修例運動,到上個月《港區國安法》的實施,令致大部分香港人的發財夢,基本上已經破滅.他們是被迫面對一些更埋身,更基本的問題:這麼辛苦賺錢供樓,到底是為了甚麼?這麼辛苦做人,到底又是為了甚麼?

一位中大學生讀了陳健文教授的《獄中書簡》之後寫了一篇讀後感.他最後的結語是://說到底,我們都要問自己,在這個時代,在一己短暫的生命裡,到底願意為「甚麼」東西,付上怎樣的代價?怎樣才能順著自己的性,在亂世中不虛此生?//

我們該擔心的,並不是人信不信,而是我們有一個怎樣的信息,是否能夠對應人們心底的需要.

近年聽到不少教會思想到怎樣去接觸「草根」階層的問題.但這個講法,會不會暗示我們在心態上已慢慢「中產化」,意思是我們不自覺地在福音的信息上,加了一層信仰以外的意識形態,以致難於跟草根階層溝通?

其實接觸「草根」對我們來說應該不是問題.因為我們或我們的上一代,不少都曾經「草根」；我們當中不少人(包括我自己)都曾經住過徙置區.是受教育的機會,令我們攀上社會的階梯,逐漸脫離所謂「草根」的生活.

其實「草根」不「草根」並不涉及是非對錯.社會的小康化,脫離貧窮是一件好事.但儘管我們的生活可以脫離「草根」,但心態上卻不一定要「中產化」(主要表現於將大部分時間精力,放在改善物質生活和社會地位的追求).

不忘記自己的根源,確實可以提醒我們一個生命的真相:就是「人赤身出於母腹,亦必赤身歸回」,無論我們有多少學位,身家(或「有身家,送學位!」:香港大部分有錢人都是博士!).

「有知識」,「沒有知識」其實都不是最重要的問題.

在使徒中有打魚的彼得,亦有讀書人保羅.而保羅更是一人包辦了過半數新約聖經的書卷.有知識的人不需要先洗腦再加入教會,我們各人都是「照自己本相」來到上帝面前,無論我們有怎樣的學歷.

我的意思是,我們所面對的,其實並不關乎「知識分子」或「草根階層」之間的選擇,而是我們是不是有一個信息!

我有一個擔心,就是我們有時會以為為了遷就「草根階層」,而過分簡化我們的信息.「草根」階層可能因為各種原因沒有機會接受高深教育,但他們卻絕非「沒腦」.事實上沒有所謂「知識分子」的虛榮,他們很多時比讀過書的人更明白事理.

但我們給他們的又是甚麼呢?我們會不會有時以對待小朋友的態度,去對待「草根」人士?這樣做一方面是策略上的錯誤,另一方面亦是對他們的不尊重.

「簡單」的信息不一定是「幼稚」的信息.但會不會我們宣講的信息,有時確實是「too simple, sometimes naive」!?而對於那些已經參加教會一段日子的信徒,我們又有沒有幫助他們離開基督道理的開端? 抑或我們總是來來去去「三幅被」,令他們感到食之無味,棄之可惜?

好幾年前在大學教職員團契有位弟兄跟我談到香港教會增長的問題:他接觸到不少教會都有明顯人數的增長,但有相當可靠研究卻指出,過去幾年,恆常參加聚會的信徒皆維持在20萬之數.當然這兩年的社會不安加上疫症,這數目恐已難以維持.我們必須警惕:一個所謂「決過志」的人一旦離開教會,要他再回頭就更加困難.讓我們不要輕忽對待一個人的信仰決定.

在我們「出去」之前,我們必須先「上去」:直到我們覺悟到那於我們真實的真理；直到我們掌握我們真正相信的信念.直到我們所傳的是生命的道理.

○隨自己的意思叫人來,他們便來到他那裡(v.13b)

這裡說,主耶穌是「隨自己的意思」叫人來,不是人「隨自己的喜好」去跟耶穌.會不會有時我們因為要遷就人的軟弱,為了得到他們的「決志」,我們就得過且過地割價傾銷,然後為了一些虛浮的數字感謝主,讚美主?

當然,為了人能夠明白,以致得到福音的好處,信仰很多時是委婉的(保羅說他對著甚麼人就做甚麼人,為要得著那些人)；但信仰對人卻並不是沒有要求.當一個人說他願意信的時候,我們是不是就如釋重負,急不及待地「收貨」呢?

「他們便來…」耶穌到底向他們說了些甚麼!?主耶穌有甚麼魅力,以致他呼召,門徒就即時撇下一切去跟從他?是不是他有俊朗的面貌,令人說服的眼神,好像今天一些出色的佈道家那樣?賽53:「他無佳形美容,我們看見他的時候,也無美貌使我們羨慕他.他被藐視,被人厭棄,多受痛苦,常經憂患.」他只是趕鬼,醫病,到曠野禱告,傳講上帝國的道理；事實是,當我們看見一個人所做的,跟他所講的一致,他的說話就自然充滿能力.那大概不是個人的魅力,而是上帝的能力.

○他就設立十二個人

我們注意到,並不是所有跟隨耶穌的,都成了那核心的十二位!這十二位乃是將來要顛覆人間政權,改變人類歷史的核心小組.主耶穌說他要將他的「教會建立在這磐石上,陰間的權柄,不能勝過他.」(太16:18)主耶穌並不迷信數字,他甘願將短短在世工作的大部分時間,投資在這十二個人身上,因為他明白到數字是虛浮的,唯有生命的改變可以持久.那主耶穌要他們做些甚麼?

要他們常與自己同在(生命的實質薰陶)

在教育制度化之前,學徒訓練(apprenticeship)是任何人想學習一門功夫,手藝的唯一途徑.在學徒訓練的過程,徒弟與師父的關係是非常密切,往往是一起工作,一起生活,在一個比較全面的處境中,而不是作為一種獨立的技巧,去運用他的所學.因此徒弟最後學到的,往往亦不止手藝技巧,更是如何去實踐,運用那些手藝技巧的心法和態度.

傳統的學徒訓練,跟現代的職業訓練,存在著非常本質上的分別.因此,一些比較特殊的專業,譬如傳統戲曲,武術,技藝和醫學,都仍然掙扎保存著學徒的傳統.事實上,在非速度的運動項目,音樂,藝術的比試中,分別往往不在於做不做到 - 能夠做到只不過是基本 - 乃在於當中有沒有生命!譬如大家都努力準確地演繹蕭邦,是甚麼令冠軍脫穎而出?

同樣,要維持教會的經常聚會和活動亦不大困難,分別是這些聚會和活動有沒有生命,能力和聖靈的臨在.「時常與主同在」,就是嘗試在一切的事上,察驗上帝善良,純全可喜悅的旨意,這是教會一切工作的動力和基礎.

要差他們去傳道

「時常與主同」在的一個自然結果,就是眼睛會變得雪亮,心靈會變得通透,以致看透一些虛假的表象,洞悉事物背後的真相.正如主耶穌在世上傳道時,他的話語充滿能力,因為他每句說話都打中人心底隱藏的實況.我們所傳的福音,是能夠穿透生命實況的兩刃利劍,抑或只是隨時可以由褲袋掏出來的得救公式?

三,  並給他們權柄趕鬼

至於趕鬼,我曾經有過一些黑暗的經驗,但我不能確定那是不是真正涉及靈界.我相信鬼靈的存在,亦相信基督有勝過鬼靈的能力；我相信那住在我們裡面的,比萬有都大.但畢竟鬼靈在現代社會,已甚少以電影《驅魔人》那類面目出現.魔鬼撒旦仍然用各樣的詭計方法威嚇人,迷惑人,令人沉溺於各樣叫人難以自拔的泥沼中,向罪惡的現實低頭；教會亦仍然有責任,有權柄將人從撒旦的圈套中拉出來.這些都表達於教會在宣教以外的具體事工,包括向失落和絕望的人伸出援手.

這十二門徒名單,跟其他福音書的記載似乎稍有出入:

太10:1-4；路6:12-16；徒1:13.

西門:彼得,磯法(約1:42；太16:18,19)

雅各,約翰:兄弟,西庇太的兒子；半尼其/雷子(路9:54)

安德烈:彼得的兄弟(可1:16; 約1:40,41)

腓力:伯賽大人,與彼得,安德烈是同城；約1:45-51

巴多羅買:跟拿但業很可能是同一個人.約翰福音名單中沒提到巴,而符類/對觀福音亦未有提過拿(約1:45-51；21:2)

馬太(大概是馬太福音的作者)

多馬(性格多疑)

. 亞勒腓的兒子雅各(可15:40)

達太(太10:3):雅各的兒子猶大(路6:16；徒1；13)

奮銳黨的西門

加略人猶大

○幾點觀察,跟我們今天的教會生活有一定的適切性:

‧ 他們當中有不少是親屬,同鄉:

雅各與約翰是兄弟,是西庇太的兒子.〈3:17〉

安德烈是彼得的兄弟(約1:40,41)

腓力是彼得,安德烈的同鄉.(約1:45-51)

巴多羅買/拿但業是腓力介紹跟隨耶穌的.(約1:45-51)

可以說在這十二人都是藤瓜相依的.其實那是很自然的事;當你真心發現一些「好ye 」,你會去找你的老友抑或「三唔識七」的陌生人?但每當群體有爭拗或涉及利益衝突,就可能引起更大紛爭.

其中一例:太20:20-24「那時,西庇太兒子的母親,同他兩個兒子上前來,拜耶穌,求他一件事.耶穌說:你要甚麼呢?他說:願你叫我這兩個兒子在你國裡,一個坐在你右邊,一個坐在你左邊.耶穌回答說:你們不知道所求的是甚麼.…那十個門徒聽見,就惱怒他們兄弟二人.」今天一般教會所涉及的紛爭之中,不少是跟這類私人關係有關.一旦有爭拗,人的黨派本性就自然表露無遺,甚至拉埋自己的家族班底,另起爐灶.這是值得我們小心留意的.

○教會是能夠容納不同性格的群體.(Refer to the list)

我相信世界上並沒有所謂「理想的性格」或「基督徒的性格」.每種性格,都有它的強處和相對的弱點甚至盲點,很多時某種性格的強處,恰好就是它弱點的所在,就像一個硬幣的兩面.例如,比較理性的人的強項當然就是講道理,但他們卻同時容易引用堂皇的理由,去合理化自己其實不那麼理性的偏執.因此,「理性」很容易同時成為一個理性的人的盲點!

相反來說,「感性」很多時成為教會有意無意標榜的特質:因為教會是一個「愛的群體」.耶穌說:「你們若彼此相愛,眾人就因此認出你們是我的門徒了.」先不說愛在這裡其實並非一種溫馨浪漫的情緒感受,過份高舉感性,很容易動輒以一種主觀的「屬靈領受」或個人的「感動」去排斥常理,叫那些稍為「有腦」的人,望而卻步.

在這裡,重要的是要有自知之明.基督徒常常說「現在活著的不再是我,乃是基督在我裡面活著.」但當我們不真正認識自己,我們的「自我」,往往就在「為基督」的名義底下為所欲為.並且因為有了這個神聖的理由,我們就更加放不下這個「自我」.結果我們愈是「為基督」,教會就愈多意氣的爭拗.

○奮銳党的西門(徒1:13)

有學者認為西門可能只是一個政見上的「激進派」；但亦有學者推斷西門成為門徒之前,大概是某個比較激進的復國主義組織的成員.無論如何,在主耶穌呼召他時,是否知道他的政見或政治背景?沒有證據顯示他不知,但亦沒有記載顯示他要西門先政治洗底才可以去跟從他.

這對我們今天有甚麼啟示?

首先,政黨其實只是社會上不同階層或利益團體,就著社會當時具體的處境,提出一些集體的訴求.因此不同的社會階層,自然就有不同的政見和訴求；不同的政治傾向和關懷,是可以並且應該彼此共融的.特別在一個民主或開放的社會(Karl Popper),這樣的訴求是合理也是合法的,並且相信當不同的階層這樣提出,為政者就能更加能夠考慮以致平衡各方面的利益,造就一個更公平的社會.

而基督徒作為一個社會公民,當然有權利參與提出相關訴求,這跟信仰並沒有直接的矛盾.你可以按著自己在社會上具體的位置,參與代表你利益或理想的政治團體:如果你是打工仔,你可能會支持職工盟或民建聯；如果你沒有基本兩餐的憂慮,而有一些較高層次的政治理想,那麼你可能會支持公民黨或者自由黨；覺得他們代表不到你,你甚至可以自己組黨.

我們會因著自己所屬的社會階層,具體關懷,而有不同的政治取向,但不同的政見並不阻礙我們的分享和溝通,因為我們知道我們只是誠實地按著自己的本位,表達自己的關懷.我們各人按著自己心中的遠象,追求一個理想和公義的社會 ─ 倘若是我們能力所及,這甚至是責無旁貸；但我們實在無須在在將政見的分歧提升到信仰矛盾的層次.

美國曾經是不少人心目中的人間樂土.她有令人羨慕的民主選舉制度和社會保障系統.但我曾經在互聯網上看見一則報導:美國國家精神健康聯會(NMHA):當年大概有4\%的美國成人,亦即八百四十萬人曾經考慮自殺,其中六十五萬人企圖自殺,三萬人成功自殺.原因複雜:精神問題,酗酒,失業,寂寞.而種族歧視,一直是困擾著這個號稱民主國家的最大問題.

瑞士:民主和富裕的理想國度,但他們同時擁有最高的自殺率及最嚴重的青少年吸毒問題:最理想和民主的制度,不能告訴一個人他生存的意義和目的.只有意義的追尋,具體的盼望,能夠給人生存的勇氣及面對死亡的力量.因此,一切的社會,政治的改革雖然有其必須性,卻不能替代人對福音的需要.

○加略人猶大

Stigma；是教會一個揮之不去的烙印,信仰大家庭中的「黑羊」；被視為基督徒的恥辱.

但他卻令我想起那位拿香膏去膏耶穌的馬利亞.(太26:6-13):「普天之下,無論在甚麼地方傳這福音,也要述說這女人所行的,作個記念.」同樣,每當我們重溫教會的起源,讓我們不要忘記,在主耶穌所揀選的十二門徒中,亦曾經有過這麼一個因為信仰不合乎自已的想法與期望,而至終出賣了耶穌的猶大.

教會久不久就會傳出一些「醜聞」:某某執事長老撿財,某某牧師好色.…我們姑且不去討論個別事件的動機和真實性,因為就算是真的,其實對信仰亦絲毫無損.正如我們不會因為有人計錯數就否定數學,警察貪污瀆職就否定整個法律制度,我們又怎能因為有基督教領袖失足,而否定信仰的真確性?除非被證明是他的信仰令他誤入歧途,否則,個別人的道德操守,並不證明或否定他所宣認的信仰.

教會有時甚至會有出賣信仰,出賣朋友,甚至出賣自己良心的叛徒,但是教會的根基,從來不是建立在這些信仰叛徒身上,而是建立在耶穌基督穩固的磐石上.

教會有時亦會有出眾的領袖甚至是觸目的明星.他們的成功叫我們感到驕傲,他們的失敗令我們覺得難堪；但這些都只不過是一些榜樣或鑑戒,叫我們有所學習或者知所警惕.而我們大多數人,都不會生活在射燈之下.我們會過著平凡而真實的生活.但信仰的意義並不在於色彩繽紛,萬眾觸目,乃是在於有上帝同在,同行的人生.

最後要說的,是這十二個人並沒有在主耶穌呼召他們之後立時變得沉實穩重,英明神武,獨當一面.他們有不少不成熟的表現,無謂的爭拗,甚至在主耶穌背後爭做老大.這些,都要在日後和主耶穌相處,相交的日子裡面,受到雕琢和校正.

聖經說教會是主耶穌基督在地上的身體,所以當我們加入教會,與主相交就不單是一件關乎個人內心的事情,而是關乎整個信仰群體的事情,在當中我們得到教導和提醒,關心和鼓勵；受過經歷磨練的,就付出一分力,幫助哪些比我們年幼和軟弱的肢體.

在今天風雨飄搖的香港,願上帝幫助大家,在自己的崗位,成為你的屬靈群體的祝福.



\newpage

\section{給求義尋道者之言}
\label{sec:515}
\textbf{林豪恩先生 - 給求義尋道者之言}
\newline
\newline
\hyperref[sec:514]{< < < PREV SERMON < < <}
~
\hyperref[sec:toc]{[返目錄]}
~
\hyperref[sec:516]{> > > NEXT SERMON > > >}


以賽亞書 51:1-6
\newline
\begin{longtable}{cl}
\hline
\hline
章節 & 經文 (和合本修訂版)\\
\hline
51:1 & \begin{tabularx}{0.7\textwidth}{X} 追求公義、 尋求耶和華的人哪，當聽從我！你們要追想自己是從哪塊磐石鑿出，從哪個巖穴挖掘而來； \end{tabularx} \\ \\ \relax
51:2 & \begin{tabularx}{0.7\textwidth}{X} 要追想你們的祖宗亞伯拉罕和生你們的撒拉；因為我選召亞伯拉罕時，他只有一個人，但我賜福給他，使他增多。 \end{tabularx} \\ \\ \relax
51:3 & \begin{tabularx}{0.7\textwidth}{X} 耶和華已經安慰錫安，安慰了錫安一切的廢墟，使曠野如伊甸，使沙漠像耶和華的園子；其中必有歡喜、快樂、感謝，和歌唱的聲音。 \end{tabularx} \\ \\ \relax
51:4 & \begin{tabularx}{0.7\textwidth}{X} 我的民哪，要留心聽我，我的國啊，要向我側耳；因為訓誨必從我而出，我必使我的公理成為萬民之光。 \end{tabularx} \\ \\ \relax
51:5 & \begin{tabularx}{0.7\textwidth}{X} 我的公義臨近，我的救恩發出。我的膀臂要審判萬民，眾海島都要等候我，倚賴我的膀臂。 \end{tabularx} \\ \\ \relax
51:6 & \begin{tabularx}{0.7\textwidth}{X} 你們要向天舉目，觀看下面的地；天必像煙雲消散，地必如衣服漸漸破舊；其上的居民也要如此死亡。惟有我的救恩永遠長存，我的公義也不廢掉。 \end{tabularx} \\ \\
[1ex]
\hline
\hline
\end{longtable}
講題:給求義尋道者之言 Words to Those Pursue Righteousness and Seek the LORD

經文:以賽亞書51:1-6

講員:林豪恩先生

場合:香港中文大學崇基學院禮拜堂主日崇拜

日期:2020年8月23日

引言:呼喚

先知以賽亞轉達耶和華的呼喚:

51:1你們這追求公義,尋求耶和華的,當聽我言!

在這裡,耶和華向誰呼喚呢?誰是追求公義,尋求耶和華的人呢?會不會是上文陷於困境中提出投訴的猶大人呢?

49:14錫安說:「耶和華離棄了我；主忘記了我.」

49:14節記載他們尋找耶和華,但似乎找不到.51:1節記載耶和華的呼喚中卻把「尋求耶和華」與「追求公義」放在一起,表示兩者相關.這種相關性,變成一項兩個方向的問題:第一,「尋求耶和華」的人有「追求公義」嗎?第二,「追求公義」的人有「尋求耶和華」嗎?誰是「追求公義,尋求耶和華」的人呢?是他們嗎?你是其中一位嗎?耶和華有甚麼說話要對這些「追求公義,尋求耶和華」的人說呢?以賽亞書51:1-6節似乎建議讀者向三個方向觀望:第一是向後望,第二是向前望,第三是向上向下望.

一．向後望

你們要追想自己是從哪塊磐石鑿出,

從哪個巖穴挖掘而來；

2要追想你們的祖宗亞伯拉罕

和生你們的撒拉；

因為我選召亞伯拉罕時,他只有一個人,

但我賜福給他,

使他增多.

當時的猶大人為所失去的悲憤迷茫,也為群體能否延續下去憂心忡忡.前路茫茫,看不見將來.既然望不見前路,就不宜盲動,耶和華建議他們不如向後望.向後望,見到甚麼呢?向後望所見到的,能夠回應他們的焦慮嗎?沒錯,向後望正正是回應他們對於能否延續下去的擔心.「你是石頭爆出來的」這句罵人的說話我們最好不要對人說,然而,他們的確本來就不是甚麼,而是耶和華呼召出來的；他們本來就不是一個群體,而只是亞伯拉罕一個人.一個人變成一個群體,是耶和華「賜福給他」的結果.處身於困境之中,感到痛苦,擔心和焦慮是人之常情.向後望可能並沒有立即令困境消失,但卻得到啟發和釋放去面對當下.他們群體之中有一個人名叫約伯,說出了向後望的領悟:

1: 21(約伯)說:「我赤身出於母胎,也必赤身歸回；賞賜的是耶和華,收取的也是耶和華.耶和華的名是應當稱頌的.」

約伯的領悟,幫助了他處身於困境的當下而沒有失去對耶和華的仰望.我們本來就沒有甚麼,也不是甚麼,若然不是蒙賜福,我們甚麼也沒有,甚麼也不是.是誰令我們由「沒有甚麼」變成「有甚麼」呢?是誰令我們「不是甚麼」變成「是甚麼」呢?越是聚焦於失去的,會越感到痛苦；轉眼思念賜予者,視野會不一樣.

舊約聖經常常重提過去,重述歷史,由此看來,猶太人的信仰是一個講求群體記憶的信仰.當他們提起十誡的時候,不只是把它們視為抽空的規條,而是群體記憶的一部分:「我是耶和華──你的神,曾將你從埃及地為奴之家領出來.」(出20:2)讀聖經時,讀者會發現出埃及的歷史,成為經常被重述.歷史,把以色列和耶和華連繫在一起.

基督徒又怎樣呢?其實基督教也是講求群體記憶的信仰,從何見得呢?不用等到聖誕節和復活節,從我們每七日就誦讀一次的使徒信經就可以看出,其中關於聖子的部分是歷史的重述:

我信我主耶穌基督,上帝的獨生子；

因聖靈感孕,由童貞女馬利亞所生；

在本丟彼拉多手下受難,被釘於十字架,

受死,埋葬；降在陰間；第三天從死人中復活；

升天,坐在全能父上帝的右邊；

在困境之下,在迷亂之中,向後望,認識自己是誰,追索自己從何而來,有助安穩自己於當下,也有助思考和計劃如何走下一步.

我做婚前輔導的時候,第一個交給準備結婚伴侶的任務就是請他們寫下十項對對方的欣賞.你猜猜這是容易還是困難的任務呢?很容易,他們不需要幾分鐘,就列出十項對對方的欣賞了.對於準備結婚的伴侶來說,這是合理的,否則就未必會結婚.然而,對於結婚之後共同生活了多年的伴侶,再要衷心寫出欣賞對方的十項,恐怕就沒有那麼容易.當他們還是高度彼此欣賞的時候,我要求他們列寫出來,好好保存.日後萬一「天氣不似預期」,就拿出來重溫,現在不滿對方的地方,其實可能當天是欣賞的地方.關係改變,不一定是對方變差了,可能只是我的要求改變了,或者我看對方的角度改變了而已.曾經何時,我們是這樣互相欣賞的.這是事實,這是歷史,站在這歷史之上,我們應該如何向前走呢?

向前望

3耶和華已經安慰錫安,

安慰了錫安一切的廢墟,

使曠野如伊甸,

使沙漠像耶和華的園子；

其中必有歡喜,快樂,感謝,

和歌唱的聲音.

當時錫安仍是廢墟,人們從表面上看不見有甚麼起色,然而,不等於耶和華沒有行動.原來,耶和華已經默默地工作,就正如一粒被泥土覆蓋的種子,已經生根發芽,只是枝葉未顯露出來而已.將來是如何的呢?他們可以盼望甚麼呢?是伊甸園,那個充滿「歡喜,快樂,感謝,和歌唱的聲音」的耶和華的園子.既然耶和華能夠把曠野和沙漠變成伊甸園,怎麼不能夠把被破壞成廢墟的錫安更新重建呢?當下是曠野,將來可以是「伊甸」；當下是沙漠,將來可以是「耶和華的園子」；當下是廢墟,將來可以是更新的錫安.各位朋友,你感到當下如何呢?你滿意當下嗎?當下令你沮喪嗎?令你焦躁嗎?我們並非否定當下,而是盼望將來.我們的盼望,在於能夠把萬物更新的上帝,祂能夠把曠野變成「伊甸」,把沙漠變成「耶和華的園子」,把被摧毀的廢墟重建過來.

潘霍華「獄中書簡」其中一封信的其中一句令我感到非常驚訝,他的大意是說:「現在,我們要開始思考和平後如何重建教會與社會」.一位被囚禁在獄中的人,不知道二次大戰何時結束,但提出思考和平後的重建,他心裡的盼望實在令人讚嘆.艱難的當下不會是永遠的,破壞了的有被復修的可能,破損了的有更新的可能.

這個星期的九號風球「海高斯」不禁令人想起2018年9月來襲的那個超強颱風「山竹」,然而,「海高斯」與「山竹」的威力真是相差太遠了.還記得「山竹」對中文大學校園造成的破壞嗎?一名在中大工作30年的工友接受訪問時用「滿目瘡痍」四個字來形容「山竹」吹襲校園後的情景,樹技折斷,有大樹被連根拔起橫躺路上,道路不通,積水處處.美麗的崇基未圓湖被逼封閉了個多月,進行清理及維修.今日,若你置身未圓湖,你只會看見湖光山色,綠葉成蔭,而不會想像得到當日颱風吹襲後的「滿目瘡痍」.

我們相信復修和更新的可能性嗎?我們心中還有盼望嗎?我們的盼望有沒有指引著我們當下的生活呢?將來的盼望能夠指導當下的生活嗎?讀到一位婚姻輔導員的分享,讓我發現就是在婚姻之中,將來的盼望對當下的狀況也是有指引功能的.

在婚前準備課程中,她常常邀請準新人共同訂出婚姻的願景(Vision),使命(Mission)和核心價值(Value),並記錄在一張精美的咭紙(簡稱「VMV 咭」)上.有一次,有一對夫婦來尋求輔導,最令輔導員感到驚訝的是他們帶同 VMV 咭前來.原來他們激烈爭吵之後,太太即時收拾行李步往大門,打算離家並離婚；丈夫急忙在門鎖放置這張VMV咭,怒氣衝天的太太呆望這咭良久後淚如雨下,默默地帶同行李箱和咭回房.在輔導中,太太表示這張咭令她回想當初彼此的約定,並發現彼此沒有完全按照核心價值而行,大家都有不足,既然內容是大家共同訂立,今天就應共同實踐,一齊追尋共同繪畫的理想. 就是這樣,他們在輔導中重建琴瑟之好.(楊麗儀:「婚姻的未來」)

基督徒向前望的時候望到甚麼呢?我們能夠盼望甚麼呢?啟示錄如此描繪終末的景象:

21:1我又看見一個新天新地,因為先前的天和先前的地已經過去了,海也不再有了.2我又看見聖城,新耶路撒冷由上帝那裡,從天而降,預備好了,就如新娘打扮整齊,等候丈夫.3我聽見有大聲音從寶座出來,說:

「看哪,上帝的帳幕在人間!

他要和他們同住,

他們要作他的子民.

上帝要親自與他們同在.

4上帝要擦去他們一切的眼淚；

不再有死亡,

也不再有悲哀,哭號,痛苦,

因為先前的事都過去了.」

5那位坐在寶座上的說:「看哪,我把一切都更新了!」

向前望所看見的,與我們的現在有甚麼關係呢?會影響我們現在如何生活?做甚麼的人嗎?讓我們繼續讀下去:

6他又對我說:「成了!我是阿拉法,我是俄梅戛；我是開始,我是終結.我要把生命的泉水白白賜給那口渴的人喝.7得勝的要承受這些為業；我要作他的上帝,他要作我的兒子.8至於膽怯的,不信的,可憎的,殺人的,淫亂的,行邪術的,拜偶像的和一切說謊話的人,他們將在燒著硫磺的火湖裡有份；這是第二次的死.」

向上向下望

向後望,向前望是時間性的,向上向下望是處境性的.向後望和向前望提供了理解和視野,然而,人們當下必然存在於某種處境中.在這種狀態中,人們如何自處呢?

6你們要向天舉目,

觀看下面的地；

天必像煙雲消散,

地必如衣服漸漸破舊；

其上的居民也要如此死亡.

惟有我的救恩永遠長存,

我的公義也不廢掉.

在婚禮上,新郎和新娘誓約:「從今天開始,無論是順境或是逆境,富有或貧窮,健康或疾病,我將永遠愛你,珍惜你直到地老天荒.」在人們的認知中,地是不會老的,天是不會荒的,因此天和地是最恆久的,最可靠的.天地是否最恆久的呢?是否最可靠的呢?耶和華呼喚說:「6你們要向天舉目,觀看下面的地；天必像煙雲消散,地必如衣服漸漸破舊；其上的居民也要如此死亡.」如果「天必像煙雲消散,地必如衣服漸漸破舊」,其上的邪惡會永遠長存嗎?不會,耶和華繼續說,永遠長存只有「我的救恩」和「我的公義」.

既然如此,當下,我們應當如何生活呢?

4我的民哪,要留心聽我,

我的國啊,要向我側耳；

因為訓誨必從我而出,

「追求公義,尋求耶和華」的人,第一項的回應是:「要留心聽我」,「要向我側耳」.聽甚麼呢?「因為訓誨必從我而出」.當下,眾聲喧嘩,你能分辨得出耶和華的訓誨嗎?兩星期前我們舉辦了網上夏令營,學生們設計了一項估歌仔遊戲,把兩首流行曲和一首詩歌一齊播放,讓參加者分辨.人們能否聽得出那首詩歌呢?這就取決於人們對這首詩歌的熟識程度.在眾說紛紜,各方或明或暗爭奪話語權的社會,我們能否分辨上主的聲音呢?也是取決於我們對上主聲音的熟識程度.

我必使我的公理成為萬民之光.

5我的公義臨近,

我的救恩發出.

我的膀臂要審判萬民,

眾海島都要等候我,倚賴我的膀臂.

「追求公義,尋求耶和華」的人,第二項的回應是:「眾海島都要等候我,倚賴我的膀臂」.等候甚麼呢?等候「公理成為萬民之光」,「公義臨近」,「救恩發出」.倚賴耶和華的膀臂做甚麼呢?「我的膀臂要審判萬民」.

總結

講道開始時問過:誰是「追求公義,尋求耶和華」的人呢?你是其中一位嗎?我們都未必能判定別人是或不是,甚至也未必能肯定自己是或不是.不過,耶和華對當時的猶大人的呼喚,值得我們參考.第一,知道自己本來無一物,若有,都是恩典；知道自己從上主而來,最終歸回上主.第二,相信上主終會撥亂反正,將萬物復原更新,以這遠象來引導我們的追求.第三,當下,聆聽上主訓誨,努力實踐於人間.



\newpage

\section{最大的權柄復和}
\label{sec:516}
\textbf{高國雄牧師 - 最大的權柄 - 復和}
\newline
\newline
\hyperref[sec:515]{< < < PREV SERMON < < <}
~
\hyperref[sec:toc]{[返目錄]}
~
\hyperref[sec:517]{> > > NEXT SERMON > > >}


講題:最大的權柄 - 復和 The Power of Reconciliation in Christ

講員:高國雄牧師

場合:香港中文大學崇基學院禮拜堂主日崇拜

日期:2020年9月6日

可能你與我一樣,發現身邊不少人都會表示覺得基督教不錯,信耶穌很有意思,也會間中參與教會崇拜團契等聚會,但一直仍然未正或表示願意成為基督徒,我常常會問這些慕道朋友當中的原因,通常回應是覺得自己未夠好,又或未預備好等等…說穿了,不少其實是因為認為若然成為「真正的信徒」,就要放棄生活中的種種「享受」,一些自己喜愛的「生活習慣」,因為覺得作為基督徒,就會有一些事情「不可以做」,例如:食煙,飲酒,講粗口等；另外亦覺得要真正投入信仰,就一定要「做」某一些事情,例如:寬恕人,愛仇敵,放低仇恨(可能就正正是我們最不想放低的人或事),總感覺認真就很傻,很蝕底(吃虧)…從這角度就會更覺得基督教信仰很離地,其價值觀與我們身處的世界實在相差得太遠,很難完全接受並遵從；曾經認識有牧者,教會為了要吸引更多人加入教會,為能更「貼近」世界而放低「要求」,這是何等可悲的事啊!我們真的要需要認真地反省.

假若有人得罪了你,你會如何回應呢?

「你敬我一尺,我敬你一丈」本來是很正面的互相尊重的教導,在我成長的經驗當中卻發現落在我們身上卻成為了「有仇不報非君子」的負面演譯,舊約聖經中「以牙還牙,以眼還眼」的教導,這「對等復仇」的原則,原為了止息仇恨,卻變成了「有仇必報」的口號及價值觀.聖經特別在詩篇多次提及要我們學習「喜愛公義」,這本是呼籲我們喜愛追求公義的事,但今天很多人喜愛的卻是要「成為公義」,正確點的表達就是只喜愛及擁抱「自己的價值觀」,並以此為「公義」的「唯一判定準則」；當我是「對」的時候,你就是「錯」；反之,當你是「對」的時候,我就會變成了「錯」,當然我不希望自己會變成「錯」,我要成為「公義」/「公正」/「對」的一方時,你必須是「錯」,這正正反映了一個極大的問題,造成人與人之間的撕裂往往不是對事件的意見或立場不同,而是在這種「敵我分明」的價值觀之下,每一個人都要自己成為「公義」,甚至是「真理」,在這種個人主義下,人的自我意識愈發膨脹,難度我們忘記了事實的真相嗎?

「沒有義人,連一個也沒有…因為世人都犯了罪,虧缺了上帝的榮耀.」羅馬書3:10,23

我們眾人皆是罪人,只有上帝是真理,公義的主,我們唯一追求的對象；我們不是要自己成為「公義」,成為上帝,我們在世目的是要追求聖潔,以致我們能夠近一點去認識,親近上帝.

那麼,耶穌基督又教導我們如何回應得罪我們的人呢?

「倘若你的弟兄得罪你,你就去,趁著只有他和你在一處的時候,指出他的錯來.他若聽你,你便得了你的弟兄；他若不聽,你就另外帶一兩個人同去,要憑兩三個人的口作見證,句句都可定準.若是不聽他們,就告訴教會；若是不聽教會,就看他像外邦人和稅吏一樣.」馬太福音18:15-17

耶穌非常直接地指出要努力去與人和好,更有足夠清楚的指引,一步一步的做法,絕不含糊；經文很清楚告訴我們,不可以以「不明白」,「不懂」為借口,如若沒有處理這問題,就只有「不願意」這原因!

耶穌繼續提到:

「我實在告訴你們,凡你們在地上所捆綁的,在天上也要捆綁；凡你們在地上所釋放的,在天上也要釋放.我又告訴你們,若是你們中間有兩個人在地上同心合意地求甚麼事,我在天上的父必為他們成全.因為無論在哪裡,有兩三個人奉我的名聚會,那裡就有我在他們中間.」馬太福音18:18-20

我們可能從來沒有留意到,耶穌在此賜予門徒/我們非常厲害的權柄,我們能「捆綁」及「釋放」,並且影響到天上!我們有否想過,這耶穌所賜的「權柄」是要我們向誰及在何時使用的呢?當然沒有理由是為滿足自己的私慾吧…其實經文接著提到的才是更重要,因為與更大的恩典,祝福相關,耶穌提到當「中間有兩個人」,「有兩三個人」,「同心合意」「奉主的名」禱告或聚會時,主就會在他們當中,能夠有主同在難道不是比權力/權柄更為重要的追求嗎?如何行使「捆綁」,「釋放」這些權柄是一種選擇,但耶穌提到的「你們」,「兩個人」,「兩三個人」是能夠運用這選擇權柄的前題,預設,亦可以說是一種必須的要求,對應經文的上文,就可以連繫成一幅圖畫.當我們信徒要行使主所賜改變世界,甚至天上的權柄以先,就必須先與我們的弟兄姊妹們結連,這種結連必需由和好開始.還記得耶穌基督在山上寶訓中的教導嗎?

「所以,你在祭壇上獻禮物的時候,若想起弟兄向你懷怨,就把禮物留在壇前,先去同弟兄和好,然後來獻禮物.」 馬太福音5:23-24

在耶穌基督眼中與弟兄姊妹和好甚至比向上帝獻祭更為緊要!

「復和」是上帝所賜極重要的禮物,亦是極大的權柄,我們需要努力去「捆綁」那位說謊之人的父─魔鬼,「釋放」被仇恨蒙蔽眼目的人；因此「復和」不是盲目不分黑白地單向寬恕對方(這只為自己感覺良好而已),因為我們仍然需要處理「罪」的問題,不過「復和」是一個以「人」為本的思考模式,目的是要先處理我們這些「罪人」與其他「罪人」的關係,才有力量一同去處理「罪」的問題,正如上帝恨惡罪惡,但深愛我們這些罪人.

「惟有基督在我們還作罪人的時候為我們死,上帝的愛就在此向我們顯明了.」羅馬書5:8

能夠願意,有動力去與人和好,最重要的認知,就是我們眾人全都是蒙恩的罪人,包括你與我,我們皆需要主的恩典,拯救,以致能夠與主復和,與自己復和,因此主更希望我們能夠與他人復和!

到這裡,我希望你能夠以行動去回應今天的信息,因為單單聽了道而不去行的是沒有意義的!

就在這刻,如果你心裡有掙扎,想起有得罪你,又或你曾得罪過的人,可能是你的同事,朋友,弟兄姊妹,甚至可能是家人,請你現在就馬上去進行尋求復和的行動,可以是打電話,發短訊,寫卡片等,直接的由「對不起」開始,因為我們都是浪子,得罪了天,又得罪了他人…完成後才回來繼續崇拜及奉獻,如果收看錄播的,就先暫停影片,回來才繼續,請現在就行動,我會為大家禱告.

請禱告…



\newpage

\section{痛與怒愛與恕}
\label{sec:517}
\textbf{林豪恩先生 - 痛與怒,愛與恕}
\newline
\newline
\hyperref[sec:516]{< < < PREV SERMON < < <}
~
\hyperref[sec:toc]{[返目錄]}
~
\hyperref[sec:518]{> > > NEXT SERMON > > >}


詩篇 103:8-13
\newline
\begin{longtable}{cl}
\hline
\hline
章節 & 經文 (和合本修訂版)\\
\hline
103:8 & \begin{tabularx}{0.7\textwidth}{X} 耶和華有憐憫，有恩惠，不輕易發怒，且有豐盛的慈愛。 \end{tabularx} \\ \\ \relax
103:9 & \begin{tabularx}{0.7\textwidth}{X} 他不長久責備，也不永遠懷怒。 \end{tabularx} \\ \\ \relax
103:10 & \begin{tabularx}{0.7\textwidth}{X} 他沒有按我們的罪待我們，也沒有照我們的罪孽報應我們。 \end{tabularx} \\ \\ \relax
103:11 & \begin{tabularx}{0.7\textwidth}{X} 天離地何等的高，他的慈愛向敬畏他的人也是何等的大！ \end{tabularx} \\ \\ \relax
103:12 & \begin{tabularx}{0.7\textwidth}{X} 東離西有多遠，他叫我們的過犯離我們也有多遠！ \end{tabularx} \\ \\ \relax
103:13 & \begin{tabularx}{0.7\textwidth}{X} 父親怎樣憐憫他的兒女，耶和華也怎樣憐憫敬畏他的人！ \end{tabularx} \\ \\
[1ex]
\hline
\hline
\end{longtable}
講題:痛與怒,愛與恕 Hurt and Anger, Mercy and Forgiveness

講員:林豪恩先生

經文:詩篇103:8-13

場合:香港中文大學崇基學院禮拜堂主日崇拜

日期:2020年9月13日

引言

請你想想一個重要的家人,例如你爸爸或媽媽,你的丈夫或太太,他／她最常表露的是甚麼情緒呢?請你想一想自己,你經常感受到的情緒又是甚麼呢?有一位心理學家寫了一本書名為「情緒四重奏」,提出四種主要的情緒,分別是「擔憂,憤怒,悲哀,喜樂」,簡稱為「憂怒哀樂」.剛才的兩個問題你的答案是甚麼呢?你想起的那位家人經常表露的情緒是甚麼?俗語說:「養兒一百歲,長憂九十九」,你的父母最常表露的情緒會否是擔憂呢?面對著因為擔憂而凡事過問,喋喋不休的父母,你最容易出現的反應會否是憤怒呢?有沒有留意聖經作者也用情緒詞彙來描述上帝呢?在你的印象中,最常用來描述上帝的情緒詞彙是甚麼?讓我們透過今天的經課進入上帝與我們的情緒世界,體會當中的情感交纏.

一. 上主的痛與怒,愛與恕

古代有位屬靈前輩,以幾句詩描述上帝的情感,收錄在詩篇103:8-13節,讓我們一齊來體會一下.

1. 耶和華的怒

103:8耶和華有憐憫,有恩典,不輕易發怒,且有豐盛的慈愛.

103:9他不長久責備,也不永遠懷怒.

「耶和華……不輕易發怒」,即是說上帝是會發怒的；「他不長久責備」,即是說上帝會責備的；又說:「也不永遠懷怒」,即是說上帝會懷怒的.

憤怒是常見的情緒,也可以變成危險和自制困境.新聞報導說:「某天凌晨3時許,一名15歲少女在家中打機,其父親帶醉返家,見女兒仍未就寢,怒火中燒,盛怒出手向她推撞,以致其失平衡身體撼向家具受傷,遇襲少女報警.警方到場調查,救護員將受傷事主送院治理,躁父涉嫌襲擊致造成實際身體傷害被捕.」發怒的人都認為自己是有理由的,如果你是案件中的父親,見到女兒半夜三更仍在打機而不去睡覺,也可能憤怒.如何處理憤怒,如何回應令你憤怒的人或事,會影響事情的發展及結果.耶和華如何處理憤怒呢?祂如何回應令祂憤怒的人呢?

2. 耶和華的恕

103:10他沒有按我們的罪過待我們,也沒有照我們的罪孽報應我們.

即是說,上帝是可以「按我們的罪過待我們」,也可以「照我們的罪孽報應我們」的.如果上帝如此回應我們,我們也無話可說.上帝沒有這樣報應我們,顯示出上帝對我們的饒恕.到底是甚麼令上帝饒恕帶著「罪過」,「罪孽」的我們呢?是甚麼情緒與憤怒互動,在上帝裡面發揮作用,並且影響了結果呢?你有沒有被饒恕的經驗呢?你猜想對方饒恕你的時候心裡面是甚麼情緒在主導呢?你有沒有饒恕別人的經驗呢?當你饒恕別人的時候,是甚麼情緒在你心中運作呢?

3. 耶和華的痛和憐恤

詩人除了提及耶和華會「發怒」,「責備」,「懷怒」之外,也指出耶和華有「恩典」和「慈愛」,並且以父親對兒女的憐恤來描述耶和華對人的感情:

103:13父親怎樣憐恤他的兒女,耶和華也怎樣憐恤敬畏他的人!

一個人又「發怒」,又「責備」,又「懷怒」,同時又有「恩典」,「慈愛」和「憐恤」,是不是很矛盾呢?是不是很難理解呢?詩人體會到,上帝就是這樣.如果你認識的上帝不是這樣,反而可能是你的認識偏差了.難道我們自己不也是這樣的嗎?你會「發怒」,「責備」,但同時又有「恩典」,「慈愛」和「憐恤」.如果你不知道自己都是這樣,可能是你對自己的認識不足；如果你真的不是兩者兼備而只有單方面,反而可能為你製造人際關係的困難.如果你是父母,你的兒女一定曾經令人發怒,你也一定曾經責備過他們；然而,發怒和責備一定不是親子關係的全部,必定還包括慈愛和憐恤.

大家有沒有想過,憤怒是如何轉變為憐恤的呢?到底是甚麼產生這樣的轉化呢?如果要明白當中的運作,我們就需要探究憤怒這種情緒.有人認為憤怒是一種「衍生情緒」(secondary emotion),是對某種基本感覺的回應,例如:失望,嫉妒,羞愧,傷痛……等.剛才提及那位父親發現女兒半夜三點仍在打機而憤怒,可能因為感到失望；有些父母接到兒女不合格的成績表而憤怒,可能是因為失望再加上羞愧.如果只是失望和羞愧,未必會轉化為憐恤；但如果那位發現女兒半夜三點仍在打機而不睡覺的女兒感覺到自己心痛,他的憤怒就有機會轉化,他的回應也有可能不同.作為父母的,從兒女哇哇一聲出生開始,一定曾經為兒女感到過心痛,可能是在兒女患病的時候,可能是在兒女受傷的時候,也可是兒女做了一些令你失望的事情的時候.詩人說:「父親怎樣憐恤他的兒女,耶和華也怎樣憐恤敬畏他的人!」耶和華對人的憐憫,隱含著父母對兒女的「心痛」.很可能是這種對我們的「心痛」,把上帝的憤怒轉化為憐恤,表現為饒恕,因此「沒有按我們的罪過待我們,也沒有照我們的罪孽報應我們」.

4. 得到愛與恕的是誰呢?

既憤怒,又心痛；既失望,又憐憫；既饒恕,又不想放下要求……如果你感到自己的情緒多元而複雜,甚至有時互相矛盾,不用感到驚訝,因為上帝與我們之間的關係也是這樣.人一方面令上帝感到憤怒,同時又令上帝心痛；一方面蒙上帝憐恤,但又感到上帝對我們有要求:

103:11天離地何等的高,他的慈愛向敬畏 (fear) 他的人也是何等的大!

103:12東離西有多遠,他叫我們的過犯離我們也有多遠!

103:13父親怎樣憐恤他的兒女,耶和華也怎樣憐恤敬畏 (fear) 他的人!

上帝的慈愛並非盲目,上帝的憐恤並非放縱,上帝的恩典帶著血與淚,上帝的寬恕並非放棄公義.這些看似互相矛盾的情感,在上帝與我們的互動中呈現,其實,也在人與人之間的關係中呈現.今天的舊約經課創世記50:15-21所記載的約瑟與兄弟們的恩怨情仇就是一個例子.

二. 約瑟與兄弟們的痛與怒,愛與恕

約瑟的恨和怒

50:15約瑟的哥哥們見父親死了,就說:「或者約瑟懷恨我們,照著我們從前待他一切的惡足足地報復我們.」

我們不在這裡重述約瑟與兄弟們的恩怨情仇,但從哥哥們承認「我們從前待他一切的惡」得知約瑟曾經被哥哥們惡待的事實.因此,哥哥們知道,「約瑟懷恨我們」是有背景的,也有理由的.是的,約瑟的確憤恨,被哥哥們出賣,怎會不憤怒呢?除了憤怒呢?會否還有其他情緒呢?或者會否有其他情緒比憤怒更藏於心底和更真摰直接呢?

約瑟的痛和淚

50:16他們就打發人去見約瑟,說:「你父親未死以先吩咐說:

50:17『你們要對約瑟這樣說:從前你哥哥們惡待你,求你饒恕他們的過犯和罪惡.』如今求你饒恕你父親上帝之僕人的過犯.」他們對約瑟說這話,約瑟就哭了.

「約瑟就哭了」,已經不是第一次,由重遇兄弟們開始,約瑟其實已經哭了幾次,只是兄弟們未必知道而已.或者你也能體會約瑟的心情,在不少狀況之下,眼淚只能暗自流下而不能對人哭,因為你感覺到在人面前哭,只會暴露自己的傷口,只會更痛.是的,哭是受傷和感到痛楚的最直接反應,大多數小孩跌倒受傷感到痛楚都會哭.長大了不哭,可能只是被成人太多次訓示「不要哭」而已.當人不再哭的時候,會否就連傷心和心痛也再感受不到呢?「約瑟就哭了」,可見約瑟除了憤怒之外,更多的可能是傷心和心痛.傷心甚麼呢?傷心自己被傷害,更傷心自己的哥哥們曾經犯上如此殘忍的罪惡；心痛甚麼呢?心痛哥哥們被自己的過犯所捆綁,不能坦然地接受寬恕,仍然生活在擔憂和恐懼之下.約瑟的確憤怒,不過,約瑟和哥哥們的故事發展到這個場景,主導的情緒不是憤怒,很可能是心傷和心痛.當人感受到憤怒之下的心傷和心痛,並且有合適的條件讓心傷和心痛表達,事情的發展可能會很不同.如何不同呢?讓我們繼續讀故事的結局:

50:18他的哥哥們又來俯伏在他面前,說:「我們是你的僕人.」

50:19約瑟對他們說:「不要害怕,我豈能代替上帝呢?

50:20從前你們的意思是要害我,但上帝的意思原是好的,要保全許多人的性命,成就今日的光景.

50:21現在你們不要害怕,我必養活你們和你們的婦人孩子.」於是約瑟用親愛的話安慰他們.

三. 耶穌的呼喚

或者,我們不是約瑟,而是約瑟的兄弟,讓我們閱讀福音經課,換另一個身分和角度思考如何在「痛與怒,愛與恕」混雜的現實中作祂的門徒.

18:21那時,彼得進前來,對耶穌說:「主啊,我弟兄得罪我,我當饒恕他幾次呢?到七次可以嗎?」

我們都很容易感受到別人加給自己的痛與怒,而且也覺得自己已經算是好人.有時也會想到,好人都應該有個限度.不過,我們也同樣容易,或者更容易忽視自己對人的虧欠,自己帶給別人的痛與怒:

18:23天國好像一個王要和他僕人算帳.

18:24才算的時候,有人帶了一個欠一千萬銀子的來.

18:25因為他沒有甚麼償還之物,主人吩咐把他和他妻子兒女,並一切所有的都賣了償還.

當我們被逼面對自己的虧欠,被逼補償帶給別人的痛和怒時,我們會深深感到自己的困境與無助,在這情況之下,我們別無他選,唯有請求寬容:

18:26那僕人就俯伏拜他,說:『主啊,寬容我,將來我都要還清.』

18:27那僕人的主人就動了慈心,把他釋放了,並且免了他的債.

當我們的困境和無助化解了之後,我們鬆一口氣對自己說:「我真是幸運!」我們沒有被那個「動了慈心」待我們的人所感動,因此,我們竟然忽視其他人的困境與無助:

18:28「那僕人出來,遇見他的一個同伴欠他十兩銀子,便揪著他,掐住他的喉嚨,說:『你把所欠的還我!』

18:29他的同伴就俯伏央求他,說:『寬容我吧,將來我必還清.』

18:30他不肯,竟去把他下在監裡,等他還了所欠的債.

或者,我們認為欠債還錢是天經地義的,對方為所虧欠我們的付上代價是應該的.我的債主動了慈心寬待我是我的幸運,或者是他的愚蠢而已.

不要把被寬容視為幸運,也不要把憐恤人的人視為愚蠢,世間若不再有寬容和憐恤,可以是很可怕的.

18:31眾同伴看見他所做的事就甚憂愁,去把這事都告訴了主人.

18:32於是主人叫了他來,對他說:『你這惡奴才!你央求我,我就把你所欠的都免了,

18:33你不應當憐恤你的同伴,像我憐恤你嗎?』

18:34主人就大怒,把他交給掌刑的,等他還清了所欠的債.

耶穌講完了故事,留下了這一句教導給門徒:

18:35你們各人若不從心裡饒恕你的弟兄,我天父也要這樣待你們了.」

或許,門徒心裡也迴盪著故事中主人對欠債者發出的問題:

18:33你不應當憐恤你的同伴,像我憐恤你嗎?』

這個問題會否也在你的心中迴盪著呢?

四. 總結

最後,讓我以一位在神學院認識的朋友的故事作結:

他的父親在逃難到香港前是大陸的望族之子,母親則是基督教的傳道人,他們的兒子會是一位怎樣的年輕人呢?這個在龍蛇雜處的調景嶺長大的年輕人,十五歲加入黑社會,成為一個逞凶鬥狠的流氓.為了混出名號,他嘗了人生中第一口嗎啡,從此走上吸毒的不歸路.

如果你是他的媽媽,你的心情會是怎樣的呢?當這位作為傳道人的媽媽向人傳福音時,人們回應:「你叫你的兒子信了以後再來跟我講吧!」「你先教好自己的兒子再來教我們吧!」相信他的媽媽羞愧得無地自容,也被兒子激得七孔生煙.

在黑幫的日子,他曾被打壞眼睛,打碎腕骨,抽斷腳筋,誰會為他的傷感到心痛呢?對手不會,「兄弟」也不會,只有他的媽媽感到心痛.被警察拘捕的時候,警員不會為他感到心痛,只有去辦理保釋手續的媽媽感到心痛.媽媽曾經送他去台灣讀書,以為遠離是非之地他就能夠重新做人,可惜,他回到香港再後繼續在吸毒與戒毒之間不斷輪迴.

他憶述,在這段黑暗的歲月中,母親無私的愛是他唯一的依靠.當他在毒海中反復掙扎時,是他母親一次一次不放棄地向神祈禱,一聲聲流著淚將他喚回.十年之後,他終於走出黑暗,第一次看到人生的希望與未來.

他後來把自己的經歷編輯成書,名為「永不放棄的愛:劉民和牧師的生命與事奉」,細述他從吸毒的流氓到傳播福音戒毒的牧師的心路歷程.

這本書不單單是劉民和牧師的見證,更是是他的媽媽的「痛與怒,愛與恕」的見證.



\newpage

\section{在承擔責任的生命中活出永生}
\label{sec:518}
\textbf{鄧瑞強博士 - 在承擔責任的生命中活出永生}
\newline
\newline
\hyperref[sec:517]{< < < PREV SERMON < < <}
~
\hyperref[sec:toc]{[返目錄]}
~
\hyperref[sec:519]{> > > NEXT SERMON > > >}


腓立比書 1:20-30
\newline
\begin{longtable}{cl}
\hline
\hline
章節 & 經文 (和合本修訂版)\\
\hline
1:20 & \begin{tabularx}{0.7\textwidth}{X} 這就是我所切慕、所盼望的：沒有一事能使我羞愧；反倒凡事坦然無懼，無論是生是死，總要讓基督在我身上照常顯大。 \end{tabularx} \\ \\ \relax
1:21 & \begin{tabularx}{0.7\textwidth}{X} 因為我活著就是基督，死了就有益處。 \end{tabularx} \\ \\ \relax
1:22 & \begin{tabularx}{0.7\textwidth}{X} 但是，我在肉身活著，若能有工作的成果，我就不知道該挑選甚麼。 \end{tabularx} \\ \\ \relax
1:23 & \begin{tabularx}{0.7\textwidth}{X} 我處在兩難之間：我情願離世與基督同在，因為這是好得無比的； \end{tabularx} \\ \\ \relax
1:24 & \begin{tabularx}{0.7\textwidth}{X} 然而，我為你們肉身活著更加要緊。 \end{tabularx} \\ \\ \relax
1:25 & \begin{tabularx}{0.7\textwidth}{X} 既然我這樣深信，就知道仍要留在世間，且與你們眾人一起存留，使你們在所信的道上又長進又喜樂， \end{tabularx} \\ \\ \relax
1:26 & \begin{tabularx}{0.7\textwidth}{X} 為了我再到你們那裡時，你們在基督耶穌裡的誇耀越發加增。 \end{tabularx} \\ \\ \relax
1:27 & \begin{tabularx}{0.7\textwidth}{X} 最重要的是：你們行事為人要與基督的福音相稱，這樣，無論我來見你們，或不在你們那裡，都可以聽到你們的景況，知道你們同有一個心志，站立得穩，為福音的信仰齊心努力， \end{tabularx} \\ \\ \relax
1:28 & \begin{tabularx}{0.7\textwidth}{X} 絲毫不怕敵人的威脅；以此證明他們會沉淪，你們會得救，這是出於神。 \end{tabularx} \\ \\ \relax
1:29 & \begin{tabularx}{0.7\textwidth}{X} 因為你們蒙恩，不但得以信服基督，而且要為他受苦。 \end{tabularx} \\ \\ \relax
1:30 & \begin{tabularx}{0.7\textwidth}{X} 你們的爭戰，就與你們曾在我身上見過、現在所聽到的是一樣的。 \end{tabularx} \\ \\
[1ex]
\hline
\hline
\end{longtable}
講題:在承擔責任的生命中活出永生  Everlasting Life in Responsibility

經文:腓立比書1:20-30

講員:鄧瑞強博士

場合:香港中文大學崇基學院禮拜堂主日崇拜

日期:2020年9月20日

各位弟兄姊妹,早安,願各人在疫情中平安.

這次疫症,維時很長,並且看來還會維持一段日子.疫症下的生活,不太正常,因死亡比平時顯得更近.我們平時都知道人會死,但很少想到會即刻死.我們平時都知道人會死,但很少想到會死的是自己.死亡是真實的,但在平時,這只是一遠距離的圖畫.在疫情下,這幅原先遠距離的圖畫,被拉近變成一幅近距離的放大圖.日日都顯示,多少人染病.日日都報導,多少人因染病而離世.死亡被拉近,被放大.當事物一近,被放大,就變得怪異,甚至恐怖.中學時代,用顯微鏡觀察小昆蟲.遠看小昆蟲,不太可怕.但在顯微鏡下,小昆蟲被放大,牠突然一動,便嚇得從顯微鏡那裡彈開來.疫症,將死亡拉近,放大,死亡便變得更加真實,更加恐怖,這使我們的生活變得不太正常.

很多時,為了實驗,會製造極端環境,看看實驗對象的反應.譬如:刻意不給水一棵植物,看看它能捱多少日子.刻意不給食物一隻昆蟲,看看它能不進食活多長時間.在疫症下,死亡好像很近,這變成某種極端環境,將生活變成一實驗場所,考驗著人性.這是一人性實驗的時刻.有一次,坐巴士,一位乘客對另一位乘客大聲呼喝,叫他戴口罩.呼喝的人,看來失去了面對另一個人應有的基本禮貌.呼喝的人,看來受到了極大衛脅,故不得不作出一種大反應.呼喝的人,明顯認為對方犯了不可原諒的錯,以致自己可理直氣壯地宣告判語.那位乘客,我不知他為何不戴口罩,是否他一時大意,是否他認為自己的自由比他人的疑慮更重要,是否他根本對這病毒不屑一顧.平時看不見的人性表達,在疫情下,在死亡的威脅下,被放大,被看得更清楚.死亡的臨近,是人性的放大鏡.

今日,我們看看保羅在他最後的人生階段裡寫下的一段文字,這段文字反映保羅正面對死亡的威脅.在死亡的陰影下,保羅反省他的一生,他一生的志業,他一生的意義.他領悟到基督復活的重要性,體會到自己活在基督裡,在必死的生命中有幸活出基督復活的真相.在基督裡,他領悟到超越生死的真理.

講道經文:腓立比書1:20－30

腓1:20　照著我所切慕,所盼望的,沒有一事叫我羞愧.只要凡事放膽,無論是生是死,總叫基督在我身上照常顯大.

腓1:21　因我活著就是基督,我死了就有益處.

腓1:22　但我在肉身活著,若成就我工夫的果子,我就不知道該挑選甚麼.

腓1:23　我正在兩難之間,情願離世與基督同在,因為這是好得無比的.

腓1:24　然而,我在肉身活著,為你們更是要緊的.

腓1:25　我既然這樣深信,就知道仍要住在世間,且與你們眾人同住,使你們在所信的道上又長進又喜樂,

腓1:26　叫你們在基督耶穌裡的歡樂,因我再到你們那裡去,就越發加增.

腓1:27　只要你們行事為人與基督的福音相稱,叫我或來見你們,或不在你們那裡,可以聽見你們的景況,知道你們同有一個心志,站立得穩,為所信的福音齊心努力.

腓1:28　凡事不怕敵人的驚嚇,這是證明他們沉淪,你們得救都是出於神.

腓1:29　因為你們蒙恩,不但得以順服基督,並要為他受苦.

腓1:30　你們的爭戰,就與你們在我身上從前所看見,現在所聽見的一樣.

腓立比書,是保羅在羅馬坐監時寫的「監獄書信」,是保羅最後期的作品之一.在這書信裡,保羅隱約感到,生命可能到盡頭了.

(1)開場白:保羅的願望

剛才讀的經文,有個小小的開場白,保羅講出他的渴望.

腓1:20　照著我所切慕,所盼望的,沒有一事叫我羞愧.只要凡事放膽,無論是生是死,總叫基督在我身上照常顯大.

簡單一節經文,總結了保羅一生的渴望和體會.保羅的渴望,是人生無憾,一生無愧.「沒有一事叫我羞愧」.如何做到?就是「凡事放膽」,充滿勇氣地面對人生,生也好,死也好,只想用整個生命來榮耀基督.

保羅的眼目,沒有注視自己的生與死,沒有計量自己的生與死帶來什麼天翻地覆的改變,他根本不注視自己,只注視基督.他渴望透過自己這必死的軀殼,能榮耀基督.

我有一位老師－－楊錫鏘牧師,這個禮拜離開了人間.他說過,這生沒有遺憾.在我看來,他真是沒有遺憾的,他回應了神的呼召,一生忠心的教導聖經.他說過:「真正讓人記得的應是上帝的恩典,而不是人的貢獻.」

(參:\url{https://christiantimes.org.hk/Common/Reader/News/ShowNews.jsp?Nid=163330 and Pid=102 and Version=0 and Cid=2011 and Charset=big5_hkscs})像保羅這裡講的一樣,他不注視自己,只注視基督.這樣的人,不可能有遺憾.

一個放下了自己,而努力榮耀基督的人,他會在人生中如何抉擇?

(2)永生－－就在深深地活在塵世裡承擔起責任

腓1:21　因我活著就是基督,我死了就有益處.

腓1:22　但我在肉身活著,若成就我工夫的果子,我就不知道該挑選甚麼.

腓1:23　我正在兩難之間,情願離世與基督同在,因為這是好得無比的.

腓1:24　然而,我在肉身活著,為你們更是要緊的.

腓1:25　我既然這樣深信,就知道仍要住在世間,且與你們眾人同住,使你們在所信的道上又長進又喜樂,

腓1:26　叫你們在基督耶穌裡的歡樂,因我再到你們那裡去,就越發加增.

開宗明義,保羅說:「我活著就是基督」,他沒有自己的人生,也不追求自己的成功,他只擁抱基督.

基督的人生,是一種怎麼樣的人生?

一種聆聽天父旨意的人生,一種順服天父帶領的人生,一種與天父連合的人生.用另一種說話說,是一種神人結合的人生,是一種進入了永恆的人生,是一種幸福無比的人生.即刻進入這人生,是夢魅以求的.

保羅知道,離開這世界,不是一種虧損,反而是一種益處,一種盈利.他說,離開世界,「就有益處」,「好得無比」.離開這世界,保羅知道,這只會讓他與基督有更深的連繫,能更深進入基督的生命裡.對飽經風霜的保羅來說,這是好得無比的選擇.

設想有一天,有天使來探訪你,帶你到天堂走一回.去到天堂,那裡有座豪宅,門口寫著你的名字,裡面約有一萬呎,有花園,有游水池,有音響室.在那裡,你可種花,可養狗,可和最好的朋友野餐.總之,是香港人夢魅以求的生活.天使說,你隨時可以來這裡.你會如何想?

保羅在思想「離世是好得無比」時,他同時在想人間的責任.1:22節說:「成就我工夫的果子」,《現代中文譯本》譯作:「多做些有益的工作」,《呂振中譯本》譯作:「多結工作的果子」.簡單說,他想起他肩頭上的責任,他想起他承擔著的十字架,他想起基督的十字架,他想起基督呼喚他背起十架.保羅有自己的渴望,但他的心被基督的呼喚充滿.飽經風霜的保羅,一方面,想放下責任；另一方面,更深的明白他塵世的責任的意義.

在1:22,保羅說:「我不知道該挑選甚麼」.面對甘甜的美果及苦澀的重擔,平凡人一定會選擇甘甜的美果.保羅不這樣選,而說「不知道該怎樣選」,他能這樣說,其實已表示他知道怎樣選.他,在基督裡,不再是一平凡人,他像基督一樣活著,選擇了為他人的福祉而活的人生,選擇了背負他人十架的命途.在1:24,保羅說:「我在肉身活著,為你們更是要緊的.」在1:25,保羅說:「我既然這樣深信,就知道仍要住在世間,且與你們眾人同住,使你們在所信的道上又長進又喜樂.」他放下了自己,他想著他人,他掛念著他人的福祉.

在上一次的沙士疫症期間,在這個禮拜堂,我們悼念過一個人物,我們悼念過謝婉雯醫生.在疫症期間,她沒有掛慮自己的安危,主動請纓,進入深切治療病房,照顧沙士病人.當一般人急於遠離沙士病房時,她選擇進入去照顧他人.她被感染,患病時,仍掛念其他病人的景況.她後來離世,她真是「香港的女兒」,真是天父的女兒,真是一個「活著就是基督」的人.她,作了重大抉擇,為了他人的痊癒,放下了自己.在此生中,活出了永恆.在死亡中,活出了基督的復活.

我們常常想,「生命」與「死亡」是兩極,「今日」與「將來」是兩個不同時段,「當下」與「永恆」是沒有關連的.好像選擇了這樣,便得放棄那樣似的.好像若保羅選擇「死亡」,選擇那好得無比的永恆時,他便不用再在塵世中糾纏下去,像徐志摩所說的那樣,「我揮一揮衣袖,不帶走一片雲彩」.又好像若保羅選擇了塵世的責任時,他便得暫時與永恆的福樂告別,只有苦澀,沒有甘甜.我們斬斷了「永恆」與「現世」的連繫,好像耶穌死了就是死了,復活是與死亡無關的另外一回事似的.

保羅卻不是這樣想的.他知道,在耶穌的死亡中,有復活的種子.在耶穌的復活中,有死亡的影子.塵世和永恆,死亡和復活,今世與來世,在耶穌身上,從來都不是斷裂的.當耶穌在塵世背起十架時,他已向人間展露天父永恆的大愛.耶穌以塵世的十架,展露神聖的永恆.永恆不在他世,永恆在此間的十架中.耶穌的死,其實是顯露真正的生.不是死亡之後,才有復活.死亡之中,已透顯出復活的生命,或者說,真正永恆的生命.保羅完全明白耶穌的死亡與復活的互相關連,明白耶穌以塵世的十架去展示永恆的光明,明白耶穌會合了人性和神性,於是,保羅選擇背起人間的十架,藉此去體會永恆生命的福樂.

保羅在思考人生時,他想到永恆的生命,他想到自己的死亡,他想到死亡後便能更深與神合一,他想到這裡,便開心不已.但當他想到這永恆生命的豐盛時,他又回頭想,這永恆生命其實已然在這塵世生命中顯露,那將來的復活生命其實就在現世,就在甘願為他人背負的十架中.原來,不是死亡後,才步入永恆；而是,今生與來世,都在永恆之中.如此,就不用離世去進入永恆；而是,深入這塵世生活中,承擔起他人的十架,在十架裡進入永恆.在背負人生的責任中,進入基督的死亡.在進入基督的死亡裡,體會復活的力量.如此,死亡與生命,今生與來世,不再是二擇其一,而是將兩者都放在基督手上,讓基督的死而復活將兩者合起來.一個進入永恆的人,就是活出基督的人,是深深地進入塵世,背起十架.

若果信耶穌,只是沉醉著天堂的福樂裡,而忘記人間的苦難,忘記在苦難中自己應盡的責任,則在保羅眼中,這人仍未明白信仰,仍未明白基督,仍未明白永生.

真正的基督信仰,是將我們帶回人間,更深刻地看到他人的需要,更勇敢地承擔責任,更切實地期盼天國在人間.

(3)鼓勵信徒,進入基督的生命裡

保羅,活出了基督的死與生,他希望信徒也能如此.

腓1:27　只要你們行事為人與基督的福音相稱,叫我或來見你們,或不在你們那裡,可以聽見你們的景況,知道你們同有一個心志,站立得穩,為所信的福音齊心努力.

腓1:28　凡事不怕敵人的驚嚇,這是證明他們沉淪,你們得救都是出於神.

腓1:29　因為你們蒙恩,不但得以順服基督,並要為他受苦.

腓1:30　你們的爭戰,就與你們在我身上從前所看見,現在所聽見的一樣.

「行事為人與基督的福音相稱」,意思是:我們要明白自己作為天國子民的身分,整個生活就要配合這身分.所謂天國子民,就是在塵世像基督一樣活著的人.這些人:堅持天父旨意,堅持聖經真理,堅持愛人如己；不喜歡暴力,不喜歡虛假,不喜歡泯滅人性.但是,如此順乎天理的事情,在人間,原來也是極其困難的.

故此,保羅鼓勵我們,要像基督的精兵,要有同一心志,要有同一目標,為基督奮鬥.「活著就是基督」,這是艱難的,但不要害怕.

1:29很重要:「因為你們蒙恩,不但得以順服基督,並要為他受苦.」蒙恩,當然能得救.但所謂得救,是生命變得像基督一樣,甘願背起人間的十架,為他人受苦,為基督受苦.唯有這樣的基督徒,在人間,像基督一樣,開展天國.

基督走過人間,為人而死,然後,復活過來,天父藉此顯明這生命的神聖.保羅,明白基督,活出基督,為他人的好處,背起他們的十架.今日,輪到我們了.保羅鼓勵我們說:「你們的爭戰,就與你們在我身上從前所看見,現在所聽見的一樣.」我們能否像保羅一樣,力挽狂瀾,活出基督呢?這意味著什麼?這意味著:更深的進入他人的苦難中,為他們的好處,盡上自己的責任.

願榮耀歸上主,願人間享和平.

(附記:本人在崇基禮拜堂的講章已結集成《荷花池畔春草綠》講道集,試讀本可參閱:\url{http://cc.cuhk.edu.hk/dev/images/content/fund-raising-campaigns/university-sermons-by-dr-sk-tang/CCC-SK-Tang-Sermon-Text-Sample-2.pdf}.較早期的崇基講章,也已結集成《敢於跟隨主》一書.有興趣的朋友可到各大基督教書局購買.)



\newpage

\section{你仗著甚麼權柄}
\label{sec:519}
\textbf{伍渭文牧師 - 你仗著甚麼權柄?}
\newline
\newline
\hyperref[sec:518]{< < < PREV SERMON < < <}
~
\hyperref[sec:toc]{[返目錄]}
~
\hyperref[sec:520]{> > > NEXT SERMON > > >}


馬太福音 21:23-32
\newline
\begin{longtable}{cl}
\hline
\hline
章節 & 經文 (和合本修訂版)\\
\hline
21:23 & \begin{tabularx}{0.7\textwidth}{X} 耶穌進了聖殿，正教導人的時候，祭司長和百姓的長老來問他：「你仗著甚麼權柄做這些事？給你這權柄的是誰呢？」 \end{tabularx} \\ \\ \relax
21:24 & \begin{tabularx}{0.7\textwidth}{X} 耶穌回答他們說：「我也要問你們一句話，你們若告訴我，我就告訴你們我仗著甚麼權柄做這些事。 \end{tabularx} \\ \\ \relax
21:25 & \begin{tabularx}{0.7\textwidth}{X} 約翰的洗禮是從哪裡來的？是從天上來的，還是從人間來的呢？」他們彼此商議說：「我們若說『從天上來的』，他會對我們說：『這樣，你們為甚麼不信他呢？』 \end{tabularx} \\ \\ \relax
21:26 & \begin{tabularx}{0.7\textwidth}{X} 若說『從人間來的』，我們又怕眾人，因為大家都認為約翰是先知。」 \end{tabularx} \\ \\ \relax
21:27 & \begin{tabularx}{0.7\textwidth}{X} 於是他們回答耶穌：「我們不知道。」耶穌也對他們說：「我也不告訴你們，我仗著甚麼權柄做這些事。」 \end{tabularx} \\ \\ \relax
21:28 & \begin{tabularx}{0.7\textwidth}{X} 「有一件事，你們的意見如何？一個人有兩個兒子。他來對大兒子說：『孩子，今天到葡萄園裡做工去。』 \end{tabularx} \\ \\ \relax
21:29 & \begin{tabularx}{0.7\textwidth}{X} 他回答：『我不去』，以後自己懊悔，就去了。 \end{tabularx} \\ \\ \relax
21:30 & \begin{tabularx}{0.7\textwidth}{X} 他來對小兒子也是這樣說。他回答：『父親大人，我去』，卻不去。 \end{tabularx} \\ \\ \relax
21:31 & \begin{tabularx}{0.7\textwidth}{X} 這兩個兒子是哪一個照著父親的意願做了呢？」他們說：「大兒子。」耶穌說：「我實在告訴你們，稅吏和娼妓倒比你們先進神的國。 \end{tabularx} \\ \\ \relax
21:32 & \begin{tabularx}{0.7\textwidth}{X} 因為約翰到你們這裡來指引你們走義路，你們卻不信他，稅吏和娼妓倒信了他。你們看見了以後，還是不悔悟去信他。」 \end{tabularx} \\ \\
[1ex]
\hline
\hline
\end{longtable}
以西結書 18:1-4,25-32
\newline
\begin{longtable}{cl}
\hline
\hline
章節 & 經文 (和合本修訂版)\\
\hline
18:1 & \begin{tabularx}{0.7\textwidth}{X} 耶和華的話又臨到我，說： \end{tabularx} \\ \\ \relax
18:2 & \begin{tabularx}{0.7\textwidth}{X} 「你們在以色列地何以有這俗語，『父親吃了酸葡萄，兒子牙齒就酸倒』呢？ \end{tabularx} \\ \\ \relax
18:3 & \begin{tabularx}{0.7\textwidth}{X} 主耶和華說：我指著我的永生起誓，你們在以色列必不再引用這俗語。 \end{tabularx} \\ \\ \relax
18:4 & \begin{tabularx}{0.7\textwidth}{X} 看哪，所有的生命都是屬我的；父親的生命怎樣屬我，兒子的生命也照樣屬我；然而犯罪的，他必定死。 \end{tabularx} \\ \\ \relax
18:5 & \begin{tabularx}{0.7\textwidth}{X} 「人若是公義，行公平公義的事： \end{tabularx} \\ \\ \relax
18:6 & \begin{tabularx}{0.7\textwidth}{X} 未曾在山上吃祭物，未曾向以色列家的偶像舉目；未曾污辱鄰舍的妻，也未曾在婦人的經期間親近她； \end{tabularx} \\ \\ \relax
18:7 & \begin{tabularx}{0.7\textwidth}{X} 未曾虧負人，而是將欠債之人的抵押品還給他；未曾搶奪人的物件，卻把食物給飢餓的人吃，把衣服給赤身的人穿； \end{tabularx} \\ \\ \relax
18:8 & \begin{tabularx}{0.7\textwidth}{X} 未曾向人取利息，也未曾索取高利，反倒縮手不作惡，在人與人之間施行誠實的判斷； \end{tabularx} \\ \\ \relax
18:9 & \begin{tabularx}{0.7\textwidth}{X} 遵行我的律例，謹守我的典章，按誠實行事；這人是公義的，必要存活。這是主耶和華說的。 \end{tabularx} \\ \\ \relax
18:10 & \begin{tabularx}{0.7\textwidth}{X} 「他若生了兒子，兒子作強盜，流人的血，作父親的雖然未犯此過，兒子卻對弟兄行了以上所說的惡，在山上吃祭物，污辱鄰舍的妻； \end{tabularx} \\ \\ \relax
18:11 & \begin{tabularx}{0.7\textwidth}{X} 見上節 \end{tabularx} \\ \\ \relax
18:12 & \begin{tabularx}{0.7\textwidth}{X} 虧負困苦和貧窮的人，搶奪別人的物件，不歸還抵押品，卻向偶像舉目，做可憎的事； \end{tabularx} \\ \\ \relax
18:13 & \begin{tabularx}{0.7\textwidth}{X} 向人取利息，索取高利，這人豈能存活呢？他不能存活。他因做這一切可憎的事，必要死亡，他的血要歸到自己身上。 \end{tabularx} \\ \\ \relax
18:14 & \begin{tabularx}{0.7\textwidth}{X} 「看哪，他若生了兒子，兒子見父親所犯的一切罪，他見了，卻不照樣去做； \end{tabularx} \\ \\ \relax
18:15 & \begin{tabularx}{0.7\textwidth}{X} 他未曾在山上吃祭物，未曾向以色列家的偶像舉目，未曾污辱鄰舍的妻； \end{tabularx} \\ \\ \relax
18:16 & \begin{tabularx}{0.7\textwidth}{X} 也未曾虧負人，未曾取人的抵押品，未曾搶奪人的物件，卻把食物給飢餓的人吃，把衣服給赤身的人穿， \end{tabularx} \\ \\ \relax
18:17 & \begin{tabularx}{0.7\textwidth}{X} 縮手不害困苦人，未曾向人索取利息或高利；反倒順從我的典章，遵行我的律例；如此，他必不因父親的罪孽死亡，定要存活。 \end{tabularx} \\ \\ \relax
18:18 & \begin{tabularx}{0.7\textwidth}{X} 至於他父親，因為施行欺壓，搶奪弟兄，在百姓中行不善，看哪，他必因自己的罪孽死亡。 \end{tabularx} \\ \\ \relax
18:19 & \begin{tabularx}{0.7\textwidth}{X} 「你們還說：『兒子為甚麼不擔當父親的罪孽呢？』兒子若行公平公義的事，謹守遵行我一切的律例，他必要存活。 \end{tabularx} \\ \\ \relax
18:20 & \begin{tabularx}{0.7\textwidth}{X} 惟有犯罪的，卻必死亡。兒子不擔當父親的罪孽，父親也不擔當兒子的罪孽。義人的善果要歸自己，惡人的惡報也要歸自己。 \end{tabularx} \\ \\ \relax
18:21 & \begin{tabularx}{0.7\textwidth}{X} 「惡人若回轉離開所做的一切罪惡，謹守我的一切律例，行公平公義的事，他必要存活，不致死亡。 \end{tabularx} \\ \\ \relax
18:22 & \begin{tabularx}{0.7\textwidth}{X} 他所犯的一切罪過都不被記念；他因所行的義，必要存活。 \end{tabularx} \\ \\ \relax
18:23 & \begin{tabularx}{0.7\textwidth}{X} 惡人死亡，豈是我所喜悅的呢？我豈不是喜悅他回轉離開所行的道而存活嗎？這是主耶和華說的。 \end{tabularx} \\ \\ \relax
18:24 & \begin{tabularx}{0.7\textwidth}{X} 至於義人，他若轉離義行而作惡，照著惡人所做一切可憎的事去做，豈能存活呢？他所行的一切義都不被記念；反而因所行的惡、所犯的罪死亡。 \end{tabularx} \\ \\ \relax
18:25 & \begin{tabularx}{0.7\textwidth}{X} 「你們卻說：『主的道不公平！』以色列家啊，你們要聽，我的道不公平嗎？你們的道不是不公平嗎？ \end{tabularx} \\ \\ \relax
18:26 & \begin{tabularx}{0.7\textwidth}{X} 義人若轉離義行而作惡，他就因這些惡而死亡。他要死在他所作的惡中。 \end{tabularx} \\ \\ \relax
18:27 & \begin{tabularx}{0.7\textwidth}{X} 惡人若回轉離開所行的惡，行公平公義的事，他必救自己的命； \end{tabularx} \\ \\ \relax
18:28 & \begin{tabularx}{0.7\textwidth}{X} 因為他省察，回轉離開所犯的一切罪過，他必要存活，不致死亡。 \end{tabularx} \\ \\ \relax
18:29 & \begin{tabularx}{0.7\textwidth}{X} 以色列家還說：『主的道不公平！』以色列家啊，我的道不公平嗎？你們的道不是不公平嗎？ \end{tabularx} \\ \\ \relax
18:30 & \begin{tabularx}{0.7\textwidth}{X} 所以，以色列家啊，我必按你們各人所做的審判你們。當回轉，回轉離開你們一切的罪過，免得罪孽成為你們的絆腳石。這是主耶和華說的。 \end{tabularx} \\ \\ \relax
18:31 & \begin{tabularx}{0.7\textwidth}{X} 你們要把所犯的一切罪過盡行拋棄，為自己造一個新的心和新的靈。以色列家啊，你們為甚麼要死呢？ \end{tabularx} \\ \\ \relax
18:32 & \begin{tabularx}{0.7\textwidth}{X} 我不喜歡有任何人死亡，所以你們當回轉，要存活！這是主耶和華說的。」 \end{tabularx} \\ \\
[1ex]
\hline
\hline
\end{longtable}
詩篇 25:1-9
\newline
\begin{longtable}{cl}
\hline
\hline
章節 & 經文 (和合本修訂版)\\
\hline
25:1 & \begin{tabularx}{0.7\textwidth}{X} 耶和華啊，我的心仰望你。 \end{tabularx} \\ \\ \relax
25:2 & \begin{tabularx}{0.7\textwidth}{X} 我的神啊，我素來倚靠你；求你不要叫我羞愧，不要叫我的仇敵向我誇勝。 \end{tabularx} \\ \\ \relax
25:3 & \begin{tabularx}{0.7\textwidth}{X} 凡等候你的必不羞愧，惟有那無故行奸詐的必要羞愧。 \end{tabularx} \\ \\ \relax
25:4 & \begin{tabularx}{0.7\textwidth}{X} 耶和華啊，求你將你的道指示我，將你的路指教我！ \end{tabularx} \\ \\ \relax
25:5 & \begin{tabularx}{0.7\textwidth}{X} 求你指教我，引導我進入你的真理，因為你是救我的神。我整日等候你。 \end{tabularx} \\ \\ \relax
25:6 & \begin{tabularx}{0.7\textwidth}{X} 耶和華啊，求你記念你的憐憫和慈愛，因為這是亙古以來所常有的。 \end{tabularx} \\ \\ \relax
25:7 & \begin{tabularx}{0.7\textwidth}{X} 求你不要記得我幼年的罪愆和我的過犯；耶和華啊，求你因你的良善，按你的慈愛記念我。 \end{tabularx} \\ \\ \relax
25:8 & \begin{tabularx}{0.7\textwidth}{X} 耶和華是良善正直的，因此，他必教導罪人走正路。 \end{tabularx} \\ \\ \relax
25:9 & \begin{tabularx}{0.7\textwidth}{X} 他要按公平引領謙卑人，將他的道指教他們。 \end{tabularx} \\ \\
[1ex]
\hline
\hline
\end{longtable}
腓立比書 2:1-13
\newline
\begin{longtable}{cl}
\hline
\hline
章節 & 經文 (和合本修訂版)\\
\hline
2:1 & \begin{tabularx}{0.7\textwidth}{X} 所以，在基督裡若有任何勸勉，若有任何愛心的安慰，若有任何聖靈的團契，若有任何慈悲憐憫， \end{tabularx} \\ \\ \relax
2:2 & \begin{tabularx}{0.7\textwidth}{X} 你們就要意志相同，愛心相同，有一致的心思，一致的想法，使我的喜樂得以滿足。 \end{tabularx} \\ \\ \relax
2:3 & \begin{tabularx}{0.7\textwidth}{X} 凡事不可自私自利，不可貪圖虛榮；只要心存謙卑，各人看別人比自己強。 \end{tabularx} \\ \\ \relax
2:4 & \begin{tabularx}{0.7\textwidth}{X} 各人不要單顧自己的事，也要顧別人的事。 \end{tabularx} \\ \\ \relax
2:5 & \begin{tabularx}{0.7\textwidth}{X} 你們當以基督耶穌的心為心： \end{tabularx} \\ \\ \relax
2:6 & \begin{tabularx}{0.7\textwidth}{X} 他本有神的形像，卻不堅持自己與神同等； \end{tabularx} \\ \\ \relax
2:7 & \begin{tabularx}{0.7\textwidth}{X} 反倒虛己，取了奴僕的形像，成為人的樣式；既有人的樣子， \end{tabularx} \\ \\ \relax
2:8 & \begin{tabularx}{0.7\textwidth}{X} 就謙卑自己，存心順服，以至於死，且死在十字架上。 \end{tabularx} \\ \\ \relax
2:9 & \begin{tabularx}{0.7\textwidth}{X} 所以神把他升為至高，又賜給他超乎萬名之上的名， \end{tabularx} \\ \\ \relax
2:10 & \begin{tabularx}{0.7\textwidth}{X} 使一切在天上的、地上的和地底下的，因耶穌的名，眾膝都要跪下， \end{tabularx} \\ \\ \relax
2:11 & \begin{tabularx}{0.7\textwidth}{X} 眾口都要宣認：耶穌基督是主，歸榮耀給父神。 \end{tabularx} \\ \\ \relax
2:12 & \begin{tabularx}{0.7\textwidth}{X} 我親愛的，這樣看來，你們向來是順服的，不但我在你們那裡，就是我現在不在你們那裡的時候更是順服的，就當恐懼戰兢完成你們自己得救的事； \end{tabularx} \\ \\ \relax
2:13 & \begin{tabularx}{0.7\textwidth}{X} 因為是神在你們心裡運行，使你們又立志又實行，為要成就他的美意。 \end{tabularx} \\ \\
[1ex]
\hline
\hline
\end{longtable}
講題:你仗著甚麼權柄?  By What Authority?

經文:馬太福音21:23-32;以西結書18:1-4,25-32;詩篇25:1-9;腓立比書2:1-13

講員:伍渭文牧師

場合:香港中文大學崇基學院禮拜堂主日崇拜

日期:2020年9月27日

引言:今天的題目取自太二十一23,完整一句是「你仗著甚麼權柄做這些事?」略去「做這些事」是避開重覆三年前鄧瑞強博士一篇同名的講章.當天講道經文來自同一福音經課,因為修訂通用經課三年一個循環.題目略去「做這些事」有另一個原因,要指出祭司長和民間長老們關心「誰給你權柄」多於「做這些事情」.「誰給你權柄」-從人而來的權柄；「做這些事」,潔淨聖殿,醫好瞎子,瘸子-從天上來的權柄.

1祭司長祇關心地上的權柄.(毋忘天上來的召命), 太二十一23

「你仗著甚麼權柄做這些事?」

(1)聖殿城的行政長官.耶路撒冷城是當時一個聖殿城特別行政區,祭司長是行政首長,從羅馬帝國得著權柄,向中央派往當地的巡撫(彼拉多)負責.聖殿城有司法權,彼拉多就不想處理耶穌,發還祭司長.「你們自己帶他去,按著你們的律法審問他吧.」(約十八30)他們聽過耶穌在迦伯農的事蹟,也派人明查暗訪作了報告.「從耶路撒冷下來的一些文士說:『他有別西卜附身,他是靠鬼魔的王驅趕鬼魔的.』」(可二22)

羅馬撫巡一定會關注地方興起的宗教運動,尤其關乎猶太人的信仰,特別這運動的領袖宣稱是猶太人的王.現在耶穌就在他們的轄區,祭司長,長老們作為聖殿城管治團隊,要為聖殿城所發生的宗教社會事件問責.「你憑甚麼權柄」是一個政治問題；但作為聖殿負責人,但他們應該關心信仰問題,他們應該問耶穌:你是誰?

(2)聖殿祭祀禱告負責人.祭司長看守管理聖殿,因為這基本的身份,才被羅馬政權委任為聖殿城的行教首長.一如在國內基督教兩會的領導,同時被委任為當地的政協.在大學我們稱concurrent,兼任.「做這些事」,甚麼事呢?耶穌不單潔淨聖殿,也在聖殿醫好瞎子和瘸子,小孩子見到這些神蹟都喊叫:「和散那歸於大衛的子孫!」(太二十一14-16) 這些神蹟說明,耶穌所依仗的權柄,是從天上而來.

為何耶穌要潔淨聖殿?.耶穌用鞭子趕出牛羊,推倒兌換銀錢之人的桌子和賣鴿子之人的凳子,並引用聖經說:「我的殿必稱為禱告的殿(賽五十六7),你們倒使它成為賊窩了.」(太二十一13) 兌換銀錢因為市面流通的錢幣有該撒的肖像,聖殿的奉獻給不應有肖像,十誡明言不能雕刻偶像.獻祭用的祭牲要更求沒有殘疾,聖殿賣出的羊應合符標準.要是在聖殿範圍外進行這些活動沒有問題,但在聖殿範圍內就影響敬拜了,聖殿是人親近神的地方,不容被商業行為干擾.

象徵行動,帶出重要信息.耶穌潔淨聖殿,推翻兌換錢幣的桌子,已成為合法使用適當暴力,對付不公義事情的圖像和經文依據.其實潔淨聖殿是象徵性行動,用來表達重要信息.施洗約翰身穿駱駝毛的衣服,腰束皮帶,吃的是蝗蟲,野蜜,以奇異的生活服飾吸引我們聽他的信息:「天國近了,你們應當悔改!」(太三1-6) 潔淨聖殿的重要信息-基督就是那聖殿.耶穌剛出來傳道不久,在約翰福音二章記載也有一次潔淨聖殿,當時群眾並沒有指責耶穌暴力,他們要了解背後的象徵含義.他們問耶穌:「你既然做這些事,還顯甚麼神蹟(記號)給我們呢?」(約二18)耶穌隨後解釋:「你們拆毀這殿,我三日內要再建立起來.」(約二19) 潔淨聖殿重點在基督論-耶穌是誰,多於教會論-教應如何運作(Christological rather than ecclesiastical).

(3) 輕忽在上主面前的身份.祭司長看見耶穌所做的事情,應該問:你是誰?很可惜他們問:你仗著甚麼權柄作這些事?你為何不按正常程序投訴,讓我們改善,甚至要求成立獨立調查委員會調查.「誰給你權柄?」牽涉行為的合法性,權力的來源,社會的穩定.

祭司長問「你仗著甚麼權柄」有另一個不能宣之於口的原因.他們懼怕耶穌太受百姓歡迎,失去了群眾,要除掉耶穌.祭司長就像陀思妥耶夫斯基的《卡拉馬佐夫兄弟》中的紅衣主教,也是審判異端的《宗教大法官》(Grand Inquisitor).一位陌生人出現在教堂面前的廣場:

「他默默在他們中間走著,流露無限同情的寧靜微笑.他的心燃燒著愛的太陽,他的眼閃耀出光明,智慧和力量的光芒,射到人們的身上,使他們的心裡湧出感激回報的愛.這位陌生人在廣場醫治了許多前來的人.人們嚷著說:“和散那!這是他,這是他自己!”大家反覆地說,”這一定就是他,除了他,不會是別人.”人們何等地激動,喊著,哭著.」

紅衣主教見到這人那麼受歡迎,全城哄動,又嫉妒又懼怕,把他下在監內.晚上進入監倉,拿著燈照照陌生人的臉龐說:「‘真是你?真是你麼?’他沒有得到回答,就急速地接著說,’別出聲,別回答吧.……我完全知道你要說的話……你為什麼到這裡來妨礙我們?……到了明天,我將審判你,把你當作一個最兇惡的異教徒放在火堆上燒死,而今天吻你腳的那些人,明天就,會在我一揮手之下,爭先恐後地跑到你的火堆前面添柴,這你知道嗎?」

是嫉妒和懼怕,使祭司長和長老們逃避問:耶穌你是誰?他們問一個權力秩序,社會穩定的問題.稍後更視耶穌為異端分子,藉巡撫彼拉多的手,除掉耶穌,罪名是自稱是猶太人的王,背叛該撒.

應用:祭司長祇關心聖殿城行政長官的身份,向巡撫問責,忽略了更重要的身份,帶領人敬拜上主的祭司長,聖殿城行政長官祇是兼任(concurrent). 我們不是祭司長,但我們是君尊的祭司,我們的召命,也是我們的正職,是帶領人敬拜上主,地上權柄給我們的職位是兼任(concurrent).職業英文(vocation)原意是召命,我們要帶住天上的召命從事地上不同的責任.我們要帶領人敬拜上主,因為我們看見了真理,我們經歷了上主的憐憫.

2天上來的權柄,使我們看見真理, 腓二1-13

「他本有神的形像……反倒虛己,取了奴僕的形像,成為人的樣式.……」

天上的權柄捅開密封的世界,讓天上的光射進這世界,讓從高處高之處,以基督的心為心看世間事物.從外太空看地球,地面祇是不同顏色,不同圖案的平面,我們都是按上主形像被造.稅吏和娼妓,羅馬的巡撫,祭司長和民間長老,都是上主所愛,同樣需要悔改才能進入上主的國.

(1) 逆轉人類命運的洗禮.祭司長問:「(耶穌)你仗著甚麼權柄作這些事?給你這權柄的是誰呢?」(太二十一23)耶穌回答:「約翰的洗禮是從那裡來的呢?是從天上來的?是從人間來的呢?」耶穌引領他們跳出思考的模式:不要單從人的權柄角度思考,永遠問:你仗著甚麼權柄?給你權柄的是誰?耶穌要他們曉得有一種更宏大,終極,天上來的柄權.

耶穌說:你看看約翰的洗禮,就知我仗著甚麼權柄.約翰在約但河跟不少人施悔改的洗禮,迎接彌賽亞的來臨,但最重要的是他跟耶穌施洗,因這洗禮顯明耶穌是誰.耶穌沒有犯罪,為何要到約但河領洗?是為了我們.彼得前書三20, 21以方舟拯救挪亞一家比喻洗禮:「當時進入方舟,藉著水得救的不多,祇有八個人.這水所表明的洗禮,現在藉著耶穌基督復活也拯救你們.」.耶穌是那方舟,被洪水擊打,裡面的人得免於難,所以教會稱大堂是船(Nave),崇基教堂的設計就像方舟.

「耶穌受了洗,隨即從水裡上來.天忽然為他開了,他就看見神的靈仿彿鴿子降下,落在他身上.從天上有聲音說:這是我的愛子,我所喜悅的.」(太三16, 17)

我所喜悅的是舊約彌賽亞的名號:「看哪,我的僕人-我所扶持所揀選,心裡所喜悅的!」(賽四十二1)聖靈降在耶穌身上,落在約但河水中,逆轉人類的命運.正教會聖洗禮文:「你聖化了約但河的水,從天降下聖靈,把潛藏在水中龍的頭擊碎了.」海在啟示錄代表罪惡,新天新地出現,「有獸(敵基督)從海中上來(十三1)；海要交出其中的死人(二十13)；海也不再有了(二十一1).」愈來覺得有海比喻罪非當貼切,大海深不可測,巨浪捲走生命.但天劍屠龍,從天降下的聖靈,, 21已把潛藏在水中龍的頭擊碎了.耶穌踏進約但河是舊的世界,踏出約但河是另一個新的世界了.

(2) 瞎子看見的新世界.約翰福音九章生來瞎眼的瞎子,被耶穌醫好後最能明白「做這些事」的含意.被耶穌醫好後,他沒有被這彩色繽紛的世界吸引而興奮,也沒有因為自己獲得光明而欣喜若狂,忙著告訴人:「我眼睛開了!我叫眼睛開了!」叫這瞎子震驚欣喜,不是他親眼見到了花草樹木,飛鳥白雲的世界,乃是他意識到彌賽亞已經來到,萬有將被更新.他的心被這認知攝住,他的感動超越了情緒的興奮.對生出來是瞎眼的他,他每天討飯不是就等候這一刻嗎?因為彌賽亞來的時候,瞎子的眼會被打開,他盼望這天已經來臨了!

在整本舊約聖經,真的沒有記載瞎眼的被醫治,醫治瞎眼得見光明是彌賽亞來臨時的記號.「那時,聾子必聽見這書上的話；瞎子的眼必從迷矇中得以看見.」(賽二十九18) 福音書記載耶穌施行的眾多神蹟,以瞎子被醫治得看見最多.他跟質問他的法利賽人說:「從創世以來,未曾聽見有人把生來是瞎子的眼睛開了.這人若不是從神來的,甚麼也不能做.」(約九32, 33)

(3) 地上權柄伏在天上權柄之下.基督被審判時,彼拉多依仗他地上的權柄說:「你不對我說話嗎?你豈不知道的我有權柄釋放你,也有權柄把你釘十字架嗎?」(約十九10)耶穌叫他思考另一種權柄-天上的權柄,所有地上的權柄,都是出自天上的權柄.「若不是從上頭賜給你的,你就毫無權柄辦我.」(約十一11) 耶穌肯定彼拉多的權柄,同時把彼拉多的權柄相對化.君王有崩逝,政黨有輪替,朝代有更迭,文化有興衰,惟有上主是不死的上帝.

(4) 天上的權柄帶來社會制度的革新.首進神國的稅吏和娼妓.祭司長和民間的長老,是以色列的代表,受上主委托,把律法和祭祀禮儀教導百姓,使他們在列邦中成為聖潔的國度,屬神的子民,宣揚上主,等候彌賽亞和神國的降臨.但他們關心地上的權柄,多於上主的召命.無花果樹被咒咀枯萎不單是神蹟(太二十一19),也是對他們的警告.無花果代表以色列,不結果子就被上主撇棄.耶穌說:「稅吏和娼妓倒比你們先進神的國.」(太二十一31) 長期被輕視的低端人口,竟然較權貴更早進入上主的國.

(5)天上的權柄使人人平等.耶穌來了,沒有改變人倫的關系,但改變了如何看待這關係.(a)父與子.舊約經課以西結書十八1-4, 25-32提到俗語「父吃萄葡,兒子牙酸」.兒子在父權至上的時代不過是父親的附屬,但新時代來臨,他們在上主面前同樣是尊貴的,每一個人都要為自己的行為負責任.「你們在以色列中,必不再有用這俗語的因由.看哪,世人都是屬我的；為父的怎樣屬我,為子的也照樣屬我；犯罪的,他必死亡.」(結十八3,4)

(b)兄與弟.耶穌和祭司長及長老對話後,說了一個大小兒子的比喻.大兒子起初拒絕進萄葡園,以後懊悔就去了.小兒子起初承諾進葡萄園但沒有行動,耶穌問:「這兩個兒子是哪一個遵行父命呢?」(太二十一31).他們答:「大兒子.」大兒子是長子,生下來就是,這身份不會改變,也沒有選擇.這是零和遊戲,長子盡取,非長子沒有地位.所以該隱忍受不到弟弟亞伯的祭物蒙上主接納,因嫉妒殺了弟弟.亦因如此,雅各兩次騙取以掃長子的名分,一次乘人之危,以一碗紅豆湯交換長子名份,並起誓不能反悔；另一次假扮以掃欺騙父親臨終的祝福.但在上主國度中,要悔改順服,才成為真正屬靈的長子,蒙受屬靈福氣,承擔使命.

耶穌來了,並沒有把聖殿城改變為人民公社,聖殿城仍然是聖殿城.奴隸制度也沒有革除,但保羅叫亞尼西謀回到奴隸主腓利門身旁.信仰使亞尼西謀「不再是奴僕,乃是高過奴僕,是親愛的弟兄.」(門16) 密封的世界已被捅開了!

耶穌來了,沒有改變洗禮作為一個津渡的禮儀.世俗也有洗禮,是社會向上提升的禮儀,接受水禮得到更高的社會地位.但基督信仰的洗禮,是橫向共融.藉著洗禮,我們成為一體:「並不分猶太人,希臘人,自主的,為奴的,或男或女,因為你們在基督耶穌裡都成為一了.」(加三28)

3天上來的權柄,使我們經歷憐憫, 詩篇經二十五7

「求你不要記念我幼年的罪愆和我的過犯.耶和華啊,求求你因你的恩惠,按你的慈愛記念我.」

權柄與權力不同,權力以利益吸引人,或以懲罰約束人,使在你權力下面的人,聽命於你,但不是由衷的服從你.但當你一朝失去權力,以前聽從你的人,會離開你.有一位大學校長在位時,很多朋友,但退休後朋友稀少,太太知道丈夫感到孤單,叫好友多來探望丈夫.權力來自權位,沒有權位,也沒有權力.人走茶涼,樹倒猢猻散.

聖經權柄一字exoursia有權能的的意思.耶穌講道不像文士,帶有權柄(能力),甚有感染力,聽道的都感受得到.教父特土良(Tertullian)說他驚訝聖經的陌生感(strangeness),像清泉解其屬靈的饑渴.希臘的哲學,多的是縝密的思辨,談論智慧.但聖經是生命之道,談到上主的憐憫和救贖.有些話是當頭棒喝:「你們親近神,神就必親近你們.有罪的人的人哪,要潔淨你們的手!心懷二意的人哪,要清潔你們的心!」(雅四8)

這種改變生命的權能,沒有保鮮期,沒有逾時失效,甚麼時候悔改,就經歷這赦罪的能力,一生伴隨著我們.每個信徒都擁有這屬天的柄柄.「凡接待他的,就是信他名的人,他就賜他們權柄,作神的兒女.」(約一12)

詩篇二十五7:「求你不要記念我幼年的罪愆和我的過犯.耶和華啊,求求你因你的恩惠,按你的慈愛記念我.」為何特別題及幼年的罪愆呢?幼年的罪愆有何特別呢?是年幼無知時犯的罪?是年幼反叛上主犯的罪?

我想,詩人在晚年還提及幼年的罪愆,是想起多年前犯下的深重罪惡,帶來的罪疚和控訴,至今仍揮之不去,撒旦就是「在我們神面前畫夜控訴我們弟兄的.」(啟十二10下) 耶穌來了,並不能回到過去,但可以抹去過去的眼淚,罪愆,和過犯.我們不再活在永遠的控訴中,每天都是新的一天.因他的本性是憐憫,而憐憫是向審判跨勝.

路德稱這勝過罪惡控訴的權柄,比肩當時一位君王.基督徒是最自由的君王,連罪都不能綑縛我們,便是最自由的君王了；路德的話使我想起聖詩「我寧願有耶穌」:

我寧願有耶穌,勝得金錢,我寧屬耶穌,勝得財富無邊；

我寧願有耶穌,勝得華宇,願主釘痕手,引導我前途.

我寧願有耶穌,勝得稱頌,我寧忠於主,成全祂的事工；

我寧願有耶穌,勝得名聲,願向主聖名永赤膽忠誠.

恩主比百合花更加美艷,遠比蜂房蜜更加可口甘甜；

副歌:勝過做君王,雖統治萬方,卻仍受罪惡捆綁；

我寧願有耶穌,勝得世界榮華,富貴,聲望.

基督徒是最自由的君王,同時也是最順從的僕人.得著這天上來的權柄,沒有取消要服從地上的權柄的責任.地上的權柄帶來諸般的身份,是我們顯揚愛心的地方.

耶穌降生,特意選擇巴勒斯坦進行人口調查,首報名上冊的歷史時刻,他順服羅馬帝國給猶太人的新稅制.他順服人的制度,納稅給政府.教會年聖誕後的主日稱聖家主日,記念約瑟和馬利亞攜同幼兒耶穌逃難埃及,逃避希律王的刀劍.為了滅絕新生王,希律把兩歲以下的嬰孩殺掉.教會記念這些殉道的嬰孩,十二月二十八日為嬰孩被殺日(Feast of Holy Innocents),為世上的孩童代禱:難民兒童,棄嬰,虐兒,童工,雛妓,兒童柺帶販賣.

結論

1)每人都有兩種權柄在他身上,兩種權柄都屬出自上主,各有不同的功能.地上的權柄給我們不同的角色,天上的權柄給我們不變的召命,成就上主隱密的計劃.

2) 這天上的權柄,使我們看見真理:每人在上主面前都是平等的,為父的,為子的,為兄的,為弟的,為奴的,自主的,男的,女的,都在上主面前極其尊貴.天上的權柄,像一襲嶄新的靈風,吹進了這密封的世界,帶來社會的革新.

3) 凡接受基督的,都擁有這屬天的權柄,使我們脫離罪的控訴,這權柄就是上主的憐憫和慈愛,終生伴隨我們.

香港通過國安法後,突然間這城市好像進入密封的世界；香港製造已變成中國製造.新冠狀病毒瘟疫席捲全球,舉世正焦慮等待疫苗面世,整個世界活在禁足與經濟復甦的拉力中.今天經文給我們安慰,有一更宏大的權柄,支撐著這世界,這權柄有真理,說明世界的去向；這權柄有憐憫,帶給我們上主的慈愛.

這世界和其上的權柄都會過去.聖殿城沒有了,但聖殿的敬拜永不止息.在地上我們有人倫的關系,但在新天新地,也不嫁不娶.我們的身體將隨年日朽壞,但藉身體為上主作的工作存到永遠.我們在地上,有不同的角色,伏在不同權柄之下,也有人伏在我們權柄之下.我們要順服天上來的權柄,尋求上主給我們智慧,學習如何使用權柄,如何順服權柄,榮神益人.

問:「你仗著甚麼權柄做這些事?」

答:「我仗著上主兒女的權柄作這些事,因我們看見了真理,因我們經歷了憐憫.」



\newpage

\section{我們看得到這希奇嗎}
\label{sec:520}
\textbf{高國雄牧師 - 我們看得到這希奇嗎?}
\newline
\newline
\hyperref[sec:519]{< < < PREV SERMON < < <}
~
\hyperref[sec:toc]{[返目錄]}
~
\hyperref[sec:521]{> > > NEXT SERMON > > >}


馬太福音 21:33-46
\newline
\begin{longtable}{cl}
\hline
\hline
章節 & 經文 (和合本修訂版)\\
\hline
21:33 & \begin{tabularx}{0.7\textwidth}{X} 「你們再聽一個比喻：有一個家的主人開墾了一個葡萄園，四周圍上籬笆，裡面挖了一個醡酒池，蓋了一座守望樓，租給園戶，就出外遠行去了。 \end{tabularx} \\ \\ \relax
21:34 & \begin{tabularx}{0.7\textwidth}{X} 收果子的時候快到了，他打發僕人到園戶那裡去收果子。 \end{tabularx} \\ \\ \relax
21:35 & \begin{tabularx}{0.7\textwidth}{X} 園戶拿住僕人，打了一個，殺了一個，用石頭打死了一個。 \end{tabularx} \\ \\ \relax
21:36 & \begin{tabularx}{0.7\textwidth}{X} 主人又打發別的僕人去，比先前更多；園戶還是照樣對待他們。 \end{tabularx} \\ \\ \relax
21:37 & \begin{tabularx}{0.7\textwidth}{X} 最後他打發自己的兒子到他們那裡去，說：『他們會尊敬我的兒子。』 \end{tabularx} \\ \\ \relax
21:38 & \begin{tabularx}{0.7\textwidth}{X} 可是，園戶看見他兒子，彼此說：『這是承受產業的。來，我們殺了他，佔他的產業！』 \end{tabularx} \\ \\ \relax
21:39 & \begin{tabularx}{0.7\textwidth}{X} 於是他們拿住他，把他扔出葡萄園外，殺了。 \end{tabularx} \\ \\ \relax
21:40 & \begin{tabularx}{0.7\textwidth}{X} 葡萄園的主人來的時候，要怎樣處置那些園戶呢？」 \end{tabularx} \\ \\ \relax
21:41 & \begin{tabularx}{0.7\textwidth}{X} 他們說：「要狠狠地除滅那些惡人，將葡萄園轉租給那些按時候交果子的園戶。」 \end{tabularx} \\ \\ \relax
21:42 & \begin{tabularx}{0.7\textwidth}{X} 耶穌對他們說：「『匠人所丟棄的石頭已作了房角的頭塊石頭。這是主所做的，在我們眼中看為奇妙。』這段經文你們從來沒有念過嗎？ \end{tabularx} \\ \\ \relax
21:43 & \begin{tabularx}{0.7\textwidth}{X} 所以我告訴你們，神的國必從你們奪去，賜給那能結果子的民。 \end{tabularx} \\ \\ \relax
21:44 & \begin{tabularx}{0.7\textwidth}{X} 誰跌在這石頭上，一定會跌得粉碎；這石頭掉在誰的身上，就要把誰壓得稀爛。」 \end{tabularx} \\ \\ \relax
21:45 & \begin{tabularx}{0.7\textwidth}{X} 祭司長和法利賽人聽見他的比喻，就看出他是指著他們說的。 \end{tabularx} \\ \\ \relax
21:46 & \begin{tabularx}{0.7\textwidth}{X} 他們想要捉拿他，但是懼怕眾人，因為眾人認為他是先知。 \end{tabularx} \\ \\
[1ex]
\hline
\hline
\end{longtable}
講題:我們看得到這希奇嗎?It is amazing in our eyes’?

講員:高國雄牧師

經文:馬太福音21:33-46

場合:香港中文大學崇基學院禮拜堂主日崇拜

日期:2020年10月5日

各位弟兄姊妹,在你們的所遇見的事情當中,有什麼事令你覺得希奇的呢?

在日新月異,幻變不斷的今天,可能很多本來應該值得希奇的事情,我們也慢慢習慣接受了,而不再有感覺,今天想與大家集中思考福音書的經課,有關葡萄園家主與園戶的比喻,請留意這個比喻在三卷符類福音書中皆有記述,這耶穌基督所說我們熟悉的比喻,今天又有沒有將會令你希奇的地方呢?

讓我們先再一次閱讀今天馬太福音的經文:

馬太福音21:33-46

33「你們再聽一個比喻:有個家主栽了一個葡萄園,周圍圈上籬笆,裡面挖了一個壓酒池,蓋了一座樓,租給園戶,就往外國去了. 34收果子的時候近了,就打發僕人到園戶那裡去收果子. 35園戶拿住僕人,打了一個,殺了一個,用石頭打死一個. 36主人又打發別的僕人去,比先前更多；園戶還是照樣待他們. 37後來打發他的兒子到他們那裡去,意思說:『他們必尊敬我的兒子.』 38不料,園戶看見他兒子,就彼此說:『這是承受產業的.來吧,我們殺他,佔他的產業!』 39他們就拿住他,推出葡萄園外,殺了. 40園主來的時候要怎樣處治這些園戶呢?」 41他們說:「要下毒手除滅那些惡人,將葡萄園另租給那按著時候交果子的園戶.」 42耶穌說:「經上寫著:匠人所棄的石頭已作了房角的頭塊石頭.這是主所做的,我們眼中看為希奇.這經你們沒有念過嗎? 43所以我告訴你們,上帝的國必從你們奪去,賜給那能結果子的百姓. 44誰掉在這石頭上,必要跌碎；這石頭掉在誰的身上,就要把誰砸得稀爛.」 45祭司長和法利賽人聽見他的比喻,就看出他是指著他們說的. 46他們想要捉拿他,只是怕眾人,因為眾人以他為先知.

故事提到一個家主,與及一個園戶；用我們比較熟悉的表達,就係一個業主,與一個租客；這位業主是一位非常好的業主,他栽了葡萄園,周圍圈上籬笆,裡面挖了一個壓酒池,蓋了一座樓,就好像一位業主把房子做了最美的裝修,購置最好的傢俬電器,之後把屋租出去,就去外國了,弟兄姊妹你覺得希奇嗎?有沒有遇到過這麼好的業主呢?或許你會認為:為租客預備最好的裝修與及最好的家電,為的是有更好的租金回報,是聰明的業主應該做的事,但細心思考,願意為租客獻上最好的不是理所當然,我當然希望如果我們當中有作業主的要善待租客,但這位葡萄園的主人─我們的主上帝,在創世時已經開始為我們預備了最美好,安心地交在我們手中管理,我們有沒有因為有這位如此愛我們的上帝而希奇呢?還是認為是理所當然的呢?

故事繼續,當時候到了,收果子的時候近了,主人打發僕人到葡萄園收果子,就正如我們大部份人每月到期要交租一樣,但當時居然發生了這事；

太21:35「園戶拿住僕人,打了一個,殺了一個,用石頭打死一個.」

相信我們沒有辦法想像,租客竟然打傷,甚至殺害業主派來收租的人,不但如此,

太21:36「主人又打發別的僕人去,比先前更多；園戶還是照樣待他們.後來打發他的兒子到他們那裡去,意思說:『他們必尊敬我的兒子.』 不料,園戶看見他兒子,就彼此說:『這是承受產業的.來吧,我們殺他,佔他的產業!』他們就拿住他,推出葡萄園外,殺了.」

很明顯作這惡不可能只是一時的衝動,而是有預謀的,為要謀財害命,所以和合本的編者分段時把這段聖經在三部符類福音中起了同樣的「凶惡園戶的比喻」為題,相當準確,聽到這凶惡園戶的事情,你覺得希奇嗎?

二年多前因為一套電視劇為香港社會帶來一句新金句:「誰大誰惡誰正確!」[1]當然可能大家與我一樣,完全沒有看過這套電視劇,但也應該在不同的媒體中見過這句金句,相信很多人或許都會認同,當然我們應該不會接受大奸大惡,殺人放火等行為,但因為在日常生活上卻不難碰上…

排隊時插隊的人仍然「兇神惡殺」,甚至謾罵回應；投訴者不論有理或沒有理也先惡形及惡言示人；不幸遇上交通意外,司機下車第一時間必先大罵對方的不是等…

讓我想起80年代中期有一首流行歌,其中歌詞提到:

「香港地不管工作與耍樂,呢一招高招梗要記得學

惡會有收穫 惡永無鋇鑊,惡做錯事也可側側膊

總之爭取優勢咪示弱,往往世事系憑邊便惡

成日惡鬥惡,成日斗有冇pok,無論爭取緊D 乜野都梗要惡

惡鬥惡,成事要靠夠pok

無論心地點樣好法,永遠個樣都梗要裝惡

冇道理亦有機可剝削,有道理亦會冇得反駁

世事確係要睇邊便惡」[2]

因為我很喜歡笑,常常都面帶笑容,曾經有不少朋友,常常有善地提醒我,不可以這樣,只要板起臉,就會顯出更有威嚴,身邊的人就會忌你三分,做事起來會更得心應手；甚至有朋友在幾個月前很認真地提我,在崇基學院將會面對很多不同的人,太友善會有機會被欺負的,要我記住:「無論心地點樣好法,永遠個樣都梗要裝惡.」我真心多謝朋友的提醒,但我仍然認為用真誠待人最重要,可能真的在世人眼中會被看為是蝕抵的事,但我們作為信徒,不就是應該更以上主的事為念嗎?聖經教導要:各人看別人比自己強,不要單顧自己的事,要愛人甚至仇敵,要學習耶穌基督犧牲的精神等等；難道不是已經非常清楚地告知我們「自身」的利益不應該是我們所追求最重要的事嗎?有所蝕低不應該係問題,問題應該係我們一切的思想及行為,有沒有彰顯出上帝的榮美,見證祂的實在.

如果我們今天見到身邊的惡人作惡不再感到希奇,詫異,甚至認為是理所當然,更可能我們已經成為一份子,真的認罪悔改,求主憐憫及赦免.

我們回到經文,耶穌說完這比喻後,就問眾人:

太21:40-41

園主來的時候要怎樣處治這些園戶呢?

他們說:「要下毒手除滅那些惡人,將葡萄園另租給那按著時候交果子的園戶.」

耶穌基督的比喻不是只說故事而已,不單只是要我們聽一個教訓那麼簡單,正如今天在這崇拜中,上主不只是要你聽一些道理,單單學習理解一些基督教的價值觀；而是要你進入上帝的話語中,認真思考應如何去回應,並以實際行動表現出上帝的信實與及你的信心.

當耶穌身邊那群認真聽這比喻的人,投入地回應表示主人應該「以牙還牙」,除掉惡人,另覓忠心,願意做應做的事的新園戶後,我們先看看祭司長及法利賽人的回應:

太21:45-46

「祭司長和法利賽人聽見他的比喻,就看出他是指著他們說的.他們想要捉拿他,只是怕眾人,因為眾人以他為先知.」

這是另一個讓我希奇的地方,其實祭司長和法利賽人怎樣「看出」耶穌是指著他們說的呢?他們曾經殺人?他們做過相類似的惡事?祭司長和法利賽人都係宗教的領袖,應該是當時社經地位非常高,道德操守也理應被要求高的一群人,聖經中亦常常描述他們是以自為優秀,比別人高一等的人；但聖經的描述,當他們一聽見耶穌的比喻,就清楚知道他們是耶穌所指那位不願立租,謀財害命的對象,馬上的對號入座反映出當時這一群領袖的心態,可能也是我們今天的心態；我們知道應該如何過我們的日子,但我們不願意,不接受,因為我只想為「自己」而活；園戶為什麼不願意立租?為什麼要把家主所派來收租的人殺害?直到殺害家主兒子時就清楚表示其原因:

太21:38

園戶看見他兒子,就彼此說:『這是承受產業的.來吧,我們殺他,佔他的產業!』

原來園戶認為自己應該要成為家主,葡萄園應該是屬於我的,我的權益比所有都大,當然包括道德價值；因此只要爭取,維護我的個人權益,我做什麼也沒有問題…

相信因為當時這一群宗教領袖心裡根本就有這個想法,所以就會馬上對號入座,從聖經中繼續描述他們接著的想法(捉拿耶穌)及往後的行動(把耶穌釘在十字架上),就更印證耶穌基督看穿了他們的心,亦看穿了我們當中很多人的心…

再次提醒大家,我們不是這世界的主人,從來沒有擁有任何的東西,我們的主上帝已經為我們預備了一切的美好,讓我們去享受,去管理,我們要很清楚自己的位置,作為管家的我們,所要求的就是忠心及良善.

讓我們再次回到經文,當群眾回答耶穌對應該如何處置這凶惡的園戶之後,耶穌又如何回應呢?

太21:42-44

耶穌說:「經上寫著:匠人所棄的石頭已作了房角的頭塊石頭.這是主所做的,我們眼中看為希奇.這經你們沒有念過嗎? 所以我告訴你們,上帝的國必從你們奪去,賜給那能結果子的百姓.誰掉在這石頭上,必要跌碎；這石頭掉在誰的身上,就要把誰砸得稀爛.」

經文中真正讓我們希奇的事出現了,對於匠人來說,他應該有非常足夠的經驗及能力,選取最合用的器材,一個建造房屋的匠人,理應清楚知道那塊石頭最合用,如果我們當中有熱愛陶瓷,雕刻等藝術創作的弟兄姊妹應該更容易體會選擇材料的重要,匠人所棄用的就等如在世人眼中沒有價值的,甚至是會做成障礙的,用一般「正常的邏輯」來說是要放棄,要除掉的；就正正不是在這些祭司長和法利賽人眼中要除掉的那位挑戰他們價值觀的耶穌嗎?

上帝讓耶穌基督成為了教會最重要的那塊房角石頭,當時建造房屋最重要的就是選定房角石並放置在最合適的位置；今天上帝就把耶穌基督放在這個位置上.

經文令我們希奇的地方,讓我們認識到上帝的愛及寬容是我們想像不到的,就正如凶惡園戶殺害主人所派的僕人後,主人仍然願意再派人去,讓園戶有改過的機會；園戶竟然仍然作惡,主人不但沒有馬上作出刑罰,更差派出自己的兒子,再給予園戶回轉的機會；其實在今天,上主同樣地不斷向我們伸出憐憫的手,等待著我們這些惡人回轉,明白祂的心意,接受祂的愛,讓我想起二段我最愛的經文:

約翰福音3:16

上帝愛世人,甚至將祂的獨生子賜給他們,叫一切信祂的,不致滅亡,反得永生.

羅馬書5:8

惟有基督在我們還作罪人的時候為我們死,上帝的愛就在此向我們顯明了.

各位朋友,各位弟兄姊妹,如果你仍然未接受,又或許已經忘卻了主的愛,請把握機會,盡快的回轉；因為耶穌基督非常清楚的表示,當主人差派兒子來到,人仍然以惡對待,甚至最後把祂殺害,不願意接受主人的美意,「上帝的國必從你們奪去」,被「跌碎」,被「砸鬼稀爛」…上帝同樣是公義主,不以有罪的為無罪,到了審判的日子,我們每一個人也要到上帝的台前,到底我們的終局所面對的是甚麼呢?完全取決於今天我們的選擇.

[1] 電視廣播有限公司(TVB無線電視)2018年12月播出劇集《大帥哥》主角狄奇的對白

[2] 1986年林子祥【最愛】專輯內第四首歌《惡鬥惡》,作曲:A. Willis,填詞:潘源良



\newpage

\section{要飯聚不犯罪}
\label{sec:521}
\textbf{林豪恩先生 - 要飯聚,不犯罪}
\newline
\newline
\hyperref[sec:520]{< < < PREV SERMON < < <}
~
\hyperref[sec:toc]{[返目錄]}
~
\hyperref[sec:522]{> > > NEXT SERMON > > >}


以賽亞書 25:1-9
\newline
\begin{longtable}{cl}
\hline
\hline
章節 & 經文 (和合本修訂版)\\
\hline
25:1 & \begin{tabularx}{0.7\textwidth}{X} 耶和華啊，你是我的神，我要尊崇你，稱頌你的名。因為你以信實忠信行遠古所定奇妙的事。 \end{tabularx} \\ \\ \relax
25:2 & \begin{tabularx}{0.7\textwidth}{X} 你使城市變為廢墟，使堅固的城荒涼，使外邦人的城堡不再為城，永遠不再重建。 \end{tabularx} \\ \\ \relax
25:3 & \begin{tabularx}{0.7\textwidth}{X} 所以，強大的民必尊敬你，殘暴之國的城必敬畏你。 \end{tabularx} \\ \\ \relax
25:4 & \begin{tabularx}{0.7\textwidth}{X} 因為你是貧寒人的保障，貧窮人急難中的保障，暴風雨之避難所，炎熱地之陰涼處。當殘暴者盛氣凌人的時候，如暴風直吹牆壁， \end{tabularx} \\ \\ \relax
25:5 & \begin{tabularx}{0.7\textwidth}{X} 如乾旱地的熱氣，你要制止外邦人的喧嚷，殘暴者的歌要停止，好像熱氣因雲的陰影而消失。 \end{tabularx} \\ \\ \relax
25:6 & \begin{tabularx}{0.7\textwidth}{X} 在這山上，萬軍之耶和華必為萬民擺設宴席，有肥甘與美酒，就是滿有骨髓的肥甘與精釀的美酒。 \end{tabularx} \\ \\ \relax
25:7 & \begin{tabularx}{0.7\textwidth}{X} 在這山上，他必吞滅那纏裹萬民的面紗和那遮蓋列國的遮蔽物。 \end{tabularx} \\ \\ \relax
25:8 & \begin{tabularx}{0.7\textwidth}{X} 他已吞滅死亡直到永遠。主耶和華必擦乾各人臉上的眼淚，在全地除去他百姓的羞辱；這是耶和華說的。 \end{tabularx} \\ \\ \relax
25:9 & \begin{tabularx}{0.7\textwidth}{X} 到那日，人必說：「看哪，這是我們的神，我們向來等候他，他必拯救我們。這是耶和華，我們向來等候他，我們必因他的救恩歡喜快樂。」 \end{tabularx} \\ \\
[1ex]
\hline
\hline
\end{longtable}
詩篇 23:1-6
\newline
\begin{longtable}{cl}
\hline
\hline
章節 & 經文 (和合本修訂版)\\
\hline
23:1 & \begin{tabularx}{0.7\textwidth}{X} 耶和華是我的牧者，我必不致缺乏。 \end{tabularx} \\ \\ \relax
23:2 & \begin{tabularx}{0.7\textwidth}{X} 他使我躺臥在青草地上，領我在可安歇的水邊。 \end{tabularx} \\ \\ \relax
23:3 & \begin{tabularx}{0.7\textwidth}{X} 他使我的靈魂甦醒，為自己的名引導我走義路。 \end{tabularx} \\ \\ \relax
23:4 & \begin{tabularx}{0.7\textwidth}{X} 我雖然行過死蔭的幽谷，也不怕遭害，因為你與我同在；你的杖、你的竿，都安慰我。 \end{tabularx} \\ \\ \relax
23:5 & \begin{tabularx}{0.7\textwidth}{X} 在我敵人面前，你為我擺設筵席；你用油膏了我的頭，使我的福杯滿溢。 \end{tabularx} \\ \\ \relax
23:6 & \begin{tabularx}{0.7\textwidth}{X} 我一生一世必有恩惠慈愛隨著我；我且要住在耶和華的殿中，直到永遠。 \end{tabularx} \\ \\
[1ex]
\hline
\hline
\end{longtable}
腓立比書 4:1-9
\newline
\begin{longtable}{cl}
\hline
\hline
章節 & 經文 (和合本修訂版)\\
\hline
4:1 & \begin{tabularx}{0.7\textwidth}{X} 我所親愛、所想念的弟兄們，你們就是我的喜樂，我的冠冕。我親愛的，你們應當靠主站立得穩。 \end{tabularx} \\ \\ \relax
4:2 & \begin{tabularx}{0.7\textwidth}{X} 我勸友阿蝶和循都基要在主裡同心。 \end{tabularx} \\ \\ \relax
4:3 & \begin{tabularx}{0.7\textwidth}{X} 我也求你這真實同負一軛的，要幫助這兩個女人，因為她們在福音上曾與我、革利免和我其餘的同工一同勞苦，他們的名字都在生命冊上。 \end{tabularx} \\ \\ \relax
4:4 & \begin{tabularx}{0.7\textwidth}{X} 你們要靠主常常喜樂。我再說，你們要喜樂。 \end{tabularx} \\ \\ \relax
4:5 & \begin{tabularx}{0.7\textwidth}{X} 要讓眾人知道你們謙讓的心。主已經近了。 \end{tabularx} \\ \\ \relax
4:6 & \begin{tabularx}{0.7\textwidth}{X} 應當一無掛慮，只要凡事藉著禱告、祈求和感謝，將你們所要的告訴神。 \end{tabularx} \\ \\ \relax
4:7 & \begin{tabularx}{0.7\textwidth}{X} 神所賜那超越人所能了解的平安，必在基督耶穌裡，保守你們的心懷意念。 \end{tabularx} \\ \\ \relax
4:8 & \begin{tabularx}{0.7\textwidth}{X} 末了，弟兄們，凡是真實的、凡是可敬的、凡是公義的、凡是清潔的、凡是可愛的、凡是有美名的，若有甚麼德行，若有甚麼稱讚，你們都要留意。 \end{tabularx} \\ \\ \relax
4:9 & \begin{tabularx}{0.7\textwidth}{X} 你們從我所學習的，所領受的，所聽見的，所看見的事，你們都要繼續去做，賜平安的神就必與你們同在。 \end{tabularx} \\ \\
[1ex]
\hline
\hline
\end{longtable}
馬太福音 22:1-14
\newline
\begin{longtable}{cl}
\hline
\hline
章節 & 經文 (和合本修訂版)\\
\hline
22:1 & \begin{tabularx}{0.7\textwidth}{X} 耶穌又用比喻對他們說： \end{tabularx} \\ \\ \relax
22:2 & \begin{tabularx}{0.7\textwidth}{X} 「天國好比一個王為他兒子擺設娶親的宴席。 \end{tabularx} \\ \\ \relax
22:3 & \begin{tabularx}{0.7\textwidth}{X} 他打發僕人去，請那些被邀的人來赴宴，他們卻不肯來。 \end{tabularx} \\ \\ \relax
22:4 & \begin{tabularx}{0.7\textwidth}{X} 王又打發別的僕人，說：『你們去告訴那被邀的人，我的宴席已經預備好了，牛和肥畜已經宰了，各樣都齊備，請你們來赴宴。』 \end{tabularx} \\ \\ \relax
22:5 & \begin{tabularx}{0.7\textwidth}{X} 那些人不理就走了，一個到自己田裡去，一個做買賣去。 \end{tabularx} \\ \\ \relax
22:6 & \begin{tabularx}{0.7\textwidth}{X} 其餘的抓住僕人，凌辱他們，把他們殺了。 \end{tabularx} \\ \\ \relax
22:7 & \begin{tabularx}{0.7\textwidth}{X} 王就大怒，發兵除滅那些兇手，燒燬他們的城。 \end{tabularx} \\ \\ \relax
22:8 & \begin{tabularx}{0.7\textwidth}{X} 於是王對僕人說：『喜宴已經齊備，只是所邀的人不配。 \end{tabularx} \\ \\ \relax
22:9 & \begin{tabularx}{0.7\textwidth}{X} 所以你們要往岔路口上去，凡遇見的，都邀來赴宴。』 \end{tabularx} \\ \\ \relax
22:10 & \begin{tabularx}{0.7\textwidth}{X} 那些僕人就出去，到大路上，凡遇見的，不論善惡都招聚了來，宴席上就坐滿了客人。 \end{tabularx} \\ \\ \relax
22:11 & \begin{tabularx}{0.7\textwidth}{X} 王進來見賓客，看到那裡有一個沒有穿禮服的， \end{tabularx} \\ \\ \relax
22:12 & \begin{tabularx}{0.7\textwidth}{X} 就對他說：『朋友，你到這裡來怎麼不穿禮服呢？』那人無言可答。 \end{tabularx} \\ \\ \relax
22:13 & \begin{tabularx}{0.7\textwidth}{X} 於是王對侍從說：『捆起他的手腳，把他扔在外邊的黑暗裡；在那裡他要哀哭切齒了。』 \end{tabularx} \\ \\ \relax
22:14 & \begin{tabularx}{0.7\textwidth}{X} 因為被召的人多，選上的人少。」 \end{tabularx} \\ \\
[1ex]
\hline
\hline
\end{longtable}
講題:要飯聚,不犯罪 Let’s Eat Together

經文:以賽亞書25:1-9;詩篇23:1-6;腓立比書4:1-9;馬太福音22:1-14

講員:林豪恩先生

場合:香港中文大學崇基學院禮拜堂主日崇拜

日期:2020年10月11日

一．引言

曾幾何時,崇拜結束後,大家到地下休息室「飲杯茶,食件餅」,閒談幾句.尚有空的就移師禮拜堂對面的教職員聯誼會,大台小台,或回應講道,或暢談信仰,或笑談風月,各自精彩.本地人俗稱這樣的一起吃飯為「飯聚」,而這種「飯聚」一直都不「犯罪」的.然而,在疫情之下,如此的「飯聚」蒙上了「犯罪」的風險.某經典日劇主題曲曾唱:「誰會珍惜,當你還擁有?當要逝去,總想挽留.」何時才能夠一起吃飯?此時此地,不單單是家庭和朋友的關注,更是社會議題,甚至是全球共同面對的難題.有沒有想過,一起吃飯,其實也是信仰的議題?今天的幾段經課,都與「飯聚」有關,經文的用詞是「筵席」.讓我們再細閱經文,思考「筵席」的信仰意義.

二．釋經及應用

1．以賽亞書25章:上帝的筵席與巴比倫的筵席

第一場筵席出現在以賽亞書25:6,雖然是一場筵席,其實隱含著兩場筵席.

「在這山上,萬軍之耶和華必為萬民用肥甘設擺筵席,用陳酒和滿髓的肥甘,並澄清的陳酒,設擺筵席.」

你有興趣參加這筵席嗎?或許你有興趣,不過有些人可能會嗤之以鼻,因為他們正在皇宮中享受盛宴,何必長途跋涉去山上「野餐」呢?在經文的歷史時期,在大河的那一邊,就是在巴比倫的皇宮,有一場宮廷盛筵正在進行中,記載在但以理書第五章:

5:1伯沙撒王為他的一千大臣擺設盛筵,與這一千人飲酒.

5:2伯沙撒在歡飲之間,吩咐人將他父尼布甲尼撒從耶路撒冷聖殿所擄掠的金銀器皿拿來,好使王與大臣,王后,妃嬪用這器皿飲酒.

5:3於是他們把聖殿,就是耶路撒冷上帝殿中所擄掠的金器皿拿來,王和大臣,王后,妃嬪就用這器皿飲酒.

巴比倫皇宮的盛宴與以賽亞書25章有甚麼關係呢?當然有關係.還記得伯沙撒王手中拿著的金銀器皿是從可而來的嗎?但以理書明言,是從耶路撒冷聖殿所擄掠而來的.以賽亞書25章的內容正是與耶路撒冷被巴比倫攻陷有關的.

3所以,剛強的民必榮耀你；

強暴之國的城必敬畏你.

4因為當強暴人催逼人的時候,

如同暴風直吹牆壁,

巴比倫這「強暴之國」如狂風掃落葉在古代近東一帶所向披靡,猶大這些小國慘遭蹂躪,連聖殿中的金銀器皿都被搶掠一空.先知用「如同暴風直吹」來形容巴比倫軍隊的攻勢.暴風過後,猶大一片狼藉,苟延殘喘,還講甚麼筵席呢?筵席,只與勝利者相配,被打敗者怎可能有筵席呢?筵席,應該在巴比倫的皇宮舉行,怎可能在頹垣敗瓦的錫安山舉杯呢?以賽亞先知卻在如此的狀況之下宣佈筵席,並說這是「奇妙的事」.

1耶和華啊,你是我的上帝；

我要尊崇你,我要稱讚你的名.

因為你以忠信誠實行過奇妙的事,

成就你古時所定的.

「奇妙的事」是甚麼事呢?就是「在這山上,萬軍之耶和華必為萬民用肥甘設擺筵席」.如此「奇妙的事」怎可能呢?以賽亞先知指出這是可能的,因為「奇妙的事」尚包括,第一,上帝會回應「強暴之國」,毀人城者其城必被毀:

2你使城變為亂堆,

使堅固城變為荒場,

使外邦人宮殿的城不再為城,

永遠不再建造.

第二,上帝沒有忘記苦難中的人:

4因為當強暴人催逼人的時候,

如同暴風直吹牆壁,你就作貧窮人的保障,

作困乏人急難中的保障,

作躲暴風之處,

作避炎熱的陰涼.

因此,耶和華在這山上設擺筵席是可能的.上帝的筵席有甚麼食物呢?有「滿髓的肥甘」和「澄清的陳酒」.不過,食物不是最重要,最重要的,也是與巴比倫皇宮的筵席的分別所在,7-8節如此說:

7他又必在這山上除滅遮蓋萬民之物和遮蔽萬國蒙臉的帕子.

8他已經吞滅死亡直到永遠.主耶和華必擦去各人臉上的眼淚,又除掉普天下他百姓的羞辱,因為這是耶和華說的.

在上帝的筵席,「耶和華必擦去各人臉上的眼淚」,「除掉普天下他百姓的羞辱」.上帝的筵席,是貧窮人的保障,苦難者的安慰,使人性尊嚴得到恢復.巴比倫皇宮的筵席是甚麼呢?他們所吃所喝的,是戰敗國心不甘情不願地進貢的；他們所用的器皿,是靠強暴從戰敗國的聖殿擄掠回來的；他們的狂歡,建築在別人的眼淚之上；他們的狂傲,令別人感到被羞辱.

「強暴之國」巴比倫被上帝審判了,巴比倫消失了.然而,巴比倫的筵席從沒消失.在歷史上,還有多少「強暴之國」用各種方法掠奪別國的來讓自己狂歡暢飲呢?還有多少有力量的人掠奪別人的來建造自己的酒池肉林呢?16世紀至19世紀長達約400年的黑奴貿易中,學者估計非洲人口損失達到三千萬,其中一半被賣到歐美做奴隸,另一半在航運中死去.就算今時今日,仍有生產糧食的人民自己得不到飽足,其實糧食去了哪裡呢?種植可可豆的人,從來都吃不起朱古力,朱古力去了哪裡呢?巴比倫的筵席雖然被興高彩烈,人人趨之若鶩,但這卻是罪惡的筵席,巴比倫的「飯聚」,其實是真正的「犯罪」.

上帝的筵席不是巴比倫的筵席,在上帝的筵席中,「各人臉上的眼淚」被擦去,「普天下他百姓的羞辱」被除掉.因此,以賽亞先知說:

25:9「看哪,這是我們的上帝；我們素來等候他,他必拯救我們.這是耶和華,我們素來等候他,我們必因他的救恩歡喜快樂.」

上帝的筵席是值得參加的,值得等待的.

2．詩篇23:在敵人面前的筵席

第二場筵席在詩篇23篇出現:

5在我敵人面前,你為我擺設筵席；

這個筵席特別的地方,是擺設「在我敵人面前」.在敵人面前擺設的筵席,到底是甚麼筵席呢?解經者有不同的解釋,今次我們參考建道神學院的教授高銘謙博士的見解.

高教授把把「擺設筵席」的傳統與以色列民四十年在曠野的生活拉上關係.他指出神為以色列民降下嗎哪,就是在沒有可能生存的地方為以色列民擺設筵席.利未記十六章描述曠野是阿撒瀉勒的地方,很多學者認為這是一種曠野的魔鬼,而當以色列民在曠野中生存,正正就是在敵人的領域當中生存,這是一種「在我敵人面前」的情況,而嗎哪便是神所擺設的筵席.這解釋得到另一篇詩篇驗證,這詩篇指出以色列民在曠野中埋怨,質疑上帝是否真的能夠擺設筵席:「他們心中試探神,隨自己所欲的求食物,並且妄論上帝說:上帝在曠野豈能擺設筵席嗎?」(詩七十八18-19) 因此,有證據顯示「擺設筵席」與曠野的經驗連在一起.「在我敵人面前,你為我擺設筵席」代表上帝能夠在曠野的敵人領域中做出不可能的神蹟.

詩篇23:4透露,詩人「行過死蔭的幽谷」.由此可見,信靠上帝的人也會陷在「死蔭幽谷」之中.在「死蔭幽谷」之中,最難想像的是「飯聚」,卻容易「犯罪」,所犯的罪就是試探上帝.在曠野漂流四十年,雖然嗎哪從天而降,然而,以色列民還是發怨言,「他們心中試探神,隨自己所欲的求食物,並且妄論上帝說:上帝在曠野豈能擺設筵席嗎?」(詩七十八18-19) 人常常在死蔭幽谷中只看見自己的困難和欠缺,而看不見上帝的作為.

此時此地,你仍相信「在我敵人面前,上帝為我擺設筵席」的可能嗎?還是充滿失望和犬儒地妄論上帝說:「上帝在曠野豈能擺設筵席嗎?」「飯聚」與「犯罪」只是一念之差而已.

3．馬太福音22章:天國的筵席

第三場筵席出現在馬太福音22章耶穌講的天國比喻裡面:

不要以為你請食飯,筵席就辦得成.有些時候,你設擺了筵度,人們都不一定赴會；甚至答應了參加的,到時也會爽約.或許你會說,你請食飯當然沒有人到,如果是王帝請又不同了.是嗎?請聽耶穌的比喻:

2「天國好比一個王為他兒子擺設娶親的筵席,

3就打發僕人去,請那些被召的人來赴席,他們卻不肯來.

耶穌用這個「王為他兒子擺設娶親的筵席」的故事來比喻天國,指出就算天國好像這個筵席那麼愉快歡樂,也有人拒絕參加的.有些拒絕進入天國的,甚至是早前預留了座位的.他們為甚麼爽約呢?

4王又打發別的僕人,說:『你們告訴那被召的人,我的筵席已經預備好了,牛和肥畜已經宰了,各樣都齊備,請你們來赴席.』

5那些人不理就走了；一個到自己田裡去；一個做買賣去；

6其餘的拿住僕人,凌辱他們,把他們殺了.

爽約不一定有甚麼特殊的原因,只是有自己的事情要忙,有自己的人生議程要追尋,而不能容納天國的議程而已.其實,不去參加「飯聚」就不去罷了,何必因此而變成「犯罪」呢?這些爽約者如何「犯罪」?他們「凌辱」,「殺害」派來提醒他們赴約的人.為甚麼他們如此做呢?可能他們的議程和天國的議程敵對,因此,他們以為消滅了那些提醒他們天國議程的人,就可以不受干擾地追逐和推進自己的議程了.

原來天國的筵席也會遇到那麼多障礙,到底誰會進去享受這個筵席呢?

8於是對僕人說:『喜筵已經齊備,只是所召的人不配.

9所以你們要往岔路口上去,凡遇見的,都召來赴席.』

10那些僕人就出去,到大路上,凡遇見的,不論善惡都召聚了來,筵席上就坐滿了客.

為甚麼是他們呢?可能他們沒有太強烈的個人議程,或者他們的議程並非與天國的議程敵對,他們有空間和有彈性去參加這個天國的筵席.

讓我們重問自己起初的問題:你有興趣參加天國的筵席嗎?你的人生議程有阻礙你去參加天國的筵席嗎?

4．腓立比書4:同工,但能同台吃飯嗎?

今天的書信經課沒有出現「筵席」,難道也與「筵席」有關係嗎?有的,關係就在於沒有提及「筵席」.為甚麼沒有「飯聚」呢?相信不是因為限聚令.那麼,到底是甚麼原因呢?大家又如何能夠協助「飯聚」再成為可能呢?

為甚麼不能「飯聚」呢?

2我勸友阿蝶和循都基,要在主裡同心.

3我也求你這真實同負一軛的,幫助這兩個女人,

友阿蝶和循都基本來應該是「筵席」的成員,保羅特別提名勸告她們,看來他們正處身於不能「同台吃飯」的狀況,或者「同台吃食各自修行」的困難.保羅指出,他們兩者並不是敵人,而是熱心的信徒:

3因為她們在福音上曾與我一同勞苦；還有革利免,並其餘和我一同做工的,他們的名字都在生命冊上.

有時就是如此無奈,雙方都是信主得救熱心事奉的主內肢體,但就偏偏不能「在主裡同心」.這些情況,可能也在你的群體中出現過.保羅提醒同伴們見到這樣的狀況不要只顧「吃花生」,而要幫助他們化解矛盾.如何幫助呢?保羅提出以下建議:

第一,在群體中散播喜樂和謙讓:

4你們要靠主常常喜樂.我再說,你們要喜樂.

5當叫眾人知道你們謙讓的心.

把喜樂和謙讓貫注於群體內部出現困難的人際關係之中,可能比扮警察查案,探法官審判,或者扮懲教人員執行刑罰適合得多,因為這不是正邪之辨,也不是敵我之分,只是內部矛盾而已.

第二,在群體中保持信靠和盼望,等候那超越人所能理解的和平來臨:

主已經近了.6應當一無掛慮,只要凡事藉著禱告,祈求,和感謝,將你們所要的告訴上帝.

7上帝所賜,出人意外的平安必在基督耶穌裡保守你們的心懷意念.

常言道,心病還需心藥醫.當肢體間有「心病」的時候,醫心者是誰呢?保羅指醒我們,能夠「保守你們的心懷意念」的是基督耶穌；那醫治心痛的良藥是「上帝所賜出人意外的平安」.

第三,見賢思齊,效法保羅,修養品德

8弟兄們,我還有未盡的話:凡是真實的,可敬的,公義的,清潔的,可愛的,有美名的,若有甚麼德行,若有甚麼稱讚,這些事你們都要思念.

9你們在我身上所學習的,所領受的,所聽見的,所看見的,這些事你們都要去行,賜平安的上帝就必與你們同在.

在「不能同台吃飯」的困難中,或者遇到群體中有人「不能同台吃飯」的狀況,我們常常思念甚麼呢?充斥腦海的會不會是誰對誰錯?以及錯甚麼呢?如果想著那些,保羅建議不如思念美好的品德,例如:「真實的,可敬的,公義的,清潔的,可愛的,有美名的」.思念這些,也行出來,情況應該會有所不同.

三．總結

到底友阿蝶和循都基後來能否一齊參加「飯聚」呢?腓立比書沒有交待,不過,我們相信這一定是保羅的期望.我們能否一齊參加「飯聚」呢?相信這也是我們的期望.今天的經課讓我們想到,主日崇拜之後的「飯聚」尚為小事,更值得參加的,是上帝所預備的筵席.上帝邀請我們一齊加入祂的筵席,這場「各人臉上的眼淚」被擦去,「普天下他百姓的羞辱」被除掉的筵席,這場不散的筵席由加入的一刻開始,可以延續至永恆.



\newpage

\section{神造光又造暗}
\label{sec:522}
\textbf{鄧瑞強博士 - 神造光,又造暗}
\newline
\newline
\hyperref[sec:521]{< < < PREV SERMON < < <}
~
\hyperref[sec:toc]{[返目錄]}
~
\hyperref[sec:523]{> > > NEXT SERMON > > >}


以賽亞書 45:1-7
\newline
\begin{longtable}{cl}
\hline
\hline
章節 & 經文 (和合本修訂版)\\
\hline
45:1 & \begin{tabularx}{0.7\textwidth}{X} 耶和華對所膏的居魯士如此說，他的右手我曾攙扶，使列國降服在他面前，列王的腰帶我曾鬆開，使城門在他面前敞開，不得關閉： \end{tabularx} \\ \\ \relax
45:2 & \begin{tabularx}{0.7\textwidth}{X} 「我要在你前面行，修平崎嶇之地。我必打破銅門，砍斷鐵閂。 \end{tabularx} \\ \\ \relax
45:3 & \begin{tabularx}{0.7\textwidth}{X} 我要將暗中的寶物和隱藏的財富賜給你，使你知道提名召你的就是我－耶和華，以色列的神。 \end{tabularx} \\ \\ \relax
45:4 & \begin{tabularx}{0.7\textwidth}{X} 因我的僕人雅各，我所揀選的以色列，我提名召你；你雖不認識我，我也加給你名號。 \end{tabularx} \\ \\ \relax
45:5 & \begin{tabularx}{0.7\textwidth}{X} 我是耶和華，再沒有別的了；除了我以外再沒有神。你雖不認識我，我必給你束腰。 \end{tabularx} \\ \\ \relax
45:6 & \begin{tabularx}{0.7\textwidth}{X} 從日出之地到日落之處使人都知道除我以外，沒有別的。我是耶和華，再沒有別的了。 \end{tabularx} \\ \\ \relax
45:7 & \begin{tabularx}{0.7\textwidth}{X} 我造光，又造暗；施平安，又降災禍；做成這一切的是我－耶和華。 \end{tabularx} \\ \\
[1ex]
\hline
\hline
\end{longtable}
講題:神造光,又造暗

經文:以賽亞書45:1-7

講員:鄧瑞強博士(香港神學院)

場合:香港中文大學崇基學院禮拜堂主日崇拜

日期:2020年10月18日

各位弟兄姊妹,早安,願各人平安.

我們活在歷史劇變的時刻裡.

因著疫情,我們進入一種生活的新常態裡.所謂「新常態」,即舊時從未有過的,一種全新的生活形態.事實的確如此,從出生的那日開始,我從未試過連續帶著口罩生活了近十個月,也從未試過當學生走過來時我叫他保持距離.

我們的社會也急劇轉變,很多人說,這是一個與舊時風貌完全不一樣的社會.

世界大局亦在劇變,中國與美國是政治和經濟上最大的兩塊板塊,我們知道,當板塊摩擦並互相壓迫時,便會地震,也可能有海嘯.政治和經濟上的地震和海嘯,不知何時發生?不知會否發生?但一旦發生,影響會很巨大.

在歷史劇變的時刻,人心惶惶之時,最需要先知的信息.先知們站在永恆的磐石上,以超越歷史巨變的遠大目光,發出洪亮的上帝之言,為洪流中顛簸不定的生命小舟,下一個大錨,將之寄托在永恆的神手中.

在舊約,以色列人歷史中最大的巨變,是主前587年,耶路撒冷被巴比倫大軍圍困,城被攻破,猶大人國破家亡,被擄他鄉.從此,生活於異域,難再見故鄉.詩篇137篇,最能表達他們的思鄉情懷.詩人說:「我們曾在異域,在巴比倫的河邊坐下,一追想故鄉便哭了.我們將琴掛在柳樹上.(在中國詩詞裡,應會掛在楊柳樹上,以表達離別之情)……我們哪有心情在異域唱故鄉的,頌讚上帝的歌呢?」對神的子民來說,這是異常傷痛的日子.故國沒有了,被巴比倫滅了.故城沒有了,被巴比倫夷平了.最重要的,是聖殿沒有了,這是他們生命的中心,是他們精神的家園,是他們夢縈魂牽的靈魂故鄉.現在,一切都沒有了.在這歷史巨變的時刻,在這歷史的淒慘時刻,先知還有什麼話可說呢?

今日的講道經文,是今日的舊約經課.我們細心聽聽先知以賽亞的信息.

經文:以賽亞書45:1－7

賽45:1　我─耶和華所膏的塞魯士；我攙扶他的右手,使列國降伏在他面前.我也要放鬆列王的腰帶,使城門在他面前敞開,不得關閉.我對他如此說:

賽45:2　我必在你前面行,修平崎嶇之地.我必打破銅門,砍斷鐵閂.

賽45:3　我要將暗中的寶物和隱密的財寶賜給你,使你知道提名召你的,就是我─耶和華,以色列的神.

賽45:4　因我僕人雅各,我所揀選以色列的緣故,我就提名召你；你雖不認識我,我也加給你名號.

賽45:5　我是耶和華,在我以外並沒有別神；除了我以外再沒有神.你雖不認識我,我必給你束腰.

賽45:6　從日出之地到日落之處使人都知道除了我以外,沒有別神.我是耶和華；在我以外並沒有別神.

賽45:7　我造光,又造暗；我施平安,又降災禍；造作這一切的是我─耶和華.

學者一般將以賽亞書分成三個部分:第一部分是第1章至39章,第二部分是第40章至55章,第三部分是第56章至66章.今日的講道經文屬於第二部分,寫作時期大約是波斯王塞魯士登基的日子至他擊敗巴比倫而釋放猶大人的日子,即主前550年至539年.主前539年,波斯滅巴比倫,新王上場,大赦天下,猶大人可以回歸家鄉.但是,寫作今天的講道經文的時候,這日子還未到,大地依然黑暗,猶大人仍在受苦的當下.就在這時期,先知發出從神而來的信息.

首先,我們看看神使用誰去改變歷史的進程.

賽45:1　我─耶和華所膏的塞魯士……

當神的子民覺得天地無光,漆黑一片時,神已準備好創造光明.猶如上古之初,創世記開場白的情節,地是空虛混沌,淵面黑暗.「淵」,即大水.混亂的大水,在近東文化,代表混亂與邪惡.就在這混沌無光,淵面黑暗之時,神的靈已運行在水面上.神在混亂的大水之上,駕馭邪惡.詩篇29:10說:「洪水泛濫之時,耶和華坐著為王」,就是同一個意思.神在黑暗與邪惡之上,掌管一切.在祂預定的時刻,便會終止黑暗,綻放光明.

「我─耶和華所膏的塞魯士」,他來的時候,便是神救贖祂子民的日子.塞魯士,一個波斯王,神膏立他.「膏立」這個字,就是「彌賽亞」這個字,即新約「基督」這個字.所謂「膏立」,就是祝聖他,使之屬於神,能被神所用.在歷史的某一時刻,神使用塞魯士為祂的工具.

賽45:1　我─耶和華所膏的塞魯士；我攙扶他的右手,使列國降伏在他面前.我也要放鬆列王的腰帶,使城門在他面前敞開,不得關閉.我對他如此說:

賽45:2　我必在你前面行,修平崎嶇之地.我必打破銅門,砍斷鐵閂.

神握著塞魯士的右手.在戰場上,右手握武器,左手握盾牌.左邊有保護,右邊容易受攻擊.神握著他的右手,即站在他的右邊,這就是保護他.塞魯士要制服列國,神「放鬆列王的腰帶」.古時,束緊腰帶,就是準備作戰.放鬆列王的腰帶,即解除列王的武裝.在神幫助下,塞魯士王能攻破城門,長驅直進,戰無不勝.

神祝聖塞魯士王,保護他,讓他戰無不勝.這樣的一個人,是一個認識神的人,敬畏神的人嗎?

賽45:5　……你雖不認識我,我必給你束腰.

原來,塞魯士並不認識神；但就算這樣,不會阻礙神使用他.

經文這樣輕描淡寫,卻開出重要的信仰思考.原來,不單神的子民屬於神,天下間的一切人也屬於神.神不單使用認識祂的人,神也可以使用不認識祂的人.神叫這人來,這人便來.叫這人去,這人便去.日光之下,萬民都在神的安排下,演活神安排好的角色.巴比倫的興起,有神的安排.巴比倫的滅亡,也有神的安排.當猶大國被巴比倫所滅,聖城被人蹂躪,聖殿被人搶掠時,神的子民可能會問:是神不存在嗎?是神不夠能力嗎?是神被巴比倫的神明擊敗嗎?先知在這裡清楚的說,一王興,出於神的旨意.一王滅,也出於神的旨意.耶和華不是被巴比倫的神明馬杜克(Marduk)所擊敗,馬杜克根本什麼都不是,在巴比倫之上作王的神明,就只有耶和華.不是耶和華被擊敗,所以巴比倫興起,消滅了猶大；而是耶和華興起巴比倫,去懲罰猶大.當時限到了時,神便興起波斯,去消滅巴比倫,去釋放神受苦的子民.當歷史發生巨變時,請不要驚惶,以為神不在,以為神看不到.神在,神看見.萬事萬物,都在神的掌控之中.

賽45:3　我要將暗中的寶物和隱密的財寶賜給你,使你知道提名召你的,就是我─耶和華,以色列的神.

關於這節經文,最新近的研究,都不能確定它的意思,學者們的看法,都只是猜想.既然大家都只是猜,我也不妨猜一猜.今日的講道經文,全段都是詩歌體.當讀希伯來文的詩歌體時,若意義不明,可看看經文是否有相關的平行經文.若有,可從平行經文中理解其含意.這裡講的「暗」,在第7節再出現.在第7節那裡,「暗」與「光」相對.如此,在第3節這裡,「寶物」和「財寶」的意思,會否指向「光明」?若如此,「暗中的寶物」,「隱密的財寶」的意思重點,便是在「暗中」的東西,如今「見光」了,被人認出是「寶物」了.黑暗終於得見光明.若意思是這樣,則這經文是講:在悠長的黑暗歷史裡,神使塞魯士成為帶來光明的人.塞魯士會否因此而體會:在冥冥之中,有一神力,讓他出現在歷史舞台上,成為帶來光明的人?

賽45:4　因我僕人雅各,我所揀選以色列的緣故,我就提名召你；你雖不認識我,我也加給你名號(或者譯:給你榮耀).

神在歷史中工作,是為了祂的僕人雅各.雅各,後來改名叫以色列.無論是雅各抑以色列,意思都指向與神立約的子民.神與某些人,有「約」的關係.這些人可說是神的代表,神的大使,他們將神的旨意活出來,神也保守他們,藉著他們建立屬神的國度.這些人的命運,在神手中.當他們守約時,神祝福他們.當他們違約時,神刑罰他們.罰惡揚善,像父親對待祂的兒女一樣.

神的子民,違反盟約,忘記神,於是,神興起巴比倫來警戒他們,他們被擄去巴比倫.到有一天,當神定的時候到了,祂就興起塞魯士,擊敗巴比倫,釋放神的子民.

賽45:5　我是耶和華,在我以外並沒有別神；除了我以外再沒有神.

賽45:6　從日出之地到日落之處使人都知道除了我以外,沒有別神.我是耶和華；在我以外並沒有別神.

歷史如此變化,是叫世人明白,善的後果是什麼,惡的後果是什麼.更重要的,是讓世人明白,歷史的軌跡,不是亂來的,而是有神在,祂主宰朝代興亡.

賽45:7　我造光,又造暗；我施平安,又降災禍；造作這一切的是我─耶和華.

這是今日講道經文的神學高潮,也是其神學結論.

在這經文裡,「光」與「平安」平行,「暗」與「災禍」平行.

「暗」連於「災禍」.

「黑暗」,在以賽亞書,是指神的子民國破家亡,被擄巴比倫的事件.最明顯的經文例子,莫如以賽亞書49:9如此說:「對那被捆綁的人說:出來吧!對那在黑暗的人說:顯露吧!」在這裡,很清楚,「黑暗」,對應於「被捆綁」.「黑暗」與「災禍」,都是指神的子民被擄巴比倫這歷史事件.

「我造光,又造暗」,很明顯,這講法來自創世記,在那裡,神在黑暗中創造光明.如今,神會再一次,在歷史的黑暗中,創造光明.這是新一回合的創造,讓歷史越過黑暗長夜,再見光明.神興起塞魯士,正正就是要改變歷史的軌跡,讓被擄的人得到自由.這生命的重拾自由,就是這經文講的「平安」.

這經文的特別之處,是刻意指出:神造光,又造暗；神施平安,又降災禍.歷史的美好時刻,來自神.歷史的艱難時刻,同樣來自神.黑暗的日子,災禍的日子,不是出於不測的命運,不是由於神被別的力量擊倒,反而是因為祂的子民與世人一樣行惡,祂就創造災禍,去糾正他們的錯誤.又正正因為祂創造惡,故此,祂就有權亦有能力去取消惡,創造光明與平安.

黑暗的日子,並沒有脫離神的掌控,也因此,當神要取消黑暗時,也不用詢問誰.先知這樣宣告,是要堅定神子民的心,叫他們在歷史的巨變時刻,不要失去對神的信靠,不用驚惶,只需好好反省自己是否行在神的旨意中.

親愛的弟兄姊妹,在歷史巨變的關頭,讓我們更敬虔,更忍耐,以信心看到「造作這一切的,是耶和華」.

願榮耀歸上主,願人間享和平.

(附記:本人在崇基禮拜堂的講章已結集成《荷花池畔春草綠》講道集,試讀本可參閱:\url{http://cc.cuhk.edu.hk/dev/images/content/fund-raising-campaigns/university-sermons-by-dr-sk-tang/CCC-SK-Tang-Sermon-Text-Sample-2.pdf}.較早期的崇基講章,也已結集成《敢於跟隨主》一書.有興趣的朋友可到各大基督教書局購買.)



\newpage

\section{請問}
\label{sec:523}
\textbf{梁遠耀牧師 - 請問……}
\newline
\newline
\hyperref[sec:522]{< < < PREV SERMON < < <}
~
\hyperref[sec:toc]{[返目錄]}
~
\hyperref[sec:524]{> > > NEXT SERMON > > >}


馬太福音 22:34-46
\newline
\begin{longtable}{cl}
\hline
\hline
章節 & 經文 (和合本修訂版)\\
\hline
22:34 & \begin{tabularx}{0.7\textwidth}{X} 法利賽人聽見耶穌堵住了撒都該人的口，他們就聚集在一起。 \end{tabularx} \\ \\ \relax
22:35 & \begin{tabularx}{0.7\textwidth}{X} 其中有一個人是律法師，要試探耶穌，就問他： \end{tabularx} \\ \\ \relax
22:36 & \begin{tabularx}{0.7\textwidth}{X} 「老師，律法上的誡命哪一條是最大的呢？」 \end{tabularx} \\ \\ \relax
22:37 & \begin{tabularx}{0.7\textwidth}{X} 耶穌對他說：「你要盡心、盡性、盡意愛主—你的神。 \end{tabularx} \\ \\ \relax
22:38 & \begin{tabularx}{0.7\textwidth}{X} 這是最大的，且是第一條誡命。 \end{tabularx} \\ \\ \relax
22:39 & \begin{tabularx}{0.7\textwidth}{X} 第二條也如此，就是要愛鄰如己。 \end{tabularx} \\ \\ \relax
22:40 & \begin{tabularx}{0.7\textwidth}{X} 這兩條誡命是一切律法和先知書的總綱。」 \end{tabularx} \\ \\ \relax
22:41 & \begin{tabularx}{0.7\textwidth}{X} 法利賽人聚集的時候，耶穌問他們： \end{tabularx} \\ \\ \relax
22:42 & \begin{tabularx}{0.7\textwidth}{X} 「論到基督，你們的意見如何？他是誰的後裔呢？」他們說：「是大衛的。」 \end{tabularx} \\ \\ \relax
22:43 & \begin{tabularx}{0.7\textwidth}{X} 耶穌說：「這樣，大衛被聖靈感動，怎麼還稱他為主，說： \end{tabularx} \\ \\ \relax
22:44 & \begin{tabularx}{0.7\textwidth}{X} 『主對我主說：你坐在我的右邊，等我把你的仇敵放在你腳下？』 \end{tabularx} \\ \\ \relax
22:45 & \begin{tabularx}{0.7\textwidth}{X} 大衛既稱他為主，他怎麼又是大衛的後裔呢？」 \end{tabularx} \\ \\ \relax
22:46 & \begin{tabularx}{0.7\textwidth}{X} 沒有一個人能回答一句話，從那日以後沒有人敢再問他甚麼。 \end{tabularx} \\ \\
[1ex]
\hline
\hline
\end{longtable}
利未記 19:1-2,15-18
\newline
\begin{longtable}{cl}
\hline
\hline
章節 & 經文 (和合本修訂版)\\
\hline
19:1 & \begin{tabularx}{0.7\textwidth}{X} 耶和華吩咐摩西說： \end{tabularx} \\ \\ \relax
19:2 & \begin{tabularx}{0.7\textwidth}{X} 「你要吩咐以色列全會眾，對他們說：你們要成為聖，因為我耶和華－你們的神是神聖的。 \end{tabularx} \\ \\ \relax
19:3 & \begin{tabularx}{0.7\textwidth}{X} 你們各人都當孝敬父母，也要守我的安息日。我是耶和華－你們的神。 \end{tabularx} \\ \\ \relax
19:4 & \begin{tabularx}{0.7\textwidth}{X} 你們不可轉向虛無的神明，也不可為自己鑄造神像。我是耶和華－你們的神。 \end{tabularx} \\ \\ \relax
19:5 & \begin{tabularx}{0.7\textwidth}{X} 「你們宰殺祭牲獻平安祭給耶和華的時候，要獻得使你們可蒙悅納。 \end{tabularx} \\ \\ \relax
19:6 & \begin{tabularx}{0.7\textwidth}{X} 這祭物要在獻的當天或第二天吃；若有剩到第三天的，就要用火焚燒。 \end{tabularx} \\ \\ \relax
19:7 & \begin{tabularx}{0.7\textwidth}{X} 第三天若再吃，這祭物是不潔淨的，必不蒙悅納。 \end{tabularx} \\ \\ \relax
19:8 & \begin{tabularx}{0.7\textwidth}{X} 吃的人必擔當自己的罪孽，因為他褻瀆了耶和華的聖物，這人必從百姓中剪除。 \end{tabularx} \\ \\ \relax
19:9 & \begin{tabularx}{0.7\textwidth}{X} 「你們在自己的地收割莊稼時，不可割盡田的角落，也不可拾取莊稼所掉落的。 \end{tabularx} \\ \\ \relax
19:10 & \begin{tabularx}{0.7\textwidth}{X} 不可摘盡葡萄園的葡萄，也不可拾取葡萄園中掉落的葡萄，要把它們留給窮人和寄居的。我是耶和華－你們的神。 \end{tabularx} \\ \\ \relax
19:11 & \begin{tabularx}{0.7\textwidth}{X} 「你們不可偷盜，不可欺騙，也不可彼此說謊。 \end{tabularx} \\ \\ \relax
19:12 & \begin{tabularx}{0.7\textwidth}{X} 不可指著我的名起假誓，褻瀆你神的名。我是耶和華。 \end{tabularx} \\ \\ \relax
19:13 & \begin{tabularx}{0.7\textwidth}{X} 「不可欺壓你的鄰舍，也不可偷盜。雇工的工錢不可在你那裡過夜，留到早晨。 \end{tabularx} \\ \\ \relax
19:14 & \begin{tabularx}{0.7\textwidth}{X} 不可咒罵聾子，也不可將絆腳石放在盲人面前。你要敬畏你的神。我是耶和華。 \end{tabularx} \\ \\ \relax
19:15 & \begin{tabularx}{0.7\textwidth}{X} 「你們審判的時候，不可不公正；不可偏護貧窮人，也不可看重有權勢人的臉，總要公平審判你的鄰舍。 \end{tabularx} \\ \\ \relax
19:16 & \begin{tabularx}{0.7\textwidth}{X} 不可在百姓中到處搬弄是非，不可陷害鄰舍的性命。我是耶和華。 \end{tabularx} \\ \\ \relax
19:17 & \begin{tabularx}{0.7\textwidth}{X} 「不可心裡恨你的弟兄；要指摘你的鄰舍，免得因他承擔罪過。 \end{tabularx} \\ \\ \relax
19:18 & \begin{tabularx}{0.7\textwidth}{X} 不可報仇，也不可埋怨你本國的子民。你要愛鄰如己。我是耶和華。 \end{tabularx} \\ \\
[1ex]
\hline
\hline
\end{longtable}
詩篇 1:1-6
\newline
\begin{longtable}{cl}
\hline
\hline
章節 & 經文 (和合本修訂版)\\
\hline
1:1 & \begin{tabularx}{0.7\textwidth}{X} 不從惡人的計謀，不站罪人的道路，不坐傲慢人的座位， \end{tabularx} \\ \\ \relax
1:2 & \begin{tabularx}{0.7\textwidth}{X} 惟喜愛耶和華的律法，晝夜思想他的律法；這人便為有福！ \end{tabularx} \\ \\ \relax
1:3 & \begin{tabularx}{0.7\textwidth}{X} 他要像一棵樹栽在溪水旁，按時候結果子，葉子也不枯乾。凡他所做的盡都順利。 \end{tabularx} \\ \\ \relax
1:4 & \begin{tabularx}{0.7\textwidth}{X} 惡人並不是這樣，卻像糠秕被風吹散。 \end{tabularx} \\ \\ \relax
1:5 & \begin{tabularx}{0.7\textwidth}{X} 因此，當審判的時候惡人必站立不住，罪人在義人的會眾中也是如此。 \end{tabularx} \\ \\ \relax
1:6 & \begin{tabularx}{0.7\textwidth}{X} 因為耶和華知道義人的道路，惡人的道路卻必滅亡。 \end{tabularx} \\ \\
[1ex]
\hline
\hline
\end{longtable}
帖撒羅尼迦前書 2:1-8
\newline
\begin{longtable}{cl}
\hline
\hline
章節 & 經文 (和合本修訂版)\\
\hline
2:1 & \begin{tabularx}{0.7\textwidth}{X} 弟兄們，你們自己知道我們來到你們那裡並不是徒然的。 \end{tabularx} \\ \\ \relax
2:2 & \begin{tabularx}{0.7\textwidth}{X} 我們從前在腓立比蒙難受辱，這是你們知道的，可是我們還是靠著神給我們的勇氣，在強烈反對中把神的福音傳給你們。 \end{tabularx} \\ \\ \relax
2:3 & \begin{tabularx}{0.7\textwidth}{X} 我們的勸勉不是出於錯誤，也不是出於污穢，也不是用詭詐。 \end{tabularx} \\ \\ \relax
2:4 & \begin{tabularx}{0.7\textwidth}{X} 但神既然認定我們經得起考驗，把福音託付我們，我們就照著傳講，不是要討人喜歡，而是要討那考驗我們的心的神喜歡。 \end{tabularx} \\ \\ \relax
2:5 & \begin{tabularx}{0.7\textwidth}{X} 因為我們從來沒有用過諂媚的話，這是你們知道的，也沒有藏著貪心，這是神可以作證的。 \end{tabularx} \\ \\ \relax
2:6 & \begin{tabularx}{0.7\textwidth}{X} 我們作為基督的使徒，雖然可以受人尊重，卻沒有向你們或向別人求榮耀， \end{tabularx} \\ \\ \relax
2:7 & \begin{tabularx}{0.7\textwidth}{X} 反而在你們當中心存溫柔，如同母親哺乳自己的孩子。 \end{tabularx} \\ \\ \relax
2:8 & \begin{tabularx}{0.7\textwidth}{X} 既然我們這樣愛你們，不但樂意將神的福音給你們，連自己的性命也樂意給你們，因為你們是我們所疼愛的。 \end{tabularx} \\ \\
[1ex]
\hline
\hline
\end{longtable}
講題:請問……

經文:馬太福音22:34-46;利未記19:1-2,15-18;詩篇1:1-6;帖撒羅尼迦前書2:1-8

講員:梁遠耀牧師

場合:香港中文大學崇基學院禮拜堂主日崇拜

日期:2020年10月25日

問安及禱告

各位弟兄姊妹平安,願主的恩惠,平安,常與你們同在.

在分享主道之先,請我們先一起禱告:「天父,願此刻僕人口中的言語,眾信徒的心思意念,蒙你悅納.奉主耶穌名求,誠心所願!」

引言

我曾在《讀者文摘》(Reader’s Digest)這本雜誌看過一篇有關運動項目的冷知識.很多運動項目,通常都是先有男性運動員參與,然後才有女性參與,所以很多時會是先有男子項目,才會有女性項目.但有小數是例外,例如韻律泳(或稱為水上芭蕾).在正規比賽時,韻律泳的運動員,是只有女性運動員參與.水上芭蕾的運動員,她們所穿的泳衣,有規定必須穿著不透視的泳衣,你知道原因嗎?這是希望展現公平比賽的原則,希望評判能夠專注在運動員的動作,而不是運動員的身體.

另一個我留意的運動項目是摔跤.摔跤是一項奧運會的比賽類別,但香港人較少留意.不同的摔跤(或稱摔角或角力)運動項目,基本上是比賽雙方,以不同的要求及技術,徒手相搏,把對方摔倒或壓下等.男子摔跤比賽,要求男性運動員於比賽前,必須剪指甲和剃鬚.要摔跤運動員於賽前剪指甲,我想大家都能理解,但你知道為甚麼要男性運動員剃鬚嗎?因為通常當男士有一兩天沒剃鬍鬚,會長出短短的鬍子(鬚根).有時這些短鬍子會比較硬身,當比賽時,雙方運動員的身體會有接觸,若運動員面上有短鬍子,在摔跤比賽時,這些短鬍子就會使對方感到不適.

有很多時,因為我們對宗教信仰的行為和教義等,會採取先接受,後理解的方式.但有時候時間久了,我們已接受,並可能已成為習慣,就沒有再進行理解,例如謝飯祈禱.請問大家曾否想過,或大家在生活實踐上,甚麼時候,會進行謝飯祈禱?早午晚三餐,我想大家就不會困惑.若進食一件三文治呢,你會祈禱嗎?退一步,吃一兩片餅干或曲奇呢?再退一步,只是喝杯咖啡,你又會祈禱嗎?我聽過有一打趣的說:「五元以上的食物才祈禱,因為一串魚蛋賣五元.」

若我們有以上的困惑,背後可能是因為我們把謝飯祈禱,視為一種要在進食前要進行宗教儀式.有沒有曾經試過因為美食當前,我們把食物放進口中,才醒起要祈禱,更彷彿遺忘了祈禱,而帶來罪疚感?為有食物而感恩,確是一件值得確定的事情,因聖經從來沒有一條誡命律例,要求我們在進食前,必須祈禱,也沒有任何一次大公會議,把飯前禱告,列為教義.但我重申,在進食前祈禱感恩,是一件美事,只是我們有這行為,久而久之,似乎有時會忘記謝飯祈禱的初衷──是感謝.若我們失去了感謝,飯前祈禱就只是一種宗教行為,而不是發自內心感恩的行動.所以,若你因為太餓了,或是一時忘記,沒有在進食前祈禱,這並不是罪,到你在進食期間才忽然想起了這件事,你可以隨時祈禱,我認為你這樣做,並沒有得罪上帝.若你在飯前已經祈禱,但在進食期間,感覺食物真的美味,或在飯後,感到滿足,也可以再次祈禱感恩呀.

了解及認識

不要讓我們的信仰,淪為只是一些宗教行為.但無可否認,信仰亦總是有一套既定的行為規範,並非完全是一種純粹個人的主觀願望.故此,我們必須先要明白規範背後的因由,才不會盲目跟從,或死守傳統.

即使上帝透過摩西吩咐以色列人,要成為神聖的子民.但在利未記19章1-2節,我們可看到上帝亦會解釋,因為祂是神聖的上帝,故此,祂的子民亦要具有這種上帝的特質.你或許認為,人類又能是神聖,怎能具有與上帝相同的特質?其實,經文中「神聖」一詞,是指「有分別」.我們的上帝是神聖的,因上帝與其他神靈或偶像,是有分別的.同樣,上帝吩咐以色列人,這群屬上帝的子民,他們是一群與其他民族有分別的人,與其他民族相比,除了在宗教祭禮的方式有所不同,在對照顧貧窮人方式,對金錢與財物的態度,人際相處的要求,對公義的追求上,也要有更高的要求,與其他民族有所不同.所以我們作為人類大家庭的一分子,我們會與一般人無異地生活,但我們也作為上帝的子民,也有我們與世界不一樣的生活態度,不可混淆,請清楚嗎?

請問你除了在主日崇拜聽道之外,還有沒有系統地去認識聖經,認識我們的信仰嗎?請問你對上一次上主日學或研讀有關聖經的書籍,是在何時?請花一點時間,去明白聖經.

默想與背誦

除了理解,傳統學習的方式,也包括背誦.對我而言,背誦其實是反覆地念誦或思考某一句子,或人與事,其實與基督教的靈修方式──默想,也是很相似的.

今天的學習,已不太強調反覆念誦,加上現代科技的幫助,很多資料已經可以在彈指之間,便能找到資料.但對我來說,因為反覆念誦,還是帶來益處的,除了能鍛鍊腦袋的記憶力,也讓人能安靜下來,也使人在思考的過程中,有更闊的內容.就好像小孩子要背誦乘數表,過程雖然有點乏味,而且計數機也是十分方便及準確,但乘數表卻是在數學學習上,必須下的功夫.

所以,對於詩篇第1篇,要得上帝的賜福,除了要理解第1節──「不從惡人的計謀,不站罪人的道路,不坐褻慢人的座位,」,詩人提醒我們不要與惡人同行(從),同站和同坐.在第2節也提醒我們,要晝夜思想上帝的話語.經常念誦和思想上帝的話,這是蒙福的.

請問你最喜愛哪幾句聖經的金句?你有經常念誦,能背出來嗎?

問問題

對於時常念誦經文,在新約時代的代表者,應屬法利賽人莫屬.他們原是一群在新舊約之間所興起的群體,他們追求愛慕和嚴謹地遵守律法,但可惜隨著時間過去,到耶穌時代,有些法利賽人漸漸變質.但法利賽人仍然是熱衷於念誦聖經,而且也會以間答方式,例如他們會彼此問答:「十誡之中,哪一條誡命較難實行,哪一條誡命又較易實行?」等等.所以作為日常學習,其中一位法利賽人,也是一位專研律法的律法師,問耶穌哪一條誡命是最大的?

作為學習討論,這是一條很好的問題.只可惜這位律法師的用意,並不是真心來向耶穌討教,聖經和合本以「試探」來翻譯描寫這位律法師的動機,而聖經現代中文譯本,以更直率的字──「陷害」來描寫律法師的動機.律法師想藉這條問題,希望為耶穌設陷阱,找耶穌說話上的把柄.結果,耶穌引用申命記6:5和利未記19:18來回應,就是人要盡心,盡性,盡意,愛主你的神,又要愛人如己.

耶穌也向法利賽人詢問有關基督論的問題,就是基督與大衛的關係.基督既是大衛王的子孫,也是大衛王的主,兩者如何平衡呢?耶穌發問之後,竟然沒有法利賽人能回應耶穌的提問.聖經說:「從那日以後,也沒有人敢再問祂甚麼.」

人常生活之中,我們會用問句的方式,來詢問或表達我們思想.例如:「你知道現在是甚麼時間嗎?」若單看字面,只是問你是否知道,你可以回答「知道」或「不知道」,便滿足了這問題.但這問題也可能是向你詢問現在是甚麼時間,詢問者可能期望你回應現在是上午11點或其他具體的時間.當然,發出這問題人,若是你的家人,正在等候你出門,你就不用回答甚麼,只要盡情穿好衣服出門便可以了,因為你的家人在催促你.律法師的問題是陷阱,希望找耶穌的錯處；耶穌也給他們問題,只是法利賽人似乎對彌賽亞已有一套既定的答案,卻沒有深思.

各位,一道好的問題,我不會說一定比答案更重要,但好問題的價值,其貢獻是十分之大.因為問題可以引導我們思考,幫助我們反省.在信仰裡,中世紀經院哲學(scholasticism)的重要人物安瑟倫(Anselmus),提出「信仰尋求理解」(faith seeking understanding)的說法,表示信仰需要有理性的.我想,提問,正是我們在理性追尋認識信仰的重要一步.

實踐

當然,聖經不單是一本用作學習及誦讀的書籍.我在神學生時間,曾在一所新界區的教會實習,而實習的其中一項內容,試探訪.有一次隨牧師去一個信徒家庭探訪,在家中,我看到有一本聖經,高高的放在一個玻璃門的組合櫃上.我感到奇怪,為何要把聖經放在這樣高的位置,十分不方便拿取.後來我詢問我的實習督導牧師,牧師表示因這家庭中的長者,是在數十年前信主,由傳統拜偶像的信仰,歸信基督.又由於他們沒有受教育的機會,所以他們不會讀聖經,但他們卻把聖經放在家中顯眼的地方,取代以往在家中放偶像的位置.

聖經並不是家中的陳設品,也不只是家中書櫃的書籍.以中華基督教會(The Church of Christ in China)為例,她有三條信仰綱要,其中一條的內容,是「接納舊新兩約聖經為上帝之言……為吾人信仰及本份無上之準則.」聖經是上帝的話,是我們生命及行為的指導.故此,閱讀聖經,識認聖經的目的,是實踐信仰.

保羅也曾向帖撒羅尼迦的信徒分享,「神既然驗中了我們,把福音託付我們,我們就照樣講,不是要討人的喜歡,乃是要討那察驗我們心的神喜歡.」保羅被上帝選召,作福音的使徒,與人分享福音.而每位信徒,我相信上帝都會給與我們恩賜,去建立教會.信徒就好像身體中的不同部分,各有不同的作用,發揮不同的功能,使身體能正常運作.

結語

信仰有知識的內容,可以學習和理解,但她不單是一套宗教知識；信仰有不同的歷程,可以一步一步進深經歷,但她不單是一種情感與情緒；信仰有行為的誡律,可以模仿和遵行,但她不單是一套行為規範.信仰仍是我們全人投入,愛主愛人.就正如主耶穌對律法誡命和先知教導的結語:「你要盡心,盡性,盡意,愛主你的神」,又「要愛人如己.」

祈禱

請我們一同祈禱:

天父上帝,你是奇妙和值得稱頌的,因為你藉聖經,承載著你的道,使人能認識你,明白你的心意.願我們能以聆聽,閱讀,念誦及研經等功夫,能進深認識你.

求你賜我們悟性,使我們能理解你的心意,更能在生活中,活出我們作為上帝子民應有的生活態度及價值.使你的旨意成就在地上,如同行在天上.

奉主名求.阿們!



\newpage

\section{信仰只說不行行嗎}
\label{sec:524}
\textbf{高國雄牧師 - 信仰只說不行,行嗎?}
\newline
\newline
\hyperref[sec:523]{< < < PREV SERMON < < <}
~
\hyperref[sec:toc]{[返目錄]}
~
\hyperref[sec:525]{> > > NEXT SERMON > > >}


馬太福音 23:1-12
\newline
\begin{longtable}{cl}
\hline
\hline
章節 & 經文 (和合本修訂版)\\
\hline
23:1 & \begin{tabularx}{0.7\textwidth}{X} 那時，耶穌對眾人和門徒講論， \end{tabularx} \\ \\ \relax
23:2 & \begin{tabularx}{0.7\textwidth}{X} 說：「文士和法利賽人坐在摩西的位上， \end{tabularx} \\ \\ \relax
23:3 & \begin{tabularx}{0.7\textwidth}{X} 所以凡他們所吩咐你們的，你們都要謹守遵行。但不要效法他們的行為，因為他們能說不能行。 \end{tabularx} \\ \\ \relax
23:4 & \begin{tabularx}{0.7\textwidth}{X} 他們把難挑的重擔捆起來，擱在人的肩上，但自己一個指頭也不肯動。 \end{tabularx} \\ \\ \relax
23:5 & \begin{tabularx}{0.7\textwidth}{X} 他們所做的一切事都是要讓人看見，所以把佩戴的經匣加寬了，衣裳的繸子加長了， \end{tabularx} \\ \\ \relax
23:6 & \begin{tabularx}{0.7\textwidth}{X} 喜愛宴席上的首座、會堂裡的高位， \end{tabularx} \\ \\ \relax
23:7 & \begin{tabularx}{0.7\textwidth}{X} 又喜歡人們在街市上向他們問安，稱呼他們拉比。 \end{tabularx} \\ \\ \relax
23:8 & \begin{tabularx}{0.7\textwidth}{X} 但你們不要接受拉比的稱呼，因為只有一位是你們的老師；你們都是弟兄。 \end{tabularx} \\ \\ \relax
23:9 & \begin{tabularx}{0.7\textwidth}{X} 也不要稱呼地上的人為父，因為只有一位是你們的父，就是在天上的父。 \end{tabularx} \\ \\ \relax
23:10 & \begin{tabularx}{0.7\textwidth}{X} 不要接受師傅的稱呼，因為只有一位是你們的師傅，就是基督。 \end{tabularx} \\ \\ \relax
23:11 & \begin{tabularx}{0.7\textwidth}{X} 你們中間誰為大，誰就要作你們的用人。 \end{tabularx} \\ \\ \relax
23:12 & \begin{tabularx}{0.7\textwidth}{X} 凡自高的，必降為卑；自甘卑微的，必升為高。 \end{tabularx} \\ \\
[1ex]
\hline
\hline
\end{longtable}
講題:信仰只說不行,行嗎?  Walk the Talk

經文:馬太福音 23:1-12

講員:高國雄牧師

場合:香港中文大學崇基學院禮拜堂主日崇拜

日期:2020年11月1日

今天是聖靈降臨後第二十二主日,亦是諸聖日,當世界因萬聖節前夕(Halloween)而去扮鬼扮馬地狂歡時,我們作信徒的不要忘記,hallows是聖人的意思,all hallows evening亦即是all saints day前夕的意思,諸聖日是記念教會史上為主殉道聖徒的節日,讓我們在今天也回想起先聖先賢的忠心服侍,提醒我們在今天這世代,仍然盡力持守信仰,致死忠心的信,求主憐憫!

在我年幼時,間唔中會聽到有人話「大隻講,口爽爽」,又或者係「講就天下無敵,做就無能為力」,來形容一些「得把口」的人,願意一生信守承諾,竭盡所能去承擔的人在這世代可能已經被認為已經絕種,難得一見的一群.

今天的經課提醒我們,福音書中曾經記載:

馬太福音12:1-12

1那時,耶穌對眾人和門徒講論,2說:「文士和法利賽人坐在摩西的位上,3所以凡他們所吩咐你們的,你們都要謹守遵行.但不要效法他們的行為,因為他們能說不能行.4他們把難挑的重擔捆起來,擱在人的肩上,但自己一個指頭也不肯動.5他們所做的一切事都是要讓人看見,所以把佩戴的經匣加寬了,衣裳的繸子加長了,6喜愛宴席上的首座,會堂裡的高位,7又喜歡人們在街市上向他們問安,稱呼他們拉比.8但你們不要接受拉比的稱呼,因為只有一位是你們的老師；你們都是弟兄.9也不要稱呼地上的人為父,因為只有一位是你們的父,就是在天上的父.10不要接受師傅的稱呼,因為只有一位是你們的師傅,就是基督.11你們中間誰為大,誰就要作你們的用人.12凡自高的,必降為卑；自甘卑微的,必升為高.

聖經福音書中常常提到要挑戰,試探,甚至攻擊耶穌的幾班人,當中的常客包括文士及法利賽人；文士在希伯來名稱字源的意思是抄寫,排列,數算,是舊約律法和經文的專家,專職講解並執行聖經教訓的人,也有時稱為律法師；法利賽人這名稱原意是「分離主義者」,在祭司體系與文士中分離出來,他們認為忠於耶和華的律法,照猶太教最高的理想與外教隔離是比任何事物都重要的,因此必須要過著比任何人皆更敬虔的生活,謹守律法與及律法一切的延伸.

原本這二批認真對待信仰的人都是值得欣賞及學習的,但歷史上見證著這些人漸漸的敗壞,最後當中不少因為走向極端,堅持了一些本應是附帶/輔助的東西,而失去本來的焦點；這就是為什麼耶穌多次直接指責他們假冒為善,假冒為善在希臘文原文中有「演戲」的意思,因此他們非常敵視耶穌,但為何耶穌在V3中提到「凡他們所吩咐你們的,你們都要謹守遵行.」呢?難道耶穌認同他們的教導嗎?

首先我們要認識V2中所提到「摩西的位」是什麼,摩西是以色列人史上最偉大的領袖,曾帶領以色列民離開埃及這為奴之地,進到上帝所賜的迦南美地當中,並按上帝心意頒布律法誡命,其教導成為了猶太人生活的典範,在當時的猶太人心目中,摩西是拉比的典範,是妥拉(Torah)的最高權威.但相信這裡所提到的不是指他們承繼了象徵意義中「典範」的角色,因為經過考古的研究,在哥拉汛和迦百農的古猶太會堂遺跡中皆發現,在會堂右方講經台階的旁邊,設有一個名為「摩西的位」的座位,供讀經後坐下作講解與教導之用；可想而知,當時耶穌很有可能直接指著這一種解經教導的形態,而不是指文士及法利賽人日常教導的內容的權威性；「所以」耶穌邀請我們「凡他們所吩咐你們的」,就是指凡依從摩西律法所引申的教導,我們都要謹慎遵守,這是理所當然的.

在這裡耶穌提醒我們,不要因人廢言,合理的教導無論出自何人都應該學習並遵守的,我們太多時候是按我們與對象的關係而選擇接受教導與否,真理不會因為出自不同人的口而有所改變.

但可惜的是耶穌卻跟著說「但不要效法他們的行為,因為他們能說不能行.他們把難挑的重擔捆起來,擱在人的肩上,但自己一個指頭也不肯動.」

道要聽,更要行,只講不做,口不對心,最終只會成為一個虛偽的人,這是何等可悲的事啊!

雅各書1:22-25提醒我們

「22但是,你們要作行道的人,不要只作聽道的人,自己欺騙自己. 23因為只聽道而不行道的,就像人對著鏡子觀看自己本來的面目, 24注視後,就離開,立刻忘了自己的相貌如何. 25惟有查看那完美,使人自由的律法,並且時常遵守的,他不是聽了就忘,而是切實行出來,這樣的人在所行的事上必然蒙福.」

為什麼有人願意去照鏡後,就離開並忘記?

為何人會選擇自己欺騙自己?

「認真就輸了…」成為這一代人常用的表達,不時就會聽到有人對承諾作這回應「講下je ,唔使太認真ge !」,人為何不願意選擇「認真」?寧願作假,或是「扮」作假又是什麼導致的呢?

或許這種表達背後有著自我形象與自信的問題,大家在學期間,相信身邊最少有一位成績超然的同學不斷表示自己沒有溫習/溫習不足,考試測驗後馬上表示自己沒有溫習那試卷內的內容…但最終卻成績極優異,明顯他們所表達的根本與事實不乎；這不想讓人知道自己曾經投入/努力,相信原因是害怕如若努力/盡力後仍然失敗,那就好像代表了自己真正失敗,寧願別人不知道自己努力過,最終如若失敗那就可以歸咎於「未盡力」；認真就輸了說到底就是怕讓人見到真正的自己,一個自己也不願意面對的自己…

另有一些人就在另一個極端,自己仍然未做到(甚至未開始做),就已經表達到已經完成,甚至完美,「講jor 等如做jor 」在我們的生活,工作間,甚至教會內其實並不罕見,冗長的會議後可能連一個行動的議決也沒有；每年新計劃的口號如何的陳腔濫調,屋苑商場的名稱如何「高,大,空」「世界之巔」「超級無敵」「宇宙最強」等；信徒們不斷地在社交媒體中引用聖經金句,發出與信仰相關的圖像,share了,like了就真的就等如行了「公義」,好了「憐憫」?真的就傳了「福音」?回應了大使命嗎?相信你我撫心自問,應該很清楚答案…

只講不做最簡單的原因不是不明白,不是不知道,而是不願意,不願意付出代價,在今日講求性價比的年代,人總期望付出低於應該付的代價,但能夠獲得應有(或更大)的回報,甚至不勞而獲,山寨版(A貨)的出現,就是為了滿足這心態,假冒為善的信徒其實就係A貨信徒,只以外表包裝「扮」真實的信仰,就算有多相似,也絕對不會等如真實的；我們不需要「似」屬靈人,學習「講」出很屬靈的話,我們要成為真正屬靈的人,請要清晰焦點與核心,敬虔的外貌與宗教的行為並不能夠使你成為主的門徒,唯有真誠認罪悔改,讓基督成為你的救主及生命的主,才是得救的唯一途徑.

接受自己的真相/真我是信仰中重要的一步,我們要相信上帝的創造及管治,當中包括萬物,當然亦包括你及我；真正的信心,需要配合行動才有意義,只睇醫生而不食藥,真的能醫好病嗎?

自我形象出現問題而導致內外不一致可以有另一邊的鐘擺情況,人就會自我膨脹.

V5他們所做的一切事都是要讓人看見,所以把佩戴的經匣加寬了,衣裳的繸子加長了,6喜愛宴席上的首座,會堂裡的高位,7又喜歡人們在街市上向他們問安,稱呼他們拉比.8但你們不要接受拉比的稱呼,因為只有一位是你們的老師；你們都是弟兄.9也不要稱呼地上的人為父,因為只有一位是你們的父,就是在天上的父.10不要接受師傅的稱呼,因為只有一位是你們的師傅,就是基督.11你們中間誰為大,誰就要作你們的用人.12凡自高的,必降為卑；自甘卑微的,必升為高.

真正的敬虔是個人與上帝建立親密的關係,刻意做出外在「敬虔的行為」讓人看見的唯一原因就是因為裡面沒有與主「真實」的關係,因為真正與主契合的人,就必會認清自己的角色及位置,施恩的主與蒙恩的罪人這一個最重要的關係,真能體會耶穌基督降卑為人,道成肉身的價值及意義.

使徒保羅曾提醒我們:

林前4:7 使你與人不同的是誰呢?你有甚麼不是領受的呢?若是領受的,為何自誇,彷彿不是領受的呢?

如果我們清楚認知,一切都是主所賜的,又有何可誇耀的呢?有沒有想過,我有機會過敬虔的生活,能有智慧,花時間去在信仰知識/生活中比別人投入更多,也全是主所賜的恩典,那麼我們將會明白何謂真正的謙卑…

世人期望卓越,隱藏弱點,爭取表現,為「成功」而盡上一切.

上帝期望我們做一個真實的人,內外一致,因為在祂面前沒有可隱藏的地方,忠於上帝給我們的一切,按我們所領受的而有相應的行動,

其實我們非常明白,「信仰只說不行,行嗎?」這問題的答案,但是我們需首先完全向主的依靠,因為信仰本身就不是靠自己力量去賺取或努力換來的東西,而是主的恩賜,所以行動的配合也由依靠施恩的主開始；接著需要勇氣,一份敢於面對真實自己的勇氣:再加上一點信心,把思想與行為連結的信心；與及一個願意與你同行,一同努力的群體,就是主的教會.

讓我們一同努力,一同在這窄路上奔跑.



\newpage

\section{那日子}
\label{sec:525}
\textbf{林豪恩先生 - 那日子}
\newline
\newline
\hyperref[sec:524]{< < < PREV SERMON < < <}
~
\hyperref[sec:toc]{[返目錄]}
~
\hyperref[sec:526]{> > > NEXT SERMON > > >}


阿摩司書 5:18-24
\newline
\begin{longtable}{cl}
\hline
\hline
章節 & 經文 (和合本修訂版)\\
\hline
5:18 & \begin{tabularx}{0.7\textwidth}{X} 想望耶和華日子的人有禍了！為甚麼你們要耶和華的日子呢？那是黑暗沒有光明的日子， \end{tabularx} \\ \\ \relax
5:19 & \begin{tabularx}{0.7\textwidth}{X} 好像人躲避獅子卻遇見熊；進房屋以手靠牆，卻被蛇咬。 \end{tabularx} \\ \\ \relax
5:20 & \begin{tabularx}{0.7\textwidth}{X} 耶和華的日子豈不是黑暗沒有光明，幽暗毫無光輝嗎？ \end{tabularx} \\ \\ \relax
5:21 & \begin{tabularx}{0.7\textwidth}{X} 「我厭惡你們的節期，也不喜悅你們的嚴肅會。 \end{tabularx} \\ \\ \relax
5:22 & \begin{tabularx}{0.7\textwidth}{X} 你們雖然向我獻燔祭和素祭，我卻不悅納，也不看你們用肥畜獻的平安祭。 \end{tabularx} \\ \\ \relax
5:23 & \begin{tabularx}{0.7\textwidth}{X} 要使你們歌唱的聲音遠離我，因為我不聽你們琴瑟的樂曲。 \end{tabularx} \\ \\ \relax
5:24 & \begin{tabularx}{0.7\textwidth}{X} 惟願公平如大水滾滾，公義如江河滔滔。 \end{tabularx} \\ \\
[1ex]
\hline
\hline
\end{longtable}
詩篇 70:1-5
\newline
\begin{longtable}{cl}
\hline
\hline
章節 & 經文 (和合本修訂版)\\
\hline
70:1 & \begin{tabularx}{0.7\textwidth}{X} 神啊，求你快快搭救我！耶和華啊，求你速速幫助我！ \end{tabularx} \\ \\ \relax
70:2 & \begin{tabularx}{0.7\textwidth}{X} 願那些尋索我命的，抱愧蒙羞；願那些喜悅我遭害的，退後受辱。 \end{tabularx} \\ \\ \relax
70:3 & \begin{tabularx}{0.7\textwidth}{X} 願那些對我說「啊哈、啊哈」的，因羞愧退後。 \end{tabularx} \\ \\ \relax
70:4 & \begin{tabularx}{0.7\textwidth}{X} 願所有尋求你的，因你歡喜快樂；願那些喜愛你救恩的，常說：「當尊神為大！」 \end{tabularx} \\ \\ \relax
70:5 & \begin{tabularx}{0.7\textwidth}{X} 但我是困苦貧窮的；神啊，求你速速到我這裡來！你是幫助我的，搭救我的；耶和華啊，求你不要耽延！ \end{tabularx} \\ \\
[1ex]
\hline
\hline
\end{longtable}
帖撒羅尼迦前書 4:13-18
\newline
\begin{longtable}{cl}
\hline
\hline
章節 & 經文 (和合本修訂版)\\
\hline
4:13 & \begin{tabularx}{0.7\textwidth}{X} 弟兄們，至於已睡了的人，我們不願意你們不知道，恐怕你們憂傷，像那些沒有指望的人一樣。 \end{tabularx} \\ \\ \relax
4:14 & \begin{tabularx}{0.7\textwidth}{X} 既然我們信耶穌死了，復活了，那些已經在耶穌裡睡了的人，神也必將他們與耶穌一同帶來。 \end{tabularx} \\ \\ \relax
4:15 & \begin{tabularx}{0.7\textwidth}{X} 我們照主的話告訴你們一件事：我們這活著還存留到主來臨的人，絕不會在那已經睡了的人之先。 \end{tabularx} \\ \\ \relax
4:16 & \begin{tabularx}{0.7\textwidth}{X} 因為，召集令一發，天使長的呼聲一叫，神的號角一吹，主必親自從天降臨；那在基督裡死了的人必先復活， \end{tabularx} \\ \\ \relax
4:17 & \begin{tabularx}{0.7\textwidth}{X} 然後我們這些活著還存留的人必和他們一同被提到雲裡，在空中與主相會。這樣，我們就要和主永遠同在。 \end{tabularx} \\ \\ \relax
4:18 & \begin{tabularx}{0.7\textwidth}{X} 所以，你們當用這些話彼此勸勉。 \end{tabularx} \\ \\
[1ex]
\hline
\hline
\end{longtable}
馬太福音 25:1-13
\newline
\begin{longtable}{cl}
\hline
\hline
章節 & 經文 (和合本修訂版)\\
\hline
25:1 & \begin{tabularx}{0.7\textwidth}{X} 「那時，天國好比十個童女拿著燈出去迎接新郎。 \end{tabularx} \\ \\ \relax
25:2 & \begin{tabularx}{0.7\textwidth}{X} 其中有五個是愚拙的，五個是聰明的。 \end{tabularx} \\ \\ \relax
25:3 & \begin{tabularx}{0.7\textwidth}{X} 愚拙的拿著燈，卻沒有帶油； \end{tabularx} \\ \\ \relax
25:4 & \begin{tabularx}{0.7\textwidth}{X} 聰明的拿著燈，又盛了油在器皿裡。 \end{tabularx} \\ \\ \relax
25:5 & \begin{tabularx}{0.7\textwidth}{X} 新郎遲延的時候，她們都打盹，睡著了。 \end{tabularx} \\ \\ \relax
25:6 & \begin{tabularx}{0.7\textwidth}{X} 半夜有人喊：『看，新郎來了，你們出來迎接他。』 \end{tabularx} \\ \\ \relax
25:7 & \begin{tabularx}{0.7\textwidth}{X} 那些童女就都起來挑亮她們的燈。 \end{tabularx} \\ \\ \relax
25:8 & \begin{tabularx}{0.7\textwidth}{X} 愚拙的對聰明的說：『請分點油給我們，因為我們的燈要滅了。』 \end{tabularx} \\ \\ \relax
25:9 & \begin{tabularx}{0.7\textwidth}{X} 聰明的回答：『恐怕不夠你我用的；你們還是自己到賣油的那裡去買吧。』 \end{tabularx} \\ \\ \relax
25:10 & \begin{tabularx}{0.7\textwidth}{X} 她們去買的時候，新郎到了。那預備好了的，與他進去共赴婚宴，門就關了。 \end{tabularx} \\ \\ \relax
25:11 & \begin{tabularx}{0.7\textwidth}{X} 其餘的童女隨後也來了，說：『主啊，主啊，給我們開門！』 \end{tabularx} \\ \\ \relax
25:12 & \begin{tabularx}{0.7\textwidth}{X} 他卻回答：『我實在告訴你們，我不認識你們。』 \end{tabularx} \\ \\ \relax
25:13 & \begin{tabularx}{0.7\textwidth}{X} 所以，你們要警醒，因為那日子，那時辰，你們不知道。」 \end{tabularx} \\ \\
[1ex]
\hline
\hline
\end{longtable}
講題:那日子 The Day

經文:阿摩司書5:18-24;詩篇70:1-5;帖撒羅尼迦前書4:13-18;馬太福音25:1-13

講員:林豪恩先生

場合:香港中文大學崇基學院禮拜堂主日崇拜

日期:2020年11月8日

引言

提到「那日子」,你會想到甚麼日子呢?由於世界各地封關防疫,復活節假期和暑假都不能外出旅行,此時此地,大人和小朋友共同期待的日子可能是「去旅行的日子」.不過,對於父母而言,對「那日子」的想像可能更多,例如:懷孕時期想像嬰兒出生的那日子；嬰兒出生後想像戒夜奶的那日子；識行識走之後想像上學的那日子；上學後想像畢業的那日子；長大成人之後想像成家立室的那日子；成家立室之後想像弄孫為樂的日子.這些都可能是人們經常期待的那日子.不過,不是所有日子都是人們期待的,有些日子倒是人們逃避的,人們甚至期望那些日子不要來,例如:大多數上班的人都不想星期一的到來；大多數學生都不想考試日的到來,大多數成人都不想衰老的日子來得那麼快；一般人都不想生命結束的那日子來得那麼快.我們逃避甚麼日子呢?我們盼望甚麼日子呢?我們希望甚麼日子快些來,又希望甚麼日子遲些來呢?這一切,一方面反映我們在期待甚麼,同時也反映我們當下的處境,讓我們透過今天的幾段經課思考自己處境及自己的期望,也思考與信仰的關係.

一.詩篇:渴望那日子快快來到

詩篇70篇的作者渴望那日子快快來到,為甚麼呢?他身處怎樣的狀況呢?他渴望甚麼呢?

70:1神啊,求你快快搭救我!耶和華啊,求你速速幫助我!

70:2願那些尋索我命的,抱愧蒙羞；願那些喜悅我遭害的,退後受辱.

70:3願那些對我說「啊哈,啊哈」的,因羞愧退後.

70:4願所有尋求你的,因你歡喜快樂；願那些喜愛你救恩的,常說:「當尊神為大!」

70:5但我是困苦貧窮的；神啊,求你速速到我這裡來!你是幫助我的,搭救我的；耶和華啊,求你不要耽延!

甚麼人渴望上主快快來到呢?如詩人一樣「困苦貧窮」的人可能是其中之一.從詩人的自述中可見,「困苦貧窮」不只是個人是智是愚,或勤或懶的問題,更有背後的社會和宗教條件.在詩人的處境中,有某種力量與他對抗,有一些人想消滅他.詩人表示自己是尋求神的,尊神為大的,暗示壓制他的對手並不尋求神,也不尊神為大.

詩人落入「困苦貧窮」的困境中,令我們想到世界上有些個人或群體,也因某種社會和宗教背景而變成「困苦貧窮」.從啟示錄我們得知第一世紀的基督徒因為信仰的緣故而被社會邊緣化.

啟示錄13章

16牠又使眾人,無論大小,貧富,自主的,為奴的,都在右手上,或是在額上,打一個印記；17這樣,除了那有印記,有獸的名或有獸名數字的,都不得買或賣.

神聽到「困苦貧窮」的呼求,神回應「困苦貧窮」的呼求,應許的那日子來了,拯救的那日子來了.如何知道神來了呢?如何得知那日子來到呢?昔日施洗約翰也問過這個問題.

馬太福音11章

2約翰在監獄裡聽見基督所做的事,就派他的門徒去,3問耶穌:「將要來的那位就是你嗎?還是我們要等候另一位呢?」4耶穌回答他們:「你們去,把所聽見,所看見的告訴約翰:5就是盲人看見,瘸子行走,痲瘋病人得潔淨,聾子聽見,死人復活,窮人聽到福音.」

如何辨認那日子來臨呢?原來「窮人聽到福音」竟然是神來到人間的其中一項記號.詩人說:「但我是困苦貧窮的」,所以詩人祈求:「神啊,求你速速到我這裡來!」因為詩人相信:「你是幫助我的,搭救我的」.「困苦貧窮」的人期望上主快來是可以理解的,對嗎?不知道你是否覺得自己「困苦貧窮」呢?無論如何,我們的世界中的確有不少「困苦貧窮」的人,教會能否把福音分享給他們呢?能否讓他們發現上主已經來臨人間呢?

二．亞摩司書:真想快來嗎?其實禍福未知

除了如詩人一樣「困苦貧窮」的人,原來還有另外一些人都期望耶和華的日子快些來到.不過,亞摩司先知卻向這些人「撥冷水」,耶和華的日子來到時,對他們到底是福是禍也說不定.為甚麼如此呢?這到底是怎樣一回事呢?

5:18想望耶和華日子的人有禍了!為甚麼你們要耶和華的日子呢?

亞摩司對他的聽眾說,你們最好期望耶和華的日子不要那麼快來到,因為對你們來說,耶和華的日子並不是你們想像一般美好的.

那是黑暗沒有光明的日子,

5:19好像人躲避獅子卻遇見熊；進房屋以手靠牆,卻被蛇咬.

5:20耶和華的日子豈不是黑暗沒有光明,幽暗毫無光輝嗎?

人人都以為耶和華的日子是光明的,原來可以是黑暗的；人人都以為耶和華的日子是拯救,原來可以是禍害.原來,耶和華的日子是因人而異的.

5:21「我厭惡你們的節期,也不喜悅你們的嚴肅會.

5:22你們雖然向我獻燔祭和素祭,我卻不悅納,也不看你們用肥畜獻的平安祭.

5:23要使你們歌唱的聲音遠離我,因為我不聽你們琴瑟的樂曲.

最令人感到意外的,是亞摩司的聽眾並不是外邦人,也不是異教徒,他們甚至是在某種意義上虔誠的耶和華敬拜者.他們勤守節期,也辦嚴肅會；他們獻燔祭和素祭,甚至慷慨地用肥畜獻的平安祭；他們鼓瑟彈琴,唱歌讚頌,所有宗教儀式都做足了,所有宗教行為都遵守了,而且他們也是如此相信,因此他們「想望耶和華日子」.亞摩司先知指出,這若不是非常嚴重的錯誤,就是非常嚴重的誤會.

5:24惟願公平如大水滾滾,公義如江河滔滔.

原來耶和華寧願要公平和公義,而不是祭祀和守節期.亞摩司的當頭棒喝不單對當時的聽眾適用,不幸地,也適用於歷史中出現過的不少信徒群體.

回想上一段詩人訴苦說自己是「困苦貧窮」的人,「困苦貧窮」這現象何以出現呢?其中可能包括個人因素,也難免有社會因素.亞摩司書明言,不公平和不公義的社會是「困苦貧窮」的背景,因此宣告上帝「惟願公平如大水滾滾,公義如江河滔滔」.

對於阿摩司的聽眾來說,那日子最好不要來得那麼快,遲些來更好,讓誤會了的人有時間改正,讓走了錯路的人有機會重回正軌.

三．馬太福音:那日子其實你們不知道,你們要知道的是要警醒

如果老師告訴學生,延遲一個月才交功課,有些學生會很興奮.興奮,不一定因為有更多時間做功課,而是不需要現在就開始做功課,可以繼續做自己喜歡做的事情.其實,老師在開學時派的「course outline」已經通知了學期尾交功課的日期,讓學生有多些時間準備,不過,大部分人都是「deadline fighters」,只會在deadline前才趕工.耶穌的天國比喻讓我們聯想到,上主彷彿向人派發了course outline,有目標,有內容,有要求,也有deadline,只是沒說清楚deadline是何年何月何日.如此的安排,對人有甚麼意義呢?人又應如何回應呢?

25:1「那時,天國好比十個童女拿著燈出去迎接新郎.

25:2其中有五個是愚拙的,五個是聰明的.

25:3愚拙的拿著燈,卻沒有帶油；

25:4聰明的拿著燈,又盛了油在器皿裡.

25:5新郎遲延的時候,她們都打盹,睡著了.

25:6半夜有人喊:『看,新郎來了,你們出來迎接他.』

25:7那些童女就都起來挑亮她們的燈.

25:8愚拙的對聰明的說:『請分點油給我們,因為我們的燈要滅了.』

25:9聰明的回答:『恐怕不夠你我用的；你們還是自己到賣油的那裡去買吧.』

25:10她們去買的時候,新郎到了.那預備好了的,與他進去共赴婚宴,門就關了.

25:11其餘的童女隨後也來了,說:『主啊,主啊,給我們開門!』

25:12他卻回答:『我實在告訴你們,我不認識你們.』

25:13所以,你們要警醒,因為那日子,那時辰,你們不知道.」

或遲或早,新郎總會來的,這方面童女們是沒有猶豫的.正如或遲或早,人生總有deadline,總有面對上主的一刻,這方面我們也不需要動搖.如果deadline訂明了日子,人們尚可以做deadline fighters.但如果「那日子,那時辰,你們不知道」,耶穌認為人們唯一應該持守的態度是「你們要警醒」.警醒是甚麼意思呢?相信警醒並非不去睡覺,警醒對阿摩司的聽從來說是從不公平和不公義的行徑中清醒過來,秉公行義,使公平和公義好像大水江河一樣充滿全地.

四．帖撒羅尼迦前書:無論遲早,最終來到

書信經課的帖撒羅尼迦前書有關於那日子的描述:

4:13弟兄們,至於已睡了的人,我們不願意你們不知道,恐怕你們憂傷,像那些沒有指望的人一樣.

4:14既然我們信耶穌死了,復活了,那些已經在耶穌裡睡了的人,神也必將他們與耶穌一同帶來.

4:15我們照主的話告訴你們一件事:我們這活著還存留到主來臨的人,絕不會在那已經睡了的人之先.

4:16因為,召集令一發,天使長的呼聲一叫,神的號角一吹,主必親自從天降臨；那在基督裡死了的人必先復活,

4:17然後我們這些活著還存留的人必和他們一同被提到雲裡,在空中與主相會.這樣,我們就要和主永遠同在.

4:18所以,你們當用這些話彼此勸勉.

那日子似乎是關於未的,那麼,對我們仍然在世間存活的人有甚麼意義呢?最後一句說:「所以,你們當用這些話彼此勸勉」,勸勉甚麼呢?4:13節指出「至於已睡了的人,我們不願意你們不知道,恐怕你們憂傷,像那些沒有指望的人一樣」.對於大部分人來說,生離死別的憂傷總會過去的,就算憂傷不滅,人生仍然要活下去,最多是帶著憂傷活下去.對於活下去的人,切身的問題是如何在等候那日子之下而活下去呢?回答這個問題,我們不妨參考這個段落之前的兩句,保羅如此教導:

4:11要立志過安靜的生活,管自己的事,親手做工,正如我們從前吩咐你們的,

4:12好使你們的行為能得外人的尊敬,同時也不依賴任何人.

總括而言,是正正當當地生活,負責任地生活,有尊嚴地生活.這會不會也可以成為耶穌吩咐門徒「你們要警醒」的註腳呢?

總結:一位中學生期待的日子

「那日子」原來可以給人眾多的想像,不同的描述.你對於「那日子」的想像又是怎樣的呢?有一位中學生對那日子的期待是這樣的:

我們期待的日子是花

每一縷芬芳都充盈我們銳敏的呼吸

把美和善良,淳樸和真誠灑進花裡

在土地的懷裡吮吸著雨露甘霖的精華

探出小巧的微笑

向太陽問好

在風雨中茁壯成長

帶來每一個空氣清新的早晨

我們期待的日子是樹

枝繁葉茂

強壯地紮根土地

讓我們像棵最偉岸高潔的樹

站成一道風景

在朝霞暮靄中

詩意地遠望

人生至美的燈塔

我們期待的日子是天空

遼闊的襟懷是飛鳥和夢想的天堂

我們的天空一碧千里

每一個善良願望在天空的溫柔中悄聲歌唱

比如風箏　比如童話　比如托著腮幫想少女的祈禱:

寒流不再來　世情不再冷漠　眼睛不再塵封

最美不過夕陽紅

讓每一步路都走出最甜蜜的回音

我們期待的日子是眼淚

在黑亮,真誠的眼裡流出

或者來自懺悔和過錯的心靈

盡情沖刷,滌蕩

讓每一雙眼睛都清亮如同童心或鏡子

每一顆心靈潔淨如同白金

每一副胸襟都廣闊如海

每一個腳印都堅實如山

我們期待的日子是詩篇,是音樂

在詩篇的聖殿裡我們無比和美

在我們的眼睛裡閃著友善的光

在我們的詩篇裡看到熱情的手

讓我們用最美麗的詩篇奏出動聽的音樂

慰籍每一顆心靈

充滿天空　舖滿大地

(來自網上中學生作品選)

或者有人會認為那位中學生過分天真和浪漫,不過,期待美好難道有甚麼妥當嗎?會否只是人們長大了,在眾多社會條件之下已經把期待美好的天真和浪漫消滅了,或者不敢再期待美好.基督徒對那日子的期待或者在某些人的眼中都是過分天真和浪漫,然而,基督徒不但維持對那日子的期待,更盼望那日子出現在當下,正如主禱文所言:「願你的旨意行地地上,如同行在天上」.在剛過去的週五學院週會,有人問演講嘉賓,學生如何善用中大這優美的校園幫助自己達至內心的平靜.講員回應說,外在環境是不重要的,真正重要的是你的心,若你修煉,你在任何環境之中也可以平靜.基督徒對那日子的盼望並不在於自己的心,也不在於修煉,我們的心和我們的修煉並不能把那日子帶來.基督徒也說在任何環境之中也可以存有盼望,因為耶穌基督道成肉身就已經把那日子帶來了,並且呼喚我們參與祂已經開展的國度,直至祂再次來臨把關於那日子的一切完滿實現.



\newpage

\section{主啊你世世代代作我們的避難所}
\label{sec:526}
\textbf{鄧瑞強博士 - 主啊,你世世代代作我們的避難所}
\newline
\newline
\hyperref[sec:525]{< < < PREV SERMON < < <}
~
\hyperref[sec:toc]{[返目錄]}
~
\hyperref[sec:527]{> > > NEXT SERMON > > >}


講題:主啊,你世世代代作我們的避難所

經文:詩篇90

講員:鄧瑞強博士

場合:香港中文大學崇基學院禮拜堂主日崇拜

日期:2020年11月15日

各位弟兄姊妹,早安,願各人平安.

今天,我要講解詩篇90篇.一般而言,讀詩篇的時候,我們習慣是將150篇詩篇看成是獨立的篇章,一篇一篇獨立地理解.這一篇詩篇講的東西,與下一篇詩篇講的東西,沒有什麼關連.事實是否是這樣?

二十世紀,聖經學者發展出一種理解聖經的方法,英文叫「Canonical Criticism」,中文叫「正典批判」,用最簡單的話來講,這方法是從經文的總體編排去理解個別的經文,而不是將經文分割開來作微觀的解剖.將這方法應用在詩篇研究上,就是從詩篇的總體編排去理解個別詩篇的內容及用意.今天,我們試試用這方法去理解詩篇90篇,看看會否令我們更深理解這詩篇?

詩篇是經過後人編輯的,經過篩選,編排,意義重組,然後,將詩篇結集成五卷.分成五卷,看來,是刻意與摩西五經對應.在中國,我們也有孔子刪詩的講法.孔子將古代流下來的3000多首詩歌,刪減並編輯為《詩三百》,成為留傳後世的《詩篇》.大家想想,經編輯後的《詩篇》,有沒有可能沒有孔子思想的貫注而將之改造成儒家思想的出口?聖經裡的詩篇,作為一經過編輯的結集,也內貫了編輯者的用心及期盼.

我先讀一讀今日的詩篇經課:詩篇90篇(詩篇卷四的第一首)

詩90:1　〔神人摩西的祈禱〕主啊,你世世代代作我們的居所(或者譯:「避難所」).

詩90:2　諸山未曾生出,地與世界你未曾造成,從亙古到永遠,你是神.

詩90:3　你使人歸於塵土,說:你們世人要歸回.

詩90:4　在你看來,千年如已過的昨日,又如夜間的一更.

詩90:5　你叫他們如水沖去；他們如睡一覺.早晨,他們如生長的草,

詩90:6　早晨發芽生長,晚上割下枯乾.

詩90:7　我們因你的怒氣而消滅,因你的忿怒而驚惶.

詩90:8　你將我們的罪孽擺在你面前,將我們的隱惡擺在你面光之中.

詩90:9　我們經過的日子都在你震怒之下；我們度盡的年歲好像一聲歎息.

詩90:10　我們一生的年日是七十歲,若是強壯可到八十歲；但其中所矜誇的不過是勞苦愁煩,轉眼成空,我們便如飛而去.

詩90:11　誰曉得你怒氣的權勢?誰按著你該受的敬畏曉得你的忿怒呢?

詩90:12　求你指教我們怎樣數算自己的日子,好叫我們得著智慧的心.

詩90:13　耶和華啊,我們要等到幾時呢?求你轉回,為你的僕人後悔.

詩90:14　求你使我們早早飽得你的慈愛,好叫我們一生一世歡呼喜樂.

詩90:15　求你照著你使我們受苦的日子,和我們遭難的年歲,叫我們喜樂.

詩90:16　願你的作為向你僕人顯現；願你的榮耀向他們子孫顯明.

詩90:17　願主─我們神的榮美歸於我們身上.願你堅立我們手所做的工；我們手所做的工,願你堅立.

要初步了解這詩篇,需要留意三點.首先,要知道這是詩篇第四卷的開頭；其次,這是全部詩篇中,唯一一次在詩篇的標題中標示「摩西」；第三,稱「摩西」為「神人摩西」,這是「神人摩西的祈禱」.

在整部詩篇的佈局中,頭三卷詩篇集中講大衛的王朝.大衛王朝每況愈下,後來的王拜偶像,行不義,終被神懲罰,被巴比倫所滅,聖殿被毀,人民被擄去巴比倫,正是國破家亡.這光景反映在詩篇第三卷的最後一首詩內,即詩篇89篇.這光景如何,我們看看詩篇89篇怎樣說.

詩89:38　但你惱怒你的受膏者,就丟掉棄絕他.

詩89:39　你厭惡了與僕人所立的約,將他的冠冕踐踏於地.

詩89:40　你拆毀了他一切的籬笆,使他的保障變為荒場.

詩89:41　凡過路的人都搶奪他；他成為鄰邦的羞辱.

詩89:42　你高舉了他敵人的右手；你叫他一切的仇敵歡喜.

詩89:43　你叫他的刀劍捲刃〔按:即刀鋒向自己〕,叫他在爭戰之中站立不住.

詩89:44　你使他的光輝止息,將他的寶座推倒於地.

詩89:45　你減少他青年的日子,又使他蒙羞.

詩89:46　耶和華啊,這要到幾時呢?你要將自己隱藏到永遠嗎?你的忿怒如火焚燒要到幾時呢?

詩89:47　求你想念我的時候是何等的短少；你創造世人,要使他們歸何等的虛空呢?

詩89:48　誰能常活免死,救他的靈魂脫離陰間的權柄呢?

詩89:49　主啊,你從前憑你的信實向大衛立誓要施行的慈愛在哪裡呢?

詩89:50　主啊,求你記念僕人們所受的羞辱,記念我怎樣將一切強盛民的羞辱存在我懷裡.

詩89:51　耶和華啊,你的仇敵用這羞辱羞辱了你的僕人,羞辱了你受膏者的腳蹤.

詩89:52　耶和華是應當稱頌的,直到永遠.阿們!阿們!〔這是編輯標記,標示卷三的結束.〕

大家都看到了,卷三的結尾,就是大衛王朝的結尾.大衛王朝的結束,標示著一個價值真空的時代的到來.再沒有王去引導社會,再沒有聖殿可讓人去寄托心靈,再沒有屬於自己的土地而被迫留放他鄉.過去的生活,一去不復返；面對前路,只是一片空白.

在這「沒有王,沒有聖殿,沒有土地」的真空時刻,編輯詩篇的聖徒們,他們想起過去,想起有一個時刻,他們同樣沒有王,沒有聖殿,沒有土地,但他們有神,有神的同在,有神給予的希望,那是摩西的時代.那個時代,神的子民在曠野.他們無飯吃,神降天糧；沒水飲,神讓磐石流出泉水；沒有土地,神讓他們望向「應許之地」,他們有前景,有希望.他們像是什麼都沒有,但他們有神,這已是最豐盛的生命.

於是,在詩篇卷三結尾的絕望裡,編輯們以摩西的詩作為卷四的開場白,去回應卷三結尾時留下的絕望.從時序而言,將摩西的詩編在大衛的詩之後,是不合時序的.但在大衛王朝崩壞之後,言說摩西,言說一沒有王朝卻有上帝的時代,是合乎神學的.來到這裡,大家應該明白,若不割裂地,獨立地看詩篇90篇,而從詩篇整體的佈局去看,則詩篇90篇並不是個人感悟生命之易逝而渴求個體幸福的詩篇,而是在國破家亡的日子裡渴求民族的復興,歷史重回正軌.

編輯者將這詩標示為「神人摩西的祈禱」,這是編輯者留下的神學路標,指引我們應如何去理解這詩篇.在整本聖經裡,稱摩西為「神人摩西」的,只有幾處,而又連於他的祈禱的,則只有一處,就是摩西帶領以色列人快將進入「應許之地」的前夕,他籌數即盡,便為他的同胞祈禱,這祈禱有警告,也帶來希望.這「神人摩西的祈禱」記載在申命記32章至33章.我們用經文互讀的方法(intertextuality),藉申命記32章至33章的經文,去理解詩篇90篇的含意.

詩90:1　〔神人摩西的祈禱〕主啊,你世世代代作我們的居所(或者譯:「避難所」).

要了解這經文,必需對照「神人摩西」在申命記33:27-29的祈禱:

申33: 27　永生的神是你的居所；他永久的膀臂在你以下.他在你前面攆出仇敵,說:毀滅吧.

申33: 28　以色列安然居住；雅各的本源獨居五穀新酒之地.他的天也滴甘露.

申33: 29　以色列啊,你是有福的!誰像你這蒙耶和華所拯救的百姓呢?他是你的盾牌,幫助你,是你威榮的刀劍.你的仇敵必投降你；你必踏在他們的高處.

這裡講到,神是我們的居所；也提及,神克服仇敵.這正好回應詩篇89篇提及的困局,在那裡,提到神懲罰祂的子民,他們被交予仇敵,國破家亡.

如此,詩篇90:1描繪的,就不是一個買不到樓的香港人的渴望:「主啊,你世世代代作我們的居所」,而是一群望著前景而感到無望的人,那些被擄異鄉的人的渴望:「主啊,你世世代代作我們的居所」.無家可歸,以神為家.

詩90:2　諸山未曾生出,地與世界你未曾造成,從亙古到永遠,你是神.

詩90:3　你使人歸於塵土,說:你們世人要歸回.

詩90:4　在你看來,千年如已過的昨日,又如夜間的一更.

詩90:5　你叫他們如水沖去；他們如睡一覺.早晨,他們如生長的草,

詩90:6　早晨發芽生長,晚上割下枯乾.

90:2講到「山」和「地」,講到「你是神」.當中提到諸山的「生出」,地和世界的「造成」,用字都是「生育」的用字,神像母親生產那樣創造大地.90:5講到生命如草,被大水沖去.表面看,這是哀歎人生的短暫.

我們對照一下「神人摩西」在申命記32章的祈禱:

申32:22　因為在我怒中有火燒起,直燒到極深的陰間,把地和地的出產盡都焚燒,山的根基也燒著了.

這裡也提到「山」和「地」,但這已不再是人生活的地方,而是遭咒詛的地方,被火焚燒的地方.

申32:3　我要宣告耶和華的名；你們要將大德歸與我們的神.

問題是:神的子民已不懂宣告神的名,以致在人生中不知所措.

申32:18　你輕忽生你的磐石,忘記產你的神.

這裡也講到神像母親生產那樣創造大地和我們,問題是我們輕忽了神,忘記了神.

申32:2　我的教訓要淋漓如雨；我的言語要滴落如露,如細雨降在嫩草上,如甘霖降在菜蔬中.

人生雖然如草,但這草被神的真理滋潤,被細雨輕撫.草,不一定是短暫的.但當草離開真理,便被「大水」沖去了.

對照著申命記32章的經文來看,我們明白到人生的短暫不是一生物學事件,而是一信仰事件.詩篇講的,其實是人的不信所導致的淒慘後果.

問題出在哪裡?

詩90:7-11點出問題所在.

這段經文是典型的平行經文:90:7對應90:11,講到「神的忿怒」；90:8對應90:10,講到「人的可悲」；90:9將「神的忿怒」與「人的可悲」合在一起.

「神的忿怒」

詩90:7　我們因你的怒氣而消滅,因你的忿怒而驚惶.

詩90:11　誰曉得你怒氣的權勢?誰按著你該受的敬畏曉得你的忿怒呢?

「人的可悲」

詩90:8　你將我們的罪孽擺在你面前,將我們的隱惡擺在你面光之中.

詩90:10　我們一生的年日是七十歲,若是強壯可到八十歲；但其中所矜誇的不過是勞苦愁煩,轉眼成空,我們便如飛而去.

「神的忿怒」與「人的可悲」

詩90:9　我們經過的日子都在你震怒之下；我們度盡的年歲好像一聲歎息.

表面看,這是講個人的罪咎體驗,講人生的短暫與悲哀.

我們對照一下「神人摩西」在申命記32章的祈禱:

申32:16　敬拜別神,觸動神的憤恨,行可憎惡的事,惹了他的怒氣.

申32:19　耶和華看見他的兒女惹動他,就厭惡他們,說:

申32:20　我要向他們掩面,看他們的結局如何.他們本是極乖僻的族類,心中無誠實的兒女.

申32:21　他們以那「不算為神」的觸動我的憤恨,以虛無的神惹了我的怒氣.我也要以那「不成子民」的觸動他們的憤恨,以愚昧的國民惹了他們的怒氣.

神的怒氣,是由於我們不再敬畏祂.我們追逐世間的成功,將世人或世物「神化」,「偶像化」,向這些偶像下拜；也因敬拜這些偶像而令生命走入歧途.

詩篇90篇講的,不是個人的迷失與悲哀,而是神整體子民的迷失與悲哀.

詩90:11　誰曉得你怒氣的權勢?誰按著你該受的敬畏曉得你的忿怒呢?

詩90:12　求你指教我們怎樣數算自己的日子,好叫我們得著智慧的心.

這裡提及「數算自己的日子」,「得著智慧的心」.問題是:數算自己的日子來作什麼?什麼才是智慧?

我們對照一下「神人摩西」在申命記32章的祈禱:

申32:6　愚昧無知的民哪,你們這樣報答耶和華嗎?他豈不是你的父,將你買來的嗎?他是製造你,建立你的.

申32:7　你當追想上古之日,思念歷代之年；問你的父親,他必指示你；問你的長者,他必告訴你.

智慧的相反便是「愚昧無知」,「愚昧無知」在於我們不再相信神是創造天地的主,神是歷史的主,神是救贖的主.我們以為可以用自己的手創造天地,可以用自己的方式扭轉歷史,可以自己救贖自己.這就是「愚昧無知」.

與「數算自己的日子」相對的,是「追想上古之日,思念歷代之年」.我們的日子短少,但神的作為世代無盡.在世世代代裡,我們見到的,是神的救贖作為,是神聆聽人的哀聲,是神解救受苦的人.如此,「數算自己的日子」,為的是看到自己的微少,而仰賴在悠長歷史中工作的神.

詩90:13　耶和華啊,我們要等到幾時呢?求你轉回,為你的僕人後悔.

詩90:14　求你使我們早早飽得你的慈愛,好叫我們一生一世歡呼喜樂.

詩90:15　求你照著你使我們受苦的日子,和我們遭難的年歲,叫我們喜樂.(《現代中文譯本》:你使我們遭受多少年日的苦難,求你也使我們享受多少年日的喜樂.)

詩90:16　願你的作為向你僕人顯現；願你的榮耀向他們子孫顯明.

「轉回」這個字,有「悔改」的含意.通常,是人向神「轉回」,詩90:3說「你們世人要歸回(按:同「轉回」)」,就有這種含意.現在,是求神「轉回」,求神「為你的僕人後悔」.

我們對照一下「神人摩西」在申命記32章的祈禱:

申32:36　耶和華見他百姓毫無能力,無論困住的,自由的都沒有剩下,就必為他們伸冤,為他的僕人後悔(按:與詩90:13中的「後悔」是同一個字).

這裡講的「神的後悔」的意思是很清楚的,是神因祂與人立的約,因祂的仁慈與憐憫,不再按罪人應得的刑罰降罰於人,而以慈愛待他們.簡單說,所謂「為祂的僕人後悔」,就是求神收回懲罰的手.

現在,要回答詩90:13的問題:「耶和華啊,我們要等到幾時呢?」這問題到底求什麼,單就詩篇90篇而論,是不清楚的.但詩篇90篇若是回應詩篇89篇的被擄結局,並在申命記32章提供的亮光下,我們知道,「耶和華啊,我們要等到幾時呢?」這問題是問:我們被擄受苦,要等多久呢?受苦的日子,好像無限長.詩人求的,是喜樂的日子,也這樣長.

詩90:17　願主─我們神的榮美歸於我們身上.願你堅立我們手所做的工；我們手所做的工,願你堅立.

這是這首詩的結尾.問題是:「我們手所做的工」是什麼樣的工,以至值得神去堅立.

我們對照一下「神人摩西」在申命記32章的祈禱:

申32:46　又說:「我今日所警教你們的,你們都要放在心上；要吩咐你們的子孫謹守遵行這律法上的話.

申32:47　因為這不是虛空,與你們無關的事,乃是你們的生命；在你們過約旦河要得為業的地上必因這事日子得以長久.」

原來,「我們手所做的工」,便是「謹守遵行這律法上的話」,如此,神的子民便能克服生命的短暫,活得長久,活出永恆的意義.

詩篇90篇,對我們有何意義呢?

這詩篇提示我們,在不如意的世界中,在短暫的人生裡,千萬不要忘記神,也不要忘記神的真理,要努力奉行.相信神的仁慈,仰賴神的解救,盼望神賜的將來,我們一起往前行.希望有一天,我們一起說,在這短暫的人生,我們活出了永恆.

願榮耀歸上主,願人間享和平.



\newpage

\section{切膚之痛}
\label{sec:527}
\textbf{王玉教授 - 切膚之痛}
\newline
\newline
\hyperref[sec:526]{< < < PREV SERMON < < <}
~
\hyperref[sec:toc]{[返目錄]}
~
\hyperref[sec:528]{> > > NEXT SERMON > > >}


馬太福音 25:31-46
\newline
\begin{longtable}{cl}
\hline
\hline
章節 & 經文 (和合本修訂版)\\
\hline
25:31 & \begin{tabularx}{0.7\textwidth}{X} 「當人子在他榮耀裡，同著眾天使來臨的時候，要坐在他榮耀的寶座上。 \end{tabularx} \\ \\ \relax
25:32 & \begin{tabularx}{0.7\textwidth}{X} 萬民都要聚集在他面前。他要把他們分別出來，好像牧人分別綿羊、山羊一般， \end{tabularx} \\ \\ \relax
25:33 & \begin{tabularx}{0.7\textwidth}{X} 把綿羊安置在右邊，山羊在左邊。 \end{tabularx} \\ \\ \relax
25:34 & \begin{tabularx}{0.7\textwidth}{X} 於是王要向他右邊的說：『你們這蒙我父賜福的，可來承受那創世以來為你們所預備的國。 \end{tabularx} \\ \\ \relax
25:35 & \begin{tabularx}{0.7\textwidth}{X} 因為我餓了，你們給我吃；渴了，你們給我喝；我流浪在外，你們留我住； \end{tabularx} \\ \\ \relax
25:36 & \begin{tabularx}{0.7\textwidth}{X} 我赤身露體，你們給我穿；我病了，你們看顧我；我在監獄裡，你們來看我。』 \end{tabularx} \\ \\ \relax
25:37 & \begin{tabularx}{0.7\textwidth}{X} 義人就回答：『主啊，我們甚麼時候見你餓了，給你吃；渴了，給你喝？ \end{tabularx} \\ \\ \relax
25:38 & \begin{tabularx}{0.7\textwidth}{X} 甚麼時候見你流浪在外，留你住；或是赤身露體，給你穿？ \end{tabularx} \\ \\ \relax
25:39 & \begin{tabularx}{0.7\textwidth}{X} 又甚麼時候見你病了，或是在監獄裡，來看你呢？』 \end{tabularx} \\ \\ \relax
25:40 & \begin{tabularx}{0.7\textwidth}{X} 王回答他們說：『我實在告訴你們，這些事你們做在我弟兄中一個最小的身上，就是做在我身上了。』 \end{tabularx} \\ \\ \relax
25:41 & \begin{tabularx}{0.7\textwidth}{X} 「王又要向那左邊的說：『你們這被詛咒的人，離開我！進入那為魔鬼和他的使者所預備的永火裡去！ \end{tabularx} \\ \\ \relax
25:42 & \begin{tabularx}{0.7\textwidth}{X} 因為我餓了，你們沒有給我吃；渴了，你們沒有給我喝； \end{tabularx} \\ \\ \relax
25:43 & \begin{tabularx}{0.7\textwidth}{X} 我流浪在外，你們沒有留我住；我赤身露體，你們沒有給我穿；我病了，我在監獄裡，你們沒有來看顧我。』 \end{tabularx} \\ \\ \relax
25:44 & \begin{tabularx}{0.7\textwidth}{X} 他們也要回答：『主啊，我們甚麼時候見你餓了，或渴了，或流浪在外，或赤身露體，或病了，或在監獄裡，沒有伺候你呢？』 \end{tabularx} \\ \\ \relax
25:45 & \begin{tabularx}{0.7\textwidth}{X} 王要回答：『我實在告訴你們，這些事你們沒有做在任何一個最小的弟兄身上，就是沒有做在我身上了。』 \end{tabularx} \\ \\ \relax
25:46 & \begin{tabularx}{0.7\textwidth}{X} 這些人要往永刑裡去；那些義人要往永生裡去。」 \end{tabularx} \\ \\
[1ex]
\hline
\hline
\end{longtable}
以西結書 34:11-16,20-24
\newline
\begin{longtable}{cl}
\hline
\hline
章節 & 經文 (和合本修訂版)\\
\hline
34:11 & \begin{tabularx}{0.7\textwidth}{X} 「主耶和華如此說：『看哪，我必親自尋找我的羊，將牠們尋見。 \end{tabularx} \\ \\ \relax
34:12 & \begin{tabularx}{0.7\textwidth}{X} 牧人在羊群四散的日子怎樣尋找他的羊，我必照樣尋找我的羊。這些羊在密雲黑暗的日子散在各處，我要從那裡救回牠們。 \end{tabularx} \\ \\ \relax
34:13 & \begin{tabularx}{0.7\textwidth}{X} 我要從萬民中領出牠們，從各國聚集牠們，引領牠們歸回故土。我要在以色列山上，在一切溪水旁邊，在境內所有可居住的地牧養牠們。 \end{tabularx} \\ \\ \relax
34:14 & \begin{tabularx}{0.7\textwidth}{X} 我要在肥美的草場牧養牠們。牠們的圈必在以色列高處的山上，牠們必躺臥在佳美的圈內，在以色列山肥美的草場上吃草。 \end{tabularx} \\ \\ \relax
34:15 & \begin{tabularx}{0.7\textwidth}{X} 我要親自牧養我的群羊，使牠們得以躺臥。這是主耶和華說的。 \end{tabularx} \\ \\ \relax
34:16 & \begin{tabularx}{0.7\textwidth}{X} 失喪的，我必尋找；被逐的，我必領回；受傷的，我必包紮；有病的，我必醫治；只是肥的壯的，我要除滅；我必秉公牧養牠們。』 \end{tabularx} \\ \\ \relax
34:17 & \begin{tabularx}{0.7\textwidth}{X} 「我的羊群哪，論到你們，主耶和華如此說：看哪，我要在羊與羊中間、公綿羊與公山羊中間施行審判。 \end{tabularx} \\ \\ \relax
34:18 & \begin{tabularx}{0.7\textwidth}{X} 你們在肥美的草場上吃草還以為是小事嗎？竟用你們的腳踐踏剩下的草；你們喝了清水，竟用你們的腳攪渾剩下的水。 \end{tabularx} \\ \\ \relax
34:19 & \begin{tabularx}{0.7\textwidth}{X} 至於我的羊，只能吃你們所踐踏的，喝你們所攪渾的。 \end{tabularx} \\ \\ \relax
34:20 & \begin{tabularx}{0.7\textwidth}{X} 「所以，主耶和華對牠們如此說：看哪，我要親自在肥羊和瘦羊中間施行審判。 \end{tabularx} \\ \\ \relax
34:21 & \begin{tabularx}{0.7\textwidth}{X} 因為你們用側邊用肩推擠一切瘦弱的羊，又用角牴撞，使牠們四散在外； \end{tabularx} \\ \\ \relax
34:22 & \begin{tabularx}{0.7\textwidth}{X} 所以，我要拯救我的群羊，牠們必不再作掠物；我也要在羊和羊中間施行審判。 \end{tabularx} \\ \\ \relax
34:23 & \begin{tabularx}{0.7\textwidth}{X} 我必在他們之上立一牧人，就是我的僕人大衛，牧養牠們；他必牧養他們，作他們的牧人。 \end{tabularx} \\ \\ \relax
34:24 & \begin{tabularx}{0.7\textwidth}{X} 我---耶和華必作他們的神，我的僕人大衛要在他們中間作王。這是我---耶和華說的。 \end{tabularx} \\ \\
[1ex]
\hline
\hline
\end{longtable}
講題:切膚之痛 (Keenly Felt Pain)

經文:馬太福音25:31-46;以西結書34:11-16,20-24

講員:王玉教授

場合:香港中文大學崇基學院禮拜堂主日崇拜

日期:2020年11月22日

I.引言

今天是基督君王主日(日即將臨前主日),我們將會從馬太福音25:31-46看一個不一樣,超乎想像的的君王形象,即一個與受苦者和被壓迫者認同的君王.

馬太福音的寫作對像,通常被學者聖經學者稱為馬太群體 (Matthean Community) ,是一班以猶太信徒為主的信仰群體,但成員也包括一些外邦信徒.基於馬太福音的猶太背景,馬太福音的作者在福音書引用或暗示了很多猶太經卷的內容,當中包括了不少先知書的內容.

本主日福音書的引題 (馬太福音25:31-46)是一段有關於「終末審判」的經文.經文不但隱晦地指向以西結書第34章中牧人與羊比喻,經文中所記載「分綿羊,山羊」和「人子在他榮耀寶座上施行審判」的一段經文更是由以西結書34:17-22發展出來.基於馬太群體的猶太背景,無用置疑,馬太群體的讀者 (或聽眾) 對猶太經卷熟識. 在誦讀這段經文時,他們必然先入為主地受到以西結書第34章影響去解讀馬太福音25:31-46的內容.稍後,我們會以以西結書第34章的視野和馬太群體身處羅馬帝國時代巴勒斯坦的處境去解讀馬太福音25:31-46,並看看兩段經文的互文性(intertextuality) 如何豐富我們對馬太福音25:31-46的解讀,嘗試就此破解馬太福音25:31-46的一些詮釋的謎團.

II.基督為終未審判中的「人子」,「牧人」及「君王」

馬太福音25:31-46 是福音書在耶穌踏上受苦之路前最後一個記載的耶穌教訓,內容帶著很強烈的終末論的元素 (eschatological features).經文開場白是:「當人子在他榮耀裡,同著眾天使來臨的時候,要坐在他榮耀的寶座上.萬民都要聚集在他面前.他要把他們分別出來,好像牧人分別綿羊,山羊一般,把綿羊安置在右邊,山羊在左邊.於是王要向他右邊的說:『你們這蒙我父賜福的,可來承受那創世以來為你們所預備的國.』」

這段經文所描述是一個「終末審判」(The Last Judgement)的情境.其中,作者以三個不同的稱謂形容耶穌基督:「人子」,「牧人」,「王」.

(1)「人子」是猶太天啟文學 ( Apocalyptic Literature) 中所期待在終末時期 (The End Time) 出現審判世上政權的彌賽亞,也是但以理書7:13-14所指的「人子」:「我在夜間的異象中觀看,見有一位像人子的,駕著天雲而來,被領到亙古常在者面前,得了權柄,榮耀,國度,使各方,各國,各族的人都事奉他.他的權柄是永遠的,不能廢去；他的國必不敗壞.」

天啟文學常見的主題﹐包括「命運逆轉」,「對行惡者的終極報應」 (參太16:27),「從混亂恢復宇宙秩序」等,都在這段經文中出現.

(2)「牧人」是用作比喻審判者進行判斷,分別義人和行不義人 (即經文所指「被詛咒的人」)的方法,有如牧人將綿羊和山羊分開兩邊一樣.這的「牧人」的比喻是取材於以西結書34:17:「我的羊群哪,論到你們,主耶和華如此說:我必在羊與羊中間,公綿羊與公山羊中間施行判斷.」及剛才所讀的20至22節.但是,以西結書所分別出來的「公錦羊與公山羊」或「肥羊和瘦羊」所指分別是只求自己利益和享樂並欺壓百姓的「以色列的牧人」及被他們所欺壓的羊,即以色列的百姓.

「牧羊人」是古代近東 (包括以色列所在地巴勒斯坦在內) 一個經常用來形容神明或君王的比喻.在此,先知以西結擴充了比喻,以「牧羊人」(眾數) 形容「以色列的君王和領袖」,以比喻形式來指斥他們的惡行.先知譴責他們自私自利,自顧自己的溫飽和財富,沒有履行責任保護弱者,反倒剝削和壓迫弱者,不能雪中送炭,反而濫用權力,恃強凌弱,橫行霸道. 他們的表現,從現代人的角度看,可以形容為反社會行為,道德意識低落,只求自己益處,不惜危害他們,甚至做一些損人不利己的事情.有別於以西結書,馬太福音用了「綿羊」和「山羊」比喻行公義的人與行不義的人.

(3)第三個對耶穌基督稱號是「君王」,一個坐在榮耀寶座上,審判萬民的君王.馬太以「萬民」一字意味著這個審判是普世的,涉及所有人.相反地,以西結書34章所提及的審判只是針對以色列民族,其領袖和平民百姓.

III.第一個謎團:誰是終末審判所暗示的壓迫者?

意識形態批判 (ideological criticism)有一個概念為「可測量的空缺」(measurable absence),[1] 所指的是作者應該說的,但又沒有說或不能說的內容. 例如:我們會期望一篇訃聞簡單提及死者的死因和去世的地點,如果訃聞沒有提及的話,這就是一個「可測量的空缺」,令人懷疑作者為何不提及.「可測量的空缺」隱藏了作者應表達但沒有表達的內容,作者不表達可能是要避免一些可能引起讀者不安或對作者不利的情況. 「可測量的空缺」是一種帶有意圖的寫作手法,它的出現令人懷疑作者不表述的動機.

馬太福音25:31-46的其中一個「可測量的空缺」,就是經文沒有提及誰是終未審判所暗示的壓迫者.我們知道耶穌認同受迫害者的遭遇,並且感受到他們在肉身上所承受的痛苦.他對義人說:「因為我餓了,你們給我吃；渴了,你們給我喝；我流浪在外,你們留我住；我赤身露體,你們給我穿；我病了,你們看顧我；我在監獄裡,你們來看我.」(vv.35-36).相反地,譴責行不義的人,餓了,沒有給他吃；渴了,沒有給他喝；他流浪在外,沒有收留他；他赤身露體,沒有給他穿；他病了或坐監,沒有來看顧他 (vv.42-43).

是誰在迫害耶穌呢?或者我們應該問「是誰在迫害這些貧瘠及軟弱的人,以至耶穌感受到切膚之痛呢?」經文並沒有透露迫害者的身份.但是相信,作者無需直接說出,馬太群體已經知道作者所暗示的是誰.

以西結書34:1-10告訴我們,上帝的子民中有「瘦弱的,沒有[人]養壯；有病的,沒有[人]醫治；受傷的,沒有[人]纏裹；被逐的,沒有[人]領回；失喪的,沒有[人]尋找；反而用強暴嚴嚴地轄制」.經文所指壓迫上帝的子民的人,正是「以色列的牧人」,即「以色列的君王和領袖」.以西結書33章記載了猶大國被滅,耶路撒冷淪陷的訊息傳到在被擄於巴比倫迦巴魯河邊 (River Chebar of Babylon)的猶太流放者,先知以西結也在其中.34章隨即以牧人的比喻將猶大國被滅的責任歸咎於「以色列君王和政治領袖」的管理不善,腐敗和暴政.

耶和華的話臨到以西結說:「人子啊,你要向以色列的牧人發預言,攻擊他們,說,主耶和華如此說:禍哉!以色列的牧人只知牧養自己.牧人豈不當牧養群羊嗎?你們吃脂油,穿羊毛,宰肥壯的,卻不牧養群羊.」(vv.1-3) 耶和華之後繼續說:「你們這些肥壯的羊,在美好的草場吃草還以為小事嗎?剩下的草,你們竟用蹄踐踏了；你們喝清水,剩下的水,你們竟用蹄攪渾了.至於我的羊,只得吃你們所踐踏的,喝你們所攪渾的.」(vv.18-19) 以西結書比喻中,虐待他們,迫害及剝削的自己的子民.

在剛才我們所讀的以西結書34:11-17中,我們看見耶和華不甘看見百姓沒有牧者的情況,四散,從各國領他們歸回故土,親自作他們的好牧人.「失喪的,我必尋找；被逐的,我必領回；受傷的,我必纏裹；有病的,我必醫治；只是肥的壯的,我必除滅,也要秉公牧養牠們.」(v.16)

為何馬太福音的作者不清楚指出?相信大家都知道在帝國的霸權下生活的民族,往往為了保護自己的群體安全,避免遭逼害,招來殺身之禍,需要逃避帝國耳目,審查,而不會直接說出帝國不喜歡聽的話,例如指出帝國就是其迫害者.對帝國的批判往往只能含沙射影地,以間接形式,統治者聽不懂,但受壓迫者聽得懂的形式來表達,這也是學者James C. Scott所指的「隱藏文本」 (hidden transcripts).[2] 以西結書34章的情況不同,因為以西結書所指責的迫害者「以色列的君王和領袖」已經被巴比倫帝國所滅或取替,作者也無須避諱.

IV.第二個謎團:誰是「一個最小的弟兄」?

從福音書我們知道:耶穌與門徒周圍傳道,四處流浪,居無定所,飽受飢餓,甚至要承受赤身露體,疾病或坐監之危機.耶穌說:「狐狸有洞,天空的飛鳥有窩,人子卻沒有枕頭的地方.」(太8:20；路9:58) 而他說這句,正正是要讓門徒清楚跟隨耶穌自己的代價.

當耶穌呼召十二門徒時,他曾經說:「人接待你們就是接待我；接待我就是接待那差我來的.人因為先知的名接待先知,必得先知所得的賞賜；人因為義人的名接待義人,必得義人所得的賞賜.無論何人,因為門徒的名,只把一杯涼水給這小子裡的一個喝,我實在告訴你們,這人不能不得賞賜.」(太10:40-42) 耶穌與門徒四處傳道,沿路上都是依靠的別人的收容和接濟.在古代(希伯來人或希臘人)的社會很看重接待旅客.當時期社會並不像現在那麼發達,可隨處找到酒店住宿.古時代旅店不多,租金亦很昂貴,設施簡陋,而且聲譽不好(如經營妓院).所以,人們出門都十分依賴被接待.保羅,提摩太,耶穌都曾經被接待過.接待是一個很普遍的現象.

早期基督徒更要面對一些猶太組織的逼迫,連使徒保羅也曾經逼迫過基督徒.及後,他成為基督徒,也要面對種種的迫害,甚至因福音的原故坐牢.羅8:35-37充分表達了保羅為主受苦的決心:「誰能使我們與基督的愛隔絕呢?難道是患難嗎?是困苦嗎?是逼迫嗎?是飢餓嗎?是赤身露體嗎?是危險嗎?是刀劍嗎?如經上所記:『我們為你的緣故終日被殺；人看我們如將宰的羊.』然而,靠著愛我們的主,在這一切的事上已經得勝有餘了.」(羅8:35-37) 可想而知,在第一世紀當基督徒所要面對的困難有多大,所需要的韌力和毅力是何等大.居無定所,依靠別人接濟,更要面對迫害的可能性.

馬太福音的寫作時間是在公元後70年至100年之間,即第一世紀的初期教會,當時期的猶太基督徒,的確要面對監禁,被虐待的遭遇 (參希10:32-34; 13:1-3).主後49年,該撒革老丟 (Claudius) 下令驅逐猶太人離開羅馬,其後羅馬王尼祿 (Nero) 更逼害早期教會.第一次的猶太羅馬戰爭(The First Jewish-Roman War) 發生於公元66年至73年,由一班猶太平民發動反羅馬起義,尼祿派兵鎮壓,事件引至耶路撒冷被攻陷,第二聖殿於公元70年被毀.在這段騷亂的日子內,不少猶太人及基督徒遭受苦害,被捆綁,受監禁,甚至成為無家可歸的人.

因此,雖然馬太沒有說明的受逼迫的處境,但是從聖經的記載,當時的歷史的處境,以及當時馬太群體對以西結書34章所指的迫害的了解,相信馬太群體很容易會聯想耶穌所指的受壓迫者就是當是受羅馬政權及一些親羅馬政權的猶太組織逼迫猶太人和基督徒!接待這一班受苦的陌生人,就是等同於接待耶穌.

耶穌設身處地地明白他們的遭遇和感受,因為他自己也是在羅馬政權下,被親羅馬政權的猶太組織逼迫而死的.因此,耶穌不願意看見有人好像他一樣受到困苦壓迫,成為政權底下的犧牲品.誰是「一個最小的弟兄」?相信在初期教會時代,這個「最小的弟兄」就是那些在羅馬政權,被帝國主義奴化,被邊緣化,受盡壓迫和剝削的人,以及堅決跟隨基督的人.

V.總結:學習主以同理心與受壓迫者同行

從以西結書34章,我們看見上帝指責「以色列的君王與領袖」,比喻他們為「壞牧人」.上帝取替了這些「壞牧人」後,自己當上了以色列的「牧人」,定意修補百姓的創傷,領以色列百姓歸回美地,安居樂業,以身作則當一個「好牧人」,並施行判斷,懲治「壞牧人」.但是,無論耶和華上帝有多不滿壓迫,剝削的行為,在以西結書34章中,受壓迫者仍然是被描述為一個被上帝拯救的對象.

到了馬太福音25:31-46,情況就不同了,作者將以西結書的「好牧人」和「分羊」的比喻加以發揮,他並不停留於形容上主對受壓迫者的拯救,而是將上主與受壓迫者劃上一個等號,並且以天啟文學的「終極審判」為論述的主軸,邀請他的讀者馬太群體(以及你與我)去反思﹐今天我們是否有參與「榮耀君王」和「好牧人」耶穌對受壓迫者的拯救行動,並且認清所服侍的對像就是耶穌所設身處地所認同的受壓迫者,甚至進深一步與受壓迫者認同,與受壓迫者為伍.他們所受的壓迫,就是主自己所受的壓迫,他們的無助,就是他的無助.弟兄姊妹,你們能否感受到耶穌所感受到的切膚之痛.

相信在座為人父母的或為人子女,應該理解「傷在你身,痛在我心」道理.父母愛子女的深,是足以令父母感受到子女所承受的痛苦.同樣地,耶穌因為深愛我們,他能夠感受我們所承受的痛苦,甚至壓迫.親愛的弟兄姊妹們,你知道耶穌有多愛你嗎?你又能否像耶穌一樣,能夠切膚之痛地認同受苦者,受壓迫者所承受的苦難呢?當我們去服侍他們時,你能否看見受苦的基督?明白到你所服侍的就是受苦的基督嗎?

仁愛會 (Missionaries of Charity) 的德蘭修女 (Mother Teresa) .在印度加爾各答 (Calcutta)最窮的地區莫地及爾(Moti Jihl)服侍窮人.德蘭修女為失學窮苦兒童辦了一所露天學校,由此開展了仁愛傳教女修會的團體.德蘭修女一生服務貧窮弱勢社群,對象包括愛滋,痳瘋患者以及街頭垂死的人群,她自己也在貧寒簡陋的環境中生活.在她的自傳中,她提及她服侍的力量完全來自耶穌基督的愛,當她看著所服侍對象時,她眼中只看到受苦的基督.

今主日的福音經課其實在說什麼呢?它所強調的就是對軟弱,無助,困苦者要有同理心 (empathy).不錯,如果我們好像耶穌一樣,曾經經歷過飢餓 (太4:1),口渴 (約19:28),流浪(太8:20),赤身露體 (約19:23-24),被羅馬兵丁羞辱,並含冤被釘十架,我想我們會更能體會被壓迫者的經歷和認同他們感受.經歷過類似的事件,的確會提升人的同理心.

讓我與大家分享一個經驗:當我初中時候我的全家移民到加拿大.自從16歲,我已搬離家獨立,自此之後也再沒有問父母拿取生活費.大學時代,我的生活開支及學費是依賴學生貸款,少量的兼職及暑期工的收入,少量的助學金及獎學金度過的.每年的學期終都是我經濟最緊絕的時候,因為大部分的收入都是在學年初取得.記得大學第三年接近學期終的時候,我的經濟又再一次出現困難.最差時期,我的銀行戶口只餘下二十多元的加幣.除了上帝外,我再沒有告訴任何人我的經濟出現了困難.我心想,只要我節衣縮食多一段時候,等兼職出糧,到了暑期找一分工作,困境就很快會過去.記得一個周五晚上,團契完畢後我返回大學宿舍.一開大門,就看見證前方的餐桌上放了一包sung 菜,內有一個這般大的肇菜和一件這般大的牛肉(以手勢表示).若我一個人吃,我想足夠我吃兩個星期.我的室友告訴我,剛才教會的一個弟兄Peter到訪,託她告訴我,他將這袋sung 菜放在這裡.我好奇地問:「為何他把sung 菜放在這裡?」可是,我的室友也百思不得其解.我以為他只是把食物暫存在我的宿舍,待他在學校溫習完,就會取回.我立即致電問他將會何時來取回這袋sung 菜.原來那袋sung 菜是他買來給我的.

我問:「為何無緣無故買袋sung 菜給我?」

他回答:「知你近來唔掂,所以買袋sung 比你囉!俾錢,我知你一定唔要ge 啦!」

我問:「你點知我唔掂ge ?」

他回答:「一見你個樣就知你唔掂啦!我以前都試過經濟出現問題,可能係咁,好容易就感覺到你唔掂!而且,見你咁耐唔出來食飯,就知你無錢啦!」

聽了他的說話,我就靜下來,連說聲「多謝!」也不能,因為我的雙眼已充滿了淚水.我好感動,為著Peter的愛心和洞察力感恩.對於我來說,Peter是我的「天使」,是上帝在我困境中差來供給我需要的使者,將祝福帶給我.也是一位以行動將基督徒愛心的精神傳遞給我看的使者,讓我更明白神的心意.

但是,我們並不一定要有過類似的經驗,才能明白,理解,同情受苦者或被壓迫者的經歷,我們只需將心比己,嘗試將自己的腳放在他的鞋內去理解他的困境,嘗試感同身受地去明白他的遭遇,關注他的需要,我們也可以明白他的困境.必須一提的是:這種態度假設了我們要先暫時放下自己的鞋,放下我們的成見和我們既有的看事物的框架,嘗試以一個站於困境內的角度看事物.

今天,我們無法時光倒流,回到第一世紀的巴勒斯坦,給飢餓或渴了的耶穌,一個麵包,一杯涼水；給流浪的耶穌,一片瓦遮頭；給赤身露體的耶穌,一件衣物遮身；或者是看顧在病榻上的耶穌,又或者試探望被監禁的耶穌.但是,我們可以做到的是明白耶穌體恤軟弱者,認同受壓迫者的心,服侍身邊有需要的人,如同服侍耶穌.我門的君王基督不是一個高高在上的帝王,而是一個與受苦者和受壓迫者身同感受的君主.

[1]  See Terry Eagleton, Criticism and Ideology (London: Verso, 1976), 72.

[2] James C. Scott, Domination and the Arts of Resistance: Hidden Transcripts. New Haven, CT: Yale University Press, 2008.



\newpage

\section{萬物朝終}
\label{sec:528}
\textbf{伍渭文牧師 - 萬物朝終}
\newline
\newline
\hyperref[sec:527]{< < < PREV SERMON < < <}
~
\hyperref[sec:toc]{[返目錄]}
~
\hyperref[sec:529]{> > > NEXT SERMON > > >}


以賽亞書 64:1-9
\newline
\begin{longtable}{cl}
\hline
\hline
章節 & 經文 (和合本修訂版)\\
\hline
64:1 & \begin{tabularx}{0.7\textwidth}{X} 願你破天而降，願山在你面前震動， \end{tabularx} \\ \\ \relax
64:2 & \begin{tabularx}{0.7\textwidth}{X} 好像火燒乾柴，又如火將水燒開，使你敵人知道你的名，列國必在你面前發顫！ \end{tabularx} \\ \\ \relax
64:3 & \begin{tabularx}{0.7\textwidth}{X} 你曾做我們不能逆料可畏的事；那時你降臨，山嶺在你面前震動。 \end{tabularx} \\ \\ \relax
64:4 & \begin{tabularx}{0.7\textwidth}{X} 自古以來，人未曾聽見，未曾耳聞，未曾眼見，除你以外，還有神能為等候他的人行事。 \end{tabularx} \\ \\ \relax
64:5 & \begin{tabularx}{0.7\textwidth}{X} 你迎見那歡喜行義、記念你道的人；看哪，你曾發怒，因我們犯了罪；這景況已久，我們還能得救嗎？ \end{tabularx} \\ \\ \relax
64:6 & \begin{tabularx}{0.7\textwidth}{X} 我們都如不潔淨的人，所行的義都像污穢的衣服。我們如葉子漸漸枯乾，罪孽像風把我們吹走。 \end{tabularx} \\ \\ \relax
64:7 & \begin{tabularx}{0.7\textwidth}{X} 無人求告你的名，無人奮力抓住你。你轉臉不顧我們，你使我們因罪孽而融化。 \end{tabularx} \\ \\ \relax
64:8 & \begin{tabularx}{0.7\textwidth}{X} 但耶和華啊，現在你仍是我們的父！我們是泥，你是陶匠；我們都是你親手所造的。 \end{tabularx} \\ \\ \relax
64:9 & \begin{tabularx}{0.7\textwidth}{X} 耶和華啊，求你不要大發震怒，也不要永遠記得罪孽；看哪，求你垂顧我們，因我們都是你的百姓。 \end{tabularx} \\ \\
[1ex]
\hline
\hline
\end{longtable}
詩篇 80:1-7,17-19
\newline
\begin{longtable}{cl}
\hline
\hline
章節 & 經文 (和合本修訂版)\\
\hline
80:1 & \begin{tabularx}{0.7\textwidth}{X} 領約瑟如領羊群的以色列牧者啊，求你側耳而聽！在基路伯之上坐寶座的啊，求你發出光來！ \end{tabularx} \\ \\ \relax
80:2 & \begin{tabularx}{0.7\textwidth}{X} 在以法蓮、便雅憫、瑪拿西面前求你施展你的大能，拯救我們。 \end{tabularx} \\ \\ \relax
80:3 & \begin{tabularx}{0.7\textwidth}{X} 神啊，求你使我們回轉，使你的臉發光，我們就會得救！ \end{tabularx} \\ \\ \relax
80:4 & \begin{tabularx}{0.7\textwidth}{X} 耶和華---萬軍之神啊，你因你百姓的禱告發怒，要到幾時呢？ \end{tabularx} \\ \\ \relax
80:5 & \begin{tabularx}{0.7\textwidth}{X} 你以眼淚當食物給他們吃，量出滿碗的眼淚給他們喝。 \end{tabularx} \\ \\ \relax
80:6 & \begin{tabularx}{0.7\textwidth}{X} 你使鄰邦因我們紛爭，我們的仇敵彼此戲笑。 \end{tabularx} \\ \\ \relax
80:7 & \begin{tabularx}{0.7\textwidth}{X} 萬軍之神啊，求你使我們回轉，使你的臉發光，我們就會得救！ \end{tabularx} \\ \\ \relax
80:8 & \begin{tabularx}{0.7\textwidth}{X} 你從埃及拔出一棵葡萄樹，趕出外邦人，把這樹栽上。 \end{tabularx} \\ \\ \relax
80:9 & \begin{tabularx}{0.7\textwidth}{X} 你在它面前清除雜物，它就深深扎根，蔓延滿地。 \end{tabularx} \\ \\ \relax
80:10 & \begin{tabularx}{0.7\textwidth}{X} 它的影子遮蔽群山，枝子好像高大的香柏樹。 \end{tabularx} \\ \\ \relax
80:11 & \begin{tabularx}{0.7\textwidth}{X} 它長出枝子，直到大海，伸展嫩枝，延到大河。 \end{tabularx} \\ \\ \relax
80:12 & \begin{tabularx}{0.7\textwidth}{X} 你為何拆毀這樹的籬笆，任憑路人摘取？ \end{tabularx} \\ \\ \relax
80:13 & \begin{tabularx}{0.7\textwidth}{X} 林中的野豬踐踏它，田裡的走獸吞吃它。 \end{tabularx} \\ \\ \relax
80:14 & \begin{tabularx}{0.7\textwidth}{X} 萬軍之神啊，求你轉回，從天上垂看觀察，眷顧這葡萄樹； \end{tabularx} \\ \\ \relax
80:15 & \begin{tabularx}{0.7\textwidth}{X} 保護你右手所栽的根，你為自己所堅固的幼苗。 \end{tabularx} \\ \\ \relax
80:16 & \begin{tabularx}{0.7\textwidth}{X} 這樹已經被火焚燒，被刀砍伐，因你臉上的怒容就滅亡了。 \end{tabularx} \\ \\ \relax
80:17 & \begin{tabularx}{0.7\textwidth}{X} 願你的手扶持你右邊的人，你為自己所堅固的人子。 \end{tabularx} \\ \\ \relax
80:18 & \begin{tabularx}{0.7\textwidth}{X} 這樣，我們就不背離你；求你救活我們，讓我們得以求告你的名。 \end{tabularx} \\ \\ \relax
80:19 & \begin{tabularx}{0.7\textwidth}{X} 耶和華---萬軍之神啊，求你使我們回轉，使你的臉發光，我們就會得救！ \end{tabularx} \\ \\
[1ex]
\hline
\hline
\end{longtable}
哥林多前書 1:3-9
\newline
\begin{longtable}{cl}
\hline
\hline
章節 & 經文 (和合本修訂版)\\
\hline
1:3 & \begin{tabularx}{0.7\textwidth}{X} 願恩惠、平安從我們的父神並主耶穌基督歸給你們！ \end{tabularx} \\ \\ \relax
1:4 & \begin{tabularx}{0.7\textwidth}{X} 我常為你們感謝我的神，因神在基督耶穌裡所賜給你們的恩惠。 \end{tabularx} \\ \\ \relax
1:5 & \begin{tabularx}{0.7\textwidth}{X} 因為你們在他裡面凡事富足，具有各種口才、各樣知識， \end{tabularx} \\ \\ \relax
1:6 & \begin{tabularx}{0.7\textwidth}{X} 正如我為基督作的見證在你們心裡得以堅固， \end{tabularx} \\ \\ \relax
1:7 & \begin{tabularx}{0.7\textwidth}{X} 以致你們在恩賜上一無欠缺，切切等候我們主耶穌基督的顯現。 \end{tabularx} \\ \\ \relax
1:8 & \begin{tabularx}{0.7\textwidth}{X} 他也必堅固你們到底，使你們在我們主耶穌基督的日子無可指責。 \end{tabularx} \\ \\ \relax
1:9 & \begin{tabularx}{0.7\textwidth}{X} 神是信實的，他呼召你們好與他兒子—我們的主耶穌基督—共享團契。 \end{tabularx} \\ \\
[1ex]
\hline
\hline
\end{longtable}
馬可福音 13:24-37
\newline
\begin{longtable}{cl}
\hline
\hline
章節 & 經文 (和合本修訂版)\\
\hline
13:24 & \begin{tabularx}{0.7\textwidth}{X} 「在那些日子、那災難以後，太陽要變黑，月亮也不放光， \end{tabularx} \\ \\ \relax
13:25 & \begin{tabularx}{0.7\textwidth}{X} 眾星要從天上墜落，天上的萬象都要震動。 \end{tabularx} \\ \\ \relax
13:26 & \begin{tabularx}{0.7\textwidth}{X} 那時，他們要看見人子帶著大能力和榮耀駕雲來臨。 \end{tabularx} \\ \\ \relax
13:27 & \begin{tabularx}{0.7\textwidth}{X} 他要差遣天使，從四方，從地極直到天邊，召集他的選民。」 \end{tabularx} \\ \\ \relax
13:28 & \begin{tabularx}{0.7\textwidth}{X} 「你們要從無花果樹學習功課：當樹枝發芽長葉的時候，你們就知道夏天近了。 \end{tabularx} \\ \\ \relax
13:29 & \begin{tabularx}{0.7\textwidth}{X} 同樣，當你們看見這些事發生，就知道那時候近了，就在門口了。 \end{tabularx} \\ \\ \relax
13:30 & \begin{tabularx}{0.7\textwidth}{X} 我實在告訴你們，這世代還沒有過去，這一切都要發生。 \end{tabularx} \\ \\ \relax
13:31 & \begin{tabularx}{0.7\textwidth}{X} 天地要廢去，我的話卻絕不廢去。」 \end{tabularx} \\ \\ \relax
13:32 & \begin{tabularx}{0.7\textwidth}{X} 「但那日子，那時辰，沒有人知道，連天上的天使也不知道，子也不知道，惟有父知道。 \end{tabularx} \\ \\ \relax
13:33 & \begin{tabularx}{0.7\textwidth}{X} 你們要謹慎，要警醒，因為你們不知道那時刻幾時來到。 \end{tabularx} \\ \\ \relax
13:34 & \begin{tabularx}{0.7\textwidth}{X} 這事正如一個人離家遠行，授權給僕人們，分派各人的工作，又吩咐看門的警醒。 \end{tabularx} \\ \\ \relax
13:35 & \begin{tabularx}{0.7\textwidth}{X} 所以，你們要警醒，因為你們不知道這家的主人甚麼時候來，是晚上，或半夜，或雞叫時，或早晨， \end{tabularx} \\ \\ \relax
13:36 & \begin{tabularx}{0.7\textwidth}{X} 免得他忽然來到，看見你們睡著了。 \end{tabularx} \\ \\ \relax
13:37 & \begin{tabularx}{0.7\textwidth}{X} 我對你們所說的話，也是對眾人說的：要警醒！」 \end{tabularx} \\ \\
[1ex]
\hline
\hline
\end{longtable}
講題:萬物朝終

經文:以賽亞書64:1-9;詩篇80:1-7,17-19;哥林多前書1:3-9;馬可福音13:24-37

講員:伍渭文牧師

場合:香港中文大學崇基學院禮拜堂主日崇拜

日期:2020年11月29日

【萬物朝終】賽六十四1-9; 詩八十1-7,17-19; 林前一3-9; 可十三24-37

1942二戰時,英軍在埃及北非亞拉莫(El Alamein)首次大敗隆美爾(Rommel)率領之德軍,之前英軍從未戰勝德軍,之後從未敗給德軍.英軍興奮極了,歡呼大戰將快結束.英首相邱吉爾(Winston Churchchill)演說提及這事時說:「這不是結束,也不是結束的開始,也許,剛剛開了個頭.」(Now this is not the end. It is not even the beginning of the end. But it is, perhaps, the end of the beginning.) 亞拉莫是End of the beginning-剛剛開了個頭, 將臨期是Beginning of the End-終末的開始.

聖經對「終」的概念有兩個論述:創造論的終結,和救贖論的終末.

創造論論述:「凡事都有定期,天下萬務都有定時,生有時,死有時,栽種有時,拔出所栽種的也有時.」(傳三1, 2)聖經智慧文學像傳道書,箴言和很多詩篇,常常提醒我們人生是有終結的.智慧文學是從創造的世界,以冷峻的眼光看紛陳世事,是現實人生的清醒劑.自然生物界,有開始也有結束,人有一生,必有一死.

杜甫《登高》:「無邊落木蕭蕭下,不盡長江滾滾來」也是終結.樹葉發芽,茁壯成長,由翠綠轉淡黃,由淡黃轉淺褐；枯葉經風一吹,蕭蕭而下,化作黃土.人年輕時,銳氣凌厲,像上游河出高山,破峽而出.壯年意氣風發,像河之中游,百川匯集,水力充沛,河道寬闊.趨近晚年,體力衰減,如臨海堆積平原的曲流,逶迤緩慢,漸漸沒入大海,消失了,這也是終結.

救贖論論述:基督道成肉身已開始了終末,因為永恆已進入時間到了.他擔負世人的罪被釘死在十架,埋葬了,第三天從死裡復活,指出死不是終點,還有身體復活,最後審判,新天新地.基督第一次來時已捅開了封閉的世界,讓我們窺探到人類最後的命運,上主永恆旨意的成就,這一切在他第二次來臨時必要成就.

終末也是終於.終於,按部就班,計劃完成.創始那一位,也是成終那一位.在那一刻,我們看到人歸塵土,百川匯海,萬葉歸根,不是宿命,是上主隱密的計劃,有其獨特的意義.終局誰主浮沉?誰能笑到最後?是那位又是始,亦是終的上主.

今天讀的經課是新的教會年第一主日的經課,共通的概念,都提及那日子和悔改:

1)可十三24:「在那些日子,那災難以後」；32:「但那日子,那時辰,沒有人知道,連天上的使者也不知道,惟有父知道.」(日子)

2)林前一8:「他必堅固你們到底,叫你們在我們主耶穌基督的日子無可責備.」(日子)

3)賽六十四1,3:「願你裂天而降,願山在你面前震動.」「求你不要大發震怒,也不永遠記念罪孽.」(悔改)

4)詩八十3:「神啊,求你使我們回轉(復興),使你的臉發光,我們便要得救.」(悔改)

究竟那日子和悔改有何關係呢?那日子重要,因為迎見永生神.悔改是靈性不斷更新,心意不斷變化,像新娘裝飾自己,迎見新郎.將臨期是基督降生的預備期,迎接世界的光,首先要除去內心的幽暗.迎接基督再來,就要站立得穩,無可責備.將臨期不單是紀念基督第一次來臨,也提醒信徒,預備迎接基督第二次來臨.

1 日子臨近

1)安慰的日子

我們先看看「那日子」.想到終末那日子,就想到民攻打民,國要攻打國,地震,海嘯,漫天烽火,滿眼廢墟,世界末日.在美國有些有線電視台全年的信息,就是未世的信息；如何曉得耶穌再來的時間,如何知道再來的徵兆.從以色列復國,歐洲聯盟的興起,尋找基督再來的預言,對號入座.但聖經談論終末,說到基督再來時,重點是叫我們儆醒,不是猜測時間和徵兆.福音經課提醒我們:「但那日子,那時辰,沒有人知道,連天上的使者也不知道,子也不知道,惟有父知道,你們要謹慎祈禱,因為你們不曉得那日期幾時來到.」(可十三32,33)

馬可福音十三章被稱為小啟示錄(Small Apocalypse),一如啟示錄,寫作目的是安慰鼓舞正受迫逼的教會,不是叫我們懼怕災難和迫害,乃曉得上主的保護.「人把你們拉去交官的時候,不要預先思慮說甚麼；到那時候,賜給你們甚麼話,你們就說甚麼.」(可十三11)

這日子並不遙遠,三十年後主後六十六年,猶太人起義,主後七十年羅馬提多將軍以武力鎮壓猶太人起義,圍城近五個月,110萬猶太人死去,聖殿徹底被毀,「沒有一塊石頭留在地上」(可十三2).「我實在告訴你們,這世代還沒有過去,這些事都要成就.」(可十三30)

2)喜樂的日子

聖經不單強調安慰,且指出這日子是喜樂的.「民要攻打民,國要攻打國,多處必有地震,饑荒,這都是災難的起頭.(災難原文作生產的苦難).」(可十三8)沒有待產的婦人單看生產之痛,母親專注的是新生命的誕生帶來的欣悅.苦難是產痛,新天新地誕生前的產痛.基督再來的日子,沒有人知道,聖經說到那日子是安慰和盼望,我們不用懼怕.經課提到那日子時,伴隨的是要悔改,預備自己迎見永生上主.

3) 回轉的日子

詩八十3:「神啊,求你使我們回轉(復興),使你的臉發光,我們便要得救.」詩篇用復興(restoration)表述悔改,指出我們遠離了原有的興旺,優美的生命.本來茁壯堅拔的生命,因為疏於料理變得萎弱不振.

我宿舍後園的白蘭花,若多天不澆水,長長的葉垂頭喪氣,沒精打彩.但當我立刻澆水,兩三小時後葉子就挺地來.有一天大吃了一驚,白蘭-樹看來特別風姿綽約,顧盼生輝,葉子特別油綠.我很不明白,問問師母,你是否施了肥?她說:是的,我前幾天使了肥.白蘭樹本來就天生麗質,奈何疏於照顧令它蒙污納垢,失去原來的亮麗.

使你的臉發光,求上主親近,驅走內心的幽暗.我們便要得救,感到自己生命的危機重重,頻於死亡邊沿,呼求上主救拔,不讓生命沉淪下去.詩人渴慕上主臨近,復興生命,向高處行,活出豐盛優美的人生.終末是那日,我們迎接新天新地；終末也是此刻,我們的生命已是新天新地,在基督裡造成的新天新地.

讓我們回轉,使你的臉發光,其實就是將臨期的主題.燃點將臨期燭光,不單迎世界的光的降生,也讓燭光掃除心中的幽暗,因罪藏在黑暗裡.王爾德(Ocar Wilde)《格雷的畫像》(The Picture of Dorian Gray)深刻指出罪藏在不為人知的黑暗中.小說的主角格雷英朗帥氣,失意的畫家裴素(Basil)見到他時振作起來,要為他畫像.格雷看到自己的肖像真的驚為天人:「希望我的外形能保持青春常駐,若有衰老,邪惡,由自己的肖像承受.我仍可保持無邪,迷人的笑容.」

有一次格雷做了一件邪惡的事,特意回到上了鎖的閣樓,看看令他心醉的畫像,竟然發現畫像中的他,嘴角有陰陰的奸笑,他赫然驚怕,真的每逢幹了一宗壞事,不單良心不安,而且鎖在閣樓沒有人知道的畫像,都會顯露出來.不過,他有恃無恐,因為靈魂的傷害不會形於色.跟著十多年,格雷生活變得放蕩,荒誕不羈,惡事做盡,以求一己肉慾的愉悅,而容貌永遠秀美青春迷人.然而,他的靈魂,像關在閣樓的俊美少年,每天愈來愈衰老,邪惡和難看.有一天他對畫他肖像的裴素說:「為了這像-我漂亮的像,我可以重新來過嗎?我愛的的肖像,我要好好做人.」

不過,肖像還是愈來愈差.為甚麼?因為格雷並不是猛然醒悟悔改,乃是為了保持肖像的完美,他仍是追求虛榮和自戀,不是悔改.最後,他竟然用刀剌死為他繪像的畫家裴素,並用同一張刀,刺破他不能接受-又老,又醜,又邪惡的肖像.在這時候,可怕的結局終於來到了.格雷突然變得瘋狂,在閣樓傳來恐佈的叫聲,令人不寒而憟,僕人急急召來警察,門開了,發現地上躺了一位又老又醜的死屍,若不是手上的戒指,根本認不出是格雷,而畫架上,出現一位秀美迷人的少年肖像.

我們親近上主,上主便親近我們,以他的臉光照我們.祭司祝福:「願耶和華賜福給你,保護你,願耶和華使他的臉光照你,願耶和華向你仰臉,賜你平安.」(民六24-26)

4)漫長的日子

基督再來那日子,沒有人知道,但我們知道,這日子是一個漫長的朝聖路,充滿試探和信心的煎熬,像耶穌被釘埋在地下的三天,期間甚至有門徒因失望重操故業,回去打魚.書信經課林前一8:「他必堅固你們到底,叫你們在我們主耶穌基督的日子無可責備.」彼得勸勉信徒:「就要愛慕那純淨的靈奶,像纔生的嬰孩愛慕奶一樣,叫你們因此漸長,以至得救.」(彼前二2)

天路客在漫長朝聖路上,有時難免迷失方向,誤入岐途,不時要調較生命的焦點.參加教會聚會,幫助我們重新聚焦.舊約經課賽六十四1-9就是聚會時的禱告,在禱告中上主臨到他們,聽他們禱告.

2 耶穌在線

1) 河海交匯

耶穌復活,四十天跟門徒談論天國的道理,之後升到父上主面前,坐在他的右邊,直到終末,被天軍天使簇擁,復臨地上.耶穌升天到父那裡,求父差遣聖靈作我們的保惠師,藉聖靈與我同在.用現在的話說,基督祇是不在線下-實體,與我們面對面.若我們從基督第二次再來的那點,逆向往後劃一條線把第一次來那點連起來,沿著這線的時間就是終末.在這線上我們可以藉著聖靈,透過聖道,聖禮,和聖教會,經歷基督的臨在,耶穌在線.

終末這條線,像頻海的一段河道,這段河道海水和河水混雜交匯.若說浩瀚無際的海水是永世,流淌不住的河水是今世.頻海的一段河道,海水和河水交匯的一段就是末世.這段河道發生了生態的變化.一些河鮮因著海水的鹽分,不能生長；亦因水的鹹度變化,出現了嶄新的河鮮.終末之前,我們不能不犯罪,耶穌來了,我們可以犯罪,也可以不犯罪.教會群體的出現就是新生事物,藉著聖洗,無論是男的,女的,為奴的,自主的,在基督裡都成為一體了.接受基督就活在終末裡,永世的權能已進入我們的生命.

2)申訴哀歌.

舊約經課賽六十四1說到上主在禱告中臨到我們.「願你裂天而降,願山在你面前震動.」這是一篇哀歌的中間,哀歌開始是六十三15:「求你從天垂顧……你的熱心和你的大能的作為在哪裡呢?」,結束是六十四12:「耶和華啊,有這些事,你還忍得住嗎?你仍靜默使我們深受苦難麼?」.

在舊約,當人向上主禱告,他們相信上主會走近禱告者,答允他們.以賽亞先知在崇拜中自覺咀唇不潔,上主就藉天使撒拉弗飛到他跟前,拿著從祭壇上取下來的紅炭,沾他的口潔淨他.

以賽亞書分三部分,三個不同歷史場景.一至三十九章處於主前八世紀,猶太國處於列強狹縫,地緣政治複雜.北有以法蓮和亞蘭(今敘利亞),亞述(今伊拉克),東有大國崛起中的巴比倫(今伊朗),南有埃及.公元前722年亞述攻陷撒馬利亞,北國以色列滅亡.

第二部分四十章至五十五章,巴比倫公元前587年攻陷耶路撒冷,聖殿被毀,猶太人被擄巴比倫.

第三部分,五十六至六十六章,記敘回歸故土在耶路撒冷的境況.回歸初期,境況堪虞.首批回歸故土的人,用心建造自己的房屋,卻對重建聖殿不起勁.「這殿仍然荒涼,你們自己還住天花板的房屋.」(該一4)舊約經課賽六十四1-9就是回歸初期的一首哀歌.

詩人看見回歸初期,聖殿仍然荒廢,百姓靈性冷淡,就對對上主發勞叨,表示不滿,這是哀歌的特色:渲洩痛苦,哀傷,不滿.詩人代表會眾詰問:「為何使我們走差離開你的道,使我們心裡剛硬,不敬畏你呢?」(賽六十三17)；也不滿「你的聖民不過暫時得這產業；我們的敵人已經踐踏你的聖所.」(賽六十三18)為何上主容許子民被擄巴彼倫,聖殿被毀壞,聖物被搶奪一空.

沒有信心的人,看見上主任由子民犯罪行惡,以致聖殿被毀,被擄異域,可能推論這位上主是惡的,因他容隱惡事而不作干預.見到聖殿被毀而沒有作為,也可推論這上帝是軟弱無力的.這正是英國懷疑主義之父休謨(David Hume)所說的,因他論說上主(talk about him).但作為禱告的哀歌,盡管表示不滿,是向上主申訴(talk to him).

魯益師提到一個黑暗殼倉門縫一線開的比喻.一線陽光透過門縫射進黑暗的殼倉,在裡面有一個人凝視光線(look at),祇見到不少微塵在光線中翻躍飛舞,不見其它.但當他踏進從門縫射進來的光線中望出去(Look along),見到另一個世界:綠樹交蔭,藍天白雲,山巒起伏；這是另一個真實的世界.崇拜中的哀歌,容許我們對上主舒洩痛苦,嗟嘆哀傷,甚至質問上主,因為我們直面上主,不是遙遙遠眺.在上主面光之中,我們才能得到光照.

3)抓住上主.

因此,在第二次求上主臨在的禱告(賽六十四1),跟第一次(賽六十三15)已經不同,他求上主裂天而降聽他禱告,想起昔日在西乃山對子民的顯現:「那時你降臨,山嶺在你面前震動.」(賽六十四3)「西乃全山冒煙,因為耶和華在火中降於山上,山上的煙氣上騰,如燒窯一般,遍山大大的震動.」(出二十18)在禱告中,他開始明白:「我們仍然犯罪,這景況已久,我們還能得救麼.我們都像不潔淨的人,所有的義都像污穢的衣服,我們都像葉子漸漸枯乾,我們的罪孽好像風把我們吹去.並且無人求告你的名,無人奮力抓住你.」(賽六十四5-7)

經文意像鮮活.罪異化我們的生命,不動聲色,生命像葉子慢慢枯乾,罪像風把我們吹散,情況非常危險,我們要奮力抓住上主,才能不被罪孽吹散.「無人奮力抓住你」也可應用在靈命的追求上.上主的感動像火炭含輝發亮(glow),使我們感到到溫暖火熱,我們要奮力抓住難得的感動,把發亮的炭扇出火焰(flame),才使感動不熄滅.

很少信徒會否認基督再來的教義,但不多愛慕基督的再臨,因我們很少經歷上主對我們顯現,或說造訪(visit)我們.我們疏於禱告,甚至停止聚會.不要停止聚會,給上主有機會造訪我們.

愛默生(Emerson)甚至說,聖禮也可以在路旁發現的:「美麗是上主的的手筆-路旁的聖禮.掛臉龐,橫天際,垂花蕊,感謝上主而來的福杯.」上主藉聖道,聖教會,甚至路旁聖禮,造訪我們,對我們顯現,堅立我們,直到見他的日子.

3誰主終末?

林前一3-9:「我常為你們感謝上主,因上主在基督裡所賜給你們的恩惠……等候我們的主耶穌基督顯現.他也必堅固你們到底,叫你們在我們主耶穌基督的日子無可指責.」

起初上帝

十月三十日校長段崇智教授,在崇基學院六十九周年校慶感恩崇拜應邀作嘉賓演講.他問大家:聖經最重要的是那一章節?是約翰福音三16「神愛世人」,是詩篇二十三篇「耶和華是我的牧者」?他說應該是「起初神」(In the beginning God).這話無似簡單,其實深邃.終末其實由起初決定.終末英文是end,有計劃和目的的意涵.

約但河的洗禮是悔改迎接神國的洗禮,耶穌沒有罪不需接受悔改的洗禮,他到約但河代替我們的罪而領洗.「他從水裡上來,就看見天裂開了,聖靈彷彿像鴿子,降在他身上.又有聲音從天上來,說,你是我的愛子,我喜悅你.」這是父上主發聲.「我喜悅你」是舊約彌賽亞的名號:「看哪,我的僕人-我所扶持所揀選,心裡所喜悅的!」(賽四十二1)當約但河水漫過其身,就像洪水漫過大地,洪水擊打方舟,在方舟裡面的挪亞一家八口,得免於難.這是父的原初救贖計劃,由基督執行完成.

天劍已屠龍.

正教會聖洗禱文提及耶穌約但河的洗禮:「你聖化了約但河的水,從天降下聖靈,把潛藏在水中龍的頭擊碎了.」這圖像活像天劍屠龍.在第二次基督再來時,新天新地出現:「海要交出其中的死人」(二十13)；「海也不再有了」(二十一1).耶穌踏進約但河是舊的世界,踏出約但河是另一個新的世界-終末的世界.終末第一點,決定了終末的第二點.已濟決定了未濟 (Already decides Not Yet).

基督已復活.

保羅跟哥林多教會說:「我當日所領受又傳給你們的:第一,就是基督照聖經所說,為我們的罪死了,而且埋葬了；又照聖經所說,第三天復活了,並且顯給磯法看,然後顯給十二使徒看；後來一時顯給五百多弟兄看,其中一大半到如今還在,卻也有已經睡了的.」(林前十五3-6)我們將要復活,因為基督已經復活.

死後將如何?我們不知道,惟有他知道,因他為我們嘗過死味.「不要懼怕,我是首先的,我是未後的……我曾死過,現在又活了,直活到永永遠遠.並且拿著死亡和陰間的鑰匙.」(啟一17, 18)基督死而復活,開啟了終末,也決定了終局.

終結與終末:喬布斯與奧古斯丁.

喬布斯傳記作者華萊士.以撒遜(Walter Isaacson)與喬布斯作了超過四十次訪談,最後一次在他家中的後花園進行,是陽光普照的下午,當時喬布斯身體很不舒服,喬布斯把話題轉到宗教和死亡.當問及他是否相信來生(afterlife),是否相信有上帝時,他說:「五十五十.這是大奧秘,人永遠不知道,但我希望死不是生命的終結,我但願累積的智慧不會隨死亡歸於無有,可以永續.但他跟著說:但可能它真的像開關制,click一下,你就離開了；停頓一秒後他笑說:所以我們蘋果產品沒有開關制的.」喬布斯對人生終點之後有沒有永生,祇能說「五十五十,但願可以永續,但可能像開關制click一下就沒有了.」

奧古斯丁經歷在基督裡的救恩,令他生命徹底改變,他把這終末的經歷寫成《懺悔錄》,並作了一個壓卷的禱告:「上帝啊,你既賜我們一切,就請你賜給我們平安,那歇息的平安,安息日的平安,沒有日落的平安.一切受造之物雖然美善,然而它們都要消逝.時限將到,美好的一切也會結束.它們有它們的清晨,各自的黃昏.

可是,第七日卻沒有日落,因為你祝聖這日,使它成為永恆的白晝.雖然創造的工作,並沒有中斷你永恆的安息,但你造物之後,看見所造皆為美好,在第七天就安息 了.在你聖言中,我們看到一個預示:我們完成今生的工作後,亦會在你懷中得著永恆生命的歇息.然而,衹有得著你的恩典,我們才能妥善完成工作.

在那永恆的安息日,你在我們裡面安息；一如現在,你在我們裡面工作.我們享受的安息本是你的,一如我們現在的工作,也是你在我們裡面完成的.可是,主阿,你永遠在工作,也永遠在安息.

你不在時間下看見,不在時間中活動,也不在時間內歇息.但我們在時間裡看見到的一切,皆你斧鑿而成,你創造時間,也在時間的盡 頭,設立安息.阿們.」 (伍渭文中譯)

結論:「萬物朝終」,但無知曉終在何日.「那日子,那時辰,沒有人知道.」迎見上主有兩個方式,一是他來,一是我去.我們不能選擇何時何處生在這世界,也不能掌控何時何處離開,這是上主的主權.但我們有自由,選擇如何活得精彩.

在結束時,我們作一個新年的禱告:「神啊,求你使我們回轉(復興),使你的臉發光,我們便要得救.」(詩八十3)

復興我們,活出優美的生命.使你的臉發光,驅走心中的幽暗.

我們要得救,經歷你的能力.啊們.



\newpage

\section{各就各位預備}
\label{sec:529}
\textbf{高國雄牧師 - 各就各位,預備……}
\newline
\newline
\hyperref[sec:528]{< < < PREV SERMON < < <}
~
\hyperref[sec:toc]{[返目錄]}
~
\hyperref[sec:530]{> > > NEXT SERMON > > >}


馬可福音 1:1-8
\newline
\begin{longtable}{cl}
\hline
\hline
章節 & 經文 (和合本修訂版)\\
\hline
1:1 & \begin{tabularx}{0.7\textwidth}{X} 神的兒子，耶穌基督福音的起頭。 \end{tabularx} \\ \\ \relax
1:2 & \begin{tabularx}{0.7\textwidth}{X} 正如以賽亞先知書上記著：「看哪，我要差遣我的使者在你面前，他要為你預備道路。 \end{tabularx} \\ \\ \relax
1:3 & \begin{tabularx}{0.7\textwidth}{X} 在曠野有聲音呼喊著：預備主的道，修直他的路。」 \end{tabularx} \\ \\ \relax
1:4 & \begin{tabularx}{0.7\textwidth}{X} 照這話，施洗約翰來到曠野，宣講悔改的洗禮，使罪得赦。 \end{tabularx} \\ \\ \relax
1:5 & \begin{tabularx}{0.7\textwidth}{X} 猶太全地和全耶路撒冷的人都出去，到約翰那裡，承認他們的罪，在約旦河裡受他的洗。 \end{tabularx} \\ \\ \relax
1:6 & \begin{tabularx}{0.7\textwidth}{X} 約翰穿駱駝毛的衣服，腰束皮帶，吃的是蝗蟲和野蜜。 \end{tabularx} \\ \\ \relax
1:7 & \begin{tabularx}{0.7\textwidth}{X} 他宣講，說：「有一位在我以後來的，能力比我更大，我就是彎腰給他解鞋帶也不配。 \end{tabularx} \\ \\ \relax
1:8 & \begin{tabularx}{0.7\textwidth}{X} 我用水給你們施洗，他卻要用聖靈給你們施洗。」 \end{tabularx} \\ \\
[1ex]
\hline
\hline
\end{longtable}
以賽亞書 40:1-11
\newline
\begin{longtable}{cl}
\hline
\hline
章節 & 經文 (和合本修訂版)\\
\hline
40:1 & \begin{tabularx}{0.7\textwidth}{X} 你們的神說：「要安慰，安慰我的百姓。 \end{tabularx} \\ \\ \relax
40:2 & \begin{tabularx}{0.7\textwidth}{X} 要對耶路撒冷說安慰的話，向它宣告，它的戰爭已結束，它的罪孽已赦免；它為自己一切的罪，已從耶和華手中加倍受罰。」 \end{tabularx} \\ \\ \relax
40:3 & \begin{tabularx}{0.7\textwidth}{X} 有聲音呼喊著：「要在曠野為耶和華預備道路，在沙漠為我們的神修直大道。 \end{tabularx} \\ \\ \relax
40:4 & \begin{tabularx}{0.7\textwidth}{X} 一切山窪都要填滿，大小山岡都要削平；陡峭的要變為平坦，崎嶇的必成為平原。 \end{tabularx} \\ \\ \relax
40:5 & \begin{tabularx}{0.7\textwidth}{X} 耶和華的榮耀必然顯現，凡有血肉之軀的都一同看見，因為這是耶和華親口說的。」 \end{tabularx} \\ \\ \relax
40:6 & \begin{tabularx}{0.7\textwidth}{X} 有聲音說：「你喊叫吧！」我說：「我喊叫甚麼呢？」凡有血肉之軀的盡都如草，他的一切榮美像野地的花。 \end{tabularx} \\ \\ \relax
40:7 & \begin{tabularx}{0.7\textwidth}{X} 耶和華吹一口氣，草就枯乾，花也凋謝。百姓誠然是草； \end{tabularx} \\ \\ \relax
40:8 & \begin{tabularx}{0.7\textwidth}{X} 草必枯乾，花必凋謝，惟有我們神的話永遠立定。 \end{tabularx} \\ \\ \relax
40:9 & \begin{tabularx}{0.7\textwidth}{X} 報好信息的錫安哪，要登高山；報好信息的耶路撒冷啊，要極力揚聲。揚聲不要懼怕，對猶大的城鎮說：「看哪，你們的神！」 \end{tabularx} \\ \\ \relax
40:10 & \begin{tabularx}{0.7\textwidth}{X} 看哪，主耶和華必以大能臨到，他的膀臂必為他掌權；看哪，他的賞賜在他那裡，他的報應在他面前。 \end{tabularx} \\ \\ \relax
40:11 & \begin{tabularx}{0.7\textwidth}{X} 他要像牧人牧養自己的羊群，用膀臂聚集羔羊，抱在胸懷，慢慢引導那乳養小羊的。 \end{tabularx} \\ \\
[1ex]
\hline
\hline
\end{longtable}
詩篇 85:1-2,8-13
\newline
\begin{longtable}{cl}
\hline
\hline
章節 & 經文 (和合本修訂版)\\
\hline
85:1 & \begin{tabularx}{0.7\textwidth}{X} 耶和華啊，你已經向你的地施恩，救回被擄的雅各。 \end{tabularx} \\ \\ \relax
85:2 & \begin{tabularx}{0.7\textwidth}{X} 你赦免了你百姓的罪孽，遮蓋了他們一切的過犯。（細拉） \end{tabularx} \\ \\ \relax
85:3 & \begin{tabularx}{0.7\textwidth}{X} 你收回所發的憤怒，撤銷你猛烈的怒氣。 \end{tabularx} \\ \\ \relax
85:4 & \begin{tabularx}{0.7\textwidth}{X} 拯救我們的神啊，求你使我們回轉，使你向我們所發的憤怒止息。 \end{tabularx} \\ \\ \relax
85:5 & \begin{tabularx}{0.7\textwidth}{X} 你要向我們發怒到永遠嗎？要將你的怒氣延留到萬代嗎？ \end{tabularx} \\ \\ \relax
85:6 & \begin{tabularx}{0.7\textwidth}{X} 你不再將我們救活，使你的百姓因你歡喜嗎？ \end{tabularx} \\ \\ \relax
85:7 & \begin{tabularx}{0.7\textwidth}{X} 耶和華啊，求你使我們得見你的慈愛，又將你的救恩賜給我們。 \end{tabularx} \\ \\ \relax
85:8 & \begin{tabularx}{0.7\textwidth}{X} 我要聽神－耶和華所說的話，因為他必應許賜平安給他的百姓，就是他的聖民；他們卻不可再轉向愚昧。 \end{tabularx} \\ \\ \relax
85:9 & \begin{tabularx}{0.7\textwidth}{X} 他的救恩誠然與敬畏他的人相近，使榮耀住在我們的地上。 \end{tabularx} \\ \\ \relax
85:10 & \begin{tabularx}{0.7\textwidth}{X} 慈愛和誠實彼此相遇，公義與和平彼此相親。 \end{tabularx} \\ \\ \relax
85:11 & \begin{tabularx}{0.7\textwidth}{X} 誠實從地而生，公義從天而現。 \end{tabularx} \\ \\ \relax
85:12 & \begin{tabularx}{0.7\textwidth}{X} 耶和華必賜福氣給我們；我們的地也要出土產。 \end{tabularx} \\ \\ \relax
85:13 & \begin{tabularx}{0.7\textwidth}{X} 公義要行在他面前，使他的腳蹤有可走之路。 \end{tabularx} \\ \\
[1ex]
\hline
\hline
\end{longtable}
彼得後書 3:8-15
\newline
\begin{longtable}{cl}
\hline
\hline
章節 & 經文 (和合本修訂版)\\
\hline
3:8 & \begin{tabularx}{0.7\textwidth}{X} 親愛的，有一件事你們不可忘記，就是：主看一日如千年，千年如一日。 \end{tabularx} \\ \\ \relax
3:9 & \begin{tabularx}{0.7\textwidth}{X} 主沒有遲延他的應許，就如有人以為他是遲延，其實他是寬容你們，不願一人沉淪，而是人人都來悔改。 \end{tabularx} \\ \\ \relax
3:10 & \begin{tabularx}{0.7\textwidth}{X} 但主的日子要像賊一樣來到；那日，天必在轟然一聲中消失，天體都要被烈火熔化，地和地上的萬物都要燒盡。 \end{tabularx} \\ \\ \relax
3:11 & \begin{tabularx}{0.7\textwidth}{X} 既然這一切都要如此消失，你們處世為人必須聖潔敬虔， \end{tabularx} \\ \\ \relax
3:12 & \begin{tabularx}{0.7\textwidth}{X} 等候並催促神的日子來到。因為在那日，天要被火燒而消滅，天體都要被烈火熔化。 \end{tabularx} \\ \\ \relax
3:13 & \begin{tabularx}{0.7\textwidth}{X} 但照他的應許，我們等候新天新地，其中有正義常住。 \end{tabularx} \\ \\ \relax
3:14 & \begin{tabularx}{0.7\textwidth}{X} 所以，親愛的，既然你們等候這些事，就要竭力使自己沒有玷污，無可指責，在主前和睦； \end{tabularx} \\ \\ \relax
3:15 & \begin{tabularx}{0.7\textwidth}{X} 並且要以我們主的容忍作為你們得救的機會，就如我們所親愛的弟兄保羅，照著所賜給他的智慧寫信給你們。 \end{tabularx} \\ \\
[1ex]
\hline
\hline
\end{longtable}
講題:各就各位,預備……

經文:馬可福音1:1-8;以賽亞書40:1-11;詩篇85:1-2,8-13;彼得後書3:8-15

講員:高國雄牧師

場合:香港中文大學崇基學院禮拜堂主日崇拜

日期:2020年12月6日

各位弟兄姊妹,今天是將臨期第二主日,你預備好未呢?又或更需要問的問題係:我們要預備甚麼呢?要做甚麼才算預備足夠呢?我是否願意成為基督的預備者呢?

福音書的經課(可1:1)就以「耶穌基督福音的起頭」作為一切內容的開始,無錯,馬可福音第一個字就是「起頭」,耶穌基督就是福音的起頭,這一宣告可以話係全馬可福音的核心,而我們有沒有為這福音的到來而做好準備呢?到底甚麼是福音呢?對現今的世代,我們一直所信奉,所傳揚的真的是大家所嚮往,所需要的「福音」嗎?耶穌是「基督」,是「上帝的兒子」對這一代的人是「福音」嗎?每年臨近聖誕,當大家也盡情「普天同慶地」消費,享受時,相信在將臨期的日子中,我們更要思考這一個問題…

馬可福音的作者刻意以回應舊約的呼籲作為預備的開始,

舊約經課以賽亞書40:1-2

你們的上帝說:「要安慰,安慰我的百姓.要對耶路撒冷說安慰的話,向它宣告,它的戰爭已結束,它的罪孽已赦免；它為自己一切的罪,已從耶和華手中加倍受罰.」

以安慰為宣告的內容:戰爭結束,罪孽蒙赦,惡已經加倍受罰；代表著一個時代的終結,一切要過去,上帝讓它重新開始,一切不再一樣；這就是先知所宣告的「福音」,我們需要這「福音」嗎?還是更多人想繼續現有的生活,或是期望回到過去,福音是要為人/世界帶來改變,不是按我們的喜好,而是按上帝的心意而行.

以賽亞書40:3-5

有聲音呼喊著:「要在曠野為耶和華預備道路,在沙漠為我們的上帝修直大道.一切山窪都要填滿,大小山岡都要削平；陡峭的要變為平坦,崎嶇的必成為平原.耶和華的榮耀必然顯現,凡有血肉之軀的都一同看見,因為這是耶和華親口說的.」

預備者的工作就是要負責把一切阻礙除去,外在的就是要除去人與上帝中的任何阻隔,內在的就是要認罪悔改,最終要讓每一個人都見證到上帝的榮耀,很清楚,很直接,很簡單的；在今天我們又如何實踐呢?我們有盡力為身邊的人除去行近上帝的阻隔嗎?盡我們的能力,方法帶領人到上帝面前嗎?還是覺得責任不在我?甚至用很多「理由」去合理化自己敵擋上帝的行動嗎?

以賽亞書40:6-11

有聲音說:「你喊叫吧!」我說:「我喊叫甚麼呢?」凡有血肉之軀的盡都如草,他的一切榮美像野地的花.耶和華吹一口氣,草就枯乾,花也凋謝.百姓誠然是草；草必枯乾,花必凋謝,惟有我們上帝的話永遠立定.報好信息的錫安哪,要登高山；報好信息的耶路撒冷啊,要極力揚聲.揚聲不要懼怕,對猶大的城鎮說:「看哪,你們的上帝!」看哪,主耶和華必以大能臨到,他的膀臂必為他掌權；看哪,他的賞賜在他那裡,他的報應在他面前.他要像牧人牧養自己的羊群,用膀臂聚集羔羊,抱在胸懷,慢慢引導那乳養小羊的.

馬可回應先知的預言,除了以賽亞先知之外,還可能有瑪拉基,馬可福音原文中(1:2)沒有直接指是以賽亞.

瑪拉基書3:1

萬軍之耶和華說:「看哪,我要差遣我的使者在我前面預備道路.你們所尋求的主必忽然來到他的殿；立約的使者,就是你們所仰慕的,看哪,快要來到.」

預備者要為主向世人喊叫,喊叫的內容是「看哪,你們的上帝!」與及「看哪,快要來到!」,原來我們作為上帝的預備者,就是要人去轉目看見上帝,世人因很多不同的原因看不見上帝,不是因為上帝隱藏了,上帝極其偉大,應該沒有可能看不見的,但人往往選擇「看不見」,明明人活在有限的能力,經歷與及軟弱的生命之中,如野地的花,經風一吹,使歸無有…相對於永恆全權的上帝,我們是何等的微小；所以人應當認識上帝永遠立定的話,等待祂大能臨到,膀臂掌權,在祂懷中被乳養；這就是要盡全力去喊叫,去呼籲的原因.

馬可福音1:2-4

正如以賽亞先知書上記著:「看哪,我要差遣我的使者在你面前,他要為你預備道路.在曠野有聲音呼喊著:預備主的道,修直他的路.」照這話,施洗約翰來到曠野…

馬可記述因回應先知們,期望我們作為讀者的也要轉目去看,上帝要親自差遣使者到人的面前,今天上帝是否仍然差遣我們到他人的面前呢?

呼喊的內容直接清晰仔細,就是要「預備主的道,修直他的路」,施洗者約翰就以行動回應上帝藉先知的呼籲,為願意悔改者進行悔改的洗禮；

馬可福音1:5

猶太全地和全耶路撒冷的人都出去,到約翰那裡,承認他們的罪,在約旦河裡受他的洗.

聖經描述「猶太全地」,「全耶路撒冷」的人都以行動回應,只要我們願意依從上帝的吩咐,為主作出喊叫,呼籲(傳揚福音),人就會回應(是回應主,不是因為我們),簡單,直接,不需要計算,思考太多,因為這往往是撒旦阻隔人認識上帝,阻礙我們為主預備的技倆.

我們如何能夠成為忠心的預備者呢?聖經接著這樣描述施洗者約翰…

馬可福音1:6-8

約翰穿駱駝毛的衣服,腰束皮帶,吃的是蝗蟲和野蜜. 他宣講,說:「有一位在我以後來的,能力比我更大,我就是彎腰給他解鞋帶也不配.我用水給你們施洗,他卻要用聖靈給你們施洗.」

與舊約中以利亞的形象很類似,帶著權柄宣告的約翰,也有一些學者認為施洗者約翰可能是愛色尼派的成員,無論衣著外表及所宣告內容也非常突出,當然聖經不是教導我們傳福音時要「標奇立異」以「嘩眾取寵」,但卻要謹記信徒要分別成聖,要記得施洗者約翰在母親伊利莎白的肚腹內就與馬利亞肚腹內的耶穌共鳴,並被聖靈充滿；施洗者約翰清楚自己的角色及位置,這也是身為預備者最首要的,他清楚要做甚麼,而盡力去完成；同時又清楚他不是甚麼,有甚麼不能/不會做,真正謙卑的人,有清晰位置並願意順服；施洗者約翰認識聖靈,並知道耶穌基督為人帶來聖靈的洗是他沒有能力做到的,但悔改的水禮就能為人迎見聖靈的洗作預備,施洗者約翰就是這樣為耶穌基督預備…

請留意:馬可福音開始第一句的描述:「上帝的兒子,耶穌基督福音的起頭.」基督福音是從何時開始的呢?不同人有不同的理解,各福音書的作者亦有不同的切入點,馬可清晰地告訴我們,他所理解基督福音的起頭就是由願意回應舊約先知的預言,以行動成為基督的預備者施洗約翰的呼喚開始.今天基督的福音在我們身邊又如何被開展的呢?有誰願意成為這個世代,與及我們身邊人的約翰呢?

我們今天要為這世代預備甚麼呢?

經課詩篇85:1-2

耶和華啊,你已經向你的地施恩,救回被擄的雅各.你赦免了你百姓的罪孽,遮蓋了他們一切的過犯.

當我們要回答「我們要預備甚麼」時,請先回想,是上帝先向我們施恩.

約翰一書4:19-20

我們愛,因為上帝先愛我們.不是我們愛上帝,而是上帝愛我們,差他的兒子為我們的罪作了贖罪祭；這就是愛.

羅馬書5:6-8

因我們還軟弱的時候,基督就按所定的日期為罪人死.為義人死,是少有的；為仁人死,或者有敢做的.惟有基督在我們還作罪人的時候為我們死,上帝的愛就在此向我們顯明了.

耶穌基督為我們犧牲,擺上予一切,我們唯一能夠做的,就是盡一切去回應上帝的恩典,除了自己回應之外,就是為身邊的人預備最有助回應的環境,空間…

經課彼得後書3:8-9

親愛的,有一件事你們不可忘記,就是:主看一日如千年,千年如一日. 主沒有遲延他的應許,就如有人以為他是遲延,其實他是寬容你們,不願一人沉淪,而是人人都來悔改.

上帝的心意希望人人都悔改,我們的心意又如何呢?如何能夠做到人人悔改呢?我們的角色,位置又是甚麼呢?主的豐盛宴席已經擺開了,並且一直在等待著人的回應,我們又如何為主預備及修直道路呢?

彼得後書3:10

但主的日子要像賊一樣來到；那日,天必在轟然一聲中消失,天體都要被烈火熔化,地和地上的萬物都要燒盡…

將臨期帶給我們其中一個很重要的信息,主會再來,必定會再來,並且必定在我們想像不到的情況下再來,所以我們自身除了要謹慎自守,警醒禱告,更要把握機會,盡預備者的職份,為身邊的人掃除阻隔,把山窪填滿,山岡削平,要下苦工,要面對壓力,甚至有一天如施洗者約翰一樣,為此獻上生命,你真的預備好了嗎?

每一次在田徑賽上聽到On your marks, get set時,那種全身每一個細胞都集中,毛孔擴張,隨時聽候「GO」就彈出去的狀態；或許可以是描繪現今作信徒如何預備回應主應有的心態的圖畫.

現在是行動的時候了,聖誕節將到,作為信徒與及追求信仰的我們,是最好的機會去行動的日子…



\newpage

\section{終極好消息}
\label{sec:530}
\textbf{王家輝牧師 - 終極好消息}
\newline
\newline
\hyperref[sec:529]{< < < PREV SERMON < < <}
~
\hyperref[sec:toc]{[返目錄]}
~
\hyperref[sec:531]{> > > NEXT SERMON > > >}


以賽亞書 61:1-4,8,11
\newline
\begin{longtable}{cl}
\hline
\hline
章節 & 經文 (和合本修訂版)\\
\hline
61:1 & \begin{tabularx}{0.7\textwidth}{X} 主耶和華的靈在我身上，因為耶和華用膏膏我，叫我報好信息給貧窮的人，差遣我醫好傷心的人，報告被擄的得釋放，被捆綁的得自由； \end{tabularx} \\ \\ \relax
61:2 & \begin{tabularx}{0.7\textwidth}{X} 宣告耶和華的恩年和我們的神報仇的日子；安慰所有悲哀的人， \end{tabularx} \\ \\ \relax
61:3 & \begin{tabularx}{0.7\textwidth}{X} 為錫安悲哀的人，賜華冠代替灰燼，喜樂的油代替悲哀，讚美為衣代替憂傷的靈；稱他們為「公義樹」，是耶和華所栽植的，為要彰顯他的榮耀。 \end{tabularx} \\ \\ \relax
61:4 & \begin{tabularx}{0.7\textwidth}{X} 他們必修造久已荒涼的廢墟，建立先前淒涼之處，重修歷代荒涼之城。 \end{tabularx} \\ \\ \relax
61:5 & \begin{tabularx}{0.7\textwidth}{X} 那時，陌生人要伺候、牧放你們的羊群；外邦人必為你們耕種田地，修整你們的葡萄園。 \end{tabularx} \\ \\ \relax
61:6 & \begin{tabularx}{0.7\textwidth}{X} 但你們要稱為「耶和華的祭司」，稱作「我們神的僕人」。你們必享用列國的財物，必承受他們的財富。 \end{tabularx} \\ \\ \relax
61:7 & \begin{tabularx}{0.7\textwidth}{X} 因為他們所受雙倍的羞辱，凌辱被稱為他們的命運，因此，他們在境內必得雙倍的產業，永遠之樂必歸給他們。 \end{tabularx} \\ \\ \relax
61:8 & \begin{tabularx}{0.7\textwidth}{X} 因為我－耶和華喜愛公平，恨惡搶奪與惡行；我要憑誠實施行報償，與我的百姓立永約。 \end{tabularx} \\ \\ \relax
61:9 & \begin{tabularx}{0.7\textwidth}{X} 他們的後裔必在列國中為人所知，他們的子孫在萬民中為人所識；凡看見他們的必承認他們是耶和華所賜福的後裔。 \end{tabularx} \\ \\ \relax
61:10 & \begin{tabularx}{0.7\textwidth}{X} 我因耶和華大大歡喜，我的心因神喜樂；因他以拯救為衣給我穿上，以公義為外袍給我披上，好像新郎戴上華冠，又如新娘佩戴首飾。 \end{tabularx} \\ \\ \relax
61:11 & \begin{tabularx}{0.7\textwidth}{X} 地怎樣使芽長出，園子怎樣使所栽種的生長，主耶和華也必照樣使公義和讚美在萬國中發出。 \end{tabularx} \\ \\
[1ex]
\hline
\hline
\end{longtable}
講題:終極好消息

經文:以賽亞書61:1-4,8,11

講員:王家輝牧師

場合:香港中文大學崇基學院禮拜堂主日崇拜

日期:2020年12月13日

經文:

61:1主耶和華的靈在我身上,因為耶和華用膏膏我,叫我報好信息給貧窮的人,差遣我醫好傷心的人,報告被擄的得釋放,被捆綁的得自由；2宣告耶和華的恩年和我們的上帝報仇的日子；安慰所有悲哀的人,3為錫安悲哀的人,賜華冠代替灰燼,喜樂的油代替悲哀,讚美為衣代替憂傷的靈；稱他們為「公義樹」,是耶和華所栽植的,為要彰顯他的榮耀.4他們必修造久已荒涼的廢墟,建立先前淒涼之處,重修歷代荒涼之城.

61:8因為我－耶和華喜愛公平,恨惡搶奪與惡行；我要憑誠實施行報償,與我的百姓立永約.9他們的後裔必在列國中為人所知,他們的子孫在萬民中為人所識；凡看見他們的必承認他們是耶和華所賜福的後裔.10我因耶和華大大歡喜,我的心因上帝喜樂；因他以拯救為衣給我穿上,以公義為外袍給我披上,好像新郎戴上華冠,又如新娘佩戴首飾.11地怎樣使芽長出,園子怎樣使所栽種的生長,主耶和華也必照樣使公義和讚美在萬國中發出.

引言

相信大家也聽過西諺有這一句說話:「No news is good news」,直譯的意思就是「沒有消息就是好消息」.此語究竟出自何方呢?坊間有不同的流傳.其中一個流傳跟英王詹姆士一世(King James I, 1566-1625)有關(他就是下令編纂《欽定版聖經》(King James Version)的那位英王),提到他在1616年一封書信中寫下這一句:「No news is better than evil news」,就是沒有消息比起壞消息更好.另一個流傳跟美國內戰或稱南北戰爭(1861-1865)有關,在參戰的將士陣亡後,其家人會收到通知,所以將士家庭都害怕得到噩耗；相反,如果沒有什麼消息,則至少還有活著的希望.

今時今日,我們也許同樣有類似的期待,每日看見新聞報導,有多少宗確診個案,然後又要再留意自己所住的大廈,屋苑或工作地點是否有累積個案.倘若出現累積的確診個案,可能就要馬上接受強制檢測,在未有結果前也會擔心自己染病.在這樣的情況下,不單是香港,甚至乎是全世界也認同”No news is good news”.儘管現實如此,然而這是消極的想法,不想得到壞消息,讓自己少一份擔憂.若從積極的方向想,我們期待什麼好消息呢?這個不單是生活經驗,更是信仰的思考.

「好消息」以信仰的語言來翻譯就是「福音」,也是我們在這段日子 ─ 將臨期,準備帶出的信息.今日我想跟大家集中以舊約經課的以賽亞書61章1至4節及8至11節作這樣的思考.

上主的受膏者

眾所周知,一個有效的訊息傳遞至少由三個部份組成:傳訊者,訊息和接收者.在以賽亞書61:1上「主耶和華的靈在我身上,因為耶和華用膏膏我」,清晰地介紹了耶和華預備了祂的傳訊者,亦是一位執行使命的僕人.他有主的靈,亦受膏立.

在舊約聖經中,受膏是一項特別的儀式,而受膏者(在聖經中稱為「彌賽亞」)是分別為聖,按上主吩咐行事.在出埃及記30:22-33記述了亞倫和他兒子受膏,成為耶和華的祭司.在撒母耳記上16:1-13記述了撒母耳膏立大衛為以色列的王.在列王記上19:15-18記述了以利亞奉上主吩咐,要膏以利沙成為先知,接續他的工作.換言之,受膏者可以是祭司,君王和先知,身份雖然有分別,但重點就是聽憑上主差遣的僕人.按照對彌賽亞身份的認識,路加福音4:18-19記述了耶穌在會堂裡,宣讀以賽亞書61:1-2,從而顯示祂同時扮演祭司,君王,先知的角色.

今日誰是上主的受膏者?誰把從上主而來的訊息傳遞出去?誰會按上主的吩咐行事?那當然是教會,亦是我和各位弟兄姊妹.我們成為了上主的受膏者,分別為聖,被主差遣.

從上主而來的訊息與使命

究竟上主有什麼訊息?受膏者受託什麼使命?以賽亞61:1下-3上清楚羅列出來:

報好信息給貧窮的人；

醫好傷心的人；

報告被擄的得釋放,被捆綁的得自由；

宣告耶和華的恩年和我們的上帝報仇的日子；

安慰所有悲哀的人,

為錫安悲哀的人,賜華冠代替灰燼,喜樂的油代替悲哀,讚美為衣代替憂傷的靈；稱他們為「公義樹」,是耶和華所栽植的,為要彰顯他的榮耀.

我們可以想像,什麼人和什麼景況需要這樣的訊息和行動?當然是貧窮的人,傷心的人,被擄,被捆綁的人,悲哀的人.普遍聖經學者也認同以賽亞書56章至66章的訊息是特別針對被擄後的以色列群體,那麼這群人所面對的處境與心境,的確跟先知所言的是吻合的.這樣的處境到新約時代,跟耶穌一起生活的猶太人,他們的遭遇也類近.在羅馬帝國強勢管治下,猶太人在反抗不成功後,形勢更惡劣.再加上貴族跟宗教領袖大致上也樂意配合帝國的管治,故此一般人的生活也極不容易.

活在今日的香港,這裡已不像以往被看為一片樂土.過往有不少人不理政治,以為憑藉個人的勤奮就可以為自己締造美好的生活.可是今日就是你不理政治,政治也滲入生活每個細節.你在社交媒體留下的一句說話或一張合照,也可能為你帶來麻煩.我們仍然可以自由呼吸,但每日不斷接收令人沮喪的消息時,卻教人感到心情沉重,甚至有快要窒息的感覺.

我們期待好消息,而且是終極的好消息!這好消息不是由誰人成為特首或美國總統,因為經驗告訴我們另一個不一定比現在這個好.這好消息亦不是那個藥廠研發出最有效用的疫苗,因為處理到一種新病毒卻不能保證將來不會出現更加新型或變種病毒.倘若人的自私,霸權,惡行不被消滅,根本沒有好消息可言.

上主的受膏者傳遞的訊息是人所面對的一切困難,人無法自救,而上主是知道的,唯有祂親自介入,親自拯救,我們才有脫困,獲救的盼望.上主的介入,以道成肉身住在我們中間,具體地跟人一起生活,明暸人的限制與掙扎,而祂亦預備了拯救的方式.這個訊息需要由我們這群屬於基督的人去傳開,亦從我們的生活去展現出來.

上主會以什麼方式解決我們眼前的困境,我沒有水晶球,所以無法預知和預測.我只知道,當面對困難時,一群人去面對總會比獨自面對更加容易.作為主的僕人,我知道亦選擇留下來繼續回應主託付的使命,但需要弟兄姊妹一起同行.

在這樣的情況下,我呼籲弟兄姊妹,若是可行請你留在香港.不是因為現在航空交通未恢復,所以你才留下,就是可能在半年或一年後,一切運作也能恢復後,請你選擇留下來.我深信上主在此時此地,仍有祂的心意在其中,有許多人需要你的援助,關懷.這個地方的未來仍然有你可以貢獻之處.

前路雖然不明朗,可是上主卻有應許.以賽亞書61:2「宣告耶和華的恩年和我們的上帝報仇的日子.」上主的受膏者就是在這樣特殊的情況下,領受了訊息和使命.

人的反應和行動

當上主的訊息傳遞出去,而受膏者亦執行使命時,百姓會有什麼反應和相應的行動呢?以賽亞書61:3下-4「他們為「公義樹」,是耶和華所栽植的,為要彰顯他的榮耀.他們必修造久已荒涼的廢墟,建立先前淒涼之處,重修歷代荒涼之城.」

身心受創的人在得到醫治,安慰後,他們回復生命的尊嚴.他們的生命能夠反映上主的榮耀,因為人本來就是按上主的形象造成,反映著上主的榮耀.不但是身心復原,更讓被毀的家園得以重建.這個獲重建的家園,並不只在硬件上得到重建,而是一些內在重要的價值也獲重建.這些價值是按上主的屬性來重建的,當中包括:公平,排除搶奪與惡行,施行報償.

許多時我們只會著眼硬件的建築,因為這是看得見,也令人感到有某種成就.尤其是以色列人對於聖殿被毀與重建聖殿也非常重視,因為聖殿象徵上主的同在.可是不能見的價值卻是整個城市,國家,民族的基石,若忽略了內裡價值的重建,就是再宏偉,華麗的建築物,只是空殼一個!今天我們著緊的又是否只在看得見的硬件,而忽略上主所看重的內在價值呢?

上主希望我們可以在祂所造和喜愛的地方生活,那並不是一個只追求物質生活,並且容許擔奪與惡行所建立的地方.那是一個使人看見上主的旨意能夠實現的地方.這樣的地方使人得到真正的滿足和歡樂,那是在地若天,是上主的國度臨在的境界.這就是上主給我們終極的好消息,也是我們需要聆聽並傳揚的訊息.求主恩助.阿們.



\newpage

\section{祂顧念窮人的卑微}
\label{sec:531}
\textbf{鄧瑞強博士 - 祂顧念窮人的卑微}
\newline
\newline
\hyperref[sec:530]{< < < PREV SERMON < < <}
~
\hyperref[sec:toc]{[返目錄]}
~
\hyperref[sec:532]{> > > NEXT SERMON > > >}


路加福音 1:46-55
\newline
\begin{longtable}{cl}
\hline
\hline
章節 & 經文 (和合本修訂版)\\
\hline
1:46 & \begin{tabularx}{0.7\textwidth}{X} 馬利亞說：「我心尊主為大； \end{tabularx} \\ \\ \relax
1:47 & \begin{tabularx}{0.7\textwidth}{X} 我靈以神我的救主為樂； \end{tabularx} \\ \\ \relax
1:48 & \begin{tabularx}{0.7\textwidth}{X} 因為他顧念他使女的卑微；從今以後，萬代要稱我有福。 \end{tabularx} \\ \\ \relax
1:49 & \begin{tabularx}{0.7\textwidth}{X} 因為那有權能的為我做了大事；他的名是聖的。 \end{tabularx} \\ \\ \relax
1:50 & \begin{tabularx}{0.7\textwidth}{X} 他憐憫敬畏他的人，直到世世代代。 \end{tabularx} \\ \\ \relax
1:51 & \begin{tabularx}{0.7\textwidth}{X} 他用膀臂施展大能；他趕散心裡妄想的狂傲人。 \end{tabularx} \\ \\ \relax
1:52 & \begin{tabularx}{0.7\textwidth}{X} 他叫有權柄的失位，叫卑賤的升高。 \end{tabularx} \\ \\ \relax
1:53 & \begin{tabularx}{0.7\textwidth}{X} 他叫飢餓的飽餐美食，叫富足的空手回去。 \end{tabularx} \\ \\ \relax
1:54 & \begin{tabularx}{0.7\textwidth}{X} 他扶助了他的僕人以色列，不忘記施憐憫， \end{tabularx} \\ \\ \relax
1:55 & \begin{tabularx}{0.7\textwidth}{X} 正如他對我們的列祖說過，『憐憫亞伯拉罕和他的後裔，直到永遠。』」 \end{tabularx} \\ \\
[1ex]
\hline
\hline
\end{longtable}
講題:祂顧念窮人的卑微

經文:路加福音1:46-55

講員:鄧瑞強博士(香港神學院)

場合:香港中文大學崇基學院禮拜堂主日崇拜

日期:2020年12月20日

各位弟兄姊妹,早安.在這疫症裡,祝願各人平安.

今日是「將臨期」的第四主日.「將臨期」是期盼的節期,期盼一些重要的事情發生.期盼,是人類的存在特質,能超出當下的現實,去想像一個還未出現的將來.我們盼望,疫症快點過去,我們再能無恐懼地生活.我們盼望,冬天的寒夜早點過去,生命在春天中再現生機.我們盼望,塵世的黑暗早點過去,天上的光明普照人間.這些都是夢,毫不現實,但人類,作為唯一懂得期盼的動物,卻敢於想像這毫不現實,卻支撐著人類走下去的遠象.我們期盼,有一天,公義是社會的根基,仁愛是人與人交往的內容,並且,神在人間行走.

這些都是夢.「將臨期」是發這些夢的日子,聖誕節是這夢境成真.

今日的講道經文,是路加福音中出名的「馬利亞頌」(Magnificat),看看夢境如何成真:路加福音1:46-55.

路1:46　馬利亞說:我心尊主為大；

路1:47　我靈以神我的救主為樂；

路1:48　因為他顧念他使女的卑微；從今以後,萬代要稱我有福.

路1:49　那有權能的,為我成就了大事；他的名為聖.

路1:50　他憐憫敬畏他的人,直到世世代代.

路1:51　他用膀臂施展大能；那狂傲的人正心裡妄想就被他趕散了.

路1:52　他叫有權柄的失位,叫卑賤的升高；

路1:53　叫飢餓的得飽美食,叫富足的空手回去.

路1:54　他扶助了他的僕人以色列,

路1:55　為要記念亞伯拉罕和他的後裔,施憐憫直到永遠,正如從前對我們列祖所說的話.

這裡有一個小女孩,她叫馬利亞,一個在猶太人中普通得不再普通的名字.她看自己是誰?她看自己為「使女」,即「女奴隸」,一個低下得不能再低下的人.她生長在窮鄉僻壤,一個叫「拿撒勒」的地方.稍有身分的猶太人,都會說,拿撒勒能出什麼好東西來?她的童年,是貧窮的.現在,她也一無所有.將來,對她而言,哪有什麼將來可言.她的過去,沒有什麼值得提起的.將來,她將湮沒在無聲無息之中.這樣的一個人,誰也看不到她.

然而,在人類的想像力以外,想像不到的事,發生在她身上.在這段經文裡,那稱為「主,神,救主,有權能的(大能者),聖者」的那位,那位超越人類想像的神,看到她,認識她,猶如在森林中看到地上的一棵小草.神要透過這柔弱的生命,創造出千古未有的神蹟.要在萬古的黑暗中,藉著她,射進從高天而來的光明.一個小女孩,點燃了人類希望的火把.

她說:

路1:46　我心尊「主」為「大」；

路1:47　我靈以「神」我的「救主」為樂.

她述說著一個超凡的體驗,一種神聖的相遇,一個深度的真相.她體會到一無以名狀的偉大生命,她稱之為「主」.這是她生命的主人,她賴以為生的倚靠對象,她生命之所托.她形容之為「大」,在這「偉大」面前,人只是卑微的蟲.

「偉大」與「微小」的相遇,結果可能只是「偉大」顯得更「偉大」,「微小」顯得更「微小」.莊子講了一個故事:世間有大鵬鳥,鳥的背像泰山,翅膀像天邊的雲,能扶搖直上九萬里,背負起蒼天.地上卻有隻小麻雀,自鳴得意,認為自己能在草叢間跳來跳去已經很了不起.莊子的按語是:「此小大之辯也」.這就是「偉大」之為「偉大」,「微小」之為「微小」的分別所在.「偉大」與「微小」的相遇,可能只成就一種哲理,卻不成就一種信仰.

卑微的馬利亞體會的,卻是深奧的神聖事件.她稱這偉大者為「神」,為救助她生命的「救主」.馬利亞與神的相遇,並不只是顯出她的卑微；與這偉大者相遇,她體會到生命的救贖,她在其中湧溢出歡樂.她繼續說:

路1:48　因為他顧念他使女的卑微；從今以後,萬代要稱我有福.

路1:49　那有權能的,為我成就了大事；他的名為聖.

這位偉大者,不只是「曬命」似的顯出他的偉大,而是以一種超乎想像的方式去顯明他的偉大.他的偉大,在於他「顧念」這位女奴隸的卑下.「顧念」的希臘文,有「注視,看得起,尊重,眷顧」的意思.至偉大者,看得起至卑賤者；猶如一個君王,向一個乞丐行禮一樣.這是超乎想像的,這是超乎人類期望的.與馬利亞相遇的,不是那種不理人死活的偉大者,而是一位顧念卑下者的,紆尊降貴的,救助世人的救主.

這位「有權能的」大能者,他不是為了弦耀他的大能而行大事,馬利亞說:他「為我成就了大事」.這位至偉大的大能者,是刻意「為了」一個卑下的人而行事；猶如一個長官,遞一杯涼水給一個坐監的人喝一樣.「有權能的」,不是為了自己做什麼,而是為了有需要的人做一點什麼.聖誕節故事,偉大的神降生在馬槽裡,是重新注解什麼叫「偉大」,什麼叫「有權能」.

馬利亞說:「他的名為聖」.神聖,在聖經裡,指向那位與「被造的萬物」全然不一樣的創造主.「神聖」,就是與萬物有別,所以聖經常常說「分別為聖」,屬於神聖的,就必與一般事物有分別.神學家巴特(Karl Barth)說:神是全然的他者(wholly other),與我們這些被造物全然不同.這講法很能講明什麼叫「神聖」.然而,如此神聖的偉大者,他的偉大和神聖,卻表現在他「看得起」,「顧念」一個至卑微的奴隸,並「為她」而有所行動.聖誕節的故事,真奇特,至高者與至卑賤者走在一起,並且,至高者為至卑賤者提供服務,猶如主耶穌為門徒洗腳一樣.

神為世上的貧窮者,為奴者,卑賤者做什麼事呢?

路1:50　他憐憫敬畏他的人,直到世世代代.

路1:51　他用膀臂施展大能；那狂傲的人正心裡妄想就被他趕散了.

路1:52　他叫有權柄的失位,叫卑賤的升高；

路1:53　叫飢餓的得飽美食,叫富足的空手回去.

神大能的膀臂,要做什麼呢?

他要將狂傲的人趕散.「趕散」這個希臘文,有用風將糠秕吹散的意思.在神眼中,狂傲人是無用的糠秕,要用風將他們篩走.在神眼中,看為重要的,是那些被狂傲人欺壓的,手無寸鐵的人,神要抬舉他們.

這位至大的神,要扭轉乾坤,要改變人間墮落的秩序.世間有悲慘的人,就因為有權柄實踐正義的人只是袖手旁觀.世間有飢餓的人,就因為資源沒有公平地分配.神要叫有權柄的失位,叫富足的空手而回；相反,神叫卑賤的升高,叫飢餓的得飽美食.在聖誕節,神要人正視的問題是:為何人間有人被欺壓,而活得毫無尊嚴?為何有人「飲飽食醉」卻有人凍臥街頭?這個世界究竟出了什麼問題?

前年,香港的貧窮人口是140.6萬,＜2019年香港貧窮情況＞報告還未發表,今年的情況看來會更差.香港人口是750萬,貧窮人口約佔18.75\%,頗嚴重.另一方面,按「地球之友」網頁的資料,香港超巿丟棄的而仍可食用的食物,每天29公噸.(\url{http://foodwaste.foe.org.hk/html/chi/topic.php?id=39})食物足夠嗎?肯定足夠.有飢餓的人嗎?大有人在.

卑微的馬利亞說,神來到人間,看到的,顧念的,不是高官巨賈,而是卑賤的和飢餓的人.神來到人間,不是住在皇宮裡,而是安放在簡陋的馬槽裡.神來到人間,是希望這個世界有一真的變革.權力的結構要改變,不是有權者享盡一切利益,而是懂得讓利,分利,將自己的權力看成為顧念弱者的責任,將自己的富足理解為幫助社會總體脫貧的工具.有些社會制度,強調把持權力.有些社會制度,強調操控金錢.神來到人間,針對的,是權力和金錢的誤用.他要叫有權柄的失位,叫富足的空手回去,因權力的誤用和金錢的誤用帶來人間極大的痛苦.聖誕節,不是只有溫馨,同時也要求社會的變革.馬利亞這裡講的,是神要進行的激烈的變革,將扭曲的權力結構和金錢結構摧毀,重建一個人性的世界.在這世界裡,再沒有被強權壓下的人,再沒有被貧窮蠶食的人.相反,地位低下的人,他們的要求會被正視；貧窮的人,他們的渴望會被聆聽.統治世間的,不再是權力和金錢,而是仁愛和同情.聖誕節,與其說是一慶祝的節日,不如說是一呼喚我們悔改,進行變革,實現人間新秩序的日子.

馬利亞的禱文,結尾兩句是:

路1:54　他扶助了他的僕人以色列,

路1:55　為要記念亞伯拉罕和他的後裔,施憐憫直到永遠,正如從前對我們列祖所說的話.

提起亞伯拉罕,就是提起神與人立的約.所謂「立約」,就是將原先毫無關係的兩方,連繫在一起,只要雙方守約,雙方就要保護對方.而神與人立的約,是一不平等的約,是對神不平等的.提摩太後書2:13說:「我們縱然失信,他仍是可信的,因為他不能背乎自己.」在神人的約裡,人不守約,人離棄神,神依然守約,用力扭轉乾坤,這就是神的「信實」.在這神與人的約裡,神永遠將人的命運連於自己的存在.神與人立約之後,眼目就不會離開人.對使人受苦的人,神必定追究他們的罪行.對受苦的人,神要扶起他們,醫治他們,堅立他們.在約裡,神要透過他的僕人,將福氣帶給世人.這就是卑微的馬利亞看到的景況,這就是聖誕節故事的內容,神來到人間,連結於人類的悲慘故事,以仁愛與正義重整世間的秩序.基於神與人的約,神來人間,看顧悲慘的人,他看到馬利亞.神注視她,一個卑微的女奴隸.神看得起她,不計較她的卑微.要透過她,讓耶穌生在世上.藉此,神抬舉卑微的人,讓他們有份與人類的救世工程.在馬利亞的故事裡,我們看到神與卑微者的親近,看到神多麼願意讓卑微者參與,透過他們的努力,去改變這個世界.

將臨期是期待的日子,這很切合當下的情境.香港人在期待,世人在期待.我們期待神聖的臨在,我們期待神來到人間.我們期待有權柄者,在嬰孩耶穌面前,明白什麼叫「偉大」.我們期待有財富者,在嬰孩耶穌面前,明白什麼叫「富有」.我們期待窮人得飽足,我們期待被剝奪人性尊嚴的人能重拾尊嚴.我們期待無權無勢者能體會自己的尊嚴,體會自己有份參與在神的救世工程中.我們期待一正義和光明的將來.簡言之,我們期待基督.

或許有人說:這一切只是痴心妄想.

將臨期就是敢於作這種痴心妄想,而聖誕節就是這種痴心妄想的實現.

在寒夜裡,請不要放棄期待.在黑暗裡,請繼續痴心妄想著光明.在看不到前景的深淵裡,請想像著神來到人間.這不可能的事件,在寒冷的聖誕節中實現了.

願榮耀歸上主,願人間享和平.



\newpage

\section{不凡的竟然如此平凡}
\label{sec:532}
\textbf{林豪恩先生 - 不凡的竟然如此平凡}
\newline
\newline
\hyperref[sec:531]{< < < PREV SERMON < < <}
~
\hyperref[sec:toc]{[返目錄]}
~
\hyperref[sec:533]{> > > NEXT SERMON > > >}


路加福音 2:22-40
\newline
\begin{longtable}{cl}
\hline
\hline
章節 & 經文 (和合本修訂版)\\
\hline
2:22 & \begin{tabularx}{0.7\textwidth}{X} 按摩西律法滿了潔淨的日子，他們就帶著孩子上耶路撒冷去，要把他獻給主。 \end{tabularx} \\ \\ \relax
2:23 & \begin{tabularx}{0.7\textwidth}{X} 正如主的律法上所記：「凡頭生的男子必歸主為聖」； \end{tabularx} \\ \\ \relax
2:24 & \begin{tabularx}{0.7\textwidth}{X} 又要照主的律法上所說，用一對斑鳩，或用兩隻雛鴿獻祭。 \end{tabularx} \\ \\ \relax
2:25 & \begin{tabularx}{0.7\textwidth}{X} 那時，在耶路撒冷有一個人，名叫西面；這人又公義又虔誠，素常盼望以色列的安慰者來到，又有聖靈在他身上。 \end{tabularx} \\ \\ \relax
2:26 & \begin{tabularx}{0.7\textwidth}{X} 他得了聖靈的啟示，知道自己未死以前必看見主所立的基督。 \end{tabularx} \\ \\ \relax
2:27 & \begin{tabularx}{0.7\textwidth}{X} 他受了聖靈的感動，進入聖殿，正遇見耶穌的父母抱著孩子進來，要照律法的規矩而行。 \end{tabularx} \\ \\ \relax
2:28 & \begin{tabularx}{0.7\textwidth}{X} 西面就把他抱過來，稱頌神說： \end{tabularx} \\ \\ \relax
2:29 & \begin{tabularx}{0.7\textwidth}{X} 「主啊，如今可以照你的話，容你的僕人安然去世； \end{tabularx} \\ \\ \relax
2:30 & \begin{tabularx}{0.7\textwidth}{X} 因為我的眼睛已經看見你的救恩， \end{tabularx} \\ \\ \relax
2:31 & \begin{tabularx}{0.7\textwidth}{X} 就是你在萬民面前所預備的： \end{tabularx} \\ \\ \relax
2:32 & \begin{tabularx}{0.7\textwidth}{X} 是啟示外邦人的光，是你民以色列的榮耀。」 \end{tabularx} \\ \\ \relax
2:33 & \begin{tabularx}{0.7\textwidth}{X} 孩子的父母因論耶穌的這些話就驚訝。 \end{tabularx} \\ \\ \relax
2:34 & \begin{tabularx}{0.7\textwidth}{X} 西面給他們祝福，又對孩子的母親馬利亞說：「這孩子被立，是要叫以色列中許多人跌倒，許多人興起；又要成為毀謗的對象， \end{tabularx} \\ \\ \relax
2:35 & \begin{tabularx}{0.7\textwidth}{X} 叫許多人心裡的意念顯露出來；你自己的心也要被劍刺透。」 \end{tabularx} \\ \\ \relax
2:36 & \begin{tabularx}{0.7\textwidth}{X} 又有位女先知，名叫亞拿，是亞設支派法內力的女兒，年紀已經老邁，從童女出嫁，同丈夫住了七年， \end{tabularx} \\ \\ \relax
2:37 & \begin{tabularx}{0.7\textwidth}{X} 就寡居了，現在已經八十四歲。她不離開聖殿，禁食祈求，晝夜事奉神。 \end{tabularx} \\ \\ \relax
2:38 & \begin{tabularx}{0.7\textwidth}{X} 正當那時，她進前來感謝神，對一切盼望耶路撒冷得救贖的人講論這孩子的事。 \end{tabularx} \\ \\ \relax
2:39 & \begin{tabularx}{0.7\textwidth}{X} 約瑟和馬利亞照主的律法辦完了一切的事，就回加利利，到自己的城拿撒勒去了。 \end{tabularx} \\ \\ \relax
2:40 & \begin{tabularx}{0.7\textwidth}{X} 孩子漸漸長大，強健起來，充滿智慧，又有神的恩典在他身上。 \end{tabularx} \\ \\
[1ex]
\hline
\hline
\end{longtable}
加拉太書 4:4-7
\newline
\begin{longtable}{cl}
\hline
\hline
章節 & 經文 (和合本修訂版)\\
\hline
4:4 & \begin{tabularx}{0.7\textwidth}{X} 等到時候成熟，神就差遣他的兒子，為女子所生，且生在律法之下， \end{tabularx} \\ \\ \relax
4:5 & \begin{tabularx}{0.7\textwidth}{X} 為要把律法之下的人贖出來，使我們獲得兒子的名分。 \end{tabularx} \\ \\ \relax
4:6 & \begin{tabularx}{0.7\textwidth}{X} 因為你們是兒子，神就差他兒子的靈進入我們的心，呼叫：「阿爸，父！」 \end{tabularx} \\ \\ \relax
4:7 & \begin{tabularx}{0.7\textwidth}{X} 可見，你不再是奴隸，而是兒子了，既然是兒子，就靠著神也成為後嗣了。 \end{tabularx} \\ \\
[1ex]
\hline
\hline
\end{longtable}
講題:不凡的竟然如此平凡 Ordinary of Extraordinary

經文:路加福音2:22-40;加拉太書4:4-7

講員:林豪恩先生

場合:香港中文大學崇基學院禮拜堂主日崇拜

日期:2020年12月27日

引言

香港有一間已經結業的旅行社,它的電視廣告歌街知巷聞,不知道大家仍否記得?歌詞頭幾句是這樣的:「食得招積,住得舒適,遊得出色,旅遊超值」.這幾句歌詞以輕鬆有趣的方式描繪出不平凡的旅行團.事實上,人們只會渴慕及追求不凡,因為平凡是不用追求的,平凡本是人的本相.人人都赤身出於母胎,也赤身歸回,在創造的高度來看,人人都是平凡的.縱使如此,人類從沒有放棄在這平凡的人生中追求不凡,個人如是,社會也如是,不用說商界,就是教育界,甚至宗教界都高舉不凡,鞭策著人們脫離平凡,向著不凡的境界奔馳.平凡,好像變成了貶義詞,自甘平凡,被視為自甘墮落.人們怎樣知道自己已經進入不凡的境界呢?最普遍及最常用的方法是與大眾比較:在吃方面,吃大眾吃的食物是平凡,吃大眾吃不到的食物是不凡；在讀書方面,在香港的處境之下,入讀本地公立學校是平凡,能夠脫離本地教育體系的才是不凡,連制定教育政策的官員也是如此為他們的孩子安排的；在機場出入境時,按照既定程序和路線的是平凡,能夠用特別通道及特別程序的就顯出不凡.綜合來說,接受一般待遇的是平凡；能夠享用特別待遇的是不凡.試想想一年前,當大眾都在街上排隊買口罩的時候,有人駕著汽車沿著排隊的人龍一直向前,停在店舖門口,售貨員把口罩拿出來交給他,然後他的汽車絕塵而去.在那一刻間,充分顯出誰是平凡,誰是不凡.在這種觀念之下,人們渴望自己與大眾不同,以此來確定和展示自己的不凡.

今天是聖誕後第一主日,是聖子耶穌露面的時候.耶穌是平凡還是不凡呢?你一定會說,耶穌當然不是凡人.相信大家也認同,耶穌不是凡人,有一首詩歌就叫做「耶穌豈是凡人」.今天,讓我們閱讀福音經課對耶穌出場的描述,加上書信經課的解釋,看看耶穌到底是不凡還是平凡?或者是如何不凡,又如何平凡?希望藉著思考耶穌的平凡與不凡,更新我們對平凡和不凡的觀念.

釋經及應用

1．「及至時候滿足,上帝就差遣他的兒子,為女子所生,且生在律法以下」(加拉太書4:4),

2:22按摩西律法滿了潔淨的日子,他們帶著孩子上耶路撒冷去,要把他獻與主,

2:23(正如主的律法上所記:「凡頭生的男子必稱聖歸主.」)

2:24又要照主的律法上所說,或用一對斑鳩,或用兩隻雛鴿獻祭.

2:39約瑟和馬利亞照主的律法辦完了一切的事,就回加利利,到自己的城拿撒勒去了.

2:40孩子漸漸長大,強健起來,充滿智慧,又有上帝的恩在他身上.

從加拉太書的介紹可見,一方面,耶穌是不凡的,因為他是上帝的兒子；另一方面,耶穌卻是平凡的,因為他和普天下每一個人一樣,為女子所生.上帝的兒子是他的身分,然而身分是看不見的；看得見的是甚麼呢?是為女子所生,即是經歷在媽媽肚子裡面懷胎十月的過程,和經歷出生及出生後的長大過程.耶穌擁有不凡的身分,與眾不同的來頭,卻沒有免去與眾生一樣的生命歷程,眾生所經歷的懷孕,出生,長大這些過程及其中的風險,耶穌都沒有幸免.我們可以假設,如果在耶穌成長的年代有甚麼傳染病在爆發的話,他和其他嬰孩和兒童一樣都處於同樣的風險之中,他的父母和其他父母都一樣地擔心和操心.

加拉太書繼續指出,耶穌生在律法之下.讓我們思想另一個層面,律法是誰發出的呢?是給誰遵守的呢?根據舊約記載,律法是上帝發出的,是給摩西帶領的以色列人遵守的.耶穌是誰呢?是上帝的兒子,與上帝同等,他要不要遵守律法呢?今天我們誦讀的路加福音經文正正是強調耶穌如何緊慎地遵守律法:第一,按摩西律法(2:22按摩西律法滿了潔淨的日子,他們帶著孩子上耶路撒冷去,要把他獻與主)；第二,正如主的律法上所記(2:23(正如主的律法上所記:「凡頭生的男子必稱聖歸主.」)；第三,又要照主的律法上所說(2:24又要照主的律法上所說,或用一對斑鳩,或用兩隻雛鴿獻祭).這幾句敘述,作者每一句都強調耶穌如其他生活在律法之下的人一樣遵守律法.耶穌作為上帝兒子的不凡身分,律法本來不是給他遵守的,他卻如平凡的世人一樣去遵守；反過來,平凡的世人是律法頒布的對象,卻常常渴望自己與眾不同,不受律法的約束,以為自己超越律法.

「2010年10月16晚21時40分許,在河北大學新區超市前,一牌照為“冀FWE420”的黑色轎車,將兩名女生撞出數米遠.被撞一陳姓女生于17日傍晚經搶救無效死亡,另一女生重傷.肇事者口出狂言:“有本事你們告去,我爸是李剛.”此後,這句話成為網友們嘲諷跋扈“官二代”的流行語.」(百度百科)

耶穌是不凡的,但卻接受與平凡的大眾一樣在律法之下生活,沒有運用特權.耶穌自己知道他是不凡的嗎?還是他根本不知道自己的來頭,因此接受平凡眾生的遭遇呢?我們可以在他接受洗禮的時候,看見他是知道自己不凡的身分的,只是他刻意與眾生看齊.

馬太福音第3章

13當下,耶穌從加利利來到約旦河,見了約翰,要受他的洗.14約翰想要攔住他,說:「我當受你的洗,你反倒上我這裡來嗎?」15耶穌回答說:「你暫且許我,因為我們理當這樣盡諸般的義.」

不凡的耶穌,視自己與平凡的人一樣,甘願接受律法的要求；我們作為平凡人的試探,卻是以為自己不凡,以為自己不受約束,不受規管,這往往是危機的開始.越成功的人,越受歡迎的人,越有權力的人,就越容易墮入這樣的試探之中.在政界,商界,甚至教育界和宗教界,都有類似的例子,以為自己超凡脫俗不需要接受規管,出來後得聲名狼藉,甚至身陷囹圄的困局.

「要把律法以下的人贖出來,叫我們得著兒子的名分.」加拉太書4:5

從外貌來看,耶穌是平凡的；從際遇來看,耶穌也是平凡的.那麼,耶穌到底是如何不凡呢?或者你會說,他是上帝的兒子,從上帝而來的,當然不凡.然而,正如剛才說過,耶穌這個身分並不能從他的外表和人生遭遇被看見.正如王子與乞丐對換了衣服和遭遇之後,人們都以為王子是乞丐,乞丐是王子,因為世界是以外表和際遇來判定不凡與平凡的.

聖經是如何判定平凡與不凡的呢?

2:25在耶路撒冷有一個人,名叫西面；這人又公義又虔誠,素常盼望以色列的安慰者來到,又有聖靈在他身上.

2:26他得了聖靈的啟示,知道自己未死以前,必看見主所立的基督.

2:27他受了聖靈的感動,進入聖殿,正遇見耶穌的父母抱著孩子進來,要照律法的規矩辦理.

2:28西面就用手接過他來,稱頌上帝說:

2:29主啊!如今可以照你的話,釋放僕人安然去世；

2:30因為我的眼睛已經看見你的救恩－

2:31就是你在萬民面前所預備的:

2:32是照亮外邦人的光,又是你民以色列的榮耀.

聖經不用外表和際遇去判定平凡與不凡,而是看這個人為世界帶來甚麼,是好處還是壞處?是建立人還是摧毀人?幫助人還是剝削人?西面知道他手上抱著這個平凡的嬰孩耶穌是不凡的,因為他為萬民帶來救恩,為外邦人帶來光,又為以色列民帶來榮耀.

3．加拉太書:「你不是奴僕,乃是兒子了」

剛才我們提過,人們不假思索地以耶穌作為上帝兒子的身分來判定耶穌的不凡,我們沒有否定耶穌是上帝的兒子這身分,只是世人習慣以外表和際遇來顯示身分,耶穌卻不是這樣.他在外表和遭遇上與世人無異,因此他作為上帝的兒子的身分並不是從外表和際遇顯示,而是從他為這個世界帶來甚麼好處來顯示.是他活出他的身分而為這個世界帶來的幸福顯示他的不平凡,這方面給世人很好的反省和學習.關於上帝的兒子這身分,加拉太書是這樣說的:

4:4及至時候滿足,上帝就差遣他的兒子,為女子所生,且生在律法以下,

4:5要把律法以下的人贖出來,叫我們得著兒子的名分.

4:6你們既為兒子,上帝就差他兒子的靈進入你們的心,呼叫:「阿爸!父!」4:7可見,從此以後,你不是奴僕,乃是兒子了；既是兒子,就靠著上帝為後嗣.

原來,不獨是耶穌擁有上帝的兒子這身分,世人也可以是上帝的兒女.得到上帝的兒女這身分的平凡人,有時仍然脫離不到世人對平凡和不凡的判定,就是外表和際遇.因此,以為自己既是上帝的兒女,就可以免受眾生的際遇,變成超人.例如患上惡病時會問:為甚麼上帝的兒女也患上如些惡劣的疾病?遇上意外時會問:為甚麼上帝不看顧他的兒女,讓他們遇上這不幸?過去一年,新型冠狀病毒肆虐,上帝的兒女會否以為自己自動免疫呢?有些人竟然真的如此相信.這種以為作為上帝的兒女這不平凡的身分就享有與眾生不同的生理結構信念,不但沒有為世界帶來好處,甚麼帶來傷害.

2月底韓國大邱新天地教會信徒無視防疫措施,進行大型宗教集會活動,導致防疫破口大開釀成大禍,其中盛傳出「趴趴走」超級傳播者(Super-spreader)信徒大嬸,隱藏自己信徒身份躲避檢疫,以及教會鼓勵信徒參與教會活動時,禁止戴口罩與眼鏡等等,都為其他民眾帶來傷害.

半年後,又有「愛第一教會」事件.其教會牧師確診新冠病毒,然而他一直相信憑藉著上帝恩典,可以克服病情,沒有主動求醫.後來被帶走時,心不甘情不願「半」戴起口罩,令緊慎防疫的民眾感到反感.第一愛教會成為了一個龐大的感染群組,並且漫延至其他群組.教會因為拒絕接受和妨礙流行病學調查,提交虛假資訊,違反《傳染病預防法》,導致疫情擴散而被提出索償46.2億韓元,即大約3千萬港元的訴訟.(The News Lens)

著名新約聖經學者賴特被問及對新冠病毒肆虐的回應時,他說:「最近重讀馬丁路德書信時,其中的一封讓我特別著迷.路德在生時代每隔幾年,就不得不應對這種事情,或是為自己,或是為鄰近鄉鎮哭號的人們,“幫幫我們吧! 我們這裡瘟疫鬧得厲害,人們正在死去,我們該怎麼辦?” 路德談到遵守有關服藥的規定,盡己所能幫助別人,如果自己可能有傳染性,就不要妨礙防疫,免得把疾病傳給別人.他很務實,實際上就是在說,我們就應這樣對付瘟疫. 我們無需在神學上對此做過深的討論.」

那麼,基督徒應該如何回應呢?他回答說:「我一開始就講到,路德已經指出,我們不能傳播感染,這是不負責任的,是對別人的生命不當一回事. 如果我們愛教會的建築物甚於愛我們的鄰里,那麼我們有禍了.」「在古代世界,基督徒走進疫區照顧患者.現在,從某種意義上說,某些基督教理想已經擴散到世界中,世界有系統地去對抗疫情及治療患者,我們應該為此感謝神.但在教會裡,從第一天起,我們就一直在做醫療,濟貧和教育方面的事情.它們深植於教會的DNA中. 因此,基督徒應該恢復這一傳統,並堅持它,而不僅僅是當瘟疫橫行的時候才做.」

上帝的兒女的是平凡的,因為他們跟隨長子耶穌的榜樣,與世人一樣接受平凡的外表和平凡的際遇；上帝的兒女是不凡的,也是因為他們跟隨長子耶穌一樣為世界帶來幫助,光明,安慰和救贖.世人不會因為上帝的兒女的外表而感動,也不會因為上帝的兒女的刀槍不入而感動；若果世人感動,是因為上帝的兒女一方面與他們同行相同的人生旅程,同時在他們有需要的時候伸出援手,彼此扶持.

三．總結

【耶穌豈是凡人】

何因被祂所惑困?感微妙我心靈更新；

耶穌祂只是凡人,因何令我心靈堅穩?

何因被祂所惑困?心靈內似春雷震撼.

耶穌祂只是凡人,竟然令世界變更.

聽,祂每一句話,能帶有大能力.

我默然地想,默想那每一句話.

何因祂敢於面對?祂行事卻出塵脫俗；

耶穌祂只是凡人,竟然為我舒懷解困!

耶穌祂顯出大愛,點燃著我心頭的光.

耶穌祂豈是凡人,欣然地復活再生.

看,祂每一舉動,除去了苦與憂.

我默然地想,十架那每一顆淚.

看,祂的眼神,神態更具寧靜.

我默然地呼,耶穌我更加想念你.



\newpage

\section{道成了肉身}
\label{sec:533}
\textbf{高國雄牧師 - 道成了肉身}
\newline
\newline
\hyperref[sec:532]{< < < PREV SERMON < < <}
~
\hyperref[sec:toc]{[返目錄]}
~
\hyperref[sec:534]{> > > NEXT SERMON > > >}


約翰福音 1:1-18
\newline
\begin{longtable}{cl}
\hline
\hline
章節 & 經文 (和合本修訂版)\\
\hline
1:1 & \begin{tabularx}{0.7\textwidth}{X} 太初有道，道與神同在，道就是神。 \end{tabularx} \\ \\ \relax
1:2 & \begin{tabularx}{0.7\textwidth}{X} 這道太初與神同在。 \end{tabularx} \\ \\ \relax
1:3 & \begin{tabularx}{0.7\textwidth}{X} 萬物都是藉著他造的，沒有一樣不是藉著他造的。凡被造的， \end{tabularx} \\ \\ \relax
1:4 & \begin{tabularx}{0.7\textwidth}{X} 在他裡面有生命，這生命就是人的光。 \end{tabularx} \\ \\ \relax
1:5 & \begin{tabularx}{0.7\textwidth}{X} 光照在黑暗裡，黑暗卻沒有勝過光。 \end{tabularx} \\ \\ \relax
1:6 & \begin{tabularx}{0.7\textwidth}{X} 有一個人，是從神那裡差來的，名叫約翰。 \end{tabularx} \\ \\ \relax
1:7 & \begin{tabularx}{0.7\textwidth}{X} 這人來是為了作見證，是為那光作見證，要使眾人藉著他而信。 \end{tabularx} \\ \\ \relax
1:8 & \begin{tabularx}{0.7\textwidth}{X} 他不是那光，而是要為那光作見證。 \end{tabularx} \\ \\ \relax
1:9 & \begin{tabularx}{0.7\textwidth}{X} 那光是真光，來到世上，照亮所有的人。 \end{tabularx} \\ \\ \relax
1:10 & \begin{tabularx}{0.7\textwidth}{X} 他在世界，世界是藉著他造的，世界卻不認識他。 \end{tabularx} \\ \\ \relax
1:11 & \begin{tabularx}{0.7\textwidth}{X} 他來到自己的地方，自己的人並不接納他。 \end{tabularx} \\ \\ \relax
1:12 & \begin{tabularx}{0.7\textwidth}{X} 凡接納他的，就是信他名的人，他就賜他們權柄作神的兒女。 \end{tabularx} \\ \\ \relax
1:13 & \begin{tabularx}{0.7\textwidth}{X} 這些人不是從血生的，不是從情慾生的，也不是從人的意願生的，而是從神生的。 \end{tabularx} \\ \\ \relax
1:14 & \begin{tabularx}{0.7\textwidth}{X} 道成了肉身，住在我們中間，充充滿滿地有恩典有真理，我們也見過他的榮光，正是父獨一兒子的榮光。 \end{tabularx} \\ \\ \relax
1:15 & \begin{tabularx}{0.7\textwidth}{X} 約翰為他作見證，喊著說：「這就是我曾說：『那在我以後來的先於我，因為在我以前，他已經存在。』」 \end{tabularx} \\ \\ \relax
1:16 & \begin{tabularx}{0.7\textwidth}{X} 從他的豐富裡，我們都領受了恩典，而且恩上加恩。 \end{tabularx} \\ \\ \relax
1:17 & \begin{tabularx}{0.7\textwidth}{X} 律法是藉著摩西頒佈的；恩典和真理卻是由耶穌基督來的。 \end{tabularx} \\ \\ \relax
1:18 & \begin{tabularx}{0.7\textwidth}{X} 從來沒有人見過神，只有在父懷裡獨一的兒子將他表明出來。 \end{tabularx} \\ \\
[1ex]
\hline
\hline
\end{longtable}
講題:道成了肉身 The Word Became Flesh

經文:約翰福音1:1-18

講員:高國雄牧師

場合:香港中文大學崇基學院禮拜堂主日崇拜

日期:2021年1月3日

耶穌基督福音的開始,福音書中有不同描述的重點,讓我們先重溫福音書記載耶穌基督傳道服務之前的內容.

馬太福音:(1:1-17)耶穌基督的家譜開始,形容祂是亞伯拉罕的後裔,大衛的子孫…； (1:18-25)主的使者在約瑟夢中出現,耶穌/以馬內利名字的由來；(2:1-23)博士來朝,獻三禮物與孩童耶穌,希律追殺2歲下男孩,約瑟一家逃到埃及(剛過去平安夜崇拜的信息)；(3:1-17)施洗者約翰出場,為耶穌洗禮；(4:1-11)耶穌受魔鬼的試探.

路加福音:(1:5-25)預言施洗者約翰的出生；(1:26-38)預言耶穌的出生；(1:39-56)馬利亞與伊利莎白相遇並稱頌主；(1:57-80) 施洗者約翰的出生並撒迦利亞對主的稱頌；(2:1-7)耶穌在伯利恆誕生；(2:8-20)天使向牧羊人報喜訊,牧羊人在馬槽遇見嬰孩耶穌；(2:21-38)耶穌受割禮及被獻呈給主；(2:39-52)少年的耶穌；(3:1-22) 施洗者約翰在約旦河為人洗禮,耶穌洗禮；(3:23-38) 耶穌基督的家譜,由約瑟倒數到亞當；(4:1-13)耶穌受魔鬼的試探.

馬可福音:(1:1-11) 施洗者約翰出場,為耶穌洗禮；(1:12-12) 耶穌受魔鬼的試探.

約翰福音在 (1:19-34)記述施洗者約翰出場,為耶穌洗禮,請留意四卷福音書相同的記述就只有約翰及洗禮的事；在這以先,真正描述由「道」開始…

道是什麼呢?

約翰福音1:1中「太初有道」中的「道」Logos λόγος,由λέγω (I say)而來,英語世界翻譯為Word,說話,也就是上帝的話語；傳教士來到中國,尋找出比Word蘊含更豐富意義,更貼近希臘文,甚至是基督教對上帝的理解,這就是「道」；中國人推崇的是天道,老子在道德經所提到的是「道可道,非常道」；道不單只是話語的意思,更是智慧,論証,甚至可以是宇宙的法則,天地之始…

道=聖經=上帝話語,帶著能力,從創世開始,上帝的啟示,律例,教訓,指引,展現,打開.

詩篇119:105你的話是我腳前的燈,是我路上的光.

約翰福音所表達的「道」又是怎樣一回事呢?

約翰福音1:1-18

1太初有道,道與上帝同在,道就是上帝.2這道太初與上帝同在.3萬物都是藉著他造的,沒有一樣不是藉著他造的.凡被造的,4在他裡面有生命,這生命就是人的光.5光照在黑暗裡,黑暗卻沒有勝過光.6有一個人,是從上帝那裡差來的,名叫約翰.7這人來是為了作見證,是為那光作見證,要使眾人藉著他而信.8他不是那光,而是要為那光作見證.9那光是真光,來到世上,照亮所有的人10他在世界,世界是藉著他造的,世界卻不認識他.11他來到自己的地方,自己的人並不接納他.12凡接納他的,就是信他名的人,他就賜他們權柄作上帝的兒女.13這些人不是從血生的,不是從情慾生的,也不是從人的意願生的,而是從上帝生的.14道成了肉身,住在我們中間,充充滿滿地有恩典有真理,我們也見過他的榮光,正是父獨一兒子的榮光.15約翰為他作見證,喊著說:「這就是我曾說:『那在我以後來的先於我,因為在我以前,他已經存在.』」16從他的豐富裡,我們都領受了恩典,而且恩上加恩.17律法是藉著摩西頒佈的；恩典和真理卻是由耶穌基督來的.18從來沒有人見過上帝,只有在父懷裡獨一的兒子將他表明出來.

耶穌基督=道

道=太初與上帝同在

道=上帝

道=創造萬物

道=創造生命之物

道=真光,照亮黑暗

道:世界不認識,不接納

接納主就可以成為上帝的兒女,因為人是從上帝所生的.

「道」成了肉身!

道成肉身就是上帝自限的表現,降卑,全因為愛.

耶穌基督是完全的人,也是完全的上帝.

上帝成為了人,不單是上帝願意明白人,而是「成為」了人.

道=耶穌基督是上帝獨一兒子,住在我們中間.

道=承傳舊約律法,表明上帝

中國人傳統,老子道德經:「道可道,非常道.」=能夠言說解釋的道,就不是真正永恆的道；宇宙本源的道,可感應可領悟,但不可解釋,因為其浩瀚,偉大…

道成肉身是有目的的:救贖,愛世人

約翰福音3:16 上帝愛世人,甚至將祂的獨生子賜給他們,叫一切信祂的,不致滅亡,反得永生!

道成肉身就是主親自道出這道來!

所以基督信仰的道不單可道,甚至必須人人去道,就是要傳這道.

反映出主的榮光,讓更多人成為主的兒女…

上帝先主動啟示自己,把自己展開,成為可親迎的上帝；因此,基督教的信仰不是自己尋找,努力的信仰…



\newpage

\section{你們將要看見天開了}
\label{sec:534}
\textbf{鄧瑞強博士 - 你們將要看見天開了}
\newline
\newline
\hyperref[sec:533]{< < < PREV SERMON < < <}
~
\hyperref[sec:toc]{[返目錄]}
~
\hyperref[sec:535]{> > > NEXT SERMON > > >}


約翰福音 1:43-51
\newline
\begin{longtable}{cl}
\hline
\hline
章節 & 經文 (和合本修訂版)\\
\hline
1:43 & \begin{tabularx}{0.7\textwidth}{X} 又過了一天，耶穌想要往加利利去。他找到腓力，就對他說：「來跟從我！」 \end{tabularx} \\ \\ \relax
1:44 & \begin{tabularx}{0.7\textwidth}{X} 這腓力是伯賽大人，是安得烈和彼得的同鄉。 \end{tabularx} \\ \\ \relax
1:45 & \begin{tabularx}{0.7\textwidth}{X} 腓力找到拿但業，對他說：「摩西在律法書上所寫的，和眾先知所記的那一位，我們遇見了，就是約瑟的兒子拿撒勒人耶穌。」 \end{tabularx} \\ \\ \relax
1:46 & \begin{tabularx}{0.7\textwidth}{X} 拿但業對他說：「拿撒勒還能出甚麼好的嗎？」腓力說：「你來看。」 \end{tabularx} \\ \\ \relax
1:47 & \begin{tabularx}{0.7\textwidth}{X} 耶穌看見拿但業向他走來，就論到他說：「看哪，這真是個以色列人！他心裡是沒有詭詐的。」 \end{tabularx} \\ \\ \relax
1:48 & \begin{tabularx}{0.7\textwidth}{X} 拿但業對耶穌說：「你從哪裡認識我的？」耶穌回答他說：「腓力還沒有呼喚你，你在無花果樹底下，我就看見你了。」 \end{tabularx} \\ \\ \relax
1:49 & \begin{tabularx}{0.7\textwidth}{X} 拿但業回答他：「拉比！你是神的兒子，你是以色列的王。」 \end{tabularx} \\ \\ \relax
1:50 & \begin{tabularx}{0.7\textwidth}{X} 耶穌回答他說：「因為我說在無花果樹底下看見你，你就信嗎？你將看見比這些更大的事呢！」 \end{tabularx} \\ \\ \relax
1:51 & \begin{tabularx}{0.7\textwidth}{X} 他又說：「我實實在在地告訴你們，你們將要看見天開了，神的使者在人子身上，上去下來。」 \end{tabularx} \\ \\
[1ex]
\hline
\hline
\end{longtable}
講題:你們將要看見天開了 You will all See Heaven Open

講員:鄧瑞強博士

經文:約翰福音1:43-51

場合:香港中文大學崇基學院禮拜堂主日崇拜

日期:2021年1月17日

各位弟兄姊妹,早安.在這疫症裡,祝願各人平安.

人類是一種不停渴望的動物.我們渴望的範圍,從生理的滿足,到心理或情感的滿足,再到精神的滿足.嬰孩會哭,渴望母親的餵哺.小孩子做很多古靈精怪的行為,渴望大人的注意.年輕人努力做出各種成績來,渴望得到別人的認同.心智成熟的,會進而渴望一個理想的世界.人真的很特別,我們雖然活在當下,活在此時此刻,卻渴望一個從未出現過的世界.希臘的柏拉圖,活在當下,卻渴望那個從未見過的,超越時間的理型世界.中國古代的《禮記》,渴望一個「人不獨親其親,不獨子其子,使老有所終,壯有所用,幼有所長,鰥寡孤獨廢疾者皆有所養」的大同世界.耶穌時代的猶太人,他們渴望,他們等待,一個能解救他們脫離人間苦難的救主；他們渴望,他們等待,這救主能開展一意想不到的天國.人類不停地在渴望.

今日,我們讀一個渴望得到初步實現的故事:約翰福音1:43-51.

約1:43　又次日,耶穌想要往加利利去,遇見腓力,就對他說:「來跟從我吧.」

約1:44　這腓力是伯賽大人,和安得烈,彼得同城.

約1:45　腓力找著拿但業,對他說:「摩西在律法上所寫的和眾先知所記的那一位,我們遇見了,就是約瑟的兒子拿撒勒人耶穌.」

約1:46　拿但業對他說:「拿撒勒還能出甚麼好的嗎?」腓力說:「你來看!」

約1:47　耶穌看見拿但業來,就指著他說:「看哪,這是個真以色列人,他心裡是沒有詭詐的.」

約1:48　拿但業對耶穌說:「你從哪裡知道我呢?」耶穌回答說:「腓力還沒有招呼你,你在無花果樹底下,我就看見你了.」

約1:49　拿但業說:「拉比,你是神的兒子,你是以色列的王!」

約1:50　耶穌對他說:「因為我說『在無花果樹底下看見你』,你就信嗎?你將要看見比這更大的事」；

約1:51　又說:「我實實在在地告訴你們,你們將要看見天開了,神的使者上去下來在人子身上.」

約1:43　又次日,耶穌想要往加利利去,遇見腓力,就對他說:「來跟從我吧.」

「又次日」,這是約翰福音開場白中「七日結構」的第四日,這「七日結構」是刻意與創世記的「七日創造的結構」相對應的.在創世記,在第四日,神創造天上的光體,從此,無論日夜,都有光體將光明帶來人間.在約翰福音,第四日,耶穌呼召人跟從他.這些人,將要成為反照神榮光的光體,將神的榮光帶來人間.

耶穌準備往加利利去,忽然遇見腓力,耶穌直接就對他說:「來跟從我吧.」在這裡,約翰福音沒有交代任何背景,沒有交代耶穌是否事前認識腓力,也沒有交代腓力的反應.總之,耶穌行過,見到人,便叫人跟從他.

你可以試試,在旺角街頭,見到人,隨便叫他跟從你,看看有何結果.

看來,約翰福音在這裡,不是記述一件生活中的小事件,而是描述一件神學上的大事件.耶穌臨在世界,就是一把大聲音,一個大呼喚,一個震動每個人的靈魂的大呼喚.這個大呼喚,呼喚每個人去跟從他.

大自然充滿呼喚.大學時代,我上過黃山.那日,日落時分,下著雨.忽然,雨停了,頃刻間,雲降在我腳下,變成一望無際的雲海.我站在雲上,夕陽在我面前,徐徐落在雲海裡,真是「何似在人間」.面對這震撼的美,我聽到它對人類的呼喚.它呼喚我們,謙虛下來,放下爭競,心懷敬畏地生活.壯麗的大地會呼喚,高超的人格也會呼喚.巍峨的高山,呼喚我們承認生命的渺小,卻不能呼喚我們跟從他.神聖的耶穌,高貴的人格,卻能呼喚我們跟從他.

耶穌來到塵世,呼喚我們放下手頭的工作,放下日日夜夜的營營役役,耶穌呼喚我們跟從他.

這呼喚真古怪,卻又令人不得不停下來沉思,沉思我們一直以來的古老渴望.那心底的渴望,從沒有消失過,但在消磨人的生活中,我們已漸漸忘記,忘記我們還有夢,忘記我們還可以有夢,忘記我們仍可期待那夢.

約1:45　腓力找著拿但業,對他說:「摩西在律法上所寫的和眾先知所記的那一位,我們遇見了,就是約瑟的兒子拿撒勒人耶穌.」

聽到耶穌那呼喚,腓力隨後的行動,是去找拿但業.拿但業是誰,經文沒有交代.腓力說:「摩西在律法上所寫的和眾先知所記的那一位,我們遇見了,就是約瑟的兒子拿撒勒人耶穌.」這說話講的是什麼?它論及一個古代留下來的夢想,這是古代聖人留在每個猶太人心靈中的渴望.幾百年來,猶太人沒有了自己的國家,失去了自由,在屈辱下生活.他們發了幾百年的夢.他們渴望,有一天,有一個人會來,扶助他們,建立一個由神掌權的國度,讓他們自由,讓他們活得有尊嚴,讓他們在神的仁愛下生活.這個夢,他們發了幾百年了.這個夢,在他們艱難的歲月中,也漸漸褪色了.依稀的,有個夢,一個古老的夢,但是,還應該繼續發這春秋大夢嗎?或許,有多少猶太人,有多少神的子民,已放棄這個夢了.忽然,聽到耶穌的呼喚,腓力回憶起這個夢.啊,他記起了,摩西在律法上寫過的,眾先知曾預言過的,那一位,那一位.他聽到耶穌的呼喚,他心靈一震,靈機一觸,他知道,是他了,是他了,是這個約瑟的兒子拿撒勒人耶穌.這個渴望,這個夢,終於實現了.

拿但業的反應,頗耐人尋味.

約1:46　拿但業對他說:「拿撒勒還能出甚麼好的嗎?」腓力說:「你來看!」

拿但業斷然拒絕相信這古老的夢會實現.猶如一個被感情傷透的人,斷然拒絕愛情的可能.拿撒勒,一塊窮鄉僻壤,與古老的夢想,怎可能有關連呢?

香港呢?香港的夢呢?會實現嗎?

或許,拿但業同樣會說:不可能.

腓力對拿但業說:「你來看!」「你來睇睇先講!」

然後,情節峰迴路轉.

約1:47　耶穌看見拿但業來,就指著他說:「看哪,這是個真以色列人,他心裡是沒有詭詐的.」

約1:48　拿但業對耶穌說:「你從哪裡知道我呢?」耶穌回答說:「腓力還沒有招呼你,你在無花果樹底下,我就看見你了.」

約1:49　拿但業說:「拉比,你是神的兒子,你是以色列的王!」

不是拿但業先看到耶穌,而是耶穌先看到他.不是我們先看見神,而是神先看見我們.不是我們先尋找神,而是神先尋找我們.耶穌看見拿但業,神看見他.

耶穌指著拿但業說:「看哪,這是個真以色列人,他心裡是沒有詭詐的.」猶太人的先祖雅各,他後來改名叫「以色列」,他的生命特徵,就是「詭詐」.他的詭詐,令他走了很多冤枉路,嚐盡風霜.如今,耶穌說,眼前的這個人,心裡沒有詭詐,是個「真以色列人」.耶穌讚他,耶穌欣賞他,耶穌看見他.耶穌看見一個人的靈魂的真正質素,耶穌從人的真我層面認識一個人.

拿但業對耶穌說:「你從哪裡知道我呢?」我們只是萍水相逢,素未謀面,為何你好像能看穿我的靈魂一樣似的?

耶穌回答說:「腓力還沒有招呼你,你在無花果樹底下,我就看見你了.」

「在無花果樹底下」,到底是什麼的一回事?

不同文化,賦予不同植物以不同含意.在古代中國,楊柳就與離別有關.最經典的詩詞是柳永的《雨霖鈴‧寒蟬淒切》:「多情自古傷離別,更那堪,冷落清秋節.今宵酒醒何處?楊柳岸,曉風殘月.」又或是蘇軾的《水龍吟》:「細看來,不是楊花,點點是離人淚.」

在猶太文化,無花果樹又代表什麼?

在聖經,無花果樹第一次出現,是創世記的第3章.那時,亞當,夏娃食了禁果,發現自己是赤身露體的,便用無花果樹的葉子,遮蓋自己的身體及羞恥.當人類不再配住在伊甸園時,無花果樹象徵著一種修補的力量.後來,無花果樹與「幸福的生活」和「期盼天國的來臨」有關.列王記上4:25:「所羅門在世的日子,從但到別是巴的猶大人和以色列人都在自己的葡萄樹下和無花果樹下安然居住.」這是猶太人的黃金歲月,最好的生活就是在葡萄樹下和無花果樹下安然居住.後來,先知彌迦預告猶大的滅國,但先知同時預告,苦難過後,天國會來臨.彌迦書4:3-4論到天國的光景,說:「……他們要將刀打成犁頭,把槍打成鐮刀.這國不舉刀攻擊那國,他們也不再學習戰事.人人都要坐在自己葡萄樹下和無花果樹下,無人驚嚇.這是萬軍之耶和華親口說的.」將要來臨的天國,其幸福生活,就是坐在自己葡萄樹下和無花果樹下,無人驚嚇.無花果樹與葡萄樹,代表著幸福,美滿,天國的實現,重返伊甸園.現在,在拿但業的故事裡,提及無花果樹.遲一點,耶穌在救贖工程完成的前夕,會談論葡萄樹,即約翰福音裡出名的「葡萄樹與枝子」的經文.

若果無花果樹代表著對天國的期盼,對重返伊甸園的渴望,則我們可以想像,拿但業在無花果樹底下,他在那裡,懷緬古時所羅門王時代最安穩的日子,同時,他思想彌迦先知的說話,他渴望先知所講的,一個神設立的國度,那時,各人將刀打成犁頭,把槍打成鐮刀；這國不舉刀攻擊那國,各人也不再學習戰事.多麼好的世界,但是,畢竟,這只是夢,一個古老的夢,一個看來不會實現的夢.

這時候,耶穌對他說:「腓力還沒有招呼你,你在無花果樹底下,我就看見你了.」耶穌說,我看見你,我看見你是一個有夢想的人,我看見你是一個對神的應許有信心的人,我看見你是一個多麼渴想伊甸園的人,當世人已失去對伊甸園的渴望,當世人已失去對神的信任,我看見你,我看見你,你依然站在信仰的高崗上,以不朽的盼望,去抵抗人間的絕望.我看見你,你是真以色列人,你是真正的神的子民,那個在無夢的世界,堅持夢想的人.我看見你.

拿但業心中的渴望,被耶穌一語道破.他心中的信仰之火,像被靈風吹拂一樣,燃聖靈魂的烈焰.他衝口而出,對耶穌說:「拉比,你是神的兒子,你是以色列的王!」耶穌的三個身份,一氣呵成地說出來.在世人眼前,耶穌是拉比,是真理的教師.就生命深處的存在本質而言,耶穌代表神,啟示神,將神的旨意活出來,是神的兒子.就歷史的命途而言,耶穌來建立神的國度,是以色列的王.

這就是拿但業在耶穌身上所見到的一切,但這就已是事實的全部了嗎?耶穌繼續說的事,超乎他的想像,超乎人類的想像.

約1:50　耶穌對他說:「因為我說『在無花果樹底下看見你』,你就信嗎?你將要看見比這更大的事」；

約1:51　又說:「我實實在在地告訴你們,你們將要看見天開了,神的使者上去下來在人子身上.」

這段經文,明顯以舊約創世記28章那個「雅各的天梯」的故事為背景.那個後來改名叫「以色列」的雅各,在伯特利這個地方,夢見一連接天地的梯子,天使在梯子上,上上落落.神就在梯子的頂端,俯視人間,祝福人間.雅各驚醒,體會到,他身處的地方,正是神的殿,天之門.

耶穌對拿但業說,我就是神的殿.從今以後,再不用什麼通天的梯子了,我在的地方,就是聖殿.在我裡面活著的,就是在神的臨在之下生活.真正的聖殿,不在這裡,不在那裡,而是在我身上.聖殿不是一個地方,而是耶穌的生命.在耶穌身上,人不單能重返伊甸,不單能看到天國的實現,更能體會天地的溝通,天人之合一.

人類是一種不停渴望的動物,最終極的渴望,是神聖的渴望,我們渴望看見天開了,我們渴望神的臨在,我們渴望天人合一.在耶穌身上,這一切都實現了.在這絕望人間,深願你能在耶穌身上,看到這永恆渴望的實現.

願榮耀歸上主,願人間享和平.



\newpage

\section{悔改相信}
\label{sec:535}
\textbf{張亦農牧師 - 悔改,相信!}
\newline
\newline
\hyperref[sec:534]{< < < PREV SERMON < < <}
~
\hyperref[sec:toc]{[返目錄]}
~
\hyperref[sec:536]{> > > NEXT SERMON > > >}


馬可福音 1:14-20
\newline
\begin{longtable}{cl}
\hline
\hline
章節 & 經文 (和合本修訂版)\\
\hline
1:14 & \begin{tabularx}{0.7\textwidth}{X} 約翰下監以後，耶穌來到加利利，宣講神的福音， \end{tabularx} \\ \\ \relax
1:15 & \begin{tabularx}{0.7\textwidth}{X} 說：「日期滿了，神的國近了。你們要悔改，信福音！」 \end{tabularx} \\ \\ \relax
1:16 & \begin{tabularx}{0.7\textwidth}{X} 耶穌沿著加利利的海邊走，看見西門和西門的弟弟安得烈在海上撒網；他們本是打魚的。 \end{tabularx} \\ \\ \relax
1:17 & \begin{tabularx}{0.7\textwidth}{X} 耶穌對他們說：「來跟從我，我要叫你們得人如得魚一樣。」 \end{tabularx} \\ \\ \relax
1:18 & \begin{tabularx}{0.7\textwidth}{X} 他們立刻捨了網，跟從他。 \end{tabularx} \\ \\ \relax
1:19 & \begin{tabularx}{0.7\textwidth}{X} 耶穌稍往前走，又見西庇太的兒子雅各和他弟弟約翰在船上補網。 \end{tabularx} \\ \\ \relax
1:20 & \begin{tabularx}{0.7\textwidth}{X} 耶穌隨即呼召他們，他們就把父親西庇太和雇工留在船上，跟從了耶穌。 \end{tabularx} \\ \\
[1ex]
\hline
\hline
\end{longtable}
講題:悔改,相信! Repent and Believe!

經課:馬可福音1:14-20

講員:張亦農牧師

場合:香港中文大學崇基學院禮拜堂主日崇拜

日期:2021年1月24日

I. 相信．悔改悔改．相信

基督徒日常傳福音的經驗:相信,悔改

馬可記述耶穌的話卻是:「悔改,相信」

「日期滿了,上帝的國近了.你們要悔改,信福音!」(可1:5)

耶穌傳道的信息不是獨創,而是建基在施洗約翰的信息之上

施洗約翰的話

「你們要悔改!因為天國近了.」(太3:2)

「宣講悔改的洗禮,使罪得赦.」(可1:4)

面對上帝國近了,指的是上帝要來臨,彰顯祂的統治.面對公義聖潔的主,人能夠甚至應該做的是作出改變.特別是除去一些主不喜悅的行為.

耶穌的信息的不同在於

悔改是不足夠的,並且要相信福音

甚麼是福音呢?

「上帝的兒子,耶穌基督福音」(可1:1)

馬可傳遞上帝的心意:人不單要悔改,更加要相信上帝藉著耶穌基督為我們預備的救恩

彼得:「除他以外,別無拯救,因為在天下人間,沒有賜下別的名,我們可以靠著得救.」(徒4:12)

II. 耶穌為甚麼往加利利

馬可福音開首提到,耶穌在約旦河受施洗約翰的洗,然後在曠野面對撒旦的試探.

約旦河耶穌受洗的地點,相信是在猶大地,近耶路撒冷的地方.

今日經課提到:「約翰下監以後,耶穌來到加利利,宣講上帝的福音」(可1:14)

耶穌為甚麼要離開猶大地,離開耶路撒冷,去到加利利傳道呢?

施洗約翰和耶穌雖然相識,但是他們不是合作同工,形成一個集團.所以當施洗約翰被希律捉拿下監,我們難以想像耶穌到加利利為的是逃避政治逼害.

加利利雖然有不少以色列人聚居,但不是宗教信仰發展蓬勃的地方,也不是猶太教重心所在.猶太教的重心仍然在猶大地的耶路撒冷,和耶路撒冷一帶.就算連施洗約翰傳道也是在耶路撒冷附近的曠野和約旦河,就可想而知.

耶穌為甚麼要到加利利宣教?是一個值得思考的問題.

當耶穌在加利利開始宣教,馬可告訴我們,第一批回應耶穌的是四個漁夫:彼得,安得烈,約翰和雅各.有趣的是,這四個人是第一次認識耶穌嗎?

約翰福音1:35-42告訴我們約翰和安得烈,甚至彼得都可能是施洗約翰的門徒.他們在耶路撒冷附近,因為施洗約翰的介紹,認識了耶穌.所以在加利利,耶穌和他們不是初相識,而是再見面.

為甚麼彼得,安得烈,約翰等會在加利利重操故業,做漁夫捕魚呢?

相信與施洗約翰的際遇有關.

「照這話,施洗約翰來到曠野,宣講悔改的洗禮,使罪得赦.猶太全地和全耶路撒冷的人都出去,到約翰那裡,承認他們的罪,在約旦河裡受他的洗.」(可1:4-5)馬可的這段話,相信反映了他宣教事業的高峰.

好景不常,因為他批評希律王,遭到政治逼害,被捕下監.(可6:17)

所謂「樹倒猢猻散」,作為施洗約翰的門生,大能當前,總有害怕和失望.逃亡,回老家避難,尤其是離開希律的勢力範圍,是最好的選擇.

若果我們從這個角度,思考耶穌為甚麼要到加利利的話,原因就可圈可點了.

耶穌到加利利,為的是要尋找那幾個失去信心的,施洗約翰的門生.

III. 門徒的新召命

耶穌要幫助彼得,安得烈,約翰和雅各重拾對宗教的熱心和追求.

耶穌的話:「你們要悔改,信福音!」,對於四個人來說,就是挑戰他們走出失意灰心,作出改變,並且相信並跟從耶穌,作他的門徒.

他們原本的師父施洗約翰不錯是曾經一時無兩,令到許多人認真思考與天父上帝的關係,願意悔改回轉,藉著領受洗禮,表達悔罪,潔淨自己.但是,總不是所有人都願意面對自己的過犯,認罪悔改.上帝國即使要來臨,也會面對拒絕和抗拒.所以施洗約翰被捕下獄.

然而上帝國的來臨並不會因著人的抗拒而停步,上帝拯救世人的心意仍舊要按著祂的計劃成就.這要成就在降生成人的耶穌基督的身上.

所以當耶穌宣告:「要悔改,信福音!」,對於他的門徒來說,是不要灰心失意,要相信耶穌基督,相信天父藉著耶穌施行拯救的心意.

「來跟從我,我要叫你們得人如得魚一樣.」(可1:17)賜給門徒新的願景.

耶穌更加應許,他們能夠在主耶穌的恩助下,成為得人漁夫,領人歸主.

IV. 悔改．相信的現代意義?

人生不如意事十常八九.

意料之外的事情常常會發生.容易讓人感到失望,灰心,甚至放棄.

但是我們的主是那位願意鼓勵人的神.昔日他願意走到那四位失意灰心的門徒那裡,呼召他們跟從他.

今日我們的主也會願意走進我們失意的人生,鼓勵我們,呼召我們跟從祂.願我們珍惜祂對我們的愛.

主也願意為我們的人生賜下遠象,讓我們知道生命努力的方向.

悔改,相信:重新振作,信靠我們的主.尋求祂的旨意,成為我們生命努力的方向.



\newpage

\section{先知文士耶穌}
\label{sec:536}
\textbf{林豪恩先生 - 先知,文士,耶穌}
\newline
\newline
\hyperref[sec:535]{< < < PREV SERMON < < <}
~
\hyperref[sec:toc]{[返目錄]}
~
\hyperref[sec:537]{> > > NEXT SERMON > > >}


講題:先知,文士,耶穌 Prophet, Scribe, Jesus

講員:林豪恩先生

場合:香港中文大學崇基學院禮拜堂主日崇拜

日期:2021年1月31日

一.引言

不久之前,網路上流傳著關於各行各業,各身分,各族群的「別人覺得我們…」六格搞笑圖文.例如有一幅圖名為「教師是這樣的人」:第一格的標題是「我朋友覺得我在做的事」,圖片是在沙灘上躺著享受日光浴；第二格的標題是「我媽覺得我在做的事」,圖片是坐在路邊行乞的乞丐；第三格的標題是「社會覺得我在做的事」,圖片是參加爭取教育經費的遊行；第四格的標題是「學生覺得我在做的事」,圖片是軍官在喝罵士兵；第五格的標題是「我覺得我在做的事」,圖片是摩西站在紅海旁邊高舉雙手號令海水分開；最後一格的標題是「我實際上正在做的事」,圖片是一位馴獸師正在訓練老虎.

不知道有沒有「基督徒是這樣的人」的類似六格搞笑圖文呢?朋友覺得基督徒是怎樣的呢?媽媽又覺得基督徒是怎樣的呢?社會又覺得基督徒是怎樣的呢?青少年又覺得基督徒是怎樣的呢?基督徒又覺得自己是怎樣的呢?最後一格,基督徒實際是怎樣的呢?

這六格圖文雖然搞笑,不過,學界的確有探討公眾如何看待基督徒的研究?崇基神學院曾經進行一項名為「新聞傳媒(報章)中的基督教形象」(2001-2011)研究,最常用來描述基督徒的詞彙是甚麼呢?是「虔誠」和「歧視」.傳媒眼中的「虔誠」是怎樣的呢?其中一個例子:例如鄧萃雯提及 TVB 其中一位老闆是「很虔誠的基督徒」,因此覺得「有點希望.做生意他們已經很有錢了,如果又將 TVB 當成是一盤純粹的生意就慘了,希望他們真的想搞好它.」(《星島》2011-5-29)從下半句描述可見,鄧氏對「虔誠」信徒的寄望,是相信他們即或面對利益關頭,仍願意抱緊理想,做良心老闆.至於傳媒同「歧視」一詞來描述基督徒,逾六成議題集中於「同性戀」及「性傾向歧視條例立法」之上.

作為基督徒的你,你覺得自己是甚麼呢?或者正在慕道的朋友,你覺得自己成為基督徒之後會變成甚麼?今天的經課提及幾種人物,包括先知,文士和耶穌.基督徒是哪一種人物呢?讓我們一齊透過經文思考這個問題,並尋求對基督徒生活的啟發.

二．申命記:先知

18:15耶和華－你的上帝要從你們弟兄中間給你興起一位先知,像我,你們要聽從他.

18:16正如你在何烈山大會的日子求耶和華－你上帝一切的話,說:『求你不再叫我聽見耶和華－我上帝的聲音,也不再叫我看見這大火,免得我死亡.』

18:17耶和華就對我說:『他們所說的是.

18:18我必在他們弟兄中間給他們興起一位先知,像你.我要將當說的話傳給他；他要將我一切所吩咐的都傳給他們.

在申命記18章18章耶和華上帝清楚定義先知的角色和任務:「我要將當說的話傳給他；他要將我一切所吩咐的都傳給他們」.先知是一位傳話人,他所傳的是上帝的話說,所要傳的對象是上帝指示要聆聽的人,在這場景是出埃及的以色列民,在另一些場景中可以是外邦某國或列國.

先知傳言並不是「講下就算」,喜歡就讚,「不喜勿插」,而是有後果的.

18:19誰不聽他奉我名所說的話,我必討誰的罪.

由此可見,先知角色相當超然.可能因此,歷世歷代都有人自稱先知,或被稱為先知.以下是其中一個例子:

大衛貝治(David Berg)就是一位引用這段聖經,自稱是先知的人.並且在一個稱為「家族」(The Family)《前名:上帝的孩子(The Children of God)》的宗教教派中,也被他的追隨者認為是一位先知,同時稱呼其為Moses David.

在一般教會,很少聽到有人自稱先知,或者被稱為先知.不過,有意見認為教會作為整體,或者基督徒作為個體具有先知的角色.例如有呼籲說:「哪裡有不公義,就要嚴正指明；哪裡有人受逼害,就要挺身而出護衛弱小者；哪裡有可改善之處,就要極力推進.」(基督教週報2031期)如果先知的角色在當下被定義為:「指斥不公義,保護弱小者,推進可改善之處」,相信沒有人會反對.不過,困難卻是如何把這些概念落實到具體事件上,因為具體狀況往往涉及判斷,而判斷則受到各種因素影響.例如,怎樣才算是不公義呢?誰是弱小者呢?因此,在現實上,不容易得到如摩西一樣的肯定,「我要將當說的話傳給他；他要將我一切所吩咐的都傳給他們」的那種來自上主的先知性宣告.使徒保羅有時也說:「我沒有主的命令……就把我的意見告訴你們」(林前7章)雖然教會或個別信徒具有先知的角色,但保羅的這份緊慎和自知是值得參考的,因為申命記這段論及先的的話也提醒:18:20若有先知擅敢託我的名說我所未曾吩咐他說的話,或是奉別神的名說話,那先知就必治死.』

基督徒會否覺得自己是上主的先知呢?就算是,也應該分辯所出之言是從上主而來還是從別處而來.

三．馬可福音:文士

1:21到了迦百農,耶穌就在安息日進了會堂教訓人.

1:22眾人很希奇他的教訓；因為他教訓他們,正像有權柄的人,不像文士.

馬可福音在第一章已經開始記述耶穌的傳道工作,並且在這個開始的時候,就在毫無背景之下把文士拿來和耶穌比較,對於不熟識來龍去脈的讀者,可能會感到文士「躺著也中槍」.馬可福音記述耶穌在安息日進了會堂講道,是假設讀者知道平日在會堂講道的是文士,因此自然地把耶穌這位外來講員和駐堂講員比較.如此比較,就正如現在一樣,通常都會把駐堂講員比下去的.話說回頭,這段經文描述的場景,其實是與申命記那段經文有關的,因為在會堂裡,有一個地方叫做「摩西的座位」,講道者就是在這個座位講解摩西律法.文士就是獻身於抄寫律法,捍衛律法和解釋律法工作的知識份子.如果要比較,現代的講道者可能類似文士多於先知,一方面沒有直接從上帝領受說話,另一方面也沒有被賦予呼風喚雨的能力,能夠做好傳承律法,捍衛律法和解釋律法這些工作已經不負所託了.因此,文士沒有權柄是合理的,有甚麼問題呢?問題就是,文士以為自己有權柄,運用這本來沒有的權柄來管轄別人,這是值得我們留心及反省的.

福音書記述耶穌多次指責文士假冒為善,很可能因為文士經常講解律法教導民眾,就好像我們今天一樣,「講左等於做左」,宣稱和真實的落差太大,人人都看到,但自己卻看不到.另外可能是文士捍衛律法心切,不能接受不守律法的狀況,不能容忍破壞律法的人,於是動用權力,甚至武力,就好像遇見耶穌之前血氣方剛的保羅追捕基督徒一樣.

如果,我們只是忠心做一個文士,傳承律法,不假設自己有權柄,會不會過份保守呢?可能會,但總比偷竊權柄並運用來打壓異見者安全得多.或許有人說,上帝的說話不是帶著權柄的嗎?這是對的,路加福音記述同一件事情時是這樣說的:「4:32他們對他的教導感到很驚奇,因為他的話裡有權柄.」讓我們留心,有權柄的是「他的話」,不是我們這些重述「他的話」的人.

接著,讓我們看看位「正像有權柄的人,不像文士」是怎樣的.

四．馬可福音:耶穌

1:23在會堂裡,有一個人被污鬼附著.他喊叫說:

1:24「拿撒勒人耶穌,我們與你有甚麼相干?你來滅我們嗎?我知道你是誰,乃是上帝的聖者.」

1:25耶穌責備他說:「不要作聲!從這人身上出來吧.」

1:26污鬼叫那人抽了一陣瘋,大聲喊叫,就出來了.

1:27眾人都驚訝,以致彼此對問說:「這是甚麼事?是個新道理啊!他用權柄吩咐污鬼,連污鬼也聽從了他.」

真正有權柄的人是令人驚訝的,因為他說話,就成真.「他用權柄吩咐污鬼,連污鬼也聽從了他.」「講左等於做左」是對我們的諷刺,但對上帝來說,「講左等於做左」是真實描述.還記得創世記第一章如何描述創造嗎?1:3上帝說:「要有光」,就有了光.」上帝是用話語創造世界的.

這個故事是如此結束的:

1:28耶穌的名聲就傳遍了加利利的四方.

或許,教會和個別信徒都想令人驚訝；或許,我們都想名聲傳遍,舉世聞名,世人慕名而來.耶穌和文士都適逢其會地在摩西的座位上教導人,耶穌令人驚訝的,是他所講的,就是實現的,不是他猜中事情如此發生,而是他把所說的話成真.或許,名聲傳遍,是因為人們渴慕那位說話成真的人.

五．總結:絕對言行一致的只有上帝,我們學習的是互相體諒

說話與實踐一致是高尚的情操,是我們應該追求的境界.對於假冒為善的人,可能是「非不能也,乃不為也」；對於真誠的人,可能是「非不為也,實不能也」.真正說話與實踐一致的,只有上帝而已,所以只有上帝可以是完全信賴的對像,也是有上帝是被讚美和敬拜的對象.今天經課的詩篇如此肯定上帝的言行一致:

111:3他所行的是尊榮和威嚴；他的公義存到永遠.

111:4他行了奇事,使人記念；耶和華有恩惠,有憐憫.

111:5他賜糧食給敬畏他的人；他必永遠記念他的約.

111:6他向百姓顯出大能的作為,把外邦的地賜給他們為業.

111:7他手所行的是誠實公平；他的訓詞都是確實的,

111:8是永永遠遠堅定的,是按誠實正直設立的.

111:9他向百姓施行救贖,命定他的約,直到永遠；他的名聖而可畏.

詩人擬人化地描述上帝所行的,他手所做的,是能夠令他的訓詞成為確實的,保證他的應許是永遠有效的.這不等於反過來說人不需要向著言行一致前進,只是人在各種限制之下,只能盡力而為而已.

如果我和你不是上帝直接把說話交給我要我宣告的先知,也不是擁有「說了就成」的權柄的耶穌,那麼,我們充其量只是一個受託傳承信仰的文土.如此,我們值得提醒自己,我們並沒有擁有其他人沒有的神喻,事實上大家都研讀著同一本聖經；我們也不擁有比其他人更大的能力和權柄,我們都是弟兄姊妹.那麼,我們如何實踐信仰呢?我們應該向甚麼方向行呢?讓我們以今天的書信經課使徒保羅提醒哥林多信徒的話來互勉:

8:1論到祭偶像之物,我們曉得我們都有知識.但知識是叫人自高自大,惟有愛心能造就人.

8:2若有人以為自己知道甚麼,按他所當知道的,他仍是不知道.

8:3若有人愛上帝,這人乃是上帝所知道的.

8:4論到吃祭偶像之物,我們知道偶像在世上算不得甚麼,也知道上帝只有一位,再沒有別的上帝.

8:5雖有稱為神的,或在天,或在地,就如那許多的神,許多的主；

8:6然而我們只有一位上帝,就是父－萬物都本於他；我們也歸於他－並有一位主,就是耶穌基督－萬物都是藉著他有的；我們也是藉著他有的.

8:7但人不都有這等知識.有人到如今因拜慣了偶像,就以為所吃的是祭偶像之物.他們的良心既然軟弱,也就污穢了.

8:8其實食物不能叫上帝看中我們,因為我們不吃也無損,吃也無益.

8:9只是你們要謹慎,恐怕你們這自由竟成了那軟弱人的絆腳石.

8:10若有人見你這有知識的,在偶像的廟裡坐席,這人的良心若是軟弱,豈不放膽去吃那祭偶像之物嗎?

8:11因此,基督為他死的那軟弱弟兄,也就因你的知識沉淪了.

8:12你們這樣得罪弟兄們,傷了他們軟弱的良心,就是得罪基督.

8:13所以,食物若叫我弟兄跌倒,我就永遠不吃肉,免得叫我弟兄跌倒了.

保羅在回應吃祭偶像之物的爭論中,他的首要考慮是:免得叫我弟兄跌倒.事情不是不可以討論,不是沒有對與錯之分；可以有合理的理解,也可以堅持這種理解；做不做某個行為,正如吃不吃祭肉,這些都是自由,不應被否定.不過,在這些事情上,保羅的心意是「造就人」,而不是傷害人.在醫療及助人專業中,有一個重要的守則是:Do No Harm,這項守則也值得基督徒考慮.

現代的基督徒有先知的角色,但不是如摩西一般上帝親自向我們交託律例典章；基督徒稱為基督徒當然是跟隨基督,但只是跟隨,而不是那位耶穌基督；基督徒就算以保存及傳承信仰為己任,但也得提醒自己保持謙虛的心,提防自己以為擁有那位啟示信仰者的權柄而帶來傷害.



\newpage

\section{眾人都在找你}
\label{sec:537}
\textbf{高國雄牧師 - 主啊:「眾人都在找你!」}
\newline
\newline
\hyperref[sec:536]{< < < PREV SERMON < < <}
~
\hyperref[sec:toc]{[返目錄]}
~
\hyperref[sec:538]{> > > NEXT SERMON > > >}


馬可福音 1:29-39
\newline
\begin{longtable}{cl}
\hline
\hline
章節 & 經文 (和合本修訂版)\\
\hline
1:29 & \begin{tabularx}{0.7\textwidth}{X} 他們一出會堂，就同雅各和約翰進了西門和安得烈的家。 \end{tabularx} \\ \\ \relax
1:30 & \begin{tabularx}{0.7\textwidth}{X} 西門的岳母正發燒躺著，就有人告訴耶穌。 \end{tabularx} \\ \\ \relax
1:31 & \begin{tabularx}{0.7\textwidth}{X} 耶穌進前拉著她的手，扶她起來，燒就退了，於是她服事他們。 \end{tabularx} \\ \\ \relax
1:32 & \begin{tabularx}{0.7\textwidth}{X} 傍晚日落的時候，有人帶著一切害病的和被鬼附的，來到耶穌跟前。 \end{tabularx} \\ \\ \relax
1:33 & \begin{tabularx}{0.7\textwidth}{X} 全城的人都聚集在門前。 \end{tabularx} \\ \\ \relax
1:34 & \begin{tabularx}{0.7\textwidth}{X} 耶穌治好了許多害各樣病的人，又趕出許多鬼，不許鬼說話，因為鬼認識他。 \end{tabularx} \\ \\ \relax
1:35 & \begin{tabularx}{0.7\textwidth}{X} 次日早晨，天未亮的時候，耶穌起來，到曠野地方去，在那裡禱告。 \end{tabularx} \\ \\ \relax
1:36 & \begin{tabularx}{0.7\textwidth}{X} 西門和同伴出去找他， \end{tabularx} \\ \\ \relax
1:37 & \begin{tabularx}{0.7\textwidth}{X} 找到了就對他說：「眾人都在找你！」 \end{tabularx} \\ \\ \relax
1:38 & \begin{tabularx}{0.7\textwidth}{X} 耶穌對他們說：「讓我們往別處去，到鄰近的鄉村，我也好在那裡傳道，因為我是為這事出來的。」 \end{tabularx} \\ \\ \relax
1:39 & \begin{tabularx}{0.7\textwidth}{X} 於是他走遍全加利利，在他們的會堂傳道，並且趕鬼。 \end{tabularx} \\ \\
[1ex]
\hline
\hline
\end{longtable}
講題:主啊:「眾人都在找你!」 My Lord:Everyone is looking for You!

講員:高國雄牧師

經文:馬可福音1:29-39

場合:香港中文大學崇基學院禮拜堂主日崇拜

日期:2021年2月7日

馬可福音記載耶穌基督的傳道服務按照祂的宣告開始「日期滿了,上帝的國近了.你們要悔改,信福音!」(可:1-15),呼召4位門徒後(西門,安德烈,雅各,約翰),先記載的神蹟就是趕鬼(可:1-21-28),然後是治好西門岳母的發燒…

馬可福音1:29-34

29他們一出會堂,就同雅各和約翰進了西門和安得烈的家. 30西門的岳母正發燒躺著,就有人告訴耶穌. 31耶穌進前拉著她的手,扶她起來,燒就退了,於是她服事他們.32傍晚日落的時候,有人帶著一切害病的和被鬼附的,來到耶穌跟前. 33全城的人都聚集在門前. 34耶穌治好了許多害各樣病的人,又趕出許多鬼,不許鬼說話,因為鬼認識他.

很明顯,福音書的作者希望展示出耶穌基督的能力,權柄及愛心,並且在一開始傳道服務時已經非常忙碌,受歡迎…

這導致

馬可福音1:35-37

35次日早晨,天未亮的時候,耶穌起來,到曠野地方去,在那裡禱告. 36西門和同伴出去找他, 37找到了就對他說:「眾人都在找你!」

「眾人都在找你!」=信眾的期望

人們為什麼要尋找耶穌呢?

我們又因為什麼原因要尋找耶穌呢?

我們是否只想找到我們自己所期望的「耶穌基督」呢?

耶穌對我們又有何期望呢?

馬可福音1:38

38耶穌對他們說:「讓我們往別處去,到鄰近的鄉村,我也好在那裡傳道,因為我是為這事出來的.」

耶穌的優先次序又是否我們的優次呢?

馬可福音1:35

次日早晨,天未亮的時候,耶穌起來,到曠野地方去,在那裡禱告.

耶穌接著的行動是…

馬可福音1:39

於是他走遍全加利利,在他們的會堂傳道,並且趕鬼.

看似與開始時一樣,但請留意「走遍全加利利」的價值及意義,背後的理念及態度…

傳道並且趕鬼,處理實際生活難處重要,但不是首要!

上主的心意 VS 我的期望,我的權益,我的價值觀,我所追求的美好

以賽亞書40:21-28

21你們豈不知道嗎?豈未曾聽見嗎?難道沒有人從起頭就告訴你們嗎?自從地的根基立定,你們豈不明白嗎?22上帝坐在地的穹窿之上,地上的居民有如蚱蜢.他鋪張穹蒼如幔子,展開諸天如可住的帳棚.23他使君王歸於虛無,使地上的審判官成為虛空.24他們剛栽上,剛種好,根也剛扎在地裡,經他一吹,就都枯乾；旋風將他們吹去,像碎秸一樣.25那聖者說:「你們將誰與我相比,與我相等呢?」26你們要向上舉目,看是誰創造這萬象,按數目領出它們,一一稱其名,以他的權能和他的大能大力,使它們一個都不缺.27 雅各啊,你為何說,以色列啊,你為何言,「我的道路向耶和華隱藏,我的冤屈上帝並不查問」?28你豈不曾知道嗎?你豈未曾聽見嗎?永在的上帝耶和華,創造地極的主,他不疲乏,也不困倦；他的智慧無法測度.

上帝由始至終也是坐著為王的上帝,擁有絕對的王權,是超越一切,及一切的主宰,全能的上帝,人無法相比,亦沒有任何東西可以向祂隱藏,上帝並不是人可以完全明白,理解的!

詩篇147:2-6,8-11

2耶和華建造耶路撒冷,聚集以色列中被趕散的人.3他醫好傷心的人,包紮他們的傷處.4他數點星宿的數目,一一稱它們的名.5我們的主本為大,大有能力,他的智慧無法測度.6耶和華扶持謙卑的人,將惡人傾覆於地.

8他用密雲遮天,為地預備雨水,使草生長在山上.9他賜食物給走獸和啼叫的小烏鴉.10他不喜悅馬的力大,不喜愛人的腿快.11耶和華喜愛敬畏他和盼望他慈愛的人.

同樣表明主上帝本為大,大有能力,智慧無法測度…

但祂不是只身處高高在上,而是滿有愛!

以賽亞書40:29-31

29疲乏的,他賜能力；軟弱的,他加力量.30就是年輕人也要疲乏困倦,強壯的也必全然跌倒.31但那等候耶和華的必重新得力.他們必如鷹展翅上騰；他們奔跑卻不困倦,行走卻不疲乏.

上帝選擇去愛人,成為我們力量真正的來源!

哥林多前書9:19-23

19我雖然是自由的,不受人管轄,但我甘心作了眾人的僕人,為贏得更多的人. 20對猶太人,我就作猶太人,為要贏得猶太人；對律法以下的人,我雖不在律法以下,還是作律法以下的人,為要贏得律法以下的人. 21對沒有律法的人,我就作沒有律法的人,為要贏得沒有律法的人；其實我在上帝面前,不是沒有律法,而是在基督的律法之下. 22對軟弱的人,我就作軟弱的人,為要贏得軟弱的人.對甚麼樣的人,我就作甚麼樣的人.無論如何我總要救一些人. 23凡我所做的,都是為福音的緣故,為要與人共享這福音的好處.

使徒保羅正正示範了如何實踐愛,並展現出真正信仰所追求的核心價值.

無論如何我總要救一些人.凡我所做的,都是為福音的緣故,為要與人共享這福音的好處.

從當日門徒向主耶穌的表達

主啊:「眾人都在找你!」

到今天我們,又或我們身邊的人都在尋找,尋問,今天你在聽道當中,當然也必然是一位尋道者；

耶穌基督就是道路,真理,生命,是唯一能去到天父上帝的路徑

我們從來沒有擁有真理,亦不代表真理,尋道者只可透過追隨主耶穌而向他人指向真理

所以人要尋找的,不應該是自我的成就,偉大的社會價值,而是單單上主本身!



\newpage

\section{登山顯容}
\label{sec:538}
\textbf{伍渭文牧師 - 登山顯容}
\newline
\newline
\hyperref[sec:537]{< < < PREV SERMON < < <}
~
\hyperref[sec:toc]{[返目錄]}
~
\hyperref[sec:539]{> > > NEXT SERMON > > >}


馬可福音 9:2-9
\newline
\begin{longtable}{cl}
\hline
\hline
章節 & 經文 (和合本修訂版)\\
\hline
9:2 & \begin{tabularx}{0.7\textwidth}{X} 過了六天，耶穌帶著彼得、雅各、約翰，領他們悄悄地上了高山。他在他們面前變了形像， \end{tabularx} \\ \\ \relax
9:3 & \begin{tabularx}{0.7\textwidth}{X} 衣服放光，極其潔白，地上漂布的人沒有一個能漂得那樣白。 \end{tabularx} \\ \\ \relax
9:4 & \begin{tabularx}{0.7\textwidth}{X} 有以利亞和摩西向他們顯現，並且與耶穌說話。 \end{tabularx} \\ \\ \relax
9:5 & \begin{tabularx}{0.7\textwidth}{X} 彼得對耶穌說：「拉比，我們在這裡真好！我們來搭三座棚，一座為你，一座為摩西，一座為以利亞。」 \end{tabularx} \\ \\ \relax
9:6 & \begin{tabularx}{0.7\textwidth}{X} 彼得不知道說甚麼才好，因為他們很害怕。 \end{tabularx} \\ \\ \relax
9:7 & \begin{tabularx}{0.7\textwidth}{X} 有一朵雲彩來遮蓋他們，又有聲音從雲彩裡出來，說：「這是我的愛子，你們要聽從他！」 \end{tabularx} \\ \\ \relax
9:8 & \begin{tabularx}{0.7\textwidth}{X} 門徒連忙向周圍觀看，不再看見任何人，只見耶穌同他們在一起。 \end{tabularx} \\ \\ \relax
9:9 & \begin{tabularx}{0.7\textwidth}{X} 下山的時候，耶穌囑咐他們說：「人子還沒有從死人中復活，你們不要把所看到的告訴人。」 \end{tabularx} \\ \\
[1ex]
\hline
\hline
\end{longtable}
講題:登山顯容 Transfiguration at the High Mountain

講員:伍渭文牧師

經文:馬可福音9:2-9

場合:香港中文大學崇基學院禮拜堂主日崇拜

日期:2021年2月14日登山顯容主日(Transfiguration Sunday)

海明威小說The Snows of Kilimanjaro(乞力馬扎羅的雪,或譯雪山盟) 一開始記載了坦桑尼亞的非洲第一高峰-乞力馬扎羅山 (菲洲語「神的殿堂」)山頂附近,有一隻凍僵的雪豹屍體,出現一個懸念:雪豹到這樣高寒的地方來尋找什麼?

1999年在世界最高的額非爾士峰腳下,發現了一具屍體,證實是英國登山運動員劍橋大學畢業生喬治.馬洛里(George Mallory).他在1922年參加登山隊,破了最高的記錄,8225公尺(26,980英呎).兩年後1924三十七歲的他,想著是最後的機會了,連同另一位登山運動員安德魯.歐文(Andrew Irvine)勇闖絕頂,在途中遇難,七十五年後才找到他的遺體.

曾有記者問馬洛里為何要冒險犯難?他答道:「因為它就在那裡」(Because it’s there).對登山運動員來說,世界上最高峰永遠是一個吸引,是體育精神的召喚.田徑運動員追求更高,更遠,更快的記錄,是跟對手公平競賽,登山是跟自己的極限競賽,也需群體協作,後勤支援,沒有其它人幫助,登上絕頂後也難保平安回來.

海明威《乞力馬扎羅的雪》給讀者宗教想像空間,雪豹可能相信上主的殿堂就在峰頂,在那裡可見到創造牠的主,真的嗎?馬洛里看到畢生努力的目標,人找到畢生奮鬥的目標是幸福的.浮士德一生尋找要攀登的高峰,甚至與魔鬼交易,出賣自己的靈魂,換取魔鬼的加持.但他一生仍在尋找中:他從哲學轉向法律,再轉向醫學和神學,最後轉向神秘的玄學.不少人以為找到了,發覺祇是鏡花水月,它根本不在那裡.

多年前,在美國一位越南船民好友曾告訴我:他認識一位北越士兵,攻陷南越西貢市時,帶著一公斤白沙糖預備送給「解放區」的親友.當然西貢甚麼都不缺,他發現被政府欺騙了.作家白樺《苦戀》有一名句:「您愛這國家,可是國家愛您嗎?」

1789年法國大革命倡導的「平等,自由,博愛」祇是一個口號.政權被推翻後,1793年發生很大的混亂和恐怖統治.不單路易十六,還有很多貴族和政敵,未經公平審訊就被送上斷頭台.一位貴族夫人被誣告,行刑前面向自由女神像說:「自由,自由,多少罪惡皆奉你的名而行.」

耶穌登上高山,因為山就在那裡；這山名叫加略,又稱各各他,是耶路撒冷城外的小山崗.他帶門徒登上高山,叫門徒定睛這矮矮的小丘,他將被掛在十架上,在山崗被舉起來,吸引萬民歸向他.我們看看今天福音經課.

1山上的榮光

過了六天,耶穌帶著彼得,雅各,約翰,暗暗的上了高山.就在他們面前變了形像,衣服放光,極其潔白,地上漂布的,沒有一個能漂得那樣白……有一朵雲彩來遮蓋他們,也有聲音從雲彩裡出來說:這是我的愛子,你們要聽從他.門徒忽然周圍一看,不再見一人,祇見耶穌同他們在那裡.下山的時候,耶穌囑咐他們說:人子還沒有從死裡復活,你們不要將所看見的告訴人.」 (可九2-9)

1)光影聲音,重要信息.耶穌高山顯容,衣服放光,並有聲音從雲彩出來說:這是我的愛子,你們要聽從他.顯容有光影有聲音,是傳遞重要信息.1972年基督教著名護教學者薛華(Francis Shaeffer)博士出版一本書名叫《他在那裡,他並不緘默》(He Is There and He Is Not Silent).上主在那裡又如何,人可以跟他接觸嗎?他有興趣和人溝通嗎?雪豹來到稱為上主殿堂乞力馬扎羅山峰頂門口又如何?造物主會理睬牠嗎?

國內著名搖滾歌手刀郎寫了《衝動的愆罰》,作為與前妻感情的墓誌銘.二十二歲的刀郎,與妻子生下女兒四十天,妻子便離開他.他尋妻十年,了無音訊.他知道妻子是在那裡,但拒絕溝通.

「那夜我喝醉了,拉著你的手,胡亂的說話,只顧著自己心中壓抑的想法,狂亂的表達, 我迷醉的眼睛, 已看不清你表情, 忘記了你當時會有怎樣的反應,我拉著手,放在我手心, 我錯誤的感覺到你也沒有生氣,所以我以為… 如果說不是老天讓緣份把我捉弄,想到你我就不會那麼心痛,就把你忘記吧,應該把你忘了,這是對衝動最好的懲罰 .」

刀郎的經驗,就像自然神論,創造的神就在那裡,但無法跟他接觸.但我們的神並非如此,他在那裡,他並不緘默.登山顯容,是上主想和門徒溝通,發出重要訊息:他要進耶路撒冷受許多的苦,被祭司長長老和文士們棄絕,最後掛在木頭上,親身擔當我們的罪.

2) 過了六天.六天前發生了甚麼事呢?彼得在該撒利亞腓立比公開認信耶穌是基督,跟著耶穌說人子要受許多的苦,並且被殺,第三天復活.彼得就拉著他勸他.「耶穌轉過頭來,看著門徒,就責備彼得說:撒旦,退我後邊吧,因為你不體貼上主的意思,祇體貼人的意思.」(可八33) 登山顯容,就要告訴門徒特別是彼得,要順服耶穌.「這是我的愛子,你們要聽他.」(九7)

3) 高山的試探.教會年曆把登山顯容放在預苦期前一個主日,是呼應周三的聖灰日.聖灰日是大齋期或預苦期的開始,教會把去年棕枝主日燒成灰燼,在該天塗在信徒額上:「你要記住,你本是塵土,仍要歸於塵土.當遠離罪,效忠基督.」塗灰顯容,指出神人之別；人是朽壞,上主不死.耶穌高山顯容,有聲音從雲彩裡出來說:「這是我的愛子.」神的兒子成為人,使人成為神的兒子(和女兒).但人登上高山,事業名氣都在峰頂,就有試探「高山變像」,當自己是上主.

耶穌在曠野受魔鬼第二,第三個試探,都在高處發生.魔鬼叫耶穌從聖殿頂跳下來,不會受傷；魔鬼帶耶穌到在最高的山,看見萬國的榮華,祇要耶穌俯伏拜一拜他,全都屬耶穌,不用受許多的苦,也不用上十字架.

前香港中文大學沈祖堯教授,有一次在校園基督徒午餐會分享:「從那麼高跳下來,怎會不受傷,除非你以為自己是上主,但確有頂尖的學者真的以為自己是上帝.」他所指是一位韓國教授,被揭發虛報自己成功首位複制羊,後來弄至聲名狼藉.沈教授說,世界級知名教授當自己是上帝不足為奇.他們容易被成功,名聲沖昏頭腦.他解釋:真的,當你在國際學術會議發表論文,巨大熒幕上單單看到你的大頭相,上千名來自世界各地的專家,屏息靜氣聽你在台上演講,之後報以熱烈掌聲.你真的感到飄飄然,大地就在腳下；再進一步,就會認為自己是上主,從高處跳下來也不會受傷.所以高山不在眼內,在心裡.

3) 高山在心中.梅頓(Thomas Merton)的屬靈自傳《七重山Seven Storey Mountian》,提及每個人心中有層層高山,攔阻我們親近上主.七重山的意念來自但丁神曲的煉獄界(地獄,煉獄,天堂三界)的七宗罪:驕傲,貪婪,貪色,暴怒,嫉妒和懶惰.七宗罪就是淨界的七重山,勝過這些試探,便臻天堂樂園.攀越七重山,便見到上主的榮光.

梅頓是一個備受愛戴,才華橫溢的詩人,政治評論家.世上的榮華,榮譽不能滿足他,心靈的空處,惟上主能填滿,他追求修道院獨處生活.他的自傳《七重山》自一九四八年出版以來即風行全球,譯成二十多種文字.《紐約時報》將之譽為二十世紀的《懺悔錄》.此書帶動美國年輕一代的屬靈覺醒,尤其響往修道院生活.

他效法奧古斯丁,作了一壓卷禱告:我喜愛獨處唯一的原因-在被造的東西中消失,向它們死去,不知道它們,因為他們提醒我離你多麼遠.它們提醒我:你和這些相距甚遠,雖然你在其中.你創造它們,它們因你而存在,它們把你遮蓋起來,讓我看不見你.我將獨自生活,離開它們.福哉獨處!(伍渭文中譯)

(That is the only reason why I desire solitude-to be lost to all created things, to die to them and to the knowledge of them, for they remind me of my distance from You. They tell me something about You: that You are far from them, even though You are in them. You have made them and Your presence sustains their being, and they hide You from me. And I would live alone, and out of them. O beata solitudo!-

-Thomas Merton, The Seven Storey Mountain. An Autobiography of Faith, (多馬斯.梅頓,七重山.信仰自傳),1998:461.

登山顯容,是上主的工作:「耶穌帶著彼得,雅各,約翰,暗暗的上了高山.」(可九2)顯容時,摩西以利亞出現,談論耶穌去世的事.不是耶穌帶領,我們不知方向；不是上主的啟示,我們不會明白.耶穌顯容的榮光,和聖道的榮光,指向福音,是白白的恩典,我們祇能憑信心接受.但接受福音後,我們要為福音發光,這是信徒的職份和責任.

2福音的榮光

1) 書信經課.林後四1-6:「我們既然蒙憐憫受了這職份,就不喪膽,乃將那些暗昧可恥的事棄絕了,不行詭詐,不謬講神的道,祇將真理表明出來,好叫在神面前把自己荐各人的良心……那吩咐光從黑暗裡照出來的神,已經照在我們心裡,叫我們得知神榮耀的光,顯在耶穌的面上.」

顯容在聽話.路加提到耶穌登山題容時是這樣說:「忽然有摩西以利亞兩個人,同耶穌說話,他們在榮光裡顯現,談論耶穌去世的事就是他在耶路撒冷將成的事.」摩西以利亞同時出現,父叫門徒聽從耶穌,表示登山顯容,重點不在顯容,乃在聽話.

在以馬忤斯路上,耶穌並沒有與兩位門徒顯出復活的肖容,反而隱藏在陌生的同路香客的身份,向兩位門徒解釋聖經.「於是從摩西和眾先知起,凡經上所指著自己的話,都給他們講解明白了.」(路二十四17)

經驗,無論是如何令你渾然忘我,也會隨時間漸漸模糊.分享個人宗教經驗,聽的人為你感恩,祇有羨慕.但分享思想,可以即時接收,而且候可以傳遞.那兩個門徒,其中一位叫革流巴,就夤夜披星帶月趕回耶路撒冷,把耶穌如何解釋聖經,基督這樣受苦又進入榮耀,傳遞給使徒,透過傳遞,他們自己也愈來愈明白.聖經的思想不同屬靈經驗,愈多看聖經,默想聖經,教導聖經,反覆細嚼經文,隨著日月推移,更能體會其意思.而屬靈經驗,日子夕遠便愈發模糊.

2)上主已發光.詩篇經課五十1, 2:「大能者神耶和華,已經發言招呼天下,從日出之地到日落之處.從全美的錫安中,神已經發光了.」山上顯容發出榮光,摩西以利亞代表律法書和先知書也發出光.聖經是我們腳前的燈,路上的光,不單相信,更要遵行,篤信力行.他不要祭牲,他要順服.「我的民啊,當聽我的話,以色列阿我要勸戒你.我是神,是你的神.我並不因你的祭物責備你,你的燔祭常在我面前.」(詩五十7,8)

3) 教導真理.舊約經課王下二1-12記述以利亞升天,以利沙繼續他的工作.除夕跟幾位因疫情今年無法回家過年的內地神學畢業生吃團年飯.飯後請國內校友分享信主的經歷.一位在昆明唸大學時,透過大學英語角的老師信主；昨天就傳來這位老師在除夕前一天,被上主接去.我安慰並鼓勵這同學:以利亞被上主接去,以利沙繼續；英語角老師被接去,英語角學生繼續,這正是王下二1-12的信息.

保羅在書信經課林後四1, 2,稱教導聖經,傳揚福音是他的職分,他不謬講神的道,祇將真理表明山來.高山顯容的榮光屬於上主的工作,但福音榮光屬於信徒的工作,傳福音的職份已托付給每一個信徒.

把光傳開.今天二月十四是正教的聖燭日.聖燭節又稱聖母行潔淨禮日,獻主節.是突顯基督為世界之光的節日.耶穌出生後四十日,馬利亞得潔淨,到聖殿燃點蜡燭祈禱.我們崇拜開始,也燃點蠟燭,象徵上主的臨在和我們的盼望,聖殿內就有長明不滅的金燈台.在聖殿西緬接過聖嬰,受聖靈感動唱出西緬頌:「主啊,如今可以照你的話,釋放僕人安然去世.」

西方教會重基督人性,聖燭日是二月二日,聖誕後十二月二十五日後四十天.東方教會強調耶穌的神性,一月六日耶穌領洗日為聖誕日,四十天就是今天二月十四日.領洗時聖父發聲,認證耶穌的神性:「你是我的愛子,我喜悅你」.(可一11)

聖燭日崇拜中有燃燭禮,象徵基督為世界之光.主禮首先以灑水和焚香祝聖蠟燭.灑水(sprinkling)除去物質之低下,焚香(censing)提升物質的尊貴,分別為聖,為上主所用；用禮儀的符號,形象地表達禱告的內容.灑水和焚香沒有能力使蜡燭聖化,是聖徒祈禱,求上主祝聖(Sanctify:).跟著依次讓信徒一個傳一個,把燭光傳開,並在教會巡行至每一角落,讓光充滿大地,信徒就是點亮世界的人.燃點蜡燭和巡遊時詩班唱Nunc Dimitti(西緬頌).拜結束時,信徒獲派洋燭回家,為主發光.

3)活出真理.保羅勸勉我們,為了傳遞福音真光的職份,我們要除去暗昧可恥的事,不行詭詐,不謬講神的道理,把自己薦與各人的良心.(林後四1,2)我們要肖似基督；人見基督,就認識父；人見到我們,也看見基督,所以我們稱為基督徒.

有人分辨聖像和偶像是不同的:聖像(icon)就像教堂的彩色玻璃,自然的光透過玻璃灑落教堂地上,莊嚴肅穆,是透光的.偶像使光停留在自己身上.一個裡外如一,沒有詭詐,沒有暗昧可恥,就能讓「基督榮耀福音的光」穿透我們的生命,使人認識基督.正教以聖徒人物繪畫的聖像,不按地上的比例,目的是引我們進入一個屬天的世界.

4) 暗暗的上了高山.要肖似基督,折射出福音榮光,就要親近上主,暗暗的上高山.高山使人變像,當自己是上主；但登高山亦能顯容-我來自塵土,也要歸回塵土.高山遠離塵囂,遠離世俗的牽挂.高山沒有網絡,可以放下手機,專心親近上主.這段時間,社交媒體食物圖片少了,湖光山色圖片多了.山川秀美,空氣清新,涼風拂臉,靈魂也清醒過來.登山臨絕頂,一覽群山小,惟上主巍峨偉大.與高山相比,人感到渺小和短暫.

耶穌選擇在高山顯容,說追求親近上主,總要付些代價.有一首聖詩「成聖須用工夫」,第一,二節:「成聖須用工夫,儆醒常禱告；時常與主親近,領受主真道,結交屬神兒女,幫助軟弱人,無論所作何事,莫忘求主恩.成聖須用工夫,世界雖忙碌,應當放下世事,單單親近主；時常仰望耶穌,更肖主模樣,親友藉你言行,可見主形像.」

戶外步行運動,我一般不帶手機,是默想祈禱的寶貴時刻,保持身體靈魂的健康.山上榮光,是基督顯容,他是父的肖像.信徒肖似基督,就能發出福音榮光.上主榮光不單在高山上彰顯；也在人世間的黑暗,苦難處彰顯；在人來人往的路上彰顯.因為他的榮光,就是他的憐憫；他以他的憐憫為他的榮耀.

3山下的榮光

「那吩咐光從黑暗裡照出來的神,已經照在我們心裡,叫我們得知神榮耀的光,顯在耶穌基督的面上.」(林後四6)

1)從黑暗出來的光.

光從黑暗裡照出來是很特別的提法.光暗不兩立,光不可能從黑暗出來,光一般是照進黑暗.從黑暗出來的光祇有兩次,一是創造:「起初神創造天地.地是空虛混沌,淵面黑暗.神的靈運行在水面,神說有光,就有了光.神看光是好的,就把光暗分開了.」(創一1-4)另一次是救贖:當耶穌掛在木頭上,從正午到申初三個小時,遍地黑暗,耶穌氣斷,埋在黑暗的地下,但翌日清晨,旭日破曉,耶穌復活.

古教會在復活前夕晚上領洗,聖洗禮儀的未段就是復活的一段,禮儀設計,是經歷從黑暗開出光明.其實每主日前夕的晚上,我們躺在黑暗中,像基督埋在地下,等候主日天亮,慶賀基督的復活.在北美數年參加教會崇拜,開始時有對答.牧師說:基督已復活!會眾答:是的,哈利路亞,他已復活了.這救贖的光,已經照在我們心裡.

2)大馬色路上的榮光.山上的榮光,祇有彼得,雅各,約翰三人看見過.彼得追憶:「他從父神得尊貴榮耀的時候,從極大榮光之中,在聲音說,這是我的途子,我所喜悅的.我們同他在聖山的時候,親自聽見這聲音從天上出來.」(彼後一17, 18)但保羅也見過基督的榮光,在往大馬色的路上,他拿著大祭司的文書正往大馬色搜捕基督徒.「忽然從天上發光,四面照著他.他就仆倒在地,聽見有聲音說,掃羅,掃羅,你為甚麼逼迫我.我說,主阿,你是誰.主說,我就是你所逼迫的耶穌.」(徒九3-5)

保羅有沒有看見基督的容貌呢?聖經說:「同行的人,站在那裡,說不出話來,聽見聲音,卻看不見人.」(徒九7)保羅與同行的人都見到光,聽到聲音,但惟有保羅仆倒在地,三天看不見東西.為甚麼?因為基督祇向保羅顯容.有解經家說保羅看見的基督,是臉相扭曲受苦的基督.掃羅不能接受耶穌是彌賽亞,而且是死在十架上的彌賽亞,上主就給他看見在十架上受苦基督的面.

3)顯在基督面上的榮光.「叫我們得知神榮耀的光,顯在耶穌基督的面上.」保羅寫這句話時,很有可能想起大馬色路上見到基督的面；但受苦扭曲的面,如何顯出神榮耀的光呢?除非上主以他的憐憫,作為他的榮耀；事實正如此.

摩西上西乃山,久久不下來,百姓按捺不住,造金牛犢代替上主敬拜.摩西下山聽到百姓圍繞金牛犢跳舞唱歌,忘形放肆,「便發烈怒,把兩塊版扔在山下摔碎了」.(出三十二19)

人不能憑自己遵守律法,摩西求上主赦免百姓的罪:「不然,求你從你所寫的冊上塗抹我的名.」(v. 32)摩西知道百姓不容易帶領,求上主同行,並顯出他的榮耀作為保證:「求你顯出你的榮耀給我看.」(出三十三18)「耶和華說:『我要顯出我一切的恩慈,在你面前經過,宣告我的名.我要恩待誰就恩待誰；要憐憫誰就憐憫誰』；又說:『你們不能看見我的面,因為人見我面不能存活.』耶和華說:『看哪,在我這裡有地方,你要站在磐石上.我榮耀經過的時候,我必將你放在磐石穴中,用我的手遮掩你,等我過去,然後我要將我的手收回,你就得見我的背,卻不得見我的面.』」(出三十三21-23)

從午正到申初遍地玄黑,因為人的罪.2014年香港展出巴洛克時期光暗大師卡拉瓦喬(Caravaggio)藏於義大利米蘭布雷拉畫廊(Pinacoteca di Brera)1606年作品《以馬忤斯的晚餐》.同場有本港藝術家以光暗為主題展出作品.其中一幅印象深刻.在潔白的熒幕,出現一條條捲動的誡命,白紙黑字.最初慢慢,後來加速.不可殺人,不可作假見證…….轉瞬間,潔白的熒幕變成玄黑；教我稍稍了解,午正的白日,如何成為黑暗.但黑亦是玄,玄妙之處是上主竟以憐憫為他的榮耀.要看上主的榮耀嗎?請看十字架上的基督,基督受苦扭曲的臉,綻發出上的主榮耀.

4)癲癇扭曲的臉相.從顯容的高山下山時,耶穌和門徒遇見一位父親,求醫治患癲癇症被鬼附身的孩子.他身體枯乾,病發時口中流沫,咬牙切齒,顯出極難看的臉相.父親說:「從小的時候,鬼屢次把他扔在火裡,水裡,要滅他,你若能作甚麼,求你憐憫我們,幫助我們.」(可九21, 22)耶穌對父親說:你若能信,在信的人,凡事都能.孩子的父親立時喊著說,我信.但我信不足,求主幫助.耶穌趕出污鬼,孩子好像死了一般,眾人多半都以為他死了.但耶穌拉著他的手,扶他起來,他就站起來了.

今天也許我們不容易看到鬼附以致臉相扭曲的人,但因著生活的重擔,壓得死去活來；沒有生命動力,沒有人生樂趣,甚至自詡是走肉行屍的人,為數卻不少.他們需要有人拉他們一把,扶他們起來,就可以站起來了.

領養兔唇娃娃.有一年校園基督徒團契到廣州退修參觀,住在白天鵝賓館.我至今難忘的一幕,是進入升降機時,同乘有三數外國家庭,拉在手的孩子趣致精靈,金頭髮,像年輕的父母.抱在懷中的年紀小,黑頭髮,全是兔唇的中國娃娃.他們來廣州辦理領養孤兒.聽他們說話,像是基督徒.

德蘭修女在加爾各答街頭,扶起躺在街頭的病患者,帶回修院醫治.頻臨死亡的,清洗他們的身體和臉龐,讓他們有尊嚴地離開.有人問她為甚麼這樣做?她說:「在受苦無助的人的臉上,我看見受苦的基督.」

總結:山上的榮光,不單在高山上；也落到山下的民宅,行人眾多的路上,也顯在每個受苦的人的臉上,因為上主憐憫他們.今天也許我們不容易看到鬼附以致臉相扭曲的人,但不少人因著生活的重擔,突然而來的橫逆,但願我們效法耶穌:「耶穌拉著他的手,扶他起來,他就站起來了.



\newpage

\section{身分考驗信息}
\label{sec:539}
\textbf{鄧瑞強博士 - 身分,考驗,信息}
\newline
\newline
\hyperref[sec:538]{< < < PREV SERMON < < <}
~
\hyperref[sec:toc]{[返目錄]}
~
\hyperref[sec:540]{> > > NEXT SERMON > > >}


馬可福音 1:9-15
\newline
\begin{longtable}{cl}
\hline
\hline
章節 & 經文 (和合本修訂版)\\
\hline
1:9 & \begin{tabularx}{0.7\textwidth}{X} 那時，耶穌從加利利的拿撒勒來，在約旦河裡受了約翰的洗。 \end{tabularx} \\ \\ \relax
1:10 & \begin{tabularx}{0.7\textwidth}{X} 他從水裡一上來，就看見天裂開了，聖靈彷彿鴿子降在他身上。 \end{tabularx} \\ \\ \relax
1:11 & \begin{tabularx}{0.7\textwidth}{X} 又有聲音從天上來，說：「你是我的愛子，我喜愛你。」 \end{tabularx} \\ \\ \relax
1:12 & \begin{tabularx}{0.7\textwidth}{X} 聖靈立刻把耶穌催促到曠野裡去。 \end{tabularx} \\ \\ \relax
1:13 & \begin{tabularx}{0.7\textwidth}{X} 他在曠野四十天，受撒但的試探，並與野獸同在一起，且有天使來伺候他。 \end{tabularx} \\ \\ \relax
1:14 & \begin{tabularx}{0.7\textwidth}{X} 約翰下監以後，耶穌來到加利利，宣講神的福音， \end{tabularx} \\ \\ \relax
1:15 & \begin{tabularx}{0.7\textwidth}{X} 說：「日期滿了，神的國近了。你們要悔改，信福音！」 \end{tabularx} \\ \\
[1ex]
\hline
\hline
\end{longtable}
講題:身分,考驗,信息 Identity, Trial and Message

講員:鄧瑞強博士

經文:馬可福音1:9-15

場合:香港中文大學崇基學院禮拜堂主日崇拜

日期:2021年2月21日

各位弟兄姊妹,早安.在這疫症裡,祝願各人平安.

今日仍是新春的日子,希望各人體會神的恩澤,生命更新.

但丁(Dante Alighieri)在他的《神曲》的開場白時這樣說:「35歲的那年,在人生路上,忽然進入了一座黑暗的森林,在其中迷途了.」35歲,但丁說,這是他人生的高峰.在這風華正茂的日子,忽然,進入了黑暗的森林,不知出路何在.不知有多少香港人,正有如此體會?

人生最大的森林,就是忘記了自己是誰,迷失在營營役役之中.

忘記了自己的真正名字.忘記了天父是怎樣叫喚我們的.忘記了要做什麼,才能讓自己體會神在萬古以先為我們確立的生命意義.忘記了,除了活著(living)以外,還有活著背後的那個深深藏著「神的印記」的存在(being).

今日,希望各人注目耶穌,在耶穌身上,找回自己的存在.

講道經文:馬可福音1:9－15

可1:9　那時,耶穌從加利利的拿撒勒來,在約旦河裡受了約翰的洗.

可1:10　他從水裡一上來,就看見天裂開了,聖靈彷彿鴿子,降在他身上.

可1:11　又有聲音從天上來,說:「你是我的愛子,我喜悅你.」

可1:12　聖靈就把耶穌催到曠野裡去.

可1:13　他在曠野四十天,受撒旦的試探,並與野獸同在一處,且有天使來伺候他.

可1:14　約翰下監以後,耶穌來到加利利,宣傳神的福音,

可1:15　說:「日期滿了,神的國近了.你們當悔改,信福音!」

經文分成三個小段落,分別討論:身分,考驗,信息.

(1)身分:馬可福音1:9－11

可1:9　那時,耶穌從加利利的拿撒勒來,在約旦河裡受了約翰的洗.

這裡,有兩點值得我們留意.

首先,是留意耶穌從何而來.

馬可在之前的經文(1:5),提到很多人去施洗約翰那裡受洗.那些人是從猶太地及耶路撒冷去的.猶太地,猶如香港的「中環」；耶路撒冷,猶如「半山區」.這裡的人,是有頭有面的人.這些地區,位處約旦河的南面.耶穌來的地方,剛好相反,位處約旦河的北面,是加利利地區的拿撒勒.《射鵰英雄傳》裡有「南帝」「北丐」,這講法頗合乎香港的地理經濟情況,「南區」是富裕的,「北區」是貧困的.這講法也剛好符合猶太地區與加利利地區的南北分野,南部的猶太地區及耶路撒冷,是高貴的,純種的,有權力的；北部的加利利地區,是低下的,混雜的,無權無勢的.神的兒子耶穌來自無權無勢之地,這是令人詫異的.

1:9這經文,第二點值得我們留意的地方是:這經文只簡單地描述耶穌在約旦河裡受了約翰的洗,並沒有交代耶穌為何要受約翰的洗.馬太福音意識到,耶穌受約翰的洗,看來不太恰當,起碼都要提個理由來吧.馬太福音給的理由,是耶穌「當盡諸般的義」.然而,馬可福音什麼也沒有交代.如此看來,在馬可福音裡,耶穌受洗,是無需什麼理由的,就是單純地與無權無勢的人認同.

可1:10　他從水裡一上來,就看見天裂開了,聖靈彷彿鴿子,降在他身上.

可1:11　又有聲音從天上來,說:「你是我的愛子,我喜悅你.」

耶穌受洗了,隨即,發生了三件事情.

一,「天裂開了」,這是神從天界闖入人間,神有所作為,要扭轉乾坤.通常,神闖進來,人面對神,是福禍相隨的.「福」者,是神看到我們的需要,加以顧念與眷佑.「禍」者,是神同時看到我們的罪惡,加以審判與懲罰.在舊約,人仰望神的面.但是,人看到神的面,卻是會死的.眷顧和審判,是相伴隨的.

神闖進人間,人間自會天翻地覆,因為黑暗會拼命拒絕光明,謊言會拼命拒絕真理,仇恨會拼命拒絕愛.在仇恨眼中,愛是不分黑白.在謊言眼中,真理是不識時務.在黑暗眼中,光明是令人討厭的過分耀眼.神闖入人間,人間便全力的要釘死祂.當邪惡「發爛渣」時,我們便知道,神闖入人間了.

問題是:當神闖入人間時,我們是站在神施恩的那一方,抑或是站在神審判的那一方呢?

二,「聖靈彷彿鴿子,降在他身上」.聖靈進到耶穌的生命裡,成為耶穌內在的生命.聖靈是耶穌生命的動力來源,也是耶穌生命的指導中心.耶穌不是為自己的喜好,目標,理想而活,耶穌在聖靈的指引下,活出天父的心意.耶穌的生活,不是為實現自己,而只為實現天父的旨意.

回看我們自己,我們到底為了什麼而活?

在神的世界裡,事情是很奇怪的.為了自己而活的人,他最後會失喪生命.為了愛他人而放下自己的人,他最後會得到生命.不斷奪取的,最後一無所得.甘心捨棄的,生命卻能豐盛而永恆.

問題是:我們有沒有時間靜下來,聆聽一下內心的聖靈的聲音?我們會否聽從聖靈的指引,放下自己?

三,「有聲音從天上來,說:『你是我的愛子,我喜悅你.』」這是天父這個「我」和耶穌這個「你」之間的「我你」對話,是父子間的私隱對話.「你是我的愛子,我喜悅你」.這是耶穌知道的,深藏內心的真正身分.這身分的認定,不關乎別人對你的看法,不關乎他人如何評價你,只關乎神對你的肯定.旁人的指指點點,並不是重點；重點是:神如何看你.而正正是這點,一個現代人,一個高度重視「自我」的人,一個追求「自我」的價值的人,最容易忘記的東西.我們活在一個「物化」的世界裡,談論的是「數字」,輕視神聖,我們不再接受由神來界定我們的身分,我們不再理會神給我們的名字,我們不願回應神的呼喚.

曾經,神給亞伯蘭一個新的名字,叫亞伯拉罕.神給雅各一個新的名字,叫以色列.神給掃羅一個新的名字,叫保羅.在這新的名字下,他們都在歷史裡留下神聖的足跡.其實,在神心中,也為我們預備了新的名字.問題是:我們更喜歡擁抱這世界給我們的名字,這名字可以是「成功者」,「卓越者」,「勝利者」,多過神給我們的名字.我們喜歡世人的稱讚,多於因神稱我們為祂的兒女而滿足.忘記了身分,忘記了名字,我們也就在黑暗的森林裡忘記了人生.

(2)考驗:馬可福音1:12－13

可1:12　聖靈就把耶穌催到曠野裡去.

可1:13　他在曠野四十天,受撒旦的試探,並與野獸同在一處,且有天使來伺候他.

一個人明白了自己的身分,便要接受這個身分的考驗.

可1:12　聖靈就把耶穌「催」到曠野裡去.

這裡的「催」字,有「強力催迫」的含意.耶穌生命裡的聖靈,催迫耶穌去接受考驗.考驗成功,神才能透過他,行在人間.考驗失敗,耶穌將一事無成.

考驗的地點,是曠野.曠野,是個一無所有的地方.沒有群眾的歡呼,沒有群眾的寄望,沒有群眾的掌聲.我們的世界,是一個「群眾」的世界.多少人的人生,活像是業餘的,或專業的youtuber,將生命的價值建立在群眾的likes之上.沒有了like,會變得無法活下去似的.吃每餐飯之前,很多人都忙得不亦樂乎,哪怕檯上只是一個普通得不再普通的漢堡包,也要從不同角度拍上一大輯相片,然後急不及待的放上網,然後期待眾人的回應,眾人的稱讚.若這些相片真的引來一大堆群眾的掌聲時,他便帶著無比滿足去吃這個包.這個包的味道如何,是不重要的.重要的,是群眾的認同.有了認同,就能獲得多一點存在的感覺.這點存在的感覺,支撐著他在下一輪飯,拍另一輪的相片.食物的相片,你來我往,於是,大家的手機,充斥著食物的相片.

聖靈催迫耶穌去到曠野.在這裡,什麼都沒有.有的,是天下間孤獨的靈魂,孤獨的我,以及隱伏在這孤獨靈魂最深處的撒旦,和隱伏在這孤獨靈魂更深處的上帝.

我們可以想像,當四野無人,麻醉我們的事物和群眾都沒有時,我們面對的,就是人性裡最純的邪惡.這邪惡,驅使我們成為一隻野獸,與野獸在一起,為慾望而生存,將生命變成純粹的慾望.

在這孤獨而無望的光景裡,耶穌渡過了他的考驗.耶穌一定是定睛注目那闖進人間的上帝.面對這位上帝時,他留意到聖靈輕柔的聲音,他甘心跟隨聖靈的指引,他放下了自我,放下了自己的規劃,放下了一切.他頭頂有天父,內心有聖靈,耶穌將自己融入神聖的生命中.在這一無所有的曠野,耶穌領悟到當外在的一切東西都沒有時那剩下來的無價的東西--他的身分,他是天父的愛子.耶穌已準備好,在人生的種種艱難中,都要守住這身分.當一個人以「神的兒子」這身分活在世上時,就無人能取消他的永恆價值.

問題是:我們能排除萬難,守住我們作為「神的兒子」,「神的女兒」這身分嗎?

(3)信息:馬可福音1:14－15

可1:14　約翰下監以後,耶穌來到加利利,宣傳神的福音,

可1:15　說:「日期滿了,神的國近了.你們當悔改,信福音!」

堅守自己身分的人,才有信息可言.

很多人沒有信息可傳,是因為在考驗中放棄了,忘記了他們在人生中的「天命」,遺忘了自己的「身分」,失去了神賜予的「名字」.一個只隨從世界的風向而飄來飄去的人,對這世界,是不會傳遞出任何信息的.

約翰,代表「舊時代」,「約翰下監以後」,即那「舊時代」過去了,「新時代」來臨了.

「日期滿了」,這裡的「日期」不是指日曆上的日期,不是指時鐘上的時間,而是指神行事的時間表.「日期滿了」,是指神行事的時機到了,神闖入人間的時機到了.神的闖入,帶來一個全新的世界,這就是「神的國」.

「神的國」近了.「近了」,是一個複雜的神學概念.「近了」,即是「未到」.但同時,「近了」,即是「近到可接觸到了」.「近了」,就是「若近若遠」.在這「若近若遠」的時間夾縫中,就是我們.我們的生命,可使天國「近」,也可使天國「遠」,視乎我們是否「悔改,信福音」.「悔改,信福音」,就是我們的視線不再注目世俗與自我,而轉而注目神；不再相信世俗與自我,而相信神.這是一種全新的人生,這是重尋天命的人生,這是找回自己真正的名字和身分的人生.這個世界,多一個這樣的人,天國就近一點.

問題是:我們真的信耶穌嗎?真的願意像耶穌一樣,經歷曠野的考驗嗎?真的願意像耶穌一樣,放下自我,擔起塵世的十架嗎?真的願意對天父說:「父啊,我願意作討你喜悅的兒女」嗎?事實上,找到這身分,生命才重獲意義.

願榮耀歸上主,願人間享和平.



\newpage

\section{我們與魔的距離}
\label{sec:540}
\textbf{林豪恩先生 - 我們與魔的距離}
\newline
\newline
\hyperref[sec:539]{< < < PREV SERMON < < <}
~
\hyperref[sec:toc]{[返目錄]}
~
\hyperref[sec:541]{> > > NEXT SERMON > > >}


講題:我們與魔的距離 The Devil in Us

講員:林豪恩先生

場合:香港中文大學崇基學院禮拜堂主日崇拜

日期:2021年2月28日

一．引言

2019年台劇「我們與惡的距離」引起廣泛關注.導演林君陽接受訪問時表示希望觀眾透過該劇思考究竟「什麼是惡」.他說:殺人是惡嗎?是,殺人可能是惡.那殺人的原因是什麼?那個原因也是惡嗎?這是可以討論的.惡裡面一定有那麼一點點善,錯的事情中一定可以找出動機,而那個原因不一定是那麼邪惡的事情.」(黃奕霖:「我們與惡的距離．一」)

我們在今天的福音經課讀到,耶穌的門徒彼得被耶穌責罵為撒旦,即是說彼得突然入了魔道,變成了魔鬼.這個奇特的片段,會否也刺激我們思考究竟甚麼是魔鬼呢?用導演林君陽的說話來問:「變成撒旦是入了魔嗎?是.但入魔的原因是甚麼?那個原因也是魔嗎?入魔的原恩惠不會有那麼一點點善呢?而那個入魔的原恩惠不會不一定是那麼邪惡的呢?

二．經文及反省

1．入魔的一點點善意

8:31從此,他教訓他們說:「人子必須受許多的苦,被長老,祭司長,和文士棄絕,並且被殺,過三天復活.」

8:32耶穌明明地說這話,彼得就拉著他,勸他.

8:33耶穌轉過來,看著門徒,就責備彼得

雖然耶穌責備彼得,但我想替彼得平反.

惻隱之心

孟子曰:「人皆有不忍人之心……今人乍見孺子將入於井,皆有怵惕惻隱之心.非所以內交於孺子之父母也,非所以要譽於鄉黨朋友也,非惡其聲而然也.由是觀之,無惻隱之心,非人也」.若我們善意地解讀彼得阻止耶穌去受苦,會不會也是惻忍之心的表現呢?對一個陌生人尚且有惻忍之心,對自己認識的人,甚至是所敬愛的人,豈不更關心和著緊嗎?在一個娛樂節目上,主持人問鄭伊健為甚麼不生孩子.鄭伊健回答說:「不生孩子是因為不想讓孩子在世間受苦受累.」你越愛惜的人,你越不想他們受苦.

未來盼望

不忍孩子受苦,除了是惻隱之心之外,更因為孩子是社會的未來.個人的籌命有限,無論如何偉大,總有離開的日子.正如傳道書所說,一代過去,一代又來,盼望和未來總是寄託在下一代身上的.「浸會大學前校長錢大康在電台節目指出,本港經過佔領中環及反修例事件,有聲音指要放棄年輕人,惟自己絕對不同意,年輕人是社會的未來,不可以放棄,亦不要只是批評,要明白年青人的需要,給予他們希望,重建對社會的信心.(信報2021年1月17日)下一代,其實是上一代的盼望所在,因此,社會總都是努力栽培下一代成長的.

彼得的未來在哪裡呢?彼得的盼望在哪裡呢?就在他的眼前,就是他眼前的耶穌.為甚麼這樣說呢?因為上文第8章29節記載耶穌問門徒:「你們說我是誰?」彼得回答他:「你是基督.」彼得獨具慧眼(馬太福音補充是天上的父啟示他的),認出耶穌是基督.認出耶穌是基督,並非等閒的事,而是彼得及整個民族的未來和盼望所在.何以見得呢?還記得馬太福音的聖誕故事中提及:「 4他(希律王)就召集了祭司長和民間的文士,問他們:「基督該生在哪裡?」祭司長和文士除了回答地方外外,還包含了盼望的所在.

5他們說:「在猶太的伯利恆.因為有先知記著:

6『猶大地的伯利恆啊,

你在猶大諸城中並不是最小的；

因為將來有一位統治者要從你那裡出來,

牧養我以色列民.』」

當時,絕大部分人都認不出耶穌就是基督,只有彼得發現這位就是大衛的苗裔,是以色列的盼望.彼得要保護耶穌,阻止耶穌受苦和被殺,不單是合情合理,而且是大仁大義.如果我們還未能理解,不妨回憶一下香港宋王台的故事.「南宋末年,蒙古軍隊大舉南侵,先後攻佔了襄樊,建康和臨安.宋軍節節敗退,皇室成員被殺被擄.蒙古大軍依舊窮追不捨,結果兩小末代皇兄弟被逼流亡.在文天祥,陸秀夫,張世傑,陳宜中等人護送下,他們經過海路先後途經泉州,潮州,惠州,廣州等地,最終逃至九龍官富場(即今九龍城附近),在此建立行宮.」(維基百科).文天祥,陸秀夫等將軍保護宋朝的苗裔被譽為中神,受後人稱頌；彼得保護耶穌,就算不被稱頌,都不應該被責難吧!

2.不體貼上帝的意思只體貼人的意思就成魔

8:33耶穌轉過來,看著門徒,就責備彼得說:「撒旦,退我後邊去吧!因為你不體貼上帝的意思,只體貼人的意思.」

受到耶穌如此回應,彼得可能感到「好心著雷劈」.其實,相信大家都體會過「好心無好報」的滋味,尤其是作為父母的.常常聽到父母對兒女說:「我做的一切都是為你好的,難道我會害你嗎?」.然而,子女的回應可能令父母感到「好心無好報」.有一位少年在網上留言區傾訴她的苦惱說:「最近真的很煩,常常會被媽媽嘮叨,被責罵.我把浴袍放在床上,媽媽又一直嘈,我不想跟她吵,但她又說:「你那張臭臉是怎樣?跟你說你就不高興?」每次都罵我房間亂七八糟,但是這幾天我確實整理了,可是她好像仍是不滿意.我不是三歲的小孩,我有想法的.連我的房間也要管,也不准我去同學家,跟朋友出去玩也不行,整天找一些強暴或一堆關於青少年的可怕新聞給我,真的很煩.」

作為旁觀者,你一定看得出那位媽媽的用意是保護她的女兒,避免她受到傷害.然而,結果卻令媽媽自己感到委屈,令女兒感到憤怒.那位少女最後說:「我不想回家了,我不想面對媽媽.有甚麼解決方法呢?」網友們的留言主要都是建議少女和媽媽互相聆聽對方,增加了解.

如果「好心」都不是關鍵,那麼,甚麼才能補救呢?沒錯,就是網友建議的「聆聽和了解」.彼得何以突然成魔呢?未必是他心腸惡毒.耶穌是這樣判斷的:「因為你不體貼上帝的意思,只體貼人的意思」.原來,是不是變成了撒旦,是不是入了魔道,分辨的標準是「體貼上帝的意思」還是「只體貼人的意思」.在人與人之間尚且要聆聽對方,希望確認「我都是為你好的」那些「好」,對方都認為是好,那就不會產生太多誤會和勉強,抱怨和憤怒.在人與上帝之間,豈不需要更多的聆聽和確認嗎?在教會歷史中,我們已經太多「為上帝的名」或「奉上帝的名」而做成的傷害了.

甚麼才是上帝的意思呢?上帝的意思是甚麼呢?在那一個片段中,耶穌是如此說的:

3.體貼上帝的意思就是跟隨基督捨己救人

8:34於是叫眾人和門徒來,對他們說:「若有人要跟從我,就當捨己,背起他的十字架來跟從我.

關於在各種歷史狀態和人生處境中如何體貼上帝的意思,聖經中記載了眾多的例子.在這個事件中,從耶穌對門徒的教導來看,體貼上帝的意思的第一步是跟隨基督.「跟隨基督」或「追隨基督」,正是潘霍華其中一本書的中文名.潘霍華經常問自己:「耶穌想要對我們說甚麼麼呢?今天祂對我們的旨意是什麼呢?在現代的世界中,祂如何幫助我們作一個好基督徒呢?」當時,耶穌對門徒的回答是:「就當捨己,背起他的十字架來跟從我」.

耶穌要我背起十字架,我的十字架是甚麼是呢?這個問題引起無限的聯想:有人以為人生的經歷和遭遇就是我們的十字架,例如:有人的丈夫脾氣急躁,有人的配偶不信主,有人碰到難相處的老闆或同事,有人是生病,失業,失婚等,我們以為這些就是「十字架」,然而,若果我們不是從自己的角度去看,而是從耶穌的角度去看,十字架所代表的可能不是這些.如果不是這些,那麼會是甚麼呢?

以賽亞書53章的僕人之歌有如此的描述:

4他誠然擔當我們的憂患,

背負我們的痛苦；

5他為我們的過犯受害,

為我們的罪孽被壓傷.

因他受的懲罰,我們得平安；

因他受的鞭傷,我們得醫治.

耶穌背起他的十字架,並不是背起自己人生中的失意和挫折,然後發奮向上,變為成功人士；耶穌的十字架,是擔當其他人的憂患,背負其他人的痛苦,為世人帶來平安和醫治；因此,耶穌的十字架之路,是捨己救人之路.耶穌呼喚他的跟隨者反省及改變人生的追求和目標.

4．生命的意義在於為他者而活

如何反省人生的追求,生命的意義呢?如果要改變,向甚麼方向改變呢?耶穌以強烈的對比要求門徒思考和選擇:

8:35因為,凡要救自己生命的,必喪掉生命；凡為我和福音喪掉生命的,必救了生命.

8:36人就是賺得全世界,賠上自己的生命,有甚麼益處呢?

8:37人還能拿甚麼換生命呢?

「要救自己生命」,要「賺得全世界」都是世界認可的目標,也是人之常情,可能就是耶穌責備彼彼所說「你只體貼人的意思」一句中「人的意思」所代表的.「人的意思」暗示著人只為自己而活,具體表現的退守立場是「要救自己生命」,進攻立場是要「賺得全世界」.耶穌呼籲他的跟隨著反省,活著的目標是否就只是令自己繼續生存呢?如果只是這樣,生活是充滿生存的焦慮,還是生活的豐富呢?如果只是這樣,生命有甚麼意義呢?有人反省之後說:「生命不在乎長短,而在乎活得有價值」.即是說,生命是可以多於生存本身的,也可以多於生存的個體本身.突破了生存本身及生存的個體本身,就會看見他人,也可以延伸至他人,耶穌所說「為我和福音喪掉生命的」必得著生命,就是為他者而活,在為他者而活之中體會到生命的價值及實在.

耶穌所說的好像很難明白,也好像曲高和寡,其實不是,人們在日常生活中也可以體會到的.人們常說「助人為快樂之本」,「施比受更為有福」大概就是這個方向吧!觀察嬰孩的成長,有沒有發現某個階段的BB總喜歡說「我來幫你啊」,甚麼也想參與.不過,因為他們能力有限,他們通常都越幫越忙,因此被大人拒絕,甚至責罵.其實,孩子很小就喜歡幫忙,也是因為他能藉著幫忙參與大人的活動,並從中感受到自我肯定,自尊和自信.就算是成人,當你幫助了一個人渡過難關,你也會感到喜悅和有意義.

人要「救自己生命」,可能只是受到威脅的反應,就正如其他動物一樣進入「fight or flight」的模式之中,然而這不應該是常態,如果是常態,就是高度焦慮的狀況,有可能誘發情緒病.另外,要透過「賺得全世界」證明生命的價值,可能只是比較和競爭氣氛過於激烈帶來的欲望,一般來說,人類都是在和平共處,彼此幫助的狀況之下更感到安全和正常的.世界上有一個非常特別的國家,冬天寒冷而且漫長,期間日照時間短而黑夜時間長,我們以為生活於灰暗的天氣之下的人應該會心情也灰暗,然而,那裡卻是世界上最快樂的國家之一.是甚麼令他們快樂呢?其中一個原因是這個只有33萬人口的小國,人民之間培養出深厚的共存意識,也高度地互相信任,96\%的冰島人深信有值得信任的人在他們需要幫助時向他們伸出援手.潘霍華用信仰的語言來解說這種狀況,他說:「因為基督擔當我們的重擔,所以我們也該擔當同伴的重擔」.

三．總結

是人是魔,可以在轉念之間.回到電視劇「我們與惡的距離」,導演林君陽說:「我們都以為自己是好人,但我們很多時候都有惡念而不自知,有些人的惡念展現出來,有些人沒有,而每個人的善念與惡念都主導了現在這個世界的面貌.如果每個人都能意識到這件事,這世界會不會好一點點?」製作人林昱伶則認為,大家都把中文劇名重點放在「惡」上,但更重要的應該是「我們」.如同林君陽所說,這世界是由「我們」所組成,劇組在劇名的字體設計上也刻意讓「我們」佔較大的比例,希望提醒觀眾,這兩個字才是劇情的核心,而非糾結於「誰是惡人」,「誰是善人」.(黃奕霖: 【我們與惡的距離．一】善惡難分界　活在灰色世界的「我們」)

看完彼得的故事,或者我們也應該把焦點轉回「我們」.說到底,正邪榮恥的終極判決不在於我們,反而是我們也是等待被判決的:

8:38凡在這淫亂罪惡的世代,把我和我的道當作可恥的,人子在他父的榮耀裡,同聖天使降臨的時候,也要把那人當作可恥的.」

讓我們一齊跟隨基督,「在我們的主耶穌同他眾聖徒來臨的時候,在我們父神面前,成為聖潔,無可指責.」(帖撒羅尼迦前書3:13)



\newpage

\section{來就光}
\label{sec:541}
\textbf{高國雄牧師 - 來就光!}
\newline
\newline
\hyperref[sec:540]{< < < PREV SERMON < < <}
~
\hyperref[sec:toc]{[返目錄]}
~
\hyperref[sec:542]{> > > NEXT SERMON > > >}


約翰福音 3:14-21
\newline
\begin{longtable}{cl}
\hline
\hline
章節 & 經文 (和合本修訂版)\\
\hline
3:14 & \begin{tabularx}{0.7\textwidth}{X} 摩西在曠野怎樣舉蛇，人子也必須照樣被舉起來， \end{tabularx} \\ \\ \relax
3:15 & \begin{tabularx}{0.7\textwidth}{X} 要使一切信他的人都得永生。 \end{tabularx} \\ \\ \relax
3:16 & \begin{tabularx}{0.7\textwidth}{X} 「神愛世人，甚至將他獨一的兒子賜給他們，叫一切信他的人不致滅亡，反得永生。 \end{tabularx} \\ \\ \relax
3:17 & \begin{tabularx}{0.7\textwidth}{X} 因為神差他的兒子到世上來，不是要定世人的罪，而是要使世人因他得救。 \end{tabularx} \\ \\ \relax
3:18 & \begin{tabularx}{0.7\textwidth}{X} 信他的人不被定罪；不信的人已經被定罪了，因為他不信神獨一兒子的名。 \end{tabularx} \\ \\ \relax
3:19 & \begin{tabularx}{0.7\textwidth}{X} 光來到世上，世人因自己的行為是惡的，不愛光，倒愛黑暗，這就定了他們的罪。 \end{tabularx} \\ \\ \relax
3:20 & \begin{tabularx}{0.7\textwidth}{X} 凡作惡的人都恨惡光，不來接近光，恐怕他的行為被暴露。 \end{tabularx} \\ \\ \relax
3:21 & \begin{tabularx}{0.7\textwidth}{X} 但實行真理的人就來接近光，為要顯明他的行為是靠神而行的。」 \end{tabularx} \\ \\
[1ex]
\hline
\hline
\end{longtable}
講題:來就光!　Come to the Light

講員:高國雄牧師

經文:約翰福音3:14-21

場合:香港中文大學崇基學院禮拜堂主日崇拜

日期:2021年3月14日

約翰福音3:14-21

14摩西在曠野怎樣舉蛇,人子也必須照樣被舉起來,15要使一切信他的人都得永生.16「上帝愛世人,甚至將他獨一的兒子賜給他們,叫一切信他的人不致滅亡,反得永生.17因為上帝差他的兒子到世上來,不是要定世人的罪,而是要使世人因他得救.18信他的人不被定罪；不信的人已經被定罪了,因為他不信上帝獨一兒子的名.19光來到世上,世人因自己的行為是惡的,不愛光,倒愛黑暗,這就定了他們的罪.20凡作惡的人都恨惡光,不來接近光,恐怕他的行為被暴露.21但實行真理的人就來接近光,為要顯明他的行為是靠上帝而行的.」

你喜愛光,還是暗呢?為什麼?為何有人喜愛黑暗呢?

舊約故事:蛇被舉起,人就得救,摩西為何需要舉蛇呢?

民數記21:4-9

4他們從何珥山起行,繞過以東地往紅海那條路走.在路上,百姓心中煩躁.5百姓向上帝和摩西發怨言,說:「你們為甚麼把我們從埃及領上來,使我們死在曠野呢?這裡沒有糧食,沒有水,我們厭惡這淡而無味的食物.」6耶和華派火蛇進入百姓當中去咬他們,於是以色列中死了許多百姓.7百姓到摩西那裡,說:「我們有罪了,因為我們向耶和華和你發怨言.求你向耶和華禱告,叫蛇離開我們.」於是摩西為百姓禱告.8耶和華對摩西說:「你要造一條火蛇,掛在杆子上.凡被咬的,一望這蛇就必存活.」9摩西就造了一條銅蛇,掛在杆子上.凡被蛇咬的,一望這銅蛇就活了.

為什麼百姓心中煩躁,向上帝和摩西及怨言?我們現在生活煩躁嗎?為什麼?

要先回想；上帝的拯救,離開埃及,過紅海,嗎哪…

「使我們死在曠野!」內容屬實嗎?真的沒有糧食,沒有水

人的耐性,自以為事,上帝的計劃與我的不同

上帝沒有帶來人的滅亡,發怨言才會…

人何時會發現自己做錯呢?罪帶來審判,懲罰,刑罰讓人知罪,願意悔改

人呼求上帝,上帝回應,施行拯救

「聯合國世界衛生組織」(WHO),會徽是由聯合國的標誌,再加一條蛇盤繞權杖的圖像湊成

耶穌被舉起,我們得永生

耶穌來,不是要定人的罪,乃是要救人,罪是因自己而定的

以弗所書2:1-10

1從前,你們因著自己的過犯罪惡而死了.2那時,你們在過犯罪惡中生活,隨從今世的風俗,順服空中掌權者的領袖,就是現今在悖逆的人心中運行的邪靈.3我們從前也都生活在他們當中,放縱肉體的私慾,隨著肉體和心中的意念去做,和別人一樣,生來就是該受懲罰的人.4然而,上帝有豐富的憐憫,因著他愛我們的大愛,5竟在我們因過犯而死了的時候,使我們與基督一同活過來-可見你們得救是本乎恩-6他又使我們在基督耶穌裡與他一同復活,一同坐在天上,7為要把他極豐富的恩典,就是他在基督耶穌裡向我們所施的恩慈,顯明給後來的世代.8你們得救是本乎恩,也因著信；這並不是出於自己,而是上帝所賜的；9也不是出於行為,免得有人自誇.10我們是他所造之物,在基督耶穌裡創造的,為要使我們行善,就是上帝早已預備好要我們做的.

罪與罰的關係,我們是因自己的過犯罪惡而死

人最大的罪,就是順服的對象錯了

上帝的愛,與人的愛不一樣,上帝竟在我們因過犯而死了的時候,使我們與基督一同活過來

唯獨恩典!!唯獨信心!!

恩,信

不是出於自己,仍是上帝所賜的

福音書中耶穌為何說起摩西舉銅蛇的故事?

約翰福音3:1-13

1有一個法利賽人,名叫尼哥德慕,是猶太人的官.2這人夜裡來見耶穌,對他說:「拉比,我們知道你是由上帝那裡來作老師的；因為你所行的神蹟,若沒有上帝同在,無人能行.」3耶穌回答他說:「我實實在在地告訴你,人若不重生,就不能見上帝的國.」4尼哥德慕對他說:「人已經老了,如何能重生呢?豈能再進母腹生出來嗎?」5耶穌回答:「我實實在在地告訴你,人若不是從水和聖靈生的,就不能進上帝的國.6從肉身生的就是肉身；從靈生的就是靈.7我說『你們必須重生』,你不要驚訝. 8風隨著意思吹,你聽見風的聲音,卻不知道是從哪裡來,往哪裡去；凡從聖靈生的也是如此.」 9尼哥德慕問他:「怎麼能有這些事呢?」10耶穌回答,對他說:「你是以色列人的老師,還不明白這些事嗎?11我實實在在地告訴你,我們所說的是我們知道的,我們所見證的是我們見過的,你們卻不領受我們的見證. 12我對你們說地上的事,你們尚且不信,若對你們說天上的事,如何能信呢? 13除了從天降下的人子,沒有人升過天.

人不是不明白,不知道,而是不願意接受,不就近上帝,因為行惡,不愛光

甚至愛黑暗,因這就定了罪

起來,回轉,就近那光,反映那光,一同光照世界!



\newpage

\section{因所受的苦難學了順從}
\label{sec:542}
\textbf{鄧瑞強博士 - 因所受的苦難學了順從}
\newline
\newline
\hyperref[sec:541]{< < < PREV SERMON < < <}
~
\hyperref[sec:toc]{[返目錄]}
~
\hyperref[sec:543]{> > > NEXT SERMON > > >}


希伯來書 5:5-10
\newline
\begin{longtable}{cl}
\hline
\hline
章節 & 經文 (和合本修訂版)\\
\hline
5:5 & \begin{tabularx}{0.7\textwidth}{X} 同樣，基督也沒有自取作大祭司的榮耀，而是在乎向他說話的那一位，他說：「你是我的兒子，我今日生了你。」 \end{tabularx} \\ \\ \relax
5:6 & \begin{tabularx}{0.7\textwidth}{X} 就如又有一處說：「你是照著麥基洗德的體系永遠為祭司。」 \end{tabularx} \\ \\ \relax
5:7 & \begin{tabularx}{0.7\textwidth}{X} 基督在他肉身的日子，曾大聲哀哭，流淚禱告，懇求那能救他免死的神，就因他的虔誠蒙了應允。 \end{tabularx} \\ \\ \relax
5:8 & \begin{tabularx}{0.7\textwidth}{X} 他雖然為兒子，還是因所受的苦難學了順從。 \end{tabularx} \\ \\ \relax
5:9 & \begin{tabularx}{0.7\textwidth}{X} 既然他得以完全，就為凡順從他的人成了永遠得救的根源， \end{tabularx} \\ \\ \relax
5:10 & \begin{tabularx}{0.7\textwidth}{X} 並蒙神照著麥基洗德的體系宣稱他為大祭司。 \end{tabularx} \\ \\
[1ex]
\hline
\hline
\end{longtable}
講題:因所受的苦難學了順從　Learning Obedience through Suffering

經文:希伯來書5:5-10

講員:鄧瑞強博士

場合:香港中文大學崇基學院禮拜堂主日崇拜

日期:2021年3月21日

各位弟兄姊妹,早安.在這疫症期間,祝願各人平安.

下禮拜五,便是主耶穌的受難日了.今日,讓我們思想主耶穌的苦難.

苦難,從來都是人類思想上的謎.在宗教哲學裡,有一種講法,說苦難的存在,是為了成就更偉大的美善,猶如畫面中的黑暗,是為了凸現更豐富的色彩.我不同意這個講法.這個講法,會合理化苦難.我看不到在戰爭中數不盡的無辜小孩子的死,對文明的進步有何必要性.在舊約,約伯記探討苦難時,從來沒有替苦難找一個合理的理由,甚至連道德的理由都找不到.苦難,是一個謎.約伯能做的,是沉默,然後,神向他顯現.這不是要合理化苦難,而是要提升人的境界.人在神聖的境界裡,才有面對苦難的力量.

在二次世界大戰時,有群猶太人,為逃避「納粹」,躲在科隆的一個地下室裡.在牆上,他們留下了這些文字.

「沒有陽光的日子,我仍相信太陽.感受不到愛的歲月,我仍相信愛.當上帝隱匿時,我仍相信上帝.」

(I believe in the sun even when it is not shining

And I believe in love even when there’s no one there

And I believe in God even when He is silent)

有人將這些文字譜成詩歌,歌名叫＜Inscription of Hope＞,這歌有廣東話版本,叫＜盼望銘刻＞.這些銘刻在牆上的文字,寫出了真信仰.「當上帝隱匿時,我仍相信上帝.」唯當我們進入神聖境界時,才能在一個毫無上帝的世界,深信上帝.唯有深信上帝的人,才能在苦難的世界中,堅持活出愛.

今日的講道經文,是今日的書信經課:希伯來書5:5－10

來5:5　如此,基督也不是自取榮耀作大祭司,乃是在乎向他說「你是我的兒子,我今日生你」的那一位；

來5:6　就如經上又有一處說:「你是照著麥基洗德的等次永遠為祭司.」

來5:7　基督在肉體的時候,既大聲哀哭,流淚禱告,懇求那能救他免死的主,就因他的虔誠蒙了應允.

來5:8　他雖然為兒子,還是因所受的苦難學了順從.

來5:9　他既得以完全,就為凡順從他的人成了永遠得救的根源,

來5:10　並蒙神照著麥基洗德的等次稱他為大祭司.

我先講耶穌的苦難經驗.

來5:7　基督在肉體的時候,既大聲哀哭,流淚禱告,懇求那能救他免死的主,就因他的虔誠蒙了應允.

「基督在肉體的時候」:這是基督的塵世經驗.他沒有特權,過一幸福快樂的生活.他來到人間,體會到人類最深刻的受苦經驗.

「大聲哀哭,流淚禱告,懇求…」,又大叫,又流淚,又禱告,又懇求,連續幾個字,便看出耶穌忍受的,是大悲苦,大艱難.

他求神救他免死.這不是普通的死,這是受盡折磨的死.幾乎所有人,都會在這種折磨中低頭.耶穌既然在苦難中懇求,就表明他有掙扎.那一夜,那苦難前的一夜,是漫漫長夜.萬般思緒,湧上心頭.這苦難真有意義嗎?難道這個墮落世界,會因一個無辜人的受苦而改變嗎?數不盡的無辜人受過苦,有人留意過他們嗎?有人為他們流淚嗎?有人流淚又如何,強暴的人會放下暴力嗎?這個世界,無一日沒有十架,無一日沒有釘在十架上的人,但世界會因此而改變嗎?為何醉生夢死的人可終日歡樂,而我竟孤獨地背起十架?這真是我的天命嗎?背負起這天命,又如何?若我逃避了,世界真有損失嗎?

天地無言,上帝隱沒,這是最淒涼的苦難.

電影《基督最後的誘惑》形象地呈現耶穌的掙扎.放下十架,過一平凡人的生活,娶妻生子,歸隱山林,有何不可?

經文說:「就因他的虔誠蒙了應允」.

「虔誠」這個字,有尊敬也有畏懼的含意,合起來,正好是「敬畏」之意.耶穌敬畏神.他懼怕苦難,但他更敬畏神.他可能對苦難的人生充滿懷疑,但他對神毫不懷疑.在苦難前夕,他向天父哀求.天父默默無言,讓他去承受.他與天父之間,有一種外人難以理解的交流.在這交流中,漸漸,聽命於天父的心,大於對自身的顧慮.漸漸,天父的無聲,變成內心裡唯一的大聲.有些人的內心,外人很難明白的.話說,多年前,在颱風日子,有人去海邊看海,卻被風扯了落海.有人經過,落水救人.最後,被風扯落海的人獲救,救人者卻遇溺死亡.這個人,當見到有人墮海遇溺時,在那瞬間,在那永恆的一剎那,他到底在想什麼?他應該有想過家人,他的家人就在身旁.他應該有衡量過自己的安危,那的確不是尋常的日子,那是狂風巨浪的日子.但是,我想,應該有一種巨大的力量抓住他,這力量是沉默的,但這力量是巨大的,巨大得讓這個人放下了自己,放下了家人,走向死亡的漩渦.或者說,應該有一種超乎尋常的,高貴的生命呼喚著他,他心悅誠服地降服,然後,他放下了恐懼,走向一個需要他幫助的人.這是超乎常理的,難以理解的.這種人與天父之間,有一種外人難以理解的交流.

經文說,因耶穌的虔誠,因耶穌對神的敬畏,他蒙了應允.「他蒙了應允」,這句子真耐人尋味.耶穌不是因為「蒙了應允」而不用受苦受死,反而正正是「蒙了應允」,而要去受苦受死.

在什麼意義上,他「蒙了應允」呢?

我想,當耶穌以敬畏的心面對天父時,他像那落海救人的人那樣,被偉大的力量抓住,被天父呼喚,被天心充溢.然後,他明白,短暫的人生,其意義在於實踐天父的心意,讓天心彰顯人間.他也明白,天父愛世人,祂的心與受苦的人同在.然後,他深刻體會到,在自己的受苦裡,正體現著天父與受苦的人的同在.再然後,他明白到,在風平浪靜的日子,無人墮海,自己能過一簡單生活,是好的.但在風雨飄搖的日子,有人墮海,則放下自己,聽從一更大的生命的呼喚,也是理所當然的.這是天父的心在人間的彰顯.這是大愛,這是救贖,這是希望.耶穌明白了,他願意這樣做.在這意義上,他這「願意」是「蒙了應允」的.

在受苦的事上,耶穌領悟到什麼?

來5:8　他雖然為兒子,還是因所受的苦難學了順從.

這節經文,就是我今日講道的題目.我覺得,這節的關鍵是一個「學」字.這不是死讀書,學習不同理論的那種「學」,而是像學徒那樣,在實踐中學習的那種「學」.耶穌不是先學習「順從」,然後去承受苦難.根本沒有這種學.實踐和領悟,是同一回事.耶穌是在實踐中,在放下自己,甘願為他人而受苦的實踐中,他領悟何謂「順從」,何謂將生命交給天父.有「順從」,而沒有承擔苦難,這「順從」是空的.耶穌這裡的「學」,是一種「知行合一」的領悟.耶穌在承擔十架中,活現了「順從」的真意.耶穌在「順從」天父時,甘心為愛而受苦.

關於基督教的信仰,我們知道得很多.但天父不要這種「知道」,這種「知道」是離地的,這種「知道」與你的生命的豐盛無關,這種「知道」也與他人的福祉無關.天父要的,是要你在實踐中磨煉你的「知道」,深化你的「知道」.在耶穌身上,我們看到,不經受十架的苦難,領悟不到「順從」的真義.

耶穌領悟了順從,那又如何?

來5:9　他既得以完全,就為凡順從他的人成了永遠得救的根源.

當甘願承受愛的十架時,耶穌的生命得以「完全」.在這裡,所謂「完全」,主要不是指一種道德上的完美,而是指一種人性生命上的圓滿實現.用哲學的術語來說,這不是一「道德學」的用字,而是一「存在論」的用字.人性生命的圓滿實現,用中國儒家的話來說,就是「盡人性」.「盡人性」者,就是「聖人」.

在這節經文裡,「順從」這個字又出現了.耶穌在背負人類苦難的過程中,領悟了順從.同樣,我們也在背負人類苦難的過程中,領悟順從.領悟了,才進入基督教所講的「救恩」裡.這是重複耶穌活著的過程.面對天父,敬畏天父,明白天父,讓天父的心抓住你.在大愛的包圍下,你甘心放下自己.在這苦難的世界,你願意為他人多走一步,背起生命的十架.在這過程中,你會像耶穌一樣,完成你人生的責任,圓滿地活出人性,領悟生命的全面意義.

在這生命裡,不再問為何有苦難,而是在苦難人生中盡力去愛.不再埋怨為何自己要受苦,而是在受苦中去安慰其他受苦的人.不再抱怨天父為何不聽祈禱,而是明白天父渴望我們在苦難的人間,透過我們承擔十架,向世人宣告平安和希望.

如此這般,耶穌成了何等樣的人物?

來5:5　如此,基督也不是自取榮耀作大祭司,乃是在乎向他說「你是我的兒子,我今日生你」的那一位；

來5:6　就如經上又有一處說:「你是照著麥基洗德的等次永遠為祭司.」………

來5:10　並蒙神照著麥基洗德的等次稱他為大祭司.

當一個人甘願被天父抓住,甘願承擔人類苦難的十架時,一個偉大的人格生命便出現於世了,他成了「大祭司」,溝通天地,站立在上帝與人類之間,將人間苦難呈現在上帝眼前,將上帝的仁愛帶進人類中間.用中國儒家的話來說,這是達致「天人合一」的聖人生命.不過,經文特別說,「基督也不是自取榮耀作大祭司,乃是在乎向他說『你是我的兒子,我今日生你』的天父」,這不是中國式的,自力的,也自負的「靠自己」去成聖.而是卑微的,順服的,由天父引導,讓天父的旨意成就在人間.

耶穌這大祭司,不是按人間的一般祭司的級數,而是按「麥基洗德」的級數.「麥基洗德」是一神秘人物.在創世記第14章,亞伯蘭(後來叫亞伯拉罕)打完一場勝仗,救回姪兒羅得後,神秘的「麥基洗德」出場,祝福亞伯蘭,也接受了亞伯蘭的奉獻.如此,「麥基洗德」作為祭司的級數,高於亞伯拉罕,高於神的子民的始祖.在希伯來書裡,「麥基洗德」代表一永恆的祭司職份,而耶穌作為人間的大祭司,是屬於這級數的,永恆的,不被替換的.

從今,當構想天人關係時,背起十架的耶穌,就總在畫面中央.

當耶穌撐起這個崩壞的世界時,在這苦難人間,我們看到天父的憐憫.

當不安籠罩大地時,讓我們像耶穌那樣祈禱,不再注視不安,而聆聽天父微聲的呼喚.當絕望瀰漫社會時,讓我們像耶穌那樣行動,不是埋怨退縮,而是承擔更大的責任.當上帝沉默隱退時,讓我們像耶穌那樣信靠,不因看不見上帝而不信,而在看不見上帝時更深的信靠.



\newpage

\section{榮耀之主謙卑服侍}
\label{sec:543}
\textbf{鄧達榮牧師 - 「榮耀之主,謙卑服侍」}
\newline
\newline
\hyperref[sec:542]{< < < PREV SERMON < < <}
~
\hyperref[sec:toc]{[返目錄]}
~
\hyperref[sec:544]{> > > NEXT SERMON > > >}


講題:「榮耀之主,謙卑服侍」　“Glorious Lord, Humbly Service”

講員:鄧達榮牧師

場合:香港中文大學崇基學院禮拜堂主日崇拜

日期:2021年3月28日

引言︰

各位親愛的弟兄姊妹,大家平安,主恩常偕.今天的主日(28/3)是棕枝主日,也就是我們教會紀念耶穌受死前最後一週的開始,我們的主耶穌騎著一匹驢駒進入耶路撒冷的時刻,群眾在耶穌前行的路上舖著自己的衣服和棕樹枝,前呼後擁地喊著說:「和散那!奉主名來的是應當稱頌的!」藉此表達出對耶穌那份尊崇和歡迎……,與此同時今天也是大齋期屬操練旅程的最後一站,緊接著我們便是進入聖週(Holy week),也就是受難週(Passion week)的日子,下一個主日我們便歡慶主耶穌的復活並所帶給我們的盼望,因此,讓我們在今天早上,藉著耶穌進入聖城的事件和腓立比書第二章內有關對耶穌生命的「讚頌」,一同認識這位「榮耀之主」,而在祂的身上又是如何作「謙卑服侍」的榜樣.

正文:

一. 耶穌這榮耀之主在群眾的歡呼中光榮地進入聖城

從馬可福音的記載,我們看到耶穌是受到極大歡呼和夾道歡迎下進入耶路撒冷,門徒在耶穌的囑咐中已預備好一隻驢駒子,並把自己的衣服搭在這驢駒上,仿如作馬鞍讓君王騎在上面,在沿途進程上,更有許多人把衣服舖在路上,還有人把田間的樹枝砍下來舖在地上,前呼後擁的群眾們都喊著說:

「和散那!奉主名來的是應當稱頌的!

那將要來的我祖大衛之國是應當稱頌的!至高無上的,和散那!」

(可11:9)

這樣的場景就仿如古代的君王在進入城裡去接受加冕,充滿光彩和榮耀.在此,首先值得我們注意的,這些人不是由政府或官方部門發動安排下出來歡迎耶穌(就仿如今天刻意安排紅地毯的歡迎儀式),而是他們自願來到這裡,發自內心的歡呼和流露他們的期盼.其次,他們所呼喊的「和散那」卻深深藏著一份「歷史」的情感因素在其中……

今天很多敬拜的詩歌內都有用「和散那」作歌詞的內容,以表達出一份讚美的呼聲,這個用語“Hosanna”往往被誤解為「讚美」的意思,但若深究下去,這個用詞不是希臘文,而是希伯來文,意思為「求來(此刻)拯救我們」,曾分別在撒下14:4和王下6:26出現,使用的處境都是百姓對他們的王呼求幫助和保護.因此,當群眾前呼後擁高喊「和散那」時,他們不單止在歌頌讚美和歡迎耶穌,而是在對他們的王呼求幫助和保護,是他們心目中上帝所差遣的「彌賽亞」,將要坐在大衛王座的君王,並要恢復以色列民族光榮的拯救者,而這份熱切期待的心情,其實可以追溯至馬加比時代百姓們對敘利亞人佔領他們土地和污辱聖殿的反抗有關,當日與敘利亞抗爭時,話說馬加比族人並沒有正規的軍服,身上只以棕樹枝識別,結果馬加比家族抗爭勝利,建立起自己的管治,因此馬加比王朝在發行的錢幣上就刻上有棕樹枝作為記念.若從這一個歷史向度去作觀察,以色列百姓舞動棕樹枝歡迎耶穌入城,就猶如歡迎期待已久的王帶領他們「搞革命」一樣,因此若參照馬太福音21:10的記載,無怪乎全城的人不論統治的羅馬官兵或是以色列民都十分被驚動並說:「這是誰?」是的,這就是「至高無上」的榮耀之主,值得我們去歡慶和擁戴祂為王,這就正如百姓所呼喊一樣:「奉主名來的是應當稱頌的!」

二. 耶穌這榮耀之主乃騎著驢駒子進入聖城

剛才我們看到耶穌進入聖城的情景,祂以「王者」的身份受到熱烈的歡呼,眾百姓在歷史的緬懷和對未來「彌賽亞」臨到的期望中,大家都拿著棕樹枝歡迎祂和呼求祂,然而,耶穌的寶座又在那裡呢?是在那榮耀的宮殿又或在宏偉的聖殿中嗎?全都不是,祂的寶座竟在那高掛的十字架上.此刻我們面對一個充滿吊詭的真理(paradoxical truth),那就是這位榮耀之主,也是上帝的兒子並要成為君王的耶穌基督,祂卻以謙卑的方式進入聖域,之後又替門徒洗腳(約13:1-20),並最終死在十架上而獲得最大的榮耀.就讓我們回到入城這一幕,在馬可的記述中,耶穌所騎的並非高頭駿馬,也不是古代君王所騎的戰馬.反之,是一匹驢駒,這正好應驗了撒迦利亞先知的預言說:「錫安哪,應當大大喜樂；耶路撒冷啊,應當歡呼.看哪,你的王來到你這裡!他是公義的,並且施行拯救,謙和地騎著驢,騎著小驢,驢的駒子.」(亞9:9)在聖經中,「馬」的形像都是「戰馬」,使我們想起的就是戰爭,征服,殺戮與死亡；反之,驢駒卻給人有平實,勸勞,溫和和服侍的形像,牠更是平民百姓日常所用的交通和運貨工具,任勞任怨,默默作出服侍……,此刻耶穌吩咐門徒去找一匹驢駒入城,正好去顯明祂帶來的並非戰爭和奪權,祂雖是君王和那真正的彌賽亞,卻是一位謙卑的榮耀之主,在地上,祂沒有政治和軍事權力,只以僕人的形像去服待,也不受到別人應有對祂的尊重,因此耶穌說:「我來了,並不是要受人的服事,乃是要服事人,並且要捨命作多人的贖價.」(可10:45),因此,讓我們再一次去細想,耶穌騎著驢駒入城的真正意義和給予我們的挑戰,就是祂本來就已經擁有一切的權柄,然而祂卻甘心放棄了,祂騎著騎駒而來,並非要在地上建立自己的國度和權力,而是要成就父神拯救世界的旨意,因著這一份的謙卑和順服,我們這位榮耀的主把那本應屬乎祂的權柄都放棄了.各位主內弟兄姊妹們,我們也能擁抱基督這份謙卑服侍的心志嗎?

三. 耶穌這榮耀之主的「虛己」,卻要成為我們謙卑服侍的榜樣和操練

各位主內弟兄姊妹,學習耶穌基督謙卑服侍的生命榜樣,是我們所有樂意跟隨耶穌(作主門徒)的終身目標, 而其中最大的挑戰,莫過於學習「放下」自己的權利,為了愛就樂意去「犧牲」自己,為了順服父神的心意就甘心去為別人「捨命」,而最著名去描述這樣服侍生命的一段偉大經文,莫過於在腓立比書第二章5-11節之內的讚頌……

「5你們當以基督耶穌的心為心:

6他本有上帝的形像,卻不堅持自己與上帝同等(的地位)；

7反倒虛己,取了奴僕的形像,成為人的樣式；既有人的樣子,

8就謙卑自己,存心順服,以至於死,且死在十字架上.

9所以上帝把他升為至高又賜給他超乎萬名之上的名,

10使一切在天上的,地上的和地底下的,因耶穌的名,眾膝都要跪下,

11眾口都要宣認:耶穌基督是主,歸榮耀給父上帝.……」

這一段又被稱為偉大的基督「虛己論」,我們又當怎樣作出正確的理解呢?在第6節首先說出基督本有上帝的形像,當中的「本有」,在希臘文中源自一個比所採用的動詞“to be”更有力的字詞,NIV的譯本將緊隨它的名詞更傳神地譯為「真本性」,亦即在成為人之前的耶穌,本身就是「真神」,擁有完全的神性(fully God),那麼「虛己」又是指甚麼,是祂倒空了自己的神性嗎?並在地上這一段的日子失去了祂的神性?非也,我們絕不可作這樣的理解,「虛己」這個字在新約中只有保羅採用過五次,其用法全是指某些東西失掉了功效,而不是指倒空某些具體的東西,換言之,這裡所指的,主耶穌不是倒空自己的神性,乃是謙卑自己,「不以自己與上帝同等為強奪的」(和修版:卻不堅持自己與上帝同等),而ASV譯本則譯為「祂使自己變得默默無聞」,因此這位充滿榮耀的主,祂擁有上帝的本性,榮耀,權柄和地位,並受天使天軍的敬拜,但道成肉身來到人世間後,主耶穌就不在這樣的榮耀中表現自己,我們可以這樣說,祂改變的是存在的方式,而非祂的本性,主耶穌這一刻暫時放棄天上存在的方式,甘心樂意地放下自己應當享有的權柄,不在榮耀中表現自己,這就是和合本「不強奪」的真正意思.

主耶穌不但「反倒虛己」,祂還「取了奴僕的形像,成為人的樣式」,祂為了完成救贖的工作,在神性上又加上人性,就是加上人性中一切的限制,噢!這榮耀的主竟成為軟弱和受盡限制的人,把不必要的擔子也加到自己的身上,而腓比立書這一段美妙的讚頌中,就是聚焦地描逑上帝的愛子耶穌基督不斷加深的「降卑」,祂不但願意離開天上的榮耀,降生在卑微的馬槽裡,更倒空自己取了奴僕的形像,成為人的樣式,經文更進一步說出祂「就謙卑自己,存心順服(即樂意放棄祂的一切權利和本有的榮耀),以至於死,且死在十字架上」,而釘十字架是昔日古羅馬時代最羞辱的死刑,專為罪大惡極又或十惡不赦的嚴重罪犯而設的酷刑.感恩我們的主竟為我們犧牲到這樣的地步,倒空自己,願意降至最深的深淵(又就是信經所提及降到陰間),經歷徹底的死亡,為的就是去拯救我們這些有罪的靈魂,重拾我們可以與上帝和好的關係.

各位主內弟兄姊妹們,我們又願否效法主耶穌這樣謙卑的生命榜樣呢?又願否以基督耶穌的心為心?雖然我們不能像主耶穌有這樣獨特「道成肉身」和「倒空」自己榮耀的歷史經驗,但我們也應有主耶穌這份謙卑和順服的心志--就是為了完成上帝所交付給我們的旨意,為了能成為上帝所合用的器皿,甘心樂意去犧牲自己應有的權利,負起艱苦的擔子,就是你和我的十字架,繼續去完成上帝所託付給我們的使命和挑戰,我們又願意嗎?

結語與生命反省:

最後,在結束今日棕枝主日信息的講道之前,讓我們有以下的總結和作出這樣生命的反省:昔日耶穌騎著驢駒以一位謙卑君王的姿態進入耶路撒冷,這份謙卑,並非只在做門面外在功夫,也不像現代的政客一樣,聘用形像顧問,為祂去度身訂造一個親民的「謙卑」形像,好贏得民眾的掌聲和選票,而我們的主耶穌確實放下祂作為榮耀之主的身段,尊貴和地位,以犧牲捨己的精神走向各各他山,完成祂在十架上救贖世人的使命,還有幾天至本週五(2/4)便是受苦節的日子,傳統稱這個星期為聖週(Holy week)又或稱為「受難週」(Passion week),是極俱莊嚴,神聖和充滿奧秘的重要教會節期,甚至應該比聖誕節期有過之而無不及.在教會傳統裡更有不少信徒學習在此受難週禁食,或在這一段期間內只吃麵包,白飯,稀粥,甚或只飲清水…….那就讓我們以簡樸,虔誠,謙卑的心懷,預備好各人的心,以摯誠的態度去迎接和記念主耶穌為我們承受的苦難,並因祂的謙卑服侍成全了圓滿救贖而發出「和散那」的歡呼和感恩,誠心所願!!



\newpage

\section{他已經復活了不在這裡}
\label{sec:544}
\textbf{高國雄牧師 - 他已經復活了,不在這裡.}
\newline
\newline
\hyperref[sec:543]{< < < PREV SERMON < < <}
~
\hyperref[sec:toc]{[返目錄]}
~
\hyperref[sec:545]{> > > NEXT SERMON > > >}


馬可福音 16:1-8
\newline
\begin{longtable}{cl}
\hline
\hline
章節 & 經文 (和合本修訂版)\\
\hline
16:1 & \begin{tabularx}{0.7\textwidth}{X} 過了安息日，抹大拉的馬利亞、雅各的母親馬利亞，和撒羅米，買了香料，要去膏耶穌的身體。 \end{tabularx} \\ \\ \relax
16:2 & \begin{tabularx}{0.7\textwidth}{X} 七日的第一日清早，太陽出來後，她們來到墳墓那裡， \end{tabularx} \\ \\ \relax
16:3 & \begin{tabularx}{0.7\textwidth}{X} 彼此說：「誰要替我們把石頭從墓門滾開呢？」 \end{tabularx} \\ \\ \relax
16:4 & \begin{tabularx}{0.7\textwidth}{X} 她們抬頭一看，看見石頭已經滾開了，原來那石頭很大。 \end{tabularx} \\ \\ \relax
16:5 & \begin{tabularx}{0.7\textwidth}{X} 她們進了墳墓，看見一個年輕人坐在右邊，穿著白袍，就很驚奇。 \end{tabularx} \\ \\ \relax
16:6 & \begin{tabularx}{0.7\textwidth}{X} 那年輕人對她們說：「不要驚慌！你們尋找那釘十字架的拿撒勒人耶穌，他已經復活了，不在這裡。來看安放他的地方。 \end{tabularx} \\ \\ \relax
16:7 & \begin{tabularx}{0.7\textwidth}{X} 你們去，對他的門徒和彼得說：『他要比你們先到加利利去，在那裡你們會看見他，正如他從前所告訴你們的。』」 \end{tabularx} \\ \\ \relax
16:8 & \begin{tabularx}{0.7\textwidth}{X} 於是她們出來，從墳墓那裡逃走，又發抖又驚訝，甚麼也沒有告訴人，因為她們害怕。 \end{tabularx} \\ \\
[1ex]
\hline
\hline
\end{longtable}
講題:他已經復活了,不在這裡.　He has been raised; He is not here.

講員:高國雄牧師

經文:馬可福音16:1-8

場合:香港中文大學崇基學院禮拜堂主日崇拜

日期:2021年4月4日

哈利路亞,主已經復活了!

今天是基督復活日崇拜,復活節是基督信仰中最重要的節日.

哥林多前書15:14「若基督沒有復活,我們所傳的便是枉然,你們所信的也是枉然.」

對,使徒保羅的提醒非常重要,我們信仰的核心就是耶穌基督從死人中復活,勝過死亡,更為信靠基督的人帶來復活的盼望,生命的力量.

因此在今天復活節的主日中,我們要再一次重述基督復活這一個最重要的神蹟,反思此神蹟在今天對我們的意義.

馬可福音16:1-8

1過了安息日,抹大拉的馬利亞和雅各的母親馬利亞並撒羅米,買了香膏要去膏耶穌的身體. 2七日的第一日清早,出太陽的時候,她們來到墳墓那裡, 3彼此說:「誰給我們把石頭從墓門滾開呢?」 4那石頭原來很大,她們抬頭一看,卻見石頭已經滾開了. 5她們進了墳墓,看見一個少年人坐在右邊,穿著白袍,就甚驚恐. 6那少年人對她們說:「不要驚恐!你們尋找那釘十字架的拿撒勒人耶穌,他已經復活了,不在這裡.請看安放他的地方. 7你們可以去告訴他的門徒和彼得,說:『他在你們以先往加利利去.在那裡你們要見他,正如他從前所告訴你們的.』」 8她們就出來,從墳墓那裡逃跑,又發抖又驚奇,甚麼也不告訴人,因為她們害怕.

福音書經課一開始提到時間是在安息日後的一天早上,首先希望與大家解釋一下,有一點很重要請留意的,有部份基督新教的教會主張必須在安息日時舉行崇拜,因為這是從舊約時代就開始,上帝在十誡中清楚表明,要求人要守安息日；請大家先認識清楚,我們今天是沒有守「安息日」的一群,除了因為耶穌基督在世時討論到有關安息日的教導時,常常提醒:「人子才是安息日的主」,所以守安息日不是要守由人所定的規條,而是要記念上帝的拯救,與及要我們休息,並有責任協助解放我們身邊仍然生活在捆鎖中的人,因此我們守安息日的方法就改變為以敬拜主,服侍人為焦點.另外就是因為基督在主日(禮拜日)復活,因此我們慶祝更重要的基督復活日(主日),在主日進行崇拜及休息.

基督復活在禮拜日的證明就是在今天這段經文所提到,安息日之後的一天,由於耶穌基督在週五的受難日(Good Friday)中被釘十字架,猶太人的傳統週五黃昏就是安息日的開始(他們一天的計算方法是由黃昏開始到另一天的黃昏,與中國人傳統日出而作,日入而息有別),所以這是為什麼不希望罪犯被釘十字架超過週五的黃昏時間(約翰福音19:31 猶太人因這日是預備日,又因那安息日是個大日,就求彼拉多叫人打斷他們的腿,把他們拿去,免得屍首當安息日留在十字架上.)

要清楚認識耶穌基督死後在第三天復活,這「第三天」的計算方法是以猶太人傳統對日子的計法,耶穌基督被釘,斷氣時候(約下午三時)的一天是安息日前的預備日(禮拜五),這是「第一天」,安息日就是「第二天」,由於按傳統猶太人在安息日中很多事務也不可作,所以耶穌基督在預備日中死去,離開安息日開始的時間實在不多,為耶穌所需要做的安葬事宜相信仍然大部份未完成,但在安息日中不能繼續,因此婦女們就要在安息日的後一天(即週六黃昏開始的這一天)才可以繼續為耶穌處理後事,當然以當時的情況一定要在有陽光的白天才可能工作,因此就能夠解釋,為何婦女們要在主日出太陽的時候(原文新約:在太陽升起非常早時)就到墳墓.而又同時可以理解到信經提到耶穌基督「第三天從死人中復活」的意思,所以要分清,耶穌基督不是在死了,三天 (72小時) 之後復活,因為如果按我們今天對時間的理解,耶穌基督由死亡到復活應該大約40小時之內.

經文中幾位愛主的婦女們在主日大清早帶著香料來到耶穌基督的墳墓前,為要希望盡最後的份(處理耶穌的屍體),相信她們帶著極度的悲傷,不捨,不明白,甚至憤怒(請不要忘記,當日耶穌基督是以革命罪/叛國罪被冤枉而處極刑的),到了墳墓時又非常擔心實際的問題(當時墓穴的設計是由大石頭擋住為墳墓之門,以幾位婦女之力應該未能把它挪開,所以V3彼此說:「誰給我們把石頭從墓門滾開呢?」),這百感交雜的狀況,我們又有沒有經歷過呢?

V5提到婦女們的心情「甚驚恐」,和合本修訂版用「很驚奇」,希臘文原文有「驚奇」,「驚訝」,可以附帶著正面或負面的反應,即是所指婦女們未必一定是負面的害怕/驚恐,而是可能對這突如其來的事情不知所措(石頭已經被滾開,墳墓內有穿白長袍青少年住著),在我們的信仰生活中又有沒有不知所措的時候呢?並且最驚恐,害怕的又會是什麼呢?

上帝的使者安慰婦女們:

V6那少年人對她們說:「不要驚恐!你們尋找那釘十字架的拿撒勒人耶穌,他已經復活了,不在這裡.請看安放他的地方. 7你們可以去告訴他的門徒和彼得,說:『他在你們以先往加利利去.在那裡你們要見他,正如他從前所告訴你們的.』」

並宣告:「他已經復活了!」

聖經中從來沒有提及這位穿白袍的青少年到底是誰,當然很自然地會被理解為天使,雖然我們無法確定他的身份,但可以很清楚的就是,他代表上帝向婦女們宣告及邀請,同時指示出耶穌已經不再在本來被安放在的位置上.

婦女們原意要到墳墓去尋找那位已然被釘在十字架上,並死了的耶穌的身體,原本帶著悲傷而來,為要盡愛/所愛的責任,我們對待信仰有時也是持有這種態度,目標或心態,「宗教皆導人向善」,「聖經/耶穌的道理很有意思」,「教會/信徒朋友曾經幫我很多,我也應該要回饋」；但上帝的宣告與邀請完全改變了婦女們原有的計劃,想法及行動,因為事實上基督真的復活了,一切不再一樣,人故有的想法,態度要改變,行為也要改變,婦女們不單不需要擔憂大石頭會搬不動,更無法/需要用香料去膏耶穌的屍體,當我們原定的一切計劃也沒法/不能實行時,我們的心情又會如何呢?

V8她們就出來,從墳墓那裡逃跑,又發抖又驚奇,甚麼也不告訴人,因為她們害怕.

婦女們的反應是發抖,驚奇(原文意思是驚駭,極度驚訝,已到達失去自制力的狀況),行動是逃跑,因為害怕(懼怕)而不對任何人說任何事；極度的震撼,難以接受,古怪難明,往往人最難接受的就是按自己邏輯無法想通,理解的事,其實背後的原因就是因為我們自以為是,自以為聰明,一切必須掌握於自己手中才是最穩妥,基督教信仰的核心就正正是要我們認識自己/人真正的位置,在上帝面前我們是如何的微小及無能,人類以為可以靠自己的能力去完全理解上帝,就是我們最愚昧的地方,所以我們要認罪,悔改,接受上帝的拯救.

V8這節經文在大部份新約學者眼中就是馬可福音最原本的結尾,現存的第9節(和合本修訂版的短結語)及9-20節(和合本修訂版的長結語及和合本中的結尾)都被視為是後期編輯加上的；因為一方面在最可靠的古抄本中沒有V9以後的內容,加上結語與前文的詞彙及風格的表達不太一致,在不同的抄本/譯本中出現不同的結尾,很可能是早期讀者覺得如果經文停在V8就完結是很奇怪及不完整,我沒有打算在此深入討論這個不同結尾所延伸相關的問題,如果我們嘗試接受如若馬可福音的作者就把經文完在16:8,那又代表著期望我們作出什麼思考呢?

上世紀八十年代開始流行的讀者反應批評理論(reader-response criticism),其中一個重點是閱讀者對文本的主動參與,意義並非存在於文本本身,而在讀者的心靈中,相信對我們作為信徒,此理論閱讀聖經亦有所幫助；如若用此角度去理解馬可福音這結尾,難道那不就是作者期望我們每一個閱讀的人要去回應,這回應不單止是要思考/想像到底在這一群非常害怕的婦女們及門徒身上將發生了什麼事?更要我們同樣因為上帝竟然愛我們到底,耶穌基督真的被釘在十架上,有人宣告基督已經復活,並且叫我們去等待與祂會面(但實際上我們仍然未見到,正如當日的婦女們一樣),在震驚中的我們,又會以什麼行動去回應呢?

馬可福音1:1 上帝的兒子,耶穌基督福音的起頭.

以教會的發展來說,福音書的寫成只是基督信仰被開展的起頭,歷世歷代信徒的回應,行動才讓福音不斷地的拓展,今天同樣,復活主日所展示的信仰,不是單單在重述二千多年前所發生的舊故事,一個已經完結的故事,而是要我們參與在其中,一同經驗生命中的震撼,改變,是一個由你我開始,延續並展現復活的大能,改變世界的故事.

延伸反思題目:

哥林多前書15:3-11

3我當日所領受又傳給你們的:第一,就是基督照聖經所說,為我們的罪死了, 4而且埋葬了；又照聖經所說,第三天復活了, 5並且顯給磯法看,然後顯給十二使徒看； 6後來一時顯給五百多弟兄看,其中一大半到如今還在,卻也有已經睡了的. 7以後顯給雅各看,再顯給眾使徒看, 8末了也顯給我看；我如同未到產期而生的人一般. 9我原是使徒中最小的,不配稱為使徒,因為我從前逼迫神的教會. 10然而,我今日成了何等人,是蒙神的恩才成的,並且他所賜我的恩不是徒然的.我比眾使徒格外勞苦；這原不是我,乃是神的恩與我同在. 11不拘是我,是眾使徒,我們如此傳,你們也如此信了.

我們在今天如何經歷復活的主呢?

這福音是如何傳給我的呢?我又如何傳出去呢?

使徒行傳10: 40-42

40第三日,神叫他復活,顯現出來； 41不是顯現給眾人看,乃是顯現給神預先所揀選為他作見證的人看,就是我們這些在他從死裡復活以後和他同吃同喝的人. 42他吩咐我們傳道給眾人,證明他是神所立定的,要作審判活人,死人的主.

誰人能夠看到復活的主呢?

為神作見證的人

與復活主同吃喝的人

傳道給眾人的人



\newpage

\section{我主我神一位誠實存疑者的認信}
\label{sec:545}
\textbf{伍渭文牧師 - 我主我神--一位誠實存疑者的認信}
\newline
\newline
\hyperref[sec:544]{< < < PREV SERMON < < <}
~
\hyperref[sec:toc]{[返目錄]}
~
\hyperref[sec:546]{> > > NEXT SERMON > > >}


講題:我主我神--一位誠實存疑者的認信　My Lord and My God - The Confession of An Honest Doubter

講員:伍渭文牧師

場合:香港中文大學崇基學院禮拜堂主日崇拜

日期:2021年4月11日

引言:

人的記憶是有選擇性的,我們記起人所作的壞事多過其好事.莎士比亞說過:「人做的惡事,比他的命長；人做的好事,跟骨頭長埋地下.」(The evil that men do lives after them; the good is oft interred with their bones.)所以我們記得:雅各騙取以掃長子名份；大衛殺人奪妻；彼得三次不認主；多馬多疑(Doubting Thomas).就像中國俗語所說,「好事不出門,醜事傳千里.」但上述例子中,多馬的多疑,並不是一樁惡行,不同殺人奪妻,欺騙,否認主.為何說多馬多疑呢?這跟耶穌復活後對門徒的顯現有關.

四福音有關耶穌復活後對門徒顯現的時序如下:抹大拉的馬利亞,撒羅米,雅各的母親馬利亞一清早帶著香膏來墓園,想膏抹耶穌屍首,發現墓石滾開了,裹屍布和頭巾像人形還放在石棺上.抹大拉馬利亞以為有人偷屍,趕緊奔告彼得和約翰.約翰跑得快,彼得隨後,進入空墳看見了就信,但不完全明白聖經的意思.(約二十8,9)馬利亞離開後,幾位婦女在墓內見到天使,極其害怕.天使囑咐她們:「快去告訴他的門徒,說他從死裡復活了,並且在你們以先往加利利去,在那裡你們要見他.」(太二十八7)

1)抹大拉後來獨自回到墓園,耶穌向她首先顯現.(約二十11-17)

2)婦女們急忙離開墳墓奔走告訴門徒,在途中耶穌對她們顯現:「願你們平安!」她們就上前抱住他的腳拜他.耶穌對他們說:「不要害怕!你們去告訴我的弟兄,叫他們往加利利去,在那裡必見我.」(太二十八9, 10)婦女們告訴門徒,他們不信.

3)以馬忤斯路上向革流巴和另一門徒顯現.他們回到耶路撒冷報告,門徒也不信.

4)耶穌向十一個門徒顯現,適時多馬不在場.耶穌責備使徒們不信,心裡剛硬.「因為他們不信那些在他復活以後看見他的人」.(可十六14)

5)隔了八天,耶穌第二次向門徒顯現,特別叫多馬伸手摸他:「不要疑惑,總要信.」多馬沒有伸手摸耶穌就說:「我的主!我的神!」(約二十27, 28)

6)第三次耶穌向門徒顯現是在加利利海.他們整晚打魚一無所獲,耶穌叫他們把網撒在右邊,網拉上來153條大魚而網沒有裂開.耶穌跟門徒共進早餐,三次問彼得:你愛我比這些更多嗎?彼得答:「主阿,你是無所不知的；你知道我愛你.」-耶穌說:「你餵養我的羊.」(約二十一15-19)

1多馬是認真的存疑者.

「那十二個門徒中,有稱為低土馬的多馬；耶穌來的時候,他沒有和他們同在.那些門徒就對他說:『我們已經看見主了.』多馬卻說:『我非看見他手上的釘痕,用指頭探入那釘痕,又用手探入他的肋旁,我總不信.』」(約二十24, 25)

1)多疑是多馬的性格.

他喜歡思考,要弄清楚信仰的原委,喜歡問問題.耶穌說他要往父那裡去,然後聖靈才來.「我往那裡去,你們知道；那條路你們也知道.」多馬對他說:「主啊,我們不知道你往那裡去,怎麼知道那條路呢?」(約十四4, 5)也許門徒同樣納悶,但他們不作聲.多馬不同,要問個究竟,他尋求合乎理性的答案.

多馬性格悲觀.猶太地區殺機四伏,為要使拉撒路復活,耶穌要赴接近耶路撒冷的伯大尼.門徒說:「拉比,猶太人近來要拿石頭打你,你還往那裡去嗎?」(約十一8)門徒叫耶穌避開死亡.但多馬較悲觀,若不能逃避死亡,願意以死相隨:「我們也去和他同死吧.」(約十一16)像叔本華(Schopenhauer)和祁克果(Kierkegaard),悲觀的人會多方思考問題.當敵人節節敗退,彼得會說:要乘勝追擊,凱旋在望了.但多馬會說:小心可能這是圈套；深入敵陣,有足夠後勤支援嗎?多馬性向存疑,但對信仰認真,祇要撥開雲霧,就義無反顧,付上生命在所不計.

2)存疑是思考.

魯益師(C.S. Lewis)地獄來鴻(Screwtape Letters)幽默吸引,藉反話說真理.內容說魔鬼給門生教路,如何誘惑信徒.第一封信魔鬼指點門生:用辯論折服信徒奪其靈魂,效果不彰而且危險；他們會愈來愈明白信仰,信得更堅定.叫信徒專注即時的官感經驗就可以了,很湊效.存疑者要問過究竟,其實是認真.1930年出版的《墓石懸迷》(Who Moved the Stone?),作者法蘭.莫理遜(Frank Morrison)是一位著名律師和專欄作家,他懷疑耶穌復活,收集受死和復活前後七天的證據,模擬在法庭辯論,希望證明耶穌沒有復活.但所有證據都指向耶穌真的復活,特別是空了的墳墓；這書最後成為護教名著.

耶穌跟撒馬利亞婦人談到崇拜時說:「你們所拜的,你們不知道；我們所拜的,我們知道,因為救恩是從猶太人出來的.」(約四22)婦人在撒馬利亞的聖殿敬拜,是人的傳統；信仰要經個人思考,不能人云亦云.

3)有思考始有根基.

1981年暑假,我在美國春田城成立了華人的崇拜聚會.開始時,我們到唐餐館吃飯,到移民學英語的學校探訪,拿起電話簿看似中國人的名字,就打電話接觸.認識一位來自香港的姊妹,丈夫是腦科醫生.她坦承到了美國就很少去教會,說是對信仰的反叛,rebound(原話).她在香港時是熱心的乖乖女.教會非常保守,看電影是愛世界,化裝是扮靚,利是錢是不勞而獲.信仰跟社會文化是割裂的.存疑幫助我們思考,沒有思考信仰不能生根,就像落在淺土和石頭上的種子,飛鳥一來就吃掉了.

4)存疑是認真.

戀愛到一個臨界點,要下決心結婚時,通常會有存疑:這真的是我終身伴侶嗎?沒有其它的選擇嗎?魯益師說他信主時,是英國最低沉和不起勁的歸依者(The most dejected and reluctant convert in all England).一九二九年經過多年上下求索,多方閱讀,師友切磋,他來到要作抉擇的一刻,知道門一推開,一跨進去,將面對一個不可估量的世界.在作信心的跳躍前,他存疑猶豫,不是不認真,乃非常認真.

5)耶穌欣賞多馬對信仰認真.

耶穌第一次向十一使徒顯現時,多馬不在場.為甚麼?極可能多馬悲傷哀慟,想一個人獨處,顯然門徒了解他,沒有怪他缺席,錯過機會,多馬回來時祇對他說:「我們已經看見主了.」多馬盼望耶穌對他顯現,像對其他門徒一樣,而且有更親蜜的接觸:「我非看見他手上的釘痕,用指頭探入他的肋旁,我總不信.」(約二十25)

過了八天耶穌第二次對門徒顯現,多馬也在其中.對多馬來說,過去八天真不好受,他應聽到不少門徒遇見耶穌的見證,但一天一天過去了,耶穌還未對自己顯現,他一定很焦慮.會不會自己開價過高,對耶穌不敬呢?喜歡思考又悲觀的多馬,也許認為耶穌拒絕他不理他了.耶穌出現了!「願你們平安!……(多馬)伸過過你的指頭來,摸我的手；伸出你的手來,探入我的肋旁.不要疑惑,總要相信!」(約二十26, 27)耶穌並沒有嚴責多馬,反而溫柔叫他不要疑惑,總要相信；摸我吧,探我吧. 此刻,多馬被耶穌的接納和愛所感動,不須要觸摸耶穌就敬拜他:「我的主!我的神!」.這是門徒對基督最高的認信:耶穌與父同尊同榮.

6)哀歌作者也是存疑者.

多馬願意相信,但他要清楚地相信.他站在信仰的立場,希望驅走疑惑,全然相信.像聖經哀歌的禱告者,把存疑帶到上主面前,對上主傾訴(Talk to God),不是談論上主(Talk at God) :「我的神!我的神!為甚麼離棄我?」(詩二十二).在哀歌,禱告的人求告上主,增強其信心,撥開雲霧,好作信心的跳躍(leap of faith).因而,耶穌向多馬顯現並叫他摸自己.

應用:在信仰和生活上,我們總有疑惑.誠實的懷疑者為上主所喜悅,因他們對信仰認真.若我們作好信心的跳躍,上主也會親近我們,顯為真實.「不要疑惑,總要相信.」

2 沒有看見就信的有福了

「你因看見了我才信,那沒有看見就信的有福了.」(約二十29)

7)看見的時代將過去.

這句話看似責備多馬,不應要求眼見才相信,但耶穌對多馬的接納和溫柔,應另有所指:有福的時代快要來到,不再憑眼見,乃憑信心.耶穌降生,受死,埋葬,復活,升天,我們看見；但升天之後直到第二次再來,就不能憑眼見,乃憑信心.信徒要透過聖靈,聖經,和聖徒,與基督相遇.這日子是有福的,因聖靈已降臨,教會已誕生.耶穌經歷十架的苦痛,就像婦人忍受生產的苦難,迎來新生命-教會的誕生.

8)不要拉住我.

當抹大拉馬利亞聽到耶穌叫她:「馬利亞」,她立刻回應「拉波尼」(亞蘭文拉比,老師)並伸手拉住耶穌.耶穌說:「不要摸我(拉住我,新譯本；Do not hold on to me, NIV；Don’t cling to me, Message),因我還沒有升上去見我的父.」(約二十16,17)耶穌豈不曾說:「等不多時,你們就不得見我；再等不多時,你們還要見我.」(約十六16)馬利亞以為耶穌離開又回來了,可以回到過去美好的時光,朝夕相處,陪伴左右.耶穌提醒她,新時代已來了,我不再是以前的拉波尼,我是萬有的主.這時候的耶穌,可以隱藏在同路人的身上,跟以馬忤斯的門徒講解聖經；也可以突然進入關閉的房子.「不要拉住我」指出世界已不一樣了.我們已進入不再憑眼見,乃憑信心的新紀元.耶穌要上升見父,直到他復臨,我們才能見到他.

耶穌公開出來傳道三年多,見過他面作他門徒的,竟然不相信他復活,墳墓空了屍首不見了,以為被偷走,傷心欲絕.兩批見證人:婦女們,及革流巴等遇見主的見證,「使徒以為是胡言,就不相信.」(路二十四11)耶穌責備門徒不接受他們的見証.「耶穌責備他們不信,心裡剛硬,因為他們不信那些在他復活以後看見他的人.」(可十六14)

事實上,憑信心的我們,比眼見到耶穌的第一代信徒,知道更多有關耶穌的救贖事件.使徒們雖然眼見過耶穌,但他們不明白也不相信耶穌復活,而且在耶穌四十天教導天國奧秘後,到升天的一刻還問:你復興以色列就在今天嗎?我們比他們更了解信仰和教義.憑信心的信徒,擁有完整的聖經,也累積了歷代信徒了解聖經的智慧,幫助我們明白和應用聖經.

9)向群體顯現.

為甚麼耶穌不直接了當向門徒顯現,而派兩批見證人跟門徒作見證呢?因為新時代來到了,不單使徒們代表教會,婦女們,革流巴等也代表教會.信徒如何接受,也要如何給予.這也說明為何耶穌要等到八天後,當多馬回到門徒當中時,才向他顯現.事實上福音書記載耶穌的顯現,除了單獨對抹大拉的馬利亞和彼得(路二十四34),都是兩個人或以上的群體.

為何要多馬等候八天之久?耶穌在其他門徒一起時才對多馬顯現,肯定多馬,不加苛責,要門徒接納並欣賞多馬存疑的積極面,同時也使用多馬信心的跳躍和認信,堅固門徒,這是聖徒的互勉(mutual consolation ).路德認為這正是教會的記號,教會是信徒的結集,在其中互相勸慰,彼此支持.耶穌升天後,沒有無縫接上聖靈降臨,還等十天,使用這空窗期,讓門徒一起懇切禱告,彼此勸慰,同心合意,等候聖靈的澆灌,教會的誕生.今天的經課,都以教會群體生活為主題.

10)今日經課都指向教會群體.

使四32-35; 信的人都是一心一意的,凡物公用.詩一三三；看哪,弟兄和睦同居是何等的善,何等的美.約一1-2: 2；信徒彼此相交；聖靈降臨前連續讀出約壹,主題是彼此相愛.約二十19-31；耶穌向信徒群體顯現.

經課使徒行傳展示門徒彼此相愛,凡物公用,充滿喜樂.百姓又希奇,驚訝,懼怕,但同時喜愛他們.(使二43；三10)封閉的世界被打開了,相比之下,帝國的羅馬和平(Pax Romana)春秋大夢,變得蒼白,軟弱無力.耶穌的跟隨者,憑信心看到,比眼睛所看到的,還要真實的世界.聖靈像風,從耶穌的墳墓,滾開了墓石,帶走死亡和絕望,並吹襲更新這世界.巴爾特(Karl Barth)說得好:「復活的信息告訴我們,我們的仇敵,罪惡,詛咒和死亡已被打敗.它們已徹徹底底沒有作為.它們也許詐作未知鹿死誰手,勝負尚未分曉,我們也不應輕敵,但已不再懼怕它們了.」在十架耶穌已一役竟其功,現在祇是清理殘餘,世界從此已不一樣.

11)我愛我主教會

香港教會這幾年因不同的政見,造成嚴重撕裂.有一次我在香港一教會講道,一位中年弟兄全場站立,不是他尊敬講員,乃表示抗議.主任牧師早跟我打了招呼,他不是抗議講員,也不是抗議教會,乃抗議當權者使用不當武力鎮壓示威者.因著不同的政見,有些信徒對教會失望,離開自己的教會,甚至不去教會.

門徒中間也有張力,因為性格和背景不同:彼得樂觀,有氣魄但衝動,多馬悲觀,謹慎多疑；馬太原是稅吏有建制背景,奮銳党西門以前推崇激進革命.他們一起被召為門徒,和睦同居,何等地善,何等地美.(詩一三三1)沒有撕裂,不須割席出逃.

宗教改革時期有兩位曠世奇材,伊斯拉姆(Erasmus)和路德(Martin Luther).伊斯拉母編輯了第一版新約希臘文聖經,歐洲王侯將相都向他請益,是備受敬重的國師.馬丁路德是一位修士和大學神學教授.兩人皆精準看出那樹枝正在腐爛敗落,但路德拿起斧頭砍掉它.路德批伊斯拉姆,應該為教會哭泣時,他竟然嘲弄教會.路德是教會忠誠的叛徒(Loyal rebel),他為教會哭泣,並改革教會,教導真理,牧養上主交托給他的羊群.

為教會哭泣使我想起《我愛我主教會》這聖詩:「我愛我主教會,主靈與人同居,救贖之主,親流寶血,建設清潔基礎,愛她我常流淚,念她我常祈求,為她事工,我更努力,直到力盡方休……耶穌為她元首,為她摯愛良友.」

總結:多馬是誠實的存疑者,但耶穌欣賞他對真理的認真,向他顯現:你可以前來摸我的釘痕,探我的肋旁,溫柔對他說:不要疑惑,祇要相信.耶穌同時趁機會帶出重要的信息:「你因看見了我才信,那沒有看見就信的有福了.」我們要進入不再憑眼見,要憑信心的新世界,憑信心,我們將看到眼睛看不見的奇妙世界.

馬利亞,不要拉住我叫我拉波尼,舊事已過,一切都變成新的了.我要上去見我的父,聖靈將降臨,教會快誕生；你們是我的見證人,是我的身體,要將福音傳遍世界；同時,餵養我的羊,建立我的教會,等候我再復臨大地.因此,講章還未結束,還有一位誠實的存疑者,稱為磯法(磐石)的西門彼得.多馬對耶穌復活存疑,彼得對自己存疑,三次否認主後,他對自己沒有信心了.

3另一位誠實的存疑者-彼得

約二十一3:「西門彼得對他們說:我打魚去.」「他們說:『我們也和你同去.』他們就出去,上了船；那一夜並沒有打著甚麼.天亮的時候,耶穌站在岸上,門徒卻不知道是耶穌.耶穌就對他們說:『小子!你們有吃的沒有?』他們回答說:沒有.耶穌說,你們把網撒在船的右邊,就必得著.」(約二十一5, 6)

12)為何彼得要去打魚?

耶穌已單獨向彼得顯現(路二十四34),耶穌對門徒顯現時,他也應該在其中.耶穌要門徒等候,聖靈快要降臨,他不相信嗎?為何要重操故業呢?耶穌問彼得三次「你愛我比這些(魚)更多嗎?」,很明顯針對他三次不認主,這失敗非常困擾彼得.他不是拒絕使徒之首的召命,但他失去自信,他懷疑自己沒有能力承擔責任,要藉打魚恢復信心.因打魚是他的本行,最有經驗,加利利海是他的地頭,沒有失過手.彼得想藉職業治療( Occupational therapy),重拾失去的自信.他沒有叫其他門徒跟他一起去打魚,所以不是重操故業,放棄原來的召命.

13)陪同支援.

門徒為何不勸阻彼得,反而一起和應他呢?因他們了解彼得,猶如了解多馬,他們陪同彼得,給予支持.互勉有時,陪同有時.肢體見面有時不須多言,陪同就足夠了.都是作家的兩父子,因為一些誤會結了心病,很久沒有見面,兒子自己的孩子也長大了,大家覺得這樣的關系不好,做兒子鼓起勇氣,約爸爸到兒時一起垂釣的地方,重遊舊地.整個下午父子坐得很近,說話很少.回家後,兩人寫下當天的日記.兒子覺得很失敗,父親跟自己沒有說過幾句話,有點自責.但同一天,父親寫下:今天樂透了,兒子陪了我整個下午,很滿足.父親需要也是最珍貴的-兒子陪同自己.

我1980年前往美國進修前一年,跟我一起在教會長大的弟兄突然感到上主離開他,不理他,他懷疑信仰.專業考試失敗,太太跟母親不和；單親媽媽覺得有人奪去兒子,加上媽媽和自己的身體都不好,這些壓力加在一起,就成為不能承受生命的重.他雖然不來崇拜,作為在教會一起成長的好友兼傳道,我每星期一次到他家附近跟他一起吃早餐,陪同他,保持聯系.若有伉儷團契聚會,邀請他夫婦參加.一年後我去了美國,後來傳來喜訊,他恢復信心,跟太太參加崇拜聚會了,並把以前的什一捐獻,全數還給上主；之後熱心在教會服事,成為信徒敬重的屬靈領袖.

14) 總結

多馬是誠實的懷疑者,耶穌接納多馬,因多馬對信仰極其認真,祇要雲霧撥開,他將義無反顧,作信心的跳躍.他認信:我的主!我的神!.彼得也是誠實的懷疑者,他懷疑自己.因為曾經三次不認主,他需要工作治療,恢復自信,打魚去了.耶穌了解彼得,吃完早飯後對彼得說:「你愛我比這些更深嗎?」彼得因為耶穌第三次對他說「你愛我嗎?」就憂愁回答:「主啊,你是無所不知的；你知道我愛你.」耶穌說:「你餵養我的羊.」(約二十一15-17)耶穌提醒彼得,是耶穌的羊,耶穌會負責,他會賜能力.憑自己,整晚一無所獲,但聽耶穌在右邊撒網,就滿載而歸.彼得明白了,他認信:「你知道我愛你,我願意牧養你的羊.」

今天,你是否對上主存疑,為何橫逆突然臨到你?你不明白,覺得上主離你很遠.耶穌說:不要疑惑,總要相信.今天,你是否對自己存疑,因為多次的失敗,覺得自己離開上主很遠,不配事奉.耶穌說:「你跟從我吧!你要牧養我的羊.」上主沒有放棄多馬,也沒有放棄彼得,也沒有放棄你.



\newpage

\section{你們看我的手我的腳}
\label{sec:546}
\textbf{鄧瑞強博士 - 「你們看我的手,我的腳」}
\newline
\newline
\hyperref[sec:545]{< < < PREV SERMON < < <}
~
\hyperref[sec:toc]{[返目錄]}
~
\hyperref[sec:547]{> > > NEXT SERMON > > >}


路加福音 24:36-43
\newline
\begin{longtable}{cl}
\hline
\hline
章節 & 經文 (和合本修訂版)\\
\hline
24:36 & \begin{tabularx}{0.7\textwidth}{X} 正說這些話的時候，耶穌親自站在他們當中，說：「願你們平安！」 \end{tabularx} \\ \\ \relax
24:37 & \begin{tabularx}{0.7\textwidth}{X} 他們卻驚慌害怕，以為所看見的是魂。 \end{tabularx} \\ \\ \relax
24:38 & \begin{tabularx}{0.7\textwidth}{X} 耶穌對他們說：「你們為甚麼驚恐不安？為甚麼心裡起疑惑呢？ \end{tabularx} \\ \\ \relax
24:39 & \begin{tabularx}{0.7\textwidth}{X} 你們看我的手和我的腳，就知道實在是我了。摸摸我，看，因為魂無骨無肉，你們看，我是有的。」 \end{tabularx} \\ \\ \relax
24:40 & \begin{tabularx}{0.7\textwidth}{X} 說了這話，他就把手和腳給他們看。 \end{tabularx} \\ \\ \relax
24:41 & \begin{tabularx}{0.7\textwidth}{X} 他們還在又驚又喜、不敢相信的時候，耶穌對他們說：「你們這裡有甚麼吃的沒有？」 \end{tabularx} \\ \\ \relax
24:42 & \begin{tabularx}{0.7\textwidth}{X} 他們給了他一片烤魚， \end{tabularx} \\ \\ \relax
24:43 & \begin{tabularx}{0.7\textwidth}{X} 他接過來，在他們面前吃了。 \end{tabularx} \\ \\
[1ex]
\hline
\hline
\end{longtable}
講題:「你們看我的手,我的腳」　See My Hands and My Feet

經文:路加福音24:36-43

講員:鄧瑞強博士

場合:香港中文大學崇基學院禮拜堂主日崇拜

日期:2021年4月18日

各位弟兄姊妹,早安.在這疫症裡,祝願各人平安.

我們都相信復活.但,復活之後,會有什麼樣的外貌的呢?我們每個人,都有不同想像.頭髮不多的人,可能想像,復活後的自己,一頭蓬勃的黑髮.身體瘦弱的人,可能想像,復活後的自己,有體操運動員的身型.至於我,其貌不揚,我就想像,復活後,最好似高牧師,氣宇軒昂.我們很難想像,復活後,禿頭的,仍然禿頭；殘缺的,仍然殘缺；其貌不揚的,仍然其貌不揚.曾經,在以色列特拉維夫的海邊,我遇到一個老伯伯,他在看海,我在看他,他手臂上刻著一列數字,很明顯,這是二戰時德國納粹刻在他身上的標記,他入過集中營,他有過不尋常的苦難經驗.在天堂,我想,只是我想,復活的他,最想移除的,就是手上這個標記,這個苦難標記.不過,又或者,他最不想移除的,就是這個標記.他的苦難,刻鑄著他的記憶,琢磨出他的一生.他這記號,包藏著他的存在最神聖的奧祕.或許,這是他生命裡向世人展示的最寶貴的歷史痕跡.

耶穌呢,當耶穌復活時,會是什麼樣的模樣的呢?

講道經文:路加福音24:36－43

路24:36　正說這話的時候,耶穌親自站在他們當中,說:「願你們平安!」

路24:37　他們卻驚慌害怕,以為所看見的是魂.

路24:38　耶穌說:「你們為什麼愁煩?為什麼心裡起疑念呢?

路24:39　你們看我的手,我的腳,就知道實在是我了.摸我看看!魂無骨無肉,你們看,我是有的.」

路24:40　說了這話,就把手和腳給他們看.

路24:41　他們正喜得不敢信,並且希奇；耶穌就說:「你們這裡有什麼吃的沒有?」

路24:42　他們便給他一片燒魚.

路24:43　他接過來,在他們面前吃了.

這就是復活的耶穌的顯現.他對門徒說:「你們看我的手,我的腳,就知道實在是我了.」我們知道,他的手,他的腳,有釘痕,有釘十字架留下的痕跡,有受過大苦的記號.復活的耶穌,沒有變得皮光肉滑,他保留了人間慘痛經歷留下的記號.

今日,我們一起看看他的手,看看他的腳.

(1)揭露世人的罪

人間多有苦難.加害者,會千方百計,將苦難的痕跡洗淨.最深刻的畫面,莫如電影《滿城盡帶黃金甲》的一幕.一場屠殺過後,轉眼間,屍體被移走,血地被洗淨,然後,瞬間鋪滿金黃色的菊花,一片太平盛世的景象.

然而,人類的良知和歷史意識,並不輕易忘記苦難,總要讓苦難留下痕跡.歷史發生過很多場大屠殺,之後,都有屠殺紀念館將這些苦難永遠記在人類的靈魂深處.中國有南京大屠殺遇難同胞紀念館,東埔寨的金邊有「吐斯廉屠殺博物館」,當然,最出名的,是以色列耶路撒冷的大屠殺紀念館.在耶路撒冷的大屠殺紀念館裡,有一個地方,叫「遇難兒童紀念館」(The Children’s Memorial),在這裡,沒有燈光,只有燭光,緩緩地,背景的廣播唸出受難兒童的名字.每一個都是人,每一個都有名字,每一個都將神的光彩帶來人間.他們在人類造成的苦難中死去了,但人類要永遠唸著他們的名字.一個名字,包藏著無限苦難.

然後,在人類苦難的紀念館裡,我們聽到「耶穌」這個名字.

如今,他永遠活著.他永遠地向我們展示著他的手和他的腳,他永遠地向我們呼喚:「看我的手,看我的腳」,看看我的手和腳所承受的苦難,千真萬確,這是人類加給另一個人類的苦難,永恆地,永恆地,不會磨滅,不會消退.他的手,他的腳,像所有屠殺紀念館一樣,永恆地展示著苦難的真實,永恆地叫人不要忘記,永恆地要人正視,這是人類加給另一個人類的苦難.復活的耶穌,沒有展示完美,健康,幸福無比的身體,他展示他的手和腳.人類的苦難,刻在他身上,也永遠地刻在上帝的身上.永恆地,上帝記得人類的殘酷.上帝以自己為苦難的紀念碑,讓人類永遠不要忘記.

在聯合國的網頁裡,有一個網頁,題目叫《大屠殺教育的意義何在?》,網頁上寫(\url{https://zh.unesco.org/news/da-tu-sha-jiao-yu-yi-yi-he-zai}):

「大屠殺歷史能讓我們明白,理應保護所有人安全和權利的社會和相關機構是何等脆弱,某些機構甚至會轉身針對我們中的特定群體.……」

「大屠殺歷史能讓我們明白,人類行為的某些方面會影響到整個社會,比如尋找替罪羊,以簡單方案敷衍複雜問題；極端暴力和濫用權力的可能性；恐懼,同儕壓力,冷漠,貪婪,怨恨等在社會政治關係中所扮演的角色.」

簡言之,紀念人類的苦難,就是讓人類認清自己的軟弱和罪性.

耶穌向我們所有人說:「來,看看我的手,看看我的腳」,是要我們認清這個事實,無論我們人類多進步,社會多繁榮,經濟多興旺,我們人性的深處,都藏著對「異於我者」的排斥,敵對,殺機.我們人性背後的殘忍,令人不寒而慄.看著耶穌的手和腳,我們知道,我們只是軟弱無比的罪人.看著耶穌的手和腳,你就會明白,你完全沒有自滿,自負,驕傲的理由.由神的身體建造起來的,不是勝利紀念碑,而是屠殺紀念館,在這裡,讓人永遠安靜,永遠沉思,永遠懺悔.

(2)與人類同受苦難(Com-passion)

復活的耶穌,永恆地保存著人間受難的標記,這不單告訴世人,世人的真面目；也同時告訴世人,上帝的真面目.

上帝是誰?

祂是那位進入人間,受盡人間苦楚的上帝.這位上帝,不是自我清高的,隔岸觀火的上帝,祂與人同受苦難.苦難的英文是passion,同受苦難就是com-passion,這就是憐憫的真義.真正的憐憫,就是進入苦境中,與受苦的人同在.有些傳道人,為了給露宿者一點溫暖,會與露宿者一起睡在街頭,這是同受淒涼,這是com-passion.

我知道,有一種「神觀」,認為真正的上帝永不受苦,受苦的,只是人性的耶穌.但明顯,復活的耶穌,已是神性的生命.復活的耶穌,要升到天上,坐在父神的右邊,成為萬有的主.這是說明,耶穌與天父為一.耶穌就是上帝,祂的身體上,仍永恆地帶著人類苦難的痕跡.看看祂的手,看看祂的腳,受苦的記號永遠烙在上帝的生命中.我們的上帝是會受苦的,並且,祂永遠與受苦的人同在.

很多人問我,世間這麼多苦難,為何神坐視不理?為何神不用神蹟變走一切苦難?言下之意,是神要為這一切負責.這種想法合理嗎?人類歷史的浩劫,都是人類自己搞出來的,人不為此負責,而要神負責,這合理嗎?

面對人類的苦難,神沒有變走苦難,但祂並不是坐視不理,祂用一種意想不到的方式採取行動.祂默默的,卑微的,來到人間,陪伴著受苦的人,成了受苦人之中最受苦的人.他陪伴孤獨的人,事實上,有誰比他更孤獨?他陪伴被社會遺棄的人,事實上,有誰比他更被社會遺棄?他在十架上,陪伴釘十架的人,事實上,有誰比他更受人間的苦楚?請不要怨上帝袖手旁觀,他進入這個人間火海,比你我進得更深.看他的手,看他的腳,有深刻的痛苦痕跡.對你而言,這是否已足夠?看看他的手,看看他的腳,你的苦難,神也有份.在你苦難的深處,他在那裡,對你而言,這是否已足夠?究竟,一個識變法術的上帝好,抑或一個陪著你同度苦難的上帝好?

來,看看這位上帝的手,看看這位上帝的腳,祂陪伴受苦的人,祂自己也受盡苦楚,在祂生命裡,永遠記得受苦的人,永遠與他們同在.這就是我們的上帝.

(3)向前展望的回憶(Remembering Looking Forward)

人類建築這麼多的「屠殺紀念館」,不只是為了回憶,不只是為了重述過去的歷史,不只是為了挑起我們的傷感,而是為了將來,為了將來不再重蹈覆轍,為了有一天人類不再需要建築另一間新的「屠殺紀念館」.

德國神學家默茨(Johann Baptist Metz)提出一個觀念,他說,基督教講的回憶,重點不是向過去回望,更重要的,是向將來展望.一種面向將來的回憶.

譬如:以聖餐來說,若我們用哥林多前書的經文,那裡說:「你們每逢吃這餅,喝這杯,是表明主的死,直等到他來.」(林前11:26)回憶主的死,不是結尾,重點是:「直等到他來」.回憶主的死,是向將來展望的.若我們用馬太福音的經文,主耶穌說:「我告訴你們,從今以後,我不再喝這葡萄汁,直到我在我父的國裡同你們喝新的那日子.」(太26:29)同樣,回憶不是結尾,重點是:「等待喝新的那日子」.回憶,紀念,為的不是抓緊過去,為的是讓我們有正確的生命面向將來.

看,耶穌的手,耶穌的腳,耶穌以自己的身體作為人類苦難的紀念館,永恆地展示人類加給無辜受害者的痛苦.耶穌告訴我們,留在上帝的永恆裡的,不是勝利者的歷史書寫,而是每一個受苦人的生命.在人間,勝利者用勝利的口吻書寫歷史.在永恆,神永遠記得的,是受害者充滿傷痕的身體.

耶穌以自己的身體書寫著人類苦難的歷史,是讓我們對現世的一切,一切經濟,社會,科技,政治的強大力量,加以深刻反省.從受苦者的角度來反省,就是反省人類為何還是不斷誤用種種力量.從受苦者的歷史來展望,就是展望那一天,一切力量,都不再是暴力,而是仁愛的力量.耶穌以苦難去反對暴力.

然後,讓我們特別留意,今日,耶穌叫我們看的手和腳,是復活耶穌的手和腳.復活這件事情,意義宏大,影響著天地萬物,古往今來,我窮一生,都未能說明於萬一.復活,不是翻生.翻生,是在自然世界重新打滾一次,只是再度紅塵.復活,不是這種塵世事件,而是永恆事件.上帝使耶穌復活.復活,是一種全新的開始,是一種全新的創造.上帝,以其創造的全能,從死亡中,創造出新生命.從虛無中,創造存在.From nothing to being.透過耶穌的復活,上帝顯明,祂創造了一個全新的世界,永恆進入了塵世,生命的光明驅散了黑暗.這是上帝對死亡,對黑暗,對暴力永恆地說「不」.死亡,No.黑暗,No.暴力,No.

神學家說,耶穌透過死亡,令「死亡」這東西死亡.看,耶穌的手,耶穌的腳,這雙手腳,銘刻著人類的苦難,這些苦難,曾令耶穌死亡.今日,我們看耶穌的手,耶穌的腳,這已是復活者的手和腳,上帝將自己生命的權能,滲透在這雙手腳裡.當復活的耶穌叫我們去看祂的手和腳時,讓我們看到,上帝記掛著塵世每個受苦的人,上帝進到苦難的無盡處,以復活的光明,推翻苦難的終極話事權.藉此,激發我們盼望,向歷史的將來盼望,盼望一個不再有苦難的將來.今日,當主耶穌邀請我們去看祂的手,祂的腳時,祂是邀請你,去參與這歷史的工程,放下仇恨,放下暴力,放下惡言惡語,放下對弱者的嘲諷,放下無意義的爭競,放下傷害自己及他人的嫉妒,放下製造人類苦難的任何武器.上帝叫我們看耶穌的手和腳,是要呼喚我們,參與建造一個無苦難的將來.既然復活的生命已顯明,故此,不要失望,天國是會降臨的.

願榮耀歸上主,願人間享和平.



\newpage

\section{存在的勇氣}
\label{sec:547}
\textbf{林豪恩先生 - 存在的勇氣}
\newline
\newline
\hyperref[sec:546]{< < < PREV SERMON < < <}
~
\hyperref[sec:toc]{[返目錄]}
~
\hyperref[sec:548]{> > > NEXT SERMON > > >}


使徒行傳 4:5-13
\newline
\begin{longtable}{cl}
\hline
\hline
章節 & 經文 (和合本修訂版)\\
\hline
4:5 & \begin{tabularx}{0.7\textwidth}{X} 第二天，官長、長老和文士在耶路撒冷聚集， \end{tabularx} \\ \\ \relax
4:6 & \begin{tabularx}{0.7\textwidth}{X} 又有亞那大祭司、該亞法、約翰、亞歷山大，和大祭司的親族都在那裡。 \end{tabularx} \\ \\ \relax
4:7 & \begin{tabularx}{0.7\textwidth}{X} 他們叫使徒站在中間，問他們：「你們憑甚麼能力，奉誰的名做這事呢？」 \end{tabularx} \\ \\ \relax
4:8 & \begin{tabularx}{0.7\textwidth}{X} 那時，彼得被聖靈充滿，對他們說：「民間的官長和長老啊， \end{tabularx} \\ \\ \relax
4:9 & \begin{tabularx}{0.7\textwidth}{X} 倘若今日我們被查問是因為在殘障的人身上所行的善事，就是這人怎麼得了痊癒， \end{tabularx} \\ \\ \relax
4:10 & \begin{tabularx}{0.7\textwidth}{X} 那麼，你們大家和以色列全民都當知道，站在你們面前的這人得痊癒，是因你們所釘在十字架、神使他從死人中復活的拿撒勒人耶穌基督的名。 \end{tabularx} \\ \\ \relax
4:11 & \begin{tabularx}{0.7\textwidth}{X} 這位耶穌是：『你們匠人所丟棄的石頭，已成了房角的頭塊石頭。』 \end{tabularx} \\ \\ \relax
4:12 & \begin{tabularx}{0.7\textwidth}{X} 除他以外，別無拯救，因為在天下人間，沒有賜下別的名，我們可以靠著得救。」 \end{tabularx} \\ \\ \relax
4:13 & \begin{tabularx}{0.7\textwidth}{X} 他們見彼得、約翰的膽量，又看出他們原是沒有學問的平民，就很驚訝，認出他們曾是跟耶穌一起的； \end{tabularx} \\ \\
[1ex]
\hline
\hline
\end{longtable}
講題:存在的勇氣　The Courage to Be

講員:林豪恩先生

經文:使徒行傳4:5-13

場合:香港中文大學崇基學院禮拜堂主日崇拜

日期:2021年4月25日

引言

「存在的勇氣(The Courage To Be)」是神學家田立克一本書的書名.為甚麼人需要存在的勇氣呢?他認為因為人生面對焦懼(anxiety)及恐懼(fear),恐懼是有具體對象,例如疾病,飢餓,人會找方法去解決,但焦懼是空無,是有限的人面對死亡的無限,是無對象的,因此也找不到方法去解決.今天,我們不是討論「存在的勇氣」這本書,只是借書名為點題來閱讀經課的四段經文.

經文及解釋

使徒行傳:耶穌死而復活是我們存在的勇氣

4:5第二天,官府,長老,和文士在耶路撒冷聚會,

4:6又有大祭司亞那和該亞法,約翰,亞歷山大,並大祭司的親族都在那裡,

4:7叫使徒站在當中,就問他們說:「你們用甚麼能力,奉誰的名做這事呢?」

4:8那時彼得被聖靈充滿,對他們說:

4:9「治民的官府和長老啊,倘若今日因為在殘疾人身上所行的善事查問我們他是怎麼得了痊癒,

4:10你們眾人和以色列百姓都當知道,站在你們面前的這人得痊癒是因你們所釘十字架,上帝叫他從死裡復活的拿撒勒人耶穌基督的名.

4:11他是你們匠人所棄的石頭,已成了房角的頭塊石頭.

4:12除他以外,別無拯救；因為在天下人間,沒有賜下別的名,我們可以靠著得救.」

這是第一個勇氣的故事.勇氣在哪裡呢?在接下來的那一節經文:

4:13他們見彼得,約翰的膽量,又看出他們原是沒有學問的小民,就希奇,認明他們是跟過耶穌的；

膽量這個字,可以譯為勇氣,自信或無懼.是甚麼狀況需要門徒鼓起勇氣回應呢?你有沒有遇過需要勇氣去面對的狀況呢?

這個段落以「第二天」開始,讓人猜想事情應該與「第一天」有關.沒錯,門徒在之前的一天被聖殿的守殿官拘捕,被祭司們扣押過夜,第二天被帶上公會審訊.何事被審問呢?:(4:7)「你們用甚麼能力,奉誰的名做這事呢?」「這事」是甚麼事呢?(4:9)彼得的回答是「在殘疾人身上所行的善事」.那又是甚麼事呢?這是一個經常在兒童主日學講的故事:

3:1申初禱告的時候,彼得,約翰上聖殿去.

3:2有一個人,生來是瘸腿的,天天被人抬來,放在殿的一個門口(那門名叫美門),要求進殿的人賙濟.

3:3他看見彼得,約翰將要進殿,就求他們賙濟.

3:4彼得,約翰定睛看他；彼得說:「你看我們!」

3:5那人就留意看他們,指望得著甚麼.

3:6彼得說:「金銀我都沒有,只把我所有的給你:我奉拿撒勒人耶穌基督的名,叫你起來行走!」

3:7於是拉著他的右手,扶他起來；他的腳和踝子骨立刻健壯了,

3:8就跳起來,站著,又行走,同他們進了殿,走著,跳著,讚美上帝.

如果我們細心閱讀,就會發現門徒被拘捕和審判,未必是醫治瘸子這行動,而是他們對行動的解說:

3:12彼得看見,就對百姓說:「以色列人哪,為甚麼把這事當作希奇呢?為甚麼定睛看我們,以為我們憑自己的能力和虔誠使這人行走呢?

3:13亞伯拉罕,以撒,雅各的上帝,就是我們列祖的上帝,已經榮耀了他的僕人耶穌；你們卻把他交付彼拉多.彼拉多定意要釋放他,你們竟在彼拉多面前棄絕了他.

3:14你們棄絕了那聖潔公義者,反求著釋放一個兇手給你們.

3:15你們殺了那生命的主,上帝卻叫他從死裡復活了；我們都是為這事作見證.

4:1使徒對百姓說話的時候,祭司們和守殿官,並撒都該人忽然來了.

4:2因他們教訓百姓,本著耶穌,傳說死人復活,就很煩惱,

4:3於是下手拿住他們；

門徒面對官府需要勇氣,因為他們宣講與官府相反的判斷.官府判定耶穌「誘惑國民,禁止納稅給凱撒,並說自己是基督,是王」(路加福音23:2),最後把他視為判亂罪的罪犯處以釘十字架死刑.現在門徒卻公開宣傳耶穌不是罪犯,而是彌賽亞,是「那聖潔公義者」,即是否定官府的判決,指責官府判斷錯誤.

官府擁有權力,也擁有武力以實施其權力.不久之前,官府就運用其權力把他們所跟隨的耶穌拘捕和處決.現在門徒再次面對官府,怎可能不恐懼呢?他們如何有如此的勇氣呢?

其實,門徒的勇氣不是一直都有的.大家還記得耶穌被拘捕時,門徒立即四散,連信誓旦旦的彼得,也三次不認主.現在門徒的勇氣,是因為重遇死而復活的耶穌而來的,在另一次面對官府的審判時,門徒如此說:

5:29彼得和眾使徒回答說:「順從上帝,不順從人,是應當的.

5:30你們掛在木頭上殺害的耶穌,我們祖宗的上帝已經叫他復活.

5:31上帝且用右手將他高舉,叫他作君王,作救主,將悔改的心和赦罪的恩賜給以色列人.

5:32我們為這事作見證；上帝賜給順從之人的聖靈也為這事作見證.」

耶穌死而復活,證實他果真是彌賽亞,對門徒來說,這是事實,他們是為這事實作見證.官府曾經恐嚇門徒,禁止他們宣揚耶穌是基督.彼得回應說:「我們所看見所聽見的,不能不說.」(4:20)原來,耶穌死而復活這事實能夠給予門徒巨大的勇氣面對壓力和恐嚇.各位弟兄姊妹,你也像當年的門徒一樣肯定耶穌已經復活嗎?耶穌死而復活這事實能夠給予你勇氣去生存及面對生存的各種挑戰嗎?

詩23篇:上主順逆相伴是我們存在的勇氣

或者有人認為,生存其實不一定需要勇氣,而是很自然的事情.對某些人在某些人生狀況來說,生活順暢,無風無浪,無驚無險,有足夠能力及資源應付生活中的大小事情,完全不用觸及勇氣這回事.詩23篇的頭幾句會不會就是這種狀況呢?

1耶和華是我的牧者,

我必不致缺乏.

2他使我躺臥在青草地上,

領我在可安歇的水邊.

3他使我的靈魂甦醒,

為自己的名引導我走義路.

不缺乏,躺臥在青草地上,安歇在溪水旁……在如此的狀況中,舒服暢快,怎會需要勇氣呢?然而,你也知道這些狀況不會是人生的全部,從小開始,你就已經遇過不少需要勇氣作決定,也需要勇氣回應的處境,有些時候,甚至是需要勇氣才能繼續走下去.詩23篇的作者也不例外:

4我雖然行過死蔭的幽谷,

也不怕遭害,

因為你與我同在；

你的杖,你的竿,都安慰我.

5在我敵人面前,你為我擺設筵席；

你用油膏了我的頭,使我的福杯滿溢.

人生路上會有「死蔭的幽谷」,也會遇到「敵人」.陷入「死蔭的幽谷」中,或者面對「敵人」的時候,「怕遭害」是一般人正常的反應吧!詩人說「不怕遭害」,甚至在敵人面前擺設筵席如此豪邁,應該不是一般人的常態,而是有神奇力量的介入.詩人深信那介入於他的人生中的力量就是耶和華上帝,無論是順或是逆,耶和華上帝都與他同在,使他能夠享受青草地,溪水旁,也有勇氣走過「死蔭的幽谷」和面對有「敵人」的處境.

有一位朋友一直健康良好,生活有序,事奉殷勤,而且好學進修,最近意外地發現自己患上癌症,以下是她分享的心路歷程:

「10月23日進行MRI,25日抽組織化驗,12月2日醫生說我患的是淋巴癌.那位醫生是基督徒,他和我交代病情後,說要挑戰我的信仰.他問我怎樣看自己患病這事.我回答說:「日光之下無新事,人生亦無常.」但以理三位朋友「即或不然」的信仰,是我追求的屬靈情操.無病無痛時說倚靠主有多難?!唯有在艱難時刻仍然信靠順服才是真的操練.未來的日子雖然不容易,我深信每一個經歷都是寶貴的,求主幫助我無論在順境或逆境中,仍然尊主為大!以神為樂!

12月28日第一次化療後覆診,因中性白血球是零而需要留院.住院期間最初幾天的中性白血球仍然是零,後來又突然升到超過正常的指數,當中體會到前人說「病人常經歷像坐過山車高高低低的滋味」.2021年1月1日在醫院安靜默想時,我禱告上帝「我不要坐在過山車上,而是上帝你陪伴我在觀望台看著這輛過山車」.我會留意自己身體的狀況,但求主賜我不受太大的牽動,因為好壞,禍福,生死,都在上帝手裡.我的主診醫生每天都來看我,每次交代完病況,總是在信仰上給我一些挑戰或提醒.他對聖經的認真,對信仰的執著,給我很多啟迪.他多次提到「病」可以使人歸正,停一停,靜一靜,使人回歸正路,歸回上主,當人專注上主,是可以將「病」放下的.能夠在醫院安靜的房間過除夕及新年,是上帝給我很特別的禮物,好像住在退修營中,過了幾天安靜反思的日子.

詩23篇的作者繼續寫道:

6我一生一世必有恩惠慈愛隨著我；

我且要住在耶和華的殿中,直到永遠.

耶和華的殿在哪裡呢?不在耶路撒冷,不在禮拜堂,而是在有耶和華以恩惠慈愛相隨相伴的人生中.

約翰一書:人們彼此相愛是我們存在的勇氣

存在的勇氣不單單從上帝的同在而來,也從我們生活的社群而來.耶穌基督不單單呼召我們個別地跟隨祂,而且呼喚我們進入憂戚與共,互相扶持的社群.約翰一書的作者如此勸勉我們:

3:16主為我們捨命,我們從此就知道何為愛；我們也當為弟兄捨命.

3:17凡有世上財物的,看見弟兄窮乏,卻塞住憐恤的心,愛上帝的心怎能存在他裡面呢?

3:18小子們哪,我們相愛,不要只在言語和舌頭上,總要在行為和誠實上.

還記得上次講道時提及過冰島這個快樂指數名列前芧的地方,調查報告顯示,冰島超過百分之九十的人相信如果他們遇到困難,總會有人會來幫助他們.試想想,如果你知道當你遇到困難時,你周圍總會有人來幫助你,你會不會更有勇氣面對生活中各種的處境呢?

剛才提及那位朋友自己一個人住,她在患病接受治療的時候也有類似的發現:

「我第一次看專科醫生時,有一位前查經組長和她丈夫要陪我,當時我說「不必啦,我自己可以呀.」然而他們盛意拳拳,我卻之不恭.平日我間中陪同事看醫生看報告,今次反過來被人陪伴,原來這感覺是很窩心的.我母會及前實習教會的弟兄姊妹都嚷著要來探望及幫忙家居清潔；還有我就讀的神學院院長,屬靈導師,論文督導,老師,同學的探望,關心及代禱；家人,特別是我妹妹常常陪伴及幫忙處理很多實務；還有很多其他弟兄姊妹,親戚朋友及鄰居,不能盡錄,他們的禱告及關懷備至,頓時令我感覺萬千寵愛在一身.」

約翰一書這段經文有幾句特別的話:

3:21親愛的弟兄啊,我們的心若不責備我們,就可以向上帝坦然無懼了.

「無懼」這個字在使徒行傳譯為「瞻量」,即是「勇氣」,描述門徒因為傳揚耶穌是基督而面對官府審判時的態度.那份連面對上帝都可以坦然的勇氣是如何來的呢?前文說是「我們的心若不責備我們」.那麼,我們的心如何能在上帝面前安穩呢?前文說:

3:18小子們哪,我們相愛,不要只在言語和舌頭上,總要在行為和誠實上.

3:19從此就知道我們是屬真理的,並且我們的心在上帝面前可以安穩.

約翰在這裡提醒我們,內心的安穩並不是來自合起眼睛的冥想,而是來自真誠的和行動上的彼此相愛,即是互相照顧.

總結:約翰福音:屬於基督是我們存在的勇氣

在「存在的勇氣」書中,田立克認為恐懼有具體對象,例如疾病,飢餓,人會找方法去解決,但焦懼是空無,是有限的人面對死亡的無限,是無對象的,因此也找不到方法去解決.如果人能夠在空無之中找到所屬,在死亡的無限之中找到歸宿,這將會為存在注入勇氣,為人生建立基礎.在空無之中,你找著甚麼呢?在死亡的無限之中,哪裡是你的歸屬呢?福音書約課中耶穌的說話,會不會成為你的引導呢?

10:11我是好牧人；好牧人為羊捨命.

10:12若是雇工,不是牧人,羊也不是他自己的,他看見狼來,就撇下羊逃走；狼抓住羊,趕散了羊群.

10:13雇工逃走,因他是雇工,並不顧念羊.

10:14我是好牧人；我認識我的羊,我的羊也認識我,

10:15正如父認識我,我也認識父一樣；並且我為羊捨命.

10:16我另外有羊,不是這圈裡的；我必須領他們來,他們也要聽我的聲音,並且要合成一群,歸一個牧人了.

10:17我父愛我；因我將命捨去,好再取回來.

10:18沒有人奪我的命去,是我自己捨的.我有權柄捨了,也有權柄取回來.這是我從我父所受的命令.」

你歸屬誰呢?歸屬於不同的對象,狀況和結果都會不同.有些對象會用威嚇來令你就範,有些總會令你感到不安全和不放心.歸屬於耶穌基督的,就是在危難的時候,祂也不會撇下你不顧,反而捨命相救.祂是空無之中你可以找到的實在,是死亡的無限之處你可以安穩的歸屬.



\newpage

\section{問道門徒}
\label{sec:548}
\textbf{王家輝牧師 - 問道門徒}
\newline
\newline
\hyperref[sec:547]{< < < PREV SERMON < < <}
~
\hyperref[sec:toc]{[返目錄]}
~
\hyperref[sec:549]{> > > NEXT SERMON > > >}


使徒行傳 8:26-40
\newline
\begin{longtable}{cl}
\hline
\hline
章節 & 經文 (和合本修訂版)\\
\hline
8:26 & \begin{tabularx}{0.7\textwidth}{X} 有主的一個使者對腓利說：「起來！向南走，往那從耶路撒冷下迦薩的路上去。」那路是曠野。 \end{tabularx} \\ \\ \relax
8:27 & \begin{tabularx}{0.7\textwidth}{X} 腓利就起身去了。不料，有一個埃塞俄比亞人，是個有大權的太監，在埃塞俄比亞女王甘大基的手下總管銀庫，他上耶路撒冷去禮拜。 \end{tabularx} \\ \\ \relax
8:28 & \begin{tabularx}{0.7\textwidth}{X} 回程中，他坐在車上，正念著以賽亞先知的書， \end{tabularx} \\ \\ \relax
8:29 & \begin{tabularx}{0.7\textwidth}{X} 聖靈對腓利說：「你去！靠近那車走。」 \end{tabularx} \\ \\ \relax
8:30 & \begin{tabularx}{0.7\textwidth}{X} 腓利就跑到太監那裡，聽見他正在念以賽亞先知的書，就說：「你明白你所念的嗎？」 \end{tabularx} \\ \\ \relax
8:31 & \begin{tabularx}{0.7\textwidth}{X} 他說：「沒有人指教我，怎能明白呢？」於是他請腓利上車，與他同坐。 \end{tabularx} \\ \\ \relax
8:32 & \begin{tabularx}{0.7\textwidth}{X} 他所念的那段經文是這樣：「他像羊被牽去宰殺，又像羔羊在剪毛的人手下無聲，他也是這樣不開口。 \end{tabularx} \\ \\ \relax
8:33 & \begin{tabularx}{0.7\textwidth}{X} 他卑微的時候，得不到公義的審判，誰能述說他的身世？因為他的生命從地上被奪去。」 \end{tabularx} \\ \\ \relax
8:34 & \begin{tabularx}{0.7\textwidth}{X} 太監回答腓利說：「請問，先知說這話是指誰，是指自己，還是指別人呢？」 \end{tabularx} \\ \\ \relax
8:35 & \begin{tabularx}{0.7\textwidth}{X} 腓利就開口，從這段經文開始，對他傳講耶穌的福音。 \end{tabularx} \\ \\ \relax
8:36 & \begin{tabularx}{0.7\textwidth}{X} 二人正沿路往前走，到了有水的地方，太監說：「看哪！這裡有水，有甚麼能阻止我受洗呢？」 \end{tabularx} \\ \\ \relax
8:37 & \begin{tabularx}{0.7\textwidth}{X} 腓利說：『你若是一心相信，就可以。』他回答：『我信耶穌基督是神的兒子。』 \end{tabularx} \\ \\ \relax
8:38 & \begin{tabularx}{0.7\textwidth}{X} 於是他吩咐把車停下來，腓利和太監二人一同下到水裡，腓利就給他施洗。 \end{tabularx} \\ \\ \relax
8:39 & \begin{tabularx}{0.7\textwidth}{X} 他們從水裡上來，主的靈把腓利提了去，太監再也看不見他了，就歡歡喜喜地上路。 \end{tabularx} \\ \\ \relax
8:40 & \begin{tabularx}{0.7\textwidth}{X} 後來有人在亞鎖都遇見腓利；他走遍那地方，在各城宣揚福音，一直到凱撒利亞。 \end{tabularx} \\ \\
[1ex]
\hline
\hline
\end{longtable}
講題:問道門徒　Questioning as disciple-making

講員:王家輝牧師

經文:使徒行傳8:26-40

場合:香港中文大學崇基學院禮拜堂主日崇拜

日期:2021年5月2日

26有主的一個使者對腓利說:「起來!向南走,往那從耶路撒冷下迦薩的路上去.」那路是曠野.27腓利就起身去了,不料,有一個衣索匹亞人,是個有大權的太監,在衣索匹亞女王甘大基的手下總管銀庫,他上耶路撒冷禮拜去了.28現在回來,在車上坐著,念先知以賽亞的書.29聖靈對腓利說:「你去!貼近那車走.」30腓利就跑到太監那裡,聽見他念先知以賽亞的書,便問他說:「你所念的,你明白嗎?」31他說:「沒有人指教我,怎能明白呢?」於是請腓利上車,與他同坐.32他所念的那段經,說:他像羊被牽到宰殺之地,又像羊羔在剪毛的人手下無聲；他也是這樣不開口.33他卑微的時候,人不按公義審判他；誰能述說他的世代?因為他的生命從地上奪去.

34太監對腓利說:「請問,先知說這話是指著誰?是指著自己呢?是指著別人呢?」35腓利就開口從這經上起,對他傳講耶穌.36二人正往前走,到了有水的地方,太監說:「看哪,這裡有水,我受洗有甚麼妨礙呢?」38於是吩咐車站住,腓利和太監二人同下水裡去,腓利就給他施洗.39從水裡上來,主的靈把腓利提了去,太監也不再見他了,就歡歡喜喜地走路.40後來有人在亞鎖都遇見腓利；他走遍那地方,在各城宣傳福音,直到凱撒利亞.

引言

置身在大學校園之中,我們很容易就會遇上尋求知識的學生和進行學術研究的學者.通過閱讀,觀察,實驗,思考等這些過程我們可以累積知識,甚至有新的發現,有新的洞見.總括而言,這些都是做學問不可或缺的元素.

懂得提問,是增進知識的關鍵,故此最有意思的教育方式並不是令學生懂得答題,而是曉得發問!我們該是學「問」而不是學「答」!胡適先生有一句座右銘「做學問要在不疑處有疑,待人要在有疑處不疑」.這樣的思想不單是對一般的教育有幫助,對信仰的培育也同樣有幫助.

雖然今日香港的中小學已離開以背誦為主的教學方式,可是為了幫助學生應付考試,熟讀標準答案和答題技巧,仍然主導著學生的學習.同樣不少華人教會也是然熱衷於鼓勵信徒背誦聖經金句,彷彿背誦多少金句就可以顯示一個信徒的靈命高低.殊不知有些時候是金句背得朗朗上口,卻不一定明白當中的真意,更可惜的是因為只記憶一,兩節金句而忽略整段經文的上文下理,結果變成曉得背,卻不明瞭當中含意.

其實初期教會早已作了示範,領人歸主,建立門徒,一切也是以提問作為開始.當然有問亦需要有答,在答問之間引發更多思考,甚至有新發現,而不是只得既定的標準答案.今天經課當中的使徒行傳8:26-40正好給我們一個特別的例子作說明.讓我們先從這段記載中的兩位主角－腓利和衣索匹亞(《和合本修訂版》譯作「埃塞俄比亞」)太監說起.

傳道者腓利執事

經文承接前面講述腓利到撒瑪利亞展開傳福音的工作(徒8:5),故此在第26節就是繼續記載腓利遇見什麼人和做了那些事.腓利是何許人?按徒6:1-7所載教會揀選七個人去管理照顧寡婦的飯食問題,這就是選立教會執事的聖經基礎.當時腓利是其中一個獲選的,然而他最強的地方原來並不是管理飯食,而是宣教的熱忱.

在腓利開始傳道工作之先,有一位同樣被選派去管理照顧寡婦飯食問題的司提反,因為在會堂傳講耶穌基督的事,結果被捉拿,在大祭司和公會面前為主作見證,最後殉道而死.徒8:1記載從這日起教會遭遇大逼迫,除使徒在耶路撒冷之外,門徒都分散在猶太和撒瑪利雅各處,這些分散的人就往各處去傳道,福音就從耶路撒冷往外傳開了.在初期教會要向猶太人以外的族群傳講上帝接納和拯救他們是一件匪夷所思的事,就好像是殺人不用填命,還要得到政府嘉許!

腓利和那些被選立為執事的人,他們有一重要生命特質就是滿有聖靈與智慧(徒6:3),他們不是按自己的喜好辦事,而是按聖靈的帶領行事,因此腓利便在撒瑪利亞人當中,甚至向外邦人傳福音.徒8:26節告訴我們腓利得到指示往曠野去,很明顯那裡不會是人煙稠密的地方,以廣東話的講法更傳神,就是「鬼影都冇」.這個情況就像我告訴弟兄姊妹,起來我們去北極傳福音吧!大家會有什麼反應?可能大家心裡會想,難道牧師叫我們向北極熊傳福音?其實北極圈也有幾百萬人口,只是按人口密度計算是相當低.與其到北極,不如到北潭涌還好,現在每逢週末都會有許多人到這些郊野地方.傳福音當然要到多人的地方去,然而腓利卻順從指示就去了.

腓利的順從令我們向自己發出一個問題,就是我們跟腓利有什麼相似與相異之處?我們知道要傳揚福音,也有這樣的心志,這是跟腓利一樣的；可是我們許多時有太多自己固有的想法,甚至蓋過了聖靈的指示,這個就跟腓利不一樣了.我估有不少弟兄姊妹也有這樣的經驗.以為向某些人傳福音較為容易,例如那些比我們年輕,看來比較容易相處的人,因此大部份教會都會著力在學生福音工作.相反,在我們的印象中,成年人是比較難信主,因為他們會固執己見,好辯駁,難說服他們,所以無謂花太多力氣.現實又如何呢?學生時代信主當然是較容易,但若沒有好好的栽培,建立,長大後就會失落了.成年才信主,要突破的限制的確較多,可是當真心回應,一般都會是堅定的相信.相信在座有不少弟兄姊妹也是有這樣的信仰歷程.

尋道者埃塞俄比亞太監

這段記載的另一主角就是埃塞俄比亞太監.他是一位有大權的官,總管銀庫,大概相當於現在的財政司長／財務大臣.我們可以想像這位太監的心田是軟土,可以接受福音的種子.我們從何得知呢?

首先若這位太監是從現代埃塞俄比亞的位置到耶路撒冷上聖殿敬拜.一程的距離是4,114公里遠,就算以現代汽車行駛,也需時84小時,即是要3至4日.若在古代,可能要花上十天至兩個禮拜才可到達.聖經學者有另外的觀點,在原文中「埃塞俄比亞」這個字原來意思可以是「深色皮膚的地方」(place of the burnt-faces).如果這位太監是來自「深色皮膚的地方」,那麼其中一個可能就是今日蘇丹共和國的地方.距離耶路撒冷2,457公里,若在古代乘馬車,大概幾日至最多一個星期已到達.按使徒行傳時代的地域概念,這個地方大概就是他們所理解的地極之處.換言之,福音由耶路撒冷傳開,到司提反殉道後,就傳至撒瑪利亞,而埃塞俄比亞太監信主,就表示福音已傳到地極.

這位太監縱然到達耶路撒冷,他也不能進入聖殿敬拜上帝,原因是他是一位太監,這個規定於申命記23:1有清楚的記載.雖然他沒有機會進入聖殿,這位太監對於信仰的追求沒有因為離開耶路撒冷而結束,當他坐在馬車上,於歸家途中時,在搖搖晃晃的情況下仍然閱讀先知書.由此可見,他的心是何等的渴慕認識上帝.估計他不曾想過,自己沒有機會在耶路撒冷進入聖殿與上帝相遇,卻在回家路途的曠野上被安排遇見上帝所差的使者.這有如先知所說的:上帝是「在曠野開道路,在沙漠開江河.」(賽43:19)的上帝.

上帝為這位太監開出一條又真又活的道路,為他乾渴如沙漠的心開出活水的江河來,使他歡歡喜喜地回家.若照今天的環境而言,領人認識上帝也不一定在教會裡,可能是在路上,公司,家庭,旅行的路途上任何的地方,只要我們隨時作好準備,上帝也可以用我們作見證.

不論這位太監是來自現今的埃塞俄比亞(該國的人也按這段聖經記載,相信他們是承繼初期教會將福音傳至外邦當中,所以他們很早已成為基督教國家),還是今日的蘇丹,非洲教會其實也是歷史相當悠久的.根據初期教父愛任紐(Irenaeus,130年-202年)說(參考Against Heresies III 12.8),這位太監後來成為向埃埃塞俄比亞傳福音的宣教士.這個國家是非洲唯一從公元300到1974年,十七個世紀之久都保持著基督教國家,並且令北非洲很多國家也受基督教的影響,直到回教的勢力進入這個地區才有變化.1966年著名佈道家葛培理在德國柏林召開世界福音大會時,開幕禮拜演講的貴賓,就是當時埃塞俄比亞的國王賽拉西亞.尋道者埃塞俄比亞太監為他自己與自己的同胞有機會認識耶穌基督的福音帶來極重要的開始.

問道

經文繼續記載了腓利與太監的對話及當中的經過.當中包括了三組的問與答的交流.

第一條問題是由腓利提出:「你所念的,你明白嗎?」太監以反問方式回答:「沒有人指教我,怎能明白呢?」

第二條問題是由太監提出:「請問,先知說這話是指著誰?是指著自己呢?是指著別人呢?」當時太監所讀的是以賽亞書53:7-8《七十士譯本》.當腓利聽到太監問及有關這段以賽亞書的經文,自然可以將有關耶穌基督的事告訴太監了.聖經學者嘗試解釋以賽亞書53:7-8為何會令太監感興趣會大聲讀出來呢?顯然在太監還未遇上腓利前,太監對耶穌的事可能也不太清楚.反而先知書這段的記載令他想起自身的狀況.因為他是女王甘大基的大臣,而作為女王重要的臣子,跟女王接觸頻繁,而為了保證王室血統純正,這些大臣的命運就是要被閹割,令他們失去使女王懷孕的能力.這位太監讀到先知所說:「他像羊被牽到宰殺之地,又像羊羔在剪毛的人手下無聲；他也是這樣不開口.」這個情景跟他無法違抗女王之命,成為被閹之人有很大的相似之處.因此他向腓利查問,先知所說的是講及先知自己的情況,還是別人呢?如果先知所講的是指別人,那就可以推斷這段說話同樣適用於這位太監身上.能夠找到經文講及自身的狀況,那種共鳴感當然強烈.許多時,弟兄姊妹靈修讀經,也會覺得某些經文是在對自己說話.尤其是準備讀神學的弟兄姊妹,每每也是因為讀到某段經文而感到是上帝透過經文的信息來呼召他／她的.

第三條問題也是由太監提出的:「看哪,這裡有水,我受洗有甚麼妨礙呢?」腓利和太監去到有水的地方,這個是十分奇特的情況.因為按照上帝使者的指示,腓利和太監在曠野地方相遇,那裡能夠找到水源並不尋常.今天我們聽到太監的提問並不感希奇,因為有人願意信主並希望接受水禮,這一切也看似是理所當然.可是在初期教會,外邦人有份於救恩已是需要辯論之事,更何況接受水禮?還有太監是外邦人再加上是被閹之人,在猶太人傳統的觀念中,根本上是絕不可能得到接納,故此太監的提問,對腓利而言也是極大的考驗.這個考驗有如我們以往從未試過上班,上學全部也在家進行一樣,開始的時候完全不知道會有什麼將會發生一般.

最終太監就了洗,而腓利就像風一樣,一下子不見了,剩下太監,經歷人生新的開始,歡歡喜喜繼續踏上歸家之路.根據初期教父帕皮亞(Papias,70年-155年)說(參考Eusebius Ecclesiastical History III 39),主繼續使用腓利傳福音直到亞細亞省份,晚年得到眾教會的敬愛.約20年以後,腓利住在該撒利亞,聖經稱他是「傳福音的腓利」(徒二十8).

總結

初期教會按照基督的吩咐要從耶路撒冷起,進而擴至猶太全地,撒瑪利亞,直到地極為主作見證.腓利的傳福音工作就是為主作見證直到地極,不過並不是由他本人到地極,而是透過他與一位太監相遇而達成.他們兩人透過問與答,由不理解經文所說的,到得著澄清與指導.他們要突破當時固有的想法,外邦人得著救恩是否可能?一個被閹之人也能夠得著救恩?這樣的人可以受洗歸主?這些在今日我們不用太多爭辯,甚至認為是理所當然的決定,在那個時候是完全不能想像的.

今日我們傳揚福音,建立門徒,並不是把一套一套標準答案,基本教義灌輸給信徒,而是鼓勵他們提問,從問題中再繼續尋索理解.基督是復活的主,今日仍然活著.我們的信仰也是活的信仰,在不同的時代,環境中,會有新的可能,新的挑戰,保持求問,尋真的心,定必有新發現,能夠繼續被建立起來.求主賜智慧,也加給我們力量.阿們.



\newpage

\section{我愛}
\label{sec:549}
\textbf{高國雄牧師 - 我愛}
\newline
\newline
\hyperref[sec:548]{< < < PREV SERMON < < <}
~
\hyperref[sec:toc]{[返目錄]}
~
\hyperref[sec:550]{> > > NEXT SERMON > > >}


講題:我愛　I Love

講員:高國雄牧師

場合:香港中文大學崇基學院禮拜堂主日崇拜

日期:2021年5月9日

如果要用一個字來表達基督信仰的核心,相信大部份人亦會認同,就是「愛」!

「愛」這個字我們每一個都曾經講過,聽過,但相信大家對「愛」皆有不同的感受與體會.

最不學術,但最多人使用的維基百科中,對「愛」有以下描述:

愛,通常多見於人或動物,是指對一個人或某物的一種強烈的情感和依戀,衍生自親人之間的強烈關愛,忠誠及善意的情感與心理狀態.包含一系列強烈和積極的情感和精神狀態,從最崇高的美德或好習慣,最深的人際關係,到最簡單的愉快.這種情感和依戀驅使那些有這種感覺的人尋求與愛的對象在身體上,智力上甚至想像上的接近.

愛最佳的定義可能是主動心動,以真心對待某個生命體或物體(可以是人,動物 ,物品,神明),使其整體得到快樂.簡而言之,愛即主動使整體得到快樂.

愛可以包括靈魂或心靈上的愛,對法律與組織的愛,對自己的愛,對食物的愛,對金錢的愛,對學習的愛,對權力的愛,對名譽的愛,對他人的愛等,數之不盡.不同人對其所接受的愛有著不同的重視程度.

愛本質上是一個抽象概念,可以體驗但卻難以用言語或文字來表達.

以上描述未必我們完全認同,但相信也有其可參考性…

當我們聽到別人對自己說「我愛你」的時候,我們會有什麼反應呢?

回想在學的日子,有一次我在茶餐廳與同學一同午餐,一個侍應到來台邊,詢問我們需要點什麼餐點,他問到:「你們要(愛)些什麼?」有同學叫A餐,有同學叫常餐,快餐等…當侍應轉身向我詢問時,我忽然調皮起來,轉頭情款款地望著該侍應:「我愛你呀…」當然他馬上被我嚇了一跳,同學們也抱腹大笑…

廣東話有其非常考妙的地方,我「要」某事物,可以用我「愛」來取代,當然我們清楚知道,愛一個A餐,一件玩具與愛一個人有很大分別,但其背後可能心態上仍然有共通之處,「我愛」等如我們想擁有,想佔據,想掌控…

耶穌基督對我們講「我愛你們」又是什麼意思呢?

在約翰福音 15:9上「我愛你們,正如父愛我一樣；」

在經課中,我們看到耶穌基督直接提出祂愛我們,而且這愛是極親密,極神聖的,因為祂以祂與天父上帝的關係作比較,正如父愛耶穌一樣,請弟兄姊妹留意,千萬不要輕看這個描述,耶穌基督與天父上帝源於一,是三位一的上帝,其緊密/親密性是沒有任何關係可以相比的,耶穌基督一開始就用了最高的例子作類比,意即耶穌基督以最大的愛愛我們,超越世間上一切所有的愛…

約翰福音 15:9下:「你們要常在我的愛裡.」

耶穌基督對我們的期望就只有一個,就是常在祂的愛裡.是期望,也是應許,更是祝福的因由及途徑,不只是在愛中,更是要常在愛中…

如何可以常在耶穌基督的愛中呢?秘訣只有一個

約翰福音 15:10「你們若遵守我的命令,就常在我的愛裡,正如我遵守了我父的命令,常在他的愛裡.」

我們想起遵守命令,常常馬上容易產生負面的感覺,守命令就等如非自願,受捆綁,沒選擇,不喜歡,愛不就應該是自由的嗎?愛一個人就應該放手讓他享自由,而不是要用規條限制他!

約翰一書 5:1-3「凡信耶穌是基督的,都是從神而生,凡愛生他之神的,也必愛從神生的.我們若愛神,又遵守他的誡命,從此就知道我們愛神的兒女.我們遵守神的誡命,這就是愛他了.」

我們與神在愛中要建立的關係,是由「信」而生的「父母與兒女」的關係,因為有此關係,才自然地產生相應的行動,遵守神的誡命就是相感應有的行動,也是愛神正常/該有的表現,所以耶穌基督也自稱祂遵守了父的命令,因面常在父的愛裡…

作為信徒,我們其實也心知理應許要好好嚴守神的誡命,但「立志為善由得我,只是行出來由不得我.」羅馬書7 :18更常是我們的寫照及掙扎的地方…能夠努力好好守誡命的秘訣,動力,方法是什麼呢?

很多人覺得守誡命很難,又或很抗拒守誡命,甚至因此拒絕信仰,是因為以為要靠一己之力,甚至誤會了要以自己的能力/行為去換取信仰,人是軟弱的,並且生處於魔鬼掌權罪的世界中,當然無法勝過…

但約翰一書繼續直接的表達5:4-6「並且他的誡命不是難守的.因為凡從神生的,就勝過世界；使我們勝了世界的,就是我們的信心.勝過世界的是誰呢?不是那信耶穌是神兒子的嗎?6這藉著水和血而來的,就是耶穌基督；不是單用水,乃是用水又用血.」

守誡命不難,因為其實不是靠著我們自身的能力去遵守,也不是要我們努力勝過世界,而是憑信依靠那位已經勝過世界的耶穌基督,祂藉水和血,釘身十架,受死,而復活,又以水和血:洗禮與聖餐,使我們進入及繼續在祂的恩,祂的愛裡,從祂而生,那就能勝過世界…

耶穌基督期望我們持守誡命,常在祂的愛裡有一個目的,這同時又是最大的命令,就是要我們彼此相愛:

約翰福音 15:12「你們要彼此相愛,像我愛你們一樣；這就是我的命令.」

約翰福音 15:13-17「人為朋友捨命,人的愛心沒有比這個大的.你們若遵行我所吩咐的,就是我的朋友了.以後我不再稱你們為僕人,因僕人不知道主人所做的事.我乃稱你們為朋友；因我從我父所聽見的,已經都告訴你們了.不是你們揀選了我,是我揀選了你們,並且分派你們去結果子,叫你們的果子常存,使你們奉我的名,無論向父求甚麼,他就賜給你們.我這樣吩咐你們,是要叫你們彼此相愛.」

耶穌基督要成為我們的朋友,我們原是被造,更因罪的原故成為了罪的奴僕,耶穌基督的拯救不單讓我們可以脫離罪,更不需要我們只轉為主的僕人,而竟然要成為我們的「朋友」,把祂所知道的告訴我們,更揀選我們去為祂結果子(差派所屬事工給我們),甚至賜與我們祈求的權柄,為要讓我們(人)能夠彼此相愛

愛是什麼呢?愛是關係,是結連,是承諾,是委身…

今天是母親節,上主為我們預備的經課,正正就是提醒我們,父母對子女的愛是關係的愛,而不是計算的愛；是連繫(血緣)的愛,是上帝所賜而不是自己去賺取的；所以作兒女的要承諾及回饋.而對耶穌基督的愛,我們又會作什麼回應呢?



\newpage

\section{沒有幻想不用絕望}
\label{sec:550}
\textbf{鄧瑞強博士 - 沒有幻想,不用絕望}
\newline
\newline
\hyperref[sec:549]{< < < PREV SERMON < < <}
~
\hyperref[sec:toc]{[返目錄]}
~
\hyperref[sec:551]{> > > NEXT SERMON > > >}


約翰福音 17:13-19
\newline
\begin{longtable}{cl}
\hline
\hline
章節 & 經文 (和合本修訂版)\\
\hline
17:13 & \begin{tabularx}{0.7\textwidth}{X} 現在我到你那裡去，我在世上說這些話，是要他們心裡充滿了我的喜樂。 \end{tabularx} \\ \\ \relax
17:14 & \begin{tabularx}{0.7\textwidth}{X} 我已把你的道賜給他們；世界恨他們，因為他們不屬世界，正如我不屬世界一樣。 \end{tabularx} \\ \\ \relax
17:15 & \begin{tabularx}{0.7\textwidth}{X} 我不求你把他們從世上接走，只求你保全他們，使他們脫離那惡者。 \end{tabularx} \\ \\ \relax
17:16 & \begin{tabularx}{0.7\textwidth}{X} 他們不屬世界，正如我不屬世界一樣。 \end{tabularx} \\ \\ \relax
17:17 & \begin{tabularx}{0.7\textwidth}{X} 求你用真理使他們成聖；你的道就是真理。 \end{tabularx} \\ \\ \relax
17:18 & \begin{tabularx}{0.7\textwidth}{X} 你怎樣差我到世上，我也照樣差他們到世上。 \end{tabularx} \\ \\ \relax
17:19 & \begin{tabularx}{0.7\textwidth}{X} 我為他們的緣故使自己分別為聖，為要使他們也因真理成聖。 \end{tabularx} \\ \\
[1ex]
\hline
\hline
\end{longtable}
講題:沒有幻想,不用絕望　No Illusions, No Despair

經文:約翰福音17:13-19

講員:鄧瑞強博士

場合:香港中文大學崇基學院禮拜堂主日崇拜

日期:2021年5月16日

各位弟兄姊妹,早安.在這疫症裡,祝願各人平安.

只要你的思想境界夠高超,有藝術家般的敏銳情感,對靈性有深度的體會,你總會覺得,你不屬於這裡.你屬於高遠的藍天,遼闊的大海,無際的星空,無拘無束的流浪,鮮花綻放的原野,滿懷溫馨的黃昏,你屬於詩,屬於美酒,屬於神話中的英雄.而人間那些平淡無味的語言,令人作嘔的謊話,扭曲的嘴臉,卑污的手段,暴力的爭奪,對你而言,就十足像看youtube時忽然爆出來的冗長廣告,令人厭煩不堪.

但,生活也有不得不如此的另一面,人像被綁在大地上的野獸,有原始的獸性,有肉體的慾望,有基本的需求.

我們就是如此奇怪的組合,是精神與泥土的結合,介乎天使與野獸之間,存於上天與下地之中.人性的這兩種向度所構成的衝突與張力,引發出人類永遠都在追問的問題:理想重要抑或現實重要?上帝口出的話重要抑或麵包重要?像耶穌那樣活出生命的重量重要抑或像烏龜那樣縮起來保命重要?不同的人,不同的時代,有不同的答案.

有些哲人,為了解決精神與物質這兩種向度的張力,會高舉一方,貶抑另一方.如哲人柏拉圖,他強調精神,對他而言,物質是變幻的,易逝的,沒有永恆意義.他將一切事物的物質性剝掉,留下的,便是純觀念,這是不變的,這是永恆的,這是理型世界的真相.個別的人,會死的,沒有永恆性.但「人」的觀念,這是不變的,這才是真理所在.這種想法發展下去,便會發展成「為了某些不變的原則」而取消「個別的具體的人」.新約的法利賽人如此,二十世紀的希特拉也如此.從靈修而言,這種想法是渴慕那純粹的靈性,而漠視塵世.這有一種柏拉圖式的喜悅,沒有了時間,只留下永恆.飛上了天上,卻失去歷史.

另一些哲人,為了解決精神與物質這兩種向度的張力,寧願放棄精神,擁抱塵世.他們認為,根本沒有那種離地的精神,有的,只是腳下的大地.唯物主義者認為,一切一切,到最後,只是物質.沒有什麼永恆的東西,一切都是變動的,都只是歷史的產物.這種想法發展下去,就只剩下對物質的追求,和在物質的追求的過程中無可避免的殘酷惡鬥.新約中的代表,便是和法利賽人相對立的撒都該人,他們的思想特徵是不相信復活,不相信永恆,有的,只是現世.他們和權貴勾結,活一種講求當下滿足的人生.從靈修而言,這種想法造就兩種可能性.若是權貴,這想法就是在歷史的過程中,力保自己的利益,剷除阻路者.若是受苦者,這想法就是推動不斷的變革.他們有的是歷史的行動,有的是現世利益的角力,沒有的,是永恆的價值.

在福音書裡,互相敵對的法利賽人和撒都該人最終攜手合作,除滅耶穌,但兩者理由不同,法利賽人是為了理想,為了不變的原則,而消滅耶穌.撒都該人是為了現實,為了實際的利益,而消滅耶穌.自命清高的人,嫌耶穌太世俗；完全世俗的人,嫌耶穌太清高.最後,他們都要置耶穌於死地.

今日,我們看看耶穌,在精神與現實,在現世與他世,在永恆與現在之間,他如何活著?他如何鼓勵跟從他的人?

今日講道經文:約翰福音17: 13-19

17:13現在我往你那裡去,我還在世上說這話,是叫他們心裡充滿我的喜樂.

17:14我已將你的道賜給他們.世界又恨他們；因為他們不屬世界,正如我不屬世界一樣.

17:15我不求你叫他們離開世界,只求你保守他們脫離那惡者.

17:16他們不屬世界,正如我不屬世界一樣.

17:17求你用真理使他們成聖；你的道就是真理.

17:18你怎樣差我到世上,我也照樣差他們到世上.

17:19我為他們的緣故,自己分別為聖,叫他們也因真理成聖.

這是主耶穌上十架前的禱文,為門徒祈禱,求天父保守他們.這禱文完全反映耶穌如何看待自己及信徒的生命.

耶穌如何看待自己及信徒的生命,我歸納為七點.

一,基督及基督的跟隨者,他們的生命不屬世界,他們是另一個世界的人.

二,但同時,他們卻也不離世界,雙腳緊緊地立足於大地之上.

三,屬於另一世界,耶穌稱之為「聖」,耶穌及其門徒是屬於神「聖」的,不是屬於世俗的,即聖經講的「分別為聖」.

四,是真理使人成聖,使人與世界有別.

五,「神聖的人」活在「不神聖的世界」,當然有衝突,有張力,耶穌稱之為「恨」,這世界恨那些不屬這世界的人.

六,不屬世界的人,留在這世界上,是因為有使命在身.天父差派不屬世界的耶穌來到這世界,是因為有使命交給耶穌去完成.耶穌差派不屬世界的門徒繼續留在這世界,同樣,是因為有使命交給他們去完成.

七,耶穌的跟隨者知道自己的身分,知道自己有任務在身,同時也知道天父的保守,故能勇於面對種種困難,「笑傲江湖」,在險惡的江湖裡,能高唱「滄海一聲笑」.耶穌稱這態度為「喜樂」.

以上各點,我每點講少許,即是說,當我講到第7點時,大家便可開心地期待,我快講完了.

(1)信徒的生命不屬這世界

很多體會到永恆的人,都不屬乎這世界的.

其中一個經典人物是古希臘的數學家阿基米德,據說,當羅馬大軍攻佔他居住的城巿時,他正在研究圓形的幾何學.一個士兵走進來,踏在他的圓圈上,他對現實的戰爭看來毫不上心,他沉思著圓形,然後對那士兵說,「請不要踏著我的圓圈」.儘管現實在戰爭中,這位天才的心思卻貫注在圓形的永恆法則裡.那士兵怒不可遏,就這樣,揮劍殺死了一個天才.

近來讀一本小書,是德國浪漫文學的一個代表人物Wackenroder的小書,書名叫《一個熱愛藝術的修士的內心傾訴》,我讀中文譯本.書中講到大畫家拉斐爾如何畫聖母像.拉斐爾說:「由於難得尋見完美的女人形體,所以我遵從降臨到我靈魂中的源於精神的形象.」這即是說,在現實中,拉斐爾根本找不到合宜的模特兒,神聖的形象只能源自天上.書中講到,拉斐爾很沮喪,因為他找不到滿意的女士形象.直到有一個晚上,聖母向他顯現.「聖母用一種無法言表的動人目光注視著他,活靈活現」.神聖的形象臨在於他,他憑此畫下一幅又一幅的聖母像.屬於這個世界的任何形象,對這大畫家而言,毫無幫助.幫助他的,是不屬乎這世界的,來自他世的,一種神聖的顯現.

天才人物,看來都是不屬這世界的,他們屬於永恆的維度,屬於聖人的團體.基督的信徒,也是天才人物,我們領悟到天父永恆的愛,聆聽到永恆的真理之言,生命進入了不滅的天國.耶穌說,我們像他一樣,不屬乎這個世界.

(2)信徒卻也不離世界

一個人,體悟永恆,也永遠留駐在永恆中,則應是哲人.一個人,體悟永恆,卻又回頭走向人間,為了要救蒼生,這就是宗教人.

大乘佛教有觀世音菩薩,在遠古以前,他已成佛,但他不留在淨土,而入凡塵,是因為他要教化世人.

耶穌並沒有將他的信徒帶離塵世,進入永恆,因為耶穌之為耶穌,總是「道成肉身」,將塵世的苦難擔在自己的肩頭上的.耶穌如此,他的信徒也如此.

若果耶穌完全遠離塵世,不吃人間煙火,法利賽人應該接納他.若果耶穌完全擁抱塵世,在物質中打滾,撒都該人應該歡迎他.但耶穌不屬世界,卻又不離世界,於是,法利賽人和撒都該人都拒絕他,除之而後快.

(3)信徒之「分別為聖」

信徒與世人有分別,而單單屬於神,這就是聖經講的「分別為聖」.

信徒最大的問題,是忘記了自己屬於誰.世上各種吸引人的東西,如:金錢,權力,性愛,時刻構成誘惑,向我們建議,你下拜我們,就會得到這世上的富貴快活.

耶穌提醒我們,我們要永遠記得自己屬誰.我們屬於神,單單敬拜祂.

(4)是真理使人成聖

如何保持自己屬於神?

堅持真理.特別是在說謊話有著數的年代,更要堅持說真話.真與假,中間有一絕對界線.

我們不屬世界,卻又活在世界中,將這生命活出來,就是:在謊話的世界中,堅持真話.這是難之又難的.當謊話夾雜著暴力與利益,則真話便意味著受苦與虧損.這是人間的十架,這是基督的十架,這是信徒的十架.信徒活在人間,就是跟從基督,背起這十架.這是難之又難的.只求慈愛的天父,減少這艱難的日子.

(5)世界恨你們

耶穌說,世界恨你們.為什麼?因為世界不喜歡真理,不喜歡光明,不喜歡上帝.

有些信徒,為免被世界恨惡,便完全循入永恆.信仰,只為靈魂的得救,無關乎塵世的責任.一旦遠離世界,不再接觸世界,不再擔起世間的什麼十字架,當然便不會和世界有任何衝突,也不會被世界恨惡,因為這樣的信仰,根本不會和現世有任何的交集.一個從未誕生在塵世的永恆耶穌,當然不需要受死.

另有一些信徒,為免被世界恨惡,便索性完全歸屬世界,成為世界的一分子.「基督徒」這身分變成只是化裝舞會的面具,不再代表生命的真正歸屬,不再有任何真正的內容.成為世界的一分子,按世界的邏輯來生活,當然不會被世界恨惡,反而得到世界的擁抱.

耶穌沒有抽離世界,也沒有變成世界的一部分.他永遠以神子的身分,行走在大地上.面對世界的恨,他沒有幻想.同時,面對世界的恨,也不會失望.耶穌明白,要將真理植根大地,沒有不用付代價的樂觀,也沒有不結果子的悲觀.沒有坦途,但也沒有絕境.耶穌走過這充滿「恨」的大地,他受盡苦楚,然後,復活的真理在他身上展現.因著耶穌,天父展現祂對大地的永恆主權.今日,世界也會恨我們,讓我們以耶穌的生命去面對,沒有幻想,但也不用絕望.

(6)耶穌差派我們到世上

耶穌留我們在世上,是要我們像他一樣,完成天國的使命.

只要大地有黑暗,便需要有人點燃燭光.只要世界有恨,便需要有人勇於去愛.只要塵世充滿著死亡的絕望,便需要有人顯明復活的希望.請不要為自己而活,請為天國的使命而活,直到天父的國度,圓滿地實現在人間.

(7)喜樂的人生

講到第7點了,各人應該很開心,體會到什麼叫喜樂.

天降大任於我們,逆境,艱難,十架,世界的誘惑,世人的恨,人性的恐懼,現實的絕望,這一切,都是真的.

我們仰賴的,是天父的保護.我們相信,天父看見.

在這難得一活的歷史時刻,就憑天父的看顧,讓我們喜樂地面對.

(8)結語

最後,我以哥林多後書6: 2-10作結,我以《現代中文譯本》讀出.

保羅說:

6:2下　聽啊,現在就是接納上帝恩惠的時刻!今天就是上帝拯救的日子!

6:3　我們不願意別人毀謗我們的事奉,所以我們也不妨礙別人.

6:4　為了在一切事上表明我們是上帝的僕人,我們忍受各種患難,貧窮,困苦.

6:5　我們曾遭受鞭打,監禁,暴民的騷擾；我們勞苦,失眠,挨餓.

6:6　我們以純潔,知識,忍耐,仁慈,更藉著聖靈,真摯的愛,

6:7　真理的信息,和上帝的大能,來表明我們是上帝的僕人.我們以正義作武器,攻擊敵人,也保衛自己.

6:8　我們受尊敬,也被侮辱；受責罵,也被稱讚.人家說我們是騙子,其實我們說的是誠實話；

6:9　說我們寂寂無聞,其實是家喻戶曉；說我們死了,其實都好好地活著.我們受刑罰,卻沒有被殺；

6:10　憂傷,卻常有喜樂；貧窮,卻使許多人富足；好像一無所有,卻樣樣都有.

願榮耀歸上主,願人間享和平.



\newpage

\section{從五旬節到聖靈降臨}
\label{sec:551}
\textbf{林豪恩先生 - 從五旬節到聖靈降臨}
\newline
\newline
\hyperref[sec:550]{< < < PREV SERMON < < <}
~
\hyperref[sec:toc]{[返目錄]}
~
\hyperref[sec:552]{> > > NEXT SERMON > > >}


使徒行傳 2:1-21
\newline
\begin{longtable}{cl}
\hline
\hline
章節 & 經文 (和合本修訂版)\\
\hline
2:1 & \begin{tabularx}{0.7\textwidth}{X} 五旬節那日到了，他們全都聚集在一起。 \end{tabularx} \\ \\ \relax
2:2 & \begin{tabularx}{0.7\textwidth}{X} 忽然，有響聲從天上下來，好像一陣大風吹過，充滿了他們所坐的整座屋子； \end{tabularx} \\ \\ \relax
2:3 & \begin{tabularx}{0.7\textwidth}{X} 又有舌頭如火焰向他們顯現，分開落在他們每個人身上。 \end{tabularx} \\ \\ \relax
2:4 & \begin{tabularx}{0.7\textwidth}{X} 他們都被聖靈充滿，就按著聖靈所賜的口才說起別國的話來。 \end{tabularx} \\ \\ \relax
2:5 & \begin{tabularx}{0.7\textwidth}{X} 那時，有從天下各國來的虔誠的猶太人，住在耶路撒冷。 \end{tabularx} \\ \\ \relax
2:6 & \begin{tabularx}{0.7\textwidth}{X} 這聲音一響，許多人都來聚集，各人因為聽見門徒用他們各自的鄉談說話，就甚納悶， \end{tabularx} \\ \\ \relax
2:7 & \begin{tabularx}{0.7\textwidth}{X} 都詫異驚奇說：「看哪，這些說話的不都是加利利人嗎？ \end{tabularx} \\ \\ \relax
2:8 & \begin{tabularx}{0.7\textwidth}{X} 我們每個人怎麼聽見他們說我們生來所用的鄉談呢？ \end{tabularx} \\ \\ \relax
2:9 & \begin{tabularx}{0.7\textwidth}{X} 我們帕提亞人、瑪代人、以攔人，和住在美索不達米亞、猶太、加帕多家、本都、亞細亞、 \end{tabularx} \\ \\ \relax
2:10 & \begin{tabularx}{0.7\textwidth}{X} 弗呂家、旁非利亞、埃及的人，並靠近古利奈的利比亞一帶地方的人，僑居的羅馬人， \end{tabularx} \\ \\ \relax
2:11 & \begin{tabularx}{0.7\textwidth}{X} 包括猶太人和皈依猶太教的人， 克里特人和阿拉伯人，都聽見他們用我們的鄉談講論神的大作為。」 \end{tabularx} \\ \\ \relax
2:12 & \begin{tabularx}{0.7\textwidth}{X} 眾人就都驚奇困惑，彼此說：「這是甚麼意思呢？」 \end{tabularx} \\ \\ \relax
2:13 & \begin{tabularx}{0.7\textwidth}{X} 還有人譏誚，說：「他們是灌滿了新酒吧！」 \end{tabularx} \\ \\ \relax
2:14 & \begin{tabularx}{0.7\textwidth}{X} 彼得和十一個使徒站起來，他就高聲向眾人說：「猶太人和所有住在耶路撒冷的人哪，這件事你們要知道，要側耳聽我的話。 \end{tabularx} \\ \\ \relax
2:15 & \begin{tabularx}{0.7\textwidth}{X} 這些人並不像你們所想的喝醉了，因為現在才早晨九點鐘。 \end{tabularx} \\ \\ \relax
2:16 & \begin{tabularx}{0.7\textwidth}{X} 這正是藉著先知約珥所說的： \end{tabularx} \\ \\ \relax
2:17 & \begin{tabularx}{0.7\textwidth}{X} 『神說：在末後的日子，我要將我的靈澆灌凡血肉之軀的。你們的兒女要說預言；你們的少年要見異象；你們的老人要做異夢。 \end{tabularx} \\ \\ \relax
2:18 & \begin{tabularx}{0.7\textwidth}{X} 在那些日子，我要把我的靈澆灌，甚至給我的僕人和婢女，他們要說預言。 \end{tabularx} \\ \\ \relax
2:19 & \begin{tabularx}{0.7\textwidth}{X} 在天上，我要顯出奇事，在地下，我要顯出神蹟，有血，有火，有煙霧。 \end{tabularx} \\ \\ \relax
2:20 & \begin{tabularx}{0.7\textwidth}{X} 太陽要變為黑暗，月亮要變為血，這都在主大而光榮的日子未到以前。 \end{tabularx} \\ \\ \relax
2:21 & \begin{tabularx}{0.7\textwidth}{X} 那時，凡求告主名的都必得救。』 \end{tabularx} \\ \\
[1ex]
\hline
\hline
\end{longtable}
講題:從五旬節到聖靈降臨　Pentecost and the Coming of the Holy Spirit

講員:林豪恩先生

經文:使徒行傳2:1-21

場合:香港中文大學崇基學院禮拜堂主日崇拜

日期:2021年5月23日

引言

事有湊巧,小弟已經第四次在聖靈降臨日講道.聖誕節的經文尚有輪替,因為有四卷福音書,而聖靈降臨日的經文一定是使徒行傳2章1-21節.如果不重複之前三次的內容,還有甚麼可以講呢?在思索的時候,看見聖靈降臨日的中英文名稱原來是不同的,中文稱為聖靈降臨日,英文稱為Day of Pentecost,第五十日.第五十日有兩種意義,其一是指耶穌基督復活後第五十日,其二是使徒行傳第2章所說的五旬節.

2:1五旬節到了,門徒都聚集在一處.

2:2忽然,從天上有響聲下來,好像一陣大風吹過,充滿了他們所坐的屋子,

2:3又有舌頭如火焰顯現出來,分開落在他們各人頭上.

2:4他們就都被聖靈充滿,按著聖靈所賜的口才說起別國的話來.

2章1節提到「五旬節」,即是Day of Pentecost,是事件的發生日期,接著2-4節描述的是事件內容,即當時聖靈來臨在門徒身上的狀況.由此可見,英文名稱Day of Pentecost標示日子,而中文名稱「聖靈降臨日」提示事件.那麼,講題名稱「從五旬節到聖靈降臨日」豈不是誤導嗎?因為,「從五旬節到聖靈降臨日」暗示著距離,其實聖靈在五旬節降臨,因此聖靈降臨日與Day of Pentecost是沒有分別的.聖靈降臨日與五旬節在時間上的確沒有差距,但又不等於在意義上沒有分別.因為五旬節是一個存在已久的猶太人節日,有其本來的起源,意義及過節行為.對於基督徒來說,五旬節的意義就是聖靈降臨,但對猶太教徒來說,五旬節仍然是五旬節.不過,兩者卻有相同之處,都是一個歷史流傳下來的節日.如何活化歷史流傳下來的節日,使之在當下仍有意義,這是值得我們互相參考和彼此啟發的.今天,讓我們以五旬節這個節日為背景,以猶太教朋友如何在非獻祭的狀況中活化五旬節為參考,閱讀聖靈降臨主日這幾段經課,思考聖靈降臨日的意義.

一．因聖靈的創造及維繫而感恩(詩篇104)

其中一段五旬節起源的經文是申命記16章9-11節:

9「你要計算七個七日:從你用鐮刀開始收割莊稼時算起,一共七個七日.

10你要向耶和華－你的神守七七節,按照耶和華－你神所賜你的福,獻上你手裡的甘心祭.

11你和你的兒女,僕婢,以及住在你城裡的利未人,在你中間寄居的和孤兒寡婦,都要在那裡,耶和華－你神選擇作為他名居所的地方,在耶和華－你神面前歡樂.

經過七七四十九日,猶太人在第五十天慶祝「五旬節」(「所以又稱為七七節」).在農業社會,五旬節是為了歡慶收成.生活於非農業社會的猶太教朋友如何記念五旬節呢?小麥,大麥,無花果,棗椰,葡萄,石榴這些農作物雖然不是自己親手種植,但無論甚麼時代,這些農作物仍是他們的主要食糧.在五旬節的時候,帶著這些農作物,或者這些農作物的製成品,例如麵餅去過節,仍是可行及有意義的.有人說,舉行安息禮拜的目的並不是要幫助離去者得享安息,因為上主已經讓他們得享安息,反而是幫助在生者得平安和盼望；猶太教的朋友帶這些農作物去過五旬節獻給上帝,也不是要供應食物給上主,而是確認食物都是從上主而來,提醒自己培養感恩的心.獻上,其實不是付出,而是為得到的而感恩.

今天,基督徒過聖靈降臨節的時候,是不是要重覆當年使徒的經歷才算是有意思呢?或者有人有如此的期望,也有如此的經歷.不過,沒有如此經歷的,是否就沒有意義呢?

正如今時今日不少猶人教朋友不是農民,但仍然可以藉著五旬節向上帝獻上農作物或製成品表達確認和感恩,沒有當年的使徒在那一天的經歷的基督徒,可否也藉著節日表達對聖靈的確認和感恩呢?今天的詩篇經課可以給我們啟發:

104:24耶和華啊,你所造的何其多!都是你用智慧造成的；遍地滿了你的豐富.

104:25那裡有海,又大又廣；其中有無數的動物,大小活物都有.

104:26那裡有船行走,有你所造的鱷魚游泳在其中.

104:27這都仰望你按時給牠食物.

104:28你給牠們,牠們便拾起來；你張手,牠們飽得美食.

104:29你掩面,牠們便驚惶；你收回牠們的氣,牠們就死亡,歸於塵土.

104:30你發出你的靈,牠們便受造；你使地面更換為新.

104:31願耶和華的榮耀存到永遠!願耶和華喜悅自己所造的!

104:32他看地,地便震動；他摸山,山就冒煙.

104:33我要一生向耶和華唱詩!我還活的時候,要向我上帝歌頌!

104:34願他以我的默念為甘甜!我要因耶和華歡喜!

詩人把我們對聖靈的認識從使徒行傳事件帶回創造的起始點,104:30「你發出你的靈,牠們便受造；你使地面更換為新」,這句詩令人們想起創世記第一章第二節說:「上帝的靈運行在水面上」,萬物的出現都和聖靈有關.在聖餐禮中我們這樣祈禱:「最恩慈的上主,我們感謝你,你賜給我們美麗的大地和海洋,豐裕的高山,平原和江河,歌唱的小鳥,可愛的花朵；為這些美好的賞賜我們讚美你」.這祈禱回應著詩104:24「耶和華啊,你所造的何其多!都是你用智慧造成的；遍地滿了你的豐富.」在聖靈降臨日,讓我們再一次確認和感謝聖靈在創造中的參與,使我們得享這一切豐富.

詩104:29詩人又說:「你掩面,牠們便驚惶；你收回牠們的氣,牠們就死亡,歸於塵土.」「你收回牠們的氣」裡面的「氣」字,就是創世記「上帝的靈運行在水面上」的「靈」字.原來,聖靈不但參與創造萬物,其實萬物的存留也在乎聖靈.當我們仍看見萬物生機勃勃,就知道聖靈仍然運行,聽到小鳥的歌唱,就提示我們為聖靈對萬物的保存和護衛而感恩.

二．因聖靈引導我們進入真理而安心(約翰福音)

聖殿仍在的時候,猶太教朋友非常重視獻祭.聖殿被毀後,五旬節失去了它大半的重要性與實際意義,因為再也沒有任何地方可以讓百姓獻祭,況且許多的猶太人也從此脫離農耕,遷移到城鎮生活,他們如何有意義地過五旬節呢?

眼見這個《希伯來聖經》明訂的節日愈來愈受到忽視,為了再次賦予它意義,好讓百姓忠心奉行,拉比賢哲們於是教導,五旬節不只是一個農業節期.拉比賢哲們根據《希伯來聖經》的記載,運用「匠心獨具的算數法」,推定從以色列人出埃及直至上帝於西乃山敵示律法為止,剛巧過了七星期,藉此將五旬節與上帝於西乃山頒賜《妥拉》的顯聖歷史連在一塊.也因此,五旬節擁有一個別名,叫作「賞賜我們《妥拉》的時刻」.因著這個新含義,五旬節重新取得在猶太生活中的重要地位.畢竟,對猶太人和猶太信仰來說,有什麼比頒賜及領受《妥拉》來得更重要?拉比賢哲們表示,五旬節所紀念的,是一件發生在猶太信仰史及世界歷史中的劃時代事件.

從使徒行傳的記載中,我們看見聖靈降臨在使徒身上而賜予給使徒施行神蹟奇事的能力,今天的福音經課讓我們看見聖靈降臨的另一重意義.

15:26但我要從父那裡差保惠師來,就是從父出來真理的聖靈；他來了,就要為我作見證.

15:27你們也要作見證,因為你們從起頭就與我同在.」

「我起先沒有將這事告訴你們,因為我與你們同在.

16:5現今我往差我來的父那裡去,你們中間並沒有人問我:『你往哪裡去?』

16:6只因我將這事告訴你們,你們就滿心憂愁.

16:7然而,我將真情告訴你們,我去是與你們有益的；我若不去,保惠師就不到你們這裡來；我若去,就差他來.

16:8他既來了,就要叫世人為罪,為義,為審判,自己責備自己.

16:9為罪,是因他們不信我；

16:10為義,是因我往父那裡去,你們就不再見我；

16:11為審判,是因這世界的王受了審判.

16:12「我還有好些事要告訴你們,但你們現在擔當不了.

16:13只等真理的聖靈來了,他要引導你們明白一切的真理；因為他不是憑自己說的,乃是把他所聽見的都說出來,並要把將來的事告訴你們.

16:14他要榮耀我,因為他要將受於我的告訴你們.

16:15凡父所有的,都是我的；所以我說,他要將受於我的告訴你們.」

耶穌對聖靈的介紹不在於賜予人行神蹟奇事的能力,而在於對人的轉化和啟迪.「16:8他既來了,就要叫世人為罪,為義,為審判,自己責備自己.」能夠叫世人承認自己做錯而責備自己,這也可算是神蹟.學生某科目考試成績不理想,通常都會解釋說:「是因為那位老師教得不好.」某人向伴侶發了脾氣,通常都會解釋說:「是因為對方激怒我.」如果找不到一個具體的推諉對象,就用最萬能的解釋:「這都是社會的錯.」無論如何,結論都是:「錯不在我.」

不肯認錯的人是不是邪惡的人呢?是不是自私的人呢?或者是,但社會心理學家認為,這只是人類的一種普遍狀況.社會心理學家艾略特阿倫森(Elliot Aronson)和卡羅爾塔夫里斯(Carol Tavris)合著一本社會心理學的經典著作,名為《錯不在我Mistakes Were Made (but not by me) Why We Justify Foolish Beliefs, Bad Decision, and Hurtful Acts》.在這本書中,作者檢示過大量案例,深入分析了人們「死不認錯」的心理,其實並不一定是邪惡或自私,而是與認知失調有關的.為了回應認知失調的狀況,人們發揮自我辯護機制去為自己的行為進行解釋,以使之與自我認知恢復協調.作者們認為自我辯護不是撒謊,而是自我蒙蔽.撒謊是自己也肯定是錯的,但卻說自己是對的.自我辯護是為自己的行為辯解,並且由衷地相信自己沒有錯.

例如一位一生精明的長者,在騙子的推銷下,買了一張天價保健床墊,最後騙子被兒子揭發了,但那位長者仍然不覺得騙子是壞人.因為她在「我很精明」和「我做了件蠢事」之間產生了「認知失調」,於是她本能地進行自我辯護.先是認為天價床墊必定具有「神奇功能」,後來認為騙子這麼做一定有苦衷……總而言之,她無法接受自己做了蠢事這個事實,她要維護自己「精明」的認知,否則認知失調會讓她非常痛苦.(\url{https://kknews.cc/psychology/ppr6qmz.html})

最近和一個查經小組讀羅馬書,作者使徒保羅用了不少篇幅探討「稱義」這個問題.在討論和思考羅馬書關於稱義的課題的過程中,我們感覺到每一個人都有一種非常強烈的「自我感覺為義」的心理需要.我們都需要證明自己是對的,而不是錯的；是值得被重視的,而不是被輕視的；是值得被接納的,而不是被拒絕的.從這個現象出發,羅馬書提出「稱義」正正回應著人們的具體處境,回應著人們心理上的基本需要.由於「自我感覺為義」是一種基本心理需要,也是一種存在的狀況,總不能以道理說服,也不宜以權威壓服,而是透過對自我坦然的自覺和勇敢地自我接受,就正如保羅自我體會:「7:19我所願意的善,我不去做；我所不願意的惡,我反而去做……24我真苦啊!誰能救我脫離這必死的身體呢?」人類的傾向是自我辯護,如保羅這種自我認識和坦然承認如何可能呢?保羅在林前說:「若不是被聖靈感動的,也沒有能說耶穌是主的.」(12:3)能夠在上主面前自我坦誠,實在有賴聖靈的幫助.

16:13只等真理的聖靈來了,他要引導你們明白一切的真理.

「明白一切的真理」不單單是理性上的認識,和修版譯為「進入一切真理」.那個字的意思的確是「進入」,也可以譯為「朝向」,「在…..裡」.「明白」這個字令人感覺是靠自己的聰明智慧和努力,「進入」則令我們想到那位引領者,聖靈就是那位引領者,引領我們進入耶穌基督所要啟示我們的.很多信徒都樂意盡最大的努力,加以「腦力」去閱讀聖經,以求明白真理；耶穌說,若要進入真理,聖靈這位導師的引導必不可少.

四．因聖靈的契合而親密(羅馬書)

拉比賢哲與猶太神秘主義者再進一步把五旬節演繹為一個上帝與猶太人「完婚」的永恆委身時刻.他們說上帝與他的百姓在逾越節時彼此誓約,然後在「俄梅珥數算期」談戀愛與訂婚,讓彼此間的信任與愛情更深刻與強烈.所以「俄梅珥數算期」不再只是用來數日子等收割,而是也代表了一段銘感於心,今人雀躍期待的時期.在這段期間裡,上帝與猶太人的關係快速地朝顛峰邁進,亦即朝向西乃山的婚禮.在這場婚禮中,上帝將《妥拉》這「結婚信物」賜給了猶太人.(以斯拉培訓網站)

聖靈降臨日的的書信經課,保羅在羅馬書也描述聖靈與信徒這種親密的契合:

8:22我們知道,一切受造之物一同歎息,勞苦,直到如今.

8:23不但如此,就是我們這有聖靈初結果子的,也是自己心裡歎息,等候得著兒子的名分,乃是我們的身體得贖.

8:24我們得救是在乎盼望；只是所見的盼望不是盼望,誰還盼望他所見的呢?

8:25但我們若盼望那所不見的,就必忍耐等候.

8:26況且,我們的軟弱有聖靈幫助；我們本不曉得當怎樣禱告,只是聖靈親自用說不出來的歎息替我們禱告.

8:27鑒察人心的,曉得勝靈的意思,因為聖靈照著上帝的旨意替聖徒祈求.

上次講道提及一位患癌的信徒的故事,崇拜結束後有朋友來告訴我他聽的時候心裡有迴響,因為他也是一位癌症患者.他憶述當初得知患病時,心情非常混亂,想到孩子仍然年幼,悲哀徬徨.走到聖壇前,在十字架之下,欲語淚先流.不知道你有沒有經歷過「欲語淚先流」的狀況呢?如保羅所說:「不曉得當怎樣禱告」.沒有言語能夠表達得到那一刻的心情,甚至也不知道要祈求甚麼.保羅的體會,正如那位朋友的經歷是:「聖靈親自用說不出來的歎息替我們禱告」.

在一個討論天人合一的哲學節目中,有教授指出中國哲學的精神是追求天人合一的境界,而基督教則沒有這種精神.他認為基督教的上帝是一位高高在上的越級存在者,而人類是有限的存在者,無限與有限兩者對立,人的位置只是遵守戒律,舊約和新約就是這些戒律.這可能是不少人對基督教神人關係的理解,然而,這卻不是基督教神人關係的全部.從今天的經課耶穌在約翰福音所講的,以及保羅在羅馬書所講的,無限的神和有限的人之間似乎並不是完然沒有互通的可能,類似中國哲人關於天人合一想像,基督徒有沒有神人契合的境界呢?透過聖靈,這是可能的,讓我們繼續在信仰的路上體會及交流.



\newpage

\section{上主的臨在}
\label{sec:552}
\textbf{蒲錦昌牧師 - 上主的臨在}
\newline
\newline
\hyperref[sec:551]{< < < PREV SERMON < < <}
~
\hyperref[sec:toc]{[返目錄]}
~
\hyperref[sec:553]{> > > NEXT SERMON > > >}


以賽亞書 6:1-8
\newline
\begin{longtable}{cl}
\hline
\hline
章節 & 經文 (和合本修訂版)\\
\hline
6:1 & \begin{tabularx}{0.7\textwidth}{X} 當烏西雅王崩的那年，我看見主坐在高高的寶座上。他的衣裳下襬遮滿聖殿。 \end{tabularx} \\ \\ \relax
6:2 & \begin{tabularx}{0.7\textwidth}{X} 上有撒拉弗侍立，各有六個翅膀：兩個翅膀遮臉，兩個翅膀遮腳，兩個翅膀飛翔， \end{tabularx} \\ \\ \relax
6:3 & \begin{tabularx}{0.7\textwidth}{X} 彼此呼喊說：「聖哉！聖哉！聖哉！萬軍之耶和華；他的榮光遍滿全地！」 \end{tabularx} \\ \\ \relax
6:4 & \begin{tabularx}{0.7\textwidth}{X} 因呼喊者的聲音，門檻的根基震動，殿裡充滿了煙雲。 \end{tabularx} \\ \\ \relax
6:5 & \begin{tabularx}{0.7\textwidth}{X} 那時我說：「禍哉！我滅亡了！因為我是嘴唇不潔的人，住在嘴唇不潔的民中，又因我親眼看見大君王－萬軍之耶和華。」 \end{tabularx} \\ \\ \relax
6:6 & \begin{tabularx}{0.7\textwidth}{X} 有一撒拉弗向我飛來，手裡拿著燒紅的炭，是用火鉗從壇上取下來的， \end{tabularx} \\ \\ \relax
6:7 & \begin{tabularx}{0.7\textwidth}{X} 用炭沾我的口，說：「看哪，這炭沾了你的嘴唇，你的罪孽便除掉，你的罪惡就赦免了。」 \end{tabularx} \\ \\ \relax
6:8 & \begin{tabularx}{0.7\textwidth}{X} 我聽見主的聲音說：「我可以差遣誰呢？誰肯為我們去呢？」我說：「我在這裡，請差遣我！」 \end{tabularx} \\ \\
[1ex]
\hline
\hline
\end{longtable}
詩篇 29:1-11
\newline
\begin{longtable}{cl}
\hline
\hline
章節 & 經文 (和合本修訂版)\\
\hline
29:1 & \begin{tabularx}{0.7\textwidth}{X} 神的子民哪，你們要將榮耀、能力歸給耶和華，都歸給耶和華！ \end{tabularx} \\ \\ \relax
29:2 & \begin{tabularx}{0.7\textwidth}{X} 要將耶和華的名的榮耀歸給他，要敬拜神聖榮耀的耶和華。 \end{tabularx} \\ \\ \relax
29:3 & \begin{tabularx}{0.7\textwidth}{X} 耶和華的聲音在眾水上，榮耀的神打雷；耶和華打雷在大水之上。 \end{tabularx} \\ \\ \relax
29:4 & \begin{tabularx}{0.7\textwidth}{X} 耶和華的聲音大有能力，耶和華的聲音滿有威嚴。 \end{tabularx} \\ \\ \relax
29:5 & \begin{tabularx}{0.7\textwidth}{X} 耶和華的聲音震碎香柏樹，耶和華震碎黎巴嫩的香柏樹。 \end{tabularx} \\ \\ \relax
29:6 & \begin{tabularx}{0.7\textwidth}{X} 他使黎巴嫩跳躍如牛犢，使西連跳躍如野牛犢。 \end{tabularx} \\ \\ \relax
29:7 & \begin{tabularx}{0.7\textwidth}{X} 耶和華的聲音使火焰分岔。 \end{tabularx} \\ \\ \relax
29:8 & \begin{tabularx}{0.7\textwidth}{X} 耶和華的聲音震動曠野，耶和華震動加低斯的曠野。 \end{tabularx} \\ \\ \relax
29:9 & \begin{tabularx}{0.7\textwidth}{X} 耶和華的聲音驚動母鹿落胎，樹林也脫落淨光。凡在他殿中的，都述說他的榮耀。 \end{tabularx} \\ \\ \relax
29:10 & \begin{tabularx}{0.7\textwidth}{X} 耶和華坐在洪水之上為王；耶和華坐著為王，直到永遠。 \end{tabularx} \\ \\ \relax
29:11 & \begin{tabularx}{0.7\textwidth}{X} 耶和華必賜力量給他的百姓，耶和華必賜平安的福給他的百姓。 \end{tabularx} \\ \\
[1ex]
\hline
\hline
\end{longtable}
羅馬書 8:12-17
\newline
\begin{longtable}{cl}
\hline
\hline
章節 & 經文 (和合本修訂版)\\
\hline
8:12 & \begin{tabularx}{0.7\textwidth}{X} 弟兄們，這樣看來，我們不是欠肉體的債去順從肉體而活。 \end{tabularx} \\ \\ \relax
8:13 & \begin{tabularx}{0.7\textwidth}{X} 你們若順從肉體活著，必定會死；若靠著聖靈把身體的惡行處死，就必存活。 \end{tabularx} \\ \\ \relax
8:14 & \begin{tabularx}{0.7\textwidth}{X} 因為凡被神的靈引導的都是神的兒子。 \end{tabularx} \\ \\ \relax
8:15 & \begin{tabularx}{0.7\textwidth}{X} 你們所領受的不是奴僕的靈，仍舊害怕；所領受的是兒子名分的靈，因此我們呼叫：「阿爸，父！」 \end{tabularx} \\ \\ \relax
8:16 & \begin{tabularx}{0.7\textwidth}{X} 聖靈自己與我們的靈一同見證我們是神的兒女。 \end{tabularx} \\ \\ \relax
8:17 & \begin{tabularx}{0.7\textwidth}{X} 若是兒女，就是後嗣，是神的後嗣，和基督同作後嗣。如果我們和他一同受苦，是要我們和他一同得榮耀。 \end{tabularx} \\ \\
[1ex]
\hline
\hline
\end{longtable}
約翰福音 3:1-17
\newline
\begin{longtable}{cl}
\hline
\hline
章節 & 經文 (和合本修訂版)\\
\hline
3:1 & \begin{tabularx}{0.7\textwidth}{X} 有一個法利賽人，名叫尼哥德慕，是猶太人的官。 \end{tabularx} \\ \\ \relax
3:2 & \begin{tabularx}{0.7\textwidth}{X} 這人夜裡來見耶穌，對他說：「拉比，我們知道你是由神那裡來作老師的；因為你所行的神蹟，若沒有神同在，無人能行。」 \end{tabularx} \\ \\ \relax
3:3 & \begin{tabularx}{0.7\textwidth}{X} 耶穌回答他說：「我實實在在地告訴你，人若不重生，就不能見神的國。」 \end{tabularx} \\ \\ \relax
3:4 & \begin{tabularx}{0.7\textwidth}{X} 尼哥德慕對他說：「人已經老了，如何能重生呢？豈能再進母腹生出來嗎？」 \end{tabularx} \\ \\ \relax
3:5 & \begin{tabularx}{0.7\textwidth}{X} 耶穌回答：「我實實在在地告訴你，人若不是從水和聖靈生的，就不能進神的國。 \end{tabularx} \\ \\ \relax
3:6 & \begin{tabularx}{0.7\textwidth}{X} 從肉身生的就是肉身；從靈生的就是靈。 \end{tabularx} \\ \\ \relax
3:7 & \begin{tabularx}{0.7\textwidth}{X} 我說『你們必須重生』，你不要驚訝。 \end{tabularx} \\ \\ \relax
3:8 & \begin{tabularx}{0.7\textwidth}{X} 風隨著意思吹，你聽見風的聲音，卻不知道是從哪裡來，往哪裡去；凡從聖靈生的也是如此。」 \end{tabularx} \\ \\ \relax
3:9 & \begin{tabularx}{0.7\textwidth}{X} 尼哥德慕問他：「怎麼能有這些事呢？」 \end{tabularx} \\ \\ \relax
3:10 & \begin{tabularx}{0.7\textwidth}{X} 耶穌回答，對他說：「你是以色列人的老師，還不明白這些事嗎？ \end{tabularx} \\ \\ \relax
3:11 & \begin{tabularx}{0.7\textwidth}{X} 我實實在在地告訴你，我們所說的是我們知道的，我們所見證的是我們見過的，你們卻不領受我們的見證。 \end{tabularx} \\ \\ \relax
3:12 & \begin{tabularx}{0.7\textwidth}{X} 我對你們說地上的事，你們尚且不信，若對你們說天上的事，如何能信呢？ \end{tabularx} \\ \\ \relax
3:13 & \begin{tabularx}{0.7\textwidth}{X} 除了從天降下的人子，沒有人升過天。 \end{tabularx} \\ \\ \relax
3:14 & \begin{tabularx}{0.7\textwidth}{X} 摩西在曠野怎樣舉蛇，人子也必須照樣被舉起來， \end{tabularx} \\ \\ \relax
3:15 & \begin{tabularx}{0.7\textwidth}{X} 要使一切信他的人都得永生。 \end{tabularx} \\ \\ \relax
3:16 & \begin{tabularx}{0.7\textwidth}{X} 「神愛世人，甚至將他獨一的兒子賜給他們，叫一切信他的人不致滅亡，反得永生。 \end{tabularx} \\ \\ \relax
3:17 & \begin{tabularx}{0.7\textwidth}{X} 因為神差他的兒子到世上來，不是要定世人的罪，而是要使世人因他得救。 \end{tabularx} \\ \\
[1ex]
\hline
\hline
\end{longtable}
講題:上主的臨在　The Presence of God

講員:蒲錦昌牧師

經文:以賽亞書6:1-8;詩篇29:1-11;羅馬書8:12-17;約翰福音3:1-17

場合:香港中文大學崇基學院禮拜堂主日崇拜

日期:2021年5月30日

一.引言:

各位崇基禮拜堂的弟兄姊妹,平安.

自從新型冠狀病毒疫情以來,我們都很珍惜能夠面對面一起崇拜的日子,也感到向其他人問安,說一聲平安是何等實在而又幸福的事情.因為生命可以轉眼就在我們面前流逝.

我特別希望鼓勵仍然以視像方式參與崇拜的弟兄姊妹,如果健康及環境許可的話,盡量親身的參與崇拜,與視像方式仍然是有分別的.崇拜是個神聖時間和空間,我們在一星期中特定的時間到特定的地方一起崇拜,這個經驗是不可代替的.大家試想想,坐在家裡崇拜,有時候難免分心,也沒理由穿著整齊才在鏡頭前面崇拜,而且會遷就自己的時間.這些都和實體崇拜不同.還有,由準備離家至到達教會崇拜,也是靈性生活一部分.大家試想一下,一步一步走上黃山看黃山與坐登山纜車看黃山,看到的黃山也會不同,因為一步一步走上黃山會花很多氣力,結果這些氣力也成為體驗的一部分.崇拜和靈性生活也一樣.沒有付出的信仰不是信仰,只是資本主義社會下提供的一種娛樂或產品.希望大家都多參與教會的實體崇拜,不在港和抱恙的自然例外.

好了.回到今主日崇拜的主題:今天是「聖靈降臨後第一主日」,也是教會年曆的「三一主日」.我們讀了四段經課,分別是賽6:1-8,詩29:1-11,羅8:12-17及約3:1-17.我會用這四段經文來與大家分享,講題是「上主的臨在」.

二.主題:

1.呼召的上主

第一段的經文是賽6:1-8:記載上主呼召以賽亞作先知--上主的代言人.聖經中有許多上主呼召人的記載,包括呼召亞伯拉罕(創12:1-3),摩西(出3-4:17). 約書亞(書1:1-9). 耶利米(耶1:3-10)等.賽6:1-8(其實到第13節才結束)是其中經典又熟悉的一段.

上主以君王之身分. 尊貴和榮耀向以賽亞顯現,我們現代人稱之為超自然的現象.6:1-4以賽亞見到上主坐在高高的寶座上,周圍有天使相伴,而且高聲呼喊:「聖哉!聖哉!聖哉!萬軍之耶和華；他的榮光遍滿全地!」除了呼喊的聲音,伴隨的是地震和殿裡滿佈的煙雲,就如摩西在西乃山上所見的:「西奈山全山冒煙,因為耶和華在火中降臨山上.山的煙霧上騰,彷彿燒窯,整座山劇烈地震動.」(出19:18)以賽亞見到的這個異象與日常生活見到的東西截然不同,是一種超越日常的經驗,是一種神秘的經驗.對於親身經驗的人而言,無比深刻,甚至影響和定義自己的一生,但又難以向他人傳述,更無法複製.除了這個超常的情境以外,跟著是以賽亞的回應和苦惱:「禍哉!我滅亡了!因為我是嘴唇不潔的人,住在嘴唇不潔的民中,又因我親眼看見大君王－萬軍之耶和華.」(賽6:5)凡體驗過上主臨在的人,他們都有很清晰而強烈的感覺,就是上主的聖潔和自己的不潔.這是對神聖的觀念最核心的部份.不單如此,以賽亞也指向他所身處的社會--一個不潔的社會.之後,便是潔淨的禮儀:天使用燒紅的炭潔淨以賽亞的嘴唇,除去他的罪,使他成為聖潔,預備好來服侍上主.(賽6:6-7).然後,上主便向以賽亞發出呼召,但是以一個問題的方式:「我可以差遣誰呢?誰肯為我們去呢?」以賽亞便爽快的回應:「我在這裡,請差遣我!」(賽6:8)這一節經文也成為了許多基督徒受到上主呼召的印證.上主藉這段經文一次又一次呼喚人來回應他. 服侍他.

各位弟兄姊妹,我們可能也有上主臨在我們生命的經驗,祂的聖潔映照出我們的不潔.上主潔淨我們的生命,使我們如沐春風,使我們再次感到輕省,不用再承受罪的重壓,使我們得以親近聖潔的上主,與主契合,像以賽亞一樣.不過,我們要謹記:不要停留在這個階段,像彼得於耶穌登山變像後想停留在山上的美好時光一樣(參路9:28-36),而要下山服侍世人.因為,上主是呼召之主,也是使命之主.祂呼召我們,差遣我們,使我們成為上主的使者和代言人,完成祂交託我們的使命,將真理和生命帶給世界上每一個需要真理和真實生命的人.

如果你是一個教師,你會如何完成上主交給你的使命呢?只需要好好教導學生,懂得判斷是非,分辨真假,做個善良的人.正如耶穌所教導的:「是,就說是；不是,就說不是.若再多說,就是出於那惡者.」(太5:37)教歷史的老師,不可以把發生過的事情當沒發生,把沒發生的事情當發生了般教學生.如果你是一個律師,你如何見證上主呢?很簡單:持守專業守則,盡力為委托你的人辯護,使他/她得到公平的審訊,正直的人不會被屈.如果你是醫生護士,你又如何服侍神呢?以最專業的精神和技術,救治病人,不考慮他/她是否富有或貧窮,善良或邪惡.做在最微小的一個人身上,就是做在主身上.這些原則,不論世界變得如何,都要堅守.

2.威榮的上主

第二段的經文是詩29篇:一首讚美上主威嚴. 榮耀和能力的詩,祂掌管整個宇宙穹蒼.這首詩分為三個部分:1,2節是呼籲人讚美上主,3-10節是讚美的內容,11節則是結束祈禱.

以下,讓我們仔細看看這篇詩篇:

「上帝的子民哪,你們要將榮耀,能力歸給耶和華,歸給耶和華!

要將耶和華的名的榮耀歸給他,要敬拜神聖榮耀的耶和華.」(詩29:1-2)呼籲以色列人讚頌上主的榮耀. 能力和神聖.

3至4節則敘述上主的能力:「耶和華的聲音在眾水上,榮耀的上帝打雷；耶和華打雷在大水之上.耶和華的聲音大有能力；耶和華的聲音滿有威嚴.」(詩29:3-4)上主的榮耀和能力不單只是在海洋湖泊或河流上打雷,而是一種在創世時運行在水面上的力量(創1:2),使象徵毀滅和死亡的洪水和海(參創6-7,啟21:1)都屈服的威力.

5至9節講到這種力量所產生的作用:

「耶和華的聲音震碎香柏樹；耶和華震碎黎巴嫩的香柏樹.

他使黎巴嫩跳躍如牛犢,使西連跳躍如野牛犢.

耶和華的聲音使火焰分岔.

耶和華的聲音震動曠野,耶和華震動加低斯的曠野.

耶和華的聲音驚動母鹿落胎,樹林也脫落淨光.凡在他殿中的,都稱說他的榮耀.」

閃電轟雷使黎巴嫩最堅固的香柏樹也震碎,黎巴嫩和西連(即黑門山,見申3:9)這些不動的高山都震動,加低斯的曠野也震動起來,一團火焰都被分開兩半,如地震使地面崩裂,土地分開.連最安靜的母鹿也驚慌落胎,翠綠的樹木片葉不留.這是一種多麼大的力量!

第10節則總結之前所顯示上主的能力和榮耀說:

「 耶和華坐在洪水之上為王；耶和華坐著為王,直到永遠.」他是戰勝洪水,永遠為王的上主,洪水都在他的腳下,動彈不得.這當然只是一個詩的意象,而不是一個事實的描述,因為沒有任何人會坐在洪水之上而不被洪水淹沒或沖走,坐在水上也沒法阻止水的流動和前進.

最後一節(11節):「耶和華必賜力量給他的百姓；耶和華必賜平安的福給他的百姓.」對以色列人的祝福語結束整首詩,相信上主會以其能力保護和照顧以色列人--他的子民.

我們今天看世界萬物為「自然界」,為有別於上主的被造物和客觀的存在,風雷雨電. 海嘯地震等都是自然現象,沒什麼希奇.舊約時代的世界和宇宙觀完全不一樣.這篇詩篇便用了象徵性的語言,描述上主在打雷中顯示自己,臨在於世界中.

我們在城市生活,狂風暴雨時躲在家中,感到很安全,在地震帶的房屋不少也有防震設備,不會倒塌.不過,如果我們走出城市,在冰封的雪山. 無垠的沙漠. 無邊無際的海洋. 衝破堤圍的巨浪和洪水. 突如奇來的海嘯和地震,火山噴薄而出的火燄,震耳欲聾的雷聲等,我們仍然會感覺到恐懼. 渺小和脆弱,也是更深刻認識上主的時刻.教會改革時的巨人馬丁路德,就是被打雷的聲音嚇倒而降服於上主腳下.

十年前,日本福島311大地震,我們在電視畫面上見到地震和海嘯的畫面,令人震驚.311之後一個月,即2011年4月11日,日本基督教協進會的總幹事上田博子牧師,與我們幾位來自南韓. 台灣,香港的教會代表,一同往訪災區,經歷著5.4級的餘震,看到香港人旅遊常去的松島被海水浸泡過後的景象,石卷市最接近海邊的房屋破爛不堪,沒水沒電,也沒有人走動,並且發出難聞的氣味,像個死城.汽車在海嘯中被沖走,有些停留在田地上,也有「飄浮」在幾呎高的矮牆上的.多麼超現實的景象!我們一團人都親眼看見了自然的力量.

我們也會在自然界的力量中看見上主的臨在和威榮嗎?

3.感人的上主

第三段經文是羅8:12-17,是保羅論到基督徒與上主關係的一段說話.保羅說,上主藉他的靈臨在我們生命之中,使我們得以進入上主的家庭,成為他的兒女,可以親切的稱上主為「阿爸,父!」而且與耶穌基督同作上主的後嗣.這是一種愛的關係,而非奴隸對主人般恐懼的關係,因為我們都成為上主的兒女了.

當聖靈的力量在我們內心工作時,他會引導我們,使我們渴慕上主,跟隨上主的指引,過聖潔的生活,並且抗拒上主不喜歡的和邪惡的事情,會按著聖靈的意思行事為人.正如羅8:9上所講:「如果上帝的靈住在你們裡面,你們就不屬肉體,而是屬聖靈了.」

當我們在屬肉體的時候,就像保羅於羅1:29-32所講:「裝滿了各樣不義,邪惡,貪婪,惡毒,滿心是嫉妒,凶殺,紛爭. 詭詐,毒恨,又是毀謗的,說人壞話的,怨恨上帝的,侮辱人的,狂傲的,自誇的,製造是非的,忤逆父母的,頑梗不化的. 言而無信的,無情無義的,不憐憫人的.他們雖知道上帝判定做這樣事的人是該死的,然而他們不但自己去行,還贊同別人去做.」只要我們看看社會上的情況,我們很容易了解和明白人類是多麼的敗壞.最明顯的一個例子,就是許多國家和國家領袖都說愛和平,但卻不斷製造更強的武器和訓練更強的軍隊.如果你問這些國家,為什麼需要這麼多武器和軍隊呢?他們會答:為了保護我的人民.結果便是各國互不信任,並且進行軍備競賽,要成為最強的那一個.這是人類的悲劇,就是假和平之名發動種種戰爭.這樣下去,人類會有真正的和平嗎?不能.只會有越來越強的武器和越來越可怕的戰爭!

我們多麼需要上主的恩典和力量來改造我們舊有的生命!就像今天所讀的第四段經文:約3:1-17,聖靈像風一樣,我們不知道他從那裡來,往那裡去,突如其來,神秘莫測.(在希臘文中,「靈」和「風」是同一個字)他會吹進我們的內心,呼喚我們要改變,要重生,好像夜裡來找耶穌有尊貴身份的尼哥德慕,聖靈就透過耶穌的話,呼喚他要重生,要由黑暗走向光明,捨棄尊貴的身份來跟從耶穌.

當我們願意歸向上主,並且,讓聖靈住進我們的內心,所產生的結果,就像保羅在加5:22-23所講:「聖靈的果子就是仁愛,喜樂,和平,忍耐,恩慈,良善,信實,溫柔,節制.這樣的事沒有律法禁止.」我們會經歷主裡的新生,並且有新生的樣式.用保羅在林後5:17-20的話:「所以,若有人在基督裡,他就是新造的人,舊事已過,都變成新的了.一切都是出於上帝；他藉著基督使我們與他和好,又將勸人與他和好的使命賜給我們.這就是:上帝在基督裡使世人與自己和好,不將他們的過犯歸到他們身上,並且將這和好的信息託付了我們.所以,我們作基督的特使,就好像上帝藉我們勸你們一般.我們替基督求你們,與上帝和好吧!」上主臨在我們的生命中,感動我們,使我們與上主和好,和他建立起父親和子女般的關係,同時,也與人和好,建立兄弟姊妹般的關係.這也是一個重生的生命,是聖靈自己的作為.

三.結語:

總括來談,上主是呼召的上主,他向我們顯現,並呼召我們完成他創造及更新這個世界的使命.上主是威榮的上主,他在世界的萬事萬物中顯示他自己,我們透過自然界可以更深刻的認識他的臨在.上主也是感動人的上主,他差遣聖靈進入我們的內心,感動我們,更新我們,使我們成為新人,並肩負和平使者的職份,在人與神和人與人之間締造和平.盼望上主保守和賜福他的教會,使我們成為他合用的器皿!阿們!



\newpage

\section{癲狂的耶穌}
\label{sec:553}
\textbf{高國雄牧師 - 癲狂的耶穌??}
\newline
\newline
\hyperref[sec:552]{< < < PREV SERMON < < <}
~
\hyperref[sec:toc]{[返目錄]}
~
\hyperref[sec:554]{> > > NEXT SERMON > > >}


馬可福音 3:20-35
\newline
\begin{longtable}{cl}
\hline
\hline
章節 & 經文 (和合本修訂版)\\
\hline
3:20 & \begin{tabularx}{0.7\textwidth}{X} 耶穌進了屋子，眾人又聚集，甚至他連飯也顧不得吃。 \end{tabularx} \\ \\ \relax
3:21 & \begin{tabularx}{0.7\textwidth}{X} 耶穌的家人聽見，就出來要拉住他，因為他們說他癲狂了。 \end{tabularx} \\ \\ \relax
3:22 & \begin{tabularx}{0.7\textwidth}{X} 從耶路撒冷下來的文士說：「他是被別西卜附身的」，又說：「他是靠著鬼王趕鬼的。」 \end{tabularx} \\ \\ \relax
3:23 & \begin{tabularx}{0.7\textwidth}{X} 耶穌叫他們來，用比喻對他們說：「撒但怎能趕出撒但呢？ \end{tabularx} \\ \\ \relax
3:24 & \begin{tabularx}{0.7\textwidth}{X} 一國若自相紛爭，那國就立不住； \end{tabularx} \\ \\ \relax
3:25 & \begin{tabularx}{0.7\textwidth}{X} 一家若自相紛爭，那家就立不住。 \end{tabularx} \\ \\ \relax
3:26 & \begin{tabularx}{0.7\textwidth}{X} 撒但若自相攻打紛爭，他就立不住，必定滅亡。 \end{tabularx} \\ \\ \relax
3:27 & \begin{tabularx}{0.7\textwidth}{X} 沒有人能進壯士家裡，搶奪他的東西；除非先綁住那壯士，否則無法搶奪他的家。 \end{tabularx} \\ \\ \relax
3:28 & \begin{tabularx}{0.7\textwidth}{X} 我實在告訴你們，世人一切的罪和一切褻瀆的話都可以得到赦免； \end{tabularx} \\ \\ \relax
3:29 & \begin{tabularx}{0.7\textwidth}{X} 凡褻瀆聖靈的，卻永不得赦免，而要擔當永遠的罪。」 \end{tabularx} \\ \\ \relax
3:30 & \begin{tabularx}{0.7\textwidth}{X} 因為他們說：「他是被污靈附身的。」 \end{tabularx} \\ \\ \relax
3:31 & \begin{tabularx}{0.7\textwidth}{X} 那時，耶穌的母親和他兄弟來，站在外邊，打發人去叫他。 \end{tabularx} \\ \\ \relax
3:32 & \begin{tabularx}{0.7\textwidth}{X} 有許多人在耶穌周圍坐著，他們就告訴他說：「看哪！你母親、你兄弟和你姊妹在外邊找你。」 \end{tabularx} \\ \\ \relax
3:33 & \begin{tabularx}{0.7\textwidth}{X} 耶穌回答他們：「誰是我的母親？誰是我的兄弟？」 \end{tabularx} \\ \\ \relax
3:34 & \begin{tabularx}{0.7\textwidth}{X} 就環視那周圍坐著的人，說：「看哪，我的母親，我的兄弟！ \end{tabularx} \\ \\ \relax
3:35 & \begin{tabularx}{0.7\textwidth}{X} 凡遵行神旨意的人就是我的兄弟姊妹和母親。」 \end{tabularx} \\ \\
[1ex]
\hline
\hline
\end{longtable}
講題:癲狂的耶穌??　Crazy Jesus??

講員:高國雄牧師

經文:馬可福音3:20-35

場合:香港中文大學崇基學院禮拜堂主日崇拜

日期:2021年6月6日

20耶穌進了屋子,眾人又聚集,甚至他連飯也顧不得吃. 21耶穌的家人聽見,就出來要拉住他,因為他們說他癲狂了. 22從耶路撒冷下來的文士說:「他是被別西卜附身的」,又說:「他是靠著鬼王趕鬼的.」 23耶穌叫他們來,用比喻對他們說:「撒旦怎能趕出撒旦呢? 24一國若自相紛爭,那國就立不住； 25一家若自相紛爭,那家就立不住. 26撒旦若自相攻打紛爭,他就立不住,必定滅亡. 27沒有人能進壯士家裡,搶奪他的東西；除非先綁住那壯士,否則無法搶奪他的家. 28我實在告訴你們,世人一切的罪和一切褻瀆的話都可以得到赦免； 29凡褻瀆聖靈的,卻永不得赦免,而要擔當永遠的罪.」 30因為他們說:「他是被污靈附身的.」31那時,耶穌的母親和他兄弟來,站在外邊,打發人去叫他. 32有許多人在耶穌周圍坐著,他們就告訴他說:「看哪!你母親,你兄弟和你姊妹在外邊找你.」 33耶穌回答他們:「誰是我的母親?誰是我的兄弟?」 34就環視那周圍坐著的人,說:「看哪,我的母親,我的兄弟! 35凡遵行上帝旨意的人就是我的兄弟姊妹和母親.」

每一位認識主耶穌基督的人,都會對主的形象有不同的想法,歷代很多不同的藝術與文學的表達,當中不少會以謙和,溫柔為主調；記得曾經睇過一部記述耶穌基督生平的電影,整部電影當中的飾演耶穌的美國白人不單外表英俊不凡,青靚白淨,文質彬彬,加上說話陰聲細氣,雖然明白電影希望讓觀眾對所表現這「耶穌基督的形象」而對基督教產生好感,我當然也認同耶穌有感性及溫柔的一面,但耶套電影全場皆以這種表達方式,實在令我覺得毛骨悚然,非常不自在,與及不斷地反思,到底我心目中耶穌基督的形成又是怎樣的呢?

弟兄姊妹,耶穌基督在我們心裡的形象是怎樣的呢?你心中的耶穌基督又是怎樣的呢?我們是否也一樣有著定型觀念/刻板印象(Stereotype),而無法立體地認識耶穌基督呢?

今天福音書的經課可能為我們展現了一個很不一樣的耶穌基督.

V20耶穌進了屋子,眾人又聚集,甚至他連飯也顧不得吃.

眾人又聚集:明顯地過去曾經聚集過,這聚集是有原因的,從可1:21-28耶穌在迦百農,在安息日趕鬼開始,到醫好西門的岳母,再醫好痳瘋帶人,癱子,手萎縮了的人…短短一章多一點的經文,除了耶穌基督講道教訓以外,已經記述多個的醫猜趕鬼的神蹟,難怪大批群眾「又」聚集在耶穌基督所進的屋內,相信當中充滿著不少身心靈需要的朋友,希望耶穌基督成為他們的幫助與拯救…當然耶穌基督也會因此非常忙碌,經文描述「甚至他連飯也顧不得吃」,原文新約用「甚至使他們不能進餐」

相信我們也可能曾經經歷過因為忙碌而忘記/沒有機會用餐,可能是因為趕功課/工作,可能是沉迷生活(遊玩,睇書,打機…)；無論任何原因,也都會被歸類為不健康,不應該,不正常,不好或不可以.如果你的家人如此,你會有什麼反應呢?關心,擔心,痛心…你會有什麼行動呢?勸喻,責罵,阻止…

或許我們會看其原因而有不同程度的回應…

聖經記述v21

耶穌的家人聽見,就出來要拉住他,因為他們說他癲狂了.

原文新約

屬於他的人(家人)聽到後,出來捉住他說:「他發瘋了!」

我們很明白耶穌基督正在做很有意義的事情,雖然連進餐的機會也沒有,我們完全體會耶穌的辛苦,亦同時體會祂對人的愛,所以這非一般人的選擇,我們也會覺得因為這很有價值而認同；同時我們也明白作為家人的憂心,所以捉住耶穌基督,想叫停,甚至以「發瘋了/癩狂了」來形容也合理,很明顯因為我們對這「發瘋了/癩狂了」理解不是說這人「精神錯亂」,「不可理喻」,「生人勿近」…而是指他的行動其他人不一樣,作出非一般人的選擇(耶穌基督當然是非一般的),結果有點過份,有點太過極致,不過仍然可以被明白及接受；甚或我們覺得這描述只是因為來自家人(比較關心)而誇張了的用語…

所以我們從來沒有耶穌基督是癲狂/發瘋的形象…

經文接著插入一段與耶穌基督身份有關的事:

v22從耶路撒冷下來的文士說:「他是被別西卜附身的」,又說:「他是靠著鬼王趕鬼的.」

別西卜:舊約列王記下1「巴力‧西卜」,非利士人崇拜的神,蒼蠅王.

一個非常嚴厲的指控!

與其他2卷符類福音書記載的不同,馬可福音表示指控耶穌,挑戰祂的權威及身份的,是從耶路撒冷下來的文士；文士是律法專家,他們熟悉傳統律法及各規條,從耶路撒冷下來,不單表達出他們尊貴,有地位的身份(耶路撒冷是當時信仰的核心地),不是一般的,更暗示他們是有目的而來的(下來,刻意衝著耶穌基督而來的)

誣衊耶穌被別西卜附身及靠鬼王趕鬼:一方面希望否定/不願意接受耶穌基督的能力來自他們所信奉的上帝(當然也完全否定耶穌基督就是上帝),同時也否定耶穌基督做的是好事,更是完全對耶穌基督的身份在猶太價值下的質疑/否定…

他們熟悉上帝的律法,甚至可以說是能代表上帝律法及信仰權威,這些來自耶路撒冷的文士對耶穌這指控,必然引來很大的震撼與回響.

面對如此嚴峻的情況,耶穌基督決斷地選擇即時回應,這也是非常值得我們深思及學習的地方,按不同的情況及嚴重程度,對信仰影響重要的事,要即時回應!

v23耶穌叫他們來,用比喻對他們說:「撒旦怎能趕出撒旦呢? 24一國若自相紛爭,那國就立不住； 25一家若自相紛爭,那家就立不住. 26撒旦若自相攻打紛爭,他就立不住,必定滅亡. 27沒有人能進壯士家裡,搶奪他的東西；除非先綁住那壯士,否則無法搶奪他的家. 28我實在告訴你們,世人一切的罪和一切褻瀆的話都可以得到赦免； 29凡褻瀆聖靈的,卻永不得赦免,而要擔當永遠的罪.」 30因為他們說:「他是被污靈附身的.」

對此誣蔑,耶穌基督先以邏輯推論來回應,用家,國作為比喻,領導者與群眾必須同一陣線,否則必定滅亡；因此沒有可以靠鬼的能力而趕鬼(因為有內部矛盾),並且指出,必須有比對方強大的力量,才可以勝過對方,有能力綁住壯士,才可擔奪他家；同樣,要趕鬼/勝過鬼,必須先有比鬼更大的力量,因此只有上帝才可以.

耶穌基督更嚴厲地回應對「耶穌是被污靈附身的.」這指控,清楚表明:

V28我實在告訴你們,世人一切的罪和一切褻瀆的話都可以得到赦免； 29凡褻瀆聖靈的,卻永不得赦免,而要擔當永遠的罪.」

馬太福音更記載:12:31-32

所以我告訴你們:「人一切的罪和褻瀆的話都可得赦免,惟獨褻瀆聖靈,總不得赦免.凡說話干犯人子的,還可得赦免；惟獨說話干犯聖靈的,今世來世總不得赦免.」

弟兄姊妹,對待耶穌基督這教導,我們必須要非常謹慎,干犯聖靈的不得赦免,不是因為上帝不愛我們,也不是祂的救恩有限；文士這種把明顯的神蹟,大家也能見證到是上帝的能力,不單不接受,更反過來說謊,誣蔑是由鬼魔而出,絕對是明知並刻意的褻瀆,而這褻瀆的背後就是代表著對上帝的態度與選擇,是不承認及反對,是傲慢與不信,要記得；上帝是輕慢不得的.

所以在這裡耶穌基督不單清晰表達祂的能力來源來不是來自鬼魔,而是能勝過鬼魔,讓鬼魔最終也要降服的上帝.

耶穌基督這超越常人/常理的能力(醫病趕鬼)也可能成為被理解為癲狂的原因,與文士(來自耶路撒冷尊貴的文士)辯論,甚至教訓他們,也可能被視為發瘋…

不過耶穌基督真正的癲狂往往不是我們所能想像得到!

V31那時,耶穌的母親和他兄弟來,站在外邊,打發人去叫他. 32有許多人在耶穌周圍坐著,他們就告訴他說:「看哪!你母親,你兄弟和你姊妹在外邊找你.」 33耶穌回答他們:「誰是我的母親?誰是我的兄弟?」 34就環視那周圍坐著的人,說:「看哪,我的母親,我的兄弟! 35凡遵行上帝旨意的人就是我的兄弟姊妹和母親.」

經文承接開首,可能是因為耶穌基督的家人知道耶穌連飯也顧不得吃的擔心,或者有其他的需要或要事,希望找到耶穌,但明顯地因為屋子內人群眾多,耶穌的母親及弟弟只可以在門外送出訊息,當人把耶穌的母親,兄弟姊妹要尋找耶穌的信息帶給耶穌時,耶穌竟然作出我們從來沒有想像過的回應,反問眾人:

「誰是我的母親?誰是我的兄弟?」

當然我相信耶穌這一刻不是不知道/忘記了誰是他肉身的家人,耶穌更加應該不會是突然失憶或急性腦退化,而是刻意說出這說話…

連家人,母親也不認,耶穌竟然這般「大逆不道」,這一種「癲狂」相信我們絕難想像,更惶論如何學效跟從…

原來對跟從主,認識主的人所要求的不單是熱心參與教會,不只是間中花一點時間及心力去服侍,不止是奉獻有餘的金錢,或者間中做一下些可能暫時世人不明白/覺得是愚蠢的事(點點的犧牲?),而是要對價值觀,人生目標等徹底的改變,上帝從來沒有叫我們忘記/離棄肉身的家庭,當孝敬父母是十誡中唯一帶應該的誡命,非常重要,聖經中亦有非常豐富教導我們如何與家人,夫婦,人倫相處之道等,所以上主看重要一個人及其家庭是無容質疑的；

記得耶穌基督在十字架上對所愛門徒的囑咐嗎?

約翰福音19:26-27耶穌見母親和他所愛的那門徒站在旁邊,就對母親說:「母親,看,你的兒子!」又對那門徒說:「看,你的母親!」從那刻起,那門徒就接她到自己家裡去了.

而耶穌基督這裡所展示的,是要我們知道上帝與我們的連結是要超乎血緣,超乎常理的,猶太人很緊張自己是選民的身份,整個民族獨特地接受主的揀選,自以為有著不可被取代的地位,耶穌基督把整個概念改變…

v35凡遵行上帝旨意的人就是我的兄弟姊妹和母親.

誰是遵行上帝旨意最極至的代表?當然就是道成了肉身,更願意因為對我們的愛而甘願釘死在十字架上的耶穌基督；因此理所當然如若我們希望成為耶穌基督的家人,與祂有最緊密的連結,就不再是依靠身份,地位,甚至血緣,而是願意成為與祂一樣價值,人生方向,盡一生遵行主旨意的人.

請記得使徒保羅曾經如此表白:

哥林多後書5:13如果我們癲狂,是為上帝；如果我們清醒,是為你們.



\newpage

\section{一切從一粒種子開始}
\label{sec:554}
\textbf{馮天聰牧師 - 一切從一粒種子開始}
\newline
\newline
\hyperref[sec:553]{< < < PREV SERMON < < <}
~
\hyperref[sec:toc]{[返目錄]}
~
\hyperref[sec:555]{> > > NEXT SERMON > > >}


馬可福音 4:26-34
\newline
\begin{longtable}{cl}
\hline
\hline
章節 & 經文 (和合本修訂版)\\
\hline
4:26 & \begin{tabularx}{0.7\textwidth}{X} 耶穌又說：「神的國如同人把種子撒在地上， \end{tabularx} \\ \\ \relax
4:27 & \begin{tabularx}{0.7\textwidth}{X} 黑夜睡覺，白日起來，這種子就發芽生長，那人卻不知道如何會這樣。 \end{tabularx} \\ \\ \relax
4:28 & \begin{tabularx}{0.7\textwidth}{X} 土地自然而然地出產五穀，先發苗，後長穗，然後穗上結成飽滿的穀子。 \end{tabularx} \\ \\ \relax
4:29 & \begin{tabularx}{0.7\textwidth}{X} 五穀熟了，就用鐮刀去割，因為收成的時候到了。」 \end{tabularx} \\ \\ \relax
4:30 & \begin{tabularx}{0.7\textwidth}{X} 耶穌又說：「我們可用甚麼來比擬神的國呢？可用甚麼比喻來說明呢？ \end{tabularx} \\ \\ \relax
4:31 & \begin{tabularx}{0.7\textwidth}{X} 它像一粒芥菜種，種在地裡的時候，雖比地上所有的種子都小， \end{tabularx} \\ \\ \relax
4:32 & \begin{tabularx}{0.7\textwidth}{X} 但種下去以後，它長起來，比各樣的菜都大，又長出大枝，以致天上的飛鳥可以在它的蔭下築巢。」 \end{tabularx} \\ \\ \relax
4:33 & \begin{tabularx}{0.7\textwidth}{X} 耶穌用許多這樣的比喻，照他們所能聽的，對他們講道； \end{tabularx} \\ \\ \relax
4:34 & \begin{tabularx}{0.7\textwidth}{X} 若不用比喻，他就不對他們講，但私下沒有人的時候，就把一切的道講給門徒聽。 \end{tabularx} \\ \\
[1ex]
\hline
\hline
\end{longtable}
講題:一切從一粒種子開始　It All Starts With A Seed

講員:馮天聰牧師

經文:馬可福音4:26-34

場合:香港中文大學崇基學院禮拜堂主日崇拜

日期:2021年6月13日

經文:馬可福音四章 26-34 節

又說:「上帝的國如同人把種撒在地上.

黑夜睡覺,白日起來,這種就發芽漸長,那人卻不曉得如何這樣.

地生五穀是出於自然的:先發苗,後長穗,再後穗上結成飽滿的子粒；

穀既熟了,就用鐮刀去割,因為收成的時候到了.」

又說:「上帝的國,我們可用甚麼比較呢?可用甚麼比喻表明呢?

好像一粒芥菜種,種在地裡的時候,雖比地上的百種都小,

但種上以後,就長起來,比各樣的菜都大,又長出大枝來,甚至天上的飛鳥可以宿在它的蔭下.」

耶穌用許多這樣的比喻,照他們所能聽的,對他們講道.

若不用比喻,就不對他們講；沒有人的時候,就把一切的道講給門徒聽.

引言

六月,對於很多香港人來說,是一個令人思潮起伏的月份.

撇開 32 年前 6 月「六四事件」引發香港人思緒的林林總總；說近一點,不論是一年前 6 月通過及實施的國家安全法,還是兩年前 6 月因「反修例風波」而出現一連串大規模的社會運動,都徹底改變了香港這座城市的面貌,亦令香港人的未來從此不再一樣.

想不到,一宗發生在台灣的兇殺案,最後會演變成接近一年的政治風暴,不但影響到香港裡面人和人,群眾與群眾,人和政府之間的關係,還震盪到中國在國際社會中與其他國家之間的關係,餘波到今天仍然未了；更想不到的是,到了上年年初,一棵微小到肉眼看不見,現時仍未知道來源的新型冠狀病毒,令疫情席捲全球,到現今已做成超過 1 億 7 千多萬宗感染個案,超過 370 萬人死亡,帶來的經濟破壞更無從估計,無數人要過著失業或停課的生活.「星星之火,可以燎原」,可說是這兩件事件最好的註腳,一開始的時候,沒有人會意料到一些似乎微不足道的事,會帶來日後巨大的轉變.

按照主日經課的編排,今天我們讀到福音書馬可福音第 4 章 26 至 34 節,在這裡,耶穌一口氣說了兩個比喻,兩個都是跟上帝的國有關的.上帝國一開始看來毫不起眼,但在人想不到的情況下,竟然出現爆炸性的發展,深深影響著世界及改變人的生命.

一. 不曉得如何這樣

我們先看看第一個比喻,一個有關種子長大的比喻,記載在 26 至 29 節.

26 又說:「上帝的國如同人把種撒在地上.27 黑夜睡覺,白日起來,這種就發芽漸長,那人卻不曉得如何這樣.28 地生五穀是出於自然的:先發苗,後長穗,再後穗上結成飽滿的子粒；29 穀既熟了, 就用鐮刀去割,因為收成的時候到了.」

在這個比喻中,耶穌說上帝的國就如同種子一樣,當撒到地上後就開始發芽成長.耶穌要表達的是, 一粒種子的發芽,生長和結實,靠的不只是撒種的人的努力而已,因為這個人連植物怎樣生長也不清楚,撒下的種子將來會怎樣,他也不能確定.一顆種子的發芽,結實,實際上是土壤自然而然地使它生長,結實；使種子生長,發芽的不是撒種的人.是誰讓種子可以順利長大結實累累,及至到了收成期,可以用鐮刀將之收割下來?其實都是上帝在背後的工作.

當然,到了今天,我們可以從植物學的角度去知道一粒種子怎樣發芽生長,又知道當中需要滿足什麼的條件及背後的原理與操作,因此可以提供植物合適的生長環境及一切所需,甚至可以透過基因工程去改造植物及種子,讓植物可以生長得更快更強壯,還可以有以往沒有的特性.但我們仍無法由無變有,不論種子本身還是我們可以提供的環境營養等等,都是上帝創造的；而且在植物生長過程中仍有不少的不確定性,人不能完全控制.農民播下種子,及每天悉心地照料.但他卻不能強迫種子發芽, 不能完全保護農作物免受突變天氣帶來的損害,更不能使植物都一定可以結出果子來.這個受造的世界提醒我們,我們是多麼需要倚賴上帝每天的供應.

正如撒種,澆灌,及收割的農夫一樣,對植物的生長過程「不曉得如何這樣」,我們對上帝國的推進

與及上帝在歷史中的工作,都係「唔知點解會搞成咁」.

經過這一年,隨著疫情發展,出現全球經濟大衰退,所造成的經濟破壞,有專家估計,可能要花很多年,才能重現轉機.不少教會及機構都出現財政困難.上年四月香港基督教機構協會進行過一次問卷調查,探討各基督教機構在疫情下的營運狀況,共 53 間機構總幹事或負責人參與.當中逾八成受訪人對事工前景感到負面或非常負面,近半更表示已進行或考慮減少事工開支,暫緩事工發展,進行特別籌募等.不少機構為應對疫情,社會運動及奉獻收縮等情況,已進行裁員及架構重整,同工自願減薪或放無薪假.以面對財政的困難.一年後的今天,情況可能更差.

全球經濟衰退不但對整個社會產生巨大影響,也大大影響教會的福音事工及普世宣教工作.疫情下各差會的宣教士差遣禮,事工分享會,短宣活動也相繼取消,國際會議也被迫延期,改為更多以網上方式舉辦會議.很多的宣教士因著國家邊境封鎖或航班取消,被逼留在工場不能回國.亦有很多的宣教士,因為疫情的關係不能夠出發.許多國家對人數聚集的限制,也直接影響到教會的聚會及福音的傳遞.在過去這一年,很多人開始討論,這次的疫情,加上香港現時的種種變化,會對教會的未來帶來什麼的影響與改變.對教會來說,這些到底是「危」還是「機」,還是兩者都是?答案在乎於我們怎樣回應.雖然我們經常會有「唔知點解會搞成咁」的感受,但上帝的話語及歷史的發展可以印證,祂仍然在掌管著歷史.我們只需倚賴上帝的恩典,按著祂給我們的領受,盡力做好自己本份就可以.

二. 小而珍貴

第二個比喻記載在馬可福音 4:30-32,同樣的比喻記載於馬太福音 13:31-32 及路加福音 13:18-19 中, 在這比喻中,耶穌用芥菜種來比喻上帝的國.

30 又說:「上帝的國,我們可用甚麼比較呢?可用甚麼比喻表明呢?31 好像一粒芥菜種,種在地裡的時候,雖比地上的百種都小,32 但種上以後,就長起來,比各樣的菜都大,又長出大枝來,甚至天上的飛鳥可以宿在它的蔭下.」

在以色列地,芥菜種大慨好像一粒鹽那麼大.卻可以長成為有如一棵灌木叢般大,有 2 至 6 呎高,以致雀鳥可以停在上面.鹽粒那麼大的芥菜種不單可以成長為一叢大的植物,而且芥菜種亦是一種有極強生命力和適應力種子.芥菜種事實上雖然並非世上最小的種子,但是在以色列地,芥菜種代表著「極小」的意思,大部分人都會嫌棄它,所以它有不起眼,不受人重視的意思.猶太人有這樣的諺語,每說到一件事情微小時,便說如「一粒芥菜種」,就等於我們所說的「芝麻綠豆」一樣.所以芥菜種不單是指它體積細小,更是指它似乎不能帶來任何實質的影響!換句話說,芥菜種的存在好像「微不足道」(insignificant),沒有任何貢獻或影響力.比喻強調上帝的國起初看起來雖然像芥菜種一樣遭人藐視,但是長起來會比各樣的菜更大,甚至長出大枝,讓天上的飛鳥都可以宿在它的蔭下.上帝的國會深深改變世界,令無數人因而得福.

這個比喻的解釋比起上一個比喻難,困難之處在於從不同的觀點釋經會帶來截而不同的結論.這個

「芥菜種的比喻」是正面的還是負面的?有解經家將這比喻作正面解釋,比喻是指上帝國的信息不斷的被傳開被接受,教會滿有上帝的祝福而大大的發展,為世界帶來好的改變；但亦有人認為這比喻更適合作負面解釋,耶穌是要警告門徒上帝國在擴張中將會出現反常的發展與變質,因有破壞及腐化的力量在教會中滋生帶來破壞,芥菜種就是這股力量的象徵.他們指出在馬太福音第 13 章的經文結構裡,「芥菜種的比喻」是置於「稗子的比喻」及「麵酵的比喻」中間.於是解經家提出,既然「撒稗子在田裡」明顯是指撒旦破壞的工作；他們亦指,在猶太人的思想中,經常把酵與罪惡的影響聯繫起來,酵跟腐爛敗壞連有關係,故此麵酵代表罪惡,所以「麵酵」和稗子一樣是不好的東西.耶穌還說要門徒防備法利賽人和撒都該人的「酵」呢(太 16:5-12)!因此「芥菜種的比喻」和「麵酵的比喻」應要看成是「稗子的比喻」的補充說明.加上耶穌在同一章裡的「撒種的比喻」裡就提過有飛鳥來吃撒在田間的種子,耶穌又說飛鳥是象徵撒旦的勢力,所以在「芥菜種的比喻」裡面所講有「天上的飛鳥來宿在它的枝上」的意思,應是指天國不斷擴展的同時,會有撒旦的勢力同時發功,在教會的發展的過程中,會有撒旦的騷擾,攻擊,誘惑,滲透在其中.

我沒採用這種對「芥菜種的比喻」負面的解釋,時間所限,我簡單提出三個質疑.

第一,「稗子的比喻」只出現在馬太福音裡,不見於其他的福音書.若然「芥菜種的比喻」和「麵酵的比喻」是「稗子的比喻」的補充說明,如此重要,為何馬可福音的作者不將「稗子的比喻」寫入書中呢?

第二,一個比喻中用到的象徵,在別的比喻中出現時並不一定代表同樣的意思.例如「飛鳥」一詞, 耶穌也用「天上的飛鳥和野地的百合花」作比喻,去教導人要仰賴上主,不要為生命憂慮.「飛鳥」不必然一定代表撒旦的勢力.在舊約提及獻祭時,若為感謝而獻上平安祭,可以用有酵的餅,與供物一同獻上,供祭司與獻祭者一同享用(利七 13-14).又如在五旬節時,將新素祭獻給上帝,要拿細麵加酵,烤成兩個搖祭的餅,當作初熟之物獻給上帝(利二十三 17).所以,「酵」也不必然是負面的東西.福音書作者在寫書時會把一些比喻編在一起,是想透過這種手法去將耶穌的一些教導以不同的面向角度去告訴讀者.

第三,在新約次經(Apocrypha)《雅各藏經》(Apocryphon of James)裡,有記載耶穌對門徒說了一段和

「芥菜種的比喻」相似的話:「天國就像在田間生長的麥穗,當它成熟後,就會把子粒散播,令田間

再次充滿麥穗,好待明年生長起來,你們也趕快去為自己收割生命的穗吧!好使國度裡也充滿了你們.」 “For the Kingdom of Heaven is like an ear of grain which sprouted in a field. And when it ripened, it scattered its fruit and, in turn, filled the field with ears of grain for another year. You also: be zealous to reap for yourselves an ear of life, in order that you may be filled with the Kingdom.” 雖然其內容跟「芥菜種的

比喻」有不同的地方,但總體來說,是想藉植物的生長表達出正面的意思,說天國的擴張是豐盛且為

人帶來祝福的.

最後,退一步來說,「芥菜種的比喻」的中心思想是想表示微不足道的事會帶來巨大的影響,都肯定上帝國是在擴張,這兩點不論正面還是負面的解釋都一樣.

那麼,「芥菜種的比喻」對今天的我們有什麼啟迪?如果我們覺得現時面對的處境十分艱難,覺得我們的信仰好像對現實沒有什麼影響力,怎去理解比喻正面解釋中所形容上帝國的威力?

我相信,上帝在危機中沒有停止工作,仍然呼召人與祂同工,即使在動盪不安的時代也是一樣.初期教會面對很大逼迫,門徒卻因而分散到各地傳講福音.如果我們再看二千多年的教會歷史,就會發現當教會面對極大的困難時,很多時亦正是福音工作最得到擴展的時間.反之,安逸的環境容易令信徒警醒心鬆懈下來,對福音使命亦失去了熱忱.就以香港教會的歷史為例.

五十年代,內地政局的變化為在華傳教運動劃上句號.49 年,中國政權易手並實施新宗教政策,成立

「三自愛國運動委員會」,由政府參與管理教會事務,並要求內地所有的教會加入.在這種政策下, 只有加入政權認可的三自教會才能生存,其他不願加入的教會,就難逃被解散的命運.很多教會及傳教士因此被逼選擇撤離內地.不少宗派的總會及差會南遷到香港,跟內地的堂會脫離關係,自立發展.另一方面,大批難民湧進香港,使香港成為一個宣教的重點.許多原來並沒有在香港開展事工的差會及傳教士撤至香港.當他們目睹大量中國難民滯港,決定在此開展新的事工,延續其對華宣教的使命.香港教會於是出現了百花齊放的多元宗派局面.

五十年代是香港教會的增長及躍進時代.有統計,1955 年香港共有 188 家堂會,信徒只有 5 萬人左

右.三年後,調查指教會增至 294 家,信徒人數超過 7 萬人,信徒的增長率接近四成.到 1962,教會

數目為 344 家,信徒更突破十萬大關.與 1955 年相比,不用十年,香港的信徒增加了一倍有多,平

均每年有 15.9\%的增長率.五十年代基督教在香港的增長,絕對是空前絕後.

除了大量宣教資源撤至香港外,因著有超過一百萬的難民湧入,激化了香港社會原有的種種問題.由於基督教界當時能夠迅速回應社會的各種需要,不僅促成了香港教會史上的增長與躍進,更大大紓解了香港社會的危機.香港基督教在香港的教育及社會福利服務方面建立的深厚基礎,便是在五十年代以後奠定的.

五十年代同時是香港教會神學教育及文字事工的突破期.神學教育方面,從國內南下在港復校或新成立的院校,到五十年代末共有十四家,有關數目在六十年代後更進一步增加.此外,國內的出版機構遷港,也促成了本地文字工作的蓬勃發展.香港一躍而成為華人神學教育及文字工作的重鎮,其服侍的對象,擴展至東南亞以至北美各地的華人教會.

不用多說明大家都知道,大家身處的中文大學,現時一同崇拜聚會的崇基學院及禮拜堂,也是在 50 年代因內地的變化而在香港植根.崇基書院是由內地第一所基督教大學聖約翰大學的校董歐偉國,另一所基督教大學廣州嶺南大學的校長李應林等人,在新中國成立後南遷到香港後,繼承國內十三間基督教大學的辦學使命,與香港的聖公會等教會團體籌辦的專上學院.歐李二人的復辦理念,得到當時聖公會主教何明華認同,加上當時的社會需要,就借助教會的龐大資源,1951 年先借用教會會堂作學院校舍,及後取得馬料水土地,以英美教會團體的資助建成 1956 年啟用的馬料水校園,即今日的崇基學院.

我想強調的是,教會不會遇到困難就停止去傳耶穌的愛,也不會受到迫害就消失.一間有見證的教會, 是在困難中仍然勇敢見證基督的福音.願上帝賜給我們智慧,面對時代處境下的逆境.

三. 你是一粒怎樣的種子?

耶穌基督用種子生長和芥菜種作比喻,來教導天國福音的本質和潛能,講的是新生命的種子,如何更新人一生的路途和走向.那麼小的種子卻可以長成好幾公尺高的樹,關鍵在於上帝在創造它時,就把芥菜有的特質及生命力都放在它的裡面.天國是有生命力,像一粒種子會發芽生長,而且從很細小發展成很大,這是天國的現象.

種子必需要有生命力.如果它沒有生命力,是粒壞種,不能發芽,沒有用處,不能活下去.而且,種子要願意犧牲自己.耶穌說:「一粒麥子不落在地裡死了,仍舊是一粒,若是死了,就結出許多子粒來.」(約十二 24)一個基督徒不但要有生命力,還要願意為他人捨己.因此我們既然已有主永恆的生命在我們裡面,我們就應該順服主的帶領,靠著聖靈好好地活出主的生命,為主作見證.

在預備這次講道時,想起《普天頌讚》裡面的〈小小水滴歌〉.

小小點點水滴,積聚成大海；小小細細砂粒,積聚成大地.

小小一分一秒,好像不希奇,倒會積成年歲,積成無盡期.

小小一點過失,會使我的靈,離開德行路程,轉向惡中行.

小小慈愛事情,小手播了種,能夠開出香花,施福千萬眾.

小小和愛事情,小小溫柔話,能使世間改善,變成天父家.

【Little Drops of Water〕

Little drops of water, little grains of sand, make the mighty ocean and the beauteous land. And the little moments, humble though they may be, make the mighty ages of eternity.

Little deeds of kindness, little words of love, make our earth an Eden, like the heaven above. So our little errors lead the soul away, from the paths of virtue into sin to stray.

Little seeds of mercy sown by youthful hands, grow to bless the nations far in heathen lands. Glory then for-ever be to God on high, beautiful and loving, to eternity.

《小小水滴歌》是由兩位教會先賢楊蔭瀏和劉廷芳在 1934   年譯出,最初刊登在一本叫《紫晶》的基

督教文藝季刊裡面.原曲本來有 6 節,翻譯成中文後把第 6 節刪去,歌詞的次序也改變了.原曲的第

5 節(即現今譯本的第 4 節)翻成中文後意思亦有改變.原文是體現了非常明顯的宣教情懷,字面意思大概是這樣的:「年輕的手種下的憐憫種子雖小,卻成長起來遠至異教徒的土地,福佑各族.」但譯者將「成長起來遠至異教徒的土地,福佑各族」轉譯為「能夠開出香花,施福千萬眾」,「異教徒的土地」一字則刪除.楊蔭瀏和劉廷芳這樣處理並非反對宣教,只是他們將這個敏感且不利於中國教會發展的詞句改動了.意思同樣是我們的好見證好行為可以為其他人帶來祝福,使人將榮耀歸與上主.

前崇基學院神學院院長邢福增教授在 2015 年一個大專培靈會中,引用芥菜種及麵酵的比喻,去勉勵

信徒在艱難的日子「深耕細作」.

他說:「耶穌希望門徒明白,天國所看重的,不是外表,而是生命的內涵；不僅是數量,更看重質素.」又表示「我們活在一個善惡並存的破碎世代,在終末審判仍未實現的當下,既擁抱著天國福音的整全價值,又經歷整全與破碎間的各種張力,委身於微小的深耕細作,在地下建立公民社會,同時也在拓展天國國度.」「天國雖不能全然在地上實現,卻不代表這只是空想的少烏托邦,門徒必須堅持天國的理想.當耶穌說:『天國好比人撒好種在田裡』時,表明天國之子就是天國在世界中讓人可以見到的天國,當魔鬼企圖以惡者之子來覆顛這世代時,天國之子卻要證明,在惡以外,仍有善,在罪以外, 仍有救恩.基督徒雖然不是完美的聖潔者,教會也不是完美的聖潔群體,卻是世界中唯一可以傳講天國福音,並且見證天國福音,讓人能夠預嚐天國的真實.

事實上,在這段艱難的日子中,令人驚訝的是,上帝激勵衪的兒女在這世代奮勇服侍.當教會及機構很多的聚會都因為疫情及限聚令暫停,但社關及社區服事沒有停下.不少信徒積極服侍社區,有人派飯,派口罩等防疫物資給社會中有需要及基層工作的人,不止是派,更加是用心地找機會跟受助者傾談,表達關心.信仰群體轉化成為不同型態地在不同地方實踐所信,信徒雖不在建築物內,卻在網絡上,在社區中,分散在不同角落去回應不同的需要.

「不以善小而不為,不以惡小而為之」,願我們彼此勸勉.

結語

福音的種子像芥菜種一樣,若是種在人的心田裡,就能發揮它的影響力.讓我們持守在我們心中真理的種子,讓真理的種子不斷的在我們心中長大,成為大樹,成為別人的祝福,天國的福音從我們個人開始,從我們生活中開始,去做一件上帝喜悅的事,去關心一個人,去服事一個人開始,去做一件別人看來微不足道的事開始,因為一個小小的行為,裡面有耶穌的愛,祝福就會臨到一個人或是一個地方.阿們!



\newpage

\section{為何耶穌可以睡得這般香甜}
\label{sec:555}
\textbf{鄧瑞強博士 - 為何耶穌可以睡得這般香甜?}
\newline
\newline
\hyperref[sec:554]{< < < PREV SERMON < < <}
~
\hyperref[sec:toc]{[返目錄]}
~
\hyperref[sec:556]{> > > NEXT SERMON > > >}


馬可福音 4:35-41
\newline
\begin{longtable}{cl}
\hline
\hline
章節 & 經文 (和合本修訂版)\\
\hline
4:35 & \begin{tabularx}{0.7\textwidth}{X} 那天晚上，耶穌對門徒說：「我們渡到對岸去吧。」 \end{tabularx} \\ \\ \relax
4:36 & \begin{tabularx}{0.7\textwidth}{X} 門徒離開眾人，耶穌已在船上，他們就請他一同去；也有別的船和他同行。 \end{tabularx} \\ \\ \relax
4:37 & \begin{tabularx}{0.7\textwidth}{X} 忽然狂風大作，波浪打入船內，以致船灌滿了水。 \end{tabularx} \\ \\ \relax
4:38 & \begin{tabularx}{0.7\textwidth}{X} 耶穌在船尾上，枕著枕頭睡覺。門徒叫醒他，說：「老師！我們快沒命了，你不管嗎？」 \end{tabularx} \\ \\ \relax
4:39 & \begin{tabularx}{0.7\textwidth}{X} 耶穌醒了，斥責那風，向海說：「住了吧！靜了吧！」風就止住，大大平靜了。 \end{tabularx} \\ \\ \relax
4:40 & \begin{tabularx}{0.7\textwidth}{X} 耶穌對他們說：「為甚麼膽怯？你們還沒有信心嗎？」 \end{tabularx} \\ \\ \relax
4:41 & \begin{tabularx}{0.7\textwidth}{X} 他們就非常懼怕，彼此說：「這到底是誰？連風和海都聽從他。」 \end{tabularx} \\ \\
[1ex]
\hline
\hline
\end{longtable}
講題:為何耶穌可以睡得這般香甜?　Why could Jesus Sleep so Sweetly?

講員:鄧瑞強博士

經文:馬可福音4:35-41

場合:香港中文大學崇基學院禮拜堂主日崇拜

日期:2021年6月20日

各位弟兄姊妹,早安.在這疫症裡,祝願各人平安.

夜闌人靜,應該睡覺的時候,很多人卻成了哲學家,思考一個嚴肅的哲學問題:「為什麼我還未睡覺?」

有人說,移民千萬不要移民去新西蘭,因為你若去了那裡,你成晚都在數羊.

我發覺有很多人,晚上睡覺時,怕睡不夠,故此,每隔十五分鐘,便會打開眼,看看鐘,看看還有多少時間睡.並且,會越看越憂心.啊,還有三個鐘可睡.隔十五分鐘,又看看鐘.隔十五分鐘,又看看鐘.隔十五分鐘,又看看鐘.心想,唉,還只有兩個鐘可睡.如此這般,一晚便過去了.

所以,現在千萬不要對人說:「昨晚睡了十個鐘,完全不知不覺,一會兒便天光了.」你這樣說,很多人會因為妒忌你而殺了你.

不過,現在科學昌明,什麼問題,科學都能提供答案.科學家已找到應付失眠的最佳方案.應付失眠的最佳方案是:多睡一點.

哲學家黑格爾(Hegel)說:「智慧的貓頭鷹只在夜間飛行」.現代人大多放棄做人,做了貓頭鷹.在夜間,思想飛行,成了充滿智慧的哲學家.不斷問:「為什麼?」「為什麼?」「為什麼還睡不著?」「為什麼今天那股票升了那麼多而我竟沒有大手購入?」「為什麼那個人今天這樣對我?」「為什麼我今天會對上司講這些廢話?」「為什麼今天這般簡單的工作我也做不來?」簡言之,十萬個為什麼.

睡覺與睡不著覺,that is the question.

今日的講道經文:馬可福音4:35-41

在這故事中,令人羨慕的,是主耶穌睡得很甜.

可4:35　當那天晚上,耶穌對門徒說:「我們渡到那邊去吧.」

可4:36　門徒離開眾人,耶穌仍在船上,他們就把他一同帶去；也有別的船和他同行.

可4:37　忽然起了暴風,波浪打入船內,甚至船要滿了水.

可4:38　耶穌在船尾上,枕著枕頭睡覺.門徒叫醒了他,說:「夫子!我們喪命,你不顧嗎?」

可4:39　耶穌醒了,斥責風,向海說:「住了吧!靜了吧!」風就止住,大大地平靜了.

可4:40　耶穌對他們說:「為什麼膽怯?你們還沒有信心嗎?」

可4:41　他們就大大地懼怕,彼此說:「這到底是誰,連風和海也聽從他了.」

首先,讓我們zoom近,去看看主耶穌.他在船尾上,枕著枕頭睡覺.對現代人來說,單單這一鏡頭,已足以叫人信耶穌.

他周圍環境,是「起了暴風,波浪打入船內,甚至船要滿了水」.身邊的人,應是呼天搶地,驚惶失措.這些無眠的人,與耶穌成了一大對比.

於是,有了一神學問題:為何有人睡得著?為何有人睡不著?

睡不著覺,有很多原因.

中國五經之首《詩篇》第一首詩＜關雎＞便討論失眠:

「關關雎鳩,在河之洲.窈窕淑女,君子好逑.……

求之不得,寤寐思服.悠哉悠哉,輾轉反側.」

你有追求的目標,你有追求的對象,但求之而不得,日思夜想,便睡不著覺,輾轉反側.

人無求,便心安；放下,就自在.有所戀,有所求,心怎得安寧.

對「失眠」反省得最深的哲學家,要算是羅馬尼亞裔旅居法國的哲學家蕭沆(Emil Cioran).有文章顯示,他失眠了五十年.他講到自己的光景:「在我年輕時,一個星期復一個星期,沒有閉過眼睛.我活在一個沒有人活過的世界裡.我感覺到,時間及其流動深刻地觸碰我的存在.時間在我裡面達到頂峰,時間在我裡面取得絕對勝利.」蕭沆以極優美的文筆,描繪在失眠時,人終極面對的,就是時間的流逝對人的絕對控制.

在蕭沆眼中,失眠是人類最重要的哲學主題,對他來說,失眠呈現人類最悲慘的現實.失眠令身體與心靈的連結解體.失眠的人,像是無靈魂的軀殼,又像是無軀殼的靈魂.在夜間,百無聊賴地思想,又像行屍走肉般活動.想思想時,覺得不清醒.想睡覺時,卻覺得過份清醒.在失眠中,蕭沆體會到人類的本質就是「時間」.當存在的一切感覺消失時,留下的,是體會到無盡的時間在無盡地延長.人類理性探討的一切意義,在失眠中失去了意義.例如:倫理學指引我們作出某些道德行動,好體會生命的意義.但在失眠中,任何行動都是無意義的.人類能掌握的確定知識,如:動物累的時候會睡覺,在失眠中,這一切確定知識也失去確定性.失眠呈現人類面對的可怕真相.這可怕的真相是:你以為事物總有其意義,但失眠是絕對無意義的.在失眠中,一個人像是孤獨地活在地獄中,每一刻的流逝不是豐富你的存在,而是折磨你的存在,迫使你發瘋.在失眠中的天問,是這痛苦何時才會停止.在失眠中,你體會到,活著像是釘在十架上,痛苦而無意義.簡言之,失眠揭露了人類存在的空洞與無意義.

我讀了很長時間的哲學,很少有哲學家,如此這般地談論失眠的.

蕭沆以「失眠」的經驗,作為一開啟人性現實的鑰匙,道盡人類的痛苦.面對如此的一個失眠的人,面對如此痛苦的慘況,我能說什麼呢?

人類活著時,向自己存在的地平線望去,便會見到死亡.我們活著,什麼都不能確定；能確定的,是我們必有一死.死,是我們最能確定的終點.這就是哲學家蕭沆在「失眠」的經驗中體會到的空洞的黑暗.失眠的人,看來活著,但他體會到的,是空洞無物,是時間在無意義地消逝,是一個人對白天時一切悲歡離合,成敗得失的無意義的追憶.這樣的經驗,就是活著去經驗空洞.當然,最後的空洞,就是我們的死亡.活著本身,就是死亡的過程.人生,是一場大失眠.

在馬可福音的故事裡,在那船上,門徒一夜無眠,他們在自己可能面對的大限面前掙扎.經文說,他們叫醒熟睡的耶穌,說:「夫子!我們喪命,你不顧嗎?你不關心嗎?你不理會嗎?」

這是一群失眠的人,他們有所掛慮.他們的掛慮,終極來說,便是死亡.面對死亡,面對生命的毫無意義,面對絕對的空洞,他們極力地行動,極力地反抗.但一切都是枉然.像失眠的人,極力地讓自己睡去,最後都是枉然一樣.他們發出生命的呼喊,像一個失眠的人,大呼:「救命!」

幸好,他們求救於一個熟睡的人.這個人會教他們安睡之道.

面對「失眠」這哲學問題,耶穌作出「神學」的回應.

可4:39　耶穌醒了,斥責風,向海說:「住了吧!靜了吧!」風就止住,大大地平靜了.

可4:40　耶穌對他們說:「為什麼膽怯?你們還沒有信心嗎?」

可4:41　他們就大大地懼怕,彼此說:「這到底是誰,連風和海也聽從他了.」

耶穌向風和海說:「住了吧!靜了吧!」風和海便平靜了.

門徒的反應是懼怕,這是什麼人,這到底是誰,連風和海也聽從他.

一般人眼中的風浪,一般人懼怕的死亡,對真正能熟睡的人而言,算不得什麼.他若像一般人那樣緊張,便不可能熟睡了.在熟睡的人的心靈裡,風和海失去力量；在這個人面前,風和海平靜了.

回到我們的神學問題:為何有人睡得著?為何有人睡不著?

關鍵的分別在哪裡?

聽聽主耶穌向我們發出的問題:「為什麼膽怯?你們還沒有信心嗎?」睡得著與睡不著,耶穌一語道破,關鍵在信心.失眠的問題,是信仰的問題.

耶穌講的信心,不是現代人強調的自信.現代人的自信,是信自己,這其實即是不信神.失眠,作為一「隱喻」,正正揭露我們自身的不足信,正正揭露我們的「自負」背後的無能,正正揭露我們的脆弱.現代社會製造出來的宗教,是教我們相信,我們是行的,是能掌控命運的,是價值的中心.失眠,讓我們看到,我們做什麼也是無用的.連睡眠也掌控不到,何來掌控自己和他人的命運.自我是價值的中心嗎?在失眠中,你就體會到生命在無意義地流逝.失眠,作為一啟示性的人類經驗,讓我們謙卑下來.

有文章說,古羅馬殘暴的凱撒尼祿,是失眠的人.近代的希特拉,也是失眠的人.每一晚,他們都與自己搏鬥.他們認為,睡不著,不是一種失敗,而是一種成功.當世人沉睡時,只有他們清醒.因為只有他們清醒,他們誤以為自己是上天派下來的人,去操控沉睡的世人.其實,在夜裡,他們不是清醒,他們只是失眠.在失眠中,他們的身體與心靈解體,虛無飄渺地活著,每一晚,發著操控世界的春秋大夢.

在這脆弱的人生裡,耶穌不操控什麼,他枕著枕頭睡覺.

睡得著或睡不著,耶穌說,分別在「信心」.耶穌講的信心,不是現代人強調的這種自信,這種對自我的執著,這種對自我能力的迷戀.耶穌講的信心,是將操控自己生死與命運的掌控權交出來,交給永恆的上帝.你握著拳頭,就無人能與你握手.放開手,神才能握著你.

每一晚,為什麼你都要親自赤膊上陣,與自己的限制,缺乏,不足,遺憾搏鬥呢?放開手,交給神,風和海就會平靜的了.

為什麼因為有人得罪了你而心懷仇恨,而不相信神會為你主持正義?為什麼因為明天可能會遇到的幽暗而滿懷憂慮,而不相信神就是光明的源頭?為什麼因為得不到的利益而不甘心,而不相信神所賜的恩典才有永恆價值?為什麼我們總是不相信神?

睡得著與睡不著,終歸而言,是信仰問題.失眠,是一溫馨提示,提示我們要好好地信任神.

耶穌對失眠的人說:「為什麼膽怯?你們還沒有信心嗎?」

最後,我想講,我們的上帝是很關心我們的失眠情況的.祂有一滿懷恩典的解決方案,就是:祂叫我們返崇拜聽道.聽道時,大部分人都睡得香甜.

祝願各人睡得好.

願榮耀歸上主,願人間享和平.



\newpage

\section{靜默有時言語有時}
\label{sec:556}
\textbf{林豪恩先生 - 靜默有時,言語有時}
\newline
\newline
\hyperref[sec:555]{< < < PREV SERMON < < <}
~
\hyperref[sec:toc]{[返目錄]}
~
\hyperref[sec:557]{> > > NEXT SERMON > > >}


耶利米哀歌 3:22-33
\newline
\begin{longtable}{cl}
\hline
\hline
章節 & 經文 (和合本修訂版)\\
\hline
3:22 & \begin{tabularx}{0.7\textwidth}{X} 因耶和華的慈愛，我們不致滅絕，因他的憐憫永不斷絕， \end{tabularx} \\ \\ \relax
3:23 & \begin{tabularx}{0.7\textwidth}{X} 每早晨，這都是新的；你的信實極其廣大！ \end{tabularx} \\ \\ \relax
3:24 & \begin{tabularx}{0.7\textwidth}{X} 我心裡說：「耶和華是我的福分，因此，我要仰望他。」 \end{tabularx} \\ \\ \relax
3:25 & \begin{tabularx}{0.7\textwidth}{X} 凡等候耶和華，心裡尋求他的，耶和華必施恩給他。 \end{tabularx} \\ \\ \relax
3:26 & \begin{tabularx}{0.7\textwidth}{X} 人仰望耶和華，安靜等候他的救恩，這是好的。 \end{tabularx} \\ \\ \relax
3:27 & \begin{tabularx}{0.7\textwidth}{X} 人在年輕時負軛，這是好的。 \end{tabularx} \\ \\ \relax
3:28 & \begin{tabularx}{0.7\textwidth}{X} 他當安靜獨坐，因為這是耶和華加在他身上的。 \end{tabularx} \\ \\ \relax
3:29 & \begin{tabularx}{0.7\textwidth}{X} 讓他臉伏於地吧！或者還會有指望。 \end{tabularx} \\ \\ \relax
3:30 & \begin{tabularx}{0.7\textwidth}{X} 讓人打他耳光，使他飽受凌辱吧！ \end{tabularx} \\ \\ \relax
3:31 & \begin{tabularx}{0.7\textwidth}{X} 主必不永遠撇棄， \end{tabularx} \\ \\ \relax
3:32 & \begin{tabularx}{0.7\textwidth}{X} 他雖使人憂愁，還要照他豐盛的慈愛施憐憫； \end{tabularx} \\ \\ \relax
3:33 & \begin{tabularx}{0.7\textwidth}{X} 他並不存心要人受苦，令世人憂愁。 \end{tabularx} \\ \\
[1ex]
\hline
\hline
\end{longtable}
講題:靜默有時,言語有時　A Time to Keep Silent, a Time to Speak

講員:林豪恩先生

經文:耶利米哀歌3:22-33

場合:香港中文大學崇基學院禮拜堂主日崇拜

日期:2021年6月27日

引言:

伍渭文牧師曾經說:「現在教會崇拜傾向抒發．．．因此任何禮序都用上音樂,帶動崇拜者的情緒.崇拜遲開始等人時彈奏鋼琴,禱告時用琴音作背景音樂,甚至靜默時,也以靜靜的琴音助攻.若視崇拜單單有抒發作用,而且著重情緒的抒發,音樂就變成有效利器了,音樂確是有效製造甚至控制情緒的.不過崇拜不單是抒發情感,乃回應上主在基督裡的救贖.聖餐拿起餅,喝這杯,是「表明主的死,直等到他來.」崇拜是信仰的表白.所以,詩歌敬拜後,不是「敬拜完畢」,乃崇拜的開始.講道重要,因為闡釋基督事件,預備教會赴羔羊的婚宴.詩篇是詩歌,但詩篇的細拉,就是靜默,因為歌唱有時,靜默有時.

有幾次在崇基禮拜堂證婚時,司琴竟然習慣地用琴音伴奏,我客氣的說:我們想聽清楚新人的誓詞,若誤會說出不願意,就貽誤終身了.有一次參加好友牧師主持的火葬禮,棺木進人火化爐時,也是唱著詩歌.之後跟好友牧師建議:讓會眾剛唸完信經「……我信身體復活,我信永生.」語音剛落,靜聽轆轆輸棺入爐的聲音最好,這是純樸厚實的律法的聲音-人因罪而死的審判,不要用琴音和歌聲伴隨；恩典與律法並舉,才是福音的信仰.

崇拜老是說過不停,彈奏不停,會產生疲勞甚至心靈頑固,而崇拜至重要是讓心靈敏銳,否則會與那聖者失諸交臂了.(伍渭文:「音樂不是禮儀的僕婢」)

崇拜中唱歌有時,靜默有時,信仰生活豈也不也是一樣嗎?讓我們以伍牧師這篇文章的建議來啟發我們閱讀今天這四段經課,思考信仰生活中「靜默有時,言語有時」.

經文一:耶利米哀歌:靜默有時(他當獨坐無言)

3:22我們不致消滅,是出於耶和華諸般的慈愛；是因他的憐憫不致斷絕.

3:23每早晨,這都是新的；你的誠實極其廣大!

3:24我心裡說:耶和華是我的分,因此,我要仰望他.

3:25凡等候耶和華,心裡尋求他的,耶和華必施恩給他.

3:26人仰望耶和華,靜默等候他的救恩,這原是好的.

3:27人在幼年負軛,這原是好的.

3:28他當獨坐無言,因為這是耶和華加在他身上的.

3:29他當口貼塵埃,或者有指望.

3:30他當由人打他的腮頰,要滿受凌辱.

3:31因為主必不永遠丟棄人.

3:32主雖使人憂愁,還要照他諸般的慈愛發憐憫.

3:33因他並不甘心使人受苦,使人憂愁.

《耶利米哀歌》哀悼耶路撒冷陷落,焚毀和人民被擄.耶路撒冷被毀不單意味著可恥的失敗及放逐.正當這一切令人難以忍受時,這事件所引發的神學思想,更是信靠神的猶太人所難以理解的.耶和華上帝不是選擇耶路撒冷來顯明祂自己嗎?何以耶路撒冷會被外邦攻破呢?是耶和華比不上外邦的神明嗎?還是耶和華放棄了祂對耶路撒冷的應許呢?就算是耶路撒冷犯罪被懲罰,那攻陷耶路撒冷的外邦國家難道比耶路撒冷更公義嗎?不義的豈可審判比他公義的呢?這書卷中的哀歌,不單哀悼耶路撒冷被毀所帶來的痛苦,並且提出質疑和控訴.

哀歌(第一章)用擬人法,把耶路撒冷比擬為一位女士來訴說哀情:

2她夜間痛哭,淚流滿頰,

在所有親愛的人中,找不到一個安慰她的.

她的朋友都以詭詐待她,

成為她的仇敵.

3猶大被擄,

遭遇苦難,多服勞役.

她住在列國中,得不著安息；

追逼她的在狹窄之地追上她.

4錫安的道路因無人前來過節就哀傷,

她的城門荒涼,

祭司嘆息,

少女悲傷；

她自己充滿痛苦.

耶路撒冷被毀後的猶大人感到自己充滿悲傷,上帝卻又容許他們悲傷,容許他們哀悼.

在困苦中,落入逆境中,或傷痛中的人們,很自己會問:是誰令耶路撒冷落得如此悲慘境況的呢?是誰令耶路撒冷的兒女生離死別,流離失所的呢?當然是巴比倫,還有那些助紂為虐的憐國.還有呢?還有那位見死不救的耶和華上帝,祂不但不解救耶路撒冷,哀歌指出祂正是耶路撒冷的痛苦的來源:

13他從高處降火進入我的骨頭,

剋制了我；

他張開網,絆我的腳,

使我退後,

又令我終日淒涼發昏.

14他用手綁我罪過的軛,

捲繞著加在我頸項上；

他使我力量衰敗.

主將我交在我不能抵擋的人手中.

15主棄絕我們當中所有的勇士,

聚集會眾攻擊我,

要壓碎我的年輕人.

主踹下少女猶大,

在醡酒池中.

16我因這些事哭泣,

眼淚汪汪；

因為那安慰我,使我重新得力的,

離我甚遠.

我的兒女孤苦,

因為仇敵得勝了.

上帝不但容許亡國的猶大人悲傷,也容許他們悲憤,甚至容許他們向祂發出控訴.

實在有太多訴說不完的哀情,有太多控訴不盡的冤屈.或者有人說,不想發生的事情既然都已經發生了,又沒有可能回到過去修改歷史,可不可以瀟灑一點面對呢?可不可以隨遇而安呢?很多人羨慕蘇軾在「定風波」所顯露的灑脫:

莫聽穿林打葉聲,何妨吟嘯且徐行.

竹杖芒鞋輕勝馬,誰怕?一蓑煙雨任平生.

料峭春風吹酒醒,微冷,山頭斜照卻相迎.

回首向來蕭瑟處,歸去,也無風雨也無晴.

如果最終能夠釋懷,往往都已經是經過年日的洗禮,也是一連串心理過程的結果.Elisabeth Kübler-Ross在1969年出版的《論死亡與臨終》(On Death and Dying)一書中提出「哀傷的五個階段」(Five Stages of Grief)被廣泛應用於多種處境,第一階段:驚嚇及否認(Shock and Denial),第二階段:憤怒(Anger),第三階段:沮喪(Depression and Detachment),第四階段:對話與討價還價(Dialogue and Bargaining),第五階段:接受(Acceptance).這過程的長短可以因人而異,也因事情而異.

灑脫地放下,說來容易,當中的難處,只有當事人最明白.還記得李清照的《聲聲慢》怎樣說嗎?

尋尋覓覓,冷冷清清,淒淒慘慘戚戚.

乍暖還寒時候,最難將息.

三杯兩盞殘淡酒,怎敵他,晚來風急.

雁過也,正傷心,卻是舊時相識.

滿地黃花堆積.憔悴損,如今有誰堪摘.

守著窗兒,獨自怎生得黑.

梧桐更兼細雨,到黃昏點點滴滴.

這次第,怎一個愁字了得!

耶利米哀歌的作者在第一章和第二章哭訴了哀情,控訴了冤屈,在第三章28節卻說:「他當獨坐無言」.獨坐無言並不是困難已解決,也不是灑脫,只是對事情的態度正在改變,回應事情的階段也在轉變.此刻,詩人呼籲獨坐無言,很可能是在哭訴和控訴之後,漸漸察覺到自己在事情中的角色,承認自己在困境中的責任,也開始能夠從自己以外的角度去看事情,例如從耶和華上帝的角度看自己的狀況.因此,在第三章詩人如此詩:

3:25凡等候耶和華,心裡尋求他的,耶和華必施恩給他.

3:26人仰望耶和華,靜默等候他的救恩,這原是好的.

3:27人在幼年負軛,這原是好的.

等候是好的,靜默是好的,幼年負軛也是好的.在哀情之中,詩人為何有如此的盼望呢?因為在向耶和華上帝哭訴和控訴的過程中,他領悟到:

3:22我們不致消滅,是出於耶和華諸般的慈愛；是因他的憐憫不致斷絕.

3:31因為主必不永遠丟棄人.

3:32主雖使人憂愁,還要照他諸般的慈愛發憐憫.

3:33因他並不甘心使人受苦,使人憂愁.

言語有時,靜默有時.眼淚流過了,控訴說過了,或者,這是「當獨坐無言」的時候了.有些時候,獨坐無言才是合適的回應,才看見黑洞盡頭的微光,才聽到心靈的微聲.

經文二:詩30篇:言語有時(好叫我的靈歌頌你,並不住聲)

30:1耶和華啊,我要尊崇你,因為你曾提拔我,不叫仇敵向我誇耀.

30:2耶和華－我的上帝啊,我曾呼求你,你醫治了我.

30:3耶和華啊,你曾把我的靈魂從陰間救上來,使我存活,不至於下坑.

30:4耶和華的聖民哪,你們要歌頌他,稱讚他可記念的聖名.

30:5因為,他的怒氣不過是轉眼之間；他的恩典乃是一生之久.一宿雖然有哭泣,早晨便必歡呼.

30:6至於我,我凡事平順,便說:我永不動搖.

30:7耶和華啊,你曾施恩,叫我的江山穩固；你掩了面,我就驚惶.

30:8耶和華啊,我曾求告你；我向耶和華懇求,說:

30:9我被害流血,下到坑中,有甚麼益處呢?塵土豈能稱讚你,傳說你的誠實嗎?

30:10耶和華啊,求你應允我,憐恤我!耶和華啊,求你幫助我!

30:11你已將我的哀哭變為跳舞,將我的麻衣脫去,給我披上喜樂,

30:12好叫我的靈歌頌你,並不住聲.耶和華－我的上帝啊,我要稱謝你,直到永遠!

詩30篇最後一節說:「好叫我的靈歌頌你,並不住聲.耶和華－我的上帝啊,我要稱謝你,直到永遠!」有些時候,獨座無言是不適合的,反而要「不住聲」.

其他人可能不明白為甚麼詩人會「歌頌你,並不住聲」,甚至「稱謝你,直到永遠」,你明白嗎?你了解嗎?還記得少林足球中這一幕嗎?

周星馳說:「做人冇夢想,同條咸魚有咩分別?」然後一個穿著恤衫,梳「大小便裝」的男子突然出現說:

「我叫做醬~~爆~~~我從來沒有放棄成為一個偉大的作曲家的理想~~~剛才聽到這位先生唱的歌~~~熱~~~情~~~奔~~~放~~~創~~~意~~~無~~~限~~~點燃起我醬爆心~~~中的一團火,我醬爆感覺到,在這個MOMENT~~~~~~要爆~~~~~~~~~~啦」於是,他用do,re,mi等音調高歌起來.

或者,其他人都不明白醬爆為甚麼在那一刻唱起歌來,因為「做人冇夢想,同條咸魚有咩分別?」這一句引爆了他心中的那一團火,就是「從來沒有放棄成為一個偉大的作曲家的理想」那一團火.那團火是在醬爆心中的,不是在其他人心中的,難怪其他人不明白,甚至視為奇怪.

有些信徒開口埋口就說「感謝主」,有些信徒祈禱總說「哈利路亞」,有些人聽到感到很奇怪,甚麼有所保留.如果順口開河,行禮如儀,假冒為善,當然是反感,但如果是發自真誠和激情呢?試想想詩人何以「歌頌上帝並不住聲」呢?看他的自述,知道他經歷過哀哭的事情,曾經披過麻衣,曾經被害流血,下到坑中.在這些狀況之中,他也經歷到上帝把他的哀哭變為跳舞,把他的麻衣變為喜樂,把他的靈魂從陰間救上來.這是詩人親身而且真實的經歷,他的感情是有這些親身經歷為基礎的.其他人不明白,但有類似經歷的人是有共鳴的.我還記得,那次我講完道之後,有肢體與我分享他在病患中如何體會到上主的恩典.我不知道他說了多少次類似的話,但他說的時候,仍是帶著感動的,聽的人也感受到他內心的那份感動.我仍記得他說了類似詩篇30篇的詩人的話:我仍活著的日子都要述說上帝的恩典.

靜默有時,說話有時,如果要說話,當說甚麼話呢?讓詩篇30篇詩人的見證成為我們的提醒.

經文三:馬可福音:靜默有時(耶穌切切地囑咐他們,不要叫人知道這事)

5:21耶穌坐船又渡到那邊去,就有許多人到他那裡聚集；他正在海邊上.

5:22有一個管會堂的人,名叫睚魯,來見耶穌,就俯伏在他腳前,

5:23再三地求他,說:「我的小女兒快要死了,求你去按手在她身上,使她痊癒,得以活了.」

5:24耶穌就和他同去.有許多人跟隨擁擠他.

5:25有一個女人,患了十二年的血漏,

5:26在好些醫生手裡受了許多的苦,又花盡了她所有的,一點也不見好,病勢反倒更重了.

5:27她聽見耶穌的事,就從後頭來,雜在眾人中間,摸耶穌的衣裳,

5:28意思說:「我只摸他的衣裳,就必痊癒.」

5:29於是她血漏的源頭立刻乾了；她便覺得身上的災病好了.

5:30耶穌頓時心裡覺得有能力從自己身上出去,就在眾人中間轉過來,說:「誰摸我的衣裳?」

5:31門徒對他說:「你看眾人擁擠你,還說『誰摸我』嗎?」

5:32耶穌周圍觀看,要見做這事的女人.

5:33那女人知道在自己身上所成的事,就恐懼戰兢,來俯伏在耶穌跟前,將實情全告訴他.

5:34耶穌對她說:「女兒,你的信救了你,平平安安地回去吧!你的災病痊癒了.」

5:35還說話的時候,有人從管會堂的家裡來,說:「你的女兒死了,何必還勞動先生呢?」

5:36耶穌聽見所說的話,就對管會堂的說:「不要怕,只要信!」

5:37於是帶著彼得,雅各,和雅各的兄弟約翰同去,不許別人跟隨他.

5:38他們來到管會堂的家裡；耶穌看見那裡亂嚷,並有人大大地哭泣哀號,

5:39進到裡面,就對他們說:「為甚麼亂嚷哭泣呢?孩子不是死了,是睡著了.」

5:40他們就嗤笑耶穌.耶穌把他們都攆出去,就帶著孩子的父母和跟隨的人進了孩子所在的地方,

5:41就拉著孩子的手,對她說:「大利大,古米!」(翻出來就是說:「閨女,我吩咐你起來!」)

5:42那閨女立時起來走.他們就大大地驚奇；閨女已經十二歲了.

5:43耶穌切切地囑咐他們,不要叫人知道這事,又吩咐給她東西吃.

福音經課記述耶穌醫治了一位患了十二年的血漏的女人,然後又把睚魯十二歲的女兒從死亡的邊緣救活回來.最令人意外的,是故事的最後一句:「5:43耶穌切切地囑咐他們,不要叫人知道這事」.按福音書的記載,耶穌不止一次吩咐人不要宣揚.很多人都為此而感到奇怪,也有各種不同的猜想.或者,就正如耶穌自己說:「我的時候還沒有到」(約7:6)

有時,不說,不是因為事情不是事實；不說,不是因為事情不是好事.有些時候,不說,只是因為不是當說的時候.現實上,無論是主動還是被動,的確有需要保持沉默的時候.

不說,不等於放棄；不說,也不等於甚麼也不可以做.約翰一書3:18說:「小子們哪,我們相愛,不要只在言語和舌頭上,總要在行為和誠實上.」書信經課所提及的就是其中一個例子.

經文四:哥林多後書:顯露有時(在這慈惠的事上)

8:7你們既然在信心,口才,知識,熱心,和待我們的愛心上,都格外顯出滿足來,就當在這慈惠的事上也格外顯出滿足來.

8:8我說這話,不是吩咐你們,乃是藉著別人的熱心試驗你們愛心的實在.

8:9你們知道我們主耶穌基督的恩典:他本來富足,卻為你們成了貧窮,叫你們因他的貧窮,可以成為富足.

8:10我在這事上把我的意見告訴你們,是與你們有益；因為你們下手辦這事,而且起此心意,已經有一年了,

8:11如今就當辦成這事.既有願做的心,也當照你們所有的去辦成.

8:12因為人若有願做的心,必蒙悅納,乃是照他所有的,並不是照他所無的.

8:13我原不是要別人輕省,你們受累,

8:14乃要均平,就是要你們的富餘,現在可以補他們的不足,使他們的富餘,將來也可以補你們的不足,這就均平了.

8:15如經上所記:多收的也沒有餘；少收的也沒有缺.

哥林多後書這裡所提及的,是保羅奔走各地教會呼籲接濟耶路撒冷信徒的事件.由於各種原因,耶路撒冷的教會及信徒生活艱難,於是保羅請哥林多這些經濟條件較好的地區的教會捐款幫助耶路撒冷的教會及信徒.

此時此地,或者有人感到甚麼事情也做不到,因此而充滿無力感.就正如校牧室每年暑假都會舉辦短宣隊,帶領同學往海外協助當地的教會及參與社區服務.今年,不要說不能成團,就是個人旅行也不方便.就在這樣的狀況下,我們招募了一隊同學,在暑期這幾個月,每個星期三到油麻地一個教會辦的扶貧機構參與和學習.同學們清晨六時半往果欄收集果商送出或拋棄的生果,下午按著每位受助者的需要包裝物資,例如有煮食設備的就送白米和冰肉等,沒有煮食設備的只能送上飯盒或乾糧.晚上,就帶同熱飯和飲品往露宿者聚居的地方,和他們分享食物.若對方樂意傾談,也可互相認識,分享生活點滴.

耶穌說:『我實在告訴你們,這些事你們既做在我這弟兄中一個最小的身上,就是做在我身上了.』(太25:40)

靜默有時,言語有時,願上主幫助我們,各按其時,成為美好.



\newpage

\section{平凡信徒肩負不平凡任務}
\label{sec:557}
\textbf{高國雄牧師 - 平凡信徒肩負不平凡任務}
\newline
\newline
\hyperref[sec:556]{< < < PREV SERMON < < <}
~
\hyperref[sec:toc]{[返目錄]}
~
\hyperref[sec:558]{> > > NEXT SERMON > > >}


馬可福音 6:1-13
\newline
\begin{longtable}{cl}
\hline
\hline
章節 & 經文 (和合本修訂版)\\
\hline
6:1 & \begin{tabularx}{0.7\textwidth}{X} 耶穌離開那裡，來到自己的家鄉；門徒也跟從他。 \end{tabularx} \\ \\ \relax
6:2 & \begin{tabularx}{0.7\textwidth}{X} 到了安息日，他在會堂裡教導人。眾人聽見，就很驚奇，說：「這人哪來這本事呢？所賜給他的是甚麼智慧？他手所做的是何等的異能呢？ \end{tabularx} \\ \\ \relax
6:3 & \begin{tabularx}{0.7\textwidth}{X} 這不是那木匠嗎？不是馬利亞的兒子雅各、約西、猶大、西門的長兄嗎？他姊妹們不也是在我們這裡嗎？」他們就厭棄他。 \end{tabularx} \\ \\ \relax
6:4 & \begin{tabularx}{0.7\textwidth}{X} 耶穌對他們說：「先知除了在本鄉、本族和自己的家之外，沒有不被尊敬的。」 \end{tabularx} \\ \\ \relax
6:5 & \begin{tabularx}{0.7\textwidth}{X} 耶穌在那裡不能行甚麼異能，不過為幾個病人按手，治好他們。 \end{tabularx} \\ \\ \relax
6:6 & \begin{tabularx}{0.7\textwidth}{X} 他也詫異他們不信。 \end{tabularx} \\ \\ \relax
6:7 & \begin{tabularx}{0.7\textwidth}{X} 他叫了十二個使徒來，差遣他們兩個兩個地出去，也賜給他們權柄制伏污靈， \end{tabularx} \\ \\ \relax
6:8 & \begin{tabularx}{0.7\textwidth}{X} 並且吩咐他們：途中不要帶食物和行囊，腰袋裡也不要帶錢，除了手杖以外，甚麼都不要帶； \end{tabularx} \\ \\ \relax
6:9 & \begin{tabularx}{0.7\textwidth}{X} 只要穿鞋子，也不要穿兩件內衣。 \end{tabularx} \\ \\ \relax
6:10 & \begin{tabularx}{0.7\textwidth}{X} 他又對他們說：「你們無論到何處，進哪家，就住在哪裡，直到離開那地方。 \end{tabularx} \\ \\ \relax
6:11 & \begin{tabularx}{0.7\textwidth}{X} 若有甚麼地方的人不接待你們，不聽你們，你們離開那裡的時候，要跺掉你們腳上的塵土，證明他們的不是。」 \end{tabularx} \\ \\ \relax
6:12 & \begin{tabularx}{0.7\textwidth}{X} 使徒就出去傳道，叫人悔改， \end{tabularx} \\ \\ \relax
6:13 & \begin{tabularx}{0.7\textwidth}{X} 又趕出許多鬼，用油抹了許多病人，治好他們。 \end{tabularx} \\ \\
[1ex]
\hline
\hline
\end{longtable}
講題:平凡信徒肩負不平凡任務　Home town disciples, mission impossible.

講員:高國雄牧師

經文:馬可福音6:1-13

場合:香港中文大學崇基學院禮拜堂主日崇拜

日期:2021年7月4日

1耶穌離開那裡,來到自己的家鄉；門徒也跟從他. 2到了安息日,他在會堂裡教導人.眾人聽見,就很驚奇,說:「這人哪來這本事呢?所賜給他的是甚麼智慧?他手所做的是何等的異能呢? 3這不是那木匠嗎?不是馬利亞的兒子雅各,約西,猶大,西門的長兄嗎?他姊妹們不也是在我們這裡嗎?」他們就厭棄他. 4耶穌對他們說:「先知除了在本鄉,本族和自己的家之外,沒有不被尊敬的.」 5耶穌在那裡不能行甚麼異能,不過為幾個病人按手,治好他們. 6他也詫異他們不信.耶穌走遍周圍鄉村教導人. 7他叫了十二個使徒來,差遣他們兩個兩個地出去,也賜給他們權柄制伏污靈, 8並且吩咐他們:途中不要帶食物和行囊,腰袋裡也不要帶錢,除了手杖以外,甚麼都不要帶； 9只要穿鞋子,也不要穿兩件內衣. 10他又對他們說:「你們無論到何處,進哪家,就住在哪裡,直到離開那地方. 11若有甚麼地方的人不接待你們,不聽你們,你們離開那裡的時候,要跺掉你們腳上的塵土,證明他們的不是.」 12使徒就出去傳道,叫人悔改, 13又趕出許多鬼,用油抹了許多病人,治好他們.

由馬可福音4:35開始,連續不斷展示出耶穌基督的能力及使命:平靜風浪,格拉森趕鬼,醫好12年經血不止的婦人,教會堂主管葉魯女兒復活,接著就是今天的經文,耶穌基督回鄉,門徒跟從.

回鄉後就在安息日時,進入會堂教導,經文描述(V2)眾人聽見,就很驚奇(原文新約:驚訝/震驚,使大吃一驚)；弟兄姊妹,我們有沒有因為上帝的話語而驚訝過?為什麼沒有/少有呢?

(V2)說:「這人哪來這本事呢?所賜給他的是甚麼智慧?他手所做的是何等的異能呢?

對耶穌基督的教導震驚,對祂所行的異能神蹟驚訝,理應所生的是敬畏,拜服之心,從而渴望著耶穌基督教導自己,幫助自己,拯救自己,正好過去在其他地方一樣…

但耶穌基督故鄉的親友卻因故有的認識,而產生不同的反應及結果…

(V3)這不是那木匠嗎?不是馬利亞的兒子雅各,約西,猶大,西門的長兄嗎?他姊妹們不也是在我們這裡嗎?」他們就厭棄他.

這不是“那”木匠嗎?…

稱呼耶穌基督為“馬利亞的兒子雅各,約西,猶大,西門的長兄”,在他們心目中耶穌沒有任何身份,是他人的附屬,這一種附屬的理解正正代表著他們輕視耶穌,更因為以為自己認識/熟悉,這往往成為追求主最大的障礙:

「聖經的教導我從小就認識了!」

「宗教都係導人向善,基督教大概也差不多…」

「我父母是基督徒,身邊很多朋友也返教會,我知道什麼是信耶穌…」

記得1996年曾經在美國亞特蘭大百週年奧林匹克運動會中參與宣教活動,帶領香港團隊有個多月的時間在當地進行文化宣教,其中最深刻的就是邀請當地人信耶穌時,每一個的回應大概也是「我家是幾代信徒,我嬰孩時已經接受洗禮了,耶穌我當然認識啦…」但當深入一些分享交流時,就發現他們大部份根本已經再沒有回教會,沒有讀聖經,甚至覺得耶穌與自己從來也沒有個人的關係…

自以為認識/熟悉也容易產生故有的印象,以過去的印象/經驗來決定所有:

例子:「三文魚有寄生蟲,日本人從來不吃三文魚魚生」,這其實是60年代以前的故事,經過科技新的發展,已經完全改變,過去十年,日本人最愛魚生除了吞拿魚,就是三文魚…

「教會教導的都是一般導人向善的大道理,牧師未講我也大概估計到內容…」

這一種態度,往往成為了我們接受信仰/信仰成長的最大阻礙.

他們就厭棄他.

他們厭棄耶穌!?(原文新約:他們就生他的氣)

他們為何生氣呢?

見到自己人(熟悉的人)“成功”不是應該慶幸,並且覺得與有榮焉,引以為傲的嗎?

人的自我,妒忌,會造成盲點…

很多人一生追求信仰,就是要看神蹟,認為只要神蹟出現,他就必會專心一意敬畏上帝；

馬可福音8:11-12法利賽人出來盤問耶穌,要求他從天上顯個神蹟給他們看,想要試探他.耶穌心裡深深嘆息,說:「這世代為甚麼求神蹟呢?我實在告訴你們,沒有神蹟給這世代看.」

耶穌基督在這裡所提到:再沒有神蹟,與及經課(V5)耶穌在那裡不能行甚麼異能,不是指耶穌基督不再施行神蹟(實際上祂繼續行神蹟,甚至行最大的神蹟:釘身十架,第三天復活),而是告訴我們,要追求上帝(上帝本身),而不是要追求神蹟(上帝的作為),焦點要清楚；經課中清楚表示他們(鄉里)因見耶穌基督的能力及智慧而驚訝,但他們沒有因此要歸榮耀予上帝,反而更質問/質疑耶穌基督能力的來源…

請留意:由於對主的不信,厭棄,生氣,帶來極嚴重的後果:

(V5)耶穌在那裡不能行甚麼異能

神蹟因此少了!人透過神蹟認識上帝,歸榮上帝的機會少了…

不但如此,他們更因此質疑上帝!

(V2)他手所做的是何等的異能呢?(原文新約:這樣的神蹟竟然透過祂的手發生!)

竟然透過祂的手!!

很多時候我們質疑上帝,是因為我們覺得事情發生得不合我們的認知,邏輯,期望,請謹記:上帝的意念高過我們的意念,就算我們不明白,甚至祂會用我們最不能接受的方法展現祂自己,但請不要懼怕,他所賜的必定是最美好的.

接著經文提到耶穌基督差派門徒兩個兩個出去,醫病趕鬼,竟然

(V8-11)吩咐他們:途中不要帶食物和行囊,腰袋裡也不要帶錢,除了手杖以外,甚麼都不要帶；只要穿鞋子,也不要穿兩件內衣.他又對他們說:「你們無論到何處,進哪家,就住在哪裡,直到離開那地方.若有甚麼地方的人不接待你們,不聽你們,你們離開那裡的時候,要跺掉你們腳上的塵土,證明他們的不是.」

我們今天遠遊絕對不會這樣,就算參與宣教活動,也應該從來沒有牧師要求我們這樣(猜想如若我這樣要求你們,你們也不會跟從…),當然我們也要認識近東遊牧民族背景中的文化,接待客旅是非常重要的事,亦是聖經中提醒必要做的事情；但這般極致的甚麼也不帶,衣服也不穿才就上路,大概也只會被認為是粗梳,“亂來”,太沒有計劃,不合格,與及不同能,但這就是耶穌基督對門徒的教導.

今天我們又如何回應主的教導呢?特別是那些我們未能明白,心底裡其實充滿不信任的託付呢?我們會視而不見?選擇自己喜愛的才接受?解釋逃避?還是願意謙卑地照單全收呢?

為何耶穌基督會這樣吩咐,而門徒又能照做呢?原因只有一個,就是耶穌基督與他們同在.

回到經文的開首,經文沒有表達為何耶穌基督要回鄉,是因為工作忙碌,要回鄉休息嗎?是因為想念家人,要回鄉探親嗎?我們無法知道,但我們清楚願意在安息日時進入會堂教導,被厭棄仍然選擇繼續(V6)走遍周圍鄉村教導人；一個非常簡單直接,又同時是最重要的答案,是因為愛,耶穌基督愛每一個人,每一個地方,當然也深愛自己的成長地,自己的家鄉,儘管不受歡迎,儘管被人厭棄,甚至令人生氣,祂仍然帶著愛去分享,更差派所愛的門徒,去完成不容易的任務.

請留心:耶穌基督因為愛仍然要走遍周圍鄉村教導人,所以祂就差遣門徒承擔承繼；今天,我們同樣地就是那麼平凡的信徒,面對著不容易的世情,肩負起這不平凡任務,求主與我們同在,讓我們有力以愛去改變世界,引領更多人歸向基督.



\newpage

\section{生命中不能承受的輕}
\label{sec:558}
\textbf{林崇智牧師 - 生命中不能承受的輕}
\newline
\newline
\hyperref[sec:557]{< < < PREV SERMON < < <}
~
\hyperref[sec:toc]{[返目錄]}
~
\hyperref[sec:559]{> > > NEXT SERMON > > >}


講題:生命中不能承受的輕　The Unbearable Lightness of Being

講員:林崇智牧師

場合:香港中文大學崇基學院禮拜堂主日崇拜

日期:2021年7月11日

各位弟兄姊妹平安,多謝高牧師邀請我回到母校崇拜中分享,今日講道的題目「生命中不能承受的輕」,是1984年捷克作家米蘭·昆德拉出版的一本小說名稱,他透過當中不同人物,表達他對生命的看法.小說提到其中一位人物名叫湯馬斯的腦外科專家,在一次與朋友閒談聊天時,他即慶講了當蘇聯獨才者史大林實行大規模整肅運動,這惡行帶來人民無盡和巨大的傷害.而當時捷共卻效法當年史太林的做法在捷克推行.他說若推行這邪惡政策的人,若看到人民所受到的痛苦結果,稍有人性的,也應將眼珠挖出來.朋友們認為湯馬斯說得很好,叫他將看法寫出來.最後湯馬斯勉為其難的寫了一篇短文交給報社,編輯也十分讚賞這文章,只是文章等不到出版,蘇聯坦克便鎮壓了捷克的布拉格之春.湯馬斯與太太流亡到瑞士,最後回到捷克重返醫院工作,醫院的領導十分歡迎他回來復職,但有一個小小手續要辦,就是湯馬斯要在一封信上個簽名,表示他是出於無知,寫了那篇批評政府的短文,他願意收回那篇文章.領導跟他說,你的專業不是寫文章,撤回這篇文章對你沒有損失,何況這封信只是內部存檔用,不會向外發表的.湯馬斯沒有甚麼反應,領導便說,這篇文章對一個腦外科專家有甚麼咁重要? 意思是收回一篇文章並沒有甚麼大不了.湯馬斯亳無表情的說了一句「再沒有比這個更重要了」.

最後腦外科專家當然不能復職,轉行做了抹窗工人.簽一個名,執筆一揮,是很輕而易舉的事.作這麼簡單的事,就能做回付上許多年,代價才能成為的腦外科醫生.不做這輕而易舉的簽名,所付的代價是從醫生專業變為抹窗工人,從經濟價值判斷,答案是很明顯的.但為何湯馬斯不能簽這個領導看為輕而易舉的事,但他卻看為再沒有比這樣更重要的事是因為人格問題,他是一個有良知,有信念的人,他可以否定並糾正自己過去所犯的錯失,但不能否定自己所做正確的事.否定自己的信念和良知,湯馬斯會一生看不起自己,無盡地譴責自己的生活下去,「承受生命中不能承受的輕」.湯馬斯一生最大的報酬就是光明正大,心安理得的生活下去.

今日的經課同樣出現兩位承擔與「生命中不能承受的輕」有關的人物,第一位是施洗約翰,一位我們熟識的人物,我們知道他是耶穌的表哥,按聖經說,他一生的使命是要在將要來臨的救世主,預備道路,叫主的百姓因罪悔改,照亮坐在黑暗中或死蔭裡的人.聖經說,當約翰逐漸長大,他身體健壯,住在曠野,穿駱駝毛的衣服,腰束皮帶,吃的是蝗蟲野蜜.及至時候滿足,他走遍約旦河一帶,在曠野大聲宣講,「天國近了,你們要悔改了,結出果子來,與悔改的心相稱」.若不悔改,上帝的審判便臨到.我現在用水給你們施行悔改的洗禮,但有一位在我以後來的,就是給他提鞋也不配.

福音書告訴我們,施洗約翰為耶穌開路作先鋒的工作,十分有果效,那時耶路撒冷,猶太和約旦河地區的人都到約翰那裡,當中甚至有不少有權位名望的法利賽人和撒都該人都來受洗.約翰在曠野傳道不足一年,從人的眼光看,他的事業如日方中,他的成就是輝煌的,他的影響力也很大.他若能繼續留在曠野為耶穌做更多預備人心的工作,不是最好嗎? 但有一天,他聽到獨才的希律王,不合法地娶了他兄弟的妻子希羅底,他是言行一致的人,便公開責備希律,要求他悔改,正如他所講的信息一樣.他仍有那麼多的事要做,有那麼多人需要他,來到曠野聽道洗禮.事實上希律王沒有阻礙他的工作,甚至看來也十分尊重他.他能對希律王的惡事隻眼開隻眼閉,忍耐一o下咪幾好,何必因小失大.結果,他對希律王的斥責帶來曠野傳道工作,一生職事的結束,甚至帶來被殺殉道,死時不到三十歲,值得嗎?

從世俗的價值判斷,對一個獨才君王批評不會帶來改變,但結果明顯是不值得,甚至愚蠢.但相信對施洗約翰,卻是換來忠於上帝的託付,一份不背棄信仰,不受自責的磊落良心.他所做的人可能是不認同,但他得到主耶穌基督的平反和肯定對他作出三個稱讚,(1)「女子所生的,沒有一個比施洗約翰大」(太1:11 A ); (2)「他是最後的一位先知,比其他的先知大得多」(路7:28); (3) 「他是點著的明燈」(約5:35),他的生命已為上帝燃點,照亮他人.

另一位展現承受「生命中不能承受的輕」的風骨氣節,是先知阿摩司.聖經說,他是屬於南面猶大國的人,但卻去了北面以色列在宗教中心伯特利作先知.他成為上帝的代言人,告訴以色列王耶羅波安二世,上帝的審判要臨到,王朝將要滅亡.這種顛覆性的說話被當時支持王室建制,在伯特利擔任大祭司的亞瑪謝所攻擊,因為皇家祭司的信息是平安平安.阿摩司的話對國家安定造成很大衛脅,對大祭司的身份也帶來衝擊,所以亞瑪謝找人到皇帝面前,抹黑阿摩司說他圖謀背叛國王,煽動民心,令國家會出現災難後果.他希望耶羅波安採取行動對付阿摩司,但亞瑪謝亦有兩手準備,若耶羅波安王不採取行動,他也同時派人去對阿摩司說,你做先知是為了找兩餐飯食je ,那麼請你返回有南面猶太國,這裡是北國以色列國,我們的事與你無關,你不可以再在伯特利講預言.若你繼續堅持下去,你的性命也可能難保.阿摩司在北部以色列不受歡迎,他的性命也確隨時有危險.事實上他的專業是牧羊人和農夫不是先知,回到南面自己的家鄉,做回自己的專業,沒有甚麼困難.為何要堅持留在不受歡迎的伯特利做先知?因為他深知一切都不是他所選的,而是耶和華帶領和呼召他到這裡,那他便不能隨意放下上帝的託付,終使不受歡迎,甚至生命有危險.

馬太7:13告訴我們,生命要常在十字路口作出艱難選擇,通往生命的門是窄的,路是小的,找到的人也少,而通往滅亡的門是寬的,路是大的,進去的人很多.要承受生命中不能承受的輕,要接受許多人不認同,不明白,不理解,要付出不輕的代價,需要持守堅定的信念,異乎尋常的信心和勇氣.那麼我們真需要在,小事中實踐,堅守信念,在禱告中經常作操練,倚靠上主的恩典和憐憫,聖靈的能力,作合乎上帝心意的事.



\newpage

\section{中間的牆沒有了}
\label{sec:559}
\textbf{鄧瑞強博士 - 中間的牆沒有了}
\newline
\newline
\hyperref[sec:558]{< < < PREV SERMON < < <}
~
\hyperref[sec:toc]{[返目錄]}
~
\hyperref[sec:560]{> > > NEXT SERMON > > >}


以弗所書 2:11-22
\newline
\begin{longtable}{cl}
\hline
\hline
章節 & 經文 (和合本修訂版)\\
\hline
2:11 & \begin{tabularx}{0.7\textwidth}{X} 所以，你們要記得：從前你們按肉體是外邦人，是「沒受割禮的」；這名字是那些憑人手在肉身上「受割禮的人」所取的。 \end{tabularx} \\ \\ \relax
2:12 & \begin{tabularx}{0.7\textwidth}{X} 要記得那時候，你們與基督無關，與以色列選民團體隔絕，在所應許的約上是局外人，而且在世上沒有指望，沒有神。 \end{tabularx} \\ \\ \relax
2:13 & \begin{tabularx}{0.7\textwidth}{X} 從前你們是遠離神的人，如今卻在基督耶穌裡，靠著他的血，已經得以親近了。 \end{tabularx} \\ \\ \relax
2:14 & \begin{tabularx}{0.7\textwidth}{X} 因為他自己是我們的和平，使雙方合而為一，拆毀了中間隔絕的牆，而且以自己的身體終止了冤仇， \end{tabularx} \\ \\ \relax
2:15 & \begin{tabularx}{0.7\textwidth}{X} 廢掉那記在律法上的規條，為要使兩方藉著自己造成一個新人，促成了和平； \end{tabularx} \\ \\ \relax
2:16 & \begin{tabularx}{0.7\textwidth}{X} 既在十字架上消滅了冤仇，就藉這十字架使雙方歸為一體，與神和好， \end{tabularx} \\ \\ \relax
2:17 & \begin{tabularx}{0.7\textwidth}{X} 並且來傳和平的福音給你們遠處的人，也傳和平給那些近處的人， \end{tabularx} \\ \\ \relax
2:18 & \begin{tabularx}{0.7\textwidth}{X} 因為我們雙方藉著他，在同一位聖靈裡得以進到父面前。 \end{tabularx} \\ \\ \relax
2:19 & \begin{tabularx}{0.7\textwidth}{X} 這樣，你們不再是外人或客旅，是與聖徒同國，是神家裡的人了， \end{tabularx} \\ \\ \relax
2:20 & \begin{tabularx}{0.7\textwidth}{X} 被建造在使徒和先知的根基上，而基督耶穌自己為房角石， \end{tabularx} \\ \\ \relax
2:21 & \begin{tabularx}{0.7\textwidth}{X} 靠著他整座房子連接得緊湊，漸漸成為在主裡的聖殿。 \end{tabularx} \\ \\ \relax
2:22 & \begin{tabularx}{0.7\textwidth}{X} 你們也靠他同被建造，成為神藉著聖靈居住的所在。 \end{tabularx} \\ \\
[1ex]
\hline
\hline
\end{longtable}
講題:中間的牆沒有了　The Wall Is Gone

經文:以弗所書2:11-22

講員:鄧瑞強博士

場合:香港中文大學崇基學院禮拜堂主日崇拜

日期:2021年7月18日

各位弟兄姊妹,早安.在這疫症裡,祝願各人平安.

讓我講一下羅密歐與茱麗葉的故事.

劇本的開頭,有一首十四行詩,交代仇恨將要譜寫出何等樣的悲劇局面.

莎士比亞這樣寫:故事發生在意大利北部的名城維羅納(Verona),那裡,有兩家威望相當的家族.兩家的舊仇引發新怨.跟著,莎士比亞的詩句原文是:「Where civil blood makes civil hands unclean」,梁實秋的翻譯是:「使得市民的血把清白的手弄髒」.我想來想去,不如將這句意譯:「本應和平地生活在一起的文明人最後卻以流血收場」.命運注定,這兩家世仇卻生下一對不幸的戀人.他們以悲慘淒涼的死亡,去掩埋兩家的爭執.這對戀人那帶著死亡標記的愛,可怕地消逝.而這兩大家族的恩恩怨怨,就以其兒女的死亡來平息.

多麼令人傷感的開場白.世事多磨,總不能好好地讓人過平安的生活.兩大家族的新仇舊恨,讓雙方的兒女不能愛,最終只能以悲劇收場.這是命運使然,還是人類永遠不願意從歷史中汲取教訓?

這齣戲劇的結尾是這樣的:最後,卡普萊(Capulet)家族死了女兒茱麗葉,蒙太古(Montague)家族死了兒子羅密歐.審理這死亡事件的親王說:「這兩家世仇在哪裡?卡普萊,蒙太古,看看你們的仇恨招來了多大的懲罰,上天藉著愛情,奪去了你們心愛的人.而我,由於忽視了你們的積怨,也失去了一對親戚,大家都受到了懲罰.」最後,兩大家族和解,但,這已是所有人受到了懲罰之下的悲愴的和解.

人與人之間,為何總是這樣?我不信你,你仇視我,然後,以不同理由,築起種種圍牆,將人與人隔開.昔日有柏林圍牆,現今仍存的,有以色列與巴勒斯坦之間的圍牆,有美國與墨西哥之間的圍牆,還有更多的圍牆,在人的思想與心靈深處.圍牆外的人看圍牆裡的人,或者圍牆裡的人看圍牆外的人,都認為對方不是人.圍牆的作用,就是將對方隔開,免得他侵害我.而就算他存在,也要藉著圍牆,令他在我的視野中消失.圍牆無處不在,多麼可悲.

小時候,很喜歡看有關美國西部牛仔的電視劇集.每齣劇集,都出現牛仔與「紅番」的衝突場面.小時候的我,被劇集的敘事角度影響,總認為「紅番」是野蠻的,無情的,不講理的.當牛仔射殺「紅番」時,會興奮莫名,覺得壞人被消滅了.一位保加利亞裔的法國思想家托多洛夫(Tzvetan Todorov),寫了一本書《美洲的征服》(The Conquest of America)研究歐洲人對美洲的征服.他提及一種「殺害他者的認知方式」(the understanding-that-kills),就是:西班牙人認為墨西哥的原住民為一種無文化的野蠻動物,這種對他人的絕對負面的評價,鼓勵了西班牙人對這些原住民採取滅絕的手段.人與人的隔離與圍牆,最終導致悲劇.近來,在加拿大,在很多為了同化原住民的寄宿學校裡,發現很多小朋友的骸骨.這些學校本身,當然也是一種圍牆.這些小朋友的死,背後有很多故事,但少不了的,是人與人之間的隔膜與仇恨.

耶穌來,要從根源上,取消人與人之間的仇恨,拆毀人與人中間的圍牆.

今日的講道經文:以弗所書2:11-22(和合本修訂版)

2:11　所以,你們要記得:從前你們按肉體是外邦人,是「沒受割禮的」；這名字是那些憑人手在肉身上「受割禮的人」所取的.

2:12　要記得那時候,你們與基督無關,與以色列選民團體隔絕,在所應許的約上是局外人,而且在世上沒有指望,沒有神.

2:13　從前你們是遠離神的人,如今卻在基督耶穌裡,靠著他的血,已經得以親近了.

2:14　因為他自己是我們的和平,使雙方合而為一,拆毀了中間隔絕的牆,而且以自己的身體終止了冤仇,

2:15　廢掉那記在律法上的規條,為要使兩方藉著自己造成一個新人,促成了和平；

2:16　既在十字架上消滅了冤仇,就藉這十字架使雙方歸為一體,與神和好,

2:17　並且來傳和平的福音給你們遠處的人,也傳和平給那些近處的人,

2:18　因為我們雙方藉著他,在同一位聖靈裡得以進到父面前.

2:19　這樣,你們不再是外人或客旅,是與聖徒同國,是神家裡的人了,

2:20　被建造在使徒和先知的根基上,而基督耶穌自己為房角石,

2:21　靠著他整座房子連接得緊湊,漸漸成為在主裡的聖殿.

2:22　你們也靠他同被建造,成為神藉著聖靈居住的所在.

這段文字,是站在「以色列人」的角度寫的.「以色列人」,自命是「神的選民」,掌控了真理,站在道德高地上,在這超凡的位置,俯視其他人.

「其他人」是什麼人?

經文說:「從前你們按肉體是外邦人,是『沒受割禮的』」.在生命裡,沒有這神聖記號.「要記得那時候,你們與基督無關,與以色列選民團體隔絕,在所應許的約上是局外人,而且在世上沒有指望,沒有神.」用現今的講法是:你們什麼也不是,沒有永恆意義,沒有生命的內容,沒有神聖,沒有將來,只是一團會活動的物質.

在很多人眼中,人不是一樣的.人分成兩批人.一批人,自命是神聖的,帶有天命的,是人類幸福的主持人,將人類帶往天國.另一批人,被視為是低下的,沒有內在意義的,是被操控的對象.在第一批人眼中,第二批人根本不算是人.而每個人都認為自己是第一批人,是帶有天命的,對方是第二批人,死不足惜.就如在希特拉眼中,猶太人不算是人,只配送入焚化爐.這兩批人之間,沒有溝通,也不可能溝通,中間是高高的圍牆,是雙方不能越過的圍牆,也是死也不願越過的圍牆.

這是我們人類的光景.

基督來到世上,要改變這光景.

2:14　因為他自己是我們的和平,使雙方合而為一,拆毀了中間隔絕的牆,而且以自己的身體終止了冤仇,

2:15　廢掉那記在律法上的規條,為要使兩方藉著自己造成一個新人,促成了和平；

2:16　既在十字架上消滅了冤仇,就藉這十字架使雙方歸為一體,與神和好,

基督以十架的寬恕與接納,消滅了人與人之間的不信,猜忌,冤仇,拆毀了人與人中間隔絕的牆,成就了真正的和平.

這對我們有什麼提示呢?

我借用意大利神學家Armido Rizzi的一些觀點來作一點說明.(以下觀點引自Armido Rizzi, “Reflections for an Anthropology of Peace,” Religious Studies Bulletin 3 no. 3 (Sep 1983): 121-129.)

Rizzi認為,有兩種關於「和平」的觀點.第一種觀點是:上天賜予人類圓滿的和平,某些被「揀選」的人領悟這和平,並受命去守護這和平.破壞和平的人,是壞人,是不義的.守護和平的人,是「選民」,當「和平」受威脅時,可以發動「聖戰」,去消滅壞人.在「聖戰」裡,自命為「選民」者,其行動必然是正義的,也有充分的合法性.在這觀點裡,為了守護和平,便需要時刻準備作戰.問題是:人人都覺得自己是「選民」,於是便時刻準備去與另一方作戰.

第二種觀點是:基於現實上不同的人生活在一起,為了安全,彼此必須訂立一社會契約.這類近盧梭(Jean-Jacques Rousseau)講的「社會契約論」.大家約定,你不攻擊我,我也不攻擊你.為了確定大家都遵守契約,於是訴之國家的權威,及法律的制裁,如此,人與人之間的行為受到規範.國與國之間,也是這樣,以國際的法律和條約,去規範彼此的行為.

Rizzi認為以社會契約去維持和平,最終也會失敗.失敗的理由有二.第一個理由,是從原則上來說的,這種社會契約的本質,只是為了防範對方的威脅,從沒有善意地接待對方的意圖.如此,「和平」只是一種策略,最終為的是自保.這是一種極度脆弱的和平.雙方的強弱一旦失去平衡,強的一方很容易便會掙脫制約,吞吃弱的一方.第二個理由,是從事實上來說的,由於現在是價值真空的時代,人與人之間,國與國之間,再沒有一更高的價值權威作衝突時的公正裁決,於是,大家要做的,是強化自己的武裝,預備自己去面對隨時出現的戰爭.

無論是採取那一種「和平」的觀點,「和平」都是脆弱的,因為這兩種觀點,都隱藏地視對方為敵人,彼此都要不斷強化自己的力量,努力震攝對方.這像兩隻凶猛的鬥狗,都準備撲向對方,現在只是各自被幼細的繩索拉住,但這繩索始終有斷的一天.

在這劍拔弩張的緊張關頭,主耶穌問了一個神學問題:可否視對方為朋友,而不是敵人?

主耶穌在路加福音第10章,講了一個「好撒馬利亞人」的故事.故事的重點:不是找出誰是朋友,然後與他「埋堆」.若如此,你只是加入一個與自己同一陣線的幫派,去對付你的敵人.主耶穌的重點是,去成為別人的朋友,那怕他是敵人.要在敵意中,創造友誼.唯有這樣,人間一切衝突的根源,才被清除.

我再讀經文:

2:14　因為他自己是我們的和平,使雙方合而為一,拆毀了中間隔絕的牆,而且以自己的身體終止了冤仇,

2:15　廢掉那記在律法上的規條,為要使兩方藉著自己造成一個新人,促成了和平；

2:16　既在十字架上消滅了冤仇,就藉這十字架使雙方歸為一體,與神和好,

主耶穌以十架受苦的身體,拆毀人與人中間隔絕的牆,終止敵對雙方的冤仇.耶穌在這裡,展示一種新人性.也在我們的生命中,創造這種新人性.這種新人性的特徵,是主動放下仇恨,甚至因此而受苦,藉此伸出友誼之手,甘願成為對方的朋友,那怕對方本是敵人.

自命為「正義的選民」去儆惡懲奸,或採納種種「和平的契約」作為手段,都不能清除人與人之間的仇恨與圍牆,因為這兩種做法,其實都暗暗地強化了敵對的意識.唯有徹底的人性改造,才能終極消除人間的敵意.這是將「仇恨」的人性,轉化為「仁愛」的人性.這不是現實的政治計算,也不是離地的空談,而是在神的救贖下人性的變革.為了「愛敵人」而受苦,這不是人力所能為,必須深深被基督的十架大愛震撼,被聖靈的大能充滿,才能有這人性的變革.

這種一種徹底的人性變革.

2:18　因為我們雙方藉著他,在同一位聖靈裡得以進到父面前.

2:19　這樣,你們不再是外人或客旅,是與聖徒同國,是神家裡的人了,

2:20　被建造在使徒和先知的根基上,而基督耶穌自己為房角石,

2:21　靠著他整座房子連接得緊湊,漸漸成為在主裡的聖殿.

2:22　你們也靠他同被建造,成為神藉著聖靈居住的所在.

首先,原先敵對的雙方,在同一位聖靈的引導下,來到神面前.一切恩怨,在至高的神面前得以揭露呈現.在神面前,無人是絕對正義的.

然後,任何人也與「聖徒同國」,成為神國的公民.因膚色,種族,語言的不同而造成的分疆分界不再需要了,我們是同一國度的人.

然後,任何人也成為「神家裡的人」,不單同國,更進而成為同一家庭的兄弟姊妹,我們是血肉相連的.

再然後,人與人的共處,成為神臨在的聖殿.我們與對方的共同生活,不是展露人性的罪性和邪惡,而是展露神的仁愛和接納.

最後,成為「神藉著聖靈居住的所在」,人與人的共在,流露出聖靈的生命.

這一切看似不可能,這一切看似離地,這一切看似天方夜譚,但在耶穌的十架中,成為我們所看見的事實.

最後,我分享克羅地亞裔美國籍神學家沃弗(Miroslav Volf)的《擁抱神學》.

在塞爾維亞與克羅地亞對敵的日子,雙方都受到很大的苦.敵對雙方,能否重建關係?如何能達致終極和平?在這背景下,沃弗思考聖經裡「浪子的故事」.在這故事裡,小兒子得罪了父親,但父親從沒有棄絕這兒子,縱使對方真的做錯了.父親總是張開手,接納小兒,施以恩典,主動恢復關係.小兒子當然要為做過的錯事付代價,但這是後話.父親首先做的,不是追討錯誤,而是擁抱.這擁抱表達了新的人性,一種「仁愛」的人性,先對他人「信任」,這樣才能「盼望」一全新的和平局面.而這就是主耶穌的十架行動對這衝突連場的世界作出的回應,這也是基督教信仰對這不幸的世界的重大意義.

沃弗的老師莫特曼(Moltmann)問沃弗:你能擁抱塞爾維亞的戰士嗎?

沃弗回答:我自己不能,但作為耶穌基督的跟從者,我應該能夠做得到.

2:14　因為他自己〔耶穌〕是我們的和平,使雙方合而為一,拆毀了中間隔絕的牆,而且以自己的身體終止了冤仇,

2:15　廢掉那記在律法上的規條〔即:超越人間種種律則〕,為要使兩方藉著自己造成一個新人,促成了和平.

願榮耀歸上主,願人間享和平.



\newpage

\section{耶穌比曹操更快}
\label{sec:560}
\textbf{林豪恩先生 - 耶穌比曹操更快}
\newline
\newline
\hyperref[sec:559]{< < < PREV SERMON < < <}
~
\hyperref[sec:toc]{[返目錄]}
~
\hyperref[sec:561]{> > > NEXT SERMON > > >}


講題:耶穌比曹操更快　Jesus is faster than Cao Cao

講員:林豪恩先生

經文:約翰福音6

場合:香港中文大學崇基學院禮拜堂主日崇拜

日期:2021年7月25日

一. 引言

某天,有位小朋友考我IQ題:

誰是跑得最快的人?

我偷偷地google一下,然後答:牙買加飛人博爾特,他跑100米只需要9.572秒.

她說:錯.

我問:怎可能錯呢?難道有人比他跑得更快嗎?

她說:有,就是曹操.

我問:為甚麼是曹操?

她說:因為「一講曹操,曹操就到」.

讀今天的福音經課時,我發現有人比曹操更快:

6:16到了晚上,他的門徒下海邊去,

6:17上了船,要過海往迦百農去.天已經黑了,耶穌還沒有來到他們那裡.

6:18忽然狂風大作,海就翻騰起來.

6:19門徒搖櫓,約行了十里多路,看見耶穌在海面上走,漸漸近了船,他們就害怕.

6:20耶穌對他們說:「是我,不要怕!」

6:21門徒就喜歡接他上船,船立時到了他們所要去的地方.

6:21節說,耶穌上船,船就立時到了他們所要去的地方.由此可見,耶穌不但自己立時到達,更能夠令同船的人到立即到達.試想想,與耶穌一起,可以省卻很多交通費和交通時間.如此美妙的境界如何可能呢?我想起叮噹(現在稱為多啦A夢)的隨意門.隨意門是最受人喜愛道具的第一名,在叮噹故事中使用率位列第三,僅次於竹蜻蜓等時光機.只要心中想著想去的地點,放置隨意門,出發地與目的地的空間就會靠攏,打開門,跨過去就到達目的地了.

二. 耶穌叮噹化

叮噹故事的確是百看不厭,它帶給觀眾的快樂正是主題曲所唱出的:

人人期望可達到,我的快樂比天高,

人人如意開心歡笑,跳進美夢尋獲美好.

爬進奇妙口袋裡,你的希望必得到,

離奇神化不可思議,心中一想就得到.

無論是我們看過叮噹卡通的現代人,或是耶穌年代的古代人,看到耶穌在今天的福音經課中所做的事情,都難免產生對叮噹的想像,就是「離奇神化不可思議,心中一想就得到」.耶穌展示的第一件道具是我們剛才提及過的類似隨意門的「立時到達船」,有了這件道具,就不再受時空的限制；第二件道具記載在上文:

6:3耶穌上了山,和門徒一同坐在那裡.

6:4那時猶太人的逾越節近了.

6:5耶穌舉目看見許多人來,就對腓力說:「我們從哪裡買餅叫這些人吃呢?」(

6:6他說這話是要試驗腓力；他自己原知道要怎樣行.)

6:7腓力回答說:「就是二十兩銀子的餅,叫他們各人吃一點也是不夠的.」

6:8有一個門徒,就是西門‧彼得的兄弟安得烈,對耶穌說:

6:9「在這裡有一個孩童,帶著五個大麥餅,兩條魚,只是分給這許多人還算甚麼呢?」

6:10耶穌說:「你們叫眾人坐下.」原來那地方的草多,眾人就坐下,數目約有五千.

6:11耶穌拿起餅來,祝謝了,就分給那坐著的人；分魚也是這樣,都隨著他們所要的.

6:12他們吃飽了,耶穌對門徒說:「把剩下的零碎收拾起來,免得有糟蹋的.」

6:13他們便將那五個大麥餅的零碎,就是眾人吃了剩下的,收拾起來,裝滿了十二個籃子.

如果你要給這件道具起名,你會叫它做甚麼呢?叮噹故事中也有類似的道具,叫做「美食餐巾」,或「美食桌布」,只要把桌布鋪在地或桌上,說出你想要吃的食物,馬上就會跑出來要吃的食物,而且味道非常好.

叮噹在我們的時代受歡迎,其實耶穌也在他的時代更受歡迎.人們看見耶穌能夠用五條魚兩個餅變出許多食物餵飽他們,就有以下的反應:

6:14眾人看見耶穌所行的神蹟,就說:「這真是那要到世間來的先知!」

6:15耶穌既知道眾人要來強逼他作王,就獨自又退到山上去了.

這反應一點也不奇怪,因為故事開始時,作者已經明言群眾正是為了耶穌的神蹟而來的:

6:1這事以後,耶穌渡過加利利海,就是提比哩亞海.

6:2有許多人因為看見他在病人身上所行的神蹟,就跟隨他.

只要有了耶穌,「你的希望必得到」；只要有了耶穌,「心中一想就得到」.教會傳講的福音是不是這樣的呢?你是不是這樣認識耶穌的呢?你是不是追隨一位這樣的耶穌呢?

日子一天一天的過去,心中想到而沒得到的越來越多,希望而得不到的也幾來起來,有沒有累積了失望呢?有沒有感到疑惑呢?有一本書暢銷書的簡介如此提問:「你是否曾經感覺到,神對很多人兌現了祂的承諾,卻似乎漏掉了你?你是否曾經疑惑過,為什麼神施展了很多奇蹟,奇蹟卻從來不曾降臨在你身上?當你發現,神好像回應了每個人的祈求,卻獨獨不回應你的時候,你要怎麼辦?」

於是,作者馬克牧師追問:「到底該怎麼祈禱,才能得到神的回應,讓你夢想成真呢?」在這本名為「勇敢告訴神,讓祂成就你的夢想」的書,作者提供了實現夢想的三大重點:(1)大膽夢想,(2)用力祈禱,(3)積極思考,並告訴讀者有效祈禱的訣竅,包括:不做模糊的祈禱,具體化祈禱的內容,找到合適的地方與專屬自己的祈禱姿勢……等等.模糊的祈禱的例子是:上帝,願你賜我一輛車.不模糊又具體的祈禱的例子是:上帝,願你賜我一輛全新BMW旗艦豪華房車THE 7,配備主動式轉向LED頭燈,頂級Nappa真皮材質豪華座椅,後座配備更軟更舒適加厚地毯 ,車身重量1,880(kg),最高時速250 (km/h),耗油量 (10.1l/100km),白色……耶穌可能會如何回應呢?

三. 一個追,一個走

6:15耶穌既知道眾人要來強逼他作王,就獨自又退到山上去了.

一個追,一個走；一個走,一個追.如此追追走走的遊戲,最終會令人憤怒的,甚至反面相向,反目成仇.就算控制得到憤怒,也會被失望侵蝕心靈,放棄期望,也失掉熱情.要正式離婚嗎?既有手續又有成本,還要向鄉親父老叔伯兄弟解釋.不如就讓大家「你有你的生活,我有我的忙碌」吧!「但假如有一日,真係係教會偶然咁參加崇拜,我地會唱下歌,聽下道,然後代禱時候已經唔知講乜ye 好.因為你會發現,我已經改變,正如我可能唔再對你有期望,或者唔覺得你會滿足我的期望.」

入教已深的信徒,或者,我們開始同意古人所說:

日出而作.

日入而息.

鑿井而飲.

耕田而食.

帝力於我何有哉.

古人所說的是遙遠的皇帝,我們所說的是天上的上帝.「帝力於我何有哉?」

四. 「他自己原知道要怎樣行」

6:5耶穌舉目看見許多人來,就對腓力說:「我們從哪裡買餅叫這些人吃呢?」(

6:6他說這話是要試驗腓力；他自己原知道要怎樣行.)

耶穌對群眾當下的需要並非視而不見,那一次,他遇到很多沒有食物的人而提供了食物給他們,只是他道成肉身來到人間的目的並不是作一個食物生產工場；很多次,他遇到患病的人而醫治了他們,只是他不是來作一個醫生.其實,在人群當中,有食物的人是可以分享食物給沒有的人的；在人群當中,有能力的人是可以照顧沒有能力的人的.這個暑假,有同學參加了社區扶貧服務.同學們都好奇,分發給街坊的食物和物資是從可而來的呢?原來,就是從社區本身而來的,有果欄檔主捐出賣相不好的生果,有餐廳送出飯盒,有街坊獻出他們的時間和心力.耶穌基督呢?他來到人間也回應人們的需要,但更重要的是完成人類自己不能完成的任務,實踐人類失落了的使命.

五. 人還能有甚麼希望呢?

如果把人類的溫飽需要和醫療需要推回給人類自己解決,豈不是如那首古詩所言:「日出而作／日入而息／鑿井而飲／耕田而食／帝力於我何有哉.」潘霍華在一九四四年六月二十七日的信提出:「基督教復活的希望與神秘論的希望不同之處,就在基督教的希望是以全新的方式把人送回世界上去生活.」這是一個怎樣的世界呢?潘霍華指出,這是一個及齡的世界,即是一個已經成長了的世界.在七月十六日的另一封信中,他說:「這及齡的世代迫著我們誠實地去認清我們與上帝面對面的處境.上帝正在教導我們,在沒有祂同在時仍要好好地生活.」甚麼時候是「沒有祂同在時」呢?或者就是有人再去尋找祂變出食物而祂走避的時候,或者就是有人再去尋找祂醫治疾病而祂離開的時候,或者就是有人要擁祂作王而祂退到山上去的時候.潘霍華繼續說:「與我們同在的上帝,就是離棄我們的上帝.這位使我們生存在這世界上,讓我們不需以祂為運作假設的那位上帝,也就是我們得以永遠站立在祂面前的上帝.」在結束這封信時,他寫道:「至此為止,我們可以說,我們所描述的及齡的世界,在於放棄對上帝的錯誤觀念,並且準備為聖經所指示的上帝而奮鬥,上帝以祂的軟弱去征服世界的強權,這必須作為我們對聖經的「屬世」的解釋的起點.」

潘霍華呼籲我們,「在沒有祂同在時仍要好好地生活.」好好地生活的力量從何而來呢?今天的書信經課以弗所書3章16節使徒保羅如此為我們祈禱:

3:16為要他按著他豐盛的榮耀,藉著他的靈,使你們內心的力量剛強起來；

3:17又要他使基督因著你們的信住在你們心裡,使你們既在愛中生根立基,

3:18-19能夠和眾聖徒一同明白基督的愛是何等的長,闊,高,深,並知道這愛是超過人的知識所能測度的,為要使你們充滿上帝一切的豐盛.

3:20上帝能照著運行在我們心裡的大能充充足足地成就一切,超過我們所求所想的.

使徒保羅說明「充充足足地成就一切,超過我們所求所想的」就是:

3:21願他在教會中,並在基督耶穌裡,得著榮耀,直到世世代代,永永遠遠.阿們!



\newpage

\section{鳥為食亡}
\label{sec:561}
\textbf{高國雄牧師 - 鳥為食亡}
\newline
\newline
\hyperref[sec:560]{< < < PREV SERMON < < <}
~
\hyperref[sec:toc]{[返目錄]}
~
\hyperref[sec:562]{> > > NEXT SERMON > > >}


約翰福音 6:24-35
\newline
\begin{longtable}{cl}
\hline
\hline
章節 & 經文 (和合本修訂版)\\
\hline
6:24 & \begin{tabularx}{0.7\textwidth}{X} 這時眾人見耶穌和門徒都不在那裡，就上了船，往迦百農去找耶穌。 \end{tabularx} \\ \\ \relax
6:25 & \begin{tabularx}{0.7\textwidth}{X} 他們在海的對岸找到他後，對他說：「拉比，你幾時到這裡來的？」 \end{tabularx} \\ \\ \relax
6:26 & \begin{tabularx}{0.7\textwidth}{X} 耶穌回答他們說：「我實實在在地告訴你們，你們找我，並不是因見了神蹟，而是因吃餅吃飽了。 \end{tabularx} \\ \\ \relax
6:27 & \begin{tabularx}{0.7\textwidth}{X} 不要為那會壞的食物操勞，而要為那存到永生的食物操勞。這食物是人子要賜給你們的，因為父神已印證了。」 \end{tabularx} \\ \\ \relax
6:28 & \begin{tabularx}{0.7\textwidth}{X} 於是他們問他：「我們該做甚麼才算是做神的工作呢？」 \end{tabularx} \\ \\ \relax
6:29 & \begin{tabularx}{0.7\textwidth}{X} 耶穌回答，對他們說：「信神所差來的，這就是神的工作。」 \end{tabularx} \\ \\ \relax
6:30 & \begin{tabularx}{0.7\textwidth}{X} 於是他們對他說：「你行甚麼神蹟，好讓我們看見而信你呢？你到底要做甚麼呢？ \end{tabularx} \\ \\ \relax
6:31 & \begin{tabularx}{0.7\textwidth}{X} 我們的祖宗在曠野吃過嗎哪，如經上寫著：『他從天上賜下糧食來給他們吃。』」 \end{tabularx} \\ \\ \relax
6:32 & \begin{tabularx}{0.7\textwidth}{X} 於是耶穌對他們說：「我實實在在地告訴你們，那從天上來的糧不是摩西賜給你們的，那從天上來的真糧是我父賜給你們的。 \end{tabularx} \\ \\ \relax
6:33 & \begin{tabularx}{0.7\textwidth}{X} 因為神的糧就是那位從天上降下來，並且賜生命給世界的。」 \end{tabularx} \\ \\ \relax
6:34 & \begin{tabularx}{0.7\textwidth}{X} 於是他們對他說：「主啊，請常常把這糧賜給我們！」 \end{tabularx} \\ \\ \relax
6:35 & \begin{tabularx}{0.7\textwidth}{X} 耶穌對他們說：「我就是生命的糧。到我這裡來的，絕不飢餓；信我的，永不乾渴。 \end{tabularx} \\ \\
[1ex]
\hline
\hline
\end{longtable}
講題:鳥為食亡 　To die for

講員:高國雄牧師

經文:約翰福音6:24-35

場合:香港中文大學崇基學院禮拜堂主日崇拜

日期:2021年8月1日

24眾人見耶穌和門徒都不在那裡,就上了船,往迦百農去找耶穌. 25既在海那邊找著了,就對他說:「拉比,是幾時到這裡來的?」 26耶穌回答說:「我實實在在地告訴你們,你們找我,並不是因見了神蹟,乃是因吃餅得飽. 27不要為那必壞的食物勞力,要為那存到永生的食物勞力,就是人子要賜給你們的,因為人子是父神所印證的.」 28眾人問他說:「我們當行甚麼才算做神的工呢?」 29耶穌回答說:「信神所差來的,這就是做神的工.」 30他們又說:「你行甚麼神蹟,叫我們看見就信你；你到底做甚麼事呢? 31我們的祖宗在曠野吃過嗎哪,如經上寫著說:『他從天上賜下糧來給他們吃.』」 32耶穌說:「我實實在在地告訴你們,那從天上來的糧不是摩西賜給你們的,乃是我父將天上來的真糧賜給你們. 33因為神的糧就是那從天上降下來,賜生命給世界的.」34他們說:「主啊,常將這糧賜給我們!」 35耶穌說:「我就是生命的糧.到我這裡來的,必定不餓；信我的,永遠不渴.

群眾尋找耶穌基督,原因何在呢?

V22第二日,站在海那邊的眾人知道那裡沒有別的船,只有一隻小船,又知道耶穌沒有同他的門徒上船,乃是門徒自己去的.

原本以為耶穌仍然在他們當中,心靈就穩妥…

(五餅二魚神蹟後)V16到了晚上,他的門徒下海邊去, 17上了船,要過海往迦百農去.天已經黑了,耶穌還沒有來到他們那裡.(V18-21耶穌在水面上行走)

為何前一晚耶穌沒有與門徒同行?

醫病趕鬼忙碌到深夜,甚至耶穌何時/如何離開了也不知道…

V25既在海那邊找著了,就對他說:「拉比,是幾時到這裡來的?」

這問題的重點在關心什麼呢?是一個簡單的關心?不明白的疑問?

因為自我所認知與實際情況有差異.

V26耶穌回答說:「我實實在在地告訴你們,你們找我,並不是因見了神蹟,乃是因吃餅得飽.

耶穌是指責他們嗎?為”搵食”而生活有錯嗎?

以耶穌基督的能力,當然隨時可以讓人吃飽(剛剛才經歷了五餅二魚的神蹟)…

記得幼年時學習過一個故事:從前有個叫龐太祖的人,不僅家財萬貫,子孫也很多,可以稱為財丁兩旺.不過,到了晚年,他卻被這些錢財折磨得無法安穩過日子.原來他的兒子和孫子們,見龐太祖年老不中用,全都盤算著如何侵占家產.每個人都私下趕製各種武器,準備等龐太祖死後互相廝殺來搶奪財產.龐太祖知道後,非常害怕到時會家破人亡,人財兩空.想來想去,他決定暗中雇用兩位挑夫,將所有的金銀挑去山上埋掉.沒想到挑夫們也是貪心的人,他們人在挑金銀,心裡卻是各自盤算要如何獨吞埋在山上的金銀.有一天早上,瘦挑夫說:「每天這樣挑擔,實在很辛苦,你先上山,我去買些飯菜,來慰勞一下,好不好?」胖挑夫點頭同意.瘦挑夫買完酒菜後,又到藥店買毒藥,放進其中一碗飯菜裡,接著到山上.瘦挑夫把有毒的飯菜端給胖挑夫說:「早上都是你在挖坑,辛苦了,請多吃一些.」胖挑夫很餓,就將飯菜吃個精光.胖挑夫吃完說:「早上都是我在挖坑,現在你先挖,我休息一下就去幫忙.」瘦挑夫就拿起鋤頭去挖坑.沒想到胖挑夫居然拿起扁擔,把瘦挑夫打死.胖挑夫正高興時,毒藥藥性發作,也一命嗚呼,這就叫做「人為財死!」,山上的鳥飛來吃兩位挑夫的軀體,吃到胖挑夫的鳥,也全部中毒死亡,這就叫做「鳥為食亡!」[1]

故事提醒我們很重要的道理,生命不應該/不可以單為飲食,一定要持守更重要的價值及意義.

V27不要為那必壞的食物勞力,要為那存到永生的食物勞力,就是人子要賜給你們的,因為人子是父神所印證的.」28眾人問他說:「我們當行甚麼才算做神的工呢?」

“做神的工”是什麼意思?為何突然提到”做工”呢?

原文新約:V27你們不要為會消失的食物工作,而是要為會繼續存留到永恆生命的食物,人子將把這賜予你們,因父上帝已蓋印.

耶穌基督清楚指出,有二種食物:會消失的食物與會繼續存留到永恆生命的食物；首先我們要知道,這二種食物也是上帝所賜的,也是由祂而來,為祝福我們生命而賜的,所以不是好與壞的分別,是全都是好的!

所以應帶著感恩的心領受就可以了；但耶穌基督教導我們只要為會繼續存留到永恆生命的食物工作,因為是上帝親自蓋印,這才是上帝的工.

V29耶穌回答說:「信神所差來的,這就是做神的工.」

信與做工的關係:接受上帝所差來的,就即是接受耶穌就是基督,是從上帝而來,祂就是上帝,這一種接納及認信是最重要/首要的…

V30他們又說:「你行甚麼神蹟,叫我們看見就信你；你到底做甚麼事呢?

群眾問這問題,好明顯在回應耶穌V26的質問:

事實上我們知道,當日群眾已經看到很多神蹟,醫病趕鬼,五餅二魚也是很厲害的神蹟,如果群眾是為餅而來,也是為神蹟/已經歷了的神蹟而來,為何還問耶穌呢?

請留心群眾所問的重要:「你行甚麼神蹟,叫我們看見就信你；你到底做甚麼事呢?」換一句話講,大概就是在問:「你會耍出什麼把戲/招數,叫我們看見就認同你有能力,你還有什麼本領呢?」

這明顯是向耶穌的能力作出挑戰,因為他們接著就提到舊約在曠野的事情:

V31我們的祖宗在曠野吃過嗎哪,如經上寫著說:『他從天上賜下糧來給他們吃.』」

群眾挑戰耶穌幾個概念:

首先:吃飽很重要,在曠野40年就是吃嗎哪吃飽的,耶穌看似輕看這重要性.

然後指出摩西比耶穌厲害,因為摩西使天降嗎哪,吃飽40年當然比五千人吃飽一次厲害得多.

所以挑戰耶穌,如何你真來最偉大的一位,就應該有更大的本領…

耶穌如何回應他們呢?

V32耶穌說:「我實實在在地告訴你們,那從天上來的糧不是摩西賜給你們的,乃是我父將天上來的真糧賜給你們. 33因為神的糧就是那從天上降下來,賜生命給世界的.」

真正的糧是由天父所賜的,並且不只會在某段時空賜給某些人,而是給世界的.

耶穌基督把我們的焦點重新調教,人很容易只睇重自己的需要,信仰追求的目的也只期望滿足自己,因此所追求的神蹟就是要看見能力(招數/把戲),為滿足自己的好奇與慾望…

到底V26耶穌所指的”神蹟”到底是什麼呢?

V34他們說:「主啊,常將這糧賜給我們!」 35耶穌說:「我就是生命的糧.到我這裡來的,必定不餓；信我的,永遠不渴.

吃飽仍然是我們作為人所需要的,不論是肉體上,還是心靈上.

耶穌基督就是生命的糧,上帝最大的神蹟!

真正我們應該要追求的神蹟,就是追求耶穌基督；

回到最重要及基本,就是我們要先祈求賜糧的主,更要努力,為要透過主的成為而看見主!

[1] 資料來自《全球華文網》\url{https://biweekly.huayuworld.org/material/92/video}



\newpage

\section{我就是生命的糧}
\label{sec:562}
\textbf{李均熊博士 - 我就是生命的糧}
\newline
\newline
\hyperref[sec:561]{< < < PREV SERMON < < <}
~
\hyperref[sec:toc]{[返目錄]}
~
\hyperref[sec:563]{> > > NEXT SERMON > > >}


講題:我就是生命的糧 　I am the bread of life

經文:約翰福音6:35,41-51

講員:李均熊博士

場合:香港中文大學崇基學院禮拜堂主日崇拜

日期:2021年8月8日

起:糾結的約翰福音

約翰福音,是一卷十分難讀通的福音書.約翰福音難讀通,主要原因是福音書的每個段落,經文都糾結多重的主題,好像一條粗壯鋼纜(cable),是由無數鋼索 (wire) 組成一樣.好像我們今日經課福音書的選段,若你仔細分辨一下內中所包含的主題,你會驚覺經文的內容是如此豐富,而且每個主題並非單一出現,而都是貫串整卷約翰福音書,要掌握約翰福音的神學信息實在不是一件易事.

看看今日經課福音書選段的主題,當中就包括有:耶穌「我就是」的宣稱,生命的糧(原文即生命麵包),引申出來的不飢不渴,吃了就不死；接著是一連串關於舊約的主題,如嗎哪,曠野,猶太人的祖先,摩西,先知書上的教導；當中還有猶太人對耶穌的議論,天上降下來相對耶穌地上的父母,差遣耶穌的父,未看見父和看見父,還有末日復活,永遠活著等等.

若你熟識約翰福音,剛列出的每一個主題都可以貫串整卷約翰福音的.耶穌「我就是」的宣稱在約翰福音出現過至少七次；與猶太人爭論的主題在約翰福音五至十章均有出現；摩西和律法的主題在第一章引言部份和第三章耶穌與尼哥德慕夜話都有談及,也散見於五,七,八,九章；這裡不飢不渴的主題,耶穌在第四章與撒馬利亞婦人談道就說過,那處是活水,這裡則是生命麵包；信耶穌是父差來的,就由第一章開始,到四,五,六,七,九,十二,以至十四至十七章耶穌的最後道別都有提及；復活主題不單在五和六章提及,更有十一章整段關於拉撒路復活的事件和對話,耶穌自己就是從死人中復活的主．．．．．．若每一個主題都仔細追溯,講解明白,就是純粹環繞今日經課的主題,我想到了晚上大家還是要坐在這裡.

承:信的人有永生

有見及此,今日我只集中6章47節「信的人有永生」這個主題去宣道.「永生」(eternal life) 在約翰福音是一個十分重要的主題.但大家首先要留意,要完全掌握「永生」這個主題,可能要同時理解第三章耶穌與尼哥德慕的對話,當中涉及第六章沒有,但又和「永生」有關的主題:神的國,重生,得救,定罪,這是今日我會從略的.但選定「永生」這個主題,是因為我懷疑華人信徒受佛道思想,又或中國文化所影響,可能都對約翰福音所講的永生有些誤解.在今日經課中,6章35及48節耶穌宣稱「我就是生命的糧」,此一宣告帶來猶太人的疑惑和議論,耶穌回應他們的時候,就說到「信的人有永生」,以至說人若吃這生命的糧,必永遠活著.正是這裡,雖然經文講到永遠不死,但永生並非指佛家所言的超脫輪迴,又或道家追求的長生不死,又或普通人想當然的指在時間上永遠的生命.若我們將永生的希臘文ζωὴν αἰώνιον(zōēn aiōnion讀zo-heyn eh-oh-nee-ohn)直譯,應是「那屬於αἰών (aiōn 讀eh-on) 的生命」(the life belongs to aiōn).大家不要聽錯,Aion (讀eh-on)並不是某間日本品牌的超級巿場或財務公司(Aeon),而是(根據黃根春教授編寫的新約希臘文詞典)可以解作「時代」或「年代」(the age).今日講「時代」可能有一些敏感,為免大家進一步誤會,我就跳過所有學者的討論,直接將結論和大家分享.

這一個中文聖經譯作「永生」的詞語,應該理解為「那屬於時代的生命」(the life belongs to the age).而這裡的「時代」,從猶太人的末世觀出發,所指的並非今天的時代,而是「那要來的時代」(the age to come).所以「永生」並不是永遠不死的生命.耶穌作為生命的糧,當我們接受並相信,就能得著這「那屬於那要來的時代的生命」(the life belongs to the ages to come).「那屬於那要來的時代的生命」,並不會因為今世的生命完結而完結,所以是「不死」(everlasting)的.但相反地,約翰福音的運用可以理解為:因為耶穌從天上來成就父的旨意,相信又接納耶穌是父差來的人,就可以在今世享有「那屬於那要來的時代的生命」.

大家還清醒嗎?不錯,若要將方便上口的「永生」一詞,改成為一串翹口的「那屬於那要來的時代的生命」,約翰福音將變得非常難(閱)讀.大家不妨之後找機會將聖經所有譯作「永生」地方,改成「那屬於那要來的時代的生命」,看看是容易明白了, 還是更加令人疑惑.但這裡就樣我(用有限時間)解開(unpack)這一個(有如數學代號mathematical notation)「永生」所代表的「那屬於那要來的時代的生命」的意思:

在耶穌當代的猶太人,他們有一套基督教所承接的末世觀,就是將世界分為兩個時代(two ages),一個是今世的時代(this age),另一個就是「那要來的時代」.「那要來的時代」,亦就是「彌賽亞的時代」the age of Messiah,又或說「神的國降臨的時代」the age of the coming of the Kingdom of God,有時亦會說成「末後的日子」the last age,也可說成「末世」eschaton.意思就是說:當彌賽亞／基督顯現的時候,一個新的時代,也是末後的時代,便要來臨,上主作王的國度要降臨,那時上主要審判人間(定世人的罪),有些人要因此滅亡,有些人卻因基督而得救,要和上主一同作王,享受這個時代的新生命,得著一個屬於這個「要來的時代的生命」(或曰永生).

你若將這一個末世的意義代入約翰福音3章16節,就會明白甚麼叫「不至滅亡,反得永生」.

以上只是一種很簡化的說法,而且在約翰福音,或之後我們基督宗教的演譯,都和這種兩個時代的末世觀有些不同.基督第一次降臨,開始了那要來的時代,但這個「今世的時代」卻未完結,要到基督的再臨,才完全實現那個將要來的彌賽亞時代.我想大家應該都聽過既濟未濟(already but not yet)的基督教末世觀,正正就是指著這個「將要來的時代」已然開始,卻仍未成全的意思.

轉:領悟「生命的糧」

對於猶太人以至約翰福音對末世的理解,如何幫助我們領悟今日經課「生命的糧」的意義?這裡「生命的糧」所聚焦的,是C. H. Dodd所謂的「實現式終末論」(realized eschatology),所關注的不是那將要來的彌賽亞時代怎樣將萬事更新(雖然這個將來的元素也是有的,約翰福音六章中就四次題到在末日使人復活),而是如何透過相信耶穌基督,使那將要來的世代的生命,在今世的日子中實現出來.「信的人就有永生」這句說話是應該如此拆解(unpack)的,「信的人就有永生」是指我們若相信耶穌是由父從天上差來完成父的旨意的基督,就可以在今世得著那將來世代才可以得著的生命.

領受耶穌生命的糧,究竟會得著一個怎樣的生命.這是一個屬天的生命,不受地上物質的限制而變得到不足:「這是從天上降下來的糧,使人吃了就不死.我就是從天上降下來生命的糧；人若吃這糧,必永遠活著.」這句說話是回應猶太人聽到耶穌說自己「是從天上降下來的糧」而產生的疑惑.猶太人因為認識耶穌的父母,就奇怪耶穌為甚麼說自己是從天降下來.情況一如耶穌和尼哥德慕對談,尼哥德慕也誤解了重生意思一樣:耶穌的意思是從天上生,尼哥德慕想的是要由回復時間的開始再生出來.再聽一次耶穌對猶太人的回答(6:44-50):「若不是差我來的父吸引人,就沒有人能到我這裡來；到我這裡來的,在末日我要使他復活.在先知書上寫著:『他們都要蒙神教導.』凡聽了父的教導而學習的,都到我這裡來.這不是說有人看見過父,惟獨從神來的,他才看見過父.」猶太人嘗試從時間的橫向歷史現實,從過去已知的經驗,從地上人間的關係去明白耶穌的本源和說話的意思,但耶穌所講的是縱向,超越歷史平凡的上面,需從上頭啟示的源頭去明白的.

猶太人從祖宗的過去,從自己的經驗出發,只知道神賜過嗎哪,使民族離開曠野的危險.但耶穌所說的是超越我們的過去的經驗,是天上的父教導和啟示源頭的自我揭示.用回我解釋的末世觀,可以說猶太人經驗和思想仍然停留在今世的地平線(Horizon)上,但耶穌要作為生命的糧,所賜的是屬於「那將要來的世代的生命」.雖然耶穌用大家熟識,物質的麵包去,去講這個生命是可以食進肚子,延續今世的經驗；但同一時間耶穌用天上的,末世的復活,「那屬於世代的」這些用語提醒猶太人,又用到自己的肉可以吃,來阻斷猶太人在地的想法,正是要從他們盼望的要來的世代去入手,要求他們去融合今世與他世的地平線(fusion of horizons),才可以領悟耶穌是生命的糧的意義.

合:今天的福音

對今日的我們,耶穌是生命的糧,以至我們吃這個糧就可以得到那要來的世代的生命,又有甚麼意義呢?若在下已經是基督徒,代表已經相信和接納耶穌是父所差來,完成父旨意的基督聖子,透過聖餐吃過耶穌的肉和飲過耶穌的血,是已經與基督結連,領受了這個天上的糧.換句話說,作為基督徒,我們已經有「永生」,已經擁有「那要來的世代的生命」.但我們會否仍然如猶太人一般,想法仍然停留在今世的生活,而沒有領悟「永生」意義呢?我們是否仍然以我們有限的經驗,以我們地上生活的理解,去試圖活出不屬於今世,那屬天的生命呢?當然我們不會好像上一段落(6:25-34)一般群眾只要求耶穌的神蹟去餵飽我們；我們當然不會當耶穌是一般地方神靈如黃大仙對我們有求必應吧(是嗎?),我們當然不會以為我們可以將耶穌好像手機般放在口袋,隨時可以應我們需要速遞我們的要求 (還是你是這樣看信耶穌?)；但我們遇事祈禱,是否仍然只期望上主以我們的方法(in our terms),以我們經驗所容許的方法,去處理我們的問題,卻沒有想到耶穌不是我們可以操控,信仰生命不屬於今世,上主參與我們生命的困難是透過祂的方法(in his terms),那屬於他世,天上的方法呢?我們所擁有的生命是史無前例的(unprecedented),所以永遠是意料之外的(unexpected),也因如此,也是不能操控的(uncontrollable).但正因如此,這「永生」才能打破歷史的常規,展現嶄新的可能.

信耶穌,得永生,並不是使我們將來能上天堂,也不是只消除我們對死亡的恐懼而永遠不死,更不是只滿足我們在地生命的需要；信耶穌,得永生,是否能帶給我們屬天的視野,或更正確地說,信耶穌,得永生,是否能使我們將地上的生命經歷,融合從天而來的生命可能呢?我們的祈望,渴求,有否只關注在地今世的生活,也同時關注那要來的世代如何實現,願上主的國度降臨,願上主的旨意行在地上如同行在天上呢?



\newpage

\section{人生的智慧}
\label{sec:563}
\textbf{林豪恩先生 - 人生的智慧}
\newline
\newline
\hyperref[sec:562]{< < < PREV SERMON < < <}
~
\hyperref[sec:toc]{[返目錄]}
~
\hyperref[sec:564]{> > > NEXT SERMON > > >}


箴言 9:1-6
\newline
\begin{longtable}{cl}
\hline
\hline
章節 & 經文 (和合本修訂版)\\
\hline
9:1 & \begin{tabularx}{0.7\textwidth}{X} 智慧建造房屋，鑿成七根柱子， \end{tabularx} \\ \\ \relax
9:2 & \begin{tabularx}{0.7\textwidth}{X} 宰殺牲畜，調好美酒，又擺設筵席， \end{tabularx} \\ \\ \relax
9:3 & \begin{tabularx}{0.7\textwidth}{X} 派遣女僕出去，自己在城中至高處呼喚： \end{tabularx} \\ \\ \relax
9:4 & \begin{tabularx}{0.7\textwidth}{X} 「誰是愚蒙的人，讓他轉到這裡來！」又對那無知的人說： \end{tabularx} \\ \\ \relax
9:5 & \begin{tabularx}{0.7\textwidth}{X} 「你們來，吃我的餅，喝我調的酒。 \end{tabularx} \\ \\ \relax
9:6 & \begin{tabularx}{0.7\textwidth}{X} 你們要離棄愚蒙，就得存活，並要走明智的道路。」 \end{tabularx} \\ \\
[1ex]
\hline
\hline
\end{longtable}
講題:人生的智慧 　The Wisdom of Life

講員:林豪恩先生

經文:箴言9:1-6

場合:香港中文大學崇基學院禮拜堂主日崇拜

日期:2021年8月15日

引言

《人生的智慧》是十九世紀德國哲學家叔本華晚年的著作,以其深邃的洞察力闡明了對人生諸多問題的見解,使其成為百年以來的哲學經典.這本書提出的「人生的智慧」是甚麼呢?我不在這裡介紹了,留待有興趣的朋友自行閱讀.不過,讓我們退一步問,智慧對於人生是否有用呢?叔本華的生平就是對這個問題的一項說明.

叔本華生於富商家庭,儘管出身顯赫,他的童年卻很不幸.父母性格不合,經常吵架.在這種家庭氛圍下長大的叔本華,從小就敏感而憂鬱.

17歲時,叔本華的父親意外身亡.後來,他和媽媽因瑣事吵了起來,媽媽一氣之下把他推下樓梯.叔本華帶著傷痛憤然離家,從此以後他再沒有踏進過家門.後來,叔本華愛上了一個演員,就在他鼓足勇氣向對方表白時,對方卻拒絕了他.那次事件傷害了年輕的叔本華,從此他再沒有談過戀愛.

父親早逝,與母親鬧翻,愛情又受挫,一連串的打擊,命運似乎要把叔本華逼到絕境.叔本華如何保護自己免於墮落呢?就是全心全意地投入尋找智慧的工作,就是他的哲學研究.由於其非凡的哲學才能,叔本華在30歲時就發表了處女作《世界是意志的表現》.他原以為,這本書可以讓自己立即名垂青史,沒想到剛出版就滯銷,原因很簡單,因為幾乎沒人能看懂.儘管出師不利,叔本華沒有思考.在63歲時,他發表了《附錄與補遺》(就是《人生的智慧》的來源).沒想到,這本書馬上就引起了轟動.叔本華終於等到了早該屬於他的榮耀.但是此時的他,早已看淡了名利,甚至連柏林皇家科學院院士的稱號,都被他婉拒.

家庭和感情都未如願,世人的反應更是不到自己控制,叔本華對智慧的追尋幫助他跨越悲劇的人生,讓他在孤獨中變得堅強.

智慧在基督徒生活中應該佔甚麼位置呢?曾經有不少人為此而爭論.有教導說:「只要研讀聖經就夠了,不需要再閱讀其它書籍!所謂的屬靈書籍,全都是人寫的,充斥著人的思想,不是來自於上帝的啟示.」在這種教導之下,有一位旅美基督徒學者如此疑惑:「我從小熱愛中國文學中的詩詞歌賦.不過,帶我信主的教授對我說,他成為基督徒以後,就再也不碰世俗的文學作品了,只讀與聖經相關的書.這讓我陷入了困惑:基督教和知識以及文學藝術的正確關係,究竟應當如何?」有人立即斷章取義地引用保羅在哥林多前書8章1節來支持反智的傾向:「知識使人自高自大,惟有愛心能造就人.」到底聖經是否反智的呢?關於智慧,聖經有甚麼教導呢?讓我們再閱讀今天經課的四段經文,思考聖經中智慧的痕跡.

一.箴言9:1-6:追尋智慧的呼籲

9:1智慧建造房屋,鑿成七根柱子,

9:2宰殺牲畜,調和旨酒,設擺筵席；

9:3打發使女出去,自己在城中至高處呼叫,

9:4說:誰是愚蒙人,可以轉到這裡來!又對那無知的人說:

9:5你們來,吃我的餅,喝我調和的酒.

9:6你們愚蒙人,要捨棄愚蒙,就得存活,並要走光明的道.

箴言第九章把智慧擬人化,像一位女主人,邀請人來赴她的宴席,即是邀請人來領受智慧,分享智慧,學習智慧.箴言是聖經智慧文學的代表,將上帝的「智慧」應用到實際生活事務中.如果你讀過箴言,你會發現書中收集了很多箴言和格言,用以教導年青人如何生活,把前人的「智慧」代代相傳,承繼不衰.生活是處境性的,智慧也要求人認識特殊的處境,制度,人物和其他受造物,並要有恰當的行動.有人認為箴言這類智慧文學所講的只是常識,並不是深奧的屬靈啟示,因此予以輕視.常識是不是與上帝無關於呢?是否只有常理以外的特殊介入才是與上帝有關的呢?對於硬要把常識與上帝二分對立的朋友,以下這個笑話可作為提醒:

在某個村落,下了一場大雨,洪水淹沒全村,一位牧師正在教堂裡祈禱,眼看洪水已經淹到膝蓋了.一個救生員駕著舢板前來,跟牧師說:「牧師,趕快上來吧!不然洪水會把你淹死的!」牧師說:「不!我深信上帝會來救我的,你先去救別人好了.」不久,洪水淹過牧師的胸口,牧師勉強站在祭壇上.這時,又一個消防員開著快艇過來,跟牧師說:「牧師,快上來,不然你真的會被淹死的!」牧師說:「不,我要守住教堂,上帝一定會來救我的.你還是先去救別人好了.」又過了一會兒,洪水已經把整個教堂淹沒了,牧師只好緊緊抓住教堂頂端的十字架.直升飛機緩緩飛過來,飛行員丟下了繩梯之後大叫:「牧師,快上來,這是最後的機會了,我們不願意見到你被淹死!!」牧師還是堅定的說:「不,上帝一定會來救我的.你還是先去救別人好了!」

最後,洪水滾滾而來,牧師被淹死了……他見到上帝後很生氣的質問:「主啊,我終生奉獻自己,誠心誠意的事奉您,為什麼你不肯救我!」上帝說:「我哪裡不肯救你?第一次,我派了舢板去救你,你拒絕了；第二次,我又派一隻快艇去,你還是不要；第三次,我以國賓的禮儀待你,派一架直升飛機來救你,結果你還是不願意.我以為你急著想要回到我身邊來呢!」

許立中先生有一本書名為《雞毛蒜皮的信仰》, 指出基督徒不能倖免於凡人的種種,當信仰落實在人生的真實場景時,有時要回歸常識.本書的英文書名就叫Common Sense Christianity.有句話說:common sense is not common.如果有人認為聖經中的智慧文學只是common sense而輕視它,不妨自我檢視一下自己是否有足夠的common sense.若有欠缺,不妨聽聽箴言作者的呼籲:「誰是愚蒙人,可以轉到這裡來!」「你們來,吃我的餅,喝我調和的酒.」

二.詩篇34:9-14:敬畏耶和華是智慧的開端

34:9耶和華的聖民哪,你們當敬畏他,因敬畏他的一無所缺.

34:10少壯獅子還缺食忍餓,但尋求耶和華的甚麼好處都不缺.

詩人呼籲敬畏耶和華,這與智慧有關係嗎?有的,關係就在剛才讀過的箴言:

9:10敬畏耶和華是智慧的開端；認識至聖者便是聰明.

開端,也可以譯為「首要」,原來,在聖經的教導中,智慧的核心是「敬畏耶和華」.詩人呼籲敬畏耶和華,與箴言教導的生活智慧的關係是怎樣的呢?可以說,這是一個硬幣的兩面.箴言從現實生活方面教導,離開愚蒙,走向智慧,就得存活,而且活得光明；詩人從宗教信仰方面教導,敬畏耶和華,那些人一無所缺.如何敬畏耶和華呢?詩人接著所說的又回到箴言式的生活智慧:

34:11眾弟子啊,你們當來聽我的話!我要將敬畏耶和華的道教訓你們.

34:12有何人喜好存活,愛慕長壽,得享美福,

34:13就要禁止舌頭不出惡言,嘴唇不說詭詐的話.

34:14要離惡行善,尋求和睦,一心追趕.」

大家都可以看到,聖經是把生活倫理上的「敬畏」和和宗教禮儀上的「敬畏」結合.智慧並不是指得著頭腦的知識,而是活出倫理的觀念.或者有人認為,這些存留下來的智慧都是傳統,傳統好像給人古老和落後的感覺.聖經有用「道路」這意象來描述傳統.「道路」是甚麼呢?道路是能夠有效地把我們從甲地帶到去到乙地去,不用後來的人重新開闢一條新的路徑.在宗教上的說法,就是沿著合乎上帝規範的道路活於世上,這就是智慧.

三.以弗所書5:15-20:聖靈是活出智慧的幫助

5:15你們要謹慎行事,不要像愚昧人,當像智慧人.

5:16要愛惜光陰,因為現今的世代邪惡.

5:17不要作糊塗人,要明白主的旨意如何.

5:18不要醉酒,酒能使人放蕩；乃要被聖靈充滿.

5:19當用詩章,頌詞,靈歌彼此對說,口唱心和地讚美主.

5:20凡事要奉我們主耶穌基督的名常常感謝父上帝.

以弗所書這段選讀呼應著箴言第九章的呼籲,拒絕愚昧,追求智慧；在具體的行為表現上,則呼應著詩篇34篇的教導,都提及「口」.詩篇34篇從負面的角度提醒:「禁止舌頭不出惡言,嘴唇不說詭詐的話」(34:13).如果嘴唇不是用來說惡言和詭詐的話,用來做甚麼呢?以弗所書從正面的角度提議:「當用詩章,頌詞,靈歌彼此對說,口唱心和地讚美主.」(5:19)

關於說話的智慧,坊間著作汗牛充棟,因此有人認為聖經的教導也沒有甚麼特別之處,「係人都知啦」.智慧這回事,與「知道不知道」不盡相同.智慧不單單是知道,智慧是知道了之後具體行為表現是否合適.坊問有一個笑話,希望是虛構,如有雷同,實屬不幸:

有一個父親,有一天聽到兒子在家講粗口,感到非常憤怒.他知道講粗口是不好的,於是盡父親的責任教訓兒子,但是,他用來教訓兒子的說話也是粗口,即是用粗口來教訓兒子不要講粗口.這些父母教導子女的矛盾情況,實在值得我們作為父母的留意,或者我們沒有用粗口教訓兒子不要講粗口,但我們可能會以發怒控制兒女的怒氣,以恐嚇回應兒女的擔心和害怕等.我們不是沒有知識,只是知道不等於做到,知識不等於智慧.如果我們真的是這樣,也不需要對自己太過失望,因為連使徒保羅也說:「立志為善由得我,只是行出來由不得我.我所願意的善,我不去做；我所不願意的惡,我反而去做.」(羅馬書7:18-19)

那麼,有甚麼可以幫助我們呢?我們有可能培養並活出智慧嗎?如何做到以弗所書5:15的勸勉「不要像愚昧人,當像智慧人」呢?使徒保羅在這裡引介能夠幫助我們的,不知道我們有沒有想過,智慧竟然是與祂有關的,保羅吩咐我們:「要被聖靈充滿」.(5:18)經文描述,被酒充滿的人會變得放蕩；我們也體會過,被怒氣充滿會罵人,被害怕充滿會尋求操控.似乎,我們都逃不過被某些有形無形的東西充滿,既然如此,使徒保羅指出,最好的就是被聖靈充滿,因為被聖靈充滿,才能夠善用自己的嘴巴.為甚麼要談嘴巴呢?其實嘴巴只不過是一個代表,雅各書就是以舌頭來代表整個人.

坊間談智慧,可能很多內容和聖經的類似,甚至相通,這是沒有問題的,因為大家都是人,大家都在思考和累積經驗如何在同一個天空之下生活,只是聖經鼓勵及帶引我們進入自己深處,對自己的實況有更深刻的體會,以至產生類似保羅的自我認識:「立志為善由得我,只是行出來由不得我.我所願意的善,我不去做；我所不願意的惡,我反而去做.」(羅馬書7:18-19)這種自我認識,刺激我們仰望蒼天,開啟我們尋找上主的道路,對我們來說,也軟化自己的心,接受聖靈的同在和啟迪.因此,聖經所講的智慧,是與聖靈相關的,也是被聖靈所成全的.

至於保羅在哥林多前書所說好像反智的說話,「我要摧毀智慧人的智慧,廢棄聰明人的聰明.」(1:19)與箴言等智慧文學的教導沒有衝突,因為箴言所教導的與保羅在哥林多前書所講的是不同範疇的事情.那麼,保羅所講的是甚麼範疇呢?讓我們讀今天的福音經課:

四.約翰福音6:51-58:看見將要來臨的世界的智慧

6:51我是從天上降下來生命的糧；人若吃這糧,就必永遠活著.我所要賜的糧就是我的肉,為世人之生命所賜的.」

6:52因此,猶太人彼此爭論說:「這個人怎能把他的肉給我們吃呢?」

6:53耶穌說:「我實實在在地告訴你們,你們若不吃人子的肉,不喝人子的血,就沒有生命在你們裡面.

6:54吃我肉,喝我血的人就有永生,在末日我要叫他復活.

6:55我的肉真是可吃的,我的血真是可喝的.

6:56吃我肉,喝我血的人常在我裡面,我也常在他裡面.

6:57永活的父怎樣差我來,我又因父活著；照樣,吃我肉的人也要因我活著.

6:58這就是從天上降下來的糧.吃這糧的人就永遠活著,不像你們的祖宗吃過嗎哪還是死了.」

約翰福音裡面耶穌是生命的糧這個主題,在過去幾個主日都一直都提到,如果大家仍記得,那是耶穌在五餅二魚事件之後回應追隨的群眾而講的.五餅二魚事件的主題是甚麼呢?是食物與生存,食物是生存的必需條件,沒有食物,就不能生存下去.耶穌沒有否認這事實,只是在這基礎之上,啟迪群眾思考生存是怎樣的呢?甚麼是生命呢?因此談到永生,上主日李均熊博士已經詳細解釋甚麼是永生.回應五餅二魚的群眾,耶穌告訴他們,你們知道吃餅和魚能夠令你們在這個世代生存,但你們知道在將要來的世代生存要靠甚麼嗎?就是耶穌自己,因為耶穌是道路,真理,生命,若不藉著祂,沒有人能到父神那裡去.

你知道群眾聽完耶穌關於我是生命的糧的講道之後有甚麼反應嗎?相信是「趕客」的效果,按約翰福音6:66記述:「從此,他門徒中有很多退卻了,不再和他同行.」保羅在哥林多前書所談及的,就是對耶穌基督十字架福音報以類似回應的人.其實,「要摧毀智慧人的智慧,廢棄聰明人的聰明」目的是清除認識耶穌的障礙,讓人們能夠進入耶穌所說的永生.

五.總結

從經課中的箴言和詩篇,我們看見智慧是和創造秩序有關的.宇宙秩序涉及受造萬物物理上和倫理上的功能和關係,也是關於創造主所設定來維繫著這些東西的法則和規範.各位弟兄姊妹,或者我們不應該期望太多違反常理的事情,雖然上帝有能力這樣做,但何必逼上帝違反自己所創造的秩序呢?現實上來看,上帝尊重宇宙秩序是常態,打破宇宙秩序的事件只是例外而已.尼布爾著名的祈禱這樣說:「親愛的上帝,請賜給我雅量從容的接受不可改變的,賜給我勇氣去改變應該改變的,並賜給我智慧去分辨什麼是可以改變的,什麼是不可以改變的.」

然而,聖經所啟示的智慧,卻又超越常理和常規,超越被造的世界.因為聖經所啟示的,除了這個被造的世界,還有創造主本身,以及被更新的世界,或者稱為在來臨中的天國.智慧教導我們如何在這個世代生活,智慧也開啟我們眺望來臨中的世代的視野.在這個世代生活使我們有使命,眺望來臨中的世代使我們不會因為這個世界的狀況而感到絕望,使我們無論處在甚麼狀況仍然可以懷著盼望.尼布爾著名的祈禱下半段如此說:

生活在每一天,享受著每一刻；

接受生活的困苦,作為通往安寧的道路.

像耶穌那樣,接受這罪惡的世界,

如其所是,非我所願；

深信只要我順服你的旨意,你將匡正萬有,

因而我今生欣悅有度,來世與你永享至福.

阿們.



\newpage

\section{你認識祂嗎}
\label{sec:564}
\textbf{張美珍博士 - 你認識「祂」嗎?}
\newline
\newline
\hyperref[sec:563]{< < < PREV SERMON < < <}
~
\hyperref[sec:toc]{[返目錄]}
~
\hyperref[sec:565]{> > > NEXT SERMON > > >}


講題:你認識「祂」嗎? 　Do you know “Him”?

講員:張美珍博士

場合:香港中文大學崇基學院禮拜堂主日崇拜

日期:2021年8月22日

大家認識我嗎?主席剛介紹我,以及在座當中有許多崇基學院的校友或老師,當然會認識我.但假若我問,除了我的名字及工作,或者宗教信仰,你認識我有幾多?相信大部份都答不出來.

大家望一眼坐在你兩旁的弟兄姊妹,你認識他/她嗎?相信除了你的伴侶,親人或好朋友,大部份都並不認識的,因為我們學院的禮拜堂並非獨立教會,除了主日崇拜之外,並沒有很多的團契或聚會,因此大家互不認識並不出奇.不過,縱使坐在你旁邊的,是你的親人或好朋友,你真的認識他/她嗎?

社交媒體面世只不過是10多年的光景,改變了全世界的溝通模式.不論Facebook或IG的用家,相信必定也有好多朋友.但是,你認識他們嗎?我現時用的頭像是一只貓,原因是假若我用自己的相片當頭像,差不多每天都有英俊瀟洒的男士要加我為朋友,我發現其中甚至有一位男士,曾經使用不到的身份及名字.網上假身份多的是,騙案(情騙財騙)時有發生!就算你以為你真的認識那人而加進朋友圈,但我聽過不只一次,有朋友正正因為在社交媒體的說話,惹來那些所謂朋友的妒忌或報復,幾乎失去工作.我不禁會問,你以為真的認識你的「朋友」嗎?

在大約4年前,我一些失去聯絡超過40年的小學同學,竟然透過社交媒體而相認.我們畢業班大概30人,陸陸續續聯絡上超過20人!大家非常興奮,組織了數次飯局,並且一起往澳門遊玩.直至前年,有同學仔在群組發放了一些與社會事件有關的言論,惹來激烈反應,唇槍舌劍,最終不歡而散.現時群組除了互祝生日之外,再沒有什麼東西說了.兩年前的社會事件所觸發的政治風波,導致很多同學,朋友甚至家庭之間出現分歧,最後演變成分裂.突然發現,過往以為自己很認識的他／她,竟然會變得很陌生.原來自己所認識的,只是某人的某一面,或是停留在某一時空下的他／她.

我因工作關係經常要往中環開會,偶爾會聽到一些在商界發生的趣事.話說有次談到某富商過身,但他的訃聞在舉殯後才刊登的.聽說那是因為避免尷尬,因為如果一早把舉辦的時間地點刊登出來,不排除會出現小三拖著孩子到靈堂哭跪,之後便是要求分身家.這些小三,可能作為死者的未亡人根本毫不知情.結婚數十載,你能夠肯定你對你的枕邊人瞭如指掌嗎?大家可能覺得這只是笑話,但我的確遇過一件真人真事:

很多年前我同一位大學同學在中國大陸旅行時,認識了一對夫婦,是來自加拿大多倫多的.後來我和外子在美國讀書時,曾經探訪他們並在他們家中小住,兩夫婦非常恩愛,住在一間獨立屋,是我們香港人極之嚮往的神仙一般的生活.誰知兩年後(那時我已前往加拿大卡加利大學讀博士課程),有一晚深夜時分收到她的電話,在另一端哭不成聲；原因是在早一天晚上,她丈夫半夜開了她的私家車,留下了一封信,與他的情人私奔,離開這個家,原來他們已經有了一個11歲的孩子.這些好像只會發生在電影情節的,從沒想過原來真的可以在現實生活中出現.她從沒想過結婚20多年的枕邊人,會突然變了一個陌生人.

言歸正傳,今日的講題,是「你認識祂嗎」?即是你對上帝的認識有多少?

去年當校牧高牧師邀請我講道時,我是有些戰兢的,因為若要用主日的經課(尤其是有四段經文)來準備講章,是頗費勁的.聖經的寫成牽涉很多時代背景,若想透過聖經去認識上帝,必須了解聖經當年成書的時代背景,並非單憑聖靈感動便能夠明白上帝的話,否則對上帝的認識便會很片面.

今日所讀的4段經文,抱歉我並沒有太多時間去研讀背景.但坦白說,這4段經文所表達作者對上帝的認識,與我今時今日對上帝的認識是有點差距.舊約經課所展現的耶和華上帝,是一位唯我獨尊的上帝--順我者生,逆我者死；祂亦是一位善惡分明--善有善報,惡有惡報--的上帝.

約書亞記廿四章:

18耶和華又把住此地的亞摩利人都從我們面前趕出去.所以,我們必事奉耶和華,因為他是我們的神.

19約書亞對百姓說:你們不能事奉耶和華；因為他是聖潔的神,是忌邪的神,必不赦免你們的過犯罪惡.

20你們若離棄耶和華去事奉外邦神,耶和華在降福之後,必轉而降禍與你們,把你們滅絕.

詩篇三十四篇:

15耶和華的眼目看顧義人；他的耳朵聽他們的呼求.

16耶和華向行惡的人變臉,要從世上除滅他們的名號.

至於新約那兩段經課,也表達了一個善惡的對立面；以弗所書的經文甚至要求信徒穿戴上帝所賜的全副軍裝,去抵擋魔鬼的詭計.不知是否年紀大心血少,對於打打殺殺的戰爭電影已經不敢再觀看,甚至新聞報道近日阿富汗局勢,在電視上看見武裝份子的模樣,我都馬上轉台..所以讀到這個全副軍裝的經文,都有少許抗拒.不過,這段經文是描述心靈的爭戰,並非真的叫大家裝扮成軍隊去打仗.

而約翰福音第六章的經課,若單從字面看,則更容易令人受驚:

56吃我肉,喝我血的人常在我裡面,我也常在他裡面.

63叫人活著的乃是靈,肉體是無益的.我對你們所說的話就是靈,就是生命.

單從字面來看,會教人摸不著頭腦,難怪當福音剛傳入中國時,會被人誤會以為基督教是一個茹毛飲血的宗教.

聖經並非一本歷史書.它所記載的,是先賢聖哲在經歷時代的變遷,人生種種考驗時與上帝的關係,來表達他們對上帝的認識.舊約反映了很多古近東及希伯來文化(原文是希伯來文),而新約則受希臘文化影響(原文是希臘文).順帶一提.黃根春教授開發的＜原文新約＞app, 就是透過原文希臘文來研讀新約聖經,希望幫助大家能夠更深入認識上帝.換句話說,不同的作者在不同的年代受著不同的文化影響,及透過其自身的經歷,對上帝都有不同的認識.

舊約所展現上帝的形象,普遍來說都是高高在上,充滿威嚴的上帝；但當進入新約時代,四福音書所展現的,是「道成肉身」,上帝化身為人,以耶穌(上帝兒子)的形象進入人群,是一位平民並滿有慈愛,到處醫病趕鬼.所以很多釋經書都說舊約著重上帝的公義及懲罰,新約著重上帝的慈愛及拯救.

而你所認識的上帝,又是怎樣的一位上帝呢?

學院的每月都有行政職員例會,開始前都有書籍分享.有次臨時召開會議,來不及安排同事分享書籍,所以我便充當了這個角色.我告訴同事說,我所分享的這一本書籍,就是由我就讀中學以來至今,對我人生影響最深的一本書--聖經.

我分享了若干經文,其中兩段是:「你們聽見有話說:『以眼還眼,以牙還牙.』只是我告訴你們,不要與惡人作對.有人打你的右臉,連左臉也轉過來由他打…」(太五:38-39)；「只是我告訴你們,要愛你們的仇敵,為那逼迫你們的禱告… 」(太五:44)我之所以分享這兩段經文,是有感於前年的社會事件造成不少家庭,工作場所,朋友圈甚至教會等等有很多撕裂和創傷,不同立場的人對對方往往充滿敵意甚至仇恨.當時在職員會上也曾引起一些討論,甚至有基督徒同事提出疑問:若是這樣,公義何在?

很明顯,他所認識的上帝,與我所認識的上帝,是有點差距.

我在學院的工作,其中一個職責是紀律委員會秘書,每次處分學生都是由我來執行.有時同學犯事後被委員會約見時,都有很多道理,說要伸張正義.就在上月,我們剛處分了一個一年級升二年級的同學.他由於要上訴而來我房間會面,一坐下來便有很多道理,又搜集了很多資料,要與我辯論.

他說他是刻意違反這些規規的,質疑有些校規根本不應該存在.很明顯,這些都是從他個人觀點出發的.我沒有跟他辯論,因為我看事物的角度與他看事物的角度根本不一樣,不可能會說服他.我只是對他說:「同學,你進入大學只是一年光景,去年因為疫情的緣故,你根本未有機會回來校園上課,你連這個校園都未看清楚,又憑什麼去挑戰學校所定下的規則呢?還有,你今年17歲,你對這個世界認識有多深?請記著你進入大學的目的是為了接受教育,不是來教育我們的!」

我在中文大學工作了快將13年,由教書到行政,接觸過不少學生,但經歷了佔中及社會事件之後,我覺得自己越來越不認識這一代的學生.前年11月,就在中文大學被黑衣人佔領之前一星期,我教一年級新生的書院通識課,就在這禮拜堂上課.該課堂的題目是「法,理,情的道德判斷」,其中一個環節是討論當時很敏感的話題--「違法達義」.我叮囑同學不要做任何違法的事；法律是每一個人的底線.在課堂完結之前的發問時間,有一位一年級的新同學出來發言,不同意我的觀點,在六百多位同學面前,把我痛罵了一頓,說我不但自己沒有參與,甚至不鼓勵同學去爭取社會公義..

最後,我如何回答這位同學呢?

我告訴這位新同學,崇基的校訓是「止於至善」.至於如何定義「善」,聖經這句話說得很清楚:

「世人哪!耶和華已指示你何為善.祂向你所要的是什麼呢?只要你行公義,好憐憫,存謙卑的心,與你的上帝同行.」(彌迦書六章八節)

行公義,若沒有憐憫心腸,這個公義只是滿足一己的慾望,沒有顧及別人的感受及考慮別人的處境；存謙卑的心,意思即是,我們還有很多不足之處,不要以為自己掌握一切真理或道理,仍需多聆聽及多認識；最後,就是與你的上帝同行--這句說話其實有些危險的,因為若你對上帝的認識是非常片面,甚至以偏概全,把個人的慾望投射至等同上帝的意思,很容易會做出傷天害理的事；正如從古到今,有許多極端宗教的恐怖活動,都是以上帝之名來進行的.

今日這個講道題目,並非要給大家一個絕對的答案來認識上帝；事實上,年紀越大,便越發覺得自己對上帝的認識只是很皮毛.今日主要是拋磚引玉,盼望弟兄姊妹透過聖經及自身的經歷,存謙卑的心以多角度來認識上帝.

至於我所認識的上帝,又是怎樣的呢?於我來說,上帝就是愛.我每次回來崇基禮拜堂守禮拜,最感動的就是望著這個十字架.它展現了基督精神的最高境界,就是無條件的愛--主耶穌為我們死在十字架上,並非為了交換什麼,而是單純的為了愛.正因這緣故,才招聚了我們每一個人回來這裡敬拜上帝.

我亦相信愛是無堅不摧的.「愛你們的仇敵,為那逼迫你們的禱告」,的確是很難做到；但消滅敵人的最佳方法,就是把他變成你的朋友.唯有愛,才能讓人類社會團結起來向前邁進；也唯有無條件的愛,才能讓別人從我們身上見到耶穌,從而認識上帝.阿們.



\newpage

\section{不要讓傳統死去}
\label{sec:565}
\textbf{關瑞文教授 - 不要讓傳統死去}
\newline
\newline
\hyperref[sec:564]{< < < PREV SERMON < < <}
~
\hyperref[sec:toc]{[返目錄]}
~
\hyperref[sec:566]{> > > NEXT SERMON > > >}


馬可福音 7:1-8,14-15,21-23
\newline
\begin{longtable}{cl}
\hline
\hline
章節 & 經文 (和合本修訂版)\\
\hline
7:1 & \begin{tabularx}{0.7\textwidth}{X} 有法利賽人和幾個從耶路撒冷來的文士聚集到耶穌那裡。 \end{tabularx} \\ \\ \relax
7:2 & \begin{tabularx}{0.7\textwidth}{X} 他們曾看見他的門徒中有人用不潔淨的手，就是沒有洗的手吃飯。 \end{tabularx} \\ \\ \relax
7:3 & \begin{tabularx}{0.7\textwidth}{X} 法利賽人和所有的猶太人都拘守古人的傳統，若不按規矩洗手就不吃飯； \end{tabularx} \\ \\ \relax
7:4 & \begin{tabularx}{0.7\textwidth}{X} 從市場來，若不洗淨也不吃飯；他們還拘守好些別的規矩，如洗杯、罐、銅器、床鋪等。 \end{tabularx} \\ \\ \relax
7:5 & \begin{tabularx}{0.7\textwidth}{X} 法利賽人和文士問他說：「你的門徒為甚麼不照古人的傳統，竟然用不潔淨的手吃飯呢？」 \end{tabularx} \\ \\ \relax
7:6 & \begin{tabularx}{0.7\textwidth}{X} 耶穌對他們說：「以賽亞指著你們假冒為善的人所預言的說得好。如經上所記：『這百姓用嘴唇尊敬我，他們的心卻遠離我。 \end{tabularx} \\ \\ \relax
7:7 & \begin{tabularx}{0.7\textwidth}{X} 他們把人的規條當作教義教導人；他們拜我也是枉然。』 \end{tabularx} \\ \\ \relax
7:8 & \begin{tabularx}{0.7\textwidth}{X} 你們是離棄神的誡命，拘守人的傳統。」 \end{tabularx} \\ \\ \relax
7:9 & \begin{tabularx}{0.7\textwidth}{X} 耶穌又說：「你們誠然是廢棄神的誡命，為要守自己的傳統。 \end{tabularx} \\ \\ \relax
7:10 & \begin{tabularx}{0.7\textwidth}{X} 摩西說：『當孝敬父母』；又說：『咒罵父母的，必須處死。』 \end{tabularx} \\ \\ \relax
7:11 & \begin{tabularx}{0.7\textwidth}{X} 你們倒說：『人若對父母說：我所當供奉你的已經作了各耳板』（各耳板就是奉獻的意思）， \end{tabularx} \\ \\ \relax
7:12 & \begin{tabularx}{0.7\textwidth}{X} 你們就容許他不必再奉養父母。 \end{tabularx} \\ \\ \relax
7:13 & \begin{tabularx}{0.7\textwidth}{X} 這就是你們藉著繼承傳統，廢了神的話。你們還做許多這樣的事。」 \end{tabularx} \\ \\ \relax
7:14 & \begin{tabularx}{0.7\textwidth}{X} 耶穌又叫眾人來，對他們說：「你們都要聽我的話，也要明白。 \end{tabularx} \\ \\ \relax
7:15 & \begin{tabularx}{0.7\textwidth}{X} 從外面進去的不能玷污人，惟有從裡面出來的才玷污人。」 \end{tabularx} \\ \\ \relax
7:16 & \begin{tabularx}{0.7\textwidth}{X} 有耳可聽的，就應當聽！ \end{tabularx} \\ \\ \relax
7:17 & \begin{tabularx}{0.7\textwidth}{X} 耶穌離開眾人，進了屋子，門徒就問他這比喻的意思。 \end{tabularx} \\ \\ \relax
7:18 & \begin{tabularx}{0.7\textwidth}{X} 耶穌對他們說：「你們也是這樣不明白嗎？難道你們不了解，凡從外面進去的不能玷污人嗎？ \end{tabularx} \\ \\ \relax
7:19 & \begin{tabularx}{0.7\textwidth}{X} 因為不是進入他的心，而是進入他的肚子，又排入廁所。」（這是說，各樣的食物都是潔淨的。） \end{tabularx} \\ \\ \relax
7:20 & \begin{tabularx}{0.7\textwidth}{X} 耶穌又說：「從人裡面出來的，那才玷污人； \end{tabularx} \\ \\ \relax
7:21 & \begin{tabularx}{0.7\textwidth}{X} 因為從人心裡發出種種惡念，如淫亂、偷盜、兇殺、 \end{tabularx} \\ \\ \relax
7:22 & \begin{tabularx}{0.7\textwidth}{X} 姦淫、貪婪、邪惡、詭詐、淫蕩、嫉妒、毀謗、驕傲、狂妄。 \end{tabularx} \\ \\ \relax
7:23 & \begin{tabularx}{0.7\textwidth}{X} 這一切的惡都是從裡面出來，且能玷污人。」 \end{tabularx} \\ \\
[1ex]
\hline
\hline
\end{longtable}
講題:不要讓傳統死去 　Tradition vs. Traditionalism

講員:關瑞文教授

經文:馬可福音7:1-8,14-15,21-23

場合:香港中文大學崇基學院禮拜堂主日崇拜

日期:2021年8月29日

今天是施行聖水禮,很感恩,我有機會見證 15 個信徒加入上主大家庭的一刻.首先,要恭喜大家!一生一世跟隨主,是人生中一個極其重要的抉擇. 從今日起, 各位就是決志以自己的生命,來見證這活生生的信仰.有 15 個生力軍加入這個見證團隊,很難得!

眼看著聖水禮的過程,讓我想起一件與水有關的往事.

十多年前,我因為漸覺自己年紀已老,身體開始出了點毛病,常常腰酸背痛, 因此立下心志,決定天天早上游水.但當時天氣很冷,我就選擇去沙田源禾路的暖水泳池操練. 就是這樣,我早上 5:30 起來,大約 6:30 到源禾路游泳池,雄心壯志去做運動.到泳池了,發現在泳池裡大部份都是老人家, 他們居然都精力充沛,往返地沿著 50 米的直池暢泳.真厲害!而我呢?因為已經十年不沾水,第一天下水禮,就決定不游直池,只游 25 米的橫池.心想,游橫池,比較短,容易應付,以免出師未捷身先死,無謂第一天就溺斃.多麼合情合理吧.誰不知,我只橫游了一半, 便聽到在水面上有幾位婆婆和伯伯,大聲咒罵:『 邊個後生仔咁唔識規矩,阻實曬,竟然打橫游!我地ge 規矩,我地ge 傳統,係打直游啊!』 當時我很尷尬,只好潛落水底,拐個彎,潛去直池的起點,只願當我的頭伸出水面時,他們認不出這個後生仔原來就是我!

我在水底時,這樣想著:「天啊!怎麼連游水都有傳統,都有傳統定下來的規矩,要人人遵守?我可以反對這個傳統嗎?」

我是研究神學的.不少人都有一個錯覺,就是在大學裡搞神學的人,都是所謂自由派,只顧走在學術的前沿,不顧傳統的.這肯定是錯覺.搞神學的人,其實就是那些看重信仰傳統的人.正是因為視信仰傳統如珠如寶,就立心要活化信仰傳統,把信仰傳統與當下處境連結,讓二千年教會傳統能活生生的向當下的人說話,讓傳統接地.

一位已故的耶魯大學的基督教學者 Jaroslav Pelikan 有句響噹噹的名言:「Tradition is the living faith of the dead, traditionalism is the dead faith of the living. And, I suppose I should add, it is traditionalism that gives tradition such a bad name.」

「傳統是已死去的前人的活潑信仰,傳統主義則是當下活著的人的僵死了的信仰.我想我應該補充一點,[傳統本身是美好的],只是傳統主義給傳統帶來了如此壞的名聲.」這句話,也是崇基學院已故前院長沈宣仁教授經常掛在口邊的一句話:「傳統是已死去的前人的活潑信仰,傳統主義則是當下活著的人的僵死了的信仰.」

我們作為信仰的見證人,其中一個任務,就是要讓人看見,我們所安身立命的信仰,是由美好的傳統所構成的,而我們卻不是傳統主義者.

讓我們今天從馬可福音第七章去看看,耶穌是如何以「傳統」來擊退「傳統主義」的；也去看看,「傳統」與「傳統主義」的微妙分別.

可 7:1-8, 14-15, 21-23

1 有法利賽人和幾個文士從耶路撒冷來,到耶穌那裡聚集.2 他們曾看見他的門徒中有人用俗手,就是沒有洗的手,吃飯.(3 原來法利賽人和猶太人都拘守古人的遺傳,若不仔細洗手就不吃飯；4 從市上來,若不洗浴也不吃飯；還有好些別的規矩,他們歷代拘守, 就是洗杯,罐,銅器等物.)5 法利賽人和文士問他說:「你的門徒為甚麼不照古人的遺傳,用俗手吃飯呢?」6 耶穌說:「以賽亞指著你們假冒為善之人所說的預言是不錯的.如經上說:這百姓用嘴唇尊敬我,心卻遠離我.7 他們將人的吩咐當作道理教導人,所以拜我也是枉然.8 你們是離棄上帝的誡命,拘守人的遺傳」… 14 耶穌又叫眾人來,對他們說:「你們都要聽我的話,也要明白.15 從外面進去的不能污穢人,惟有從裡面出來的乃能污穢人.」… 21 因為從裡面,就是從人心裡,發出惡念,苟合,22 偷盜,凶殺,姦淫,貪婪,邪惡,詭詐,淫蕩,嫉妒,謗讟,驕傲,狂妄.23 這一切的惡都是從裡面出來,且能污穢人.」

讀這段經文,第一個問題要問的是,「究竟耶穌是否一個反傳統的人呢?」答案很清楚,當然不是.耶穌絕對不是一個背離傳統的人.這點從他如何斥責法利賽人和文士可以清楚看得出.第 8 至 9 節記載了耶穌如何斥責法利賽人和文士. 聖經和合本修訂版的翻譯比較準確,是這樣的:『你們是離棄上帝的誡命,拘守人的傳統.』耶穌又說:『你們誠然是廢棄上帝的誡命,為要守自己的傳統.』在第 6,7 兩節,耶穌更引用了猶太人看為寶貴傳統的一段話,即是以賽亞先知的說話,來譴責法利賽人及文士假冒為善.如果耶穌是藐視傳統的,那他這樣引用古人傳統的話來斥責對手,就顯得太自相矛盾了.

那麼,耶穌是怎樣看待傳統的呢?有不少釋經學者研究過這個議題.

原來這場有關進食前是否應該洗手的辯論,是關乎利未記和民數記裡有關禮儀上 的潔淨和不潔淨的條例.有關這些潔淨條例的傳統,在猶太人的圈子裡,歷來有不同的詮釋,眾說紛紜.簡單來說,一派說從外邊進入人體內的東西,如果沒有經過嚴格的潔淨儀式,是可以使人不潔淨的.相反,另外一派則說,入口的東西不能使人不潔,反而是從人裡面出來的東西,才有使大地和人不潔的可能性.這即是說,人才是污染的源頭.況且,若以最嚴謹的態度來理解聖經,根本在其中就找不到「進食」本身可使人不潔的傳統.

如此看來,當耶穌在 15,18-20 節重覆說「15 從外面進去的不能污穢人,惟有從裡面出來的乃能污穢人」時,顯明了兩點:1.耶穌沒有輕看有關潔淨律法的摩西傳統；2. 然而,耶穌對這些潔淨律法提出了更具說服力的解釋.因此,耶穌才會在 21-23 節進一步說清楚,「因為從裡面,就是從人心裡,發出惡念,苟合,偷盜,凶殺,姦淫,貪婪,邪惡,詭詐,淫蕩,嫉妒,謗讟,驕傲,狂妄.這一切的惡都是從裡面出來,且能污穢人.」

總結而言,耶穌不單只沒有輕視傳統,反而是以更嚴謹的力度來詮釋傳統,以致讓跟隨者更能活出上帝的誡命.

耶穌是嚴守傳統的人,卻不是傳統主義者-traditionalist.相反,法利賽人與文 士就肯定是傳統主義者的佼佼者了.我們可以從耶穌如何跟他們交手,明白什麼是傳統主義-traditionalism.

說法利賽人與文士所仗賴的是傳統主義而不是遵守傳統,起碼有四個互相交織的理由.

首先,第 7 節說,法利賽人與文士「將人的吩咐當作道理教導人」,並且,他們也把這個人的遺傳,說成是上帝的誡命,強逼每個人遵守.事實上,從歷史文獻,我 們可以知道,該洗手才可進食的潔淨禮儀,猶太教要到一世紀末之後才要求每個猶太 人遵守的.在耶穌時期,只有祭司才被要求遵守.法利賽人與文士的確是「將人的吩 咐當作道理教導人」.

那麼,為什麼他們要強加人的遺傳在每個人身上呢?這個答案構成了第二個理由證明他們是傳統主義者.為什麼要「強加」呢?因為他們要假借上帝之名,借人的傳 統來控制他人,透過上帝之名來譴責他人,貶低他人,尤其是要用之以排斥與自己立場不同的人.就是因為耶穌的名聲已蓋過他們的,而耶穌所持的立場與他們相異,法利賽人和文士就處處乘機借上帝之名去攻擊耶穌,排斥異己.借信仰傳統來攻擊異己,是假冒為善的傳統主義.在今天撕裂的社會裡面,這點尤為重要.在撕裂的社會裡,我們各自有自己的立場,有不同的爐堆,圍爐取暖.在爐火燒得熱騰騰之際,總是覺得上帝站在我們一方,而對方呢?對方就是妖魔化身.如果我們也是這樣,就要回憶林肯一句話:「My concern is not whether God is on our side; my greatest concern is to be on God’s side, for God is always right.」 「我在乎的並非上帝是否在我們這邊； 我最在意的,是我是否真的站在上帝那邊.因為 [只有] 上帝是對的.」事事以為自己所信奉的就是真理,以人言代替神言,以致把異己排斥,是傳統主義者的特徵.

第三,法利賽人和文士的傳統主義,其實是要以傳統來鞏固那不對等的權力關係.他們以傳統之名來責罵耶穌和門徒,就是要顯示,「我們法利賽人是高人一等, 你們是低端人口.我說了便是,我教導的你們要聽.」傳統主義的第三個特徵,就是把傳統轉為權力工具.

傳統主義的第四個特徵,可以從其後果看到.9-13 節是這樣的:

耶穌又說:「你們誠然是廢棄上帝的誡命,為要守自己的傳統.摩西說:『當孝敬父母』； 又說:『咒罵父母的,必須處死.』你們倒說:『人若對父母說:我所當供奉你的已經作了[奉獻]』 … 你們就容許他不必再奉養父母.這就是你們藉著繼承傳統,廢了上帝的話.你們還做許多這樣的事.」

這聽來荒謬得很.不過,傳統主義的荒謬之處,正正就是違反人性地墨守成規,把傳統僵化,使本來要讓人活得更好的傳統,變成害人命的枷鎖,吃人的禮教.耶穌的提醒很重要:人才是污染的源頭,身外之物不是.

讓我們總結一下什麼是傳統主義:1. 將人的吩咐當作真理教導人, 2. 借信仰傳統來攻擊異己,3. 把傳統轉為權力工具, 4. 把本來以人益處為本的傳統僵化為害人的枷鎖.

我們的信仰,是由上帝的話設立的,也經過二千年來先賢們對上帝的話的理解與實踐精煉出來的.我們承繼了豐厚而寶貴的傳統.我們信耶穌的,無論是剛剛洗禮,或是資深信徒,都要竭力避免誤墮傳統主義的圈套.

第一, 我們有責任仔細而嚴謹查考上帝的話語,以致懂得分辨,不會誤把人的吩咐當作真理奉行.

第二, 我們必須提高自覺,倘若我們發覺自己有意無意間借信仰教導來排斥與自己立場相異的人,我們必須回轉.

第三, 當我們發覺自己以認識多點信仰而自覺高人一等,該作領導,我們也必須回轉.信仰是讓人謙卑,不會叫人自高自大的.

第四, 如果我們發現,遵守某些傳統結果會對人無益,甚至會傷害人,我們就必須知道,我們已經不是真的在遵守傳統.相反,我們只是成為了傳統主義者.我們 必須回轉.

今天洗禮的也好,我們資深信徒也好,我們要謹記,我們是活潑信仰的見證人,是信仰傳統的守護者,萬萬不可衰化為傳統主義者.Pelikan 的提醒很重要:「傳統是已死去的前人的活潑信仰,傳統主義則是當下活著的人的僵死了的信仰.我想我應該補充一點,[傳統本身是美好的],只是傳統主義給傳統帶來了如此壞的名聲.」

至於到底我當年是否應該遵守游泳池的傳統規矩,就留待大家去思考了.阿們.



\newpage

\section{把孩子的餅給狗}
\label{sec:566}
\textbf{張亦農牧師 - 把孩子的餅給狗}
\newline
\newline
\hyperref[sec:565]{< < < PREV SERMON < < <}
~
\hyperref[sec:toc]{[返目錄]}
~
\hyperref[sec:567]{> > > NEXT SERMON > > >}


講題:把孩子的餅給狗 　Give kids’ bread to dog

經文:馬可福音7:24-37

講員:張亦農牧師

場合:香港中文大學崇基學院禮拜堂主日崇拜

日期:2021年9月5日

政治正確Political Correctness相信今時今日大家不再陌生的觀念.政治正確好可能是今日身處香港社會的你和我需要留心和學習的意識；就算去到世界各地生活,都要了解當地社會民情民風,禁忌界線,要懂得謹言慎行,才是明哲保身之道.今日三代經課的福音,馬可福音7:24-37節裡面,我們留意到耶穌基督說了一句政治唔正確的話.到底佢講左甚麼呢?

馬可福音提到耶穌前往外邦人聚居的城巿推羅,但聖經沒有告訴我們他去那裡的原因.

雖然馬可提到耶穌「不願意人知道」(24節),卻也沒有告訴我們原因.

馬可也指出耶穌根本想隱藏,也隱藏不了,很快就被發現.其中一個發現他的是一個希臘女人,馬可告訴我們這女人「屬敘利亞的腓尼基族」(25節)

耶穌講了一句極具冒犯的話:「拿孩子的餅丟給小狗吃是不妥的.」(27節)

曾經流傳民國時期,上海外灘有公園張貼有「華人與狗不得內進」的告示牌,被視為對中國人的侮辱.當然,歷史學家經過多番考證,否定了這個告示牌的存在.把人和狗類比,對中國人而一種侮辱,對外邦人也是一種侮辱.耶穌講呢句話,是一句政治不正確的話.

狗在舊約和新約中都是有一種負面的認知,認為牠們是污穢不潔的.甚至耶穌也曾提醒:「不要把聖物給狗」(太7:6),保羅也曾提到「應當防備犬類」(腓3:2)用普通話來形容:「不是好東西」.

耶穌的話,明眼人一聽就知道,他把以色列人比作人,把外邦人比作狗.把外邦人比喻作狗,在外邦人耳中肯定不會是好事,耶穌為要在一個外邦人聚居的城鎮中講出這番話,真的有點找死的味道.

耶穌為甚麼要打這個比喻呢?他沒理由不知道他講的那句話,用今日的講法,是帶有種族歧視成份.對那個一心想求耶穌救治她女兒的希臘婦人是一個實實在在的侮辱.

耶穌在推羅,按常理應該是暫住在猶太人的家中,當這個希臘婦人來求耶穌的時候,在耶穌身邊的除了門徒之外,相信還有其他猶太人.

若果耶穌真心認同這一句話,照理,他也不會後來願意為婦人的女兒趕鬼.所以,我個人比較傾向相信耶穌講這句話,是專登引述,反映猶太人的想法,認為上帝的恩典只屬於以色列人--上帝的選民.而耶穌為那婦人的女兒趕鬼,就是用行動表達了他不認同猶太人的狹隘想法.在耶穌眼中,不論是猶太人,還是外邦人,只要願意憑信心去尋求仰望上帝,天父也必會垂憐他們.

種族歧視一直是西方社會極其關注的事情.因為美國接二連三發生白人警察過度執法造成多位無攻擊性的黑人死亡,過去這一兩年引發世界各地出現Black Lives Matter大型示威活動,要對族種歧視說不.可惜的事,種族歧視並沒有煙消雲散.因為美國政客的政治操作,Covid-19常被稱為中國病毒(China Virus),結果在美國甚至在大都會紐約出現針對東亞人士的Asian Hate Crime.讓人傷心的是,在紐約,閉路電視影到1名黑人毒打1個65歲的亞裔婦人.黑人常常聲稱自己深受種族歧視之苦,卻也用種族歧視的態度對待無辜的亞裔人士.

孔子話:「己所不欲,勿施於人」,主耶穌基督也告訴我們:「無論何事,你們想要人怎樣待你們,你們也要怎樣待人,因為這就是律法和先知的道理.」(太7:12)

有時,我們也需要這份醒覺.有一次,我在一個與外國華人教會弟兄姊妹查經聚會,在言談中,我用左「鬼佬」來形容西方白人.後來,有一位當地教會的弟兄,短訊給我,提醒我要避免使用「鬼佬」這具冒犯性的用語.我好感謝該位弟兄好言相勸,不竟生活在一個華人佔90\%+的社會,未必為意到自己的說話會傷到人的心.

約3:16提醒我們「上帝愛世人」,我們的天父不單愛以色列人,也愛外邦人.確實福音書中,也多次提及耶穌幫助醫治外邦人的事蹟:他曾經去撒馬利亞與一個婚姻關係一團糟的婦人談道,因羅馬軍長的信心而治好他僕人的病,又醫治了撒瑪利亞人的痳瘋病.在今日的馬可福音,耶穌同樣為這個希臘婦人的女兒趕鬼.

其實耶穌引述了這一句現在看來甚具冒犯的話回應那個希臘婦人的請求,確實有點把這婦人推埋牆角的味道.婦人可以因為耶穌的話感覺受侮辱而含怒而去.又可以因為認為耶穌的拒絕而無奈離去.

婦人沒有離開,因為她知道耶穌就是她女兒唯一的希望.所以,她堅決抓緊這個機會,沒有放手,並且回應了耶穌的話.婦人用自己養育女兒的經驗,指出小孩子吃東西,難免天一半地一半,若果家中有狗的話,狗也會把子孩子掉在地上的食物吃了.即使主人看見,也不會阻止甚至責罰自己的狗.

婦人的答覆,沒有質疑上帝揀選以色列人作選民,但也相信上帝不會完全置外邦人於不顧,特別是那些從心仰望尋求祂的外邦人,上帝也會憐憫他們.這份認信得到了耶穌的肯定,讓耶穌拯救了婦人的女兒.

這位希臘婦人的經驗告訴我們:當我們認清事實,只有耶穌基督能夠幫助到我們,是我們唯一的出路的時候,我們才會全心全意去尋求耶穌的幫助,仰望他的憐憫.

當我們生命陷於絕境時,我們對主耶穌基督的經歷特別來得深刻和實在.

使徒彼得曾經清楚地宣告:「除他以外,別無拯救,因為在天下人間,沒有賜下別的名,我們可以靠著得救.」(徒4:12)係信仰路上,唯有我們認定耶穌是我們唯一的拯救,我們才會打從心底全然仰望和倚靠耶穌基督,與上帝建立更密切的關係.求主幫助我們.



\newpage

\section{你介意甚麼}
\label{sec:567}
\textbf{李均熊博士 - 你介意甚麼?}
\newline
\newline
\hyperref[sec:566]{< < < PREV SERMON < < <}
~
\hyperref[sec:toc]{[返目錄]}
~
\hyperref[sec:568]{> > > NEXT SERMON > > >}


馬可福音 8:27-38
\newline
\begin{longtable}{cl}
\hline
\hline
章節 & 經文 (和合本修訂版)\\
\hline
8:27 & \begin{tabularx}{0.7\textwidth}{X} 耶穌和門徒出去，往凱撒利亞‧腓立比附近的村莊去。在路上，他問門徒：「人們說我是誰？」 \end{tabularx} \\ \\ \relax
8:28 & \begin{tabularx}{0.7\textwidth}{X} 他們對他說：「是施洗的約翰；有人說是以利亞；又有人說是先知中的一位。」 \end{tabularx} \\ \\ \relax
8:29 & \begin{tabularx}{0.7\textwidth}{X} 他又問他們：「你們說我是誰？」彼得回答他：「你是基督。」 \end{tabularx} \\ \\ \relax
8:30 & \begin{tabularx}{0.7\textwidth}{X} 於是耶穌切切地囑咐他們不可對任何人說起他。 \end{tabularx} \\ \\ \relax
8:31 & \begin{tabularx}{0.7\textwidth}{X} 從此，他教導他們說：「人子必須受許多的苦，被長老、祭司長和文士棄絕，並且被殺，三天後復活。」 \end{tabularx} \\ \\ \relax
8:32 & \begin{tabularx}{0.7\textwidth}{X} 耶穌明白地說了這話，彼得就拉著他，責備他。 \end{tabularx} \\ \\ \relax
8:33 & \begin{tabularx}{0.7\textwidth}{X} 耶穌轉過來看著門徒，斥責彼得說：「撒但，退到我後邊去！因為你不體會神的心意，而是體會人的意思。」 \end{tabularx} \\ \\ \relax
8:34 & \begin{tabularx}{0.7\textwidth}{X} 於是他叫眾人和門徒來，對他們說：「若有人要跟從我，就當捨己，背起自己的十字架來跟從我。 \end{tabularx} \\ \\ \relax
8:35 & \begin{tabularx}{0.7\textwidth}{X} 因為凡要救自己生命的，必喪失生命；凡為我和福音喪失生命的，必救自己的生命。 \end{tabularx} \\ \\ \relax
8:36 & \begin{tabularx}{0.7\textwidth}{X} 人就是賺得全世界，賠上自己的生命，有甚麼益處呢？ \end{tabularx} \\ \\ \relax
8:37 & \begin{tabularx}{0.7\textwidth}{X} 人還能拿甚麼換生命呢？ \end{tabularx} \\ \\ \relax
8:38 & \begin{tabularx}{0.7\textwidth}{X} 凡在這淫亂罪惡的世代，把我和我的道當作可恥的，人子在他父的榮耀裡與聖天使一同來臨的時候，也要把那人當作可恥的。」 \end{tabularx} \\ \\
[1ex]
\hline
\hline
\end{longtable}
講題:你介意甚麼? 　What do you mind?

經文:馬可福音8:27-38

講員:李均熊博士

場合:香港中文大學崇基學院禮拜堂主日崇拜

日期:2021年9月12日

引言

今天經課內容仍然十分豐富.耶穌問門徒人說他是誰,再問門徒他們說他是誰,表達了當時眾人和門徒對耶穌身份的了解.當彼得承認耶穌就是基督,就是猶太人切望的彌賽亞,耶穌便明明預告他的受苦受難,是馬可福音三次耶穌預告他的受難的第一次.因著這個預告,經文說彼得將耶穌拉著和責備他!但「耶穌轉過來看著門徒,斥責彼得說:「撒旦,退到我後邊去!因為你不體會神的心意,而是體會人的意思.」之後耶穌教導門徒若要跟從他,就要捨己,背起自己的十字架來跟從他,又謂人若賺得全世界卻賠上生命,有甚麼益處呢?耶穌和門徒的對話,每一個段落都包含深刻的神學,使我們明白怎樣為之作門徒.不過,我一如平常的只選擇其中一個段落去同大家宣講,就著一個焦點去反思我們今日的景況.今日我選擇同大家分享ge 主題係耶穌和彼得的對答,甚至可以說他們在互相責備(廣東話:「互片」).究竟他們的對話,對我們有甚麼意義呢?讓我們先禱告.【祈禱】

起:彼得責備耶穌

若然你問我,為何選擇分享耶穌與彼得的對話,而不是集中講耶穌對門徒的吩咐?可能是因為一方面我有些害怕耶穌的要求,而另一方面我比較同情彼得.試想想:耶穌呼召彼得,跟從他,一路走來,耶穌都顯出救世彌賽亞的能力,治病,趕鬼,帶著權柄的教導……凡此種種,都令人對這個彌賽亞能救我們脫離這黑暗邪惡世代大有信心,但在這個大希律為凱撒奧古士督建有神殿的凱撒利亞腓立比,這位我們承認為彌賽亞基督的耶穌,卻說他必須受許多的苦,被棄絕,被殺(應該聽到這裡,大家已經聽不進三天後復活這句話,更可況大家根本不明白耶穌在說甚麼)……試想想,你是彼得,你會有何反應?

彼得是誰?我們可能有一個印象認為彼得乃是門徒的代表,為大家說話；但馬可福音卻令我們看見,「你是基督」可能是代表門徒的宣稱,但彼得接著的說話就顯示他是一個立體的角色,並非純粹門徒的集體意識.正是他豐富的性格特質,令讀者,起碼我,能認同彼得.甚麼?我們的主要受苦受死?甚麼?我們的主不是要申張正義,驅趕奸邪,帶來光明,使天國降臨的嗎?於是彼得拉著耶穌,責備他:「主啊,大家對你期望這麼高,為甚麼要受苦受死,你就用你治病趕鬼的能力,成就復興我們民族的大義啦!」我們對耶穌,是否也是這樣期望?我們是否也希望受苦不是「必須」?大能上帝為何甘願受苦?永生的神為甚麼要經歷死亡?這樣真的可以改變世間的強權?這樣真的可以救這荒謬的世代重現光明嗎?

耶穌轉過來,看著的不是彼得,而是門徒.若然作我的門徒來跟從我,就要捨己,背起屈辱的十字架來跟從我.耶穌看見的是門徒,不是彼得.彼得抗拒受苦,期盼光明,渴求救贖,追求上帝的榮耀.但門徒卻要捨去世上的生命,承受屈辱,必須受苦,並要在這淫亂罪惡的世代,不以受苦受死的基督為可恥!

我們都是彼得.但我們可以成為基督的門徒嗎?

承:跟從與歸邊(take side)

當耶穌轉過來看著門徒,就斥責彼得:「撒旦,退到我後邊去!因為你不體會神的心意,而是體會人的意思.」這一句看似嚴厲的斥責,可能需要仔細解釋一下.

首先,耶穌稱呼彼得為撒旦,他又怎可能仍是他的門徒?其實有幾個可能性.當然在我們心目中,撒旦就是黑暗之子,邪惡化身(demonic),魔鬼代表；但撒旦在聖經時期可以純粹解作敵對者,控訴者.雖然我們不排除這裡有可能指彼得不自知地被鬼王附身或控制,但更可能的解釋是耶穌斥責彼得乃抵擋耶穌使命之人,是耶穌工作的對頭和絆腳石,質疑和對抗上帝的道,而非邪靈或鬼王(king of devil)化身.

若是這樣,「退到我後邊去」亦非如趕逐惡鬼般要彼得離開.這裡亦可以有兩種理解.其一,就是叫彼得不要再阻擋耶穌,不要成為上主工作的絆腳石；另外,「到我後邊去」就是要彼得重新再跟從自己,英文:back into line.耶穌是叫彼得不要因為一時迷失,脫離了門徒之道,而一直走迷,而是回復初心,重回跟在耶穌背後的門徒之路.

最需我們細想的,可能是這一句「因為你不體會神的心意,而是體會人的意思.」很多英文譯本都類似NIV將這句譯成:You do not have in mind the concerns of God, but merely human concerns.好像說到彼得想法上出現問題,就如說彼得在心思中究竟關心甚麼.中文和合本說彼得「不體貼神的意思,只體貼人的意思」；和修將之改譯為「不體會神的心意,而是體會人的意思」.另外一些中譯本將「體貼」譯成「思念」(新譯本),「意念」(呂譯),「思想」(標準譯本)等等.黃根春教授推動的《原文新約》則將這個動詞譯作「專注於」,在《新約希臘文詞典》進一步指這字可解作「衡量；老是想著；不斷考慮」等意思,或總結為「不斷地認真考慮著某些事情」.換句話說,照字面解釋,這句意思是說彼得整天花心思想著的並不是神的事情,而是人的事情,所以彼得才會生出敵擋耶穌工作的念頭,阻止耶穌踏上這條受苦的道路.但在2015年一位新約學者發表了一份專文,透過比較當代的希臘文學,他認為只有RSV能譯出這句說話要表達的意思,就是You are not on the side of God but of me.這反映出彼得的問題在於他忠誠的對像錯置了,他站在人的那邊,而不是站在神的那邊,那是歸邊(take side)的問題,委身於誰(commitment)的問題.字面上的「體貼」,「思念」,「不斷地認真考慮的事情」,所反映的不是思想的問題,而是忠試委身於誰的問題.

綜合來說,耶穌面對彼得的責備,並非因為他每時每刻都想著人的需要,人的反應,人的要求,而是當他阻礙耶穌的路徑,抗拒基督必須受苦的前路,反映他並沒有站在耶穌的一邊,沒有忠於他作為耶穌門徒的召命,委身的只是人的想法和渴求,而非上主受苦的十字架.

耶穌轉身看到的是離開門徒之道的彼得,但耶穌眼中仍是他心中的門徒,於是耶穌對他說:抗拒這條受苦之路的彼得,退開,不要阻擋我的路,成為我絆腳石,卻要回到門徒的隊伍中,重新跟隨著我.因為你並不站在上帝的那邊,卻站在人的想法和要求的那邊,你重新想想你效忠的對像吧!

轉:馬可群體的訊息

不知從這個角度去思考彼得和耶穌的對答,你會有怎樣的反應,但對當時讀者來說,可能有更深刻的領受.馬可福音第一代的讀者,聖經學者認為,很大可能是身處羅馬,而且曾經歷過尼祿王對基督徒殘酷的逼害.相傳尼祿為要打造一個盛世羅馬,於是派軍隊一把火將舊羅馬城燒毀.雖然現今史家都懷疑此說,但無論事實如何,終究認為尼祿放火的傳言已經四起,當時的尼祿就找了基督徒,一群當時受羅馬人所          討厭,又不承認凱撒的神性,卻拜這個被針十字架的叛國賊耶穌的信徒,去背這個黑鍋,以至對基督徒展開殘酷的逼害(維基百科「尼祿」條目:「於是他下令逮捕基督徒,公開地將他們殘酷折磨--釘十字架,披獸皮讓惡狗咬死,將他們釘上柱子做為蠟燭.」).當時的逼迫有可能令很多信徒失去信心,甚至如彼得在耶穌受審當晚不承認他一樣,沒有承認自己是主的門徒.所以彼得其實就是當時讀者可認同的對像.馬可福音也可能透過講述彼得的失迷和耶穌挽回,去幫助當時因為受苦而一時放棄跟從耶穌的信徒門,能回到站在耶穌一邊,跟從耶穌的路徑上,並重新背起自己的十字架.

馬可福音中的彼得並不是一個門徒領袖的化身,卻是和你和我一樣,會在苦難面前害怕,甚至有不少時刻會「轉tie 」/「歸錯邊」,將自己的忠誠放錯在人的祈求之上.在承認耶穌是基督的當兒,卻抗拒受苦的道路；在耶穌受苦的當兒,三次不承認他是自己的主!在尼祿的殘酷逼害下,信徒有會否將自己對主的忠誠拋棄,靠著安全的一邊站立,看著弟兄受苦而心裡退縮呢?(這令我想起《沉默》這套電影)耶穌當然要斥責這樣的信徒:你這是在敵擋神,做了撒旦的幫兇了.但耶穌卻要我們回到他身邊,重新歸隊,不要在將忠誠壓在人的那邊.因為「凡為我和福音失喪生命的,必救自己的生命.人就是賺得全世界,卻賠上自己的生命,有甚麼益處呢?」

合:你介意甚麼?

今日我們作為耶穌門徒,並沒有如當年羅馬信徒飽受殘害和逼迫.但面向今日的境況,又有誰能說自己可以毫不猶疑地向耶穌歸邊,因為跟從耶穌而捨己,又能背起屈辱的十字架而不退縮?你究竟介意甚麼( What do you really mind)?我說我能認同彼得,是因為面對可能的受苦,我也會猶疑.當我的忠誠受到測試時,我又會體貼人的意思,還是會體貼神的意思呢?彼得抗拒基督的受苦,甚至三次不認主,使我知道,作為人的我是有機會迷失,退縮,甚至可能會背叛我的主.那時,彼得被基督的斥責和挽回,卻是我能重新面對自己,重新再走門徒之路,再有勇氣背上十字架的動力和提醒了.



\newpage

\section{們在路上議論的是甚麼}
\label{sec:568}
\textbf{高國雄牧師 - 「你們在路上議論的是甚麼?」}
\newline
\newline
\hyperref[sec:567]{< < < PREV SERMON < < <}
~
\hyperref[sec:toc]{[返目錄]}
~
\hyperref[sec:569]{> > > NEXT SERMON > > >}


馬可福音 9:30-37
\newline
\begin{longtable}{cl}
\hline
\hline
章節 & 經文 (和合本修訂版)\\
\hline
9:30 & \begin{tabularx}{0.7\textwidth}{X} 他們離開那地方，經過加利利；耶穌不願意人知道， \end{tabularx} \\ \\ \relax
9:31 & \begin{tabularx}{0.7\textwidth}{X} 因為他正教導門徒說：「人子將要被交在人手裡，他們要殺害他；被殺以後，三天後他要復活。」 \end{tabularx} \\ \\ \relax
9:32 & \begin{tabularx}{0.7\textwidth}{X} 門徒卻不明白這話，又不敢問他。 \end{tabularx} \\ \\ \relax
9:33 & \begin{tabularx}{0.7\textwidth}{X} 他們來到迦百農。耶穌在屋裡問門徒說：「你們在路上議論的是甚麼？」 \end{tabularx} \\ \\ \relax
9:34 & \begin{tabularx}{0.7\textwidth}{X} 門徒不作聲，因為他們在路上彼此爭論誰最大。 \end{tabularx} \\ \\ \relax
9:35 & \begin{tabularx}{0.7\textwidth}{X} 耶穌坐下，叫十二個使徒來，說：「若有人願意為首，他要作眾人之後，作眾人的用人。」 \end{tabularx} \\ \\ \relax
9:36 & \begin{tabularx}{0.7\textwidth}{X} 於是耶穌領一個小孩過來，讓他站在門徒當中，又抱起他來，對他們說： \end{tabularx} \\ \\ \relax
9:37 & \begin{tabularx}{0.7\textwidth}{X} 「凡為我的名接納一個像這小孩子的，就是接納我；凡接納我的，不是接納我，而是接納那差我來的。」 \end{tabularx} \\ \\
[1ex]
\hline
\hline
\end{longtable}
講題:「你們在路上議論的是甚麼?」 　“What were you arguing about on the way?”

經文:馬可福音9:30-37

講員:高國雄牧師

場合:香港中文大學崇基學院禮拜堂主日崇拜

日期:2021年9月19日

30他們離開那地方,經過加利利；耶穌不願意人知道, 31因為他正教導門徒說:「人子將要被交在人手裡,他們要殺害他；被殺以後,三天後他要復活.」 32門徒卻不明白這話,又不敢問他.33他們來到迦百農.耶穌在屋裡問門徒說:「你們在路上議論的是甚麼?」 34門徒不作聲,因為他們在路上彼此爭論誰最大. 35耶穌坐下,叫十二個使徒來,說:「若有人願意為首,他要作眾人之後,作眾人的用人.」 36於是耶穌領一個小孩過來,讓他站在門徒當中,又抱起他來,對他們說: 37「凡為我的名接納一個像這小孩子的,就是接納我；凡接納我的,不是接納我,而是接納那差我來的.」

耶穌有什麼不想讓任何人知道的呢?

不單如此,更要因此而離開那地…

其重要性,隱密性可想而知

「人子將被殺…要復活.」在今天可能講/聽得很習慣,但對當時門徒來說,何等震撼!!!

不過在馬可福音記載,這是第二次耶穌談及此事,上主日講道內容就是耶穌第一次對受死復活的預言(可8:31「人子必須受許多的苦,被長老,祭司長和文士棄絕,並且被殺,三天後復活.」)非常相近的內容,預言的再次被宣告,明顯是耶穌很想門徒們對此事有更深入的理解.

以上種種,反映出耶穌很緊張這事:再次預言,不讓任何人知道,離開那地…

門徒的反應呢?

不明白!!!

真令人瞎眼!為什麼仍然不明白呢?

不單止不明白,知道自己不明白但又不敢問,為什麼呢?

在我們的經驗中,又有沒有相類似的情況呢?

回想第一次預言時…可8:32「耶穌明白地說了這話,彼得就拉著他,責備他.」…彼得被耶穌責備

或許因為害怕導致不敢問,關係緊張,恐懼等蓋過真正的需要…

自我價值觀,自身經驗讓人產生盲點,特別在信仰,這有可能超乎常理/甚至邏輯的事上,更容易成為障礙.

在耶穌二次發出將被殺,復活的預言中間,門徒中間發生了什麼事?

可9:2-8記述耶穌登山顯榮

可9:9-13下山的時候,耶穌囑咐他們說:「人子還沒有從死人中復活,你們不要把所看到的告訴人.」 門徒將這話存記在心,彼此議論「從死人中復活」是甚麼意思.他們就問耶穌:「文士為甚麼說以利亞必須先來?」(瑪拉基書4:5-6)耶穌說:「以利亞的確先來復興萬事.經上不是指著人子說,他要受許多的苦和被人輕慢嗎? 我告訴你們,以利亞已經來了,他們任意待他,正如經上指著他說的.」

可9:14-29

耶穌趕出的,是門徒在一般情況下趕不到的鬼,門徒被耶穌再一次展現的大能驚嚇到…

接著就是今天的經文.

以上2段記載有何相類同的地方呢?

耶穌基督能力的彰顯,超出門徒所經歷的,所明白的,產生疑惑

讓我們再次回到經文中一個非常重要,要處理的問題

耶穌的權柄及能力已經清晰可見,但所教導的卻非常難接受

反映出的就是門徒的不明白,不敢問,到底當時門徒是在什麼狀況中…

與我們今天在屬靈的路上所經歷的又是什麼呢?

我們又在什麼狀況當中呢?

我們又有什麼不明白,不敢問的事情及原因呢?

當門徒在這狀況持續了一段時間

V33當他們來到迦百農.耶穌在屋裡問門徒說:「你們在路上議論的是甚麼?」

耶穌向門徒問了一個正正回應當時的狀況的問題

馬可的記載,明顯地耶穌不是不知道/聽不到他們議論的內容

而是希望門徒能有所反省

上帝今天可能最關心我們的,就是我們今天到底最關心的人和事是什麼…這不但透視出我們現在所持有的價值觀,更反映出我們現在與上帝的關係.

因為V34門徒不作聲,因為他們在路上彼此爭論誰最大.

門徒的不出聲,當然很可能就是覺得自己在爭論不應該爭議的話題,一般理解為:「門徒應該要謙卑,作僕人,不應該追求權/位.」當然這對我們有反思的意義,但如果以爭取排名,追求卓越的心,成為努力的動力,也可能有其價值…

相信耶穌這裡不是不希望我們「成功」,而是希望我們能夠找到真正成功的秘訣.

所以V35耶穌坐下,叫十二個使徒來,說:「若有人願意為首,他要作眾人之後,作眾人的用人.」

耶穌沒有反對我們要努力,目標要「為首」,即是追求「成功」,「卓越」,成為眾人的領袖是應該的,耶穌基督的跟隨者不應該過著Hea的生活,不應該以「不爭名逐利」為原因/理由而懶散做事,甚至覺得教會群體就應該接納慢進度,低效率,低回報的計劃及行動模式…

耶穌提到「為首」是一種選擇,由一個願意的心開始,願意帶來實踐行動,就是「在眾人裡成為最低微的,而且是眾人的僕人」(原文新約)

耶穌用例子作更深入的說明:

v36於是耶穌領一個小孩過來,讓他站在門徒當中,又抱起他來,對他們說: 37「凡為我的名接納一個像這小孩子的,就是接納我；凡接納我的,不是接納我,而是接納那差我來的.」

小孩子在門徒眼中可能是被輕看的,低微的,人總想要接待大人物,因為會以此為榮,猶太人文化中的接待更是如此.

如果接待小孩子=接待耶穌=接待上帝

我們必須學習看到,抱起那一些需要被接待的微小

也要願意成為微小/接受自己是微小的,而被其他人接待

我們可能一直在思考及討論如何成為「大/首」的問題

不斷彰顯榮耀,最大的耶穌卻教導我們學習如何能夠成為「小」

求主打開我們受教的心!



\newpage

\section{擔當不起}
\label{sec:569}
\textbf{林豪恩先生 - 擔當不起}
\newline
\newline
\hyperref[sec:568]{< < < PREV SERMON < < <}
~
\hyperref[sec:toc]{[返目錄]}
~
\hyperref[sec:570]{> > > NEXT SERMON > > >}


民數記 11:4-6,10-16,24-29
\newline
\begin{longtable}{cl}
\hline
\hline
章節 & 經文 (和合本修訂版)\\
\hline
11:4 & \begin{tabularx}{0.7\textwidth}{X} 他們中間的閒雜人動了貪慾的心；以色列人又再哭著說：「誰給我們肉吃呢？ \end{tabularx} \\ \\ \relax
11:5 & \begin{tabularx}{0.7\textwidth}{X} 我們記得在埃及的時候，不花錢就可以吃魚，還有黃瓜、西瓜、韭菜、蔥、蒜。 \end{tabularx} \\ \\ \relax
11:6 & \begin{tabularx}{0.7\textwidth}{X} 現在我們的精力枯乾了。除了這嗎哪以外，在我們眼前甚麼都沒有。」 \end{tabularx} \\ \\ \relax
11:7 & \begin{tabularx}{0.7\textwidth}{X} 嗎哪好像芫荽子，看上去如同樹脂的樣子。 \end{tabularx} \\ \\ \relax
11:8 & \begin{tabularx}{0.7\textwidth}{X} 百姓四處走動撿取嗎哪，把它用磨磨碎或用臼搗成粉，在鍋中煮了做成餅，滋味好像油烤餅的滋味。 \end{tabularx} \\ \\ \relax
11:9 & \begin{tabularx}{0.7\textwidth}{X} 夜間露水降在營中，嗎哪也隨著降下。 \end{tabularx} \\ \\ \relax
11:10 & \begin{tabularx}{0.7\textwidth}{X} 摩西聽見百姓家家戶戶在帳棚門口哀哭。因此， 耶和華的怒氣大大發作，摩西看了也不高興。 \end{tabularx} \\ \\ \relax
11:11 & \begin{tabularx}{0.7\textwidth}{X} 摩西對耶和華說：「你為何苦待僕人？我為何不在你眼前蒙恩，竟把這眾百姓的擔子加在我身上呢？ \end{tabularx} \\ \\ \relax
11:12 & \begin{tabularx}{0.7\textwidth}{X} 這眾百姓豈是我懷的胎，豈是我生下來的呢？你竟對我說：『把他們抱在懷裡，如養育之父抱著吃奶的嬰孩，一直抱到你起誓應許給他們祖宗的土地去。』 \end{tabularx} \\ \\ \relax
11:13 & \begin{tabularx}{0.7\textwidth}{X} 我從哪裡拿肉給這眾百姓吃呢？他們都向我哭著說：『給我們肉吃！』 \end{tabularx} \\ \\ \relax
11:14 & \begin{tabularx}{0.7\textwidth}{X} 我不能獨自帶領這眾百姓，這對我太沉重了。 \end{tabularx} \\ \\ \relax
11:15 & \begin{tabularx}{0.7\textwidth}{X} 如果你這樣待我，倒不如立刻把我殺了吧！我若在你眼前蒙恩，求你不要讓我再受這樣的苦。」 \end{tabularx} \\ \\ \relax
11:16 & \begin{tabularx}{0.7\textwidth}{X} 耶和華對摩西說：「你要從以色列的長老中為我召集七十個人，就是你所認識，作百姓的長老和官長的，領他們到會幕，使他們和你一同站在那裡。 \end{tabularx} \\ \\ \relax
11:17 & \begin{tabularx}{0.7\textwidth}{X} 我要在那裡降臨，與你說話，把降給你的靈分給他們。他們就和你分擔帶領百姓的擔子，免得你獨自承擔。 \end{tabularx} \\ \\ \relax
11:18 & \begin{tabularx}{0.7\textwidth}{X} 你要對百姓說：『你們要為了明天使自己分別為聖，你們將有肉吃。因你們哭著說：誰給我們肉吃呢？我們在埃及多麼好！這聲音傳到了耶和華的耳中，所以他必給你們肉吃。 \end{tabularx} \\ \\ \relax
11:19 & \begin{tabularx}{0.7\textwidth}{X} 你們不只吃一天、兩天、五天、十天、二十天， \end{tabularx} \\ \\ \relax
11:20 & \begin{tabularx}{0.7\textwidth}{X} 而是整整一個月，直到肉從你們的鼻孔噴出來，使你們厭惡。因為你們厭棄那住在你們中間的耶和華，在他面前哭著說：我們為何出埃及呢？』」 \end{tabularx} \\ \\ \relax
11:21 & \begin{tabularx}{0.7\textwidth}{X} 摩西說：「跟我在一起的百姓，步行的男人就有六十萬，你還說：『我要把肉賜給他們，使他們可以整整吃一個月。』 \end{tabularx} \\ \\ \relax
11:22 & \begin{tabularx}{0.7\textwidth}{X} 難道宰了羊群牛群，就夠給他們嗎？或者把海中所有的魚都捕來，就夠給他們嗎？」 \end{tabularx} \\ \\ \relax
11:23 & \begin{tabularx}{0.7\textwidth}{X} 耶和華對摩西說：「耶和華的膀臂豈是縮短了嗎？現在你要看我的話向你應驗不應驗。」 \end{tabularx} \\ \\ \relax
11:24 & \begin{tabularx}{0.7\textwidth}{X} 摩西出去，把耶和華的話告訴百姓，並從百姓的長老中召集七十個人來，叫他們站在會幕的四圍。 \end{tabularx} \\ \\ \relax
11:25 & \begin{tabularx}{0.7\textwidth}{X} 耶和華在雲中降臨，對摩西說話，把降給他的靈分給那七十個長老。靈停在他們身上的時候，他們就說預言，以後卻沒有再說了。 \end{tabularx} \\ \\ \relax
11:26 & \begin{tabularx}{0.7\textwidth}{X} 但有兩個人仍在營裡，一個名叫伊利達，一個名叫米達。他們本是在那些登記的人中，卻沒有到會幕那裡去。靈停在他們身上，他們就在營裡說預言。 \end{tabularx} \\ \\ \relax
11:27 & \begin{tabularx}{0.7\textwidth}{X} 有一個年輕人跑來告訴摩西說：「伊利達和米達在營裡說預言。」 \end{tabularx} \\ \\ \relax
11:28 & \begin{tabularx}{0.7\textwidth}{X} 嫩的兒子約書亞，年輕時就作摩西的助手，說：「請我主摩西禁止他們。」 \end{tabularx} \\ \\ \relax
11:29 & \begin{tabularx}{0.7\textwidth}{X} 摩西對他說：「你為我的緣故嫉妒嗎？惟願耶和華的百姓都是先知，願耶和華把他的靈降在他們身上！」 \end{tabularx} \\ \\
[1ex]
\hline
\hline
\end{longtable}
講題:擔當不起 　Too Heavy for Me

經課:民數記11:4-6,10-16,24-29

講員:林豪恩先生

場合:香港中文大學崇基學院禮拜堂主日崇拜

日期:2021年9月26日

一. 引言

(網絡新詩)

做人,要做一個有擔當的人,

有擔當的人,

在事情沒來的時候不害怕,

有擔當的人,

在事情來了的時候不退縮.

做一個有擔當的人,

既可以讓自己遇事不驚,

也可以讓別人心境安穩.

有擔當的人是一棵大樹,

灑下片片陰涼,

自己卻從來默默無語.

有擔當的人是一座山,

讓別人覺得踏實可靠,

不至於一陣風來,

就轟然倒塌.

如果能夠,相信大家都想做一個有擔當的人.然而,人生中卻又難免有令人感到擔當不起」的時候.曾經有令你感到擔當不起的事情嗎?或者,經歷過擔當不起的日子嗎?這個問題讓人感到沉重,為了令氣氛輕鬆一點,我們借用一位網友的問題來進入話題:他問:「女朋友說我對她的愛太沉重了,她擔當不起,想要暫時分開一段時間,我該怎麼辦呢?」請大家給這位男孩子提供一些建議.

以下是其中一些網友回應如下:

「要學會給愛空間,不然會適得其反,你認為你愛他,可是他會看作是自己的負擔.」

「大概是你對她太好了,而她又不能給予你等量的愛.」

原來,一個人與另一個人的關係,都可以令人感到沉重得擔當不起.那麼,一個人與更多人的關係,又會是一種怎樣的狀況呢?今天的舊約經課中,以色列的偉大先知,帶領以色列人出埃及的神人摩西,也說出這句話:「擔當不起」.是甚麼事情或甚麼狀況令神人摩西也嘆氣「擔當不起」呢?

二.經文

未能滿足被帶領者的要求

11:4他們中間的閒雜人大起貪慾的心；以色列人又哭號說:「誰給我們肉吃呢?

11:5我們記得,在埃及的時候不花錢就吃魚,也記得有黃瓜,西瓜,韭菜,蔥,蒜.

11:6現在我們的心血枯竭了,除這嗎哪以外,在我們眼前並沒有別的東西.」

11:10摩西聽見百姓各在各家的帳棚門口哭號.耶和華的怒氣便大發作,摩西就不喜悅.

11:11摩西對耶和華說:「你為何苦待僕人?我為何不在你眼前蒙恩,竟把這管理百姓的重任加在我身上呢?

11:12這百姓豈是我懷的胎,豈是我生下來的呢?你竟對我說:『把他們抱在懷裡,如養育之父抱吃奶的孩子,直抱到你起誓應許給他們祖宗的地去.』

11:13我從哪裡得肉給這百姓吃呢?他們都向我哭號說:『你給我們肉吃吧!』

11:14管理這百姓的責任太重了,我獨自擔當不起.

從事件可見,摩西夾在群眾與耶和華之間,成為了磨心.他帶領以色列人出埃及走向應許之地,他是以色列人的領袖,以色列人是他所要照顧的群眾,然而,現在他卻沒有能力滿足以色列人的要求,因此感到非常沉重.你有沒有經驗過作為照顧者的沉重呢?不要以為要照顧一間大學才會感到沉重,就算照顧一個學院,或者一個項目,或者一個家庭,甚或一個人,都可以令人感到沉重得擔當不起的.你要照顧的對象有各方面,而且是連續不斷的需要,當你看到被照顧對象或要照顧的事情的需要未能被滿足的狀況,你也會感到壓力的.

或者有人說,群眾的要求都是過份的,他們貪得無厭,旨在製造麻煩.讓我們代入群眾的角色想一想,你被要求跟大隊離開住了幾百年的地方,放下工作和房屋,所追求的是甚麼呢?豈不是那被描繪為流奶與蜜的應許之地嗎?現在,不要說沒有流奶與蜜,連「魚,黃瓜,西瓜,韭菜,蔥,蒜」這些日常普通食物也沒有,只有被稱為嗎哪的東西.不知道嗎哪好不好吃,然而,就算好吃,一日三餐,一復一日都吃同一樣食物,你會有甚麼感覺呢?

從另一方來看,如果你是摩西,你聽到百姓的哭號,你有甚麼感受呢?你有甚麼想法呢?或者你會認為百姓的期望過分而感到憤怒；或者你會感到百姓的哭號成為你的壓力而感到沮喪；或者你會認為自己滿足不到百姓的需要而感到挫敗；或者其實你自己也對現況感到失望,對前路茫茫而感到迷失和無力.於是乎,你會感受到摩西嘆息說擔當不起.

又未能向任務交付者報告完成任務

擔當不起,不單單是因為未能滿足被照顧者的期望,同時也是未能向任務委託者交待.我們作為職員要向上司交待,別以為我們的上司不用向人交待,就是校長,也要向公眾交待.

摩西呢?他受耶和華上帝委派去執行這項出帶領以色列人埃及這項任務.現在出師未捷民先亂,他如何向耶和華上帝交待呢?

當群眾哭號之聲響徹全營的時候,耶和華上帝有甚反應呢?

11:10摩西聽見百姓各在各家的帳棚門口哭號.耶和華的怒氣便大發作.

你試想像一下,圍繞著你,是遠是近,都回盪著哭號聲；同一時間,當你向上望,耶和華卻怒氣大發,你如何是好呢?若你打算向耶和華說:「不要憤怒,我立即去安撫百姓,令他們不再抱怨,不再投訴.」你認為自己能夠說句話便叫百姓收聲嗎?或者你向百姓說:「乖乖不要哭,我立即向耶和華要求賜下你們想吃的食物給你們.」你認為耶和華會照單全收嗎?摩西知道他不能控制群眾,也不能控制上帝,他夾在中間,雙方都不是他可以滿足的,也不是他有能力處理的.

11:10摩西就不喜悅.

11:11摩西對耶和華說:「你為何苦待僕人?我為何不在你眼前蒙恩,竟把這管理百姓的重任加在我身上呢?

11:12這百姓豈是我懷的胎,豈是我生下來的呢?你竟對我說:『把他們抱在懷裡,如養育之父抱吃奶的孩子,直抱到你起誓應許給他們祖宗的地去.』

11:13我從哪裡得肉給這百姓吃呢?他們都向我哭號說:『你給我們肉吃吧!』

11:14管理這百姓的責任太重了,我獨自擔當不起.

於是,摩西決定退出.或許,夾在中間的人,除了摩西之外,還有不少都卻步或退出.

有人提醒,當下全球教會正面臨老齡化,空巢期,斷層,領袖同工荒,牧師難尋……教會到底出了什麼問題呢?其中一個情況可能是領袖不易為,承擔越多,要求越多,壓力越大,誰能擔當得起呢?

大家都說有能力又肯承擔的領袖不容易找,教會界也不例外.有一則聘牧笑話如下:

聘牧委員會經過研究,候選人名單和主要特點有:施洗約翰自稱開路先鋒,不修邊幅,喜歡野餐,講道容易傷害聽眾感情.彼得蒙召清楚,脾氣不好,三次不認主.保羅其貌不揚,有博士學位,不圓滑,會發脾氣,講道長,人際關係欠佳.提摩太年輕能講道,基層服事,個性懦弱,恐難面對大場面.耶穌講道冒犯人,爭議太大,得罪過長老等人.唯有加略人猶大,能力強,當家理財,關懷得體,行政出色,茲決定禮聘為牧師.

苦不堪言求解脫

若然可以全身而退,何妨瀟灑走一回.讓我們回到摩西的處境,他有甚麼退路嗎?後面對他恨之入骨的埃及法老,四週包圍著對他抱怨投訴的百姓,上有正在發怒的耶和華,根本無路也退.因此,他對耶和華作出如此的要求:

11:15我若在你眼前蒙恩,求你立時將我殺了,不叫我見自己的苦情.」

有時候,當事人在困境中真的會感覺到退無可退,在苦不堪言的狀況之下,唯有求徹底解脫.

2020年9月5日,葵涌發生倫常慘案!一名中度智障青年,因年滿二十一歲須離開院舍,由親人接返家照顧僅四日,其母疑感心力交瘁,在壓力「爆煲」下,有人前晚刻意與智障兒同房入睡,並趁其他家人熟睡,狠心用膠台布紮成的繩索將智障兒勒死,再割腕企圖自盡,圖與子共赴黃泉.警方到場以涉嫌謀殺罪,拘捕涉案母親並送院治療.在鄰居眼中,女疑兇是一名「好媽媽」,持家有道,任勞任怨照顧三名兒女,坊鄰對慘劇均感驚訝惋惜及難以置信.《星島日報》

照顧別人其實是相當吃力的,而這種不堪重負的感覺非常普遍.照顧者可能會發現生活中充滿挑戰,感到有壓力和孤獨.如果已經有心理健康問題的人,成為照顧者可能會更影響他們的狀況.

求助與分擔

摩西進退維谷,苦不堪言,求死解脫.其實,很多人求死並不是渴望死亡,只是想解決問題,耶和華當然知道摩西的狀況及困難,因此對於摩西的要求「求你立時將我殺了」,當然不會傻得照單全收.人在極端狀態之下所說的話,都未必適宜按字面理解,更不適宜按字面行動的.例如兩夫婦吵架時,其中一個衝口而出說:「你既然咁唔鍾意我,咁離婚囉!」另一個如果按字面行動,就恨錯難返了.說出這樣的話,反映出的是那一刻大家相處得非常困難,渴望關係能夠改善,希望問題可以解決.耶和華並非只聽到摩西激動地求死的說話,也聽到他所表達的困難,他沒有把焦點放在「求你立時將我殺了」這訴求,他聽得出摩西的困難其實是前一句「11:14管理這百姓的責任太重了,我獨自擔當不起.」如果是責任太重獨自擔當不起而求死,那麼,不用獨自擔當,有人分擔責任,就不用求死了,對嗎?

11:16耶和華對摩西說:「你從以色列的長老中招聚七十個人,就是你所知道作百姓的長老和官長的,到我這裡來,領他們到會幕前,使他們和你一同站立.

11:17我要在那裡降臨,與你說話,也要把降於你身上的靈分賜他們,他們就和你同當這管百姓的重任,免得你獨自擔當.

11:24摩西出去,將耶和華的話告訴百姓,又招聚百姓的長老中七十個人來,使他們站在會幕的四圍.

11:25耶和華在雲中降臨,對摩西說話,把降與他身上的靈分賜那七十個長老.靈停在他們身上的時候,他們就受感說話,以後卻沒有再說.

11:26但有兩個人仍在營裡,一個名叫伊利達,一個名叫米達.他們本是在那些被錄的人中,卻沒有到會幕那裡去.靈停在他們身上,他們就在營裡說預言.

能夠達到這個效果,必需要有兩項修件:第一,摩西願意與人分享權力；第二,有人願意與摩西一齊承擔責任.

關於第一項,經文沒有記載摩西對耶和華提出的解決方法的感受和想法,只是記載他按著耶和華的說話去執行.耶和華對摩西說:「我要把降於你身上的靈分賜他們,他們就和你同當這管百姓的重任,免得你獨自擔當.」對於有些人來說,「他們就和你同當這管百姓的重任」都未必能接受,因為如此一來就失去了自己的獨特權力和位置,就算死也不讓別人與自己同站同一級台階；有些人樂於接受有人分擔重任,不用獨自擔當那麼辛苦,然而,考慮到「我要把降於你身上的靈分賜他們」就有保留了,因此這樣就會失去了自己的獨特身分和超然角色.讓人分擔重擔,對有些人來說是不容易的,有沒有見過有教會一年五十二個主日都是同一位牧師講道呢?會不會是因為教會只有一個牧師呢?通常都不是,這些教會往往有幾十位牧者.幸好摩西沒有這些猶豫,不過,他的跟隨著卻另有憂慮:

11:27有個少年人跑來告訴摩西說:「伊利達,米達在營裡說預言.」

11:28摩西的幫手,嫩的兒子約書亞,就是摩西所揀選的一個人,說:「請我主摩西禁止他們.」

11:29摩西對他說:「你為我的緣故嫉妒人嗎?惟願耶和華的百姓都受感說話!願耶和華把他的靈降在他們身上!」

有些時候,當事人的確是想解決問題,也廣納解決問題的方法,而沒有考慮太多個人的得失,反而是跟隨者敏感計較,甚至興風作浪,唯恐世界不亂.雖然跟隨產生了嫉妒的反應,幸好摩西沒有和應,有容仍大,無欲則剛,讓事情能夠按著耶和華的建議解決.

回想那個葵涌倫常慘案,那位母親很可能也有「獨自擔當不起」的沉重壓力.社會大眾事後會問:「為甚麼不找人幫手?」如此的詰問說就容易,事實是艱難的.社會大眾對這些特殊人士投以了多少偏見呢?有沒有聽過有聲音說這些人是社會的包袱呢?設身處地想一想,這些父母對於自己生了這樣的孩子有甚麼感受呢?可能會感到不幸,也能感到羞恥,當然也有擔心和無力.父母調教自己的觀念,感受和情緒已經不容易,如果還要面對其他人的冷眼和厭煩,那還有勇氣和自尊去向別人和社會要求甚麼嗎?如果社會能夠去除有需要人士這些心理負擔,事情可能會容易些.

第二項條件是有人願意分擔.在摩西的處境,已經有七十個長老存在,現在似乎是進行權責再分配的安排.經文沒有記載長老們的想法和感受,我們只知道事情就按著耶和華的吩咐進行,解決了摩西「「11:14管理這百姓的責任太重了,我獨自擔當不起」的困難.

讓我們又回到葵涌倫常慘案,如果社會上有完善的照顧制度,讓這些有特殊需要成員的家庭不用過分擔心,這悲劇可能不會發生.其實,社會上一定有不能照顧自己的人,可能是先天缺憾所致,也可能是後天意外產生,更多只是一般人年紀大身體機能退化或者患痛而失去自我照顧能力.如果這些照顧只落在一位家人身上,或者只落在一個家庭身上,難免會有如摩西的呼喊:「獨自擔當不起」.我們的社會需要有七十位長老,分擔個別人士的重擔.從創造的角度來看,每一個人都是上主的兒女,即是說大家都是兄弟姐妹.如果我們能夠建立社群精神,視個別人士的困難為社群的困難,一齊尋求解決方法,我們的生存狀況可以有所不同.



\newpage

\section{人不可分開}
\label{sec:570}
\textbf{高國雄牧師 - 人不可分開}
\newline
\newline
\hyperref[sec:569]{< < < PREV SERMON < < <}
~
\hyperref[sec:toc]{[返目錄]}
~
\hyperref[sec:571]{> > > NEXT SERMON > > >}


馬可福音 10:2-16
\newline
\begin{longtable}{cl}
\hline
\hline
章節 & 經文 (和合本修訂版)\\
\hline
10:2 & \begin{tabularx}{0.7\textwidth}{X} 有法利賽人來問他說：「男人休妻合不合法？」意思是要試探他。 \end{tabularx} \\ \\ \relax
10:3 & \begin{tabularx}{0.7\textwidth}{X} 耶穌回答他們說：「摩西吩咐你們的是甚麼？」 \end{tabularx} \\ \\ \relax
10:4 & \begin{tabularx}{0.7\textwidth}{X} 他們說：「摩西准許寫了休書就可以休妻。」 \end{tabularx} \\ \\ \relax
10:5 & \begin{tabularx}{0.7\textwidth}{X} 耶穌對他們說：「摩西因為你們的心硬，所以寫這誡命給你們。 \end{tabularx} \\ \\ \relax
10:6 & \begin{tabularx}{0.7\textwidth}{X} 但從起初創造的時候，神造人是造男造女。 \end{tabularx} \\ \\ \relax
10:7 & \begin{tabularx}{0.7\textwidth}{X} 因此，人要離開他的父母，與妻子結合， \end{tabularx} \\ \\ \relax
10:8 & \begin{tabularx}{0.7\textwidth}{X} 二人成為一體。既然如此，夫妻不再是兩個人，而是一體的了。 \end{tabularx} \\ \\ \relax
10:9 & \begin{tabularx}{0.7\textwidth}{X} 所以，神配合的，人不可分開。」 \end{tabularx} \\ \\ \relax
10:10 & \begin{tabularx}{0.7\textwidth}{X} 他們到了屋裡，門徒又問他這事。 \end{tabularx} \\ \\ \relax
10:11 & \begin{tabularx}{0.7\textwidth}{X} 耶穌對他們說：「凡休妻另娶的，就是犯姦淫，辜負他的妻子； \end{tabularx} \\ \\ \relax
10:12 & \begin{tabularx}{0.7\textwidth}{X} 妻子若離棄丈夫另嫁，也是犯姦淫了。」 \end{tabularx} \\ \\ \relax
10:13 & \begin{tabularx}{0.7\textwidth}{X} 有人帶著小孩子來見耶穌，要他摸他們，門徒就責備那些人。 \end{tabularx} \\ \\ \relax
10:14 & \begin{tabularx}{0.7\textwidth}{X} 耶穌看見就很生氣，對門徒說：「讓小孩到我這裡來，不要阻止他們，因為在神國的正是這樣的人。 \end{tabularx} \\ \\ \relax
10:15 & \begin{tabularx}{0.7\textwidth}{X} 我實在告訴你們，凡要接受神國的，若不像小孩子，絕不能進去。」 \end{tabularx} \\ \\ \relax
10:16 & \begin{tabularx}{0.7\textwidth}{X} 於是他抱著小孩子，給他們按手，為他們祝福。 \end{tabularx} \\ \\
[1ex]
\hline
\hline
\end{longtable}
講題:人不可分開 　Let no one separate

經文:馬可福音10:2-16

講員:高國雄牧師

場合:香港中文大學崇基學院禮拜堂主日崇拜

日期:2021年10月3日

2有法利賽人來問他說:「人休妻可以不可以?」意思要試探他. 3耶穌回答說:「摩西吩咐你們的是甚麼?」 4他們說:「摩西許人寫了休書便可以休妻.」 5耶穌說:「摩西因為你們的心硬,所以寫這條例給你們； 6但從起初創造的時候,上帝造人是造男造女. 7因此,人要離開父母,與妻子連合,二人成為一體. 8既然如此,夫妻不再是兩個人,乃是一體的了. 9所以,上帝配合的,人不可分開.」 10到了屋裡,門徒就問他這事. 11耶穌對他們說:「凡休妻另娶的,就是犯姦淫,辜負他的妻子； 12妻子若離棄丈夫另嫁,也是犯姦淫了.」

在這神聖莊嚴又漂亮的崇基禮拜堂,除了每主的崇拜,學期中的學生週會及音樂會,相信最多人認識這裡就是由婚禮開始…

多年來主持不少的婚禮,當中很多有趣的事情,有時甚至是令我覺得匪夷所思,最常見的就是有新人問可否改動婚禮程序,當然當中所選用的經文及詩歌也有不少可協調的空間(這也是令同事們常常煩惱的事),我也大都願意溝通處理,但試過有新人希望改動證婚中的誓詞,立約的內容,例如刪去「無論苦樂禍福,富貴貧窮,平安疾病,我都一樣敬愛你…」因為覺得只提一同享福才夠浪漫,誓詞大煞風景；又有試過交換信物時不願意說:「…表示將我所有,都奉獻給你…」認為現代夫婦,婚後也要留空間,財產給自己,保障自己才是最重要等…

甚至有新人直接與我講:「婚禮只是我們倆人的事,與其他所有人無關,所以不需要理會其他人.」相信這表達應該包括作為證婚牧師的我…

回應這些情況我就會問:「你們選擇在教堂中以基督教儀式舉行婚禮,你除了期望拍下美麗的照片,有美好的回憶之外,對你們未來的婚姻,你有什麼期望呢?你們的婚姻中需要上帝嗎?」

在未進入婚姻之前,首先我們要記得基督徒舉行婚禮的意義是什麼:婚禮在天主教會的傳統7項聖禮(sacrament)中,是上帝的奧秘有形的展現,所以在信仰群體當中的婚禮不只是程序與儀式,而是上帝的臨在,雖然我們基督新教只承認聖水禮及聖餐禮,但我們仍然可以在婚禮中體會上帝的神聖,恩典的臨在,所以每次婚禮中的宣告,我也會講:「婚禮是神聖的,是莊嚴的,是為人帶來歡愉的…」

而在每次婚禮的禮儀在最後,牧師會為新人禱告祝福,然後就會宣告:「倆人由現在開始,成為正式夫婦,上帝所配合的,人不可分開…」

當然這宣告的權能來自上帝及聖經,今天的經課正正就提到這教導重要的來源及因由.

首先我們要問,為什麼耶穌會提到「上帝所配合的,人不可分開」這論述呢?

可10:2有法利賽人來問耶穌說:「人休妻可以不可以?」意思要試探他.

因為人明顯違背上帝創世時的教導:

創世記2:24

因此,人要離開父母,與妻子連合,二人成為一體.

已經成為一體的倆人當然不可以分開

申命記24:1-2

人若娶妻以後,見她有甚麼不合理的事,不喜悅她,就可以寫休書交在她手中,打發她離開夫家.婦人離開夫家以後,可以去嫁別人.

先不討論法利賽人為何/如何以此問題試探耶穌,我相信其實不是摩西主動容許人休妻,而是我們很清楚,當時在以色列民當中已經有休妻的事情發生,並且應該也廣泛,甚至一般情況已是慣常及認為可接受,所以摩西照顧人的軟弱,又要要求休妻的男人必須負起責任:寫休書,讓被休的妻得自由,得生活的保障…

我們現在身處的現況又是如何呢?

網上面書中一個名為「我老婆嫁左比Mirror導致婚姻破裂關注組」近日成為熱話,不少「鏡粉」丈夫在當中投訴自己太太的瘋狂行為,該關注組成立僅5日,成員人數就已經突破24萬.

當然這關注組應該大概只是以輕鬆,開玩笑的態度面對身邊另一半的追星行為,但以「導致婚姻破裂」等字眼,而又帶來這麼多的反應及回響,相信也或多或少反映出現代香港人婚姻關係脆弱的程度及情況.

根據香港政府統計處《香港的女性及男性主要統計數字Women and Men in Hong Kong Key Statistics 2020年版》,至1991年至2019年香港離婚宗數不斷增加,由1991年的6,295宗,上升至2019的21,157宗,30年之內,增加了超過3.3倍.2019年結婚44,247宗,離婚21,157宗,即離婚率達48\%,即差不多每2對新人結婚,就會有1對夫婦離婚.

根據中國民政部婚姻的數據,全國離婚率接近40\%；歐洲有幾個國家離婚率超過60\%,最高的比利時更超過70\%,即每10對新人結婚,就會有7對夫婦離婚…

我們身處在一個甚麼樣的世代呀?!

使徒保羅對婚姻中的夫婦有著很重要的教導

以弗所書5:22-33

22你們作妻子的,當順服自己的丈夫,如同順服主. 23因為丈夫是妻子的頭,如同基督是教會的頭；他又是教會全體的救主. 24教會怎樣順服基督,妻子也要怎樣凡事順服丈夫. 25你們作丈夫的,要愛你們的妻子,正如基督愛教會,為教會捨己. 26要用水藉著道把教會洗淨,成為聖潔, 27可以獻給自己,作個榮耀的教會,毫無玷污,皺紋等類的病,乃是聖潔沒有瑕疵的. 28丈夫也當照樣愛妻子,如同愛自己的身子；愛妻子便是愛自己了. 29從來沒有人恨惡自己的身子,總是保養顧惜,正像基督待教會一樣, 30因我們是他身上的肢體. 31為這個緣故,人要離開父母,與妻子連合,二人成為一體. 32這是極大的奧祕,但我是指著基督和教會說的. 33然而,你們各人都當愛妻子,如同愛自己一樣.妻子也當敬重她的丈夫.

首先我們明白,實際上不是每人也會進入婚姻,同時也要記得獨身是恩賜,是上帝所賜的,請所有弟兄姊妹也為婚姻中的弟兄姊妹守望,禱告,支援,我們是一個群體/一個基督的身體,我們需要大家的幫助.

二人成為一體,是極大的奧秘,是上帝的奧秘,也基督與教會關係最好的演繹,要知道從始至終,人也不能完全明白這奧秘,直到基督耶穌來到,福音完全的展開才會被完全透視.

保羅向以弗所教會教導有關夫婦關係的核心,就是丈夫要愛妻子,就如基督愛教會,妻子要順服丈夫,就如教會順服基督.

夫婦的關係,就是上帝福音展現及記號,也就是說,眾人見到我們的夫婦關係,就會見到上帝與教會,就會見到福音.

所以婚姻關係就是福音盟約的實體展現的最好例子.

作為丈夫或太太的,我們比較喜歡留心上帝向對方的教導,手指著對方比較習慣及舒服,但請專心回應上帝給我們自己位份的愛心信件:

上帝向太太們的忠告,「要完全向丈夫的順服」

首先我們要必須清楚:上帝有全權去定義人的關係,特別是祂所配合的夫婦關係.

但向別人順服,向丈夫完全的順服,好像已經不合潮流,現今世代為弱勢的婦女爭取應有的權益,在父權社會中的確是非常重要,但要知道聖經所教導的,與平權問題無關的.

這教導重點是讓我們認清上帝為人所賜予的角色,順服就是回應該角色/身份應有及最好的態度.

在教會內也有不同的角色,信徒要學習順服教會領袖/牧者是理所當然的,因為不同的角色,位置,就會帶來不同的功能,不同的問責對像,這就是肢體各按其職,成為基督的身體的首要必須認同的.

但這絕不代表領袖比其他信徒高級,或某人比另一人重要,他在上帝眼中更有價值,更被重視,這是絕對錯誤的!

因為我們每一個人在上帝眼中也是最重要的,是上帝用重價所買贖的,上帝把順服放在當中成為關係建立的基礎,我們未能接受要完全的順服,問題不在此教導當中,而是在於人的軟弱,人對所授予的權力的濫用,與及人的自以為事…

要記得耶穌基督在客西馬尼園中的禱告,祂向我們完全展示了何謂順服,耶穌基督順服於父的旨意,以致於死.難道我們就因此認為天父比耶穌重要嗎?

三一(Trinity)的信仰清楚教導我們父,子,靈同榮同尊,而耶穌基督親身示範,成為我們最好的榜樣.我們完全明白,很多時候,要完全順服是極難,上帝給予作為太太這功課,當然絕對不是要輕看其地位,而是要在夫婦關係中彰顯基督.

要讓人能夠在「太太向丈夫順服」當中看見「教會順服基督」,而歸榮耀予上帝,並得著福音,這才是當中的價值所在.

所以問題的核心好大程度就是因為我們不願意服順,這種心態,就是認為人際間的關係必須由我們自己話事,我要全權選擇我所愛的,我要全權掌控我生命中的一切,當然包括我的婚姻,這種想法及態度由舊約開始,直到現在也一樣!

同樣聖經也有對做丈夫的忠告:

要記得:當中的要求更多,更豐富(經文中更詳細)

重點在於:要愛妻子,就如基督愛教會.

如果丈夫與妻子的關係,就是基督與教會的關係的展現,除非基督離棄教會,又或教會不再願意作基督的新婦,否則這盟約必然繼續生效,因為這配合就不可以分開!

如若我們作為丈夫的愛太太的程度與基督愛教會程度一樣,婚姻關係必然沒有問題,因為上帝是守約施慈愛的主,所以基督忠於教會,無論教會多麼軟弱,有無數的問題,基督仍然忠於對教會(人)的約,而繼續用無盡的愛去愛,甚至願意犧牲自己!

基督專心愛教會

丈夫,請專心愛你的妻子,

照顧妻子如同照顧/愛自己的身體

能夠成為「上帝所配合的」最重要的是我們願意接受,要記得,夫婦關係的出現不是由人所決定,即是由上帝所設立的,

創2:18「耶和華上帝說:「那人獨居不好,我要為他造一個配偶幫助他.」

是上帝主動親自設立,使二人成為一體；所以主權在上帝,恩約在我們當中

各位弟兄姊妹,

我是特別對願意認真回應上帝的道的你所講的

無論是誰,請停止一切繼續破壞家庭關係的行為,這不合上帝的心意.

如若有與他人曖昧,甚至出軌,請馬上停止.

如若你認為離婚是現在考慮中的一種選擇,只需要對自己負責任,請停一停,要記得你們的結合是上帝所配合的.

請為人在婚姻中的軟弱,跌倒禱告,認罪及悔改.

我不是說問題會馬上解決,關係會馬上改善,不再需要為夫婦感情問題而掙扎,而是期望我們能夠先去尋求上帝,先面向賜我們婚姻的上帝,讓我們可以面向對方,在屬靈的夫婦當中可以體會三一上帝的特質,夫婦2人加上帝就會成為一個三角形,夫婦愈親近,就能與上帝更近；與上帝愈親近,夫婦2人也會愈近,等邊三角形是最平衡亦是最穩固,你要決定過一個有上帝的模式的生活,還是靠自己過著沒有上帝的模式的生活而繼續掙扎呢?

我也願意與大家一同努力,有需要請來找我及師母,我們願意為大家守望禱告.

要記得接著的經文提到

可10:13-16

13有人帶著小孩子來見耶穌,要耶穌摸他們,門徒便責備那些人. 14耶穌看見就惱怒,對門徒說:「讓小孩子到我這裡來,不要禁止他們；因為在上帝國的,正是這樣的人. 15我實在告訴你們,凡要承受上帝國的,若不像小孩子,斷不能進去.」 16於是抱著小孩子,給他們按手,為他們祝福.

門徒責備帶小孩來的人,可能因為覺得小孩不可愛,不成熟,成為阻礙,但耶穌不單接受,更抱起,按手,祝福,並宣告在上帝國的,就正正是這樣像小孩的人,或許我們覺得自己及伴侶也不可愛,不成熟,甚至成為別人的阻礙,但只要我們擁有單純的心,單單仰望主,主必然接納我們,求主賜福,誠心所願!



\newpage

\section{近代中國三位傑出的教育家}
\label{sec:571}
\textbf{吳梓明教授 - 近代中國三位傑出的教育家}
\newline
\newline
\hyperref[sec:570]{< < < PREV SERMON < < <}
~
\hyperref[sec:toc]{[返目錄]}
~
\hyperref[sec:572]{> > > NEXT SERMON > > >}


阿摩司書 5:6-7
\newline
\begin{longtable}{cl}
\hline
\hline
章節 & 經文 (和合本修訂版)\\
\hline
5:6 & \begin{tabularx}{0.7\textwidth}{X} 要尋求耶和華，就必存活，免得他在約瑟家如火發出，焚燒伯特利，無人撲滅。 \end{tabularx} \\ \\ \relax
5:7 & \begin{tabularx}{0.7\textwidth}{X} 你們這使公平變為茵蔯，將公義丟棄於地的人哪！ \end{tabularx} \\ \\
[1ex]
\hline
\hline
\end{longtable}
詩篇 90:12-17
\newline
\begin{longtable}{cl}
\hline
\hline
章節 & 經文 (和合本修訂版)\\
\hline
90:12 & \begin{tabularx}{0.7\textwidth}{X} 求你指教我們怎樣數算自己的日子，好叫我們得著智慧的心。 \end{tabularx} \\ \\ \relax
90:13 & \begin{tabularx}{0.7\textwidth}{X} 耶和華啊，我們要等到幾時呢？求你轉回，憐憫你的僕人們。 \end{tabularx} \\ \\ \relax
90:14 & \begin{tabularx}{0.7\textwidth}{X} 求你使我們早早飽得你的慈愛，好叫我們一生一世歡呼喜樂。 \end{tabularx} \\ \\ \relax
90:15 & \begin{tabularx}{0.7\textwidth}{X} 求你照著你使我們受苦的日子，和我們遭難的年歲，使我們喜樂。 \end{tabularx} \\ \\ \relax
90:16 & \begin{tabularx}{0.7\textwidth}{X} 願你的作為向你僕人們顯現，願你的榮耀向他們子孫顯明。 \end{tabularx} \\ \\ \relax
90:17 & \begin{tabularx}{0.7\textwidth}{X} 願主－我們神的恩寵歸於我們身上。願你堅立我們手所做的工，我們手所做的工，願你堅立。 \end{tabularx} \\ \\
[1ex]
\hline
\hline
\end{longtable}
約翰一書 4:7-12,18-19
\newline
\begin{longtable}{cl}
\hline
\hline
章節 & 經文 (和合本修訂版)\\
\hline
4:7 & \begin{tabularx}{0.7\textwidth}{X} 親愛的，我們要彼此相愛，因為愛是從神來的。凡有愛的都是由神而生，並且認識神。 \end{tabularx} \\ \\ \relax
4:8 & \begin{tabularx}{0.7\textwidth}{X} 沒有愛的就不認識神，因為神就是愛。 \end{tabularx} \\ \\ \relax
4:9 & \begin{tabularx}{0.7\textwidth}{X} 神差他獨一的兒子到世上來，使我們藉著他得生命；由此，神對我們的愛就顯明了。 \end{tabularx} \\ \\ \relax
4:10 & \begin{tabularx}{0.7\textwidth}{X} 不是我們愛神，而是神愛我們，差他的兒子為我們的罪作了贖罪祭；這就是愛。 \end{tabularx} \\ \\ \relax
4:11 & \begin{tabularx}{0.7\textwidth}{X} 親愛的，既然神這樣愛我們，我們也要彼此相愛。 \end{tabularx} \\ \\ \relax
4:12 & \begin{tabularx}{0.7\textwidth}{X} 從來沒有人見過神，我們若彼此相愛，神就住在我們裡面，他的愛在我們裡面得以完滿了。 \end{tabularx} \\ \\ \relax
4:13 & \begin{tabularx}{0.7\textwidth}{X} 因為神將他的靈賜給我們，由此我們知道我們是住在他裡面，而他也住在我們裡面。 \end{tabularx} \\ \\ \relax
4:14 & \begin{tabularx}{0.7\textwidth}{X} 父差子作世人的救主，這是我們所看見並且作見證的。 \end{tabularx} \\ \\ \relax
4:15 & \begin{tabularx}{0.7\textwidth}{X} 凡宣認耶穌為神兒子的，神就住在他裡面，而他也住在神裡面。 \end{tabularx} \\ \\ \relax
4:16 & \begin{tabularx}{0.7\textwidth}{X} 我們知道並且深信神是愛我們的。神就是愛，住在愛裡面的就是住在神裡面；神也住在他裡面。 \end{tabularx} \\ \\ \relax
4:17 & \begin{tabularx}{0.7\textwidth}{X} 由此，愛在我們裡面得以完滿：我們可以在審判的日子坦然無懼，因為基督如何，我們在這世上也如何。 \end{tabularx} \\ \\ \relax
4:18 & \begin{tabularx}{0.7\textwidth}{X} 在愛裡沒有懼怕；完滿的愛把懼怕驅逐出去，因為懼怕裡含著懲罰，懼怕的人在愛裡尚未得到完滿。 \end{tabularx} \\ \\ \relax
4:19 & \begin{tabularx}{0.7\textwidth}{X} 我們愛，因為神先愛我們。 \end{tabularx} \\ \\
[1ex]
\hline
\hline
\end{longtable}
馬可福音 12:28-34
\newline
\begin{longtable}{cl}
\hline
\hline
章節 & 經文 (和合本修訂版)\\
\hline
12:28 & \begin{tabularx}{0.7\textwidth}{X} 有一個文士來，聽見他們的辯論，知道耶穌回答得好，就問他說：「誡命中哪一條是第一呢？」 \end{tabularx} \\ \\ \relax
12:29 & \begin{tabularx}{0.7\textwidth}{X} 耶穌回答：「第一是：『以色列啊，你要聽，主—我們的神是獨一的主。 \end{tabularx} \\ \\ \relax
12:30 & \begin{tabularx}{0.7\textwidth}{X} 你要盡心、盡性、盡意、盡力愛主—你的神。』 \end{tabularx} \\ \\ \relax
12:31 & \begin{tabularx}{0.7\textwidth}{X} 第二是：『要愛鄰如己。』再沒有比這兩條誡命更大的了。」 \end{tabularx} \\ \\ \relax
12:32 & \begin{tabularx}{0.7\textwidth}{X} 那文士對耶穌說：「好，老師，你說得對，神是一位，除了他以外，再沒有別的了； \end{tabularx} \\ \\ \relax
12:33 & \begin{tabularx}{0.7\textwidth}{X} 並且盡心、盡智、盡力愛他，又愛鄰如己，要比一切燔祭和祭祀好得多。」 \end{tabularx} \\ \\ \relax
12:34 & \begin{tabularx}{0.7\textwidth}{X} 耶穌見他回答得有智慧，就對他說：「你離神的國不遠了。」從此以後，沒有人敢再問他甚麼。 \end{tabularx} \\ \\
[1ex]
\hline
\hline
\end{longtable}
講題:近代中國三位傑出的教育家 　Three Outstanding Educators in Modern China

經文:阿摩司書5:6-7;詩篇90:12-17;約翰一書4:7-12,18-19;馬可福音12:28-34

講員:吳梓明教授

場合:香港中文大學崇基學院禮拜堂主日崇拜

日期:2021年10月10日

各位早晨!多謝高牧師的邀請,讓我有機會與崇基禮拜堂眾弟兄姊妹一同聚集,一同敬拜上帝.祝願 上主的恩惠及慈愛常與你們同在!

今年是崇基學院創校70周年,在高牧師的領導下,我們設計了一個「香江薪傳」校史本地考察計劃,與崇基的師生校友們一同反思基督教大學在香港承傳與蛻變的課題.因為疫情的緣故,我們不能夠大事宣揚,所以參加的人數不多,但參加者都是非常認真投入地思考.我們知道,崇基是一所基督教大學,我們承傳著國內13所基督教大學的教育傳統; 我們是「崇奉基督」,以博愛為基礎,致力提供「文武(身體智慧)雙全,德才(道德才華) 兼備,靈命俱長,達至全人教育」.我們記得崇基先賢,創校第二任院長為我們定下「德,智,體,群,靈」五育,凌道揚院長說:『知識必須要加上靈性的教育,才能達到「止於至善」的境界』.甚麼是靈性教育? 原來崇基所著重的是一種具「基督精神」的教育,是具有「基督的大愛,犧牲,和服務精神」的教育.

今日的主日講題是:「近代中國三位傑出的教育家」.我想和大家介紹近代中國三位傑出的教育家,他們有一共通之處就是他們都是出身自教會大學,他們不一定是基督徒,但是他們受過教會大學的基督教教育的薰陶,能夠被上帝使用,因為在他們的身上,我們可以見到他們的生命是洋溢著「基督大愛,犧牲和服務的精神」.

第一位傑出的教育家是陶行知(1891-1946),陶行知在國內被尊崇為「近代中國偉大的人民教育家」.很少人知道陶行知是一位基督徒,他是1911年在南京金陵大學畢業(就是我們所說國內13所基督教大學之一的金陵大學),在大學期間信主的.陶行知立志要「讀書為救國」,當畢業後獲得一份獎學金,他決定要去美國讀政治學,準備回國後為新政府服務.不過,1915年的夏天,他在美國日內瓦湖畔參加了基督教青年會YMCA主辦的一個夏令營.在營中他得到一個特別的召喚, 上帝呼召他轉讀教育,要他透過公立教育-面向新中國,為上帝作工.陶行知隨即轉往哥倫比亞大學師範學院攻讀教育學博士學位,成為了杜威的學生,也是胡適在哥倫比亞大學的同學.

回國後,陶行知協助北京大學校長蔡元培先生(1868-1940)推動新中國教育的改革,積極推行一系列的教育改革運動,包括『平民教育運動』,『農村教育運動』,『普及教育運動』等.陶行知更獨力主張

『愛的教育』,為中國教育創造新文化.

1927 年,陶行知創辦一所實驗學校.曉莊師範學校雖然不是一所基督教學校,卻是注入了基督教信仰的元素:『愛』.陶行知說:

『曉莊是從愛裡產生出來的.沒有愛便沒有曉莊…愛之所在即曉莊所在』.『曉莊是從愛裡產生出來的』,這樣的思想是從那裡來的呢?對中國傳統教育文化來說,它不會是來自中國的人文哲學或甚麼主義思想吧；『愛的教育』就是一種新文化教育思維,這種新的文化就是來自基督教的信仰.陶行知必定是熟識《聖經》的教導,正如我們剛才所讀的新約經課《約翰壹書》4:7-8 所說:『愛是從神而來的! 凡有愛心的,都是由神而生,並且認識上帝.沒有愛心的,就不認識神; 因為神就是愛』.

華中師範大學周洪宇教授曾多次斷言陶行知後來放棄了基督教信仰,但他也不得不承認「愛的教育」是來自基督教的.他也坦言說:

『(陶行知)所奉行的「愛滿天下」及其偉大的犧牲精神,是與基督教的博愛主張和耶穌「捨己為人」的救世精神有著某種思想淵源』.國內的學者雖然不是基督徒,卻不得不承認在陶行知的教育思維中可以見到基督教信仰的影兒,他們仍然尊崇他為「近代中國偉大的人民教育家」.

「愛的教育」其實就是「中國版」道成肉身的基督教教育.南京金陵協和神學院汪維藩教授曾說過:「中國基督教」必須是「道(在中國) 成了肉身的基督教」.基督教必須要認同中國的歷史文化傳統,與中國人民共同經歷苦難,在歷盡滄桑的苦難中孕育出一種具中國特色的基督教神學來,這才配稱為「中國基督教」.同樣地,「中國化」的基督教教育也必須是道(在中國)成了肉身的基督教教育.陶行知所主張的「愛的教育」其實就是「道(在中國)成了肉身的基督教教育」一個典型的例子,是最具「基督精神」,並且最能為中國人民接納的基督教教育.

第二位傑出的教育家是魯潔教授(1930-2020),魯潔教授1980年代初擔任南京師範大學教育系系主任(1981-1986),她1994年創立了南京師範大學道基督教育研究所,成為近代中國度基督教育研究的一個重要的基地.魯潔在1947年入讀金陵女子大學(也是一所教會大學),是吳貽芳校長在金女大最後一批學生.魯潔雖然不是基督徒,但她亦坦言她是在一個充滿基督教氛圍的環境中長大的(她身邊的親戚中,大部分的長輩都是基督徒),她全家均出身自教會學校,接受過兩代的基督教教育,包括:她的父親[魯繼曾]就讀於之江大學,母親[楊進美]就讀於上海聖瑪利亞女子中學,姐姐[魯明]是在上海中西女中,哥哥[魯平](曾是港澳辦主任,基本法起草委員會秘書)畢業於上海聖約翰大學,而[魯潔]本身則是畢業於金陵女子大學.原來魯潔的父親更是陶行知在哥倫比亞大學的同學,同是讀教育學的.當魯潔在南京師範大學做教育系系主任的時候,她必定熟識南師大校園內所豎立陶行知的肖像,並且認同陶行知所主張的「愛的教育」和「以人為本」的教育理想.南師大校園內刻上陶行知的兩句名言:「捧著一顆心來,不帶半根草去」,和「千教萬教,教人求真,千學萬學,學做真人」,成為了南京師範大學推動

「人本」教育的目標和口號.

當被問及自己的教育理念的根源時,魯潔這樣回答說:『對於一個從事教育學專業的人來說,真正讓我體悟到教育的真正意義的,不是書本,而是…我所敬愛的老師們,曾經為我「帶過路」的同學們…有

我曾經用自己的心靈與之相互溝通的學生們…對於他們,我只能用自己的生命…去解讀他們的生命,在我慢慢地讀懂他們時,也在慢慢地讀懂了一本用生命為我打開的教育學』.事實上,魯潔在南師大所提倡的「人本」(以人為本)的教育,就是承傳了她從金女大的老師和同學們身上所學到金女大的校訓:「厚生」(Abundant Life =豐盛的生命)和「以學生為本」的教育,加上陶行知的「愛的教育」,這些教育理念對近代中國教育的影響確實是深遠的.

1980年代,鄧小平在國內推行改革開放的政策,教育界亦是百花齊放之時,魯潔把握機會,提出了她對中國教育的診斷.她說:

『「文革十年」是中國經歷一場「去文化」,「去教育」的年代.因此,從1980年代開始,中國教育必須進行另一場新的啟蒙運動,是「一種文化教育的復興運動」』.魯潔教授最值得敬佩的地方,不僅是因為她能夠對過往30年來中國教育所採用的教條化的馬克思主義和蘇聯教育模式提出批判,更加精彩的地方是她能夠在馬克思主義的哲學思想中, 重新尋獲教育與人的主體性問題的邏輯起點,她宣稱說:『教育是要引導人去創造一個新世界(走馬克思主義的路),同時(教育)也要引導人在創造世界的過程中提升自己(兼容人文主義的精神)』.魯教授為中國的道基督教育重新確立一個基本的使命,她說:『道基督教育的基本使命就是要引導人走上「成人之道」,去做成一個人』；又說:『教育的原點是「育人」,教育的旨歸是「成人」-使人得以做成一個人』.陶行知被尊崇為「偉大的人民教育家」,魯潔卻被譽為是「中國社會主義教育家」;不過,魯潔教授這種「成人」的教育,其實就是用另一種的方式承傳了金陵女子大學的「厚生-以人為本」的教育和陶行知所主張的「愛的教育」,它也確實是與我們所說「具基督精神」的教育思維有相近似的思想淵源.換句話說,魯潔是與陶行知的教育思想相近,同樣是為教會大學的「厚生」,「全人教育」,和「愛的教育」找到了一條新的出路.

第三位傑出的教育家是章開沅教授(1926-2021),我特別懷念魯潔教授和章開沅教授,因為我認識他們已有30年,魯潔教授是在去年[2020年]12月25日聖誕日離世,而章開沅教授剛剛在今年5月28日離世,短短的半年間痛失了兩位我至敬愛的老前輩.章教授曾是南京金陵大學的學生,他在金陵大學只停留了兩年(1946-1948),但估不到自此以後他卻是與教會大學結下了不解之緣.當時因為國共內戰,他輾轉間被

送往接近武漢的中原大學,(奇怪的是,1951年中原大學被新政府實施的院校調整,結果是與華中大學[另一所教會大學]合併,後來成為了今日的華中師範大學).也因為這個緣故,他認識了曾是華中大學首位華人校長韋卓民,與他共事多年,讓他能能夠從另一個角度深入些認識教會大學,也能進一步與教會大學「續緣」(re-connected).

1984年,章教授剛剛58歲時,他獲政府委任為華中師範大學的校長,第二年1985年初,美國普林斯頓大學劉子健教授回國訪問章教授, 他們兩位史學家都是1940年代中國教會大學的學生-劉子健教授是燕京大學,章教授是金陵大學的學生.在兩人談話間,劉教授突然提及研究中國教會大學的必要性及迫切性,並且他對章教授說:『我看你就是帶頭研究的最適合的人選… (三個原因) 第一,你原來是教會大學的學生,現在的工作單位前身也是教會大學,對教會大學有較多親身的體驗.第二,你現在是大學校長,對學校有總體的把握,比一般人更有利於研究教會大學.第三,特別是作為歷史學者,你既有這方面的專業訓練,更有責任來填補中國近代史研究中這塊空缺』.這真是一個突如其來(從上天而來)的挑戰(呼召),[奇怪嗎?] 章教授卻是欣然接受了.

1986年(章教授60歲的時候),他本來是一位著名的辛亥革命史專家,現在搖身一變,成為在國內推動中國教會大學史研究的領導人.30 年前,中國教會大學史研究是一門「險學」(危險的學科,學術界的禁區),章教授卻為我們衝破禁區,使教會大學史研究成為中國學術界認可的一門新學科(國內的大學開始重現基督教研究中心,宗教研究中心等學術單位).後來章教授因為專心研究他在金陵大學歷史系的老師貝德士(Miner Searle Bates)的檔案文獻,他更成為了南京大屠殺歷史研究的專家,也為中國教會大學的歷史研究作了許多「活化」(我們學者稱之為「處境化」,「中國化」)的工作和見證,使它成為南京大屠殺歷史研究一個重要的部分.

2017 年,崇基校牧室主辦了一次校史考察團,我們有幸到訪了華中師範大學章開沅教授.當時章教授已經是 91 歲的老人家,但在分享會上,他還是活力充沛地為教會大學作見證,告訴我們:『基督教就是愛的宗教』,教會學校讓他學會了『教育的靈魂就是愛,沒有愛的教育就是沒有靈魂的教育』.當時,坐在他旁邊的正是華中師範大學的黨委書記,但章教授仍然無懼地與我們分享過去30 年來他是怎

樣推動中國教會大學史的研究,並亦指出教會大學的基督教教育是與近年中國高等教育界所流行的「人文素質教育」及「通識教育」頗有相通之處.換句話說,章教授也確認了「教會大學的基督教教育」委實是近代中國教育思潮中一種另類的「通識教育/博雅教育」,也可算是近年流行的「人文素質教育」的一部分吧.

在今年 7 月校史考察團的總結分享會上,我們一位 50 年代的校友(大師兄)提醒我們崇基學院是「崇奉基督」,我們著重的「基督精神」,是要在學生的生命中表現出「基督的愛與服務」的精神.剛才提到三位傑出的教育家,他們都是出身自教會大學,受過教會大學的基督教教育的薰陶,更能被上帝使用,因為從他們的身上可以見證到他們是承傳著基督教大學所信奉的「基督精神」,也為教會大學作出了美好的見證.因為他們所宣揚的不僅是「愛的教育」,更是「厚生=豐盛生命」的教育; 他們的獻身不僅是為社會主義的中國服務,也是為近代中國社會/教育思維拓展出一套「以人為本,以愛為基」的道基督教育來.我也發覺,若果我們是從這樣的角度看教會大學的基督教教育的話,今日我們所讀的經課便會顯得更有意思的.

今日的舊約經課,阿摩司先知勸籲以色列人要尋求耶和華,就必存活(摩 5:4).尋求耶和華是甚麼意思?原來有三個重要的元素:

「 尋求公義(Righteousness) , 尋求善 (Goodness) , 懂得愛/ 施恩

(Love/Grace)」(摩 5:7,14,15).新約書信約翰壹書 4:8 更清楚的說:

「神就是愛」,尋求耶和華就是尋求愛,「因為愛是從神而來,凡有愛心的,都是由神而生,並且認識神」(約壹 4:7).我剛才提及的三位教育家,都有一個共同點,是主張「愛的教育」,陶行知說:『曉莊是從愛裡產生出來的.沒有愛便沒有曉莊…愛之所在即曉莊所在』.魯潔教授也強調要愛我們的學生,幫助我們的學生成人(教人求真,學做真人)； 章開沅教授更確認說:『基督教就是愛的宗教』.我們剛剛所讀的新約福音書信也說明了耶穌的看法,當文士追問耶穌誡命中那一條是最緊要的時候,耶穌回答說:『第一要緊的是:你要盡心,盡性,盡意,盡力,愛主你的上帝; 其次就是說:要愛人如己』(可 12:30-31).換句話說,基督教就是關乎「愛神愛人」的宗教.

1950   年代,教會大學在國內消失了,取而代之是馬克思主義及蘇

聯教育制度的模式,但到了80年代,隨著國內改革開放政策的推行,

過往教會大學所推行的通識教育,博雅教育,包括以人為本的教育,愛與服務的教育,在國內又陸續重現出來.可以說,過往教會大學所推行的教育傳統現在是以另類(蛻變)的方式被承傳下來了!這豈不也是一種

「道成了肉身」(蛻變出來)的基督教教育嗎?

崇基學院創校至今已經踏進 70 周年,我們一直是承傳著13 所教會大學的基督教教育傳統,我們也是確信,「基督教就是愛的宗教」,「基督教教育就是愛的教育 - 包括愛上帝,愛國家,愛人民」.我們是「崇奉基督」,宣揚及實踐基督的大愛與服務的精神,為學生提供「文武(身體智慧)雙全,德才(道德才華) 兼備,靈命俱長, 達至全人的教育」.我們是從事教育工作者,我們相信,過去的先賢,我們的老師,教授們,同學校友們,也是傳承著基督大愛與服務的精神,並且發揚光大,就像剛才提及的三位傑出教育家一樣.教育是屬於播種的工作,我們也是播種了「基督的大愛和服務的精神」,崇基學院過去的 70 年是這樣,今天我們仍然是這樣,我們要努力地播種這種具「基督精神」的教育種子,並且祈求上主繼續使用我們,並加上(上主的)祝福.就好像今日我們所誦讀的詩篇經課 90 篇最後一節(第 17 節)所說的,我們不斷的禱告,就是:『主啊!願你堅立我們手所作的工,我們手所作的工,願你堅立』.啊們!



\newpage

\section{人生或苦不亢不卑}
\label{sec:572}
\textbf{林豪恩先生 - 人生或苦 不亢不卑}
\newline
\newline
\hyperref[sec:571]{< < < PREV SERMON < < <}
~
\hyperref[sec:toc]{[返目錄]}
~
\hyperref[sec:573]{> > > NEXT SERMON > > >}


以賽亞書 53:3-12
\newline
\begin{longtable}{cl}
\hline
\hline
章節 & 經文 (和合本修訂版)\\
\hline
53:3 & \begin{tabularx}{0.7\textwidth}{X} 他被藐視，被人厭棄；多受痛苦，常經憂患。他被藐視，好像被人掩面不看的一樣，我們也不尊重他。 \end{tabularx} \\ \\ \relax
53:4 & \begin{tabularx}{0.7\textwidth}{X} 他誠然擔當我們的憂患，背負我們的痛苦；我們卻以為他受責罰，是被神擊打苦待。 \end{tabularx} \\ \\ \relax
53:5 & \begin{tabularx}{0.7\textwidth}{X} 他為我們的過犯受害，為我們的罪孽被壓傷。因他受的懲罰，我們得平安；因他受的鞭傷，我們得醫治。 \end{tabularx} \\ \\ \relax
53:6 & \begin{tabularx}{0.7\textwidth}{X} 我們都如羊走迷，各人偏行己路；耶和華使我們眾人的罪孽都歸在他身上。 \end{tabularx} \\ \\ \relax
53:7 & \begin{tabularx}{0.7\textwidth}{X} 他被欺壓受苦，卻不開口；他像羔羊被牽去宰殺，又像羊在剪毛的人手下無聲，他也是這樣不開口。 \end{tabularx} \\ \\ \relax
53:8 & \begin{tabularx}{0.7\textwidth}{X} 因受欺壓和審判，他被奪去，誰能想到他的世代呢？因為他從活人之地被剪除，為我百姓的罪過他被帶到死裡。 \end{tabularx} \\ \\ \relax
53:9 & \begin{tabularx}{0.7\textwidth}{X} 他雖然未行殘暴，口中也沒有詭詐，人還使他與惡人同穴，與財主同墓。 \end{tabularx} \\ \\ \relax
53:10 & \begin{tabularx}{0.7\textwidth}{X} 耶和華的旨意要壓傷他，使他受苦。當他的生命作為贖罪祭時，他必看見後裔，他的年日必然長久。耶和華所喜悅的事，必在他手中亨通。 \end{tabularx} \\ \\ \relax
53:11 & \begin{tabularx}{0.7\textwidth}{X} 因自己的勞苦，他必看見光就心滿意足。因自己的認識，我的義僕使許多人得稱為義，他要擔當他們的罪孽。 \end{tabularx} \\ \\ \relax
53:12 & \begin{tabularx}{0.7\textwidth}{X} 因此，我要使他與位大的同份，與強盛的均分擄物。因為他傾倒自己的生命，以致於死，也列在罪犯之中。他卻擔當多人的罪，為他們的過犯代求。 \end{tabularx} \\ \\
[1ex]
\hline
\hline
\end{longtable}
講題:人生或苦 不亢不卑　Suffering and the Struggle for Dignity

講員:林豪恩先生

經文:以賽亞書53:3-12

場合:香港中文大學崇基學院禮拜堂主日崇拜

日期:2021年10月17日

甲. 引言

很多人都相信並且盼望「好人有好報,壞人有惡報」,你也是其中之一嗎?如何得知「善有善報,惡有惡報」呢?其實,我們比較容易見得到的,是「善報」和「惡報」,至於當事人是「善」或是「惡」則不容易看到.於是乎,人們傾向從逆向判斷,即是說,有「惡報」的應該是「壞人」,有「善報」的應該是「好人」.最常見的「祖母社會教育」是這樣的:「乖孫,你看到這個被警察拘捕的壞人嗎?你千萬不要學他啊!」「麻麻,為甚麼說這個人是壞人呢?」「乖孫,這個人被拘捕,一定是做了壞事.你明白嗎?不明白也不重要,最重要的是不要學他.」於是,乖孫自小就被教導從看見的推論看不見的:如果有人被懲罰,一定是這個人犯了事.再把這概念推廣應用於人生際遇,大概就是約伯的朋友對遭逢災禍的約伯的詰問:「請你追想:無辜的人有誰滅亡?正直的人何處被剪除?」( 約伯記4:7) 視不好的狀況為罪的懲罰,連耶穌的門徒也是這樣想的:「耶穌往前走的時候,看見一個生來就失明的人.門徒問耶穌:「拉比,這人生來失明,是誰犯了罪?是這人還是他的父母呢?」(約翰福音9:1-2)見到人有「好報」,就認為他是好人；見到人有「惡報」,就肯定他是壞人.有時,我們也會這樣來理解自己的際遇:當自己順風順水,春風得意,人見人愛,車見車載,無病無痛,精神爽利的時候,就以為自己正在被「好報」；當自己諸事不順,困難重重,身體不適,精神不振,不受歡迎,處處碰壁的時候,就感到自己被「惡報」,免不了心感不安,自我懷疑一番.

今天的舊約經課以賽亞書提到有一個人,從他的遭遇來看,按我們的邏輯來說,肯定是被「惡報」.然而,以賽亞先知的看法似乎與我們不同,因此,讓我們來細察一下.

乙.經文及內容

一.遭逢苦難的人

這個人的遭遇和境況是這樣的:

他被藐視,被人厭棄；多受痛苦,常經憂患.他被藐視,好像被人掩面不看的一樣,人們也不尊重他.(53:3)

他擔當憂患,背負痛苦；受責罰,被擊打苦待.(53:4)

他受害,被壓傷.受懲罰,受鞭傷(53:5)

他被欺壓受苦,他被牽去宰殺.(53:7)

這個人的遭遇,可以稱得上是劫難,或者有旁觀者理解為「惡報」.若我們從旁觀者的角度轉為代為當事人的角色,我們可能感到命途多舛,也感到身體的痛,以及心靈的苦.遭遇劫難,陷入苦境的人,會如何看自己呢?又會如何反應呢?一般人可能有以下兩種取向:

二.怨天尤人,憤世嫉俗

其中一種反應是初而怨天尤人,繼而憤世嫉俗.是誰令我如此不幸和痛苦的呢?是別人的錯,是社會的錯,是天地沒有公理.因此,我有權憤怒,正如先知約拿所說:「我發怒以至於死,都合乎理!」社會不公平,人們排斥我,令我利益受損,令我生意失敗,令我家破人亡,令我苦不堪言,我是不是要讓社會知道他們做得太過份呢?是不是要社會知道如此對待我的後果呢?是不是要以行動逼社會償還欠我的一切呢?

某月某日,在某省某市一名6旬獨居退休老人,疑因不滿遭受社會遺棄,持棍在校園外襲擊小學生洩憤,導致12名兒童和2路人受傷.

又某月某日,在某大城市的某著名小學附近,畢業慶典之後,有人揮刀亂砍,導致2名男童死亡,1男童和1女性受傷.砍人者來自某省某市,畢業於該省科技大學,據悉來這大城市找工作無著,心生不忿,懷藏菜刀,行兇報復.

三.自怨自艾,自卑自憐

另一種反應可能是自怨自艾,自卑自憐.一般人都認為:好人好者,怎會被苦待呢?行得正企得正,怎會被刑罰呢?如果你沒有問題,怎會被厭棄呢?那麼多人都沒有事,只有你一個有事,一定是你有問題啦.以賽亞書53章的群眾對著那位不幸受苦的人也是這樣判斷的:「我們卻以為他受責罰,被上帝擊打苦待了」(53:4)那個人如此糟糕,肯定犯了天理不容的罪惡,連老天爺都看不過眼.用兩個字來總結:活該(抵死).

接連的打擊,當事人可能也會自我懷疑:那麼多人都成功,為何只有我一個失敗呢?其他人都快樂,為何只有我獨自痛苦憂悶?為甚麼其他人都順順利利,只有我那麼多困難呢?世界上那麼多人,為甚麼選中我來欺壓和苦待呢?這些惡劣的遭遇臨到我,會不會是我真的做錯甚麼呢?我陷入如此的困境,會不會是我一直都行錯路呢?會不會是我前世罪孽深重,今世要被刑罰?我無能為力逃離這困境,我的存在成為人們的笑話,我繼續生存下去還有甚麼價值呢?

四.受苦,可以是為他人

以賽亞先知指出,當時的人對著那位遭遇災難的受苦者是這樣判斷的:

53:4 卻以為他受責罰,是被神擊打苦待.

不過,遭遇災難是不是等於被懲罰呢?就算接受懲罰,是不是由於做錯事呢?一般來說,或許是；但對這個例子來說,以賽亞先知清楚地宣告:不是.

53:4他誠然擔當我們的憂患,背負我們的痛苦；

原來,有一種苦,是為他人而受的苦.

53:5他為我們的過犯受害,為我們的罪孽被壓傷.

原來,有一種受罰,並不是因為自己犯錯,反而是因為他人犯錯.

53:5 因他受的懲罰,我們得平安；因他受的鞭傷,我們得醫治.

原來,有些情況,接受某種打擊,承受某種懲罰,是會為他人帶來平安和醫治的.

五.公平與受苦

或許有人說:不公平,真是不公平.他們的過犯,為甚麼不是他們自己受害,而要我來受害?為甚麼不是他們被自己的罪孽壓傷自己,而是壓傷我呢?他們的憂患,為甚麼不是他們自己擔當,而是要我來擔當呢?他們的痛苦,為甚麼不是他們自己背負,而是要我來背負呢?

或許,這是不公平的.你知道如此不公平的情況為甚麼會出現嗎?原因是他們自己不能:

53:6我們都如羊走迷,各人偏行己路；

走迷了路的羊自己不能回到羊園,偏行己路的人自己沒能力重回正軌.

今年9月28日,土耳其50歲男子穆特魯跟朋友一直喝酒喝到醉,隨後走失遊蕩晃進樹林.穆特魯的妻子好幾個小時都聯繫不上人,就跟友人決定報警.救援人員與義工進入樹林搜索,搜索過程中也有當地居民加入.好幾個小時過去,天快光了,救援隊也找不到迷路者.直到一名救援人員喊出「穆特魯」的名字,在搜索隊中有人回應「我在這裡」.原來酒醉了的穆特魯在樹林漫遊的時候,加入了搜尋隊,他不知道人們尋找的正是他自己,不知不覺地參加了這項自己找自己的烏龍任務.

穆特魯醉酒而勞師動眾,最後要不要補償呢?新聞沒有跟進.不過,無論要補償或者不需要補償,失而復得,平安回家團聚,都是大團圓結局.如羊走迷,偏行己路,可以令當事人身陷險境,不但當事人受傷害,關心他的人都都會受苦.因此,有人「如羊走迷,偏行己路」而不能自救時,總要有人回應的.

這兩個星期颱風襲港,令人想起有些人喜歡在打風時出海尋求刺激.有救護員說:「每次打風期間,從新聞報道看到有人在風球下出海,「成班師兄弟都嬲都震」,希望市民可以諒解救護員出海救人,等於「半隻腳踏入鬼門關」.」

縱使憤怒,如果真的有出海的人遇險,大家真的認為不用拯救嗎?或者有電視機旁邊的觀眾說得出「由他們死吧」這樣的話,但如果遇險者是家人或朋友,反應就不同了.就算是「半隻腳踏入鬼門關」,仍是要出發救人的.事實上,也的確有救援人員在各種拯救行動中失去生命.

六.人們的判斷可能改不了,但最重要的是知道自己沒有做錯甚麼

救人而犧牲的,人們會肯定和嘉許.但不是所有為人承受痛苦的,都會被認同和感激.對於以賽亞書53章描述的那位為「如羊走迷,偏行己路」的人的受苦者,那些「如羊走迷,偏行己路」的人是如此回報他的:

53:7他像羔羊被牽去宰殺,又像羊在剪毛的人手下無聲,

53:8因受欺壓和審判,他被奪去,從活人之地被剪除,被帶到死裡.

有些時候,人們的判斷可能改不了,但重要的是知道自己沒有做錯甚麼.

53:9他未行殘暴,口中也沒有詭詐

能夠如此自己我肯定,才能心安理得,坦然無所愧咎.

七.更重要的是知道自己在做甚麼

53:6耶和華使我們眾人的罪孽都歸在他身上.

53:10耶和華的旨意要壓傷他,使他受苦.當他的生命作為贖罪祭,

在逆境之中,在受苦的時候,一方面,知道自己沒有做錯,另一方面,也知道是為了他人的幸福,當事人才能夠避開憤世嫉俗或自卑自憐的兩極,

53:7他被欺壓受苦,卻不開口；

默然承受,就好像在洪水中被驚濤駭浪衝擊的方舟,無論雨怎麼打,始於盛載著方舟內的眾生,保護著大地的苗蕊.

八.世界並不終結於受苦

53:11因自己的勞苦,他必看見光就心滿意足.因自己的認識,我的義僕使許多人得稱為義,他要擔當他們的罪孽.

53:12因此,我要使他與位大的同份,與強盛的均分擄物.因為他傾倒自己的生命,以致於死,也列在罪犯之中.他卻擔當多人的罪,為他們的過犯代求.

苦並不是終結,可能是人生百味之一.對一般人來說,不會是「以致於死」,可能只是某種放棄,某程度的犧牲,或者某方面的限制.我有一位小學同學,工作地方人人吸煙,他也是煙不離手.同學們約飯聚時,他每十五分鐘便要離開一會.有一次飯聚,有其他同學離開一會時,他仍然留在飯局.我們好奇問原因,他說已經不吸煙了.故事是這樣的:當他的孩子出生後,他決定不在家吸煙,相信大家都知道原因.當他的孩子懂事後,告訴他身上有煙味,於是他決定戒煙,原因也是明顯的.改變一種嗜好,尤其是會上癮的嗜好,不是容易的事.這位爸爸可能也掙扎過:為甚麼要限制我的自由?為甚麼要放棄我的喜好?戒煙期間煙癮發作的時候,也可能會問:為甚麼我要令自己如此辛苦?答案?不用多說,相信大家是明白的.

有些受苦的時候,我們覺得完全是不幸,是沒有意義的,也看不見對其他人有甚麼幫助,反而是感到連累了其他人.曾經有一位醫生來學校分享,他少年時患骨癌,最後要切割一條腳保命.過程所受的身心痛苦及手術後的適應艱辛可想而之.後來,他當上醫生,在公立醫院服務.有一次,另一個病房有一位病人腳部肌肉被病茵入侵,需要切足保命.那位病人不能接受失去一隻腳,寧死也不接受手術.病房的醫生和護士都說服不了他.這位醫生知道後,他去探望那位病人,了解他的情況,給他看他也是沒有一隻腳,向他分享自己切足的故事.這次探訪之後,那位病人告訴主診醫生,他願意接受手術.

丙.總結

世間有苦難,但苦難不是結局；人生有困境,但困境不是終結.以賽亞書這位受苦者所受的苦不是徒然的,他承受痛苦,為世間開拓可能性,為人們帶來更新和盼望.

清代詩人龔自珍的其中一首詩如此寫道:

浩蕩離愁白日斜,吟鞭東指即天涯.

落紅不是無情物,化作春泥更護花.

好一句「落紅不是無情物,化作春泥更護花」.既然如此,就不用憤世,也不需自憐.

書信經課希伯來書

5:1凡從人間挑選的大祭司,是奉派替人辦理屬上帝的事,為要獻上禮物和贖罪祭.

5:2他能體諒那愚蒙的和失迷的人,因為他自己也是被軟弱所困.

這大祭司是誰呢?

希伯來書

4:14我們既然有一位已經升入高天尊榮的大祭司,就是上帝的兒子耶穌,便當持定所承認的道.

4:15因我們的大祭司並非不能體恤我們的軟弱.他也曾凡事受過試探,與我們一樣,只是他沒有犯罪.

4:16所以,我們只管坦然無懼地來到施恩的寶座前,為要得憐恤,蒙恩惠,作隨時的幫助.



\newpage

\section{成聖的因與果}
\label{sec:573}
\textbf{蘇成溢牧師 - 成聖的因與果}
\newline
\newline
\hyperref[sec:572]{< < < PREV SERMON < < <}
~
\hyperref[sec:toc]{[返目錄]}
~
\hyperref[sec:574]{> > > NEXT SERMON > > >}


希伯來書 9:11-14
\newline
\begin{longtable}{cl}
\hline
\hline
章節 & 經文 (和合本修訂版)\\
\hline
9:11 & \begin{tabularx}{0.7\textwidth}{X} 但現在基督已經來到，作了已實現的美事的大祭司，經過那更大更全備的帳幕，不是人手所造，也不是屬於這世界的； \end{tabularx} \\ \\ \relax
9:12 & \begin{tabularx}{0.7\textwidth}{X} 他不用山羊和牛犢的血，而是用自己的血，只一次進入至聖所就獲得了永遠的贖罪。 \end{tabularx} \\ \\ \relax
9:13 & \begin{tabularx}{0.7\textwidth}{X} 若山羊和公牛的血，以及母牛犢的灰，灑在不潔的人身上，尚且使人成聖，身體潔淨， \end{tabularx} \\ \\ \relax
9:14 & \begin{tabularx}{0.7\textwidth}{X} 何況基督的血，他藉著永遠的靈把自己無瑕疵地獻給神，更能洗淨我們的良心，除去致死的行為，好事奉那位永生的神。 \end{tabularx} \\ \\
[1ex]
\hline
\hline
\end{longtable}
講題:成聖的因與果　To sanctify and to be sanctified

經文:希伯來書9章11-14節

講員:蘇成溢牧師

場合:香港中文大學崇基學院禮拜堂主日崇拜

日期:2021年10月31日

1 引言

1.1 希伯來書的背景和思想重點

希伯來書大概成書於主後70年前後,是一封寫給陷於困境中基督徒的信.當時他們受到反對者越來越大的壓迫,有放棄信仰的危機.作者要堅定他們的信心,向他們指出,耶穌基督是上帝真確和完整的啟示,要認識以下三方面對耶穌基督的信仰重點:1)耶穌是上帝的兒子,祂受苦順服,但祂是聖子,比一切的先知和天使,甚至比摩西更高.2)上帝立耶穌作為永遠的祭司,高過一切的祭司.3)藉著耶穌,人能從罪惡,恐懼,死亡中得拯救；因祂是永遠的大祭司,更獻上自己,超越了一切猶太的禮儀和祭司所獻上的祭牲畜,永遠有效.

1.2 因此,凡在基督耶穌裡的人,因有基督大祭司所作的赦罪工作(獻上自己為贖罪祭),生命便能完全得潔淨,更新,成聖.

反省:各位基督徒,昔日你接受了洗禮,被稱為基督徒,請問,今日的你,有耶穌在你生命裡頭嗎?你又在耶穌基督裡面嗎?

希伯來書9章11-14節 (有關贖罪的律例,參考利未記16章)

11-12節 大祭司為百姓贖罪,必須帶血進入至聖所

(「贖」是個商業用語,有如今天的擔保人/抵押品.)根據利未記25: 23-28 的說法,若有以色列人沒有錢,他可以把他的產業抵押,到禧年前,無需贖回,產業便物歸原主.但他也可以隨時贖回的.方法是找一個富有的親戚代他付錢,就可贖回；又或當自己儲蓄夠錢,自行贖回.至於那些無地無產業的窮人,可以賣身當僱工服侍主人.若要贖身,可按利未記25: 39-45 所講述的方法進行:1)由親戚付款代贖回；2)自行付款為自己贖身以回復自由.

以上所講述的,是商業買賣形式的抵押/贖回行為.然而,人類因罪和惡行所要承受的代價,根不能自行以金錢代替,為自己買贖.希伯來書作者告訴我們,是那位被上帝所立的永遠的大祭司,祂自己成為贖罪祭,獻上自己的流的血,一次過成全了為全人類的救贖.(有如一個人/機構/公司欠了借貸款,有人為他/他們一次過付款,所有欠債一筆勾銷.)例子:百多年前在福建福州衛理公會第一個信徒,本是一個被通緝的罪犯.一天他因要逃避追查,躲進一間教堂.當時他聽到台上牧師在講道,講述耶穌釘十字架的情況.牧師說,在耶穌釘十架時,在他的左右兩邊分別有兩位強盜,最終一位接受耶穌得拯救,另一位未有接受.那份在教堂後聽到的通緝犯,受聖靈感動問自己,我既是個有罪的人,那麼我是那位接受耶穌還是不接受耶穌的罪犯呢?

13節 山羊的血是為百姓贖罪(利未記16: 15),公牛的血是為自己和家族贖罪(利未記16:11),母牛犢的灰用作為人去除污穢用之調和的水(民數記19:2-10).

這些都是舊約律例用作潔淨罪污,使人成為聖潔的禮儀工具.關鍵之處並不在於這些禮儀和工具,乃在於生命的罪污能否得潔淨,有罪的人可否成為聖潔?

14節 基督的血(聖子),藉永遠的靈(聖靈),獻給上帝(聖父):救贖世人是三一上帝的作為,使人心潔淨,惡行去除,並能事奉上帝.

基督自己成為贖價,洗淨了我們身心一切的污穢,罪惡,使我們成聖.舊約的成聖是用血(利未記8: 30)；使以色列民成為了聖潔的子民(利未記19: 2；20: 26).新約的成聖是因基督先成全了救贖,用真理的道(約翰福音17: 17-19)使人在基督裡(哥林多前書1: 2),藉聖靈(哥林多前書6: 11)和以水和道(以弗所書5: 26)去使生命成聖.

反省:基督徒的生命從滿身罪污到成為聖潔,全靠耶穌基督一次,完全的獻上自己所成全的救恩,將我們的生命徹底地轉換成為新的生命.今天的你和我,已經不再是舊的我,這是何等奇妙的事.

3 總結和反省

3.1 因擁有聖潔的身份,我們應曉得感恩.「從我感恩因為我快樂,滿足,改為我快樂因為我感恩.」

3.2 成聖之人的聖潔身份,不是為了自己的榮譽,高於別人,乃是為事奉上帝,祝福他人.(有如食肆的餐具已完全清潔消毒,可以放心承載食物,供人使用.)



\newpage

\section{活在山崩石裂時}
\label{sec:574}
\textbf{林豪恩先生 - 活在山崩石裂時}
\newline
\newline
\hyperref[sec:573]{< < < PREV SERMON < < <}
~
\hyperref[sec:toc]{[返目錄]}
~
\hyperref[sec:575]{> > > NEXT SERMON > > >}


馬太福音 13:1-10
\newline
\begin{longtable}{cl}
\hline
\hline
章節 & 經文 (和合本修訂版)\\
\hline
13:1 & \begin{tabularx}{0.7\textwidth}{X} 就在那天，耶穌從房子裡出來，坐在海邊。 \end{tabularx} \\ \\ \relax
13:2 & \begin{tabularx}{0.7\textwidth}{X} 有一大群人到他那裡聚集，他只好上船坐下，眾人都站在岸上。 \end{tabularx} \\ \\ \relax
13:3 & \begin{tabularx}{0.7\textwidth}{X} 他用比喻對他們講了許多話。他說：「有一個撒種的出去撒種。 \end{tabularx} \\ \\ \relax
13:4 & \begin{tabularx}{0.7\textwidth}{X} 他撒的時候，有的落在路旁，飛鳥來把它們吃掉了。 \end{tabularx} \\ \\ \relax
13:5 & \begin{tabularx}{0.7\textwidth}{X} 有的落在土淺的石頭地上，因為土不深，很快就長出苗來， \end{tabularx} \\ \\ \relax
13:6 & \begin{tabularx}{0.7\textwidth}{X} 太陽出來一曬，因為沒有根就枯乾了。 \end{tabularx} \\ \\ \relax
13:7 & \begin{tabularx}{0.7\textwidth}{X} 有的落在荊棘裡，荊棘長起來，把它擠住了。 \end{tabularx} \\ \\ \relax
13:8 & \begin{tabularx}{0.7\textwidth}{X} 又有的落在好土裡，就結出果實，有一百倍的，有六十倍的，有三十倍的。 \end{tabularx} \\ \\ \relax
13:9 & \begin{tabularx}{0.7\textwidth}{X} 有耳的，就應當聽！」 \end{tabularx} \\ \\ \relax
13:10 & \begin{tabularx}{0.7\textwidth}{X} 門徒進前來問耶穌：「對眾人講話，為甚麼用比喻呢？」 \end{tabularx} \\ \\
[1ex]
\hline
\hline
\end{longtable}
講題:活在山崩石裂時　When the Great Buildings Fall

講員:林豪恩先生

經文:馬太福音13:1-10

場合:香港中文大學崇基學院禮拜堂主日崇拜

日期:2021年11月14日

引言

座落於美國紐約的世界貿易中心由6座大樓組成,自1973年啟用以來,每天約五萬人在裡面工作,並有超過20萬名遊客出入.為甚麼有那麼多遊客出入呢?因為6座大樓中有兩座高達110層,分別稱為「北塔」和「南塔」,曾經是世界上最高的建築物,成為紐約著名地標之一,遊客可以在觀光層眺望,欣賞紐約市的美景.世貿中心能夠在紐約市拔地而起殊不簡單,建造這些摩天大樓的想法最早在1943年就有人提出,1962年開始設計,1966年動工,1973年落成啟用,是很多人多年投入的成果.無論是建築師,建築工人,在裡面上班的人,絡繹不絕的遊人都不會想到,這宏偉壯觀,高聳入雲的雙塔,竟會在瞬刻之間倒塌.2001年9月11日上午,一架民航機撞入北樓的北面.17分鐘後,另一架民航機撞入南樓.在燃燒約56分鐘後,南樓倒塌.北樓在燃燒了約102分鐘後也倒塌下來.雙塔被撞毀倒塌的畫面,不單震撼了美國,也震撼了整個世界.事隔二十年,相信不少目擊者回想起來,仍然歷歷在目,仍然感到震撼.讓我們借用這份震撼感來閱讀今天的福音經課.

大殿宇也有倒下時

13:1耶穌從殿裡出來的時候,有一個門徒對他說:「夫子,請看,這是何等的石頭!何等的殿宇!」

13:2耶穌對他說:「你看見這大殿宇嗎?將來在這裡沒有一塊石頭留在石頭上,不被拆毀了.」

耶路撒冷聖殿被拆毀得「沒有一塊石頭留在石頭上」,猶太人聽到會引起很大的反應.那種反應不但是當下的震撼,也是歷史的傷痕.對猶太人來說,耶路撒冷聖殿是甚麼呢?在他們心目中有甚麼位置呢?

根據聖經記載,第一座耶路撒冷聖殿在大衛王的時代已經開始預備材料,由繼位的所羅門王動工興建.因此,當時的聖殿又稱所羅門聖殿.聖殿落成後,隨即成為朝見耶和華上帝的敬拜場所,也成了以色列民族的象徵.我們可以從詩篇122篇感受到他們對耶路撒冷聖殿的感情:

122:1人對我說:我們往耶和華的殿去,我就歡喜.

122:2耶路撒冷啊,我們的腳站在你的門內.

122:4眾支派,就是耶和華的支派,上那裡去,按以色列的常例稱讚耶和華的名.

122:5因為在那裡設立審判的寶座,就是大衛家的寶座.

122:9因耶和華－我們上帝殿的緣故,我要為你求福!

公元前586年,巴比倫王國尼布甲尼撒二世摧毀耶路撒冷,聖殿被毀.猶太人為此而感到非常哀痛和羞辱.

那都是幾百年前的歷史了.耶穌和門徒當時所見到的聖殿,是失而復得的另一座.

公元前539年,波斯帝國推翻了巴比倫王國.波斯王古列(Cyrus the Great)頒旨,讓猶太人在公元前537年回耶路撒冷重建聖殿.按以斯拉尼希米記所述,重建過程非常艱辛,攔阻處處,幾經辛苦,第二座聖殿於公元前515年完工.又幾百年後,在公元前37年,聖殿在大希律王的時代得以擴建.

門徒對耶穌說:「夫子,請看,這是何等的石頭!何等的殿宇!」門徒對聖殿的讚譽和崇敬,很可能也代表著當時普遍人對聖殿的敬意.

耶穌似乎沒有分享這份對聖殿的敬意,更語不驚人誓不休地撥門徒冷水:「你看見這大殿宇嗎?將來在這裡沒有一塊石頭留在石頭上,不被拆毀了.」

耶穌如此說,很可能傷害了普遍猶太人的感情,也令猶太領袖們反感,因為聖殿不單單是一座建築物,更承載著宗教和民族感情,也代表著權力和領導.如果可以撇開感性,純理性來看,人們都知道如何宏偉的建築物也會有倒下的一天.就如我們現在身處的這個禮拜堂,總有一天會拆卸的,對嗎?又例如你現在所住的大廈,就算一百年內不拆卸,二百年後也不會存在吧!因此,問題不是建築物最後會不會被拆毀,而是在甚麼時候和甚麼狀況之下拆毀.正如人人都知道人總有一死,但在甚麼時候及怎樣的狀況之下死去,那才是問題的所在.當時,門徒也懷著類似的疑問.

亂世多迷惑

13:3耶穌在橄欖山上對聖殿而坐.彼得,雅各,約翰,和安得烈暗暗地問他說:

13:4「請告訴我們,甚麼時候有這些事呢?這一切事將成的時候有甚麼預兆呢?」

13:5耶穌說:「你們要謹慎,免得有人迷惑你們.

再一次,門徒的反應代表著普遍人的心聲.如果事情總會發生,人們當然想洞悉先機,趨吉避凶.如果事先知道紐約世貿中心雙子塔甚麼時候倒塌,人們當然不會照常去上班,遊客也不會去遊覽.現實上,就是不知道.就算有人說知道,也不會有多少人相信.就正如新冠病毒兩年前在人類之間傳播,造成全球疫情,人們只能視之為突如其來.這些突如其來的事情,是不是完全在人類的預計之外呢?又不是.從歷史中,人類知道瘟疫總會有的,只是無論科技如何進步,也沒有辦法知道瘟疫在甚麼時候及甚麼情況之下來襲.

就算不知道具體時間,有蛛絲馬跡可以推敲嗎?這也是門徒想知道的:「有甚麼預兆呢?」有的,但通常都是事情發生之後,重組案情的時候才可能組織到出來的.事後孔明,當然人人都可以變成專家.

耶穌沒有回答門徒關於時間和預兆的問題,反而吩咐門徒「你們要謹慎,免得有人迷惑你們.」當人人都說自己有一個預兆的時候,當人人都說自己掌握到事情的來龍去脈的時候,者當人人都說有內幕消息的時候,正是反映出人心最徬徨無助,社會最混亂不安,也是迷惑人者最有利去迷惑人的時候.如果過去兩三年你在這個城市生活,你一定接收過海量這些資訊,可能你自己也發放過這些資訊,甚至也因這些資訊而有所行動.亂世中的迷惑是甚麼呢?耶穌提出了兩種可能性,正好是人們常常陷入的兩種極端,其一是藥石亂投,焦慮地四處找救生圈；另一個極端是因強烈的無力感而投降,苟且偷生.

亂世迷惑的一端:拯救來了

13:6將來有好些人冒我的名來,說:『我是基督』,並且要迷惑許多人.

『我是基督』是甚麼意思呢?基督即是當時人們盼望的彌賽亞,即是上帝差派來的拯救者.在人們的盼望中,彌賽亞有三種身分,分別是祭司,先知和君王,只要彌賽亞來臨,生活上,宗教上及政治上的困境就得到解決,理想的國度就會出現.

有資料說,古代猶太人記戴有29人自稱是基督,近百年來有1100人自稱基督.

其中一個例子是在公元132年,有一位擁有武裝力量的猶太人將領被稱為彌賽亞君王.許多猶太人以為,這個人就是拯救他們脫離羅馬帝國,終止他們長期以來所受壓迫的人.可是,這位被封為彌賽亞君王的將領在戰鬥中一敗塗地,成千上萬的猶太人因此賠上了性命.

自古至今,人們都盼望有一個人來臨,解決所有的困境；從東至西,人們都期待有一種方法出現(社會層面稱之為制度),消除所有的痛苦.生活在二十一世紀的我們,何嘗不是如此盼望呢?

每一次有人出來宣稱他是真命天子,為民請命,帶領我們往理想的世界,我們都雙眼發光,寄以厚望,把我們僅有的雙手奉上.一次又一次,我們逃不過治亂興衰這歷史循環.在這循環中,多少性命被犧牲?多少聖靈被塗炭?多少家園被摧毀?多少痛苦白白受呢?

我們並不是在譴責某些別有用心的人,只是怪責自己既簡單又無知,一次又一次地寄予盼望,一次又一次地被欺騙.有人被網上戀人騙取家財,是騙子太邪惡嗎?當然是.但另一方面,也可能反映出我們對感情的渴求而已.有人被假醫生騙財,病人也騙,是騙子太沒良心嗎?當然是.但另一方面,也可能反映出我們對健康的渴求而已.我們還有很多不同的渴求,難度都是不應該的嗎?當然不是,還存活的人就有渴求,只是死人才沒有渴求.

我們的渴求是繼續存在的,自稱基督的也是繼續出現的,在這樣的情況中,耶穌並不是要求門徒清除渴求,也不是消滅自稱基督的,而是吩咐門徒在這樣的狀況中要謹慎,免得被迷惑.

4.亂世迷惑的另一端:世界末日了

13:7你們聽見打仗和打仗的風聲,不要驚慌.這些事是必須有的,只是末期還沒有到.

13:8民要攻打民,國要攻打國；多處必有地震,饑荒.

亂世迷惑的一端是自認拯救者群魔亂舞,亂世迷惑的另一端是人們無助無力,苟且偷生,束手等待那日來到.既然「民要攻打民,國要攻打國」,我們這些蟻民豈能螳臂擋車呢?既然「多處必有地震,饑荒」,我們這些渺小如微塵的人類豈能對抗天地風雲呢!

在教會歷史中,有多少傳道者都借用這兩節經文來宣揚末日信息,警告人們時日無多,要信耶穌就立即舉手.他們圖文並茂,影音並用,繪聲繪影地指出,世界上的戰爭越來越多,地震次數也越來越頻密,證明末日也越來越近.既然末日越來越近,世界將會過去,人生就沒有甚麼再需要執著,世界也沒有甚麼再值得投入了.最極端的版本是,變賣財產,放下工作,一齊去指定地方,等候被接升天.

5.是令你受苦的災難,還是令你期待的盼望?

這都是災難(災難:原文是生產之難)的起頭. (現中:這些事的發生正像產婦陣痛的開始一樣)

災難原文是「生產之痛苦」,比喻極大的痛苦,好像生產一樣.你知道你媽媽生你有多痛嗎嗎?婦產科專科醫生表示用數字分級法,0為不痛,10為最痛,7-10級為重度疼痛,人們說生仔是10級痛.你能夠忍受到甚麼程度的痛呢?在某地有一間醫院與一間中學合辦「分娩疼痛模擬體驗」活動.在醫護人員陪同下,學生在腹部貼上電極片,體驗儀會輸出電磁波,模擬陣痛過程.痛楚的級數從1到10級,大家猜一猜學生大多能忍受到甚麼級別的痛楚呢?有些學生起初表情輕鬆,當醫護人員接駁儀器通電,該學生立即露出痛苦的表情,輾轉反側大叫起來,更哭起來,在旁醫護人員即時拍拍學生的肩膀安撫及遞上紙巾.有的學生就躺在梳化上,體驗時面容扭曲,更典床典蓆及手腳亂舞.結果顯示,大多數學生只是忍受到1至3級的痛楚程度而已.

我們好像越講越遠,差些忘記了本來的主題,門徒和耶穌本來是在談論耶路撒冷聖殿的.我們曾經提及,聖殿是一座建築物,但它不單單是一座建築物,它承載著民族的歷史和人民的感情；它代表著宗教和政治制度.聖殿被拆毀,意味著民族歷史和人民感情被傷害,宗教和政治制度被推反和改變.這必然是混亂的時候,也是痛苦的時候.耶穌提醒門徒,在混亂和痛苦的時候,不要被迷惑,不要藥石亂投,也不要坐以待斃.我們有沒有發現耶穌用「生產之痛苦」來比喻這種日子的用意呢?是的,「生產的痛苦」是真實的,痛的程度是非常嚴重的,但在生產的時候,沒有人懷著坐以待斃的態度的,相反,所有人都懷著喜悅的盼望,因為將有嬰孩要出生,人們都在迎接新生命的來臨.如果我們用這個角度來聆聽耶穌和門徒這段對話,我們就未必用來進行未日恐嚇,反而可以用來揭露在混亂和不安中的盼望.

前瞻:

活在山崩石裂時,如果不藥石亂投,盲目地向自稱拯救者投誠；又不坐以待斃,還可以怎樣呢?我們可以參考耶穌接著向門徒說的話:

13:9但你們要謹慎；因為人要把你們交給公會, 並且你們在會堂裡要受鞭打,又為我的緣故站在諸侯與君王面前,對他們作見證.

13:10然而,福音必須先傳給萬民.

用使徒保羅的說法,無論得時不得時,都要作見證,無論對甚麼人,都要作見證.因為耶穌基督的心懷,掛念著萬民是否得到福音,這也應該是跟隨耶穌基督的人的關注.



\newpage

\section{你是猶太人的王嗎}
\label{sec:575}
\textbf{盧龍光牧師 - 你是猶太人的王嗎?}
\newline
\newline
\hyperref[sec:574]{< < < PREV SERMON < < <}
~
\hyperref[sec:toc]{[返目錄]}
~
\hyperref[sec:576]{> > > NEXT SERMON > > >}


約翰福音 18:33-37
\newline
\begin{longtable}{cl}
\hline
\hline
章節 & 經文 (和合本修訂版)\\
\hline
18:33 & \begin{tabularx}{0.7\textwidth}{X} 於是彼拉多又進了總督府，叫耶穌來，對他說：「你是猶太人的王嗎？」 \end{tabularx} \\ \\ \relax
18:34 & \begin{tabularx}{0.7\textwidth}{X} 耶穌回答：「這話是你說的，還是別人論到我時對你說的呢？」 \end{tabularx} \\ \\ \relax
18:35 & \begin{tabularx}{0.7\textwidth}{X} 彼拉多回答：「難道我是猶太人嗎？你的同胞和祭司長把你交給我。你做了甚麼事呢？」 \end{tabularx} \\ \\ \relax
18:36 & \begin{tabularx}{0.7\textwidth}{X} 耶穌回答：「我的國不屬於這世界；我的國若屬於這世界，我的部下就會為我戰鬥，使我不至於被交給猶太人。只是我的國不屬於這世界。」 \end{tabularx} \\ \\ \relax
18:37 & \begin{tabularx}{0.7\textwidth}{X} 於是彼拉多對他說：「那麼，你是王了？」耶穌回答：「是你說我是王。我為此而生，也為此來到世界，為了給真理作見證。凡屬真理的人都聽我的話。」 \end{tabularx} \\ \\
[1ex]
\hline
\hline
\end{longtable}
講題:你是猶太人的王嗎?　Are you the King of the Jews?

經文:約翰福音18:33-37

講員:盧龍光牧師

場合:香港中文大學崇基學院禮拜堂主日崇拜

日期:2021年11月21日

今天這個主日是將臨期前的最後一個主日,下個主日,我們將進入迎接耶穌來臨的將臨期,就是教會年曆新一年的開始,信徒們要透過閱讀聖經去自我省悟,警醒,悔改,以迎接喜樂與平安.

迎接耶穌來臨,是慶祝和紀念耶穌在二千多年前降生在人間；另一方面是警醒,準備好迎接耶穌的第二次的忽然來臨.

然而,這一切的核心問題,就是耶穌究竟是誰?他過去為什麼要來到世上?在世上做了什麼?影響了什麼?他是誰?他的降生值得我們慶祝嗎?另一方面,他所宣告的忽然再來是可信的嗎?歷史會真的斷裂嗎?

今天的主日,教會的福音經課是約翰福音18章33-37節,內容是關於耶穌在被捕之後,被押解到羅馬總督彼拉多面前受審的對話；首先,彼拉多問耶穌:你是猶太人的王嗎?也就是在質問:耶穌,你是誰?

彼拉多的問題令我們追問:

為什麼他會這樣問呢?

耶穌做了什麼事情會令人控訴他呢?

耶穌如何回答這問題?

彼拉多這樣問,明顯地是因為他知道猶太人對耶穌的控訴:耶穌是叛亂分子,在羅馬的統治下宣稱自己是猶太人的王,然而彼拉多並不相信,認為這個控訴沒有根據,這只是猶太人的領袖想利用他除去他們自己猶太人中的異己.在這裡我們見到一個吊詭的情形,猶太人本應不甘心被羅馬人統治,羅馬人是他們的敵人,應該去尋找他們的領袖,擁戴為王,去推翻羅馬人的統治；而羅馬的統治者彼拉多,應該敏感於猶太人的陰謀,將所有推翻羅馬統治的陰謀在萌牙的時候就鎮壓下來；但是在這裡,是猶太人的領袖主動向羅馬的統治者控訴有可能做王的猶太人,羅馬的統治者卻不相信,打算釋放這個被猶太人控訴為叛亂分子的人!為什麼會是這樣呢?因為其實猶太人自己並不相信耶穌的政治野心,但是耶穌在百姓中因所行的神蹟和教導的吸引力,正在挑戰他們的宗教權力,耶穌是危害他們權力和利益的人,是異己分子,面對異己的自然反應,就是除之而後快.出賣和殺害威脅自己權力的異己,比消滅真正的敵人還要狠毒!甚至利用敵人的力量去殺害威脅自己的權力和利益的人.這正是聖經所揭露人性的恐怖和罪惡,是耶穌和我們一樣去面對人性的醜惡!耶穌是誰?對彼拉多來說是猶太人中間被排斥的異己,是一個犧牲品.

耶穌做了什麼事情會令猶太人這樣控訴他呢?在約翰福音第二章13-25節記載了耶穌潔淨聖殿,挑戰了祭司長的權威；在五章2-30節,耶穌在安息日治病,挑戰了猶太人對於律法的解釋(另參九:13-34),並且宣稱自己和上帝的親密關係,突顯出他們和上帝之間的疏離；在六章1-15節,耶穌餵飽五千人；在八章1-11節,質問他們誰有資格定那正在犯姦淫時被他們捉的女人的罪；在十一章38-53節,耶穌叫拉撒路復活,吸引了很多人跟從他,威脅著猶太人領袖的權力和利益,但卻美其名為「若讓他這樣做,人人都要信他；羅馬人也要來毀滅我們的聖殿和我們的民族」.

耶穌是誰?他為百姓做了善良的事,吸引了百姓,但卻被視為對聖殿和民族做成威脅的危險人物!

耶穌如何回答彼拉多關於他是否猶太人的王?他沒有直接回答,而是反問彼拉多是他自己說的?還是別人說的?耶穌迫彼拉多自己回答耶穌是誰的問題!當彼拉多逃避回答,說是耶穌的同胞將他交給自己,並且質疑耶穌究竟做了什麼事而得罪了他們?耶穌也不直接回答他,而是兩次指出:自己的國不屬於這世界,自己是王,是為了真理作見證的王!

耶穌是誰?當我們準備迎接聖誕節的時候,不要被聖誕的歡樂氣氛所迷惑,要反思自己怎樣回答這問題.

他被當權者視為異己而被陷害,被當作解除對宗教和人民安全構成威脅的犧牲品,被視為對自己的權力和利益有損害的競爭對手,抑或是滿足自己的興趣和利益的工具?耶穌不是現世政治的王,我們是否企圖強迫他為王?(6:15)

對你來說,耶穌是誰?他是否真理的王?我們生命之主?

當迎接新的一年即將來臨的時候,準備耶穌隨時再降臨的時候,我們是否準備好回應,耶穌是誰的問題.



\newpage

\section{積極的等待}
\label{sec:576}
\textbf{林津牧師 - 積極的等待}
\newline
\newline
\hyperref[sec:575]{< < < PREV SERMON < < <}
~
\hyperref[sec:toc]{[返目錄]}
~
\hyperref[sec:577]{> > > NEXT SERMON > > >}


路加福音 21:25-36
\newline
\begin{longtable}{cl}
\hline
\hline
章節 & 經文 (和合本修訂版)\\
\hline
21:25 & \begin{tabularx}{0.7\textwidth}{X} 「日月星辰要顯出預兆，地上的邦國也有困苦，因海中波浪的響聲而惶惶不安。 \end{tabularx} \\ \\ \relax
21:26 & \begin{tabularx}{0.7\textwidth}{X} 人想到那要臨到世界的事，就都嚇得魂不附體，因為天上的萬象都要震動。 \end{tabularx} \\ \\ \relax
21:27 & \begin{tabularx}{0.7\textwidth}{X} 那時，他們要看見人子帶著能力和大榮耀駕雲來臨。 \end{tabularx} \\ \\ \relax
21:28 & \begin{tabularx}{0.7\textwidth}{X} 一有這些事，你們就當挺身昂首，因為你們得救贖的日子近了。」 \end{tabularx} \\ \\ \relax
21:29 & \begin{tabularx}{0.7\textwidth}{X} 耶穌對他們講了一個比喻說：「你們看無花果樹和各樣的樹， \end{tabularx} \\ \\ \relax
21:30 & \begin{tabularx}{0.7\textwidth}{X} 樹葉一長出來，你們看了自然就知道夏天近了。 \end{tabularx} \\ \\ \relax
21:31 & \begin{tabularx}{0.7\textwidth}{X} 同樣，當你們看見這些事發生，就知道神的國近了。 \end{tabularx} \\ \\ \relax
21:32 & \begin{tabularx}{0.7\textwidth}{X} 我實在告訴你們，這世代還沒有過去，一切都要發生。 \end{tabularx} \\ \\ \relax
21:33 & \begin{tabularx}{0.7\textwidth}{X} 天地要廢去，我的話卻絕不廢去。」 \end{tabularx} \\ \\ \relax
21:34 & \begin{tabularx}{0.7\textwidth}{X} 「你們要謹慎，免得被貪食、醉酒和今生的憂慮壓住你們的心，那日子就忽然臨到你們， \end{tabularx} \\ \\ \relax
21:35 & \begin{tabularx}{0.7\textwidth}{X} 如同羅網一樣，因為那日子要臨到所有居住在地面上的人。 \end{tabularx} \\ \\ \relax
21:36 & \begin{tabularx}{0.7\textwidth}{X} 你們要時時警醒，常常祈求，使你們能逃避這一切要來的事，得以站立在人子面前。」 \end{tabularx} \\ \\
[1ex]
\hline
\hline
\end{longtable}
講題:積極的等待　Wait actively in the Lord

經文:路加福音21:25-36

講員:林津牧師

場合:香港中文大學崇基學院禮拜堂主日崇拜

日期:2021年11月28日

引言:         現代人都不喜歡等待,也的確少了等待.

但其實等待是十分重要的,如果沒有等待,一切都會變得沒有價值.

背景:         主耶穌在最後的日子向門徒預言末世的事

今天是將臨期的開始,如果今晚零晨就是世界的終結,你今天會做甚麼呢?

一      照常生活

二      甚麼都不做

三      盡情享受

四      做最有意義的事情

我們可能都想知道末世是在甚麼時候,

但是我們永遠不會知道,所以重點不在於時間,而是在於我們的心態.

反應:         面對末世,其他人會被嚇得魂不附體.

我們作為信徒卻當挻身昂首,因為我們有得救贖的盼望.

不單只有平安的心,也要在世人眼前作見證.

樹葉:         知道上帝的國近了,並且相信上帝的應許都要成就.

上帝的國所強調的不是在地域上,而是在人心的歸屬,你願意嗎?

除了個人之外,上帝的國也包括其他人,你願意成為別人的幫助嗎?

謹慎:         貪食(由慾念開始)

醉酒(失去自我意識)

憂慮(失去對上帝的信心)

總結:         等待中有很多事情可以做,我們必須積極行動,

最後還要時時警醒,常常祈求,因為上帝必會引領.



\newpage

\section{等一個人}
\label{sec:577}
\textbf{林豪恩先生 - 等一個人}
\newline
\newline
\hyperref[sec:576]{< < < PREV SERMON < < <}
~
\hyperref[sec:toc]{[返目錄]}
~
\hyperref[sec:578]{> > > NEXT SERMON > > >}


西番雅書 3:14-20
\newline
\begin{longtable}{cl}
\hline
\hline
章節 & 經文 (和合本修訂版)\\
\hline
3:14 & \begin{tabularx}{0.7\textwidth}{X} 錫安哪，應當歌唱！以色列啊，應當歡呼！耶路撒冷啊，應當滿心歡喜快樂！ \end{tabularx} \\ \\ \relax
3:15 & \begin{tabularx}{0.7\textwidth}{X} 耶和華已經免去對你的審判，趕出你的仇敵。以色列的王－耶和華在你中間；你必不再懼怕災禍。 \end{tabularx} \\ \\ \relax
3:16 & \begin{tabularx}{0.7\textwidth}{X} 當那日，必有話對耶路撒冷說：「不要懼怕！錫安哪，不要手軟！ \end{tabularx} \\ \\ \relax
3:17 & \begin{tabularx}{0.7\textwidth}{X} 耶和華－你的神在你中間大有能力，施行拯救。他必因你歡欣喜樂，他在愛中靜默，且因你而喜樂歡呼。 \end{tabularx} \\ \\ \relax
3:18 & \begin{tabularx}{0.7\textwidth}{X} 我要聚集那些因無節期而愁煩的人，他們曾遠離你，是你的重擔和羞辱。 \end{tabularx} \\ \\ \relax
3:19 & \begin{tabularx}{0.7\textwidth}{X} 那時，看哪，我必對付所有苦待你的人，拯救瘸腿的，召集被趕出的；那些在全地受羞辱的，我必使他們得稱讚，享名聲。 \end{tabularx} \\ \\ \relax
3:20 & \begin{tabularx}{0.7\textwidth}{X} 那時，我必領你們回來，召集你們；我使你們被擄之人歸回的時候，必使你們在地上的萬民中享名聲，得稱讚；這是耶和華說的。」 \end{tabularx} \\ \\
[1ex]
\hline
\hline
\end{longtable}
講題:等一個人　Waiting for Someone

講員:林豪恩先生

經文:西番雅書3:14-20

場合:香港中文大學崇基學院禮拜堂主日崇拜

日期:2021年12月12日

一.引言

今天是將臨期第三主日,又稱為「喜樂主日」.將臨與喜樂有關係嗎?有人說,等待是美麗的,尤其是當你有等待的對象,在等待的過程中,往往編織浪漫的憧憬.幾年前有一套電影名為「等一個人咖啡」,其中一首主題曲名為「等一個人」,散發著類似的情懷:

天空中是否有顆星星守護我

只是我沒有發覺

人群中是否有個肩膀願為我

擋住最寒冷的冬天

誰會溫柔而堅決 帶著未來和我遇見

每個人心裡都有一個 保留著的空位

等一個人出現 擁抱著我的世界

站在愛的面前 所有寂寞防備 一瞬間崩潰

等一個人深陷 在我胸口左邊

一夜之間讓心慌 讓心痛 讓心碎

都成為紀念

等一個人成為 我生命中的永遠

幸福多麼迂迴 為何沒有發現 愛就在身邊

等一個人帶回 大雨後的藍天

從今之後每一天 每一月 每一年 都心甘情願

歌曲中的那一個人,就是電影「等一個人咖啡」中李思螢的等候.你又在等甚麼呢?你是否也是在等一個人呢?還是等某些事情,某些東西呢?你等待的來到了沒有呢?

這個學期的成長小組,學生們選擇了「迷茫」為主題.他們的解釋是:中學時等升大學,等到了,卻並沒有想像那麼滿足和快樂.這種類似的狀況,原來生活中經常出現,在未等到時,以為等到就滿足；但到出現時,卻有不外如是的感慨.這就是他們的「迷茫」.

你也有過類似的「迷茫」嗎?如果你看過電影「等一個人咖啡」,你會知道李思螢這位大學生的憧憬最後都是「迷茫」.這不只是少女情懷總是詩的情況,當年林子祥也如是:

在世間尋覓愛侶 尋獲了但求共聚

然而共處半生都過去 我偏偏又後悔

別了她原為了妳 留住愛亦留住罪

誰料伴你的心今已碎 卻有她在夢裡

為何離別了 卻願再相隨

為何能共對 又平淡似水

問如何下去 為何猜不對

何謂愛 其實最愛只有誰

有時候,我們以為等到一個人出現,我們所有的問題就得到解決.有時候,我們卻經驗到,那個人縱使出現了,但並沒有帶來「從此以後,公主和王子就快快樂樂地住在一起了」這童話式的結局,反而可能又是不外如是的失望.於是乎,有些人無奈地接受「沒有期望,就沒有失望」的信念,不再談論等待,也不再等待甚麼了.

或許,西番雅書論及的以色列人也有這樣的想法,在歷史的洪流中,不但不能力挽狂瀾,更身不由己地隨波逐流.失去的已經太多了,世界已經改變,還要等待嗎?還應該要等待甚麼嗎?

二.經文及內容

1.等失而復得／等令失而復得的上主(西番雅書)

3:14錫安哪,應當歌唱!以色列啊,應當歡呼!耶路撒冷啊,應當滿心歡喜快樂!

3:15耶和華已經免去對你的審判,趕出你的仇敵.以色列的王－耶和華在你中間；你必不再懼怕災禍.

3:16當那日,必有話對耶路撒冷說:「不要懼怕!錫安哪,不要手軟!

3:17耶和華－你的上帝在你中間大有能力,施行拯救.他必因你歡欣喜樂,他在愛中靜默,且因你而喜樂歡呼.

3:18我要聚集那些因無節期而愁煩的人,他們曾遠離你,是你的重擔和羞辱.

3:19那時,看哪,我必對付所有苦待你的人,拯救瘸腿的,召集被趕出的；那些在全地受羞辱的,我必使他們得稱讚,享名聲.

3:20那時,我必領你們回來,召集你們；我使你們被擄之人歸回的時候,必使你們在地上的萬民中享名聲,得稱讚；這是耶和華說的.」

西番雅書先知呼籲在異邦生活的以色列人,不要習慣,不要放棄,要保存等待的意志.滿心歡喜快樂,不在於目前的狀況,而在於等待將臨的狀況.將臨的是甚麼狀況呢?先知呼籲當時的以色列人等甚麼呢?

三章二十節說:「那時,我必領你們回來,召集你們；我使你們被擄之人歸回的時候,必使你們在地上的萬民中享名聲,得稱讚；這是耶和華說的.」

失去的,可以重新得到嗎?離開了,可以重回故土嗎?憑甚麼呢?無兵無將,無刀無槍,有何力量令帝國屈服呢?或許他們不是不想歸回故土,只是根本想像不到有甚麼方法可以歸回故土.畢竟,他們並不是歡歡喜喜地移民外國,現實是他們在不情不願之下被逼離開家園.無可奈何地離開,寄居外地的滋味是怎樣的呢?從他們留下的詩歌來看,很可能對故土懷著百般不捨的心情,百感交雜.詩137

1我們在巴比倫河邊,

坐在那裡,追想錫安,就哭了.

最近,有一位移居他方的才子在外地畫寫生畫時如此感懷:「人間的離合悲歡有如夕照,是短暫的,然而黑夜與白晝的輪迴是永恆的.許多年來在香港寫生無數,要以這幾張心情最複雜.許多事情年少時不明白,辜負了許多恩典,長大後方有感悟,譬若再識這片生育我們的海山.人生每值霞彩漫處,回眸方知情深.」

我們可以想像,當時離開耶路撒冷在外邦生活的以色列人的心情也是複雜的.一出一入,一次遷徙,當中擾動的情緒,如打翻的五味架,酸甜苦辣咸,一齊湧上心頭,無以名狀.有時,我們自己也承受不起翻滾的情緒,為了生活,把一切都壓下去,打落門牙和血吞.想不到,我們以為這一切都已成定局,無力扭轉乾坤,也看不到有甚麼改變的可能性的時候,先知卻重提上主的應許:「我必領你們回來,召集你們；我使你們被擄之人歸回.」失而復得,在人不能,但在上主是可能的.

上主不但知道當時的以色列人懷念耶路撒冷這個他們好幾個世代生活的家園,更關心祂的子民在這個遷徙過程中的複雜心情,主動回應他們的需要.

Ａ.安全

3:15耶和華已經……趕出你的仇敵.以色列的王－耶和華在你中間；你必不再懼怕災禍.

這是安全的提供.上主知道祂的子民在那些年間擔驚受嚇,身體,心靈和財產都受到威脅,上主也知道祂的子民不單單需要藏身之所,也需要安全地生活,需要一處不被威嚇的生活空間.

Ｂ.從罪疚釋放

3:15耶和華已經免去對你的審判

這是罪疚感的釋放.有時,罪疚感來自犯錯；有時,罪疚感來自辜負了甚麼；有時,罪疚感來自自覺背棄了理想……無論來源是甚麼,上主知道他們的沉重,幫助他們從罪疚感中釋放出來,可以問心無愧地生活.

Ｃ.羞辱變成肯定

3:19那時,看哪,我必對付所有苦待你的人,拯救瘸腿的,召集被趕出的；那些在全地受羞辱的,我必使他們得稱讚,享名聲.

這是對羞辱感的醫治.被逼離開的人,充滿挫敗,被羞辱,也自覺羞辱.上主明白他們的狀況,回應他們的羞辱感,讓他們重新得到肯定,可以挺胸昂首前行.

西番雅先知勉勵當時失落的上主子民,縱使自己力有不逮,看不到出路,仍然要懷著盼望地等待.要等待的,不是耶路撒冷,而是那位使他們失而復得,把安全,釋然,肯定重新賜予的上主.

2.等撥亂反正的上主(路加福音)

西番雅書的時代過去了,以色列人回到故土又生活了幾百年,在施洗約翰的年代,他們又等候甚麼呢?

3:7約翰對那出來要受他洗的眾人說:「毒蛇的孽種啊,誰指示你們逃避那將要來的憤怒呢?

3:8你們要結出果子來,和悔改的心相稱.不要自己心裡說:『我們有亞伯拉罕為祖宗.』我告訴你們,上帝能從這些石頭中給亞伯拉罕興起子孫來.

3:9現在斧子已經放在樹根上,凡不結好果子的樹就砍下來,丟在火裡.」

3:10眾人問他:「這樣,我們該做甚麼呢?」

3:11約翰回答:「有兩件衣裳的,就分給那沒有的；有食物的,也該這樣做.」

3:12也有稅吏來要受洗,對他說:「老師,我們該做甚麼呢?」

3:13約翰對他們說:「除了規定的數目,不要多收.」

3:14也有士兵問他說:「我們該做甚麼呢?」約翰說:「不要勒索任何人,也不要敲詐人；自己有糧餉就該知足.」

3:15百姓期待基督的來臨；他們心裡猜測,或許約翰是基督.

3:16約翰對眾人說:「我是用水給你們施洗,但有一位能力比我更大的要來,我就是給他解鞋帶也不配.他要用聖靈與火給你們施洗.

3:17他手裡拿著簸箕,要揚淨他的穀物,把麥子收在倉裡,把糠用不滅的火燒盡.」

3:18約翰又用許多別的話勸百姓,向他們傳福音.

當天文台發出颱風警告時,大家都在猜想到底會不會掛八號風球.當時,在網絡上流傳著幾句俏皮話:「整個香港,都在等颱風；就像一個初戀的少女,等待著男友,怕他不來,又怕他亂來!」

從路加福音的記載,我們看見在施洗約翰的年代,人們似乎也有類似的心情:

3:15百姓期待基督的來臨；他們心裡猜測,或許約翰是基督.

百姓都在等待基督,怕他不來,又怕他亂來(亂來的意思是不符合百姓的期望).因為他們都相信基督是大有能力的,祂有能力解決他們的問題,滿足人們的需要,同時,他們又知道基督不是人們能夠限制和控制的.約翰說基督「手裡拿著簸箕,要揚淨他的穀物,把麥子收在倉裡,把糠用不滅的火燒盡」,這情況真令人感到矛盾.到底期望祂來,還是請祂不要來呢?這意向涉及人們如何了解自己.到底我會是他收在倉裡麥子?還是祂用不滅的火燒盡的穀糠呢?我又如何知道自己是麥子還是穀糠呢?這真是一個令人不安的問題.

你又如何呢?如果基督是施洗約翰所講「手裡拿著簸箕,要揚淨他的穀物,把麥子收在倉裡,把糠用不滅的火燒盡」,你仍等候基督來臨嗎?你認為自己是麥子,還是穀糠呢?

當時的人又認為自己是麥子?還是穀糠呢?我們可以從施洗約翰的說話探出端倪來:

3:8不要自己心裡說:『我們有亞伯拉罕為祖宗.』我告訴你們,上帝能從這些石頭中給亞伯拉罕興起子孫來.

當時的人似乎認為出生比人生更決定性,出生得好,才能得獎；出生得不好,就只是雜草.即是說,一個人怎樣做人,如何處事也無關重要,最重要的是他的位置.我們又可以從施洗約翰的說話猜測,當時的人似乎是「有兩件衣裳的,都不分給那沒有的」；「有食物的,也不分給沒有的」；當時的稅吏「除了規定的數目,還會多收」；當時的士兵「除了自己有糧餉,還會勒索敲詐人」.這些都只是當時的人的信念,當時的人的表現嗎?難道在歷史中不經常是這樣的嗎?難道在現代社會,沒有人如此對待別人嗎?沒有人如此誤用自己的權力和武力嗎?正正是這樣,要來的基督「手裡拿著簸箕,要揚淨他的穀物,把麥子收在倉裡,把糠用不滅的火燒盡」,祂來撥亂反正,教導人如何做人,如何對待別人.有人以為,基督的審判是嚴厲的,然而,基督的要求其實是簡單的,相信一般人都能夠明白,而且也認為是應該的:

對擁有資源的人來說:「有兩件衣裳的,就分給那沒有的；有食物的,也該這樣做.」

對擁有權力的人來說:「除了規定的數目,不要多收.」

對擁有武力的人來說:「不要勒索任何人,也不要敲詐人；自己有糧餉就該知足.」

有人說:等待不是為了等愛情,而是準備更好的自己.借用這格式來說:等一個人不單單是基督,而是準備一個迎見基督的自己.

三.結語:

等待可以是喜樂的,如果你所等待的是你所心愛的,使徒保羅在腓立比書也如此勸勉信徒:

4:4你們要靠主常常喜樂.我再說,你們要喜樂.

為甚麼可以喜樂呢?

4:5要讓眾人知道你們謙讓的心.主已經近了.

如何可以喜樂呢?

4:6應當一無掛慮,只要凡事藉著禱告,祈求和感謝,將你們所要的告訴上帝.

4:7上帝所賜那超越人所能了解的平安,必在基督耶穌裡,保守你們的心懷意念.



\newpage

\section{絕望者的希望}
\label{sec:578}
\textbf{鄧瑞強博士 - 絕望者的希望}
\newline
\newline
\hyperref[sec:577]{< < < PREV SERMON < < <}
~
\hyperref[sec:toc]{[返目錄]}
~
\hyperref[sec:579]{> > > NEXT SERMON > > >}


路加福音 1:46-55
\newline
\begin{longtable}{cl}
\hline
\hline
章節 & 經文 (和合本修訂版)\\
\hline
1:46 & \begin{tabularx}{0.7\textwidth}{X} 馬利亞說：「我心尊主為大； \end{tabularx} \\ \\ \relax
1:47 & \begin{tabularx}{0.7\textwidth}{X} 我靈以神我的救主為樂； \end{tabularx} \\ \\ \relax
1:48 & \begin{tabularx}{0.7\textwidth}{X} 因為他顧念他使女的卑微；從今以後，萬代要稱我有福。 \end{tabularx} \\ \\ \relax
1:49 & \begin{tabularx}{0.7\textwidth}{X} 因為那有權能的為我做了大事；他的名是聖的。 \end{tabularx} \\ \\ \relax
1:50 & \begin{tabularx}{0.7\textwidth}{X} 他憐憫敬畏他的人，直到世世代代。 \end{tabularx} \\ \\ \relax
1:51 & \begin{tabularx}{0.7\textwidth}{X} 他用膀臂施展大能；他趕散心裡妄想的狂傲人。 \end{tabularx} \\ \\ \relax
1:52 & \begin{tabularx}{0.7\textwidth}{X} 他叫有權柄的失位，叫卑賤的升高。 \end{tabularx} \\ \\ \relax
1:53 & \begin{tabularx}{0.7\textwidth}{X} 他叫飢餓的飽餐美食，叫富足的空手回去。 \end{tabularx} \\ \\ \relax
1:54 & \begin{tabularx}{0.7\textwidth}{X} 他扶助了他的僕人以色列，不忘記施憐憫， \end{tabularx} \\ \\ \relax
1:55 & \begin{tabularx}{0.7\textwidth}{X} 正如他對我們的列祖說過，『憐憫亞伯拉罕和他的後裔，直到永遠。』」 \end{tabularx} \\ \\
[1ex]
\hline
\hline
\end{longtable}
講題:絕望者的希望　The Hope Of The Hopeless

講員:鄧瑞強博士

經文:路加福音1:46-55

場合:香港中文大學崇基學院禮拜堂主日崇拜

日期:2021年12月19日

各位弟兄姊妹,早安.疫症仍在,祝願各人平安.

上兩個月,觀賞了一齣話劇,就是蔡錫昌先生編導,白耀燦先生主演的《詩聖杜甫》.其中一幕,講到杜甫在成都浣花溪畔住的茅屋,被八月的狂風捲走了屋頂.杜甫的長詩《茅屋為秋風所破歌》記述這事.其中一段是這樣寫的:「床頭屋漏無乾處,雨腳如麻未斷絕.自經喪亂少睡眠,長夜沾濕何由徹?」講到風雨飄搖,長夜漫漫,屋頂被吹走了,屋內四處漏水,雨像麻線一樣從屋頂流下來,可以想像杜甫也同樣心亂如麻.杜甫說,自從安史之亂以來,他已睡得很少,再加上現在床濕被冷,長夜不能成眠.詩歌簡單幾句,道出他生命的窮困,濟世的天才未能為世所用.他蜷縮在茅寮之中,內心掛念的,卻是戰禍中悲慘的人民.睡不著,因自身的冰凍,更因人民的慘痛.長夜漫漫,是屋漏雨滴的長夜漫漫,也是苦難人間的長夜漫漫.

在長夜漫漫裡,杜甫有他的「將臨期」的渴望,他在黑暗中盼望光明.這首《茅屋為秋風所破歌》最後一段是這樣寫的:「安得廣廈千萬間,大庇天下寒士俱歡顏,風雨不動安如山!」在苦難人間,在屋漏更兼連夜雨的日子,杜甫惦念的,是千萬和自己一樣受苦的人.他渴望,有千萬間寬敞的房屋,能讓貧寒的人安居樂業.這些房屋,能在風雨中,安穩如山,屋頂不會被揭去,裡面的人獲得蔭庇.在等待光明的「將臨期」裡,杜甫不單渴望,也有決志.這首詩的結束是這樣說的:「嗚呼!何時眼前突兀見此屋,吾廬獨破受凍死亦足!」這是杜甫的悲願,他說,若果有日真有如此多的屋能讓人遮風擋雨,則就算唯獨自己的茅屋被風吹破,自己凍死,亦無憾矣.這個悲願,道出詩聖杜甫之為「聖」的理由.

一般人,遭逢時勢不佳,處境淒涼之時,多數會自憐自怨,控訴上蒼.唯有人間的聖者,在自身遭逢橫逆之時,能超越自身的痛苦,去擔負他人的痛苦,去盼望那不再有痛苦的明天.猶如耶穌在十架上,受著極大苦楚時,仍記掛著身邊的人,仍鼓舞著旁邊同樣被釘十架而將死的囚徒.

在《詩聖杜甫》這話劇中,有一幕,李白與杜甫相遇,詩仙與詩聖對飲.詩仙李白超越於人間的苦難之上,活出世的人生.詩聖杜甫雖有超世的靈魂,卻深入於人間的苦難之中.仙者,脫俗飄逸.聖者,悲憫沉重.聖者,在長夜漫漫,與千萬悲痛的靈魂同在,盼望著塵世不可企及的明天.聖者,情願自己一人受苦,期望蒼生得到解救.聖者,孤獨地背負起人類的十架,在黑暗深處,指出光明之所在.

只要我們的私心仍在,則人性的黑暗仍籠罩我們.只要一個社群與另一社群仍暴力相逼,則社會的黑暗仍瀰漫大地.只要國家與國家仍武力相向,則歷史仍在漫漫長夜中.我們仍在黑暗中,我們仍等待著光明.「將臨期」提示的盼望與等待,對我們而言,仍是重要的提示.

在耶穌出生前夕,以色列人在歷史的長夜中等待,他們等待一個被稱為「彌賽亞」的聖者.「彌賽亞」是希伯來文,希臘文是「基督」,其意思是「受膏立的人」.受膏立的人,可以是先知,祭司,或君王.以色列人等待的,最主要是君王,一位救主形態的君王,一位以神的真理統治的君王.他們期望這君王能帶領他們驅逐羅馬人,建立一正義的國度,過安樂的日子.他們稱在這新國度下的生命狀態為「Shalom」(和平,平安).當耶穌出生時,天使唱歌說:「在至高之處,榮耀歸與神；在地上,平安歸與他所喜悅的人」.「平安」,指的就是活在神的統治裡,過整全的幸福生活.這不是普普通通的開心,而是一個新時代的誕生.

當時,以色列人被外族統治了幾百年,生命困苦.每一代,他們都期望救主君王會來.幾百年過去了,他們的期待,從未止息.耶穌出生前夕,施洗約翰來了,以色列人問他:「那將要來的,是你嗎?」每一個帶來希望的人,他們都問:「你是基督嗎?你是救主君王嗎?」耶穌誕生前夕,他們的期待,不是我們點幾枝「將臨期蠟燭」輕輕鬆鬆便可以體會的.他們期待的,是一個改變天地的聖者,是一個改變他們人間命運的救主王者,是一個改變歷史軌跡的關鍵時刻.這不是一般的等待.我們在「將臨期」的等待,可能只是等待佈置家中的聖誕樹,可能只是等待朋友聚舊的聖誕聚餐,可能只是等待一個溫馨而令人感動的平安夜聚會.我們的等待,大概類似於上菜前等開飯.以色列人的等待,卻像是無期徒刑的囚犯等待著皇恩大赦,是生死的等待,是在絕望中等待著真正的希望.這才是「將臨期」應有的等待.

在這種期待下,天使對一個婢女馬利亞說,你的等候沒有落空,救主君王要來了,他卻是要由你生下來.我們聽聽這婢女的心聲.

今日的講道經文:路加福音1:46-55

路1:46　馬利亞說:我心尊主為大；

路1:47　我靈以神我的救主為樂；

路1:48　因為他顧念他使女的卑微；從今以後,萬代要稱我有福.

路1:49　(原文有「因為」)那有權能的,為我成就了大事；他的名為聖.

路1:50　他憐憫敬畏他的人,直到世世代代.

路1:51　他用膀臂施展大能；那狂傲的人正心裡妄想就被他趕散了.

路1:52　他叫有權柄的失位,叫卑賤的升高；

路1:53　叫飢餓的得飽美食,叫富足的空手回去.

路1:54　他扶助了他的僕人以色列,

路1:55　為要記念亞伯拉罕和他的後裔,施憐憫直到永遠,正如從前對我們列祖所說的話.

馬利亞這首頌歌,英文稱為Magnificat,中文一般稱為「尊主頌」,這是一首講述「終於等到救主君王」的讚美詩.

開首兩節經文,是一對平行句,點出這頌歌在頌讚誰.

路1:46　馬利亞說:我心尊主為大；

路1:47　我靈以神--(這神與我何干?)我的救主--為樂.

很平行的句子:「我心」對「我靈」；「主」對「神」(「我的救主」是對「神」的進一步說明)；「尊主為大」對「以神為樂」.在這裡,馬利亞以三個稱呼,來描繪她頌讚的對象:主,神,救主.「主」,是生命之所屬.「神」,是敬畏的對象.「救主」,在這頌歌裡,指的不是一般的救主,更是「我的」救主,是將我從逆境中搶救過來的恩人.尊主,或神,或救主為大,這是人之常情.但以主,或神,或救主為樂,則是少有的.世人很少以神為樂的,我未聽過有人以土地神,或天后,或黃大仙為樂的.以神為樂,則這位神不單救助我們,也與我們有一種親切的友誼.

跟著,要說明頌讚神的理由.「尊主頌」用兩段文字來說明.先講第一段文字.

路1:48　因為他顧念他使女的卑微；從今以後,萬代要稱我有福.

路1:49　(原文有「因為」)那有權能的,為我成就了大事；他的名為「聖」.

路1:50　他憐憫敬畏他的人,直到世世代代.

頭兩句又是很平行的句子:「使女的卑微」(即:至卑微者的卑微)對「有權能者的大事」(即:至大能者的大能)；「稱我有福」對「神名為聖」.越強調神的無限偉大,越顯出人的無限渺小.無限偉大者,如何面對無限渺小者呢?1:48開頭的動詞「顧念」很醒目.「顧念」,是神向下顧視,垂念卑微的人.

杜甫之為「詩聖」而不是「詩仙」,因為他以偉大的心靈,顧念淒涼的蒼生.在「尊主頌」裡,神被稱為「聖」.我們知道,聖經中「聖」(holy)這個字,專指神與萬物的不同.神是「創造主」,萬物只是「被造之物」,「創造主」與「被造物」之間,有一明確界線.「聖」這個字,明確標示這條界線.在「尊主頌」裡,神被稱為「聖」.這聖者如何為聖?卑微的馬利亞說,這聖者之偉大,正正表現於他顧念卑微的世人.

在歷史的黑暗長河裡,卑微的馬利亞,引頸張望,終於體會到了神的顧念.救主君王真的要來,意想不到的,是這君王要透過馬利亞這卑微的身軀誕生而來.無限偉大者要透過我們這無限渺小者去成就他的救贖.神是聖者,他的顧念,他的重視,使我們成為有福的人.卑微的婢女,艱苦的生命,由於神的顧念,體會到生命被肯定,體會到一種莫名的價值,她稱這為「福」.無論一個人多麼卑微,他仍可以不單單等待神的來臨,還可以參與神的計劃.1:50總結這神人關係,神對人,是憐憫；人對神,是敬畏.神憐憫敬畏他的人,這是神人的連繫,天人之相接,這是聖誕故事永恆的主題.

馬利亞頌讚神的第一個理由,是無限偉大的神,在他心中,竟有我們的位置.他要做的事情,竟預了我們一份.神甘願與我們連在一起.

頌讚神的理由,有第二段文字說明,又是3節經文.

路1:51　他用膀臂施展大能；那狂傲的人正心裡妄想就被他趕散了.

路1:52　他叫有權柄的失位,叫卑賤的升高；

路1:53　叫飢餓的得飽美食,叫富足的空手回去.

今次,先總結,後說明.1:51總結說,我們頌讚神的理由,就在於他的大能,正是用來驅散心懷霸道的狂傲人.在現實世界,我們最怕孔武有力的人,就如「大雄」最怕「技安」.霸道者常能大小通吃,其他人無法反抗.在馬利亞的頌歌裡,大能的神,不是霸王.他的能力,用在克服人間的霸王,回復正義.

1:52及1:53是交叉平行,第1節的頭,等於第2節的尾；第1節的尾,等於第2節的頭,即:「有權柄的失位」等於「富足的空手回去」；「卑賤的升高」等於「飢餓的得飽美食」.這種交叉平行體,或許是以文字的結構去表達「在前的,將要在後；在後的,將要在前」的乾坤大變.政治上有權勢的,再不能魚肉人民了.經濟上富足的,再不能恃強凌弱了.政治上無聲無息的,他們的尊嚴被肯定.經濟上貧乏的,他們的需要被滿足.現實政治和現實經濟,講的都是實力.但神建造的新世界,講的,是強者承擔弱者,偉大者顧念卑微者.正如後來主耶穌說的:「你們中間誰願為大,就必作你們的用人；誰願為首,就必作你們的僕人」(太20:26-27).

問題是:我們在「將臨期」裡,是否期待這個新世界?

或許,我們真正期待的,是在政治上或經濟上成為強者,勝利者,能壓在其他人的頭上.若如此,我們是在「將臨期」的期待裡,作出「反期待」.我們在期待中,期待著一種我們不願意它出現的東西.我們可能實質期待成為強者,輕視各種弱者；我們可能實質希望成為眾僕人的主人,而不甘心成為服事人的僕人.我們期待著的世界,正是反抗著神的新世界.很可能,我們是在「將臨期」裡,期望神的世界永不來臨.

馬利亞頌讚神的第二個理由,是神要撥亂反正,顛倒一個被顛倒的世界.一個本應是人性的世界,被人的罪顛倒成一個非人性的世界,神來,是將之再顛倒,重建一個充滿人性的世界.馬利亞看到,這個世界,在救主君王的救贖中來臨.今日,「將臨期」要等待的,正是這一全新世界的完全來臨.

「尊主頌」的結尾,點出神要完成的救贖,正正是他自古以來的作為.

路1:54　他扶助了他的僕人以色列,

路1:55　為要記念亞伯拉罕和他的後裔,施憐憫直到永遠,正如從前對我們列祖所說的話.

在遠古,神作出承諾,他會改變這世界,有一天,他會在人間統治,與人同在,神人連合.然後,他揀選亞伯拉罕及其子孫,透過他們成就這事.如今,他要透過耶穌,成就這事.

後來,耶穌出生了.後來,耶穌在十架上完成了他的一生.後來,耶穌要升天了.他的門徒,一群以色列人,很詫異.為何這個世界還沒有改變?救主君王為何不在這世界作王?神的寶座為何不臨現人間?在歷史長夜裡,以色列人不斷地期待,期待一位救主君王,期待他即時全面的改變這個世界.如今,令他們詫異的,是耶穌將升天了,這世界看起來並沒有不同.以色列人並沒有想過,救主君王來了,這世界卻看來沒有改變.在他們的神學裡,救主君王既來了,就進入永世.故此,在耶穌升天前夕,門徒問了一個以色列人必會問的問題:「主啊,你復興以色列國就在這時候嗎?」新天新地,是否此刻出現?耶穌的回答,按神學而言,是這樣的:「真的,天上地下所有的權柄都賜給我了,我真是救主君王,我是萬有的主.這世界表面看來沒有改變,但其底層的邪惡力量已被我克服了.然而,歷史還未終結,新天新地還未完全實現,你們留在世上,有你們要做的事情,就是:彰顯我的真理,顯明我的管治,活出新世界的生活方式.然後,我會再來,新天新地便會實現了.」

以色列的門徒,第一次理解到,救主君王的來臨,並不是帶來一完成了的「製成品的天國」,而是給予一顆「天國的種子」.或許有人會說,耶穌的降生,如此過了一生,是一種失敗,因為這世界根本沒有改變過.真的沒有改變過嗎?畢竟,耶穌留下了「天國的種子」,這是天地改變的基礎.或許有人會問,為何耶穌不一次過改變天地,一了百了呢?我想,神希望我們的活著,有一種重量,能為永恆的目標而努力,能為新天新地作出一分貢獻,而不是常常只望坐享其成.或許,唯有這樣,我們才能體會我們真的參與在神的計劃中,體會神人合作對我們生命的重大意義.

如此,「將臨期」的等待與期望,就不是百無聊賴地讓時間流逝,而是在有限的光陰裡,擔負起自己生命的十字架.這是一種有所貢獻的等待,是讓那將要來的新天新地成為生命當下的遠象,那被盼望的永恆成為當下奮鬥的動力.

在凋零破落的人生裡,杜甫在高山上呼喊:「嗚呼!何時眼前突兀見此屋,吾廬獨破受凍死亦足!」獻出了自己,讓天國臨現,死而無憾矣.

願榮耀歸上主,願人間享和平.



\newpage

\section{這就是要給你們的記號}
\label{sec:579}
\textbf{高國雄牧師 - 這就是要給你們的記號}
\newline
\newline
\hyperref[sec:578]{< < < PREV SERMON < < <}
~
\hyperref[sec:toc]{[返目錄]}
~
\hyperref[sec:580]{> > > NEXT SERMON > > >}


路加福音 2:1-20
\newline
\begin{longtable}{cl}
\hline
\hline
章節 & 經文 (和合本修訂版)\\
\hline
2:1 & \begin{tabularx}{0.7\textwidth}{X} 在那些日子，凱撒奧古斯都降旨，叫全國人民都登記戶籍。 \end{tabularx} \\ \\ \relax
2:2 & \begin{tabularx}{0.7\textwidth}{X} 這第一次登記戶籍是在居里扭作敘利亞總督的時候行的。 \end{tabularx} \\ \\ \relax
2:3 & \begin{tabularx}{0.7\textwidth}{X} 眾人各歸各城，辦理登記。 \end{tabularx} \\ \\ \relax
2:4 & \begin{tabularx}{0.7\textwidth}{X} 約瑟也從加利利的拿撒勒城上猶太去，到了大衛的城名叫伯利恆，因為他是大衛家族的人， \end{tabularx} \\ \\ \relax
2:5 & \begin{tabularx}{0.7\textwidth}{X} 要和他所聘之妻馬利亞一同登記戶籍。那時馬利亞已經懷孕。 \end{tabularx} \\ \\ \relax
2:6 & \begin{tabularx}{0.7\textwidth}{X} 他們在那裡的時候，馬利亞的產期到了， \end{tabularx} \\ \\ \relax
2:7 & \begin{tabularx}{0.7\textwidth}{X} 就生了頭胎的兒子，用布包起來，放在馬槽裡，因為客店裡沒有地方。 \end{tabularx} \\ \\ \relax
2:8 & \begin{tabularx}{0.7\textwidth}{X} 在伯利恆的野外有牧羊人，夜間值班看守羊群。 \end{tabularx} \\ \\ \relax
2:9 & \begin{tabularx}{0.7\textwidth}{X} 有主的一個使者站在他們旁邊，主的榮光四面照著他們，牧羊人就很懼怕。 \end{tabularx} \\ \\ \relax
2:10 & \begin{tabularx}{0.7\textwidth}{X} 那天使對他們說：「不要懼怕！看哪！因為我報給你們大喜的信息，是關乎萬民的： \end{tabularx} \\ \\ \relax
2:11 & \begin{tabularx}{0.7\textwidth}{X} 因今天在大衛的城裡，為你們生了救主，就是主基督。 \end{tabularx} \\ \\ \relax
2:12 & \begin{tabularx}{0.7\textwidth}{X} 你們要看見一個嬰孩，包著布，臥在馬槽裡，那就是給你們的記號。」 \end{tabularx} \\ \\ \relax
2:13 & \begin{tabularx}{0.7\textwidth}{X} 忽然，有一大隊天兵同那天使讚美神說： \end{tabularx} \\ \\ \relax
2:14 & \begin{tabularx}{0.7\textwidth}{X} 「在至高之處榮耀歸與神！在地上平安歸與他所喜悅的人！」 \end{tabularx} \\ \\ \relax
2:15 & \begin{tabularx}{0.7\textwidth}{X} 眾天使離開他們，升天去了。牧羊人彼此說：「我們往伯利恆去，看看所成的事，就是主所告訴我們的。」 \end{tabularx} \\ \\ \relax
2:16 & \begin{tabularx}{0.7\textwidth}{X} 他們急忙去了，找到馬利亞和約瑟，還有那嬰孩臥在馬槽裡。 \end{tabularx} \\ \\ \relax
2:17 & \begin{tabularx}{0.7\textwidth}{X} 他們看見，就把天使論這孩子的話傳開了。 \end{tabularx} \\ \\ \relax
2:18 & \begin{tabularx}{0.7\textwidth}{X} 聽見的人都詫異牧羊人對他們所說的話。 \end{tabularx} \\ \\ \relax
2:19 & \begin{tabularx}{0.7\textwidth}{X} 馬利亞卻把這一切的事存在心裡，反覆思考。 \end{tabularx} \\ \\ \relax
2:20 & \begin{tabularx}{0.7\textwidth}{X} 牧羊人回去了，因所聽見所看見的一切事，正如天使向他們所說的，就歸榮耀於神，讚美他。 \end{tabularx} \\ \\
[1ex]
\hline
\hline
\end{longtable}
講題:這就是要給你們的記號　This will be a sign for you

經文:路加福音2:1-20

講員:高國雄牧師

場合:香港中文大學崇基學院禮拜堂平安夜燭光崇拜

日期:2021年12月24日

1在那些日子,凱撒奧古斯都降旨,叫全國人民都登記戶籍. 2這第一次登記戶籍是在居里扭作敘利亞總督的時候行的. 3眾人各歸各城,辦理登記. 4約瑟也從加利利的拿撒勒城上猶太去,到了大衛的城名叫伯利恆,因為他是大衛家族的人, 5要和他所聘之妻馬利亞一同登記戶籍.那時馬利亞已經懷孕. 6他們在那裡的時候,馬利亞的產期到了, 7就生了頭胎的兒子,用布包起來,放在馬槽裡,因為客店裡沒有地方.

8在伯利恆的野外有牧羊人,夜間值班看守羊群. 9有主的一個使者站在他們旁邊,主的榮光四面照著他們,牧羊人就很懼怕. 10那天使對他們說:「不要懼怕!看哪!因為我報給你們大喜的信息,是關乎萬民的: 11因今天在大衛的城裡,為你們生了救主,就是主基督. 12你們要看見一個嬰孩,包著布,臥在馬槽裡,那就是給你們的記號.」 13忽然,有一大隊天兵同那天使讚美上帝說:

14「在至高之處榮耀歸與上帝!在地上平安歸與他所喜悅的人!」

15眾天使離開他們,升天去了.牧羊人彼此說:「我們往伯利恆去,看看所成的事,就是主所告訴我們的.」 16他們急忙去了,找到馬利亞和約瑟,還有那嬰孩臥在馬槽裡. 17他們看見,就把天使論這孩子的話傳開了. 18聽見的人都詫異牧羊人對他們所說的話. 19馬利亞卻把這一切的事存在心裡,反覆思考. 20牧羊人回去了,因所聽見所看見的一切事,正如天使向他們所說的,就歸榮耀於上帝,讚美他.

21滿了八天,他們就給孩子行割禮,又給他起名叫耶穌；這是他還沒有在母腹裡成胎以前天使所起的名.

在我們的生活當中,有很多不同的記號,在那些記號代表著很多不同的意義,在當中一般來講可以分為:

指示(正確的方向)

提醒(重新想起已遺忘的)

警告(危險時的幫助)

記號的例子:男女洗手間,昇降機,港鐵

小心地滑,小心花生,小心熱水

交通路牌:不准駛入,停,讓(圓形:命令,三角:警告)

沙灘紅旗,颱風暴雨

除了在生活中一般的記號以外,上帝有時候也會給我們一些很特別或者是超自然的情境,作為我們的記號,經文中上帝給牧羊人就有以下的記號:

「榮光」:在夜晚發出的大光,象徵上帝的臨在,光明必勝過黑暗.

「天使」:代表上帝宣告一個關乎萬民的信息

「嬰孩」:新生,純潔無暇,為人帶來盼望

「包著布,臥在馬槽」:謙卑,道成肉身

「天兵加入敬拜讚美」:以讚美歸榮回應主的宣告

「在至高之處榮耀歸與上帝!在地上平安歸與他所喜悅的人!」

這震撼的畫面,帶來拯救,盼望的信息,等待中的彌賽亞的世代要來臨…

從另一角度,上帝給馬利亞及約瑟的記號又是什麼呢?

「牧羊人的到來」:帶來天使對孩子的論述,對耶穌的身份(基督,救主)的確立.

「為耶穌進行割禮」:盡應盡的義,歸屬上帝,選民的重要記號.

「耶穌」一個拯救的名字:上帝所賜的名字,就是拯救.

上帝給人的記號:就是盼望的記號,等待的記號,獨特又神聖的記號.

上帝給我們的記號是什麼?又或我們今天帶著的是什麼記號呢?

基督徒生命中應該要帶著的是什麼記號呢?

「主耶穌基督」就是我們的記號:

羅馬書13:12-14

12黑夜已深,白晝將近.所以我們該除去暗昧的行為,帶上光明的兵器. 13行事為人要端正,好像在白晝行走.不可荒宴醉酒；不可好色淫蕩；不可紛爭嫉妒. 14總要披戴主耶穌基督,不要只顧滿足肉體,去放縱私慾.

「十字架」就是我們的記號:

加拉太書2:19-20

我已經與基督同釘十字架, 20現在活著的不再是我,乃是基督在我裡面活著；並且我如今在肉身活著,是因信上帝的兒子而活；他是愛我,為我捨己.

到底我們今天所仍然持有的,是極微強甚至差不多熄滅,還是要盡力點在台上能為世人帶來盼望的光?

是羞辱不願讓人知道,還是得勝高舉的十字架?

我們又是否願意「把這一切的事存在心裡,反覆思考.」



\newpage

\section{你們為甚麼找尋我呢}
\label{sec:580}
\textbf{高國雄牧師 - 你們為甚麼找尋我呢?}
\newline
\newline
\hyperref[sec:579]{< < < PREV SERMON < < <}
~
\hyperref[sec:toc]{[返目錄]}
~
\hyperref[sec:581]{> > > NEXT SERMON > > >}


路加福音 2:41-52
\newline
\begin{longtable}{cl}
\hline
\hline
章節 & 經文 (和合本修訂版)\\
\hline
2:41 & \begin{tabularx}{0.7\textwidth}{X} 每年逾越節，他父母都上耶路撒冷去。 \end{tabularx} \\ \\ \relax
2:42 & \begin{tabularx}{0.7\textwidth}{X} 當他十二歲的時候，他們按著過節的規矩上去。 \end{tabularx} \\ \\ \relax
2:43 & \begin{tabularx}{0.7\textwidth}{X} 守滿了節期，他們回去，孩童耶穌仍舊在耶路撒冷。他的父母並不知道， \end{tabularx} \\ \\ \relax
2:44 & \begin{tabularx}{0.7\textwidth}{X} 以為他在同行的人中間，走了一天的路程才在親屬和熟悉的人中找他， \end{tabularx} \\ \\ \relax
2:45 & \begin{tabularx}{0.7\textwidth}{X} 既找不著，就回耶路撒冷去找他。 \end{tabularx} \\ \\ \relax
2:46 & \begin{tabularx}{0.7\textwidth}{X} 過了三天，他們發現他在聖殿裡，坐在教師中間，一面聽，一面問。 \end{tabularx} \\ \\ \relax
2:47 & \begin{tabularx}{0.7\textwidth}{X} 凡聽見他的人都對他的聰明和應對感到驚奇。 \end{tabularx} \\ \\ \relax
2:48 & \begin{tabularx}{0.7\textwidth}{X} 他父母看見就很驚奇。他母親對他說：「我兒啊，為甚麼對我們這樣做呢？看哪，你父親和我很焦急，到處找你！」 \end{tabularx} \\ \\ \relax
2:49 & \begin{tabularx}{0.7\textwidth}{X} 耶穌對他們說：「為甚麼找我呢？難道你們不知道我應當在我父的家裡嗎？」 \end{tabularx} \\ \\ \relax
2:50 & \begin{tabularx}{0.7\textwidth}{X} 他所說的這話，他們不明白。 \end{tabularx} \\ \\ \relax
2:51 & \begin{tabularx}{0.7\textwidth}{X} 他就同他們下去，回到拿撒勒，並且順從他們。他母親把這一切的事都存在心裡。 \end{tabularx} \\ \\ \relax
2:52 & \begin{tabularx}{0.7\textwidth}{X} 耶穌的智慧和身量，並神和人喜愛他的心，都一齊增長。 \end{tabularx} \\ \\
[1ex]
\hline
\hline
\end{longtable}
講題:你們為甚麼找尋我呢?　Why were you searching for me?

講員:高國雄牧師

經文:路加福音2:41-52

場合:香港中文大學崇基學院禮拜堂主日崇拜

日期:2021年12月26日

你們為甚麼找尋我呢?

路加福音2:41-52

41每年到逾越節,他父母就上耶路撒冷去. 42當他十二歲的時候,他們按著節期的規矩上去. 43守滿了節期,他們回去,孩童耶穌仍舊在耶路撒冷.他的父母並不知道, 44以為他在同行的人中間,走了一天的路程,就在親族和熟識的人中找他, 45既找不著,就回耶路撒冷去找他. 46過了三天,就遇見他在殿裡,坐在教師中間,一面聽,一面問. 47凡聽見他的,都希奇他的聰明和他的應對. 48他父母看見就很希奇.他母親對他說:「我兒!為甚麼向我們這樣行呢?看哪,你父親和我傷心來找你!」 49耶穌說:「為甚麼找我呢?豈不知我應當以我父的事為念嗎?」  50他所說的這話,他們不明白. 51他就同他們下去,回到拿撒勒,並且順從他們.他母親把這一切的事都存在心裡. 52耶穌的智慧和身量,並上帝和人喜愛他的心,都一齊增長.

一直以來對於耶穌成長的故事,都是非常吸引人的,過去有著不少書籍,電影皆是以此為題材,嘗試去估計耶穌在出來傳道以前的生活態況,更有不少陰謀論的人士或作家,會有很多大膽的假設,為耶穌出來傳道之前的生活,產生了眾多的傳說,當中可能有不少很令人注目,很有話題,甚至不少人因此帶來可觀的收益…

但對有關耶穌生平,可靠的記載,我們就必須回到聖經.

福音書中的記載,我們不難發現,在印象中,由在伯利恆的馬槽誕降的嬰孩耶穌,按著馬上要逃到埃及(因希律要追殺2歲以下的)這一幕之後,畫面就馬上一轉,就是30歲出來傳道,中間好像空白的廿多年.

除此之外,唯獨就是只今天這段福音書的經課:路加福音2:41-52的記述…

耶穌在世上33年的人生,聖經中基本上只記錄了耶穌出生時及傳道的3年,按比例大概只有祂人生的十分之一…為什麼福音書的作者們會這樣選擇的呢?

難道耶穌整個成長階段不重要嗎?青少年的耶穌是一個怎樣的人?當中真的沒有其他值得記錄的事情嗎?

現今對青少年階段的成長有非常多的研究,我過去廿多年亦主力在青少年群體中服侍,常常思考著,如果有更多青少年的耶穌事跡的記載,我們就更能好好仿傚學習了…

根據最多人參考(但未必最準確)的維基百科:

青少年(adolescence) 是身體和心理發展的過渡階段,此一階段通常發生在青春期到法定成年期間.青少年通常是指15歲左右的人,但其應有的身體,心理或文化表現可能會更早開始和更遲結束.例如,青春期現在通常在前青少年期的階段開始,此一情況在女性中尤其普遍.認知發展(特別是大腦前額葉皮質所負責的控制衝動)和身體體型發育(特別是男性)可持續至20多歲.因此,年齡只是青少年的粗略象徵,學者們發現很難就青少年的精確定義達成共識.要對青少年的全面了解,需要從各個角度入手,包括心理學,生物學,神經認知科學,歷史,社會學,教育和人類學.宏觀所有角度,一般認為青少年是兒童和成年之間的過渡期,其文化目的是孩子要對成年人這一角色進行準備.[1]

雖然這不是一個很嚴謹的定義,但也清楚指出,可以理解「青少年」是一個過渡階段,過渡時期.

猶太人到今天仍然為到了13歲的男孩子舉行成人禮(部份傳統女孩子12歲時舉行),象徵由那時開始,該兒童就需要自己承擔責任,反之,未進行成人禮前,父母要為兒童的一切負責.

這習俗其實是與舊世界的社會很一致的,舊世界一般以農耕,畜牧,遊牧為主的社會結構來講,人生的階段從群體中的角色,簡單分類就只有3個,由一出生開始,就是一個「被人照顧的人」；成年後,有了照顧自己的能力及生產力,就成為「照顧他人的人」；到了年老,失出照顧自己的能力及生產力時,就再次成為「被照顧的人」；如果從這角度去理解,我們當然就不會特別關注所謂「過渡性」的階段及時期,因此在舊世界中,就沒有發展出對「青少年」的理解,研究；可以想像這是因為在今天的社會中,由於原本應該已經「有照顧自己的能力及生產力」,並需要「照顧他人」的人不實踐/被不實踐,因此才產生很多「青少年問題」,這真值得反思.

耶穌基督生於二千多年前的猶太地,其「青少年階段」的成長是否完全沒有可參考,學習的地方嗎?

經文路加福音2:41-43記述:

41每年到逾越節,他父母就上耶路撒冷去. 42當他十二歲的時候,他們按著節期的規矩上去. 43守滿了節期,他們回去

首先要留意,約瑟與馬利亞對主是忠誠的,也會按猶太人信仰的習慣,敬虔的過活,所以每年也會在猶太人最重要的節期─逾越節上耶路撒冷朝聖.就如往年一樣,朝聖完後,使舉家回拿撒勒…

路加福音2:43-45

孩童耶穌仍舊在耶路撒冷.他的父母並不知道, 44以為他在同行的人中間,走了一天的路程,就在親族和熟識的人中找他, 45既找不著,就回耶路撒冷去找他.

我們當中有作為父母的,請問你們有沒有試過把孩子丟失?現在大家可能能夠比較專心看顧孩童,在我成長的年代,小朋友蕩失路是非常常見的,我也清楚記得,幼年時在新年期間全家總動員在家附近花了大半天尋找走失了的弟弟,當年只有幾歲的我,也極奇擔心,花盡能量,還記得找到弟弟後的那份喜悅,與及虛脫般的情況…

相信約瑟與馬利亞由於過去的經驗(每年也帶著耶穌上耶路撒冷),加上耶穌應該也是聽教聽話的乖孩子(讓父母放心),所以在他們舉家回程時,沒有留意到耶穌沒有跟著大隊一同回去,過了一天,才開始在同行的親屬中尋找耶穌,仍然找不著,唯有回程…

試想像如若你的子女走失了,你要走回之前行過的路去找他,並且知道可能在一天前已走失了,你的心情會如何呢?除了擔心恐懼之外,可能還會忿怒與自責…

路加福音2:46-48

46過了三天,就遇見他在殿裡,坐在教師中間,一面聽,一面問. 47凡聽見他的,都希奇他的聰明和他的應對. 48他父母看見就很希奇.他母親對他說:「我兒!為甚麼向我們這樣行呢?看哪,你父親和我傷心來找你!」

當約瑟與馬利亞回到耶路撒冷(應該又走了一天路程,可能在路上也找了一整天),耶穌走失這事情已經過了三天,在聖殿中終於找到走失了的耶穌,當時祂坐在教師中間,不單聆聽,更能提出問題對應,所以讓在場的教師們希奇,這裡表現出耶穌的聰明智慧,及對信仰的熱情與熟悉…

聖經記述約瑟與馬利亞看見就很希奇,相信這「希奇」與教師們的「希奇」有著很大的分別,當然約瑟與馬利亞也可能看到耶穌的對答如流,但更重要的是:他們「看見」了尋找多日的耶穌,一個仍然未成年的兒子,一個仍然應該由他們作為父母而負責的兒子…他們更希奇為何兒子沒有按往常一樣跟隨他們回鄉,而是留在耶路撒冷,從母親的一句「我兒!為甚麼向我們這樣行呢?看哪,你父親和我傷心(非常憂慮地)來找你!」,我們可以完全想像得到他們當時的心情,身為父母的對子女的憂心及掛念,責備加上講出感受,這種語氣也是非常合理的表達…怎料,耶穌的回應竟然是:

路加福音2:49-50

49耶穌說:「為甚麼找我呢?豈不知我應當以我父的事為念嗎?」50他所說的這話,他們不明白.

耶穌這時回應的說話,當然作為父母的約瑟與馬利亞不會明白,因為在他們心目中,耶穌就是那位剛因不聽話而走失了的兒子,在責備,教導後的回應就理應是認錯及改過,耶穌的反問,好像約瑟與馬利亞回來尋找耶穌是不應該的,做錯的是父母,對未成年的人來說,被照顧是父母的責任,耶穌這個反問,在人情上很難說得過去,如果我們暫時忘記這是耶穌親口的說話,而放在一般的家庭中,這表現一定被定性為「駁咀」,不認錯,反叛,如果在華人圈子,可能更被認為是「冇大冇細」,「大逆不道」,當然很難明白…

耶穌在這裡所提到,是生命優次的問題,地上的父母是上帝所賜,亦是很重要的,但相對比於在天上的父上帝,祂才是我們真正的父親,唯一的信靠；同樣地,地上的家庭很重要,但與上帝的家(敬拜上帝的殿)相比,我們當然更應留在其中…

耶穌在此向地上的父母(約瑟與馬利亞)分享祂的價值觀並藉此向我們作出教導,也讓我們被提醒,每一個人都需要主的教導,上帝也會透過我們(這些作主門徒的人)去向人展述真理,當然對象也包括我們所愛的家人,所敬重的長輩,這或許需要很大的勇氣與信心,求主加力,並給予智慧.

請留意,事情發展不是耶穌就打發父母回鄉,自己留下在聖殿學習,成長,而是:

路加福音2:51

51他就同他們下去,回到拿撒勒,並且順從他們.

耶穌乖乖地跟從父母回去,並順從他們；耶穌選擇了順服,留在耶路撒冷,留在聖殿應該是最好的,是最應當的,但作為未成年的耶穌,祂學習了體貼父母的心,就算在客觀上耶穌所論述的完全正確,但在行動上「選擇」了順服；對,這就是耶穌所作出的示範,對青少年的弟兄姊妹,與及對我們也是非常重要的示範:就算我們認為我們是站在對的位置,也要學習體諒,明白對方的關注與感受,為對方放低當下可以放低的(耶穌沒有放低對上帝的尋求,只是在這年歲時放低留下的堅持),才會因此得著更多,更能步向更成熟.

路加福音2:51

他母親把這一切的事都存在心裡.

馬利亞的反應,讓我們想起當天使向她顯現,要揭示聖靈感孕時,記述了她就反覆思想…無論在生命中,在信仰上,當我們經歷不明白時,我們很容易有很大的情緒,甚至只懂不斷追問,希望盡快得到答覆,馬利亞所選擇的「放在心裡,反覆思想」,相信可以是我們的好榜樣,就是因為不能夠完全明白,而需要/願意繼續思考,才能不斷地記住上帝的作為,與主能夠有更緊密的團契；要記得,這事以後,相信到馬利亞完全明白耶穌基督的使命,應該還有廿多年,我真敬佩馬利亞對主的那份堅毅,讓我們也努力學習.

經文描述,這事件過後,有什麼因此而發生呢?

路加福音2:52

52耶穌的智慧和身量,並上帝和人喜愛他的心,都一齊增長.

耶穌成長了,路加福音的作者特意記述少年耶穌這事跡,並刻意地讓我們體會此事對耶穌的成長有著關鍵的意義.

路加福音2:40

40孩子漸漸長大,強健起來,充滿智慧,又有上帝的恩在他身上.

原來在整件事件前,也描述了嬰孩耶穌如何漸漸長大…

對比一下故事前後v40及v52這二節經文,同樣提到耶穌的身量/強健,智慧,與及上帝的恩/喜愛有所增長,在事件過後,v52唯獨多了一個描述,就是「人喜愛耶穌的心」也增長了,這點是值得我們留意的,特別是在我們當中的青少年,要特別留意；我們或許對耶穌有一個的印象,祂是反建制,推翻傳統,強硬,不會討好人的革命者；我也常常教導弟兄姊姊,我們處世的目的,不是要做一個人見人愛的人,因為耶穌在世時也非常惹人討厭(因為人要選擇作惡,耶穌基督的義就成為了人討厭祂的原因),但這裡卻刻意描述人喜愛祂的心也因此增長了…

正如或許我們記得,耶穌曾說過:

路加福音14:26-27

「無論甚麼人到我這裡來,若不愛我勝過愛自己的父母,妻子,兒女,兄弟,姊妹,甚至自己的性命,就不能作我的門徒. 27凡不背著自己的十字架來跟從我的,也不能作我的門徒.」

無錯,從這一個角度,作門徒是激進的,是與世界不一樣的,是要付上最大的代價的,甚至要預備為此放棄一切的,才算合格；但耶穌在世時也示範了,要按身邊/人的需要用愛去回應,由體貼門徒的限制而多番對他們不明白的信仰/比喻作講解,到以食物餵飽群眾,因愛而向病人及被鬼附的所伸出醫冶的手,在十架上不單接納身旁同釘願意悔改的,更把母親馬利亞托付給愛徒,這種種相信當中充滿了愛,體恤,人情世故；我們千萬不要理解錯誤:對,上帝是恨惡罪惡,但卻完全接納罪人的；我們不可以為討好他人而放棄信仰應該要持守的價值觀,但同時要謹記,我們也有責任去實踐愛與和平.

[1] \url{https://zh.wikipedia.org/wiki/}\%E9\%9D\%92\%E5\%B0\%91\%E5\%B9\%B4



\newpage

\section{道成了住在我們中間的肉身}
\label{sec:581}
\textbf{高國雄牧師 - 道成了住在我們中間的肉身}
\newline
\newline
\hyperref[sec:580]{< < < PREV SERMON < < <}
~
\hyperref[sec:toc]{[返目錄]}
~
\hyperref[sec:582]{> > > NEXT SERMON > > >}


約翰福音 1:10-18
\newline
\begin{longtable}{cl}
\hline
\hline
章節 & 經文 (和合本修訂版)\\
\hline
1:10 & \begin{tabularx}{0.7\textwidth}{X} 他在世界，世界是藉著他造的，世界卻不認識他。 \end{tabularx} \\ \\ \relax
1:11 & \begin{tabularx}{0.7\textwidth}{X} 他來到自己的地方，自己的人並不接納他。 \end{tabularx} \\ \\ \relax
1:12 & \begin{tabularx}{0.7\textwidth}{X} 凡接納他的，就是信他名的人，他就賜他們權柄作神的兒女。 \end{tabularx} \\ \\ \relax
1:13 & \begin{tabularx}{0.7\textwidth}{X} 這些人不是從血生的，不是從情慾生的，也不是從人的意願生的，而是從神生的。 \end{tabularx} \\ \\ \relax
1:14 & \begin{tabularx}{0.7\textwidth}{X} 道成了肉身，住在我們中間，充充滿滿地有恩典有真理，我們也見過他的榮光，正是父獨一兒子的榮光。 \end{tabularx} \\ \\ \relax
1:15 & \begin{tabularx}{0.7\textwidth}{X} 約翰為他作見證，喊著說：「這就是我曾說：『那在我以後來的先於我，因為在我以前，他已經存在。』」 \end{tabularx} \\ \\ \relax
1:16 & \begin{tabularx}{0.7\textwidth}{X} 從他的豐富裡，我們都領受了恩典，而且恩上加恩。 \end{tabularx} \\ \\ \relax
1:17 & \begin{tabularx}{0.7\textwidth}{X} 律法是藉著摩西頒佈的；恩典和真理卻是由耶穌基督來的。 \end{tabularx} \\ \\ \relax
1:18 & \begin{tabularx}{0.7\textwidth}{X} 從來沒有人見過神，只有在父懷裡獨一的兒子將他表明出來。 \end{tabularx} \\ \\
[1ex]
\hline
\hline
\end{longtable}
講題:道成了住在我們中間的肉身 　The Word became flesh and dwelt among us

經文:約翰福音1:10-18

講員:高國雄牧師

場合:香港中文大學崇基學院禮拜堂主日崇拜

日期:2022年1月2日

10他在世界,世界也是藉著他造的,世界卻不認識他. 11他到自己的地方來,自己的人倒不接待他. 12凡接待他的,就是信他名的人,他就賜他們權柄作上帝的兒女. 13這等人不是從血氣生的,不是從情慾生的,也不是從人意生的,乃是從上帝生的.

14道成了肉身,住在我們中間,充充滿滿地有恩典有真理.我們也見過他的榮光,正是父獨生子的榮光. 15約翰為他作見證,喊著說:「這就是我曾說:『那在我以後來的,反成了在我以前的,因他本來在我以前.』」 16從他豐滿的恩典裡,我們都領受了,而且恩上加恩. 17律法本是藉著摩西傳的；恩典和真理都是由耶穌基督來的. 18從來沒有人看見上帝,只有在父懷裡的獨生子將他表明出來.

去年2021年第一個主日,我以「道成了肉身」為題,分享了約翰福音第一章開始的內容,回顧了各福音書對「耶穌基督的福音」的開始,及約翰獨特的表達,亦即是今天講道的上半部份…

讓我們先作簡單的重溫:

1太初有道,道與上帝同在,道就是上帝. 2這道太初與上帝同在. 3萬物是藉著他造的；凡被造的,沒有一樣不是藉著他造的. 4生命在他裡頭,這生命就是人的光. 5光照在黑暗裡,黑暗卻不接受光.6有一個人,是從上帝那裡差來的,名叫約翰. 7這人來,為要作見證,就是為光作見證,叫眾人因他可以信. 8他不是那光,乃是要為光作見證. 9那光是真光,照亮一切生在世上的人.

耶穌基督=道

道=太初與上帝同在

道=上帝

道=創造萬物

道=創造生命之物

道=真光,照亮黑暗

道:世界不認識,不接納

接納主就可以成為上帝的兒女,因為人是從上帝所生的.

「道」成了肉身!

道成肉身就是上帝自限的表現,降卑,全因為愛.

耶穌基督是完全的人,也是完全的上帝.

上帝成為了人,不單是上帝願意明白人,而是「成為」了人.

今天的經課,集中在當耶穌道成肉身來到後,人的反應與選擇.

V10「他在世界」:道成肉身的核心,一個在世的「上帝」,弟兄姊妹,如果要你去想像,上帝出現在這世上,應該是怎樣的呢?君臨天下般的降臨?會滿足人一切的期望?

V10「世界也是藉著他造的,世界卻不認識他.」,一個重要的事實,世界是上帝所造的,按他心意,按他計劃,全權在他手中；但世界上的人(我們)卻不認識(承認)他,不單如此

V11「他到自己的地方來,自己的人倒不接待他.」,約翰指出一個要點,我們本是屬主的,由他所創造,當耶穌這道成了肉身來到,我們卻不認識,甚至不接待,抗拒,到底為什麼呢?

V14「道成為了肉身,住在我們當中,充充滿滿地有恩典有真理.」

我們是如何理解的呢?

如果有好朋友需要來你家暫住,借宿一宵,你會如何接待呢?

相信也會很殷勤,會盡力提供所有可能的協助及服侍,煮最好的飯菜,擺上最好的飲品,坐低傾很多很多的心事,一定會「客客氣氣」,遇見小問題(甚至不小的問題)也會容忍,因為很清楚這只是暫住,只是非常短暫地影響生活,是可接受,上帝是這樣道成肉身在你的世界上嗎?

又或如果你家中需要聘請傭工,與你同住,情況又會是怎樣的呢?

盡量安排分開的空間,分開生活及用餐,盡量在視線範圍內不要見到,需要時才呼喚出來,互不影響生活,因為你們是在不同世界的人?道成肉身的上帝在你身上是否如這樣的嗎?

又或試想像由於種種原因,由今天開始,你要與家人同住,在同一空間內一同生活,祖父母,外祖父母,父母,外父,外母,老爺,奶奶,新抱,女婿等等的生活相處又是如何的呢?相見好,同住難的道理,是否如道成肉身的上帝與我們的關係呢?

請留意:耶穌基督的示範是由出生開始的,我們可能從來沒有問過這個問題,為什麼上帝要選擇透過「聖靈感孕」這方式,讓耶穌由馬利亞的懷胎中出生,由嬰孩開始?

讓我們反思一下,我們心目中「道成肉身」的耶穌,不就是那位在加利利一帶傳道宣講,帶著上帝的能力醫病趕鬼,最後進入耶路撒冷,被人出賣,釘在十字架上…

如果道成肉身只有這些內容,為什麼上帝不直接把「道」降在/臨在成年人耶穌(甚至30歲時的耶穌),又或者讓耶穌直接在30歲的時空中在穿越出現,不是就已經很足夠了嗎?

耶穌基督的道成肉身,不是暫住,是完全,完整的臨在同住,所以神學上的認知,耶穌基督是完全的神,也是完全的人,不是一半一半,亦不是超人；在人看來,他與當時的人完全一樣地生活及成長,經歷同樣的種種,當然包括生命中的喜怒哀樂,當然同樣面對家庭,朋輩的影響及壓力,也當然看到當時政治經濟所帶來對人的壓迫,困苦,不會因為耶穌是神而減少了他是一個「完全,完整的人」的任何部份.

這一種道成肉身的價值,超越了我們眼中故有,超人式,「神仙降凡塵」形態的理解,因為耶穌基督道成肉身這種拯救的方式必包含著極大的愛,而這愛帶來實質的行動,而不止停留在理論與神學思考層面.

耶穌基督就是這樣道成了肉身,來到我們中間,住在我們中間,從而反映了上帝的榮美,將上帝表明出來；今天作為尋求信仰,基督的追隨者的我們又如何學習,並效法去住在世界,住在世人的中間,以彰顯基督呢?

我們又帶著什麼態度而去住呢?

像久違了的旅行,入住最好的酒店,渡假村,暫住,只關心自身最大的享受?

或如移居海外,逃避現況,只以暫住的心過活?

我們到底帶著什麼的「道」,成了怎樣的「肉身」,住在什麼人當中呢?



\newpage

\section{天上來的眾聲}
\label{sec:582}
\textbf{李均熊博士 - 天上來的眾聲}
\newline
\newline
\hyperref[sec:581]{< < < PREV SERMON < < <}
~
\hyperref[sec:toc]{[返目錄]}
~
\hyperref[sec:583]{> > > NEXT SERMON > > >}


講題:天上來的眾聲　Voices from heaven

講員:李均熊博士

場合:香港中文大學崇基學院禮拜堂主日崇拜

日期:2022年1月9日

哲理:聲從何來?

你有沒有聽過這一個經典的哲學問題:在遙遠杳無人煙的原始森林中,大樹倒下,究竟有沒有聲音發出呢?你可能即時認為,當然有,大樹倒下必然發出轟隆巨響!但你試想深一層,沒有人聽見,那其實是不是響聲?極其量都只是空氣劇烈的振動而矣.你或許知道,這個哲學問題核心其實是:究竟現實世界是外在存在,還是所謂的現實只是透過人的感觀,而其實存在我們的大腦皮層呢?這個從佛家的色即是空,到最新上影的Matrix (廿二世紀殺人網絡) 電影系列重啟 reboot,其實都和這個問題有關.

經課:天上來的聲音

今天的經課的幾段經文,都有提到聲音.以賽亞書論到上主發言,「雅各啊,創造你的耶和華；以色列啊,造成你的那位,現在如此說:你不要害怕!因為我救贖了你.我曾提你的名召你,你是屬我的……我要從北方說:交出來!對南方說:不要扣留……」這裡不單論到上主發言,更說到上主提名呼召,並對北方南方宣告命令.但上主的聲音是如何的呢?祂說話的聲線是怎樣的呢?詩篇29篇正好說明這個問題:

3耶和華的聲音發在水上；

榮耀的神打雷,

耶和華打雷在大水之上.

4耶和華的聲音大有能力；

耶和華的聲音滿有威嚴.

5耶和華的聲音震破香柏樹；

耶和華震碎黎巴嫩的香柏樹.

6他也使之跳躍如牛犢,

使黎巴嫩和西連跳躍如野牛犢.

7耶和華的聲音使火焰分岔.

8耶和華的聲音震動曠野；

耶和華震動加低斯的曠野.

9耶和華的聲音驚動母鹿落胎,

樹木也脫落淨光.

凡在他殿中的,都稱說他的榮耀.

上主的聲音如同打雷,大有能力,大有威嚴!

為甚麼舊約經課會提到上主的說話和聲音呢?皆因福音經課今日是顯現期第一週,教會傳統提醒我們紀念和思考主基督耶穌受洗的意義.特別今天福音經課選讀的是路加福音3章的一段:「21眾百姓都受了洗,耶穌也受了洗.他正禱告的時候,天開了,22聖靈降在他身上,形狀彷彿鴿子；又有聲音從天上來,說:「你是我的愛子,我喜愛你.」

疑問:究竟耶穌聽到甚麼?

「有聲音從天上來」,那會是甚麼聲音呢?那當然是上主的聲音.上主在舊約發言,呼召祂的僕人以色列,在新約祂的呼召也臨到基督耶穌身上.上主的聲音是如何呢,就是如雷貫耳,大有能力,滿有威嚴!

但你仔細再看路加的經文,耶穌是何時聽到這個聲音?是在受洗後,禱告時；換句話說,並不是有其他人聽到這個聲音,然候作出紀錄報導；並沒有立場新聞的直播或者眾新聞中國組專業的分析報導,我們又怎知這個聲音是如何,內容又是甚麼呢?若果你進一步比較對觀福音,更為詳細記載基督耶穌受洗的馬太福音,是完全不一樣的說法,馬太3章16-17節是這樣報導:「16耶穌受了洗,隨即從水裡上來.天忽然為他開了,他就看見神的靈彷彿鴿子降下,落在他身上.17從天上有聲音說:「這是我的愛子,我所喜悅的.」

不單天上的聲音內容是「這是我的愛子」而不是「你是我的愛子」；聽到聲音的時機也從耶穌禱告轉為「隨即從水裡上來」,即受洗後一上水面,就有天上來的聲音了.從敘事學的角度去分析,馬太的觀點是從一個第三方見證人的角度去敘述耶穌受洗的情況；而路加福音,則是從耶穌的自身視角去講述基督受洗的事件.

啟示:主觀經歷與客觀事件

耶穌如何聽到從天上來的聲音?一個是透過第三者的見證,一個卻是耶穌本身的經驗.這就回到我一開始的哲學問題,究竟耶穌所經驗天上來的聲音,是祂自身的感觀經驗,即只有祂聽到,是發生在祂內在的感受；還是外在的,可以由第三方的見證所可以證實的事件?你可能會問,這又有甚麼分別呢?若從這個角度去思想耶穌受洗的意義,我們可以說:馬太是邀請我們去見證基督的受洗,讓我們認識耶穌身份,是那神的愛子的客觀事實；但路加卻讓我們進入耶穌的意識中,去經驗三位一體的奧秘:聖靈的降臨,聖父的聲音,聖子的肯定,都是從耶穌的角度去體會.

回到一開始的哲學問題,我們其實可以用一種所謂明辨式的實體論(或譯批判實存論critical realism)作出回應,即是外在的事件和內在的經歷是同等重要,缺一不可；沒有外來的刺激,沒有事件的發生,一切盡都枉然；但沒有我們投入的經歷,沒有內在的體會,信仰人生也沒有任可價值.這是我們基督徒所應該把持對上主啟示的態度

所謂上主的啟示,我想今天我們未必會親耳聽到上主的聲音.在座若申請來崇基學院神學院讀書,我們一定會與你會面,有很大可能會面的第一條問題就是問你,你為甚麼想入神學院讀書.若然你說,因為有一天你祈禱時聽到神的聲音叫你來讀神學,與你會面的老師可能就會你眼望我眼,然候要你進一步講清楚,你所謂的「聽到」神的聲音是甚麼意思.若果你堅持聽到神在你耳邊和你說話,攪不好我們會請你和專業的心理醫療人員(即係心理醫生)見見面.

我們現代人並不會認為神會在我們耳邊發出物理的空氣振動,然後將聲波傳入我們耳中.通常我們說「聽到」神和我們說話,大部份時間所指的,可能是我們感受到一些訊息,好像神和我們講了甚麼似的.但我們和神的關係並不是一種純粹心理活動,宗教經驗,必須有外在的明証才可以確立我們所相信的是真確的.換句話說,耶穌是神的兒子,是基督,不是祂說了算的.我們感恩馬太透過不同人的見證去確認耶穌的確是神的兒子,耶穌的受洗就是一個確立他是神的兒子的,可被見證的事件.所以你說你聽到神叫你讀神學,我們也不是必然收你做學生,還要了解你的心志,學識程度,家庭和教會背景等等眾多的因素,看你是否適合此時此刻進入神學院讀書.

但反過來說,只看重眼前的事實,事事講求客觀,可被証實的事情,對我們的信仰可能也未必足夠.剛過去學期我教的一科處境中的希伯來聖經(Hebrew Bible in Contexts)就有一個非基督徒的同學,在課堂上表達過他不是不相信神,或基督教的實在,但他就是沒有內心的「感覺」,令他成為一個基督徒.我和班裡面的同學都為他禱告,希望他能早日有所體會,以至成為基督徒.但其實他要怎樣的經歷,才認為有所經歷呢?其實我們都有很多經歷,但我們是否相信神透過這是事件和我們交流,溝通呢?上主的確與我們說話,我們內心是否相信上主會與我們溝通,用不同的方法使我們感受祂的同在和預備呢?正如詩篇所言:諸天述說神的榮耀,穹蒼傳揚祂的手段:這日到那日發出言語,這夜到那夜傳出知識(詩19:1-2).耶穌也曾說:若是我們閉口不說,石頭也會呼叫起來,為上主作見證(路19:40).那位同學所需要的,或許是開放心靈,讓自己的意識能與上主接通,相信上主透過他生活所經歷的和他說話,使他能感受神的愛.這個星期你關注的是甚麼呢?在疫情眾多新聞中,你又有沒有在心靈中聽到天上來的聲音?我就像聽到上主用箴言的說話對坐在高位的人說:「挖陷坑的,自己必陷在其中；滾石頭的,石頭反滾在他身上.」(箴26:27)

洗禮:裡外如一的見證

其實教會的聖禮,無論是洗禮,聖餐禮,甚至婚禮,或是喪禮,都是要幫助我們從日常的事件看到背後屬靈的意義,也將屬靈的心靈連結現實的水,酒,餅,又或婚嫁,又或死亡.聖與俗並非二分,乃是互為表裡.我們的靈性若沒有每日實在的生活作為基礎和內容,那是空虛的靈性；但每日的實在沒有靈性的參悟和投入,那只有無情的實在.

基督的洗禮亦是預表我們的洗禮.福音書展示的基督的洗禮,不單是讓我們見證祂與上主的關係,也是讓我們透過基督去經歷這種關係.這也是我們的洗禮,一方面在教會,在眾人面前,見證我們與上主的關係:從水裡面上來的一刻,我們就神的兒子,不再在罪中,是教會的一員,是屬基督的子民.但我們內心又有沒有經歷這一種關係,感受到我們是上主的愛子,體會基督捨身救贖的深恩,和聖靈不住的提醒與感通?洗禮以後的生活,我們是否仍裡外如一的見證上主的深恩呢?或者反過來說,在變化莫測,甚至顯得有些荒謬的世情中,我們有沒有踏實的活出我們的信仰,又同時在黑暗世代中不失那看不見的盼望和信心呢?



\newpage

\section{獨上高樓望盡天涯路}
\label{sec:583}
\textbf{鄧瑞強博士 - 獨上高樓,望盡天涯路}
\newline
\newline
\hyperref[sec:582]{< < < PREV SERMON < < <}
~
\hyperref[sec:toc]{[返目錄]}
~
\hyperref[sec:584]{> > > NEXT SERMON > > >}


創世記 12:1-3
\newline
\begin{longtable}{cl}
\hline
\hline
章節 & 經文 (和合本修訂版)\\
\hline
12:1 & \begin{tabularx}{0.7\textwidth}{X} 耶和華對亞伯蘭說：「你要離開本地、本族、父家，往我所要指示你的地去。 \end{tabularx} \\ \\ \relax
12:2 & \begin{tabularx}{0.7\textwidth}{X} 我必使你成為大國，我必賜福給你，使你的名為大；你要使別人得福。 \end{tabularx} \\ \\ \relax
12:3 & \begin{tabularx}{0.7\textwidth}{X} 為你祝福的，我必賜福給他；詛咒你的，我必詛咒他。地上的萬族都必因你得福。」 \end{tabularx} \\ \\
[1ex]
\hline
\hline
\end{longtable}
講題:獨上高樓,望盡天涯路　Calling from Above, Unknown Way Ahead

經文:創世記12:1-3

講員:鄧瑞強博士

場合:香港中文大學崇基學院禮拜堂主日崇拜

日期:2022年1月16日

各位弟兄姊妹,早安.在這疫症裡,祝願各人平安.

這些日子以來,我越來越為弟兄姊妹是否能建立一穩固的信仰根基而焦急.信仰要穩固,必須對整本聖經所描述的神的救贖作為有深刻了解.崇基禮拜堂採用的經課系統,已大致涵蓋整本聖經的重要經文.這經課系統的好處,是根據主耶穌的生,死,復活,再來而編排,用這經課系統的經文來講道,肯定能幫助弟兄姊妹了解主耶穌的福音.這是很重要的.唯這經課系統稍稍忽略聖經正典如何藉著一卷又一卷書,連續地,一步一步地向人透露神的救贖作為.為了讓弟兄姊妹能連貫地明白聖經,從今日起,我打算逐卷逐卷順序來講,即:順序地,每次講一卷書的某章節,順道交代一下這卷書整體的神學用心.如此,六十六卷書,要講六十六個禮拜才能講完.人生有限,世事無常,看看能否有朝一日真的能講完.我希望這種講法,能讓弟兄姊妹對整本聖經有一較全面的認識.

現在是一年之初,我就以聖經的第一卷書《創世記》作開始.

中國人的儒家文化,建基於四書五經.剛好,聖經新約有四卷福音書,舊約有五經(猶太人稱之為Torah),這四書五經,也正是我們基督信仰的根基.

舊約五經之首《創世記》開頭的第一句經文是:「起初,神創造天地」.這句經文是整本聖經的基石,是支撐起整個信仰大廈的地基.這個「起初」,固然可以是「時間」之始,但更是指萬物的頭等要旨,萬物的根基.在宇宙間,具終極意義的頭等大事,便是:神創造天地.「神創造天地」的重要性,在其指出我們生命的本相,指出終極價值之所在,指出萬物之所依.

既然是「神創造天地」,我們的本相便是一有限的被造之物,這是對人的本相的終極認知.神與人之間,有一明確的界線.人類的災劫,便是越過這界線,自命為神.任何人類神化自己,都總帶來歷史的災難.

既然是「神創造天地」,我們便應知道永恆的價值並不握在人的手裡,終極的價值在神那裡.什麼時候我們絕對化我們手上的東西,我們便走上魔道.

既然是「神創造天地」,我們便應知道我們本體上不是孤島,我們不是獨個兒存在的,我們需要身外的他者,更需要那無限的他者(神).任何人漠視他人,漠視永恆的上帝,以為是自己創造整個天地,則失去人性.

「神創造天地」是首要的真理,托住人性,托住世界,是形而上的「起初」,是萬物的終極基礎.但是,但是,現實又如何?《創世記》的第2節經文便說:「地是空虛混沌,淵面黑暗」.現實的世界,離開神創造的原意很遠.一片混亂,一片黑暗.為何搞成這樣?《創世記》娓娓道來.

未講人性的墮落故事前,《創世記》先交代神的回應方案.1章3節說:「神說:『要有光』,就有了光.」在人類的黑暗長夜裡,「光」從何來?「光」來自神的說話.這不關乎人類的力量,也不關乎人類的知識,只關乎神話語的創造能力.神的說話,其實就是神內在本性的流露.神若說:「我愛你」,你就永恆地被神的愛所愛,沒有任何東西能將你和神的愛隔絕,連死亡都不能.這就是神的說話的大能.神的話語(道)後來成了肉身,來到人類中間,這就是神子耶穌.神的說話,就是光明.問題是:人類從來都不喜歡聆聽,人類喜歡發號施令.我們不喜歡聆聽異己者(異於自己者)的聲音,而最大的異己者便是神,我們不喜歡聆聽神的聲音.離開神的話語,我們看不見光.

第3章開始,《創世記》將人類的故事娓娓道來,看看人做了什麼,將大地搞得如此混亂黑暗.第3章講到,人不喜歡聆聽神的命令,我們喜歡自己作主,喜歡訂立自己的規則,滿足自己的欲望,於是,亞當夏娃違反神的禁令,吃了禁果.這禁果代表的,正是人不應越過的,神與人之間應有的界線.就像《2001太空漫遊》(2001: A Space Odyssey)那齣電影裡,那猩猩拋起那骨頭--那代表智慧,力量,權力的骨頭,便進化成為統治大地的主人.人,吃過禁果,獨立於神,便失去伊甸園,在大地飄盪.

第4章,是兄弟該隱與亞伯的故事.應是手足情深的故事,卻變成兄弟相殘的故事.像人間很多夫妻的故事,本應是相愛的故事,倒頭來卻成了你死我亡的故事.《創世記》第4章講出這人間的悲劇.大佬該隱不滿意弟弟亞伯事事好過自己,就乾脆把他殺了.一個人,頭頂上沒有神,身邊漠視他人,不敬天,不愛人,就淪落到這個地步.人自以為神,按自己的法則而行,倒頭來就只剩下暴力的法則.我強,就殺了你.大地變成黑暗的森林.《創世記》講遠古的故事,讓我們看見自己.

6章2節:「神的兒子們看見人的女子美貌,就隨意挑選,娶來為妻.」這節經文的解釋,要多少,有多少.但其總意是:天人之間應有的界線,完全被取消了.於是,人的邪惡,像洪水般瀰漫大地.最後,神行動了.神以洪水淨化這個邪惡的世界.唯有挪亞一家,及方舟裡的動物,能保存性命.

我們以為,洪水之後,人類會回到神的身邊,聆聽神的話,重返伊甸園.事實不是這樣的.第一次世界大戰後,我們想,這樣特大的災難完了,人類應該學懂了.事實不是這樣的.更瘋狂的第二次世界大戰很快便爆發了.哲學家黑格爾說:「人類從歷史中學到的唯一教訓,就是沒有從歷史中吸取到任何教訓.」

11章3-4節:「〔有一群人〕他們彼此商量說:『來吧!我們要做磚,把磚燒透了.』他們就拿磚當石頭,又拿石漆當灰泥.他們說:『來吧!我們要建造一座城和一座塔,塔頂通天,為要傳揚我們的名,免得我們分散在全地上.』」這就是著名的巴別塔故事.「巴別」即後來的「巴比倫」.巴比倫後來滅了猶大國,摧毀了耶路撒冷的聖殿.巴別塔的故事,即後來的巴比倫帝國的故事.巴別塔的故事講什麼呢?講到人造磚起通天塔.人懂得造磚,是天大的事.人不再倚靠大自然提供的石頭來建築,人能自己造出建築材料來.這是科技,這是人的力量.他們要造一通天塔,要走上神壇,通向天上,向世人宣告他們的統治.

故事來到這裡,人的科技越來越利害了,統合大地的力量越來越強了,走上神壇也走得越來越高了,「起初,神創造天地」這句話的回響越來越小了.忘記了神,忘記了聆聽他人,忘記了人的有限,於是,天地昏暗,空虛混亂.

11章5節:「耶和華降臨,要看看世人所建造的城和塔.」神臨在,變亂了人的口音,保存了人說話的多樣式,保存了人存在的多樣式,保存了人的生命尊嚴.

問題來了,人製造出來的黑暗與混亂,沒完沒了.神要如何處理?一次又一次的介入?抑或,在黑暗的大地上,找一個願意點燈的人?誰願意在漆黑的世界,舉起一把火把?

如此,便進入了第12章的故事情節,即今日的講道經文:創12: 1-3.

12:1　耶和華對亞伯蘭說:「你要離開本地,本族,父家,往我所要指示你的地去.

12:2　我必叫你成為大國.我必賜福給你,叫你的名為大；你也要叫別人得福.

12:3　為你祝福的,我必賜福與他；那咒詛你的,我必咒詛他.地上的萬族都要因你得福.」

有一天,在大地上,有一個人,忽然領悟一種天命.

他忽然對生活的一切有一種深刻的疑惑.「本地」,是他生活的領域,是他生活的疆界,他忽然有一個問號,這就是人的生命的全部範圍嗎?人活的固有領域,就是生命意義的依歸嗎?人能想像一塊更大的土地,去發現生命更大的意義嗎?「本族」,是他的親族,是他的血脈,他忽然有一個問號,人自身的親族就是人關懷的全部範圍嗎?人的關懷,能及於天地間所有的人嗎?人能關懷非我族類的人嗎?「父家」,是自己的家人,他忽然有一個問號,我真的只屬於一個家庭嗎?難道不能四海一家嗎?我身邊的敵人,難道不能看成是我的兄弟姊妹嗎?我真的要像該隱,將本是手足的亞伯看成是敵人嗎?天地間,什麼是「本地」?什麼是「本族」?什麼是「父家」?哪片大地,不可以是吾家?哪個民族,不可以是我的親族?哪個人,不可以是我的家人?大地啊,為何變得楚河漢界,你不能踏過來,我無法踏過去?人類啊,為何要生死相逼,不是你死,就是我亡?身邊的人啊,為何要親疏有別,不能平等相待?

在人間,某一天,忽然有這麼的一個人,質疑自己一直以來的生活的正當性,質疑人類一直以來的生活的慣常想法,質疑世人一直以來無法擺脫的,卻是置人於仇恨深淵的習性.他懷疑過去的一切,扣問永恆的意義,追尋人類未曾踏足過的心靈世界.他想像,能容納一切人的大地,能撇開人的種族而看到那皮膚裡面的人性,能視一切人如家人.他想像,一個從未見過的世界.他想像,一個人類的伊甸園,在那裡,人與人,人與動物,人與神,像家人一樣,相親相愛.

那是遙不可及的夢嗎?或許是.但他確實聽到那聲音,那來自上天的聲音,那來自永恆的呼喚.他心中清楚地聽到:「來吧,離開你固有的一切,本地,本族,父家,往一個全新的領域去吧.」

這裡有一個人,一個孤獨的人,一個感悟到自己是那獨特的一個人,他登上了靈魂的高樓,望向前面,望向遠方,望盡天涯路.這是不能完成的夢嗎?前路茫茫,真的要踏上去嗎?

他聽到那從天而來的呼喚,他領悟到,若不走上這茫茫天涯路,一切都會依舊,世人依舊,舊古以來的歷史依舊.萬古以來的歷史,是殘殺的歷史,是不讓他人站在大地上的歷史,是視家人如敵人的歷史.要這歷史繼續一萬年嗎?

他聽到那從天而來的呼喚,他領悟到,他必須踏出去.千里之行,始於他真的走出去.那夢,看來在他有生之年,不能實現,但,就算不能實現,也要踏出去,也要保存這夢.有這夢,才能在漆黑的大地,保存一點光明.

他知道,他要以全新的生命存在,開展一全新的國度.這國度的存在,不為自己的福祉,而為萬國萬族的福祉.如何開始呢?他要以信心--那份對終極的天地之大愛的信心--走出自己熟悉的一切.他知道自己要往何方嗎?不知道.但他有夢,他有靈魂的呼喚的引導,他有對蒼生的大愛.

太空人岩士唐作為人類第一人踏足月球這全新領域時說:「這是一個人的一小步,卻是人類的一大跳躍」.亞伯蘭(後來改名叫亞伯拉罕)離開他固有的一切的一小步,真的是人類精神的一大跳躍.生活的慣性,不再有主宰的力量.歷史的固定軌跡,不再有獨裁的力量.歷史自身的發展,不再成為人的救贖力量.歷史的發展,常常是人的力量的越發利害,暴力的越發強大,從該隱殺亞伯的兄弟仇殺,發展到巴比倫帝國對猶大國的民族滅絕.總要有人打破這循環,總要有人走出這歷史的輪迴,總要有人質問固有的一切.亞伯拉罕不將救贖放在歷史的固有發展,他將救贖寄托在從天而來的仁愛天命.這不是歷史的邏輯,這是信仰的聆聽.他不再是歷史發展下的一顆螺絲釘,他獨上高樓,成為信心之父.

他聽到從天而來的呼喚,說:

12:2　「我必叫你成為大國.我必賜福給你,叫你的名為大；你也要叫別人得福.

12:3　為你祝福的,我必賜福與他；那咒詛你的,我必咒詛他.地上的萬族都要因你得福.」

要祝福世人,就要以全新的生命存在,開展一全新的國度.這國度要慢慢開展,首先,要有一群人,甘心活出一種特殊的天命,在古時,這就需要亞伯拉罕有後裔,新生命從一家族開始,然後影響其他人,《創世記》後面的故事會發展這情節,也多次探討沒有後裔的危機.其次,一群人要活出一種特殊的生活,讓他人看見,就要有一具體的土地.土地問題,在下一次講《出埃及記》時會探討.第三,新生命的核心,是這人的生命連於神的生命.留心這句經文:「為你祝福的,我必賜福與他；那咒詛你的,我必咒詛他」.神與這人連成一線,禍福與共.簡單說,神要透過一群人的全新生命,活出神聖,好讓人斬斷罪惡的循環,展現永恆的意義.

在人類歷史漆黑的長夜裡,亞伯拉罕從歷史的固網中走出來.歷史的固網,聖經用字是「chronos」.亞伯拉罕從「chronos」走向「kairos」,而「kairos」就是人回應神的時刻--那具有永恆意義的時刻.

上天揀選了亞伯拉罕,亞伯拉罕回應了上天的呼喚,他要成為天命的承擔者.人性的邪惡會在歷史中不斷活出它的殘忍,而天命的承擔者要以仁愛對抗這人性的殘酷.天降大任於人,他的人生將面對各種艱難.無視這一切艱難,亞伯拉罕--作為領受天命的第一人--毅然離開固有的一切,尋索著生命的永恆意義.

不知道有多少人會有亞伯拉罕的那種靈魂觸動?

不知道有多少人會像亞伯拉罕那樣毅然起行?

不知道有多少人能像亞伯拉罕那樣以信仰活到底?

五經的故事如何發展下去呢,下個月第三個禮拜,我在這裡講《出埃及記》.約定你,下個月見.

願榮耀歸上主,願人間享和平.



\newpage

\section{他在那裡他並不鍼默}
\label{sec:584}
\textbf{伍渭文牧師 - 他在那裡,他並不鍼默}
\newline
\newline
\hyperref[sec:583]{< < < PREV SERMON < < <}
~
\hyperref[sec:toc]{[返目錄]}
~
\hyperref[sec:585]{> > > NEXT SERMON > > >}


講題:他在那裡,他並不鍼默　He Is There and He Is Not Silent

講員:伍渭文牧師

經文:詩篇19

場合:香港中文大學崇基學院禮拜堂主日崇拜

日期:2022年1月23日

引言

2015年7月14美國太空總署(NASA)耗資\$720百萬,穿越太空9年(2006.9.1發射)的「新視野號」(New Horizons),終於到達冥王星,傳回高清照片.冥王星離太陽系最遠(30億英里),248 年環繞太陽一周.當太空總署同人站立歡呼時,他們也知道,新視野號也有一天步航海家(Voyager)1號同一命運,消失在太空中,失去聯系.2013年9月13日NASA確認,航海家(Voyager)1號已經飛出太陽系,進入星際空間36年後已失聯.航海家一號是存的,但人沒辦法聯系.對昂貴的飛行物失聯我們感到可惜,但若失聯的是你所愛的人,而他在那裡,但拒絕跟你聯系,你就會感到莫大的痛苦了.

國內著名歌手刀郎的《沖動的懲罰》,稱為寫給妻子的墓誌銘,埋葬跟她的感情.二十二歲的刀郎,跟前妻生下女兒剛40天, 妻子就離他而去.尋妻十年,毫無音訊,最後他寫下這首歌.

「那夜我喝醉了, 拉著你的手,胡亂的說話, 只顧著自己心中壓抑的想法, 狂亂的表達,我迷醉的眼睛, 已看不清你表情, 忘記了你當時會有怎樣的反應,我拉著手, 放在我手心, 我錯誤的感覺到你也沒有生氣,所以我以為… …如果說不是老天讓緣份把我捉弄,想到你我就不會那麼心痛, 就把你忘記吧,應該把你忘了, 這是對衝動最好的懲罰 .」 尋找所愛的人,而遭沉默回絕,痛苦莫過於此.

今日講題他在那裡,他並不沉默,是來自法蘭西.薛華(Francis Schaeffer,1912-84)所寫的一本書He Is There and He Is Not Silent(中譯:祂存在而且並不沉默).這書1972年出版,三十年後2001年再版.

今天經課也說明他在那裡,他並不沉默-上主有話要說.詩篇經課:他透過諸天來述說,也透過律法來跟我們說話(詩十九)；舊約經課:藉文士以斯拉教導律法,對我們說話(尼八1-3,5,6,8-10)；福音經課:在末日差遣父懷中獨生子,他是上主心中至重要的話,也是律法的總歸,告訴我們:「主的靈在我身上,因為他用膏膏我,叫我傳福音給貧窮的人,差遣我報告被擄的得釋放,瞎眼的得看見,叫受壓制的得自由,報告神悅納人的禧年.」(路四18, 19)「今天這經應驗在你們耳中了.」(路四21)書信經課:聖靈在教會賜下使徒,先知,教師等恩賜,教導聖經,對我們說話.(林前十二29)

薛華他在那裡,他並不沉默這書,來自他的公開講演「Possible Answers to the Basic Philosophical Questions」(基本哲學問題的可能答案).他說哲學有三大基本問題:

形而上問題 -受造世界背後有造物主嗎?

道德問題-行為可有對錯(善惡)標準嗎?

知識論問題-如何得到真確的知識?

今天詩篇經課詩十九,正應對以上的問題.

1)諸天述說神的榮耀,1-5節.「穹蒼傳揚他的手段.這日到那日發出言語,這夜到那夜傳出知識.」大自然以無言之言跟我們說話.

2)律法定對錯,6-11節.「耶和華的典章真實,全然公義.」(9節)上主藉律法跟我們說話,律法有對錯的標準.

3) 12-14節談到罪與赦免.哲學知識論討論如何得到真確的知識.蘇格拉底說:知識的目的是「認識自己」(Know Thyself ).奧古斯丁說最重要的知識就是「認識自己是一個罪人,且藉悔改認罪,得到上主的赦免.」「求你攔阻僕人,不犯任意妄為的罪,不容這罪轄制我,我便完全,免犯大罪.」(詩十九13)

1諸天述說神的榮耀(詩十九1-6)

上帝存在嗎?絕大部分的人都相信有造物主.「神的永能和神性是明明可知的,雖是眼不能見,但藉著所造之物就可以曉得,叫人無可推諉.」(羅一20)但造物主與我無關.他好像鐘表匠,上了發條,鐘表就自行運轉,世界自動駕駛(auto-pilot).上帝存在,但聯系不到,或他沉默回絕,我們稱此為自然神論(Deism),自然神論與無神論祇是一步之遙.

無言之語.諸天無言無語,也無聲音可聽,是無言之語.一位同是作家兒子和父親,因為對一些事情產生很大的意見,之後不多聯系.不竟是父子,兩人還是常常想念對方,也想辦法打開心結.有一天兒子主動約父親一起往小時候常到的湖垂釣,但坐得很接近但說話不多.父親在日記中寫下:今天真失敗,我兒子很少跟我說話,雖然他坐在我旁,就是那麼專心釣魚.但兒子同一天的日記說:「我樂透了,父親整個下午陪著我,雖然沒有多說話,但我感到他的陪同.兒子重視藏在心中的意念,父親看重是宣諸於口中的言語.

靜默讚嘆.晨昏開晝夜,黑白舉乾坤；潮汐隨月轉,四季依日運.河漢渺渺,巨流滔滔.大自然鬼斧神工,秀麗宏偉,令人看得口呆目瞪.面對恢弘,讚嘆驚訝；面對奧秘,祇有靜默.

經文運用的意象,強調創造主宏大中的細致的保治(providence):日夜穩定節奏:「這日到那日,這夜到那夜」(二節)；普及眾生:「他的量帶通遍天下,他的言語傳到地極.」(四節)；恩澤萬民:「沒有一物被隱藏不得他的熱氣.」(六節)是的,「日頭照好人,也照歹人；降雨給義人,也給不義的人.」(太五45)

有人問愛恩斯坦:「有關宇宙的問題,你最想提問甚麼?」他說:「我最想知道這宇宙是否對我們友善?」相對宇宙,人確是渺小和軟弱,不要說甚麼外星人入侵,就是一塊隕石擊中地球,引起巨大海嘯,我們也遭滅頂.

輕輕呵護.從保治角度看,無言無語,就像為父母或祖父母,給睡著的兒孫,輕輕蓋上被褥,不作一聲,祇有滿足的微笑,怕弄醒他們.1528年路德寫【小問答】,教導孩童明白基本義理,談到創造,都是從上主保治看顧入手.

「我信上帝,全能的父,創造天地的主.

這是什麼意思? 我信上帝造我和萬物；賜我身體,靈魂,耳目,百體,理性和感官,並且繼續保存.衪又賜我衣屜,飲食,屋宇,家庭,妻子,兒女,土地,牲畜,和我擁有的一切.為養活我身體生命,衪每日豐富地供給我一切所需用的.無論有什麼危險,兇惡,衪護庇,看顧,保佑我.這都是出於衪為父,為神的良善和憐憫,並不因我有甚麼功勞或配得的.為這一切的事,我有責任感謝,頌揚,事奉順從衪.」

諸天像一幅圖像,誘發人思考,有效傳遞信息.出色的廣告,大部分利用圖像,畫作或攝影作品,然後一,二行文字,誘發讀者的想像力.不竟,現代資訊太多,讓人資訊疲勞.諸天是圖像也是符號,有揭示同時又遮蓋實有的功用,一如所有藝術作品,能誘發追求更大的揭示的渴望.

西箱記崔鶯鶯把手帕掉落地下,是一視覺符號,張生看見拾起手帕,是揭示同時是遮蓋.拿起手帕後,張生思前想後,究竟小姐有甚麼啟示呢?我今晚要到西廂打探打探.無言無語,也無聲音可聽,很要命的.

在【他在那裡,他並不沉默】書中,提及一位青年想不通為何世界那麼多苦難,心靈因此受到很大的折磨,有一次來到一個懸崖,本想跳下去,但被山川的秀麗攝住了,大自然無言無語,但壯麗的景色本身就是信息,造物主留下了痕跡,他應是美善的造物主.

日月雙輝.詩篇十九篇說日間太陽,詩篇第八篇提及夜間的月亮,兩篇合在一起更能體會上帝造化的奇妙.1969年七月三十日,有史以來全球最多觀眾屏息對著電視機,觀賞兩位太空人巴茲.奧爾德林(Buzz Aldrin) 與尼爾.阿姆斯特朗( Neil Armstrong)首次踏足月球.1968年聖誕前夕,太陽神八號環繞月球時,數位太空人因為在太空誦讀創世記,而被無神論者麥達琳.默里.歐黑爾(Madalyn Murray O’Hair)控告太空總署(最後法院不受理) ,奧爾德林受著很大的壓力,要他放棄在太空讀出有關創造的經文.他如何處理呢?

他很有智慧說:『休斯頓,這是雄鷹(Eagle) 機長說話.我邀請大家有片刻靜默,Over.我邀請每位聽著,無論你身在何處,你是何人,請現在就剛才數小時眼睛所見到的一切,作出你認為合適的感恩.』跟著他對自己唸出約翰福音15:5 (我是萄葡樹:你們你們是枝子,離了我你就不能作甚麼……),領受聖餐.在回程進入大氣層之際,奧爾德林讀出詩篇第八篇第三,四節: 我觀看你指頭所造的天,並你所陳設的月亮星宿,便說,人算甚麼,你竟願念他!世人算甚麼,你竟眷顧他!

他在那裡,他並不鍼默.上主藉大自然對我們說話.更進一步,他藉摩西在西乃山頒下律法對我們說話；律法是上主的訓詞和命令.時候到了他更藉他的獨生愛子對我們說話.

2 耶和華的律法全備(詩十九7-11)

這段經文用了六個不同名詞分三組形容律法:1)律法全備,法度確定；2)訓詞正直,命令清潔；3)道理潔淨,典章真實.每對是平行的,第一與第三對類似,用以突出中間的訓詞和命令.第一對:律法全備,法度確定,指律法的本身內容完整(七節),成為標竿(法度確定)；第三對:道理潔淨,典章真實,指律法是公平制度(九節),並無瑕疵(道理潔淨).法度確定,典章真實,在社群中成為重要制度.

中間第二對,指出律法是上主的訓詞和命令(八節),不是人為制定(positive law),是自然法則(natural law),與人的良心呼應.所以律法能甦醒人心,使愚人有智慧,能明亮人的眼目,律法使人明白對與錯.外在律法宣示,內在良心和應,所以是全然公義.像精金般寶貴,像蜜般甘甜.

諸天滋生萬物,供養我們；諸天無言無語,像不作聲的父母,把被褥輕輕覆蓋在睡著的孩子身上.律法是訓詞和命令,像父母對兒女百般叮嚀,甚至大聲訓斥警告,不要誤入岐途.

康德說:「有兩件事物我愈是思考愈覺神奇,心中也愈充滿敬畏,那就是我頭頂上的星空與我內心的道德準則.」上主的永能和神性,是明明可知；物理世界背後有造物主,道德世界背後有對與錯.這都是不証自明,與生俱來.

但對選民來說,律法不單是判定道德是非對錯,擁有律法是選民的記號.在西乃山上主藉摩西頒下律法,割禮和律法是選民的記號.「如今你們若實在聽從我的話,遵守我的約,就要在萬民中作屬我的子民,因為全地都是我的.你們要歸我作祭司的國度,為聖潔的國民,這些話你要告訴以色列人.」(出十九5, 6)律法小部分是敬拜的禮儀,大部分是民事,如何如人相處.如收割不可拾起掉在地上的麥穗,讓有需要的人拾起來裹腹.(申二十四19；得二3)

今天舊約經課尼八章,文士以斯拉宣讀並教道律法書,讓被擄回歸的以色列毋忘選民的召命.書信經課提到教會的教導恩賜,話語的教導沒有停止,聖靈在教會賜下使徒,先知,教師等恩賜,教導聖經.(林前十二29)

信仰群體若勵行律法,就給未信的人看到盼望.去年因疫情晚上六時後食肆不能堂食,工人下班拿著外賣,趁熱坐在沒有瓦遮頭的公園進食.有教會開放地方,讓他們有尊嚴好好享用晚餐,並供應熱茶,使用衛生間.這就是遵行律法,行公義,好憐憫.

律法有三個用途:1)第一用途,亦稱政治用途.律法像籬笆(fence),律法針對人的罪性發出警告:不可殺人,不可偷盜,不可作假見証,不可姦淫,這四條誡命,在不同的社會已成為刑法了.2)第二用途,亦稱稱義用途,律法像鏡子(mirror),使我們看到自己的虧欠,驅使我們尋求基督的拯救,藉信稱義.3)第三用途,亦稱成聖用途,律法像蹊徑(road).蒙上主赦免後,遵守誡命,顯出對人的愛,也作為對上主的感載.

律法像籬笆,建立人與人的行為現範.一九九三年全球宗教領袖召開「世界宗教議會」(Parliament of World Religions)百週年大會,通過由基督教界發起的「全球倫理宣言」,並得到各學界,政界有識之士的響應,一九九六年成立「互動委員會」(Interaction Council),推廣全球倫理教育.根據十誡中不可殺人,不可偷盜,不可作假見證,不可姦淫四條古訓,作為全球倫理的綱領,藉此建立非暴力,公平經濟秩序,兼容愛真理,男女平等的新文化.強調倫理責任,凡人皆當遵守.全球倫理糅合孔子所倡議的「己所不欲,勿施於人」及基督所說:「你想人怎樣待你,你便怎樣待人」的金科玉律.二零零零年秋在聯合國召開的「宗教與世界和平大會」重申全球倫理宣言,強調個人對他人的倫理責任,平衡過去聯合國人權宣言強調個人權益的不足.上主藉大自然眷顧眾生,也藉律法維持社會的公平穩定,這是普遍的恩典,上主讓「日頭照好人,也照歹人；降雨給義人,也給不義的人.」(太五45)

律法第二用途像鏡子,指出人的罪性.哲學很想知道的第三個是知識論問題:人如何獲得真確的知識.但最真確也是最重要的知識是甚麼呢?就是人皆有罪,上主願意赦免人的罪.這是至重要的信息,大自然說不清,律法講不透,祇有父懷裡的獨生子,將此信息顯明.就如在前面引用奧古斯丁說:最重要的知識就是「認識自己是一個罪人,且藉悔改認罪,得到上主的赦免.」

3 求你不容這罪轄制我(詩十九12-14)

律法能甦醒人心,也能明亮人的眼晴,也使愚人得智慧,但人不能憑自已遵守律法.詩人坦白承認他任意妄為,而且陷溺罪中不能自拔.「誰能知道自己的過錯呢?願你赦免我隱而未現的過錯.求你攔阻僕人,不犯任意妄為的罪,不容這罪轄制我.」(詩十九12,13)

詩人呼求上帝的拯救.「耶和華我的磐石,我的救贖主阿,願我口中的言語,心中的意念,在你面前蒙悅納.」(詩十九14)這經文是牧師講道前喜歡用來禱告的經文.有些牧師更分折經文:「願僕人口中的言語,弟兄姊妹心中的意念,蒙你悅納.」其實意念與言語互為表裡,宣諸於口為言語,藏諸心內為意念；有時候言語不慎,暴露內在的意念,說漏了嘴.「願我口中的言語,心中的意念,在你面前蒙悅納」是心口一致,別無虛言,詩人真心求上主赦免.

律法就像鏡子,使人認識自己無能為力,呼求上帝的拯救.但誰人能夠「口中的言語,心中的意念,在神面前蒙悅納呢?」除了這一位.誰能登耶和華的聖山呢?「就是手潔心清,不向虛妄,起誓不懷詭詐的人.」(詩二十四4)也是這一位.

詩篇二十四篇傳統被視為約櫃重返聖殿的詩篇.以色列人因為拜偶像,上主讓約櫃被非利士人擄去.詩二十四篇是慶賀約櫃重返聖殿的詩篇.第一,二節讚美上主之後,出現禮儀對答:朝聖者問:「誰能登耶和華的聖山?誰能站在他的聖所?」(v. 3) 然後祭司答:「就是手潔心清,不向虛妄,起誓不懷詭詐的人.」(.4) 我們可以想像,這班要登山進入聖殿的朝聖者,中間有抬著約櫃的人,他們因為約櫃引帶,轉叩問為凱旋的歡呼:「眾城門哪,你們要抬起頭來!永久的門戶,你們要被舉起!那榮耀的王將要進來」(v.7)

祇有一位是「手潔心清,不向虛妄,起誓不懷詭詐」,但他已成為贖罪祭,他的血已塗在約櫃的頂成為施恩座.施恩座要進入耶路撒冷至高之錫安山的金頂聖殿；被殺的羔羊,要登上至高上主的寶座,坐在其右.因被殺的羔羊,已成為掌管萬有的主.「榮耀的王是誰呢?萬軍之耶和華,他是榮耀的王!」(v. 19) 從來沒有人被諮詢過就生下來,也是除了這一位,但他答應了.這位是父上主心中的話語或道,這話語太初與上帝同在,這道成為肉身,名叫耶穌.

今天是顯現後第三主日,一月六日顯現日記念耶穌在約但河接受了水禮,代替人接受了審判,約但河水漫過其身,像昔日洪水漫過大地.所以,彼得前書三20-21就以洪水漫過方舟比喻信徒的洗禮:”借著水得救不多,只有八個人,這水所表明的洗禮,借著基督的復活,也拯救你們.”約旦河流淌的水,帶給我們歡樂,因我們被潔淨.在顯現日耶穌洗禮的禱文中,提及耶穌在約旦河潔淨他的新婦教會.”今天教會與他天上的配偶結合了,因為基督在約旦河洗去了他的罪惡；博士們帶著禮物迅速前來,奔赴王家的婚宴；變水為酒,使賓客歡心,哈利路亞.”

顯現節是歡樂的節日,我們的聖帷是白色或金色,是聖誕期的高潮和終結.屈梭多模(金口約翰John Chrysostom)主教說:”今天基督進入了約旦河,以洗滌世人的罪過,今天上主的聲音,響徹水面,你是我的愛子,我喜悅你.今天,聖靈借鴿子降臨在水面上,報告人類永遠滅亡的災禍已停止.」”(每日頌禱[Liturgia Horarum]一,頁308.)

領洗後,耶穌公開傳道.今天福音經課路四14-21記述耶穌受洗後首次在會堂讀出當天經課-賽六十一1, 2.耶穌的講道非常短但很震撼:「今天這經應驗在你們以耳中了.」(路四21),以賽亞先知書五百多年的預言,經過兩約之間四百年之沉默,終於當下實現了.詩十九13詩人祈求:「求你攔阻僕人,不犯任意妄為的罪,不容這罪轄制我.」耶穌說:「差遣我報告被擄的得釋放.」(路四19)

律法第三用途-律法像路徑,行出成聖的路.得到赦免的人,他口中的言語,心中的意念,便成為他人的祝福.他心被恩感,行事為人遵從聖經的教訓,律法就像路徑,行出成聖的果子.一個藉聖經像鏡子,看到自己的本相:任意妄為,被罪轄制,心魔盤踞,口中言語,心中意念都不蒙上主悅納,但當他向上主尋求憐憫,便蒙赦免.得蒙赦免後,他口中的言語,心中意念分蒙上主悅納,分別為聖,為上主所用.

述說神的憐憫.諸天述說神的榮耀,但祇罪得赦免的人,能述說神的憐憫.約櫃存放的是第二塊法版,第一塊因為拜金牛犢摔碎了.摩西上西乃山,久久不下來,他們按捺不住,造金牛犢代替上主敬拜.摩西下山聽到百姓圍繞金牛犢跳舞唱歌,忘形放肆,「便發烈怒,把兩塊版扔在山下摔碎了」.(出三十二19)

人不能憑自己遵守律法,摩西求上主赦免百姓的罪:「不然,求你從你所寫的冊上塗抹我的名.」(v. 32)摩西知道百姓不容易帶領,求上主同行,並顯出他的榮耀作為保證:「求你顯出你的榮耀給我看.」(出三十三18)「耶和華說:『我要顯出我一切的恩慈,在你面前經過,宣告我的名.我要恩待誰就恩待誰；要憐憫誰就憐憫誰』；又說:『你們不能看見我的面,因為人見我面不能存活.』耶和華說:『看哪,在我這裡有地方,你要站在磐石上.我榮耀經過的時候,我必將你放在磐石穴中,用我的手遮掩你,等我過去,然後我要將我的手收回,你就得見我的背,卻不得見我的面.』」(出三十三21-23)

摩西求見上主的榮耀,上主顯示他的憐憫,好像牛頭不搭馬嘴,其實憐憫就是上主的榮耀.基督在十字架的死,就是上主顯出他榮耀的時候.當上主的榮耀經過時「你就得見我的背」是甚麼意思呢?為何見到上主的背,就是見到他的榮耀?因為沒有人能見到上主的本體,我們一進入上主赤裸裸的聖潔,就消滅了；就如細菌見到激光極大的熱能,紛飛煙滅.上主的背就是耶穌,從來沒有人見過上主,祇有在父懷裡的獨生子將他表明出來.上主的榮耀,就是從午正到申初遍地玄黑；基督掛在木頭上就創顯出上主的榮耀,他的榮耀就是他的憐憫.

諸天述說神的榮耀,但祇罪得赦免的人,能述說神的憐憫；也祇有曾被律法定罪,後得著釋放的人,能把律法解釋得清楚.

總結:

今天誰在你心中說話,甚麼在你心中說話?在你心底至深處,被甚麼意念盤踞,揮之不去?今天詩篇給我安慰的話:雖然我們任意妄為,被罪轄制,心魔盤踞,若誠心尋求上主,必得釋放,罪得赦免.而且他還要使用我們的生命:口中的言語,心中的意念,蒙上主使用.請你看清楚,上主就在那裡,上主並不鍼默,祇要你謙卑聆聽,全心順服,必有所聞.



\newpage

\section{一時的光輝未必是永恆}
\label{sec:585}
\textbf{林豪恩先生 - 一時的光輝未必是永恆}
\newline
\newline
\hyperref[sec:584]{< < < PREV SERMON < < <}
~
\hyperref[sec:toc]{[返目錄]}
~
\hyperref[sec:586]{> > > NEXT SERMON > > >}


路加福音 4:16-29
\newline
\begin{longtable}{cl}
\hline
\hline
章節 & 經文 (和合本修訂版)\\
\hline
4:16 & \begin{tabularx}{0.7\textwidth}{X} 耶穌來到拿撒勒，就是他長大的地方。在安息日，照他素常的規矩進了會堂，站起來要念聖經。 \end{tabularx} \\ \\ \relax
4:17 & \begin{tabularx}{0.7\textwidth}{X} 有人把以賽亞先知的書交給他，他就打開，找到一處寫著： \end{tabularx} \\ \\ \relax
4:18 & \begin{tabularx}{0.7\textwidth}{X} 「主的靈在我身上，因為他用膏膏我，叫我傳福音給貧窮的人；差遣我宣告：被擄的得釋放，失明的得看見，受壓迫的得自由， \end{tabularx} \\ \\ \relax
4:19 & \begin{tabularx}{0.7\textwidth}{X} 宣告神悅納人的禧年。」 \end{tabularx} \\ \\ \relax
4:20 & \begin{tabularx}{0.7\textwidth}{X} 於是他把書捲起來，交還給管理人，就坐下。會堂裡的人都定睛看他。 \end{tabularx} \\ \\ \relax
4:21 & \begin{tabularx}{0.7\textwidth}{X} 耶穌對他們說：「你們聽見的這段經文，今天已經應驗了。」 \end{tabularx} \\ \\ \relax
4:22 & \begin{tabularx}{0.7\textwidth}{X} 眾人都稱讚他，並對他口中所出的恩言感到驚訝；他們說：「這不是約瑟的兒子嗎？」 \end{tabularx} \\ \\ \relax
4:23 & \begin{tabularx}{0.7\textwidth}{X} 耶穌對他們說：「你們一定會用這俗語向我說：『醫生，你醫治自己吧！我們聽見你在迦百農所做的事，也該在你自己的家鄉做吧。』」 \end{tabularx} \\ \\ \relax
4:24 & \begin{tabularx}{0.7\textwidth}{X} 他又說：「我實在告訴你們，沒有先知在自己家鄉被人接納的。 \end{tabularx} \\ \\ \relax
4:25 & \begin{tabularx}{0.7\textwidth}{X} 我對你們說實話，在以利亞的時候，天閉塞了三年六個月，遍地有大饑荒，那時，以色列中有許多寡婦， \end{tabularx} \\ \\ \relax
4:26 & \begin{tabularx}{0.7\textwidth}{X} 以利亞並沒有奉差往她們中任何一個人那裡去，只奉差往西頓的撒勒法一個寡婦那裡去。 \end{tabularx} \\ \\ \relax
4:27 & \begin{tabularx}{0.7\textwidth}{X} 在以利沙先知的時候，以色列中有許多痲瘋病人，但除了敘利亞的乃縵，沒有一個得潔淨的。」 \end{tabularx} \\ \\ \relax
4:28 & \begin{tabularx}{0.7\textwidth}{X} 會堂裡的人聽見這些話，都怒氣填胸， \end{tabularx} \\ \\ \relax
4:29 & \begin{tabularx}{0.7\textwidth}{X} 就起來趕他出城。他們的城造在山上；他們帶他到山崖，要把他推下去。 \end{tabularx} \\ \\
[1ex]
\hline
\hline
\end{longtable}
講題:一時的光輝未必是永恆　A Moment of Glory may not Last

講員:林豪恩先生

經文:路加福音4:16-29

場合:香港中文大學崇基學院禮拜堂主日崇拜

日期:2022年1月30日

一.引言

你聽過「一剎那的光輝並不代表永恆」這句說話嗎?來源是這樣的:1985年《勁歌金曲》第二季金曲季選,蔡楓華在節目中報導張國榮演唱會消息時,說了一句「一時ge 光輝未必係永恆ge 」,報章寫成文字時變成「一剎那的光輝並不代表永恆」.隨後傳媒大肆報導蔡楓華的言論是針對張國榮,結果令他失去了《勁歌金曲》主持的席位.事後蔡楓華在電台節目中解釋:「那句話,只是有感而發,對藝人生涯的感慨,並非針對任何人.」

「一時ge 光輝未必係永恆ge 」豈止是對藝人生涯的感慨,難道世事不也是類近嗎?

二.經文及反省

一時的光輝

4:16耶穌來到拿撒勒,就是他長大的地方.在安息日,照他平常的規矩進了會堂,站起來要念聖經.

4:17有人把先知以賽亞的書交給他,他就打開,找到一處寫著說:

4:18主的靈在我身上,因為他用膏膏我,叫我傳福音給貧窮的人；差遣我報告:被擄的得釋放,瞎眼的得看見,叫那受壓制的得自由,

4:19報告神悅納人的禧年.

4:20於是把書捲起來,交還執事,就坐下.會堂裡的人都定睛看他.

4:21耶穌對他們說:「今天這經應驗在你們耳中了.」

4:22眾人都稱讚他,並希奇他口中所出的恩言

那一次,耶穌在拿撒勒的會堂裡讀經,得到眾人的稱讚.其實,耶穌這次是載譽而歸,因為他回拿撒勒之前,名聲已經傳遍加利利,拿撒勒的鄉親當然也有所聞.

4:14耶穌滿有聖靈的能力,回到加利利；他的名聲就傳遍了四方.

4:15他在各會堂裡教訓人,眾人都稱讚他.

耶穌的名聲傳遍四方,這名聲可以維持多久呢?眾人都稱讚他,這稱讚又會持續至何時呢?

「可以紅幾耐?」這是網上討論區討論Mirror的題目.有人說:我會永遠記得他們；有人說:三兩年；又有人說:與四大天王相比,簡直不入流,根本不知道他們為甚麼會紅,就算紅,都只會是曇花一現.

「可以紅幾耐?」他們自己又怎樣說呢?

「將來是好或是壞,根本無法控制.但我地可以控制的是,此刻有沒有盡力去做想做的事.」

然後旁白說:

「將來可以紅幾耐,行幾遠,解散會否,無人知,阿祖叫佢地好好珍惜,好好享受當下.壞與更壞ge 現實社會下,對於我地,亦是如此.」

「捧得紅你都可以踩低你」

男團MIRROR出道以來獲大班粉絲追捧,當中以姜濤人氣最勁,在Instagram擁有6位數字追隨者,平日去到邊都被大批「姜糖」包圍.不過,追捧的熱鬧的背後,存在著暗潮.某天,有網民在討論區開Post,貼出姜濤在Telegram(TG)群組發給粉絲的長文.姜濤在文中表示粉絲的「私追」行為影響了拍攝工作,同時透露曾有粉絲上樓敲門,對其家人造成極大困擾和擔憂.姜濤對粉絲的指責,令部分「姜糖」非常傷心,即時在社交網取消關注偶像.有粉絲指出,姜濤走紅全靠一班「姜糖」支持,並揚言「姜糖捧得紅你,都可以踩低你」.

耶穌可以受歡迎幾耐呢?

4:28會堂裡的人聽見這話,都怒氣滿胸,

4:29就起來攆他出城(他們的城造在山上)；他們帶他到山崖,要把他推下去.

「姜糖」說,「姜糖捧得紅你,都可以踩低你」.借用這句話來描述當年耶穌的處境,可以這樣說:群眾捧得紅耶穌,都可以踩低耶穌.「姜糖」不滿,是因為姜濤發文指責他們的追星行為影響到工作和家人；拿撒勒的鄉親當下反面,所為何事呢?到底發生了甚麼事,令鄉親們的態度馬上逆轉呢?

大家都知道,人們對載譽歸來的自己人是有期望的.如果有位中大校友成為了歌星,回校參加校慶活動,不是唱首歌助慶也不肯吧?如果有人發了達衣錦還鄉,鄉親們總應該比外人有更多好處吧!

京東集團的劉東強,把生意直接帶回了老家,在家鄉創造了2萬得就業崗位,和近20億得稅收!亞里巴巴的馬雲發達之後,捐錢到哪裡呢?當然是他的家鄉杭州.馬雲對杭州的捐款巨大,杭州市政府在亞里巴巴集團成立20周年的時候,舉行授予馬雲“功勳杭州人”榮譽稱號儀式.至於馬雲後來的遭遇,那是後話了.

耶穌在外面趕鬼醫病傳道,造福遠近城鄉,現在回到拿撒勒,鄉親們說:「這不是約瑟的兒子嗎?」按載譽歸來的理解,鄉親們這句話不是貶低耶穌,而是指出耶穌是自己人,鄉親們對耶穌作為自己人有甚麼期望呢?

4:23耶穌對他們說:「你們必引這俗語向我說:『醫生,你醫治自己吧!我們聽見你在迦百農所行的事,也當行在你自己家鄉裡』」.

這是合理期望啊!既然你有能力幫助外人,更應該幫助自己人.現在,我們終於等到你回來了,你可以開始施展能力了.

怎知,耶穌斷言拒絕.他舉出兩個例子,表明「益外人,唔益自己人」.第一個例子:

4:25我對你們說實話,當以利亞的時候,天閉塞了三年零六個月,遍地有大饑荒,那時,以色列中有許多寡婦,

4:26以利亞並沒有奉差往她們一個人那裡去,只奉差往西頓的撒勒法一個寡婦那裡去.

第二個例子

4:27先知以利沙的時候,以色列中有許多長大痲瘋的,但內中除了敘利亞國的乃縵,沒有一個得潔淨的.」

如果我們把自己代入當時在現場耶穌鄉親的角色,相信你會體會到他們的反應.這個世侄不但亳無血濃於水的鄉親,連基本禮數也沒有,用我們的俗語說,是「手指拗出唔拗入」.「手指拗出唔拗入」,從古至今,從東到西,特別在某些敏感時刻,很容易引起反應,甚至會被批評和攻擊.

做自己想做的事

回到人們評論Mirror的話題:「可以紅幾耐?」他們的回答是:「將來是好或是壞,根本無法控制.但我地可以控制的是,此刻有沒有盡力去做想做的事.」

耶穌如何回答呢?

4:18主的靈在我身上,因為他用膏膏我,叫我傳福音給貧窮的人；差遣我報告:被擄的得釋放,瞎眼的得看見,叫那受壓制的得自由,

4:19報告神悅納人的禧年.

演藝人當然努力維持自己可以繼續「紅」下去,但「可以紅幾耐?」的確無法控制,因此,只能「盡力去做想做的事」.對耶穌來說,他的名聲可以傳播多遠,稱讚可以維持多久,都不是他要經營的,他知道自己人生在世所為所事,他活著有使命,有方向.嚴格上來說,他不是Mirror所說的「盡力去做想做的事」,他是盡力去做他的天父交託給他的事.他有沒有自己想的事呢?不知道,如果有,都統攝在他被捕前的禱告之中:「父啊!你若願意,求你將這杯撤去；然而,不是照我的意願,而是要成全你的旨意.」(路加福音22:42)

上帝的子民應該做的事

耶穌由被捧紅,到被踩低這現象,有兩點啟發值得教會及基督徒留意:

第一,一般的教導就如Mirror的領悟一樣,既然我們不能控制其他人歡迎自己與否,所以就「盡力去做想做的事」吧!「盡力去做想做的事」未必是教會和基督徒的追求.教會存在的意義,以及教會應該如何存在,不是教會自己想怎樣就怎樣,而是學效耶穌基督的榜樣,做天父所交託的事.基督徒也如是,保羅在哥林多後書第5章14-15說:「原來基督的愛激勵我們；因我們這樣斷定,一人既替眾人死,眾人就都死了；並且他替眾人死,是叫那些活著的人不再為自己活,乃為替他們死而復活的主活.」

原來,不少基督徒都知道要做天父所交託的事,有些認真的基督徒,更因為擔心自己每一天的生活不是在聽從天父的心意.於是,同學們計劃用一次團契聚會來討論「上帝的心意」和「自己的心意」兩者互不相容的問題.我好奇地問他們,具體上是甚麼事情令你們懷疑自己不是在聽從上帝的心意呢?他們回答說,例如升大學時讀甚麼科目?爭取好成績,「靚CV」可以嗎?應不應該和這個人拍拖 ?畢業了,繼續讀書還是找工作?如果讀書,留在本地還是出外?讀甚麼?如果工作,找甚麼工作……如果我所選的,不是上帝所要我選的,那麼,我不是違背了上帝的心意嗎?

我一方面欣賞他們的虔誠,同時也挑戰他們:如果是這些事情,把「上帝的心意」和「自己的心意」對立,會不會是偽命題呢?為甚麼是偽命題呢?若這個命題要成立,就要在剛才提及的事情確立有「上帝的心意」和「自己的心意」兩樣東西,果真有這兩樣東西存在嗎?在這些事情上,你們或許感覺到有所謂「自己的心意」,然而,也在這些事情上,有與「自己的心意」相對的「上帝的心意」嗎?就算有,你們如何得知,如何肯定呢?如果無從得知或無從肯定,那麼,探究「上帝的心意」和「自己的心意」在剛才提及的人生處境中的對立,就未必有現實意義了.同學們感到很驚訝,問:那麼,我們做基督徒不需要追求上帝的心意嗎?聽到這裡,你會如何回答這個問題呢?

經過一輪討論之後,同學們肯定作為基督徒仍是要追求上帝的心意,不過,不再擔心自己選科時會否違背上帝的心意,也不再擔心物色拍拖對象時會否違背上帝的心意,也不再擔心選擇工作時會否違背上帝的心意,因為他們想到聖經上也有人問類似的問題:

彌迦書6:6「我朝見耶和華,在至高神面前跪拜,當獻上甚麼呢?」

先知回答:

8世人哪,耶和華已指示你何為善.他向你所要的是甚麼呢?只要你行公義,好憐憫,存謙卑的心與你的神同行.

一般來說,我們不用擔心畢業後所選擇的工作有否違背上帝的心意,要擔心的是我們在工作中是向著「行公義,好憐憫」的方向,還是剝削和壓制的方向；要擔心的是在工作中「存謙卑的心與你的神同行」,還是目中無神,也目中無人,損人利己地追求滿足一己的野心.一位學生努力讀書,考入醫學院,為了想實踐濟世為懷的抱負；另一位學生也努力讀書,考入醫學院,為了名成利就,傲視同儕.分別不在於讀書的態度,也不在於所選的工作,而在於所追求的是甚麼.最後,同學們在團契聚會中辦了一次人生拍賣會,讓大家透過拍賣活動探討甚麼值得追求,甚至才有價值.

第二,從耶穌的遭遇來看,當教會及基督徒個人在地上盡力活出使命的時候,也不一定不受歡迎的.現實是,可能受歡迎,也可能不受歡迎；可能這時受歡迎,可能那時不受歡迎.甚麼時候受歡迎呢?就是世界的價值與信仰的價值一致的時候,例如當社會都認為要善待難民的時候,教會及教會機構接待難民,人們就會稱讚教會和基督徒有愛心.但當社會對難民的態度改變,排斥難民的時候,教會繼續幫助難民,就會被批評了.

世界的趨勢是會改變.大家還記得,不久之前世界曾經大力推動全球化,世界上每一個人都是地球村的一份子,人們高唱We are the world, We are the children．歌曲尚未被遺忘,大家己經說要自己優先,要再強大了.大家都聽過德蘭修女在印度的服務及建立的「仁愛傳教修女會」,濟弱扶傾,照顧社會上最有需要的人,曾經得到不少印度民眾的稱讚.但幾個月前,印度政府運用「國外捐款管理法」以避免外部勢力影響內部安全為由,取消「仁愛傳教修女會」接收國外捐款的資格.經過交涉之後,「仁愛傳教修女會」恢復接收國外捐款的資格,但同時增力了限制.

世界改變之快,教會如果要追,恐怕永遠都只會落後於形勢.如果教會要得到世界的稱讚,世界可以把教會捧紅,也可以把教會踩低,令教會存在得非常艱難.

三.總結

一時的光輝未必是永恆,教會可以如何自處呢?保羅在哥林多前書13章的教導值得我們反省和學習.

13:1我若能說萬人的方言,並天使的話語,卻沒有愛,我就成了鳴的鑼,響的鈸一般.

13:2我若有先知講道之能,也明白各樣的奧祕,各樣的知識,而且有全備的信,叫我能夠移山,卻沒有愛,我就算不得甚麼.

13:3我若將所有的賙濟窮人,又捨己身叫人焚燒,卻沒有愛,仍然與我無益.

曾經何時,一眾教會都以既羨慕又批評的目光望著那所晶瑩剔透,崇拜轉播到全世界的教會；曾經何時,一眾基督徒都以既感恩又熱情地期待那位躋身國際流行樂壇的師母；曾經何時,一眾教牧都既焦慮又渴求地去到那所人多財大地廣的教會朝聖,今日,他們的狀況如何呢?那些曾經的成就,又算得了甚麼呢?如果輝煌,受歡迎,業績都不是教會的追求,那麼,教會應該追求甚麼呢?保羅繼續說:

13:4愛是恆久忍耐,又有恩慈；愛是不嫉妒；愛是不自誇,不張狂,

13:5不做害羞的事,不求自己的益處,不輕易發怒,不計算人的惡,

13:6不喜歡不義,只喜歡真理；

13:7凡事包容,凡事相信,凡事盼望,凡事忍耐.

13:8愛是永不止息.

雖然一時的光輝未必是永恆,但使徒保羅提醒我們,世間上仍有永恆的事物,那就是愛.



\newpage

\section{來跟從耶穌}
\label{sec:586}
\textbf{高國雄牧師 - 來!跟從耶穌}
\newline
\newline
\hyperref[sec:585]{< < < PREV SERMON < < <}
~
\hyperref[sec:toc]{[返目錄]}
~
\hyperref[sec:587]{> > > NEXT SERMON > > >}


路加福音 5:1-11
\newline
\begin{longtable}{cl}
\hline
\hline
章節 & 經文 (和合本修訂版)\\
\hline
5:1 & \begin{tabularx}{0.7\textwidth}{X} 耶穌站在革尼撒勒湖邊，眾人擁擠他，要聽神的道。 \end{tabularx} \\ \\ \relax
5:2 & \begin{tabularx}{0.7\textwidth}{X} 他見有兩隻船靠在湖邊，打魚的人卻離開船，洗網去了。 \end{tabularx} \\ \\ \relax
5:3 & \begin{tabularx}{0.7\textwidth}{X} 有一隻船是西門的，耶穌就上去，請他把船撐開，稍微離岸，就坐下，在船上教導眾人。 \end{tabularx} \\ \\ \relax
5:4 & \begin{tabularx}{0.7\textwidth}{X} 他講完了，對西門說：「把船開到水深的地方下網打魚。」 \end{tabularx} \\ \\ \relax
5:5 & \begin{tabularx}{0.7\textwidth}{X} 西門說：「老師，我們整夜勞累，並沒有打著甚麼。但依從你的話，我就下網。」 \end{tabularx} \\ \\ \relax
5:6 & \begin{tabularx}{0.7\textwidth}{X} 他們下了網，圈住許多魚，網險些裂開， \end{tabularx} \\ \\ \relax
5:7 & \begin{tabularx}{0.7\textwidth}{X} 就招手叫另一隻船上的同伴來幫助。他們就來，把魚裝滿了兩隻船，船甚至要沉下去。 \end{tabularx} \\ \\ \relax
5:8 & \begin{tabularx}{0.7\textwidth}{X} 西門‧彼得看見，就俯伏在耶穌膝前，說：「主啊，離開我，我是個罪人。」 \end{tabularx} \\ \\ \relax
5:9 & \begin{tabularx}{0.7\textwidth}{X} 他和一切跟他一起的人對打到了這一網的魚都很驚訝。 \end{tabularx} \\ \\ \relax
5:10 & \begin{tabularx}{0.7\textwidth}{X} 他的夥伴西庇太的兒子雅各、約翰，也是這樣。耶穌對西門說：「不要怕！從今以後，你要得人了。」 \end{tabularx} \\ \\ \relax
5:11 & \begin{tabularx}{0.7\textwidth}{X} 他們把兩隻船靠了岸，就撇下所有的，跟從了耶穌。 \end{tabularx} \\ \\
[1ex]
\hline
\hline
\end{longtable}
講題:來!跟從耶穌　Come, follow Jesus!

講員:高國雄牧師

經文:路加福音5:1-11

場合:香港中文大學崇基學院禮拜堂主日崇拜

日期:2022年2月6日

近日最火紅的女子組合collar,成為很多人追的新星

牧師朋友,在fb差不多每天發文支持,亦同時分享牧師領也叫collar,與教友科普一下,同時在娛樂中輕鬆一下,尋找到自己喜愛的,也是非常好的事；不過同時又會有朋友覺得疑惑,私下亦有不少討論,到底信徒追星有沒有問題；本人也欣賞Mirror,特別是當中姜濤的勵志故事,我也會花時間去留心,聽下新歌,欣賞舞蹈,相信只要沒有把任何人或事高舉,其位置過於基督,就可以了.

剛過去的初三,在電視的新聞中見到有人按傳統會到沙田車公廟拜神祈福,也有接受記者訪問,當中一位被訪者的回應令我非常錯愕.

市民:「我來了十年,雖然自己信耶穌,我初一至初十當做善事,奉獻給車公,平時多去D 教堂okge .」[1] Now tv善信趁初三赤口到車公廟參拜,有人祈求疫情減退\url{https://news.now.com/home/local/player?newsId=465270}

這到底是什麼的價值觀呢?我們也是否大概如此呢?

信仰的要求到底是怎樣的呢?我們又怎樣認識耶穌?如何跟從他的呢?

路加福音5:1-11

1耶穌站在革尼撒勒湖邊,眾人擁擠他,要聽上帝的道. 2他見有兩隻船灣在湖邊；打魚的人卻離開船洗網去了. 3有一隻船是西門的,耶穌就上去,請他把船撐開,稍微離岸,就坐下,從船上教訓眾人. 4講完了,對西門說:「把船開到水深之處,下網打魚.」 5西門說:「夫子,我們整夜勞力,並沒有打著甚麼.但依從你的話,我就下網.」 6他們下了網,就圈住許多魚,網險些裂開, 7便招呼那隻船上的同伴來幫助.他們就來,把魚裝滿了兩隻船,甚至船要沉下去. 8西門‧彼得看見,就俯伏在耶穌膝前,說:「主啊,離開我,我是個罪人!」 9他和一切同在的人都驚訝這一網所打的魚. 10他的夥伴西庇太的兒子雅各,約翰,也是這樣.耶穌對西門說:「不要怕!從今以後,你要得人了.」 11他們把兩隻船攏了岸,就撇下所有的,跟從了耶穌.

故事開始由眾人擁擠耶穌開始:

擁擠的目的:聽上帝的道

你追求信仰,你來到教會,目的是什麼呢?

你參與崇拜,你接受洗禮,尋求的又是什麼呢?

為什麼眾人,認為耶穌能宣講上帝的道呢?

路4:14-15耶穌滿有聖靈的能力,回到加利利；他的名聲就傳遍了四方. 他在各會堂裡教訓人,眾人都稱讚他.

4:31-37耶穌趕鬼

v37於是耶穌的名聲傳遍了周圍地方.

4:38-41耶穌治病

v40日落的時候,凡有病人的,不論害甚麼病,都帶到耶穌那裡.耶穌按手在他們各人身上,醫好他們.

眾人看見耶穌的作為,所以認定他所宣講的就是上帝的道,所以擁擠他,希望可以領受.

我們今天是否同樣如此呢?

你看到基督的作為嗎?他在你身上顯過神蹟嗎?你有期望過領受上帝的道嗎?

當眾人在擁擠耶穌,要聽上帝的道的時間,主角西門與及同伴,卻與眾人不一樣,在旁邊洗網(v2),聆聽上帝的道不重要嗎?為何有這與眾不同的選擇呢?

自己的事(搵食的事)比起上帝的道更重要嗎?我們的選擇又是怎樣的呢?

或許這是漁夫的習慣,一面洗網, 一面聆聽上帝的話語,也可以吧…?

耶穌沒有放過洗網中的西門

借船形成講道台,為要得著西門

先教導眾人,滿足眾人的需要

道講完了,精彩故事才正式開始

4講完了,對西門說:「把船開到水深之處,下網打魚.」 5西門說:「夫子,我們整夜勞力,並沒有打著甚麼.但依從你的話,我就下網.」

耶穌向西門發出了一個極奇怪的要求

今天如若耶穌同樣向你發出這要求,你的反應是什麼呢?你又會如何回應他呢?

我們又嘗試想像一下V5當時西門的語氣,想法及目的是什麼呢?他這番話的對象真的只是耶穌嗎?

無論如何,最終西門用行動回應了,這是很重要的.

請謹記

馬太福音21:28-31

又說:「一個人有兩個兒子.他來對大兒子說:『我兒,你今天到葡萄園裡去做工.』 他回答說:『我不去』,以後自己懊悔,就去了. 又來對小兒子也是這樣說.他回答說:『父啊,我去』,他卻不去.1你們想,這兩個兒子是哪一個遵行父命呢?」

行動比一切重要,行動帶來生命的改變,更帶來意想不到的結果

V6-7他們下了網,就圈住許多魚,網險些裂開, 便招呼那隻船上的同伴來幫助.他們就來,把魚裝滿了兩隻船,甚至船要沉下去.

革尼撒勒湖是西門打魚的地盤,相信他也是打魚專家,應該非常熟悉在那裡,甚麼時間打魚最有效益

對比於V5整夜勞力,並沒有打著甚麼

這是超常理/常識/自然的神蹟

我們常常忽略/忘記問一個很重要的問題:耶穌為什麼行這神蹟?

神蹟一定是為人而行的,但不一定是為滿足人的需要而行的.

西門需要打很多魚嗎?這答案同時是「對」及「不對」

「對」是因為他是漁夫,打魚為業的人,能夠多打一些魚獲,當然是很重要的事,也可能是他的人生目標/每天工作的目的.這也在他整夜打不到魚後的心情反應可以想像得到,也非常合情合理.

但同時又「不對」,是因為當時他的需要明顯不是如此,不止不是西門當下的期望,甚至他從來沒有想過會如此發生,正如剛剛我們所想像的情況,很大可能根本上西門完全沒有信過耶穌有能力行這事.

西門需要這神蹟嗎?答案是肯定的

如果神蹟只單為滿足人當下的需要,其意義就只會成為人的工具:人有當下的需要,神蹟發生；人沒有當下的需要/需要已經被滿足,神蹟不會發生；當神的作為單為人自身的需要服務,神就成為人的奴僕,這只是阿拉丁神燈中的「燈神」,而不是我們所信奉的上帝.

西門的回應,反應出這神蹟對他的實在意義:

V8西門‧彼得看見,就俯伏在耶穌膝前,說:「主啊,離開我,我是個罪人!」

俯伏,完全的降服,投降,接受自己的限制/軟弱

這神蹟對西門真正的價值,神蹟就是要讓人看見上帝的作為,這作為單為了彰顯上帝的榮耀,使人歸榮耀給上帝

由V5「夫子(老師)」到V8「主」,特顯出西門對耶穌-上帝的身份的發現/認知.

路4:38-39 耶穌醫好了西門的岳母

耶穌滿足了西門一家當下的需要(耶穌願意),但與這打魚帶來的影響相差很遠.

最大的神蹟是生命的改變,生命的轉向

這神蹟使西門發現自己(是個罪人),發現耶穌(是主),這發現震撼了他整個生命,亦因此改變了他整個人生

這神蹟帶來了生命的選擇

繼續如過往一樣?打魚洗網打魚洗網,間中借船給耶穌講道用?甚至邀請耶穌加入成為船隊的一員?

耶穌呼召的是你,他最想得著的是你

你會如何作選擇呢?如何回應主在你身上所行的神蹟呢?



\newpage

\section{無形事物的有形記號}
\label{sec:587}
\textbf{林豪恩先生 - 無形事物的有形記號}
\newline
\newline
\hyperref[sec:586]{< < < PREV SERMON < < <}
~
\hyperref[sec:toc]{[返目錄]}
~
\hyperref[sec:588]{> > > NEXT SERMON > > >}


耶利米書 17:5-10
\newline
\begin{longtable}{cl}
\hline
\hline
章節 & 經文 (和合本修訂版)\\
\hline
17:5 & \begin{tabularx}{0.7\textwidth}{X} 耶和華如此說：「倚靠人，以血肉為膀臂，心中離棄耶和華的，那人該受詛咒！ \end{tabularx} \\ \\ \relax
17:6 & \begin{tabularx}{0.7\textwidth}{X} 他必像沙漠裡的矮樹，不見福樂來到；他要住在曠野乾旱之處，無人居住的鹽地。 \end{tabularx} \\ \\ \relax
17:7 & \begin{tabularx}{0.7\textwidth}{X} 倚靠耶和華、以耶和華為他所仰賴的，那人有福了！ \end{tabularx} \\ \\ \relax
17:8 & \begin{tabularx}{0.7\textwidth}{X} 他必像樹栽於水旁，在河邊扎根，炎熱來到，毫不察覺，葉子仍必青翠；在乾旱之年，一無掛慮，並且結果不止。 \end{tabularx} \\ \\ \relax
17:9 & \begin{tabularx}{0.7\textwidth}{X} 「人心比萬物都詭詐，壞到極處，誰能識透呢？ \end{tabularx} \\ \\ \relax
17:10 & \begin{tabularx}{0.7\textwidth}{X} 我－耶和華是鑒察人心，考驗人肺腑的，要按各人所行的和他做事的結果報應他。」 \end{tabularx} \\ \\
[1ex]
\hline
\hline
\end{longtable}
講題:無形事物的有形記號　A Visible Sign of the invisible

講員:林豪恩先生

經文:耶利米書17:5-10

場合:香港中文大學崇基學院禮拜堂主日崇拜

日期:2022年2月13日

一.引言

在日常生活中,充滿著以有形的東西作為無形的事物為記號的例子.明天是情人節,如果你的情人要求你拿出對他的愛來看看,你會如何回應呢?你能夠把愛好像一件物件一樣拿出來觀看嗎?會不會你能夠拿得到出來的,只是代表愛的東西呢?

婚禮進行時,通常有一個交換戒指的環節.牧師手持戒指,解釋戒指的意義時會這樣說:「這對戒指是圓的,代表沒有結束的永恆,它的材質是黃金做成的,代表純淨,聖潔的.當你們彼此交換戒指時,我們祈禱,你們的愛也天長地久,純淨以及聖潔.願神藉著這個婚姻的記號,賜福你們的婚姻.」有時,新郎和新娘也彼此說:「我將這戒指戴在你的手指上,從今天開始,直到永遠,做為我們婚姻的記號.戴著這戒指表明我們今天在上帝及眾人面前承諾的婚姻約定.」由此可見,無論是主持婚禮的牧師還是結婚的新郎新婚,對這對戒指都有共識,那是一個婚姻的記號.承諾,永恆,純淨,聖潔是無形無體不能觸摸的,但戒指這個記號則是有形有體可觸摸的.看見這個有形有體可觸摸的戒指,就令人想起它代表的婚姻中的承諾,永恆,純淨和聖潔.

著名教父奧古斯丁將聖禮界定為神聖事物的可見記號 (a visible sign of a sacred thing) 和不可見的恩典的可見形式 (a visible sign of an invisible grace),例如聖餐中的餅與酒,就是指向基督的身體和血這神聖事物的標記.信仰之中充滿著類似的表述,以看得見的事物來表達看不見的境界,今天的經課就有一些例子.

二.經課及思考

1.耶利米書17:5-10

耶利米書17章5節開始由敘事文變成詩歌,我們今天讀的5-10是這首詩的其中一段.古今中外,詩歌是最常以有形可見的實體東西來表達無形無體的抽象情感或人生道理的文學體裁.或借景抒情,或以物寓意.讓我們一齊來再讀這首詩歌.

17:5耶和華如此說:倚靠人血肉的膀臂,心中離棄耶和華的,那人有禍了!

17:6因他必像沙漠的杜松,不見福樂來到,卻要住曠野乾旱之處,無人居住的鹼地.

17:7倚靠耶和華,以耶和華為可靠的,那人有福了!

17:8他必像樹栽於水旁,在河邊扎根,炎熱來到,並不懼怕,葉子仍必青翠,在乾旱之年毫無掛慮,而且結果不止.

17:9人心比萬物都詭詐,壞到極處,誰能識透呢?

17:10我－耶和華是鑒察人心,試驗人肺腑的,要照各人所行的和他做事的結果報應他.

詩歌也是由文字組成,詩歌當中也有很多細節.或許我們習慣咬文嚼字地閱讀聖經,錙銖必較,因為我們被教導聖經中每一個字都是有神的啟示.然而,如果我們以同樣的態度閱讀詩歌,就有機會弄巧反拙,捉錯用神.例如讀李白的靜夜詩:

床前明月光

疑是地上霜

舉頭望明月

低頭思故鄉

學生:老師,請問第一個字「床」是甚麼床?

老師:是甚麼床不重要,這首詩的重點是……

學生:老師,不能說不重要,這是全首詩的第一個字,連第一個字都沒有搞清楚,你如何說不去呢?

老師:那麼,應該是木床,因為李白那個年代的床應該是用木造的.

學生:老師,那麼,即是說你也不肯定.算了,就當是木床.那麼,這張床的尺寸如何呢?甚麼顏色呢?放在甚麼位置呢……

如果李白看到後人如此閱讀他的靜夜詩,不知應該生氣還是發笑.

讓我們回到耶利米的詩歌:

17:5耶和華如此說:倚靠人血肉的膀臂,心中離棄耶和華的,那人有禍了!

學生:老師,第一句耶穌華「說」是如何說?用甚麼語言說?是有物理學上的聲波震動嗎?

老師:你問的問題都很有意思.但未必是這首詩歌的重點,不如我們讓詩歌本身帶領我們閱讀,看看有甚麼發現,好嗎?

若我們繼續讀下去,我們會發現詩歌引導讀者問的問題很可能在最後一句:「那人有禍了」,於是問的問題應該是:有禍是怎樣的呢?因為下文按這主題發展.

17:6因他必像沙漠的杜松,不見福樂來到,卻要住曠野乾旱之處,無人居住的鹼地.

禍並不是可以拿出來觀看和觸摸的東西,於是,詩歌借用有形有體可見的事物來表述.詩歌描繪一幅圖畫,圖畫的主角是一棵矮樹,在荒漠曠野,乾旱無水,土壤貧瘠.詩歌邀請讀者想像自己就是那棵矮樹,置身於如此那般的環境,會有甚麼感受呢?苦悶?缺乏?孤獨?隔絕?無盼望?詩歌就是要留下空間給讀者去想像,有禍就是如此惡劣的狀況.

接著,詩歌以相反的狀況作為對比:

17:7倚靠耶和華,以耶和華為可靠的,那人有福了!

按詩歌結構和脈絡,應該暗示讀者如此問:有福又是怎樣的呢?福也不是一件可以拿出來觀看和觸摸的東西,於是,詩歌又借用有形有體可見的事物來表述.

17:8他必像樹栽於水旁,在河邊扎根,炎熱來到,並不懼怕,葉子仍必青翠,在乾旱之年毫無掛慮,而且結果不止.

詩歌展示另一幅圖畫,主角又是一棵樹,卻處身於不同的環境.今次讀者見到水源不絕的河流,聽到流水淙淙的聽音,看見根深葉茂的樹.再細看,有青翠的葉子,豐盈的果實.詩歌又邀請讀者想像自己就是那棵栽圖畫中的樹,有甚麼感受呢?豐富供應?受滋養?無憂無慮?有成果?詩歌也留下空間給讀者去想像,有福是何等舒暢的狀況.

9-10節是一個小結:

詩人問:

17:9人心比萬物都詭詐,壞到極處,誰能識透呢?

因為人心詭詐,所以正邪難辯,正如Killer遊戲(狼人殺)一樣,殺手盡力假扮天使,令你死了也不知道死在誰手上.耶和華會不會被人欺騙呢?會不會把禍福錯配呢?這的確是令人擔心的.詩歌如何回應呢?

17:10我－耶和華是鑒察人心,試驗人肺腑的,要照各人所行的和他做事的結果報應他.

詩人以耶和華鑒察人心,試驗人肺腑作為保證,不會被欺騙,也不會錯配禍福,為這個小段作結.詩歌讀完了,不過,又未是離開的時候.詩歌最重要的是詩境,如果讀者能夠投入詩境,穿梭兩幅圖畫感受其意境,並停留在小結所展示的耶和華的保障,讓那份保障進入心中,讓心安妥,帶著這份安妥離開詩境,進入生活.

2.詩篇第1篇

無獨有偶,詩篇第一篇又是一首類似的詩歌,只是禍福的次序倒轉:

1:1不從惡人的計謀,不站罪人的道路,不坐褻慢人的座位,

1:2惟喜愛耶和華的律法,晝夜思想,這人便為有福!

有了讀上一首耶利米的詩歌的經驗,我們可以假設,詩篇第一篇的重點都是最後那句「這人便為有福」,暗示讀者應該問的問題是:有福是怎樣的呢?讀下去,就知道我們有沒有猜錯.

1:3他要像一棵樹栽在溪水旁,按時候結果子,葉子也不枯乾.凡他所做的盡都順利.

是熟識的畫面,耶利米詩歌的樹似乎又再出現.有水源,有茂葉,有果子.詩境相似,讀者的感受也相似,豐富供應,受滋養,無憂無慮,有成果.

接著,詩歌展示相反的狀況:

1:4惡人並不是這樣,

詩歌又暗示讀者問:不是這樣,惡人又是怎樣呢?

1:4乃像糠秕被風吹散.

今次的圖畫沒有出現另一棵樹,而是被風吹散的糠秕.詩歌需要讀者運用想像力,想像在打穀的過程中,風把不要的糠秕吹走.又邀請讀者投入感情,想像自己就是那在打穀場上被風吹散的糠秕,被放棄,被視為垃圾.

5-6節是總結:

1:5因此,當審判的時候,惡人必站立不住；罪人在義人的會中也是如此.

好像上一首詩一樣,提示讀者問:怎樣知道呢?怎能肯定呢?詩人也是以耶和華「知道」來回應:

1:6因為耶和華知道義人的道路；惡人的道路卻必滅亡.

詩歌讀完了,不過,又未是離開的時候.你進入了這首詩歌的詩境了嗎?詩歌邀請你感受栽在溪水旁的樹木的滋養和安穩,也邀請你感受被風吹散的糠秕的飄泊流落,直至感受到義人與惡人的不同結局,選擇去惡近義,就能夠與詩人共融,同感一靈了.

3.路加福音

詩境能夠啟發想像,滋潤心靈.不過,人並不只是活在精神世界,卻是活在具體的物質世界.在福音經課中,耶穌對人們的具體處境和生存需要作出回應:

第一項對比

6:20「貧窮的人有福了!因為神的國是你們的.

6:24但你們富足的人有禍了!因為你們已經受過安慰.

第二項對比

6:21現在飢餓的人有福了!因為你們將得飽足.

6:25你們現在飽足的人有禍了!因為你們將要飢餓.

第三項對比

現在哭泣的人有福了!因為你們將要歡笑.

你們現在歡笑的人有禍了!因為你們將要哀慟哭泣.

第四項對比

6:22人為人子的緣故憎恨你們,拒絕你們,辱罵你們,把你們當惡人除掉你們的名,你們就有福了!

6:23在那日,你們要歡欣雀躍,因為你們在天上的賞賜是很多的；他們的祖宗也是這樣待先知的.

6:26人都說你們好的時候,你們有禍了!因為他們的祖宗也是這樣待假先知的.」

耶穌直接提及人們具體的狀況,例如溫飽,哭笑,獎罰.然而,都是只是應許,即是不在當下出現.在當下體會到的,看得見的是:飢餓的人手中仍然沒有食物,飽足的人台上仍然有吃不完的；哭泣的人抹到自己的眼淚,歡笑的人聽到自己的笑聲.這看不見,摸不到的應許,其記號又是甚麼呢?耶穌又留下甚麼給人們作為保證呢?

6:17耶穌和他們下了山,站在一塊平地上；在一起的有許多門徒,又有許多百姓從全猶太和耶路撒冷,並推羅,西頓的海邊來,

6:18都要聽他講道,又希望耶穌醫治他們的病；還有被污靈纏磨的,也得了醫治.

6:19眾人都想要摸他,因為有能力從他身上發出來,治好了他們.

耶穌對自己的應許的保證是讓人們預嘗應許實現的滋味,人們可以摸到耶穌身上發出來的能力,可以見到折磨人的疾病得醫治,見到纏磨人的污靈被驅趕.天國不只是詩境,而是可觸摸可看見的實體；天國雖然尚未完全實現,但可以初嘗.

教會延續耶穌的使命,為天國作見證.教會沒有能力把天國實現,但教會可以讓人們預嘗天國的滋味,而不只是宣揚一個死後的世界.

4.哥林多前書

基督教信仰雖然有很多無形事物的有形記號,但並不是「本來無一物,何處惹塵埃」的宗教,基督教信仰終究是有形有體的,使徒保羅曾為此努力爭辯.

15:12既傳基督是從死裡復活了,怎麼在你們中間有人說沒有死人復活的事呢?

15:13若沒有死人復活的事,基督也就沒有復活了.

15:14若基督沒有復活,我們所傳的便是枉然,你們所信的也是枉然；

15:15並且明顯我們是為上帝妄作見證的,因我們見證上帝是叫基督復活了.若死人真不復活,上帝也就沒有叫基督復活了.

15:16因為死人若不復活,基督也就沒有復活了.

15:17基督若沒有復活,你們的信便是徒然,你們仍在罪裡.

15:18就是在基督裡睡了的人也滅亡了.

15:19我們若靠基督只在今生有指望,就算比眾人更可憐.

15:20但基督已經從死裡復活,成為睡了之人初熟的果子.

保羅指出,基督不但在人世界生活的時候有形有體,就是死後仍是有形有體,因為復活的基督展示過形體才離開門徒.人們現在能夠觸摸,能夠看見自己有形有體,將來呢?人死如燈滅嗎?不是,基督就是人們「初熟的果子」,人們會步履基督的腳步,有形有體地復活.保羅強調,這種有形有體的信仰是獨特的,是基督徒信仰不至徒然的基礎.有形有體的信仰對生活與倫理的建立也是重要的.如果有興趣思考和探索,讓我們一齊去請教我們當中的神學教授鄧瑞強博士.



\newpage

\section{浪花淘盡英雄}
\label{sec:588}
\textbf{鄧瑞強博士 - 浪花淘盡英雄}
\newline
\newline
\hyperref[sec:587]{< < < PREV SERMON < < <}
~
\hyperref[sec:toc]{[返目錄]}
~
\hyperref[sec:589]{> > > NEXT SERMON > > >}


出埃及記 14:10-18
\newline
\begin{longtable}{cl}
\hline
\hline
章節 & 經文 (和合本修訂版)\\
\hline
14:10 & \begin{tabularx}{0.7\textwidth}{X} 法老逼近的時候，以色列人舉目，看哪，埃及人追來了，就非常懼怕，以色列人向耶和華哀求。 \end{tabularx} \\ \\ \relax
14:11 & \begin{tabularx}{0.7\textwidth}{X} 他們對摩西說：「難道埃及沒有墳地，你要把我們帶來死在曠野嗎？你為甚麼這樣待我們，將我們從埃及領出來呢？ \end{tabularx} \\ \\ \relax
14:12 & \begin{tabularx}{0.7\textwidth}{X} 我們在埃及豈沒有對你說過，不要攪擾我們，讓我們服事埃及人嗎？因為服事埃及人總比死在曠野好。」 \end{tabularx} \\ \\ \relax
14:13 & \begin{tabularx}{0.7\textwidth}{X} 摩西對百姓說：「不要怕，要站穩，看耶和華今天向你們所要施行的拯救，因為你們今天所看見的埃及人必永遠不再看見了。 \end{tabularx} \\ \\ \relax
14:14 & \begin{tabularx}{0.7\textwidth}{X} 耶和華必為你們爭戰，你們要安靜！」 \end{tabularx} \\ \\ \relax
14:15 & \begin{tabularx}{0.7\textwidth}{X} 耶和華對摩西說：「你為甚麼向我哀求呢？你吩咐以色列人往前走。 \end{tabularx} \\ \\ \relax
14:16 & \begin{tabularx}{0.7\textwidth}{X} 你舉手向海伸杖，把水分開。以色列人要下到海中，走在乾地上。 \end{tabularx} \\ \\ \relax
14:17 & \begin{tabularx}{0.7\textwidth}{X} 看哪，我要任憑埃及人的心剛硬，他們就跟著下去。我要在法老和他的全軍、戰車、戰車長身上得榮耀。 \end{tabularx} \\ \\ \relax
14:18 & \begin{tabularx}{0.7\textwidth}{X} 我在法老和他的戰車、戰車長身上得榮耀的時候，埃及人就知道我是耶和華。」 \end{tabularx} \\ \\
[1ex]
\hline
\hline
\end{longtable}
講題:浪花淘盡英雄 Waves over the Egyptians

經文:出埃及記14:10-18

講員:鄧瑞強博士

場合:香港中文大學崇基學院禮拜堂主日崇拜

日期:2022年2月20日

各位弟兄姊妹,早安.在這疫症裡,祝願各人平安.

創世記講到,在人類漫長的歷史黑夜裡,神揀選點燈的人.神揀選亞伯拉罕,亞伯拉罕也願意回應神.他放下他的過去,離開本地,本族,父家,告別過去的成敗得失,跟隨著他的天命,他向前走,但卻不知要走向何方.在大地上,一切已知的領域,都不是他嚮往的地方.他的心靈,想望著一無以名狀的世界,一個神明與人同在的世界.世人熟悉的德蘭修女也是這樣,離開自己的故鄉,去到一處被世人遺忘的窮鄉僻壤,抱起被疾病與苦難折磨到瀕臨死亡的人,擁抱他們,照料他們,讓他們體會生命真正的尊嚴.一個如此的病人說:「我過去活像街上的一頭牲畜,但在你們的愛與關懷下,我死的時候,卻像天使一樣.」(In the Heart of the World (Novato: New World Library, 1997), p.64.)在這裡,我們見到神的愛在人間帶來的改變.孔子的學生曾子說:「士不可以不弘毅,任重而道遠.仁以為己任,不亦重乎?死而後已,不亦遠乎?」他說:真正的人志向不可以不遠大,意志不可以不堅強,因為責任重大,要走的路遙遠.要將天地間的大愛實踐於世,你說,這責任重大不重大?背負此重任直至死時方休,你說,此路途遙遠不遙遠?亞伯拉罕的一生,都在這漫長的道路上.

要實踐天道,除了自己,還需要一群人.中國的孔子,透過教學,凝聚了一群人.道,在學問的傳承中.遠方的亞伯拉罕,透過生育,形成一個民族.道,在耶和華的信仰裡,不是透過教學而傳遞,而是透過生命歷史中的神人相遇,神人糾纏,神人摔交而被看見.道,在耶和華的信仰裡,是在人生旅程中一次又一次的險境裡,神的應許與人的疑惑,神的顯現與人的不信,神的鼓勵與人的放棄的種種張力中,被人經驗.在耶和華的信仰裡,道,最後是要「成肉身」的,在歷史中顯現.亞伯拉罕不能透過教學而凝聚一群人,他需要透過後裔繁衍的歷史過程,在其中的種種危機與轉折裡,看見神的真實.以色列的神,需要在歷史的曲折中被看見.對亞伯拉罕而言,最大的危機,是「無後」.就神學而言,這危機意味著在歷史的過程中,看不見神的真實.神與人交往的故事,是曲折而神奇的.亞伯拉罕的正室不生育,他便找妾侍.妾侍生了男孩,但神說,這不算是他的正式後裔.到後來,正室生了,但神說,這男孩要拿來獻祭.亞伯拉罕拿他來獻祭,神又說,不用了,因你獻祭的心已是上好的祭.亞伯拉罕的後裔的最大問題,就是如何延續後裔.這問題兜兜轉轉,重重複複.最後,一個家族成型了.後來,他們遇上荒年,失收無糧,便去仍有糧食的埃及保命.這就是出埃及記的開始了.

出埃及記1:7:「以色列人生養眾多,並且繁茂,極其強盛,滿了那地.」時光飛逝,斗轉升移.亞伯拉罕的後裔,已成了一人丁鼎盛的民族.要成為在黑暗的歷史裡點燈的一群人,他們要在歷史中與神摔交,從而看到神的真實.

到了歷史的某一刻,形勢逆轉.出埃及記1:8-9:「有不認識約瑟的新王起來,治理埃及,對他的百姓說:看哪,這以色列民比我們還多,又比我們強盛.」一個新朝代的興起,新王走上歷史的舞台.他不認識這民族的過去,他也不與他們友善,他只看到這強大的民族潛在的威脅,他要壓制他們.這個民族成為了奴隸,出生的男丁被殺死,他們的歷史前景暗淡無光.如何點燈?光明在哪裡?在絕望的時刻,希望在哪裡?原來,要成就一個民族為神的子民,歷史必需走這著棋.

希望的開端,是兩個卑微的女人,是收生婆,有名有姓,一個叫「施弗拉」,意思是「靚」,我叫她「靚姑」,另一個叫「普阿」,意思是「光彩」,我叫她「彩姨」.靚姑和彩姨有好生之德,這麼低微的人,竟敢違反法老殺男嬰的法令,偷偷保留了以色列的男嬰.

之後,一個以色列男嬰出生,他母親藏不到他,便將男嬰放在一「蒲草箱」裡,放在公主來洗澡的河邊.這「蒲草箱」的「箱」字,便是挪亞方舟的「方舟」那個字.挪亞方舟是在人生的洪濤大浪中救助人的工具,如今,這方舟,這箱,救活了一個重要人物.公主來河邊洗澡,見到這嬰孩,可憐他,收養他.天意真是深妙無窮,要被消滅的嬰孩,竟住在消滅者的深宮之中.出埃及記2:10:「公主給孩子起名叫摩西,意思說:因我把他從水裡拉出來.」摩西這個名字,意思就是「拉出來」,「救出來」的意思.他要將他的同胞從水深火熱中拉出來.

上天的佈局已完成,就留待時間去發酵.

時光飛逝,摩西長大了,因抱打不平而殺了一個埃及人,逃命到米甸.有一天,他命中注定的天命來了.

出3:2:「耶和華的使者從荊棘裡火焰中向摩西顯現.摩西觀看,不料,荊棘被火燒著,卻沒有燒毀.」

這是一團不滅的火.

神藉一團不滅的火呼喚摩西,大部分人生命中的這一團火已經熄滅.神在這團火中向摩西說話.只要生命的火還在,生命裡冥冥之中的天命仍會被喚醒.

出3:7:「耶和華(對摩西)說:我的百姓在埃及所受的困苦,我實在看見了；他們因受督工的轄制所發的哀聲,我也聽見了.……

出3:10:故此,我要打發你去見法老,使你可以將我的百姓以色列人從埃及領出來.」

神讓摩西知道,他的同胞的苦情.神讓摩西知道,上天聽到這些哀聲.然後,神讓摩西知道,在人間苦情與上天聆聽的這種天人交接中,摩西的天命,就是代表上天去回應人間的哀求.我們可能已對人間的苦情麻木,故無動於衷.就或看見人間的苦情,我們可能意會不到有一位神在聆聽,故只覺生命的無奈.就或意會到神在聆聽,我們可能意識不到與自己有什麼關係,只會讓一切聽天由命.神直接說:這與你有關,這是你的天命,這就是你存在於歷史的這個時刻的意義所在.

神在歷史中存在.祂在歷史的某個艱難時刻,借歷史的契機而說話.而剛剛,人正是要在歷史的這個艱難時刻,逃避歷史的責任.就在這神與人糾結的歷史過程中,人看見神的真實.神就在這裡,因為祂說話,因為祂呼召人,因為祂呼召人去回應歷史的重負.在生命的火中,人面對神.

摩西問神:你叫什麼名字?

出3:14:神對摩西說:「我是自有永有的」(Ehyeh-Asher-Ehyeh).

神的名字,在希伯來文來說,基本上是一「be」動詞,其意思不是「自有永有」.「自有永有」的含意,基本上已脫離了塵世的歷史,指向一無時間的永恆,或抽離時間的寂然不動.但Ehyeh-Asher-Ehyeh的含意,明確指向時間中的「be」的創生力,它是在時間中讓萬物生成的偉大力量,即中國人講的「天地之大德曰生」的那個「生」.有了這「生」的力量,才能在「無」創造「有」,在「死路」創造出「出路」,在歷史的困局中創造出歷史的生機.神不是「自有永有」的存在於某種永恆的空間裡,祂是內在地存在於我們的歷史裡,祂在時間中存在,祂在人類歷史的前面,不斷為我們的歷史創造意想不到的將來.

神回答摩西,在你的將來,在你同胞的將來,在人類的將來,全在我的創造生機中.你不用以為是你用自己的力開自己的路,不,是我以我的be為你開路,你只是be的器皿.

在這不滅的火中,摩西聆聽到神的呼喚,明白自己身處歷史某一時刻的存在意義,領悟到天命,這天命關乎同胞的苦難,關乎人類歷史永恆的主題「奴役與自由」,關乎人類的歷史前程.摩西,將成為人類歷史永恆的巨人.

摩西去見法老,希望法老讓以色列人放假,到曠野敬拜神.法老當然拒絕.於是,十種災禍臨到埃及.每一個災,針對的,都是埃及的神明.例如:埃及人敬拜尼羅河,於是,尼羅河的水變血,河神死了.埃及人敬拜蛙神,於是,青蛙成為他們的災劫.埃及人敬拜太陽神,於是,有三天的黑暗.埃及人敬拜法老,於是,埃及人未來的法老--法老長子的死亡.

終於,法老趕他們走.

法老,以及法老的神明,在聖經文學裡,代表著奴役人類的強大魔力.耶和華神,以及祂的代表摩西,正面與這魔力交鋒.奴役的力量,自由的力量,在那歷史時刻,在以色列人的心靈中相遇激盪.以色列人慣於被奴役,他們雖然叫苦連天,但從未想像過自由.自由,並不是他們的想像力能夠達到的境界.但是,在摩西的奮鬥中,在神的「be」所創造的全新將來中,他們第一次想像自由.

以色列人,400年的奴役,竟然還可得嚐自由,他們帶著行裝,浩浩蕩蕩,試著踏出自由的一步.畢竟,這是從未試過的一步.但,法老,怎會如此仁慈.他悲哀過後,定一定神,想起國家的經濟命脈,沒有了奴隸,就像今日沒有了電力,這國家如何運作下去?埃及帝國的成功,少不免有人淪為奴隸.這是成王敗寇的現實,以色列人的行動是反現實.法老想通了,發出追兵,捉拿他們.

以色列人來到紅海邊,埃及的追兵也差不多到了.前無去路,後有追兵.

這就到了今日的講道經文:出埃及記14:10-18

14:10　法老臨近的時候,以色列人舉目看見埃及人趕來,就甚懼怕,向耶和華哀求.

14:11　他們對摩西說:「難道在埃及沒有墳地,你把我們帶來死在曠野嗎?你為甚麼這樣待我們,將我們從埃及領出來呢?

14:12　我們在埃及豈沒有對你說過,不要攪擾我們,容我們服事埃及人嗎?因為服事埃及人比死在曠野還好.」

14:13　摩西對百姓說:「不要懼怕,只管站住!看耶和華今天向你們所要施行的救恩.因為,你們今天所看見的埃及人必永遠不再看見了.

14:14　耶和華必為你們爭戰；你們只管靜默,不要作聲.」

14:15　耶和華對摩西說:「你為甚麼向我哀求呢?你吩咐以色列人往前走.

14:16　你舉手向海伸杖,把水分開.以色列人要下海中走乾地.

14:17　我要使埃及人的心剛硬,他們就跟著下去.我要在法老和他的全軍,車輛,馬兵上得榮耀.

14:18　我在法老和他的車輛,馬兵上得榮耀的時候,埃及人就知道我是耶和華了.」

多麼震憾的畫面,多麼令人感動的畫面.無數手無寸鐵的人,無數裝備精良的士兵,這將是一幅什麼樣的畫面?慣作奴隸的,他們的心仍是奴隸的心,他們解決問題的方式,是做回一個奴隸.多麼淒涼!自由何等重價!摩西說,這是一場神魔大戰,是讓人自由的神與迫人為奴的魔的大戰.在這歷史時刻,在時間的深處,在「be」的根源那裡,我們會看到人類的救贖,我們會看見神.神說,向前走吧!水,會為你們分開.水,會為你們回流.走過去,就是自由.退回頭,就是奴隸.天地的主,要為人間定睛義,故此,向前走吧!

中國水墨畫家島子,畫過一幅《過紅海》的大型水墨畫,以前我在崇基校園見過,不知現在是否仍在?島子說,他畫這畫時,有的是「一種感動,一種內心很沉痛的東西,就是我們的歷史,中國的近代的歷史,我們中國還沒有過紅海,所以內心很沉重,非常沉痛的感情.……正因這個感動,所以畫了這幅《過紅海》,那個是鋪天蓋地的巨浪壓過來,霎時把法老的追兵全部都撲到底下,它是用大幅度的潑墨和破墨完成.然後,手持神杖的摩西在這邊已經上岸了,在呼召他的子民趕快走,在畫面底部用很小筆觸,用焦墨畫的.你仔細看,這個巨浪是一個字母“v”,victor.這個作品揭示的是現實經驗,也是歷史真相.是對社會的觀察,思考,對救贖的盼望,從而給出一個判斷.因為神的啟示從來都不是空穴來風.」(\url{https://www.cco.cuhk.edu.hk/chaplaincy/chi/cccfestival/2012/03/}\%E4\%B8\%AD\%E5\%9C\%8B\%E8\%97\%9D\%E8\%A1\%93\%E5\%AE\%B6\%E5\%B3\%B6\%E5\%AD\%90\%E8\%88\%87\%E4\%BC\%8D\%E6\%B8\%AD\%E6\%96\%87\%E7\%89\%A7\%E5\%B8\%AB\%E5\%B0\%8D\%E8\%AB\%87\%EF\%BC\%9A\%E5\%9F\%BA\%E7\%9D\%A3\%E6\%95\%99\%E8\%97\%9D\%E8\%A1\%93\%E8\%88\%87/)

神在歷史裡,一步一步地建立一個群體,這個群體要將光明帶來人間.神首先要做的,是讓這群人體會真正的自由.走過紅海,這群奴隸,終於呼吸到真正的自由空氣.你會走過紅海嗎?

願榮耀歸上主,願人間享和平.



\newpage

\section{戰勝魔鬼之道}
\label{sec:589}
\textbf{高國雄牧師 - 戰勝魔鬼之道}
\newline
\newline
\hyperref[sec:588]{< < < PREV SERMON < < <}
~
\hyperref[sec:toc]{[返目錄]}
~
\hyperref[sec:590]{> > > NEXT SERMON > > >}


路加福音 4:1-13
\newline
\begin{longtable}{cl}
\hline
\hline
章節 & 經文 (和合本修訂版)\\
\hline
4:1 & \begin{tabularx}{0.7\textwidth}{X} 耶穌滿有聖靈，從約旦河回來，聖靈把他引到曠野， \end{tabularx} \\ \\ \relax
4:2 & \begin{tabularx}{0.7\textwidth}{X} 四十天受魔鬼的試探。在那些日子，他沒有吃甚麼，日子滿了，他餓了。 \end{tabularx} \\ \\ \relax
4:3 & \begin{tabularx}{0.7\textwidth}{X} 魔鬼對他說：「你若是神的兒子，叫這塊石頭變成食物吧。」 \end{tabularx} \\ \\ \relax
4:4 & \begin{tabularx}{0.7\textwidth}{X} 耶穌回答說：「經上記著：『人活著，不是單靠食物。』」 \end{tabularx} \\ \\ \relax
4:5 & \begin{tabularx}{0.7\textwidth}{X} 魔鬼又領他上了高山，霎時間把天下萬國都指給他看， \end{tabularx} \\ \\ \relax
4:6 & \begin{tabularx}{0.7\textwidth}{X} 對他說：「這一切權柄和榮華我都要給你，因為這原是交給我的，我願意給誰就給誰。 \end{tabularx} \\ \\ \relax
4:7 & \begin{tabularx}{0.7\textwidth}{X} 你若在我面前下拜，這一切都歸你。」 \end{tabularx} \\ \\ \relax
4:8 & \begin{tabularx}{0.7\textwidth}{X} 耶穌回答他說：「經上記著：『要拜主—你的神，惟獨事奉他。』」 \end{tabularx} \\ \\ \relax
4:9 & \begin{tabularx}{0.7\textwidth}{X} 魔鬼又領他到耶路撒冷去，叫他站在聖殿頂上，對他說：「你若是神的兒子，從這裡跳下去！ \end{tabularx} \\ \\ \relax
4:10 & \begin{tabularx}{0.7\textwidth}{X} 因為經上記著：『主要為你命令他的使者保護你； \end{tabularx} \\ \\ \relax
4:11 & \begin{tabularx}{0.7\textwidth}{X} 他們要用手托住你，免得你的腳碰在石頭上。』」 \end{tabularx} \\ \\ \relax
4:12 & \begin{tabularx}{0.7\textwidth}{X} 耶穌回答他說：「經上說：『不可試探主—你的神。』」 \end{tabularx} \\ \\ \relax
4:13 & \begin{tabularx}{0.7\textwidth}{X} 魔鬼用完了各樣的試探，就離開耶穌，再等時機。 \end{tabularx} \\ \\
[1ex]
\hline
\hline
\end{longtable}
講題:戰勝魔鬼之道 How to defeat the devil

經文:路加福音4:1-13

講員:高國雄牧師

場合:香港中文大學崇基學院禮拜堂主日崇拜

日期:2022年3月6日

講道大綱

1耶穌被聖靈充滿,從約旦河回來,聖靈將他引到曠野, 2四十天受魔鬼的試探.那些日子沒有吃甚麼；日子滿了,他就餓了. 3魔鬼對他說:「你若是上帝的兒子,可以吩咐這塊石頭變成食物.」 4耶穌回答說:「經上記著說:『人活著不是單靠食物,乃是靠上帝口裡所出的一切話.』」 5魔鬼又領他上了高山,霎時間把天下的萬國都指給他看, 6對他說:「這一切權柄,榮華,我都要給你,因為這原是交付我的,我願意給誰就給誰. 7你若在我面前下拜,這都要歸你.」 8耶穌說:「經上記著說:當拜主－你的上帝,單要事奉他.」9魔鬼又領他到耶路撒冷去,叫他站在殿頂上,對他說:「你若是上帝的兒子,可以從這裡跳下去； 10因為經上記著說:主要為你吩咐他的使者保護你；11他們要用手托著你,免得你的腳碰在石頭上.」12耶穌對他說:「經上說:『不可試探主－你的上帝.』」 13魔鬼用完了各樣的試探,就暫時離開耶穌.

信仰疑難:試探 vs 試練

耶穌面對試探的意義

耶穌如何面對試探

得勝秘訣

試探一:身體的需要 vs 証明自己的身份

路4:2-4

那些日子沒有吃甚麼；日子滿了,他就餓了. 3魔鬼對他說:「你若是上帝的兒子,可以吩咐這塊石頭變成食物.」 4耶穌回答說:「經上記著說:『人活著不是單靠食物,乃是靠上帝口裡所出的一切話.』」

回應:

申命記8:1-5

1「我今日所吩咐你的一切誡命,你們要謹守遵行,好使你們存活,人數增多,可以進去得耶和華向你們列祖起誓應許的那地. 2你要記得,這四十年耶和華─你的上帝在曠野一路引導你,是要磨煉你,考驗你,為要知道你的心如何,是否願意遵守他的誡命. 3他磨煉你,任你飢餓,將你和你列祖所不認識的嗎哪賜給你吃,使你知道,人活著,不是單靠食物,乃是靠耶和華口裡所出的一切話. 4這四十年,你身上的衣服沒有穿破,你的腳也沒有腫. 5你心裡要知道,耶和華─你的上帝管教你,像人管教兒女一樣.

要對準焦點:飲食雖然是必須,但不可單靠!

試探二:追求世上權柄,榮華 vs 拜服,跟從上帝

路4:5-8

5魔鬼又領他上了高山,霎時間把天下的萬國都指給他看, 6對他說:「這一切權柄,榮華,我都要給你,因為這原是交付我的,我願意給誰就給誰. 7你若在我面前下拜,這都要歸你.」 8耶穌說:「經上記著說:當拜主－你的上帝,單要事奉他.」

回應:

申命記6:10-15

10「耶和華－你的上帝必領你進他向你列祖亞伯拉罕,以撒,雅各起誓要給你的地.那裡有又大又美的城鎮,不是你建造的； 11有裝滿各樣美物的房屋,不是你裝滿的；有挖成的水井,不是你挖的；有葡萄園,橄欖園,不是你栽植的；你吃了而且飽足. 12你要謹慎,免得你忘記領你從埃及地為奴之家出來的耶和華. 13你要敬畏耶和華－你的上帝,事奉他,奉他的名起誓. 14不可隨從別神,就是你們四圍民族的眾神明, 15因為在你中間的耶和華－你的上帝是忌邪的上帝,恐怕耶和華－你上帝的怒氣向你發作,把你從地上除滅.

上帝的價值觀應該是我們唯一認同的價值觀

試探三:身份的質疑,試探上帝 vs 自我實現,顯現榮耀

路4:9-12

9魔鬼又領他到耶路撒冷去,叫他站在殿頂上,對他說:「你若是上帝的兒子,可以從這裡跳下去； 10因為經上記著說:主要為你吩咐他的使者保護你；11他們要用手托著你,免得你的腳碰在石頭上.」12耶穌對他說:「經上說:『不可試探主－你的上帝.』」

詩篇91: 11-12

11因他要為你命令他的使者,在你所行的一切道路上保護你.

12他們要用手托住你,免得你的腳碰在石頭上.

回應:

申命記6:16

16「你們不可試探耶和華－你們的上帝,像你們在瑪撒那樣試探他.

出埃及記17:1-8

1以色列全會眾遵照耶和華的吩咐,從汛的曠野一段一段地往前行.他們在利非訂安營,但百姓沒有水喝. 2百姓就與摩西爭鬧,說:「給我們水喝吧!」摩西對他們說:「你們為甚麼與我爭鬧呢?你們為甚麼試探耶和華呢?」 3百姓在那裡口渴要喝水,就向摩西發怨言,說:「你為甚麼把我們從埃及領出來,使我們和我們的兒女,以及牲畜都渴死呢?」 4摩西就呼求耶和華說:「我要怎樣對待這百姓呢?他們差一點就要拿石頭打死我了.」 5耶和華對摩西說:「你帶著以色列的幾個長老,走在百姓前面,手裡拿著你先前擊打尼羅河的杖,去吧! 6看哪,我要在何烈的磐石那裡,站在你面前.你要擊打磐石,水就會從磐石流出來,給百姓喝.」摩西就在以色列的長老眼前這樣做了. 7他給那地方起名叫瑪撒(「瑪撒」意思是「試探」.),又叫米利巴(「米利巴」意思是「爭鬧」.),因為以色列人在那裡爭鬧,並且試探耶和華,說:「耶和華是否在我們中間呢?」

民數記20:1-13

1正月間,以色列全會眾到了尋的曠野；百姓住在加低斯.米利暗死在那裡,也葬在那裡. 2會眾沒有水,就聚集反對摩西和亞倫. 3百姓與摩西爭鬧,說:「我們恨不得與我們的弟兄一同死在耶和華面前. 4你們為甚麼領耶和華的會眾到這曠野,使我們和我們的牲畜都死在這裡呢? 5你們為甚麼領我們從埃及上來,把我們帶到這壞的地方呢?這地方不能撒種,沒有無花果樹,葡萄樹,石榴樹,也沒有水喝.」6 摩西,亞倫離開會眾面前,到會幕的門口,臉伏於地；耶和華的榮光向他們顯現. 7耶和華吩咐摩西說: 8「你拿著杖去,和你的哥哥亞倫召集會眾,在他們眼前吩咐磐石湧出水來,水就會從磐石流出,給會眾和他們的牲畜喝.」9於是摩西遵照耶和華所吩咐他的,從耶和華面前拿了杖去. 10摩西和亞倫召集會眾到磐石前.摩西對他們說:「聽著,你們這些悖逆的人!我們要叫這磐石流出水來給你們嗎?」 11摩西舉起手來,用杖擊打磐石兩下,就有許多水流出來,會眾和他們的牲畜都喝了. 12但是耶和華對摩西,亞倫說:「因為你們不信我,沒有在以色列人眼前尊我為聖,所以你們必不能領這會眾進入我所要賜給他們的地去.」 13這就是米利巴(「米利巴」意思是「爭鬧」.)水,因以色列人與耶和華爭鬧,耶和華在他們面前顯為聖.

無論任何因由,也不可以試探上帝

結語:

耶穌受試探的價值及意義

耶穌被聖靈充滿,聖靈引到曠野 > 魔鬼出現試探耶穌

3:14-15耶穌帶著聖靈的能力回到加利利,他的名聲傳遍了四方.他在各會堂裡教導人,眾人都稱讚他.

魔鬼不會罷休,不會停手,直到最後.

不要怕

約翰福音16:33

在世上你們有苦難,但你們要有勇氣,我已經勝過世界.

勝過試探的關鍵:上帝的話語,就是只有上帝的話語!!!



\newpage

\section{雞舍狐狸}
\label{sec:590}
\textbf{李均熊博士 - 雞舍狐狸}
\newline
\newline
\hyperref[sec:589]{< < < PREV SERMON < < <}
~
\hyperref[sec:toc]{[返目錄]}
~
\hyperref[sec:591]{> > > NEXT SERMON > > >}


路加福音 13:31-35
\newline
\begin{longtable}{cl}
\hline
\hline
章節 & 經文 (和合本修訂版)\\
\hline
13:31 & \begin{tabularx}{0.7\textwidth}{X} 就在那時，有幾個法利賽人來對耶穌說：「離開這裡到別處去吧，因為希律想要殺你。」 \end{tabularx} \\ \\ \relax
13:32 & \begin{tabularx}{0.7\textwidth}{X} 耶穌對他們說：「你們去告訴那個狐狸：『你看吧，今天明天我趕鬼治病，第三天我的事就成了。』 \end{tabularx} \\ \\ \relax
13:33 & \begin{tabularx}{0.7\textwidth}{X} 雖然這樣，今天明天後天我必須向前走，因為先知是不可能在耶路撒冷之外被害的。 \end{tabularx} \\ \\ \relax
13:34 & \begin{tabularx}{0.7\textwidth}{X} 耶路撒冷啊，耶路撒冷啊，你常殺害先知，又用石頭打死那奉差遣到你這裡來的人。我多少次想聚集你的兒女，好像母雞把小雞聚集在翅膀底下，可是你們不願意。 \end{tabularx} \\ \\ \relax
13:35 & \begin{tabularx}{0.7\textwidth}{X} 看吧，你們的家要被廢棄。我告訴你們，你們絕不會再見到我，直到你們說：『奉主名來的是應當稱頌的！』」 \end{tabularx} \\ \\
[1ex]
\hline
\hline
\end{longtable}
講題:雞舍狐狸

經文:路加福音13章31-35節

講員:李均熊博士

日期:2022年3月13日

基督面對祂的仇敵的威嚇,但祂仍一心要完成上主的使命...

路加福音13章31-35節 31就在那時,有幾個法利賽人來對耶穌說:「離開這裡到別處去吧,因為希律想要殺你.」32耶穌對他們說:「你們去告訴那個狐狸:『你看吧,今天明天我趕鬼治病,第三天我的事就成了.』33雖然這樣,今天明天後天我必須向前走,因為先知是不可能在耶路撒冷之外被害的.34耶路撒冷啊,耶路撒冷啊,你常殺害先知,又用石頭打死那奉差遣到你這裡來的人.我多少次想聚集你的兒女,好像母雞把小雞聚集在翅膀底下,可是你們不願意.35看吧,你們的家要被廢棄.我告訴你們,你們絕不會再見到我,直到你們說:『奉主名來的是應當稱頌的!』」

當耶穌直斥要殺他的希律是狐狸,他轉頭卻形容自己好像母雞!

狐狸與母雞,那一個著數,那一個蝕底?大家有沒有聽過一句英語的說法: Fox in the henhouse；姑且譯它做雞舍狐狸(亦即今天的講題).這句說話可以解作「不安好心」:喜獵食雞肉的狐狸拜訪雞舍,母雞自然凶多吉少.中國其實亦有這類的說法,不知大家聽過黃鼠狼向雞拜年,歇後語就是「不安好心」.Fox in the henhouse這句說話原初亦可能是寫成Fox guarding the hen house.讓狐狸守雞舍,自是所托非人也.總而言之,將雞放在狐狸面前,大家可以想像總不會有甚麼好結果的了.

但耶穌一方面稱呼希律是狐狸,但隨即說自己好像母雞,用意又可在?

經文記載法利賽人去向耶穌「報信」,表面上好像好心通知,但我看他們也只是黃鼠狼拜年——不安好心.前兩章(第11 章),耶穌才斥責法利賽人有禍了,說他們是不顯露的墳墓,洗淨杯盤的外面,裡面卻滿了貪婪和邪惡；自此,法利賽人伺機要抓耶穌的把柄.現在他們收到消息希律想要殺害耶穌,自然用這個藉口趕走這個惱人的耶穌:你好快D 過主啦,唔係希律就要殺你!所以耶穌對他們的回答一點不感激,更一點也不客氣:你話希律要殺我,你就同我去講比呢隻狐狸聽……也不知法利賽人是否真有膽去告訴希律耶穌的說話.

大家要留意,千萬不要把希律說成「老狐狸」.老狐狸代表老奸巨滑,狡計多端.但耶穌這裡稱希律做狐狸,看來並沒有表示希律是甚麼詭計多端的惡人.在路加福音中,希律最多也只是如狐狸般四處破壞,是農戶最討厭的動物.狐狸在當時文化中亦可代表表面好像擁有威權,其實沒有實質權力,正如中文「狐假虎威」一樣,只是作威作福.耶穌稱呼希律作狐狸,自是代表他自然不會懼怕這個作威作福,並無實權的討厭傢伙.只於耶穌要法利賽人告訴「那個狐狸」的說話,大意其實是耶穌自己正忙著治病趕鬼的事工,原全不會理會這個所謂希律想要殺他的威脅.但是,耶穌不日就會繼續向前行,真如法利賽人所言離開加利利地區,前往耶路撒冷.耶穌要表明並不是希律又或法利賽人的脅迫和加害令祂要去逃避,而是因為要完成上主的召命,就是先知的工作,故此祂仍要前行.上主的心意並非要祂在加利利北方外邦人之境遇害,而是要在耶路撒冷完成基督的工作,在那裡受苦受死,完成救恩.

耶穌基督我們的主並不因迫害和威脅而退縮,但卻為著完成上主的使命,而選擇前行,並不示弱,也並不「送頭」.作為祂的們徒,我們是否知道表面上我們受逼迫,但所有逼迫我們的勢力其實都是狐狸,都是沒有實質權力,因為真正的權柄其實由始至終都掌管在我們的主手上.正如今日預苦期第二主日經課的詩篇27篇啟應一樣:

1耶和華是我的亮光,是我的拯救, 我還怕誰呢? 耶和華是我生命的保障, 我還懼誰呢?

2那作惡的就是我的仇敵, 前來吃我肉的時候就絆跌仆倒.

3雖有軍隊安營攻擊我,我的心也不害怕； 雖然興起戰爭攻擊我,我仍舊安穩.

不錯,預苦期是要我們為著耶穌的受苦受死作預備,但讓我們緊記,耶穌受死並不是因為世上的政權加害於祂,而是祂為著成就上主的救恩甘心擺上自己,行走這一條住耶路撒冷的十架路.故此,作為基督的們徒,我們可以不卑不亢,為著完成上主的召命能不懼困阻.問題是,甚麼是我們要完成的上主的使命?在這個預苦期,就讓我們思考我們作為們徒的召命,尋著上主給我們的目標,我們就能面對強權而不怕.

但今天經課讓我們看見,耶穌不怕希律,稱呼他作狐狸,並不是因為祂能對抗希律或他背後所代表的羅馬強國的極權統治.剛好相反,耶穌比擬自己是母雞,是狐狸的獵物,而不是會抗衡狐狸力量的獅子或巨鷹(羅馬帝國的旗幟就是一隻巨鷹!).我們不要以為耶穌不懼怕希律,就代表祂會對抗他的勢力.不對,基督的的確抗衡希律的殘忍或帝國的壓迫,但祂用的不是赤裸的權力,而是溫柔的招聚,將祂所愛的擁在懷中,就如母雞把小雞聚集在翅膀底下.基督展示的,是無能者的大能,並非武力的對抗當時羅馬的政權.耶穌不懼怕這個政權,但祂也不是要推翻這個政權.政權的力量是屬於上主的,政權的興衰都全然掌管在我們的主手中.但基督所做的,就是以溫柔的力量,去保護擁抱受驚受害的人.

可惜的是,人們通常追求的不是這樣一個如母雞的基督；人們祈盼的是那個猶大的獅子,駕著雲彩降臨的人子.你可明白為甚麼耶路撒冷的兒女,總是不願意聚在母雞的翅膀底下.耶穌甚至不是一隻公雞；公雞面對狐狸巨鷹至少都可以反擊自保(大家有看過殘忍的鬥雞嗎?).但耶穌比喻自己為母雞,把小雞罩在牠的翅膀下.你有沒有玩過一個遊戲叫「麻鷹捉雞仔」.做母雞的其實是以自己的身軀擋著麻鷹的攻擊!母雞是做好犧牲自己的準備.甚至有人想像母雞可能是想用自己去滿足狐狸的食慾,以求狐狸放過自己的小雞……我們喜見耶穌稱呼希律做狐狸,我們又體會耶穌為了保護耶路撒冷的子民願意犧牲自己,卻見耶路撒冷中人拒絕耶穌張開的翅膀,甚至真正殺害自己的不是那隻狐狸,而是耶路撒冷的子民?

耶穌不怕狐狸,卻甘願為母雞招聚小雞.教會可以學效基督不懼強權,視權勢為狐狸嗎?教會又真的可以如基督像母雞保護小雞般,保護在制度下受壓迫的人嗎?3月11日《時代論壇》正好有一篇報導(參下)展示教會保護受苦群眾的可能性.在預苦期的第二主日,讓我們學習耶穌拒絕強權的操弄,專心完成上主的使命；但同時願意學習基督溫柔的樣式,如母雞照顧小雞般守護受苦受害的人.

教會如何可以像基督母雞形像保護在制度下受壓受傷的人?

原報導網址: \url{https://christiantimes.org.hk/Common/Reader/News/ShowNews.jsp?Nid=168195 and Pid=102 and Version=0 and Cid=2143 and Charset=big5_hkscs}

美國教會籌款為窮人註銷醫療債務

時代論壇 / 2022年3月11日

美國聯合基督教會(United Church of Christ)在二月十四日情人節宣佈,於每年一度的「慈善星期二」(Giving Tuesday)中籌得二十萬美元,為該教會總部所在俄亥俄州的居民註銷了三千三百萬美元的醫療債務,超過一萬個家庭受惠.自二零一九年起,該教會已替俄亥俄州的低下階層註銷超過一億美元的醫療債務.

高昂的醫療費用及龐大的醫療債務是困擾美國家庭的其中一個問題.據史丹福大學經濟學家估計,全美目前的醫療債務達一千四百億美元.而美國的人均醫療開支亦遠高於其他發達國家,比第二位德國的六千美元高近八成,約一萬零五百美元.

聯合基督教會兩年前起,透過非牟利慈善組織「永別醫療債務」(RIP Medical Debt,下稱「永別」)替俄亥俄州最貧窮的人──包括收入在美國聯邦政府貧窮線以下兩倍,或是面臨破產等──「償還」債務.「永別」由兩位曾任職醫療債務收數公司管理層的人於二零一四年成立,該組織從債主收購債務,然後將它們註銷.由於很多醫療債務的債務人最終未能還款,債主便將這些債務捆綁出售給追數公司.「永別」表示,每一美元可購買約一百美元的債務.至去年六月,該組織透過從不同組織和人士獲得的捐款,成功在全美註銷了四十五億美元的醫療債務. 聯合基督教會牧師翰林(Sekinah Hamlin)表示,「醫療債務仍然是美國民眾破產和被追數的主要原因,這亦導致有些人寧願選擇放棄治療.醫療保健是一項人權,沒有人應該因為貧窮而無法得到醫治.」

「永別」在註銷債務時不會附加任何條件,教會亦不知道受惠者的名字.但組織會向因教會捐款而受惠的人寄出信件,告訴他們「註銷債務所需的資金,是由聯合基督教會慷慨提供.」

其中一位受惠者表示,「我一開始以為這是一場騙局……我從未聽說過這種事情……我很高興,也很震驚,有人能為我做這麼好的事情.」

聯合基督教會公義及本地教會事工副總主任布萊克蒙(Traci Blackmon)牧師表示,希望教會的努力能推動醫療系統改革,「這個行動的公開性質能讓更多人看見我們醫療系統的內在性不公；它針對最沒有能力反抗的人,以及這個國家從病患身上賺取巨額盈利是一件何等荒謬的事.這些都是大多數媒體鮮有主動報道的事情.」

布萊克蒙又表示,「在這個思索愛的日子,我想起了保羅對腓立比教會說的話:『所以,在基督裡若有任何勸勉,若有任何愛心的安慰,若有任何聖靈的團契,若有任何慈悲憐憫,你們就要意志相同,愛心相同,有一致的心思,一致的想法,使我的喜樂得以滿足.凡事不可自私自利,不可貪圖虛榮；只要心存謙卑,各人看別人比自己強.各人不要單顧自己的事,也要顧別人的事』.謝謝所有伸出援手的人.」



\newpage

\section{我聖爾亦當聖}
\label{sec:591}
\textbf{鄧瑞強博士 - 我聖.爾亦當聖}
\newline
\newline
\hyperref[sec:590]{< < < PREV SERMON < < <}
~
\hyperref[sec:toc]{[返目錄]}
~
\hyperref[sec:592]{> > > NEXT SERMON > > >}


利未記 19:1-18
\newline
\begin{longtable}{cl}
\hline
\hline
章節 & 經文 (和合本修訂版)\\
\hline
19:1 & \begin{tabularx}{0.7\textwidth}{X} 耶和華吩咐摩西說： \end{tabularx} \\ \\ \relax
19:2 & \begin{tabularx}{0.7\textwidth}{X} 「你要吩咐以色列全會眾，對他們說：你們要成為聖，因為我耶和華－你們的神是神聖的。 \end{tabularx} \\ \\ \relax
19:3 & \begin{tabularx}{0.7\textwidth}{X} 你們各人都當孝敬父母，也要守我的安息日。我是耶和華－你們的神。 \end{tabularx} \\ \\ \relax
19:4 & \begin{tabularx}{0.7\textwidth}{X} 你們不可轉向虛無的神明，也不可為自己鑄造神像。我是耶和華－你們的神。 \end{tabularx} \\ \\ \relax
19:5 & \begin{tabularx}{0.7\textwidth}{X} 「你們宰殺祭牲獻平安祭給耶和華的時候，要獻得使你們可蒙悅納。 \end{tabularx} \\ \\ \relax
19:6 & \begin{tabularx}{0.7\textwidth}{X} 這祭物要在獻的當天或第二天吃；若有剩到第三天的，就要用火焚燒。 \end{tabularx} \\ \\ \relax
19:7 & \begin{tabularx}{0.7\textwidth}{X} 第三天若再吃，這祭物是不潔淨的，必不蒙悅納。 \end{tabularx} \\ \\ \relax
19:8 & \begin{tabularx}{0.7\textwidth}{X} 吃的人必擔當自己的罪孽，因為他褻瀆了耶和華的聖物，這人必從百姓中剪除。 \end{tabularx} \\ \\ \relax
19:9 & \begin{tabularx}{0.7\textwidth}{X} 「你們在自己的地收割莊稼時，不可割盡田的角落，也不可拾取莊稼所掉落的。 \end{tabularx} \\ \\ \relax
19:10 & \begin{tabularx}{0.7\textwidth}{X} 不可摘盡葡萄園的葡萄，也不可拾取葡萄園中掉落的葡萄，要把它們留給窮人和寄居的。我是耶和華－你們的神。 \end{tabularx} \\ \\ \relax
19:11 & \begin{tabularx}{0.7\textwidth}{X} 「你們不可偷盜，不可欺騙，也不可彼此說謊。 \end{tabularx} \\ \\ \relax
19:12 & \begin{tabularx}{0.7\textwidth}{X} 不可指著我的名起假誓，褻瀆你神的名。我是耶和華。 \end{tabularx} \\ \\ \relax
19:13 & \begin{tabularx}{0.7\textwidth}{X} 「不可欺壓你的鄰舍，也不可偷盜。雇工的工錢不可在你那裡過夜，留到早晨。 \end{tabularx} \\ \\ \relax
19:14 & \begin{tabularx}{0.7\textwidth}{X} 不可咒罵聾子，也不可將絆腳石放在盲人面前。你要敬畏你的神。我是耶和華。 \end{tabularx} \\ \\ \relax
19:15 & \begin{tabularx}{0.7\textwidth}{X} 「你們審判的時候，不可不公正；不可偏護貧窮人，也不可看重有權勢人的臉，總要公平審判你的鄰舍。 \end{tabularx} \\ \\ \relax
19:16 & \begin{tabularx}{0.7\textwidth}{X} 不可在百姓中到處搬弄是非，不可陷害鄰舍的性命。我是耶和華。 \end{tabularx} \\ \\ \relax
19:17 & \begin{tabularx}{0.7\textwidth}{X} 「不可心裡恨你的弟兄；要指摘你的鄰舍，免得因他承擔罪過。 \end{tabularx} \\ \\ \relax
19:18 & \begin{tabularx}{0.7\textwidth}{X} 不可報仇，也不可埋怨你本國的子民。你要愛鄰如己。我是耶和華。 \end{tabularx} \\ \\
[1ex]
\hline
\hline
\end{longtable}
講題:我聖,爾亦當聖

經文:利未記19:1-18

講員:鄧瑞強博士

日期:2022年3月20日

創世記講到,在黑暗的長夜,神揀選了亞伯拉罕作點燈的人.要燈火長明,需要一個有延續性的群體,將燈火一代又一代的傳下去.亞伯拉罕的後裔所形成的民族,扮演著這個角色.要成為背負天命的民族,以色列人一定要學曉,在黑暗中堅持光明.

人類歷史的巨大黑暗,其中之一,是人性的奴役.以色列民族落入這巨大的黑暗中,他們在埃及地,成為了法老的奴隸.他們手持的火把,看來要熄滅了.神卻在不滅的火中,呼喚摩西.神要摩西去解放這群奴隸.摩西面對法老,法老有的是兵馬,摩西有什麼,他什麼也沒有,有的,只是神在他內心燃點的,對自由的渴望.法老有的是權勢,摩西無權無勢,有的只是上帝.摩西有的,是上帝的承諾,是上帝承諾下對歷史前景的寄望,是對自由的想像.於是,摩西帶領以色列人離開埃及,追求自由,書寫歷史的新的一頁,也塑造了一個有深刻歷史經驗的民族.要塑造一個群體,要塑造一個民族,要有典範性的經驗.一群奴隸,在神的救贖下,成為自由群體,這歷史經驗深深地刻在以色列人的靈魂裡,幾千年來,他們守逾越節,重述這歷史故事,回憶這動人的經驗.無數次,在歷史長河裡,他們成了奴隸,在巴比倫手下,在波斯手下,在希臘手下,在羅馬手下,在希特拉手下,他們仍堅持守逾越節,重述出埃及的故事,一次又一次,這賦予了他們不死的信念,在奴隸的人生中,想像自由,渴望自由,相信自由.這是一種信念,這信念塑造了一個信仰的民族.他們相信,人生的無奈,不是生命的全部,生命的底蘊,是上帝的救贖和同在；歷史的前景,在上帝的應許之中.就這樣,一個信仰的民族,立足於大地之上.

救贖的經驗,塑造了神的子民,但,這經驗本身,不足以成為生活的南針.這救贖經驗如何變為生活的內容呢?生活的內容如何反映上帝的救贖呢?上帝會親自教導他們.

出了埃及,踏過紅海,以色列人進入了西奈半島的曠野,他們告別了奴役,獲得了自由,發現自己卻是活在荒蕪的土地上.他們沒有什麼東西可吃,沒有水可飲,只仰賴上天的供應.天父憐憫他們,有充足的供應.但,他們竟眷戀舊時的生活,常常想回埃及.人,真是活在根深柢固的奴役力量裡,想做奴隸的心,強得不得了.要讓人做一個真正自由的人,真的很難.

在西奈半島的曠野,這群人真的有很多東西要學.曠野,是嚴苛的人性學校,神的子民要在這裡努力學習,學習在一無所有的環境中信任神,相信神的看顧；學習拒絕奴役的誘惑,背負著自由的十架；學習從神而來的全新事物,做人類中完全不一樣的民族,一群屬於神的子民.利未記,就是神的子民在曠野這所學校中,要學習的主要內容.

利未記這卷書的位置,位於摩西五經的中間,其信息也是五經的中心.但很可惜,很多信徒對利未記一點興趣也沒有,便白白錯失了五經的中心思想.在利未記裡,神教導祂的子民,如何成為獨一無二的群體.利未記有一相當簡單而平行的結構,講四個主題.不計書卷的結語,利未記的頭和尾,講祭禮與節期,其用意是教導我們,透過一些神聖禮儀,將生活與日子聖化,好讓我們體會神在人間.接著,開頭的第二大段和尾二段落,講祭司的設位和職責,這些人物是守護人,守護著神的子民的聖潔.接著,前面的第三大段落及尾三段落,講神的子民的聖潔性,這是每一位神的子民的生活操守.然後,整卷書的中心,是第16章,講「贖罪日」,在神面前,沒有人在道德及聖潔上站立得住.人若真的學曉在神面前虛心反省,乞求寬恕,則人間的仇恨便可消除了.「贖罪日」,凸顯神和人之間的無限差異.神是創造主,人是被造物,這是最重要的存在秩序,這秩序若被破壞,則有災難的後果.整卷利未記的教育規劃,是以「贖罪日」為中心,教人尊敬神聖；最外圍是祭禮與節期,讓人活在神聖的氛圍中；輔以祭司,作神聖的守護人；最後,是教育的成果,就是神子民的聖潔性,這聖潔性,表現在崇高的操守上.

今日的講道經文,是被稱為「聖潔法典」的經文,這經文教導神的子民活出聖潔:利未記19: 1-18.

為方便大家理解這段經文的結構,我先作少許說明.這段經文分成上下兩大段落.上段以「我是耶和華－你們的神」作為小分段標記,有4個這種標記.下段以「我是耶和華」作小分段標記,同樣有4個這種標記.我一會兒讀的時候,會叫大家留意一下,這樣,這段經文的結構便分明了.

19:1　耶和華對摩西說:

19:2　「你曉諭以色列全會眾說:你們要聖潔,因為我耶和華－你們的神是聖潔的.

19:3　你們各人都當孝敬父母,也要守我的安息日.我是耶和華－你們的神.

19:4　你們不可偏向虛無的神,也不可為自己鑄造神像.我是耶和華－你們的神.

19:5　你們獻平安祭給耶和華的時候,要獻得可蒙悅納.

19:6　這祭物要在獻的那一天和第二天吃,若有剩到第三天的,就必用火焚燒.

19:7　第三天若再吃,這就為可憎惡的,必不蒙悅納.

19:8　凡吃的人必擔當他的罪孽；因為他褻瀆了耶和華的聖物,那人必從民中剪除.

19:9　「在你們的地收割莊稼,不可割盡田角,也不可拾取所遺落的.

19:10　不可摘盡葡萄園的果子,也不可拾取葡萄園所掉的果子；要留給窮人和寄居的.我是耶和華－你們的神.

19:11　「你們不可偷盜,不可欺騙,也不可彼此說謊. 19:12　不可指著我的名起假誓,褻瀆你神的名.我是耶和華.

19:13　「不可欺壓你的鄰舍,也不可搶奪他的物.雇工人的工價,不可在你那裡過夜,留到早晨.

19:14　不可咒罵聾子,也不可將絆腳石放在瞎子面前,只要敬畏你的神.我是耶和華.

19:15　「你們施行審判,不可行不義；不可偏護窮人,也不可重看有勢力的人,只要按著公義審判你的鄰舍.

19:16　不可在民中往來搬弄是非,也不可與鄰舍為敵,置之於死.我是耶和華.

19:17　「不可心裡恨你的弟兄；總要指摘你的鄰舍,免得因他擔罪.

19:18　不可報仇,也不可埋怨你本國的子民,卻要愛人如己.我是耶和華.

聖潔法典的第一天條,是:

19:2　「你們要聖潔,因為我耶和華－你們的神是聖潔的.」(文理譯本:「我乃耶和華,爾之上帝,我聖,爾亦當聖.」這就成了今日的講道題目.)

我們要聖潔,因為我們的神是聖潔的,理由簡單明白,我們屬於誰,就像誰.「聖潔」,重點在於「有分別」,所以舊約常常講「分別為聖」.神與萬物有別,故我們不能敬拜萬物.我們敬拜神,也要與世人有別,不能按世人的標準來生活.以下講的,就是神的子民與別不同的生活方式.

19:3　「你們各人都當孝敬父母,也要守我的安息日.我是耶和華－你們的神.」

「孝敬」,原文即「敬畏」.在原文,先講「母」,後講「父」.人一出生,被母親乳養,這是一個人的生命所依.對生命的根源,要學懂「敬畏」.連這敬畏的心都沒有,就目中無人了.目中無人的人,不可能敬拜神.

敬畏「母和父」,連於遵守神的安息日.安息日,重點不在休息,而在:放下我們塵世的擔心,好好仰望神,信靠神.我們很易將安息日看成是「休息」的日子,然後說,休息,是為了走更遠的路.這是思念成功,多過思念上帝.這已不能從世人中分別出來了,這已違反了聖潔.

19:4　你們不可偏向虛無的神,也不可為自己鑄造神像.我是耶和華－你們的神.

這是重申「聖潔」的重要性.將生命寄托於被造之物,已不聖潔了. 19:5　你們獻平安祭給耶和華的時候,要獻得可蒙悅納.

19:6　這祭物要在獻的那一天和第二天吃,若有剩到第三天的,就必用火焚燒.

19:7　第三天若再吃,這就為可憎惡的,必不蒙悅納.

19:8　凡吃的人必擔當他的罪孽；因為他褻瀆了耶和華的聖物,那人必從民中剪除.

19:9　「在你們的地收割莊稼,不可割盡田角,也不可拾取所遺落的.

19:10　不可摘盡葡萄園的果子,也不可拾取葡萄園所掉的果子；要留給窮人和寄居的.我是耶和華－你們的神.

這裡提到「平安祭」,這種祭,不是規定要獻的,是自願獻的.若你犯罪,在舊約,要獻贖罪祭,這是規定的,不得不獻的.但平安祭,並不是規定的.若你一生都不獻平安祭,是沒有問題的.平安祭,是感恩的祭,這種祭的特點,是人可吃用獻上的祭物.

這裡提到:這些祭物只能吃兩天,不能留到第三天吃.這是防止人獻了祭,卻又將祭物私自留下,讓自己人慢慢享用.世人會如此做,但聖潔的人不能如此.在這裡,獻祭的人要學的功課,就是:既然是自願獻的祭,就當甘心與外人分享.

下面的話,也強調分享.「在你們的地收割莊稼,不可割盡田角,也不可拾取所遺落的.不可摘盡葡萄園的果子,也不可拾取葡萄園所掉的果子；要留給窮人和寄居的.」在生活裡,自己夠了,就當滿足了,要讓不夠的人有機會執拾零碎,不要賺盡,不要取盡,讓他人有生機.在新聞中聽到,有些茶餐廳沒得賺,也繼續營業,為的是讓手停口停的員工有工可返,有啖飯可吃.願神祝福這些老闆.

講到這裡,講完「聖潔法典」的上半段,重點是活出生命的聖潔,因為救贖我們的神是聖潔的.我們要與別不同,因為我們的神是與別不同的.

中國人說,「內聖」「外王」,有內在的聖潔,就會外顯為王道.有內心的敬虔,就會有外在的公義,仁愛.若果上一段講的是「敬畏神」,如今這一段便是講愛人如己.這大段也是分為4小分段.

19:11　你們不可偷盜,不可欺騙,也不可彼此說謊.

19:12　不可指著我的名起假誓,褻瀆你神的名.我是耶和華.

人與人相處,要誠實,不要貪圖別人的東西.一貪,你便會害人的了.大衛王一貪戀別人的妻子,便設計殺了別人的丈夫.貪,是萬惡之根.

在古代,很多時,東西是不是你的,最後要以起誓作準.故此,這裡說到,不要為了貪圖某些東西,而在起誓時發假誓.貪念,最後演變為冒犯神.

19:13　不可欺壓你的鄰舍,也不可搶奪他的物.雇工人的工價,不可在你那裡過夜,留到早晨.

19:14　不可咒罵聾子,也不可將絆腳石放在瞎子面前,只要敬畏你的神.我是耶和華.

以前,是日薪的,一日出一日的糧,工人靠這糧,讓家人一日一日的活.你扣壓工資,會害到人家當日無飯食.不要欺壓弱者,窮人,不要欺負聾人,盲人.敬畏神的人,當對人心存尊敬.不要單單尊敬大人物,也要尊敬小人物,因為:神看重他們.想想:就是因為神看重以色列這班卑微的奴隸,神才派摩西來解救他們.我們的神看重弱者,所以我們也要看重他們.

19:15　「你們施行審判,不可行不義；不可偏護窮人,也不可重看有勢力的人,只要按著公義審判你的鄰舍.

19:16　不可在民中往來搬弄是非,也不可與鄰舍為敵,置之於死.我是耶和華.

判斷對錯,不能偏私,不能親疏有別,不能為了自己的潛在利益而作虛假判斷,不能因為你不同意某些人的意見而判定他們一定錯誤.判斷,一定要保持正義.這是很難的,因為:我們有很多偏見,也有很多利益考慮.經文說到,你錯誤地判斷,會害死人的.故此,要力求正義.

19:17　「不可心裡恨你的弟兄；總要指摘你的鄰舍,免得因他擔罪.

19:18　不可報仇,也不可埋怨你本國的子民,卻要愛人如己.我是耶和華.

這裡提到:一,不要心中有恨.有恨,就有殺機.二,盡力指正身邊的人的錯誤,這是救他免除大惡.三,不要埋怨,因為:怨生恨,恨生殺.四,重中之重:愛人如己,或更正確的翻譯:愛鄰舍如同自己.

之前講過的,不貪婪,不欺壓,不偏私,不懷恨,仍不夠,敬畏神的人,會表現出對人的最大尊敬,就是:愛鄰舍如同自己.如此,才活出神的子民完全獨一無二的聖潔.

結語:

這「聖潔法典」,始於敬畏神,終於愛人如己.能敬畏神,則承認自身的卑下與有限；能愛人如己,則與世人有別,活出聖潔.我們要聖潔,因為耶和華－我們的神是聖潔的. 　　願榮耀歸上主,願人間享和平.

...Copyright (C) 本網站上所有視頻／音頻／文字版權歸【葡萄樹傳媒】所有,轉載請註明出處! \url{https://www.vinemedia.org/course/bible-teaching/sunday-service-1/ye-shall-be-holy-for-i-the-lord-your-god-am-holy/} .



\newpage

\section{吉甲的覺醒進入應許之地}
\label{sec:592}
\textbf{伍渭文牧師 - 吉甲的覺醒—進入應許之地}
\newline
\newline
\hyperref[sec:591]{< < < PREV SERMON < < <}
~
\hyperref[sec:toc]{[返目錄]}
~
\hyperref[sec:593]{> > > NEXT SERMON > > >}


講題:吉甲的覺醒—進入應許之地

講員:伍渭文牧師

經課:路加福音15

日期:2022年3月27日

引言:每一架飛機在燃油記錄儀中間有一標誌—Point of No Return 折返點,顯示燃油用了過半,若要回航就在此刻,過了這點就不夠燃油返回起航的機場了.

出埃及是一個屬靈生命的航程,像肉身生命的航程,沒有折返點.人可以回到過去的地點,但不能回到過去的時間.出埃及另一端的機場是迦南流奶與蜜之地,這是上主應許給亞伯拉罕的.創十二章上主呼召七十五歲亞伯蘭離開吾珥,帶同妻子,羅得和財物去到迦南地.

創十五章到了迦南地,上主再次向亞伯蘭顯現,進一步告訴他:「你的後裔必寄居別人的地,又服事那地的人,那地的人要苦待他們四百年……到了第四代,他們必回到這地,因為亞摩利人(西方人,迦南人統稱)的罪孽,還沒有滿盈.」(創十五13-16)跟著上主與亞伯蘭立約:「日落天黑,不料有冒煙的爐並燒著的火把,從那些肉塊經過.當日耶和華與亞伯蘭立約.」(創十五17, 18)這是立約儀式:行經切開的祭牲,以生命保證不會爽約,否則像祭牲剖開.到了創十七,亞伯蘭改名並行割禮.「你們所有的男子,都要受割禮,這是我與你並你的後裔所立的約,是你們所當遵守的.」(創十七10)

先知何西亞用婚姻比喻這約:「我要永遠聘娶你,以正義,公平,慈愛,憐憫聘娶你；以忠實聘娶你,使你認識我是上主.」(何2:21-22).但以色列人對這婚姻不忠誠,忘記自己是上主的約民,多次離開上主,拜偶像.因而北國以色列被亞述國所滅,南國猶大被巴比倫所滅,備受亡國之痛.

放下行李.以色列人的失敗,就如不少婚姻的失敗,其中一個原因,是忘記起航後是沒有折返點.有些人結婚時,行李仍沒有放下,五年,十年仍未放下行李,隨時可以離開,結果實現了「自我實現的預言」(Self-fulfilled prophesy).

在吉甲行割禮就是約民身份的覺醒.上主與先祖亞伯拉罕立約,賜他後裔應許之地,此刻以色列人已渡過約但河,立足應許之地,他們要決志從此結束飄流的生活,做上主的子民.

今天福音經課路十五章的浪子,也醒悟自己的身份,回到父親家中.「我要起來,到我父親那裡去,向他說:父親,我得罪了天,又得罪了你,從今以後,我不配稱為你的兒子.把我當作一個雇工吧!」(路十五18, 19)今日的題目「吉甲的覺醒」,是約民身份的覺醒,屬靈生命的覺醒.

1 吉甲的覺醒—認清身份,不再飄流

「你製造火石刀,第二次給以色列人行割禮.約書亞就製造火石刀,在除皮山那裡給以色列人行割禮.」(書五2,3)「耶和華對約書亞說,我今日將埃及的羞辱從你們身上輥去了.因此那地方名叫吉甲,直到今日.」(書五9)

約書亞叫所有在曠野出生的以色列人,在吉甲進行割禮並遵守逾越節,再次認定他們的身份,結束過去飄流的生活,回歸祭司國度,屬神的子民,征戰就在前面,要為真道打美好的仗.

割禮是聖洗

聖經稱割禮就是聖洗. 「你們在他裡面,也受了不是人手所行的割禮.你們既受洗與他一同埋葬,也就在此與他一同復活.你們從前在過犯,和未受割禮的肉體中死了.神赦免了你們一切過犯,便叫你們與基督一同活過來.」(西二11, 12)」

a聖洗是開竅,教會稱領了洗的人蒙光照(Enlightened),得著靈性的覺醒；領洗時燃點一枝蠟燭,象徵蒙了光照,也願意把真光傳遍天下.吉甲原文是輥開,割禮就是把陽皮輥開,所以吉甲另一名字叫除皮山.輥開也是開竅,蒙光照.

b蒙光照是明白真理,靈性復活.未蒙光照前,我不曉得為何來到世上,我信有造物主,但不知他是誰,如何找到他；人死後甚麼都沒有,人生有何意義呢?人生沒有上主,在世確實沒有指望.年紀愈大,更體會「沒有指望」這話的重量.但現在我知離世是進入更美的家鄉,如葛培理所說:我沒有死,祇是換了地址.蒙光照是靈性復活.信主後我真的經歷了與基督一同活過來.中學二年級信耶穌後,經驗生命的改變.一般中四(Form 4)才成為領袖生,但中三已被選為領袖生,全校祇有兩位,一男一女.班主任問我:「你父母給你的家庭教育很成功很難得.」我說:「是耶穌改變了我.」

聖洗是誓師

a回絕黑暗.教會稱聖洗是入門禮儀(Rite of initiation),在聖洗禮序中領洗者要立誓回絕魔鬼.主禮:「我問你是否決意應許拒絕魔鬼各樣行為,丟棄世上虛浮之榮華,除去一切貪慾,與肉體之邪情.」在古教會領洗者此時轉身向西(魔鬼墮落方向)；用肢體語言回絕魔鬼,然後轉過頭來面向聖壇–基督復活的東面.

b誓師入旅.吉甲在耶利哥東北一英哩.在此行割禮有開赴戰場前誓師的意味.事實上,聖洗就有歸入基督軍旅誓師的意味.每一個領洗者已被徵入伍,進行屬靈的爭戰.公禱書聖洗禮文:「我等接納此人進入基督之教會,將十字聖記號畫其額上為記,表明其自今以後,必真信奉釘十字架之基督,在人面前承認基督不以為恥,在基督旗下,向罪惡世俗魔鬼奮勇爭戰,一生盡忠,事奉基督.」

輥開埃及的羞辱

割禮或聖洗是「承認基督不以為恥」,書五9:「今日將埃及的羞辱從你們中間身上輥去了.」甚麼是埃及的羞辱呢?為奴是埃及的羞辱嗎?但此刻他們已離開為奴之地,來到吉甲.埃及的羞辱應指埃及人嘲笑上帝,無能力引領百姓進入迦南地.拜金牛犢帶來上主的震怒,摩西求上主不要消滅以色列人,寧可把自己名字從生命冊塗抹,其中一個理由是不想使埃及人議論上帝,離間上主與選民的關係.「為甚麼使埃及人議論說,他領他們出去,是要降禍與他們,把他們殺在山中.」(出三十二12)

撒旦要敵擋的是上主,人根本不是他的對手；項莊舞劍,志在沛公.「約伯敬畏神豈是無故呢?……他手所做的都蒙你賜福……你且伸手毀他一切所有的,他必當面棄掉你.」(伯一9-11)

魔鬼叫人埋怨上主,懷疑上主,離間上主與人的關係,目的是敵對上主,因為上主極愛世人,甚至犧牲獨生愛子拯救世人.

馬丁路德1537年一篇約翰福音的講章中,提及教父一則軼聞,說明道成為肉身是何等敬虔的奧秘.古教會崇拜認信,詠唱(chanting)尼西亞信經到「他道成肉身」(Et homo factus est)時,會眾都同時曲膝下跪,但有一個人卻若無其事站在那裡,也不脫帽,以示崇敬,木若柱子.魔鬼實在看不過眼,趨前重重擊打他一下,頓時金星直冒.魔鬼詛咒他說:「願地獄吞噬你,你這蠢驢.若上主成為我們一樣的天使(魔鬼乃墮落的天使),會眾歌頌『上主成為天使』,我不單下跪,且要五體投地!是的,我會匍匐爬行十肘(五米).你這可惡的人……你聽到上主沒有成為天使,卻成為人,跟你一樣,還像木杉和石柱呆呆站著.」

出了埃及後以色列人祇在曠野繞圈,繞了四十年,這羞辱來自以色列人不信和埋怨,但也算在上帝的帳上.信徒的行為,可以榮神益人,亦可羞辱上主,把基督重釘十字架.但當信徒認定是上主子民的身份,讓心靈不再隨波逐流,就把埃及的羞辱輥去了.

小結:人生旅程沒有折返點,朝聖旅程也沒有折返點.我們可以回到過去的地點,但不能回到過去的時間.吉甲的覺醒是放下行李,拿起屬靈的兵器,參與屬靈的征戰.進入應許之地是宣揚上主的律法,叫罪惡滿盈的亞摩利人,及時悔改,逃過審判.

( 1) 吉甲的覺醒—認清身份,不再飄流 )

2) 吉甲的覺醒: 宣講信息,使人得福

「約書亞靠近耶利哥的時候舉目觀看,不料有一個人手裡有拔出來的刀,對面站立.約書亞到他那裡,問他說,你是幫助我們呢,是幫助我們敵人呢?他回答說,不是的,我來是作耶和華軍隊的元帥.約書亞就俯伏在地下拜說.我主有甚麼話吩咐僕人.」(書五13, 14)

站在上主那一邊

在人看來,敵人當前行割禮是不可思議的,傷口需時間復元,萬一敵人趁機攻過來怎辦?耶利哥城離吉甲祇一英哩.行割禮後才開始戰爭說明這是上主的戰事,不是以色列人的戰爭,是以色列人為上主爭戰.第五章13-15節出現一位持刀的人.約書亞問:「你是幫助我們呢,是幫助我們敵人呢?」這人答:「我來是要作耶和華軍隊的元帥.」

這是上主對付罪惡的戰爭,不是國族與國族的戰爭.「耶和華你的神將這些國民從你面前攆出以後,你們心裡不可說,耶和華將我領進來得這地,是因我的義,其實耶和華將他們從你面前趕出去,是因為他們的惡.你進去得他們的地,並不是因你們的義,也不是因你們心裡正直,乃是因這些國民的惡,耶和華你的神將他們從你面前趕出去,又因耶和華要堅定他向你列祖亞伯罕,以撒,雅各起誓所應許的話.」(申九4, 5)

1861年美 國南北戰爭開打初期北軍節節敗退,有人建議林肯總統帶領全體士兵禱告,求神站在北軍這邊,贏得戰事.林肯說他不會那樣禱告:「我從不求神站在我們這邊,只問我們是否站在神那邊?」深信只要北軍沒有違反聖經真理,上主會幫助北軍.林肯這番話對當下的戰事很有啟發,他對戰爭帶來的苦難的話也發人深省:「毋懷苦毒,廣施慈愛；趨仁就義,擇善固執,如承天意；努力完成手中的工作,包紮國家的傷口,照顧為國出征而捐軀者的遺孀和孤兒—竭盡所能,謀求在我們和列國中間的公義和永久的和平.」[1]

上主在戰爭

渡河與圍城.渡約但河時,是上主使水斷流.抬約櫃祭司把腳站在河水裡,平時漲過兩岸的水就斷了,水立成壘,以色列人行過干地,約櫃離開,河水復流,一如昔日出埃及過紅海的情況.這事使迦南人聞而喪膽.攻打耶利哥城,繞城七天,每天一次,聽祭司吹號角,不作一聲.第七天繞城七次後才大聲呼喊,城牆就塌下來了.這是耶和華的爭戰.

滅絕Herem, Anathema.約書亞記中記載很多戰爭,得勝的以色列人把整個城的人和牲畜殺掉,從現代人看來是野蠻的種族清洗,信徒提起都有點尷尬.但約書亞記不是現代社會學,人類學的書,其寫作角度是神學的,強調上帝對罪惡的決絕. 迦南人的宗教文化極其邪淫迷信,廟宇常駐男,女祭司,藉行淫進行宗教禮儀,獸交很普遍,又把兒童作祭牲獻祭.上主要清除這些罪惡,保護選民不受異教影響信仰.但選民仍然受當地異教文化影響拜偶像,禁之不絕,像以色列王嗎哪西:「在耶和華的兩院(外院,內院)中為天上的萬象築壇,並使他的兒子經火,又觀兆,用法術,立交鬼的,和行巫術的.多行耶和華眼中看為惡的事,惹動他的怒氣.」(王下二十一5, 6)

「一切人口,盡行殺絕」也是當時文化流行誇張的話.申命記七章說要把迦南人滅絕淨盡,但又說不可與他們結親,既然滅絕何來結親?「滅絕」一詞,希伯來文Herem,希臘文翻作Anathema,強調信仰純正,與錯謬,罪惡決絕.宗教改革時期,羅馬天主教會列舉教義,若不遵守,便被Anathema咒逐,詛咒並逐出教會.

因為罪惡滿盈,耶利哥城在號角聲歡呼聲中不攻而陷,全城被除滅.同樣當以色列人亞干偷了上主庫中的財物和收藏異教的衣飾,以色列人眾不敵寡,艾城一役慘敗,亞干全家被除滅.很明顯,約書亞記的戰爭角度是屬靈的.對罪絕不能妥協,上主是聖潔的,不得輕慢,無論以色列人或迦南人,一視同仁,與種族無關.

上主在拯救

弱勢妓女喇合與大城基遍勇士.事實上弱勢妓女喇合幫助探子全家得救；大城基遍勇士眾多,但聽見耶和華為以色列人行了奇事大事,棄假從真,用計謀說是遠方來的,慕名歸順,被接納為聖殿劈柴挑水；原居耶路撒冷的耶布斯人沒有被逐,與以色列人和平共存.事實上,後來我們看到不少迦南人散居在各支派境內.

以色列人進佔應許迦南地,是上主對居民的審判,也是對他們的拯救.喇合就看到以色列上帝才是真實可靠的神.「耶和華你們上帝,本是上天下地的上帝……要救活我的父母,弟兄,姐妹,和一切屬他們的,拯救我們性命不死.」(書二11, 13)

喇合認識迦南宗教諸偶像的虛妄,以色列的上帝才是天地的主.王爾德說:「我們都躺在陰溝裡,但有些人仰視晚上的星空.」.喇合不單得救,而且成為耶穌的先祖,上主並不偏待人.

進入迦南應許之地有宣教召命,所以離開埃及後,以色列人首先被引領來到西乃山,與上主立約,接受並學習律法.西乃山立約時,上主清楚指出他們有宣教的使命:「如今你們若實在聽從我的話,遵守我的約,就要在萬民中作屬我的子民,因為全地都是我的.你們要歸我作祭司的國度,為聖潔的國民,這些話你要告訴以色列人.」(出十九5, 6)

《耶穌愛迦南人》(Jesus Loves Canaanites)一書的作者勞謝爾(Randal Rauser (2 Cups Press, 2021)指出:「我們沒有藉口行使用暴力,合理化優越主義,或妖魔化局外人.反之,我們得到嚴厲的警告,我們屬靈的自覺和聖潔乏善可陳,也提醒我們局外人的聖潔,可作他山之石,幫助我們靈性重整革新.」[2]

4) 吃迦南地的土產

「逾越節的次日,他們吃了那地的出產……第二日嗎哪就止住了.」(書五11, 12)以色列人不單把信仰,律法傳給迦南人,叫他們悔改,棄假從真,此後迦南地就是他們的家,與迦南人生活在同一天空下.吃那地的出產,表示要為那城求平安,祈求風調雨順；關心農業灌溉系統,水利工程,糧食價格.信仰與生活不可分割,因為應許之地就是活出信仰之地.

我們要為居住的城平安,更要把上主的律法的精神滲入公共論壇.當公共空間充斥著俗世思想,自我主義,消費主義便成為金科玉律；當彼岸的光微弱暗淡,此岸的黑暗就被奉為理所當然.柏拉圖早就說過:「良善的人對政治冷漠,就要給比自己差的人管治,作為懲罰.」

(1 吉甲的覺醒—認清身份,不再飄流)

2  吉甲的覺醒: 宣講信息,使人得福)

吉甲的覺醒:甚麼是應許之地? 應許之地就是西乃敬拜之地 土地是當下敏感議題,有人為要奪回昔日「應許之地」,恢復舊日王國,興起戰爭,令聖靈塗炭.應許之地其實是西乃敬拜之地；何處有敬拜的生活,何處就是應許之地.出埃及有兩個原因: (1) 到西乃曠野敬拜上主.(2) 進入迦南應許之地.上主藉摩西在西乃山頒下律法,律法涵蓋敬拜,刑法,民法,倫理.使徒約翰的徒孫里昂主教愛任紐(Irenaeus)說:「活著的人是上主的榮耀,但人的生命離不開對上主的遠像(異像). 」(The glory of God is the living man, but the life of man is the vision of God.) (抗異端書4,20)

為何說活著的人是上主的榮耀呢?因為活著的人才能活出信仰,行公義好憐憫,藉行為歸榮耀給上主.為何說「但人的生命離不開對上主的遠像(異像)」呢?因為沒有上主就沒有靈性,當上主成為我遠象(異象)時,才有真正的滿足.榮休教宗拉辛格(Ratzinger)談到敬拜時說: 「敬拜存在目的是把這遠像(異象)傳遞,並使生命彰顯上主的榮耀.」 (Cult exists in order to communicate this vision and to give life in such a way that glory is given to God.) Ratzinger(2000) ,18. 崇拜是在主日擦新,調較對上主的異象,以致平日藉好行為榮耀上主.

所以西乃山並非中途竭息之所,乃安頓心靈應許之地.沒有西乃的敬拜,踏足之土如浮沙,所見之地像蜃樓,索然無趣,出埃及是為了來到西乃山.事實上,以色列被譽為經中之經的五經,不在應許之地所獲得,是在西乃曠野頒下的.那裡有敬拜,那裡就是應許之地.無論生活在何方,祇要我們敬拜上主,愛慕親近上主,那裡就是應許之地.而敬拜愛慕上主,能幫助我們放下行李,勇往直前,心靈不再遊蕩.

新愛牽魂,回絕舊戀

能夠放下行李走天路,因為是被上帝的愛所吸引.湯馬斯.錢伯(Thomas Chamber)有一篇著名講章冠名:The Expulsive Power of A New Affection 直譯「新歡排舊愛」,但這譯名會引來用情不專的聯想,且譯「新愛牽魂,回絕舊戀」.錢伯很聰明,把約壹二15「不要愛世界,和世界上的事.人若愛世界,愛父的心就不在他裡面了.」倒轉過來.能夠排拒世界的戀慕,就是對父的敬愛.地心吸力永遠想把飛機拉下來,惟有氣流動力使飛機在空中不墮下來.

苦難中的樂歌.在今天幾段經課中都流露對上主的愛.詩三十二1-11.「你是我藏身之處,你必保佑我脫離苦難,以得救的樂歌,四面環繞我.」(v.7)「惟獨倚靠耶和華的,必有慈愛四面環繞他.」(v.10)這篇認罪詩在預苦期必讀,但詩篇的氣氛不嚴苛決絕,強調上主慈愛願意赦免,並以得救樂歌環繞我們.

父親的饒恕.路十五1-3, 11b-32.浪子覺醒他是父親的兒子,不再以擁有多少家產定義自己.「父親,我得罪了天,又得罪了你.從今以後,我不配稱為你的兒子,把我當作一個雇工罷.」(路十五18, 19)相隔甚遠,他父親看見就動了慈心,跑去抱著他的頸項,連連與他親嘴.請注意,兒子見到父親這樣,沒有提及「把我當作雇工罷」了.是父親的愛,不見父親的蔭庇,令浪子受感動.使我想起一首歌「今要主自己」:

「前所要是祝福,今所要是主；前所憑是感覺,今靠主言語；

前所要是恩賜,今要賜恩主；前我尋求醫治,今要主自己.

前用己力苦試,今完全信靠；前要一半救恩,今要全得救.」

基督的愛.書信經課林後五16-21開始用「所以」,所以是上文所說:「原來基督的愛激勵我們.因我們想一人既替眾人死,眾人就都死了.並且他替眾人死,是叫那些活著的人,不再為自己而活,乃為替他們死而復活的主活.」(vv. 14, 15) 是愛驅策保羅一生忠於職分使人與基督和好.(林後五19)

今天也是愛,使我們回到吉甲,結束屬靈生命的飄流,讓基督的愛把我們吸引,心靈安頓,扎根永恆.

結論:朝聖旅程一如生命的旅程,沒有不能返回點.我們可以回到過去的地點,不能回到過去的時間.朝聖的征途是漫長的,有很多學習.上主呼召亞伯蘭,跟他立約也要一步一步向他啟示.亞伯拉罕創十二章蒙召離開吾珥,離開時還不知往那裡去.十五章在迦南地寄居,上主向其顯現說這就是賜給你和你後裔的地,然後立約:把祭牲切開,火爐從中間經過,以生命保證不會爽約.十七章才改名叫亞伯拉罕並守割禮.這是漸進的啟示,需要忍耐學習.他也迷失過,為了保存生命,欺騙埃及法王(創十二)和基拉耳王亞比米勒(創二十)說妻子撒拉是自己的妹子.但他始終沒有忘記身上的割禮(洗禮),回轉後,繼續先行.

預苦期是領洗的預備期,也是被逐和自我放逐離開教會的人,決心結束在外信仰飄流,回歸信仰的日子.吉甲的覺醒是屬靈身份的覺醒.讓我們重溫領洗的禮文:「我等接納此人進入基督之教會,將十字聖記號畫其額上為記,表明其自今以後,必真信奉釘十字架之基督,在人面前承認基督不以為恥,在基督旗下,向罪惡世俗魔鬼奮勇爭戰,一生盡忠,事奉基督.」 [1] “With malice toward none; with charity for all; with firmness in the right, as God gives us to see the right, let us strive on to finish the work we are in; to bind up the nation’s wounds; to care for him who shall have borne the battle, and for his widow and his orphan—to do all which may achieve and cherish a just and lasting peace among ourselves and with all nations.”  [2] “We have not a pretext for violence, a justification for triumphalism, or a precedent to demonize outsiders . Rather, we have a sober warning about our own lack of spiritual awareness and holiness and a call to recognize holiness in the outsider to our community even as we seek to addresss spiritual waywardness within our ranks.”  Randal Rauser, Jesus Loves Canaanites: Biblical Genocide in the Light of Moral Intuition . (Canada: 2 Cups Press, 2021), p. 271



\newpage

\section{十二門徒容易得知心一個也難求}
\label{sec:593}
\textbf{林豪恩先生 - 十二門徒容易得知心一個也難求}
\newline
\newline
\hyperref[sec:592]{< < < PREV SERMON < < <}
~
\hyperref[sec:toc]{[返目錄]}
~
\hyperref[sec:594]{> > > NEXT SERMON > > >}


講題:十二門徒容易得　知心一個也難求

講員:林豪恩先生

日期:2022年4月3日

讀這個故事的時候,大家會發現事件中有兩個突出的角色:一個是耶穌的門徒猶大,另一個是接待耶穌的馬利亞.由於這兩個角色的戲份相當重,讀者很自然地會把他們對比,於是地得出下列幾項教導.

1. 猶大貪心,馬利亞慷慨

馬利亞用作膏抹耶穌的香膏值三十兩銀子,這的確是一個大數目.如果真的是這個數目,有人估計大約是當時一個普通工人的一年薪金.馬利亞不是富貴人家,她把「一斤極貴的真哪噠香膏」用在耶穌身上,就算不是傾盡所有,都很可能是她財產的重要部份.與馬利亞對比的角色是門徒猶大,作者直接稱「他是個賊,帶著錢囊,常取其中所存的」.我們可以從猶大和馬利亞的對比推論:猶大貪心,馬利亞慷慨.這個推論給我們的反省是:我是貪心還是慷慨呢?於是給我們的教導就是:不要學習猶大的貪心,要學習馬利亞的慷慨.

2. 猶大自大,馬利亞謙卑

聖經說,馬利亞「抹耶穌的腳,又用自己頭髮去擦」.你可以想像馬利亞面對耶穌時的身體姿態是如何的?很可能是下跪在耶穌之前.門徒猶大呢?他出言質問馬利亞的行動,可以想像他是站著的,自以為義地伸出指責的指頭.從這項對比,我們可以得到的感覺是:馬利亞是謙卑的,而猶大是自大的.這個感覺給我們的反省是:我是自大還是謙卑呢?於是給我們的教導是:不要學猶大的自大,要學馬利亞的謙卑.

3. 猶大抽離,馬利亞親近

第三項我們可以從中觀察到的是,馬利亞親近耶穌,門徒猶大與耶穌保持距離.如果你要把這個故事用圖畫表達,你的畫可能像現存的大多數一樣,畫面出現的是馬利亞抹耶穌的腳,遠處站著指責別人的猶大.這對比給我們的印象是:馬利亞親近耶穌,猶大遠離耶穌.於是給我們的反省是:我親近耶穌還是遠離耶穌呢?給我們的教導是:要學馬利亞一樣親近耶穌,不要學猶大一樣遠離耶穌.

4.誰能明白他?

我們得出上述的反省及教導,可能因為把焦點集中在馬利亞和猶大身上.不過,我們尚未觀察耶穌的反應.如果耶穌才是事件的當事人,當事人自己如何看待這事情呢?耶穌當時是如何回應的呢?

12:7耶穌說:「由她吧!她是為我安葬之日存留的.

12:8因為常有窮人和你們同在,只是你們不常有我.」

耶穌沒有反駁猶大提出的意見和理由.他沒有好像我剛剛所做的,指責猶大的不是,也沒有證明馬利亞做的合理,只是說:「由她吧!」聽到耶穌說「由她吧!」,我們能夠想像當時耶穌的心情是怎樣嗎?我們知道耶穌處於怎樣的狀況嗎?我們知道耶穌下一步是怎樣的路嗎?

上文11:53指出:「從那日起,他們就商議要殺耶穌……那時,祭司長和法利賽人早已吩咐說,若有人知道耶穌在哪裡,就要報明,好去拿他.」

下文12:10馬上報導:「但祭司長商議連拉撒路也要殺了；」

「殺」這個字在上下文出現,包圍著目前這個故事.殺機四伏,人都將要死了,有甚麼想念呢?臨死之時,甚麼才重要呢?

二十多年前有一部電影名叫《情約今生》(Meet Joe Black)上影,這部電影是改篇自1934年的舞台劇電影《死神度假》(Death Takes a Holiday).故事的開始是這樣的:有一位富可敵國,在商場上吒叱風雲大亨,在他的六十五歲生日之前,遇上死神來取其性命.大亨為了換取多幾天的壽命,提出招待死神在人世間「度假」,同時答應讓死神與他如影隨形但不可以揭露死神的身分,直至死神假期結束帶他走.生命已走到盡頭,這位有著人世間最大的財富和最大的權勢的大亨當天立即做的事情是甚麼呢?他召集家人共進晚餐.第二天,他又召集家人共進晚餐.家人都感到奇怪,但因為與死神有秘密協議,不能揭露死神這位客人的身分,也不能透露自己即將被死神帶走.看來,生命在倒數的日子,就算是吒叱風雲的大亨,最珍惜的仍然是與自己有關係的人；最後的日子,渴望與自己所愛的人及愛自己的人共同渡過.

第五波新冠疫情奪去不少長者性命,令人惋惜.然而,更傷感的是,家人和離世者不能如常地渡過最後的日子,甚至沒有機會道別.一般情況下,如果病人有生命危險,醫院都會通知家人,讓家人來與臨別者見最後一面.那個時刻,家人會帶甚麼去見最後一面呢?難道是成績表嗎?難道是金錢嗎?難道是「分身家」的文件嗎?當然不會,帶自己去已經足夠了.在這一刻,見到家人不離不棄,感到自己被珍惜,被懷念,知道自己心中有他們,也知道他們心中有自己,平安道別,釋懷離去,就是最珍貴的了.

12:7耶穌說:「由她吧!她是為我安葬之日存留的.

12:8因為常有窮人和你們同在,只是你們不常有我.」

雖然耶穌已經明言自己會被殺害,門徒也知道猶太人要拿石頭打他,同時,作者也刻意在前後文突出殺機,但在那一刻,誰會認真看待耶穌正在一步一步地走向死亡呢?誰會體諒耶穌的處境和心情呢?誰會珍惜耶穌呢?

5. 耶穌是我的知心友,我是耶穌的知心友嗎?

體諒耶穌的處境和心情?對門徒來說,耶穌是彌賽亞,是來拯救他們的,怎會需要體諒呢?珍惜耶穌?對門徒來說,耶穌是滿有能力的,門徒還有很多想望等待耶穌實現,從沒有想像過以「珍惜」這角度來對耶穌.如果門徒沒有這樣的想法,也是可以明白的,難道我們不也是一樣嗎?詩歌常唱:「神是我們的力量,是我們的盾牌」,即是說,神幫助我們戰勝敵人,神為我們擋風擋雨擋子彈.神是甚麼呢?神可能被我們視為一種能夠保護我們,滿足我們需要的工具.

讓我們問一問自己,我們又把神當作甚麼呢?或者有人會說:「我不會把神當作工具,我把神當作朋友.」曾幾何時,的確有一首本地詩歌名為「耶穌是我的知心友」,歌詞的第一和第二段及副歌是這樣的:

耶穌是我的知心友,他知道我深透,

細心帶領我向天路走,他牽扶著我手.

我向他傾訴我憂與愁,我信他必去困憂,

永記他因我將寶血流,我靈得拯救.

耶穌是我的知心友,明我祈求或苦與憂,

神到處將我施恩救,慈愛永遠伴我走.

在這世界道路多歧途,誰明白我主恩高,

曾為我要被死煎與熬,神奇恩我仰慕.

把耶穌當作朋友,會不會貶低耶穌抬高自己呢?其實,這份友誼是耶穌自己主動建立的:

「以後我不再稱你們為僕人,因僕人不知道主人所做的事.我乃稱你們為朋友；因我從我父所聽見的,已經都告訴你們了.」(約翰福音15:15)

朋友關係應該是雙向的,「耶穌是我的知心友」,但我是耶穌的知心友嗎?「他知道我深透」,但我有興趣知道他嗎?或許,有人已經把系統神學中的「神論」讀得滾瓜爛熟,不過,就算如此,也不等於就是朋友；就正如有人熟讀了你的履歷,就是你的朋友了嗎?

詩歌提及「耶穌是我的知心友」,因為「我向他傾訴我憂與愁」,他「明我祈求或苦與憂」.從另一個方向來看,作為朋友,我們有理會他的憂與愁嗎?我們與他分憂嗎?神也有憂愁嗎?聖經中有這些片段,例如:

創世記 6:5「耶和華見人在地上罪惡很大,終日所思想的盡都是惡,6耶和華就後悔造人在地上,心中憂傷.」

馬太福音26:36耶穌同門徒來到一個地方,名叫客西馬尼,就對他們說:「你們坐在這裡,等我到那邊去禱告.」37於是帶著彼得和西庇太的兩個兒子同去,就憂愁起來,極其難過,38便對他們說:「我心裡甚是憂傷,幾乎要死……」

耶穌在拉撒路家裡吃過這餐飯,第二天就上耶路撒冷了.這次進入耶路撒冷,是踏上死亡之路.在這路上,有誰同行呢?有誰相伴呢?感到孤寂,不一定是身邊沒有人,而是沒有人洞悉你的處境,沒有人明白你的心情.就正如那才提及的電影《情約今生》中的大亨,身邊當然環繞著很多人,催逼他開會,向他提交公司發展大計,建議收購合併計劃,這些都不是不好,也不是沒有道理,只是不切合處境和心境.就在進入耶路撒冷的前夕,竟然有一偶爾拜訪的女子,而不是朝夕共對的門徒,彷彿洞悉自己的命途,拿香膏來膏抹自己,相信耶穌那一刻的感受是強烈的,他深深感到馬利亞的行動「是為我安葬之日存留的」.

今天(4月3日)是預苦期第五主日,距離耶穌在客西馬尼被拘捕的日子越來越近了,也距離耶穌被侮辱,被枉法,被釘死在十字架上的時候越來越近了.如果我們喜歡唱「耶穌是我的知心友」這首詩歌,作為耶穌的朋友,讓我們與耶穌一同上路,與他同行這些艱難的日子.

預苦期是記念耶穌的苦路的日子,耶穌所受的苦,其實是世人的苦.當我們記念耶穌被拘捕,被侮辱,被枉法的時候,也記念世界上眾多被拘捕,被侮辱,被枉法的受苦者；當我們與耶穌同行步向死亡之路的時候,也讓我們與世間眾多因身體狀況或戰火無情中步向死亡的人同行.



\newpage

\section{人生如驢驢如人生}
\label{sec:594}
\textbf{梁元生教授 - 人生如驢,驢如人生}
\newline
\newline
\hyperref[sec:593]{< < < PREV SERMON < < <}
~
\hyperref[sec:toc]{[返目錄]}
~
\hyperref[sec:595]{> > > NEXT SERMON > > >}


路加福音 19:28-40
\newline
\begin{longtable}{cl}
\hline
\hline
章節 & 經文 (和合本修訂版)\\
\hline
19:28 & \begin{tabularx}{0.7\textwidth}{X} 耶穌說完了這些話，就走在前面，上耶路撒冷去。 \end{tabularx} \\ \\ \relax
19:29 & \begin{tabularx}{0.7\textwidth}{X} 快到伯法其和伯大尼，在名叫橄欖山的地方，他打發兩個門徒， \end{tabularx} \\ \\ \relax
19:30 & \begin{tabularx}{0.7\textwidth}{X} 說：「你們往對面村子裡去，進去的時候會看見一匹驢駒拴在那裡，是從來沒有人騎過的，把牠解開，牽來。 \end{tabularx} \\ \\ \relax
19:31 & \begin{tabularx}{0.7\textwidth}{X} 若有人問為甚麼解開牠，你們就這樣說：『主要用牠。』」 \end{tabularx} \\ \\ \relax
19:32 & \begin{tabularx}{0.7\textwidth}{X} 被打發的人去了，所遇見的正如耶穌對他們所說的。 \end{tabularx} \\ \\ \relax
19:33 & \begin{tabularx}{0.7\textwidth}{X} 他們解開驢駒的時候，主人問他們：「為甚麼解開驢駒？」 \end{tabularx} \\ \\ \relax
19:34 & \begin{tabularx}{0.7\textwidth}{X} 他們說：「主要用牠。」 \end{tabularx} \\ \\ \relax
19:35 & \begin{tabularx}{0.7\textwidth}{X} 他們把驢駒牽到耶穌那裡，把自己的衣服搭在上面，扶耶穌騎上。 \end{tabularx} \\ \\ \relax
19:36 & \begin{tabularx}{0.7\textwidth}{X} 他前進的時候，眾人把衣服鋪在路上。 \end{tabularx} \\ \\ \relax
19:37 & \begin{tabularx}{0.7\textwidth}{X} 他將近耶路撒冷，正下橄欖山的時候，一大群門徒因所見過的一切異能，都歡呼起來，大聲讚美神， \end{tabularx} \\ \\ \relax
19:38 & \begin{tabularx}{0.7\textwidth}{X} 說：「奉主名來的王是應當稱頌的！在天上有和平；在至高之處有榮光。」 \end{tabularx} \\ \\ \relax
19:39 & \begin{tabularx}{0.7\textwidth}{X} 人群中有幾個法利賽人對耶穌說：「老師，責備你的門徒吧！」 \end{tabularx} \\ \\ \relax
19:40 & \begin{tabularx}{0.7\textwidth}{X} 耶穌回答：「我告訴你們，若是這些人閉口不說，石頭也要呼叫起來。」 \end{tabularx} \\ \\
[1ex]
\hline
\hline
\end{longtable}
講題:人生如驢,驢如人生

經文:路加福音19: 28-40

講員:梁元生教授

經文:路加福音19: 28-40

日期:2022年4月10日

「耶穌說完這話,就在前面走,上耶路撒冷去.將近伯法其和伯大尼,在一座山名叫橄欖山那裡,就打發兩個門徒,說:你們往對面村子裡去,進去的時候,必看見一匹驢駒三在那裡,是從來沒有人騎過的,可以解開牽來.若有人問為甚麼解牠,你們就說:主要用牠.打發的人去了,所遇見的正如耶穌所說的.他們解驢駒的時候,主人問他們說:解驢駒做甚麼?他們說:主要用牠.他們牽到耶穌那裡,把自己的衣服搭在上面,扶著耶穌騎上.走的時候,眾人把衣服鋪在路上.將近耶路撒冷,正下橄欖山的時候,眾門徒因所見過得的一切異能,都歡樂起來,大聲讚美神,說:奉主名來的王是應當稱頌的!在天上有和平,在至高之處有榮光.眾人中有幾個法利賽人對耶穌說:夫子,責備你的門徒吧!耶穌說:若是他們閉口不說,這些石頭必要呼叫起來.」

一. 前面的話

各位主內弟兄姊妹,各位線上朋友:

感謝主讓我們在這疫症期間仍然能夠有主日崇拜,能夠有線上聚會,這是主給我們這一代人特別的恩典.雖然如此,我站在講臺上,看著滿堂的空凳,心裡還是感到很大唏噓.記得陸機文賦中有言:課虛無以責有,叩寂寞而求音.我希望今天我所要講的並非虛無縹緲的事,我也希望這寂寞的大堂和空虛的雲端背後,會有一點反響和回音.因為我今天要分享的題目,是有關“人生”的,不論或長或短,這是人人都擁有的,故此也是所有人皆會面對的.

屈子楚辭中有這樣一句話:“悲晨曦之易夕,感人生之長勤.”用現代白話來翻譯,可以這樣說,“我哀歎光陰流逝的迅速,感歎一生的勞碌工作.”在疫情之中反省人生,更加感受工作之沉重和心靈之重壓,以及光陰之急速和生命之短暫.我很樂意從自己七十多年的經歷和大家回顧人生,但今天不會這樣做,而是借鏡他者只經驗.今天所讀的福音經課,即路加福音十九章28-40節的那一段,記載著耶穌最後的一段人生路.由伯法其和伯大尼的橄欖山出發,只需再走一段路,就是耶路撒冷.這段路程,是耶穌的最後一段人生路.耶穌進城之後,就是受審和被釘十字架,結束了祂在世三十三年的人生.這是一段十分弔詭的文字:耶穌榮耀進城,進城之前受到夾道歡迎,但進城之後,耶穌就走向死亡,背著十字架,走向觸髏地.今天就讓我們和耶穌一同走這最後一程人生路,先從他騎驢進城開始.

這段經文展示出一幅熱鬧的場景.在這段描述中,我們看到有著不同身份和扮演著不同角色的一群人,有夾道歡呼的眾人,路人,也有在旁邊觀看熱鬧的閒人；當中有自以為義的法利賽人.當然,還有簇擁著耶穌歡然大叫“和撒哪,和撒哪,奉主名來的,是應當稱頌的!”的門徒和跟隨者.

此外,我想應該還有有驢子的主人(根據聖經注釋家言,主人這個詞在這裡是多數的,即驢子的主人可能多過一個人)和他的驢子.今天,我們在這裡用這段經文來思考人生,喘摸人生的意義,我建議我們暫且不看這些形形式式的人,而是注目看那一頭驢子.以驢子看人生!

二. 人生如驢

人生如驢.許多人聽到這個比喻,不禁歎息和唏噓.像驢子的人生有甚麼好?

在一般人心目中,驢子是蠢笨的.以前我們會常聽到罵人的話:你這蠢驢!近日疫情襲港,常常關在家中,百無聊賴,看多了電視劇集.在一部中國古裝片和一部韓劇中,我都聽到“你這蠢驢!”這句罵人的話,而且都是罵“老公”的.除了少數男人會羡慕“潘,驢,鄧,小,閑”中的驢子,一般人對驢子的看法都比較負面.驢子除了愚蠢之外,還很執拗!驢子代表固執,一成不變,不夠靈活,不肯變通.舊時中國有句老話,是貶視湖南人的,就說這人是“湖南驢子”,意思是說這是個非常固執,不易變通,難以相處的人.

臺灣的宋楚瑜是湖南人,數次出選競逐總統之位,幾次都明知不會成功,仍然一意孤行,不願妥協,臺灣報章稱他為“湖南驢子”.

驢子除了愚蠢,固執之外,還代表著“辛苦命”!驢子的工作是拖車,是犁田,是負重,是運輸,是長時間的工作,是忙碌,勞累,以及十分辛苦的生涯.我們常聽到人說一句廣東話“捱驢仔”,就是“辛苦命”的意思.

講IQ,驢子屬於蠢鈍之列,講性格,屬於固執,頑梗之徒,講生活,講生命,屬於勞碌命和辛苦命.是的,愚笨,固執,辛苦.這就是驢子人生(donkey’s life)!

然而,聖經中的驢子所寓意的人生,卻和我們剛才所說的並不一樣.今天讓我們從聖經中看幾匹驢子,並且從這些驢子的故事思考一下其對人生的意義.

三. 耶穌的驢:“主要用他!”

讓我們先來看剛才讀過一遍的福音經課:路加福音第十九章28-40節.這是棕枝主日的經文,主要是描述主耶穌和門徒進入耶路撒冷聖城是的情景.這情景顯得甚為熱鬧,氣氛也好像非常歡樂.人們一邊拿著棕樹枝,夾道歡迎耶穌,一邊呼叫說:“和撒哪,和撒哪,奉主名來的,是應當稱頌的!”

在耶穌和門徒“榮耀進城”之前,聖經也記載了一段發生在橄欖山下的小插曲,那就是耶穌吩咐兩個門徒去尋找進城的坐騎.經文這樣說:(耶穌)打發兩個門徒,說:“你們往對面村子裡去,進去的時候,必看見一匹驢駒拴在那裡,是從來沒有人騎過的,可以解開牽來.若有人問為甚麼解牠,你們就說:“主要用牠.”

這段經文對驢子的描述是怎樣的?第一,“從來沒有人騎過”,第二,“被拴(束綁)在那裡”,第三,需要“解開,牽來”,然後最後說:“主要用牠”.我們一般讀經史知識注意最後那句話:主要用牠.“主要用牠”當然是一個基督徒人生最重要的使命和意義.但我們從這段經文可以看到“為主所用”之前,需要經過“三部曲”,即是三重的考驗.首先,是從來沒有被人騎過的驢子,即是要一直保持赤字之心,或者如聖經他處所記,是“初熟的果子”,意思是有著純潔的心靈,聖潔的生活；其次,要“解開”,不再“被拴住”.意思是要脫離罪惡的捆綁和世俗的纏累.想想有哪些事,那些人,是我們放不下,解不開的,都要先處理好；第三,要“牽來”,即是跟隨門徒來到耶穌跟前,而不是盲目亂奔,成為狂野蠻驢.從解開,不再被拴住,到被牽來,我們看到一放一收,所謂張弛有道的生活藝術.解開了,卻不放縱；離開舊主人,卻有新的牽引者.

最後是“主要用牠”!這句話在這段經文中前後出現兩次,顯示出其重要性.一個基督徒生命的真正意義和價值,就是能夠為主所用.縱然,他的用途或可用之處可能只有短暫的一刻時光,和驢子負者耶穌進城一樣,只走了一段很短的路程,而不是像一些門徒那樣跟隨著耶穌走過大半生.雖然如此,耶穌騎驢進城,得著眾人夾道歡迎,就像迎接君王一樣,使耶穌得到應得的榮耀.這一刻的光陰,這一段短距離的路程,對耶穌,對驢子,都是生命中極為重要,和值得銘記的.今天我們來看驢子的故事,有沒有反思:在我們的人生當中,有沒有這一刻,有沒有這一段路,是耶穌和我共同走過的?

而在與主同行時,讓祂騎在背上,把祂高舉,使祂得著榮耀!祂要用我,以我為榮.若然的話,我就不枉此生.

四. 亞伯拉罕的驢:為主作證

另外,聖經其他地方也提到兩匹驢子:一在創世記二十二章1-8節,另一處在民數記二十二章:21-33節.在創世記二十二章那裡記載的是上帝考驗亞伯拉罕的故事.第2節說:“上帝說:你要帶你的兒子,就是你所愛的獨子以撒,往摩利亞地去,在我只是你的一座山上,把他獻為燔祭.亞伯拉罕清早起來,預備了驢,帶著跟他一起的兩個僕人和他的兒子以撒,劈好了燔祭的柴,就起身往上帝指示他的地方去了.”這段經文中的故事和故事中的主角,都是為人所熟悉的.主角是兩父子,即亞伯拉罕和他的兒子以撒,故事是父親要把兒子帶到山上,殺死獻祭,到了最緊張的一刻被叫停,耶和華神為他預備了獻祭用的羊羔,保留了以撒的性命.在這個緊張刺激的故事中,大家注目於亞伯拉罕對神的信心和順服,對幼童以撒無助和無知的同情,以及對上帝的及時介入,預備和解救的鬆一口氣.

我們讀過這段經文的,都會記得這個動人和震撼的故事,都會記住“耶和華以勒”(耶和華必預備)的箴言和教訓.然而,在人生的舞臺上,除了這些主角,除了感動人心的故事,以及震撼的情節和畫面,還有許多配角,許多旁觀者,許多不為人注意的閒人.在這個故事裡,有兩個隨行的僕人,還有那一匹驢子!

聖經說:“亞伯拉罕清早起來,預備了驢,帶著跟他一起的兩個僕人和他的兒子以撒,劈好了燔祭的柴,就起身往上帝只是他的地方去了.到了第三日,亞伯拉罕舉目遙望那地方.亞伯拉罕對他的僕人說:“你們和驢留在這裡,我和孩子要去那裡敬拜,然後回到你們這裡來.”這段經文告訴我們,僕人和驢子都沒有在獻祭的現場.但他們都參與了這場重要的祭祀大典,而且分別擔任一些預備的工作.首先,是在亞伯拉罕家裡準備燔祭所用的柴.我猜想僕人可能會幫忙劈柴,然後又驢子背負著這些燔祭的柴禾其他用品,清早出發.聖經說他們行了三天的路程,才到了上帝指示的地方.路途上沒有甚麼驚險,也沒有意外.但三天的路程,而且走的是陡峭不平的山路.相信這段路不會是舒適和歡愉的享受之旅.到了第三日,亞伯拉罕舉目遙望那地方,對僕人說:“你們和驢留在這裡,我和孩子要去那裡敬拜,然後回到你們這裡來.”換言之,驢子和僕人到了目的地,只在遠處遙望,看著一個意想不到的緊張場面在眼前展開:亞伯拉罕手裡拿著火與刀,把燔祭的柴堆放在兒子以撒身上.來到那地方,築起祭壇,把柴擺好,把兒子綁起放置柴上,然後亞伯拉罕伸手拔刀,要殺兒子.就在這電光石火之間,有耶和華的使者說:“亞伯拉罕,不要下手!”並且在灌木叢中,預備好一隻公綿羊,作為獻祭的祭品.成為“在耶和華的山上,祂必預備”的明證.僕人和驢子雖然不在事發現場,但他們在遙遠可見的地方守候,看見這一切事情發生,也就成為“耶和華必預備”的神跡的第一手見證人.

我們的人生,難免有時扮演配角,或是旁觀者的角色.但能夠為主作證,使神跡奇事,和主的承諾—“只要有信心,神必有預備”的承諾,得以流傳及宣揚開去,也就不枉此生.不一定要做亞伯拉罕,也不一定要做以撒,做僕人,或做驢子,也可以找到生存的意義.

五. 巴蘭的驢:不肯前行

民數記二十二章:21-33節那裡,提到另外一匹驢子,先知巴蘭的驢子.聖經這樣記載:“巴蘭早晨起來,備了驢,就和摩押的官員一同去了.上帝因他去就怒氣發作；耶和華的使者站在路中間敵對他.他騎著驢,有兩個僕人跟隨他.驢看見耶和華的使者站在路中間,手裡有拔出來的刀,就離開了路,岔入田間.巴蘭就打驢,要牠回到路上.耶和華的使者站在葡萄園的窄路上,這邊有牆,那邊也有牆.驢看見耶和華的使者,就往牆擠去,把巴蘭的腳擠到牆上,巴蘭再打驢；耶和華的使者有往前去,站在狹窄的地方,那裡左右都無路可轉.驢看見耶和華的使者,就伏在巴蘭底下.巴蘭怒氣發作,用杖打驢.耶和華使驢開口,對巴蘭說:我向你做了甚麼,你竟打我這三次呢?”

這是一匹有趣而且奇怪的會說話的驢子.故事說牠的主人巴蘭騎著牠到摩押去,招來上帝的怒氣,派使者在路上攔阻他.上面的經文記載著巴蘭和驢子出行的過程,有著非常生動的描寫:(巴蘭)騎著驢,兩個僕人跟隨他.驢看見耶和華的使者站在路中間,手裡有拔出來的刀,就離開了路,岔入田間.巴蘭就打驢,要牠回到路上.耶和華的使者站在葡萄園的窄路上,這邊有牆,那邊也有牆.驢看見耶和華的使者,就往牆擠去,把巴蘭的腳擠到牆上,巴蘭再打驢；耶和華的使者有往前去,站在狹窄的地方,那裡左右都無路可轉.

驢看見耶和華的使者,就伏在巴蘭底下.

巴蘭怒氣發作,用杖打驢.耶和華使驢開口,對巴蘭說:我向你做了甚麼,你竟打我這三次呢?巴蘭對驢說:因為你戲弄我,我恨不得手中有刀,現在就把你殺了.驢對巴蘭說:我不是你從小直到今天所起的驢嗎?我平時有這樣待過你嗎?巴蘭說:沒有.

從巴蘭和驢子的對答中,我們看到二者的不同.驢子三次不肯往前行,因為牠看見神的使者拔刀在前面阻擋；先知巴蘭卻看不見前面有神的使者,三次出力打驢,並令牠受傷.關於墮落先知巴蘭和驢的故事,聖經學者解釋不多.不過,我對這驢子及其寓意的人生卻是極為同情的.幾十年來從小至今一直忠心服侍主人,未嘗違背,從不忤逆.但落得如此下場,被冤屈,被打罵!原因只有一個:主人巴蘭的墮落,眼睛蒙蔽,服侍假神,看不見上帝的使者,看不到刀和死亡在前面等待著他.我們的人生也背負著很多的傳統和習慣(conventions),會蒙蔽我們的眼睛,使我們看不見神的使者,就如巴蘭先知一樣.驢子雖然受冤屈,雖然被責打,但牠保持眼睛明亮,定睛神的使者,至終化險為夷,平安通過.

六. 結語

棕枝主日,是君王榮耀進城的日子,也是耶穌步向死亡的開始.是誰陪伴著祂走這段最後的人生路?一匹驢子!一匹被解開,被捆鎖,被挑選的驢子.“主要用牠”,那是基督徒人生的座右銘!

人生如驢,如亞伯拉罕的驢,背負祭品,背負柴枝,走三天的路,為了見證一件奇跡:“耶和華必預備!”

巴蘭的驢,寓意忠心服侍的人生亦難免受冤,被屈,受責打,被欺負；只要眼目定睛於主,最終可以笑到最後,連主人也被感動而改變.

人生如驢,驢如人生!

原來聖經中的每一篇章,每一個人,每一件事,每棵植物,每只動物,都對人生有所啟示.願大家多讀聖經.因為聖經的話,使我們腳前的燈,路上的光!



\newpage

\section{路漫漫其曲折兮}
\label{sec:595}
\textbf{鄧瑞強博士 - 路漫漫其曲折兮}
\newline
\newline
\hyperref[sec:594]{< < < PREV SERMON < < <}
~
\hyperref[sec:toc]{[返目錄]}
~
\hyperref[sec:596]{> > > NEXT SERMON > > >}


民數記 14:1-10
\newline
\begin{longtable}{cl}
\hline
\hline
章節 & 經文 (和合本修訂版)\\
\hline
14:1 & \begin{tabularx}{0.7\textwidth}{X} 全會眾大聲喧嚷，那夜百姓哭號。 \end{tabularx} \\ \\ \relax
14:2 & \begin{tabularx}{0.7\textwidth}{X} 以色列眾人向摩西和亞倫發怨言，全會眾對他們說：「我們寧願死在埃及地，寧願死在這曠野！ \end{tabularx} \\ \\ \relax
14:3 & \begin{tabularx}{0.7\textwidth}{X} 耶和華為甚麼要把我們領到那地，讓我們倒在刀下呢？我們的妻子和孩子必成為擄物。我們回埃及去豈不更好嗎？」 \end{tabularx} \\ \\ \relax
14:4 & \begin{tabularx}{0.7\textwidth}{X} 他們彼此說：「我們不如選一個領袖，回埃及去吧！」 \end{tabularx} \\ \\ \relax
14:5 & \begin{tabularx}{0.7\textwidth}{X} 摩西和亞倫在以色列全會眾面前臉伏於地。 \end{tabularx} \\ \\ \relax
14:6 & \begin{tabularx}{0.7\textwidth}{X} 窺探那地的人中，嫩的兒子約書亞和耶孚尼的兒子迦勒撕裂衣服， \end{tabularx} \\ \\ \relax
14:7 & \begin{tabularx}{0.7\textwidth}{X} 對以色列全會眾說：「我們所走過、所窺探之地是極美之地。 \end{tabularx} \\ \\ \relax
14:8 & \begin{tabularx}{0.7\textwidth}{X} 耶和華若喜愛我們，就必領我們進入那地，把這流奶與蜜之地賜給我們。 \end{tabularx} \\ \\ \relax
14:9 & \begin{tabularx}{0.7\textwidth}{X} 但你們不可背叛耶和華，也不要怕那地的百姓，因為他們是我們的食物。保護他們的已經離開他們，耶和華卻與我們同在。不要怕他們！」 \end{tabularx} \\ \\ \relax
14:10 & \begin{tabularx}{0.7\textwidth}{X} 當全會眾正說著要拿石頭打死他們的時候，耶和華的榮光在會幕中向以色列眾人顯現。 \end{tabularx} \\ \\
[1ex]
\hline
\hline
\end{longtable}
講題:路漫漫其曲折兮

經文:民數記 14: 1-10

講員:鄧瑞強博士

日期:2022年4月17日

我講道的經文是在民數記,但今日是復活節,我不得不講幾句關於主耶穌復活的話.

主耶穌復活,是宇宙的頭等大事；復活節,記念的是新世界的誕生.其重要性猶如宇宙大爆炸,大爆炸之後,出現的,是一個全新的世界.主耶穌復活,不只是說,有個人,死了之後,又復活了.然後,世界又如常繼續運作下去.我們常常聽到這種古怪新奇的新聞,比如說:有個人醫生斷定已死亡,但過了一會,忽又復甦過來.聽到這類新聞,我們嘖嘖稱奇,然後,繼續吃飯,飲酒,打麻雀.主耶穌復活不是這種無傷大雅的新聞,這是宇宙的頭等大事.主耶穌的復活,克復了人類最大的恐懼和敵人——死亡.死亡,將人的存有,化作虛無.死亡,比任何人都強大.無論你是誰,都要在死亡面前低頭.它嘲笑你,戲弄你,刪除你.死亡背後,是罪的權勢.羅馬書6:23說:「罪的工價乃是死」.死亡,是罪的後果與代價.死亡背後,是罪的力量.罪的背後,是魔鬼,牠是一切邪惡的總匯.主耶穌的復活,是克服死亡,罪惡,魔鬼,是克服了一連串取消生命的力量.主耶穌的復活,是神的光明驅散一切黑暗.猶如天地初開時那神聖的一刻,在那神聖的一刻,神說:「要有光」,便有了光,然後,將黑暗驅散了.

主耶穌復活後,整個宇宙已不再一樣了,你和我都不再一樣了.死亡,罪惡,魔鬼,已失去它們終極的力量.它們的鎖鍊,已被主的復活打碎了.我們自由了.我們是新的創造.哥林多後書5:17說:「若有人在基督裡,他就是新造的人(可譯為「新的創造」),舊事已過,都變成新的了.」在復活的基督裡,我們進入了一個全新的人生,是散發著光明的人生,是具有永恆意義的人生.

復活節,是新宇宙誕生的紀念日.希望我們明白這全新世界,全新生命的意義.不要只是慶祝幾天,然後,生命又回到黑暗與死亡的權勢下.正如:不要入到一隨便你吃的美食天堂,只讚歎幾聲,然後,什麼也沒有吃便離開了,餓著肚子,又埋怨自己在捱餓.

回到今日講道的書卷:民數記.

我們溫習一下之前的經卷.創世記講到,神揀選一個人,作歷史的點燈人,在人間留住光明.這人是亞伯拉罕,神與他立約.這約有三個元素,一,是無數後裔；二,是歸屬神的生命質素；三,是土地.為何需要有無數後裔?因為一個人力量單薄,歷史太黑暗,需要一整個民族來燃亮光明.沒有足夠後裔,無法形聖民族.在創世記裡,亞伯拉罕,以撒,雅各的最大考驗,是如何生仔的問題.到了出埃及記時,後裔的問題已解決,他們繁衍甚多,甚至成為埃及人頭痛的問題.現在,重點在如何建立一個民族?要建立一個民族,必須有一全民族刻骨銘心的經驗,這將要成為民族的印記,成為他們世世代代不滅的集體回憶.出埃及的經驗,是神給這民族的難忘經驗.整整一個民族,因著歷史的種種因由,成了苦不堪言的奴隸.然後,在種種神蹟下,他們過紅海,出埃及,成了全新的自由人.他們以逾越節去紀念這新的創造,這是他們的復活節.主耶穌的受死與復活,也是發生在逾越節的日子裡.直到今日,猶太人的逾越節與基督教的受難節與復活節,常常是在同一日子.以色列民族,是一特殊民族,以回憶一救贖事件,作為他們永恆的身分象徵.

民族的記號有了,跟著要處理的,是在這民族裡陶造一「歸屬神的生命質素」.這就是利未記的主題.很繁瑣的誡命與規條,是讓他們與別不同.如何讓「歸屬神的生命質素」深深地銘刻在這民族的生命裡呢?如何讓「歸屬神的生命質素」成為他們靈魂的紋理呢?神讓他們在曠野裡受考驗.曠野,是一雕刻生命的地方.曠野,什麼也沒有,最能考驗人是否全心倚靠神.《民數記》在希伯來文聖經裡,經卷的名稱的意思,就是＜在曠野＞.今日,我們看看這民族在曠野受的考驗.

離開埃及,進入迦南地之前,中間經過曠野.在曠野,人的需要變得更強烈.人性受到的考驗,變得更明顯.在曠野,無食物吃,無水可飲,人變得暴燥,進而挑戰摩西的領導,這是經常發生的事情.《民數記》不斷記載這些事件.後來,耶穌出道時,刻意進入這曠野,經受這種沒有食物的考驗.然後,讓人明白:「人活著不是單靠食物,乃是靠神口裡所出的一切話.」在人性經受磨難時,多少人能依靠神口裡所出的話而克服艱難呢?曠野,的確是人性最大的考驗場所,因此,也同時是雕琢「歸屬神的生命質素」的最佳地方.

好快,以色列民便去到迦南地的邊界.他們決定派探子去探探虛實,他們十二個支派,每支派派一領袖,潛入迦南地,看看那地如何,是好是壞；看看那地的人如何,是強是弱.探子回來報信,那地是美地,是物產豐茂之地.但同時,那地的人高大強悍.他們的結論是:「我們所窺探,經過之地是吞吃居民之地,我們在那裡所看見的人民都身量高大.我們在那裡看見亞衲族人,就是偉人；他們是偉人的後裔.據我們看,自己就如蚱蜢一樣；據他們看,我們也是如此.」(民13:32-33)

探子如此回報,當然帶來災難性的後果.

今日的講道經文,接續著上述故事,請看:民數記14:1-10.

14:1　當下,全會眾大聲喧嚷；那夜百姓都哭號.

14:2　以色列眾人向摩西,亞倫發怨言；全會眾對他們說:「巴不得我們早死在埃及地,或是死在這曠野.

14:3　耶和華為什麼把我們領到那地,使我們倒在刀下呢?我們的妻子和孩子必被擄掠.我們回埃及去豈不好嗎?」

14:4　眾人彼此說:「我們不如立一個首領回埃及去吧!」

14:5　摩西,亞倫就俯伏在以色列全會眾面前.

14:6　窺探地的人中,嫩的兒子約書亞和耶孚尼的兒子迦勒撕裂衣服,

14:7　對以色列全會眾說:「我們所窺探,經過之地是極美之地.

14:8　耶和華若喜悅我們,就必將我們領進那地,把地賜給我們；那地原是流奶與蜜之地.

14:9　但你們不可背叛耶和華,也不要怕那地的居民；因為他們是我們的食物,並且蔭庇他們的已經離開他們.有耶和華與我們同在,不要怕他們!」

14:10　但全會眾說:「拿石頭打死他們二人.」忽然,耶和華的榮光在會幕中向以色列眾人顯現.

讓我先講講這段經文的文字表達.

14:1:「那夜百姓都哭號」,那夜是長夜,而哭號將希望之旅變成絕望之路.

14:2:全會眾說:「巴不得我們早死在埃及地,或是死在這曠野.」一是死在那裡,一是死在這裡.他們認為,這是橫又死,豎又死的光景.

14:3:他們說:「耶和華為什麼把我們領到那地…」,一問為什麼,就是質疑,就是不信任,問題便開始.妻子問丈夫,為什麼你的衫上有條這麼長的頭髮,問題便開始.

14:3:他們說:「我們回埃及去豈不好嗎?」這是經過思考,比較,抉擇的結果.比起上來,神的路不可取,回埃及更可取.去得一塊不可能得到的應許之地是愚蠢的行徑,回到埃及做奴隸更穩當.神信不過,信自己好過.

14:4　眾人彼此說:「我們不如立一個首領回埃及去吧!」人性的悲劇莫大於此.神立一領袖,帶領他們出埃及.他們寧願立一領袖,帶領他們回埃及.剛剛出黑暗入光明,現在是出光明入黑暗.剛剛才度過復活節,現在卻甘心擁抱死亡.無怪乎他們哭了一個長夜.人性如此.我們也如此.

14:5　「摩西,亞倫就俯伏在以色列全會眾面前.」摩西,亞倫不是俯伏在神面前,故此,這應不是向神認罪,而是以謙卑的生命勸阻眾人,讓眾人也謙卑下來.或者,這是對會眾的選擇表達哀傷.會眾對神絕望,而摩西,亞倫對他們的絕望表達哀傷.

14:6　「窺探地的人中,嫩的兒子約書亞和耶孚尼的兒子迦勒撕裂衣服」,約書亞和迦勒也曾一起去窺探那地,他們有不同的看法.他們撕裂衣服,明確表達他們的憤慨.

14:7　他倆說:「我們所窺探,經過之地是極美之地.」眾人說,回埃及是好的.約書亞和迦勒卻說,我們前面要去的地,是極好的.面對同一現實,這兩個人有不同的價值判斷.現實歸現實,如何看待這現實的價值判斷是另一回事.這兩人作出不一樣的選擇.

14:8　這兩人說:「耶和華若喜悅我們,就必將我們領進那地,把地賜給我們；那地原是流奶與蜜之地.」重點不在艱難,重點在我們是否討神喜悅.

14:9　這兩人繼續說:「但你們不可背叛耶和華,也不要怕那地的居民；因為他們是我們的食物,並且蔭庇他們的已經離開他們.有耶和華與我們同在,不要怕他們!」

約書亞和迦勒鼓勵會眾「不要怕」,理由有三:一,「他們是我們的食物」,這是逆反群眾的看法.群眾認為,對方太強,他們是巨人,自己是「小薯仔」.其他探子說:「據我們看,自己就如蚱蜢一樣；據他們看,我們也是如此.」這是巨人與侏儒的對決,必敗無疑.但約書亞和迦勒卻不同意,誰是巨人,誰是蚱蜢,還未知曉.事情的成敗,不能如此草率下判斷.二,「蔭庇他們的已經離開他們」,這是一種神學的判斷.在當時,每個民族都有自己的守護神.成與敗,視乎這守護神「好不好打」.約書亞和迦勒認為,耶和華最好打,故這些神明會離開他們,保護不了他們了.這是對所信的耶和華的能力的宣認.三,「有耶和華與我們同在」,其他的神明會離棄他們的子民,但耶和華不會離棄祂的子民.這是對耶和華的應許的信任.神不會撇下我們為孤兒的.

就憑這三個理由,就憑這些信念,約書亞和迦勒鼓勵會眾「不要怕」.相反,不相信神的同在,不相信神的能力,而要回埃及,就是對神的背叛了.

14:10　「但全會眾說:『拿石頭打死他們二人.』忽然,耶和華的榮光在會幕中向以色列眾人顯現.」會眾的人生方向,已決定好了.他們選擇不信神,不相信生命的美地.他們要打死阻礙他們回埃及的人,他們做奴隸的心很堅決.在這時刻,神的榮耀顯現了,對人的反叛作出仲裁.神讓群眾講的說話成真,他們不得入迦南地,無法在人世享受那生命的福樂,他們要在曠野漂流40年,至終真的死在這曠野地.人的不信,造成自己生命的悲劇.

整個故事,對我們的人生有何提示呢?

這是一群「捉到鹿不懂脫角」的人,一班被幸福看中卻最後主動放棄幸福的人,一群在恩典中卻又踐踏恩典的人.他們原本在出埃及,入迦南的途上,卻最後演變成拒絕入迦南,而返回埃及的故事.一樁救贖事件的反敘述,反出埃及(anti-Exodus)的敘述.

因著現實的艱難,他們不信神的承諾,不信神的大能,不信神真是神.他們憑眼見的現實而作出判斷,無視上帝應許的前景.他們的生命,有一種根深柢固的奴性.他們覺得,做一個奴隸是最安穩的事.他們寧願做一隻快樂的豬,好過個一個要冒險的人.誰阻礙他們做奴隸,他們就剷除誰.他們的肉身出了埃及,但他們的靈魂卻甘心在埃及.他們經歷過救贖,但他們卻輕視救贖.他們經歷過復活節,但他們卻眷戀那永不超生的黑暗.

真無辦法,人性的墮落如此深不可測.

約書亞和迦勒試圖力挽狂瀾.他們言說希望.他們認為艱難的現實不應限定靈魂的高度.艱難的現實只是走上帝道路途中遇到的挑戰,而不是注定走上帝道路的不可能.上帝的真實不是由順風順水來證明,上帝的真實由逆境中懷著希望地前行來體會.在群眾說「不」時,約書亞和迦勒懷著信仰地說「是」.

誰對誰錯,最後,上帝顯現,作出裁判.不信的人,按其不信,領受無福樂的人生.信靠神的人,按其信仰,在前行中,終於抵達神應許之地.歷史的黑暗力量,強大得不得了.在這考驗人性的曠野,原先應是點燈的民族,其中放棄光明的人竟佔多數.唯有少數的人,在曠野裡,陶造出「歸屬神的生命質素」,他們無懼艱難,堅持信仰,堅持盼望,堅持在人間傳遞光明.你又如何呢?前行而迎向巨人,抑或回頭而甘心做一隻蚱蜢?

願榮耀歸上主,願人間享和平.



\newpage

\section{超越多馬的信心}
\label{sec:596}
\textbf{高國雄牧師 - 超越多馬的信心}
\newline
\newline
\hyperref[sec:595]{< < < PREV SERMON < < <}
~
\hyperref[sec:toc]{[返目錄]}
~
\hyperref[sec:597]{> > > NEXT SERMON > > >}


講題:超越多馬的信心

講員:高國雄牧師

日期:2022年4月24日

前言:

近日被同學問到,對現時流行樂壇中最火紅的12「鏡仔」有幾多認識…當我不單能數出12人名字,甚至連12人各自的粉絲團的名字後；我反問他,作為信徒,對耶穌的12門徒有多少認識呢?對藝人/明星/偶像的一切我們會努力認識,對耶穌與及祂身邊的人呢?

十二位門徒各自精彩:使徒:ἀπόστολος,被派遣出去的人

福音書中的馬太10:1-4,馬可3:13-19,路加6:12-16也有12人的名單:頭一個叫西門(又稱彼得/磯法),還有他弟弟安得烈,西庇太的兒子雅各和雅各的弟弟約翰(半尼其(雷的兒子)),腓力和巴多羅買/拿但業,多馬和稅吏馬太/利未,亞勒腓的兒子雅各,和達太(雅各的兒子猶大),激進黨的西門,還有出賣耶穌的加略人猶大(後被馬提亞取代) (太10:1-4)

約翰福音沒有列出使徒的名單,而用”十二” δώδεκα 我們對12使徒認識不深,可能因為福音書記述只集中在幾位使徒身上.“出埸”比較多:約翰,彼得,雅各；負面:賣主的猶大,多馬,激進黨/奮銳黨的西門(?)

奮銳黨(未必是政治組織):第四大派(法利賽,撒都該,愛色尼)

對多馬的認識:

今天經文的主角多馬,大家對他有什麼認識呢?除了今天的經課之外,大家記得福音書中其他對多馬的記載嗎?稱為低土馬(孖仔)的多馬 太,可,路:只在12使徒名單中出現,約翰福音中有3段與多馬相關的記錄:約11:1-16拉撒路死了,耶穌對門徒說他睡了,要去叫醒他:

12門徒就說:「主啊,他若睡了,就會好的.」

13耶穌說這話是指拉撒路死了,他們卻以為他是指通常的睡眠.

14於是耶穌就明白地告訴他們:「拉撒路死了.

15為了你們的緣故,我不在那裡反而歡喜,為要使你們信.現在我們到他那裡去吧.」

16於是那稱為低土馬的多馬對其他的門徒說:「我們也去和他同死吧!」

怨言?悔氣說話?少信(拉撒路已死,沒有冒險的價值)

信心建基在耶穌的邀請,就算有危險,也願前行.

約14:1-6

1「你們心裡不要憂愁；你們信上帝,也當信我.

2在我父的家裡有許多住處；若是沒有,我就早已告訴你們了.我去原是為你們預備地去方.

3我若去為你們預備了地方,就必再來接你們到我那裡去,我在哪裡,叫你們也在哪裡.

4我往哪裡去,你們知道那條路.」

5多馬對他說:「主啊,我們不知道你去哪裡,怎麼能知道那條路呢?」

6耶穌對他說:「我就是道路,真理,生命；若不藉著我,沒有人能到父那裡去…

直接的問題:不知道耶穌的去向,又如何知道那條路呢?

這是一個非常好,亦很重要的問題

多馬應該是一個真誠,坦率,理智,不盲從,重思考的人

我們不要忘記就是因為多馬提出這問題,耶穌才回答出對我們很重要的金句:

「我就是道路,真理,生命；若不藉著我,沒有人能到父那裡去」

但是在歷史上,在大部份人心中,他是懷疑者,不信者

被冠以懷疑者多馬的名號

義大利著名畫家卡拉瓦喬世界級名作《懷疑者多馬》* 聖經的標題(編輯加上的) 耶穌向疑惑的多馬顯現,多馬不信

耶穌和多馬,向多馬顯現

彷彿多少對多馬的表現感負面

回到經文 約20:19-31

19那日(就是七日的第一日)晚上,門徒因怕猶太人,所在的地方門都關了.耶穌來,站在當中,對他們說:「願你們平安!」

20說了這話,他把手和肋旁給他們看.門徒一看見主就喜樂了.

21於是耶穌又對他們說:「願你們平安!父怎樣差遣了我,我也照樣差遣你們.」

22說了這話,他向他們吹一口氣,說:「領受聖靈吧!

23你們赦免誰的罪,誰的罪就得赦免；你們不赦免誰的罪,誰的罪就不得赦免.」

24那十二使徒中,有個叫低土馬的多馬,耶穌來的時候,他沒有和他們在一起.

25其他的門徒就對他說:「我們已經看見主了.」多馬卻對他們說:「除非我看見他手上的釘痕,用我的指頭探入那釘痕,用我的手探入他的肋旁,我絕不信.」

26過了八日,門徒又在屋裡,多馬也和他們在一起.門都關了,耶穌來,站在當中,說:「願你們平安!」

27然後他對多馬說:「把你的指頭伸到這裡來,看看我的手；把你的手伸過來,探入我的肋旁.不要疑惑,總要信!」

28多馬回答,對他說:「我的主!我的上帝!」

29耶穌對他說:「你因為看見了我才信嗎?那沒有看見卻信的有福了.」

30耶穌在他門徒面前另外行了許多神蹟,沒有記錄在這書上.

31但記載這些事是要使你們信耶穌是基督,是上帝的兒子,並且使你們信他,好因著他的名得生命.

多馬一句:「我絕不信!」,成為我們對他的懷疑者的認知.經文提到,門關上,耶穌顯現,帶著傷口,門徒見到,多馬卻沒有在一起…試代入多馬心情,一個真誠,坦率,理智,不盲從,重思考的人,當然一切要理性分析(這和我們在大學內所學習的一樣).所以「我絕不信!」有一個前設,就是「除非我看見他手上的釘痕,用我的指頭探入那釘痕,用我的手探入他的肋旁…」.

我欣賞多馬,理性分析,親身體會都很重要,我們不要只成為宗教狂熱者,更不可盲從附和；更重要的是耶穌不單顯現,更向門徒吹氣授權「領受聖靈吧!你們赦免誰的罪,誰的罪就得赦免；你們不赦免誰的罪,誰的罪就不得赦免.」作為十二使徒中一份子的多馬,一方面對這極要的事件要查證,更不能接受這錯失.過了八天,這八天到底多馬是怎樣過的呢?心情又如何呢?同樣的情況,門關上,這次多馬與門徒們在一起,耶穌顯現…弟兄姊妹請謹記,耶穌很愛多馬(甚至可說特別愛),特別再次為他顯現.不單特別為他再次顯現,更主動回應多馬早前在門徒中的提出的查證要求

V27耶穌對多馬說:「把你的指頭伸到這裡來,看看我的手；把你的手伸過來,探入我的肋旁.不要疑惑,總要信!」

總要信vs我絕不信

耶穌的愛何等的大,請留意多馬的回應

V28多馬回答,對他說:「我的主!我的上帝!」

多馬沒有馬上伸手去探(到最後都沒有),因為不再需要.他作出信仰上最重要的宣告:承認耶穌基督是主,是上帝!要一個著重理性分析,著重思考的人在信仰上的降服絕對不是容易的事,相信這次多馬對主的降服成為往後他忠心事奉主重要的記號,教會傳統多馬後來就去到印度傳道,並被矛刺死而殉道.多馬因為基督的顯然而完全降服,沒有過多的懷疑(沒有伸手作確認),耶穌基督接著的回應:

V29耶穌對他說:「你因為看見了我才信嗎?那沒有看見卻信的有福了.」

請留心,耶穌所指的”那沒有看見卻信的人”是誰呢?門徒親眼看見了復活的主的顯現,多馬在第一次顯現時不在,耶穌基督在一星期後特別為他再次顯現,上帝特別賜福多馬,但他們全都是”看見了”的人.”那沒有看見卻信的人”所指的就是往後的信徒,亦即是指今天的我們.一般人的理解,想法,好可能是如果能親眼見證,才是最實在的,信仰也一樣,不少人未能全身投入信仰,其原因或許是因為未來很實在地經歷,如果耶穌基督今天在我面對顯現,我一定完全信服.

耶穌的教導卻不一樣,主所要求我們的是信心.信是什麼呢?

希伯來書11:1 信就是所望之事的實底,是未見之事的確據.

信就是對所盼望之事有把握,對未見之事有確據.(和修版).理性尋真是非常重要的,因為對任何知識的追求,就算與信仰相關也不可盲目,迷信,但「認識上帝」,這認識是指與上帝溝通,交往,建立關係,所能靠的,就只有信:

希伯來書11:6 人非有信,就不能得上帝的喜悅；因為到上帝面前來的人必須信有上帝,且信他賞賜那尋求他的人.

這關係是蒙福的關係,是得生命的秘訣,正如經課提到:約翰福音20:30-31耶穌在他門徒面前另外行了許多神蹟,沒有記錄在這書上.但記載這些事是要使你們信耶穌是基督,是上帝的兒子,並且使你們信他,好因著他的名得生命.基督昔日在復活後顯現時留下傷口的痕跡,為了讓我們認得祂,也是為了記念我們在罪可被赦免,這記號在我們與基督再見的日子仍然會見到.

今天我們雖然無法用肉眼看得見祂,但復活的主仍然在我們心裡活著,這記號仍然在我們身上顯為榮耀.阿們.



\newpage

\section{生亦何歡死亦何苦}
\label{sec:597}
\textbf{林豪恩先生 - 生亦何歡,死亦何苦?}
\newline
\newline
\hyperref[sec:596]{< < < PREV SERMON < < <}
~
\hyperref[sec:toc]{[返目錄]}
~
\hyperref[sec:598]{> > > NEXT SERMON > > >}


講題:生亦何歡,死亦何苦?

講員:林豪恩先生

日期:2022年5月1日

一.引言

還記得《倚天屠龍記》嗎?如果要選最轟烈的一幕,你會選哪一個片段呢?有些人選六大派圍攻光明頂那一幕.當時,少林,武當,峨眉,崑崙,崆峒,華山六派高手齊集光明頂,明教中人與他們戰至最後一刻,生命危在旦夕之際,突然停下一同打坐,誦念禱文.他們當時誦念的禱文是這樣的:

「焚我殘軀,熊熊聖火.生亦何歡,死亦何苦?為善除惡,唯光明故.喜樂悲愁,皆歸塵土.憐我世人,憂患實多!憐我世人,憂患實多!」

生存,從來都是人們的本能；死亡,從來都是人們的威嚇.何以有人慨嘆:「生亦何歡,死亦何苦」呢?寫下這禱文的人,很可能深刻思考過生存的意義,然後對於生死建立了某種觀念和態度.有人認為,人類是唯一會思考生存意義的動物.你問過自己這個問題嗎?讀大學的時候的某個黃昏,我與我的「蛇」(非正式宿生)坐在宿舍房間的地板上,關了燈,在暗光之下,談起一個奇怪的話題:「為甚麼我們仍然生存下去?」我已經忘記了自己講過甚麼,但仍記得他說:「因為我的媽媽仍然在世,我要照顧她.」你呢?你生存下去的原因又是甚麼呢?今天經課的四段經文,都涉及生與死,讓我們再讀經文,思考對我們的啟發.

二.經文及思考

1.詩篇30篇:我被害流血,下到坑中,有甚麼益處呢?

詩篇30篇的作者如絕大部分人一樣,面對死亡威嚇的時候都不想死,而盡力求生,例如:

30:2耶和華－我的上帝啊,我曾呼求你,你醫治了我.

30:3耶和華啊,你曾把我的靈魂從陰間救上來,使我存活,不至於下坑.

30:8耶和華啊,我曾求告你；我向耶和華懇求,說:

30:9我被害流血,下到坑中,有甚麼益處呢?

為甚麼呢?

30:11你已將我的哀哭變為跳舞,將我的麻衣脫去,給我披上喜樂,

詩人把死亡與「哀哭」和「麻衣」相配,把生存和「跳舞」和「喜樂」相連,這也是人間的共識,生,總是喜慶的,死,總是哀傷的.

絕大部分人都不想死,然而,如果要為生存下去找一個理由,那會是甚麼呢?詩人如此向耶和華祈求說:

30:9我被害流血,下到坑中,有甚麼益處呢?塵土豈能稱讚你,傳說你的誠實嗎?

詩人的答案是「稱讚你」和「傳說你的誠實」.如果是你,你的答案是甚麼呢?

沒有答案一點也不奇怪,一般來說人們是不會思考生存下去的理由的,因為生存是順理成章的,詩人順境時也是一樣:

30:6至於我,我凡事平順,便說:我永不動搖.

30:7耶和華啊,你曾施恩,叫我的江山穩固；

在甚麼狀況之下,人們才會思考生存的意義呢?就如詩人一樣,在「被害流血,下到坑中」的處境,即是感到生命危在旦夕之時.正所謂,「誰會珍惜,當你還擁有,將要逝去總想挽留」.在生死邊緣之時,人們就會想:「如果今次死唔去,我就會……」.不要以為如此激烈和震撼的時刻是罕有的,可能有不少人都經驗過.有一位大家尊敬的弟兄,意想不到地發現腦腫瘤.第一次手術算是成功,但也知道腫瘤未能完全清除,有復發機會.雖然病患未能根治,威脅沒有完全解除,但他也為自己的生命能夠存留而感恩.在手術前後,相信他深入地思考過人生,坦誠地與上主溝通.手術後,他除了繼續在工作上造福人群之外,也調整自己的工作節奏,投入事奉,培育了一代又一代的兒童.相信不少受教於他的兒童(有些現在已經長大成人)及家長,都感謝他和懷念他.

可能你活了二十年,或者三十年；五十年,或者七十年……繼續活下去,所為何事呢?當你為自己尋找答案的時候,詩人提出的「稱謝上主」和「傳頌上主的誠實」會不會成為你的參考呢?

2.約翰福音21章:耶穌說這話是指著彼得要怎樣死,榮耀上帝

有生必有死.古人說:「人生七十古來稀」,現在壽命延長,人們說:「長命百二歲」,即是說,無論活多少年都總會離開.死,有人認為有輕重之分.司馬遷《報任少卿書》有云:「人固有一死,或重於泰山,或輕於鴻毛,用之所趨異也.」福音經課約翰福音21章以耶穌與門徒的溫馨故事開始,記述復活的耶穌親臨打魚的門徒當中,並且預備烤魚作早餐給門徒吃,但卻以談論使徒彼得的死為結束.耶穌三次問彼得「你愛我嗎?」彼得三次回覆「我愛你」.如果抽空來看,以為是浪漫愛情故事.放回故事脈絡,原來又是生死攸關的事情.當彼得三次確認對耶穌的愛之後,下文是這樣的:

21:18我實實在在地告訴你,你年少的時候,自己束上帶子,隨意往來；但年老的時候,你要伸出手來,別人要把你束上,帶你到不願意去的地方.」

21:19(耶穌說這話是指著彼得要怎樣死,榮耀上帝.)說了這話,就對他說:「你跟從我吧!」

有一個關於彼得殉道的傳說是這樣的:那時,尼祿皇帝展開了羅馬帝國第一次的大迫害,起初,許多在彼得身邊的基督徒都勸他趕快走避,他一開始不願意離開這些羊群,但後來沒有辦法了,只好一人逃離,當他到了那路上時,突然看到不遠處有個人走來,他認出就是主耶穌基督,他問道:“主啊!你要到哪裡去?”耶穌對他說:“彼得,你離棄了我的羊,現在我要到羅馬再釘一次十字架!”彼得頓時捶胸大哭,立刻又轉回羅馬,不久被捕,他要求被倒釘十字架,因為他自認不配與主同樣的死法.

詩人啟發我們問,生存下去所為何事呢?彼得的故事邀請我們思考,若然非死不可,死,又可以有甚麼價值呢?對詩篇30篇的作者來說,生存是呈現上主的恩典；約翰福音的作者指出,彼得的死是榮耀上帝.正如保羅在羅馬書14:8說:「我們若活著,是為主而活；若死了,是為主而死.所以,我們或活或死總是主的人.」

3.啟示錄5章:曾被殺的羔羊是配得權柄,豐富,智慧,能力,尊貴,榮耀,頌讚的.既然有生必有死,因此很多人認為緊要的不是不死,而是「好死」.善寧會於2004年對公眾人士進行了問卷調查,了解他們對「好死」的想法.結果發現,最重要的因素:「無肉體痛苦」(如:「死前唔駛病痛長時間折磨」)；其次是「留者善別」,意指對於遺留下來親人的需要是否能妥善照顧(如:知道自己死唔會拖累家人)；其三是「去者善終」,乃有關臨終病人的個人心理需要,是否安詳地離逝等(如:死前仍然能夠得到其他人尊重).

如果離開的時候有齊上述三項,算是「好死」了.如果「天氣不似預期」,那又怎樣呢?

彼得被拘捕,被判刑,被倒針十字架而死,在當時的社會看來,一定是羞辱；在身體來說,是痛苦.啟示錄第5章也提到一位「被殺過的」羔羊:5:6我又看見寶座與四活物,並長老之中有羔羊站立,像是被殺過的,有七角七眼,就是上帝的七靈,奉差遣往普天下去的.

這羔羊如何被殺的呢?

5:9他們唱新歌,說:你配拿書卷,配揭開七印；因為你曾被殺,用自己的血從各族,各方,各民,各國中買了人來,叫他們歸於上帝.

啟示錄中這位被殺的羔羊就是耶穌基督,如果用「好死」的標準來看,耶穌基督應該不算「好死」.他被判有罪,列在囚犯之中,被釘十字架而死.不過,啟示錄卻指出:

5:11我又看見且聽見,寶座與活物並長老的周圍有許多天使的聲音；他們的數目有千千萬萬,

5:12大聲說:曾被殺的羔羊是配得權柄,豐富,智慧,能力,尊貴,榮耀,頌讚的.

5:13我又聽見在天上,地上,地底下,滄海裡,和天地間一切所有被造之物,都說:但願頌讚,尊貴,榮耀,權勢都歸給坐寶座的和羔羊,直到永永遠遠!

原來,在某種歷史空間之中「不得好死」的,在永恆之中卻「配得權柄,豐富,智慧,能力,尊貴,榮耀,頌讚的」,因為:

5:9他們唱新歌,說:你配拿書卷,配揭開七印；因為你曾被殺,用自己的血從各族,各方,各民,各國中買了人來,叫他們歸於上帝.

在人世界被否定的,可以在永恆中被肯定.被殺的羔羊如是,使徒彼得如是,或許,敬拜被殺的羔羊,承繼使徒傳流下來的信仰的人士,也應該有如此的視野.在天地之間,有超越生與死的；一切價值與意義,不在變動的歷史中,而在永恆之主的心中.

三.總結:使徒行傳9章:為誰奔波為何活?

我們都存在於生與死之間,營營役役正是我們的生活寫照.我們每日奔波忙碌,勞苦不息,有些只是為求生存,有些則奮力為生存增加一些甚麼.使徒保羅稱為掃羅時期的前半生也是如此奔波:

9:1掃羅仍然向主的門徒口吐威嚇凶煞的話,去見大祭司,

9:2求文書給大馬士革的各會堂,若是找著信奉這道的人,無論男女,都准他捆綁帶到耶路撒冷.

如此奔波勞碌所為何事呢?為生命增添了甚麼價值呢?他後來在腓立比書有這樣的反省:

3:6就熱心說,我是逼迫教會的；就律法上的義說,我是無可指摘的.

3:7只是我先前以為與我有益的,我現在因基督都當作有損的.

保羅這種反省,並不是來自他自己的,而是從上而來的光照:

9:3掃羅行路,將到大馬士革,忽然從天上發光,四面照著他；

9:4他就仆倒在地,聽見有聲音對他說:「掃羅!掃羅!你為甚麼逼迫我?」

9:5他說:「主啊!你是誰?」主說:「我就是你所逼迫的耶穌.

9:6起來!進城去,你所當做的事,必有人告訴你.」

生與死是自古以來與人類這種動物不可分割的話題,其實,人們對於生與死可以控制的並不太多.我們出生來到世上,並不需要得到我們的同意；我們死亡離去之時,我們不同意也不能改變甚麼.屬於我們的,是在生與死之間的日子.有些人會問:如何過這些日子才有意義?怎樣的人生才有價值?或者,答案未必盡在人間,正如保羅的人生一樣,唯有在從上而來的光照之下,才能分辨甚麼是有意義,甚麼是有價值.然而,尋找人生意義的人絕不孤單,因為人世間有不少這樣的同路人.

有一個關於美國的第34任總統艾森豪威爾的傳說是這樣的:

他曾經和朋友一起玩打牌遊戲時,就時常抱怨自己抓的牌很爛,有時因為自己的牌爛而大發脾氣.有次他母親就教訓他道:「既然要打牌,你就必須用手中的牌打下去,不管牌是好是壞,好運氣是不可能都讓你碰上的!」

疫情之下離世的人不少,其中一位引起公眾關注,他叫羅偉祥.童年時,受疾病所影響,羅先生除了手腳不靈活外,更不能坐立.他每天只可以困在家裡,用僅有的力量,把身體貼在地上爬出門外看天.直至12歲時,羅先生才有機會在香港紅十字會甘迺迪中心讀書,經同儕間的互相鼓勵和支持,羅先生自始建立積極堅毅的人生態度,及後他更參與志願人道工作,為傷殘人士爭取權益,又協助他們改善生活.路向四肢傷殘人士協會主席嚴楚碧在羅偉祥的安息禮拜表示,祥仔臨終前表示「仲有好多ye 未做」,仍希望可為殘疾者做得更多.

艾森豪威爾的媽媽的話還未說完,下文如下:

「人生就和這打牌一樣,發牌就是上帝,不管你名下的牌是好是壞,你都必須拿著,你都必須面對,你能做的,就是讓浮躁的心情平靜下來,然後認真對待,把自己的牌打好,力爭達到最好的效果,這樣打牌,這樣對待人生才會有意義!」



\newpage

\section{彌賽亞牧人}
\label{sec:598}
\textbf{李均熊牧師 - 彌賽亞牧人}
\newline
\newline
\hyperref[sec:597]{< < < PREV SERMON < < <}
~
\hyperref[sec:toc]{[返目錄]}
~
\hyperref[sec:599]{> > > NEXT SERMON > > >}


講題:彌賽亞牧人

講員:李均熊牧師

日期:2022年5月8日

1. 引言:甚麼是彌賽亞牧人——牧人形象的再思

今日福音經課中,猶太人問耶穌,如果他是基督—彌賽亞—就明明地告訴他們.但若我們從約翰福音第五章一直看到這裡,的而且確耶穌的行動而見證他就是彌賽亞／基督,而有些猶太人亦因此想置耶穌於死地.所以耶穌的回答十分反映當時處境:我已經告訴你們,你們不信.你要叫耶穌再說甚麼呢?而這些猶太人的不信,證明他們不是耶穌的羊.因為「我的羊聽我的聲音,我也認識他們,他們也跟著我.」約翰福音中的耶穌在這裡有論到賜下永生給跟隨他的人,大家可以翻查之前我曾經講過關於永生的理解,這裡不贅.但與今日我們主題更息息相關的,是耶穌論到是父把羊賜給他,以至誰也不能從他父手裡把他們奪去.這句「誰也不能把他們奪去」其實總共出現兩次,顯見耶穌要強調這點.

究竟誰要將羊奪去呢?原來在古代牧養,並不如現代一般容易.古代的牧羊的地方,其實就是野外之地,遇見兇猛野獸是十分普遍的.童話故事說到牧童見到豺狼,其實已經小兒科.舊約阿摩司書一段話正好讓我們瞭解牧養的凶險.阿摩司3:12節論到耶和華如此說:「牧人怎樣從獅子口中搶回兩條腿或耳朵的一小片,住撒瑪利亞的以色列人得救也是如此,不過搶回床的一角和床榻的靠枕而已.」牧人所面對的不是豺狼,而是獅子!在撒母耳記上17:34-37中,當大衛要去對付巨人歌利亞時,向掃羅介紹自己,就是一個可以對付熊和獅子的牧羊人!

這就引入我今日想同大家分享的,關於彌賽亞牧人的形象問題.我們心目中,究竟牧人的形象是甚麼呢?當然我們城巿人大多都沒有牧羊經驗,聖經常說到牧人和羊,我們都只靠想像力推斷經文的意思.我們或許覺得牧羊人只是底下階層,只是帶著羊兒找水草,閒來就吹o下簫,曬太陽,是十分沉悶的工作.加上我們現今都稱呼教會聖職人員叫牧師,又十分想當然的認為牧師就是文靜,有愛心,溫柔等等.這其實都不是聖經牧人想表達的形象!就如我剛提到,牧人並不是弱不禁風,閒時無聊的人.他們要能搏擊大型的野獸,才能守護窮山惡水所可能遇到的凶險.在古代的米索匹亞,埃及,希臘等文化,牧人是用來表示君王的.聖經多處提到牧人,都在表達他們應該有的強而有力的領導力,正如詩篇80篇提到:領約瑟如領羊群的以色列牧者啊,求你側耳而聽!在基路伯之上坐寶座的啊,求你發出光來!在以法蓮,便雅憫,嗎哪西面前,求你施展你的大能,拯救我們.

從這個強悍的領導者的形象出發,我們或許能更加了解今日經課中選讀的詩篇23篇.

2. 牧養的危機

首先,讓我們先看看詩篇23篇第5和6節(按原文特色和字詞順序再譯):

你在我面前設置盛筵,

~~~~~~~~對正我眾仇敵；

你用油膏了我的頭,

~~~~~~~~我的杯滿溢

此一詩句我想大家留意一下,傳統譯為「在我敵人面前,你為我擺設筵席」,其實大大地削弱了這一詩句的威力.原句動詞設置放在最頭,而敵人放在最尾.又再想像一下,你仍然是一隻羊,貪食易驚.牧人在你面前擺設了一些甘甜的乾草,分量多過你所要的,當你一路吃一路抬頭,忽然見到正對著你的是隻野狼!但為食的你真的好少理,因為有強悍的牧人「頂住」嘛.牧人,詩中的主語「你」,正對敵人,在危險當下,就是這麼親近,又是這麼可靠,當我驚覺危險臨近,那時以經不成問題了,因為你已經預備盛筵【激鬼死隻狼!!】.

其實詩篇23篇整首詩並不是在風和日麗的日子,牧人與羊群閒逸地在青草地,溪水旁靜靜的打發日子.這兩節突兀的出現,曾經困擾很多解經家.若果詩篇23篇只是說說羊兒安舒的日子,偶然經過死蔭幽谷,那這裡提到仇敵環伺就與全詩顯得格格不入.但有猶太解經家指出,全詩其實是詩人在危急中的呼求,並非安逸中的祈禱.

3. 危急中的引導

大家不況再看一看第二和第三節的翻譯:

在茂密草原,祂讓我躺下【回味】,

~~~~~~~~下於靜水之間,祂帶領我

我的生命由祂挽回,

~~~~~~~~祂引導我在正確路徑上,皆為祂的名.

這節的翻譯有幾點要大家留意.與平時希伯來字詞排列並不一樣,開始三句所有動詞都是在句尾,故詩句所強調的是句中的名詞而非動作；而那些地方都是與羊的飲食有關.但同一時間,所有動詞都指向牧羊人的帶領.連第三節開始的一句,其實都和牧羊人的工作有關.和合本譯作:「他使我的靈魂甦醒」,其實不太恰當.若這一句仍然是講到羊和牧羊人的關係,那麼羊又可來「靈魂」去給牧羊人「甦醒」?更直接的翻譯,是牧羊人「挽回」羊的「生命」.牧人很多時要找回走失的羊,將牠帶回羊群中.所以牧人的工作就是帶領,挽回,陪伴,緊追,又領羊往正確的路徑.

但更值得留意的,是詩句動詞的時態.理論上,幾個動詞都是指仍未發生的狀況,而非回憶曾發生的事情或當下發生中的景象.若詩人—詩中的「我」—正在急難中,他要表達並不是他曾享過的風光,而是確信強捍的牧人能引領他走到靜水邊享受水草,甚至現在迷失的自己,他也想信上主作為牧者必然挽回他的「生命」,再引導他行在正途上.

4. 危急中的同在

我們繼續看看第四節:

縱然我走在死蔭的谷中,

我不懼凶險；

皆因你陪伴我--

~~~~~~~~你的杖與你的竿,

~~~~~~~~~~~~~~~~它們,它們勉勵我

你比較二,三節,和這一節,你會否發現一個很突然的轉變?

不錯,就是原本詩人稱牧人為他,在這裡突然改口叫牧人為你.

為什麼會有這樣的轉變?

之前是確信「他」會引導自己脫離危險,想像脫險後的種種可能；但在死蔭幽谷中,自己才驚覺我需要的是你,不是那些草地溪水,而是牧人本身.

當中詩中有一個轉折地方,就是全句的焦點只有「我」.唯有「我」意識自己身處凶險,才會明白皆因你的同在陪伴我才會不驚懼.

我們可時從經歷上主由相對於我的「他」而變成我的「你」??

試想像,為什麼羊會落入幽谷呢?這有兩個可能性,一是自己與羊群走散了,一是牧人知道幽谷之後有更好的草原,於是帶領羊經過死蔭的幽谷.兩個可能性,牧羊人的枚和竿都成為羊的安慰,勉勵.使羊可以走下去,離開幽谷.

什麼是竿和杖?它們怎樣安慰,勉勵呢?

上主的竿,是用以擊打的短棍,正是強捍的牧人對付熊人獅子的武器.但自己走散了,牧人也會責打,領回.

而牧人的杖,則是長長的,可以高舉過頭的標記.牧人領路,你遠遠仍見到牧人的蹤影.好多時候我們看圖劃,杖的前端是呈彎鉤狀的,也是領回走散羊群的工具.

最後一節,是當在急難中的詩人回望一生時,發現追趕他的原來是化妝的祝福.

原來是良善與慈愛緊追著我,

~~~~~~~~在我一生的日子；

我且要回到上主殿中,

~~~~~~~~長久的日子.

緊追,原意是追趕.又想像你是一隻羊.容易受驚的你,突然覺得好像有些東西追著你,於是你發足狂奔.但最後,「原來是良善與慈愛」追你呀!是牧人在追你回去羊群中呢?在你生命的每一天中,你是否知道好牧人其實一路緊追著你呢?你還「行」到那裡呢?

做一隻羊,是要經常回到主的身邊的

5. 主乃吾牧我無缺

詩人用上一個「變重負面」(double negative)的字句,正是要帶出「一無所缺」並不是靜態的「什麼也有」.不是因為我們什麼也有,銀行的儲備夠我們食「過世」,所以變得「一無所缺」.從整篇詩篇的意像與佈局,這個「一無所缺」是一種動態的狀況.甚麼是「動態的一無所缺」呢?是剛剛發現米缸沒有米,就收到公司寄來的出糧支票；明天要考試,今日就收到老師發放的必考題示；這一刻還擔心考不上大學,下一刻就收到入學的通知書；獨居的你剛好扭到腳踝不良於行,朋友就「突然」致電慰問候,正好解決你即是的困境；病到五顏六色,室友剛好煲了老火湯問你要不要?；中國人所講的「及時雨」,「救急扶危」,「挽狂瀾於既倒,扶大廈之將傾」,「絕處逢生」；保羅所說的:「我們四面受敵,卻不被困住；心裡作難,卻不至失望；遭逼迫,卻不被乏棄；打倒了,卻不至死亡.」

不是什麼也有,你會覺得「一無所缺」；是上帝及時,剛好的補足,使你覺得「一無所缺」.

在不足的當下感受充足的供應,這才使人驚覺:我真的一無所缺.

為什麼上帝不一次過供應我們一生所需要的嗎哪?當你以為什麼也有時,你還會意識自己只是一隻無用的羊嗎?無時無刻需要引導和帶領嗎?

以上我所舉的例子都或多或小是我聽過或我自己親身體會的類似經歷；但我們永遠都只有回頭去看,才知道:原來一直追著我們的困難,真是上帝化了裝的祝福；原來使我們受苦的杖和竿,才是真正的安慰,勉勵.

6. 結語:我是誰的羊?

不過,這篇詩篇提醒我們,最重要的,或者說對我們最難做到的,不是牧人對我們有多好,而是你願意做一隻羊嗎?約翰福音中耶穌回答猶太人,說他們不是他的羊,所以不信.但本乎羊的習性,可能是猶太人自以為知道怎樣分辨彌賽亞,又或拒絕跟隨耶穌,不願走入基督彌賽亞牧人的羊圈而已.

猶太人圍著他,說:「你叫我們猶疑不定到幾時呢?你若是基督,就明明地告訴我們.」耶穌回答說:「我已經告訴你們,你們不信.我奉我父之名所行的事可以為我作見證；我的羊聽我的聲音,我也認識他們,他們也跟著我.我又賜給他們永生；他們永不滅亡,誰也不能從我手裡把他們奪去.我父把羊賜給我,他比萬有都大,誰也不能從我父手裡把他們奪去.」

做一隻羊代表什麼呢?代表你放下領導權,甘心被帶領.

你願意做一隻看似愚蠢的羊,放下自我,讓牧人引領你路嗎?這樣,會否你成為父所賜給彌賽亞牧人的羊,將會永不失落?

在蘇格蘭的高地,羊很容易在群石之間迷失,而走到一些地方牠們不能夠走出來.那些山間的野草特別香甜,有時羊可以一躍十二尺去吃那些牠們喜歡的草堆.但之後牠們卻不能夠跳回來.牧人會聽到牠們受困的哀鳴,不過牧人會不理會等那隻羊,等牠受困多數日,直致牠吃盡那處的草地.牧人要等到那隻羊餓到發軟蹄,才會用繩套將那隻羊套回來.有人問那些蘇格蘭牧人,為什麼不立即救回那隻迷失的羊?原來牠們太易受驚,以致愚蠢到牧人施救時躍下山崖跌死.

為什麼上帝不立時施救?在不合適的時機施救我們可能就會拒絕牧人的繩套以自尋死路.

但又是什麼令我們可以在死蔭幽谷裡堅持下去?是什麼是我們在眾敵環繞的處景中乃能心存盼望?

耶穌是好牧人.

但我們願意做相信祂,跟從祂的羊嗎?



\newpage

\section{生路如此美好何解選擇滅亡}
\label{sec:599}
\textbf{鄧瑞強博士 - 生路如此美好,何解選擇滅亡?}
\newline
\newline
\hyperref[sec:598]{< < < PREV SERMON < < <}
~
\hyperref[sec:toc]{[返目錄]}
~
\hyperref[sec:600]{> > > NEXT SERMON > > >}


講題:生路如此美好,何解選擇滅亡?

經文:申命記10:12-22及30:5-20

講員:鄧瑞強博士

日期:2022年5月15日

今日,我要總結摩西五經的講道了.

人犯罪,令世界失去秩序,泯滅人性,也遺忘神聖.神定意救贖世人,但祂沒有選擇一種直接了當的方式,祂選擇了曲折迂迴的方式,祂選擇了與人在歷史中同行,讓人在歷史中將祂的光明傳播開來.神這種選擇最終使祂行在人間,流出承擔人類罪惡的鮮血,在歷史中留下救贖的記號.

神選擇與人在歷史中同行,神人合作地在黑暗的大地上傳送光明,這是漫長的過程.神人合作,有點似玩「二人三足」的遊戲.神就著人,步伐緩慢了.人要跟著神的節奏,也需要努力的學習.在這神人互動中,神的光明,透過人,在大地上蔓延開來.

神選擇了亞伯拉罕,亞伯拉罕也選擇了神.神要求亞伯拉罕大膽地信任神,大膽地為神獻上一切.亞伯拉罕信任地回應神,在他的信仰行動中,他看到神的同行,看到神的救贖,看到歷史絕境中從神而來的希望.在漆黑的夜空下,神叫亞伯拉罕抬頭看看那無盡的星星,這是他手中的火把所點燃的無盡火光,將散佈在大地的每個角落上.在慘淡的人生裡,神叫亞伯拉罕想像有一天萬國都要因他得福,亞伯拉罕想像到,在暗淡的歷史長河的盡頭那裡,大地瀰漫著天國的福樂.

亞伯拉罕在每天都是絕境的光景下,大無畏地活出對神的信靠,對前景的希望.他對神的信仰,塑造了一個家族.這個家族,發展出一個民族.這個民族,以信仰神,敬畏神作為其民族精神.這精神,成為火把,在黑暗的人間,燃點光明.這是創世記的主題.

神要透過一個民族救贖人間,必先讓這民族深刻地經驗救贖.神要透過這個民族傳遞光明,必需讓這民族明白光明的生活大道.在歷史的曲折迂迴中,這民族成了埃及的奴隸.神選擇摩西,帶領他們離開埃及.如此,他們經歷了一段刻骨銘心的歷史,經歷了改變他們命運的救贖事件.在這救贖事件的過程中,神頒佈誡命與律法,讓這群人明白如何在生活中活出光明.如此,世上不單有一以信仰神為其內在精神的民族,如今,也能明確知道這信仰的具體內容.這是出埃及記及利未記的主題.

要讓這民族實在地連繫於歷史,要讓光明的生活方式在塵世中具體可見,神要賜一塊土地給他們,如此,這民族才不是大地上的遊魂,而是具體可見的神的子民.這就是民數記和申命記的主題.信仰的光明,若要點亮漆黑的大地,就必須具體可見.但現今的文化,是將信仰劃定為私人領域的東西.很多基督徒的信仰,只是一種內心的信念,對黑暗的大地,並不提供半點光明.這是一種沒有土地的信仰,香港人叫這做「離地」.神賜土地給祂的子民,因信仰根植大地.貼近大地,才能將光明傳遞.暗淡的歷史,才有光明的前景.

摩西五經的三個主題:亞伯拉罕有無數的後裔可組成一民族,這民族以遵守律法及活出光明去凸顯他們與神的關係,他們有一活在其上的土地好具體地祝福世界,已完全呈現出來了.來到申命記的經文,神所應許的土地,已遙遙在望了.

摩西,帶領著這一大群人,離開埃及,經過曠野,到達神應許之地的邊界.他老了,他知道自己去不到這應許之地了.他望向前面,他望到神所應許之地,他望到這民族可能有的前景.再望遠一點,他望到歷史前面的希望,他望到人類終極的救贖.但,這老人家,也回憶.一路走來,曲曲折折.他想到這群人一直以來由朝到晚因著種種慾望得不到滿足而發的埋怨,他想到這群人因不信神而敬拜自造的金牛,他想到他的親信不信任他而想奪權,他想到這群人因面對的逆境而絕望,他想到這群人甘願放棄自由而想返回埃及做奴隸.前塵往事,每一幕都叫這老人家擔心,每一幕都讓他流淚.唉,這漫長的路!唉,這群軟弱的人!

在申命記,這位即將踏出歷史舞台的摩西,在神所應許的土地的邊界,一方面充滿盼望,一方面充滿憂心.如何向這群人,講最後一番話呢?

我們看看摩西的一段話:申命記30:5-10

特別留意,除了第8節,每節都強調「耶和華—你的神」要成就的事情.

30:5「耶和華－你的神必領你進入你列祖所得的地,使你可以得著；〔這是「土地」這主題〕又必善待你,使你的人數比你列祖眾多.〔這是「後裔」這主題〕

30:6耶和華－你神必將你心裡和你後裔心裡的污穢除掉,〔字面意義是「心受割禮」,受了割禮,便打上了屬於神的記號〕好叫你盡心盡性愛耶和華－你的神,使你可以存活.〔這是「與神的特殊關係」這主題,保持這關係,就得生命〕

30:7耶和華－你的神必將這一切咒詛加在你仇敵和恨惡你,逼迫你的人身上.

30:8你必歸回,聽從耶和華的話,遵行祂的一切誡命,就是我今日所吩咐你的.

30:9-10你若聽從耶和華－你神的話,謹守這律法書上所寫的誡命律例,又盡心盡性歸向耶和華－你的神,祂必使你手裡所辦的一切事,並你身所生的,牲畜所下的,地土所產的,都綽綽有餘；因為耶和華必再喜悅你,降福與你,像從前喜悅你列祖一樣.」〔再次強調「與神的特殊關係」這主題,但這回強調人的回應,若人謹守律法,保持這神人關係,就得福氣〕

最怪的是第8節:「你必歸回,聽從耶和華的話……」.「歸回」指什麼?「歸回」這字,可譯作「再次」.若譯作「再次」,即「你必再次聽從耶和華的話」,則這意味著這群人曾不聽從耶和華的話了.其實,譯作「歸回」可能更恰當,因這更能對應上文30:3的講法.申30:3說:「那時,耶和華－你的神必憐恤你,救回〔與「歸回」同一個字〕你這被擄的子民；耶和華－你的神要回轉〔與「歸回」又是同一個字〕過來,從分散你到的萬民中將你招聚回來.」如此,「歸回」這字,便直指後來猶太人被擄巴比倫的事件了.我們要明白,摩西五經在歷史流傳的過程中,猶太的神學家會將對後來事件的信仰反省,加一點點入摩西五經中,好讓摩西五經能對身處後來事件中的人產生意義.摩西五經留有很多後來者加插注腳的痕跡.若申命記30:3及30:8是指向後來猶太人被擄巴比倫的事件,則透過摩西的口,一方面言說希望,講出「耶和華—你的神」必為祂的子民成就救贖；另一方面卻表達失望,神的子民必會像他們過去所作的一樣離棄神,以致有一天,得到的土地會再次失去,成為無家可歸的人.摩西,這位老人家,在神所賜的土地的邊界,代表著猶太歷代的先知,發出信仰的悲歌.他看到神將要賜給祂的子民的土地,他看到神的子民可以有的美好前景,但他也看透世情,看透人性,看到神的子民總會離棄神,也看到神的子民離棄神之後神另一回合的救贖工程.

我們看看摩西的另一段話:申命記30:11-14

30:11「我今日所吩咐你的誡命不是你難行的,也不是離你遠的；

30:12不是在天上,使你說:『誰替我們上天取下來,使我們聽見可以遵行呢?』

30:13也不是在海外,使你說:『誰替我們過海取了來,使我們聽見可以遵行呢?』

30:14這話卻離你甚近,就在你口中,在你心裡,使你可以遵行.」

摩西這番話,強調神的誡命合乎人性,合乎情理,若人不遵守,就不是「做不到」的問題,而是「不願做」的問題.不願遵行神的誡命,有千百種理由.禮拜日返教會崇拜神,容不容易做,容易得很.但不願意這樣做的,也有千百種理由.我昨晚去酒吧飲酒飲得太夜太醉,故今天醒不來了.我今天要和家人上茶樓,畢竟很久沒見過他們了.我今天要做功課,明天要交了.「不願做」,總有理由的.但,就因為傳遞光明的人,最終沒有傳遞到光明,黑暗的力量總是很大.

我們看看摩西的再另一段話:申命記30:15-20

聽我讀的時候,請留意「今日」這字眼,及「福與禍」,「生與死」的強烈對比,也留意「土地」與「遵守神的誡命」的連繫

30:15「看哪,我今日將生與福,死與禍,陳明在你面前.

30:16〔原文這裡有:「我今日」〕吩咐你愛耶和華－你的神,遵行祂的道,謹守祂的誡命,律例,典章,使你可以存活,人數增多,耶和華－你神就必在你所要進去得為業的地上賜福與你.

30:17倘若你心裡偏離,不肯聽從,卻被勾引去敬拜事奉別神,

30:18我今日明明告訴你們,你們必要滅亡；在你過約旦河,進去得為業的地上,你的日子必不長久.

30:19我今日呼天喚地向你作見證；我將生死禍福陳明在你面前,所以你要揀選生命,使你和你的後裔都得存活；

30:20且愛耶和華－你的神,聽從祂的話,專靠祂；因為祂是你的生命,你的日子長久也在乎祂.這樣,你就可以在耶和華向你列祖亞伯拉罕,以撒,雅各起誓應許所賜的地上居住.」

經文十分強調「今日」,你明天的命運,由你「今日」的抉擇所造成.我在神學院教書,近年多了年輕人來讀神學,他們說,教會的明天看來更艱難,但就因為艱難,就越需要有人去承托.為了教會的明天,他們今日便放下平常的工作,去神學院讀書,好服事將來的教會.今日的抉擇,影響著明天.

「生與死」的對比,「土地」,「遵守神的誡命」這些主題交織在一起,讓我們再次返回創世記開頭的場景.在那裡,神創造了伊甸園,一塊讓人得福的土地.園中有分別善惡樹,考驗人是否真心遵守神的誡命.遵守的,便能留在這土地上,得到生命與福樂.不遵守的,便被逐離這土地,面對死亡的恐懼.如今,「今日」,在神的子民眼前的,就是這塊神應許的土地——迦南地.若守神的誡命,這群人就能安居這裡,永享福樂.神的誡命難守嗎?不,一點不難.但最後呢?他們又再重複亞當夏娃的抉擇,他們選擇不服從神.同一情景,不斷重複.人的墮落與神的救贖在歷史中不斷循環,歷史依然暗淡,直至人學曉堅持光明.今日,若有一樂土在我們前面,但需要我們遵守神的誡命才可進去,我們又如何抉擇呢?生與死,福與禍,人生的兩條大路,常常在我們前面,我們又如何抉擇呢?摩西很「肉緊」,很擔心,知道這群人最後一定會離棄神,但仍苦口婆心,對他們說,這不是可以「玩玩下」的抉擇,這是生與死的抉擇.選對了,就是生；選錯了,便是死.真要好好抉擇呀!

30:19寫道:「我今日呼天喚地向你作見證……」,摩西呼喚天地,為這神與人之間的生死盟約作證.我們知道,神必然守約.遵守祂誡命的,神必賜生命.離棄神的,必自招死亡.天與地會見證著,在黑暗的大地上,誰活出光明?

來到這裡,讓我們看看神的誡命的具體內容.請大家評論一下,神的誡命是否合乎人性,合乎情理,也切實可行?

讓我們看看:申命記10:12-22

10:12「以色列啊,現在耶和華－你神向你所要的是甚麼呢?只要你敬畏耶和華－你的神,遵行祂的道,愛祂,盡心盡性事奉祂,

10:13遵守祂的誡命律例,就是我今日所吩咐你的,為要叫你得福.10:14看哪,天和天上的天,地和地上所有的,都屬耶和華－你的神.

10:15耶和華但喜悅你的列祖,愛他們,從萬民中揀選他們的後裔,就是你們,像今日一樣.

10:16所以你們要將心裡的污穢除掉〔字面意義是「心受割禮」〕,不可再硬著頸項.

10:17因為耶和華－你們的神－祂是萬神之神,萬主之主,至大的神,大有能力,大而可畏,不以貌取人,也不受賄賂.

10:18祂為孤兒寡婦伸冤,又憐愛寄居的,賜給他衣食.

10:19所以你們要憐愛寄居的,因為你們在埃及地也作過寄居的.

10:20你要敬畏耶和華－你的神,事奉祂,專靠祂,也要指著祂的名起誓.

10:21祂是你所讚美的,是你的神,為你做了那大而可畏的事,是你親眼所看見的.

10:22你的列祖七十人下埃及；現在耶和華－你的神使你如同天上的星那樣多.」

是否很簡單,就是敬神愛人而已.不敬拜偶像,不向任何勢力低頭.不歧視任何人,不因私利而害人,具體地照顧一切弱者.多麼合情合理.問題是,這麼簡單的事,我們能做得到嗎?

申命記的結尾是「開放」的,生與死,容讓人去抉擇.你會如何抉擇呢?

願榮耀歸上主,願人間享和平.



\newpage

\section{開始與終結}
\label{sec:600}
\textbf{馮天聰牧師 - 開始與終結}
\newline
\newline
\hyperref[sec:599]{< < < PREV SERMON < < <}
~
\hyperref[sec:toc]{[返目錄]}
~
\hyperref[sec:601]{> > > NEXT SERMON > > >}


約翰福音 14:23-29
\newline
\begin{longtable}{cl}
\hline
\hline
章節 & 經文 (和合本修訂版)\\
\hline
14:23 & \begin{tabularx}{0.7\textwidth}{X} 耶穌回答他說：「凡愛我的人就會遵守我的道，我父也會愛他，並且我們要到他那裡去，與他同住。 \end{tabularx} \\ \\ \relax
14:24 & \begin{tabularx}{0.7\textwidth}{X} 不愛我的人就不遵守我的道。你們所聽見的道不是我的，而是差我來之父的。 \end{tabularx} \\ \\ \relax
14:25 & \begin{tabularx}{0.7\textwidth}{X} 「我還與你們在一起的時候，已對你們說了這些事。 \end{tabularx} \\ \\ \relax
14:26 & \begin{tabularx}{0.7\textwidth}{X} 但保惠師，就是父因我的名所要差來的聖靈，他要把一切的事教導你們，並且要使你們想起我對你們所說的一切話。 \end{tabularx} \\ \\ \relax
14:27 & \begin{tabularx}{0.7\textwidth}{X} 我留下平安給你們，我把我的平安賜給你們。我所賜給你們的，不像世人所賜的。你們心裡不要憂愁，也不要膽怯。 \end{tabularx} \\ \\ \relax
14:28 & \begin{tabularx}{0.7\textwidth}{X} 你們聽見我對你們說過，我去了還要回到你們這裡來。你們若愛我，就會因我到父那裡去而喜樂，因為父比我大。 \end{tabularx} \\ \\ \relax
14:29 & \begin{tabularx}{0.7\textwidth}{X} 現在事情還沒有發生，我預先告訴你們，使你們在事情發生的時候會信。 \end{tabularx} \\ \\
[1ex]
\hline
\hline
\end{longtable}
使徒行傳 16:9-15
\newline
\begin{longtable}{cl}
\hline
\hline
章節 & 經文 (和合本修訂版)\\
\hline
16:9 & \begin{tabularx}{0.7\textwidth}{X} 夜間，有異象向保羅顯現。有一個馬其頓人站著求他說：「請你過來，到馬其頓來幫助我們！」 \end{tabularx} \\ \\ \relax
16:10 & \begin{tabularx}{0.7\textwidth}{X} 保羅既看見這異象，我們就立即設法往馬其頓去，認為神呼召我們傳福音給那裡的人。 \end{tabularx} \\ \\ \relax
16:11 & \begin{tabularx}{0.7\textwidth}{X} 我們從特羅亞開船，直行駛到撒摩特喇，第二天到了尼亞坡里； \end{tabularx} \\ \\ \relax
16:12 & \begin{tabularx}{0.7\textwidth}{X} 從那裡來到腓立比，就是馬其頓這一帶的一個重要城市，也是羅馬的駐防城。我們在這城裡住了幾天。 \end{tabularx} \\ \\ \relax
16:13 & \begin{tabularx}{0.7\textwidth}{X} 在安息日，我們出城門，到了河邊，知道那裡有一個禱告的地方，我們就坐下來對那些聚會的婦女講道。 \end{tabularx} \\ \\ \relax
16:14 & \begin{tabularx}{0.7\textwidth}{X} 有一個賣紫色布的婦人，名叫呂底亞，是推雅推喇城的人，素來敬拜神。她在聽著，主就開導她的心，使她留心聽保羅所講的話。 \end{tabularx} \\ \\ \relax
16:15 & \begin{tabularx}{0.7\textwidth}{X} 她和她一家都領了洗，就求我們說：「你們若以為我是真心信主的，請到我家裡來住。」於是她堅決請我們留下。 \end{tabularx} \\ \\
[1ex]
\hline
\hline
\end{longtable}
講題:開始與終結

講員:馮天聰牧師

經文:詩篇67;約翰福音14:23-29;使徒行傳16:9-15

日期:2022年5月15日

引言:

全球疫情肆虐超過兩年,隨著疫情放緩,過去半年不少國家開始陸續開關,放寬入境的限制,讓非公民可以進入.對於很多喜歡旅行的人來說,這真是一個好消息.對於一些必需要到不同國家公幹的人來說,當然亦是一個喜訊.還有一類人會對通關非常雀躍的,就是一群宣教士.我們都知道,在疫情下,全球超過100個國家和地區頒布不同程度的隔離政策,進行全面或局部的「封城」或「禁足」命令,以阻止疫情擴散,因此很多宣教工場都被迫關閉,不少宣教士都因為避免染疫的考量,被差會及教會徵召回到自己的國家,之後無法再次回到工場；或是一些想離開但被迫滯留的宣教士,一直等待可以回家的一天.

我們需要為這些宣教事工祈禱,讓差會及眾宣教士在疫情下有智慧去回應,能繼續將福音帶給未信主的人.

自創世以來,上帝一直使用不同的人,去向世人啟示自己.以詩篇67篇為例:

詩篇67篇

1願上帝憐憫我們,賜福與我們,用臉光照我們,(細拉)

2好叫世界得知你的道路,萬國得知你的救恩.

3上帝啊,願列邦稱讚你!願萬民都稱讚你!

4願萬國都快樂歡呼；因為你必按公正審判萬民,引導世上的萬國.(細拉)

5上帝啊,願列邦稱讚你!願萬民都稱讚你!

6地已經出了土產；上帝－就是我們的上帝要賜福與我們.

7上帝要賜福與我們；地的四極都要敬畏他!

有人說詩篇

67篇是宣教的詩篇,甚至有人說是主禱文的翻版,因為全篇表明上帝國度與旨意的彰顯.這首詩篇應是每年七月猶太人守住棚節唱的詩歌,那是以色列人的新年,慶祝豐收的日子.這歌帶領守節的百姓向上帝呼求憐憫和賜福,但要注意這詩不是求上帝幫忙解決他們的需要,而是呼求「好叫世界得知你的道路,萬國得知你的救恩」,詩歌帶領百姓用上帝的心思看世界,因為上帝的心意是要用以色列作世界的光,使萬民都得著救恩.這篇詩篇使以色列人明白上帝的憐憫,賜福,光照,乃是為了讓他們可以成為祂的見證,成為一個使人認識上帝的群體；以致於世界能知道上帝的道路,萬國能曉得上帝的救恩,叫萬國萬民能一起來讚美主,領受上帝的祝福.一開始,上帝要使用以色列這個民族,去向萬國見證上帝.

然而,以色列人一直叫上帝失望.歷史中他們不斷離開上帝,遺忘上帝的教導.即使經過上帝多番的提醒與管教,甚至用上亡國與回歸作警告,他們也屢勸不改.所以,為了拯救世人,上帝開始了新的工作,差派耶穌道成肉身在我們中間,以言以身教去啟示上帝.可惜的是,人仍是拒絕耶穌,還把他殺死了.

在耶穌要走上十字架的道路上前,他在約翰福音14章23至29節留低安慰門徒及我們的說話:

約翰福音14章23至29節

23耶穌回答說:「人若愛我,就必遵守我的道；我父也必愛他,並且我們要到他那裡去,與他同住.24不愛我的人就不遵守我的道.你們所聽見的道不是我的,乃是差我來之父的道.25「我還與你們同住的時候,已將這些話對你們說了.26但保惠師,就是父因我的名所要差來的聖靈,他要將一切的事指教你們,並且要叫你們想起我對你們所說的一切話.27我留下平安給你們；我將我的平安賜給你們.我所賜的,不像世人所賜的.你們心裡不要憂愁,也不要膽怯.28你們聽見我對你們說了,我去還要到你們這裡來.你們若愛我,因我到父那裡去,就必喜樂,因為父是比我大的.29現在事情還沒有成就,我預先告訴你們,叫你們到事情成就的時候就可以信.

25-29節主要是再次去勸勉和鼓勵門徒,其中帶出聖靈保惠師的工作.耶穌即將離去,到26節中耶穌提到保惠師,就是上帝因耶穌的名而差來的聖靈,聖靈將一切的事指教門徒,保惠師這一個詞彙是源自希臘的法庭中,為申辯者作陳述,幫助和安慰的人,因此,在這裡,聖靈是門徒在世遇到困難,甚至在被控訴時的鼓勵,勸勉和申辯者,聖靈亦是在耶穌離去之後,耶穌親自臨在的表徵,令門徒隨時想起耶穌的教導.在這個令人傷感和哀傷的時刻,耶穌將要受苦,黑暗好像掌權的時刻,耶穌已預言了他的復活,他以他的平安賜給他們,這平安不是從世界上的物質和權勢中得到,反而是在受苦,困難和離別中所領受的.

在死亡臨近時,耶穌再一次鼓勵門徒不要為這一切必要來的事而憂愁和膽怯,他的離去,他地上工作的完結是他再來的第一步.他地上生命的完結將為人帶來新的生命新的開始.換言之,這個終結將為人類帶來一個新的篇章.故此耶穌勉勵門徒要為他所成就的事快樂,耶穌預示了他的復活,因他從死裡復活,成就了他的使命,而門徒就因此而信他所說的一切話.

約翰福音14章23-29節這段經文使我們看到耶穌的離去沒有即時帶來一個美好完全的境況,門徒在世上仍要倚靠聖靈活下去,這個世界的王好像仍然掌權,但耶穌再三鼓勵和勸勉他們要持守他的道和他的命令,堅持直到他的再來.

耶穌升天後,在五旬節聖靈果然如耶穌應許一樣降臨在信徒身上,自此他們便積極將福音帶到不到的地方.而後來以保羅為首的宣教隊伍,在使徒行傳16章開始了第二次的宣教旅程.這次旅程的目的本來是要培植及堅固第一次宣教旅程中所建立的教會的信徒.宣教隊伍出發後,先朝西南方向行走.可是,聖靈禁止他們,不要「他們在亞細亞講道」(6節).既然西南的路線被封,他們便轉向北面,直至他們抵達「每西亞的邊界」(7節).從這裡他們嘗試向北前行,進入庇推尼,怎料「耶穌的靈卻不許」(7節).究竟聖靈怎樣禁止他們呢?原因又是甚麼呢?聖經並沒有講明.

無論怎樣,宣教士發現西南和向北的路都被阻塞了,惟一剩下敞開的方向是西北,他們順勢到特羅亞.最後,上帝就透過晚間的異象,讓保羅的宣教隊伍清楚知道,祂的計劃是要他們渡過愛琴海到馬其頓去傳福音.兩次消極的「不准」和一次積極的「呼喚」構成了人的籌算與上帝的計劃的互動:

使徒行傳16章9至15節

9在夜間有異象現與保羅.有一個馬其頓人站著求他說:「請你過到馬其頓來幫助我們.」10保羅既看見這異象,我們隨即想要往馬其頓去,以為上帝召我們傳福音給那裡的人聽.11於是從特羅亞開船,一直行到撒摩特喇,第二天到了尼亞坡里.12從那裡來到腓立比,就是馬其頓這一方的頭一個城,也是羅馬的駐防城.我們在這城裡住了幾天.13當安息日,我們出城門,到了河邊,知道那裡有一個禱告的地方,我們就坐下對那聚會的婦女講道.14有一個賣紫色布疋的婦人,名叫呂底亞,是推雅推喇城的人,素來敬拜上帝.她聽見了,主就開導她的心,叫她留心聽保羅所講的話.15她和她一家既領了洗,便求我們說:「你們若以為我是真信主的,請到我家裡來住」；於是強留我們.

第二次宣教旅程的最突出之處是福音第一次擴展到歐洲大陸.那時代亞洲與歐洲的分隔並不顯著,因為全都在羅馬帝國的版圖之內,像是從一個省份到另一省份.但宣教士渡過愛琴海的北面時確實是離開了亞洲,進入了歐洲大地.這個區域性的轉變主因就是眾所周知的馬其頓呼聲!

宣教隊決定前往馬其頓,並非保羅個人的決定,而是宣教隊伍的決議.「保羅既看見這異象,我們隨即想要往馬其頓去,以為上帝召我們傳福音給那裡的人聽」(10節).「以為」這字有從各方不同的資料推斷出來的結論的意思.《和合本》翻成「以為」未能表達得清晰,其他中文及英文譯本都指出是有「認為」《和合本修訂版》,「認定」《新譯本》,「推斷」《呂振中譯本》,「推知」《思高譯本》,「確認」(KJV),「作出結論」(RSV,NIV,NASB)及「被說服」(NRSV)的意思.他們是有商有量,最後才得到一致的信念,認定上帝呼召他們到馬其頓去宣教.

宣教隊伍從特羅亞開船順風而行,只需兩天行程就到撒摩特利(11節).抵達後,從尼坡尼步行約十哩就到達腓立比,開展了歐洲的福音工作.敘事者自此加入了這個宣教隊伍；路加於此開始用第一身眾數「我們」的稱呼；他極有可能就是那位在夜間異象中,向保羅呼求來馬其頓幫助的那個人.

當保羅一夥人順服聆聽聖靈的引領到達腓立比時,他們見證到上帝奇妙的工作.保羅來到腓立比沒有發現猶太會堂,但卻在河邊遇上一群敬虔的婦女,其中一人是呂底亞.呂底亞是一位賣紫色布疋的商人,生活應該頗為富裕,因為紫色布疋是當時名貴的衣料,只有有錢人才能穿紫色衣服.

全世界大約有200多個國家,而我們都知道國旗是國家的象徵,有的國旗圖樣複雜,有的則是簡單大方,不過在這麼多樣式,顏色不同的國旗中,卻獨獨有一種顏色幾乎沒出現過,那就是紫色!直到近代一些比較新成立的國家才有有使用紫色這個顏色.包括尼加拉瓜共和國,多米尼克,與及1931年至1939年西班牙第二共和國時期的國旗曾一度更改為紅,黃,紫三色,後來才改為今日紅黃兩色.這是由於,過去紫色染料非常珍貴.

據說,在新約時代,紫色染料需要一滴一滴地從當時的腓尼基城邦的泰爾境內獨有的一種甲殼類昆蟲蒐集所製成,非常珍貴.由於紫色染料稀有而珍貴,所以羅馬帝國383年時曾頒布敕令,限制民間進行紫色染色工作,並將其稱作「帝王紫」.布商製作紫色染料費時又費力.所以紫色染料比等重的黃金還貴.倘若你想要一磅,那就得拿三磅黃金來換,算一算,價格差不多等於現今5萬美元.總之,當時只有有財有勢的人才買得起紫色染料.

因此,在當時從事這類買賣的都是大富商,包括呂底亞在內.路加亦形容呂底亞是「素來敬拜上帝」(14節)；路加形容那些信奉多神的外邦人,後來接受以色列一神信仰的人為「敬虔人」.我們可從她經常在安息日到禱告的地方聚會知道她對上帝的渴慕.原來猶太人最少要有十位成年男丁才可以在一個地方組織會堂,否則只可成立為一個禱告的地方.這也反映當時的腓立比以外邦人為主,猶太人數量並不多.上帝預備了呂底亞的心,她聽了保羅,主就開導她的心,使她留心聽保羅所講的話.她和她一家都領了洗.

呂底亞和家人信主後,力邀保羅和他的同工在她家中停留幾天.她擁有接待幾位成員的地方,家境必定相當；她能夠作主請客,這表明她可能是單身的婦女,她沒有結婚或丈夫已離世.呂底亞信主後,成為宣教隊伍的接待家庭.保羅後來在腓立比被囚獲釋後,第一時間也是到了呂底亞的家中作客,可見呂底亞這位歐洲第一位信徒,應是腓立比教會成立之初的具影響力的信徒之一.呂底亞信主後,一直在經濟上支持保羅的宣教工作後來提到腓立比教會他的描述.保羅在腓立比書中提到:「我初傳福音,離了馬其頓的時候,論到接受的事,除了你們以外,並沒有別的教會供給我.就是我在帖撒羅尼迦,你們也一次兩次的,打發人供給我的需用」(腓四15~16).腓立比教會如此供應保羅,相信是呂底亞的影響和資助.

今日來到腓立比也會到呂底亞教堂參觀,其中教堂外的浸池相傳是她當日受浸時同一條河.教會後來選擇在這裡建堂,並建造了這個具備東正教特色,十字型的戶外浸池來紀念第一位女信徒及她的熱心.

保羅在異象中聽到馬其頓的呼聲,當他到了腓立比,首先去到的地方是河邊的禱告會,而福音就從那裡開始,從這個祈禱會,福音自亞洲傳到了歐洲.

這跨洲際的宣教工作是在聖靈的帶領中進行,聖靈禁止保羅原先的計畫,讓他改變行進的方向.對我們來說,我們也有很多自己想做的事,可是上帝也有禁止的時候,但在禁止中祂會讓我們知道,接下來祂對我們的心意是什麼.

這兩年多的全球大疫症帶來的不單只是人命損失,短期的經濟破壞和衰退,而是更廣泛及劇烈的全球性社會及文化關係的改變.

現在也是當今大國在政治,經濟及文明上進行全球博弈的時期,東西方政治開始進入冷戰期,國際經濟體也正在進行脫鉤及「去全球化」的行動,民族與文明之間更見敵我分明,互不相讓.新冠肺炎疫症在這時候出現也發揮了催化劑的作用,讓這些社會,經濟及文化的改變大幅加劇.

「我們的世界從此不再一樣.」已是大家的共識.面對急劇和大規模的改變,教會同樣無可避免要面對.我們過去熟悉的一切也很可能不會回復到過去的樣子了!我們今天極可能是面對一個「典範轉移」(Paradigm Shift)的時刻,我們面對的是不單是翻天覆地的改變,更是全面而恆久的變化.過去我們熟悉的世界,無論是政治經濟秩序,社會文化現象,人際關係準則都在改變.擺在我們面前的是一個新的世界,我們的挑戰就是如何在這陌生的新環境中依靠聖靈回應上帝給我們的使命.

今天我們會面對怎樣的挑戰呢?又怎樣去回應呢?

1.人與人之間的矛盾衝突加劇

新冠肺炎蔓延至全球大流行,各國為求自保,實施隔離封關,限制出入境往來從而控制疫情.一些國家又搶購醫療物資,甚至互相猜疑.同時,不同族裔之間的矛盾升級,保護主義及民族主義再次抬頭,種族問題隨之而來.加上對防疫措施的不同理解,更激化了人與人之間的矛盾.

回應:保持良善

2.世界權力的失衡,全球化崩裂與產業鏈轉移帶來的貧窮問題

因為有疫苗而經濟可以重啟,人民可以自由的往來之後,那麼原本國力,經濟能力就比較強大的國家將更快的復甦,全球的國力與權力的失衡將會越來越明顯,也就是說富強者更富,更快速向前；而醫療體系不發達的國家將更難在短時間內恢復重啟,甚至整個國家要從癱瘓再起來要一個相當的時日.考慮到勞工成本,一些高勞工成本的製造業逐漸轉移至低勞工成本的國家生產.去年開始的中,美貿易戰,加上新冠肺炎疫情,西方國家紛紛鼓勵,甚至資助企業將一些具戰略性的產業(例如醫療產品)撤回本國,或分散在共他國家生產,加速了產業鏈轉移.

回應:樂於分享

3.精神健康與靈命需要

在抗疫期間,很多人面對工作,經濟,生活等問題而煩躁不安,加上禁足在家引發個人的精神,心理,靈命,以致夫婦,家庭,學習等問題,令更多人思考人生的意義及心靈健康的重要性.

回應:學習主耶穌走遍各城各鄉去關懷那裡的人,讓他們得到身,心,社,靈的幫助.

結語:

啟示錄21章10節及21章22節至22章5節

10我被聖靈感動,天使就帶我到一座高大的山,將那由上帝那裡,從天而降的聖城耶路撒冷指示我.

22我未見城內有殿,因主上帝－全能者和羔羊為城的殿.23那城內又不用日月光照；因有上帝的榮耀光照,又有羔羊為城的燈.24列國要在城的光裡行走；地上的君王必將自己的榮耀歸與那城.25城門白晝總不關閉,在那裡原沒有黑夜.26人必將列國的榮耀,尊貴歸與那城.27凡不潔淨的,並那行可憎與虛謊之事的,總不得進那城；只有名字寫在羔羊生命冊上的才得進去.

2在河這邊與那邊有生命樹,結十二樣果子,每月都結果子；樹上的葉子乃為醫治萬民.3以後再沒有咒詛；在城裡有上帝和羔羊的寶座；他的僕人都要事奉他,4也要見他的面.他的名字必寫在他們的額上.5不再有黑夜；他們也不用燈光,日光,因為主上帝要光照他們.他們要作王,直到永永遠遠.

在啟示錄寫作的時代,那時的信徒受迫害相當嚴重,隨時都會有被逮捕的可能,他們的聚會甚至被逼到必須進入墓穴中,或是到深山人煙稀少的地方去聚會.但在信仰的支持下,他們還是照樣聚會,但他們在期盼一個新的開始,就是有上帝永遠同在的世界,他們深信在這種環境之下,生命會充滿著喜悅,活力.看我們的時代,豈不是也在期盼這種沒有死亡,悲傷,哭泣,或痛苦的時代臨到我們當中?我們豈不也是經常生活在這種痛苦與不安的世代中?

在聖經最後一章啟示錄第廿二章裡,作者敘述邪惡的力量被制服之後,一個如同創造之初那種嶄新的天地出現,呈現出一片得勝的榮耀景象.在新天新地的耶路撒冷裡,有生命水湧流成生命河,在河的兩邊有生命樹.生命樹這主題出現在《創世記》第二章和第三章,但是因為人類犯罪,使得人從此與伊甸園中的生命樹隔絕.而生命樹在《啟示錄》22章重新出現,代表著上帝給人的永恆生命,而生命樹所結的果子,代表著豐盛生命的刻畫.我們對這上帝給人的永恆應許放在我們的心中,成為我們努力的一個目標.讓我們知道,原來我們將來要去的這個歸宿是如此完美,我們的終結是美好的開始,就以以終為始(begin with the end in mind),讓這份盼望鼓勵我們今天持續奔跑,保持定向,於是我們就能夠持守到底.



\newpage

\section{無情世界有情天}
\label{sec:601}
\textbf{林豪恩先生 - 無情世界有情天}
\newline
\newline
\hyperref[sec:600]{< < < PREV SERMON < < <}
~
\hyperref[sec:toc]{[返目錄]}
~
\hyperref[sec:602]{> > > NEXT SERMON > > >}


使徒行傳 16:16-34
\newline
\begin{longtable}{cl}
\hline
\hline
章節 & 經文 (和合本修訂版)\\
\hline
16:16 & \begin{tabularx}{0.7\textwidth}{X} 後來，我們往那禱告的地方去時，有一個被占卜的靈附身的使女迎面走來，她使用法術使她的主人們發了大財。 \end{tabularx} \\ \\ \relax
16:17 & \begin{tabularx}{0.7\textwidth}{X} 她跟隨保羅和我們，喊著說：「這些人是至高神的僕人，對你們傳講救人的道路。」 \end{tabularx} \\ \\ \relax
16:18 & \begin{tabularx}{0.7\textwidth}{X} 她一連好幾天這樣喊叫，保羅就心中厭煩，轉身對那靈說：「我奉耶穌基督的名吩咐你從她身上出來！」那靈立刻出來了。 \end{tabularx} \\ \\ \relax
16:19 & \begin{tabularx}{0.7\textwidth}{X} 使女的主人們見發財的指望沒有了，就揪住保羅和西拉，拉他們到市上去見官； \end{tabularx} \\ \\ \relax
16:20 & \begin{tabularx}{0.7\textwidth}{X} 又帶他們到行政官長們面前，說：「這些騷擾我們城的，他們是猶太人， \end{tabularx} \\ \\ \relax
16:21 & \begin{tabularx}{0.7\textwidth}{X} 竟傳佈我們羅馬人所不可接受、不可遵守的規矩。」 \end{tabularx} \\ \\ \relax
16:22 & \begin{tabularx}{0.7\textwidth}{X} 群眾就一齊起來攻擊他們。官長們吩咐撕開他們的衣裳，用棍子打； \end{tabularx} \\ \\ \relax
16:23 & \begin{tabularx}{0.7\textwidth}{X} 打了許多棍，就把他們下在監裡，囑咐獄警嚴緊看守。 \end{tabularx} \\ \\ \relax
16:24 & \begin{tabularx}{0.7\textwidth}{X} 獄警領了這樣的命令，就把他們下在內監，兩腳拴在木架上。 \end{tabularx} \\ \\ \relax
16:25 & \begin{tabularx}{0.7\textwidth}{X} 約在半夜，保羅和西拉正在禱告，唱詩讚美神，眾囚犯也側耳聽著的時候， \end{tabularx} \\ \\ \relax
16:26 & \begin{tabularx}{0.7\textwidth}{X} 忽然，地大震動，甚至監牢的地基都搖動了，監門立刻全開，眾囚犯的鎖鏈也都解開了。 \end{tabularx} \\ \\ \relax
16:27 & \begin{tabularx}{0.7\textwidth}{X} 獄警一醒，看見監門全開，以為囚犯已經逃走，就拔刀要自殺。 \end{tabularx} \\ \\ \relax
16:28 & \begin{tabularx}{0.7\textwidth}{X} 保羅大聲呼叫：「不要傷害自己！我們都在這裡。」 \end{tabularx} \\ \\ \relax
16:29 & \begin{tabularx}{0.7\textwidth}{X} 獄警叫人拿燈來，就衝進去，戰戰兢兢地俯伏在保羅和西拉面前。 \end{tabularx} \\ \\ \relax
16:30 & \begin{tabularx}{0.7\textwidth}{X} 然後獄警領他們出來，說：「二位先生，我必須做甚麼才可以得救？」 \end{tabularx} \\ \\ \relax
16:31 & \begin{tabularx}{0.7\textwidth}{X} 他們說：「當信主耶穌，你和你一家都必得救。」 \end{tabularx} \\ \\ \relax
16:32 & \begin{tabularx}{0.7\textwidth}{X} 他們就把主的道講給他和他全家的人聽。 \end{tabularx} \\ \\ \relax
16:33 & \begin{tabularx}{0.7\textwidth}{X} 當夜，就在那時候，獄警把他們帶去，洗他們的傷；他和他所有的家人立刻都受了洗。 \end{tabularx} \\ \\ \relax
16:34 & \begin{tabularx}{0.7\textwidth}{X} 於是獄警領他們上自己的家裡去，給他們擺上飯。他和全家的人，因為信了神，都滿心喜樂。 \end{tabularx} \\ \\
[1ex]
\hline
\hline
\end{longtable}
講題:無情世界有情天

經文:使徒行傳16:16-34

講員:林豪恩先生

經文:使徒行傳16:16-34

日期:2022年5月29日

引言

常言道:「人非草木,孰能無情」.為甚麼會有人無情地對待其他人呢?難道可以對他人的感受和痛苦視而不見嗎?使徒行傳16章16-34節就記載了一個這樣的故事.保羅和同伴到腓立比城傳福音,突然被人捏造罪名,煽動群眾攻擊,並且「未見官先打八十大板」,被陷冤獄.在人類歷史中,保羅的遭遇並不獨特,只是不幸遭遇的其中一個例子.無情地對待保羅的人,或許也不是特別邪惡,只是平常人在某些情況之下的反應.是甚麼因素令人突然變得無情呢?可能在某些狀狀之下他們不把人當作有感覺和有尊嚴的人.那麼,他們把人當作甚麼呢?

一.使女主人的無情:把人看成是利益的來源還是阻礙

16:16後來,我們往那禱告的地方去.有一個使女迎著面來,她被巫鬼所附,用法術,叫她主人們大得財利.16:17她跟隨保羅和我們,喊著說:「這些人是至高上帝的僕人,對你們傳說救人的道.」16:18她一連多日這樣喊叫,保羅就心中厭煩,轉身對那鬼說:「我奉耶穌基督的名,吩咐你從她身上出來!」那鬼當時就出來了.16:19使女的主人們見得利的指望沒有了,便揪住保羅和西拉,拉他們到市上去見首領；16:20又帶到官長面前說:「這些人原是猶太人,竟騷擾我們的城,16:21傳我們羅馬人所不可受不可行的規矩.」

使女的主人們為甚麼「揪住保羅和西拉,拉他們到市上去見首領」呢?因為「見得利的指望沒有了」(16:19).俗語說:「阻人發達猶如殺人父母」.第一個令人對人無情的因素是「把人看作利益的來源還是阻礙」,如果你是他的利益來源他就維護你,如果你與他的利益不相干就不理會你,如果是他的利益的阻礙就要清除你.當既得利益者的利益受都影響的時候,就算只是他們的個人利益受損,但為了清除阻礙,也可能捏造罪名陷害無辜.在使徒行傳中,使女的主人們不惜以公共罪名對付保羅和同伴:

第一,「這些人原是猶太人」,把個人事件提升至群族層次.

第二,「竟騷擾我們的城」,把個人事件提升至城市的公共安全層次.

第三,「傳我們羅馬人所不可受不可行的規矩」,把個人事件提升至帝國層次.

他們無限上綱地對付保羅,源頭其實是利益.為了利益,甚麼也做得出.有人認為,世界就是這樣,你不能夠期望世界講情義,要講情義,只能夠回到家中,家人才會以情相待.現實是,如果不把人看作有血有肉人,只當作利益的中介,為了利益,家人之間也可以大打出手；為了利益,有家人互相告上法庭.為了利益,有家人從此成為陌路人,甚至變成世仇.「你點可以對家人咁無情?」「因為我根本沒有把他當作有感情的人,他們只是我的利益的渠道.」有利益衝突的個人可能被利益蒙蔽,然而,一般相信群眾的眼睛是雪亮的,群眾力量可以平衡個人的無情嗎?可以為保羅他們主持公道嗎?

二.群眾的無情:從眾,把人看成是破壞分子

16:22眾人就一同起來攻擊他們.

在某些情況之下,群眾力量未必能夠保護受害者,當時腓立比城的群眾就是一個不幸的例子.他們不是那使女的主人,沒有直接經濟損失,他們為甚麼落井下石,一起攻擊受害者呢?很可能群眾也沒有把保羅和西拉等人當作有感情的人看待.群眾把他們當作甚麼呢?

如果有機會訪問當時參與攻擊保羅團隊的腓立比人,例如問他們:

「你們攻擊的人叫甚麼名字?」他們很可能答不出,最多只會說:「他們是外來人.」

「那麼你們為甚麼攻擊他們呢?」「因為他們來搞亂我們的城市.」

「他們如何搞亂你們的城市呢?」「他們做了甚麼事情呢?」「他們來你們這個城市的目的是甚麼呢?」

所有這些問題,可能他們都回答不到.因為他們聽到的只是利益受損的主人們的口號:「有外來人搞我們呀!」「他們破壞我們的社會秩序呀!」「他們威脅我們的帝國價值呀!」

在這種情況之下,群眾把保羅看作破壞分子.因此,對破壞分子,怎會以情相待呢?

如果追問下去「你怎知道他們是破壞分子呢?」他們可能回答「人人都說他們破壞分子,怎可能不是呢?」

這個令人放棄個人判斷能力及無情地對待別人的狀況可以稱為「從眾」,或者「羊群效應」,即是其他人做甚麼,我就做甚麼.人們以為「從眾」是安全的,無論是對與錯,把自己藏在群眾中,做大眾做的事,不做大眾不做的事,在這些狀況之下,比突出自己安全得多.或許那天在腓立比城起哄的人群中,有人見到保羅和西拉等人被攻擊,仍有惻隱之心,但為甚麼沒有挺身阻止呢?這樣的情況我們在新聞報導中也會見到,例如有人被侵犯,圍觀者眾,

但卻無人出手救助.

當群眾起哄的時候,最後的期望唯有是維護法紀的官長.官長帶著捕快來到,明察秋毫,主持公道,分清善惡.然而,有時現實並不符合期望.

三.官長的無情:把人看作治安的妨礙

官長吩咐剝了他們的衣裳,用棍打；

16:23打了許多棍,便將他們下在監裡,囑咐禁卒嚴緊看守.16:24禁卒領了這樣的命,就把他們下在內監裡,兩腳上了木狗.

在保羅的事件中,原來官長對被公審中的人也是無情的.官長出現不但不是救助受害者,反而不問緣由舉棍便打.官長何以如此無情呢?在官長的眼中,很可能沒有把保羅和同伴視為有血有肉有感受的人.那麼,官長視他們為甚麼呢?很可能視他們為亂象出現的源頭.官長的關注是治安,最重要是不出亂子.為了平亂,為了他管治地區的治安,順從民意清除保羅這些外人,比調查案件分清對錯有利得多.

禁卒會不會對保羅等這幾位無辜者有情呢?他們會回答:「當然不會.」「他們都是人,而且被打,你們也看見他們身體的損傷呢!有些傷口仍流著血呢?」禁卒可能回答:「官長吩咐我們監禁他們,不讓他們逃走,這是我們要執行的任務.」官長命令,執行任務,是在二次大戰之後的戰爭罪行審判法庭中,經常聽到的辯解理由之一.「你們為甚麼對無辜者施放毒氣毒殺?」「那是長官的命令,我們只是執行任務.」

四.制度的無情:把人看作執行任務的工具或對象

禁卒不但沒有把保羅他們視為有感受的人,只視為執行任務的對像,他們也沒有把自己視為有感情的人,只視為執行任務的工具,因此,當任務失敗時,他們也就沒有其他價值了.

16:25約在半夜,保羅和西拉禱告,唱詩讚美上帝,眾囚犯也側耳而聽.16:26忽然,地大震動,甚至監牢的地基都搖動了,監門立刻全開,眾囚犯的鎖鍊也都鬆開了.16:27禁卒一醒,看見監門全開,以為囚犯已經逃走,就拔刀要自殺.

任務失敗就拔刀自殺,真是悲哀.自殺或者不是禁卒自己的意願,但想到任務失敗之後也會被官長算帳,就自行了斷好了.無論是自願還是被逼,都反映出禁卒自視或被視為執行任務的工具.一個有血有肉有感情的人,自視或被視為工具,實在可惜.如何可以去工具化,恢復人性呢?

五.受屈者有情:同是天涯淪落人,把人自作自己一樣生命寶貴的人

16:28保羅大聲呼叫說:「不要傷害自己!我們都在這裡.」

如果你是被囚者,尤其是在叫天不應叫地不聞的冤獄的被囚者,監獄的門打開了,自然的反應就是走.至於那位拔刀要自殺的禁卒呢?這種不分是非助紂為虐的獄卒,也應有此報吧!他該不該死,那是他與他的長官之間的瓜葛而已.

保羅卻沒有走,更以行動制止禁卒自殺,為甚麼呢?很可能保羅看禁卒為一個有生命的人,不只是執行任務的工具,也不只是囚禁他們碌奪他們自由的勢力.如果禁卒也如自己一樣是有生命的人,禁卒的生命也是珍貴的.如果禁卒如自己一樣是有感情的人,禁卒與他的家人的感情也是珍貴的.若自殺死亡,不單單禁卒失去自己的生命,更可能是他的父母失去兒子,他的太太失去丈夫,他的兒女失去父親,他的朋友失去好友.

六.有道者有情:把人看作自己一樣在上帝眼中珍貴的聖靈

16:29禁卒叫人拿燈來,就跳進去,戰戰兢兢地俯伏在保羅,西拉面前；16:30又領他們出來,說:「二位先生,我當怎樣行才可以得救?」16:31他們說:「當信主耶穌,你和你一家都必得救.」16:32他們就把主的道講給他和他全家的人聽.

保羅和同伴視禁卒為真實的人,如自己一樣有血有肉有感情的人.他們知道自己是上主所創造和珍惜的人,也承認禁卒全為每一個人也上主創造無珍惜的人.因此,他們自己如何領受上主的恩典,如何得救,也就與禁卒全家分享這恩典,盼望禁卒全家如自己一樣得救.保羅和同伴對禁卒有情,很可能是因為他們經驗到上主對他們有情.就是這份情,開啟人心,幫助人能夠以情相待.

七.禁卒有情:互相承認和滿足身體需要

16:33當夜,就在那時候,禁卒把他們帶去,洗他們的傷；他和屬乎他的人立時都受了洗.16:34於是禁卒領他們上自己家裡去,給他們擺上飯.他和全家,因為信了上帝,都很喜樂.

保羅他們對禁卒展示的情幸運地也恢復了禁卒和家人的情,他們也以情相待保羅.禁卒不再把保羅他們只視為執行任務的對象,而是把他們視為如自己一樣的人.作為有血有肉的人,大家都有共同的基本需要,就好像這一幕,保羅他們承認禁卒的生存需要,也回應他們的靈性需要；而禁卒承認保羅他們治療傷口的需要和飲食的需要,互通有無,互相滿足,這是多麼美好的圖畫.這景況縱使是出現在監牢裡面,也彷彿是天國一樣.這一幕,讓我們想像,天國並不一定要在高貴宏偉的地方出現,而在人們以天國價值彼此相待的處境中出現.

總結

當晚,冰冷的監牢充滿了人情味.或許,第二天當他們步出監牢時,發現腓立比城雖然在陽光照耀下仍然冰冷.的確,世界沒有因為一小點的人情而被溶化.如此,可以怎樣呢?最近在網絡上讀到的這個故事.

「戰爭時要帶住家人逃難,一D 都唔容易,有位烏克蘭人需要帶成家人逃難,但係佢無jor 最重要ge 交通工具——車.arm arm 由防空洞出黎ge 佢,就見到附近舖頭,停jor 一架冇人但插住匙ge 車.心急ge 佢係附近等jor 兩個鐘頭,諗住等到車主出黎,但最後冇人出現,於是佢就決定上車睇一睇.

就好似電影入面ge 情節咁,上jor 架陌生ge 車,通常都會開個雜物箱望下,而佢就係箱入面,搵到一個電話號碼.帶住屋企人逃難ge 佢,就打jor 去呢個電話,希望打俾車主,同佢講自己偷jor 佢架車::「對唔住,我偷jor 你架車,用黎救我家人離開…」

車主:「感謝上帝!你唔使擔心,我有4架車,我用架SUV載走jor 屋企人,然後為其他車入滿油,放低車匙同我電話號碼,再放係唔同地方.」

估唔到原來呢架車,係車主特登放係街邊,留比有需要ge 人去用.而且仲唔止一架,係幾架咁多.

車主仲話同事主講「計埋你,而家全部車都有人Call我喇!希望和平之後,我地可以相見,保重!」(100毛)

人總是盼望,也需要有情相待而不是無情以對,就讓我們在無情的世界留下有情的痕跡.



\newpage

\section{清楚嗎}
\label{sec:602}
\textbf{姚志豪牧師 - 清楚嗎?}
\newline
\newline
\hyperref[sec:601]{< < < PREV SERMON < < <}
~
\hyperref[sec:toc]{[返目錄]}
~
\hyperref[sec:603]{> > > NEXT SERMON > > >}


講題:清楚嗎?

講員:姚志豪牧師

日期:2022年6月5日

大約2,3年前,一次睡醒之後,我慣常打開聖經看也顯得模糊,我不斷擦眼晴,還是沒有變化,於是就放在一旁,不理它,如是者,一次一次也是如此,向別人請教,這叫老花(遠視).我很難接受模糊的視力,因為我從來沒有近視(只有20度).模糊了,看事物不清楚,就是等於苦嗎?

根據教會傳統,聖靈降臨節亦是耶穌基督「祂受害以後,用許多確據向使徒顯明自己是活著的,在40天之中向他們顯現,並講說上帝國的事.」1:3,再過10天,(50日)後,與門徒一同聚集,聖靈的降臨(又稱五旬節),聖靈的降臨也是耶穌基督對門徒的承諾.與門徒3年的有血有肉耶穌,他的所言所行,有時都不太明白,清楚,對門徒而言尚且模糊不清,何況是一個觸摸不到的聖靈呢?特別當聖靈降臨之後說出不同的語言,似乎產生更多模糊.使2:17記述,這是先知約珥所說的,在末後的日子,就是兒女,少年,老人也要說預言,異象及異夢,大概也是不太清楚的境象.

保羅於林前13:12我們現在是對著鏡子觀看,模糊不清([13.12]「模糊不清」:原文直譯「如同猜謎」.)；到那時,就要面對面了.我如今所認識的有限,到那時就全認識,如同主認識我一樣.

明顯地,要告知當下的信徒一個事實,我們的現在是真實的,但又非真實的全部.事實上,我們知道是有限的,至少下一刻是如何,我們可以預計,卻不是確切的肯定,當中可能有很多時導致所預計的有所偏差.譬如:旅行預先看資料,到達現場的實景,大概你就會明白,我們所知不是真實的全部,有時更加會模糊不清的.回看這一千多日(3年)的香港,試問2019年之前的香港,誰會知道今天是如此呢?大概只有水晶球才預知,這點模糊,有時更令人迷失.

事實上,聖靈顯現的模糊,不是祂獨有,今天讀到耶穌基督於約翰福音的叮囑,也叫幾位門徒摸不著頭腦的.在約14:8之前,多馬與耶穌基督的對話已可見.當耶穌基督正安慰門徒,將要與他們分離,是要回去父家為眾人預備地方,而耶穌基督自己就是那條道路,真理及生命,並且承諾一定會接門徒一起去,亦明確告知,他們一定會在那裡.耶穌基督更講:我往哪裡去,你們知道那條路.14:4.多馬卻不會做:「主啊,我們不知道你去哪裡,怎麼能知道那條路呢?」好明顯對於他來說:是模糊的,都不清楚.耶穌基督再補充:「我就是道路,真理,生命；若不藉著我,沒有人能到父那裡去.」這個時候,另一個門徒腓力再問:「主啊,將父顯給我們看,我們就知足了.」顯明,他也感到耶穌基督所講的不清楚.耶穌基督反問:我與你們在一起這麼久了,你還不認識我嗎?看見我的就是看見了父,你怎麼還說「將父顯給我們看」呢?」耶穌基督的反問,假設著門徒是知道,認識的.

不論是多馬,腓力,甚至乎彼得(約13:36):「主往哪裡去?」耶穌回答說:「我所去的地方,你現在不能跟我去,後來卻要跟我去.」之後(14:22)的猶大:「主啊,為甚麼要向我們顯現,不向世人顯現呢?」似乎在講這四個門徒,其實同樣不太清楚耶穌基督所講的,為什麼?一方面可能如耶穌基督回答腓力(14:9)「我與你們同在這樣長久,你還不認識我嗎?」即門徒一直沒有意識耶穌基督所作,或覺得是理所當然,沒去思考,沒去在意；另一個可能性,就是耶穌基督刻意沒有說得清楚,其實只是一種開式方式的表達手法(OPEN ENDING).沒有說得清楚就是給予人一種自由去思考,去想像.事實上,如果你見過有些抽象派的畫,或是一些沙龍的相片,你就會明白,所謂「模糊」所給予人一種無限的想像及說不出的美.想像:亦代表不只停留在某一原點,亦是原點以外更多的可能,亦不是只見到無可能,特別在一些苦無出路之時,為人有空間創造出更多的可能.

想像也是超越人世間一直被束縛的思考.中文大學哲學系退休教授張燦輝:他引伽利略被教皇逐出教會之際,說過的一番話:「你可以不准我出聲,燒光我的書,不准我與任何人說話,不准我做任何事,但卻不能禁止我在夜間仰望星空.」就是表達我們既有的限制以外的想像空間.

耶穌基督不是說得不清不楚,只是視(門徒)我們為有血有肉會思考的人,不是一個只遵從指令的機械人,只會鸚鵡學舌,人肉錄迫音機；又或只當我們是一塊木頭搬來搬去.因此當門徒仍然不明白的時候,他就說:(14:11)「你們當信我,我在父裡面,父在我裡面；即或不信,也當因我所做的事信我.」換言之,你不明白,不清楚,也看一下我所做的事.祂做了什麼事?

約翰13章,耶穌基督為門徒洗腳,叫眾人彼此相愛,表徵著有權力的也當卑微地服侍人；約翰15章,提到上主與人是葡萄樹及枝子的關係,常時連結,並2次提到要彼此相愛(15:12,17),而且是會為朋友捨命,並且為受人無理恨惡(15:25).15:18世人若恨你們,你們知道(

或譯:該知道),恨你們以先已經恨我了.

耶穌基督為何會這樣做?祂當時在一個甚麼的處境呢?就是一個羅馬人統治猶太人的年代,而猶太的宗教領袖表面受制於他們,但權力上,猶太人的領袖也是舉足輕重,到左耶穌釘十字架時,猶太群眾要釋放巴拉巴,其實彼拉多本要釋放耶穌基督,最終也敵不過群眾的壓力,完全反映到這種的權力的角力.在這樣一個處境下,耶穌基督的教導一點也不模糊,就是叫門徒要從有權力的服侍無權的,甚至乎罪人,他們看不順眼的人.這個不單是對門徒的挑戰,也是對當權者的挑戰.因為從來只是卑微的服侍在位,怎會反過來呢?!服侍無權的,不單是對人的權力放下身段的挑戰,更是體驗到一個人對另一個的包容,接納,愛.

耶穌基督所做不單如此,還會為朋友捨命,至於捨不捨?怎樣捨?幾時捨?耶穌基督只留給我們(門徒)去思考,去選擇.事實上,每個人對信仰的理解的不同的,但重點是信仰的本質就是不只為了自己,而是為了別人,不是一個人,而是一班人聚在一起.龍應台女士(HKU醫生畢業禮的演講):一個人所抵抗的以及所堅持的,匯成一個總體,就叫做「信仰」.但是信仰,依靠的不是隆重的大聲宣告；信仰深藏在日常生活的細節裡,信仰流露在舉手投足之間最尋常最微小的決定裡.今天,你要為何服侍?你要捨棄的是?就是可以從你的想像,從你生活小節中而生,當你一直繼續堅持,將其集合埋,就是一份信仰所展現的力量.

耶穌基督所講,所做,就是超越了我們習慣的思想,釘死在框框裡,理所當然的一套,EG.PATCH ADAM.門徒所發問源於不清楚,或許只是指向更大的可能.

有一次,可4:10無人的時候,跟隨耶穌的人和十二個門徒問他這比喻的意思.

可4:11 耶穌對他們說:「上帝國的奧祕只叫你們知道,若是對外人講,凡事就用比喻.」我想起葉紹鈞《以畫為喻》,有時需要用「實用性的畫」或者「非實用性的畫」去講述一些道理.簡單而言,有些時候就是要直腸直肚講出怎樣去做「實用性的畫」,譬喻,你去操作一些機械,你就需要一些清楚的說明書(lego實用畫)每一個步驟畫清楚,斷不會叫你思考下:按下這個紅色掣／藍掣,你自己猜一下吧!但有些時間,不用說得如此清楚,就如要表達更多的情感,空間,想像,

大概就是要「非實用性的畫」.

信仰,有些時候,就是不能言傳,是需要想像,而比喻,正正就是你可以在你有的基礎上,不斷的想像.如果我問:耶穌是什麼樣子的?大概會有一個較西化的形象,因為一般人眼中的耶穌基督是從西方傳統而來,而西方傳統,又有說從裹屍布的耶穌基督影像而來,無論如何……無人真正見過耶穌基督,所以對於耶穌,不同的地方也可以有不同的想像.本是模糊不清的耶穌基督,一下子因著不同人,民族的想像多左好多可能,豐富左我們的信仰.我相信,不是因為商業決定而衍生不同款的耶穌基督,而是因為不同人的理解,去想像不同的耶穌,目的,就是讓對所想像的群體帶來一點安慰.

另一個想像:十字蓮花(崇基圖書館)

「大秦景教流行中國碑」上的十字架,聶斯脫里十字架(英語:Nestorian cross),即景教十字架,是一種源自歷史上的東方教會及今日的東方亞述教會的十字架.東方教會於唐朝時傳入中國被稱為景教.十字架是基督教的象徵.日本專家佐伯好朗,認為十字架上周圍的「珍珠」是代表教會信眾心目中信條的化身,2004年年曾任崇基學院客席教授顧衛民先生,素有研究亞洲基督教藝術史,他認為,「火」代表著信徒傳教的熱忱,下面飛雲,白雲,相信道教及伊斯蘭教的藝術手法,而蓬花就是佛教的象徵.

崇基學院前院長梁元生博士說:頂端刻上的蓮花座與十字架圖案,除了是基督教來華最古老的圖像見證外,亦象徵基督教與中國文化的相遇和會通.貫穿《十字蓮花》全書的,就是會通這兩個文化傳統,溝通中西文明的人文關懷.而事實上,更有人想像,蓮花上的十字架,就是代表著耶穌基督於壓迫他的當權者手下,仍出於污泥而不染,那種純潔.

十字蓮花／不同耶穌又好,想像,給予人很多的空間,不是死的信仰,而是活的展現.

當耶穌基督再講:世人不會接受祂,就是因為看不到,不認識祂(14:17).因此,祂會差遣聖靈(勸慰者)安慰我們,而且更講明聖靈更會指示我們,使門徒記起祂一切的話(14:27),而祂賜下聖靈,就是安慰,並且留下平安:「我將我的平安賜給你們.我所賜的,不像世人所賜的.你們心裡不要憂愁,也不要膽怯.」14:27因為耶穌基督,深明他們所遭遇的更難,壓迫只也不少,只因「世人若恨你們,你們知道(或譯:該知道),恨你們以先已經恨我了.」(15:18)

平安+勸慰—指向困難,不安,恐懼.但耶穌基督卻3次承諾他會回來(3,18,28).耶穌所謂不清不楚,可能只是我們對祂不清楚,或許是祂為我們留下這一個彈性.讓好似模糊不清,有時反而才是最美,所謂畫人不用把腸子畫出來.人生有點迷糊,就是冒險,才有驚喜；或許迷糊才渴望有造物主帶領,指引.

可能你會反問:清清楚楚不好嗎?創世故事告訴我們,人類太清楚(只有一種語言),「他們想要做的任何事,就沒有甚麼可攔阻他們了(現:為所欲為).」創11:6.使徒行傳2:6,聖靈充滿讓人說起不同的語言,那種不清楚,大家也感到納悶,(現:興奮)(原文:混亂／驚奇,興奮).人看似不好,有時只是我們不知道.

或許,所謂迷糊不清,正是曾經清楚／以為清楚,但當人生看得太清楚,或許使自己更難以入目,好像現在高科技高清／4D電視一樣.不清楚,才叫人有想像,以有動力繼續前行,因眼前所見,不是永恆不變,大概只是模糊不清蘊藏著無限的可能.



\newpage

\section{我們因信稱義嗎}
\label{sec:603}
\textbf{高國雄牧師 - 我們因「信」稱義嗎?}
\newline
\newline
\hyperref[sec:602]{< < < PREV SERMON < < <}
~
\hyperref[sec:toc]{[返目錄]}
~
\hyperref[sec:604]{> > > NEXT SERMON > > >}


羅馬書 5:1-5
\newline
\begin{longtable}{cl}
\hline
\hline
章節 & 經文 (和合本修訂版)\\
\hline
5:1 & \begin{tabularx}{0.7\textwidth}{X} 所以，我們既因信稱義，就藉著我們的主耶穌基督得以與神和好。 \end{tabularx} \\ \\ \relax
5:2 & \begin{tabularx}{0.7\textwidth}{X} 我們又藉著他，因信得以進入現在所站立的這恩典中，並且歡歡喜喜盼望神的榮耀。 \end{tabularx} \\ \\ \relax
5:3 & \begin{tabularx}{0.7\textwidth}{X} 不但如此，就是在患難中也是歡歡喜喜的，因為知道患難生忍耐， \end{tabularx} \\ \\ \relax
5:4 & \begin{tabularx}{0.7\textwidth}{X} 忍耐生老練，老練生盼望， \end{tabularx} \\ \\ \relax
5:5 & \begin{tabularx}{0.7\textwidth}{X} 盼望不至於落空，因為神的愛，已藉著所賜給我們的聖靈，澆灌在我們心裡。 \end{tabularx} \\ \\
[1ex]
\hline
\hline
\end{longtable}
講題:我們因「信」稱義嗎?

講員:高國雄牧師

經文:羅馬書5:1-5

日期:2022年6月12日

羅馬書5:1-5

1我們既因信稱義,就藉著我們的主耶穌基督得與上帝相和.

2我們又藉著他,因信得進入現在所站的這恩典中,並且歡歡喜喜盼望上帝的榮耀.

3不但如此,就是在患難中也是歡歡喜喜的；因為知道患難生忍耐,

4忍耐生老練,老練生盼望；

5盼望不至於羞恥,因為所賜給我們的聖靈將上帝的愛澆灌在我們心裡.

*因信稱義是基督新教信仰核心中的核心

宗教改革:回應當時情況而高舉「憑著信心,被上帝判為無罪」VS 傳統教會(羅馬天主教,東方正教會):「得蒙救贖需要加上善功」

人需要稱義的原因:原罪需解決

馬丁路德:贖罪券事件

馬丁路德的經歷

1505一次接近死亡邊緣(險些被閃電擊中)的經驗而回應委身

1512年他獲得神學博士學位,成為維滕貝格大學的教授

同時是聖奧斯定修會的修士,監督好幾家修院,並在教堂中當司祭

羅馬中1:17

因為上帝的義,正在這福音上顯明出來；這義是本於信以致於信.如經上所記,義人必因信得生.

馬丁路德:我極其盼望能理解保羅的『羅馬書』,一切都很順利,唯一不能解決的就是『上帝的義』這個用語.我一直以為它是指上帝大公無私,凡行不義者祂必懲罰的意思.我雖是一位無可指責的修道士,但在上帝面前我仍是良心不安的罪人,我也不知道我的善行能否平息祂的怒氣.因此我不愛這位公義而震怒的上帝,相反地我恨惡祂,並向祂發怨言.但我還是緊抓住保羅的話,誓要弄清楚祂的意思.

我日以繼夜地思索『上帝的義』的問題,直到我看出它與『義人必因信得生』的關係.不久我便明白到正是靠著『上帝的義』,上帝通過恩典和憐憫使我們因信稱義.我立時感到自己已得重生,樂園的門大開,讓我進入.《聖經》的一切話語有了新意義,從前『上帝的義』令我聞而生厭,現在它使我在更大的愛中有說不出的甜蜜.保羅這段話使我進入天堂之門.

如果你真相信基督是你的救主,你就立刻有了一位恩典的上帝,因為信心帶領你進入並打開上帝的心思意念,你就能看見完全的恩典和豐富的愛.在信心裡看上帝,你就看見他為父,為友的心腸,其中沒有憤怒,全是恩典.那看上帝是憤怒的,沒有正確地看見上帝,而只是看見一層簾子,好比有一層烏雲遮住了他的臉」.[1]

馬丁路德翻譯德文聖經時加上“唯獨”

五個唯獨

唯獨信心(sola fide),唯獨聖經(sola scriptura),唯獨恩典(Sola gratia),唯獨基督(Solo Christo),唯獨上帝的榮耀(Soli Deo gloria)

何謂稱義?

羅馬書4章

1如此說來,我們的祖宗亞伯拉罕憑著肉體得了甚麼呢?

2倘若亞伯拉罕是因行為稱義,就有可誇的；只是在上帝面前並無可誇.

3經上說甚麼呢?說:「亞伯拉罕信上帝,這就算為他的義.」

4做工的得工價,不算恩典,乃是該得的；

5惟有不做工的,只信稱罪人為義的上帝,他的信就算為義.

6正如大衛稱那在行為以外蒙上帝算為義的人是有福的.

7他說:得赦免其過,遮蓋其罪的,這人是有福的.8主不算為有罪的,這人是有福的.9如此看來,這福是單加給那受割禮的人嗎?不也是加給那未受割禮的人嗎?因我們所說,亞伯拉罕的信,就算為他的義,

10是怎麼算的呢?是在他受割禮的時候呢?是在他未受割禮的時候呢?不是在受割禮的時候,乃是在未受割禮的時候.

11並且他受了割禮的記號,作他未受割禮的時候因信稱義的印證,叫他作一切未受割禮而信之人的父,使他們也算為義；

12又作受割禮之人的父,就是那些不但受割禮,並且按我們的祖宗亞伯拉罕未受割禮而信之蹤跡去行的人.13因為上帝應許亞伯拉罕和他後裔,必得承受世界,不是因律法,乃是因信而得的義.

14若是屬乎律法的人才得為後嗣,信就歸於虛空,應許也就廢棄了.

15因為律法是惹動忿怒/叫人受刑的；哪裡沒有律法,那裡就沒有過犯.

16所以人得為後嗣是本乎信,因此就屬乎恩,叫應許定然歸給一切後裔；不但歸給那屬乎律法的,也歸給那效法亞伯拉罕之信的.

17亞伯拉罕所信的,是那叫死人復活,使無變為有的上帝,他在主面前作我們世人的父.如經上所記:「我已經立你作多國的父.」

18他在無可指望的時候,因信仍有指望,就得以作多國的父,正如先前所說:「你的後裔將要如此.」

19他將近百歲的時候,雖然想到自己的身體如同已死,撒拉的生育已經斷絕,他的信心還是不軟弱；

20並且仰望上帝的應許,總沒有因不信心裡起疑惑,反倒因信心裡得堅固,將榮耀歸給上帝,

21且滿心相信上帝所應許的必能做成.

22所以,這就算為他的義.

23「算為他義」的這句話不是單為他寫的,

24也是為我們將來得算為義之人寫的,就是我們這信上帝使我們的主耶穌從死裡復活的人.

25耶穌被交給人,是為我們的過犯；復活,是為叫我們稱義.

人在上帝面前並無可誇,包括亞伯拉罕

做工的得工價,不算恩典,乃是該得的,得恩典,蒙上帝賜福不能以任何「工」換取

所以亞伯拉罕蒙主賜福是在受割禮之前

人不是因為有律法,守律法而得稱為義

亞伯拉罕的唯一:他在無可指望的時候,因信仍有指望,就得以作多國的父

「算為他義」的這句話不是單為亞伯拉罕寫的,也是為我們將來得算為義之人寫的

耶穌基督的受死,復活,是為叫我們稱義.

羅5:1稱義,因此能夠與上帝相和

*這才是信仰的真正核心

V2因信進入所站的這恩典中

站在不同的位置,新的身份與生活方式,看世界的角度,眼光不再一樣

V2歡歡喜喜盼望上帝的榮耀

人生追求的目標不再一樣

V3-5不但如此,就是在患難中也是歡歡喜喜的；因為知道患難生忍耐,忍耐生老練,老練生盼望；盼望不至於羞恥,因為所賜給我們的聖靈將上帝的愛澆灌在我們心裡.

患難成為價值,成為生命中重要的意義

因信稱義在我們身上的意義:

不應只是神學的議題

基督信仰最公平,最好的地方,就是一切由信實的上主所應許賜的恩開始

到生命中的經歷,願意放低自己/自我,以全人的降服作回應

信心VS行為

雅各書2

14我的弟兄們,若有人說自己有信心,卻沒有行為,有甚麼益處呢?這信心能救他嗎?

15若是弟兄或是姊妹,赤身露體,又缺了日用的飲食；

16你們中間有人對他們說:「平平安安地去吧!願你們穿得暖,吃得飽」,卻不給他們身體所需用的,這有甚麼益處呢?

17這樣,信心若沒有行為就是死的.18必有人說:「你有信心,我有行為；你將你沒有行為的信心指給我看,我便藉著我的行為,將我的信心指給你看.」

19你信上帝只有一位,你信的不錯；鬼魔也信,卻是戰驚.

20虛浮的人哪,你願意知道沒有行為的信心是死的嗎?

21我們的祖宗亞伯拉罕把他兒子以撒獻在壇上,豈不是因行為稱義嗎?

22可見,信心是與他的行為並行,而且信心因著行為才得成全.

23這就應驗經上所說:「亞伯拉罕信上帝,這就算為他的義.」他又得稱為上帝的朋友.

24這樣看來,人稱義是因著行為,不是單因著信.

25妓女喇合接待使者,又放他們從別的路上出去,不也是一樣因行為稱義嗎?

26身體沒有靈魂是死的,信心沒有行為也是死的.

雖然馬丁路德早期甚至懷疑雅各書不是使徒所寫,是“稻草書信”,後來卻完全接受了

信心沒有行為是死的

亞伯拉罕:信心是與他的行為並行,而且信心因著行為才得成全

身體沒有靈魂是死的,信心沒有行為也是死的.

信心,行為是錢幣的兩面

整全的信心必須同時帶著相應的行為

有意義的行為必須內含上主的信實與及我們對主的信心

以弗所書2:

8-10

你們得救是本乎恩,也因著信；這並不是出於自己,乃是上帝所賜的；

也不是出於行為,免得有人自誇.

我們原是他的工作,在基督耶穌裡造成的,為要叫我們行善,就是上帝所預備叫我們行的.

[1] Roland Bainton, 1950  and  1978,”Here I Stand: A Life of Martin Luther”, Abingdon Press.(翻譯參古樂人,陸中石:這是我的立場,道聲出版社)



\newpage

\section{你們好好選擇吧}
\label{sec:604}
\textbf{鄧瑞強博士 - 你們好好選擇吧}
\newline
\newline
\hyperref[sec:603]{< < < PREV SERMON < < <}
~
\hyperref[sec:toc]{[返目錄]}
~
\hyperref[sec:605]{> > > NEXT SERMON > > >}


講題:你們好好選擇吧

經文:約書亞記24:1-28

講員:鄧瑞強博士

日期:2022年6月19日

各位弟兄姊妹,早安.在這疫症裡,祝願各人平安.

小孩子常常問大人:「為什麼?」小孩子問:「為什麼非洲人跑得那麼快?」大人答:「若果後面有獅子追著你,你也會跑得這麼快.」小孩子問:「為什麼進入到教堂要安靜?」大人答:「不要妨礙別人睡覺嘛.」小孩子問:「為什麼人要打仗?」大人不懂回答了.「為什麼人會受苦?」「為什麼人會死?」「為什麼人要存活下去?」這些問題,大人越來越不懂回答了.

「為什麼我要存活下去?」「存活於世,對我有什麼意義?」這些問題,或遲或早,我們總會問.特別是遭逢極大困境時,如:天災人禍,連吃的東西都沒有的時候；或戰火連天,連安睡一覺的寧靜都沒有的時候,我們便會問:「存活於世,對我有什麼意義?」當我太太在我身邊離世時,我心中也曾問:「存活下去,還有什麼意義嗎?」人存活下去,需要理由,需要意義.神學家田立克(Paul Tillich)指出,現代人充滿憂慮,有嚴重的精神問題,是因現代人比以前的人更容易失落存活的意義.找不到意義,人就可能放棄活下去.無怪乎文學家卡謬(Albert Camus)說:「只有一個真正嚴肅的哲學問題,就是自殺問題」.所謂自殺問題,即存在的意義問題.

一個人,仰視夜空,念及前塵往事,念及天地悠悠,總會思索生命的意義.摩西五經,是人思考存在意義的開始點.亞伯拉罕,在天地間,把握到一種「天命」,一種上天與他之間的連繫.這天命要貫注人間,不能單靠一個人,更要靠一群人.要成就一群明白天命的人,他們便要親身經歷上天的「好生之德」.神帶領他們離開埃及,不做奴隸,做一個自由人,這經驗塑造了這群人的靈魂.然後,神在曠野,將天命的生活內容告訴他們,他們要在人間,將神的仁愛大道向世人活現出來.要這樣做,他們需要一塊土地,成為歷史裡的一個持存群體,活現天道,改變人間.這就是摩西五經對存在意義的一步一步的思考.

藉著摩西五經的思路來思考,一個人明白到,存活於世,不只是尋索自身創造出來的意義,更重要的是連於上帝,從神那裡領受天命.不只是一個人好好地活,更是塑造一個群體.這群體不是為自身的利益而活,它要向世人展示生命的正道.這群體也不是只為當下而活,它負起了歷史的重負,要在悠悠的歷史裡,在每一個明天裡,向世人展現具有永恆意義的生活方式.透過這種生活,將人間建設成天國,將荒涼的大地建設成樂園.

在摩西五經的結尾,在申命記的最後篇章裡,將要離世的摩西,遙望那塊神應許的土地,那塊讓神的子民安居的土地,那塊讓神的子民向世界展示神的仁愛大道的土地,面對這土地,摩西可能在想,這群人的將來,將會如何?那土地的將來,將會如何?這是一未知數.在此,摩西死了.歷史的將來會如何,這未知數,永遠就是信仰的關鍵所在.神的子民,能懷著信仰,完成天命嗎?

帶著這個問題,我們來到今日要講的《約書亞記》.

一般弟兄姊妹閱讀《約書亞記》時,大多以「個人靈修」的方式讀,看著以色列人進入迦南地,看著他們與不同的民族開戰,當中有成有敗,從中領悟這些戰爭背後的靈性教訓.這種「個人靈修」的方式,有其好處.但在《約書亞記》裡,有很多神叫以色列人滅絕其他民族的經文,遇到這些經文,一般弟兄姊妹便不知如何處理.若將《約書亞記》的文體,看成是歷史事件的如實記錄,則這些經文真的很難處理.

學術上,理解《約書亞記》,是另一回事.有學者將《約書亞記》的最後編撰,定在後來猶大國被巴比倫滅國的時期.大家要明白,每個時代都會重寫歷史的.秦朝寫的秦朝歷史,與漢朝寫的秦朝歷史,與今人寫的秦朝歷史,呈現出來的意義,會很不一樣的.不斷書寫歷史,是要借鑑歷史的智慧,去面對當前的難關.《約書亞記》記述的事情,大約是主前1200年的事情,這日期只是一推斷,無法確定.約600年後,猶太人面對史無前例的劫數,猶大國被滅,猶太人被擄去巴比倫,時為主前586年.國破家亡,痛定思痛,猶太人重新回顧歷史,要去找出失去土地,被擄他鄉的神學理由.這時,他們回顧那段進入迦南地的歷史,那段600年前的往事,藉這回顧,反思他們當下失去土地的神學理由.在巴比倫,在被擄之地,猶太神學家最後編撰《約書亞記》.如此,《約書亞記》就不是歷史事件的如實記錄,而是回答眾多神學問題的神學作品.《約書亞記》的文體不是如實記錄的歷史,而是承載受苦者面對失去家園苦況的神學反思.它要回答一連串的神學問題:為什麼會亡國?為什麼會失去土地?這土地不是神應許給祂子民安居樂業之地嗎?神不是這土地的保護神嗎?失去這土地是否意味著神的無能呢?失去這土地是永遠不能改變的命運嗎?在失去土地的絕境中,神的子民還有希望嗎?《約書亞記》回答這些問題.

被擄去巴比倫的猶太人,在編撰《約書亞記》時,第一個反省的神學問題是:在歷史裡,那群在曠野流浪,無家可歸的人,靠什麼進入神應許的土地?他們回想那段歷史,那群人由曠野進入迦南地,必須跨越約旦河.《約書亞記》頭五章記述的大事,就是渡河故事.《約書亞記》3:15說,那時的河水「漲過兩岸」,是發大水的時候.這些記述,讓我們回到近東文化對大水的恐懼.在挪亞時期,這大水足以毀滅蒼生.神的救贖,就是駕馭這大水.全本聖經的第2句經文創世記1:2說:「神的靈運行在水面上」.唯有神的靈,這創發生命的靈,能駕馭這毀滅生命的大水.在《約書亞記》裡,以色列人是怎樣過河的呢?怎樣克服這大水的呢?怎樣橫越阻止人進入樂土的限制呢?他們沒有築起行軍用的浮橋,他們只是派祭司抬著約櫃,站在河水之上.如此,水便斷流了.從無家可歸之地,踏過這凶險的河水,便能進入建立家園的土地,其中靠的,是約柜.約柜裡面,有十誡.十誡代表的,是神與人之間的盟約.這盟約,這十誡,這神的律法,就是創發生命的靈,足以駕馭大水,駕馭人間的凶險,能讓人踏入神聖的土地.

踏過生命的限制,跨過人間的凶險,迎來的,是佔據了土地的迦南人.信仰從來不是一勞永逸的.信了耶穌,便等著上天堂,這真是人類精神上的一大發明.主耶穌叫我們背起十架來跟隨祂,這是一天又一天的為神的真理而奮鬥.以色列人過了約旦河,面對的是迦南人,迦南文化,迦南的神明.面對強大壓力,以色列人倚靠的,是與他們立約的神.倚靠的,是對神的信任,對神人之約中人應盡的責任的承擔,對神的十誡與律法的恭敬奉行.

失去土地,被擄去巴比倫的猶太人,第二個反省的神學問題是:我們實在是如何得到土地的?藉此,他們思考,他們如今是如何失去土地的?去到《約書亞記》第6章,是第一場「戰役」.這「戰役」,我特別用引號括著它,因它根本是一場非一般,非人間,非常態的「戰役」.這場「戰役」是如何打的?這場「戰役」根本沒有打.以色列人只是列隊巡遊,像宗教的儀仗隊,圍著耶利哥城巡遊.祭司抬著約柜,走在隊伍前頭,一邊行,一邊吹著號角.一天繞城一次,繞了六天.第七天,繞城七次,號角齊鳴,人民呼喊,然後,耶利哥城便崩塌了.在這裡,大家看到的,不是戰鬥,而是神聖的宗教禮儀.

弟兄姊妹,請問:你們認為這是一場什麼樣的「戰役」?從中得到什麼教訓?耶利哥之戰,是以色列人獲取土地的第一場「戰役」,也是以後一切「戰役」的典範.這場「戰役」最觸目的地方是:人並沒有真的打,人的角色只是敬拜神.到了下一場戰役,人試圖以自己的力去打時,便輸得一敗塗地.原來,得到安身立命的土地,靠的不是武器,而是靠對神的敬拜.以上這句話,希望普京能明白.土地,完全不是靠武力搶過來.主耶穌後來在登山寶訓裡,總結這關於土地的教訓,他說:「溫柔的人有福了,因為他們必承受地土」.土地,是神的賞賜,是神給人的禮物.在土地的事上,人的角色是信靠神,聽命於神,敬拜神.神才是土地的主人,祂勝過一切土地神,祂能任意分派土地,給聽從祂命令的人.

被擄去巴比倫的猶太神學家,藉這耶利哥之戰的故事,思考到:土地是屬於神的,一切都是神的禮物.神是大地的主,其他的「土地神」皆屬虛無.他們猶太人是與神有約的民族,應當敬拜神,聽從神的誡命.如此,他們整理出被擄去巴比倫的神學理由.這理由很簡單:被擄,失去土地,不是因為神是假的,而是因為神是真的；不是因為神無力量保護,而是因為祂有全然的大能,以致當祂的子民失信於祂時,祂有能力將土地這禮物取回.問題的核心,是神的子民違背了約,離棄了真神,走向那虛無的,迦南的土地神.這土地神名叫巴力,它主宰風調雨順,它應許富足豐盛.在我們今天,這巴力就是股神,經濟之神,金錢的神.人服膺金錢的法則,便離棄神的誡命,那些叫我們照顧弱者的誡命.人敬拜巴力,便會用金錢的暴力,去搶奪土地.當人視土地為自己的戰利品時,神便收回祂的土地了.

被擄去巴比倫的猶太神學家,痛定思痛,認定一切問題的根源,便是生命跪拜了虛無的巴力.巴力是迦南文化的核心,是宗教,經濟,政治的合體.解決問題的答案,便是:徹底與巴力劃清界線.這即是:從存在的根源深處,完全清除迦南文化的所有影響.用主耶穌的講法講,人不是拜上帝,便是拜瑪門(金錢).要清心敬拜上帝,便要徹底擺脫金錢的操控法則.用一種「戰爭」的語言來表達,便成為:徹底消滅迦南人.在這裡,我要再說一次,《約書亞記》的文體不是如實記錄的歷史,而是承載受苦者面對失去家園苦況的神學反思.作為神學反思,所謂「徹底消滅迦南人」,就不是什麼屠殺,而是在生命裡,來一次大掃除,將世間一切的虛假力量,消滅淨盡.耶利哥的經典「戰爭」,已明白提示,這是一場敬拜.唯純潔的敬拜,將生命裡的雜質完全清除.

反省至此,被擄去巴比倫的猶太人,明白到失去家園,問題不在神,問題在人.不是神無能,而是人不守約.這被擄他方的困局,是死局嗎?不是.還有希望嗎?有.希望何在?在神.神要求我們什麼?神要求我們守約,遵守律法,單單敬拜祂.我們做得到嗎?在《約書亞記》結束的章節裡,編撰《約書亞記》的神學家,將這問題舖開來,讓所有讀者自行回答.道理已講明白了,生路與死路也清楚了,人實際上又會作何選擇呢?

我們看看今日的講道經文,是《約書亞記》結束的經文:約書亞記24:

1-28,這是典型的古代近東的盟約體裁.由於經文頗長,我只能讀重點經文.

(一)盟約的開場白

24:1　約書亞將以色列的眾支派聚集在示劍,召了以色列的長老,族長,審判官,並官長來,他們就站在神面前.

(二)盟約背後的歷史:因有這歷史因由,才有這盟約

24:2　約書亞對眾民說:「耶和華－以色列的神如此說:『古時你們的列祖,就是亞伯拉罕和拿鶴的父親他拉,住在大河〔按:幼發拉底河〕那邊事奉別神,

24:3　我將你們的祖宗亞伯拉罕從大河那邊帶來,領他走遍迦南全地,又使他的子孫眾多,把以撒賜給他；「大河那邊」,即幼發拉底河的東邊,即「巴比倫」.猶太人編撰《約書亞記》時,正是在巴比倫為奴.這裡暗示有一天,會離開巴比倫.

24:4　又把雅各和以掃賜給以撒,將西珥山賜給以掃為業；後來雅各和他的子孫下到埃及去了.

24:5　我差遣摩西,亞倫,並照我在埃及中所行的降災與埃及,然後把你們領出來.………

24:13　我賜給你們地土,非你們所修治的；我賜給你們城邑,非你們所建造的.你們就住在其中,又得吃非你們所栽種的葡萄園,橄欖園的果子.』

這裡提到埃及,以色列人曾在此地為奴.他們最後離開埃及,進入迦南.他們期望,有一天,他們同樣會離開巴比倫.

在地理位置上,以色列夾在兩大帝國中間,東邊有巴比倫,南邊有埃及,以色列位於這兩大帝國貿易路線的中間,是亞洲和非洲貿易的必經之地.歷來是兵家必爭之地.以色列有時和巴比倫結盟,去反埃及,在古代,這意味和巴比倫的神明結盟.有時,他們和埃及結盟,去反巴比倫,這時,他們和埃及的神明結盟.當以色列人為了自身的安危和富足去和別的神明打交道時,也意味著他們的生活摻雜了別國的風俗,從宗教上講,這是破壞了與耶和華訂立的盟約；從生活上講,是離棄了神的律法.結果是:他們失去土地,在異地為奴.他們的盼望是:回歸純正的信仰,重回安居的樂土.

(三)盟約的條款:立約者的義務,立約者發誓遵守

24:14　「現在你們要敬畏耶和華,誠心實意地事奉他,將你們列祖在大河那邊和在埃及所事奉的神除掉,去事奉耶和華.

24:15　若是你們以事奉耶和華為不好,今日就可以選擇所要事奉的:是你們列祖在大河那邊所事奉的神呢?是你們所住這地的亞摩利人的神呢?至於我和我家,我們必定事奉耶和華.」

透過約書亞的口,對神的子民作出最大的諷刺和警告.你們有選擇的機會的,但選來選去,你們總是選擇別人的神明.約書亞代表著真正的神的子民,作出斬釘截鐵的決定:「至於我和我家,我們必定事奉耶和華.」

24:16　百姓回答說:「我們斷不敢離棄耶和華去事奉別神；

24:17　因耶和華－我們的神曾將我們和我們列祖從埃及地的為奴之家領出來,在我們眼前行了那些大神蹟,在我們所行的道上,所經過的諸國,都保護了我們.

24:18　耶和華又把住此地的亞摩利人都從我們面前趕出去.所以,我們必事奉耶和華,因為他是我們的神.」

(四)違約的咒詛

24:19　約書亞對百姓說:「你們不能事奉耶和華；因為他是聖潔的神,是忌邪的神,必不赦免你們的過犯罪惡.

24:20　你們若離棄耶和華去事奉外邦神,耶和華在降福之後,必轉而降禍與你們,把你們滅絕.」

24:21　百姓回答約書亞說:「不然,我們定要事奉耶和華.」

(五)立約的證人,立約文件的安放

24:22　約書亞對百姓說:「你們選定耶和華,要事奉他,你們自己作見證吧!」他們說:「我們願意作見證.」………

24:25　當日,約書亞就與百姓立約,在示劍為他們立定律例典章.

24:26　約書亞將這些話都寫在神的律法書上,又將一塊大石頭立在橡樹下耶和華的聖所旁邊.

24:27　約書亞對百姓說:「看哪,這石頭可以向我們作見證；因為是聽見了耶和華所吩咐我們的一切話,倘或你們背棄你們的神,這石頭就可以向你們作見證.」

那些失去土地的猶太神學家,向我們呈現他們的深刻反省.人在大地上張羅,謀算最大利益.這裡有著數,便拜這裡的神,和它勾結,活它喜悅的方式.那裡有著數,便拜那裡的神,和它勾結,活它喜悅的方式.一生只在利益中打滾,違忘了真神.最後,成了被奴役的人,或在這裡被奴役,或在那裡被奴役.《約書亞記》呼喚我們,作明智的抉擇.各位弟兄姊妹,你們會如何抉擇呢?



\newpage

\section{離開的和留下的}
\label{sec:605}
\textbf{林豪恩先生 - 離開的和留下的}
\newline
\newline
\hyperref[sec:604]{< < < PREV SERMON < < <}
~
\hyperref[sec:toc]{[返目錄]}
~
\hyperref[sec:606]{> > > NEXT SERMON > > >}


列王記下 2:1-14
\newline
\begin{longtable}{cl}
\hline
\hline
章節 & 經文 (和合本修訂版)\\
\hline
2:1 & \begin{tabularx}{0.7\textwidth}{X} 耶和華要用旋風接以利亞升天的時候，以利亞與以利沙從吉甲往前行。 \end{tabularx} \\ \\ \relax
2:2 & \begin{tabularx}{0.7\textwidth}{X} 以利亞對以利沙說：「耶和華差遣我往伯特利去，你可以留在這裡。」以利沙說：「我指著永生的耶和華，又指著你的性命起誓，我必不離開你。」於是二人下到伯特利。 \end{tabularx} \\ \\ \relax
2:3 & \begin{tabularx}{0.7\textwidth}{X} 在伯特利的先知的門徒出來，到以利沙那裡，對他說：「耶和華今日要接你的師父離開你，你知不知道？」他說：「我知道，你們不要作聲。」 \end{tabularx} \\ \\ \relax
2:4 & \begin{tabularx}{0.7\textwidth}{X} 以利亞對以利沙說：「耶和華差遣我往耶利哥去，你可以留在這裡。」以利沙說：「我指著永生的耶和華，又指著你的性命起誓，我必不離開你。」於是二人到了耶利哥。 \end{tabularx} \\ \\ \relax
2:5 & \begin{tabularx}{0.7\textwidth}{X} 在耶利哥的先知的門徒來靠近以利沙，對他說：「耶和華今日要接你的師父離開你，你知不知道？」他說：「我知道，你們不要作聲。」 \end{tabularx} \\ \\ \relax
2:6 & \begin{tabularx}{0.7\textwidth}{X} 以利亞對以利沙說：「耶和華差遣我往約旦河去，你可以留在這裡。」以利沙說：「我指著永生的耶和華，又指著你的性命起誓，我必不離開你。」於是二人一同往前行。 \end{tabularx} \\ \\ \relax
2:7 & \begin{tabularx}{0.7\textwidth}{X} 有五十個先知的門徒同去，遠遠地站在他們對面；他們二人在約旦河邊站住。 \end{tabularx} \\ \\ \relax
2:8 & \begin{tabularx}{0.7\textwidth}{X} 以利亞捲起自己的外衣，用來打水，水就左右分開，二人走乾地過去。 \end{tabularx} \\ \\ \relax
2:9 & \begin{tabularx}{0.7\textwidth}{X} 過去之後，以利亞對以利沙說：「我未被接去離開你以前，你要我為你做甚麼，只管求。」以利沙說：「願感動你的靈雙倍感動我。」 \end{tabularx} \\ \\ \relax
2:10 & \begin{tabularx}{0.7\textwidth}{X} 以利亞說：「你求的是一件難事。我被接去離開你的時候，你若看見我，就必得著；若不然，就得不著了。」 \end{tabularx} \\ \\ \relax
2:11 & \begin{tabularx}{0.7\textwidth}{X} 他們邊走邊說話的時候，看哪，有火馬和火焰車出現，把二人隔開，以利亞就乘旋風升天去了。 \end{tabularx} \\ \\ \relax
2:12 & \begin{tabularx}{0.7\textwidth}{X} 以利沙看見，就呼叫說：「我父啊！我父啊！以色列的戰車騎兵啊！」以利沙不再看見他的時候，就把自己的衣服撕為兩片。 \end{tabularx} \\ \\ \relax
2:13 & \begin{tabularx}{0.7\textwidth}{X} 他拾起以利亞身上掉下來的外衣，回去站在約旦河邊。 \end{tabularx} \\ \\ \relax
2:14 & \begin{tabularx}{0.7\textwidth}{X} 他用以利亞身上掉下來的外衣打水，說：「耶和華－以利亞的神在哪裡呢？」打水之後，水也左右分開，以利沙就過去了。 \end{tabularx} \\ \\
[1ex]
\hline
\hline
\end{longtable}
路加福音 9:51-62
\newline
\begin{longtable}{cl}
\hline
\hline
章節 & 經文 (和合本修訂版)\\
\hline
9:51 & \begin{tabularx}{0.7\textwidth}{X} 耶穌被接上升的日子將到，他決定面向耶路撒冷走去。 \end{tabularx} \\ \\ \relax
9:52 & \begin{tabularx}{0.7\textwidth}{X} 他打發使者在他前頭走；他們進了撒瑪利亞的一個村莊，要為他作準備。 \end{tabularx} \\ \\ \relax
9:53 & \begin{tabularx}{0.7\textwidth}{X} 那裡的人不接待他，因為他面向著耶路撒冷去。 \end{tabularx} \\ \\ \relax
9:54 & \begin{tabularx}{0.7\textwidth}{X} 他的門徒雅各和約翰看見了，就說：「主啊！你要我們吩咐火從天上降下來，燒滅他們嗎？」 \end{tabularx} \\ \\ \relax
9:55 & \begin{tabularx}{0.7\textwidth}{X} 耶穌轉身責備兩個門徒。 \end{tabularx} \\ \\ \relax
9:56 & \begin{tabularx}{0.7\textwidth}{X} 於是他們就往別的村莊去了。 \end{tabularx} \\ \\ \relax
9:57 & \begin{tabularx}{0.7\textwidth}{X} 他們在路上走的時候，有一個人對耶穌說：「你無論往哪裡去，我都要跟從你。」 \end{tabularx} \\ \\ \relax
9:58 & \begin{tabularx}{0.7\textwidth}{X} 耶穌對他說：「狐狸有洞，天空的飛鳥有窩，人子卻沒有枕頭的地方。」 \end{tabularx} \\ \\ \relax
9:59 & \begin{tabularx}{0.7\textwidth}{X} 他又對另一個人說：「來跟從我！」那人說：「主啊，容許我先回去埋葬我的父親。」 \end{tabularx} \\ \\ \relax
9:60 & \begin{tabularx}{0.7\textwidth}{X} 耶穌對他說：「讓死人埋葬他們的死人，你只管去傳講神的國。」 \end{tabularx} \\ \\ \relax
9:61 & \begin{tabularx}{0.7\textwidth}{X} 又有一人說：「主啊，我要跟從你，但容許我先去辭別我家裡的人。」 \end{tabularx} \\ \\ \relax
9:62 & \begin{tabularx}{0.7\textwidth}{X} 耶穌對他說：「手扶著犁向後看的人，不配進神的國。」 \end{tabularx} \\ \\
[1ex]
\hline
\hline
\end{longtable}
加拉太書 5:1-26
\newline
\begin{longtable}{cl}
\hline
\hline
章節 & 經文 (和合本修訂版)\\
\hline
5:1 & \begin{tabularx}{0.7\textwidth}{X} 基督釋放了我們，為使我們得自由。所以要站穩了，不要再被奴隸的軛挾制。 \end{tabularx} \\ \\ \relax
5:2 & \begin{tabularx}{0.7\textwidth}{X} 我—保羅告訴你們，你們若受割禮，基督就對你們無益了。 \end{tabularx} \\ \\ \relax
5:3 & \begin{tabularx}{0.7\textwidth}{X} 我再指著凡受割禮的人確實地說，他有義務遵行全部的律法。 \end{tabularx} \\ \\ \relax
5:4 & \begin{tabularx}{0.7\textwidth}{X} 你們這要靠律法稱義的是與基督隔絕，從恩典中墜落了。 \end{tabularx} \\ \\ \relax
5:5 & \begin{tabularx}{0.7\textwidth}{X} 至於我們，我們是靠著聖靈，憑著信心，等候所盼望的義。 \end{tabularx} \\ \\ \relax
5:6 & \begin{tabularx}{0.7\textwidth}{X} 因為在基督耶穌裡，受割禮不受割禮都沒有功效，惟獨使人發出仁愛的信心才有功效。 \end{tabularx} \\ \\ \relax
5:7 & \begin{tabularx}{0.7\textwidth}{X} 你們向來跑得好，誰攔阻了你們，使你們不順從真理呢？ \end{tabularx} \\ \\ \relax
5:8 & \begin{tabularx}{0.7\textwidth}{X} 這樣的勸導不是出於那召你們的。 \end{tabularx} \\ \\ \relax
5:9 & \begin{tabularx}{0.7\textwidth}{X} 一點麵酵能使全團都發起來。 \end{tabularx} \\ \\ \relax
5:10 & \begin{tabularx}{0.7\textwidth}{X} 我在主裡深信你們必不懷別樣的心；但騷擾你們的，無論是誰，必須承受懲罰。 \end{tabularx} \\ \\ \relax
5:11 & \begin{tabularx}{0.7\textwidth}{X} 弟兄們，我若仍舊傳割禮，為甚麼還受迫害呢？若是這樣，十字架絆倒人的地方就沒有了。 \end{tabularx} \\ \\ \relax
5:12 & \begin{tabularx}{0.7\textwidth}{X} 恨不得那騷擾你們的人把自己閹割了。 \end{tabularx} \\ \\ \relax
5:13 & \begin{tabularx}{0.7\textwidth}{X} 弟兄們，你們蒙召是要得自由；只是不可把這自由當作放縱情慾的機會，總要用愛心互相服侍。 \end{tabularx} \\ \\ \relax
5:14 & \begin{tabularx}{0.7\textwidth}{X} 因為全部律法都包括在「愛鄰如己」這一句話之內了。 \end{tabularx} \\ \\ \relax
5:15 & \begin{tabularx}{0.7\textwidth}{X} 你們要謹慎，你們若相咬相吞，恐怕要彼此消滅了。 \end{tabularx} \\ \\ \relax
5:16 & \begin{tabularx}{0.7\textwidth}{X} 我說，你們要順著聖靈而行，絕不可滿足肉體的情慾。 \end{tabularx} \\ \\ \relax
5:17 & \begin{tabularx}{0.7\textwidth}{X} 因為肉體的情慾和聖靈相爭，聖靈和肉體相爭，這兩個彼此敵對，使你們不能做所願意做的。 \end{tabularx} \\ \\ \relax
5:18 & \begin{tabularx}{0.7\textwidth}{X} 但你們若被聖靈引導，就不在律法之下。 \end{tabularx} \\ \\ \relax
5:19 & \begin{tabularx}{0.7\textwidth}{X} 情慾的事都是顯而易見的；就如淫亂、污穢、放蕩、 \end{tabularx} \\ \\ \relax
5:20 & \begin{tabularx}{0.7\textwidth}{X} 拜偶像、行邪術、仇恨、紛爭、忌恨、憤怒、自私、分派、結黨、 \end{tabularx} \\ \\ \relax
5:21 & \begin{tabularx}{0.7\textwidth}{X} 嫉妒、醉酒、荒宴等類。我從前告訴過你們，現在又告訴你們，做這樣事的人必不能承受神的國。 \end{tabularx} \\ \\ \relax
5:22 & \begin{tabularx}{0.7\textwidth}{X} 聖靈的果子就是仁愛、喜樂、和平、忍耐、恩慈、良善、信實、 \end{tabularx} \\ \\ \relax
5:23 & \begin{tabularx}{0.7\textwidth}{X} 溫柔、節制。這樣的事沒有律法禁止。 \end{tabularx} \\ \\ \relax
5:24 & \begin{tabularx}{0.7\textwidth}{X} 凡屬基督耶穌的人，是已經把肉體與肉體的邪情私慾同釘在十字架上了。 \end{tabularx} \\ \\ \relax
5:25 & \begin{tabularx}{0.7\textwidth}{X} 我們若靠著聖靈而活，也要靠著聖靈行事。 \end{tabularx} \\ \\ \relax
5:26 & \begin{tabularx}{0.7\textwidth}{X} 不要貪圖虛名，彼此惹氣，互相嫉妒。 \end{tabularx} \\ \\
[1ex]
\hline
\hline
\end{longtable}
詩篇 16:1-11
\newline
\begin{longtable}{cl}
\hline
\hline
章節 & 經文 (和合本修訂版)\\
\hline
16:1 & \begin{tabularx}{0.7\textwidth}{X} 神啊，求你保佑我，因為我投靠你。 \end{tabularx} \\ \\ \relax
16:2 & \begin{tabularx}{0.7\textwidth}{X} 我曾對耶和華說：「你是我的主，我的福氣惟獨從你而來。」 \end{tabularx} \\ \\ \relax
16:3 & \begin{tabularx}{0.7\textwidth}{X} 論到世上的聖民，他們是尊貴的人，是我最喜悅的。 \end{tabularx} \\ \\ \relax
16:4 & \begin{tabularx}{0.7\textwidth}{X} 追逐別神的，他們的愁苦必增加；他們所澆奠的血我不獻上，我嘴唇也不提別神的名號。 \end{tabularx} \\ \\ \relax
16:5 & \begin{tabularx}{0.7\textwidth}{X} 耶和華是我的產業，是我杯中的福分；我所得的，你為我持守。 \end{tabularx} \\ \\ \relax
16:6 & \begin{tabularx}{0.7\textwidth}{X} 用繩量給我的地界，坐落在佳美之處；我的產業實在美好。 \end{tabularx} \\ \\ \relax
16:7 & \begin{tabularx}{0.7\textwidth}{X} 我要稱頌那指引我的耶和華，在夜間我的心腸也指教我。 \end{tabularx} \\ \\ \relax
16:8 & \begin{tabularx}{0.7\textwidth}{X} 我讓耶和華常在我面前，因他在我右邊，我就不致動搖。 \end{tabularx} \\ \\ \relax
16:9 & \begin{tabularx}{0.7\textwidth}{X} 因此，我的心歡喜，我的靈快樂；我的肉身也要安然居住。 \end{tabularx} \\ \\ \relax
16:10 & \begin{tabularx}{0.7\textwidth}{X} 因為你必不將我的靈魂撇在陰間，也不讓你的聖者見地府。 \end{tabularx} \\ \\ \relax
16:11 & \begin{tabularx}{0.7\textwidth}{X} 你必將生命的道路指示我。在你面前有滿足的喜樂，在你右手中有永遠的福樂。 \end{tabularx} \\ \\
[1ex]
\hline
\hline
\end{longtable}
講題:離開的和留下的

講員:林豪恩先生

經文:列王記下2:1-14;路加福音9:51-62;加拉太書5:1-26;詩篇16:1-11

日期:2022年6月26日

一．離開的以利亞和留下的以利沙(列王記下第2章)

踏入七月,標誌著大學學年的結束,八月又是新一個學年的開始.校園那清雅的荷塘,蒼鬱的老樹,翠瓦丹柱的亭台,兩側樹木成行的大道,幽靜的庭院,芳馨襲人的中藥園,著名的雕塑,得獎的建築,以及其他或人工或天然的景物,凡此種種,讓人感覺到桃花依舊,一切如常.其實,學年交接,正是人事交接的季節,或退休的,或轉職的,每年這個時候都有教職員離開.雖不能說人面全非,也總是教人依依不捨.離開,可能一百個人有一百種心情,若然好像以利亞一樣,或許也是完美的狀況.

1．有徒弟千里相送,緊隨不離,既是尊師重道,也是師生情深

2:1耶和華要用旋風接以利亞升天的時候,以利亞與以利沙從吉甲前往.

2:2以利亞對以利沙說:「耶和華差我往伯特利去,你可以在這裡等候.」以利沙說:「我指著永生的耶和華,又敢在你面前起誓,我必不離開你.」於是二人下到伯特利.

2:6以利亞對以利沙說:「耶和華差遣我往約旦河去,你可以在這裡等候.」以利沙說:「我指著永生的耶和華,又敢在你面前起誓,我必不離開你.」於是二人一同前往.

2．有徒弟繼承衣缽,完全承接了自己的心志和能力

2:7有先知門徒去了五十人,遠遠地站在他們對面；二人在約旦河邊站住.

2:8以利亞將自己的外衣捲起來,用以打水,水就左右分開,二人走乾地而過.

2:9過去之後,以利亞對以利沙說:「我未曾被接去離開你,你要我為你做甚麼,只管求我.」以利沙說:「願感動你的靈加倍地感動我.」

2:10以利亞說:「你所求的難得.雖然如此,我被接去離開你的時候,你若看見我,就必得著；不然,必得不著了.」

2:11他們正走著說話,忽有火車火馬將二人隔開,以利亞就乘旋風升天去了.

2:12以利沙看見,就呼叫說:「我父啊!我父啊!以色列的戰車馬兵啊!」以後不再見他了.於是以利沙把自己的衣服撕為兩片.

2:13他拾起以利亞身上掉下來的外衣,回去站在約旦河邊.

2:14他用以利亞身上掉下來的外衣打水,說:「耶和華－以利亞的上帝在哪裡呢?」打水之後,水也左右分開,以利沙就過來了.

徐志摩說:「悄悄的我走了,正如我悄悄的來；我揮一揮衣袖,不帶走一片雲彩.」以利亞比徐志摩更瀟灑,因為以利亞有從天而降的火車火馬開路,然後乘風升天而去.徐志摩揮一揮衣袖,不帶走一片雲彩；以利亞卻留下他的外衣.不要輕看這件外衣,這件外衣可說是繼承的標誌.離去後的繼承若處理不善,可以留給後人無盡苦惱.武俠小說倚天屠龍記中的明教教主楊頂天突然死亡,來不及選立繼位人,結果導致明教內部群龍無首,四分五裂,相互鬥爭,明教從此衰落.若有城中富豪離世而未乃時安排好遺產分配,也難免引起後人之間的糾紛.現實上,離開,並不是如詩人揮一揮手那麼瀟瀟的.

以利亞既能瀟灑揮別,又能安心離去,因為有人繼承他的衣缽.這位繼承人所求的是甚麼呢?竟然是「願感動你的靈加倍地感動我」,以利亞別無所求了.

瀟灑離開的,可能是如徐志摩或以利亞這些高人；對我們凡夫俗子來說,能夠安心離開,已經是難得的了.離開得好的,成為了留下者的祝福.以利沙「2:14用以利亞身上掉下來的外衣打水,說:「耶和華－以利亞的上帝在哪裡呢?」打水之後,水也左右分開,以利沙就過來了.」以利沙不但在心志上承繼了以利亞的使命,更接過以利亞的外衣這個外在記號,並且在打水行動中顯出上帝的確認.離開得如此完美,以利亞夫復何求呢?

二．離開的耶穌和留下的門徒(路加福音第9章)

不過,如以利亞那樣既瀟灑又安心地離開的人可能並不多,就算是耶穌,也沒有如此優厚的待遇.

9:51耶穌被接上升的日子將到,

耶穌將要離開了,能夠瀟灑和安心嗎?看來未必.

他就定意向耶路撒冷去,

9:52便打發使者在他前頭走.他們到了撒馬利亞的一個村莊,要為他預備.

9:53那裡的人不接待他,因他面向耶路撒冷去.

耶穌的日子已經不多了,撒馬利亞村民不要那麼「小氣」可以嗎?還放不下撒馬利亞還是耶路撒冷的爭議嗎?算了吧,你不介意,不等於別人不介意；你可以放下,不等於別人都可以放下.不接待,對耶穌來說意味著不歡送,因為這次接待,就是最後一次接待了.人家不歡送就不歡送吧,你糾纏也沒有意思,因為你的路仍然要你繼續走下去.

9:54他的門徒雅各,約翰看見了,就說:「主啊,你要我們吩咐火從天上降下來燒滅他們,像以利亞所做的嗎?」

9:55耶穌轉身責備兩個門徒,說:「你們的心如何,你們並不知道.

9:56人子來不是要滅人的性命,是要救人的性命.」說著就往別的村莊去了.

不提以利亞尤自可,一提以利亞,可能令耶穌心情複雜.想到以利亞離開的時候,有一位明白師父心意又有能力承繼師父衣缽的門徒,完美交接,可能會羨慕以利亞.自己的門徒呢?外人不懂自己就算了,連門徒都不懂自己,怎能叫耶穌放心得下安心離開呢?唉,日夜相隨,言教身教三年,只是學到以天火來消滅不歡迎自己的人,叫為師的情何以堪呢?如果師父離去之後,門徒被賦予以利沙一樣的能力,將會造福世人還是危害世界呢?

撒馬利亞的村莊不接待,就往別的村莊去吧!自己的門徒教而不善,要另擇他人嗎?不要忘記,作者提示「耶穌被接上升的日子將到」,即是已經沒有時間換人了.耶穌並沒有因為自己即將離開而「放軟手腳」,也沒有進入「半退休狀態」,他珍惜最後機會,在離開之前對跟隨者盡力而教,教得多少就多少.

9:57他們走路的時候,有一人對耶穌說:「你無論往哪裡去,我要跟從你.」

9:58耶穌說:「狐狸有洞,天空的飛鳥有窩,只是人子沒有枕頭的地方.」

9:59又對一個人說:「跟從我來!」那人說:「主,容我先回去埋葬我的父親.」

9:60耶穌說:「任憑死人埋葬他們的死人,你只管去傳揚上帝國的道.」

9:61又有一人說:「主,我要跟從你,但容我先去辭別我家裡的人.」

9:62耶穌說:「手扶著犁向後看的,不配進上帝的國.」

門徒明白嗎?馬可福音的作者把答案告訴了大家,書中多次寫著「門徒仍是不明白」.面對這種狀況,怎麼能夠既瀟灑又安心地離開呢?

在座中或許有朋友能夠明白耶穌的心情.你要離開了,人人都說一句all the best,難道人人都明白你的心情嗎?或許你如耶穌一樣,心中仍有無限牽掛,仍有很多人放不下,有很多事未畫上完整句號,事實如此,又可以如何呢?這就是我們的狀況,接受吧!耶穌也沒有免卻離開的牽絆,我們怎能免卻呢?有些事情若然未能在離開之前處理好,就離開之後跟進吧!耶穌已經預備好「售後服務」,由自己跟進也好,由聖靈跟進也好.離開,並不是結局.不過,在這一刻,就算既瀟灑又安心,是前行的時候,就要起步.讓我們看看使徒保羅對起程前行的人有甚麼忠告.

三．向後看的和向前看的(加拉太書第5章)

5:1基督釋放了我們,叫我們得以自由.所以要站立得穩,不要再被奴僕的軛挾制.

5:2我－保羅告訴你們,若受割禮,基督就與你們無益了.

5:3我再指著凡受割禮的人確實地說,他是欠著行全律法的債.

5:4你們這要靠律法稱義的,是與基督隔絕,從恩典中墜落了.

5:5我們靠著聖靈,憑著信心,等候所盼望的義.

5:6原來在基督耶穌裡,受割禮不受割禮全無功效,惟獨使人生發仁愛的信心才有功效.

5:7你們向來跑得好,有誰攔阻你們,叫你們不順從真理呢?

5:8這樣的勸導不是出於那召你們的.

5:9一點麵酵能使全團都發起來.

5:10我在主裡很信你們必不懷別樣的心；但攪擾你們的,無論是誰,必擔當他的罪名.

5:11弟兄們,我若仍舊傳割禮,為甚麼還受逼迫呢?若是這樣,那十字架討厭的地方就沒有了.

5:12恨不得那攪亂你們的人把自己割絕了.

5:13弟兄們,你們蒙召是要得自由,只是不可將你們的自由當作放縱情慾的機會,總要用愛心互相服事.

5:14因為全律法都包在「愛人如己」這一句話之內了.

5:15你們要謹慎,若相咬相吞,只怕要彼此消滅了.

5:16我說,你們當順著聖靈而行,就不放縱肉體的情慾了.

5:17因為情慾和聖靈相爭,聖靈和情慾相爭,這兩個是彼此相敵,使你們不能做所願意做的.

5:18但你們若被聖靈引導,就不在律法以下.

5:19情慾的事都是顯而易見的,就如姦淫,污穢,邪蕩,

5:20拜偶像,邪術,仇恨,爭競,忌恨,惱怒,結黨,紛爭,異端,

5:21嫉妒,醉酒,荒宴等類.我從前告訴你們,現在又告訴你們,行這樣事的人必不能承受上帝的國.

5:22聖靈所結的果子,就是仁愛,喜樂,和平,忍耐,恩慈,良善,信實,

5:23溫柔,節制.這樣的事沒有律法禁止.

5:24凡屬基督耶穌的人,是已經把肉體連肉體的邪情私慾同釘在十字架上了.

5:25我們若是靠聖靈得生,就當靠聖靈行事.

5:26不要貪圖虛名,彼此惹氣,互相嫉妒.

保羅在加拉太書第5章提出了長篇的勸勉,歸納來說,就是不要再後看了,向前看吧!割禮已經是過去事,在基督裡的恩典和自由是前面的事；情慾是後面的事,所顯現的是姦淫,污穢,邪蕩,拜偶像,邪術,仇恨,爭競,忌恨,惱怒,結黨,紛爭,異端,嫉妒,醉酒,荒宴等,希望把這些都視為過去了；前面有甚麼等著我們呢?是聖靈所結的果子,就是仁愛,喜樂,和平,忍耐,恩慈,良善,信實,溫柔,節制,這些都在我們前頭等著我們去實現.既然是離開的時候了,就朝向恩典和自由吧!既然起行了,就追求聖靈所結的果子吧!

三．此心安處是吾鄉(詩篇第16篇)

離開習慣的環境,難免感到不安；放下身份和地位,難免感到失落.朝向未知的將來,怎可能不感到內心忐忑呢?詩人的話能否成為我們的安慰呢?

16:1上帝啊,求你保佑我,因為我投靠你.

16:2我的心哪,你曾對耶和華說:你是我的主；我的好處不在你以外.

16:3論到世上的聖民,他們又美又善,是我最喜悅的.

16:4以別神代替耶和華的,他們的愁苦必加增；他們所澆奠的血我不獻上；我嘴唇也不提別神的名號.

16:5耶和華是我的產業,是我杯中的分；我所得的,你為我持守.

16:6用繩量給我的地界,坐落在佳美之處；我的產業實在美好.

16:7我必稱頌那指教我的耶和華；我的心腸在夜間也警戒我.

16:8我將耶和華常擺在我面前,因他在我右邊,我便不致搖動.

16:9因此,我的心歡喜,我的靈快樂；我的肉身也要安然居住.

16:10因為你必不將我的靈魂撇在陰間,也不叫你的聖者見朽壞.

16:11你必將生命的道路指示我.在你面前有滿足的喜樂；在你右手中有永遠的福樂.

或在這處,或者那處,我們的心安於何處呢?蘇軾有一位好友名叫王鞏(字定國),因為受到使蘇軾遭殺身之禍的「烏臺詩案」牽連,被貶謫到地處嶺南荒僻之地的賓州.王定國受貶時,其歌妓柔奴毅然隨行到嶺南.公元1083年(元豐六年)王鞏北歸,出柔奴(別名寓娘)為蘇軾勸酒.蘇問及廣南風土,柔奴答以「此心安處,便是吾鄉」.蘇軾聽後,大受感動,寫了《定風波.南海歸贈王定國侍人寓娘》一詞:

常羨人間琢玉郎.天教分付點酥娘.盡道清歌傳皓齒.風起.雪飛炎海變清涼.

常常羨慕這世間如玉雕琢般丰神俊朗的男子,就連上天也憐惜他,贈予他柔美聰慧的佳人與之相伴.人人稱道那女子歌聲輕妙,笑容柔美,風起時,那歌聲如雪片飛過炎熱的夏日使世界變得清涼.

萬里歸來顏愈少.微笑.笑時猶帶嶺梅香.試問嶺南應不好.卻道.此心安處是吾鄉.

你從遙遠的地方歸來卻看起來更加年輕了,笑容依舊,笑顏裡好像還帶著嶺南梅花的清香；我問你:「嶺南的風土應該不是很好吧?」你卻坦然答道:「心安定的地方,便是我的故鄉.」

無論是離開的,還是留下的,你的心安定於甚麼呢?詩人的領悟是:

16:8我將耶和華常擺在我面前,因他在我右邊,我便不致搖動.

16:9因此,我的心歡喜,我的靈快樂；我的肉身也要安然居住.



\newpage

\section{痊癒始於放下身份}
\label{sec:606}
\textbf{林豪恩先生 - 痊癒始於放下身份}
\newline
\newline
\hyperref[sec:605]{< < < PREV SERMON < < <}
~
\hyperref[sec:toc]{[返目錄]}
~
\hyperref[sec:607]{> > > NEXT SERMON > > >}


列王記下 5:1-14
\newline
\begin{longtable}{cl}
\hline
\hline
章節 & 經文 (和合本修訂版)\\
\hline
5:1 & \begin{tabularx}{0.7\textwidth}{X} 亞蘭王的元帥乃縵在他主人面前是一個偉大的人，得王的喜悅，因為耶和華曾藉他使亞蘭人得勝。他雖然是大能的勇士，卻染上了痲瘋。 \end{tabularx} \\ \\ \relax
5:2 & \begin{tabularx}{0.7\textwidth}{X} 亞蘭人成群出征的時候，從以色列地擄了一個小女孩，她就服事乃縵的妻子。 \end{tabularx} \\ \\ \relax
5:3 & \begin{tabularx}{0.7\textwidth}{X} 她對女主人說：「我希望主人去見撒瑪利亞的先知，他必能治好主人的痲瘋。」 \end{tabularx} \\ \\ \relax
5:4 & \begin{tabularx}{0.7\textwidth}{X} 乃縵去告訴他主人說，從以色列地來的女孩如此如此說。 \end{tabularx} \\ \\ \relax
5:5 & \begin{tabularx}{0.7\textwidth}{X} 亞蘭王說：「你可以去，我也會送信給以色列王。」於是乃縵手裡帶十他連得銀子、六千舍客勒金子和十套衣裳去了。 \end{tabularx} \\ \\ \relax
5:6 & \begin{tabularx}{0.7\textwidth}{X} 他帶著這信給以色列王，說：「現在你接到這信，看哪，我派臣僕乃縵到你這裡來，你要治好他的痲瘋。」 \end{tabularx} \\ \\ \relax
5:7 & \begin{tabularx}{0.7\textwidth}{X} 以色列王讀了信就撕裂衣服，說：「我豈是神，能使人死使人活呢？這人竟派人來，叫我治好一個人的痲瘋。你們要知道，看，這人是找機會來跟我吵架的。」 \end{tabularx} \\ \\ \relax
5:8 & \begin{tabularx}{0.7\textwidth}{X} 神人以利沙聽見以色列王撕裂衣服，就派人到王那裡，說：「你為甚麼撕裂衣服呢？讓那人到我這裡來，他會知道以色列中有先知。」 \end{tabularx} \\ \\ \relax
5:9 & \begin{tabularx}{0.7\textwidth}{X} 於是乃縵帶著車馬到了以利沙的家，站在門前。 \end{tabularx} \\ \\ \relax
5:10 & \begin{tabularx}{0.7\textwidth}{X} 以利沙派一個使者，對乃縵說：「去，在約旦河中沐浴七次，你的肉就必復原，你會得潔淨。」 \end{tabularx} \\ \\ \relax
5:11 & \begin{tabularx}{0.7\textwidth}{X} 乃縵卻發怒走了。他說：「看哪，我以為他必定會出來，到我這裡，站著求告耶和華－他神的名，在患處上搖手，治好這痲瘋。 \end{tabularx} \\ \\ \relax
5:12 & \begin{tabularx}{0.7\textwidth}{X} 大馬士革的亞瑪拿河和法珥法河豈不比以色列的一切水更好嗎？我難道不可以在那裡沐浴而得潔淨嗎？」於是他生氣，轉身走了。 \end{tabularx} \\ \\ \relax
5:13 & \begin{tabularx}{0.7\textwidth}{X} 他的僕人近前來，對他說：「我父啊，先知若吩咐你做一件大事，你豈不做嗎？何況是吩咐你去沐浴，得潔淨呢？」 \end{tabularx} \\ \\ \relax
5:14 & \begin{tabularx}{0.7\textwidth}{X} 於是乃縵下去，照著神人的話，在約旦河裡浸了七次。他的肉復原，好像小孩的肉，他就潔淨了。 \end{tabularx} \\ \\
[1ex]
\hline
\hline
\end{longtable}
講題:痊癒始於放下身份

講員:林豪恩先生

經文:列王記下5:1-14

日期:2022年7月3日

題.元帥染上痲瘋

很多吸引的故事都有一個難題或挑戰要克服,列王記下5:1-14這個故事也一樣,而且起頭就直接指出難題是甚麼.

5:1亞蘭王的元帥乃縵在他主人面前是一個偉大的人,得王的喜悅,因為耶和華曾藉他使亞蘭人得勝.他雖然是大能的勇士,卻染上了痲瘋.

很見「雖然...卻」這句式結構,就知道事有蹊蹺.暗示前面提及乃縵的所有功積和優勢,都無補於事,只因為他染上了痲瘋.在乃縵的年代,處理痲瘋病人就好像當年處理新冠病人類似,一定要隔離.不但病人要被隔離,病人所接觸過的物件和住所都要被隔離.就算是元帥,是大能的勇士,「中招」之後,都變為無兵司令,生人勿近.大家有沒有發現這個故事非常弔詭,乃縵染上痲瘋病是要被隔離,但是,乃縵的痲瘋病得以痊癒,竟然是一個打破隔離,恢復接觸的過程.乃縵的痊癒要打破甚麼隔離呢?要超越的不是物理空間的隔離,因為在這個故事中醫生和病人是沒有身體接觸過的.原來,要超越的竟然是身份的隔離.身份會造成隔離的嗎?超越身份的隔離與痊癒有關係嗎?

一．放得下亞蘭人與以色列人的敵我身份嗎?

故事很直接指出難題,也很直接提出解決的方法.

5:2亞蘭人成群出征的時候,從以色列地擄了一個小女孩,她就服事乃縵的妻子.

5:3她對女主人說:「我希望主人去見撒瑪利亞的先知,他必能治好主人的痲瘋.」

正如很多現實情況一樣,問題不是沒有方法,而是如何可能用得到呢?例如若有人餓死,只要有人把剩餘的食物給飢餓的人就解決問題了,但是,如何做到呢?經文指出「亞蘭人成群出征的時候,從以色列地擄了一個小女孩」,即是亞蘭和以色列剛打過仗,兩國關係很可能處於交戰狀況.乃縵元帥很可能就是亞蘭打敗以色列的功臣.被擄到乃縵家裡的這位以色列小女孩,得知乃縵染上痳瘋病,如何回應呢?如果小女孩以自己是以色列國民,乃縵是亞蘭國民的身份來看待事情,即是自己和乃縵是敵我關係.眼前染上大痳瘋的是頭號敵人,真是大快人心,就算不能「開香檳」慶祝,也心裡暗喜.令人驚訝的是這位以色列小女孩竟然提供醫治亞蘭元帥乃縵的方法.怎可能呢?看來這位小女孩突破了以敵我身份來看待自己和乃縵.那麼,在她面前這位乃縵是誰呢?這個問題可以問一問在座中的醫護人員,在手術台上接受手術的人是誰呢?是富商?是高官?是清潔工人?是露宿者?都不是,只是一位準備接受手術的病人.感染了痳瘋的乃縵是誰呢?對這位以色列小女孩來說,他是一個痳瘋病人,而剛好她有幫助這位病人痊癒的方法,於是,她就把這方法告訴這位病人,就是那麼簡單.

乃縵呢?如果他視自己是亞蘭國的元帥,視小女孩是以色列國的戰擄,以敵對的身份來看待小女孩的行動,就會順理成章地假設敵人知道自己染上大痳瘋,一定落井下石,怎可能提供醫治方法呢?就算提供醫治方法,都很可能是陷阱.乃縵能夠突破這種敵我身份的觀念嗎?能夠相信對方提出的方法嗎?能夠接受對方的幫助嗎?

你還記得武俠小說倚天屠龍記中峨眉派掌門滅絕師太是如何喪命的嗎?當時,六大派武林人士朝廷擒捕,囚於萬安寺.明教中人往救,主事者下令放火燒塔,一舉消滅這批武林高手.火勢從下而上,眾人走頭無路,張無忌要求他們從塔上跳下來,運用乾坤大挪移在地面承接,讓跳塔者安然著地.最後剩下滅絕師太和徒弟周芷若.火勢越來越猛烈,已經燒到身邊,在這危急時刻,滅絕師太抱住周芷若跳下.差不多跌到地面,滅絕師太雙臂托起周芷若,自己加速下跌.張無忌見此情形想去救她,然而滅絕師太卻堅持自己是正派人士,不接受魔教恩惠,因此拼盡餘力將張無忌一掌打開,自己重重摔在地上,頓時脊骨斷裂而死去.

二．放得下勝利者與失敗者的身份嗎?

乃縵是誰呢?是一個染上痳瘋祈求痊癒的病人,如果單純以一個病人的身份,對於小女孩提供的醫治方法,一試也無妨,畢竟沒有其他辦法.然而,他是亞蘭國的元帥,向敵國求醫,茲事體大,他不敢輕率,否則醫了疾病而被國王砍頭,就得不償失了.於是向國王請示.

5:4乃縵去告訴他主人說,從以色列地來的女孩如此如此說.

亞蘭王又如何看待乃縵治病的需要呢?

5:5亞蘭王說:「你可以去,我也會送信給以色列王.」於是乃縵手裡帶十他連得銀子,六千舍客勒金子和十套衣裳去了.

5:6他帶著這信給以色列王,說:「現在你接到這信,看哪,我派臣僕乃縵到你這裡來,你要治好他的痲瘋.」

如果以滅絕師太的觀念來回應,當然是寧死也不向敵人救助.如果以敵我的戰略關係處理,不妨將計就計,批准乃縵去,如果能夠痊癒回來,就救回自己一位元帥；如果客死異鄉或無功而還,就為下次出戰提供藉口.我們不知道亞蘭王的心思意念,只見他寫給以色列王的信十足是吩咐下屬一樣,「你要治好他的痲瘋」.畢竟上一場仗亞蘭是戰勝國,以色列王是他的手下敗將.亞蘭王似乎以勝利者的身份看待以色列王,命令這位戰敗者執行任務.越有身份的人就越難放下身份,尤其是勝利的一方,既然佔盡優勢,當然是「有風駛盡帆」.失勢那方呢?他們又能夠放下戰敗者的身份去考慮事情嗎?其實也是不容易.

5:7以色列王讀了信就撕裂衣服,說:「我豈是神,能使人死使人活呢?這人竟派人來,叫我治好一個人的痲瘋.你們要知道,看,這人是找機會來跟我吵架的.」

以色列王和亞蘭王果然都是政治領袖,他們一致地以國王的身份來看待乃縵醫治痳瘋病的事情,也以戰勝者和戰敗者的身份來看待自己.從以色列王以及戰敗者的身份,醫治乃縵的痳瘋病不是一件醫療事件,而是一件政治事件,甚至是一件軍事事件,目的是「找機會來跟我吵架」.雖然乃縵突破了敵我身份帶著大筆醫藥費去求醫,然而,如果其他重要人士不能放下身份去回應他的需要,他的痊癒只是鏡花水月.就在這危急關頭,推倒之際,事情又有轉機.

三．放得下命令者與接受命令者的身份嗎?

5:8神人以利沙聽見以色列王撕裂衣服,就派人到王那裡,說:「你為甚麼撕裂衣服呢?讓那人到我這裡來,他會知道以色列中有先知.」

5:9於是乃縵帶著車馬到了以利沙的家,站在門前.

5:10以利沙派一個使者,對乃縵說:「去,在約旦河中沐浴七次,你的肉就必復原,你會得潔淨.」

5:11乃縵卻發怒走了.他說:「看哪,我以為他必定會出來,到我這裡,站著求告耶和華－他神的名,在患處上搖手,治好這痲瘋.

5:12大馬士革的亞嗎哪河和法珥法河豈不比以色列的一切水更好嗎?我難道不可以在那裡沐浴而得潔淨嗎?」於是他生氣,轉身走了.

我們曾經以為乃縵放下了國族間的敵我身份,只視自己為一位病人,去以色列求醫.原來,有些身份仍然是不容易放得下的.乃縵可能想:「如果要講勢力,現在是我們亞蘭的國王吩咐你這個手下敗將醫好我,你豈能不醫?如果要講人情,我已經帶了大筆醫療費用來,你醫好我後,一世都不用憂柴憂米.我們是戰勝國都那麼有誠意帶大筆錢財來,你竟然戲弄我?你不要忘記,我就算感染痳瘋,仍然是一位元帥,是我下命令吩咐人行動,不是你下命令叫我行動,何況是如此無聊的行動.」士可殺,不可辱,乃縵生氣,轉身就走.唉,此次一走,不但乃縵不能痊癒,更有可能讓亞蘭王找到藉口再次出兵,聲討以色列戲弄亞蘭元帥的罪.列王記的作者真是講故事高手,讀者剛剛以為事情解決了,怎知危機又起,如何是好呢?俗語說,高手在民間,今次化解這危機的是一位無名無姓的僕人.

5:13他的僕人近前來,對他說:「我父啊,先知若吩咐你做一件大事,你豈不做嗎?何況是吩咐你去沐浴,得潔淨呢?」

5:14於是乃縵下去,照著神人的話,在約旦河裡浸了七次.他的肉復原,好像小孩的肉,他就潔淨了.

乃縵不能放下發號命名者與接受命令者的身份,拒絕被吩咐行動；而他最後的說話:「大馬士革的亞嗎哪河和法珥法河豈不比以色列的一切水更好嗎?」顯示他再次披上國族及敵對身份.僕人向乃縵進言指出,暗示乃縵應該看以利沙是一位先知,不用視他為以色列國王之下的一位官員.先知並不是國家的代表,而是神的代表.現實上,以色列的先知經常與國王意見相左,甚至譴責國王.所以,不需要以國族身份來看待先知.至於先知的行徑,從來都是古靈精怪,出乎意料之外的,不過,既然先知是神的代表,就不需要用人間的身份來比較誰是命令者,誰是接受命令者了.

乃縵真幸運,有一位民間高手級別的僕人隨行,為他拆了這個炸彈.如果身邊的人都是慣於發號師令的軍官,或者立場先行的官員,可能又引發一場戰爭.

四．總結:能夠接受眾人在神裡面的身份嗎?

5:15乃縵帶著所有跟隨他的人,回到神人那裡,站在他面前,說:「看哪,我知道,除了以色列,全地沒有神.現在請你收下僕人的禮物.」

病醫好了,當然要給醫藥費,此乃正常行為.不過,故事的發展告訴我們,以利沙不肯收這筆醫藥費.以利沙是先知,並不是醫生,當然不收醫藥費,免得人們誤會了先知是醫生.先知的行動的回報已經顯出,就是乃縵的領悟:「看哪,我知道,除了以色列,全地沒有神.」換句話說,即是全地的人,都在一位神之下.全地的人都在一位神之下是人最原始和最基本的身份,其他的都是處境性的,都是建構出來的.

痳瘋病製造了人與人之間的隔離,其實,人的各種身份所產生的隔離效果更加嚴重及廣泛.唯有一種身份能夠打破隔離,就是上帝的兒女.如果世人能夠放下自願或被加上的各種身份,而以上帝的兒女這身份來互相看待,世界會否不再一樣呢?



\newpage

\section{好人當賊辦}
\label{sec:607}
\textbf{李均熊牧師 - 好人當賊辦}
\newline
\newline
\hyperref[sec:606]{< < < PREV SERMON < < <}
~
\hyperref[sec:toc]{[返目錄]}
~
\hyperref[sec:608]{> > > NEXT SERMON > > >}


阿摩司書 7:7-17
\newline
\begin{longtable}{cl}
\hline
\hline
章節 & 經文 (和合本修訂版)\\
\hline
7:7 & \begin{tabularx}{0.7\textwidth}{X} 他又指示我一件事，看哪，主手拿鉛垂線，站立在依鉛垂線建好的牆邊。 \end{tabularx} \\ \\ \relax
7:8 & \begin{tabularx}{0.7\textwidth}{X} 耶和華對我說：「阿摩司，你看見甚麼？」我說：「鉛垂線。」主說：「看哪，我要在我子民以色列中吊起鉛垂線，不再寬恕他們。 \end{tabularx} \\ \\ \relax
7:9 & \begin{tabularx}{0.7\textwidth}{X} 以撒的丘壇必荒涼，以色列的聖所必荒廢；我要起來用刀攻擊耶羅波安的家。」 \end{tabularx} \\ \\ \relax
7:10 & \begin{tabularx}{0.7\textwidth}{X} 伯特利的祭司亞瑪謝派人到以色列王耶羅波安那裡，說：「阿摩司在以色列家中圖謀背叛你，他所說的一切話，這地不能承擔； \end{tabularx} \\ \\ \relax
7:11 & \begin{tabularx}{0.7\textwidth}{X} 因為阿摩司這樣說：『耶羅波安必被刀殺，以色列百姓必被擄，離開本地。』」 \end{tabularx} \\ \\ \relax
7:12 & \begin{tabularx}{0.7\textwidth}{X} 於是亞瑪謝對阿摩司說：「你這先見哪，要逃到猶大地，在那裡過活，在那裡說預言； \end{tabularx} \\ \\ \relax
7:13 & \begin{tabularx}{0.7\textwidth}{X} 卻不要在伯特利再說預言，因為這裡有王的聖所，有王的宮殿。」 \end{tabularx} \\ \\ \relax
7:14 & \begin{tabularx}{0.7\textwidth}{X} 阿摩司對亞瑪謝說：「我原不是先知，也不是先知的門徒；我是牧人，是修剪桑樹的。 \end{tabularx} \\ \\ \relax
7:15 & \begin{tabularx}{0.7\textwidth}{X} 耶和華帶領我，叫我不再牧放羊群，對我說：『你去向我子民以色列說預言。』 \end{tabularx} \\ \\ \relax
7:16 & \begin{tabularx}{0.7\textwidth}{X} 「現在你要聽耶和華的話。你說：『不要向以色列說預言，也不要向以撒家傳講。』 \end{tabularx} \\ \\ \relax
7:17 & \begin{tabularx}{0.7\textwidth}{X} 所以耶和華如此說：『你的妻子要在城中作妓女，你的兒女要倒在刀下；你的地必有人用繩子量了瓜分，你自己必死在不潔淨之地；以色列百姓必被擄，離開本地。』」 \end{tabularx} \\ \\
[1ex]
\hline
\hline
\end{longtable}
講題:好人當賊辦

經課:阿摩司書7:7-17

講員:李均熊牧師

日期:2022年7月10日

起:好人當賊辦?

皇家祭壇伯特利的祭司亞瑪謝,派人與當時的北國耶戶王朝,史家稱之為耶羅波安二世的以色列王如此說:「阿摩司在以色列家中圖謀背叛你,他所說的一切話,這地不能承擔；因為阿摩司這樣說:『耶羅波安必被刀殺,以色列百姓必被擄,離開本地.』」(摩7:10)亞瑪謝引述阿摩司的說話所言非虛.因為在摩7:9的經文的確記載阿摩司接收從上主而來的信息正是:「以撒的丘壇必荒涼,以色列的聖所必荒廢；我要起來用刀攻擊耶羅波安的家.」阿摩司可能以某種說法在以色列宣講,而亞瑪謝將「用刀攻擊耶羅波家安的家」,聽成為王上本人「耶羅波安必被刀殺」,還不趕快派人轉告給王上聽.的確,這樣對王上的凶言和對國家預告可能的災難,會做成人心惶惶,令以色列國民「不能承擔」,但亞瑪謝說阿摩司「圖謀背叛」,就是好人當賊辦了.

我們從經文中看到,阿摩司得見上主審判的異象時,並不是幸災樂禍.在7:1-6節中,阿摩司見到第一個異象是蝗蟲之災,他見到這異象第一反應是向上主求饒:「主耶和華啊,求你赦免；因為雅各弱小,他怎能站立得住呢?」上主也因為如此,所以改變心意,免去這災.阿摩司見到第二個會燒盡產業的火之異象,也一樣向上主求饒:「主耶和華啊,求你止息；因為雅各弱小,他怎能站立得住呢?」可見若阿摩司要圖謀不軌,大可以讓上主施行審判,不用為以色列求饒.你說這是不是好人當賊辦?

承:先知與(建制)祭司的衝突

但先知與祭司的衝突並不是表面這麼簡單.從亞瑪謝口中,我們看到他關心的不是將要來的災禍,又或考慮阿摩司是上主的先知前來預警,讓國民能逃避災禍；亞瑪謝的報告中,並沒有說阿摩司乃本於上主的異象去宣告,也沒有說到阿摩司其他控告領袖道德敗壞的,欺壓貧困的說話,而是關心說話者是否有政治動機,國家會否落入不穩定的局面.政治先行,穩定為先,就是這些建制祭壇的祭司關心的事情?是否從上主而來的信息,先知分析本國的欺壓問題卻漠不關心?

亞瑪謝眼中的阿摩司是何許人?「你這先見哪,要逃到猶大地,在那裡過活,在那裡說預言；卻不要在伯特利再說預言,因為這裡有王的聖所,有王的宮殿.」用廣東話講:呢位師傅,你返你原居地啦,係個度搵食,係個度講預言啦；唔好係呢度伯特利講乜ye 預言,呢度係皇家重地,官家ge 地方.但阿摩司回應卻明言,他不是一般的先知,亦不是跟從那一個派別出來混日子.他原本是牧戶,是專業人士,真要找餬口也不用靠說甚麼預言.乃上主帶領他,阿摩司才向以色列說預言.這是召命,是神聖的,是從上主而來的,並不是出於政治,餬口(搵食),又或要令到人心惶惶.

轉:無人代求下的兵凶勢危

上主要審判,看來必先差遣使者好作出警告.阿摩司書第三章就明言先知的任務:

3:3二人若不同心,豈能同行呢?

3:4獅子若無獵物,豈會在林中咆哮呢?少壯獅子若無所得,豈會從洞裡吼叫呢?

3:5若未設圈套,雀鳥豈能陷入地上的羅網呢?羅網若無所得,豈會從地上翻起呢?

3:6城中若吹角,百姓豈不戰兢嗎?災禍若臨到一城,豈非耶和華所降的嗎?

3:7主耶和華不會做任何事情,除非先將奧祕指示他的僕人眾先知.

3:8獅子吼叫,誰不懼怕呢?主耶和華既已說了,誰能不說預言呢?

有因必有果,兩個人若不是相約好,又怎會在道上同行?(唔係講追女子……)獅子吼叫,城中吹角,必然有因.主上主不會做任何事情,除非先將奧秘指示他的僕人眾先知.上帝要我講,唔通我唔講囉咩?當然我還可以選擇點樣講,但你卻叫我「不要向以色列說預言,也不要向以撒家傳講.」問題就不在我身上了,而是你藐視了上主的警告,阻外上主的作為了!

阿摩司其實也早料到此著,上主不也曾透露:

4:6「我使你們在每一座城裡牙齒乾淨,使你們各處的糧食缺乏,你們仍不歸向我.」這是耶和華說的.

4:7「在收割的前三個月,我不降雨在你們那裡,我降雨在這城,不降雨在那城；這塊地有雨,那塊無雨的地就必枯乾.

4:8兩三城的人擠到一個城去找水喝,卻喝不足,你們仍不歸向我.」這是耶和華說的.

4:9「我以焚風和霉爛攻擊你們,你們許多的菜園,葡萄園,無花果樹,橄欖樹屢屢被剪蟲所吃,你們仍不歸向我.」這是耶和華說的.

4:10「我降瘟疫在你們中間,如在埃及的樣子；用刀殺戮你們的年輕人和你們遭擄掠的馬匹,營中臭氣撲鼻,你們仍不歸向我.」這是耶和華說的.

4:11「我傾覆你們,如同神從前傾覆所多瑪,蛾摩拉一樣；你們好像從火中搶救出來的一根柴,你們仍不歸向我.」這是耶和華說的.

上主三番四次的警告你們不聽,到先知明明白白的宣告審判,你們又嗤之以鼻.亞瑪謝啊,你知不知到阿摩司其實也不斷向上主求饒,現在你連這個好人都當賊辦,那你叫上主情可以堪?

其實阿摩司書七至九章阿摩司總計看見有五個審判的異象,但當中加插了這一段亞瑪謝與阿摩司的對話,有聖經學者認為這是後來的加插,但亦有聖經學者認為這個短篇的加插,正要解釋為甚麼阿摩司在之後所見的異象再無為以色列求饒.先知不再求饒,乃因建制祭司要求他走路(廣東話:過主),因此先知也不為以色列求饒了!沒有先知的求饒,建制祭司所面對的正是戰爭的悲劇:妻子在大城中作妓女,兒女倒在刀下,自己的地被人瓜分,自己又要死在不屬於上主的的不潔之地.亞瑪謝亦一語成讖:以色列百姓必被擄,離開本地.(阿摩司幾時講過會咁?)

合:面對無可避免的災禍,神子民仍有盼望嗎?

先知為何要宣告上主的審判?我們基督徒可能以為只要悔改就能免卻禍患.但悔改並不是容易的.除了要承認過錯,也要承擔罪責,回頭行善,正如阿摩司的勸告:

5:14你們要尋求良善,不要尋求邪惡,就必存活.這樣,耶和華－萬軍之神必照你們所說的與你們同在.

5:15要恨惡邪惡,喜愛良善,在城門口秉公行義；或者耶和華－萬軍之神會施恩給約瑟的餘民.

但阿摩司書充滿的卻多是審判,除了九章的終末復興,全書只有三節勸導以色列人尋求上主(5:4,6,14).先知預告的滅亡不單近在眉睫,更是不可逆轉,無可避免.現在加上領袖事事政治化,保平穩的建制思維窒息了任何對時局的批評,甚至連為國民求饒的先知也要趕絕.這個社會要被消滅只是時間問題.

但既然審判無可避免,悔改遙遙無期,那麼先知為何仍要花費唇舌,極力揚聲,迫使人聽這些凶言惡語?當然一方面這是召命,不能不說；另一方面,就是大家明白一切有因,災禍背後並不是人力的不濟,而是上主的審判和責罰.面對必然的審判,先知並不是選擇避之則吉,又或幸災樂禍,而是帶著悲憫,去代禱,去勸告,去同行.這會否也是我們可以在現時可以做的事情?

最後,面對無可避免的災禍,神子民仍有盼望嗎?約翰巴頓在他的阿摩司書的神學一書中所言:

「所以說到底,還是有希望的.

但這並不在於期望上帝立即停止災難,而是相信在災難的另一邊仍然有未來.

這種進路並不認為災難是可以避免的,但確實相信上帝可能在它之後(或通過它)還有仁慈的目的.

這本質上是何西阿書最後版本的論證……這也是阿摩司書所傳達的信息.

但可能連阿摩斯本人也沒有對這種思維方式視而不見.

如果是這樣的話,從某種意義上說,我們是兩者兼得:完全沒有希望的信息本身仍然包含著希望的種子!

這並不是因為先知認為人們會聽從他的警告並悔改；

這是因為先知相信上主最終是仁慈的,儘管他的人民仍將不得不經歷可怕的苦難,一種現在無法避免的苦難,卻朝向——也許——以某種難以想像的形式終須出現的重建.」

“So in the end, there is some hope, after all.

But it lies not in an immediate expectation that God will avert disaster, but in a sense of trust that on the other side of disaster there may still be a future.

Such an approach does not think that the disaster can be avoided, but does trust that God may still have gracious purposes beyond it (or through it).

This is essentially the argument of the book of Hosea in its finished form, and it is also……the message that the book of Amos communicates.

But it could be that even Amos himself was not blind to this way of thinking.

If that is the case, we can in a sense have our cake and eat it too: the totally unhopeful message still contains within itself the seeds of hope!

This is not because the prophet thought people would heed his warnings and repent; it was because the prophet believed that Yhwh was ultimately merciful, though his people would still have to pass through dire suffering, a suffering that could not now be averted, on the way – perhaps – to eventual restoration in some as yet unimaginable form.”

(John Barton, The Theology of the Book of Amos, p.99)

這正是阿摩司第九章的信息:

9:8看哪,主耶和華的眼目察看這有罪的國度,要把它從地面上滅絕,卻不將雅各家滅絕淨盡.」這是耶和華說的.

我們可依靠的不是自己無定向的悔改心思,而是上主的對立約子民的憐憫.惟有審判後的餘民仍可成為散播盼望的餘種.在面臨困苦災難的時間,你仍有盼望,信靠上主,仍可勇敢前行嗎?



\newpage

\section{將樂園變成了地獄}
\label{sec:608}
\textbf{鄧瑞強博士 - 將樂園變成了地獄}
\newline
\newline
\hyperref[sec:607]{< < < PREV SERMON < < <}
~
\hyperref[sec:toc]{[返目錄]}
~
\hyperref[sec:609]{> > > NEXT SERMON > > >}


士師記 19:1-30
\newline
\begin{longtable}{cl}
\hline
\hline
章節 & 經文 (和合本修訂版)\\
\hline
19:1 & \begin{tabularx}{0.7\textwidth}{X} 當以色列中沒有王的時候，有一個利未人寄居以法蓮山區的邊界，他娶了一個猶大伯利恆的女子為妾。 \end{tabularx} \\ \\ \relax
19:2 & \begin{tabularx}{0.7\textwidth}{X} 這妾對丈夫生氣，離開丈夫，回到猶大伯利恆的父家，在那裡住了四個月。 \end{tabularx} \\ \\ \relax
19:3 & \begin{tabularx}{0.7\textwidth}{X} 她的丈夫起來，帶著一個僕人、兩匹驢跟著她去，要用好話勸她回來。女子就帶丈夫進到父親家裡。女子的父親看見了他，就歡歡喜喜地迎接他。 \end{tabularx} \\ \\ \relax
19:4 & \begin{tabularx}{0.7\textwidth}{X} 這岳父，就是女子的父親，留他住了三天。他們在那裡吃喝，住宿。 \end{tabularx} \\ \\ \relax
19:5 & \begin{tabularx}{0.7\textwidth}{X} 第四日，他們清早起來，利未人起身要走，女子的父親對女婿說：「先吃點東西，加添心力，然後你們才走。」 \end{tabularx} \\ \\ \relax
19:6 & \begin{tabularx}{0.7\textwidth}{X} 於是二人坐下，一同吃喝。女子的父親對那人說：「請你答應再住一夜，使你的心舒暢。」 \end{tabularx} \\ \\ \relax
19:7 & \begin{tabularx}{0.7\textwidth}{X} 那人起身要走，他岳父挽留他，他就留下，在那裡又住了一夜。 \end{tabularx} \\ \\ \relax
19:8 & \begin{tabularx}{0.7\textwidth}{X} 第五日，他清早起來要走，女子的父親說：「來，請加添心力，留到太陽偏西吧。」於是二人一同再吃。 \end{tabularx} \\ \\ \relax
19:9 & \begin{tabularx}{0.7\textwidth}{X} 那人同他的妾和僕人起身要走，但他岳父，就是女子的父親，對他說：「看哪，太陽下山，天快晚了，你們再住一夜吧。看哪，太陽偏西了，就在這裡住宿，使你的心舒暢，明天你們一早起來上路，回你的帳棚去。」 \end{tabularx} \\ \\ \relax
19:10 & \begin{tabularx}{0.7\textwidth}{X} 那人不願再住一夜，就備上兩匹驢，帶著他的妾起身走了，來到耶布斯的對面，耶布斯就是耶路撒冷。 \end{tabularx} \\ \\ \relax
19:11 & \begin{tabularx}{0.7\textwidth}{X} 將近耶布斯的時候，太陽快下山了，僕人對主人說：「來吧，我們進這耶布斯人的城，在這裡住宿。」 \end{tabularx} \\ \\ \relax
19:12 & \begin{tabularx}{0.7\textwidth}{X} 主人對他說：「我們不可進入外邦人的城，那不是以色列人的地方，我們越過這裡到基比亞去吧。」 \end{tabularx} \\ \\ \relax
19:13 & \begin{tabularx}{0.7\textwidth}{X} 他又對僕人說：「來，讓我們到基比亞或拉瑪的一個地方住宿。」 \end{tabularx} \\ \\ \relax
19:14 & \begin{tabularx}{0.7\textwidth}{X} 於是他們越過那裡往前走，將到便雅憫的基比亞的時候，太陽已經下山了。 \end{tabularx} \\ \\ \relax
19:15 & \begin{tabularx}{0.7\textwidth}{X} 他們進入基比亞要在那裡住宿。他來坐在城裡的廣場上，但沒有人接待他們到家裡住宿。 \end{tabularx} \\ \\ \relax
19:16 & \begin{tabularx}{0.7\textwidth}{X} 看哪，晚上有一個老人從田間做工回來。他是以法蓮山區的人，寄居在基比亞；那地方的人是便雅憫人。 \end{tabularx} \\ \\ \relax
19:17 & \begin{tabularx}{0.7\textwidth}{X} 老人舉目看見那過路的人在城裡的廣場上，就說：「你從哪裡來？要到哪裡去？」 \end{tabularx} \\ \\ \relax
19:18 & \begin{tabularx}{0.7\textwidth}{X} 他對他說：「我們從猶大的伯利恆過來，要到以法蓮山區的邊界去。我是那裡的人，去了猶大的伯利恆，現在要到耶和華的家去，卻沒有人接待我到他的家。 \end{tabularx} \\ \\ \relax
19:19 & \begin{tabularx}{0.7\textwidth}{X} 其實我有飼料草料可以餵驢，我和你的使女，以及與我們在一起的僕人都有餅有酒，甚麼都不缺。」 \end{tabularx} \\ \\ \relax
19:20 & \begin{tabularx}{0.7\textwidth}{X} 老人說：「願你平安！你所需用的我都會給你們，只是不可在廣場上過夜。」 \end{tabularx} \\ \\ \relax
19:21 & \begin{tabularx}{0.7\textwidth}{X} 於是老人領他到家裡，餵上驢。他們洗了腳，就吃喝起來。 \end{tabularx} \\ \\ \relax
19:22 & \begin{tabularx}{0.7\textwidth}{X} 他們心裡歡樂的時候，看哪，城中的無賴圍住房子，連連叩門，對老人，這家的主人說：「把那進你家的人帶出來，我們要與他交合。」 \end{tabularx} \\ \\ \relax
19:23 & \begin{tabularx}{0.7\textwidth}{X} 這家的主人出來對他們說：「弟兄們，不要做這樣的惡事。這人既然進了我的家，你們就不要做這樣可恥的事。 \end{tabularx} \\ \\ \relax
19:24 & \begin{tabularx}{0.7\textwidth}{X} 看哪，我有個女兒還是處女，還有這人的妾，我把她們領出來任由你們污辱她們，就照你們看為好的對待她們吧！但對這人你們不要做這樣可恥的事。」 \end{tabularx} \\ \\ \relax
19:25 & \begin{tabularx}{0.7\textwidth}{X} 那些人卻不肯聽從他。那人抓住他的妾，把她拉出去給他們。他們強姦了她，整夜凌辱她，直到早晨，天色快亮才放她走。 \end{tabularx} \\ \\ \relax
19:26 & \begin{tabularx}{0.7\textwidth}{X} 到了早晨，婦人回來，仆倒在留她主人住宿的那人的家門前，直到天亮。 \end{tabularx} \\ \\ \relax
19:27 & \begin{tabularx}{0.7\textwidth}{X} 早晨，她的主人起來開了門，出去要上路。看哪，那婦人，他的妾倒在屋子門前，雙手搭在門檻上。 \end{tabularx} \\ \\ \relax
19:28 & \begin{tabularx}{0.7\textwidth}{X} 他對婦人說：「起來，我們走吧！」婦人卻沒有回應。那人就將她馱在驢上，起身回自己的地方去了。 \end{tabularx} \\ \\ \relax
19:29 & \begin{tabularx}{0.7\textwidth}{X} 到了家裡，他拿刀，抓住他的妾，把她的屍身切成十二塊，分送到以色列全境。 \end{tabularx} \\ \\ \relax
19:30 & \begin{tabularx}{0.7\textwidth}{X} 凡看見的人都說：「自從以色列人離開埃及地上來，直到今日，像這樣的事還沒有發生過，也沒有見過。大家應當想一想，商討一下再說。」 \end{tabularx} \\ \\
[1ex]
\hline
\hline
\end{longtable}
講題:將樂園變成了地獄

經課:士師記19:1-30

講員:鄧瑞強博士

日期:2022年7月17日

各位弟兄姊妹,早安.在這疫症裡,祝願各人平安.

1983年諾貝爾文學獎得主William

Golding的得獎作品《蒼蠅王》(Lord of the Flies)講述一群乘搭飛機逃避戰禍的孩子,因飛機被擊落,流落孤島的故事.他們都是小孩子,應是純真無邪的,這個孤島看來也似人間樂土.純真的小孩,在美麗的小島上,如果他們能活下去,這裡會否就是人間天堂?童話一定會這樣寫的,但一本直視人性的小說便沒有這種美麗的結局.故事發展下去,我們看到,人性裡本有的不同元素,會在不同孩童身上發展出來.有些小孩較理智,講規矩；有些小孩較嗜肉,要獵殺動物,較用原始的力量去解決問題.久而久之,孩童們便分成兩派,一派希望用道理來解決問題,一派卻以暴力來解決問題.最後,是一派獵殺另一派.將一個可能是樂園的地方,變成了徹頭徹尾的地獄.可惜啊,真可惜啊,人類常常將樂園變成地獄!

當領受了十誡的以色列人,進入神所應許的樂土時,我們自然會有一童話式的想像:好了,艱苦的旅程終於完成了,純良的神的子民,終於能快快樂樂地生活下去了.很多返教會的人,常常有相似的想像:好了,在人世間,終於找到完美的群體了,這裡有一群道德高尚的基督徒,在這裡以愛相連,沒有邪惡,真是人間天堂啊.當然,我們知道,這只是幻想.

約書亞死後,以色列人在迦南地,其實是一片亂局.他們散居在迦南地,與不同的民族混居.我們常常想像,這是敵我同在,我生你死的局面,其實不是這樣的,只要看一看我下次將講的《路得記》開頭的幾節經文,你會看到,就算是在士師的年代,當猶大地饑荒時,猶太人是會去摩押人那裡求助的,而摩押人並沒有拒絕,也沒有「你死你事」的敵意的.

在士師的年代,以色列人和不同民族比鄰而居,時而和睦共處,時而衝突相爭,這其實只是人類的常態.所謂士師,其實即民族英雄,或者,更似江湖大佬,在「自己人」被別人欺負時,為自己民族出頭,討回公道.不少士師亦正亦邪,例如士師基甸,他擊敗米甸人,但他卻以戰利品造了一個用來求神問卜的神像(呂振中譯本:8:27:「神諭像」),這導致以色列人後來拜這偶像.又例如士師參孫,力拔山兮氣蓋世,推倒神廟,與非利士人同歸於盡,但他生前的生活,看來糜爛不堪.士師只是一時的英雄,不是供人效法的聖人.

講《士師記》的人,大多會提到書中的某個循環規律,就是:一,神的子民拜偶像,離棄真神；二,耶和華發怒,懲罰他們,讓外族人欺凌他們；三,他們向神求救；四,神興起士師,解救他們.然後,過一段太平日子,之後,士師死了,他們又再拜偶像,離棄真神,這循環又再來一次.這其實也是我們信仰生活的寫照.生活安定,然後,慢慢忘記神,這招來了罪的苦果,我們便認罪求救,然後,神派來救星,危機解決了.一段安定的生活後,這循環又再來一次了.

今日,我們看一段《士師記》的經文,看看他們如何將神應許的樂園變成了人間地獄:士師記19:1—30

19:1當以色列中沒有王的時候,有住以法蓮山地那邊的一個利未人,娶了一個猶大伯利恆的女子為妾.

「當以色列中沒有王的時候」,這是一個「無王管」的時代,人間秩序凋零.但經文這樣說,即意味著:它等待著一個王者的來臨.這王者來臨時,能重建大地,帶來和平.《士師記》的內容記述的年代,約為主前1300-1050年,這年代當然不是精確的.但其最後編輯成書的年代,應是400多年後的事了.我不是隨便說的,19:1的前兩節經文18:30清楚提及北國的被擄,和合本這裡「遭擄掠」的用字,在舊約裡,差不多就是「被擄」的專用字眼了.事實上,以色列人最「無王管」的時代,就是後來北國以色列被擄到亞述,及南國猶大被擄到巴比倫的日子.在那些日子,他們渴望一個彌賽亞式的君王的來臨,這期望一直延伸到新約時代.問題只是,這救主王者將來自何方?這不是一雞毛蒜皮的事,在混亂的時代,誰能確保自己作為王者的正統性及合法性,是關係重大的.在中國的三國時代,魏蜀吳其中一個爭論的地方,是誰才是漢王室的正統繼承者.

「有住以法蓮山地那邊的一個利未人,娶了一個猶大伯利恆的女子為妾」,從地點的描述來看,這明顯是站在南國猶大伯利恆的位置來說的,北國的以法蓮被描述成「那邊」.在前一章,提及一見利忘義的利未人,這利未人幫殺人放火的北方支派「但」做祭司,從前文可看,這裡的利未人不是好東西.這裡看到一種南北對峙.伯利恆是後來的大衛王的家鄉,明顯,經文暗示:大衛王的血脈,才是彌賽亞式的君王的正統.

19:2妾行淫離開丈夫,回猶大的伯利恆,到了父家,在那裡住了四個月.

呂振中譯本譯:「那妾惱怒丈夫」,這翻譯更可取.因為若那妾為了「行淫」而離開丈夫,就不會回娘家.以色列人有強烈的「榮辱」(Honour and Shame)文化,若女兒不守婦道,會敗壞家聲,她的家庭就算不拿石頭打死她(中國人是「浸豬籠」了),也不可能容她回來.故此,只能說,不知何因,那妾回了娘家,但應與「行淫」無關.有學者說,是因那男人有家暴,這當然是一種猜測.

19:3她丈夫起來,帶著一個僕人,兩匹驢去見她,用好話勸她回來.女子就引丈夫進入父家.她父見了那人,便歡歡喜喜地迎接.

在古代男人至上的文化裡,能對妻妾說「好話」(直譯是「向她的心說話」),顯示這男人犯錯而氣走這妾的機會甚大.

19:10那人不願再住一夜,就備上那兩匹驢,帶著妾起身走了,來到耶布斯的對面(耶布斯就是耶路撒冷.)

19:11臨近耶布斯的時候,日頭快要落了,僕人對主人說:「我們不如進這耶布斯人的城裡住宿.」

19:12主人回答說:「我們不可進不是以色列人住的外邦城,不如過到基比亞去」；

19:13又對僕人說:「我們可以到一個地方,或住在基比亞,或住在拉瑪.」

19:14他們就往前走.將到便雅憫的基比亞,日頭已經落了；

19:15他們進入基比亞要在那裡住宿,就坐在城裡的街上,因為無人接他們進家住宿.

這裡,刻意提及「耶布斯」(耶布斯就是耶路撒冷,即大衛王朝的京城),但他們不選擇留在這裡,他們離開後來大衛王朝的根據地,後果堪虞.最後,仍有兩個地方可選擇,一是基比亞,一是拉瑪.基比亞是後來以色列人第一位君王掃羅的家鄉.拉瑪是後來的最後一位士師撒母耳的家鄉,他先膏立掃羅,但其後放棄掃羅,膏立大衛.那人最後選擇被棄的君王掃羅的家鄉基比亞.

19:16晚上,有一個老年人從田間做工回來.他原是以法蓮山地的人,住在基比亞；那地方的人卻是便雅憫人.………

19:22他們心裡正歡暢的時候,城中的匪徒圍住房子,連連叩門,對房主老人說:「你把那進你家的人帶出來,我們要與他交合.」

19:23那房主出來對他們說:「弟兄們哪,不要這樣作惡；這人既然進了我的家,你們就不要行這醜事.

19:24我有個女兒,還是處女,並有這人的妾,我將她們領出來任憑你們玷辱她們,只是向這人不可行這樣的醜事.」

19:25那些人卻不聽從他的話.那人就把他的妾拉出去交給他們,他們便與她交合,終夜凌辱她,直到天色快亮才放她去.

19:26天快亮的時候,婦人回到她主人住宿的房門前,就仆倒在地,直到天亮.

這裡提及基比亞,那裡住著便雅憫人.剛好,掃羅正是住在基比亞的便雅憫人.在這裡發生的事,讓我們重溫創世記裡最淒涼的事.創世記19章,記載了所多瑪城的故事.有兩個天使,晚上來到所多瑪,被羅得接待.所多瑪城的人,連老帶少,重重圍著羅得的房子.他們要羅得交出那兩人,任他們魚肉.羅得打算交出兩個女兒.在千鈞一髮間,天使救了他們.之後,便是天火焚城的故事了.除了諾亞方舟講的洪水毀滅大地的故事外,天火毀滅所多瑪城是神對犯罪的世人最嚴厲的審判了.

我們以為,在神子民的歷史裡,不可能發生像所多瑪的邪惡事件.但我們的想法是錯的,就是在神子民的地方,也會發生天下間最駭人聽聞的惡事.那些便雅憫人,終夜凌辱一個弱質女流,將她凌辱至死.這是一個什麼樣的世界!

來到這裡,我們看看神的救贖故事由人來主導的話,發展到什麼地步.神揀選亞伯拉罕,讓他的後裔成為強大民族,為要將光明帶進黑暗的人間.神帶領他們脫離奴役的枷鎖,成為自由的人.神帶領他們克服困難,踏進被稱為「流奶與蜜」的豐盈土地.神希望他們不要學迦南人,不要陷進他們的黑暗,要點燃光明,要成為在墮落世界中的聖潔力量.但,結果如何?他們比迦南人更迦南人,比黑暗更黑暗,比墮落更墮落.他們將好好的一片樂土,變成人間地獄.上帝看著,我想,祂不只是默默流淚,而是號咷大哭,泣不成聲.人類啊人類,我們什麼時候,才能止住不當的慾望,放下操控弱者的權力,不再只留心自己的得益,而學懂十字架的捨己.

不要單單說那些神的子民,我們看看自己.在教會裡,不是安全得很的嗎?不是的.有個網站,叫＜離教者之家＞.我不是要大家做個「離教者」,我只是想大家聆聽一下離教者的故事.在這網站裡,很多「離教者」的見證.他們的故事,像那位被凌辱的妾的故事一樣,都是受害者的故事.有一個「離教者」的見證是這樣的:「當你不再去教會或者對信仰有了動搖,教友關心的只是你是否離教,而不是你這個人.似乎他們忘記了你是一個有著血肉之軀的人,有思想有感情的人,你只是他們為主做工的道具而已.………

最具有諷刺意味的是,一個偽善的基督徒,長相英俊,以信仰的名義接近敬虔的姊妹,玩弄姊妹的感情,導致該女孩精神失常,至今無法正常生活和工作.一次聚會碰到,我上前問他心裡是否有愧,之前這麼欺騙教會裡的姊妹能對得起上帝嗎?他承認自己過去是做錯了,但他說事情已經過去了,上帝已經原諒他了,現在他的事業和感情都很棒,他說自己已經過上了理想中的生活,他快要結婚了,他說真的很感謝上帝賜給他的一切.………」(\url{https://exchristian.hk/home/article/show/323})

上帝聽著,我想,祂依舊泣不成聲.

19:27早晨,她的主人起來開了房門,出去要行路,不料那婦人仆倒在房門前,兩手搭在門檻上；

19:28就對婦人說:「起來,我們走吧!」婦人卻不回答.那人便將她馱在驢上,起身回本處去了.

19:29到了家裡,用刀將妾的屍身切成十二塊,使人拿著傳送以色列的四境.

19:30凡看見的人都說:「從以色列人出埃及地,直到今日,這樣的事沒有行過,也沒有見過.現在應當思想,大家商議當怎樣辦理.」

這裡拋出一個問題:「大家商議當怎樣辦理?」

故事發展下去,是以色列聯軍去伐罪.結果是:幾乎將「便雅憫」支派滅絕.便雅憫支派的女人和孩童都死光,只剩下600個男人.忽然,以色列人又覺得十二個支派少了一個支派不太好,於是又想辦法找老婆給這600個便雅憫人.他們的方法是:誰無參加伐罪聯軍的,便將那城的人全殺掉,只留處女.於是,他們出發到「基列˙雅比」進行屠城,留下400處女.但,還欠200個女子.於是,他們強搶示羅的民女去湊數.其實,如此解決問題的方式,是再一次的胡亂屠殺和淫辱婦女,將之前的罪惡再倍大,更多人被殺,更多婦女被害.

這就是《士師記》給我們看的終極圖畫:不接待外人,淫亂,任意的殺戮,視生命如物,以更罪惡的手法處理罪惡.他們對付迦南人,從來沒有這般齊心,這般徹底,但殺害自己人,卻手起刀落.對不參與的,也絕不留情.我希望,在這面歷史的鏡前,我們看到自己的本相.

《士師記》的寫作,不單記述以色列人還沒有王之前的亂局,其在400多年後最後的編輯定稿,更是在面對被亞述及巴比倫滅國後的一片混亂時,反省造成這一片混亂的背後的原因.猶大的神學家將歷史反省後的神學思考擺在我們面前,看,是我們的手,將神所創造的美好世界作「反創造」的處理,將秩序變成「混亂」,將樂園變成地獄.看,就算將樂園交給我們,就算伊甸園的大門重新向我們開放,就算神應許的福樂向我們傾倒,我們會將之搞成什麼樣子?不要埋怨世界的黑暗,是我們將手中的火把熄滅.

無怪乎,保羅千叮萬囑,提醒我們信仰的意義:

羅12:2不要效法這個世界,只要心意更新而變化,叫你們察驗何為神的善良,純全,可喜悅的旨意.

林後5:17若有人在基督裡,他就是新造的人(新的創造),舊事已過,都變成新的了.

在絕望的光景裡,《士師記》留下希望.希望在哪裡?在彌賽亞式的君王的來臨.世上不同的人都會逐鹿王位.誰是天命所歸的人?經文嚴厲批評掃羅王及後來的北國以色列的不敬神與不愛人.《士師記》將希望指向像大衛那樣的王.大衛雖壞,但始終有一敬畏上帝的心.歷史一路發展下去,我們知道,大衛王朝最終都倒下.然後,我們問:在大衛的子孫裡,會否有一救世的王者來臨,來拯救這黑暗的蒼生?

這位王者的確來了,大衛的子孫耶穌來了,改變了這個世界.問題是:我們是活出新世界的和平與福樂,抑或重複舊世界的墮落與黑暗?

願榮耀歸上主,願人間享和平.



\newpage

\section{讓我們一同禱告}
\label{sec:609}
\textbf{高國雄牧師 - 讓我們一同禱告}
\newline
\newline
\hyperref[sec:608]{< < < PREV SERMON < < <}
~
\hyperref[sec:toc]{[返目錄]}
~
\hyperref[sec:610]{> > > NEXT SERMON > > >}


路加福音 11:1-13
\newline
\begin{longtable}{cl}
\hline
\hline
章節 & 經文 (和合本修訂版)\\
\hline
11:1 & \begin{tabularx}{0.7\textwidth}{X} 耶穌在一個地方禱告。禱告完了，有個門徒對他說：「主啊，求你教導我們禱告，像約翰教導他的門徒一樣。」 \end{tabularx} \\ \\ \relax
11:2 & \begin{tabularx}{0.7\textwidth}{X} 耶穌對他們說：「你們禱告的時候，要說：『父啊，願人都尊你的名為聖；願你的國降臨； \end{tabularx} \\ \\ \relax
11:3 & \begin{tabularx}{0.7\textwidth}{X} 我們日用的飲食，天天賜給我們。 \end{tabularx} \\ \\ \relax
11:4 & \begin{tabularx}{0.7\textwidth}{X} 赦免我們的罪，因為我們也赦免凡虧欠我們的人。不叫我們陷入試探。』」 \end{tabularx} \\ \\ \relax
11:5 & \begin{tabularx}{0.7\textwidth}{X} 耶穌又對他們說：「你們中間誰有一個朋友半夜到他那裡去，對他說：『朋友！請借給我三個餅； \end{tabularx} \\ \\ \relax
11:6 & \begin{tabularx}{0.7\textwidth}{X} 因為我有一個朋友旅途中來到我這裡，我沒有東西招待他。』 \end{tabularx} \\ \\ \relax
11:7 & \begin{tabularx}{0.7\textwidth}{X} 那人在裡面回答：『不要打擾我，門已經關了，孩子們也同我在床上了，我不能起來給你。』 \end{tabularx} \\ \\ \relax
11:8 & \begin{tabularx}{0.7\textwidth}{X} 我告訴你們，雖不因他是朋友起來給他，也會因他不顧面子地直求，起來照他所需要的給他。 \end{tabularx} \\ \\ \relax
11:9 & \begin{tabularx}{0.7\textwidth}{X} 我又告訴你們，祈求，就給你們；尋找，就找到；叩門，就給你們開門。 \end{tabularx} \\ \\ \relax
11:10 & \begin{tabularx}{0.7\textwidth}{X} 因為凡祈求的，就得著；尋找的，就找到；叩門的，就給他開門。 \end{tabularx} \\ \\ \relax
11:11 & \begin{tabularx}{0.7\textwidth}{X} 你們中間作父親的，誰有兒子求魚，反拿蛇當魚給他呢？ \end{tabularx} \\ \\ \relax
11:12 & \begin{tabularx}{0.7\textwidth}{X} 求雞蛋，反給他蠍子呢？ \end{tabularx} \\ \\ \relax
11:13 & \begin{tabularx}{0.7\textwidth}{X} 你們雖然不好，尚且知道拿好東西給兒女，何況天父，他豈不更要把聖靈賜給求他的人嗎？」 \end{tabularx} \\ \\
[1ex]
\hline
\hline
\end{longtable}
講題:讓我們一同禱告

經課:路加福音11:1-13

講員:高國雄牧師

日期:2022年7月24日

1耶穌在一個地方禱告；禱告完了,有個門徒對他說:「求主教導我們禱告,像約翰教導他的門徒.」

2耶穌說:「你們禱告的時候,要說:我們在天上的父:願人都尊你的名為聖.願你的國降臨；願你的旨意行在地上,如同行在天上.3我們日用的飲食,天天賜給我們.4赦免我們的罪,因為我們也赦免凡虧欠我們的人.不叫我們遇見試探；救我們脫離凶惡.」

5耶穌又說:「你們中間誰有一個朋友半夜到他那裡去,說:『朋友!請借給我三個餅；

6因為我有一個朋友行路,來到我這裡,我沒有甚麼給他擺上.』

7那人在裡面回答說:『不要攪擾我,門已經關閉,孩子們也同我在床上了,我不能起來給你.』

8我告訴你們,雖不因他是朋友起來給他,但因他情詞迫切地直求,就必起來照他所需用的給他.

9我又告訴你們,你們祈求,就給你們；尋找,就尋見；叩門,就給你們開門.

10因為,凡祈求的,就得著；尋找的,就尋見；叩門的,就給他開門.

11你們中間作父親的,誰有兒子求餅,反給他石頭呢?求魚,反拿蛇當魚給他呢?

12求雞蛋,反給他蠍子呢?

13你們雖然不好,尚且知道拿好東西給兒女；何況天父,豈不更將聖靈給求他的人嗎?」

禱告對你來說是什麼呢?

大部份字典類的解釋,大概就是:「祈禱就是向神說話,獻上感恩及祈求幫助.」

在神學上,在教會歷史中對禱告可能有很多不同的定義,不同的演繹,先聖先賢對禱告也會有很個人不同的體會及理解,最重要的是:禱告對你來說是什麼呢?

對我來說,每次聚會時(特別是在聚餐前),很多時候被邀請帶領禱告,除了相信是一份尊重之外,也可能是因為對禱告有不同的理解而產生,甚至連未信主的朋友也會覺得牧師帶領禱告比較合適(更好).

耶穌所教導的禱告非常重要,亦是我們最熟悉的,就是在每個主日一同念誦的,很多教會學校,從少就教導同學背誦主禱文,多年前曾經聽過有學校舉行「鬥快背誦主禱文比賽」,亦親身體會過小學生們非常流暢地背誦“救我們脫離空殼”,當然這內容也非常值得反省,但這錯誤更值得我們反思,應該如何對待“禱文”

聖經中曾經記述不少不同人對上帝禱告的禱文,歷世歷代也有很多非常精彩及值得我們細心的學習的禱文.

聖法蘭西斯(St. Francis of Assisi)的禱文

主啊!使我作您和平之子

在憎恨之處~~播下您的愛

在傷痕之處~~播下您寬恕

在懷疑之處~~播下信心

在絕望之處~~播下您的盼望

在幽暗之處~~播下您的光明

在憂愁之處~~播下您的歡愉

在赦免時~~我們便蒙赦免

在捨去時~~我們便有所得

迎接死亡時~~我們便進入永生

主啊~~使我少為自己求

少求受安慰~~但求安慰人

少求被瞭解~~但求瞭解人

少求愛~~但求全心付出愛

尼布爾的靈靜禱文

上帝啊!懇求你賜我恩典,安詳接受不能改的事件；鼓我義念,勇敢改變可以改的事件；贈我慧劍,能夠分辨接受還是改變.

讓我每天實在過活,時刻安然享受；讓我在艱困的路上,內心不失平安；讓我師法祂,接納這罪孽的人世間；讓我信任你,萬物萬事會變成完善.

但願我順服你的旨意,好叫我此生悠然快樂,來世更與你同享極樂.阿們.

除了歷代精彩禱文之外,市面上有很多教人禱告的書籍及網頁,但請先問,為何要學習?

因為不懂其方式?怕錯誤?怕上帝不聆聽?

有人常常會將禱告視為一種「魔法」.

甚至是通靈的咒語,通靈是沒有錯的(當然不是民間宗教那一種),因為禱告就是透過主賜我們內在的靈,與上帝的靈溝通,所以禱告的對象不是向空間說話；

而其魔法就是因著主的應許

馬太福音18:19-20我又告訴你們,若是你們中間有兩個人在地上同心合意地求甚麼事,我在天上的父必為他們成全.因為無論在哪裡,有兩三個人奉我的名聚會,那裡就有我在他們中間.

有些人認為,如果我們禱告時沒有說對話,擺對姿勢或特別的順序,神就不會垂聽和應允我們的禱告.這種看法是完全不符合聖經的.

今天的經課,路加福音11章就是主耶穌教導門徒如何禱告的主禱文,首先必須要認清,主禱文不是“主自己的禱文”

V2耶穌說:「你們禱告的時候,要說…

所以此主禱文是主教導我們,又特別為我們所設計,賜給我們的禮物

為何耶穌要教導門徒如何禱告呢?

我們又如何學習禱告的呢?

以賽亞書1:15 你們舉手禱告,我必遮眼不看,就是你們多多地祈禱我也不聽,你們的手都滿了殺人的血.

約翰一書5:14-15我們若照他的旨意求甚麼,他就聽我們,這是我們向他所存坦然無懼的心.

既然知道他聽我們一切所求的,就知道我們所求於他的,無不得著.

門徒請求耶穌教他們如何禱告,因為:施洗約翰有教導他的門徒如何禱告；耶穌常常禱告；門徒覺得禱告重要,亦想有所進步.

我們有這羨慕,追求的心嗎?還是禱告只期望滿足自己的需要?

接著耶穌就以我們所熟悉的「主禱文」教導門徒,主禱文內容非常豐富,過去我在教會曾經以1個學期(12堂)分享當中的重點,相信今天沒有辦法在這短短的十多分鐘內詳細講解,只想特別提出當中幾點,作為大家的提醒及鼓勵.

我們在天上的父:願人都尊你的名為聖.

以上帝為開始點,不單是稱呼,而是歸榮,更是整個禱告的核心.

願你的國降臨；願你的旨意行在地上,如同行在天上.

先求上帝的國:

馬太福音6:33你們要先求他的國和他的義,這些東西都要加給你們了.

願你的旨意行在地上,如同行在天上.

祈求上帝旨意成就(如何成就?)

我們日用的飲食,天天賜給我們.

為生命所需祈禱,認定上帝的應許及恩典

赦免我們的罪,因為我們也赦免凡虧欠我們的人.

求赦罪(因為只有手潔心清的人才可以到上帝面前),並承諾要以同樣的愛愛人

不叫我們遇見試探；救我們脫離凶惡.

求拯救,以活在主內

有了「標準」,可跟從的學習內容後,耶穌以比喻,補充當中最重要的部分:

首先:作為遊牧民族的後代,接待是猶太人生活中其中最重要的功課,

早期教會同樣看重:

希伯來書13:2不可忘記用愛心接待客旅；因為曾有接待客旅的,不知不覺就接待了天使.

路加用連續幾段經文教導真正接待的價值及意義:

特別由9:43-45耶穌再次預言自己將會被殺害之後；

9:46-48對門徒爭大的問題的回應:凡為我名接待這小孩子的,就是接待我；

9:51-56對不接待耶穌及門徒的撒瑪利亞村莊:他的門徒雅各,約翰看見了,就說:「主啊,你要我們吩咐火從天上降下來燒滅他們,像以利亞所做的嗎?」

耶穌轉身責備兩個門徒,說:「你們的心如何,你們並不知道.

人子來不是要滅人的性命,是要救人的性命.」說著就往別的村莊去了

10:1-12差派70個門徒去被接待(學習被接待):你們要住在那家,吃喝他們所供給的,因為工人得工價是應當的；不要從這家搬到那家.

8無論進哪一城,人若接待你們,給你們擺上甚麼,你們就吃甚麼.

10:25-37你想,這三個人哪一個是落在強盜手中的鄰舍呢?」

他說:「是憐憫他的.」耶穌說:「你去照樣行吧.」:好撒瑪利亞人的比喻,「愛鄰舍」的真正意義,就是要學習被接待,我就是那個遍體麟傷需要鄰舍的人.

10:38-42馬大,馬利亞接待耶穌:馬利亞已經選擇那上好的福分,是不能奪去的.

然後到今天的經課的內容

朋友的請求,情詞迫切地直求；人必然應該盡能力回應其需要

11:9-10

我又告訴你們,你們祈求,就給你們；尋找,就尋見；叩門,就給你們開門.

因為,凡祈求的,就得著；尋找的,就尋見；叩門的,就給他開門.

上主對禱告的回應就是如此直接,但為什麼我們常常覺得未能經歷呢?

11:11-12

你們中間作父親的,誰有兒子求餅,反給他石頭呢?求魚,反拿蛇當魚給他呢?

求雞蛋,反給他蠍子呢?

耶穌以一般能經驗的父親與兒子的關係作為例子,上主回應的直接,就如我們也必會盡力直接回應所愛的人的祈求；但問題會否是出於兒子的祈求呢?我們作兒子的,是在求「餅」,「魚」及「雞蛋」?還是只是每天不斷的求「石頭」,「蛇」及「蠍子」呢?

放心,上主非常清楚什麼是對我們最好的…

11:13你們雖然不好,尚且知道拿好東西給兒女；何況天父,豈不更將聖靈給求他的人嗎?

要深信上主給我們必定是最好的,儘管…

以作者路加醫生的角度,路加福音可理解為使徒行傳前傳,而使徒行傳其實就是聖靈行傳

耶穌在此清晰表達,真正最好的,我們要好好祈求的內容,就是聖靈,就是主的臨在,與我們同在的聖靈,就是禱告追求的核心…



\newpage

\section{煥然一新}
\label{sec:610}
\textbf{王玉教授 - 煥然一新}
\newline
\newline
\hyperref[sec:609]{< < < PREV SERMON < < <}
~
\hyperref[sec:toc]{[返目錄]}
~
\hyperref[sec:611]{> > > NEXT SERMON > > >}


歌羅西書 3:1-11
\newline
\begin{longtable}{cl}
\hline
\hline
章節 & 經文 (和合本修訂版)\\
\hline
3:1 & \begin{tabularx}{0.7\textwidth}{X} 所以，既然你們已經與基督一同復活，就當求上面的事；那裡有基督，坐在神的右邊。 \end{tabularx} \\ \\ \relax
3:2 & \begin{tabularx}{0.7\textwidth}{X} 你們要思考上面的事，不要思考地上的事。 \end{tabularx} \\ \\ \relax
3:3 & \begin{tabularx}{0.7\textwidth}{X} 因為你們已經死了，你們的生命與基督一同藏在神裡面。 \end{tabularx} \\ \\ \relax
3:4 & \begin{tabularx}{0.7\textwidth}{X} 基督是你們的生命，他顯現的時候，你們也要與他一同在榮耀裡顯現。 \end{tabularx} \\ \\ \relax
3:5 & \begin{tabularx}{0.7\textwidth}{X} 所以，要治死你們在地上的肢體；就如淫亂、污穢、邪情、惡慾和貪婪—貪婪就是拜偶像。 \end{tabularx} \\ \\ \relax
3:6 & \begin{tabularx}{0.7\textwidth}{X} 因這些事，神的憤怒必臨到那些悖逆的人。 \end{tabularx} \\ \\ \relax
3:7 & \begin{tabularx}{0.7\textwidth}{X} 當你們在這些事中活著的時候，你們的行為也曾是這樣的。 \end{tabularx} \\ \\ \relax
3:8 & \begin{tabularx}{0.7\textwidth}{X} 但現在你們要棄絕這一切的事，就是惱恨、憤怒、惡毒、毀謗和口中污穢的言語。 \end{tabularx} \\ \\ \relax
3:9 & \begin{tabularx}{0.7\textwidth}{X} 不要彼此說謊，因為你們已經脫去舊人和舊人的行為， \end{tabularx} \\ \\ \relax
3:10 & \begin{tabularx}{0.7\textwidth}{X} 穿上了新人，這新人照著造他的主的形像在知識上不斷地更新。 \end{tabularx} \\ \\ \relax
3:11 & \begin{tabularx}{0.7\textwidth}{X} 在這事上並不分希臘人和猶太人，受割禮的和未受割禮的，未開化的人、西古提人、為奴的、自主的；惟獨基督是一切，又在一切之內。 \end{tabularx} \\ \\
[1ex]
\hline
\hline
\end{longtable}
講題:煥然一新

經課:歌羅西書3:1-11

講員:王玉教授

日期:2022年7月31日

I.引言

相信在座應該有不少人都聽過一句頗流行的勵志短句「態度決定高度」(Attitude determines altitude).這句短句可理解成「決定你人生高度的,不是你的才能,而是你的態度.」(It’s your attitude, not your aptitude that will determine your altitude.)[1] 的縮短版.這句說話出自已故美國作家,勵志演說家,兼銷售員Zig Ziglar.他用了很多的勵志短句來激勵銷售員和企業家去追求一個成功的人生,鼓倡成功的關鍵其實很簡單,就是要有積極的心態.無可口非這句說話有它的道理,但是它強調的是個人的成就和得失.

相反地,今天的新約書信歌羅西書3:1–11及其後的12至15節所提倡的是一種很不一樣的人生態度,經文強調的是「個人外在行為及內心素質的更新決定群體關係的轉化」,這種的轉化包括了個人生命及群體關係轉化的層面.經文重視的並不是一個以個人的成就與得失為目的,而是如何由個人內在素質的更新達致一個更包容,更有愛與平等的群體.講題「煥然一新」所指的正這種由個人轉化至群體轉化的願景.全段經文的精髓位於5a, 9b至10節的一個「換衫」的比喻:「所以,要治死你們在地上的肢體…,因為你們已經脫去舊人和舊人的行為,穿上了新人,這新人照著造他的主的形像在知識上不斷地更新.」

II.「穿上新人」的理據:基督復活與個人生命的轉化(3:1–4)

首先,1至4節提供了「穿上新人」的理據,因為我們已經與勝過世上一切權勢並且與上帝得享榮耀的基督同死同復活,我們的生命也「與基督一同藏在上帝裡面」,意思是我們的生命一同在上帝的蔭蔽/保護 (κέκρυπται; be concealed)下.將來當基督的生命呈現於 (φανερωθῇ; be revealed)我們的生命中,我們也將與他一同得榮耀.這意味著基督的死和復活是必然地帶來基督追隨者的生命轉化,讓我們更像基督,讓我們的生命更能活出基督的生命.基督的死和復活是具有一種強烈的道德感染力,足以轉化其跟隨者的生命,令人的內心得到更新,啟發人追求更崇高的道德價值,更美善,更公義,更和諧的社會.

耶穌基督本是一個義人,卻無辜受迫害,並甘願為愛,擺上自己的生命,死在十架.義人受苦之所以能夠帶來崇高的道德力量是因為它能夠產生一種替代的經驗 (vicarious experience).

所謂替代的經驗是指一個人代替或代表你承受了一些後果,處境,感受,或痛苦,而這一切是你感到你本應承擔或分擔的.

一個義人無辜受害激發歷代千千萬萬的基督徒繼續跟隨基督,追求愛,和平,平等,及公義.這就是義人受苦所帶來的道德力量,能夠感染及轉化生命.這種「生命轉化」並不是在我們死後發生的,而是從你願意接受基督,將你的生命交給基督,依靠他,並跟隨他的引導過生活的那一刻開始.從那一刻開始,你已經「復活」了,得著「永生」了.因此,復活的果效不是來世才能經驗到的,而是今世(此時此刻)已經呈現於信徒生命的轉化.

哥林多後書5:14-15說:「原來基督的愛激勵我們；因我們這樣斷定,一人既替眾人死,眾人就都死了；並且他替眾人死,是叫那些活著的人不再為自己活,乃為替他們死而復活的主活.」一個曾經被基督的愛激勵過的人,被基督的死,受苦,復活感動過的人,是會被他的愛感染和轉化的,人生因此變得不一樣,因為這人再不會為自己而活,而是為「替他死而復活的主活」.歌羅西書3:11指出,這種生命的轉化是普世性,無論任何種族,任何傳統,任何社會階層,基督的生命都能夠觸及,都能夠轉化.

III.脫去「舊人與舊人的行為」:行為上的轉化(3:5–9)

5至15節指出基督徒生命的轉化包含了三個層面.

5至9節指出的第一層面是行為上的轉化.一開始,經文提醒基督徒要去「思考上面的事,不要思考地上的事」,又或者說「去思考一些上帝所重視,有永恆價值的事,而不是一些人所重視,短暫無益,甚至是帶有傷害性的事.」這是對我們一個很大的提醒.弟兄姊妹,到底你的心思意念是充斥著「地上的事」,還是「上面的事」呢?在5節,作者舉出一些關注「地上的事」的意念的例子—「淫亂,污穢,邪情,惡慾和貪婪—貪婪就是拜偶像.」在8至9節,作者繼續舉出其他對人,對群體帶有破壞性,傷害性的脾氣和行為—「惱恨,憤怒,惡毒,毀謗和口中污穢的言語」,以及「說謊」.這些例子反映了人的罪性,人為了個人的慾望和利益,主動去做一些傷害其他人的事.

弟兄姊妹,到底什麼是「罪」?誰是「罪人」?罪是否僅僅指一些主動傷害的行為?或僅僅是個人的行為問題呢?

希臘東正教,英國區的大主教尼基塔斯·盧利亞斯 (Nikitas Lulias) 在4月份所舉行的「Faith in Ukraine」(信念在烏克蘭)的跨宗教聚會中曾經分享過一個故事.[2]

他說:「幾年前,一位老婦人去找她的神父懺悔,對他說:『神父,我這輩子沒有偷竊,我沒有拿走任何不屬於我的東西.我從來沒有欺騙過我的丈夫,未有過外遇.我從來沒有想要任何不屬於我的東西.我沒有做那些我不應該做的事情.』理所當然地,她坐下來等待神父告訴她,『做得好,你這善良又忠心的僕人.』

但神父一臉悲傷地看著她,說:『我很遺憾地告訴你,你可能會直接下地獄,因為生命不是在乎你不應該做或不做的事情.生命是在乎你做什麼,與及時做好這一刻.親愛的姐妹兄弟,生命是關於我們必須做的事情.因為冷漠可能是所有罪惡中最大的.』」

或許你會覺得這個故事比較極端,但是類似的故事聖經中也有.路加福音18:18-30(太19:16-30;

可10:17-31)講到「一個官問耶穌說:『善良的老師,我該做甚麼事才能承受永生?』耶穌對他說『……

誡命你是知道的:「不可姦淫；不可殺人；不可偷盜；不可作假見證；當孝敬父母.」』那人說:『這一切我從小都遵守了.』最後耶穌吩咐他變賣一切所有的,分給窮人,來跟從我.他聽見這些話,就很憂愁,因為他很富有.另外,路加福音16章講述了「財主和拉撒路」的故事,財主天天經過在他門口討飯的拉撒路,但是他漠視他的需要,漠視他的生命危在旦夕,沒有給他充飢的食物.最後兩個都死了,財主在陰間受苦,而拉撒路就在亞伯拉罕的懷裡得著安慰.接著,耶穌說:「有錢財的人進神的國是何等的難哪!駱駝穿過針眼比財主進神的國還容易呢!」(vv.24-25)

我想透過大主教的故事及聖經的故事指出一個對「罪」的非常普遍的理解,就是「罪為行為上的偏差」.在福音書中,「罪人」一詞是有特定的意思,指一群生活方式和行為偏離當時主導社會制度和其主導道德價值,期望的人.這包括了不應做的反去做 (sin of commission) 及應做的不去做(sin of omission).問耶穌如何得永生的那個官只能停留在一個律法主義式的信仰,關注不去做不應做的.而路加福音16章中的財主則未能去做應做,去關心一個生命危在旦夕的乞丐.

罪固然之可以是行為上的偏差,但是罪也可以形容一種人存在的境況或社會結構性邪惡產生的倫理道德問題.人性軟弱,恐懼,內疚(不論是主觀感覺或客觀標準衡量),又或是歌羅西書3:8所提及的「惱恨,憤怒,惡毒」等等；往往令人與人之間的關係疏離 (estranged),異化 (alienated).父親在一時「火遮眼」下,掌摑了兒子,足以令父子關係從此破裂.人因為做了內疚的事情,無法原諒自己,無法抬起頭做人,影響了與家人和朋友的關係,與人疏遠.因為他無法原諒自己,他亦覺得別人無法,也不應原諒自己.

社會結構問題邪惡產生的倫理道德問題,更使「罪」變得避無可避,即令人無法選擇「應做的事」或不得不選擇做「不應做的事」.人基於生活處境或制度問題,在工作上,社會中,家庭中都機會做一些身不由己,逼於無奈的事情.例如:父母因為要養家活口,一日工作十多小時,逼於無奈疏忽照顧子女.法官要依法判刑,但無權去質疑法例的正當性.就算存疑,職責上也得去執行.紀律部隊人員或軍人要執行上級指示,無論他是否同意或明白所執行的職責.

列王記下5章講述了亞蘭王的元帥乃縵染上了痲瘋,去求先知以利沙醫治,他按先知要求,在約旦河中沐浴七次,痲瘋就痊癒了.他最後歸依耶和華,但因為工作關係,無法不對亞蘭神明臨門叩拜.「乃縵向以利沙請求:『惟有一件事,願耶和華饒恕你僕人:我主人進臨門廟在那裡叩拜的時候,他總是扶著我的手,所以我也在臨門廟叩拜.我在臨門廟叩拜的這事,願耶和華饒恕你僕人.』以利沙對他說:『你平安地回去吧!』」(vv.18-19)

兩害相權原則 (Lesser of two evils principle)指出,有時候人被迫在兩個道德都不對的選擇上作決定,你可以選的不是錯或對,你只能選擇「無咁錯」.例如:說謊言和拯救生命之間.一個人應該說謊來挽救一條生命嗎?納稅德國很多波蘭人就選擇了說謊匿藏猶太人.生活在當下的香港,無論是疫情,還是社會處境,我相信我們經常都有左右為難,不能兩全其美,迫於選擇「無咁錯」.

最後一點,社會結構問題邪惡往往令人在不知情下參與了剝削性的事情.例如我們所穿的棉質衣服或所吃的肉類,可能是透過剝削工人權益,動物權益,甚至是強迫勞動而得出的,而我們都在不知情下使用了.2019年底,我去了巴峇島參觀了「貓糞咖啡」(civet coffee) 的制作場.「貓糞咖啡」,由印尼棕櫚果子狸排泄的咖啡豆製成,譽為世界上最貴的咖啡.導賞員強調制作過程沒有涉及虐待動物,我當時還無知地相信她.後來作一些調查,原來果子狸是被困在狹窄的籠子裡,在一個極不人道的環境下生活,強迫吃大量不能消化的咖啡豆,有些甚至嚴重受傷.我在不知情下或太天真下參與了虐待野生動物.

上述的例子告訴我們「罪」並不僅是一種個人「行為上的偏差」,更是一種存在的境況或社會結構性邪惡,有時候令我們避無可避.只要你和我生存在這個世界裡,我們永遠離不開罪的境況.使徒保羅在羅馬書3:10精闢地指出了要害:「沒有義人,連一個也沒有.」他又在羅馬書7:19-24說:「我所願意的善,我不去做；我所不願意的惡,我反而去做.如果我去做我不願意做的,就不是我做的,而是住在我裡面的罪做的.我覺得有個律,就是我願意行善的時候,就有惡纏著我.因為,按著我裡面的人,我喜歡神的律,但我看出肢體中另有個律和我內心的律交戰,把我擄去,使我附從那肢體中罪的律.我真苦啊!誰能救我脫離這必死的身體呢?」使徒保羅分享他的心路歷程,他在罪的境況掙扎,令他感到無助,沮喪,但他仍然「感謝神」,因為他深信他能夠倚靠主耶穌基督去面對罪的困境(v.25).

IV.「穿上新人」由內在素質的更新開始 (3:12)

歌羅西書3:12指出了基督徒生命的轉化的第二層面,就是內在素質的更新:「所以,你們既是上帝的選民,聖潔,蒙愛的人,要穿上憐憫,恩慈,謙虛,溫柔和忍耐.」原來基督徒生命的轉化不僅是外在行為的轉化,只「脫去舊人和舊人的行為」,更要進一步「穿上新人」,去更新內心.換言之,單是避免做「不應做的事」是不足的,我們還去有內在素質的更新,才能做「應做的事」.箴言4:23提醒我們:「你要保守你心,勝過保守一切,因為一生的果效是由心發出.」(和合本)

人去做「應做的事」必須發自內心,出自內心的更新﹐這種內心的更新並不是一次性的,而是多次的,就像電腦軟件的更新一樣,是需要不繼進行的,才能培養出「憐憫,恩慈,謙虛,溫柔和忍耐」的心.

初中時期是我人生其中一段最動蕩不安的時期,中二後,我媽媽再婚,家姐,二哥和我隨同母親和繼父移民加拿大.移民的頭一年,家裡經常有衝突和摩擦,吵架變成了家常便飯,吵架聲浪之大,未入屋已聽到,每周總有幾回,而且涉及家裡所有成員.當時期家人的關係,家裡的氣氛和局面,以「惱恨,憤怒,惡毒,毀謗和口中污穢的言語」來形容,絕對不過分.每次聽到家中成員無辜受罵,我都會氣憤無比,加入參與罵戰.大家可以想像,情況不但沒有改善,而且還變本加厲.

感恩的是同年也是我決志信主的一年.決志信主之後的我,被聖靈提醒,每次家裡有罵戰,無論內心有多憤怒,裡面都有一個微少的聲音提醒不能加入罵戰.記得有一次家裡又有罵戰,而且越罵越厲害,越罵越大聲,我開始感到按捺不住,之後我一手拿起聖經和詩集,將自己反鎖在房間中,向上帝祈禱.我求主幫助我,因為我實在無法再忍受下去,我感到很無助,很孤單,無出路.祈禱後,我隨意翻開詩集,翻到詩歌《我心靈得安寧》,輕聲地唱起來.

有時享平安,如江河平又穩,有時悲傷來似浪滾.

不論何環境,主已教導我說:我心靈得安寧,得安寧.

我心靈得安寧,我心靈得安寧,得安寧.[3]

一邊唱,眼淚一邊流下來,房門外仍然聽到聲嘶力竭的罵戰,但是內心卻有一分出人意外的平靜,特別是唱到「基督深知我,無依靠之苦境,甘流血救贖我,賜安寧」那一句時.這次之後,家裡再有吵架,我彷彿有了特異功能,不再感到憤怒,而且我還能以平靜的心,溫柔的聲線,耐心地作「和事佬」.我不但見證了耶穌如何更新我的內心,往後日子裡,我也見證了耶穌如何更新了其他家人的內心(有的比我的改變更為奇妙),及最終家人的關係也得以修復.

在主裡,我們一同學習「穿上憐憫,恩慈,謙虛,溫柔和忍耐」的功課.

V.外在群體關係的轉化 (3:13–15)

我剛才分享的個人經驗正好能夠帶出基督徒生命的轉化的第三層面,就是「由個人外在行為及內心素質的更新轉化成群體關係的更新」.正如歌羅西書3:13–15所描述的情況一樣:「倘若這人與那人有嫌隙,總要彼此容忍,彼此饒恕；主怎樣饒恕了你們,你們也要怎樣饒恕人.除此以外,還要穿上愛心,因為愛是貫通全德的.你們要讓基督所賜的和平在你們心裡作主,也為此蒙召,歸為一體.你們還要存感謝的心.」這就是群體關係更新轉化後的一個情境.弟兄姊妹,你有沒有想過你個人的更新改變可以促成你身邊的人的更新改變呢?「罪」對人類最大的破壞,是關係性的破裂,令人與自己,人與人,人與群體,人與其他受造物,人與上帝之間的關係性的破裂.

耶穌的死與復活對我們有何意義呢?耶穌的死,讓人在上帝的眼中,一切所犯的都可以一筆勾銷.令生命可以Take

Two!這不代表人不用為所作的承受後果或懲罰；而是無論你的過去是怎樣,將來要面對的是什麼,你已經被上帝無條件接納了.而這種無條件接納,足以感化人對自己接納,學習被接納,及給予其他人接納,使人復和.保羅在哥林多後書5:17-18用了一個類以「換衫」的比喻,說:「若有人在基督裡,他就是新造的人,舊事已過,都變成新的了.一切都是出於上帝；他藉著基督使我們與他和好,又將勸人與他和好的使命賜給我們.」

要和好,我們必須學會寬恕.寬恕一個人其實是讓自己得著自由,擺脫仇恨,擺脫憤怒,擺脫自憐,重拾去愛的能力,活出真正的自己,自由地過最充實的生活.基督的死不但令我們的罪得赦,更使我們成就了人與上帝,人與自己,人與人,人與世界,群體與群體之間的復和的中介.

VI.

總結

從歌羅西書3:1–15我們看到基督徒生命轉化的三個不同層面:

(1) 個人外在行為上的轉化 ,即「脫去舊人和舊人的行為」；

(2) 個人內在素質的更新,即「穿上新人」,這兩者在基督徒生命轉化的過程是缺一不可的.有了個人外在與內在的轉化更新,才能夠達至第三個層面

(3) 群體關係的轉化更新.基督徒生命轉化是一個過程,有時候是突變,有時候是漸變；有時候改變是細微,有時候是明顯的.無論如何,我們總要更新轉化.讓我們明白我們個人的更新轉化是足以影響群體的更新轉化.所以,千其不要少看自己的.相信在座不少人都聽過「蝴蝶效應」(butterfly effect):一隻蝴蝶在南美洲拍翼,會引致北美洲產生一場風暴.同樣地,就算你力薄弱如蝴蝶,上帝都可以使用你產生蝴蝶效應!

[1] Zig Zilar, See You at the Top (Gretna: Pelican, 1986), 292.

[2] “Faith in Ukraine,” YouTube, April 12, 2022, \url{https://www.youtube.com/watch?v=FFimBYHxsVo} , 45:35–46:50.

[3] 生命聖詩342首；詞:Horatio Gates Spafford (1873)；曲:Philip Paul Bliss (1876).



\newpage

\section{攝禮歸仁}
\label{sec:611}
\textbf{林豪恩先 - 攝禮歸仁}
\newline
\newline
\hyperref[sec:610]{< < < PREV SERMON < < <}
~
\hyperref[sec:toc]{[返目錄]}
~
\hyperref[sec:612]{> > > NEXT SERMON > > >}


以賽亞書 1:10-20
\newline
\begin{longtable}{cl}
\hline
\hline
章節 & 經文 (和合本修訂版)\\
\hline
1:10 & \begin{tabularx}{0.7\textwidth}{X} 所多瑪的官長啊，你們要聽耶和華的言語！蛾摩拉的百姓啊，要側耳聽我們神的教誨！ \end{tabularx} \\ \\ \relax
1:11 & \begin{tabularx}{0.7\textwidth}{X} 耶和華說：「你們許多的祭物於我何益呢？公綿羊的燔祭和肥畜的油脂，我已經膩煩了；公牛、羔羊、公山羊的血，我都不喜悅。 \end{tabularx} \\ \\ \relax
1:12 & \begin{tabularx}{0.7\textwidth}{X} 「你們來朝見我，誰向你們的手要求這些，使你們踐踏我的院宇呢？ \end{tabularx} \\ \\ \relax
1:13 & \begin{tabularx}{0.7\textwidth}{X} 不要再獻無謂的供物了，香是我所憎惡的。我不能容忍行惡又守嚴肅會：初一、安息日和召集的大會。 \end{tabularx} \\ \\ \relax
1:14 & \begin{tabularx}{0.7\textwidth}{X} 你們的初一和節期，我心裡恨惡，它們成了我的重擔，擔當這些，令我厭煩。 \end{tabularx} \\ \\ \relax
1:15 & \begin{tabularx}{0.7\textwidth}{X} 你們舉手禱告，我必遮眼不看，就算你們多多祈禱，我也不聽；你們的手沾滿了血。 \end{tabularx} \\ \\ \relax
1:16 & \begin{tabularx}{0.7\textwidth}{X} 你們要洗滌、自潔，從我眼前除掉惡行；要停止作惡， \end{tabularx} \\ \\ \relax
1:17 & \begin{tabularx}{0.7\textwidth}{X} 學習行善，尋求公平，幫助受欺壓的，替孤兒伸冤，為寡婦辯護。」 \end{tabularx} \\ \\ \relax
1:18 & \begin{tabularx}{0.7\textwidth}{X} 耶和華說：「來吧，我們彼此辯論。你們的罪雖像硃紅，必變成雪白；雖紅如丹顏，必白如羊毛。 \end{tabularx} \\ \\ \relax
1:19 & \begin{tabularx}{0.7\textwidth}{X} 你們若甘心聽從，必吃地上的美物； \end{tabularx} \\ \\ \relax
1:20 & \begin{tabularx}{0.7\textwidth}{X} 若不聽從，反倒悖逆，必被刀劍吞滅；這是耶和華親口說的。」 \end{tabularx} \\ \\
[1ex]
\hline
\hline
\end{longtable}
生

講題:攝禮歸仁

經課:以賽亞書1章10-20節

講員:林豪恩先生

日期:2022年8月7日

一．克己復禮為仁

各位弟兄姊妹,當你讀舊約聖經時,有沒有看到不少內容都和宗教禮儀有關,例如敬拜耶和華的帳幕和約櫃要如何造,祭司的服飾如何,燔祭,素祭,平安祭,贖罪祭,贖愆罪要如何獻,安息日要如何過,還有仔細的祭司工作指引,以及長長的祈禱文等,你知道為甚麼如此多禮儀及如此仔細嗎?就算不知道原因,你也會想像得到禮儀對以色列人敬拜耶和華的重要性.奇怪的是,到了主前第八世紀,先知們,例如阿摩司,以賽亞和彌迦等,卻對宗教禮儀卻另有觀點,他們對當時人們的宗教禮儀行為予以否定,甚至嚴厲批評.例如剛才誦讀的以賽亞書就是其中一段.這底是怎樣一回事呢?讓我們再細閱經文一齊思考.

11耶和華說:

你們所獻的許多祭物與我何益呢?

公綿羊的燔祭和肥畜的脂油,

我已經夠了；

公牛的血,羊羔的血,公山羊的血,

我都不喜悅.

12你們來朝見我,

誰向你們討這些,

使你們踐踏我的院宇呢?

13你們不要再獻虛浮的供物.

香品是我所憎惡的；

月朔和安息日,並宣召的大會,

也是我所憎惡的；

作罪孽,又守嚴肅會,

我也不能容忍.

14你們的月朔和節期,我心裡恨惡,

我都以為麻煩；

我擔當,便不耐煩.

15你們舉手禱告,我必遮眼不看；

就是你們多多地祈禱,我也不聽.

如果我們是當時的以色列人,聽到先知如此嚴厲的批評,很可能會感到委屈,甚至憤怒.追本尋源,是誰要我們獻祭的?是誰要我們守安息日的?是誰要我們守節期的?是誰要我們祈禱的?這一切,豈不是耶和華吩咐的嗎?豈不是律法要求的嗎?當我們嚴格遵守所有宗教禮儀的時候,你竟然嚴厲地批評我們,說甚麼耶和華「不喜悅所」燔祭(11),「恨惡」供物香品(12),「不能容忍」嚴肅會(13),對節期「都以為麻煩」(14),對禱告「遮眼不看,也不聽」(15)!我們不做,就等於不遵守律法,我們做了,又被批評,你叫我們如何是好呢?

各位弟兄姊妹,我們可以如何理解這種關於禮拜的掙扎呢?古代以色列是一個非常重視禮儀的民族,剛好在古代的東方,也有一個非常重視禮儀的民族,雖然兩者的禮儀範疇有所不同,例如前者重視的是宗教禮儀,後者重視的是人倫禮儀,但兩者在禮儀方面的經驗似乎有些可以互相借鏡的地方.今天,讓我們大膽嘗試以這個東方古老民族在禮儀方面的經驗來閱讀以賽亞先知留下的這段說話,或許能夠啟發我們理解古代以色列人在禮儀的掙扎.在座之中很多學者和專家對這個題目有更廣博和精深的認識,請恕小弟班門弄斧,還望賜教指正.

這個東方的古老民族有「禮儀三百,威儀三千」的說法,因此被稱為「禮儀之邦」.或許有人認為禮儀只是表面行為,但你聽過先賢的心聲嗎?他們認為「克己復禮為仁,一日克己復禮,天下歸仁焉.」(論語顏淵)(語譯:克制自己的私慾,使言行舉止合乎禮節,這就是仁.哪一天人們能克己復禮,天下就能歸服仁德.)你有沒有想過,禮儀除了個人修養及社會穩定的效果之外,更能把人提升至崇敬的境界.他們又說:「禮生敬,敬生情,敬情生神聖」.借助這種對禮儀的理解,或許能夠幫助我們明白古代以色列的宗教禮儀的「初心」.「禮生敬,敬生情,敬情生神聖」,一切的禮儀,無論是獻祭,節期,集會,禱告等都希望能引發神聖之情,朝向神聖之主.試想想,我們進入教堂參加主日崇拜,崇拜的程序也是一種禮儀,或坐下或站立,或安靜或頌唱,或低頭或仰望,這些集體行動都是朝向「禮生敬,敬生情,敬情生神聖」的境界.

你可以想像沒有禮序的集體崇拜會是怎樣的嗎?大家記得使徒保羅在哥林多前書如何批評哥林多教會的聚會情況嗎?

「你們聚會的時候彼此分門別類...在你們中間不免有分門結黨的事...你們聚會的時候,算不得吃主的晚餐；因為吃的時候,各人先吃自己的飯,甚至這個飢餓,那個酒醉.」(11:18-21)

保羅對這混亂的狀況非常憤怒,指責哥林多教會說:「你們要吃喝,難道沒有家嗎?還是藐視神的教會,叫那沒有的羞愧呢?我向你們可怎麼說呢?可因此稱讚你們嗎?我不稱讚!(11:22)並且斬釘截鐵地說:「你們聚會不是受益,乃是招損.」(11:17)用四個字來形容哥林多教會的狀況,可以說是「禮崩樂壞」,即然禮己崩,何以生「敬」呢?沒有「敬」,何以生「神聖」呢?因此,禮儀之於宗教,並不是可有可無的.既然如此,為甚麼主前八世紀的先知又如此嚴厲地批評時人的禮儀行為呢?

二．吃人的禮教

以賽亞書1:15節指出問題所在:「你們的手都滿了殺人的血.」原來當時那些嚴守禮儀的人士,被指責為滿手鮮血的人,怪不得他們所行的禮被嚴厲地批評.先知批評主前八世紀的以色列行禮如儀的人,會否令我們想起東方這個禮儀之邦也有一段嚴厲地批評禮教的時期呢?古代的耶和華敬拜者被稱為滿手鮮的人,而東方這個禮儀之邦傳統遺留下來的禮儀則被批評為「吃人的禮教」.想不到,千里之外,千年之隔,竟然點點相通.

魯迅在《狂人日記》裡說:「古來時常吃人,我也還記得,可是不甚清楚.我翻開歷史一查,這歷史沒有年代,歪歪斜斜的每頁上都寫著『仁義道德』幾個字.我橫豎睡不著,仔細看了半夜,才從字縫裡看出字來,滿本都寫著兩個字是『吃人』!」

後來,吳虞接續魯迅的口吻,發表〈吃人與禮教〉一文說:「到了如今,我們應該覺悟!我們不是為君主而生的!不是為聖賢而生的!也不是為綱常禮教而生的!什麼「文節公」呀,「忠烈公」呀,都是那些吃人的人設的圈套,來誑騙我們的!我們如今應該明白了!吃人的就是講禮教的!講禮教的就是吃人的呀!」

在如此對照之下,我們有一點頭緒為甚麼先知說耶和華「都不喜悅所」燔祭(11),「恨惡」供物香品(12),「不能容忍」嚴肅會(13),對節期「都以為麻煩」(14),對禱告「遮眼不看,也不聽」(15),因為禮儀行動已經遠離了它的初心,不但沒有朝向神聖,反而與神聖背道而馳.

3．攝禮歸義,攝義歸仁

耶和華上帝拒絕當時的以色列人那些宗教禮儀,祂想要的是甚麼呢?先知繼續說:

16你們要洗濯,自潔,

從我眼前除掉你們的惡行,

要止住作惡,

有人指出,禮教本來並不吃人.先賢的「初心」是「攝禮歸義」,禮依於義而立,義是禮的實質,禮是義的表現.「義」是甚麼呢?我們常常為「義」配上一個「正」字,成為「正義」這個詞,義即是「正當性」.以賽亞先知指責的耶和華敬拜者,已經失去了行為的正當性,他們的舉動被稱為「作惡」,他們的行為被稱為「惡行」.「攝禮歸義」要求的是行為的正當性.怎樣的行為算是正當的呢?以賽亞先知呼籲耶和華的敬拜者「17學習行善,尋求公平」,這就是一個參考的準則.

如果我們再問為甚麼要「行善,公平」呢?「行善,公平」是為了甚麼了?難道我們只是為了正當性而做正當的事嗎?原來「攝禮歸義」不是終點,先賢還提出「攝義歸仁」.「義」是有所指向的,義的關注並不是義本身,而是「仁」.「仁」又是甚麼呢?你有沒有聽過:「仁者愛人」(孟子．離婁下)這句說話?所以,「仁」和「愛」兩個字常常配在一起,成為「仁愛」一詞.大家有沒有發現,耶和華要求敬拜祂的人「學習行善,尋求公平」指向甚麼呢?原來也是指向仁愛.仁愛的具體行動是怎樣的呢?對甚麼人仁愛呢?為甚麼是他們呢?

解救受欺壓的；

給孤兒伸冤,

為寡婦辨屈.

先知在這裡提出了三種人,其一是受欺壓的,其二是孤兒,其三是寡婦,他們代表著古代以色列社會中弱小得不能保護自己而需要別人伸出援手的人.對於當時的人來說,「學習行善,尋求公平」的具體行動是怎樣的呢?就是「解救受欺壓的」,「給孤兒伸冤」,「為寡婦辨屈」.

你認為當時那些敬拜耶和華的人寧願遵守11-15節提及的那些繁複的宗教禮儀,還是執行17節提及的三個行動呢?答案看明顯,他們選擇前者.為甚麼呢?試想想,如果不遵守安息日,會立即被發然而且可能被懲罰,但不解救受欺壓的,會受懲罰嗎?如果不獻祭,就可能被抽出來批評,但不給孤兒伸冤又有甚麼後果呢?如果人人守節期你不守,就可能被排斥,但不為寡婦辨屈又有甚麼問題呢?

如果讀聖經的時候,很容易就把自己代入了先知的角色,義正詞嚴地加入指責的行列.現實是,我們可能比較接近被指責的人；撫心自問,我們可以明白和體諒被指責的人的.為甚麼呢?就換個時空以現在的教會作為一個敬拜耶和華的群體來想想,如果教會不舉辦主日崇拜,會不會遭人離棄呢?但教會不參與解救受欺壓的,應該還是平安無事的.如果教會平安夜沒有崇拜,聖誕節沒有活動,受苦節沒有默想聚會,復活節沒有慶祝,教會就會被批評不像教會；然而,教會不「給孤兒伸冤」,不「為寡婦辨屈」,誰人會理會呢?若我們將心比己,就會明白當時的以色列人的選擇了.他們的選擇,大多數的教會都可以明白的,因為如果要在遵守宗教禮儀和「解救受欺壓的」兩者之間作選擇,大家都的選擇都差不多.或者有人會問:「為甚麼一定要二擇其一,Why not both?」當然可以兩者兼顧,也應該兩者兼顧,只是現實上呈現的狀況往往不是如此理想.耶和華上帝如何回應當時以色列人不理想的現實狀況呢?

4．過而能改,善莫大焉

18耶和華說:

你們來,我們彼此辯論.

你們的罪雖像硃紅,必變成雪白；

雖紅如丹顏,必白如羊毛.

我們明白自己的限制,知道自己的軟弱,我們怎敢接受耶和華的挑戰,有甚麼可以與祂辯論呢?我們自知「罪像硃紅,紅如丹顏」,還有甚麼話好說呢?只有慚愧地等待批評.怎知,祂竟然說,「必變成雪白,必白如羊毛」,這真是出乎意料之外.原來,耶和華上帝提供給我們的機會比我們想像的要豐厚很多,祂常常準備接受我們改過.

以賽亞先知勸諫古代的以色列人「學習行善」,剛巧崇基學院的校訓「止於至善」也有一個「善」字.何謂善?有很多種說法.古書上有一句話可以參考:「過而能改,善莫大焉」(左傳).改過,原來也是善.

就算我們立志為善,都只是在通往善的路上跌跌撞撞,誰能夠宣稱完成了「解救受欺壓的；給孤兒伸冤,為寡婦辨屈」這愛鄰舍的使命呢?想不到的是,上帝道成肉身,親自來承擔我們的軟弱,接過了我們所做不到的事情,祂說:

「18主的靈在我身上,因為他用膏膏我,叫我傳福音給貧窮的人；差遣我報告:被擄的得釋放,瞎眼的得看見,叫那受壓制的得自由,19報告上帝悅納人的禧年.」(路加福音4章)

以賽亞先知之後幾百年,上帝關心的仍然是「被擄的,瞎眼的,受壓制的,貧窮的人」,就是那些最缺乏愛的人,社會中弱小得不能保護自己而需要別人伸出援手的人.祂不只是透過先知告訴敬拜祂的群體如此行,更是透過耶穌基督親自做,使我們在這條路上不致孤單,而是有肯定,有榜樣,有同行.願上主幫助我們.



\newpage

\section{帶來分裂的信仰}
\label{sec:612}
\textbf{高國雄牧師 - 帶來分裂的信仰?}
\newline
\newline
\hyperref[sec:611]{< < < PREV SERMON < < <}
~
\hyperref[sec:toc]{[返目錄]}
~
\hyperref[sec:613]{> > > NEXT SERMON > > >}


路加福音 12:49-56
\newline
\begin{longtable}{cl}
\hline
\hline
章節 & 經文 (和合本修訂版)\\
\hline
12:49 & \begin{tabularx}{0.7\textwidth}{X} 「我來是要把火丟在地上，假如已經燒起來，不也是我所希望的嗎？ \end{tabularx} \\ \\ \relax
12:50 & \begin{tabularx}{0.7\textwidth}{X} 我有當受的洗還沒有受，在這事完成之前，我是多麼地焦急！ \end{tabularx} \\ \\ \relax
12:51 & \begin{tabularx}{0.7\textwidth}{X} 你們以為我來是要使地上太平嗎？不！我告訴你們，是使人紛爭。 \end{tabularx} \\ \\ \relax
12:52 & \begin{tabularx}{0.7\textwidth}{X} 從今以後，一家五個人將要紛爭，三個和兩個相爭，兩個和三個相爭： \end{tabularx} \\ \\ \relax
12:53 & \begin{tabularx}{0.7\textwidth}{X} 父親和兒子相爭，兒子和父親相爭；母親和女兒相爭，女兒和母親相爭；婆婆和媳婦相爭，媳婦和婆婆相爭。」 \end{tabularx} \\ \\ \relax
12:54 & \begin{tabularx}{0.7\textwidth}{X} 耶穌又對眾人說：「你們看見西邊起了雲彩，就說：『要下大雨了』，果然就有； \end{tabularx} \\ \\ \relax
12:55 & \begin{tabularx}{0.7\textwidth}{X} 起了南風，你們就說：『要燥熱了』，也就有了。 \end{tabularx} \\ \\ \relax
12:56 & \begin{tabularx}{0.7\textwidth}{X} 假冒為善的人哪，你們知道分辨天地的氣象，怎麼不知道分辨這是甚麼時代呢？」 \end{tabularx} \\ \\
[1ex]
\hline
\hline
\end{longtable}
講題:帶來分裂的信仰?

講員:高國雄牧師

經文:路加福音12:49-56

日期:2022年8月14日

49「我來要把火丟在地上,倘若已經著起來,不也是我所願意的嗎?

50我有當受的洗還沒有成就,我是何等地迫切呢?

51你們以為我來,是叫地上太平嗎?我告訴你們,不是,乃是叫人紛爭.

52從今以後,一家五個人將要紛爭:三個人和兩個人相爭,兩個人和三個人相爭；53父親和兒子相爭,兒子和父親相爭；母親和女兒相爭,女兒和母親相爭；婆婆和媳婦相爭,媳婦和婆婆相爭.」54耶穌又對眾人說:「你們看見西邊起了雲彩,就說:『要下一陣雨』；果然就有.

55起了南風,就說:『將要燥熱』；也就有了.

56假冒為善的人哪,你們知道分辨天地的氣色,怎麼不知道分辨這時候呢?」

前幾天新生來中大認識書院,不同人有不同的考慮因由:喜愛,避開,害怕,擔心選擇錯誤,最有趣是未揀先後悔

面對信仰,上帝的話語,我們又是否同樣用這「選擇」的心及態度呢?

基督教信仰,耶穌基督的教訓=地上太平的福音?

v51你們以為我來,是叫地上太平嗎?我告訴你們,不是,乃是叫人紛爭.

我們可以如果理解?

這與一般人理解很不一樣

以「宗教都是導人向善」來認識基督教有問題嗎?

向善是使人願意去行善,或是以為人可以因此而過更美善的生活呢?

這世上有一些「律」其實是上帝所定的:

做任何事都有「代價」

行善的得獎賞

作惡的要受罰

不會以有罪為無罪

「律」可以是規範,可以是要求,亦可以是期望…

你對信仰/上帝的期望是什麼呢?

有限的人,以自身有限的角度去理解上帝是可明白的,但不應該是唯一的,因為我們的上帝是主動啟示自己的上帝,我們更應從上帝的啟示來認識祂

耶穌基督以比喻教導12:35-40

35「你們腰裡要束上帶,燈也要點著,

36自己好像僕人等候主人從婚姻的筵席上回來.他來到,叩門,就立刻給他開門.

37主人來了,看見僕人警醒,那僕人就有福了.我實在告訴你們,主人必叫他們坐席,自己束上帶,進前伺候他們.

38或是二更天來,或是三更天來,看見僕人這樣,那僕人就有福了.

39家主若知道賊甚麼時候來,就必警醒,不容賊挖透房屋,這是你們所知道的.

40你們也要預備；因為你們想不到的時候,人子就來了.」

作為等待主人的僕人/守門者,其責任在那?如何能夠成為有福的人呢?

警醒!

因為想不到

人生有很多想像/估計不到的事情及時刻,這就是「意外」,意料之外

我們很不想/不喜歡/不願意接受「意外」

人總想掌握一切,控制一切

信仰=趨吉避凶?

信耶穌=得「保佑」?

「主必看顧」又會是什麼意思呢?

彼得前書3:12

–

因為,主的眼看顧義人；主的耳聽他們的祈禱.惟有行惡的人,主向他們變臉.

因主的眼睛(看著)義人,他的耳(聽著)他們的祈禱,但他的面(臨在)卻針對行惡的人.(原文新約)

12:41

彼得說:「主啊,這比喻是為我們說的呢?還是為眾人呢?」

彼得此問所關心的是什麼呢?

自持是門徒就能「與眾不同」?能逃過所要求/期望的責任?

我們是否也是希望祝福臨在自己,要求就放他人身上?

信徒期望有「特別」的待遇?特別的「保守」?

上帝對我特別嚴格?

因為上帝愛我們,亦因為我們屬於祂,有著不同的關係及身份

路加福音12:42-48

42主說:「誰是那忠心有見識的管家,主人派他管理家裡的人,按時分糧給他們呢?

43主人來到,看見僕人這樣行,那僕人就有福了.

44我實在告訴你們,主人要派他管理一切所有的.

45那僕人若心裡說:『我的主人必來得遲』,就動手打僕人和使女,並且吃喝醉酒；

46在他想不到的日子,不知道的時辰,那僕人的主人要來,重重地處治他,定他和不忠心的人同罪.

47僕人知道主人的意思,卻不預備,又不順他的意思行,那僕人必多受責打；

48惟有那不知道的,做了當受責打的事,必少受責打；因為多給誰,就向誰多取；多託誰,就向誰多要.」

耶穌沒有直接回答問題

「忠心有見識」的管家

vs

等待主人的僕人

主人的指派,是要求,是責任,是期待,也是恩典/蒙祝福的好機會

自以為正確,放縱的行為,必受刑罰

知道而不以行動回應,必受責打

不知者,罪比較少,必少受責打

所以主人的指派,是恩典/蒙祝福的好機會,也是要求,是責任,是期待,亦是降服在上帝的律當中.

因此耶穌提到這世代的事實,一個在上帝的律當中的事實

路加福音12:49-56

49「我來要把火丟在地上,倘若已經著起來,不也是我所願意的嗎?

50我有當受的洗還沒有成就,我是何等地迫切呢?

51你們以為我來,是叫地上太平嗎?我告訴你們,不是,乃是叫人紛爭.

52從今以後,一家五個人將要紛爭:三個人和兩個人相爭,兩個人和三個人相爭；53父親和兒子相爭,兒子和父親相爭；母親和女兒相爭,女兒和母親相爭；婆婆和媳婦相爭,媳婦和婆婆相爭.」54耶穌又對眾人說:「你們看見西邊起了雲彩,就說:『要下一陣雨』；果然就有.

55起了南風,就說:『將要燥熱』；也就有了.

56假冒為善的人哪,你們知道分辨天地的氣色,怎麼不知道分辨這時候呢?」

人懂得看雲彩,而知天色變

見南風,而知要燥熱

上帝透過在大自然中的啟示,讓我們認識祂/祂的律

我們是懂得的

但為何仍然扮作不知道如何分辨這世代呢?

耶穌嚴厲地講出事實:這世代不再一樣

家人相爭,地上沒有和平是必然發生的事情

V49耶穌要把火丟在地上

火有時被理解為災禍

路加所提到的火,必然是使徒行傳/教會行傳的主角:聖靈

聖靈的臨在的先決條件就是耶穌「當受的洗成就」,也就是基督要走上十架,受死,埋葬,復活,升天…

領受福氣,領受救恩是要付上代價的

我們的罪就歸在耶穌基督身上

耶穌基督就是最大的贖罪祭

同時,今天我們也要付上代價/責任

耶穌基督沒有因為要討好群眾而避重就輕

我們也不要因為怕別人對信仰負面/感受不好而淡化信仰中對我們的要求

我們太喜愛扮無知

面對終末來臨,不可以再充耳不聞

面對福音,更要勇敢面對,作出應有的選擇

路是窄的,門也是窄的,但上帝恩典是無窮大的

請同來奮力,我們一起攜手,同進天國!



\newpage

\section{耶穌教導的餐桌禮儀}
\label{sec:613}
\textbf{高國雄牧師 - 耶穌教導的餐桌禮儀}
\newline
\newline
\hyperref[sec:612]{< < < PREV SERMON < < <}
~
\hyperref[sec:toc]{[返目錄]}
~
\hyperref[sec:614]{> > > NEXT SERMON > > >}


講題:耶穌教導的餐桌禮儀

講員:高國雄牧師

日期:2022年8月28日

1安息日,耶穌到一個法利賽人的首領家裡去吃飯,他們就窺探他.

7耶穌見所請的客揀擇首位,就用比喻對他們說:

8「你被人請去赴婚姻的筵席,不要坐在首位上,恐怕有比你尊貴的客被他請來；

9那請你們的人前來對你說:『讓座給這一位吧!』你就羞羞慚慚地退到末位上去了.

10你被請的時候,就去坐在末位上,好叫那請你的人來對你說:『朋友,請上座.』那時,你在同席的人面前就有光彩了.

11因為,凡自高的,必降為卑；自卑的,必升為高.」

12耶穌又對請他的人說:「你擺設午飯或晚飯,不要請你的朋友,弟兄,親屬,和富足的鄰舍,恐怕他們也請你,你就得了報答.

13你擺設筵席,倒要請那貧窮的,殘廢的,瘸腿的,瞎眼的,你就有福了!

14因為他們沒有甚麼可報答你.到義人復活的時候,你要得著報答.」

很多人認為信仰是大道理,聖經只記載與生死存亡有關的事情.

常常聽到教導說耶穌基督的教訓就簡單一個「愛」字成為總結就足夠了.

甚至可能因此大家對經文的細節沒有興趣…

同時,反過來另有人就會鑽牛角尖,把聖經中解釋道理的細節放大,不斷研究探討,以致最終可能無法明白及捉緊上主教導的核心.

今天經課所提到的:

耶穌要教導的,期望我們能學習得到的,真的就只是生活智慧?社交禮儀?

為何耶穌要教導我們這些在宴會中的潛規則?

首先,耶穌是要回應那要窺探他的法利賽人首領而講這一個比喻.

這比喻的場景是真實的情況,即是真人真事改編而成的.

婚宴中座席的首末=在婚宴中的地位

自以為重要,自視過高=先取羞辱

「面子是他人比的,架卻是自己丟的」

人總是應該要學習為免丟架…

想起聖經

腓立比書2:3只要存心謙卑,各人看別人比自己強

耶穌基督好像在教導我們待人處事要以謙卑的態度,真的是如此「簡單」/「生活化」嗎?

比喻繼續

V12你擺設午飯或晚飯,不要請你的朋友,弟兄,親屬,和富足的鄰舍,恐怕他們也請你,你就得了報答.

投資望回報是理所當然的!

耶穌卻教導要選擇向沒有能力回報的施予…

v13-14你擺設筵席,倒要請那貧窮的,殘廢的,瘸腿的,瞎眼的,你就有福了!因為他們沒有甚麼可報答你.到義人復活的時候,你要得著報答.

我們不是不要求回報,追求回報是應該的.

而是要求那真正有價值,對永恆生命有價值的回報!

V15-24

15同席的有一人聽見這話,就對耶穌說:「在上帝國裡吃飯的有福了!」

16耶穌對他說:「有一人擺設大筵席,請了許多客.

17到了坐席的時候,打發僕人去對所請的人說:『請來吧!樣樣都齊備了.』

18眾人一口同音地推辭.頭一個說:『我買了一塊地,必須去看看.請你准我辭了.』

19又有一個說:『我買了五對牛,要去試一試.請你准我辭了.』

20又有一個說:『我才娶了妻,所以不能去.』

21那僕人回來,把這事都告訴了主人.家主就動怒,對僕人說:『快出去,到城裡大街小巷,領那貧窮的,殘廢的,瞎眼的,瘸腿的來.』

22僕人說:『主啊,你所吩咐的已經辦了,還有空座.』

23主人對僕人說:『你出去到路上和籬笆那裡,勉強人進來,坐滿我的屋子.

24我告訴你們,先前所請的人沒有一個得嘗我的筵席.』」

由赴宴到設宴:

耶穌分別用宴席主人,赴宴者,相對的二個角度(對立角度)來演繹.

耶穌先以赴宴者,參與宴席的人的態度開始,然後以宴席主人(設宴者)的角色,更深演繹其關注的…

耶穌這比喻,是比喻什麼呢?誰是設宴者?誰是赴宴者?

詩篇23:5在我敵人面前,你為我擺設筵席

啟示錄19:6-9羔羊的婚宴

6我聽見好像群眾的聲音,眾水的聲音,大雷的聲音,說:哈利路亞!因為主－我們的上帝,全能者作王了.7我們要歡喜快樂,將榮耀歸給他.因為,羔羊婚娶的時候到了；新婦也自己預備好了,8就蒙恩得穿光明潔白的細麻衣.9天使吩咐我說:「你要寫上:凡被請赴羔羊之婚筵的有福了!」又對我說:「這是上帝真實的話.」

耶穌基督就是宴席的主人,為我們設宴!

*經課中所教導的不單只是一般赴宴,設宴的智慧或提醒而是以比喻,講出設宴的主的價值觀及態度,祂的心腸,與及我們這些被邀請赴宴的人該如何回應

12耶穌又對請他的人說:「你擺設午飯或晚飯,不要請你的朋友,弟兄,親屬,和富足的鄰舍,恐怕他們也請你,你就得了報答.

13你擺設筵席,倒要請那貧窮的,殘廢的,瘸腿的,瞎眼的,你就有福了!

14因為他們沒有甚麼可報答你.到義人復活的時候,你要得著報答.」

基督設宴的目的不是要我們有什麼回報,即單單是恩典的施予,其實我們根本就無法回報!

17到了坐席的時候,打發僕人去對所請的人說:『請來吧!樣樣都齊備了.』

18眾人一口同音地推辭.

大好的宴席擺在面前,只要接受就可以,但仍然有很多人(原來是大部份人)會用很多不同的藉口拒絕,實在是一個很可惜的真相…

22僕人說:『主啊,你所吩咐的已經辦了,還有空座.』

23主人對僕人說:『你出去到路上和籬笆那裡,勉強人進來,坐滿我的屋子.

24我告訴你們,先前所請的人沒有一個得嘗我的筵席.』」

上主的恩典豐厚,從未想像過可以有機會接受的人,也會被接納!

上主眼中的尊貴(宴席坐上位的貴賓)與世人所理解的不盡一樣,所以必需以謙卑的心,因為能參與宴席全是恩典,獎賞全在於主.

各人總有很多不同的情況,理由,原因拒絕/被拒絕,最終決定接受邀請而赴宴,能夠享受宴席,才是最重要的關鍵



\newpage

\section{千里共嬋娟}
\label{sec:614}
\textbf{林豪恩先生 - 千里共嬋娟}
\newline
\newline
\hyperref[sec:613]{< < < PREV SERMON < < <}
~
\hyperref[sec:toc]{[返目錄]}
~
\hyperref[sec:615]{> > > NEXT SERMON > > >}


腓利門書 1:1-21
\newline
\begin{longtable}{cl}
\hline
\hline
章節 & 經文 (和合本修訂版)\\
\hline
1:1 & \begin{tabularx}{0.7\textwidth}{X} 為基督耶穌被囚的保羅，同弟兄提摩太，寫信給我們所親愛的同工腓利門、 \end{tabularx} \\ \\ \relax
1:2 & \begin{tabularx}{0.7\textwidth}{X} 亞腓亞姊妹，和我們的戰友亞基布，以及在你家裡的教會。 \end{tabularx} \\ \\ \relax
1:3 & \begin{tabularx}{0.7\textwidth}{X} 願恩惠、平安從我們的父神和主耶穌基督歸給你們！ \end{tabularx} \\ \\ \relax
1:4 & \begin{tabularx}{0.7\textwidth}{X} 我在禱告中記念你的時候，常為你感謝我的神， \end{tabularx} \\ \\ \relax
1:5 & \begin{tabularx}{0.7\textwidth}{X} 因聽說你對眾聖徒的愛心，和你對主耶穌的信心。 \end{tabularx} \\ \\ \relax
1:6 & \begin{tabularx}{0.7\textwidth}{X} 願你與人分享信心的時候，能產生功效，讓人知道我們所行的各樣善事都是為基督做的。 \end{tabularx} \\ \\ \relax
1:7 & \begin{tabularx}{0.7\textwidth}{X} 弟兄啊，由於你的愛心，我得到極大的快樂和安慰，因為眾聖徒的心從你得到舒暢。 \end{tabularx} \\ \\ \relax
1:8 & \begin{tabularx}{0.7\textwidth}{X} 雖然我靠著基督能放膽吩咐你做該做的事， \end{tabularx} \\ \\ \relax
1:9 & \begin{tabularx}{0.7\textwidth}{X} 可是像我這上了年紀的保羅，現在又是為基督耶穌被囚的，寧可憑著愛心求你， \end{tabularx} \\ \\ \relax
1:10 & \begin{tabularx}{0.7\textwidth}{X} 就是為我在捆鎖中所生的兒子阿尼西謀求你。 \end{tabularx} \\ \\ \relax
1:11 & \begin{tabularx}{0.7\textwidth}{X} 從前他與你沒有益處，但如今與你我都有益處。 \end{tabularx} \\ \\ \relax
1:12 & \begin{tabularx}{0.7\textwidth}{X} 我現在打發他回到你那裡去，他是我心肝。 \end{tabularx} \\ \\ \relax
1:13 & \begin{tabularx}{0.7\textwidth}{X} 我本來有意將他留下，在我為福音所受的捆鎖中替你伺候我。 \end{tabularx} \\ \\ \relax
1:14 & \begin{tabularx}{0.7\textwidth}{X} 但不知道你的意見，我不願意這樣做，好使你的善行不是出於勉強，而是出於自願。 \end{tabularx} \\ \\ \relax
1:15 & \begin{tabularx}{0.7\textwidth}{X} 他暫時離開你，也許是要讓你永遠得著他， \end{tabularx} \\ \\ \relax
1:16 & \begin{tabularx}{0.7\textwidth}{X} 不再是奴隸，而是高過奴隸，是親愛的弟兄；對我確實如此，何況對你呢！無論在肉身或在主裡更是如此。 \end{tabularx} \\ \\ \relax
1:17 & \begin{tabularx}{0.7\textwidth}{X} 所以，你若以我為同伴，就接納他，如同接納我一樣。 \end{tabularx} \\ \\ \relax
1:18 & \begin{tabularx}{0.7\textwidth}{X} 他若虧負你，或欠你甚麼，都算在我的賬上吧， \end{tabularx} \\ \\ \relax
1:19 & \begin{tabularx}{0.7\textwidth}{X} 我必償還。這是我—保羅親筆寫的。我並不用對你說，甚至你自己也虧欠我呢！ \end{tabularx} \\ \\ \relax
1:20 & \begin{tabularx}{0.7\textwidth}{X} 弟兄啊，希望你使我在主裡因你得益處，讓我的心在基督裡得到舒暢。 \end{tabularx} \\ \\ \relax
1:21 & \begin{tabularx}{0.7\textwidth}{X} 我寫信給你，深信你必順服，知道你所要做的，必過於我所說的。 \end{tabularx} \\ \\
[1ex]
\hline
\hline
\end{longtable}
講題:千里共嬋娟

經文:腓利門書1:1-21

講員:林豪恩先生

日期:2022年9月4日

引言:圓月念遠人

中秋圓月,最容易引起對親友的思念.望著今年的明月,誰會出現於你的思念中呢?是移居了往地球的另一邊生活的親友?還是留下在地球那麼的家人?或者是因為疫情而分隔各地不能共聚的同伴?舉頭望明月,我們心裡面都有同一心願,就是蘇東坡在水調歌頭所唱的「但願人長久,千里共嬋娟」.

當年,在羅馬被囚的保羅也想著一位遠方的人,他是歌羅西的大戶腓利門.於是保羅寫了一封信給他,就是我們今天書信經課所讀的腓利門書.為甚麼保羅想著千里之外的腓利門呢?保羅這封給腓利門的信想表遠甚麼呢?這封信對二千年後的讀者又有甚麼啟發呢?

一．但願人長久,千里共嬋娟

蘇東坡以「但願人長久,千里共嬋娟」的祝願結束他的詞,保羅以類似的祝願開始他的信:

1:1為基督耶穌被囚的保羅,同兄弟提摩太寫信給我們所親愛的同工腓利門,

1:2和妹子亞腓亞並與我們同當兵的亞基布,以及在你家的教會.

1:3願恩惠,平安從上帝我們的父和主耶穌基督歸與你們!

若你想起遠方的某人,你想起他們的甚麼呢?保羅想起腓利門,想到的是腓利門的愛心和信心:

1:4我禱告的時候提到你,常為你感謝我的上帝；

1:5因聽說你的愛心並你向主耶穌和眾聖徒的信心.

1:6願你與人所同有的信心顯出功效,使人知道你們各樣善事都是為基督做的.

1:7兄弟啊,我為你的愛心,大有快樂,大得安慰,因眾聖徒的心從你得了暢快.

二．人有悲歡離合,月有陰晴圓缺,此事古難全

相信保羅並不是要美化腓利門,正如蘇東坡的感悟:「人有悲歡離合,月有陰晴圓缺,此事古難全」.世上沒有完美的人,腓利門也不例外.不過,想念對方的好,還是記住對方的不好,這是可以選擇的.如果想念對方的好,會讓他把這些好發揮得更好,豈不是美事嗎?保羅曾經在寫給哥林多教會的信指出,信仰的美德有三,其一是信心,其二是盼望,他認為永不止息的是愛心.腓利門的愛心本來就被肯定,保羅盼望腓利門的愛心更長,更闊,更深,更具體地展現於信中所談及事情之上.

1:8我雖然靠著基督能放膽吩咐你合宜的事,

1:9然而像我這有年紀的保羅,現在又是為基督耶穌被囚的,寧可憑著愛心求你,

1:10就是為我在捆鎖中所生的兒子阿尼西謀求你.

1:11他從前與你沒有益處,但如今與你我都有益處.

1:12我現在打發他親自回你那裡去；他是我心上的人.

1:13我本來有意將他留下,在我為福音所受的捆鎖中替你伺候我.

1:14但不知道你的意思,我就不願意這樣行,叫你的善行不是出於勉強,乃是出於甘心.

從保羅用那麼多個人情感用語來看,這似乎是一件極困難的事情.這事情何以如此困難呢?保羅信中提及那位名叫阿尼西謀的人,原來是從腓利門那裡逃走的奴隸.信中暗示這位奴隸在主人家裡做了一些錯事,逃走到一個大城市,希望能夠在人群中隱姓埋名而不被發現.我們不知道阿尼西謀如何認識保羅,只從信中得知阿尼西謀受到保羅的影響,悔改歸主(第10節)；並受保羅鍾愛,成為保羅心上的人(第12節)；他參與了保羅的福音工作,以致保羅有意把他留下侍侯,作他「親愛忠心的弟兄」(第13節).

三．在不完美的世界中締造美好

對於現代人來說,奴隸制度是罪惡的.現代人可能想,如果保羅要求腓利門釋放奴隸,還阿尼西謀自由,那才是應有之義.然而,保羅竟然如此安排:

1:17你若以我為同伴,就收納他,如同收納我一樣.

保羅決定把阿尼西謀這位奴隸送回去他的奴隸主那裡,這樣做豈不是送羊入虎口?在保羅的年代,奴隸被視為奴隸主的擁有物之一,是奴隸主可以任意處置的.對於一個犯錯逃走的奴隸,回到奴隸主那裡會被怎樣對待呢?難道保羅這自由人完全不知道奴隸的慘況嗎?難道保羅只是口中說公義而完全莫視奴隸制度的罪惡嗎?如果保羅也認為奴隸制度是不公義的,為甚麼不起而推倒這不公義的制度呢?改變制度,豈不是更徹底地解決問題嗎?

在現代社會,這些詰問顯得義正詞嚴,難以反駁.不過,當我們以超越時空的眼睛來看歷史中的事情,看容易忽略了歷史的限制.事實上,除了科幻片中能夠超越時空的人,所有人都只能活在他所屬的時間和空間之內,受到他那個時代的條件限制.有些制度是經過長時間形成,也可能需要長時間才能被改變.如果我們這些二千多年之後的人質問保羅為甚麼不推翻不義的奴隸制度,未免過於超越時空了.直至最近一二百年,透過很多代人的努力,甚至兵戎相見,奴隸制度才能被廢除.不過,不義的制度暫時不能推翻,是不是就沒有事情可以做呢?在制度的限制之下,還有甚麼可以做呢?以下是當時保羅提出的:

1:15他暫時離開你,或者是叫你永遠得著他,

1:16不再是奴僕,乃是高過奴僕,是親愛的兄弟.在我實在是如此,何況在你呢!這也不拘是按肉體說,是按主說.

奴隸主與奴隸,在奴隸制度下是人與物的關係.當時保羅要求腓利門把阿尼西謀視為一個人,即是轉化成人與人的關係.

第二,奴隸主與奴隸,在制度下是身分尊卑的關係.當時保羅要求腓利門視阿尼西謀為兄弟,即是轉化為平等的關係.

奴隸主與奴隸,在制度下是工具的關係.當時保羅要求腓利門視阿尼西謀為兄弟,即是轉化為共享的關係.

保羅指出,這些要求「也不拘是按肉體說,是按主說」,即是說,平面來看,大家作為人,本應如此相待；但敬拜上主的人,多了一個從上來看的角度,在這個角度來看,大家都是上主的兒女,在上主的眼中本是如此.

在古代奴隸制度之下,奴隸主以他們的權力及財力把人「買起」,把人變成他們的奴隸.現在,奴隸制度已經被廢除,然而,以權力及財力把人「買起」的情況仍然存在嗎?最近有些人被騙去泰國,柬埔寨,緬甸工作,到Po 後被禁錮強逼從事非法活動.旁觀者或者會認為這些人太過無知,不過,這是不是先以財力為餌,後以武力為權來「買起」人來奴役的現象呢?

強逼從事非法活動本身就是非法的,然而,沒有犯法的又如何呢?曾經有一間工廠,有幾十萬員工在廠內工作,某一年內有18位員工在工廠內跳樓自殺.國立中正大學勞工關係學系教授朱柔若(Chu, Jou-Juo)進行了一項研究,揭示該工廠的非人化的管理策略,其中三項如下:

「低無長高」:所謂的「低無長高」基本上是指「低底薪」搭配同意「無彈性加班」下,透過「超長工時」「高耗損體力」的勞動過程,賺取高於底薪兩倍的工資水平.

「羞辱式管理」(另一種說法叫做「屌人文化」):所謂「屌人文化」是指讓工廠的生產管理體制,管理者經常用粗暴的辱罵對員工進行管理,「被屌」是珠三角工人的流行用語(香港人都會明白),被罵的人則稱「挨屌」.據說,使用這類策略的最主要目的,乃在於期待喚起「被屌者」的個人自尊心,進而產生即時行為嚇阻的效果.

「狼性管理」:指其以反人性的管理手段,非法搜查,拘禁,和毆打,不惜壓榨員工經濟價值來追求效率至上的單一目標；在此過程之中,管理者因其奴化員工人本價值而被狼化.

旁觀者或者會說,君不見每天大量人等著進入這工廠工作嗎?如果這工廠如此非人化,怎會有那麼多人渴望進去呢?其實,這也是無奈.當時員工之間流傳著這樣的說法:「一跳保全家,死要死在富士康」.這個順口溜側面地透露出連環跳樓員工的處境及悲哀,要自己死去,才能養家.到底人是甚麼呢?人有甚麼需要呢?心理學家Abraham Maslow 認為人有五方面的需要:(一)生理需要(Physiological needs),(二)安全需要(Safety needs),(三)社交需要(Love and belonging needs),(四)尊嚴需要(Esteem needs),(五)自我實現需要(Self-actualization).離鄉別井長期在工廠生活的工人,他們被視為甚麼呢?是「人力」,專業名詞稱為「人力資源」.這個名詞有沒有令你想起從前用牛耕種的「牛力」,和用馬拉車的「馬力」呢?或許,在現代社會中還隱藏不少被非人性地對待的狀況而沒有引起公眾關注.或者又有人說,不喜歡就不幹吧!現實上,雖然再沒有奴隸制度,但都不能說離開就離開的.

就是以今日的眼光來看待保羅在腓利門書所提出的,在當時其實是非常革命性的.這件事情對腓利門的要求是不容易的,也是艱難的挑戰.能不能把一個奴隸當作一個與自己一樣的人?這不單是挑戰腓利門的個人道德及靈性,更是挑戰著腓利門如何面對他處身的社會.試想想,腓利門收納一個逃走的奴隸而不予嚴懲,其他奴隸主會如何看呢?腓利門視一個奴隸為兄弟,其他奴隸主會怎樣反應呢?其他奴隸主會不會指責腓利門不合禮法,及破壞制度呢?

如果我們從這種處境閱讀保羅寫給腓利門的信,就會感到雖然簡短,但字字都有重量,正是簡短而絕不簡單.不過,保羅並不是站於道德高地,隔岸觀火地指指點點,他是把自己投入了去成就這場改革:

1:18他若虧負你,或欠你甚麼,都歸在我的帳上；

1:19我必償還,這是我－保羅親筆寫的.

保羅在做甚麼呢?在人間來看,是做「架兩」.在信仰來看,是為一個奴隸代贖.保羅為一個奴隸代贖,令人想到基督為世人代贖.從基督代贖的高度來看,腓利門和阿尼西謀一樣都是蒙救贖的人,因此保羅向他指出在信仰上的實況:

1:19我並不用對你說,連你自己也是虧欠於我.

保羅在哥林多前書曾經說:

4:7使你與人不同的是誰呢?你所有的有哪一個不是領受的呢?若是領受的,為何自誇,彷彿不是領受的呢?

在耶穌基督的福音裡面,腓利門和阿尼西謀有甚麼分別呢?

四．但願人長久,千里共神恩

千百年來那一輪明月寄託了多少文人墨客的情深.「但願人長久,千里共嬋娟」的意思是:只希望自己思念的人平安長久,不管相隔千山萬水,都可一起看到明月皎潔美好的樣子.

當我們仰望上主的時候,會否同樣有「但願人長久,千里共神恩」的情懷呢?會不會希望自己思念的人都得到上主的保守和恩典呢?不管相隔千山萬水,都可以一起看到上主創造的美好呢?在時間上,我們的日子或長或短；在空間上,我們的所在地或遠或近；在社會制度上,我們的身份或尊或卑；在經濟上,我們或富或貧,當我們共同仰望同一個「嬋娟」,這些差異的意義可以不同,這些差異所產生的效果也可能不同.「人有悲歡離合,月有陰晴圓缺,此事古難全.」差異仍然繼續存在,人間不是天上宮闕,然而,人與人的關係可以更有情有義.



\newpage

\section{混沌中的盼望}
\label{sec:615}
\textbf{李均熊牧師 - 混沌中的盼望}
\newline
\newline
\hyperref[sec:614]{< < < PREV SERMON < < <}
~
\hyperref[sec:toc]{[返目錄]}
~
\hyperref[sec:616]{> > > NEXT SERMON > > >}


耶利米書 4:11-12,22-28
\newline
\begin{longtable}{cl}
\hline
\hline
章節 & 經文 (和合本修訂版)\\
\hline
4:11 & \begin{tabularx}{0.7\textwidth}{X} 那時，必有話對這百姓和耶路撒冷說：「來自曠野光禿高地的熱風吹向我的百姓，不是為簸揚，也不是為揚淨。 \end{tabularx} \\ \\ \relax
4:12 & \begin{tabularx}{0.7\textwidth}{X} 又有一陣比這更大的風向我颳來；現在，我要向他們宣讀我的判決。」 \end{tabularx} \\ \\ \relax
4:13 & \begin{tabularx}{0.7\textwidth}{X} 看哪，他必如雲湧上；他的戰車如旋風，他的馬比鷹更快。我們有禍了！我們敗落了！ \end{tabularx} \\ \\ \relax
4:14 & \begin{tabularx}{0.7\textwidth}{X} 耶路撒冷啊，你當洗去心中的惡，使你可以得救。惡念在你裡面要存到幾時呢？ \end{tabularx} \\ \\ \relax
4:15 & \begin{tabularx}{0.7\textwidth}{X} 有聲音從但傳出，有災禍從以法蓮山傳來。 \end{tabularx} \\ \\ \relax
4:16 & \begin{tabularx}{0.7\textwidth}{X} 你們當傳給列國，看哪，要向耶路撒冷報告：「有圍攻的人從遠方來到，向猶大的城鎮大聲喊叫。 \end{tabularx} \\ \\ \relax
4:17 & \begin{tabularx}{0.7\textwidth}{X} 他們包圍耶路撒冷，好像看守田園的，因為它背叛了我。這是耶和華說的。 \end{tabularx} \\ \\ \relax
4:18 & \begin{tabularx}{0.7\textwidth}{X} 你的作風和行為招惹這事；這是你罪惡的結果，實在是苦，刺透了你的心！」 \end{tabularx} \\ \\ \relax
4:19 & \begin{tabularx}{0.7\textwidth}{X} 我的肺腑啊，我的肺腑啊，我心疼痛！我的心在我裡面煩躁不安。我不能靜默不言，因我已聽見角聲和打仗的喊聲。 \end{tabularx} \\ \\ \relax
4:20 & \begin{tabularx}{0.7\textwidth}{X} 毀壞的信息不斷傳來，因為全地荒廢。我的帳棚忽然毀壞，我的幔子頃刻破裂。 \end{tabularx} \\ \\ \relax
4:21 & \begin{tabularx}{0.7\textwidth}{X} 我看見大旗，聽見角聲，要到幾時呢？ \end{tabularx} \\ \\ \relax
4:22 & \begin{tabularx}{0.7\textwidth}{X} 「我的百姓愚頑，不認識我；他們是愚昧無知的兒女，有智慧行惡，沒有知識行善。」 \end{tabularx} \\ \\ \relax
4:23 & \begin{tabularx}{0.7\textwidth}{X} 我觀看地，看哪，地是空虛混沌；我觀看天，天也無光。 \end{tabularx} \\ \\ \relax
4:24 & \begin{tabularx}{0.7\textwidth}{X} 我觀看大山，看哪，盡都震動，小山也都搖來搖去。 \end{tabularx} \\ \\ \relax
4:25 & \begin{tabularx}{0.7\textwidth}{X} 我觀看，看哪，無人；空中的飛鳥也都躲避。 \end{tabularx} \\ \\ \relax
4:26 & \begin{tabularx}{0.7\textwidth}{X} 我觀看，看哪，肥田變為荒地；所有城鎮在耶和華面前，因他的烈怒都被拆毀。 \end{tabularx} \\ \\ \relax
4:27 & \begin{tabularx}{0.7\textwidth}{X} 耶和華如此說：「全地必然荒涼，我卻不毀滅淨盡。 \end{tabularx} \\ \\ \relax
4:28 & \begin{tabularx}{0.7\textwidth}{X} 因此，地要悲哀，天上也必黑暗；因為我言已出，我意已定，必不改變，也不由此轉回。」 \end{tabularx} \\ \\
[1ex]
\hline
\hline
\end{longtable}
詩篇 51:1-10
\newline
\begin{longtable}{cl}
\hline
\hline
章節 & 經文 (和合本修訂版)\\
\hline
51:1 & \begin{tabularx}{0.7\textwidth}{X} 神啊，求你按你的慈愛恩待我！按你豐盛的憐憫塗去我的過犯！ \end{tabularx} \\ \\ \relax
51:2 & \begin{tabularx}{0.7\textwidth}{X} 求你將我的罪孽洗滌淨盡，潔除我的罪！ \end{tabularx} \\ \\ \relax
51:3 & \begin{tabularx}{0.7\textwidth}{X} 因為我知道我的過犯；我的罪常在我面前。 \end{tabularx} \\ \\ \relax
51:4 & \begin{tabularx}{0.7\textwidth}{X} 我向你犯罪，惟獨得罪了你，在你眼前行了這惡，以致你責備的時候顯為公義，判斷的時候顯為清白。 \end{tabularx} \\ \\ \relax
51:5 & \begin{tabularx}{0.7\textwidth}{X} 看哪，我是在罪孽裡生的，在我母親懷胎的時候就有了罪。 \end{tabularx} \\ \\ \relax
51:6 & \begin{tabularx}{0.7\textwidth}{X} 你所喜愛的是內心的誠實；求你在我隱密處使我得智慧。 \end{tabularx} \\ \\ \relax
51:7 & \begin{tabularx}{0.7\textwidth}{X} 求你用牛膝草潔淨我，我就乾淨；求你洗滌我，我就比雪更白。 \end{tabularx} \\ \\ \relax
51:8 & \begin{tabularx}{0.7\textwidth}{X} 求你使我得聽歡喜快樂的聲音，使你所壓傷的骨頭可以踴躍。 \end{tabularx} \\ \\ \relax
51:9 & \begin{tabularx}{0.7\textwidth}{X} 求你轉臉不看我的罪，塗去我一切的罪孽。 \end{tabularx} \\ \\ \relax
51:10 & \begin{tabularx}{0.7\textwidth}{X} 神啊，求你為我造清潔的心，使我裡面重新有正直的靈。 \end{tabularx} \\ \\
[1ex]
\hline
\hline
\end{longtable}
提摩太前書 1:12-17
\newline
\begin{longtable}{cl}
\hline
\hline
章節 & 經文 (和合本修訂版)\\
\hline
1:12 & \begin{tabularx}{0.7\textwidth}{X} 我感謝那賜給我力量的我們的主基督耶穌，因為他認為我可信任，派我服事他。 \end{tabularx} \\ \\ \relax
1:13 & \begin{tabularx}{0.7\textwidth}{X} 我從前是褻瀆、迫害、侮慢神的人；然而我還蒙了憐憫，因為我是在不信、不明白的時候做的。 \end{tabularx} \\ \\ \relax
1:14 & \begin{tabularx}{0.7\textwidth}{X} 而且我們的主的恩典格外豐盛，使我在基督耶穌裡有信心和愛心。 \end{tabularx} \\ \\ \relax
1:15 & \begin{tabularx}{0.7\textwidth}{X} 這話可信，值得完全接受：「基督耶穌到世上來是要拯救罪人」，而在罪人中我是個罪魁。 \end{tabularx} \\ \\ \relax
1:16 & \begin{tabularx}{0.7\textwidth}{X} 然而，我蒙了憐憫，好讓基督耶穌在我這罪魁身上顯明他完全的忍耐，給後來信他得永生的人做榜樣。 \end{tabularx} \\ \\ \relax
1:17 & \begin{tabularx}{0.7\textwidth}{X} 願尊貴、榮耀歸給永世的君王，那不朽壞、看不見、獨一的神，直到永永遠遠。阿們！ \end{tabularx} \\ \\
[1ex]
\hline
\hline
\end{longtable}
路加福音 15:1-10
\newline
\begin{longtable}{cl}
\hline
\hline
章節 & 經文 (和合本修訂版)\\
\hline
15:1 & \begin{tabularx}{0.7\textwidth}{X} 許多稅吏和罪人都挨近耶穌，要聽他講道。 \end{tabularx} \\ \\ \relax
15:2 & \begin{tabularx}{0.7\textwidth}{X} 法利賽人和文士私下議論說：「這個人接納罪人，又同他們吃飯。」 \end{tabularx} \\ \\ \relax
15:3 & \begin{tabularx}{0.7\textwidth}{X} 耶穌就用比喻對他們說： \end{tabularx} \\ \\ \relax
15:4 & \begin{tabularx}{0.7\textwidth}{X} 「你們中間誰有一百隻羊，失去其中的一隻，不把這九十九隻留在曠野，去找那失去的羊，直到找著呢？ \end{tabularx} \\ \\ \relax
15:5 & \begin{tabularx}{0.7\textwidth}{X} 找到了，他就歡歡喜喜地把羊扛在肩上。 \end{tabularx} \\ \\ \relax
15:6 & \begin{tabularx}{0.7\textwidth}{X} 他回到家裡，請朋友和鄰舍來，對他們說：『你們和我一同歡喜吧，我失去的羊已經找到了！』 \end{tabularx} \\ \\ \relax
15:7 & \begin{tabularx}{0.7\textwidth}{X} 我告訴你們，一個罪人悔改，在天上也要這樣為他歡喜，比為九十九個不用悔改的義人歡喜還大呢！」 \end{tabularx} \\ \\ \relax
15:8 & \begin{tabularx}{0.7\textwidth}{X} 「同樣，哪一個婦人有十塊錢，若失落一塊，不點上燈，打掃屋子，細細地找，直到找著呢？ \end{tabularx} \\ \\ \relax
15:9 & \begin{tabularx}{0.7\textwidth}{X} 找到了，她就請朋友和鄰舍來，對她們說：『你們和我一同歡喜吧，我失落的那塊錢已經找到了！』 \end{tabularx} \\ \\ \relax
15:10 & \begin{tabularx}{0.7\textwidth}{X} 我告訴你們，一個罪人悔改，神的使者也是這樣為他歡喜。」 \end{tabularx} \\ \\
[1ex]
\hline
\hline
\end{longtable}
講題:混沌中的盼望

經文:耶利米書4:11-12,22-28;詩篇51:1-10;提摩太前書1:12-17;路加福音15:1-10

講員:李均熊牧師

日期:2022年9月11日

起:又審判?

作為神學院教授舊約的講師,我曾說過在崇基禮拜堂講道會盡量解說舊約的經課.但上月和今天我負責宣講的舊約經課都是環繞先知審判的信息.對,又是審判.看舊約的先知書,我們印象中經文都像是充滿對猶太人的審判.尤其是今日經課中耶利米書的經文,「不是為簸揚,也不是為揚淨.」,甚至不是只針對敗懷的領袖,而似要滅絕所有的以色列人.讀得太多這類審判的信息,會否令人變得麻木?好像父母親見到我們就是責備我們:怎麼不早些起床?怎麼衣著這麼不整齊?怎麼說話這麼沒分寸．．．慢慢我們便覺父母煩厭,甚至覺得父母只是挑剔,懷疑我們在父母心中的位置,慢慢也不再聽他們的嚕囌.更何況聖經的審判針對的是以色列人,我們或會問:與我何干?反而新約福音經課,耶穌接納罪人,甚至撇下九十九隻羊,去尋找一隻迷羊呢!今日可能有些基督徒,就如早期教會被定為異端的馬吉安一般,放棄了舊約,認為舊約中的神是一個殘忍暴戾的神,我們相信的卻是新約中一位慈愛溫柔的耶穌.

今日不去討論新舊約的關係,只要你細心讀完整本舊約,你就會找到先知書中上主對祂子民安慰的信息.例如耶利米書,由30章至33章就有所謂「安慰之書」(book of consolation).但怎說也好,不知你讀到好像今日經課的經文,會否覺得「這話很難,誰聽得進呢?」

就讓我們嘗試再仔細閱讀今日的經課,在審判的宣告中尋找一絲的盼望,尋找看先知書的動力.

承:熱風的逆創造

先旨聲明(大家可能也應該知道),先知書的審判信息是有其歷史的脈絡,我們讀的時候不能抽離當中的歷史背景,免得張冠李戴,將十分帶針對性的判詞,想像性凶暴的神要滅絕世界的宣告.例如,在耶利米書裡面,常常提到北方來的「災禍與大毀滅」(4:6,6:22).熟認聖經歷史背景的朋友都會知道,這是指公元前六世紀興起的新巴比倫帝國,對當時的巴勒斯坦地的軍事威脅,肯定不是位處北方的俄羅斯呀之類國家啦．．．故此耶利米書很多的審判之言是針對公元前六世紀的猶大國,主要是針對由約西亞王到西底家王任期內國內所出現的弊端而發.當時政治領袖,宗教,甚至民生都因巴比倫掘起所引發國際局勢的風起雲湧,而變得極度混亂.但最主要的問題是信仰上主的人民再不相信上主的保護,不順從上主的律法,因為時勢艱難人人變得自私自利,又嘗試找方法自保,而不再忠於上主,正如今日經課其中一節:

「我的百姓愚頑,不認識我；他們是愚昧無知的兒女,有智慧行惡,沒有知識行善.」4:22

他們「有智慧行惡」,就是說到他們想盡辦法枉法營私,又以混合宗教的敬拜期望得到更多更大的庇佑.但今日的經課集中講對上主子民的「判決」,而非對他們罪行的剖析,所以這方面就從略了.

既然經文是講判決,我們就仔細看看這個判決是甚麼:

那時,必有話對這百姓和耶路撒冷說:「來自曠野光禿高地的熱風(ruah)吹向我的百姓,不是為簸揚,也不是為揚淨.又有一陣比這更大的風(ruah)向我颳來；現在,我要向他們宣讀我的判決.」4:11-12

「我觀看地,看哪,地是空虛混沌(tohu wabohu)；我觀看天,天也無光.

我觀看大山,看哪,盡都震動,小山也都搖來搖去.

我觀看,看哪,無人(adam)；空中的飛鳥也都躲避.

我觀看,看哪,肥田變為荒地；所有城鎮在耶和華面前,因他的烈怒都被拆毀.

耶和華如此說:「全地必然荒涼,我卻不毀滅淨盡.

因此,地要悲哀,天上也必黑暗；因為我言已出,我意已定,必不改變,也不由此轉回.」(4:23-28)

今日經課選讀的經文略去了4:13-21節關於先知的內心回應,將焦點集中在上主的判決中.心水清的弟兄姊妹或會認得,耶利米書這裡的經文所說的,其實是上主將其創世的工作逆轉.在創世記第一章,神的靈運行在水面上,因為地是空虛混沌,淵面黑暗,然候神依次造出光,讓旱地露出,使天空充滿飛鳥,有用自己形象造人；第二章則造了亞當,將他安置在伊甸園內修理看守.當中神的靈——ruch elohim中「靈」與11,12節的「風」是同一個名詞； 23節空虛混沌——tohu wabohu與創世記1:2 的「空虛混沌」都是tohu wabohu.上主造光,現在卻無光了.上主造人,名叫亞當adam去修理看守,但上主的判決卻使肥田變荒地,無人——沒有adam了,飛鳥也避.所以上主在耶利米書的判決,並不是歷史性,而是宇宙性的,上主原本生育萬物的ruah,靈／風,因為以色列人的罪而成為熱風,大風,使到上主創造的秩序逆轉為tohu wabohu空虛混沌,所創造的光,人,和因而產生的肥田與城鎮都被拆毀.甚至被造的地要悲哀,天上變成灰黑!上主說全地必然荒涼!

這是一個宇宙性的判決,並非歷史性的判決.這是甚麼意思呢?意思就是,不是歷史的偶然,又或大國政治軍事的角力,以至以色列人遭遇國家的劫難.非也,那是創造天地的主,使大地有氣息,使人類有生機,使萬物成為可能的上主,要將他的創造的工逆轉,當以色列人看見這些兆頭時,候知道是宣稱創造這個宇宙的上主的審判,而不是時運或政治而遇到這些災禍.

轉:「我卻不毀滅淨盡」4:27

這段信息看來完全絕望的滅世判決,如何能看到一絲的盼望?的確,先知的宣講,他用的比喻,意像和語言所談及的上主的毀滅性審判,讀的人只會感到絕望,像是一點可能改變的希望也沒有.甚至悔改也顯得太遲,無法逆轉將要面對可怕的結局.但是,若你仔細看完每一句說話,你就會找到,任何一段審判的宣告,都留有一句半句的「空轉子」loophole(通常譯「漏洞」),所謂網開一面也:

「全地必然荒涼,我卻不毀滅淨盡.」4:27

上主之言自我解構(deconsturct itself),一方面要全地必然荒涼,另方面卻不毀滅淨盡!有些學者可能會認為這只是後期文士因為不能接受上主全然滅絕所以附加上這一小句；但若我們搜尋一下這個「不毀滅淨盡」,除了在這裡之名,在耶利米書出現起碼有其他四次:

「你們要上去毀壞它的葡萄園,但不可毀壞淨盡,只可除掉其枝子,因為不屬耶和華.」5:10

「就是在那些日子,我也不會將你們毀滅淨盡.這是耶和華說的.」5:18

「我的僕人雅各啊,不要懼怕；以色列啊,不要驚惶；因我從遠方拯救你,從被擄之地拯救你的後裔；雅各必回來得享平靜安逸,無人能使他害怕.這是耶和華說的.因我與你同在,要拯救你,也要將那些國滅絕淨盡,就是我趕你去的那些國；卻不將你滅絕淨盡,倒要從寬懲治你,但絕不能不罰你.這是耶和華說的.」30:11

「我的僕人雅各啊,不要懼怕!因我與你同在.我要將那些國滅絕淨盡,就是我趕你去的那些國；卻不將你滅絕淨盡,倒要從寬懲治你,但絕不能不罰你.這是耶和華說的.」46:28

同樣的說法也散見其他的先知書,甚至有學者將之稱為餘民神學——就是在上主滅絕審判的當下,先知同時宣告上主必留有餘種,若用今日的說法,就是上主的恩典要保留起來,將來臨到那些倖存者中間.因此,這些經文對被擄流散的一代有特別的意義,亦可能因為如此,這些凶惡的判語才會流傳下來.耶利米書的對應讀者自然是公元前六世紀的猶大國領袖與人民,但同時間也是為亡國被擄到巴比倫和流散不同國家的猶太人而被保存和傳播的,否則當猶大國被巴比倫毀滅後,後世若對耶利米書內容毫無共鳴的話,經文便不會被閱讀或流傳下來.

對我們而言(正如上個月我曾講過),上主的審判若然無何避免,而且是絕對的消滅,那對聽眾來說有何意義?但是,先知書,又或上主的審判,是帶有盼望的,但正如耶利米書30:11或46:28所言,上主不將以色列人滅絕淨盡,但絕不能不罰!刑罰和審判一方面的確使犯罪者知罪,但另一方面那是罪的後果,是不能避免的.當我們面對凶惡,當我們承擔罪責,那是痛苦而絕望的；但這一切混沌過去之後,才有更新的可能,才有悔改回轉的意義,才有重新出發的盼望.

或許有我最近小小的犯錯經驗,作一個類比……

我的經驗與以色列的離開上主以至亡國的刑罰自然不能作比較,但我想指出的是(正如我上月曾說過的):真正的盼望,是在災難的另一邊,但審判是要我們明白災難的因由,以及上主的心意,以致我們可以悔改前行,而不是在罪中原地踏步.悔改是有多個向度,其中之一是悔懊錯誤的過去,但另一個向度則是向前的,轉向神,重新的跟從,為判決後的將來體察上主的思典.

合:賞月念神恩

最後,容我幫充一點.我們不能忽略今日經課的經文所論到的審判是宇宙性的審判.這提醒我們每論順逆,我們並不是被命運播弄,也不是歷史物理偶然的力量,使我們遭遇今日所遭遇的.災困之後仍可重生,全因宇宙是掌握在創造的上主手上.祂造光,也造暗；祂賜災禍,也施平安.今日經課的確是一個沉重的審判,但也叫我們知道,那是創造的上主的工作,而非北方之王,亞述,巴比倫,波斯,或地上任何政權所做成的.耶利米書的後段,所謂「安慰之書」,就有這兩節的提醒:

耶和華使太陽白晝發光,按定例使月亮和星辰照耀黑夜,又攪動大海,使海中波浪澎湃,萬軍之耶和華是他的名,他如此說:「這些定例若能在我面前廢掉,以色列的後裔才會在我面前斷絕,永遠不再成國.這是耶和華說的.」31:35-36

耶和華如此說:「除非我沒有立白日黑夜之約,也未曾安排天和地的定例,

否則我不會棄絕雅各的後裔和我僕人大衛的後裔,使大衛的後裔不再治理亞伯拉罕,以撒,雅各的後裔.我必使他們被擄的人歸回,也必憐憫他們.」33:25-26

今日經課中,上主的宇宙性審判(cosmic judgement)反而彰顯了祂創世的能力,惟有祂的掌權是永久,只有祂能使我們重新再造.以賽亞書40章對著被擄流散的子民,也是這樣提醒他們:

你們要向上舉目,看是誰創造這萬象,按數目領出它們,一一稱其名,以他的權能和他的大能大力,使它們一個都不缺.

雅各啊,你為何說,以色列啊,你為何言,「我的道路向耶和華隱藏,我的冤屈神並不查問」?

你豈不曾知道嗎?你豈未曾聽見嗎?永在的神耶和華,創造地極的主,他不疲乏,也不困倦；他的智慧無法測度.

疲乏的,他賜能力；軟弱的,他加力量.

就是年輕人也要疲乏困倦,強壯的也必全然跌倒.

但那等候耶和華的必重新得力.他們必如鷹展翅上騰；他們奔跑卻不困倦,行走卻不疲乏.賽40:26-31

大自然的規律秩序,非以人的意志可轉移.只要仰觀日夜更替,俯視潮夕漲退,我們就知道造物主仍然掌權.只要我們回歸這位造物之主,祂永遠帶來新的可能,甚至今日造就新的創造:「所以,若有人在基督裡,他就是新造的人,舊事已過,都變成新的了.」林後5:17

今日仍是中秋追月夜,讓我們以賞月去記念神的恩典,相信在祂裡面永遠有新的可能.

我們在混沌中仍有盼望,全因我們可透過基督耶穌轉回到創天造地的主,經歷祂的更新和再造.



\newpage

\section{比亞伯拉罕更亞伯拉罕}
\label{sec:616}
\textbf{鄧瑞強博士 - 比亞伯拉罕更亞伯拉罕}
\newline
\newline
\hyperref[sec:615]{< < < PREV SERMON < < <}
~
\hyperref[sec:toc]{[返目錄]}
~
\hyperref[sec:617]{> > > NEXT SERMON > > >}


路得記 1:16-18
\newline
\begin{longtable}{cl}
\hline
\hline
章節 & 經文 (和合本修訂版)\\
\hline
1:16 & \begin{tabularx}{0.7\textwidth}{X} 路得說：「不要勸我離開你，轉去不跟隨你。你往哪裡去，我也往哪裡去；你在哪裡住，我也在哪裡住；你的百姓就是我的百姓；你的神就是我的神。 \end{tabularx} \\ \\ \relax
1:17 & \begin{tabularx}{0.7\textwidth}{X} 你死在哪裡，我也死在哪裡，葬在哪裡。只有死能使你我分離；不然，願耶和華重重懲罰我！」 \end{tabularx} \\ \\ \relax
1:18 & \begin{tabularx}{0.7\textwidth}{X} 拿娥米見路得決意要跟自己去，就不再對她說甚麼了。 \end{tabularx} \\ \\
[1ex]
\hline
\hline
\end{longtable}
講題:比亞伯拉罕更亞伯拉罕

經文:路得記1:1,6-18;2:1-13;3:1-13;4:1-6,13-17

講員:鄧瑞強博士

日期:2022年9月18日

各位弟兄姊妹,早安.在這疫症裡,祝願各人平安.

我今日講《路得記》.關於《路得記》,首先希望大家知道,希伯來聖經的經卷編排,和我們的聖經的經卷編排,是不同的.在希伯來聖經裡,《士師記》之後,不是《路得記》,而是《撒母耳記》.至於《路得記》,在希伯來聖經裡,是放在經卷結集的第三部分,即所謂「聖卷」(Kethuvim, Writings)部分(我稱之為「雜項」的部分).在其中,《路得記》與《雅歌》,《耶利米哀歌》,《傳道書》,《以斯帖記》又組成一獨特的單元,稱之為「五小卷」(Megillot, Five Scrolls),分別在五個節期中誦讀.《路得記》是在五旬節時誦讀的,慶祝豐收.但基督教聖經將《路得記》抽了出來,放在《士師記》及《撒母耳記》中間,將之看成是歷史書的其中一卷.經卷位置的不同編排,改變了經卷的性質,也左右了讀者對經卷的理解.《士師記》講的,是「無王管」下的社會的一片混亂.《撒母耳記》講的,是王權的設立,藉此建立社會的秩序.基督教聖經將《路得記》放在《士師記》及《撒母耳記》中間,試圖交代從「無王管」到「有王管」的歷史過渡.

一般而言,歷史都是帝王將相的歷史,唯有這些大人物,才有資格創造歷史,扭轉歷史,甚至改寫歷史.《路得記》作為一種歷史書寫,講的卻是小人物的故事,在「沒人管」的混亂時代,正氣的小人物依舊可以織出美麗的人生.正所謂「仗義每多屠狗輩」,仁與義在一些小人物身上滲透出來.不要少看這些仁義的小人物,他們在冥冥之中,創造出歷史的新軌跡,以致《撒母耳記》後來講的王者,也與這些小人物扣上關係.大時代的波動,我們這些小人物管不了.能做的,是活出仁義.就讓仁義的小種子,在冥冥之中,結出意想不到的果子.

路1:1當士師秉政的時候,國中遭遇饑荒.在猶大的伯利恆,有一個人帶著妻子和兩個兒子往摩押地去寄居.

一個簡單的開場白,交代一個時代的種種困局.「士師秉政的時候」,是一個「無王管」的混亂時代.「伯利恆」,希伯來文的意思是「糧食之鄉」,如今,「糧食之鄉」無糧食.有一個人,他的名字叫「以利米勒」,意思是「我的神是王」,但他的神看來保不住他,他要移民保命.他帶著妻子和兩個兒子,移民去摩押.摩押,在舊約傳統裡,是一個對以色列人不懷好意的地方.看看這一家人,生於亂世,「找不到吃的」,神不保佑,只能逃命.

後來,主人公以利米勒死了.他兩個兒子,娶了摩押女子為妻,後來,連這兩個兒子也死了.這是悲劇中的悲劇,一切經濟命脈都斷了.將來如何,不要談將來了.男的死盡,也絕了後,只留下婆媳三人.前路茫茫,以利米勒的妻子拿俄米,只好回老家.

路1:6她就與兩個兒婦起身,要從摩押地歸回；因為她在摩押地聽見耶和華眷顧自己的百姓,賜糧食與他們.

路1:7於是她和兩個兒婦起行離開所住的地方,要回猶大地去.

路1:8拿俄米對兩個兒婦說:「你們各人回娘家去吧.願耶和華恩待你們,像你們恩待已死的人與我一樣!

路1:9願耶和華使你們各在新夫家中得平安!」於是拿俄米與她們親嘴.她們就放聲而哭,

路1:10說:「不然,我們必與你一同回你本國去.」

路1:11拿俄米說:「我女兒們哪,回去吧!為何要跟我去呢?我還能生子作你們的丈夫嗎?

路1:12我女兒們哪,回去吧!我年紀老邁,不能再有丈夫；即或說,我還有指望,今夜有丈夫可以生子,

路1:13你們豈能等著他們長大呢?你們豈能等著他們不嫁別人呢?我女兒們哪,不要這樣.我為你們的緣故甚是愁苦,因為耶和華伸手攻擊我.」

路1:14兩個兒婦又放聲而哭,俄珥巴與婆婆親嘴而別,只是路得捨不得拿俄米.

拿俄米是好人,她希望媳婦不要跟著她,這樣她們才有機會找到新丈夫,建立新家庭,開展新生活.她們若去到猶大地,文化不通,宗教不同,很難組織新家庭,也很難有新生活的可能.再者,拿俄米認為自己「很衰敗」,在1:13裡,她說:「我為你們的緣故甚是愁苦,因為耶和華伸手攻擊我」；或者說:「我想起你們就愁了,我幫不到你們,因為神苦待我.」拿俄米照顧不到自己,也看來照顧不到兩個媳婦,媳婦跟著她,只會「攬著一齊死」.

有一個新抱離開了拿俄米,另一個新抱路得卻作出了非凡的決定.

路1:15拿俄米說:「看哪,你嫂子已經回她本國和她所拜的神那裡去了,你也跟著你嫂子回去吧!」

路1:16路得說:「不要催我回去不跟隨你.你往哪裡去,我也往那裡去；你在哪裡住宿,我也在那裡住宿；你的國就是我的國,你的神就是我的神.

路1:17你在哪裡死,我也在那裡死,也葬在那裡.除非死能使你我相離!不然,願耶和華重重地降罰與我.」

路1:18拿俄米見路得定意要跟隨自己去,就不再勸她了.

路得明確拒絕她婆婆的建議,她作出自己人生的抉擇.婆婆叫她「回去」,她卻堅決「前行」.因對婆婆的愛,她大膽前行,以婆婆的家為家,以婆婆的國為自己的身分認同所在,以婆婆的神為自己敬拜的對象,甘心改變過去一生的價值取向.她為這堅決的抉擇,作一類乎婚姻盟約的誓詞:「從今以後,生死相連」.這盟誓更是以耶和華的名義作的,對路得而言,她的命運已交給這新的神明.

我們說,亞伯拉罕是信心之父,但看看路得這位女子,她作出更偉大的決定.亞伯拉罕有所行動,起碼都因聽到神的呼喚.路得,什麼也沒有聽到,她的行動只出於對她婆婆的愛,一種平凡不過的人間情懷.亞伯拉罕離開吾珥,離開父家,離開過去的一切,但其實,他帶著很多財富和人馬.路得,離開摩押,離開父家,就真的離開過去的一切,她沒有帶著什麼.亞伯拉罕往迦南走去,這是他不認識的地方,但他有強大的人馬作後盾,起碼,他仍有點「抓拿」.路得往猶大地走去,這是她不認識的地方,沒有男人作後盾,對一個外地女子而言,不單沒有「抓拿」,更是危險異常的.路得,因為人間的愛,不怕冒險,無懼艱難,捨棄一切,大膽向全新的世界走去.亞伯拉罕起碼仍講著自己的語言,保持著自己生活的傳統,敬拜著熟悉的耶和華.但路得,卻要適應全新的文化,面對全新的人,學習敬拜一全新的神.簡單來說,路得要放下過去的一切,做一個全新的人.你們說,她的決心,她的抉擇,她的信念,是否比亞伯拉罕要冒更大的險?是否比亞伯拉罕捨棄更多?是否比亞伯拉罕更勇於向未知的領域前進?是否比亞伯拉罕對耶和華有更大的信任?是否比亞伯拉罕更亞伯拉罕?這一切,竟不是出於回應神的呼喚,而只是出於一個小人物對另一個小人物的愛.不需要神的顯現,只需要在平凡生活中仁愛地活.小人物的仁愛,平凡而正氣的生活,默默地改變著歷史的軌跡.如此,為何小看自己小人物的角色呢?

香港的貧窮婦女,有些以拾取紙皮為生.貧窮的路得,以在田間拾取收割人剩下的麥粒為生.

路2:3路得就去了,來到田間,在收割的人身後拾取麥穗.她恰巧到了以利米勒本族的人波阿斯那塊田裡.………

路2:8波阿斯對路得說:「女兒啊,聽我說,不要往別人田裡拾取麥穗,也不要離開這裡,要常與我使女們在一處.

路2:9我的僕人在那塊田收割,你就跟著他們去.我已經吩咐僕人不可欺負你；你若渴了,就可以到器皿那裡喝僕人打來的水.」

路2:10路得就俯伏在地叩拜,對他說:「我既是外邦人,怎麼蒙你的恩,這樣顧恤我呢?」

路2:11波阿斯回答說:「自從你丈夫死後,凡你向婆婆所行的,並你離開父母和本地,到素不認識的民中,這些事人全都告訴我了.

路2:12願耶和華照你所行的賞賜你.你來投靠耶和華－以色列神的翅膀下,願你滿得他的賞賜.」

貧窮的外族女子,在田間,與本地的一個財主相遇.這段邂逅,有很多東西可以說,但時間關係,只能點到即止.財主波阿斯稱呼路得為「女兒」,可能顯示兩者年紀有一段差距.波阿斯先對路得說話,並且說得比路得多很多,內容是善意的,但又保持合宜的距離.他叫路得不要往別處去拾取麥穗,可見他想幫助路得,或許,也想見路得多幾面.他叫路得與他的使女與僕人在一起,顯然,他看路得為自己人,不是外人,甚至有喝僕人打來的水的權利.他知道路得離開自己的父母和家鄉,卻捨不得離開自己的婆婆,情與義都有了.波阿斯由衷地祝福這女子,這女子投靠在神的翅膀下,願她滿得賞賜.特別留意在神翅膀下得「覆庇」這主題.

這位財主波阿斯,是一好人.他欣賞像路得這類有情有義的女子,在他能力範圍內,他盡量提供幫助.這些幫助,不影響他自己的生活,也不用改變他自己的人生.無可厚非,他是好人,社會上多幾個這樣的人,已經很好了.

路得的積極性漸漸改變想法負面的婆婆.婆婆也積極起來,對幸福的前景,有了新的希望.路得的婆婆知道波阿斯對路得有一種安份而純正的好感,便想,可否,將之推進一步?

路3:6路得就下到場上,照她婆婆所吩咐她的而行.

路3:7波阿斯吃喝完了,心裡歡暢,就去睡在麥堆旁邊.路得便悄悄地來掀開他腳上的被,躺臥在那裡.

路3:8到了夜半,那人忽然驚醒,翻過身來,不料有女子躺在他的腳下.

路3:9他就說:「你是誰?」回答說:「我是你的婢女路得.求你用你的衣襟遮蓋我,因為你是我一個至近的親屬.」

路3:10波阿斯說:「女兒啊,願你蒙耶和華賜福.你末後的恩比先前更大；因為少年人無論貧富,你都沒有跟從.

路3:11女兒啊,現在不要懼怕,凡你所說的,我必照著行；我本城的人都知道你是個賢德的女子.

路3:12我實在是你一個至近的親屬,只是還有一個人比我更近.

路3:13你今夜在這裡住宿,明早他若肯為你盡親屬的本分,就由他吧!倘若不肯,我指著永生的耶和華起誓,我必為你盡了本分,你只管躺到天亮.」

依從婆婆的建議,路得更大膽地行動.一個外邦女子,晚間去到打麥場,臥在一男子腳下,其結果可大可小,禍福未卜,但路得就是敢於行動.那晚,那個男人,飲醉了酒,在四野無人的大麥田上,睡著了.這個女子,掀開那男人下身的被,躺在那裡.那男人醒了,surprise,有個女子.男人問:「你是誰?」女子答:「我是你的婢女路得.」這是拉關係,波阿斯從來只是叫路得與自己的使女在一起,從來沒有說路得是自己的婢女.路得說:「求你用你的衣襟遮蓋我」,這其實是用波阿斯先前的講法和挑戰波阿斯,先前波阿斯說:「你來投靠耶和華－以色列神的翅膀下,願你滿得他的賞賜」.「翅膀」這字與「衣襟」這字,其實是同一個字(kanaph).路得其實是對波阿斯說:與其你求耶和華覆庇我,不如你自己覆庇我.而「用衣襟遮蓋」一個人,意味著娶她為妻.兩者可能有年齡差距,但路得明示,我不介意.路得這主動要求,於波阿斯不是件易事.幫助人,而不用改變自己,相對容易.幫助人,而需要整個人「落水」,則很難.捐錢幫助露宿者,易.與露宿者一起生活,難.

路得的要求有道理嗎?路得說:「因為你是我一個至近的親屬」.「至近的親屬」是一個特別的用字(gaal),是指宗族中最有義務購買或贖回死者土地的親人.如此,死者的土地才能保存在自己的宗族內.這裡,應只涉及土地的買贖,不應涉及婚姻.但路得將土地的責任,連繫到婚姻的可能.路得,這位至情,至義,至勇的女子,迫使世上的好人,思考一下,你幫人可以幫到幾盡?只是隔空地讓一個窮人拾幾粒米,抑或讓這窮人進入自己的人生裡?是不「落水」的關懷,抑或是「濕晒身」的參與?

面對這一奇女子,波阿斯下了誓言,親自落水去愛她.但,路得提出的買贖死者土地的要求,在法理上,有另一人更有資格去做,故此,必先擺平這人才行.

路4:1波阿斯到了城門,坐在那裡,恰巧波阿斯所說的那至近的親屬經過.波阿斯說:「某人哪,你來坐在這裡.」他就來坐下.

路4:2波阿斯又從本城的長老中揀選了十人,對他們說:「請你們坐在這裡.」他們就都坐下.

路4:3波阿斯對那至近的親屬說:「從摩押地回來的拿俄米,現在要賣我們族兄以利米勒的那塊地；

路4:4我想當贖那塊地的是你,其次是我,以外再沒有別人了.你可以在這裡的人面前和我本國的長老面前說明,你若肯贖救贖,若不肯贖就告訴我.」那人回答說:「我肯贖.」

路4:5波阿斯說:「你從拿俄米手中買這地的時候,也當娶死人的妻摩押女子路得,使死人在產業上存留他的名.」

路4:6那人說:「這樣我就不能贖了,恐怕於我的產業有礙.你可以贖我所當贖的,我不能贖了.」

其實,買地歸買地,法理上,是不用娶死者遺孀的,只有同住的親兄弟,才有這種要求.波阿斯和那人都不是死者同住的親兄弟,故根本沒什麼婚姻的義務的.明顯,波阿斯「打橫來講」,將這婚姻義務加了上去,如此,那人便不買這地了.也如此,波阿斯明正言順地,買了那地,娶了路得.於路得,於波阿斯,於拿俄米,這是從未想過的全新的開始,源自一比亞伯拉罕更亞伯拉罕的信仰抉擇.

路4:13於是,波阿斯娶了路得為妻,與她同房.耶和華使她懷孕生了一個兒子.

路4:14婦人們對拿俄米說:「耶和華是應當稱頌的!因為今日沒有撇下你,使你無至近的親屬.願這孩子在以色列中得名聲.

路4:15他必提起你的精神,奉養你的老,因為是愛慕你的那兒婦所生的.有這兒婦比有七個兒子還好!」

路4:16拿俄米就把孩子抱在懷中,作他的養母.

路4:17鄰舍的婦人說:「拿俄米得孩子了!」就給孩子起名叫俄備得.這俄備得是耶西的父,耶西是大衛的父.

他們是大衛的曾祖父,大衛是猶大人理想的王者,猶如我國的帝堯帝舜,是民族命運的樞紐人物.

結語:《路得記》講的,是平凡人的生活,在他們生活中,神沒有明顯的介入.這些平凡的人,按其良知,盡心去愛,真誠行義,放膽去構想美好的將來,無懼拒絕地要求人擔起額外的責任,勇毅地向渴想的領域前進.這是平凡不過的生活,卻結出意想不到的歷史果子.各位弟兄姊妹,不要少看你平凡的美善.在冥冥之中,神使用我們這些平凡的人,引導著歷史的前進.

願榮耀歸上主,願人間享和平.



\newpage

\section{那日之後}
\label{sec:617}
\textbf{伍渭文牧師 - 那日之後}
\newline
\newline
\hyperref[sec:616]{< < < PREV SERMON < < <}
~
\hyperref[sec:toc]{[返目錄]}
~
\hyperref[sec:618]{> > > NEXT SERMON > > >}


講題:那日之後

講員:伍渭文牧師

日期:2022年9月25日

引言:倪匡前一段時間辭世,崇基基督教文節一視頻火熱,在視頻中他說:達摩面壁九年,我用十年在美國看聖經,得到一結論:「聖經說今生很短暫,來生很長,不值得為短暫的今生,勞勞碌碌.」其實聖經還提醒我們,要積財寶在天,因為時候一到,那日之後,地上所有財富終歸無有,要善用地上的財富,賙濟窮人,好叫我們能進入將來永存的帳幕.

教會採用之修訂通用經課(Revised Common Lectionary)又稱三代經題,分三年依次每年連續讀完一卷符類福音書.馬太強調耶穌是教師,宣講天國度理;馬可指出耶穌的權能,驅魔醫病,服事有需要的人;路加顯明耶穌的憐憫,關心貧窮的人和弱勢群.

約翰福音每年在預苦期和復活期誦讀.今年是路加年(丙年),請讀經:「有一個財主,穿著紫袍和細麻布衣服,天天奢華宴樂.又有一個討飯的,名叫拉撒路,渾身生瘡,被人放在財主門口…….」(路十六19-31)

I.【財主和拉撒路】

我們先看在這比喻的背景.之前耶穌講了一個「不義(思高譯本,不忠信)的管家」的比喻(路十六1-13).主人把錢給管家放貸,管家不誠實,被人告發,主人定意回來解僱管家.管家知道可以支配的金錢很快歸零,失去工作如何謀生?耕田沒有力氣,討飯又害羞.在主人回來前,快快利用職權善待債仔,建立關係,謀個門路,將來好辦事.在這點上,耶穌讚賞不義管家今世之子,較光明之子更加聰明.他懂得利用「那不義(屬世)的錢財(NIV, worldly wealth)結交朋友,到了錢財無用的時候,他們可以接你們到永存的帳幕裡去.」 (9節)

另一方面,耶穌警誡我們不可效法不忠信的管家,因為錢財的誘惑,失去主人的信任,失去受托管理更多錢財的機會.「人在最少的事上忠心,在大事上也忠心……倘若你們在屬世的錢財上不忠心,誰還把那真實的錢財托付你們呢?」(10-12節)跟著耶穌作了一結論:「一個僕人不事奉兩個主,不是惡這個愛那個,就是重這個輕那個,你們不能又事奉神,又事奉瑪門.」(13節)法利賽人是貪愛錢財的,他們聽見這一切話,就嗤笑耶穌..跟著耶穌講了財主與拉撒路的比喻.

這比喻接續不義(不忠信)管家的教訓,管家貪財不誠實是事奉瑪門,觸犯了禁止的罪(sin of commission)；財主愛財,炫耀財富多於榮耀上主,不肯施捨,也是愛瑪門,觸犯了缺失的罪(sin of omission).財主是反面教材,說明沒有利用今世的錢財,接我們到永存的帳幕；這些錢財到了那日,一切歸零.

聖經沒有抑富重貧

比喻把財主與討飯的並列,一個是極富有的財主,住在有圍牆(gated)的大宅.另一個是極須照顧的病人.財主死後在陰間受苦,貧窮者死後在天堂享福.這是肯定貧窮,貶抑富財嗎?不是,因為聖經中敬畏上主的不少是富人:亞伯蘭的金,銀,牲畜極多(創十三2)；約伯是東方最富有的(伯一1)；所羅門王富甲一方；以斯帖是皇后；腓利門是奴隸主(亞尼西母主人)；呂底亞是高檔布商;耶穌葬在財主約瑟的新墓園.

另一方面,聖經說懶惰使人貧窮.箴言十4:「懶散的手,使人貧窮」；傳道書十18:「因人懶惰,房頂塌下;因人手懶,房屋滴漏」.

筵席是天國的記號

財主的問題不是擁有財富,乃沒有用財富憐恤需要幫助的人.「穿著紫色袍和細麻布衣服,天天奢華宴樂.」(路十六19)紫色布用極貴的顏料染成,他連細麻布的內衣也是紫色.他的大宅築有圍牆,重視氣派.討飯的拉撒路渾身生瘡,被人放在財主門口.「又有一個討飯的,名叫拉撒路,渾身生瘡,被人放在財主的門口,要得財主桌子掉下來的零碎充饑.」(20,21節)這是當時社會保障風俗,讓不能照顧自己也沒有人照顧的貧乏者,被帶到大戶人家門前,富人可以讓他進來大宅的柴房或小屋,叫僕婢照顧他們.聖經常用筵席表示天國,耶穌與稅吏並罪人坐席.法利賽人和文士就向耶穌的門徒發怨言說:「你們為甚麼和稅吏,並罪人一同吃飯呢?」(路五30)路十四章說天國像擺設大筵席,人人獲邀參加；十五章浪子歸家父親設宴,大兒子生氣不參加.崇拜領聖餐是參加筵席,基督復臨時我們赴羔羊的筵席.但在比喻中,財主拒絕拉撒路一同坐席.

A.

滿身生瘡被視為不潔

財主並沒有運用其財富治好其病.經文特別提及狗餂拉撒路的瘡:「並且狗來餂他的瘡.」似乎祇有狗才關心他的疾病和痛楚,狗成了他的動物醫生,給他安慰,成為他的朋友.上主並沒有離開拉撒路,上主遺派這條狗陪伴他,減輕痛楚.

正如詩篇經課所說:「他……賜食物給饑餓的……扶起被壓下的人……扶持孤兒和寡婦,卻使惡人的道路彎曲.」(詩一四六7-9)財主除了把桌子掉下來給狗吃的給他一份,沒有做甚麼.天天祇顧奢華宴樂.,沒有把拉撒路渾身是瘡的痛苦放在心上,因而比喻中財主沒有名字,不被上主紀念.相對好撒馬利亞人的故事,作為過客的好撒馬利亞人,見到有客旅遭強盜打到半死,雖然趕路,但立刻用隨身的油和酒包裹傷口,把傷者扶上自己的牲口帶往旅館,交付所需費用,囑咐店主好好照顧,回程再來看傷者,並補上所需之費用.一位完全不相識的過客,尚且如此悉心照顧不幸的人.

B.「要得財主桌子掉下來的零碎充饑」

財主視他為犬類的外邦人.迦南女子求耶穌醫好被鬼附的女兒,耶穌說:「不好拿兒女的餅,丟給狗吃.」,婦人說狗也吃主人從桌子掉下來的碎渣兒.耶穌大讚婦人信心並醫她其女兒.(太十五21-29)

C.

亞伯拉罕的懷中

拉撒路死後「被天使帶去放在放在亞伯拉罕的懷中」(路十六22)顯示他參加亞伯拉罕的筵席.猶太人用餐時身體依向隔離的胸懷,最後晚餐時「有一個門徒(約翰),是耶穌所愛的,側身挨近耶穌的懷裡.」反而財主一個人在陰間受苦.

D.

拉撒路—恩慈,憐恤

財主在陰間求憐恤:「我祖亞伯拉罕哪,可憐我吧.」(路十五24)但上主很早就把憐恤給了他,拉撒路就是上主給財主的憐恤,但他拒絕拉撒路.八福中提到憐恤人的人有福了,因為他們必蒙憐恤.若財主憐恤拉撒路,進而幫助其它有需要的人,慢慢明白這樣用錢比天天奢華宴樂來得有意義,身體也健康.為何財主天天奢華宴樂?會不會是心靈不滿足的一種補償?上主希望拉撒路的出現,財主學曉憐憫的美德；心中沒有上主的話語,就如污鬼被逐後,心靈打掃干淨,引來更多污鬼.(太十二45)所以湯馬斯.錢伯(Thomas Chamber)講論約翰一書二15「不要愛世界和世界上的事.人若愛世界,愛父的心就不在他裡面了.」把講章冠名:The Expulsive Power of A New Affection(新的愛慕的排拒能力).學曉憐憫和分享,也學曉敬虔和知足,明白「沒有帶甚麼來世上來,也不能帶甚麼走.」(提前六7)上主把一些需要幫助的人帶到我們面前,其實是我們的天使；憐恤人的人有福了,因為他們必蒙憐恤.

但,財主幡然醒悟悔改了嗎?沒有悔意的財主.財主生前把拉撒路排除筵席之外,在陰間見到拉撒路仍全無愧疚.他沒有說:「拉撒路,真抱歉,我後悔生前沒有好好照顧你.」他還像生前隨意指喚人替他辦事,當他是婢僕差遣服事自己,顯然他沒有真正的悔改.財主祇想脫離苦況,並沒有真正悔改.所以亞伯拉罕指出:「你我之間,有深淵限定」.

II【有深淵限定】

你我之間,有深淵限定.

「可憐我吧,打發拉撒路來,用指頭尖蘸點水,涼涼我的舌頭.亞伯拉罕說:「兒阿,你該回想你生前享過福,拉撒路也受過苦……你我之間,有深淵限定,以至人要從這邊過到你那邊,是不能的.」(路十六26)

深淵英文是Chasm分裂之意,從來沒有用這字描述天堂與地獄上下的鴻溝.這字用來表述平面的巨大裂縫,完全不同甚至相反的人生觀,價值觀,世界觀.

財主說:「打發拉撒路到我父家去,因為我還有五個兄弟,他可以對他們作見証,免得他們也來到這痛苦的地方.」亞伯拉罕說:「他們有摩西和先知的話可以聽從.他(財主)說,我祖亞伯拉罕哪,不是的,若有一個人從死裡復活的,到他們那裡去的,他們必要悔改.亞伯拉罕說,若不聽從摩西和先知的話,就是有一個從死裡復活的,他們也是不聽從.」(路十六29-31) 事實確如此,耶穌使馬大弟弟拉撒路復活.「猶太人見了耶穌所作的事,就多有信他的.」(約十一45)但祭司長和法利賽人就有不同反應:「若這樣由著他,人人都要信他.」(約十一.48)他們從那時開始定意捉拿耶穌.人硬了心,一切要依我的旨意而行,就算死人從死裡復活,也無補於.你看信與不信的深淵限定,何等的大.深淵限定,勢不兩立被視為浪漫主義先驅的詩人作家威廉.布萊克(William Blake)[1]寫了《天堂與地獄締婚》(The Marriage of Heaven and Hell).他認為人與生俱來就有地獄的熾欲,同時亦希望在地上建立天堂,這是上帝的計劃.[2]透過發展,適應,和昇華,去蕪存菁,邪惡終蛻變成善良.魯益斯(C.S.Lewis)寫了《大離婚》(The Great Divorce)回應.他指出自然的發展,似枝幹岐出的樹木,而不似百川共匯的河流,

時間不能醫治邪惡,咒詛要慢慢解除,人要改邪歸正.

此書透過一夢境,說到離世後人坐著巴士開往地獄的種種見聞.聖靈和魔鬼,嚮往天堂的人和嚮往地獄的人,都是書中的角色.

魯益斯說得好:「最後祇有兩類人,一類對上主說:

願你的旨意成就.

另一類人,是上主對他們說:

願你的旨意成就.」往地獄的人,都是自己的選擇,

沒有自己選擇就沒有地獄,因為凡尋找就尋見,凡叩門的就開門.」[3]

往地獄的人不想往天堂,因為水土不服.有些人行為怪異恐怖,普通人看來的地獄行為,竟當作天堂的樂趣:

縱毒,殘暴,扭曲的性行為…….有些人白天行為像天堂,

晚上行為像地獄;更有一些披上天堂的外衣,內心是地獄.表面道貌岸然,骨子裡錢財和權力的心魔,盤踞其思想,蠶食其良心.他們追求「願我的旨意成就」,

排拒甚至敵對「願你(上主)的旨意成就」.

在書中,說到人離世後,人才看清楚人世間種種遭遇背後的原因.明白世上的痛苦,有時是偽裝的祝福,看來是享福,其實是迷惑.不是嗎?財主在世,天天奢華宴樂,迷失自己；忘記那日一到,一切歸零.

時間不能醫治邪惡,教育也不能醫治邪惡.最近柬甫寨西城人口販賣新聞,使人想起赤柬管治下,用暴力推行極端政治運動,造成六百萬人口中二百萬遇難.執政的核心全是留法高級知識分子,教育不能把地獄變成天堂.1789年法國大革命,革命黨人也多飽學之士,他們設計斷頭台,處決路易十六,也處決不少政見不同的人.貴族羅蘭夫人涉嫌同情皇室被推上斷頭台,臨刑前朝著廣場上的自由雕像鞠躬並說:「自由,自由,古今多少之罪惡,假你之名而行.」政治的狂熱像魔咒,令這些人幹著地獄的事,作為通往天堂的道路.

Yehiel Dinur (耶喜.狄魯)是納粹大屠殺(holocaust) 的幸存者,在審訊納粹頭號戰犯Adolf Eichmann(阿道夫.艾希孟)中作証.狄魯直視被防彈玻璃後邊艾希孟,就像曾經被牽到宰殺之地的的羊面對屠夫.突然,狄魯飲泣,不久暈過去,不是因為勾起傷痛或忿怒暈倒.他後來解釋:「那一刻我突然驚慄起來,我是懼怕自己,我可以成為他一樣,做出這樣可怕的事.這人是惡魔?是瘋子?還是令人感到心寒–是正常的?艾希孟就在我們的心中.」

小結:深淵限定是大決裂(Chasm),不會殊途同歸.往天堂之路的人,對上主說:「願你的旨意成就」；往地獄之路的人,是自己的選擇,上主對他們說:「願你的旨意成就.」

III【天堂與地獄】

被天使帶去

「後來那討飯的死了,被天使帶去放在亞伯拉罕的懷裡.財主也死了,並且埋葬了.他在陰間受痛苦.」(路十六22)離世時拉撒路有天使迎接,而財主氣一斷,甚麼都沒有,而且在陰間受苦.

這段寶貴的經文,給我們窺探到那日之後(Hereafter),令信徒面對死亡時,有著極大的安慰.一位太太臨終時感到懼怕,因為死亡來找她.丈夫緊握太太的手說:不用怕,是天使來接你歸家了.立時太太平靜安詳下來,天使接我們歸家,何等安慰.

著名布道家葛培理2018安息了.媒體在他安息禮拜舉行前一段時間,不斷回顧他的生平和講道,其中有一段提到他如何面對自己的離世.他說:「有一天,有人會說葛培理死了,離開我們了.其實他沒有死,只是改了一個地址.」葛培理說得對,他沒有死,他活在另一個世界–新界(New Territories).新界就是天上榮耀的教會(Church Triumphant),尚在人間的是作戰的教會(Church Militant);那裡有歷代離世的聖徒(聖者)和眾天使,環繞在上主寶座,不停頌讚,也為地上作戰的教會禱告.他們也參與地上教會的崇拜,正如信經所說:

我信聖徒相通.

公禱書殯葬禮儀出殯往墳場安葬時有辭靈禮,其禮文[4]:

主禮:上主的聖者(徒),請前來扶助他；上主的天使,請來迎接他.請接納他的靈魂,轉獻於至高者的台前.

領:基督既召喚了他,願主收納他,並願天使領他到聖祖的懷抱中.

眾:請接納他的的靈魂,轉獻於至高者的台前.

領:主,求你賜給他永遠的安息,並求你以永遠的光明照耀他.

眾:請接納他的的靈魂,轉獻於至高者的台前.

甚麼是天堂?

除了亞伯拉罕的懷裡(筵席).

聖經還用後多意象形容天堂.樂園:「耶穌對他說,我實在告訴你,今天你要同我在樂園裡了.」(路二十三43)；國度:「你們小群,不下懼怕,因為你們的父,樂意把國賜給你們.」(路十二32)；生命的冠冕:「你務要至死忠心,我就賜給你生命冠冕.」(啟二10).天堂是是信心終點,盼望的目標:「我現在被澆奠,我離世的時候到了,那美好的仗我已經打過了,當跑的路我已經跑盡了.」(提後四6-8,腓三13, 14)

天堂是一個空間.永存的帳幕:「在我父的家裡有很多住.」(約十四2)「因為他等候那一座有根基的城,就是神所經營所建造的.」(來十一10)新天新地:「但我們照他的應許,盼望新天新地,有義者居在其中.」(彼後三13)

總的來說:天堂是與上主親密無間的關係,耶穌同在就是天堂.在天堂可以面對面見到上主,真知道他.「親愛的弟兄阿,我們現在是上主的兒女,將來如何,還未顯明,但我們知道主若顯現,我們必要像他.因為必得見他的真體.」

陰間(Sheol)是甚麼?

(1)

陰間Sheol遺骸埋葬的地方.舊約出現65次,但意義還不清楚,King

James譯作墳墓,地獄,坑洞,要看上下文了解:「等到安息在塵土中,這指望必下到陰間的門閂那裡了.」(伯十七16)

「我說,正在我中年之日,必進入陰間之門,我餘下的年歲不得享受.」(賽三十八10)人離世後,下到陰間,即埋葬在地下.

陰間是所有亡者的去處.

「誰能常活免死救他的靈魂脫離陰間的權柄呢?」(詩八十九48).「因為你必不將我的靈魂撇在陰間,也不叫你的聖者見朽壞.」(詩十六10)彼前三19說基督「他藉這靈,曾去傳道給那些在監獄裡的靈聽.」

陰間是亡者的去處,

無論公義和邪惡的人,也要經歷死亡；陰間在這意義,是監獄.但基督到這監獄宣告他已勝過死亡,因為「陰間的權柄不能勝過教會.」(太十六18)耶穌之前的聖徒離世後,落在陰間光明之處,

在亞伯拉罕的懷中,等候彌賽亞的來臨；我們是回顧,他們是等候.不悔改的人,像比喻中的財主,落在陰間的火焰中.

信經:基督下陰間(Harrowing

of

Hell).我們每次崇拜認信基督曾降至陰間,加爾文認為基督極度謙卑,從至高降到至卑,也認同人間的諸般悲慘經驗.路德強調凱旋得勝,降至陰間是要廢掉陰間的權柄.正教會(Orthodox

Church)有大星期六崇拜,紀念基督下到陰間(看圖一),雙手握著亞當與夏娃,把他們(人類)從陰間救拔出來.

正教神聖偉大的星期六崇拜

(凌晨一時,牧區教堂星期五晚上),其禮文如下:

尊貴的約瑟,從十架取下至純潔的身體,用潔白細麻和香料包裹,放進新墳墓.

榮耀歸於父……

不朽的生命,紆尊降陰間,

閃耀的神性,把地獄滅絕,

你把死人從陰間帶走.

天軍高聲歌頌:

生命賜予者,基督我主,榮耀歸給你.……

婦女攜香膏,迎向墓穴來

天使站墳前,大聲高喊叫:

沒藥是為死人,基督不知朽壞.[5]

IV那日之後

凡事都有限期

在超市買食物,聰明的消費者不單看價錢,也看標簽標示的成份,更重要是產品時限.食物都列出安全食用的時限.路加十六章財主與拉撒路.不義的管家及十五章浪子的三個比喻,

都談及擁有錢財的時限.

(1)路十五11-32浪子分家產,浪費父親的錢財.「小兒子就把他一切所有的,都收拾起來,往遠方去了,在那裡任意放蕩,浪費貲財.」(路十五13)浪費也有時限,

時限一到便成貧困.

他快快分了父親的家產,急急揮霍與娼妓即時的歡樂中.

與之相反,舊約經課(耶三十二1-3a,

6-15)記錄耶利米用金錢買地,為將來神的計劃投資.耶利米因說預言(也是兇言):巴比倫人入侵,是對子民拜偶像的懲罰,國民將被擄,但他日將回歸故土而被拘禁.「那時巴比倫王的軍隊圍困耶路撒冷,先知耶利米囚在護衛兵的院內,在猶大王的宮中.」(耶三十二2)

但他用錢買了一塊地,並將地契封存說:「將來在這地必有人再買房屋,田地和葡萄園.」(耶三十二15)耶利米使用金錢為將來打算,不像浪子追求即時肉欲的歡愉；這地是被擄回歸的信心,天國大業的投資.巴比倫人入侵,戰爭把所有今世的財富瞬間歸於無有,但投資在神國的事業不永恆價值.

(2)路十六1-12不義管家(不忠信)浪費主人的錢財,主人要免去管家的職務,管家也失去管理錢財的機會,遑論將來受托更重的責任

(路十六13).

機會也有時限,現在有的機會,將來不一定有.

(3)財主與拉撒路.財主浪費自己的錢財.他沒有好好賙濟拉撒路,失去積財寶在天的機會,死後下到陰間受苦.

回眸過去,我們真的浪費不少時間,也錯過很多機會,但這些都可以贖回來,浪子可以回頭.過去是不義的管家,可以回轉求上主給予機會.但財主與拉撒路的比喻,是不能逆轉的死亡.那日之後,決定了永遠的命運.

耶穌曾跟祇顧屯積財富的財主說:「無知的人哪,今夜必要你的靈魂;你所預備的要歸誰呢?」(路十二20)那日是何日?前有一日.但我們不用對那日的來臨患得患失,因為在今世我們己經歷到那一日了,窺探到將要來的那一天.

那日是快樂日:藉著重生的洗與聖靈的更生,我們已經出死入生了.(多三5)聖詩快樂日:「(1)前有一日我已決定,信靠耶穌救主我神,那時心中快樂歡欣,願將此樂普告萬民.(2我真快樂與主立約,我願全心愛主忠誠,當我進入上主聖殿,歡樂歌聲充滿會中.(3)救贖大恩今已完成,主為我君,我為主民,我願跟從救主引領,傳揚福音,無上歡欣.副歌:快樂日,快樂日,耶穌洗淨我眾罪孽.主教導我儆醒禱告,快樂日,快樂日,耶穌洗淨我眾罪孽.」

快樂日使我想起信主的那天.1963年我中學的佈道會決志信耶穌,講道經文徒四12:「除他以外,別無拯救,因為在天下人間,沒有賜下別的名,我們可以靠著得救.」坐在回家巴士頂層,感到從未有過的快樂.我問邀請我出席佈道會的同學說:這快樂會否有一天消失?

那日之前.生命的終點,是生命圓滿的彩帶,行經時間,完成上主交付的工作.保羅說:「那美好的仗我已經打過了,當跑的路我已經跑盡了,所信的道我已經守住了.」(提後四7)

那日之後,我們都被帶走,被天使帶走,帶進亞伯拉罕的懷中.

面對面迎見我們的救贖主.今生很短暫,來生很長,不值得為短暫的今生,勞勞碌碌.願我們行經時間,成就上主的旨意,善用今世的財富,進入永存的帳幕.阿們.

[1] 布萊克名句:To see a world in a grain of sand, And a heaven in a wild flower, Hold infinity in the palm of your hand And eternity in an hour .一沙見世界,一花窺天堂.手心握無限,須臾納永恆.(譯者不詳)

[2] 「耶穌全是德行,但他率性而行,不是按儀文章則.耶穌曾說:我看見天使,張開兩手擁抱烈焰,在火中升天,像以利亞.備注:這天使現變成撒旦-我的特別朋友.」(Jesus was all virtue, and acted from impulse, not from rules…When he had so spoken, Ibeheld the Angel, who strectched out his arms, embracing the flame of fire and he was consumed, and rose as Elizah. Note:–This Angel , who is now become a devil, is my particular friend.) William Blake, The Marriage of Heaven and Hell (Digireads.com Publishing, 2021), p.11.

[3] “There are only two kinds of people in the end: those who say to God, “Thy will be done,” and those to whom God says, in the end, “Thy will be done.” All that in Hell, choose it. Without that self-choice there could be no Hell. No soul that seriously and constantly desires joy will ever miss it. Those who seek find. To those who knock it is opened.” C.S. Lewis, The Great Divorce in The C.S. Lewis Signature Classics (HarperCollins, 2017), p. 506.

[4] 台灣聖公會出版之《公禱書各種聖禮之施行,其它禮儀與教會慶典及全體詩篇》《辭靈禮》.

[5] Holy and Great Saturday: “Noble Joseph, taking down Thy most pure body from the Tree, wrapped it in clean linen with sweet spices, and he laid it in a new tomb. Glory be to the Father… Going down to death, O Life immortal, Thou hast slain hell with the dazzling light of Thy divinity. And when Thou has raised up the dead from their place beneath the earth, all the powers of heaven cried aloud: ‘Giver of Life, O Christ our God, glory to Thee.’ …. The Angel stood by the tomb, and to the woman bearing spices he cried aloud: ‘Myrrh is fitting for the dead, but Christ has shown Himself a stranger to corruption.”(The Lenten Triodion, translated by Mother Mary and Archimandrite Kallistos Ware, 1977, 622.



\newpage

\section{當你凝視深}
\label{sec:618}
\textbf{林豪恩先生 - 當你凝視深淵}
\newline
\newline
\hyperref[sec:617]{< < < PREV SERMON < < <}
~
\hyperref[sec:toc]{[返目錄]}
~
\hyperref[sec:619]{> > > NEXT SERMON > > >}


講題:當你凝視深淵

講員:林豪恩先生

日期:2022年10月5日

引言

有人說,經濟有週期,經濟擴張或繁榮時期之後,就會進入收縮或衰退時期,盛極就會衰退,谷底又會反彈.各位朋友,你認為現在處於上行週期還是下行週期呢?你認為高峰過了沒有?還是已經沉到谷底呢?有人建議,升市莫估頂,跌市莫估底.估頂可能是無傷大雅的,因為經濟上行時歌舞昇平,賣樓賣早了,最多只是賺少了.但跌市時,人心虛怯,低處還更低,萬丈深淵望不見底.有時個人的際遇也類似,有起有跌,有順有逆.話雖如此,但跌進深淵之時,連最豁達瀟灑的人也難以超脫.

蘇東坡詩詞文章所表現的豁達灑脫為後人樂道,然而,當他身陷逆境,處於人生低谷時,也難掩坎坷鬱結的心情.他留下了著名的書法作品《寒食帖》訴說了坎坷時鬱結的心情,不知道你有沒有經驗過類似的狀況呢?

自我來黃州,已過三寒食.(宋神宗元豐二年,蘇東坡受誣告以詩諷時政,險遭斬頭,後來他在弟弟蘇轍,曹太后和其他大員力保下,改謫任黃州團練副使.團練副使只是個虛職,收入讓蘇東坡一家捉襟見肘過活猶然不足.)

年年欲惜春,春去不容惜.

今年又苦雨,兩月秋蕭瑟.(春雨綿綿連二個月,好像把秋天的蕭瑟都帶回來)

臥聞海棠花,泥污燕支(*胭脂)雪.

闇中偷負去,夜半真有力.(夜夜綿綿雨,屋邊的海棠花落污泥,胭脂映雪般的春顏,青壯的活力,一夕間被誰偷走了呢?落污泥的海棠就是蘇東坡自己.)

何殊病少年,病起鬚已白.(冷雨催人病,病中的蘇東坡見自己壯年生白髮,更覺悲愴.)

春江欲入戶,雨勢來不已.

小屋如漁舟,濛濛水雲裡.(真是「屋漏偏逢連夜雨」)

空庖煮寒菜,破灶燒濕葦.(病中,雨中的日子,空蕩的廚房破灶,蘆葦溼透,連火也燒不起來)

那知是寒食,但見烏銜紙(*紙).

君門深九重,墳墓在萬里.

也擬哭塗(*途)窮,死灰吹不起.(窮途無門,真是失望至極)

1.當你凝視深淵

個人順逆有時,蘇東坡的《寒食帖》是在哀自身；城市也有治亂興衰,耶利米哀歌就是悲我城:

1:1先前滿有人民的城,現在何竟獨坐!先前在列國中為大的,現在竟如寡婦；先前在諸省中為王后的,現在成為進貢的.

那城,應該是耶路撒冷,從前的繁華衰落了；從前的地位失去了；從前的榮耀被奪走了.

1:2她夜間痛哭,淚流滿腮；在一切所親愛的中間沒有一個安慰她的.她的朋友都以詭詐待她,成為她的仇敵.

以為可以依靠的,竟然落井下石.曾經是朋友的,包括埃及,亞述,以東,摩押,與亞捫,不但沒有加以援助,有些甚至趁火打劫,把你本來有的資源和地位都搶了過去.

1:3猶大因遭遇苦難,又因多服勞苦就遷到外邦.她住在列國中,尋不著安息；追逼她的都在狹窄之地將她追上.

人民流離失所,如風浪中的小舟,任風雨吹打,隨時被這無情的風浪淹沒.

1:4錫安的路徑因無人來守聖節就悲傷；她的城門淒涼；她的祭司歎息；她的處女受艱難,自己也愁苦.

1:5她的敵人為首；她的仇敵亨通；因耶和華為她許多的罪過使她受苦；她的孩童被敵人擄去.

1:6錫安城的威榮全都失去.她的首領像找不著草場的鹿；在追趕的人前無力行走.

錫安悲傷,歎息,受艱難,愁苦的時候,竟然看見「仇敵亨通」,怎不叫人悲憤呢?在「無力行走」的時候仍被敵人追趕,怎不叫人絕望呢?

不知道蘇東坡寫《寒食帖》的時候,有沒有想起誣告他的人呢?自己飢寒交逼的時候,他們是否正在大魚大肉呢?自己「屋漏偏逢連夜雨」的時候,他們是否正在高床軟枕呢?自己病臥在床的時候,他們是否正在意氣風發呢?

自己受傷是痛,但見到傷害自己的人快樂,就更加苦不堪言,內心難平,這種痛苦會打擊人們的心智,可以扭曲人們的良知.

2.深淵也凝視著你

在這坎坷鬱結之中,詩篇的作者提出如下的勸導:

37:1不要為作惡的心懷不平,也不要向那行不義的生出嫉妒.

自己愁苦,仇敵亨通,怎可能不心懷不平呢?「殺人放火金腰帶,修橋補路無屍骸」,怎不叫人不嫉妒呢?

37:2因為他們如草快被割下,又如青菜快要枯乾.

常言道,「善有善報,惡有惡報.若然不報,時辰未到」,問題是「要等到幾時呢?」「快被割下」,「快要枯乾」,到底有多快呢?處於水深火熱的人,一日也嫌長,因為一日也是很難過的.這些痛苦,相信詩人是明白的,就算如此,他仍然作出類似的提醒:

37:7不要因那道路通達的和那惡謀成就的心懷不平.

37:8當止住怒氣,離棄忿怒；不要心懷不平,以致作惡.

身體受傷令人感到痛,這是名副其實的「傷痛」；失去重要的人或珍貴的事物也令人感到傷痛,被欺壓也令人感到傷痛,失望也令人傷痛.不過,人們對這些傷痛的直接反應,卻未必是傷痛,更常見到的是憤怒.怒火中燒時候的行動,就很容易「你做初一,我做十五」,甚至替天行道.

還記得鄧博士講士師記時的那個故事嗎?有一位住在以法蓮山地的利未人娶了一個猶大伯利恆的女子為妾,他帶著妾回家的途中在基比亞過夜,結果妾被基比亞人終夜凌辱致死.基比亞人的行為令人髮指,如何回應呢?那人帶著屍體回家,用刀將妾的屍身切成十二塊,使人拿著傳送以色列的四境.「故事發展下去,是以色列聯軍去伐罪.結果是:幾乎將「便雅憫」支派滅絕.便雅憫支派的女人和孩童都死光,只剩下600個男人.忽然,以色列人又覺得十二個支派少了一個支派不太好,於是又想辦法找老婆給這600個便雅憫人.他們的方法是:誰無參加伐罪聯軍的,便將那城的人全殺掉,只留處女.於是,他們出發到「基列˙雅比」進行屠城,留下400處女.但,還欠200個女子.於是,他們強搶示羅的民女去湊數.其實,如此解決問題的方式,是再一次的胡亂屠殺和淫辱婦女,將之前的罪惡再倍大,更多人被殺,更多婦女被害.」

37:9因為作惡的必被剪除；

「殺人放火金腰帶,修橋補路無屍骸」是令人義憤難平的.不過,自己因此也幹「殺人放火」的事情,是無補於事的,只是令世上出現更多自己不願意看見的惡事.惡事始於是惡事,不會因為人人都如此做就變為好事；也不因為做了惡事的人逍遙法外而令惡事變為好事.尼采曾說:「與怪物戰鬥的人,應當小心自己不要成為怪物.當你凝視深淵,深淵也凝視著你.」

在上帝審判未臨到之前,等待「作惡的必被剪除」的日子,我們又可以如何呢?如果不以惡報惡,可以怎樣呢?詩人指出:

37:3你當倚靠耶和華而行善,住在地上,以他的信實為糧；

回應忿怒,詩人呼籲我們:

37:4又要以耶和華為樂,他就將你心裡所求的賜給你.

回應心懷不平和疾妒,詩人建議我們:

37:5當將你的事交託耶和華,並倚靠他,他就必成全.

37:6他要使你的公義如光發出,使你的公平明如正午.

或許我們會問:「要等都幾時呢?」時間並不由我們控制,但詩人仍然忠告我們等候:

37:7你當默然倚靠耶和華,耐性等候他；

37:9惟有等候耶和華的必承受地土.

總結:我們有本分當盡

等候可以令人感到焦慮和煎熬,既然不能以惡報惡,又不能扭轉乾坤,有甚麼好做呢?有「講股佬」說:「最好甚麼也不做,以免做得多輸得多.」甚麼也不做,百無聊賴其實也是一種痛苦,很容易令人思想和情感都沒有依託,精神極度空虛無聊.因此又有人說:「生意淡薄,不如賭博」.這樣,只會陷入更深的深淵而不能自跋.保羅在提摩太後書提醒我們:

1:7因為上帝賜給我們,不是膽怯的心,乃是剛強,仁愛,謹守的心.

1:8你不要以給我們的主作見證為恥,也不要以我這為主被囚的為恥；總要按上帝的能力,與我為福音同受苦難.

1:12為這緣故,我也受這些苦難.然而我不以為恥；因為知道我所信的是誰,也深信他能保全我所交付他的,直到那日.

心理治療學派 – 意義治療與存在主義分析(Existential Psychoanalysis)的創辦人維克多·弗蘭克(Viktor Emil

Frankl是一位二次大戰時猶太人大屠殺的倖存者,他提出了一條公式,D=S-M (Despair equals suffering minus meaning),絕望是受苦但找不到意義.保羅的勸勉正是對抗絕望的出路,表明不以受苦為恥,因為他清楚知道他是為福音受苦.人們的確不知道深淵有多深,不知道何時谷底反彈,保羅的態度值得我們參考,「上帝賜給我們,不是膽怯的心,乃是剛強,仁愛,謹守的心」.

在下行的週期中,是否甚麼也不做就是最好呢?耶穌在路加福音所講的比喻指出,其實我們在任何時候都有應分做的事情:

17:7你們誰有僕人耕地或是放羊,從田裡回來,就對他說:『你快來坐下吃飯』呢?

17:8豈不對他說:『你給我預備晚飯,束上帶子伺候我,等我吃喝完了,你才可以吃喝』嗎?

17:9僕人照所吩咐的去做,主人還謝謝他嗎?

17:10這樣,你們做完了一切所吩咐的,只當說:『我們是無用的僕人,所做的本是我們應分做的.』」

普頌735「本分當盡」是我們經常在主日崇拜中唱的詩歌,我們如此唱:

我有本分當盡,我有真神當尊,

我有永恆靈魂當救,離罪與主相親.

我願奉獻此生,成就主派之工,

將我才能不住施展,勤懇勉力至終.

求主使我謹慎,如在主前為人,

常作準備,將我責任一一向主面陳.

助我警醒禱告,終身靠主引導,

恪守基督仁愛之道,向主永歸榮耀.

這樣生活,是不是需要超凡的信心呢?

17:5使徒對主說:「求主加增我們的信心.」

17:6主說:「你們若有信心像一粒芥菜種,就是對這棵桑樹說:『你要拔起根來,栽在海裡』,它也必聽從你們.

「若有信心像一粒芥菜種」,就可以叫一棵桑樹自己連根拔起移植到海裡.其實,耶穌描述僕人所應盡的本分,都只不過是日常工作而已,又不是「叫一棵桑樹自己連根拔起移植到海裡」的超自然神奇事情,何需要談有沒有信心呢?要有信心才做,沒有信心就不做,那可能是藉口而已.我們還是學耶穌比喻中僕人的態度吧:『我們是無用的僕人,所做的本是我們應分做的.』上主自會在祂認為合適的時候作出回應.願主幫助我們.



\newpage

\section{王者的隕落}
\label{sec:619}
\textbf{鄧瑞強博士 - 王者的隕落}
\newline
\newline
\hyperref[sec:618]{< < < PREV SERMON < < <}
~
\hyperref[sec:toc]{[返目錄]}
~
\hyperref[sec:620]{> > > NEXT SERMON > > >}


撒母耳記上 13:5-14
\newline
\begin{longtable}{cl}
\hline
\hline
章節 & 經文 (和合本修訂版)\\
\hline
13:5 & \begin{tabularx}{0.7\textwidth}{X} 非利士人集合，要與以色列人作戰。他們有戰車三萬輛，騎兵六千，士兵像海邊的沙那樣多。他們上來，在伯‧亞文東邊的密抹安營。 \end{tabularx} \\ \\ \relax
13:6 & \begin{tabularx}{0.7\textwidth}{X} 以色列人見自己危急，軍隊被圍攻，百姓就藏在山洞、叢林、巖隙、地窖和深坑中。 \end{tabularx} \\ \\ \relax
13:7 & \begin{tabularx}{0.7\textwidth}{X} 有些希伯來人過了約旦河，逃到迦得和基列地。掃羅還在吉甲，所有的人都戰戰兢兢地跟隨他。 \end{tabularx} \\ \\ \relax
13:8 & \begin{tabularx}{0.7\textwidth}{X} 掃羅照著撒母耳所定的日期等了七日。但是，撒母耳還沒有來到吉甲，百姓就離開掃羅散去了。 \end{tabularx} \\ \\ \relax
13:9 & \begin{tabularx}{0.7\textwidth}{X} 於是掃羅說：「把燔祭和平安祭帶到我這裡來。」掃羅就獻上燔祭。 \end{tabularx} \\ \\ \relax
13:10 & \begin{tabularx}{0.7\textwidth}{X} 他剛獻完燔祭，看哪，撒母耳就到了。掃羅出去迎接他，向他問安。 \end{tabularx} \\ \\ \relax
13:11 & \begin{tabularx}{0.7\textwidth}{X} 撒母耳說：「你做了甚麼事啊？」掃羅說：「因為我見百姓離開我散去，你又不照所定的日期來到，而且非利士人已在密抹集合； \end{tabularx} \\ \\ \relax
13:12 & \begin{tabularx}{0.7\textwidth}{X} 我說：『現在非利士人已經下到吉甲來攻擊我，可是我還沒有向耶和華禱告。』所以我就勉強獻上燔祭。」 \end{tabularx} \\ \\ \relax
13:13 & \begin{tabularx}{0.7\textwidth}{X} 撒母耳對掃羅說：「你做了糊塗事了，沒有遵守耶和華－你神吩咐你的命令。不然，耶和華會在以色列中堅立你的國度，直到永遠。 \end{tabularx} \\ \\ \relax
13:14 & \begin{tabularx}{0.7\textwidth}{X} 現在你的國度必不長久。耶和華已經尋著一個合他心意的人，立他作百姓的君王，因為你沒有遵守耶和華所吩咐你的。」 \end{tabularx} \\ \\
[1ex]
\hline
\hline
\end{longtable}
講題:王者的隕落

經文:撒母耳記上13:5-14

講員:鄧瑞強博士

日期:2022年10月16日

今日,我們一起讀《撒母耳記上》的一個小片段.《撒母耳記》原文不分上下兩卷,上接《約書亞記》,《士師記》,下開《列王記》.這四卷書可以看成是一整體,以信仰的眼睛,細看歷史的興亡.《士師記》講的,是「無王管」的混亂時期.《撒母耳記》講的,是君王制度的設立.《列王記》講的,是君王制度的失敗,終結於猶大國被巴比倫所滅,神的子民被擄巴比倫.在猶大國被巴比倫所滅之後,猶大的神學家坐在巴比倫的河邊,作神學的沉思.他們想起,古時,約書亞帶領他們那些一無所有的先祖進入迦南地,然後,他們散亂地在這裡生活了很久.再然後,他們設立了君王制度,開始時,真是一段輝煌的歲月.再然後,經過四百多年的盛衰,最後,國破家亡,長恨綿綿.耶路撒冷變成荒地,莊嚴的聖殿一片凋零,人去城空,只留下一些野狗在白骨之間徘徊.那些猶大的神學家在孤寂中,在虛無中,拿起筆來,記下歷史沉重地犁出來的深坑,在這些令人難忘的痕跡裡,挖掘出盛衰興亡背後的天理,好讓自己的民族,在那空洞得像活在真空裡無法呼吸的日子,明白自己走到這絕地,自有其緣由,也因為有跡可尋,才能回頭逆行,重返故鄉.

西漢的司馬遷,寫信給「任安」這位朋友,說明自己寫《史記》的用意,是「究天人之際,通古今之變,成一家之言」.歷史的發展,總有些測不透的因素,這是天意.事情的成敗,人的抉擇也很重要,這是人意.歷史的興亡,是天意與人意的際會,故謂「天人之際」.歷史家的史識,就是透視這天人之際的軌跡,明乎其理,則能明白古今之變的通則,故謂「通古今之變」.若這識見能幫助活在歷史洪流裡的人有所領悟,以天地之正氣,作美善的抉擇,引歷史回正途,則這歷史家便能「成一家之言」. 在巴比倫的河邊,失去了身分,沒有了前景的猶大神學家,正是要整理這幾百年來的歷史,他們也是要「究天人之際,通古今之變,成一家之言」,為的是從歷史汲取教訓,好能在沒有了歷史的絕望中,可以走出黑暗,再寫歷史.他們以摩西五經裡神的誡命為天意,特別是《申命記》裡詳盡的法典,正是神對人明確的要求.這班猶大神學家,以《申命記》法典裡神的要求,審視人的回應,好整理出這幾百年來歷史興亡的道理,於是,寫下了《約書亞記》,《士師記》,《撒母耳記》,《列王記》.現今學者稱這四卷歷史書為「申命記學派」(Deuteronomistic School)的歷史書寫.申命記神學的重點,就是主耶穌所講的:「你要盡心,盡性,盡力,盡意愛主－你的神；又要愛鄰舍如同自己.」(路10:27)人能盡心愛神,敬畏神,就明白自己並不是創造歷史的主人,而是讓天意實現的僕人.如何可見天意的實現,就是愛鄰舍如同自己.一個君王能如此,就不會成為霸主.一個人能如此,就不會為權為利而失去人性.

今日的講道,是看看以色列人第一個王掃羅的隕落.

撒上9:2說:「…掃羅,又健壯,又俊美,在以色列人中沒有一個能比他的,身體比眾民高過一頭.」掃羅系出名門,有錢有面,高大靚仔.一見到他,人人都喜愛他.他看來高人一等,是天生的王者.人喜愛他,神揀選他.他成為以色列人第一個王.

在古代,要證明自己是上天揀選的王,就要證明自己是蒙神祝福的.如何證明?在戰場上獲勝,便是神祝福的證明.戰爭,充滿不確定的因素,沒有神的祝福,無法獲勝.

我們看看今日的講道經文:撒母耳記上13:5—14,看看掃羅王打的一場戰爭.

撒上13:5非利士人聚集,要與以色列人爭戰,有車三萬輛,馬兵六千,步兵像海邊的沙那樣多…

撒上13:6以色列百姓見自己危急窘迫,就藏在山洞,叢林,石穴,隱密處,和坑中.

撒上13:7有些希伯來人過了約旦河,逃到迦得和基列地.掃羅還是在吉甲,百姓都戰戰兢兢地跟隨他.

撒上13:8掃羅照著撒母耳所定的日期等了七日.撒母耳還沒有來到吉甲,百姓也離開掃羅散去了.

撒上13:9掃羅說:「把燔祭和平安祭帶到我這裡來.」掃羅就獻上燔祭.

撒上13:10剛獻完燔祭,撒母耳就到了.掃羅出去迎接他,要問他好.

撒上13:11撒母耳說:「你做的是甚麼事呢?」掃羅說:「因為我見百姓離開我散去,你也不照所定的日期來到,而且非利士人聚集在密抹.

撒上13:12所以我心裡說:恐怕我沒有禱告耶和華.非利士人下到吉甲攻擊我,我就勉強獻上燔祭.」

撒上13:13撒母耳對掃羅說:「你做了糊塗事了,沒有遵守耶和華－你神所吩咐你的命令.若遵守,耶和華必在以色列中堅立你的王位,直到永遠.

撒上13:14現在你的王位必不長久.耶和華已經尋著一個合他心意的人,立他作百姓的君,因為你沒有遵守耶和華所吩咐你的.」

掃羅面對一場戰爭,形勢很壞,對方人強馬壯,以色列人怕得要命.從當時的文化看,地上的戰爭涉及天上的神明.先獻祭,後作戰,才能保證自己得到神明的保佑.

掃羅約好撒母耳來獻祭,但他遲遲不來,掃羅便私自獻了一他不當獻的祭.祭剛獻完,撒母耳就到了.撒母耳問:「你做的是甚麼事呢?」

掃羅答得很詳細.他的回答涉及四個理由,個個都重要.

一,「因為我見百姓離開我散去」,這是第一個理由.軍心動搖,逃兵極多,百姓失去作戰的勇氣和信心,在這光景下,是否應該做點什麼,去振奮士氣呢?

二,「你也不照所定的日期來到」.掃羅明說,這局面的出現,不是他造成的,而是撒母耳不按時間出現.責任在撒母耳,而不在掃羅自己.

三,「而且非利士人聚集在密抹」.敵人集結兵力,已殺到埋身.我能坐以待斃嗎?百姓散去,你又未到,只能靠我自己.做點什麼去穩定民心,有什麼問題呢?

四,第四個理由很「屬靈」的:「所以我心裡說:恐怕我沒有禱告耶和華.非利士人下到吉甲攻擊我,我就勉強獻上燔祭.」其實掃羅明白,打仗不是自己打的,任何地上的打仗都涉及天上的神明.敵人已殺到埋身,必須應戰.應戰前,必須獻祭.掃羅的想法很正路.撒母耳未到,能怎麼樣呢?只能「勉強」獻祭.所謂「勉強」,即不是真想這樣做,但為顧全大局,身不由己,勉為其難去做.

想想這局面,自己人散去,敵人殺到埋身,幫助你的人未到,但你又知某些事情必須先做,只是你沒有資格做.這時,你能怎麼樣?

我好同情掃羅.我比他更軟弱,我比他更驚慌,我比他更衝動.

我想,一般人在掃羅的處境,也會這樣做.

關於這段經文,有本書這樣講(Nico ter Linden,荷蘭牧者,書名叫The Story Goes…):「若這是一般人的故事,掃羅做的,全對.但,這是一信仰故事.在信仰裡,重要的,不是恐懼,而是信靠.只要你是王者,你就必須在群眾面前,顯出更大的信靠.一個領袖,應將眼光全放在神,而不是放在仇敵之上.」

簡單說,按人的一般智慧來說,掃羅做的,合乎人之常情,人之常理,但卻不合乎信仰的法度,不合乎《申命記》的要求.以色列的王者,不是普通人,他是讓神的旨意,神的祝福臨現大地的僕人.他要對神有絕對的信任,否則,他就會強行將自己的統治意志加在百姓身上,漸漸遠離天意.他要對神有絕對的敬畏,否則,他很易成為操控一切的霸主,成為亂世的魔王.他要遵從不應踰越的神聖界線,這界線一越過,他便會自以為自己是決定歷史軌跡的神明,這終將歷史帶向歧途.

信仰的人生,是一敬畏的人生.合乎人之常情常理,不表示合乎神的心意.神若要你「咩都唔好做,企定定係度」,你就只能「咩都唔好做,企定定係度」.神若叫你不要吃那樹上的果子,縱使在你看來,這果子與別的果子根本沒有分別,但神的命令就是不應踰越的神聖界線,你就只能不吃.主耶穌若叫你寬恕仇敵,則縱使在你看來,這要求多麼不合常情常理,你就只能照做.

在南非的歷史裡,白人加害黑人,挖出了充滿仇恨的歷史深坑,看來,仇恨會將他們的歷史前景湮沒,但有人聽從了神那不合情理的和解與寬恕的命令,於是,看來無法避免的仇殺停止了,歷史轉入正途.

以色列的第一個君王掃羅,在那歷史的轉折點,他只看到危機,看到大軍壓境,看到強大的敵人.又看到自己的人民怯懦,看到他們退縮,看到他們散去.掃羅在巨大的壓力下,叫天不應,叫地不聞,他就大膽越過神聖的界線,獻上一他沒有資格獻的祭.他明明知道獻祭的重要性,就這而言,他明明知道歷史的成敗,關鍵在神.在神面前,他就應該停下來,讓神主導這絕望的局面.但他就是按捺不住,做出一些他不應做的事,如此舉動,表明他想將事情的結果操控在自己手中.他原先是明白神的介入的重要性的,但他的舉動表明,他的介入比等待神的介入更重要.當人以為能操控歷史的軌跡時,便將歷史帶離正道了.

早前去學生家,她家養貓.地上放著貓糧,但有我這陌生人在,貓不太安心吃.於是,我懷著善意,拿起貓糧,放到貓的嘴邊,讓牠安心地吃,但自此,貓就隱藏到角落裡去,不再出來了.

其實,只要我安靜不動,時間一過,貓便會安心地去吃東西的,但我按捺不住,越過了某種界線,貓的平靜歷史便被我打破了.

掃羅一獻完祭,撒母耳就到了.撒母耳一口氣的說:「你做了糊塗事了,沒有遵守耶和華－你神所吩咐你的命令.若遵守,耶和華必在以色列中堅立你的王位,直到永遠.現在你的王位必不長久.耶和華已經尋著一個合他心意的人,立他作百姓的君,因為你沒有遵守耶和華所吩咐你的.」

「遵守耶和華所吩咐你的」,這意念貫穿這番話.就因掃羅在壓力下做不到,他的王位便完了.

總結:

司馬遷寫《史記》,寫到西楚霸王項羽.他被困垓下,在烏江自刎.霸王自嘆說:「此天之亡我,非戰之罪也」.司馬遷卻如此評論他:「自矜功伐,奮其私智,而不師古,謂霸王之業,欲以力征經營天下,五年,卒亡其國,身死東城,尚不覺寤,而不自責,過矣.乃引『天亡我,非用兵之罪也』,豈不謬哉!」這位霸主,毫無德行,以力稱霸,到死時,仍不覺悟；還說「天亡我,非用兵之罪也」,竟還推說這只是天意弄人,與他的行動無關.司馬遷說:他的講法,「豈不謬哉」!其實,一切都是他自己作的孽,離開了天意,他便要承受惡果.歷史的成敗,在乎人是否遵從天意.

《撒母耳記》是一系列歷史書寫的一個環節,這一系列歷史書寫,是審視君王在歷史轉折關頭的反應.成敗得失,就在乎這君王是否遵守耶和華所吩咐的法度.這法度的主旨,是敬神和愛人.敬神,人就懂得約束和自限.愛人,人就將天意實現人間.真正的王者,是甘心的僕人,讓上天的仁愛,遍灑人間.唯這些王者,能開萬世之太平.

對很多人來說,此刻是歷史的關鍵時刻.但願,人能從歷史中汲取教訓,敬神愛人,讓歷史穩步前進.

願榮耀歸上主,願人間享和平.



\newpage

\section{上帝啊開恩可憐我這個罪人}
\label{sec:620}
\textbf{高國雄牧師 - 「上帝啊,開恩可憐我這個罪人!」}
\newline
\newline
\hyperref[sec:619]{< < < PREV SERMON < < <}
~
\hyperref[sec:toc]{[返目錄]}
~
\hyperref[sec:621]{> > > NEXT SERMON > > >}


路加福音 18:9-14
\newline
\begin{longtable}{cl}
\hline
\hline
章節 & 經文 (和合本修訂版)\\
\hline
18:9 & \begin{tabularx}{0.7\textwidth}{X} 耶穌向那些自以為義而藐視別人的人講了這比喻： \end{tabularx} \\ \\ \relax
18:10 & \begin{tabularx}{0.7\textwidth}{X} 「有兩個人上聖殿去禱告，一個是法利賽人，一個是稅吏。 \end{tabularx} \\ \\ \relax
18:11 & \begin{tabularx}{0.7\textwidth}{X} 法利賽人獨自站著，自言自語地禱告說：『神啊，我感謝你，我不像別人勒索、不義、姦淫，也不像這個稅吏。 \end{tabularx} \\ \\ \relax
18:12 & \begin{tabularx}{0.7\textwidth}{X} 我每週禁食兩次，凡我所得的都獻上十分之一。』 \end{tabularx} \\ \\ \relax
18:13 & \begin{tabularx}{0.7\textwidth}{X} 那稅吏遠遠地站著，連舉目望天也不敢，只捶著胸，說：『神啊，開恩可憐我這個罪人！』 \end{tabularx} \\ \\ \relax
18:14 & \begin{tabularx}{0.7\textwidth}{X} 我告訴你們，這人回家去比那人倒算為義了。因為凡自高的，必降為卑；自甘卑微的，必升為高。」 \end{tabularx} \\ \\
[1ex]
\hline
\hline
\end{longtable}
講題:「上帝啊,開恩可憐我這個罪人!」

經文:路加福音18章9至14節

講員:高國雄牧師

日期:2022年10月23日

9耶穌向那些仗著自己是義人,藐視別人的,設一個比喻,

10說:「有兩個人上殿裡去禱告:一個是法利賽人,一個是稅吏.

11法利賽人站著,自言自語地禱告說:『上帝啊,我感謝你,我不像別人勒索,不義,姦淫,也不像這個稅吏.

12我一個禮拜禁食兩次,凡我所得的都捐上十分之一.』

13那稅吏遠遠地站著,連舉目望天也不敢,只捶著胸說:『上帝啊,開恩可憐我這個罪人!』

14我告訴你們,這人回家去比那人倒算為義了；因為,凡自高的,必降為卑；自卑的,必升為高.」

請留意:比喻的對象是認識上帝,會上聖殿禱告,參與宗教活動的「信徒」

法利賽:這個名詞源於希伯來語的,意思是「分離」,指一些為保持純潔而與俗世保持距離,願意分別為聖的人,給我們的印象是嚴守律法/規約的宗教領袖,也會被認為是宗教的偽善者.

稅吏:代羅馬政府向猶太人收取稅捐,被其他猶太人視為叛徒,以及榨取百姓金錢的人,會被禁止參加宗教活動.

「上帝啊,我感謝你!」一句信徒在禱告時最常講得說話,你會為什麼發出「感謝」呢?為了「成功」?為了「順利」?為了「自我感覺良好」?

耶穌在這比喻中加入「站著,自言自語地禱告說」,真是非常大的諷刺,你的禱告只是「自言自語」嗎?

在殿裡站著,就應該是向上帝禱告的時候,正如我們上教會一樣,但為何只有「自言自語」呢?

這法利賽人的「感謝」內容:「我不像別人勒索,不義,姦淫,也不像這個稅吏.

我一個禮拜禁食兩次,凡我所得的都捐上十分之一.」

再翻譯出心底的說話,看似就是要公開宣告:「我不像別人一樣道德敗壞,我是個義人!」

和合本用了4次「我」這個字…

「我不像別人」:「別人」原文是「其餘的人」,有剩餘,餘下的意思,這描述為要特顯自己的優越.

耶穌又如何描述另一位主角呢?

V13那稅吏遠遠地站著,連舉目望天也不敢…

連行近,舉目也不敢,這是不是我們面向上主的心情及行動呢?

尋求上主,向主禱文,理應最重要的就是要先走近,再仰望；這看似是祈求的根本,也是我們的日常…但耶穌基督以稅吏作為主角,其心況及表現,又對我們有什麼提醒呢?

V13只捶著胸說:『上帝啊,開恩可憐我這個罪人!』

認清自己的本相,角色及位置,承認自己的限制,過犯,真心的認罪悔改,這才是真正走近上主的開始…

當你願意認罪悔改時,你就是與上主最親近的時候.(當耶穌呼召西門彼得時,西門發現自己是罪人,請耶穌離開他,耶穌卻呼召他同行)

一個簡單,真誠,信仰中最重要的禱告:

『上帝啊,開恩可憐我這個罪人!』

是信靠的開始,也是蒙恩的時候.

請回想耶穌釘十架時的,身旁的強盜對主的祈求…

v14「這人回家去比那人倒算為義了」

「家」,地方/建築物,也可以是指「聖殿/聖所」,就是那一個接納自己的地方及群體,也是上帝的管治與國度內.

願意悔改的「罪人」

vs

自以為是的「義人」

我們全都是「罪人」,生活在罪的世界中,是軟弱,不能自救的,這就是始祖犯罪後,「人」的根本.

我們誰都可以因著「信」而被稱為「義」,因為我們都可以成為蒙恩的「罪人」,因著上主的信實,與及我們的認信.

耶穌用這比喻的最重要的目的是什麼呢?

V9「耶穌向那些仗著自己是義人,藐視別人的,設一個比喻…」

這世界很多人自以為義,在科技網絡中更顯明出來,我們不難看到很多人在網上成為評論家,以不斷指出他人的錯處來抬高自己,以為自己站在「道得高地」,甚至連事情的根本也沒有攬清楚,就先去評論(其實是論斷).

來到教會,參與宗教活動的人,可能更甚,拿起聖經,教會的傳統規約等,就好像手執「上方寶劍」一樣,甚至自以為自己能夠代表上帝,去審判世界…

V14因為,凡自高的,必降為卑；自卑的,必升為高.

在上帝的眼中,祂所關心的與我們不一樣

路加福音14同樣使用這一句

耶穌教導赴宴要不揀首位,要揀末位,被邀請上座而得光彩…

馬太福音23也出現這經文

描述文士及法利賽人能說不能行(講一套做一套),又喜歡在宴會中坐首席及被人稱呼為老師,所以耶穌以此例子來教導,誰要為大,就要作眾人的用人…

被上主看重的人,真正在信仰中的謙卑,是由清楚自己的本相及位份開始…

路加福音的作者在這裡不單是記述耶穌基督的教訓及道理,而是要先劇透,為19章要出場的撒該作引子,因為耶穌即將要進入他的家,今天主要進入你的家,你預備好了嗎?



\newpage

\section{何明華贈予崇基的三份禮物}
\label{sec:621}
\textbf{吳梓明教授 - 何明華贈予崇基的三份禮物}
\newline
\newline
\hyperref[sec:620]{< < < PREV SERMON < < <}
~
\hyperref[sec:toc]{[返目錄]}
~
\hyperref[sec:622]{> > > NEXT SERMON > > >}


約拿書 3:1-10
\newline
\begin{longtable}{cl}
\hline
\hline
章節 & 經文 (和合本修訂版)\\
\hline
3:1 & \begin{tabularx}{0.7\textwidth}{X} 耶和華的話第二次臨到約拿，說： \end{tabularx} \\ \\ \relax
3:2 & \begin{tabularx}{0.7\textwidth}{X} 「起來，到尼尼微大城去，把我告訴你的信息向其中的居民宣告。」 \end{tabularx} \\ \\ \relax
3:3 & \begin{tabularx}{0.7\textwidth}{X} 約拿就照耶和華的話起來，到尼尼微去。尼尼微是一座極大的城，約有三天的路程。 \end{tabularx} \\ \\ \relax
3:4 & \begin{tabularx}{0.7\textwidth}{X} 約拿進城，走了一天，宣告說：「再過四十天，尼尼微要傾覆了！」 \end{tabularx} \\ \\ \relax
3:5 & \begin{tabularx}{0.7\textwidth}{X} 尼尼微人就信服神，宣告禁食，從最大的到最小的都穿上麻衣。 \end{tabularx} \\ \\ \relax
3:6 & \begin{tabularx}{0.7\textwidth}{X} 這消息傳到尼尼微王那裡，他就從寶座起來，脫下朝服，披上麻布，坐在灰中。 \end{tabularx} \\ \\ \relax
3:7 & \begin{tabularx}{0.7\textwidth}{X} 他叫人通告尼尼微全城，說：「王和大臣有令，人、畜、牛、羊都不可嘗任何東西，不可吃，也不可喝水。 \end{tabularx} \\ \\ \relax
3:8 & \begin{tabularx}{0.7\textwidth}{X} 人與牲畜都要披上麻布，切切求告神。各人要回轉離開惡道，離棄自己掌中的殘暴。 \end{tabularx} \\ \\ \relax
3:9 & \begin{tabularx}{0.7\textwidth}{X} 誰知道神也許會回心轉意，不發烈怒，使我們不致滅亡。」 \end{tabularx} \\ \\ \relax
3:10 & \begin{tabularx}{0.7\textwidth}{X} 神察看他們的行為，見他們離開惡道，神就改變心意，原先所說要降與他們的災難，他不降了。 \end{tabularx} \\ \\
[1ex]
\hline
\hline
\end{longtable}
詩篇 90:1-6,12-17
\newline
\begin{longtable}{cl}
\hline
\hline
章節 & 經文 (和合本修訂版)\\
\hline
90:1 & \begin{tabularx}{0.7\textwidth}{X} 主啊，你世世代代作我們的居所。 \end{tabularx} \\ \\ \relax
90:2 & \begin{tabularx}{0.7\textwidth}{X} 諸山未曾生出，地與世界你未曾造成，從亙古到永遠，你是神。 \end{tabularx} \\ \\ \relax
90:3 & \begin{tabularx}{0.7\textwidth}{X} 你使人歸於塵土，說：「世人哪，你們要歸回。」 \end{tabularx} \\ \\ \relax
90:4 & \begin{tabularx}{0.7\textwidth}{X} 在你看來，千年如已過的昨日，又如夜間的一更。 \end{tabularx} \\ \\ \relax
90:5 & \begin{tabularx}{0.7\textwidth}{X} 你叫他們如水沖去，他們如睡一覺。早晨，他們如生長的草； \end{tabularx} \\ \\ \relax
90:6 & \begin{tabularx}{0.7\textwidth}{X} 早晨發芽生長，晚上割下枯乾。 \end{tabularx} \\ \\ \relax
90:7 & \begin{tabularx}{0.7\textwidth}{X} 我們因你的怒氣而消滅，因你的憤怒而驚惶。 \end{tabularx} \\ \\ \relax
90:8 & \begin{tabularx}{0.7\textwidth}{X} 你將我們的罪孽擺在你面前，將我們的隱惡擺在你面光之中。 \end{tabularx} \\ \\ \relax
90:9 & \begin{tabularx}{0.7\textwidth}{X} 我們經過的日子，都在你震怒之下，我們度盡的年歲，好像一聲嘆息。 \end{tabularx} \\ \\ \relax
90:10 & \begin{tabularx}{0.7\textwidth}{X} 我們一生的年日是七十歲，若是強壯可到八十歲；但其中所矜誇的不過是勞苦愁煩，轉眼即逝，我們便如飛而去。 \end{tabularx} \\ \\ \relax
90:11 & \begin{tabularx}{0.7\textwidth}{X} 誰曉得你怒氣的權勢？誰因著敬畏你而曉得你的憤怒呢？ \end{tabularx} \\ \\ \relax
90:12 & \begin{tabularx}{0.7\textwidth}{X} 求你指教我們怎樣數算自己的日子，好叫我們得著智慧的心。 \end{tabularx} \\ \\ \relax
90:13 & \begin{tabularx}{0.7\textwidth}{X} 耶和華啊，我們要等到幾時呢？求你轉回，憐憫你的僕人們。 \end{tabularx} \\ \\ \relax
90:14 & \begin{tabularx}{0.7\textwidth}{X} 求你使我們早早飽得你的慈愛，好叫我們一生一世歡呼喜樂。 \end{tabularx} \\ \\ \relax
90:15 & \begin{tabularx}{0.7\textwidth}{X} 求你照著你使我們受苦的日子，和我們遭難的年歲，使我們喜樂。 \end{tabularx} \\ \\ \relax
90:16 & \begin{tabularx}{0.7\textwidth}{X} 願你的作為向你僕人們顯現，願你的榮耀向他們子孫顯明。 \end{tabularx} \\ \\ \relax
90:17 & \begin{tabularx}{0.7\textwidth}{X} 願主－我們神的恩寵歸於我們身上。願你堅立我們手所做的工，我們手所做的工，願你堅立。 \end{tabularx} \\ \\
[1ex]
\hline
\hline
\end{longtable}
哥林多前書 13:4-12
\newline
\begin{longtable}{cl}
\hline
\hline
章節 & 經文 (和合本修訂版)\\
\hline
13:4 & \begin{tabularx}{0.7\textwidth}{X} 愛是恆久忍耐；又有恩慈；愛是不嫉妒；愛是不自誇，不張狂， \end{tabularx} \\ \\ \relax
13:5 & \begin{tabularx}{0.7\textwidth}{X} 不做害羞的事，不求自己的益處，不輕易發怒，不計算人的惡， \end{tabularx} \\ \\ \relax
13:6 & \begin{tabularx}{0.7\textwidth}{X} 不喜歡不義，只喜歡真理； \end{tabularx} \\ \\ \relax
13:7 & \begin{tabularx}{0.7\textwidth}{X} 凡事包容，凡事相信，凡事盼望，凡事忍耐。 \end{tabularx} \\ \\ \relax
13:8 & \begin{tabularx}{0.7\textwidth}{X} 愛是永不止息。先知講道之能終必歸於無有；說方言之能終必停止；知識也終必歸於無有。 \end{tabularx} \\ \\ \relax
13:9 & \begin{tabularx}{0.7\textwidth}{X} 我們現在所知道的有限，先知所講的也有限， \end{tabularx} \\ \\ \relax
13:10 & \begin{tabularx}{0.7\textwidth}{X} 等那完全的來到，這有限的必消逝。 \end{tabularx} \\ \\ \relax
13:11 & \begin{tabularx}{0.7\textwidth}{X} 我作孩子的時候，說話像孩子，心思像孩子，意念像孩子；既長大成人，就把孩子的事丟棄了。 \end{tabularx} \\ \\ \relax
13:12 & \begin{tabularx}{0.7\textwidth}{X} 我們現在是對著鏡子觀看，模糊不清；到那時，就要面對面了。我如今所認識的有限，到那時就全認識，如同主認識我一樣。 \end{tabularx} \\ \\
[1ex]
\hline
\hline
\end{longtable}
路加福音 10:30-35
\newline
\begin{longtable}{cl}
\hline
\hline
章節 & 經文 (和合本修訂版)\\
\hline
10:30 & \begin{tabularx}{0.7\textwidth}{X} 耶穌回答：「有一個人從耶路撒冷下耶利哥去，落在強盜手中。他們剝去他的衣裳，把他打個半死，丟下他走了。 \end{tabularx} \\ \\ \relax
10:31 & \begin{tabularx}{0.7\textwidth}{X} 偶然有一個祭司從那條路下來，看見他就從另一邊過去了。 \end{tabularx} \\ \\ \relax
10:32 & \begin{tabularx}{0.7\textwidth}{X} 又有一個利未人來到那裡，看見他，也照樣從另一邊過去了。 \end{tabularx} \\ \\ \relax
10:33 & \begin{tabularx}{0.7\textwidth}{X} 可是，有一個撒瑪利亞人路過那裡，看見他就動了慈心， \end{tabularx} \\ \\ \relax
10:34 & \begin{tabularx}{0.7\textwidth}{X} 上前用油和酒倒在他的傷處，包裹好了，扶他騎上自己的牲口，帶他到旅店裡去，照應他。 \end{tabularx} \\ \\ \relax
10:35 & \begin{tabularx}{0.7\textwidth}{X} 第二天，他拿出兩個銀幣來，交給店主，說：『請你照應他，額外的費用，我回來時會還你。』 \end{tabularx} \\ \\
[1ex]
\hline
\hline
\end{longtable}
講題:何明華贈予崇基的三份禮物

經文:約拿書3:1-10；詩篇90:1-6,12-17；哥林多前書13:4-12；路加福音10:30-35

講員:吳梓明教授

日期:2022年10月30日

各位早晨!多謝高牧師的邀請,讓我再有機會與崇基禮拜堂眾弟兄姊妹一同聚集,一同敬拜上帝.

祝願上主的恩惠及慈愛常與你們同在!

去年因為疫情的緣故,打擾了我們很多活動,但也讓我們可以用兩年時間來慶祝崇基學院創校70周年的紀念.去年在高牧師的領導下,我們設計了一個十分獨特的校史考察計劃,出版了一本「香江薪傳 – 中國基督教大學在香港的承傳與蛻變」(梁智仁等,2022),大家有興趣可以向高牧師索取,今年我們也繼續與崇基的師生校友們一同以「何明華與崇基」為題,探索崇基學院的先賢為崇基存留了怎麼樣的傳統和教育使命,所以今天我亦想與大家分享這個題目:「何明華會督贈予崇基的三份禮物」.

何明華會督(或稱為主教)是崇基的創校先賢之一,甚至很多老前輩如李應林校長,謝昭杰先生等也確認,說:「崇基學院得以…創立,何明華會督實在居功至偉」.(李應林,1951,頁1-2)

讓我今日簡單地介紹何明華會督贈送給崇基學院的三份禮物,至今我們仍是念念不忘的.第一件是何會督為崇基設計的「人生哲學」課程.

「人生哲學」– 靈性的關注與修為

創校之初,何會督倣效國內基督教大學的經驗,擬定「人生哲學」課程為崇基學生的必修課程之一,他亦特別安排在一年級第一個學期有三次親臨授課,還有四年級的「柏拉圖」哲學課,也是由他教授的,由此可見何會督對「人生哲學」課程的重視.(謝昭杰,1951,頁2-3及7；詹肇凱,1952,頁2)

我的主修是「宗教教育」,所以我特別關注「教育」,也特別欣賞何會督在他繁忙的時間中還能抽出時間親自教授「人生哲學」和「柏拉圖」(一年級及四年級)的課,這對學生的成長,特別在基礎教育上,確是十分重要的.最近我聽聞一位我們宗教系畢業的校友在一所聖公會中學擔任校長多年,今年八月她接受《教聲》訪問時亦透露自己擔任校長時仍是負責教授中一的宗教科,目的是要認識新同學,建立友好互信的關係,這樣,六年後她便可以說自己是親身教過全校的學生了.(《教聲》,2022,頁4-5)

我很欣賞這樣的教學態度,這位校長就像何會督一樣,也承傳著我們崇基獨特的教育精神,不單是盡心地教導我們的學生,更是很重視低年級學生所需要的基礎教育.

何會督的「人生哲學」,顧名思義,就是直指人生的哲學,它是從哲學(理性)的角度談論人生,思考人生的問題.何會督認為「人生哲學」就是一種「對話」(dialogue)

和「會談」(conversation),他在崇基學院第一屆畢業崇拜時對學生說:『崇基學院…的「人生哲學」…是要(讓)每一個學生和歷代中西哲人的會談,也與現代思潮和不同的人生觀對話』.[2]

這種「對話」和「會談」的教學模式成為了今日崇基學院

(不論是形式或非形式的)

通識教育課程中一種很獨特的教學模式,是每一個崇基學生(特別是我們校友們可以作見證說),是我們一生受用的.原來這種「對話」和「會談」的教學模式,早在1950年代的崇基學院已經出現了.我們要多謝何會督!

我最敬愛的老師之一是沈宣仁教授(1931-2004),他是在1962年加入崇基學院宗教及哲學系(簡稱「宗哲系」)[3]

任教,他當時亦致力於協助學院在既有的「人生哲學」課程的基礎上,制定「綜合基本課程」(Integrative Basic Studies,簡稱為「IBS」),[4]

1967年獲學院通過,1968年9月正式推行.我有幸在當年進入崇基學院,成為第一批接受IBS通識教育課程的學生,也見證到IBS是如何從「人生哲學」課程演進過來.沈教授也像何會督一樣,他願意親自教授一年級的「大學之理念」和「思想方法」等基礎課程.[5]

崇基學院在2016年成立「宣仁通識教育資源中心」(中心設在崇基圖書館旁邊),並推崇沈教授為「崇基學院通識教育之父」.若果沈教授仍是活著的話,我相信他會說:假如我(沈教授)是崇基學院通識教育之父,創始人(我們的祖父)應該是何明華會督.沈教授曾經講過一個親身的故事,原來在1953年他還未遠赴美國深造之前,他曾親自造訪何明華,在會督府詳談他的未來計劃,他們又在祈禱室一起跪下禱告,沈教授憶述說那次的禱告改變了他的一生.[6]

沈教授大學的時候是主修英文,副修哲學,後來在美國他轉讀西方哲學及神學,又在歐柏林學院和芝加哥大學學會了甚麼是博雅教育,回到崇基後便開始教授宗教,哲學和通識教育等課程,更是將何明華的人生哲學提升至「綜合基本課程」(IBS).沈教授亦坦言說他是深受何明華會督的感染和影響的,所以我們可以說:「何明華是崇基學院通識教育之先祖」.

「人生哲學」課程有另一個特色,是讓學生可以接觸基督教的新舊約聖經.《聖經》不單是代表西方文化,更重要的是引導我們思考人生和靈性的問題.或許有人會問:「靈性」到底是甚麼?簡單來說,靈性就是關乎人生的思考,更是尋求超越人生的局限.一般人相信「全人教育」是「德,智,體,群,美 - 五育」,而靈性教育是屬於「第六育」.北大蔡元培校長曾主張要「以美育代宗教教育」,但(崇基學院第二任院長)凌道揚院長卻有另一種看法,他說:「崇基的(教育)主旨是德,靈,智,體,群五育」.(凌道揚,1955,頁3-4)換句話說,凌院長是主張以「靈性教育」代替「美育」的,或者更清楚些,應是在知識之外,必須加上靈性的思維.凌院長也曾說過:「崇基…的教育方針,不僅是要使同學在知識方面獲得高度的成就,同時,也要使每一個同學,在靈性方面獲得充分的修養.」換句話說,知識和靈性兩方面配合並進,才能達到我們校訓所說『止於至善』的境界.(凌道揚,1959,頁1)去年在崇基學院70周年校慶時,我們的校董會主席陳德霖博士提出了12字真言,即是說:「(崇基的教育主旨是要學生達到)

文武雙全,德才兼備,靈性修為」(鍾啟然等,2021,頁1-7)[7],正好就是承傳了何明華會督的「人生哲學」和凌道揚院長所講的教育方針,突顯出崇基學院通識教育的獨特精神,是著重「全人教育」,兼顧學生對「人生」與「靈性」的思考和修為,也不要忘記創校先賢留給我們的教育主旨,是要「在知識之外,加上靈性的思維」.這就是何明華給我們「人生哲學」的第一份禮物.

服務精神 – 培訓謙卑的學習

何明華給我們的第二份禮物是「服務精神」.「服務學習」不單是崇基學院的教育特色之一,其實也是何會督留給崇基另一份珍貴的禮物.

何會督是聖公會維多利亞教區的主教,但最為人敬重的是他對香港社會戰後重建的貢獻.除了在1951年與歐偉國先生,李應林校長一起創建崇基學院之外,何會督亦創立聖雅各福群會,將「福群會」(“settlement”=即是要進入服務的社群)的概念帶進香港,何會督也在紅磡創立「聖匠職業訓練學校」和「聖匠社區中心」(為同一個社群提供多樣化的服務,也是一種創新的思維).柴灣的「明華大廈」是以「明華」為名,亦是因為應用了他另一個創新的「社區屋村」(“communal estate”)的概念.另一個例子是「何文田村」.還有,如「童膳會」(提供學童午膳),「香港小童群益會」,「航海學校」,「勞工子弟學校」等,這些都是何明華會督積極回應香港戰後重建所作出的貢獻.

事實上,更重要的貢獻是何會督將這種基督教「社會服務」的精神注入在崇基的教學課程之中.大家有否留意?

崇基是香港最早成立「社會服務(社工)系」的專上學院之一,崇基的課程設計最具特色的是它將「宗教教育」與「社會工作」結合(共融)在一起.自1956年,崇基已經設有「宗教教育及社會服務學系」,提供了「宗教教育及社會工作文憑課程」,教師包括芮陶菴博士(Dr. Andrew T. Roy),馬鴻述教授,黃壽林先生等.[8]

這種將「宗教教育與社工訓練」融洽在一起是當時香港社會所獨有的.剛於前年(2020年8月31日)我們的大師兄溫漢璋先生(84歲)離世,他就是「宗教教育及社會工作文憑課程」第一屆1959年的畢業生,他的一生的確是活出了「崇基服務精神」一個最佳美的模範.[9]原來何明華在崇基學院第一屆的畢業崇拜中已經教導了我們的大師兄,大師姐甚麼是崇基最具特色的「服務精神」.何會督說:一般人以為基督徒生活就是「服務的生活」,其實是錯誤的,他說基督徒生活應該是「友誼的生活」.因為耶穌說:「以後我不再稱你們為僕人,……我乃稱你們為朋友」.(約翰十五章十五節)

(「僕人」重點在服務,「朋友」是生命的交流)何會督說:「友誼」(friendship) 是包括「會談」(conversation)- 是人與人之間心靈的「對話」(dialogue),「服務」不一定會含有「生命的分享」或「心靈對話」的成份,有時反而會變成為權力與傲慢的工具.(何明華,1955b,頁45)所以何會督認為,有基督精神的「服務」必須是發自內心,有「生命的分享」和「心靈的對話」.

當年的校牧龐德明牧師也曾解釋,說:『崇基學院作為一所基督教大學,是負有三個主要的目的:

其一是研究高深的學問；

其二是發展高貴的品德；其三是發展服務的精神.』他繼續解釋說:『這三個目的不是基督教大學所獨有的,任何一間有價值的大學都是要達到這些目的的.但在尋求達到這些教育目的時,基督教大學卻表現出它有獨特的方法和態度來,即是尋求真理時要有「信仰」,發展道德時要有「謙卑」,服務時要有「愛心」』.(龐德明,1954,11-12)

這也就是何會督建議我們要將「服務」變成「友誼」,要用愛心活出崇基獨具特色的「服務」精神來.

今年五月,我們訪問了一位老校友-

陳佐才法政牧師.為什麼要訪問陳佐才法政牧師?他不單只是我們的大師兄(1958年英文系畢業),原來他中學時已認識何明華,是何明華引導他選擇進入崇基讀書的,後來按立他做牧師,更教導他認識甚麼是崇基所提倡的服務.當訪問陳法政時,他說:『何會督相信的上帝是一位「樂意擺上自己的上帝」(Self-giving

God);聖經說:「上帝愛世人,甚至將祂的獨生子賜給他們,叫一切信祂的,不至滅亡,反得永生」.(約3:16)「Self-giving」常被譯作「捨己,犧牲,或是否定自我」,但何會督認為Self-giving並沒有「否定自己」的意思,Self-giving是「giving yourself = 將你所有的,好的東西擺上」』.「服務也是這樣」,是「發自內心,願意將你所有的美好的東西拿出來」,與有需要的人分享.這就是陳佐才法政牧師與我們的分享.

有些同學會問:「服務學習」為什麼需要有愛心的呢?原來「愛心的服務」是一種很特別的,另類的學習.

我們必須要認識我們服務的對象不僅是需要金錢或是物質上的東西,更也是需要受到別人的尊重,接納,和愛護.

今天我們所讀到的福音經課是一個大家都會熟識的故事,就是耶穌所講的「一個善撒瑪利亞人」的故事.耶穌說:有一個人從耶路撒冷落去耶利哥去.途中遇到強盜,他們剝去他的衣服,把他打個半死,並丟下他走了.當時有一祭司從這條路下來,看見他就轉向另一邊走了.又有一個利未人,來到這個地方,也照樣從另一邊過去了.唯有一個撒瑪利亞人,來到這裡,看見他,就動了慈心,上前用油和酒倒在他的傷處,包裹好了,扶他騎上自己的牲口,帶到店裡去照料他.第二天,他拿出兩錢銀子來,交給店主說:

你且照料他,若需額外費用,我回來必定還你.…

奇怪嗎?這個善撒瑪利亞人並沒有告訴那店主他是誰,他沒有借此機會宣揚自己,也沒有期望受傷者將來會給他甚麼回報,起碼會記得他是誰,將來有機會要報答他等.這個就是何明華會督,他將自己所有的給予崇基,他的「服務」並不是為炫耀自己,也不是期望將來得到甚麼回報,或是要建立自己的名聲,權位,而只是把自己所有的(能力所及的),用愛心去回應身邊有需要的人,這就是何會督「服務精神」的一個典範.

2013年,在梁元生教授(院長)領導下,崇基學院創立了中大唯一的一所「服務學習中心」.(《崇基學院概覽》,2013,頁99)梁院長的在中心開幕典禮上亦重新確認「服務學習」是國內13所基督教大學的辦學傳統之一,他引用燕京大學的校訓:「因真理,得自由,以服務」,說明崇基是繼承著像燕京大學等13所基督教大學同樣的教育使命.(麥嘉輝,2014；許開明,2019)陶兆銘教授擔任中心主任的時候,也曾借用《新約聖經》《哥林多前書》13:13所說:『如今常存的有信,有望,有愛這三樣,其中最大的是愛』,說明崇基的服務學習中心亦是承傳著基督的博愛精神,實踐於服務學習之中.(梁智仁等,2022,頁121)

大家或許也知道: 今個學年(2022-2023)中文大學開始推行全體學生均須參與服務學習計劃(Service Learning Program),現在中大所有學生於2022-2025(四年內)學年間必須修讀一個學分的服務學習課程.或許有人會認為學分太少,但亦聊勝於無!其實重點不在於學分多少,時間多少,而是「服務學習」背後的「精神」,我們能學會多少.我相信這個計劃更重要的意義,是讓我們崇基學院的師生不可忘記何明華會督所贈予我們的第二份禮物 -要我們學習「活出謙卑,獻出愛心,擺上自己的服務精神」.第三份禮物是甚麼?

大家估到嗎?

何明華贈予我們的第三份禮物是「崇基的校徽」.

崇基校徽 – 中西文化的共融

近日很多校友討論中大的校徽,他們特別關注中大校徽所表達出創校先賢的教育使命和遠象.

我們作為崇基人,其實也應該同樣地關注我們「崇基的校徽」.原來崇基校徽是何明華會督賜給我們一份非常寶貴的禮物.這個校徽也是表達了崇基的創校先賢留給我們重要的遠象和教育使命.

創校之時,何會督悉心設計了這個「崇基校徽」,這個校徽現時(70年後的今天)仍然是高高掛在中環會督府旁的一棟高樓之上,這棟高樓叫做「霍約瑟會督紀念大樓」(大家可以參看我們今年6月拍到的照片,若有機會去中環的時候,不妨花些時間去看看,已經有70年的歷史了),這棟大樓是何明華在1952開始建造,目的為借給崇基上課之用,直至崇基學院找到永久的校園,即是崇基搬入來馬料水之後,才歸還給聖公會的.(所以話:崇基的創校,何會督真是居功至偉!)

我們對「崇基校徽」認識有多少呢?很多崇基學生,校友們只知道我們的校徽是選自唐代《大秦景教流行中國碑》上的徽號,但他們卻不知道為什麼何明華選用這個徽號作為崇基的校徽.何會督曾經就崇基校徽作過這樣的描述,他說:「(大秦景教)碑上的十字架下採用了中國兩大宗教的象徵:『蓮花的清淨』和『雲形的神秘』…

我希望這從『蓮花』浮現的十字架和從『雲』中吐露的火焰,會成為崇基學院的校徽」.(何明華,1955,5)[10]

原來『蓮花』和『浮雲』是代表(或是融合)中國的兩大宗教 – 佛教和道教的象徵,原來何會督是相信基督教與中國宗教文化兩者不是對立的,更是可以共生共存,融合為一的,他亦是希望崇基學院能夠肩負起這份教育使命,不僅是溝通中西文化的交流,更是能夠致力達至中西文化的共融.

今天的舊約經課提醒我們「上帝愛世人」,祂的慈愛遍及全世界,各民各族,上帝甚至差派先知約拿往尼尼微城(尼尼微城人是代表外邦人,在猶太人眼中是不能被接納的外邦人,但上帝卻要差約拿去尼尼微城),要宣示上帝「救恩的普世性」.何會督在崇基的校徽上擺上唐代《大秦景教》十架的徽號,正是提醒我們基督教的福音已經在唐朝傳入中國,上帝的愛早已經在唐代臨到中國.何會督曾被稱為是「紅色主教」,或許有很多人會批評他,但我相信何會督是不會介意的,因為何會督知道是上帝遣派他來中國,來香港,要他學習尊重中國的歷史文化,他不單只是愛我們香港,更也是愛我們中國人,包括國內13億的中國人.

原來何明華是特別重視宗教間或中西文化間的對話與交流的,尤其是基督教與中國宗教文化的共融,何會督曾經說過:「人生就是人與『上帝』,人與『經驗』,人與『人』不斷的會談(對話,在對話當中),你要有可堅持的見解,然後他人才能和你會談.同樣,你也必須尊重對方的見解,虛懷若谷地去聆聽他人的意見,然後會談才有實現的可能」.(何明華,1955b,頁44)這是何明華會督在1955年說的.

在會談或對話的過程中,我們必須要學習謙卑,正如今日所念的哥林多前書中保羅所說的,保羅說:「我如今所知道的有限,到那時候就全知道,如同主知道我一樣」(哥林多前書13:12).同樣地,「十架」與「蓮花」和「雲火」只是象徵著不同文化對上帝,對人生的局部理解,三者並不衝突,還是可以互為補足的.崇基作為一所在中國的基督教高等學府,不僅是可以發揚基督十架的精神,我們更需要學習謙卑,尊重中國的文化,努力達至中西文化的和諧與共融.

我在崇基讀書的時候便認識容啟東校長(1960-1975年出任崇基學院校長一職);

原來他在就職演講詞中也有談及崇基學院的教育使命和兩大特色,其中第一個特色是:「東西文化的交流」.(容啟東,1960,頁6)

他也是認同何會督的看法,希望我們透過崇基學院所提供的基督教高等教育,能夠探索中西文化的不同,可以進一步達至如校歌所說「(我們要尋找)中西(的)結晶,增益(世界的)文明」.今日當我們看見世界上國與國之間的衝突,戰爭,不同的文化與文明或意識形態的對立時,我會更加懷念昔日何明華會督的睿智.原來何會督是透過崇基學院校徽的設計,贈送予崇基的師生們一種項崇高的教育使命,要我們透過基督教高等教育的實踐,達至「中西文化的共融,增益世界文明」.這就是何會督給我們崇基學院的第三份禮物 –「尋求中西文化共融的教育使命」,多謝何會督!

今日當我們慶祝崇基學院71周年校慶之時,也讓我們謹記何會督贈予我們的三份禮物.盼望我們能夠繼續傳承何會督為我們樹立這三個崇基的優良傳統,大家記得嗎?

分別是:

(一)何會督的人生哲學課程,「在知識之外,加上靈性的修為」；(二)崇基學院服務學習的傳統:讓我們一起「活出謙卑,獻出愛心,擺上自己」的服務學習精神來；(三)校徽:「尋求中西文化的共融,達至增益世界文明的教育使命」.

何會督是在1975年4月21日離世.翌日的南華早報有這樣一段的描述:「何明華(會督)…

不單建立了一個教區,更締造了香港社會的良心…

他的夢想是一個仁愛的社會,…

藉著教會謙卑的服侍,彰顯基督的愛在現代的社會中.」

崇基的校名是帶有「崇奉基督」的教育使命,「崇奉基督」,就是要活出,彰顯基督的愛在現代的社會中.何會督不單是贈予我們「人生哲學」的課程,教導我們如何與中西哲人,與中西文化中不同的世界觀,人生觀會談和對話.他也是以身作則地「活出謙卑,獻出愛心,擺上自己」的獨具崇基特色的服務學習精神!更重要的是何會督也贈予我們必須致力尋求達至中西文化溝通和共融的教育使命.

祈求上主祝福我們,幫助我們延續崇基先賢傳承給我們的教育使命,讓我們一起低頭祈禱說:「主啊!求你堅立何明華會督和崇基創校先賢們手所作的工!也幫助我們繼續承傳他們留給我們的遠象和教育使命；願主也堅立我們手所作的工,我們手所作的工,願你堅立」(詩篇90:17).啊們!



\newpage

\section{女人永遠只是男人的某某}
\label{sec:622}
\textbf{林豪恩先生 - 女人永遠只是男人的某某?}
\newline
\newline
\hyperref[sec:621]{< < < PREV SERMON < < <}
~
\hyperref[sec:toc]{[返目錄]}
~
\hyperref[sec:623]{> > > NEXT SERMON > > >}


路加福音 20:27-38
\newline
\begin{longtable}{cl}
\hline
\hline
章節 & 經文 (和合本修訂版)\\
\hline
20:27 & \begin{tabularx}{0.7\textwidth}{X} 有些撒都該人來見耶穌。他們說沒有復活這回事，於是問耶穌： \end{tabularx} \\ \\ \relax
20:28 & \begin{tabularx}{0.7\textwidth}{X} 「老師，摩西為我們寫下這話：『某人的哥哥若死了，有妻無子，他該娶哥哥的妻子，為哥哥生子立後。』 \end{tabularx} \\ \\ \relax
20:29 & \begin{tabularx}{0.7\textwidth}{X} 那麼，有兄弟七人，第一個娶了妻，沒有孩子死了。 \end{tabularx} \\ \\ \relax
20:30 & \begin{tabularx}{0.7\textwidth}{X} 第二個、 \end{tabularx} \\ \\ \relax
20:31 & \begin{tabularx}{0.7\textwidth}{X} 第三個也娶過她；同樣地，七個人都娶過她，沒有留下孩子就死了。 \end{tabularx} \\ \\ \relax
20:32 & \begin{tabularx}{0.7\textwidth}{X} 後來，那婦人也死了。 \end{tabularx} \\ \\ \relax
20:33 & \begin{tabularx}{0.7\textwidth}{X} 那麼，在復活的時候，那婦人是哪一個的妻子呢？因為他們七個人都娶過她。」 \end{tabularx} \\ \\ \relax
20:34 & \begin{tabularx}{0.7\textwidth}{X} 耶穌對他們說：「這世代的人有娶有嫁， \end{tabularx} \\ \\ \relax
20:35 & \begin{tabularx}{0.7\textwidth}{X} 惟有配得那要來的世代和從死人中復活的人不娶也不嫁。 \end{tabularx} \\ \\ \relax
20:36 & \begin{tabularx}{0.7\textwidth}{X} 因為他們不能再死，和天使一樣；既然是復活的人，他們就是神的兒子。 \end{tabularx} \\ \\ \relax
20:37 & \begin{tabularx}{0.7\textwidth}{X} 至於死人復活，摩西在《荊棘篇》上就指明了，他稱主是亞伯拉罕的神，以撒的神，雅各的神。 \end{tabularx} \\ \\ \relax
20:38 & \begin{tabularx}{0.7\textwidth}{X} 神不是死人的神，而是活人的神，因為對他來說，人都是活的。」 \end{tabularx} \\ \\
[1ex]
\hline
\hline
\end{longtable}
講題:女人永遠只是男人的某某?

經文:路加福音20章27至38節

講員:林豪恩先生

日期:2022年11月6日

一．引言

2022年9月13日,伊朗首都德黑蘭22歲女子阿米尼(Mahsa Amini)因為違反頭巾法而被道德警察逮捕.三日後,她在被拘留期間死亡,當局宣稱她的死亡與心臟病有關,但人們懷疑她的死另有原因.年輕女孩因為頭巾沒戴好就死去,引發了部分伊朗民眾再次關注女性的處境,他們上街要求「正義,自由,不要強制頭巾」,至今,事件仍在發展中.

在另外一個地方,2021年8月政權易手,新政府發佈了各種限制婦女自由的法令,例如禁止她們在沒有男性監護人的情況下從事政府工作,接受中等教育,以及旅行超過45英里(72公里).

這些要有男性監護人的限制都只是施加於女人,而沒有施加於男人,令人想起文獻記載古時對女性「三從四德」的要求.「三從」,指在家從父,出嫁從夫,夫死從子.從這些要求所來看,暗示女人是甚麼呢?女人必須是男人的某某嗎?耶穌時代的撒都該人似乎也懷著類似的觀念,在今天的福音經課中,我們就聽見他們問:「當復活的時候,她是哪一個的妻子呢?」撒都該人的觀念似乎認為,女人的身分是男人的妻子,在生是；如果有將來,將來都只有是.這樣的講法,不禁令人疑惑,難道女人永遠都只是男人的某某?

二．照顧還是擁有?

讓我們回到經文細讀撒都該人提問的背景及具體內容:

20:27撒都該人常說沒有復活的事.有幾個來問耶穌說:

20:28「夫子!摩西為我們寫著說:『人若有妻無子就死了,他兄弟當娶他的妻,為哥哥生子立後.』

20:29有弟兄七人,第一個娶了妻,沒有孩子死了；

20:30第二個,第三個也娶過她；

20:31那七個人都娶過她,沒有留下孩子就死了.

20:32後來婦人也死了.

20:33這樣,當復活的時候,她是哪一個的妻子呢?因為他們七個人都娶過她.」

撒都該人提出這個問題,是否真正關心那個女人呢?這也值得懷疑,因為經文提示撒都該人來問耶穌這個問題,用意是藉著這個問題指出復活的荒謬,證實他們否定復活的立場.不過,真心也好,假意也好,問題往往反映出問問題者的假設或信念.在撒都該人的眼中,似乎女人就只是某人的妻子,他們對女人的這種觀念是有根據的,他們所說的摩西律法記載在申命記25章:

25:5「弟兄同居,若死了一個,沒有兒子,死人的妻不可出嫁外人,她丈夫的兄弟當盡弟兄的本分,娶她為妻,與她同房.

25:6婦人生的長子必歸死兄的名下,免得他的名在以色列中塗抹了.

25:7那人若不願意娶他哥哥的妻,他哥哥的妻就要到城門長老那裡,說:『我丈夫的兄弟不肯在以色列中興起他哥哥的名字,不給我盡弟兄的本分.』

我們懷疑撒都該人的問題並不是出於關心那位女人,翻查摩西律法,相關的條文明言這種安排的目的是「25:6婦人生的長子必歸死兄的名下,免得他的名在以色列中塗抹了」,這就是「盡弟兄的本分」的意思.雖然條文沒有明言照顧寡嫂是目的,但在當時的社會處境中卻產生了照顧寡嫂的功能.

試想像在「在家從父,出嫁從夫,夫死從子」的古代社會,女性的身分依附於男性,弟娶寡嫂這種安排就給予了寡嫂身分,現實上產生了照顧寡嫂的效果.在古代社會,人們關注的是存傳宗接代,其實,在女人來說,她們盼望的豈不也是生孩子嗎?你還記得信心之父亞伯拉罕的妻子撒拉生下以撒之前最苦惱的是甚麼事嗎?對了,就是未能生孩子.你有沒有留意(撒母耳記上1:8)撒母耳的媽媽哈拿「哭泣,不吃飯,心裡愁悶」,所謂何事呢?就是因為未能生孩子.摩西律法寫著「她丈夫的兄弟當盡弟兄的本分,與她同房.婦人生的長子必歸死兄的名下」.當寡嫂成功生孩子,為亡夫續後,承接亡夫的田地資產,寡嫂生活就有所依靠.現實一點來想,在古代社會,對一個女人來說,你認為丈夫重要還是兒子重要呢?我們可以從撒母耳的媽媽哈拿的故事得到一點啟發.當哈拿「哭泣,不吃飯,心裡愁悶」的時候,她丈夫以利加拿對她說:「哈拿啊,你為何哭泣,不吃飯,心裡愁悶呢?有我不比十個兒子還好嗎?」你猜猜哈拿的丈夫這些甜言蜜語能否令哈拿重現笑顏呢?後來,哈拿終於笑逐顏開了,「哈拿禱告說:我的心因耶和華快樂；我的角因耶和華高舉.我的口向仇敵張開；我因耶和華的救恩歡欣.」(2:1)發生甚麼事呢?因為哈拿終於生了一個兒子,就是撒母耳!

摩西律法中亡夫的其中一位兄弟要娶寡嫂為妻這安排一方面是「盡弟兄的本分,為兄弟留後」,同時為那位喪夫的寡婦賦予身分,令她可以繼續在社會中正常地生活,這種安排發揮了家人之間互相照顧的功能.但撒都該人可能想得太遠了,當他們問:「這個女人是誰的妻子?」他們似乎把男人視為擁有者,把女人視為被擁有品.按著撒都該人對女人的觀念,這個女人被七個男人擁有過,她到底是屬於誰的呢?他們的問題令他們困惑,因為他們自己深陷於對女人的照顧責任和擁有權的迷思之中.相信沒有人否定人是需要互相照顧的,不單是女人需要被照顧,男人也需要被照顧呢!但照顧責任不等於擁有權,被照顧者也不等於附屬於照顧者.如果有人以照顧你之名,把你視為他的擁有物,要你依附於他,取消你的獨立自主,可能你會感到不妥當.若用我們現在的觀念來照顧寡嫂,還給她如男人一樣的財產擁有權就可以了.這樣,她就可以繼承家庭的資產,有資源繼續生活下去.

某些狀況之下,個人未必有足夠資源,因此,仍是需要互相照顧.例如不幸患重病要長期治療,治療費用從何而來呢?在某些地方,病人因此就傾家蕩產,或者負債累累.在香港的公立醫院,住院收費好像仍然是一百元一日,你知道為甚麼是一百元嗎?我曾經在一個醫療講座聽到,原來那一百元是一日三餐的大約成本.那麼,醫療費用呢?原來醫療是免費的.參加者都感到驚訝,為甚麼會是免費的呢?講者解釋說,當時香港跟隨英國的醫療政策,他們認為國家有責任照顧人民的健康,因此國家負起了人民的醫療費用,沒有人會因為沒有錢而得不到治療.錢從何來呢?當然是從人民而來.這樣的制度,讓人民參與互相照顧.人民負起了互相照顧的責任,然而每個人,無論是男還是女,不會因此歸附於另一個人,仍然是獨立及自主的.

三．「當下投射終末」還是「終末塑造當下」?

對於撒都該人的假設性問題,耶穌如何回應呢?

20:34耶穌說:「這世界的人有娶有嫁；

20:35惟有算為配得那世界,與從死裡復活的人也不娶也不嫁；

20:36因為他們不能再死,和天使一樣；既是復活的人,就為上帝的兒子.

20:37至於死人復活,摩西在荊棘篇上,稱主是亞伯拉罕的上帝,以撒的上帝,雅各的上帝,就指示明白了.

20:38上帝原不是死人的上帝,乃是活人的上帝；因為在他那裡,人都是活的.」

順著撒都該人的問題提及的在生狀況及復活狀況,耶穌用了「這世界」和「那世界」來區分.撒都該人本來就不相信有「那世界」,不過,如果真的有「那世界」的話,「那世界」會是怎樣的呢?按撒都該人的推想,「那世界」只不過是「這世界」的延續.如果,女人在「這世界」的身分是男人的妻子,在「那世界」都仍然只是男人的妻子.如此推想也難怪,人們都習慣以現在推想將來,在宗教上以人類社會狀態投射神靈世界.

民間信仰神明多不勝數,誰比較大,誰管什麼你搞得清楚嗎?有人認為,「玉皇大帝」原來不是最大,祂算是神界的CEO,其上還有董事長太上老君及董事會.關於土地問題,應該找誰呢?土地公相當於村里長,城隍爺是省市長,而地后則是國土部部長.農曆年底為甚麼要拜灶君呢?人間有賄賂官員的現象,人們對天上神明也行這一套.因為灶神監察人間善惡,並在年底稟報天庭.農曆年底送神日,人們很謹慎的祭拜灶君,供奉肉食,糖餅,茶酒,水果等等,還要把糕點,糖漿,糖粉之類的甜食塗在灶君圖像的嘴上,希望灶君不在玉帝那裡講自己家的壞話,以避免上天的責罰.

耶穌卻指出,「那世界」並不是「這世界」的延續,「這世界」和「那世界」不是時間上的不同,而是性質上的不同.耶穌指出,配得那世界,與從死裡復活的人,是不能再死的人,都是活著的人.「那世界」既然沒有死亡,就完全沒有兄弟死去留下「寡嫂」這現象；既然沒有死亡,也就不存在在生兄弟娶寡嫂為妻「為兄弟留後」這需要.「那世界」是一個不再一樣的世界,到底我們要用「這世界」去投射「那世界」,還是用「那世界」來指導「這世界」呢?即是說,我們要用暫時的當下去投射永恆,還是以永恆來塑造暫時的當下呢?

耶穌指出,「那世界」「從死裡復活的人也不娶也不嫁」,那麼,到時女人的身分如何定義呢?她不再是男人的女兒,不再是男人的妻子,不再是男人的母親,她是誰呢?這些問題,為甚麼不反過來問男人?男人的身分又如何定義呢?回答可以很簡單:男人就是男人!那麼,為甚麼不能用相同的回答來定義女人呢?女人就是女人!耶穌提示,在「那世界」其實不需要把男人和女人對立,男人不是女人的某某,女人也不是男人的某某,因為耶穌說:「既是復活的人,就為上帝的兒子(女)」.在「那世界」,沒有人被其他人所擁有,因為所有人都屬於上帝.

撒都該人以「這世界」,尤其是在他們當時的歷史時空中的婚姻狀況去投射「那世界」,以為「那世界」也是如此；耶穌提醒我們,反而應該以「那世界」為願景去指導「這世界」前進.聖經也有以婚姻來描述上帝與人的關係,例如保羅在以弗所書論都婚姻時,大家看見字裡行間的夫妻關係和神人關係的互動:

5:21又當存敬畏基督的心,彼此順服.

5:22你們作妻子的,當順服自己的丈夫,如同順服主.

5:23因為丈夫是妻子的頭,如同基督是教會的頭；他又是教會全體的救主.

5:24教會怎樣順服基督,妻子也要怎樣凡事順服丈夫.

5:25你們作丈夫的,要愛你們的妻子,正如基督愛教會,為教會捨己.

最後,保羅清楚指出:

5:31為這個緣故,人要離開父母,與妻子連合,二人成為一體.

5:32這是極大的奧祕,但我是指著基督和教會說的.

保羅強調「與妻子連合,二人成為一體」這種「這世界」可以體會的親密,其實是「那世界」基督與教會,或說神與人的親密的預嚐.到底是怎樣的呢?那麼多親密的需要和渴求,如何可能被滿足呢?保羅說,這是極大的奧祕.詩篇23篇的詩人說:「耶和華是我的牧者,我必不致缺乏.」不致缺乏,可以譯為「沒有其他需要了」.奧古斯丁曾說:「上帝造人時,在人的靈魂裡留下一個像上帝的形狀一樣的破洞,因此只有當人找到上帝時,才能彌補心靈的空缺.」在「那世界」,我們對親密關係的需求,都將會一無所缺.

馬丁路德金牧師展示出以「那世界」來塑造「這世界」的見證,他看見遠象:「那時,上帝的所有兒女,黑人和白人,猶太教徒和非猶太教徒,基督教徒和天主教徒,都將手攜手,合唱一首古老的黑人靈歌:「終於自由啦!終於自由啦!感謝全能的上帝,我們終於自由啦!」因為他對「那世界」的熱愛,成為了改變「這世界」的動力,推動著「這世界」向著「那世界」前進.



\newpage

\section{看哪我造新天新地}
\label{sec:623}
\textbf{李均熊牧師 - 看哪!我造新天新地}
\newline
\newline
\hyperref[sec:622]{< < < PREV SERMON < < <}
~
\hyperref[sec:toc]{[返目錄]}
~
\hyperref[sec:624]{> > > NEXT SERMON > > >}


以賽亞書 65:17-25
\newline
\begin{longtable}{cl}
\hline
\hline
章節 & 經文 (和合本修訂版)\\
\hline
65:17 & \begin{tabularx}{0.7\textwidth}{X} 「看哪，我造新天新地！從前的事不再被記念，也不被人放在心上； \end{tabularx} \\ \\ \relax
65:18 & \begin{tabularx}{0.7\textwidth}{X} 當因我所造的歡喜快樂，直到永遠；看哪，因為我造耶路撒冷為人所喜，造其中的居民為人所樂。 \end{tabularx} \\ \\ \relax
65:19 & \begin{tabularx}{0.7\textwidth}{X} 我必因耶路撒冷歡喜，因我的百姓快樂，那裡不再聽見哭泣和哀號的聲音。 \end{tabularx} \\ \\ \relax
65:20 & \begin{tabularx}{0.7\textwidth}{X} 那裡沒有數日夭折的嬰孩，也沒有壽數不滿的老人；因為百歲死的仍算孩童，未達百歲而亡的算是被詛咒的。 \end{tabularx} \\ \\ \relax
65:21 & \begin{tabularx}{0.7\textwidth}{X} 他們建造房屋，居住其中，栽葡萄園，吃園中的果子； \end{tabularx} \\ \\ \relax
65:22 & \begin{tabularx}{0.7\textwidth}{X} 並非造了給別人居住，也非栽種給別人享用；因為我百姓的日子必長久如樹木，我的選民必享受親手勞碌得來的。 \end{tabularx} \\ \\ \relax
65:23 & \begin{tabularx}{0.7\textwidth}{X} 他們必不徒然勞碌，所生產的，也不遭災害，因為他們和他們的子孫都是蒙耶和華賜福的後裔。 \end{tabularx} \\ \\ \relax
65:24 & \begin{tabularx}{0.7\textwidth}{X} 他們尚未求告，我就應允；正說話的時候，我就垂聽。 \end{tabularx} \\ \\ \relax
65:25 & \begin{tabularx}{0.7\textwidth}{X} 野狼必與羔羊同食，獅子必吃草，與牛一樣，蛇必以塵土為食物；在我聖山的遍處，牠們都不傷人，也不害物；這是耶和華說的。」 \end{tabularx} \\ \\
[1ex]
\hline
\hline
\end{longtable}
講題:看哪!我造新天新地

經文:以賽亞書65章17至25節

講員:李均熊牧師

日期:2022年11月13日

1.引言:新天新地的想像

過去當我在崇基禮拜堂宣道時,遇到的舊約經課都是環繞審判的信息；今天否極泰來,終於可以看到上主給與祂子民的應許和盼望.「看哪!我造新天新地!」不知大家想像的這個新天新地是一個怎樣的地方?可能我們受啟示錄21章的影響,會想像新天新地就是黃金街,碧玉城,不再有死亡,哭泣,痛苦.又有生命泉可以白白取用.但你仔細看以賽亞書中的經文,在這個新天新地裡,大家並不是升仙脫離人世；反之,這裡的人都會死亡,只是沒有壽數不足而去逝的人而矣(20節).大家不是再不吃人間煙火,大家仍要食仍要住,要耕種工作,只是人民的物質生活得到保障(21-23節),也可能因為如此,住在這新天新地的人才都充滿歡喜快樂,而不再聽見哭泣的聲音(18-19節),所以這也不算甚麼新奇事情.經文明言,這個新天新地是上主所造的耶路撒冷,是一個人間的城巿,並從天降下的新耶路撒冷.這個新天新地之所以是新,並不是上主要推倒祂第一個創造,然後重新再造一個天地；這個天地之所以是新,是因為「從前的事不再被記念,也不再被追想」(17節)了.這個天地之所以新,是住在耶路撒冷的居民在心態上看它是嶄新(just anew in mindset)而矣,而非上主再次從無造有(creatio ex nihilo),造一些新奇事物出來.我們可能覺得,唯一比較神奇的,是最後一節「豺狼必與羊羔同食；獅子必吃草與牛一樣；塵土必作蛇的食物.在我聖山的遍處,這一切都不傷人,不害物.」(25節)但若你回想創世記第一章,上主所造的天地本身就是所有動物吃青草,而人就吃有結種子的菜蔬和有核的果子(創1:29-30),蛇就真的是吃塵土(創3:14),是到了洪水之後,人才會吃肉(創9:3).所以上主只是回復創造本來的設計,而非造甚麼我們認為神奇的事情.

2.釋經:為何提起新天新地?

雖然以賽亞書這個新天新地不像啟示錄那個基督再來所打造的那樣金碧輝煌,但對於我們今日活在幻變無常的世界,無天生活充滿焦慮與不確定,這個新天新地也不賴呀!可惜,這個上主應許到今天仍未實然,感覺上甚至遙遙無期.是甚麼時候上主才會造這個新天新地呢?為何應許遲遲未可實現呢?要解答這個問題,其實首先要問:為甚麼以賽亞書要在這裡提起這個新天新地呢(why bring it up)?若我們細看這段經課的上文下理,你就會發現上主要造的這個新天新地是有一個很重要的背景的.這段經文並不是獨立存在在以賽亞書65章中的,若你細看之前8-16節的經文,你或會驚覺兩段經文之間竟有如此之多的呼應地方.

賽65章17-25節與8-16節的關係

讓我用下表去展示今日經課經文和前段8-16節經文的呼應之處:

賽65章17-25節

賽65章8-16節

植物的意像

22節因為我民的日子必像樹木的日子

8節葡萄中尋得新酒,

人就說:不要毀壞,

因為福在其中.

動物的意像

25節野狼必與羔羊同食,獅子必吃草,與牛一樣,塵土必作蛇的食物；在我聖山的遍處,這一切都不傷人,不害物……

10節沙崙平原必成為羊群的圈,亞割谷成為牛群躺臥之處,都為尋求我的民所得.

可喜可樂vs.悲痛而哀哭

18-19節你們當因我所造的永遠歡喜快樂；因我造耶路撒冷為人所喜,造其中的居民為人所樂.我必因耶路撒冷歡喜,因我的百姓快樂……

14節我的僕人因心中高興歡呼,你們卻因心中憂愁哀哭,又因心裡憂傷哀號.

足享壽數vs.屈身被殺

20節其中必沒有數日夭亡的嬰孩,也沒有壽數不滿的老者；因為百歲死的仍算孩童,未達百歲死的罪人算被詛咒的.

12節我要命定你們歸於刀下,都必屈身被殺……15節主耶和華必殺你們……

未求告就應允vs.我呼喚你們不回應

24節他們尚未求告,我就應允；正說話的時候,我就垂聽……

12節因為我呼喚,你們沒有答應；我說話,你們沒有聽從……

從前的事不再被記念vs.

你們必留下自己的名(用作賭咒)

17節看哪,我造新天新地!從前的事不再被記念,也不再追想……

15節你們必留下自己的名,為我選民指著賭咒……

其他呼應:

21節吃……23節子孫……後裔

13節飢餓……吃9節後裔……繼承

從此表可見,17至25節的應許,其實是在回應8至16節的處境.那其實是甚麼的境況?

經文的背景——聖民重建的假象(參賽56-59章)

大部份學者都認為由以賽亞書56章至66章以賽亞書結束為止,其關注的歷史時空已經由40-55章那處關於巴比倫帝國末期,轉而為波斯時期,聖殿已被重建的日子中.因為在56章開始,經文就論到上主的殿必稱為萬民禱告的殿(賽56:7).在40-55章中,子民在帝國交替的日子,又灰心又惶恐,不知如何面對前路的時候,經文的信息就應許塞魯士(古列)就是上主的受膏者(賽45:1),要帶來復興和釋放(49:8-9).現在塞魯士已經收服巴比倫,波斯當政,但所應許的復興和釋放卻仍沒有實現,人民的生活仍是苦不堪言.面對這樣的困境,經文焦點就放在回應應許為何遲遲未實現.其原因可以在56-59章可以看到,正如59:1-2所言:「耶和華的膀臂並非縮短,不能拯救,耳朵並非發沉,不能聽見,但你們的罪孽使你們與神隔絕；你們的罪惡使他掩面不聽你們.」換句話說,上主並不是不能實現祂的應許,而是因為子民中的罪孽使到神與他們隔絕.若我們從56-66章那同心圓般的經文編排結構來看,今日經課的經文正相對於56-59章解釋應許未能實現的原因.

從56-59章的經文來看,上主應許為何未能實現呢?若是因為子民的罪孽將神隔開,那又是怎樣的罪孽呢?綜合56-59章的經文來看,猶太人當時的罪狀主要有三:一,就是雖有聖殿可以獻祭,有機會盡上一切宗教禮儀,守安息日和禁食會,但人民心裡面卻只是想著如何欺壓窮人,「搵人著數」,正如58:3-4所言:「看哪,你們禁食的日子仍求利益,勒逼人為你們做苦工.你們禁食,卻互相爭競,以凶惡的拳頭打人.」(我稱這為「齋口唔齋心」)其次,他們雖然重建了聖殿,但竟仍有人恢復舊有的迦南宗教,崇拜令他們亡國的偶像,以為它們可以帶來復興!(參57:3-13)最後,亦都是最致命的罪狀,就是社會中充滿貪婪,欺壓,謊言和強暴,不公不義無日無之,但竟然無人理會,也無人為此代求,「無一人按公義告狀,無一人憑誠實辯白；都倚靠虛妄,說謊言.所懷的是毒害；所生的是罪孽.」(59:4)

在波斯時期,猶太人能夠回歸故土,重建聖殿,但聖殿卻只空有其表,因為上主期望的聖民卻沒有因為聖殿的重建,而沒有跟著重建起來.沒有聖民的重建,所有回歸和復興都只是假像

.正因如此,所應許的復興和更新並沒有如期出現,人民仍面對惡劣的社會境況.

應許仍在嗎?(參賽60-62章)

但若只是如此,我們就真的只有灰心失望.但由56-66章的經文編排的同心圓結構來看,上主的應許仍然是有效的,因為由60-62章的核心部份,上主在之前所言的釋放,復興仍然存在,不過需要等待上主的僕人去宣告和實現.大家熟識的,耶穌在會堂曾宣讀的經文,就在這一個核心之內:「主耶和華的靈在我身上；因為耶和華用膏膏我,叫我傳好信息給謙卑的人,差遣我醫好傷心的人,報告被擄的得釋放,被囚的出監牢；報告耶和華的恩年,和我們神報仇的日子；安慰一切悲哀的人.」(賽61:1-2)

3.釋義:應許延遲,也是審判延遲

在這個大背景下看賽65:8-25,我希望你看到,其實應許延遲,但同時也表示上主審判延遲.

子民:尋求我的

應許的確是延遲了,但不代表上主忘記了祂的子民.因為在子民當中,仍有願意尋求上主的,上主就要使他們領受這個上主的應許.在第9-10節中,上主就宣告:「我必從雅各中領出後裔,從猶大中領出承受我眾山的.我的選民必承受；我的僕人要在那裡居住.沙崙平原必成為羊群的圈；亞割谷必成為牛群躺臥之處,都為尋求我的民所得.」又在22節中這樣說:「他們建造的,別人不得住；他們栽種的,別人不得吃；因為我民的日子必像樹木的日子；我選民親手勞碌得來的必長久享用.」

我們是否仍然保有盼望?我們是否願意在艱難,甚至虛謊暴力的日子去尋求上主?我們願意作尋求上主的選民,在應許和盼望好像失落時仍相信上主仍會工作?

你們:離棄上主的

但我們不要忘記,上主應許實現被阻礙的原因,可能就是因為我們沒有回應上主,正如11節所言:「但你們這些離棄耶和華,忘記我的聖山,給時運擺筵席,給天命盛滿調和酒的,我要命定你們歸在刀下,都必屈身被殺；因為我呼喚,你們沒有答應；我說話,你們沒有聽從；反倒行我眼中看為惡的,揀選我所不喜悅的.」

「你們」是誰?我們會否手指向外,以為就是別人?面對社會不公不義的時候,我們有沒有回應上主的呼喚?上主要求我們照顧有需要的人時,我們有沒有聽從?我們所選擇的,是上主所喜悅的,還是祂看為可惡的呢?攪不好,「我們」就是那些離棄上主的「你們」,走去尋求時運和天命,以為這些非神的力量會為我們帶來生機,誰不知這些只會使我們倒下來呢?究竟上主應許的延遲,會否與我有關?

僕人的角色:不全然毀滅的原因

在整個8-25節的大段落的一開始,上主宣告:「葡萄中尋得新酒,人就說:不要毀壞,因為福在其中.我因我僕人的緣故也必照樣而行,不將他們全然毀滅.」原來應許延遲,並不是因為離棄上主的人「拖後腿」.上主要履行祂的應許,你認為真會被這些不忠不信的人阻礙嗎?反倒是因為上主的僕人仍在守望,仍在工作,上主才寬容忍耐,延遲的其實是祂的審判,而非祂的應許.

正如彼後3:9所說的:「主所應許的尚未成就,有人以為他是耽延,其實不是耽延,乃是寬容你們,不願有一人沉淪,乃願人人都悔改.」

4.今日意義:我們能為社群帶來新天新地嗎?

究竟,我們是離棄上主,令上主應許耽延的人?還是我們正尋求上主,可以帶來新天新地的選民,甚至是上主的僕人呢?



\newpage

\section{一子錯}
\label{sec:624}
\textbf{鄧瑞強博士 - 一子錯}
\newline
\newline
\hyperref[sec:623]{< < < PREV SERMON < < <}
~
\hyperref[sec:toc]{[返目錄]}
~
\hyperref[sec:625]{> > > NEXT SERMON > > >}


講題:一子錯

經文:撒母耳記下11:1-4,13:7,10-14,21-22,15:13-14,16:20-22

講員:鄧瑞強博士

日期:2022年11月20日

各位弟兄姊妹,早安.在這疫症裡,祝願各人平安.

看電視劇,已看夠宮廷裡的殘酷鬥爭.在權力底下,誰人生,誰人死；幾時生,幾時死；為什麼生,為什麼死,很多時,連當事人都說不出一個所以然來.這些宮廷裡的黑暗,也發生在以色列人最顯赫的帝王大委身上.聖經的歷史書寫,不是帝王自己的歷史書寫,而是上帝視角下的歷史書寫,將帝王的黑暗,暴露在光天化日之下.聖經的歷史書寫,是要讓人明白,就算是最隱密的事情,神的眼睛仍然明察秋毫,祂要讓事情得到正義的了斷.一個人若不正視自己的黑暗,將引發更深的黑暗,聖經的歷史書寫,就是叫那些在黑暗中的人,及早回頭,走向光明.

以色列民族的歷史,去到某個階段,需要建立王權制度.但,人真能善用權力嗎?聖經的歷史書寫,明明白白,表明人多麼傾向誤用權力.聖經肯定王權制度,同時,聖經批判王權制度.要維護秩序,運用權力是有需要的；但當權力變成終極力量時,就可能連秩序都失去了.今日,我們看宮廷裡連場的權力遊戲,神將人的劣根性暴露在我們面前,讓我們自己決定,是延續他們的遊戲,抑或放下屠刀,走向光明.

今日講《撒母耳記下》,講道的經文分佈得頗零散,我不得不如此,因為我想將事情的起頭,中間的轉折,事情導致的結局,作一整體而扼要的說明.

我們逐段經文來看.

第一段經文:撒母耳記下11:1-4

11:1過了一年,到列王出戰的時候,大衛又差派約押,率領臣僕和以色列眾人出戰.他們就打敗亞捫人,圍攻拉巴.大衛仍住在耶路撒冷.

11:2一日,太陽平西,大衛從床上起來,在王宮的平頂上遊行,看見一個婦人沐浴,容貌甚美,

11:3大衛就差人打聽那婦人是誰.有人說:「她是以連的女兒,赫人烏利亞的妻拔示巴.」

11:4大衛差人去,將婦人接來；那時她的月經才得潔淨.她來了,大衛與她同房,她就回家去了.

大衛王,一個王朝的奠基人,一個所謂最好的君王,當他的元帥約押在前線衝鋒陷陣時,他在王宮裡睡到日上三竿.他看上一個女人,這女人是他手下猛將烏利亞的妻子,他忍不住內心的慾望,他將她召來,他與她同房.聖經只是白描,對這些事沒有加以道德判斷,但看在我們眼裡,已覺得這皇帝很不妥當.

故事發展下去,是這女人懷了孕.為了掩飾自己的罪行,大衛王從前線召回這女人的丈夫,希望他與妻子同房,如此,大衛王作過的醜事便能瞞天過海.誰知,人算不如天算,這男人竟因有戰事在身,不肯與妻子同房.大衛王無計可施,只能用計殺了他.於是,他命令前線的元帥,將這人派往前線最危險的必死之地.

最後,大衛王奸計得逞,借刀殺了這忠勇之士,奪了他的妻子.大衛王以為這事,除了自己和元帥約押知道以外,無人知道.但是,神看見,神知道.大衛王以為自己能操控這事的發展,事實上,邪惡的發展超出他的預期.

描寫完大衛王的惡事後,《撒母耳記下》平行地描繪他的長子暗嫩的惡事.

暗嫩戀慕同父異母的妹妹她瑪.為了得到她,暗嫩裝病,求父親叫妹妹拿東西來給他吃.

13:7大衛就打發人到宮裡,對她瑪說:「你往你哥哥暗嫩的屋裡去,為他預備食物.」

奸險的大衛不知他兒子的奸險,他不知道,邪惡已生出邪惡.他作的惡,已生出惡果.

13:10暗嫩對她瑪說:「你把食物拿進臥房,我好從你手裡接過來吃.」她瑪就把所做的餅拿進臥房,到她哥哥暗嫩那裡,

13:11拿著餅上前給他吃,他便拉住她瑪,說:「我妹妹,你來與我同寢.」

13:12她瑪說:「我哥哥,不要玷辱我.以色列人中不當這樣行,你不要做這醜事；

13:13你玷辱了我,我何以掩蓋我的羞恥呢?你在以色列中也成了愚妄人.你可以求王,他必不禁止我歸你.」

13:14但暗嫩不肯聽她的話,因比她力大,就玷辱她,與她同寢.

大衛覬覦別人的妻子,大衛的兒子卻覬覦大衛的女兒.大衛在外的醜事,報應在他的家裡.在這裡,傷害人的,是大衛家中的長子.受害的,是他家中的女兒.大衛那為所欲為的性格,他的長子完全承繼.悲劇在家中慢慢發酵,大衛終有一天要親自體會自己種下的種子所帶來的苦果.

家中發生這等醜事,身為家長的大衛,如何處理?

13:21大衛王聽見這事,就甚發怒.

13:22押沙龍並不和他哥哥暗嫩說好說歹；因為暗嫩玷辱他妹妹她瑪,所以押沙龍恨惡他.

大衛王只是發怒,他事實上什麼也沒有做.這個父親或許慣於縱容他的兒女,或許他個性優柔寡斷,不過,看來這些都不是原因,我想,是因為他自己其身不正,他沒有一個穩固的立足點,可指責他人.一個為了奪人妻子而殺人的人,斷沒有任何理由,去指責一個罪行比自己輕的人.罪的結果是縱容罪.

罪的更新得牢固了,它會生發另一條枝幹,另一個主角出場了.她瑪的同母的親哥哥押沙龍,他將一切記在心中.表面上,押沙龍和他的大哥暗嫩仍稱兄道弟,但心中已盤算著總會行的一步.每一日的相遇,每一次的應對,每一步的讓對方疏於防避,就是為了將來的一天.這是什麼樣的家庭生活?表面的親近隱藏著內裡的殺機,正如大衛的情慾伴隨著謀殺.

兩年後,當人人都忘掉暗嫩強姦她瑪的事件時,當人人都以為一切已復歸平靜時,押沙龍就趁機殺了暗嫩.一方面,替妹妹的受辱討一公道；另方面,也替自己在爭奪王位一事上剷除最大的障礙.

大衛這兩個兒子,反映著大衛的人格特質.大衛的長子暗嫩,繼承了父親的為求目的不擇手段.暗嫩走著大衛同一條軌跡,但大衛殺了人,暗嫩卻被殺.大衛另一個兒子押沙龍,繼承了父親的陰險.他殺了哥哥,主導著一場大衛王朝的「玄武門之變」.這條路走下去,必然就是篡位.

15:13有人報告大衛說:「以色列人的心都歸向押沙龍了!」

15:14大衛就對耶路撒冷跟隨他的臣僕說:「我們要起來逃走,不然都不能躲避押沙龍了；要速速地去,恐怕他忽然來到,加害於我們,用刀殺盡合城的人.」

大衛,一代英雄人物,能指點江山,誰想到,會有這種下場.被兒子押沙龍奪權,追殺,他要漏夜逃命,何等淒涼.戰勝巨人,力拔山兮的大衛,如今像過街老鼠,落泊潦倒.世事真難料.

16:20押沙龍對亞希多弗說:「你們出個主意,我們怎樣行才好?」

16:21亞希多弗對押沙龍說:「你父所留下看守宮殿的妃嬪,你可以與她們親近.以色列眾人聽見你父親憎惡你,凡歸順你人的手就更堅強.」

16:22於是人為押沙龍在宮殿的平頂上支搭帳棚；押沙龍在以色列眾人眼前,與他父的妃嬪親近.

押沙龍看來已得到大衛的江山了,他問軍師亞希多弗,現在應做什麼?軍師亞希多弗對他說,要公開地親近你父親的妃嬪.這步棋是「將」對方「軍」的妙著.親近上一任君王的妃嬪,作為一公開表演,是表明他的東西現在是我的了,這是公開的宣示主權,像新的獅王趕走舊的獅王後,便能佔有整個雌獅群一樣.這行動當然是對上一任君王大衛的公開羞辱,這必然招致大衛憎惡押沙龍,而兩父子的關係從此就誓成水火.敵對關係一形成,軍民就必須靠邊站,站這邊的,便以消滅那邊的為己任.在榮辱文化裡,被辱的人應努力挽回尊嚴,若大衛無法這樣做,則證明他的勢已去,靠押沙龍這邊站的人,就更能一往情深地押注在押沙龍身上.如此,押沙龍的王位,便坐得更穩了.

故事當然還有下文,但我就在此打住了.我們發現,這些歷史的畫面是如此荒謬,超乎小說家的想像力.真實的歷史情節,總是比我們能想像的,更不可思議.

十日前,有人在海底發現美國太空穿梭機「挑戰者」號的碎片.1986年,這隻超高科技的巨無霸穿梭機升空,透過電視轉播,全世界的人都看著它.起飛後73秒,它就爆炸解體.這隻穿梭機,每一個環節,部件,都是高科技弟兄姊妹產品.然而,它的失敗,卻只是源自一個低級錯誤,一個合成橡膠製的「Ｏ形環」因發射前周遭環境的低溫而失去了彈性,令助推器燃燒時的高溫能外泄,便造成如此巨大的災難.巨無霸的倒下,不需要偉大的原因；微少的差錯,就足以置命.

大衛宮廷裡連串的悲劇,源自大衛的一個低級錯誤.他百無聊賴,卻貪戀一個女子.為了得到她,他殺了她的丈夫.唉,如此一個人物,怎會走出一步這麼差的棋?為了貪吃一隻卒,他敗了一盤棋.唉,真是一子錯,滿盤皆落索.人在做,神在看.神沒有出手,因為人的邪惡,自會結出邪惡的苦果.人怎能以為,手所作的惡事,是沒有後果的.

大衛的長子暗嫩走大衛的老路,他強姦他的妹妹.大衛作為父親,沒有什麼反應.他讓邪惡生根.

大衛另一個兒子押沙龍為親妹妹報仇雪恥,大衛作為父親,也沒有什麼大反應.如此,權力的誤用與暴力的仇殺,已在家族中根深蒂固.

然後,是兒子暴力地推翻父親,兒子淫辱父親的妃嬪,將暴力與淫亂推向極致.大衛終於吃到他種下的果子.

原先,神看著人間充滿邪惡,就選召以色列民族,去開展天國工程.以色列民族進入迦南地,神希望他們建立一敬拜神的國度,但很快,他們就與迦南人同化,拜迦南的神.《士師記》記載,他們的強暴與淫亂,比起迦南人,有過之而無不及.為了確立秩序,他們設立君王制度.《撒母耳記上》記載,他們第一個王掃羅,因不遵守神的吩咐,很快便被大衛取代了.今日講的《撒母耳記下》,我們看到這個王朝是如何展開的.這一幕,光天化日之下,淫辱他人的妻妾,又使整幅圖畫回到《士師記》的結尾.歷史發展了這麼久了,最令人矚目的畫面,依然是權力的誤用,依然是暴力,依然是將對方物化,將之變成是自己發洩的對象.王權制度有其用處,但怎能保證君王不誤用權力而淫辱他人?

聖經的歷史書寫,並沒有給出什麼答案,它只是讓這人性的事實,呈現在人眼前.

當我們看到這些事實時,我們自然明白,君王與人民都敬畏神,是多麼寶貴.

我們看到,天國降臨,是多麼重要的一件事.歷史將會發展下去,神呼召祂的子民,為天國的實然而繼續努力.



\newpage

\section{黑夜到底多長}
\label{sec:625}
\textbf{葉菁華教授 - 黑夜到底多長?}
\newline
\newline
\hyperref[sec:624]{< < < PREV SERMON < < <}
~
\hyperref[sec:toc]{[返目錄]}
~
\hyperref[sec:626]{> > > NEXT SERMON > > >}


馬太福音 24:36-44
\newline
\begin{longtable}{cl}
\hline
\hline
章節 & 經文 (和合本修訂版)\\
\hline
24:36 & \begin{tabularx}{0.7\textwidth}{X} 「但那日子，那時辰，沒有人知道，連天上的天使也不知道，子也不知道，惟有父知道。 \end{tabularx} \\ \\ \relax
24:37 & \begin{tabularx}{0.7\textwidth}{X} 挪亞的日子怎樣，人子來臨也要怎樣。 \end{tabularx} \\ \\ \relax
24:38 & \begin{tabularx}{0.7\textwidth}{X} 在洪水以前的那些日子，人照常吃喝嫁娶，直到挪亞進方舟的那日， \end{tabularx} \\ \\ \relax
24:39 & \begin{tabularx}{0.7\textwidth}{X} 不知不覺洪水來了，把他們全都沖去。人子來臨也要這樣。 \end{tabularx} \\ \\ \relax
24:40 & \begin{tabularx}{0.7\textwidth}{X} 那時，兩個人在田裡，一個被接去，一個被撇下。 \end{tabularx} \\ \\ \relax
24:41 & \begin{tabularx}{0.7\textwidth}{X} 兩個女人推磨，一個被接去，一個被撇下。 \end{tabularx} \\ \\ \relax
24:42 & \begin{tabularx}{0.7\textwidth}{X} 所以，你們要警醒，因為不知道你們的主哪一天來到。 \end{tabularx} \\ \\ \relax
24:43 & \begin{tabularx}{0.7\textwidth}{X} 你們要知道，一家的主人若知道晚上甚麼時候有賊來，就必警醒，不讓賊挖穿房屋。 \end{tabularx} \\ \\ \relax
24:44 & \begin{tabularx}{0.7\textwidth}{X} 所以，你們也要預備，因為在你們想不到的時候，人子就來了。」 \end{tabularx} \\ \\
[1ex]
\hline
\hline
\end{longtable}
講題:黑夜到底多長?

經文:馬太福音24:36-44

講員:葉菁華教授

日期:2022年11月27日

馬太福音24:36-44

36「但那日子,那時辰,沒有人知道,連天上的使者也不知道,子也不知道,惟獨父知道.37挪亞的日子怎樣,人子降臨也要怎樣.38當洪水以前的日子,人照常吃喝嫁娶,直到挪亞進方舟的那日；39不知不覺洪水來了,把他們全都沖去.人子降臨也要這樣.40那時,兩個人在田裡,取去一個,撇下一個.41兩個女人推磨,取去一個,撇下一個.42所以,你們要警醒,因為不知道你們的主是哪一天來到.43家主若知道幾更天有賊來,就必警醒,不容人挖透房屋；這是你們所知道的.44所以,你們也要預備,因為你們想不到的時候,人子就來了.」

引言

人們對未來的不同態度

樂觀,「好友」,明天會更好

悲觀,「淡友」,明天會更差

香港基督教協進會(Hong Kong Christian Council)與香港中文大學崇基學院神學院(Divinity School of Chung Chi College, CUHK)合作進行一項主題是市民對基督教的觀感(public perception of Protestantism)的研究,委託中大香港亞太研究所(Hong Kong Institute of Asia-Pacific Studies)進行的一項全港調查(survey):

香港市民對前景

「幾樂觀」加「非常樂觀」:23\%

一般:27\%

「幾唔樂觀」加「非常唔樂觀」:50\%

特別是:

學歷和收入較高的人士傾向於不樂觀

基督徒也傾向不樂觀

不知道世界其他地方居民對前景如何,但相信也難以樂觀

全球氣候暖化,加上就快三年的新冠疫情,再加上俄羅斯發動對烏克蘭的戰爭

導致糧食,能源價格飆升,因而通貨膨脹加劇；但不少地方經濟倒退,引發滯脹(stagflation)的危機

晶片供應短缺和大國博奕,更加深經濟危機

社會政治方面,有些地方政局混亂,人民惶恐不安；有些地方威權甚至極權統治,人民常受監控和壓制；有些地方的疫情和防疫政策,令人民飽受困苦.

或者,無論遠近,每天都有荒謬的事情,禮崩樂壞,扭曲人性,違反公義,持強凌弱,顛倒是非,指鹿為馬

馬太福音的處境

(1)公元66年:猶太人與羅馬人戰爭

(2)公元70年:

耶路撒冷圍城,數十萬人死亡

全城破壞,聖殿被毀,猶太人民族浩劫

(3)馬太群體

跟從耶穌基督,跟從那被羅馬帝國用十字架酷刑對待的危險人物

基督徒中的猶太人與外邦人關係緊張,屢有衝突

(4)不安的狀況

(5)基督徒問:主何時再來

(6)太24:3 門徒問耶穌,「你來臨和世代的終結有甚麼預兆呢?」

(7)24章整章回答這個問題

7民要攻打民,國要攻打國,多處必有饑荒,地震.8這都是災難的起頭.9那時,人要使你們陷在患難裡,也要殺害你們；你們又要為我的名被萬民憎恨.10那時,會有許多人跌倒,也會彼此陷害,彼此憎恨

……14這天國的福音要傳遍天下,對萬民作見證,然後終結才來到

到時終末時候會有大災難,有假先知

當你們看見這一切,就知道那時候近了

馬太24:36-44

(1)但那日子,那時辰,沒有人知道,連天上的使者也不知道,子也不知道,惟獨父知道.

(2)「所以,你們要警醒,因為不知道你們的主是哪一天來到……你們也要預備,因為你們想不到的時候,人子就來了.」

(3)主再來(Christ’s second coming)的不確定

無人知,不用預測世界末日

只有父知道:耶穌基督的人性

(4)只知不確定,容易以為未有來那麼快來,以至你仍可生活如常,不受影響

(5)馬太強調我們不能掉以輕心,不能滿不在乎,因為不確定,所以主隨時會來,而且時間可以好近,所以我們要警醒等候

(6)人子降臨時就如挪亞時代的洪水,人們生活如常,但一向做的事情會突然中斷,災難會突然臨到

(7)兩個人在田裡工作,兩個女人在推磨,取去一個,撇下一個

「被提」(rapture):

有些人很關心被提,神學上有爭論,是災前,災中,災後被提

太關注被提,基督徒容易流於自我中心,甚至流於個人化,只是關心如果今日主再來,我是否會被提,忘記經文關心的是世界和歷史

被提與撇下,是主再來的審判

審判世人的主正是為世人承受審判的主

經文沒有說被撇下的那些人犯了甚麼罪,所以我們不能說「唔關我事」

若看下一章(太25)的終極審判,似乎主耶穌所關心的,不是信主與不信主,而是有沒有服事主,透過對於弱勢人士的服事

(8)賊的比喻

盜賊不會預先話你知幾時來

總要時常準備

結語

(1)將臨期(Advent)

基督再來,帶來天國的最終臨在,爾國降臨,爾旨得成

(2)講道題目:「黑夜到底多長?」(How Long?)

來自《教會根基歌》(The Church’s One Foundation)其中一節

信徒警醒心焦問,「黑夜到底多長?」yet saints their watch are keeping, their cry goes up, “How long?”

荒謬的事情,禮崩樂壞,扭曲人性,違反公義,顛倒是非等,要到何時?

詩歌說:「但是悲泣變歌聲,轉瞬便見晨光」and soon the night of weeping shall be the morn of song

但轉瞬(soon)似乎有排,我們未見晨光

(3)主基督必要再來,必會審判,但不知何時,隨時會來

已知的未知 “known unknown”

(4)在充滿危機和荒謬的世界中,馬太要我們持守信仰,警醒等候我們基督；要警醒,準備

警醒:心思意念超越現世,放在來臨中的上帝；我們的盼望

(5)12黑夜已深,白晝將近.我們就當脫去暗昧的行為,帶上光明的兵器.13行事為人要端正,好像行在白晝.作鹽作光,對抗黑暗及敗壞的權勢

跟從主,服事主,包括弱勢人士

那些是弱勢人士?我們如何服事他們?

(6)黑夜已深,白晝將近.

冬至(Winter Solstice)一陽生



\newpage

\section{那盡了責的見證人}
\label{sec:626}
\textbf{關瑞文教授 - 那盡了責的見證人}
\newline
\newline
\hyperref[sec:625]{< < < PREV SERMON < < <}
~
\hyperref[sec:toc]{[返目錄]}
~
\hyperref[sec:627]{> > > NEXT SERMON > > >}


馬太福音 3:1-12
\newline
\begin{longtable}{cl}
\hline
\hline
章節 & 經文 (和合本修訂版)\\
\hline
3:1 & \begin{tabularx}{0.7\textwidth}{X} 在那些日子，施洗的約翰出來，在猶太的曠野宣講： \end{tabularx} \\ \\ \relax
3:2 & \begin{tabularx}{0.7\textwidth}{X} 「你們要悔改！因為天國近了。」 \end{tabularx} \\ \\ \relax
3:3 & \begin{tabularx}{0.7\textwidth}{X} 這人就是以賽亞先知所說的：「在曠野有聲音呼喊著：預備主的道，修直他的路。」 \end{tabularx} \\ \\ \relax
3:4 & \begin{tabularx}{0.7\textwidth}{X} 這約翰身穿駱駝毛的衣服，腰束皮帶，吃的是蝗蟲和野蜜。 \end{tabularx} \\ \\ \relax
3:5 & \begin{tabularx}{0.7\textwidth}{X} 那時，耶路撒冷、全猶太和全約旦河地區的人，都到約翰那裡去， \end{tabularx} \\ \\ \relax
3:6 & \begin{tabularx}{0.7\textwidth}{X} 承認他們的罪，在約旦河裡受他的洗。 \end{tabularx} \\ \\ \relax
3:7 & \begin{tabularx}{0.7\textwidth}{X} 約翰看見許多法利賽人和撒都該人也來受洗，就對他們說：「毒蛇的孽種啊，誰指示你們逃避那將要來的憤怒呢？ \end{tabularx} \\ \\ \relax
3:8 & \begin{tabularx}{0.7\textwidth}{X} 你們要結出果子來，和悔改的心相稱。 \end{tabularx} \\ \\ \relax
3:9 & \begin{tabularx}{0.7\textwidth}{X} 不要自己心裡說：『我們有亞伯拉罕為祖宗。』我告訴你們，神能從這些石頭中給亞伯拉罕興起子孫來。 \end{tabularx} \\ \\ \relax
3:10 & \begin{tabularx}{0.7\textwidth}{X} 現在斧子已經放在樹根上，凡不結好果子的樹就砍下來，丟在火裡。 \end{tabularx} \\ \\ \relax
3:11 & \begin{tabularx}{0.7\textwidth}{X} 我是用水給你們施洗，叫你們悔改；但那在我以後來的，能力比我更大，我就是給他提鞋子也不配，他要用聖靈與火給你們施洗。 \end{tabularx} \\ \\ \relax
3:12 & \begin{tabularx}{0.7\textwidth}{X} 他手裡拿著簸箕，要揚淨他的穀物，把麥子收在倉裡，把糠用不滅的火燒盡。」 \end{tabularx} \\ \\
[1ex]
\hline
\hline
\end{longtable}
講題:那盡了責的見證人

經文:馬太福音3章1至12節

講員:關瑞文教授

日期:2022年12月4日

馬太福音3章1至12節

1在那些日子,施洗的約翰出來,在猶太的曠野宣講:2“你們要悔改!因為天國近了.”3這人就是以賽亞先知所說的:“在曠野有聲音呼喊著:預備主的道,修直他的路.”4這約翰身穿駱駝毛的衣服,腰束皮帶,吃的是蝗蟲和野蜜.5那時,耶路撒冷,全猶太和全約旦河地區的人,都到約翰那裡去,6承認他們的罪,在約旦河裡受他的洗.7約翰看見許多法利賽人和撒都該人也來受洗,就對他們說:“毒蛇的孽種啊,誰指示你們逃避那將要來的憤怒呢?8你們要結出果子來,和悔改的心相稱.9不要自己心裡說:‘我們有亞伯拉罕為祖宗.’我告訴你們,神能從這些石頭中給亞伯拉罕興起子孫來.10現在斧子已經放在樹根上,凡不結好果子的樹就砍下來,丟在火裡.11我是用水給你們施洗,叫你們悔改；但那在我以後來的,能力比我更大,我就是給他提鞋子也不配,他要用聖靈與火給你們施洗.12他手裡拿著簸箕,要揚淨他的穀物,把麥子收在倉裡,把糠用不滅的火燒盡.”

在這個主日,我們讀到馬太福音3章1至12節,要聆聽上主透過施洗約翰的事蹟,有什麼話要對我們說.

我們要看看,施洗約翰這位具有超凡勇氣並且很大膽的見證人,如何在人生的高低起跌裡,始終堅守作為見證人的本份,始終不忘信仰的初心.

我相信,坐在這裡的弟兄姊妹,有不少是大學老師,研究生,大學生,大學畢業生,即使不是人生大贏家,都應該是社會的中上階層,或者是準中上層.對於我們這群人,或許施洗約翰的事蹟,會對我們的人生構成一點挑戰.

讓我們回到經文去.

馬太福音安排了一個極有排場的鏡頭角度來讓施洗約翰出場的.

首先,在第3節,馬太福音作者提醒讀者們,施洗約翰殊不是普通人,他是萬眾期待的,是幾百年前,以賽亞先知所預告要出現的大人物.

第4節就描述到施洗約翰的服飾打扮和行徑,說他「身穿駱駝毛的衣服,腰束皮帶,吃的是蝗蟲和野蜜.」

大家千萬不要誤會,以為施洗約翰因為貧窮或者樸素,才這樣穿,這樣吃.

事實上,不說可能你不知,按列王記下1章8節,這種服飾打扮,正是他們的已故民族英雄,人民革命家,舊約偉大先知以利亞的標志性服飾打扮.

所以,隨後在馬太福音11章和17章,我們就看見耶穌直接說,施洗約翰就是以利亞.

還有,施洗約翰絕對是萬人迷,粉絲絕對比姜濤的多何止千百倍.

第5節,「那時,耶路撒冷,全猶太和全約旦河地區的人,都到約翰那裡去

…

在約旦河裡受他的洗.」

第7節甚至提到,連福音書大反派法利賽人和撒都該人,都要加入這個「粉絲團」!

其實,馬太也的確把施洗約翰的出場,搞得非常高調.

第1節那句「在那些日子,施洗的約翰出來」當中的「出來」這個詞組,在原文是與第13節介紹耶穌出場時的同一個詞組.

在新約時代,這個詞組通常是用來表示一個神聖的或尊貴的人物大駕光臨.

所以,如果這要拍成電影,這大概是一個君臨天下的震撼場面,背景音樂應該可以選用貝多芬那首降E大調第五鋼琴協奏曲,作品73,《皇帝協奏曲》.

的確,施洗約翰真的是一個非凡成功的人物.

就是連一世紀猶太著名歷史學家約瑟夫,也在他的歷史著作中提到施洗約翰,並稱呼他是個正直的人,是一名有無數跟隨者的出色演說家.

雖說英雄莫問出處,可是,你知道施洗約翰前半生的經歷嗎?你想知道他的出身嗎?

按一些歷史研究推斷,施洗約翰出身非常平凡,是個「唔多識野ge 鄉下仔」.他是個鄉間祭司的兒子.

他出身貧窮,受教育不多,絕對不是社會精英.

當日的祭司行業,很有階級制度.作為鄉下祭司階層,他的家庭,被有錢人或在大城市裡的祭司看不起,是被剝削的祭司階層.

按當時的歷史記載,這種低端祭司,無收入穩定的工作,甚至在同階層的祭司,有餓死的情況.

有研究聖經歷史的學者認為,施洗約翰有可能曾經因在大城市失意受欺負而加入了在鄺野(沙漠)的隱修團體,後來才再重出江湖,向百姓宣揚天國的審判即將來臨.

換言之,施洗約翰是出身寒微,自少貧窮,被人輕看,曾經在大城市打拼卻失敗,結果潦倒落魄而退隱江湖.

施洗約翰的前半生,有點類似像我這一代戰後出生的第一代香港人.

我們小時侯家貧,父母是從內地走難而來的難民,無論本來在內地有多富貴,走難來到香港,都是在工廠「打份牛工」,「餐搵餐食餐餐清」.

我媽媽本來就是中山石岐大地主家族的嬌嬌女,我爸爸就是順德大鋼鐵和印刷企業大老闆的長子.他們二人都是名符其實的富二代.不過同樣都是一夜成為難民,偷渡來到香港,展開「捱世界」的人生.

雖然有些當代的難民結果走上了類似李嘉誠的發達之路,大部分卻長時間留在草根階層,做其半斤八兩被人剝削的「打工仔」.

幼年時我們都是住在徙置區,70年代才遷入公屋,以「屋村仔」的身份長大.

不知在座的弟兄姊妹,有多少是這種出身的人呢?

不過,正如剛才所說,後期的施洗約翰,已經成為一個成功人物,人民心目中的英雄,眾人的偶像—萬人空巷排著隊要前來接受他的洗禮,加入他的群體.

我們香港人,也是像施洗約翰一樣,先苦後甜,先貧後富.這是香港的歷史.

在六七暴動後,當時的香港殖民政府痛定思痛,大灑金錢發展本地房屋,醫療,教育,社會福利,公共建設,土地,廉潔運動等等.這不但讓留下來的人開始安居樂業,生兒育女,疏遠政治鬥爭和意識形態,還成為了後來六,七十年代因全球化為香港帶來經濟奇跡的良好基礎.七十年代十年之間,香港一夜發財,其亞洲四小龍地位奠定下來了,即使期間發生了股災,也沒有給香港經濟帶來致命打擊.1965年至1980年,香港的國民生產總值年均增長率為6\%.到1987年,香港的人均國民生產總值為8070美元,高於紐西蘭,西班牙和愛爾蘭,只比英國低約五分之一.這時的香港,幾乎沒有人失業.消費主義抬頭,物質生活躍升,機會處處,遍地黃金.

這一代人,只要循規蹈矩,稍加努力,循著社會階梯,把握機會向上流,生活必然改善.

還有,若然你的智力比常人稍勝,努力稍多,又懂得把握機會,真的不會太難就出人頭地,成為中產階級,專業人士,甚至成為大學教授,或者從商而發達.你能夠達到的成就,擁有的財富,遠遠超過你們上一代,也為你們的下一代創造福氣.

還有,如果你能夠建立良好的人際網絡,累積更多社會資本,你的前途就更有把握了.所以如果你加入了類似教會般的互助互愛群體,你就更容易扶著社會階梯的把手向上直爬.這就是為什麼現在在教會裡的中產階層人數比總人口中的中產階層人數要多.

一般人,如果人生是先苦後甜,先貧後富,大概都會竭力不讓自己走回頭路.尤其到人生收成期,都會盡量保住自己已經擁有的,為下一代著想,會事事小心,即使不趨炎附勢,也不敢隨便見義勇為.這應該是人之常情,你與我都可能都是這樣的人.是嗎?

施洗約翰卻是例外.他受上主所托,見證天國福音,斥責世人罪惡,勸人悔改.

就算是遇到有權有勢的人,也都「冇面俾」照樣斥責,甚至明知要付上代價,也都無畏無懼,堅持說真話,見義勇為.

第7節:約翰看見許多法利賽人和撒都該人也來受洗,就對他們說:「毒蛇的孽種啊,誰指示你們逃避那將要來的憤怒呢?8你們要結出果子來,和悔改的心相稱.9不要自己心裡說:『我們有亞伯拉罕為祖宗.』我告訴你們,上帝能從這些石頭中給亞伯拉罕興起子孫來.10現在斧子已經放在樹根上,凡不結好果子的樹就砍下來,丟在火裡.

嘩!法利賽人和撒都該人是「毒蛇的孽種」,「要求他們結出果子來,和悔改的心相稱」,「警告他們,若不結好果子就砍他們下來,丟在火裡.」

真的罵過痛快啊!不過,你知道法利賽人和撒都該人是何許人嗎?

撒都該人是當時猶太教的四大派別之一,另外三大派別為法利賽人,愛色尼派和奮銳黨.

撒都該人是圍繞著祭司長為中心的教派,在宗教主導的猶太社會裡,祭司是社會文化等的核心階層,撒都該人自然就是社會的領袖.

作為猶太教領袖,撒都該人在當日是有權有勢有地位並且有錢的.

他們不只是宗教領袖,更是因為宗教領袖的身份,實行以夷制夷政策的羅馬帝國,把巴勒斯坦猶太地區的政治行政權,交予他們.所以,他們是非常有權勢的,是代表著統治者,真的「一隻手指尾就可以摙死你」施洗約翰.施洗約翰罵他們是毒蛇的孽種,絕對是「膽生毛」,「唔識寫個死字」.

至於法利賽人,相信我們不會感到太陌生.雖然他們人數不多,又不似撒都該人有正統的宗教領袖地位,但他們的影響力是很大的,也廣受百姓尊敬.為甚麼呢?因為他們在宗教上是很熱誠的,也是摩西律法專家,並者非常之有理想,他們教導百姓如何按處境遵守摩西的法律,他們不斷地抄寫律法,研究及應用律法,要建立一群上帝喜悅的子民.他們不只空口說白話,為到實踐宗教理想,他們過著並不富足的生活,他們更加身體力行,並且東奔西跑,常到猶太人的聚會點,即是福音書中的會堂,去教導並傳揚律法.

施洗約翰得罪法利賽人,其實是冒很大的險的.因為法利賽人在百姓心中有很權威的地位,若然法利賽人號召百姓不要聽從施洗約翰的道理,這條出來要跟隨施洗約翰的長龍,分分鐘可以走了一半人.並且,他辛苦賺來的美好聲譽,很可能會嚴重受損.

施洗約翰的膽子真的很大,與這兩批人作對,等於得罪統治者,得罪可毀他聲譽的人,換來的下場,可能是把他打回原形,本來已到人生收成期,得來不易,結果被絕對清零.

事實上,在稍後的日子裡,施洗約翰因為繼續堅持講真話,見義為之,堅持實踐上帝所交付的天國使命,就拿起勇氣,直接斥責統治者希律犯了奪人妻子的罪,結果他被收監,最後被斬頭而死.

我相信施洗約翰不會不知道,凡事折衷一點,妥協一點,人生會好過一點,至少可以為他避開一些劫數.

正所謂君子不立危牆,退一步海闊天空.

不過,在施洗約翰的生命裡,什麼都可以退,信仰原則就不可退,半步都不可退.

他從一個窮小子,在時代裡搖身一變成為大人物,結果又被打回原形.

他因為立心要在彎曲勃謬的世代裡堅持上主所交託的使命,在大時代裡最後變得一無所有,兩手空空而終結一生.

這樣去做天國的見證人,值得嗎?這樣盡責做上帝的見證人,是忠心還是愚蠢呢?

我相信他被斬首後回天家之路上,是微笑地對著自己說,「感謝主,你讓我成為那盡了責的見證人.」

他那無畏無懼的勇氣是從哪裡來的呢?看看第11節,謎底就解開了.

原來這個成功人士,享有高度聲譽的施洗約翰,看自己根本不配得到那麼多跟隨者,和那麼高的地位.

他說,「11我是用水給你們施洗,叫你們悔改；但那在我以後來的,能力比我更大,我就是給他提鞋子也不配,他要用聖靈與火給你們施洗.」

「我就是給他提鞋子也不配」,

施洗約翰自覺自己連替耶穌挽鞋都不配.他擁有的,都只是因為上帝的恩典.

他有,但他不擁有,因為一切都是從上主而來,他把從上主而來的都歸給主.人到無求品自高,事能知足心常泰.

弟兄姊妹,原則上,每個基督徒都是施洗約翰,無後顧之憂,方能見義勇為.

在將臨期,在這「準備期」與「等待期」裡,上帝要我們知道,我們都應該是連替主挽鞋都不配的施洗約翰.

我們都應該是施洗約翰,不戀棧所擁有的,時刻要見義勇為.不過這真是很難的,因為當人擁有得越多,就會越懂得計算,就會越事事小心,凡事折衷一點,妥協一點.

我們都是人生贏家.

也許,在大時代裡,我們都需要一首詩歌提醒我們:

〈願那靈火復興我〉

願那復興我的靈火,重新再來一次復興我,

起初的信心何處失落,我願悔改重新再得著,

在明媚的春光裡,在幸福的生活中,

我們陶醉,我們歡樂,早把主恩典忘記.

願那復興我的靈火,重新再來一次復興我,

起初的希望何處失落,我願悔改重新再得著,

在穩妥的事業裡,在迷人的笑聲中,

我們遊戲,我們沉迷,早把主恩典忘記.

願那復興我的靈火,重新再來一次復興我,

起初的愛心何處失落,我願悔改重新再得著,

在艱苦的歲月裡,在百般的試煉中,

我們禱告,我們警醒,毋把主恩典忘記.

願那復興我的靈火,重新再來一次復興我,

起初的信望愛何處失落,我願悔改現在就得著.



\newpage

\section{凡不因我跌倒的就有福了}
\label{sec:627}
\textbf{姚志豪牧師 - 凡不因我跌倒的就有福了!}
\newline
\newline
\hyperref[sec:626]{< < < PREV SERMON < < <}
~
\hyperref[sec:toc]{[返目錄]}
~
\hyperref[sec:628]{> > > NEXT SERMON > > >}


講題:凡不因我跌倒的就有福了!

經文:詩146:5-10,賽35:1-10；太11:2-11,雅5:7-10

講員:姚志豪牧師

日期:2022年12月11日

凡不因我跌倒的就有福了!(那對我不懷疑是多麼有福啊!) Blessed is he who takes no offense at me

幸福是不斷懷疑,不斷拷問與揚棄,是一種「通過苦惱的歡欣」,而不是簡單的信仰.而這種「高層次的懷疑」,原是對自己而不是任何人,任何權勢,任何意識形態教條的忠誠,對自己也不是指外在的名利權,而是內在的良知.以信仰某種價(上帝,和平,學術)為藉口逃避(懷疑)是愚蠢的,因為你逃不了.由懷疑,拷問,苦腦,揚棄而導致的認知的改變,會有歡欣,但只是過程,而未必是歸宿.

「人類對於外在世界和自我的認識,都是透過對一切知識,觀念,經驗的不斷懷疑才能獲得.意見不同的各種對話,真是每一次都長知識.」——李怡《重拾優雅》

懷疑:總是被人有種傾向負面的感覺,香港的教育(中小幼)制度,好少鼓勵因懷疑去問,最好就是坐定定,跟住照做,有好成績,大概就是好學生的先決條件.倘若提出疑問(過多),總是比人有種攪事的感覺.然而,有位老牧者同我分享,所謂學問,學問,要學就要問；有懷疑就更加要問.

將臨期是聖誕前有既四個主日,新教PROTESTANT傳統每個主日也有一個主題:重生,聖經,施洗約翰(第三主日),馬利亞.施洗約翰亦是今天經文所提到對耶穌產生懷疑的那位.今天的講題是選自和合本修定版的太11:6「凡不因我跌倒的就有福了!」,現修版:「那對我不懷疑是多麼有福啊!」跌倒/懷疑又好,耶穌講這句說話,正代表著,信他的人對他懷疑是正常不過的事,一點也不出奇!為何一個歡天喜地迎接聖誕的來臨,提到懷疑?

經文一開始就提到施洗約翰是位在囚人士,(太4:12)就是當耶穌被試探40天後,就聽到約翰下監,什麼事?之後第十四章1-13節提到,就是因為他批評希律王的婚姻.這個希律王叫安提帕,是大希律的兒子.他跟妻子離婚,然後娶了兄弟的妻子希羅底,這是很不道德的,更犯左猶太法規.由於安提帕喜歡希羅底,希羅底又喜歡他,就將批評他的施洗約翰拉去坐監,他更下令斬了約翰頭.皇帝都敢批評,以下犯上,惹來殺生之禍是必然的.而今天讀的經文太11章,就是在他未被殺之前,他聽過耶穌的事,就派自己的門徒去見耶穌,問:「那將要來的是你嗎?還是我們等候別人呢?」我想:監內無理由只聽到這事,不過,明顯這事也是約翰一直也有留意著,所以約翰知道有一拯救者要來,自己也會想?耶穌是否就是要等果位彌賽亞(拯救者)?於是就派門徒人去見耶穌,但顯明他也是充滿疑問.因為他要門徒問:「那將要來的是你嗎?還是我們等候別人呢?」,有趣的是,耶穌沒有直接的回答,耶穌的回答是猶太人耳熟能詳拯救者會做的事:「就是瞎子看見,瘸子行走,長大痲瘋的潔淨,聾子聽見,死人復活,窮人有福音傳給他們.」為什麼會這麼說?如果作為一個猶太的老師,大概對舊約以賽亞書35:5-6也會知道的「那時,瞎子的眼必睜開；聾子的耳必開通.那時,瘸子必跳躍像鹿；啞巴的舌頭必能歌唱.在曠野必有水發出；在沙漠必有河湧流.」詩146:「耶和華開了瞎子的

眼睛；耶和華扶起被壓下的人.耶和華喜愛義人.」經文中,我們會見到好多的顛覆記述:太11章,瞎眼可見,痲瘋的潔淨,聾子聽見,這班全是罪人,只因猶太人認為,要不你犯罪,要不你父母,否則怎麼會有病?窮人有福音傳給他們,正是代表,窮人被拒門外囉,就好似(華人與狗不得內進),倒過來說,福音窮人點會有份呢?若果他們被醫治,即代表罪可被赦免,邊個有權赦罪?在他們心目中,一個不像君王的耶穌,說出這顛覆了傳統的想法,點會不被質疑呢?

《循規踰矩》有罪?無罪?——法庭講求證據與你們所想的見證是不同的.我們不關心你知多少關於基督教的東西,談及多少信仰的事情,參與多少宗教活動,亦不關心你那些以口,以筆,以音樂,以藝術來進行的敬拜.我們只會懼怕那些會放下畫筆,離開舒適的教堂,押上自己生命,如像那一位危險人物耶穌一樣,去顛覆世界,煽動人心,去衝擊我們和諧,挑戰我們主權的人.直至你與耶穌一樣,危險到一個地步,寧願放棄自己生命,膽敢挑戰這個不可動搖的系統,直至你明知沒有可能,仍願意以你微弱的性命來成為我們的一根刺,在此之前,你都不會成為我們的敵人,朋友!

一個來自猶太小城拿撒勒的木匠的兒仔,破壞傳統,亂說話,他能厲害到哪裡去?還不是搗蛋?

福音書,以賽亞及詩篇的對比(瘸子–跳躍；聾子—開通；瞎子—睜開；曠野—有水；沙漠—有河湧流……),正講出,當慣常遇上異常,會造成差異.就好像世界盃,日本兩度勝出歐洲強隊德國及西班牙,慣常支持者的差異當然是失望；然而被看低的一隊的差異,是反映出全國的歡謄.亞洲/香港人,也感到爭光了.或許,這種所謂的差異,其實是幾十年前日本球隊己開始培養/蘊藏了,只是作為強者沒有察覺.而那些被視為邊緣的一群同樣也一直祈盼,甚至祈求上主,這種轉變的來臨,同樣那高高在上的不以為意,亦視而不見.然而,上主,一直也在看顧,也在聆聽祈求.當有一天所謂差異來臨之時,只是上主恢復應有的平等世界,對強者不習慣是因為他以為自己永遠也是強者,忘記左在一切權勢之上,還有以愛及公義掌管世界的上主.

生活要讓差異進入,甚至乎要製造差異,可能會叫人不習慣,因為差異叫人見到其他/無限可能,差異引發人思考,看法/傳統/習慣是否只得一種,別無選擇?差異亦讓人開啟另一扇窗:就是希望.對於一直被看輕,甚至被視罪人,被排擠在外的小眾,耶穌對約翰門徒的回應,豈不是挽回他們生命一點希望,嘗到人性應有的尊嚴嗎?

當約翰的門徒問耶穌,路加7:18-26補充,原來當時耶穌正在醫治人,換言之,即不是只有講,沒有做.或許門徒如此快的離開,不知是否已釋除/加多了一些疑慮?事實上,門徒一句都無回應,但無論如何,相信他們都好想將這種懷疑/看見盡快告知監裡的約翰.耶穌陣他走之前再補一句話:「凡不因我跌倒的就有福了!(那對我不懷疑是多麼有福啊!)」

進來崇基,因著懷疑,最初感到好困擾,為何與我在教會所認識的信仰有很大的差異,我形容教會好比在馬場上跑圈要戴上眼罩的馬,目標清晰向前；而神學院就是除了眼罩的馬,四野遼闊,起初不知方向,後來卻感到這種自由.因為開闊緊我視野,我就發現自己是渺滄海之一栗,視野開闊使我開闊胸襟,當胸襟被開闊,你個人自然會謙卑.當我經驗過懷疑之後,我感到是有福的.

當帶著疑惑的門徒走了之後,耶穌就轉身向另一班群眾說……..從前你們出去是想看什麼?你們從前出到曠野是要看甚麼呢?要看風吹動的蘆葦嗎?你們出去到底是要看甚麼?要看穿細軟

衣服的人嗎?那穿細軟衣服的人是在王宮裡.

當時印在希律年代的錢幣上的就是蘆葦,而細軟(華麗,於宮庭宴樂渡日的衣服).因此,耶穌問你們既然出來,要看的是什麼?你們有沒有找錯地方?今天,我們所找的信仰又是一種什麼價值的信仰?是有顛覆世界力量的福音?還是只講「幸福音」?

8歲時開始去教會,每年傳統的平安夜聽到的信息大概也叫人安心快樂,因平安的來臨.當2019年的12月,我再看同一段經文的平安夜晚,就是不一樣.神學訓練叫我思考,這種安舒,潔淨的馬槽,似乎對於我認識聖經當下的世界不太真實,馬槽的真身,其實是因為他們一家要報名上冊沒有地方下的選擇,馬利一亞已急住生產(路2:1-7),即使動物排泄馬槽也無法.生下以後,當夜(太2:13)因為嬰孩耶穌一家被希律政權追殺,以致被迫逃往埃及.聖誕的故事豈有如此的浪漫?浪漫大概就是與商業掛了勾.然而,正是我有當下處境讀進這經文之時,我才發現,真正的平安,不是無事無幹節日的期待,有禮物,大餐的平平安安,而是在急難,壓迫中,仍感受到上主同在,這種想像/讀入,會讓人疑惑/懷疑嗎?

有個潮語叫:「懷疑人生」是指因遇到重大打擊或認為事情發展不符合常理,令人生觀產生巨大改變.懷疑,疑惑,不是要停止我們思考信仰的原因,相反,正是因為疑惑,令人繼續去尋問,就好像約翰一樣.肉身被囚的約翰,最大的囚禁不是肉身,或許是傳統/思想的限制.孫寶玲牧師說:囚牢是通往福音的一扇門,在等待的囚獄,跌倒和困惑是免不了的.然而,這是福音呼喚的開始,或許,在囚的我們,才開始覺知瞎眼,瘸腿,耳聾,貧乏,才真正渴望瞎子看見,瘸子行走,長大痲瘋潔淨,聾子聽見,死人復活,窮人有福音傳給他們的喜樂.

最後:我引用雅各書:「弟兄們哪,你們要忍耐,直到主來.看哪,農夫忍耐等候地裡寶貴的出產,直到得了秋雨春雨.你們也當忍耐,堅固你們的心,因為主來的日子近了.」5:7-8可能,當我們看不見土下的種子變化時,我們也會產生疑問,甚至覺得白費心機,但我們會否相信,他正在花芽,生根,有一天,那根甚至乎迫開壓迫它的水泥而出呢!

詩146篇:他為受屈的伸冤,賜食物與飢餓的.耶和華釋放被囚的,耶和華扶起被壓下的人.使惡人的道路彎曲.最終,如賽35:7發光的沙( – 或譯:蜃樓)要變為水池；乾渴之地要變為泉源.在野狗躺臥之處,必有青草,蘆葦,和蒲草.或許我們要更深的思考是:為何要疑惑上主,疑惑他就是那真正能顛覆人傳統,思考,使黑暗,混亂的世代,重現秩序,重見光明的那一位.

將臨第二主日點的(喜樂之燭),不是因有慶祝節目來臨而喜樂,而是在黑暗中存有盼望而喜樂,就如捷克捷克共和國的第一位民選總統哈維爾說:希望不是相信事情會好起來的信念,而是不管結果如何,仍確信事情是有意義的.



\newpage

\section{噢流星雨}
\label{sec:628}
\textbf{鄧瑞強博士 - 噢,流星雨}
\newline
\newline
\hyperref[sec:627]{< < < PREV SERMON < < <}
~
\hyperref[sec:toc]{[返目錄]}
~
\hyperref[sec:629]{> > > NEXT SERMON > > >}


馬太福音 1:18-25
\newline
\begin{longtable}{cl}
\hline
\hline
章節 & 經文 (和合本修訂版)\\
\hline
1:18 & \begin{tabularx}{0.7\textwidth}{X} 耶穌基督降生的事記在下面：他母親馬利亞已經許配給約瑟，還沒有迎娶，馬利亞就從聖靈懷了孕。 \end{tabularx} \\ \\ \relax
1:19 & \begin{tabularx}{0.7\textwidth}{X} 她丈夫約瑟是個義人，不願意當眾羞辱她，想要暗地裡把她休了。 \end{tabularx} \\ \\ \relax
1:20 & \begin{tabularx}{0.7\textwidth}{X} 正考慮這些事的時候，忽然主的使者在約瑟夢中向他顯現，說：「大衛的子孫約瑟，不要怕，把你的妻子馬利亞娶過來，因她所懷的孕是從聖靈來的。 \end{tabularx} \\ \\ \relax
1:21 & \begin{tabularx}{0.7\textwidth}{X} 她將要生一個兒子，你要給他起名叫耶穌，因他要將自己的百姓從罪惡裡救出來。」 \end{tabularx} \\ \\ \relax
1:22 & \begin{tabularx}{0.7\textwidth}{X} 這整件事的發生，是要應驗主藉先知所說的話： \end{tabularx} \\ \\ \relax
1:23 & \begin{tabularx}{0.7\textwidth}{X} 「必有童女懷孕生子；人要稱他的名為以馬內利。」（以馬內利翻出來就是「神與我們同在」。） \end{tabularx} \\ \\ \relax
1:24 & \begin{tabularx}{0.7\textwidth}{X} 約瑟醒來，就遵照主的使者的吩咐把妻子娶過來； \end{tabularx} \\ \\ \relax
1:25 & \begin{tabularx}{0.7\textwidth}{X} 但是沒有和她同房，直到她生了兒子，就給他起名叫耶穌。 \end{tabularx} \\ \\
[1ex]
\hline
\hline
\end{longtable}
講題:噢,流星雨

經文:馬太福音1:18-25

講員:鄧瑞強博士

日期:2022年12月18日

當夜空漆黑,大地昏暗時,抬頭望天,偶然見到一顆流星,已開心不已,若然見到流星雨,則更喜出望外.流星雨裡的眾多流星,都是從外太空同一個源頭來的,它們以平行的軌道,闖向地球,與大氣層摩擦,產生平行的,令人驚嘆的光的軌跡.外太空的東西,闖向地球,我們便見到流星.類比地說,當神聖世界的東西,闖入我們的世界,與我們世界慣常的事物產生摩擦時,也會產生耀眼的,令人驚嘆的光的軌跡,我們稱之為神蹟,這是從神聖世界而來的流星.在馬太福音的聖誕故事裡,神聖世界密集地闖進我們的世界,譬如:聖靈闖進馬利亞的生命,讓她懷孕；天使闖進約瑟的夢境,讓他接納馬利亞；天上有特別的星光,引領東方的博士；神聖的指引,也闖進博士們的夢境裡；然後,天使不斷指引約瑟當行的事.神聖世界的事物,不斷闖入我們的世界,劃破人間漆黑的夜空,閃耀出聖誕故事特有的光芒.在漆黑穹蒼下的我,只能讚嘆地說:噢,流星雨!好美的流星雨!

今日的講道,與你們分享這流星雨景象的一個片段.我想,你們也會讚嘆:噢,好美的流星雨!

今日的講道經文是:馬太福音1:18-25

太1:18耶穌基督降生的事記在下面:他母親馬利亞已經許配了約瑟,還沒有迎娶,馬利亞就從聖靈懷了孕.

太1:19她丈夫約瑟是個義人,不願意明明地羞辱她,想要暗暗地把她休了.

太1:20正思念這事的時候,有主的使者向他夢中顯現,說:「大衛的子孫約瑟,不要怕!只管娶過你的妻子馬利亞來,因她所懷的孕是從聖靈來的.

太1:21她將要生一個兒子,你要給他起名叫耶穌,因他要將自己的百姓從罪惡裡救出來.」

太1:22這一切的事成就是要應驗主藉先知所說的話,

太1:23說:必有童女懷孕生子；人要稱他的名為以馬內利.(以馬內利翻出來就是「神與我們同在」.)

太1:24約瑟醒了,起來,就遵著主使者的吩咐把妻子娶過來；太1:25只是沒有和她同房,等她生了兒子,就給他起名叫耶穌.

第一幕:馬太福音1:18

太1:18耶穌基督降生的事記在下面:他母親馬利亞已經許配了約瑟,還沒有迎娶,馬利亞就從聖靈懷了孕.

馬太的聖誕故事,入局得很快.聖靈闖入一個女子的生命裡,她就從聖靈懷了孕.這不是一個普通的懷孕故事,這是神聖介入人間的故事.神聖介入人間,就是將新的元素帶入人類歷史裡,改變歷史的軌跡,以致歷史能邁向神聖的終點.

歷史會陷入低潮,會墮落,會沉入黑暗的時刻裡,歷史之上的神聖上帝會按其旨意在關鍵時機裡介入人間,讓人間在漆黑裡仍看到盛載著神蹟的流星,免得人類在漫長的黑暗中走向虛無.

馬利亞懷著孕,有一天,會生下耶穌.耶穌生命的底蘊,耶穌生命的內在精神,耶穌要開顯的生命內容,卻是聖靈.聖靈是神的靈,是生命的靈,是天地間的正氣.「靈」與「氣」,在聖經的用字裡,是同一個字.耶穌乃神聖之正氣的精華,也將以天地的正氣驅散人間的黑暗.聖誕節,紀念這天地的正義貫注人間.

馬利亞已許配了約瑟,還未迎娶,卻懷了孕,於是,將皮球交給了約瑟,看看約瑟如何處理.

第二幕:馬太福音1:19

太1:19她丈夫約瑟是個義人,不願意明明地羞辱她,想要暗暗地把她休了.

不要忽略經文這裡對約瑟的描述:「他是個義人」.這裡點出了馬太福音的一個信仰主題:成為一個義人.馬太福音最出名的「登山寶訓」討論「八福」,其中第四福是:「飢渴慕義的人有福了!因為他們必得飽足」(太5:6)；第八福是:「為義受逼迫的人有福了!因為天國是他們的」(太5:10).耶穌在「登山寶訓」裡說:「我告訴你們,你們的義若不勝於文士和法利賽人的義,斷不能進天國」(太5:20).

馬太福音的結尾部分(第25章),講到對人類的終極判決,就是分判義人和惡人.義人可進天國.在馬太福音第25章的經文那裡,義人問:「主啊,我們什麼時候見你餓了,給你吃,渴了,給你喝?什麼時候見你作客旅,留你住,或是赤身露體,給你穿?又什麼時候見你病了,或是在監裡,來看你呢?」主回答說:「我實在告訴你們,這些事你們既做在我這弟兄中一個最小的身上,就是做在我身上了.」(太25:37-40)能不能進天國,在這段經文裡,重點在義人的行為.

如何成為義人,是馬太福音信息的重點.在馬太福音裡,耶穌兩次提到:「我喜愛憐恤,不喜愛祭祀」(太9:13,12:7).一個人如何能勝過文士和法利賽人的義,就是對人的憐恤,悲憫.

回到我們的聖誕故事:約瑟,對著這懷了孕的未婚妻,他如何做?

一般猶太男人,懷疑自己的女人紅杏出牆,自己「戴了綠帽」,可按傳統律法,公開休棄妻子,讓妻子當眾受辱,甚至被石頭打死.約瑟卻一反常態,選擇保護這女子的名聲,暗暗地休了她,保留她的生機.

在馬太福音裡,耶穌在登山寶訓中常常如此說:「傳統的律法如此如此說,我卻這樣說」,而耶穌說法的重點,是對人的憐恤,悲憫.聖誕故事中的約瑟,其實是提前地實踐耶穌的教訓.他不是按律法傳統寫下的文字,來處理他的未婚妻,而是憐恤她,向她展示一種純粹的悲憫.對人的憐恤,大於跟從律法傳統寫下的無血無肉的死文字.

第三幕:馬太福音1:20-23

太1:20正思念這事的時候,有主的使者向他夢中顯現,說:「大衛的子孫約瑟,不要怕!只管娶過你的妻子馬利亞來,因她所懷的孕是從聖靈來的.

太1:21她將要生一個兒子,你要給他起名叫耶穌,因他要將自己的百姓從罪惡裡救出來.」

太1:22

這一切的事成就是要應驗主藉先知所說的話,

太1:23說:必有童女懷孕生子；人要稱他的名為以馬內利.(以馬內利翻出來就是「神與我們同在」.)

神聖的事物又介入人間,改變歷史的軌跡了.天使闖入約瑟的夢境.在古代,夢是上天向人啟示的中介.天使對約瑟說:「不要怕」.這是神聖人物與世人相遇時的典型開場白.神聖的闖入,打破了人間的常態,引發驚訝,是可以預期的.天使說「不要怕」,讓人平靜下來,迎接新的事態.天使向約瑟解釋,馬利亞所懷的胎的不尋常.她要生的兒子,要給他起名叫耶穌.「耶穌」這個名字,即舊約的「約書亞」.耶穌要做的事,與約書亞所做的事,有相似性.約書亞是摩西的接班人.摩西無法領神的子民進入神所應許的「迦南地」,這事是由約書亞完成的.當日,在「迦南地」的邊界,神的子民大部分都說,不進「迦南地」了,因打不過迦南人.他們決意回埃及,重新成為法老王的奴隸,回到歷史的黑暗裡.是約書亞,獨排眾議,力挽狂瀾,將神的百姓,從黑暗帶到光明,將他們從罪惡的權勢中救出來.耶穌,將會以新的方式,再進行約書亞那力挽狂瀾的工程,將自己的百姓從罪惡裡救出來.

回到聖誕故事:天使對約瑟說,這一切事,是要應驗以賽亞先知的一段經文.這經文出自以賽亞書7:14,該經文的背景是:敘利亞和以色列聯手攻打猶大,猶大的先知以賽亞說,有一年輕女子懷孕生子,當那孩子還未懂得作道德判別時,敘利亞和以色列便會被消滅,猶大便會得救.那孩子有一富有象徵意義的名字,人給他起名叫「以馬內利」,其意思是:神與我們同在.如今,以賽亞先知的信息,在耶穌身上,以一種全新面貌進一步地應驗.童女馬利亞,會生下耶穌,他的出生,意味著神的子民面對的危機會終極地獲得解決.若你深入探究:在黑暗中的人類,憑什麼獲得終極的光明?你會發現,其終極理由是:神願意與人同在,光明願意照入黑暗,神聖世界願意闖進人間.就因如此,我們才能看見流星.

從聖誕故事看,我們的生活世界,不是一封閉的世界,而是一向神聖世界開放的世界.神聖世界也不是一封閉的世界,而是一向人間世界開放的世界.塵世的事務,我們自有當盡的責任.我們要力圖成為一義人,但因著塵世的種種原因,我們總有力不能勝的地方.黑暗的世界,有種種黑暗的勢力,人在其中,無法突破.好在,我們的世界不是只有我們,我們的世界向上帝開放.神聖的正氣闖入人間,生下耶穌,以世間本來沒有的神聖恩典,以人類自身無法亮起的光明,驅散人間的黑暗,將人類從黑暗中救出來.聖誕節,是一天人交疊的圖畫,因上天的正氣的介入,人能秉承這正氣,作一正義的人.聖誕故事,呼喚我們去明白天心,去成為義人,去悲憐蒼生.

第四幕:馬太福音1:24-25

太1:24約瑟醒了,起來,就遵著主使者的吩咐把妻子娶過來；

太1:25只是沒有和她同房,等她生了兒子,就給他起名叫耶穌.

約瑟,順服天命,娶了馬利亞,讓這女子有家.他給孩子起名,就是公開地接納這孩子為他兒子,讓這孩子進入約瑟的族譜內,成為大衛的後裔.

結語:

整個故事,是神聖世界不斷介入人類世界的故事,上帝不斷地介入人間.我們的生活,被上帝廣闊無邊的憐憫包圍著.黑暗的大地,被外來的神聖恩典不斷闖入.噢,流星雨!好美的流星雨!

願榮耀歸與在至高之處的神,願平安歸與神所喜悅的人.



\newpage

\section{在沒有聖景之中尋找聖誕意義}
\label{sec:629}
\textbf{林豪恩先生 - 在沒有聖景之中尋找聖誕意義}
\newline
\newline
\hyperref[sec:628]{< < < PREV SERMON < < <}
~
\hyperref[sec:toc]{[返目錄]}
~
\hyperref[sec:630]{> > > NEXT SERMON > > >}


講題:在沒有聖景之中尋找聖誕意義

講員:林豪恩先生

日期:2022年12月25日

一．引言

你有沒有見過崇基學院禮拜堂前草地陳設的聖景呢?根據校友們說,早於六十年代,禮拜堂外便擺放馬槽聖景,每年由將臨期開始擺放至元旦後.參照耶穌出生的情景,聖景中有一間以芒草搭建的小屋,屋內有馬槽,也曾經有耶穌,瑪利亞,約瑟,牧羊人和其小羊,博士等人物,亦會掛上一顆星,象徵引領博士前往耶穌出生之地的景星「伯利恆之星」.根據崇基校園通訊所述,馬槽小屋由物業管理處搭建,用作屋頂的芒草就是於中大校園收集.後來因為校園芒草漸少,加上收割芒草的安全問題,因此無法再搭建聖景.

這個有馬槽,耶穌,瑪利亞,約瑟,牧羊人和其小羊,博士等人物的聖景,是綜合了馬太福音和路加福音的聖誕故事.博士和那顆星是來自馬太福音的；牧羊人是來自路加福音的.基督徒常常問,這是否合乎聖經呢?這個聖景,是有經文支持的.那麼,現在不lor 設聖景,是不是不合乎聖經呢?又不是,因為今年聖誕節主日的福音經課所講述的聖誕故事,就是沒有我們熟識的聖景；又或者說,是一個不一樣的聖景.這個沒有聖景的聖誕故事,仍然有聖誕意義嗎?如有,是怎樣的聖誕意義呢?讓我們以經文互讀的方法,把約翰福音1:1-14節和創世記互相參照,尋找沒有聖景的聖誕意義.

東方博士在哪裡尋找到生下來作猶太人之王的基督呢?馬太福音指向伯利恆的小屋；曠野的牧羊人又在哪裡尋找到救主基督呢?路加福音指向馬槽.約翰福音從另一個角度啟示我們,基督不在伯利恆小屋,也不在馬槽.那麼,基督在哪裡呢?

二．內容

1.在萬物之上尋找基督

1:1太初有道,道與上帝同在,道就是上帝.

1:2這道太初與上帝同在.

1:3萬物是藉著他造的；凡被造的,沒有一樣不是藉著他造的.

約翰福音指出1:1太初有道,你會想到甚麼呢?熟識舊約的讀者會想起創世記一章一節:

1:1起初,上帝創造天地.

約翰福音的「太初」與希臘文創世記的「起初」是同一個字.約翰福音提示我們,回到「起初」尋找基督,回到萬物被創造之前尋找基督.基督在約瑟和馬利亞之先,在希律王和東方博士之前,在牧羊人與馬槽之外,在天使與那星之上.基督在萬物出現之前已經存在,約翰福音開宗明義指出基督就是上帝,萬物都是祂創造的.然而,人如何能在祂的創造中尋找基督呢?

詩篇十九篇的詩人就是透過上主創造的天地感悟上主的同在:

1諸天述說上帝的榮耀,穹蒼傳揚他手的作為.

2這日到那日發出言語,這夜到那夜傳出知識.

歷世歷代有不少人也被大自然的浩瀚或美景震撼心靈而感悟上主,也有人透過研究受造物的奇妙而步向上主.世界頂尖遺傳工程科學家法蘭西斯‧柯林斯 Francis S. Collins本來是一位無上帝論的科學家,經過30年的探索,在宇宙論,演化論和基因研究中,看見了上帝的容貌.探索自然的科學,不是信仰的衝突,而是榮耀與讚美造物主的語言.

2.在黑暗中尋找光

1:5光照在黑暗裡,黑暗卻不接受光.

「黑暗卻不接受光」或譯「黑暗卻沒有勝過光」,熟識創世記的讀者,就會想到創造的開始:

創世記

1:2地是空虛混沌,淵面黑暗；上帝的靈運行在水面上.

1:3上帝說:「要有光」,就有了光.

1:4上帝看光是好的,就把光暗分開了.

在光出現之前,遍地黑暗,上帝說「要有光」,然後,「光照在黑暗裡」.約翰福音指出現實狀況,雖然「光照在黑暗裡」,「黑暗卻不接受光」,不過,縱使不接受光,「卻沒有勝過光」.現實上黑暗雖然不接受光,但創世記指出「上帝看光是好的,就把光暗分開了」.從此之後世上除了有黑暗,也有光明,人也能分辨得出黑暗和光明.約翰福音提示我們基督道成肉身臨在世間,是光明再次照射在黑暗裡,是創世記的重現,是重新創造.約翰提示我們,在哪裡尋找到基督呢?在光明之中.或許有人說,世界黑暗,縱使如此,約翰呼籲人在黑暗裡尋找光明.的確,很多人都不喜歡黑夜,越黑的夜越冷,越孤單,伸手不見五指,看不清前方的路,就越迷茫.似乎一切有關黑夜的東西都成了被嫌棄的對象,可是看到過妙不可言的星光的人,就知道越黑的夜,越能看到最美的星空.有人說,越黑暗的地方越需要光；越黑暗的地方,發光的生命越顯得珍貴.很多人都是在人生的艱難時候尋找上帝,並且相信上帝的.這是如何發生的呢?因為感到失望不安的人們,發現身邊基督徒的仍然有平安和盼望,於是進一步去認識給予基督徒平安和盼望的上帝.

3.在物化了的人類中尋找生命

1:4生命在他裡頭,這生命就是人的光.

「生命在他裡頭」,又提醒熟識創世記的讀者想起生命的來源:

創世記

2:7耶和華上帝用地上的塵土造人,將生氣吹在他鼻孔裡,他就成了有靈的活人,名叫亞當.

第一次創造的時候,上帝為出自塵土的人賦予生命靈氣.可惜人類的發展過程中,往往自我貶為塵土,互相把對方物化,彼此變為工具.一個人,就好像一座機器中的一粒小螺絲.或許,你已經很少在人身上看到光輝,更多時候,反而看到醜惡和黑暗.或者你曾經聽過有人慨嘆,生存有如行屍走肉,或者投訴有只有生存而沒有生活.這些慨嘆和投訴,都反映出人不滿足於被物化,內心一直在追求另一個層次的東西.約翰福音指出,基督來臨,就是再次創造,就正如當初賦予聖靈給人一樣,把物化了的人重新活化,把人生命裡的光重新點燃.你又尋找甚麼呢?是不是尋找一個有生命的自己呢?

4.在人間尋找上帝

1:14道成了肉身,住在我們中間,充充滿滿地有恩典有真理.我們也見過他的榮光,正是父獨生子的榮光.

天地萬物都是基督所創造的,即是說,基督在萬物之外,人有沒有可能在世間尋找到基督呢?按約翰福音所說,是有可能的.因為「道成了肉身,住在我們中間」.其實,基督本來就與人一齊,並且知道人的需要,關心人的需要,也滿足人的需要.創世記告訴讀者,本來是這樣的:

創世記

2:18耶和華上帝說:「那人獨居不好,我要為他造一個配偶幫助他.」

2:19耶和華上帝用土所造成的野地各樣走獸和空中各樣飛鳥都帶到那人面前,看他叫甚麼.那人怎樣叫各樣的活物,那就是牠的名字.

2:20那人便給一切牲畜和空中飛鳥,野地走獸都起了名；只是那人沒有遇見配偶幫助他.

2:21耶和華上帝使他沉睡,他就睡了；於是取下他的一條肋骨,又把肉合起來.

2:22耶和華上帝就用那人身上所取的肋骨造成一個女人,領她到那人跟前.

上帝就好像父母關懷兒女一樣關懷人,與人同在,明白人的需要,也滿足人的需要,特意為人預備同伴.既然上帝與人那麼親近,為甚麼人尋找不到上帝呢?

5.在世界尋找失散的兒女

1:10他在世界,世界也是藉著他造的,世界卻不認識他.

1:11他到自己的地方來,自己的人倒不接待他.

1:12凡接待他的,就是信他名的人,他就賜他們權柄作上帝的兒女.

1:13這等人不是從血氣生的,不是從情慾生的,也不是從人意生的,乃是從上帝生的.

約翰福音揭示,「他在世界,世界也是藉著他造的,世界卻不認識他」,不但如此,「他到自己的地方來,自己的人倒不接待他」.怎會這樣的呢?又讓我們想起創世記:

3:8天起了涼風,耶和華上帝在園中行走.那人和他妻子聽見上帝的聲音,就藏在園裡的樹木中,躲避耶和華上帝的面.

原來,人躲避上帝.躲避慣了,就越來越陌生；越來越陌生,就越來越拒絕；越來越拒絕,就越來越容不下上帝在世界之中,也容不下上帝在自己的生命之中.

3:9耶和華上帝呼喚那人,對他說:「你在哪裡?」

今天,當人說找不到上帝的時候,耶穌基督來臨人間,或許仍然在對人說創世記那一句話:「你在哪裡?」

三．總結

這就是約翰福音的聖誕故事,是一個道成肉身來尋找兒女的故事,是一個等候兒女回應的的故事.

或者有人會問,上帝在天上開聲叫就可以了,何必要降生成人,進入這個污煙瘴氣的世界呢?

從前有一位首相,他周遊列國時受了基督福音的感化作了基督徒.回國以後,他告訴國王他找到了那位來到世界拯救罪人的救主,並且作了基督徒.皇帝問他說:『我若要作什麼事,只要吩咐僕人,他就去作.那位萬王之王的上帝,可用一句話救人,何必親自成為人的身體呢?』首相便求皇帝賜他一天時間,然後回答.

回家之後,首相叫來一個精巧木匠,要他照著皇帝周歲孩子之形像,刻個木偶,穿上美衣,好好裝飾明天備用.

第二天皇帝和首相坐在船裡,要他回答昨天的問題.這時宮女來了,懷中抱著木偶,站在岸上,皇帝以為真是他的兒子,伸手去接,宮女早已領了首相吩咐,故意失手,木偶掉到水裡去了.

皇帝立時就要跳下水去,救他兒子.宰相攔阻他說:『王阿,你用不著親自下去呀,吩咐我去還不夠嗎?』王回答說:『不行不行,我要自己去救,別人我不放心,這是父愛!”首相說:『全能的主所以不肯只用一句話救世人,必要親自成為人形,惟一原因,也就是愛!』

在沒有東方博士,也沒有牧羊人的聖誕故事中,我們尋找到道成肉身來尋找我們的上帝.



\newpage

\section{迷霧中如何前行}
\label{sec:630}
\textbf{林豪恩先生 - 迷霧中如何前行?}
\newline
\newline
\hyperref[sec:629]{< < < PREV SERMON < < <}
~
\hyperref[sec:toc]{[返目錄]}
~
\hyperref[sec:631]{> > > NEXT SERMON > > >}


傳道書 3:1-11
\newline
\begin{longtable}{cl}
\hline
\hline
章節 & 經文 (和合本修訂版)\\
\hline
3:1 & \begin{tabularx}{0.7\textwidth}{X} 凡事都有定期，天下每一事務都有定時。 \end{tabularx} \\ \\ \relax
3:2 & \begin{tabularx}{0.7\textwidth}{X} 生有時，死有時；栽種有時，拔出有時； \end{tabularx} \\ \\ \relax
3:3 & \begin{tabularx}{0.7\textwidth}{X} 殺戮有時，醫治有時；拆毀有時，建造有時； \end{tabularx} \\ \\ \relax
3:4 & \begin{tabularx}{0.7\textwidth}{X} 哭有時，笑有時；哀慟有時，跳舞有時； \end{tabularx} \\ \\ \relax
3:5 & \begin{tabularx}{0.7\textwidth}{X} 丟石頭有時，撿石頭有時；懷抱有時，不抱有時； \end{tabularx} \\ \\ \relax
3:6 & \begin{tabularx}{0.7\textwidth}{X} 尋找有時，失落有時；保存有時，拋棄有時； \end{tabularx} \\ \\ \relax
3:7 & \begin{tabularx}{0.7\textwidth}{X} 撕裂有時，縫補有時；沉默有時，說話有時； \end{tabularx} \\ \\ \relax
3:8 & \begin{tabularx}{0.7\textwidth}{X} 喜愛有時，恨惡有時；戰爭有時，和平有時。 \end{tabularx} \\ \\ \relax
3:9 & \begin{tabularx}{0.7\textwidth}{X} 這樣，做事的人在他所勞碌的事上得到甚麼益處呢？ \end{tabularx} \\ \\ \relax
3:10 & \begin{tabularx}{0.7\textwidth}{X} 我觀看神給世人的擔子，使他們在其中勞苦： \end{tabularx} \\ \\ \relax
3:11 & \begin{tabularx}{0.7\textwidth}{X} 神造萬物，各按其時成為美好，又將永恆安放在世人心裡；然而神從始至終的作為，人不能測透。 \end{tabularx} \\ \\
[1ex]
\hline
\hline
\end{longtable}
講題:迷霧中如何前行?

講員:林豪恩先生

經文:傳道書3:1-11

一.

有神也不等於知曉前路

2023將會是怎樣的一年呢?疫情會怎樣呢?經濟會怎樣呢?樓市會怎樣呢?很多人都喜歡對新一年做各種預測,因為人人都希望能夠行在可預知及可預備的路上.我有一位朋友做投資工作,他說在2021年底,絕大部分的金融機構都預測2022年恆生指數最低是22000點,最高是28000點.預測準確度有多少呢?大家看一看自己的強積金過去一年的回報就知道了.基督徒呢?基督徒有神同在,是不是萬事都在預期之內呢?有些基督徒的確是這樣相信的.當你驚訝地跟他們說:「唉,想不到俄羅斯竟然攻打烏克蘭!」他們回應說:「有甚麼值得驚訝?耶穌不是一早說好了嗎?民要攻打民,國要攻打國.(馬太福音24:7)」世上萬事似乎都給他們看透了.剛才我們誦讀的舊約經課傳道書的著名詩句,很多人都喜歡引用這些詩句來證明一切都在預期之內,世事都在掌握之中.

「唉,新冠疫情,真係死了不少人!」

3:2生有時,死有時；栽種有時,拔出所栽種的也有時；

「炮火轟隆,一下子就炸毀了整座建築物!」

3:3殺戮有時,醫治有時；拆毀有時,建造有時；

「男朋友要分手.」

3:4哭有時,笑有時；哀慟有時,跳舞有時；

3:5拋擲石頭有時,堆聚石頭有時；懷抱有時,不懷抱有時；

「丟失了iPhone,慘過風丟失錢包.」

3:6尋找有時,失落有時；保守有時,捨棄有時；

「我和太太在冷戰.」

3:7撕裂有時,縫補有時；靜默有時,言語有時；

3:8喜愛有時,恨惡有時；爭戰有時,和好有時.

朋友,你還煩惱甚麼呢?你還執著甚麼呢?傳道者不是已經說了:「凡事都有定期,天下萬務都有定時」嗎?作為基督徒,你還有甚麼是不知道的呢?你還有甚麼是預測不到的呢?你還為甚麼感到迷茫呢?

如果我們如此應用傳道者的詩句,傳道者可能提出抗議.這些詩句是表達一切都在測期之內,一切都在掌握之中嗎?還是剛剛相反,一切都似懂非懂,似乎在掌握之中,但卻又是掌握之外.且看傳道者的感慨:

3:9這樣看來,做事的人在他的勞碌上有甚麼益處呢?

3:10我見上帝叫世人勞苦,使他們在其中受經練.

3:11上帝造萬物,各按其時成為美好,又將永生安置在世人心裡.然而上帝從始至終的作為,人不能參透.

「萬物各按其時成為美好」,這是對誰來說的?只有誰知道呢?是上帝.人呢?傳道者指出,「然而上帝從始至終的作為,人不能參透」.人不能參透,可以不是甚麼大問題,如果人只是如一隻豬,有食物就食,有水就喝,有一日就過一日.但人不能夠做一隻快樂的豬,因為上帝「又將永生安置在世人心裡」.「永生」,和修本譯為「永恆」,有英文聖經譯為「a sense of past and future」,即是思前想後的本能.人會思前想後,希望能夠找出世事的規律,人生的軌跡,和未來的趨勢.人不是不知道生有時死有時,而是不知道何時生何時死,也不知道死得舒服還是辛苦,不在預期的時間和方式死去是很難受的；人也不是不知道股市有升有跌,痛苦之處是測不透何時升何時跌；人人都知道每個星期都有幾次搞珠搞出六個數字,飲恨的是估不中這六個數字.

我們能夠因此就反過來說人對前路一無所知,完全沒有預測能力嗎?如此說又未免過激了.正如傳道者所說,日光之下無新事,人大概仍是覺得自己對過去未來有某程度的洞見,然而,卻又測不準,看不透,也保證不了.因此,無論是基督徒與否,大家的處境都是相同的,都是在知與不知之間的迷霧中前行.在如此的狀態中前行,有甚麼錦囊呢?

二.

仰望穹蒼,俯瞰大地帶來的安穩

以下是詩人提供給我們的錦囊,他建議我們仰望穹蒼,俯瞰大地:

8:3我觀看你指頭所造的天,並你所陳設的月亮星宿,

8:4便說:人算甚麼,你竟顧念他!世人算甚麼,你竟眷顧他!

8:5你叫他比天使微小一點,並賜他榮耀尊貴為冠冕.

8:6-8你派他管理你手所造的,使萬物,就是一切的牛羊,田野的獸,空中的鳥,海裡的魚,凡經行海道的,都服在他的腳下.

8:9耶和華－我們的主啊,你的名在全地何其美!

生活在不碓定的狀況,行走在看不清前景的迷霧之中,很多人都會感到不安,很多人都希望找到安全和穩妥的感覺.你有沒有經驗過這種安全和穩妥的感覺呢?中學時曾經讀過巴金的一篇短文,名為「繁星」,想不到在人生的某個階段再讀,更有感受.其中兩段如此說:「我愛月夜,但我也愛星天.從前在家鄉,六七月間的夜晚,在庭院中納涼時,我最愛看天空中密密麻麻的繁星.看著那星天,我就會忘掉一切,彷彿就回到了母親的懷裡／／如今在海上,每晚我和繁星相對,我把它們一個個認得很熟了.我躺在艙面上,仰望著天空,深藍色的天空裡,正懸著無數半明半昧的星.船在動,星也在動；它們掛得那麼低,真是搖搖欲墜呢!漸漸地我的眼睛模糊了,我就像看見無數的螢火蟲在我周圍飛舞.海上的夜是柔和的,是靜寂的,是夢幻的.我望著那許多認識的星,我彷彿看見它們在眨眼,我彷彿聽見它們在低聲說話.這時候,我真忘掉了一切.在星的懷抱中,我微笑著,我沉睡著.我覺得自己是一個小孩子,現在睡在母親的懷裡了.」

巴金的安全和穩妥感覺來自哪裡呢?他說了兩次類似的話:「彷彿就回到了母親的懷裡」,「我覺得自己是一個小孩子,現在睡在母親的懷裡了」.詩篇131:2詩人也有類似的感覺:「我使我心安穩平靜,好像母親懷中斷奶的孩子；我的心在我裡面如同斷過奶的孩子.」詩人得到這分安全感的途徑和巴金類似,仰望穹蒼,俯瞰大地,心曠神怡,豁然開朗.這不單是詩人和巴金,也是很多人經驗過的,或許你也經驗過,正如詩歌「默默回望」所唱:「望見星星心裡感安寧」,然後可以進入「回望主恩說不清」的狀態.

三.

安穩在於上主的同在

詩人「觀看你指頭所造的天,並你所陳設的月亮星宿」,就感到釋懷,令他感悟到「人算甚麼,你竟顧念他!世人算甚麼,你竟眷顧他!」啟示錄給予的錦囊向我們保證,那位創造地的上主,並不只是高高在上,而同時也是與我們同在的上主.

21:1我又看見一個新天新地；因為先前的天地已經過去了,海也不再有了.

21:2我又看見聖城新耶路撒冷由上帝那裡從天而降,預備好了,就如新婦妝飾整齊,等候丈夫.

21:3我聽見有大聲音從寶座出來說:「看哪,上帝的帳幕在人間.他要與人同住,他們要作他的子民.上帝要親自與他們同在,作他們的上帝.

21:4上帝要擦去他們一切的眼淚；不再有死亡,也不再有悲哀,哭號,疼痛,因為以前的事都過去了.」

21:5坐寶座的說:「看哪,我將一切都更新了!」又說:「你要寫上；因這些話是可信的,是真實的.」

21:6他又對我說:「都成了!我是阿拉法,我是俄梅戛；我是初,我是終.

在迷霧中前行的我們,難免跌跌碰碰,有眼淚,有死亡,有悲哀,哭號,疼痛,這一切,上主都知道,因此祂要擦去我們一切的眼淚,代表著祂親自的安慰,同時,祂應計把一切都更新.

巴金想像自己「回到了母親的懷裡」就感到安穩,如果那不是想像,而是真實呢?啟示錄宣告,「看哪,上帝的帳幕在人間.他要與人同住,他們要作他的子民.上帝要親自與他們同在,作他們的上帝.」那將會是怎樣的狀況呢?那時的政治,經濟,社會狀態是如何的呢?人與人之間還會互相傷害嗎?我們不知道.啟示錄只讓我們知道的狀況是:「不再有死亡,也不再有悲哀,哭號,疼痛,因為以前的事都過去了.」怎有可能呢?為何不可能呢?因為上主應許「看哪,我將一切都更新了!」有可能嗎?當耶路撒冷被巴比倫摧毀之後,人人都說:『看哪,我們的骨頭枯乾了,我們的指望失去了,我們滅絕淨盡了!』(以西結書37:11).上主給以西結看見異象:『看哪,我必使氣息進入你們裡面,你們就要活過來.6我要給你們加上筋,長出肉,又給你們包上皮,使氣息進入你們裡面,你們就要活過來；你們就知道我是耶和華.』」創造天地萬物的上主,怎不可能把萬物更新,重新創造呢?

四.

迷霧不阻善行

終有一天,上主會把萬物更新；終有一天,我們不在在迷霧中前行.使徒保羅曾說:(哥林多前書13章)「12我們如今彷彿對著鏡子觀看,模糊不清,到那時就要面對面了.我如今所知道的有限,到那時就全知道,如同主知道我一樣.」我們現今模糊不清,就不需要假扮知道一切,也不用尋找和依賴那些自吹知道一切的人,卻需要學習如何在模糊不清中生活.如何在模糊不清的狀況前行呢?我們不需要祈福黨的保證,不需要保藥黨的安慰,我們有我們當做的事情,就是福音經課耶穌基督留給我們的行動錦囊:

25:31「當人子在他榮耀裡,同著眾天使降臨的時候,要坐在他榮耀的寶座上.

25:32萬民都要聚集在他面前.他要把他們分別出來,好像牧羊的分別綿羊山羊一般,

25:33把綿羊安置在右邊,山羊在左邊.

25:34於是王要向那右邊的說:『你們這蒙我父賜福的,可來承受那創世以來為你們所預備的國；

25:35因為我餓了,你們給我吃,渴了,你們給我喝；我作客旅,你們留我住；

25:36我赤身露體,你們給我穿；我病了,你們看顧我；我在監裡,你們來看我.』

25:37義人就回答說:『主啊,我們甚麼時候見你餓了,給你吃,渴了,給你喝?

25:38甚麼時候見你作客旅,留你住,或是赤身露體,給你穿?

25:39又甚麼時候見你病了,或是在監裡,來看你呢?』

25:40王要回答說:『我實在告訴你們,這些事你們既做在我這弟兄中一個最小的身上,就是做在我身上了.』

25:41王又要向那左邊的說:『你們這被咒詛的人,離開我!進入那為魔鬼和他的使者所預備的永火裡去!

25:42因為我餓了,你們不給我吃,渴了,你們不給我喝；

25:43我作客旅,你們不留我住；我赤身露體,你們不給我穿；我病了,我在監裡,你們不來看顧我.』

25:44他們也要回答說:『主啊,我們甚麼時候見你餓了,或渴了,或作客旅,或赤身露體,或病了,或在監裡,不伺候你呢?』

25:45王要回答說:『我實在告訴你們,這些事你們既不做在我這弟兄中一個最小的身上,就是不做在我身上了.』

25:46這些人要往永刑裡去；那些義人要往永生裡去.」

耶穌基督這段說話足夠具體,聽了不會說不知道要做甚麼.「弟兄中一個最小的餓了,給他吃,渴了,給他喝；作客旅,留他住；赤身露體,給他穿；病了,看顧他；在監裡,去看他.」有人說,宗教信仰很神秘,很抽象,很虛幻,很「離地」；事實是,宗教信仰有神秘的部分,也有具體清楚的地方,例如耶穌基督在這段說話提出的行為就夠清楚具體了.我們看不透,測不準來年會怎樣,但不影響我們去做耶穌基督吩咐我們做的這些事情.正如詩歌【我知誰掌管明天】所唱:

我不知明天將如何,每一天只為主活,

我不借明天的陽光,因明天或不晴朗,

我不要為將來憂慮,因我信主的應許,

我今天要與主同行,因祂知前面路程.

有許多未來的事情,我現在不能識透,

但我知誰掌管明天,我也知誰牽我手.



\newpage

\section{看哪我的僕人}
\label{sec:631}
\textbf{李均熊牧師 - 看哪!我的僕人!}
\newline
\newline
\hyperref[sec:630]{< < < PREV SERMON < < <}
~
\hyperref[sec:toc]{[返目錄]}
~
\hyperref[sec:632]{> > > NEXT SERMON > > >}


講題:看哪!我的僕人!

講員:李均熊牧師

經文:以賽亞書 42

日期:2023年1月8日

耶和華的僕人

42:1看哪,我的僕人,我所扶持,所揀選,心裡所喜悅的!我已將我的靈賜給他,他必將公理傳給萬邦.

42:2他不喧嚷,不揚聲,也不使街上聽見他的聲音.

42:3壓傷的蘆葦,他不折斷；將殘的燈火,他不吹滅.他憑信實將公理傳開.

42:4他不灰心,也不喪膽,直到他在地上設立公理；眾海島都等候他的訓誨.

42:5那創造諸天,鋪張穹蒼,鋪開地與地的出產,賜氣息給地上眾人,賜生命給行走其上之人的神耶和華如此說:

42:6「我－耶和華憑公義召你,要攙扶你的手,保護你,要藉著你與百姓立約,使你成為萬邦之光,

42:7開盲人的眼,領囚犯出監獄,領坐在黑暗中的出地牢.

42:8我是耶和華,這是我的名；我必不將我的榮耀歸給別神,也不將我所得的頌讚歸給雕刻的偶像.

42:9看哪,先前的事已經成就,現在我要指明新事,告訴你們尚未發生的事.

起:耶穌以賽亞預言的上主僕人嗎?

今天福音經課焦點乃是紀念耶穌基督的受洗.而先知書的經課則選上以賽亞書42:1-9作對應.以賽亞書42:1-9實為以賽亞書四首僕人之歌的第一首(其餘三首分別為49:1-7;

50:4-11和信徒較熟識的52:13-53:12).耶穌基督的受洗標誌他地上事奉的開始,而聖教會將以賽亞書的經文與福音書對應,正表示耶穌就是那一個先知書所說的「僕人」,耶穌的事奉正好成就了舊約的「預言」!

當然我也相信耶穌就是以賽亞書所說的僕人,但是先知書是特定「預言」幾百年後一個叫耶穌的人物,還是這段經文是指向其他人的呢?又或說,若先知書指定要應驗這個「預言」就是／只是耶穌,那對我們今日的信徒又有甚麼關係?

承:以賽亞書中所論的上主僕人是誰?

以賽亞書經文所說的「僕人」究竟是誰?其實歷代學者對此皆有不同的見解,基督教會傳統當然認為這是指著基督耶穌,但從經文的上文下理看,近代聖經學者就提出有不同的可能性.例如在賽44:28提到波斯君王居魯士乃上主的牧人,而居魯士平定四方,不就是經課一開始所說「將公理傳給萬邦」嗎?但一位外邦君王又怎會成為上主的僕人,所以學者轉而看其他經文,例如之前一章的賽41:8-10就明言「惟你以色列—我的僕人,雅各—我所揀選的,我朋友亞伯拉罕的後裔」,那麼「雅各」就是上主的僕人.但「以色列／雅各」是指創世記裡的先祖雅各,還是以此名字代表當時的以色列人呢?若我們再看以賽亞書論到僕人的經文,上主的僕人就更為明顯.若我們看後小小,在42:19-20就說「誰比我的僕人眼瞎呢?誰比我所差遣的使者耳聾呢?誰瞎眼像那獻身給我的人?誰瞎眼像耶和華的僕人呢?看見許多事卻不領會,耳朵開通卻聽不見.」這段會否令你想起以賽亞書第6章上主對以賽亞去宣告的「禍音」:你去告訴這百姓說:你們聽是要聽見,卻不明白；看是要看見,卻不曉得.要使這百姓心蒙脂油,耳朵發沉,眼睛昏迷；恐怕眼睛看見,耳朵聽見,心裡明白,回轉過來,便得醫治.以賽亞書43:8-10就有相似的說法:你要將有眼而瞎,有耳而聾的民都帶出來!…耶和華說:你們是我的見證,我所揀選的僕人.

承:第二以賽亞與以色列被擄

若綜合這些經文來看,以賽亞書這部份所論的上主僕人就是上主的子民以色列.不錯,這些以色列人就是上主所揀選雅各的後裔.但是,這些以色列民因為犯罪(「耳聾眼瞎」)而受上主審判,經歷亡國被擄之苦,但現在卻要被上主使用,成為祂的僕人,正如42:6所言:我－耶和華憑公義召你,要攙扶你的手,保護你,要藉著你與百姓立約,使你成為萬邦之光!

以賽亞書42章的前文後理,相信是由40章至55章這十六章組成的大單元,現代學者多數稱之為「第二以賽亞書」.整個背景是論到身處異邦,國破家亡的以色列人,應該如何理解上主在歷史和世界中的作為.耶和華／上主是以色列的上帝,為何祂容許國破家亡和國民被擄到異邦的事情發生?是否祂敵不過外邦的假神?以色列人又應該如何面對他身處的處境?

以色列人亡國被擄,是因為他們心靈遠離上主,在社會中不單不顧念體恤孤兒寡婦,只顧自己發財,而且以為只要聖殿不倒和獻祭不停,上主必然保佑國家,免受外敵欺凌；於是上主使用亞述和巴比倫等帝國去侵略他們,分散他們,使他們國破家亡.這些所有的審判和懲罰,其實國中的先知一早就指出,並構成以賽亞書的頭39章(「先前的事」,賽42:9及48:3-6).但現在一切都過去,上主透過當時波斯帝國居魯士王的勢力,要使以色列的命運復興,使祂的子民回歸故土.這個復興救贖計劃的視野則構成餘下27章的以賽亞書(「新事」,賽42:9及48:6-8節).

這個歷史背景讓我們明白到,受審判亡國被擄失落的以色列人是如何成為上主之僕.若果你是公司老闆,員工下屬工作表現完全不達你的標準,你會如何處置?很多人會將他／她開除了事.但上主這位老闆卻沒有放棄這個員工,卻透過歷煉期望他／她能達到工作的標準,成為理想的員工.其實,上主不視我們為祂的員工,而是視我們為祂的拍檔,並非因為祂子民的不足就放棄他們.

在上世紀60年代之前,美國底特律地的汽車工業的生產量是首屈一指的,但到了後來日本汽車的生產量已經超越美國,而美國人也寧願買日本製的汽車而不買本土產品.為甚麼會這樣呢?80年代末就有人提出這是因為日本用上一種叫作「精益生產」(lean production)的管理方法.長話短說,精益生產其中一個與我們今天經課有關的理念,就是部件生產商生產出現不理想,日本的汽車公司並不是替換了生產商了事,而是成立精益小組,協助和訓練生產商和生產隊伍生產不用多餘加工和零瑕疵的部件,從而將生產線裝配效益最大化.車廠視部件生產商為伙伴,而不是除時置換的外判公司.

正如第6節我－耶和華憑公義召你,要攙扶你的手,上主在子民失敗和失落的時候並非放棄他們,而是幫助他們成功,期望他們成為他的理想僕人.亡國被擄一方面是對子民的叛逆的審判,但也是更新煉淨上主子民的機會,使上主揀選的子民,能真正成為上主理想的僕人.

轉:理想之僕

究竟上主理想的僕人有甚麼特質呢?

以賽亞書42:1-9節,論到上主的僕人如何將公理傳給外邦:「他不喧嚷,不揚聲,也不使街上聽見他的聲音.壓傷的蘆葦,他不折斷；將殘的燈火,他不吹滅.他憑真實將公理傳開.他不灰心,也不喪膽,直到他在地上設立公理,海島都等候他的訓誨.」可見這個作萬邦之光的僕人,並不是以君臨天下的救世主心態出現,卻是默默地去幫助受壓受傷的人,去看見上主的幫助和工作.因為這位僕人並非「惦晒」才去傳遞上主救恩的信息,因為他也經歷亡國被擄的創傷(trauma),而同時也經歷從上主而來的復興.

七個「不」

若你數一數,在2-4節其實出現了七個「不」.這七個「不」可分成三組,代表著上主僕人三種特質.

「他不喧嚷,不揚聲,也不使街上聽見他的聲音」並不是,上主僕人並不會開佈道會,又或不會大聲疾呼叫人悔改,這裡並不是說上主之僕如何運用他的傳理策略.「喧嚷」原文意思是在「困苦中求助」(cry out for help in distress)意思.所以僕人不使街上聽見他的是甚麼「聲音」呢,我認為是他因痛苦而有的哀鳴(lament),呻吟(moaning)或呼喊(crying out).換句話說,這一組三個「不」,是說到上主之僕並不因痛苦而「伊伊喔喔」(murmuring),他不會因經歷苦難而喋喋不休(今日所謂的「放負」).上主之僕不記苦楚,在困難中默默承擔.

第二組「不」是「壓傷的蘆葦,他不折斷；將殘的燈火,他不吹滅.」蘆葦是代表埃及(36:9/19:6),而燈火是代表巴比倫(43:16-17).兩者都是曾經壓迫過以色列的強國,但這兩個國家將會受到波斯帝國的攻擊,同樣受到被消滅的威協.但上主僕人卻不會因此落井下石,卻能不計前嫌,且能溫柔面對這些曾對付過自己的國家.

第三組「不」是「他不灰心,也不喪膽……」最後,

上主的僕人並不會灰心喪膽,在亡國被擄的處境下,在困難的境況中,他仍能堅持成就上主的召命.

所以上主理想之僕有三個特質:一,不記苦楚,二,不記前嫌,三,不會放棄.但經歷創傷,甚至是別人加害的創傷,上主的理想僕人如何可以做到不去記念而能堅持不休呢?

主賜靈力

在42:1-9節中,原文其實出現了四個「賜」字.一個說到上主不會「賜」他的榮耀給別神(第8節,和合本譯「不『將』我的榮耀歸給別神」)；另一個是上主「賜」下僕人作「萬邦之光」(第6節,和合本譯作「『使』你成為萬邦之光」)；其中值得留意的,是兩個論到上主「賜」下祂的靈(第1節「我已將我的靈賜給他」和第5節「賜氣息…又賜生命」原文只得一個「賜」字).上主的靈代表著上主的能力臨到僕人的身上,使他縱然不揚聲不喧嚷,仍能達成上主的心意!上主僕人的姿態好像處於被動(不喧嚷,不揚聲),但可以努力的堅持(不灰心,不喪膽),皆因上主的僕人不單首先領受上主的「公理」(第1,3,4節皆有出現,原文可解作「秉公辦理」,「申訴得直」)(我們的希望?),更因為得到了上主的匡扶(第6節,「我……要攙扶你的手」).是上主屬靈屬天的能力醫治了僕人的創傷,是上主的公理使僕人面對加害者能不計前嫌,是上主的扶助使僕人能努力不懈,勇敢前行.

而在這段經文的第7節,指出上主要僕人作萬邦之光,其目的是要「開瞎子的眼,領被囚的出牢獄,領坐黑暗的出監牢.」「瞎子」並非因眼疾不能見的人,而是因為他活在古代作為監倉的地洞,那裡暗無天日,進去的人就如瞎子不見物一樣.故此,這裡的三句,其實都是指使人從監禁人的牢獄中得到自由,就如被擄的國家從被擄中得回自主的權力一樣.這不正正是上主僕人以色列的親身經歷?要作萬邦之光,自己要先得釋放見光明,才可領人出牢獄捆鎖.

靈恩/秉公行義:皇室身份!

是甚麼人能得賜這靈力?是甚麼人被要求秉公行義的呢?

在古代近東「君權神授」的處境中,是皇室人員才會得這靈力,去代表上主在世間秉公行義,為被屈者伸冤.因此,這是表示上主的理想僕人,其實是有著皇室的身份.這皇室身份標示以色列若跨越創傷,成為上主理想之僕,就能成為上主管治這世間的代理人,能代表上主實現祂在世的旨意(「願祂的旨意行在地上如同在天上」).

其實先知經課所言的僕人,是一個上主理想的僕人.當然祂期望在被擄苦難創傷洗禮的以色列人能承擔此使命,但很可惜,歷史的實然卻叫我們看到這並未完全實現在以色列人身上.但這個理想之僕的特質,卻完全顯示在耶穌基督的身上:皇室身份,披帶上主的靈恩,不念苦楚,愛仇敵,愛到底,彰顯甚麼是公義和真理,最終成就了上主救贖的使命.耶穌基督應驗了先知的說法,並非因為先知在預言一個獨特的人物.先知指示的是上主要扶持祂的子民,成為祂理想的僕人.但這個理想,終於在神子耶穌身上得著完全的實現.

合:如何成為萬邦之光?

正是因為先知所言的是個理想的化身,雖然在耶穌身上得著完全的實現,但今日的我們,也可以成就這一個理想,成為上主僕人.跟從基督這個理想之僕的腳蹤,我們也能成為上主理想僕人.

這一刻,我們或可反思自問,我們的軟弱和破碎如何可以成就上主的僕人呢?我們今天需要再進一步培訓成為上主理想的僕人,還是我們靠上主所賜的靈力而已經超越我們的苦難,創傷,仇恨,軟弱和不足,領受皇室的身份,憑信實將公理傳開,使自己成為上主理想僕人,成為萬邦之光?

[Matthew Perry (扮演著名美劇Friends 中Chandler Bing一角的演員)自傳作為例證]



\newpage

\section{如此風流人物}
\label{sec:632}
\textbf{鄧瑞強博士 - 如此風流人物}
\newline
\newline
\hyperref[sec:631]{< < < PREV SERMON < < <}
~
\hyperref[sec:toc]{[返目錄]}
~
\hyperref[sec:633]{> > > NEXT SERMON > > >}


申命記 17:16-19
\newline
\begin{longtable}{cl}
\hline
\hline
章節 & 經文 (和合本修訂版)\\
\hline
17:16 & \begin{tabularx}{0.7\textwidth}{X} 只是王不可為自己加添馬匹，也不可為加添馬匹使百姓回埃及去，因耶和華曾對你們說：『不可再回那條路去。』 \end{tabularx} \\ \\ \relax
17:17 & \begin{tabularx}{0.7\textwidth}{X} 王不可為自己多立妃嬪，免得他的心偏離；也不可為自己多積金銀。 \end{tabularx} \\ \\ \relax
17:18 & \begin{tabularx}{0.7\textwidth}{X} 他登了國度的王位之後，要在利未家的祭司面前，將這律法書為自己抄寫一份在書卷上。 \end{tabularx} \\ \\ \relax
17:19 & \begin{tabularx}{0.7\textwidth}{X} 這書要存在他那裡，他一生的年日要誦讀，好使他學習敬畏耶和華－他的神，謹守遵行這律法書上的一切話和這些律例， \end{tabularx} \\ \\
[1ex]
\hline
\hline
\end{longtable}
講題:如此風流人物

經文:申命記17:16-19;列王記上3:1,9:15-21,24,10:26-11:9

講員:鄧瑞強博士

日期:2023年1月15日

今日,我們來到《列王記》的經文.《列王記》原本不分上下卷,但因篇幅長,為方便起見,分成上下兩卷.記述的內容:起自大衛王崩,所羅門繼位；至猶大國亡,被巴比倫所滅.亡國後,猶大的神學家在異鄉,以事後回顧的歷史高度,反思亡國的神學理由.故事的背景是大衛的興起,他成為王,是出於神的揀選,神的恩典,神又以立約的形式去確立這揀選的恩典.大衛是天命所歸的人,同時,他也要承擔起天命的託負,就是神人盟約裡人應盡的責任,這些責任都寫在《申命記》裡了.這王朝之後開展的歷史,卻是人慢慢忘記天命的歷史.在君王的私慾裡,在人民的墮落中,在國際形勢的壓迫下,神的子民逐漸與世俗同化,拜世俗的偶像,倚靠世俗的軍事同盟,而忘記了自身作為神的子民的獨特身分,忘記了自身獨特的天命,忘記了神,於是,神就讓巴比倫消滅他們.當失去了身分,失去了天命,失去了上帝之後,神的子民變得清醒.他們在心靈的廢墟中,重新思考這一切往事,然後,寫下《列王記》,訴說一個民族墮入虛無的理由,這理由正正就在於他們忘記了神的約.忘記了神的約,也即是忘記了人生在世的意義就是將神的旨意成就在人間.

我們簡短地重溫一下《申命記》對王的基本要求——申命記17:16-19:

申17:16只是王不可為自己加添馬匹,也不可使百姓回埃及去,為要加添他的馬匹,因耶和華曾吩咐你們說:「不可再回那條路去.」

申17:17他也不可為自己多立妃嬪,恐怕他的心偏邪；也不可為自己多積金銀.

申17:18他登了國位,就要將祭司利未人面前的這律法書,為自己抄錄一本,

申17:19存在他那裡,要平生誦讀,好學習敬畏耶和華－他的神,謹守遵行這律法書上的一切言語和這些律例.

「王不可為自己加添馬匹」:這裡的「馬匹」是戰爭工具,王不應窮兵黷武,不應以武力壓迫他人.

「不可使百姓回埃及去,為要加添他的馬匹」,可以理解成:不應為了強化王的權力或武力,而使百姓淪為奴隸,這樣做,就如將百姓帶返埃及一樣.

「不可為自己多立妃嬪」:以前多立妃嬪,多與政治的聯盟有關.妃嬪太多,聯盟太多,就會令王的眼光集中於軍事保障,而忘記上帝的仁義.

「不可為自己多積金銀」:過份積累財富,過份看錢份上,當然也會令人走向權力與暴力,而忘卻天命.

作為一個王,他最重要的,是敬畏神,讓神的旨意成就在人間,讓天國臨到大地,讓百姓活在神的恩典裡.

《列王記》的起點是所羅門王的登基,我們看看所羅門王是否符合《申命記》所講的王的標準.這些標準以三樣具體事物表示:馬匹(特別是埃及的馬匹),妃嬪,財富.很快,我們便發現,所羅門王完全不合格.《列王記》以所羅門王作為一具體例子,說明猶大國至終的消亡,在於其開始之時,已出現嚴重的朽爛.

我們看看《列王記上》對所羅門一生功業的描述.

所羅門王登基後,逐一消滅那些會影響他王位的人,之後,正式介紹他的第一句經文,便是以下這經文.

王上3:1所羅門與埃及王法老結親,娶了法老的女兒為妻,接她進入大衛城,直等到造完了自己的宮和耶和華的殿,並耶路撒冷周圍的城牆.

約略地看這經文,已大感不妙.所羅門的第一項政績,便是與法老結盟.在聖經文學裡,法老很多時代表奴役.接著,他將法老的女兒,穿堂入室地接入大衛城,讓她進駐王權的中心.再看看所羅門建造工程的排位,自己的宮竟排在耶和華的殿之前.這些寫法,都是不祥之兆.

王上9:15所羅門王挑取服苦的人,是為建造耶和華的殿,自己的宮,米羅,耶路撒冷的城牆,夏瑣,米吉多,並基色.

王上9:16先前埃及王法老上來攻取基色,用火焚燒,殺了城內居住的迦南人,將城賜給他女兒所羅門的妻作妝奩.

王上9:17所羅門建造基色,下伯‧和崙,

王上9:18巴拉,並國中曠野裡的達莫,

王上9:19又建造所有的積貨城,並屯車和馬兵的城,與耶路撒冷,黎巴嫩,以及自己治理的全國中所願建造的.

王上9:20至於國中所剩下不屬以色列人的亞摩利人,赫人,比利洗人,希未人,耶布斯人,

王上9:21就是以色列人不能滅盡的,所羅門挑取他們的後裔作服苦的奴僕,直到今日.

王上9:24法老的女兒從大衛城搬到所羅門為她建造的宮裡.那時,所羅門才建造米羅.

一個王是否守約,履行天命,第一個標準是看他如何處理「財富」的事情.以上經文,牽涉到所羅門如何累積「財富」的情況,其中最引人注目的字眼是「積貨城」,這字眼曾出現在《出埃及記》,那時,埃及法老要以色列人做奴隸,就是要他們蓋「積貨城」.如今,世界輪流轉,所羅門慢慢變成了法老,在他底下,有奴隸為他蓋「積貨城」.一個要從法老的奴役下走出來的民族,如今,竟然變成了奴役別人的民族.在神子民的國度,竟然有奴隸.在這段文字的結尾部分,刻意再提及「法老的女兒」.看來,所羅門已與法老的血脈,連成一體.

王上10:26所羅門聚集戰車馬兵,有戰車一千四百輛,馬兵一萬二千名,安置在屯車的城邑和耶路撒冷,就是王那裡.

王上10:27王在耶路撒冷使銀子多如石頭,香柏木多如高原的桑樹.

王上10:28所羅門的馬是從埃及帶來的,是王的商人一群一群按著定價買來的.

王上10:29從埃及買來的車,每輛價銀六百舍客勒,馬每匹一百五十舍客勒.赫人諸王和亞蘭諸王所買的車馬,也是按這價值經他們手買來的.

一個王是否守約,履行天命,第二個標準是看他如何處理「馬匹」的事情.「馬匹」,就是軍事工具,就是當時的軍火.所羅門大量屯積馬匹,屯積軍火,還做起軍火商,發軍火財.經文特別提到,這些是埃及的馬匹.又提到埃及.埃及的馬匹,是令人成為奴隸的軍事力量,所羅門現在成了這力量的中間人.這中介生意當然會讓所羅門賺很多錢,也令所羅門變得更強大,但卻令他與那位帶領無能者出埃及的神越來越遠了.

王上11:1所羅門王在法老的女兒之外,又寵愛許多外邦女子,就是摩押女子,亞捫女子,以東女子,西頓女子,赫人女子.

王上11:2論到這些國的人,耶和華曾曉諭以色列人說:「你們不可與她們往來相通,因為她們必誘惑你們的心去隨從她們的神.」所羅門卻戀愛這些女子.

王上11:3所羅門有妃七百,都是公主；還有嬪三百.這些妃嬪誘惑他的心.

王上11:4所羅門年老的時候,他的妃嬪誘惑他的心去隨從別神,不效法他父親大衛誠誠實實地順服耶和華－他的神.

王上11:5因為所羅門隨從西頓人的女神亞斯她錄和亞捫人可憎的神米勒公.

王上11:6所羅門行耶和華眼中看為惡的事,不效法他父親大衛專心順從耶和華.

王上11:7所羅門為摩押可憎的神基抹和亞捫人可憎的神摩洛,在耶路撒冷對面的山上建築邱壇.

王上11:8他為那些向自己的神燒香獻祭的外邦女子,就是他娶來的妃嬪也是這樣行.

王上11:9耶和華向所羅門發怒,因為他的心偏離向他兩次顯現的耶和華－以色列的神.

一個王是否守約,履行天命,最後一個標準是看他如何處理「妃嬪」的事情.所羅門的妃嬪實在多得驚人,有1000人,世上竟有如此風流人物.若他每晚見一個,也要兩年零九個月才見得完.這些婚姻,十居其九都是為了建立國際關係而締結的政治婚盟.這些女子來的時候,也帶著她們的神明到來.政治的聯合無可避免地伴隨著宗教的混和,智慧如所羅門,最後也離棄神.這妃嬪名單,以法老的女兒為首,這當然標誌著所羅門越來越與法老連結了.在他治下,百姓的生活光景,彷彿又回到埃及一樣,而這正是最違反君王天職的事情.

所羅門王死後,猶大國便一分為二,分成北部的以色列國和南部的猶大國.分裂的理由,與出埃及的理由一樣,就是被奴役的人不想再做奴隸.在所羅門治下,人民吃了大苦,苦中求變,於是,所羅門一死,以色列便有十個支派脫離大衛王朝,另立以色列國去了.

透過這樣的描述,身處巴比倫的神學家剖視猶大亡國的神學理由.理由很簡單,就是王不成王.王,應是神子民的典範,領受上天的恩典,在人間,作神旨意的僕人.神透過王,去管治大地,讓世人回復人性,不被物化而變成奴隸,以致人人能盡其性.猶大國的敗亡,在於帝王的骨子裡,有變成法老的基因.帝王傾向與各方政治結盟,多過與神結盟；傾向倚靠武力,多過實踐仁愛；傾向累積財富,多於讓人安息.這不是實踐上帝的天心,而是根本地反上帝.這傾向一路發展下去,最後,便是猶大的亡國,一切化為烏有.

若我要講猶大亡國的神學理由,來到這裡,我大概可以結束了.但我更感興趣的,是那些流放他鄉的猶太人,他們是如何構想民族的將來的.在寫作及編輯《列王記》時,這些神學家在字裡行間,留下了他們對將來的不同構思.我要再說一次,《列王記》背後,有一群神學家,他們在反思歷史的同時,構想著將來.這些構想,不是明明白白地說出來的,而是隱藏在歷史敘事的字裡行間之中的.這話從何說起.大家有沒有留意,我今日講的所羅門,與大家平時對所羅門的觀感不大一樣,有點格格不入似的.一般而言,我們欣賞所羅門王,《列王記》有別的經文顯示,所羅門懂得向神求智慧；他又以智慧判案,最出名的案件是兩個母親爭一個嬰孩的故事,他以智慧伸張了公義；他又為神建聖殿,榮耀神的名.這一切都甚好,贏盡世人掌聲.但,為何《列王記》卻同時記下我今日講道用的那些非常負面的經文?我猜想(純粹猜想),這是由於寫作及編輯《列王記》的神學家們,代表著不同群體,他們對於如何開展將來的歷史,各有不同想法.例如:將來應以某種方式延續大衛王朝的統治嗎?有人明顯認為不應該,這些人可能是掃羅王的後人,更有可能是同被擄去異鄉的北國以色列人.他們在《列王記》留下的文字,是強烈認為所羅門及其繼位者正是王國分裂的始作俑者,他們是民族的千古罪人.

另一相關問題是:若不由大衛王朝統治,將來的猶大民族仍保留君王體制嗎?明顯,有人不單反對大衛王朝,連君王體制也劇烈反對,這些人認為君王的權力集中正是所有問題的根源.這些人對所羅門的王權的描述,非常負面.但另外有些人,卻認為沒有君王體制就會進入無政府狀態,這只會造成混亂,於是,他們留下了對所羅門王相當正面的描述.

但如何在堅持君王體制的同時卻能避免君王的墮落呢?有人便提出「智慧君王論」,這些人在《列王記》裡特別強調所羅門王的智慧.「智慧君王論」有別於傳統的「盟約君王論」.「盟約君王論」就是說,神揀選某人,與他訂立盟約,於是,這人便成了天命所歸的君王.問題是:大衛王朝裡的所有君王,原則上都是天命所歸的君王,但不行天命的,卻佔多數.這些人提出「智慧君王論」,就是反對傳統「天命所歸」那種君權神授的想法,而強調君王本身的智慧,當然,最高的智慧是敬畏耶和華.敬畏神的人,才能成為君王.來到這裡,我們見到《列王記》背後,有不同群體,他們對於如何運作將來的世界,有不同的構想.他們將這些不同想法,寫入《列王記》的歷史敘述中,在字裡行間,隱隱現出對將來的豐富想像.

過往我們讀聖經,習慣認為聖經在具體問題上會給我們單一的答案,但現在的聖經研究,卻重視聖經本身的「多聲道」.聖經的文字,是不同群體,不同觀點的平等申述.不同群體各自表述他們的看法,不曾試圖去取消他人的聲音,用俄國思想家巴赫金(Bakhtin)的講法,這就是「眾聲喧嘩」(heteroglossia).面對歷史的前景,經文容讓不同的想像.不同群體可以探索不同的可能性,而無須扼殺任何聲音.同時,我們看到,縱使《列王記》對歷史前景有不同的想像,但總意卻是在探索:人如何能將神的旨意更好地實現在大地上,好使天國早日臨現人間.今日,我們如何能在歷史的錯綜複雜中保持敬畏神的心,同時又保持對不同聲音的開放性?

願榮耀歸上主,願人間享和平.



\newpage

\section{天國近了悔改吧}
\label{sec:633}
\textbf{高國雄牧師 - 天國近了,悔改吧!}
\newline
\newline
\hyperref[sec:632]{< < < PREV SERMON < < <}
~
\hyperref[sec:toc]{[返目錄]}
~
\hyperref[sec:634]{> > > NEXT SERMON > > >}


馬太福音 4:12-23
\newline
\begin{longtable}{cl}
\hline
\hline
章節 & 經文 (和合本修訂版)\\
\hline
4:12 & \begin{tabularx}{0.7\textwidth}{X} 耶穌聽見約翰下了監，就退到加利利去； \end{tabularx} \\ \\ \relax
4:13 & \begin{tabularx}{0.7\textwidth}{X} 後來離開拿撒勒，往迦百農去，住在那裡。那地方靠海，在西布倫和拿弗他利地區。 \end{tabularx} \\ \\ \relax
4:14 & \begin{tabularx}{0.7\textwidth}{X} 這是要應驗以賽亞先知所說的話： \end{tabularx} \\ \\ \relax
4:15 & \begin{tabularx}{0.7\textwidth}{X} 「西布倫，拿弗他利，沿海的路，約旦河的東邊，外邦人的加利利— \end{tabularx} \\ \\ \relax
4:16 & \begin{tabularx}{0.7\textwidth}{X} 那坐在黑暗裡的百姓看見了大光；坐在死蔭之地的人有光照耀他們。」 \end{tabularx} \\ \\ \relax
4:17 & \begin{tabularx}{0.7\textwidth}{X} 從那時候，耶穌開始宣講，說：「你們要悔改！因為天國近了。」 \end{tabularx} \\ \\ \relax
4:18 & \begin{tabularx}{0.7\textwidth}{X} 耶穌沿著加利利海邊行走，看見兩兄弟，就是那叫彼得的西門和他弟弟安得烈，正往海裡撒網；他們本是打魚的。 \end{tabularx} \\ \\ \relax
4:19 & \begin{tabularx}{0.7\textwidth}{X} 耶穌對他們說：「來跟從我，我要叫你們得人如得魚一樣。」 \end{tabularx} \\ \\ \relax
4:20 & \begin{tabularx}{0.7\textwidth}{X} 他們立刻捨了網，跟從他。 \end{tabularx} \\ \\ \relax
4:21 & \begin{tabularx}{0.7\textwidth}{X} 耶穌從那裡往前走，看見另外兩兄弟，就是西庇太的兒子雅各和他弟弟約翰，同他們的父親西庇太在船上補網，耶穌就呼召他們。 \end{tabularx} \\ \\ \relax
4:22 & \begin{tabularx}{0.7\textwidth}{X} 他們立刻捨了船，辭別父親，跟從了耶穌。 \end{tabularx} \\ \\ \relax
4:23 & \begin{tabularx}{0.7\textwidth}{X} 耶穌走遍加利利，在各會堂裡教導人，宣講天國的福音，醫治百姓各樣的疾病。 \end{tabularx} \\ \\
[1ex]
\hline
\hline
\end{longtable}
講題:天國近了,悔改吧!

經文:馬太福音4:12-23

講員:高國雄牧師

經文:馬太福音4:12-23

日期:2023年1月22日

中國人在新年傳統會說好多吉利的說話,拜年的時間,你說得最多的是那些呢?如果我們細心留意,用心理解,就會發覺當中可以大概分為幾大類:

祈求保持/存,不變”Keep”:身體健康,青春常駐,身壯力健,龍馬精神,出入平安等…

一些期盼,希望獲取更多(特別是物資上的)”Wish/Take”:恭喜發財,一本萬行,步步高升,心想事成,花開富貴等…

對於未來一年,我們對信仰又有怎樣的期盼呢?

保持?不變已經非常滿足?上帝主權的挑戰?

獲取更多?想要獲取的原因?

中國傳統仍有另一種智慧,新年賀詞還有:一元復始,萬象更新…

年輕時總是不明白,這與我的生活有什麼關係?

不單是去舊迎新,更是復始,更新,有一個回到從前的意思”ReStart/Reset”

所以中國人今天開始是癸卯年(有說是由立春那日開始),由甲子,乙丑到壬戌,癸亥,天干與地支配合共60年,之後又由甲子開始,重點在於重新開始,是一個循環,有回到最初,返到本源…這或許才是每年我們應該有的最大願望!

今天福音書的經課

馬太福音4:12-23

12耶穌聽見約翰下了監,就退到加利利去；13後又離開拿撒勒,往迦百農去,就住在那裡.那地方靠海,在西布倫和拿弗他利的邊界上.14這是要應驗先知以賽亞的話,15說:西布倫地,拿弗他利地,就是沿海的路,約旦河外,外邦人的加利利地－－16那坐在黑暗裡的百姓看見了大光；坐在死蔭之地的人有光發現照著他們.17從那時候,耶穌就傳起道來,說:「天國近了,你們應當悔改!」

耶穌基督所傳的道,從此記載中,非常直接清晰,不轉彎沒角,不婉轉修飾…

極重要的背景:「天國近了!」

其實好直接,我會用廣東俗語「gwo 頭近啦/那邊近了」

可能你會馬上反應…新年流流,仲要在大年初一,簡單幾句輕鬆愉快,互相祝賀的好說話就夠了,「gwo 頭近啦/那邊近了」太不吉利了.

但認真思考一下,人生命有限制(時限),是所有人也經歷及要接受的,新年互相祝賀的說話,其實正正就是道出在限制中希望超越/擺脫,從而過得更好…

有關「天國」是一個不容易理解的概念,簡單來說

天國是以上帝的價值而運作的地方/場所

天國是上帝掌權的國度

天國不等如天堂/樂園,雖然好有可能到末日審判之後,完全了的天國才會實現

「天國近了!」「gwo 頭近啦/那邊近了」

不是說我們各人快要離開世界

(雖然人出生後就是向著死亡例數,每過一年,就是離見主面近一年)

我們正在走向世界/自己的終末

又或以被動式(很多年前本人在基督徒學生福音團契所參與的營會題目)；看明日向我們走來!

耶穌基督以這極具震撼,壓迫感,同時又非常真實的描述來成為傳道服務開始的宣告背景,實在值得我們深思及學習…

「時日無多」是鐵一般的事實,我今天的信息不是要掃大家的興,而是要在重要日子,道出生命中重要的道理…

「天國近了,你們應當悔改!」

今天(在新的一年)我們首要做的,就是要悔改,回轉,是”ReStart/Reset”,更要是”Transform”

耶穌基督不單只留下宣告的內容,更馬上化成行動

4:18耶穌在加利利海邊行走,看見弟兄二人,就是那稱呼彼得的西門和他兄弟安得烈,在海裡撒網；他們本是打魚的.19耶穌對他們說:「來跟從我,我要叫你們得人如得魚一樣.」20他們就立刻捨了網,跟從了他.21從那裡往前走,又看見弟兄二人,就是西庇太的兒子雅各和他兄弟約翰,同他們的父親西庇太在船上補網,耶穌就招呼他們,22他們立刻捨了船,別了父親,跟從了耶穌.

作出呼召與接受呼召

當耶穌基督宣告「天國近了,你們應當悔改!」的同時,他就去到有需要的人當中作出呼召/邀請,這邀請帶來生命及生活的轉變.

所以首先要問,我們在座當中,是否每一位也已經回應呼召你的主,改變了你的生命及生活方式,放下了你的網來跟從他呢?

我更要問當中已經信息多年的弟兄姊妹,你所尋找的信仰,是Keep,Take還是Transform的信仰呢?你願意成為別人的祝福,向身邊親友良朋發出這呼召呢?

與新年慣常的”客套”祝賀不一樣,信仰中的期盼是應該要帶來行動/改變的,因為呼召是有力量的,我們已經接受主的呼召,就應該可以有力量傳下去.

4:23耶穌走遍加利利,在各會堂裡教訓人,傳天國的福音,醫治百姓各樣的病症.

這段經文以此描述作結,今天我也以此作為對大家的祝願:

「經歷主的恩典,呼召,我們就得著力量,以新的生命去走遍主所帶領到的各處,在教會/信仰學習的人當中以聖經彼此服侍,傳揚上帝要掌權這大好的消息,醫治每一個受傷的身心靈.」

祝大家:「新年蒙恩,靈命長進,實踐真理,活證主道!」



\newpage

\section{天堂的位置為誰而留}
\label{sec:634}
\textbf{林豪恩先生 - 天堂的位置為誰而留}
\newline
\newline
\hyperref[sec:633]{< < < PREV SERMON < < <}
~
\hyperref[sec:toc]{[返目錄]}
~
\hyperref[sec:635]{> > > NEXT SERMON > > >}


講題:天堂的位置為誰而留?

講員:林豪恩先生

日期:2023年1月29日

引言

看到這個題目,有人會聯想到從前某位公眾中物的故事:某年她出席公開場合時,有感而發提到:「我今朝有去望彌撒.今朝ge 福音有真福八端.真福八端ge 第八端,有人話好arm 我.第八端就係:為義被難真有福,因為係天堂已經預jor 個位比我.」雖然,「係天堂已經預jor 個位比我」這句話引起不少討論,但我想起的,是另一個故事與我的工作比較有關係的故事,這個故事你很可能聽過:

有一位傳道人,他壽終了；剛巧同一天有一位巴士司機也去世了,但是,傳道人卻下了地獄,而巴士司機卻上了天堂,為甚麼呢?傳道人氣憤的跑去跟上帝理論說:「上帝啊,你很不公平,為什麼我那麼忠心地傳講你的信息,卻下了地獄,那巴士司機開車橫衝直撞卻上了天堂,難道您沒有觀察到嗎?」上帝說:「我看見啊!但是你每次在講道時,台下的教友全都睡著了,而巴士司機每次載著教友時,全車的教友都在虔誠的祈禱,所以….」

天堂的位置為誰而留呢?在今天經課的詩篇,我們讀到詩人提出了類似的問題:

15:1耶和華啊,誰能寄居你的帳幕?誰能住在你的聖山?

我們不知道天堂有沒有留位制度,不敢隨便判定天堂留了位給誰,沒有留位給誰；我們也不是意圖為進天堂設定三個問題,或四條定律；我們只是期望透過今天經課的幾段經文,一齊觀察和思考,在天堂會和誰一齊的呢?在天堂不一定見到我,但一定會見到上帝,對嗎?如果在天堂日日和上帝一齊,上帝喜歡見到怎樣的人呢?反過來問,怎樣的人和上帝一起會感到快樂,而不是辛苦呢?

可能不是這樣的人

1.彌迦先知指出不是那些一面犯罪一面獻祭的人:

6:6我朝見耶和華,在至高上帝面前跪拜,當獻上甚麼呢?豈可獻一歲的牛犢為燔祭嗎?

6:7耶和華豈喜悅千千的公羊,或是萬萬的油河嗎?我豈可為自己的罪過獻我的長子嗎?為心中的罪惡獻我身所生的嗎?

這是反問句,用來表達否定的意思.變為直述句可以如下:

「我朝見耶和華,在至高上帝面前跪拜,當獻上甚麼呢?不可獻一歲的牛犢為燔祭.耶和華不喜悅千千的公羊,不喜悅萬萬的油河.我不可為自己的罪過獻我的長子,不可為心中的罪惡獻我身所生的.」

絕大部分宗教信徒,都以為透過宗教行為或儀式,獻上貴重或大量的供物,就可修復和保證自己和神明之間的關係.如果一隻牛犢不足夠,我可以獻千千的公羊；如果萬萬的油河不足夠,我還可以好像異教徒一樣做得更徹底,獻上自己的長子.正如現代人所說,錢能夠解決的問題,就不是問題了.彌迦先知指出,上帝不接受用錢來解決問題；如果用錢就可以解決問題,那麼我們這些沒有錢的人就一個問題都解決不到.幸好上帝不接受用千千公羊,萬萬油河來解決問題,否則我們這些連一隻牛犢都沒有的就注定被上帝拒諸門外了；又幸好上帝不接受用獻長子解決問題,否則我們這些只生了一個女兒,或者沒有生子女的就永不超生了.

先知們都一致地否定以獻祭來取悅上帝,從撒母耳開始,已經如此說:

15:22撒母耳說:「耶和華喜愛燔祭和祭物,豈如喜愛人聽從他的話呢?看哪,聽命勝於獻祭,順從勝於公羊的脂肪.」

以賽亞更對此大加抨擊:

11耶和華說:「你們許多的祭物於我何益呢?公綿羊的燔祭和肥畜的油脂,我已經膩煩了；公牛,羔羊,公山羊的血,我都不喜悅.

15你們舉手禱告,我必遮眼不看,就算你們多多祈禱,我也不聽；你們的手沾滿了血.

有些宗教信徒以為,香火鼎盛,受供奉的神明就喜悅.先知們指出,耶和華上帝不受香火奉承,也不被供奉賄賂.彌迦揭示獻祭不能掩蓋上帝的眼晴令祂看不到罪惡；以賽亞斷言宗教行為不能洗淨沾滿血的手.正如彼拉多判處耶穌基督死刑之前金盆洗手,但都洗不走把無罪者判為有罪的徇私枉法,以後人們每個主日誦讀使徒信經時,仍然訴說耶穌基督在彼拉多手下受難.先知們指出,一面流人血,一面進行宗教行為的人,只會令上帝感到厭煩和憤怒；其實,這也是一般人的共識.試想想,口蜜腹劍,笑裡藏刀的人,難道你會喜愛嗎?

2.使徒保羅在哥林多前書指出另一些離上帝越走越遠的人,就是那些自恃聰明智慧的人:

1:19就如經上所記:我要滅絕智慧人的智慧,廢棄聰明人的聰明.

1:20智慧人在哪裡?文士在哪裡?這世上的辯士在哪裡?上帝豈不是叫這世上的智慧變成愚拙嗎?

有人認為保羅反智,有人因此而推崇反智.保羅真的反智嗎?保羅繼續說:

1:21世人憑自己的智慧,既不認識上帝,上帝就樂意用人所當作愚拙的道理拯救那些信的人；這就是上帝的智慧了.

1:22猶太人是要神蹟,希臘人是求智慧,

1:23我們卻是傳釘十字架的基督,在猶太人為絆腳石,在外邦人為愚拙；

對於「世人自己的智慧」來說,耶穌事件是一宗醜聞.耶穌徹頭徹尾是一個失敗者,出師未捷身先死,自以為聰明的人怎可能跟隨如此一個失敗者呢?對於當時的社會來說,耶穌是一個重犯,所犯的罪嚴重得被判以釘十字架這極端羞辱的刑罰,有體面的人怎可能信仰一個重犯呢?對於要求神蹟的人來說,耶穌簡直是一個令人失望的笑話.馬太福音記載耶穌被打十字架時,「39從那裡經過的人譏笑他,搖著頭,40說:「你這拆毀殿,三日又建造起來的,救救你自己吧!如果你是上帝的兒子,就從十字架上下來呀!」41眾祭司長,文士和長老也同樣嘲笑他,說:42「他救了別人,不能救自己.他是以色列的王,現在從十字架上下來,我們就信他.」44和他同釘的強盜也這樣譏諷他.」釘十字架的基督,對猶太人來說的確是絆腳石；對智慧人來說是愚拙.

懷著這些希臘人和猶太人想法的人,相信現在也有不少.他們期望一位戰無不勝的基督,一位點石成金的基督,一位比成功人士更成功的基督.他們繼續尋找基督,但仍未尋找得到.未必是基督拒絕這些朋友,有可能是這些朋友離基督卻越行越遠.

是怎樣的人呢?

1.行公義和好憐憫的人

我們看到了彌迦先知所說的令上帝厭煩和憤怒的人,也看到了保羅所說的那些離上帝越走越遠的人.那麼,若要正面回答「15:1耶和華啊,誰能寄居你的帳幕?誰能住在你的聖山?」他們又提供了甚麼參考呢?

詩人的問題其實是設問,他自己是有答案的.且看他提供的答案:

15:2就是行為正直,做事公義,心裡說實話的人.

15:3他不以舌頭讒謗人,不惡待朋友,也不隨夥毀謗鄰里.

15:4他眼中藐視匪類,卻尊重那敬畏耶和華的人.他發了誓,雖然自己吃虧也不更改.

15:5他不放債取利,不受賄賂以害無辜.行這些事的人必永不動搖.

有沒有發現詩人對自己的問題的答案和彌迦先知所說的類似?只是先知在這裡說的較為精簡及概括:

6:8世人哪,耶和華已指示你何為善.他向你所要的是甚麼呢?只要你行公義,好憐憫,存謙卑的心,與你的上帝同行.

彌迦先知指出,與上帝同行的人是行公義的.關於公義,詩人的描述是:「行為正直,做事公義,心裡說實話的人」,詩人舉出的例子是:「他不放債取利,不受賄賂以害無辜」.另外,與上帝同行的人是憐憫的人,關於憐憫,詩人舉的例子有:「他不以舌頭讒謗人,不惡待朋友,也不隨夥毀謗鄰里.15:4他眼中藐視匪類,卻尊重那敬畏耶和華的人.他發了誓,雖然自己吃虧也不更改.」彌迦先知又說,與上帝同行的人是存謙卑的心的,存謙卑的心是怎樣的呢?這特質帶我們回到保羅在哥林多前書的教導.

2.存謙卑的心的人(哥林多前書)

1:26弟兄們哪,想一想你們的蒙召,按著人的觀點,有智慧的不多,有能力的不多,有尊貴地位的也不多.

1:27但是,上帝揀選了世上愚拙的,為了使有智慧的羞愧；又揀選了世上軟弱的,為了使強壯的羞愧.

1:28上帝也揀選了世上卑賤的,被人厭惡的,以及那一無所有的,為要廢掉那樣樣都有的,

1:29使凡血肉之軀的,在上帝面前,一個也不能自誇.

1:30但你們得以在基督耶穌裡是本乎上帝,他使基督成為我們的智慧,成為公義,聖潔,救贖.

1:31如經上所記:「要誇耀的,該誇耀主.」

有智慧的人的容易恃才傲物,強壯的人傾向自命不凡,樣樣都有的人可能空目一切,他們都分享著同一種特質,就是保羅所說的「自誇」；同時,他們都欠缺同一樣的特質,就是彌迦先知所說的「謙卑」.其實,這種謙卑的特質,不是上帝要求我們的,而是上帝自己展示的,保羅在腓立比書說:

2:6他本有上帝的形像,卻不堅持自己與上帝同等；

2:7反倒虛己,取了奴僕的形像,成為人的樣式；既有人的樣子,2:8就謙卑自己,存心順服,以至於死,且死在十字架上

我們回望自己的信仰歷程,「自誇」的時候,往往就是遠離上帝的時候；甚麼時候謙卑,就是靠近上帝的時候.不接受耶穌基督的朋友可能有千百樣理由,很可能其中一種就是自以為是的自恃.跟隨耶穌基督的朋友也可能有千百樣理由,其中一種應該是願意在上帝面前謙卑.

3.八福中的謙卑人

我們今天的福音經課提及的八福,我們都習慣地以「得福」的角度去讀,看看我們能夠得到甚麼「福」.如果用來回答詩人「15:1耶和華啊,誰能寄居你的帳幕?誰能住在你的聖山?」我們會發現除了得到甚麼福之外,更重要的可能是做怎樣的人.

5:3虛心的人有福了!因為天國是他們的.

5:4哀慟的人有福了!因為他們必得安慰.

5:5溫柔的人有福了!因為他們必承受地土.

5:6飢渴慕義的人有福了!因為他們必得飽足.

5:7憐恤人的人有福了!因為他們必蒙憐恤.

5:8清心的人有福了!因為他們必得見上帝.

5:9使人和睦的人有福了!因為他們必稱為上帝的兒子.

5:10為義受逼迫的人有福了!因為天國是他們的.

按八福的描述,和上帝／天國有關的是怎樣的人呢?

虛心的人／哀慟的人／溫柔的人／飢渴慕義的人／憐恤人的人／清心的人／使人和睦的人／為義受逼迫的人

這人是誰呢?不正正就是耶穌基督自己嗎?

2:6他本有上帝的形像,卻不堅持自己與上帝同等；

2:7反倒虛己,取了奴僕的形像,成為人的樣式；既有人的樣子,

2:8就謙卑自己,存心順服,以至於死,且死在十字架上

所以保羅勸勉我們:(腓立比書)

2:5你們當以基督耶穌的心為心:

總結

或者,很多人都想知道「天堂的位置為有沒有為我而留」,又知道進入天堂要滿足甚麼準則.大家有沒有感覺得到今天的經課提示我們從另一個角度思考這個問題,與其問上帝有沒有為我留位,不如問自己,我在那個場景會是怎樣的?那個是上帝的場,即是說我要和祂一起,我會感到舒服嗎?我會感到快樂嗎?我們沒有懷疑那是一個美好的場景,但是,如果我和這個場景格格不入,就未必是一件美事.最近有新聞報導有一位家長交付校長二萬元,要求學校收錄她的兒子.對這位家長的原居地來說,這可能是常態,但當她處身於另一個地方時,就顯得格格不入,並且為她帶來痛苦.人人都想進天國,但如果某人和天國的價值和品質相距太遠,可能會感到格格不入,感到自己是陌生人；如果某人和天國的價值和品質接近,可能會感到舒服,有一種終於來到自己的地方的釋然和自在.願上主幫助我們.



\newpage

\section{如何可以勝過文士及法利賽人的義}
\label{sec:635}
\textbf{高國雄牧師 - 如何可以勝過文士及法利賽人的義?}
\newline
\newline
\hyperref[sec:634]{< < < PREV SERMON < < <}
~
\hyperref[sec:toc]{[返目錄]}
~
\hyperref[sec:636]{> > > NEXT SERMON > > >}


馬太福音 5:13-48
\newline
\begin{longtable}{cl}
\hline
\hline
章節 & 經文 (和合本修訂版)\\
\hline
5:13 & \begin{tabularx}{0.7\textwidth}{X} 「你們是地上的鹽。鹽若失了味，怎能叫它再鹹呢？它不再有用，只好被丟在外面，任人踐踏。 \end{tabularx} \\ \\ \relax
5:14 & \begin{tabularx}{0.7\textwidth}{X} 你們是世上的光。城造在山上是不能隱藏的。 \end{tabularx} \\ \\ \relax
5:15 & \begin{tabularx}{0.7\textwidth}{X} 人點燈，不放在斗底下，而是放在燈臺上，就照亮一家的人。 \end{tabularx} \\ \\ \relax
5:16 & \begin{tabularx}{0.7\textwidth}{X} 你們的光也要這樣照在人前，叫他們看見你們的好行為，把榮耀歸給你們在天上的父。」 \end{tabularx} \\ \\ \relax
5:17 & \begin{tabularx}{0.7\textwidth}{X} 「不要以為我來是要廢掉律法和先知。我來不是要廢掉，而是要成全。 \end{tabularx} \\ \\ \relax
5:18 & \begin{tabularx}{0.7\textwidth}{X} 我實在告訴你們，就是到天地都廢去，律法的一點一畫也不能廢去，直到一切都實現。 \end{tabularx} \\ \\ \relax
5:19 & \begin{tabularx}{0.7\textwidth}{X} 所以，無論誰廢掉這誡命中最小的一條，又教導人也這樣做，他在天國裡要稱為最小的。但無論誰遵行並如此教導人的，他在天國裡要稱為大。 \end{tabularx} \\ \\ \relax
5:20 & \begin{tabularx}{0.7\textwidth}{X} 我告訴你們，你們的義若不勝過文士和法利賽人的義，絕不能進天國。」 \end{tabularx} \\ \\ \relax
5:21 & \begin{tabularx}{0.7\textwidth}{X} 「你們聽過有對古人說：『不可殺人』；『凡殺人的，必須受審判。』 \end{tabularx} \\ \\ \relax
5:22 & \begin{tabularx}{0.7\textwidth}{X} 但是我告訴你們：凡向弟兄動怒的，必須受審判；凡罵弟兄是廢物的，必須受議會的審判；凡罵弟兄是白痴的，必須遭受地獄的火。 \end{tabularx} \\ \\ \relax
5:23 & \begin{tabularx}{0.7\textwidth}{X} 所以，你在祭壇上獻祭物的時候，若想起有弟兄對你懷恨， \end{tabularx} \\ \\ \relax
5:24 & \begin{tabularx}{0.7\textwidth}{X} 就要把祭物留在壇前，先去跟弟兄和好，然後來獻祭物。 \end{tabularx} \\ \\ \relax
5:25 & \begin{tabularx}{0.7\textwidth}{X} 你同告你的冤家還在路上，就要趕快與他講和，免得他把你送交給法官，法官交給警衛，你就下在監裡了。 \end{tabularx} \\ \\ \relax
5:26 & \begin{tabularx}{0.7\textwidth}{X} 我實在告訴你，就是有一個大文錢還沒有還清，你也絕不能從那裡出來。」 \end{tabularx} \\ \\ \relax
5:27 & \begin{tabularx}{0.7\textwidth}{X} 「你們聽過有話說：『不可姦淫。』 \end{tabularx} \\ \\ \relax
5:28 & \begin{tabularx}{0.7\textwidth}{X} 但是我告訴你們：凡看見婦女就動淫念的，這人心裡已經與她犯姦淫了。 \end{tabularx} \\ \\ \relax
5:29 & \begin{tabularx}{0.7\textwidth}{X} 若是你的右眼使你跌倒，就把它挖出來，丟掉。寧可失去身體中的一部分，也不讓整個身體被扔進地獄。 \end{tabularx} \\ \\ \relax
5:30 & \begin{tabularx}{0.7\textwidth}{X} 若是你的右手使你跌倒，就把它砍下來，丟掉。寧可失去身體中的一部分，也不讓整個身體下地獄。」 \end{tabularx} \\ \\ \relax
5:31 & \begin{tabularx}{0.7\textwidth}{X} 「又有話說：『無論誰休妻，都要給她休書。』 \end{tabularx} \\ \\ \relax
5:32 & \begin{tabularx}{0.7\textwidth}{X} 但是我告訴你們：凡休妻的，除非是因不貞的緣故，否則就是使她犯姦淫了；人若娶被休的婦人，也是犯姦淫了。」 \end{tabularx} \\ \\ \relax
5:33 & \begin{tabularx}{0.7\textwidth}{X} 「你們又聽過有對古人說：『不可背誓，所起的誓總要向主謹守。』 \end{tabularx} \\ \\ \relax
5:34 & \begin{tabularx}{0.7\textwidth}{X} 但是我告訴你們：甚麼誓都不可起。不可指著天起誓，因為天是神的寶座。 \end{tabularx} \\ \\ \relax
5:35 & \begin{tabularx}{0.7\textwidth}{X} 不可指著地起誓，因為地是他的腳凳；也不可指著耶路撒冷起誓，因為耶路撒冷是大君王的京城。 \end{tabularx} \\ \\ \relax
5:36 & \begin{tabularx}{0.7\textwidth}{X} 又不可指著你的頭起誓，因為你不能使一根頭髮變黑變白。 \end{tabularx} \\ \\ \relax
5:37 & \begin{tabularx}{0.7\textwidth}{X} 你們的話，是，就說是；不是，就說不是。若再多說，就是出於那惡者。」 \end{tabularx} \\ \\ \relax
5:38 & \begin{tabularx}{0.7\textwidth}{X} 「你們聽過有話說：『以眼還眼，以牙還牙。』 \end{tabularx} \\ \\ \relax
5:39 & \begin{tabularx}{0.7\textwidth}{X} 但是我告訴你們：不要與惡人作對。有人打你的右臉，連另一邊也轉過去由他打。 \end{tabularx} \\ \\ \relax
5:40 & \begin{tabularx}{0.7\textwidth}{X} 有人想要告你，要拿你的裡衣，連外衣也由他拿去。 \end{tabularx} \\ \\ \relax
5:41 & \begin{tabularx}{0.7\textwidth}{X} 有人強迫你走一里路，你就跟他走二里。 \end{tabularx} \\ \\ \relax
5:42 & \begin{tabularx}{0.7\textwidth}{X} 有求你的，就給他；有向你借貸的，不可推辭。」 \end{tabularx} \\ \\ \relax
5:43 & \begin{tabularx}{0.7\textwidth}{X} 「你們聽過有話說：『要愛你的鄰舍，恨你的仇敵。』 \end{tabularx} \\ \\ \relax
5:44 & \begin{tabularx}{0.7\textwidth}{X} 但是我告訴你們：要愛你們的仇敵，為那迫害你們的禱告。 \end{tabularx} \\ \\ \relax
5:45 & \begin{tabularx}{0.7\textwidth}{X} 這樣，你們就可以作天父的兒女了。因為他叫太陽照好人，也照壞人；降雨給義人，也給不義的人。 \end{tabularx} \\ \\ \relax
5:46 & \begin{tabularx}{0.7\textwidth}{X} 你們若只愛那愛你們的人，有甚麼賞賜呢？就是稅吏不也是這樣做嗎？ \end{tabularx} \\ \\ \relax
5:47 & \begin{tabularx}{0.7\textwidth}{X} 你們若只請你弟兄的安，有甚麼比別人強呢？就是外邦人不也是這樣做嗎？ \end{tabularx} \\ \\ \relax
5:48 & \begin{tabularx}{0.7\textwidth}{X} 所以，你們要完全，如同你們的天父是完全的。」 \end{tabularx} \\ \\
[1ex]
\hline
\hline
\end{longtable}
講題:如何可以勝過文士及法利賽人的義?

講員:高國雄牧師

經文:馬太福音5:13-48

日期:2023年2月5日

13「你們是世上的鹽.鹽若失了味,怎能叫它再鹹呢?以後無用,不過丟在外面,被人踐踏了.

14你們是世上的光.城造在山上是不能隱藏的.

15人點燈,不放在斗底下,是放在燈臺上,就照亮一家的人.

16你們的光也當這樣照在人前,叫他們看見你們的好行為,便將榮耀歸給你們在天上的父.」

17「莫想我來要廢掉律法和先知.我來不是要廢掉,乃是要成全.

18我實在告訴你們,就是到天地都廢去了,律法的一點一畫也不能廢去,都要成全.

19所以,無論何人廢掉這誡命中最小的一條,又教訓人這樣做,他在天國要稱為最小的.但無論何人遵行這誡命,又教訓人遵行,他在天國要稱為大的.

20我告訴你們,你們的義若不勝於文士和法利賽人的義,斷不能進天國.」

什麼是文士及法利賽人的義呢?

文士:是指那些有抄寫才能,負責抄寫聖經,律法條文,整理資料等.後來發展成為學者階層,負責抄寫猶太律法,並在聖殿和會堂中向百姓教導和講解.

法利賽:最初由一群反對猶太人與希臘同化,對抗外邦異俗的人,他們嚴守摩西律法,目的在保全國家精神之完整,所以此黨乃由熱烈的愛國精神及宗教虔誠而產生；及後因其要旨為牢守摩西律例淪為形式主義,以為只要遵守律法者,就可以得到上主賜福.

所以他們的使命就是要:守護律法,維護真理,並自以為是上帝的代言人.

在福音書內,耶穌對法利賽人的指責非常嚴厲,稱他們為「假冒為善,粉飾的墳墓(太23:27)」,和魔鬼的兒子(約8:44)

V20耶穌指出「義」是進入天國的必備條件,這不是一般的「義」,文士及法利賽人的「義」也不足夠!!

但這句話不是說我們要靠「(自己)義的行為」賺取進入天堂的門票,而是指「義」和「上帝的律法」在基督徒生命中的功能.

耶穌的義

V17莫想我來要廢掉律法和先知.我來不是要廢掉,乃是要成全.

耶穌基督完全順服上帝的律法,藉此為那些信靠祂的人獲得了義.(羅馬書中的因信稱義)

耶穌基督的義是種類上(而不是程度上)勝過了法利賽人的義.

我們在基督裡的人是從上帝的律法中被拯救了出來,因此我們不是靠遵行律法而得救的.

而是被上帝拯救,而去遵行上帝的律法,以遵行律法來作為敬拜拯救我們的上帝的方式/途徑.

所以相對於期望以遵行律法而獲取上帝的拯救/繼續維持在上帝的子民身份中(文士及法利賽人),更為超越,更加優勝.

因此耶穌更以實際例子去說明對我們「義」的要期望/要求:

V21-26

21「你們聽見有吩咐古人的話,說:『不可殺人』；又說:『凡殺人的難免受審判.』

22只是我告訴你們,凡向弟兄動怒的,難免受審判；凡罵弟兄是拉加的,難免公會的審判；凡罵弟兄是魔利的,難免地獄的火.

23所以,你在祭壇上獻禮物的時候,若想起弟兄向你懷怨,

24就把禮物留在壇前,先去同弟兄和好,然後來獻禮物.

25你同告你的對頭還在路上,就趕緊與他和息,恐怕他把你送給審判官,審判官交付衙役,你就下在監裡了.

26我實在告訴你,若有一文錢沒有還清,你斷不能從那裡出來.」

超越/勝過的「義」例子一:由「不可殺人」,延伸至不可對弟兄動怒/咒罵.

不單如此,更要提升層次,要求我們要主動與他人和好.

與人和好比一切重要,甚至比獻祭(宗教禮儀上的「義」)重要.

與人和好,同時也能解決地上公義要解決的問題,甚至可免審判(公義下需受罰的代價).

V27-30

27「你們聽見有話說:『不可姦淫.』

28只是我告訴你們,凡看見婦女就動淫念的,這人心裡已經與她犯姦淫了.

29若是你的右眼叫你跌倒,就剜出來丟掉,寧可失去百體中的一體,不叫全身丟在地獄裡.

30若是右手叫你跌倒,就砍下來丟掉,寧可失去百體中的一體,不叫全身下入地獄.」

超越/勝過的「義」例子二:由「不可姦淫」,延伸至不同對婦女動淫念.

要堅決拒絕罪惡,由萌芽時就開始.

真正的「義」,不單處理罪的行為,更要由源頭開始.

V31-32

31「又有話說:『人若休妻,就當給她休書.』

32只是我告訴你們,凡休妻的,若不是為淫亂的緣故,就是叫她作淫婦了；人若娶這被休的婦人,也是犯姦淫了.」

超越/勝過的「義」例子三:不要以為寫休書後休妻,就是已經盡了「義」,在律法上可能已經滿足,但耶穌對此不滿意,休妻要有根有據有理由.

做影響他者的決定時,不單要滿足基本要求(要做的),而是要讓原定規條的原意得滿足(應做的),否則「罪」就增加了.

V33-37

33「你們又聽見有吩咐古人的話,說:『不可背誓,所起的誓總要向主謹守.』

34只是我告訴你們,甚麼誓都不可起.不可指著天起誓,因為天是上帝的座位；

35不可指著地起誓,因為地是他的腳凳；也不可指著耶路撒冷起誓,因為耶路撒冷是大君的京城；

36又不可指著你的頭起誓,因為你不能使一根頭髮變黑變白了.

37你們的話,是,就說是；不是,就說不是；若再多說就是出於那惡者

.」

超越/勝過的「義」例子四:不可自以為義,不要以為自己可靠,「修飾」很可能就是為要「掩飾」.

V38-42

38「你們聽見有話說:『以眼還眼,以牙還牙.』

39只是我告訴你們,不要與惡人作對.有人打你的右臉,連左臉也轉過來由他打；

40有人想要告你,要拿你的裡衣,連外衣也由他拿去；

41有人強逼你走一里路,你就同他走二里；

42有求你的,就給他；有向你借貸的,不可推辭.」

超越/勝過的「義」例子五(更高的例子):真正的義,不是維護/悍衛自己權益,為證明自己有理.

「義」所要滿足的對象是上帝,讓上帝滿意才是義的存在價值.所以彰顯上帝的愛,這種類的義比單單計算公平/公允的義更有價值,更為超越.

V43-47

43「你們聽見有話說:『當愛你的鄰舍,恨你的仇敵.』

44只是我告訴你們,要愛你們的仇敵,為那逼迫你們的禱告.

45這樣就可以作你們天父的兒子；因為他叫日頭照好人,也照歹人；降雨給義人,也給不義的人.

46你們若單愛那愛你們的人,有甚麼賞賜呢?就是稅吏不也是這樣行嗎?

47你們若單請你弟兄的安,比人有甚麼長處呢?就是外邦人不也是這樣行嗎?

超越/勝過的「義」例子六(最高的例子):公義常被認為就是好善惡惡.

但耶穌教導我們,不是要以惡對惡,更要以愛勝惡,因為真正的公平公允只在上帝手中.

最大最高的「義」,就是耶穌基督道成肉身的愛.

惟有基督在我們還作罪人的時候為我們死,上帝的愛就在此向我們顯明了.(羅5:8)

V48所以,你們要完全,像你們的天父完全一樣.

真正能夠勝過文士及法利賽人的義,就只有最終極的追求,進入天國的條件,也同樣地向著這方向追求及邁進,方向正確比一切都重要…

記得一定要加上實際的行動!

V13-16

13「你們是世上的鹽.鹽若失了味,怎能叫它再鹹呢?以後無用,不過丟在外面,被人踐踏了.

14你們是世上的光.城造在山上是不能隱藏的.

15人點燈,不放在斗底下,是放在燈臺上,就照亮一家的人.

16你們的光也當這樣照在人前,叫他們看見你們的好行為,便將榮耀歸給你們在天上的父.」

努力成為世上的鹽,世上的光；一同見證主的豐富與美善,止於至善.



\newpage

\section{大人物與小薯仔}
\label{sec:636}
\textbf{鄧瑞強博士 - 大人物與小薯仔}
\newline
\newline
\hyperref[sec:635]{< < < PREV SERMON < < <}
~
\hyperref[sec:toc]{[返目錄]}
~
\hyperref[sec:637]{> > > NEXT SERMON > > >}


列王記下 5:1-15
\newline
\begin{longtable}{cl}
\hline
\hline
章節 & 經文 (和合本修訂版)\\
\hline
5:1 & \begin{tabularx}{0.7\textwidth}{X} 亞蘭王的元帥乃縵在他主人面前是一個偉大的人，得王的喜悅，因為耶和華曾藉他使亞蘭人得勝。他雖然是大能的勇士，卻染上了痲瘋。 \end{tabularx} \\ \\ \relax
5:2 & \begin{tabularx}{0.7\textwidth}{X} 亞蘭人成群出征的時候，從以色列地擄了一個小女孩，她就服事乃縵的妻子。 \end{tabularx} \\ \\ \relax
5:3 & \begin{tabularx}{0.7\textwidth}{X} 她對女主人說：「我希望主人去見撒瑪利亞的先知，他必能治好主人的痲瘋。」 \end{tabularx} \\ \\ \relax
5:4 & \begin{tabularx}{0.7\textwidth}{X} 乃縵去告訴他主人說，從以色列地來的女孩如此如此說。 \end{tabularx} \\ \\ \relax
5:5 & \begin{tabularx}{0.7\textwidth}{X} 亞蘭王說：「你可以去，我也會送信給以色列王。」於是乃縵手裡帶十他連得銀子、六千舍客勒金子和十套衣裳去了。 \end{tabularx} \\ \\ \relax
5:6 & \begin{tabularx}{0.7\textwidth}{X} 他帶著這信給以色列王，說：「現在你接到這信，看哪，我派臣僕乃縵到你這裡來，你要治好他的痲瘋。」 \end{tabularx} \\ \\ \relax
5:7 & \begin{tabularx}{0.7\textwidth}{X} 以色列王讀了信就撕裂衣服，說：「我豈是神，能使人死使人活呢？這人竟派人來，叫我治好一個人的痲瘋。你們要知道，看，這人是找機會來跟我吵架的。」 \end{tabularx} \\ \\ \relax
5:8 & \begin{tabularx}{0.7\textwidth}{X} 神人以利沙聽見以色列王撕裂衣服，就派人到王那裡，說：「你為甚麼撕裂衣服呢？讓那人到我這裡來，他會知道以色列中有先知。」 \end{tabularx} \\ \\ \relax
5:9 & \begin{tabularx}{0.7\textwidth}{X} 於是乃縵帶著車馬到了以利沙的家，站在門前。 \end{tabularx} \\ \\ \relax
5:10 & \begin{tabularx}{0.7\textwidth}{X} 以利沙派一個使者，對乃縵說：「去，在約旦河中沐浴七次，你的肉就必復原，你會得潔淨。」 \end{tabularx} \\ \\ \relax
5:11 & \begin{tabularx}{0.7\textwidth}{X} 乃縵卻發怒走了。他說：「看哪，我以為他必定會出來，到我這裡，站著求告耶和華－他神的名，在患處上搖手，治好這痲瘋。 \end{tabularx} \\ \\ \relax
5:12 & \begin{tabularx}{0.7\textwidth}{X} 大馬士革的亞瑪拿河和法珥法河豈不比以色列的一切水更好嗎？我難道不可以在那裡沐浴而得潔淨嗎？」於是他生氣，轉身走了。 \end{tabularx} \\ \\ \relax
5:13 & \begin{tabularx}{0.7\textwidth}{X} 他的僕人近前來，對他說：「我父啊，先知若吩咐你做一件大事，你豈不做嗎？何況是吩咐你去沐浴，得潔淨呢？」 \end{tabularx} \\ \\ \relax
5:14 & \begin{tabularx}{0.7\textwidth}{X} 於是乃縵下去，照著神人的話，在約旦河裡浸了七次。他的肉復原，好像小孩的肉，他就潔淨了。 \end{tabularx} \\ \\ \relax
5:15 & \begin{tabularx}{0.7\textwidth}{X} 乃縵帶著所有跟隨他的人，回到神人那裡，站在他面前，說：「看哪，我知道，除了以色列，全地沒有神。現在請你收下僕人的禮物。」 \end{tabularx} \\ \\
[1ex]
\hline
\hline
\end{longtable}
講題:大人物與小薯仔

經文:列王記下5:1-15

講員:鄧瑞強博士

日期:2023年2月19日

今日,我們來看《列王記下》的經文.

大家可能讀過《萬曆十五年》(1587, A Year of No Significance)這本書,作者為歷史學家黃仁宇,他在這書中表達一種「大歷史觀」,以宏觀的視野,討論事情如此發展的背後,有著龐大的制度之制約.今天,我講《列王記下》的一個小故事,我視之為一種「微歷史」,作者透過一「微敘事」,去舒發他對歷史前景的看法.

《列王記》寫成之時,猶大國已被巴比倫所滅,他們被擄去異鄉,他們的歷史看來驟然中斷.那些猶太神學家,在巴比倫的土地上,他們思考一些嚴肅的歷史神學問題.如:在強大的巴比倫帝國面前,猶太人渺小得可憐,但,若這歷史的渺小者仍獲得上天眷顧的話,則他們在強大者面前能扮演什麼角色?這些角色有何重要性?這些微小角色會改變歷史的軌跡嗎?這些微小角色能對人類的救贖作出貢獻嗎?

猶太的神學家透過歷史書寫,透露其信仰信念.這些在巴比倫的奴隸,因著其不死的信仰,無視自己的卑微,在心靈裡,進行著改變天地的沉思.

他們在《列王記下》,寫下一個小故事.我們不用質疑這故事的歷史性,但寫歷史的人,他在千萬故事中挑選這故事,書寫這故事,重點就在於他要透過這故事,去想像歷史的將來.

這故事記載在《列王記下5:1-15》,我們先看《列王記下》5:1-8:

5:1亞蘭王的元帥乃縵在他主人面前為尊為大,因耶和華曾藉他使亞蘭人得勝；他又是大能的勇士,只是長了大痲瘋.

5:2先前亞蘭人成群地出去,從以色列國擄了一個小女子,這女子就服事乃縵的妻.

5:3她對主母說:「巴不得我主人去見撒馬利亞的先知,必能治好他的大痲瘋.」

5:4乃縵進去,告訴他主人說,以色列國的女子如此如此說.

5:5亞蘭王說:「你可以去,我也達信於以色列王.」於是乃縵帶銀子十他連得,金子六千舍客勒,衣裳十套,就去了；

5:6且帶信給以色列王,信上說:「我打發臣僕乃縵去見你,你接到這信,就要治好他的大痲瘋.」

5:7以色列王看了信就撕裂衣服,說:「我豈是上帝,能使人死使人活呢?這人竟打發人來,叫我治好他的大痲瘋.你們看一看,這人何以尋隙攻擊我呢?」

5:8神人以利沙聽見以色列王撕裂衣服,就打發人去見王,說:「你為甚麼撕了衣服呢?可使那人到我這裡來,他就知道以色列中有先知了.」

故事一開始,便是一位大人物的出場.

5:1亞蘭王的元帥乃縵在他主人面前為尊為大,因耶和華曾藉他使亞蘭人得勝；他又是大能的勇士,只是長了大痲瘋.

這大人物叫乃縵,是「亞蘭王的元帥」,亞蘭即敘利亞,這人是軍隊的元帥.他在王面前「為尊」「為大」,耶和華使他「得勝」,他又是「大能的勇士」.所有的形容詞都與「偉大」有關.上天眷顧他,帝王看重他,他自己又有才能,真是人間極品.可惜,他「長了大痲瘋」,他面對的問題,也是大問題.

跟著,一個「小薯仔」出場.

5:2先前亞蘭人成群地出去,從以色列國擄了一個小女子,這女子就服事乃縵的妻.

這裡有一個小女子,原文真是有「小」這個形容詞的,她無名無姓,是被擄來在異鄉服事敵人的一個奴隸.

最偉大的大人物,最卑微的小薯仔,何等大的對比.

當《列王記》寫成之時,它的讀者的處境其實正是這樣.在強大的帝國面前,自己只是無名無姓的小薯仔.問題是:這小薯仔在歷史裡真的是完全無角色可言嗎?猶太的神學家在巴比倫,思考著自己在歷史中的角色.

《列王記》沒有對這「小女子」多費筆墨,她實在太微小了.現按人之常情作一般的推論.這小女子是以色列人,因打敗仗,被擄去服事敵人,她的心總有不甘.身為奴隸,她的苦境是由乃縵的軍隊造成的.乃縵若病死了,雖則未必能改變自己的命運,但起碼能吐一口怨氣.在這情況下,她靜看事情發展便是最好的了.但,在無人期望她發言的情況下,她竟然發言.她為什麼這麼蠢,出言獻計呢?若因她的建議,事情反而變得更壞的話,她的命運會更淒慘.這個小女子,因她的信仰,竟大膽發言.

5:3她對主母說:「巴不得我主人去見撒馬利亞的先知,必能治好他的大痲瘋.」

請弟兄姊妹留意,這是對上帝何等大的信仰,這是對敵人何等大的善意.她相信上帝的醫治大能,故她發言.她心中沒有敵人,故她向敵人作出建議.在戰爭的世界,這信仰發言其實是創造復和的契機.乃縵乃敵方主帥,若被己方的上帝醫治,人間的戰爭或許就會平息.大家要明白,古時的戰爭,開戰前要稟告神明,地上的戰爭,同時也是天上神明之間的戰爭.若敵方的神明醫好自己,則可視為敵方伸出了友誼之手,戰爭平息的機會大增.想想,若普京患了絕症,烏克蘭的主教為他按手祈禱而他痊癒,則和平的機會會否大增?

寫作《列王記》的神學家,他們被擄巴比倫,在強國下當奴隸,但他們相信,上帝是真實的,其恩典是宏大的,宏大得一切敵人也是其施恩的對象.在強權歷史裡,猶太人是微小得可憐的,只像一個無名無姓的小女子,但她的信仰宣告是偉大的,能創造人間的復和,醫治人間的敵意,顯揚上帝乃萬國之上的恩典主宰.這信仰宣告締造的復和,就是卑微的神的子民對世界歷史和人類救贖的貢獻.

信仰的善意,總被人性的卑污所污染.

5:4乃縵進去,告訴他主人說,以色列國的女子如此如此說.

5:5亞蘭王說:「你可以去,我也達信於以色列王.」於是乃縵帶銀子十他連得,金子六千舍客勒,衣裳十套,就去了；

5:6且帶信給以色列王,信上說:「我打發臣僕乃縵去見你,你接到這信,就要治好他的大痲瘋.」

5:7以色列王看了信就撕裂衣服,說:「我豈是上帝,能使人死使人活呢?這人竟打發人來,叫我治好他的大痲瘋.你們看一看,這人何以尋隙攻擊我呢?」

原本,可以是復和的好事,但亞蘭王卻搞局,將之變成發動戰爭的藉口.我想起周星馳那齣《唐伯虎點秋香》,那個「寧王」為了殺「華太師」,總會找一些不是理由的理由去「發爛渣」.亞蘭王有心「發爛渣」,是很難搞的.以色列王也很差勁,他的信仰比一個小女子還差.

先知出場了.

5:8神人以利沙聽見以色列王撕裂衣服,就打發人去見王,說:「你為甚麼撕了衣服呢?可使那人到我這裡來,他就知道以色列中有先知了.」

我們看看這故事如何發展下去,《列王記下》5:9-12:

5:9於是,乃縵帶著車馬到了以利沙的家,站在門前.

5:10以利沙打發一個使者,對乃縵說:「你去在約旦河中沐浴七回,你的肉就必復原,而得潔淨.」

5:11乃縵卻發怒走了,說:「我想他必定出來見我,站著求告耶和華－他上帝的名,在患處以上搖手,治好這大痲瘋.

5:12大馬士革的河亞罷拿和法珥法豈不比以色列的一切水更好嗎?我在那裡沐浴不得潔淨嗎?」於是氣忿忿地轉身去了.

乃縵大陣仗的來到以利沙那裡,怎料以利沙只派一個僕人出來招呼他,吩咐他去約旦河沐浴七次,如此,他的大痲瘋便得醫治.

乃縵的反應非常大.經文說:「乃縵卻發怒走了,說:我想他必定出來見我,站著求告耶和華他神的名,在患處以上搖手,治好這大痲瘋.」乃縵期望先知出來迎接他,在他患處上作下法,想不到竟遭到如此冷漠對待.乃縵繼續說:「大馬士革的河亞罷拿和法珥法豈不比以色列的一切水更好嗎?我在那裡沐浴不得潔淨嗎?」於是氣忿忿地轉身去了.

大人物乃縵有一種傲氣或霸氣,他習慣大場面,忍受不到被人當作小人物般對待.他忍受不到先知不出來見見他；他忍受不到先知竟不作一場大法事,不即席行一個大神蹟去醫治他；他忍受不到只需要在約旦河沐浴七次這等小兒科的小事.他覺得他的偉大被差辱,他覺得「丟臉」,覺得不受重視,於是忿然離去.

若事情真的這樣完了,則那小女孩的建議不但得不到好結果,反招來大禍.

我們看看這故事的結局,《列王記下》5:13-15:

5:13他的僕人進前來,對他說:「我父啊,先知若吩咐你做一件大事,你豈不做嗎?何況說你去沐浴而得潔淨呢?」

5:14於是乃縵下去,照著神人的話,在約旦河裡沐浴七回；他的肉復原,好像小孩子的肉,他就潔淨了.

5:15乃縵帶著一切跟隨他的人,回到神人那裡,站在他面前,說:「如今我知道,除了以色列之外,普天下沒有上帝.現在求你收點僕人的禮物.」

當乃縵忿然離場時,有另一個小薯仔介入.這也是一個無名的奴僕,他對乃縵說,若先知叫你做一些高難度的大事,你才得醫治,你粉身碎骨都會完成.多難的事,都難不到你.既然大事難不到你,怎能讓小事難到你?若視去約旦河裡沐浴為羞辱而不為,豈不被這小事難到你嗎?你這偉大的人,怎能被這小事難到?

這小薯仔的高見,竟改變了乃縵.於是,乃縵去約旦河,沐浴七回.結果是:「他的肉復原,好像小孩子的肉,他就潔淨了」.留意這裡的形容:他的肉,好像小孩子的肉.他的肉身,回到「小孩子的樣式」.在這裡,原文真是有「小」這個形容詞,與小女孩的「小」,同一字根.在這裡,我們看到信仰的神蹟,當一個不同凡響的大人物,願意聆聽小女孩的信仰宣告,願意走出自己的領域,願意跨越敵意的圍牆,願意離開自己宗教,文化,民族的界線而進入神國度的界線,願意長途跋涉來找先知,願意放下自己的成見而謙卑地遵照先知的吩咐,當他面對上帝的恩慈而一切都願意時,他痊癒了,回復到「小孩子的樣式」,回復到不受戰爭污染的最簡樸人性.

當人謙虛下來,體會到上天的恩典時,人只能讚嘆.

5:15乃縵帶著一切跟隨他的人,回到神人那裡,站在他面前,說:「如今我知道,除了以色列之外,普天下沒有上帝.現在求你收點僕人的禮物.」

乃縵如今承認,萬國之上,只有一位上帝.作奴隸的小女孩,因著她的信仰宣告,叫一個人間的大元帥,認信耶和華為獨一的神.現在,他自稱是「僕人」,是奴僕,並向神獻上感恩的禮物.

在強國底下,我們大部分人只是小人物,我們有的,是信仰.在巴比倫作奴隸的猶太神學家,藉著以上這「微敘事」,再次肯定他們的信仰.他們宣告:神不單在猶大地作王,神也在亞蘭及巴比倫作王,神是萬國之王.縱使在強權的歷史下,他們成了奴隸,看來失去了身分,但他們透過書寫《列王記》的故事,叫同胞不要忘記自己的真正身分,無論去到何處,無論在現世有何身分,他們的真正身分,是神的子民.神的子民的角色,是無論自己多微小,都要宣告信仰.這信仰能穿越一切的邊界,能醫治人類的頑疾,能改變人心.這是我們的歷史角色,這是我們對人類救贖的貢獻.終有一天,世人會宣認神的主權,在那刻,人類會復和,戰爭會止息,天國會臨在人間.這是我們這些小人物在信仰裡的期望.

親愛的弟兄姊妹,不要小看自己「小薯仔」的角色.

願榮耀歸上主,願人間享和平.



\newpage

\section{走不易的道路}
\label{sec:637}
\textbf{陳智衡牧師 - 走不易的道路}
\newline
\newline
\hyperref[sec:636]{< < < PREV SERMON < < <}
~
\hyperref[sec:toc]{[返目錄]}
~
\hyperref[sec:638]{> > > NEXT SERMON > > >}


路加福音 9:51-62
\newline
\begin{longtable}{cl}
\hline
\hline
章節 & 經文 (和合本修訂版)\\
\hline
9:51 & \begin{tabularx}{0.7\textwidth}{X} 耶穌被接上升的日子將到，他決定面向耶路撒冷走去。 \end{tabularx} \\ \\ \relax
9:52 & \begin{tabularx}{0.7\textwidth}{X} 他打發使者在他前頭走；他們進了撒瑪利亞的一個村莊，要為他作準備。 \end{tabularx} \\ \\ \relax
9:53 & \begin{tabularx}{0.7\textwidth}{X} 那裡的人不接待他，因為他面向著耶路撒冷去。 \end{tabularx} \\ \\ \relax
9:54 & \begin{tabularx}{0.7\textwidth}{X} 他的門徒雅各和約翰看見了，就說：「主啊！你要我們吩咐火從天上降下來，燒滅他們嗎？」 \end{tabularx} \\ \\ \relax
9:55 & \begin{tabularx}{0.7\textwidth}{X} 耶穌轉身責備兩個門徒。 \end{tabularx} \\ \\ \relax
9:56 & \begin{tabularx}{0.7\textwidth}{X} 於是他們就往別的村莊去了。 \end{tabularx} \\ \\ \relax
9:57 & \begin{tabularx}{0.7\textwidth}{X} 他們在路上走的時候，有一個人對耶穌說：「你無論往哪裡去，我都要跟從你。」 \end{tabularx} \\ \\ \relax
9:58 & \begin{tabularx}{0.7\textwidth}{X} 耶穌對他說：「狐狸有洞，天空的飛鳥有窩，人子卻沒有枕頭的地方。」 \end{tabularx} \\ \\ \relax
9:59 & \begin{tabularx}{0.7\textwidth}{X} 他又對另一個人說：「來跟從我！」那人說：「主啊，容許我先回去埋葬我的父親。」 \end{tabularx} \\ \\ \relax
9:60 & \begin{tabularx}{0.7\textwidth}{X} 耶穌對他說：「讓死人埋葬他們的死人，你只管去傳講神的國。」 \end{tabularx} \\ \\ \relax
9:61 & \begin{tabularx}{0.7\textwidth}{X} 又有一人說：「主啊，我要跟從你，但容許我先去辭別我家裡的人。」 \end{tabularx} \\ \\ \relax
9:62 & \begin{tabularx}{0.7\textwidth}{X} 耶穌對他說：「手扶著犁向後看的人，不配進神的國。」 \end{tabularx} \\ \\
[1ex]
\hline
\hline
\end{longtable}
講題:走不易的道路

經文:路加福音9:51-62

講員:陳智衡牧師

日期:2023年2月26日

路加福音九51節開始,至路十九44節,是整本路加福音的中央部分.上一個大段是路四14開始,講述耶穌呼召門徒,行神蹟奇事,宣講天國真理.而在下一個大段,是由路十九45開始,那是講述耶穌基督在耶路撒冷的受難和復活.而中間這大段,路加福音九51節開始這,那是講述在位的猶太人領袖,他們不斷地抗拒和挑戰耶穌,而耶穌做的,就是不斷呼召上帝的子民悔改歸主.而今日這段經文,是這中央部分耶穌進入耶路撒冷聖殿作為開始.

「耶穌被接上升的日子將到,他決定面向耶路撒冷走去」(路九51).這次耶穌往耶路撒冷之旅,是耶穌準備受死的救贖旅程.耶穌沒有拒絕撒瑪利亞人,但撒瑪利亞人卻是拒絕接待耶穌.為什麼呢?原因是,當時撒瑪利亞人本身對猶太人就有排斥,凡是往耶路撒冷的猶太人朝聖者,撒瑪利亞人都會這樣態度對待他們.

門徒有見及此,就跟耶穌說:「主啊!你要我們吩咐火從天上降下來,燒滅他們嗎?」(路九55)門徒這樣的反應,亦是被兩族的多年仇恨所影響.

「你無論往哪裡去,我都要跟從你.」(路九57)相信耶穌聽到這句說話會很疑惑,因為那人根本不知道耶穌去那裡.耶穌上耶路撒冷,是以死進行對世人的救贖.

耶穌叫第二個想跟隨他,而那人的回應是「主呀,容許的先回去埋葬我的父親」(路九60).而耶穌的回應是「讓死人埋葬他們的死人,你只管去傳講神的國」.(路九61).

第三個人跟耶穌說「容許我先去辭別我家裡的人」(路九61).那是引用舊約列王記上十九20節,以利亞呼召以利沙跟隨他的時候,以利沙卻說「請你讓我先與父母吻別,然後我就跟隨你.」(王上十九20).而當時以利亞是容許以利沙可以這樣做,與父母吻別才回到以利亞身邊學習和服侍神.

耶穌說:「手扶著犁向後看的人,不配進神的國」,「不配進」原文意思是「不合用」,意思是「對上帝的國是不適用」.

耶穌認為的重點,並非是人一般所理解的,所以人才會錯覺以為耶穌叫價大高,其實並不是.

某程度上,因為我們與耶穌的視角不同,我們低估了耶穌的要求.

基督信仰,其實是要付代價的.

作為耶穌門徒,亦是有著基督信仰的人.我們會因為基督信仰而甘願放棄做某些事情,又或因為基督信仰而甘願做一些自己未必喜歡的事情.嚴格來說,基督徒的行徑是不合乎常理的,早期教會信徒因信仰被族群斷絕關係,被社會大眾唾棄.



\newpage

\section{無水之道}
\label{sec:638}
\textbf{李均熊牧師 - 無水之道}
\newline
\newline
\hyperref[sec:637]{< < < PREV SERMON < < <}
~
\hyperref[sec:toc]{[返目錄]}
~
\hyperref[sec:639]{> > > NEXT SERMON > > >}


出埃及記 17:1-7
\newline
\begin{longtable}{cl}
\hline
\hline
章節 & 經文 (和合本修訂版)\\
\hline
17:1 & \begin{tabularx}{0.7\textwidth}{X} 以色列全會眾遵照耶和華的吩咐，從汛的曠野一段一段地往前行。他們在利非訂安營，但百姓沒有水喝。 \end{tabularx} \\ \\ \relax
17:2 & \begin{tabularx}{0.7\textwidth}{X} 百姓就與摩西爭鬧，說：「給我們水喝吧！」摩西對他們說：「你們為甚麼與我爭鬧呢？你們為甚麼試探耶和華呢？」 \end{tabularx} \\ \\ \relax
17:3 & \begin{tabularx}{0.7\textwidth}{X} 百姓在那裡口渴要喝水，就向摩西發怨言，說：「你為甚麼把我們從埃及領出來，使我們和我們的兒女，以及牲畜都渴死呢？」 \end{tabularx} \\ \\ \relax
17:4 & \begin{tabularx}{0.7\textwidth}{X} 摩西就呼求耶和華說：「我要怎樣對待這百姓呢？他們差一點就要拿石頭打死我了。」 \end{tabularx} \\ \\ \relax
17:5 & \begin{tabularx}{0.7\textwidth}{X} 耶和華對摩西說：「你帶著以色列的幾個長老，走在百姓前面，手裡拿著你先前擊打尼羅河的杖，去吧！ \end{tabularx} \\ \\ \relax
17:6 & \begin{tabularx}{0.7\textwidth}{X} 看哪，我要在何烈的磐石那裡，站在你面前。你要擊打磐石，水就會從磐石流出來，給百姓喝。」摩西就在以色列的長老眼前這樣做了。 \end{tabularx} \\ \\ \relax
17:7 & \begin{tabularx}{0.7\textwidth}{X} 他給那地方起名叫瑪撒，又叫米利巴，因為以色列人在那裡爭鬧，並且試探耶和華，說：「耶和華是否在我們中間呢？」 \end{tabularx} \\ \\
[1ex]
\hline
\hline
\end{longtable}
講題:無水之道

講員:李均熊牧師

經文:出埃及記17:1-7

日期:2023年3月12日

起:求水那有錯?

17:1以色列全會眾遵照耶和華的吩咐,從汛的曠野一段一段地往前行.他們在利非訂安營,但百姓沒有水喝.

以色列人離開為奴的埃及,進入了自由的曠野.他們是按著上主的吩咐,一段段的在曠野往前走的.但曠野雖然帶來自由,但亦充滿危機和不確定.當以色列人行到利非訂,百姓沒有水喝.

百姓渴了.在曠野之地會眾沒有水喝,向上主發出要求,何錯之有?

承:米利巴與瑪撒

17:2百姓就與摩西爭鬧,說:「給我們水喝吧!」摩西對他們說:「你們為甚麼與我爭鬧呢?你們為甚麼試探耶和華呢?」

17:3百姓在那裡口渴要喝水,就向摩西發怨言,說:「你為甚麼把我們從埃及領出來,使我們和我們的兒女,以及牲畜都渴死呢?」

17:4摩西就呼求耶和華說:「我要怎樣對待這百姓呢?他們差一點就要拿石頭打死我了.」

沒有水而向上主求水喝,自然不過.但經文明確的給我們看到,百姓的要求,卻變成向摩西——他們的領袖——爭鬧的理由.而摩西更進一步指出,百姓與摩西爭鬧,骨子裡其實是在試探上主.

爭鬧和試探,成為為事件發生地方命名的依據:米利巴,就是爭鬧；瑪撒,就是試探.

爭鬧(argue),用今日的語言來說,就是埋怨(grumble),控訴(complaint),責罵(chide),「嗌交」(quarrel),甚至初而口角,繼而動武(from verbal quarrl to physical fight).不是嗎?「他們差一點就要拿石頭打死我了」——摩西不是看見他們每個都拿起石頭準備向他丟過去嗎?

沒有水而求水,之所以變成爭鬧,並不是危機如何艱巨,困難如何凶險.以色列人不是看見過上主在出埃及施行的神蹟奇事嗎?過紅海後以色列人不是高唱上主是戰士嗎?在書珥的曠野上主不是令苦水變甜嗎?在以琳他們不是享受過十二股水泉,七十棵棕樹的美景嗎?在汛的曠野他們不是享用嗎哪和鵪鶉直到如今嗎?為甚麼這刻發現無水會演變成埋怨爭鬧?

百姓對摩西不是只要求「給我們水喝」,而是質疑摩西「為甚麼把我們從埃及領出來,使我們和我們的兒女,以及牲畜渴死呢?」他們不領情不特止,更質疑摩西的動機,認為他們領袖的出發點不是要解放他們免受奴役,而是要引他們到無水之地將他們害死!「給你害死啦!」

你不要奇怪以色列百姓為何會如此.出埃及記只是將人性的一面向我們展露.人就是會忘恩負義,無論經歷幾多恩典,遇到困難問題,也會質疑身邊的一切,甚至演變成埋怨和投訴.不是嗎?

因為這樣,這裡就叫做米利巴.

摩西進一步使我們意識到,這裡百姓並不是針對他們的領袖,而是領袖背後代表的上主:「你們為甚麼試探耶和華?」.百姓其實是在「試探」上主.甚麼是「試探」?這裡的「試探」不是test,並不是出一道難題看你能否解決,而是trick,製造一個陷阱,看一下你會不會失足掉下去.就如今日我們常常遇到的網絡釣魚(phishing),詐騙電話(scam call)等等.又如魔鬼試探耶穌,叫祂自找麻煩,從聖殿頂跳下去,「看一下神會不會幫助你!」耶穌駁斥魔鬼引用的申命記6:16「不可試探主—你的神」,下半句就是「像你們在瑪撒那樣試探他.」

以色列人怎樣「試探」上主呢,經課結束的說話可能就是線索:百姓「試探耶和華,說:耶和華是否在我們中間呢?」百姓並不是只是口渴,所以要求水喝,而是「啍,看看山你是否真的做得到!」百姓是否遇見無水之境,就想試一下引上主「現形」?甚至可能當上主「現形」的時候,就找個百寶袋將祂封印起來,到有甚麼需要,就隨時可以找祂出來幫手!

情形就好似你遇到一個精靈,想用牠的法力去幫你做事,於是乎就「試探」(trick)這個精靈:「你真的可以鑽進這盞那麼小的神燈嗎?我不信!」,於是精靈為了證明牠可以的,就化作一陣青煙,鑽進那個神燈裡,然候你就將牠封印起來,到要用牠的時候,就擦三下神燈,牠就要出來滿足你的願望一樣.

上主要隨時滿足我的需要,不然跟祂走曠野路有什麼好處?不如回埃及地去,起碼不會餓死,渴死.

因為這樣,所以這裡就叫做瑪撒.

轉:上主站那裡?

17:5耶和華對摩西說:「你帶著以色列的幾個長老,走在百姓前面,手裡拿著你先前擊打尼羅河的杖,去吧!

17:6看哪,我要在何烈的磐石那裡,站在你面前.看我啊,站著在你面前,在何烈的磐石之上.你要擊打磐石,水就會從磐石流出來,給百姓喝.」摩西就在以色列的長老眼前這樣做了.

但說到底,百姓的要求是合理的,於是上主就吩咐摩西,帶著以色列的幾個長老,走在百姓前面.大家不要忘記,每個百姓手上拿著石頭.這個舉動,是要鍛煉摩西,等他壯膽?還是要提醒一班凶神惡煞的以色列人,摩西只是其中一個代表上主的領導,其實還有一大班長老都是代表上主,這是集體領導,針對摩西乃愚蠢之舉?

關於領導問題,今日暫且按下不表.

但上主同時叫摩西手裡拿著先前擊打尼羅河的杖,更可能就是提醒這一班忘恩負義的以色列人,上主一直沒有離開以色列人,透過可見的杖顯出祂的神蹟奇事.昔日擊打尼羅河的杖,會否就今日的十字架,又或聖餐的酒與餅?

上主跟著的吩咐其實我們更要留心,我嘗試按原文再翻譯一下:

看我啊,站著在你面前那處,在何烈的磐石之上.

你要擊打磐石,從那裡水就會流出來,百姓就可以喝.

若我們照字面意思去理解,摩西會見到上主就在磐石的上面,那麼摩西擊打磐石,便是朝著上主打上去了?

無論這個解釋是否合理,這個擊打磐石而出水的意像,可能是預苦期選讀這一段舊約的原由.這塊磐石,早期教會認為就是基督的預表,在結合猶太人認為這塊磐石在此之後就在曠野一路上跟著以色列人走的早期傳統,保羅在林前10:1-4就這樣寫道:「弟兄們,我不願意你們不知道,我們的祖宗從前都在雲下,都從海中經過,都在雲裡,海裡受洗歸了摩西,並且都吃了一樣的靈糧,也都喝了一樣的靈水,所喝的是出於跟隨著他們的靈磐石；那磐石就是基督.但他們中間多半是神不喜歡的人,所以倒斃在曠野裡了.」

那靈磐石是跟隨著他們的靈磐石,而那磐石就是基督.磐石被擊打而出水,成為一個異象,讓我們省察我們今日井旁婦人期望能得到的活水,是透過基督的受苦而成存的.

合:不要留在瑪撒米利巴

17:7他給那地方起名叫瑪撒,又叫米利巴,因為以色列人在那裡爭鬧,並且試探耶和華,說:「耶和華是否在我們中間呢?」

可以見到,在預苦期選讀這段舊約經課,其正面目的是回響井旁婦人所得到的活水,乃透過那靈磐石所預表的基督被擊打而能夠賜給我們.但整段舊約經課的焦點,正正是17:7要我們記得的,就是那地方的名字:名叫瑪撒,又叫米利巴.詩篇的選讀也不厭其煩地提醒我們:「手你們不可硬著心,像當日在米利巴,就是在曠野的瑪撒.那時,你們的祖宗試我探我,並且觀看我的作為.四十年之久,我厭煩那世代,說:這是心裡迷糊的百姓,竟不曉得我的作為!所以,我在怒中起誓,說:他們斷不可進入我的安息.」之前提到的申6:16也說「不可試探主—你的神,像你們在瑪撒那樣試探他.」我也提過保羅在林前10引用以色列人在曠野的經歷,告訴我們:這些事都是我們的鑒戒,使我們不要貪戀惡事,像他們貪戀過的一樣.

為何要將事發地點改名做米利巴和瑪撒,每次提起都不厭其煩叫我們不要重蹈覆轍,好像當日在曠野的百姓般呢?

在基督徒的日常生活裡,遇到困難是必然的.人生本來就充滿危機與困難,基督徒並不能幸免.但我們透過基督成為上主的兒女,能因耶穌的救恩而得享活水天糧,並不代表我們就會常存感恩,或是面對困厄時能真心倚靠神.米利巴提醒我們,以色列人出埃及過紅海,得解放得自由,並未換來他們對上主的感恩,一遇危難就會去與帶領他們,施恩給他們的主去爭鬧.而瑪撒就提醒我們,米利巴背後可能就是百姓想透過這般撒野,去誘使trick

上主去滿足他們的需求.問題不在我們的渴求是否正當,又或上主能否又或會否滿足我們的渴求；問題是我們想透過我們的方法去促(迫?)使上主為我們服務,而不是我們跟隨上主引領,使人尊主為聖,讓上主旨意在地若天.

這是人性的軟弱,就算我們如何有信心,在面對危難困厄時也有機會落入米利巴和瑪撒的錯誤中而不自知.要常常憶起米利巴和瑪撒,提醒自己不可對上主無理要求.是我們遵主指使行曠野路,而不是在曠野路搵上主著數(要上主為我們服務).

預苦期是紀念耶穌受苦與受死的日子.摩西擊打磐石而出水,能成為提醒我們上主的受苦使我們得活水的意像.但發生這事的地方名叫米利巴,即是瑪撒,就要提醒我們不要以為上主施與我們生命的活水,我們就可以對上主強作要求,要祂為我們的需要服務.

不過,今日書信經課也提醒我們,我們還作罪人的時候主耶穌就替我們死,而顯明神對我們的愛.所以,雖然以色列人忘恩負義,又要試探祂,只想上主為他們服務,上主仍然透過磐石被擊打而出水,去滿足他們在無水之時的渴求.因此讓我們不要忘記,今天我們能如井旁婦人得活水而不渴,領受上主因受苦而賜下給我們生命,並不是因我們的義,而是上主全然的接納和恩典.



\newpage

\section{家譜會說話}
\label{sec:639}
\textbf{鄧瑞強博士 - 家譜會說話}
\newline
\newline
\hyperref[sec:638]{< < < PREV SERMON < < <}
~
\hyperref[sec:toc]{[返目錄]}
~
\hyperref[sec:640]{> > > NEXT SERMON > > >}


講題:家譜會說話

經文:歷代志上1:1,28,3:1-10,17-19,4:1,6:1,31-33,49-50,8:1-2,9:17-19,39-40

講員:鄧瑞強博士

日期:2023年3月19日

各位弟兄姊妹,早安.

今日,我們來看《歷代志上》的經文.在希伯來聖經裡,《歷代志上下》原是一卷.

未入正題之前,我要交代一下上兩個月講的《列王記》與今日開始講的《歷代志》的神學差異.《列王記》結束的經文,是講到猶大人被擄巴比倫.寫作的時候,猶大人仍在巴比倫.《列王記》的寫作,主要是失敗後的賽後檢討,追尋亡國的神學理由,所以我早前以所羅門王為例,講到所羅門王因失去王者的風範,已預示了猶大國的消亡.今日要講的《歷代志》,它結束的經文,是講到波斯王塞魯士的興起,他宣告猶大人可回歸故土,重建家園.故此,《歷代志》的寫作,重點就不在失敗後的賽後檢討,而在如何重整隊伍,準備下一輪的賽事.若檢討失敗,則重點在指出錯誤,嚴責失職.若準備下一輪的賽事,則重點不在嚴責,而在鼓勵,鼓勵各人重整隊形,鼓勵各人守好崗位.由於《歷代志》的重點在鼓勵人,讓人重整隊形,故此讀《歷代志》時,你會發覺大衛好像從沒有淫人妻子,殺人丈夫似的,而所羅門的一生,都是那麼光明神聖.從《歷代志》的角度看,猶大人亡國,「輸波」已「輸到不能再輸」,現在有機會回歸重建,有機會再戰江湖,就不好再埋怨那些主將(大衛,所羅門)的失職了,不如倒過來,強調主將的強項,讓這些主將支撐起整支球隊,看看能否打好以後的賽事.

要重整球隊,就要點將,排陣,抓緊前進的重心.《歷代志》的開場白,就是一連九章的家譜經文.書寫家譜,就是點將,排陣,講明日後發展的重心.這種家譜神學,上承創世記的家譜神學,下開馬太福音和路加福音的家譜神學,但《歷代志》令人矚目的地方,是它的家譜經文篇幅特別長.

講起家譜,想起有一次,白耀燦老師帶我們的團友到龍躍頭的松嶺鄧公祠作古蹟導賞,祠堂裡,有一牌位,上有龍頭,原來是宋高宗之女因戰亂南下,輾轉嫁了鄧氏族人.家譜是會講故事的,它訴說著王室與香港的一段歷史姻緣.

美國的NBA有所謂名人堂,誰能入這名人堂,當然需要經過嚴格的甄選程序.這人有什麼卓越的成就,他對美國籃球運動有什麼非凡的貢獻,他的名字被永遠地紀念會對後來者帶來什麼積極的影響?在高手雲集的NBA,要回答這些問題,不是件易事.誰被紀念,誰被遺忘,誰被留下名字去啟發後來者的夢想,在歷史的建構裡,這從來都是件複雜的事情.

《歷代志上》的家譜,不是一自然家譜,而是一神學家譜,解決嚴肅的神學問題.《歷代志》面對的神學問題是:被擄去巴比倫的猶大人回歸故土後,如何重建社群?這社群要具備「神的子民」的氣質,免得重蹈覆轍,再次滅亡.

一大批難民一樣的猶大人回歸故土,要組織他們,家譜的書寫是首要大事.學者認為,家譜有幾項重要功能.一,法律的功能:處理誰可分配土地的問題,在香港,這就是誰有「丁權」的問題.對散沙一樣的難民來說,這是最基本的問題.二,社會的功能:處理社會上的階級和職業等問題,哪種行業最尊貴?由誰去承擔?三,政治的功能:處理誰可合法管治的問題.四,心理的功能:處理如何建構一個身分認同的問題.這群人將是一群什麼樣的人,以什麼身分去凝聚他們.以上種種問題很重要,影響著這群人將來的命運,塑造著他們的歷史角色,也判斷出他們能否成為上帝救贖世界的器皿.《歷代志上》建構一神學家譜,一併處理以上問題.

由於《歷代志上》的家譜很長,我們只能抽一些段落來看.

1:1亞當生塞特……

我們知道,亞當生了該隱和亞伯.因該隱殺了亞伯,該隱就失去進入這家譜的資格.塞特的後裔,卻開出亞伯拉罕的譜系,這是神救贖世人所揀選的譜系.《歷代志上》1:1的經文已顯示,這不是自然的譜系,而是神救贖計劃的譜系.

1:28亞伯拉罕的兒子是以撒,以實瑪利.

我們知道,以實瑪利才是長子,但他是按亞伯拉罕的意願而不是按神的應許而生的,在家譜裡,他退而為次.但將他列入家譜,明顯,是《歷代志上》有一天下一家的精神.它嘗試將非以色列人納入家譜中.或許,猶大人經過被擄後,發展出這種天下一家的想法.他們去到巴比倫後,他們發現,神不單是猶大地的神,也是巴比倫地的神.這位神,在巴比倫帝國之上掌權,主宰著巴比倫的命運.有了這一認識,他們擴大了他們的神觀,神不只是以色列人的神,也是萬民之上的神.這種超越民族的一神觀,讓以色列人更明白他們作為祝福萬國萬民的角色.

3:1大衛在希伯崙所生的兒子記在下面:長子暗嫩是耶斯列人亞希暖生的.次子但以利是迦密人亞比該生的.……

3:4這六人都是大衛在希伯崙生的.大衛在希伯崙作王七年零六個月,在耶路撒冷作王三十三年.

3:5大衛在耶路撒冷所生的兒子是示米亞,朔罷,拿單,所羅門.這四人是亞米利的女兒拔書亞生的.……

3:9這都是大衛的兒子,還有他們的妹子她瑪,妃嬪的兒子不在其內.

3:10所羅門的兒子是羅波安……〔按:在這裡,所羅門成了最重要的大衛之子〕

這是大衛在耶路撒冷作王前和後的兩列後裔名單.值得留意的是3章5節.這裡提到的女人名字是「拔書亞」,她的名字在《撒母耳記下》和《列王記上》都是「拔示巴」.改變了寫法,是否想隱藏她與大衛之間的醜事?這女人的名字,在《歷代志》只出現一次,且是作為一個母親的名字出現,明顯,是想淡化她的事蹟.再者,在別處地方(撒下5:14),我們知道大衛有示米亞,朔罷,拿單這三個兒子,但從來無提及「拔示巴」是他們的母親.按《撒母耳記下》的記載,大衛與「拔示巴」生的兒子,除了夭折的那個孩子之外,便是所羅門.在這裡,所羅門忽然排到第四,是令人驚奇的.但這種寫法,就將所羅門作為大衛和「拔示巴」姦情的頭號確鑿證據隱藏,這有利大衛和所羅門的名聲.這家譜,重點不在嚴責罪惡,而是為民族日後的存亡找些關鍵人物,以他們的偉大事蹟,凝聚散沙一盤的人.也為日後的社群秩序,確立合法的管治人物.

3:17耶哥尼雅被擄.他的兒子是撒拉鐵……

3:19毗大雅的兒子是所羅巴伯……

這是猶大亡國後,在被擄之人中,大衛後裔的名單.請留意,亡國後,大衛的後裔作為管治者的正統血脈,並沒有散失.不但沒有散失,反被強化.在群龍無首的情況下,有人可合法管治,是關鍵的.3:19提到所羅巴伯,在日後我要講的《以斯拉記》裡,這人就是帶領回歸的猶大人重建社群的領袖.明顯,《歷代志》強調大衛的家系,他們是王室正統,他們可合法地管治從流散中回來的人.這當然會強化後來的彌賽亞信仰,從神而來的救主,一定是大衛的子孫.

第4章開始,便以十二支派來記載家譜.

4:1猶大的兒子是法勒斯,希斯崙,迦米,戶珥,朔巴.

在這家譜裡,十二支派以猶大支派排頭位,這當然也不是自然次序.按出生次序,十二支派應以「呂便」為首.但《歷代志》的家譜,重點不在出生次序,而在重建一個新秩序.合法的管治者是大衛的家系,而大衛出自猶大支派.猶大支派也是南國猶大的主要成員.可以看到,這家譜的設計,指向猶大國.猶大國作為一地理位置,在回歸者的新秩序中,這將是世界的中心.

我們追問:猶大作為世界的中心,意義何在?

回答這問題,我們要繼續看這家譜的設計,我們來到這家譜的中心部分.

6:1利未的兒子是革順,哥轄,米拉利.

家譜的中心是利未支派.若這家譜的功能是用來分地,則不應列出利未支派,因為利未支派是服務聖殿的,心在聖殿,家在聖殿,無地分給他們.但這家譜是用以重整這群曾流亡的人日後的人生的,則將利未放在家譜的中央,用意便明顯了,就是要他們像利未人,心在聖殿,家在聖殿,以聖殿為人生的中心.一個民族心中若沒有聖殿,便什麼也沒有了.

6:31約櫃安設之後,大衛派人在耶和華殿中管理歌唱的事.

6:32他們就在會幕前當歌唱的差,及至所羅門在耶路撒冷建造了耶和華的殿,他們便按著班次供職.

6:33供職的人和他們的子孫記在下面:哥轄的子孫中有歌唱的希幔.希幔是約珥的兒子；約珥是撒母耳的兒子……

重中之重,就是聖殿中的約櫃,裡面有神與人立約的法板.流亡回歸的猶大人,要重建身分,要凝聚成一個民族,其身分認同的核心,就是神人之約.他們永遠要記得,他們的真正身分是「神的子民」.若失落這身分,就沒有所謂前景,沒有所謂民族的天命,歷史也將永遠遺忘他們.

以上這段經文,又提到大衛和所羅門.來到這裡,我們終於知道,為何大衛和所羅門的劣跡被洗掉,他們的家系被抬舉,理由是:大衛設計好聖殿的崇拜,而所羅門建造了聖殿.猶大支派,大衛,所羅門,猶大國之所以被重視,是因為他們建構,保護,支撐著聖殿的崇拜,而這是神的子民一切希望之所依.

在《歷代志上》第6章有長長的名單,是服務聖殿者的名單.好長好長的名單,記著很多小人物的名字,他們列在這家譜的中央位置,成了後世景仰的人物.有一天,在這聖殿裡服務的人,包括:崇拜主席,領詩,司琴,詩班,傳譯同工,接待同工,司事同工,音響同工等等,都會寫入這長長的名單內,成了永被記念的人物.

6:49亞倫和他的子孫在燔祭壇和香壇上獻祭燒香,又在至聖所辦理一切的事,為以色列人贖罪,是照上帝僕人摩西所吩咐的.

6:50亞倫的兒子是以利亞撒……

這裡提及聖殿裡的另一事奉者:祭司.他們作為神與人之間的中介,溝通著天地.最重要的,他們向人間展示上帝無限的仁慈.憑這上天的仁慈,絕望的人才有力氣走向明天.也因這上天的仁慈,神的子民對世界有宣講的內容.

十二支派的名單,我們來到壓尾部分.

8:1便雅憫的長子比拉,次子亞實別,三子亞哈拉,

8:2四子挪哈,五子拉法.……

壓尾的支派是便雅憫支派.以它來壓尾,是因為南國猶大的組成部分,正是猶大支派和便雅憫支派.頭和尾,都指向猶大國.猶大國的中心,便是聖殿.聖殿就是神的子民存在的基礎.鄧麗君有首歌,叫做《忘記他》,歌詞是:「忘記他,等於忘掉了一切,等於將方和向拋掉,遺失了自己.」神的子民若忘記聖殿,就是這樣,等於忘掉了一切,等於將方和向拋掉,遺失了自己.

9:17守門的是沙龍,亞谷,達們,亞希幔,和他們的弟兄；沙龍為長.

9:18從前這些人看守朝東的王門,如今是利未營中守門的.

9:19可拉的曾孫,以比雅撒的孫子,可利的兒子沙龍,和他的族弟兄可拉人都管理使用之工,並守會幕的門.他們的祖宗曾管理耶和華的營盤,又把守營門.

我們猜想不到,連看守聖殿門口的守衛,也有一系列名單,列在這家譜內.我不期然的想起,詩篇84:10的經文:「在你的院宇住一日,勝似在別處住千日；寧可在我上帝殿中看門,不願住在惡人的帳棚裡.」

我們最後要看的一小段家譜,是提到便雅憫支派不得不提的一個顯赫家族:掃羅的家譜.

9:39尼珥生基士；基士生掃羅；掃羅生約拿單,麥基舒亞,亞比拿達,伊施‧巴力.

9:40約拿單的兒子是米力‧巴力；米力‧巴力生米迦.

掃羅是以色列人的第一個王.這裡有他的家譜,說明這王族的影響力從未消失,他們是管治權的潛在競爭者.但我們留意到,在這家譜裡,不斷出現帶有「巴力」這外邦神明的名字.家譜這樣記載,讓神的子民明白,掃羅家族已離開了聖殿的純正敬拜,失去了統治的合法性了.

結語:面對著一盤散沙的群眾,寫作《歷代志》的神學家,藉這家譜書寫,首要解決誰去管治的問題.他們賦予大衛家族以管治的合法性.接著,他們處理凝聚群眾的身分認同問題.他們以聖殿為精神家園,以神人之約為身分的核心,從而讓失去身分的人重拾神子民的身分.大衛家族之所以有管治的合法性,正是因為他們維護了聖殿的敬拜.至於社會上的階級和職業等問題,就教導群眾以服務聖殿的精神來考量他們的工作選擇.費盡苦心,《歷代志》的神學家在廢墟中,力圖重建一個新的世界.

若你身處一荒涼世界,《歷代志》在什麼意義下幫助你重建一信仰人生呢?

願榮耀歸上主,願人間享和平.



\newpage

\section{人在此身中}
\label{sec:640}
\textbf{林豪恩先生 - 人在此身中}
\newline
\newline
\hyperref[sec:639]{< < < PREV SERMON < < <}
~
\hyperref[sec:toc]{[返目錄]}
~
\hyperref[sec:641]{> > > NEXT SERMON > > >}


羅馬書 8:6-11
\newline
\begin{longtable}{cl}
\hline
\hline
章節 & 經文 (和合本修訂版)\\
\hline
8:6 & \begin{tabularx}{0.7\textwidth}{X} 體貼肉體就是死；體貼聖靈就是生命和平安。 \end{tabularx} \\ \\ \relax
8:7 & \begin{tabularx}{0.7\textwidth}{X} 因為體貼肉體就是與神為敵，對神的律法不順服，事實上也無法順服。 \end{tabularx} \\ \\ \relax
8:8 & \begin{tabularx}{0.7\textwidth}{X} 屬肉體的人無法使神喜悅。 \end{tabularx} \\ \\ \relax
8:9 & \begin{tabularx}{0.7\textwidth}{X} 如果神的靈住在你們裡面，你們就不屬肉體，而是屬聖靈了。人若沒有基督的靈，就不是屬基督的。 \end{tabularx} \\ \\ \relax
8:10 & \begin{tabularx}{0.7\textwidth}{X} 基督若在你們裡面，身體就因罪而死，靈卻因義而活。 \end{tabularx} \\ \\ \relax
8:11 & \begin{tabularx}{0.7\textwidth}{X} 然而，使耶穌從死人中復活的神的靈若住在你們裡面，那使基督從死人中復活的，也必藉著住在你們裡面的聖靈使你們必死的身體又活過來。 \end{tabularx} \\ \\
[1ex]
\hline
\hline
\end{longtable}
以西結書 37:1-14
\newline
\begin{longtable}{cl}
\hline
\hline
章節 & 經文 (和合本修訂版)\\
\hline
37:1 & \begin{tabularx}{0.7\textwidth}{X} 耶和華的手按在我身上。耶和華藉著他的靈帶我出去，把我放在平原中，平原遍滿骸骨。 \end{tabularx} \\ \\ \relax
37:2 & \begin{tabularx}{0.7\textwidth}{X} 他使我從骸骨的四圍經過，看哪，平原上面的骸骨甚多，看哪，極其枯乾。 \end{tabularx} \\ \\ \relax
37:3 & \begin{tabularx}{0.7\textwidth}{X} 他對我說：「人子啊，這些骸骨能活過來嗎？」我說：「主耶和華啊，你是知道的。」 \end{tabularx} \\ \\ \relax
37:4 & \begin{tabularx}{0.7\textwidth}{X} 他又對我說：「你要向這些骸骨說預言，對它們說：枯乾的骸骨啊，要聽耶和華的話。 \end{tabularx} \\ \\ \relax
37:5 & \begin{tabularx}{0.7\textwidth}{X} 主耶和華對這些骸骨如此說：『看哪，我必使氣息進入你們裡面，你們就要活過來。 \end{tabularx} \\ \\ \relax
37:6 & \begin{tabularx}{0.7\textwidth}{X} 我要給你們加上筋，長出肉，又給你們包上皮，使氣息進入你們裡面，你們就要活過來；你們就知道我是耶和華。』」 \end{tabularx} \\ \\ \relax
37:7 & \begin{tabularx}{0.7\textwidth}{X} 於是，我遵命說預言。正說預言的時候，有響聲，看哪，有地震；骨與骨彼此接連。 \end{tabularx} \\ \\ \relax
37:8 & \begin{tabularx}{0.7\textwidth}{X} 我觀看，看哪，骸骨上面有筋，長了肉，又包上皮，只是裡面還沒有氣息。 \end{tabularx} \\ \\ \relax
37:9 & \begin{tabularx}{0.7\textwidth}{X} 耶和華對我說：「人子啊，你要說預言，向風說預言。你要說，耶和華如此說：氣息啊，要從四方而來，吹在這些被殺的人身上，使他們活過來。」 \end{tabularx} \\ \\ \relax
37:10 & \begin{tabularx}{0.7\textwidth}{X} 於是我遵命說預言，氣息就進入骸骨，骸骨就活過來，並且用腳站起來，成為極大的軍隊。 \end{tabularx} \\ \\ \relax
37:11 & \begin{tabularx}{0.7\textwidth}{X} 他對我說：「人子啊，這些骸骨就是以色列全家。他們說：『看哪，我們的骨頭枯乾了，我們的指望失去了，我們滅絕淨盡了！』 \end{tabularx} \\ \\ \relax
37:12 & \begin{tabularx}{0.7\textwidth}{X} 所以你要說預言，對他們說，主耶和華如此說：我的子民，看哪，我要打開你們的墳墓，把你們帶出墳墓，領你們進入以色列地。 \end{tabularx} \\ \\ \relax
37:13 & \begin{tabularx}{0.7\textwidth}{X} 我的子民哪，我打開你們的墳墓，把你們帶出墳墓時，你們就知道我是耶和華。 \end{tabularx} \\ \\ \relax
37:14 & \begin{tabularx}{0.7\textwidth}{X} 我必將我的靈放在你們裡面，你們就要活過來。我把你們安置在本地，你們就知道我---耶和華說了這話，就必成就。這是耶和華說的。」 \end{tabularx} \\ \\
[1ex]
\hline
\hline
\end{longtable}
約翰福音 11:1-45
\newline
\begin{longtable}{cl}
\hline
\hline
章節 & 經文 (和合本修訂版)\\
\hline
11:1 & \begin{tabularx}{0.7\textwidth}{X} 有一個患病的人，名叫拉撒路，住在伯大尼，就是馬利亞和她姐姐馬大的村莊。 \end{tabularx} \\ \\ \relax
11:2 & \begin{tabularx}{0.7\textwidth}{X} 這馬利亞就是那用香膏抹主，又用頭髮擦他腳的；患病的拉撒路是她的弟弟。 \end{tabularx} \\ \\ \relax
11:3 & \begin{tabularx}{0.7\textwidth}{X} 姊妹兩個就打發人去見耶穌，說：「主啊，你所愛的人病了。」 \end{tabularx} \\ \\ \relax
11:4 & \begin{tabularx}{0.7\textwidth}{X} 耶穌聽見後卻說：「這病不至於死，而是為了神的榮耀，為要使神的兒子藉此得榮耀。」 \end{tabularx} \\ \\ \relax
11:5 & \begin{tabularx}{0.7\textwidth}{X} 耶穌素來愛馬大和她妹妹，以及拉撒路。 \end{tabularx} \\ \\ \relax
11:6 & \begin{tabularx}{0.7\textwidth}{X} 他聽見拉撒路病了，仍在原地住了兩天， \end{tabularx} \\ \\ \relax
11:7 & \begin{tabularx}{0.7\textwidth}{X} 然後對門徒說：「我們再到猶太去吧！」 \end{tabularx} \\ \\ \relax
11:8 & \begin{tabularx}{0.7\textwidth}{X} 門徒對他說：「拉比，猶太人近來要拿石頭打你，你還再到那裡去嗎？」 \end{tabularx} \\ \\ \relax
11:9 & \begin{tabularx}{0.7\textwidth}{X} 耶穌回答：「白天不是有十二小時嗎？人若在白天行走，就不致跌倒，因為他看見這世上的光。 \end{tabularx} \\ \\ \relax
11:10 & \begin{tabularx}{0.7\textwidth}{X} 人若在黑夜行走，就會跌倒，因為他沒有光。」 \end{tabularx} \\ \\ \relax
11:11 & \begin{tabularx}{0.7\textwidth}{X} 耶穌說了這些話，隨後對他們說：「我們的朋友拉撒路睡了，我去叫醒他。」 \end{tabularx} \\ \\ \relax
11:12 & \begin{tabularx}{0.7\textwidth}{X} 門徒就說：「主啊，他若睡了，就會好的。」 \end{tabularx} \\ \\ \relax
11:13 & \begin{tabularx}{0.7\textwidth}{X} 耶穌說這話是指拉撒路死了，他們卻以為他是指通常的睡眠。 \end{tabularx} \\ \\ \relax
11:14 & \begin{tabularx}{0.7\textwidth}{X} 於是耶穌就明白地告訴他們：「拉撒路死了。 \end{tabularx} \\ \\ \relax
11:15 & \begin{tabularx}{0.7\textwidth}{X} 為了你們的緣故，我不在那裡反而歡喜，為要使你們信。現在我們到他那裡去吧。」 \end{tabularx} \\ \\ \relax
11:16 & \begin{tabularx}{0.7\textwidth}{X} 於是那稱為低土馬的多馬對其他的門徒說：「我們也去和他同死吧！」 \end{tabularx} \\ \\ \relax
11:17 & \begin{tabularx}{0.7\textwidth}{X} 耶穌到了，知道拉撒路在墳墓裡已經四天了。 \end{tabularx} \\ \\ \relax
11:18 & \begin{tabularx}{0.7\textwidth}{X} 伯大尼離耶路撒冷不遠，約有六里路。 \end{tabularx} \\ \\ \relax
11:19 & \begin{tabularx}{0.7\textwidth}{X} 有好些猶太人來看馬大和馬利亞，要為她們弟弟的緣故安慰她們。 \end{tabularx} \\ \\ \relax
11:20 & \begin{tabularx}{0.7\textwidth}{X} 馬大聽見耶穌來了，就出去迎接他；馬利亞卻仍然坐在家裡。 \end{tabularx} \\ \\ \relax
11:21 & \begin{tabularx}{0.7\textwidth}{X} 馬大對耶穌說：「主啊，你若早在這裡，我弟弟就不會死了。 \end{tabularx} \\ \\ \relax
11:22 & \begin{tabularx}{0.7\textwidth}{X} 我也知道，即使現在，你無論向神求甚麼，神也必賜給你。」 \end{tabularx} \\ \\ \relax
11:23 & \begin{tabularx}{0.7\textwidth}{X} 耶穌對她說：「你弟弟會復活的。」 \end{tabularx} \\ \\ \relax
11:24 & \begin{tabularx}{0.7\textwidth}{X} 馬大對他說：「我知道在末日復活的時候，他會復活。」 \end{tabularx} \\ \\ \relax
11:25 & \begin{tabularx}{0.7\textwidth}{X} 耶穌對她說：「復活在我，生命也在我。信我的人雖然死了，也必復活。 \end{tabularx} \\ \\ \relax
11:26 & \begin{tabularx}{0.7\textwidth}{X} 凡活著信我的人必永遠不死。你信這話嗎？」 \end{tabularx} \\ \\ \relax
11:27 & \begin{tabularx}{0.7\textwidth}{X} 馬大對他說：「主啊，是的。我信你是基督，是神的兒子，就是那要臨到世界的。」 \end{tabularx} \\ \\ \relax
11:28 & \begin{tabularx}{0.7\textwidth}{X} 馬大說了這話就回去，叫她妹妹馬利亞，私下說：「老師來了，他在叫你。」 \end{tabularx} \\ \\ \relax
11:29 & \begin{tabularx}{0.7\textwidth}{X} 馬利亞聽見了，急忙起來，到耶穌那裡去。 \end{tabularx} \\ \\ \relax
11:30 & \begin{tabularx}{0.7\textwidth}{X} 那時，耶穌還沒有進村子，仍在馬大迎接他的地方。 \end{tabularx} \\ \\ \relax
11:31 & \begin{tabularx}{0.7\textwidth}{X} 那些同馬利亞在家裡安慰她的猶太人，見她急忙起來，出去，就跟著她，以為她要往墳墓那裡去哭。 \end{tabularx} \\ \\ \relax
11:32 & \begin{tabularx}{0.7\textwidth}{X} 馬利亞到了耶穌那裡，看見他，就俯伏在他腳前，對他說：「主啊，你若早在這裡，我弟弟就不會死了。」 \end{tabularx} \\ \\ \relax
11:33 & \begin{tabularx}{0.7\textwidth}{X} 耶穌看見她哭，並看見與她同來的猶太人也哭，就心裡悲嘆，又甚憂愁， \end{tabularx} \\ \\ \relax
11:34 & \begin{tabularx}{0.7\textwidth}{X} 就說：「你們把他安放在哪裡？」他們對他說：「主啊，請你來看。」 \end{tabularx} \\ \\ \relax
11:35 & \begin{tabularx}{0.7\textwidth}{X} 耶穌哭了。 \end{tabularx} \\ \\ \relax
11:36 & \begin{tabularx}{0.7\textwidth}{X} 猶太人就說：「你看，他多麼愛他！」 \end{tabularx} \\ \\ \relax
11:37 & \begin{tabularx}{0.7\textwidth}{X} 其中有人說：「他既然開了盲人的眼睛，難道不能叫這人不死嗎？」 \end{tabularx} \\ \\ \relax
11:38 & \begin{tabularx}{0.7\textwidth}{X} 耶穌又心裡悲嘆，來到墳墓前。那墳墓是個穴，有一塊石頭擋著。 \end{tabularx} \\ \\ \relax
11:39 & \begin{tabularx}{0.7\textwidth}{X} 耶穌說：「把石頭挪開！」那死者的姐姐馬大對他說：「主啊，他現在必定臭了，因為他已經死了四天了。」 \end{tabularx} \\ \\ \relax
11:40 & \begin{tabularx}{0.7\textwidth}{X} 耶穌對她說：「我不是對你說過，你若信就必看見神的榮耀嗎？」 \end{tabularx} \\ \\ \relax
11:41 & \begin{tabularx}{0.7\textwidth}{X} 於是他們把石頭挪開。耶穌舉目望天，說：「父啊，我感謝你，因為你已經聽了我。 \end{tabularx} \\ \\ \relax
11:42 & \begin{tabularx}{0.7\textwidth}{X} 我知道你常常聽我，但我說這話是為了周圍站著的眾人，要使他們信是你差了我來的。」 \end{tabularx} \\ \\ \relax
11:43 & \begin{tabularx}{0.7\textwidth}{X} 說了這些話，他大聲呼叫說：「拉撒路，出來！」 \end{tabularx} \\ \\ \relax
11:44 & \begin{tabularx}{0.7\textwidth}{X} 那死了的人就出來了，手腳都裹著布，臉上包著頭巾。耶穌對他們說：「解開他，讓他走！」 \end{tabularx} \\ \\ \relax
11:45 & \begin{tabularx}{0.7\textwidth}{X} 於是來看馬利亞的猶太人中，有很多人見了耶穌所做的事，就信了他。 \end{tabularx} \\ \\
[1ex]
\hline
\hline
\end{longtable}
講題:人在此身中

講員:林豪恩先生

經文:羅馬書8:6-11;詩篇130;以西結書37:1-14;約翰福音11:1-45

日期:2023年3月26日

一．為身體而糾結

(羅馬書)

有一樣事物,你不可以沒有「它」,但有了「它」,「它」可以令你一生都糾結不休.你想到這是甚麼嗎?你可能想到很多事物,今次我想講的是身體.試想一想,你覺得身體(包括家人的身體)帶給你快樂多些,還是擔驚受怕多些?有一首舊歌叫〈為甚麼〉,道出了不少人對身體的慨歎:

第一種是對身體老化的驚嚇:

為甚麼雙鬢斑白光彩消失面容上

為甚麼齒髮俱落一張怪模樣

第二種是身體患病帶來的困擾:

為甚麼竟會生病輾轉反側在床上

病榻中許我一問怎可永無恙

第三種是身體死亡的無奈:

為甚麼凄冷孤寂輕飄飄像無力

為甚麼不見光亮飄渺沒形象

在目前如何風光仍是泥中葬

沒法牽走一根線那許依戀臭皮囊

人們對身體的態度其實是相對矛盾的,何以見得呢?歌詞雖然用「臭皮囊」這個棄之也不可惜的名稱來稱呼身體,但卻又說「依戀」.對一般人來說,身體令人又愛又憂.如果加上宗教因素呢?在宗教來看,身體的糾結會變得簡單些,還是更加複雜呢?

最近,你有沒有聽過關於「攝理教」?可能因為藝人方力申女朋友的緣故,「攝理教」成為了熱話.大家有沒有留意到「攝理教」其實與身體的糾結非常突出?網上串流平台一套以南韓邪教和異端教派為題的紀錄片《以神之名:信仰的背叛》揭露,那些教主通常神化自己,在他們那裡,越自認神聖的,就越身體性的.為甚麼這樣說呢?因為那些神聖的教主對女性的身體最有興趣,神聖與身體結合,裡面的人稱之為救贖,外面的人指責為性侵.網絡新聞指出攝理教教主鄭明析於2009年2月於韓國因性侵被判刑10年,並於2018年2月18日刑滿獲釋.其後再次於2022年10月4日因涉嫌強姦和性騷擾被警方逮捕.因此,我們可以看見宗教與身體的糾結,越自認神聖的,就越身體性的.

基督徒又如何看待身體呢?基督徒都會為身體而糾結嗎?如果有不為身體糾結的基督徒,修為實在高超得令人敬佩,就是使徒保羅,也感到掙扎.你還記得我們剛剛讀過保羅在羅馬書的經文嗎?

8:6體貼肉體的,就是死；體貼聖靈的,乃是生命,平安.

8:7原來體貼肉體的,就是與上帝為仇；因為不服上帝的律法,也是不能服,

8:8而且屬肉體的人不能得上帝的喜歡.

8:9如果上帝的靈住在你們心裡,你們就不屬肉體,乃屬聖靈了.人若沒有基督的靈,就不是屬基督的.

8:10基督若在你們心裡,身體就因罪而死,心靈卻因義而活.

8:11然而,叫耶穌從死裡復活者的靈若住在你們心裡,那叫基督耶穌從死裡復活的,也必藉著住在你們心裡的聖靈,使你們必死的身體又活過來.

這段文字中,不斷重覆「肉體」和「身體」,又提到「死」及「活」,你能馬上疏理出當中的關係嗎?當你讀經文的時候,有沒有感覺到,肉體好像與死有關,身體似乎與活有關?然而,我們下一個問題自然會問:「肉體」和「身體」有甚麼不同呢?是同一樣東西?還是兩樣東西呢?從字義及用法來說,兩個字有交互重疊的地方,可以交替互用,因此,單從用字來說,是難以分辨兩者的差異.那麼,我們不妨按作者在某此篇幅的用法來看,例如在這小段落中,我們看見「肉體」是與上帝／聖靈相對的:

8:6體貼肉體的,就是死；體貼聖靈的,乃是生命,平安.

8:7原來體貼肉體的,就是與上帝為仇；因為不服上帝的律法,也是不能服,

8:8而且屬肉體的人不能得上帝的喜歡.

「身體」是與基督／聖靈相關的:

8:10基督若在你們心裡,身體就因罪而死,心靈卻因義而活.

8:11然而,叫耶穌從死裡復活者的靈若住在你們心裡,那叫基督耶穌從死裡復活的,也必藉著住在你們心裡的聖靈,使你們必死的身體又活過來.

因此,我們可以假設保羅在這裡對肉體和身體的描述並不是純粹物質性的,而是信仰性的,是相對於與上帝的關係而定義的.讓我們想像,身體作為生存狀態,當與上帝為敵的時候,可以墮落為行屍走肉的「肉體」,當與基督同在的時候,就提升為「身體」.

無論是希臘哲學,或者是眾多宗教,都有一種想法,就是永存的靈魂暫寄身於身體之中.受這種想法影響,有基督徒也認為只有人的靈才能與上主溝通,而上主也只透過人的靈與人接觸.不過,又有人提問,有沒有一個沒有身體而只有靈的人呢?相信大家都沒有見過.反過來說,人們會指責某些人只有身體而沒有靈魂.我們見過的人都是有身體的,就是耶穌基督死後復活,都是有身體的.因此,有人認為人不在身體之外,人在此身中.身體為甚麼如此令人糾結呢?因為人與自己的關係,與別人的關係,及別人對待自己,都離不開身體,都與身體有關.接著,讓我們透過今日經課的其他三段經文思考這三者與上帝的關係.

與自我在上帝之中的關係

(詩篇)

130:1耶和華啊,我從深處向你求告!

130:2主啊,求你聽我的聲音!願你側耳聽我懇求的聲音!

130:3主－耶和華啊,你若究察罪孽,誰能站得住呢?

130:4但在你有赦免之恩,要叫人敬畏你.

130:5我等候耶和華,我的心等候；我也仰望他的話.

130:6我的心等候主,勝於守夜的,等候天亮,勝於守夜的,等候天亮.

130:7以色列啊,你當仰望耶和華!因他有慈愛,有豐盛的救恩.

130:8他必救贖以色列脫離一切的罪孽.

130:1詩人在深處向耶和華求告,求甚麼呢?130:6詩人徹夜等候,等候甚麼呢?130:7等待慈愛和恩典.耶和華的慈愛和恩典對詩人有甚麼作用呢?130:4赦免罪孽.由此推測,詩人的深處很可能就是表達陷於自己的罪孽之中的狀況.陷於自己的罪孽之中可以是非常痛苦和惡劣的自我關係.

2023年3月15日,明報有一段新聞的標題是這樣的:「獨留孫女車上墮地亡,外公聞緩刑哭聲未止,求情稱「完全接受唔到」自己」

「完全接受唔到自己」是非常痛苦及自責的狀況,同時顯現於身體和心靈.你還記得大衛殺人丈夫奪人妻子的故事嗎?相傳他在這件事之後寫的詩篇如此訴說:

51:3因為我知道我的過犯；我的罪常在我面前.

是甚麼令他意識到「我的罪常在我面前」呢?你猜會不會是被他所殺的烏利亞現身索命的幻覺呢?

32:3我閉口不認罪的時候,因終日呻吟而骨頭枯乾.

4黑夜白日,你的手壓在我身上沉重；我的精力耗盡,如同夏天的乾旱

從這些詩句中,你能想像得到罪孽深重的人的身體狀況嗎?他的面容,他的精神,他的肌肉,他的動力,他的腸胃,他的自我形象等,都不會是你羨慕的狀況.這就是整個人的狀況,再分不開身體和心靈.面對「完全接受唔到自己」這種與自己分離的痛苦和自責,如何能夠再次活過來呢?大衛的詩篇是這樣說的:

7求你用牛膝草潔淨我,我就乾淨；求你洗滌我,我就比雪更白.

8求你使我得聽歡喜快樂的聲音,使你所壓傷的骨頭可以踴躍.

「得聽歡喜快樂的聲音」,「骨頭可以踴躍」,代表著一個人的心理狀況及身體狀況,也就是一個人的自我狀況得以融洽,由「完全接受唔到自己」改善為「接受到自己」.如此的改變,並不是靠自己的幻想,詩人的經驗從自我與上主的關係有關的,在上主的恩典和慈愛之中,詩人經驗到上主的赦免及接納,轉化成為對自己的赦免與接納.

相信我們都能夠明白那位外公的內疚,如果這不幸事件發生在我們身上,我們也會耿耿於懷.縱使法庭不把我們監禁,我們也接受不了自己,甚至內疚終生.誰能幫助我們脫離種痛苦的自我關係呢?誰把我們從過去釋放出來呢?但願這位外公能夠經驗到上主的恩典和慈愛,經驗到上主的赦免及接納,轉化成為對自己的赦免與接納.

與他者在上帝之中的關係(以西結書)

37:1耶和華的靈降在我身上.耶和華藉他的靈帶我出去,將我放在平原中；這平原遍滿骸骨.

37:2他使我從骸骨的四圍經過,誰知在平原的骸骨甚多,而且極其枯乾.

37:3他對我說:「人子啊,這些骸骨能復活嗎?」我說:「主耶和華啊,你是知道的.」

37:4他又對我說:「你向這些骸骨發預言說:枯乾的骸骨啊,要聽耶和華的話.

37:5主耶和華對這些骸骨如此說:『我必使氣息進入你們裡面,你們就要活了.

37:6我必給你們加上筋,使你們長肉,又將皮遮蔽你們,使氣息進入你們裡面,你們就要活了；你們便知道我是耶和華.』」

37:7於是,我遵命說預言.正說預言的時候,不料,有響聲,有地震；骨與骨互相聯絡.

37:8我觀看,見骸骨上有筋,也長了肉,又有皮遮蔽其上,只是還沒有氣息.

37:9主對我說:「人子啊,你要發預言,向風發預言,說主耶和華如此說:氣息啊,要從四方而來,吹在這些被殺的人身上,使他們活了.」

37:10於是我遵命說預言,氣息就進入骸骨,骸骨便活了,並且站起來,成為極大的軍隊.

37:11主對我說:「人子啊,這些骸骨就是以色列全家.他們說:『我們的骨頭枯乾了,我們的指望失去了,我們滅絕淨盡了.』

37:12所以你要發預言對他們說,主耶和華如此說:『我的民哪,我必開你們的墳墓,使你們從墳墓中出來,領你們進入以色列地.

37:13我的民哪,我開你們的墳墓,使你們從墳墓中出來,你們就知道我是耶和華.

37:14我必將我的靈放在你們裡面,你們就要活了.我將你們安置在本地,你們就知道我－耶和華如此說,也如此成就了.這是耶和華說的.』」

這是以西結先知看見的「枯骨復甦」異象.上一段我們用羅馬書經文思考過「人在此身中」的「人身」作為個體與自己及上主的關係,現在讓我們用「枯骨復甦」異象來思考「人身」作為群體和上主的關係.

既然平原滿佈骸骨,有沒有想像過耶和華建立一個骸體兵團,故事剪接如下:

7於是,我遵命說預言.正說預言的時候,不料,有響聲,有地震；骨與骨互相聯絡.10骸骨便活了,並且站起來,成為極大的軍隊.(以下是創作)骸骨兵團一步一步操向敵軍,不怕痛,不流血,敵軍無以還擊,落荒而逃.骸骨兵團成功驅逐侵略者,奪回耶路撒冷.

你認為這個策略如何呢?敵人被趕走了,耶路撒冷奪回了,那又怎樣呢?這就是以色列人的盼望嗎?其實,有時人們的確會陷入類似的迷惑中,以為驅逐了敵人,奪回了土地就成功,最後,得到甚麼呢?沒有身體沒有靈魂的骸骨?講遠一點,對著我們不願接受的東西,我們不顧一切地驅趕清除,對著我們渴望的東西,我們拼命地搶奪,最後,我們得到甚麼呢?沒有人會稱骸骨為人,能夠稱之為人的是以下描述的:

37:8我觀看,見骸骨上有筋,也長了肉,又有皮遮蔽其上,只是還沒有氣息.

37:9主對我說:「人子啊,你要發預言,向風發預言,說主耶和華如此說:氣息啊,要從四方而來,吹在這些被殺的人身上,使他們活了.」

有骨,有筋,有肉,還有從創造主而來的氣息,我們才感到這是人.我們所要的是甚麼呢?就是覺得自己是人.我們所要的是甚麼呢?難道只是權力,土地和資源嗎?就算大地在你腳下,但只有你一個人,難道你就感到滿足了嗎?

在以西結的年代,以色列人極度絕望,他們說:『我們的骨頭枯乾了,我們的指望失去了,我們滅絕淨盡了.』他們的盼望是甚麼呢?甚麼異象能夠讓他們繼續有盼望呢?

群體瓦解,個人得以苟存的以色列人,仍然感到絕望.對群體瓦解的傷痛,此時此地的人們是感受得到的.人們不知道怎樣訴說這種群體瓦解的傷痛,大家不是仍然活著嗎?何事傷春悲秋呢?俗語有話:「畫公仔唔使畫出腸」,又或者不方便「畫出腸」.還有盼望嗎?甚麼是盼望呢?

耶和華藉由「枯骨復甦」異象賜給經歷群瓦解的以色列人盼望:

37:14我必將我的靈放在你們裡面,你們就要活了.我將你們安置在本地,你們就知道我－耶和華如此說,也如此成就了.」

重建群體,這是以西結年代的以色列人的盼望；或許,也是此時此地的人的盼望.有一位心理學家(Harry Stack Sullivan)指出,我們唯有在與他人的關係中才能學會成為一個個體,並以此持續運作.

通過他者而關涉的自我與上主的關係(約翰福音)

11:1有一個患病的人,名叫拉撒路,住在伯大尼,就是馬利亞和她姊姊馬大的村莊.

11:2這馬利亞就是那用香膏抹主,又用頭髮擦他腳的；患病的拉撒路是她的兄弟.

11:3她姊妹兩個就打發人去見耶穌,說:「主啊,你所愛的人病了.」

11:4耶穌聽見,就說:「這病不至於死,乃是為上帝的榮耀,叫上帝的兒子因此得榮耀.」

11:5耶穌素來愛馬大和她妹子並拉撒路.

11:6聽見拉撒路病了,就在所居之地仍住了兩天.

11:7然後對門徒說:「我們再往猶太去吧.」

11:8門徒說:「拉比,猶太人近來要拿石頭打你,你還往那裡去嗎?」

11:9耶穌回答說:「白日不是有十二小時嗎?人在白日走路,就不致跌倒,因為看見這世上的光.

11:10若在黑夜走路,就必跌倒,因為他沒有光.」

11:11耶穌說了這話,隨後對他們說:「我們的朋友拉撒路睡了,我去叫醒他.」

11:12門徒說:「主啊,他若睡了,就必好了.」

11:13耶穌這話是指著他死說的,他們卻以為是說照常睡了.

11:14耶穌就明明地告訴他們說:「拉撒路死了.

11:15我沒有在那裡就歡喜,這是為你們的緣故,好叫你們相信.如今我們可以往他那裡去吧.」

11:16多馬,又稱為低土馬,就對那同作門徒的說:「我們也去和他同死吧.」

11:17耶穌到了,就知道拉撒路在墳墓裡已經四天了.

11:18伯大尼離耶路撒冷不遠,約有六里路.

11:19有好些猶太人來看馬大和馬利亞,要為她們的兄弟安慰她們.

11:20馬大聽見耶穌來了,就出去迎接他；馬利亞卻仍然坐在家裡.

11:21馬大對耶穌說:「主啊,你若早在這裡,我兄弟必不死.

11:22就是現在,我也知道,你無論向上帝求甚麼,上帝也必賜給你.」

11:23耶穌說:「你兄弟必然復活.」

11:24馬大說:「我知道在末日復活的時候,他必復活.」

11:25耶穌對她說:「復活在我,生命也在我.信我的人雖然死了,也必復活；

11:26凡活著信我的人必永遠不死.你信這話嗎?」

11:27馬大說:「主啊,是的,我信你是基督,是上帝的兒子,就是那要臨到世界的.」

11:28馬大說了這話,就回去暗暗地叫她妹子馬利亞,說:「夫子來了,叫你.」

11:29馬利亞聽見了,就急忙起來,到耶穌那裡去.

11:30那時,耶穌還沒有進村子,仍在馬大迎接他的地方.

11:31那些同馬利亞在家裡安慰她的猶太人,見她急忙起來出去,就跟著她,以為她要往墳墓那裡去哭.

11:32馬利亞到了耶穌那裡,看見他,就俯伏在他腳前,說:「主啊,你若早在這裡,我兄弟必不死.」

11:33耶穌看見她哭,並看見與她同來的猶太人也哭,就心裡悲歎,又甚憂愁,

11:34便說:「你們把他安放在哪裡?」他們回答說:「請主來看.」

11:35耶穌哭了.

11:36猶太人就說:「你看他愛這人是何等懇切.」

11:37其中有人說:「他既然開了瞎子的眼睛,豈不能叫這人不死嗎?」

11:38耶穌又心裡悲歎,來到墳墓前；那墳墓是個洞,有一塊石頭擋著.

11:39耶穌說:「你們把石頭挪開.」那死人的姊姊馬大對他說:「主啊,他現在必是臭了,因為他死了已經四天了.」

11:40耶穌說:「我不是對你說過,你若信,就必看見上帝的榮耀嗎?」

11:41他們就把石頭挪開.耶穌舉目望天,說:「父啊,我感謝你,因為你已經聽我.

11:42我也知道你常聽我,但我說這話是為周圍站著的眾人,叫他們信是你差了我來.」

11:43說了這話,就大聲呼叫說:「拉撒路出來!」

11:44那死人就出來了,手腳裹著布,臉上包著手巾.耶穌對他們說:「解開,叫他走!」

11:45那些來看馬利亞的猶太人見了耶穌所做的事,就多有信他的；

今天的福音經課是一個較為長篇的故事,由於時間關係,我們不再重讀這個故事.如果用「拉撒路的死和復活」作為標題,某程度上可以概括這個故事的內容了.

〈為甚麼〉這首歌最後一段哀嘆:

為甚麼凄冷孤寂輕飄飄像無力

為甚麼不見光亮飄渺沒形象

在目前如何風光仍是泥中葬

沒法牽走一根線那許依戀臭皮囊

如果再給一次機會,「臭皮囊」再次活起來,又如何呢?有一位作者認為,「每個成年人,都是劫後餘生的人」.曾幾何時,我們都曾立願:「如果我今次大能不死,我會怎樣怎樣.」「如果我今次過得到,我會怎樣怎樣.」我們不知道拉撒路有沒有立過類似的願,不過,拉撒路生命的take two,可以在我們思考這個問題時作為啟發.在這事情上,我們沒有聽到拉撒路的心聲,不過,卻聽到給予他take two的耶穌的說話:

11:4耶穌聽見,就說:「這病不至於死,乃是為上帝的榮耀,叫上帝的兒子因此得榮耀.」

這榮耀的具體顯現是怎樣的呢?

11:45那些來看馬利亞的猶太人見了耶穌所做的事,就多有信他的；

拉撒路這個人,這個身體,已經不單單是一個個體,他這個人身的存在,是為他人而存在,可以說是一個活在上主和他人之中的人身,也是一個為上主和他人而活的人身.拉撒路這個人身活著的意義,連祭司長也看見:約翰福音12章

9有許多猶太人知道耶穌在那裡,就來了,不但是為耶穌的緣故,也是要看他從死裡所復活的拉撒路.10但祭司長商議連拉撒路也要殺了；11因有好些猶太人為拉撒路的緣故,回去信了耶穌.

正如〈為甚麼〉這首歌唱出,身體變化的過程令人充滿糾結.使徒保羅也認同這種糾結,因此勸勉我們接受聖靈同在,把身體獻上當作活祭；詩人見證在上主的恩典和慈愛中如何化解自身的糾結,帶來生命的整合；以西結先知向我們揭示人在群體中的盼望；拉撒路的「返生」指引出人身在世的方向.想到這裡,我記起一位在座的弟兄曾經跟我說過:「患上這不治之症而仍然能夠存活的日子,都是為見證上主而活的日子.」



\newpage

\section{主的復活}
\label{sec:641}
\textbf{李均熊牧師 - 主的復活}
\newline
\newline
\hyperref[sec:640]{< < < PREV SERMON < < <}
~
\hyperref[sec:toc]{[返目錄]}
~
\hyperref[sec:642]{> > > NEXT SERMON > > >}


約翰福音 20:1-18
\newline
\begin{longtable}{cl}
\hline
\hline
章節 & 經文 (和合本修訂版)\\
\hline
20:1 & \begin{tabularx}{0.7\textwidth}{X} 七日的第一日清早，天還黑的時候，抹大拉的馬利亞來到墳墓，看見石頭已從墳墓挪開了， \end{tabularx} \\ \\ \relax
20:2 & \begin{tabularx}{0.7\textwidth}{X} 就跑來見西門‧彼得和耶穌所愛的那個門徒，對他們說：「有人從墳墓裡把主移走了，我們不知道他們把他放在哪裡。」 \end{tabularx} \\ \\ \relax
20:3 & \begin{tabularx}{0.7\textwidth}{X} 彼得和那門徒就出來，往墳墓去。 \end{tabularx} \\ \\ \relax
20:4 & \begin{tabularx}{0.7\textwidth}{X} 兩個人同跑，那門徒比彼得跑得快，先到了墳墓， \end{tabularx} \\ \\ \relax
20:5 & \begin{tabularx}{0.7\textwidth}{X} 低頭往裡看，看見細麻布還放在那裡，只是沒有進去。 \end{tabularx} \\ \\ \relax
20:6 & \begin{tabularx}{0.7\textwidth}{X} 西門‧彼得隨後也到了，進了墳墓，看見細麻布放在那裡， \end{tabularx} \\ \\ \relax
20:7 & \begin{tabularx}{0.7\textwidth}{X} 又看見耶穌的裹頭巾沒有和細麻布放在一起，是另在一處捲著。 \end{tabularx} \\ \\ \relax
20:8 & \begin{tabularx}{0.7\textwidth}{X} 然後先到墳墓的那門徒也進去，他看見就信了。 \end{tabularx} \\ \\ \relax
20:9 & \begin{tabularx}{0.7\textwidth}{X} 他們還不明白聖經所說耶穌必須從死人中復活的意思。 \end{tabularx} \\ \\ \relax
20:10 & \begin{tabularx}{0.7\textwidth}{X} 於是兩個門徒回自己的住處去了。 \end{tabularx} \\ \\ \relax
20:11 & \begin{tabularx}{0.7\textwidth}{X} 馬利亞卻站在墳墓外面哭。她哭的時候，低頭往墳墓裡看， \end{tabularx} \\ \\ \relax
20:12 & \begin{tabularx}{0.7\textwidth}{X} 看見兩個天使穿著白衣，在安放耶穌身體的地方坐著，一個在頭，一個在腳。 \end{tabularx} \\ \\ \relax
20:13 & \begin{tabularx}{0.7\textwidth}{X} 天使對她說：「婦人，你為甚麼哭？」她對他們說：「因為有人把我主移走了，我不知道他們把他放在哪裡。」 \end{tabularx} \\ \\ \relax
20:14 & \begin{tabularx}{0.7\textwidth}{X} 說了這些話，她轉過身來，看見耶穌站在那裡，卻不知道他是耶穌。 \end{tabularx} \\ \\ \relax
20:15 & \begin{tabularx}{0.7\textwidth}{X} 耶穌問她：「婦人，你為甚麼哭？你找誰？」馬利亞以為他是看園子的，就對他說：「先生，若是你把他移了去，請告訴我，你把他放在哪裡，我去把他移回來。」 \end{tabularx} \\ \\ \relax
20:16 & \begin{tabularx}{0.7\textwidth}{X} 耶穌對她說：「馬利亞。」馬利亞轉過身來，用希伯來話對他說：「拉波尼！」（「拉波尼」就是老師的意思。） \end{tabularx} \\ \\ \relax
20:17 & \begin{tabularx}{0.7\textwidth}{X} 耶穌對她說：「不要拉住我，因為我還沒有升上去見我的父。你到我弟兄那裡去告訴他們，我要升上去見我的父，也是你們的父，見我的神，也是你們的神。」 \end{tabularx} \\ \\ \relax
20:18 & \begin{tabularx}{0.7\textwidth}{X} 抹大拉的馬利亞就向門徒報信：「我已經看見了主。」她又把主對她說的話告訴他們。 \end{tabularx} \\ \\
[1ex]
\hline
\hline
\end{longtable}
講題:主的復活

經文:約翰福音20章1至18節

講員:李均熊牧師

日期:2023年4月9日

起:相認場景

約翰福音有很多段記載都用上希臘劇場常用的手法,就是所謂的相認場景.在這些場景中,很多時台上角色不能相認眼前人,但台下觀眾卻明確知道這一個眼前人是誰.台下觀眾比台上的角色知道的更多,台下觀眾為台上的角色或是緊張,或是失笑.

的確,是馬利亞不識眼前人,而我們卻是清清楚楚知道與馬利亞說話的人就是復活的主耶穌基督.

可能正因為這樣,歷來聖經詮釋者對發現空墳墓的馬利亞的評價都頗為負面,認為她雖見到耶穌,卻不認到耶穌,因乃其信心不足.但我作為旁觀者,反而覺得馬利亞更接近我.試想想,二,三天前才親眼在十架上看著耶穌斷氣,現在一大早去到墳墓只想看看屍身狀況,但卻發現所愛之主軀體不知所蹤,相信普通人,甚至應該任何人,在那一刻都會迷惘不知所措.再突然間有兩三人出現在眼前,我反而覺得馬利亞沒有被嚇走,仍然堅持要求眼前人將耶穌找回來,其對主的忠誠和意志比我們這些任何人都更堅定.

承:重複:你為甚麼哭

作為觀眾的我們,自然會看得出馬利亞之所以認不出耶穌,除了是沒有預計兩日前受苦受難的主會復活以外,也是因為心靈憂傷,淚眼迷糊,而看不情眼前人就是耶穌.

天使和復活的主耶穌,都說了同一段說話:「婦人,你為甚麼哭呢?」

天使和耶穌好像提醒馬利亞,痛苦要過去,也已過去了,不要再停留在哀傷和疼痛中.過度的哀傷和悲痛不單不必要,也會影響我們對於當下生命的景況作出錯誤的評估.

轉:重複:馬利亞轉身

同樣重複出現在今天復活經課的,就是每每與眼前人對答之後,馬利亞都「轉過身來」:

13天使對她說:「婦人,你為甚麼哭?」她對他們說:「因為有人把我主移走了,我不知道他們把他放在哪裡.」14說了這些話,她轉過身來,看見耶穌站在那裡,卻不知道他是耶穌.

16耶穌對她說:「馬利亞.」馬利亞轉過身來,用希伯來話對他說:「拉波尼!」(「拉波尼」就是老師的意思.)

若果轉個身來直接解作身體轉動,我們會看到一個在團團轉的馬利亞,第一次轉身可以看成她從面對墓穴內的兩個白衣人轉向墓外,從而發現另一人站在穴口；但第二次轉身,不就從對她說話的耶穌轉回面向墓穴內的兩人,那她又怎會發現和她說話的就是復活的主耶穌呢?

其實中文「回過神來」(英文或譯come to her senses)可能會譯得更貼切(可能第一次仍可譯轉過身來?).馬利亞並非只是將身體轉來轉去,而是回過神來,將心思而不是身體轉回(turn her mind but not her body),專注和注視起眼前境況(focusing and paying attention on the present situation),開始弄清發生甚麼事情(and start to comprehending what is going on).可能因為有人和她對話,甚至耶穌直呼其名,都令馬利亞回過神來,而不是陷入自己的哀傷而與現況脫節了.

合:去我兄弟姊妹那裡,去告訴他們

最後,耶穌叫馬利亞不要拉住自己,而是派她去到「我的弟兄(姊妹)那裡去」,去他們中間做甚麼呢?就是去向他們「宣告」(原文和第11節中馬利亞站著「哭」,和第16節馬利亞說出拉波尼時的「轉過來」,都是從屬主動詞的分詞participle)耶穌的去向,就是「要升上去見我的父,也是你們的父,見我的神,也是你們的神.」

馬利亞從目睹耶穌的受難而帶來的痛苦和哀傷,失望和迷網,透過耶穌的說話,慢慢回過神來,終於見到在她面前的空墳墓所代表的是主耶穌基督的復活.但遇見復活的主並非沉醉於喜樂的時候,耶穌更要差派她去到弟兄姊妹當中,見證主的復活所帶來的應許和新的可能性.



\newpage

\section{是結束也是開始}
\label{sec:642}
\textbf{鄧瑞強博士 - 是結束,也是開始}
\newline
\newline
\hyperref[sec:641]{< < < PREV SERMON < < <}
~
\hyperref[sec:toc]{[返目錄]}
~
\hyperref[sec:643]{> > > NEXT SERMON > > >}


歷代志下 36:5-23
\newline
\begin{longtable}{cl}
\hline
\hline
章節 & 經文 (和合本修訂版)\\
\hline
36:5 & \begin{tabularx}{0.7\textwidth}{X} 約雅敬登基的時候年二十五歲，在耶路撒冷作王十一年。他行耶和華－他神眼中看為惡的事。 \end{tabularx} \\ \\ \relax
36:6 & \begin{tabularx}{0.7\textwidth}{X} 巴比倫王尼布甲尼撒上來攻擊他，用銅鏈鎖著他，要把他帶到巴比倫去。 \end{tabularx} \\ \\ \relax
36:7 & \begin{tabularx}{0.7\textwidth}{X} 尼布甲尼撒又將耶和華殿裡的一些器皿帶到巴比倫，放在巴比倫自己的宮裡。 \end{tabularx} \\ \\ \relax
36:8 & \begin{tabularx}{0.7\textwidth}{X} 約雅敬其餘的事和他所行可憎的事，以及發生在他身上的事，看哪，都寫在《以色列和猶大列王記》上，他兒子約雅斤接續他作王。 \end{tabularx} \\ \\ \relax
36:9 & \begin{tabularx}{0.7\textwidth}{X} 約雅斤登基的時候年八歲，在耶路撒冷作王三個月十天，他行耶和華眼中看為惡的事。 \end{tabularx} \\ \\ \relax
36:10 & \begin{tabularx}{0.7\textwidth}{X} 過了一年，尼布甲尼撒王差遣人將約雅斤和耶和華殿裡寶貴的器皿帶到巴比倫，然後立約雅斤的叔父西底家作猶大和耶路撒冷的王。 \end{tabularx} \\ \\ \relax
36:11 & \begin{tabularx}{0.7\textwidth}{X} 西底家登基的時候年二十一歲，在耶路撒冷作王十一年。 \end{tabularx} \\ \\ \relax
36:12 & \begin{tabularx}{0.7\textwidth}{X} 他行耶和華－他神眼中看為惡的事，沒有謙卑聽從耶利米先知所傳達耶和華的話。 \end{tabularx} \\ \\ \relax
36:13 & \begin{tabularx}{0.7\textwidth}{X} 尼布甲尼撒王曾叫他指著神起誓，他卻背叛，硬著頸項，內心頑固，不歸向耶和華－以色列的神。 \end{tabularx} \\ \\ \relax
36:14 & \begin{tabularx}{0.7\textwidth}{X} 眾祭司長和百姓也多多犯罪，效法列國一切可憎的事，玷污耶和華在耶路撒冷分別為聖的殿。 \end{tabularx} \\ \\ \relax
36:15 & \begin{tabularx}{0.7\textwidth}{X} 耶和華－他們列祖的神因為愛惜自己的百姓和居所，一再差遣使者去警戒他們。 \end{tabularx} \\ \\ \relax
36:16 & \begin{tabularx}{0.7\textwidth}{X} 他們卻嘲笑神的使者，藐視他的話，譏誚他的先知，以致耶和華向他的百姓大發烈怒，甚至無法可救。 \end{tabularx} \\ \\ \relax
36:17 & \begin{tabularx}{0.7\textwidth}{X} 所以，耶和華使迦勒底人的王來攻擊他們，在他們聖殿裡用刀殺了他們的壯丁，不憐憫他們的少男少女、老人長者。耶和華把所有的人都交在他手裡。 \end{tabularx} \\ \\ \relax
36:18 & \begin{tabularx}{0.7\textwidth}{X} 他把神殿裡一切的大小器皿與耶和華殿裡的財寶，以及王和眾領袖的財寶，全都帶到巴比倫去。 \end{tabularx} \\ \\ \relax
36:19 & \begin{tabularx}{0.7\textwidth}{X} 迦勒底人焚燒了神的殿，拆毀耶路撒冷的城牆，用火燒了城裡所有的宮殿，毀壞了城裡一切寶貴的器皿。 \end{tabularx} \\ \\ \relax
36:20 & \begin{tabularx}{0.7\textwidth}{X} 凡脫離刀劍的倖存者，迦勒底王都擄到巴比倫去，作他和他子孫的僕婢，直到波斯國興起。 \end{tabularx} \\ \\ \relax
36:21 & \begin{tabularx}{0.7\textwidth}{X} 這就應驗耶和華藉耶利米的口所說的話：地得享安息；在荒涼的日子，地就守安息，直到滿了七十年。 \end{tabularx} \\ \\ \relax
36:22 & \begin{tabularx}{0.7\textwidth}{X} 波斯王居魯士元年，耶和華為要應驗藉耶利米的口所說的話，就激發波斯王居魯士的心，使他下詔書通告全國，說： \end{tabularx} \\ \\ \relax
36:23 & \begin{tabularx}{0.7\textwidth}{X} 「波斯王居魯士如此說：耶和華－天上的神已將地上萬國賜給我，又委派我在猶大的耶路撒冷為他建造殿宇。你們中間凡作他子民的可以上去，願耶和華－他的神與他同在。」 \end{tabularx} \\ \\
[1ex]
\hline
\hline
\end{longtable}
講題:是結束,也是開始

經文:歷代志下36:5-23

講員:鄧瑞強博士

日期:2023年4月16日

今日,我們來看《歷代志下》的經文.在希伯來聖經裡,《歷代志》是最後的一卷書,這與我們的聖經裡的經卷編排是不一樣的.《歷代志》書寫的歷史,明顯先於《以斯拉記》及《尼希米記》,但在希伯來聖經裡,《歷代志》卻排在《以斯拉記》及《尼希米記》之後.如此說來,經卷的編排,不是按時序的,而是按神學的.希伯來聖經的編排方式,反映以色列人的神學；我們聖經的編排方式,反映我們的神學.

我們今日的講道經文,是《歷代志下》結束的段落.在希伯來聖經裡,這就是壓尾的經文.以色列人以這段經文作為全本聖經的結束,一方面,是以此作為以色列民族過去的歷史的蓋棺定論；另一方面,在歷史的終結處,思索如何可能有歷史的前景.

今日的講道經文,是《歷代志下》36:5-23.這段經文是快鏡,迅速交代猶大國的滅亡.我們小段小段的看,先看《歷代志下》36:5-8:

36:5約雅敬登基的時候年二十五歲,在耶路撒冷作王十一年,行耶和華－他上帝眼中看為惡的事.

36:6巴比倫王尼布甲尼撒上來攻擊他,用銅鍊鎖著他,要將他帶到巴比倫去.

36:7尼布甲尼撒又將耶和華殿裡的器皿帶到巴比倫,放在他神的廟裡.

36:8約雅敬其餘的事和他所行可憎的事,並他一切的行為,都寫在以色列和猶大列王記上.他兒子約雅斤接續他作王.

數句經文,描繪了約雅敬王的過場.這段經文交代了三件事.

一,約雅敬王「行耶和華－他上帝眼中看為惡的事」,這解釋了他在歷史舞台上匆匆謝幕的理由.不用說其他的事了,就這一件事已是致命的了.

二,緊接著約雅敬行惡的描述之後,便是「巴比倫王尼布甲尼撒上來攻擊他,用銅鍊鎖著他,要將他帶到巴比倫去」.讀起來,約雅敬的行惡,即時招來了神的懲罰.在這裡,經文讓我們看到,神使用的歷史工具,可以是巴比倫王尼布甲尼撒.神的舞台,不再只是局限在猶太民族及巴勒斯坦地,而是擴大到萬國萬民.神召喚來使用的人,不單只是神的子民,也可以是神的子民的對頭人.

約雅敬王被擄去巴比倫,用古時的事件來類比,就是重返埃及為奴.曾經,以色列人離開埃及,經歷神的解救,是要在神應許的土地上,建立一和平與公義的國度.他們立了王,建立了一個小小的王國.歷史證明,這個應該有神聖意味的國,卻與任何一個墮落的國無異.神的子民所作的,是神眼中看為惡的事.他們建立的王國,實質上是欺壓人民的王國.在所羅門王治下,竟有奴隸,與埃及沒有分別.就生活的內容而言,他們看似逃離了埃及,實質是建造了另一個埃及.他們沒有離開過埃及,神就藉尼布甲尼撒王的手,帶他們重返埃及.

三,經文刻意描寫「尼布甲尼撒將耶和華殿裡的器皿帶到巴比倫」.大家還記得嗎,《歷代志》特別看重「聖殿」,「聖殿」是神子民生命及身分的核心.神的子民在人間最重要的工作,便是建造「聖殿」,在其中敬拜.當神的榮光臨在聖殿中時,世界就找到意義的根源.但神的子民忘了實踐神的正義,忘了將神的旨意行在地上,如同行在天上,忘了建造聖殿的天命,於是,聖殿的器皿便被拿走了.猶如燈臺上的燈不再發光,便被拿走一樣.

我們看《歷代志下》36:9-10:

36:9約雅斤登基的時候年八歲,在耶路撒冷作王三個月零十天,行耶和華眼中看為惡的事.

36:10過了一年,尼布甲尼撒差遣人將約雅斤和耶和華殿裡各樣寶貴的器皿帶到巴比倫,就立約雅斤的叔叔西底家作猶大和耶路撒冷的王.

經文更短,約雅斤王過場得更快.他沒有在神面前,思考他父親的命運.他重複他父親的三件事情.他「行耶和華眼中看為惡的事」,他被擄去巴比倫,聖殿的器皿被拿走.

我們看《歷代志下》36:11-14:

36:11西底家登基的時候年二十一歲,在耶路撒冷作王十一年,

36:12行耶和華－他上帝眼中看為惡的事.先知耶利米以耶和華的話勸他,他仍不在耶利米面前自卑.

36:13尼布甲尼撒曾使他指著上帝起誓,他卻背叛,強項硬心,不歸服耶和華－以色列的上帝.

36:14眾祭司長和百姓也大大犯罪,效法外邦人一切可憎的事,污穢耶和華在耶路撒冷分別為聖的殿.

西底家是猶大國最後一個王了,這裡的話,是對猶大亡國的蓋棺定論.國之將亡,西底家竟仍然「行耶和華－他上帝眼中看為惡的事」.先知有沒有勸導他?有,但西底家王並沒有謙卑下來,即是說,他行惡是義無反顧的.這裡也指出西底家的政治失誤,他背叛奉神的名與巴比倫王立的盟約.如此,就是一雙重背叛,背叛盟友巴比倫王,也背叛神.

在上有政治上背信棄義的王,在下便有宗教上背信棄義的祭司和百姓.經文說,祭司和百姓「污穢耶和華在耶路撒冷分別為聖的殿」,用今天的情境來說,就是基督徒在教會裡光天化日下作羞辱基督的事情.這是完全違背自己的身分的.如此,猶大國由上至下,在信仰上與道德上凋零,敗局已定.

我們看《歷代志下》36:15-16:

36:15耶和華－他們列祖的上帝因為愛惜自己的民和他的居所,從早起來差遣使者去警戒他們.

36:16他們卻嘻笑上帝的使者,藐視他的言語,譏誚他的先知,以致耶和華的忿怒向他的百姓發作,無法可救.

約雅敬,約雅斤,西底家,一連三代君王,都行「耶和華眼中看為惡的事」,神才終極刑罰猶大國.三代過去了,都沒有改變,神才下定決心處置他們,這是因為神「愛惜自己的民和他的居所」,這就是舊約常說的:「上帝不輕易發怒,且有豐盛的慈愛」.就因為「愛惜」,神的審判多麼遲緩,也多番差遣先知去警戒他們.這裡的用字很特別,神「從早起來」差遣使者去警戒他們,這「從早起來」有急不及待的含意.在善惡的事上,神比人更緊張.人在惡事上徘徊,不急於改正,神比我們更焦急,希望我們回轉.

面對先知的警戒,經文用三個動詞來表達神的子民的反應:嘻笑,藐視,譏誚,與新約裡猶太人對耶穌的反應一樣.如此,神的忿怒發出了,因為猶太人無藥可救.

這兩節經文總括了猶大國的歷史和命運.從天的角度來看,神已極有憐憫,也多番透過先知發出上天的信息；從人的角度來看,猶大國冥頑不靈,無藥可救；從地的角度來看,猶太人將無地立足,被擄他鄉.

我們看《歷代志下》36:17-21:

36:17所以,耶和華使迦勒底人的王來攻擊他們,在他們聖殿裡用刀殺了他們的壯丁,不憐恤他們的少男處女,老人白叟.耶和華將他們都交在迦勒底王手裡.

36:18迦勒底王將上帝殿裡的大小器皿與耶和華殿裡的財寶,並王和眾首領的財寶,都帶到巴比倫去了.

36:19迦勒底人焚燒上帝的殿,拆毀耶路撒冷的城牆,用火燒了城裡的宮殿,毀壞了城裡寶貴的器皿.

36:20凡脫離刀劍的,迦勒底王都擄到巴比倫去,作他和他子孫的僕婢,直到波斯國興起來.

36:21這就應驗耶和華藉耶利米口所說的話:地享受安息；因為地土荒涼便守安息,直滿了七十年.

歷史的興亡,表面有其種種原因,背後卻是神的作為.將歷史的悲劇訴之於上帝,對很多人來說,就是上帝要為這些悲劇負責,上帝是邪惡的.猶太人不是這樣思考的.他們認為,若這歷史的悲劇不訴之於上帝的話,則這悲劇便已是徹底的黑暗與絕望.此話何解?想想,若歷史的悲劇出於命運的無情之力,則怎能知道命運會在下一步給我們光明?若歷史的悲劇來自事件本身的因果,而事情已到了無法翻盤的局面,則事件本身從何來一扭轉乾坤的「因」?若歷史的悲劇源於人自身的無用,則最弱勢的人何來戰勝巨人的力量?對脆弱的猶太人來說,將歷史的悲劇歸因於人間的任何原因,只會帶來絕望.他們拒絕絕望.他們將歷史的悲劇歸因於上帝,是因為當人間的一切只帶給他們更深的黑暗時,只有上帝能給他們帶來光明.一般人認為,將悲劇訴之於上帝,只會顯出上帝是邪惡的.猶太人不是這樣思考的,當他們將悲劇訴之於上帝時,他們才見到光明的可能.猶太的神學家在這裡落下沉重的一筆:「耶和華使迦勒底人的王來攻擊他們,在他們聖殿裡用刀殺了他們的壯丁,不憐恤他們的少男處女,老人白叟.耶和華將他們都交在迦勒底王手裡.」

神興起迦勒底人,即巴比倫人.經文提到,巴比倫人是「不憐恤」的.「憐恤」這個字,與上文提及的神的「愛惜」那「愛惜」的用字,是同一個字.巴比倫人「不憐恤」,好對照神的「憐恤」.但神縱然「憐恤」,猶太人卻毫不稀罕,毫不珍惜.巴比倫人的「不憐恤」,是徹底的「不憐恤」.他們屠殺猶太人,從少到老.他們掠去聖殿及貴族的財寶.他們破壞聖殿,城牆,宮殿.未被殺的,他們擄去巴比倫作奴隸.如此,猶大地一片荒涼,又回到創世記開頭的經文,「地是空虛混沌」,等待著神的重新創造.猶太人成了奴隸,又回到出埃及記的經文,等待著神新一回合的出埃及救贖作為.一切又回到虛無的太初.

然而,就是到了這絕望的深淵,已看到神創造與救贖的苗頭.第20節提及一個時機,神會興起新的帝國,到時,黑暗的歲月會過去.第21節提及一個時限,70年,這是一段讓罪惡在猶太地平息下來的時間.在創世記裡,神創造的高峰是「安息」,諷刺的是,只有當猶大亡國時,這創造的高峰才能重現.過了這罪惡止息的70年,便會有新的開始了.

我們看希伯來聖經結尾的兩節經文,《歷代志下》36:22-23:

36:22波斯王塞魯士元年,耶和華為要應驗藉耶利米口所說的話,就激動波斯王塞魯士的心,使他下詔通告全國,說:

36:23「波斯王塞魯士如此說:耶和華－天上的上帝已將天下萬國賜給我,又囑咐我在猶大的耶路撒冷為他建造殿宇.你們中間凡作他子民的,可以上去,願耶和華－他的上帝與他同在.」

波斯王塞魯士元年,神創造及救贖的新開始.神興起巴比倫,神同樣也可興起波斯.萬國的王,都由神主宰其興亡.塞魯士王下詔書,宣稱天上的神明祝福他.當然,他不是敬拜耶和華的,但他覺得萬國的神明也祝福他,讓他得到天下,故他會順帶提及耶和華的名字.在這詔書中,提及三點.

一,塞魯士王肯定上天對他的揀選,對猶太讀者來說,這就是耶和華以他來開啟一個新的機遇.

二,塞魯士王讓猶太人回耶路撒冷建聖殿,對塞魯士來說,他只是寄望萬國的神明祝福他.對猶太讀者來說,這就是耶和華透過他來進行一次新的創造.重建聖殿,就是重建一個神人合一的世界.《歷代志》最重視「聖殿」,重建聖殿,就是重建世界,也重建沒有了身分的猶太人,重建一新的人類.這是新的創造.我們在這裡,聽見神說:「要有光」,就有了光.

三,塞魯士王讓猶太人自由地離開被擄的地方,重返故鄉.最後的23節最後的一個字是「可以上去」,一般的英文翻譯是「and let him go up」.在《出埃及記》3:8,神向摩西顯現,對摩西說:「我下來是要救他們脫離埃及人的手,領他們出了那地,到美好,寬闊,流奶與蜜之地…」,「領他們出」與「可以上去」乃同一個字.《歷代志》最後一個字,就是出埃及事件的用字.神要帶領奴隸離開為奴之地,進入神應許之地.猶太讀者讀到《歷代志》的結尾,便意識到這是新的出埃及事件.

對猶太人來說,亡國被擄,是他們民族歷史的斷裂與終結.事實上,很多民族被征服及同化後,便在歷史上消失了.歷史本身若果是一嚴密的因果系列,而歷史的終結猶如人的個體的死亡,而死亡本身又沒有令人復生的因,則歷史的終結只會帶來沒有歷史前景的絕望.

但猶太民族沒有視歷史本身是一嚴密的因果系列,而視歷史為一向神開放的系統,神在歷史之上掌管著歷史,神在歷史的將來引導著歷史的前進.歷史的斷裂不源自歷史內部的因果,而來自神收回祂的憐恤.歷史斷裂後,能否再現歷史的前景,猶如耶穌死後,能否超越死亡的因果之力而復活,全賴神的介入與恩典.猶大亡國後,他們的歷史前景消失了.但,神介入,祂興起新的王,揭開歷史新的一頁,在混沌的世界,重新創造天地,再次救贖絕望的人.

就因為猶太民族倚靠歷史以外,歷史以上,歷史的前方的神,縱使他們在黑暗的深淵,仍以信仰的眼睛,看到由神所開展的,由神所應許的,由神以祂權能的命令所創造出來的一個想也想像不到的新世界.從聖經神學來講,這就是由「神的應許」所開啟的全新的將來.從系統神學來講,這就是終末論裡期待的神所創造的全新天地.新的世界是會出現的,但不由人的強力帶來,而由神的恩典帶來.強權並不創造歷史,歷史由神主宰.神是萬王之上的王,是歷史斷裂與延續的主宰,是新天地的創造者,是受苦者的救贖主.到了這裡,我們明白猶太人為何將《歷代志》放在希伯來聖經終結的地方,因為《歷代志》的結尾,講的是猶大國的結束,同時,言說創造與救贖的新的開始.

若你身處黑暗的深淵,願你聽到神說:「光明,出現吧」!



\newpage

\section{如果耶穌顯現你認得出他嗎}
\label{sec:643}
\textbf{林豪恩先生 - 如果耶穌顯現,你認得出他嗎?}
\newline
\newline
\hyperref[sec:642]{< < < PREV SERMON < < <}
~
\hyperref[sec:toc]{[返目錄]}
~
\hyperref[sec:644]{> > > NEXT SERMON > > >}


路加福音 24:13-35
\newline
\begin{longtable}{cl}
\hline
\hline
章節 & 經文 (和合本修訂版)\\
\hline
24:13 & \begin{tabularx}{0.7\textwidth}{X} 同一天，門徒中有兩個人往一個村子去；這村子名叫以馬忤斯，離耶路撒冷約有二十五里。 \end{tabularx} \\ \\ \relax
24:14 & \begin{tabularx}{0.7\textwidth}{X} 他們彼此談論所發生的這一切事。 \end{tabularx} \\ \\ \relax
24:15 & \begin{tabularx}{0.7\textwidth}{X} 正交談議論的時候，耶穌親自走近他們，和他們同行， \end{tabularx} \\ \\ \relax
24:16 & \begin{tabularx}{0.7\textwidth}{X} 可是他們的眼睛模糊了，沒認出他。 \end{tabularx} \\ \\ \relax
24:17 & \begin{tabularx}{0.7\textwidth}{X} 耶穌對他們說：「你們一邊走一邊談，彼此談論的是甚麼事呢？」他們就站住，臉上帶著愁容。 \end{tabularx} \\ \\ \relax
24:18 & \begin{tabularx}{0.7\textwidth}{X} 兩人中有一個名叫革流巴的回答：「你是在耶路撒冷的旅客中，惟一還不知道這幾天在那裡發生了甚麼事的人嗎？」 \end{tabularx} \\ \\ \relax
24:19 & \begin{tabularx}{0.7\textwidth}{X} 耶穌對他們說：「甚麼事呢？」他們對他說：「就是拿撒勒人耶穌的事。他是個先知，在神和眾百姓面前，說話行事都大有能力。 \end{tabularx} \\ \\ \relax
24:20 & \begin{tabularx}{0.7\textwidth}{X} 祭司長們和我們的官長竟把他解去，定了死罪，釘在十字架上。 \end{tabularx} \\ \\ \relax
24:21 & \begin{tabularx}{0.7\textwidth}{X} 但我們素來所盼望要救贖以色列民的就是他。不但如此，這些事發生到現在已經三天了。 \end{tabularx} \\ \\ \relax
24:22 & \begin{tabularx}{0.7\textwidth}{X} 還有，我們中間的幾個婦女使我們驚奇：她們清早去了墳墓， \end{tabularx} \\ \\ \relax
24:23 & \begin{tabularx}{0.7\textwidth}{X} 不見他的身體，就回來告訴我們，說她們看見了天使顯現，說他活了。 \end{tabularx} \\ \\ \relax
24:24 & \begin{tabularx}{0.7\textwidth}{X} 又有我們的幾個人往墳墓那裡去，所發現的正如婦女們所說的，只是沒有看見他。」 \end{tabularx} \\ \\ \relax
24:25 & \begin{tabularx}{0.7\textwidth}{X} 耶穌對他們說：「無知的人哪，先知所說的一切話，你們的心信得太遲鈍了。 \end{tabularx} \\ \\ \relax
24:26 & \begin{tabularx}{0.7\textwidth}{X} 基督不是必須受這些苦難，然後進入他的榮耀嗎？」 \end{tabularx} \\ \\ \relax
24:27 & \begin{tabularx}{0.7\textwidth}{X} 於是，他從摩西和眾先知起，凡經上所指著自己的話都給他們作了解釋。 \end{tabularx} \\ \\ \relax
24:28 & \begin{tabularx}{0.7\textwidth}{X} 他們走近所要去的村子，耶穌好像還要往前走， \end{tabularx} \\ \\ \relax
24:29 & \begin{tabularx}{0.7\textwidth}{X} 他們卻強留他說：「時候晚了，天快黑了，請你同我們住下吧。」耶穌就進去，要同他們住下。 \end{tabularx} \\ \\ \relax
24:30 & \begin{tabularx}{0.7\textwidth}{X} 坐下來和他們用餐的時候，耶穌拿起餅來，祝福了，擘開，遞給他們。 \end{tabularx} \\ \\ \relax
24:31 & \begin{tabularx}{0.7\textwidth}{X} 他們的眼睛開了，這才認出他來。耶穌卻從他們眼前消失了。 \end{tabularx} \\ \\ \relax
24:32 & \begin{tabularx}{0.7\textwidth}{X} 他們彼此說：「在路上他和我們說話，給我們講解聖經的時候，我們的心在我們裡面豈不是火熱的嗎？」 \end{tabularx} \\ \\ \relax
24:33 & \begin{tabularx}{0.7\textwidth}{X} 於是他們立刻起身，回耶路撒冷去，看見十一個使徒和與他們正在一起的人聚集在一處， \end{tabularx} \\ \\ \relax
24:34 & \begin{tabularx}{0.7\textwidth}{X} 說：「主果然復活了，已經顯現給西門看了。」 \end{tabularx} \\ \\ \relax
24:35 & \begin{tabularx}{0.7\textwidth}{X} 於是，兩個人把路上所遇到，和耶穌擘餅的時候怎麼被他們認出來的事，都述說了一遍。 \end{tabularx} \\ \\
[1ex]
\hline
\hline
\end{longtable}
講題:如果耶穌顯現,你認得出他嗎?

經文:路加福音24:13-35

講員:林豪恩先生

日期:2023年4月23日

一.引言

有沒有聽過有人說:「如果要我信耶穌,除非耶穌出來見我.」其實,事情沒有那麼簡單.試想想,就算耶穌現在就出現,你肯定這就是耶穌嗎?就算他說自己就是耶穌,你就就信他是嗎?不要說二千多年後的你和我,我們剛才讀的音經課顯示,就算是耶穌時代那兩位在以馬忤斯路上的門徒,都認不出復活的耶穌呢!

我們再讀這個故事,並不是為了耶穌突然出現在作準備,而是用來思信仰的元素.有些人喜歡把信仰約簡至「一點通」:「不要又這又那,只要告訴我唯一重要的是甚麼就夠了.」「唔好同我講咁多,最緊要係靈定唔靈!」；「唔使睇咁多書既,那些都是人的思想,睇聖經已經足夠!」是否可以簡化至「一招了」呢?從以馬忤斯路上的兩個門徒如何認出耶穌這個片段來看,信仰之路並不是一招了,而是一套組合拳.讓我們細看這個片段,窺探這套組合拳包含的招式.

二.經文及內容

1.疑問和傾談／理性討論

24:13同一天,門徒中有兩個人往一個村子去；這村子名叫以馬忤斯,離耶路撒冷約有二十五里.

24:14他們彼此談論所發生的這一切事.

24:15正交談議論的時候,耶穌親自走近他們,和他們同行,

24:16可是他們的眼睛模糊了,沒認出他.

為甚麼那兩位門徒認不出耶穌呢?作者提出的理由是「他們的眼睛模糊了」.解決「眼睛模糊」的困難,用的不是眼鏡或眼藥水.用甚麼呢?

24:17耶穌對他們說:「你們一邊走一邊談,彼此談論的是甚麼事呢?」他們就站住,臉上帶著愁容.

24:18兩人中有一個名叫革流巴的回答:「你是在耶路撒冷的旅客中,惟一還不知道這幾天在那裡發生了甚麼事的人嗎?」

耶穌用的方法是和他們傾談,讓他們表達所見所聞,所思所想.對於他們的疑惑,你認為耶穌有沒有答案呢?但耶穌不心急講答案,反而引導他們繼續講.

24:19耶穌對他們說:「甚麼事呢?」他們對他說:「就是拿撒勒人耶穌的事.他是個先知,在神和眾百姓面前,說話行事都大有能力.

24:20祭司長們和我們的官長竟把他解去,定了死罪,釘在十字架上.

24:21但我們素來所盼望要救贖以色列民的就是他.不但如此,這些事發生到現在已經三天了.

24:22還有,我們中間的幾個婦女使我們驚奇:她們清早去了墳墓,

24:23不見他的身體,就回來告訴我們,說她們看見了天使顯現,說他活了.

24:24又有我們的幾個人往墳墓那裡去,所發現的正如婦女們所說的,只是沒有看見他.」

對於信仰,或多或少,每個人都有所聽聞,也有自己的想法和疑惑.在信仰裡面,連耶穌也接受這種探索的過程,參與在聆聽與傾談之中,這是值得我們反省的.有人投訴基督徒傳福音時不聆聽別人心聲,不理會別人的狀況,硬銷福音.其實,不單單在傳福音時,就是在信仰的任何一個階段,都需要傾談討論,才能互相明白,互相啟發,一齊成長.

甚麼地方是不許問,不許講的呢?有一位攝理教的離教者指出:「他們會「教育」信徒不可在網上觀看任何反對該教的事,否則「靈魂就會中槍」及「落地獄」.除此之外,他們亦會用不同理由說服教徒,令他們逐漸「斷六親」.」

有學者(Ronald Enroth)的研究指出,異端用來掌控成員的其中一項工具是思想控制.他們不鼓勵批判性思考,禁止任何批評教主或教會的言論.

耶穌接受門徒表達疑惑,也和門徒傾談他們的疑惑.傾談思考,理性討論,是信仰其中一項重要元素.崇基禮拜堂不但鼓勵人思考和討論,甚至連講道都希望會眾參與.你有沒有留意到報告事項中有「圓桌講壇」事項?崇拜講員邀請大家一齊預先閱讀經課,一齊分享領受,希望講道成為聖靈在會眾群體的互動與回應的成果.

傾談和討論是信仰不可或缺的元素,不過,有甚麼可以作為有益於信仰成長的傾談基礎呢?有甚麼可以啟發傾談者呢?我們繼續讀這個故事.

2.聖經

24:25耶穌對他們說:「無知的人哪,先知所說的一切話,你們的心信得太遲鈍了.

24:26基督不是必須受這些苦難,然後進入他的榮耀嗎?」

24:27於是,他從摩西和眾先知起,凡經上所指著自己的話都給他們作了解釋.

24:28他們走近所要去的村子,耶穌好像還要往前走,

24:29他們卻強留他說:「時候晚了,天快黑了,請你同我們住下吧.」耶穌就進去,要同他們住下.

耶穌用甚麼來為門徒解惑呢?聖經.「從摩西和眾先知起,凡經上所指著自己的話都給他們作了解釋」.聖經在基督信仰的重要性,相信無人置疑,差不多所有教會的信仰大網也列出以聖經為信仰依據.

有些人認為,聖經是一本自明的書,只要你肯讀,你就會明白；有些群體甚至以實踐聖經字面意思而感到自己鶴立雞群.早期教會有一位著名教父名叫俄利根,他對信仰非常忠誠,在羅馬帝國的壓逼之下寧死也不低頭.或許少年當時太年輕,也或許受力求按字面實踐的解經方法影響,傳說他讀到馬太福音19章12節「因為有人生來是閹人,也有被人閹的,並有為天國的緣故自閹的；這話誰能領受,就可領受.」於是,他舉刀自閹.

個別人如何運用聖經只是個人的事情,但群體如何運用聖經就足以留下洗不清的污點.最近聽到一位在英國留學的校友分享,歐洲年輕學人之間隱約地有鄙視基督教的氣氛,因為他們認定聖經曾經被用來支持帝國殖民主義.

相信故事中那兩位門徒都讀過聖經,耶穌並不是叫他們再讀一次,而是向他們「解釋」他們所讀的,聖經是需要解釋的.沒有基督徒否認聖經是信仰的重要元素,除此之外,第三項又是甚麼呢?

3.聖禮

24:30坐下來和他們用餐的時候,耶穌拿起餅來,祝福了,擘開,遞給他們.

24:31他們的眼睛開了,這才認出他來.耶穌卻從他們眼前消失了.

門徒邀請這位路上遇到的人一齊用餐,本來門徒是主人,耶穌是客人.突然間,耶穌變成了主人,他「拿起餅來,祝福了,擘開,遞給他們」,如此特別的場面,不正是耶穌和門徒共進逾越節晚餐的情況嗎?看到耶穌這個舉動,門徒的眼睛開了,終於認出耶穌來.

以後,「拿起餅來,祝福了,擘開,遞給他們」這一系列舉動是甚麼呢?它是教會聖餐的重要部分.更重要的是,在這個禮儀中,主客對換了:當我們思考討論信仰和閱讀聖經時,我們常常認為自己是主,而被討論或被閱讀的是客.但在聖禮時,我們變成了客人,耶穌成為了主人.聖禮揭示了信仰中的性質,就是上主是賜予者,我們是接受者.如果信仰只是人的思考和解釋,那就不是基督教信仰了.耶穌和以馬仵斯路上的門徒這個片段啟發我們,在基督教信仰中,上主是賜予的主,我們是接受者.聖禮這個字是奧秘的意思,最偉大的奧秘是耶穌基督道成肉身,是本來不可見的恩典的可見形式,這是上主對世人的賜予.

聖禮是信仰的第三項重要元素,第四項又是甚麼呢?

4.感動/情感

討論,聖經,聖禮,似乎都是外在的元素,但信仰不能沒有信仰者的回應,信仰者的回應也是構成信仰的重要元素之一.

24:32他們彼此說:「在路上他和我們說話,給我們講解聖經的時候,我們的心在我們裡面豈不是火熱的嗎?」

在這裡,我們看到門徒的情感反應,他們感到自己「心裡火熱」.

有人指出,人類的行動不少是被情感推動的,也往往是回應情感的需要.理由呢?往往是用來作為行動的解釋.

有同學想跟女朋友分手,來找我商量用甚麼理由,他在網上找的有以下幾個:

我覺得我們個性不合

我不想傷害你,我是壞人

我想先認真讀書

我配不上你,我不夠好

大家覺得哪個理由最好呢?我問他為甚麼要分手,他說:「我不喜歡她了.」

我們不要低估情感的力量,對信仰亦然.試想想人們成為基督徒的過程是怎樣的?思考討論,閱讀聖經,參加聚會等都有幫助,但總離不開當事人情感的層面.例如有人被審判的懼怕推動,有人被自己的罪疚感推動,有人被上主的愛所感動,有人被天堂的美好吸引...丹麥神學家和哲學家齊克果認為宗教是充滿激情的委身,如此的委身,其實是一種對客觀上的不確定性的冒險,他稱之為「信心的跳躍」(leap of faith).

最近與學生們籌備夏令營,他們把夏令營的主題名為「感動／麻木」,因為他們意識到感動在信仰的重要性,同時也意識到日常生活令他們麻木.這個主題令我想起一首他們不認識的古舊詩歌,其中一段如此唱:

願那復興我的靈火,重新再來一次復興我,

起初的信心何處失落,我願悔改重新再得著.

在明媚的春光裡,在幸福的生活中,

我們陶醉,我們歡樂,早把主恩典忘記.

在以馬忤斯路上的門徒的「心裡火熱」,是信仰第四項重要元素.離開以馬忤斯的路,我們看見第五項:

5.群體

24:33於是他們立刻起身,回耶路撒冷去,看見十一個使徒和與他們正在一起的人聚集在一處,

24:34說:「主果然復活了,已經顯現給西門看了.」

24:35於是,兩個人把路上所遇到,和耶穌擘餅的時候怎麼被他們認出來的事,都述說了一遍.

認出耶穌,感到心裡火熱,那兩位門徒的馬上折返找其他門徒.為甚麼呢?就是急不及待與大家分享剛才與耶穌的相遇.這片段讓我們想到,群體也是信仰中重要的元素.關於群體於信仰的意義的探討非常豐富,有一間教會的異象宣言如此精簡地說:

神透過祂的話,呼召一個人,改變一個人；唯有連結於神的話,和謙卑禱告,才能成為主的門徒；勇於面對改變,預備自己領受召命,才能回應不斷變化的世界；信仰群體能協助這信念得到印證,守望,互勉和鼓勵,彼此連結於同一價值,成為一個整體,對世界展示出福音更大的力量.

(基督教中華宣道會恩友堂:「我們的異象」)

人們常說「愛」是基督教的重要詞彙,耶穌基督也吩咐門徒彼此相愛,試想想,如果不在群體中,又如何實踐彼此相愛呢?

三.總結

總括而言,以馬忤斯路上的片段啟發我們思考基督教信仰的五項元素:一,尋問；二,聖經；三,聖禮；四,情感；五,群體.

或許有人立即想到,改教運動不是提出「唯獨」嗎?「唯獨」一樣,清楚簡單很多呢!沒錯,改教運動的確提出過「唯獨」,但卻不是一樣,其實也是一個組合:分別有「唯獨恩典」,「唯獨信心」,「唯獨聖經」,「唯獨基督」和「唯獨上帝的榮耀」.

個別教會用來建立信仰的根基,也是一個組合:例如聖公會的三重法碼是聖經,傳統,理性；衛斯理約翰的四大神學支柱是聖經,傳統,理性,經驗.這些種種,與以馬忤斯路上的故事給我們的啟發指向同一個方向,基督教信仰的根基是一個多元素的組合,願我們在各種元素中滋養自己,建立生命.



\newpage

\section{同羊唔同命}
\label{sec:644}
\textbf{林豪恩先生 - 同羊唔同命}
\newline
\newline
\hyperref[sec:643]{< < < PREV SERMON < < <}
~
\hyperref[sec:toc]{[返目錄]}
~
\hyperref[sec:645]{> > > NEXT SERMON > > >}


講題:同羊唔同命

講員:林豪恩先生

日期:2023年4月30日

一.引言

剛才大家聆聽經課的時候,有沒有留意幾段經文都提及某種動物.是甚麼動物呢?沒錯,是羊.有一首詩歌叫「我是主的羊」:

主領我到青草地,安歇在溪水旁；

黃昏時主與我一路同行,

牧場上主的羊都得飽足心快暢,我是主的羊.

人們以為歌詞中有三種動物:鵝,豬,羊,其實只有羊.大家都知道羊是用來借喻信徒,以羊與牧羊人的關係來比喻信徒和上主的關係.雖然同是主的羊,但每隻羊都可能有不同的際遇.有一次,有學生來問我:「他積極事奉,我也積極事奉,為甚麼他有拖拍,而我沒有?」或許因為我們看到的羊都好似一樣,於是以為我們這些被比喻為主的羊的信徒也是一個樣子,有著相同的際遇.今天的幾段經課讓我們發現有不同的羊.

二.內容

1.青草地溪水旁的羊

最典型的羊出現在詩篇23篇,差不多令人以為「主的羊」就是這樣的:

23:1耶和華是我的牧者,我必不致缺乏.

23:2他使我躺臥在青草地上,領我在可安歇的水邊.

23:3他使我的靈魂甦醒,為自己的名引導我走義路.

我們可以稱這為「青草地溪水旁的羊」.「青草地溪水旁的羊」給人們的印象是生活得無憂無慮,悠閒自得,令人想到福音單張上常說的話:「相信主耶穌能使我們擁有新的生命,有好的生命,有健康的身體,快樂的心靈,深處的平安以及和諧的人際關係,進而又能幫助別人擁有他所擁有的,那時賺多少錢就是其次了,因為賺錢本來就是為了求取保障,謀求幸福,而這個生命已經帶來了真正的健康,快樂,平安以及美滿的人生.」「青草地溪水旁的羊」和「健康,快樂,平安以及美滿的人生」,這兩個畫面是互相配合的,也是我們所羨慕的.在上主的群羊中,你是不是一隻這樣的羊呢?或者,你有沒有經驗過如此的狀態呢?

2.豐盛生命的羊

詩篇23篇讓我們繪畫出「青草地溪水旁的羊」,約翰福音的比喻又讓我們繪畫怎樣的圖畫呢?

10:9我就是門；凡從我進來的,必然得救,並且出入得草吃.

10:10盜賊來,無非要偷竊,殺害,毀壞；我來了,是要叫羊得生命,並且得的更豐盛.

我們暫且稱這為「豐盛生命的羊」.「豐盛生命的羊」的形象是怎樣的呢?很可能是肥肥白白,精神飽滿,健壯有力的.換了是人呢?福音單張指出一個「豐盛生命」的人有四個可度量的指標:

第一,一個「豐盛生命」的人與神有美好的關係.

第二,一個生命豐盛的人,一定會與人有美好的關係.

第三,一個生命豐盛的人,與世界物質有美好的關係.

第四,一個生命豐盛的人,與自己有美好的關係.

在上主的群羊中,你是不是一隻這樣的羊呢?或者,你有沒有經驗過如此的狀態呢?

3.在地若天的羊

使徒行傳雖然沒有提及羊,但如果我們把經文所描述的繪畫成圖畫,那會是一幅怎樣的圖畫呢?大家一齊來想像一下.

2:42都恆心遵守使徒的教訓,彼此交接,擘餅,祈禱.

2:43眾人都懼怕；使徒又行了許多奇事神蹟.

2:44信的人都在一處,凡物公用,

2:45並且賣了田產,家業,照各人所需用的分給各人.

2:46他們天天同心合意恆切地在殿裡,且在家中擘餅,存著歡喜,誠實的心用飯,

2:47讚美上帝,

人們各盡所能,各取所需,彼此互通有無,愉快共享.這是甚麼狀況呢?是不是類似人類一直渴望的「烏托邦」?彷如人們所能想像的天堂一樣美好呢?我們暫且稱他們為「在地若天的羊」.「在地若天」的生活,當然令人羨慕,經文如此指出其他人也想加入這樣的群體:

得眾民的喜愛.主將得救的人天天加給他們.

「烏托邦」這個詞出於希臘語:οὐ,意為「不」,和τόπος,意為「地方」；這個詞現在托寓那些不可能出現的理想社會和事物.是的,那種理想狀況是不可能出現的,天是天,地是地,若非掌管天地的上主的恩典,在地的人怎能嘗到在天的滋味呢?當時這群信徒,就是在上主的特別恩典之中,經驗到在地若天的狀況,變成在地若天的羊.

在上主的群羊中,你是不是一隻這樣的羊呢?或者,你有沒有經驗過如此的狀況呢?如果你不是上述這些羊之一,也不用質疑或否定他們的經驗和狀況.就正如我不是「高富帥」,但也不用否定這個世界有「高富帥」的存在.不過,也不用因為自己不是「高富帥」而自卑或自憐,這個世界不是只有「高富帥」才有價值的.在「主的羊」之中,也不是只有上述的羊,還有其他在不同狀況之下的羊,都是「主的羊」.其他羊是怎樣的呢?

4.死蔭幽谷中的羊

23:4我雖然行過死蔭的幽谷,也不怕遭害,因為你與我同在；你的杖,你的竿,都安慰我.

「死蔭幽谷」用來比喻人生的患難和危險,令人感到自己隨時會跌進黑暗無助之中.死蔭幽谷對很多人來說是可怕的,所以我們盡量避開,不會自招麻煩.但最叫人頭痛的,就是死蔭幽谷有時突然出現,我們不禁問:「我已忠心跟隨你,為何會有死蔭幽谷迎面而來呢?」

大家有沒有讀過盧雲神父的書呢?他廿五歲成為神父,曾任教於美國聖母院大學,耶魯大學和哈佛大學.一生著作等身,是二十世紀有名的靈修作家.他在《心靈愛語》書中自述:「一九八七年十二月至一九八八年六月期間,是我生命中最艱難的一個時期.那時我極其痛苦,甚至無法肯定能否繼續自己的人生.我的一切都在瓦解:自信,生活工作的力量,被愛的感覺,痊癒的盼望,對上主的信心……無一倖免.我作為書寫靈性生活的作家,被大家視為愛主,為人帶來盼望的人,如今竟死氣沉沉,暗無天日...正當身邊所有人一再向我肯定他們愛我,關心我,欣賞我,甚至欽佩我,我卻只覺得自己一無是處,人見人惡.正當人們溫暖地擁抱我,我卻只覺人生就是無盡的痛苦,毫無值得活下去的理由.正當我終於找到一個家,我卻感到無家可歸.正當我被稱讚深具靈性洞見,我卻感到信心全失.正當別人感激我引導他們更親近上主,我卻感到上主撇棄了我...苦痛令我徹底癱瘓,夜不能寢,不能自己地哭了又哭,任何安慰說話都聽不進去,也再沒有心思為別人分憂.我變得食不下咽,對音樂,藝術,甚至自然之美亦興致索然.世界只剩漆黑.」

在上主的羊群中,你是否這樣的羊呢?或者曾否身同感受呢?有些人的事情與盧雲的不一樣,例如突然面對自己或家人的危疾,或者事業上的挫折,或者親密關係的失去...只要是經歷過死蔭幽谷,都在某程度上體會到盧雲這段自述的狀況.

5.在敵人面前的羊

詩篇23篇不但提及在死蔭幽谷的羊,更有在敵人面前的羊:

23:5在我敵人面前,你為我擺設筵席；你用油膏了我的頭,使我的福杯滿溢.

有人說:「我沒有敵人,也沒有仇恨」,但不等於沒有人仇恨你,也不等於你不需要面對敵人.時代論壇24/4/2023報導,剛過去的四月廿一日,數以千計的斯里蘭卡人於在首都科倫坡組成人鏈,要求政府認真調查二零一九年的復活節恐襲.二零一九年四月廿一日復活節,當地三座教堂──兩座天主教,一座新教,以及三間酒店和另外兩個處所發生自殺式炸彈襲擊,近二百七十人死亡,五百人受傷.科倫坡佩雷拉神父(Fr. Julian Patrick Perera)表示,天主教會不僅是為天主教徒,也是為「全體斯里蘭卡人民」尋求正義,又表示除非找到真正的罪魁禍首,否則斯里蘭卡永遠無法成為一個自由國家.但他亦同時強調,基督徒必須懷有饒恕別人的心,「我們想饒恕,但我們也要知道需要饒恕的是誰……因著上主的恩典,我們已取得了重大進展,並將繼續奮鬥下去」.

6.與強盜為伴的羊

約翰福音10:1「我實實在在地告訴你們,人進羊圈,不從門進去,倒從別處爬進去,那人就是賊,就是強盜.」我們以為羊圈裡面全部都是純潔小綿羊,冷不提防,裡面竟然有賊和強盜.我們可能從沒想過自己是「與強盜為伴的羊」.時代論壇2023年4月4日報導:「播道總會三月底發聲明指,該會牧師ＸＸＸ多次邀請教友參與投資活動,涉及巨大金額,曾承諾歸還款項,可是至今仍未歸還,...該會牧職部及執行委員會議決,褫奪ＸＸＸ的牧師職銜.」這群羊,可能損失了金錢才知道內中有強盜.

7.忍受冤屈苦楚的羊

彼得前書讓我們想像到還有一些「忍受冤屈苦楚的羊」:

2:19倘若人為叫良心對得住上帝,就忍受冤屈的苦楚,這是可喜愛的.

2:20你們若因犯罪受責打,能忍耐,有甚麼可誇的呢?但你們若因行善受苦,能忍耐,這在上帝看是可喜愛的.

彼得前書的讀者在加入這個群體的時候,可能沒有想到竟然要「忍受冤屈的苦楚」.如果我們是他們,我們會不會投訴:「你的福音單張不是這樣說的,你嫌疑違反商品說明條例,現在我們要求退貨,並且要求賠償.」所有向我們傳教的人,都說這是「福音」.是福還是禍,還要看你在甚麼時空之下做基督徒.

「敞開的門」2023報告指出,去年全球有5,621基督徒被殺害,2,110教會遭襲擊,4,542基督徒被逮捕.在阿富汗,新政權上台後,大多數基督徒試圖離開該國,或者躲藏起來.化名加山的秘密信徒說:「塔利班正在挨家挨戶搜查尋索我們,如果他們找到了我們,就會當場把信徒殺死.」

局勢如風雲變色,誰想像得到藍天白雲轉眼間被狂風暴雨取代呢?

龔立人教授最近一篇文章論及「救贖痛苦」:「簡單來說,救贖痛苦指痛苦為他人和自己帶來救贖...每一個人都可以成為基督苦難的分擔者,當一個人願意分擔基督救世苦難的同時,其他人也能藉此獲得救贖.他成為了基督在世上的「共同救贖者」.」在這個概念次下,他想起父母.他說:「我們不需浪漫化父母一切行為.事實上,以「為你好」之名的父母之愛可以是專制和侵犯.留意這虛假的愛之外,我們體會父母的救贖痛苦.每次為朋友的父母之安息禮拜負責講道時,我都有機會聆聽朋友對父母付出的感謝,以致他們可以有今日的人生和生活.這是父母的救贖痛苦.經歷痛苦的父母不一定會看他們的付出是痛苦.為了孩子有好的生活,父母已不覺自己受的苦是一回事了.作為父親的我明白這道理.」另外,他想起那些為社會不公義挺身而出,並可能因此付出痛苦的人.他們或許獲得讚賞,但也可能因此而被攻擊或監禁.

其實,那隻承受痛苦帶來救贖的羊是誰呢?彼得指出:

2:22他並沒有犯罪,口裡也沒有詭詐.

2:23他被罵不還口；受害不說威嚇的話,只將自己交託那按公義審判人的主.

2:24他被掛在木頭上,親身擔當了我們的罪,使我們既然在罪上死,就得以在義上活.因他受的鞭傷,你們便得了醫治.

是的,主的羊有青草地溪水旁,有豐盛生命,甚至有在地若天的預嘗；彼得不是報喜不報憂,他明言主的羊也會有冤屈和苦楚,就是如此,他的勸勉是:

2:21你們蒙召原是為此；因基督也為你們受過苦,給你們留下榜樣,叫你們跟隨他的腳蹤行.

三.總結:仰望牧人的羊

還記得開始時提過那位同學的問題嗎?「他積極事奉,我也積極事奉,為甚麼他有拖拍,而我沒有?」類似的問題可以無限延伸:「我返教會,他也返教會,為甚麼他無憂無慮,而我卻事事不順?」「他信主,我也信主,為甚至他生活優裕,而我卻捉襟見肘?」...最後,經課提醒我們,與其把眼目注視於其他羊,與其他羊比較,不如轉眼仰望那位牧羊人:

詩篇23:6我一生一世必有恩惠慈愛隨著我；我且要住在耶和華的殿中,直到永遠.

彼得前書2:25你們從前好像迷路的羊,如今卻歸到你們靈魂的牧人監督了.

約翰福音

10:10我來了,是要叫羊得生命,並且得的更豐盛.

10:11我是好牧人；好牧人為羊捨命.



\newpage

\section{若不藉著我沒有人能到父那裡去}
\label{sec:645}
\textbf{高國雄牧師 - 若不藉著我,沒有人能到父那裡去}
\newline
\newline
\hyperref[sec:644]{< < < PREV SERMON < < <}
~
\hyperref[sec:toc]{[返目錄]}
~
\hyperref[sec:646]{> > > NEXT SERMON > > >}


約翰福音 14:1-14
\newline
\begin{longtable}{cl}
\hline
\hline
章節 & 經文 (和合本修訂版)\\
\hline
14:1 & \begin{tabularx}{0.7\textwidth}{X} 「你們心裡不要憂愁；你們信神，也當信我。 \end{tabularx} \\ \\ \relax
14:2 & \begin{tabularx}{0.7\textwidth}{X} 在我父的家裡有許多住處；若是沒有，我就早已告訴你們了。我去原是為你們預備地方去。 \end{tabularx} \\ \\ \relax
14:3 & \begin{tabularx}{0.7\textwidth}{X} 我若去為你們預備了地方，就必再來接你們到我那裡去，我在哪裡，叫你們也在哪裡。 \end{tabularx} \\ \\ \relax
14:4 & \begin{tabularx}{0.7\textwidth}{X} 我往哪裡去，你們知道那條路。」 \end{tabularx} \\ \\ \relax
14:5 & \begin{tabularx}{0.7\textwidth}{X} 多馬對他說：「主啊，我們不知道你去哪裡，怎麼能知道那條路呢？」 \end{tabularx} \\ \\ \relax
14:6 & \begin{tabularx}{0.7\textwidth}{X} 耶穌對他說：「我就是道路、真理、生命；若不藉著我，沒有人能到父那裡去。 \end{tabularx} \\ \\ \relax
14:7 & \begin{tabularx}{0.7\textwidth}{X} 既然你們認識了我，也會認識我的父。從今以後，你們就認識他，並且已經看見他了。」 \end{tabularx} \\ \\ \relax
14:8 & \begin{tabularx}{0.7\textwidth}{X} 腓力對他說：「主啊，將父顯給我們看，我們就知足了。」 \end{tabularx} \\ \\ \relax
14:9 & \begin{tabularx}{0.7\textwidth}{X} 耶穌對他說：「腓力，我與你們在一起這麼久了，你還不認識我嗎？看見我的就是看見了父，你怎麼還說『將父顯給我們看』呢？ \end{tabularx} \\ \\ \relax
14:10 & \begin{tabularx}{0.7\textwidth}{X} 我在父裡面，父在我裡面，你不信嗎？我對你們所說的話不是憑著自己說的，而是住在我裡面的父在做他的工作。 \end{tabularx} \\ \\ \relax
14:11 & \begin{tabularx}{0.7\textwidth}{X} 你們要信我，我在父裡面，父在我裡面；即使不信，也要因我所做的工作信我。 \end{tabularx} \\ \\ \relax
14:12 & \begin{tabularx}{0.7\textwidth}{X} 我實實在在地告訴你們，我所做的工作，信我的人也要做，並且要做得比這些更大，因為我到父那裡去。 \end{tabularx} \\ \\ \relax
14:13 & \begin{tabularx}{0.7\textwidth}{X} 你們奉我的名無論求甚麼，我必成全，為了使父因兒子得榮耀。 \end{tabularx} \\ \\ \relax
14:14 & \begin{tabularx}{0.7\textwidth}{X} 你們若奉我的名向我求甚麼，我必成全。」 \end{tabularx} \\ \\
[1ex]
\hline
\hline
\end{longtable}
講題:若不藉著我,沒有人能到父那裡去

經文:約翰福音14:1-14

講員:高國雄牧師

日期:2023年5月7日

1「你們心裡不要憂愁；你們信上帝,也當信我.2在我父的家裡有許多住處；若是沒有,我就早已告訴你們了.我去原是為你們預備地方去.3我若去為你們預備了地方,就必再來接你們到我那裡去,我在哪裡,叫你們也在那裡.4我往哪裡去,你們知道；那條路,你們也知道.」5多馬對他說:「主啊,我們不知道你往哪裡去,怎麼知道那條路呢?」6耶穌說:「我就是道路,真理,生命；若不藉著我,沒有人能到父那裡去.7你們若認識我,也就認識我的父.從今以後,你們認識他,並且已經看見他.」8腓力對他說:「求主將父顯給我們看,我們就知足了.」9耶穌對他說:「腓力,我與你們同在這樣長久,你還不認識我嗎?人看見了我,就是看見了父；你怎麼說『將父顯給我們看』呢?10我在父裡面,父在我裡面,你不信嗎?我對你們所說的話,不是憑著自己說的,乃是住在我裡面的父做他自己的事.11你們當信我,我在父裡面,父在我裡面；即或不信,也當因我所做的事信我.12我實實在在地告訴你們,我所做的事,信我的人也要做,並且要做比這更大的事,因為我往父那裡去.13你們奉我的名無論求甚麼,我必成就,叫父因兒子得榮耀.14你們若奉我的名求甚麼,我必成就.」

在今天復活期的第5主日,我們繼續體會基督復活的豐富,首先讓我們一同進入福音書的經文內容:約翰福音第14章,經文一開始的時候,耶穌直接向他們表示「V1你們心裡不要憂愁」,這裡所用的憂愁內容是指苦惱,不安,在經文一開始的時候耶穌就以這一個情緒去描述當時門徒的狀況,在此先要安撫門徒苦惱,不安,憂愁的情緒,到底發生什麼事呢?

到底當時門徒的心境又是怎樣的呢?

各位弟兄姊妹,我們今天是否也在憂愁,苦惱,不安的狀態當中?所謂何事呢?

在經文中耶穌接著去題到「V1你們信上帝,也當信我.」一個很簡單的邏輯,直線的描述推論,信上帝的人就應該相信耶穌.

在我們生命中所遇見的人,可能也表達他們「信有上帝」,可能他們表達只相信世上有主宰,有一位創造主在冥冥中管理一切的存在,但是他們未必(甚至很多時候他們不會)相信耶穌就是哪一位上帝,為什麼呢?

因為很多時候我們對上帝的形象有著我們自己的理解,期盼,甚至這形象是由我們自己腦袋中幻想而塑造出來的,所以當我們打開聖經,當我們回到教會,當我們開始認真認識歷史中的耶穌,又聽到耶穌走上十字架,為我們釘身,為我們受死埋葬並且復活等…很多時候我們沒有辦法完全理解明白,甚至以我們的邏輯智慧覺得有一點愚蠢,有一點太兒戲,有一點我們不能接受,不能明白的地方,所以很多人會選擇相信地上有主宰,但選擇不相信耶穌就是那位主宰…

在約翰福音當耶穌要安慰門徒不要憂愁,不要苦惱的時候,耶穌直接提出「你們信上帝,也當信我!」這可能也是當時門徒當中最大的疑惑,因為他們見證著耶穌的所作所為,他們有數年的時間與耶穌一同生活,一定見證過耶穌的能力,耶穌說話的權柄等等…

但是與他們心目中天父上帝的形象又或者是他們心目中彌賽亞拯救主的形象還有一點差距.所以他們憂愁,接著耶穌繼續描述,耶穌說他要回去(他要離開他們)回到父那裡,這是有原因的,經文當中第2節很清楚「我去原是為你們預備地方,我若去為了你們預備地方,就必再來接你們到我那裡去.」耶穌在這裡說明他要離開,請緊記這段經文是耶穌在離世前,被捉拿前與門徒相聚的最後分享,他預言自己將要被出賣,將要離開.

我們非常清楚耶穌非常明白門徒心裡憂傷其中一個要點就是懼怕失去,懼怕失去耶穌,懼怕未來的日子當耶穌真的被害離開的時候,他們便失去方法.

他們的憂愁,恐懼,苦惱令他們離開信仰,令他們失去生活/生存的力量.

耶穌在這裡清楚描述他離開是有原因的:「他要為我們預備地方」這描述對我們來說是什麼意義呢?

耶穌繼續在世上不是就能夠繼續行神蹟,能夠繼續醫病趕鬼,行更多奇事難道不是更好嗎?難道不是更能夠幫助更多人嗎?難道不是更有價值更有意義嗎?

很多人的理解,所想像「好的宗教信仰」的價值與意義不就是不斷做好事,不斷行善,不斷幫助人,盡最大的時空努力去做最大的事,性價比最高難道不是應該的嗎?耶穌那個時候按我一般我們的理解大概33歲左右,傳道服務也不過3年左右,如果他在地上多留幾十年甚至最少多留幾個月,難道不是更有價值嗎?

期望宗教導人向善這道理難道有錯嗎?

基督教不是旨在導人向善,而是要導人向上帝!

耶穌選擇要離開,耶穌選擇接受被出賣,耶穌選擇走上十字架,為要為我們預備更好的地方,有關這個「地方」當然有很多我們不能明白的事情,到底是那裡呢?有人會理解為耶穌在天上(天堂)為我們預備/佈置最美的住處,讓我們死後就可以到天堂有美好的居住場所,這個理解當然有它的意義,但是另一個理解方向,就是耶穌要成就救恩為我們預備一個能夠與上帝永遠同在的場所,真正的豐富,真正美好的家鄉,就是與上帝一同生活,一同作王的地方,正如當日創世時那伊甸園的情況一樣.

所以耶穌要去為我們預備地方,那個「去預備」就是他自己走上十字架,祂自己獻上他的生命,他的受死埋葬復活升天為我們帶來重新的力量,讓我們可以透過他潔淨一切的罪以致能夠返回最初上帝創造的本意,我們能夠同主永遠同在,那就是最美的地方.

所以耶穌說「我在那裡,叫你們也在那裡,我往那裡去你們知道那條路.」耶穌非常清晰他的目的是期望我們與祂永遠同在,所以他去預備地方的目的就是要讓我們也有一天能夠完全進入他裡面,與主相遇這也就是福音最核心的要點.

聖經繼續提到門徒多馬對耶穌說:「主啊,我不知道你往那裡怎麼能知道那條路呢?」

這是一個很有趣的提問,耶穌只表達他去預備地方但他沒有清晰講是哪裡,接著耶穌的描述「我在那裡叫你們也在那裡,我往那裡去,你們知道那條路.」

耶穌在這裡很清楚表達,他認為門徒應該完全明白他將要走的那條路是怎樣的路,他期望門徒能夠往同一個目的地,就是耶穌為他們預備的那美地,但是多馬在這裡的疑問其實是很容易理解的,當我們不知道目的地又怎知道該行那條路呢?

在我們日常生活中是常常能夠經歷到的,今天我們只要打開手機,打開地圖,打開導航,第一個動作必然是輸入目的地,我地我們就能夠按照指引清楚到達,當日的問題同樣,如果沒有目的地,沒有地址,沒有目標,我們又如何知道那條道路呢?

這也可能是我們信徒活命的寫照,很多時候我們好像很迷茫,好像未來目標不太清晰,特別是當我們看見身邊很多風浪的時間,我們就會懼怕,想起昔日彼得在船上看見耶穌在海面上行走,彼得也望著耶穌想走過去,原本因為有清晰的目標路徑,他真的能在海面上行走,但是當他看見的是風浪大作的時候,他就懼怕,人就沉下去…

在我們信仰成長的路可能也是一樣,我們望見身旁很大的風浪,我們就會懼怕,當我們懼怕時,我們就會沉下去.

很多時候我們也會因為很懼怕兜圈子,很懼怕走錯路,很想所有事情很清晰才踏出那一步:分享一個經驗,我休息的時候大多會去外邊旅遊,要去一個新地方,打開地圖尋找路線,輸入正確目的地名稱,當然是很習慣；但往往發生當路線不太清晰,又或當那些導航有出現誤差的時候,我會很驚惶,很懼怕,甚至有時會馬上停下,等待重新計劃路線才行動,當然我不是推薦大家以後盲目亂走,但是其實很多時候只要我們繼續走上一小段,導航又再會重新定了位置,就能夠有更新,更清晰,更能夠直達目標的路徑出現在我們面前.

所以耶穌的回應,就是耶穌一句非常精彩並且很多信徒成為金句的名言「V6-7我就是道路,真理,生命,若不藉著我沒有人能到父那裡去.你們若認識了我,也就認識我的父,從今以後,你們認識他,並且已經看見他了.」

這裡很清晰,耶穌就是那道路,目的地可能暫時我們沒有辦法完全清楚細節,這是我們的限制,但是我們很清楚我們要踏上的這條道路就是耶穌基督本身,當我們能夠清楚認識:我們只要透過耶穌這條道路,就能夠連接上帝,我們安然憑信心倚靠踏上這條道路,我們必會到達將要去的目的地,與創造主同在,與創造主重新聯結,正如當我們只要清晰登上一輛合適的巴士,只要安心信靠,我們必然到達那應該要到的目的地,這就是耶穌向門徒的應許.

接著另一位門徒腓力就對耶穌問:「求主將父顯給我們看,我們就知足了.」腓力這個要求好像與前面耶穌一直的描述沒有多大的關係,為何腓力在這裡突然問這一個問題呢?

很明顯第一:他現在不是非常滿足(仍然未知足)

第二:他表示他期望要看天父所顯出的榮耀,因為仍然未看到.

這是因為剛剛耶穌說:「你們若認識我(耶穌),也就認識我父,從今以後你們認識,並且已經看見他了.」耶穌清晰描述在耶穌身上他們應該就夠看見上帝.

腓力的描述好像斬釘截鐵地表達出:「我看不見!」

真的是很可悲的事情,跟隨耶穌的門徒,與耶穌同住多年,竟然在此時(耶穌將要離開)仍然沒有辦法透過耶穌看見上帝…

耶穌如何回答他呢?第9節耶穌對他說:「腓力,我與你們在一起這麼久了,你還不認識我嗎?」耶穌在這裡心痛地反問為什麼到現在還未能真正認識耶穌呢?

是認識得不夠深?是耶穌所行D 神蹟奇事仍然未能打動他們的心?還是正如一開始我所提到他們對上帝所差來彌賽亞的期望,與他們所經歷的有分別,有落差?

相信他們更從沒有想過耶穌行這些神蹟奇事,耶穌所施展的大能就是表明他與天父上帝的關係,他們沒有因為看見耶穌就看見上帝,所以耶穌接著說:「人看見了我,就是看見了父；你怎麼說『將父顯給我們看』呢?10我在父裡面,父在我裡面,你不信嗎?我對你們所說的話,不是憑著自己說的,乃是住在我裡面的父做他自己的事.11你們當信我,我在父裡面,父在我裡面；即或不信,也當因我所做的事信我.」

這裡耶穌明明白白地直接描述他與父的關係,「父在他裡面,他在父裡面」,當耶穌反問腓力為何這麼多年仍然沒有清楚耶穌的角色,身份,當腓力期望「父的顯現」.

這體會很多時候與我們經歷實很相近,我們所見的實在太片面,有很多人說如果耶穌就像我身邊顯然,耶穌在我身上行神蹟,我就會好好的跟隨主了…

但事實上,在我們身邊每天仍然有很多的神蹟出現,在我們身上不斷地也有很多上帝所造美妙的事情,但是我們很多人也仍然未能夠完全跟從,倚靠.

耶穌在這裡說明就算在理解上沒辦法完全明白,但是從耶穌所行的所作的也應該要信,12節:「我實實在在地告訴你們,我所做的事,信我的人也要做,並且要做比這更大的事,因為我往父那裡去.13你們奉我的名無論求甚麼,我必成就,叫父因兒子得榮耀.14你們若奉我的名求甚麼,我必成就.」

耶穌很有意義的延伸,因他從父那裡得著能力,他在父裡面,父在他裡面,所以他有能力作大事,而今天耶穌要在門徒裡面,門徒也要在耶穌裡面,所以門徒要承繼耶穌這能力,甚至要做更大的事.

回到這段經文開首,當耶穌提及他們心裡憂愁,苦惱,不安的原因,在前一段經文第13章末段,36-38節耶穌與西門彼得的一段對話中,「西門彼得問耶穌說:主往去哪裡?耶穌回答說:我所去D 地方,你現在不能去,以後卻要跟我去；彼得說,主呀,我為什麼現在不能跟你去?我願意為你捨命,耶穌回答,你願意為我捨命嗎?我實實在在地告訴你,雞叫以先,你要三次不認我.」原來約翰福音的記載14章耶穌自述是道路,真理,生命這段經文之前,就是直接向著彼得預言,他會3次不認主,因為彼得在問耶穌要去哪裡後,他揚言必會跟從到底,就算捨命也願意,我相信當時彼得這番話是真心的,到最後他三次不認主,不完全是因為貪生怕死,而是因為當時他們對耶穌有獨特的期望,正如早前與大家分享過他們期望耶穌成為他們當時的拯救者,當時他們被政權欺壓下他們需要人幫他們從捆綁中釋放他們,所以他們問耶穌你要走的是哪條路時,這個問題相信不是問耶穌要走的是哪個方向(地域上),而是問耶穌將要如何走他人生的路,直接的表達,相信就是要問耶穌你的政治取向,走哪一條路,就是問你將要作甚麼事呢…

接著耶穌以這一段的回應門徒:14:1-14節整段的回答,說明一件很重要的事,耶穌所緊張的,是門徒未來將要承斷,而作比耶穌所作更大的事.

雖然在肉身來說,耶穌將要離世,肉身上我們無辦法再見到耶穌,但是我們很清楚耶穌這裡所說的就是要期望與門徒同在(在裡面),當耶穌說若不是藉著我沒有人能到父那裡去,是因為它與父原為一,他要與我們合一.

約翰福音往後的3章經文中不斷提到耶穌離開後,聖靈要來到,耶穌是真葡萄樹與我們要連結,更同時要提醒門徒要合一,從整段的脈絡,我們再一次去理解耶穌是道路,真理,生命,就不應該單理解為耶穌(耶穌的救贖)只是一個工具,單單明白我們能透過這工具就可以到父那裡去是不足夠的.

耶穌這段離別之話的重點是耶穌要與我們同在,我們的主就是那位願意並喜愛與人同在的上帝!

雖然感受上,生活形態上將會不一樣,正如門徒曾經歷的地上道成肉身的耶穌,將要離開,當聖靈降臨在他們身上,門徒就再一次經歷上帝的同在,正如耶穌回到天上與父同在一樣的同在.

弟兄姊姊,今天我們生命中經歷的可能有不一樣,信仰路成長的軌跡也有所不同,甚至出現斷斷續續的情況,但是我們要清楚知道主願意/期望/喜愛/要與我們同在.

所以當我們在每一次崇拜,當主席與大家問安「願主與你們同在」,當我們回應「也與你同在」的時候,我們是否真實體會上主透過聖靈臨在我們當中呢?

當我們在每一次領受聖餐,一同吃餅喝杯的時候,我們有沒有體會到耶穌基督,我們的主上帝就再一次進入我們生命當中?我們再一次與主一同連結,體會這同在的美好?

弟兄姊妹,信仰在每一個人心目中可能有不同的理解與期盼,不同的狀態可能也是我們無法想像的,但是我們要清晰,只有一條道路,只有一位從上主以來的耶穌能夠成為我們進入永生的唯一途徑,你們是否願意一同在這路上與上帝連結,尋找到生命的的磐石,看得見上帝的榮耀呢?



\newpage

\section{為那不可知的神作見證}
\label{sec:646}
\textbf{李均熊牧師 - 為那不可知的神作見證}
\newline
\newline
\hyperref[sec:645]{< < < PREV SERMON < < <}
~
\hyperref[sec:toc]{[返目錄]}
~
\hyperref[sec:647]{> > > NEXT SERMON > > >}


使徒行傳 17:22-31
\newline
\begin{longtable}{cl}
\hline
\hline
章節 & 經文 (和合本修訂版)\\
\hline
17:22 & \begin{tabularx}{0.7\textwidth}{X} 保羅站在亞略巴古當中，說：「諸位雅典人！我看你們凡事很敬畏鬼神。 \end{tabularx} \\ \\ \relax
17:23 & \begin{tabularx}{0.7\textwidth}{X} 我到處走走的時候，仔細觀察你們所敬拜的，發現一座壇，上面寫著『獻給未識之神明』。你們所不認識而敬拜的，我現在向你們宣告： \end{tabularx} \\ \\ \relax
17:24 & \begin{tabularx}{0.7\textwidth}{X} 他是創造宇宙和其中萬物的神；他既是天地的主，就不住在人手所造的殿宇裡， \end{tabularx} \\ \\ \relax
17:25 & \begin{tabularx}{0.7\textwidth}{X} 也不用人手去服侍，好像缺少甚麼似的；自己倒將生命、氣息、萬物賜給萬人。 \end{tabularx} \\ \\ \relax
17:26 & \begin{tabularx}{0.7\textwidth}{X} 他從一人造出萬族，居住在全地面上，並且預先定準他們的年限和所住的疆界， \end{tabularx} \\ \\ \relax
17:27 & \begin{tabularx}{0.7\textwidth}{X} 為要使他們尋求神，或者可以揣摩而找到他，其實他離我們各人不遠。 \end{tabularx} \\ \\ \relax
17:28 & \begin{tabularx}{0.7\textwidth}{X} 我們生活、行動、存在都在於他。就如你們的詩人也有人說：『我們也是他所生的。』 \end{tabularx} \\ \\ \relax
17:29 & \begin{tabularx}{0.7\textwidth}{X} 既然我們是神所生的，就不應該以為神的神性像人用手藝和心思所雕刻的金、銀、石像一般。 \end{tabularx} \\ \\ \relax
17:30 & \begin{tabularx}{0.7\textwidth}{X} 世人蒙昧無知的時候，神並不追究，如今卻吩咐各處的人都要悔改。 \end{tabularx} \\ \\ \relax
17:31 & \begin{tabularx}{0.7\textwidth}{X} 因為他已經定了日子，要藉著他所設立的人按公義審判天下，並且使他從死人中復活，給萬人作可信的憑據。」 \end{tabularx} \\ \\
[1ex]
\hline
\hline
\end{longtable}
詩篇 66:8-20
\newline
\begin{longtable}{cl}
\hline
\hline
章節 & 經文 (和合本修訂版)\\
\hline
66:8 & \begin{tabularx}{0.7\textwidth}{X} 萬民哪，你們當稱頌我們的神，使人得聽讚美他的聲音。 \end{tabularx} \\ \\ \relax
66:9 & \begin{tabularx}{0.7\textwidth}{X} 他使我們的性命存活，不叫我們的腳搖動。 \end{tabularx} \\ \\ \relax
66:10 & \begin{tabularx}{0.7\textwidth}{X} 神啊，你曾考驗我們，你熬煉我們，如煉銀子一樣。 \end{tabularx} \\ \\ \relax
66:11 & \begin{tabularx}{0.7\textwidth}{X} 你使我們進入羅網，把重擔放在我們身上。 \end{tabularx} \\ \\ \relax
66:12 & \begin{tabularx}{0.7\textwidth}{X} 你使人坐車軋我們的頭；我們經過水火，你卻使我們到豐富之地。 \end{tabularx} \\ \\ \relax
66:13 & \begin{tabularx}{0.7\textwidth}{X} 我要帶著燔祭進你的殿，向你還我的願， \end{tabularx} \\ \\ \relax
66:14 & \begin{tabularx}{0.7\textwidth}{X} 就是在急難時我嘴唇所發的、口中所許的。 \end{tabularx} \\ \\ \relax
66:15 & \begin{tabularx}{0.7\textwidth}{X} 我要將肥牛的燔祭和公羊的香祭獻給你，又要把公牛和公山羊獻上。（細拉） \end{tabularx} \\ \\ \relax
66:16 & \begin{tabularx}{0.7\textwidth}{X} 敬畏神的人哪，你們都來聽！我要述說他為我所做的事。 \end{tabularx} \\ \\ \relax
66:17 & \begin{tabularx}{0.7\textwidth}{X} 我曾用口求告他，我的舌頭也稱他為高。 \end{tabularx} \\ \\ \relax
66:18 & \begin{tabularx}{0.7\textwidth}{X} 我若心裡注重罪孽，主必不聽。 \end{tabularx} \\ \\ \relax
66:19 & \begin{tabularx}{0.7\textwidth}{X} 但神實在聽見了，他留心聽了我禱告的聲音。 \end{tabularx} \\ \\ \relax
66:20 & \begin{tabularx}{0.7\textwidth}{X} 神是應當稱頌的！他沒有推卻我的禱告，也沒有使他的慈愛離開我。 \end{tabularx} \\ \\
[1ex]
\hline
\hline
\end{longtable}
彼得前書 3:13-22
\newline
\begin{longtable}{cl}
\hline
\hline
章節 & 經文 (和合本修訂版)\\
\hline
3:13 & \begin{tabularx}{0.7\textwidth}{X} 你們若熱心行善，有誰會害你們呢？ \end{tabularx} \\ \\ \relax
3:14 & \begin{tabularx}{0.7\textwidth}{X} 即使你們為義受苦，也是有福的。不要怕人的威嚇，也不要驚慌； \end{tabularx} \\ \\ \relax
3:15 & \begin{tabularx}{0.7\textwidth}{X} 只要心裡奉主基督為聖，尊他為主。有人問你們心中盼望的理由，要隨時準備答覆； \end{tabularx} \\ \\ \relax
3:16 & \begin{tabularx}{0.7\textwidth}{X} 不過，要以溫柔、敬畏的心回答。要存無虧的良心，使你們在何事上被毀謗，就在何事上使那些凌辱你們在基督裡有好品行的人自覺羞愧。 \end{tabularx} \\ \\ \relax
3:17 & \begin{tabularx}{0.7\textwidth}{X} 神的旨意若是要你們因行善受苦，這總比因行惡受苦好。 \end{tabularx} \\ \\ \relax
3:18 & \begin{tabularx}{0.7\textwidth}{X} 因為基督也曾一次為罪受苦 ，就是義的代替不義的，為要引領你們到神面前。在肉體裡，他被治死；但在靈裡，他復活了。 \end{tabularx} \\ \\ \relax
3:19 & \begin{tabularx}{0.7\textwidth}{X} 他藉這靈也曾去向那些在監獄裡的靈傳道， \end{tabularx} \\ \\ \relax
3:20 & \begin{tabularx}{0.7\textwidth}{X} 就是那些從前在挪亞預備方舟、神容忍等待的時候不信從的人。當時進入方舟，藉著水得救的不多，只有八個人。 \end{tabularx} \\ \\ \relax
3:21 & \begin{tabularx}{0.7\textwidth}{X} 這水所預表的洗禮，現在藉著耶穌基督的復活拯救你們，不是除掉肉體的污穢，而是向神懇求有無虧的良心。 \end{tabularx} \\ \\ \relax
3:22 & \begin{tabularx}{0.7\textwidth}{X} 耶穌已經到天上去，在神的右邊，眾天使、有權柄的、有權能的都服從了他。 \end{tabularx} \\ \\
[1ex]
\hline
\hline
\end{longtable}
約翰福音 14:15-21
\newline
\begin{longtable}{cl}
\hline
\hline
章節 & 經文 (和合本修訂版)\\
\hline
14:15 & \begin{tabularx}{0.7\textwidth}{X} 「你們若愛我，就會遵守我的命令。 \end{tabularx} \\ \\ \relax
14:16 & \begin{tabularx}{0.7\textwidth}{X} 我要求父，父就賜給你們另外一位保惠師，使他永遠與你們同在。 \end{tabularx} \\ \\ \relax
14:17 & \begin{tabularx}{0.7\textwidth}{X} 他就是真理的靈，是世人不能接受的。因為他們既看不見他，也不認識他；你們卻認識他，因他常與你們同在，也要在你們裡面。 \end{tabularx} \\ \\ \relax
14:18 & \begin{tabularx}{0.7\textwidth}{X} 我不會撇下你們為孤兒，我必到你們這裡來。 \end{tabularx} \\ \\ \relax
14:19 & \begin{tabularx}{0.7\textwidth}{X} 再過不久，世人不再看見我，你們卻會看見我，因為我活著，你們也要活著。 \end{tabularx} \\ \\ \relax
14:20 & \begin{tabularx}{0.7\textwidth}{X} 到那日，你們就會知道我在父裡面，你們在我裡面，我也在你們裡面。 \end{tabularx} \\ \\ \relax
14:21 & \begin{tabularx}{0.7\textwidth}{X} 有了我的命令而又遵守的人，就是愛我的；愛我的人，我父要愛他，我也要愛他，並且要親自向他顯現。」 \end{tabularx} \\ \\
[1ex]
\hline
\hline
\end{longtable}
講題:為那不可知的神作見證

講員:李均熊牧師

經文:使徒行傳17:22-31;詩篇66:8-20;彼得前書3:13-22;約翰福音14:15-21

日期:2023年5月14日

起:今天福音經課帶出的主題是甚麼?

五旬節快到了,也是傳統教會慶祝聖靈降臨的日子.初看這星期的福音經課,主耶穌就預告聖靈的降臨,那今天的宣講應順理成章講聖靈吧?但你若仔細看看其他的經課,你就會察覺無論是使徒行傳17章,詩篇66的選段,和書信經課的彼得前書3章的選段,都不約而同的論到見證上主的主題.使徒行傳記載保羅要向凡事敬畏鬼神的雅典人論說那位他們的未識之神明就是創造和維繫萬有的上主；詩篇中詩人經過上主的考驗,在羅網,重擔,水火之後,回到上主聖殿中獻祭還願時,就召喚「敬畏神的人哪,你們都來聽!」,因為他「要述說他為我所做的事.」；在彼得前書中,經文鼓勵我們在受苦中仍要行善,不單「有人問你們心中盼望的理由,要隨時準備答覆」,而且「要以溫柔,敬畏的心回答」；更重要的,就是用我們的好行為去使毀謗,凌辱我們的人「自覺羞愧」!這些內容都是我們教會傳統所謂的做見證.向未信主的人傳福音是做見證,在教會分享自己的困難並上主如何使我們經過是感恩見證,甚至在批評我們信仰的人面前做好見證(但在教會我們通常是反面的指責人沒有用好行為「做好見證」).

那麼,今天的福音經課和這些其他經課,主題就不是一致了?

承:見證——使人知道祂

但當我仔細去看今日的福音經課,我最後的結論是,四段主日經課其實都是教導我們如何去見證復活的主,而且福音經課是見證的基礎原理,其他三段則是不同處境的發揮.為甚麼這樣說呢?

大家不妨細讀一次今日的福音經課,短短6節的經文,其實包含很多主題:一頭一尾關於愛我就遵守命令的主題,接著論到保惠師,就是真理的靈的主題,間中穿插的又有「常在」的主題.但其中一個交織著的主題,就是「看見／看不見」,「認識／不認識」這些與「知道」有關的題旨.

以住我也指出這其實是約翰福音的特色,就是每個約翰福音的段落,經文都糾結多重的主題,好像一條粗壯鋼纜(cable),是由無數鋼索 (wire) 組成一樣.

若果其他經課選段都是和見證上主有關,那約翰福音的這一段福音經課放在「見證上主」的脈絡內,便可以這樣理解:因我們和主耶穌相互的愛,父就賜給我們另一個保惠師,又叫真理的理,使他永遠與我們同在,使基督可以常在我們裡面.但「再過不久,世人不再看見我,你們卻會看見我」,而且「他們既不看見他,也不認識他」,而我們卻認識他,因他常與我們同在,也在我們裡面.配合其他經課選段,這一天的經課傳統讓我們去思想,你可以並要見證上主,因為你看得見他,又認識他,惟有這樣,你才可以讓看不見又不認識他的世人明白他是怎樣的一方愛世人的主.見證是源與我們看見上主,不單看見,更是認識.

轉:見證不是宣傳

根據約翰福音的講法,其實見證不是單分享自己的經歷；而是透過自己的經歷,去讓別人認識耶穌／保惠師／上主是一個怎樣的神.所以我們必須要先從自身的經驗中看到上主的實在,並明白祂是怎樣的神,我們才能在不認識祂的人面前見證祂.很多時我們只分享自己(主觀的)經驗,甚至強調一些「得勝」的經驗,卻沒有令人透過我們經歷看見神,並更認識神是一位怎樣的神,或神如何在我們生命中使我們得見祂的實在,我們其實還是未作出見證,只是個人經驗分享而矣,就好像買藥的人,只求宣傳藥的效能,卻不是使人明白那是怎樣的藥,為何有效.我們在廣告洗腦的網絡時代,只知道甚麼是宣傳.但宣傳不是見證,或者不是今日福音經課所導引出來見證的原則和基礎.若我們只宣傳上主「得唔得」,也只會吸引只要靈驗經歷的人,卻不是見證賜恩主是誰,因認識他而接受他而成為與上主建立關係的兒女.

合:見神才可以見證祂

我們以為「得勝」的經歷才可以拿來見證神；又或當我們仍在逆境掙扎中,我們會覺得生命沒有甚麼「好」見證能榮耀神.有些時候我們甚至感到我們這般主觀的經驗,未必可以見證上主的實在.這些都是因為我們將見證的焦點放錯在「宣傳」神的靈驗上,而非讓人透過我們的經驗看見神的實在,明白和認識這位世人所不認識甚至誤解了的神.

在我們無論或順或逆的經驗中,若我們「見」到世人不能見的神,「認識」世人不能認識的主,那透過我們分享,別人能得見基督的實在,認識這一位神(的某些屬性),這才是真實的見證.見證的焦點是使人得見神,認識神,明白神,並因此而接受祂,愛祂和敬拜祂,而不是我們的經歷的好與壞,又或這個神是否一個可以替我消災解難的精靈(認得米利巴,就是瑪撒嗎?).



\newpage

\section{奉誰的旨}
\label{sec:647}
\textbf{鄧瑞強博士 - 奉誰的旨?}
\newline
\newline
\hyperref[sec:646]{< < < PREV SERMON < < <}
~
\hyperref[sec:toc]{[返目錄]}
~
\hyperref[sec:648]{> > > NEXT SERMON > > >}


講題:奉誰的旨?

經文:以斯拉記1:1-4,7,4:1-7,18-24,5:1-6,6:1-15,19-22;7:21-23

講員:鄧瑞強博士

日期:2023年5月21日

在猶太人的歷史裡,國破家亡,京城敗落,聖殿被毀,被擄去巴比倫,是無法忘懷的事情.在巴比倫的河邊,他們為亡國前的腐敗生活懊悔.在巴比倫的土地上,先知以西結在異象中看到,神的寶座從耶路撒冷移到巴比倫.這是一個重大的信念,神不單在猶大地掌權,耶和華也在看似沒有祂位置的領域上掌權,耶和華在敵對祂子民的空間上掌權.太陽在黑雲之上依然燦爛,光明在黑暗之上仍舊照耀.這一信念,讓猶大人懷著盼望,在異鄉裡安然生活.

時移勢異,乾坤轉動.有點出人意外,巴比倫竟然急速敗亡.波斯帝國踏入歷史舞台,一統江山.今日講道要講的《以斯拉記》,就是要交代在這新的歷史契機下,猶太人的信仰反省.

(1)《以斯拉記》的開場白:神的旨意在王的旨意之上

波斯王塞魯士一踏進舞台,便施行新政,實施懷柔政策,讓被擄的各族人民重回故土.在大英博物館裡,有一藏品叫「Cyrus Cylinder」(塞魯士圓筒),一個手掌般大的石圓筒,上面寫著塞魯士王的諭令:「我招聚所有這些人民,讓他們回歸故土……

願我送返各神廟裡的所有神明……

祝福我長壽,述說我的好行為……

我令全地和平生活…」(…I collected together all of their people and returned them to their settlements…

May all the gods that I returned to their sanctuaries…

ask for a long life for me, and mention my good deeds…

I have enabled all the lands to live in peace.)(\url{https://shorturl.at/bhnAY})

1:1波斯王塞魯士元年,耶和華為要應驗藉耶利米口所說的話,就激動波斯王塞魯士的心,使他下詔通告全國說:

1:2「波斯王塞魯士如此說:『耶和華天上的上帝已將天下萬國賜給我,又囑咐我在猶大的耶路撒冷為他建造殿宇.

1:3在你們中間凡作他子民的,可以上猶大的耶路撒冷,在耶路撒冷重建耶和華－以色列上帝的殿(只有他是上帝).願上帝與這人同在.

1:4凡剩下的人,無論寄居何處,那地的人要用金銀,財物,牲畜幫助他,另外也要為耶路撒冷上帝的殿甘心獻上禮物.』」……

1:7塞魯士王也將耶和華殿的器皿拿出來,這器皿是尼布甲尼撒從耶路撒冷掠來,放在自己神之廟中的.

這裡,有幾點可以留意:

一,塞魯士王下詔,讓猶太人回歸,並在耶路撒冷建殿.建殿是合法的,在塞魯士王眼中,這不單不會威脅波斯帝國的安定,反讓帝國得到福祉.

二,他將巴比倫王尼布甲尼撒從耶路撒冷掠來的聖殿之物歸還,對猶太人而言,這是恢復古老的聖殿；從信仰而言,這就是天地的復修,是救恩的新天地的開始.

三,留意,塞魯士本身及帝國並沒有給予實質的物資,在神學上,這是關鍵的.在近東的建殿文化裡,誰出資,誰就是恩主,誰的神明就住在這殿中.明顯,波斯的神明與耶和華的殿無關.

四,最富有神學深意的,是猶大神學家領悟到,在塞魯士王的詔書背後,是神的「詔書」,神的旨意.明顯,耶和華是萬王之上的王,祂主宰王的心.能看到神的手主宰著歷史的演變,對信徒而言,是最大的鼓舞.

五,整段經文的重點是「建殿」.建殿,是肯定神的主權,並在祂的主權下下拜.建殿,也是沒有了身分的猶太人重建「神的子民」這身分的關鍵行動.人一定有所「拜」,拜什麼,便是他的身分.

(2)阻滯

4:1猶大和便雅憫的敵人聽說被擄歸回的人為耶和華－以色列的上帝建造殿宇,

4:2就去見所羅巴伯和以色列的族長,對他們說:「請容我們與你們一同建造；因為我們尋求你們的上帝,與你們一樣.自從亞述王以撒哈頓帶我們上這地以來,我們常祭祀上帝.」

4:3但所羅巴伯,耶書亞,和其餘以色列的族長對他們說:「我們建造上帝的殿與你們無干,我們自己為耶和華－以色列的上帝協力建造,是照波斯王塞魯士所吩咐的.」

4:4那地的民,就在猶大人建造的時候,使他們的手發軟,擾亂他們；

4:5從波斯王塞魯士年間,直到波斯王大流士登基的時候,賄買謀士,要敗壞他們的謀算.

這裡,有幾點可以留意:

一,「猶大和便雅憫」,這是南國猶大的支派.這裡,以文學手法表達,北國以色列的支派無份建殿,因為他們不在耶路撒冷的聖殿敬拜神.

二,想建殿,也不是容易的,他們面對「敵人」.這些「敵人」表明上不似敵人,他們願意一起建殿.但留意,這些人不是猶太人,是亞述王帶到那地方的,他們很可能是混種的人,有著混種的信仰,敬拜著混種的神明.

三,所羅巴伯等領導,不容他們建殿,因為這時是猶太人重建身分的關鍵時刻.一摻雜別的神明,慢慢便會被別的信仰「同化」,如此,便會失去信仰的純正,也就失去神子民的純正身分,也就忘記活在世上的天命.

四,《以斯拉記》這卷書,分上下兩部分.上半部,重點在重建聖殿,講求的是保持聖殿的「分別為聖」；下半部,重點在重建人民,講求的是保持人民的「分別為聖」.「建殿」其實即是「建人」.下半部的經文,來得很決絕,任何娶了外邦女子的猶太人,要與妻子離婚,也要離開她們生的子女.這是在關鍵時刻,以最極端的方式,保持身分的純正.

五,那些「敵人」被拒絕,老羞成怒,真的成了敵人,擾亂他們的工程.

六,4:5提及「大流士」王,按《以斯拉記》6:15的記述,聖殿在「大流士」王第6年建成.

(3)上奏

4:6在亞哈隨魯才登基的時候,上本控告猶大和耶路撒冷的居民.

4:7亞達薛西年間,比施蘭,米特利達,他別,和他們的同黨上本奏告波斯王亞達薛西.本章是用亞蘭文字,亞蘭方言.

亞哈隨魯王和亞達薛西王,是大流士王之後的王,按事實,在他們執政時期,聖殿早已建成.故此,隨後的文字的重點,在於表達一種信仰洞見.書寫的文字刻意轉用「亞蘭文」,明顯,吸引讀者留意人間的王的諭旨.誰下終極的「旨」?

(4)亞達薛西王的諭旨:反對建殿

4:18〔王的覆諭:〕「你們所上的本,已經明讀在我面前.

4:19我已命(tǝēm,「旨」)人考查,得知此城古來果然背叛列王,其中常有反叛悖逆的事.

4:20從前耶路撒冷也有大君王統管河西全地,人就給他們進貢,交課,納稅.

4:21現在你們要出告示命(tǝēm,「旨」)這些人停工,使這城不得建造,等我降旨(tǝēm).

4:22你們當謹慎,不可遲延,為何容害加重,使王受虧損呢?」

4:24於是,在耶路撒冷上帝殿的工程就停止了,直停到波斯王大流士第二年.

主題由「建殿」換成了「建城」.「城」是軍事設施,政治意味加強了.在這段諭旨裡,「宗教議題」換成了「政治議題」,在王眼中,這是很敏感的.如此,亞達薛西王下了「聖旨」,禁止猶太人建殿.背後的動機很明顯:「為何容害加重,使王受虧損呢?」

4:24的經文,忽然又返回之前的大流士王.

(5)找到先王塞魯士的諭旨,大流士王下旨:興建聖殿

5:1那時,先知哈該和易多的孫子撒迦利亞奉以色列上帝的名向猶大和耶路撒冷的猶大人說勸勉的話.

5:2於是撒拉鐵的兒子所羅巴伯和約薩達的兒子耶書亞都起來動手建造耶路撒冷上帝的殿,有上帝的先知在那裡幫助他們.

5:3當時河西的總督達乃和示他•波斯乃,並他們的同黨來問說:「誰降旨讓你們建造這殿,修成這牆呢?」

5:4我們便告訴他們建造這殿的人叫什麼名字.

5:5上帝的眼目看顧猶大的長老,以致總督等沒有叫他們停工,直到這事奏告大流士,得著他的回諭.

5:6河西的總督達乃和示他•波斯乃,並他們的同黨,就是住河西的亞法薩迦人,上本奏告大流士王.

主題又回到「聖殿」.對敵人來說,這是政治的「城」的問題；對猶太人來說,這只是「殿」的問題.因著先知的鼓勵,猶太人並沒有停止建殿.其中的信仰思想是很清晰:敬拜上帝,重建身分,這是頭等要事,不應因外力終止.

為建殿的事,猶太地的總督上奏大流士王,如此,就造成人間不同的王之間的「聖旨對決」.

6:1於是大流士王降旨(tǝēm,「旨」),要尋察典籍庫內,就是在巴比倫藏寶物之處；

6:2在米底亞省亞馬他城的宮內尋得一卷,其中記著說:

6:3「塞魯士王元年,他降旨(tǝēm,「旨」)論到耶路撒冷上帝的殿,要建造這殿為獻祭之處,堅立殿的根基.殿高六十肘,寬六十肘,

6:4用三層大石頭,一層新木頭,經費要出於王庫；

6:5並且上帝殿的金銀器皿,就是尼布甲尼撒從耶路撒冷的殿中掠到巴比倫的,要歸還帶到耶路撒冷的殿中,各按原處放在上帝的殿裡.」

找到先王塞魯士的諭旨記錄,對比第1章的寫法,這裡的諭旨多了一些條款,包括「經費要出於王庫」.

之前,強調聖殿不容外人插手,強調「分別為聖」,故此沒有這條款.

如今,強調先王有旨,強調先王強力支持建殿.如何支持?甚至動用國庫.

回到塞魯士的聖旨,是要讀者們記起:在塞魯士的聖旨之上,是上帝的聖旨.人間的王,有不同聖旨,有支持建殿的,有反對建殿的.但最終,重要的,是誰的旨意終極地指引著天地?

6:7〔大流士王降旨:〕「不要攔阻上帝殿的工作,任憑猶大人的省長和猶大人的長老在原處建造上帝的這殿.

6:8我又降旨(tǝēm,「旨」),吩咐你們向猶大人的長老為建造上帝的殿當怎樣行,就是從河西的款項中,急速撥取貢銀作他們的經費,免得耽誤工作.

6:9他們與天上的上帝獻燔祭所需用的公牛犢,公綿羊,綿羊羔,並所用的麥子,鹽,酒,油,都要照耶路撒冷祭司的話,每日供給他們,不得有誤；

6:10好叫他們獻馨香的祭給天上的上帝,又為王和王眾子的壽命祈禱.

6:11我再降旨(tǝēm,「旨」),無論誰更改這命令,必從他房屋中拆出一根樑來,把他舉起,懸在其上,又使他的房屋成為糞堆.

6:12若有王和民伸手更改這命令,拆毀這殿,願那使耶路撒冷的殿作為他名居所的上帝將他們滅絕.我－大流士降這旨意(tǝēm,「旨」),當速速遵行.」

這聖旨的特別之處,是有「咒詛」條款.誰阻礙建殿,死無葬身之地.

6:14……他們遵著以色列上帝的命令(tǝēm,「旨」)和波斯王塞魯士,大流士,(按:這裡應有亞哈隨魯),亞達薛西的旨意(tǝēm,「旨」),建造完畢.

6:15大流士王第六年,亞達月初三日,這殿修成了.

大家看到,《以斯拉記》的作者,是完全知道列王的歷史次序的,也知道聖殿是在大流士王年間建成的.但在建殿的記述中,加插後來的亞達薛西王的反對,這是要凸顯:建殿的事不可能一帆風順,人間的聖旨也不斷轉變,在這歷史過程中,神的子民應看到「誰的旨意主宰人間」,「誰掌管明天」.

建殿完成,6:19即刻返回「希伯來文」的書寫.「希伯來文」,是書寫神的旨意的文字.

6:19正月十四日,被擄歸回的人守逾越節.

6:20原來,祭司和利未人一同自潔,無一人不潔淨.……

並為自己宰逾越節的羊羔.

6:21從擄到之地歸回的以色列人和一切除掉所染外邦人污穢,歸附他們,要尋求耶和華－以色列上帝的人都吃這羊羔,

6:22歡歡喜喜地守除酵節七日；因為耶和華使他們歡喜,又使亞述王的心轉向他們,堅固他們的手,作以色列上帝殿的工程.

聖殿建成,猶太人守逾越節.逾越節是慶祝「離開為奴之地的埃及」的節日.聖殿建成,就是神的旨意行在地上,如同行在天上.這是創造了新的天地,讓人得解放,享自由.

(6)亞達薛西王的轉變(7:12-26):重新下旨,支持建殿

經文在7:12,又轉用「亞蘭文」書寫.

7:21「我－亞達薛西王又降旨與河西的一切庫官,說:『通達天上上帝律法的文士祭司以斯拉,無論向你們要什麼,你們要速速地備辦……

7:23凡天上之上帝所吩咐的,當為天上上帝的殿詳細辦理.為何使忿怒臨到王和王眾子的國呢?」

之前,為了自己的利益而反對建殿的亞達薛西王,如今改變了.他敬畏神的心,大於他對利益的考量.他的改變越戲劇性,就越發顯明神的主權在人間.

(7)總結

建造敬拜神的聖殿,是人生首要大事.建造好聖殿,才能建造好人生.其中的過程是艱難的,有敵人,也有不同的「聖旨」.我們的信仰是:神的旨意才是終極的,祂的旨意會成就在地上,如同成就在天上.在「聖旨」混雜的人間,你聽從誰的「旨」呢?

願榮耀歸上主,願眾人享平安.



\newpage

\section{領受聖靈吧}
\label{sec:648}
\textbf{高國雄牧師 - 領受聖靈吧!}
\newline
\newline
\hyperref[sec:647]{< < < PREV SERMON < < <}
~
\hyperref[sec:toc]{[返目錄]}
~
\hyperref[sec:649]{> > > NEXT SERMON > > >}


講題:領受聖靈吧!

經文:約20:21-23

講員:高國雄牧師

日期:2023年5月28日

記得有一次有一位學生與我討論,分享到關於華人教會對聖靈的理解,他觀察發現華人教會群體很少提及聖靈,甚至不少弟兄姊妹參與教會十數年,從來沒有在講道中聽到有關聖靈的教導.「聖靈降臨節」完全沒有在他們的信仰生活中出現過,華人教會對「聖靈」的迴避可能是由於過去靈恩運動的影響,曾經因為多次靈恩運動當中有一些信仰群體比較個人,又或比較熱烈的表達,讓很多不明所意的弟兄姊妹感到困惑(甚至覺得恐怖難明),很多時候有關聖靈感動這一個課題,只會在信徒間私下討論分享,因為好像是非常個人的事情,與教會整體很難連上密切的關係,甚至我聽過有牧者向弟兄姊妹教導,我們只要勤讀聖經,真理只在聖經當中,不要有太多個人的感受體會與想法,不要擅自用個人的感受而去解釋信仰,當然我明白面對信仰要謹慎,要認真,但個人經歷這層面是必須的,今天我們(教會)在世上也是被聖靈一直帶領與感動的活著的.

當我再作更深一層的思考,發現不單是華人教會,其實在教會的傳統中,很多時候我們也沒有花太多時間心力去理解研究聖靈今天的工作,正如我們每一次在崇拜講道後同誦讀的「使徒信經」,「信經」的出現是為了要回應當日教會出現異端的問題,所以你會發現使徒信經花了最多的內容在描述耶穌基督,由他的出生到他釘十字架,受死,埋葬,復活升天,並且將要再來,很仔細D 描述；對創造萬物的天父也很直接地表示出他就是創造天地的主,這認知我們也很清晰,但對於三一上帝中聖靈這個位格,使徒信經只有「我信聖靈」這幾個字,甚至比起接著描述的教會「我信聖而公之教會」所描述更少,因為對教會也表達出教會「聖」而「公」的特色,反之只有「我信聖靈」這四個字,我們真的很難由此作出更深入的討論及研究.

當然早期教會也因為對聖靈這位格(或者只是能力)而作出很多不同的討論,甚至由此產生不同的異端,在這裡我不會詳述.因此今天的經課:林前12:11已經非常清楚

林前12:11:「這一切都是由惟一的,同一位聖靈所運行,隨著自己的旨意分給各人的.」

但當我們很清晰地表達我們信的是三位一體的上帝,是聖父,聖子,聖靈,但我們真的很少探討聖靈的特色及工作,對聖靈的認識也很缺乏,各位弟兄姊妹回想一下你過去的禱告中,有多少比例是向「天父上帝」,「阿爸父」祈求?有多少時間是向耶穌基督我們的救主祈求?又有多少時間向著聖靈呢?

靈是風,亦是氣,代表著生命與氣息.

今天詩篇經課

詩104:29你收回牠們的氣,牠們就死亡,歸於塵土.

當中所提到當「氣」被收回,生命就完結.

而「靈」的出現早在創造就開始

創2:7「耶和華上帝用地上的塵土造人,將生命之氣吹進他的鼻孔,這人就成了有靈的活人」

復活的主再一次使人回復與上帝的關係

今天的經課

約20:21-23

21於是耶穌又對他們說:「願你們平安!父怎樣差遣了我,我也照樣差遣你們.」

22說了這話,他向他們吹一口氣,說:「領受聖靈吧!

23你們赦免誰的罪,誰的罪就得赦免；你們不赦免誰的罪,誰的罪就不得赦免.」

耶穌向門徒吹氣,門徒所領受的不單是力量,不單是一份感覺/感動,而是上帝聖靈自己,得著的赦罪權柄是因為有主的同在.

使徒行傳中聖靈震撼地降臨:

使2:1-4

1五旬節那日到了,他們全都聚集在一起.

2忽然,有響聲從天上下來,好像一陣大風吹過,充滿了他們所坐的整座屋子；

3又有舌頭如火焰向他們顯現,分開落在他們每個人身上.

4他們都被聖靈充滿,就按著聖靈所賜的口才說起別國的話來.

從天而來的巨響,好像大風,充滿房間,聖靈的舌頭如火焰落在各人身上,聖靈充滿,發出方言…

弟兄姊妹,請問你羨慕這個景象嗎?今天就是聖靈降臨主日,你有想像過今天聖靈要同樣地臨到我們當中嗎?你是期望還是害怕呢?

當我們繼續往下看這一段經文時,你會發現門徒突如其來的說方言的能力帶來很大的迴響,由於當時很多來自各地不同的猶太人(猶橋),他們來到耶路撒冷,可能是因為他們要去朝聖,由於他們來自不同的地方,他們有自己不同的語言(或者這樣說他們有不同的母語),當聖靈降臨在門徒身上,他們就宣告,而來自不同地方的猶太人他們竟然聽到自己不同的母語(鄉談)…

各位弟兄姊妹請問你認為門徒有能力說出不同的語言,比起剛剛的境象,「響聲從天上來好像大風,有聖靈的舌頭如火焰在他們的頭上」,到底哪一個畫面更震撼呢?

如果這裡的目的是要展現上帝的能力,如果聖靈降臨是單單為要展現上主的能力,更震撼的畫面(例如醫病趕鬼等等)的出現難度不比門徒有說不同語言(說不同鄉談)來得更精彩,更吸引嗎?對上帝來說應該沒有什麼困難…

就如我們今天,我國祈求聖靈給我們不同的能力,就如經課哥林多前書當中提到聖靈顯在各人身上,賜給我們智慧的言語,賜給我們知識的言語,也賜給我們信心,更有醫病的恩賜,行異能,作先知,辨別諸靈等等恩賜…

當我們祈求期盼著這些恩賜出現在我們身上,我們是為甚麼目的的呢?

保羅提醒哥林多教會肢體雖多,仍是一個身子,我們有很多不同能力/恩賜,不同的差異,但是都是受同一位聖靈的洗,一切都是由同一位聖靈隨祂的意思賜給我們的.

使徒行傳中聖靈降臨的時候,清楚讓我們知道聖靈的能力,最大最重要的能力,其實就是讓人溝通的能力,就是讓人能夠無阻隔地把信仰傳開的能力,這也同樣回應到保羅提醒哥林多教會一樣,各人從聖靈中領受不同的恩賜,就是要合一地見證同一位基督.

今天是聖靈降臨主日,我們預備好接受聖靈嗎?就是接受更多/更大的能力,以致能更有效地與人溝通,傳揚福音,以致更能夠與身邊未信主(不同語言)的朋友溝通,亦更能夠與我們的上主溝通,這也是信徒必須的,因為沒有靈,就沒有生命…

當日耶穌向門徒吹氣,與及在五旬節聖靈的降臨,你準備好了嗎?

讓我們領受聖靈吧!



\newpage

\section{三一萬能俠}
\label{sec:649}
\textbf{林豪恩先生 - 三一萬能俠}
\newline
\newline
\hyperref[sec:648]{< < < PREV SERMON < < <}
~
\hyperref[sec:toc]{[返目錄]}
~
\hyperref[sec:650]{> > > NEXT SERMON > > >}


講題:三一萬能俠

講員:林豪恩先生

日期:2023年6月4日

一.三一主日

哥林多後書13:14:願主耶穌基督的恩惠,上帝的慈愛,聖靈的感動常與你們眾人同在!

馬太福音28:19:所以,你們要去,使萬民作我的門徒,奉父,子,聖靈的名給他們施洗.

大家有沒有留意到經課重覆出現聖父,聖子,聖靈?三一主日是聖靈降臨期的第一個主日.主後1334年,教會規定以此主日作為紀念三位一體的教義.有人以為三位一體是神學家想出來的概念,其實不盡然如此,早在教會初期,三位一體的這個結構就已經出現.例如保羅的書信中,常常出現:「願主耶穌基督的恩惠,天父上帝的慈愛,聖靈的感動常與你們同在」這一句.

今日的講題「三一萬能俠」是漫畫改編成的電視動畫.「三一萬能俠」由三架飛機合體組成,分別為:「1號–空戰用機」,「2號–陸戰地底用機」,「3號–海戰用機」.因此,「三一萬能俠」涵蓋了海陸空三個領域.

信義會監督張振華牧師在《教會年曆實用指南》解釋三一主日的意義說:「我們認識三位一體上帝,就解答了我們人類三個必然詢問自己的問題:我從哪裡來?我怎樣過一個稱為義的生活?最後,就是我的將來如何?」借用電視動畫「三一萬能俠」的概念來思考這三個問題,三位一體的上帝也涵蓋了過去,現在和將來三個時空.

二.我從哪裡來?

首先讓我們思考第一個問題:「我從哪裡來?」,這個問題提示我們回望「過去」.讓我們透過創世記1:1-2:4尋找啟發.

經文相當長,剛才讀過就不再重覆了.請大家特別留意這個段落的最後一句,和修版譯為:「這就是天地創造的來歷.」「來歷」這個字有幾個翻譯,分別為:後代,家譜,來歷...等.「這是ＸＸ的來歷」這格式在創世記中出現了多次,例子有:

6:9這是挪亞的後代.

10:1這是挪亞的兒子閃,含,雅弗的後代.

25:19這是亞伯拉罕的兒子以撒的後代.

創世記起首這段落就是要告訴讀者天地萬物,包括人類的來歷.其實,來歷重要嗎?或者「我從哪裡來?」這個問題重要嗎?有人未必認為重要.齊豫唱出三毛的作品「橄欖樹」如此說:

不要問我從哪裡來我的故鄉在遠方

為什麼流浪流浪遠方

流浪

為了天空飛翔的小鳥為了山間輕流的小溪

為了寬闊的草原流浪遠方

流浪

還有還有為了夢中的橄欖樹

橄欖樹

不要問我從哪裡來我的故鄉在遠方

為什麼流浪為什麼流浪遠方

為了我夢中的橄欖樹

不要問我從哪裡來!好像很灑脫,然而,真的可以永遠無根地流浪嗎?有一首詩歌如此回應:

不要像流雲,流浪天邊,不要像浮萍漂浮水面,

它們不問方向如何分辨,剎那間,蹤影就消失不見.

你可知道在那永恆世界,信主的人享受永生無限；

你我生命不是雲煙,不是為偶然出現.

「我從哪裡來?」這個問題關乎過去,然而,過去與現在和將來是不容易割斷的.盧雲神父寫了一本書,書名叫《浪子回頭》,是他默想林布蘭同名油畫而來的感悟.油畫描繪的是聖經「浪子回頭」的故事.盧雲觀賞油畫,發現自己的心靈和那個浪子沒分別,一直想出人頭地,想向全世界證明自己,因此不斷打拼,與別人競爭,心靈惶恐不安,失落沮喪,沒有生活在愛中,在家裡的安全感.他如離開父家的浪子一樣,都是在遠方過著孤獨的生活.其實他深深渴望回家,投靠在天父的懷抱,尋回他的身份,他的真我.

「我從哪裡來?」不是關於地點的問題,而是關於歸屬的問題.盧雲後來發現,沒有離家的大兒子其實也如小兒子一樣,沒有生活在愛中,在家裡的安全感,都好像是在遠方過著孤獨的生活.流浪在外的小兒子渴望回家,在家的大兒子其實也渴望回家.回家,與地點無關,如盧雲神父所說:深深渴望回家,投靠在天父的懷抱,尋回身份和真我.你從哪裡來的呢?你的歸屬在哪裡呢?

三.我怎樣過一個稱為義的生活?

第二個問題是:「我怎樣過一個稱為義的生活?」詩篇第八篇提示我們從自己如何看待自己開始:

第一,不宜妄自尊大

8:3我觀看你指頭所造的天,並你所陳設的月亮星宿,

8:4便說:人算甚麼,你竟顧念他!世人算甚麼,你竟眷顧他!

望星空穹蒼,念天地悠悠,很容易產生「人算甚麼」的感悟.轉個頭,回到人間,卻往往不接受自己算不得甚麼,因此常常妄自尊大.相信大家都同意,妄自尊大並不是過一個稱為義的生活的合適態度,反而可能造成災難和悲劇.電視劇《秦始皇》的主題曲對狂傲的表現描繪得非常突出:

大地在我腳下

國計掌於手中

哪個再敢多說話

這些妄自尊大的人,怎會讓別人說話呢?秦始皇採納李斯的建議,史書只保留秦國史書,其他國家的史書都焚毀；《詩篇》,《書經》及諸子百家之書只有博士官職可以保留,民間的都限期交出燒毀.《秦始皇》主題曲接著的幾句是:

夷平六國是誰

哪個統一稱霸

誰人戰績高過孤家

稱霸,顯戰績,一將功成的代價是甚麼呢?你知道誰在烏克蘭東部巴赫穆特作戰嗎?俄羅斯僱傭兵首領普里戈任五月二十四日表示,瓦格納集體招募了5萬名囚犯到巴赫穆特作戰,當中兩成人戰死,另陣亡的合約僱傭兵也有約一萬人,但稱烏克蘭也有5萬至7萬人受傷,5萬人死亡.他抱怨說,「戰爭突顯社會不平等,窮人兒子的屍體從前線被送回來,精英階層的孩子則在陽光下『搖著屁股』」.

第二,不宜妄自菲薄

人類不宜妄自尊大,然而,也不宜妄自菲薄.為甚麼呢?因為:

第一,人雖渺小,卻是上帝顧念的對象.

8:4便說:人算甚麼,你竟顧念他!世人算甚麼,你竟眷顧他!

第二,人雖如嬰孩,上帝仍賜予能力

8:2你因敵人的緣故,從嬰孩和吃奶的口中,建立了能力,使仇敵和報仇的閉口無言.

第三,人雖微小,上帝仍賜予身分位置

8:5你叫他比天使(上帝)微小一點,並賜他榮耀尊貴為冠冕.

第四,人雖不算甚麼,上帝仍委以重任

8:6-8你派他管理你手所造的,使萬物,就是一切的牛羊,田野的獸,空中的鳥,海裡的魚,凡經行海道的,都服在他的腳下.

人類不宜妄自尊大,也不宜妄自菲薄,那麼,人應該怎樣生活呢?寫創世記釋經書的舊約學者鄺炳釗博士指出,人類既然被上帝所重視所愛護,也應該自重自愛.以下是自重自愛的人的其中幾項特點《從聖經看如何活得更快樂》:

他們感到自由,因為能夠接納自己,明白世界上沒有完美的人,所以願意承認自己的缺點,也能夠欣賞自己的優點.

他們有信心,有安全感,不用整天與人比較.

他們曉得健康地愛護並且照顧自己,不會為了討好別人而令自己心裡充滿憎恨和不快.

他們樂意接受稱讚,欣然地說「謝謝」.

他們常常是好父母,接納子女的優點和弱點,尤其在子女失敗的時候,不會因此而釀成自卑和羞恥.

人們喜歡和自重自愛的人在一起,因為他們不會太敏感,也因為他們不怕取笑自己,而且願意接受別人無傷大雅的玩笑.

我怎樣過一個稱為義的生活?在個人層面,詩人提示我們不宜妄自尊大,也不宜妄自菲薄,最合適的是自重自愛.在群體層面,讓我們看看保羅對哥林多信徒的忠告:

13:11末了,弟兄們,願你們喜樂.要追求完全；要接受鼓勵；要同心合意；要彼此和睦.如此,慈愛和平的神必與你們同在.

保羅勸勉哥林多信徒,第一,同心合意；第二,彼此和睦.現實是「一人有一個夢想,兩人熱愛漸迷惘,三人有三種愛找各自理想」,怎可能「同心合意」呢?世界充滿混亂和衝突,「彼此和睦」怎可能呢?甚至有人語帶諷刺說,除非有人操控其他人的思想,控制所有人的行為,就有可能「同心合意」,「彼此和睦」了.有一位校友離港進修,所讀的課程名為「衝突與和平」,曾經赴盧旺達工作.我問他,推動和平的工作如何?他說非常艱難,但世界上仍有不少人為此而努力.下個學年,他將會再入學校,繼續做「衝突與和平」的研究.

追求「同心合意」,「彼此和睦」,或許不應著眼於內容,即是說不是要所有人都同意同一內容,接受同一行為規條；反而是接受「不同」為現實,培養能夠幫助我們邁向「同心合意」,「彼此和睦」的態度和價值觀.在這裡,保羅提出了這種態度的基礎,這個基礎是關乎我們如何看待其他人的.讓我們再聽一次保羅對哥林多教會的忠告:

13:11末了,弟兄們,願你們喜樂.要追求完全；要接受鼓勵；要同心合意；要彼此和睦.如此,慈愛和平的神必與你們同在.

保羅如何稱呼哥林多信徒呢?對於哥林多教會來說,保羅是使徒,是屬靈父親,然而,這些都不是最重要的.最後,保羅稱收信人為「弟兄姊妹」.弟兄姊妹是互相之間最基本的身分.互道弟兄姊妹,讓我們意識到大家是平等的,沒有誰高級過誰；以弟兄姊妹互稱,讓我們感到大家是同根同源的同伴,而不是互相排斥的異己.人們平等相待是很困難的,這是事實,因為人人都想自己高於其他人,也想自己的民族優於其他民族.然而,聖子耶穌基督已經為世人立下榜樣:「他本有神的形像,卻不堅持自己與神同等；反倒虛己,取了奴僕的形像,成為人的樣式」(腓立比書2:6-7).對於第二個問題:「我怎樣過一個稱為義的生活?」保羅給我們的啟發是建立互視為弟兄姊妹的平等態度和價值.

四.我的將來如何?

最後,讓我們透過馬太福音28:16-20節思考「我的將來如何?」這個問題.

28:18耶穌進前來,對他們說:「天上地下所有的權柄都賜給我了.」既然如此,門徒可以合理地問自己的將來,擁有「天上地下所有的權柄」的耶穌應該不會不知道的,對嗎?不過,耶穌並沒有揭示他們的個人前程,或許就正如那首舊歌Whatever will be, will be所唱的:

當我還是個小女孩

我問媽媽:”我將來會變成什麼樣子呢?”

“我會變得漂亮嗎?”

“我會變得富有嗎?”

她這麼回答我

順其自然吧

該來的總會到來

我們無法預知未來

順其自然吧

該來的總會到來

現在我有自己的孩子

他們問我:”我長大後會成為怎樣的人”

“我會變得帥氣嗎?”

“我會變得富有嗎?”

我溫柔地回應他們

順其自然吧

該來的總會到來

我們無法預知未來

順其自然吧

該來的總會到來

耶穌沒有預告前程,只告訴門徒現在要做的事,就是「使萬民作我的門徒,奉父,子,聖靈的名給他們施洗.凡我所吩咐你們的,都教訓他們遵守」.對於將來,耶穌只應許「我就常與你們同在」,這應許也是在約翰福音所解釋的聖靈要來作信徒保衛師的應許.憑著這應許,有詩人寫了【我知誰掌管明天】這首詩歌,讓我們以此詩歌作為結束及互勉:

我不知明天將如何,每一天只為主活,

我不借明天的陽光,因明天或不晴朗,

我不要為將來憂慮,因我信主的應許,

我今天要與主同行,因祂知前面路程.

有許多未來的事情,我現在不能識透,

但我知誰掌管明天,我也知誰牽我手.

我不知明天將如何,或遭遇貧苦飢餓,

但那位看顧麻雀者,祂必然也看顧我,

祂是我旅途的良伴,縱遭遇各樣災害,

我救主必與我同在,祂寶血把我遮蓋.

有許多未來的事情,我現在不能識透,

但我知誰掌管明天,我也知誰牽我手.



\newpage

\section{lekh-lekhah你要去}
\label{sec:650}
\textbf{李均熊牧師 - lekh-lekhah 你要去!}
\newline
\newline
\hyperref[sec:649]{< < < PREV SERMON < < <}
~
\hyperref[sec:toc]{[返目錄]}
~
\hyperref[sec:651]{> > > NEXT SERMON > > >}


講題:lekh-lekhah 你要去!

講員:李均熊牧師

日期:2023年6月11日

引言:甚麼是lekh-lekhah

從今天福音經課來看,今天宣講可以有兩個主題:一是關於呼召,二是關於信心.我選擇了舊約經課中從亞伯拉罕的生命來宣講關於信心的題目.亞伯拉罕被稱為信心之父,不單是因為我們憑信得稱為他的後嗣,更因為他信心的表現可以成為我們在人生旅程中實踐信心的榜樣.但今日的題目中有一個音譯詞lekh-lekhah,其實是出於創12:1經文,上主對還未改名為亞伯拉罕的亞伯蘭的呼召:「你要離開本地,本族,父家,往我所要指示你的地去」,則要解釋一下. lekh-lekhah 是希伯來文(לֶךְ־לְךָ)的音譯.和合本將此翻譯為「你要離開」,大部份英譯本會譯作Go out或Go forth,英皇欽定譯本King James Version 則譯作Get thee out.這個字不是容易翻譯的,因為這個字有兩個部份,lekh比較易理解,乃是行走的命令式,「去!」Go 就可以了.但經文這裡卻加了lekhah,乃是一個前置詞再加一個「你」,可以解作「給你」(to you),「為你」(for you),「被你」(by you)等等的可能性.為甚麼要拿這個lekh-lekhah去討論呢?是因為這樣組成的希伯來文在聖經其實只出現過兩次,一次在這裡,另一次就在創世記22:2,當神要亞伯拉罕獻以撒時,「神說:你帶著你的兒子,就是你獨生的兒子,你所愛的以撒,往摩利亞地去,在我所要指示你的山上,把他獻為燔祭.」在「往摩利亞地去」就包含了lekh-lekhah了.換句話說,這組希伯來文成為了貫串亞伯拉罕生平的關鍵字了.

起:亞伯拉罕是不問情由的相信?

當你看到亞伯拉罕生命中這兩處獨特的經驗時,不知你會不會問:當上帝對亞伯蘭說話,說:「你要離開本地,本族,父家,往我所要指示你的地去.」(創12:1)亞伯拉罕就照著耶和華的吩咐去行──當中沒有記載亞伯蘭的反應,沒有爭扎,沒有懷疑.為什麼沒有記載呢?是因為亞伯拉罕的「信心」中完全沒有懷疑嗎?那會是一個怎樣的信心?同樣的情形出現在與十二章遙遙相對的廿二章,上帝對亞伯拉罕說:「你帶著你的兒子,就是你獨生的兒子,你所愛的兒子以撒,往摩利亞地去,在我所要指示你的山上,把他獻為燔祭.」(創22:2)當中亦沒有記載亞伯拉罕的反應,思想,爭扎和懷疑.為什麼沒有呢?是亞伯拉罕的「信心」中完全沒有懷疑嗎?但亞伯拉罕怎可以沒有反應的呢?若亞伯拉罕是一個有感情,有意識的人,對於上主這樣的吩咐怎可以無動於衷.

這兩段中所用的字眼,都是容易令人有情緒反應的字眼.「本地,本族,父家」,從一個地理的位置「本地」,進到一個與自己有親密關聯的「父家」,若要離開,怎可能沒有感情上的反應?而「你的兒子,你獨生的兒子,你所愛的兒子」,每進一層關係,就加深一份亞伯拉罕和兒子的親密感,亞伯拉罕焉能沒有反應?然而,亞伯拉罕二話不說就「順服」上帝的吩咐.在今日看來,這是否出於一種盲目的信心?究竟亞伯拉罕的信心,是怎樣的一種信心?

承:亞伯拉罕的一生

亞伯拉罕的一生中,是否都沒有懷疑呢?兩個 lekh-lekhah 之間,就是亞伯拉罕實踐信心的生命旅程,猶太釋經提出亞伯拉罕一生遇過十次考驗,中世紀的解經家就認為在這兩個 lekh-lekhah 之間.離開父家之然是一個大考驗,在離開父家到達迦南,亞伯拉罕不單經歷饑荒,在逃避饑荒下埃及時,妻子就被帶入法老的皇宮；之後為了解救姪兒羅得,他不得已與迦南諸王爭戰起來.而「在這些事以後」,即是在這些考驗以後,上主說要賞賜亞伯拉罕(15:1),亞伯拉罕的回應是:「我既無子,你還賜我甚麼呢?」(15:2)；上主說要賜「這地」(即迦南地)為亞伯拉罕的產業,他卻問:「我怎能知道必得這為業呢?」(15:7-8)上主並沒有因為他的疑問而責怪亞伯拉罕,反而以立約的方法,去肯定祂自己曾說過的話(15:9-21)凡此種種,「亞伯蘭信耶和華,耶和華就以此為他的義.」(15:6)

可見,當上主與亞伯蘭立約時,亞伯蘭就發出了疑問(創15).懷疑,不單為上主所容許,之後帶來的信心,上主更會加以肯定.

另外一次記載上主與亞伯拉罕立約,就是當他九十九歲的時候,並以割禮為記(拉比們認為以九十九歲高齡要承受割禮,也算是亞伯拉罕十大考驗之一),他的名字亦由亞伯蘭改為亞伯拉罕,因他要成為多國之父.在此,當上主應許他的妻子要為他生一個孩子,亞伯拉罕卻沒有「即時」的相信,卻用「喜笑」作為回應(創17:17).當中亞伯拉罕不單帶著懷疑,而且亦含有情感的反應.在亞伯拉罕為所多瑪向上主求情的事情上(18:20-33),亞伯拉罕的討價還價:「無論善惡,你都要剿滅嗎?假若那城裡有五十個義人,你還剿滅那地方嗎?不為城裡這五十個義人饒恕其中的人嗎?將義人與惡人同殺,將義人與惡人一樣看待,這斷不是你所行的.審判全地的主豈不行公義嗎?」(創18:23-25)這表現出亞伯拉罕的信心並非盲目和缺乏反省的,也顯出其勇毅和堅定的一面.

亞伯拉罕的信並不是不問情由的相信,教會也應該鼓勵信徒對信仰發問,甚至勇敢的面質上帝──這就是「以信求知」的態度吧!

轉:懷疑甚麼?相信甚麼?

但今日書信經課,卻好像不是這樣說亞伯拉罕.羅馬書4:20-22說亞伯拉罕「仰望神的應許,總沒有因不信心裡起疑惑,反倒因信心裡得堅固,將榮耀歸給神,且滿心相信神所應許的必能做成.所以,這就算為他的義.」

在我看來,這裡經文說到亞伯拉罕不會疑惑的,乃是因為相信神的應許,但在亞伯拉罕的生命裡,就面對的處境仍是有疑惑的,這就是我仍為亞伯拉罕的信心是我和你都可以學習的.甚麼是上主給亞伯拉罕的應許呢,這我們就要回到今日的舊約經課了.

在上主呼召亞伯拉罕lekh-lekhah的時候,中世紀的猶太拉比Rashi(Rabi Salomon Isaacides 1040-1105)建議直接將lekhah譯為”for your benefit and for your good”,意即lekh-lekhah就是「為你的利益和你的好處,你要去!」因為接著的12:2-3節經文明言上主的呼召是帶著應許的:「我必叫你成為大國.我必賜福給你,叫你的名為大；你也要叫別人得福.」這一個解釋聚焦的,就是亞伯拉罕相信上主而踏上這一個成為流散異地的不確之旅,其中很大因素就是相信上主應許的實現(羅馬書4章則將焦點放在獻以撒上主仍能實現應許之上).

上主的應許,其實正正對應古代行旅的種種挑戰和困難.古代旅行並不如今天的方便,除了因為路上的危險而帶來人丁和金錢的損失,更重要的是你不能掌握你將會面對的是甚麼人和事,而且上主只說「往我所要指示你的地去」,對亞伯拉罕而言,這個不確定的旅程帶來的是可等大的挑戰和焦慮!而且對亞伯拉罕來說,離開本地本族父家,就沒有族群的保護,更重要的,是失去了古代人最重視的名聲.學者一直都認為,古以色列人所謂永生,就是子嗣繁衍,家族的名聲代代的傳承下去.但離開了本地本族父家,你的名字就不會有人認識和紀念了.

但在上主的呼召中,上主應許使亞伯拉罕的名為大,並使他成為大國,地上萬族也因他而得福.在一切的不確定中,上主應許必賜福給他.所以lekh-lekhah,上主呼召他離開自己的地方,去開展不確定的將來,是要「為你的好處」.亞伯拉罕踏上這不確的行旅,流散的旅程,並不是因為一切疑慮消取,才有不問情由的信心.亞伯拉罕的信心,是雖然前面充滿不確與焦慮,但因為相信上主的應許,又相信上主的信實必然實現他的應許,「仰望神的應許,總沒有因不信心裡起疑惑,反倒因信心裡得堅固,將榮耀歸給神,且滿心相信神所應許的必能做成.所以,這就算為他的義.」

合:在無可指望的時候,因信仍有指望

羅馬書4:18節是這樣肯定亞伯拉罕的信心的:「他在無可指望的時候,因信仍有指望,就得以作多國的父.」亞伯拉罕並非沒有經歷過困難絕望,他也並非盲阿Q的相信明天會更好,因為上主給祂的召命就是要他經歷前路的不確的.亞伯拉罕會懷疑現實的可能性,會質問上主的誠信公義為何,但他在無可指望的時候,仍有指望,是因信上主的應許必要實現,是相信上主所發的言語必然成就,是相信上主必然忠於祂所立的約.亞伯拉罕的指望是在於那愛他應許他能戰勝困難的主,他的信心是投放在上主的誠信中,以至在環境與現實證據下像是沒有盼望的時候,仍有盼望.

當環境令我們產生懷疑,當人生行旅充望不確定,甚至是困難重重,令我們感到沒有出路,再無指望時,我們因信那應許之主,可以仍有指望.雖然世事不確,令我們懷疑人生,但主對我們的愛和應許,因主復活大能,我們可以不用懷疑.並不是我們信心的高與低使我們有盼望,而是主立約的誠信是我們有盼望.

這就是亞伯接罕的信心,也可以是我們對主耶穌的信心.



\newpage

\section{在淤泥上植一池蓮花}
\label{sec:651}
\textbf{鄧瑞強博士 - 在淤泥上植一池蓮花}
\newline
\newline
\hyperref[sec:650]{< < < PREV SERMON < < <}
~
\hyperref[sec:toc]{[返目錄]}
~
\hyperref[sec:652]{> > > NEXT SERMON > > >}


尼希米記 13:10-31
\newline
\begin{longtable}{cl}
\hline
\hline
章節 & 經文 (和合本修訂版)\\
\hline
13:10 & \begin{tabularx}{0.7\textwidth}{X} 我發現利未人當得的份無人供給他們，甚至供職的利未人與歌唱的都各奔回自己的田地去了。 \end{tabularx} \\ \\ \relax
13:11 & \begin{tabularx}{0.7\textwidth}{X} 我就斥責官長說：「你們為何不顧神的殿呢？」於是我召集利未人，使他們在自己的崗位上供職。 \end{tabularx} \\ \\ \relax
13:12 & \begin{tabularx}{0.7\textwidth}{X} 猶大眾人就把五穀、新酒和新油的十分之一送入庫房。 \end{tabularx} \\ \\ \relax
13:13 & \begin{tabularx}{0.7\textwidth}{X} 我派示利米雅祭司、撒督文士和利未人毗大雅作司庫管理庫房，副手是哈難；哈難是撒刻的兒子，撒刻是瑪他尼的兒子；這些人都是忠實的，他們的職務是分派他們弟兄所當得的份。 \end{tabularx} \\ \\ \relax
13:14 & \begin{tabularx}{0.7\textwidth}{X} 我的神啊，求你因這事記念我，不要塗去我為神的殿與其中的禮儀所獻的忠心。 \end{tabularx} \\ \\ \relax
13:15 & \begin{tabularx}{0.7\textwidth}{X} 那些日子，我在猶大見有人在安息日踹醡酒池，搬運禾捆馱在驢上，又把酒、葡萄、無花果和各樣的擔子在安息日扛入耶路撒冷，我就在他們賣食物的那日警戒他們。 \end{tabularx} \\ \\ \relax
13:16 & \begin{tabularx}{0.7\textwidth}{X} 有一些住在城裡的推羅人也把魚和各樣貨物運進來，甚至在耶路撒冷，在安息日賣給猶大人。 \end{tabularx} \\ \\ \relax
13:17 & \begin{tabularx}{0.7\textwidth}{X} 我就斥責猶大的貴族，對他們說：「你們怎麼會做這惡事干犯安息日呢！ \end{tabularx} \\ \\ \relax
13:18 & \begin{tabularx}{0.7\textwidth}{X} 你們祖先豈不是這樣做，以致我們的神使一切災禍臨到我們和這城嗎？你們竟干犯安息日，使憤怒越發臨到以色列！」 \end{tabularx} \\ \\ \relax
13:19 & \begin{tabularx}{0.7\textwidth}{X} 安息日前一日黃昏的時候，我吩咐人把耶路撒冷城門鎖上；我又吩咐，不過安息日不准開門。我也派幾個僕人在城門口站崗，免得有人在安息日挑擔子進城。 \end{tabularx} \\ \\ \relax
13:20 & \begin{tabularx}{0.7\textwidth}{X} 於是商人和販賣各樣貨物的人，有一兩次在耶路撒冷城外過夜。 \end{tabularx} \\ \\ \relax
13:21 & \begin{tabularx}{0.7\textwidth}{X} 我警告他們說：「你們為何在城牆前過夜呢？若再這樣，我必下手辦你們。」從此以後，他們在安息日就不再來了。 \end{tabularx} \\ \\ \relax
13:22 & \begin{tabularx}{0.7\textwidth}{X} 我吩咐利未人潔淨自己來守城門，使安息日分別為聖。我的神啊，求你因這事記念我，照你豐盛的慈愛憐憫我。 \end{tabularx} \\ \\ \relax
13:23 & \begin{tabularx}{0.7\textwidth}{X} 那些日子，我又看見猶太人娶了亞實突、亞捫和摩押的女子為妻。 \end{tabularx} \\ \\ \relax
13:24 & \begin{tabularx}{0.7\textwidth}{X} 他們的兒女，一半說亞實突話，或其他種族的方言，不會說猶大話。 \end{tabularx} \\ \\ \relax
13:25 & \begin{tabularx}{0.7\textwidth}{X} 我就斥責他們，詛咒他們，打了他們幾個人，拔下他們的鬍鬚，叫他們指著神起誓：「你們不可把自己的女兒嫁給外邦人的兒子，也不可為自己和兒子娶他們的女兒。 \end{tabularx} \\ \\ \relax
13:26 & \begin{tabularx}{0.7\textwidth}{X} 以色列王所羅門不也在這樣的事上犯罪嗎？在許多國家中並沒有一位王像他，蒙他神喜愛，神立他作王治理全以色列。然而，連他也被外邦女子引誘犯罪。 \end{tabularx} \\ \\ \relax
13:27 & \begin{tabularx}{0.7\textwidth}{X} 我們豈能聽憑你們行這一切大惡，娶外邦女子干犯我們的神呢？」 \end{tabularx} \\ \\ \relax
13:28 & \begin{tabularx}{0.7\textwidth}{X} 以利亞實大祭司的孫子，耶何耶大的一個兒子是和倫人參巴拉的女婿，我就把他從我這裡趕出去。 \end{tabularx} \\ \\ \relax
13:29 & \begin{tabularx}{0.7\textwidth}{X} 我的神啊，求你記得他們的罪，因為他們玷污了祭司的職分，違背祭司和利未人的約。 \end{tabularx} \\ \\ \relax
13:30 & \begin{tabularx}{0.7\textwidth}{X} 這樣，我潔淨他們，使他們脫離屬外邦人的一切；我又分派祭司和利未人的班次，使他們各盡其職， \end{tabularx} \\ \\ \relax
13:31 & \begin{tabularx}{0.7\textwidth}{X} 按定期奉獻木柴和初熟的土產。我的神啊，求你記念我，施恩於我。 \end{tabularx} \\ \\
[1ex]
\hline
\hline
\end{longtable}
講題:在淤泥上植一池蓮花

講員:鄧瑞強博士

經文:尼希米記13:10-31

日期:2023年6月18日

今日講道會講《尼希米記》,《尼希米記》已是舊約歷史記述的尾聲.以色列人的歷史,一路走來,是「神對人的恩典」及「人對神的背叛」交織而成的長卷.昔日,神將火把交給亞伯拉罕,要他將光明帶進黑暗的人間.後來,他的子孫面對當時最強大的黑暗力量,就是奴役人的埃及.神聽到他們的哀聲,差派摩西去解救他們.幾經波折,他們來到迦南地.神希望他們在這裡安居樂業,以神美善的律法,耕耘一片人間樂土.在《士師記》,我們看到,以色列人將這裡搞得烏煙瘴氣.神容許他們設立君王制度,希望社會從始井井有條.到所羅門王時代,極盡繁榮,但也就在這時候,所羅門王變成另一個埃及法老,奴役人去建造「積貨城」.自始王國分裂,分成北面的以色列國及南面的猶大國.時移世異,國際政治興起不同的霸主,以色列國和猶大國為了對應這些霸主,玩起國際權力的遊戲來,慢慢放棄了自己的信仰價值.君王離棄上帝,人民的生活就越發腐敗.最後,便走向敗亡.以色列國亡於亞述,猶大國亡於巴比倫.人民失去家園,失去身分,成了流落異鄉的nobody.神再次憐憫這群人,興起波斯帝國,波斯的塞魯士大帝容讓猶太人回歸故土.《尼希米記》講的,便是猶太人回歸故土,重建家鄉,重覓身分的故事.問題是:這些歷盡劫難的人,能否從歷史中吸取教訓?

我們看第一段經文:尼希米記13:10-14

13:10我見利未人所當得的分無人供給他們,甚至供職的利未人與歌唱的俱各奔回自己的田地去了.

13:11我就斥責官長說:「為何離棄上帝的殿呢?」我便招聚利未人,使他們照舊供職.

13:12猶大眾人就把五穀,新酒,和油的十分之一送入庫房.

13:13我派祭司示利米雅,文士撒督,和利未人毗大雅作庫官管理庫房；副官是哈難.哈難是撒刻的兒子；撒刻是瑪他尼的兒子.這些人都是忠信的,他們的職分是將所供給的分給他們的弟兄.

13:14我的上帝啊,求你因這事記念我,不要塗抹我為上帝的殿與其中的禮節所行的善.

《尼希米記》這卷書的上半部分,講重建耶路撒冷的城牆,這是重建神的子民身分的硬件.《尼希米記》下半部分,講重建猶太人的信仰人格,這是重建神的子民身分的軟件.我們發覺,建造硬件是相對容易的,塑造人性卻是艱難得多.

在尼希米心中,人的罪,社會的墮落,神子民身分的失落,最核心的問題,在於人不再敬畏神.人不再敬畏神,就失去超越的視野,失去對他者的絕對尊重,失去對天命的領悟,失去那要在世上成就永恆價值的獨特身分的自覺.敬畏神,對尼希米來說,就表現在維護敬拜神的聖殿體制之上.

在13:10-14的經文裡,我們看到:無人供應服事聖殿的利未人,無人支援頌讚神的詩班員,無人顧念聖殿的運作.一言以蔽之,人只顧自己塵世的生活,無心思索永恆的事情.人心遠離了上帝,失去了神聖,人性裡的醜惡就更容易爆發出來,也就更輕易拿起武器,殘害身邊的人.有人說,沒有了上帝,人就可任意妄為.的確,人若不敬畏上帝,還有什麼東西可限制人的自私,貪婪,暴力?

是人性的敗亡,導致民族的敗亡.要避免猶太人重蹈覆轍,首先要重建他們神的子民的身分.這身分需要的,是敬畏神的心.敬畏神的心,體現在尊重聖殿的禮儀上.尼希米用心地重整聖殿的體制.在這段落的結尾,是尼希米的祈禱.

13:14我的上帝啊,求你因這事記念我,不要塗抹我為上帝的殿與其中的禮節所行的善.

用心聽這祈禱,會感到一陣心酸.為什麼尼希米祈求神記念他一個人?仿佛,天下間,只有他,孤獨一人,記掛著聖殿的事.其他人,在做什麼?或許,只謀求自己的福樂.看來,只有尼希米,在人性傾頹之際,孤獨地,力挽狂瀾.

我們看第二段經文:尼希米記13:15-22

13:15那些日子,我在猶大見有人在安息日醡酒,搬運禾捆馱在驢上,又把酒,葡萄,無花果,和各樣的擔子在安息日擔入耶路撒冷,我就在他們賣食物的那日警戒他們.

13:16又有泰爾人住在耶路撒冷；他們把魚和各樣貨物運進來,在安息日賣給猶大人.

13:17我就斥責猶大的貴冑說:「你們怎麼行這惡事犯了安息日呢?

13:18從前你們列祖豈不是這樣行,以致我們上帝使一切災禍臨到我們和這城嗎?現在你們還犯安息日,使忿怒越發臨到以色列!」

13:19在安息日的前一日,耶路撒冷城門有黑影的時候,我就吩咐人將門關鎖,不過安息日不准開放.我又派我幾個僕人管理城門,免得有人在安息日擔什麼擔子進城.

13:20於是商人和販賣各樣貨物的,一兩次住宿在耶路撒冷城外.

13:21我就警戒他們說:「你們為何在城外住宿呢?若再這樣,我必下手拿辦你們.」從此以後,他們在安息日不再來了.

13:22我吩咐利未人潔淨自己,來守城門,使安息日為聖.我的上帝啊,求你因這事記念我,照你的大慈愛憐恤我.

一種信仰人格,就有一種信仰態度.一種信仰態度,就有一種經濟行為.若果在敬拜神的日子裡,做生意可賺錢,則當時不敬神的猶太人會問:那為何不做生意賺錢,而去敬拜神呢?他們在安息日,大做買賣.

安息日是什麼日子?

在《創世記》,神造萬物,造完後,神都看為是「好」的.在第七日,神什麼都沒有造,祂停了下來,安息了,然後,神賜福這日,將之祝聖為「聖」日.在人間,被造的萬事萬物,最多都只是「好」的,並沒有是「聖」的.「聖」,是專屬於神的.在創造過程裡,神沒有祝聖任何「實物」,免得人去爭這些聖物.就像《奪寶奇兵》電影系列裡的情節,各人不斷去爭奪約櫃,聖杯.在創造過程裡,神沒有祝聖任何「實物」,神祝聖人無法儲存在夾萬裡的時間,神祝聖「安息日」.安息日是聖的,是屬於神的.神在這聖日臨在,親近敬拜祂的人.

舊約學者布魯格曼(Walter Brueggemann)寫了一本書,書名叫《Sabbath as Resistance》,安息日是對塵世價值觀的一種反抗.安息日反抗什麼呢?它反抗不專心敬拜神的一切行為.它反抗我們一邊崇拜,一邊打whatsapp,或處理購物清單,或處理公事.安息日,重點不在放假休息,而在放下塵世的一切煩慮,專心親近神,專心體會神的臨在,專心沐浴在神的愛裡.但,人就是不專心,在親近神的時刻,卻同時用種種方法,努力去控制生活上的各樣事情.布魯格曼在《Sabbath as Resistance》這本書裡面說:守安息日的時候,卻同時在做買賣,就像與愛人溫馨時,卻不斷看錶一樣.

尼希米時代那些回歸故土的猶太人,他們就是那樣不重視神在安息日的臨在.他們在安息日做買賣,或許那些貨物有時限,要及時買賣,或許他們急不及待要賺更多的錢,或許他們要買些東西作當下享受之用,或許,或許,總有諸多藉口.要在親近神之時做買賣；要在敬拜神之時,同時想賺得世界；要在以神為最高價值時,同時不忘高舉其他價值,要這樣做,總有無限理由.

安息日,要求我們停下來,面對神.這是反抗我們用盡時間,甚至侵佔神聖時間,去操控這個世界.當人不理會神聖的界線而瘋狂追求利益時,便會淪為貪婪的奴隸.為了重建猶太人在安息日裡的信仰生活,為了重整他們的經濟行為,為了重整他們的人性,尼希米不惜在安息日前後,關上城門,截斷物流通道.在這段落的結尾,再次是尼希米的祈禱.

13:22……

我的上帝啊,求你因這事記念我,照你的大慈愛憐恤我.

同樣,看來,天下間,只有尼希米一人,期望著這些回歸的猶太人,在安息日,專心仰望神的臨在.但,看來,他的同胞,在安息日,仍記掛著做買賣.

我們看第三段經文:尼希米記13:23-31

13:23那些日子,我也見猶大人娶了亞實突,亞捫,摩押的女子為妻.

13:24他們的兒女說話,一半是亞實突的話,不會說猶大的話,所說的是照著各族的方言.

13:25我就斥責他們,咒詛他們,打了他們幾個人,拔下他們的頭髮,叫他們指著上帝起誓,必不將自己的女兒嫁給外邦人的兒子,也不為自己和兒子娶他們的女兒.

13:26我又說:「以色列王所羅門不是在這樣的事上犯罪嗎?在多國中並沒有一王像他,且蒙他上帝所愛,上帝立他作以色列全國的王；然而連他也被外邦女子引誘犯罪.

13:27如此,我豈聽你們行這大惡,娶外邦女子干犯我們的上帝呢?」

13:28大祭司以利亞實的孫子,耶何耶大的一個兒子是和倫人參巴拉的女婿,我就從我這裡把他趕出去.

13:29我的上帝啊,求你記念他們的罪；因為他們玷污了祭司的職任,違背你與祭司利未人所立的約.

13:30這樣,我潔淨他們,使他們離絕一切外邦人,派定祭司和利未人的班次,使他們各盡其職.

13:31我又派百姓按定期獻柴和初熟的土產.我的上帝啊,求你記念我,施恩與我.

語言,是一個人的身分認同的核心元素.在這裡,特別提及所羅門王娶外族女子為妻,因而轉向敬拜她們的神明的事.如此,凸顯所謂不再講猶大語言,正是遺忘神的子民的身分問題.尼希米在這裡最擔心的,是異族通婚令猶太人失去了言說自己的傳統語言的能力.他們的傳統語言,承載著他們祖先的信仰故事,貫注了上帝對他們民族的叮嚀,也銘刻著民族的苦難記憶.失去了語言,就將遺忘身分.假以時日,就不再記起自己作為神的子民的身分了.

現代社會是語言氾濫的社會,我們掌握致富的經濟語言,致勝的競爭語言,達致自我滿足的心理技術語言,這是我們在社會裡能互相溝通的語言.熟習這些語言之時,我們可能會漸漸忘記上帝的語言,上帝的話語.久而久之,我們會與世俗同化,忘記了自己的真正身分.

在這段落的結尾,再一次是尼希米的祈禱.

13:31……

我的上帝啊,求你記念我,施恩與我.

看來,尼希米,仍是孤獨一人,在黑暗的歷史長河邊,拿著燈籠,為神的子民的將來,憂心忡忡.

來到這裡,我們看見,被擄回歸的猶太人,並沒有珍惜神的恩典.一轉眼,他們便重蹈祖先的罪行.他們漠視聖殿,輕忽安息日,遺忘自己的獨特身分.若不是神超乎想像的眷顧,他們早已在歷史中湮沒了.

這般暗淡的局面,猶太人不忍將這放在他們聖經的結尾,而將《歷代志下》放在結尾,成為最後的經卷.《歷代志下》的結束經文,是波斯的塞魯士王宣佈大赦天下,被擄的猶太人可以回歸,開展他們信仰的新一頁.如此,就能以開展前景的希望,作為他們聖經終結的經文.

我當然希望眾人的歷史的終結,不是蒼白暗淡,而是充滿從天而來的恩典.問題是:我們尊重聖殿中的崇拜嗎?在安息日裡,我們能重視神的臨在嗎?在歷史的巨變中,我們能守護神的子民的身分嗎?

願榮耀歸上主,願眾人享平安.



\newpage

\section{真是我的最佳選擇嗎}
\label{sec:652}
\textbf{高國雄牧師 - 真是我的最佳選擇嗎?}
\newline
\newline
\hyperref[sec:651]{< < < PREV SERMON < < <}
~
\hyperref[sec:toc]{[返目錄]}
~
\hyperref[sec:653]{> > > NEXT SERMON > > >}


馬太福音 10:24-39
\newline
\begin{longtable}{cl}
\hline
\hline
章節 & 經文 (和合本修訂版)\\
\hline
10:24 & \begin{tabularx}{0.7\textwidth}{X} 「學生不高過老師，僕人不高過主人。 \end{tabularx} \\ \\ \relax
10:25 & \begin{tabularx}{0.7\textwidth}{X} 學生所遭遇的與老師一樣，僕人所遭遇的與主人一樣，也就夠了。既然有人罵一家的主人是『別西卜』，更何況他的家人呢？」 \end{tabularx} \\ \\ \relax
10:26 & \begin{tabularx}{0.7\textwidth}{X} 「所以，不要怕他們，因為掩蓋的事沒有不顯露出來的，隱藏的事也沒有不被知道的。 \end{tabularx} \\ \\ \relax
10:27 & \begin{tabularx}{0.7\textwidth}{X} 我在暗中告訴你們的，你們要在明處說出來；你們耳中所聽的，要在屋頂上宣揚出來。 \end{tabularx} \\ \\ \relax
10:28 & \begin{tabularx}{0.7\textwidth}{X} 那殺人身體但不能滅人靈魂的，不要怕他們；惟有那能在地獄裡毀滅身體和靈魂的，才要怕他。 \end{tabularx} \\ \\ \relax
10:29 & \begin{tabularx}{0.7\textwidth}{X} 兩隻麻雀不是賣一銅錢嗎？你們的父若不許，一隻也不會掉在地上。 \end{tabularx} \\ \\ \relax
10:30 & \begin{tabularx}{0.7\textwidth}{X} 就是你們的頭髮也都數過了。 \end{tabularx} \\ \\ \relax
10:31 & \begin{tabularx}{0.7\textwidth}{X} 所以，不要懼怕，你們比許多的麻雀還貴重！」 \end{tabularx} \\ \\ \relax
10:32 & \begin{tabularx}{0.7\textwidth}{X} 「所以，凡在人面前認我的，我在我天上的父面前也必認他； \end{tabularx} \\ \\ \relax
10:33 & \begin{tabularx}{0.7\textwidth}{X} 凡在人面前不認我的，我在我天上的父面前也必不認他。」 \end{tabularx} \\ \\ \relax
10:34 & \begin{tabularx}{0.7\textwidth}{X} 「你們不要以為我來是帶給地上和平，我來並不是帶來和平，而是刀劍。 \end{tabularx} \\ \\ \relax
10:35 & \begin{tabularx}{0.7\textwidth}{X} 因為我來是要叫『人與父親對立，女兒與母親對立，媳婦與婆婆對立。 \end{tabularx} \\ \\ \relax
10:36 & \begin{tabularx}{0.7\textwidth}{X} 人的仇敵就是自己家裡的人。』 \end{tabularx} \\ \\ \relax
10:37 & \begin{tabularx}{0.7\textwidth}{X} 愛父母勝過愛我的，不配作我的門徒；愛兒女勝過愛我的，不配作我的門徒。 \end{tabularx} \\ \\ \relax
10:38 & \begin{tabularx}{0.7\textwidth}{X} 不背自己的十字架跟從我的，不配作我的門徒。 \end{tabularx} \\ \\ \relax
10:39 & \begin{tabularx}{0.7\textwidth}{X} 得著性命的，要喪失性命；為我喪失性命的，要得著性命。」 \end{tabularx} \\ \\
[1ex]
\hline
\hline
\end{longtable}
羅馬書 6:3-11
\newline
\begin{longtable}{cl}
\hline
\hline
章節 & 經文 (和合本修訂版)\\
\hline
6:3 & \begin{tabularx}{0.7\textwidth}{X} 難道你們不知道，我們這受洗歸入基督耶穌的人，就是受洗歸入他的死嗎？ \end{tabularx} \\ \\ \relax
6:4 & \begin{tabularx}{0.7\textwidth}{X} 所以，我們藉著洗禮歸入死，和他一同埋葬，是要我們行事為人都有新生的樣子，像基督藉著父的榮耀從死人中復活一樣。 \end{tabularx} \\ \\ \relax
6:5 & \begin{tabularx}{0.7\textwidth}{X} 我們若與他合一，經歷與他一樣的死，也將經歷與他一樣的復活。 \end{tabularx} \\ \\ \relax
6:6 & \begin{tabularx}{0.7\textwidth}{X} 我們知道，我們的舊人和他同釘十字架，使罪身滅絕，叫我們不再作罪的奴隸， \end{tabularx} \\ \\ \relax
6:7 & \begin{tabularx}{0.7\textwidth}{X} 因為已死的人是脫離了罪。 \end{tabularx} \\ \\ \relax
6:8 & \begin{tabularx}{0.7\textwidth}{X} 我們若與基督同死，我們信也必與他同活， \end{tabularx} \\ \\ \relax
6:9 & \begin{tabularx}{0.7\textwidth}{X} 因為知道基督既從死人中復活，就不再死，死也不再作他的主了。 \end{tabularx} \\ \\ \relax
6:10 & \begin{tabularx}{0.7\textwidth}{X} 他死了，是對罪死，只這一次；他活，是對神活著。 \end{tabularx} \\ \\ \relax
6:11 & \begin{tabularx}{0.7\textwidth}{X} 這樣，你們也要看自己對罪是死的，在基督耶穌裡對神卻是活的。 \end{tabularx} \\ \\
[1ex]
\hline
\hline
\end{longtable}
講題:真是我的最佳選擇嗎?

講員:高國雄牧師

經文:馬太福音10:24-39;羅馬書6:3-11

日期:2023年6月25日

24「學生不高過老師,僕人不高過主人.

25學生所遭遇的與老師一樣,僕人所遭遇的與主人一樣,也就夠了.既然有人罵一家的主人是『別西卜』,更何況他的家人呢?」

26「所以,不要怕他們,因為掩蓋的事沒有不顯露出來的,隱藏的事也沒有不被知道的.

27我在暗中告訴你們的,你們要在明處說出來；你們耳中所聽的,要在屋頂上宣揚出來.

28那殺人身體但不能滅人靈魂的,不要怕他們；惟有那能在地獄裡毀滅身體和靈魂的,才要怕他.

29兩隻麻雀不是賣一銅錢嗎?你們的父若不許,一隻也不會掉在地上.

30就是你們的頭髮也都數過了.

31所以,不要懼怕,你們比許多的麻雀還貴重!」

耶穌基督安慰門徒,不要怕.

弟兄姊妹,在我們的生活中,有多少讓你懼怕的事情呢?我們最懼怕的是什麼呢?

當耶穌向門徒說:「不要怕!」.到底當日門徒正處於什麼狀況呢?因什麼原因而害怕呢?

耶穌接著表達:「因為掩蓋的事沒有不顯露出來的,隱藏的事也沒有不被知道的.」當時是否有一些事情?他們很害怕被顯露/不被顯露.很害怕,有些事情沒有人知道.需要在明處說出來,需要在屋頂上宣揚.

相信正是回應當時門徒最大的擔憂,以及最懼怕的事情.

我們是一群信仰基督的人,我們懼怕宣揚基督嗎?

我們懼怕,因為宣揚基督而受傷害嗎?

甚至如果有一天,因為宣揚基督,而招致殺身之禍,我們要如何自處呢?

V28「那殺人身體但不能滅人靈魂的,不要怕他們；惟有那能在地獄裡毀滅身體和靈魂的,才要怕他.」

地上很多不同的權勢.

歷世歷代也有很多信徒因為信仰的緣故殉道.

今天我們又如何回應這個呼召呢?

早期教會的神學家特土良曾說:「殉道者的血是教會的種子」.

耶穌的安慰:如果父不容許,沒有任何權勢可以打倒我們.

上帝把我們的頭髮也數過,我們在他眼中非常非常貴重,因此主不會輕易讓我們陷在難處當中.

這裡清楚表達:作為信徒,所面對的敵人只有一個,就是那能毀滅身體和靈魂的,就是那敵擋基督的魔鬼撒旦.

32「所以,凡在人面前認我的,我在我天上的父面前也必認他；

33凡在人面前不認我的,我在我天上的父面前也必不認他.」

「在人面前承認信仰」,在平穩安定的日子中,當然好像沒什麼特別.

在開放的社會中,在基督信仰被認同的地方/國家政權中,當然信仰是比較能夠輕鬆直接表達.

公開的認信,公開的洗禮,公開的崇拜…在我們經歷上好像是理所當然的.

但當我們回想,歷代教會所遭遇到的迫害.

當我們環顧,現時世界各地仍然有很多人,沒有辦法自由地表達信仰.

在人面前認主,可以就成為生命中最大的掙扎.

34「你們不要以為我來是帶給地上和平,我來並不是帶來和平,而是刀劍.

35因為我來是要叫『人與父親對立,女兒與母親對立,媳婦與婆婆對立.36人的仇敵就是自己家裡的人.』

很多人形容基督教是一個「愛」的宗教,基督徒更是很多人眼中為世界帶來和平的群體,當然這描述基本上是非常正確的.

上帝愛世人,他的愛遍佈全地.

但今天的經課中耶穌所表達的,卻是另一個向度.

V34「你們不要以為我來是帶給地上和平,我來並不是帶來和平,而是刀劍…」

耶穌要帶來刀劍,為什麼呢?

刀劍要對付的,就是那敵擋上帝的勢力,就是那惡者撒旦.

很多人以為,有基督徒的世界,就會有和平,而那「和平」,就是在地上的和平,就是「和諧」,這和諧就等如沒有爭鬥,沒有衝突.

當然,我不是說我們作為信徒,要刻意引起衝突.但是原來當我們認清這世代的真相.今天這世代,仍然被那迷惑普天下,空中掌權者的首領的影響當中.惡人惡事隨處可見,敵擋基督的事情在我們身旁不斷發生.所以,耶穌說:「他來不是為地上帶來和平,而是刀劍.」而且很清楚地表達,就是在最近的人身上,可能就在我們家人當中.

我們要清晰,這裡耶穌沒有叫我們不去孝順父母,關愛家人.

耶穌要我們知道,正如上一段所講,作為信徒最大的敵人,就只有那惡者撒旦.

當我們發現知道身邊的人是敵擋基督的時間,我們就要立定我們的信仰立場.就算最終要動起刀劍,也不遺餘力.

因為這是根本的問題.

當然實際處理手法很多,但這裡耶穌所指出的是一個本質的問題.

亦是我們作信徒是否堅持信仰,是否把信仰捉得緊的問題.

這是生命中最重要的優先次序問題.

37愛父母勝過愛我的,不配作我的門徒；愛兒女勝過愛我的,不配作我的門徒.

38不背自己的十字架跟從我的,不配作我的門徒.

39得著性命的,要喪失性命；為我喪失性命的,要得著性命.」

所以耶穌好像非常極端地表示:「愛父母勝過愛我的,不配作我的門徒；愛兒女勝過愛我的,不配作我的門徒.」

父母愛兒女,子女愛父母,好像是人性很重要,很基本的價值.

但是愛主,才應該是信徒最重要的焦點.如果要選擇,如果真的需要選擇,我們應該清楚知道我們要如何選擇.

正如聖經所說,不背自己的十字架跟從我的,不配作我的門徒.我們自己的十字架是什麼呢?

如果有一天我們所愛的人與物,成為了敵擋信仰,成為了我們與上主隔離的因由…

為此我們背上了自己的十字架.我們又是否為主願意揹起?

所以最後耶穌提到.

得著性名的,要喪失性命?

為我名喪失性命的,要得著姓名.

真正的生命在哪裡呢?我們看準.我們真的要看得準.

羅馬書6

3難道你們不知道,我們這受洗歸入基督耶穌的人,就是受洗歸入他的死嗎?

4所以,我們藉著洗禮歸入死,和他一同埋葬,是要我們行事為人都有新生的樣子,像基督藉著父的榮耀從死人中復活一樣.

5我們若與他合一,經歷與他一樣的死,也將經歷與他一樣的復活.

6我們知道,我們的舊人和他同釘十字架,使罪身滅絕,叫我們不再作罪的奴隸,

7因為已死的人是脫離了罪.

8我們若與基督同死,我們信也必與他同活,

9因為知道基督既從死人中復活,就不再死,死也不再作他的主了.

10他死了,是對罪死,只這一次；他活,是對上帝活著.

11這樣,你們也要看自己對罪是死的,在基督耶穌裡對上帝卻是活的.

信耶穌,洗禮加入教會,就是歸入基督,歸入他的死,埋葬,復活.

同死,同活.

對罪死,對主活.

這就是我們生命中應做的選擇!

【面對抉擇】

1)環境諸多變遷,誰可以看見,明天的發展和際遇?

人智慧有限,難作選擇,前路一片茫茫如沒路牌!

2)人心思可變更,誰可以看透,人心中的構思和計劃?

人你我互謀,難作審斷,前路一片茫茫如沒路牌!

副歌:

面對抉擇面臨困惑,心底裡幾番衝動:

人去亦難留更難,常退避沒有用!

今天我有主耶穌的足印,給我作指引!

面對艱難我誓要跟從,原因主智慧無限.

實在我不配為我作選擇,原因我已屬上主所有.



\newpage

\section{命運就算顛沛流離}
\label{sec:653}
\textbf{蒲錦昌牧師 - 命運就算顛沛流離}
\newline
\newline
\hyperref[sec:652]{< < < PREV SERMON < < <}
~
\hyperref[sec:toc]{[返目錄]}
~
\hyperref[sec:654]{> > > NEXT SERMON > > >}


講題:命運就算顛沛流離

經文:路得記1:1-22,4:13-22

講員:蒲錦昌牧師

日期:2023年7月2日

一,前言:

我們讀歷史書往往看到大人物如帝王將相等對歷史的影響,小人物卻常常被忽略.究竟他們對歷史的走向是否有影響力呢?今天我們讀的路得記,只記載一個家庭的故事:兩婆媳拿俄米和路得移民和返鄉的故事,再加上一個配角波阿斯.為什麼這卷書會被放進聖經之中呢?她們的故事有什麼價值和對我們今天的生活有什麼啟示呢?現在讓我們來看看.

二,主題:

1.

顛沛流離的時代

得1:1介紹拿俄米和路得活在「士師統治的時候」.我們讀士師記後段,大概便會了解 拿俄米和路得時的以色列人社會是怎樣一個狀況.士師記到中段士師基甸便越來越差,他的兒子亞比米勒把兄弟七十人全部殺害以奪得統治以色列人的地位,殘忍而暴戾.士師參孫則只按著自己的喜好,慾望和情緒行事,絲毫沒有政治領袖的特質.到士師記

19章記載一個利未人和他的妾侍在路經便雅憫支派的地方,竟被同為以色列人的便雅憫人暴力對待,妾侍也死掉.這個利未人把妻子的屍體切成十二塊,分別送到以色列全地各支派中,結果全部以色列人聚集向便雅憫人大興問罪之師,差不多把他們滅族,最後又擔心十二支派從始只剩十一個,於是為便雅憫男人強搶基

列雅比和示羅的女子為妻.這就是士師記的結束.這是多麼混亂的一個場面!

士師記的最後一節說:「那時,以色列中沒有王,各人照自己眼中看為對的去做.」(士 21:25)換而言之,這是一個十二支派同盟政治制度崩潰的時代,一個價值顛倒的時代,一個暴力充斥的時代,一個人和人互相難以信任的時代,也是一個不平靜和顛沛流離的時代.

2.

苦命人拿俄米

路得記開首部分(1:1-5)記載一個伯利恆家庭的故事.因著飢荒,丈夫以利米勒帶著妻子拿俄米和兩個兒子瑪倫和基連到摩押地生活,兩個兒子分別娶了摩押女子俄珥巴和路得為妻.但是,三個男

人後來都死了,剩下三個女人,並無子嗣.拿俄米三人都變成了寡婦,一無所有,只有絕望.從拿俄米後來的話推測:「我滿滿地出去,耶和華使我空空地回來.」(得 1:21上)大概他們原來是一個比較富裕的家庭,屬於有經濟條件移民者,但是,現在卻變得一無所有.

當拿俄米知道家鄉情況好轉,而她也沒有留戀摩押地的原因,於是便決定起程回鄉.

按照那時的習俗,家庭中的長輩是有責任照顧後輩的,而後輩則需要服從長輩的按照那時的習俗,家庭中的長輩是有責任照顧後輩的,而後輩則需要服從長輩的教導和安排.拿俄米既然沒法給兩個媳婦帶來什麼希望和好日子,於是便叫她們教導和安排.拿俄米既然沒法給兩個媳婦帶來什麼希望和好日子,於是便叫她們不要再跟著她,找一不要再跟著她,找一主主好人家改嫁便算.俄珥巴見拗好人家改嫁便算.俄珥巴見拗不過婆婆,唯有含淚和拿俄不過婆婆,唯有含淚和拿俄米吻別,只有路得堅決跟從拿俄米回伯利恆.米吻別,只有路得堅決跟從拿俄米回伯利恆.

當拿俄米回到家鄉伯利恆時,已是十多二十年之後,大概她潦倒和衰老的樣子令當拿俄米回到家鄉伯利恆時,已是十多二十年之後,大概她潦倒和衰老的樣子令全城的婦女們都吃驚,不禁相問:「這個人真的是拿俄米嗎?」(得全城的婦女們都吃驚,不禁相問:「這個人真的是拿俄米嗎?」(得1:19)拿俄米)拿俄米的回答反映了她對自己人生的評價:的回答反映了她對自己人生的評價:

「不要叫我拿娥米,

要叫我瑪拉,

因為全能者使我受盡了苦.

我滿滿地出去,

耶和華使我空空地回來.

耶和華使我受苦,

全能者降禍於我.

你們為何還叫我拿娥米呢?」(得1:20-21)

她不再想別人叫她拿俄米(希伯來文是甜的意思),要人叫她瑪拉(希伯來文是苦的意苦的意思).她的生命太苦了!不單只這樣,她還認為自己有有這樣的這樣的遭遭遇遇是上帝是上帝有意降禍予她,令她人生這麼苦.有意降禍予她,令她人生這麼苦.

人間充滿了痛苦,包括令人失望的環境,戰爭的殘酷,醜陋的人性,不幸的遭遇等等.有時候,我們都和拿俄米一樣,覺得父母為我們改了一個好名字,但人生等等.有時候,我們都和拿俄米一樣,覺得父母為我們改了一個好名字,但人生卻是另外一回事.不少人都和拿俄米有同感,都覺得自己是苦命人,被上帝和命卻是另外一回事.不少人都和拿俄米有同感,都覺得自己是苦命人,被上帝和命運捉弄.運捉弄.

3.

同行者路得同行者路得

路得記另一個主角是路得.在路得記的敘述中,經常連著「摩押女子」這個稱號來叫她.根據申命記23章章3至至4節說:「亞捫人或摩押人不可入耶和華的會；甚至到第十代,也永不可入耶和華的會.至到第十代,也永不可入耶和華的會.

因為你們出埃及的時候,他們沒有拿食物和水在路上迎接你們,並且雇了美索不達米亞的毗奪人,比珥的兒子巴蘭來詛咒你.」換句話說,以色列人視摩押人為世仇,摩押人永遠也不能入籍成為以色列人一份子,因為以色列人在由埃及往迦南建立民族和國家的歷史中,摩押人並沒有歡迎他們的到臨.那麼,可想而知,路得陪伴拿俄米在伯利恆生活,將會遇上多大的歧視和艱難!

路得在面對前路要作出抉擇的時候,她說了一番驚天動地的話:

「不要勸我離開你,

轉去不跟隨你.

你往哪裡去,

我也往哪裡去；

你在哪裡住,

我也在哪裡住；

你的百姓就是我的百姓；

你的上帝就是我的上帝.

你死在哪裡,

我也死在哪裡,葬在哪裡.

只有死能使你我分離；

不然,願耶和華重重懲罰我!」(得1:16-17)

她以自己的性命起誓,願意陪伴拿俄米到天涯海角,而且至死不渝.她的誓言比很多戀人還要堅定和偉大.

當我們閱讀短短四章的路得記時,我們會發現,路得對拿俄米的話往往言聽計從,而且是拿俄米先說而路得作出回應.路得有古代傳統女性的美德:服從.但是,她也作過一次主動,在路得記二章一開始,路得便向拿俄米建議:「讓我到田裡去拾取麥穗,我在誰的眼中蒙恩,就跟在誰的身後.」(得2:2)大概她們二人的生活已捉襟見肘,不能不工作了.拿俄米會不會對這麼丟臉的狀況感到難堪,也難以開口叫路得做最卑下的工作?而路得為了生活,為了她婆婆的三餐,竟然主動的表示願意做以色列人社會中最卑賤的工作:跟在奴僕後面拾取田間剩餘的麥穗.

還有,路得這樣做不單只會被人看輕,而且還要承受一定的風險,就像古代沒有男人保護的女人會被他人佔便宜一樣.男人保護的女人會被他人佔便宜一樣.有有田田地地的的富人波阿斯勸告路得:「女兒啊,聽我說,不要到別人田裡去拾取麥穗,也不要離開這裡,要緊跟著我的女僕們.你要看好我的僕人正在哪塊田收割,就跟著女僕們去.我已經吩咐僕人不可侵犯你.」(路2:8-9)又或者拿俄米提醒路得要格外留心:「女兒啊,你要跟著他的女僕出去,免得你在別人的田間受人騷擾.」(路2:22)原來無依無靠的寡婦和外邦人,尤其是受到鄙視的外邦人,是很容易被他人騷擾和侵犯的.

我們想一下,路得的遭遇和拿俄米其實很相似:她也是寡婦,也像拿俄米一樣失去丈夫.路得選擇陪伴拿俄米回她的家鄉伯利恆,她也要要在異地居住,變成異鄉人.至於她在伯利恆初期的生活,可能比拿俄米初到摩押時更糟.不過,在拿俄米喊苦的時候,我們卻沒有在聖經中聽到路得喊苦的聲音.為什麼呢?

我想先跟大家跟大家講一個故事.在二戰時的集中營裡,住了一個心理學家叫維克多．法蘭克(Victor E. Frankl).他在集中營裡發現一件很奇怪的事情:不少年青和身體強壯的人捱不過集中營艱難的日子,反而有些老弱的捱過了.當他研究這些成功捱過的人,他發現他們都有一個共通點:他們都有一些重要的事要做,不論是與未婚妻結婚,又或者有分散了的親人等待重聚.而法蘭克自己也有一件重要的事等著他去完成:他冒險收藏和保存自己著作的手稿,希望有一天可以出版.他用這個發現來開創一種新的心理治療法,叫意義治療法,以幫助有心理問題者找到生存下去的動力.他有一句名言:「人只要找到意義,就什麼都可以承受.」

路得大概以照顧婆婆拿俄米為她最重要的人生目標,因而有動力面對困難,並且能忘記自己的痛苦,超越自己的煩惱,專注於他人身上.她無疑有傳統女性的犧牲精神,處處為他人設想.

各位弟兄姊妹,我們的人生可能也有不少苦,像拿俄米一樣,會覺得世界對自己不公平,甚至認為是上主苦待自己.但是,請大家想一想,留意一下周圍發生的事,會不會在你的家人,朋友,同事或弟兄姊妹中間,找到一個像路得的同行者?他/她雖然和你的處境相同,不過,他/她不單只沒有埋怨,反而默默的與你同行.如果我們身邊找不到像路得這樣的人,我們是否可以想想,自己是否可以成為路得一樣的同行者,與痛苦的人同行,使他們在痛苦中嘗到一點的甜蜜?你也可以得一樣的同行者,與痛苦的人同行,使他們在痛苦中嘗到一點的甜蜜?你也可以因著找到人生目標和意義而提升自己面對痛苦的能力.

4.

上主的作為上主的作為

在短短四章的路得記中,上主並沒有出現,只有被人提起及談及.

上主究竟在那裡呢?他有沒有參與和影響人類的歷史呢?他叫拿俄米受苦嗎?是他安排路得這個外邦女子成為拿俄米的同行者,改變了拿俄米的命運嗎?是他安排路得這個外邦女子成為拿俄米的同行者,改變了拿俄米的命運嗎?是他安排路得遇上波阿斯,使他們成婚和有後嗎?正如傳道書所講:上主的作為,人不能參透,尤其是在當時.

當我們跳過中間的部分直接看路得記四章時,我們發覺路得嫁了富有的好人波阿斯,而且生子俄備得,更成為後世大衛王的祖先.多麼神奇!上主藉兩個藉藉無名的婦女,改變了以色列整個民族的命運.由士師時代分裂的十二支派,到統一的王國；由崩壞的制度及社會環境,到新的政治制度和人民生活.原來在最難堪的時代,上主仍然賜下人間的溫暖和愛,使我們感受到生命的價值和意義.在二戰的集中營裡,我們仍然看到在極度飢餓中願意與人分享食物的人,顯示出人類的尊貴和有上主的形象.上主在亂世中仍為人開出路,使人類不至絕望.路得記以死亡開始,以新生結束；以絕望開始,以希望結束,正正說明這樣一個道理.

我們雖然不能完全明白上主的作為,但是,讓我們以信心倚靠他,深信他必成就他在歷史中的美意,因為他是人類的同行者.上主差遣愛子耶穌基督道成肉身,成為人類的一份子,飽嘗人間的痛苦,就是要與人同在同行,又為我們的生命開新的出路.他的名是以馬內利(上主與我們同在)!

三,結語:

最後,請讓我以香港歌星李克勤填詞的「紅日」來結束今日的信息:

「命運就算顛沛流離,

命運就算曲折離奇,

命運就算恐嚇著你做人沒趣味,

別流淚心酸,更不應捨棄,

我願能一生永遠陪伴你.

一生之中兜兜轉轉那會看清楚,

徬徨時我也試過獨坐一角像是沒協助,

在某年那幼小的我,

跌倒過幾多幾多落淚在雨夜滂沱.

一生之中彎彎曲曲我也要走過,

從何時有你有你伴我給我熱烈地拍和,

像紅日之火燃點真的我,

結伴行千山也定能踏過.」

在充滿苦難的人生中,上主為我們預備了同行者,使我們有力量走下去.他又讓我們成為別人的同行者,使我們的生命活得更加有意義和充實.盼望我們都能夠結伴同行,在人生的旅途中看到希望,嘗到甘甜!



\newpage

\section{識見僕人柔和謙卑}
\label{sec:654}
\textbf{李均熊牧師 - 識見僕人．柔和謙卑}
\newline
\newline
\hyperref[sec:653]{< < < PREV SERMON < < <}
~
\hyperref[sec:toc]{[返目錄]}
~
\hyperref[sec:655]{> > > NEXT SERMON > > >}


講題:識見僕人．柔和謙卑

講員:李均熊牧師

經文:創世記24

日期:2023年7月9日

引言:我的軛是容易的

今天我都會集中講舊約的經課,但在我進入舊約經課之前,先讓我解釋一下福音經課中耶穌說他的軛是「容易」的是甚麼意思.軛是套在牛身上用作耕種的器具,可能大家都了解,但怎樣謂之「容易」的軛呢?其實「容易」應該是指「合適」(fit),試想像作為木匠之子的耶穌,他其實是在推銷自己所做的軛,為牛度身訂造,配合牛身軀線條,所以牠們拉起這個合適的軛時,就很「容易」發力,用起上來「輕省」,也不會因不合適而刮傷牛身.跟從耶穌,耶穌就會將合適的軛加在我們身上,我們雖然承擔使命,作祂的僕人,但卻並非受勞役,被驅使,而是輕省的.

起:亞伯拉罕僕人最大的難題是甚麼?

今日的舊約經課,就是講到亞伯拉罕要求他的僕人負起一個對他合適的軛,就是為他的兒子以撒回到他的「鄉下」去「討老婆」:

37我主人叫我起誓說:『不要為我兒子娶我所居住的迦南地的女子為妻.38你要往我父家,我本族那裡去,為我的兒子娶妻.』(和修本)

經課沒有交代的,是這個僕人如何找到現在這個利伯加的呢?我們可以看看創24:10-22的記載:

10那僕人從他主人的駱駝中取了十匹駱駝,他手中也帶著他主人各樣的貴重物品離開,起身往美索不達米亞去,到了拿鶴的城.11傍晚時,眾女子出來打水,他就讓駱駝跪在城外的水井旁.12他說:「耶和華—我主人亞伯拉罕的神啊,求你施恩給我的主人亞伯拉罕,讓我今日就遇見吧!13看哪,我站在井旁,城內居民的女子們正出來打水.14我向哪一個少女說:『請你放下水瓶來,給我水喝』,她若說:『請喝!我也給你的駱駝喝』,願她作你所選定給你僕人以撒的妻.這樣,我就知道你施恩給我的主人了.」

15話還沒說完,看哪,利百加肩頭上扛著水瓶出來.利百加是彼土利所生的；彼土利是亞伯拉罕的兄弟拿鶴妻子密迦的兒子.16那少女容貌極其美麗,是未曾與人親近的童女.她下到井旁,打滿了瓶子的水,就上來.17僕人跑上前去迎著她,說:「請你讓我喝你瓶子裡的一點水.」18少女說:「我主請喝!」就急忙拿下瓶子托在手上,給他喝水.19那少女給他喝足了,又說:「我也為你的駱駝打水,直到駱駝喝足了.」20她就急忙把瓶子裡的水倒在槽裡,又跑到井旁打水,為所有的駱駝打了水.21那人定睛看著少女,一句話也不說,要知道耶和華是否使他的道路亨通.

22駱駝喝足了,那人就拿出一個比加重的金環,一對十舍客勒重的金手鐲,23說:「請告訴我,你是誰的女兒?你父親家裡有沒有地方可以讓我們過夜?」24少女說:「我是密迦為拿鶴生的兒子彼土利的女兒.」25又說:「我們家裡有充足的乾草和飼料,也有住宿的地方.」26那人就低頭向耶和華敬拜,27說:「耶和華—我主人亞伯拉罕的神是應當稱頌的,因他不斷以慈愛信實待我主人.至於我,耶和華一路引領我,直到我主人的兄弟家裡.」(和修本)

看來事情在上主的引導下十分順利,這個僕人用上一個「考驗」,試出這個勤奮好客的女子,最後竟是亞伯拉罕兄弟的孫女兒利伯加!

不過,但若我們回看亞伯拉罕的要求是甚麼:

1亞伯拉罕年紀老邁,耶和華在一切事上都賜福給他.2亞伯拉罕對他家中管理他一切產業最老的僕人說:「把你的手放在我大腿底下.3我要叫你指著耶和華—天和地的神起誓,不要為我兒子娶我所居住的迦南地的女子為妻.4你要往我的本地本族去,為我的兒子以撒娶妻.」5僕人對他說:「如果那女子不肯跟我來到這地,我必須把你的兒子帶回到你出來的地方嗎?」6亞伯拉罕對他說:「你要謹慎,不可帶我兒子回那裡去.7耶和華—天上的神曾帶領我離開父家和本族的地,對我說話,向我起誓說:『我要將這地賜給你的後裔.』他要差遣使者在你面前,你就可以從那裡為我兒子娶妻.8倘若那女子不肯跟你來,我叫你起的誓就與你無關了,只是你不可帶我的兒子回到那裡去.」9僕人就把手放在他主人亞伯拉罕的大腿底下,為這事向他起誓.(和修本)

如何令利伯加自己(如亞伯拉罕一般?)離開家庭「下嫁」自己的小主?這才是亞伯拉罕的僕人真正的困難,也是這一個老僕人合適這個「軛」的地方.當我們十分心急想知道利伯加的反應,中途卻殺出她的哥哥:

28那少女跑去,把這些話告訴她母親家裡的人.29利百加有一個哥哥,名叫拉班,拉班就跑到外面井旁那人那裡.30當他看見金環和戴在他妹妹手上的金鐲,又聽見他妹妹利百加說的話:「那人如此對我說」,他就來到那人面前,看哪,他還站在井旁的駱駝旁邊,31就對他說:「你這蒙耶和華賜福的人,請進來吧!為甚麼站在外面?我已經收拾了房屋,也為駱駝預備了地方.」32那人就進了拉班的家.拉班卸了駱駝,用飼料餵牠們,拿水給那人和隨從他的人洗腳,33把食物擺在他面前,請他吃.他卻說:「我不吃,等我把我的事情說完了再吃……(和修本)

利伯加的「跑」——看來她也是樂意(eagar)的吧?她不是問題吧?但拉班的跑:因為他看到金環金鐲和她妹妹的說話(事實上她講了什麼?),見到(和修版有原文的「看哪!」,強調拉班即時見到的事情)D 駱駝(眾數)【參和修本】!蒙福(是否只看到物質的豐富),對駱駝的「殷勤」——拉班會不會只是見錢開眼,甚至會否當妹妹是「搖錢樹」?看來拉班才是這個僕人的對手.(我們後來對拉班的認識,就知他「數口精」)

承:僕人的「說詞」:他如何覆述事情始末

面對這個可能見錢開眼的拉班,亞伯拉罕必須在不受任何他的「好處」之前,道明來意,所以要「把事情說完了再吃」.首先,這個僕人投其所好,介紹自己主人時並不是說上主要祝福他子孫繁多,或別人會因他得福等等,而是純物質的:

「我是亞伯拉罕的僕人.35耶和華大大地賜福給我主人,使他發達,賜給他羊群,牛群,金銀,奴僕,婢女,駱駝和驢.36我主人的妻子撒拉年老的時候為我主人生了一個兒子；我主人把他一切所有的都給了他.

而且那個將會成為利伯加丈夫的人,不單是主人年老所生的心干寶貝,更已承受主人的一切.然候,僕人再說:

37我主人叫我起誓說:『不要為我兒子娶我所居住的迦南地的女子為妻.38你要往我父家,我本族那裡去,為我的兒子娶妻.』39我對我主人說:『恐怕那女子不肯跟我來.』40他就說:『我所事奉的耶和華必要差遣他的使者與你同去,使你的道路亨通,你就可以在我父家,我本族那裡,為我的兒子娶妻.41只要你到了我本族那裡,我叫你起的誓就與你無關.他們若不把女子交給你,我叫你起的誓也與你無關.』

當中,僕人將原本與主人的對話,作出微妙的改動.原本亞伯拉罕叫僕人去自己的「出生地」(moladeti中文譯本族)去為以撒娶妻,但僕人向拉班卻兩次說他主人要他回父家和族人(mispahti中文都是譯本族)那裡去,令拉班意識到這次僕人回鄉娶妻並不是家家戶戶也合適.當然他主人和他所說關於上主怎樣帶領他離開父家和本族和向他起誓這些細節省略,免得令拉班誤會亞伯拉罕的神怎麼只揀選亞伯拉罕而不是他的家族等等事情而節外生枝.

然候,他就覆述他如何遇到他的妹妹利伯加:

42「我今日到了井旁,就說:『耶和華—我主人亞伯拉罕的神啊,願你使我所行的道路亨通.43看哪,我站在井旁,對哪一個出來打水的女子說:請你讓我喝你瓶子裡的一點水,44她若說:你只管喝,我也為你的駱駝打水；願那女子作耶和華給我主人兒子所選定的妻子.』

45「我心裡的話還沒有說完,看哪,利百加肩頭上扛著水瓶出來,下到井旁打水.我對她說:『請你給我水喝.』46她就急忙從肩頭上拿下瓶子來,說:『請喝!我也給你的駱駝喝.』我就喝了；她也給我的駱駝喝了.47我問她說:『你是誰的女兒?』她說:『我是彼土利的女兒,彼土利是密迦和拿鶴生的兒子.』我就把環子戴在她鼻子上,把鐲子戴在她雙手上.48然後我低頭向耶和華敬拜,稱頌耶和華—我主人亞伯拉罕的神,因為他引導我走合適的道路,使我得著我主人兄弟的孫女,給我主人的兒子為妻.

若你對比之前經文記載的細節,你會發現亞伯拉罕僕人的覆述其實是要說服拉班,整件事並非任何人為刻意安排,而是純粹上主的引導.僕人自然沒有說他首先見到一個女子「容貌極其俊美」,又看來還未嫁人才上去問話,又或者自己如何不確定,要觀察這個女子的言行等等,這些僕人思量考據的細節略掉；而是心裡的話還沒說完,就看到利伯加,讓拉班感到利伯加就是上主指派唯一選擇.向著拉班交待,僕人當然是先問清楚這個女子姓甚名誰,是誰家姑娘,肯定是親戚上頭,才將貴重的鼻環手鐲送給她,而不是心急送禮之後才問人家是誰.最後,僕人在拉班面前所覆述他對上主的感恩的禱告,並非單單上主在路上引領自己找到主人兄弟的家,要不然為甚麼自己不直接走到主人的兄弟家中,而要在井旁攪這麼多小動作,而是「上主引領我使我得著我主人兄弟的孫女,給我主人的兒子為妻」.換句話說,拉班可要知道,是上主引導我找到利伯加,而不是找到你們這一家.僕人的覆述,更加側重上主的帶領和引導,而沒有一丁點兒是為測試主人的未來新袍的為人如何.

轉:識見僕人——靠上主,也要發揮識見

最後,亞伯拉罕的僕人發出一個看似一般,卻是令人無可推搪的要求:

49現在你們若願以慈愛誠實待我主人,就告訴我；若不然,也告訴我,使我可以或向左,或向右.

比起中性的說:若個你同意(if you consent)……,這個要求一方面期望拉班一如上主般恩待他的主人,就如上主在引導他找到利伯加一樣；另一方面,若果拉班不同意,就是不願以慈愛誠實待我家主人而便變得無情無義了!最後一句「若不然,也告訴我,使我可以或向左,或向右.」用廣東話翻出來就是「你無誠意ge 就出聲,唔好拖住我,我早D 搵個第二個啦(非誠勿擾,我就好另謀高就).」

都這麼說了,拉班那可以拒絕?

50拉班和彼土利回答說:「這事乃出於耶和華,我們不能向你說好說歹.51看哪,利百加在你面前,可以將她帶去,照著耶和華所說的,給你主人的兒子為妻.」

聽到這個回答,僕人並沒有好像剛遇到利伯加般,急急送出禮物.面對可能見錢開眼的拉班,僕人要讓拉班知道現在不是做買賣,而是一切的確出於上主,於是先向上主俯伏祈禱,然候才將金器,銀器,衣服,相信是作為聘禮送給利伯加,然候才將「寶物」(貴重禮物)送給拉班和利伯加母親,然候才吃喝.

過了一晚,僕人便要帶著利伯加回到主人身邊.但拉班卻要留他們多十日,也不知拉班是否想要多些著數,最後只有問問利伯加.最後,利伯加爽快的答「我去」,可見利伯加最終也能自主命運,才沒有夜長夢多,亞伯拉罕吩咐僕人至關重要的事情,也就真的辦妥.

整件事情的記載,讓我們親身經歷上主是如何的引導亞伯拉罕的僕人,這個僕人也處處顯出他如何敬虔地依靠上主.但在關鍵時刻,這個僕人也隨機應變,發揮他的識見和智慧.面對拉班這個對手,他非常巧妙地使拉班不得不承認這是上主的工作,而非利益考慮的婚嫁交易.亞伯拉罕給這個老練的僕人的軛,確是合適的軛,容易而輕省.

合:柔和謙卑,負主的軛

今日舊約經課,正好讓我們明白崇基學院神學院常說所要培養的識見僕人的意思.上主並非直接干預這個世界,祂創造我們,使我們和他一同合作管理這個大地.而基督徒的生活,就是一種神人協作的生活.我們要學主耶穌的樣式,就是柔和謙卑,就是不靠自己,而學習每事都要依靠上主,並且進到祂的安息中；但我們也不忘自己的身份角色,要好好培養我們的識見,好好負起主耶穌的合適,容易而輕省的軛,讓主可以使用我們作祂的僕人,成就祂的旨意.



\newpage

\section{荒謬而諧趣}
\label{sec:655}
\textbf{鄧瑞強博士 - 荒謬而諧趣}
\newline
\newline
\hyperref[sec:654]{< < < PREV SERMON < < <}
~
\hyperref[sec:toc]{[返目錄]}
~
\hyperref[sec:656]{> > > NEXT SERMON > > >}


講題:荒謬而諧趣

經文:以斯帖記1:1,10-22,2:16-17,21-23,3:1-6,4:7-16,5:12-14,6:1-9,12-13;7:1-10

講員:鄧瑞強博士

日期:2023年7月16日

基督徒很習慣一個嚴厲的上帝,很不習慣一個搞笑的上帝.上帝扮法官,信徒覺得很自然.上帝扮小丑,行一步,跌兩步,信徒總覺得格格不入.靈修大師盧雲(Nouwen)寫了一本小書,書名是《羅馬城的小丑戲》(Clowning in Rome).有一次,盧雲住了在羅馬一段較長時間,在這座天主教總部所在的城巿裡,盧雲參觀了很多宏偉的教堂,也見過很多教會高層,然後,他在這本小書的＜前言＞中說:羅馬這個大城巿,像一個大馬戲團,最吸引他的,不是馴獸師或空中飛人,而是小丑.在這馬戲團裡,像馴獸師或空中飛人等這些本領高強的人,能不斷製造種種精神緊張.唯有小丑,卻讓人在緊張中綻放笑容.就是這笑容,讓平凡的人,在困局中,不失希望.

今日,我們看《以斯帖記》.基督徒習慣將《以斯帖記》看成是嚴肅的歷史作品,但在希伯來聖經裡,《以斯帖記》是與詩歌智慧文學連在一起的,又特別與《雅歌》,《路得記》,《耶利米哀歌》,《傳道書》合成一獨特組合,我們稱之為「五小卷」(megilloth).五小卷都和節日有關.《以斯帖記》用來解釋一個節日(普珥節)的起源.「普珥節」是一個狂歡的節日,在這節日裡,猶大人會上演一些搞笑的話劇,也會做很多破格的事情,例如:要醉酒,醉到剛好不能正常運用理性的程度,好能用心去直接感受人間的歡樂.在希伯來聖經的一本具有學術嚴謹性的英譯本《The Jewish Study Bible》裡,在《以斯帖記》導言的第2段便說:《以斯帖記》最好以喜劇讀之.香港人越來越不懂得笑,越來越失去笑的本能,每天都神經兮兮.今日,我們看看《以斯帖記》這齣喜劇,它教我們如何在一「不好笑」的世界裡笑.

(第一幕)亞哈隨魯王和他的佞臣

1:1亞哈隨魯作王,從印度直到古實,統管一百二十七省.

「古實」,即埃塞俄比亞.波斯帝國的版圖,由南亞到東非,是當時最大的帝國.亞哈隨魯王,管治這麼大的帝國,我們想,他應是個很精明的人物.

1:10第七日,亞哈隨魯王飲酒,心中快樂,就吩咐在他面前侍立的七個太監……

1:11請王后瓦實提頭戴王后的冠冕到王面前,使各等臣民看她的美貌,因為她容貌甚美.

1:12王后瓦實提卻不肯遵太監所傳的王命而來,所以王甚發怒,心如火燒.

1:13-14那時,在王左右常見王面,國中坐高位的,有波斯和米底亞的七個大臣……..都是達時務的明哲人.按王的常規,辦事必先詢問知例明法的人.王問他們說:

1:15「王后瓦實提不遵太監所傳的王命,照例應當怎樣辦理呢?」

亞哈隨魯王和群臣尋歡作樂,真是「沒什麼好玩的了」,叫深宮裡的王后出來,讓人欣賞.王后不出來,皇帝氣得跳起來,問身邊的佞臣,如何辦理?

這群「笨蛋」的回答真的笑死人,大家當心笑死.我將經文以意譯來讀.

1:16米母干…說:「王后這樣做,會害死我們這些男人.

1:17我們的老婆,若知道皇帝叫王后來,王后都敢不來,以後,我們這些男人還有站的位置嗎?

1:18整個波斯帝國,若個個女人學王后這樣做,則會「亂套了」.

1:19故此,皇帝啊,要振一下夫綱,不能讓王后想不來就不來,而是著她永遠不用來.之後,找個聽話的人做老婆.

1:20如此,全國所有老婆都不敢不聽老公話.」

1:21王和眾首領都以米母干的話為美,王就照這話去行,

1:22於是,在全國裡,男人以後在家中話事,說話說了算.

故事一開場,便見到土包一般的亞哈隨魯,一個反應古怪的王后,一群笑死人的妄臣.在歡宴與華麗背後,我們看到的,是無比的空虛.偉大的波斯帝國背後,是何等的荒誕.

(第二幕)新王后以斯帖,末底改

亞哈隨魯王挑選新王后,選中以斯帖.

2:16亞哈隨魯王第七年十月……,以斯帖被引入宮見王.

2:17王愛以斯帖過於愛眾女…….王就把王后的冠冕戴在她頭上,立她為王后……

跟著,是引介以斯帖的堂兄末底改的事蹟,他有恩於亞哈隨魯王,這是故事的重大伏筆.

2:21當那時候,末底改坐在朝門,王的太監中有兩個守門的,辟探和提列,惱恨亞哈隨魯王,想要下手害他.

2:22末底改知道了,就告訴王后以斯帖.以斯帖奉末底改的名,報告於王；

2:23究察這事,果然是實,就把二人掛在木頭上,將這事在王面前寫於歷史上.

(第三幕)末底改得罪哈曼

跟著的經文出現《以斯帖記》的一個大壞蛋,名叫哈曼.猶大人在普珥節讀《以斯帖記》時,一讀到哈曼,便會蹬腳或發出噓聲,好掩沒哈曼的名字.

3:1這事以後,亞哈隨魯王抬舉亞甲族哈米大他的兒子哈曼,使他高升,叫他的爵位超過與他同事的一切臣宰.

3:2在朝門的一切臣僕都跪拜哈曼,因為王如此吩咐；惟獨末底改不跪不拜.

3:3在朝門的臣僕問末底改說:「你為何違背王的命令呢?」

3:4他們天天勸他,他還是不聽

……

因他已經告訴他們自己是猶大人.

3:5哈曼見末底改不跪不拜,他就怒氣填胸.

3:6他們已將末底改的本族告訴哈曼；他以為下手害末底改一人是小事,就要滅絕亞哈隨魯王通國所有的猶大人,就是末底改的本族.

哈曼是亞甲族人.亞甲族是何許人?亞甲族是亞瑪力人亞甲王的後人.當摩西帶領以色列人出埃及時,在曠野,出來截擊以色列人的第一個民族,就是亞瑪力人.亞瑪力人是攔阻神的子民得救贖的邪惡代表.《出埃及記》17:14,16記載,神說:「我要將亞瑪力的名號從天下全然塗抹了」,經文又說:「耶和華已經起了誓,必世世代代和亞瑪力人爭戰.」後來,在《撒母耳記上》15章記載,掃羅王與亞瑪力人的亞甲王爭戰時,放過了亞甲王,就是這事件,令神最終放棄掃羅而另立大衛.亞瑪力人亞甲族的哈曼在《以斯帖記》出現,代表著攔阻神的子民得救贖的邪惡代表又再次出現.這就是驚悚片的橋段,不知何解,殭屍或暴龍總在身邊忽然出現.

哈曼要所有人向他叩頭,偏偏末底改不跪不拜,他的理由是:他是猶大人.用這個民族身分去面對亞瑪力人哈曼,劇情便變成「神的子民」與敵擋神的「邪惡子民」的相遇.

面對末底改不肯叩頭,哈曼的反應是誇張的.他不單要殺害末底改,也不單要殺害末底改一家,更要殺害由印度到埃塞俄比亞的所有猶大人.真誇張!這誇張的情節讓我們看到,卑微的猶大民族,要在這龐大的波斯帝國裡生存,保存自己的身分,並不是一件易事.

(第四幕)末底改求救於堂妹以斯帖王后

4:7末底改將自己所遇的事

……

4:8又將……要滅絕猶大人的旨意交給哈他革,要給以斯帖看……,並囑咐她進去見王,為本族的人在王面前懇切祈求.

4:9哈他革回來,將末底改的話告訴以斯帖；

4:10以斯帖就吩咐哈他革去見末底改,說:

4:11「王的一切臣僕和各省的人民都知道有一個定例:若不蒙召,擅入內院見王的,無論男女必被治死；除非王向他伸出金杖,不得存活.現在我沒有蒙召進去見王已經三十日了.」

4:12人就把以斯帖這話告訴末底改.

4:13末底改託人回覆以斯帖說:「你莫想在王宮裡強過一切猶大人,得免這禍.

4:14此時你若閉口不言,猶大人必從別處得解脫,蒙拯救；你和你父家必致滅亡.焉知你得了王后的位分不是為現今的機會嗎?」

4:15以斯帖就吩咐人回報末底改說:

4:16「你當去招聚書珊城所有的猶大人,為我禁食三晝三夜,不吃不喝；我和我的宮女也要這樣禁食.然後我違例進去見王,我若死就死吧!」

末底改求救於堂妹以斯帖王后,催促她去見王,求王解救.有幾點值得留意.一,王后是必須行動的,沒權選擇的.若她向王求情,還有生存機會；若不行動,所有猶大人必死,王后也自身難保.二,不知何解,末底改卻認為猶大人必蒙解救.全卷《以斯帖記》沒有提過耶和華或神的名字,沒有提過神的任何救贖事件,也未提及過神與其子民的約.末底改在書中只有民族意識,絲毫沒有半點信仰內涵,奇怪的是,就算當事人對神毫無意識,讀者卻又好像看到神的手無處不在.三,以斯帖要求猶大人為她禁食三天,這看似是一宗教行動,但事實是這禁食沒有伴隨應有的尋求耶和華的祈禱和舉動.學者David

Lambert說,以斯帖的禁食是一種「喜劇性的禁食」(comic

fast),提供一種由悲情轉化為歡笑的轉折點.

到現時為止,看不到什麼信仰英雄,只有平凡的人,在無常的世局裡,為了生存,作出必要的行動.然後,訴之那測不透的天命,期望解救.

(第五幕)大壞蛋哈曼的詭計

5:12哈曼又說:「王后以斯帖預備筵席,除了我之外不許別人隨王赴席.明日王后又請我隨王赴席；

5:13只是我見猶大人末底改坐在朝門,雖有這一切榮耀,也與我無益.」

5:14他的妻細利斯和他一切的朋友對他說:「不如立一個五丈高(約75呎)的木架,明早求王將末底改掛在其上,然後你可以歡歡喜喜地隨王赴席.」哈曼以這話為美,就叫人做了木架.

香港的樓宇,樓底一般約10呎,哈曼造的木架,有7層樓那麼高.將人釘在上面或吊在上面,地面的人可能連看都看不清.若要羞辱敵人,有需要造一個這麼高的木架嗎?情節越誇張,效果越滑稽.猶如一個小孩子,拖著十尺大刀,要與一隻蚊決鬥,是否喜劇十足?

(第六幕)諸多偶然,冥冥之中

6:1那夜王睡不著覺,就吩咐人取歷史來,念給他聽.

6:2正遇見書上寫著說:王的太監中有兩個守門的,辟探和提列,想要下手害亞哈隨魯王,末底改將這事告訴王后.

歷史有無數偶然,也有無數偶然交織而成的必然.哈曼實施詭計的前一晚,亞哈隨魯王失眠.無事做,偶然讀起歷史來.偶然讀到末底改通風報信的事蹟.

6:3王說:「末底改行了這事,賜他什麼尊榮爵位沒有?」伺候王的臣僕回答說:「沒有賜他什麼.」

末底改救了王,作了這等大事,竟然沒有什麼賞賜,這是一不可能的偶然.

6:4王說:「誰在院子裡?」(那時哈曼正進王宮的外院,要求王將末底改掛在他所預備的木架上.)

6:5臣僕說:「哈曼站在院內.」王說:「叫他進來.」

6:6哈曼就進去.王問他說:「王所喜悅尊榮的人,當如何待他呢?」哈曼心裡說:「王所喜悅尊榮的,不是我是誰呢?」

6:7哈曼就回答說:「王所喜悅尊榮的,

6:8當將王常穿的朝服和戴冠的御馬,

6:9都交給王極尊貴的一個大臣,命他將衣服給王所喜悅尊榮的人穿上,使他騎上馬,走遍城裡的街市,在他面前宣告說:王所喜悅尊榮的人,就如此待他.」

在這靜夜,大壞蛋竟還走來王宮,急不及待的求王成全他的詭計,這是多麼的不可能.當很多「剛好如此」拼在一起時,便會構成笑話.以前的「老夫子」漫畫,題目常常是「無獨有偶」.當一隻貴婦狗的毛,與拖著牠的貴婦的頭髮一模一樣時,便會引人發笑.話說,在蘇聯時代,有個美國人去莫斯科旅行,在克里姆林宮門口,和當地人閒聊.很快,他們便談到自由.美國人說,我們最大的自由,便是在白宮門口,大罵美國政府的不是.那蘇聯人拍手回應說,真巧,我們最大的自由,竟和你們一樣,便是在克宮門口,大罵美國政府的不是.「剛好如此」,「無獨有偶」,是笑話的來源.大壞蛋剛好出現在一剛好的時刻,成了笑料.他以為王要提拔他,豈料,王要提拔的,正是大壞蛋想弄死的人.

(第七幕)給蕉皮人踩,自己卻踩到蕉皮

7:1王帶著哈曼來赴王后以斯帖的筵席.

7:2這第二次在酒席筵前,王又問以斯帖說:「王后以斯帖啊,你要什麼,我必賜給你；你求什麼,就是國的一半也必為你成就.」

7:3王后以斯帖回答說:「我若在王眼前蒙恩,王若以為美,我所願的,是願王將我的性命賜給我；我所求的,是求王將我的本族賜給我.

7:4因我和我的本族被賣了,要剪除殺戮滅絕我們.我們若被賣為奴為婢,我也閉口不言；但王的損失,敵人萬不能補足.」

7:5亞哈隨魯王問王后以斯帖說:「擅敢起意如此行的是誰?這人在哪裡呢?」

7:6以斯帖說:「仇人敵人就是這惡人哈曼!」哈曼在王和王后面前就甚驚惶.

7:7王便大怒,起來離開酒席往御園去了.哈曼見王定意要加罪與他,就起來,求王后以斯帖救命.

7:8王從御園回到酒席之處,見哈曼伏在以斯帖所靠的榻上；王說:「他竟敢在宮內,在我面前凌辱王后嗎?」這話一出王口,人就蒙了哈曼的臉.

7:9伺候王的一個太監名叫哈波拿,說:「哈曼為那救王有功的末底改做了五丈高的木架,現今立在哈曼家裡.」王說:「把哈曼掛在其上.」

7:10於是人將哈曼掛在他為末底改所預備的木架上.王的忿怒這才止息.

大壞蛋以為王及王后看重他,邀請他赴宴.那知,劇情180度逆轉.猶大人的敵人,如今,成了王的敵人.劇情需要,王「離開酒席往御園去了」.這時,大壞蛋向王后求情,卻不經意地伏在王后靠著來吃東西的榻上.剛好在這時,王又回來,以為大壞蛋想污辱王后.其實,眾目睽睽,這怎麼可能!如此這般,土包一般的亞哈隨魯隨即怒從心上起.之後,大壞蛋便掛在自己預備的,高75呎的木架上.這種情節,十足像是他放下蕉皮,等人踩上,到頭來,卻是自己踩上了,翻了個觔斗.

總結:

《以斯帖記》記述的世事,反覆無常,毫無道理.它向世人點出,在波斯帝國的大系統下,事情原來如此空洞.

大壞蛋,總會有的,但到頭來,只是個栽一大觔斗的小丑.

《以斯帖記》全書沒有提及神,明顯,它是向認為世界如此荒謬而神根本不存在的人說話.的確,我們看不到神,但世事的不可預料正造成很多奇妙的巧合,神在我們理性無法掌控的這些巧合中,創造出生機.

理性框架以外的生機,就是喜劇,就是笑話.笑話總是在人的理性所掌控的常態外突破出全未想過的新事,笑話總是在現狀之外思考,笑話讓人走出既定的框框,想像一從未想過的世界.

當亞伯拉罕100歲時,他的妻子撒拉已90歲,神卻應許撒拉要生一個兒子,好實現神的救贖計劃.撒拉卻笑出來,這不是因為能懷孩子而喜悅地笑,而是覺得神講了一個大笑話而笑.從理性看,從任何角度看,90歲了,懷孩子是不可能的,一切只是笑話.最後,她生了以撒.「以撒」這名字的意思,便是「他笑」.不可能的,從未想過的,成為可能.原先那看來不可能的「笑話」,最後,真的能叫人滿足地笑.

在這個不好笑的世界裡,神創造笑話,讓我們懂得笑.

願上主得著榮耀,頌讚；願眾人得享喜樂,平安.



\newpage

\section{人在旅途做夢時}
\label{sec:656}
\textbf{林豪恩先生 - 人在旅途做夢時}
\newline
\newline
\hyperref[sec:655]{< < < PREV SERMON < < <}
~
\hyperref[sec:toc]{[返目錄]}
~
\hyperref[sec:657]{> > > NEXT SERMON > > >}


創世記 28:10-14
\newline
\begin{longtable}{cl}
\hline
\hline
章節 & 經文 (和合本修訂版)\\
\hline
28:10 & \begin{tabularx}{0.7\textwidth}{X} 雅各離開別是巴，往哈蘭去。 \end{tabularx} \\ \\ \relax
28:11 & \begin{tabularx}{0.7\textwidth}{X} 到了一個地方，因為已經日落，就在那裡過夜。他拾起那地方的一塊石頭枕在頭下，就躺在那地方。 \end{tabularx} \\ \\ \relax
28:12 & \begin{tabularx}{0.7\textwidth}{X} 他做夢，看哪，一個梯子立在地上，梯子的頂端直伸到天；看哪，神的使者在梯子上，上去下來。 \end{tabularx} \\ \\ \relax
28:13 & \begin{tabularx}{0.7\textwidth}{X} 看哪，耶和華站在梯子上面，說：「我是耶和華—你祖父亞伯拉罕的神，以撒的神。你現在躺臥之地，我要將它賜給你和你的後裔。 \end{tabularx} \\ \\ \relax
28:14 & \begin{tabularx}{0.7\textwidth}{X} 你的後裔必像地上的塵沙，必向東西南北開展；地上萬族必因你和你的後裔得福。 \end{tabularx} \\ \\
[1ex]
\hline
\hline
\end{longtable}
馬太福音 13:24-43
\newline
\begin{longtable}{cl}
\hline
\hline
章節 & 經文 (和合本修訂版)\\
\hline
13:24 & \begin{tabularx}{0.7\textwidth}{X} 耶穌又設個比喻對他們說：「天國好比人撒好種在田裡， \end{tabularx} \\ \\ \relax
13:25 & \begin{tabularx}{0.7\textwidth}{X} 在人睡覺的時候，他的仇敵來，把雜草撒在麥子裡就走了。 \end{tabularx} \\ \\ \relax
13:26 & \begin{tabularx}{0.7\textwidth}{X} 到長苗吐穗的時候，雜草也顯出來。 \end{tabularx} \\ \\ \relax
13:27 & \begin{tabularx}{0.7\textwidth}{X} 地主的僕人進前來對他說：『主人，你不是撒好種在田裡嗎？哪裡來的雜草呢？』 \end{tabularx} \\ \\ \relax
13:28 & \begin{tabularx}{0.7\textwidth}{X} 主人回答他們：『這是仇敵做的。』僕人對他說：『你要我們去拔掉嗎？』 \end{tabularx} \\ \\ \relax
13:29 & \begin{tabularx}{0.7\textwidth}{X} 主人說：『不必，恐怕拔雜草，也把麥子連根拔出來。 \end{tabularx} \\ \\ \relax
13:30 & \begin{tabularx}{0.7\textwidth}{X} 讓這兩樣一起長，等到收割。當收割的時候，我會對收割的人說，先把雜草拔出來，捆成捆，留著燒，把麥子收在我的倉裡。』」 \end{tabularx} \\ \\ \relax
13:31 & \begin{tabularx}{0.7\textwidth}{X} 他又設個比喻對他們說：「天國好比一粒芥菜種，有人拿去種在田裡。 \end{tabularx} \\ \\ \relax
13:32 & \begin{tabularx}{0.7\textwidth}{X} 它原比所有的種子都小，等到長起來，卻比各樣的菜都大，且成了樹，以致天上的飛鳥來在它的枝上築巢。」 \end{tabularx} \\ \\ \relax
13:33 & \begin{tabularx}{0.7\textwidth}{X} 他又對他們講另一個比喻：「天國好比麵酵，有婦人拿來放進三斗麵裡，直到全團都發起來。」 \end{tabularx} \\ \\ \relax
13:34 & \begin{tabularx}{0.7\textwidth}{X} 這都是耶穌用比喻對眾人說的話，不用比喻，他就不對他們說甚麼。 \end{tabularx} \\ \\ \relax
13:35 & \begin{tabularx}{0.7\textwidth}{X} 這是要應驗先知所說的話：「我要開口說比喻，說出從創世以來所隱藏的事。」 \end{tabularx} \\ \\ \relax
13:36 & \begin{tabularx}{0.7\textwidth}{X} 當時，耶穌離開眾人，進了屋子。他的門徒進前來，說：「請把田間雜草的比喻講給我們聽。」 \end{tabularx} \\ \\ \relax
13:37 & \begin{tabularx}{0.7\textwidth}{X} 他回答：「那撒好種的就是人子， \end{tabularx} \\ \\ \relax
13:38 & \begin{tabularx}{0.7\textwidth}{X} 田地就是世界，好種就是天國之子，雜草就是那惡者之子， \end{tabularx} \\ \\ \relax
13:39 & \begin{tabularx}{0.7\textwidth}{X} 撒雜草的仇敵就是魔鬼，收割的時候就是世代的終結，收割的人就是天使。 \end{tabularx} \\ \\ \relax
13:40 & \begin{tabularx}{0.7\textwidth}{X} 正如把雜草拔出來用火焚燒，世代的終結也要如此。 \end{tabularx} \\ \\ \relax
13:41 & \begin{tabularx}{0.7\textwidth}{X} 人子要差遣他的使者，把一切使人跌倒的和作惡的從他國裡挑出來， \end{tabularx} \\ \\ \relax
13:42 & \begin{tabularx}{0.7\textwidth}{X} 丟在火爐裡，在那裡要哀哭切齒了。 \end{tabularx} \\ \\ \relax
13:43 & \begin{tabularx}{0.7\textwidth}{X} 那時，義人要在他們父的國裡發出光來，像太陽一樣。有耳的，就應當聽！」 \end{tabularx} \\ \\
[1ex]
\hline
\hline
\end{longtable}
講題:人在旅途做夢時

講員:林豪恩先生

經文:創世記28:10-14;馬太福音13:24-43

日期:2023年7月23日

一．香港人的夢

一般香港普羅市民發甚麼夢呢?有一位著名Samuel Chan的有這樣的觀察:「我身邊有一位朋友,期期都買六合彩,開八百萬又好金多寶一億半又好,一定買.「八百萬都可以買房子住.」講到真的會中獎那般.其實很多人,打從心底覺得自己一定不會中,但都堅持要買.他對我說:「你以為我買六合彩真的以為是為了錢嗎?我也知道一定不會中獎,我買六合彩是買一個作夢的時間.」什麼夢?「這個夢厲害了.從你塞一張二十塊錢給櫃台,跟櫃台小姐姐說『麻煩電腦票』,夢就開始了,直到到開獎那一刻才醒過來.每個人買完六合彩之後,一離開投注站就開始想『呀,如果我中了頭獎,那些錢會怎麼用呢?』「最要緊最要緊,是買房.買一間給自己住,我喜歡禮頓山,買一間太古城的送給爸爸媽媽,他們喜歡太古城的間格.再多買兩三套房收租金.」講到很有計劃,好像是真的.然後他又回到現實.「唉,其實我都知道沒那麼容易的,發財的事,你想就可以了嗎?作夢就可以.不過這二十塊錢就是買這兩,三天的時間.作夢一樣那麼開心.」

這個城市,很多人工作辛苦,朝十晚十為兩餐,倒頭來可能收入不過兩萬,給了家用夠生活,就不夠供房子.而六合彩有一種魔力,小小的夢想,為城市人大大的打氣.今日還在上班?不打緊,買金多寶,真命天子可能是你.

這不是賭博,只是買一個夢,作幾天的夢,不中獎的話又再出發.你可以說香港人要靠六合彩完夢是一種悲哀,現實是,幾乎每個香港人都作過這發財的夢.」(\url{https://www.britannia-study.com/en/article/}\%E7\%99\%BC\%E9\%81\%94\%E5\%A4\%A2\%E7\%9A\%84\%E5\%A4\%A2/)

二．雅各在伯特利做的夢

今天的舊約經課,也記載了一個夢.是誰的夢呢?是怎樣的夢呢?這個夢對我們有甚麼啟發呢?

28:10雅各出了別是巴,向哈蘭走去；

28:11到了一個地方,因為太陽落了,就在那裡住宿,便拾起那地方的一塊石頭枕在頭下,在那裡躺臥睡了,

當時的雅各很可能還是一個少年人.這個男孩子為甚麼獨自一人上路呢?為甚麼露宿荒野,以石頭為枕呢?

雅各一個人在途上,其實是逼不得已的.雅各本來生活在一家四口,令人想到溫馨的小康之家.他的父親是以撒,母親是利百加,有一個孿生兄長叫以掃.現實上,這個家庭似乎並不溫馨,反而是充滿著明爭暗鬥.有一天,雅各以一碗紅豆湯收買了哥哥的長子名分,並在媽媽的安排下假扮哥哥騙取了父親對長子的祝福.事件被揭發後,以掃盛怒.在風頭火勢的情況下,媽媽恐怕以掃會殺死雅各報仇,於是安排雅各逃走,去哈蘭投靠自己的娘家.於是,雅各展開了一個人在途上的逃難之旅.

一個人在途上是怎樣的心境呢?中學範文有郁達夫的「一個人在途上」,他引述盧騷晚年的作品《孤獨散步者的夢想》的頭幾句話:「自家除了己身以外,已經沒有弟兄,沒有鄰人,沒有朋友,沒有社會了,自家在這世上,像這樣的,已經成了一個孤獨者了.」文章結束時,他如此表達自己的狀況:「現在去北京遠了,去龍兒更遠了,自家只一個人,只是孤零丁的一個人.在這裡繼續此生中大約是完不了的飄泊.」

雅各一個人在逃命的路上,還有夢可以造嗎?生存已經是壓倒性的掙扎,還可以有甚麼夢呢?行行重行行,這位少年拖著疲累的身軀倒下了.我們不知道在荒野以石頭枕首的雅各睡得好不好,然而,他當晚的確造了一個夢:

28:12夢見一個梯子立在地上,梯子的頭頂著天,有上帝的使者在梯子上,上去下來.

28:13耶和華站在梯子以上,說:「我是耶和華－你祖亞伯拉罕的上帝,也是以撒的上帝；我要將你現在所躺臥之地賜給你和你的後裔.

28:14你的後裔必像地上的塵沙那樣多,必向東西南北開展；

這個夢,對於在逃亡中的雅各來說是一個美夢:第一,釋放了生存的焦慮,保證了不會被殺害也至於死去.第二,不單止不會死,更有機會結婚生子,傳宗接代.第三,不單止後續有人,更能開枝散葉,後代多得地上的塵沙那樣；第四,不單止有人,更重要的是有地,這真是回應了當時要寄人籬下的雅各的需要.一個夢,就正如香港人造中六合彩的發達夢,只要中了六合彩,所有需要都滿足了,所有問題都解決了.

三．夢想達到又如何?

現實上又是怎樣呢?就算美夢實現了,又如何呢?不少人都以為中六合彩就解決一切問題,一切困難都自動消失.網媒文章「六合彩中頭獎是福是禍?盤點各地幸運兒悲慘結局!」指出,「原來中彩票後的下場並非大家幻想中般美好,外國有些幸運兒更在中獎後遇到被殺害,破產等下場.」「2006年,美國佛羅里達州一位貨車司機中了3,100萬美元的彩票巨獎.後來,一名女子表示想為他寫自傳,於是開始接近中獎男子.其後,二人逐漸發展成為了男女朋友.及後,中獎男子被發現躺臥在朋友的寓所後花園,證實死亡.讓男子為甚麼死亡呢?後來,他的女朋友被判一級謀殺和終身監禁,因為她為了搶奪男友的獎金而殺害了他.」

又有英國一名男子19歲就獲得970萬英鎊獎金,正所謂「有錢就任性」,開名車,穿名牌服飾,食得招積,住得舒適,遊得出色,經常在牙買加,巴巴多斯和波多黎各巴努斯等豪華目的地與朋友聚會.有錢就容任性之外,有錢就生活不檢點,縱慾無道,最後他的妻子受不了跟他離婚,並帶走女兒.過了十年荒誕生活,至2012年,他的錢花光了,被迫開始過領取救濟金的生活.」(\url{https://www.edigest.hk/206747/?utm_campaign=ED_ContentCopy and utm_source=Web-inventory and utm_medium=Content-Copy_ED})

雅各的後代,要人有人,要地有地,然而,自己十二兄弟卻打生打死,內憂外患,打了無數大大小小的戰爭,死傷無數.雅各的夢實現了,他的性命得保,開技散葉,有人有地,那又怎樣呢?這個美夢到底所為何事呢?

周星馳在「少林足球」有一句經典對白︰「做人如果無夢想,同條鹹魚有咩分別呀?」其實,不需要擔心人們不造夢,正如之前提及,香港普羅市民都有一個中六合彩的發達夢；不少公司都有一個要超越對手,傲視同儕的夢；不少學校都有一個排名上升,成為地區第一,甚至世界第一的夢；不少國家都有一個要強要大的復興帝國夢.我們要擔心的反而是,這些夢所為何事?這些夢實現了又如何?

四．雅各夢的意義

讓我們回到雅各的夢,如果雅各的夢只停留在土地和後代,其實只是一個好像香港普羅市民的發達夢,一個解決個人焦慮和滿足個人需要的夢,一個令自己揚眉吐氣的出人頭地夢.如此,這個上帝所賜予的夢和我們的白日夢有甚麼分別呢?這個上帝揀選的人和其他人有甚麼分別呢?雅各的後代和其他民族有甚麼分別呢?上帝賜給雅各這個夢所為何事呢?讓我們再進入雅各的夢境仔細查探.

28:12夢見一個梯子立在地上,梯子的頭頂著天,有上帝的使者在梯子上,上去下來.

28:13耶和華站在梯子以上,說:「我是耶和華－你祖亞伯拉罕的上帝,也是以撒的上帝；我要將你現在所躺臥之地賜給你和你的後裔.

28:14你的後裔必像地上的塵沙那樣多,必向東西南北開展；地上萬族必因你和你的後裔得福.

我們剛才刻意忽略了最後那句,結果雅各夢就變得平凡無奇.最後那句「地上萬族必因你和你的後裔得福」不是可有可無,而是雅各夢的遠景和意義所在.如果雅各夢沒有這遠象,就根本不值得後人當作經典來研讀.雅各夢不是保證我們上帝必然救我們不死,而是啟發我們生存的目的和意義；雅各夢也不是保證我們必然興旺強大,而是挑戰我們興盛強大到底所為何事?就算個人爭奪得武林至尊寶座,就算學校爬上了排名的頂峰；就算國家戰鬥力無人匹敵,又如何呢?對於「地上萬族」,能夠「必因你和你的後裔得福」嗎?

雅各的夢在逃亡中被賜予,或許生存的焦慮蓋過了夢的理想,自身的安危逼走了對其他人的關切.然而,上帝賜予雅各這夢,回應雅各當下的需要還算是小事,上帝有更大的關切,就是雅各自身以外的「地上萬族」的福祉.今天,強大的聲音也正在催促及誘導我們造滿足自身需要的夢,造解決自身焦慮的夢,造令自己舒暢和縱慾的夢.這些夢,最好沒有對他者的關切,也不要有對其他民族和其他國家的關切,更不要有對地球,對世界的關切.雅各夢的遠像和意義還有著意嗎?作為信仰上帝的人,如何看待這夢呢?還要追尋嗎?值得追尋嗎?讓我們閱讀福音經課來尋找啟迪.

五．天國夢

耶穌向門徒講了一個天國夢:

13:24耶穌又設個比喻對他們說:「天國好像人撒好種在田裡,

13:25及至人睡覺的時候,有仇敵來,將稗子撒在麥子裡就走了.

13:26到長苗吐穗的時候,稗子也顯出來.

天國在夢中是完美的,然而,在看得見的狀態卻是摻雜的.

13:27田主的僕人來告訴他說:『主啊,你不是撒好種在田裡嗎?從哪裡來的稗子呢?』

13:28主人說:『這是仇敵做的.』僕人說:『你要我們去薅出來嗎?』

完美的天國夢被摻雜,怎能夠忍受呢?越熱切尋夢的就越要把摻雜的清除.

13:29主人說:『不必,恐怕薅稗子,連麥子也拔出來.

天國的主人竟然能夠忍受完美被稗子摻雜,為甚麼呢?他不是稗子麥子不分,也不是認可稗子,他忍受稗子摻雜他的天國,完全是因為他關切麥子福祉,他不忍心為了除掉稗子而傷害麥子,他清楚知道天國夢的目的和意義並不是建功立業,而是慈愛和公義.

13:41人子要差遣使者,把一切叫人跌倒的和作惡的,從他國裡挑出來,

13:42丟在火爐裡；在那裡必要哀哭切齒了.

13:43那時,義人在他們父的國裡,要發出光來,像太陽一樣.有耳可聽的,就應當聽!」

天國夢是有價值導向的,慈愛的性質保障無辜者不受傷害,公義的性質保證善有善報,惡有惡報.天國夢雖然看不到,因為看得到的都是已經被摻雜了的狀況,然而,那位賜予夢的上主,卻保證天國夢是真實的,終有一天會實現的.在旅途中的教會值得繼續追尋這夢嗎?在旅途中的信徒值得追尋這夢嗎?值得呼籲在旅途中的其他人一齊追尋這夢嗎?耶穌告訴我們這是值得的,因為追尋這夢的最終被肯定,他們知道所獻上的一切對這個夢的實現來說或許微不足道,然而,耶穌基督保證:

13:43那時,義人在他們父的國裡,要發出光來,像太陽一樣.有耳可聽的,就應當聽!」



\newpage

\section{來到了尼波山}
\label{sec:657}
\textbf{梁元生教授 - 來到了尼波山}
\newline
\newline
\hyperref[sec:656]{< < < PREV SERMON < < <}
~
\hyperref[sec:toc]{[返目錄]}
~
\hyperref[sec:658]{> > > NEXT SERMON > > >}


講題:來到了尼波山

講員:梁元生教授

日期:2023年7月30日

一.前言

上個禮拜助理校牧林豪恩先生和大家講人生旅途做夢時.今日我想繼續講人生旅途.人生旅途中有多個驛站,而每個驛站都很可能是一個人生中的轉捩點.

今日要和大家講道的題目是關於尼波山(Mount Nebo)的.而尼波山對很多人來說,可以說是人生旅途的最後一個驛站.熟悉聖經史地的當然都會聽過或者甚至去過尼波山,因為現代的聖地旅遊都會把它看作一個重要的景點.尼波山位於今日中東約旦國內,

離首都安曼約有三十餘哩,尼波山的山勢不算高峻,只有六百多米(約2700英尺),和我們對面的馬鞍山差不多高度.但正如中國古語有雲:山不在高,有龍則靈.尼波山在聖經傳統中,相當有名,有著許多關於耶和華神的故事和傳說.現時山上有紀念摩西的碑石,也有一具銅蛇的雕塑,是一位現代義大利藝術家的作品；山上還有廢棄的古老教堂和修道院,是中世紀的建築.這些人為的添加物品,無論是建築物或是藝術雕塑,都是以前所沒有的.舊約時代的尼波山,大概就是一座自然的光禿的山丘,和蓖斯迦山(Mt.Pisgah)山脈相連,雙峰相對.我自己雖然沒有去過尼波山,但今天且容許我權充導遊,和大家來到尼波山,看看尼波山的面貌和風景.

先來看看尼波山的地圖,看它坐落何方:

下面帶大家看看尼波山的實景:尼波山 (Mount Nebo)(2706呎).

二．摩西來到尼波山:偉人之“憾”

聖經上提到尼波山時,多次都會提及一個人物—-摩西.上帝在尼波山和他說話,與他檢視過去,也向他展示未來.在舊約聖經《申命記》三十四章:1－4節那裡,對這次神人之會有相當清晰的記載:

「摩西從摩押平原,登尼波山,上了那與耶利哥相對的毗斯迦山頂.耶和華把基列全地,直到但,拿弗他利全地,以法蓮,嗎哪西的地,猶大全地,直到西海,南地,和棕樹城耶利哥的平原,直到瑣珥,都指給他看.耶和華對他說:這就是我向亞伯拉罕,以撒,雅各起誓應許之地,說:我必將這地賜給你的後裔.現在我使你眼睛看見了,你卻不得過那邊去.」

聖經中的尼波山是一個神人相會或神人相遇的地方(a place of encountering God),在那裡,人得見上帝,並且得著上帝的曉諭和指示.正因如此,歷代許多基督教的聖者(如教宗保祿二世)都曾來到尼波山,要尋求一見神的面,或聽到祂的諭旨,和得到祂的啟示.但更多的人,像摩西那樣,來到尼波山,是來到了人生旅途的最後一個驛站,人生的最後一個關口.尼波山象徵著人生的終站或者是人生的尾段:到了尼波山,意思是人生的旅程將要終結,死亡快要來到.到了尼波山,亦即到了晚年,到了老年,到了垂暮之年.在此時此刻,應該做些甚麼?難道還充滿那少年的幻想和青年的追夢時刻麼?

到了這裡,是時候對自己的人生作全盤的總結與檢討,是時候重溫上帝的教誨和應許,因為到了尼波山,你知道這段人生路將要結束,同時你將要面對另一個全新的旅程,一個進到另一世界的新征程!在這裡,你會有甚麼想法?你會作何準備?

對於摩西來說,尼波山既是與神交會的地方,也是他人生旅途的一個柺點,也可以說是他的人生之中最重要的轉捩點.

不但如此,對摩西來說,尼波山還隱含著另一層重要的意義.對他而言,尼波山是可以看見,但又不能過去的山.“可以見”但“不能過”!摩西一生大部份的時候都聽從神的吩咐,青年時期他和有權有勢的埃及法老進行抗爭,帶領以色列人做過一連串轟轟烈烈的大事:過紅海,越沙漠,經曠野,攀高山,戰蠻王...等等,經歷漫長的四十年,最終來到神所應許的目的地—流奶與蜜的迦南美地.在以色列人的歷史中,他是一個不可多得的領袖,是一個極受尊崇的偉大人物.是偉人,是先知!然而,耶和華神把他帶到尼波山上,讓他看見應許之地,卻對他說:你只能望見,但不得過去.你猜摩西心頭上是怎麼樣的感覺?有人也許會覺得這是一樁悲劇,也有人會為付出了這麼多的摩西而憤憤不平?

可望而不可即!可以見而不能過!這是否人生最大的遺憾?盡了一切的努力,而未能看到願望實現,是否一樁悲劇?是否終生的遺憾?

讓我告訴你:我們都會遇到尼波山.人生未必事事如意,更未必事事如願.人生有許多未圓的夢,人生有許多想做卻做不了的事.然而,如果我們能像摩西在尼波山上那樣,我們就是非常幸福和蒙恩的了.因為神使我們“看見”,意思是祂向我們展示了祂的信實,兌現了祂的承諾和應許.對摩西來說,甚麼是實實在在的盼望?不就是在尼波山上所見到的一切嗎?!摩西雖然不得過去約旦河彼岸的應許地,但他清楚看見神應許的兌現.他親自體驗到和他說話的那一位正是賜他十誡,法版,以嗎哪,鵪鶉,雲柱,火柱帶領他和以色列人走過四十年的曠野路.尼波山上和他說話的那一位,

既要提醒摩西其應許和承諾不變,但也要指出摩西曾經犯下的不信之罪.使摩西明白,他的人生為何充滿祝福,但犯了罪也必須接受咒詛,承受懲罰.明白到這一點,因此他不再奮力前行,而欣然接受上帝的安排,孤身一人獨自上了高山,在尼波山上找到自己最終的歸宿.那最終的孤獨(the ultimate solitude),變成為與神同在的愉悅(happiness of being together with God).摩西不用跨越約旦河和穿過耶利哥,尼波山變成他與神永遠同在的地方.偉人之“憾”,至終乃得無憾!

三．約書亞來到尼波山:天人之“諾”

從我們的角度看,摩西的終身遺憾是看見而不得親臨應許之地.這是整代出埃及的以色列人的悲哀,流浪曠野四十年,雖有雲柱火柱看顧隨行,但他們信心不足,受異族文化和宗教的影響,膜拜金牛,崇奉巴力,惹怒了耶和華.

所以,幾乎所有追隨摩西出埃及,過紅海的以色列人都不能進入上帝應許的迦南流奶與蜜之地.除了約書亞和迦勒!約書亞是以法蓮支派嫩族的後人,出埃及時仍算年青,由他繼承摩西,

帶領以色列人的第二代,或許還有許多第三代,第四代的人,他們是在出埃及之後在曠野出生和長大的以色列人.摩西代表的是年老的一代,來到尼波山時已經是垂垂老矣,走到人生最後的一段路,尼波山前,千思纏繞,百感交集；回顧多於前瞻,感恩中常帶愧疚.

我們在座中也有不少步入晚年的長者,多少能體諒摩西在尼波山的心境.即是:一切繁華將會過去,曾經擁有的權力也終歸於虛空,就如摩西在詩篇九十篇中祈禱所說的話.“我們一生的年日是七十歲,若是強壯可到八十歲,但其中可矜誇的不過是勞苦愁煩,轉眼成空,我們便如飛而去....求你指教我們怎樣數算自己的日子,好叫我們得著智慧的心.”我們當中有摩西那一代的人,生命中或許有過叱吒與輝煌的時候,但到了尼波山,即是晚景殘年之際,都會有一種覺悟:即人之作為遠不如神之應許來得可靠.

然而,來到尼波山前以約書亞為首的第二代和第三代人,

他們的心情和想法和摩西可能是很不一樣的.第二代或第三代人,他們的心態是怎樣的呢?我們來看約書亞記一章:1-6節:“耶和華的僕人摩西死了以後,耶和華曉諭摩西的幫手,嫩的兒子約書亞,說:我的僕人摩西死了,先在你要起來,和眾百姓過這約旦河,往我所要賜給以色列人的地去.凡你們腳掌所踏之地,我都照著我所應許摩西的話賜給你們了.從曠野和這黎巴嫩,直到伯拉大河,赫人的全地,

誘導大海日落指出,都要做你們的境界.你平生的日子,必無一人能在你面前站立得住.我怎樣與摩西同在,也必照樣與你同在.我必不撇下你,也不丟棄你.你當剛強壯膽,因為你必使這百姓承受那地為業,就是我向他們列祖起誓應許賜給他們的地.”

耶和華神要摩西停止腳步,卻命令約書亞繼續前行.過了尼波山,還要克服那阻擋前路的約旦河.約旦河與死海相連,許多聖經學者皆以約旦河比喻生命的終結和死亡的界線.然而,對於信仰群體來說,約旦河水攔不住向前的步伐

,死海寒波也不能隔絕上帝對其子民的應許.只有犯罪(如以色列人在曠野的生活和行為)和不信,才會封鎖那通往應許之路.耶和華神在尼波山,在摩西死後,特別鄭重的對約書亞說:“你要起來,和眾百姓過這約旦河,往我所要賜給你們的地去.”這是上帝向約書亞,也是向以色列人的應許.為了使約書亞更加剛強

,神更進一步清楚地強調:“

你當剛強壯膽.我怎樣與摩西同在,也必照樣與你同在.”這是神在尼波山向約書亞所許的承諾.同樣地,面對滔滔滾滾的約旦河水,堅固不墜的耶利哥城,以及那些並不友善的蠻族異邦,過河之後的以色列人並不是全是一帆風順,而是處處挑戰,困難重重的境況,面對這些困難和挑戰,約書亞也向耶和華神作出相對的承諾.那就是,無論境況是順是逆:“至於我和我的家,我們必定事奉耶和華!“(書24:15)

下了尼波山,約書亞和他帶領的信仰群體,帶著上帝的應許和彼此作出的承諾,繼續前行.過河,攻城,平亂,在應許之地安頓下來,成家,安族,到最後建城,立國‧‧‧.當然,

那是另一新篇章和另一段歷史了!

四.結語:我們來到尼波山

按照教會年曆提供的經文,今天我應該講的聖經人物是雅各.但我思量再三,選擇要講的是摩西和約書亞,原因是我近一兩年來,身體逐漸衰壞,病痛纏繞,對於進入人生晚年,有比較強烈的感受.摩西到了尼波山的情景和心態,時常在腦海中回盪,思前想後,既後悔以前的錯失,就如摩西的不信,又恐時不我予,渡河之日將至,難以彌補一生之遺憾.我想在座不少弟兄姊妹也漸漸踏進老年,愈來愈接近尼波山界.是回顧人生,檢討生命的時候了.首先,是和耶和華相會,交談,重新得著祂的應許和承諾,再求進一步在人生旅途上前行.然而,我今天也講約書亞,他代表著我們這個信仰群體的第二代,和其他年輕的一代.

我要提醒他們:他們也需要經過尼波山,在這山上重燃對上帝的愛火和信心,得著上帝的重新承諾和鼓勵:你當剛強壯膽.然後上路.在這裡,你和約書亞一樣,應當也作出你們向神的承諾:就是:至於我和我的一家,都必事奉耶和華.

我們都活此一生,旅途上我們都會遇到我們的尼波山.它提醒我們人生須有上帝作主,人生須有上帝同行,始能無憾,始能看到應許,始能欣然上路,始能毅然渡河.生命的充實和豐足,其意義從尼波山的反省或可略窺一二.

最後,我想以幾行詩句作為今日講道的結束:

陟彼高崗,我心其惶；惡水堅城,前路阻且長.

尼波遠眺,我心其傷；應許美地,可望不可即.

下彼高崗,我心其懽；水分城塌,護我抵家鄉.

尼波山兮,人生之驛旅；上帝之帳篷兮,伴我在人間!

用現代白話翻譯出來如下:

來到了高山之上,我的心滿是惶恐；

前路險阻,是洶洶的約旦河水,和堅固不墮的耶利哥城；

在尼波山巔遠望,我的心帶著憂傷；

上帝所應許的美麗地方,

我只能遠望,卻不能夠親身抵達!

從山上下來,我的心滿是歡慰；

河水分開,城牆崩塌,上帝帶領我到達迦南家鄉.

尼波山啊尼波山,人生旅途上的驛站；

上帝的帳篷啊,在這人間之旅,賴你永遠與我相隨.



\newpage

\section{你們給他們吃吧}
\label{sec:658}
\textbf{高國雄牧師 - 你們給他們吃吧!}
\newline
\newline
\hyperref[sec:657]{< < < PREV SERMON < < <}
~
\hyperref[sec:toc]{[返目錄]}
~
\hyperref[sec:659]{> > > NEXT SERMON > > >}


馬太福音 14:13-21
\newline
\begin{longtable}{cl}
\hline
\hline
章節 & 經文 (和合本修訂版)\\
\hline
14:13 & \begin{tabularx}{0.7\textwidth}{X} 耶穌聽到了，就從那裡上船，私下退到荒野的地方去。眾人聽到後，從各城來，步行跟隨他。 \end{tabularx} \\ \\ \relax
14:14 & \begin{tabularx}{0.7\textwidth}{X} 耶穌出來，見有一大群人，就憐憫他們，治好了他們的病人。 \end{tabularx} \\ \\ \relax
14:15 & \begin{tabularx}{0.7\textwidth}{X} 傍晚的時候，門徒進前來，說：「這地方偏僻，而且時候已經晚了，請叫眾人散去，他們好進村子，自己買些食物。」 \end{tabularx} \\ \\ \relax
14:16 & \begin{tabularx}{0.7\textwidth}{X} 耶穌對他們說：「不用他們去，你們給他們吃吧！」 \end{tabularx} \\ \\ \relax
14:17 & \begin{tabularx}{0.7\textwidth}{X} 門徒說：「我們這裡只有五個餅、兩條魚。」 \end{tabularx} \\ \\ \relax
14:18 & \begin{tabularx}{0.7\textwidth}{X} 耶穌說：「拿過來給我。」 \end{tabularx} \\ \\ \relax
14:19 & \begin{tabularx}{0.7\textwidth}{X} 於是他吩咐眾人坐在草地上，就拿著這五個餅和兩條魚，望著天祝福，擘開餅，遞給門徒，門徒又遞給眾人。 \end{tabularx} \\ \\ \relax
14:20 & \begin{tabularx}{0.7\textwidth}{X} 他們都吃，並且吃飽了。門徒把剩下的碎屑收拾起來，裝滿了十二個籃子。 \end{tabularx} \\ \\ \relax
14:21 & \begin{tabularx}{0.7\textwidth}{X} 吃的人中，男的約有五千，還不算婦女和孩子。 \end{tabularx} \\ \\
[1ex]
\hline
\hline
\end{longtable}
講題:你們給他們吃吧!

講員:高國雄牧師

經文:馬太福音14:13-21

日期:2023年8月6日

13耶穌聽到了,就從那裡上船,私下退到荒野的地方去.眾人聽到後,從各城來,步行跟隨他.

14耶穌出來,見有一大群人,就憐憫他們,治好了他們的病人.

15傍晚的時候,門徒進前來,說:「這地方偏僻,而且時候已經晚了,請叫眾人散去,他們好進村子,自己買些食物.」

16耶穌對他們說:「不用他們去,你們給他們吃吧!」

17門徒說:「我們這裡只有五個餅,兩條魚.」

18耶穌說:「拿過來給我.」

19於是他吩咐眾人坐在草地上,就拿著這五個餅和兩條魚,望著天祝福,擘開餅,遞給門徒,門徒又遞給眾人.

20他們都吃,並且吃飽了.門徒把剩下的碎屑收拾起來,裝滿了十二個籃子.

21吃的人中,男的約有五千,還不算婦女和孩子.

4卷福音書都有記述,極少見

從經文脈絡,馬太福音的表達:V13耶穌聽到的,就是施洗約翰的死訊.

約翰是耶穌的表哥,更是他的先鋒(傳道的第一個伙伴).

這訊息讓耶穌選擇要「私下退到荒野」

所以「事情發生在荒野」是耶穌的選擇,也是他的需要,這成了這神蹟的重要背景及必需的源由…

耶穌因需要:退到荒野

眾人因需要:到荒野找耶穌

相遇由憐憫開始,神蹟亦由此而發生…

V

15傍晚的時候,門徒進前來,說:「這地方偏僻,而且時候已經晚了,請叫眾人散去,他們好進村子,自己買些食物.」

一個十分合理(甚至可能是日常的習慣),很正常不過的處理方案…

地方偏僻V13

馬可福音6:31他就說:「你們來,同我私下到荒野的地方去歇一歇.」這是因為來往的人多,他們連吃飯的時間也沒有.

因耶穌的選擇,在門徒合情合理的請求之下…

卻得著意想不到的回應

V16耶穌對他們說:「不用他們去,你們給他們吃吧!」

卻遇上一個”無理/無知/天真/naïve”的回應…

上帝曾否對我有類似的要求?

v17門徒說:「我們這裡只有五個餅,兩條魚.」

路加福音9:13下

他們說:「我們不過有五個餅,兩條魚,若不去為這許多人買食物就不夠.」

馬可福音6:37下

門徒對他說:「我們要拿兩百個銀幣去買餅給他們吃嗎?」

約翰福音6:

7腓力回答他:「就是兩百個銀幣的餅也不夠給他們每人吃一點點.」

門徒只有”事實的陳述”

五餅二魚的神蹟就接著發生…

反思:我們今天讀這段經文的意義…

五餅二魚能夠餵飽五千人?耶穌超然的大能力?明白耶穌有多「厲害」?

神蹟的參與?我的份?

我們可以如何學習?

神蹟的重點:

請注意:這神蹟與我們所認識的大部份神蹟不同,因為按我們的理解,這是”非必要的神蹟”,與其他醫病,趕鬼,叫死人復活的不一樣…

耶穌看似沒有必要行這神蹟,甚至門徒亦沒有期望過耶穌在此行神蹟,那麼,這神蹟的價值及意義就變得更為重要!

先回想一下,有沒有相類似的情況…

約翰福音2:1-11

迦拿婚宴,把水變成酒,一個使人盡慶的神蹟,一個表示出耶穌願意與人同歡的神蹟,一個彰顯耶穌榮耀的神蹟,亦是一切神蹟的開始(作為記號的起點(原文新約))…

對我們有什麼意義…

門徒所看見的,與耶穌所看見的有什麼分別?

今天我們又看見什麼呢?

耶穌所看見的…我們看到嗎?

今天我們所關心的是“資源”?

考慮的是性價比?

耶穌的吩咐:「你們給他們吃吧!」,對我們來說是何等的沉重呢?

耶穌遞給我們的,我們如何遞出去呢?

上帝所賜的,是足夠,並有餘的,你願意參與其中,使更多人因你的回應而得福嗎?



\newpage

\section{時不利兮}
\label{sec:659}
\textbf{林豪恩先生 - 時不利兮}
\newline
\newline
\hyperref[sec:658]{< < < PREV SERMON < < <}
~
\hyperref[sec:toc]{[返目錄]}
~
\hyperref[sec:660]{> > > NEXT SERMON > > >}


講題:時不利兮

講員:林豪恩先生

日期:2023年8月13日

一.引言

十七歲的時候,你在做甚麼呢?有沒有想過,這可能是死亡的時候呢?剛才誦讀的舊約,有一位十七歲的少年,名叫約瑟.就在十七歲的時候,他的十個哥哥密謀要殺死他.這件事令人想起秀茂坪童黨燒屍案,當年有一位與約瑟差不多年幾的16歲男孩被謀殺.

案發在1997年,當時我仍住在秀茂坪,距離案發現場約800米.16歲少年陸志偉和另外15名年齡介乎14至17歲的少年男女經常一齊玩耍,一齊出入,因此被稱為童黨.5月14日晚上,陸志偉被他的朋友殺死.相隔幾千年的兩件事件都有殺機,後者果真成為了謀殺案,前者則變化為企圖謀殺案,是甚麼做成這分別呢?兩個處境都動了殺機,群情洶湧,可謂大勢已去,大局已定,局中的個別人可以有甚麼反應呢?個別人的反應有甚麼效果呢?

二.積怨的背景

第一:約瑟將他哥哥們的惡行報給他們的父親

37:2約瑟十七歲與他哥哥們一同牧羊.他是個童子,與他父親的妾辟拉,悉帕的兒子們常在一處.約瑟將他哥哥們的惡行報給他們的父親.

第二:父親愛約瑟過於愛他的眾子

37:3以色列原來愛約瑟過於愛他的眾子,因為約瑟是他年老生的；他給約瑟做了一件彩衣.

37:4約瑟的哥哥們見父親愛約瑟過於愛他們,就恨約瑟,不與他說和睦的話.

第三:難道你真要管轄我們嗎?

37:5約瑟做了一夢,告訴他哥哥們,他們就越發恨他.

37:6約瑟對他們說:「請聽我所做的夢:

37:7我們在田裡捆禾稼,我的捆起來站著,你們的捆來圍著我的捆下拜.」

37:8他的哥哥們回答說:「難道你真要作我們的王嗎?難道你真要管轄我們嗎?」他們就因為他的夢和他的話越發恨他.

秀茂坪童黨經常聚集在任職清潔工人的輕度智障男子陳木清的住所,陳木清經常遭到童黨欺負,陸志偉勸他報警求助,童黨認為陸志偉出賣他們,被童黨視為背叛者.

三.殺機

37:13以色列對約瑟說:「你哥哥們不是在示劍放羊嗎?你來,我要打發你往他們那裡去.」約瑟說:「我在這裡.」

37:14以色列說:「你去看看你哥哥們平安不平安,群羊平安不平安,就回來報信給我」；於是打發他出希伯崙谷,他就往示劍去了.

37:15有人遇見他在田野走迷了路,就問他說:「你找甚麼?」

37:16他說:「我找我的哥哥們,求你告訴我,他們在何處放羊.」

37:17那人說:「他們已經走了,我聽見他們說要往多坍去.」約瑟就去追趕他哥哥們,遇見他們在多坍.

37:18他們遠遠地看見他,趁他還沒有走到跟前,大家就同謀要害死他,

1997年5月14日晚上,各被告以陳木清的名義誘騙陸志偉到陳木清的住所,陸志偉到達現場後隨即被各被告輪流拳打腳踢,以鐵通,摺椅等武器毆打陸志偉,並強迫他吃煙頭,倒豎蔥用頭撞地,撞牆,墊凳底,昏倒後淋凍水等…

陸志偉被打至奄奄一息…

失救致死.當天凌晨時分,童黨成員將陸志偉的屍體放進紙皮箱,運到馬路斜對面空置待拆的第32座放火燒屍,並向屍體潑硫酸使容貌更難辨認,眾人離開後再度返回燒屍現場,發現屍體仍然未被燒盡,於是再度放火燒屍,然後用垃圾膠袋裝好死者骨灰拋棄到垃圾站.

四.集體殺人案中的個人反應

14人認為死者背叛他們,敵愾同仇,誓要教訓死者,但對於是否要致對方於死地,到底有沒有人有其他想法的呢?如果有不同想法,在群情洶湧之下,如何自處呢?在毆打的整個晚上,現場的那14個人,到底有沒有人有不同的想法呢?如果有,在似乎大局已定的情況下,可以怎樣呢?

回望謀殺約瑟的案件,十位兄長都被約瑟向父親揭發惡行,或許因此都受過父親的責罰；十位兄長都見到父親對約瑟偏心,敵愾同仇；十位兄長都受不了約瑟的自視過人,對約瑟的反感和恨意已經把十兄弟連成一線.十兄弟出去放牧的日子,聲討約瑟可能已經是大家的共同話題.無論是認真,還是發洩,大家都已經有共識,就是要把約瑟除之而後快.

37:19彼此說:「你看!那做夢的來了.

37:20來吧!我們將他殺了,丟在一個坑裡,就說有惡獸把他吃了.我們且看他的夢將來怎麼樣.」

看來他們比秀茂坪童黨更有計劃,已經設計好處理屍體的方法.

敵愾同仇,群情洶湧的氣氛之下,置身其中的人的言論很容易被激化,說話和行動都會越演越烈.在這樣激烈高昂的群體行動之中,每一個人都衷心地認同行動的嗎?每一個人都義無反顧地去執行嗎?在秀茂坪童黨案中,有一個人提早離開現場,在受害人垂死沒有再予以毆打.另一個人的律師代表求情時指出,被告本來是好女孩,受男朋友影響才犯下此案,案法時曾經阻止其他被告毆打死者.雖然這些個別反應都未能改變結局,但卻令人看見在瘋狂的群體行動中的個體並不是鐵餅一塊.或許言詞上贊同,但若真的化為行動,仍然會有不同的回應.

眾兄弟已經計劃好了如何殺人滅屍,這個時候,大哥呂便出聲了:

37:21呂便聽見了,說:「我們不可害他的性命」；

37:22又說:「不可流他的血,可以把他丟在這野地的坑裡,不可下手害他.」

作者告訴讀者,「呂便的意思是要救他脫離他們的手,把他歸還他的父親.」可能有讀者認為這位大哥太過軟弱,作為大哥,既然不同意殺死約瑟,為甚麼不斥責弟弟們的行動?為甚麼不撥亂反正?既然想救約瑟,為甚麼不放他回家,反而建議「把他丟在這野地的坑裡」?這些質問都是合理的,只是事情就是如此進行.縱使呂便大哥的回應沒有徹底解救約瑟,但就是因為他的回應,拖延了約瑟的死亡時間,把約瑟的生命存留多一刻.

37:23約瑟到了他哥哥們那裡,他們就剝了他的外衣,就是他穿的那件彩衣,

37:24把他丟在坑裡；那坑是空的,裡頭沒有水.

把約瑟的性命存多多一刻會有甚麼事情發生,之後如何,或許呂便自己都未想得到.

37:25他們坐下吃飯

吃飯的時候,大家或許仍然在言語上嘲諷約瑟,以洩累積多年的憤怒.然而,我們可以想像,呂便心裡應該有點焦急,腦海中不斷思索,下一步可以怎樣.原來,不只他一個人在思索,還有另一個人似乎也是飯不知味,他雖然在和大家一齊食飯,但卻是眼看四方,似乎在搜索甚麼似的.終於,給他看見了:

37:25舉目觀看,見有一夥米甸的以實瑪利人從基列來,用駱駝馱著香料,乳香,沒藥,要帶下埃及去.

37:26猶大對眾弟兄說:「我們殺我們的兄弟,藏了他的血有甚麼益處呢?

37:27我們不如將他賣給以實瑪利人,不可下手害他；因為他是我們的兄弟,我們的骨肉.」眾弟兄就聽從了他.

以實瑪利人商隊的出現,給猶大製造了機會,讓他找到替代殺死約瑟的行動.然而,讀者可能會質疑,如果猶大真的如此重視兄弟情,為甚麼不對兄弟們曉以大義呢?為甚麼不幫助大家和好呢?這些問題也是合理的,但事情就是如此發展下去.

37:28有些米甸的商人從那裡經過,哥哥們就把約瑟從坑裡拉上來,講定二十舍客勒銀子,把約瑟賣給以實瑪利人.他們就把約瑟帶到埃及去了.

五.總結:時不利兮,或許只能見步行步

時不利兮,古有項羽拼死一戰,玉石俱焚；秀茂坪童黨案中無人能夠阻止慘案發生,所有人都悲劇收場；約瑟謀殺案卻被化解為企圖謀殺案,雖然被賣去埃及為奴,前途不明,生死未卜,但最低限度約瑟沒有在那一刻被殺死.生命仍然存留,就仍然有可能性.或者讀者仍然困惑不解,呂便和猶大既然想救約瑟,為甚麼不救到底呢?或者,在那種一面倒激烈高昂的氣氛之下,呂便和猶大只能見到行步,阻止約瑟死在兄弟手上,但同時也未有辦法化解大家對約瑟的恨意,未有辦法化解大家要約瑟消失的躁動.這個絕對不是美好的故事,也沒有甚麼榜樣值得仿傚.說到底,謀殺案怎樣說也是罪惡,約瑟後來的成就,絕對不可以合理他謀殺的罪惡.罪惡無處不在,跨越時空,跨越地域,有時收歛一點,有時猖狂,或者讀者可以藉著這個故事思考如何在罪惡的環境之中自處,如何能夠令邪惡不那麼猖獗.

這不是故事的結束,可以說只是故事的開始.在這個企圖謀殺的片段中,讀者看不見上帝,上帝似乎隱藏了.37章以約瑟的父親喪子之痛作結,38章耶和華這個名字出了,同時,劇情也出現了猶大連喪兩子,在那一刻,或者猶太終於體會到父親喪子之痛的滋味.39章開始是約瑟的故事,故事開首作者就強調「耶和華與他同在」.在接下來約瑟的故事中,讀者可以看見上帝在約瑟生命中的工作.

有時,或許我們感到時不利兮,但希望我們不要把這作為故事的結束,讓故事繼續發展下去.有時,我們未必能夠看見上帝的介入,然而,有一天上帝會顯現他的工作給我們知道.



\newpage

\section{婦人你的信心很大}
\label{sec:660}
\textbf{高國雄牧師 - 婦人,你的信心很大!}
\newline
\newline
\hyperref[sec:659]{< < < PREV SERMON < < <}
~
\hyperref[sec:toc]{[返目錄]}
~
\hyperref[sec:661]{> > > NEXT SERMON > > >}


馬太福音 15:21-30
\newline
\begin{longtable}{cl}
\hline
\hline
章節 & 經文 (和合本修訂版)\\
\hline
15:21 & \begin{tabularx}{0.7\textwidth}{X} 耶穌離開那裡，退到推羅、西頓境內。 \end{tabularx} \\ \\ \relax
15:22 & \begin{tabularx}{0.7\textwidth}{X} 有一個迦南婦人從那地方出來，喊著說：「主啊，大衛之子，可憐我！我女兒被鬼纏得很苦。」 \end{tabularx} \\ \\ \relax
15:23 & \begin{tabularx}{0.7\textwidth}{X} 耶穌卻一言不答。門徒進前來，求他說：「這婦人在我們後頭喊叫，請打發她走吧。」 \end{tabularx} \\ \\ \relax
15:24 & \begin{tabularx}{0.7\textwidth}{X} 耶穌回答：「我奉差遣只到以色列家迷失的羊那裡去。」 \end{tabularx} \\ \\ \relax
15:25 & \begin{tabularx}{0.7\textwidth}{X} 那婦人來拜他，說：「主啊，幫幫我！」 \end{tabularx} \\ \\ \relax
15:26 & \begin{tabularx}{0.7\textwidth}{X} 他回答：「拿孩子的餅丟給小狗吃是不妥的。」 \end{tabularx} \\ \\ \relax
15:27 & \begin{tabularx}{0.7\textwidth}{X} 婦人說：「主啊，不錯，可是小狗也吃牠主人桌上掉下來的碎屑。」 \end{tabularx} \\ \\ \relax
15:28 & \begin{tabularx}{0.7\textwidth}{X} 於是耶穌回答她說：「婦人，你的信心很大！照你所要的成全你吧。」從那時起，她的女兒就好了。 \end{tabularx} \\ \\ \relax
15:29 & \begin{tabularx}{0.7\textwidth}{X} 耶穌離開那地方，來到靠近加利利的海邊，就上山坐下。 \end{tabularx} \\ \\ \relax
15:30 & \begin{tabularx}{0.7\textwidth}{X} 有一大群人到他那裡，帶著瘸子、盲人、肢殘的、聾啞的，和好些別的病人，都放在他腳前，他就治好了他們。 \end{tabularx} \\ \\
[1ex]
\hline
\hline
\end{longtable}
講題:婦人,你的信心很大!

講員:高國雄牧師

經文:馬太福音15:21-30

日期:2023年8月20日

21耶穌離開那裡,退到推羅,西頓境內.

22有一個迦南婦人從那地方出來,喊著說:「主啊,大衛之子,可憐我!我女兒被鬼纏得很苦.」

23耶穌卻一言不答.門徒進前來,求他說:「這婦人在我們後頭喊叫,請打發她走吧.」

24耶穌回答:「我奉差遣只到以色列家迷失的羊那裡去.」

25那婦人來拜他,說:「主啊,幫幫我!」

26他回答:「拿孩子的餅丟給小狗吃是不妥的.」

27婦人說:「主啊,不錯,可是小狗也吃牠主人桌上掉下來的碎屑.」

28於是耶穌回答她說:「婦人,你的信心很大!照你所要的成全你吧.」從那時起,她的女兒就好了.—馬太福音15:21-28

24耶穌從那裡起身,往推羅境內去,進了一家,他不願意人知道,卻隱藏不住.

25立刻有一個婦人,她的小女兒被污靈附著,一聽見耶穌的事,就來俯伏在他腳前.

26這婦人是希臘人,屬敘利亞的腓尼基族.她求耶穌從她女兒身上趕出那鬼.

27耶穌對她說:「讓孩子們先吃飽,拿孩子的餅丟給小狗吃是不妥的.」

28婦人回答:「主啊,桌子底下的小狗也吃小孩子的碎屑呀!」

29耶穌對她說:「憑著這句話,你回去吧,鬼已經離開你的女兒了.」

30她就回家去,見小孩子躺在床上,鬼已經出去了.—馬可福音7:24-28

在一般人心目中,基督教信仰最多表達的就是愛心.

除了愛心這個很大的題目之外,學習信心的功課也常常在不同的信仰群體中.

很多基督徒對信心的學習視為生命中最重要的功課.

聖經中耶穌基督亦常常為門徒的信心有不同的表達.

大多時候會教訓他們小信,甚至問門徒為甚麼沒有信心?

今天這一段經文描述耶穌對這一個婦人的表述:一個非常重要的表述.

耶穌大大稱讚這婦人:「你的信心很大.」

這是我希望在開始時大家要記住的.

我們可能也期望自己能夠成為信心大的人.

到底在甚麼情況?甚麼境遇之下?我們的信心可以比現在大一點呢?

從經文開始:耶穌與門徒到推羅西頓的地方.

當時有一位迦南人.

馬可福音7:26這婦人是希臘人,屬敘利亞的腓尼基族.

是一位外邦人的婦女.

他一直跟著耶穌跟著門徒們,不斷的喊叫,懇求.

期望耶穌能夠幫助他.

首先我們要理解一下這婦人對耶穌的稱呼.

主呀!

大衛之子!

「主呀」這個稱呼可能也是我們容易比較明白的.

當時的社會中主僕的關係,會對所尊崇的稱呼為主.

但「大衛之子」是一個好很獨特的稱呼/描述.

大衛之子對猶太人來說非常清晰.

除了是指出君王後裔的身份,這也是對彌賽亞的稱號.

彌賽亞就是基督.

對猶太人而言,是他們整個民族的期盼,救贖的希望.

但對外國人而言,其實應該沒有什麼關係.

從這一段經文我們看出,這婦人有對耶穌獨特的認識.

雖然它不是猶太人,但他以猶太人期盼的拯救彌賽亞來形容耶穌,稱呼耶穌.

希望求耶穌可憐是因為他的女兒被污鬼纏繞了很久.

這個看似直接簡單的描述,但仿佛又看到深層的意義.

婦人期盼耶穌可憐的不是他的女兒,而是她自己.

很明顯,他作為母親為著女兒的苦.他也受了很多苦.

為著女兒因被污鬼纏繞的問題,相信這位母親也長期在困苦患難當中.

需要主的憐憫.

所以他求耶穌可憐他.

其實這也是一個很重要,亦很簡單的禱告.

婦人就是這樣直接地向耶穌乞求開恩憐憫.

我們都要學習.

V23耶穌基督卻的一言不發.我們可能不太明白其原因,可能也沒法明白.

或許耶穌希望給與門徒以及這婦人一點空間.

耶穌希望給予我們今天一點空間.當我們去到主面前懇求.可能那一刻我們經歷是主的沉默,可能那一刻我們要學習安靜等待.

而那一刻的等待,可能就是我們信心可以萌芽的時間.

V23門徒進前來,求他說:「這婦人在我們後頭喊叫,請打發她走吧.」

門徒的反應,既直接又真實且合理…

V24耶穌回答:「我奉差遣只到以色列家迷失的羊那裡去.」

讓我們難明困惑的:耶穌是否冷酷無情?上帝不是愛世人的嗎?

上帝的主權,上帝的揀選,上帝的節奏與次序

羅馬書1:16我不以福音為恥；這福音本是上帝的大能,要救一切相信的,先是猶太人,後是希臘人.

羅馬書2:

9他要把患難,困苦加給一切作惡的人,先是猶太人,後是希臘人；

10卻把榮耀,尊貴,平安加給一切行善的人,先是猶太人,後是希臘人.

11因為上帝不偏待人.

上帝的公平,上帝的計劃…

V24耶穌回答:「我奉差遣只到以色列家迷失的羊那裡去.」

這說話既是向婦人,同時也是回應門徒的.

25那婦人來拜他,說:「主啊,幫幫我!」

婦人再三懇求,耶穌接著回應.

26他回答:「拿孩子的餅丟給小狗吃是不妥的.」

容易誤會的經文:

猶太人習慣用狗來形容外邦人,早期有輕看的意味,後來成為習慣語.(有點類似香港人稱呼外國人為「鬼佬」)

但耶穌不是用他們慣常用的「狗」(在外邊,與人不同,不被看重),而是用「小狗」,即家犬,是寵物,是家中的一份子.

所以智慧的婦人聽得明天耶穌所描述的,更按主所給予的信心回應…

27婦人說:「主啊,不錯,可是小狗也吃牠主人桌上掉下來的碎屑.」

婦人深信上帝恩典豐厚,並且愛及萬人…

28於是耶穌回答她說:「婦人,你的信心很大!照你所要的成全你吧.」從那時起,她的女兒就好了.

耶穌不單馬上拯救婦人的女兒脫離困苦,更發出對這婦人最大的稱讚…

我們對主的恩典認識有多少呢?

我們是否完全信靠主的恩典夠用呢?

上帝所看重的是什麼呢?

你信上帝嗎?



\newpage

\section{施弗拉和普亞}
\label{sec:661}
\textbf{李均熊牧師 - 施弗拉和普亞}
\newline
\newline
\hyperref[sec:660]{< < < PREV SERMON < < <}
~
\hyperref[sec:toc]{[返目錄]}
~
\hyperref[sec:662]{> > > NEXT SERMON > > >}


講題:施弗拉和普亞

講員:李均熊牧師

日期:2023年8月27日

究竟是那一個法老?

歷史學家一直以來都很想知道以色列人在埃及為奴和出埃及的確實年期.究竟是那一個埃及法老的時代呢?現在普遍都認為是蘭塞二世(約主前十三世紀).不過,從另一個角度讀出埃及記,你可能會發現:聖經經文的作者根本不關心這些我們現代人所緊張的歷史問題.你可有發覺,當你細讀聖經所載由創世記到列王記的所謂歷史,埃及的王不是叫「法老」,就是直接叫「埃及王」,直到列王記上11:40節我們才有一個具名的法老王,叫「示撒」.有了法老的名字,歷史學家就可以透過眾多的埃及學的考古資料,透過這些「外証」去追溯這段時期的歷史事件究竟是如何發生.但出埃及記經文作者,則較少理會這些歷史學家的渴求,就是叫埃及王做法老:亞伯蘭在法老面前認撒萊做妹妹；約瑟下埃及成了法老的宰相；法老又吩咐民眾浸死以色列人的男嬰；摩西去見法老,要他放走以色列人；所羅門最後娶了法老的女兒……慢著……收養摩西的不也是法老的女兒嗎?那摩西和所羅門說不好有姻親關係了?此法老不同彼法老；但若不說清楚,所羅門就成了摩西的後父了.

施弗拉和普阿的角色

那麼經文作者關心甚麼呢?帝皇將相是歷史演義的慣常主角,但細看出埃及記第一及二章,不計開始時的族長,故事中只有三個人是記名的:施弗拉,普阿和摩西.摩西是當然的第一男主角,但其他配角繼施弗拉和普阿這兩個收生婆後,還有藏了孩子三個月的利未人媽媽,追著看弟弟命運的姐姐,大發慈悲的法老女兒和最後與摩西在井旁相遇的米甸祭司的七個女兒.這十二個女角,就是保住摩西生命的關鍵人物.但這眾多的人物,就只有施弗拉和普阿,名流千古.

施弗拉和普亞,她們其實有甚麼特別之處,可以流芳百世?有人說她們是第一對「公民抗命」的人民英雄,亦有人會疑惑她們為何說謊反得祝福.這些對她們行為的詮釋,看來只在乎詮釋者是宗教左派還是右派了.但神卻懶理你是左派抑或右派,「收生婆敬畏神,神便叫她們成立家室」.甚麼?只是成家立室?不錯,施弗拉和普亞只想「有啖安樂茶飯」,成家立室.她們並沒有甚麼政治的計算,也不理會這個或那個法老的手段.她們純綷因為「敬畏神」,就不照埃及王的吩咐去做,讓男孩活著.

她們面對埃及王的質問,也答得巧妙.有人問收生婆為何散謊,但你又怎知她們在哄騙這個法老王呢.用今日的說話,她們只是「躺平」一下.可能只是因為搭共車,慢了半拍,「收生婆還沒有到,她們已經生產了.」這個說自己如何用巧計的法老(出1:10),對收生婆的作為也沒有辦法了.

大歷史的諷刺

以色列民族的起始,充滿對大歷史的反諷:沒有法老可以留名,反而留名的是兩個敬畏神的收生婆；法老害怕戰爭和政治的反抗,收生婆只想成家立室；男人的計謀多多,暴力兇殺；但都不及眾女士合演的一場好戲精彩.教會內不單每一個人都重要,當中每一件小事也都不可忽略.是不願意的人,不顯眼的事,保存了生命,改變了民族的歷史.



\newpage

\section{焚而不毀的荊棘}
\label{sec:662}
\textbf{伍渭文牧師 - 焚而不毀的荊棘}
\newline
\newline
\hyperref[sec:661]{< < < PREV SERMON < < <}
~
\hyperref[sec:toc]{[返目錄]}
~
\hyperref[sec:663]{> > > NEXT SERMON > > >}


講題:焚而不毀的荊棘

講員:伍渭文牧師

日期:2023年9月3日

引言:夏威夷毛伊島(Maui)最近發生山火,近一百人罹難,經此一役,島上居民若再見到火燒著的荊棘,第一時間一定驚惶逃命.在希臘神話中,普羅米修斯(Prometheus)盜取天火,帶給人類文明,火帶給人類幸福.但對曾為埃及王子,熟讀戰爭歷史的摩西來說,火更代表毀滅:古代用火攻城掠地,用火堅壁清野,火燒連環船,都是利用火作為利器—毀滅的利器.但這時候的摩西,對火有不一樣的懷想.

1 神聖的火The Holy Fire

摩西終年一百二十,前四十年是埃及王子,次四十年是曠野牧人,最後四十年帶領以色列人出埃及.約四十歲時,他看見埃及人欺負希伯來人,觀看沒有人就把埃及人打死,藏在沙土裡.(出二11, 12) 他立刻逃到米甸的曠野,娶了米甸女子西坡拉(Zipporah)並在米甸曠野為岳父葉忒羅看羊.此刻摩西不再是埃及王子了,他離開王宮後,開始以亞伯拉罕的後裔的身份,思考自己的命途.

摩西看見焚而不毀的荊棘,同時也聽到上主的聲音.「不要近前來.當把你腳上的鞋脫下來,因為你所站之地是聖地.」(出三5) 火代表上主的臨在,這是聖地.曾在美國洛杉磯參加埃及正教科普特教會(Coptic Church)崇拜,領聖餐時信徒都脫了鞋趨近聖桌.

2 立約的火The Covenant Fire

曠野牧羊的摩西,離開王宮的權力富貴,臣僕簇擁,來到孤寂荒涼的曠野,曠野沒有地標的,天上北斗千萬年就在那裡為旅客導航.星空恢宏渺渺,埃及帝國祇是滄海的一粟,埃及國祚祇是蜉蝣一現；他想起上主曾應許先祖,後裔將如星宿繁多.上主說:「我是你父的上帝,是亞伯拉罕的上帝,以撒的上帝,雅各的上帝.」(出三6) 他想起立約的火把.上主和先祖亞伯拉罕立約當晚:「日落天黑,不料有冒煙的爐並燒著的火把從那些肉塊中經過.」(創十五17) 這是立約儀式,立約雙方行經切開一半的祭牲,表示若然爽約,就像祭牲被剖開,以生命為證,立約是切約 (cut the covenant).這約是用上主的獨生子的生命來保證,領聖餐時:耶穌說這是我的身體為你們而捨的.

3 祭壇的火The Altar Fire

荊棘中的火代表上主的臨在,在會幕和聖殿,有祭壇不滅的火,是第一次獻燔祭時,從天而降.每天早晚獻上祭牲的脂油,使火不熄滅.人若獻祭,須在獻壇點火,否則視為獻凡火遭上主擊殺,祭司亞倫兩個兒子拿達和亞比戶因獻凡火死了.(利十1, 2) 獻祭須從這祭壇點火,表示人要藉基督,才能來到父面前.基督就是那燔祭:「正如基督愛我們,為我們捨了自己,當作馨香的供物和祭物,獻與上帝.」(弗五2) 祭壇不滅的火所預表的基督,帶給我們不滅的永生.

4 荊棘的火The Bush Fire

「我要過去看這大異象,荊棘為何沒有被燒壞呢?」(出三3) 立約的火把,祭壇不滅的火,代表上主的臨在,聖潔和超越,屬於另一個世界.但摩西看到這超越神聖的火,陪伴卑下朽壞的荊棘,火不燒毀荊棘,而且讓荊棘在火中閃閃生輝；此刻荊棘經歷不朽壞,焚而不毀.而自己朽壞的身體,隨歲月漸漸衰弱,八十歲的牧羊人,已不及盛年時的埃及王子當年勇了.但神聖的火伴隨住容易朽壞的荊棘；你不須懼怕法老,我必與你同在.

摩西在焚而不毀的荊棘中還看見甚麼?他看見耶和華的使者.

耶和華的使者.

釋經學者認為耶和華的使者,就是道成肉身前的基督.一般被造的天使都有名字:向馬利亞報喜訊是加百列；天使長是米迦勒；守護約櫃是基路伯；在上主寶座旁侍立是撒拉弗.但耶和華的使者沒有名字.「耶和華的使者從荊棘裡火焰中向摩西顯現.」(出三2) 從來沒有人見過上主,祇有懷裡的獨生子將他表明出來.「(上主)那獨一不死,住在人不能靠近的光裡,是人未曾看見,也不能看見的,(基督)要將他表明出來.」(提前六16)

獻以撒.

耶和華使者在舊約出現,都與預表基督有關:亞伯拉罕獻以撒,手起刀落的一剎那,耶和華的使者出現:「耶和華的使者從天上呼叫他說……你不可在這童子身上下手,一點不可害他.」(創二十二11, 12) 但父為我們的罪,獻上他的獨生愛子作為贖罪祭.

伯特利天梯.

雅各冒認哥哥以掃,騙取父親臨終的祝福,離家出走,在伯特利作了一個夢,見到天梯,神的使者上去下來.「夢見一個梯子立在地上,梯子的頭頂著天,有神的使者在梯子上,上去下來.」(創二十八12) 上主沒有定雅各的罪,反而應許保護他,成就列祖的約:「我必與你同在,你無論往那裡去,我必保佑你,領你歸回這地.」(創二十八51) 耶穌在約翰福音說:「你們將要看見天開了,神的使者上去下來在人子身上.」(約一15)

因為耶和華的使者,因為道成肉身的基督,我們朽壞凡軀,得著不滅之身.

奧本海默.最近上畫的電影《奧本海默》(Oppenheimer)講述原子彈試爆故事.他是原子彈之父,又稱美國的普羅米修斯(American Prometheus).按鈕試爆成功,全場靜止,時間停住,烈焰滔天,從此世界已不一樣.奧本海默跟著說:「Now I Am Become Death, the Destroyer of Deaths (我已成為死神,世界的毀滅者).」這話引自印度古詩《薄伽梵歌》(Bhagavad-Gita).

但上主無窮的能力,保護著荊棘；他的能力用來建立,不是毀滅.

願你堅立.

摩西在野外晚上放羊,每次仰望星空,都會想起上主對先祖的應許:後裔如海沙之多,星宿之繁茂.但在他心底深處,他還是嚮往不朽的生命.詩九十篇是摩西回顧一生的詩篇,這主題哀嘆人生的短暫,行經時間,渴望被上主紀念:「你叫他們如水沖去；他們如睡一覺.早晨如生長的草,早晨發芽生長,晚上割下枯乾.……願你堅立我們手所作的工；我們手所作的工,願你堅立.」」(詩九十5, 6, 17)

喬布斯.

傳記作者華特.艾薩克森(Walter Isaacson)與喬布斯(Steve Job)作了超過四十次訪談,最後一次在他家中的後花園進行,是陽光普照的下午,當時喬布斯身體很不舒服,喬布斯把話題轉到宗教和死亡.當問及他是否相信死後有永生(afterlife),是否相信有上帝時,他說:「五十五十.這是大奧秘,人永遠不知道,但我希望死不是生命的終結,我但願累積的智慧不會隨死亡歸於無有,可以永續.但他跟著說:但可能它真的像開關制,click一下,你就離開了；停頓一秒後他笑說:所以我們蘋果產品沒有開關制的.」

跟喬布斯不同,摩西知道上主會堅立他手所作的工,因為在那一天,摩西看見上主的火,沒有焚毀荊棘,本來遇風折枝,遇火成灰的荊棘,竟然絲毫無損.其實每個人都希望凍齡,逆齡,返老還童,希望不朽,希望永生.

5 朽壞的荊棘The Mortal Bush

但人生必要面對身體自然的朽壞, 李察.尼德姆(Richard Needham)令人深思的人生七步曲:spills(嬰孩吐奶), drills(少年運動), thrills(求偶興奮), bills(成年帳單), ills(老年多病), pills(藥不離身), and wills(立下遺囑).

死亡的破壞.

保羅田立克(Paul Tillich)在一篇名為「死亡的破壞」(Destruction of Death)的講章說:「死亡不單是著名古老比諭所說像較剪,把生命的線剪斷.死亡更像貫穿生命其中的一條線,編織進我們的存在,從開始到結束.」身體老化,凡軀朽壞,我們都要死去的事實,每一刻都影響著我們的身體和靈魂.

年輕財主.

渴望永生,是生命不能承受的重.聖經記載一位富有的年輕人,公開在路上跑來跪在耶穌前面求問:「良善的夫子,我當作甚麼事才可以承受永生?」(可十17) 這人認為永生是可以作一些事來承受的；第二,承受表示永生是人死後的事情.耶穌回答:「為何你稱我是良善的?除了上主沒有一位是良善的.誠命你是曉得的:不可殺人,不可奸淫,不可偷盜,不可作假見証…….年輕財主答:這一切從少都遵守了.」(可十18-20)

讀出律法時,耶穌沒有提及另一塊法版:「除了我以外,你不可有別的神.」因為財主的神就是他的財富.耶穌說:「你還缺少一件,去變賣你所有的,分給窮人,就必有財寶在天上,你還要來跟從我.」永生不能承受,是藉悔改才得到的恩典；永生不是死後的事,跟隨基督就是永生.今天福音就解釋如何跟從耶穌.

6 為我喪掉生命Losing Life for Me

「若有人要跟從我,就當捨己,背起他的十字架來跟從我.因為凡要救自己生命的,必喪掉生命；凡為我喪掉生命的,必得著生命.人若賺得全世界,賠上自己的生命,有甚麼益處呢?人還能拿甚麼換生命呢?」(太十六24-26) ”[i]

點著的洋燭.

人生像點著的洋燭,愈來愈短；過一天就賠上一天,喪掉一天.我們可以為基督喪掉這(一天)生命,也可以為自己喪掉這(一天)生命.凡要救自己生命的,必喪掉生命:凡為自己,不是為基督用這一生,就喪掉這一生,虛度這一生.但為基督過這一生,必得著這一生.

永生與不朽.

聖經說永生,罕談不朽,因為祇有上主獨一不死,與他結連就是永生,這關係連死亡也不能間斷.對牙痛的人來說,永恆是可怕的.躺在床上全身癱瘓,靠管導餵養,不能說話,永恆更是夢魘.所以當人吃了禁果後,被趕出樂園,恐怕他們再吃生命樹的果子(創三22),永遠活著,帶著罪永遠活著多痛苦.

格式化的生命.

朽壞死亡有它的功用.死亡為人生作了限制,就像電腦的「格式化」(Formating),沒有格式化的硬碟不能使用.[1]老師給學生作業也設定「死線」(Deadline),沒有死線,學員失去用功的目標和動力.修讀學位也有完成學程的年限.不設年限的學程不能鼓勵學習,看不到終點,也失去衝過畢業終點的想像力.沒有終點,也失去過程；過程是起點和終點的連線.生命有終結,才迫使我們珍惜光陰,作出選擇,體會懊悔與得失,認識智慧和愚昧.生命的終點,是生命圓滿的彩帶,行經時間,完成上主交付的工作.保羅說:「那美好的仗我已經打過了,當跑的路我已經跑盡了.所信的道我已經守住了.」(提後四7)

今生很短.

倪匡在一次聚會中說:「達摩面壁九年悟出道理,我在美國十年讀聖經,認識真諦:「今生很短,來生很長,不值得今生勞勞碌碌,要為來生打算.」

上星期有一天旁晚逛大衛新商場,看見將開業買衣服的店寫了宣傳口號:Life Is Too Short to Wear Boring Clothes (人生苦短,何必穿著沉悶的衣服呢?) 我在臉書即興寫了下:是的,除了披戴基督,世上所有衣飾皆沉悶,人生真的太短暫.

牟敦梅頓 (Thomas Merton)的一生,可說明除了披戴基督,始得滿足,其它衣飾如糠秕,看似金光華麗,經風一吹了無痕.他是一個備受愛戴,才華橫溢的詩人,政治評論家.世上的榮華,榮譽不能滿足他,心靈的空處惟上主始能填滿,他追求修道院獨處生活.他的自傳《七重山》(The Seven Storey Mountain) 自一九四八年出版以來即風行全球,譯成二十多種文字.《紐約時報》將之譽為二十世紀的《懺悔錄》. 此書帶動美國年輕一代的屬靈覺醒,尤其響往修道院生活.)他效法奧古斯丁的懺悔錄,作了壓卷的禱告:

【福哉獨處!】

我喜愛獨處唯一的原因——

在被造的東西中消失,

向它們死去,不知道它們,

因為它們提醒我離你多麼遠.

它們提醒我:

你和這些相距甚遠,

雖然你在其中.

你創造它們,它們因你而存在,它們把你遮蓋起來,

讓我看不見你.

我將獨自生活,離開它們.

福哉獨處![2]

荊棘冠冕.

耶穌憑甚麼要求我們跟隨他呢?請看太十六21:「從此,耶穌才指示門徒,他必須上耶路撒冷去,受長老,祭司長,文士許多的苦,並且被殺,第三日復活.」荊棘焚而不毀,因基督把荊棘將做成冠冕,戴在頭頂,走上十字架,第三天復活,帶給我們永生.

7 遠象的火The Visonary Fire

焚而不毀不是火與荊棘和平共存,互不相干,生活在平行的空間；乃火在荊棘中,荊棘在火中；像燒至火紅的鐵,鐵在火中,火在鐵中,被火滲透的荊棘,此刻成為照明的火把,滿有遠象的火把.

在荊棘焚而不毀的異象中,上主第一次稱以色列人為「我的百姓」.「我的百姓在埃及所受的困苦,我實在看見了他們受督工的轄制所發的哀聲,我也聽見了.我原知道他們的痛苦,我下來是要救他們脫離埃及人的手,領他們出了那地,到美好寬闊,流奶與蜜之地,就是到迦南人,赫人,亞摩利人,比利洗人,希未人,耶布斯人之地.……故此我要打發你去見法老,使你可以將我的百姓以色列人從埃及領出來.」(出三7-10)

進入迦南應許美地前,以色列人要在西乃山敬拜上主,與上主立約,領受律法.「如今你們若聽從我的話,遵守我的約,就要在萬民中作屬我的子民,因為全地都是我的.你們要歸我作祭司的國度,為聖潔的國民.」(出十九4-6)

很明顯,進迦南應許之地有宣教的使命,上主賜福給亞伯拉罕,目的是別人因他而得福,因此他們進入迦南先要在西乃山立約學習上主的律法.以色列人要成為火把,照亮迦南人,離開虛妄的偶像,轉向真神.

迦南人的宗教文化極其邪淫迷信,廟宇常駐男,女祭司,藉行淫進行宗教禮儀,獸交很普遍,又把兒童作祭牲獻祭.上主要清除這些罪惡,一方便保護選民不受異教影響信仰,另一方便把信仰傳給迦南人.

喇合就看到以色列上帝才是真實可靠的神.「耶和華你們上帝,本是上天下地的上帝……要救活我的父母,弟兄,姐妹,和一切屬他們的,拯救我們性命不死.」(書二11, 13) 喇合認識迦南宗教諸偶像的虛妄,以色列的上帝才是天地的主.

王爾德說:「我們都躺在陰溝裡,但有些人仰視晚上的星空.」喇合不單得救,而且成為耶穌的先祖,上主並不偏待人.

8 燃點靈火The Spiritual Fire

我們的本體也像朽壞的荊棘,遇風枝折,遇火成灰；但被聖靈的火陶造後,也能成為照亮這世界的火把.五旬節像火降在門徒頭上的聖靈,也降在我們的身上.今天書信羅十二9-11:「愛人不可虛假.惡要厭惡；善要親近.愛弟兄,要彼此親熱；恭敬人,要彼此推讓.殷勤,不可懶惰；要心裡火熱,常常服事主.

詩篇經課指出生命須經熬煉,向上主敞開:「耶和華啊,求你察看我,試驗我,熬煉我的肺腑心腸.因為你的慈愛常在我眼前,我也按你的真理而行.」(詩二十六2, 3)

受真理熬煉陶造的生命,漸漸成為照亮人的火把,照亮面對黑暗的人,這火把給人溫暖,伴隨著人看來卑微,被人忽視,被邊緣化的人.

在正教的信仰傳統,明天九月四日記念焚而不毀荊棘節慶,馬利亞就是焚而不毀的荊棘[3].朽壞的身體被上主使用；聖靈感孕誕下聖嬰,仍保有其童貞.信徒也懷有基督,也能藉生命顯揚基督的香氣. 結論:我們本是朽壞的荊棘,但基督把我們戴在頭上,作為冠冕,走上十字架；聖靈在我們心中燃點,聖道陶造我們的生命,成為照亮他人的火把.為人的需要燃點生命,就是為耶穌喪掉生命,就必得著生命,經歷真正滿足的生命.

人若賺得全世界,賠上自己的生命,有甚麼益處呢?人還能拿甚麼換生命呢?人若賺得全世界,賠上自己的生命,有甚麼益處呢?人還能拿甚麼換生命呢?

[1] 新購買的磁碟在使用之前,要能讓作業系統認得它,先寫入一些磁性的記號到磁碟上儲藏資料每一個扇區(Disk sector),便可在該作業系統下取用磁碟上的資料,這個動作就稱為格式化.

[2] That is the only reason why I desire solitude— to be lost to all created things, to die to them and to the knowledge of them, for they remind me of my distance from You. They tell me something about You: that You are far from them, even though You are in them. You have made them and Your presence sustains their being, and they hide You from me. And I would live alone, and out of them. O beata solitudo! –Thomas Merton, The Seven Storey Mountain. An Autobiography of Faith (多馬斯.梅頓,七重山.信仰自傳),1998:461.

[3] 有關聖像(Icon)用重疊的鑽石,各有四角.紅色代表火在外,在內綠色上主之母(Theotokos)馬利亞手抱聖嬰.八個角象徵屬天:挪亞一家八口進方舟；嬰孩第八天行割禮；六日創造,第七日安息日結束是受造世界；主日是基督復活的第八天,開始新造的世界

...Copyright (C) 本網站上所有視頻／音頻／文字版權歸【葡萄樹傳媒】所有,轉載請註明出處! \url{https://www.vinemedia.org/course/bible-teaching/sunday-service-1/the-burning-bush/} .



\newpage

\section{逾越抑保護}
\label{sec:663}
\textbf{李均熊牧師 - 逾越抑保護}
\newline
\newline
\hyperref[sec:662]{< < < PREV SERMON < < <}
~
\hyperref[sec:toc]{[返目錄]}
~
\hyperref[sec:664]{> > > NEXT SERMON > > >}


講題:逾越抑保護? Passover or Protect?

講員:李均熊牧師

場合:香港中文大學崇基學院禮拜堂主日崇拜

日期:2023年9月10日

引言:新學期,新開始

有人說現今的中國人有兩個新年,一個新曆新年和一個農曆新年,所以新年立願可以做兩次.但其實現今社會不只有兩個新年,不同公司甚至政府都有自己的「財政年度」.大學新的學年也正好在上個星期開始了.每一個新開始,我們都可以回顧過去,展望將來.

聖經記載以色列人也有不止一個的新年.現今的猶太人新年叫作Rosh Hashanah,那其實是民數記廿九章所記的吹角節,就是猶太人秋天慶祝的新年,乃是猶太曆的七月,約西曆九至十月間的日子,今年(2023)就在下星期五9月15至9月17星期日.但我們基督徒可能更加熟識,就是聖經所記載以色列人在春天慶祝的宗教信仰上的新年.在出12:1-2記載:「耶和華在埃及地曉諭摩西,亞倫說:『你們要以本月為正月,為一年之首……』」所以以色列人出埃及,吃逾越節晚餐的那一個月,才是猶太曆的一月.

現代猶太人的曆法就如中國人的農曆,混合了陰曆和陽曆,兩者均同時依靠月亮圓缺和不同時期的日照時間來推算日子,所以是比較複雜的.但出埃及記這裡的記載,是以色列人由奴役的生活轉為自由生活的新開始,也提醒我們作為神子民,我們的日子安排不再是依從日月星晨的自然現象,也不是依從農耕經濟生活周期的要求,日子乃由神的命令作定位,信仰成為生活的新坐標.這裡再一次叫我們反思,崇奉基督只是令我們多了星期日返教會的活動,是我們生活可有可無的點綴；還是我們日常生活行事為人的定位和依據?

今日舊約經課正好讀到這個逾越節的設立,就讓這個星期,作為我們新學年的開始,細想這段經文的意義.

起:逾越抑保護?

出埃及記12章1-13節記載上主吩咐摩西逾越節應如何遵守.

我們基督徒相信,耶穌就是那一隻沒有殘疾的羔羊,為救我們,已變成為那永遠屬天的逾越節祭物,他血已經為我們獻上了,使我們可以從罪的困鎖中釋放,生命可以自由和滿有恩典的去活.加上逾越節經過很多世代的流變,今日的慶祝方式和當日第一個逾越節的舉行方法可能已經有很大的出入,所以這裡所記載的很多細節對於今日的我們基督徒可能沒有太大意義.但我們或可以留意第12和13節:羊羔的血成為記號,而滅命的使者見到那一家有這個記號,就會「越過」那一家.所以我們叫這一個節令為「逾越節」(Passover).只不過,代表「逾越」的那一個希伯來文Pesah,亦是今日猶太人稱呼自己節期的字,本意其實不是代表「越過」.其基本的用意乃指「保護」.這個意思,在以賽亞書31:5就十分之明顯:「雀鳥怎樣搧翅覆雛,萬軍之耶和華也要照樣保護耶穌撒冷,他必保謢拯救,要越門保守.」什麼越門保守,字根就是Pesah.將Pesah翻譯成「越門保守」,根本不配合上文講到雀鳥用翅旁覆蓋幼雛的舉動.此句應該翻譯為「就如雀鳥搧翅覆雛,萬軍之耶和華也會圍護耶路撒冷,他必圍護拯救,保護搭救.」

承:主動保守的記念

若果我們將保護套用在出埃及記十二13,意思就並不是上帝越過血在門上的一家,乃是「我一見這血,就保護你們」是上帝主動的保護,使災禍離開我們.這和詩篇廿三篇「我們雖然行為死蔭幽谷,也不怕遭害,因為你與我同在」,和主禱文「救我們脫離兇惡」都有異曲同功的意思.

危難和兇險基本上時時在我們身邊出現,但是上帝主動的保護和搭救,我們才不至滅亡.很多時我們皆沒有用心去留意上帝的搭救和保護,以為事事順利,生活平安是理所當然的.【鎖車例子:下意識之下,我自然就鎖了車,但也常回頭想,我究竟鎖了沒有?】新的學年,讓我們先回顧一下過去日子,我們怎樣避過了生活的兇險,我們要提醒自己,使我們免受患難的是上主主動的幫助和保守,我們其實都要學習去記念.

轉:敘事儀式化

雖然今天我們未必跟從經文的內容照做,但正如兩星期前我的講道,讀聖經不單要留意經文說了甚麼,也留意經文怎樣說.若我們仔細去看經文的舖排,出埃及記由12章到13章是指引和敘事相間的.

由出埃及記開始(甚至可以說從創世記開始)到出埃及記11章,都是一段一段故事的敘述,但到了出12章,敘述卻變成上主透過摩西向以色列人吩咐遵守逾越節的律例.猶太人很早就留意這一點.猶太人雖然稱呼五經為「律法書」,但要延至出埃及記12章,才看到上主吩咐律法!

這樣將敘事變成儀式,又或所謂儀式化敘事,目的不單要神的子民清楚知道他們得救贖得自由的歷史事件是如何發生,更要神的子民世世代代透過儀式去銘記塑造他們民族生命的基礎事件,更透過儀式化使敘述成為不是單一,一次過的事件,乃是神子民世世代代都可以透過儀式所包含的文化象徵,重新體會這個事件的效果和意義.換句話說,經文並不是簡單要神的子民記得上主的拯救,更要使這一份記念帶著一份儀式感,使這一個神子民的歷史經驗,化成每年充滿象徵意義的民族生活文化中的儀式.出埃及記12章的表達形式,將民族的基礎敘事,慶祝儀式和它們的象徵意義,都結合起來(出12:14).

近來很多網絡文章都說我們生活需要「儀式感」.若你Google「儀式感」都會彈出內容大同小異的網頁.通常這些網頁也會引用法國經典文學《小王子》的內容解釋甚麼是儀式,以下我就直接引用其中一篇網頁文章的內容:

儀式感意思是什麼呢?探討儀式感定義前先從經典童話《小王子》談起

狐狸對小王子說:「如果你下午四點鐘來,從三點鐘開始,我就開始感覺到快樂,時間越臨近,我就越來越感到快樂.到了四點鐘的時候,我就會坐立不安,我發現了幸福的價值,但是如果你隨時什麼時候來,我就不知道在什麼時候準備好迎接你的心情了,……這應當有一定的儀式.」

「儀式是什麼?」小王子問道.

「這也是經常被遺忘的事情.」狐狸說,「它就是使某一天與其他日子不同,使某一時刻與其他時刻不同.」

大家一定沒想過,儀式感定義是來自《小王子》中的狐狸說的話,他想傳達的儀式感意思是賦予生活中一些看似瑣碎的細節特殊的意義,為平淡的日子增添幾分詩意,讓生活不只是湊合過日子.

大部份這些網頁的內容其實是參考Karan Johnson在BBC網站於2021年10月25譔寫的一篇題為〈日常儀式的驚人力量:為什麼生活需要有儀式感〉(崇拜之後大家不妨也搜查一下).文章說到心理學家將儀式定義為「一組預先設定好的象徵性動作序列,通常以形式和重覆為特徵,缺乏直接的效用目的」.文章又說,「儀式的作用遠不止幫助我們實現價值觀.它們也可能讓我們不那麼焦慮.」

「儀式性的做法有助於給不確定的未來帶來一定程度的可預測性.科學家稱,它們讓人類大腦相信存在穩定和可預測性,是『對抗不確定性和焦慮的習慣性緩衝器』.」

合:重尋儀式的安全感

儀式或是儀式感對現代人來說其實並不陌生,只是很多活動我們不再稱呼它為儀式而矣,也未必當成儀式去看待.節日,生日的飯聚,甚至開學禮,開學崇拜,宿舍的學Sem飯等等,又甚至是起床的刷牙洗面,也可以做得很有儀式感.有無試過因為遲了起床,沒有完成平時的起床「儀式」便出門口,整天生活就好像不太完整.(疫情三年沒有了O camp?)每主日返教會崇拜,每年的聖誕復活節,甚或今天經課的逾越節,就是透過一些重覆的儀式,讓我們的生活文化與信仰象徵有機會結合,讓有時平淡,有時混亂的生活,添上意義和方向,重回穩定,可測,安全的感覺,成為「對抗不確定性和焦慮的習慣性緩衝器」.

逾越節可能只是一年慶祝一次,又或我們基督徒未必會慶祝逾越節,但上主的保護和救贖,需要我們時常的紀念.新學期的開始,讓我們重拾失落的儀式感,再生活中發現上主的保護帶領,或者透過每天的讀經祈禱靈修,每星期的主日崇拜,每年的節日省思,深化基督對我們生命的影響,在暗昧不明的日子,又或有時焦慮不確的生活中,使我們找到安息與鼓勵,又或召命和方向,幫助我們每天勇敢的前行.



\newpage

\section{主啊如若弟兄得罪我}
\label{sec:664}
\textbf{高國雄牧師 - 主啊,如若弟兄得罪我}
\newline
\newline
\hyperref[sec:663]{< < < PREV SERMON < < <}
~
\hyperref[sec:toc]{[返目錄]}
~
\hyperref[sec:665]{> > > NEXT SERMON > > >}


馬太福音 18:21-35
\newline
\begin{longtable}{cl}
\hline
\hline
章節 & 經文 (和合本修訂版)\\
\hline
18:21 & \begin{tabularx}{0.7\textwidth}{X} 那時，彼得進前來，對耶穌說：「主啊，我弟兄得罪我，我當饒恕他幾次呢？到七次夠嗎？」 \end{tabularx} \\ \\ \relax
18:22 & \begin{tabularx}{0.7\textwidth}{X} 耶穌說：「我告訴你，不是到七次，而是到七十個七次。 \end{tabularx} \\ \\ \relax
18:23 & \begin{tabularx}{0.7\textwidth}{X} 因為天國好像一個王要和他僕人算賬。 \end{tabularx} \\ \\ \relax
18:24 & \begin{tabularx}{0.7\textwidth}{X} 他開始算的時候，有人帶了一個欠一萬他連得的僕人來。 \end{tabularx} \\ \\ \relax
18:25 & \begin{tabularx}{0.7\textwidth}{X} 因為他沒有甚麼償還之物，主人下令把他和他妻子兒女，以及一切所有的都賣了來償還。 \end{tabularx} \\ \\ \relax
18:26 & \begin{tabularx}{0.7\textwidth}{X} 那僕人就俯伏向他叩頭，說：『寬容我吧，我都會還你的。』 \end{tabularx} \\ \\ \relax
18:27 & \begin{tabularx}{0.7\textwidth}{X} 那僕人的主人就動了慈心，把他釋放了，並且免了他的債。 \end{tabularx} \\ \\ \relax
18:28 & \begin{tabularx}{0.7\textwidth}{X} 那僕人出來，遇見一個欠他一百個銀幣的同伴，就揪著他，扼住他的喉嚨，說：『把你所欠的還我！』 \end{tabularx} \\ \\ \relax
18:29 & \begin{tabularx}{0.7\textwidth}{X} 他的同伴就俯伏央求他，說：『寬容我吧，我會還你的。』 \end{tabularx} \\ \\ \relax
18:30 & \begin{tabularx}{0.7\textwidth}{X} 他不肯，卻把他下在監裡，直到他還了所欠的債。 \end{tabularx} \\ \\ \relax
18:31 & \begin{tabularx}{0.7\textwidth}{X} 同伴們看見他所做的事就很悲憤，把這一切的事都告訴了主人。 \end{tabularx} \\ \\ \relax
18:32 & \begin{tabularx}{0.7\textwidth}{X} 於是主人叫了他來，對他說：『你這惡奴才！你央求我，我就把你所欠的都免了； \end{tabularx} \\ \\ \relax
18:33 & \begin{tabularx}{0.7\textwidth}{X} 你不應該憐憫你的同伴，像我憐憫你嗎？』 \end{tabularx} \\ \\ \relax
18:34 & \begin{tabularx}{0.7\textwidth}{X} 主人就大怒，把他交給司刑的，直到他還清了所欠的債。 \end{tabularx} \\ \\ \relax
18:35 & \begin{tabularx}{0.7\textwidth}{X} 你們各人若不從心裡饒恕你的弟兄，我天父也要這樣待你們。」 \end{tabularx} \\ \\
[1ex]
\hline
\hline
\end{longtable}
講題:主啊,如若弟兄得罪我…… Lord, If My Brother Sins Against Me……

講員:高國雄牧師

場合:香港中文大學崇基學院禮拜堂主日崇拜

經文:馬太福音18:21-35

日期:2023年9月17日

馬太福音 18:21-35

21那時,彼得進前來,對耶穌說:「主啊,我弟兄得罪我,我當饒恕他幾次呢?到七次可以嗎?」22耶穌說:「我對你說,不是到七次,乃是到七十個七次.23天國好像一個王要和他僕人算帳.24才算的時候,有人帶了一個欠一千萬銀子的來.25因為他沒有甚麼償還之物,主人吩咐把他和他妻子兒女,並一切所有的都賣了償還.26那僕人就俯伏拜他,說:『主啊,寬容我,將來我都要還清.』27那僕人的主人就動了慈心,把他釋放了,並且免了他的債.

28「那僕人出來,遇見他的一個同伴欠他十兩銀子,便揪著他,掐住他的喉嚨, 說:『你把所欠的還我!』29他的同伴就俯伏央求他,說:『寬容我吧,將來我必還清.』30他不肯,竟去把他下在監裡,等他還了所欠的債.31眾同伴看見他所做的事就甚憂愁,去把這事都告訴了主人.32於是主人叫了他來,對他說:『你這惡奴才!你央求我,我就把你所欠的都免了,33你不應當憐恤你的同伴,像我憐恤你嗎?』34主人就大怒,把他交給掌刑的,等他還清了所欠的債.35你們各人若不從心裡饒恕你的弟兄,我天父也要這樣待你們了.」

饒恕很難

信徒特別要有好的「道德」行為,要行多幾多步?七次的意義:完美的「七」七十個七=490次嗎?比喻的重點及意義:王有理嗎?不是已經免了債嗎?重新追討?上帝會如此「不合理」嗎?惡奴才,惡在那裡?

V33「應當」:主禱文:「免我們的債,如同我們免了人的債…」這「如同」,是因果?是換取?約翰一書4:19「我們愛,因為上帝先愛我們.」羅馬書5:8「惟有基督在我們還作罪人的時候為我們死,上帝的愛就在此向我們顯明了」上帝主動先付出,先愛我們,先主動為我們預備「赦罪的恩典」

我們要接受,就要回應,要認罪,悔改…(路加福音15)浪子回頭的比喻中,浪子的認罪自述:「父親!我得罪了天,又得罪了你, 從今以後,我不配稱為你的兒子,把我當作一個雇工吧.」承認過失,立志改變行為,才是完整/真誠/真實的悔改.當然我們也得承認,我們是罪人,在罪的世界中…

羅馬書7:18-19「我也知道,住在我裡面的,就是我肉體之中,沒有善.因為立志為善由得我,只是行出來由不得我.我所願意的善,我不去做；我所不願意的惡,我反而去做.」

但重點在:V35你們各人若不從心裡饒恕你的弟兄,我天父也要這樣待你們了.「從心裡」就是核心.

請再一次留意經文

26那僕人就俯伏拜他,說:『主啊,寬容我,將來我都要還清.』27那僕人的主人就動了慈心,把他釋放了,並且免了他的債.祈求寬容的是時間,而不是債務內容…免去的是全部,是一切的債.上主所賜的恩典是無價的,不可換取,不能比較,只需接受.真心的接受,必須要內含對恩典的認同.上主願意赦免是因為愛,我們要接受赦免,也接受了愛,並認同要以愛去赦免,所以也必要以此對待他人.再回想一開始時,彼得的提問:「主啊,如若弟兄得罪我,我當……」人有時難免會自以為義,很容易只看見別人的過犯,弱點,做得不完美的地方…今天的投訴文化,自身權益的爭取,繼上輕易的惡意批評…希望當下一次我們再向上帝問同一個問題時,願意先自省,學習路加福音中的稅吏

18:13那稅吏遠遠地站著,連舉目望天也不敢,只捶著胸,說:『上帝啊,開恩可憐我這個罪人!』



\newpage

\section{大能的摩西與無能的基督}
\label{sec:665}
\textbf{林豪恩先生 - 大能的摩西與無能的基督}
\newline
\newline
\hyperref[sec:664]{< < < PREV SERMON < < <}
~
\hyperref[sec:toc]{[返目錄]}
~
\hyperref[sec:666]{> > > NEXT SERMON > > >}


講題:大能的摩西與無能的基督 The Powerful Moses and the Powerless Christ

講員:林豪恩先生

場合:香港中文大學崇基學院禮拜堂主日崇拜

日期:2023年10月1日

1. 耶穌被質疑(馬太福音)

21:23耶穌進了殿,正教訓人的時候,祭司長和民間的長老來問他說:「你仗著甚麼權柄做這些事?給你這權柄的是誰呢?」

按福音書記載,耶穌不只一次被質疑.從開始傳道至被釘十字架,耶穌都不斷被質疑,要求證明自己.

你曾經被質疑嗎?質疑你的能力?質疑你的身分?你通常有甚麼感受呢?感到憤怒?感到委屈?你通常是如何回應呢?反駁?自衛?呈強?還是逃走呢?

2. 摩西被質疑(出埃及記)

17:1以色列全會眾都遵耶和華的吩咐,按著站口從汛的曠野往前行,在利非訂安營.百姓沒有水喝,

17:2所以與摩西爭鬧,說:「給我們水喝吧!」摩西對他們說:「你們為甚麼與我爭鬧?為甚麼試探耶和華呢?」

17:3百姓在那裡甚渴,要喝水,就向摩西發怨言,說:「你為甚麼將我們從埃及領出來,使我們和我們的兒女並牲畜都渴死呢?」

不知道你有沒有感到奇怪?百姓沒有水,竟然要摩西給他們水.如果我們現在沒有水,要求你給我們水,你會怎樣回應?

「無水呀?開水喉啦!」

「水喉沒有水出來呀!有就唔使叫你啦!」

「既然水喉都沒有水出來,你怎可能要求我給你們水呢?你把我當作甚麼呢?你要水,去找那個林豪恩啦,因為他叫做阿水.」(林豪恩:對不起,我雖然叫做阿水,但都不可以無中生水,今次真係幫唔到你!)

在曠野的以色列人沒有水,就來質問摩西,為甚麼呢?可能認為摩西要為現在的危機負責,也可能他們認為摩西有能力無中生水.他們要求摩西無中生水,你認為合理嗎?不久之前,摩西舉起手中的仗,紅海的水便立即分開.又不久之前,沒有食物,摩西就叫天降嗎哪.現在沒有水,去找摩西,豈不是合理期望嗎?這個大能的摩西有甚麼是做不到的呢?這個大能的摩西有甚麼問題是解決不了的呢?

3. 大能的摩西

17:5耶和華對摩西說:「你手裡拿著你先前擊打河水的杖,帶領以色列的幾個長老,從百姓面前走過去.

17:6我必在何烈的磐石那裡,站在你面前.你要擊打磐石,從磐石裡必有水流出來,使百姓可以喝.」摩西就在以色列的長老眼前這樣行了.

摩西果然不負眾望,他在以色列的長老眼前,擊打磐石,磐石裡就有水流出來.摩西的大能又一次被證實,滿足民眾的能力又一次被肯定.因此,後世的以色列人一直視摩西為英雄.如果有拯救者,那人一定似摩西,是摩西再臨.

基督徒也被呼籲學習摩西.不少書籍名為「摩西領袖學」,有些文章名為「向摩西學習談判技巧」,有些會課程名為「從摩西看僕人的訓練及成長」,有些講座名為「向摩西學習健康長壽」,有些訓練名為「效法摩西的信心」.摩西的確是信徒的模範,他靈力充沛,信心巨大.大家都想為上帝做好見證,希望能夠透過自己見證上帝的大能.例如當法老質疑摩西及摩西的上帝,摩西就向法老顯示上帝給予他的神力,把手上的杖變蛇,又把蛇變回手杖；更進一步,摩西向法老宣告的災禍不是「得個講字」,而是會兌現的.如果我們的信仰和能力都像摩西一樣就好了,不至於「講就天下無敵,做就有心無力」.

4. 耶穌對質疑的奇怪回應

耶穌被質疑「仗著甚麼權柄做這些事」是甚麼事呢?按上文來理解,很可能是指耶穌潔淨聖殿的行動.如果是摩西式的回答,很可能說:「當然是靠著大能的上帝的權柄!」然後展示一下作為證明.

耶穌沒有這樣做,卻跟他們繞圈子,最後竟然說:「我也不告訴你們我仗著甚麼權柄做這些事.」

如果你真是有能力的,請你出示證明都合理呢!然而,耶穌面對質疑,或者要求證明他的身分,從開始至被釘在十字架上,他的回應到令人感到奇怪,令人覺得他一直拒絕顯示能力,迴避證明自己.

耶穌開始傳道時,魔鬼挑戰他:

「你若是上帝的兒子,可以吩咐這些石頭變成食物.」(馬太福音4:3)

「你若是上帝的兒子,可以在(殿頂)跳下去,因為經上記著說:主要為你吩咐他的使者用手托著你,免得你的腳碰在石頭上.」(馬太福音4:6)

耶穌沒有把石頭變成食物,也沒有在(殿頂)跳下去,他似乎拒絕用這些方法去回應對他的身分的質疑.快要死了,把握最後機會展示一下能力吧!耶穌被釘在十字架上時,旁邊的強盜嘲笑他:

27:40「你這拆毀聖殿,三日又建造起來的,可以救自己吧!你如果是上帝的兒子,就從十字架上下來吧!」

耶穌就是死在十字架上,也不從十字架下來以證實自己是上帝的兒子.

下面的官長嘲諷他:

27:42「他救了別人,不能救自己.他是以色列的王,現在可以從十字架上下來,我們就信他.

耶穌就是在最後一刻,都沒有從十字架下來滿足那些人對於相信他是以色列的王的要求.

27:43他倚靠上帝,上帝若喜悅他,現在可以救他；因為他曾說:『我是上帝的兒子.』」

上帝任由耶穌死去,也沒有救他滿足人們對於上帝喜悅他的證明.

耶穌以如此無能的方式回應挑戰,以如此無能的方式回應質疑,可見他和摩西的分別何等巨大.期待著摩西重臨的以色列人如何可能接受耶穌就是那位拯救者呢?

基督徒,你們又在期待甚麼呢?你期待基督教領袖如摩西一樣大有能力嗎?如果像耶穌一樣,你能夠接受嗎?基督徒,你對自己又有甚麼期待呢?你期待自己如摩西一樣大有能力見證上帝的大能嗎?如果你的遭遇和狀況如耶穌一樣,你會否覺得自己很失敗,失見證,不配做基督徒呢?

5. 無能的基督

使徒保羅有另一種看法,基督看似無能地回應挑戰,看似軟弱地回應質疑,並不是失敗,也不是失見證,反而正正是作見證.他說:

2:6他本有上帝的形像,不以自己與上帝同等為強奪的；

2:7反倒虛己,取了奴僕的形像,成為人的樣式；

2:8既有人的樣子,就自己卑微,存心順服,以至於死,且死在十字架上.

很多人一生都在為了證明自己有用,有能力而掙扎,內心忐忑不安,夜不安枕.很多人為了證明自己比人強而隨意貶低他人,指責他人無用,無能力,製造很多傷害.有些教會為要證明自己做得好而想盡辦法增加人數和奉獻數字；有些學校為了要成為名校而把跟不上的學生請走；有些醫院為了減低死亡率而拒收臨危病人...無奈又痛苦的現實是,大家都渴求被肯定,被讚賞,被吹捧.更無奈及更痛苦的是,那些肯定我們,讚賞我們,吹捧我們的,是可以隨時「搬龍門」的.正如某位明星的歌迷所說:「我們捧得你上去,也可以把你拉下來!」對著那個隨時移位的龍門,你只有疲於奔命.

耶穌被世人視為無用,無能力,無成就,無地位,就算如此,他也不去用世人的標準證明自己有用,有能力,有成就,有地位,因為他清楚知道最終只有誰才有資格判斷怎樣才是有用和有成就.

2:9所以,上帝將他升為至高,又賜給他那超乎萬名之上的名,

2:10叫一切在天上的,地上的,和地底下的,因耶穌的名無不屈膝,

2:11無不口稱「耶穌基督為主」,使榮耀歸與父上帝.

6. 總結:當以基督耶穌的心為心

因此,使徒保羅勸勉信徒:

2:5你們當以基督耶穌的心為心.

很多人以為,教會要夠強大,才是好見證.的確,人們是會被強大的教會吸引的.很多人以為,成功才是好見證,的確,人們是會被有成功人士演講的聚會吸引的.很多人以為,如果國家元首是基督徒,或者基督徒晉身政府高層,就是好見證,的確,基督徒這身分可能會得到多一點正面形象及支持的.不過,這可能只是一剎那的光輝,你可能還記得,美國某教會短時間內變身世界知名的大教會,主日崇拜向全球百多個國家轉播,但幾年後就申請破產了.

「有多年基督教出版經驗的資深文字工作者吳思源形容香港曾是「整個世界華文基督教出版的中心」...香港成為當時華文出版的唯一自由窗口.隨著第一代信徒成長,吳思源指當時青年信徒求知若渴,對屬靈書籍需求大,基督教出版業七,八十年代出現「蓬勃期」,當時在外國唐人街或華人教會買一本書,「十居其八」都是香港出版...自從九十年代末開始,閱讀人口漸漸減少.近三年的疫情,帶來的影響更是非常巨大,書店面對很大的困難,步入「衰退期」.」(時代論壇)

曾幾何時,這個城市何止是「整個世界華文基督教出版的中心」,更是「ＸＸ中心」,「ＸＸ中心」.即是說,這個城市曾經「威」過的.基督教出版事業的變化,某程度上也反映出這個城市眾多方面的變化,包括教會.我們失敗了嗎?我們失見證了嗎?我們要祈求上主幫助我們重奪「ＸＸ中心」,「ＸＸ中心」的地位才算是好見證嗎?或許有人正為此而祈求和努力,各人有不同的領受.然而,如果事情並不如理想,時勢並不配合,「不強奪」,「虛己」,「卑微」,「順服至死」,對耶穌來說,也是好見證,對我們跟隨耶穌基督的人,可值得參考.



\newpage

\section{對話中的十誡}
\label{sec:666}
\textbf{李均熊牧師 - 對話中的十誡}
\newline
\newline
\hyperref[sec:665]{< < < PREV SERMON < < <}
~
\hyperref[sec:toc]{[返目錄]}
~
\hyperref[sec:667]{> > > NEXT SERMON > > >}


講題:對話中的十誡 10 Commandments in Dialogue

講員:李均熊牧師

場合:香港中文大學崇基學院禮拜堂主日崇拜

日期:2023年10月8日

起:十條誡命有長短

今日舊約經課選讀仍然是出埃及記,而我們進入西奈山,聆聽上主賜下的,傳統稱為「十誡」的經文.這個被稱為「十誡」的段落,應是其中一段舊約最為人熟識的經文之一.但有時越熟識的事物,我們對它的了解並不一定是最全面；更可能因為自以為熟識,就沒有再細心參詳,被先入為主的感覺所主導.

首先讓我們略為對經文進行一些觀察.

八條以不可為始(不x13),二條以當為始

六條沒有其他枝節(短)

一條帶著應許:當孝敬父母

一條帶著原因:守安息日

二條帶著警告?必懲罰/必不以他為無罪

十條誡命如何分?例: 那句算為第一誡?而且不可貪戀出現兩次

承: 對話中的「十誡」

只是從表面的觀察,就提醒我們兩件事,這段經文並不是簡單的十句說說,而且十誡並不是「鐵板一塊」,不同群體可以有不同的閱讀.這個不同閱讀的情況並不只出現在猶太教又或公教與新教的分別,其實在聖經內我們已經能看到這種不同詮釋不同理解的可能.因為今日經課選讀的(所謂)「十誡」,不單只在出埃及記出現,也在申命記出現,其實也以不同形式出現在聖經的不同地方.先不說出埃及記和申命記中的分別,當福音書記載一個少年的官進見耶穌,耶穌就引用了十誡(太19:18-19; 可10:19; 路18:20).當中耶穌不但只引用後半部份的誡命,引用的次序在路加的經文和馬太,馬可的經文也是不同的.路加中的耶穌是先引不可姦淫,才引不可殺人,和我們熟識的不可殺人不可姦淫的次序有所不同.當中原因可能是路加作者引用的希臘文聖經的順序是不可姦淫先行.而且耶穌引用後半也可能和當時人背誦十誡是由短句開始有關.

說回「十誡」在出埃及記和申命記的兩個不同版本,其不同處不是簡單的一兩處……

學者當然提出很多不同理論去解釋十誡在出埃及記和申命記中的不同,但我們只要留心經文的上文下理,就是在經文的世界內也是可以理解的.經文記載的十誡並不是像校規(宿規)只是列出一條一條抽離的誡命條例,在出埃及記乃是出現在西乃山上主在以色列人中顯現的對話情節中,而在申命記則是出現在摩西經過四十年之後向將要入迦南而又沒有經歷過西乃神顯的新生代重申上主說話的場境中.

轉:十「誡」?

凡此種種,其實都告訢我們十誡並非一成不變的清規誡律(比對:校規),而是有結構和文學性的說話,會因應對話對像的處境而有不同「說法」.

「誡律」對你來說是甚麼?我們可能認為我們不會因為嚴守誡律而成為律法主義,但誡律總會令人聯想起不問情由的遵從,甚至是高壓權力的強制措施.在今日自由社會的處境下,誡律更會被視為管束.誡律就是防止你做錯事,卻不是告訴你怎樣做才為對.當然誡律是籬笆,邊防,謹守遵行可免於我們行差踏錯,然而,今日我們看誡律是否會想起負面的事情多於正面的事情呢?

但出埃及記二十章的經文,細心看其實是上主的吩咐,教誨,叮嚀,是在對話中上主表達期望和要求有個性(personal)的說話,乃上主位格(person)的彰顯.十誡在英文有另一個更好的翻譯,就是decalogue (Greek: deka logous,出34:28 Hebrew: ʿĂseret haDībbərōt ),中譯其實可以是十言(10 words),而非10  Commandments,以至我們對神的說話產生不必要的誤解.

合:喜樂中對話

十誡是處於神人盟約的對話中(出19-24 ),從記載的背景來看,也是因為上主帶領祂子民由為奴之家出來後將要進入新生活的吩咐.上主的「十言」,不是沒有斟酌空間,要求不問情由地接受,甚至盲目使用的條文,而是在對話中生成,須隨處境的轉換而再思如何實踐.作為神的話(在出及記,真是文義上(literally)神的說話),其價值對今天有十分重要的指導教誨作用.[猶太人稱之為「妥拉」 Torah,今天(2023年10月7日黃昏至8日黃昏)其實就是猶太教住棚節最後一天,稱為Simchat Torah「在妥拉中喜樂」的意思.] 上主的說話並非無情的清規誡律,而是上主有情的叮嚀吩咐,我們會否帶著高興快樂去接受?上主十言也須要我們用心思聆聽和回應,在對話中認識上主心意,用心思考如何踐行祂的要求.



\newpage

\section{天地間的一場豪賭}
\label{sec:667}
\textbf{鄧瑞強博士 - 天地間的一場豪賭}
\newline
\newline
\hyperref[sec:666]{< < < PREV SERMON < < <}
~
\hyperref[sec:toc]{[返目錄]}
~
\hyperref[sec:668]{> > > NEXT SERMON > > >}


約伯記 1:1-22
\newline
\begin{longtable}{cl}
\hline
\hline
章節 & 經文 (和合本修訂版)\\
\hline
1:1 & \begin{tabularx}{0.7\textwidth}{X} 烏斯地有一個人名叫約伯。這人完全、正直、敬畏神、遠離惡事。 \end{tabularx} \\ \\ \relax
1:2 & \begin{tabularx}{0.7\textwidth}{X} 他生了七個兒子，三個女兒。 \end{tabularx} \\ \\ \relax
1:3 & \begin{tabularx}{0.7\textwidth}{X} 他的家產有七千隻羊，三千匹駱駝，五百對牛，五百匹母驢，並有許多僕婢。這人在東方人中為至大。 \end{tabularx} \\ \\ \relax
1:4 & \begin{tabularx}{0.7\textwidth}{X} 他的兒子按著日子各在自己家裡擺設宴席，派人去請他們的三個姊妹來，與他們一同吃喝。 \end{tabularx} \\ \\ \relax
1:5 & \begin{tabularx}{0.7\textwidth}{X} 宴席的日子過了，約伯派人去叫他們自潔。他清早起來，按著他們眾人的數目獻燔祭，因為他說：「恐怕我的兒子犯了罪，心中背棄神。」約伯常常這樣行。 \end{tabularx} \\ \\ \relax
1:6 & \begin{tabularx}{0.7\textwidth}{X} 有一天，神的眾使者來侍立在耶和華面前，撒但也來在其中。 \end{tabularx} \\ \\ \relax
1:7 & \begin{tabularx}{0.7\textwidth}{X} 耶和華對撒但說：「你從哪裡來？」撒但回答耶和華說：「我從地上走來走去，在那裡往返。」 \end{tabularx} \\ \\ \relax
1:8 & \begin{tabularx}{0.7\textwidth}{X} 耶和華對撒但說：「你曾用心察看我的僕人約伯沒有？地上再沒有人像他那樣完全、正直、敬畏神、遠離惡事。」 \end{tabularx} \\ \\ \relax
1:9 & \begin{tabularx}{0.7\textwidth}{X} 撒但回答耶和華說：「約伯敬畏神，豈是無故呢？ \end{tabularx} \\ \\ \relax
1:10 & \begin{tabularx}{0.7\textwidth}{X} 你豈不是四面圈上籬笆圍護他和他的家，以及他一切所有的嗎？他手所做的都蒙你賜福，他的家產也在地上增多。 \end{tabularx} \\ \\ \relax
1:11 & \begin{tabularx}{0.7\textwidth}{X} 但你若伸手毀他一切所有的，他必當面背棄你。」 \end{tabularx} \\ \\ \relax
1:12 & \begin{tabularx}{0.7\textwidth}{X} 耶和華對撒但說：「看哪，凡他所有的都在你手中；只是不可伸手加害於他。」於是撒但從耶和華面前退出去。 \end{tabularx} \\ \\ \relax
1:13 & \begin{tabularx}{0.7\textwidth}{X} 有一天，約伯的兒女正在他們長兄的家裡吃飯喝酒， \end{tabularx} \\ \\ \relax
1:14 & \begin{tabularx}{0.7\textwidth}{X} 有報信的來見約伯，說：「牛正耕地，母驢在旁邊吃草， \end{tabularx} \\ \\ \relax
1:15 & \begin{tabularx}{0.7\textwidth}{X} 示巴人忽然闖來，把牲畜擄去，並用刀殺了僕人；惟有我一人逃脫，來報信給你。」 \end{tabularx} \\ \\ \relax
1:16 & \begin{tabularx}{0.7\textwidth}{X} 他還說話的時候，又有人來說：「神從天上降下火來，把羊群和僕人都吞滅了；惟有我一人逃脫，來報信給你。」 \end{tabularx} \\ \\ \relax
1:17 & \begin{tabularx}{0.7\textwidth}{X} 他還說話的時候，又有人來說：「迦勒底人分成三隊忽然闖來，把駱駝擄去，並用刀殺了僕人；惟有我一人逃脫，來報信給你。」 \end{tabularx} \\ \\ \relax
1:18 & \begin{tabularx}{0.7\textwidth}{X} 他還說話的時候，又有人來說：「你的兒女正在他們長兄的家裡吃飯喝酒， \end{tabularx} \\ \\ \relax
1:19 & \begin{tabularx}{0.7\textwidth}{X} 看哪，有狂風從曠野颳來，襲擊房屋的四角，房屋倒塌在年輕人身上，他們就都死了；惟有我一人逃脫，來報信給你。」 \end{tabularx} \\ \\ \relax
1:20 & \begin{tabularx}{0.7\textwidth}{X} 約伯就起來，撕裂外袍，剃了頭，俯伏在地敬拜， \end{tabularx} \\ \\ \relax
1:21 & \begin{tabularx}{0.7\textwidth}{X} 說：「我赤身出於母胎，也必赤身歸回；賞賜的是耶和華，收取的也是耶和華。耶和華的名是應當稱頌的。」 \end{tabularx} \\ \\ \relax
1:22 & \begin{tabularx}{0.7\textwidth}{X} 在這一切的事上，約伯並沒有犯罪，也不以神為狂妄。 \end{tabularx} \\ \\
[1ex]
\hline
\hline
\end{longtable}
講題:天地間的一場豪賭 Who will Win?

講員:鄧瑞強博士

經文:約伯記1:1-22

場合:香港中文大學崇基學院禮拜堂主日崇拜

日期:2023年10月15日

各位弟兄姊妹,早安.願神的恩典常與你們同在.

我已經講完五經及歷史書,今日,我們來到基督教聖經「詩歌智慧書」的部分.「詩歌智慧書」以《約伯記》開首,這是一卷典型的智慧文學.未講約伯記前,讓我分享一下智慧文學的特點.智慧文學不同我們之前看的五經及歷史書,五經及歷史書重視神與祂的子民的立約,重視神的應許,重視神的救贖.在智慧文學裡,以上這些典型的信仰元素都沒有了,而訴之人類的智慧,叫人用心探討這個生活世界的客觀現象,叫人明白自己在這生活世界裡當盡的本分,當然,這是叫人謙虛下來,明白自己的有限,好面對真正無限的上帝.

舉個例子,讓大家明白現在要說的「智慧文學」與我們習慣的「救贖文學」的分別.若你在埃及做奴隸,天天苦不堪言,天天祈求神解救,在「救贖文學」裡,神會聽禱告,神會派摩西來解救,會有神蹟,紅海會分開,你會得到自由.但在「智慧文學」裡,你當然可以祈禱,但不要期望神蹟會發生.「智慧文學」會教你,在埃及做奴隸的生存智慧,如:不要埋怨,因為埋怨後你仍是要做奴隸的.不要對管工口出惡言,禍從口出,之後你會做得更辛苦.要勤力一點,若因此得讚賞,你可能會輕省一點,但世事也不一定如此的,可能你越勤力,做得越多.你要做的事,要用心做.但之後的結果,你是預計不到的.你要明白,人生有種種限制,知道你自己不是神,要懂得敬畏真正的神.

再舉個例子,若一個人患絕症,在「救贖文學」裡,他通常首先是認罪,因為按照「約的神學」,不聽從神的吩咐,會帶來禍患.認罪後,祈求神蹟醫治,因為神救贖受苦的人.但「智慧文學」會告訴你,患絕症的人要快一些找醫生.若他找一個壞醫生去醫,那是他愚昧.若他找一個好醫生去醫,那是他有智慧.若他無錢找一個好醫生,則他要反省:為何平時他不好好賺錢?為何他不買一份好的保險?他若說:我已經很勤力工作,奈何很多地方仍不及人,故賺不到多少錢?「智慧文學」會說:世界就是這樣的,你只能接受生命的限制.在這限制下,請不要胡作非為,因這終究是無智慧的,要好好敬畏神.基督徒很習慣「救贖文學」,較不明白「智慧文學」.主耶穌說:「日頭照好人,也照歹人」,這是典型的「智慧文學」.不會因為一個人壞,太陽便不照他的.若這壞人後來做了很多壞事,我們也不能怪責神為何做一個會照壞人的太陽.神造了太陽,它照好人壞人,太陽就是這樣,世界就是這樣,絕不會因你嫉惡如仇的想法改變.智慧文學,要我們放下神必然會介入而顯神蹟的那種「救贖思考」,而要我們真切地,智慧地面對這蒼茫世界,腳踏實地的努力前行.在這裡,能依賴的,是人生的常理；生活要求的,是我們當盡的責任；在這一切背後,是敬畏神的心.

明白了「智慧文學」的基本特徵後,我們進入約伯記的經文裡了.

今日的講道經文:約伯記1:1-22

1:1　烏斯地有一個人名叫約伯；那人完全正直,敬畏上帝,遠離惡事.

1:2　他生了七個兒子,三個女兒.〔按:數字加起來是「十」,有圓滿的意思〕

1:3　他的家產有七千羊,三千駱駝,五百對牛,五百母驢,並有許多僕婢.這人在東方人中就為至大.

1:4　他的兒子按著日子各在自己家裡設擺筵宴,就打發人去,請了他們的三個姊妹來,與他們一同吃喝.

1:5　筵宴的日子過了,約伯打發人去叫他們自潔.他清早起來,按著他們眾人的數目獻燔祭；因為他說:「恐怕我兒子犯了罪,心中棄掉上帝.」約伯常常這樣行.

故事的主角是約伯,他完全正直,敬畏上帝,遠離惡事.這是他的人格素質,在宗教敬虔上,在倫理操守上,他都無可指摘.接著,便講他的外在際遇,他有圓滿的兒女群,有豐足的家業,在東方人中他為至幸福的人.他的兒女關係和諧,約伯作為一家之長,守護著家族的敬虔.他德福兼備,無人能及.這樣的生活,能永無止盡地延續下去,你說多好!

1:6　有一天,上帝的眾子來侍立在耶和華面前,撒旦也來在其中.

1:7　耶和華問撒旦說:「你從哪裡來?」撒旦回答說:「我從地上走來走去,往返而來.」

1:8　耶和華問撒旦說:「你曾用心察看我的僕人約伯沒有?地上再沒有人像他完全正直,敬畏上帝,遠離惡事.」

1:9　撒旦回答耶和華說:「約伯敬畏上帝,豈是無故呢?

1:10　你豈不是四面圈上籬笆圍護他和他的家,並他一切所有的嗎?他手所做的都蒙你賜福；他的家產也在地上增多.

1:11　你且伸手毀他一切所有的；他必當面棄掉你.」

1:12　耶和華對撒旦說:「凡他所有的都在你手中；只是不可伸手加害於他.」於是撒旦從耶和華面前退去.

在天上,有上帝主持的會議,這是古代近東文化常有的畫面.經文說:「撒旦也來在其中」.關於撒旦,在這裡,只是扮演一個小人物的角色.有學者說,「撒旦」這名字的字根,意思是「走來走去」.他走來走去幹什麼?他收集人類的情報,向上帝匯報.在這裡,撒旦不是後來的大魔王,這觀念還未發展出來,他只是小小的情報員而已.這個撒旦,只是上帝之下的其中一個小被造物.他要向上帝負責.上帝發問,他就要回答.他想做的事,若沒有上帝允許,也不能亂來.其實,無論是那時的小情報員撒旦,抑或是今日的大魔王撒旦,他的角色,基本上沒有什麼變動.他不是主,也不是為所欲為的.

耶和華問撒旦說:「你曾用心察看我的僕人約伯沒有?地上再沒有人像他完全正直,敬畏上帝,遠離惡事.」在這裡,上帝重複第1節對約伯的讚許,約伯真是倫理上,敬虔上無可指摘的人.上帝問情報員撒旦,你的情報如何?

1:9　撒旦回答耶和華說:「約伯敬畏上帝,豈是無故呢?」

不知是否出於撒旦的妒忌,抑或出於撒旦對人性觀察得夠深入,他問了一個很重要的宗教問題.人敬畏上帝,豈是無故呢?人相信神,不可能背後沒有原因!約伯的敬虔,與他世上的福樂,一定有關連!人的敬虔,到頭來是否只是一種利益的盤算?若拿走約伯的福樂,他還會相信神嗎?拿走多少,他才不再相信呢?有時,我會問學生(當然,這純粹只是用於刺激他們思考的),若沒有天堂,你們還會不會信耶穌?信上帝這回事,到底有多少是因為得到福樂?我們著眼上帝,抑或只是著眼福樂?我們終極追求的,是上帝,抑或福報?撒旦真是一個利害的宗教哲學家.

1:11　撒旦對上帝說:「你且伸手毀他一切所有的；他必當面棄掉你.」

在這裡,撒旦「大」上帝:「你這個世上最完美的敬虔者,若失去一切世上的福樂,必當著你面咒罵你.」

上帝如何接招呢?上帝能如何接招呢?

上帝剛剛才讚約伯是世上最完美的敬虔者,難道上帝信不過他嗎?若上帝拒絕撒旦的測試,即是連最敬虔的人,上帝都信不過,那上帝還有希望嗎?你想想:若你認為有一個人,是你世上最好的朋友,十世難求.有人挑戰你說:「你敢不敢打賭,用十萬元就能叫他出賣你?」你會怎說呢?十萬元,你當然敢賭這一局.若這人又說:「你敢不敢打賭,用一億元就能叫他出賣你?」一億元,你還敢賭嗎?十億元呢?你還能信任你的最好朋友嗎?我們可能收手了.上帝又如何呢?祂能信任祂的朋友約伯嗎?撒旦「曬冷」了,上帝呢?好驚訝,上帝都「曬冷」了.上帝接受了天地間的這場豪賭.祂選擇相信人.若世上最敬虔的人都信不過,這世間還有希望嗎?我們常常說,信仰是人選擇相信神,但事實是,信仰是神選擇相信人.若神根本不相信人,這個世界轉眼間便消失了.

以上的情節,是「日光之上」的事,是「天界」的事,對人類來說,這是隱秘的.苦難問題,終究是理性不可解的.對於活在「日光之下」的人類來說,我們知道的,面對的,就是當下的現實.這殘酷的現實人生,就其隱秘的意義而言,是人向天地萬物表明:是上帝對,抑或是撒旦對?人活在世上,到底是否只是為了物質的富貴?抑或,就算沒有了物質的富貴,我們仍然能在上帝那裡找到存在的意義?令我常常心靈震動的是:上帝竟然這樣大膽,信任我們!

「日光之上」的事,對人是隱秘的.約伯能見到的,是「日光之下」的事情.

1:13　有一天,約伯的兒女正在他們長兄的家裡吃飯喝酒,

1:14　有報信的來見約伯,說:「牛正耕地,驢在旁邊吃草,

1:15　示巴人忽然闖來,把牲畜擄去,並用刀殺了僕人；惟有我一人逃脫,來報信給你.」

1:16　他還說話的時候,又有人來說:「上帝從天上降下火來,將群羊和僕人都燒滅了；惟有我一人逃脫,來報信給你.」

1:17　他還說話的時候,又有人來說:「迦勒底人分作三隊忽然闖來,把駱駝擄去,並用刀殺了僕人；惟有我一人逃脫,來報信給你.」

1:18　他還說話的時候,又有人來說:「你的兒女正在他們長兄的家裡吃飯喝酒,

1:19　不料,有狂風從曠野颳來,擊打房屋的四角,房屋倒塌在少年人身上,他們就都死了；惟有我一人逃脫,來報信給你.」

這就是真實的人間.

一浪緊接一浪的災難.先是人禍(示巴人來搶掠),跟著是天火(應是閃電),接著是人禍(迦勒底人來搶掠),再接著是暴風.人間的悲劇就發生在我們身邊.近來,我們見到暴風的威力.在電視上,我們也見到人禍帶來的悲慘.在以色列,一群人在音樂會裡唱歌跳舞,轉瞬間,子彈橫飛,下一刻,家人已永別,家園已被毀.這就是約伯的經驗,苦難來時,根本來不及思考,痛苦得無可言喻.

賭局來到這裡,撒旦這樣出牌,看來,他已穩操勝券.在如此苦難的人間,人還能相信神嗎?

1:20　約伯便起來,撕裂外袍,剃了頭,伏在地上下拜,

1:21　說:「我赤身出於母胎,也必赤身歸回；賞賜的是耶和華,收取的也是耶和華.耶和華的名是應當稱頌的.」

1:22　在這一切的事上約伯並不犯罪,也不以上帝為愚妄(或譯:也沒有埋怨上帝).

難以置信,真難以置信,約伯接受一切現實,現實就是這樣的.我空空而來,也必空空而去.生命本來就是這樣的.我擁有的東西,是神賜予的,我無權也無法永遠擁有.我失去了這些東西,這是神收回,我無權也無法拒絕.人生在世,本是如此.這就是「智慧文學」的特質.世事本來如此,不按我們主觀意志改變.你接受也好,你不接受也好,都要學曉放手.在這之上,約伯仍不犯罪,仍不埋怨上帝,仍稱頌上帝,仍堅持世上的大苦不阻礙人對神的敬畏.在深深的苦難裡,約伯堅持,人生終歸是有意義的,美善是存在的,上帝是可信賴的.

看著這一幕,我不知上帝會如何想?

或許,祂也驚訝,一個人竟然能忠誠到這個程度.其實,後來的耶穌,在苦難的深淵,也表達出這種對神的完全信賴.我們,作為相信基督的人,作為基督的跟隨者,是否也能在苦難人間,同樣忠誠地信靠神?

願上主得著榮耀,頌讚；願眾人在苦難人間,忠誠地敬畏神.



\newpage

\section{博雅教育與中國基督教大學的故事}
\label{sec:668}
\textbf{吳梓明教授 - 博雅教育與中國基督教大學的故事}
\newline
\newline
\hyperref[sec:667]{< < < PREV SERMON < < <}
~
\hyperref[sec:toc]{[返目錄]}
~
\hyperref[sec:669]{> > > NEXT SERMON > > >}


講題:博雅教育與中國基督教大學的故事 Liberal Arts and the Story of China’s Christian Universities

講員:吳梓明教授

場合:香港中文大學崇基學院禮拜堂主日崇拜

日期:2023年10月22日

多謝高牧師的邀請,讓我再有機會與崇基禮拜堂眾兄弟姊妹一同聚集,一同敬拜上帝. 祝願 上主的恩惠及慈愛常與我們同在!

五年前崇基校史考察活動開始由一班崇基校友與校牧室,及學院的通識教育中心一起聯合主辦,後來更有全國基督教大學同學會(CCUAA) 加入一齊協辦.五年來我們每年出版一本校史考察與學生的專題研習報告,由於有通識教育中心的參與,我們的學生研習亦增加了不少有關基督教大學與博雅教育/或通識教育課題的研習,逐漸地發覺原來博雅教育的課題不單只是我們老師,學生們,甚至是校友們 (包括浸會,嶺南大學等其他院校的校友們) 也對這課題甚感興趣.今年在全國基督教大學同學會慶祝70周年的晚宴上也是以「基督教大學與博雅教育」為主題,還有,今年5月我和許美基督教授一同編輯的一本書剛好出版,也是與博雅教育有關的,題目就是: “Liberal Arts and the Legacy of China’s Christian Universities”. [1] 今早我想與大家介紹這本書,特別是我們所認識的博雅教育和它與中國基督教大學有關的故事.

這本書是我與許美基督教授(Prof. Ruth Hayhoe)及她的一位博士生合編的書.在1972年,當我剛從崇基畢業並開始在聖傑靈女子中學教書的時候,許美德女士正在協恩中學教書,我們是在香港教會學校宗教教育協會(Hong Kong Association for Religious Education in Christian Schools)上合作認識的.1978年,許美德往英國倫敦大學教育學院讀博士,後來我也去倫大讀博士.她是鑽研比較教育,專攻中國高等教育的發展,我是從事宗教教育與中國教會大學史研究的.這本書的出版一方面是慶祝我們相識50年的友誼,另一方面也是見證著我們過去30多年來與國內學者們在中國高等教育史上共同研究,默默耕耘的成果.

我在1989年參加了第一屆中國教會大學史國際學術研討會上,認識華中師範大學校章開沅教授,開始了我們對四九年前中國教會大學史的研究.1991年第二次的會議是在南京師範大學舉行,許美基督教授剛巧在國內工作,擔任加拿大駐華大使館教育,科學與文化參贊(1989-1991),她也來參加了南京的會議.30年後,到了2019年,我們又召開了一個慶祝研討會舉辦30周年的大會,許美基督教授已經返回加拿大,但她派了她的博士生 (即後來的牟樂平博士) 參加這個會議,後來我們三人決定一起編訂這本書Liberal Arts and the Legacy of China’s Christian Universities,總結學者們過去30年來對中國基督教大學史研究的成果.

這本書其中一個特色是我們搜集了13篇具代表性的文章,能夠概括了過去30多年來國內的學者,包括來自上海復旦大學的徐以驊教授,華中師範大學的章開沅教授,他的弟子馬敏教授,中國人民大學的劉賢副教授,淅江大學的張凱教授,山東大學的劉家峰教授,福建師範大學的朱峰教授,和吳巍巍教授等,在教會大學史研究的成果和心得.較為明顯的發現是他們能夠將教會大學與近代中國的高等教育研究連在一起,尤其是在大學教育的現代化 (國學研究,宗教研究方法,宗教比較等),國際化 (擴闊世界視野,結合中西文化教育),跨學科 (文理兼容,引進新學科如心理學,社會學,人類學等跨學科的研究方法),跨文化 (促進中西文化的交流與共融),和高雅品德與人文素質 (即全人教育) 的追求等,這些均是與教會大學所倡導的博雅教育有關的.

「博雅教育」

甚麼是「博雅教育」? 一般學者談及香港的博雅教育時,多以為「博雅教育」是源自嶺南大學,他們只說對一半.原因是1995年當陳坤耀教授出任嶺南學院校長時,他為快將升格為大學的嶺南學院未來發展目標定位為「博雅大學」[2]. 陳教授更進一步闡釋「博雅」一詞為: 「學識廣博,生活高雅」[3]. 但事實上,博雅教育是20世紀初期國內13所教會大學創立時,已將歐美(主要是美國的) Liberal Arts College 的教育模式帶來中國的,尤其是北京的燕京大學和廣州的嶺南大學的表現得尤為突出.中國基督教大學博雅教育的特色在於其「跨學科,跨文化,推行雙語教學,兼重文理與全人教育,關顧個人成長和品德發展等」[4].1950年開始,國內教會大學的師生遷移來香港,崇基學院是在1951年創立,目的是為要保存國內13所基督教大學的博雅教育模式,讓這些教育傳統能夠繼續傳承下去.除了剛剛所說的「跨學科,跨文化,推行雙語教學等」外,在課程上較獨特的是設立哲學和宗教系,社工系(社會服務),音樂系,和體育系等,這些就是燕京大學和嶺南大學等基督教大學所保留獨特的博雅教育傳統.崇基學院的創校校長是李應林博士,他本來是廣州嶺南大學的校長,因此,早期崇基學院的教職員大部分都是來自廣州嶺南大學的,包括李應林,何明華,謝昭杰,曾昭森,黃壽林,謝扶雅,麥健增,容啟東等.[5] 他們不僅是幫助開拓崇基的校務及行政事宜,還有協助學院設立那些獨具特色的哲學宗教系 (編寫及教授「人生哲學」課程,後來發展成為「綜合基本課程」和「通識教育課程」),社工系 (包括「宗教教育與社會服務」的課程,是香港大專院校最早提供,也是獨具一格的社會服務課程) [6],和體育系 (我還記得我入讀崇基那年,在校園內有一位體育科老師是李小洛先生,他就是李應林校長的兒子),音樂系方面則是由燕京大學的音樂系系主任及資深教授范天祥教授協助創立的.

崇基創校之初,何明華會督倣效國內基督教大學的經驗 (包括廣州嶺南大學的經驗,因為何會督曾任教於廣州嶺南大學協和神學院,並兼任神學院校董),設計「人生哲學」課程為崇基學生的必修課程之一.[7] 何會督的「人生哲學」,顧名思義,就是直指人生的哲學,它是從哲學(理性)的角度探討人生,思考人生的問題.何會督認為「人生哲學」就是一種「對話」(dialogue) 和「會談」(conversation),他對學生說:『崇基…的「人生哲學」…是要(讓)每一個學生和歷代中西哲人的會談,也與現代思潮和不同的人生觀對話』.[8] 這種「對話」和「會談」的教學模式成為了今日崇基學院(更是香港中文大學的) 通識/博雅教育課程中一種很獨特的教學模式.原來今日香港中文大學通識教育所突顯的「與人文對話」和「會談」的教學模式,早在1949年前的中國基督教大學及在1950年代開始的崇基學院已經出現了.

我最敬愛的一位老師沈宣仁教授(1931-2004) 在1962年加入崇基學院時[9],即致力協助學院在既有的「人生哲學」課程基礎上,製訂「綜合基本課程」(Integrative Basic Studies,簡稱為「IBS」),[10] 崇基學院在2016年成立「宣仁通識教育資源中心」(中心設在崇基圖書館旁邊),並推崇沈教授為「崇基學院通識教育之父」.若果沈教授仍是活著的話,我相信他會說: 假如我 (沈教授) 是崇基學院通識教育之父,我們的祖父(通識教育的創始人)應該是何明華會督.沈教授亦曾坦言說他是深受何明華會督的感染和影響的,所以我們可以說:「沈教授是崇基學院通識教育之父,何明華是我們學院通識教育之祖父」[11]

我和許美基督教授在學術研究上有一個共通點,是喜歡從人物敘事(narratives)入手尋找歷史故事.許教授最初是翻譯馬相伯的書,知道馬相伯 (Ma Xiangbo, 1840-1939)是在中國創辦過三間大學,分別是: 震旦大學(Aurora/ Zhendan Academy, 1903),復旦大學(Fudan University, 1905),和天主教輔仁大學(Fu Jen Catholic University, 1925). 後來我們發現原來馬相伯是最早將希臘文,拉丁文和希臘經典文學 (跨經典文學研究) 放進中國的大學教育課程中,也是最早將「博雅教育」的概念帶來中國的.

另一位重要的人物是燕京大學的創校校長- 使徒雷登 (John Leighton Stuart, 他有一位兄長使徒華林, Warren H. Stuart也曾是另一所教會大學「(杭州)之江大學」的校長, 1916-1922).使徒雷登是在杭州出生(1876年),長大,是一位中國通,他對中國最大的貢獻是打造了一間燕京大學出來.燕京大學其實亦是一間典型的博雅教育的大學.使徒雷登在1928年創立了「哈佛燕京學社」,不僅是推動了燕京大學的國學研究,也為中西文化的交流建構了一個學術平台,讓中西方的學術研究人員能夠有雙向的交流,亦讓許多燕京大學的學生可以往哈佛大學深造,同時亦在燕大培訓了不少西方著名的漢學家.換言之,燕京大學在不斷進行「中國化」的同時,也在致力提升中國現代教育的「國際化」.

使徒雷登在校園中更是營造了一個最適合推廣博雅教育的校園環境.他親自帶領師生在校園內同時使用中文及英語,靈活地運用雙語溝通.這種「雙語並用」的校園生活模式,不僅是有助提升了燕京大學的國際形象,也促進了中西文化的交流和共融.

許美基督教授曾撰寫過一本書: Portraits of Influential Chinese Educators [12],描述謝希德(1921-2000) 教授的故事.謝教授是一位固體物理學家,中國科學院院士,1988年獲第三世界科學院院士,她曾在上海復旦大學擔任副校長(1978—1983)和校長(1983—1988)之職達10年.後來我也補充發現了原來謝希基督教授的父親是謝玉銘(1893-1986) 教授,亦是物理學家,早年在北京通州協和大學 (即燕京大學的前身)畢業,赴美國深造後便回燕京大學任教,擔任燕京大學物理系系主任 (1929-1932).所以謝希基督教授從小便在燕京大校園內生活,長大,認識使徒雷登,也深受燕大博雅教育校園的薰陶.加上她多年在美國麻省史密斯學院(Smith College)校園生活的影響,當謝教授擔任上海復旦大學校長之後,她也自然地將博雅教育的特色- 包括自由開放的學術風氣,雙語的校園生活,兼容中西文化及具國際視野等元素引入上海復旦大學.[13]

自從1950年代開始,國內的「博雅教育」隨著中國基督教大學的消逝而被忘卻了,取而代之的是蘇聯的「專才- 科技訓練」的教育模式.至1980年代,隨著國內改革開放政策的推行,高等教育政策又改變了,為補專才教育的不足,大陸高校重新引進「通識教育」,創新「人文素質教育」等課程,甚至有些大學增設「博雅教育學院」,全面推行現代化的「博雅教育」.[14]今日在國內,除了北京大學,上海復旦大學外,還有許多近年新建,標榜博雅教育模式的大學/學院如: 「寧波諾丁漢大學」(University of Nottingham, Ningbo China, 2004),「北京師範大學- 香港浸會大學聯合國際學院」(Beijing Normal University – Hong Kong Baptist University  United International College, 2005),「廣州中山大學博雅學院」(Boya (Liberal Arts) College at Sun Yatsen University in Guangzhou, 2009),「上海紐約大學」(New York University Shanghai, 2012),「崑山杜克大學」(Duke Kunshan University of Wuhan China, 2014).我們也可以說: 昔日中國教會大學的博雅教育傳統也再次以「蛻變」,復活過來,是以不一樣的方式被傳承下來了.

結論 – 三個比喻或象徵

最後,我會用三個比喻或是象徵來總結我對中國基督教大學和博雅教育的一些看法,是我累積了30多年來在中國基督教大學歷史研究中所體會到的,與各位同道們一起分享吧!.

1.「一粒麥子的比喻」: 2014年,在我發表一篇文章 ──〈中國教會大學的興起與發展〉[15]── 的總結時,我引用了耶穌的一句話說: 『一粒麥子,若不落在地裡死了,仍舊是一粒,若是死了,就結出許多子粒來』(約翰福音十二24).[16] 這句話或許也是一個比喻,它說出了一個重要的真理,就是: 一粒麥子,若是要發揮其潛在的生命力,結出果實來,就必須要「被埋在地裡」,「死了」,才能「結出許多子粒來」. 事實上,基督教教育在中國的歷史中,曾經歷過「調適」和「處境化」的演變過程中,我們或許也可以看見到,它曾是「被埋在地裡」,「死了」,現在卻已經「結出許多子粒來」了.

2.「落花生的比喻」: 20世紀初,中國教會一位史學家王治心教授在討論本色化運動時,曾用過一個象徵比喻 -「落花生」.[17]「花生」在中國本是稱為「洋花生」,或是稱為「東洋果」,但現在已沒有人這樣稱呼它了,甚至也不知道它原來是從外國帶回來中國栽種的.「本色化」(Indigenization)一詞也有人翻譯成「本位化」,「本地化」或「本土化」,意思就是在特定的文化環境裡的具體化,成為「土生土長」的.同樣地,當基督教教育進入中國的土壤中,它也必須要找到落腳點與生根點,埋藏在中國的地土裡,吸取中國土壤的養分,才可以達成「調適」和「處境(中國)化」的過程,最終成為真真正正是「土生土長」的中國基督教教育來.

3.「毛毛蟲的蛻變」:兒時也常聽到一個有關毛毛蟲的故事,唱著:「點蟲蟲,蟲蟲飛」等.為什麼蟲蟲會飛呢?原來那條毛蟲經過「蛻變」之後,便會變成一隻蝴蝶飛走了. 當我們見到一隻美麗的蝴蝶在花園中翩翩飛舞時,我們卻忘掉了原來它是一條從不起眼的毛毛蟲蛻變出來的,原來它也是經歷了一個奇妙的「蛻變」過程,才變成今天一隻美麗的蝴蝶. 基督教高等教育在中國的歷史中,或許也像毛毛蟲一樣,從一條不起眼的毛毛蟲在不斷地蛻變中,其間也許有人會以為它已經從歷史中消逝了(像毛毛蟲死去了一樣),誰知它已經變成了一隻美麗的蝴蝶,在花園中翩翩飛舞著,它或許也已成為另一個「生命蛻變」故事的主角了.

這就是我們的新書 liberal arts and the Legacy of China’s Christian Universities 想與大家分享的故事.

盼望上帝祝福我們,幫助我們延續崇基先賢傳承給我們的教育使命及博雅教育的傳統,我們同心的禱告是 :「主啊! 求你堅立崇基創校先賢們,和中國基督教大學的先賢們手所作的工! 也幫助我們繼續承傳他們所留給我們的教育遠象和使命. 主啊! 願你也堅立我們手所作的工,我們手所作的工,願主堅立」(詩篇90:17).啊們!

參考文獻

〈人生哲學課–崇基通識的由來〉,刊載於《崇基校園通訊》第58期,卷1 (2022年1月20日),頁24.

王治心 .(1925).〈中國本色教會的討論〉,頁11-16,載於《青年進步》,第79冊.

何明華.(1955a).〈專載: 何明華: 霍約瑟堂奉獻典禮演講詞〉,《崇基校刊》,第6期,1955年1月20日,頁5-6.

何明華.(1955b).〈會談(英文題目為 “Conversation”)〉,何明華會督在崇基學院第一屆畢業崇拜上講詞,1955年6月23日.中譯本節錄,刊載於崇基學院校史檔案館:《哲懷深致立崇基──創校文物展(1951-1963)展覽小冊》,附錄六,頁44-45.

胡志偉編 (2017).《開創以傳承 – 嶺南大學香港復校》,香港: 中華書局.

梁智仁,魏智樂,彭穎研,李穎妍(2021).《香江薪傳 - 中國基督教大學在香港的承傳與蛻變》.香港: 香港中文大學崇基學院.

陳沛瑜,彭穎研,呂俊傑,劉昀善,楊芷欣(2022).《明華予崇基 會督與香江 – 香港中文大學崇基學院通識教育校史考察計劃暨學生專題研習報告》.香港: 香港中文大學崇基學院.

陳坤耀 (2014).〈香港人對博雅教育仍多誤解〉,詳見網上文章: \url{https://w.w.w.}《灼見名家》,2014年9月4日.

詹肇凱.(1952).〈我從安德生教授獲得了人生哲學的概念〉,刊於《崇基校刊》第2期,卷1(1952年5月10日),頁2-3.

謝昭杰.(1951).〈校務報告〉,刊於《崇基校刊》,第1期,卷1 (1951年12月),頁2-3.

楊福家.(2017).〈我與博雅教育的相遇〉,擷自: \url{https://www.sohu.com/a/134400160_507651} ,2017-04-16.

Hayhoe, Ruth. (2006). Portraits of Influential Chinese Educators. Hong Kong: Comparative Education Research Centre, University of Hong Kong and Springer.

Ng, Peter Tze Ming. (2014) “Rise and Development of Christian Higher Education in China”, in Carpenter, Joel et al. (eds.) Christian Higher Education: A Global Reconnaissance. Grand Rapids, Michigan: William B. Eerdmans Publishing Company.

Ng, Peter Tze Ming, Leping Mou and Ruth Hayhoe (eds.) Liberal Arts and the Legacy of China’s Christian Universities. Singapore: Springer Nature Pte Ltd., 2023.

[1]   參Peter Tze Ming Ng, Leping Mou and Ruth Hayhoe (eds.) Liberal Arts and the Legacy of China’s Christian Universities. Singapore: Springer Nature Pte Ltd., 2023.

[2]     參《開創以傳承 – 嶺南大學香港復校》,香港: 中華書局,2017,頁76.

[3]     參陳坤耀:〈香港人對博雅教育仍多誤解〉,詳見網上文章: \url{https://w.w.w.}《灼見名家》,2014年9月4日.

[4]     參Ng et al., 2023年書.

[5]     參呂羅中:〈嶺南與崇基的淵源與情緣: 他們是「崇基人」,也是「嶺南人」〉.梁智仁等編:《香江薪傳: 中國基督教大學在香港的承傳與蛻變》,香港中文大學崇基學院通識教育校史考察學生專題研習報告,2021,頁222-235.

[6]    自1956年,崇基已經設有「宗教教育及社會服務學系」,提供了「宗教教育及社會工作文憑課程」,教師包括芮陶菴博士(Dr. Andrew T. Roy),馬鴻述教授,黃壽林先生等.有關「宗教教育及社會工作文憑課程」之介紹,參《崇基校刊》,第9期.1956年11月刊,頁13.「宗教教育及社會服務學系」教職員名單可參《崇基校刊》,第15期.1958年5月1日,頁13.

[7]     參謝昭杰,1951,頁2-3及7; 詹肇凱,1952,頁2.

[8]   何會督解釋「對話」是接近希腊語的思維,「會談」是較接近拉丁語的思維.「對話」容易產生對立和爭辯的結果,「會談」較側重生命的交流與分享,所以何會督選擇了用「會談」作第一次畢業演講的題旨. 參何明華,1955b,頁44.

[9]   崇基學院於1959年設立「神學及宗教教育系」,1962年改名為「宗教及哲學系」,至1968年又改稱為「哲學及宗教學系」(簡稱「哲宗系」),系內再分哲學,宗教知識,和神學三組.(梁智仁等合編,2021,頁61)

[10] 原來沈宣仁教授畢業於菲律賓基督教大學時是主修英文,副修哲學,在美國深造時,卻剛巧是在與早年謝昭杰主任所讀同一間以博雅教育馳名的歐柏林學院修讀碩士,然後轉往芝加哥大學繼續攻讀神學博士的.美國的歐柏林學院和芝加哥大學同是重視博雅和通才教育的,因此,謝昭杰先生和沈宣仁教授亦自然地將西方的博雅/通識教育引進崇基來.在教學期間,沈教授曾擔任哲宗系系主任(1973-74,1975-76),崇基學院通識教育主任(1977-90),崇基學院文學院院長(1969-71,1976-77),香港中文大學文學院院長(1986-89)和崇基學院院長(1990-94)等要職.期間,沈教授亦為「綜合基本課程」編寫了一些教材,如《大學教育與大學生》(1973),《論文的註釋及參考書目》(1976),《讀和寫:大學之鑰初輯》(1977)等.

[11]   亦可參閱吳梓明:〈何明華贈予崇基的三份禮物〉,載於《明華予崇基 會督與香江》,香港中文大學崇基學院通識教育校史考察計劃暨學生專題研習報告,2022年,頁270-271.

[12]   參Hayhoe, Ruth. Portraits of Influential Chinese Educators. Hong Kong: Comparative Education Research Centre, University of Hong Kong and Springer, 2006.

[13]   1985-88年,謝校長亦兼任上海復旦大學美國研究中心所長,她率先在中國國內打破綜合大學只有文科,理科的前蘇聯模式,根據上海復旦大學的條件增設了技術科學,生命科學,管理科學等5個學院.1987年6月接受美國紐約州立大學奧爾巴尼分校授予名譽博士學位時,謝希德被《今日美國》報社記者採訪時被稱為是「中國的哈佛大學校長」.

[14]   楊福家教授也曾是上海復旦大學的校長,現任寧波諾丁漢大學校長.他曾這樣描述自己的經歷,說: 『我第一次瞭解真正的博雅教育是在2012年,那時我去訪問了一所美國的博雅學院,就是當年謝希德老校長曾去念書的史密斯學院(Smith College).美國的這些博雅學院,是以本科生教育為主的.提起美國的大學,我們往往會想到哈佛,耶魯,但其實在美國本科教育的排名上,一些小而精,專注博雅教育的文理學院(Liberal Arts College)甚至要排在哈佛耶魯之前.那次從美國參觀回來,我便體會到博雅教育的好處.於是回來後我就發起了博雅教育的會議,到這次是第三屆』.擷自: \url{https://www.sohu.com/a/134400160_507651} 《我與博雅教育的相遇》2017-04-16 20:49.他所提及美國麻省史密斯學院正是謝希基督教授曾在美國接受教育的學院.

[15]  可參閱Ng 2014.

[16] 參Ng 2014,頁89.

[17] 參王治心 1925,頁11-16.



\newpage

\section{為福音作見證}
\label{sec:669}
\textbf{蘇成溢牧師 - 為福音作見證}
\newline
\newline
\hyperref[sec:668]{< < < PREV SERMON < < <}
~
\hyperref[sec:toc]{[返目錄]}
~
\hyperref[sec:670]{> > > NEXT SERMON > > >}


講題:為福音作見證 To be A Witness for the Gospel of Jesus

講員:蘇成溢牧師

場合:香港中文大學崇基學院禮拜堂主日崇拜

日期:2023年10月29日

1. 前言

見證是基督徒熟悉的用語, 教會也常常鼓勵信徒要為所信的道作見證. 對今日大多數的基督徒來說, 作見證就是, 向人傳講所信的福音內容, 口述耶穌基督為人成就的救贖故事. 然而, 按新約聖經原文, 見證一字 (martus) 並非單作口傳, 乃指以生命為信仰作證, 甚至作出犧牲. 因此, 在早期教會, 見證的人也成為了殉道者 (martyr).

*根據教會歷史所載, 士每拿主教坡旅甲 (Polycarp 69-155 AD), 這位被詳細紀錄如何殉道的見證人/殉道者, 他為了堅守所信的福音 (相傳他是使徒約翰的門生), 被捆綁在柱上燒死的.

*中國基督教第一位殉道者名叫車錦江. 他在1856年初聞福音, 於香港由理雅各牧師施洗. 從1857-1860年他多次領人到香港受洗. 由於車錦江因熱心傳福音, 努力領人信主, 被當時廣東博羅的鄉紳禁止和威嚇, 最終因誓死護教被殺害, 碎屍投入東江. 1961年中華基督教會為紀念他殉道一百周年, 在現今黃大仙東頭村道興建了錦江紀念禮拜堂.

2. 帖撒羅尼迦教會

根據使徒行傳16-17章所載, 帖撒羅尼迦是使徒保羅在歐洲宣教的第二站. 第一站在腓立比, 雖有成效卻因助人驅走邪靈, 被誣告收監,結果要離開. 之後保羅到了帖撒羅尼迦, 連續三個安息日在當地猶太會堂傳講耶穌基督的福音, 有不少人相信, 卻引起當地猶太人的攻擊威嚇, 又被迫離開.

剛建立不久的幼嫩的帖撒羅尼迦教會, 雖有信心和愛心, 能勇敢面對逼迫, 且有美好的德行. 然而, 教會面對的挑戰也不少, 包括: 持續受迫害; 有假冒的信徒混入教會, 製造分化; 也有信徒在試探中失腳; 和有信徒對基督再來時的盼望抱有懷疑等. 為此保羅特寫信去勉勵和教導他們, 堅固他們的信仰.

*從保羅和初期教會信徒的經歷, 說明了為福音作見證, 是隨時隨地要付出, 會受苦, 甚至以生命作代價的. 反觀今天, 不少基督徒所傳講的, 都是信者有好處得獎賞的信息, 忽略了耶穌基督所教導要捨己的吩咐, 更少提見證與殉道的關係.

3. 帖撒羅尼迦前書2章1-8節

1節 保羅在帖撒羅尼迦雖有被趕走的經歷, 但確信自己所作的工 (見證) 是有果效的

2節 在受迫害, 面對如戰爭般的處境, 仍能奮勇見證福音 (“強烈反對中” (和修本); “極大的奮鬥 “ (呂振中本); “艱巨的格鬥中” (思高本)

3-4節 保羅強調, 他所見證的福音, 不是那些以宗教為名為己謀利的騙子, 乃是上帝所認定和託付的

* 作為福音作見證的人, 由於知道自己是上帝所認定 (approved RSV) 和託付的 (entrusted RSV), 故此一切為福音作見證的事, 都必須向主問責 (accountable to God).

5-6節 保羅表明, 為福音作見證, 全是以遵主旨意, 討衪喜悅而作的; 並非為個人名利或聽從某些人的心意而行

7-8節 保羅進一步指出, 他用溫柔的態度對待帖撒羅尼迦教會, 不是用權威的方式去命令. 傳福音時他會如母親般去愛人, 甚至不惜付上生命, 領人歸於基督.

* 透過這八節經文所載有關保羅在帖撒羅尼迦教會的經歷和教導, 我們認識到, 對為福音作見證的人, 要有以下的要求: 1) 對主絕對忠心; 2) 對人以愛以誠; 3) 對己虛懷自律.

4. 總結與反省

以生命見證勝於語言技巧

*英國藝術評論家拉斯金 (John Ruskin 1819-1900) 一次在傍晚時分在家與友人共聚, 言談大家提到 “如何能看出誰是基督徒?” 這時拉斯金望向窗外遠處, 見到遠處正有點燈人正手拿火把, 在山坡上點亮路燈. 由於天色已晚, 看不到點燈人, 但由於手握的火把和點亮的燈, 就知道點燈人正沿山而上. 拉斯金立時有所領悟, 對友人說: “基督徒就是這樣, 人未必認識他, 甚至不能分辨誰是誰不是; 但, 人能從他們在人生路途上所點亮的燈.

*多年前我曾與一批香港及內地中學生的遊學團, 有兩星期到北美遊學交流. 在行程的最後一天, 一位來自內地的中三學生, 主動到我面前, 向我提問. 她說: “我從未聽過基督教, 也沒有見過牧師. 但經過了這兩星期行程, 我很想知道基督教. 可否請你用幾個關鍵詞, 讓我認識基督教.

換言之, 為福音作見證不在於口講甚麼, 乃在乎活出甚麼.



\newpage

\section{誰是叛教者}
\label{sec:670}
\textbf{林豪恩先生 - 誰是叛教者}
\newline
\newline
\hyperref[sec:669]{< < < PREV SERMON < < <}
~
\hyperref[sec:toc]{[返目錄]}
~
\hyperref[sec:671]{> > > NEXT SERMON > > >}


講題:誰是叛教者? The Apostate(s)?

講員:林豪恩先生

場合:香港中文大學崇基學院禮拜堂主日崇拜

日期:2023年11月5日

引言:

還記得年初上映的Netflix紀錄片《以神之名:信仰的背叛》嗎?這紀錄片上映當時震驚全亞洲,因為:「這部紀錄片系列聚焦於四名自稱先知的韓國領袖人物,探索令人膽寒的真實故事,揭發盲目信仰的黑暗面.」紀錄片的名稱直指這些宗教領袖背叛信仰,其實背叛信仰是怎樣一回事呢?如果信奉同一位神的各方互相指責背叛信仰,如何判斷哪一方背叛信仰,哪一方沒有背叛呢?有沒有一些判斷準則可供參考呢?例如在今天的福音經課,我們聽到耶穌指責文士和法利賽人背叛上帝,扭曲信仰.反過來說,文士和法利賽人一早已經指責耶穌褻瀆上帝和不遵守祖宗傳下來的規條.要分辨誰是誰非,或許要仔細考察指責內容,正如《以神之名:信仰的背叛》這紀錄影片,用了不少時間去考察及追蹤,拍攝那四名自稱先知的宗教領袖的相關片段,把他們的行為和操守展示人前.因此,讓我們再讀經文,細閱耶穌指責文士和法利賽人背叛信仰的內容.

一. 耶穌指責法利賽人叛教(馬太福音)

23:1那時,耶穌對眾人和門徒講論,

23:2說:「文士和法利賽人坐在摩西的位上,

23:3凡他們所吩咐你們的,你們都要謹守遵行.

從耶穌對文士和法利賽人的指責來看,能夠被指責為背叛信仰的,似乎不是一般人,而是有某種條件的人,就正如這紀錄片所揭示的,這些人都是領袖.不是說一般人不可能背叛信仰,只是一般人人微言輕,沒有影響力.位高權重的人就不同了,他們不但有影響力,更有強制別人的能力,例如文士和法利賽人擁有宣講律法的資格和權威,他們和其他信眾的地位和權力並不對等,就正如紀錄片中那些人宗教領袖和信眾的權力不對等的狀況類似.在他們的群體內,他們擁有懲罰不跟隨者的權力,最少也令不合作者感到害怕和威脅.在這種地位和權力差異的情況之下,有機會令宗教領袖產生以下耶穌指責文士和法利賽人的特徵:

1. 律人以嚴

23:4他們把難擔的重擔捆起來,擱在人的肩上,

不是說律人以嚴的都是背叛信仰者,然而對信眾施以過分壓力,甚至如《以神之名》紀錄片中剝削信眾的狀況,都非常有可疑.為甚麼這些狀況能夠出現呢?很可能這些背叛信仰的領袖說話義正詞嚴,形象正義凜然,彷如是上帝的代表,正義的化身.因此,他們要求信眾做甚麼,信眾不敢不從.

2. 待己以寬

23:4但他們自己一個指頭也不肯動.

可能他們自覺身分特殊,有別於一般人；他們認為自己的角色是教導其他人,監察其他人,然後對違規者予以糾正或施以懲罰.因此,對他們來說,規矩是定給其他人守的.

3. 要令人覺得自己重要

23:5他們一切所做的事都是要叫人看見,所以將佩戴的經文做寬了,衣裳的繸子做長了,

23:6喜愛筵席上的首座,會堂裡的高位,

23:7又喜愛人在街市上問他安,稱呼他拉比 .

耶穌這項指責會不會有商榷餘地呢?身為宗教領袖,讓人知道,讓人看見,進身高位,豈不是對宣揚宗教有利嗎?耶穌豈不是在馬太福音5:15說過「人點燈,不放在斗底下,而是放在燈臺上,就照亮一家的人」嗎?那麼,應當如何考量這種狀況呢?其實耶穌這句話未說完的,下一句說:5:16「你們的光也要這樣照在人前,叫他們看見你們的好行為,把榮耀歸給你們在天上的父.」第一,問題不是要不要讓人看見,而是看見甚麼.耶穌認為值得被人看見的是「好行為」；當時的宗教領袖要人看見的是甚麼呢?第二,讓人看見的目的是甚麼呢?耶穌的目的是能夠令人「把榮耀歸給在天上的父」,當時的宗教領袖叫人把榮耀歸給誰呢?

上述幾節經文中耶穌對文士和法利賽人背叛信仰的指責只是引言,具體內容記載在23:7節之後.不過,故事還未結束,耶穌和法利賽人互指對方為叛徒這種狀況延續至跟隨耶穌基督的人和法利賽人之間的互判,我們剛才讀的書信經課貼撒羅尼迦前書的下文可以感受到這種張力:

二．叛教的保羅對叛教的法利賽人

2:15這猶太人殺了主耶穌和先知,又把我們趕出去.他們不得上帝的喜悅,且與眾人為敵,

2:16不許我們傳道給外邦人使外邦人得救,常常充滿自己的罪惡.上帝的忿怒臨在他們身上已經到了極處.

使徒保羅指斥猶太人背叛上帝,相反,對猶太人來說,保羅才是叛徒,要除之而後快.有一次保羅口耶路撒冷聖殿過節,被同來過節的外邦猶太人認出,於是號召人群捉拿保羅,當時:

(使徒行傳)

21:30合城都震動,百姓一齊跑來,拿住保羅,拉他出殿,殿門立刻都關了.

21:31他們正想要殺他,有人報信給營裡的千夫長說:「耶路撒冷合城都亂了.」

其實,保羅也不否認自己的改變,保羅被捕後向猶太人講述自己「棄暗投明」的經歷:

22:3保羅說:「我原是猶太人,生在基利家的大數,長在這城裡,在迦瑪列門下,按著我們祖宗嚴緊的律法受教,熱心事奉上帝,像你們眾人今日一樣.

22:4我也曾逼迫奉這道的人,直到死地,無論男女都鎖拿下監.

就是在往大馬士革追捕基督徒的路上,保羅經歷耶穌向他顯現,吩咐他向外邦人傳福音.於是,保羅由猶太人的英雄變為叛徒,猶太人對他咬牙切齒:

22:22眾人聽他說到這句話,就高聲說:「這樣的人,從世上除掉他吧!他是不當活著的.」

猶太人指控保羅背叛,保羅指控猶太人背叛,我們又問那個問題:如何判斷背叛與否呢?我們不妨把耶穌對文士和法利賽人作為宗教領袖的指責,對比一下保羅作為使徒在帖撒羅尼迦前書的自述,看看有沒有甚麼發現.

1:對待自己

2:9弟兄們,你們記念我們的辛苦勞碌,晝夜做工,傳上帝的福音給你們,免得叫你們一人受累.

馬太福音23:4「但自己一個指頭也不肯動」對「辛苦勞碌,晝夜做工」

2:對待信仰

2:10我們向你們信主的人,是何等聖潔,公義,無可指摘,有你們作見證,也有上帝作見證.

馬太福音23:23「你們將薄荷,茴香,芹菜獻上十分之一,那律法上更重的事,就是公義,憐憫,信實,反倒不行了.」

3:對待信眾

2:11你們也曉得,我們怎樣勸勉你們,安慰你們,囑咐你們各人,好像父親待自己的兒女一樣,

馬太福音23:4他們把難擔的重擔捆起來,擱在人的肩上.」

4:帶領信眾朝往的目的地

2:12要叫你們行事對得起那召你們進他國,得他榮耀的上帝.

馬太福音23:13「你們正當人前,把天國的門關了,自己不進去,正要進去的人,你們也不容他們進去」.

5．給予信眾的生命

2:13為此,我們也不住地感謝上帝,因你們聽見我們所傳上帝的道就領受了；不以為是人的道,乃以為是上帝的道.這道實在是上帝的,並且運行在你們信主的人心中.

馬太福音23:15「你們走遍洋海陸地,勾引一個人入教,既入了教,卻使他作地獄之子.」

總結:耶穌教導門徒如何避免跌入法利賽人背道的危險

其實在馬太福音耶穌指責文士和法利賽人的經文,中間也插入了對門徒的吩咐,教導他們如何避免重蹈文士和法利賽人的覆轍.使待保羅的表現似乎正是耶穌的教導的實踐.

1. 平等相待

23:8但你們不要受拉比的稱呼,因為只有一位是你們的夫子；你們都是弟兄.

2. 同奉天父

23:9也不要稱呼地上的人為父,因為只有一位是你們的父,就是在天上的父.

3. 同師基督

23:10也不要受師尊的稱呼,因為只有一位是你們的師尊,就是基督.

4. 謙虛服務

23:11你們中間誰為大,誰就要作你們的用人.

23:12凡自高的,必降為卑；自卑的,必升為高.

宗教領袖按著耶穌這幾項吩咐而行,宗教群體應該會相對安全.如果反其道而行,就可以產生危險,正如紀錄片《以神之名:信仰的背叛》所揭露的其中一個原名叫基督教福音宣教會的攝理教,教主鄭明析自稱「彌賽亞」.由此可見,他和信眾並非「弟兄」的平等關係,他自視為師,自視為父,凌駕於弟兄之上,形成為所欲為的條件.如果門徒互為弟兄,同奉天父,同尊基督為師,就有機會避免有人被捧上凌駕其他人的位置.

背叛信仰的現象在歷史中不斷發生,有時出現在好像紀錄片《以神之名》那些新興教派,有時可以出現在主流教會之內.讓我們學習誠實現對教會歷史,成為我們的警惕,以耶穌基督的教導指導我們的個人操守及團體運作.



\newpage

\section{事奉神諗清楚好喎}
\label{sec:671}
\textbf{李均熊牧師 - 事奉神?諗清楚好喎}
\newline
\newline
\hyperref[sec:670]{< < < PREV SERMON < < <}
~
\hyperref[sec:toc]{[返目錄]}
~
\hyperref[sec:672]{> > > NEXT SERMON > > >}


講題:事奉神?諗清楚好喎! Serving God? You Better Think Twice!

講員:李均熊牧師

場合:香港中文大學崇基學院禮拜堂主日崇拜

日期:2023年11月12日

起:你要事奉神?諗清楚喎……

今日舊約經課是約書亞記最後一章.經文的處境是以色列人由曠野漂流到現在可以在應許地安居,而他們的領袖約書亞亦走到人生最後一程,向這一班神的子民作最後的宣講.經課中沒有選用的4-13節乃係約書亞轉述上主的說話,提到以色列人由先祖在「大河」那邊,到下埃及,出埃及,過紅海而入到曠野的歷史,和上主指出他們後來到約旦河東與亞摩利人和進入迦南和各個民族爭戰的事蹟,以色列的種種遭遇都是上主的作為:「我將你們的祖宗……帶出來…….我差遣摩西和亞倫……把你們領出來…我領你們祖宗出埃及……領你們到約旦河東……我把他們交在你們手中,你們就得了他們的地為業……我賜給你們的地,不是你們開墾的；我賜給你們的城鎮,不是你們建造的.你們卻住在其中,又得吃那不是你們栽植的葡萄園和橄欖園的果子……」說完這一些上主的「提醒」之後,約書亞便要求以色列人作出選擇:

[14] 「現在你們要敬畏耶和華,誠心實意地事奉他,將你們列祖在大河那邊和在埃及所事奉的神除掉,去事奉耶和華. [15] 若是你們以事奉耶和華為不好,今日就可以選擇所要事奉的:是你們列祖在大河那邊所事奉的神呢?是你們所住這地的亞摩利人的神呢?至於我和我家,我們必定事奉耶和華.」

當中一句——至於我和我家,我們必定事奉耶和華——都是很多信徒熟識,甚至引用的「金句」!但若你仔細留意上文下理或當中可能的語氣,或許你會感到一絲約書亞的「不滿」:「你地揀D A貨啦,我就會揀真ye 啦!」(有無可能潛台詞係「你地咪咁蠢啦」?)

點解會咁樣覺得,你再看下去,當以色列人異口同聲說:我們絕不離棄上主,約書亞連隨說:

[19] 約書亞對百姓說:「你們不能事奉耶和華；因為他是聖潔的神,是忌邪的神,必不赦免你們的過犯罪惡. [20] 你們若離棄耶和華去事奉外邦神,耶和華在降福之後,必轉而降禍與你們,把你們滅絕.」

以色列人一而再堅定地說:不,我們要事奉上主,約書亞的回應也好像不太積極和相信:你們選擇上主,要事奉祂,你們自己作個證吧!(係先好講呀!)

我和我家,我們必定事奉耶和華,但係你要事奉神?諗清楚好喎!

承:點解約書亞要咁講?

點解約書亞會咁講呢?

縱觀以色列人由出埃及到入迦南,他們的往積(track record)都係十分之差:一出埃及在紅海遇到法老軍兵便質疑摩西要帶他們出曠野受死；一到曠野就埋怨無水無食物；遇見應許地要對付的是一群巨人就懦弱退縮,仲話寧願唔入迦南,於是上主要懲罰他們在曠野漂流40年 才可以進入應許之地.以色列人一路以來的表現,也難怪約書亞對他們是否會專一事奉上主,其間不會背叛他們的神, 心存疑慮.在聖經以後的記載,又的而且確給約書亞「講中」,以色列人因為不專一跟從上主而遭到亡國之災.可能值得過世了在天家的約書亞感到安慰的,是起碼在現場立石為証的以色列沒有背棄上主,正如約書亞記最後也說:「約書亞在世的日子和他死了以後,那些知道耶和華為以色列所做一切事的長老還在世的時候,以色列人事奉耶和華.」(書24:31)

轉:對主忠誠——真係no matter what?

(要諗清楚喎….諗D 乜?)

可見在平順日子,在應許之地,要立志忠信事奉上主,表面看來應該沒有甚麼問題.約書亞所質疑的,可能是在往後的日子,我們的際遇並不會一帆風順,在危難困難時,我們是否仍然會持守信仰,一心跟隨這個曾帶我們出死入生的神呢?

更何況:

上主是聖潔忌邪(jealous)的,所以任何將祂的愛看成廉價恩典, 用祂的愛作幌子隨意犯罪,都會惹祂厭惡.

上主是聖潔忌邪的,所以祂有超越我們凡俗心思的旨意,任何自把自為按私意事奉都會帶來祂的審判 你 斜 有好多 零 全部 兩只 快.

上主是聖潔忌邪的,所以任何替代祂,濫用祂或以為可以將祂私有化,要祂隨傳隨到為自己服務,都會惹祂忿怒.

上主要求的是我們要忠心於祂,無論任何境況,或順或逆,我地都要放下自我,以佢為主為首,no matter what!

但只要稍有偏離,我們就有機會越行越遠,我們必須盡早悔改回頭.否則,一如約書亞所說:你們若離棄耶和華去事奉外邦的神明,耶和華再降福之後,必轉而降禍給你們,把你們滅絕!今日我們或者不會拜黃大仙,財神,觀音,但這裡所說外邦的神明,用現代的說法乃是說我們以非神,代替真神；以不是上主來代替上主,無論那是金錢,愛情,權力,自己,名譽,家庭,成就.

合:一生的抉擇

第二次世界大戰時期,對抗納粹的德國神學家潘霍華(Dietrich Bonhoeffer)在他的著作「追隨基督」(The Cost of Discipleship)有一句說話:當基督呼召一個人,祂是在呼召他來為祂死　(When Christ calls a man, He bids him come and die).

事奉上主不是一時一刻感覺良好的決定.乃是無時無刻,一生的抉擇.Whenever, whereever, no matter what!

你要事奉神,真係諗清楚好喎……



\newpage

\section{啊上帝真忙}
\label{sec:672}
\textbf{鄧瑞強博士 - 啊上帝真忙}
\newline
\newline
\hyperref[sec:671]{< < < PREV SERMON < < <}
~
\hyperref[sec:toc]{[返目錄]}
~
\hyperref[sec:673]{> > > NEXT SERMON > > >}


詩篇 82:1-8
\newline
\begin{longtable}{cl}
\hline
\hline
章節 & 經文 (和合本修訂版)\\
\hline
82:1 & \begin{tabularx}{0.7\textwidth}{X} 神站立在神聖的會中，在諸神中施行審判。 \end{tabularx} \\ \\ \relax
82:2 & \begin{tabularx}{0.7\textwidth}{X} 你們審判不秉公義，抬舉惡人的臉面，要到幾時呢？（細拉） \end{tabularx} \\ \\ \relax
82:3 & \begin{tabularx}{0.7\textwidth}{X} 當為貧寒的人和孤兒伸冤，為困苦和窮乏的人施行公義。 \end{tabularx} \\ \\ \relax
82:4 & \begin{tabularx}{0.7\textwidth}{X} 當保護貧寒和貧窮的人，救他們脫離惡人的手。 \end{tabularx} \\ \\ \relax
82:5 & \begin{tabularx}{0.7\textwidth}{X} 他們愚昧，他們無知，在黑暗中走來走去；地的根基都搖動了。 \end{tabularx} \\ \\ \relax
82:6 & \begin{tabularx}{0.7\textwidth}{X} 我曾說：「你們是諸神，都是至高者的兒子。 \end{tabularx} \\ \\ \relax
82:7 & \begin{tabularx}{0.7\textwidth}{X} 然而，你們要死去，與世人一樣，要仆倒，像任何一位王子一般。」 \end{tabularx} \\ \\ \relax
82:8 & \begin{tabularx}{0.7\textwidth}{X} 神啊，求你起來審判全地，因為你必得萬國為業。 \end{tabularx} \\ \\
[1ex]
\hline
\hline
\end{longtable}
講題:啊,上帝真忙! Oh, How Busy God is!

經文:詩篇82:1-8

講員:鄧瑞強博士

場合:香港中文大學崇基學院禮拜堂主日崇拜

日期:2023年11月19日

各位弟兄姊妹,早安.願神的恩典常與你們同在.

今日,我和大家分享「詩歌智慧文學」的第二卷書《詩篇》.在實體聖經裡,詩篇大概位於聖經的中央位置.除了幾篇「智慧詩」外,詩篇不算是智慧文學.智慧文學不涉及救贖這主題,但詩篇卻大量祈求神的解救.詩篇是詩歌,用於敬拜禮儀裡,是唱出來的祈禱文.詩歌文字的特點,是平行,對比,採用很多象徵,聯想,文字是押韻的,但這在翻譯過程中就無法保留了.

作為祈禱文,詩篇有不同類型,最常見的,有三類.一,讚美詩:讚美神的偉大與榮耀.二,祈求詩(或者叫「哀歌」):人生有苦情,有危機,自然地,便會求神解救.有群體的祈求詩,也有個體的祈求詩.三,感恩詩:苦難過後,雨過天青,便感謝神的解救.祈求詩和感恩詩,都會涉及救贖這主題.

詩人將生命的所有實況,喜樂的,悲傷的,平安的,恐懼的,期盼的,絕望的,無論什麼際遇,都像祭物一樣,獻呈在神面前.無論境況順逆,詩人都滿懷信心,仰望神.一代又一代的詩人,他們浸淫在神的恩典裡,讓人生的一切經驗,釀成豐富生命的美酒.在這些祈禱裡,就算是痛不欲生的經驗,仍在信仰中體會到顧念蒼生的上帝的臨在.耶穌在十架上,在極度悲苦之中,他的呼喊:「我的神!我的神!為什麼離棄我?」這呼喊本身就是詩篇的禱文,出於詩篇22:1.死前,耶穌說:「父啊!我將我的靈魂交在你手裡.」這信靠神的話,也是詩篇的禱文,出於詩篇31:5.連神的兒子,在最深刻的痛苦裡,他的禱文,都是出自詩篇.詩篇,可說是將人和神扣得最緊的經文.詩篇,從人角度而言,是詩人懷著喜怒哀樂的不同情緒,真誠地走向神.詩篇,從神角度而言,是神親自表述,祂以悲憫擁抱著悲苦的人間.在這些禱文裡,人和神,心和心連在一起.無怪乎靈修大師Thomas Merton(梅頓)說:詩篇可理解為耶穌向天父的祈禱.人以詩篇來祈禱,就是在基督裡,以基督的話來祈禱.詩篇,是人和神最深邃的交流.

今日,和大家分享詩篇82篇.這首詩難於歸類,因其內容有多種元素,故此,一般學者將之歸類為「綜合詩」.這首詩的內容,包含:對神至高權柄的尊崇,神作為王對天地的審理,神親自宣告的神諭,先知式的批判,對人間苦難的控訴,對美好世界的嚮往,也有禮儀性的祈求,集「王者—先知—禮儀」的元素於一身.

今日的講道經文:詩篇82篇

亞薩的詩.

詩82:1　上帝站在有權力者的會中,在諸神中行審判,

詩82:2　說:你們審判不秉公義,徇惡人的情面,要到幾時呢?(細拉)

詩82:3　你們當為貧寒的人和孤兒伸冤；當為困苦和窮乏的人施行公義.

詩82:4　當保護貧寒和窮乏的人,救他們脫離惡人的手.

詩82:5　你們仍不知道,也不明白,在黑暗中走來走去；地的根基都搖動了.

詩82:6　我曾說:你們是神,都是至高者的兒子.

詩82:7　然而,你們要死,與世人一樣,要仆倒,像王子中的一位.

詩82:8　上帝啊,求你起來審判世界,因為你要得萬邦為業.

這首詩被稱為「亞薩的詩」.亞薩是利未族人,其世系都是服事聖殿的人.亞薩是詩班長,或是領唱者.通常,服事聖殿的人,都強調神在聖殿中,坐著為王,因聖殿裡的約櫃之上,有施恩座,這就是神的座位.服事聖殿的人,也重視秩序,正如我們的崇拜,也有不同的環節,也有其禮儀秩序.每樣物品如何擺,每句話如何講,皆有其規定.在聖殿中事奉,一板一眼,絲毫不得馬虎.上帝的王權,天地的秩序,在詩篇82篇裡,是兩大主題.

詩篇82篇的經文,有工整的平行結構:

第1節及第8節:你會見到「審判」這字眼,首尾呼應,講至高神的審判

第2-4節:講「審問」的案情；第6-7節作對應:講「判罪」的議決

中間第5節:點出問題的核心,講明「判決」的理據

詩82:1　上帝站在有權力者的會中,在諸神中行審判.

經文採用神話式畫面,模擬天上開神明大會,耶和華上帝主持會議.或者說,是上帝開庭,審判各方神明.在神話結構中,天上的眾神明,要為人間的福樂負責.明顯,人間有禍患,上帝要找相關的神明來審問.用現今的話來說,這些小神明的「KPI」不行,上帝找他們來問責.

詩82:2　說:你們審判不秉公義,徇惡人的情面,要到幾時呢?(細拉)

詩82:3　你們當為貧寒的人和孤兒伸冤；當為困苦和窮乏的人施行公義.

詩82:4　當保護貧寒和窮乏的人,救他們脫離惡人的手.

至高神直斥一眾直接掌控人間事務的小神明.可以想像,上帝是拍著檯說的.上帝斥責他們什麼呢?他們審判時,不按公義.人間的審判,最重要的責職,就是主持公道,正直地評道理.若應公道時不公道,應評理時不評理,則蒼世將如何,大地將如何?此外,這些小神明又徇惡人的情面,偏幫惡人,有些聖經譯本按字面直譯作「抬舉惡人的面」,對惡人給足面子.當神明都偏幫惡人時,這世界的亂局可以想像,或者說,這世界的亂局不可以想像.在這裡,上帝一聲歎息,這「要到幾時呢?」(How Long?)在第2節的結尾,有一音樂符號「細拉」,我們已不知其真正意義,或許,這是一聲歎息的延長符號(fermata):「唉.........」

第3-4節列出很多小神明應該照顧的人:貧寒的人,無依無靠的孤兒,被欺壓的困苦人,一無所有的窮乏人.小神明應該照顧這些人,但偏偏不理這些人,卻關照惡人.真是顛倒是非.上帝要這些小神明,滿足四項「KPI」指標,完成四項工作(經文用四個動詞來表達):為弱者「伸冤」,為被欺壓者「施行公義」,「保護」貧寒者,「解救」無助者免受欺凌.這些應有的行動,小神明們無一樣做到.「要到幾時」(How Long)這些小神明們才會停止助紂為虐?

這三節經文,真令人感動.神像先知一樣,親自在天地間吶喊.祂像一個荒原上的孤獨者,在天地間的眾神明面前,為人間的最弱者伸冤.當人間的弱者連為自己的痛苦都無法發出呼聲時,神為他們發出響遍宇宙的最大吶喊.我覺得,將生命交給這樣的神,是人生最幸福的事.

詩82:5　你們(應譯為「他們」)仍不知道,也不明白,在黑暗中走來走去；地的根基都搖動了.

這節經文是詩篇82篇的中心經文,位於這詩篇平行結構的正中位置.這經文點出人間正義的重要性,是上帝判案時的理據所在.

經文由「你們」轉為「他們」,如此,對象就不單是天上的小神明,也關乎他們地上的代理人.在神話佈局中,天上與人間是連成一體的,天上的神明主宰著人間的事情,神明的決定靠賴一些代理人去執行.如今,上帝一併向天上的小神明和他們人間的代理人發言.

為何正義不彰,是非顛倒?原因是他們不知道,不明白.他們不知道什麼?他們不知道正義是天地立穩的根基.他們不明白什麼?他們不明白正義是支撐天地的棟梁.無論是上天的小神明,或他們地上的代理人,竟領悟不到,對人間弱者的憐愛,正是天地間的大道.他們的無知,是對人間終極真相的絕對無知.經文說,他們「在黑暗中走來走去」,像鬼鬼祟祟般,在背後做些見不得光的事情.由於應該主持正義的,沒有主持正義,於是,「地的根基都搖動了」,天地崩壞,秩序瓦解.再沒有一節經文,像這經文一樣,講出人間正義的重要性.沒有了仁愛,失去了正義,天地便分崩離析了,又回到太初的混沌.

詩82:6　我曾說:你們是神,都是至高者的兒子.

詩82:7　然而,你們要死,與世人一樣,要仆倒,像王子中的一位.

第2-4節講出了案情,權力偏袒惡人,弱者反被欺壓.第5節講出判案時的理據,正義是支撐天地的棟梁,不義令天地崩壞.第6-7節便宣判刑罰.

對象由「他們」又轉回「你們」,神向小神明們宣判刑罰.天上的小神明,由於他們失職,沒有好好管理人間,反令人間大亂,故此,他們失去其不朽的生命.他們失去了神聖,失去將來,不再被人記念,更沒有人向他們膜拜.神明般的存在,被貶為凡人,失去了存在的光采,在歷史裡被遺忘.這是何等大的失落!

原來,神聖和正義之間,有一種微妙的關連.堅持正義者,才配稱神聖.神聖者,必堅守正義.失去正義,就失去神聖.在主禱文裡,我們說,「我們在天上的父,願人都尊你的名為聖」.天上的父配稱為聖,是因祂主持人間的正義.在苦難的人生,有神主持正義,這是人間的希望.

詩82:8　上帝啊,求你起來審判世界,因為你要得萬邦為業.

來到這首詩的結尾,是一祈求.在這裡,你會發現,其他神明退場了,只留下真正的神,那位歷史的主,那位世事的審判者,那位主持正義的神.詩人求神起來,起來,就是有所行動.詩人求神行動,執行判決,回復人間正義.神如此行,因為祂要得萬邦為業,這正是主禱文所求的,願神的國降臨,全地回到祂的王權之下.這是對天國降臨的終極期盼,這是對大地與上帝復和的終極期待,這是信徒對歷史的發展的終極期許.

這首詩是一首很簡短的詩.但在短短八節經文裡,卻看到上帝忙得不得了.祂為人間秩序的分崩離析而操心,祂召開會議,或者說,開庭審議.祂肯定細心聆聽過貧寒人的苦情,用心體會過被欺壓者的歎息.祂將人間的不義,一一陳明在眾神明面前.祂斥責這些小神明,點出問題的嚴重性.維護不義,就是動搖天地的棟梁.正義,就是神聖的根本.失去正義,就算連神明,都要貶為凡人.而人類的希望,就在於正義的神起來行動,掌管萬邦,讓天國降臨.天地不是無情的,上帝仁愛的情眷顧著受苦的人.天地不是無義的,上帝正直的義默默地扭轉著乾坤.讓我們信靠上帝,懷著希望,好好地活,好好地活出仁愛的正義.

願上主得著榮耀,頌讚；願眾人得享喜樂,平安.



\newpage

\section{主啊我甚麼時候見到你呢}
\label{sec:673}
\textbf{高國雄牧師 - 主啊!我甚麼時候見到你呢?}
\newline
\newline
\hyperref[sec:672]{< < < PREV SERMON < < <}
~
\hyperref[sec:toc]{[返目錄]}
~
\hyperref[sec:674]{> > > NEXT SERMON > > >}


馬太福音 25:31-46
\newline
\begin{longtable}{cl}
\hline
\hline
章節 & 經文 (和合本修訂版)\\
\hline
25:31 & \begin{tabularx}{0.7\textwidth}{X} 「當人子在他榮耀裡，同著眾天使來臨的時候，要坐在他榮耀的寶座上。 \end{tabularx} \\ \\ \relax
25:32 & \begin{tabularx}{0.7\textwidth}{X} 萬民都要聚集在他面前。他要把他們分別出來，好像牧人分別綿羊、山羊一般， \end{tabularx} \\ \\ \relax
25:33 & \begin{tabularx}{0.7\textwidth}{X} 把綿羊安置在右邊，山羊在左邊。 \end{tabularx} \\ \\ \relax
25:34 & \begin{tabularx}{0.7\textwidth}{X} 於是王要向他右邊的說：『你們這蒙我父賜福的，可來承受那創世以來為你們所預備的國。 \end{tabularx} \\ \\ \relax
25:35 & \begin{tabularx}{0.7\textwidth}{X} 因為我餓了，你們給我吃；渴了，你們給我喝；我流浪在外，你們留我住； \end{tabularx} \\ \\ \relax
25:36 & \begin{tabularx}{0.7\textwidth}{X} 我赤身露體，你們給我穿；我病了，你們看顧我；我在監獄裡，你們來看我。』 \end{tabularx} \\ \\ \relax
25:37 & \begin{tabularx}{0.7\textwidth}{X} 義人就回答：『主啊，我們甚麼時候見你餓了，給你吃；渴了，給你喝？ \end{tabularx} \\ \\ \relax
25:38 & \begin{tabularx}{0.7\textwidth}{X} 甚麼時候見你流浪在外，留你住；或是赤身露體，給你穿？ \end{tabularx} \\ \\ \relax
25:39 & \begin{tabularx}{0.7\textwidth}{X} 又甚麼時候見你病了，或是在監獄裡，來看你呢？』 \end{tabularx} \\ \\ \relax
25:40 & \begin{tabularx}{0.7\textwidth}{X} 王回答他們說：『我實在告訴你們，這些事你們做在我弟兄中一個最小的身上，就是做在我身上了。』 \end{tabularx} \\ \\ \relax
25:41 & \begin{tabularx}{0.7\textwidth}{X} 「王又要向那左邊的說：『你們這被詛咒的人，離開我！進入那為魔鬼和他的使者所預備的永火裡去！ \end{tabularx} \\ \\ \relax
25:42 & \begin{tabularx}{0.7\textwidth}{X} 因為我餓了，你們沒有給我吃；渴了，你們沒有給我喝； \end{tabularx} \\ \\ \relax
25:43 & \begin{tabularx}{0.7\textwidth}{X} 我流浪在外，你們沒有留我住；我赤身露體，你們沒有給我穿；我病了，我在監獄裡，你們沒有來看顧我。』 \end{tabularx} \\ \\ \relax
25:44 & \begin{tabularx}{0.7\textwidth}{X} 他們也要回答：『主啊，我們甚麼時候見你餓了，或渴了，或流浪在外，或赤身露體，或病了，或在監獄裡，沒有伺候你呢？』 \end{tabularx} \\ \\ \relax
25:45 & \begin{tabularx}{0.7\textwidth}{X} 王要回答：『我實在告訴你們，這些事你們沒有做在任何一個最小的弟兄身上，就是沒有做在我身上了。』 \end{tabularx} \\ \\ \relax
25:46 & \begin{tabularx}{0.7\textwidth}{X} 這些人要往永刑裡去；那些義人要往永生裡去。」 \end{tabularx} \\ \\
[1ex]
\hline
\hline
\end{longtable}
講題:主啊!我甚麼時候見到你呢?

經文:馬太福音25:31-46

講員:高國雄牧師

場合:香港中文大學崇基學院禮拜堂主日崇拜

日期:2023年11月26日

基督君王日

31「當人子在他榮耀裡,同著眾天使來臨的時候,要坐在他榮耀的寶座上. 32萬民都要聚集在他面前.他要把他們分別出來,好像牧人分別綿羊,山羊一般, 33把綿羊安置在右邊,山羊在左邊. 34於是王要向他右邊的說:『你們這蒙我父賜福的,可來承受那創世以來為你們所預備的國. 35因為我餓了,你們給我吃；渴了,你們給我喝；我流浪在外,你們留我住； 36我赤身露體,你們給我穿；我病了,你們看顧我；我在監獄裡,你們來看我.』 37義人就回答:『主啊,我們甚麼時候見你餓了,給你吃；渴了,給你喝? 38甚麼時候見你流浪在外,留你住；或是赤身露體,給你穿? 39又甚麼時候見你病了,或是在監獄裡,來看你呢?』 40王回答他們說:『我實在告訴你們,這些事你們做在我弟兄中一個最小的身上,就是做在我身上了.』

41「王又要向那左邊的說:『你們這被詛咒的人,離開我!進入那為魔鬼和他的使者所預備的永火裡去! 42因為我餓了,你們沒有給我吃；渴了,你們沒有給我喝； 43我流浪在外,你們沒有留我住；我赤身露體,你們沒有給我穿；我病了,我在監獄裡,你們沒有來看顧我.』 44他們也要回答:『主啊,我們甚麼時候見你餓了,或渴了,或流浪在外,或赤身露體,或病了,或在監獄裡,沒有伺候你呢?』 45王要回答:『我實在告訴你們,這些事你們沒有做在任何一個最小的弟兄身上,就是沒有做在我身上了.』 46這些人要往永刑裡去；那些義人要往永生裡去.」

今天是基督君王日,提醒我們耶穌基督作為救恩的泉源.

基督君王日是教會年的最後一個主日,這節期的代表顏色為白色,焦點為慶祝救主耶穌基督為王,並表達基督再來的威榮.這主日上承「聖靈降臨日」後漫長的「常年期」,下接新一個教會年的新開始:「將臨期」.

今天是宣告基督王權的全面展現,成全完備的救恩的大日子；同時也是對基督將要來的大君王的預表.

1919年,第一次世界大戰結束,戰勝國無視戰敗國的民族利益,強行重新規劃新世界的地圖.

戰後部份歐洲國家經濟萎縮,因而產生民族之間的仇視.

在這樣的情勢下,民族主義高漲,希特拉,墨索利尼等領袖堀起,被視為民族的救星.

由於他們強調自己民族的優越性,所以轉而欺壓其他民族.

有見於這類極端民族主義的情緒日漸高漲,教宗庇護十一世於1925年12月11日遂以《設立基督君王日》通諭(QuamPrimas)宣布訂定每年十月最後一個主日為「耶穌基督君王禮」.

盼望透過這個教會的節日,提醒信徒救主並非俗世裡任何一位魅力領袖,而是內心柔和謙卑的耶穌基督.

祂與屬世的領袖大不相同,祂甘願為世人得救的緣故而犧牲自己.深願世人認清基督為世界的主,尊崇基督為王,停止干戈,使世人得享和平幸福.

1969年,在教宗保祿六世重新修訂的羅馬普通曆法中,這節日被更改為「常年期」的最後一個主日.

教會把基督君王日,放在教會年的最後一個主日慶祝,作為教會年的終向,以示萬物都要在基督權下,獲得圓滿的救贖.

所以今天就是「基督君王日」.

信徒實在很容易忘記「基督才是我們生命的王」

當「權力等同真理」,「誰大誰惡誰就正確」變成常態的時候,人就任意而行.

正如約翰福音9:10

彼拉多對耶穌說:「難道你不知道我有權柄釋放你,也有權柄把你釘十字架嗎?」

惡人以為自己擁有權力就能主宰一切.

約翰福音18:36

耶穌基督清楚表示:「我的國不屬於這世界」.

祂不是倚仗權力,聲望及財富來統治的；

祂不是一位高高在上,威嚴地發號施令的君王,而是一位服侍人的君尊義僕.

基督君王最具體的表現是在十字架上,履行與天父一致的心意,愛,寬恕並拯救願意認罪尋求赦免的人.

基督君王日同時是指向這位將要回來的君王耶穌基督—將要回來審判萬民的君王,祂是超過世上一切所有的國度和王權.

講到末日審判,主耶穌明說有些人要蒙主賜福,要承受永生；

但有些人要被咒詛,要離開主往永刑裡去.到底祂是在講誰呢?

是指那些沒有信主的人嗎?

那為何都是主的羊,平日裡一起生活一起吃草,也被同一個牧人眷養,最後結果卻有這麼大的差別?

V31-33節

當人子在他榮耀裡,同著眾天使降臨的時候,要坐在他榮耀的寶座上.萬民都要聚集  在他面前.他要把他們分別出來,好像牧羊的分別綿羊山羊一般,把綿羊安置在右邊,山羊在左邊.

主再來時將施行審判,「他要把他們分別出來」

這工作是極其重要也無人能取代.

若綿羊山羊的區分,如果只是看他們是否信耶穌上教堂做禮拜,現在誰都可以看出來,何必要在未後主再來才能揭曉答案?

因此我們有理由相信,無論綿羊還是山羊,平日裡都在同一個牧場生活,一起吃草一起生活,也被同一個牧人牧養.

正如前面幾個比喻所講,無論是領一千兩千或五千的僕人,平日都一起同工,服事同一位主人；

十個童女都聚在一處等候新郎的來臨,分別只在有的預備有的沒預備.

所以我們明白,這幾個比喻其實都不是在講教外之人,而是指教會裡的人,即我們.

從耶穌的教導可以歸納出真正屬祂的羊有三個特徵:

第一,羊會認識他的主,

第二,羊會聽主的聲音,

第三,羊會跟從他的主.

約翰福音10章,耶穌說:「我是好牧人；我認識我的羊,我的羊也認識我,我的羊聽我的聲音,我也認識他們,他們也跟著我.我又賜給他們永生；他們永不滅亡,誰也不能從我手裡把他們奪去.

約翰福音14章

耶穌說:「你們若愛我,就必遵守我的命令.有了我的命令又遵守的,這人就是愛我的」

真正屬祂的羊必定愛祂,必定樂意遵守祂的命令,因此看他的行為如何便是最好的證明.

在巴勒斯坦一帶,綿羊和山羊經常在同樣的草地上同時出現.農夫,牧場主人和牧羊人通常允許牠們在同樣的草場上吃草和走動.這兩種動物都是常在一起的.

把牠們分散在不同的草場,是不必要.

除非把綿羊和山羊分開,否則牠們自由地閒逛,混合和交往.

況且,牠們不會打架,不會嫉妒,也不會排斥對方.

誰是綿羊?誰是山羊?上帝清楚知道.

時間和行為的考驗會證明,最後的結局將會把他們分別出來.

你的基督信仰是與基督建立關係,或只是沒有基督的宗教行為而已?

你是實踐真理的信徒,或只是名義上的教徒?

返到經文的重點

重點在於是否「見到」.

我們常常期望見到主,見到祂的作為,見到祂的心意…

上帝期望我們見到小弟兄,見到人的需要,見到我們應該要有的行動…

上帝何時才在我們身上作王?

就是我們見到的時候!



\newpage

\section{要不要照愚昧人的愚妄話回答他}
\label{sec:674}
\textbf{鄧瑞強博士 - 要不要照愚昧人的愚妄話回答他?}
\newline
\newline
\hyperref[sec:673]{< < < PREV SERMON < < <}
~
\hyperref[sec:toc]{[返目錄]}
~
\hyperref[sec:675]{> > > NEXT SERMON > > >}


箴言 26:1-12
\newline
\begin{longtable}{cl}
\hline
\hline
章節 & 經文 (和合本修訂版)\\
\hline
26:1 & \begin{tabularx}{0.7\textwidth}{X} 愚昧人得尊榮不相宜，正如夏天落雪，收割時下雨。 \end{tabularx} \\ \\ \relax
26:2 & \begin{tabularx}{0.7\textwidth}{X} 詛咒不會無故臨到，正如麻雀掠過，燕子翻飛。 \end{tabularx} \\ \\ \relax
26:3 & \begin{tabularx}{0.7\textwidth}{X} 鞭子是為打馬，轡頭是為勒驢，刑杖正是為打愚昧人的背。 \end{tabularx} \\ \\ \relax
26:4 & \begin{tabularx}{0.7\textwidth}{X} 不要照愚昧人的愚昧話回答他，免得你與他一樣。 \end{tabularx} \\ \\ \relax
26:5 & \begin{tabularx}{0.7\textwidth}{X} 要照愚昧人的愚昧話回答他，免得他自以為有智慧。 \end{tabularx} \\ \\ \relax
26:6 & \begin{tabularx}{0.7\textwidth}{X} 藉愚昧人的手寄信的，就像砍斷雙腳，喝下殘暴。 \end{tabularx} \\ \\ \relax
26:7 & \begin{tabularx}{0.7\textwidth}{X} 箴言在愚昧人的口中，正如瘸子的腳懸空無用。 \end{tabularx} \\ \\ \relax
26:8 & \begin{tabularx}{0.7\textwidth}{X} 將尊榮給愚昧人的，就像石頭綁在彈弓上。 \end{tabularx} \\ \\ \relax
26:9 & \begin{tabularx}{0.7\textwidth}{X} 箴言在愚昧人的口中，好像荊棘刺入醉漢的手。 \end{tabularx} \\ \\ \relax
26:10 & \begin{tabularx}{0.7\textwidth}{X} 雇愚昧人的，與雇過路人的，就像弓箭手射傷任何人。 \end{tabularx} \\ \\ \relax
26:11 & \begin{tabularx}{0.7\textwidth}{X} 愚昧人重複做愚昧之事，就如狗轉過來吃自己所吐的。 \end{tabularx} \\ \\ \relax
26:12 & \begin{tabularx}{0.7\textwidth}{X} 你看見自以為有智慧的人嗎？愚昧人比他更有指望。 \end{tabularx} \\ \\
[1ex]
\hline
\hline
\end{longtable}
講題:要不要照愚昧人的愚妄話回答他? Answer or Answer not a Fool According to His Folly?

經文:箴言26:1-12

講員:鄧瑞強博士

場合:香港中文大學崇基學院禮拜堂主日崇拜

日期:2023年12月3日

各位弟兄姊妹,早安.願神的恩典常與你們同在.

我們今日看聖經裡「智慧文學」的第二卷書《箴言》.《箴言》大量對比智慧與愚昧,義人與惡人,努力向善與自甘墮落的兩種人生.道理講得明明白白,但自甘墮落的人依然很多.在智慧與愚昧的對比中,《箴言》講出智慧的總綱.箴1:7:「敬畏耶和華是知識的開端,愚妄人藐視智慧和訓誨.」箴9:10:「敬畏耶和華是智慧的開端,認識至聖者便是聰明.」原來,智慧的基礎不在「智能」,而在「敬畏」.問題不關乎「智商」,而關乎人生善惡的抉擇.對智慧文學而言,愚昧不單是「蠢」的問題,也是善惡抉擇的「倫理」問題,這抉擇又關乎人生的終極信仰問題.大賊通常都是很聰明的,說其愚昧,問題不在智能,而在他們選擇了邪惡的路.如此選擇,又反映出他們的終極信仰.讀智慧文學,要明白:智慧或愚昧,與倫理抉擇,及神學信念,是三者合一的.

先看兩節《箴言》經文,這是今日講道的重點:

箴26:4　不要照愚昧人的愚妄話回答他,恐怕你與他一樣.

箴26:5　要照愚昧人的愚妄話回答他,免得他自以為有智慧.

《箴言》是否很有趣?到底,要不要照愚昧人的愚妄話回答他呢?

當你在茶樓,正開餐時,旁邊的人問你:「你在開餐嗎?」要答他抑或不答他呢?

我上網尋找如何回答愚蠢問題時,發現很多是教男人如何回答女人的蠢問題的.女人問男人的最蠢問題,是:「你是否覺得我好蠢?」當然,我們知道,面對這問題,一定要記得箴言的告誡,千萬不要回答,一答,便比問的人更蠢.

另有一條蠢問題,被喻為人類史上最難答的問題之一,老婆問老公:「我和你媽掉到水裡,你會先救誰?」男人為這問題,煩惱了很久.在網上,終於有人破局.老公以一個問題回答:「你愛我嗎?」以問題回答問題,很有猶太民族的色彩.老婆答:「我當然愛你啊.」老公便說:「愛我就不要問我這個問題.」老婆便追問:「若我不愛你呢!」她這一問便中伏了.老公便答:「那我為何要救你?」這一問一答,獲得網民喝采.據說,終於找到方法,解決這個世紀難題.

我們回到經文.箴26:4-5同時列出看來相反的告誡,到底用意何在?

有人說:這是因為世事無常,根本沒有一本通書可看到老,列出看來相反的告誡,就是要人隨機應變,以智慧去判斷,各獨特處境應如何對應.

有拉比說:在不同領域裡,有不同對應方式.在日常領域裡,智慧人就不要理會愚昧人了,免得將自己降至和愚昧人同一水平.但在關乎聖經或真理的情境時,就必須回答,免得愚昧人以為自己掌握了真理,這是很危險的.

又有人說:這不關乎情境,而關乎問問題的那人的質素.若他愚昧得無藥可救,就不要和他糾纏了.若他仍有開竅的機會,就應盡力幫助他.

究竟考慮的重點在人的質素,抑在處境的性質,抑或根本次次不同?

這兩節經文,到底用意何在?

學者發現,這兩節經文其實是一個獨特段落的組成部分.要理解這兩節經文的真義,必須從這經文段落的總體去理解才行.箴26:4-5是箴26:1-12的組成部分.這段經文前面的段落(即第25章)講社會倫理,後面段落(即26:13-16)講懶惰,而箴26:1-12這段落集中講「愚昧」,與前文和後理都不同,獨立出來,自成一段.

我們看看全段的經文:箴26:1-12(留意:差不多每節都提到愚昧)

箴26:1　夏天落雪,收割時下雨,都不相宜；愚昧人得尊榮也是如此.

箴26:2　麻雀往來,燕子翻飛；這樣,無故的咒詛也必不臨到.

箴26:3　鞭子是為打馬,轡頭是為勒驢；刑杖是為打愚昧人的背.

箴26:4　不要照愚昧人的愚妄話回答他,恐怕你與他一樣.

箴26:5　要照愚昧人的愚妄話回答他,免得他自以為有智慧.

箴26:6　藉愚昧人手寄信的,是砍斷自己的腳,自受損害.

箴26:7　瘸子的腳空存無用；箴言在愚昧人的口中也是如此.

箴26:8　將尊榮給愚昧人的,好像人把石子包在機弦裡.

箴26:9　箴言在愚昧人的口中,好像荊棘刺入醉漢的手.

箴26:10　雇愚昧人的,與雇過路人的,就像射傷眾人的弓箭手.

箴26:11　愚昧人行愚妄事,行了又行,就如狗轉過來吃牠所吐的.

箴26:12　你見自以為有智慧的人嗎?愚昧人比他更有指望.

1-3節成一小段落:

箴26:1　夏天落雪,收割時下雨,都不相宜；愚昧人得尊榮也是如此.

這節經文討論合不合宜的問題:「夏天落雪,收割時下雨,都不相宜」.然後作一平行的論斷:愚昧人與尊榮,也不相宜.

箴26:2　麻雀往來,燕子翻飛；這樣,無故的咒詛也必不臨到.

這節經文討論事情不會無緣無故臨到的,鳥兒在天空飛,不會無端落下.同理,無緣無故的咒詛,也不會無端的臨到人.言下之意,尊榮不會臨到愚昧人,但咒詛或惡運臨到愚昧人,就不是無緣無故的.

箴26:3　鞭子是為打馬,轡頭是為勒驢；刑杖是為打愚昧人的背.

鞭打驢馬,是因為牠們頑固,不聽教,同理,愚昧人像頑固的動物,不打不聽教.1-3節合起來,意思便是:愚昧人與尊榮不相配；相配的,是刑杖,而刑杖臨到愚昧人,不是無緣無故的,是因他們冥頑不靈.

6-8節又自成一小段落,對比著1-3節,而以4-5節為中心.

箴26:6　藉愚昧人手寄信的,是砍斷自己的腳,自受損害.

靠愚昧人做跑腿去幹某些事情,就像砍斷自己的腳一樣,自尋煩惱,因為那事情是不會完成的了.這節經文對應第3節.第3節講愚昧人因愚昧受刑杖,如今,講到信任愚昧人的,連自己也受害.愚昧人的禍害,不只自作自受,也令人受害.

箴26:7　瘸子的腳空存無用；箴言在愚昧人的口中也是如此.

瘸子的腳空存無用,有用的東西,被廢了武功,同理,智慧的箴言,對愚昧人而言,也是毫無用處.這節對應第2節,愚昧人受禍,真的不是無緣無故的.他們視智慧如無物,走滅亡之路.

箴26:8　將尊榮給愚昧人的,好像人把石子包在(譯作:「放在」更好)機弦裡.

這節經文很難解,參照箴26:10後,將「把石子放在機弦裡」解作胡亂地發射,傷害更多的人.這節對應第1節,重複「尊榮」與「愚昧人」不配.

來到這裡,可看到愚昧人真是無可救藥.傷害自己,傷害他人,視智慧如無物,故此,智慧人應避之則吉,不要回應他們,免得「蠢埋一堆」.這解釋了26:4的勸戒:「不要照愚昧人的愚妄話回答他,恐怕你與他一樣」.但第5節又如何解釋呢?我們繼續看經文的發展.

9-11節又自成一小段落,強化之前的討論.

箴26:9　箴言在愚昧人的口中,好像荊棘刺入醉漢的手.

「箴言在愚昧人的口中」與第7節的講法一樣,但第7節只是講箴言於愚昧人無用,現在卻講箴言會傷害愚昧人,看來,這是因為愚昧人錯解箴言,誤用箴言.其實,錯解聖經以致誤導自己,誤導他人,其造成的傷害也是很大的.

箴26:10　雇愚昧人的,與雇過路人的,就像射傷眾人的弓箭手.

這經文綜合第6及8節.靠愚昧人去做事的,必招損失.其損失之大,猶如將箭交給一個亂射的弓箭手,導致死傷無數.早前,某啤酒廠的某愚昧員工向釀酒原料撒了一泡尿,據報,翌日開巿,酒廠的巿值便蒸發了接近90億港元.雇用這愚昧員工的,其帶來的損失,難以估計.

箴26:11　愚昧人行愚妄事,行了又行,就如狗轉過來吃牠所吐的.

愚昧如醉酒的人,總說自己無醉.蠢人做蠢事,總能提出一些蠢的理由.真愚昧的人,是那些不知道自己是多麼愚昧的人.他們很難打破愚昧的循環.

講到這裡,看來仍只是重複之前的論點,找不到解釋第5節的線索.我們看看最後一節經文:箴26:12.

箴26:12　你見自以為有智慧的人嗎?愚昧人比他更有指望.

忽然來一破格的經文.之前的經文,多是平行句,而第4及5節是告誡,而這第12節,卻是一「比較」經文:Ａ比Ｂ好.在比較下,看不到有什麼優點的愚昧人,還不算是最差.比愚昧人更差的,是那些「自以為有智慧的人」.原先一大段經文,講了這麼長,以為只是數落愚昧人,教人不要做愚昧人,卻忽然來個急轉彎,指出有人比愚昧人更差,就是那些「懶醒」的人.聽說,有個少年人,向別人誇說自己懂聖經原文,其他人當然很欣賞他,真是年少出狀元.然後,這少年人慢慢拿出一本英文聖經來,教導其他人「原文」是怎樣寫的.真如箴言所講:「愚昧人比他更有指望」.

原來,第12節才是高潮.回頭望第5節,發現第12節正正對應第5節,兩節都出現「自以為有智慧」的字眼.於是,解開第5節的線索出現了.為什麼「要照愚昧人的愚妄話回答他」,是因為唯有這樣,才能叫醒那沉睡的人,才能讓那些「自以為有智慧」的人知道自己根本沒有智慧.智慧人若不點醒他們,則他們會自鳴得意,自以為是,而這比單純的愚昧更沒有希望.

第4-5節的意思,忽然豁然開朗.第4-5節的意思,原來是:「不要照愚昧人的愚妄話回答他」,不要與愚昧人打交道,近墨者黑,恐怕你會成為與他一樣愚昧.但,縱使有這危機,智慧人仍要跨越這勸戒,要豁出去,努力「照愚昧人的愚妄話回答他」,要與愚昧人打交道,為的是開他們的竅,免得他們在封閉的自我中「自以為有智慧」,而這比單純的愚昧更差.當智慧人放下身段,與愚昧人打交道時,可能真會變得像愚昧人,但,若這能帶一點智慧給愚昧人,叫愚昧人不那麼「自以為有智慧」,這已是一比較好的結果了.真智慧人,不怕落坑渠,去救助那些掉落坑渠而不自知的盲目人.這才是智慧文學講的「智慧」.

有人以耶穌作說明例子.他滿有智慧,是真理的化身.但為了與人打交道,他來到人間,並作出了最愚昧的抉擇,走上十架.藉此,讓在黑暗中「自以為有智慧」的人,看到自己的無知,從而走出盲目的自負.十架,是神的真智慧.

講到這裡,大家有沒有發現,當我們讀到智慧人與愚昧人對比的經文時,大家自自然然便代入智慧人的角度去讀.如:箴26:4-5:「不要照愚昧人的愚妄話回答他,恐怕你與他一樣.要照愚昧人的愚妄話回答他,免得他自以為有智慧.」讀到這經文時,我們心中自然有一問題:「那麼,我應照愚昧人的愚妄話回答他嗎?」我們自自然然便採用智慧人的視角,心中思想的,是回答抑或不回答愚昧人的問題.你可能沒有發現,由始至終,你就是那「自以為有智慧」的人.你一直認為自己有智慧,想著如何去應對身邊的愚昧人.但,這樣想的人,比愚昧人更沒有希望.

你們說:《箴言》是否很有趣呢?當你以為掌握其中的智慧時,來一個大逆轉,才發覺自己愚昧到極點.

願各人努力領悟神的智慧,記住:真正的智慧人,是明知與愚昧人打交道會變成愚昧,但仍甘於進入愚昧人當中,救助他們.

願上主得著榮耀,頌讚；願眾人得享智慧,平安.



\newpage

\section{黎緊咁即係未到}
\label{sec:675}
\textbf{李均熊牧師 - 黎緊……咁即係未到?}
\newline
\newline
\hyperref[sec:674]{< < < PREV SERMON < < <}
~
\hyperref[sec:toc]{[返目錄]}
~
\hyperref[sec:676]{> > > NEXT SERMON > > >}


以賽亞書 40:1-31
\newline
\begin{longtable}{cl}
\hline
\hline
章節 & 經文 (和合本修訂版)\\
\hline
40:1 & \begin{tabularx}{0.7\textwidth}{X} 你們的神說：「要安慰，安慰我的百姓。 \end{tabularx} \\ \\ \relax
40:2 & \begin{tabularx}{0.7\textwidth}{X} 要對耶路撒冷說安慰的話，向它宣告，它的戰爭已結束，它的罪孽已赦免；它為自己一切的罪，已從耶和華手中加倍受罰。」 \end{tabularx} \\ \\ \relax
40:3 & \begin{tabularx}{0.7\textwidth}{X} 有聲音呼喊著：「要在曠野為耶和華預備道路，在沙漠為我們的神修直大道。 \end{tabularx} \\ \\ \relax
40:4 & \begin{tabularx}{0.7\textwidth}{X} 一切山窪都要填滿，大小山岡都要削平；陡峭的要變為平坦，崎嶇的必成為平原。 \end{tabularx} \\ \\ \relax
40:5 & \begin{tabularx}{0.7\textwidth}{X} 耶和華的榮耀必然顯現，凡有血肉之軀的都一同看見，因為這是耶和華親口說的。」 \end{tabularx} \\ \\ \relax
40:6 & \begin{tabularx}{0.7\textwidth}{X} 有聲音說：「你喊叫吧！」我說：「我喊叫甚麼呢？」凡有血肉之軀的盡都如草，他的一切榮美像野地的花。 \end{tabularx} \\ \\ \relax
40:7 & \begin{tabularx}{0.7\textwidth}{X} 耶和華吹一口氣，草就枯乾，花也凋謝。百姓誠然是草； \end{tabularx} \\ \\ \relax
40:8 & \begin{tabularx}{0.7\textwidth}{X} 草必枯乾，花必凋謝，惟有我們神的話永遠立定。 \end{tabularx} \\ \\ \relax
40:9 & \begin{tabularx}{0.7\textwidth}{X} 報好信息的錫安哪，要登高山；報好信息的耶路撒冷啊，要極力揚聲。揚聲不要懼怕，對猶大的城鎮說：「看哪，你們的神！」 \end{tabularx} \\ \\ \relax
40:10 & \begin{tabularx}{0.7\textwidth}{X} 看哪，主耶和華必以大能臨到，他的膀臂必為他掌權；看哪，他的賞賜在他那裡，他的報應在他面前。 \end{tabularx} \\ \\ \relax
40:11 & \begin{tabularx}{0.7\textwidth}{X} 他要像牧人牧養自己的羊群，用膀臂聚集羔羊，抱在胸懷，慢慢引導那乳養小羊的。 \end{tabularx} \\ \\ \relax
40:12 & \begin{tabularx}{0.7\textwidth}{X} 誰曾用手心量諸水，用手虎口量蒼天，用升斗盛大地的塵土，用秤稱山嶺，用天平稱岡陵呢？ \end{tabularx} \\ \\ \relax
40:13 & \begin{tabularx}{0.7\textwidth}{X} 誰曾測度耶和華的靈，或作他的謀士指教他呢？ \end{tabularx} \\ \\ \relax
40:14 & \begin{tabularx}{0.7\textwidth}{X} 他與誰商議，誰教導他，以公平的路指示他，將知識傳授與他，又將通達的道指教他呢？ \end{tabularx} \\ \\ \relax
40:15 & \begin{tabularx}{0.7\textwidth}{X} 看哪，列國都像水桶裡的一滴，又如天平上的微塵；看哪，他舉起眾海島，好像舉起極微小之物。 \end{tabularx} \\ \\ \relax
40:16 & \begin{tabularx}{0.7\textwidth}{X} 黎巴嫩不夠當柴燒，其中的走獸也不夠作燔祭。 \end{tabularx} \\ \\ \relax
40:17 & \begin{tabularx}{0.7\textwidth}{X} 列國在他面前如同不存在，在他看來微不足道，只是虛空。 \end{tabularx} \\ \\ \relax
40:18 & \begin{tabularx}{0.7\textwidth}{X} 你們究竟將誰比神，用甚麼形像與他相較呢？ \end{tabularx} \\ \\ \relax
40:19 & \begin{tabularx}{0.7\textwidth}{X} 至於偶像，匠人鑄造它，銀匠用金子包裹它，又為它鑄造銀鏈。 \end{tabularx} \\ \\ \relax
40:20 & \begin{tabularx}{0.7\textwidth}{X} 沒有能力捐獻的人，就挑選不易朽壞的木頭，為自己尋找巧匠，豎立不會倒的偶像。 \end{tabularx} \\ \\ \relax
40:21 & \begin{tabularx}{0.7\textwidth}{X} 你們豈不知道嗎？豈未曾聽見嗎？難道沒有人從起頭就告訴你們嗎？自從地的根基立定，你們豈不明白嗎？ \end{tabularx} \\ \\ \relax
40:22 & \begin{tabularx}{0.7\textwidth}{X} 神坐在地的穹窿之上，地上的居民有如蚱蜢。他鋪張穹蒼如幔子，展開諸天如可住的帳棚。 \end{tabularx} \\ \\ \relax
40:23 & \begin{tabularx}{0.7\textwidth}{X} 他使君王歸於虛無，使地上的審判官成為虛空。 \end{tabularx} \\ \\ \relax
40:24 & \begin{tabularx}{0.7\textwidth}{X} 他們剛栽上，剛種好，根也剛扎在地裡，經他一吹，就都枯乾；旋風將他們吹去，像碎秸一樣。 \end{tabularx} \\ \\ \relax
40:25 & \begin{tabularx}{0.7\textwidth}{X} 那聖者說：「你們將誰與我相比，與我相等呢？」 \end{tabularx} \\ \\ \relax
40:26 & \begin{tabularx}{0.7\textwidth}{X} 你們要向上舉目，看是誰創造這萬象，按數目領出它們，一一稱其名，以他的權能和他的大能大力，使它們一個都不缺。 \end{tabularx} \\ \\ \relax
40:27 & \begin{tabularx}{0.7\textwidth}{X} 雅各啊，你為何說，以色列啊，你為何言，「我的道路向耶和華隱藏，我的冤屈神並不查問」？ \end{tabularx} \\ \\ \relax
40:28 & \begin{tabularx}{0.7\textwidth}{X} 你豈不曾知道嗎？你豈未曾聽見嗎？永在的神耶和華，創造地極的主，他不疲乏，也不困倦；他的智慧無法測度。 \end{tabularx} \\ \\ \relax
40:29 & \begin{tabularx}{0.7\textwidth}{X} 疲乏的，他賜能力；軟弱的，他加力量。 \end{tabularx} \\ \\ \relax
40:30 & \begin{tabularx}{0.7\textwidth}{X} 就是年輕人也要疲乏困倦，強壯的也必全然跌倒。 \end{tabularx} \\ \\ \relax
40:31 & \begin{tabularx}{0.7\textwidth}{X} 但那等候耶和華的必重新得力。他們必如鷹展翅上騰；他們奔跑卻不困倦，行走卻不疲乏。 \end{tabularx} \\ \\
[1ex]
\hline
\hline
\end{longtable}
講題:黎緊……咁即係未到?

經文:以賽亞書40:1-31

講員:李均熊牧師

場合:香港中文大學崇基學院禮拜堂主日崇拜

日期:2023年12月10日

起:黎緊/未到——將臨期的張力

將臨,簡單一點來說,就是「黎緊」.係教會裡面,「黎緊」講緊ge 會係聖誕節,當然亦都係指會再黎ge 耶穌基督.但我不知你們有否在等人時,等了很久,等到有些不耐煩,便致電對方,詢問對方的位置；當對方答案是「黎緊」,你會有何感受……又或者,當你在餐廳叫菜,等了很久也未起菜,若你詢問侍應時他的回答是「黎緊」,你又會如何?「黎緊」,即係未到,在等待中,等待有時令人焦慮煩躁,「黎緊」一句說話反而可能會令你無名火起.無了期的等待甚至令人失去興趣,失去希望,又原本興奮的期待變成「到jor 先至算」ge 無奈.有時可能因為等待一些影響重大的人物或事情來臨,但卻遲遲未到未發生,「黎緊」未必能緩解當中的焦慮和緊張,反而加劇負面的想法:「會唔會出jor 事?」「會唔會黎唔到?」「遲jor 件事會唔會有問題有咩後果?」……有些人,多數是小朋友或年青人吧,會為期待聖誕而感到興奮,甚至可能一早買定禮物準備交換.但會否有些人(好似我咁?),年過中年,年年難過年年過,感覺聖誕已沒有年少時慶祝的氣氛；甚至有些人(我估係讀緊或畢業於中大咁有人文氣息ge 大學ge 人掛)會感到聖誕已被消費主義的商家騎劫,將臨期的儀式成了例行公事,而對聖誕「黎緊」已經失去任何興趣.若將這種心情轉為我們等候那說好了要再來但不不知不覺說了差不多二千年要再來的耶穌基督,究竟我們又懷著甚麼心情呢?是興奮?是焦躁?是無奈?還是在現實種種失望之下,已經失去任何盼望的動力?

承:我喊叫甚麼?

以賽亞書40章,一開始正正是一種說「黎緊」但反而顯出「未到」的張力.

[3] 有人聲喊著說: 在曠野預備耶和華的路, 在沙漠地修平我們神的道. [4] 一切山窪都要填滿, 大小山岡都要削平； 高高低低的要改為平坦, 崎崎嶇嶇的必成為平原. [5] 耶和華的榮耀必然顯現； 凡有血氣的必一同看見； 因為這是耶和華親口說的.

這一段正正是說:上主要顯現了,祂「黎緊」啦,好好為祂預備啦!但接著回應的人是怎樣的態度?

[6] 有人聲說:你喊叫吧!有一個說:我喊叫甚麼呢?說:凡有血氣的盡都如草；他的美容都像野地的花. [7] 草必枯乾,花必凋殘,因為耶和華的氣吹在其上；

那是怎樣的態度呢?「我有甚麼好喊叫呢?」講D 乜ye ?有咩好講?上主「黎緊」有乜咁巴閉?關我乜事?點解我要關心?「人只不過好似草一樣,好似野地的花,過幾天這些野花野草就會枯萎；更可況上主的氣(原文可解作「沙漠吹來的熱風」)一吹,草更必枯乾,花更必凋殘!」

我們仔細再看這一段以賽亞書40章的經文,其實當中最少起碼出現三把聲音.一開始那聲音宣告說:「你們的神說:『你們要安慰,安慰我的百姓……』」這一把聲音宣稱「你們的神」,語氣帶著權威,但也於人一種置身事外的感覺.

接著第3節:「有人聲喊著說【BTW,原文是沒有人,只是有把聲音喊著說】……」其實是第二把聲音,因為他稱呼上主為「我們的神」,可見是與受話者屬同一個群體,而不是第一把聲音的宣稱上主為「你們的神」那種抽離語氣 [我們學聖經的希伯來文的,通常會花大量時間分辨語句中的第一人稱我或者第二人稱你或者第三人稱他,在這裡便派上用場了.].

第6節的那把聲音說【原文也是沒有「人」的】:「你喊叫吧」,也不知是出自第一把聲音還是第二把聲音,又或是第三把聲音；但接著回應的「有一個說:我喊叫甚麼」(和合本)就肯定是這段經文的另一把聲音了.接著是頹然的一句「凡有血氣的盡都如草；他的美容都像野地的花.草必枯乾,花必凋殘,因為耶和華的氣吹在其上.」以至第二把聲音(會說「我們的神」的那聲音)要鼓勵他:「百姓誠然是草.草必枯乾,花必凋殘,惟有我們神的話必永遠立定.」不錯,百姓的確如草一般會枯乾,但你要將焦點放在神的話,神的應許上面,那是永遠立定的.

經文餘下的部份,則回復第一把呼召人「你們的神說要安慰百姓」的聲音繼續呼召:「報好信息給錫安的啊,你要登高山；報好信息給耶路撒冷的啊,你要極力揚聲.揚聲不要懼怕,對猶大的城邑說:看哪,你們的神!主耶和華必像大能者臨到；他的膀臂必為他掌權.他的賞賜在他那裡；他的報應在他面前. 他必像牧人牧養自己的羊群,用膀臂聚集羊羔抱在懷中,慢慢引導那乳養小羊的.」

今日讓我們集中了解一下,那把頹喪,灰心,失望的聲音,為何會聽到上主「黎緊」的時候,會有這樣的反應.

轉:面對歷史轉接的無力感

這裡先處理一個confusion位【令人混亂的地方】——若你參考和合本修定本,第6節中的「有一個說」是寫作「我說:我喊叫甚麼」.這是因為不同抄本有不同人稱,可能代表古代抄寫的文士對這經文有兩方面的理解,其一,就是第一把聲音(你們的神)是上主的使者／天使,向著先知宣講,而第二把聲音就是先知本人,宣告要為將臨的上主在難行崎嶇不平的曠野預備道路,而對「你喊叫吧」(第一把聲音或是另一把上主使者的聲音)的回應者,就是當時群體中「有一個【原文是他】」代表性的聲音吧,帶著灰心失望的語氣回應著「我喊叫甚麼」.另一種理解,第二把聲音可能都仍是上主的使者,而要到回應「你喊叫把」的要求時,先知作為第一身才以失望的語調說:「我喊叫甚麼」.

怎說也好,是群體中的一位,還是先知本人,在當時面對有人高呼上主「黎緊」的心態,並不是引頸以待,又或興奮莫名,而是反高潮的:「我喊叫甚麼?!」的回應.[就好似教會攪聖誕嘉年華,牧師對會友說,嘉年華會「黎緊」,出去宣傳一下,叫周圍的街坊來參加.但你的反應是(甚至可能是牧師自己的反應):有乜好宣傳呀?講咩ye 嘉年華會呀?仲搞嘉年華會,有冇搞錯呀,21世紀都過jor 23年喇……]

是甚麼因素引至這樣的反應呢?這可能要從當時可能的讀者群體的心理狀態去理解.如果你從以賽亞書的脈絡來看,40章是進入另一個歷史時空下的宣講:以賽亞書40章前,大部份的經文信息都是宣告對以色列的審判的信息,到40章突然宣告:「要對耶路撒冷說安慰的話,又向她宣告說,她爭戰的日子已滿了；她的罪孽赦免了；她為自己的一切罪,從耶和華手中加倍受罰.」聖經學者一般都認為以賽亞書40章的生成背景,是以色列亡國被擄的中期,可能是巴比倫帝國末期,波斯帝國漸取代巴比倫成為當時的霸主的歷史時期.這是一個歷史轉折的時期,原本消滅以色列國家的強大帝國開始衰落,同時間一個更強大的勢力:瑪代波斯,正漸漸興起.當中帝君古列／居魯士(Cyrus the Great)究竟是以色列民族的解放者,還是壓迫者?對這個歷史轉折,在以賽亞書40章後的經文是十分正面的(44:28; 45:1).但已然亡國被擄超過四十多年的以色列人來說,他們等待的解救已經太久太耐,可能大部份以色列人,甚至先知【未必是以賽亞本人,而是當時在被擄群體宣講工作的一位沒有記名的先知】自己也對上主「黎緊」沒有任何預期與盼望.根據之後以賽亞的經文,更可能有一部份以色列人認為他們敬拜的上主根本就是被巴比倫神明打敗了,甚至轉而去拜巴比倫的神明；又或因為他們以為令色列人亡國的罪孽太大,以至這個神已經放棄了祂的子民,不再愛顧這些曾是祂子民的以色列人.沒有了期盼,甚至對上主失望,又或對自己民族失望,都造成了當時人聽到有宣告說上主「黎緊」而頹然和提不起勁,「黎到至算啦」!

以色列人面對歷史轉折ge 無力感,使佢地對上主「黎緊」ge 宣告毫無興趣.今日處身種種不安與無奈的我們,我們對於耶穌基督將臨的應許會否也失去興趣,還是仍然熱切期待?

合:turn our frustration into curiosity

今年是馬可年,福音經課都會選用馬可福音.今天就是馬可福音的起頭,亦是宣告「神的兒子,耶穌基督福音的起頭」.馬可福音引用的就是以賽亞書40章的經文,指出施洗約翰就是那位「預備主的道,修直他的路」的那一位.「照這話,約翰來了,在曠野施洗,傳悔改的洗禮.」今日是將臨期第二主日,福音經課都會將焦點放在施洗約翰的身上.因為他就是基督的先鋒,提醒那不再期待上主再臨的世代,要預備迎見那用聖靈給我們施洗的那一位.

耶穌基督已經黎jor ,已經證明上主和先知的說話是真實會發生的.所以今天讓我們將心裡的疑惑轉為積極的預備.「預備主的道,修直他的路」古有施洗者約翰,今日在等候上主再臨的將臨期,我們又應如何預備呢?

若果你面對今日生活的種種挑戰和困難,心情仍像被擄的以色列人一樣,感到無奈沮喪又徬徨,內心被負面的焦慮和抑鬱佔據, 讓我們聽聽呼召我們去安慰百姓的那聲音的鼓勵:

[9] 報好信息給錫安的啊, 你要登高山； 報好信息給耶路撒冷的啊, 你要極力揚聲. 揚聲不要懼怕, 對猶大的城邑說: 看哪,你們的神! [10] 主耶和華必像大能者臨到； 他的膀臂必為他掌權. 他的賞賜在他那裡； 他的報應在他面前. [11] 他必像牧人牧養自己的羊群, 用膀臂聚集羊羔抱在懷中, 慢慢引導那乳養小羊的.

不知近日當大家打開電視新聞,每日去看到的消息,是否對你來說都可能是壞消息；又或者你感覺自己只係一個小信徒,沒有發言權,沒有有能力去改變任何事.但你又會否同時看到上主的話語係立定ge ,從古時所預備的真確如實發生了.上帝唔係「黎緊」,而係「黎jor 」,而且會「再黎」,上帝再黎到的時候就要不再是一個軟弱的嬰兒,上帝再黎的時候就會為我地爭戰,得著勝利,同時間又將我地好似小羊咁懷抱係佢ge 心懷,安慰我們!

在友人的面書看到,剛過去的星期五(2023年12月8日),台灣中央研宄院社會學研究所的「當香港研究走向全球」學術研討會中,李靜君,一名著名的旅美香港社會學家(現為加利福尼亞大學洛杉磯分校社會學教授),說了這麼一句話:當前,研究香港一定會遇到不少frustrations,但要「turn our frustrations into intellectual curiosity」.或者讓我轉換一下 講法:當前面對香港生活氣氛一定會遇到不少沮喪,但要將沮喪轉為信仰的期盼.讓我們在這個將臨期,這個等待慶祝耶穌 曾經來過世上的聖誕的日子,成為等待,尋找和見證上主奇妙的作為的機會.



\newpage

\section{猜猜我是誰}
\label{sec:676}
\textbf{姚志豪牧師 - 猜猜我是誰?}
\newline
\newline
\hyperref[sec:675]{< < < PREV SERMON < < <}
~
\hyperref[sec:toc]{[返目錄]}
~
\hyperref[sec:677]{> > > NEXT SERMON > > >}


講題:猜猜我是誰?

講員:姚志豪牧師

場合:香港中文大學崇基學院禮拜堂主日崇拜

日期:2023年12月17日

女騙徒: 「婆婆,把首飾給我,讓我為你一家人祈福吧!」

婆婆: 「我們一家人都很幸福.」

「猜猜我是誰」係一種始於2000年代尾,活躍於中國大陸,香港,台灣,以至海外華人社區ge 電話詐騙行為.呢種騙案中,騙徒會隨機打一堆電話號碼,叫接聽電話ge 人估下係邊個打黎,如果對方講jor 某個名之後,就會冒認成其人,吹水吹一輪之後後收線.隔一段日子後,騙徒又會致電事主,然後冒認身份表示自己遇上jor 意外,叫對方匯款去騙徒的銀行帳戶.單係2010年亞運會期間,廣州當局已經截獲62宗類似的電話騙案.除jor 華人地區之外,日本亦有類似詐騙手段,被稱為「我啦我啦詐騙」.

明報(2023年6月16日星期五):警方偵破「猜猜我是誰」詐騙集團,拘捕5名男子,據了解包括負責向事主收款的港大經濟系學生,他未有收取報酬.警方估計集團共牽涉16宗案件,最大一筆騙款為20萬,合共涉及110萬.「猜猜我是誰」我相信即使我們沒有被騙,也會經驗過.

不過,「猜猜我是誰?」在我小時候,其實是一件既溫情,又叫人驚喜的經驗.當時用的不是智能電話,是固網的有線電話,有朋友/親戚打來,不是現在那樣有來電顯示,結果,有時一家人一齊估對方是誰,估中的興奮,叫人有一種期待.

【差派被質疑】

福音的經課,提到施洗約翰的到來,當時宗教領袖,顯明就是估下佢是何方神聖.經文交待施洗約翰是上主差派來要為那光作見證,但講到明,他不是那光,因此,當時的宗教領袖(耶路撒冷的猶太人),就差派人去見JN要問佢,即係邊個?不過你吾好以為他們想估中,反而吾想估中的,點解?(1:24-25).

四本福音書從不同的角度記述了這件事(施洗約翰)的出現,路加福音(3:1-2)更補充了更多客觀的資料:凱撒提庇留在位第十五年(於主後十四年至卅七年正式任羅馬皇帝,他在位第十五年,相當於主後廿六年.這是施洗約翰開始傳道的年分,此時約翰大約是三十歲.)//本丟‧彼拉多作猶太總督(主耶穌是在他任內被審問並判釘十字架)//希律作加利利分封的王(『希律』指希律安提帕(Herod Antipas) ,分封王---是指羅馬帝國內較小籓屬國的統治者.從主前四年至主後卅九年,長達四十三年.殺害施洗約翰的就是他.)//他兄弟腓力作以土利亞和特拉可尼地區分封的王(『腓力』(Herod Philip)是大希律的另一個兒子；他所分領的國土在加利利的東北面,惟境內猶太人不多.),呂撒聶作亞比利尼分封的王(『呂撒聶』(Lysanias)的領土在黑門山以北,黎巴嫩山脈以東,與猶太人幾無關係.)//亞那(岳父)和該亞法(女婿)作大祭司(兩位都是羅馬政府所公認的大祭司,猶太的大祭司既是宗教領袖,又是與羅馬政府交涉事務的政治領袖.當時的猶太公會相當於猶太族的最高法庭,大祭司便是公會的主持人,權勢很大.當年耶穌被無辜定罪處死,就是由該亞法主謀煽動猶太人,迫使羅馬巡撫下令執行的.當施洗約翰開始傳道的時候,正值該亞法作大祭司,但因亞那仍舊舉足輕重,故路加提到兩個人的名字.)那時,撒迦利亞的兒子約翰在曠野裡,上帝的話臨到他.即羅馬人統治的年代.

路加補充了這個重要時間性的資料,正正就道出,當時約翰在一個怎樣管治下的處境.是一個被監管且高壓下.即使宗教領袖,都不一定會站宗教人物「約翰」的一邊,更加質疑他的身份,而呢班權力人物與約翰及耶穌的死甚有關係.

這班宗教領袖是咩人?點解可有咁大權力?他們是來自猶太人議會,即猶太人政治及宗教最高領導核心(中央政府權力核心),差派人查問JN的身份,點解要咁問JN?就是JN為人洗禮.對於猶太人最高權力,就是問緊你有咩權力咁做?換言之,即一般人不可以這樣做,再講得白D,你吾係我們權力核心的人,無我地批准,是不可這樣做.事實上,對於猶太人而言,只有外邦人歸入猶太教是需要接受潔淨儀式(不等於洗禮)；另某些宗派都有施行潔淨禮,是他們(宗教有權力人士)進行的,而且不止一次.一方面,JN幫耶穌洗禮,不是幫外邦人皈信猶太教(耶穌就是猶太人)施洗,更加不只簡單的潔淨儀式,而是悔改的洗禮,對猶太人而言,他們根本不需要悔罪洗禮(他們已分別為聖).不難想像的是,約翰超越緊佢地的權力,挑戰緊他們的特權,(悔罪)甚至乎覺得影射他們.當擁有權力的人,感到一個無權者的威脅,他們就是腎上線上升,甚至乎用盡方法把他壓下去.試問:他們仲點會想估中約翰所引導出來的,就是那位彌賽亞(拯救者)呢?估中左,豈不是會動搖/威脅自己的地位?難怪(太2:4)當希律聽到猶太人之王降生:「希律王聽見了,就心裡不安；耶路撒冷全城的人也都不安.」2:16,而且,還打算差人將伯利恆城裡和四境所有的男孩,凡兩歲以內的,都殺盡了.

對於約翰被上主差派作見證,傳好消息,應該是喜樂(期待已久),就好似將臨期第三主日點起的「粉紅色」蠟燭,就是「喜樂之燭」.但對某些人來是卻是壞消息,不單身份被質,更認為是煽動民眾.我想:這就是信仰的本質:不是要衝著權力而來,而是在人中間彰顯緊人性/神聖,「光照在黑暗裡,黑暗卻沒有勝過光.」(約1:5)表明信仰的內涵.因此,1:22,果班被差派來祭司及利未人問約翰(估估下估下,應該又驚驚地):「請告訴我們,你到底是誰,好讓我們回覆派遣我們來的人.你自己說,你是什麼人?」(好交差)一個是上主差派,VS而一班是有權力的人士差派來的,形成了JC出場前一種的張力.

事實上,福音書的約翰,只是延續舊約所讀以賽亞先知講述為關乎全人類的好消息作好準備.

【差派是報好消息】

上主差派先知以賽亞,就是向受壓迫的人帶來希望,向貧窮的人傳佳音.他差遣我醫治傷心的人；差遣我去安慰悲傷的人；賽61:1-----被擄的,得釋放；被囚的,得自由.差遣我宣告:上主拯救他子民的恩年(ACCEPTABLE)；好明顯全部都是佳音/好消息/GOODNEWS.

〈GOODNEWS〉

值得我們深思是:(福音GOOD NEWS)這個字的意思?對貧窮人甚麼是的佳音?當然不只是留於口說「耶穌愛你」,「愛」其實是好抽象的,有沒有具體的事其實很難明白?我成日覺得係學校做宗教科老師好難做,這頭說:主耶穌好愛你,轉頭問你重咩吾交功課.耶穌愛你,但我要罰你.吾知點算!GOOD NEWS好消息:是以JC釘十架對我們的拯救為核心,怎樣將這份無條件的犧牲轉化日常具體的行動?每個主日,崇拜結束前差遣我們返去生活崗位做咩呢?就是報好消息GOODNEWS.我想:對於一個交吾切功課的人,好消息又是甚麼?對於一個寒冷的人,好消息又是甚麼?對於一個沒得食的人,好消息又是甚麼?對於一個受排擠的人,好消息又是甚麼?對於一個坐冤獄的人?對於一個沒被公正對待的人?對於一個孤單的人,好消息又是甚麼?如果上主是愛,但在你生命無做過少任何野?大概你會覺得不是太真實,如果JC不用我們咁具體地去愛,他自己就不需要餵飽人,不需要醫好人,不需要陪伴人,不需要為你犧牲,GOODNEWS豈只留於形而上的愛?更加不需要生在充滿張力的人世間.:EG.教會講道的一杯水.就是對你接觸的人他的需要而回應,即對人產生感覺,除去冷漠,從而讓其他人明白基督徒不是得個講字.

〈恩年〉

因此,先知以賽亞又提到上主拯救他子民的恩年(61:2):學者認為正是JC路4:19「並宣告主拯救他子民的恩年(禧).」而這個恩年(禧),利25:8-10「你們要計算七個七年,一共四十九年.25:9 那一年的七月十日,就是贖罪日,你們要派人到全國各地吹號.五十年分別出來,向全體居民宣佈這一年為自由年,為禧年.在這一年,所有賣了的產業要歸還給原主或原主的後代；從前被販賣為奴的要得釋放,回到自己的家.」

可能你會問,我幾時有囚禁過任何人需要釋放呢?今天,在我們崗位上(可控的範圍內),對人的態度上,言語上的標籤(性小眾,低端人口),限制,是否已囚禁了多年呢?有時,我們對自己也標籤左自己,有時精神囚禁不會遜色於肉身的囚禁,甚至有過之而無不及.上主差遣我地出去,就是要釋放人,是上主所悅納的.正如「賽49:9 我要對被俘的說:你們自由了!對黑暗中的人說:你們重見光明吧!」釋放人就是放過自己,釋放自己也是釋放人.EG.二戰時期,被囚禁的德國潘霍華牧師,他怎樣理解將臨期?監獄的牢房,人們在其中等待,期盼,而牢房通往自由之門必須從外面開啟,這是個頗能貼切的形容將臨期的意象.

先知的記述,就是表達上主期望將不好轉為好,因此以賽亞書61章3提到:為錫安悲哀的人,賜華冠代替灰燼,喜樂的油代替悲哀,讚美為衣代替憂傷的靈；當講完這些轉變後,就話要好似上主親手種植的樹,稱他們為「公義樹」,是耶和華所栽植的,為要彰顯他的榮耀.8因為我－耶和華喜愛公平,恨惡搶奪與惡行( [ 61.8] 「恨惡搶奪與惡行」:有古卷是「恨惡搶奪又獻燔祭」.)；我要憑誠實施行報償,與我的百姓立永約.10我因耶和華大大歡喜,我的心因上帝喜樂；因他以拯救為衣給我穿上,以公義為外袍給我披上,好像新郎戴上華冠,又如新娘佩戴首飾.11地怎樣使芽長出,園子怎樣使所栽種的生長,主耶和華也必照樣使公義和讚美在萬國中發出.

短短8節經文,出現了4次公義/公道,而且61:-2為被擄得釋放,被綑綁的自由,顯明也是一個公義的意象.

【差派:種植公義的種子】

公義「樹」,有咩特性,所謂十年樹木,即不是一時間會出現,是由一粒種子開始,並且於生長的過程會遇左遭遇到很多困難,甚麼乎有一天會因強風而倒下,又有機會比人鋸,然而當樹根找得越深,就是找得越緊.樹最初是種子的階段是見不到的,慢慢長大不單觸目,更替人遮風擋雨.公義是需要培養的,公義就是當眾人仍沉默/沒動靜之時,你仍願意表達一點聲音/做一些事情,即使很微弱,即使被壓下去,那棵公義樹都不易跌倒,只因是上主親手種的.今天,我們被栽種的是怎麼樣的種子?或是我們會選擇甚麼種子栽種我們的心中?環顧四周,不公不義之事往往就是近在咫尺,工作?家庭?身處之地?甚而我們信仰對事情理解.

持續兩個多月(8/10/2023)的以(哈以)巴衝突:事件一出,不少人會因為「以色列」的名字而盲撐,原來當我地細心留意下,巴勒斯坦長期遭到不公平的對待.二戰後,德國殺害猶太人後,以色列復國呼聲高,對於當時未有實際邊界的以色列國來說,以色列及巴勒斯坦人都生活在此地上.不論在1947年《聯合國大會181號的議決》(巴45\%,耶城共同,以55\%)或是1993年《奧斯陸協議》(巴18\%,共同22\%,以60\%),巴勒斯坦(小)與以色列(多)所分得的地非常不公道,而且當以色列因經濟起飛,吸引不少以色人進入該地以後,迫使巴勒斯坦人離開,而且更不讓他回到居住之地,其實是長期受壓.研究巴勒斯坦人神學學者:”voices fom the margin”,使巴勒斯坦人想起:當耶和華命定以色列人進入迦南美地之時,他們同樣壓迫緊迦南人離開,只有以色人的視野去睇聖經,對被壓迫的迦南人公道嗎?可以把巴勒斯坦人放在中心非邊緣的位置去睇聖經嗎?EG.JC的出身的相片.

詩85:13,「公義要走在上主前面,為他預備道路.」因此,我們不難意會,施洗約翰為耶穌基督開路,亦即是說,當有公義的出現之後,上主隨之而來.EG.就好似,一對結婚的新人步入教堂,大家最先見到的不是這對新人,而是花仔花女,他們好像為新開路似的,隨之而來,必定是那位粉墨登場的新娘子.將臨期,再次提醒我們自己有這花仔花女角色.

總結:

施洗約翰被質疑,JC更會被質疑,因為他更顛覆當下宗教/政權的慣性.《拉比給我們生命的九堂課》「人可能只有宗教信仰而不認識神嗎?人可能認識神而不屬於任何一個組織的宗教嗎?」「上帝不在教堂或會堂裡,不在聖書中,不在虔誠的信徒的心中,乃在人際之間.p.98當我地要「猜猜佢是誰?」之時,正代表我們不一定估中,甚至乎會偏差,因為他是超越我們所能理解,所知道的,但我們儘管去猜,儘管去試,說不定,當有日揭曉之時,你中了JACKPOT.



\newpage

\section{從上帝而來的沒有任何事情是不可能的}
\label{sec:677}
\textbf{高國雄牧師 - 從上帝而來的,沒有任何事情是不可能的!}
\newline
\newline
\hyperref[sec:676]{< < < PREV SERMON < < <}
~
\hyperref[sec:toc]{[返目錄]}
~
\hyperref[sec:678]{> > > NEXT SERMON > > >}


路加福音 1:26-38
\newline
\begin{longtable}{cl}
\hline
\hline
章節 & 經文 (和合本修訂版)\\
\hline
1:26 & \begin{tabularx}{0.7\textwidth}{X} 到了第六個月，天使加百列奉神的差遣往加利利的一座城去，這城名叫拿撒勒， \end{tabularx} \\ \\ \relax
1:27 & \begin{tabularx}{0.7\textwidth}{X} 到一個童女那裡，她已經許配大衛家的一個人，名叫約瑟；童女的名字叫馬利亞。 \end{tabularx} \\ \\ \relax
1:28 & \begin{tabularx}{0.7\textwidth}{X} 天使進去，對她說：「蒙大恩的女子，你好，主和你同在！」 \end{tabularx} \\ \\ \relax
1:29 & \begin{tabularx}{0.7\textwidth}{X} 馬利亞因這話就很驚慌，又反覆思考這樣問候是甚麼意思。 \end{tabularx} \\ \\ \relax
1:30 & \begin{tabularx}{0.7\textwidth}{X} 天使對她說：「馬利亞，不要怕，你在神面前已經蒙恩了。 \end{tabularx} \\ \\ \relax
1:31 & \begin{tabularx}{0.7\textwidth}{X} 你要懷孕生子，要給他起名叫耶穌。 \end{tabularx} \\ \\ \relax
1:32 & \begin{tabularx}{0.7\textwidth}{X} 他將要為大，稱為至高者的兒子；主神要把他祖先大衛的王位給他。 \end{tabularx} \\ \\ \relax
1:33 & \begin{tabularx}{0.7\textwidth}{X} 他要作雅各家的王，直到永遠；他的國沒有窮盡。」 \end{tabularx} \\ \\ \relax
1:34 & \begin{tabularx}{0.7\textwidth}{X} 馬利亞對天使說：「我沒有出嫁，怎麼會有這事呢？」 \end{tabularx} \\ \\ \relax
1:35 & \begin{tabularx}{0.7\textwidth}{X} 天使回答她說：「聖靈要臨到你身上；至高者的能力要庇蔭你，因此，那要出生的聖者要稱為神的兒子。 \end{tabularx} \\ \\ \relax
1:36 & \begin{tabularx}{0.7\textwidth}{X} 況且，你的親戚伊利莎白，就是那素來稱為不生育的，在年老的時候也懷了男胎，現在懷孕六個月了。 \end{tabularx} \\ \\ \relax
1:37 & \begin{tabularx}{0.7\textwidth}{X} 因為，出於神的話，沒有一句不帶能力的。」 \end{tabularx} \\ \\ \relax
1:38 & \begin{tabularx}{0.7\textwidth}{X} 馬利亞說：「我是主的使女，願意照你的話實現在我身上。」於是天使離開她去了。 \end{tabularx} \\ \\
[1ex]
\hline
\hline
\end{longtable}
講題:從上帝而來的,沒有任何事情是不可能的!

經文:路加福音1:26-38

講員:高國雄牧師

場合:香港中文大學崇基學院禮拜堂主日崇拜

日期:2023年12月24日

路加福音1:26-38

26到了第六個月,天使加百列奉上帝的差遣往加利利的一座城去(這城名叫拿撒勒), 27到一個童女那裡,是已經許配大衛家的一個人,名叫約瑟.童女的名字叫馬利亞； 28天使進去,對她說:「蒙大恩的女子,我問你安,主和你同在了!」 29馬利亞因這話就很驚慌,又反覆思想這樣問安是甚麼意思. 30天使對她說:「馬利亞,不要怕!你在上帝面前已經蒙恩了. 31你要懷孕生子,可以給他起名叫耶穌. 32他要為大,稱為至高者的兒子；主上帝要把他祖大衛的位給他. 33他要作雅各家的王,直到永遠；他的國也沒有窮盡.」 34馬利亞對天使說:「我沒有出嫁,怎麼有這事呢?」 35天使回答說:「聖靈要臨到你身上,至高者的能力要蔭庇你,因此所要生的聖者必稱為上帝的兒子 . 36況且你的親戚伊利莎白,在年老的時候也懷了男胎,就是那素來稱為不生育的,現在有孕六個月了. 37因為,出於上帝的話,沒有一句不帶能力的.」 38馬利亞說:「我是主的使女,情願照你的話成就在我身上.」天使就離開她去了.

今天是將臨期第四主日,也是平安夜的日子.在聖誕這節期當中,我們可能有很多慶祝的活動,與家人相聚的時間.

當我們回到聖經當中,有關耶穌基督出生的故事.

我們好像很熟悉,但是當中有很多重要的事情我們會遺忘,甚至我們刻意忘記.

基督的誕降是上帝極大的奧秘.

在今天經課中:羅馬書提醒我們.

這奧秘是要顯明出來.

而耶穌基督的,但正正就是奧秘顯明彰顯的重要點.

這道成了肉身,道就是耶穌基督,就是上帝.

路加福音記載了耶穌出生前一個很重要的片段.

這裡提到當時馬利亞已經許配了給約瑟.

聖經描述天使對馬利亞說:蒙大恩的女子,我問你安,主和你同在了.

瑪利亞的反應是很驚慌.

又反覆思想,這是什麼意思?

可能我們不太明白到底馬利亞為什麼驚惶?

這很明顯,天使對馬利亞的問安是非常態的,特別是「蒙大恩的女子」,這描述是非常特別的.

從原文看,這個蒙恩者的描述是一個對尊貴者,被寵愛的意思.

所以這裡天使就是向馬利亞表述,他特別特別特別被上帝寵愛.

弟兄姊妹你期望被上帝寵愛嗎?

我們知道上帝愛我們每一個人.

但當「特別被寵愛」在我們身上發生象徵著什麼呢?

馬利亞被上帝寵愛,上帝揀選了瑪利亞.

因此在他身上行了獨特的事,奇妙的事.

更是那奧秘的事.

天使的回應:「不要怕.」

其實懼怕是現代人常常出現的情感與對事物的反應.

在耶穌出生的故事當中,無論在馬太福音或是路加福音中也充滿了懼怕的情感.

當天使天軍向牧羊人報喜訊,在中大光出現.

牧羊的人就很懼怕.

當約瑟發現仍未迎娶的未婚妻馬利亞有了身孕.

他很懼怕.

當希律王從東方來的智者口中知道將會有新生的王.

他很懼怕.

當希律王的懼怕導致很多孩童帶來殺身之禍.

全城也懼怕.

懼怕往往是我們現代人生活好像不能逃脫的.

我們也常活在當中.

所以當馬利亞聽到天使代表著上帝表示對他特別特別特別寵愛.

他當然非常懼怕.

所以天使對馬利亞的安慰:就是由「不要怕」開始.

因為他將要收到一個十分震撼,會令他萬分懼怕的消息…

馬利亞要懷孕生子.

但當時她仍未正式出嫁給約瑟,如何會懷孕呢?

相信這是歷世歷代最難明白的奧秘之一.

今天我們在教會,神學上習慣用「聖靈感孕」來描述這事件,我們好像就明白了,解開了這疑惑,但其實一個人怎能直接從聖靈中懷孕呢?

請再仔細閱讀經文:天使在此所描述的

30天使對她說:「馬利亞,不要怕!你在上帝面前已經蒙恩了. 31你要懷孕生子,可以給他起名叫耶穌. 32他要為大,稱為至高者的兒子；主上帝要把他祖大衛的位給他. 33他要作雅各家的王,直到永遠；他的國也沒有窮盡.」

馬利亞的回應:

34馬利亞對天使說:「我沒有出嫁,怎麼有這事呢?」

很明顯馬利亞清楚認知將要發生的是什麼事.

但不能接受,因為覺得不可能…

35天使回答說:「聖靈要臨到你身上,至高者的能力要蔭庇你,因此所要生的聖者必稱為上帝的兒子 .

聖靈臨在是什麼意思/意義呢?

我們作為基督徒,渴望聖靈的臨在.

下主日將在本堂進行的洗禮中,透過聖禮,聖靈要臨在每一位受禮者當中.

聖靈在是我們所期盼的.

或許我們一向認為聖靈降臨必定是歡喜快樂,開心雀躍,興奮莫名的事情.

但在馬利亞身上,聖靈要降臨卻在憂慮,懼怕驚惶的時候,因為實在很難接受.

聖靈的臨在,在路加福音中一個非常重要的課題及核心.

路加福音大概有11段經文中提到聖靈(15次).

當中提到聖靈臨在/充滿的對象與場景包括:施洗者約翰,其父母撒迦利亞,伊利沙白,耶穌受割禮後上耶路撒冷時遇見的西面,今日經課的主角馬利亞,當然包括我們很熟悉的:耶穌接受約翰洗禮時聖靈如鴿子般降在耶穌身上.

論到聖靈,路加福音10:21

正當那時,耶穌被聖靈感動而歡喜快樂,說:「父啊,天地的主,我感謝你!因為你把這些事向聰明智慧的人隱藏起來,而向嬰孩啟示出來.父啊,是的,因為你的美意本是如此. 22一切都是我父交給我的.除了父,沒有人知道子是誰；除了子和子所願意啟示的人,沒有人知道父是誰.」

聖靈會令人歡喜快樂的原因就是:

上帝會透過聖靈而啟示自己,我們就能透過聖靈而認識上帝更多.

因此聖靈最重要的能力就是溝通的能力,人與上帝溝通的能力,更提高人與人溝通的能力.

11章更提到

11: 9我又告訴你們,祈求,就給你們；尋找,就找到；叩門,就給你們開門. 10因為凡祈求的,就得著；尋找的,就找到；叩門的,就給他開門. 11你們中間作父親的,誰有兒子求魚,反拿蛇當魚給他呢? 12求雞蛋,反給他蠍子呢? 13你們雖然不好,尚且知道拿好東西給兒女,何況天父,他豈不更要把聖靈賜給求他的人嗎?

聖靈是上帝給人最好的禮物!

也是我們人生應該追求的最重要/唯一目標.

天使對馬利亞以另一個見證作為印證:就是有關他親戚伊利莎白在年老的時候也懷了男胎.

從而引證:出於上帝的話,沒有一句不能不帶著能力的.

理性的信仰思考,神學研究,透過邏輯思維推論等能夠為我們帶來很多衝擊,有時開拓我們對上帝的理解.

但同時也有機會狹窄了我們對上帝的認知,甚至因此使沒有辦法透徹地進入上帝當中.

回想一下,我們的屬靈生命,回想一下我們的成長.

聖靈的臨在,不是我們用理性思考就能夠完成理解的.

當我們經歷時,更多的情況是因懼怕而沒有辦法接受.

有時我們期盼上帝以我們渴求的內容回應我們.

更多時候,我們根本沒有完全擺上信心,因為我們只期盼上帝完成我所計劃的…

就例如在疫情時,會眾能夠保持教會人數,奉獻,能夠繼續進行崇拜,已經是難能可貴的.

在我們的心底裡,是否其實是對上帝沒有太大的期望呢?

又或我們根本沒有認識到上帝的計劃及心意呢?

「在人不能,在上帝凡事都能」這一金句,我們很熟悉.但到底我們是否真的相信呢?

路1:38

馬利亞說:「我是主的使女,情願照你的話成就在我身上.」

降服與接受就是對聖靈臨在最好的回應!

因為從上帝而來的,沒有任何事情是不可能的!



\newpage

\section{應許與期望}
\label{sec:678}
\textbf{李均熊牧師 - 應許與期望}
\newline
\newline
\hyperref[sec:677]{< < < PREV SERMON < < <}
~
\hyperref[sec:toc]{[返目錄]}
~
\hyperref[sec:679]{> > > NEXT SERMON > > >}


馬太福音 1:18-25
\newline
\begin{longtable}{cl}
\hline
\hline
章節 & 經文 (和合本修訂版)\\
\hline
1:18 & \begin{tabularx}{0.7\textwidth}{X} 耶穌基督降生的事記在下面：他母親馬利亞已經許配給約瑟，還沒有迎娶，馬利亞就從聖靈懷了孕。 \end{tabularx} \\ \\ \relax
1:19 & \begin{tabularx}{0.7\textwidth}{X} 她丈夫約瑟是個義人，不願意當眾羞辱她，想要暗地裡把她休了。 \end{tabularx} \\ \\ \relax
1:20 & \begin{tabularx}{0.7\textwidth}{X} 正考慮這些事的時候，忽然主的使者在約瑟夢中向他顯現，說：「大衛的子孫約瑟，不要怕，把你的妻子馬利亞娶過來，因她所懷的孕是從聖靈來的。 \end{tabularx} \\ \\ \relax
1:21 & \begin{tabularx}{0.7\textwidth}{X} 她將要生一個兒子，你要給他起名叫耶穌，因他要將自己的百姓從罪惡裡救出來。」 \end{tabularx} \\ \\ \relax
1:22 & \begin{tabularx}{0.7\textwidth}{X} 這整件事的發生，是要應驗主藉先知所說的話： \end{tabularx} \\ \\ \relax
1:23 & \begin{tabularx}{0.7\textwidth}{X} 「必有童女懷孕生子；人要稱他的名為以馬內利。」（以馬內利翻出來就是「神與我們同在」。） \end{tabularx} \\ \\ \relax
1:24 & \begin{tabularx}{0.7\textwidth}{X} 約瑟醒來，就遵照主的使者的吩咐把妻子娶過來； \end{tabularx} \\ \\ \relax
1:25 & \begin{tabularx}{0.7\textwidth}{X} 但是沒有和她同房，直到她生了兒子，就給他起名叫耶穌。 \end{tabularx} \\ \\
[1ex]
\hline
\hline
\end{longtable}
以賽亞書 7:1-17
\newline
\begin{longtable}{cl}
\hline
\hline
章節 & 經文 (和合本修訂版)\\
\hline
7:1 & \begin{tabularx}{0.7\textwidth}{X} 烏西雅的孫子，約坦的兒子，猶大王亞哈斯在位的時候，亞蘭王利汛和利瑪利的兒子以色列王比加上來攻打耶路撒冷，卻不能攻取。 \end{tabularx} \\ \\ \relax
7:2 & \begin{tabularx}{0.7\textwidth}{X} 有人告訴大衛家說：「亞蘭與以法蓮已經結盟。」王的心和百姓的心就都顫動，好像林中的樹被風吹動一樣。 \end{tabularx} \\ \\ \relax
7:3 & \begin{tabularx}{0.7\textwidth}{X} 耶和華對以賽亞說：「你和你的兒子施亞‧雅述要出去，到上池的水溝盡頭，往漂布地的大路上，迎見亞哈斯， \end{tabularx} \\ \\ \relax
7:4 & \begin{tabularx}{0.7\textwidth}{X} 對他說：『你要謹慎，要鎮定，不要害怕，不要因利汛和亞蘭，以及利瑪利的兒子這兩個冒煙火把的頭所發的烈怒而心裡膽怯。 \end{tabularx} \\ \\ \relax
7:5 & \begin{tabularx}{0.7\textwidth}{X} 因為亞蘭、以法蓮，和利瑪利的兒子設惡謀要害你，說： \end{tabularx} \\ \\ \relax
7:6 & \begin{tabularx}{0.7\textwidth}{X} 我們要上去攻擊猶大，擾亂它，攻破它來歸我們，在其中立他比勒的兒子為王。 \end{tabularx} \\ \\ \relax
7:7 & \begin{tabularx}{0.7\textwidth}{X} 主耶和華如此說：這事必站立不住，也不得成就。 \end{tabularx} \\ \\ \relax
7:8 & \begin{tabularx}{0.7\textwidth}{X} 因為亞蘭的首都是大馬士革，大馬士革的領袖是利汛；六十五年之內，以法蓮必然國破族亡， \end{tabularx} \\ \\ \relax
7:9 & \begin{tabularx}{0.7\textwidth}{X} 以法蓮的首都是撒瑪利亞；撒瑪利亞的領袖是利瑪利的兒子。你們若是不信，必站立不穩。』」 \end{tabularx} \\ \\ \relax
7:10 & \begin{tabularx}{0.7\textwidth}{X} 耶和華又吩咐亞哈斯： \end{tabularx} \\ \\ \relax
7:11 & \begin{tabularx}{0.7\textwidth}{X} 「你向耶和華－你的神求一個預兆：在陰間的深淵，或往上的高處。」 \end{tabularx} \\ \\ \relax
7:12 & \begin{tabularx}{0.7\textwidth}{X} 但亞哈斯說：「我不求；我不試探耶和華。」 \end{tabularx} \\ \\ \relax
7:13 & \begin{tabularx}{0.7\textwidth}{X} 以賽亞說：「聽啊，大衛家！你們使人厭煩豈算小事，還要使我的神厭煩嗎？ \end{tabularx} \\ \\ \relax
7:14 & \begin{tabularx}{0.7\textwidth}{X} 因此，主自己要給你們一個預兆，看哪，必有童女懷孕生子，給他起名叫以馬內利。 \end{tabularx} \\ \\ \relax
7:15 & \begin{tabularx}{0.7\textwidth}{X} 到他曉得棄惡擇善的時候，他必吃乳酪與蜂蜜。 \end{tabularx} \\ \\ \relax
7:16 & \begin{tabularx}{0.7\textwidth}{X} 因為在這孩子還不曉得棄惡擇善之先，你所憎惡的那兩個王的土地必被撇棄。 \end{tabularx} \\ \\ \relax
7:17 & \begin{tabularx}{0.7\textwidth}{X} 耶和華必使亞述王臨到你和你的百姓，並你的父家，自從以法蓮脫離猶大的時候，未曾有過這樣的日子。 \end{tabularx} \\ \\
[1ex]
\hline
\hline
\end{longtable}
講題:應許與期望 Promise and Expectation

經文:馬太福音1:18-25；以賽亞書7:1-17

講員:李均熊牧師

場合:香港中文大學崇基學院禮拜堂平安夜燭光崇拜

日期:2023年12月24日

引言:應許與期望的複雜關係

今天平安夜的宣講,我選擇了「應許與期望」作為講題.我自己一路預備的時候,發覺這個題目其實是十分複雜的.有應許自然有期望,但有時候,有應許你卻沒有期望:例如區議會選舉,有很多應許,但期望卻十分低.但有時候,即使無應許,你仍有期望.情況如看電影,有些時不是因為套電影好看不好看,而是看你對那套電影期望有多高.如果因為電影宣傳或別人的推薦,都應許電影會十分好看,你帶著好高期望去觀看,你就會發覺不外如是,可能就覺得不好看；反之,完全無期望,例如陪伴家人或朋友去看一套電影,但卻因為無期望反而你發覺好好看.

對於過聖誕節,你又有沒有期望呢?人越來越大,可能就越來越沒有期望了.對於聖誕節,可能又是那些party,又是那些火雞,又是那些經文,又是那樣的報佳音.但小朋友可能仍有一定的期望.如果再加上父母或家人「應許」了買甚麼禮物,小朋友的期望可能更高. 不過若你應許一些很特別的事情,但最後卻不能實現,帶來失望,那麼以後家人對你的應許,就會降低期望,甚至沒有期望.所以可能今天你很少再應許家人或別人什麼事情,又或很少會應許小孩子送他任何東西,因為我們常常害怕自己做不到,買不到,所以我就不要讓人失望了. 遂漸地不單沒有期望, 更加沒有了應許.

可見應許與期望的互動關係並不是簡單.

今日讓我們回看聖經對耶穌出生的記載,特別是馬太福音的記載,看看當時人是怎樣的期望,上主的應許又如何回應.

馬太的藍色聖誕

大家常聽說「白色聖誕」,想像一個漫天細雲下的聖誕節會十分浪漫.但大家有沒有聽過「藍色聖誕」(blue Christmas)?有些西方教會會在一年黑夜最長的一天,亦即冬至日,會舉辦聚會去向失去所愛或在哀中的人致敬.也有些研究指出節日更容易令人感到抑鬱,焦慮,疲倦,哀傷,甚至有些心理學家叫這些為Christmas blue(聖誕愁思).今年在加沙或者烏黑蘭當中的戰爭,亦令到今年聖誕添上一抹藍色.

其實聖經記載耶穌的出生,並不是完全歡歡喜喜的.我們了解的聖誕故事,有時從教會「聖景」擺設(nativity scene),會以為東方博士獻禮物,天使天軍和牧羊人朝見臥在馬糟的嬰孩耶穌是同時出現,其實這些擺設只是將路加的記載和馬太的記載結合了.但若我分開兩卷福音書來仔細查看耶穌出生的記載,我們就或者會留意到,只有路加福音會認為耶穌出生是一個要普天同慶的大喜訊息(路1:14, 28, 42, 47, 48, 57; 二10, 14….),但若你轉到馬太福音的記載,你會見到的是約瑟的忐忑不安快,甚至是合城的不安,逃亡,希律的凶殺和母親的悲哭(太2:3, 13, 18).

為何馬太福音和路加福音記載耶穌出生的調子會這樣不同?這可能是兩者的寫作對像不同所致.路加福音的寫的對像,極有可能是一些操希臘語的外邦信徒,或更可能是一群希臘化的猶太人.用回我們今日的主題,基督,亦即猶太人聖經所預言的彌賽亞,對這些人來說,其實沒有甚麼期望,對他們來說基督降生也並不是實現了甚麼的應許.所以,路加福音聚照耶穌作為全人類的救主,係為你而生,能拯救你幫助你,故此是一大喜的信息.無期望反而更覺驚喜!

但馬太福音的成書對像又是誰呢?現在聖經的研究認為馬太的第一代讀者很大可能是經歷過公元70 年羅馬軍隊鏟平了第二聖殿的猶太信徒,他們和約瑟和馬利亞一家,與耶路撒冷,又或伯利恆的母親們一樣,都經歷不安,或哀傷,或逃亡的信徒群體.他們本來可能對彌賽亞的應許十分期待,但所經歷的卻是令他們懷疑究竟這個應許有多真實的群體.高期望反而更失望,原初滿滿盼望變成不再期望,甚至質疑經書的應許是否會實現.所以,馬太福音記載基督的降生,正是發生在不安,殺戮和逃難的當下,正如他們所經歷的處境一般.亦正是在這艱難時刻,耶穌基督就照著先知的預言在伯利恆誕生了.

天使的聲音(太1:18-25)

樣我們仔細看看,究竟馬太福音的記載如何對應(address)當時歷盡苦劫的猶太信徒群體.

當你細讀今日的選讀的馬太經文,你會發覺整段記載只有天使的對白,而我們以為的主角約瑟和馬利亞,其實經文沒有記載過他們任何一句的說話.那麼天使那句說話是甚麼呢?

「大衛的子孫約瑟,不要怕,把你的妻子馬利亞娶過來,因她所懷的孕是從聖靈來的.21她將要生一個兒子,你要給他起名叫耶穌,因他要將自己的百姓從罪惡裡救出來.」

好了,究竟天使的聲音是否停在這裡?若我們繼續讀下去:

22這一切的事成就是要應驗主藉先知所說的話,23說:必有童女懷孕生子；人要稱他的名為以馬內利.(以馬內利翻出來就是「上帝與我們同在」.)

這句說話,有多少人會覺得仍是天使的說話?(我年少時看,常以為這一句也是天使的宣告.為何不可呢?天使進一步指出預言應驗,那怕約瑟你不能不信!)

其實這句說話和之前的說是有些矛盾的,若然天使叫約瑟將要生的孩子叫作耶穌,怎麼之後又說人要稱他的名為「以馬內利」?如何兩句說話都是天使對約瑟的宣告,那約瑟可能會問:究竟嬰孩要叫什麼名字?是耶穌,還是以馬內利?那很大可能這句說話並非天使的宣告一部份.

若然22-23節不是天使的宣告,那是誰的聲音呢?這就要回看馬太福音寫作的時候,可能因為在經歷好多的事情後,他們才終於明白到,這個耶穌就係經書預言,就是應許的那一位「以馬內利」.換句話說,是在編修耶穌生平時候,這一句是後來加插入去,為要提醒讀的人,這個耶穌就是「以馬內利」,就是神與我們同在了.耶穌的就是在困難的境況之下出生,就是那刻上帝的應許便應驗了,但那是在後來回想才有這樣的意識.如果馬太福音的讀者,是公元70年後那些經歷各樣,壓迫,創傷,戰亂,流亡,活在惶恐不安中的信徒群體中,這一句後來的加插正正是在回應他們「以馬內利」是真實的,神的應許沒有失落,耶穌基督就是神與我們同在的明證.

「以馬內利」的應許(賽七1-17)

「以馬內利」的應許其實源於以賽亞對當時猶大國亞哈斯王的宣告:

10耶和華又曉諭亞哈斯說:11「你向耶和華－你的上帝求一個兆頭:或求顯在深處,或求顯在高處.」12亞哈斯說:「我不求；我不試探耶和華.」13以賽亞說:「大衛家啊,你們當聽!你們使人厭煩豈算小事,還要使我的上帝厭煩嗎?14因此,主自己要給你們一個兆頭,必有童女懷孕生子,給他起名叫以馬內利.

「以馬內利」乃對大衛王朝的(政治性)應許

當時亞蘭王(古敘利亞)利汛和以色列王比加聯合起來出兵猶大國,並非要消滅猶大,而是要推翻亞哈斯的大衛王朝(設謀要害你),另立一個傀儡君王(在其中立他比勒的兒子為王).這正是亞哈斯的關注,並非國家的安全,而是關乎自己和國家王朝的興衰榮辱.

亞蘭國聯合以色列國是屬於抗拒亞述的勢力.亞哈斯則要在臣服亞述而失去自主國權,與加入聯軍卻九死一生中作出抉擇.而顯然亞哈斯王傾向臣服保命.亞蘭與以色列的聯合陣線自然不能讓猶大國在他們後方歸向敵人亞述的一方,於是才有聯手進攻,期望推翻大衛王朝,建立一個會和他們合作的傀儡君王.故此對亞哈斯來說,保著王位為先,不斷侵吞鄰國的亞述反而成了盟友.

上主的信息指出利迅和比加的計謀必不能成功,因為就算猶大不向亞述臣服,亞蘭國與以色列國的滅亡也是指日可待,故不用擔心.只要亞哈斯願意信靠上主而非其他國家勢力,在謹慎鎮定下守護猶大國,王權自可以站立得穩.

但當以賽亞要求亞哈斯求一個記號(兆頭),以作為上主應許保護大衛王朝的明證,但亞哈斯卻說:「我不求,我不試探耶和華」.那是什麼意思?表面上亞哈斯好像有堂皇的理由「不試探上主」,但背後其實是拒絕上主介入他的國防軍事和與強國的外交工作.從當其時事情結局和列王記／歷代志的記載來看,亞哈斯其實已定意尋求亞述的幫助,他可能認為自己有法子面對當前危機,上主的介入亳無意義.

這個又回到我一開始所講,應許,期望的複雜關係.亞哈斯的回應,其中可能有以下的原因:

Case 1: 我唔再信佢, 信佢係無用ge /佢唔可信/佢都唔惦 (失望)

Case 2: 我唔需要佢, 我自己攪得掂 (無望)

Case 3: 呢D 事唔關佢事, 等我地自己做ye  (無望)

Case 4: 求都無用, 佢都唔理呢D 國際大事ge  (失望)

亞哈斯最大的危機是王位不保,在上主保證下仍以亞述作他解危方案,正顯示他相信自己的謀略,勝於被動地依靠上主來得實際.亞哈斯過於自信?以為上主應許太遙遠?太虛無漂眇?實際一點,親自解決危機勝於被動等待?

對於平民百姓,受帝國壓迫,可能因迫於無奈而對上主失望；對於帝王將相,以為自己可以解決任何問題,好像不需任何應許,自然對上主也不再有期望.

結語:應許實現,但不是以你的期望實現.你會點?

讓我們再聽一次天使對約瑟的宣告:

「大衛的子孫約瑟,不要怕,把你的妻子馬利亞娶過來,因她所懷的孕是從聖靈來的.21她將要生一個兒子,你要給他起名叫耶穌,因他要將自己的百姓從罪惡裡救出來.」

不錯,不是所有人對上主的應許失去期望.當時仍有一些猶太人,對上主會派遣彌賽亞來復興萬事的應許充滿期望.他們或會字面地認為推翻帝國壓迫的,將會是大衛王朝之後,大衛的子孫,帶領他們軍事起義推翻暴政.所以約瑟作為大衛的兒子,他的子裔必然就是一個帶領群眾的強人.那知他卻死生十字架上,破碎了這些人對應許的期望.但馬太福音,在面對一切都變失敗的處境,卻要讓他們重新看見,耶穌其實並不是約瑟的兒子.天使的宣告,告訴我們,馬利亞所懷的孕,是從聖靈來的.今日沒有時間再細看馬太福音開始的一大段家譜,但在家譜的結束,馬太福音同樣告訢我們,耶穌或許是約瑟的兒子,大衛的子孫,但1:16卻說: 約瑟…馬利亞的丈夫,那稱為基督的是從馬利亞生的.

耶穌家譜中有幾個女性:他瑪,路得和烏利亞的妻,加上未婚懷孕的馬利亞,都是有失身份,緋聞纏身的女性角色.馬太福音是否告訢我們,以馬內利,神的親臨,並不是我們所期望的,由正統,從政治去實現的,卻是由你所想不到的過程,而實現的更是超出你的期望.神並不是派遣一個由人間王朝後裔所生的人物,帶來政治的復興.上主所實現的應許,是自己親臨到人間,在苦難和壓迫中與人同在.最後是這個復活主帶我們進入末世的天國,所要復興的是一個永恆的國度!

這樣,應許是錯落了嗎?還是只是不按你的期望去實現呢?

面對耶穌基督降生,「以馬內利」上帝與我們同在的應許,我們是已經因為生活的打壓,無奈地放棄期望?還是已經滿足現狀,以為我們「攪得掂」就不再期望?抑或其實我們對於聖誕不是沒有期望,但上主給我們實現的應許,並不是按著我們要求和想像來實現的呢?但你又有否想過,上主所應許的,比我們期望的或沒有期望的,更加美好,更值得擁有?



\newpage

\section{至人只是常}
\label{sec:679}
\textbf{林豪恩先生 - 至人只是常}
\newline
\newline
\hyperref[sec:678]{< < < PREV SERMON < < <}
~
\hyperref[sec:toc]{[返目錄]}
~
\hyperref[sec:680]{> > > NEXT SERMON > > >}


路加福音 2:22-40
\newline
\begin{longtable}{cl}
\hline
\hline
章節 & 經文 (和合本修訂版)\\
\hline
2:22 & \begin{tabularx}{0.7\textwidth}{X} 按摩西律法滿了潔淨的日子，他們就帶著孩子上耶路撒冷去，要把他獻給主。 \end{tabularx} \\ \\ \relax
2:23 & \begin{tabularx}{0.7\textwidth}{X} 正如主的律法上所記：「凡頭生的男子必歸主為聖」； \end{tabularx} \\ \\ \relax
2:24 & \begin{tabularx}{0.7\textwidth}{X} 又要照主的律法上所說，用一對斑鳩，或用兩隻雛鴿獻祭。 \end{tabularx} \\ \\ \relax
2:25 & \begin{tabularx}{0.7\textwidth}{X} 那時，在耶路撒冷有一個人，名叫西面；這人又公義又虔誠，素常盼望以色列的安慰者來到，又有聖靈在他身上。 \end{tabularx} \\ \\ \relax
2:26 & \begin{tabularx}{0.7\textwidth}{X} 他得了聖靈的啟示，知道自己未死以前必看見主所立的基督。 \end{tabularx} \\ \\ \relax
2:27 & \begin{tabularx}{0.7\textwidth}{X} 他受了聖靈的感動，進入聖殿，正遇見耶穌的父母抱著孩子進來，要照律法的規矩而行。 \end{tabularx} \\ \\ \relax
2:28 & \begin{tabularx}{0.7\textwidth}{X} 西面就把他抱過來，稱頌神說： \end{tabularx} \\ \\ \relax
2:29 & \begin{tabularx}{0.7\textwidth}{X} 「主啊，如今可以照你的話，容你的僕人安然去世； \end{tabularx} \\ \\ \relax
2:30 & \begin{tabularx}{0.7\textwidth}{X} 因為我的眼睛已經看見你的救恩， \end{tabularx} \\ \\ \relax
2:31 & \begin{tabularx}{0.7\textwidth}{X} 就是你在萬民面前所預備的： \end{tabularx} \\ \\ \relax
2:32 & \begin{tabularx}{0.7\textwidth}{X} 是啟示外邦人的光，是你民以色列的榮耀。」 \end{tabularx} \\ \\ \relax
2:33 & \begin{tabularx}{0.7\textwidth}{X} 孩子的父母因論耶穌的這些話就驚訝。 \end{tabularx} \\ \\ \relax
2:34 & \begin{tabularx}{0.7\textwidth}{X} 西面給他們祝福，又對孩子的母親馬利亞說：「這孩子被立，是要叫以色列中許多人跌倒，許多人興起；又要成為毀謗的對象， \end{tabularx} \\ \\ \relax
2:35 & \begin{tabularx}{0.7\textwidth}{X} 叫許多人心裡的意念顯露出來；你自己的心也要被劍刺透。」 \end{tabularx} \\ \\ \relax
2:36 & \begin{tabularx}{0.7\textwidth}{X} 又有位女先知，名叫亞拿，是亞設支派法內力的女兒，年紀已經老邁，從童女出嫁，同丈夫住了七年， \end{tabularx} \\ \\ \relax
2:37 & \begin{tabularx}{0.7\textwidth}{X} 就寡居了，現在已經八十四歲。她不離開聖殿，禁食祈求，晝夜事奉神。 \end{tabularx} \\ \\ \relax
2:38 & \begin{tabularx}{0.7\textwidth}{X} 正當那時，她進前來感謝神，對一切盼望耶路撒冷得救贖的人講論這孩子的事。 \end{tabularx} \\ \\ \relax
2:39 & \begin{tabularx}{0.7\textwidth}{X} 約瑟和馬利亞照主的律法辦完了一切的事，就回加利利，到自己的城拿撒勒去了。 \end{tabularx} \\ \\ \relax
2:40 & \begin{tabularx}{0.7\textwidth}{X} 孩子漸漸長大，強健起來，充滿智慧，又有神的恩典在他身上。 \end{tabularx} \\ \\
[1ex]
\hline
\hline
\end{longtable}
講題:至人只是常 The Extraordinary in the Ordinary

經課:路加福音2:22-40

講員:林豪恩先生

場合:香港中文大學崇基學院禮拜堂主日崇拜

日期:2023年12月31日

你認為自己是平凡人,還是不平凡的人呢?有人認為平凡人貌不驚人,不平凡的人氣宇軒昂,你是否同意呢?《菜根譚》一書中有這樣一句話:「濃肥辛甘非真味,真味只是淡；神奇卓異非至人,至人只是常.」意思是真正的美味無需添加過多的調味料,所謂能人異士往往看上去平淡無奇.最神奇的可能是能夠在尋常中看見不尋常,在平淡無奇中發現偉大.今天是聖誕後第一主日,耶穌基督出場了,到底是氣宇軒昂,鶴立雞群的神奇嬰孩,還是貌不驚人普通ＢＢ呢?

一．如一般按摩西律法行事的嬰兒

2:22按摩西律法滿了潔淨的日子,他們就帶著孩子上耶路撒冷去,要把他獻給主.

2:23正如主的律法上所記:「凡頭生的男子必歸主為聖」；

2:24又要照主的律法上所說,用一對斑鳩,或用兩隻雛鴿獻祭.

你如何看得出與眾不同的重要人物呢?例如當你在機場辦理登機手續時,和你一齊排長龍的,你不會認為是與眾不同的重要人物.然而,當你見到有人被專車送來,有專人開路,不但不用排隊,更有特別通道出入,你會猜想這位不是普通人了.我們認為接受一般待遇,按一般規矩行事的是一般人；相反,不平凡的人以特別待遇顯示出他們的不平凡.

按此常識推測,耶穌應該是平凡的普通人,因為他的待遇是一般人的待遇,按規矩而行:「按摩西律法」,「正如主的律法上所記」,「又要照主的律法上所說」,所獻的祭,是一般平民百姓等級的祭品.看來,這個耶穌只是萬千嬰兒的其中一個,是平凡人中的平凡人.正如耶穌的其中一個門徒拿但業曾經說過:「拿撒勒還能出甚麼好的嗎?」

開始時,我問大家,你認為自己是平凡人還是不平凡的人?我相信大部分人都認為自己是平凡人,望左望右,自己和其他人沒有甚麼大分別?然而,每個人底裡都覺得自己是獨特的,不是任何人的其中一個.或者,這是人,既是和其他人一樣的人,又不是任何人的其中一個.問題是,誰能夠在眾裡尋得你呢?誰能夠在你的平凡中發現你的不平凡呢?當年在聖殿中,誰能夠從從一般的嬰兒中發現聖嬰呢?甚麼樣的人能夠從平凡中看見不凡呢?

諸葛亮在《出師表》中如此描述自己:「臣本布衣,躬耕於南陽,苟全性命於亂世,不求聞達於諸侯.先帝不以臣卑鄙,猥自枉屈,三顧臣於草廬之中,諮臣以當世之事.」諸葛亮這位能人異士藏身於布衣之內,在耕者群中,住草廬之下.劉備發掘諸葛亮出山相助,扭轉了自己的頹勢,邁進新的境界.

二．在平常中看見不平常的西面和亞拿

當年在耶路撒冷聖殿,的確有人能夠從芸芸眾嬰兒中看得出耶穌的獨特:

2:25那時,在耶路撒冷有一個人,名叫西面；這人又公義又虔誠,素常盼望以色列的安慰者來到,又有聖靈在他身上.

2:26他得了聖靈的啟示,知道自己未死以前必看見主所立的基督.

2:27他受了聖靈的感動,進入聖殿,正遇見耶穌的父母抱著孩子進來,要照律法的規矩而行.

2:28西面就把他抱過來,稱頌神說:

2:29「主啊,如今可以照你的話,容你的僕人安然去世；

2:30因為我的眼睛已經看見你的救恩,

2:31就是你在萬民面前所預備的:

2:32是啟示外邦人的光,是你民以色列的榮耀.」

2:36又有位女先知,名叫亞拿,是亞設支派法內力的女兒,年紀已經老邁,從童女出嫁,同丈夫住了七年,

2:37就寡居了,現在已經八十四歲.她不離開聖殿,禁食祈求,晝夜事奉神.

2:38正當那時,她進前來感謝神,對一切盼望耶路撒冷得救贖的人講論這孩子的事.

後世讀者已經得知耶穌的來龍去脈,言行際遇及信仰意義,對於這些對嬰兒耶穌的講論感到理所當然,沒有特別反應,然而,當時他的父母卻沒有這種後來的視角,他們的反應是怎樣的呢?

2:33孩子的父母因論耶穌的這些話就驚訝.

三．能夠從平凡中看出不平凡的人

當時人人都看見嬰兒耶穌,難道人人都看見聖嬰?正如今天慶祝聖誕的人,難道人人都看見救主嗎?我們常常以為看見或看不見決定於被觀看的對象,對象展現出他的不平凡,我才能看見；對象沒有特別,我能看見甚麼呢?現實又不一定是這樣的,看見或看不見,有時是和觀察者有關的.

何以聖殿那裡多人,只有西面和亞拿兩人看見呢?有甚麼特別讓他看得出嬰孩耶穌的不平凡呢?路加福音指出,第一,他「公義」的人；第二,他是「虔誠」的人；第三,他「素常盼望以色列的安慰者來到」；第四,「又有聖靈在他身上」.至於亞拿,路加福音指出她也是一位渴慕神的人:「她不離開聖殿,禁食祈求,晝夜事奉神」.公義虔誠的人渴望神,他們能夠看見道成肉身的神—就算只是嬰兒耶穌,正是「你們祈求,就給你們；尋找,就尋見；叩門,就給你們開門」.

有一次,同學們討論天堂地獄的話題,他們列出多種案例,詢問誰上天堂,誰不能上天堂.我建議他們想一想,如果在天堂就是和神一起及由神管治,甚麼人享受這樣的狀況呢?對於那些根本不想與神一齊,又不想被神管治的人,在天堂未必感到快樂的,因此,天堂未必是他們的天堂.

基督徒都是平凡的人,但基督徒又是不平凡的人.基督徒的不平凡,不是因為信了耶穌變超人,而是好像西面和亞拿一樣,從平凡人中發現耶穌的不平凡.

四．應用:洗禮

今天,我們為發現和確認耶穌的不平凡的新人舉行水禮.你們不要忘記,就算接受了水禮,仍然是平常的人；然而,你又是不平凡的人,因為你進入了一個新的國度,一個不平常的群體.這樣的事情是如何發生的呢?

在耶穌的年代,男嬰接受割禮.宗教改革家約翰加爾文對割禮和洗禮的關係如此論述:「割禮和洗禮都是恩約的外在記號和印記.雖然它們在外表上有不同,但其實質是一樣的,都是代表上帝的應許.上帝藉著這兩個記號和印記,在不同的時代告訴選民,他們已經被收納進入上帝的家中,成為上帝國度的一分子,得以享受天國的福分.換句話說,割禮可以說是洗禮的一個雛形.古時的以色列人從這個雛形中,得到我們從洗禮所得到的相同的屬靈的應許.它們是神與我們立約的記號,在我們身上打上成為上帝家人的印記.(4.16.2-3)關於嬰兒洗禮,約翰加爾文認為,因上帝的這個記號,就如把文件蓋上封印,傳給嬰孩,證實這個賜給敬虔父母的應許,並宣告它獲得確認:主不但要作他的上帝,也同樣要作他後裔的上帝；他不只要向他們彰顯自己的慈愛與恩典,也要向他的後裔彰顯這恩典,直到千代 .父母該孩子加入恩典的群體,希望孩子知道在起伏不定的人生路上,有一位眷顧他和憐憫他的上帝,並且協助他成為這位上帝眷顧及憐憫的人．今天,我們在此舉行水禮．無論是成人或是兒童,都透過水禮成為上帝家的人,成為上帝眷顧和憐憫的人.

五．覆歸平常

2:39約瑟和馬利亞照主的律法辦完了一切的事,就回加利利,到自己的城拿撒勒去了.

2:40孩子漸漸長大,強健起來,充滿智慧,又有神的恩典在他身上.

水禮之後,你們就回到自己的家中,或學校,或工作場所,過著平常的生活,但不要忘記你們蒙受上帝恩典的人,活在恩典之中,是上帝家中的成員.願你們好像嬰孩耶穌一樣,漸漸長大,強健起來,充滿智慧,又有上帝的恩典在你們身上」.



\newpage

\section{傳道者約翰}
\label{sec:680}
\textbf{高國雄牧師 - 傳道者約翰}
\newline
\newline
\hyperref[sec:679]{< < < PREV SERMON < < <}
~
\hyperref[sec:toc]{[返目錄]}
~
\hyperref[sec:681]{> > > NEXT SERMON > > >}


講題:傳道者約翰 John the Preacher

講員:高國雄牧師

場合:香港中文大學崇基學院禮拜堂主日崇拜

日期:2024年1月7日

約翰的形象

可1:6約翰穿駱駝毛的衣服,腰束皮帶,吃的是蝗蟲,野蜜.

與耶穌的關係

表哥(相差半年左右)(路加福音1)

名字同樣是上帝/天使改的

路1:13-15

13天使對他說:「撒迦利亞,不要害怕,因為你的祈禱已經被聽見了.你的妻子伊利莎白要給你生一個兒子,你要給他起名叫約翰. 14你必歡喜快樂；有許多人因他出世,也必喜樂. 15他在主面前將要為大,淡酒濃酒都不喝,從母腹裡就被聖靈充滿了.

太1:20-21

有主的使者向他(約瑟)夢中顯現,說:「大衛的子孫約瑟,不要怕!只管娶過你的妻子馬利亞來,因她所懷的孕是從聖靈來的. 21她將要生一個兒子,你要給他起名叫耶穌,因他要將自己的百姓從罪惡裡救出來.」

施洗者約翰

為耶穌預備道路,為悔罪者施洗

馬可福音1:

1上帝的兒子,耶穌基督福音的起頭. 2正如先知以賽亞書上記著說:看哪,我要差遣我的使者在你前面,預備道路.3在曠野有人聲喊著說:預備主的道,修直他的路.4 照這話,約翰來了,在曠野施洗,傳悔改的洗禮,使罪得赦. 5猶太全地和耶路撒冷的人都出去到約翰那裡,承認他們的罪,在約旦河裡受他的洗.

先鋒

先知約翰

成為上帝的傳聲筒

獨特形象的需要及原因?

吸引人的佈道方式

使人願意到來認罪悔改

約翰的使命

可1:7-8

7他傳道說:「有一位在我以後來的,能力比我更大,我就是彎腰給他解鞋帶也是不配的. 8我是用水給你們施洗,他卻要用聖靈給你們施洗.」

約1:15約翰為他作見證,喊著說:「這就是我曾說:『那在我以後來的,反成了在我以前的,因他本來在我以前.』」

約1:29-34

29次日,約翰看見耶穌來到他那裡,就說:「看哪,上帝的羔羊,除去世人罪孽的! 30這就是我曾說:『有一位在我以後來,反成了在我以前的,因他本來在我以前.』 31我先前不認識他,如今我來用水施洗,為要叫他顯明給以色列人.」 32約翰又作見證說:「我曾看見聖靈,彷彿鴿子從天降下,住在他的身上. 33我先前不認識他,只是那差我來用水施洗的,對我說:『你看見聖靈降下來,住在誰的身上,誰就是用聖靈施洗的.』 34我看見了,就證明這是上帝的兒子.」

為耶穌施洗

可1:9-11

9那時,耶穌從加利利的拿撒勒來,在約旦河裡受了約翰的洗. 10他從水裡一上來,就看見天裂開了,聖靈彷彿鴿子,降在他身上. 11又有聲音從天上來,說:「你是我的愛子,我喜悅你.」

被聖靈充滿

耶穌認同,認證他的召命

太3:13-17

13當下,耶穌從加利利來到約旦河,見了約翰,要受他的洗. 14約翰想要攔住他,說:「我當受你的洗,你反倒上我這裡來嗎?」 15耶穌回答說:「你暫且許我,因為我們理當這樣盡諸般的義.」於是約翰許了他. 16耶穌受了洗,隨即從水裡上來.天忽然為他開了,他就看見上帝的靈彷彿鴿子降下,落在他身上. 17從天上有聲音說:「這是我的愛子,我所喜悅的.」

約翰的獨特之處

定位清晰

四福音同樣記述

馬可福音:「福音的起頭」

耶穌承接了約翰這福音的棒

可1:14-15

14 約翰下監以後,耶穌來到加利利,宣傳上帝的福音, 15說:「日期滿了,上帝的國近了.你們當悔改,信福音!」

約1:35-37

35再次日,約翰同兩個門徒站在那裡. 36他見耶穌行走,就說:「看哪,這是上帝的羔羊!」 37兩個門徒(其中一位就是安得烈)聽見他的話,就跟從了耶穌.

盡心服侍至終

馬可福音6

17先是希律為他兄弟腓力的妻子希羅底的緣故,差人去拿住約翰,鎖在監裡,因為希律已經娶了那婦人. 18約翰曾對希律說:「你娶你兄弟的妻子是不合理的.」 19於是希羅底懷恨他,想要殺他,只是不能； 20因為希律知道約翰是義人,是聖人,所以敬畏他,保護他,聽他講論,就多照著行,並且樂意聽他. 21有一天,恰巧是希律的生日,希律擺設筵席,請了大臣和千夫長,並加利利作首領的. 22希羅底的女兒進來跳舞,使希律和同席的人都歡喜.王就對女子說:「你隨意向我求甚麼,我必給你.」 23又對她起誓說:「隨你向我求甚麼,就是我國的一半,我也必給你.」 24她就出去對她母親說:「我可以求甚麼呢?」她母親說:「施洗約翰的頭.」 25她就急忙進去見王,求他說:「我願王立時把施洗約翰的頭放在盤子裡給我.」 26王就甚憂愁；但因他所起的誓,又因同席的人,就不肯推辭, 27隨即差一個護衛兵,吩咐拿約翰的頭來.護衛兵就去,在監裡斬了約翰, 28把頭放在盤子裡,拿來給女子,女子就給她母親. 29約翰的門徒聽見了,就來把他的屍首領去,葬在墳墓裡.

施洗者成為殉道者

這就是傳道者約



\newpage

\section{始料不及}
\label{sec:681}
\textbf{李均熊牧師 - 始料不及}
\newline
\newline
\hyperref[sec:680]{< < < PREV SERMON < < <}
~
\hyperref[sec:toc]{[返目錄]}
~
\hyperref[sec:682]{> > > NEXT SERMON > > >}


講題:始料不及 Unexpected

講員:李均熊牧師

場合:香港中文大學崇基學院禮拜堂主日崇拜

日期:2024年1月14日

起:撒母耳的錯聽 (Samuel’s mishearing)

今天顯現期第二週舊約經課,是年青時候的撒母耳因為錯聽晚間的呼喚,以為是來自他所服侍的老祭司以利,而差點錯過上主的啟示.

當時的撒母耳可能已經十三,四歲,而以利已經年老,眼目昏花,看不分明(應該是白內障吧).守護聖殿長明燈的責任,就落在撒母耳身上.所以和合本第三節經文,(可能參考英王欽定譯本)沒有譯出撒母耳其實是睡在上主的殿中.和修本則讓我們清楚看到:上帝的燈還沒有熄滅,撒母耳睡在耶和華的殿內,上帝的約櫃就在那裡.在夜半撒母耳聽到呼叫他的聲音,年青的撒母耳就以為老祭司以利有甚麼需要,即時應了句「我在這裡!」,就一支箭的跑到以利面前問他,你叫我嗎?我在這裡.不知以利會否欣賞這個孩子的孝心,就對他說,我沒有叫你,你回去睡吧.撒母耳見以利睡在被窩好地地的,也安心回去殿中再睡.

沒多久,撒母耳又聽到呼喚他的聲音,就又起來去到以利處問他是否叫他,今次以利真係感到撒母耳的孝心,就加了句:我兒,我沒有叫你,回去睡吧.到了第三次,撒母耳不厭其煩仍然起來去到以利身邊問他是否叫他,反而以利感到奇怪,怎麼這個小子就算多麼孝順,也不會這麼一晚三次的來找他吧,於是意識到可能睡在聖殿中的小子聽到的是上主的呼喚,於是就教他,再聽到聲音呼叫他,就回應說:上主,請說,僕人敬聽.

又再一次,撒母耳聽到呼叫他的聲音,他也忘記要先講出上主名號才作回應,就直接應了:請說,僕人敬聽.撒母耳聲到的,卻是上主要貫徹祂在之前對以利家的判語,「叫聽見的人都必耳鳴」,以至第二天撒母耳也不敢將上主的默示告訴以利,要以利賭咒叫撒母耳不可隱瞞,他才將一切話告訴以利.

撒母耳以為年老的祭司以利在半夜需要他的服侍,所以不斷呼叫他；但始料不及,原來是上主有話要向他說.更始料不及的,是上主向他所說的話,竟是對自己服侍的老祭司一家的凶言判語.

承:拿但業被看見

顯現期第二週福音經課,是腓力叫拿但業來看耶穌,而他反而早被耶穌所看見.

之前施洗者約翰的門徒安得烈和另一個約翰的門徒,聽見約翰稱耶穌為上帝的羔羊,就去跟著耶穌；耶穌看見他們跟著,就問他們要甚麼,他們問了拉比,在哪裡住,耶穌就叫他們來看.然候,安得烈就領他哥哥西門彼得去見耶穌.接著,這一次,是耶穌主動去加利利,找到腓力,就叫他來跟從他.就像安得烈找西門彼得去見耶穌,腓力也找著拿但業,對他說他見到律法和先知所寫的那一位,就是約瑟的兒子拿撒勒人耶穌.怎知拿但業卻說:拿撒勒還能出甚麼好的嗎?(就好像有人在崇基學院神學院上完我的課,找來友人說:崇基神學院有個舊約講師好勁架,怎知他會說:崇基還能出甚麼好東西)……

腓力堅持:你來看!耶穌看見拿但業就指著他說了句話:看哪,這是真以色列人,他心裡沒有詭詐的.拿但業就覺得很奇怪,就問:你從哪裡知道我呢?耶穌的回答仍是使人一頭霧水:腓力還沒有招呼你,你在無花果樹底下,　我就看見你了——但因為這一句話,拿但業就和之前輕蔑的態度 180度的改變過來:拉比,你是上帝的兒子,你是以色列的王!這回到耶穌覺得奇怪:因為我說在無花果樹底下看見你,你就信嗎 ! ?

拿但業以為拿撒勒來的耶穌不是甚麼好東西,只是腓力堅持要他過去看；但始料不及,原來被看見的不是耶穌,而是在無花果樹下的拿但業被看見.拿但業也真是始料不及,拿撒勒來的竟是神的兒子,以色列的王!

轉:我們有限的視野

解經家並不明白耶穌的那句關於「在無花果樹下就看見你」為何對拿但業有這樣大的影響.從猶太拉比認為無花果樹下是人學習律法的地方,到撒迦利亞書3:8-10因為同時提到「大衛苗裔」和「在那日你們各人要請鄰舍坐在葡萄樹和無花果樹下」,因而是與彌賽亞相關,都有學者認為是這裡無花果樹下的意思.究竟是甚麼原因令到拿但業的態度有這樣的改變,我們可能永遠都不會知曉.但拿但業對拿撒勒人的偏見(prejudice)卻是顯而易見的.拿但業對拿撒勒人的偏見有可能是產生於當時社會對北方小鎮的一般見解.受人重視的人物,通常若不是由中部比較農產經濟富庶的地區,就是由南部猶大或耶路撒冷等政治中心而來的.像拿撒勒這些落後的窮鄉僻壤,會有甚麼好東西?

撒母耳起初並不知道叫喚他的是上主的聲音,不全然是因為撒母耳年幼資歷淺,亦有可能是時代給予的負條件.經文充分描寫了當時聆聽上主聲音的負條件:言語稀少不常有默示v.1 /以利老眼昏花看不分明v.2/撒母耳還未認識耶和華未(what?)未得默示(耶和華的話也未曾向他啟示) v.7.

其實我們都有我們的偏見,若稱我們有偏見有些負面,或者稱之為前設了解(pre-understanding)吧,若同時比較一下撒母耳的處境,我們了解面前的事情都會有一定假設,有一定的因為社會,文化,成長,生活而對事情的先存的看法和想法,而這些都會造成我們回應上主召喚的反應.上主很多時候的召請,都是我們始料不及的.其中的原因當然是上帝獨行奇妙的事,和我們平常所經驗的會有很大距離；但上主的召請,會不會同時顛覆我們對人事的偏見,對信仰,對人生一些前設理解,或對自己的生活,又或這個時代某些事情有可能發生,或沒有可能發生的既定看法?

合:信仰仍可使我門詫異嗎?

究竟耶穌的思典仍否使我們感覺詫異?會否我們對現況太多的限制,以至失去信仰那奇妙的感覺?在2024年,你準備在你始料不及的處境下與神相遇嗎?你會否跨過你的預設的框框,又或不按社會的偏見,時代的標準或別人的期望來選擇和生活,去體會上主對你的召喚與看見?



\newpage

\section{命限恩典義}
\label{sec:682}
\textbf{鄧瑞強博士 - 命限,恩典,義}
\newline
\newline
\hyperref[sec:681]{< < < PREV SERMON < < <}
~
\hyperref[sec:toc]{[返目錄]}
~
\hyperref[sec:683]{> > > NEXT SERMON > > >}


傳道書 3:1-14
\newline
\begin{longtable}{cl}
\hline
\hline
章節 & 經文 (和合本修訂版)\\
\hline
3:1 & \begin{tabularx}{0.7\textwidth}{X} 凡事都有定期，天下每一事務都有定時。 \end{tabularx} \\ \\ \relax
3:2 & \begin{tabularx}{0.7\textwidth}{X} 生有時，死有時；栽種有時，拔出有時； \end{tabularx} \\ \\ \relax
3:3 & \begin{tabularx}{0.7\textwidth}{X} 殺戮有時，醫治有時；拆毀有時，建造有時； \end{tabularx} \\ \\ \relax
3:4 & \begin{tabularx}{0.7\textwidth}{X} 哭有時，笑有時；哀慟有時，跳舞有時； \end{tabularx} \\ \\ \relax
3:5 & \begin{tabularx}{0.7\textwidth}{X} 丟石頭有時，撿石頭有時；懷抱有時，不抱有時； \end{tabularx} \\ \\ \relax
3:6 & \begin{tabularx}{0.7\textwidth}{X} 尋找有時，失落有時；保存有時，拋棄有時； \end{tabularx} \\ \\ \relax
3:7 & \begin{tabularx}{0.7\textwidth}{X} 撕裂有時，縫補有時；沉默有時，說話有時； \end{tabularx} \\ \\ \relax
3:8 & \begin{tabularx}{0.7\textwidth}{X} 喜愛有時，恨惡有時；戰爭有時，和平有時。 \end{tabularx} \\ \\ \relax
3:9 & \begin{tabularx}{0.7\textwidth}{X} 這樣，做事的人在他所勞碌的事上得到甚麼益處呢？ \end{tabularx} \\ \\ \relax
3:10 & \begin{tabularx}{0.7\textwidth}{X} 我觀看神給世人的擔子，使他們在其中勞苦： \end{tabularx} \\ \\ \relax
3:11 & \begin{tabularx}{0.7\textwidth}{X} 神造萬物，各按其時成為美好，又將永恆安放在世人心裡；然而神從始至終的作為，人不能測透。 \end{tabularx} \\ \\ \relax
3:12 & \begin{tabularx}{0.7\textwidth}{X} 我知道，人除了終身喜樂納福，沒有一件幸福的事。 \end{tabularx} \\ \\ \relax
3:13 & \begin{tabularx}{0.7\textwidth}{X} 並且人人吃喝，在他的一切勞碌中享福，這也是神的賞賜。 \end{tabularx} \\ \\ \relax
3:14 & \begin{tabularx}{0.7\textwidth}{X} 我知道神所做的都必存到永遠；無所增添，無所減少。神這樣做，是要人在他面前存敬畏的心。 \end{tabularx} \\ \\
[1ex]
\hline
\hline
\end{longtable}
講題:命限,恩典,義 Limitation, Gift, and Rightness

經文:傳道書3:1-14

講員:鄧瑞強博士

場合:香港中文大學崇基學院禮拜堂主日崇拜

日期:2024年1月21日

各位弟兄姊妹,早安.願神的恩典常與你們同在.

我們今日繼續看智慧文學.我將分享傳道書的智慧.傳道書是舊約智慧文學裡神學思想發展到最高峰的作品,是一位智者對塵世事情的深刻領悟.

先用兩節經文作為開場白:

1:2　傳道者說:虛空的虛空,虛空的虛空,凡事都是虛空.

1:14　我見日光之下所做的一切事,都是虛空,都是捕風.

「虛空的虛空」,這種句式,是希伯來文表達最高級意思的句式,如:King of Kings,Song of Songs.「虛空的虛空」,即最虛空,虛空到極點.希伯來文「虛空」這個字,原始的含意是「氣」(空氣的氣).氣,不是完全沒有東西,不是「無」,只是捉不到而已.有些人理解這「虛空」為虛無主義或悲觀主義的nothingness或meaningless,我認為不對.「虛空」是氣,它不是什麼也沒有,也不是完全沒有意義,只是你要捉住它,掌控它,就只會徒勞無功.

傳道書這「虛空」的講法,類似佛家講的「空」.佛家講的「空」也不是指什麼東西也沒有,而只是指一切事物皆依因緣(conditional causality)而生,緣聚則在,緣散則滅,背後沒有一獨立的實體,此即佛教講的「諸法無我」.「法」,即事物.「諸法」,即一切事物,皆互相依存,沒有一獨立的「我」性,故曰「諸法無我」.這不等於「無」,而只是沒有不變的實體而已.根本沒有實體而你硬要認為某事物是實體,就是執著.「健康」本來就不是一樣不變的「實體」,要依據很多因素才能出現,你硬要永遠擁有它,便是執著.有執著,就有煩惱.

1:14　我見日光之下所做的一切事,都是虛空,都是捕風.

傳道者說,人間生活的實相,都是「虛空」.我再說,「虛空」不是什麼也沒有,不是做什麼也無意義.只是你硬要去掌控像「空氣」一樣的東西,便是這裡說的「捕風」,就是試圖去追逐風,去抓緊風.氣,或風,雖是有某些東西存在,但它的本質就是不能被抓住.硬要去抓住風,就是硬要去執著無法執著的東西,當然會帶來虛空的感受.

在第14節,我們讀到「日光之下」這字眼,再一次提醒我們智慧文學的用心所在.智慧文學只討論「日光之下」可見的事情,不討論「日光之上」不可見的領域.例如:智慧文學會討論「日光之下」人有肥有瘦這現象,會討論肥人宜做什麼,瘦人宜做什麼；但不會討論「為何神造我,不吃飯也肥；造她,卻任吃也瘦,如此這般,神是否不公平」這等問題.智慧文學會討論:病了,看哪個醫生較好；但不會討論:「為何我這般敬虔,又生活健康,卻得絕症」這問題；也不會回答:「神為何不聽我禱告,為何不醫治我」這等問題.智慧文學討論「日光之下」的生活現實,不討論「日光之上」的救贖問題.當然,救贖問題也是重要的.沒有救贖,人會感到失落,但從智慧文學的角度講,人不正視「日光之下」的現實,產生的問題也一大堆.如:應放手的時候,你不放手,與其說神不救你,不如說,你不正視現實,自造麻煩.

在youtube裡看到,有些人在洞口很窄的地洞裡,或開了個窄口的葫蘆瓜裡,放些食物,就能引到猴子來,猴子引手入洞口取食物,由於牠們緊握拳頭不放,牠們無法把手收回來,於是,被人捉著了.若果猴子有智慧文學,那麼,牠們會讀到,智慧文學沒有教牠們祈禱,求神解救,而是教牠們放手.這是日光之下完全可以理解,完全可以實踐的道理.

接著,我們看看第3章的一段經文,在這裡,傳道者解釋:「為何日光之下所做的一切事,都是虛空,都是捕風?」這段經文是智慧文學最深刻的神學反省.

3:1　凡事都有定期,天下萬務都有定時.

3:2　生有時,死有時；栽種有時,拔出所栽種的也有時；

3:3　殺戮有時,醫治有時；拆毀有時,建造有時；

3:4　哭有時,笑有時；哀慟有時,跳舞有時；

3:5　拋擲石頭有時,堆聚石頭有時；懷抱有時,不懷抱有時；

3:6　尋找有時,失落有時；保守有時,捨棄有時；

3:7　撕裂有時,縫補有時；靜默有時,言語有時；

3:8　喜愛有時,恨惡有時；爭戰有時,和好有時.

這裡有十四對「定時」,a time for everything,涵蓋生命的種種活動.當然,還可以不斷加上不同的「定時」的.炒樓炒到識飛有時,炒到訓街有時.聽道時有所得益有時,睡到打呼嚕有時.

在這裡,傳道者描繪神所造的世界的本相.這個世界,凡事都有一個「定時」.定時,講的就是「時間性」,Temporality.這是哲學上的一個大課題.時間性的特性,就是流變性,變幻不居.萬物在變,這即是佛家所說的「諸行無常」.「行」者,動也.一切事物在變動,而沒有常態.萬物的流變性,對人的感受來說,有好有不好.如:世上有戰爭,但好在戰爭有時限,有其定時,和平會到來,這對人來說,是好的,起碼知道苦難不會永續.但和平亦有其定時,到了某個時刻,又會發生戰爭,這對人來說,是不好的,因人人都想和平永續.若人想抓住和平,便發覺像捕風一樣,抓不住,捉不到.就是這種捉不到,人感到不安.

中國的易經,講變易,很早便察覺生命在恆常變化中.易經以六十四卦,象徵人生的不同境遇.每一卦又分六爻,代表每個境遇的變化.易經試圖抓住變化中的規律,其中一個規律,就是物極必反.如:泰極否來,好到盡頭便是壞.按照這規律做人的話,則事情就算很順利,也不要去到盡,凡事留有餘地,若去到盡,事情便會向相反方向變化.易經認為,萬事在變,但變化中仍有規律可尋.傳道書卻認為,萬事在變,各有其定時,而其定時的規律,卻無跡可尋.你以為「否」極就會「泰」來,壞到盡頭便是好,但事實可以是,你一生只見到壞,而不見好.壞有時,好有時,但這個「有時」卻掌管在神的手裡,人永遠無法參透.

這一切「有時」,就是人的命限.命限不同命運,命限即生命的限制.我唱歌不及高牧師好聽,這是我的命限.而命運指的,卻是生命的終極軌跡,被一種無形之力決定.命運是無法改變的,古希臘的悲劇,講的是命運.悲劇人物伊底帕斯(Oedipus)注定會殺父娶母,則無論他如何逃避,他最終就是殺父娶母.傳道書講的「定時」,是一種命限,而不是命運.生有時,死有時,人有生就有死,死是我們的限制,有人說,這是大限.但這命限,並不決定我們如何生,如何死,這是我們真正可以選擇的.人的煩惱,就是不接受自己的限制,卻想如神一樣,操控人生.但,我們怎能克服萬物的「定時」呢?你努力練球,想做球王,但下一刻,就斷了阿基里斯腱,正是「斷腱有時」,你能如何呢?

3:9　這樣看來,做事的人在他的勞碌上有什麼益處呢?

「益處」,指賺到的利潤.世事常變,人努力做事,到頭來,能賺到什麼呢?你努力賺錢,來一場金融風暴,就可能失去所有.你可能努力經營一個幸福的家庭,忽然來一場急病,就可能與所愛的人永別.在變幻不定的人生裡,到頭來,能賺到什麼呢?

3:10　我見上帝叫世人勞苦,使他們在其中受經練.

「上帝叫世人勞苦」:「叫」這個字,就是「賜予」這個字.人生的變幻不居,無法掌控,原來是神賜給人的.這是神創造的美意,為的是使人在其中受經練.所謂「受經練」,其實即:使人深深地經驗時間性下的流變,讓人體會生命的命限.知道限制,從而不去操控,其實是神給人的禮物.

生有時,死有時,這是命限,在神看來,這是好的,但人偏偏看為不好,不能接受死亡的突然其來.於是,我們試圖做上帝,操控自己的命運,這其實就是試圖去掌控無法掌控的「氣」,試圖去捕風,這當然無法成功,而虛空由此而來.

為何人總是有這衝動,有這企圖?

3:11　上帝造萬物,各按其時成為美好,又將永生〔原文是永遠〕安置在世人心裡.然而上帝從始至終的作為,人不能參透.

「神造萬物,各按其時成為美好」:在神的創造裡,每一樣事物的「其時／定時」,都是美好的,合宜的.在這「時間性」的世界裡,神卻同時將「永恆」安置在世人心裡,神將「永恆性」(Eternity)賜給人.「永恆性」(Eternity),哲學上的另一個大課題.一條狗,沒有永恆性,牠在這變易的時間性世界裡生活,沒有憂慮的,牠不用想著去買哪一份保險.但人不同,人兼具Temporality與Eternity,在「時間性」的世界,人強求「永恆性」,虛空之感由此而來.

凡事都有定時,這本是神美好的安排.但人的「永恆性」抗拒「時間性」的變易.在神眼中,生有時,死有時,是很好的.但人不想死,尋求長生不死.這裡講的,不是「信耶穌,得永生」的那種「永生」.智慧文學不討論這「日光之上」的永生；智慧文學討論的,是日光之下,現實人生裡,人人都不想死的欲望.現代人食很多健康食品(不是自然的那種,是精煉出來的,有不少是藥丸狀的),做很多檢查,做大量運動,為的是多少滿足長生不死的欲望.在必死的命限中,尋求永恆不死,這只是捕風,終究不會成功,故此,人感到虛空.

「神從始至終的作為,人不能參透」:由於人以「永恆性」嘗試去抓住不變的東西,好安頓自己怕變的心,故此,人就無法理解神為何如此創造；也不能參透這個充滿「定時」,不可測的世界,到底有什麼好.很多人說:神一創造,便是天堂,那麼,多好.其實,他的意思是說:將人造成天使,一開始,便是永恆的,多好.事實上,天使是羨慕做人的,多少傳統故事,講仙女下凡,為的只是感覺一下短暫,卻具體的人間愛情.我常常想,當我吃雪糕時,天使看到「流口水」.

一個具有「永恆性」的人,卻活在「時間性」的塵世裡,他的心無法在變幻的世界中找到安頓,但也無法像天使般甩掉塵世.他恆常在不滿足的虛空中.人無法忍受這種折磨,於是,他一是放棄「時間性」,一是放棄「永恆性」.

古希臘的柏拉圖,他不能忍受變幻的現實,他選擇放棄塵世的「時間性」.他認為,一切變易的,都是假的.真的東西,必定是永恆的.於是,他追求不變的理型(idea),漠視易變的塵世.他輕視身體,否定慾望,強調心靈對永恆理型的直觀.他活著,只是心靈的活著.死,就是心靈的解脫.現代有些靈修,偏向這種想法,遠離塵世,追求一種心靈的寂靜.內心越平安,就越發對現實反感.看不到,神造這變幻的世界,原是好的.

另有一些人,卻選擇放棄「永恆性」,只浸淫在「時間性」中.哲學家祈克果(Kierkegaard)用西班牙的風流人物「唐璜」(Don Juan)來形容這類人.莫扎特有齣歌劇Don Giovanni就是講「唐璜」這個人物.色情狂唐璜將自己變成慾望的化身,他的生命沒有任何「永恆性」的元素,他的所謂存在,其實就只是一充滿慾望的軀殼不斷地尋求滿足.現代那些不斷地打機的人,接近這種生命形態,不斷地在刺激下作反應,不再有什麼永恆的想像.

柏拉圖是「純心靈」的活著,追求永恆.唐璜是「純肉體」的活著,追求當下的快感.他活得醉生夢死,墮落在流逝的感覺中.

柏拉圖幸福嗎?他的柏拉圖愛情令他經歷不到人間真實愛情的波瀾壯闊,他浪費了塵世.唐璜幸福嗎?他令自己活得像野獸,他浪費了人性.

人能否兼顧「時間性」和「永恆性」,兼顧「肉體」和「心靈」,而不感到虛空?

3:12　我知道世人,莫強如終身喜樂行善；

3:13　並且人人吃喝,在他一切勞碌中享福,這也是上帝的恩賜.

3:14　我知道上帝一切所做的都必永存；無所增添,無所減少.上帝這樣行,是要人在他面前存敬畏的心.

人生的一切,本來就是神的「恩賜」,禮物.神造的世界,變幻不居,卻是他永恆的美意,無所增添,無所減少.神造人有永恆性,自然會對有「定時」,有「命限」的人生不滿足,無論人走向極端的永恆性,自視為神,或走向極端的時間性,自貶為獸,都無法擺脫虛空之感.這是神巧妙的設計,「時間性」使人感到不滿足,「永恆性」使人尋求真正永恆的神.

神這樣行,是叫人敬畏他,尊重他,感謝他.當人面向神,相信神,接受生命的定時與命限都是禮物時,虛空之感便會被克服.原來,我們可以擁抱生命的有限,在敬畏神之心境下,不必以自己的角度去評價事物的好壞.我們可以不視「生」為祝福,視「死」為咒詛,而可以視「生」為祝福,視「死」也是祝福.若生與死,苦與樂,哭與笑都是神的禮物,那麼,我們就不必須硬要永遠擁有生命,快樂,歡笑,不必須硬要抓緊一些根本不能抓緊的東西.一切的限定,一切的際遇,都是好的,都可以開心擁抱的,都是神獨特的禮物,塑造我們成獨特的人.

若能接受生命的一切為禮物,便能作合宜的事.這就是我的講題的第三部分:「義」.「義」者,宜也.人生作何事最宜?

3:12　我知道世人,莫強如終身喜樂行善；

3:13　並且人人吃喝,在他一切勞碌中享福,這也是上帝的恩賜.

「莫強如」:字面直譯是「沒什麼好得過」.其中「好」(tov)字,就是神創造萬物時,神看是「好的(tov)」的那個「好」字.最好的事,就是終身喜樂行善.「行善」的「善」(tov)字,又是神看是「好的」的那個「好」字.行善,在這裡,與喜樂平行,是指一種生活態度,並沒有任何道德含意.在這裡,行善的意思是:看生活的一切是好的,以此心態好好地生活.最合宜的好生活,就是「人人吃喝,在他一切勞碌中享福」.「享福」的「福」(tov)字,又是「好」那個字.所謂「享福」,即享受生命中的好.三次用「好」(tov)這個字,是希望我們看到神所看的,看到這世界及生活的美好.

縱使生命中的一切皆有定時,皆是命限,有成有敗,有愛有恨,有聚有散,但若以敬畏神的心視之,看一切皆為禮物,則我們便可好好地享受這人生.一切都是神的禮物,白白的得來,也當白白地捨去.生命中的一切,可好好享受之,卻千萬不能據為己有.當人將生命,健康,財富,成就,甚至兒女據為己有時,則會千方百計握緊它們,但,這一切皆有定時,皆像風,根本無法掌控.強行去捕風,到頭來,便是虛空.

親愛的弟兄姊妹,好好享受人生,這是神的禮物.既是禮物,就好好和他人分享.生的喜悅,是禮物；死的恐懼,也是禮物.無牽無掛的,活一場好活.

最後,以主耶穌的一段話作結束.

「你們看那天上的飛鳥,也不種,也不收,也不積蓄在倉裡,你們的天父尚且養活牠.你們不比飛鳥貴重得多嗎?你們哪一個能用思慮使壽數多加一刻呢?所以,不要憂慮說,吃什麼?喝什麼?穿什麼?這都是外邦人所求的.你們需用的這一切東西,你們的天父是知道的.你們要先求他的國和他的義,這些東西都要加給你們了.所以,不要為明天憂慮,因為明天自有明天的憂慮；一天的難處一天當就夠了.」(太6: 26-27, 31-34)

願上主得著榮耀,頌讚；願眾人得享喜樂,平安.



\newpage

\section{顯現有時隱藏有時}
\label{sec:683}
\textbf{林豪恩先生 - 顯現有時,隱藏有時}
\newline
\newline
\hyperref[sec:682]{< < < PREV SERMON < < <}
~
\hyperref[sec:toc]{[返目錄]}
~
\hyperref[sec:684]{> > > NEXT SERMON > > >}


講題:顯現有時,隱藏有時 A Time to Reveal, A Time to Reserve

講員:林豪恩先生

場合:香港中文大學崇基學院禮拜堂主日崇拜

日期:2024年1月28日

今天是顯現期第四主日.顯現期有多長呢?可以短至兩個主日,可以長至六個主日,視乎該年復活節的早與遲而定.顯現期是關於甚麼的呢?這節期的焦點是耶穌顯明自己.大家還記得過去三個主日所讀的福音經課嗎?

顯現期第一主日,即是耶穌受洗日,當日的經課如此說:「那時,耶穌從加利利的拿撒勒來,在約旦河裡受了約翰的洗.他從水裡一上來,就看見天裂開了,聖靈彷彿鴿子降在他身上.又有聲音從天上來,說:「你是我的愛子,我喜愛你.」」.(馬可福音1章9-11節)

顯現期第二主日福音經課記載拿但業遇見耶穌之前說:「拿撒勒還能出甚麼好的嗎?」(1:46).遇到耶穌之後,拿但業對耶穌說:「拉比!你是上帝的兒子,你是以色列的王.」(約翰福音1:49).

顯現期第三主日就記述門徒對這位「上帝的愛子,以色列的王」的回應.耶穌對西門和安得列說:「來跟從我,我要叫你們得人如得魚一樣.」他們立刻捨了網,跟從他.」(馬可福音1:17-18)

今天是顯現期第四主日,耶穌向誰顯現呢?或者說,誰發現耶穌的身分呢?這次顯現,產生了甚麼效果呢?這種顯現的現象對我們的信仰生活有甚麼啟發和意義呢?

一．上主向所有人顯現

1:21他們到了迦百農,耶穌就在安息日進了會堂教導人.

1:22他們對他的教導感到很驚奇,因為他教導他們正像有權柄的人,不像文士.

當時的猶太人每逢安息日到會堂聚會,類似我們每逢主日去禮拜堂參加崇拜；他們的文士在會堂教導；類似我們的傳道人在崇拜講道.所以,當耶穌也去會堂教導時,人們自然地把耶穌當作一位文士.坦白說,人們對於這些日復一日,年復一年的講道,還會抱甚麼期望呢?好像每天回家吃飯,通常都是家常便飯.要驚喜?就要出外吃.所以小朋友特別期待出外吃飯.在迦百農會堂裡的人們,或許也持著類似的態度來聽耶穌講道.這次,他們卻感到驚奇.令他們驚奇的,未必是講道的內容,作者也沒有記下內容.令人們感到驚奇的是耶穌這個人,經文指出,「因為他教導他們正像有權柄的人,不像文士」.這種情況對我們有甚麼啟發呢?你還記得基督教信仰令你感到驚奇的時刻嗎?

在座當中有誰曾經在教會學校讀書呢?香港有很多教會學校,不少都是受歡迎的學校.或是幼稚園,或是小學,或是中學甚至大學,相信不少人都曾經置身其中.在教會辦的學校讀書,如果不是每星期聽一次講道,最少都每年聽一次.學生甚麼時候感到驚奇呢?有一位學生的經驗是這樣的:在早會,聖經課,福音週,聽講道,他如果不是遊魂,就是水過鴨背.基督教信仰甚麼時候引起他的關注呢?有一天,他的校服褲「爆tie 」,在非常尷尬之際,有位老師拿來一條運動褲給他替換,然後替他修補「爆tie 」的西褲.有一年,班主任要求同學寫月記,每一次他收到班主任的回覆,篇幅如果不比他所寫的為多,也跟他所寫的差不多.這位同學對講道的內容沒有興趣,但卻對基督徒老師的善良及誠意感到驚奇.再過一些年日,他身邊有些要好的同學成為了基督徒,但那些基督徒同學仍然和他做好朋友,一齊讀書,一齊玩耍,令他感到基督徒的友善及接納.

後來,有人送了一本書給他,講述基督徒如何理解潮流和文化,以及如何在複雜的現實中活出真誠和有意義的人生.這本書開了他的眼界,原來基督徒不是善良的傻瓜,而是對時代和處境有精闢的分析.他們嚮往美善,追求人生意義.對一個開始尋找前路,思考人生,和討論社會議題的年輕人,基督徒的人生及追求令他感到吸引,因為這些人不像電視劇中那些追逐名利或者爭奪武林霸主的人.

各位弟兄姊妹,各位朋友,你還記得你在甚麼時候開始留意基督信仰嗎?信仰令你感到驚奇的是甚麼呢?驚奇,或者是只一種吸引,只是一個開始,不代表和你產生關係,也不代表改變了你甚麼.當時那位學生雖然欣賞基督徒的善良及友好,但他仍然是旁觀者,和他個人的生命仍未有關係.他心裡想,若要把自己投向上主,除非上主向他顯現,或者讓他感到上主的實在.

二．上主向個別人士顯現

1:23當時,會堂裡有一個污靈附身的人,他在喊叫,

1:24說:「拿撒勒人耶穌,你為甚麼干擾我們?你來消滅我們嗎?我知道你是誰,你是上帝的聖者.」

1:25耶穌斥責他說:「不要作聲,從這人身上出來吧!」

1:26污靈使那人抽了一陣風,大聲喊叫,就出來了.

污靈附身的人所喊叫的內容:「你是上帝的聖者」突出了顯現期的主題.不知道大家會不會問,為甚麼耶穌基督是上帝的聖者這個身分只有這個污靈附身個人發現,而其他人沒有發現呢?相信大家都知道,決定對方是甚麼,不單單取決於對方,也取決於觀察者本身.當天,在迦百農會堂,人人都遇到同一個耶穌,但發現耶穌是上帝的聖者就只有這個人.客觀上來說,耶穌向所有人顯現,但主觀上來說,只有那個人發現耶穌的顯現.發現耶穌顯現和沒有發現耶穌顯現是很不同的狀況.你可以從經文中看見,發現耶穌顯現的那個人是不能沒有反應的,而且反應相當強烈.這種反應改變了這個人的狀況,影響了他的生命,這種改變是其他人沒有的.旁觀者只是由感到驚奇,進而感到驚訝.他們只是把看見的事情視為談論的話題.這種談論,和他們的生命是沒有關係,對他們的生命也沒有甚麼影響.但是,經驗耶穌顯現的那個人就不同了,污靈離開了他,表示他的生命改變了,或多或少和之前不再一樣了.

當我們讀到這個片段的時候,催促我們問自己:在自己和上主的關係上,我是親身經歷上主顯現的人?還是在旁邊「食花生」的人?或許每個人經體上主顯現的狀況都不一樣,彼得前書說,上主多次多方曉諭列祖,即是說上主的顯現不限於迦伯農會堂的例子.上主是如何向你顯現的呢?剛才提及的那位同學仍未能信靠上主的原因,就是感到上主既隱藏又遙遠.他不是不同意基督教的教導,甚至覺得做基督徒也是好事,但總是感到欠缺甚麼似的,未能把自己交托給上主.不知道你有沒有類似的掙扎呢?

有一晚,他留校溫習後下山回家.同學各走各路,他一個人在的路上,抬頭望見藍黑的天空,皎潔的月亮和若隱若現的星星,又引起他無限的想像.他向長空發問:「如果天地之間真有上主,為何不顯現讓我能夠信靠呢?」上主不負有心人,夜空彷彿有聲音在他心中回應:「你已經尋找了一段日子,為何仍猶豫不決呢?」他就把這一刻當作上主向他顯現,於是他做了一個信靠上主的決定,把自己交托給上主.回家後,他打電話告訴同學他成為基督徒了.他們問他是不是剛剛參加佈道會,他說只是在回家的路上做了這個決定.

三．你向上主顯現自己嗎?

剛才讀福音經課時,我們說耶穌向一個污靈附身的人顯現,其實這種說法值得商榷.事實上,耶穌向所有人顯現,為甚麼只有那位污靈附身的人有回應呢?從另一個角度來看,耶穌的顯現逼使這個人同時向耶穌顯現自己.當上主顯現時,如果有人同時向上主顯現自己,改變就有可能發生.對於沒有向上主顯現自己的人,沒有甚麼發生是可以理解的.客觀上來說,上主常在；但經驗上來說,人與上主互相顯現是一個特殊的時候,是一種契機,而未必是常態.有流行歌唱「神的不在場證明」,其實在日常生活經驗中,神不在場才是常態.然而,人與上主那種特殊的互相顯現時刻,對靈性生命來說卻是關鍵性的.雖然是不期而遇,但卻是可以重覆的,每一次獨特而深刻的互相顯現,都可能為靈命成長帶來突破和更新.

還記得那位把那個晚上的經驗當作上主向他顯現的年輕人嗎?那次顯現帶給他改變,因為他把自己交托給上主.後來,他順利進入大學,就在展開人生新一頁的時候,他遇到了十幾年來最嚴重的挫折.那個晚上,他又遙望長空,自慚形穢,感到自己一無是處,對自己非常失望.他問上主,為甚麼自己如此不濟?如此無用?如此令人討厭?憑甚麼再站起來呢?那一次,當他把最不濟的自己向上主顯露的時候,上主又不負他.他彷彿感到上主告訴他:「其實你有多不濟,有多令人討厭,有多少缺點和不足,我一早就知道,而且比你知道得更清楚.但是,我都照樣愛你,為甚麼你不能愛自己呢?」他感到震撼.然後,使徒保羅的說話在他心中響起:「惟有基督在我們還作罪人的時候為我們死,上帝的愛就在此向我們顯明了.(羅馬書5:8)同時,耶穌對彼得說的話也浮現在他的腦海中:「但我已經為你祈求,使你不至於失了信心.你回頭以後,要堅固你的弟兄.」(路加福音22:32)那一晚,他又再一次把自己交托給上主.

四．總結

有人以為舉手決志,獲取了天堂的入場券票就可以「收工」.其實,那不是結束,而只是開始.有人認為基督徒的生命是靈修旅程.靈修是甚麼呢?是每天早上讀一段經文,唱一首詩歌,祈一次禱嗎?一位研究靈修的神父葛羅謝爾(Fr. Benedict Groeschel)在其著作《煉淨光明合一:靈性成長心理學Spiritual Passages: The Psychology of Spiritual Development》一書中如此定義靈修:

「或許對靈修生活最好的描述是:完全地回應上主的呼召.然而,靈修生活並不是把自己關起來,它是融合在人的生活各種複雜的行為中,一連串不斷成長且協調的回應.把靈修生活和一個人日常生活分離是自古以來不斷發生的錯誤,當一個人決定回應內在經驗到的上主的呼召,並努力地讓這呼召成為行動和抉擇的中心,他或她才真正堪稱過著靈修生活的人.於是,這個呼召便成了使回應者得以整合的要素,靈修生活也就成為一生的事情.」

在靈修路上,靈修者偶爾經驗與上主互相顯現的特別時刻,這些時刻,有甚麼重要性呢?葛羅謝爾神父以焦慮為例子如此解釋:「由奧古斯丁都佛落依德,在他們之前還有使徒保羅,偉大的心理學家們認知到人類內心有種可怕的焦慮,對於徒勞無功和勞苦困頓的恐懼,由嬰孩時期開始便縈繞不去.當我們意識到自己人格中的靈性成分,我們日益察覺到對於生命超越層面的意義的可能,而這種覺察開始讓焦慮得以緩和.」原來,當靈修者向上主顯露他的焦慮的時候,上主向他顯現祂的靈性,這種互動就有機會幫助他以靈性來超越焦慮.

葛羅謝爾神父提醒,靈修路上沒有武林至尊,以為自己以「高人」的,可以是危險的.在靈修路上越有成就的人,就越懂得用高言大志和複雜細緻的理論來解釋自己的行動.葛羅謝爾神父指出,那些「高人」以為自己已經克服了自我中心的傾向,以為自己已經到達不受約束的地步,實際上,他可能正在放縱自己,並且自欺欺人.葛羅謝爾神父沒有指名道姓,讀者可能想起著名神學家約翰尤達(John Howard Yoder, 1927-1997).

尤達性侵事件發生於上世紀七十至八十年代,當時他以研究一種嶄新的基督教性倫理學說為名,性侵害逾百女性.尤達任教的神學院院長收到不少投訴,有些來自尤達家人.與此同時,尤達事奉的教會也關注事件,院長和牧師力勸尤達十多年,也是徒勞無功.因為尤達「德高望重」,辯才了得,跟他們拉布,耍花槍,帶他們遊花園,令處理事件的人士身心俱疲.差不多二十年之後,教會改變以前為挽救他們的明星神學家,以及家醜不出外揚的處理手法,對尤達的性罪行做一個全面徹底的調查.

耶利米先知曾經說過:「人心比萬物都詭詐,壞到極處,誰能識透呢?」(17:9)然後他回答說:「耶和華是鑒察人心,考驗人肺腑的」.在教會年曆,現在是顯現節期,耶穌向我們顯現,但願我們也向耶穌顯現我們自己,讓互相顯現的時刻成為我們更新及成長的契機.



\newpage

\section{以利亞已經來了}
\label{sec:684}
\textbf{李均熊牧師 - 以利亞已經來了}
\newline
\newline
\hyperref[sec:683]{< < < PREV SERMON < < <}
~
\hyperref[sec:toc]{[返目錄]}
~
\hyperref[sec:685]{> > > NEXT SERMON > > >}


馬可福音 9:2-13
\newline
\begin{longtable}{cl}
\hline
\hline
章節 & 經文 (和合本修訂版)\\
\hline
9:2 & \begin{tabularx}{0.7\textwidth}{X} 過了六天，耶穌帶著彼得、雅各、約翰，領他們悄悄地上了高山。他在他們面前變了形像， \end{tabularx} \\ \\ \relax
9:3 & \begin{tabularx}{0.7\textwidth}{X} 衣服放光，極其潔白，地上漂布的人沒有一個能漂得那樣白。 \end{tabularx} \\ \\ \relax
9:4 & \begin{tabularx}{0.7\textwidth}{X} 有以利亞和摩西向他們顯現，並且與耶穌說話。 \end{tabularx} \\ \\ \relax
9:5 & \begin{tabularx}{0.7\textwidth}{X} 彼得對耶穌說：「拉比，我們在這裡真好！我們來搭三座棚，一座為你，一座為摩西，一座為以利亞。」 \end{tabularx} \\ \\ \relax
9:6 & \begin{tabularx}{0.7\textwidth}{X} 彼得不知道說甚麼才好，因為他們很害怕。 \end{tabularx} \\ \\ \relax
9:7 & \begin{tabularx}{0.7\textwidth}{X} 有一朵雲彩來遮蓋他們，又有聲音從雲彩裡出來，說：「這是我的愛子，你們要聽從他！」 \end{tabularx} \\ \\ \relax
9:8 & \begin{tabularx}{0.7\textwidth}{X} 門徒連忙向周圍觀看，不再看見任何人，只見耶穌同他們在一起。 \end{tabularx} \\ \\ \relax
9:9 & \begin{tabularx}{0.7\textwidth}{X} 下山的時候，耶穌囑咐他們說：「人子還沒有從死人中復活，你們不要把所看到的告訴人。」 \end{tabularx} \\ \\ \relax
9:10 & \begin{tabularx}{0.7\textwidth}{X} 門徒將這話存記在心，彼此議論「從死人中復活」是甚麼意思。 \end{tabularx} \\ \\ \relax
9:11 & \begin{tabularx}{0.7\textwidth}{X} 他們就問耶穌：「文士為甚麼說以利亞必須先來？」 \end{tabularx} \\ \\ \relax
9:12 & \begin{tabularx}{0.7\textwidth}{X} 耶穌說：「以利亞的確先來復興萬事。經上不是指著人子說，他要受許多的苦和被人輕慢嗎？ \end{tabularx} \\ \\ \relax
9:13 & \begin{tabularx}{0.7\textwidth}{X} 我告訴你們，以利亞已經來了，他們任意待他，正如經上指著他說的。」 \end{tabularx} \\ \\
[1ex]
\hline
\hline
\end{longtable}
講題:以利亞已經來了 Elijah Has Come

講員:李均熊牧師

經文:馬可福音9:2-13

場合:香港中文大學崇基學院禮拜堂主日崇拜

日期:2024年2月11日

可9章

耶穌改變形像

2過了六天,耶穌帶著彼得,雅各,約翰暗暗地上了高山,就在他們面前變了形像,3衣服放光,極其潔白,地上漂布的,沒有一個能漂得那樣白.4忽然,有以利亞同摩西向他們顯現,並且和耶穌說話.5彼得對耶穌說:「拉比,我們在這裡真好!可以搭三座棚,一座為你,一座為摩西,一座為以利亞.」6彼得不知道說甚麼才好,因為他們甚是懼怕.7有一朵雲彩來遮蓋他們；也有聲音從雲彩裡出來,說:「這是我的愛子,你們要聽他.」8門徒忽然周圍一看,不再見一人,只見耶穌同他們在那裡.9下山的時候,耶穌囑咐他們說:「人子還沒有從死裡復活,你們不要將所看見的告訴人.」10門徒將這話存記在心,彼此議論「從死裡復活」是甚麼意思.11他們就問耶穌說:「文士為甚麼說以利亞必須先來?」12耶穌說:「以利亞固然先來復興萬事；經上不是指著人子說,他要受許多的苦被人輕慢呢?13我告訴你們,以利亞已經來了,他們也任意待他,正如經上所指著他的話.」

起:為何以利亞要來?

今日是年初二,中國人的農曆新年第二天.先恭祝大家新年快樂,龍馬精神,主恩常在,喜樂平安.

今日也是教會年曆的顯現期最後一個主日,也是紀念耶穌登山變像,顯現榮耀的主日.下星期就進入預苦期／大齋期,今天也是我們為面對基督的受苦而作出預備.

耶穌登山改變形像,顯得潔白發光,充滿榮耀,又有代表舊約兩大傳統的摩西和以利亞顯現.可是,在場的門徒「不知道說甚麼才好,因為他們甚是懼怕」.最後雲彩遮蓋他們,有聲音說「這是我的愛子,你們要聽他」,那是在馬可福音所記載耶穌的生平,第二次有天上聲音宣告耶穌是「我的愛子」(第一次是耶穌受洗時).在下山時耶穌和他一同登山的彼得,雅各和約翰說到人子還沒從死裡復活,就不要將所看見的告訴別人.經文卻仍有後續:

門徒將這話存記在心,彼此議論「從死裡復活」是甚麼意思.11他們就問耶穌說:「文士為甚麼說以利亞必須先來?」12耶穌說:「以利亞固然先來復興萬事；經上不是指著人子說,他要受許多的苦被人輕慢呢?13我告訴你們,以利亞已經來了,他們也任意待他,正如經上所指著他的話.」

為甚麼摩西和以利亞會顯現,為甚麼之後門徒會提起以利亞必須先來?舊約經課提到以利亞直接升天,不見死亡,所以猶太人在後來都相信他要再來,特別在瑪拉基書4章這樣說以利亞:

「看哪,耶和華大而可畏之日未到以前,我必差遣先知以利亞到你們那裡去.6他必使父親的心轉向兒女,兒女的心轉向父親,免得我來咒詛遍地.」

承:猶太人的盼望

其實記載以利亞生平的列王記中的以利亞,和新約時期猶太人心目中的以利亞,是很不一樣的.若仔細看列王記的經文,不受我們新約教會福音書對以利亞正面形像影響,其實以利亞是個自以為是,自把自為,不順從神旨,甚至情緒不穩定,令人不知所措,不被信任,所以孤僻離群,也不放別人在眼裡的問題人物來的!

但這可能是以色列人亡國初期,按申命記內中概念而寫成的舊約歷史書一貫對先知的不信任.先知的確是宣告神旨的人,但他們是否值得信任,根據申命記的吩咐,是在乎他所言是否實現(申18:22),也不是任何人行神蹟傳異夢的就是先知(申13:1-5).所以在列王記中出現的先知好壞參半,例如王上13章另一處說到一位老先知如何欺騙一名神人以至神人被獅子咬死的記載可見一斑.

但在猶太人流散後期,可能是先知言論集結,成書和流傳,先知的地位變得十分重要,從波斯時期成書的歷代志的最後評語(代下13:15-16, 21-22),以色列國滅亡,以色列人流散,並不是不聽申命記上主的吩咐,而是不聽如耶利米等先知的警告,而引至上主的忿怒向他的百姓發作.

以利亞作為神偉大先知,要來復興萬事,作上主的彌賽亞先鋒,便在希羅時期猶太人的著作中甚為顯著.在猶太人的著作中,以利亞不單是彌賽亞先鋒,更是滿有智慧的律法師傅,有時會幻化成不同人物如白髮長者或妓女,為人排難解紛,揭示天上的秘密,甚至以利亞會親臨在每次猶太人的割禮,或在逾越節和猶太人一起共進逾越節的晚餐.

所以以利亞不單是我們熟悉的舊約先知,更是猶太人的守護天使,協助猶太人重建家國,成就神的旨意.對我們來說,以利亞使人看見盼望的實現,是成就盼望的先驅,是守護人心中盼望的使者.

轉:得見耶穌的榮耀——為悲傷作預備 (進入預苦期)

在耶穌進入受苦之前,門徒能得見耶穌登山變像的榮耀,意識摩西和以利亞對彌賽亞企盼的守護,是十分重要的.

當面對耶穌的受苦受死,可以想像到門徒心裡面一定充滿疑惑,甚至乎會感到失望.怎樣在失望,甚至乎在崩潰的邊緣,仍能保持盼望,正因為他們預先見過耶穌的榮耀,知道摩西和以利亞會守護這個盼望.耶穌登山變像之後,門徒仍在思想「從死裡復活」的意思,不正好是為耶穌受死的時刻來臨的時候,幫助門徒先有一個概念.在驚訝,恐懼 ,慌亂 ,失望的同時,能仍有一絲的盼望, 好像 書信經課 之後的經文:

6那吩咐光從黑暗裡照出來的上帝,已經照在我們心裡,叫我們得知上帝榮耀的光顯在耶穌基督的面上.7我們有這寶貝放在瓦器裡,要顯明這莫大的能力是出於上帝,不是出於我們.8我們四面受敵,卻不被困住；心裡作難,卻不致失望；9遭逼迫,卻不被丟棄；打倒了,卻不致死亡.

合:在受苦時不至失望

盼望是危險的東西,盼望更是脆弱的東西(hope is a dangerous thing, yet hope is even more fragile ).

今日有甚麼幫助我們仍有盼望?

我們基督徒見證基督復活的大能,是否也能靠著聖靈的幫助,在充滿失望不安的時代,守護上主用愛同行釋放我們綑鎖的期望?

新春祝賀 促進盼望



\newpage

\section{重返伊甸}
\label{sec:685}
\textbf{鄧瑞強博士 - 重返伊甸}
\newline
\newline
\hyperref[sec:684]{< < < PREV SERMON < < <}
~
\hyperref[sec:toc]{[返目錄]}
~
\hyperref[sec:686]{> > > NEXT SERMON > > >}


講題:重返伊甸 Back to Eden

經文:雅歌4:1-7,12-16,5:6-8,10-16,8:5-7

講員:鄧瑞強博士

場合:香港中文大學崇基學院禮拜堂主日崇拜

日期:2024年2月18日

各位弟兄姊妹,早安.願神的恩典常與你們同在.今天是年初九,還算是新春的日子,祝願各人身體健康,心靈平安.

今天的講道,講《雅歌》,這是聖經「智慧文學」的最後一卷書了.智慧文學不祈求神蹟,不理會超越的世界,而是腳踏實地,就著可見的現實,努力活一滿足的人生.我們討論過,《約伯記》探討苦難的問題,苦難的來源無法可知,人能作的,是在苦難中不去懷疑神的美善；生命的正直,也不由福與禍去決定.《箴言》探討在這善惡混雜的世界裡,如何以敬畏耶和華的心去持守智慧,並以智慧去開導愚昧.《傳道書》探討人生的虛空,提示我們不要執著變幻的世物,要將人生看成是神的禮物,懷著感恩的心去接受,也懷著分享禮物的心去送贈.白白地得來,也白白地捨去.如此,過一自由而愉悅的人生.如此這般,我們探討了幾個人生的大課題:苦難的無法理喻與忠信的難能可貴,愚昧的無藥可救與智慧的春風化雨,執著的徒勞無益與放下的消遙自在.今日,透過《雅歌》,我們探討另一個人生大課題:愛情.

我們先看看《雅歌》這愛情故事的場景.若你回家細心閱讀《雅歌》,你會發現,每一章,都有大量香花芳草.香柏樹,乳香岡,沒藥,沉香,香花畦,香草臺,香草山,四處都是.若你懂得聞一聞經文,每一頁,都滲透出幽幽的清香,瀰漫你的思緒,充實你的靈魂.由於時間關係,我們無法一一細閱.我們看其中一段:《雅歌》4:12-16,每句經文,不是花香,就是流水.

歌4:12 我妹子,我新婦,乃是關鎖的園,禁閉的井,封閉的泉源.

歌4:13 你園內所種的結了石榴,有佳美的果子,並鳳仙花與哪噠樹.

歌4:14 有哪噠和番紅花,菖蒲和桂樹,並各樣乳香木,沒藥,沉香,與一切上等的果品.

歌4:15 你是園中的泉,活水的井,從黎巴嫩流下來的溪水.

歌4:16 北風啊,興起!南風啊,吹來!吹在我的園內,使其中的香氣發出來.願我的良人進入自己園裡,吃他佳美的果子.

簡單一看,有美人,有多層次的花香樹香,有靜靜的流水,有愛情,是否已美得無話可說?妹子,新婦,就是美人.「關鎖的園,禁閉的井,封閉的泉源」就是美人的貞潔.各式各樣的香花芳草,就是美人的高潔清雅.從男子的角度看,這女子就是「園中的泉,活水的井,從黎巴嫩流下來的溪水」.以色列地是熾熱乾旱之地,黎巴嫩有雪山,流下來的冰水,真讓人心曠神怡.「願我的良人進入自己園裡,吃他佳美的果子」,這是美人對男子說的,經文的意義充滿禪意,只能意會,不能言傳.

以芳草喻美人,中國的大詩人屈原是個中高手.其《九歌•山鬼》的開首便是:「若有人兮山之阿,被薜荔兮帶女蘿」.美人,在山的深處,身披薜荔,腰束女蘿.薜荔,女蘿,都是香草.又說:「山中人兮芳杜若,飲石泉兮蔭松柏,君思我兮然疑作」.美人像杜若一樣芬芳,飲石泉之水,居松柏之間,何等高雅脫俗.「君思我兮」,美人猜想著愛人是否思念著自己,「然疑作」,「然」者,確是這樣；「疑」者,未必是這樣.「作」者,產生這些念頭.「君思我兮然疑作」,美人心神不寧,一時覺得對方思念著她,一時覺得對方忘記了她.在這齷齪的塵世,高潔的美人孤獨異常.

我想,若由屈原去翻譯《雅歌》,則《雅歌》將會是一篇何等絕倫的美文?

回到《雅歌》,美人與她的愛人,香花,芳草,流水,花園,這根本就是伊甸園的畫面.自從人類離開了伊甸園之後,在人類的夢裡,在詩人的想像裡,在聖人的祈禱裡,常常回到這靈魂的故鄉.在這齷齪的塵世,我們渴慕那純潔的淨土.

世上有不少關於美人的藝術品,其中一件是《米洛的維納斯》(Venus de Milo),一座斷了臂,半祼的女神維納斯的大理石雕塑,展現出女性美麗的胴體.很多小型的石膏擺設,以此為模型.另一件藝術品是波提切利(Botticelli)的畫作《維納斯的誕生》(The Birth of Venus),不少人認為這畫作是世上十大名畫之一,全裸的女神維納斯在一巨型的貝殼中誕生.左邊的風神將鮮花吹向她,右邊的春天女神準備為她蓋上繡滿花朵的斗篷,中間的維納斯純潔地凝望著遠方.我講了這麼多,我相信你們能欣賞裸體所呈現的純真的美.

現在,我們看看《雅歌》以文字描繪的裸體——《雅歌》4:1-7.

歌4:1 〔原文有「看」〕我的佳偶,你甚美麗!〔原文有「看」〕你甚美麗!你的眼在帕子內好像鴿子眼〔原文沒有這「眼」字〕.

〔你的眼,露在紗外,像鴿子.根據《雅歌》2:14的講法,鴿子藏身在陡峭的岩洞上,好美,但看不到全貌.意指:美人的面藏在面紗裡,晶瑩的眼睛吸引人想看到她隱藏著的美貌.〕

你的頭髮如同山羊群臥在基列山旁.〔黑山羊其實不是「臥在」基列山旁,而是緩緩地從山上走下來,意指:美人的頭髮一卷一卷地垂下來,從頭直到肩膊.〕

歌4:2 你的牙齒如新剪毛的一群母羊,洗淨上來,個個都有雙生,沒有一隻喪掉子的.〔牙齒白而整齊.〕

歌4:3 你的唇好像一條朱紅線；你的嘴也秀美.你的兩太陽在帕子內,如同一塊石榴.〔「兩太陽」即面頰,紅撲撲的,像擘開為兩半的石榴.正是:唇紅齒白,兩頰緋紅.〕

歌4:4 你的頸項好像大衛建造收藏軍器的高臺〔頸挺而直〕,其上懸掛一千盾牌,都是勇士的籐牌〔掛著一環又一環的項鏈〕.

歌4:5 你的兩乳好像百合花中吃草的一對小鹿,就是母鹿雙生的.〔小鹿,嫩草,百合花,給我們一種芬芳而柔軟的感覺.〕

歌4:6 我要往沒藥山和乳香岡去〔沒藥山,乳香岡,是美人的身體外形〕,直等到天起涼風,日影飛去的時候回來〔男子要登沒藥山,乳香岡,直到清晨涼風吹拂的時候〕.

歌4:7 我的佳偶,你全然美麗,毫無瑕疵!

欣賞完女性的美,我們看看男性的身體——《雅歌》5:10-16.

由於時間關係,我不仔細解釋了.基本上,是展示一精壯的身體.

歌5:10 我的良人白而且紅,超乎萬人之上.

歌5:11 他的頭像至精的金子；他的頭髮厚密纍垂,黑如烏鴉.

歌5:12 他的眼如溪水旁的鴿子眼〔原文沒有這「眼」字〕,用奶洗淨〔眼白很白〕,安得合式.〔即是指眼睛黑白分明〕

歌5:13 他的兩腮如香花畦,如香草臺；他的嘴唇像百合花,且滴下沒藥汁.

歌5:14 他的兩手好像金管,鑲嵌水蒼玉；他的身體如同雕刻的象牙,周圍鑲嵌藍寶石.

歌5:15 他的腿好像白玉石柱,安在精金座上；他的形狀如黎巴嫩,且佳美如香柏樹.

歌5:16 他的口極其甘甜；他全然可愛.耶路撒冷的眾女子啊,這是我的良人；這是我的朋友.

《雅歌》純潔地展示身體與性,是要教我們如何在「日光之下」過一整全而滿足的人生.

《雅歌》教我們正視我們是「身體與心靈」的統合體(psychomatic unity).身體沒有本然的邪惡,心靈沒有本然的聖潔.邪惡與聖潔,在乎「身體與心靈」這統合體是否敬神和愛人.敬神和愛人,是心靈的事,也必須透過身體來表達.正如最高貴的神人聯合,也要透過吃一塊聖餐餅,喝一杯聖餐酒這身體行動來領會.人與人的愛,不能不透過身體的親密接觸來表達.在婚姻關係裡,夫婦的性愛,是以身體的連合去表達心靈的愛,其實,這是最富有靈性的行動.

教宗若望保祿二世(John Paul II)寫了一本巨著《男和女:天主的傑作—身體神學》(Man and Woman He Created Them: A Theology of the Body),他說,由於「道成了肉身」,我們便能合法地透過「肉身」去了解「道」.他用這原則去理解神對人的創造.神造人之初,認為人「獨居不好」,為何不好?因為人的身體與其他動物的身體不同,無法有共同「體」驗,無法分享生命,因此,人感到無比孤獨.於是,神為人造一個配偶,造出男與女.人有身體,身體有性的欲望,性的欲望驅使人走向他人.因身體可以接觸,性的接觸讓人與他人有生命的融合.性的融合,按若望保祿二世的講法,有聖禮的性質.身體的合一,是一可見的神聖標記,表達出上帝內在的不可見的,位格的愛.男與女的身體融合,作為聖禮,讓人類參與上帝的愛的奧祕,藉這聖禮,男與女共同展現出「神的形象」.換句話說,透過身體的性與愛的結合,上帝不可見的愛具體地呈現人間.

這是很弔詭的.我們常以為,性和身體,是阻礙人接近神的元凶.事情剛好相反,正因性和身體,人才能體會上帝最內在的愛,也將這愛具體展現.想想,若我們沒有身體,沒有性,則我們只是一個又一個孤立的靈魂,不可能有生命的接觸和融合,也無法透過身體呈現愛.身體和性,令人更明白神.

如此,我們就能將身體,性(sexuality),靈性(spirituality)結合起來.真正的性愛,是基於雙方真誠地接納對方.接納對方,始於接納對方的身體,也敢於向對方展示自己的身體.視身體如禮物,送贈對方；也視對方的身體,是愛的餽贈.這愛的交流,有付出,有犧牲,也有愉悅和滿足.性的融合,需要人尊重對方,也學習敢於向對方展示自己的私隱,這是位格與位格相遇的契機.在其中,人學習愛,也學習被愛,這愛帶來身體上的滿足,也帶來人格上的圓滿.人能在愛中,活出生命的整全,活現上帝的內在生命,這就是靈性.

令人整全的性接觸,具靈性意義.反過來說,靈性上的缺失,會造成聖的障礙.靈性上的缺失,各式各樣,如:恐懼他人,沒有安全感,不能完全接納他人,不能完全接納自己,不能完全相任他人,不能完全相任自己,不能完全愛惜他人或自己,,物化他人或自己,沒有足夠的愛與被愛的能力,自慚形穢,因不相任而來的羞恥,壓力過大,自我刑罰,或刑罰對方....... 林林總總,歸根究柢,從靈性上講,是不相信神對自己的無條件接納,也無法無條件地接納對方.靈性上的缺失,阻礙性的滿足；性的障礙,也影響靈性的圓滿.如何打破這惡性循環,需要對神,對人有更大的開放與信任.

性與靈性,有一種連繫性.有一個字眼,常常出現在性愛文學和靈修文學中,這個字眼就是「ecstasy」(神魂超拔).在性愛文學裡,這是指忘我的高潮.在靈修文學裡,這是指神人超越的聯合.若以存在主義的哲學來分析,「ecstasy」就是指人走出自己的生命圍牆,走向他人.「ecstasy」的前綴「ec」即「ex」,即「走出去」的意思,英文「exodus」就是「走出」埃及這塊「土地」的意思.「ecstasy」的「stasy」演變成英文的「stasis」,「stand」這些字,這些字都是「固定物」的意思.人的身體就是人的固定圍欄,性的欲望令人意圖走出圍欄,愛是雙方圍欄的開放,性的聯合便是成功走出圍欄,被接納進入他人的圍欄內,其高峰是「ecstasy」,真正走出自己而成功地走向他人,生命找到超越自己的圓滿意義.

講了一大堆話,讓我回到《雅歌》的經文.《雅歌》教我們正視身體,不要物化身體,不要善價而沽,而是將身體視為位格相遇的中介,透過身體,去接觸他人.《雅歌》教我們正視性,性驅使身體與身體融合,在純潔的性行為中,雙方接納對方那具有身體的人格,走出自身的圍欄,體會在關係中呈現的愛,這愛正是神的形象.當人在愛的關係中活出神的形象時,人才能找到生命的圓滿.

在伊甸園,夏娃與亞當違反神的禁令後,發現自己赤身露體,感到羞恥.《雅歌》卻展示,只要真愛在,夫婦能展露身體,彼此接納,彼此開放,並透過身體的連合,去體會神創造男與女的美意.在純淨的愛中,男與女能赤祼相對,並不羞恥.這是修復吃禁果後的羞恥,人在愛中,一步一步地重返伊甸.

社會心理學家佛洛姆(Erich Fromm)寫了一本書《愛的藝術》(The Art of Loving),他指出,人類懼怕不被人愛,但更懼怕去愛人,因為愛意味走出自己的圍牆,走出安全區,要冒險,要付出,要犧牲,卻可能得不到什麼回應,甚或被傷害.愛的行動,是向他人赤誠相向,若得不到回應,則只會令自己成了毫無保護的脆弱的人.人怕冒險,怕受傷,故現代多了很多將自己封閉在密室裡的孤獨靈魂.

《雅歌》講的愛,不是純肉慾的快感,《雅歌》描繪的愛,是走出圍牆,敢於冒險的愛.我們看《雅歌》5:6-8.

歌5:6 我給我的良人開了門；我的良人卻已轉身走了.他說話的時候,我神不守舍；我尋找他,竟尋不見；我呼叫他,他卻不回答.

歌5:7 城中巡邏看守的人遇見我,打了我,傷了我；看守城牆的人奪去我的披肩.

歌5:8 耶路撒冷的眾女子啊,我囑咐你們:若遇見我的良人,要告訴他,我因思愛成病.

這是美人的獨白,她尋覓她的愛人.《雅歌》很前衛,很女性主義,主動去愛,敢於犧牲的,是女性.這女子打開門,走出圍牆,越出安全區,四處尋覓她的愛人,她遇到壞人,被凌辱,受傷,被人奪去掩蓋身體的圍巾.縱使這樣,她仍堅持去愛.在這齷齪的塵世,高潔的美人孤獨異常.

在《創世記》伊甸園的故事裡,夏娃與亞當違反神的禁令後,他們隱藏自己,是神四處尋覓他們.整本聖經66卷書,就是神不斷尋覓罪人的漫長故事.神也走出祂的安全區,耶穌也要受人鞭打,最後犧牲在十架上.這是真愛的故事.在《雅歌》裡,這女子受盡苦難,為的是愛.故此,《雅歌》不是粉紅色地只描繪愛情的浪漫,也用紫色,黑色重筆地呈現愛情的苦難深度.當女子仿效上帝,冒險地去尋覓,受苦地去愛時,人類又再進一步接近伊甸園了.

最後,堅貞的愛獲得苦難後的成果.我們看《雅歌》8:5-7.

歌8:5 那靠著良人從曠野上來的是誰呢?〔這句話是旁白:倚偎著愛人,從艱難之曠野地上來的,是誰呢?〕

我在蘋果樹下叫醒你.你母親在那裡為你劬勞；生養你的在那裡為你劬勞.

〔這是美人的話,場景回到伊甸園,在蘋果樹下,美人叫醒她的愛人.叫醒他,做什麼呢?「你母親在那裡為你劬勞；生養你的在那裡為你劬勞」,是重複的,婉轉的說,你的母親在那裡受孕.美人口中說的,是男子的母親.言下之意,想著的,卻是自己.〕

歌8:6 求你將我放在你心上如印記,帶在你臂上如戳記.因為愛情如死之堅強,嫉恨如陰間之殘忍；所發的電光是火焰的電光,是耶和華的烈焰.

歌8:7 愛情,眾水不能熄滅,大水也不能淹沒.若有人拿家中所有的財寶要換愛情,就全被藐視.

在世上,死亡的力量,是無堅不摧的.在世上,也只有一種力量,能與之匹敵,就是愛情的力量.「愛情如死之堅強」.在人間,陰間的鍥而不捨,是無人能及的.在人間,也真只有一種力量,能與之相提並論,那就是愛情的力量.「愛情的嫉恨如陰間之殘忍」.在滅世的洪水下,無一物能存活,但,「愛情,眾水不能熄滅,大水也不能淹沒」.

在古老的伊甸園,當人類犯罪後,便得面對死亡的威脅.如今,在堅貞的愛情裡,人找到克服死亡的力量.經過艱苦的旅程,人終於重返伊甸了.當然,這只是「日光之下」人能獲得的成果,至於「日光之上」的永恆生命,就需要神的救贖去獲取了.

願上主得著榮耀,頌讚；願眾人得享愛情,福樂.新年快樂.



\newpage

\section{止於至善}
\label{sec:686}
\textbf{郭志榮牧師 - 止於至善}
\newline
\newline
\hyperref[sec:685]{< < < PREV SERMON < < <}
~
\hyperref[sec:toc]{[返目錄]}
~
\hyperref[sec:687]{> > > NEXT SERMON > > >}


馬太福音 5:17-20
\newline
\begin{longtable}{cl}
\hline
\hline
章節 & 經文 (和合本修訂版)\\
\hline
5:17 & \begin{tabularx}{0.7\textwidth}{X} 「不要以為我來是要廢掉律法和先知。我來不是要廢掉，而是要成全。 \end{tabularx} \\ \\ \relax
5:18 & \begin{tabularx}{0.7\textwidth}{X} 我實在告訴你們，就是到天地都廢去，律法的一點一畫也不能廢去，直到一切都實現。 \end{tabularx} \\ \\ \relax
5:19 & \begin{tabularx}{0.7\textwidth}{X} 所以，無論誰廢掉這誡命中最小的一條，又教導人也這樣做，他在天國裡要稱為最小的。但無論誰遵行並如此教導人的，他在天國裡要稱為大。 \end{tabularx} \\ \\ \relax
5:20 & \begin{tabularx}{0.7\textwidth}{X} 我告訴你們，你們的義若不勝過文士和法利賽人的義，絕不能進天國。」 \end{tabularx} \\ \\
[1ex]
\hline
\hline
\end{longtable}
講題:止於至善 In Pursuit of Excellence

經文:馬太福音5:17-20

講員:郭志榮牧師

場合:香港中文大學崇基學院禮拜堂主日崇拜

日期:2024年2月25日

一．經文:太五17-20

17 莫想我來要廢掉律法和先知.我來不是要廢掉,乃是要成全.

18 我實在告訴你們,就是到天地都廢去了,律法的一點一畫也不能廢去,都要成全.

19 所以,無論何人廢掉這誡命中最小的一條,又教訓人這樣做,他在天國要稱為最小的.但無論何人遵行這誡命,又教訓人遵行,他在天國要稱為大的.

20 我告訴你們,你們的義若不勝於文士和法利賽人的義,斷不能進天國.

二．前言

弟兄姊妹,主內平安,新年蒙恩!

回到崇基禮拜堂這個熟悉的地方,感覺到很親切.當年我就是在這裡讀書,也是在那個時候有朋友帶我返教會信主,轉眼間已經過了三十五年.

當年初信的時候,就有人教我,要讀懂和讀通《約翰福音》和《羅馬書》.我當時選擇了先讀《羅馬書》,因為我是一個好喜歡抽象思考的人,而《羅馬書》當中就包括了很多金句,都是一些重要的信仰命題,而保羅也在這書信中一步步推論,所以「arm 哂我」.特別是開頭八章我讀了好幾次,直到我以為自己已經明白為止.這給我一個印象(現在知道是錯誤的),就是基督教的信仰是強調福音,強調因信稱義而貶低甚至否定律法.

帶著這個誤解,當讀到〈登山寶訓〉,發現耶穌說祂要「成全律法」的時候,就感覺到好疑惑.《羅馬書》不是說人根本不能倚靠律法和行為稱義,為何耶穌說「律法的一點一畫」都要「成全」,「不能廢去」呢?是保羅曲解了耶穌的教導?當然不會.

另一方面,從《福音書》的記載中可以看到,耶穌的確似乎有做過一些不遵守《舊約》律法規定的事情,例如在安息日治病,又例如,不按照律法要求的方式洗手.[1]既然耶穌似乎也沒有嚴格遵守律法的規條,祂又怎樣把律法「成全」呢?

坊間聽過的一個解釋,就是說因為《馬太福音》是寫給猶太人看的,所以便加入這些內容.[2]但這樣解釋,是否意味著猶太人信的福音和外邦人所信的是不同的嗎?當然不是!就算猶太人和外邦人在信仰的實踐上因應文化不同而有分別,但信仰最核心的內容也不應互相矛盾.

而假若繼續再細閱〈登山寶訓〉接下來的經文時,大家更會發覺耶穌提出的要求——例如不僅不可殺人和不可姦淫,連向弟兄發怒和動淫念也不可以；又例如有人打你的右臉,連左臉也要轉過來由他打——根本極難完全實現.曾經有一位未信的朋友就問過我,這怎麼可能做得到?信耶穌的要求實在太高了.是否真的如此?

耶穌說要「成全律法」,應該怎樣理解呢?讓我們今天去看看這段經文,希望讓大家對耶穌「成全律法」的教導有更深的理解.

經文來自《馬太福音》第五章17-20 節,今日主要講第17 和第20 節.經文已經顯示在電幕上,讓我們一同誦讀.

三．釋經

A. 為何耶穌要提出「成全律法」這個講法?

要解釋這段經文,第一個要處理的問題,就是為何耶穌要提出祂來是要「成全律法」呢?祂這麼說,到底要回應甚麼問題呢?

耶穌在17節中一開始就說:「莫想我來要廢掉律法和先知.」耶穌這麼說,明顯暗示有人會認為祂的出現是要「廢掉律法和先知」.

首先,「律法和先知」指甚麼?在耶穌的年代,《舊約》的正典雖然仍未完全定下來,但已有基本的範圍.大體而言,「律法和先知」這個組合就是指整卷《舊約》——「律法」指《摩西五經》,是猶太拉比心目中《舊約》經卷的核心,而「先知」就是泛指其餘的經卷.[3]

耶穌澄清,祂的到來所要成就的,並不是要「廢掉」,即不是要「破壞」,「拆毀」和「放棄」《舊約》所寫的內容.[4]問題是:耶穌要回應誰?誰會誤會或指控耶穌要廢掉「律法和先知」呢?答案可以在結尾的第20節中找到提示,這些人就是「文士和法利賽人」.

B. 甚麼是「文士和法利賽人的義」?

在第20節中提到,耶穌向跟隨祂的人說:「你們的義若不勝於文士和法利賽人的義,斷不能進天國.」那麼,甚麼是「文士和法利賽人的義」呢?[5]

「文士」就是那些仔細研究和解讀「律法」(即《摩西五經》)的「專家」.這些「專家」當時把「律法」的教導,拆解和演繹成為可以具體應用到各種生活處境的誡命規條,共有613 條.[6] 這613 條誡命規條極之繁瑣,條文盡可能清楚地訂明,在不同的情況下,一個人應該做以及不能做的事情.例如,在安息日原則上不可以醫治病人,除非病人有生命危險；而且同時訂明,遇上有生命危險的病人時,醫生只能做最基本的搶救,保住病人的性命,而不能在安息日把病人治癒.[7]

那麼,「法利賽人」呢?「法利賽人」的意思就是「分離者」,[8]意思就是嚴守這些誡命規條而與一般人分別出來的人.[9]

「文士和法利賽人」認為,在這613 條的誡命規條中,一個人如果能夠符合到超過一半,即遵守的比違犯的多,便算是「義人」,否則被視為「不義」.這就是「文士和法利賽人的義」.

耶穌對「文士和法利賽人」如此演繹「律法」的教導不以為然,自然也招致「文士和法利賽人」的評擊,指控耶穌要「廢掉律法」.

C. 「文士和法利賽人的義」有甚麼問題?

「文士和法利賽人」把生活的準則巨細無遺地定得清清楚楚,讓大家可以跟著條文實踐,不好嗎?有甚麼問題呢?

人常常有傾向,想把規矩定得清楚一點.就以我所事奉的堂會為例,要求崇拜的服飾要莊重整潔.間中就會有弟兄姊妹問,可否定一個詳細一點的指引.既然人有這樣的傾向,「文士和法利賽人」的做法有甚麼問題呢?

大家先要弄清楚一個更基本的問題:「律法」的性質是甚麼?

「律法」就是神與人立約的條款,是在生命約下如何實踐信仰的指引.這個指引包含了兩個向度,就是怎樣與神和與人相處,而「愛」——或更準確地說,是「行動之愛」就是與神和與人相處的最高原則.所以,聖經教導大家,「愛神」和「愛人如己」兩條誡命就是懸掛起整個「律法」的「總綱」,也就是說,這兩條「愛」的誡命是詮釋一切別的誡命的前設或依據.

「文士和法利賽人」對「律法」繁瑣的演繹,把焦點錯放在禮儀,以及一切實際上無關痛癢的生活和行為規條之上,完全沒有「導人向愛」,沒有讓遵守這些禮儀和規條的人感知到神的「行動之愛」,更遑論要把這「行動之愛」實踐出來.「文士和法利賽人」就這樣偏離了兩條「律法的總綱」,也就是離開了對「律法」正確的詮釋和引伸.這樣的信仰流於空洞,沒有靈性,因為他們所跟隨的只是自己所訂下來的一堆規條,死的文字,而並非跟隨著一位有位格,有情意的永活之道.

另一方面,「文士和法利賽人」的做法表面上很嚴格,但實質上是把「律法」的要求降低了.

有神學家提出一個好精彩的見解,就是「文士和法利賽人」把「律法」本來更重要和更高的要求擱在一旁,而以他們自己演繹出來的誡命規條取而代之.雖然這些條文極之繁瑣,但卻使「律法」變得更為容易執行.試想想,要做到「愛人如己」困難?還是不在安息日按升降機困難?當然是前者.所以,這其實是一種變相「折衷／讓步」(compromise)的做法.[10] 《馬太福音》第五章接下來就記載了耶穌一連串的教導,把天國應有的要求,和那些「古人的話」——亦即「文士和法利賽人」的教導——作對比,正好說明了他們的「讓步」.

再者,「文士和法利賽人」也並非要求人要符合全部613 條誡命規條,而是只需要做到不少於一半便可以,於是就更進一步降低「義」的門檻.耶穌說要勝過「文士和法利賽人」的義這個門檻才能「進天國」,意思就是說:一個人若只達到「文士和法利賽人的義」根本就還沒有真正遵行神的旨意,亦即還沒有降服和活在神的主權之下,當然仍然留在天國之外.

D. 耶穌怎樣「成全律法」?

解釋了勝過「文士和法利賽人」的義之後,便可以回到一開始提出的問題:耶穌說要「成全律法」,到底應該怎樣理解呢?

「成全」這個字的基本意義就是「填補」,「倒滿」的意思,[11]引伸出來,就可以解作「完成」,或者更貼切解作「圓滿」.[12]而「圓滿」這個動詞的希臘文形式,顯示「圓滿律法」是一個一次性的行動,[13] 也就是耶穌「來」所要完成的事.從上述的討論看,要勝過「文士和法利賽人」的義,就是要從「律法」的詮釋,引伸和實踐三個角度去理解,而「圓滿律法」也包括了這三個層面的內涵.

1. 「律法」的詮釋

首先,「圓滿律法」的第一個層次,就是要正確地詮釋「律法」.耶穌作為神成了肉身的聖道甘願俯就卑微而「來」到世上,最後主動走上十字架上將自己的生命獻上,當中所彰顯出來的「行動之愛」,正正是正確地詮釋「律法」的關鍵.換言之,要判斷某一個對「律法」的詮釋是否正確,必要條件就是視乎這個詮釋是否與神的「行動之愛」相符,又或者是否指導人以「行動的愛」來愛神和愛人.

例如,「文士和法利賽人」經常評擊耶穌不守安息日的規條,但耶穌卻指出「安息日是為人設立的,人不是為安息日設立的.」(可二27)也就是說,神設立安息日不是為了限制人,規定人在安息日最多可以走多遠的路,或禁止煮食,而是要使人得著生命的好處,讓人在開展每週的工作前,可以先放下世務專心親近神,從而思想和預備好怎樣在未來一週中敬虔度日.

神為人設立安息日,本身就是神愛人的行動.人願意在安息日暫時放下世務,也是愛神,渴慕神的表現.而耶穌在安息日治病,使人的身體或心靈恢復過來,正是「行動之愛」的實踐,同時也讓那些得到醫治的和在場見證著事件發生的人,認識到神對人的大愛並因此而歸榮耀給神.

2. 「律法」的引伸

耶穌說祂要「圓滿律法」,除了要詮釋出「律法」的基本含義外,也要從這基本含義作進一步的引伸,使「律法」的教導所涵蓋的範圍就更全面,內容也更完備.

《馬太福音》的作者在五章21-48 節中,就分別用了六個人際之間相處的例子,來示範「律法」的引伸到底是怎樣的一回事.例如,由「不可殺人」引伸出「不可向弟兄動怒」的教訓.又例如,由「不可姦淫」進一步引伸出,其實連「看見婦女就動淫念」也是不可以的.

耶穌如此「圓滿律法」所引伸出來的教導,正好說明了神對人的要求極高.這是因為人是按著神的形像而受造,神當然要求人活出能反照出神形像的生命樣式.這正好與「文士和法利賽人」把「律法」要求降低的做法,形成強烈的對比.

3. 「律法」的實踐

耶穌不僅對「律法」作正確的詮釋和進一步的引伸,讓人明白和掌握「律法」真正的含義和要求,因為這樣並不足夠,並未「圓滿」.人就算完全明白了「律法」,即神的話語,或說聖經的所有教訓,也是行不出來的.因為人的罪性使人自我,自保,加上信心不足,明知是對的也不願實行.

例如,我堂會中負責婚前輔導的同工就教導夫婦,如果伴侶令自己非常生氣,自己正要爆發,忍不住要向伴侶表達自己的不滿時,可以在責罵對方的說話前先加上「你個傻豬豬」這個稱謂,整個氣氛就會逆轉過來.但問題是,雖然明知這樣說可以得到美好的果效,緩和緊張的氣氛,但夫婦有誰能真正約束自己的怒氣,並且願意放下自我,以愛平息紛爭?

耶穌深知人本性的限制,知道人就算已掌握「字句的話語」,也要與祂同行,並在親近祂的禱告過程中得力,這樣才能把「律法」的要求逐一實踐出來.耶穌「圓滿律法」的最後一個層次,就是要以祂在十字架上所賜下的救贖之恩,使願意降服和跟隨祂的人得著實踐「律法」的動力和能力.

以上三個層面結合起來,才是耶穌「圓滿律法」完整的意思.

四．引伸應用

最後,還有一個問題要進一步澄清,就是耶穌所「圓滿」的「律法」要求極高,到了一個程度好像信徒一生也無法完全百分百達到.那信徒應怎麼辦呢?這麼高的「律法」的要求有何作用?

一個傳統的觀點就是,「律法」就是使人知罪.大家可以把耶穌所「圓滿」的「律法」要求,視為「理想」(ideal)的典範.而人達不到這個「理想」的典範,就是要讓人知道自己跟神所訂下的標準相去甚遠,因而認識到自己有罪的真相,也認識到自己無力脫離罪牽引的事實,於是就知道自己需要尋求耶穌所成就的赦罪之恩,[14]因為若不是神的赦免,根本沒有人能逃避在神的審判下滅亡的厄運.

然而,這個講法只講述了事情的上半部分.有神學家就指出,一個信徒倘若把這些「律法」的要求淡化,單單視為只有讓人知罪和尋求神赦罪的功能,而並非真的要實踐出來,在這樣理解下的信仰只會是在追求「廉價的恩典」.這樣的信徒只希望神赦免自己,不在烈怒中審判自己,但卻沒有打算改變自己的生命品格和言行表現,更不用說要為事奉神而付上代價.[15]

沒錯,這些「律法」的「理想」典範,對於有限的人來說,很可能終其一身都無法完全達到.但即使是這樣,這些「律法」的要求,卻是跟隨神的人需要窮一生的精力去追求的「理想」目標.

換言之,今天一個信徒或許發覺自己與這個「理想」的典範仍有極大的距離,但耶穌為要「圓滿律法」而賜下的「成長之恩」,可以加力給他,幫助他朝向這個「理想」目標進發,因此知道只要願意倚靠聖靈的幫助,明天,後天,大後天,每天都可以再走近這個目標多一點.這就是一個沒有終結的,生命不斷成長的,邁向成聖與稱義的歷程.

總結一下,耶穌所說的「圓滿律法」,既是有要求,而且要求極高.但同時也是有恩典的,包括了容許人會犯錯的「赦罪之恩」,也包括「成長之恩」,幫助人去追求實踐「律法」的要求!

最後,當「文士和法利賽人」把「律法」演繹為誡命規條,總是會問某件事「做唔做得」.這是問錯了問題,因為「律法」中那個「理想」的典範其實是一個信徒要倚靠基督而不斷追求的目標.所以,要問的問題不是「做唔做得」,而是怎樣可以做得更好,更更好,更更更好.如果要說明這個追求過程的終點,或許可以借用我們崇基學院的校訓來表達,就是「止於至善」.

書目

Abernathy, David. An Exegetical Summary of Matthew 1–16, Exegetical Summaries.Dallas, TX: SIL International, 2013.

Arndt, William, Frederick W. Danker, et al. A Greek-English Lexicon of the New Testament and Other Early Christian Literature. Chicago: University of ChicagoPress, 2000.

Barclay, William. The Gospel of Matthew, Third Ed., The New Daily Study Bible.Edinburgh: Saint Andrew Press, 2001.

Boring, M. Eugene “The Gospel of Matthew,” in New Interpreter’s Bible, ed. Leander E. Keck, vol. 8. Nashville: Abingdon Press, 1994–2004.

Blomberg, Craig. Matthew, vol. 22, The New American Commentary. Nashville: Broadman  and  Holman Publishers, 1992.

Bromiley, Geoffrey W. The International Standard Bible Encyclopedia, Rev. ed. Wm. B. Eerdmans, 1979–1988.

Carson, D. A. “Matthew,” in The Expositor’s Bible Commentary: Matthew, Mark, Luke, ed. Frank E. Gaebelein, vol. 8. Grand Rapids, MI: Zondervan Publishing House, 1984.

Davies, W. D. and Dale C. Allison Jr., A Critical and Exegetical Commentary on the Gospel according to Saint Matthew, vol. 1, International Critical Commentary. London; New York: T and T Clark International, 2004.

Freedman, David Noel. The Anchor Yale Bible Dictionary. New York: Doubleday, 1992.

Hagner, Donald A. Matthew 1–13, vol. 33A, Word Biblical Commentary. Dallas: Word, Incorporated, 1993.

Hauerwas, Stanley. Matthew, Brazos Theological Commentary on the Bible. Grand Rapids, MI: Brazos Press, 2006.

Hendriksen, William and Simon J. Kistemaker, Exposition of the Gospel According to Matthew, vol. 9, New Testament Commentary. Grand Rapids: Baker Book House, 1953–2001.

Küng, Hans. On Being a Christian, trans. Edward Quinn. New York: Doubleday, 1984.

Luz, Ulrich. Matthew: A Commentary, ed. Helmut Koester, Hermeneia—a Critical and Historical Commentary on the Bible. Minneapolis, MN: Augsburg, 2001.

Morris, Leon. The Gospel according to Matthew, The Pillar New Testament Commentary. Grand Rapids, MI; Leicester, England: W.B. Eerdmans; Inter-Varsity Press, 1992.

Nolland, John. The Gospel of Matthew: A Commentary on the Greek Text, New International Greek Testament Commentary. Grand Rapids, MI; Carlisle: W.B. Eerdmans; Paternoster Press, 2005.

Peluso-Verdend, Gary. “Theological Perspective on Matthew 13:31–33, 44–52,” in Feasting on the Word: Preaching the Revised Common Lectionary: Year A, ed. David L. Bartlett and Barbara Brown Taylor, vol. 3. Louisville, KY: Westminster John Knox Press, 2011.

[1] William Barclay, The Gospel of Matthew, Third Ed., The New Daily Study Bible (Edinburgh: Saint Andrew Press, 2001), 147.

[2] 同上.

[3] W. D. Davies and Dale C. Allison Jr., A Critical and Exegetical Commentary on the Gospel according to Saint Matthew, vol. 1, International Critical Commentary (London; New York: T and T Clark International, 2004), 484.

[4] William Arndt et al., A Greek-English Lexicon of the New Testament and Other Early Christian Literature (Chicago: University of Chicago Press, 2000), 521–522.

[5] 原文在「文士和法利賽人」後沒有「義」字,「義」是按著「你們的義」中的「義」字而補上.

[6] William Barclay, The Gospel of Matthew, Third Ed., The New Daily Study Bible (Edinburgh: Saint Andrew Press, 2001), 149.

[7] 同上.

[8] William Arndt et al., A Greek-English Lexicon of the New Testament and Other Early Christian Literature (Chicago: University of Chicago Press, 2000), 1049.

[9] William Barclay, The Gospel of Matthew, Third Ed., The New Daily Study Bible (Edinburgh: Saint Andrew Press, 2001), 149.

[10] Hans Küng, On Being a Christian, trans. Edward Quinn (New York: Doubleday, 1984), 201-211.

[11] William Arndt et al., A Greek-English Lexicon of the New Testament and Other Early Christian Literature (Chicago: University of Chicago Press, 2000), 828.

[12] Donald A. Hagner, Matthew 1–13, vol. 33A, Word Biblical Commentary (Dallas: Word, Incorporated, 1993), 105.

[13] 「圓滿」這動詞用了active, aorist和infinitive的形式.Aorist tense在這裡應不是指過去發生的事,而描述一個一次性,未有指定時間的動作.

[14] 同上.

[15] 同上.



\newpage

\section{智慧抑兆頭}
\label{sec:687}
\textbf{李均熊牧師 - 智慧抑兆頭}
\newline
\newline
\hyperref[sec:686]{< < < PREV SERMON < < <}
~
\hyperref[sec:toc]{[返目錄]}
~
\hyperref[sec:688]{> > > NEXT SERMON > > >}


哥林多前書 1:18-25
\newline
\begin{longtable}{cl}
\hline
\hline
章節 & 經文 (和合本修訂版)\\
\hline
1:18 & \begin{tabularx}{0.7\textwidth}{X} 因為十字架的道理，在那滅亡的人是愚拙，在我們得救的人卻是神的大能。 \end{tabularx} \\ \\ \relax
1:19 & \begin{tabularx}{0.7\textwidth}{X} 就如經上所記：「我要摧毀智慧人的智慧，廢棄聰明人的聰明。」 \end{tabularx} \\ \\ \relax
1:20 & \begin{tabularx}{0.7\textwidth}{X} 智慧人在哪裡？文士在哪裡？這世上的辯士在哪裡？神豈不是已使這世上的智慧變成愚拙了嗎？ \end{tabularx} \\ \\ \relax
1:21 & \begin{tabularx}{0.7\textwidth}{X} 既然世人憑自己的智慧不認識神，神就本著自己的智慧樂意藉著人所傳愚拙的話拯救那些信的人。 \end{tabularx} \\ \\ \relax
1:22 & \begin{tabularx}{0.7\textwidth}{X} 猶太人要的是神蹟，希臘人求的是智慧， \end{tabularx} \\ \\ \relax
1:23 & \begin{tabularx}{0.7\textwidth}{X} 我們卻是傳被釘十字架的基督，這對猶太人是絆腳石，對外邦人是愚拙； \end{tabularx} \\ \\ \relax
1:24 & \begin{tabularx}{0.7\textwidth}{X} 但對那蒙召的，無論是猶太人、希臘人，基督總是神的大能，神的智慧。 \end{tabularx} \\ \\ \relax
1:25 & \begin{tabularx}{0.7\textwidth}{X} 因為，神的愚拙總比人智慧；神的軟弱總比人強壯。 \end{tabularx} \\ \\
[1ex]
\hline
\hline
\end{longtable}
講題:智慧抑兆頭? Wisdom or Sign?

經文:哥林多前書1:18-25

講員:李均熊牧師

場合:香港中文大學崇基學院禮拜堂主日崇拜

日期:2024年3月3日

[18] 因為十字架的道理,在那滅亡的人是愚拙,在我們得救的人卻是上帝的大能. [19] 就如經上所記: 「我要摧毀智慧人的智慧, 廢棄聰明人的聰明.」 [20] 智慧人在哪裡?文士在哪裡?這世上的辯士在哪裡?上帝豈不是已使這世上的智慧變成愚拙了嗎? [21] 既然世人憑自己的智慧不認識上帝,上帝就本著自己的智慧樂意藉著人所傳愚拙的話拯救那些信的人. [22] 猶太人要的是神蹟,希臘人求的是智慧, [23] 我們卻是傳被釘十字架的基督,這對猶太人是絆腳石,對外邦人是愚拙； [24] 但對那蒙召的,無論是猶太人,希臘人,基督總是上帝的大能,上帝的智慧. [25] 因為,上帝的愚拙總比人智慧；上帝的軟弱總比人強壯.

上帝愚拙的智慧

熟識新約聖經或者保羅書信的人都會知道哥林多教會充滿是非.恩賜多多,但同時問題多多.哥林多前書開首沒有幾句話,保羅就指斥哥林多教會結黨分裂,這個分裂有可能是源於有一黨人誇讚自己是「智慧人」,又認為保羅並沒有甚麼「智慧」或「智慧的言語」,去批評論斷保羅,說他不是甚麼好東西,因而分裂了教會,甚至使人誤解耶穌基督的福音.於是保羅用了足足四章的經文為自己辯護,指出他傳的福音,並不用上智慧的言語,因為耶穌的福音是關乎耶穌在十字架上成就的救恩,「在猶太人為絆腳石,在外邦人為愚拙」.

我不是研究新約的,但也知道今日的書信經課中常出現的「智慧」並不是普通含義下的智慧,不是說你擁有多少博士學位,或道德化了的所謂「做人的道理」,智慧只是教你如何做個好人或道德的人等等.哥林多前書中常常提到的「智慧」,希臘文sophia,總共在哥林多前書出現了十七次之多,其中十六次在1-3章,冠絕新約.學者對這個「智慧」的理解都有不同,有認為是諾斯底主義的雛型(proto-Gnosticism),又或希羅哲士的詭辯精英主義 (Graeco-Roman rhetorical elitism),亦可能是希臘化的猶太人所講的智慧傳統(Hellenistic Jewish wisdom tradition).其中希臘化的猶太人的智慧傳統,最可能呼應今日舊約的十誡經課和詩篇選讀.因為猶太人的律法,或稱妥拉(Torah),在新約時期又或更後的時間,漸漸轉化成智慧的教導,正如申命記4:5-6「我照著耶和華－我上帝所吩咐的將律例典章教訓你們,使你們在所要進去得為業的地上遵行.所以你們要謹守遵行；這就是你們在萬民眼前的智慧,聰明.他們聽見這一切律例,必說:『這大國的人真是有智慧,有聰明!』」

但無論這個智慧在哥林多前書是何所指,都不能叫我們認識上帝更加多,也不能使我們得到救恩.反而保羅認為,上帝的智慧,是被人當作愚拙的道理,並且上帝「樂意用人所當作愚批的道理拯救那些信的人」.那是甚麼道理呢,就是十字架的道理,「在那滅亡的人為愚拙,在我們得救的人,卻為上帝的大能」.

今日我們崇拜望著一個大大的十字架,可能大家身上也會有十字架的飾物.但大家或多或少都應該知道,十字架對哥林多人來說,是殘忍的刑具,最羞辱的酷刑,是羅馬帝國為叛國犯所預備的.因為一個人死在十字架上大家就可以得救,真是荒謬絕倫,愚蠢至極(背後原理是怎樣的,大家就要看羅馬書或加拉太書了).但正如接續今日經課林前一章的經文:26弟兄們哪,可見你們蒙召的,按著肉體有智慧的不多,有能力的不多,有尊貴的也不多.27上帝卻揀選了世上愚拙的,叫有智慧的羞愧；又揀選了世上軟弱的,叫那強壯的羞愧.28上帝也揀選了世上卑賤的,被人厭惡的,以及那無有的,為要廢掉那有的,29使一切有血氣的,在上帝面前一個也不能自誇.30但你們得在基督耶穌裡是本乎上帝,上帝又使他成為我們的智慧,公義,聖潔,救贖.31如經上所記:「誇口的,當指著主誇口.」

上帝就是樂意用人所看不起的,成就了不起的救恩.

被誤解的兆頭

另一個在林前1:22出現的詞語「神蹟」,也正好呼應今日福音經課中,猶太領袖詢問耶穌:你能顯甚麼神蹟給我們看,表明你可以做這些事呢?我的講題不用「神蹟」這個字,免得你以為我們是討論耶穌那些超自然的「神蹟」.這裡以英文Sign去解釋可能更為適合.Sign的中文可譯做記號,指示,跡象,徵兆,或基督徒常說的兆頭.字典解說是 Something showing that something else exists or might happen or exist in the future (表示其他事情存在或可能發生或將來會出現的事情).林前1:22論到猶太人要求sēmeion,並不是說祂們要求超自然的神怪事情,或要求耶穌如魔術師在袖裡拿出兔子,猶太人渴求看到時代的徵兆,想知道那將要來的世代來了沒有.可能猶太人在今世經歷太多苦痛吧,所以極其渴望進入末世／那將要來的世代(the age to come).如何得知末世／他世而臨,就是需要一些指示,徵兆,明證才感安心吧?

但一個被針十字架的彌賽亞,對猶太人來說也是愚蠢致極,匪夷所思,又怎可能是新時代來臨的徵兆／兆頭.當耶穌說「你們拆毀這殿,我三日內要再建立起來.」,猶太領袖們自然聽得一頭霧水,還是門徒在主耶穌死而復活之後才明白耶穌以他的身體為殿.怪就只可以怪耶穌有甚麼事情都不會直講,未到時候所謂的兆頭就算顯現出來你也無從得知.

反轉期望的應許

在我看來,說穿了,人就是期望掌握人生,得到指引,在茫茫世情中明白處世之道.無論是希臘人的智慧,或猶太人的妥拉,又或他們追求的兆頭,都是希望在世情多變的不知所措之下找到一些安心的明證.如果尋求智慧和兆頭是為了緩解內心的焦慮和不安,耶穌的在十字架,又怎樣可以成為上帝的大能呢?保羅不也說傳釘十字架的基督「在猶太人為絆腳石,在外邦人為愚拙」?

上帝的智慧,就是要反轉人間的期望.人之所以焦慮和不安,在上帝眼中,可能就是因為我們所追求和渴望的都是錯重點.我們以為知道考試答案,就可以考好一百份嗎?我們以為知道明天是陰是晴,就可以安心快活嗎?那知道考試的目的不是要我們擁有一百份,而是預備自己接受艱難的挑戰；知道明天陰晴只可以幫助我們選揀行裝,但路徑仍有無窮變數.當我們的期望不是投注在世人的智慧,自己的能力,外在的轉變,良好的兆頭,對不可操控的掌握,而是曾經受苦受死,且死在十字架上的我們的主耶穌.因為祂受死卻復活過來,證明祂是神的大能,正如保羅在林前2:9-12對哥林多人的保証:

9. 上帝為愛他的人所預備的是眼睛未曾看見,耳朵未曾聽見,人心也未曾想到的.

10只有上帝藉著聖靈向我們顯明了,因為聖靈參透萬事,就是上帝深奧的事也參透了.

11除了在人裡頭的靈,誰知道人的事?像這樣,除了上帝的靈,也沒有人知道上帝的事.

12我們所領受的,並不是世上的靈,乃是從上帝來的靈,叫我們能知道上帝開恩賜給我們的事.

現在的我怎樣面對明日未知的不安和事情不可掌控的焦慮?當然我作決定前仍會祈求上主的兆頭和智慧,不安和焦慮也是叫人預備好自己面對挑戰的內在動力.但明白到上主自有主意,祂又會將最好的給我,甚至連祂獨生愛子也願意為我死在十字架上,我還有甚麼需要怕,需要擔心的呢?



\newpage

\section{上帝愛世人}
\label{sec:688}
\textbf{高國雄牧師 - 上帝愛世人}
\newline
\newline
\hyperref[sec:687]{< < < PREV SERMON < < <}
~
\hyperref[sec:toc]{[返目錄]}
~
\hyperref[sec:689]{> > > NEXT SERMON > > >}


約翰福音 3:16
\newline
\begin{longtable}{cl}
\hline
\hline
章節 & 經文 (和合本修訂版)\\
\hline
3:16 & \begin{tabularx}{0.7\textwidth}{X} 「神愛世人，甚至將他獨一的兒子賜給他們，叫一切信他的人不致滅亡，反得永生。 \end{tabularx} \\ \\
[1ex]
\hline
\hline
\end{longtable}
講題:上帝愛世人 God So Loved the World

講員:高國雄牧師

經文:約翰福音3:16

場合:香港中文大學崇基學院禮拜堂主日崇拜

日期:2024年3月10日

約翰福音3:16

16「上帝愛世人,甚至將他的獨生子賜給他們,叫一切信他的,不致滅亡,反得永生.

基督教是「愛」的宗教相信沒有人會異議.

這與今天一般人所認識,所體會的愛,有什麼分別呢?

中國人的「愛」字,戰國時出現「金文」

這個張口的人,用右手撫著心以表示“張口告人,心裡喜歡”的意思,這就是“愛”這個會意字.

秦始皇併吞六國以後統一使用的的小篆“愛”,字腳卻多了一隻向下的足,以表示“愛”是一種行為與行動,[1]

現今世代的愛有何特徵?

愛=感覺良好(好Feel)?

愛=自我滿足?投訴文化,鍵盤戰士?

愛=佔有?

愛=非理性?

約翰福音3:14-15

14摩西在曠野怎樣舉蛇,人子也必照樣被舉起來, 15叫一切信他的都得永生.

摩西舉蛇,銅蛇的杆

民數記21:4-9

4他們從何珥山起行,繞過以東地往紅海那條路走.在路上,百姓心中煩躁. 5百姓向上帝和摩西發怨言,說:「你們為甚麼把我們從埃及領上來,使我們死在曠野呢?這裡沒有糧食,沒有水,我們厭惡這淡而無味的食物.」 6耶和華派火蛇進入百姓當中去咬他們,於是以色列中死了許多百姓. 7百姓到摩西那裡,說:「我們有罪了,因為我們向耶和華和你發怨言.求你向耶和華禱告,叫蛇離開我們.」於是摩西為百姓禱告. 8耶和華對摩西說:「你要造一條火蛇,掛在杆子上.凡被咬的,一望這蛇就必存活.」 9摩西就造了一條銅蛇,掛在杆子上.凡被蛇咬的,一望這銅蛇就活了.

認罪 > 悔改 > 仰望

人犯罪 > 定罪 > 應受公義下的審判 > 死

以弗所書2:1-3

1從前,你們因著自己的過犯罪惡而死了. 2那時,你們在過犯罪惡中生活,隨從今世的風俗,順服空中掌權者的領袖,就是現今在悖逆的人心中運行的邪靈. 3我們從前也都生活在他們當中,放縱肉體的私慾,隨著肉體和心中的意念去做,和別人一樣,生來就是該受懲罰的人.

基督教的愛:我們可以有所仰望

約翰福音3:1-15

1有一個法利賽人,名叫尼哥德慕,是猶太人的官. 2這人夜裡來見耶穌,說:「拉比,我們知道你是由上帝那裡來作師傅的；因為你所行的神蹟,若沒有上帝同在,無人能行.」 3耶穌回答說:「我實實在在地告訴你,人若不重生,就不能見上帝的國.」 4尼哥德慕說:「人已經老了,如何能重生呢?豈能再進母腹生出來嗎?」 5耶穌說:「我實實在在地告訴你,人若不是從水和聖靈生的,就不能進上帝的國. 6從肉身生的就是肉身；從靈生的就是靈. 7我說:『你們必須重生』,你不要以為希奇. 8風隨著意思吹,你聽見風的響聲,卻不曉得從哪裡來,往哪裡去；凡從聖靈生的,也是如此.」 9尼哥德慕問他說:「怎能有這事呢?」 10耶穌回答說:「你是以色列人的先生,還不明白這事嗎? 11我實實在在地告訴你,我們所說的是我們知道的；我們所見證的是我們見過的；你們卻不領受我們的見證. 12我對你們說地上的事,你們尚且不信,若說天上的事,如何能信呢? 13除了從天降下,仍舊在天的人子,沒有人升過天.14摩西在曠野怎樣舉蛇,人子也必照樣被舉起來, 15叫一切信他的都得永生.

人需要愛,愛就是盼望

約翰福音3:16-21

16「上帝愛世人,甚至將他的獨生子賜給他們,叫一切信他的,不致滅亡,反得永生. 17因為上帝差他的兒子降世,不是要定世人的罪,乃是要叫世人因他得救. 18信他的人,不被定罪；不信的人,罪已經定了,因為他不信上帝獨生子的名. 19光來到世間,世人因自己的行為是惡的,不愛光,倒愛黑暗,定他們的罪就是在此. 20凡作惡的便恨光,並不來就光,恐怕他的行為受責備. 21但行真理的必來就光,要顯明他所行的是靠上帝而行.」

以弗所書2:4-10

4然而,上帝有豐富的憐憫,因著他愛我們的大愛, 5竟在我們因過犯而死了的時候,使我們與基督一同活過來—可見你們得救是本乎恩— 6他又使我們在基督耶穌裡與他一同復活,一同坐在天上, 7為要把他極豐富的恩典,就是他在基督耶穌裡向我們所施的恩慈,顯明給後來的世代. 8你們得救是本乎恩,也因著信；這並不是出於自己,而是上帝所賜的； 9也不是出於行為,免得有人自誇. 10我們是他所造之物,在基督耶穌裡創造的,為要使我們行善,就是上帝早已預備好要我們做的.

上帝的愛,就是拯救的恩典

單單出於上帝

羅馬書5:8

8惟有基督在我們還作罪人的時候為我們死,上帝的愛就在此向我們顯明了.

約翰一書4:9

10不是我們愛上帝,乃是上帝愛我們,差他的兒子為我們的罪作了挽回祭,這就是愛了.

本乎恩的愛

上帝因愛而有所行動,道成了肉身

我們愛,又會否行動呢?

因著信的愛

約翰一書1:9

9我們若認自己的罪,上帝是信實的,是公義的,必要赦免我們的罪,洗淨我們一切的不義

約翰一書3:16-18

16主為我們捨命,我們從此就知道何為愛；我們也當為弟兄捨命. 17凡有世上財物的,看見弟兄窮乏,卻塞住憐恤的心,愛上帝的心怎能存在他裡面呢? 18小子們哪,我們相愛,不要只在言語和舌頭上,總要在行為和誠實上.

For God so loved the world

(TEXT: Frances Townsend, MUSIC: Alfred B. Smith)

For God so loved the world,

He gave His only Son,

to die on Calvary's tree

from sin to set me free.

Some day He's coming back,

what glory that will be!

Wonderful His love to me!



\newpage

\section{我見主坐在高高的寶座上}
\label{sec:689}
\textbf{鄧瑞強博士 - 我見主坐在高高的寶座上}
\newline
\newline
\hyperref[sec:688]{< < < PREV SERMON < < <}
~
\hyperref[sec:toc]{[返目錄]}
~
\hyperref[sec:690]{> > > NEXT SERMON > > >}


以賽亞書 6:1-13
\newline
\begin{longtable}{cl}
\hline
\hline
章節 & 經文 (和合本修訂版)\\
\hline
6:1 & \begin{tabularx}{0.7\textwidth}{X} 當烏西雅王崩的那年，我看見主坐在高高的寶座上。他的衣裳下襬遮滿聖殿。 \end{tabularx} \\ \\ \relax
6:2 & \begin{tabularx}{0.7\textwidth}{X} 上有撒拉弗侍立，各有六個翅膀：兩個翅膀遮臉，兩個翅膀遮腳，兩個翅膀飛翔， \end{tabularx} \\ \\ \relax
6:3 & \begin{tabularx}{0.7\textwidth}{X} 彼此呼喊說：「聖哉！聖哉！聖哉！萬軍之耶和華；他的榮光遍滿全地！」 \end{tabularx} \\ \\ \relax
6:4 & \begin{tabularx}{0.7\textwidth}{X} 因呼喊者的聲音，門檻的根基震動，殿裡充滿了煙雲。 \end{tabularx} \\ \\ \relax
6:5 & \begin{tabularx}{0.7\textwidth}{X} 那時我說：「禍哉！我滅亡了！因為我是嘴唇不潔的人，住在嘴唇不潔的民中，又因我親眼看見大君王－萬軍之耶和華。」 \end{tabularx} \\ \\ \relax
6:6 & \begin{tabularx}{0.7\textwidth}{X} 有一撒拉弗向我飛來，手裡拿著燒紅的炭，是用火鉗從壇上取下來的， \end{tabularx} \\ \\ \relax
6:7 & \begin{tabularx}{0.7\textwidth}{X} 用炭沾我的口，說：「看哪，這炭沾了你的嘴唇，你的罪孽便除掉，你的罪惡就赦免了。」 \end{tabularx} \\ \\ \relax
6:8 & \begin{tabularx}{0.7\textwidth}{X} 我聽見主的聲音說：「我可以差遣誰呢？誰肯為我們去呢？」我說：「我在這裡，請差遣我！」 \end{tabularx} \\ \\ \relax
6:9 & \begin{tabularx}{0.7\textwidth}{X} 他說：「你去告訴這百姓說：『你們聽了又聽，卻不明白；看了又看，卻不曉得。』 \end{tabularx} \\ \\ \relax
6:10 & \begin{tabularx}{0.7\textwidth}{X} 要使這百姓心蒙油脂，耳朵發沉，眼睛昏花；恐怕他們眼睛看見，耳朵聽見，心裡明白，回轉過來，就得醫治。」 \end{tabularx} \\ \\ \relax
6:11 & \begin{tabularx}{0.7\textwidth}{X} 我就說：「主啊，這到幾時為止呢？」他說：「直到城鎮荒涼，無人居住，房屋空無一人，土地極其荒蕪； \end{tabularx} \\ \\ \relax
6:12 & \begin{tabularx}{0.7\textwidth}{X} 耶和華將人遷到遠方，國內被撇棄的土地很多。 \end{tabularx} \\ \\ \relax
6:13 & \begin{tabularx}{0.7\textwidth}{X} 國內剩下的人若還有十分之一，也必被吞滅。然而如同大樹與橡樹，雖被砍伐，殘幹卻仍存留，聖潔的苗裔是它的殘幹。」 \end{tabularx} \\ \\
[1ex]
\hline
\hline
\end{longtable}
講題:我見主坐在高高的寶座上 I Saw the Lord Sitting on a High and Lofty Throne

經文:以賽亞書6:1-13

講員:鄧瑞強博士

場合:香港中文大學崇基學院禮拜堂主日崇拜

日期:2024年3月17日

各位弟兄姊妹,早安.願神的恩典常與你們同在.

之前幾次講道,我們讀詩歌智慧文學,今天,我們開始看舊約的另一部分:先知書.由於聖經被釘裝成一本書,我們以為聖經是一本「書」.其實,聖經不是一本書,而是66卷獨立書卷的結集,裡面包含不同經卷的獨特視角.有些經卷從「神人之約」的角度看問題；有些經卷從「大衛王朝」的延續看事情；有些經卷講「祭司神學」,講禮儀,講秩序；有些經卷卻講「先知神學」,講悔改,講遠象；智慧文學講神創造的秩序,這是人人都可參透的智慧；先知文學講神在歷史中的作為,這就只有信仰的心靈才能領悟.弟兄姊妹有時很困惑,為何一本書有看來不太協調的不同向度的,這是因為聖經不是一本書,而是眾多作品的結集.

中國的《四庫全書》,一個中國傳統典籍的結集,以四個部分來分類:經,史,子,集.舊約的經卷也大概可作如此分類.經,是摩西五經.史,是約書亞記到以斯帖記的一系列書卷.子,是先知諸子.集,是約伯記到雅歌的一系列詩歌智慧文學.上幾個月,我們看詩歌智慧文學,你會覺得很有趣味,飽含哲理,令人心曠神怡.今日開始,我們進入「子」的系列,會聽到不同先知對罪惡的譴責,對歷史的反思,對前景的想像.你們會開始討厭我,因為無人喜歡別人揭露自己的罪.先知文學連耶利米哀歌共有17卷,我大致要一年半時間才能講完,你們討厭我的日子將會很長.

今日,我們看以賽亞書6:1-13,先看6:1-4.

賽6:1　當烏西雅王崩的那年,我見主坐在高高的寶座上.他的衣裳垂下,遮滿聖殿.

賽6:2　其上有撒拉弗侍立,各有六個翅膀:用兩個翅膀遮臉,兩個翅膀遮腳,兩個翅膀飛翔；

賽6:3　彼此呼喊說:聖哉!聖哉!聖哉!萬軍之耶和華；他的榮光充滿全地!

賽6:4　因呼喊者的聲音,門檻的根基震動,殿充滿了煙雲.

烏西雅王,在位大概50年,是猶大國在位最長的其中一位君主.那時,國際新興的勢力亞述,由於內亂,無暇東征西討.烏西雅王得以享受長時間的安定,他在位時,是一和平安穩的盛世.烏西雅王死後,新興的勢力亞述挑戰傳統的勢力埃及,爭做武林盟主,猶大國夾在亞述與埃及中間,其艱難可想而知.想想:獅子和老虎打架,一隻小雞站在中間,其凶險不言而喻.事實上,烏西雅王死後大概20年,比猶大國強大得多的以色列國,便被亞述所滅.亞述也曾圍攻耶路撒冷,多得一個神蹟,可能是瘟疫,才令亞述軍大敗而回.

「烏西雅王崩的那年」,是極盛時代轉向衰敗的轉折點,戰爭的鼓聲,隱約可聞.猶大人眼看堂堂一個姊妹國以色列被滅,內心一定會想,會否輪到自己?就在這個歷史的拐變點,神興起先知以賽亞.

先知是信仰敏銳的人,當人人都沈醉在人間的歌舞昇平裡時,他卻看到上帝的神聖世界；當人人都對罪惡習以為常之時,他卻對一個老闆欺騙了一個婆婆一毫子而耿耿於懷；當人人都在宗教裡尋求安穩和慰藉時,他卻以上帝之名挑戰人去承擔歷史的重負；當人人都活得像泥漿一樣時,他呼喚人活得有靈魂；當人人都對前景憂心忡忡時,他卻宣告意想不到的希望.猶太神學家Abraham Heschel說:當人人舒適安逸地在睡覺時,先知在夜間驚叫,因為他感覺到那從天而來的爆炸.(Abraham Heschel, The Prophets, vol.1, p.16)

賽6:1　當烏西雅王崩的那年,我見主坐在高高的寶座上.他的衣裳垂下,遮滿聖殿.

在人間,一個偉大王者離去.以賽亞卻看到另一個王者,那掌管人間歷史的神.當光輝的時代結束時,人最需要知道的,是誰掌管明天.對猶大人來說,烏西雅王崩了,盛世也可能結束了,明天將會如何?以賽亞卻在聖殿裡,在這神人相遇之地,看到在高高的寶座上,永活的神坐在為王.明天,在祂的引導之下.以賽亞有的,是一種超越的視野.他看到神,他從神的角度看人間,這是先知獨特的視角.

神坐著,祂的四圍有撒拉弗侍立.撒拉弗可能是神獸,當然可以理解為天使,其字根的意義為「燃燒著的東西」,他們的外形說不定像火,他們一刻間用炭火去潔淨以賽亞.這些撒拉弗呼喊:「聖哉!聖哉!聖哉!萬軍之耶和華；他的榮光充滿全地!」「聖哉!聖哉!聖哉!」是對神的聖潔的至高頌讚.奇怪的是,如此神聖的上帝,祂的榮光不是閃耀在宇宙太空中,而是充滿全地.我們活在一個被上帝的榮耀包圍的大地之上.歷史充滿危機,上帝的榮耀卻像母親的胸懷,擁抱著我們.你所在之地,被神的榮光包圍著.

賽6:5　那時我說:「禍哉!我滅亡了!因為我是嘴唇不潔的人,又住在嘴唇不潔的民中,又因我眼見大君王－萬軍之耶和華.」

賽6:6　有一撒拉弗飛到我跟前,手裡拿著紅炭,是用火剪從壇上取下來的,

賽6:7　將炭沾我的口,說:「看哪,這炭沾了你的嘴,你的罪孽便除掉,你的罪惡就赦免了.」

面對神,以賽亞有強烈的感受.神學家Rudolf Otto寫了一本書《The Idea Of The Holy》,描繪人面對神之時的感受.最基本的感受,是體會到一種無以名狀的無窮偉大,Rudolf Otto稱之為「Mysterium」.在這無窮偉大之下,人感到無限卑微；不單如此,人也體會到這無限者對我們的存在的審視,人不期然感到戰驚恐懼,不知所措,Rudolf Otto稱之為「Tremendum」.奇妙的是,人同時又感到這無限者有一種無窮的吸引力,祂充滿著安慰的力量,人被吸引,來到祂面前,然後體會被無條件接納,Rudolf Otto稱之為「Fascinans」.

以賽亞體會到無限偉大的神就在眼前,然後,他有這兩種感受:戰驚恐懼,無條件被接納.他說:「禍哉!我滅亡了!因為我是嘴唇不潔的人,又住在嘴唇不潔的民中,又因我眼見大君王－萬軍之耶和華.」他體會到自己的不潔,體會到自己活在神的審判之下.在神面前,他戰驚恐懼.現在,很多信徒已沒有了這種戰驚恐懼的心,神像令人開心的玩具,多過是令人敬畏的無限神祕.

之後,撒拉弗用祭壇上的炭火沾以賽亞的口,宣告他的罪被赦免了.一個自覺被審判的人,忽然間,他的罪便被赦免了.在神面前,敬畏神,自覺有罪而認罪的人,體會到這無條件的赦免.先知對罪敏感,也對神的恩典敏感.相比而言,我們可能已對罪麻木,也對恩典麻木.明知是謊言,也照樣講,這是對罪麻木.手中所擁有的,不願意分享,這是對恩典麻木.

賽6:8　我又聽見主的聲音說:「我可以差遣誰呢?誰肯為我們去呢?」我說:「我在這裡,請差遣我!」

「我可以差遣誰呢?誰肯為我們去呢?」這裡的「我們」,當然與「三位一體」沒有什麼關係.這是天庭會議,神在策劃著人間歷史的演變.神將要藉著國際的大變局,來對祂的子民猶大國的罪作一清算.神要向祂的子民宣佈審判.祂問:「誰願意承擔這發言人的職責呢?」

以賽亞看來並不知道要為神發什麼言,但他意識到猶大民族的前途,端賴有人向他們宣講神的信息.先知是奇怪的人,當人人都不願意為歷史的將來負上責任時,他們卻甘心背起時代的十架.對他們而言,信仰不是等上天堂,而是為神的真理實現在人間而獻上自己.這不是追求個人的幸福,而是追求天國的實現.以賽亞回應神,他說:「我在這裡,請差遣我!」在一個對真理無興趣的世界言說真理,真是一個重負.

賽6:9　他說:「你去告訴這百姓說:你們聽是要聽見,卻不明白；看是要看見,卻不曉得.

賽6:10　要使這百姓心蒙脂油,耳朵發沉,眼睛昏迷〔盲目〕；恐怕眼睛看見,耳朵聽見,心裡明白,回轉過來,便得醫治.」(留意:心,耳,眼,盲目vs.看見,眼,耳,心,他們對神完全盲目.)

這段經文,新約引用了六次,這包括兩次不完整的引用(太13:14-15; 可4:12, 8:18〔局部引用〕; 路8:10〔局部引用〕; 約12:39-40; 徒28:26-27).在福音書裡,耶穌在人間,講神的審判與恩典,大部分在場的人,就是不信與拒絕.這段經文,大致而言,有兩種解釋向度.一是強勢向度,神就是要使百姓心蒙脂油,耳朵發沉,眼睛昏迷,以致他們拒絕先知的信息,好降禍給他們.強調神的主權及預定論的神學家,傾向這觀點.另一種解釋向度,多少認為這經文表達的,是神對人間罪惡現實的哀歎.持這觀點的人指出,這段經文其實是詩歌體裁,不能完全理解為事實描繪,其中舒發情感的元素不可忽視.神眼見人的墮落已是一結構性的墮落,你越叫他們不要犯罪,他們越是犯罪給你看.正如廣東俗語所講:「阿崩叫狗,越叫越走」,上帝也徒呼奈何.現在,很多父母叫子女不要打機,他們越是打給你看,若打不到機,他們寧願死給你看.面對冥頑不靈的人,父母無辦法,上帝也無辦法,只能讓他們承受苦果.如此這般,就不是神令他們盲目,而是他們定意要盲目,如此,神也只得無奈接受這現實,讓他們繼續盲目下去,好讓他們自招的苦果,能令他們猛然省悟.

不知以賽亞心情如何?原來神要他講的信息,是人人都討厭的信息.是要叫貪錢的人不要再貪錢,叫吃得太多的人不要吃得太多,叫縱慾的人不要再縱慾,這是貪錢,貪吃,縱慾的人最不想聽到的信息.當整個社會都是這樣的人,你能講些什麼?為了喚醒世人,先知將成為一個孤獨的人.

賽6:11　我就說:「主啊,這到幾時為止呢?」他說:直到城邑荒涼,無人居住,房屋空閒無人,地土極其荒涼.

賽6:12　並且耶和華將人遷到遠方,在這境內撇下的地土很多.

賽6:13　境內剩下的人若還有十分之一,也必被吞滅,像栗樹,橡樹雖被砍伐,樹墩子卻仍存留.這聖潔的種類在國中也是如此.

要向一群不會回應的人,傳講神的說話,猶如對著一群睡著覺的學生努力講書一樣,花心機卻注定徒勞無功.以賽亞問:「主啊,這到幾時為止呢?」神說:直到學校倒塌,變成頹垣敗瓦之時；直到猶大國被滅,城邑荒涼,人民被擄之日.這又回到太古之初的虛無了.看來,人性不會變,罪性不會改,不到絕境,人是不會回頭的了.徹底拆毀了,才有重生的可能.

最後一節經文,是亡國後的希望.樹砍了,還留有餘幹.這餘幹仍能發新芽,神的子民裡總還有一些願意悔改的人,他們成了日後的神的新子民.

從烏西雅王崩的那年,到猶大國真的覆亡,其實還有約150年.以賽亞沒有親眼看著猶大國的毀滅,但他的信息一直傳留在猶大人當中.當猶大人被擄到巴比倫之時,他們回憶起,以賽亞曾經警告過他們的先祖,要聽從神的命令,否則國破家亡,如今,一切已成事實.在絕望中,他們記起,以賽亞留下希望,民族會重生.就因先知這句話,他們在異鄉,堅強地站了起來.

今日,先知在哪裡呢?在這聖殿裡,還有人看見神坐在高高的寶座上嗎?在神的臨在裡,還有人意識到自己的罪嗎?還有人意識到社會的罪嗎?還有人體會到能在神面前存活已是恩典嗎?還有人聽到從天而來的召喚嗎?還有人願意背起時代的責任嗎?或許,在這甘於平庸的世界,已沒有先知了.

願上主得著榮耀,頌讚；願眾人得享喜樂,平安.



\newpage

\section{揮舞棕枝的激情}
\label{sec:690}
\textbf{林豪恩先生 - 揮舞棕枝的激情}
\newline
\newline
\hyperref[sec:689]{< < < PREV SERMON < < <}
~
\hyperref[sec:toc]{[返目錄]}
~
\hyperref[sec:691]{> > > NEXT SERMON > > >}


講題:揮舞棕枝的激情 The Passion of Waving Palm Branches

講員:林豪恩先生

場合:香港中文大學崇基學院禮拜堂主日崇拜

日期:2024年3月24日

引言

今天是棕技主日.這主日是紀念當日群眾以棕樹枝歡迎耶穌進入耶路撒冷的事件.傳統上,信徒在棕枝主日首先齊集在教堂外,各人手持棕枝或樹枝參與巡行禮,即是圍繞教堂外面一同巡行,群眾一邊行走,一邊高唱聖詩,並高呼基督為王.大家如果仍記得剛才誦讀的福音經課,就會想起經文所記述的情況.

一．揮舞棕枝的激情(馬可福音)

11:1耶穌和門徒快到耶路撒冷,來到伯法其和伯大尼,在橄欖山那裡.耶穌打發兩個門徒,

11:2對他們說:「你們往對面村子裡去,一進去的時候會看見一匹驢駒拴在那裡,是從來沒有人騎過的,把牠解開,牽來.

11:3若有人對你們說:『為甚麼做這事?』你們就說:『主要用牠,但會立刻把牠牽回到這裡來.』」

11:4他們去了,看見一匹驢駒拴在門外街道上,就把牠解開.

11:5在那裡站著的人,有幾個說:「你們解開驢駒做甚麼?」

11:6門徒照著耶穌的話說,那些人就任憑他們牽去了.

11:7他們把驢駒牽到耶穌那裡,把自己的衣服搭在上面,耶穌就騎上.

11:8有許多人把衣服鋪在路上,還有人把田間的樹枝砍下來鋪上.

11:9前呼後擁的人都喊著說:「和散那!奉主名來的是應當稱頌的!

11:10那將要來的我祖大衛之國是應當稱頌的!至高無上的,和散那!」

11:11耶穌到了耶路撒冷,進入聖殿,看了周圍的一切.天色已晚,他就和十二使徒出城,往伯大尼去.

在剛才的片段中,你看見棕樹枝嗎?可能有人看見,有人看不見,因為馬可福音只提及樹枝,而約翰福音則明言這是棕樹枝.當時人們拿著棕樹枝迎接耶穌進入耶路撒冷有甚麼含意呢?和現在信徒拿著棕樹枝在教堂外巡行是不是一樣的意思呢?讓我們講幾個背景故事,或許能夠幫助我們探究當時人們拿棕樹枝迎接耶穌入城的想法和期望.

關於棕樹枝的故事可以追溯100多年之前.《馬加比一書》記載了一件對猶太人非常重要的事件.當時猶太人被西流古王朝統治,到安提亞古四世時期,在猶大地推行去猶太化,全面希臘化的政策.有猶太人抗拒,安排亞古四世的軍隊入城,殺死抗拒者,褻瀆聖所,禁止猶太人獻祭,並放置偶像及猶太人視為不潔之物在聖殿中.在一處鄉郊地方,有一個民間祭司家庭起而反抗,後來被稱為「馬加比革命」.馬加比兄弟前仆後繼,到第三任領袖兼大祭司約拿單領導時,終於打敗外敵,驅逐了佔領耶路撒冷的外邦人.猶太人潔淨聖城,掃除污穢之物.《馬加比一書13:51》:「他們帶頌讚,棕櫚枝,琴瑟,鐃鈸,弦樂器,詩歌和頌詞,進入堡壘,因為大敵已經被粉碎,從以色列中除去.」百多年之後,當耶穌進入耶路撒冷的時候,揮動棕樹枝迎接他的百姓心中會不會想著這段歷史呢?會不會盼望著百多年前猶太民族這光輝的一頁呢?

馬加比革命令猶太人擺脫希臘人的統治而恢復猶太王國,耶穌時代第一世紀猶太人活在羅馬統治之下,他們會否也渴望脫離外邦統治者呢?事實上,耶穌時期的猶大地區充滿動盪不安.大家還記得,耶穌出年時是希律王管治,但在耶穌被判死刑的時候,卻是羅馬總督彼拉,為甚麼是希律王呢?彼拉多取代希律王的過程,正正反映出當地局勢的動盪.耶穌出年時的希律王被稱為大希律,他管治猶大附近地區四十年.希律王駕崩,燃起了不少猶太人的期望,各地不時出現自稱彌賽亞的軍事領袖,帶領猶太民眾對抗羅馬政府.羅馬政府把大希律的管轄地分為三分交給他的三個兒子,繼承了耶路撒冷兒子無力處理這些衝突,被羅馬當局流放,羅馬派總督直接管治耶路撒冷和猶大地區.在羅馬統治的這些年間,以整個大敘利亞地區來計,被判為反羅馬動亂而被釘死在十字架上的人,是以千為單位的.

當時,各地的群眾湧往耶路撒冷過逾越節.耶路撒冷聚集如此多人,對管治者來說是一個非常敏感及危險的時刻,只要有人振臂高呼,隨時就有事情發生.就在這個時候,耶穌騎上驢駒進城.人們看到,會想到甚麼呢?讓我們講另外一個背景故事:大約500年前,有一位先知撒迦利亞曾預言:「錫安的民哪,應當大大喜樂；耶路撒冷的民哪,應當歡呼.看哪,你的王來到你這裡!他是公義的,並且施行拯救,謙謙和和地騎著驢,就是騎著驢的駒子.」(九 9)先知預告,騎著驢來的是王.當時的猶太人心目中的王,是一位怎樣的王呢?如果他們在想像一位政治軍事領袖,帶領他們驅逐外邦人,恢復猶大王國,不正正呼應著他們的民族感情和盼望嗎?

馬可福音指出「前呼後擁的人都喊著說:「和散那!奉主名來的是應當稱頌的!那將要來的我祖大衛之國是應當稱頌的!至高無上的,和散那!」」(11:9-10) 大衛之國是甚麼?無論說的人有沒有動機,聽在某些人心裡,一定感到不妥.路加福音記下了「人群中有幾個法利賽人對耶穌說:「老師,責備你的門徒吧!」」(19:39)為甚麼呢?可能法利賽人覺察到他們的言行其實非常危險,公然高呼王者來臨是意圖謀反的死罪.然而,這正正反映著當時不少人的心聲和期望,期望上主的拯救者來到,救他們脫離外族統治,恢復猶大王國.在當時群眾的心目中,如果有任何人自稱基督,都應該是這樣的一位領袖.對一位勝利者基督的盼望,連耶穌的門徒也不例外.你還記得這個片段嗎?

在往凱撒利亞腓立比的路上,耶穌問門徒:「人們說我是誰?」門徒回答說:「是施洗的約翰；有人說是以利亞；又有人說是先知中的一位.」耶穌又問他們:「你們說我是誰?」彼得回答他:「你是基督.」(馬可8:27-29) 然後呢?如果你記得下文,答對了問題的彼得不但沒有100分,反而被耶穌斥責.為甚麼呢?

8:31從此,他教導他們說:「人子必須受許多的苦,被長老,祭司長和文士棄絕,並且被殺,三天後復活.」

8:32耶穌明白地說了這話,彼得就拉著他,責備他.

8:33耶穌轉過來看著門徒,斥責彼得說:「撒旦,退到我後邊去!因為你不體會上帝的心意,而是體會人的意思.」

從這事件可以想到,當時所有人都認為基督是帶領猶太人復國的勝利者,絕不可能是失敗者和受苦者.

二．The Passion of Waving Palm Branches

今天講題的英文是The Passion of Waving Palm Branches.Passion可以解作激情,剛才我們感受到人們迎接耶穌入城的激情.一語相關,The Passion,卻是指耶穌受難.人們的passion和耶穌的Passion是截然不同的,甚至指向相反的方向.當你讀今天的福音經課時,在你心中所產生的是群眾的激情,還是想到耶穌的受難呢?如果沒有想到,也不要緊,因為經課的其他三段經文作出了充分的提示.

以賽亞書:受苦的僕人

50:4主耶和華賜我受教者的舌頭,使我知道怎樣用言語扶助疲乏的人.主每天早晨喚醒,喚醒我的耳朵,使我能聽,像受教者一樣.

50:5主耶和華開啟我的耳朵,我並未違背,也未退後.

50:6人打我的背,我任他打；人拔我兩頰的鬍鬚,我由他拔；人侮辱我,向我吐唾沫,我並不掩面.

50:7主耶和華必幫助我,所以我不抱愧.我硬著臉面好像堅石,也知道我必不致蒙羞.

50:8稱我為義的與我相近；誰與我爭論,讓我們來對質；誰與我作對,讓他近前來吧!

50:9看哪,主耶和華必幫助我,誰能定我有罪呢?

詩篇:受苦的僕人

31:9耶和華啊,求你憐憫我,因為我在急難之中；我的眼睛因憂愁而昏花,我的身心也已耗盡.

31:10我的生命為愁苦所消耗,我的年歲為嘆息所荒廢；我的力量因我的罪孽衰敗,我的骨頭也枯乾.

31:11我因所有的敵人成了羞辱,在我鄰舍跟前更加羞辱；那認識我的都懼怕我,在街上看見我的都躲避我.

31:12我被遺忘,如同死人,無人記念；我好像破碎的器皿.

31:13我聽見許多人的毀謗,四圍盡是驚嚇；他們一同商議攻擊我,圖謀害我的性命.

31:14耶和華啊,我仍要倚靠你；我說:「你是我的上帝.」

31:15我終生的事在你手中,求你救我脫離仇敵的手和那些迫害我的人.

腓立比書:受苦的僕人

2:5你們當以基督耶穌的心為心:

2:6他本有上帝的形像,卻不堅持自己與上帝同等；

2:7反倒虛己,取了奴僕的形像,成為人的樣式；既有人的樣子,

2:8就謙卑自己,存心順服,以至於死,且死在十字架上.

2:9所以上帝把他升為至高,又賜給他超乎萬名之上的名,

2:10使一切在天上的,地上的和地底下的,因耶穌的名,眾膝都要跪下,

2:11眾口都要宣認:耶穌基督是主,歸榮耀給父上帝.

如果迎接耶穌進城的群眾盼望的是勝利的基督,經課引導我們去認識的卻是受苦的基督.

總結:

透過今天棕枝主日的經課,讓我們作以下幾點反省:

今天,我們仍然可以揮動棕樹枝迎接耶穌基督.但耶穌基督入城之後發生在他身上的事情,提示我們所迎接的不是勝利的基督,而是受苦的基督.因為接下來,耶穌基督被門徒出賣(欺騙,背叛,憤怒,悲傷),被拘捕(失去自由),被門徒四散離棄(孤苦無援),被官兵毆打嘲諷(身體痛楚和心靈痛楚),被假證供判處死刑(屈枉正直),被天父離棄(可15:34我的上帝!我的上帝!為甚麼離棄我?),最後死在十字架上(失去生命).這是接下來聖週的內容,我們可以不急於慶祝復活節,而先與耶穌基督共情同感共度這受苦的聖週嗎?

其實,耶穌基督所受的苦,正是人間眾苦的寫照.被背叛的傷痛和憤怒,自由受限制的困頓,孤獨無援的恐懼,被批評嘲諷的難受,身體病患帶來的痛苦,各種的心靈憂傷,被不公平對待的委屈等,任何一個人或多或少總會嘗過.有人埋怨上帝為甚麼不創造一個無痛的世界,怪責上帝不制止人們互相傷害.上帝的回應卻是親臨有痛的世界,進人互相傷害的人間,和萬千受苦者一同受苦.

我們要如彼得和當時的群眾一樣單單接受勝利的基督,還是也接受受苦的基督呢?

我們要「身光頸靚」成功神學基督教,還是要投入人間疾苦,在污泥濁水中和受苦的人一齊同行的基督教呢?

我們自己要做無痛的基督徒,還是要做就算受苦都跟隨基督的基督徒呢?

受苦不一定等於蒙羞(以賽亞書50:7主耶和華必幫助我,所以我不抱愧.我硬著臉面好像堅石,也知道我必不致蒙羞.)

受苦不一定等於有罪(以賽亞書50:8稱我為義的與我相近；誰與我爭論,讓我們來對質；誰與我作對,讓他近前來吧!50:9看哪,主耶和華必幫助我,誰能定我有罪呢?)

受若不一定等於失去依靠(詩篇31:14耶和華啊,我仍要倚靠你；我說:「你是我的上帝.」31:15我終生的事在你手中,求你救我脫離仇敵的手和那些迫害我的人.)

受苦可以不失盼望(腓立比書2:9所以上帝把他升為至高,又賜給他超乎萬名之上的名,2:10使一切在天上的,地上的和地底下的,因耶穌的名,眾膝都要跪下,2:11眾口都要宣認:耶穌基督是主,歸榮耀給父上帝.).



\newpage

\section{願你們平安}
\label{sec:691}
\textbf{張漢強牧師 - 「願你們平安」}
\newline
\newline
\hyperref[sec:690]{< < < PREV SERMON < < <}
~
\hyperref[sec:toc]{[返目錄]}
~
\hyperref[sec:692]{> > > NEXT SERMON > > >}


馬太福音 28:1-15
\newline
\begin{longtable}{cl}
\hline
\hline
章節 & 經文 (和合本修訂版)\\
\hline
28:1 & \begin{tabularx}{0.7\textwidth}{X} 安息日過後，七日的第一日，天快亮的時候，抹大拉的馬利亞和另一個馬利亞來看墳墓。 \end{tabularx} \\ \\ \relax
28:2 & \begin{tabularx}{0.7\textwidth}{X} 忽然，地大震動；因為有主的一個使者從天上下來，把石頭滾開，坐在上面。 \end{tabularx} \\ \\ \relax
28:3 & \begin{tabularx}{0.7\textwidth}{X} 他的相貌如同閃電，衣服潔白如雪。 \end{tabularx} \\ \\ \relax
28:4 & \begin{tabularx}{0.7\textwidth}{X} 看守的人嚇得渾身顫抖，甚至和死人一樣。 \end{tabularx} \\ \\ \relax
28:5 & \begin{tabularx}{0.7\textwidth}{X} 天使回應婦女說：「不要害怕！我知道你們是尋找那釘十字架的耶穌。 \end{tabularx} \\ \\ \relax
28:6 & \begin{tabularx}{0.7\textwidth}{X} 他不在這裡，照他所說的，他已經復活了。你們來！看看安放他的地方。 \end{tabularx} \\ \\ \relax
28:7 & \begin{tabularx}{0.7\textwidth}{X} 快去告訴他的門徒，說他已從死人中復活了，並且要比你們先到加利利去，在那裡你們會看見他。看哪！我已經告訴你們了。」 \end{tabularx} \\ \\ \relax
28:8 & \begin{tabularx}{0.7\textwidth}{X} 婦女們急忙離開墳墓，又害怕，又大為歡喜，跑去告訴他的門徒。 \end{tabularx} \\ \\ \relax
28:9 & \begin{tabularx}{0.7\textwidth}{X} 忽然，耶穌迎上她們，說：「平安！」她們就上前抱住他的腳拜他。 \end{tabularx} \\ \\ \relax
28:10 & \begin{tabularx}{0.7\textwidth}{X} 耶穌對她們說：「不要害怕！你們去告訴我的弟兄，叫他們往加利利去，在那裡會見到我。」 \end{tabularx} \\ \\ \relax
28:11 & \begin{tabularx}{0.7\textwidth}{X} 她們去的時候，看守的兵有幾個進城去，把所發生的事都報告祭司長。 \end{tabularx} \\ \\ \relax
28:12 & \begin{tabularx}{0.7\textwidth}{X} 祭司長和長老聚集商議，就拿許多銀錢給士兵， \end{tabularx} \\ \\ \relax
28:13 & \begin{tabularx}{0.7\textwidth}{X} 說：「你們要這樣說：『夜間我們睡覺的時候，他的門徒來把他偷去了。』 \end{tabularx} \\ \\ \relax
28:14 & \begin{tabularx}{0.7\textwidth}{X} 若是這話被總督聽見，有我們勸他，保你們無事。」 \end{tabularx} \\ \\ \relax
28:15 & \begin{tabularx}{0.7\textwidth}{X} 士兵收了銀錢，就照所囑咐他們的去做。這話就在猶太人中間流傳，直到今日。 \end{tabularx} \\ \\
[1ex]
\hline
\hline
\end{longtable}
講題:「願你們平安」 “Peace be with you”

經文:馬太福音28:1-15

講員:張漢強牧師

場合:香港中文大學崇基學院禮拜堂主日崇拜

日期:2024年3月31日



\newpage

\section{領受聖靈就會看見}
\label{sec:692}
\textbf{高國雄牧師 - 領受聖靈,就會看見}
\newline
\newline
\hyperref[sec:691]{< < < PREV SERMON < < <}
~
\hyperref[sec:toc]{[返目錄]}
~
\hyperref[sec:693]{> > > NEXT SERMON > > >}


約翰福音 20:19-31
\newline
\begin{longtable}{cl}
\hline
\hline
章節 & 經文 (和合本修訂版)\\
\hline
20:19 & \begin{tabularx}{0.7\textwidth}{X} 那日（就是七日的第一日）晚上，門徒因怕猶太人，所在的地方門都關了。耶穌來，站在當中，對他們說：「願你們平安！」 \end{tabularx} \\ \\ \relax
20:20 & \begin{tabularx}{0.7\textwidth}{X} 說了這話，他把手和肋旁給他們看。門徒一看見主就喜樂了。 \end{tabularx} \\ \\ \relax
20:21 & \begin{tabularx}{0.7\textwidth}{X} 於是耶穌又對他們說：「願你們平安！父怎樣差遣了我，我也照樣差遣你們。」 \end{tabularx} \\ \\ \relax
20:22 & \begin{tabularx}{0.7\textwidth}{X} 說了這話，他向他們吹一口氣，說：「領受聖靈吧！ \end{tabularx} \\ \\ \relax
20:23 & \begin{tabularx}{0.7\textwidth}{X} 你們赦免誰的罪，誰的罪就得赦免；你們不赦免誰的罪，誰的罪就不得赦免。」 \end{tabularx} \\ \\ \relax
20:24 & \begin{tabularx}{0.7\textwidth}{X} 那十二使徒中，有個叫低土馬的多馬，耶穌來的時候，他沒有和他們在一起。 \end{tabularx} \\ \\ \relax
20:25 & \begin{tabularx}{0.7\textwidth}{X} 其他的門徒就對他說：「我們已經看見主了。」多馬卻對他們說：「除非我看見他手上的釘痕，用我的指頭探入那釘痕，用我的手探入他的肋旁，我絕不信。」 \end{tabularx} \\ \\ \relax
20:26 & \begin{tabularx}{0.7\textwidth}{X} 過了八日，門徒又在屋裡，多馬也和他們在一起。門都關了，耶穌來，站在當中，說：「願你們平安！」 \end{tabularx} \\ \\ \relax
20:27 & \begin{tabularx}{0.7\textwidth}{X} 然後他對多馬說：「把你的指頭伸到這裡來，看看我的手；把你的手伸過來，探入我的肋旁。不要疑惑，總要信！」 \end{tabularx} \\ \\ \relax
20:28 & \begin{tabularx}{0.7\textwidth}{X} 多馬回答，對他說：「我的主！我的神！」 \end{tabularx} \\ \\ \relax
20:29 & \begin{tabularx}{0.7\textwidth}{X} 耶穌對他說：「你因為看見了我才信嗎？那沒有看見卻信的有福了。」 \end{tabularx} \\ \\ \relax
20:30 & \begin{tabularx}{0.7\textwidth}{X} 耶穌在他門徒面前另外行了許多神蹟，沒有記錄在這書上。 \end{tabularx} \\ \\ \relax
20:31 & \begin{tabularx}{0.7\textwidth}{X} 但記載這些事是要使你們信耶穌是基督，是神的兒子，並且使你們信他，好因著他的名得生命。 \end{tabularx} \\ \\
[1ex]
\hline
\hline
\end{longtable}
講題:領受聖靈,就會看見 Receive the Holy Spirit, and You Will See

經文:約翰福音20:19-31

講員:高國雄牧師

場合:香港中文大學崇基學院禮拜堂主日崇拜

日期:2024年4月7日

約20:19-31

19那日(就是七日的第一日)晚上,門徒所在的地方,因怕猶太人,門都關了.耶穌來,站在當中,對他們說:「願你們平安!」 20說了這話,就把手和肋旁指給他們看.門徒看見主,就喜樂了. 21耶穌又對他們說:「願你們平安!父怎樣差遣了我,我也照樣差遣你們.」 22說了這話,就向他們吹一口氣,說:「你們受聖靈! 23你們赦免誰的罪,誰的罪就赦免了；你們留下誰的罪,誰的罪就留下了.」

24那十二個門徒中,有稱為低土馬的多馬；耶穌來的時候,他沒有和他們同在. 25那些門徒就對他說:「我們已經看見主了.」多馬卻說:「我非看見他手上的釘痕,用指頭探入那釘痕,又用手探入他的肋旁,我總不信.」 26過了八日,門徒又在屋裡,多馬也和他們同在,門都關了.耶穌來,站在當中說:「願你們平安!」 27就對多馬說:「伸過你的指頭來,摸我的手；伸出你的手來,探入我的肋旁.不要疑惑,總要信!」 28多馬說:「我的主!我的上帝!」 29耶穌對他說:「你因看見了我才信；那沒有看見就信的有福了.」

30耶穌在門徒面前另外行了許多神蹟,沒有記在這書上. 31但記這些事要叫你們信耶穌是基督,是上帝的兒子,並且叫你們信了他,就可以因他的名得生命.

顯現的意義

為何「信經」沒有提到?不重要?

顯現是復活的重要部份

肉身復活的示範,基督的特徵,釘痕乃完美的記號

能與門徒日常生活,又同時擁有超越的能力(隨時出現及離開)

有明顯的使命:顯現的對象,升天的安排…

堅固使徒的信心,向門徒傳遞使命,就是為了今天的你和我

“多疑的多馬”?

“多馬不信”, “多馬釋疑”(和合本)

“耶穌向疑惑的多馬顯現”(新譯本)

“耶穌和多馬”(新標點和合本)

“耶穌向多馬顯現”

由“未能經歷”開始V24

理性的要求,合情合理,比其他門徒更認真更謹慎V25

耶穌為多馬顯現,耶穌特意向多馬顯現,耶穌極愛多馬V26

耶穌沒有責備多馬,反要堅固他的信心V27

多馬以信心回應(沒有真正伸手!)V28

多馬比其他門徒少了福份嗎?V29

“沒有看見的” > 就是今天的我們!!

門徒領受了聖靈,領受了赦罪的權柄,領受了生命

多馬等了一星期,我們呢?

我們今天看見嗎?

啟真道博士Dr. LG. Kilborn

(7/4/2024)129歲生日

1955年出任崇基學院校董會主席

1956年為崇基學院新校舍奠基

1960年受聘於崇基學院任副院長,不久成為組建委員會的主席,並引導制定了三個獨立學院(新亞書院,聯合書院和崇基學院)的學術標準.

為禮拜堂籌款

崇基學院校董會會議室:啟真道室



\newpage

\section{復活基督的顯現}
\label{sec:693}
\textbf{李均熊牧師 - 復活基督的顯現}
\newline
\newline
\hyperref[sec:692]{< < < PREV SERMON < < <}
~
\hyperref[sec:toc]{[返目錄]}
~
\hyperref[sec:694]{> > > NEXT SERMON > > >}


路加福音 24:36-48
\newline
\begin{longtable}{cl}
\hline
\hline
章節 & 經文 (和合本修訂版)\\
\hline
24:36 & \begin{tabularx}{0.7\textwidth}{X} 正說這些話的時候，耶穌親自站在他們當中，說：「願你們平安！」 \end{tabularx} \\ \\ \relax
24:37 & \begin{tabularx}{0.7\textwidth}{X} 他們卻驚慌害怕，以為所看見的是魂。 \end{tabularx} \\ \\ \relax
24:38 & \begin{tabularx}{0.7\textwidth}{X} 耶穌對他們說：「你們為甚麼驚恐不安？為甚麼心裡起疑惑呢？ \end{tabularx} \\ \\ \relax
24:39 & \begin{tabularx}{0.7\textwidth}{X} 你們看我的手和我的腳，就知道實在是我了。摸摸我，看，因為魂無骨無肉，你們看，我是有的。」 \end{tabularx} \\ \\ \relax
24:40 & \begin{tabularx}{0.7\textwidth}{X} 說了這話，他就把手和腳給他們看。 \end{tabularx} \\ \\ \relax
24:41 & \begin{tabularx}{0.7\textwidth}{X} 他們還在又驚又喜、不敢相信的時候，耶穌對他們說：「你們這裡有甚麼吃的沒有？」 \end{tabularx} \\ \\ \relax
24:42 & \begin{tabularx}{0.7\textwidth}{X} 他們給了他一片烤魚， \end{tabularx} \\ \\ \relax
24:43 & \begin{tabularx}{0.7\textwidth}{X} 他接過來，在他們面前吃了。 \end{tabularx} \\ \\ \relax
24:44 & \begin{tabularx}{0.7\textwidth}{X} 耶穌對他們說：「這就是我從前和你們同在時所告訴你們的話：摩西的律法、先知的書，和《詩篇》上所記一切指著我的話都必須應驗。」 \end{tabularx} \\ \\ \relax
24:45 & \begin{tabularx}{0.7\textwidth}{X} 於是耶穌開他們的心竅，使他們能明白聖經， \end{tabularx} \\ \\ \relax
24:46 & \begin{tabularx}{0.7\textwidth}{X} 又對他們說：「照經上所寫的，基督必受害，第三天從死人中復活， \end{tabularx} \\ \\ \relax
24:47 & \begin{tabularx}{0.7\textwidth}{X} 並且人們要奉他的名傳悔改、使罪得赦的道，從耶路撒冷起直傳到萬邦。 \end{tabularx} \\ \\ \relax
24:48 & \begin{tabularx}{0.7\textwidth}{X} 你們就是這些事的見證。 \end{tabularx} \\ \\
[1ex]
\hline
\hline
\end{longtable}
講題:復活基督的顯現 Appearance of the Risen Christ

經文:路加福音24章36至48節

講員:李均熊牧師

場合:香港中文大學崇基學院禮拜堂主日崇拜

日期:2024年4月14日

[36] 正說這話的時候,耶穌親自站在他們當中,說:「願你們平安!」 [37] 他們卻驚慌害怕,以為所看見的是魂. [38] 耶穌說:「你們為甚麼愁煩?為甚麼心裡起疑念呢? [39] 你們看我的手,我的腳,就知道實在是我了.摸我看看!魂無骨無肉,你們看,我是有的.」 [40] 說了這話,就把手和腳給他們看. [41] 他們正喜得不敢信,並且希奇；耶穌就說:「你們這裡有甚麼吃的沒有?」 [42] 他們便給他一片燒魚.  [43] 他接過來,在他們面前吃了.

[44] 耶穌對他們說:「這就是我從前與你們同在之時所告訴你們的話說:摩西的律法,先知的書,和詩篇上所記的,凡指著我的話都必須應驗.」 [45] 於是耶穌開他們的心竅,使他們能明白聖經, [46] 又對他們說:「照經上所寫的,基督必受害,第三日從死裡復活, [47] 並且人要奉他的名傳悔改,赦罪的道,從耶路撒冷起直傳到萬邦. [48] 你們就是這些事的見證.

起:身體,不適

今個星期人感到不適.星期初先是感到倦怠,喉嚨痕癢.到星期中間倦怠依然,但喉嚨痕癢卻變成咳嗽,而且晚間睡覺時比日間咳嗽更利害.雖是幾天之間,咳嗽而經減少並消失,但那種好像有螞蟻在咽喉間爬行的感覺其實十分難受,恨不得能將手掌放入喉嚨內給抓癢.另一個問題是倦困.早上因為晚上咳嗽已經感到睡眠不足,很想午間來個少睡補眠,可惜仍要上課；到下午上完一輪課後,人就感到有氣無力,連咳也好像有心無力般.所以今次的講章也特別遲才完成,為此要向傳譯的姊妹們致歉.

當你平時健康無病無痛時,你可能沒有留意身體的機能,你自然的呼吸,走路,吃喝,睡訓,甚至大小二便,都沒有特別的意識.但當你像我人到中年,又或身體不適,甚至因各種遭遇身體漸漸不聽使喚,不知你會否幻想靈魂離開身體,不會因疾病感到不適,也不會肚餓,肚痛,內急,痕癢,疲倦,疼痛?作為一個經常要坐在案頭讀書寫字的神學院老師,久坐一會就希望這個身體,這個背脊,這雙長腿,這個屁股可以消失,只是剩下一個清明的腦袋在思考,不需感到肚餓疲倦.現在人年紀大了,不像後生時可以連續幾個鐘頭集中精神去讀書,有時打開書本看了一兩頁,眼皮就跌下來了.眼皮自閉癥這個癥狀好像是沒有藥醫的.

承:肉身的顯現

今天的福音經課剛好是耶穌復活後向門徒顯現的記載,路加福音的經文再三強調耶穌向門徒確定復活後的祂不是鬼魂,而是有身體的人,更特意向門徒要了一條魚,好像就是要證明祂是一個不折不扣身體健全的人.

在聖誕節時我們會慶祝基督道成了肉身,而在復活節我們再一次慶祝耶穌肉身的復活.基督教其實是一個十分重視肉身的宗教.但有些宗教看身體,是煩惱的源頭,慾望的源頭,罪惡的源頭.正所謂色即是空,空即是色,甚至希望透過修練能夠超脫這個身體,又或超脫這個身體所帶來的煩惱.但你看,耶穌已經戰勝了死亡,脫離了痛苦,再沒有了眼淚；但祂復活之後,卻仍然帶著那個有釘痕,傷口,需要飲食的身體.為什麼呢?耶穌既然從死裡復活了,為什麼仍然要帶著這個身體呢?

有學者認為,路加福音在這裡特別提到耶穌復活強調自己的身體,向門徒證明自己是有肉身的人,是要回應當時群體當中一類諾斯底主義,認為正正是因為身體帶來邪惡,身體是不好的東西,所以若不苦待自己希望能使自己靈魂得救,就會是放縱自己感觀的享受,因為身體在上帝救贖的計劃內並沒有任何重要的位置.故此,路加福音在耶穌復活顯現的當兒,再一次強調耶穌戰勝死亡,同時亦救贖了肉身,神所賜的身體帶來的感觀和慾求是個好東西,是可以被救贖,也應該被救贖的生命不可分割的部份.

轉:復活主提供的新經驗引發的領悟

姑勿論路加福音寫作的本意,他所身處的群體面對當時的思潮是什麼,今日福音經課都提醒我們,上帝的救贖是包括這一個在我們身上會漸漸衰殘的身體.魯益師(C. S. Lewis)曾經留下一段說話:「不是一種抽象的東西叫作『人性』要去救贖.而是你……你的靈魂,甚至是(從某個角度來看我們還需要多了解的)你的身體,都是朝向那崇高而神聖之處的.」(“It is not an abstraction called Humanity that is to be saved. It is you, … your soul, and, in some sense yet to be understood, even your body, that was made for the high and holy place.”)

但若我們再仔細看一看今日的經課,耶穌肉身的復活,是十分不容易被認出的一件事情.

大家可以留意,這一段經文經常出現的一個動詞,就是「看」:

[36] 正說這話的時候,耶穌親自站在他們當中,說:「願你們平安!」 [37] 他們卻驚慌害怕,以為所看見的是魂. [38] 耶穌說:「你們為甚麼愁煩?為甚麼心裡起疑念呢? [39] 你們看我的手,我的腳,就知道實在是我了.摸我看看!魂無骨無肉,你們看,我是有的.」 [40] 說了這話,就把手和腳給他們看.

若我們只抽這一段經文來看,好像十分合理:門徒見到從死裡復活的主,以為耶穌是鬼魂,大家應該都可以理解.但是,大家若讀讀之前的記載,就會覺得門徒反應的遲鈍和奇怪.門徒晚間一起聚集時,早上已經離開往以馬忤斯的革流巴和另一個門徒剛好匆匆歸來,走到他們中間,他們就彼此述說耶穌如何先後向他們顯現.就算彼得遇見主時間太匆匆不太了解情況,但革流巴他們二人差不多和耶穌同行了半天,更一同坐席用餐,當他們向門徒述說耶穌向他們顯現的時候,應該已經有一定程度對耶穌復活後的情況有些微了解,甚麼耶穌再次向大家顯現,大家就會驚惶失惜,以為耶穌是鬼魂呢?

OK,可能你會覺得我批評門徒當時不應如此反應有些過分了；門徒突然見到在他們親眼看見死在十架上的耶穌在他們面前顯現,有些慌失失實為人之常情.不過,路加福音中的經文不單強調耶穌肉身的復活,更多關注就是這一個身體復活的主並不容易被認出.在往以馬忤斯路上的兩個人,不就是和復活主談了大半天,仍然不知祂就是復活的基督嗎?

有可能門徒心目中的耶穌,仍然是之前見到受盡酷刑苦頭已經面目全非的耶穌,但現在眼前突然出然的人物青靚白淨,認不出祂就是受死復活的主也不為過.我們是否對我們所信仰的基督,帶著太多前理解/偏見,受我們經驗的限制,因而看不見耶穌復活的身體呢?

耶穌復活後的身體,是要我們重新認識的身體.耶穌的肉身,但祂卻看似再不受時空限制.個同時間,祂的肉身仍然帶著受苦受傷的印記.不錯,耶穌的救贖也是包括我們的身體,但我們這個身體也是限制我們經驗的媒介,當中我們的感觀和我視野的限制,並不能百分百使我們弄清事情的意義(路加福音中耶穌顯現也強調祂要向門徒解釋經上指著祂的內容是甚麼意思,但今日就不多談了).為甚麼耶穌復活的身體不是完好無缺的身體呢?可能耶穌就是讓門徒可以認得出祂,而要保留自己的釘痕手和肋旁傷.復活的主的顯現,不單要讓我們知道祂是真的復活了,也同時要提供新的角度給我們了接受苦和復活的生命是如何一回事(約一3:2「親愛的,我們現在是神的兒女,將來如何還未顯明.我們所知道的是:基督顯現的時候,我們會像他,因為我們將見到他的本相.」).復活是帶著新的生命,去理解過去的傷痕和遭遇.在復活生命的領悟下,我們才能看得清這個將殘肉體本相的真義,也可以在不限於自己理解和經驗的限制,看出這被救贖的身體可能帶來的新意義.

合:見證--不是得個講字

我們大學人通常被批評活在象牙塔,滿口理論,堅離地.基督徒信耶穌也常被批評得個講字,尤其那些口口聲聲要信仰要反思,要讀神學的人,更是不解世情.但當復活的主耶穌顯現的最後,強調門徒就是這些事的見證.唯有經歷過耶穌的復活,才能真正見證這些事.就好比我們說見證一場球賽,不會就只是看著電視,甚至聽電台足球旁述,就說自己見證了那場球賽.就算你不是落場踢的選手,也起碼是親身入場,親身經歷球賽,才好說你見證那場球賽.

我們要見證的復活主,是抽像理念,只存在於思想世界,還是我們有血有肉的生命曾經經歷有血有肉的耶穌?



\newpage

\section{捨去與重拾}
\label{sec:694}
\textbf{高國雄牧師 - 捨去與重拾}
\newline
\newline
\hyperref[sec:693]{< < < PREV SERMON < < <}
~
\hyperref[sec:toc]{[返目錄]}
~
\hyperref[sec:695]{> > > NEXT SERMON > > >}


約翰福音 10:11-18
\newline
\begin{longtable}{cl}
\hline
\hline
章節 & 經文 (和合本修訂版)\\
\hline
10:11 & \begin{tabularx}{0.7\textwidth}{X} 「我是好牧人，好牧人為羊捨命。 \end{tabularx} \\ \\ \relax
10:12 & \begin{tabularx}{0.7\textwidth}{X} 雇工不是牧人，羊不是他自己的，他一看見狼來，就撇下羊群逃跑；狼抓住羊，把牠們趕散。 \end{tabularx} \\ \\ \relax
10:13 & \begin{tabularx}{0.7\textwidth}{X} 雇工逃走，因為他是雇工，對羊毫不關心。 \end{tabularx} \\ \\ \relax
10:14 & \begin{tabularx}{0.7\textwidth}{X} 我是好牧人；我認識我的羊，我的羊也認識我， \end{tabularx} \\ \\ \relax
10:15 & \begin{tabularx}{0.7\textwidth}{X} 正如父認識我，我也認識父一樣；並且我為羊捨命。 \end{tabularx} \\ \\ \relax
10:16 & \begin{tabularx}{0.7\textwidth}{X} 我另外有羊，不屬這圈裡的，我必須領牠們來，牠們也要聽我的聲音，並且要合成一群，歸一個牧人。 \end{tabularx} \\ \\ \relax
10:17 & \begin{tabularx}{0.7\textwidth}{X} 為此，我父愛我，因為我把命捨去，好再取回來。 \end{tabularx} \\ \\ \relax
10:18 & \begin{tabularx}{0.7\textwidth}{X} 沒有人奪去我的命，是我自己捨的；我有權捨棄，也有權再取回。這是我從我父所受的命令。」 \end{tabularx} \\ \\
[1ex]
\hline
\hline
\end{longtable}
講題:捨去與重拾    Lay It Down and Take It Up

經文:約翰福音10章11至18節

講員:高國雄牧師

場合:香港中文大學崇基學院禮拜堂主日崇拜

日期:2024年4月21日

11我是好牧人；好牧人為羊捨命. 12若是雇工,不是牧人,羊也不是他自己的,他看見狼來,就撇下羊逃走；狼抓住羊,趕散了羊群. 13雇工逃走,因他是雇工,並不顧念羊. 14我是好牧人；我認識我的羊,我的羊也認識我, 15正如父認識我,我也認識父一樣；並且我為羊捨命. 16我另外有羊,不是這圈裡的；我必須領他們來,他們也要聽我的聲音,並且要合成一群,歸一個牧人了. 17我父愛我；因我將命捨去,好再取回來. 18沒有人奪我的命去,是我自己捨的.我有權柄捨了,也有權柄取回來.這是我從我父所受的命令.」

「斷,捨,離」

人生不斷經歷「捨棄」與「被捨棄」

不願意,不明白,不相信,不接受

捨棄(Give up/Abandon)≠放棄

主動,自願犧牲,為了更大的目標/得著

有什麼我願意為主捨棄的呢?

經課前

約翰福音10:10盜賊來,無非要偷竊,殺害,毀壞；我來了,是要叫羊得生命,並且得的更豐盛.

耶穌教導後

約翰福音10:19-21

19 猶太人為這些話又起了紛爭. 20內中有好些人說:「他是被鬼附著,而且瘋了,為甚麼聽他呢?」 21又有人說:「這不是鬼附之人所說的話.鬼豈能叫瞎子的眼睛開了呢?」

超出預期

不合理,不可能,無法接受

起了紛爭

同樣的情況曾經出現…

約翰福音7:40-44

40眾人聽見這話,有的說:「這真是那先知.」 41有的說:「這是基督.」但也有的說:「基督豈是從加利利出來的嗎? 42經上豈不是說『基督是大衛的後裔,從大衛本鄉伯利恆出來的』嗎?」 43於是眾人因著耶穌起了紛爭. 44其中有人要捉拿他,只是無人下手.

曾因耶穌的身份(先知?基督?)起紛爭

計劃要下毒手

約翰福音9:16

16法利賽人中有的說:「這個人不是從上帝來的,因為他不守安息日.」又有人說:「一個罪人怎能行這樣的神蹟呢?」他們就起了紛爭.

再次為耶穌的身份(從上帝而來VS罪人)起紛爭

回到今天經課約翰福音10:11-18

耶穌這次如何論述自己的身份,以致又引起紛爭呢?

猶太人眼中牧人的角色

好牧人VS不好/一般/壞牧人

不好/一般/壞牧人 >> 雇工

關係,角色(Rode/relationship)≠崗位/工作內容(position)

耶穌這好牧人的特色

羊是自己的

顧念羊(無論什麼情況)

認識羊

互相認識

建立羊的群體

願意為羊捨命

真的為羊捨了命!

好牧人捨命,從狼群中拾回了迷失的羊

耶穌基督捨命,從邪惡世代中魔鬼手上拾回了失落了的我們

更讓我們重拾了愛

我們也要效法主,捨棄而重拾身邊的弟兄

約翰一書3:16-18

16主為我們捨命,我們從此就知道何為愛；我們也當為弟兄捨命. 17凡有世上財物的,看見弟兄窮乏,卻塞住憐恤的心,愛上帝的心怎能存在他裡面呢? 18小子們哪,我們相愛,不要只在言語和舌頭上,總要在行為和誠實上.



\newpage

\section{往加沙的路上}
\label{sec:695}
\textbf{林豪恩先生 - 往加沙的路上}
\newline
\newline
\hyperref[sec:694]{< < < PREV SERMON < < <}
~
\hyperref[sec:toc]{[返目錄]}
~
\hyperref[sec:696]{> > > NEXT SERMON > > >}


講題:往加沙的路上   On the Road to Gaza

講員:

場合:香港中文大學崇基學院禮拜堂主日崇拜

日期:2024年4月28日

一．踏上往加沙之路

使徒行傳8:26如此為這個故事做背景介紹:

有主的一個使者對腓利說:「起來!向南走,往那從耶路撒冷下迦薩的路上去.」

迦薩是甚麼地方呢?如果用另一個翻譯—加沙,相信大家都聽過了.加沙位於巴勒斯坦西岸,連接以色列及埃及的沿海地帶.為甚麼加沙最近被世界所關注呢?因為去年10月,武裝分子從加沙出來襲擊以色列南部,擄走人質,隨後以色列軍隊攻入加沙,開始了這場到今日仍未結束的悲劇.

如果今日有人邀請去加沙,你去不去呢?不用擔心,因為一般人根本不可能去加沙.自從2007年開始,以色列嚴密封鎖加沙.另一邊的埃及亦關閉加沙南部的邊界.加沙變成一個生活艱苦的地方,居民只能依靠國際救濟；一個衝突頻繁的地方.沒有人能夠自由離開這個地方,一般人也不會進入這個地方.不過,曾幾何時,加沙是一處自由出入的地方,甚至是一條來往南北常用的通道,人們稱之為加沙走廊.

二．往加沙的路,曾經是尋道者的通道

8:27有一個衣索匹亞人,是個有大權的太監,在衣索匹亞女王甘大基的手下總管銀庫,他上耶路撒冷禮拜去了.

8:28現在回來,在車上坐著,念先知以賽亞的書.

為甚麼有位衣索匹亞太監在這條路上出現呢?衣索匹亞,和修版聖經譯為埃塞俄比亞,其實並不等於今日的埃塞俄比亞共和國,可能是舊約中古實的所在地,位於埃及的南方及蘇丹的北部.相傳列王記上10章那位來拜訪所羅門王的示巴女王就是從這地區而來.

經文指出「他上耶路撒冷禮拜去了.現在回來」.為甚麼有人從衣索匹亞去耶路撒冷禮拜呢?幾百年前,猶大以色列地區曾被亞述和巴比倫入侵,戰亂的時候,人們被迫逃離家園,有些人南下來到此地定居,把猶太教信仰帶到此處.這人坐馬車來回,應該不是普通人,經文指出是一位「有大權的」,在朝廷掌管財政的高官.

現在的加沙被封鎖為一個禁區,無人能夠自由進出.當時,加沙是一條走廊,方便人們從非洲東部來往耶路撒冷朝聖.兩種狀況作為對比,你認為加沙作為朝聖者的通道,還是被封鎖成為禁區,更能夠成為人們的祝福呢?看著現在被封鎖的加沙,不禁令人惋惜.或者,加沙仍然期盼自己可以再次成為尋道者的通道.

三．往加沙的路,曾經是傳道者的通道

8:29聖靈對腓利說:「你去!貼近那車走.」

8:30腓利就跑到太監那裡,聽見他念先知以賽亞的書,便問他說:「你所念的,你明白嗎?」

8:31他說:「沒有人指教我,怎能明白呢?」於是請腓利上車,與他同坐.

8:32他所念的那段經,說:他像羊被牽到宰殺之地,又像羊羔在剪毛的人手下無聲；他也是這樣不開口.

8:33他卑微的時候,人不按公義審判他；誰能述說他的世代?因為他的生命從地上奪去.

8:34太監對腓利說:「請問,先知說這話是指著誰?是指著自己呢?是指著別人呢?」

8:35腓利就開口從這經上起,對他傳講耶穌.

故事中另外一位出現在往加沙的路上的是腓利.腓利在往加沙的路上出現,遠因是使徒行傳8:1交代的背景:「耶路撒冷的教會大遭逼迫,除了使徒以外,門徒都分散在猶太和撒馬利雅各處.」近因是上文所說主的使者的吩咐.腓利這次下加沙,經文記載只做了一件事,就是剛才讀過的內容,向那位從衣索匹亞來的人傳講耶穌.

現在,奔跑於往加沙路上的人所為何事呢?殺戮嗎?戰鬥嗎?報復嗎?曾經何時,往加沙的路,是傳道的路,是拯救的路.兩種狀態相比,哪一種更為美好呢?往加沙的路,何時可以再成為傳道者的通道呢?

四．往加沙的路,曾經是受洗歸主的驛站

8:36二人正往前走,到了有水的地方,太監說:「看哪,這裡有水,我受洗有甚麼妨礙呢?」

8:38於是吩咐車站住,腓利和太監二人同下水裡去,腓利就給他施洗.

今時今日,如果有兩個來自不同地方,不同種族的人在往加沙的路上偶然相遇,會產生甚麼結果呢?當年那位從衣索匹亞來的人和腓利在往加沙的路上相遇,締造出既神聖又充滿生機的時刻.神聖,是因為水禮的舉行.水禮是神聖的,其信仰意義及對受洗者的意義都非常豐富；水禮同時是生機的象徵.保羅在羅馬書6:4說:「所以,我們藉著洗禮歸入死,和他一同埋葬,原是叫我們一舉一動有新生的樣式,像基督藉著父的榮耀從死裡復活一樣.」

今日,加沙是一個甚麼地方呢?外面的人不想進去,裡面的人想走也逃不出來.加沙變成了人間地獄,燃料,食物,食水與醫療資源嚴重短缺,供電不足,醫院,污水處理與海水淡化廠等基礎建設皆停擺.其實,加沙有教會,有基督徒,本來也有神聖而充滿生機的水禮舉行.現在,是誰把神聖變為地獄呢?是甚麼把生機變為死亡呢?加沙這條走廊,以及在其上的教會,何時才能恢復生機?何時才能讓神聖再顯現呢?

五．來回而不佔領加沙

8:39從水裡上來,主的靈把腓利提了去,太監也不再見他了,就歡歡喜喜地走路.

8:40後來有人在亞鎖都遇見腓利；他走遍那地方,在各城宣傳福音,直到凱撒利亞.

有些人踏上往加沙的路,為了攻佔加沙；有些人踏上往加沙的路,為了把加沙據為已有.當年那位被經文指明「有大權的」衣索匹亞大臣只是經過加沙,沒有運用他的財力和權力強奪加沙；當年那位來無蹤去無影,充滿神力異能的腓力只是踏上往加沙的路,沒有運用他的神力異能控制加沙.如果有權力和能力的人士不強奪加沙,讓它開放成為來往南北的安全通道,讓尋道者和傳道者通過和相遇,這會不會比現在的狀況更美好呢?

這個世界上,除了加沙被稱為走廊,大家還想到有甚麼地方都扮演著類似的橋樑角色嗎?當世界各地的人才,資金,貨物,資訊都能夠自由出入,這地就成為文化交流,知識傳遞,民族共融的多元並且充滿生機的地方.或者,有些地方能夠對當地居民及世界的作出貢獻,就是繼續成為開放及自由的通道,讓不同的人在這裡相遇,併發出美麗及創新的景象.

四．反省

使徒行傳告訴我們,加沙曾經是尋道者和傳道者的通道.讓我們以通道,或者道路這個主題作一點信仰思考.

1．作為道路的耶穌基督(約翰福音)

14:1「你們心裡不要憂愁；你們信上帝,也當信我.

14:2在我父的家裡有許多住處；若是沒有,我就早已告訴你們了.我去原是為你們預備地方去.

14:3我若去為你們預備了地方,就必再來接你們到我那裡去,我在哪裡,叫你們也在那裡.

14:4我往哪裡去,你們知道；那條路,你們也知道.」

14:5多馬對他說:「主啊,我們不知道你往哪裡去,怎麼知道那條路呢?」

14:6耶穌說:「我就是道路,真理,生命；若不藉著我,沒有人能到父那裡去.

這是耶穌和門徒分離之前給門徒的安慰.門徒不感到被安慰,因為門徒根本不知道那條路,也沒有看見那條路.就正我們有時遇到的狀況,看不見有路.耶穌指出祂自己就是這條路.還記得拿但業初遇耶穌的時候,他們的對話嗎?

1:49拿但業說:「拉比,你是上帝的兒子,你是以色列的王!」

1:50耶穌對他說:「因為我說『在無花果樹底下看見你』,你就信嗎?你將要看見比這更大的事」；

1:51又說:「我實實在在地告訴你們,你們將要看見天開了,上帝的使者上去下來在人子身上.」

那條路在哪裡呢?耶穌指出那條路在「人子身上」.接著,讓我們以耶穌基督作為道路的主題思考一下教會及信徒的使命.

2．作為道路的教會及信徒

有一首詩歌名為【流通管子】(生命聖詩369),給我們以下的提示:

我要頌揚寶貴救主,你的愛使我歸你,

你潔淨我又充滿我,使我成流通管子.

主啊,我願成為管子,讓你奇妙大能力；

充滿我心,藉我流出,每時每刻不停止.

我願倒空讓主充滿,成為主潔淨器皿；

我無力量,惟主加添,能完成主的命令.

主啊,我願成為管子,讓你奇妙大能力；

充滿我心,藉我流出,每時每刻不停止.

3．作為道路成就別人的人生

有一首歌名為「橋」(頌主聖詩533),或許也是一個引發我們思考的例子:

父母親從小撫育我,栽種愛惜幼樹苗,維護照料牽我手；

父母親延伸天父愛,甘心獻身像道橋,禍困凶險我不怕過.

父母親年紀雖漸老,一生獻身像道橋,

容讓你我在上面輕輕鬆鬆再走過,再經過,不分風雨,渡過!

迦薩之路,成就了腓利和衣索匹亞大臣的信仰故事；耶穌作為道路,成就了人與天父的和好；教會和信徒會否也可以用「一條路」來想像自己,延續耶穌基督的使命呢?作為個人,如果我們也用「一條路」來想像自己,我們想要成就甚麼呢?



\newpage

\section{常在主愛中}
\label{sec:696}
\textbf{高國雄牧師 - 常在主愛中}
\newline
\newline
\hyperref[sec:695]{< < < PREV SERMON < < <}
~
\hyperref[sec:toc]{[返目錄]}
~
\hyperref[sec:697]{> > > NEXT SERMON > > >}


講題:常在主愛中   Always in His Love

講員:高國雄牧師

場合:香港中文大學崇基學院禮拜堂主日崇拜

日期:2024年5月5日

約翰福音 15:9-17

9我愛你們,正如父愛我一樣；你們要常在我的愛裡.10你們若遵守我的命令,就常在我的愛裡,正如我遵守了我父的命令,常在他的愛裡.11「這些事我已經對你們說了,是要叫我的喜樂存在你們心裡,並叫你們的喜樂可以滿足.12你們要彼此相愛,像我愛你們一樣；這就是我的命令.13人為朋友捨命,人的愛心沒有比這個大的.14你們若遵行我所吩咐的,就是我的朋友了.　15以後我不再稱你們為僕人,因僕人不知道主人所做的事.我乃稱你們為朋友；因我從我父所聽見的,已經都告訴你們了.16不是你們揀選了我,是我揀選了你們,並且分派你們去結果子,叫你們的果子常存,使你們奉我的名,無論向父求甚麼,他就賜給你們.17我這樣吩咐你們,是要叫你們彼此相愛.」

V9我愛你們,最直接簡單的描述

因由:正如父愛我一樣,曾經歷過的,體會到的

你們要常在我的愛裡

「你們要與我的愛保持聯合」(原文新約中譯本)

V10遵守命令,「如若」,帶應計/祝福的命令

遵守耶穌基督的命令=與主的愛保持聯合

耶穌基督自己就是最好的示範

V11發出命令使耶穌基督「喜樂」,又因此這「喜樂」就存在我們心裡,並且是完全的「喜樂」.

耶穌基督的命令是什麼呢?

V12彼此相愛

耶穌愛我們,就是最好的示範(又一次)

遵守主的命令 > 成為朋友

耶穌基督是如何遵守天父的命令呢?愛我們到底,為我們捨命

V13人為「朋友」捨命

V14基督要成為我們的朋友…若…

V15-16成為耶穌朋友的好處

朋友 X僕人

知道主所做的事

收到一切由天父來的信息

被主揀選

為主結果子,有成果,有傳遞

能靠主名求天父賜一切

V17 這就是命令,是吩咐「你們要彼此相愛」

約翰一書 5:4-5

4因為凡從上帝生的,就勝過世界；使我們勝了世界的,就是我們的信心.5勝過世界的是誰呢?不是那信耶穌是上帝兒子的嗎?

「常在主愛中,與主一起勝過世界!」



\newpage

\section{栽在溪水旁}
\label{sec:697}
\textbf{李均熊牧師 - 栽在溪水旁}
\newline
\newline
\hyperref[sec:696]{< < < PREV SERMON < < <}
~
\hyperref[sec:toc]{[返目錄]}
~
\hyperref[sec:698]{> > > NEXT SERMON > > >}


詩篇 1:1-6
\newline
\begin{longtable}{cl}
\hline
\hline
章節 & 經文 (和合本修訂版)\\
\hline
1:1 & \begin{tabularx}{0.7\textwidth}{X} 不從惡人的計謀，不站罪人的道路，不坐傲慢人的座位， \end{tabularx} \\ \\ \relax
1:2 & \begin{tabularx}{0.7\textwidth}{X} 惟喜愛耶和華的律法，晝夜思想他的律法；這人便為有福！ \end{tabularx} \\ \\ \relax
1:3 & \begin{tabularx}{0.7\textwidth}{X} 他要像一棵樹栽在溪水旁，按時候結果子，葉子也不枯乾。凡他所做的盡都順利。 \end{tabularx} \\ \\ \relax
1:4 & \begin{tabularx}{0.7\textwidth}{X} 惡人並不是這樣，卻像糠秕被風吹散。 \end{tabularx} \\ \\ \relax
1:5 & \begin{tabularx}{0.7\textwidth}{X} 因此，當審判的時候惡人必站立不住，罪人在義人的會眾中也是如此。 \end{tabularx} \\ \\ \relax
1:6 & \begin{tabularx}{0.7\textwidth}{X} 因為耶和華知道義人的道路，惡人的道路卻必滅亡。 \end{tabularx} \\ \\
[1ex]
\hline
\hline
\end{longtable}
約翰福音 17:6-19
\newline
\begin{longtable}{cl}
\hline
\hline
章節 & 經文 (和合本修訂版)\\
\hline
17:6 & \begin{tabularx}{0.7\textwidth}{X} 「你從世上賜給我的人，我已把你的名顯明給他們。他們本是你的，你把他們賜給我，他們也遵守了你的道。 \end{tabularx} \\ \\ \relax
17:7 & \begin{tabularx}{0.7\textwidth}{X} 現在他們知道，你所賜給我的一切都是從你那裡來的； \end{tabularx} \\ \\ \relax
17:8 & \begin{tabularx}{0.7\textwidth}{X} 因為你所賜給我的話，我已經賜給他們，他們也領受了，又確實知道，我是從你出來的，並且信你差了我來。 \end{tabularx} \\ \\ \relax
17:9 & \begin{tabularx}{0.7\textwidth}{X} 我為他們祈求，不為世人祈求，卻為你所賜給我的人祈求，因他們本是你的。 \end{tabularx} \\ \\ \relax
17:10 & \begin{tabularx}{0.7\textwidth}{X} 凡是我的都是你的，你的也是我的，並且我因他們得了榮耀。 \end{tabularx} \\ \\ \relax
17:11 & \begin{tabularx}{0.7\textwidth}{X} 我到你那裡去；我不再留在世上，他們卻在世上。聖父啊，求你因你的名，就是你所賜給我的名，保守他們，使他們像我們一樣合而為一。 \end{tabularx} \\ \\ \relax
17:12 & \begin{tabularx}{0.7\textwidth}{X} 我與他們同在的時候，我奉你的名，就是你所賜給我的名，保守了他們，我也護衛了他們；其中除了那滅亡之子，沒有一個滅亡的，好使經上的話得以應驗。 \end{tabularx} \\ \\ \relax
17:13 & \begin{tabularx}{0.7\textwidth}{X} 現在我到你那裡去，我在世上說這些話，是要他們心裡充滿了我的喜樂。 \end{tabularx} \\ \\ \relax
17:14 & \begin{tabularx}{0.7\textwidth}{X} 我已把你的道賜給他們；世界恨他們，因為他們不屬世界，正如我不屬世界一樣。 \end{tabularx} \\ \\ \relax
17:15 & \begin{tabularx}{0.7\textwidth}{X} 我不求你把他們從世上接走，只求你保全他們，使他們脫離那惡者。 \end{tabularx} \\ \\ \relax
17:16 & \begin{tabularx}{0.7\textwidth}{X} 他們不屬世界，正如我不屬世界一樣。 \end{tabularx} \\ \\ \relax
17:17 & \begin{tabularx}{0.7\textwidth}{X} 求你用真理使他們成聖；你的道就是真理。 \end{tabularx} \\ \\ \relax
17:18 & \begin{tabularx}{0.7\textwidth}{X} 你怎樣差我到世上，我也照樣差他們到世上。 \end{tabularx} \\ \\ \relax
17:19 & \begin{tabularx}{0.7\textwidth}{X} 我為他們的緣故使自己分別為聖，為要使他們也因真理成聖。 \end{tabularx} \\ \\
[1ex]
\hline
\hline
\end{longtable}
講題:栽在溪水旁   Planted by the River

經文:詩篇1:1-6；約翰福音17:6-19

講員:李均熊牧師

場合:香港中文大學崇基學院禮拜堂主日崇拜

日期:2024年5月12日

起:希伯來詩的結構欣賞

今天的啟應詩篇,是大家可能十分熟悉的詩篇第一篇.很多學者都認為這詩是在詩篇書卷形成時最後添加的詩篇,好作為整卷詩篇的引言.所以大家或會留意到詩篇第一篇既不是讚美詩,也不是哀求詩,乃是一首叫人喜愛律法的智慧詩,充滿著教導的氣氛.就讓我們今天一同看看這首詩教導我們甚麼道理.

用詩篇作為宣講經文其實十分困難,因為我們單靠閱讀中文翻譯,可能不能領會原本希伯來詩的妙韻.所以通常我都會嘗試重新翻譯經文,並且加入一些註腳好使會眾留意一些只靠中文譯本可能未必看到的細節.詩篇第一篇作為一首智慧詩,它的內容就大玩文字遊戲,從希伯來原文的音,形,義都有不同程度的重複和呼應.今天我也不會逐點細數,大家就試細閱登在週刊的譯文和註腳,看會否有所領會.歡迎大家有機會再交流.

雖然如此,只要大家掌握兩個要點,或許更能對詩篇經文有所領會.這兩個要點,就是要掌握希伯來詩篇每篇的結構,和所運用的意像.

希伯來詩有些如我們中國人的古詩,是十分講求格律的,所以每篇詩都有工整的結構.透過仔細閱讀經文字眼的重複和句子的編排,我們就可以大概掌握一首希伯來詩的結構,從而可以透過結構所聚焦的焦點,來了解詩的主題和修辭目的.

例如詩篇第一篇,首兩節和第5節就有多番的重複.比較突出的是前置詞「在」(in)的重複,在詩一共8次出現的「在」,首兩節便出現了5次,而第5節則出現了2次.其次詩第1節中出現的人物:「惡人」,「罪人」也均在第5節再出現,第1節的「道路」則在第6節重複(參週刊倒影結構中不同字體的標示).(第5節的「站立」乃指「起立」,不是第1節的「站」字).從這些重複中,再以符合古詩的格律要求,將九行詩均分,便可分成三行詩句一組,我們就可以將詩一分為工整的三個段落(參平行結構).由此,我們看到當中首,尾段落提出教導(或說智慧人的觀察),而中間則以一段以「樹」的明喻作為解說(illustration).其實我們也可以用一種倒影結構的方法去分析這詩(參週刊),而無論詩的外圍和核心部份,都展示了詩要表達的要點.

承:意像的妙用

另一個幫助我們了解詩篇的方法,就是細想詩人運用的意像(imagery)和比喻(metaphor).詩篇語言通常用上非字面的手法,透過很多濃縮的以文字生成的圖像,和民族中家傳戶曉的典故,來表達豐富的內容.拆解這些意像比喻,現成用語和民族典故,就可以大致了解詩篇要引導我們進入的境界為何.

就以今日的詩篇第一篇來看,這首智慧詩明顯用上了一些植物的比喻來作意境的傳遞.詩第3節就用上一棵樹栽種在溪水旁的圖畫,來帶出詩的要點.當一個人喜愛上主的訓誨,晝夜思想的時候,他的福份就好像一棵樹栽在溪水旁,按時候結出果子,葉子也不會枯黃.這個比喻最後卻以一個原全相反的圖畫:糠秕被風吹散,來對比惡人的結局.

大家可以在腦內畫出兩幅由詩詞拼湊出來的圖畫,去了解喜愛上主訓誨的義人,和反過來拒絕上主教誨的惡人的不同結局.第3節原文用上12個字去描繪的那棵樹,表達出有福的那個人豐富的生命,然而只用4個字就說完的惡人,則完全內容欠奉,甚至糠秕不能作為任何動詞的主語(subject),只能作為隨風飄散的物件.這個意像,也暗合第1 節教導我們,不要跟從惡人和罪人而動作多多,又行又站又坐,因為這只會帶來散亂的人生.不過後來,詩人在第5節下半節卻省去動詞,這種由第一段落到第三段落動詞的多寡運用,成為詩一結構幫助我們聚焦詩中要點的一大特色.

惡人說夠了,也不值得多說,就讓我們細看這個有福人如何「像一棵樹栽在溪水旁」.這個比喻也用在耶利米書17:7-8:

倚靠耶和華,以耶和華為可靠的, 那人有福了! 他必像樹栽於水旁, 在河邊扎根, 炎熱來到,並不懼怕, 葉子仍必青翠, 在乾旱之年毫無掛慮, 而且結果不止.」

但若仔細對比兩段經文,就更見詩一用詞的別具用心.大家不知有沒有留意我的翻譯裡並沒有用上溪水,而是用上「流水」.因為原文所用的字眼,是指一個人工做成的水道(channel),而非自然因雨水和地形而產生的水溪,那可是耶17:8所用的「河邊」一字.加上「栽」字所指刻意的「種植」,這棵樹並不是自然而生,而是受到呵護照顧的果樹,所以它才會「按時候結出果子」.

「按時候結出果子」原文應該翻成為「在它的時候給出果子」,一方面這裡其實用上了前置詞的「在」字,呼應詩一首尾段落的重複用語；另一方面,比較耶17:8的「結果」是用「製造」to make／「結出」to yield,而這裡用上「給予」to give一字,令人想起創世記的典故:上主創造天地時,第三天造樹木結果子(創1:11-12),好可以到了創造人類時宣告將果子「給」人作食物(創1:29)!所以當你因喜愛上主的訓誨而日以繼夜地,不間斷地,連續地,「在」祂的訓誨中,你就會好像被悉心安置(「栽」)在有水道灌溉(「流水旁」)的果樹(這又會否暗指生在伊甸園的河水旁的生命樹的典故(創2:8-15)?),因而會「在」適當的時候,為別人供應你果實,就如上主賜人食物般祝福別人.當然這個有福的人自己不會枯乾(葉子不枯乾,參耶17:8,那裡集中講倚靠上主的人,自己能渡過難關),但他生命產生的成果,卻是在社群中給予他人的恩賜,而非單為自己所用.「凡他所做的盡都順利」——這裡其實也有雙重意思,「他」可以是有福的人,「所做的盡都順利」；但「它」也可以是指那棵樹,那此句就可譯做「它所有的出產(果和葉)都繁茂」all it produces will flourish.

轉:在那裡的省思

當你細想那有福的人如何像一棵樹栽於水旁的這個意像,你就會更明白詩一要教導我們的是甚麼了.一棵樹之所以在社群中對人對己有益繁茂,皆因它能栽於流水旁,它能長期不間段的受到照料,灌溉,這就回應了詩一結構首尾呼應的兩個段落的重點:有福的人,他「在」那裡呢?他經常接觸,流年,相伴的是甚麼人和事呢?

他不會在「當神無到」的「惡人」中,他們以自己的計謀行事,常常圖謀不軌,卻將神排除於生命之外.

他也不會在離開神不行神心意的「罪人」中,他們或許浸醉在罪之樂,也可能「應做的不做,不應做的反去做」,並且樂此不疲,毫無悔意.

他們更不會在褻慢人中間,就是那些用言語製造紛爭(箴22:10「趕出褻慢人,爭端就消除, 紛爭和羞辱也必止息.」),或心驕氣傲而看不起別人的人(箴21:24 「心驕氣傲的人,名叫褻慢,他行事狂妄,都出於驕傲.」).用今日的用語,就是那些「有毒的人」toxic people,諸事八卦,流言蜚語,抱怨不停,情緒勒索；又或為恐天下不亂,散播懷疑,動不動製造紛爭的人 ；又或那些看不起人,常用言語貶低別人,無論這是因為他想抬高自己,還是想令人難受,他就是不斷的對別人不留餘地地「窒,串,潤,踩」(sorry, 講東話較傳神).

這些惡人最後都「不在」了.都像糠秕,隨風而散.對比會結果的樹,它們不單毫無用處,也毫無質量,只有隨風飄動,隨水漂流,永遠站不住腳,在上主審判時更抬不起頭,不能起立.在義人的聚會中,你應該找不到這些惡人罪人和toxic people 的.

那麼有福的人應該「在」那裡呢?還用說,當然是「在」主的訓誨中!那甚麼是「在」主的訓誨中呢?就是晝夜思想,任何時候,不間斷,都在誦讀(思想的本意,有人認為可解作「反芻」的動作)上主的訓誨,甚至愛上它,以反思上主的訓誨為樂.不過大家要留意,古代並沒有私下閱讀書卷的可能,所有的誦讀都是在社群中發生的,誦讀上主的訓誨更多時候是群體活動,而不是私有化上主的說話.

留意我不用「律法」去翻譯Torah妥拉一字,因為上主的「妥拉」,與我們今日所理解的「律法」相去甚遠.舊約所謂的妥拉,可以是指摩西在西奈山所頒佈神的要求律例,也可以是指摩西五經以色列民族的基礎歷史,更可以指整本希伯來聖經,但在詩篇更多時是描述上主的叮嚀,吩咐,要求,甚至詩篇五卷也視為神的妥拉／訓誨；故舊約學者今天多建議把妥拉譯作訓誨instruction,或乾脆直接音譯作妥拉torah.到了新約,耶穌基督更是成全這些上主要求妥拉／訓誨的那位,從他身上和他給我們的道(今日福音經課),我們就更能明白和實踐上主的妥拉／訓誨.「在」上主的妥拉,與其說是默誦聖經的經文,不如說是回想父親的叮嚀和吩咐,或反思基督如何成全上主的要求,以致我們可以學效跟從.

已故的楊錫鏘牧師曾經指出,詩篇第一篇作為詩篇的引言,也是我們讀詩篇的心法,那就是要像一棵樹栽在溪水旁,如樹根沈浸在水中一般,自己也要沈浸在詩篇的意境中,不斷重複誦讀,也不斷思考詩中的意思,看一次是不夠的,要看上過千百次你才會開始領會詩篇的意境,讓詩篇將你的信仰帶入新境界.

合:義人的抉擇

初看我們會以為詩篇第一篇是要我們看惡人和義人的兩個結局,但其實兩者根本無可比擬.而詩一作為一首智慧詩,其實它教導我們的是:我們若想成為義人或有福的人,其實是要作出智慧的選擇(wise choice),就是「在」那裡的抉擇:你要處身在何方?你要與甚麼人同行?你生活以甚麼為伴?當你知道惡人的結局,又領會喜愛上主訓誨的人如何好像一棵樹置在流水旁,自己自然不枯乾,更能在適當時候能成為別人的祝福幫助,你會怎樣選擇你的生活,和與你共渡生活的人呢?

詩篇第一篇(新譯:平衡結構)

1 真有福啊!這一個人!他------(1)~~~~~~~~不行(2)在(3)惡人的汁落中(4)；

在罪人的道路中(5),不站；~~~~~~~~並在褻慢人的座位中,不坐

2 相反地在上主的訓誨中是他至愛(6);在祂的訓誨中(7)他默誦(8)日以繼夜

3 就絛一棵樹(9)~~~~~~~~栽於流水(10)旁

在它的時候結(11)果子,~~~~~~~~葉子不枯乾(12),

凡他(它)所做的盡都順利.(13)

4 不是這樣的是惡人:~~~~~~~~相反地他像糠秕~~~~~~~~被風吹散(14).

5 因此惡人必站立不住在審判中；~~~~~~~~罪人------在義人的會中.

6 因為上主認識(15)義人的道路；~~~~~~~~但是惡人的道路是滅亡•

--------------------------------------------

1: 原文乃押強:asg-rei ha-ish ae-ser.

2: 原文動詞的格式強調動作的不連貫.

3: 全詩共出現前置詞達8次之多.在這裡均譯作「在…」.

4:「在惡人的計謀中」,與 v5「在義人的會中」產生呼應.留意「在...計謀中」與「在...會中」的形似.

5: 另一個與後面呼應的字詞(v6 義人的道路,惡人的道路).

6: 留意由上文三個不連貫卻遞進的動作,轉變成為此節一個名詞(原文沒有「是」),以及不間斷的動作(參註8),去作為如何「在…訓誨中」的狀態.

7: 留意「訓誨」(原文:Torah)的重覆,傳統譯「律法」.

8: 原文有「默想」和「低聲誦念」二重意思(或透過低聲誦念去默想),並且動詞的格式強調動作的不間斷.

9: 「像一棵樹」與下文「像糠秕」的原文乃十分相似的.

10: 留意「栽」與「流水」(channel of water)皆為人工而非自然的安排.

11: 原文為「給予」.

12: 「枯乾」的字根與「愚蠢」的字根是用上一樣的宇母的.

13: 這句其實有雙重意思,「他」可以是有福的人,「所做的盡都順利」；但「它」也可以是指那棵樹,那此句就可譯做「它所有的出產(果葉)都繁茂」all it produces will flourish.

14: 在以十二個字去形容樹後,原文只有四個字去形容糠秕.而且鐮秕並不能成為任何動詞的主語,而只能作為吹散的對象.

15: 原文有「深入的關係而了解」的意思.

--------------------------------------------

<引子>

真有福啊!這一個人!

他不行,在惡人的汁謀

<「在」哪裡?>

~~~~ 在罪人的道路,不站

~~~~ 並在褻慢人的座位,不坐

相反地在上主的訓海,是他至愛

在祂的訓誨,他默誦------日以繼夜

<比喻一>

就像一棵樹,栽於溪水旁,

在它的時候結果子.

~~~~ 葉子不枯乾,

<小結>

~~~~~~~~ 凡他(它)所做的盡都順利.

<小引>

~~~~~~~~ 不是這樣的是惡人:

<比喻二>

~~~~ 相反地像糠秕

它們被風吹散.

<「在」那裡!>

因此惡人必站立不住,在審判時；

罪人~~~~ 在義人的會中.

<結語>

因為上主深深認識義人的道路；

但惡人的道路是滅亡.



\newpage

\section{誰是真先知}
\label{sec:698}
\textbf{鄧瑞強博士 - 誰是真先知}
\newline
\newline
\hyperref[sec:697]{< < < PREV SERMON < < <}
~
\hyperref[sec:toc]{[返目錄]}
~
\hyperref[sec:699]{> > > NEXT SERMON > > >}


耶利米書 28:1-17
\newline
\begin{longtable}{cl}
\hline
\hline
章節 & 經文 (和合本修訂版)\\
\hline
28:1 & \begin{tabularx}{0.7\textwidth}{X} 當年，就是猶大王西底家登基第四年五月，押朔的兒子基遍人哈拿尼雅先知，在耶和華的殿中當著祭司和眾百姓的面對我說： \end{tabularx} \\ \\ \relax
28:2 & \begin{tabularx}{0.7\textwidth}{X} 「萬軍之耶和華－以色列的神如此說：我已經折斷巴比倫王的軛。 \end{tabularx} \\ \\ \relax
28:3 & \begin{tabularx}{0.7\textwidth}{X} 二年之內，我要將巴比倫王尼布甲尼撒從這地擄掠到巴比倫的器皿，就是耶和華殿中的一切器皿，都帶回此地。 \end{tabularx} \\ \\ \relax
28:4 & \begin{tabularx}{0.7\textwidth}{X} 我又要將約雅敬的兒子猶大王耶哥尼雅和被擄到巴比倫所有的猶大人帶回此地，因為我要折斷巴比倫王的軛。這是耶和華說的。」 \end{tabularx} \\ \\ \relax
28:5 & \begin{tabularx}{0.7\textwidth}{X} 耶利米先知當著祭司和站在耶和華殿裡眾百姓的面，對哈拿尼雅先知說： \end{tabularx} \\ \\ \relax
28:6 & \begin{tabularx}{0.7\textwidth}{X} 「阿們！願耶和華如此行，願耶和華實現你所預言的話，將耶和華殿中的器皿和所有被擄去的人從巴比倫帶回此地。 \end{tabularx} \\ \\ \relax
28:7 & \begin{tabularx}{0.7\textwidth}{X} 然而我在你和眾百姓耳中所要說的話，你應當聽。 \end{tabularx} \\ \\ \relax
28:8 & \begin{tabularx}{0.7\textwidth}{X} 從古以來，在你我以前的眾先知，向多國和大邦說預言，論到戰爭、災禍、瘟疫的事。 \end{tabularx} \\ \\ \relax
28:9 & \begin{tabularx}{0.7\textwidth}{X} 至於那預言平安的先知，到先知的話應驗的時候，人就知道他真是耶和華所差來的。」 \end{tabularx} \\ \\ \relax
28:10 & \begin{tabularx}{0.7\textwidth}{X} 哈拿尼雅先知就取下耶利米先知頸項上的軛，把它折斷。 \end{tabularx} \\ \\ \relax
28:11 & \begin{tabularx}{0.7\textwidth}{X} 哈拿尼雅又當著眾百姓的面說：「耶和華如此說：二年之內我必照樣從列國的頸項上折斷巴比倫王尼布甲尼撒的軛。」耶利米先知就離開了。 \end{tabularx} \\ \\ \relax
28:12 & \begin{tabularx}{0.7\textwidth}{X} 哈拿尼雅先知折斷耶利米先知頸項上的軛以後，耶和華的話臨到耶利米，說： \end{tabularx} \\ \\ \relax
28:13 & \begin{tabularx}{0.7\textwidth}{X} 「你去告訴哈拿尼雅說，耶和華如此說：你折斷木軛，卻換來鐵軛！ \end{tabularx} \\ \\ \relax
28:14 & \begin{tabularx}{0.7\textwidth}{X} 萬軍之耶和華－以色列的神如此說：我已將鐵軛加在這些國的頸項上，使他們服事巴比倫王尼布甲尼撒。他們總要服事他，我也把野地的走獸給了他。」 \end{tabularx} \\ \\ \relax
28:15 & \begin{tabularx}{0.7\textwidth}{X} 於是耶利米先知對哈拿尼雅先知說：「哈拿尼雅啊，你應當聽！耶和華並沒有差遣你，你竟使這百姓倚靠謊言。 \end{tabularx} \\ \\ \relax
28:16 & \begin{tabularx}{0.7\textwidth}{X} 所以耶和華如此說：看哪，我要把你從地面上除掉，你今年必死，因為你向耶和華說了叛逆的話。」 \end{tabularx} \\ \\ \relax
28:17 & \begin{tabularx}{0.7\textwidth}{X} 這樣，哈拿尼雅先知當年七月間就死了。 \end{tabularx} \\ \\
[1ex]
\hline
\hline
\end{longtable}
講題:誰是真先知?   Who is the Real Prophet?

經文:耶利米書28:1-17

講員:鄧瑞強博士

場合:香港中文大學崇基學院禮拜堂主日崇拜

日期:2024年5月19日

各位弟兄姊妹,早安.願神的恩典常與你們同在.

今日講耶利米書.主前627年,神選召耶利米為先知.那一年,亞述帝國最後一位偉大王者「亞述巴尼拔」(Ashurbanipal)逝世,巴比倫乘勢崛起,形成一種三雄鼎立的局面:傳統的勢力埃及,新興的勢力巴比倫,和衰敗中的亞述.主前605年,埃及決心爭做霸主,便聯同亞述的殘餘力量,在迦基米施(Carchemish)與正在冒起的巴比倫一決雌雄,巴比倫大勝.巴比倫尼布甲尼撒王稱霸武林的時代,正式來臨.

猶大國與埃及接壤,因地緣因素,它無法不靠近埃及.但一靠近埃及,便惹怒巴比倫.主前597年,巴比倫連廢兩位猶大王,將猶大王擄去當俘虜,同時擄走一批精英,包括但以理,以西結等人,也掠走一批聖殿器皿,這就等於將戰敗國的神明擄走.巴比倫改立西底家為傀儡王,這西底家是末代君王,後來因反叛巴比倫,巴比倫大軍圍城.在位11年,耶路撒冷陷落,聖殿被毀,猶大國破家亡,人民被擄他方.耶利米書用「空虛混沌」(耶4:23)來形容這劫後景象,「空虛混沌」這用字曾經出現在創世記1:2裡:「地是空虛混沌,淵面黑暗」,那是形容神創造前的空洞荒涼.亡國被擄,對猶大人來說,是整個世界回到太古前的荒蕪.先知耶利米就是生活於這時代,他親眼目睹這一切.

今日的講道經文:耶利米書28:1-17

28:1　當年,就是猶大王西底家登基第四年五月,基遍人押朔的兒子,先知哈拿尼雅,在耶和華的殿中當著祭司和眾民對我說:

28:2　「萬軍之耶和華─以色列的神如此說:我已經折斷巴比倫王的軛.

28:3　二年之內,我要將巴比倫王尼布甲尼撒從這地掠到巴比倫的器皿,就是耶和華殿中的一切器皿都帶回此地.

28:4　我又要將猶大王約雅敬的兒子耶哥尼雅〔即:約雅斤〕和被擄到巴比倫去的一切猶大人帶回此地,因為我要折斷巴比倫王的軛.這是耶和華說的.」

28:5　先知耶利米當著祭司和站在耶和華殿裡的眾民對先知哈拿尼雅說:

28:6　「阿們!願耶和華如此行,願耶和華成就你所預言的話,將耶和華殿中的器皿和一切被擄去的人從巴比倫帶回此地.

28:7　然而我向你和眾民耳中所要說的話,你應當聽.

28:8　從古以來,在你我以前的先知,向多國和大邦說預言,論到爭戰,災禍,瘟疫的事.

28:9　先知預言的平安,到話語成就的時候,人便知道他真是耶和華所差來的.」

28:10　於是,先知哈拿尼雅將先知耶利米頸項上的軛取下來,折斷了.〔按:之前耶利米作了一先知式的象徵行動,將軛套在頸上,象徵要服從巴比倫；如今,哈拿尼雅作一相反的象徵行動,將軛取走折斷,象徵要脫離巴比倫的轄制了.〕

28:11　哈拿尼雅又當著眾民說:「耶和華如此說:二年之內我必照樣從列國人的頸項上折斷巴比倫王尼布甲尼撒的軛.」於是先知耶利米就走了.

28:12　先知哈拿尼雅把先知耶利米頸項上的軛折斷以後,耶和華的話臨到耶利米說:

28:13　「你去告訴哈拿尼雅說,耶和華如此說:你折斷木軛,卻換了鐵軛!

28:14　因為萬軍之耶和華─以色列的神如此說:我已將鐵軛加在這些國的頸項上,使他們服事巴比倫王尼布甲尼撒,他們總要服事他；我也把田野的走獸給了他.」

28:15　於是先知耶利米對先知哈拿尼雅說:「哈拿尼雅啊,你應當聽!耶和華並沒有差遣你,你竟使這百姓倚靠謊言.

28:16　所以耶和華如此說:看哪,我要叫你去世,你今年必死,因為你向耶和華說了叛逆的話.」

28:17　這樣,先知哈拿尼雅當年七月間就死了.

我們仔細看看這個故事.

28:1　當年,就是猶大王西底家登基第四年五月......

西底家為什麼會坐在王位,就因為4年前,巴比倫擄走前任君王約雅斤,而立他作傀儡王.西底家王第4年,離開亡國,還有7年.那當然不會是風平浪靜的年代.巴比倫的壓迫,無日無之.作為傀儡國,要納的稅,一定重得難以負荷.前任的君王被擄走,真是顏面無存.但最令人耿耿於懷的,是聖殿的器皿被奪走.聖殿是大能的神的居所,聖殿的器皿是神聖的,人藉這些器皿去進行信仰活動,與神連繫.對猶大人來說,聖殿的一切代表著神的臨在,神的祝福,神的眷顧.在聖殿那裡,猶大人尋索生命的意義,寄托人生的希望.聖殿的器皿被掠去,他們感到一種無名的空洞.

28:1　當年...... 基遍人押朔的兒子,先知哈拿尼雅,在耶和華的殿中當著祭司和眾民對我說:

28:2　「萬軍之耶和華─以色列的神如此說:我已經折斷巴比倫王的軛.

28:3　二年之內,我要將巴比倫王尼布甲尼撒從這地掠到巴比倫的器皿,就是耶和華殿中的一切器皿都帶回此地.

28:4　我又要將猶大王約雅敬的兒子耶哥尼雅〔即:約雅斤〕和被擄到巴比倫去的一切猶大人帶回此地,因為我要折斷巴比倫王的軛.這是耶和華說的.」

在混亂的歲月,是有很多先知發言的.有人說應這樣行,有人說應那樣行.個個都是奉耶和華的名發言的,個個都有深厚的神學知識,個個都帶著牧養的情懷.這裡就有一個,他叫哈拿尼雅.他在聖殿,奉神的名,當著祭司和民眾,公開地向耶利米發言.他的發言,開頭和結尾都是:神要折斷巴比倫王加給猶大的軛.嘩,這番話多麼中聽.這是每個猶大人的心底話.誰願意活在巴比倫的強權下呢?

哈拿尼雅說,兩年之內,神要將巴比倫王尼布甲尼撒從聖殿掠走的一切器皿帶回耶路撒冷.嘩,這先知真有信心,兩年之內,這事必成.這種話不是隨口能說的,說兩年內事必成,必須對事情有相當把握.或許,他是根據一種堅實的聖殿神學說的.很久以前,猶大人和非利士人打仗,被非利士人搶走約櫃,後來,非利士人還不是乖乖的將約櫃送回.誰能奪去聖殿的東西而不招禍的呢?聖殿的東西,是耶和華的東西,耶和華一定會親手奪回來,以顯示祂的威榮.或許,哈拿尼雅如此相信.

哈拿尼雅又說,神也會將被擄去的約雅斤王和其他猶大精英帶回來.哈拿尼雅如此相信,或許,是根據一種大衛王室神學而說的.神與大衛立下永遠的約,神祝福大衛的王系血脈,故此,大衛王位的正統繼承人,不會流落異鄉,必會回歸故土的.再者,看來哈拿尼雅深信出埃及的神學.按這神學,神能帶領被擄的人歸回.哈拿尼雅講法的每一點,都有其信仰根據.可以想像,哈拿尼雅會挑戰耶利米:「耶利米啊,耶利米,你為何這般小信,為何對神這般無信心,神要消滅巴比倫,要取回聖殿的器皿,要帶回被擄的人,有多難呢?」哈拿尼雅講的,完全是猶大人的渴望.他講的東西,完全建基於猶大人的信仰.他表達出來的信心,完全是眾人的模範.面對個人的危難,或者面對社會的困局,不少人像哈拿尼雅,憑信心,按其神學知識,說:「沒事的,沒事的,神會救助我們」.

我們看看耶利米的反應.

28:5　先知耶利米當著祭司和站在耶和華殿裡的眾民對先知哈拿尼雅說:

28:6　「阿們!願耶和華如此行,願耶和華成就你所預言的話,將耶和華殿中的器皿和一切被擄去的人從巴比倫帶回此地.

28:7　然而我向你和眾民耳中所要說的話,你應當聽.

28:8　從古以來,在你我以前的先知,向多國和大邦說預言,論到爭戰,災禍,瘟疫的事.

28:9　先知預言的平安,到話語成就的時候,人便知道他真是耶和華所差來的.」

聽到哈拿尼雅的話,耶利米很冷靜,他說:「阿們!願耶和華如此行.」哈拿尼雅講的,都是好事,哪個猶大人會反對這信息呢?當四面楚歌時,哈拿尼雅的話,的確能振奮人心.特別是要團結人心,抵抗外敵時,哈拿尼雅的話很有鼓舞作用.耶利米希望他的話成真.然後,耶利米對哈拿尼雅說,以前的先知講的,多是審判的信息.這些信息能叫人及早悔改,及時回頭.現在,你講的,是平安的信息,是真是假,就要「放長雙眼」了.耶利米只是希望哈拿尼雅的話成真,但是,這是否真的是來自神的安慰之言,耶利米有所保留,走著瞧吧!

之後,耶和華的話臨到耶利米.我們聽聽耶和華說什麼?

28:13　「你去告訴哈拿尼雅說,耶和華如此說:你折斷木軛,卻換了鐵軛!

28:14　因為萬軍之耶和華─以色列的神如此說:我已將鐵軛加在這些國的頸項上,使他們服事巴比倫王尼布甲尼撒,他們總要服事他；我也把田野的走獸給了他.」

耶和華講什麼?一語道盡,耶和華叫猶大人向巴比倫投降.

o下,這是什麼話?猶大民族的主保神,大能的耶和華,叫自己子民向他人低頭.自古以來,耶和華只有幫助猶大人脫離為奴之境,從未有過叫猶大人甘心去作奴隸的.耶和華講的話,完全脫離哈拿尼雅的神學框架,也脫離我們習以為常的信仰框架.耶和華如此說:「哈拿尼雅折斷了耶利米的木軛,我卻換上鐵軛!」哈拿尼雅這麼有信心,相信神的解救.一般人,包括我們,也會以為:神必定會應允他的渴求.但事實是:神加重猶大人的苦難.想想:你祈求神解救,神不單不解救,反加重艱難,這是否超出你的信仰框架?我們習慣救贖神學,神像黃大仙,一定會救你.我們不習慣創造神學,神是萬有的主,祂行事不按你的渴望,而只按祂自己的美意.想想:天父沒有救耶穌脫離十架的苦難,而是按祂自己的美意,讓他受苦,好成就超出一己之外的,關乎全人類的,涉及悠長歷史的天國工程.我們習慣只看到自己,但神處理的,是歷史的大局.

耶和華說:猶大人要好好服事巴比倫王尼布甲尼撒.28:14最後的短句很有趣:「我也把田野的走獸給了他」.這牽涉神創造的主權,祂把萬物交給誰,就交給誰.在這裡,神講得很清楚,祂把一切生命交給了尼布甲尼撒,由他主宰生死.順從他的,會得到安寧.但後來西底家也是聽了哈拿尼雅那一派先知的說話,抵抗巴比倫,至終導致亡國.

耶和華的話,與猶大人慣常的信仰想法,很不一樣.按猶大人慣常的上帝觀,君王觀及民族觀,耶和華是猶大人的上帝,祂設立猶大君王來管治,猶大人應該只服事猶大君王,這等於服從耶和華.服事巴比倫王,即等於服從巴比倫的神明,這是猶大人做不到的.但是,現在,耶和華叫他們服事異教的巴比倫王,這超出了他們慣常的信仰框架.想想:耶和華若叫你去服事反耶和華的君王,你都會覺得匪夷所思.

經文來到這裡,就涉及耶利米書研究的一個重大課題:如何辨別真先知和假先知?答案是:沒有方法,沒有簡易的方法,沒有一成不變的方法.但參照哈拿尼雅和耶利米的這個對峙個案,我們得出一些想法.哈拿尼雅的信仰,是一種框架式的信仰.按神過去的作為,歸納出某些信念框架,然後,按這些框架作出推論.很有經文根據的,很「屬靈」的,很有信心的.這種信仰模式很「穩陣」,但就將神困在一個框框裡.面對全新形勢,這個「框框裡的神」能做的,是將全新形勢塞回舊框架裡,事情要回復到原先想好的框框裡.這種信仰,看起來很「屬靈」,實質是守舊,是將上帝和信仰封閉在一安全密室裡.這樣的一位神,無法作出什麼新事來,也開不出從未有過的將來.想想這個思想遊戲:設想——我們有幸,在天上,與未成肉身的耶穌對話.耶穌說,我要成肉身,住在人間,與人同在.若我們只有一陳舊的上帝觀,我們會勸說他,千萬不要這樣做.神人有別,無限的神要成為人,會貶損神性.這對人也無好處,我們不需要有人性軟弱的神,我們需要無窮大能的神.耶穌又說,他成肉身後,會在十架上死去.我們又會反對,這般羞辱的死,羞辱了神的榮耀,人不會相信這種羞辱的神的.......若以我們的框架去信仰,神就什麼也做不到.當然,也就無法在羞辱的死亡中,開出人類過去從未想像過的復活生命.哈拿尼雅的信仰,將神限死在一固定的框架裡.

耶利米卻走出了舊的框框,面對活著的上帝,也面對歷史的現實.因人的罪,耶和華可以放棄與祂有永恆之約的大衛王朝.因人的罪,耶和華可以驅逐祂的子民離開那應許給他們的地.神可以與猶大人為敵,而使用猶大人的敵人巴比倫為僕人.神可以放棄祂在耶路撒冷居所的聖殿,而將寶座遷去巴比倫.神可以使用巴比倫人,去推進歷史,而不一定要使用猶大人.這種種想法,都是突破性的.要等到國破家亡後,猶大人才能大徹大悟——原來,神不單是猶大人的神,更是萬國萬族的神；被擄去巴比倫,這不是神的離棄,因為神也可以在巴比倫與他們同在；再者,猶大人可以不再以本位為中心去思考神的救贖,神救贖的,是萬國萬族.就是這種信仰,開啟了猶大人的心胸和眼界,支撐著他們越過全新的歷史變局,越過了前所未有的苦難.

28:15　於是先知耶利米對先知哈拿尼雅說:「哈拿尼雅啊,你應當聽!耶和華並沒有差遣你,你竟使這百姓倚靠謊言.

28:16　所以耶和華如此說:看哪,我要叫你去世,你今年必死,因為你向耶和華說了叛逆的話.」

28:17　這樣,先知哈拿尼雅當年七月間就死了.

經文是有點黑色幽默的.哈拿尼雅啊,你講到猶大國會平安大吉,不會亡國,看,連你自己也不能自保.耶利米原先說,放長雙眼看看結局,怎知,不到兩個月,哈拿尼雅便死了.

我們習慣在框架下信仰,這給我們一種安全感.病了,我們求神醫治,這是人之常情.問題不在這「常情」,問題在我們如何信.或許,我們持守著一框架——神是憐憫我的,必定會醫治我；若祂不醫治我,則多數是我犯了罪.若我認罪,神會赦免我,便會醫治我了.若最後神都不醫治我,則祂不愛我,或者,祂根本無能,如此,祂根本不可信.我們常常在這些框架裡走來走去,像哈拿尼雅.耶利米的上帝會說:我是健康人的主,也是病人的主.我使人健康,是希望他能為天國做點事情.我不醫治一個病人,也是希望他能為天國做點事情.在光明中,我眷顧你；在黑暗中,我也眷顧你.詩篇139:12說:「黑暗和光明,在你看都是一樣.」唯有這種闊大的信仰心靈,能支撐我們,度過難關.

願神賜福你們.也願榮耀歸於聖父,聖子,聖靈.



\newpage

\section{當遵守律法與追求美善不可兼得的時候}
\label{sec:699}
\textbf{林豪恩先生 - 當遵守律法與追求美善不可兼得的時候}
\newline
\newline
\hyperref[sec:698]{< < < PREV SERMON < < <}
~
\hyperref[sec:toc]{[返目錄]}
~
\hyperref[sec:700]{> > > NEXT SERMON > > >}


講題:當遵守律法與追求美善不可兼得的時候 When You Cannot Be Lawful and Good at the Same Time

講員:林豪恩先生

場合:香港中文大學崇基學院禮拜堂主日崇拜

日期:2024年6月2日

引言

最近拜讀白耀燦師兄的大作《歷史戲劇作品集:詩比歷史更真實?》,啟發了我重思今次講道的形式.於是,我把原本正在草擬的論說式講章放棄,嘗試運用想像力,把論說轉化為對話.不過,一個人讀對白,一定有不足之處,或者會混淆人物,希望各位包容.剛才讀福音經課的時候,大家都發現經文以對話和行動為主,讓我們以這些對話為基礎,發掘事件中人物在事件中的思考及掙扎.就讓我們從2:24法利賽人對耶穌所說的話開始.

第一回合:耶穌和法利賽人關於門徒被指責違反律法

2:24法利賽人對耶穌說:看哪,他們在安息日為甚麼做不可做的事呢?(原文的意思是不合乎律法,和修:不合法的事；NIV: unlawful)

門徒搶先回應:律法專家,請問我們做了甚麼不合律法的事呢?

法利賽人:你們做過甚麼事,你們心知肚明,還要問我們嗎?你們經過麥地的時候,他們掐了麥穗來吃.

門徒:掐麥穗來吃又犯甚麼律法呢?

法利賽人:你們的老師沒有教你們嗎?在安息日工作是犯法的.難道你沒有讀過出埃及記20章記載的十戒嗎?關於安息日是這樣寫的:「當記念安息日,守為聖日.六日要勞碌做你一切的工,但第七日是向耶和華─你的神當守的安息日.這一日你和你的兒女,奴僕,婢女,牲畜,以及你城裡

寄居的客旅,都不可做任何的工.因為六日之內,耶和華造天,地,海和其中的萬物,第七日就安息了；所以耶和華賜福與安息日,定為聖日.」(8-11)

門徒:我們吃麥穗等於工作嗎?難道你們在安息日不吃東西的嗎?

法利賽人:不要狡辯!安息日不是不可以吃東西,只是不可以工作.連路加福音的作者也承認你們不但「摘了麥穗」,還「用手搓著吃」(6:1).又摘又搓,如此多工序,還不是違反了「安息日不可做任何的工」的律法嗎?

門徒:請問律法師們,如果不把麥穗用手搓去外殼才吃,難道要連殼一齊吃?

法利賽人:你們怎樣吃我們不理.但他們在安息日工作就是違反律法.

門徒:因為肚餓,我們才摘麥穗來呢.馬太福音12:1指出說明:「他的門徒餓了,就摘麥穗來吃.」肚餓要吃東西,這是人的基本需要啊!

法利賽人:你們肚餓還是不肚餓我們不理,但違反律法就是違反律法,這是無可爭辯的事實.

就在這個時候,耶穌發言了.

耶穌:律法專家們,我告訴你們一件事,請你們判斷一下,好嗎?

2:25耶穌對他們說:「經上記著大衛和跟從他的人缺乏飢餓之時所做的事,你們沒有念過嗎?

法利賽人:你指的是甚麼事?你想說甚麼?

2:26他當亞比亞他作大祭司的時候,怎麼進了上帝的殿,吃了陳設餅,又給跟從他的人吃.這餅除了祭司以外,人都不可吃(not lawful for any but the priests to eat).」

耶穌:請你們判斷一下大衛吃陳設餅這件事.

法利賽人:大衛是我們的民族英雄,是耶和華選立的君王,是以色列人的驕傲,你怎可以抹黑大衛?

耶穌:這件事記載在撒母耳記上21章,你們一定讀過,怎麼說我抹黑大衛呢?你們不會認為律法只是平民百姓要導守,王帝不用遵守吧!

法利賽人:胡說八道!律法是人人都要遵守的,無分貴賤.

門徒:老師,你的意思是不是說,大衛因為肚餓吃了陳設餅,就算是違反規定的,但沒有人指責他.即是說不能夠用「違反律法就是違反律法」一刀切去判斷所有事情,而是人們會考慮當時的處境,以及當事人的動機和目的.對嗎?

法利賽人:「但是,違反律法就是違返律法.」

這個時候,耶穌又出聲了:

2:27又對他們說:「安息日是為人設立的,人不是為安息日設立的.

2:28所以,人子也是安息日的主.」

門徒:「老師,你的意思是不是想他們在執行律法前想一想,到底著緊的是甚麼?著緊的是安息日律法,還是活生生,有血有肉的人呢?他們一直以為不犯安息日律法,就不會得罪上帝.但是,卻會因此扭曲了安息日的意義.你曾經教導我們,設立安息日的意義在於讓上主的創造,包括人,動物和大地都有安息的機會,與上主一齊欣賞創造的美意.從這個角度來看,安息日是讓人得到釋放來到上主跟前,讓人的需要得到滿足；安息日並不是要為人製造囚牢,把人囚禁起來.

門徒:我記起你曾經問我們,安息日的主角是誰呢?是安息日律法嗎?我們都以為是,但你告訴我們安息日的主角不是律法,而是上主本身.如果因為執行安息日律法而令人認為上主是冷酷無情的,吹毛求疵的,視人的需要於不顧的,如此,只會令人對上主敬而遠之,甚或心生厭惡,這又豈是上主所願的呢?」

各位弟兄姊妹,你們又認為如何呢?另一方面,你們猜想法利賽人會不會被說服呢?欲知後事如何,請聽下回分解.

第二回合:耶穌親自在安息日治好病人

3:1耶穌又進了會堂,在那裡有一個人枯乾了一隻手.

各位弟兄姊妹,讓我們進入場景,你聽不聽得到會堂的一角有兩個人在細聲交談.

群眾Ａ:這個人受傷之後都很少出現了,今天知道耶穌來了,所以他也來.相信他來是請求耶穌醫治他.

群眾Ｂ:但今天是安息日,如果耶穌醫好他,就犯安息日律法了.

群眾Ａ:在安息日醫病是犯律法的嗎?

群眾Ｂ:又不可以這樣說,在安息日醫治急病是容許的.但醫治的不是急病就違反了安息日不可以工作的律法了.這個人只是手受傷,應該死不了的.

群眾Ａ:你看看,法利賽人都來了!今天我們這個小會堂真熱鬧,好戲在後頭.你看,法利賽人盯著耶穌,相信他們正在看耶穌會不會醫好那個人,如果耶穌出手,他們就找到耶穌違反安息日律法的證據了.

群眾Ｂ:其實這個人都好慘,新約外典「多馬福音」說,這個人本來是泥水匠.他剛才暗暗地對耶穌說:『先生,請你幫助我,因為我的家庭需要我用這隻手來養活他們.』」他一家大小都等著他工作的收入來生活.如果耶穌回應他的需要,就應該醫好他；但醫好他,就犯安息日律法.如何是好呢?

群眾Ａ:喂,暫停討論,來看,耶穌有行動了!

3:3耶穌對那枯乾一隻手的人說:「起來,站在當中.」

3:4又問眾人說:「在安息日行善行惡,救命害命,哪樣是可以的呢?」他們都不作聲.

耶穌:各位鄉親,今日是安息日,律法禁止在安息日工作,連醫病也被指違反律法.請問各位,醫好這個人的手,減少他的痛苦,讓他可以工作養家,這樣做,是好事還是壞事呢?

耶穌:反過來說,如果可以減輕他的痛苦而不做,可以醫好他的手而不做,讓他和一家大小繼續貧病交迫,是好事還是壞事呢?

耶穌:那麼,你們認為在安息日行善合乎律法,還是行惡合乎律法呢?救命合乎律法,還是害命合乎律法呢?

耶穌:大家為甚麼不回答呢?律法師,你們精通律法,一定知道甚麼是善,甚麼是惡,然而,你們竟然對著一個貧病交迫的人無動於衷?甚至要用安息日律法阻止別人伸出援手?你們所信的上主是怎樣的上主呢?你們把那位有恩典有憐憫的上主收藏去了哪裡呢?

3:5耶穌怒目周圍看他們,憂愁他們的心剛硬,就對那人說:「伸出手來!」他把手一伸,手就復了原.

耶穌:這位兄弟,伸出手來!你的手已經康復,可以去工作養活你的家人了.

剛才全場都屏息以待,耶穌行動之後,大家才鬆一口氣.那兩個本來在交談的人忍不住又說話了:

群眾Ａ:嘩,剛才那一幕好刺激啊!耶穌那一番話,真是大快人心!法利賽人一直要求我們守律法,守不到就責罰,我以為敬拜耶和華就是等於守律法.

群眾Ｂ:你看,法利賽人走了,不知道他們有甚麼打算呢?

第三回:為堅持美善而付出代價

3:6法利賽人出去,同希律一黨的人商議怎樣可以除滅耶穌.

後來發生的事情,大家都知道了.鏡頭一轉,有一位年輕法利賽人也曾經不辭勞苦地追捕違反律法的人.後來,他竟然也成為了被指責違反律法,由捕快變成被追捕的對象.讓我們詢問這位朋友,這是甚麼緣故呢?

4:5我們原不是傳自己,乃是傳基督耶穌為主,並且自己因耶穌作你們的僕人.

4:6那吩咐光從黑暗裡照出來的上帝,已經照在我們心裡,叫我們得知上帝榮耀的光顯在耶穌基督的面上.

原來你明白了律法不是上主,耶穌基督才是上主.但是,你如此行,會被堅持律法的人逼害呢!你如何自處呢?

4:8我們四面受敵,卻不被困住；心裡作難,卻不致失望；

4:9遭逼迫,卻不被丟棄；打倒了,卻不致死亡.

4:10身上常帶著耶穌的死,使耶穌的生也顯明在我們身上.

4:11因為我們這活著的人是常為耶穌被交於死地,使耶穌的生在我們這必死的身上顯明出來.

佩服!請問你何來有這種能耐的呢?

4:7我們有這寶貝放在瓦器裡,要顯明這莫大的能力是出於上帝,不是出於我們.

保羅先生,謝謝你留下書信給我們,讓我們看見你的榜樣.但願你的見證成為我們的鼓勵.



\newpage

\section{深淵裡的拯救}
\label{sec:700}
\textbf{李均熊牧師 - 深淵裡的拯救}
\newline
\newline
\hyperref[sec:699]{< < < PREV SERMON < < <}
~
\hyperref[sec:toc]{[返目錄]}
~
\hyperref[sec:701]{> > > NEXT SERMON > > >}


詩篇 130:1-9
\newline
\begin{longtable}{cl}
\hline
\hline
章節 & 經文 (和合本修訂版)\\
\hline
130:1 & \begin{tabularx}{0.7\textwidth}{X} 耶和華啊，我從深處求告你！ \end{tabularx} \\ \\ \relax
130:2 & \begin{tabularx}{0.7\textwidth}{X} 主啊，求你聽我的聲音！求你側耳聽我懇求的聲音！ \end{tabularx} \\ \\ \relax
130:3 & \begin{tabularx}{0.7\textwidth}{X} 耶和華啊，你若究察罪孽，主啊，誰能站得住呢？ \end{tabularx} \\ \\ \relax
130:4 & \begin{tabularx}{0.7\textwidth}{X} 但在你有赦免之恩，要叫人敬畏你。 \end{tabularx} \\ \\ \relax
130:5 & \begin{tabularx}{0.7\textwidth}{X} 我等候耶和華，我的心等候；我也仰望他的話。 \end{tabularx} \\ \\ \relax
130:6 & \begin{tabularx}{0.7\textwidth}{X} 我的心等候主，勝於守夜的等候天亮，勝於守夜的等候天亮。 \end{tabularx} \\ \\ \relax
130:7 & \begin{tabularx}{0.7\textwidth}{X} 以色列啊，你當仰望耶和華，因耶和華有慈愛，有豐盛的救恩。 \end{tabularx} \\ \\ \relax
130:8 & \begin{tabularx}{0.7\textwidth}{X} 他必救贖以色列脫離一切的罪孽。 \end{tabularx} \\ \\
[1ex]
\hline
\hline
\end{longtable}
講題:深淵裡的拯救 Deliverance from Depths

經課:詩篇130:1-9

講員:李均熊牧師

場合:香港中文大學崇基學院禮拜堂主日崇拜

日期:2024年6月9日

起:詩一百三十篇主題是甚麼?

上次我在這裡宣講信息時,分享了如何研讀詩篇.正好今天和下一次我在這裡講道的啟應詩篇都會是詩篇裡面的上行之詩.所以今天和下一次講道,我都會選用經課中的啟應詩篇作為宣講信息的經文.

希望大家都記得,上一次我說到研讀詩篇需要沉浸在詩篇的經文內,透過不繼頌讀去發現詩篇的文學結構和投入在詩篇的意像語言中.但講道時間有限,我也不會細細分析經文的細節.剛好很多年前李思敬牧師在中神期刊中重譯過今日啟應詩篇的詩篇一百三十篇,所以我就借用了李思敬牧師的翻譯,來和大家細想詩篇一百三十篇的信息.李思敬牧師的翻譯如下:

詩的結構相對簡單,就是四個段落,每段落有兩行詩句,詩句有時是兩小句,有時是三小句.除了工整原則,所以八行分四段落外,你也可以發現一些字詞的分佈去支持這樣的分段.你可以留意一下上主(即和合本譯作「耶和華」的地方)一詞在每個段落都出現一次；而且第一,三段的主語的文法人稱不是「我」便是「你」,第二,四段的主語則是環繞「你」和「他」,這種文法上人稱的轉換,就是不懂希伯來文的讀者均可輕易察覺的分段提示.

這樣的分段帶出四個主題:第一段是詩人向上主呼求；第二段詩人說到罪與赦免；第三段主題則是等候；第四段焦點卻是上主的「解救」(和合本將同一個字根名詞譯「救恩」,動詞譯「救贖」).初看整首詩的主題可以連貫成詩人陷入如深淵中的境況中,他向上主呼求,而在他的等候中,向以色列他的同胞確認上主有充裕的解救(和合本:有豐盛的救恩).但若你真沉浸在這詩篇的經文中,你就會發覺有些地方好像不怎麼對勁.第二個段落所論到的上主究察罪孽和赦免之恩其實好像離了題,和詩人呼求,等候和確信上主解救有甚麼關係?

由中世紀開始,大公教會禮儀傳統上,會認為詩篇一百三十篇是七首「懺悔詩」(Penitential Psalms)其中之一(七首分別是詩篇六,卅二,卅八,五十一,一百零二,一百三十和一百四十三篇).很多釋經書也環繞「懺悔」這個主題去解釋詩篇一百三十篇.但若以沉浸式方法去細讀這首詩,不單大部份內容均與悔罪無關,而且你也找不到經文中有任何具體的懺悔,認罪,沒有承認任何罪行,也沒有因為懺悔而有的內疚和痛心(試比較其他的「懺悔詩」如詩篇五十一或一百四十三篇).詩篇一百三十篇所謂的「懺悔」,有的只有一句設問,就是主若究察罪孽誰能站得住,承受得到?這句似是詩人對人生的領會,多於他對自己的自身罪行有任何懺悔.而且詩人下一句就已經承認上主有赦免之恩；看來這個段落論說上主之恩多於詩人的認罪懺悔.

承: 能呼求上主全是恩典

若本詩第二段不是關於認罪懺悔,那麼細想這一段和整首詩就不會格格不入了.當詩人身在深淵般的處境,他便向上主呼求,懇求他的主垂聽他的聲音；但突想到自己其實並不配來到主面前呼求,因為自己本身就是一個罪人,「若你究察罪孽,誰能站得住呢?」若真要計算,我們又怎會有資格來到上主面前祈求幫助呢?若上主真要斤斤計較,沒有人可以在他面前承受他的審判,又有任何人可以來到主前作任何懇求呢?若上主有本小氣簿,記下所有我的過錯,我的名字應該寫滿這本簿子!但詩人轉念一想,不錯,若上主計較,沒有人可以承受,但就是因為在他有赦免之恩,所以我們才可以來到上主面前,才會在困難深處向祂祈求幫助解救呀!所以在第三段詩人轉向專心的等候,並且在第四段向他的以色列同胞確認在上主中有信實慈愛,在他又有充裕豐富的救恩和解救.就是這三個「在」(with):在你卻有赦免(v.4),在上主有信實(v.7),在祂有充裕的解救(v.7),成為在深淵中呼求的基礎.

(試想像:你因為經營不善,公司陷入債務「深淵」,當絕望的時候,想到可以向父親救助「解救」,當即想到自己其實仍欠父親一大筆債務,但你知道父親並不是斤斤計較的人,於是硬著頭皮向父親發出請求,現在等候他的回覆.同時你向公司員工保證,父親不單有「豐盛」的資金,更是講親情講義氣的,故此就一起「仰望」他,解救因過失而臨危的公司….)

轉:詩人的困境是甚麼?

李思敬牧師的分析中,也認為詩篇一百三十篇並非一首懺悔詩.他的文章主要追綜這首詩中的一些公式用語,在其他詩篇的用法,從而探究這首詩的主題.當中有兩句說話(v. 2 「我懇求的聲音」, v.5 「祂的話我正期待」),均以不同形式出現在詩篇內的其他的經文中.大家可以在週刊有關這詩的英文翻譯中找到相關的經文.

我提過研讀詩篇除了分析結構外,更需要投入詩的意境.公式用語的使用,就是在濃縮精煉的詩歌語言中,引發熟悉這些公式傳統的讀者產生聯想.我們未必如第一代讀者般,熟悉這些公式用語,但我們也可以透過串珠(cross-refrence)查看這兩句公式用語在不同詩篇的引用,來解答經文中(因為詩歌的濃縮語言而)沒有解答的問題如:究竟詩人在怎樣的深淵?他的懇求是甚麼?他等候的又是甚麼?為甚麼他最後可以確信上主有豐盛的救恩?

例如在第三段落中,究竟詩人等候甚麼?其實在這一段落中是有答案的:「祂的話我正期待」(和合本譯:我也仰望他的話).但這又代表甚麼呢?李思敬從這一句公式用語在詩篇一一九篇四次的出現(v. 74, 81, 114, 147),綜合出這句話,代表詩人不是單純等待上主的答覆,或神諭oracle,或一般的妥拉等等,其實是等待上主幫助他脫離惡人的攻擊加害.

原來在詩篇一一九篇,一篇不斷強調上主的話語如何美好的一篇詩,細讀之下你會領略到詩篇一一九篇內的詩人正面對惡人的逼迫,攻擊和羞辱,而上主的話語卻成為他活命的力量,使他深信上主必實現祂的應許,解救他脫離惡人的圍困.詩篇一一九篇共有22x8=174節經文,今日就不多討論.這裡只抽兩段來回應詩篇一百三十篇詩人究竟深入怎樣的深淵:

其中一段,和合本標題就是「求主解救」:

119:81我心渴想你的救恩,仰望你的應許(或譯:你的話我正期待).

119:82我因盼望你的應許眼睛失明,說:你何時安慰我?

119:83我好像煙薰的皮袋,卻不忘記你的律例.

119:84你僕人的年日有多少呢?你幾時向逼迫我的人施行審判呢?

119:85不從你律法的驕傲人為我掘了坑.

119:86你的命令盡都誠實；他們無理地逼迫我,求你幫助我!

119:87他們幾乎把我從世上滅絕,但我沒有離棄你的訓詞.

119:88求你照你的慈愛將我救活,我就遵守你口中的法度.

當中詩人的困境就是被驕傲人陷害,被人無理逼迫,「幾乎把我從世上滅絕」!

另外一段就更能和詩篇一百三十篇呼應,和合本標題都是「求主解救」:

119:145耶和華啊,我一心呼籲你；求你應允我,我必謹守你的律例!

119:146我向你呼籲,求你救我!我要遵守你的法度.

119:147我趁天未亮呼求；我仰望了你的言語(或譯:你的話我正期待).

119:148我趁夜更未換將眼睜開,為要思想你的話語.

119:149求你照你的慈愛聽我的聲音；耶和華啊,求你照你的典章將我救活!

119:150追求奸惡的人臨近了；他們遠離你的律法.

119:151耶和華啊,你與我相近；你一切的命令盡都真實!

119:152我因學你的法度,久已知道是你永遠立定的.

這一段說詩人趁天未亮呼求,趁夜更未換就思想,期待上主的話,因為「追求奸惡的人臨近了」,甚麼是惡人呢:「他們遠離你的律法」.149節「求你照你的慈愛聽我的聲音；耶和華啊,求你照你的典章將我救活!」中心思想也和詩篇一百三十篇相約.

詩一百三十篇第2節有另一句也在其他詩篇出現的公式語:「我懇求的聲音」,若在其他詩篇出現時,詩人懇求的都是和他面對的「惡人」有關.例如「我懇求的聲音」出現了兩次的詩篇廿八篇,用李思敬文章的說法:「作者向上主的祈禱非常清晰:他懇求上帝報應惡人,因為這樣就等如幫助和拯救了上主的子民.」其他如詩卅一篇或一百四十篇,當中詩人懇求的聲音都包括求主救自己「脫離仇敵的手和那些逼迫我的人」(31:15)或「求你不要遂惡人的心願；不要成就他們的計謀」(140:8)

雖然我們其實也不知詩人真確的敵人和仇敵是誰,那些「惡人」夠竟是社會上的人物還是詩人身邊的同事或親人,但這也是我們能沉浸在詩篇經文的原因,若然詩篇出現的惡人敵人仇人都是實在的人物,那我們就未必能將詩篇經文代入我們今日的處境了.但起碼從這些公式用語中,我們了解到詩人所陷入的深淵,並不是甚麼財困或疾病,而是壞人的加害和不從神律的惡人的逼迫.可能正因如此,當詩人要求上主去對付罪人惡人時,其實也想到自己也並不比這些罪人和惡人好,因為上主「若究察罪孽,誰能站得住呢?」,我們能面對敵人和惡人,並不是我們有甚麼義舉,而只是上主有赦免之恩,願意解救祂自己的立約的子民(慈愛原文乃指立約的信實),也要救贖祂的子民脫離罪惡／罪孽.

(詩129-131上行之詩的流散處境:提供一個歷史向度,讓我們探索詩篇一百三十篇中的困境是甚麼,當中涉乃的「惡人」會是誰——下回再論)

合:等候與盼望

綜合來說,詩篇一百三十篇是一個人因為周遭惡人和不義的人的攻擊和圍困,故此深感生命落入深淵危險中,而發聲向上主懇求解救.但我們又憑甚麼可以向上主呼求呢?我們比我們的敵人更有義嗎?不是的,「你若究察罪孽,誰能站得住呢?」但因為「在」上主有赦免之恩,「在」祂有充裕的救恩,更重要的是「在」上主向子民守約的信實和慈受,故此我們能夠在深淵中呼求,在黑夜中等候,在困境中期待.不單我們如此,我們更要在神子民中,鼓勵身邊人一同仰望上主,在這個惡人當道,不容我們發聲的黑暗世代中(!),向上主呼求,和期待上主在深淵中救我們脫離一切的罪孽和凶惡.

最後,詩篇一百三十篇還有一個意像,我還未細講:那就是詩人等候上主的解救,勝於守夜者等候天亮的意境.大家試過捱夜未?詩人說等候上主要「勝於」守夜的等候天亮!那是怎樣的一種決心和毅力?或者大家這個星期就投入去這個意境中,沉浸在這首詩篇裡,細味如何在艱難日子仍可等候上主的恩典.



\newpage

\section{無處話淒涼}
\label{sec:701}
\textbf{鄧瑞強博士 - 無處話淒涼}
\newline
\newline
\hyperref[sec:700]{< < < PREV SERMON < < <}
~
\hyperref[sec:toc]{[返目錄]}
~
\hyperref[sec:702]{> > > NEXT SERMON > > >}


講題:無處話淒涼 Hope in Gloom

經課:耶利米哀歌1:1-5,2:20-21,3:22-27,4:9-10,5:21-22

講員:鄧瑞強博士

場合:香港中文大學崇基學院禮拜堂主日崇拜

日期:2024年6月16日

各位弟兄姊妹,早安.願神的恩典常與你們同在.

戰爭,從來都是悲慘的.但很不幸,縱使明知是悲慘,人類仍要戰爭.

有一首漢代的樂府詩《戰城南》:

「戰城南,死郭北,野死不葬烏可食.為我謂烏︰『且為客豪.野死諒不葬,腐肉安能去子逃?』」

戰士在城南打仗,城北有很多戰死沙場的人.暴屍荒野,烏鴉來啄食.詩人對烏鴉說,可否先為戰死者哀悼,他們的屍體一定無法安葬,腐肉又怎能逃得過你們的啄食呢?

戰爭剛結束,遍地屍骸,無處話淒涼.

在耶利米時代,巴比倫大軍圍攻耶路撒冷,圍城1年半,於主前587年攻下耶路撒冷,維時400多年的大衛王朝終結,猶大亡國,人民被擄.今日,我們看《耶利米哀歌》,這是劫後餘生,沒有被擄去巴比倫的耶路撒冷人的哀歌.他們坐在廢墟之上,故國不堪回首,眼下無盡淒涼.透過這哀歌,他們沒有掩飾的描繪人間的悲劇,表達無盡的哀愁,也流著眼淚,無語問蒼天.

先看《耶利米哀歌》1:1-5:

1:1　先前滿有人民的城,現在何竟獨坐!先前在列國中為大的,現在竟如寡婦；先前在諸省中為王后〔公主〕的,現在成為進貢〔奴隸〕的.

1:2　她夜間痛哭,淚流滿腮；在一切所親愛的中間沒有一個安慰她的.她的朋友都以詭詐待她,成為她的仇敵.

1:3　猶大因遭遇苦難,又因多服勞苦就遷到外邦.她住在列國中,尋不著安息；追逼她的都在狹窄之地將她追上.

1:4　錫安的路徑因無人來守聖節就悲傷；她的城門淒涼；她的祭司歎息；她的處女受艱難,自己也愁苦.

1:5　她的敵人為首；她的仇敵亨通；因耶和華為她許多的罪過使她受苦；她的孩童被敵人擄去.

1:1　先前滿有人民的城,現在何竟獨坐!先前在列國中為大的,現在竟如寡婦；先前在諸省中為王后〔公主〕的,現在成為進貢〔奴隸〕的.

詩人將耶路撒冷比喻為一個女子.按原文結構,1:1的開始是「何竟 坐 孤獨」.寂寞的女子,唉聲歎氣,坐在那裡,一個人.她在做什麼呢?她想起前塵往事,對比當下的淒涼.想起城巿以前的光景,「滿有人民」,現在只得孤身一人.我想起亡國被囚的李後主,他也有相似的感歎.「多少恨,昨夜夢魂中.還似舊時游上苑,車如流水馬如龍.花月正春風.」過去越繁華,當下越淒涼.詩人講到猶大國以前的尊貴,現在竟似一個失去倚靠的寡婦.以前像公主,現在只是服事他人的奴隸.李後主的《浪淘沙》,也有相似意境.「簾外雨潺潺,春意闌珊；羅衾不耐五更寒.夢裡不知身是客,一晌貪歡.獨自莫憑闌,無限江山,別時容易見時難!流水落花春去也,天上人間!」孤寂冰冷,只有在夢中,才能忘記自己已是囚徒,暫享片刻的歡愉.想起舊時的江山,已是永別.一切已過去,永不再返.很傷感,很淒涼.

1:2　她夜間痛哭,淚流滿腮；在一切所親愛的中間沒有一個安慰她的.她的朋友都以詭詐待她〔按:特別是她的盟友埃及,在巴比倫圍城時,為了自保,沒有幫助她〕,成為她的仇敵.

1:3　猶大因遭遇苦難,又因多服勞苦就遷到外邦.她住在列國中,尋不著安息；追逼她的都在狹窄之地將她追上.

1:3指出,留在廢墟中的人,生活太苦,不得不流落異鄉.但生活的苦難,仍不斷纏著他們.這些是戰爭後的難民,因無法在原地生活,飄泊他方.

1:4　錫安的路徑因無人來守聖節就悲傷；她的城門淒涼；她的祭司歎息；她的處女受艱難,自己也愁苦.

留意一下這經文的用字:悲傷,淒涼,歎息,艱難,愁苦.整座城巿,代表精神的祭司,代表帶來下一代希望的處女,全都在歎息中.

1:5　她的敵人為首；她的仇敵亨通；因耶和華為她許多的罪過使她受苦；她的孩童被敵人擄去.

敵人成了她的頭,主宰了她的命運.她的孩童,這是城巿的將來,也沒有了.這一切,全因猶大自身的罪,神因她的罪審判他們.

我們看看第2章,詩人描繪慘況,也將這慘況陳示在神面前.

2:20　耶和華啊,求你觀看!見你向誰這樣行?〔新譯本:你曾這樣對待過誰呢?〕婦人豈可吃自己所生育手裡所搖弄的嬰孩嗎?祭司和先知豈可在主的聖所中被殺戮嗎?

2:21　少年人和老年人都在街上躺臥；我的處女和壯丁都倒在刀下；你發怒的日子殺死他們.你殺了,並不顧惜.

戰爭就是如此淒慘,到處都是屍體,而母親要吃自己的嬰孩.詩人要求神看看,「耶和華啊,求你觀看!」這裡有一種很微妙的心情.一方面,若神是歷史主宰,則這一切都關祂的事.詩人似是向神投訴.另一方面,詩人是將現實帶到神面前,請神去看.若詩人根本不相信神能做些什麼,則叫神看也於事無補.詩人默然地,對於神能處理這災難現實,有一種信任.

面對慘痛的現實,詩人寄望上帝的仁慈,同時,在神面前呈示這些慘況,事實上是叫人類正視自己的責任.說到底,是人類的自私自利,搞成這個局面.

我們看看第3章,這是《耶利米哀歌》的中心部分.

3:22　我們不致消滅,是出於耶和華諸般的慈愛；是因他的憐憫不致斷絕.〔新譯本:「耶和華的慈愛永不斷絕,他的憐憫永不止息.」按:如此翻譯,「不致消滅」的,就不是「我們」,而是神的慈愛.〕

3:23　每早晨,這都是新的；你的誠實極其廣大!

3:24　我心裡說:耶和華是我的分,因此,我要仰望他.

在絕望深淵,詩人靈光一閃,找到在苦難中存活下去的勇氣.在人間,流水落花春去也,什麼也抓不住.但,有些東西是永恆的,是歷史的浪花淘不盡的,是支撐著人類努力生存下去的力量.永不斷絕的,是神的慈愛；永不止息的,是祂的憐憫.歷史,不是苦難的輪迴.神的恩典,神的憐憫,能打破這惡性循環,人間是有新的可能性的.新的一天,可以不是昨天苦難的延續,而是神的光明朗耀的一天.「神的誠實」,即神守約的決心,即神憐憫人的意志,是牢不可破的.

詩人說:「耶和華是我的分」,所謂「分」,即祖先傳下來的產業.耶和華是我們的信仰前輩傳給我們的產業.這產業不是我們賺來的,是一分恩典,若我們好好珍惜,這產業能再傳下去.在苦難深淵裡的人,覺得一無所有,但事實不是這樣的,耶和華本身就是我們的產業.祂是存在,光明,意義的源頭,有了祂,人就有依傍,能堅強地活下去.在艱難中的弟兄姊妹,抓緊神,堅強地活下去.

3:25　凡等候耶和華,心裡尋求他的,耶和華必施恩給他.

3:26　人仰望耶和華,靜默等候他的救恩,這原是好的.

3:27　人在幼年負軛,這原是好的.

這3節經文的開頭,希伯來文都是「tov」(good, 好).3:25原文是這樣的:「好 耶和華 對那些等候祂的 對那些生命 尋求祂的」.在苦難中,詩人講到,不是全部東西都是壞的,有些東西是好的.等候神,是好的.等候神的救恩,是好的.人在幼年負軛,是好的.相反來說,不等候神,反去靠賴偶像,去改變風水,努力改運,這是不好的.不等候神的救恩,反而放棄希望,怨天尤人,這是不好的.至於幼年負軛,這是好的,這與孟子講的道理是一致的.孟子說:「故天將降大任於是人也,必先苦其心志,勞其筋骨,餓其體膚,空乏其身,行拂亂其所為,所以動心忍性,曾益其所不能.」艱難的日子,令人更謙和,也更堅強的面對人生.在困苦的歲月裡,黑暗居多,靈光很少；負面的思想多,信仰的心靈少.當永恆之光一閃而過時,詩人卻抓住了,他找到冥冥中上帝的恩情,他靠這恩情,努力地活下去.

殘酷的事實當然不會因此而改變,我們看《耶利米哀歌》第4章.

4:9　餓死的不如被刀殺的,因為這是缺了田間的土產,就身體衰弱,漸漸消滅.

4:10　慈心的婦人,當我眾民被毀滅的時候,親手煮自己的兒女作為食物.

為何《哀歌》不停在第3章?剛才,才走到信仰的高峰,為何忽然卻又回到現實的低谷?有了信仰,不就是將艱難一筆勾銷了嗎?原來,真正的信仰並不會取消現實的艱難,而只是給我們力量在艱難的現實中實現永恆的意義.在客西馬尼園,在十架前夕,耶穌向天父祈禱,祈禱後,他明白了天命,抓住了天命,現實的苦難不會因此而消失,只是耶穌有了力量去面對苦難,在苦難中展示人的生存的神聖意義.

我們看看《耶利米哀歌》的結尾.

5:21　耶和華啊,求你使我們向你回轉,我們便得回轉.求你復新我們的日子,像古時一樣.

5:22　你竟全然棄絕我們,向我們大發烈怒?

〔新譯本:耶和華啊!除非你真的棄絕我們,你向我們所發的怒氣難收,不然,求你使我們歸向你,好使我們回轉過來；求你更新我們的日子,像古時一樣.〕

無論怎樣翻譯,《哀歌》的最後一節經文,仍是懸疑未決:神是否棄絕我們,向我們大發烈怒?用這種較負面的經文作經卷結束,在聖經裡是少見的.或許,現實的痛苦太強烈,並不容許輕言希望.但又不是不講希望,希望的信息在第3章這經卷的中心,只是這希望不是「啊,無事了」這種輕省的希望.面對艱難,有些信徒懷著巨大的信心祈禱,以為祈禱完後,艱難便消失了.當艱難沒有消失,不少信徒便失望了,覺得神好無用,於是,放棄神了.其實,不是神棄絕人,而是人在絕望中棄絕神.有些信徒,卻能像耶穌,祈禱後,便倚靠神的力量,面對艱難,開出新生命的前景.縱使艱難仍在,這些人找到在現實以外及感覺以外的永恆上帝.就因為上帝,他們有力量,在絕境中前行.

願上主得著榮耀,頌讚；願眾人得享喜樂,平安.



\newpage

\section{老師我們喪命你不顧嗎}
\label{sec:702}
\textbf{高國雄牧師 - 老師,我們喪命,你不顧嗎?}
\newline
\newline
\hyperref[sec:701]{< < < PREV SERMON < < <}
~
\hyperref[sec:toc]{[返目錄]}
~
\hyperref[sec:703]{> > > NEXT SERMON > > >}


馬可福音 4:35-41
\newline
\begin{longtable}{cl}
\hline
\hline
章節 & 經文 (和合本修訂版)\\
\hline
4:35 & \begin{tabularx}{0.7\textwidth}{X} 那天晚上，耶穌對門徒說：「我們渡到對岸去吧。」 \end{tabularx} \\ \\ \relax
4:36 & \begin{tabularx}{0.7\textwidth}{X} 門徒離開眾人，耶穌已在船上，他們就請他一同去；也有別的船和他同行。 \end{tabularx} \\ \\ \relax
4:37 & \begin{tabularx}{0.7\textwidth}{X} 忽然狂風大作，波浪打入船內，以致船灌滿了水。 \end{tabularx} \\ \\ \relax
4:38 & \begin{tabularx}{0.7\textwidth}{X} 耶穌在船尾上，枕著枕頭睡覺。門徒叫醒他，說：「老師！我們快沒命了，你不管嗎？」 \end{tabularx} \\ \\ \relax
4:39 & \begin{tabularx}{0.7\textwidth}{X} 耶穌醒了，斥責那風，向海說：「住了吧！靜了吧！」風就止住，大大平靜了。 \end{tabularx} \\ \\ \relax
4:40 & \begin{tabularx}{0.7\textwidth}{X} 耶穌對他們說：「為甚麼膽怯？你們還沒有信心嗎？」 \end{tabularx} \\ \\ \relax
4:41 & \begin{tabularx}{0.7\textwidth}{X} 他們就非常懼怕，彼此說：「這到底是誰？連風和海都聽從他。」 \end{tabularx} \\ \\
[1ex]
\hline
\hline
\end{longtable}
講題:「老師,我們喪命,你不顧嗎?」 “Teacher, do you not care that we are perishing?”

講員:高國雄牧師

經文:馬可福音4:35-41

場合:香港中文大學崇基學院禮拜堂主日崇拜

日期:2024年6月23日

35當那天晚上,耶穌對門徒說:「我們渡到那邊去吧.」

36門徒離開眾人,耶穌仍在船上,他們就把他一同帶去；也有別的船和他同行.

37忽然起了暴風,波浪打入船內,甚至船要滿了水.

38耶穌在船尾上,枕著枕頭睡覺.門徒叫醒了他,說:「夫子!我們喪命,你不顧嗎?」

39耶穌醒了,斥責風,向海說:「住了吧!靜了吧!」風就止住,大大地平靜了.

40耶穌對他們說:「為甚麼膽怯?你們還沒有信心嗎?」

41他們就大大地懼怕,彼此說:「這到底是誰,連風和海也聽從他了.」

當耶穌吩咐你要行動,但在月黑風高的晚上,在危險滿佈的時間…

你會如何回應呢?

「我們渡到那邊去吧.」主要與我們一同冒險

面對風浪,你會選擇…認命?逃避?埋怨?求救?

「夫子!我們喪命,你不顧嗎?」

向老師的呼救 VS 對耶穌的認識/期望

耶穌醒了:耶穌睡著?耶穌仍然睡著?

原來耶穌在斥責我們之前,先斥責(命令)風

風浪也願意馬上聽命而安靜,我們呢?

40耶穌對他們說:「為甚麼膽怯?你們還沒有信心嗎?」

「我們應該要有信心」???

過去的經歷…耶穌的同在…耶穌的使命

41他們就大大地懼怕,彼此說:「這到底是誰,連風和海也聽從他了.」

我們因什麼而懼怕?

風浪?耶穌的斥責?

耶穌的能力(太厲害)?這不是我們所期待的嗎?

耶穌與門徒忙碌了一整天,教導了幾個重要的比喻(撒種,芥菜種)

道與天國

耶穌接著就用行動展現出來

邀請同往…目的地就是格拉森

可5:21耶穌又坐船渡到對岸…

目標就只有一個…去拯救那一個在受苦,被捆綁中的弟兄

耶穌因顧念格拉森的弟兄,而願意與門徒冒這風險,怎可能會不顧所愛的門徒呢?

「老師,我們喪命,你不顧嗎?」

這是”少信”的寫照,也是要我們經歷信心的開始



\newpage

\section{在絕望孤單中偶然遇到你}
\label{sec:703}
\textbf{林豪恩先生 - 在絕望孤單中偶然遇到你}
\newline
\newline
\hyperref[sec:702]{< < < PREV SERMON < < <}
~
\hyperref[sec:toc]{[返目錄]}
~
\hyperref[sec:704]{> > > NEXT SERMON > > >}


講題:在絕望孤單中偶然遇到你 Meeting You by Chance in Despair and Loneliness

講員:林豪恩先生

場合:香港中文大學崇基學院禮拜堂主日崇拜

日期:2024年6月30日

引言

從前在中大讀書時,認識一位讀工程的師兄,原來他會寫作詩歌.後來,他的其中一首詩歌走紅,歌名為「偶然遇上的驚喜」,其中幾句歌詞如下:

能在絕望孤單中偶然遇到你,是我一生中最大的福氣.

來日禍困迎面,滿途挑戰,能贈我勇氣,完全是你.

能在絕望孤單中偶然遇到你,的確是人生的福氣.剛才聆聽經課的時候,有沒有聽到幾個類似的故事呢?

一．疾病令人絕望孤單(馬可福音)

馬可福音記述了一個因為疾病而陷入絕望孤單的故事:

5:25有一個女人,患了十二年的血漏.

郭詠觀醫生在「病患者的心理健康」一文中指出　:「重病會令人感覺哀傷,驚慌,憂慮或惱怒,因為患者會覺得對自己的身體和整體情況失去了控制,並可能認為自己無能為力,患者也會有與親友疏離和寂寞的感受.」

故事中的女人患的是血漏,當時被視為不潔,被煎熬和隔離十二年之久,絕望孤單可想而知.禍不單行的是:

5:26在好些醫生手裡受了許多的苦,又花盡了她所有的,一點也不見好,病勢反倒更重了.

或許她已經不再把希望寄託在人身上,也再沒有錢財求醫.然而,最近有一個人出現,又燃起了她的希望.準確地說,這不是一個人,因為人不可能有「我只摸他的衣裳,就必痊癒」的能力.這次,她找對了對象,不花一分一亳,頑疾就醫好了.

5:30耶穌頓時心裡覺得有能力從自己身上出去,就在眾人中間轉過來,說:「誰摸我的衣裳?」

疾病都已經醫好了,為甚麼耶穌不讓她靜靜地離去呢?疾病醫好,不等於脫離絕望孤單的狀況.耶穌在人群中宣布這個人已經痊癒,宣告她蒙受上主眷顧及恩惠,鼓勵社群接納她,也鼓勵她再投入社群,離開絕望孤單的境況.

二．死亡令人絕望孤單(馬可福音)

故事中還有另一個人「快要死」的人.這令我們想到,死亡會不會是另一種令人絕望孤單的狀況呢?

5:22有一個管會堂的人,名叫睚魯,來見耶穌,就俯伏在他腳前,

5:23再三地求他,說:「我的小女兒快要死了,求你去按手在她身上,使她痊癒,得以活了.」

雖然垂死是小女孩,但故事的主角是她的爸爸.「白頭人送黑頭人」更是悽愴.河南有一位36歲男士的女兒5歲時被診斷出癌症.夫婦帶女兒到北京治療,花盡儲蓄但無法挽救女兒性命.女兒落葬後,父親轉工到墓園工作,為了「陪伴」女兒.每天擦拭及親吻女兒的墓碑,再與女兒聊天,在其墓碑前擺放了各種玩偶,水果和鮮花.(晴報2021/02/02)

5:38他們來到管會堂的家裡；耶穌看見那裡亂嚷,並有人大大地哭泣哀號,

5:41就拉著孩子的手,對她說:「大利大,古米!」(翻出來就是說:「閨女,我吩咐你起來!」)

5:42那閨女立時起來走.他們就大大地驚奇；閨女已經十二歲了.

面對死亡,有時可以奇蹟地起死回生,然而,更多情況是一去不返,販夫走卒如是,帝王將相也如是.

三．終結的絕望孤單(撒母耳記下)

歷史上,成者為王,敗者為寇.掃羅和三個兒子戰死沙場,被非利士人割下首級侮辱.有人偷去他的冠冕,送到他的對手大衛那裡.人都有一死,人死的時候,能否給他一點尊嚴呢?大衛立即撕裂衣服,並作哀歌,弔掃羅和他兒子約拿單.

1:19歌中說:以色列啊,你尊榮者在山上被殺!大英雄何竟死亡!

1:20不要在迦特報告；不要在亞實基倫街上傳揚；免得非利士的女子歡樂；免得未受割禮之人的女子矜誇.

1:21基利波山哪,願你那裡沒有雨露!願你田地無土產可作供物!因為英雄的盾牌在那裡被污丟棄；掃羅的盾牌彷彿未曾抹油.

1:22約拿單的弓箭非流敵人的血不退縮；掃羅的刀劍非剖勇士的油不收回.

1:23掃羅和約拿單－活時相悅相愛,死時也不分離－他們比鷹更快,比獅子還強.

1:24以色列的女子啊,當為掃羅哭號!他曾使你們穿朱紅色的美衣,使你們衣服有黃金的妝飾.

1:25英雄何竟在陣上仆倒!約拿單何竟在山上被殺!

人人都知道死亡無可逃避,能否有尊嚴地步過最後一程呢?德蘭修女有一次在污水渠檢到一名男子,把他帶回修院.他滿身蛆,而且很髒.修女們花了三小時為他清潔,把蛆統統拿掉.他死時面露微笑,有尊嚴,被愛,被關懷,而且與神一起,很安詳.他說:我像野獸般在街上生活,但卻像天使般死去.

四．以群體回應絕望孤單(哥林多後書)

當年那位女士和少女在絕望孤單之中遇到耶穌,得到拯救.其他人又怎樣呢?保羅寫給哥林多信徒的書信,提示我們,人與人之間的互相幫助也是拯救絕望孤單的途徑.

8:7你們既然在信心,口才,知識,熱心,和待我們的愛心上,都格外顯出滿足來,就當在這慈惠的事上也格外顯出滿足來.

上星期我們在柬埔寨,同學把香港和柬埔寨比較之後說:世界真是不公平的.如果我出生在柬埔寨,我的成長就不一樣了.我和柬埔寨人有甚麼分別呢?分別只是我出生在香港,他們出生在柬埔寨.當年,保羅就勸勉條件較好的哥林多信徒向耶路撒冷信徒伸出援手:

8:13我原不是要別人輕省,你們受累,

8:14乃要均平,就是要你們的富餘,現在可以補他們的不足,使他們的富餘,將來也可以補你們的不足,這就均平了.

人人都有條件自力更生當然理想,但際遇難料,如果落入不能夠照顧自己的狀況又如何呢?藝人鄭俊弘與何雁詩夫婦接受訪問,公開兒子被診斷患罕見遺傳病天使綜合症.何雁詩稱身為父母,最擔心的是,「我們可以照顧他多久?我們將來年老後他會如何呢?」這不單單是鄭俊弘與何雁詩夫婦的擔心,也是普天下有特殊孩子的父母的擔心.

8:9你們知道我們主耶穌基督的恩典:他本來富足,卻為你們成了貧窮,叫你們因他的貧窮,可以成為富足.

保羅指出,互通有無,互相幫助,不但是人類自身的需要,更是主耶穌基督的榜樣.「能在絕望孤單中偶然遇到你,是我一生中最大的福氣.」對患血漏病的女人以及少女的父親,那個「你」是耶穌；對掃羅和約拿單,那個「你」是大衛；對耶路撒冷的信徒,那個「你」是哥林多的弟兄姊妹.在座各位,你會否就是其中一位別人在絕望孤單中偶然遇到的「你」呢?



\newpage

\section{上主啊求你憐憫我們}
\label{sec:704}
\textbf{李均熊牧師 - 上主啊!求你憐憫我們}
\newline
\newline
\hyperref[sec:703]{< < < PREV SERMON < < <}
~
\hyperref[sec:toc]{[返目錄]}
~
\hyperref[sec:705]{> > > NEXT SERMON > > >}


詩篇 123:1-4
\newline
\begin{longtable}{cl}
\hline
\hline
章節 & 經文 (和合本修訂版)\\
\hline
123:1 & \begin{tabularx}{0.7\textwidth}{X} 坐在天上的主啊，我向你舉目。 \end{tabularx} \\ \\ \relax
123:2 & \begin{tabularx}{0.7\textwidth}{X} 看哪，僕人的眼睛怎樣仰望主人的手，婢女的眼睛怎樣仰望女主人的手，我們的眼睛也照樣仰望耶和華－我們的神，直到他憐憫我們。 \end{tabularx} \\ \\ \relax
123:3 & \begin{tabularx}{0.7\textwidth}{X} 耶和華啊，求你憐憫我們，憐憫我們！因為我們受盡了藐視。 \end{tabularx} \\ \\ \relax
123:4 & \begin{tabularx}{0.7\textwidth}{X} 我們受盡了安逸人的譏誚和驕傲人的藐視。 \end{tabularx} \\ \\
[1ex]
\hline
\hline
\end{longtable}
講題:上主啊!求你憐憫我們 Have Mercy on Us, O Lord

經文:詩篇 123:1-4

講員:李均熊牧師

場合:香港中文大學崇基學院禮拜堂主日崇拜

日期:2024年7月7日

上行之詩

1 向你我抬起我的眼睛 (我-你)

~~~那坐在天上的!

2 看哪!如眾僕的眼睛

~~~對著他們主人的手

如使女僕的眼睛

~~~對著她主母的手

照檬我們的眼睛對著上主我們的神 (我們-上主)

~~~直到祂施恩於我們

3 施恩於我們!上主阿!施恩於我們 (我們-上主)

~~~因我們極致充滿了藐視

4 極致地充滿了是我們的心靈

~~~安逸人的譏誚

~~~~~~驕傲人的藐視

起: 詩123的文學特色

今天的啟應詩篇是詩篇123篇.4節經文共6行的詩句,看來應十分容易處理.但這麼短小精悍的詩篇,其實反而十分難掌握.詩人將要說的意念濃縮成只有三兩句說話,而且多個用語重重覆覆(眼睛,手,施恩／憐憫,極致／極點,飽受／充滿,藐視),有話卻不說盡,自然令人費解.用我們一直以來欣賞聖經詩篇的方法,就是沈浸在詩的結構與意象中,詩123對讀詩的人也是不容易分解.

初看6行詩句,用工整原則,若不是3行一段分2個段落,就會是2行一段分三個段落.但這詩的2到4行,聖經分節的第2節,用和合本的翻譯來讀:「僕人的眼睛怎樣望著主人的手,婢女的眼睛怎樣望著主母的手,我們的眼睛也照樣望著耶和華我們的　神,直到他向我們施恩.」這是一句完整的句子,第2行和第3行的喻象,要到第4行才看到比喻的喻意,並產生一種延後懸疑的文學效果,是詩篇文學絕皆的詩句.因此本詩不能機械地將行數分割為等份.以第2節第2至4行三行詩句為一段,則本詩第1節一行為一段,詩3和4節第5,6行兩行詩句作為結束,全詩便共分為一行,三行,二行不規則的三個段落了.這個分段可以從詩句的文法主語變化得到確定:第一行作為引言詩句主語是「我」(第一人稱單數),對話的對像是你——那坐在天上的(參週刊).而第二段比喻部份則詩句主語是「我們」(第一人稱眾數),上主在此是第三人稱——我們的神,對話對像應該也是「我們」.最後總結一段主語仍是「我們」(第一人稱眾數),但上主則成為呼求的對像,回到第一段落一般,成為對話的對像.

詩的喻象部份也是十分耐人尋味:究竟僕人使女望著主人主母的手,實際上在期待著甚麼呢?剛才說過這一段運用了延後懸疑的效果,用上兩個平行喻象才說出喻意.兩個喻象也充滿文學色彩,不單用上男僕女僕,細看更是男僕為眾數,女僕為單數的文法變化,有解家認為這是代表全面,完整,但大部份學者則認為這種陰陽單眾的變化純粹文學裝飾,沒有任何修辭意義.這種陰陽單眾的文法變化也出現在最後一段,動詞飽受／充滿在第3節(第5行)主語是「我們」(第一人稱眾數)；在第4節(第6行)則為「我們的心靈」,文法上是第三人稱陰性單數主語.這可以是詩人連結各個段落的文學手法,沒有甚麼意義上的加增.

說了這麼多,無非是想大家看到,這一首短少的詩篇,其實是一篇文學手法十分高超的作品,大家懂得欣賞的話,確能引人入勝,讓讀者沈浸體會詩的意境.

但在意義上,這首詩對我們現代人其實仍然十分費解.就算你看上千次萬次,不明白就是不明白.我已說過,僕人望著主人的手其實是期望甚麼呢?我們不太知道.更令現今讀者費解的,可能就是詩人祈求上主施恩憐憫的原因.我們或會想像,一個人可能身患重疾,落入懷人圈套,被惡人迫害,生命有危險等等重大危機,我們會呼求上主施恩憐憫.但看詩最後一段,詩人是因為飽受／心靈充滿了別人的藐視和譏誚,而向上主呼求憐憫.究竟這個飽受藐視,譏誚的境況是一個怎樣的處境,令到詩人要這樣呼求天上的主?

承:為何事祈求施恩憐憫?——羞辱的社會代價

當然,今日我們也會緊張我們的名聲名譽.我們也知道社交媒體有時做成的網絡欺凌,可以致人與死.但今日亦有人淡薄名利,沒有Facebook account,對於別人的閒話置之不理,自信清者自清,只要自重自量,自然不怕別人的貶抑.不過今日有那麼健康的自我意識其實並不容易,尤其年輕人今時今日的自我形像實在建立得不太完整(大家應該看看《反轉腦朋友2》Inside Out 2).在自由社會,我們能開放,真誠地建立自我的健康形像意識,自然不會理會別人的所謂「藐視」.

不過,若你看電影間諜片之類看得多的話,有些極權社會要毀滅一個人,並非一定要取其性命,而是誣陷這個人,使之身敗名裂.有些極權社會更會進行社會評級,作為操控人民的手段.就是簡單的信貸評級(credit score),若管理不善也會帶來很多個人在社會的問題.可見,詩篇這裡所說的藐視譏誚,並不是單純的對個人心理健康的挑戰.當中其實是涉及名譽羞辱的社會性功能,信譽,名聲,這些其實是一些重要的社會資本,當一個人,或一個群體,在社會中被羞辱,或今天常說的污名化(stigmatization),這個人或這些群體在社會中將會面對極大的困難,甚至會帶來自身的危險(南亞,中亞或中東的「名譽殺人」(honor killing),互聯網上各種的文化戰,和宗教,種族,甚至性別仇恨).

究竟詩篇123所要展示的處境是一個怎樣的處境.要理解詩篇123單純的文學結構分析可能未必足夠,我們也要進入詩篇成文的社會文化處境,才可以了解詩篇對當時讀者的意義,從而理解在千多二千年後對我們的意義.

近百年來,人類學家或歷史學家的研究都指出古代世界,特別生成聖經經文的西亞(近東)或地中海區域,榮辱文化(honor-shame culture)對當時的社會是極重要的運作原則.榮譽,關乎個人社會的地位聲望,對男性而言,通常與他的勇氣,忠誠,可信等特質相關；而女性的名聲更是影響女性在社會能否被接納的重要因素,當中甚至操控著女性的舉止,衣著,作息,或可以出入的場所,可以交住的人士等等.一個人若被羞辱,大多數會被主流社會排斥,男的可能不能參與任何社會活動, 無論是工作機會或宗教祭典,女的甚至會被仇視,以至生命受到威脅.正如我先前說到,社會信譽在今日社會也是十分重要的社會資本,但若放在古代,例如詩篇123猶太讀者身處的境況,其影響和嚴重性可以放大十倍百倍也不為過.

轉:上行之詩的流散背景

但對於猶太讀者,為甚麼會認同因為人的藐視而要向上主呼求施恩,乃是因為詩篇123所展示的不是一個個人的境況,而更可能是在被擄流散的處境,整個民族在外邦帝國統治下,都成了藐視的對像.

詩123是詩篇第五卷其中一首上行之詩.在詩篇中,標題為上行之詩的共有十五首,由詩篇120首到134首.上月我講道的經文詩130篇也是其中一首.十五首上行之詩在詩篇第五卷為一個特別的集結,其中各詩之間有很多微妙的關係,例如今日經文詩123開首「我向你舉目」,就和大家熟識的詩121開首「我向山舉目」異曲同工,其他有關聯的字眼也包括「坐在天上」,若比較另一首上行之詩,詩132:12,14裡坐著其實是指王座,故有學者譯詩123:1為「那坐天上寶座的一位」.上行之詩也常用重覆,遞進,呼應等等文學手法,詩123也與詩131在結構十分相似,等等等等.

但更重要的,是為甚麼這十五首詩會稱為「上行之詩」而結集在一起?

學者認為這些詩篇稱為「上行之詩」有幾個可能性,其一說法是聖殿有十五級梯級,祭司每上行一級便念誦一首；另一說法是聖殿處於耶路撒冷山丘之上,人向聖殿朝聖便是上行之旅.這兩個說法相信都是源於猶太拉比後期的解釋,指出十五首詩都是朝聖者在旅程中唱頌的詩歌,和耶路撒冷的聖殿息息相關.另一個比較切合舊約神學的說法,則是這個「上行」是回應歷代志下36:23節,在波斯帝君居魯士元年下令讓猶太人可以「上去」耶路撒冷重建聖殿的詔書:這十五首詩是反映被擄流散的以色列人,回歸耶路撒冷境況而收集的詩篇,其歷史讀者的處境也和向耶路撒冷聖殿朝聖有關,但他們正是身處帝國統治,民族流散不安的境況,雖然有部份人可以回歸故土重建,但國民身份仍是屈辱被擄之民,在帝國中只是被統治的亡國遺民.

說到這裡,不知大家是否開始明白詩篇123(和其他上行之詩)的主題,其實都是與上主的子民後被擄流散身份有關.在波斯時期,猶太人雖然能回歸重建,在新建之聖殿中朝聖朝拜,但心繫家國的猶太遺民並未忘記亡國的屈辱,在帝國下仍是被譏誚藐視的民族.

合:僕人眼望之手——等候上主作為

了解了最後段落,讓我們回到中間段落的意象喻意.要掌握詩篇123的比喻,也需要我們對古代社會文化有所認識.今天我們或許家中有傭工幫助我們家居的照顧,但古代主僕關係和今日的僱傭關係是完全不一樣的.在古時的奴隸制度中,僕人婢女地位極低,主人視他們為家庭的資本,且擁有對奴婢的生殺大權(這很小使用,不竟奴僕是家中財產).

奴婢仰望主人的手,有可能是因為受罰,等代處置,但這個解釋並不配合詩篇123的意境.更有可能,是僕人使女在等候主人主母的指示,雖然詩人可能心中有一個主人的具體行動,例如用手指的指示動作來表示賦予權力或授予恩惠,但對「仰望主人的手」意象的最可能的喻意是僕人在等待主人的某個行動(任何行動).在讀者所處的情況下,他們耐心等待上主對他們因別人的輕蔑和嘲笑而遭受的不公正苦難作出干預,任何來自上主活動的跡象都會是一個好兆頭.奴隸們可以等待幾個小時以尋求他們主人的一個跡象:拍手召喚奴隸預備食物,或揚手叫僕人收拾行裝或物品以便出行,於是僕人就起身,並配合行動.上主我們神的手的任何動作都是祂的施恩,就如主人用手暗示僕人行動的時刻已經到來.

在這段經文中,我們看到了一個深刻的教訓.就像僕人期待主人的行動一樣,我們也耐心等待神的行動.無論我們的處境如何,無論我們面對的困難有多大,我們都應該保持信心,相信神會在適當的時候採取行動.神的手的任何動作,無論多麼微小,都是一個好兆頭,因為它意味著神正在行動,正在為我們的利益工作.

詩篇123開始於個人向天舉目祈求,但這個個人所以代表的是整個神子民群體——我們——的失落和凌辱.詩篇123篇向上主呼求的憐憫,是以色列人對民族流散,國亡家恨的解救.詩篇123篇祈求上主施恩,也是我們祈望上主回應時代的被擄,在強權的帝國下重拾人性的尊嚴和自由.我們仰望上主的作為,等候上主復興的兆頭,我們又預備好配合,成就我們作主僕的使命嗎?



\newpage

\section{盤子上的頭和籃子上的餅}
\label{sec:705}
\textbf{林豪恩先生 - 盤子上的頭和籃子上的餅}
\newline
\newline
\hyperref[sec:704]{< < < PREV SERMON < < <}
~
\hyperref[sec:toc]{[返目錄]}
~
\hyperref[sec:706]{> > > NEXT SERMON > > >}


講題:盤子上的頭和籃子上的餅 The Head on the Plate and the Bread in the Basket

講員:林豪恩先生

場合:香港中文大學崇基學院禮拜堂主日崇拜

日期:2024年7月14日

引言一．盤子上的頭

剛才誦讀的福音書經課是一個西洋藝術史常出現的聖經故事.現在投影出來的是意大利畫家卡拉瓦喬(Michelangelo Merisi da Caravaggio) 晚年(1607年)的作品.畫家把自己的頭像繪成施洗者約翰的頭,讓劊子手遞交給莎樂美,於是這幅畫被稱為「莎樂美接收施洗者約翰的頭( Salome Receives the Head of Saint John the Baptist」.這就是今天講題的前半「盤子上的頭」.

在畫裡,希羅底的女兒從劊子手手上接過施洗者約翰的頭顱,眼神往畫框外望去,像是不願面對約翰一般.莎樂美身後的老婦人,則凝望著約翰的頭顱,後人推測她應為莎樂美的母親希羅底.現在,「莎樂美接收施洗者約翰的頭」這幅名畫收藏於位於倫敦的英國國家畫廊 (The National Gallery).

引言二．籃子上的餅

講題的下半部「籃子上的餅」在哪裡呢?就在「盤子上的頭」的下文:

6:30使徒聚集到耶穌那裡,將一切所做的事,所傳的道全告訴他.

6:31他就說:「你們來,同我暗暗地到曠野地方去歇一歇.」這是因為來往的人多,他們連吃飯也沒有工夫.

6:32他們就坐船,暗暗地往曠野地方去.

6:33眾人看見他們去,有許多認識他們的,就從各城步行,一同跑到那裡,比他們先趕到了.

6:34耶穌出來,見有許多的人,就憐憫他們,因為他們如同羊沒有牧人一般,於是開口教訓他們許多道理.

6:35天已經晚了,門徒進前來,說:「這是野地,天已經晚了,

6:36請叫眾人散開,他們好往四面鄉村裡去,自己買甚麼吃.」

6:37耶穌回答說:「你們給他們吃吧.」門徒說:「我們可以去買二十兩銀子的餅給他們吃嗎?」

6:38耶穌說:「你們有多少餅,可以去看看.」他們知道了,就說:「五個餅,兩條魚.」

6:39耶穌吩咐他們,叫眾人一幫一幫地坐在青草地上.

6:40眾人就一排一排地坐下,有一百一排的,有五十一排的.

6:41耶穌拿著這五個餅,兩條魚,望著天祝福,擘開餅,遞給門徒,擺在眾人面前,也把那兩條魚分給眾人.

6:42他們都吃,並且吃飽了.

6:43門徒就把碎餅碎魚收拾起來,裝滿了十二個籃子.

這個又是經常出現在西洋藝術史的聖經故事,現在投影出來的是義大利畫家喬瓦尼弗朗科(Giovanni Lanfranco)1620-1623年的作品,名為《麵包和魚的奇跡》.畫中耶穌身穿紅色長袍,頭部發光,一手指著裝滿麵包的籮筐,一手邀請眾人前來分享.無論老人孩子,還是男人女人都充滿喜悅之情,祂在群眾的簇擁下,使整個畫面呈現出一片歡騰的景象.

福音書作者把這兩個故事放在一起成為上下文,大家有沒有發現強烈的對比效果呢?藉著這對比,福音書作者想告訴讀者甚麼信息呢?這對比對我們的信仰有甚麼啟發呢?

對比一．氣氛和狀況:恐怖VS喜悅

首先,讓我們感受一下這兩個故事帶出的氣氛和狀況.「盤子上的頭」展示的狀況給你甚麼感覺呢?你可能有多種感覺,其中一種可能是殘酷和恐怖,這種感覺也符合經文的描述:

6:25(希羅底的女兒)就急忙進去見王,求他說:「我願王立時把施洗約翰的頭放在盤子裡給我.」

6:27(王)隨即差一個護衛兵,吩咐拿約翰的頭來.護衛兵就去,在監裡斬了約翰,

6:28把頭放在盤子裡,拿來給女子,女子就給她母親.

你還記得那是甚麼場合嗎?

6:21有一天,恰巧是希律的生日,希律擺設筵席,請了大臣和千夫長,並加利利作首領的.

宮廷生日宴會必然有大量美食,盤子本來應該用來盛載甚麼的呢?當然是美食.如果你是座上客,看見捧出來的盤子盛載著的不是食物,而是人頭,你會感到驚嚇,毛骨悚然嗎?那是多麼殘忍和恐怖的場景,這個殘忍和恐怖的場景出現在希律的宮廷.

對比之下,「籃子上的餅」展示的狀況又給你甚麼感覺呢?這個畫面,在曠野的五餅二魚故事出現.

6:42他們都吃,並且吃飽了.

6:43門徒就把碎餅碎魚收拾起來,裝滿了十二個籃子.

雖然只是民間的餅和魚,不是宮廷美食,但大家都有得吃,而且吃得飽,還把碎餅碎魚都收拾起來.這個畫面給你甚麼感覺呢?會不會是有福同享,有餅食餅,有魚食魚那種共享的感覺呢?會不會是有一餐「安樂茶飯」食的感覺呢?

「盤子上的頭」和「籃子上的餅」對比之下所顯示的氣氛和狀況,你寧願要哪一種呢?福音書所描述的這兩種狀況絕非罕有,在現實中經常出現.殘酷和殺戮只是在希律的宮廷出現嗎?其實在世界各地的宮廷也出現.另一方面,群眾互相分享,一齊「食餐安樂茶飯」也在世界各處可見.雖然兩者都是現實,但讀者從這個對比中充分感受到,心中渴望的是「籃子上的餅」的喜悅氣氛,而不是「盤子上的頭」的恐怖氣氛.這種對比也邀請耶穌基督的跟隨者反省,我們在自己的影響範圍,例如社會,工作間,家庭,所製造的氣氛和狀況是哪一種呢?哪一種是屬於希律宮廷的?哪一種是屬於耶穌基督的呢?

對比二．價值和觀念:視人命如草芥VS對蒼生憐憫

為甚麼會出現兩種如此強烈對比的情況呢?了解這兩種狀況背後的原因,或者能夠幫助我們遠離希律宮廷的恐怖,而能夠學效耶穌創造和平共享的喜悅.首先讓我們看看「盤子上的人頭」是如何出現的.

在希律的宮廷,沒有道德,只有慾望,連兄弟的妻子也搶過來；

6:17先是希律為他兄弟腓力的妻子希羅底的緣故,差人去拿住約翰,鎖在監裡,因為希律已經娶了那婦人.

第二,在希律的宮廷,不講對錯,只有仇恨,得罪自己的人就要除之而後快；

6:18約翰曾對希律說:「你娶你兄弟的妻子是不合理的.」

6:19於是希羅底懷恨他,想要殺他,只是不能；

第三,在希律的宮廷,輕視教養,只講手段,連自己的女兒也拖進殺人的圈套；

6:22希羅底的女兒進來跳舞,使希律和同席的人都歡喜.王就對女子說:「你隨意向我求甚麼,我必給你.」

6:23又對她起誓說:「隨你向我求甚麼,就是我國的一半,我也必給你.」

6:24她就出去對她母親說:「我可以求甚麼呢?」她母親說:「施洗約翰的頭.」

第四,在希律的宮廷,草菅人命,視人命為兒戲,宮廷的玩樂.

6:22希羅底的女兒進來跳舞,使希律和同席的人都歡喜.王就對女子說:「你隨意向我求甚麼,我必給你.」

6:23又對她起誓說:「隨你向我求甚麼,就是我國的一半,我也必給你.」

6:25她就急忙進去見王,求他說:「我願王立時把施洗約翰的頭放在盤子裡給我.」

第五,在希律的宮廷,不守正義,只顧面子,明知是義人也要殺；

6:20因為希律知道約翰是義人,是聖人,所以敬畏他,保護他,聽他講論,就多照著行,並且樂意聽他.

6:26王就甚憂愁；但因他所起的誓,又因同席的人,就不肯推辭,

6:27隨即差一個護衛兵,吩咐拿約翰的頭來.護衛兵就去,在監裡斬了約翰,

「籃子上的餅」所傳遞出來的喜悅共享又是如何出現的呢?參與的人的價值觀和希律宮廷的人有甚麼不同呢?

相對於希律宮廷視人命如草芥,耶穌自己辛苦仍然憐憫和珍惜人；

6:31他就說:「你們來,同我私下到荒野的地方去歇一歇.」這是因為來往的人多,他們連吃飯的時間也沒有.

6:32他們就坐船,私下往荒野的地方去.

6:33眾人看見他們走了,有許多認識他們的,就從各城步行,一同跑到那裡,比他們先趕到了.

6:34耶穌出來,見有一大群的人,就憐憫他們,因為他們如同羊沒有牧人一般,於是開始教導他們許多事.

第二,相對於希律宮廷那「盤子上的頭」隱喻著把人視為食物,耶穌提供食物滿足人的需要

6:35天已經晚了,門徒進前來,說:「這是野地,天已經晚了,

6:36請叫眾人散開,他們好往四面鄉村裡去,自己買甚麼吃.」

6:37耶穌回答說:「你們給他們吃吧.」門徒說:「我們可以去買二十兩銀子的餅給他們吃嗎?」

6:38耶穌說:「你們有多少餅,可以去看看.」他們知道了,就說:「五個餅,兩條魚.」

6:39耶穌吩咐他們,叫眾人一幫一幫地坐在青草地上.

6:40眾人就一排一排地坐下,有一百一排的,有五十一排的.

6:41耶穌拿著這五個餅,兩條魚,望著天祝福,擘開餅,遞給門徒,擺在眾人面前,也把那兩條魚分給眾人.

是甚麼製造出殘酷恐慌狀況呢?是甚麼製造喜悅共享的狀況呢?其實都離不開相關的人,尤其是持份者的價值觀.希律宮廷的人把其他人當作滿足自己的工具,視人命如草芥；五餅二魚事件中的人憐憫愛惜蒼生,回應和滿足他人的需要.對比清楚顯示出這兩種對人的不同態度,邀請耶穌基督的跟隨著反省,自己對人的態度所反映的是哪一種價值呢?自己追求的是哪一個國度呢?

總結:

其實,「盤子上的頭」和「籃子上的餅」兩個故事都有共同之處,就是和飲食有關.在希律宮廷的本來是美食當前的生日宴會,可惜最後變成令人作嘔的場面；在曠野的群眾本來欠缺食物,最後變成人人「有安樂茶飯食」的歡宴.希律王宮是世界權力的象徵,曠野的耶穌呢?彼拉多審問耶穌說:「這樣,你是王嗎?」耶穌回答說:「你說我是王.我為此而生,也為此來到世間」(約翰福音18:37)希律宮廷展示的王權是「吃人」的權力,而耶穌所展示的王權是為他人獻出自己的權力:

14:22耶穌拿起餅來,祝福了,就擘開,遞給他們,說:「你們拿去,這是我的身體.」

14:23他又拿起杯來,祝謝了,遞給他們；他們都喝了.

14:24耶穌對他們說:「這是我立約的血,為許多人流出來的.」

各位跟隨耶穌基督的弟兄姊妹,願上主保守我們不要誤入了希律宮廷的歧途,求上主帶領我們跟隨耶穌基督的榜樣.



\newpage

\section{這些骸骨能復活嗎}
\label{sec:706}
\textbf{鄧瑞強博士 - 這些骸骨能復活嗎}
\newline
\newline
\hyperref[sec:705]{< < < PREV SERMON < < <}
~
\hyperref[sec:toc]{[返目錄]}
~
\hyperref[sec:707]{> > > NEXT SERMON > > >}


講題:這些骸骨能復活嗎? Can These Bones Live?

講員:鄧瑞強博士

場合:香港中文大學崇基學院禮拜堂主日崇拜

日期:2024年7月21日

各位弟兄姊妹,早安.願神的恩典常與你們同在.

上兩個月,我們講到耶利米書和耶利米哀歌.在猶大國亡國前夕,先知耶利米叫人及早悔改,但人總是不聽,最後,巴比倫大軍攻入耶路撒冷,猶大亡國,人民被擄.耶利米哀歌描繪亡國時的淒涼,慘不忍睹.

今日,我們一起看以西結書.在猶大國亡國前十年,先知以西結已經被擄去巴比倫.當時,巴比倫擄走猶大國的約雅斤王,作為人質,然後立西底家作為傀儡王.先知以西結就是與約雅斤王一起被擄走的精英,他是在巴比倫聽到猶大能民的講述,才知道猶大國已亡.

先知以西結本身是祭司,很熟悉聖殿,很熟悉聖殿裡至聖所內約櫃上的「施恩座」,很熟悉獻祭禮儀,很重視聖潔.以西結書和這些事件相關.

在以西結書的開頭,先知在巴比倫地,看見能飛行移動的「施恩座」,看見神在巴比倫上空掌權.原來,耶路撒冷的聖殿雖然被毀,「施恩座」仍在,變得會飛行,能無處不在.耶和華不單沒有被巴比倫的神明擊敗,更在巴比倫上空掌權.這是猶大神學家典型的思考,在絕望之時,他們不是放棄耶和華,覺得祂無用,然後轉去敬拜別人的神明,而是擴闊他們的神觀,耶和華不只是猶大地的神,更是跟著祂的子民,去到巴比倫,在那裡坐著為王.在艱難中,猶大神學家不是放棄信仰,而是擴大信仰,好使新的信仰觀念能承托他們去面對逆境.信仰的創造力成了猶大民族堅毅地面對艱難的根據.

又是先知又是祭司的以西結,以贖罪日的贖罪禮儀去理解亡國被擄的事件.在贖罪禮儀裡,會用兩隻公山羊.有一隻公山羊被殺為祭,而另一隻公山羊,大祭司會按手在牠頭上,將人民的罪歸到這隻代罪羊身上,然後將牠放生到曠野去.在以西結眼中,在亡國事件裡,有人死去,有人被帶到外地,正是一種贖罪祭,讓土地和人民得以潔淨.

以西結書的結尾部分,先知看到歷史的終極遠象,是聖殿的重建,作為大地的中心.

可以看到,以西結書是以祭司的視角看世界局勢,看事情的核心,看歷史的遠象.以西結的一生,是在兵荒馬亂的國際情景下,身處不潔的巴比倫,卻努力完成祭司的身分,呼籲人好好思考聖潔.

今日,我們看以西結書37:1-14.

結37:1　耶和華的靈〔原文是手〕降在我身上.耶和華藉他的靈帶我出去,將我放在平原中；這平原遍滿骸骨.

結37:2　他使我從骸骨的四圍經過,誰知在平原的骸骨甚多,而且極其枯乾.

結37:3　他對我說:「人子啊,這些骸骨能復活嗎?」我說:「主耶和華啊,你是知道的.」

結37:4　他又對我說:「你向這些骸骨發預言說:枯乾的骸骨啊,要聽耶和華的話.

結37:5　主耶和華對這些骸骨如此說:『我必使氣息進入你們裡面,你們就要活了.

結37:6　我必給你們加上筋,使你們長肉,又將皮遮蔽你們,使氣息進入你們裡面,你們就要活了；你們便知道我是耶和華.』」

結37:7　於是,我遵命說預言.正說預言的時候,不料,有響聲,有地震；骨與骨互相聯絡.

結37:8　我觀看,見骸骨上有筋,也長了肉,又有皮遮蔽其上,只是還沒有氣息.

結37:9　主對我說:「人子啊,你要發預言,向風發預言,說主耶和華如此說:氣息啊,要從四方〔原文是風〕而來,吹在這些被殺的人身上,使他們活了.」

結37:10　於是我遵命說預言,氣息就進入骸骨,骸骨便活了,並且站起來,成為極大的軍隊.

結37:11　主對我說:「人子啊,這些骸骨就是以色列全家.他們說:『我們的骨頭枯乾了,我們的指望失去了,我們滅絕淨盡了.』

結37:12　所以你要發預言對他們說,主耶和華如此說:『我的民哪,我必開你們的墳墓,使你們從墳墓中出來,領你們進入以色列地.

結37:13　我的民哪,我開你們的墳墓,使你們從墳墓中出來,你們就知道我是耶和華.

結37:14　我必將我的靈放在你們裡面,你們就要活了.我將你們安置在本地,你們就知道我─耶和華如此說,也如此成就了.這是耶和華說的.』」

結37:1　耶和華的靈〔原文是手〕降在我身上.耶和華藉他的靈帶我出去,將我放在平原中；這平原遍滿骸骨.

結37:2　他使我從骸骨的四圍經過,誰知在平原的骸骨甚多,而且極其枯乾.

這是亡國後,猶大民族的寫照,像枯骨,沒有生機,沒有希望.

結37:3　他對我說:「人子啊,這些骸骨能復活嗎?」我說:「主耶和華啊,你是知道的.」

「這些骸骨能復活嗎?」在絕望深處,還有希望嗎?

先知說:「主耶和華啊,你是知道的.」希望,只建基於神.

在真正的絕望之處,像死亡這存在的虛無之境,人是無能為力的.以西結沒有一開口便回答「有希望」,他反對輕言希望,反對「無根據的希望」,反對阿Ｑ精神.同時,以西結也不輕言絕望,他知道神乃萬有的存在根源,他反對「不信神的絕望」.絕處是否有生機,端賴神的恩典.

結37:4　他又對我說:「你向這些骸骨發預言說:枯乾的骸骨啊,要聽耶和華的話.

結37:5　主耶和華對這些骸骨如此說:『我必使氣息進入你們裡面,你們就要活了.

結37:6　我必給你們加上筋,使你們長肉,又將皮遮蔽你們,使氣息進入你們裡面,你們就要活了；你們便知道我是耶和華.』」

神的說話帶著權能,向枯骨說,活過來,他們就要活過來.個體或社群,在歷史的流逝中,都會遭逢絕境.是神的說話,神的真理,神的作為,帶來全新的創造,這使人知道我們自己不是歷史的主人,暴力的力量也不是歷史的主人,耶和華才是萬物的主人.

結37:7　於是,我遵命說預言.正說預言的時候,不料,有響聲,有地震；骨與骨互相聯絡.

結37:8　我觀看,見骸骨上有筋,也長了肉,又有皮遮蔽其上,只是還沒有氣息.

神如何吩咐以西結,以西結便照樣做.他向枯骨講說神的話,枯骨便長了皮肉,只欠氣息.

結37:9　主對我說:「人子啊,你要發預言,向風發預言,說主耶和華如此說:氣息啊,要從四方〔原文是風〕而來,吹在這些被殺的人身上,使他們活了.」

結37:10　於是我遵命說預言,氣息就進入骸骨,骸骨便活了,並且站起來,成為極大的軍隊.

經文明言,這些枯骨是「被殺者」的枯骨,他們是巴比倫戰爭下的亡魂.神對以西結說,吩咐「四風」,吹在這些被殺的人身上,使他們復活.為什麼提到「四風」,這是什麼來的?學者指出,要理解這些經文,需要參照巴比倫最著名的史詩《埃努瑪•埃利什》(Enuma Elish).這些人敗在巴比倫手下,按當時的宗教文化來說,他們是被巴比倫的大神馬杜克(Marduk)擊敗的.在神話裡,馬杜克用什麼收服敵人的呢,就是用「四風」.馬杜克以「四風」作網,將敵人圍困.在中國的《西遊記》裡,佛祖是用「五指山」將孫悟空困住.各處神明,困住敵人,有不同板斧.

現在,神叫以西結呼喚「四風」,將「四風」變成生氣,讓被殺者復活.這是神學上最複雜的一種敘事.將巴比倫大神克服猶大人的手段,變成是耶和華手上的工具,就像將「西毒歐陽鋒」用來殺人的「蛤蟆功」,變成救助被「蛤蟆功」所殺的人的方法.

以西結書這樣說,其實就是廢了巴比倫大神的武功,他那用來殺敗猶大人的「四風」,在耶和華手中,成了祝福猶大人的工具.如此,沒有力量可與耶和華的創造大能匹敵.「四風」失去了克制猶大人的魔力,反成了讓猶大人復活的工具.馬杜克(Marduk)失去了主場,巴比倫現在成了耶和華話事的地方.耶和華的武功真高強,敵人最大的殺著,卻成就了自己的救贖行動.

以西結宣告神的話,巴比倫大神手中的武器「四風」便聽命,成了生命氣息,讓枯骨復活,且成了「極大的軍隊」.這軍隊,在世上,為神的聖潔與神聖而奮鬥.

結37:11　主對我說:「人子啊,這些骸骨就是以色列全家.他們說:『我們的骨頭枯乾了,我們的指望失去了,我們滅絕淨盡了.』

結37:12　所以你要發預言對他們說,主耶和華如此說:『我的民哪,我必開你們的墳墓,使你們從墳墓中出來,領你們進入以色列地.

結37:13　我的民哪,我開你們的墳墓,使你們從墳墓中出來,你們就知道我是耶和華.

結37:14　我必將我的靈放在你們裡面,你們就要活了.我將你們安置在本地,你們就知道我─耶和華如此說,也如此成就了.這是耶和華說的.』」

這幾節經文,是這個故事的神學結語了.

真正的希望,不是根據看到的景象來判斷的.眼睛看到的,只是事物本身內含的所有可能性.若看到事物本身內藏的生機,那當然能鼓舞人積極前進.但事實上,很多時,事物本身已沒有生機,這時,如何積極也沒有什麼用的.就像那些枯骨說,我們沒有希望了.事實上,它們真的沒有希望的.死亡,本身並不內藏復活的可能性.猶大這樣的小國,亡在超級霸主巴比倫手下,也看不到什麼生機.

信仰卻不是憑眼見的.信仰就是不將生命交給現實,而是交給上帝.基督教的上帝,最終顯明在耶穌的生命裡,他死後復活了,一方面顯明死亡與絕望的真實,另方面,卻顯明神有創造生命的大能,他不被死亡困住,卻利用了死亡這絕境,叫人認識到死亡終極而言沒有話事權,神有克服死亡的大能.正如我們的信經說:「我信身體復活,我信永生」.

真正的希望,不是看看現實能提供什麼出路,而是相信神,相信祂的說話,相信祂對我們的愛,在神裡,找到希望的根據.真正的希望,不是看到絕境本身有出路,而是看到絕境真是絕境,但它並不決定我們的命運,唯有上帝才能決定我們的命運.當猶大人亡國,成了別人的奴隸,變得一無所有時,唯有這種不看現實,只看上帝的非一般希望,支撐他們度過難關.願你們仰望神,同樣找到這種希望.

願上主得著榮耀,頌讚；願眾人得享喜樂,平安.



\newpage

\section{你看見什麼}
\label{sec:707}
\textbf{高國雄牧師 - 你看見什麼}
\newline
\newline
\hyperref[sec:706]{< < < PREV SERMON < < <}
~
\hyperref[sec:toc]{[返目錄]}
~
\hyperref[sec:708]{> > > NEXT SERMON > > >}


講題:你看見什麼? What Do You See?

講員:高國雄牧師

場合:香港中文大學崇基學院禮拜堂主日崇拜

日期:2024年7月28日

1這事以後,耶穌渡過加利利海,就是提比哩亞海.

2有許多人因為看見他在病人身上所行的神蹟,就跟隨他.

3耶穌上了山,和門徒一同坐在那裡.

4那時猶太人的逾越節近了.

5耶穌舉目看見許多人來,就對腓力說:「我們從哪裡買餅叫這些人吃呢?」

( 6他說這話是要試驗腓力；他自己原知道要怎樣行.)

7腓力回答說:「就是二十兩銀子的餅,叫他們各人吃一點也是不夠的.」

8有一個門徒,就是西門‧彼得的兄弟安得烈,對耶穌說:

9「在這裡有一個孩童,帶著五個大麥餅,兩條魚,只是分給這許多人還算甚麼呢?」

10耶穌說:「你們叫眾人坐下.」原來那地方的草多,眾人就坐下,數目約有五千.

11耶穌拿起餅來,祝謝了,就分給那坐著的人；分魚也是這樣,都隨著他們所要的.

12他們吃飽了,耶穌對門徒說:「把剩下的零碎收拾起來,免得有糟蹋的.」

13他們便將那五個大麥餅的零碎,就是眾人吃了剩下的,收拾起來,裝滿了十二個籃子.

14眾人看見耶穌所行的神蹟,就說:「這真是那要到世間來的先知!」

15耶穌既知道眾人要來強逼他作王,就獨自又退到山上去了.

16到了晚上,他的門徒下海邊去,

17上了船,要過海往迦百農去.天已經黑了,耶穌還沒有來到他們那裡.

18忽然狂風大作,海就翻騰起來.

19門徒搖櫓,約行了十里多路,看見耶穌在海面上走,漸漸近了船,他們就害怕.

20耶穌對他們說:「是我,不要怕!」

21門徒就喜歡接他上船,船立時到了他們所要去的地方.

今天我們又要再次經歷五餅二魚的神蹟,你是否興奮雀躍,充滿著期待呢?

相信我們當中不少人對此神蹟十分熟悉,前主日林先生也有分享這神蹟(盤子上的頭和籃子上的餅 ),我在這講壇亦曾經與大家分享過,到底仍然有什麼值得再次分享呢?

首先,今天的經課是約翰福音,約翰福音是一本很獨特的福音書,與其餘3卷福音書有很大的分別,馬太,馬可及路加福音被歸類為符類福音/共觀福音(Synoptic Gospels),她們有近似的記述次序及編排結構,一般認為她們有著相當的關聯,我們亦習慣在合適時把她們相互參考.

但約翰福音不單是在福音書中最後成書的(~90AD),更很清楚地自我表達:

約20:30-31耶穌在門徒面前另外行了許多神蹟,沒有記在這書上.但記這些事要叫你們信耶穌是基督,是上帝的兒子,並且叫你們信了他,就可以因他的名得生命.

約翰就是要為一個獨特的目標而選取某些有關耶穌基督的事蹟,甚至可以理解為用另一角度去補充當時流行福音書未有展述的觀點；在細心閱讀約翰福音後,我們不難發現其獨特之處,由「太初有道」取代「耶穌誕降」,到「為門徒洗腳」取代「最後晚餐」,足以看出其選材的獨特及意義.

當我們集中認識今天的經課,五餅二魚的神蹟,請先留意,約翰福音中記載了耶穌行了7個(或8個)神蹟,分別是:

1.約2:1-11:迦拿婚宴中把水變成酒(第一個神蹟)

2.約4:46-54:醫治大臣之子

3.約5:1-18:醫治畢士大池邊38年的癱子

4.約6:6-14:五餅二魚吃飽五千人(太14:13-21,可6:30-44,路9:10-17)

5.約6:16-21:耶穌在海面上行走(太14:22-33,可6:45-52)

6.約9:醫好生來瞎眼的

7.約11:叫拉撒路復活

*8.約21:復活的主向西門彼得顯現,啟示他打了153條魚,網沒有破

請留意,除了今天的經課中的2個接續相連的神蹟,其他的全部都是約翰福音的獨特/唯一記載,而五餅二魚這神蹟,更是四福音中唯一都共同有記載的神蹟,可以理解為約翰是刻意,用心地選擇把這神蹟再重述,相信值得我們細心留意約翰想向我們所展示的…

今天想與大家思考,我們所「看到」的是什麼,又因此有什麼反應呢?

當一打開五餅二魚的經課,我們看到什麼:熟悉?沒有新意?耶穌很厲害?…反應:把耳關上?不再期待?

門徒的眼目,所睇見的

V7腓力回答說:「就是二十兩銀子的餅,叫他們各人吃一點也是不夠的.」

腓力看到與反應:二十兩銀子(200日的人工)資源不足(現實),限制,太難(人太多,就算有…也不能)…不作任何行動,不願意/覺得需要作任何行動.

V8有一個門徒,就是西門‧彼得的兄弟安得烈,對耶穌說: 9「在這裡有一個孩童,帶著五個大麥餅,兩條魚,只是分給這許多人還算甚麼呢?」

安得烈看到與反應:部份資源,五個大麥做的餅,跨大其不足,資源沒有用處…但也願意獻上.想起這表達:「就是西門‧彼得的兄弟安得烈」,安得烈是一個怎樣的人?沒有領導的身份?低調專心?不要忘記,按約翰福音,安得烈是跟從耶穌的第一個門徒,西門是他介紹給耶穌的,這小孩也是…

眾人呢?他們又看到什麼?

V2有許多人因為看見他在病人身上所行的神蹟,就跟隨他.

眾人看到與反應:耶穌所行的神蹟,跟隨.

V14眾人看見耶穌所行的神蹟,就說:「這真是那要到世間來的先知!」 15耶穌既知道眾人要來強逼他作王,就獨自又退到山上去了.

眾人看到與反應:耶穌所行的神蹟,認信…強逼耶穌做不是他願意/計劃的事.

耶穌呢?

耶穌的眼目一直看到什麼呢?

V5耶穌舉目看見許多人來,就對腓力說:「我們從哪裡買餅叫這些人吃呢?」

看到:人的需要,要馬上回應這需要…

腓力回應了資源不足,充滿限制…

但其實耶穌所期望他回應的是什麼呢?

V5耶穌舉目看見許多人來,就對腓力說:「我們從哪裡買餅叫這些人吃呢?」V6他說這話是要試驗腓力；他自己原知道要怎樣行.

耶穌清楚知道如何處理資源的問題,而是希望邀請腓力參與在這次「分享」當中.

看到:腓力可以幫助(他是當地人,熟悉環境)

耶穌不是需要門徒幫助,上帝清楚要做什麼,只是我們不知道…耶穌想門徒有份參與,想腓力有所成長,邀請腓力一同經歷這神蹟.

V10耶穌說:「你們叫眾人坐下.」

看到:安得烈的投入,是可以有意義回報的,也邀請他一同經歷這神蹟.

V12他們吃飽了,耶穌對門徒說:「把剩下的零碎收拾起來,免得有糟蹋的.」 13他們便將那五個大麥餅的零碎,就是眾人吃了剩下的,收拾起來,裝滿了十二個籃子.

耶穌讓眾人一同看到:人能體會超乎想像的豐富,更要珍惜主的恩典.

故事還未完結…

V16到了晚上,他的門徒下海邊去, 17上了船,要過海往迦百農去.天已經黑了,耶穌還沒有來到他們那裡. 18忽然狂風大作,海就翻騰起來. 19門徒搖櫓,約行了十里多路,看見耶穌在海面上走,漸漸近了船,他們就害怕.

耶穌看到在狂風中的門徒需要祂,就以最快的方式去拯救他們(心靈上的拯救).

門徒看到與反應:狂風大作,以為風浪中的耶穌是鬼魅…害怕

V20耶穌對他們說:「是我,不要怕!」 21門徒就喜歡接他上船,船立時到了他們所要去的地方.

門徒看到與反應:看到是耶穌…喜歡接納.

五餅二魚神蹟後,眾人看到什麼呢?

6:22第二日,站在海那邊的眾人知道那裡沒有別的船,只有一隻小船,又知道耶穌沒有同他的門徒上船,乃是門徒自己去的. 23然而,有幾隻小船從提比哩亞來,靠近主祝謝後分餅給人吃的地方. 24眾人見耶穌和門徒都不在那裡,就上了船,往迦百農去找耶穌. 25既在海那邊找著了,就對他說:「拉比,是幾時到這裡來的?」 26耶穌回答說:「我實實在在地告訴你們,你們找我,並不是因見了神蹟,乃是因吃餅得飽.

30他們又說:「你行甚麼神蹟,叫我們看見就信你；你到底做甚麼事呢? 31我們的祖宗在曠野吃過嗎哪,如經上寫著說:『他從天上賜下糧來給他們吃.』」 32耶穌說:「我實實在在地告訴你們,那從天上來的糧不是摩西賜給你們的,乃是我父將天上來的真糧賜給你們. 33因為上帝的糧就是那從天上降下來,賜生命給世界的.」

34他們說:「主啊,常將這糧賜給我們!」 35耶穌說:「我就是生命的糧.到我這裡來的,必定不餓；信我的,永遠不渴. 36只是我對你們說過,你們已經看見我,還是不信.

41 猶太人因為耶穌說「我是從天上降下來的糧」,就私下議論他, 42說:「這不是約瑟的兒子耶穌嗎?他的父母我們豈不認得嗎?他如今怎麼說『我是從天上降下來的』呢?」

52因此,猶太人彼此爭論說:「這個人怎能把他的肉給我們吃呢?」

60他的門徒中有好些人聽見了,就說:「這話甚難,誰能聽呢?」 61耶穌心裡知道門徒為這話議論,就對他們說:「這話是叫你們厭棄嗎?

66從此,他門徒中多有退去的,不再和他同行. 67耶穌就對那十二個門徒說:「你們也要去嗎?」 68西門‧彼得回答說:「主啊,你有永生之道,我們還歸從誰呢? 69我們已經信了,又知道你是上帝的聖者.」

一個神蹟,吸引成千上萬的群眾.

一篇信息,最後只餘下12門徒.

今天我們看到什麼?

人的需要?上帝的大能?上帝的計劃?

看到又為我們帶來什麼行動?

主動付出?喜歡接納?

一個看似沒有必須的神蹟,耶穌就是因為看到當下的需要,就回應…

我們能夠戴上這對上帝的眼鏡嗎?



\newpage

\section{心靈和誠實}
\label{sec:708}
\textbf{陳永業傳道 - 心靈和誠實}
\newline
\newline
\hyperref[sec:707]{< < < PREV SERMON < < <}
~
\hyperref[sec:toc]{[返目錄]}
~
\hyperref[sec:709]{> > > NEXT SERMON > > >}


約翰福音 4:19-26
\newline
\begin{longtable}{cl}
\hline
\hline
章節 & 經文 (和合本修訂版)\\
\hline
4:19 & \begin{tabularx}{0.7\textwidth}{X} 婦人對他說：「先生，我看你是一位先知。 \end{tabularx} \\ \\ \relax
4:20 & \begin{tabularx}{0.7\textwidth}{X} 我們的祖宗在這山上敬拜神，你們倒說，應當敬拜的地方是在耶路撒冷。」 \end{tabularx} \\ \\ \relax
4:21 & \begin{tabularx}{0.7\textwidth}{X} 耶穌對她說：「婦人，你要信我。時候將到，你們敬拜父，既不在這山上，也不在耶路撒冷。 \end{tabularx} \\ \\ \relax
4:22 & \begin{tabularx}{0.7\textwidth}{X} 你們所敬拜的，你們不知道；我們所敬拜的，我們知道，因為救恩是從猶太人出來的。 \end{tabularx} \\ \\ \relax
4:23 & \begin{tabularx}{0.7\textwidth}{X} 時候將到，現在就是了，那真正敬拜父的，要用心靈和誠實敬拜他，因為父要這樣的人敬拜他。 \end{tabularx} \\ \\ \relax
4:24 & \begin{tabularx}{0.7\textwidth}{X} 神是靈，所以敬拜他的必須用心靈和誠實敬拜他。」 \end{tabularx} \\ \\ \relax
4:25 & \begin{tabularx}{0.7\textwidth}{X} 婦人對他說：「我知道彌賽亞—就是那稱為基督的—要來；他來了，會把一切的事都告訴我們。」 \end{tabularx} \\ \\ \relax
4:26 & \begin{tabularx}{0.7\textwidth}{X} 耶穌對她說：「我就是，正在跟你說話呢！」 \end{tabularx} \\ \\
[1ex]
\hline
\hline
\end{longtable}
講題:心靈和誠實 In Spirit and In Truth

經文:約翰福音4:19-26

講員:陳永業傳道

場合:香港中文大學崇基學院禮拜堂主日崇拜

日期:2024年8月4日

今天跟大家分享這段經文是一個井旁婦人遇見耶穌的事蹟.一個井旁的婦人在街上打水,那時候是二千多年前的耶穌時代,她在井旁打水遇到一個人問她:你打水可不可以給我一杯水來喝?大家覺得是不是一件很平常的事?其實很平常的,你在中大校園坐下,突然有路人問你:崇基禮拜堂在哪裡?你都很容易告訴他,大家只是聊聊天,之後再說多兩句,聊天很平常,但其實我告訴你,這個相遇這件事在耶穌時代是極不尋常的事,為什麼?

有幾個不尋常的,第一個不尋常就是一個男人和一個女人聊天,現今社會當然沒問題了,我們現在坐在一起,男女坐在一起沒有問題的,可以握握手,但當時男女授受不親,就好像我們中國人的背景,我們早幾代都已經很不同了,即是男女有別,是不會溝通什麼,而且當時男性的地位很高,而女性很卑微的,雖然一個男人是不會跟一個女人說話的.第二個不尋常就是一個撒瑪利亞婦人和一個猶太人聊天,因為他們是世仇,弟兄姊妹有沒有不喜歡的人?可能多多少少都有,但這個不單止是普通交往那種不喜歡,而是一種好根深蒂固的仇恨,是歧視,為什麼呢?

在耶穌時代,即是說二千幾年前,耶穌出世之前七百多年前,以色列已經分開了南北兩國,再早一些,我們知道以色列在聖經裡面說第一個皇帝叫掃羅,第二個是牧羊的大衛,第三個他的兒子很有智慧的所羅門.之後就沒有了,就分開南北兩國,南國就是承傳著所羅門血統的嫡系皇朝叫猶大國,上面叫以色列國,南北兩國開始分開,去到耶穌出世七百多年前也經歷很多代了,不過就給在上面有一個很強勁的帝國叫亞述帝國去殲滅了,殲滅了就沒有了自己的地方,沒有自己的家園,然後就受人統治,被人流放,演變到他們猶太人的血統慢慢被同化,被侵入了.就等於我們今天有很多人去移居海外,然後他們結婚生子,娶了一個外國人為妻,或者嫁了一個外國人,然後生了混血小孩.今天我們都挺喜歡混血兒,男就都挺帥的,女就很漂亮的,輪廓好深,可以做模特兒.第一,當時猶太人很不喜歡自己的血統不純正,視為不潔,他們覺得自己的身份是上帝的選民,絕對不可以侵入外邦人的血統.

於是他們世世代代就跟撒瑪利亞人有一個很大的洪溝,看不起他們,覺得他們是異族,所以不屑跟他們交往,另一方面,撒瑪利亞人有他們自己居住的地方,有他們自己敬拜的地方,因為他們都是有猶太人血統,都是敬拜耶和華上帝,但是他們敬拜的地方就不會去耶路撒冷,因為他們知道猶太人不喜歡他們,於是他們找了另一個山叫基利心山,因為他們說這座山更有歷史,更有傳統,比起錫安山,比起耶路撒冷的聖殿更有歷史價值,因為摩西曾經在這裡祝福我們十二支派,包括我們,所以他們就去基利心山敬拜,猶太人就到耶路撒冷,大家河水不犯井水.

還有一個不尋常,就是時間.不尋常的地方就是在日頭太陽正午十二點的時間,這個婦人走來打水是很不尋常的,因為是自找麻煩的.香港現在三十幾度都融化了,但是當時如果你有機會去到聖地的中東地方,是四十幾度的高溫,所以沒有人在這個時間打水的.今日我們打開水喉就有水不覺得什麼叫辛苦,但大家如果大家試過制水,托著的水桶去樓下拿水就很辛苦了.當時的人要在城外的井打水,打完之後又拿回去,一來一回是很辛苦的.可能你從家中去到城外,可能半個小時路程,要托著水桶,還要不止一次,可能一天去幾次,所以沒有人在這個時間打水,通常是一早清晨太陽未去到正午的時候,或者是日落西山,下午三,四點才去,沒有人在這個時間去,但這個女人為什麼要在這個時間去?她自找麻煩嗎?想做健身?想出身汗?她其實是不想見人,大家如果知道這個故事的發展,或者知道這個撒瑪利亞人的背景,你會知道其實她是一個很不想見人,那個生命很複雜的,她的家庭背景很複雜,不想別人見到她的時候,怕大家指指點點或者帶著一些不好的眼光,所以她情願辛苦一點都不想見人,所以很不尋常的.

剛才經文那裡都有提到,在第四節那裡,經文說:耶穌離開了猶大去到加利利,從南面上去到北面,必須要經過撒瑪利亞,但其實很多釋經研究的學者說其實不是必須經過撒瑪利亞,而是耶穌必須經過的,其他人其實就不會這樣走,因為由南去北就經過撒瑪利亞,剛才我說他們是不想來往的,所以很多猶太人是寧願「兜路」,有其他路線是可以上加利利,是不需要經過撒瑪利亞,但是為什麼耶穌說要必須去撒瑪利亞?我們看回事情的發展,其實耶穌是故意去找她,去找這個撒瑪利亞婦人,所以你試想一下,其實那些門徒跟著耶穌去的,那些門徒跟著耶穌,耶穌說:我們走這邊,就是經過撒瑪利亞的,如果你是其中一個門徒,你又是猶太人,你就心裡面有很多疑惑,耶穌為什麼要去這裡?我們通常走這邊,為什麼你要走那邊?但他們不敢問,因為耶穌是他們的老師,他說去就去就跟著,所以他們心裡面非常多疑惑,甚至乎心裡面是好忐忑,心裡面說:「有冇搞錯?」

但耶穌定意要經過那裡,定意要去到這個井旁,定意要去找這個撒瑪利亞婦人,耶穌親自尋找她,過往在聖經裡面都有耶穌親自去找一些人,好像格拉森被鬼附那個人,都是外邦人,所以耶穌親自去找這個撒瑪利亞婦人是有目的,有衪心意的,耶穌去到這個撒瑪利亞婦人那裡就問他:「你可不可以給水我喝?」然後那個撒瑪利亞婦人說:「你是猶太人,為什麼問我拿水喝?」所以她覺得耶穌很奇怪,我們向沒有來往,然後耶穌就說:「如果你知道上帝的恩賜,跟你說給我喝水那個是誰,你會立刻給他,你會求他……」然後這個女人就更加疑惑,更加出奇,耶穌你說你有水喝,你有活水,你又何必問我拿水喝?你想怎樣?是不是騙案?這個井又深,你又沒有打水的器具,在哪裡找到你所講的活水?然後接住問耶穌:「這井是我們祖宗雅各留給我的,他是十二支派的老祖宗,難道你比他大?耶穌說︰「凡喝這個井的水,還要再喝,但是我給你的水是永遠不渴,這個水是直湧到在你裡面,成為泉源,直湧到永生,就是使你能永遠不渴.」

撒瑪利亞婦人聽完這句話之後,她轉念了,她由充滿疑惑突然間變成趨之若鶩,彷彿有人說中了她心裡面的渴望︰「竟然你有永遠不渴的水,可不可以給我?」這個婦人在想什麼?因為如果有人可以給她永遠不渴的水,可以不用再口乾,不用再渴,即是說她可以不用再來打水,不用再見人,多好的事情!她正正是因為不想見人,所以十二點去打水,如果她說有水可以永遠不渴的,可以要就好了,立刻有一種好渴望去求耶穌︰「請你將那水給我喝.」

耶穌怎樣回應?耶穌說:「好的,我給你,但是你叫你丈夫來喝」,她呆了和尷尬地回答︰「我沒有丈夫.」於是耶穌再追問,再進深的問:「是的,你說得對,你是沒有丈夫,你以前有過五個丈夫,現在跟你在一起的那個也不是你丈夫.」這個女人突然間呆了,為什麼祂會知道我身世?為什麼祂知道我所有事情?難道我的身世一直傳到耶路撒冷,傳到所有人都知道?祂怎麼會「起我底」?

今天說「起底」很容易,我們上網輸入一個名字,你會知道很多資料,輸入一個關鍵詞就很容易找到,一點都不出奇的.但那個時候好奇怪,這個女人很詫異,為什麼耶穌和她素未謀面,他竟然知道我一切,故事再發展下去,就是剛才我們讀的經文:「先生,我看得出來你是一個先知!」當她知道站在她面前的,是一個高人的時候,就立刻就問一個她打從心底裡面一直都在想,一直在問的問題,就是你們猶太人就說在耶路撒冷,錫安山那裡敬拜神,我們的祖先就說我們在基利心山這裡敬拜神,究竟哪裡才是真正敬拜神的地方?

她在問好像一個地點的問題,但其實不是的.你試想想,如果今天你完了崇拜之後走到街上,然後見到一個你看他就知道智慧老人的人,有白頭髮,有些鬍鬚穿著長衣服,不是傻的,是一個智慧老人,你會問他什麼問題?很難得平時在的街坊見不到,你會問他什麼問題?你在想你心底最想問的問題是什麼?「高人,想問那個樓市跌價就現在這樣嗎?谷底了未?可以入市嗎?幾時回升?」或者問︰「高人,我旁邊那個是否我結婚的對象?如果是就不要拖了……」什麼問題你會問?那婦人很難得見到一個先知,見到一個高人,她竟然問一個在哪裡敬拜神的問題,有沒有那麼敬虔?想深一層,其實她在問的,不是一個地點的問題.

你想一想,其實這個婦人一而再再而三的有不同的婚姻的關係.我以前小時候讀主日學的時候,讀到這個井旁婦人的時候,我會覺得這個井旁婦人都不是很正經的人,甚至乎可能是一個妓女,或者是一些很不受歡迎的人,當我們再細心的看,人大了經歷多了,我們又知道其實一個人有她的背景有很多原因,即是我們不可以從我們的主觀角度,一看她就是這樣,有一些很直接的批評或者評論.當我們繼續去從聖經時代及文化去看的時候,其實我們會知道其實這個女人與其你說她是一個很不正經的女人,不如你說她其實是一個很破碎的人,因為在當時的社會制度裡,男性的地位很高,女性是需要男性的依靠,以至她在社會才有地位,即是那個女人其實是很被動的,而當時休妻這件事是很平常的,你試想一個父系的社會,任何原因都可以休妻的,即是說你今天煮的飯sung 不好吃,這樣可以休妻了；你的臉多了些皺紋,可以休妻了；你為甚麼頂撞我?可以休妻了!這個女人一次又一次可能被丈夫拋棄休妻,又或者可能當時這些男人很容易會病死的,當一個丈夫死了,因為在當時的社會制度裡面,要有身份就必須要再嫁人,你試想,無論甚麼原因都好,她嫁完一次又一次,嫁了五次,她的生命很破碎,她在問:我們在哪裡敬拜神?其實她在問哪裡可以找到真正的依靠?人靠不住了,我想靠神,但我們又說這座山,你又說另一座山,究竟神在哪裡?究竟真正神同在的地方在哪裡?其實她在問這個問題.

我們今天有沒有問這個問題?你真正的依靠在哪裡?我們心裡面最深的需要是甚麼?我們回來教會都是很想尋找一個生命的倚傍,這個女人是一個破碎的生命,很想見到高人就問這個很重要的問題,但耶穌怎麼回答她呢?耶穌在21節說︰「婦人,你要相信我,時候將到,你們敬拜神的不在基利心山,也不在耶路撒冷」在哪裡呢?在這裡,是我!(耶穌) 你要來找我,我才是你所找的依附對象,是真正敬拜的對象.耶穌說你們所拜的你們不知道,我們所敬拜的我們知道,因為救恩是在猶太人那裡出,其實他在說自己.

23節說:「時候將到,現在就是了,那些真正拜父的要用心靈和誠實去敬拜他,因為父是要這樣的人去敬拜他.」時候將到其實這裡出現了兩次,在23節那裡說過,還有剛才再之前21節那裡有提過,這兩次的經文都說到時候將到,其實聖經學者再看清楚,其實這個時候這個字在整卷的約翰福音,不是在說一般時間,而是在說耶穌基督來到這個世界,彌賽亞來臨那個的時間,耶穌將要被高舉,將要被釘在十字架上,復活那個時間將到,所以耶穌知道婦人需要的對象,需要依靠的,就是他自己本人.耶穌就告訴這個婦人知道,是的,真正要去敬拜父的人,真正依靠的對象,你要用心靈和誠實去敬拜衪.

我們用心靈誠實這個字我們都常常都說的,都是說我們要崇拜主的時候,我們要用心靈和誠實,心靈就是我們要用心,清心,盡力去敬拜神,誠實就更加要是一個切切實實,我們剛才提到我們祈禱最後說「阿們」或「誠心所願」,意思都是我們要真誠實意的去敬拜主.

但其實耶穌在這裡說的,是有更深的一個層次.耶穌說你要用心靈誠實去敬拜他,其實這個字本身,「心靈」這個字是沒有「心」字的,在原文裡面只有「靈」這個字,而在約翰福音裡面多數出現這個「靈」字,是指「聖靈」,所以這裡除了是指人的心靈之外,其實耶穌想說的是「聖靈」.另一方面就是「誠實」,誠實這個字更加可以解釋為「真理」,真理的意思,究竟什麼是真理呢?今天這個世界還有沒有真理的?有的,就是在耶穌那裡,耶穌說:「衪就是道路,真理,生命.」原文及英文譯本之前有一個冠詞「the」字,意思是唯一的,是那條道路,唯一通向永生的道路,衪就是真理,所以這裡說是按著真理的意思,其實就是在說耶穌基督衪自己,所以整句說話的意思,其實用心靈誠實去敬拜就是說當我們要敬拜父的時候,我們要在聖靈的裡面和耶穌基督的裡面去敬拜天父,這裡就出現了聖父聖子聖靈三一的關係的裡面,就是我們要去敬拜神,我們要透過聖靈和耶穌基督,在耶穌基督和聖靈的裡面去敬拜天父.

那什麼是在耶穌基督和聖靈裡面?這裡說一個很重要的道理就是聖靈的工作,聖靈其中一個很重要的工作就是叫我們明白真理,明白耶穌基督的心腸,明白耶穌基督的救恩,我們回看其實聖經裡說到有很多的使徒跟從耶穌的人,就像是彼得,雅各,約翰,他們跟從耶穌三年多了,但是有時他們對耶穌所作的比是一知半解,不是很知耶穌的心腸,就好像就這次事件那這樣,為什麼耶穌要去撒馬利亞?為什麼要找這個婦人?他們充滿疑惑,他們不知道也不明白,耶穌做了五餅二魚,然後再用同一方式去到外邦人的面前,再將七個餅幾條魚分給四千人,但他們還是不明白.

當耶穌復活升天時,門徒還在問耶穌得國的時候,就是現在嗎?你復興以色列國的時候是不是現在?他們還在想他們心裡面的問題,他們不明白耶穌的心腸,直至什麼時候呢?直至到使徒行傳聖靈降臨之後,他們明白了,他們恍然大悟將舊約和新約的道理貫穿,把耶穌基督就是彌賽亞這個信息到處跟人宣講,帶著能力的,帶著信心權柄的,和五千人信主,三千人信主,即使是被捆鎖,被人打,被人恐嚇,他們仍然決心的要去傳揚耶穌基督的救恩.他們可以大膽說︰「順從神不順從人是應該的!」然後被人打了,趕出會堂.他們仍然有這種勇氣和信心,就是聖靈的工作和能力.

弟兄姊妹,今天我們帶著這個福音在我們生命裡面,我們有沒有這種信心權柄勇氣,在聖靈裡面,並且我們知道我們所信的核心就是耶穌基督衪自己,就好像這個婦人見到彌賽亞,見到耶穌那一日,她知道衪是又真又活的神,就是我們一直等著的彌賽亞,就是她一直尋找的倚靠.

在耶穌基督和聖靈裡面敬拜天父的意思就是這份關係了.今天我們敬拜神,我們回來聚集是很好的,但是回來教會的意思和目的,你真的是知道嗎?是不是就是來大家互相認識,聯誼一下,然後去喝茶?不是的,而是我們一起聚集,一起敬拜,是一起和這個三一的神維繫那種關係,大家彼此的守望.所以這份關係對上帝那種敬畏,我們真是怎樣去尊主為大,讓耶穌作王,在這個心裡面,是這個敬拜的生命先是我們崇拜裡面最核心要說的,我們要以基督為中心,去屈膝俯伏這位信實不變的上帝.敬拜就是一個俯伏,降服的意思,降服的意思不單是心態的,不單是動作裡面的俯伏,而是我們從頭到尾承認耶穌為主,衪是主,我們是衪的兒女,我們聽衪話,降服在衪那裡.

這是一個生命的態度,不只是在教會裡面,是教會以外.在我們的社群當中,在我們的家庭裡面,在我們的職場,在我們人與人之間的關係裡面,我們能不能夠活出一個基督的信仰,一種喜樂雀躍的人生,我們信耶穌的人很棒的,能笑得出也睡得著,常常可以面帶笑容,笑看這個風雲.是可以的,我們不是「阿Q精神」來的,而是我們真實的知道我們所敬拜的是誰,而是我們真實知道誰看著我們.

所以弟兄姊妹,今天我們跟弟兄姊妹一齊來敬拜神,我們都一起找到我們真正的依靠,但我很想告訴你知,你找到了,但你為什麼這麼多時候都這麼多煩慮呢?我們還有很多不安或者不開心,我們可以試想一下我們有沒有真正的交託和依靠.所以弟兄姊妹,我們一起渴望上帝常常在我們一起,跟我們一起去經歷衪,就像這個撒馬利亞婦人一樣,有一顆很雀躍的心,我們有耶穌在我們心裡面,我們就可以周圍跟人去說的,是很興奮的事情.

我們嘗試撫心自問,第一,我有沒有這個很雀躍的心?第二,我的心是否真的可以這樣信靠耶穌,可以釋懷,可以完全的信賴衪,完全的交託給衪?這個是我們真正得著耶穌,得著這個敬拜的生命的一個很大的一個回報和喜樂,我們同心一齊祈禱︰

在這個時間,我也很想大家一齊安靜地去默禱,我們問問自己,我們是否真的找到我們生命的依靠?信仰對我們的生命有沒有這種雀躍的影響力? 今天耶穌就透過這段經文告訴你,衪就是你真正所依靠的對象,你找對了.所以祈求主給我們可以能明白你,識得向你禱告,我們可以全然的將我們生命一切的事情都交託在這位信得過的主耶穌裡面.衪是那位創造天地萬物的神,去愛我們,衪有能力幫助我們,保護我們,縱然我們今天很多的人或事情都不能夠依靠的時候,我們嘗試學習完全的相信,完全的信靠.主啊,我們在這裡,求你憐憫,求你聖靈工作,我們都會祈禱祈求主耶穌基督再一次的興起我們,復興我們,更新我們,讓我們真的可以全然的交託,讓我們心裡面得著真正的平安喜樂和盼望.

聖靈祈求你來到我們每一個人的生命當中,來到我們的生活當中幫助我們,提醒我們,鼓勵我們.主啊,我們一群願意渴慕你信靠你的人,願意我們都能夠找到我們生命最重要的依靠.多謝主,多謝你讓我們一起去聽你的話語,願我們能夠信靠順服,我們這樣的禱告,奉耶穌基督聖名祈求,阿們!

祝福大家!



\newpage

\section{義人的好處}
\label{sec:709}
\textbf{義人的好處 - 李均熊牧師}
\newline
\newline
\hyperref[sec:708]{< < < PREV SERMON < < <}
~
\hyperref[sec:toc]{[返目錄]}
~
\hyperref[sec:710]{> > > NEXT SERMON > > >}


詩篇 34:1-22
\newline
\begin{longtable}{cl}
\hline
\hline
章節 & 經文 (和合本修訂版)\\
\hline
34:1 & \begin{tabularx}{0.7\textwidth}{X} 我要時時稱頌耶和華，讚美他的話常在我口中。 \end{tabularx} \\ \\ \relax
34:2 & \begin{tabularx}{0.7\textwidth}{X} 我的心必因耶和華誇耀，謙卑的人聽見就喜樂。 \end{tabularx} \\ \\ \relax
34:3 & \begin{tabularx}{0.7\textwidth}{X} 你們要和我一同尊耶和華為大，讓我們一同高舉他的名。 \end{tabularx} \\ \\ \relax
34:4 & \begin{tabularx}{0.7\textwidth}{X} 我曾尋求耶和華，他就應允我，救我脫離一切的恐懼。 \end{tabularx} \\ \\ \relax
34:5 & \begin{tabularx}{0.7\textwidth}{X} 仰望他的人，就有光榮；他們的臉必不蒙羞。 \end{tabularx} \\ \\ \relax
34:6 & \begin{tabularx}{0.7\textwidth}{X} 這困苦人呼求，耶和華就垂聽，救他脫離一切的患難。 \end{tabularx} \\ \\ \relax
34:7 & \begin{tabularx}{0.7\textwidth}{X} 耶和華的使者在敬畏他的人四圍安營，要搭救他們。 \end{tabularx} \\ \\ \relax
34:8 & \begin{tabularx}{0.7\textwidth}{X} 你們要嘗嘗主恩的滋味，便知道他是美善；投靠他的人有福了！ \end{tabularx} \\ \\ \relax
34:9 & \begin{tabularx}{0.7\textwidth}{X} 耶和華的聖民哪，你們當敬畏他，因敬畏他的一無所缺。 \end{tabularx} \\ \\ \relax
34:10 & \begin{tabularx}{0.7\textwidth}{X} 少壯獅子尚且缺食忍餓，但尋求耶和華的甚麼好處都不缺。 \end{tabularx} \\ \\ \relax
34:11 & \begin{tabularx}{0.7\textwidth}{X} 孩子們哪，來聽我！我要將敬畏耶和華的道教導你們。 \end{tabularx} \\ \\ \relax
34:12 & \begin{tabularx}{0.7\textwidth}{X} 有誰喜愛生命，愛慕長壽，得享美福？ \end{tabularx} \\ \\ \relax
34:13 & \begin{tabularx}{0.7\textwidth}{X} 你要禁止舌頭不出惡言，嘴唇不說詭詐的話。 \end{tabularx} \\ \\ \relax
34:14 & \begin{tabularx}{0.7\textwidth}{X} 要棄惡行善，尋求和睦，一心追求。 \end{tabularx} \\ \\ \relax
34:15 & \begin{tabularx}{0.7\textwidth}{X} 耶和華的眼目看顧義人，他的耳朵聽他們的呼求。 \end{tabularx} \\ \\ \relax
34:16 & \begin{tabularx}{0.7\textwidth}{X} 耶和華向行惡的人變臉，要從地上除滅他們的名字。 \end{tabularx} \\ \\ \relax
34:17 & \begin{tabularx}{0.7\textwidth}{X} 義人呼求，耶和華聽見了，就拯救他們脫離一切患難。 \end{tabularx} \\ \\ \relax
34:18 & \begin{tabularx}{0.7\textwidth}{X} 耶和華靠近傷心的人，拯救心靈痛悔的人。 \end{tabularx} \\ \\ \relax
34:19 & \begin{tabularx}{0.7\textwidth}{X} 義人多有苦難，但耶和華救他脫離這一切， \end{tabularx} \\ \\ \relax
34:20 & \begin{tabularx}{0.7\textwidth}{X} 又保護他全身的骨頭，連一根也不折斷。 \end{tabularx} \\ \\ \relax
34:21 & \begin{tabularx}{0.7\textwidth}{X} 惡必害死惡人，恨惡義人的，必被定罪。 \end{tabularx} \\ \\ \relax
34:22 & \begin{tabularx}{0.7\textwidth}{X} 耶和華救贖他僕人的性命，凡投靠他的，必不致定罪。 \end{tabularx} \\ \\
[1ex]
\hline
\hline
\end{longtable}
講題:義人的好處 The Righteous May See Good

經文:詩篇34:1-22

講員:李均熊牧師

場合:香港中文大學崇基學院禮拜堂主日崇拜

日期:2024年8月11日

起:詩34作為一首詩

今日的啟應詩篇是詩32篇,全詩共22節,今天選讀1至8節 ；下週會選讀9至14節,而餘下的15至22節則在再下一週讀完.今天啟應開始於詩人的讚美和稱頌上主的邀請,而完於一節很美麗的詩句,和合本譯作「你們要嘗嘗主恩的滋味,便知道他是美善；投靠他的人有福了!」若你細看9-14節,主題「敬畏上主」則呼之欲出,並完於另一句和合本譯得朗朗上口的經文:「要離惡行善,尋求和睦,一心追趕」.最後15至22節著眼於上主對義人的幫助,和合本所譯的第17節也容易記憶:「義人呼求,耶和華聽見了,便救他脫離一切患難!」

當然這樣按主題,又配合其他經課經文,對教會崇拜甚有助益.例如今日的福音經課論到耶穌乃生命的糧,正好暗合「要嘗嘗主恩的滋味」一語.但作為研究舊約的學者,我認為若整體地研讀此詩,我們可能更能欣賞這詩的獨特之處.其實這首詩是一首字母詩,或稱之為離合字母詩(acrostic alphabetical poem),即詩每句的起首字詞都以字母的順序選字,若以英文為例子,則第一行要用a起首的英文字作為句首開始,第二行以b,第三行以c,如此類推,直到26個字母都用完.希伯來文共有22個字母, 故此,這詩共有22節,一節一行,並不需要在分行上再劃分.離合詩的格式有時帶有完全,完備,詩34當中重覆了「所有」一詞共8次之多,配合起來這首詩可能要表達其中一個意思就是第10節下半節和合本所譯「但尋求耶和華的,什麼好處都不缺」.為了將希伯來文中那個「所有」還原,我重譯為「但尋求上主的,他們所有好處都不缺」.「所有」,「一切」,「全部」的好處,都是屬於尋求上主的人!

除了詩34是一首離合詩,這詩另一個特點是將某些字詞多次重複使用.剛說了「所有」(all)一詞出現了8次之多,另一個明顯的重複當然是「上主」/「耶和華」 (the Lord),共出現16次.其他重複出現的有「聽」hear/listen,「敬畏」fear,「好」good,「惡」evil/bad,都出現了4次,「解救」delivery有3次,連其他意義相近的字詞「拯救」save(出現2次),「搭救」rescue和「救贖」redeem,這組詞彙出現了可計7次之多.其他出現2次的字詞順序有「性命」soul,「尋求上主」(seek the Lord 此詩也出現另一個相類字眼「尋找和平」search salom/peace),患難troubles(都是連於解救delivery),看see (都是去看好see good),投靠trust(原意是「找著庇護」find refuge),「不缺」no want/not lack,「做」do(一次do good,一次do evil),「破碎」break/broken(一次心破碎broken heart,一次骨頭連一根也不破碎bone…not one of them is broken)和「受刑罰」to be condemned .

這詩的兩個特點:離合字母詩和重複詞語的出現,使我們不能將之劃分成三首詩,否則我們就會錯失了此詩前後呼應的信息.

如果整首詩來看,我們就會觀察到,詩的受話者addressee 並不是上主,因此這首詩不是一首祈禱詩,而是談論上主如何對待「義人」的一首詩,而且透過上主如何解救「義人」,去教導一起敬拜的「你們」為何／如何去敬畏上主,詩人稱呼他的受話者的「你們」為「祂的聖民」,可見與詩人一起敬拜的,本來就與詩中的上主關係非淺,詩人也如箴言開首般稱呼受話者為「孩子」/「兒子」,可見總的來說,這詩的體栽是一首「智慧教導詩」,而不是一般的祈禱或感恩的詩歌.

承:詩的結構——誰是「義人」?

當中的「義人」一詞其實也有重複出現,但有時這名詞以眾數出現(15節),有時為單數名詞(19,21節),不過大部份時間「義人」在詩的15節以後的下半部份,均以代名詞的形式出現(17節原句並無「義人」一字,是翻譯者後來加上,否則此處的「他們」就不知可所指吧).

在15至18節中的4行詩句中,義人以眾數和第三身眾數代名詞「他們」出現.在另譯這詩的過程中,為凸顯這個眾數代名詞,我都硬譯為「他們」／「人們」(英文用those).而在19至21節中的3行詩句中,義人則以單數名詞和代名詞「他」出現.從代名詞眾數/單數這個角度出發,大家可能會發現8至11節4行詩句中第二身代名詞是眾數的「你們」,而12至14節3行詩句中第二身代名詞則是「單數」的「你」(留意此三句上主一次也沒有出現).也為因如此,將整詩分成7個段落(見附表),行數以4句和3句相間出現,而最後一行自成一段.

第一段4行,是詩人「我」向「你們」發出讚揚上主的邀請,原因此段落最後一句:「我曾尋求上主,祂就應允我,從我所有的恐懼祂解救我.」接著第二段3句將焦點放在「他們」身上:他們仰望他,他們就發光,他們的臉必不蒙羞.但他們是誰?在這段結束便揭曉:就是敬畏他的人們(眾數).但此段中間的困苦人(單數)又是誰呢,就著「祂從他所有的患難拯救了他」與第一段結束「從我所有的恐懼祂解救了我」相若的句式,這個困苦人會不會就是第一段的「我」?同樣相似的句式也出現詩的下半部分的第五段17節和第六段19節,詩人「我」作為敬畏上主的困苦人,其實是進一步代表著義人「他們」,因「我」也就是同獲解救的「義人」.若貫串整首詩篇前後呼應的重複,上半部分詩人所訴說的經歷,不單是他個人的經歷,亦是敬畏上主的人的經歷,也代表下半部分一眾義人.由始至終,這首詩所分享給「你們」的其實都環繞「義人」在上主面前得蒙解救的經歷,因為「我」所親身經歷的就是義人的經歷.

綜合來說,詩中所說的義人是一個/一群怎樣的人呢?他們可能就是第2節所說的一群謙卑的人,不過「貧窮」/「貧乏」可能更接近詩人的原意,因而更能對應「他們一無所缺」,「所有好處都不缺」.但這群詩人自己要代表的義人,也是困苦的,經歷一切(所有)的恐懼和患難,就尋求上主和仰望上主,作為敬畏上主的人,便投靠上主,亦即在上主中尋找庇護,在後半段明顯被稱為義人的一群人,看來他們的困苦,恐懼和患難,以至內心破碎,心靈被壓破,這一切所有都是出於另一群「行惡」(做惡do evil,相對做好do good)的人所做成,但「惡必害死奸人,而恨惡義人的必受刑罰」；也因此,詩人自己,同時邀請「你們」一同頌讚神恩,而詩人又教導「你們」敬比上主之道.

轉:「義人」的好處與凶險

若全詩由「我」與「他」到「他們」都是見證「義人」如何經歷上主,那詩人在詩中間邀請「你們嘗嘗和看看上主的好」,並向「你們」力證敬畏上主的人「所有好處都不缺」,就必然和詩中「義人」關係密切.

詩篇34篇其實並沒有言明「好處」是甚麼,只是說到尋求上主的人「所有好處都不缺」,不單呼應敬畏上主的人們一無所缺,並且以「小壯獅子尚且缺食忍餓」來對比.接著問誰人喜愛生命,愛慕在世有多些日子,為的就是可以看多些「好」(通常和合本會譯see good為「享受」,例如在傳道書多次出現「看為好」也是譯為「享受」),而返過來,怎樣才有多些日子在世上,就是要「做好」——就是禁止舌頭從惡,亦即嘴唇被免詭詐的話——這可能就是那些「做惡」的人的所作所為.

但若是以義人的經驗去評估詩篇中間所強調的「好」,那「義人的好處」會否就是上主將義人所有的恐懼,患難,解救/拯救出來?從詩34的起承轉合的布局,在開首「我」在分享義人得解救的經歷,並引導「你們」去看上主的好,其實最終是在記念義人因惡人的攻擊所受的困苦和患難.惡人的攻擊可能就是口舌的詭詐中傷,使到義人心破碎,靈受壓,令他們窮乏而謙卑,無時無刻的讚美其實源於無時無刻的呼求!

合:「他們」的經驗與「你們」的堅持

讀詩篇或聖經可以跟著經文思路順勢地讀,但更多時候我喜歡逆向思考地讀經文,即認為經文強調某一個方面,其實是多表經文作者或讀者缺乏某一個方面.當詩人提到敬畏中的人「所有都不缺」,其實正表示詩人或讀者正處於缺乏之中；當詩人說到義人必得解救,其實表示經文背後的世界中的詩人或讀者正處於一個義人受逼迫困難危險的處境.當期望上主要對付惡人,說到惡必害死惡人,則經文背後世界正是惡人當道的世界:越是缺乏,便在詩文中越發祈求.所以從我一路分析詩34篇的主題是詩人教導讀者敬拜讚頌和敬畏上主,並以「義人的好處」去作為動機,會否讀者的世界其實正處於一個惡人當道,義人受苦,但上主好像解救無期,以至神的聖民對上主灰心失望,甚至走向惡人的道路,放棄義人之道呢?

其實我在預備這篇講章時,我所感受的並不是見義人的好處而興奮讚美,反而沉重地覺得在我們經驗的世界中,少見義人得解救的經歷,卻見惡人的詭詐控訴至使追求和平公義的義人下在困苦危險之中.不知在我對詩篇34篇的分析中,你又會感受到在帶出義人的好處時,其實我的潛台詞乃是義人其實正正活在困乏患難之中,心靈破碎受壓的境況.

彼得前書3:10-12正正以義人受苦仍要行好的角度引用之詩.若由3:7-17節來看,彼得前書的上文下理清楚交待了這首詩的主題乃是「為義行善而受苦」:

“總而言之,你們都要同心,彼此體恤,相愛如弟兄,存憐憫和謙卑的心. 不要以惡報惡,以辱罵還辱罵,倒要祝福,因為你們正是為此蒙召的,好使你們承受福氣. 因為經上說: 「凡要愛惜生命, 享受好日子的人, 要禁止舌頭不出惡言, 嘴唇不說詭詐的話. 也要棄惡行善, 尋求和睦,一心追求. 因為主的眼看顧義人, 他的耳聽他們的祈禱； 但主向行惡的人變臉.」 你們若熱心行善,有誰會害你們呢? 即使你們為義受苦,也是有福的.不要怕人的威嚇,也不要驚慌； 只要心裡奉主基督為聖,尊他為主.有人問你們心中盼望的理由,要隨時準備答覆； 不過,要以溫柔,敬畏的心回答.要存無虧的良心,使你們在何事上被毀謗,就在何事上使那些凌辱你們在基督裡有好品行的人自覺羞愧. 上帝的旨意若是要你們因行善受苦,這總比因行惡受苦好.” 彼得前書 3:8-17 (和合本修訂本)

我們對詩34的背景並不是十分清楚,但對彼得前書的寫作背景有卻有比較合理的推論.彼得前書成書於羅馬君主尼祿誣告基督徒有分火燒羅馬城,並驅趕基督徒離開羅馬城,至使基督徒成為不受歡迎的人物,那時作基督徒就要流離失所,被人輕看迫害.彼得前書的作者正要鼓勵作為基督的門徒,在受苦中仍要堅持行善,因為基督也是這樣輕看苦難,以祂的義的代替我們的不義,並在艱難的時刻堅持所信.正如彼得前書結束前的鼓勵:

“所以,你們要謙卑服在上帝大能的手下,這樣,到了適當的時候,他必使你們升高. 你們要將一切的憂慮卸給上帝,因為他顧念你們. 務要謹慎,要警醒.因為你們的仇敵魔鬼,如同咆哮的獅子,走來走去,尋找可吞吃的人. 你們要用堅固的信心抵擋他,因為知道你們在世上的眾弟兄也正在經歷這樣的苦難. 那賜一切恩典的上帝曾在基督裡召了你們,得享他永遠的榮耀,在你們暫受苦難之後,必要親自成全你們,堅固你們,賜力量給你們,建立你們.” 彼得前書 5:6-10 和合本修訂本

不知彼得前書的那隻獅子是否受詩34那隻少壯獅子所啟發,但無論如何,彼得前書作者在信徒艱難的處境引用詩34,成為了我們理解此詩的最佳註解.

作為一首字母離合詩,詩篇34篇可能暗示讀者,要走完所有字母,才能掌握詩文的真理.我們又能否在惡人當道,義人受苦的世代,長久的堅持作個敬畏上主,擇善固執,追求和平,以信靠,呼求,尋找上主,以得見「義人的好處」呢?你又能相信上主會解救,保護那些仍在為義受苦,受逼迫的「義人們」呢?



\newpage

\section{千古興亡多少事悠悠}
\label{sec:710}
\textbf{鄧瑞強博士 - 千古興亡多少事,悠悠}
\newline
\newline
\hyperref[sec:709]{< < < PREV SERMON < < <}
~
\hyperref[sec:toc]{[返目錄]}
~
\hyperref[sec:711]{> > > NEXT SERMON > > >}


但以理書 7:1-14
\newline
\begin{longtable}{cl}
\hline
\hline
章節 & 經文 (和合本修訂版)\\
\hline
7:1 & \begin{tabularx}{0.7\textwidth}{X} 巴比倫王伯沙撒元年，但以理在床上做夢，腦中看見異象，就記錄這夢，述說其中的大意。 \end{tabularx} \\ \\ \relax
7:2 & \begin{tabularx}{0.7\textwidth}{X} 但以理說：我在夜間的異象中觀看，看哪，天上有四風，突然颳在大海之上。 \end{tabularx} \\ \\ \relax
7:3 & \begin{tabularx}{0.7\textwidth}{X} 有四隻巨獸從海裡上來，牠們各不相同： \end{tabularx} \\ \\ \relax
7:4 & \begin{tabularx}{0.7\textwidth}{X} 頭一個像獅子，有鷹的翅膀；我正觀看的時候，牠的翅膀被拔去，牠從地上被扶起來，用兩腳站立，像人一樣，還給了牠人的心。 \end{tabularx} \\ \\ \relax
7:5 & \begin{tabularx}{0.7\textwidth}{X} 看哪，另有一獸如熊，就是第二獸，半身側立，口裡的牙齒中有三根獠牙。有人吩咐這獸說：「起來，吞吃許多的肉。」 \end{tabularx} \\ \\ \relax
7:6 & \begin{tabularx}{0.7\textwidth}{X} 其後，我觀看，看哪，另有一獸如豹，背上有四個鳥的翅膀；這獸有四個頭，還給了牠權柄。 \end{tabularx} \\ \\ \relax
7:7 & \begin{tabularx}{0.7\textwidth}{X} 其後，我在夜間的異象中觀看，看哪，第四獸可怕可懼，極其強壯，有大鐵牙，吞吃嚼碎，剩下的用腳踐踏。這獸與前面所有的獸不同，牠有十隻角。 \end{tabularx} \\ \\ \relax
7:8 & \begin{tabularx}{0.7\textwidth}{X} 我正思考這些角的時候，看哪，其中又長出另一隻小角；先前的角中有三隻角在它面前連根被拔出。看哪，這角有眼，像人的眼，有口說誇大的話。 \end{tabularx} \\ \\ \relax
7:9 & \begin{tabularx}{0.7\textwidth}{X} 我正觀看的時候，有寶座設立，上面坐著亙古常在者。他的衣服潔白如雪，頭髮如純淨的羊毛。寶座是火焰，其輪為烈火。 \end{tabularx} \\ \\ \relax
7:10 & \begin{tabularx}{0.7\textwidth}{X} 有火如河湧出，從他面前流出來；事奉他的有千千，在他面前侍立的有萬萬；他坐著要行審判，案卷都展開了。 \end{tabularx} \\ \\ \relax
7:11 & \begin{tabularx}{0.7\textwidth}{X} 於是我觀看，因這角說誇大的話，我正觀看的時候，那獸被殺，身體被毀，扔在火中焚燒。 \end{tabularx} \\ \\ \relax
7:12 & \begin{tabularx}{0.7\textwidth}{X} 其餘的獸，權柄都被奪去，生命卻得以延續，直到所定的時候和日期。 \end{tabularx} \\ \\ \relax
7:13 & \begin{tabularx}{0.7\textwidth}{X} 我在夜間的異象中觀看，看哪，有一位像人子的，駕著天上的雲而來，被領到亙古常在者面前。 \end{tabularx} \\ \\ \relax
7:14 & \begin{tabularx}{0.7\textwidth}{X} 他得了權柄、榮耀、國度，使各方、各國、各族的人都事奉他。他的權柄是永遠的，不能廢去，他的國度必不敗壞。 \end{tabularx} \\ \\
[1ex]
\hline
\hline
\end{longtable}
講題:千古興亡多少事,悠悠 Rises and Falls

經文:但以理書7:1-14

講員:鄧瑞強博士

場合:香港中文大學崇基學院禮拜堂主日崇拜

日期:2024年8月18日

各位弟兄姊妹,早安.願神的恩典常與你們同在.

今日,我們一起看但以理書.

先交代一下但以理書的性質.希伯來聖經將經卷分成三大部分:一,摩西五經；二,先知書；三,著作(Writings).在希伯來聖經裡,但以理書不屬於「先知書」,而屬於「著作」.由於但以理書第7章講到一從天而來的「人子」,及第12章講到死人復活,基督教很容易將之理解成是對耶穌的預言,故此基督教聖經將但以理書歸在「先知書」類別之下.不同的神學,對聖經正典裡同一經卷應放在哪裡,歸屬何類,有不同處理.這些編排,又影響後來的讀者對該經卷的理解.若干年後,若再出現另一個「馬丁路德」時,可能又會將聖經正典裡經卷的次序及歸類,重新整理一次.

就文學類型來說,但以理書前半部屬於宮廷傳奇,後半部屬於天啟文學.但以理書不屬於先知書,故請不要胡亂猜想,書中究竟預言些什麼?其實,就算但以理書屬先知書,也是沒有什麼預言將來的成分的,因為舊約的先知,他們的言論,重點根本就不在預言將來,而是糾正神的子民當下的錯誤,叫他們回歸神.一個典型的先知故事,是大衛王做錯事,殺人奪妻,先知拿單就去和他講「耶和華如此說......」,指出他的錯誤.先知說話的重點,就是在此,不是預告將來,而是糾正當下.

但以理書1:1寫「猶大王約雅敬在位第三年,巴比倫王尼布甲尼撒來到耶路撒冷,將城圍困」,這是主前605年,巴比倫王尼布甲尼撒剛剛登上歷史舞台,圍困耶路撒冷,後來取代了當時亞述的霸主地位,成了耶路撒冷的宗主國.

去到但以理書第11章時,書中講到希臘西流古王朝(Seleucid Empire,塞琉古王朝)「安提阿哥四世」(Antiochus IV Epiphanes)對耶路撒冷的鐵腕統治,「安提阿哥四世」推行文化改造,要將猶大地希臘化.在聖殿裡,「安提阿哥四世」安放不應安放的偶像,他安放宙斯神像在其中,並獻豬肉為祭,這引發猶大民族著名的「馬加比革命」,爭取猶大地的獨立和宗教純潔.第11章結尾,講到「安提阿哥四世」的死亡,時為主前164年.學者相信,這就是但以理書最後成書的年代.

簡單說,但以理書的內容由主前605年到主前164年,橫跨441年.441年是多長時間,我好奇計一計,由現在倒數441年,是1583年,是明朝萬曆年間,距離黃仁宇教授寫的《萬曆十五年》只差4年.

我剛才說,這書寫於「馬加比革命」的時代.「安提阿哥四世」雖然已死,但抗拒希臘化這事是否成功,猶大地能否獨立,聖殿能否保持純潔,一切尚未可料.在這其間,很多對抗,很多戰事,很多傷亡.

我們想像一下,這時,有些猶大神學家要寫一點東西,來鼓舞未知前景的猶大人.如何開始?他們想起400多年前一個類似的大時代,那時,巴比倫王尼布甲尼撒興起,欺壓他們的祖先.他們失去了土地,生活在異鄉.關於在巴比倫的生活,祖先流傳下來幾個故事,講到有幾個信仰英雄,抵住了巴比倫帝國的異教壓力,成就了流芳百世的信仰傳奇.這就是大家常常聽到的,但以理書的故事,即:但以理的三個朋友因不拜尼布甲尼撒的金像而被扔進火yiu 裡卻安然無恙的故事,但以理因堅持向神祈禱而被扔進獅子坑裡卻獲救的故事.這些古代傳奇,很能激勵寫作但以理書的當代人,即「馬加比革命」的當代人,鼓勵他們,堅持信仰,抵擋希臘的帝國暴力和偶像文化.

由於是複述古代的宮廷傳奇,但以理書由2:4b開始,便用了巴比倫宮廷用的亞蘭文.這些傳奇故事去到第6章尾便結束.在我看來,第7章是以歷史神學來總結之前的故事,並開展第8章打後的討論.由於第7章有這轉折功能,讓它顯得有點怪異.其文字仍舊是用亞蘭文,看得出它與2-6章的宮廷傳奇是相關的.但第7章的文體,卻轉用了「天啟文學」(apocalyptic literature)的文體,這配合8-12章的天啟情節.第8章打後,經文便轉用希伯來文了.

「天啟文學」是「舊約時代」到「新約時代」這中間過渡時期興起的文體,與《魔戒》一書的文體相似,有神界人物,魔界人物,人類要在歷史旅程中承擔起艱巨的任務,在力不能勝時會有神界人物的介入,在歷史的終點,邪惡會被克服.天啟文學強調的,是神掌權,祂主宰著歷史的進程,至終,祂的國度會來臨.很多人讀但以理書,一讀到天使,魔君,便以為是預言末世,其實不是的,這只是天啟文學的表達手法而已.它是鼓勵「馬加比革命」的當代人,面對邪惡的暴力,不要怕,要持守純正的信仰,因為神在萬有之上掌權.

我們今日一起讀的經文,是但以理書第7章的經文.文字是亞蘭文,文體是天啟文學.我們看:但以理書7:1—14

但7:1　巴比倫王伯沙撒元年,但以理在床上做夢,見了腦中的異象,就記錄這夢,述說其中的大意.

但7:2　但以理說:我夜裡見異象,看見天的四風陡起,颳在大海之上.

但7:3　有四個大獸從海中上來,形狀各有不同:

但7:4　頭一個像獅子,有鷹的翅膀；我正觀看的時候,獸的翅膀被拔去,獸從地上得立起來,用兩腳站立,像人一樣,又得了人心.

但7:5　又有一獸如熊,就是第二獸,旁跨而坐,口齒內啣著三根肋骨.有吩咐這獸的說:「起來吞吃多肉.」

但7:6　此後我觀看,又有一獸如豹,背上有鳥的四個翅膀；這獸有四個頭,又得了權柄(呂振中譯本:「又有權柄給了牠.」).

但7:7　其後我在夜間的異象中觀看,見第四獸甚是可怕,極其強壯,大有力量,有大鐵牙,吞吃嚼碎,所剩下的用腳踐踏.這獸與前三獸大不相同,頭有十角.

但7:8　我正觀看這些角,見其中又長起一個小角；先前的角中有三角在這角前,連根被牠拔出來.這角有眼,像人的眼,有口說誇大的話.

但7:9　我觀看,見有寶座設立,上頭坐著亙古常在者.他的衣服潔白如雪,頭髮如純淨的羊毛.寶座乃火焰,其輪乃烈火.

但7:10　從他面前有火,像河發出；事奉他的有千千,在他面前侍立的有萬萬；他坐著要行審判,案卷都展開了.

但7:11　那時我觀看,見那獸因小角說誇大話的聲音被殺,身體損壞,扔在火中焚燒.

但7:12　其餘的獸,權柄都被奪去,生命卻仍存留,直到所定的時候和日期.

但7:13　我在夜間的異象中觀看,見有一位像人子的,駕著天雲而來,被領到亙古常在者面前,

但7:14　得了權柄,榮耀,國度,使各方,各國,各族的人都事奉他.他的權柄是永遠的,不能廢去；他的國必不敗壞.

但7:1　巴比倫王伯沙撒元年,但以理在床上做夢,見了腦中的異象,就記錄這夢,述說其中的大意.

但7:2　但以理說:我夜裡見異象,看見天的四風陡起,颳在大海之上.

伯沙撒是巴比倫的末代君王.1:1講到巴比倫尼布甲尼撒的興起.7:1講到巴比倫的末代君王.藉此,將巴比倫時期的故事括起來.之後,便是巴比倫滅亡後的歷史發展了.

上次,講「以西結書」時講過,「四風」是巴比倫主神「Marduk(馬杜克)」的武器,用來克服造成混亂的海神「Tiamat(提亞馬特)」的.如今,耶和華在「Marduk(馬杜克)」之上,取用了他的武器,去克服大海.在聖經語言裡,「大海」常代表政治暴力引發的混亂.

但7:3　有四個大獸從海中上來,形狀各有不同:

但7:4　頭一個像獅子,有鷹的翅膀；我正觀看的時候,獸的翅膀被拔去,獸從地上得立起來,用兩腳站立,像人一樣,又得了人心.

「獸」代表暴力的帝國.在歷史上,獸接續而來.

第一隻獸,是有翅膀的獅子,這是巴比倫常見的神獸,明顯代表巴比倫.留意,這獸的翅膀是被拔去的.有一比巴比倫更強大的力量,駕馭著牠.

但7:5　又有一獸如熊,就是第二獸,旁跨而坐,口齒內啣著三根肋骨.有吩咐這獸的說:「起來吞吃多肉.」

第二隻獸,像熊,力大無窮,一般認為是指米底亞(有時譯為「瑪代」).米底亞很多時與波斯聯合起來,並稱「米底亞」「波斯」.留意,這獸是「被吩咐」的,有一比牠更強大的力量,駕馭著牠.

但7:6　此後我觀看,又有一獸如豹,背上有鳥的四個翅膀；這獸有四個頭,又得了權柄(呂振中譯本:「又有權柄給了牠.」).

第三隻獸,像豹,行動迅速,一般認為是指波斯.留意,這獸的權柄是「被給予牠」的,有一比牠更強大的力量,駕馭著牠.

但7:7　其後我在夜間的異象中觀看,見第四獸甚是可怕,極其強壯,大有力量,有大鐵牙,吞吃嚼碎,所剩下的用腳踐踏.這獸與前三獸大不相同,頭有十角.

但7:8　我正觀看這些角,見其中又長起一個小角；先前的角中有三角在這角前,連根被牠拔出來.這角有眼,像人的眼,有口說誇大的話.

第四隻獸,無法形容似什麼,只形容牠「甚是可怕」.「這獸與前三獸大不相同」,因為前三獸巴比倫,米底亞,波斯全是「亞洲帝國」,講的都是閃族語系的語言.但希臘卻是「歐洲帝國」第一次統治著這些古老的亞洲帝國.希臘這怪獸,「頭有十角」,指其武力十分強大.

在複雜的權力鬥爭中,有一「小角」興起來,這就是希臘帝國下的「西流古王朝」的「安提阿哥四世」.這就是寫作但以理書時面對的惡王,他為了鎮壓耶路撒冷,殺了很多人.「馬加比革命」對抗的,就是他.「這角有眼」,這形像十足是「魔戒」裡「索倫之眼」(Eye of Sauron)之形像.

但7:9　我觀看,見有寶座設立,上頭坐著亙古常在者.他的衣服潔白如雪,頭髮如純淨的羊毛.寶座乃火焰,其輪乃烈火.

但7:10　從他面前有火,像河發出；事奉他的有千千,在他面前侍立的有萬萬；他坐著要行審判,案卷都展開了.

一講到猶大人當下面對的暴君,經文便即刻轉往天上的神.在天啟文學裡,描繪神在天上,是很具體的.祂坐在寶座上,有代表審判的烈火,也有審案的案卷.天啟文學揭露宇宙真相,神在高天之上,審判人間的不義.

但7:11　那時我觀看,見那獸因小角說誇大話的聲音被殺,身體損壞,扔在火中焚燒.

但7:12　其餘的獸,權柄都被奪去,生命卻仍存留,直到所定的時候和日期.

這隻窮凶極惡,自命不凡的野獸,最後被殺,有一比牠更強大的力量,駕馭著牠；牠面對的,是審判的火.

但7:13　我在夜間的異象中觀看,見有一位像人子的,駕著天雲而來,被領到亙古常在者面前,

但7:14　得了權柄,榮耀,國度,使各方,各國,各族的人都事奉他.他的權柄是永遠的,不能廢去；他的國必不敗壞.

這是歷史的遠象.既然人間的帝國,總是變成怪獸,那麼,人類的希望就只能寄望從天而來的人物.這人物不是怪獸的模樣,像人子,帶來人性化的管治.靠這從天而來的人物,最後,開展永恆的天國.到了新約,當耶穌到來時,他引用這經文指向自己,祂的救贖,就是實現天國.

今日講的這段短短的天啟文學的經文,是對悠悠歷史的一個神學反省.一個帝國接一個帝國,用的都是暴力,從來沒有一個帝國以仁愛去愛惜百姓.很多人說,在現實政治講仁愛,在現實世界講仁慈,肯定會被怪獸吞了,與自殺無異.從天啟文學來說,這只是因為我們看不到真相,才認為仁愛不可行.天啟文學開我們的眼睛,讓我們能看見,在人間之上,還有審視萬物的神.祂看著一切,審判著一切,並在某時某刻,主持人間正義,帶來歷史的改變.在神的引領下,歷史是有終極意義的.歷史的當下,縱有艱難,但不用灰心,神會千方百計,讓天國最終實現的.

我們的問題,是見人用暴力,自己便用暴力,將自己變成怪獸一部分,不相信神的仁愛.我們的另一問題,是不太相信神引領著歷史的前景,只關心自己眼前短暫的幸福.當個體的人生脫離了歷史的總體意義時,就像一個人在漆黑的太空中飄浮,注定是孤獨而失落的.很多人活得不開心,因為他只關心自己的開心,拒絕與歷史的總體意義有任何關連.猶大神學家寫這本但以理書,就是要鼓勵在「安提阿哥四世」壓迫下的猶大人,不要逃避現實,不用因當下的逆境而絕望,也不要只顧著自己的幸福,在歷史的這一刻,堅信上帝,盡上自己的努力,使自己有分於天國的建立.願這一切,也成為對我們的鼓勵.

願上主得著榮耀,頌讚；願眾人得享喜樂,平安.



\newpage

\section{靈與肉}
\label{sec:711}
\textbf{林豪恩先生 - 靈與肉}
\newline
\newline
\hyperref[sec:710]{< < < PREV SERMON < < <}
~
\hyperref[sec:toc]{[返目錄]}
~
\hyperref[sec:712]{> > > NEXT SERMON > > >}


講題:靈與肉 Spirit and Flesh

講員:林豪恩先生

場合:香港中文大學崇基學院禮拜堂主日崇拜

日期:2024年8月25日

一．靈與肉的困惑

如果你是宗教聚會講員,有沒有想過講道之後的結果會是這樣呢?

6:66從此,他門徒中多有退去的,不再和他同行.

令人驚訝的是,這位講員並不是無名小卒,而是耶穌.耶穌講道都產生「趕客」效果,到底發生甚麼事呢?

6:60他的門徒中有好些人聽見了,就說:「這話甚難,誰能聽呢?」

原來耶穌講了人們聽不下的說話.到底是甚麼話如此難聽呢?

6:56吃我肉,喝我血的人常在我裡面,我也常在他裡面.

耶穌竟然叫人吃他的肉,喝他的血.面對如此驚嚇的言論,輕鬆的回應可能是:「傻的嗎?」批評的回應可能是:「還不是邪教嗎?」耶穌這幾句話,不但引起當時猶太人反感,後來基督教傳入中國時也成為猜疑和攻擊的話柄.

如此容易引起誤會,甚至「趕客」的說話,何以耶穌在那個場合發表呢?到底這些說話有甚麼意思呢?歷代讀者各有見解,今次我們從兩方面思考:第一,說話通常都有前文後理,追蹤說話所在的脈絡,避免斷章取義或望文生義.第二,把說說放回脈絡中來理解.

二．靈與肉的脈絡及意思

耶穌關於「吃我肉,喝我血」的講論,可以追溯到那一件事呢?

1．五餅二魚帶出食物這個主題

這是7月28日的福音經課,耶穌以五個餅二條魚餵飽5000人的事件.這件事帶出了「食物」這個主題.

2．食物這個主題帶出了「必壞的」和「永生的」兩種食物

耶穌施行了五餅二魚神蹟,及引了一大群Fans.8月4日的福音經課指出耶穌和門徒渡海去迦百農避開群眾,但群眾窮追不捨.還記得耶穌如何回應嗎?

6:26耶穌回答說:「我實實在在地告訴你們,你們找我,並不是因見了神蹟,乃是因吃餅得飽.

6:27不要為那必壞的食物勞力,要為那存到永生的食物勞力,就是人子要賜給你們的,因為人子是父上帝所印證的.」

3．帶出猶太人和耶穌討論甚麼是永生的食物

談及永生的食物,猶太人立即想到摩西叫天降下的嗎哪, 在8月11日的課音經課耶穌澄清:

6:49你們的祖宗在曠野吃過嗎哪,還是死了.

6:50這是從天上降下來的糧,叫人吃了就不死.

甚麼是「從天上降下來的糧」呢?有甚麼食物「叫人吃了就不死」的呢?

4．耶穌表明這種糧食就是他自己

8月18日的福音經課耶穌如此說:

6:51我是從天上降下來生命的糧；人若吃這糧,就必永遠活著.我所要賜的糧就是我的肉,為世人之生命所賜的.」

從這個脈絡來看,自從五餅二魚開始,話題一直環繞著「食物」這個主題.人們對於「五餅二魚」類的食物既認識又需要,因為人們以此來支撐肉身的生命；談到靈性的生命,又有甚麼食物可以支撐永恆的生命呢?

因此在今天的福音經課,耶穌指出:

6:58這就是從天上降下來的糧.吃這糧的人就永遠活著,不像你們的祖宗吃過嗎哪還是死了.」

支撐永恆生命的食物必然來自天上那永恆者,而不是來自地上會腐壞的食物,這是可以明白的,只是,當時的聽眾想像不到怎樣可能,也無法接受.

對於「吃肉,喝血」的理解,可能懂中文的人會容易接受一點,因為中文有四字詞語謂「血肉長城」.大家都明白意思不是用血和肉作為建築材料去建造長城,而是表達有人犧牲性命來抵抗外敵拯救家園.會不會耶穌的意思也是獻出自己的血肉,即是犧牲自己的生命來拯救人呢?這個想法,約翰福音的作者是同意的:

11:49內中有一個人,名叫該亞法,本年作大祭司,對他們說:「你們不知道甚麼.

11:50獨不想一個人替百姓死,免得通國滅亡,就是你們的益處.」

11:51他這話不是出於自己,是因他本年作大祭司,所以預言耶穌將要替這一國死；

11:52也不但替這一國死,並要將上帝四散的子民都聚集歸一.

使徒保羅也同意,他在以弗所書說:

2:13從前你們是遠離上帝的人,如今卻在基督耶穌裡,靠著他的血,已經得以親近了.

2:14因為他自己是我們的和平,使雙方合而為一,拆毀了中間隔絕的牆,而且以自己的身體終止了冤仇,

2:16既在十字架上消滅了冤仇,就藉這十字架使雙方歸為一體,與上帝和好.

使徒保羅也用基督耶穌的「血」和「身體」這些詞彙為媒介,成就「得以(和上帝)親近」,「與上帝和好」的效果,這些效果在今天的福音經課耶穌也說:

6:54吃我肉,喝我血的人就有永生,在末日我要叫他復活.

如果覺得「永生」這個概念太抽象,末日太遙遠,下句平衡句子的補充是:

6:56吃我肉,喝我血的人常在我裡面,我也常在他裡面.

「常在我裡面,我也常在他裡面」類似的句子,讀過約翰福音的朋友會感到熟識,耶穌在其他場合也講過,有人認為這是約翰福音的重要主題之一.

三．對靈與肉的回應

當時聽完耶穌講道的人,有失望而離開,但也有人繼續跟隨:

6:67耶穌就對那十二個門徒說:「你們也要去嗎?」

6:68西門‧彼得回答說:「主啊,你有永生之道,我們還歸從誰呢?

是甚麼做成離開的和留下的這兩批人的分別呢?原因可能是大家對耶穌懷著不同的期望.彼得向耶穌尋求的是「永生之道」,耶穌提供的剛好就是「永生之道」,互相配合.對另一群人來說,他們向耶穌尋求的是「吃餅得飽」,但耶穌提供的卻是「永生之道」,兩者並不配合.俗語說,因誤會而結合,因了解而分開.耶穌似乎也不勉強,因此問留下的人:「你們也要去嗎?」

在跟隨耶穌的路上,類似的問題也值得我們不時問問自己,我向耶穌尋求甚麼呢?如果我們所期望的和耶穌提供的不同,我們如何回應呢?曾經同路的人或許已經離去,如果耶穌問你:「你也要去嗎?」你會如何回答呢?

四．靈與肉的應用

今天的福音經課指出耶穌分享了他的身體,這意味著分享甚麼呢?對於跟隨耶穌的人有甚麼啟發呢?我們可以從下列三方面思考:

分享了關係

有人認為耶穌死在十字架上成就救恩,我得到天堂的入場券,這就是信仰的全部.原來耶穌所賜予的比這些更豐富.耶穌期望的是「人常在我裡面,我也常在他裡面」如此密切的關係.這種關係如何可能呢?如何建立呢?最近讀靈修學和心理學教授葛羅謝爾神父的書,書名為《煉淨,光明,合一:靈性成長的心理學》.

書中描屬靈修之路的三個階段,分別是煉淨之路,光明之路,以及合一之路.作者介紹了幾位經驗到「與天主合一的境界」的先賢,例如:聖文德,聖女大德蘭,聖十字若望等.讀到先賢們與上主如此深入及親密的記錄,今人感到羨慕,也燃起期盼.原來我們的先賢經驗的信仰是如此豐富和深刻,「人常在我裡面,我也常在他裡面」是真實地有人經驗到的.耶穌基督分享他的身體,就是分享了關係,如果用血肉這詞彙來講,可說是一種血肉相連的關係.

分享了價值觀

因五餅二魚吃飽的群眾聚焦在飽腹的糧食,對耶穌提出的永生糧食不感興趣,這是可以理解的,我們何嘗不是這樣呢?常言道「耶穌俾飯你食呀?」不過,如果把人生壓縮到只剩下吃餅得飽,也不是大家想要的情況.世界上的確有人捱饑抵餓,然而,也有研究指出,捱饑抵餓的問題不在於欠缺糧食,而是分配不均.糧食是看得見,數得到的物質；分配則由看不見,也數到的力量影響,例如:價值觀.今天的書信經課,使徒保羅帶我們進入價值和品格的範疇來回應福音經課:

6:10我還有末了的話:你們要靠著主,倚賴他的大能大力作剛強的人.

6:12因我們並不是與屬血氣的爭戰,乃是與那些執政的,掌權的,管轄這幽暗世界的,以及天空屬靈氣的惡魔爭戰.

6:14所以要站穩了,用真理當作帶子束腰,用公義當作護心鏡遮胸,

6:15又用平安的福音當作預備走路的鞋穿在腳上.

6:16此外,又拿著信德當作盾牌,可以滅盡那惡者一切的火箭；

6:17並戴上救恩的頭盔,拿著聖靈的寶劍,就是上帝的道；

6:18靠著聖靈,隨時多方禱告祈求；並要在此警醒不倦,為眾聖徒祈求,

英文諺語說:You are what you eat,可以理解為吃甚麼,就會像甚麼.以耶穌的身體為食糧的人,又似甚麼呢?保羅指出這些依靠耶穌,以他為主的人會培育出以下特質:真理,公義,和平,信德,救恩(安全感),為眾聖徒祈求(關懷他人).如果人間充滿這些特質,或許糧食分配的狀況得以改變,吃餅得飽不會再成為人們壓倒性的追求,有空間發現自己及他人的需要.

3．分享了使命

既然這幾個主日的福音經課的題材都環繞著食物,大家知不知道耶穌的食物是甚麼呢?

福音書記述耶穌也吃當時的食物,但他自己是這樣說的:(約翰福音)

4:32耶穌說:「我有食物吃,是你們不知道的.」

4:33門徒就彼此對問說:「莫非有人拿甚麼給他吃嗎?」

4:34耶穌說:「我的食物就是遵行差我來者的旨意,做成他的工.

耶穌的食物就是「遵行差我來者的旨意,做成他的工」,可以理解為當耶穌遵行天父的旨意,完成天父的託付時,就得到靈性的飽足.耶穌呼籲人以他為食物,會不會同時呼籲人「遵行差我來者的旨意,做成他的工」呢?耶穌清楚知道「做成他的工」是甚麼意思.對我們來說,「做成他的工」又是意味著甚麼呢?使徒保羅在以弗所書的說這值得我們參考:

2:10我們原是他的工作,在基督耶穌裡造成的,為要叫我們行善,就是上帝所預備叫我們行的.

保羅的說話啟發我們,「做成他的工」原來不是要勞役我們,而是呼喚我們活出上主創造的美好,因慈相待,彰顯良善,這就是上主「預預備叫我們行的」.



\newpage

\section{得居聖山者之品行}
\label{sec:712}
\textbf{李均熊牧師 - 得居聖山者之品行}
\newline
\newline
\hyperref[sec:711]{< < < PREV SERMON < < <}
~
\hyperref[sec:toc]{[返目錄]}
~
\hyperref[sec:713]{> > > NEXT SERMON > > >}


詩篇 15:1-5
\newline
\begin{longtable}{cl}
\hline
\hline
章節 & 經文 (和合本修訂版)\\
\hline
15:1 & \begin{tabularx}{0.7\textwidth}{X} 耶和華啊，誰能寄居你的帳幕？誰能居住你的聖山？ \end{tabularx} \\ \\ \relax
15:2 & \begin{tabularx}{0.7\textwidth}{X} 就是行為正直、做事公義、心裡說實話的人。 \end{tabularx} \\ \\ \relax
15:3 & \begin{tabularx}{0.7\textwidth}{X} 他不以舌頭讒害人，不惡待朋友，也不隨夥毀謗鄰舍。 \end{tabularx} \\ \\ \relax
15:4 & \begin{tabularx}{0.7\textwidth}{X} 他眼中藐視匪類，卻尊重那敬畏耶和華的人。他發了誓，雖然自己吃虧也不更改。 \end{tabularx} \\ \\ \relax
15:5 & \begin{tabularx}{0.7\textwidth}{X} 他不放債取利，不受賄賂以害無辜。做這些事的人必永不動搖。 \end{tabularx} \\ \\
[1ex]
\hline
\hline
\end{longtable}
講題:得居聖山者之品行 The Character of Those Who May Dwell with the Lord

經課:詩篇15:1-5

講員:李均熊牧師

場合:香港中文大學崇基學院禮拜堂主日崇拜

日期:2024年9月1日

1. 引言:居主聖山者的品行

今天啟應的詩篇是詩篇十五篇,相對於上次我分析的詩篇三十四篇,可謂十分簡單直接.5節詩句,除第一節/第一句引子用上一個「誰」的問題的平行句作開始外,主體都是用上三疊遞進three-fold progression的格式,即每句環繞同一意念,以三短句構成一句.此詩文字直白,並沒有用上任何的比喻,典故,意象,並不需要讀者太多運用想像力.所以我也不會如上次一般拐彎抹角,就直接討論得居主聖山者品行上的三個範疇,以三點式作為講道大綱.

開始之先,要對這首詩的主題作一回應.研究詩篇的學者,大都定性這首詩為聖殿的進門禮儀詩(Entrance Liturgy),即認為這是一首朝聖者要進入聖所必先經歷的儀式.進門前祭司會用問句詢問進入聖所的朝聖者誰人可進入聖所,入門者或祭司則回應當中進入聖所者的品行要求.所以學者多以詩論敬拜與敬拜者道德的角度討論本詩.

但我看這種解釋有些困難.對比同被認為是進門禮詩的詩篇24篇,詩篇24篇的提問是:誰能登耶和華的山?誰能站在他的聖所?(詩24:3)詩篇15篇的起首明明問:

15:1 耶和華啊,誰能寄居你的帳幕? 誰能居住你的聖山?

「誰能居住在你的聖山中」,不是問誰何進入聖殿或會幕.誰會居住在聖山會幕的呢?並不是來了又去的朝聖者,更可能是在委身會幕供職的祭司.若以今日的語言講,是誰可以委身事奉上主,在基督的教會服事呢?選擇居於聖山的,是願意呈獻自己全心全意服事上主的人,而非特來拜拜的一般朝聖者.當然事奉或英文Service,都指向「敬拜」worship,但我認為這詩的對像不是對只是每一週來禮拜堂參加敬拜的信徒,而是對願意委身於服事上主的人來說的.

詩第二節便直接講到這些願居聖山,願意委身服事上主者品行的三個面向:

15:2 就是行為正直,做事公義, 心裡說實話的人.

「行為正直」:行為原文是「行」,可解作日常的行為,平時習慣.而「正直」不是「曲直」的直,而是解作「完全／全備」,即於創世記17章當上主與亞伯拉罕立永約時,要他先「在我面前做完全人」.但經文中的「完全」這個字,並非指「十全十美」(perfect),而是一個禮儀用詞,最常出現的是在利未記,如利1:3「他的供物若以牛為燔祭, 就要在會幕門口獻一隻沒有殘疾的公牛, 可以在耶和華面前蒙悅納」中的「沒有殘疾」,應視為「祭品無瑕疵」,應用在人身上,應循「可蒙悅納」的方向去思考.

「做事公義」:若「行」是日常和平時的表現,「做事」則是在特定的工作,職務上的表現,用現代角度看可以是我們的工作job或職業career.而「公義」在舊約出現多與「公平」共用,在古代君王所立的銘刻中主要是指「社會公義」(social justice),特別是照顧孤兒寡婦,或在社會中無權無勢,貧苦無依的人.

「心裡說實話」:這是關於「說話」的範疇.「實話」原文就是我們常說的「阿們」,誠實,真心,希臘文譯為alethea(女孩常取名字),是指真理,真實.而「心裡」今日會視為情感或思想的所在,但古代「心」卻更多是指「決斷／決策／決定」的地方,而非一般的心思意念:「你要保守你心,勝過保守一切,因為一生的果效是由心發出」(箴言4:23)乃是指做人要有好的決斷力,因為任何一個決定都有可能影響一生.

之後的每一節,就剛好可以圍繞著居聖山者品行上這三個範疇發揮(我也未深究這是否這詩原本的設計,但起碼我認為可以這樣解釋):第3節關於決定如何運用言語,第4節關於日常習慣中如何蒙主接納,和第5節在職務工作上如何實踐社會公義上.

2. 品行第一個範疇:言語／說話

15:3 他不以舌頭讒害人, 不惡待朋友, 也不隨夥毀謗鄰舍.

若此句是討論居聖山者在「言語上」的品行,那中間一小句「不惡待朋友」,雖然沒有說明怎樣「不惡待」或「不作」(參英譯nor does或參另譯:不作——對他的朋友——壞事),但或許就是指用言語對朋友做壞事／惡事了.(不知何解KJV將原文通常解作「朋友」的譯成neighbor,反而下半節原文通常解作「鄰舍」的譯成friend?)

至於「讒害」/「讒謗」原文其實從名詞「腳」一字而來,動詞在聖經另外一些地方是解「派出探子四處窺探」,如創世記42章中約瑟指控去埃及買糧的兄弟為「奸細」或在約書亞記6章中約書亞派兩個探子去窺探耶利哥城；在何西亞書11:3則比喻上帝「教行路」.只有撒下19:27用於散播謠言(那處的謠言是關於瘸腿的米非波設,可能經文是一語相關,謠言米非波設其實是假裝瘸腿,其實是可以四圍行走的),而詩篇15節就是第二次用在「舌頭」上.大家或者可以形象化的想像一下,「以舌頭讒害人」或「用舌頭散播謠言」就是一條舌頭有一對腳,如探子／奸細一般在大街小卷上四圍行走,窺探私隱.

整句的意思就是居聖山者不會口舌招尤,搬弄事非；又或者透過窺探和散播別人(朋友／鄰舍)私隱,製造謠言,惡意中傷,在背後說朋友鄰舍的壞話.這樣刻意為之自然罪大惡極,但有時教會內也會因為關心別人,以至周圍討論某人某家庭的困難,不知不覺成為散播私隱,或將錯誤的假消傳播的人,無心卻成為製造傷害的元兇.總而言之,就是事奉上主(牧師?),於教會互相稱呼對方第兄姊妹,也必須小心說話,尤其是別人的事情,在談論時就要小心處理了.或者,得居聖山者若愛護友朋鄰里,都會是在實際行動上,而非只是在言語談論上.

3. 品行第二個範疇:日常所行

接著的一節, 我將之歸類為居聖山者之日常所行,或可指委身敬拜事奉上主之人的生活態度:

15:4 他眼中藐視匪類, 卻尊重那敬畏耶和華的人. 他發了誓,雖然自己吃虧也不更改.

他尊重敬畏上主的人,而在他眼中卑鄙的人(參另譯)則會藐視.「藐視」簡單點說就是輕視,不理會,和尊重(原文字根從「重」)代表的重視剛好相反.那經文中的「匪類」／卑鄙的人(參另譯)又會是甚麼人呢?其實此字有「厭惡」／「厭棄」despiesd或「拒絕」rejected之意,根據利未記26:43-44對違反上主立約之咒詛的總結:「他們離開這地,地在荒廢無人的時候就要享受安息.並且他們要服罪孽的刑罰,因為他們厭棄了我的典章.心中厭惡了我的律例,我卻要為他們的緣故記念我與他們先祖所立的約,他們的先祖是我在列邦人眼前從埃及地領出來的,為要做他們的神．我是耶和華.」上主並不會厭棄祂的子民,反而是上主之民屢次厭棄上主.所以無論是正是反,得居上主聖山的人,日常的態度就是重視上主,「當神有到」,而拒絕上主的人他都「費事采佢」(不作理會).

但這種人會如何表現呢?「他發了誓,雖然自己吃虧也不更改.」發誓自然是對上主發誓,與人立約,但可能在某些情況之下,他與人訂定之約反而對自己不利,例如……,但正因他尊重上主,認真看待他在上主面前之約,所以寧願「蝕底」,也不違背在上主面前的立誓.

我初也以為「他發了誓,雖然自己吃虧也不更改」 ,就是在社會規範和法理依據之外,仍是一個有道德的人,不自私自利,不是每事都只講求自己的權益,而甘願為別人「蝕底」.但再看經文的字眼,這個能居聖山的人,就是第2節中的「完全人」,生活行為「可蒙悅納」,就不單是道德高尚,寧願吃虧也不負人；而是因信仰上主,尊崇上主,是寧負自己也不負上主,上主對這種人怎能不悅納呢?這可謂完美演譯羅馬書12:1活祭的意思 :「將身體獻上,當做活祭,是聖潔的,是神所喜悅的,你們如此事奉乃是理所當然的.」

4. 品行第三個範疇:工作／職務

若然詩15:2 做事公義是能居聖山者的品行,那此詩最後一節則能反映這人在工作上的態度:

15:5(上) 他不放債取利, 不受賄賂以害無辜.

此句原文正好是一「給」give(參另譯:金錢他不給出),一「取」take(參另譯:賄賂……他不收取),所以無論是施或受,在職務行使上面都要實踐公義.若果公義或社會公義是代表照顧社會的孤寡,或在社會中無權無勢,貧苦無依的人,那「不受賄賂以害無辜」當然是行使社會公義,但不放債取利,又和社會公義有甚麼關係?五經律法清楚表明放債取利是不容許的,例如出埃及記22:25「我民中有貧窮人與你同住,你若借錢給他,不可如放債的向他取利.」為甚麼呢?這要從古代債務問題說起.

要知道古代並無今日比較完善的銀行和金融體系.特別在一些農村社會,大部份按著祖地的分配,自耕自養,自給自足,沒有必要借貸,也不會有太多金錢的積蓄,更勿論放債修飾等今日視為平常的商業活動.但正因為沒有金融保險等等現代視為平常的保障系統,遇有意外不幸或乾旱天災,以至帶來農產失收等困難日子,一個家庭就可能面對財務上的困境.唯一可做的就是投靠別人,等有能力的人申出援手以渡時艱.(一時的失收可能帶來幾年的失收,因為務農的會在今年的收成預留下年播種的種子,固一年的失收可能引至緊接數年都未能耕種的情況,唯一出路是向別家借種子或借貸金錢購買種子,這樣或可走出困境.)若然這時利「為債利」而給出金錢,這不單不是對窮困無依者帶來幫助,這根本就是見到別人在困難中謀取利益.

因此,不給出金錢為債利不是居上主聖山者在職務上不能做銀行金融工作,而是在任何職務上,都要為著社會公義而行,以協助窮困無依者,而非趁機會壓榨需要幫助的人.

5. 結語:做這些事的人必永不動搖.

詩篇15篇,最後一節還有一句短句:

15:5(下) 做這些事的人必永不動搖.

但其實除了第4節上半節提到正面要做的事外,其餘詩篇15篇提到的都是負面的「不」.舊約倫理好像沒有太多正面的建議,要我們在生活工作如何去「做」,每次提起上主的要求,都是以「不可」為主,例如十誡也是以「不可」為主.又其今日正向心理學提醒我們不要講太多負面的事情,教導小孩也應該多說「該做」甚麼,小些禁制,多些鼓勵；十誡或詩篇15篇的「不可」就變得不合時宜.但猶太人的智慧告訴我們,上主的「不可」是為祂的子民設置的籬笆,為我們的行為準則設置的界限.在這界限內,可以做的事情,比不可以做的事情自然多十倍佰倍.但在每個情況每個處境,應如何合適有創意地實踐信仰和道德,那要記載這些,十誡可能就佰變成佰誡,仟誡了.所以,詩篇15篇提到的是居主聖山者不應做的事情,那也只有六,七,八項,但今日如何實踐愛神愛鄰舍愛社會上困苦的大罪,則要我們按每天多變的處境好好思考和實踐了.而最終,活在上主的籬笆中,生活就可以穩妥,不動搖了(最後也是一個「不」).

我們有時講求敬拜祈禱讀經,都可能只求個人的敬虔,自己安心,返教會變成追求上主祝福自己,但是詩篇15篇提醒我們,敬拜上主,委身上主的人,並不是追求自我利益或私德,而是因為敬畏上主,重視上主的約,願意在關愛友鄰,和在社會中實踐公義.



\newpage

\section{睇}
\label{sec:713}
\textbf{姚志豪牧師 - 睇}
\newline
\newline
\hyperref[sec:712]{< < < PREV SERMON < < <}
~
\hyperref[sec:toc]{[返目錄]}
~
\hyperref[sec:714]{> > > NEXT SERMON > > >}


講題:睇 Look

講員:姚志豪牧師

場合:香港中文大學崇基學院禮拜堂主日崇拜

日期:2024年9月8日

情人眼裡出西施,狗眼看人低,……「睇」人有很多種視野,只是我們怎樣去選擇,或許選擇的背後,正反映我們心裡的價值,而呢D價值實在是受很多野影響的,難怪有說:「眼睛是心靈的叛徒.」

人究竟會點樣睇一個人?如果他不說話以先,相信第一睇就是睇外觀,打扮.所謂人靠衣裝.苦心自問:你要去一個地方,若果接待你的人打扮得企理一點,談吐溫文一點,大概印象也會好一點.細個的時候,爸爸叮囑我:「要簽好你的名字,因為日後當你見工寄求職信,未見其人,先見其字.」(當然時代轉變,今天未必完全合用).不過大家都知,這個其實很表面的,但又係使我們好容易有一種價值判斷在內.外表的觀感確實係一種真實的感覺,只是這不是整個人的全部.有時,更是金玉其外,敗絮其中.

箴言是智慧文學,不是某某人的智慧,其實是生活的集體智慧的結集.因此不是一家之言而成,就是人與人之間的日常.(現中)箴22:1-2「信譽比財富寶貴；名望比金銀可羨.貧窮人和有錢人有一處相同,二者都是上主所造.」經文提到,點睇一個人呢?內裡/待人接物比你所擁有的寶貴.有些人一生很多朋友,可能是以財富地位來結交的,也有是從人品來建立的.或許香港咁急速的城市,使我們沒有時間看得深入,只能看表面,甚至會考慮到往後的利益而建立的關係.久而久之,不經覺間,人靠衣裝彷佛是以貌取人的糖衣.

雅各書2章,一開始就提到:不憑外貌睇人(待人).(現中)就加多句「就不可憑著人的外表,用不同的態度對待人.」態度:就是同一件事,因著你的身份,對你的語氣,言詞,面部表情都不一樣.EG.打條呔VS 著短褲就已經大不同.跟住就更仔細描述處境出來:「若有一個人戴著金戒指,穿著華麗的衣服,進入你們的會堂,又有一個窮人穿著骯髒的衣服也進去,而你們只看重那穿華麗衣服的人,說:「請坐在這裡」,又對那窮人說:「你站在那裡」,或「坐在我腳凳旁」；(現中)2:4你就是偏心,是惡意歧視人.」

事實上,雅各書的背景:當時讀者一般就是貧農受剝削和欺詐的一群,有學者更指雅各書是針對當時流動性高的中產人士,因此,雅1:27「在父上帝眼中(點睇),那純潔沒有缺點的虔誠便是:照顧苦難中的孤兒寡婦和保守自己不受世界的腐化(和:不沾染世界).」而且更深一層指出,2:6「你們(富有)卻羞辱貧窮的人.欺壓你們(貧窮),拉你們到公堂去的,不就是這些富有的人嗎?」5:4「你們剝削在田裡替你們做工的工人,剋扣他們的工錢,他們在控訴你們呢!那些收割工人喊冤的聲音已經進到主—萬軍的統帥的耳朵了.」所以作者要好坦白/寫下書信去提醒,要留心,信仰所關心的是甚麼?點樣睇一個人,從而點樣對待.有沒有不經意受世界污染(馮:拒絕這世界污染) .

換言之,不是說點樣靠外表睇人咁簡單,「以外貌取人」更多一層意思,就是一方面不要睇低果D外表貧寒的；另方面,更不要被用華麗衣裝所遮掩住壓迫人的面具,要留心,有沒有已經被污染到不捨得除下這面具!《穿起西裝不會變高層,返了教會不會變好人》EG.教會咳嗽的婆婆,搬到下一條街.

馬可福音再進一步指出這種以貌/種族來睇人的現象隨處可見.馬可福音的讀者以外邦人群體為主,因此作為猶太 人的耶穌正穿越猶太人所鄙視的外邦地區(推羅.西頓),並且接觸緊一位受鬼附的女孩身上,耳聾結舌的人(宗教上污穢).對於分別為聖的猶太人傳統來說:兩者都應該被隔絕,甚至乎當時要放他們在一些樹林.耶穌卻主動與他們相遇,表達一種很強烈的訊息,外邦人沒有被排除以外.EG.(午夜男).

不過,可7:27耶穌對她說:「讓孩子們先吃飽,拿孩子的餅丟給小狗吃是不妥的.」以「小狗」喻「外邦人」是十分不妥當的.就是我們難以接受為何耶穌會這樣睇外邦人?歷來有少學者嘗試去解釋:

1) 耶穌在考驗這位婦人的信心.最終外邦被邊緣的聲音,能勇敢地對權力說真話.

2)一些人認為耶穌最初的拒絕(第27節)是帶著調皮的神情說出的,他是在給這位女性一個表達他知道她心中所懷信仰的機會,然後再心甘情願地治癒她的女兒.

3) 亦有認為耶穌這樣說,是以拉比的爭辯方式來諷刺法利賽人的態度,並為周圍的人和女人提供教訓.我們無法知道他的面部表情或語調,但我們知道他對請求者的關係模式,始終充滿溫柔和感人

4)有很多人嘗試幫耶穌解釋.基督論:如果我地話JC是完全既人,又是完全的神,所以我們會說耶穌係有情感的,因此會喜,怒,哀,樂,即是話被人世間(猶太文化)價值觀影響是有可能.事實上,耶穌不是一次顯出人性的張力「苦杯挪開」但神性的力量卻改變了事情及人性的睇法.

不論是基於哪一個原因,最終的結果不是以猶太人的睇法去拒絕這位婦人.而是看見了耶穌衝破/超越了人世的價值而說:你回去吧!鬼已離開你的女兒了!(可7:29).

當我們把馬可第七章整章去細讀,你會發現可7:1-8,14-15,21-23 (上主日經課),正是要表達出那固有的猶太潔淨規條,點樣「睇」一個人潔吾潔淨,以致點樣限制左一個人,有無視一個人為人,而非種族,類別.耶穌就是睇穿左經學教師(文士)奉行猶太飲食條例,就是要有效地將猶太人與鄰邦異教徒區別開來,就是「聖」「俗」二分,以這種視野,規條,傳統去「睇」一個有血有肉的人.因此故事的顛覆性力量遠超出傳統的界限突破.耶穌的教導已經廢除了潔淨(猶太人)和不潔(外邦人)之間的傳統區分.

共觀福音太15:28都有記載這故事,耶穌:「你的信心救了你」如果這句說話去神話化,耶穌說了一句「人話」今天這句說話亦是很有價值.EG.肥/女人40路/范柏依/殘奧,是人規限左自己?定是自己規限左自己?

究竟係咩既框框阻礙我們睇一個人?傳統,習慣,文化,傳媒,價值,身份…….相反這個框框又點樣框住左我自己睇自己呢?我既價值是被甚麼來定斷呢?EG.點樣看----樹

「賢愚共一丘」-----當我們面對赤裸裸的BB「出生」,或是一點也帶不走面對「死亡」------詩146:4「他的氣一斷,就歸回塵土,他所打算的,當日就消滅了.」會否讓我們反思:我赤身出於母胎,也必赤身歸回意義?我地睇吾睇見:人的原(根)本才是反映上主特質呢?你睇------上主照住佢自己的形象而創造左你.



\newpage

\section{性價比與等價交換}
\label{sec:714}
\textbf{高國雄牧師 - 「性價比」與「等價交換」}
\newline
\newline
\hyperref[sec:713]{< < < PREV SERMON < < <}
~
\hyperref[sec:toc]{[返目錄]}
~
\hyperref[sec:715]{> > > NEXT SERMON > > >}


馬可福音 8:27-38
\newline
\begin{longtable}{cl}
\hline
\hline
章節 & 經文 (和合本修訂版)\\
\hline
8:27 & \begin{tabularx}{0.7\textwidth}{X} 耶穌和門徒出去，往凱撒利亞‧腓立比附近的村莊去。在路上，他問門徒：「人們說我是誰？」 \end{tabularx} \\ \\ \relax
8:28 & \begin{tabularx}{0.7\textwidth}{X} 他們對他說：「是施洗的約翰；有人說是以利亞；又有人說是先知中的一位。」 \end{tabularx} \\ \\ \relax
8:29 & \begin{tabularx}{0.7\textwidth}{X} 他又問他們：「你們說我是誰？」彼得回答他：「你是基督。」 \end{tabularx} \\ \\ \relax
8:30 & \begin{tabularx}{0.7\textwidth}{X} 於是耶穌切切地囑咐他們不可對任何人說起他。 \end{tabularx} \\ \\ \relax
8:31 & \begin{tabularx}{0.7\textwidth}{X} 從此，他教導他們說：「人子必須受許多的苦，被長老、祭司長和文士棄絕，並且被殺，三天後復活。」 \end{tabularx} \\ \\ \relax
8:32 & \begin{tabularx}{0.7\textwidth}{X} 耶穌明白地說了這話，彼得就拉著他，責備他。 \end{tabularx} \\ \\ \relax
8:33 & \begin{tabularx}{0.7\textwidth}{X} 耶穌轉過來看著門徒，斥責彼得說：「撒但，退到我後邊去！因為你不體會神的心意，而是體會人的意思。」 \end{tabularx} \\ \\ \relax
8:34 & \begin{tabularx}{0.7\textwidth}{X} 於是他叫眾人和門徒來，對他們說：「若有人要跟從我，就當捨己，背起自己的十字架來跟從我。 \end{tabularx} \\ \\ \relax
8:35 & \begin{tabularx}{0.7\textwidth}{X} 因為凡要救自己生命的，必喪失生命；凡為我和福音喪失生命的，必救自己的生命。 \end{tabularx} \\ \\ \relax
8:36 & \begin{tabularx}{0.7\textwidth}{X} 人就是賺得全世界，賠上自己的生命，有甚麼益處呢？ \end{tabularx} \\ \\ \relax
8:37 & \begin{tabularx}{0.7\textwidth}{X} 人還能拿甚麼換生命呢？ \end{tabularx} \\ \\ \relax
8:38 & \begin{tabularx}{0.7\textwidth}{X} 凡在這淫亂罪惡的世代，把我和我的道當作可恥的，人子在他父的榮耀裡與聖天使一同來臨的時候，也要把那人當作可恥的。」 \end{tabularx} \\ \\
[1ex]
\hline
\hline
\end{longtable}
講題:「性價比」與「等價交換」 “Cost Performance Ratio” and “Exchange of Equal Value”

講員:高國雄牧師

經文:馬可福音 8:27-38

場合:香港中文大學崇基學院禮拜堂主日崇拜

日期:2024年9月15日

馬可福音 8:27-38

27耶穌和門徒出去,往凱撒利亞．腓立比附近的村莊去.在路　上,他問門徒:「人們說我是誰?」

28他們對他說:「是施洗的約翰；有人說是以利亞；又有人說是先知中的一位.」

29他又問他們:「你們說我是誰?」彼得回答他:「你是基督.」

30於是耶穌切切地囑咐他們不可對任何人說起他.

31從此,他教導他們說:「人子必須受許多的苦,被長老,祭司長和文士棄絕,並且被殺,三天後復活.」

32耶穌明白地說了這　話,彼得就拉著他,責備他.

33耶穌轉過來看著門徒,斥責彼得說:「撒旦,退到我後邊去!因為你不體會上帝的心意,而是體會人的意思.」

34於是他叫眾人和門徒來,對他們說:「若有人要跟從我,就當捨己,背起自己的十字架來跟從我.

35因為凡要救自己生命的,必喪失生命；凡為我和福音喪失生命的,必救自己的生命.

36人就是賺得全世界,賠上自己的生命,有甚麼益處呢?

37人還能拿甚麼換生命呢?

38凡在這淫亂罪惡的世代,把我和我的道當作可恥的,人子在他父的榮耀裡與聖天使一同來臨的時候,也要把那人當作可恥的.」

「人們說我是誰?」

心理學中面對身份危機時,在我們心底最常有的疑問,這疑問對自身的定位有著重大的意義.

耶穌也問,為什麼他要問呢?

耶穌當然非常清楚自己的身份,他是怕門徒不知道,不認識…

「是施洗的約翰；有人說是以利亞；又有人說是先知中的一位.」

施洗的約翰:形象鮮明,不怕強權,伸張正義,上帝的代表/代言,叫人悔改,實為先鋒.

以利亞:偉大的先知,為上帝出戰,戰勝別神(巴力,亞舍拉),大有力量.

很明顯,當時群眾認定耶穌能回應他們的期盼,他們所認知當下的需要:滿有能力,願意正面面對強權/惡勢力,戰鬥到底的一位「拯救者」.

耶穌沒有對眾人對他的理解表示贊同與否,只轉向問門徒,很明顯,他期望門徒有不同的答案.

我們作為追求信仰的人,與一般坊眾,對耶穌是誰這問題,是否也有著不同的答案呢?

彼得回答他:「你是基督.」

彼得答出一個「驚人」的答案…

馬太福音16:17 耶穌回答他說:「約拿的兒子西門,你是有福的!因為這不是屬血肉的啟示你的,而是我在天上的父啟示的. 18我還告訴你,你是彼得,我要把我的教會建造在這磐石上,陰間的權柄不能勝過它. 19我要把天國的鑰匙給你,凡你在地上所捆綁的,在天上也要捆綁；凡你在地上所釋放的,在天上也要釋放.」

耶穌是基督,不是約翰,不是以利亞,也不只是一位先知,一位偉人…

基督=彌賽亞=猶太人一直等待的那位

耶穌是基督,是拯救

於是耶穌切切地囑咐(原文:命令)他們不可對任何人說起他.

一個當時不可對任何人說起的身份,時候未到…

V31從此,他教導他們說:「人子必須受許多的苦,被長老,祭司長和文士棄絕,並且被殺,三天後復活.」

從此…耶穌向門徒所教導的重點不再一樣…

基督=必須受許多的苦,被長老,祭司長和文士棄絕,並且被殺,三天後復活.

這是耶穌所理解,並將要做的事:耶穌版的「我是誰」

門徒(那位答中「耶穌是基督」的彼得)不明白,不接受,不相信,覺得沒有可能,覺得耶穌錯了,因而「責備」耶穌…

V32耶穌明白地說了這話,彼得就拉著他,責備他. 33耶穌轉過來看著門徒,斥責彼得說:「撒旦,退到我後邊去!因為你不體會上帝的心意,而是體會人的意思.」

彼得大錯特錯,更被斥責(責備同一字)為「撒旦」(敵對者)…

這錯誤非常嚴重,是致命性的錯誤!

高的性價比:今天很多人追求的生活目標

以前只講是否「抵」,

現在集中高性能(多功能),低價格(少付出)

心態:不勞而獲vs 一勞永逸

不是偷懶,不是貪心,只是「合理」

追求高性價比=我有理,我合理,應該的,聰明的…

對信仰呢?在信仰內的價值觀,是否同樣的追求呢?

我們常常努力地計算,我在信仰中所付出的,是否與所得到/獲得的合比例(我期望的比例)

決志 > 得天堂入場券 ?

返教會 > 生活順利 ?

奉獻金錢 > 賺回更多 ?

如果經驗與期望不同,就埋怨「基督教信仰/教會」性價比低?甚至要重新考慮,重新選擇?

近年來另一常提到的就是「等價交換」

記得廿年前的日本漫畫「鋼之練金術師」:

「人沒有犧牲的話就什麼也得不到,為了得到某些東西,就必須付出同等的代價,這就是煉金術的「等價交換」原則.

沒有伴隨傷痛的教訓根本沒甚麼意義,因為人類若是不作任何犧牲就不能得到任何收穫.但是當忍過這份痛楚,並且戰勝它之後,人就將獲得不屈服於任何事物的堅強的心.沒錯,如鋼鐵一般的心.」

耶穌說:「V35因為凡要救自己生命的,必喪失生命；凡為我和福音喪失生命的,必救自己的生命. 36人就是賺得全世界,賠上自己的生命,有甚麼益處呢? 37人還能拿甚麼換生命呢?」

耶穌提到人最需要得到的,人必須要有「生命」才可以繼續為人,沒有任何東西可以換取,全世界也不可以,只有捨棄「生命」,才能夠獲取「生命」!

人生中最重要的就是得到「生命」,這亦是信仰最重要的追求,所以信仰中的等價交換,不是我用什麼/做什麼而可以換取,「生命」的等價就只有「生命」

唯一能夠等價交換的就是以耶穌基督的「生命」,來換取我們的「生命」!

羅馬書6:23因為罪的工價乃是死；惟有上帝的恩賜,在我們的主基督耶穌裡,乃是永生.

耶穌活的生命為我們捨棄,我們必死的生命因此得生!

我們的性價比追求自己賺得最多,耶穌基督的性價比就是為你我付出最多!

V34於是他叫眾人和門徒來,對他們說:「若有人要跟從我,就當捨己,背起自己的十字架來跟從我.

這就是耶穌對我們的吩咐,期望.

保羅的回應:加拉太書2:20現在活著的不再是我,乃是基督在我裡面活著；並且我如今在肉身活著,是因信上帝的兒子而活；他是愛我,為我捨己. 21我不廢掉上帝的恩；如果義是藉著律法而獲得,那麼基督就白白死了.



\newpage

\section{苦苦單戀}
\label{sec:715}
\textbf{鄧瑞強博士 - 苦苦單戀}
\newline
\newline
\hyperref[sec:714]{< < < PREV SERMON < < <}
~
\hyperref[sec:toc]{[返目錄]}
~
\hyperref[sec:716]{> > > NEXT SERMON > > >}


講題:苦苦單戀 A Bitter Love

經文:何西阿書4:11-14及11:8-11

講員:鄧瑞強博士

場合:香港中文大學崇基學院禮拜堂主日崇拜

日期:2024年9月22日

各位弟兄姊妹,早安.願神的恩典常與你們同在.

過去幾個月,我們一起看以賽亞書,耶利米書,耶利米哀歌,以西結書,但以理書,習慣上,我們稱這幾卷書為「大先知書」.今日,我們看何西阿書.由何西阿書開始,到舊約聖經最後一卷書瑪拉基書,總共有12卷書,習慣上,我們稱之為「十二小先知書」.稱為為「小」,大概是因為它們每卷的篇幅較小,感覺上也覺得它們沒那麼重要.這稱呼和感覺,是有點誤導的.事實上,在希伯來聖經裡,這12卷先知書是合在一起而成一整體的,書名就叫做「十二先知書」,各卷不是獨立的.它的總長度,比以西結書還要長.12個小矮人疊起來,差不多可與偉大的以賽亞比高.有時,讀這12卷先知書,你會覺得某些書卷「無頭無尾」,這是因為我們將這些書卷從一整體裡切割開來所致.這情況類似:將一戲曲切割成12段,某些段落能自成體系,成為折子戲；但某些段落因為沒有了前因,也失落了後果,看起來便有點突兀了.

這12卷先知書的編排,不是按時序的,看來是按文學結構或神學思想而編排的.根據什麼原則來編排,在學術上,仍在探索中.12卷先知書的成書年代,不完全明確,粗略看,可能橫跨400年,由以色列國和猶大國仍然健在的日子,到經歷以色列國被亞述所滅,猶大國被巴比倫所滅；然後,在波斯帝國容許下,猶大人回耶路撒冷重建；到最後的寫作日子,可能去到希臘帝國的早期.在這悠長的400年裡,不是只有那些赫赫有名的「大先知」在指引歷史,原來還有無數「小先知」,伴隨著神的子民,在種種人生小事上,講述神的要求.這400年間,當然不只有這12個先知,抽12個出來而成一卷書,明顯與神的子民的數目「12」有關.如此編排,其背後的神學用意便很明顯了:在悠長歲月裡,神與祂的子民同在,無論他們在歷史裡經歷什麼風風雨雨,神都和他們同行.神安排了很多名不經傳的小僕人,在種種小事上,作出微小的提示.無論你在何方,神的僕人總在身邊.神與我們一起,越過驚濤駭浪的歲月.

雖然「十二先知書」是連成一體的,但為方便講道,我仍將會逐卷來講.今日講何西阿書.先知何西阿生活的年代,比先知以賽亞還要早.他生活於北國以色列最鼎盛的年代,卻目擊以色列人最終消亡在亞述的手上.他看到以色列人在鼎盛時期的奢華,看到亞述帝國興起時以色列人的人心惶惶.在奢華的年代,人顧住宴樂而忘記神.在敗落的年代,人顧住用自己的方式自保而忘記神.無論平安抑或危難,神總是記掛著祂的子民；但無論平安抑或危難,祂的子民總是忘記神.先知好傷感,他也體會到神的傷感.先知無法壓抑自己,他盡情傾吐傷感,那種與神同感的傷感.他呼喚人活得清醒,恢復靈性,回歸對神的忠誠.

我們看看先知何西阿怎樣描寫神的子民的靈性光景.

我們看何西阿書4:11-14.讀的時候,留意「淫」這個字.在全本聖經裡,這是「淫」及相關意象出現頻率最高的段落.

何4:11　姦淫和酒,並新酒,奪去人的心(take away one’s understanding).

何4:12　我的民求問木偶,以為木杖能指示他們；因為他們的淫心使他們失迷,他們就行淫離棄神,不守約束,

何4:13　在各山頂,各高岡的橡樹,楊樹,栗樹之下,獻祭燒香,因為樹影美好.所以【因此之故】,你們的女兒淫亂,你們的新婦行淫.

何4:14　你們的女兒淫亂,你們的新婦行淫,我卻不懲罰她們；因為你們自己離群與娼妓同居,與妓女一同獻祭.這無知的民(people without understanding)必致傾倒.

讀完這段經文,我想,大家都會同意,這段經文的主旨,一言以蔽之:「淫」,這就是神的子民的光景.

這段短短的經文,有一簡單的平行結構:

頭尾平行:4:11講到姦淫「奪去人的心」,人不再清醒；4:14下半節講到淫亂使人成為「無知的民」.不再清醒,無知,都是活得毫無靈性.對神無知,對生活的意義無知,對自己的人性無知,只追求當下的歡愉,活得像動物一樣.這就是「淫」,這是「對神無知的人」的總體描述.

接著,是4:12-13上半節與4:14上半節的平行.4:12-13上半節講到「求問木偶」,在樹下燒香.4:14上半節講到他們「與妓女一同獻祭」.這是對神無知的人的宗教表達,他們四處拜不同的神明.淫亂,就是不忠誠.終極的不忠誠,就是耶和華的子民根本不信耶和華,而投靠別的神明.本來,神的子民有耶和華,已一無所缺.「耶和華是我的牧者,我必不致缺乏」.有了創造萬有的耶和華,神的子民還欠什麼呢?有了掌握萬有權柄的耶穌,基督徒還欠什麼呢?但神的子民就是偏偏不信神,內心充滿各種不安,於是,去到不同神廟,都入去拜一拜以求心安；每周又看看星座運程,以求趨吉避凶；閒來又擺些風水陣,好得些橫財.用盡各種方法去求心安,但偏偏就是不將生命交給神.這就是「淫」,這根本就是不信神.

這段平行經文的核心,就是4:13下半節:「【因此之故】,你們的女兒淫亂,你們的新婦行淫.」這裡描繪生活總體的腐爛,已爛到入骨髓.古時,是男性主導的社會,有嚴格的榮辱(honour and shame)結構.男性要守護女性的貞潔,這是家族的榮譽所在.若家中的女兒,甚或自己的妻子淫亂,就意味無法守住生活的基本秩序,家庭結構的瓦解,傳統價值的失落,家族的尊嚴蕩然無存.在先知眼中,因神的子民不好好信靠神,生命的方方面面都崩壞解體了.

當亞述的威脅越來越大時,以色列人一方面繼續夜夜笙歌,若不及時行樂,他們怕將來想行樂也沒有機會了；另方面他們胡亂地尋求幫助,為求心安,拜不同的神明,倚靠不同的勢力.但真奇怪,以色列人偏偏就是不信任耶和華,這令神很傷心.

耶和華與祂的子民,有一種盟約關係,若神的子民忠誠於這關係,則神會帶領他們度過種種難關.而何西阿是第一個先知,以性和婚姻的意象去表述這盟約關係.神是丈夫,祂的子民就是祂的妻子,祂深愛祂的妻子,忠誠於這婚姻關係.開了這先河後,後來不少先知都繼續沿用婚姻意象去表述神人關係.耶穌來世,也將自己看為是新郎,而信徒群體就是新娘.到最後,在啟示錄裡,耶穌要和信徒群體完成一「羔羊的婚禮」,成就神人的永久連繫.

在何西阿書裡,神以忠誠的愛,去愛祂的子民.祂的子民卻以不忠的「淫」,去回報祂的神.一忠誠,一淫亂,在何西阿書裡,成了極大的對比.為了讓何西阿體會「忠誠」遇上「淫亂」的痛苦,神叫何西阿娶一個淫婦為妻.這女子以「淫亂」為本性,先知卻以忠貞的愛待她.可以想像,這關係有多複雜.這女子一次又一次的紅杏出牆,先知一次又一次的接納寬恕.先知終於明白,面對死不悔改的人,無條件的接納是多麼艱難,等待她回轉是多麼渺茫,獨自一人忠貞地愛是多麼叫人絕望.在這不忠貞的人間,先知堅持愛,終於,他明白這就是神的心.神在高天,默默地注視世人,注視祂的子民,他們快樂時忘記祂,天天領受祂的恩典卻總是忘恩負義,艱難時卻尋求各種偶像的幫助.神依戀地望著世人,世人卻貪婪地注視世界.一個苦苦地單戀的上帝,祂看著祂的子民漠視祂,離棄祂,愛著別的東西.一個苦戀的上帝,一個被人離棄的上帝,一個流淚的上帝.

我們聽聽神的心聲.

何11:8　以法蓮哪,我怎能捨棄你?以色列啊,我怎能棄絕你?我怎能使你如押瑪?怎能使你如洗扁?我回心轉意,我的憐愛大大發動.

「以法蓮哪,以色列啊」,像是丈夫呼喚著不貞的妻子.妻子已離去,但這苦苦單戀的上帝卻說:「我怎能捨棄你?我怎能棄絕你?」聽起來,像是一個鍾情的男子,不顧愛人千次萬次的離棄及出賣,仍保持著忠貞的愛情.無論你對我怎樣,我對你的愛一如初戀.這個男子,就算遭愛人遺棄,仍孤獨地,奮力擔起愛的十架.

「我怎能使你如押瑪?怎能使你如洗扁?」押瑪和洗扁,在所多瑪,蛾摩拉附近,這是古時天火懲罰的兩座城巿.神壓抑著祂無限的毀滅能力,承受著被愛人背叛的屈辱.沒有恨,只有愛.經文說:「我回心轉意,我的憐愛大大發動」.你們不回心轉意,我自己回心轉意.你們不願意愛我,我主動愛你們.你們不願意回到我身邊,我主動去到你們身邊.在我們的世界,被愛人拋棄後,愛變成恨的例子很常見,但神拒絕這樣做.祂反轉曹操的講法,祂說:「寧教天下人負我,我絕不負天下人.」神決心去愛,就愛到底.

在古代,其實多是被遺棄的怨婦,等候遲遲不歸的丈夫,那有被遺棄的丈夫,等候根本不願歸的妻子的道理.妻子棄夫而去,早就被休了.妻子紅杏出牆,在古代以色列,早就被石頭打死了.如今,卻是超乎想像的,大能的神自甘卑微,忍受羞辱,沒有休棄淫亂的妻,沒有置她死地,卻成為一等待愛人回歸的怨夫,祂堅持愛,堅持得令人吃驚.

何11:9　我必不發猛烈的怒氣,也不再毀滅以法蓮.因我是神,並非世人,是你們中間的聖者；我必不在怒中臨到你們.

在我們的想像裡,創造萬有的神疾惡如仇,世人對祂不忠,祂便降災懲罰.在極端的情況下,神的確會這樣做的,但這絕不是神的本性.若神真要懲治罪惡,下一刻這世界便歸於無有了.神的本性是愛,祂寂寞地堅持愛,縱使我們不理祂,祂仍選擇苦苦單戀下去.當世人以恨還愛,以暴力取代和平,以決裂代替和解時,神寧願彎下腰,替人洗腳,創造仁愛.這就是我們相信的神,這就是真正的聖者,祂向我們展現一種毫無理由的愛.祂愛我們,不為什麼,只因祂是神,我們是祂的戀人.縱使我們背叛祂,祂仍堅持愛我們.沒有理由,祂就是愛我們.

何11:10　耶和華必如獅子吼叫,子民必跟隨他.他一吼叫,他們就從西方急速而來.

何11:11　他們必如雀鳥從埃及急速而來,又如鴿子從亞述地來到.我必使他們住自己的房屋.這是耶和華說的.

獅子,是亞述的神獸,以色列國最終被亞述毀滅,像被獅子吞噬一樣.但耶和華像獅子王,祂吼叫,宣示自己對萬有的主權.祂吼叫,呼喚祂的子民從被擄之地歸回.祂吼叫,像痴情的丈夫,呼籲祂的愛人返回家園.經文最後說:「這是耶和華說的」,這是神以其永恆的神性下終極的決定.神好像是說:我知我的子民水性楊花,毫不忠誠；我知我的子民追逐世間的榮華富貴,一點不放我在眼內；我知我的子民完全不渴望生命的永恆意義,活得朝生暮死；我知我的子民崇拜我的時候,漫不經心.我知我的子民根本一點不愛我.這一切,我都知道.但我依然愛他們,這是我說的.

各位弟兄姊妹,神這樣說了,你又如何說呢?

願上主的大愛臨到你,願你回應上主的愛,得享滿足,平安.



\newpage

\section{信心這回事}
\label{sec:716}
\textbf{關瑞文教授 - 信心這回事}
\newline
\newline
\hyperref[sec:715]{< < < PREV SERMON < < <}
~
\hyperref[sec:toc]{[返目錄]}
~
\hyperref[sec:717]{> > > NEXT SERMON > > >}


雅各書 5:13-20
\newline
\begin{longtable}{cl}
\hline
\hline
章節 & 經文 (和合本修訂版)\\
\hline
5:13 & \begin{tabularx}{0.7\textwidth}{X} 你們中間若有人受苦，他該禱告；有人喜樂，他該歌頌。 \end{tabularx} \\ \\ \relax
5:14 & \begin{tabularx}{0.7\textwidth}{X} 你們中間若有人病了，他該請教會的長老們來為他禱告，奉主的名為他抹油。 \end{tabularx} \\ \\ \relax
5:15 & \begin{tabularx}{0.7\textwidth}{X} 出於信心的祈禱必能救那病人，主必叫他起來；他若犯了罪，也必蒙赦免。 \end{tabularx} \\ \\ \relax
5:16 & \begin{tabularx}{0.7\textwidth}{X} 所以，你們要彼此認罪，互相代求，使你們得醫治。義人祈禱所發的力量是大有功效的。 \end{tabularx} \\ \\ \relax
5:17 & \begin{tabularx}{0.7\textwidth}{X} 以利亞與我們是同樣性情的人，他懇切地祈求不要下雨，地上就三年六個月沒有下雨。 \end{tabularx} \\ \\ \relax
5:18 & \begin{tabularx}{0.7\textwidth}{X} 他又禱告，天就降下雨來，地就有了出產。 \end{tabularx} \\ \\ \relax
5:19 & \begin{tabularx}{0.7\textwidth}{X} 我的弟兄們，你們中間若有人迷失了真理而有人使他回轉， \end{tabularx} \\ \\ \relax
5:20 & \begin{tabularx}{0.7\textwidth}{X} 這人該知道，使一個罪人從迷途中回轉，會從死亡中把他的靈魂救回來，而且遮蓋許多的罪。 \end{tabularx} \\ \\
[1ex]
\hline
\hline
\end{longtable}
講題:信心這回事 This Thing Called Faith

經文:雅各書5:13-20

講員:關瑞文教授

場合:香港中文大學崇基學院禮拜堂主日崇拜

日期:2024年9月29日

弟兄姊妹,早晨好!很高興今日又來到大家當中,和大家一起崇拜.今日我們要看雅各書第5章13節至20節.其實整本雅各書都相當重要,當中談到一個很重要的信仰話題,就是究竟甚麼是「真正」的信心?

究竟甚麼是真正的信心呢?閱讀雅各書這段第5章的經文,或許會讓我們聯想起,信心就是帶著力量的禱告.當我們的心懷著對上帝全然的信靠,當我們向祂禱告時,上帝必然會因著我們大有信心而垂聽我們的祈求.

不是嗎?翻開雅各書第5章13-18節,經文這樣教導我們:你們中間有受苦的嗎?他就應該禱告.接著,經文繼續說道:你們中間有病了的嗎?他就該請教會的長老來,他們可以奉主的名用油抹他,為他禱告.出於信心的祈禱,必能救那病人,主必叫他起來；他若犯了罪,也必蒙赦免.這段經文清楚地告訴我們,出於信心的禱告擁有醫治的大能,因為主必叫他起來,並且赦免他的罪.第16節更進一步強調:所以你們要互相代求,使你們可以得醫治.義人祈禱所發的力量是大有功效的.經文還引用了先知以利亞的例子.以利亞憑著對上帝堅定的信心,懇切地禱告,求上帝使天閉塞不下雨.三年零六個月,地土乾旱,滴雨未降.之後,他又再次向神祈求,天就降下雨來,地也生出土產.這真是太神奇了!

由此看來,似乎只要我們有信心,上帝就會答應我們一切的祈求,甚至連呼風喚雨都能做到,因為上帝是全能的上帝,如果我們全然信靠,就天衣無縫,萬事俱備,不欠東風!

我記得很久以前,我曾經聽過一位資深的傳道人分享他一次去醫院探訪的經歷.當天,他有一位姊妹教友的媽媽病得很嚴重,醫生說媽媽大概很快就會離世.但這位姊妹還有很多話未跟媽媽說,很不願意見到媽媽就這樣離去,沒機會與媽媽訴說最後的話.於是就請傳道人跪下,為她媽媽禱告,求上帝多給她媽媽一晚的時間,讓她可以在媽媽離世之前,可以最後一次傾訴.當時傳道人心裡面想,其實她媽媽已經昏迷,並且生命跡象都非常微弱,醫生亦都說她媽媽不會再醒過來,所以他根本不相信禱告之後她媽媽會醒過來和姊妹說話.但是眼見這位姊妹如此渴求,如此有信心,就和她一起跪下,為她媽媽和她們兩母女禱告,懇求天父可以讓她媽媽多留一天.意想不到的是,在他們禱告之後大約一小時,媽媽真的醒過來了,以致在媽媽過身之前,她們兩母女真的可以互訴衷腸.親眼見證這個神蹟,當時傳道人心裡面好自責,怪責自己為甚麼那麼沒有信心,也責怪自己的信心比不上這位姊妹.他當時說如果讓他再來一次,他一定會帶著更大的信心,相信上帝的神蹟,為這位姊妹和她媽媽禱告.

然而,信心果真是如此嗎?信心是否等同心想事成的萬能鑰匙?是否每次我們為垂危病人禱告,都必能起死回生?如果是這樣,是不是要說,信心可以使人長命百歲,就地永生?不會是這般吧!但雅各卻的確在第15節那裡說出於信心的禱告,「必」能救那病人,主「必」叫他起來.難道雅各在信口雌黃,假傳聖旨?也不會吧!那麼,有關信心的事情,到底雅各書想要告訴我們什麼呢?

究竟甚麼是信?你可能會想起我們每個禮拜都背誦的《使徒信經》:「我信上帝,全能的父......我信我主耶穌基督,上帝獨生的子......我信聖靈,我信聖而公之教會,我信聖徒相通,我信罪得赦免,我信身體復活,我信永生.」當我們以為我們已經答對了得到100分的時候,雅各書在2章19節卻這樣回應:「你信神只有一位,你信得很好；連鬼魔也信,卻怕得發抖.」雅各這個回應真是令人太詫異了!

然而,其實當我們了解雅各書的背景的時候,我們就不會太過詫異於他這樣的回應.雅各書是一本成書非常早期的一卷聖經書卷.相信大家都知道,在相當早期的教會,保羅向外邦人的宣教運動尚未成形,所以教會的弟兄姊妹主要都是猶太背景的基督跟隨者,而雅各當然也是一位擁有濃厚猶太傳統背景的信徒.有不少釋經學者都認為,寫成雅各書主要部分的那位作者,就是耶穌的兄弟雅各,他亦是在耶穌受死,復活,升天之後十年之內成為了耶路撒冷教會的領袖.所以,可以這樣說,雅各書是猶太基督跟隨者對著猶太基督跟隨者說話的.這一點對我們要把握雅各書當中所講的「信心」是很重要的,是很關鍵的.

有關「信」這回事,猶太人的思想和用語習慣到底是怎樣的呢?在猶太人的用語裡面,「信」是“Emunah”.“Emunah”所講的絕對不是就等於頭腦裡面的信念.我意思是,當猶太人講到「信」的時候,其實他們不是在講「我信上帝只有一位」,「我信天父」,「我信主耶穌基督」,「我信聖靈」.在猶太人的思想和用語習慣裡面,「信」其實是一種行動,是身體力行的一種行動.所以,有學者指出,雖然「信」這個字是名詞,但是在猶太傳統用語裡面,「信」這個字,是指涉行動的,因此可以用動詞的概念來理解這個「信」字.所以,雅各這樣說:

2.14我的弟兄們,若有人說自己有信心,卻沒有行為,有甚麼益處呢?這信心能救他嗎

2. 20你這虛浮的人哪,你願意知道沒有行為的信心是沒有用( [根本唔係野來的],[邪惡])的嗎?

2.24這樣看來,人稱義是因著行為,不是單因著信.

弟兄姊妹,我們必須留意到,在雅各書裡面,真正的信心其實是行動.但我們必須追問,所謂行動,雅各所指的是甚麼行動呢?

我記得,差不多二十多年前,有一晚我去了一次此生難忘的宵夜.當時我是神學樓學生宿舍的舍監,經常和學生一起去宵夜,因此養成了一種宵夜的興趣.有一晚,「gwo 條宵夜筋又有發作」,所以「夜媽媽」11時許,即興和家人開車朝向沙角村的那個大排檔吃宵夜.不到10分鐘,車就停泊在沙角村街市那道側門前.下車時,見到一位年事已高的老婆婆,她一隻眼睛嚴重發炎,有水滲出.彎著早已佝僂的腰,當時她正在街市外面,垃圾桶旁,拾取棄置地上那些白天街市沒有賣完的霉霉爛爛的蔬菜及碎肉.見此,我的心突然沉了下來.「朱門酒肉臭,路有凍死骨!」我不愁溫飽,不飢不餓,即興去大魚大肉,而婆婆即使肚餓也只能吃這些爛菜碎肉.當時我走過去婆婆身邊,勸她不要吃這些蒼蠅亂飛的食物.她仍然彎著腰,緩緩地抬起頭來,面露著一種令我內疚的笑容,對我說:「這些還可以吃,還可以吃.」那輕描淡的抬頭,那老人家虛弱的眼神,那句「還可以吃,還可以吃」,剎那間就化為一種「他者」的呼喚；這張滿滿皺紋的面容,剎那間就化為「他者」的敲問!呼喚什麼?敲問什麼?是那張受苦的面容,在敲問為什麼在這面容下面的上帝形象和樣式,會扭曲得如此沉默無聲.頃刻,我的心不只突然感到很難過,更是為到自己來吃宵夜的決定感到很羞恥.在資本主義的遊戲規則下,正當我們享受著職業帶來的物質生活時,想不到原來自己已間接成為了社會資源分配不均的幫兇,間接讓有需要的人連溫飽都成為每天的難題,讓人受苦.朱門酒肉臭,路有凍死骨.當晚,我只敢卑卑微微地走近婆婆身旁,問候她的近況,了解她的需要,最後給了她我身上所有的錢.當時,我絕對不是出於「助人為快樂之本」的心,更加不是想起「施比受更有福」；相反,我心裡只充滿了「欠債還錢,天公地道」之感受.其實,所欠婆婆的,還有很多呢!我立刻打消了宵夜的念頭,掉頭離開了.往後的日子,若不是為了公事,我再也沒有去吃宵夜了!

弟兄姊妹,我們必須留意到,在雅各書裡面,真正的信心其實是行動.所謂行動,雅各所指的是甚麼行動呢?

2.15-17我的弟兄們,若有人說自己有信心,卻沒有行為,有甚麼益處呢?這信心能救他嗎?若是弟兄或是姊妹沒有衣服穿,又缺少日用的飲食；你們中間有人對他們說:「平平安安地去吧!願你們穿得暖,吃得飽」,卻不給他們身體所需要的,這有甚麼益處呢?信心也是這樣,若沒有行為是死的.

1.27在上帝—我們的父面前,清潔沒有玷污的虔誠就是看顧在患難中的孤兒寡婦.

雅各提醒我們,在沙角村街市遇到婆婆,被她那受苦的面容呼喚,而開口問候,以及掏出自己腰包給她身上所有的錢,這絕對不是一種慈善行為,也不是因為信了耶穌,就努力聽從上帝的教導,學習服侍有需要的人；見到弟兄或姊妹沒有衣服穿,缺少日用的飲食,就將我們所有的分給他們,這絕對不是一種慈善行為,也不是因為信了耶穌,就努力聽從上帝的教導,學習服侍有需要的人；見到在患難中的孤兒寡婦,就去照顧,這絕對不是一種慈善行為,也不是因為信了耶穌,就努力聽從上帝的教導,學習服侍有需要的人.這全都因為看見受苦者的面容,就看見那扭曲了的上帝形象.

這好比孟子曾說︰「…今人乍見孺子將入於井,皆有怵惕惻隱之心,非所以內交於孺子之父母也,非所以要譽於鄉黨朋友也,非惡其聲而然也.由是觀之,無惻隱之心,非人也 …」什麼意思呢?孟子在說︰任何人突然看到小孩子要掉到井裡,都會感到驚嚇和同情,情不自禁地伸手援救.這並不是因為要得到什麼好處,或避免任何損失,而是因為人皆有惻隱之心,這種不忍之心是發自內心的,當下被天呼喚要成為一個天所設定的真人.

如果所謂「信」就是與上帝修復關係,那麼,被那埋藏在「他者」的受苦面容皮骨底下的上帝呼喚及拷問,以致給婆婆身上所有錢,以及給予衣服及食物有需要的人,並且照顧在患難中的孤兒寡婦,本身就是所謂「信」.所謂「信」就是與「他者」上帝相遇時的行動.這些都不是「信」所產生出來的行動,這些行動本身就是「信」,這是雅各要告訴我們的.

所以讀到第5章13節開始,雅各說:你們中間若有人受苦,就該禱告；你們中間若有人病了,就應該請教會的長老來為他禱告,奉主的名為他抹油.出於信心的祈禱,必能救那病人,主必叫他起來；他若犯了罪,也必蒙赦免.這裡的重點不是在講論祈禱的威力有多大,而是在提醒我們,當我們遇到受苦,患病,被罪惡纏繞的人的時候,我們就是遇見那扭曲了面容的上帝,而且這面容呼喚我們要作出回應,作出行動,這就是與上帝重建關係,這才是真正的「信」.所以雅各說:

• 2.17「信心也是這樣,若沒有行為是死的.」

• 2.18「但是有人會說:『你有信心,我有行為.』你把沒有行為的信心給我看,我就藉著我的行為,把我的信心給你看.」

• 2.22-24「可見,信心是與他的行為並行,而且信心因著行為才得成全.這就應驗了經上所說:『亞伯拉罕信上帝,這就算為他的義.』他又得稱為上帝的朋友.這樣看來,人稱義是因著行為,不是單因著信.」

雅各不是要辯論到底人稱義是「因信稱義」還是「因行為稱義」,他只是想要提醒我們,「信」就是那些被「他者」上帝面容呼喚而產生的行動.

雅各書很豐富,當中也提到其它「信」的行動,譬如要勒住舌頭,不為富不仁,不追求這世界所追求的等等,很值得我們慢慢細讀.

但總結而言,我們可以問自己:真的「我信上帝只有一位,信耶穌得永生」.不過,我們的「信」,是真的「信」嗎?

2.18「但是有人會說:『你有信心,我有行為.』你把沒有行為的信心給我看,我就藉著我的行為,把我的信心給你看.」



\newpage

\section{人算甚麼}
\label{sec:717}
\textbf{李均熊牧師 - 人算甚麼}
\newline
\newline
\hyperref[sec:716]{< < < PREV SERMON < < <}
~
\hyperref[sec:toc]{[返目錄]}
~
\hyperref[sec:718]{> > > NEXT SERMON > > >}


詩篇 8:1-9
\newline
\begin{longtable}{cl}
\hline
\hline
章節 & 經文 (和合本修訂版)\\
\hline
8:1 & \begin{tabularx}{0.7\textwidth}{X} 耶和華－我們的主啊，你的名在全地何其美！你將你的榮耀彰顯於天。 \end{tabularx} \\ \\ \relax
8:2 & \begin{tabularx}{0.7\textwidth}{X} 你因敵人的緣故，從孩童和吃奶的口中建立了能力，使仇敵和報仇的閉口無言。 \end{tabularx} \\ \\ \relax
8:3 & \begin{tabularx}{0.7\textwidth}{X} 我觀看你手指所造的天，並你所陳設的月亮星宿。 \end{tabularx} \\ \\ \relax
8:4 & \begin{tabularx}{0.7\textwidth}{X} 人算甚麼，你竟顧念他！世人算甚麼，你竟眷顧他！ \end{tabularx} \\ \\ \relax
8:5 & \begin{tabularx}{0.7\textwidth}{X} 你使他比神微小一點，賜他榮耀尊貴為冠冕。 \end{tabularx} \\ \\ \relax
8:6 & \begin{tabularx}{0.7\textwidth}{X} 你派他管理你手所造的，使萬物， \end{tabularx} \\ \\ \relax
8:7 & \begin{tabularx}{0.7\textwidth}{X} 就是一切的牛羊、田野的牲畜、空中的鳥、海裡的魚，凡游在水裡的， \end{tabularx} \\ \\ \relax
8:8 & \begin{tabularx}{0.7\textwidth}{X} 都服在他的腳下。 \end{tabularx} \\ \\ \relax
8:9 & \begin{tabularx}{0.7\textwidth}{X} 耶和華－我們的主啊，你的名在全地何其美！ \end{tabularx} \\ \\
[1ex]
\hline
\hline
\end{longtable}
講題:人算甚麼 What Is Man?

經文:詩篇8:1-9

講員:李均熊牧師

場合:香港中文大學崇基學院禮拜堂主日崇拜

日期:2024年10月6日

起:上主的威榮

今日啟應詩篇,是很多人喜愛的詩篇第八篇.但是詩篇第八篇的主旨是甚麼呢?詩人所探索的是甚麼題目呢?大部份研讀這詩的,都會被中間問句「人算甚麼,你竟顧念他?世人算什麼,你竟眷顧他?」吸引,不其然就會以為這詩定是講論人的存在價值,繼而大談存在的危機,存在主義等等,變得很哲學.接著第6-8節:「你派他管理你手所造的,使萬物,就是一切的牛羊,田野的獸,空中的鳥,海裡的魚,凡經行海道的,都服在他的腳下.」又不其然想現代人處於全球暖化,生態危機的困境中,轉而覺得這詩是講人對護理大地的責任,然候就開始討論人在自然生態環境中的角色和使命,變得很環保.但是這詩的主題真是環繞人的存在價值和環保角色嗎?

當我們細閱這詩是,不要忘記這是詩篇內的一首讚美詩,我們不要忽略開頭幾節對上主的讚頌.可以留意詩開首與結尾詩句是完全一樣的:「耶和華我們的主啊,你的名在全地何其美!」李思敬選用「威嚴」而不用「美」字,反映的可能就是這詩不是因為大自然的美麗而想到上主的「美」(可惜的是李思敬省去了「何其」那設問口吻,因為這個設問正好呼應中間「人算甚麼」的設問).

詩人並不是向第三者陳述上主的威榮,而是向著上主「你」而發出.所以這詩並非一般的讚美詩.詩篇內的讚美詩一般都會有頌讚上主的邀請,而後內容必然帶著原因.用詩篇最短的一首讚美詩117篇作例子:「萬國啊,你們都當讚美耶和華!萬民哪,你們都當頌讚他!因為他向我們大施慈愛；耶和華的誠實存到永遠.你們要讚美耶和華!」傳統讚美詩編排必有的讚美的邀請和讚美的原因便顯而易見.所以詩篇第8篇並不是詩人觀星時見到滿天星星而有感而發,然後邀請人一同頌讚上主(這其實更似是啟蒙運動和浪漫主義的產物).詩人是面向上主,因他所敬拜的上主無可比擬的偉大而用上天際萬像去作出比較.若是這樣,第2節出現的「孩童和吃奶的口中」去說出上主的威榮,其實是相當另人費解.可能詩人在比擬上主的威榮時,想到上主的能力並不是赤裸的權力,好像極權國家的領導,要人民透過洗腦或脅逼而要讚美國家一般.不,上主的威榮是可以透過小孩,甚至還在吃奶的嬰兒的口而傳頌.孩童與嬰兒的純真才是可以真實,毫不虛假的展現上主的威榮所在!

承:古代敬拜者的視角

要明白詩篇8篇的主旨,或許讓先我們心神穿越回到古代,投入一位進入聖殿敬拜的古以色列人,他的思想和視野是如何的,才有可能明白詩篇8篇的境界.

作為古代敬拜上主的人,一定不會關注甚麼人的存在價值或環保議題.古人生活在窮山惡水的環境中,大自然並不是一件賞心樂事,在假日郊遊享受的地方(「唔係講玩」).現代人以為大自然是美麗安舒的地方,這是因為我們根本沒有直接接觸大自然.我們今日城巿人所接觸的,極其量只能算是「郊外」,是一處被城巿化下仍被保留的「郊野公園」(有政客還說要用郊野公園土地興建樓宇).真正的野外,是不適而人居住的地方.若你沒有足夠的裝備,在「大自然」下你可能二,三天便吃不消了.古語有云:「天地不仁,以萬物為芻狗」芻狗即用草紮成的狗,用來祭祀焚燒,故此在凶險的大自然底下,人又怎會有價值(雖然中國這句古語其實是指眾生皆平等)?在聖經出現「某某算甚麼」的意思,通常就是某某其實沒有甚麼價值(例如王下8:13當以利沙說哈薛要作亞蘭王時,哈薛便說:「你僕人算什麼,不過是一條狗……」).詩篇這一個反問,不是問「人是甚麼」(What is Man),而是人「算得上」甚麼(What does a person “count”?),人真是一文不值worthless .作為古以色列人,面向自然穹蒼,其實不會如嬌生慣養的現代人,想到的並不是上主的「美麗」,而是凶險的大自然乃是臣服於威嚴的上主,因為那是由上主創造,祂說有就有的世界,是祂「指頭」的玩物.但就是這個反差,更顯這詩頌讚的這個上主的獨特性.「人算甚麼,你竟派他治理你手的作為!」

在加上古人活在一個多神宗教的文化中,詩人正是要透過宣告自己所敬拜的獨一真神的與眾不同,所以才有詩篇8篇的內容.詩中開始講到上主的敵人,那會是誰呢?熟識周遭神話的古代人,可能正在想其他神明都有一個宿敵nemesis,都要戰勝別個勢力才可宣稱對宇宙的王權,例如位於北敘利亞的古烏加列Ugarit文明的神話中,風暴之神「巴力」Baal 因要掌控天庭,就要戰勝海洋邪神「閻王」Yamｍ.但詩篇8篇的上主,並沒有與之匹敵的勢力.上主的能力,能「消滅仇敵還擊」(見李思敬翻譯),換句話說就算上主有敵人,這個敵人就根本連還擊的機會也沒有!

這個熟識古代神話的詩人,可能也同時想到這些神話是如何描寫人的「沒有」價值.古代的神話,大都是寫成「人」在神明面前只是奴隸,其價值只在為神明服務,按時提供祭品.在古代沒有民主,人民只能服從在上掌權的君王,他們對神明的觀念,就是一個絕對的君上,所有臣下都是他的奴僕,如君主的神明有絕對生殺之權,人只能仰望君主或神明的施恩.古代巴比倫的神話Atra-Hasis中,人,只是因為一些低等的神明不願服侍高級的神明,從而造出來的奴隸.最後統粹因為這些作為奴隸的人太嘈吵,所以神明就降下洪水去消滅他們!

轉:王權代理人

好了,若你今日進入一個古代敬拜者的心靈,他所身處的文化環境,只會告訢你在充滿凶險勢力的宇宙中,「人」是沒有價值的,神明只會當人是奴隸,大自然充滿災禍凶險,生活在世只是在逆境中求生存.不過,你卻又不是一般的古代人,因為你敬拜的是耶和華上主,以色列的神,獨一真主,無可比擬,不單沒有其他勢力可以威脅祂,祂更順服了凶惡的大地,只需用指頭便安排了顯示人命運的天像,更令你驚訝的,就是毫無價值的人,祂竟然會記念,會眷顧,「你又使他僅次於神明(很多譯本將之譯為天使,但在古人眼中,可能神明才是比較的對像),卻以榮耀尊貴為他冠冕」,若上主係君上,我們不單不是臣僕,我們被加以皇室「冠冕」,並且上主還派我們去治理祂所臣服的大地萬物!你來到以色列的聖殿,並不是單單讚美這一個威榮的上主,你其實在宣告耶和華上主,我們的主,是獨一至高,無可比擬的上帝,但作為敬拜祂的子民,也會是與眾不同,是皇家子弟,能協助至高神治理祂的國度!

合:回到現在

今日和古代世界已經大大不同.今人的世界宇宙觀也和昔日古人完全不一樣.因為歷史的發展,文藝復興,宗教改革,啟蒙運動,以至工業革命,可能也因為基督教對西方文明的影響,人對自然的觀感與才有今日的改變.

今日我們不用再面對窮山惡水,能在相對安全的環境生活,能夠發揮和欣賞人的價值,從歷史的角度看,是因為基督教的價值改變了人們對世界和人生的想法.人生而有價值,平等和自由,我們又會改造環境讓它更適宜人類居住,可以說都是基督教文明所帶來的發展.但我們卻慢慢忘記我們對這個宇宙的敬畏,也慢慢忘記了我們所敬拜的上主的獨特和無可比擬!詩篇8篇的主旨並不是天象,人類或受造界,這些都是映襯上主的威榮.詩篇8篇的主旨,是在再困難重重的人生中有順服宇宙百物和危險勢力,那有能力的「上主,我們的主」.

不過,無論古代或現今,我相信讀這篇詩的人都會體會到,經文展示的不單是一位在危險困難的環境中,展現威榮能力而獨特的神,更加是祂與人的獨特關係,使我們驚異.古代一般的宗教文化將人置於一個被神明奴役的身份,故此詩篇第8篇獨排眾議,告訢我們,我們所敬拜的獨一威榮的上主,不單順興宇宙百物,更視我們為治理大地的伙伴partner,祂在世的代理agent.今日我們身處的危險環境可能不再是窮山惡水,但今日我們或許要居安思危,也要留心身處「世界」其實仍然危險,但更要注目我們所信仰的耶穌基督,因著祂復活的能力為我們帶來平安.今日我們又會否令我們獨排眾議,活出上帝兒子的獨特身份,並領受與上主同工的召命呢?



\newpage

\section{錢財是永生的敵人嗎}
\label{sec:718}
\textbf{林豪恩先生 - 錢財是永生的敵人嗎}
\newline
\newline
\hyperref[sec:717]{< < < PREV SERMON < < <}
~
\hyperref[sec:toc]{[返目錄]}
~
\hyperref[sec:719]{> > > NEXT SERMON > > >}


馬可福音 10:17-31
\newline
\begin{longtable}{cl}
\hline
\hline
章節 & 經文 (和合本修訂版)\\
\hline
10:17 & \begin{tabularx}{0.7\textwidth}{X} 耶穌剛上路的時候，有一個人跑來，跪在他面前，問他：「善良的老師，我該做甚麼事才能承受永生？」 \end{tabularx} \\ \\ \relax
10:18 & \begin{tabularx}{0.7\textwidth}{X} 耶穌對他說：「你為甚麼稱我是善良的？除了神一位之外，再沒有善良的。 \end{tabularx} \\ \\ \relax
10:19 & \begin{tabularx}{0.7\textwidth}{X} 誡命你是知道的：『不可殺人；不可姦淫；不可偷盜；不可作假見證；不可虧負人；當孝敬父母。』」 \end{tabularx} \\ \\ \relax
10:20 & \begin{tabularx}{0.7\textwidth}{X} 他對耶穌說：「老師，這一切我從小都遵守了。」 \end{tabularx} \\ \\ \relax
10:21 & \begin{tabularx}{0.7\textwidth}{X} 耶穌看著他，就愛他，對他說：「你還缺少一件：去變賣你所有的，分給窮人，就必有財寶在天上；然後來跟從我。」 \end{tabularx} \\ \\ \relax
10:22 & \begin{tabularx}{0.7\textwidth}{X} 他聽見這話，臉就變了色，憂憂愁愁地走了，因為他的產業很多。 \end{tabularx} \\ \\ \relax
10:23 & \begin{tabularx}{0.7\textwidth}{X} 耶穌看了看周圍，對門徒說：「有錢財的人進神的國是何等的難哪！」 \end{tabularx} \\ \\ \relax
10:24 & \begin{tabularx}{0.7\textwidth}{X} 門徒對他的話非常驚奇。耶穌又對他們說：「孩子們，要進神的國是何等的難哪！ \end{tabularx} \\ \\ \relax
10:25 & \begin{tabularx}{0.7\textwidth}{X} 駱駝穿過針眼比財主進神的國還容易呢！」 \end{tabularx} \\ \\ \relax
10:26 & \begin{tabularx}{0.7\textwidth}{X} 門徒就更為驚訝，彼此對問：「這樣，誰能得救呢？」 \end{tabularx} \\ \\ \relax
10:27 & \begin{tabularx}{0.7\textwidth}{X} 耶穌看著他們，說：「在人不能，在神卻不然，因為在神凡事都能。」 \end{tabularx} \\ \\ \relax
10:28 & \begin{tabularx}{0.7\textwidth}{X} 彼得就對他說：「看哪，我們已經撇下一切跟從你了。」 \end{tabularx} \\ \\ \relax
10:29 & \begin{tabularx}{0.7\textwidth}{X} 耶穌說：「我實在告訴你們，凡為我和福音撇下房屋，或是兄弟、姊妹、父親、母親、兒女、田地， \end{tabularx} \\ \\ \relax
10:30 & \begin{tabularx}{0.7\textwidth}{X} 沒有不在今世得百倍的，就是房屋、兄弟、姊妹、母親、兒女、田地，並且要受迫害，在來世得永生。 \end{tabularx} \\ \\ \relax
10:31 & \begin{tabularx}{0.7\textwidth}{X} 然而，有許多在前的，將要在後；在後的，將要在前。」 \end{tabularx} \\ \\
[1ex]
\hline
\hline
\end{longtable}
講題:錢財是永生的敵人嗎? Is Wealth the Enemy of Eternal Life?

經文:馬可福音10:17-31

講員:林豪恩先生

場合:香港中文大學崇基學院禮拜堂主日崇拜

日期:2024年10月13日

一. 看得起你才邀請你

還記得福音書的那個故事嗎?這個人滿懷期望,興高采烈地跑過來找耶穌,結果卻是垂頭喪氣地離開.望著他無奈地離開的身影,你有甚麼感覺呢?會不會有點同情這個人呢?會不會覺得那位老師要求過分呢?由於那位老師是耶穌,於是信徒都為他辯護.辯護的主要策略是「質疑對手」,指出做成「憂憂愁愁地走了」這結局的責任在於這個人自己,而不在耶穌.你有沒有聽過下列的說法呢?

指責這個人稱呼耶穌為「良善的夫子」(17)是貶低耶穌.在這個人心目中,耶穌只是個很好的老師,卻不是基督.辯護者順便「抽水」說:這等於今天有人認為耶穌是絕頂好人,卻不承認祂是神的兒子.

指責這個人問耶穌「我當做什麼事才可以承受永生?」(17)表示他以為靠行為得救.

指責這個人表面上很不錯,但只做好了十誡後面對人的部分,卻沒有做好前面對神的部分.指責者順便「抽水」說:這也是現代人失誤之處,只著重對人做善事,熱心公益,造福社會,慷慨解囊,一生正直,然而對神方面卻一片空白.

指責這個人心目中以錢財比神更為重要；他以瑪門為神,他沒有尊耶和華為獨一的真神.

如此質疑這個人,是否合符耶穌對他的觀感呢?讓我們看看耶穌如何評價這個人.當耶穌聽到他說:「夫子,這一切(律法)我從小都遵守了.」(10:20)作者記下耶穌的反應:「耶穌看著他,就愛他」(10:21)

為耶穌辯護的人負評這個人,但其實耶穌欣賞這個人；為耶穌辯護的人指責他咎由自取,但其實耶穌喜愛他.耶穌應該遇過很多人,但不是對每個人他都發出跟隨自己的邀請.耶穌向這個人發出邀請,可能是耶穌喜愛他,看得起他,認為他合適.

這個人「憂憂愁愁地走了」這個畫面難免令人感到可惜.不過,也不需要判定這就是這個人一生的結局.我們怎知道下一刻會怎樣呢?事實上,不少人都曾經和耶穌失諸交臂,但在人生路上某一刻又重遇.就算是舉過手決過志的基督徒,也可能在人生某一刻走開,然後又在另一刻重投上主懷抱.正所謂,一剎那的光輝不代表永恆,一剎那的走開也不代表永遠.此刻「憂憂愁愁地走了」,可能創造了空間,慢慢想清楚人生是甚麼,甚麼才重要；想清楚對自己來說跟隨耶穌是甚麼意思.(聽說有人就是喜歡崇基禮拜堂有類似的空間)

二．永生,得救,進上帝的國

這個人向耶穌請教如何承受永生,引發了接下來耶穌和門徒的討論.大家有沒有留意耶穌的用詞改變了?

10:23耶穌周圍一看,對門徒說:「有錢財的人進上帝的國是何等地難哪!」

10:24門徒希奇他的話.耶穌又對他們說:「小子,倚靠錢財的人進上帝的國是何等地難哪!

10:25駱駝穿過針的眼,比財主進上帝的國還容易呢.」

耶穌改用了「進上帝的國」來承接關於永生的話題.這個人所尋求的永生和耶穌所說進上帝的國會不會是相同的事情呢?如果是,會不會啟發我們對「永生」更豐富的想像呢? 另外,有沒有留意門徒用甚麼詞彙來談論這個話題呢?

10:26門徒就分外希奇,對他說:「這樣誰能得救呢?」

門徒名之為「得救」,這個人名之為「永生」,耶穌稱之為「進上帝的國」.如果我們用耶穌的表達來閱讀這個故事,又會有甚麼啟發呢?

三．上帝的國與這世界

「上帝的國」相對於甚麼呢?耶穌被彼拉多審問時說過:「我的國不屬這世界」(約翰福音18:36).我們可以假設「上帝的國」相對於「這世界」.如果我們以「這世界」來理解故事中的這個人,可能對他有不同的觀感.關於「永生」,我們會說:「你有沒有永生?」／「我有永生」,彷彿永生是我們的擁有品之一.

擁有在這世界是一個重要的主題,這世界鼓勵人們擁有,人們也為了擁有而努力奮鬥,例如擁有財富,學歷,權力,名譽,甚至擁有全世界.這世界認為擁有得越多就越成功,擁有累積起來就成為雄厚的資本.資本越來,人就越安全,也帶來越大的收益.有同學說,在名校就讀的中學生,已經懂得累積資本,他們的資本是甚麼呢?就是簡稱為CV的履歷表,他們稱這種操作為「砌靚張CV」.甚麼人懂得「砌靚張CV」呢?越有條件的人,越懂得「砌靚張CV」.如果我們用這個角度來看故事中來訪耶穌的這個人,他的CV相當不錯.

從他和耶穌的對話中,可見他的CV「遵守校規」部分必定A+.另外一個部分如果詳細列出,恐怕要加幾頁紙.那是甚麼?是資產,經文指出他的產業很多.張CV已經咁靚,還未足夠嗎?在資本世界,有足夠這個概念嗎?你做足100分,競爭對手挑戰無限分.永無止境的競爭和比較,不要說前線職員感到壓力,就是高層也如履薄冰.

可能故事中這個人都認同自己的CV已經好靚,但好像還欠一樣,就是永生.你有沒有感到奇怪,那麼有條件的人都需要宗教?擁有越多的人,越需要宗教,因為他們需要安全和穩定.大家知不知道有些富豪擁有自己的私家宗教場所?擁有私家聖職人員?

雖然相隔二千年,我們或多或少都明白故事中這個人的處境.他所承受的價值觀,我們也承受.他其實是這個世界的精英,這個世界教育出來的模範.作為父牧師長的不妨自問,如果我們的兒女學生能夠有這個人的成就,已經心滿意足了,對嗎?做成這個人憂憂愁愁地離開的畫面,不是這個人不夠好,反而顯示出這個世界和上帝的國兩者之間的差異實在太大.喝這個世界的奶長大,以這個世界的價值觀為血肉的人,怎可能不對於耶穌所說的上帝國感到陌生,感到難以接受呢?因此,耶穌慨嘆地對門徒說:

10:24「倚靠錢財的人進上帝的國是何等地難哪!」

上帝的國是怎樣的呢?和這世界有甚麼不同呢?在這個段落中,耶穌是這樣說的:

10:29耶穌說:「我實在告訴你們,人為我和福音撇下房屋,或是弟兄,姊妹,父母,兒女,田地,

10:30沒有不在今世得百倍的,就是房屋,弟兄,姊妹,母親,兒女,田地,並且要受逼迫,在來世必得永生.」

不知道你會否疑問:要一百間房屋尚好,但要一百個弟兄,姊妹,母親,兒女做甚麼呢?對於現代城市人來說,兒女一個都嫌多,一百個簡直是一個惡夢.其實,一百只是表示數量龐大的約數,重要的是耶穌對關係的重新定義:(馬可福音)

3:32有許多人在耶穌周圍坐著,他們就告訴他說:「看哪,你母親和你弟兄在外邊找你.」

3:33耶穌回答說:「誰是我的母親?誰是我的弟兄?」

3:34就四面觀看那周圍坐著的人,說:「看哪,我的母親,我的弟兄.

3:35凡遵行上帝旨意的人就是我的弟兄姊妹和母親了.」

由此可以想像,進入上帝的國,意味著進入導行上帝旨意的大家庭.如此的大家庭會是一種怎樣的狀況呢?使徒行傳記載曾經類似的狀況:

2:44信的人都在一處,凡物公用,

2:45並且賣了田產,家業,照各人所需用的分給各人.

最早期的信徒群體曾經體驗過人間天堂般的狀況,然而,這種狀況一瞬即逝,就正如很多美好的事物,無論是日出的壯麗,或日落的瑰麗,都一瞬即逝,但卻動人心弦,令人響往.人總是渴望留住美好一刻變為擁有品.死纏不放,如果不是徒勞無功,就可能是痛苦的開始.這世界畢竟不是上帝的國,勉強把這世界變成天堂,可能弄巧成拙,為蒼生帶來痛苦.

10:26門徒就分外希奇,對他說:「這樣誰能得救呢?」

10:27耶穌看著他們,說:「在人是不能,在上帝卻不然,因為上帝凡事都能.」

有些事情,如果人類能夠承認自己的限制,尊重和信任上帝,反而能夠減少個人痛苦和人間災難.

四．上帝國的局部體現

耶穌邀請人們進入上帝的國,然而人們仍然生活在這世界之中.如何在這世界體現上帝的國呢?這個故事揭示出這世界和上帝國的不同,在於這世界推崇擁有和累積,上帝國的精神是分享和分擔.福音書中有一個例子可以作為說明,還記得耶穌離世前在十字架上作出的交托嗎?(約翰福音)

19:26耶穌見母親和他所愛的那門徒站在旁邊,就對他母親說:「母親,看,你的兒子!」

19:27又對那門徒說:「看,你的母親!」從此,那門徒就接她到自己家裡去了.

孟子在描述他的理想社會時指出「老吾老以及人之老,幼吾幼以及人之幼」.其實,是不是有一百個母親並不重要,重要的是把眼前有需要人視為家人般照顧.

不同的社會設立不同的制度照顧有需要的人,有些較為全面,有些相對不足.然而,就是在不足的狀況之下,仍可看見人們守望相助的美好畫面.

大約在十六層公共屋村落成不久的那個年代,有一位小孩子入學讀書,對著學生手冊要填寫父母的資料而感到迷茫.第一日送他上學及接他放學的,是對面單位在附近上學的姐姐.放學後,他就去隔鄰的聽障婆婆家吃午餐.後來,這位小孩子懂得煮即食麵,就自己解決午餐了.某年,他的班主任知道他吃即食麵作午餐,於是放學後留下他做一點事情.在做那一點事之前帶同他去學校附近的餐廳吃午餐.老師離開學校時,經過街市,順道買一些肉,吩咐他帶回家做晚餐.時光飛逝,這位孩子要考會考了.那段日子,有些鄰居夜夜笙歌,那孩子擔心睡得不好而沒精神考試.走廊另一端的獨居婆婆把門銷交給他,對他說如果你那邊太吵睡不著,就過來我那邊睡.

雖然這個世界的主旋律仍舊是擁有及累積,但也有不少人樂意以分享及分擔回應別人的需要.這些人所作的或許在這世界來看微不足道,但卻讓有需要的人淺嚐上帝國的美善,也在這世界中為上帝的國畫上美麗的一筆.



\newpage

\section{因真理得自由以服務}
\label{sec:719}
\textbf{伍渭文牧師 - 因真理,得自由,以服務}
\newline
\newline
\hyperref[sec:718]{< < < PREV SERMON < < <}
~
\hyperref[sec:toc]{[返目錄]}
~
\hyperref[sec:720]{> > > NEXT SERMON > > >}


馬可福音 10:35-45
\newline
\begin{longtable}{cl}
\hline
\hline
章節 & 經文 (和合本修訂版)\\
\hline
10:35 & \begin{tabularx}{0.7\textwidth}{X} 西庇太的兒子雅各和約翰進前來，對耶穌說：「老師，我們無論求你甚麼，願你為我們做。」 \end{tabularx} \\ \\ \relax
10:36 & \begin{tabularx}{0.7\textwidth}{X} 耶穌對他們說：「要我為你們做甚麼？」 \end{tabularx} \\ \\ \relax
10:37 & \begin{tabularx}{0.7\textwidth}{X} 他們對他說：「在你的榮耀裡，請賜我們一個坐在你右邊，一個坐在你左邊。」 \end{tabularx} \\ \\ \relax
10:38 & \begin{tabularx}{0.7\textwidth}{X} 耶穌對他們說：「你們不知道所求的是甚麼。我所喝的杯，你們能喝嗎？我所受的洗，你們能受嗎？」 \end{tabularx} \\ \\ \relax
10:39 & \begin{tabularx}{0.7\textwidth}{X} 他們對他說：「我們能。」耶穌對他們說：「我所喝的杯，你們要喝；我所受的洗，你們也要受。 \end{tabularx} \\ \\ \relax
10:40 & \begin{tabularx}{0.7\textwidth}{X} 可是坐在我的左右，不是我可以賜的，而是為誰預備就賜給誰。」 \end{tabularx} \\ \\ \relax
10:41 & \begin{tabularx}{0.7\textwidth}{X} 其餘十個門徒聽見，就對雅各和約翰很生氣。 \end{tabularx} \\ \\ \relax
10:42 & \begin{tabularx}{0.7\textwidth}{X} 耶穌叫了他們來，對他們說：「你們知道，外邦人有君王作主治理他們，有大臣操權管轄他們。 \end{tabularx} \\ \\ \relax
10:43 & \begin{tabularx}{0.7\textwidth}{X} 但是在你們中間，不可這樣。你們中間誰願為大，就要作你們的用人； \end{tabularx} \\ \\ \relax
10:44 & \begin{tabularx}{0.7\textwidth}{X} 在你們中間誰願為首，就要作眾人的僕人。 \end{tabularx} \\ \\ \relax
10:45 & \begin{tabularx}{0.7\textwidth}{X} 因為人子來，並不是要受人的服事，乃是要服事人，並且要捨命作多人的贖價。」 \end{tabularx} \\ \\
[1ex]
\hline
\hline
\end{longtable}
講題:因真理,得自由,以服務 Freedom Through Truth For Service

經文:馬可福音10:35-45

講員:伍渭文牧師

場合:香港中文大學崇基學院禮拜堂主日崇拜

日期:2024年10月20日

今天的題目來自燕京大學的校訓,亦來自今天福音經課:「你們知道,外邦人有尊為君王的,治理他們,有大臣操權管束他們.祇是在你們中間,誰願為大,就必作你們的用人.在你們中間,誰願為首,就必作眾人的僕人.因為人子來,並不是要受人的服事,乃是要服事人,並且要捨命作多人的贖價.」(可十42-45) 這段經文也是貫穿馬可福音的重要主題.

我們崇拜讀的經課又稱三代經題,因三年一循環,讀畢三卷重點不同的符類福音: 馬太強調耶穌是教師,宣講天國度理; 馬可指出耶穌的權能, 驅魔醫病, 服事有需要的人; 路加顯明耶穌的憐憫, 關心貧窮的人和弱勢群體.今年是馬可年.馬可福音篇幅最小 (可23頁, 路40,太37,約29),語句濃縮(compact),鮮活生動(vivid),鋪排有序(orderly).例如作者馬可有意在八,九,十章,三度提及耶穌對門徒預告: 他要上耶路撒冷,被文士,長老和祭司長們棄絕,並且被殺,第三天復活.此後耶穌就騎著驢駒子進入耶路撒冷,十一至十六終章,記錄耶穌的在耶路撒冷所受諸般的苦難.

讓我們看看這三段上耶路撒冷受死的預告和門徒的反應,來了解聖經對服事的教導.我們看看這地圖三處耶穌跟門徒預告他要是耶路撒冷受死的地點.由北到南:該撒利亞腓立比,加利利湖旁迦伯農,耶利哥往耶路撒冷路上.

1)八27-38地點:該撒利亞腓立比.善用人生服事他者.

耶穌問門徒,人說我是誰?彼得回答說:你是基督.「耶穌就禁戒他們不要告訴人.從此他教訓他們說,人子必須受許多的苦,被長老祭司長和文士棄絕,並且被殺,過三天復活.耶穌明明的說這話,彼得就拉著他勸他……耶穌轉過來,看著門徒,就責備彼得說:撒旦,退我後邊去吧,因為你不體貼神的意思,祇體貼人的意思.若有人要跟從我,就當捨己,背起他的十字架來跟從我……凡為我和和福音喪掉生命的必救了生命.」(八30-34)

該撤利亞腓立比.巴勒斯坦最高的山是終年積雪的黑門山,耶穌曾在其中較低的山峰顯容.積雪溶解南流至很多巖洞和泉眼的該撤利亞腓立比.這裡建有眾多希臘,羅馬的神盦,香客絡驛不絕,他們來此尋來人生的意義,求神祇指點迷津.在這裡,舉頭高山白雪鎧鎧無瑕,低首泉湧清澈涼凜,叫人想起自己的罪污,和靈魂的救贖；在這裡,人求神問卜,如何過一個有意義的人生.

為何耶在此問門徒:「人說我是誰?」因為尋道者最焦慮是找到救贖的方法.薛華(Francis Schaeffer)【他在那裡,他並不沉默】(He Is There and He Is Not Silent)這書,來自他的公開講演「Possible Answers to the Basic Philosophical Questions」(基本哲學問題的可能答案).他說哲學有三大基本問題:形而上問題 :受造世界背後有造物主嗎?道德問題:行為可有對錯(善惡)標準嗎?知識論問題:如何得到真確的知識?尋道者不難相信有造物主；亦不難相信行為有對錯；但如何得到真理,就不容易了.  耶穌藉彼得的認信,確認他就是彌賽亞,為人捨命的救世主.

古教父奧古斯丁(A.D. 354-430)悔改後,寫了《懺悔錄》,他說最重要的知識,就是「認識自己是一個罪人,且可藉悔改認罪,得到上主的赦免.」

彼得不想耶穌赴義,因為他以為神的國就是用武力結束羅馬鐵蹄統治,若耶穌死了,這運動就完了.耶穌趁機會解釋:跟隨他的,也要背起十字架捨棄生命.耶穌釘十架是為人的罪而捨命,跟隨耶穌也要背起十架捨命,這捨命是感戴耶穌代贖之恩,而為他者的需要用我們的一生.

毋忘死亡(Mememto Mori).我們每天都喪掉生命,像燃點中的洋燭,愈來愈短.捨去生命,沒有選擇；但我們可以選擇為了甚麼捨去生命,如何捨去生命.為基督捨去生命,為基督用我們在世的時間,才經歷有意義,滿足的人生.「凡為我和福音喪掉生命的,必救了生命.人就是賺得全世界,賠上自己的生命有甚麼益處呢.人還能拿甚麼換生命呢?」(可八35, 36)

倪匡在  一聚會中說,達摩面壁九年,我用十年在美國勤讀聖經,了解一個真理:今生很短很短,來生很長很長,不值得為了短暫的今世勞勞碌碌,忽視來生.

聯合國第二任秘書長韓馬紹(Dag Hammarskjold)說得好:「我們不需要尋求死亡,死亡會找上門的,但我們要找到往死亡的路,使死亡顯得有意義 ……命運本身不能逃避,也不需響往,它不是沒有原委的隱密,因為它說明這世界和人類歷史的軌跡,都有意義.」

德國宣教士史懷哲Albert Schweitzer才華橫溢,他是出色的風琴演奏家,二十五歲成為神學院講師,三十一歲出版【歷史耶穌的追尋】(The Quest of Historical Jesus).為了回應宣教的感動,二十九歲重回學校,用六年時間研習醫學,奉派到菲洲醫療宣教,被稱為菲洲之父.他經常思考耶穌所說:「我們不應該只為自己而活.」認為幸福的人,應該造福他人；擺脫苦惱的人,應該去減輕他人的苦惱；世上一切的不幸,應該由大家共同來承擔.史懷哲說:人生最大的滿足,是服事有需要的人.

如何為上主捨命?在教會歷史中,有信徒以守獨身,到修道院與世隔絕,專心禱告回應耶穌為我們捨命.今天也有信徒以敬拜讚美,來供奉上主.其實上主並不需萬萬油河的,祇想我們行公義,好憐憫,存謙卑的心,與上主同行.捨己是為他者的需要而活,服事他們.

中古修道院的興起,其實是針對社會的需要,服事社會.十三世紀初,歐洲大學興起,帶來護教解惑的需要,道明會(Dominican Order, Order of Preachers) 1216年應運而生,多馬斯.亞奎亞(Thomas Aquinas)是其姣姣代表.1209,城市興起,社會貧富懸殊,托缽(Mendicant)方濟會(Franciscan Order)成立,扶貧濟世.1650年法國耶穌會會士麥道爾(Medaille)組織年輕婦女創立麥道爾聖約瑟姊妹會(Congregation of the Sisters of St. Joseph of Medaille):「去到城市,分散進入鄰舍,找出他們的需要,招募樂於行善的平信徒一起同工.」她們開始照顧孤寡,濟貧治病,推廣教育.1789年法國大革命爆發,反宗教情緒高漲,院舍土地被強行沒收,社群被驅散,有些人下在獄中,甚至上了斷頭臺.[1]

在該撒利亞腓立比耶穌第一次預告他要受死,為我們受死,叫我們不用懼怕死亡,終身患得患失過日子；但毋忘死亡,人生短暫,好好利用日子為上主工作,愛上主所愛的人.

2)九30-32 地點: 加利利,近迦百農(加利利海北端).大的要服事小的, overpower or empower.

耶穌在黑門山高山顯容,下山途中醫好癲癇病的孩子.「經過加利利……教訓門徒說:人子將要被交在人手裡,他們要殺害他,被殺以後,過三天他要復活.門徒卻不明白這話,又不敢問他.」(可九30-32)

來到迦百農一屋裡,耶穌問門徒:「在路上議論的是甚麼?門徒不作聲,因為他們在路上彼此爭論誰為大.」耶穌說為首先的,他必作眾人末後的,作眾人的用人.「於是領過一個小孩子來,叫他站在門徒中間,又抱起他來,對他們說,凡為我名,接待一個像這小孩子的就是接待我,凡接待我的,不是接待我,乃是接待那差我來的.」(可九36, 37) 抱起小孩,可翻作雙手放在孩子背後.有人把孩子帶到耶穌面前,使他親近耶穌,甚至聆聽耶穌講道.今屆金書獎(教會及佈道事工)得主是一位崇基神學院校友,作品【讓孩子爬進耶穌懷中】,提及信仰的培育,始於幼年,就在家中.

孩子也可以是需要照顧的家人,像馬可第九章記述耶穌黑門山(麓)登山顯出榮耀的臉容.下山途中遇啞吧鬼附身小孩,發作時像害癲癇,污鬼把他摔倒,口吐流沫,咬牙切齒,身體枯乾,孩子顯出的是使人懼怕扭曲的臉容.耶穌驅逐污鬼,使孩子痊愈.「夫子,我帶了我的兒子到你這裡來.」(可九17)這「孩子」亦可能是家中年邁需要照顧的長輩；服事從家中開始,不用到橋底,也不用到非洲.

帝國藉武力東征西伐,建立殖民地,宣示其文化,強徵賦稅.耶穌降生在伯利恆,因為約瑟和馬利亞要回鄉為交稅「報名上冊」.羅馬士兵可要求猶太人無償跟他背負行李,走一哩路程.(太五41)這是帝國邏輯,小的要服事大.在現代辦公室,也有下級職員自願為上級主管沖咖啡,甚至買早餐.世俗社會擁有權力,用來壓制(overpower)小的,但耶穌說,大的要為小的充權(empower).服事是為他者的需要,為首的要服事小的.這觀念顛覆了帝國的權力觀.

服事他者,為他者充權,其實為自己的屬靈生命充權.八福中第五福也說:「憐恤人的人有福了,因為他們必蒙憐恤.」前香港特區律政司長黃仁龍有一次在崇基周會講演說:當律政司長前,他每個月盡可能參加有關團體舉辦的活動,探訪天橋底的無家者,免得自己在富裕安舒的生活中,迷失自己.當了律政司長,「我每次看到堆在桌面一大堆文件,都看到背後很多不同的臉相,滿臉愁容和焦慮.」耶穌說:「凡為我和福音捨去生命的(用去光陰),必救了生命.」(可八35)

3)十32-45地點:往耶路撒冷途中,接近耶利哥.

第三段將要被殺的預告前,特別描述耶穌的心情:「他們行路上耶路撒冷去,耶穌在前頭走,門徒就希奇,跟隨的人也害怕.耶穌又叫過十二個門徒來,把自己將要遭遇的事,告訴他們.」

耶穌行在前面很特別,一般有門徒在前面開路,防避潛在兇險,保護老師.耶穌知道此次斯人一去不復還.他在伯大尼使拉撒路復活後,祭司長該亞法已通過猶太人公會(Sanhedrin)要處決耶穌.耶穌與門徒退到約但河東南的曠野退隱,避開風頭.

耶穌對耶路撒冷並不陌生,十二歲那年,他在聖殿三天專心耹聽律法師講解聖經,連父母離開也不知道,因為他以父的事為念.作為成年猶太人,每年幾次都上聖殿敬拜.

但今次不同,他要赴耶路撒冷受死,不是一般的死亡,是被出賣,公開被羞辱,被鞭打,背著沉重的十架赴城外各各他小山丘,釘死在其上.十字架的刑罰是羅馬人非常不人道的刑具,有羅馬公民權的免釘十字架.二零零四年二月二十五日聖灰日米路吉遜(Mil Gelson) 導演的《受難曲》(The Passion of the Christ) 上影,因為太血腥定為三級片,其實導演祇是寫實,十架的刑罰就是那麼血腥殘酷.

《馬利亞的泣訴》.香港電台一分鐘閱讀曾經介紹當代愛爾蘭文學巨人柯姆．托賓(Colm Tóibín)所著《馬利亞的泣訴》(The Testament of Mary).全書以第一身書寫,這位全程目睹兒子慘痛遭遇的母親痛不欲生,她的聲音終於被聽見:「我看見我的兒子被綁起來而且流著血,我聽見他的哭喊……」作者說:當長釘刺穿耶穌手臂時,兩千年來不曾有人聽過她(馬利亞)的一聲尖叫.聖母的沉默,至今才有人書寫她的泣訴.

為了迎接這苦路,耶穌與門徒到約但河東南面的曠野避靜,時候到了,他昂首闊步,直面耶路撒冷,走在前面,門徒感到希奇,隨行的人也害怕.莫非,神的國快將降臨?我們將得著榮耀了!他們真的這樣想,路十九11:「他們以為神的國快要顯出來.」

門徒兩次聽到耶穌宣告他要往耶路撒冷受死,並沒有醒覺,反面談論誰為大,誰可坐在耶穌左邊和右邊.他們是否太過份呢?我們觀察到耶穌並沒有嚴嚴責備他們,原因是他們以為耶穌的上耶路撒冷要建立彌賽亞的國度,推翻羅馬鐵蹄統治,像公元前167至37年,猶太人推翻希臘帝國統治,建立馬迦比王朝.革命原因是希臘強權強行推動其宗教文化,掠奪聖殿聖物,並在祭壇獻上被視為污穢的豬.

門徒爭論為大,是想更多付出,率先領前赴義；坐在耶穌左右,是希望被耶穌信任重用.耶穌並沒有嚴厲責備他們,反而溫柔教導他們,使他們慢慢明白,他們也會步其後塵,為信仰澆奠鮮血.成聖的路和明白真理的路是漫長的,上主對我們忍耐,我們也要對別人和自已忍耐,但明白真理後就要義無反顧.

「耶穌說,你們不知道所求的是甚麼.我所喝的杯你們能喝嗎?我所受的洗你們能嗎?他們說,我們能.耶穌說,我所喝的杯,你們也要喝,我所受的洗你們也要受.祇是坐在我我的左右,不是我可以賜的,乃是為誰豫備的,就賜給誰.那十個門徒聽見,就惱怒雅各,約翰.」(可十38-40)

耶穌十二使徒,除了約翰被放逐拔摩海島終其一生,其他都殉道.雅各是第一位為信仰被殺的.「希律(亞基帕)王下手苦害教會中幾個人.用刀殺了約翰的哥哥雅各.他見猶太人喜歡這事,又去捉拿彼得.」(使十二1-3) 古代歷史記載,他們有的死于刀劍,有的被火焚燒,有的鞭打至死,有的用鐵叉叉死,或釘十字架(彼得倒釘,安得烈橫釘),或沉于大海,被剝去皮肉,用鐵鋸鋸死,用石頭打死,放在嚴寒之處凍死,不給飲食饑渴而死,或砍去手腳,赤身露體,受盡淩辱而死.

形勢是嚴峻的,十架就在前面,但耶穌還要教訓門徒:誰願為首,就必作眾人的僕人.「耶穌叫他們來,對他們說,你們知道,外邦人有尊為君王的,治理他們.有大臣操權管束他們.祇是在你們中間不是這樣.你們中間誰願為大就必作眾人的用人,在你們中間,誰願為首,就必作眾人的僕人.因為人子來,並不是要受人的服事乃是要服事人,並且要捨命,作多人的贖價.」(可十42-45)

再細看今天的題目,【因真理,得自由,以服事】.作為大學校訓,因真理得自由,可以解釋為:因知識的啟蒙,擺脫愚昧的自由,以學識服務社會.但馬可所說的,是因信稱義帶給我們的自由.「因為人子來,並不是要受人的服事,乃是要服事人,並且要捨命作多人的贖價.」(可十42-45) 這自由來自基督代贖,不被定罪的釋放.

路德把這自由比肩君王的自由.路德在1520年【論基督徒的自由】開始的一段說得好:「因著信,基督徒是全然自由的萬人之主,不受任何人管轄；因著愛,基督徒是全然順服的萬人之僕,受一切的管轄.」當人不被罪束縛時,能束縛我們?到上主面前,唯獨信心,唯獨恩典；但心被恩感,回到人面前是愛,回應上主赦罪之恩之愛,服事他者的愛.為弱小充權的愛.

在改教運動前夕,服事並不是完全自由的,服事是一種功德,藉此得到上主接納,罪得赦免.沒有人行經時間,有絕對把握能進天堂,終生都在誠惶誠恐中,面對最後的審判.

君王治理,大臣操權,說的是外在的約束；為首作僕人是內在,自發,甘心.出於愛的行為,超越律法要求,這愛來自對基督捨命的回應.來到上主面前,唯獨信心,罪得赦免.若罪不能轄制我們,我們的自由何其大,路德把這自由比肩君王的自由,沒有人攔阻君王的自由.心被救贖之恩感動,驅使我們去到人群當中,為他者的需要顯出愛.服事是為他者的需要,大的要服事小的,服事是對基督捨命的愛的回應.

1)服事是為他者的需要而活.我們每天都喪掉生命,捨去生命,沒有選擇；但我們可以選擇為了甚麼捨去生命,如何捨去生命.為基督捨去生命,為基督用我們在世的時間,這才經歷有意義,滿足的人生.因此,活著為耶穌才是有意義的生命.

2) 服事是為弱者充權,大的要服事小的.我想起那鬼附的小孩,病發時像發癲癇的小孩,誰人服事他呢?「我帶了我的兒子到你這裡來」(可九17) ,像很多令人敬佩的父母,悉心照顧有成長障礙的孩子；家人中也有年邁的長輩像孩子,需要我們照顧.

大的要服事小的,我也想到耶穌抱起那小孩,教導門徒要接待小子如同接待他.這小孩是蒙福的,因有人帶他親近耶穌,甚至被耶穌使用當實物教材.我們在世行經時間,都擁有諸般的關係和責任,用愛心盡上本份,就是服事,不一定要到菲洲和橋底.服事他者,從自己的家開始.

3) 服事是對上主救贖之恩的回應.基督信仰一路來都有服務的傳統.香港基督教服務處1952年成立,官網介紹:「一直為有需要人士提供適切,專業,真誠的優質服務,對弱勢及被忽略的社群尤為關注,並以社會仁愛公義,人人全面發展為願景,為市民及社會締造希望,倡導公義,牽引共融.」

耶穌在世有三重職事,除了宣講(福音的宣講),還有對人的服事:驅魔,醫病.宣講是聽得見的福音,驅魔,醫病是看得見的福音.驅魔是靈體的撥亂反正,醫病是肉體的撥亂反正,都是天國降臨的記號,兩者不可分割.

服事他者也是對上主真信心的表示.馬太二十五章最後審判時,不是口說主阿!主阿!就可以進天堂,乃我餓了你給我吃,我渴了你給我喝,我下在獄中你來探我,那些以行動來服事他者顯出對上主信心的人,沒有行為的信心是死的.

在結束時問問:那個門徒為大?那個門徒在耶穌得榮耀時,列在耶穌的左右?耶穌被高舉在各各他的十架時,是他得榮耀時刻,在那一刻,憐憫向審判跨勝了.在他兩旁的強盜,都清楚聽到十架七言中的召喚:「父阿,赦免他們!因為他們所做的,他們不曉得.」(路二十三34)

其中一個強盜悔改:「耶穌阿,你得國降臨的時侯,求你.記念我.」耶穌說:「我實在告訴你:今日你要同我在樂園裡了.」(路二十三43)每一位悔改信徒,悔改的一刻憑信與基督進入樂園了,現在亦憑信與基督坐在父上主的右邊,將來與基督一同審判世界.(弗二4-7；林前六2, 3)

信耶穌那一刻,每個信徒都列在耶穌左右.「因著信,基督徒是全然自由的萬人之主,不受任何人管轄；因著愛,基督徒是全然順服的萬人之僕,受一切的管轄.」

「因著信,基督徒是全然自由的萬人之主,不受任何人管轄；因著愛,基督徒是全然順服的萬人之僕,受一切的管轄.」

[1] 現全球五十國共有14,000姊妹會會員,過半在美國,2,000在法國.1937在紐約洲西創立麥道爾大學.2017年其屬下一間醫療機構「聖猶大醫院」(St. Jude Hospital)獲選全美婦女海選最佳醫院 (2017 Women’s Choice Award, American Best Hospitals).服事是為他者的需要而活.



\newpage

\section{聖靈引導教會進入真理}
\label{sec:720}
\textbf{吳梓明教授 - 聖靈引導教會進入真理}
\newline
\newline
\hyperref[sec:719]{< < < PREV SERMON < < <}
~
\hyperref[sec:toc]{[返目錄]}
~
\hyperref[sec:721]{> > > NEXT SERMON > > >}


約伯記 42:1-610-17
\newline
\begin{longtable}{cl}
\hline
\hline
章節 & 經文 (和合本修訂版)\\
\hline
42:1 & \begin{tabularx}{0.7\textwidth}{X} 約伯回答耶和華說： \end{tabularx} \\ \\ \relax
42:2 & \begin{tabularx}{0.7\textwidth}{X} 「我知道，你萬事都能做；你的計劃不能攔阻。 \end{tabularx} \\ \\ \relax
42:3 & \begin{tabularx}{0.7\textwidth}{X} 誰無知使你的旨意隱藏呢？因此我說的，我不明白；這些事太奇妙，是我不知道的。 \end{tabularx} \\ \\ \relax
42:4 & \begin{tabularx}{0.7\textwidth}{X} 求你聽我，我要說話；我問你，求你讓我知道。 \end{tabularx} \\ \\ \relax
42:5 & \begin{tabularx}{0.7\textwidth}{X} 我從前風聞有你，現在親眼看見你。 \end{tabularx} \\ \\ \relax
42:6 & \begin{tabularx}{0.7\textwidth}{X} 因此我撤回，在塵土和爐灰中懊悔。」 \end{tabularx} \\ \\ \relax
42:7 & \begin{tabularx}{0.7\textwidth}{X} 耶和華對約伯說話以後，耶和華就對提幔人以利法說：「我的怒氣向你和你兩個朋友發作，因為你們議論我，不如我的僕人約伯說的正確。 \end{tabularx} \\ \\ \relax
42:8 & \begin{tabularx}{0.7\textwidth}{X} 現在你們要為自己取七頭公牛、七隻公羊，到我的僕人約伯那裡去，為自己獻上燔祭，我的僕人約伯就為你們祈禱。我必悅納他，不按你們的愚妄處置你們。你們議論我，不如我的僕人約伯說的正確。」 \end{tabularx} \\ \\ \relax
42:9 & \begin{tabularx}{0.7\textwidth}{X} 於是提幔人以利法、書亞人比勒達、拿瑪人瑣法遵照耶和華所吩咐的去做，耶和華就悅納約伯。 \end{tabularx} \\ \\ \relax
42:10 & \begin{tabularx}{0.7\textwidth}{X} 約伯為他的朋友祈禱。耶和華就使約伯從苦境中轉回，並且耶和華賜給他的比他從前所有的加倍。 \end{tabularx} \\ \\ \relax
42:11 & \begin{tabularx}{0.7\textwidth}{X} 約伯的兄弟、姊妹，和以前所認識的人都來到他那裡，在他家裡跟他一同吃飯。他們因耶和華所降於他的一切災禍，都為他悲傷，安慰他。每人送他一塊可錫塔和一個金環。 \end{tabularx} \\ \\ \relax
42:12 & \begin{tabularx}{0.7\textwidth}{X} 這樣，耶和華後來賜福給約伯比先前更多。他有一萬四千隻羊，六千匹駱駝，一千對牛，一千匹母驢。 \end{tabularx} \\ \\ \relax
42:13 & \begin{tabularx}{0.7\textwidth}{X} 他也有七個兒子，三個女兒。 \end{tabularx} \\ \\ \relax
42:14 & \begin{tabularx}{0.7\textwidth}{X} 他給長女起名叫耶米瑪，次女叫基洗亞，三女叫基連哈樸。 \end{tabularx} \\ \\ \relax
42:15 & \begin{tabularx}{0.7\textwidth}{X} 在全地的婦女中找不著像約伯的女兒那樣美貌的。她們的父親使她們在兄弟中得產業。 \end{tabularx} \\ \\ \relax
42:16 & \begin{tabularx}{0.7\textwidth}{X} 此後，約伯又活了一百四十年，得見他的四代兒孫。 \end{tabularx} \\ \\ \relax
42:17 & \begin{tabularx}{0.7\textwidth}{X} 這樣，約伯年紀老邁，日子滿足而死。 \end{tabularx} \\ \\
[1ex]
\hline
\hline
\end{longtable}
使徒行傳 2:1-21
\newline
\begin{longtable}{cl}
\hline
\hline
章節 & 經文 (和合本修訂版)\\
\hline
2:1 & \begin{tabularx}{0.7\textwidth}{X} 五旬節那日到了，他們全都聚集在一起。 \end{tabularx} \\ \\ \relax
2:2 & \begin{tabularx}{0.7\textwidth}{X} 忽然，有響聲從天上下來，好像一陣大風吹過，充滿了他們所坐的整座屋子； \end{tabularx} \\ \\ \relax
2:3 & \begin{tabularx}{0.7\textwidth}{X} 又有舌頭如火焰向他們顯現，分開落在他們每個人身上。 \end{tabularx} \\ \\ \relax
2:4 & \begin{tabularx}{0.7\textwidth}{X} 他們都被聖靈充滿，就按著聖靈所賜的口才說起別國的話來。 \end{tabularx} \\ \\ \relax
2:5 & \begin{tabularx}{0.7\textwidth}{X} 那時，有從天下各國來的虔誠的猶太人，住在耶路撒冷。 \end{tabularx} \\ \\ \relax
2:6 & \begin{tabularx}{0.7\textwidth}{X} 這聲音一響，許多人都來聚集，各人因為聽見門徒用他們各自的鄉談說話，就甚納悶， \end{tabularx} \\ \\ \relax
2:7 & \begin{tabularx}{0.7\textwidth}{X} 都詫異驚奇說：「看哪，這些說話的不都是加利利人嗎？ \end{tabularx} \\ \\ \relax
2:8 & \begin{tabularx}{0.7\textwidth}{X} 我們每個人怎麼聽見他們說我們生來所用的鄉談呢？ \end{tabularx} \\ \\ \relax
2:9 & \begin{tabularx}{0.7\textwidth}{X} 我們帕提亞人、瑪代人、以攔人，和住在美索不達米亞、猶太、加帕多家、本都、亞細亞、 \end{tabularx} \\ \\ \relax
2:10 & \begin{tabularx}{0.7\textwidth}{X} 弗呂家、旁非利亞、埃及的人，並靠近古利奈的利比亞一帶地方的人，僑居的羅馬人， \end{tabularx} \\ \\ \relax
2:11 & \begin{tabularx}{0.7\textwidth}{X} 包括猶太人和皈依猶太教的人， 克里特人和阿拉伯人，都聽見他們用我們的鄉談講論神的大作為。」 \end{tabularx} \\ \\ \relax
2:12 & \begin{tabularx}{0.7\textwidth}{X} 眾人就都驚奇困惑，彼此說：「這是甚麼意思呢？」 \end{tabularx} \\ \\ \relax
2:13 & \begin{tabularx}{0.7\textwidth}{X} 還有人譏誚，說：「他們是灌滿了新酒吧！」 \end{tabularx} \\ \\ \relax
2:14 & \begin{tabularx}{0.7\textwidth}{X} 彼得和十一個使徒站起來，他就高聲向眾人說：「猶太人和所有住在耶路撒冷的人哪，這件事你們要知道，要側耳聽我的話。 \end{tabularx} \\ \\ \relax
2:15 & \begin{tabularx}{0.7\textwidth}{X} 這些人並不像你們所想的喝醉了，因為現在才早晨九點鐘。 \end{tabularx} \\ \\ \relax
2:16 & \begin{tabularx}{0.7\textwidth}{X} 這正是藉著先知約珥所說的： \end{tabularx} \\ \\ \relax
2:17 & \begin{tabularx}{0.7\textwidth}{X} 『神說：在末後的日子，我要將我的靈澆灌凡血肉之軀的。你們的兒女要說預言；你們的少年要見異象；你們的老人要做異夢。 \end{tabularx} \\ \\ \relax
2:18 & \begin{tabularx}{0.7\textwidth}{X} 在那些日子，我要把我的靈澆灌，甚至給我的僕人和婢女，他們要說預言。 \end{tabularx} \\ \\ \relax
2:19 & \begin{tabularx}{0.7\textwidth}{X} 在天上，我要顯出奇事，在地下，我要顯出神蹟，有血，有火，有煙霧。 \end{tabularx} \\ \\ \relax
2:20 & \begin{tabularx}{0.7\textwidth}{X} 太陽要變為黑暗，月亮要變為血，這都在主大而光榮的日子未到以前。 \end{tabularx} \\ \\ \relax
2:21 & \begin{tabularx}{0.7\textwidth}{X} 那時，凡求告主名的都必得救。』 \end{tabularx} \\ \\
[1ex]
\hline
\hline
\end{longtable}
馬可福音 10:46-52
\newline
\begin{longtable}{cl}
\hline
\hline
章節 & 經文 (和合本修訂版)\\
\hline
10:46 & \begin{tabularx}{0.7\textwidth}{X} 他們到了耶利哥。耶穌同門徒並許多人離開耶利哥的時候，有一個討飯的盲人，是底買的兒子巴底買，坐在路旁。 \end{tabularx} \\ \\ \relax
10:47 & \begin{tabularx}{0.7\textwidth}{X} 他聽見是拿撒勒的耶穌，就喊了起來，說：「大衛之子耶穌啊，可憐我吧！」 \end{tabularx} \\ \\ \relax
10:48 & \begin{tabularx}{0.7\textwidth}{X} 有許多人責備他，不許他作聲，他卻越發喊著：「大衛之子啊，可憐我吧！」 \end{tabularx} \\ \\ \relax
10:49 & \begin{tabularx}{0.7\textwidth}{X} 耶穌就站住，說：「叫他過來。」他們就叫那盲人，對他說：「放心，起來！他在叫你啦。」 \end{tabularx} \\ \\ \relax
10:50 & \begin{tabularx}{0.7\textwidth}{X} 盲人就丟下衣服，跳起來，走到耶穌那裡。 \end{tabularx} \\ \\ \relax
10:51 & \begin{tabularx}{0.7\textwidth}{X} 耶穌回答他說：「你要我為你做甚麼？」盲人對他說：「拉波尼，我要能看見。」 \end{tabularx} \\ \\ \relax
10:52 & \begin{tabularx}{0.7\textwidth}{X} 耶穌對他說：「你去吧！你的信救了你。」盲人立刻看得見，就在路上跟隨耶穌。 \end{tabularx} \\ \\
[1ex]
\hline
\hline
\end{longtable}
講題:聖靈引導教會進入真理 Holy Spirit Will Lead Us to All Truths

經文:詩篇126篇1至6節；約伯記42章1至6及10至17節；使徒行傳2章1至21 節；馬可福音10章46至52 節

講員:吳梓明教授

場合:香港中文大學崇基學院禮拜堂主日崇拜

日期:2024年10月27日

我們低頭祈禱: 願我口中的言語,眾弟兄姊妹的心懷意念,都在上主面前得蒙悅納.(詩19: 14) 啊們!

各位親愛的弟兄姊妹主內平安.多謝高牧師的邀請,讓我再有機會與崇基禮拜堂眾兄弟姊妹一同聚集,一同敬拜上帝. 祝願 上主的恩惠及慈愛常與我們同在!

今早福音經課記載一個瞎子得醫治的故事.當耶穌在耶利哥城的時候,有一個瞎子知道他經過,就喊著說:『大衛的子孫耶穌啊,可憐我吧!』.有些人責備他,不准他開聲,耶穌卻停下來,問他說:『你要我為你作甚麼呢?』.瞎子說:『夫子!我要能看見!』.耶穌於是讓他能夠看見! 瞎子立刻跟隨了耶穌.這個故事十分簡單,瞎子知道自己需要甚麼,他大聲呼喊,求耶穌幫助,耶穌於是幫助他,讓他能夠看見.今日我們來到聖堂,我們需要甚麼呢? 我們會否放膽開聲祈求上帝的幫助,說:『我要能看見!』.今日講道的題目是:『聖靈引導教會進入真理!』.耶穌曾應許說:『只等真理的聖靈來了,他要引導你們明白一切的真理』(約16:13).我們能否放膽,祈求上主施恩,開啟我們的心眼,讓我們也可以看見更多,更大的真理嗎?

我最近思考一個問題:到底我們需要明白甚麼真理呢?《使徒行傳》記載,五旬節聖靈降臨帶來教會很大的改變,引導教會明白許多真理! 初期教會保留著猶太教的傳統–嚴守割禮和摩西律法,歧視外邦人,認為外邦人是不信上帝的人,是無得救的.聖靈卻是引導門徒明白『上帝愛世人(是愛世上所有的人)』.耶穌曾吩咐門徒要往普天下傳福音給「萬民」(包括外邦人) 聽,可惜門徒只聚集在耶路撒冷,並未有向外邦人傳福音.終於耶穌親自出馬找保羅,說:『掃羅! 掃羅! 你為什麼逼迫我?』.保羅說:『主啊! 你是誰?』.耶穌說:『我就是你所逼迫的耶穌!』.耶穌隨即呼召保羅,要他作「外邦人的使徒」(徒9:1-22; 羅11:13),將福音傳給外邦人; 上帝亦讓彼得見到異象,引導彼得將福音傳給哥尼流,帶領哥尼流全家得救.連彼得也開口說:『我真的看出上帝是不偏待人的』(徒10:34).結果,初期教會在第一次耶路撒冷大會上,正式准許外邦人入教.(徒15:1-21)這就是教會需要明白的一個重要的真理,今天或許我們心目中也有一些「外邦人」,我們會認為他們是教會之外的人,我們不需要關心,也不會主動向他們傳福音(因為他們是「教外人」).但原來上帝是不偏待人,祂也是愛那些「我們不喜歡的外邦人」.我們需要祈求耶穌『開啟我們的心眼,讓我們能看見!』,也祈求聖靈引導我們,明白這個重要的真理!

我有一個親身的經歷,當我祈求上主幫助的時候,聖靈就像耶穌一樣,打開我的心眼,讓我看見更多,更大的真理.香港回歸中國已經27年了! 我發覺很多弟兄姊妹仍然在掙扎,到底怎樣才能夠「又愛國又愛教」呢? 我感謝上帝,過去30多年我有很多機會往返國內訪問,旅遊,遇到很多愛主的弟兄姊妹.我也有機會在英國,美國等地方生活,認識不少英國,美國的基督徒.我發現,絕大多數的英國或美國的基督徒都是愛自己的國家的,不管當權政府(或許是英國保守黨,工黨,或是美國民主黨,共和黨等),也不管當權政府是否真心信靠上帝,有否按聖經的教導去執政,他們也很少因為信仰不同而反對他們的國家政府的.我因此發現,原來上帝愛世人,也是包括所有的中國人.我也學會了,原來「愛國」和「愛教」兩者並不是對立的,「愛國」和「愛教」是可以共融共存的.我開始與外地的弟兄姊妹分享我在香港一些獨特的經驗:我兩個月大已經被父母帶往聖公會諸聖堂接受嬰兒洗禮.所以我自幼便接受了聖公宗的宗教教育,我的「愛國」的觀念是從宗教教育(「愛教」)培育出來的.聖公宗早禱崇拜是採用《公禱書》,我們崇拜中誦念〈為國家禱文〉(所以我從少便懂得愛國,並且要為國家祈禱).1997年前,我們會為英國祈禱(“God Bless the Queen”),因為香港是英國殖民地; 但1997年後,香港已經回歸中國,成為中華人民共和國香港特別行政區,我們在崇拜中,也自然地轉向為中國祈禱.奇怪嗎? 我們誦念的是同一篇〈為國家禱文〉,但禱告的對象已經改變了的.原來基督教教育是可以中國化的,〈為國家禱文〉是一個很好的例子,由「為英國祈禱」轉變為「為中國祈禱」,這就是「基督教教育的中國化」,是從「愛教」伸展至「愛中國」.我也發覺,原來「愛國」不僅是一種身份的改變,更也是心態上的改變,表達我們對國家的歸屬,是一種委身的表達,像我們愛父母,愛家庭一樣.感謝主! 靠賴上帝聖靈的幫助,打開了我的心眼,讓我看見更多,更大的真理.

我們剛剛讀到舊約經課約伯記42章,說: 『約伯回答耶和華說:我知道你萬事都能作,你的旨意不能攔阻…這些事太奇妙,是我不知道的.』(約伯記 42:1-3) 約伯經歷到從苦難帶來的祝福,他知道上帝萬事都能作,實在太奇妙了,連約伯也不曉得為什麼會是這樣的.讓我想起南京金陵協和神學院一位老前輩汪維藩教授,他在1988年曾應邀來香港參加中文大學的學術研討會.汪教授以〈毀滅,反思,再生– 文革後基督教與中國文化之關係〉為題在大會上發言.[1] 文革時期國家和教會均一起經歷大災難,但汪教授卻有一個新思維,建議我們學習從苦難中找到「再生」的希望.他指出文革十年,中國教會雖未能夠煉成精金,卻也學會了與中國人民一同經歷苦難,『凡事與弟兄相同』[2],在苦難中更可將信仰提升,尋覓一套具中國特色的基督教神學來.他還見證說:文革時期教會所獲得的(祝福)或許會比失去的更多.[3] 奇怪嗎?「苦難可以變為祝福!」; 汪教授引用聖經一句話說:『道成了肉身,住在我們中間』[4],原來今日教會 (所見證的「道」) 必須是「在這裡」- 即是在中國「(道)成肉身」–換句話說,教會必須認同中華民族的文化傳統,與中國人民共同經歷苦難,歷盡滄桑,才能在苦難中衍生出一個具有中國特色的教會與中國基督教神學來.[5] 「(道)如何能在中國成為肉身」,能夠從苦難中看到上帝的祝福呢? 上帝啊! 求你賜下聖靈,開啟我們的心眼,讓我們看到更多,更大的真理.

後來,我又發現,《普天頌讚》原來也有三首聖詩是與愛國有關的,包括290首:「為國祈禱歌」,290[1]首:「主佑中華歌」和291首「主佑我國歌」.這三首歌本是屬於18-19世紀的聖詩,後來被翻譯成為中文,分別由燕京大學宗教學院的劉廷芳教授和趙紫宸教授,及香港文藝出版社的社長黃永熙博士翻譯的.三位均是我十分敬重的老前輩,他們所翻譯的聖詩不僅是幫助我們學習怎樣為國家祈禱,更是見證了「基督教聖詩的中國化」.譬如: 劉廷芳提醒我們:『先賢之神(原來上主也是我們國家的先賢之神),(祂的)聖手無限權威… (在為國家祈禱時,我們也要)求主掌管,監護,支持和引導』; 趙紫宸也確認說:『中華立國,至今四千餘年(我們相信國家的歷史也在上帝的掌管中)…前代歷有聖賢,(我們也須)懇求(上主)引導,懇求(上主)振發中華,常使中華(我們的國家),得蒙(上主)恩澤無涯』.黃永熙博士曾是我們崇基學院音樂系系主任,也曾擔任崇基合唱團指揮,在1974年(文革尚未結束之時),黃博士翻譯「主佑我國歌」,他這樣祈禱說:『求主(祝)佑我(的)國家,黑夜風雨交加,國(家根)基(仍是)鞏固; 每當風雨陰霾,(上帝仍是)掌管風浪(的)主宰,求(主)施(行)大能慈愛,保我國土… (他還說)期待上主宏恩…不單(保)佑我(國家)土地,更求主賜恩慈,遍及(惠及)其他(國家),遠達萬邦,(促進)人類相助相望…天下一家』.這是黃永熙博士為國家的祈禱,我們能否回應說:「誠心所願」嗎?

今年暑假期間,我有幸獲邀往江西省基督教兩會主辦的一個中青年教牧同工的培訓班授課,題目是:『講好中國基督教的故事』.我亦趁這機會學習認識基督教在南昌的故事,我發現,原來天主教傳教士利瑪竇曾經在南昌停留一段時間,他尊重中國文化,願意與士大夫,和佛教僧侶交往,頗受中國官紳的重視;利瑪竇曾寫過一本書:《交友論》(1595年寫成).利瑪竇在中國生活了13年,終於發現「友誼」在中國文化中佔有一個獨特的地位.他明白,若要基督教被中國人民接納的話,傳教士必須要表明西方文化(包括基督教)是不會與中國的社會文化秩序對立的.利瑪竇的發現是十分重要的,傳教士若是想基督教能夠被中國人接受的話,他們必須尊重中國文化,與中國人建立真正的友誼,這就是利瑪竇撰寫《交友論》一書的基本立論.我有機會參觀南昌市一間天主教堂,發現天主教堂內設有一個展覽館,記述利瑪實在華的事蹟,他們正在努力嘗試重新認識利瑪竇「交友之道」,實踐利瑪竇的「文化適應策略」,有學者更稱之為「利瑪竇規矩」,是很有意思的.

南昌市原來也是共產黨起義的一個重要市鎮,市內有一間博物館,是特別記念1927年8月1日「八一南昌起義」的第一槍(砲); 當時,共產黨借用了一所教會學校的校舍 (地址是:子固路83號),原來就是當日南昌市中華聖公會的「宏道堂」及「宏道中學」[6],成為了革命黨起義時的臨時指揮部,由賀龍將軍革命軍駐zhak 的地方.當年宏道堂主任牧師兼宏道中學校長是劉屏庚牧師,當賀龍將軍到達南昌時,他騰空自己的牧師樓給賀龍將軍居住了10天(1927年7月27日到8月5日),策劃「八一起義」. 8月1 日凌晨,賀龍將軍與劉伯承和周逸群三人一起站在牧師樓小樓前台階上指揮武裝起義.臨別時,劉牧師更特別祝福賀龍將軍,祝賀他早日凱旋歸來.賀龍將軍後來成為新中國十大開國功臣之一,中央政府更將宏道堂和宏道中學保留成為「賀龍指揮部舊址(博物館)」.博物館籌建之時,劉牧師亦將賀龍將軍贈送給他的物品,連同起義時國民革命軍第20軍會議室的傢俱(本屬牧師的私人物品)共31套92件一同捐獻給紀念館.[7]

我也有機會參觀「八一南昌起義博物館」,「賀龍總指揮部舊址(博物館)」,在賀龍指揮部博物館外牆大門的門楣上仍然保留了一塊四個字的石刻:「宏毅道成」,奇怪嗎? 這石刻是宏道堂和宏道中學建立時表達出『堅持不懈,宏揚基督福音的信心(和毅力),終必能(道成肉身)結成仁義果實的意思』,這也剛巧是突顯出共產黨起義時的革命精神與基督教的宏道精神相配合,所以能夠被保留下來;博物館的二樓亦「還原貌恢復地」保留了學校的禮堂,見證著當年的基督教團體,青年會,福音醫院,學校,教會牧者,信徒等也有參與和協助共產黨在南昌起義的活動;禮堂內還仍掛著學校當年「宏道小學」的校歌和校訓,校歌寫著:『人能宏道是良圖,具此好基礎; 品格端須有,技術不可無,願為青年之楷模,入世皆坦途』.禮堂兩旁還掛有幾張1935年時祝賀宏道堂和建校10周年紀念的題詞,包括黃光斗的題詞:「聖道宏敷」,和韓仁敦的題詞:「宏開道域」[8] ;還有創立中華基督教會的誠靜怡牧師在宏道小學創校10周年(1935年)紀念時贈送給學校的題詞:「啟發天真」,亦是很有意思的.[9]

我獲邀在南昌市一所聖經學校內負責暑期研習班的講課,從全球地域化的角度重新檢視基督教的歷史,講好「中國基督教會歷史的故事」,我也深刻地被年青教牧同工們學道的熱誠所感動.很多人以為當傳教士離開中國的時候,上帝也離開了中國.不過,當我們重新審視教會歷史的時候,我們卻發現更多,更大的真理.事實上,不管過去的75年來中國發生了甚麼事情,我們看到上帝並沒有離開中國.從一些簡單的統計數字便可得知,1949年中國基督新教徒的人數只是70萬,但2018年國家的統計數字是3,800萬,當然我們也知道國家的統計數字是很保守的,這個數字並沒有將未登記的地下教會和家庭教會計算在內,有些人會估計,中國境內的基督徒人數應該是接近 7,000-8,000 萬.換句話說,70多年間有超過 50-100倍的收成,這豈不是一個神蹟嗎? 證明上帝真的沒有離開中國,上帝的聖靈仍在中國,奮興中國教會,更創造了「教會增長」的奇蹟.

現在回頭看,新中國成立75年來的改變,其實也可以說是另一個神蹟.中國有14億的人口,要脫貧真的不容易.縱然中國在經濟發展上仍是遇到不少的問題,但按照世界銀行的說法,今天中國已經是一個上中等收入的國家和世界第二大經濟體;目前中國是全球貿易第一大國,成為世界第一大出口國和第二大進口國.當今世界仍是在大動亂之中,沒有一個國家可以避免世界動亂帶來的影響,但我們仍然能夠見到國家的進步,它不單是可以獨立自主,還能穩步地發展,更能負起大國的責任,支持公義,敢於為正義發聲,亦致力尋求國與國之間和平共處的方案,帶來人類未來世界和平的希望.這也是值得我們為國家而感到自豪的.剛才提到趙紫宸的祈禱是:『懇求(上主)引導,懇求振發中華,常使中華(我們的國家),得蒙(上主)恩澤無涯』;黃永熙博士亦教導我們祈禱說:『求主(祝)佑我(們的)國家… 不單(保)佑我(國家)土地,更賜下恩慈,惠及其他(國家),遠達萬邦,(促進)人類相助相望…天下一家』.

感謝主! 今日我們來到崇基的禮拜堂,求上帝差派真理的聖靈,進入我們的內心,打開我們的心眼,讓我們見到更多,更大的真理.剛才我們已經提出有三個我們需要明白的真理,包括: (一) 上帝是愛世上所有的人,包括國內的同胞,上帝是真的愛中國和大陸的同胞的.(二) 「(道)在中國成為肉身」,原來我們也要與國內的同胞一同經歷苦難,愛我們的國家,像愛我們的父母一樣; (三) 我們也要學習從苦難中看到上帝的祝福,在磨鍊中發展出一套具有中國特色的基督教神學來.上帝的聖靈能夠開啟我們的心眼,讓我們見到更多,更大的真理.我發現還有第四個真理,就是「基督教教育也可以中國化的」,總結有以下四點:(一)按我個人的經驗,我的「愛國教育」是從「基督教教育」而來,基督徒不單需要愛教,也必須要愛國.我發覺我愈是愛教,也會愈更愛國的,因為我每天都會誦念〈為國家禱文〉,誦唱〈為國祈禱歌〉; (二)「九七」之後,我們是繼續誦念《公禱書》〈為國家禱文〉,但是我們的禱告對象改變了,是由「為英國祈禱」轉變為是為「為中國祈禱」,這不僅是從「基督教教育」發展出來的「愛國教育」,更也是見證了「基督教教育的中國化」! (三) 不論是在國內或是在香港,基督教教育必須發展新的課程,學習「講好中國基督教會的故事」,剛才提到南昌的故事便是一個例子,講好基督教的故事就是要將基督教的歷史放回中國歷史中;另一方面,當談及中國基督教的故事時,我們不應太偏重傳教士的故事,潛意識地將基督教變成「洋教」,反之,我們應該更多地發掘華人信徒在中國歷史中的故事.剛才提到南昌宏道堂的劉屏庚牧師便是一個例子.我們經常聽到很多關於八一南昌起義的故事,其實在教會內我們也應讓信徒更多認識華人信徒也有為國家作了不少有貢獻的事蹟.這些故事就是今天教會最需要的「中國化基督教教育」的課程. (四) 「基督教教育」就是關乎上帝作為和大愛的故事.基督教教育不僅是教聖經,神學,它更需要是見證和宣揚上帝的作為和大愛.今日中國需要的基督教教育,就是能夠見證和宣揚上帝作為的基督教教育.我們要見證的其實就是「基督教教育的中國化」,亦是一種「中國化的基督教教育」.

原來「基督教中國化」必須是從「教育」開始,我們可以教育年青的一代,讓他們能從教會歷史中看到上帝的作為和大愛,看到上帝能夠將苦難變成祝福.今年的巴黎奧運,喚起了我們香港人有一份「愛香港」的心,也讓我想到:其實我們也可以透過「基督教教育的中國化」,不單只是喚醒信徒一份「愛香港」的心,更也是有一份「愛中國」的心.

今日我們來到崇基的禮拜堂,盼望我們也能一起禱告,祈求上帝真理的聖靈,進入我們的內心,打開我們的心眼,好讓我們可以見到更多,更大的真理.願上主親自祝福大家! 啊們! (誠心所願!)

[1]     參 汪維藩:〈毀滅,反思,再生– 文革後基督教與中國文化之關係〉,見刊於《金陵神學誌》,復刊後第9期,1988年11月,頁80-81.

[2]     參《新約聖經》《希伯來書》2 :17.

[3]     他這樣說: 『歷史永遠是上帝手中的一把兩刃利劍.祂藉長達十年的痛苦,使我們這些早已接受基督為救主的罪人,進入那種為自己的民族得拯救甘作代罪羔羊的靈性境界.上帝又藉著這長達十年的苦難,使我們的民族和人民,逐步理解並接近那位背負人類罪孽的羔羊基督』. 參汪維藩,見前註2引文,頁81.

[4]     參《新約聖經》《約翰福音》1:14. 汪先生亦是引用拿撒勒一座教堂地下室祭壇大理石上刻著的一句話: 「道在這裡成為肉身」. 可參見汪維藩:〈道在這裡成了肉身〉,載於《金陵神學誌》,總第17期,1992年12月,頁38.

[5]     參同上註引文,頁39.

[6]     1922年初建時稱為「宏道中學」,1925年後改稱為「宏道小學」.

[7]     (內容參見《黨史學習教育官網》,\url{http://dangshi.people.com.cn/n1/2022/0208/c436975-32347514.html}).

[8]     黃,韓二是當代有名的鄉紳,題詞是為表達他們對宏道堂與學校的美好祝願.「聖道宏敷」,意思就是: 聖道廣傳至高峰,振興神威達顯赫.讀者亦可參閱南昌中華聖公會宏道堂1935年出版之《南昌宏道堂落成十周年紀念刊》,據聞已屬古文物收藏於露天市集古書善本內(244988),亦深具研究價值云.

[9]     (內容參見《黨史學習教育官網》,\url{http://dangshi.people.com.cn/n1/2022/0208/c436975-32347514.html}).



\newpage

\section{上帝的國離我有多遠}
\label{sec:721}
\textbf{高國雄牧師 - 上帝的國離我有多遠}
\newline
\newline
\hyperref[sec:720]{< < < PREV SERMON < < <}
~
\hyperref[sec:toc]{[返目錄]}
~
\hyperref[sec:722]{> > > NEXT SERMON > > >}


講題:「上帝的國離我有多遠?」 “How far is the Kingdom of God from me?”

講員:高國雄牧師

場合:香港中文大學崇基學院禮拜堂主日崇拜

日期:2024年11月03日

各位在信仰的路上尋求的朋友,弟兄姊妹,你覺得自己離上帝的國有多遠呢?

馬可福音12:28-34

28有一個文士來,聽見他們辯論,曉得耶穌回答的好,就問他說:「誡命中哪是第一要緊的呢?」 29耶穌回答說:「第一要緊的就是說:『以色列啊,你要聽,主－我們上帝是獨一的主. 30你要盡心,盡性,盡意,盡力愛主－你的上帝.』 31其次就是說:『要愛人如己. 』再沒有比這兩條誡命更大的了.」 32那文士對耶穌說:「夫子說,上帝是一位,實在不錯；除了他以外,再沒有別的上帝； 33並且盡心,盡智,盡力愛他,又愛人如己,就比一切燔祭和各樣祭祀好的多.」 34耶穌見他回答的有智慧,就對他說:「你離上帝的國不遠了.」從此以後,沒有人敢再問他甚麼.

文士=知識份子,抄寫人員,解經家,律法師,律法監護人…

「曉得耶穌回答的好」,覺得自己有能力/資格去評論耶穌?

馬太福音22:34-40

34法利賽人聽見耶穌堵住了撒都該人的口,他們就聚集. 35內中有一個人是律法師,要試探耶穌,就問他說: 36「夫子,律法上的誡命,哪一條是最大的呢?」 37耶穌對他說:「你要盡心,盡性,盡意愛主－你的上帝. 38這是誡命中的第一,且是最大的. 39其次也相倣,就是要愛人如己. 40這兩條誡命是律法和先知一切道理的總綱.」

V35試探耶穌…

「誡命中哪是第一要緊的呢?」

為何有這問題?期望得到什麼答案?試探點何在?

第一要緊=比較不同律法>>借此指控耶穌看輕律法

耶穌回答申命記6:4-9

4「以色列啊,你要聽!耶和華－我們上帝是獨一的主. 5你要盡心,盡性,盡力愛耶和華－你的上帝. 6我今日所吩咐你的話都要記在心上, 7也要殷勤教訓你的兒女.無論你坐在家裡,行在路上,躺下,起來,都要談論. 8也要繫在手上為記號,戴在額上為經文； 9又要寫在你房屋的門框上,並你的城門上.」

Shama(你要聽),猶太人每日的禱文

Tefillin護經匣綁在額上,Mezuzah門框經文盒

利未記19:18不可報仇,也不可埋怨你本國的子民,卻要愛人如己.我是耶和華.

第一要緊?

愛上帝,愛人

是1條還是2條誡命?

二合一的誡命

32那文士對耶穌說:「夫子說,上帝是一位,實在不錯；除了他以外,再沒有別的上帝； 33並且盡心,盡智,盡力愛他,又愛人如己,就比一切燔祭和各樣祭祀好的多.」

文士對耶穌的肯定/贊許

更再加例子以深化內容

34耶穌見他回答的有智慧,就對他說:「你離上帝的國不遠了.」從此以後,沒有人敢再問他甚麼.

耶穌的回應

「你離上帝的國不遠了.」

是讚賞嗎?是我們所期望耶穌對我的描述嗎?…真的是讚賞嗎?

眾人的反應

為何「不敢再問」?因為十分滿意?因為無話可說?

回想文士為何與耶穌開啟這對話

28有一個文士來,聽見他們辯論,曉得耶穌回答的好…

馬可福音12:18-27

18撒都該人常說沒有復活的事.他們來問耶穌說: 19「夫子,摩西為我們寫著說:『人若死了,撇下妻子,沒有孩子,他兄弟當娶他的妻,為哥哥生子立後.』 20有弟兄七人,第一個娶了妻,死了,沒有留下孩子. 21第二個娶了她,也死了,沒有留下孩子.第三個也是這樣. 22那七個人都沒有留下孩子；末了,那婦人也死了. 23當復活的時候,她是哪一個的妻子呢?因為他們七個人都娶過她.」 24耶穌說:「你們所以錯了,豈不是因為不明白聖經,不曉得上帝的大能嗎? 25人從死裡復活,也不娶也不嫁,乃像天上的使者一樣. 26論到死人復活,你們沒有念過摩西的書荊棘篇上所載的嗎?上帝對摩西說:『我是亞伯拉罕的上帝,以撒的上帝,雅各的上帝.』 27上帝不是死人的上帝,乃是活人的上帝.你們是大錯了.」

耶穌直接指出:「你們所以錯了,豈不是因為不明白聖經,不曉得上帝的大能嗎?… 上帝不是死人的上帝,乃是活人的上帝.你們是大錯了.」

撒都該人是一群拒絕文士們所累積的口傳律法,自稱只接受律法書上所寫的,否認身體復活,不相信天使和靈魂的存在.

文士見到耶穌如何用智慧/聖經知識,辯倒他們其中一群可能的敵對者.

以為耶穌在他們一方?

認同耶穌解釋聖經/信仰的智慧?

耶穌也認同文士所提出內容的價值及意義

但明顯,最終文士也得不到耶穌「完全」的認同.

馬可福音12:38-40

38耶穌在教訓之間,說:「你們要防備文士；他們好穿長衣遊行,喜愛人在街市上問他們的安, 39又喜愛會堂裡的高位,筵席上的首座. 40他們侵吞寡婦的家產,假意作很長的禱告.這些人要受更重的刑罰!」

因為文士只在知識/頭腦上理解,認識上帝,這頂多是離上帝的國不遠.

行得近一點上帝的國,當然可以是一件美事,也是信仰中應該有的進程.

但這完全不足夠!

亦不是目的地

我們要進到上帝的國,而不是只走近上帝的國!!!

不少追求信仰的人焦點模糊,以為信仰只是學道理,學知道

更期望用我的「智慧」來「理解」上帝,甚至以為可以用我的「智慧」完全「理解」上帝,這是錯了,更是大錯特錯了!

信仰要有參與的部份,上帝的國有我份,是因為我參與在其中.

在地上生活時就可以,以基督的心共建天國.

主禱文:願你的國降臨,願你的旨意行在地上,如同行在天上…

這就是進入天國的關鍵

弟兄姊妹,上帝的國離我們有多遠呢?我們是否願意,又準備好更進去呢?



\newpage

\section{寡婦和信仰}
\label{sec:722}
\textbf{李均熊牧師 - 寡婦和信仰}
\newline
\newline
\hyperref[sec:721]{< < < PREV SERMON < < <}
~
\hyperref[sec:toc]{[返目錄]}
~
\hyperref[sec:723]{> > > NEXT SERMON > > >}


講題:寡婦和信仰 Widow and Faith

講員:李均熊牧師

場合:香港中文大學崇基學院禮拜堂主日崇拜

日期:2024年11月10日

起:寡婦和信仰

今日可以說是寡婦主日,除了這段時間都是選讀希伯來書的書信經課,其他經課選段都是和寡婦有關.舊約經課主要選讀路得記中拿俄米對路得的一段說話,和這兩個寡婦——路得和拿俄米——的結局.當中的曲折,有機會再和大家分享.啟應詩篇146:9的一句:「耶和華保護寄居的,扶持孤兒和寡婦,卻使惡人的道路彎曲」,道出聖經一貫論到上主如何顧念社會中沒有依靠的人:寄居的,孤兒和寡婦.福音經課更是耶穌稱讚一個寡婦奉獻的兩個小錢:「比眾人所投的更多」.

我今日想集中宣講的,卻不是「寡婦的信仰」,而是「寡婦和信仰」.我不是想大家留意「寡婦的信仰」,所以不會講拿俄米或路得如何在沒有依靠之下仍努力掙扎,最後如何得到上主的祝福；也不是想討論耶穌所稱讚的這個寡婦如何透過將自己生活所有的奉給上主,而成為我們信仰學效的榜樣.我講道的題目不是「寡婦的信仰」,卻是「寡婦和信仰」,正是希望大家反思一下,正因為社會中有著寡婦和一些沒有依靠的人,怎樣影響我們的信仰.

承:寡婦的處境

舊約申命記的律法,多數將寡婦和孤兒,寄居者,劃為同一類人:

10章17-18

17因為耶和華－你們的神是萬神之神,萬主之主,是偉大,強有力,可畏的神,不看人的情面,也不受賄賂.18他為孤兒寡婦伸冤,愛護寄居的,賜給他衣食.」

24章17-21

17不可對寄居的和孤兒屈枉正直,也不可拿寡婦的衣服作抵押.18要記得你曾在埃及作過奴僕,耶和華－你的神從那裡救贖了你,所以我吩咐你遵行這事.19「你在田間收割莊稼,若忘了一捆在田間,就不要再回去拿,要留給寄居的,孤兒和寡婦；好讓耶和華－你的神在你手裡所做的一切,賜福給你.20你打了橄欖樹,枝上剩下的不可再打,要留給寄居的,孤兒和寡婦.21你摘葡萄園的葡萄,掉落的不可拾取,要留給寄居的,孤兒和寡婦.22你要記得你曾在埃及地作過奴僕,所以我吩咐你遵行這事.」

他們和一般的窮人是有所分別的.甚至有時寄居的,孤兒和寡婦,和利未人看成同一類:

26: 12-13

12每逢第三年,就是捐十分之一的那年,你從你一切土產中取了十分之一,要分給利未人,寄居的,孤兒和寡婦,使他們在你的城鎮中可以吃得飽足.13你又要在耶和華－你神面前說:『我已將聖物從家裡拿出來,給了利未人,寄居的,孤兒和寡婦,是遵照你吩咐我的一切命令.你的命令,我沒有違背,也沒有忘記.』

為甚麼申命記將寡婦和利未人相提並論呢?因為在申命記內利未人是住在聖殿供職的人士,是沒有祖宗地業可以繼承的,正如一些寡婦和孤兒,以及寄居者,在社會中沒有屬於自己地業可以耕作,自給自足.原來當舊約提到寡婦時,其實也有三種說法:一種說法是「死人的妻」,重點是在沒有生育能繼承祖宗產業子裔就成了寡婦,這些寡婦是可以由丈夫的兄弟再娶,目的是「為死人立後」,以至祖宗的產業可以有人繼承.路得就是這樣的寡婦；另一種說法是「作為妻子的寡婦」(原文只是「妻子-寡婦」),這說法表達的是這個寡婦已透過兒子全面承繼先夫祖宗的產業,列王記上11:26就是這樣稱呼耶羅波安的母親.最後,就是一般的「寡婦」,就是那些因種種理由沒有祖宗產業可以承繼.故此利未人,寄居的,孤兒和寡婦,都是那些在以色列人社會中沒有祖宗產業,對將來也沒有保障的一群人.

我想不用我們穿越二,三千年,我們也可以想到這些寡婦就是無依無靠,生活與未來也沒有保障的人.用今日的說法,就是那「跌出了社會安全網」的人.現代社會通常透過社會福利,建立社會的安全網,去照顧那些在社會中無依無靠的人.但在昔日以色列社會以氏族為社會結構核心的古代世界,沒有或不可能繼承丈夫祖宗產業的,就不屬於家族的成員,族群是沒有責任繼續照顧這個寡婦的,他們有困難時家族也可以不為他們出頭.當然出於良心仍然照料這些寡婦的家族仍是會有,但社會或多或少出現對欺壓孤兒寡婦的一類事件,故此在出埃及記中,上主就明明的吩咐:

22:21-24

21不可虧待寄居的,也不可欺壓他,因為你們在埃及地也作過寄居的.22不可苛待寡婦和孤兒；23若你確實苛待他,他向我苦苦哀求,我一定會聽他的呼求,24並要發烈怒,用刀殺你們,使你們的妻子成為寡婦,兒女成為孤兒.」

正因為上主對這些無權無依者的關心,慢慢聖殿就發展成主要照顧這些寡婦孤兒的核心.在古代如是,而在今日社會,教會也常充當這種角色.在社會未全面發展,社會福利欠缺的時候,教會就成為那些無依無靠者的保護者和供應者,因為我們的主「為孤兒寡婦伸冤,愛護寄居的,賜給他衣食.」我也不必重複多講上世紀四,五,六十年代教會來港的難民派物資,辦學校,和開創各樣社會福利的事業.就是到今日,我們又會否見到社會中的寡婦呢?我的意思是,那些今日無權無勢,無依無靠,甚至有時跌出政府或社會設置的安全網的人群.我心目中今日的所謂「寡婦」,並不是那些窮困的人,而是被社會忽略,失去尊嚴,也失去親友依靠的人.很多無家者,他們並非沒有金錢,大多數是因為種種原因,沒有了人和人之間的連系；獨居的長者,沒有了家人,孤獨的面對生活種種的困難和挑戰；基層的清潔工,他們或因工作的性質被人看輕,在超時工作,極低工資的處境下艱辛努動的人.今日你看到我們身邊的「寡婦」嗎?

可以告訴你今日有不同的教會在協助和關心這些「寡婦」.我認識有傳道人,並不是在教會工作,而是充當地產經紀,為無家者配對,連繫幾個無家者好去找合適的住房租住,不單解決他們無家可歸的困境,更為他們重新連結社區和人群；我認識有傳道人,為獨居長者配對維修技工,而那些維修技工就是一些無家者,使老人家有人探望,使無家者感到有所貢獻；亦有教會為清潔工服務,開放他們的地方,使他們休息用膳時不用留在垃圾站或公廁,能有地方可去,又教他們申展運動,又為他們預備有營養的午餐,使他們感受到人的尊重.崇基學院也有服務學習計劃,也是很好的機會讓我們的同學看到今日社會中的「寡婦」,使無依者可以連系社區,和如何使他們有尊嚴地生活.

今日教會能成為無依無靠的「寡婦」的保護者和供應者.昔日,就是聖殿社群和當中領袖要承擔的角色.

轉／合:耶穌看見寡婦

最後,讓我們回到福音經課在聖殿中的一幕.越來越多新約學者關心的,不只是寡婦在不足中仍投入自己養生的二個小錢,而是這個寡婦這行動背後的處境.有人認為耶穌怎會欣賞一個兩餐不繼的窮寡婦,把用來生活的小錢投入聖殿的錢庫內呢?以耶穌在福音書內的表現,祂豈不應該關心寡婦的生活所需多於她投入多少金錢呢?近幾十年來,就有學者留意這一段的一上文下理,和耶穌對寡婦的稱讚格格不入.在接著的馬可福音13章,就是耶穌預言聖殿將要被毀,「將來在這裡沒有一塊石頭留在石頭上,不被拆毀了.」(可13:2)所以任為聖殿投入多少金錢都是無補於事,故此耶穌並非關心投入庫內多少的問題.而且今日經課其實可分為兩段,在耶穌對錢庫坐者看眾人怎樣投錢入庫之前,大家不要忽略耶穌其實正指斥那些文士「他們好穿長衣遊行,喜愛人在街市上問他們的安,又喜愛會堂裡的高位,筵席上的首座.他們侵吞寡婦的家產,假意作很長的禱告.這些人要受更重的刑罰!」這一些在聖殿中熟讀聖經,應該按聖經吩咐照顧寡婦的文士,現在卻是侵吞無權無勢無依無靠無人為他們發聲的寡婦的家產的人!所以有些聖經學者認為,耶穌的一句:「但這寡婦是自己不足,把她一切養生的都投上了.」語氣並非讚賞,而是心中悲憤的一種控訴!宗教領袖並沒有按他們熟識的律法要求,保障這些不受保障的人的生活,更在侵吞她們的產業之後,宗教領袖仍要求她們奉獻金錢,害她們連養生的兩個小錢都要花在今日華麗卻將要消失的聖殿上?豈有此理(有無攪錯／How can this be so)???

不過,經文並沒有透露寡婦在想甚麼,也沒有記錄耶穌說話時的神情語氣.是寡婦心甘情願,為著上主而放棄生活所需,將產生的二個小錢都要奉獻給上主?耶穌的一句「這窮寡婦投入庫裡的,比眾人所投的更多」在我看來的確帶著一點點讚嘆之情.但誰知道,是不是寡婦面對著身邊侵吞她產業的文士,想到還要將養生的都要投入而 心裡愁苦呢?以至耶穌義憤填膺,哀歎「(啍!)他們都是自己有餘,拿出來投在裡頭；但這寡婦是自己不足,把她一切養生的都投上了(有無攪錯!)!」

福音書記載的,是一段段關於耶穌的敘述,有時耶穌會明確叫我們做甚麼不做甚麼,但大部份情況並沒有明文要我們要怎樣想怎樣做.福音書是要樣我們透過耶穌的眼光看世界,有時留白沒有言明,就是讓我們帶著自己的生命進入耶穌的世界去思考人生.經課今天或許告訴我們「寡婦的信仰」,帶我們看一個弱少困苦的生命,如何仍可以敬虔,奉獻和犧牲；但我願大家思想「寡婦和信仰」,當我們進入耶穌的世界,我們看到我們身邊的寡婦,又看到寡婦的處境嗎?我們看到這個身邊不因生活壓迫和困難,仍堅守信仰,甘願奉獻和擺上所有的弟兄或姊妹,並學習自己也成為這樣的人；我們又會否同時看到身旁困苦無助,無權無勢,無依無靠的眾生,讓我們細心思考,我們的主耶穌基督今日會期望我們如何回應?



\newpage

\section{耶和華的日子臨近了}
\label{sec:723}
\textbf{鄧瑞強博士 - 耶和華的日子臨近了}
\newline
\newline
\hyperref[sec:722]{< < < PREV SERMON < < <}
~
\hyperref[sec:toc]{[返目錄]}
~
\hyperref[sec:724]{> > > NEXT SERMON > > >}


約珥書 1:4-2:32
\newline
\begin{longtable}{cl}
\hline
\hline
章節 & 經文 (和合本修訂版)\\
\hline
1:4 & \begin{tabularx}{0.7\textwidth}{X} 剪蟲吃剩的，蝗蟲來吃；蝗蟲吃剩的，蝻子來吃；蝻子吃剩的，螞蚱來吃。 \end{tabularx} \\ \\ \relax
1:5 & \begin{tabularx}{0.7\textwidth}{X} 醉酒的人哪，要清醒，要哭泣；好酒的人哪，都要為甜酒哀號，因為酒從你們的口中斷絕了。 \end{tabularx} \\ \\ \relax
1:6 & \begin{tabularx}{0.7\textwidth}{X} 有一隊蝗蟲，強盛且不可數，上來侵犯我的地；牠的牙齒如獅子的牙齒，如母獅的大牙。 \end{tabularx} \\ \\ \relax
1:7 & \begin{tabularx}{0.7\textwidth}{X} 牠毀壞我的葡萄樹，撕裂我的無花果樹，剝光又丟棄，使枝條露白。 \end{tabularx} \\ \\ \relax
1:8 & \begin{tabularx}{0.7\textwidth}{X} 你要像童女腰束麻布，為她年少時的丈夫哀號。 \end{tabularx} \\ \\ \relax
1:9 & \begin{tabularx}{0.7\textwidth}{X} 耶和華的殿中斷絕素祭和澆酒祭，事奉耶和華的祭司都悲哀。 \end{tabularx} \\ \\ \relax
1:10 & \begin{tabularx}{0.7\textwidth}{X} 田荒涼，地悲哀；因為五穀毀壞，新酒枯竭，新的油也缺乏。 \end{tabularx} \\ \\ \relax
1:11 & \begin{tabularx}{0.7\textwidth}{X} 農夫啊，你們要慚愧；修整葡萄園的啊，你們要哀號；因為大麥、小麥與田間的莊稼全都毀了。 \end{tabularx} \\ \\ \relax
1:12 & \begin{tabularx}{0.7\textwidth}{X} 葡萄樹枯乾，無花果樹衰殘，石榴樹、棕樹、蘋果樹，田野一切的樹木都枯乾；眾人的喜樂盡都消逝。 \end{tabularx} \\ \\ \relax
1:13 & \begin{tabularx}{0.7\textwidth}{X} 祭司啊，當束上麻布痛哭；事奉祭壇的啊，要哀號；事奉我神的啊，你們要來，披上麻布過夜，因為在你們神的殿中不再有素祭和澆酒祭了。 \end{tabularx} \\ \\ \relax
1:14 & \begin{tabularx}{0.7\textwidth}{X} 你們要使禁食的日子分別為聖，宣告嚴肅會，召集長老和這地所有的居民來到耶和華－你們神的殿，向耶和華哀求。 \end{tabularx} \\ \\ \relax
1:15 & \begin{tabularx}{0.7\textwidth}{X} 哀哉，這日子！因為耶和華的日子臨近，好像毀滅從全能者來到。 \end{tabularx} \\ \\ \relax
1:16 & \begin{tabularx}{0.7\textwidth}{X} 糧食不是在我們眼前斷絕了嗎？歡喜快樂不是從我們神的殿中止息了嗎？ \end{tabularx} \\ \\ \relax
1:17 & \begin{tabularx}{0.7\textwidth}{X} 種子在土塊下朽爛，倉荒涼，廩破壞，因為五穀枯乾了。 \end{tabularx} \\ \\ \relax
1:18 & \begin{tabularx}{0.7\textwidth}{X} 牲畜哀鳴，牛群混亂，因無草場，羊群也受苦。 \end{tabularx} \\ \\ \relax
1:19 & \begin{tabularx}{0.7\textwidth}{X} 耶和華啊，我向你求告，因為有火吞噬野地的草場，火焰燒盡田野的樹木。 \end{tabularx} \\ \\ \relax
1:20 & \begin{tabularx}{0.7\textwidth}{X} 田野的走獸切慕你，因為溪水乾涸，火吞噬了野地的草場。 \end{tabularx} \\ \\ \relax
2:1 & \begin{tabularx}{0.7\textwidth}{X} 你們要在錫安吹角，在我的聖山發出警報。這地所有的居民要發顫，因為耶和華的日子快到，已經臨近了。 \end{tabularx} \\ \\ \relax
2:2 & \begin{tabularx}{0.7\textwidth}{X} 那是黑暗、陰森的日子，是密雲、烏黑的日子，如同黎明籠罩山嶺。有一隊蝗蟲，又大又強，自古以來沒有像這樣的，以後直到萬代也必沒有。 \end{tabularx} \\ \\ \relax
2:3 & \begin{tabularx}{0.7\textwidth}{X} 牠們前面有火吞噬，後面有火焰燒盡。牠們未到以前，地如伊甸園，過去以後，卻成了荒涼的曠野，沒有一樣能躲避牠們。 \end{tabularx} \\ \\ \relax
2:4 & \begin{tabularx}{0.7\textwidth}{X} 牠們形狀如馬，奔跑如戰馬。 \end{tabularx} \\ \\ \relax
2:5 & \begin{tabularx}{0.7\textwidth}{X} 響聲如戰車在山頂上跳動，如火焰吞噬碎秸，好像強大的軍隊擺陣備戰。 \end{tabularx} \\ \\ \relax
2:6 & \begin{tabularx}{0.7\textwidth}{X} 在牠們面前，萬民傷慟，臉都變色。 \end{tabularx} \\ \\ \relax
2:7 & \begin{tabularx}{0.7\textwidth}{X} 牠們如勇士奔跑，如戰士攀登城牆，各行於自己的道路，不亂隊伍； \end{tabularx} \\ \\ \relax
2:8 & \begin{tabularx}{0.7\textwidth}{X} 牠們並不彼此推擠，各行於自己的大道，衝過防禦，並不停止。 \end{tabularx} \\ \\ \relax
2:9 & \begin{tabularx}{0.7\textwidth}{X} 牠們蹦上城，跳上牆，爬上房屋，從窗戶進來，如同盜賊。 \end{tabularx} \\ \\ \relax
2:10 & \begin{tabularx}{0.7\textwidth}{X} 在牠們面前，地動天搖，日月昏暗，星宿無光。 \end{tabularx} \\ \\ \relax
2:11 & \begin{tabularx}{0.7\textwidth}{X} 耶和華在他的軍旅前出聲，他的隊伍龐大，遵行他命令的強盛。耶和華的日子大而可畏，誰能當得起呢？ \end{tabularx} \\ \\ \relax
2:12 & \begin{tabularx}{0.7\textwidth}{X} 然而你們現在要禁食，哭泣，哀號，一心歸向我。這是耶和華說的。 \end{tabularx} \\ \\ \relax
2:13 & \begin{tabularx}{0.7\textwidth}{X} 你們要撕裂心腸，不要撕裂衣服。歸向耶和華－你們的神，因為他有恩惠，有憐憫，不輕易發怒，有豐盛的慈愛，並且會改變心意，不降那災難。 \end{tabularx} \\ \\ \relax
2:14 & \begin{tabularx}{0.7\textwidth}{X} 誰知道他也許會回心轉意，留下餘福，就是獻給耶和華－你們神的素祭和澆酒祭。 \end{tabularx} \\ \\ \relax
2:15 & \begin{tabularx}{0.7\textwidth}{X} 你們要在錫安吹角，使禁食的日子分別為聖，宣告嚴肅會。 \end{tabularx} \\ \\ \relax
2:16 & \begin{tabularx}{0.7\textwidth}{X} 聚集百姓，使會眾自潔；召集老年人，聚集孩童和在母懷吃奶的；使新郎出內室，新娘離開洞房。 \end{tabularx} \\ \\ \relax
2:17 & \begin{tabularx}{0.7\textwidth}{X} 事奉耶和華的祭司要在走廊和祭壇間哭泣，說：「耶和華啊，求你顧惜你的百姓，不要使你的產業受羞辱，在列國中成為笑柄。為何讓人在萬民中說『他們的神在哪裡』呢？」 \end{tabularx} \\ \\ \relax
2:18 & \begin{tabularx}{0.7\textwidth}{X} 耶和華為自己的地發熱心，憐憫他的百姓。 \end{tabularx} \\ \\ \relax
2:19 & \begin{tabularx}{0.7\textwidth}{X} 耶和華應允他的百姓說：「看哪，我要賞賜你們五穀、新酒和新的油，使你們飽足，我必不再使你們受列國的羞辱。 \end{tabularx} \\ \\ \relax
2:20 & \begin{tabularx}{0.7\textwidth}{X} 我要使北方來的隊伍遠離你們，將他們趕到乾旱荒蕪之地：前隊趕入東海，後隊趕入西海；臭氣上升，惡臭騰空。耶和華果然行了大事！ \end{tabularx} \\ \\ \relax
2:21 & \begin{tabularx}{0.7\textwidth}{X} 「土地啊，不要懼怕，要歡喜快樂，因為耶和華行了大事。 \end{tabularx} \\ \\ \relax
2:22 & \begin{tabularx}{0.7\textwidth}{X} 田野的走獸啊，不要懼怕，因為曠野的草已生長，樹木結果，無花果樹、葡萄樹也都效力。 \end{tabularx} \\ \\ \relax
2:23 & \begin{tabularx}{0.7\textwidth}{X} 「錫安的民哪，你們要歡喜，要因耶和華－你們的神快樂；因他賞賜你們合宜的秋雨，為你們降下甘霖，秋雨和春雨，和先前一樣。 \end{tabularx} \\ \\ \relax
2:24 & \begin{tabularx}{0.7\textwidth}{X} 「禾場充滿五穀，池中漫溢新酒和新的油。 \end{tabularx} \\ \\ \relax
2:25 & \begin{tabularx}{0.7\textwidth}{X} 我差遣到你們中間的大軍隊，就是蝗蟲、蝻子、螞蚱、剪蟲，那些年間所吃的，我要補還給你們。 \end{tabularx} \\ \\ \relax
2:26 & \begin{tabularx}{0.7\textwidth}{X} 「你們必吃得飽足，讚美耶和華－你們神的名，他為你們行了奇妙的事。我的百姓不致羞愧，直到永遠。 \end{tabularx} \\ \\ \relax
2:27 & \begin{tabularx}{0.7\textwidth}{X} 你們必知道我是在以色列中，又知道我是耶和華－你們的神，沒有別的。我的百姓不致羞愧，直到永遠。」 \end{tabularx} \\ \\ \relax
2:28 & \begin{tabularx}{0.7\textwidth}{X} 「以後，我要將我的靈澆灌凡有血肉之軀的。你們的兒女要說預言，你們的老人要做異夢，你們的少年要見異象。 \end{tabularx} \\ \\ \relax
2:29 & \begin{tabularx}{0.7\textwidth}{X} 在那些日子，我要將我的靈澆灌我的僕人和婢女。 \end{tabularx} \\ \\ \relax
2:30 & \begin{tabularx}{0.7\textwidth}{X} 「我要在天上地下顯出奇事，有血，有火，有煙柱。 \end{tabularx} \\ \\ \relax
2:31 & \begin{tabularx}{0.7\textwidth}{X} 太陽要變為黑暗，月亮要變為血，這都在耶和華大而可畏的日子未到以前。 \end{tabularx} \\ \\ \relax
2:32 & \begin{tabularx}{0.7\textwidth}{X} 那時，凡求告耶和華名的就必得救；因為照耶和華所說的，在錫安山，在耶路撒冷將有逃脫的人。凡耶和華所召的，都在餘民之列。」 \end{tabularx} \\ \\
[1ex]
\hline
\hline
\end{longtable}
講題:耶和華的日子臨近了 The Day of the Lord is Near

經文:約珥書1:4-2:32

講員:鄧瑞強博士

場合:香港中文大學崇基學院禮拜堂主日崇拜

日期:2024年11月17日

各位弟兄姊妹,早安.願神的恩典常與你們同在.

今日,我講12小先知書的第二卷:約珥書.講先知書時,我通常會交代一下時代背景,講一些年份,弟兄姊妹聽到這些資料,很頭痛,我今日不講了,因為約珥書本身沒有給我們任何時代背景的提示,而學者們對這問題也沒有什麼確定看法.

約珥書的篇幅很短,但先知的心靈很廣闊.他從大自然的災禍中,從世間的人禍中,看到上帝的天意.從上帝的天意裡,看到人類應負的責任.大地,人類,上帝,有一種緊密的聯繫,人從來都不是孤零零地走在無情的大地之上的.大地感應人的善惡,上帝透過大地發出祂的警告.人在大地的警告中,審視自己,面向上帝.若然人甘願謙卑下來,走回正途,則上帝會以仁愛回應,讓大地與人重拾光明.

我們看看約珥書的經文.

珥1:4 剪蟲剩下的,蝗蟲來吃；蝗蟲剩下的,蝻子來吃；蝻子剩下的,螞蚱來吃.

珥1:5 酒醉的人哪,要清醒哭泣(醒過來,哭吧)；好酒的人哪,都要為甜酒哀號,因為從你們的口中斷絕了.

剪蟲剩下的,蝗蟲來吃；蝗蟲剩下的,蝻子來吃；蝻子剩下的,螞蚱來吃.簡單幾隻字,便交代了一場浩劫,是蝗災.

最近期的蝗災,是2020年從東非蔓延到印度的蝗災.這些蝗蟲,可以在3個月內繁殖20倍,密度達到每平方公里8000萬隻,每日吃掉3萬5千人一天的食糧.(參考:\url{https://www.bbc.com/future/article/20200806-the-biblical-east-african-locust-plagues-of-2020})

好好的一個猶大國,好好的一個地方,一場蝗災,轉瞬間,便成為荒土.世情的變化,真是預想不到.本來,可能是等著豐收的,但一下子,便要準備捱飢抵餓.

珥1:5 酒醉的人哪,要清醒哭泣〔醒過來,哭吧〕；好酒的人哪,都要為甜酒哀號,因為從你們的口中斷絕了.

面對困境,對上天毫無感覺的人,應該醒一醒了.

經文繼續描述災情.

珥1:11 農夫啊,你們要慚愧〔失望〕；修理葡萄園的啊,你們要哀號；因為大麥小麥與田間的莊稼都滅絕了.

珥1:12 葡萄樹枯乾；無花果樹衰殘.石榴樹,棕樹,蘋果樹,連田野一切的樹木也都枯乾；眾人的喜樂盡都消滅.

珥1:13 祭司啊,你們當腰束麻布痛哭；伺候祭壇的啊,你們要哀號；事奉我上帝的啊,你們要來披上麻布過夜,因為素祭和奠祭從你們上帝的殿中斷絕了.

珥1:14 你們要分定禁食的日子,宣告嚴肅會,招聚長老和國中的一切居民到耶和華─你們上帝的殿,向耶和華哀求.

珥1:15 哀哉!耶和華的日子臨近了.這日來到,好像毀滅從全能者來到.

珥1:16 糧食不是在我們眼前斷絕了嗎?歡喜快樂不是從我們上帝的殿中止息了嗎?

大麥,小麥,葡萄樹,無花果樹,石榴樹,棕樹,蘋果樹,連田野一切的樹木也都枯乾了,一切食物都被蝗蟲食光了,經濟活動停頓了,民間的歡樂沒有了.

從大自然的災劫裡,先知看到,這是「耶和華的日子臨近」.這就是今日講道的主題.耶和華的日子,是永恆的上帝明明白白的介入人類歷史的日子.耶穌誕生的日子,就是這類日子.這些日子,讓我們明白,人類歷史不是純粹由人類的手打造出來的.這些日子,讓我們停下來,看看我們將歷史搞到一個什麼樣子.上帝介入歷史,帶來歷史的轉折,好讓人思考生命應有的新意.先知約珥說:「哀哉!耶和華的日子臨近了.這日來到,好像毀滅從全能者來到.」(珥1:15)在先知眼中,今次的蝗災,是上帝審判的作為.大自然的災劫,背後是上帝的懲罰.

上帝的臨在,總有兩面.對邪惡的人來說,上帝的臨在是審判.對謙虛向善的人,卻是安慰與拯救.如此說來,你可能說,那就請上帝不要來到我人生裡了,免受審判.但我要問,為何不正正因為上帝要來到你的人生裡,你就及早回歸正途?

人類的知識越來越利害了,我們有生物工程,可改變物種的基因,說不定,有一天,我們能製造出某種超級物種來.現在,我們能遠距離控制傳呼機,令它們爆炸,說不定,有一天,你的手機會出現什麼情況.人類的操控能力越來越利害,或許會以為,從今以後,人類能掌握歷史的未來了.但在先知眼中,若人自以為是而作惡,則大自然會作出警告,並面對上帝的臨在.上帝的臨在,是要人反省:人究竟是在行善,抑或作惡呢?

約珥書第1章的經文,是叫人認真看看世事的變化,從而領悟上帝的臨在.第2章的經文,是叫人細心聽聽上帝的要求.

珥2:12 耶和華說:雖然如此,你們應當禁食,哭泣,悲哀,一心歸向我.

珥2:13 你們要撕裂心腸,不撕裂衣服.歸向耶和華─你們的上帝；因為他有恩典,有憐憫,不輕易發怒,有豐盛的慈愛,並且後悔不降所說的災.

珥2:14 或者他轉意後悔,留下餘福,就是留下獻給耶和華─你們上帝的素祭和奠祭,也未可知.

上帝不希望人在驕傲中墮落,祂呼喚人謙卑,悔改.悔改有用嗎?有用的.這與上帝的性情有關.上帝不是法官,不是按你犯的錯要你一一十足補償.祂是慈父,祂基本的態度就是接納你,祂是在接納你的心情下對你的錯作出警戒.祂的審判,終歸而言,不是叫你受罰,而是叫你走上正途.慈父的懲戒,都是為兒女好的,不是要摧毀他們的.

約珥書講到上帝的性情,「祂有恩典,有憐憫,不輕易發怒,有豐盛的慈愛」.在古時,在西乃山上,上帝向摩西頒佈十誡,而以色列民就在山下拜金牛贖.下山後,摩西很忿怒,便摔碎寫著十誡的法版.後來,上帝叫摩西再上西乃山,將十誡再寫在法版上,這時,上帝說:「耶和華,耶和華,是有憐憫有恩典的神,不輕易發怒,並有豐盛的慈愛和誠實.」(出34:6)縱使祂的子民遠離祂,上帝仍按自己的恩典親近祂的子民.

在另一卷小先知書約拿書裡,也出現過同一表達.那時,上帝叫先知約拿去亞述的尼尼微城宣告審判.尼尼微城全城悔改,上帝就後悔,不降災與他們.這時,先知約拿便大發脾氣.他對上帝說:「...我知道你是有恩典,有憐憫的神,不輕易發怒,有豐盛的慈愛,並且後悔不降所說的災,所以我急速逃往他施去.」(拿4:2)約拿早知上帝就是這樣的上帝,故他逃避向尼尼微城宣告上帝的信息,尼尼微城是以色列的敵人,約拿很不希望上帝放過他們.明顯,就算不是上帝的子民,就算是上帝子民的敵人,只要他們悔改,上帝依然善待他們,這是上帝的本性.面對惡人,上帝就算警戒他們,上帝的本意仍是希望他們回歸正道而得享福樂.面對大自然的災害,面對社會的邪惡,先知約珥呼喚人悔改,而他這樣積極呼籲,而不是心灰意冷,是他相信人生的底蘊,是那極端希望我們在美善中得福樂的大愛.浪子只要肯回頭,便會看到,他的父親早已在遠遠的地方跑過來,擁抱他.

珥2:15 你們要在錫安吹角,分定禁食的日子,宣告嚴肅會.

珥2:16 聚集眾民,使會眾自潔:招聚老者,聚集孩童和吃奶的；使新郎出離洞房,新婦出離內室.

珥2:17 事奉耶和華的祭司要在廊子和祭壇中間哭泣,說:耶和華啊,求你顧惜你的百姓,不要使你的產業受羞辱,列邦管轄他們.為何容列國的人說:「他們的上帝在哪裡」呢?

先知繼續呼籲不同的人悔改.老的,嫩的,都是悔改.有趣的是,要「使新郎出離洞房,新婦出離內室」.舊約對新婚夫婦,很多時都很優待,連打仗都不用打.但在悔改的事上,他們也不能置身事外.在這裡,祭司求上帝顧念他們,提及他們的「產業受羞辱」,「列邦管轄他們」.看來,他們不單面對大自然災害,還面對人禍,面對戰敗後被列強凌辱的淒涼.在天災,人禍下,人很易以為上帝離棄了他們.但祭司的信念是:不要讓人以為上帝離棄了我們.這是一種何等大的信心.

當人願意悔改時,上帝便顯明祂的慈父性格了.

珥2:18 耶和華就為自己的地發熱心,憐恤他的百姓.

珥2:19 耶和華應允他的百姓說:我必賜給你們五穀,新酒,和油,使你們飽足；我也不再使你們受列國的羞辱.

珥2:18是天地改變的轉折點,因百姓的悔改,上帝便轉向他們,使祂的子民從蝗災中恢復過來,也使他們不再受列強欺凌.

我們今日讀的最後經文,是上帝再次明顯的臨在,這次,祂以恩慈憐憫祂的子民.

珥2:30 「在天上地下,我要顯出奇事,有血,有火,有煙柱.

珥2:31 日頭要變為黑暗,月亮要變為血,這都在耶和華大而可畏的日子未到以前.

珥2:32 到那時候,凡求告耶和華名的就必得救；因為照耶和華所說的,在錫安山,耶路撒冷必有逃脫的人,在剩下的人中必有耶和華所召的.」

在這裡,很多上天下地的怪異事件,這是聖經文學的表達手法,表明上帝的臨在.耶和華的日子再次出現,這是上帝顯明祂的大能的特別日子,在這裡,特別提到祂要照顧在戰亂中倖存的人.世事紛亂,人心惶惶,上帝臨在,為要解救痛苦的人.

來到新約時代,「耶和華的日子」演變成「耶穌再來的日子」.有了這日子,人面對歷史的將來,就不是空白一片,而是面對上帝的來臨.現在,可能有天災,有人禍,迎向我們的,卻不是無意義的黑暗一片,而是上帝帶來的光明,故此,無論遇到什麼絕境,都不應絕望而自殺,因生命的底蘊,不是無意義的空白,而是那十足支撐我們的仁愛天父.因為有「耶和華的日子」,有「耶穌再來的日子」,有這充滿上帝光明的明天,我們今日就不要胡亂地活,不要自求當下快樂,不要以不義待人,而以正直的生命,迎向來臨中的上帝.

對邪惡的人來說,「耶穌再來的日子」是審判之日.對謙虛而努力向善的人來說,「耶穌再來的日子」是生命得以圓滿之時.面對天災,人禍,先知約珥呼籲神的子民,去思考「耶和華的日子」.今日,或許,你覺得人生艱難,前路茫茫,我邀請你,去思考「耶穌再來的日子」,希望你明白,歷史的將來,仍由仁愛的上帝掌管,因而請你懷著希望,努力生活,活出仁愛.

願上主的大愛臨到你,願你回應上主的愛,得享滿足,平安.



\newpage

\section{在亂世中再思基督君王的意義}
\label{sec:724}
\textbf{林豪恩先生 - 在亂世中再思基督君王的意義}
\newline
\newline
\hyperref[sec:723]{< < < PREV SERMON < < <}
~
\hyperref[sec:toc]{[返目錄]}
~
\hyperref[sec:725]{> > > NEXT SERMON > > >}


講題:在亂世中再思基督君王的意義 Reflecting the Meaning of Christ the King in Turbulent Times

講員:林豪恩先生

場合:香港中文大學崇基學院禮拜堂主日崇拜

日期:2024年11月24日

一. 亂世

本主日是基督君王日,宣道講題是「在亂世中再思基督君王的意義」.這些年頭世界局勢如何呢?今年的聯合國大會期間,有本地報章發表評論如下:

「今年的「聯大」非比尋常,俄烏衝突沒有停息跡象,以色列與哈馬斯激戰未止,又開啟與黎巴嫩的新戰線,俄羅斯還威脅使用核武器.中美關係究竟是後冷戰還是新冷戰難以界定,這個世界冷戰與熱戰硝煙瀰漫.而「聯大」仍然是各說各話,全球迫在眉睫的氣候變化,以及發展中國家無以發展的問題,被與會國代表互相指摘,謾罵,威脅之聲淹沒.」【明報社評2024年9月30日】

二. 亂世中人們的盼望:強大的王

常言道,亂世出英雄,或許處於亂世的人們,心中會盼望著偉大領袖的出現,強而有力地保護人民.從今天的舊約經課中,我們也可以看見類似的盼望.

撒母耳記下 23

23:3以色列的上帝,以色列的磐石曉諭我說:那以公義治理人民的,敬畏上帝執掌權柄,

23:4他必像日出的晨光,如無雲的清晨,雨後的晴光,使地發生嫩草.

公義正真,能夠撥亂反正的君王是亂世中人民的盼望,然而,就連大衛自己也如此說:

23:5我家在上帝面前並非如此

如果連古代以色列人心中的偉大君王大衛都承認自己並非公義之君,不是仁愛之王,那麼,以色列的第一位君王,被他取代的掃羅王呢?撒母耳記上15:10-11節如此記載耶和華上帝的評價:

耶和華的話臨到撒母耳說:「我立掃羅為王,我後悔了；因為他轉去不跟從我,不遵守我的命令.」

既然掃羅令耶和華後悔,為甚麼當初又立他為王呢?古代以色列的君王制度是如何出現的呢?就是因為亂世,在亂世中人民盼望偉大的君王.撒母耳記上12:12透露了以色列君王制度出現的背景:

你們見亞捫人的王拿轄來攻擊你們,就對我說:『不,要有一個王治理我們.』

其實撒母耳當時已經苦口婆心向以色列人說明立王的後果:

8:11「管轄你們的王必這樣行:他必派你們的兒子為他趕車,跟馬,奔走在車前；

8:12又派他們作千夫長,五十夫長,為他耕種田地,收割莊稼,打造軍器和車上的器械；

8:13必取你們的女兒為他製造香膏,做飯烤餅；

8:14也必取你們最好的田地,葡萄園,橄欖園賜給他的臣僕.

8:15你們的糧食和葡萄園所出的,他必取十分之一給他的太監和臣僕；

8:16又必取你們的僕人婢女,健壯的少年人和你們的驢,供他的差役.

8:17你們的羊群他必取十分之一,你們也必作他的僕人.

既然王帝帶來奴役,當時的人為甚麼仍然要立王呢?都是因為亂世,在亂世中的百姓如此說:

8:19「不然!我們定要一個王治理我們,

8:20使我們像列國一樣,有王治理我們,統領我們,為我們爭戰.」

亂世中的百姓期待強大的君王,其實是盼望甚麼呢?從撒母耳時代的以色列人來看,他們盼望的是安全和保障,這樣的盼望是完全可以理解的.聖經作者如何評價呢?北國19個王,全部被評為「行耶和華眼中看為惡的事」；南國20個王,超過一半都被評為「行耶和華眼中看為惡的事」.

亂世中的人們盼望強大的君王是出於現實需要,這種需要是可以理解的.然而,真正能夠帶領人們跨越亂世,而不是為人們帶來奴役的君王何處尋呢?

三. 基督為王:不屬於這世界的王

歷世歷代,多少英雄為了爭奪君王寶座而打生打死?多少百姓為了各路英雄的君王夢而戰死沙場?古詩有云:一將功成萬骨枯.一千年後,又有一個被傳為君王的人物出現,就連羅馬帝國派到猶大地的總督也起了疑問.今天的福音經課約翰福音18章33節記下了彼拉多的疑問:

18:33彼拉多又進了衙門,叫耶穌來,對他說:「你是猶太人的王嗎?」

耶穌沒有直接回答是還是不是,卻反問彼拉多:

18:34「這話是你自己說的,還是別人論我對你說的呢?」

如此耐人尋味的回答,令彼拉多更疑惑:

18:35彼拉多說:「我豈是猶太人呢?你本國的人和祭司長把你交給我.你做了甚麼事呢?」

彼拉多可能想,如果耶穌是猶太人的王,猶太人應該不會把他捉拿來給自己.如果耶穌根本沒有動機要做猶太人的王,就不用勞煩自己這位總督大人出手.耶穌不直接解答彼拉多的疑惑,反而繼續帶他去「遊花園」.

18:36耶穌回答說:「我的國不屬這世界；我的國若屬這世界,我的臣僕必要爭戰,使我不至於被交給猶太人.只是我的國不屬這世界.」

究竟耶穌有沒有回答彼拉多的問題呢?可以說沒有,也可以說有.「你是猶太人的王嗎?」可以說不是,又可以說是.如果用「這世界」來說,耶穌不是猶太人的王；如果用「我的國」來說,耶穌一定是猶太人的王.其實,耶穌的回答可以繼續擴展下去:「你是羅馬人的王嗎?」可以說不是,又可以說是.如果用「這世界」來說,耶穌不是羅馬人的王；如果用「我的國」來說,耶穌一定是羅馬人的王,因此後來耶穌對彼拉多說:

19:11耶穌回答他:「若不是從上頭賜給你的,你就毫無權柄辦我」

到底耶穌是不是君王呢?彼拉多到最後到仍然搞不清楚,在18:37他仍然問這個問題:

18:37彼拉多就對他說:「這樣,你是王嗎?」

彼拉多搞不清楚,因為彼拉多只有這個世界的帝王概念,而不能想像上帝國的君王是怎樣的.上帝國的君王和這個世界的帝王有甚麼不同呢?

18:37耶穌回答說:「你說我是王.我為此而生,也為此來到世間,特為給真理作見證.凡屬真理的人就聽我的話.」

彼拉多是不會明白,也不會接受,因為彼拉多的世界是權力鬥爭的世界,勝者為王,敗者為寇,拳頭就是真理,被拘捕帶到自己面前的這個人,如此軟弱無力,眾叛親離的耶穌,哪一點似王者呢?失敗者怎可能代表真理呢?如果世界不接受耶穌為王,不認同耶穌所見證的真理,也不用驚訝,在拳頭為真理的世界,人人都盼望強大的君王,怎可能認同耶穌呢?除非他們的心中另有所想.心中所想的是甚麼,就有可能認同耶穌這位君王呢?

四. 終末的王:以愛為管治的王／為子民自我犧牲的王

今天經課中的啟示錄1章給我們提示:

1:4但願從那昔在,今在,以後永在的上帝,和他寶座前的七靈,

1:5並那誠實作見證的,從死裡首先復活,為世上君王元首的耶穌基督,有恩惠,平安歸與你們!他愛我們,用自己的血使我們脫離罪惡,

1:6又使我們成為國民,作他父上帝的祭司.但願榮耀,權能歸給他,直到永永遠遠.阿們!

啟示錄1章5節指出耶穌這位君王的兩項特徵:

第一:「他愛我們」.撒母耳先知清楚指出,以掃羅為代表的這個世界的君王是以權力來管治,以武力來令人臣服,以奴役來對待人民；但上帝國的君王耶穌基督是用愛來管治,用愛來感召子民,用愛來對待子民.如果人們心中日思夜想的都是權力,武力和奴役,無論是王帝或是平民百姓,都不可能認同耶穌基督為君王；然而,那些心中渴望愛,盼望世間多點愛,也希望人人以愛相待的人,都會被耶穌基督這位君王吸引,向祂走近,跟隨祂的道路.

第二:「用自己的血使我們脫離罪惡」.相對於這個世界奴役人民來服侍自己的王帝,犧牲人民的生命來建立王國的王帝,耶穌基督反其道而行,為人民犧牲自己的生命,救人民脫離罪惡,為人民的福祉而不是為自己.在這個世界中信奉「人不為己天誅地滅」的人懷疑耶穌是可以理解的,他們的字典裡面沒有為人犧牲這個詞彙.然而,那些厭惡互相撕殺,受夠了你贏我輸這種零和遊戲,希望人間有雙贏,盼望人與人之間互相幫助的人,就可能被耶穌基督這位君王吸引,向祂走近,跟隨祂的道路.

五．總結:昔在,今在,以後永在的王

今天是「基督君王日」,提醒我們耶穌基督才是昔在,今在,永在的真正君王,他所展示的才是真正的君王風範.其實這節期的歷史不算悠久.它起源於亂世,上世紀初歐洲戰爭之後,多國經濟萎縮,人民生活困苦,國家民族之間因戰爭而結下的仇恨情緒瀰漫.在這樣的情勢下,民族主義高漲,人民期待及崇拜魅力領袖,視之為民族的救星.由於他們強調自己民族的優越性,轉而欺壓其他民族.有見於這類極端民族主義情緒日益高漲,教宗庇護十一世於1925年12月11日宣布《設立基督君王日》,盼望透過這個教會的節日,深願世人認清基督為世界的主,尊崇基督為王,停止干戈,使世人得享和平幸福.(香港聖樂促進會通訊2021年7月)

耶穌基督是昔在,今在,以後永在的王.昔日,以色列人遇到亂世:

12:12你們見亞捫人的王拿轄來攻擊你們,

於是要求立王:

就對我說:『我們定要一個王治理我們.』

也改變不到一個事實:

其實耶和華－你們的上帝是你們的王.

現在,無論人們拒絕或受耶穌基督為王,啟示錄揭示他們終有一天都要面對耶穌基督為王這個事實:

1:7看哪,他駕雲降臨!眾目要看見他,連刺他的人也要看見他；地上的萬族都要因他哀哭.這話是真實的.阿們!

1:8主上帝說:「我是阿拉法,我是俄梅戛,是昔在,今在,以後永在的全能者.」



\newpage

\section{將臨期的意義}
\label{sec:725}
\textbf{李均熊牧師 - 將臨期的意義}
\newline
\newline
\hyperref[sec:724]{< < < PREV SERMON < < <}
~
\hyperref[sec:toc]{[返目錄]}
~
\hyperref[sec:726]{> > > NEXT SERMON > > >}


講題:將臨期的意義 The Meaning of Advent

講員:李均熊牧師

場合:香港中文大學崇基學院禮拜堂主日崇拜

日期:2024年12月1日

起

新年快樂!

新年快樂?李牧師你是否弄錯了?聖誕還沒過,今日才是十二月一日,還有整整一個月才是新曆新年,為甚麼這個時候說新年快樂?

我沒有弄錯.因為今日是教會年曆的第一週,可以說是教會年的新年.若你參加崇基禮拜堂的聚會超過一年或以上,或本身參加奉行教會年曆的禮儀教會,你就知道每年的十一月尾或十二月頭,或正確點說是聖誕節前的第四個主日,就是教會年曆的開始,亦即教會的新年,也是我們進入將臨期的第一主日.

崇拜開始時我們已經點燃了將臨節花環(Advent Wreath)上四支蠟燭的第一支.將臨期的四個主日,會燃點四支蠟燭,直至聖誕主日點燃中間白色代表基督降生的蠟燭.每支都代表一種信徒等待慶祝基督降生的應有心情.第一支紫色的蠟燭代表「盼望」,第二支仍是紫色的蠟燭代表「等候」,第三支則是代表「喜樂」的粉紅色蠟燭,最後一支則回復紫色的蠟燭,代表「和平」.這四個主題,除了「喜樂」之燭用上粉紅色,其他三支蠟燭都是紫色的.紫色在教會年曆中主要出現在預備基督受死復活的四旬期／預苦期／大齋期(Lent),基調就是省思和悔罪.但紫色也出現在將臨期,故此當我們等候慶祝基督降生時,也提醒我們要省思悔罪.

承

但為甚麼一個預備聖誕普天同慶的日子,要有悔罪之心呢?

我們以為教會常常強調罪(guilt),對非信徒來說會否像是教會常常「情緒勒索」,要令到信徒常常感到內疚(guilty),以控制信徒循規蹈矩呢?若你因返教會以至內心常感到愧疚自責,我想你是誤解了悔罪(repent)的意思.不錯,悔罪是有一百八十度改變你的行為,從罪惡的生活中回轉的意思.但更多時候,悔罪是我們留意我們的生活行為,有沒有偏差,從上主的心意或作門徒的路徑中「偏行己路」,以至和上主的道越行越遠.所以悔罪並不需要每次都一百八十度轉變,若是如此你不就在原地團團轉?每次教會提出我們要省思悔罪,都是我們細察會否因為生活的種種世情以至走差了,迷失了,以至失落了領受恩典的機會.悔罪就好像一艘在暴風雨中航行的小船,要不斷留意是否偏離航道,並不斷修正航行的方向.

不錯,不斷留意自己生活的方向,就是今日經課提醒我們要「警醒」的意思.我們要「時時警醒,常常祈求」,免得因為種種外在的引誘(「貪食,醉酒」)或內在的不安(「今生的思慮累住你們的心」),因而迷失,偏行己路,遠離了迎接上主的方向,以至當上主再來時,未能夠「站立在人子面前」.

轉

將臨期的意義是甚麼呢?將臨期是有兩層意義的.最首要的當然是我們為著慶祝救主耶穌基督的降生而預備教會,預備自己.但將臨期的第一和第二週,也提醒我們基督要再次臨到人間,「那時,他們要看見人子有能力,有大榮耀駕雲降臨.一有這些事,你們就當挺身昂首,因為你們得贖的日子近了.」我們的救主耶穌基督要再一次臨到我們中間,驅走一切邪惡,更新一切.

但我們等了二千多年了,基督再臨好像遙遙無期.對於路加福音的讀者,那更是帶來極大的困惑.基督在離世前,不是預言了聖殿被毀,耶路撒冷的陷落嗎?當時跟隨耶穌的人,以為他的再來是即時的,是他復活之後很快會發生的事情,甚至認為其預兆就是聖殿和耶路撒冷的陷落.學者研究認為,在馬可福音成書時,當時讀者仍未經歷公元七十年羅馬軍隊進入耶路撒冷燒毀希律所修的聖殿,故此仍然認為基督的再來就是聖殿被毀之後.但到路加福音成書時,聖殿已經被羅馬軍隊無情的摧毀,但基督好像還未回來復興以色列,耶穌說人子會顯現的應許好像不會實現一般.你可以想像當時的門徒是多麼的失望,多麼的困惑!就是到了今天,我們是否仍會相信,基督真的會再來嗎?今日種種的黑暗邪惡,上主真的關心嗎?耶穌真的會回來驅走一切黑暗邪惡,帶來新的創造嗎?

經課中耶穌這一段,在馬太,馬可,路加福音都有記載,就是對於天上顯出異像,地上各國會有困苦,而人子帶著大榮耀駕雲降臨的應許,並且我們以無花果樹學比方,留心上帝國臨近的兆頭,知道天國必然成就.不過,若你比較今年路加年中路加福音的經課,和馬太或馬可福音相應的經文,路21:34-36一段則是路加獨有的,取代了馬太和馬可福音強調那日子那時辰沒有人知道的段落.

合

面對因為主再來遲遲未實然而帶來信徒的失望和困惑,路加額外提醒我們,要謹慎,時時警醒禱告,因為那日子會忽然臨到,是你不能預期的.若那時你生命的船偏離了航道,那時主再來時你卻離開了祂,那時要轉舵就太遲了.

所以讓將臨期的燭光照亮我們,使我們仍有盼望,仍願等候,不要因為上主時間和我們心目中的時間不同而失去信心.又讓將臨期的紫色,提醒我們,要謹慎警醒,要留心會否因為上主的應許和我們期望的不同,以至生活的方向變成追求今生的享樂和舒適,或因為種種生活的困難思慮煩憂,使我們失去依靠上主的信心,偏行己路,漸漸遠離上主的恩典和慈愛.

今日就讓我們以聖餐桌前的靜默,以悔罪的心,檢視和修正我們生命的方向,以迎接基督的再來.



\newpage

\section{在盼望中邁向更美好的新一天}
\label{sec:726}
\textbf{高國雄牧師 - 在盼望中邁向更美好的新一天}
\newline
\newline
\hyperref[sec:725]{< < < PREV SERMON < < <}
~
\hyperref[sec:toc]{[返目錄]}
~
\hyperref[sec:727]{> > > NEXT SERMON > > >}


路加福音 3:1-6
\newline
\begin{longtable}{cl}
\hline
\hline
章節 & 經文 (和合本修訂版)\\
\hline
3:1 & \begin{tabularx}{0.7\textwidth}{X} 凱撒提庇留在位第十五年，本丟‧彼拉多作猶太總督，希律作加利利分封的王，他兄弟腓力作以土利亞和特拉可尼地區分封的王，呂撒聶作亞比利尼分封的王， \end{tabularx} \\ \\ \relax
3:2 & \begin{tabularx}{0.7\textwidth}{X} 亞那和該亞法作大祭司。那時，撒迦利亞的兒子約翰在曠野裡，神的話臨到他。 \end{tabularx} \\ \\ \relax
3:3 & \begin{tabularx}{0.7\textwidth}{X} 他就走遍約旦河一帶地方，宣講悔改的洗禮，使罪得赦。 \end{tabularx} \\ \\ \relax
3:4 & \begin{tabularx}{0.7\textwidth}{X} 正如以賽亞先知書上所記的話：「在曠野有聲音呼喊著：預備主的道，修直他的路！ \end{tabularx} \\ \\ \relax
3:5 & \begin{tabularx}{0.7\textwidth}{X} 一切山窪都要填滿；大小山岡都要削平！彎彎曲曲的地方要改為筆直；高高低低的道路要改為平坦！ \end{tabularx} \\ \\ \relax
3:6 & \begin{tabularx}{0.7\textwidth}{X} 凡有血肉之軀的，都要看見神的救恩！」 \end{tabularx} \\ \\
[1ex]
\hline
\hline
\end{longtable}
講題:在盼望中邁向更美好的新一天  Better Days—There Is Hope

講員:高國雄牧師

經文:路加福音3:1-6

場合:香港中文大學崇基學院禮拜堂主日崇拜

日期:2024年12月8日

路加福音3:1-6

1凱撒提庇留在位第十五年,本丟‧彼拉多作猶太巡撫,希律作加利利分封的王,他兄弟腓力作以土利亞和特拉可尼地方分封的王,呂撒聶作亞比利尼分封的王,　2亞那和該亞法作大祭司.那時,撒迦利亞的兒子約翰在曠野裡,上帝的話臨到他.3他就來到約旦河一帶地方,宣講悔改的洗禮,使罪得赦.4正如先知以賽亞書上所記的話,說:在曠野有人聲喊著說:預備主的道,修直他的路!5一切山窪都要填滿；大小山岡都要削平!彎彎曲曲的地方要改為正直；高高低低的道路要改為平坦!6凡有血氣的,都要見上帝的救恩!

上主要來臨!

將臨期為紀念慶祝基督的第一次降臨；

期待迎接基督於末世第二次的光榮降臨.

我們今天帶著什麼的心情去等待呢?

近日大熱電影「破.地獄」中提到一些理念

「人一出世就在倒數.」

「有機會來到這個世界已經賺了,何必介懷甚麼時候下車,不如好好欣賞沿途的風景.」

「不止死人要超渡,生人也需要破地獄,生人都有好多地獄.」

當中不乏很有意義的思考及真理

今天我們燃點起的是「信心」之燭

對上主的信心能夠為我們帶來盼望

耶穌基督的先鋒施洗者約翰就是要在那邪惡,是非扭曲的世代向人呼喚:要回轉

5一切山窪都要填滿；大小山岡都要削平!彎彎曲曲的地方要改為正直；高高低低的道路要改為平坦!6凡有血氣的,都要見上帝的救恩!

這就是信心的宣告,這宣告就為人帶來盼望.

施洗約翰的下場就是因指出希律的惡行而被斬頭…

在人看來,這世界仍然在「惡」當中,甚至以為這「惡者」已經得勝.

但耶穌基督的誕降,為人類帶來了盼望

約翰福音1:14道成了肉身,住在我們中間,充充滿滿地有恩典有真理,我們也見過他的榮光,正是父獨一兒子的榮光.

上帝的道成了肉身,成為人的拯救

這拯救不單是我們死後可以上「天堂」,而是我們現在可以在「天國」,就是在上帝主權下生活,每天經歷主的同在,提醒,鼓勵,支持,引導,安慰,教訓,擁抱與接納…

在地上就得活出天國,體驗一點天國的美好,以致我們更能確信,更有盼望地過每一天的生活.

Better Days,不是客觀環境的一天比一天好,這世界因要走向滅亡,所以只會更壞…

但Better Days,係因為There is Hope,因為盼望,我們一天比天活得更好,活得更有價值,有意義,更精彩,直到見主的面!



\newpage

\section{就這樣習以為常地活了一生}
\label{sec:727}
\textbf{鄧瑞強博士 - 就這樣習以為常地活了一生}
\newline
\newline
\hyperref[sec:726]{< < < PREV SERMON < < <}
~
\hyperref[sec:toc]{[返目錄]}
~
\hyperref[sec:728]{> > > NEXT SERMON > > >}


阿摩司書 4:6-5:24
\newline
\begin{longtable}{cl}
\hline
\hline
章節 & 經文 (和合本修訂版)\\
\hline
4:6 & \begin{tabularx}{0.7\textwidth}{X} 「我使你們在每一座城裡牙齒乾淨，使你們各處的糧食缺乏，你們仍不歸向我。」這是耶和華說的。 \end{tabularx} \\ \\ \relax
4:7 & \begin{tabularx}{0.7\textwidth}{X} 「在收割的前三個月，我不降雨在你們那裡，我降雨在這城，不降雨在那城；這塊地有雨，那塊無雨的地就必枯乾。 \end{tabularx} \\ \\ \relax
4:8 & \begin{tabularx}{0.7\textwidth}{X} 兩三城的人擠到一個城去找水喝，卻喝不足，你們仍不歸向我。」這是耶和華說的。 \end{tabularx} \\ \\ \relax
4:9 & \begin{tabularx}{0.7\textwidth}{X} 「我以焚風和霉爛攻擊你們，你們許多的菜園、葡萄園、無花果樹、橄欖樹屢屢被剪蟲所吃，你們仍不歸向我。」這是耶和華說的。 \end{tabularx} \\ \\ \relax
4:10 & \begin{tabularx}{0.7\textwidth}{X} 「我降瘟疫在你們中間，如在埃及的樣子；用刀殺戮你們的年輕人和你們遭擄掠的馬匹，營中臭氣撲鼻，你們仍不歸向我。」這是耶和華說的。 \end{tabularx} \\ \\ \relax
4:11 & \begin{tabularx}{0.7\textwidth}{X} 「我傾覆你們，如同神從前傾覆所多瑪、蛾摩拉一樣；你們好像從火中搶救出來的一根柴，你們仍不歸向我。」這是耶和華說的。 \end{tabularx} \\ \\ \relax
4:12 & \begin{tabularx}{0.7\textwidth}{X} 「因此，以色列啊，我要如此對待你；因為我要這樣對待你，以色列啊，你當預備迎見你的神。」 \end{tabularx} \\ \\ \relax
4:13 & \begin{tabularx}{0.7\textwidth}{X} 看哪，那創山，造風，將其心意指示人，使晨光變幽暗，踩行在地之高處的，他的名是耶和華－萬軍之神。 \end{tabularx} \\ \\ \relax
5:1 & \begin{tabularx}{0.7\textwidth}{X} 以色列家啊，聽我為你們所作的哀歌： \end{tabularx} \\ \\ \relax
5:2 & \begin{tabularx}{0.7\textwidth}{X} 「以色列民跌倒，不得再起；躺在地上，無人扶起。」 \end{tabularx} \\ \\ \relax
5:3 & \begin{tabularx}{0.7\textwidth}{X} 主耶和華如此說：「以色列家的城派出一千，只剩一百；派出一百，只剩十個。」 \end{tabularx} \\ \\ \relax
5:4 & \begin{tabularx}{0.7\textwidth}{X} 耶和華向以色列家如此說：「你們要尋求我，就必存活。 \end{tabularx} \\ \\ \relax
5:5 & \begin{tabularx}{0.7\textwidth}{X} 不要往伯特利尋求，不要進入吉甲，也不要過到別是巴；因為吉甲必被擄走，伯特利必歸無有。」 \end{tabularx} \\ \\ \relax
5:6 & \begin{tabularx}{0.7\textwidth}{X} 要尋求耶和華，就必存活，免得他在約瑟家如火發出，焚燒伯特利，無人撲滅。 \end{tabularx} \\ \\ \relax
5:7 & \begin{tabularx}{0.7\textwidth}{X} 你們這使公平變為茵蔯，將公義丟棄於地的人哪！ \end{tabularx} \\ \\ \relax
5:8 & \begin{tabularx}{0.7\textwidth}{X} 那造昴星和參星，使死蔭變為晨光，使白晝變為黑夜，召喚海水、使其傾倒在地面上的，耶和華是他的名。 \end{tabularx} \\ \\ \relax
5:9 & \begin{tabularx}{0.7\textwidth}{X} 他快速摧毀強壯的人，毀滅就臨到堡壘。 \end{tabularx} \\ \\ \relax
5:10 & \begin{tabularx}{0.7\textwidth}{X} 你們怨恨那在城門口斷是非的，憎惡那說正直話的。 \end{tabularx} \\ \\ \relax
5:11 & \begin{tabularx}{0.7\textwidth}{X} 所以，因你們踐踏貧寒人，向他們勒索糧稅；你們雖建造石鑿的房屋，卻不得住在其內；雖栽植美好的葡萄園，卻不得喝其中所出的酒。 \end{tabularx} \\ \\ \relax
5:12 & \begin{tabularx}{0.7\textwidth}{X} 我知道你們的罪過何其多，你們的罪惡何其大；你們迫害義人，收受賄賂，在城門口屈枉貧窮人。 \end{tabularx} \\ \\ \relax
5:13 & \begin{tabularx}{0.7\textwidth}{X} 所以智慧人在這樣的時候必靜默不言，因為這是險惡的時候。 \end{tabularx} \\ \\ \relax
5:14 & \begin{tabularx}{0.7\textwidth}{X} 你們要尋求良善，不要尋求邪惡，就必存活。這樣，耶和華－萬軍之神必照你們所說的與你們同在。 \end{tabularx} \\ \\ \relax
5:15 & \begin{tabularx}{0.7\textwidth}{X} 要恨惡邪惡，喜愛良善，在城門口秉公行義；或者耶和華－萬軍之神會施恩給約瑟的餘民。 \end{tabularx} \\ \\ \relax
5:16 & \begin{tabularx}{0.7\textwidth}{X} 因此，主耶和華－萬軍之神如此說：「在一切的廣場上必有哀號的聲音；在各街市上必有人說：『哀哉！哀哉！』他們叫農夫來哭號，叫善唱哀歌的來舉哀； \end{tabularx} \\ \\ \relax
5:17 & \begin{tabularx}{0.7\textwidth}{X} 各葡萄園都有哀號的聲音，因為我必從你中間經過。」這是耶和華說的。 \end{tabularx} \\ \\ \relax
5:18 & \begin{tabularx}{0.7\textwidth}{X} 想望耶和華日子的人有禍了！為甚麼你們要耶和華的日子呢？那是黑暗沒有光明的日子， \end{tabularx} \\ \\ \relax
5:19 & \begin{tabularx}{0.7\textwidth}{X} 好像人躲避獅子卻遇見熊；進房屋以手靠牆，卻被蛇咬。 \end{tabularx} \\ \\ \relax
5:20 & \begin{tabularx}{0.7\textwidth}{X} 耶和華的日子豈不是黑暗沒有光明，幽暗毫無光輝嗎？ \end{tabularx} \\ \\ \relax
5:21 & \begin{tabularx}{0.7\textwidth}{X} 「我厭惡你們的節期，也不喜悅你們的嚴肅會。 \end{tabularx} \\ \\ \relax
5:22 & \begin{tabularx}{0.7\textwidth}{X} 你們雖然向我獻燔祭和素祭，我卻不悅納，也不看你們用肥畜獻的平安祭。 \end{tabularx} \\ \\ \relax
5:23 & \begin{tabularx}{0.7\textwidth}{X} 要使你們歌唱的聲音遠離我，因為我不聽你們琴瑟的樂曲。 \end{tabularx} \\ \\ \relax
5:24 & \begin{tabularx}{0.7\textwidth}{X} 惟願公平如大水滾滾，公義如江河滔滔。 \end{tabularx} \\ \\
[1ex]
\hline
\hline
\end{longtable}
講題:就這樣習以為常地活了一生  A Life like This

經文:阿摩司書4:6-5:24

講員:鄧瑞強博士

場合:香港中文大學崇基學院禮拜堂主日崇拜

日期:2024年12月15日

各位弟兄姊妹,早安.願神的恩典常與你們同在.

今日,我們一起看阿摩司先知書.基督教很多討論經濟正義和社會正義的文章,會引用阿摩司書.南美洲那些強調救助窮人的神學,會引用阿摩司書.美國的民權領袖馬丁路德金那著名的演說《我有一個夢》,也引用阿摩司書.阿摩司先知批判的,是社會上,經濟上的一切不義,他的遠象是建立一個不會欺凌弱小,尊重人性,資源共享的社會,而這社會之所以可能,建基於人對上帝的無限敬畏.

阿摩司書1:1說到阿摩司活在南國猶大烏西雅王,北國以色列耶羅波安二世的年代,大約是公元前744年左右,那時是猶大國和以色列國極鼎盛的時代.阿摩司先知有點怪,他本是猶大國的人,神卻吩咐他去以色列國傳講神的信息.以色列國有個祭司亞瑪謝,就對阿摩司很不客氣,說:「喂,你撈過界了,要搵食,唔該你行遠一點,請你返回你的鄉下去吧.」(摩7:12意譯)在人間,傳講神的說話,只要你夠認真,很多時人都會叫你「行遠一點」,不要在這裡講.或許,有一天,我都要「行遠一點」,返鄉下去.

為什麼上帝要調派一個南國的猶大先知去北國以色列傳道呢?那時,是以色列國的盛世,有錢人夜夜笙歌,不亦樂乎.或許,在這美好世界裡,已很難找到一個向同胞宣講神的審判的先知了.或許,沒有對的人可用,神就抽調阿摩司去,因為要傳遞的信息緊迫.以色列人完全沒有察覺,由那時的盛世,到被亞述完全殲滅,只是二十多年的光景.由極盛到無有,只需二十多個年頭.上帝要向他們,發出一個緊急的警報.前面是懸崖了,要緊急調頭.人心的崩壞,與社會的衰亡,息息相關.在這存亡之秋,先知呼喚人心悔改,敬神愛人,或許,這才能扭轉國家的命運.

我們看看今天的講道經文,先看:阿摩司書4:6—11

我讀經文的時候,請大家留意不同的災禍,和重複的句子.

摩4:6 我使你們在一切城中牙齒乾淨,在你們各處糧食缺乏,你們仍不歸向我.這是耶和華說的.

摩4:7 在收割的前三月,我使雨停止,不降在你們那裡；我降雨在這城,不降雨在那城；這塊地有雨,那塊地無雨；無雨的就枯乾了.

摩4:8 這樣,兩三城的人到一城去找水,卻喝不足；你們仍不歸向我.這是耶和華說的.

摩4:9 我以旱風,霉爛攻擊你們,你們園中許多菜蔬,葡萄樹,無花果樹,橄欖樹都被剪蟲所吃；你們仍不歸向我.這是耶和華說的.

摩4:10 我降瘟疫在你們中間,像在埃及一樣；用刀殺戮你們的少年人,使你們的馬匹被擄掠,營中屍首的臭氣撲鼻；你們仍不歸向我.這是耶和華說的.

摩4:11 我傾覆你們中間的城邑,如同我從前傾覆所多瑪,蛾摩拉一樣,使你們好像從火中抽出來的一根柴；你們仍不歸向我.這是耶和華說的.

在社會極繁盛之時,先知宣告五類即將來臨的災禍,這大致涵蓋人類生活的種種危機.

「各處糧食缺乏」:在農業社會,上天不賜福,大地不出糧食,有錢都無用.

「沒有足夠的水可喝」:天不降雨,生命即時受到威脅,也令農業失收.

「農業失收」:不夠雨水,大地炎熱,蟲害增生,導致沒有收成,這是經濟崩潰.

「瘟疫和戰爭」:有了「沙士」和「Covid」的經驗,我們知道瘟疫的可怕.再加上國與國的戰爭,真是死傷無數.

「城巿衰敗」:在天災人禍下,繁華的都市也只能苟延殘喘,失去前景.

以上講的,是經濟,民生,國際衝突,城巿衰敗的危機,在興盛的日子,阿摩司講這些東西,是叫人居安思危.先知向昔日的以色列國傳講神的信息,對今日的我們來說,或許,會覺得有點遙遠.但,真的是這樣嗎?近來,我們城巿負資產的數字越來越多,政府財赤也越來越大,也有其危機.講瘟疫,這個我們較有親身體會.至於戰爭,好似好遠,但又可以好近.我們周邊國家的軍備支出,近年都大幅飆升,國與國之間的磨擦越來越多,體諒越來越少.至於這個城市,前景如何,則見仁見智.

人是脆弱的,城市是脆弱的,國家是脆弱的,只要上帝不祝福,以色列國如何富強也無補於事.

我們應留意到,以上短短經文,五次講到「你們仍不歸向我」,這就點出問題的核心.就算是飢餓,旱乾,經濟崩壞,戰爭,城巿的逐漸衰敗,一切就在眼前,但人心依舊死不悔改,對生命的永恆意義,對人性裡應有的仁愛,對有限的人無法測透的上帝,以色列人毫不放在心上.這正是問題所在:他們對神毫不敬畏,他們仍不歸向神,他們離神好遠.

等一會,我們會讀另一段經文,那時,我們會看到以色列人具體的錯誤.現在,先知指出他們問題的核心:就算死到臨頭,以色列人仍是沒有半點敬畏上帝的心.人若不歸向神,沒有半點敬畏上帝的心,則人終歸就不會考慮什麼價值與道德了.沒有永恆,沒有上帝,沒有上帝對人的終極期望與審判,則還有什麼力量可阻止邪惡的人行邪惡的事呢?所以,有人說,若沒有上帝的話,則人可以任意而行了.事實上,以色列人的罪,就是有權者對無權者的任意欺凌,這使得神不得不作出警告了.

剛才我們讀的經文,神五次說「你們仍不歸向我」,接著,是五次「這是耶和華說的」.我們看到,神多麼認真,多麼鄭重地警告墮落的人,多麼渴望人的回轉.真的到懸崖了,還不剎車!但人怎願意聽神的要求呢?

我們看看以色列人的生活,看看你們是否同意神要警告他們.

阿摩司書5:7,10—15:我們逐節經文來看

摩5:7 你們這使公平變為茵蔯〔苦的〕,將公義丟棄於地的.........

他們藐視公平和公義.他們漠視公平,使弱者飽嚐苦澀.他們踐踏公義,使之如地上的糞土.他們做了什麼?

摩5:10 你們怨恨那在城門口責備人的〔指出人錯誤的〕,憎惡那說正直話的〔實話實說的〕.

城門口,是當時的民間法庭,是做判決的地方.他們怨恨那些主持公道,指責他們不公不正的人；他們憎惡那些講事實的人.換句話說,他們希望假話當真話,從中得到有利自己的判決.

摩5:11 你們踐踏貧民,向他們勒索麥子；你們用鑿過的石頭建造房屋,卻不得住在其內；栽種美好的葡萄園,卻不得喝所出的酒.

有錢人用種種方法,欺壓窮人,令窮人食無可食,居無可居,飲無可飲.無論窮人如何辛勞去賺取生活,都依然活得淒涼.

摩5:12 我知道你們的罪過何等多,你們的罪惡何等大.你們苦待義人,收受賄賂,在城門口屈枉窮乏人.

在城門口,握有裁判權的,有地位的人,收取賄賂,屈枉無錢的正直人.應講法律的地方,不講法律,則有權有勢的人便能胡作非為,窮人只能任他們魚肉,有冤無路訴.

摩5:13 所以通達人見這樣的時勢必靜默不言,因為時勢真惡.

醒目的人都不再出聲了,沉默是金,因為時勢真惡.我信耶穌信了幾十年,以前很少聽這經文,近這些年,聽這經文真是聽多了.

摩5:14 你們要求善,不要求惡,就必存活.這樣,耶和華─萬軍之神必照你們所說的與你們同在.

摩5:15 要惡惡好善,在城門口秉公行義；或者耶和華─萬軍之神向約瑟的餘民施恩.

又再提及「城門口」,又再提及「公義」,看來,先知很著緊法律的公正執行.在「時勢真惡」的日子,先知呼籲人,要求善,不要求惡,要惡惡好善,這才有生機,社會才有出路,國家才有希望.在是非不分,真假難辨的時代,要堅持美善,是很難的.但,這正是神的子民的責任.

近幾年,有本很暢銷的基督教書籍,書名叫《時勢好惡,做基督徒好難》.作者梁柏堅在書中第一章如此寫:「當阿摩司說,『時勢真惡.你們要求善,不要求惡,就必存活』,我們看見一個天秤擺在面前,天秤的兩端,分別有來自天國的善和地上的惡,陳明生死與禍福.......揀選生命......我們從此有了不一樣的視點,相信世間有另一套真實......當人能擺脫權力的枷鎖,不隨波逐流,不以世界行惡的方式積聚財富,加官晉爵,而能跟隨天國的良善,俯就卑微的人,行事公正無私,為無權無勢的人抵擋貪婪者的獠牙,拭乾傷心者的眼淚,天國就離我們不遠.」(頁26)

我們可能覺得好奇怪,以色列人都是敬拜耶和華的人,他們的信仰為何沒有更正他們的生活?理由是:他們的信仰錯了.他們相信一個不斷賜福他們,讓他們豐足的神明,而不相信一個要求他們行為正直的上帝.他們開心的接受神給他們的權力,卻不接受神要求的道德理念.他們祈求神讓他們做買賣時掠奪最大的盈利,卻不尊重神要求的經濟正義.他們相信一個能被他們操控的神,卻不接受神是他們的主.他們願意奉行各種禮儀,但對神沒有半點敬畏之情.他們有宗教,卻沒有上帝.注目自己,而沒有仰望蒼天.有表面的敬虔,卻沒有與信仰相稱的道德人格.他們如此敬拜上帝,倒頭來,便越執著於自己的利益和享樂,越遠離上帝,越遺忘上帝.他們不想改變,他們寧願習以為常地活其一生.

最後,先知叫他們正視自己的信仰.

摩5:21 我厭惡你們的節期〔三大節期:逾越節,五旬節,住棚節；我們今日的三大節典,可能是聖誕節,受難節,復活節〕,也不喜悅你們的嚴肅會〔節期裡的宗教聚會〕.

摩5:22 你們雖然向我獻燔祭和素祭,我卻不悅納,也不顧你們用肥畜獻的平安祭；

摩5:23 要使你們歌唱的聲音遠離我,因為我不聽你們彈琴的響聲.

摩5:24 惟願公平如大水滾滾,使公義如江河滔滔.

他們守的節期,參與的嚴肅會,獻的燔祭,素祭,平安祭,歌唱的聲音,彈琴的響聲,七項核心的宗教活動,神全部不喜悅.當宗教活動失去了敬畏上帝的心,只淪為開開心心的社交活動,上帝怎會悅納呢?上帝要的,是有道德內容的信仰,這信仰流露在正直的生活裡,這生活呈現出對上帝的敬畏.上帝要的,是這種信仰.

上帝希望:公平如大水滾滾,公義如江河滔滔.在曠野地方,最需要的,是長流不息的活水,唯有這活水,才能令乾旱的大地孕育出生命.在人心枯燥,社會沒有生氣的時刻,我們需要這長流不息的活水.上帝渴望:公平如大水滾滾,公義如江河滔滔.但願我們,在艱難時刻,努力活出美善,活出公平,活出公義.但願弟兄姊妹們,能不同流合污,能成為社會的清泉,孕育生命.

願你聆聽上主的要求,願你回應上主的呼喚,過一美善而正直的人生.



\newpage

\section{為這個世界帶來一個新生命}
\label{sec:728}
\textbf{林豪恩先生 - 為這個世界帶來一個新生命}
\newline
\newline
\hyperref[sec:727]{< < < PREV SERMON < < <}
~
\hyperref[sec:toc]{[返目錄]}
~
\hyperref[sec:729]{> > > NEXT SERMON > > >}


講題:為這個世界帶來一個新生命  Bringing a New Life to This World

講員:林豪恩先生

場合:香港中文大學崇基學院禮拜堂主日崇拜

日期:2024年12月22日

一. 引言:新生命帶來的突破

今天有十位成人及兩位兒童將會接受洗禮.有人把洗禮視為自己的第二個生日,代表一個新生命的開始.新生命往往帶給人盼望,希望由壞變好,希望狀況有突破.還記得武俠小說倚天屠記中,張無忌的出生,如何改變金毛獅王謝遜嗎?

話說金毛獅王謝遜,張翠山和殷素素流落到冰火島.殷素素十月懷胎,嬰兒快要出生的時候,張翠山和殷素素不是感到興奮,反而是憂心忡忡.因為他們發現金毛獅王謝遜的臉色越來越兇狠,時時拿著屠龍刀狂揮狂舞,咒天罵人,擔心他發狂起來傷害孩子.有一天,謝遜突然大吼:「張三豐不是東西,他的弟子張翠山更加不是東西,讓我捏死他的老婆再說!」謝遜奔進殷素素臨盆的山洞,張翠山力阻謝遜傷害快將分娩的殷素素.謝遜揮舞屠龍刀,刀花如雪,寒光閃閃,張翠山心中一酸,想起今日和殷素素同時畢命,竟不能見一眼那未出世的孩兒.突然間「哇」的一聲,內洞中傳出一響嬰兒的哭聲.謝遜大吃一驚,站著不動.張翠山和殷素素凝視著這初生的嬰兒,知道只要謝遜一刀下來,夫妻倆連著嬰兒便同時送命.突然之間,謝遜良知激發,狂性消失,頭腦清醒過來,問:「是男孩還是女孩?」張翠山道:「是個男孩.」謝遜道:「很好.剪了臍帶沒有?」張翠山道:「要剪臍帶嗎?啊,是的,是的,我倒忘了.」謝遜倒轉長劍,將劍柄遞了過去.張翠山接過長劍,割斷了嬰兒的臍帶,這時方始想起,謝遜已然迫近身邊,可是他居然並不動手,心中奇怪,回頭望了他一眼,只見謝遜臉上充滿關切之情,竟似要插手相助一般.謝遜又道:「你有沒燒了熱水,給嬰兒洗一個澡?」謝遜主動奔出去燒水,過了一陣,端了一盆熱水進來,張翠山便替嬰兒洗澡.謝遜聽得嬰兒哭聲洪亮,問道:「孩兒像媽媽呢還是像爸爸?」張翠山微笑道:「還是像媽媽多些,不大肥,是張瓜子臉.」謝遜歎了口氣,低聲道:「但願他長大之後,多福多壽,少受苦難.」

一個新生命的出現,原來能夠為另一個人帶來改變,並且為艱難的狀況帶來突破.

二．彌迦書:等待那生產的婦人生下子來

今天的舊約經課也有一個等待新生命的故事:

5:2伯利恆的以法他啊,你在猶大諸城中為小,將來必有一位從你那裡出來,在以色列中為我作掌權的；他的根源從亙古,從太初就有.

5:3耶和華必將以色列人交付敵人,直等那生產的婦人生下子來.

公元前七二二年,北國以色列被亞述滅國,南國猶大的先知彌迦看見猶大社會腐敗不公,相信上主的懲罰終必到來,猶大能逃北國以色列的覆轍.在審判之中,彌迦先知宣告盼望.盼望是甚麼呢?5:3「直等那生產的婦人生下子來」——盼望就是一個新生命的出現.倚天屠龍記中的新生命把惡人的戾氣化為祥和,把殺機轉為生機；彌迦先知預告的新生命又會帶來甚麼改變呢?

第一,5:3「其餘的弟兄必歸到以色列人那裡」——從流散裂國變為回歸重聚.

第二,5:4「他必起來,倚靠耶和華的大能,並耶和華－他上帝之名的威嚴,牧養他的羊群.他們要安然居住」——從群羊無首流離失所變為有人照顧安居樂業.

第三,5:5「這位必作我們的平安」——從動盪不安轉變為和平共存.

「生產的婦人」是誰呢?「生下子來」的孩子又是誰呢?歷史上曾經出現過很多似是而非的人物,引起人們無限盼望,又令人們無盡失望.

三．路加福音:萬眾期望集於一身的新生命

今天的福音經課出現了兩位懷胎的婦人,會不會是他們其中一位呢?那孩子會不會就是其中一個呢?

1:39那時候,馬利亞起身,急忙往山地裡去,來到猶大的一座城；

1:40進了撒迦利亞的家,問伊利莎白安.

1:41伊利莎白一聽馬利亞問安,所懷的胎就在腹裡跳動.

1:42伊利莎白且被聖靈充滿,高聲喊著說:「你在婦女中是有福的!你所懷的胎也是有福的!

1:43我主的母到我這裡來,這是從哪裡得的呢?

1:44因為你問安的聲音一入我耳,我腹裡的胎就歡喜跳動.

1:45這相信的女子是有福的!因為主對她所說的話都要應驗.」

從伊利莎白的說話可以知道,預言中的婦人不是他,而是來探望她的馬利亞；那個盼望中的新生命是馬利亞所懷的胎.何以是馬利亞所懷的胎呢?伊利莎白說「主對她所說的話都要應驗」.主對馬利亞說了甚麼話呢?

1:30「馬利亞,不要怕,你在神面前已經蒙恩了.

1:31你要懷孕生子,要給他起名叫耶穌.

1:32他將要為大,稱為至高者的兒子；

主神要把他祖先大衛的王位給他.

1:33他要作雅各家的王,直到永遠；

他的國沒有窮盡.」

如果是此時此地的父母聽到這番預告自己的孩子的成就,會有甚麼反應呢?相信大多數都會說:「不用客氣了,只要孩子健康成長,將來生活愉快就可以了.」一個被如此預告的新生命出現,很可能令各路人馬產生不同的聯想,引起各個群體不同的期望或擔憂.看著被各種猜測懷疑的眼光投射在孩子身上,被各種人物的眾多期待撕裂的孩子,最心痛的是誰呢?在聖殿裡,西面曾經對孩子的母親馬利亞說:

2:34「這孩子被立,是要叫以色列中許多人跌倒,許多人興起；又要成為毀謗的對象,

2:35叫許多人心裡的意念顯露出來；你自己的心也要被劍刺透.」

每一個人都對這個新生命有期望,然而,有誰明白他呢?有誰願意認識這個新生命來到這個世界的使命及意義呢?書信經課希伯來書指出指引我們從這個角度來理解這個新生命:

四．希伯來書:帶來新生命的新生命

10:5所以基督到世上來的時候,就說:上帝啊,祭物和禮物是你不願意的；你曾給我預備了身體.

10:6燔祭和贖罪祭是你不喜歡的.

10:7那時我說:上帝啊,我來了,為要照你的旨意行；我的事在經卷上已經記載了.

10:8以上說:「祭物和禮物,燔祭和贖罪祭,是你不願意的,也是你不喜歡的(這都是按著律法獻的)」；

10:9後又說:「我來了為要照你的旨意行」；可見他是除去在先的,為要立定在後的.

10:10我們憑這旨意,靠耶穌基督,只一次獻上他的身體,就得以成聖.

或許當時的人期待這個新生命「使以色列再強大」,然而,驍勇善戰的君王不是沒有,民族英雄也不是沒有,唯獨沒有「我來了為要照你的旨意行」的人.這個新生命是一個真正活出上主心意的生命,也由於活出上主心意而能夠為其他人帶來新生命的生命.

五．洗禮帶來的新生命

有人說洗禮是新生命的誕生,使徒保羅在羅馬書6:4的確如此說:

「所以,我們藉著洗禮歸入死,和他一同埋葬,原是叫我們一舉一動有新生的樣式,像基督藉著父的榮耀從死裡復活一樣.」

無論你是今天受洗,或是已經受洗,接受洗禮的人將會開展新生命,這個新生命的確是賜給你的,然而,也是賜給這個世界的.你會被這個新生命改變,這個世界也會因為多了一個新生命而被改變.或者有人說,我已經一把年紀了,還能談甚麼新生命嗎?願使徒保羅在哥林多後書的話成為我們的勉勵:

4:16外體雖然毀壞,內心卻一天新似一天.



\newpage

\section{渴求的平安在地上屬於人的}
\label{sec:729}
\textbf{高國雄牧師 - 渴求的平安在地上屬於人的!}
\newline
\newline
\hyperref[sec:728]{< < < PREV SERMON < < <}
~
\hyperref[sec:toc]{[返目錄]}
~
\hyperref[sec:730]{> > > NEXT SERMON > > >}


講題:渴求的平安在地上屬於人的!On Earth Peace Among Those Whom He Favors!

講員:高國雄牧師

場合:香港中文大學崇基學院禮拜堂平安夜燭光崇拜

日期:2024年12月24日

(第一部份)

可以差不多說全人類都慶祝聖誕節

今天在香港的我們,又正為什麼事情慶祝呢?

人尋求生命,生活所需

每天努力「搵食」

聖誕願望:

證券公司期盼股票上漲

2023澳洲一名13歲少女寫下聖誕願望清單,列出一共43件物品,總值約3萬港元,而「金錢」則寫了兩次.

你又為什麼而在這節期慶祝呢?

你的聖誕願望呢?

(第二部份)

馬太福音2

1在希律作王的時候,耶穌生在猶太的伯利恆.有幾個博學之士從東方來到耶路撒冷,說: 2「那生下來作猶太人之王的在哪裡?我們在東方看見他的星,特來拜他.」 3希律王聽見了,就心裡不安；耶路撒冷全城的人也都不安. 4他就召集了祭司長和民間的文士,問他們:「基督該生在哪裡?」 5他們說:「在猶太的伯利恆.因為有先知記著:6『猶大地的伯利恆啊,你在猶大諸城中並不是最小的；因為將來有一位統治者要從你那裡出來,牧養我以色列民.』」7於是,希律暗地裡召了博學之士來,查問那星是甚麼時候出現的, 8就派他們往伯利恆去,說:「你們去仔細尋訪那小孩子,找到了就來報信,我也好去拜他.」 9他們聽了王的話就去了.

從東方來的博士(星象家)到來不是敬拜希律王,而是要拜那位更大的「新生王」

希律雖然存心不良,但卻指引了博士們對的方向

上帝在我們想像不到的地方,也要賜下恩典!

(第三部份)

路加福音2:8-14(原文新約)

8在那地方有些牧羊人逗留在野外,在晚上看守他們的羊群.

9有主的天使站在他們旁邊,主的榮光照亮他們四周,他們感到大大的懼怕和恐懼.

10天使對他們說:「不要懼怕,看哪!我要告訴你們好消息,使整個民族獲得大喜悅.

11因為今天在大衛的城裡為你們生了救主,就是主基督.

12這是給你們的記號:你們將會找到一個嬰兒,被包在布裡,並躺臥在飼料糟裡.」

13突然,一群天上的軍隊到來與天使讚美上帝說:

14「榮耀在高天上屬於神的,渴求的平安在地上屬於人的.」

天軍與天使的讚美,宣告這世界的實在

耶穌所賜的平安

約翰福音14:27

我留下平安給你們,我把我的平安賜給你們.我所賜給你們的,不像世人所賜的.你們心裡不要憂愁,也不要膽怯.

人類在地上所渴求的,真正心底裡的需要,就是要在地上得到平安

這渴求在地上能完全被滿足,就是因為完全反映了上帝在天上的榮耀

是一對,同時出現的



\newpage

\section{東方博士的結局正邪交戰的序幕}
\label{sec:730}
\textbf{郭志榮牧師 - 「東方博士」的結局——正邪交戰的序幕}
\newline
\newline
\hyperref[sec:729]{< < < PREV SERMON < < <}
~
\hyperref[sec:toc]{[返目錄]}
~
\hyperref[sec:731]{> > > NEXT SERMON > > >}


講題:「東方博士」的結局——正邪交戰的序幕 The Outcome Brought by The Magi from The East –  A Prelude to The War between Good and Evil

講員:郭志榮牧師

場合:香港中文大學崇基學院禮拜堂主日崇拜

日期:2024年12月29日

一．前言

弟兄姊妹,主內平安.雖然今年的聖誕節已經過了,但今天仍在聖誕節期之中,所以還可以祝大家聖誕快樂!

若要用一個形容詞來形容聖誕節,大家會選擇哪一個形容詞?在大家心中,聖誕故事的氣氛,總是離不開「平安」,「喜樂」,「歡欣」,「普世歡騰」,「充滿盼望」等等.

然而,這並非聖誕故事的全部.就以《馬太福音》為例,東方博士的拜訪並非在他們獻上黃金,乳香,沒藥等禮物之後就大團圓結局,而是仍有「下文」.這「下文」正好突顯出聖誕故事有不平安,不喜樂,不歡欣的一面.

今天跟大家分享的經文,就是〈東方博士〉故事的下集.經文來自《馬太福音》二章16-18節,今天所用的是《和合本》,讓我們一同誦讀.

二．經文:太二16-18(和合本)

16  希律見自己被博士愚弄,就大大發怒,差人將伯利恆城裡並四境所有的男孩,照著他向博士仔細查問的時候,凡兩歲以裡的,都殺盡了.

17  這就應了先知耶利米的話,說:

18  在拉瑪聽見號咷大哭的聲音,是拉結哭她兒女,不肯受安慰,因為他們都不在了.

三．釋經

A. 太二16上:「希律見自己被博士愚弄,就大大發怒,......」

經文一開始提到:「希律見自己被博士愚弄,就大大發怒.」到底發生了甚麼事情,令希律如此憤怒?而那些「博士」又是甚麼人呢?

首先跟大家交代一下事件發展的背景.《馬太福音》第二章一開始就提到,耶穌在伯利恆誕生後,有幾位來自東方的博士到達耶路撒冷(1節).

「東方」指甚麼地方呢?「東方」所指的可能是波斯,巴比倫一帶,也可能是阿拉伯或埃及等地——位於南方的阿拉伯或埃及的旅客會從東面的路進入猶大地,所以也是「來自東方」.[1]

那麼「博士」又指甚麼人呢?「博士」一詞是原本用來指那些瑪代波斯的祭司,後來卻泛指那些在異教社會中擁有某些特別知識或能力的人,如那些能解讀天文現象或夢境的「智者」,或那些占卜算命,施行法術的術士——後者的用法(即「術士」)通常帶點負面的含義,暗示他們只是招搖撞騙的騙子.[2]

幾位東方博士聲稱他們觀測到不尋常的天文現象,看到一顆很特別的星,特別到一個程度令他們認定這顆星是「猶太人之王」誕生的標記.他們於是長途拔涉來到耶路撒冷,希望朝拜這位非凡的王者(2節).

大家可能會問,幾位東方博士怎樣從「星的出現」而推斷出「猶太人之王誕生」的結論呢?波斯的占星術有這樣記載嗎?機會極微!幾位東方博士對「猶太人之王誕生」的認知,更可能是從散居在他們當中的猶太人那裡獲得的.[3] 《死海古卷》就收錄了一些猶太人的文獻,當中把《民數記》廿四章17節中「有星出於雅各,有杖要興於以色列」解讀為彌賽亞降臨的預言,說明早已有類似的觀點在某些猶太社群當中流傳.[4] 幾位從事解讀天文現象的東方博士,很可能就是聽過類似的說法.

幾位東方博士在耶路撒冷到處打聽「猶太人之王」的事,最終傳到希律的耳中,觸動了他的神經,因而使他「心裡不安」.「心裡不安」本來是形容把水「搖動」或「攪動」的動作,在這裡引伸出來,就是指希律被傳聞所「攪擾」,使他感到不安,憂愁,煩燥,驚恐,內心難以平靜.[5]

希律徵詢了祭司長和文士的意見,得知《彌迦書》預告了彌賽亞將在伯利恆出生,然後就秘密召見幾位東方博士.希律一方面從東方博士那裡「套料」,找出這位「猶太人之王」大概甚麼時候出生,另一方面派幾位博士到伯利恆,叫博士找到「猶太人之王」後要馬上回來通知他.東方博士似乎沒有懷疑希律的動機,答應了希律,卻因為在夢中得到指示而在朝拜聖嬰後秘密繞道離開.

可以想像得到,本來已「心裡不安」的希律,等了一段時間卻杳無音訊,必然是愈等愈著急,愈等愈煩燥.最後見到幾位東方博士沒有回來,覺得被他們愚弄,因此就「大大發怒」.「發怒」這個字本身已經是「狂怒」的意思,[6] 還要加上「大大」這個副詞,表示希律已經「狂怒」到情緒失控的狀態.

這就是事件發展的背景,也就是希律「大大發怒」的原因.

B. 太二16下:「......差人將伯利恆城裡並四境所有的男孩,照著他向博士仔細查問的時候,凡兩歲以裡的,都殺盡了.」

經文跟著說:「(希律)差人將伯利恆城裡並四境所有的男孩,照著他向博士仔細查問的時候,凡兩歲以裡的,都殺盡了.」

情緒失控,變得瘋狂的希律於是就派人前往伯利恆,把伯利恆和周邊地區中所有兩歲或以下的男孩都屠殺了.

大家有沒有想過:為何是「兩歲以下」?

東方博士由觀測到不尋常的天文現象到決定起行,再長途跋涉來到耶路撒冷,再到希律見不到東方博士回來而動殺機,前後應該已經過了一段時間,最少也有數個月,或許接近一年.[7] 既然希律已經「向博士仔細查問」了「那星是甚麼時候出現」,所以希律大可以把屠殺的範圍縮小,只幹掉那些有「嫌疑」的男嬰便可.[8] 例如,剛出世沒多久的嬰孩,「星」出現的時候他們還未出生,照計不會是那位「猶太人之王」；「星」出現之前就已經出生的,也不符合要求.

但瘋狂的希律卻「寧枉勿縱」,把要殺害嬰孩的範圍擴大至所有「兩歲以下」的「男孩」,而非只是「嬰兒」,地區也擴闊至伯利恆的「四境」.[9]

希律屠殺男孩的決定不僅是兇殘和血腥,也是非常弔詭,很矛盾.大家可以想想:到底在希律心目中,是否相信這些不明來歷的東方博士所講,彌賽亞已經降生呢?希律是否真正相信彌賽亞的應許?

假若希律自己相信彌賽亞的應許,並相信彌賽亞已經降生,而既然彌賽亞是神所差遣來作王,希律憑甚麼可以敵擋神的受膏者?他不怕惹神的怒氣而自招滅亡嗎?

事實上,希律一開始時並沒有派人跟隨或跟蹤著幾位博士,可見希律雖然「心裡不安」,但他對幾位東方博士的說話最多只是半信半疑.耶路撒冷的祭司長,文士以至「合城的人」對傳言更是「冷處理」,他們當中似乎沒有人打算到伯利恆看個究竟.他們最擔心的,主要是希律會否有過激的反應.

另一方面,假若希律並不相信彌賽亞的應許,不相信那些東方博士的說話,他何必要大費周章,派人去伯利恆殺害那些無辜的男孩呢?對於不相信有彌賽亞應許的人來說,幾位博士大概只是招搖撞騙的術士,對他們所說的話需要那麼認真和緊張嗎?所以,希律屠殺男孩的決定看來並不合理.

按常理來說,希律反而應該要問,幾位招搖撞騙的術士散佈謠聞,是否別有所圖?是否要煽動猶太人反抗希律的政權?假若如此,要對付的應該是這些圖謀不軌的術士,騙子!但希律也沒有這樣做.

大家最後一定會問:那麼,希律為何動了殺機?

答案實際上可以在經文中找到,關鍵就是「被愚弄」這個字.「愚弄」這個字只有在《福音書》中出現,除了兩次之外,全都是用來形式耶穌被「戲弄」的情境.[10] 「愚弄」也有「嘲弄」,「恥笑」的成份.[11] 《馬太福音》的讀者當然知道,那些「東方博士」並沒有「愚弄」希律,但希律卻是這樣理解.經文中的「見」,即「理解」之意.[12]

「心裡不安」的希律,心情本來已經難以平伏,現在又覺得被這幾個人「愚弄」了.在希律心中,他甚至可能會認定天下人都正在背後「恥笑」他.如此大的羞辱,希律怎會受得了?於是他老羞成怒,情緒失控,下令以屠殺男孩來洩憤,也藉此威嚇那些背後「恥笑」他的人.

C. 太二17-18:「這就應了先知耶利米的話,說:在拉瑪聽見號咷大哭的聲音,是拉結哭她兒女,不肯受安慰,因為他們都不在了.」

經文跟著就是馬太對這次屠殺男孩事件的註釋和評論:「這就應了先知耶利米的話,說:在拉瑪聽見號咷大哭的聲音,是拉結哭她兒女,不肯受安慰,因為他們都不在了.」

大家可能會問:這是否意味,馬太認為這次屠殺事件是神所命定,所計劃的呢?馬太當然不是這個意思!

《和合本》把17節的經文譯作「應了先知耶利米的話」,同一個希臘文動詞,為甚麼是譯作「應了」而不是好像《馬太福音》一章22節和二章15節那樣譯「應驗」呢?背後其實是大有文章的.

在一章22節和二章15節「應驗」的那個動詞前有一個連接詞「ινα」.這連接詞就好像英文的「in order that」那樣,表示連接詞後面的內容是前者的「目的」,[13] 所以,《和合本》把那裡的動詞譯作「應驗」,說明上文所發生的事件是依循著神的計劃而行,目的是要應驗《聖經》的話語,實現神在《聖經》中的應許.

這裡17節則用了另一個連接詞「τότε」,可解作「那時」或「於是」(in that case)之意,是沒有包含任何「目的」關係的含義.[14] 馬太特意要提醒讀者,無辜男孩的屠殺絕非神的目的,也不是神所命定或計劃的事情,而是出於希律邪惡的意念,所以希律要為到他自己的惡行負上全部的責任.[15] 同樣的情況也出現在《馬太福音》廿七章9節,馬太指猶大以三十塊錢出賣耶穌,也是「應了先知耶利米的話」,原文也是用了「τότε」而不是「ινα」.[16]

18節引用了《耶利米書》卅一章15節的經文,這節經文本來要表達甚麼呢?

耶利米先知要對耶利米時代的人說:連已死去並埋葬在拉瑪附近的拉結——她是雅各最愛的妻子,所以也是以色列的民族之母——也要為到當時被殺和被擄的以色列民發出哀哭的聲音.拉結「不肯受安慰」,表達出那是持續和深沉的傷痛,以致別人安慰的說話都聽不進去,終日以眼淚洗面,哀哭無法止息.[17]

拉瑪正好在伯利恆的附近,在伯利恆周邊發生的悲劇震動了那裡每一個人的心靈,連葬在那裡的拉結也要再一次哀哭.馬太在這裡引用《耶利米書》這一節經文,正是向那些無辜被害的男孩表達同樣的哀痛.[18] 換句話說,馬太要指出,為這些男孩哀哭的不只是他們的父母或他們的家人,以色列的先祖以至整個民族也要為他們哭泣!

四．引伸應用

A. 預示了敵擋天國的邪惡勢力

聖誕節記念主耶穌的降生,為幽暗的世界帶來光明,帶來盼望.

但馬太藉著〈東方博士〉這個聖誕故事告訴《馬太福音》的讀者,即使主耶穌已經降生,罪的權勢仍沒有想過放棄敵擋神.罪的權勢試圖藉著希律瘋狂的行動,阻止神藉著成為肉身的聖道在世上施行救恩,妨礙神國度的開展和實現.

這權勢要竭盡所能去破壞神的計劃,要繼續迷惑人心,扭曲人生命應有的形像樣式,要繼續使人成為它奴役人的差役和工具,同時也要使人成為被奴役,被拆騰的受害者,生命變得肢離破碎,最後在沉淪中步向滅亡.

〈東方博士〉這個「不一樣」的聖誕故事正要說明,那「正邪屬靈爭戰」的序幕,正隨著耶穌的降生而正式揭開.這是一場「屬靈爭戰」,因為它的戰場就在人的心靈之中,是天國與罪的權勢兩套不同價值的正面交鋒.

這正預示了,跟隨和事奉主耶穌基督的信徒,自然也要面對罪權勢的攻擊,因為「僕人不能大於主人」(約十五20).所以,信徒們也要參與在天國與罪權勢的屬靈爭戰之中.這個世界充滿了歪理和謊言,要迷惑人心,當中所發生的事情,即使是信徒也感到困惑,迷惘和無所適從,找不到出路.所以,信徒務要與基督同一陣線,努力活出並傳揚基督的福音真理,並要在做人處事上堅守天國的屬靈價值.

B. 人的不安與不信所帶來的惡果

上面提到,希律所作的,無論從甚麼角度看,也是極不合理的.為何希律會這樣?

希律並非天生就是一個瘋子.假若從希律的政績來評價他的一生,可以說希律曾經也是一個精明的政治家.他取得了羅馬帝國的信任而獲冊封為「猶大地的王」(King of Judea),統治的範圍廣闊,包括了撒馬利亞和敘利亞的土地.[19]

希律的父親是以土買人(即以東人的後代),他的母親是南面的阿拉伯人.但他卻歸信了猶太教,至少表面上遵行猶太的律法.希律也娶了猶太人馬加比王朝的公主米利暗(Mariamme)為妻,以確立自己統治的立法性.他擴建了聖殿,也為猶大地帶來穩定發展和繁榮的時代.[20]

希律最愛米利暗這位妻子,他卻在妹妹撒羅米(Salome)不斷的中傷和唆擺下,親自處決了自己的愛妻,把本來已經到手的幸福親手毀掉.宮庭的鬥爭沒有因此而止息.希律連自己的至愛也親手殺害,他自此以後就變本加厲,變得更加猜疑和更加可怕.希律完全被自己不安的情緒所掌控了,使他最後十年的統治陷在一片風腥血雨和混亂之中.希律愈是不安,愈是要設法以強權和殘暴的手法來加強操控,但這反而使情況變得更加糟糕.在這混亂的十年中,他竟然先後殺死了自己三個兒子.[21] 已經瘋狂到這個地Po ,希律下令屠殺無辜的男孩,一點也不出奇.

神創造人,讓人有自由意志作出抉擇.渴求永恆的人因自己的有限而感到不安,而心裡不安的人愈是希望可以通過掌控自己的命運以解除不安的困擾,愈是發覺自己的有限.解決之道是應該是信靠神,讓神來掌管自己的人生,但人選擇了不信,最終陷在罪的權勢之中,最終甚至作出種種瘋狂的行為,因而招致滅亡.這就是這個世界充滿了罪惡和混亂的根源.

或許大家心中會問,為何神不去救這些無辜遇害的男孩?經文沒有明確交代,神有沒有差派天使向這些男孩的父母報夢.問題是,即使有,又有多少人會相信而馬上收拾行裝逃走?不過這並非重點,更根本的問題是,除非能改變人因不信而陷在罪中的事實,否則這個邪惡的世界是不會改變,而無辜的人被害的事情,只會繼續在世界不同的角落發生.

神按著祂的應許,讓祂的聖道臨降人間,就是要打破罪權勢對人的捆鎖.

C. 神記念無辜的人

有學者嘗試從《聖經》以外的歷史文獻或記錄,尋找這次屠殺事件的外證,但暫時仍未找到,也很可能永遠都找不到.這是因為伯利恆只是一個細小的鄉鎮,人口並不多.學者一般估計,被殺害的男孩不會超過二十人,[22] 可能不足十個.相對於希律其他的殺戳,這顯然是小巫見大巫,沒有特別記載這件事也不足為奇.

屠殺男孩的事件大概發生在希律在生最後一年,頂多兩年,因為耶穌在主前四或五年誕生,而希律就在主前四年去世.希律後來得了大病,在臨終前就捉拿了大批猶太人的領袖,並命令他的手下在他死去後處決他們,因為希律認為人民不會為自己的死而哀號,但他至少要令大家不能因他的離世而歡欣.[23]

然而,《聖經》作為神的話語卻特意要把這批男孩遇害的事記錄下來,不只讓後世的人知道他們無辜被殺,要為他們哀哭,也要表明,神是記念他們的.這些神所記念的男孩今生的生命雖然短暫,但他們必然會在基督再來的時候復活,被引領到永世之中.

不僅如此,假若大家翻開馬太所引用《耶利米書》的經文,並且繼續讀下去,就會看到緊接在下面的,正是神對遭遇大能的以色列民的應許:

耶和華如此說: 「不要出聲哀哭, 你的眼目也不要流淚； 因你的辛勞必有報償, 他們必從仇敵之地歸回.」這是耶和華說的.「你末後必有指望, 你的兒女必回到自己的疆土.」這是耶和華說的.

(耶卅一16-17,和修版)

馬太要讓《馬太福音》的讀者知道,神的應許必然會在這個充滿苦難的世界中實現.

書目

1. Arndt, William et al., A Greek-English Lexicon of the New Testament and Other Early Christian Literature. Chicago: University of Chicago, 2000.

2. Blomberg, Craig. Matthew, vol. 22, The New American Commentary. Nashville: Broadman  and  Holman, 1992.

3. Carson, D. A. “Matthew,” in The Expositor’s Bible Commentary: Matthew, Mark, Luke, ed. Frank E. Gaebelein, vol. 8. Grand Rapids: Zondervan, 1984.

4. Clements, R. E. Jeremiah, Interpretation, a Bible Commentary for Teaching and Preaching. Atlanta: J. Knox, 1988.

5.Davies, W. D. and Dale C. Allison Jr., A Critical and Exegetical Commentary on the Gospel according to Saint Matthew, vol. 1, International Critical Commentary. London: T and T Clark, 2004.

6. Feinberg, Charles L. “Jeremiah,” in The Expositor’s Bible Commentary: Isaiah, Jeremiah, Lamentations, Ezekiel, ed. Frank E. Gaebelein, vol. 6. Grand Rapids: Zondervan, 1986.

7. Freedman, David Noel et al. The Anchor Yale Bible Dictionary. New York: Doubleday, 1992.

8. Hagner, Donald A. Matthew 1–13, vol. 33A, Word Biblical Commentary. Dallas: Word, Incorporated, 1993.

9. Hare, Douglas R. A. Matthew, Interpretation, a Bible Commentary for Teaching and Preaching. Louisville: John Knox, 1993.

10. Holladay, William Lee. Jeremiah 2: A Commentary on the Book of the Prophet Jeremiah, Chapters 26–52, ed. Paul D. Hanson, Hermeneia—a Critical and Historical Commentary on the Bible. Minneapolis: Fortress, 1989.

11. Lenski, R. C. H. The Interpretation of St. Matthew’s Gospel. Minneapolis: Augsburg Publishing, 1961.

12. Louw, Johannes P. and Eugene Albert Nida, Greek-English Lexicon of the New Testament: Based on Semantic Domains. New York: United Bible, 1996.

13. Luz, Ulrich. Matthew 1–7: A Commentary on Matthew 1–7, ed. Helmut Koester, Rev. ed., Hermeneia—a Critical and Historical Commentary on the Bible. Minneapolis: Fortress, 2007.

14. Morris, Leon. The Gospel according to Matthew, The Pillar New Testament Commentary. Grand Rapids: W.B. Eerdmans, 1992

15. Nolland, John. The Gospel of Matthew: A Commentary on the Greek Text, New International Greek Testament Commentary. Grand Rapids: W.B. Eerdmans, 2005.

注解

[1] Donald A. Hagner, Matthew 1–13, vol. 33A, Word Biblical Commentary (Dallas: Word, Incorporated, 1993), 27.

[2] John Nolland, The Gospel of Matthew: A Commentary on the Greek Text, New International Greek Testament Commentary (Grand Rapids: W.B. Eerdmans, 2005), 108.

[3] D. A. Carson, “Matthew,” in The Expositor’s Bible Commentary: Matthew, Mark, Luke, ed. Frank E. Gaebelein, vol. 8 (Grand Rapids: Zondervan, 1984), 85.

[4] 同上,86.

[5] William Arndt et al., A Greek-English Lexicon of the New Testament and Other Early Christian Literature (Chicago: University of Chicago, 2000), 990-991.

[6] Johannes P. Louw and Eugene Albert Nida, Greek-English Lexicon of the New Testament: Based on Semantic Domains (New York: United Bible Societies, 1996), 761.

[7] Willoughby C. Allen, A Critical and Exegetical Commentary on the Gospel according to S. Matthew, International Critical Commentary (New York: C. Scribner’s Sons, 1907), 15.

[8] Leon Morris, The Gospel according to Matthew, The Pillar New Testament Commentary (Grand Rapids: W.B. Eerdmans, 1992), 45.

[9] Craig Blomberg, Matthew, vol. 22, The New American Commentary (Nashville: Broadman  and  Holman, 1992), 68.

[10] 參Logos的Bible Word Study.

[11] William Arndt et al., A Greek-English Lexicon of the New Testament and Other Early Christian Literature (Chicago: University of Chicago, 2000), 323.

[12] 同上,279-280.

[13] William Arndt et al., A Greek-English Lexicon of the New Testament and Other Early Christian Literature (Chicago: University of Chicago, 2000), 475-477.

[14] William Arndt et al., A Greek-English Lexicon of the New Testament and Other Early Christian Literature (Chicago: University of Chicago, 2000), 1012-1013.

[15] R. C. H. Lenski, The Interpretation of St. Matthew’s Gospel (Minneapolis: Augsburg Publishing, 1961), 82.

[16] 同上.

[17] John Nolland, The Gospel of Matthew: A Commentary on the Greek Text, New International Greek Testament Commentary (Grand Rapids: W.B. Eerdmans, 2005), 124-125.

[18] 同上.

[19] David Noel Freedman et al., eds., The Anchor Yale Bible Dictionary (New York: Doubleday, 1992), 161-169.

[20] 同上.

[21] 同上.

[22] Donald A. Hagner, Matthew 1–13, vol. 33A, Word Biblical Commentary (Dallas: Word, Incorporated, 1993), 37.

[23] David Noel Freedman et al., eds., The Anchor Yale Bible Dictionary (New York: Doubleday, 1992), 161-169.



\newpage

\section{在我們當中住下的道}
\label{sec:731}
\textbf{高國雄牧師 - 在我們當中住下的「道」}
\newline
\newline
\hyperref[sec:730]{< < < PREV SERMON < < <}
~
\hyperref[sec:toc]{[返目錄]}
~
\hyperref[sec:732]{> > > NEXT SERMON > > >}


講題:在我們當中住下的「道」 “LOGOS” that Dwells among Us

講員:高國雄牧師

場合:香港中文大學崇基學院禮拜堂主日崇拜

日期:2025年1月5日

今日與大家談「道」

Λόγος(LOGOS)

約翰福音1:1-2

1太初有道,道與上帝同在,道就是上帝. 2這道太初與上帝同在.

中文聖經把Λόγος翻譯成「道」

相比起英文Word,「道」所包含的意思更豐富

WIKIPEDIA

邏各斯(古希臘語:λόγος)希臘語的原義是「話語」(英語:word)的意思,但是它包含了許多層面的意義.一方面它代表了語言,演說,交談,故事,原則等意涵,另一方面,它也代表了理性,思考,計算,關係,因果,類推等等.邏各斯通常使用在哲學之中,在希臘語中,一般的日常語言,或交談,是用另一個字古希臘語:λέξις(lexis)來代表,它與邏各斯都起源於相同的字根古希臘語:λέγω(legō),意思是交談,計算,說話或演說.

英語:logic(邏輯)的字根來自於邏各斯,各種學科的字尾(-logy)

它譯為拉丁語:ratio,後轉為法語:raison,最後成為英文中理智(reason)與理性(Rationality)的字根.

「道」一個非常精彩的翻譯

中國人的「道」,可以是道理,法則,描述形容整個世界

周易「一陰一陽之謂道」

道德經「道可道,非常道.名可名,非常名.無名天地之始；有名萬物之母.故常無欲,以觀其妙；常有欲,以觀其徼.此兩者,同出而異名,同謂之玄.玄之又玄,眾妙之門.」

今天教會:

講道,宣道,證道,傳道(人),「領受聖道」,聖禮聖道…

道是什麼呢?

道理,教訓?

「道」在創世時就發出

詩篇33:6

6諸天藉耶和華的命(logos)而造；萬象藉他口中的氣而成.

創世記1:1-3

1起初,上帝創造天地. 2地是空虛混沌,淵面黑暗；上帝的靈運行在水面上.3上帝說:「要有光」,就有了光.

「道」帶著能力,創造的能力

約翰福音4:49-50

49那大臣對他說:「先生,求你趁著我的孩子還沒有死就下去吧.」 50耶穌對他說:「回去吧,你的兒子會活!」那人信耶穌所說的話(logos),就回去了.

「道」帶著能力,醫治的能力,改變的能力

今天的經課

約翰福音1:14 道成了肉身,住在我們中間,充充滿滿地有恩典有真理.我們也見過他的榮光,正是父獨生子的榮光.

信仰中最重的,這「道」不單是「天道」,「神道」,而是成了肉身的「道」,更選擇住在我們當中.

人有不同的回應

約翰福音1:15約翰為他作見證,喊著說:「這就是我曾說:『那在我以後來的,反成了在我以前的,因他本來在我以前.』」

約翰福音1:10-11

10他在世界,世界也是藉著他造的,世界卻不認識他. 11他到自己的地方來,自己的人倒不接待他.

生命中不同的選擇,帶來不同的結果

約翰福音1:12-13

12凡接待他的,就是信他名的人,他就賜他們權柄作上帝的兒女. 13這等人不是從血氣生的,不是從情慾生的,也不是從人意生的,乃是從上帝生的.

「道」成為了唯一的連繫,使人與上帝能重建關係,更因此人能領受「上帝的兒女」的身份.

上帝的道,成為了通往上帝的道路

約翰福音14:6

耶穌對他說:「我就是道路,真理,生命；若不藉著我,沒有人能到父那裡去.」



\newpage

\section{從耶穌受洗說起}
\label{sec:732}
\textbf{林豪恩先生 - 從耶穌受洗說起}
\newline
\newline
\hyperref[sec:731]{< < < PREV SERMON < < <}
~
\hyperref[sec:toc]{[返目錄]}
~
\hyperref[sec:733]{> > > NEXT SERMON > > >}


講題:從耶穌受洗說起 Reflecting on Jesus' Baptism

講員:林豪恩先生

場合:香港中文大學崇基學院禮拜堂主日崇拜

日期:2025年1月12日

引言:

本主日是顯現後第一主日,也是耶穌受洗日.顯現期以耶穌受洗開始,以登山變像結束,首尾兩個主日的顏色都是白色,中間的是綠色,標示出這兩個主日的特別.

路加福音對耶穌受洗的描述是這樣的:

3:21眾百姓都受了洗,耶穌也受了洗.正禱告的時候,天就開了,

3:22聖靈降臨在他身上,形狀彷彿鴿子；又有聲音從天上來,說:「你是我的愛子,我喜悅你.」

耶穌受洗的事件大家都耳熟能詳,其中幾個元素相信大家都看得到,包括河水,聖靈及天上來的聲音.今次,讓我們在幾段經課的引導之下再思考這幾項元素的意義及我們的啟發.

一．河水與歸屬

路加福音這段記述耶穌受洗的文字雖然沒有提及河水,但讀者必定假設有河水在其中,因為前文指出洗禮的地方是約旦河.耶穌洗禮的河水讓我們想到甚麼呢?有沒有印象我們剛才讀過的舊約經課以賽亞書也提及河水呢?

43:2你從水中經過,我必與你同在；你渡過江河,水必不漫過你；你從火中行過,必不被燒,火焰也不著在你身上.

以賽亞書43章2節提及的河水和耶穌受洗事件的河水,兩種場景給人的感覺很不一樣.在耶穌受洗的約旦河,人們是自願及帶著盼望踏入河水中接受洗禮；但在以賽亞書中,人們是無奈地及驚恐地被逼走入河水的.經文把水與火並排:「從水中經過,從火中行走」,喻意艱辛凶險的路程.這是描述以色列人被打敗,家園被毀,被擄於外邦的慘痛經歷.國破家亡,作為戰敗者以戰俘的身分存活於打敗自己的國家,會是怎樣的心境呢?會否懷念故國,追憶往事呢?古代被滅國後,被擄於北宋的李後主曾經如此哀嘆:

春花秋月何時了,往事知多少? 小樓昨夜又東風,故國不堪回首月明中!

亡國被擄的悲哀何時完結呢?對將來的絕望何時終止呢?李後主的詞繼續訴說:

問君能有幾多愁?恰似一江春水向東流!

對於沒有信仰的人,世事只是成王敗寇.但對於敬拜耶和華的色列人卻是雙重加擊.戰敗代表甚麼呢?是不是耶和華沒有能力保護我們呢?如果耶和華有能力但不保護我們,是不是代表他已經放棄我們了呢?思前想後,看不見摩西再世,也看不見大衛再臨.不禁念天地之悠悠,獨愴然而涕下.

當以色列人陷入水火之中,感到被拋棄,孤苦無依的時候,耶和華透過先知以賽亞向他們說:

43:1雅各啊,創造你的耶和華,以色列啊,造成你的那位,現在如此說:你不要害怕!因為我救贖了你.我曾提你的名召你,你是屬我的.

43:4因我看你為寶為尊；又因我愛你,所以我使人代替你,使列邦人替換你的生命.

按人的期望,如果我是屬你的,如果你愛我,你就不會讓我落在如此艱難的處境中.現實是,陷入困境有很多原因,可能是被逼害,也可能是自作自受；可能有跡可尋,也可能不知何故.耶和華對身陷困境的以色列人的回應是:我沒有放棄你們,我仍然愛你們.以賽亞書的河水並不浪漫,反而是艱苦的經驗.踏過艱辛,被擄的以色列人在茫茫天地之中知道自己情歸何處,確認自己並不是沒有人要的孤兒.「沒有人要我」,是人生最淒涼的狀況之一.曾經有失戀的同學來哭訴:「佢唔要我呀!我以後點算呀?」在婚姻中感到被遺棄者的淒涼詰問也是:「點解佢唔要我?」工作中的人有時也有「公司幾時唔要我」的恐懼,有些老人家也擔心「仔女會唔會唔要自己」.但上主的應許是「必不撇下你,也不丟棄你,不要懼怕,也不要驚惶」(申命記31:8)

河水這項元素提示我們參考經課中這段以賽亞書思考耶穌的洗禮,他並不是如其他人一樣要洗淨自身的罪污,反而是展示著如何在天地間確認自己的歸屬,在時空中肯定的自己的身分.何以如此說呢?讓我們繼續閱讀耶穌受洗的經文.

二．聖靈與接納

3:21眾百姓都受了洗,耶穌也受了洗.正禱告的時候,天就開了,

3:22聖靈降臨在他身上,形狀彷彿鴿子；又有聲音從天上來,說:「你是我的愛子,我喜悅你.」

除了河水之外,耶穌受洗的另一項元素就是聖靈降臨.聖靈降臨帶我們去讀新約經課使徒行傳:

8:14使徒在耶路撒冷聽見撒馬利亞人領受了上帝的道,就打發彼得,約翰往他們那裡去.

8:15兩個人到了,就為他們禱告,要叫他們受聖靈.

8:16因為聖靈還沒有降在他們一個人身上,他們只奉主耶穌的名受了洗.

8:17於是使徒按手在他們頭上,他們就受了聖靈.

聖靈降臨是上帝同在的證明,代表著上帝接納這人,確認這人歸屬自己.這裡已經是使徒行傳第八章,聖靈降臨不是新鮮的事情,新鮮的是聖靈降臨在撒馬利亞人身上.我們作為局外人,對於撒馬利亞人獲得上帝的接納不會感到新鮮,但對於當時的猶太人,卻是出乎他們意料之外.以下這段記載在路加福音的事件反映出當時猶太人和撒馬利亞人的交惡狀況:

9:51耶穌被接上升的日子將到,他決定面向耶路撒冷走去.

9:52他打發使者在他前頭走；他們進了撒瑪利亞的一個村莊,要為他作準備.

9:53那裡的人不接待他,因為他面向著耶路撒冷去.

9:54他的門徒雅各和約翰看見了,就說:「主啊!你要我們吩咐火從天上降下來,燒滅他們嗎?」

從新約經課使徒行傳這段經文再看耶穌受洗時聖靈降臨,會帶給我們甚麼啟發呢?會不會也是天父對耶穌的接納的宣告呢?從經課使徒行傳的事件來看,聖靈降臨代表著上帝的接納,甚至是接納那些我們不接納的人,或者是接納那些我們應為上帝不應接納的人.當你面對上主的時候,你會否擔心上主不接納你呢?會否感到自己不值得被上主接納呢?有沒有想過這些擔心和感覺可能不是來自上主,而是來自我們在人間經驗的不被接納的經驗.上主是怎樣說的呢?讓我們再讀耶穌受洗的第三個元素.

三．愛子與喜悅

3:22聖靈降臨在他身上,形狀彷彿鴿子；又有聲音從天上來,說:「你是我的愛子,我喜悅你.」

「你是我的愛子,我喜悅你」這是多麼重要的一句說話,很多人一生的渴望就是聽到這句說話.

近期的話題電影「破地獄」中有一個家庭有一位父親有一子一女.女兒郭文玥感到父親重男輕女,在父親心目中自己身為女兒身是「污穢」的.這個標籤如影隨形般束縛著她,成為她揮之不去的心結,認定父親輕視和拒絕她.兒子郭志斌其實卻感覺不到父親特別重視他,反而感到父親對他特別嚴厲,在儀式上有何表現不足,往往被嚴厲斥責.他認定父親看不起他,認為他無用.這位兒子及女兒心底裡的渴望甚麼呢?會不會都是這一句「你是我的愛子／愛女,我喜悅你」呢?最後,郭文玥的心結是如何解開的呢?就是讀到父親在遺書中告訴她:「玥」的意思是「珍寶」,「文玥」即是「文哥的珍寶」.原來父親一直把自己視作珍寶,這是一直沒有說出口的愛意.

能否不需要等到遺書才能得到這種肯定呢?可否不需要透過遺書來表達愛意呢?最近我收到亞洲家庭治療學院的一篇文享文章,其中有幾段是這樣寫的:

「踏進三十歲那一年,我幾乎雙失.我很想知道自己的人生究竟發生了什麼事?於是乎我去上一個心理學家辦的課程.導師一開始就說:「你要你的人生好事業好愛情好,就要跟父母關係好!」我心想「今次大件事了!」跟媽媽的關係好得像閨密一樣無所不談,但跟爸爸……基本上是「零」溝通.我一直十分討厭他的冷漠,無理,自私,野蠻與及對媽媽的不尊重,又對我們兄妹三人愛理不理.課堂上,導師邀請我分享與爸爸的關係,那一刻我泣不成聲,剎那間小時侯扭著要爸爸抱的畫面浮現在眼前.原來我潛意識仍像個小女孩般渴望爸爸的愛,意識上卻對他不斷批判並一直疏遠他.

課程有一個練習,要我們打電話去多謝一個很重視的人.或許是上天的安排,那天媽媽不方便,那我只好打給爸爸!坦白說,我真的不知如何開始,也想不到要多謝他甚麼.我故意先找其他話題開腔.

「我們打算和媽媽去日本旅行,你一唔一齊去呀?」

爸爸:「唔去啦,日本人殺jor 咁多中國人,我不去呢D鬼地方,你地自己去啦!」

冷不防他這樣回應,我呆了一會,定神過後,生硬地說:「我同妹妹都你一齊去呀!」

爸爸:「我唔去啦!不如去夏威夷啦,你媽媽一直都想去既!」

我又呆了.為什麼爸爸會這樣說呢?他什麼時候會這樣關心媽媽呢?從小到大,絕大部份時間只是目睹他們兩人在吵吵鬧鬧,是不是我誤會了什麼呢?我鼻子一酸,忍不住在抽泣,但又怕爸爸聽到,那一刻真的很混亂.趁著未泣不成聲我趕緊說:「爸爸,多謝你……多謝你一直痛錫我地.」話未完已淚流滿面,忍不住哭了出聲.

爸爸:「你係邊度呀?唔好喊啦,佢唔arm 你就算啦,搵過第個啦!不要緊嫁!你係邊度呀?我來搵你啦.」那一刻,是我人生第一次感受到他真誠的關心,來自那個我曾經討厭,既冷漠又自私的爸爸.我已經泣不成聲,只盡量忍住抽泣聖潔尾:「我無事!只係想說聲多謝!我收線啦!」

幾年後在我的婚禮上,我能夠用最真誠的心去答謝爸爸媽媽.而這一次,泣不成聲的是爸爸.」

渴望被接納,渴望被珍惜相信是每一個人與生俱來的基本需要,現實上卻不容易滿足.有時,人們為了證明自己值得被接納和珍惜,努力表現自己,甚至排斥他人,但仍然患得患失,掙扎不安.來到上主面前,我們不需要證明自己值得,因為上主只想對我們說:「你是我的愛子,我喜悅你.」

四．總結:我們都是上帝的兒女

「你是我的愛子,我喜悅你.」這句說話只是對耶穌說的,還是都對我們說的呢?

其實人類出現的時候,上帝已經對人類說:「我喜悅你!」(創世記):

1:27上帝就照著自己的形像造人,乃是照著他的形像造男造女.

1:31上帝看著一切所造的都甚好.

「你是我的愛子,我喜悅你」這句說話帶我們回到詩篇經課:

29:1上帝的眾子啊,你們要將榮耀,能力歸給耶和華,歸給耶和華!

29:2要將耶和華的名所當得的榮耀歸給他,以聖潔的妝飾敬拜耶和華.

經課提示我們,「你是我的愛子,我喜悅你」這句說話不單是對耶穌說的,不單是對你說的,還對甚麼人說的呢?讓我們用以下這個小故事作結.

有一位拉比向學生提出一個問題:“你們怎麼知道黑夜何時結束,白天何時開始嗎?”

"是太陽出來的時候嗎?

“不是的,那並不是黑夜結束和白天開始的時刻.”

另一位學生:“那是當你看著一棟建築物,能分辨出它的顏色的時候嗎?”

“不是的,那也不是黑晚結束和白天開始的時刻.”

學生們繼續提出答案,但都不正確.最後,他問拉比:“請告訴我們,怎樣才能知道黑晚何時結束,白天何時開始.”

拉比回答說,“當你看到一個陌生人的臉,能夠在他的臉上看到他是你的兄弟姐妹的時候,就是黑晚結束,白天開始的時候了.若然你看不到他都是上帝的兒女,無論光線多麼明亮,仍然是黑夜.”各位,當你照鏡的時候,你看到自己是上帝的兒子嗎?你聽到上帝對你說:「你是我的愛子／女,我喜悅你」嗎?黑夜已經過去,白日已經開始,願我們都看到自己及其他人都是上帝喜悅的兒女.



\newpage

\section{相煎何太急}
\label{sec:733}
\textbf{鄧瑞強博士 - 相煎何太急}
\newline
\newline
\hyperref[sec:732]{< < < PREV SERMON < < <}
~
\hyperref[sec:toc]{[返目錄]}
~
\hyperref[sec:734]{> > > NEXT SERMON > > >}


俄巴底亞書 1:1-21（和合本修訂版）
\newline
\begin{longtable}{cl}
\hline
\hline
章節 & 經文 (和合本修訂版)\\
\hline
1:1 & \begin{tabularx}{0.7\textwidth}{X} 俄巴底亞所見的異象。論以東說：我從耶和華那裡聽見信息，並有使者被差往列國去，說：起來吧，一同起來與以東爭戰！ \end{tabularx} \\ \\ \relax
1:2 & \begin{tabularx}{0.7\textwidth}{X} 看哪，我要使你在列國中為最小，被人大大藐視。 \end{tabularx} \\ \\ \relax
1:3 & \begin{tabularx}{0.7\textwidth}{X} 你狂傲的心欺騙了你，你住在巖穴，居所在高處，心裡說：「誰能把我拉下來到地上呢？」 \end{tabularx} \\ \\ \relax
1:4 & \begin{tabularx}{0.7\textwidth}{X} 你雖如鷹高飛，在星宿之間搭窩，我必從那裡拉你下來。這是耶和華說的。 \end{tabularx} \\ \\ \relax
1:5 & \begin{tabularx}{0.7\textwidth}{X} 盜賊若來到你那裡，小偷夜間來到，豈不是只偷他們所需要的嗎？摘葡萄的若來到你那裡，豈不留下幾串嗎？你竟全然滅絕！ \end{tabularx} \\ \\ \relax
1:6 & \begin{tabularx}{0.7\textwidth}{X} 以掃遭到搜查，他隱藏的寶物竟被尋出！ \end{tabularx} \\ \\ \relax
1:7 & \begin{tabularx}{0.7\textwidth}{X} 與你結盟的都驅趕你，直到邊界，與你和好的欺騙你，勝過你，吃你飯的人設下圈套陷害你---他卻毫無聰明。 \end{tabularx} \\ \\ \relax
1:8 & \begin{tabularx}{0.7\textwidth}{X} 到那日，我豈不從以東除滅智慧人？從以掃山除滅聰明人？這是耶和華說的。 \end{tabularx} \\ \\ \relax
1:9 & \begin{tabularx}{0.7\textwidth}{X} 提幔哪，你的勇士必驚惶，以致以掃山的人都被殺戮剪除。 \end{tabularx} \\ \\ \relax
1:10 & \begin{tabularx}{0.7\textwidth}{X} 因你向兄弟雅各施暴，你必蒙羞，永被剪除。 \end{tabularx} \\ \\ \relax
1:11 & \begin{tabularx}{0.7\textwidth}{X} 當陌生人擄掠雅各的財物，當外邦人進入他的城門，為耶路撒冷抽籤分取財物的日子，你竟站在一旁，像與他們同夥。 \end{tabularx} \\ \\ \relax
1:12 & \begin{tabularx}{0.7\textwidth}{X} 你兄弟遭難的日子，你不該瞪著眼看；猶大人被滅的日子，你不該幸災樂禍；他們遭難的日子，你不該說狂傲的話。 \end{tabularx} \\ \\ \relax
1:13 & \begin{tabularx}{0.7\textwidth}{X} 我子民遭災的日子，你不該進他們的城門；他們遭災的日子，你不該瞪著眼看他們受苦；他們遭災的日子，你不該伸手搶他們的財物。 \end{tabularx} \\ \\ \relax
1:14 & \begin{tabularx}{0.7\textwidth}{X} 他們遭難的日子，你不該站在岔路口剪除他們逃脫的人，你不該交出他們的倖存者。 \end{tabularx} \\ \\ \relax
1:15 & \begin{tabularx}{0.7\textwidth}{X} 耶和華的日子臨近萬國；你所做的，人也必向你照樣做，你的報應必歸到自己頭上。 \end{tabularx} \\ \\ \relax
1:16 & \begin{tabularx}{0.7\textwidth}{X} 你們在我聖山怎樣喝了苦杯，萬國必照樣不停地喝，且喝且吞，他們就必歸於無有。 \end{tabularx} \\ \\ \relax
1:17 & \begin{tabularx}{0.7\textwidth}{X} 但在錫安山必有逃脫的人，那山必成為聖；雅各家必得原有的產業。 \end{tabularx} \\ \\ \relax
1:18 & \begin{tabularx}{0.7\textwidth}{X} 雅各家必成為大火，約瑟家成為火焰；以掃家必如碎秸，遭燃燒，被吞滅，以掃家必無倖存者。這是耶和華說的。 \end{tabularx} \\ \\ \relax
1:19 & \begin{tabularx}{0.7\textwidth}{X} 他們必得尼革夫和以掃山，得謝非拉，非利士人之地，他們必得以法蓮地和撒瑪利亞地，得便雅憫和基列； \end{tabularx} \\ \\ \relax
1:20 & \begin{tabularx}{0.7\textwidth}{X} 被擄的以色列大軍必得迦南人的地，直到撒勒法，在西法拉被擄的耶路撒冷人必得尼革夫的城鎮。 \end{tabularx} \\ \\ \relax
1:21 & \begin{tabularx}{0.7\textwidth}{X} 必有一些解救者上到錫安山，審判以掃山，國度就歸耶和華了。 \end{tabularx} \\ \\
[1ex]
\hline
\hline
\end{longtable}
講題:相煎何太急 Hey, He is Your Brother

經文:俄巴底亞書1:1-21(和合本修訂版)

講員:鄧瑞強博士

場合:香港中文大學崇基學院禮拜堂主日崇拜

日期:2025年1月19日

各位弟兄姊妹,早安.願神的恩典常與你們同在.

今日講舊約最短的一卷書,全書只有1章,只有21節經文,就是俄巴底亞書.「俄巴底亞」的意思,就是「耶和華的僕人」,一個很謙卑的名字.從內容看,俄巴底亞書是對猶大國的國破家亡,作一信仰回應.公元前586年,巴比倫大軍消滅猶大國,摧毀耶路撒冷聖殿,對猶大民族而言,這是他們從未經歷過的大劫.香港一個小小的城市,經歷少許風雨,已經造成很大的震盪,想想當時的猶大國,經歷毀滅性的重創,其悲痛是無法形容的.戰亂過後,未死的人多少會有創傷後遺症.有些人意志消沉,有些人對前景絕望,有些人迷失了自己的身分,最嚴重的,是很多人失去了安身立命的基礎,那個曾承諾保護他們的耶和華,還信得過嗎?信仰崩潰,價值崩潰,人便失去存活的意義.

在歷史的大劫裡,猶大民族在消亡的邊緣,畢竟,猶大民族是很弱小的民族,隨時會在歷史中消失.事實上,很多弱小的民族,被強大的帝國蹂躪後,便在歷史中消失了.就在這關鍵時刻,猶大民族湧現不少出色的神學家,他們領悟天理,以從天而來的視野教人重新理解苦難,以從神而來的遠象指導同胞看到黑暗前面的光明,以耶和華的信仰重新拼湊大劫下的歷史碎片.現實,像丟落地上的一大幅拼圖,混亂不堪,這一小片拼圖與那一小片拼圖沒有接合,看不到任何有意義的畫面.但這些偉大的神學家,以超然的信仰視角,將這些亂七八糟的拼圖,一塊一塊的拼合起來,讓人看到全幅畫面,一幅從神而來的,充滿意義的畫面.當人看到破碎畫面背後那全新而完整的意義時,人就找到應對逆境,存活下去的勇氣.先知俄巴底亞,就是其中一位偉大的神學家,以短短21節經文,支撐起一個衰敗的民族.

今日,不少基督徒的信仰,主要是為了將來上天堂.但在聖經裡,先哲寫下的信仰文字,往往寫於生死存亡的關頭,給人存活下去的力量,激發人走下去的勇氣,讓人意識到一個人立在天地間的存在意義.

我們看看俄巴底亞書這21節經文的巨大信仰力量.

由於《新標點和合本》在俄巴底亞書的翻譯上,有少許沙石,我今天採用《和合本修訂版》的經文.

俄巴底亞書大致可分成兩部分.前半部分1-16節是神審判邪惡,這邪惡的代表就是「以東」.後半部分17-21是受苦者猶大的復興,猶大民族雖然弱小,但命不該絕,走下去,會見到光明.

(1)開場白

俄1 俄巴底亞所見的異象.

我們從耶和華那裡得到消息,有使者被差往列國去:「起來吧,我們要起來與以東爭戰!」主耶和華論以東如此說......

耶和華在天上決議,派使者向以東宣戰,又派出戰士(經文中的「我們」),與以東作戰.神要恢復人間正義,難免一場正邪之戰.

在辦公室,有人被欺凌,最怕是老闆不知道,沒有及時作出處理.在人間,一個民族欺凌另一個民族,先知說,萬有之主是知道的,祂會適時作出行動.單就這一點,已叫被欺凌的猶大在逆境中保持希望.

(2)心高氣傲的以東

俄2 看哪,我要使你在列國中為最小,被人大大藐視.

俄3 你狂傲的心欺騙了你,你住在巖穴,居所在高處,心裡說:「誰能把我拉下來到地上呢?」

俄4 你雖如鷹高飛,在星宿之間搭窩,我必從那裡拉你下來.這是耶和華說的.

以東這民族,其城堡建在高山上,如鷹巢,經文說他「如鷹高飛,在星宿之間搭窩」.提外話,在德國最南部地區,某座山上有個地方叫「鷹巢」,是希特勒的別墅.這居高臨下的位置,象徵以東的心高氣傲.

以東以為他們的城堡難以攻破,覺得很安全,便大膽地欺凌弱小的猶大.他們心裡說:「誰能把我拉下來到地上呢?」神就說:「我必從那裡拉你下來.」神不喜悅人驕傲,要將以東拉下來,成為人人藐視的國.

短短的經文,提升被壓下的猶大民族的心情.將你壓下來的,無論他有多大本事,神會將他壓下來,挫敗他的傲氣.總有一天,猶大民族將可以抬起頭來.

(3)審判以東不顧手足之情

俄5 盜賊若來到你那裡,小偷夜間來到,豈不是只偷他們所需要的嗎?摘葡萄的若來到你那裡,豈不留下幾串嗎?你竟全然滅絕!

俄6 以掃遭到搜查,他隱藏的寶物竟被尋出!

俄7 與你結盟的都驅趕你,直到邊界,與你和好的欺騙你,勝過你,吃你飯的人設下圈套陷害你——他卻毫無聰明〔按:這一小節應接去下一段落〕.

第5節經文,點出以東的罪狀,他貪得無厭.小偷來,豈不只是偷所需要的；摘葡萄的來,豈不留下幾串?盜亦有道,賊也懂適可而止,留有餘地.但以東,當欺凌猶大時,竟不留餘地,要「全然滅絕」.正因如此,以東遭到報應時,也被別人搶走一切,變得一無所有.

特別留意,第6節稱以東為「以掃」.以掃是以東民族的先祖,而以掃是雅各的哥哥,雅各是猶大民族的先祖.以東和猶大原來是兄弟民族,有著同一血緣,也有著相近的信仰.先知在這裡,其實是控訴以東竟不念親情,本是同更新,相煎何太急.理應顧念手足而不顧,當然罪加一等.

由於以東不念手足之情,將有一天,他自己信任的「手足」,也會出賣他.第7節經文說,與以東結盟的,和好的,一齊吃飯的,這些人應是他可信任的手足,但全都出賣他.他對兄弟不義,有一天,他會吃盡同樣的苦果.神不希望這個世界,淪落成出賣手足的不義世界.

(4)審判以東不施援手

俄7 ......他卻毫無聰明.

俄8 到那日,我豈不從以東除滅智慧人?從以掃山除滅聰明人?這是耶和華說的.

俄9 提幔哪,你的勇士必驚惶,以致以掃山的人都被殺戮剪除.

俄10 因你向兄弟雅各施暴,你必蒙羞,永被剪除.

俄11 當陌生人擄掠雅各的財物,當外邦人進入他的城門,為耶路撒冷抽籤分取財物的日子,你竟站在一旁,像與他們同夥.

留意第11節,講到巴比倫攻陷耶路撒冷時的光景,財物被掠,城門被攻破,入侵者抽籤分贓.這時,以東在場,在犯罪現場,他們是在場的,卻沒有幫助自己的手足.

一個有能力幫助自己的兄弟,卻在緊急關頭沒有提供幫助的人,神將廢去他的「武功」.第7-8節經文講到,以東將失去聰明,即戰爭的謀略.第9節經文講到,他將失去戰爭的勇氣.頭腦沒有謀略,手腳沒有勇力,就像被廢武功的人.一個有力量而不用於行善的人,他將失去力量.

(5)審判以東助紂為虐

以下一段經文,不斷出現重複的字眼,盡數以東的不義.

俄12 你兄弟遭難的日子〔按,原文是:你兄弟的日子,他遭難的日子〕,你不該瞪著眼看；猶大人被滅的日子,你不該幸災樂禍；他們遭難的日子,你不該說狂傲的話.

俄13 我子民遭災的日子,你不該進他們的城門；他們遭災的日子,你不該瞪著眼看他們受苦；他們遭災的日子,你不該伸手搶他們的財物.

俄14 他們遭難的日子,你不該站在岔路口剪除他們逃脫的人,你不該交出他們的倖存者.

第12節開頭說「你兄弟的日子」,什麼日子呀?一連講了七次,就是「遭難」,「被滅」,「遭災」的日子.然後,一連八次說:「你不該...」,「你不該...」,「你不該...」.先知好像也對我們說:當你身邊的弟兄姊妹不開心時,你不該忽略他的需要；當你身邊的弟兄姊妹遇上信仰危機時,你不該吝嗇你對他的鼓勵；當你身邊的弟兄姊妹心中痛苦時,你不該輕視他的痛苦.當猶大遭逢劫難時,以東的反應,完全不合格.幾節經文,暴露我們對他人苦難的漠視.人類應是同胞,手足,很多時,我們辜負了手足之情.

最嚴重的,是第14節:「他們遭難的日子,你不該站在岔路口剪除他們逃脫的人,你不該交出他們的倖存者.」這就不只是「食花生」,而是參與屠殺的行動了.在戰場中虎口逃生的人,以東竟攔截他們,將他們再送回虎口.以東已不顧親情到極點,一個活口也不讓留下.

神的審判是這樣的:

俄15 耶和華的日子臨近萬國；你所做的,人也必向你照樣做,你的報應必歸到自己頭上.

俄16 你們在我聖山怎樣喝了苦杯,萬國必照樣不停地喝,且喝且吞,他們就必歸於無有.

大家要留意,當先知一講到「耶和華的日子」,他就是要將我們的視點,從塵世的角度帶到神的角度,從當下的視野帶到永恆的視野,從歷史當下的一點帶到歷史終末時的全面圖畫.先知希望我們的視野,不要局限在當下的一個丁點裡,先知希望我們看到神眼中的歷史和世界.用我們豆大的眼光看歷史和世界,很多時會失望.但從永恆的角度看,看到的圖畫大很多,意義也豐富很多.神在萬國之上,在歷史之上,祂在俯視塵世,在祂認為最好的時機裡,祂會報應以東.第15節經文說,以東向人怎樣做,就受到相應的懲罰.猶大人,在受苦的當下,每日只感覺到苦,是看不到任何有意義的畫面的.先知的說話,開他們的眼睛,讓他們知道,神是會主持正義的.加害人的人,是會受到相應懲罰的.

在「耶和華的日子」裡,神不單處置以東,也審判萬國.萬國必為他們做的事情,承擔相應的結果.天地間,有神在主持正義.有神主持公道,受苦者才能保持希望.

真正的正義,不單是惡人受罰,也必須包含受苦的弱者的復興,這就是俄巴底亞書後半部分17-21節的重點.

(6)命運的逆轉:猶大的復興

俄17 但在錫安山必有逃脫的人,那山必成為聖；雅各家必得原有的產業.

錫安山,雅各家,全都指猶大民族.他們命不該絕,留有餘種.不單留有餘種,還會復興.從失去家園,到最後,會再次重建家園.

俄18 雅各家必成為大火,約瑟家成為火焰；以掃家必如碎稭,遭燃燒,被吞滅,以掃家必無倖存者.這是耶和華說的.

毀滅別人家園的,將失去家園；吞滅人的,反被吞滅；不讓人有倖存者的,自己反而沒有了倖存者.這是命運的逆轉,或者說,是不義的世界恢復正義.

俄19 他們必得尼革夫和以掃山,得謝非拉,非利士人之地,他們必得以法蓮地和撒瑪利亞地,得便雅憫和基列；

俄20 被擄的以色列大軍必得迦南人的地,直到撒勒法,在西法拉被擄的耶路撒冷人必得尼革夫的城鎮.

失去土地的,會重獲土地；掠奪土地的,必失去土地.不同的民族在執位,因為神在扭轉乾坤.

俄巴底亞書的結語是:

俄21 必有一些解救者上到錫安山,審判以掃山,國度就歸耶和華了.

俄巴底亞書的開頭,講到神派戰士去與以東爭戰,克服不義,這裡講到「解救者」,作一對應.對弱者來說,最要緊的,是有人解救他們.至於欺壓人的,會受到審判.最後,人看到的歷史終點是:神作王,統治萬有,神國在人間.神國降臨的期盼,燃點起受苦者的希望.對前景有希望,人才能在痛苦的逆境中活下去.

短短21節經文,討論了很多主題,包括:神在天庭上作出了正義的決定,祂有適切的行動,會與邪惡爭戰,祂令驕傲的人謙卑下來,作出公義的審判,讓邪惡的人得到相應的刑罰,人與人之間應有兄弟情誼,手足之間不應互相出賣,要有同理心去體會他人的痛苦,要及時伸出援手,弱者會留有餘種,有神定好的耶和華的日子,有終末的遠象,有救贖者的來臨,弱者的命運會被扭轉,弱者有復興的前景,能重建家園,最後是神親自管治,讓大地回復正義.

看,短短的經文,承載著這麼多觀念,這些都是人類期盼的天國的相關觀念.這些觀念,支撐著受苦的猶大民族,讓他們對民族的明天,正義的前景,歷史的將來保持希望.若沒有希望,他們早就在歷史中消亡了.有了希望,他們會努力活下去.他們的存活,訴說著神的救贖故事,傳遞著天國的聖火,讓一代又一代受苦的人,在苦難中,看到天國的希望.

耶穌呼召我們,將天國的聖火交給我們,我們能好好傳遞這聖火嗎?

願平安臨到你們.願榮耀歸給聖父,聖子,聖靈的名,阿們.



\newpage

\section{耶穌講道}
\label{sec:734}
\textbf{李均熊牧師 - 耶穌講道}
\newline
\newline
\hyperref[sec:733]{< < < PREV SERMON < < <}
~
\hyperref[sec:toc]{[返目錄]}
~
\hyperref[sec:735]{> > > NEXT SERMON > > >}


講題:耶穌講道 The Preaching of Jesus

講員:李均熊牧師

場合:香港中文大學崇基學院禮拜堂主日崇拜

日期:2025年1月26日

起:路加福音中耶穌講道

有些人喜歡稱今日的經課為耶穌第一次講道.但經課的頭兩節,清楚的讓我們知道耶穌其實已經在加利利地區「在各會堂裡教訓人」,所以理論上路加記載的並非耶穌初試啼聲的第一篇講章,因為耶穌已經是一位廣受人歡迎的講員,因為他在各會堂內教訓人,眾人都稱讚他.不過,若我們比較其他福音書,路加的記載,是將耶穌呼召門徒,和傳道治病趕鬼的工作ministry放在後面,而把耶穌進會堂講道教訓人,作為耶穌受洗之後的首要記載.相信這樣安排的目的,是以這篇講章作為耶穌在世服事的綱領,甚至可以說是路加福音先要說明這福音書要記載耶穌一生的焦點.所以路加福音中耶穌講道值得我們細心了解一下.

承:耶穌如何講道

不過要理解耶穌講道的信息,我們也許要先理解當時人是如何講道的.我們今日的崇拜都是大家坐在下面,講員通常在高出來的講壇講道.但第一世紀會堂的講道卻是調轉的,是先有人站起來讀經,然候講員是坐下來講道的.今日我們的講道是單向性,但當時的會堂的講道,更似一個經文辯論會,其他人和講員是互有答辯的.下週經課就會讓我們看到,耶穌進一步與眾人答辯,也因此激發眾人由稱讚變成想將耶穌推下山崖的反應.不過既然教會傳統將路加福音的這一段經文分為兩段經課,那麼眾人對耶穌講道的反映,今天就按下不表了.

耶穌這次講道的場景是拿撒勒的會堂,也值得我們了解一下.我們對第一世紀猶太人會堂的布置和安排,所知其實十分有限.近代學者都提醒我們,不要用今日已具規模的會堂作類比.一來昔日會堂未必有一個恆常的標準,而是各處鄉村各處例(different places have different customs)；拿撒勒亦只是一個不顯眼的鄉村小鎮,當中的會堂可能只是規模十分細小的地方,純粹是鄉親父老在安息日聚集的一個地方,平日也可能用來相聚,傾談,商討鄉事的地方.甚至有學者推斷拿撒勒的會堂並沒有一整套完整的Torah托拉/摩西五經,有的經卷就可能只是我們今天福音經課中,有人交給耶穌的以賽亞書.至這也可能不是整卷的以賽亞書,而可能是半份或三分之一份的以賽亞書.大家要知道在公元第一世紀,大部份人都不懂寫字的社會,要抄寫一份經書,所需要的人力和成本是十分高昂的.耶穌講道的場境,亦並非如一些解經家所說,會堂裡有人讀完了律法書經課,然候耶穌才起來讀相配的先知書經課；這只是今日會堂的安排,昔日可能就是有人從藏經閣內拿出大家珍而重之的以賽亞手卷,交給耶穌,耶穌就找到要讀的經文,然候就讀起經來.

轉:耶穌講乜道?

那究竟耶穌講道的內容是甚麼呢?經課內我們只見耶穌站起來讀經,便坐下來.當大家眼定定看著耶穌,期望他講道解經的時候(「會堂裡的人都定睛看他」),耶穌就只說了一句:「今天這經應驗在你們耳中了.」大家就好像很滿足的,又稱讚他,「稀奇他口中所出的恩言」又說:「這不是約瑟的兒子嗎?」(22節在今天的經課是沒有的)!

試想像今日讀經員讀完經,大家唱完榮耀頌後,主席介紹講員李均熊牧師上台,點知李牧師只說了一句話(可能是:「你們就照著行吧」),就下台了,大家會有何反應?(可能ＡＡＡ[Alternative Activities to Assembly週會替代活動])的同學會鼓起掌來,估不到這麼快捱完個講道!於是乎眾人稱讚李牧師,又說:這不是神樓的舍監嗎……)

事情當然不會這麼簡單.

當然有機會路加福音的記載從略了耶穌的講解,因為整個路加福音就是這段經文的解讀；又或是期實下星期的經課才是耶穌的講解,路加福音的記載亦使大家焦點集中於聽眾前後不同的反應,而非耶穌的講道.但更有可能就是當時大家講道聽道的方式和今日是截然不同的.

今日大家坐在這裡聽道,都可能期望講員解釋一下經文的歷史背景,字義用法,或講明經文的意義和應用.原因其實是我們對聖經這本在二千多年前生成的文獻著實有太多不認識的地方,不單經文的語言是古代的語言,現代社會和古代社會也著實有太多差異,牧師傳道唯有進神學院讀三四年書,成為一位能溝通古今的講員.但當時的人基本上對經文瞭如指掌,故此如前述,就算沒有一本正式的經卷在前,古代人其實都能把經文記憶起來.而且猶太人的解釋聖經的方法是用另外的經文去註解要認識的經文,所以耶穌的講道的信息就在他所讀的經文內.

所以不講大家大概也不知道,耶穌宣讀的經文,並不是一段以賽亞書的經文,而是結合了兩段的經文,一段是今日我們聖經的以賽亞書61章1-2節:

61:1 主耶和華的靈在我身上,因為耶和華用膏膏我,叫我報好信息給貧窮的人,差遣我醫好傷心的人,報告被擄的得釋放,被捆綁的得自由(古希臘文譯本譯「瞎眼的得看見」)；61:2宣告耶和華的恩年(dekton)和我們的神報仇的日子；安慰所有悲哀的人,

另一段是以賽亞書58:5-7節:

58:5 這豈是我所要的禁食,為人所用以刻苦己心的日子嗎?我難道只是叫人如蘆葦般低頭,鋪上麻布和灰燼嗎?你能稱此為禁食,為耶和華所悅納的日子(dekten)嗎?58:6我所要的禁食,豈不是要你鬆開兇惡的繩,解開軛上的索,使被欺壓的得自由,折斷一切的軛嗎?58:7豈不是要你把食物分給飢餓的人,將流浪的窮人接到家中,見赤身的給他衣服遮體,而不隱藏自己避開你的骨肉嗎?

大家可以對比一下路加福音(以希臘文)記載耶穌宣讀的內容[粗體字就是引用的字眼]:

4:18主的靈在我身上,因為他用膏膏我,叫我傳福音給貧窮的人；差遣我報告:被擄的得釋放,瞎眼的得看見,叫那受壓制的得自由,4:19報告神悅納人的禧年(dekton).

很大可能是耶穌用賽58:5-7節的經文,註解賽66:1-2節的經文.

大家就這樣看應該會一頭霧水,正因為我們沒有當時人對以賽亞書的熟識,耶穌的以經解經對我們而言是沒有意思的.但當時在羅馬帝國統治下(因此成為?)窮鄉僻壤的拿撒勒的鄉親父老,卻聽得出耶穌的宣告是甚麼意思:那受膏者／彌賽亞來到要報告的上帝的禧年(陽性受格詞dekton)/悅納人的日子(陰性受格詞dekten),並不是來報復使猶太人傷心悲哀的外族統治者(省略的部份),是要帶來社會上貧窮,受欺壓的人得自由.但觀乎耶穌引用賽58:5-7的經文,欺壓人的人不單只是外族帝國,可能就是在坐的鄉親父老!要成就上主的鿋年,「豈不是要你把食物分給飢餓的人,將流浪的窮人接到家中,見赤身的給他衣服遮體,而不隱藏自己避開你的骨肉嗎?」(經文沒有明言的部份)耶穌的宣道,其實是很「有骨」的,很不明顯地將賽58章上主控訴當時的猶太人「齋口唔齋心」,口說要敬虔,要婪食,卻並不照顧貧窮有需要的人的意思,隱藏在賽61章彌賽亞禧年的經文內.所以拿撒勒的鄉親父老聽了即時好歡喜,因為耶穌宣告彌賽亞的時代要應驗了,上主要報復統治我們,使我們成為窮鄉僻壤的羅馬人了.拿撒勒的鄉親父老以為彌賽亞要來是要推翻羅馬帝國,但耶穌要帶來的彌賽亞時代,是基督要來關心和幫助貧苦大眾,使受欺壓的人得到釋放和自由,這就是路加福音被稱為貧窮者的福音的原因.

其實下星期經課,正是耶穌看見他們喜歡的反應,就知道他們不明白耶穌講道的意思.於是要進一步解釋一下:

4:23 耶穌對他們說:「你們必引這俗語向我說:『醫生,你醫治自己吧!我們聽見你在迦百農所行的事,也當行在你自己家鄉裡』」；4:24又說:「我實在告訴你們,沒有先知在自己家鄉被人悅納的.4:25我對你們說實話,當以利亞的時候,天閉塞了三年零六個月,遍地有大饑荒,那時,以色列中有許多寡婦,4:26以利亞並沒有奉差往她們一個人那裡去,只奉差往西頓的撒勒法一個寡婦那裡去.4:27先知以利沙的時候,以色列中有許多長大痲瘋的,但內中除了敘利亞國的乃縵,沒有一個得潔淨的.」4:28會堂裡的人聽見這話,都怒氣滿胸,4:29就起來攆他出城(他們的城造在山上)；他們帶他到山崖,要把他推下去.

可見,如果當時拿撒勒的鄉親父老真是明白彌賽亞／基督並不是要推翻羅馬帝國,而是要幫助不認識上主恩典的外邦人明白上帝的禧年,至於身邊的孤兒寡婦和有需要的人,是那些領受上主恩典的人都有份要去照顧的.若然大家都了解作基督的門徒不是受人服待,而是要如基督一樣去犧牲服待別人,應該也未必這麼高興.

合:

下星期的經課,我就不在這裡多講了.其實今天的經課還有很多地方仍需要深思,例如究竟耶穌是用甚麼語言去宣道?因為大家都見到,希伯來文的舊約經文和希臘文的舊約經文其實有一些差別(中譯希伯來文「被捆綁的得自由」古希臘文譯本卻譯「瞎眼的得看見」?),這些差別為何會出現?在耶穌的時候,大家又會否留意這些差別?路加福音是否刻意用上古希臘文的經文,是因為路加福音是寫給只認識希臘文的信徒看的嗎?那這段經文對這些後來的信徒又有甚麼意義?這些問題都不在這裡一一解答了.

今日我們聽道讀經可能只觸及表面的皮毛,或許,最後我邀請大家繼續思考,就是我真的明白耶穌作為基督——彌賽亞的身份和顯現的目的嗎?究竟我心目中所認識的耶穌基督,是一位怎樣的耶穌基督呢?



\newpage

\section{心想事成}
\label{sec:735}
\textbf{梁遠耀牧師 - 心想事成}
\newline
\newline
\hyperref[sec:734]{< < < PREV SERMON < < <}
~
\hyperref[sec:toc]{[返目錄]}
~
\hyperref[sec:736]{> > > NEXT SERMON > > >}


路加福音 4:21-30
\newline
\begin{longtable}{cl}
\hline
\hline
章節 & 經文 (和合本修訂版)\\
\hline
4:21 & \begin{tabularx}{0.7\textwidth}{X} 耶穌對他們說：「你們聽見的這段經文，今天已經應驗了。」 \end{tabularx} \\ \\ \relax
4:22 & \begin{tabularx}{0.7\textwidth}{X} 眾人都稱讚他，並對他口中所出的恩言感到驚訝；他們說：「這不是約瑟的兒子嗎？」 \end{tabularx} \\ \\ \relax
4:23 & \begin{tabularx}{0.7\textwidth}{X} 耶穌對他們說：「你們一定會用這俗語向我說：『醫生，你醫治自己吧！我們聽見你在迦百農所做的事，也該在你自己的家鄉做吧。』」 \end{tabularx} \\ \\ \relax
4:24 & \begin{tabularx}{0.7\textwidth}{X} 他又說：「我實在告訴你們，沒有先知在自己家鄉被人接納的。 \end{tabularx} \\ \\ \relax
4:25 & \begin{tabularx}{0.7\textwidth}{X} 我對你們說實話，在以利亞的時候，天閉塞了三年六個月，遍地有大饑荒，那時，以色列中有許多寡婦， \end{tabularx} \\ \\ \relax
4:26 & \begin{tabularx}{0.7\textwidth}{X} 以利亞並沒有奉差往她們中任何一個人那裡去，只奉差往西頓的撒勒法一個寡婦那裡去。 \end{tabularx} \\ \\ \relax
4:27 & \begin{tabularx}{0.7\textwidth}{X} 在以利沙先知的時候，以色列中有許多痲瘋病人，但除了敘利亞的乃縵，沒有一個得潔淨的。」 \end{tabularx} \\ \\ \relax
4:28 & \begin{tabularx}{0.7\textwidth}{X} 會堂裡的人聽見這些話，都怒氣填胸， \end{tabularx} \\ \\ \relax
4:29 & \begin{tabularx}{0.7\textwidth}{X} 就起來趕他出城。他們的城造在山上；他們帶他到山崖，要把他推下去。 \end{tabularx} \\ \\ \relax
4:30 & \begin{tabularx}{0.7\textwidth}{X} 他卻從他們中間穿過去，走了。 \end{tabularx} \\ \\
[1ex]
\hline
\hline
\end{longtable}
講題:心想事成 May All Your Wishes Come True

經文:路加福音4:21-30

講員:梁遠耀牧師

場合:香港中文大學崇基學院禮拜堂主日崇拜

日期:2025年2月2日

在新春佳節,分享主道前,首先讓我恭祝各位蛇年蒙福,靈巧像蛇,純良如鴿,出入平安.

大家不知有沒有見過或聽過,有些小朋友或大朋友,當他們向人拜年祝賀時,把不同的祝福語,琅琅上口,說出十多個四個字的祝福語.其實我小時候就想,那用說這麼多,只要一句:心想事成,都應該是已經可以包括全部的祝賀語了.

但我後來也明白,大家都只是喜歡聽「好入耳 / arm 聽」的說話,至於能否達成,則不是重點.因為如果我們相信這此祝賀語真的成真,照理我們在大年初一起床後要做的第一件事,就立即衝去廁所,或在梳妝台,或總之要找一面鏡,然後對著鏡上反映的自己,誠心祝賀自己:恭喜發財,青春常駐,準時收工,百事可樂,不勞而獲,成為世一,食極唔肥等等.我們有時真的只是喜歡聽「好入耳 / arm 聽」的說話,但我們有沒有為達成心願而做好預備.

若大家上週有參與崇基禮拜堂的主日崇拜,無論是現場出席還是在網上參與,聽到講員李均熊牧師(Rev. Dr. Li Kwan Hung Leo)的講道,你有聽到尾的話,今天應該會幾開心.開心的原因,因為李牧師上週講道的內容,已經講解了今天的新約經課,所以講道時間可以縮短5分鐘.李牧師在上星期已解釋了,當時拿撒勒會堂的群眾,他們起初對耶穌所讀的以賽亞書內容後,起初十分喜悅,因為這群會堂的會眾,很可能是誤會了耶穌宣讀以賽亞書61:1的意思,以為是「好入耳 / arm 聽」,講述彌賽亞要來使我得自由,使我得醫治,使我得幫助.但當耶穌表示這段以賽亞書的所宣告的禧年,連非只向猶太人宣告,同時非猶太人亦包容在內,也是上主所看顧的子民,就好像昔日的先知以利亞(Prophet Elijah)和先知以利沙( Prophet Elisha)的事蹟,上主看顧普世全人類.群眾聽到耶穌的講解後,立即轉變態度,由欣賞,肯定,變為討厭,甚至是憎厭,要動手對付耶穌,因為耶穌的講解「不好入耳 / 唔arm 聽」.

在很多年前,我以教牧人員的身份,參加一個聯校的中學畢業禮.當日畢業禮致勉勵詞的講員是蔡元雲醫生(Dr. Choi Yuen Wan Philemon).當日蔡醫在致勉勵詞開首,面帶笑容,祝賀同學們畢業了,這是同學們人生的高峰.當時大家聽到都在微笑,但蔡醫的面容立即變得嚴肅,說:你知不知,祝賀你今天到達高峰,是一項咒詛.因為這表示了今天到達高峰,明天走向低谷.蔡醫其實是想勉勵畢業同學,要不斷努力,不要自滿.

祝賀別人的學業,事奉達到高峰,可能很多人都會欣然接受,很少人會再細想其意思,總之「好入耳 / arm 聽」就好了.

有一個故事,兩個朋友在早上碰面,朋友甲一臉憔悴,沒精打采,朋友乙便走上前去關心朋友甲.乙問候甲:「你為何這樣憔悴?」甲回答:「我昨夜夢見摩西,他向我展示十誡,我看見後就覺得好難遵行,所以整夜失睡.」乙卻展靈笑容,說:「真巧,我昨夜也夢見摩西,同樣都向我展示十誡,我細心去睇,感到十分放鬆和愉悅.」甲向乙問:「為何你看到十條誡命,還會感到放鬆和愉悅?」乙回答:「因為在刻上十誡的石塊上,還有一行細字,你有看到嗎?」甲問:「是甚麼字,我沒細看.」乙說:「任擇五條.」信仰是沒有選擇,不能只要祝福,而沒有責任,因為信仰是要改變生命,就必然帶來祝福與責任.

在我小時候,有一牌子的朱古力,是一種一粒粒朱古力豆,但這牌子的朱古力豆,表面會有一層糖衣包裹.小時欣,這牌子的朱古力豆,在電視廣告裡,經常會重覆一句口號:「只溶係口,唔溶係手」(只溶在口,不溶在手),意思是這種朱古力豆,不會弄污我們的手.但我們的信仰,又是否「只溶係口,唔溶係手」(只溶在口,不溶在手)?當聖經提示我們要關心貧窮人時,我的信仰行動是否只是向對方表示會為你祈禱,求主供應,「只溶係口,唔溶係手」?當要關心有家歸不得,又或無家可歸的人,我的信仰行動是否只是祝福對方平安,願你吃得飽,穿得暖,「只溶係口,唔溶係手」?當要關心有身心障礙的人,我的信仰行動是否只說外面很危險,不要出來,你留在家,「只溶係口,唔溶係手」呢?

耶穌在家鄉拿撒勒會堂(Synagogue in Nazareth)的宣講,提醒上主的子民,不要把信仰上主視為只求獲得祝福和保守,忘記我們也要努力去迎接福氣.在主裡的心想事成,我們並不是把手交疊胸前,然後等待上主來完成的旁觀者,而是上主讓我們有機會參與在祂救恩計劃中的參與者.

在今天的舊約經課,是上主呼召耶利米,要他宣講「拔出,拆毀,毀壞,傾覆,又要建立,栽植」的訊息.這是「不好入耳 / 唔arm 聽」的訊息,我想耶利米都會明白,這是一份吃力不討好的事奉.故此,耶利米以自己年紀輕為理由,向上主推辭.但上主吩咐耶利米,請他不要擔心自己的年紀輕而不受人重視,因為上主必與耶利米同在.

童話故事《國王的新衣》裡,上至達官貴人,下至平民百姓,都只說國王「好入耳 / arm 聽」的說話,但對於國王沒穿外衣的醜態,則置之不理,就只是故事後來的童子,以孩子純真,率直和誠實的態度,指出了國王沒穿外衣的醜態.

我見過很多基督徒在婚禮上,都會選讀稱為「愛篇」的哥林多前書13:1-13,即今日的新約書信經課.可能因為經文不斷反覆出現「愛」字,大家就很喜歡.但我們只要細心看看就知,因為哥林多前書的其中一個主要重點,是處理哥林多教會的紛爭,教會內的信徒互相排斥,所以保羅就勸勉哥林多的信徒要以愛來互相連繫,而不是分享男女或婚姻情愛的甜蜜.雖然如此,「愛篇」仍然可能是基督徒婚禮上,最常被選讀的經文,因為「好入耳 / arm 聽」.

人喜歡聽「好入耳 / arm 聽」的話,其實也無可厚非.雖然我們明白我們被批評,也可以是我們改進的動力,經常被人批評,就總不是味兒.

例如若我今天祝高牧師身體健康,相信都是合宜和「好入耳 / arm 聽」.但隨後若在崇拜後,我們去餐廳食飯,侍應說今天有芋頭扣肉(Braised Pork Belly with Taro),要不要試?要心想事成,我就會請高牧師不要食了,但可以點,我可以打包外賣.

在蛇年,讓我講一個蛇的故事.有兩條蛇在對話,甲蛇向乙蛇詢問:「乙,你知不知我是否毒蛇」.乙蛇反問甲蛇:「你竟然連自己是否毒蛇都不知道.」甲蛇回答:「我想再確定.」乙蛇再反問:「你發生了甚麼事,要再確定自己是否毒蛇?」甲蛇回答:「因為我剛才不小心,自己咬到自己條利 / 舌頭.」

我們是上主的子民,當我們誦念《主禱文》時,我們祈願天父國度降臨,願天父旨意行在地上如同行在天上.在我們的生活上,又有否遵行上主的旨意而行?

我鼓勵大家訂立新年的靈命目標,無論是在讀經,祈禱,服侍還是傳福音上,都能持續成長.我們要相信,無論面對甚麼挑戰,上主都與我們同在,且必會帶領我們走過一切困難,成就祂的美意.在這新的一年裡,希望大家都立志遵行主道,讓我們的生命充滿信心與盼望,心想事成.



\newpage

\section{呼召}
\label{sec:736}
\textbf{李均熊牧師 - 呼召……}
\newline
\newline
\hyperref[sec:735]{< < < PREV SERMON < < <}
~
\hyperref[sec:toc]{[返目錄]}
~
\hyperref[sec:737]{> > > NEXT SERMON > > >}


講題:呼召…… Calling...

講員:李均熊牧師

場合:香港中文大學崇基學院禮拜堂主日崇拜

日期:2025年2月9日

起:是關於呼召不是?

今天的經課,無論是舊約經課(賽6:1-8),還是福音經課(路5:1-11),通常都被稱為呼召敘述(call narrative)．和合本的標題就直接用上「主呼召以賽亞作先知」,或「耶穌呼召頭幾個門徒」.但我們若仔細看經文,其實無論上主向以賽亞,或耶穌對彼得,都沒有直接的「呼召」.若你比對路加福音和馬可或馬太福音,路加只記載耶穌和彼得說「不要怕!從今以後,你要得人了.」而並沒有馬太或馬可福音記載耶穌簡單直接呼召彼得和他的兄弟(西門和安得烈):「來跟從我,我要叫你們得人如得魚一樣.」(太4:19；可1:17)其實路加福音的記載,和馬太和馬可福音的記載是這樣不同,有些人甚至認為路加福音記載的是馬可或馬太福音呼召彼得後的另一件事件.加上路加福音的記載,和約翰福音在最後所記載的,們徒在耶穌復活後在岸邊與他相遇的那次事件(約21章)是這樣的相似,就更加使路加福音的記載撲朔迷離,我們更難確定其實路加記載這段經文的重點是否就是耶穌呼召彼得去跟從他.

另一邊廂,以賽亞書第6章,記載一個「我」,相信就是以賽亞,在聖殿窺見上主的榮耀,並聽到上主說「我可以差遣誰呢?誰肯為我們去呢?」就急不及待回應「我在這裡,請差遣我!」你細心想一下,這並非一個直接的呼召,更似是以賽亞聽到上主不知和誰(「我們」)在討論事情,以賽亞就插口要「我們」差遣自己.這個「我們」是誰呢?根據古代近東的傳統,學者有理由相信這是上主在天庭與眾天軍天使開會的情景來的(就如中國宮庭片常看到的皇帝上朝會見眾卿),這些情景在約伯記頭一,二章就有出現過；列王記上22章亞哈詢問先知米該亞應否出戰亞蘭,米該亞就形容了一幕上主在天庭與天軍開會的境況(王上22:19-22米該雅說:「你要聽耶和華的話!我看見耶和華坐在寶座上,天上的萬軍侍立在他左右……」)；耶利米書也記載上主警告那些假先知說:有誰站在耶和華的會中得以聽見並會悟他的話呢?(耶23:18),就是他們沒有到上主的天庭會議開會,又怎會傳出上主真實的言語呢:「我沒有打發那些先知,他們竟自奔跑；我沒有對他們說話,他們竟自預言.他們若是站在我的會中,就必使我的百姓聽我的話,又使他們回頭離開惡道和他們所行的惡.」(耶23:21-22)所以其實以賽亞書6章記載的,實質是上主萬軍之耶和華與眾卿家商討派誰去向以色列宣告審判的信息,這誤闖入這個天庭會議,敬陪末席,旁聽會議的以賽亞,連議程是甚麼也不知道,就亂舉手說「我在這裡,請差遣我」!但當他聽清楚那個工作要求時(今日經課沒有使用的經文):「你去告訴這百姓說:『你們聽了又聽,卻不明白；看了又看,卻不曉得.』要使這百姓心蒙油脂,耳朵發沉,眼睛昏花；恐怕他們眼睛看見,耳朵聽見,心裡明白,回轉過來,就得醫治.」以賽亞就慌張起來,連忙問道:「主啊,這到幾時為止呢?」他得到的回應是:直到城邑荒涼,無人居住,房屋空閒無人,地土極其荒涼.並且耶和華將人遷到遠方,在這境內撇下的地土很多.境內剩下的人若還有十分之一,也必被吞滅,像栗樹,橡樹雖被砍伐,樹墩子卻仍存留.這聖潔的種類在國中也是如此.然候,就沒有了然候——以賽亞書6章就突然完在此處,沒有再記載以賽亞的回應.你說,這算那門子的呼召敘述?

承:以賽亞 vs. 彼得

但今日舊約經課和福音經課,卻有著一些特殊的連結.

你會發現,兩處經文都展現了神性的顯現:

在以賽亞書,以賽亞見到了上主的榮耀,並天使撒拉弗震撼的呼喊；

在路加福音,彼得見到了耶穌神聖不尋常的一面,能使他們兩首船都載滿了魚,甚至船要沉下去.

兩處經文,局中人都情緒波動:

彼得和伙伴們整夜勞力,筋疲力竭,而且失望地一無所得,但突然經歷神奇巨變,「他和一切同在的人都驚訝這一網所打的魚.」

以賽亞則在經歷了長治50多年的烏西雅王的駕崩,可能心裡愁悶,對未來國運疑惑,突在聖殿親見神聖榮耀,於是大呼「禍哉!我滅亡了!」

兩處經文,當局者都有些語無倫次:

以賽亞在毫無頭緒,一無所知下高呼「我在這裡,請差遣我」；

彼得突逢奇遇,卻在眾人都想留住耶穌之際,反說:「主啊,離開我」.

最主要的連結,就是兩處經文的當事人,都驚覺自己的罪與不配:

以賽亞眼見大君王萬軍之上主,就說自己是「嘴唇不潔的人,又住在嘴唇不潔的民中」；

西門彼得則俯伏在耶穌膝前,稱自己是個「罪人」.

每個人經歷上主,體會信仰的神聖,可能都各有不同:以賽亞可能作為宮庭先知憂心國事,退到聖殿沉思神子民的未來,就在退修的當刻經驗上主的榮耀；彼得和同伴卻在日常工作的期間,在尋常工作的場所遇到不尋常的基督.我們有時會在生活中感到迷失迷惘,有時會退下來思索,有時則在日常中更敏銳上主的聲音.但無論是何人,在甚麼境況,我們就是不能無動於衷,更大可能,是因為與主的相遇,我們更能看見自己,醒察自身,而體會自己的不配和不足.

轉:不在乎「呼」,而在乎「應」

今日無論是舊約的經課或是福音書的記載,細看並不似呼召的經文,但結局都是經歷有著神聖體驗的當事人,作出了回應.彼得和同伴最後「就撇下所有的,跟從了耶穌.」以賽亞也回應「我在這裡,請差遣我!」最後以賽亞書6章整章結束我們或許看不見以賽亞的回應,但我認為以賽亞書接著的經文,就是以賽亞的回應:他在國家危難的當前,規勸君王,發出先知的宣告,為神的子民帶來信靠的依據和盼望.

不知在座或網上參與崇拜的,有沒有人正積極尋求上主的心意,現在是神學院召生的季節,可能你也會問上主有否呼召我為上帝國擺上甚麼?上主有沒有呼召我呢?我如何得知上主的呼召呢?今日的經課其實提醒我們,跟從耶穌,作上主之僕,重點其實不在上主或耶穌的「呼」,而是在乎我們如何去「應」!就算天上有聲音叫你現在就放下一切,獻身傳道,報名神學院,終身事奉神,但你卻無動放衷,抽離無感,直行直過,那怎樣的呼召也是無用無效的.但在亳無準備之下,甚至灰心失意的時候,給你遇上了威榮的上帝,奇妙的主,就算沒有明確的呼召,你也願意回應這個值得你擺上一切,跟從到底的耶穌基督.那呼召不呼召,又是那麼重要嗎?

合:顯現期的呼召

不知大家有沒有留意,顯現後的這幾個主日經課,特別是舊約的選段,都和差遣相關的:顯現後第三主日耶穌宣讀的經文:主的靈在我身上,因為他用膏膏我,叫我傳福音給貧窮的人；差遣我報告:被擄的得釋放,瞎眼的得看見,叫那受壓制的得自由,報告神悅納人的禧年；上星期舊約經課,就是耶利米書1:4-10節:耶利米說,耶和華的話臨到我說:我未將你造在腹中,我已曉得你；你未出母胎,我已分別你為聖；我已派你作列國的先知.今日經課的以賽亞書6章和路加福音5章,都可以說圍繞上主的差遣和召命.顯現期是我們在主耶穌誕生之後,思想基督在世顯現他的工作的節期.但上主藉基督的顯現,又能否召請你去為上主擺上自己,在生活中見證榮耀他呢.上主的呼召,並不是只向要入神學院讀神學,將來出來在教會工作事奉的人的.上主的顯現和呼召,是對一切相信他,願意一生作他的們徒的.你又有何回應?



\newpage

\section{愛到底}
\label{sec:737}
\textbf{高國雄牧師 - 愛到底}
\newline
\newline
\hyperref[sec:736]{< < < PREV SERMON < < <}
~
\hyperref[sec:toc]{[返目錄]}
~
\hyperref[sec:738]{> > > NEXT SERMON > > >}


講題:愛到底 L-O-V-E

講員:高國雄牧師

場合:香港中文大學崇基學院禮拜堂主日崇拜

日期:2025年2月16日

福音書的場景,眾人對耶穌的期望

路加福音6:17-19

17耶穌和他們下了山,站在一塊平地上；同站的有許多門徒,又有許多百姓,從猶太全地和耶路撒冷,並泰爾,西頓的海邊來,都要聽他講道,又指望醫治他們的病； 18還有被污鬼纏磨的,也得了醫治. 19眾人都想要摸他；因為有能力從他身上發出來,醫好了他們.

聽道,得醫治(身心靈),得能力

與我們新年祈福很相似!

耶穌祝福他們什麼呢?

新年論「福氣」

生命中不同的境遇:

路加福音6:20-23

20耶穌舉目看著門徒,說:你們貧窮的人有福了!因為上帝的國是你們的.21你們飢餓的人有福了!因為你們將要飽足.你們哀哭的人有福了!因為你們將要喜笑.22「人為人子恨惡你們,拒絕你們,辱罵你們,棄掉你們的名,以為是惡,你們就有福了! 23當那日,你們要歡喜跳躍,因為你們在天上的賞賜是大的.他們的祖宗待先知也是這樣.

貧窮,飢餓,哀哭,被恨惡,被拒絕,被辱罵,名字被棄掉,被以為是惡.

貧窮:客觀的,相對的,社會性

「上帝的國是你們的」:屬於上帝的,上帝接納的,最豐富的,進入上帝的國是人生最重要的.

飢餓:實在的,當下的,最容易解決的,也是最危急的.

「將要飽足」:上帝所賜的是最合適,最對應當下的.

哀哭:當下的,悲憫的,心底裡的需要

「將要喜笑」:上帝要改變你,你的狀況,與及你的心

被恨惡,被拒絕,被辱罵,名字被棄掉,被以為是惡(被認為是沒有價值):為耶穌的原故.

在天上有大賞賜.

路加福音6:24-26

24但你們富足的人有禍了!因為你們受過你們的安慰.25你們飽足的人有禍了!因為你們將要飢餓.你們喜笑的人有禍了!因為你們將要哀慟哭泣.26「人都說你們好的時候,你們就有禍了!因為他們的祖宗待假先知也是這樣.」

新年論「禍患」

富足,飽足,喜笑

富足VS貧窮?

已經經歷一切你們所配得的安慰(原文新約:路6:24)

飽足VS飢餓?

將會飢餓

喜笑VS哀哭

將會悲傷和哭泣

「人都說你們好的時候,你們就有禍了!」

「人說你們好」是什麼時候呢?風光時?得意時?自我陶醉時?沽沽自喜時?…忘記上帝時?

可選擇VS不可選擇

個人情況VS社會眼光與待遇

耶穌實在是回應當下的需要,人就是來到主前,因為體會/知道自己有需要…

耶穌因為愛,揀選門徒,回應群眾…

耶穌是否就是單要滿足我們的需要?

使徒保羅道出了信仰的關鍵

哥林多前書 15:12-20

12既傳基督是從死裡復活了,怎麼在你們中間有人說沒有死人復活的事呢? 13若沒有死人復活的事,基督也就沒有復活了. 14若基督沒有復活,我們所傳的便是枉然,你們所信的也是枉然； 15並且明顯我們是為上帝妄作見證的,因我們見證上帝是叫基督復活了.若死人真不復活,上帝也就沒有叫基督復活了. 16因為死人若不復活,基督也就沒有復活了. 17基督若沒有復活,你們的信便是徒然,你們仍在罪裡. 18就是在基督裡睡了的人也滅亡了. 19我們若靠基督只在今生有指望,就算比眾人更可憐.20但基督已經從死裡復活,成為睡了之人初熟的果子.

當時在教會群體當中有人不信「死人會復活」

其實有這些人是不難理解的,我們今天在信仰追求的路上,也會遇到一些不明白,沒有經歷過,存疑的事情,但我們可以選擇用信心去接受,接受自己的有限；當然也有人選擇逃避,跳過,甚至對抗…

但保羅在此提到的關鍵內容,就是必須要接受有死人復活的事,才能信基督耶穌的復活,以致我們信主的人有一天也會一同復活…否則,一切都沒有價值與意義.

很多人來到教會只求得到福氣,覺得信仰導人向善,能(在道德上)教好自己/家人/兒女,做一個好人…

人很容易失去焦點

信仰的真正核心不是要得到當下的利益,世界的認同與稱讚,我們所追求的不是要成為世人眼中的好人.

而是以上帝的眼光去看世界,用上帝的愛的對待身邊人.

接受我們總覺得很難接受的事情.

剛剛過去的元宵節與及情人節,雖然商業成份味道比較重,但也有值得我們思考的地方.

情人節先讓我想起一些舊歌

L-O-V-E Love (溫拿樂隊1974,曲詞:黃霑)

L-O-V-E, Love La love, l-o-v-e 愛啦...

I will not give you riches. 我不會給你財富

Riches are not my wishes. 財富不是我的願望

I'll give you the vow of my life. 我會給你我一生的承諾

L-O-V-E, Love La love, l-o-v-e 愛啦...

L is for life and living. L 代表生活和活著

O is the oath of giving. O 是給予的誓言

V and E, vow of eternity. V 和 E 是永恆的誓約

現代愛情故事(1991)

別離沒有對錯 要走也解釋不多

現代說永遠已經很傻…

你我情如路半經過

深知道再愛痛苦必多

願你可輕輕鬆鬆放低我

剩了些開心的追憶送走我

皆因了解之後認清楚

離別時笑笑明晨剩我一個

瀟洒裡也會記起當初

若你的心中孤單再找我

若你的心窩中空虛再找我

不必痛苦當憶起我

情人節(Valentine's Day)來源的傳說

根據學術最不嚴謹的維基百科

來源1:「公元268年,羅馬皇帝克勞狄二世禁止年輕男子結婚,他認為未婚男子可以成為更優良的士兵.聖瓦倫丁違反了皇帝的命令,秘密為年輕男子主持婚禮,引起皇帝不滿,結果被收監,據說瓦倫丁於公元269年2月14日被處決.」為紀念聖瓦倫丁精神,人們將每年2月14日定為紀念日,成了後來的「情人節」.

來源2:瓦倫丁是早期基督徒.當時羅馬帝國是迫害基督徒的,因此在帝國境內信基督教意味著危險,甚至死刑.瓦倫丁為了掩護其他殉教者,暴露了自己的身分而被軟禁.在軟禁期間中,法官與他討論他的信仰時,討論到了耶穌的合法性.法官對瓦倫丁進行了考驗,並將法官收養的失明女兒帶到了他面前,法官說如果瓦倫丁成功地恢復了女兒的視力,他會幫瓦倫丁做任何事.於是瓦倫丁向神祈禱,把手放在她的雙眼上,女兒的視力就恢復了.法官信了耶穌,並透過禁食和祈禱,釋放了他管轄下的所有基督徒囚犯.

瓦倫丁後來因繼續傳福音而再次被捕,隨後他被送往羅馬皇帝那裡.瓦倫丁試圖說服皇帝接受基督教.皇帝拒絕接受基督教並判處瓦倫丁死刑,命令瓦倫丁放棄他的信仰.瓦倫丁拒絕了,並於269年2月14日被處決.傳說在瓦倫丁行刑的早晨,他給法官女兒寫了一封情意綿綿的告別信,信的落款是「來自你的瓦倫丁」.

原來情人節是選擇的故事

是為愛而選擇犧牲的故事

這正正同樣地就是耶穌基督的愛的故事

羅馬書5:7-8

7為義人死,是少有的,為仁人死,或者有敢作的.8 惟有基督在我們還作罪人的時候為我們死,上帝的愛就在此向我們顯明了.

基督在世時不是單回應當下群眾的需要,更讓人知道生命中更大的需要─我們需要生命,需要復活,永恆的生命.因此祂願意為人走上十架,為我們流出寶血,並且按照應許,耶穌基督從死人中復活.

這就是一個愛到底的故事!



\newpage

\section{斬斷延續怨恨的鎖鏈}
\label{sec:738}
\textbf{鄧瑞強博士 - 斬斷延續怨恨的鎖鏈}
\newline
\newline
\hyperref[sec:737]{< < < PREV SERMON < < <}
~
\hyperref[sec:toc]{[返目錄]}
~
\hyperref[sec:739]{> > > NEXT SERMON > > >}


約拿書 3:1-4:11
\newline
\begin{longtable}{cl}
\hline
\hline
章節 & 經文 (和合本修訂版)\\
\hline
3:1 & \begin{tabularx}{0.7\textwidth}{X} 耶和華的話第二次臨到約拿，說： \end{tabularx} \\ \\ \relax
3:2 & \begin{tabularx}{0.7\textwidth}{X} 「起來，到尼尼微大城去，把我告訴你的信息向其中的居民宣告。」 \end{tabularx} \\ \\ \relax
3:3 & \begin{tabularx}{0.7\textwidth}{X} 約拿就照耶和華的話起來，到尼尼微去。尼尼微是一座極大的城，約有三天的路程。 \end{tabularx} \\ \\ \relax
3:4 & \begin{tabularx}{0.7\textwidth}{X} 約拿進城，走了一天，宣告說：「再過四十天，尼尼微要傾覆了！」 \end{tabularx} \\ \\ \relax
3:5 & \begin{tabularx}{0.7\textwidth}{X} 尼尼微人就信服神，宣告禁食，從最大的到最小的都穿上麻衣。 \end{tabularx} \\ \\ \relax
3:6 & \begin{tabularx}{0.7\textwidth}{X} 這消息傳到尼尼微王那裡，他就從寶座起來，脫下朝服，披上麻布，坐在灰中。 \end{tabularx} \\ \\ \relax
3:7 & \begin{tabularx}{0.7\textwidth}{X} 他叫人通告尼尼微全城，說：「王和大臣有令，人、畜、牛、羊都不可嘗任何東西，不可吃，也不可喝水。 \end{tabularx} \\ \\ \relax
3:8 & \begin{tabularx}{0.7\textwidth}{X} 人與牲畜都要披上麻布，切切求告神。各人要回轉離開惡道，離棄自己掌中的殘暴。 \end{tabularx} \\ \\ \relax
3:9 & \begin{tabularx}{0.7\textwidth}{X} 誰知道神也許會回心轉意，不發烈怒，使我們不致滅亡。」 \end{tabularx} \\ \\ \relax
3:10 & \begin{tabularx}{0.7\textwidth}{X} 神察看他們的行為，見他們離開惡道，神就改變心意，原先所說要降與他們的災難，他不降了。 \end{tabularx} \\ \\ \relax
4:1 & \begin{tabularx}{0.7\textwidth}{X} 這事令約拿大大不悅，甚至發怒。 \end{tabularx} \\ \\ \relax
4:2 & \begin{tabularx}{0.7\textwidth}{X} 他就向耶和華禱告，說：「耶和華啊，這不就是我仍在本國的時候所說的嗎？我知道你是有恩惠，有憐憫的神，不輕易發怒，有豐盛的慈愛，並且會改變心意，不降那災難。我就是因為這樣，才急速逃往他施去的呀！ \end{tabularx} \\ \\ \relax
4:3 & \begin{tabularx}{0.7\textwidth}{X} 耶和華啊，現在求你取走我的性命吧！因為我死了比活著更好。」 \end{tabularx} \\ \\ \relax
4:4 & \begin{tabularx}{0.7\textwidth}{X} 耶和華說：「你這樣發怒，對嗎？」 \end{tabularx} \\ \\ \relax
4:5 & \begin{tabularx}{0.7\textwidth}{X} 約拿出城，坐在城的東邊，在那裡為自己搭了一座棚。他坐在棚子的蔭下，要看看城裡會發生甚麼事。 \end{tabularx} \\ \\ \relax
4:6 & \begin{tabularx}{0.7\textwidth}{X} 耶和華神安排了一棵蓖麻，使它生長高過約拿，影子遮蓋他的頭，使他免受苦難；約拿因這棵蓖麻大大歡喜。 \end{tabularx} \\ \\ \relax
4:7 & \begin{tabularx}{0.7\textwidth}{X} 次日黎明，神卻安排一條蟲來咬這蓖麻，以致枯乾。 \end{tabularx} \\ \\ \relax
4:8 & \begin{tabularx}{0.7\textwidth}{X} 太陽出來的時候，神安排炎熱的東風，太陽曝曬約拿的頭，使他發昏，他就為自己求死，說：「我死了比活著更好！」 \end{tabularx} \\ \\ \relax
4:9 & \begin{tabularx}{0.7\textwidth}{X} 神對約拿說：「你因這棵蓖麻這樣發怒，對嗎？」他說：「我發怒以至於死，都是對的！」 \end{tabularx} \\ \\ \relax
4:10 & \begin{tabularx}{0.7\textwidth}{X} 耶和華說：「這棵蓖麻你沒有為它操勞，也不是你使它長大的；它一夜生長，一夜枯死，你尚且愛惜； \end{tabularx} \\ \\ \relax
4:11 & \begin{tabularx}{0.7\textwidth}{X} 何況這尼尼微大城，其中不能分辨左右手的就有十二萬多人，還有許多牲畜，我豈能不愛惜呢？」 \end{tabularx} \\ \\
[1ex]
\hline
\hline
\end{longtable}
講題:斬斷延續怨恨的鎖鏈 Break the Chain of Resentment

經文:約拿書3:1-4:11

講員:鄧瑞強博士

場合:香港中文大學崇基學院禮拜堂主日崇拜

日期:2025年2月23日

各位弟兄姊妹,早安.願神的恩典常與你們同在.

人類社會難免衝突,某些人欺壓另一些人,被欺壓的人作出反抗.為了合理化欺壓,欺壓者會稱呼對方為惡棍；為了合理化反抗,被欺壓者會稱呼對方為邪魔.一旦污名化已成形,雙方各自對對方懷有成見,演變成正邪不兩立,則很難復和.

我們今日講約拿書,內容涉及以色列與亞述的世代恩怨.這恩怨如何能化解呢?約拿書寫成的年代,應是猶太人被擄歸回的年代.你們聽我講舊約,常常提及猶太人被擄歸回的事件.這是他們民族歷史的分水嶺.他們國家被強國消滅後,人民被擄他方,猶太人生活的世界,就完全不能回到從前了.以前,有自己的國家,有自己的民族身分,有正常的宗教生活；國破後,成了別人的奴隸,沒有了國家,沒有了身分,沒有了安身立命的宗教.昔日的社會,已一去不返了.在國際舞台上,霸主興起又沒落,到了波斯稱霸的年代,猶太人竟能從被擄之地回歸故土,重建家園.經過這些風雨之後,倖存的猶太人,要重新認識自己到底是誰,也要重新確立自己存活的意義.我們相信,約拿書便是在這處境下寫成的.

翻查猶太人的歷史,在列王記下14:25裡,提及一個先知,正是亞米太的兒子約拿,這個約拿是北國以色列耶羅波安二世時的先知.他顯然是一個逢迎君王的先知,他沒有像同期的先知阿摩司那樣,批評君王的窮奢極侈,魚肉人民.以色列國卻在這些歌功頌德的啦啦隊掌聲中,日漸消亡.以色列國最後被亞述所滅.順帶一提,猶大國後來被巴比倫所滅.

斗轉星移,這位約拿先知二百多年之後,波斯已取代了亞述和巴比倫,成了當代霸主.猶太人可以從被擄之地歸回,回到昔日家園.唉,一片荒涼景象,正是無處話淒涼.

搞到如此田地,能怪誰呢?有恨嗎?恨,總是有的.恨誰?恨亞述,恨巴比倫,也應該恨自己吧.如何看這些統治自己的霸權?是敵人?但,不也正是新興起的霸主波斯,讓自己從被擄之地歸回?能愛他們嗎?面對這些外來的強大對手,又恨又愛,又愛又恨.面對他們,我是誰?我們和他們日後如何相處?上帝的心意又是如何?這一切,都需要重新思考.

有神學家寫下約拿書,刺激他的同胞對以上問題作一點信仰思考.約拿書是聖經中最有卡通feel,最有繪本feel的一卷書.一條大魚吞了一個人,一個人在大魚肚裡感恩地祈禱,在繪本裡,這些情節是完全可以想像的.我們習慣了將聖經看為是千篇一律的,嚴肅的事實命題大全,但其實聖經的文體是很多元化的.類乎繪本的文體,也可以用來啟發信仰.或者,一般人看繪本,比看講道理的論文,更能得到信仰的啟迪.

用誰來當繪本的主角呢?作者用了古時那位先知約拿來當主角.這個古時的約拿,是以色列國一個壞君王的啦啦隊隊長,很有民族感情,但不知他是否夠長命,看到以色列國被亞述吞滅.這有點像:在足球場上,這啦啦隊隊長大聲喊著:「我們醒,我們勁」,但轉頭,便被對方「勁炒十蛋」.場面是有點諷刺,有點滑稽的.這個人物,便成了約拿書「戲仿/惡搞」的對象,他帶給讀者一點歡樂.或者,沉重的國仇家恨,只能在神聖的歡樂中才能化解.

我們看看這本最有繪本feel的書.

拿1:1 耶和華的話臨到亞米太的兒子約拿,說:

拿1:2 「你起來往尼尼微大城去,向其中的居民呼喊,因為他們的惡達到我面前.」

拿1:3 約拿卻起來,逃往他施去躲避耶和華；下到約帕,遇見一隻船,要往他施去.他就給了船價,上了船,要與船上的人同往他施去躲避耶和華.

典型的先知故事,開頭通常就是一個蒙召的經歷.題外話,整本約拿書,從來沒有稱呼約拿為先知.這本書看似是先知書,但「惡搞」的成分居多.在這開場白裡,耶和華叫約拿去亞述首都尼尼微,宣講審判的信息.若果我們熟讀傳統的先知書,我們會預期:蒙召的人會自覺不配,推說自己擔當不了這麼大的任務.如:先知以賽亞會說:「禍哉!我滅亡了!因為我是嘴唇不潔的人,又住在嘴唇不潔的民中.」(賽6:5)又或者像先知耶利米說:「主耶和華啊,我不知怎樣說,因為我是年幼的.」(耶1:6)有了這些慣性預期,便看到約拿的反應多麼出人意外.神叫他去東面的尼尼微,他就走去西面的他施.他施在西班牙的南端,是地中海極西之處,當時人認為這就是世界的西端.

情節假定我們知道,約拿是個民族主義者,他恨透滅了自己國家的亞述,他想亞述被神消滅,而不是獲救.經文仔細地說,「他給了船價」,這個人付錢去逃避一份工.這有點似:老闆叫他去傾生意,他就丟下錢,走了去游水.約拿像其他先知那樣,自覺卑微,不勝重負嗎?沒有,他自覺必會成功,但他就是不希望自己工作成功,他「劈炮唔撈」.約拿的反應,顛覆了我們對蒙召者的想像.

跟著的情節,大家可能知道.約拿上了船,神便興風作浪.水手們各自求神拜佛,約拿卻在船艙裡睡覺.在聖經裡,除了耶穌,只有約拿在暴風下的船上睡覺.水手們掣籤,看看這風雨因誰而起.啊,原來是約拿.水手們問他:你是誰?你做了什麼?約拿的答案耐人尋味,他說:「我是希伯來人.我敬畏耶和華－那創造滄海旱地之天上的上帝.」(拿1:9)你們說,這約拿是敬畏耶和華的人嗎?老闆叫你工作,你顧住打機.同事問你,「你搞邊科?」你答:我是敬畏老闆的人.這應引來一場大笑.

跟著,水手們將約拿拋到海裡,風浪止息了.神安排一條大魚吞了約拿,他在魚腹裡三日三夜.

來到這裡,若有弟兄姊妹堅持約拿書講的,是確實的神蹟事件,不是什麼像繪本般的故事,當然是可以的.約拿書的文學性質是什麼,是有討論餘地的.

約拿在魚肚裡,作了什麼?他作了一感恩的祈禱,他感謝神救他出深淵.一個人在魚肚裡,好平靜地祈禱,這畫面好有繪本feel.最形象的描繪,是「海草纏繞我的頭」(拿2:5).這像是一個人正跌落懸崖,卻說,我見到一隻紅蜻蜓飛過,很超現實,也有點搞笑.面對如此大的懲罰,約拿半句悔改的話也沒有,他看來像大無畏的大英雄.接著,「耶和華吩咐魚,魚就把約拿吐在旱地上.」(拿2:10)

拿3:1 耶和華的話二次臨到約拿說:

拿3:2 「你起來!往尼尼微大城去,向其中的居民宣告我所吩咐你的話.」

拿3:3 約拿便照耶和華的話起來,往尼尼微去.這尼尼微是極大的城,有三日的路程.

拿3:4 約拿進城走了一日,宣告說:「再等四十日,尼尼微必傾覆了!」

耶和華再叫約拿去尼尼微,宣講審判的信息.貫穿這城,需要三日的路程,但約拿進城只走了一日,便宣告信息了.約拿是否很「hea」?校長叫你巡三間班房,你巡了一間,便打報告,這是否很「hea」?約拿的信息,原文只有五個字,「再等」「四十」「日」「尼尼微」「必傾覆了」.這是聖經裡最短的信息之一.約拿沒有針對任何人發言,沒有指出尼尼微人的任何罪行,也沒有要求任何具體的改變,他看來只是很隨意的,講了五個字,便收工了.

在正常情況下,先知講的信息,都是沒有人聽的.摩西向法老宣告神的要求,法老當耳邊風.耶利米叫西底家王投降,西底家也沒聽從.亞述是窮凶極惡之人,而約拿又毫無激情,講得極「hea」,我們想,必定沒有什麼效果喇.區區五個字,能帶來什麼轉變呢?這次,又出乎我們的想像.

拿3:5 尼尼微人信服上帝,便宣告禁食,從最大的到至小的都穿麻衣〔或譯:披上麻布〕.

拿3:6 這信息傳到尼尼微王的耳中,他就下了寶座,脫下朝服,披上麻布,坐在灰中.

拿3:7 他又使人遍告尼尼微通城,說:「王和大臣有令,人不可嘗甚麼,牲畜,牛羊不可吃草,也不可喝水.

拿3:8 人與牲畜都當披上麻布；人要切切求告上帝.各人回頭離開所行的惡道,丟棄手中的強暴.

拿3:9 或者上帝轉意後悔,不發烈怒,使我們不致滅亡,也未可知.」

完全不可思議,由下到上,全民悔改.他們穿上麻衣,全面禁食及禁水.從來沒有先知,要求人禁食時同時禁飲水的.這次更奇怪,連動物也披麻,禁食禁水.想像一下,若他們有貓狗,那些貓狗也穿上麻衣.更難想像,動物如何禁食禁水?這些畫面,若你不是古板的人,應覺得好笑.

在聖經中,有眾多先知,唯有這位最不順服的約拿,他的hea講獲得最全面的順服,獲得最大的回響.

拿3:10 於是上帝察看他們的行為,見他們離開惡道,他就後悔,不把所說的災禍降與他們了.

約拿最不想見到的局面,出現了,神寬恕尼尼微.

約拿像小孩子一樣,發脾氣.

拿4:1 這事約拿大大不悅,且甚發怒,

拿4:2 就禱告耶和華說:「耶和華啊,我在本國的時候豈不是這樣說嗎?我知道你是有恩典,有憐憫的上帝,不輕易發怒,有豐盛的慈愛,並且後悔不降所說的災,所以我急速逃往他施去.

拿4:3 耶和華啊,現在求你取我的命吧!因為我死了比活著還好.」

在古代,先知以利亞曾經在神面前求過死.不過,那次,以利亞講的話,無人聽,還被人追殺.今次,約拿講的,全部人都聽,效果好得不得了.那次,是考試零分,故求死.今次,是考試一百分,但不喜歡那獎品,又求死.零分,求死；一百分,又求死.可以肯定的是:上帝會被激死.

最有趣的,是約拿投訴的理由.我意譯一下4:2:「耶和華啊,我『早就話jor 你喇』,你耳朵軟,心腸軟,容易變,沒有一點堅持,所以我不理你.」言下之意,是神太慈愛,祂不應如此慈愛.

拿4:4 耶和華說:「你這樣發怒合乎理嗎?」

這個「理」字,不單只有理性上「合理」的意思,更似中文「天理」的「理」,有道德含意,美好人格的含意,令人愉悅的含意.神問約拿的話,用廣東話來表達:「你這樣發怒,arm 唔arm 呢?」這個「arm 」字,包含理性上的對,道德上的善,情境上的恰當.

拿4:5 於是約拿出城,坐在城的東邊,在那裡為自己搭了一座棚,坐在棚的蔭下,要看看那城究竟如何.

再次,約拿像小孩子一樣,「粒聲唔出」,看事情發展.難道,他期望上帝改變心意,消滅尼尼微?

拿4:6 耶和華上帝安排一棵蓖麻,使其發生高過約拿,影兒遮蓋他的頭,救他脫離苦楚；約拿因這棵蓖麻大大喜樂.

神令一棵蓖麻長大,為約拿遮陰,約拿因這棵蓖麻「大大」喜樂.「大大」這字的原文,就是形容尼尼微是「大城」的那個「大」字.這喜樂真大.約拿因這棵小小的蓖麻,就能有大大的喜樂,真是很易「哄」.

拿4:7 次日黎明,上帝卻安排一條蟲子咬這蓖麻,以致枯槁.

拿4:8 日頭出來的時候,上帝安排炎熱的東風,日頭曝曬約拿的頭,使他發昏,他就為自己求死,說:「我死了比活著還好!」

上帝安排一條蟲咬那蓖麻,那蓖麻便死了.又是一很有繪本feel的畫面.這麼小事,約拿又大發脾氣,又想死.

拿4:9 上帝對約拿說:「你因這棵蓖麻發怒合乎理嗎?」他說:「我發怒以至於死,都合乎理!」

拿4:10 耶和華說:「這蓖麻不是你栽種的,也不是你培養的；一夜發生,一夜乾死,你尚且愛惜；

拿4:11 何況這尼尼微大城,其中不能分辨左手右手的有十二萬多人,並有許多牲畜,我豈能不愛惜呢?」

神問約拿,你因這棵蓖麻發怒,「arm 唔arm 呀?」約拿說:「我arm 晒.」神沒有責罵約拿,反而像老爺爺「哄」小孩子般「哄」他.那蓖麻不是你種的,你尚且愛惜它.何況這許多人,他們像無知的小孩般,分不清是非黑白,也還有許多牲畜,我們怎能不愛惜呢?

我們很多人像約拿,很愛惜身邊的小盆栽,貓貓狗狗,但對心中的敵人,卻要除之而後快.神也很難和我們講道理,對於消滅敵人一事,我們認為沒有道理可講.神讓我們坐下來,一起看繪本.有一個不聽話的小孩,爸爸叫他去東,他就去西,因為他不想去救他的敵人,他想他們死.然後,大魚吞了他......

在這個故事裡,沒有壞人,只有蠢人.有叫動物著麻服,要牠們禁水禁食的尼尼微人.也有約拿,一棵小植物就叫他開心得手舞足蹈,也因這小植物的枯死而要生要死.其實,壞人不太壞,但蠢人就真的很蠢,他們誓要對方死.對方不死,惱得自己想死.但,有需要如此嗎?這樣「arm 唔arm 呢」?這繪本的劇情,基本上,是大團圓結局的.約拿的工作,超額完成.邪惡的尼尼微人,原來並不太壞,尼尼微至終沒有遭天譴.神愛約拿,也愛尼尼微城的人.在人間,人人都可以是天父的兒女.人能從這繪本,得到應得的智慧嗎?這繪本的結尾,最後是開放式的,人會忽然再想消滅對方,抑或,保留大團圓到結尾?

願神祝福你們.



\newpage

\section{主的榮光}
\label{sec:739}
\textbf{李均熊牧師 - 主的榮光}
\newline
\newline
\hyperref[sec:738]{< < < PREV SERMON < < <}
~
\hyperref[sec:toc]{[返目錄]}
~
\hyperref[sec:740]{> > > NEXT SERMON > > >}


講題:主的榮光 The Glory of the Lord

講員:李均熊牧師

場合:香港中文大學崇基學院禮拜堂主日崇拜

日期:2025年3月2日

起:彼得和同伴睡著了?

當我預備今天經課講道時,我發現原來我同樣在上年的「顯現後最後主日——耶穌登山變像日」(2024年2月11日)講道,講道題目為〈以利亞已經來了〉.那次講道我將焦點放在馬可和馬太福音記載門徒在耶穌變像顯現之後問耶穌「文士為甚麼說以利亞必須先來?」,然候講到猶太人的盼望,並結論說我們應在進入預備預苦期前好好守護盼望.但今天福音經課選取的路加福音卻沒有這一段.這可能因為路加福音的宣講對像主要是操希臘語的猶太人,或甚至是活在希臘文化的外邦人,所以傳統猶太人所關心的「文士為甚麼說以利亞必須先來?」這個問題就從略.

反之,路加以下的節情只有路加福音記載:

9:32 彼得和他的同伴都打盹,既清醒了,就看見耶穌的榮光,並同他站著的那兩個人

為甚麼獨有路加記載彼得和門徒睡著了的細節,而其他對觀的馬太,馬可卻沒有記載?

首先,路加福音說彼得和他的同伴都「打盹」,即代表他們是怎樣的狀況呢?「打盹」即是「很想睡但還未睡著」,而NIV譯為very sleepy,也可能給人一個感覺是彼得和同伴只是很想睡覺,只是仍苦苦撐著,不讓自己入睡.但轉為直譯的KJV卻譯為heavy with sleep；NASB呢則譯作overcome with sleep；新近一些英譯本也多譯為in a deep sleep(Holman Christian Standard Bible), sound asleep (Contemporary English Version),黃根春老師編的《原文新約》也譯作「沉睡了」……若真要硬譯希臘文的意思,或許可以說彼得和同伴「給睡眠負累」(to be burdened by sleep).那究竟是他們很想訓覺,還是訓得很lum ,兩者都是解得通的.不過新近期的字典都比較傾向他們是睡著了(to be in a sstat of deep sleep),起碼肯定他們對身邊的事情是迷迷糊糊,否則也不用再加一句「既清醒了」,字典的解釋是「從睡眠中完全醒過來」(to awaken completely from sleep)之意.

大家細心想想,其實加入這個門徒睡著了的情節(無論是迷迷糊糊的打盹還是真的沉沉睡著了),不是會令讀者產生一個疑問, 耶穌一路祈禱,一路面容改變,身體發光,又和兩位同樣在榮光裡顯現的人物交談(我也時常被人問到,他們怎知那兩人一個是摩西,一個是以利亞呢?),究竟門徒是否親眼目睹? 路加福音這樣寫,不是使彼得和門徒成了不可靠的沉睡證人?!

承:彼得和門徒的「懵懂」

路加福音其實和其他福音書一樣,以彼得為我們作門徒的代表,去顯出我們對耶穌作為基督的身份是如何充滿誤解,又是多麼糊塗muddled多麼懵懂befuddled.馬太,馬可和路加一樣,在記載耶穌帶門徒上山祈禱變像前,都記載耶穌問門徒群眾認為自己是誰.他們答了一般人的見解:有人說是已被希律斬首的施洗約翰；有人說是以利亞或古時的一個先知又活了.但耶穌續問門徒認為自己是誰,彼得就回答說:是神所立的基督.之後,耶穌便明明的告訴門徒,人子必須受許多的苦,被長老,祭司長,和文士棄絕,並且被殺,第三日復活.但路加福音卻沒有如馬太和馬可一樣,這兩卷福音書均記載彼得拉著耶穌,又勸他,叫他萬不可如此,而耶穌說撒旦,退我後邊去(這句話相信大家應很熟識).

在馬太和馬可福音,當耶穌談到他的受死和受苦,彼得就已經跳出來說耶穌萬不可如此.顯出彼得對耶穌作為基督的身份只是一知半解.但路加卻沒有加插這個彼得對主受苦受死的回應,並不是路加認為彼得完全明瞭耶穌作為基督的意義.更可能是,路加要顯出彼得連耶穌講甚麼,可能他也聽不進去,故此路加要將彼得(和門徒)的「糊塗」muddled「懵懂」befuddled,放在耶穌登山顯榮光的時候,才將之表露出來.

轉:要留住主的榮光因為不懂主的身份

很多解經家都留意到,路加福音的整段登山變像經文其實可分為兩個段落,並由彼得與門徒由沉睡到清醒為轉接點,可見路加福音的作者是刻意加入這個情節.上半段關於耶穌登山變像,並與摩西和以利亞兩位偉大先知代表同顯榮光的傳統,相信是初代信徒之間流傳甚廣,又十分熟識的耶穌事蹟.正是三卷對觀福音的資料來源.馬太和馬可都表明耶穌是在三位門徒面前變了形像,所以彼得,雅各,約翰是以目擊証人的角色見證主的榮光.主基督在榮耀中顯現,在路加福音中彼得和門徒不單沒有親眼去見證,反而落在沉睡中.

當他們「從睡眠中完全醒過來」真正看到主的榮光,已經是二位帶著榮光的人物要和耶穌分離的時候.那時,彼得對耶穌說,要搭三座棚,把耶穌,摩西和以利亞留在這個山上,留住這個時刻.當然路加不忘加一句「他卻不知道所說的是甚麼」.當摩西和以利亞與耶穌將要在耶路撒冷將要成的事,彼得卻開倒車留戀住棚節以搭棚來記念昔日上主在曠野的同在和供應?就在彼得說出這些完全誤解耶穌作為基督的說話,就有一朵雲(和合本譯作雲彩常害我想像登山變像是在大白天發生……)遮蓋他們.那是遮蓋主耶穌的榮光,卻同時是神顯現的密雲,使人心生懼怕.而從這神顯現的密雲中,不是眼見的榮光,而是神的聲音:這是我的兒子,我所揀選的,你們要聽他.聲音住了,就只見耶穌一人在那裡.沒有了榮光,沒有了古時的偉人,就只有平凡但我們卻要聽他話的耶穌.

合:為預苦期預備

路加福音要展現的,可能是另一種我們對耶穌作為基督身份的誤解,就是期望基督的榮光可以留在世上,甚至完全忘記耶穌就是那位受苦的基督.有時信徒會否太過渴求要見到基督的榮耀,追求特殊的宗教經,而忽略了作門徒的代價?路加福音的作者卻透過彼得和同伴的沉睡,解構了這一種渴望.因為當你要找緊這一個主的榮光時,你卻可能會忘記了主耶穌的說話.

在路加記載的耶穌登山變像的尾聲,是密雲中有聲音說:你們要聽他.但耶穌在整段記載中,都沒有說過一句話.那我們要聽他甚麼呢,不就是上山八日前,耶穌說過的話,就是要天天背起十字架來跟從主.

23 耶穌又對眾人說:「若有人要跟從我,就當捨己,天天背起他的十字架來跟從我.24因為,凡要救自己生命的,必喪掉生命；凡為我喪掉生命的,必救了生命.25人若賺得全世界,卻喪了自己,賠上自己,有甚麼益處呢?」

主的榮光可以在我們預苦期中思想基督的受苦時仍不失盼望.但主的榮光又會否成為我們留戀的對像,以為有主在我們就不需要付上作門徒的代價?當主耶穌的榮光被神的密雲遮蔽,其實並不是取去我們盼望的動力,而是要我們更專注去聽從主的聲音而已.



\newpage

\section{勝過人生三大試探}
\label{sec:740}
\textbf{高國雄牧師 - 勝過人生三大試探}
\newline
\newline
\hyperref[sec:739]{< < < PREV SERMON < < <}
~
\hyperref[sec:toc]{[返目錄]}
~
\hyperref[sec:741]{> > > NEXT SERMON > > >}


講題:勝過人生三大試探 Overcome the Three Major Temptations in Life

講員:高國雄牧師

場合:香港中文大學崇基學院禮拜堂主日崇拜

日期:2025年3月9日

大齋(Lent)一詞源自古英語"Lencten",意指「春天」.由聖灰日(Ash Wednesday即大齋首日)開始.

整個大齋期是四十天(不包括主日).因為每個主日都視作一個「小救主復活日」.而在救主復活日前主日稱為「棕樹主日」,又是「聖週」(Holy Week)的開始.

主後一百年左右的基督教文獻中已有提及有關大齋期的活動.早期教會很多準備在復活日受洗的人都會進行禁食祈禱來作好準備受洗之心.而一般信眾多會透過禁食,自我反省與懺悔認罪,再一次把自己獻上當作活祭,與上帝復和.

大齋期所用的節期顏色是紫色,象徵懺悔認罪.有些教會在此期間不擺設鮮花,甚或不舉行婚禮或其他慶典.

大齋期的目的是更新和強化基督徒的屬靈生命.

禁食,讀上帝的話語——聖經或屬靈書藉,行動.

目的是給予我們有機會與上帝復和,與自己復和,又能夠與別人復和.

路加福音 4:1-13

1耶穌被聖靈充滿,從約旦河回來,聖靈將他引到曠野,2四十天受魔鬼的試探.那些日子沒有吃甚麼；日子滿了,他就餓了.3魔鬼對他說:「你若是上帝的兒子,可以吩咐這塊石頭變成食物.」4耶穌回答說:「經上記著說:『人活著不是單靠食物,乃是靠上帝口裡所出的一切話.』」5魔鬼又領他上了高山,霎時間把天下的萬國都指給他看,6對他說:「這一切權柄,榮華,我都要給你,因為這原是交付我的,我願意給誰就給誰.7你若在我面前下拜,這都要歸你.」8耶穌說:「經上記著說:當拜主－－你的上帝,單要事奉他.」9魔鬼又領他到耶路撒冷去,叫他站在殿頂上,對他說:「你若是上帝的兒子,可以從這裡跳下去； 10因為經上記著說:主要為你吩咐他的使者保護你；11他們要用手托著你,免得你的腳碰在石頭上.」12耶穌對他說:「經上說:『不可試探主－－你的上帝.』」13魔鬼用完了各樣的試探,就暫時離開耶穌.

前言

試探:魔鬼引誘?不應該出現?壞事?

過去一直的教導:「我們要學習如果避開試探,如何勝過試探」,「靠著聖經就可以勝過試探,所以要多讀聖經」…

到底是否真的如此呢?

V1耶穌被聖靈充滿,從約旦河回來,聖靈將他引到曠野,2四十天受魔鬼的試探.

試探因聖靈充滿,被聖靈引導

V9魔鬼又領他到耶路撒冷去,叫他站在殿頂上,對他說:「你若是上帝的兒子,可以從這裡跳下去； 10因為經上記著說:主要為你吩咐他的使者保護你；11他們要用手托著你,免得你的腳碰在石頭上.」

魔鬼引用詩篇91:11-12

魔鬼也用聖經

物資的試探

V3魔鬼對他說:「你若是上帝的兒子,可以吩咐這塊石頭變成食物.」

食物的試探=貧吃(七宗罪之一)的提醒?食得太豐富?有信徒因此覺得食得太好是罪過…這段聖經真的是如此教導嗎?

食物是身體的需要,最基本的慾望…

V1聖靈將他引到曠野,2四十天受魔鬼的試探.那些日子沒有吃甚麼；

耶穌在曠野禁食四十天,請留意:聖經沒有說耶穌是禁食禱告,只是禁食,這可能與曠野(出埃及時)四十年呼應的體驗,請留意:在曠野四十年的以色列人也有上帝所賜的嗎哪與鵪鶉作食物,上帝沒有忽視人需要食物這實況,向上帝求食物,在主禱文內也有教導(我們日用的飲食,求父今日賜給我們)…

所以耶穌當下的需要期望被滿足是合情合理的…

V4耶穌回答說:「經上記著說:『人活著不是單靠食物,乃是靠上帝口裡所出的一切話.』」

申命記8:1-20

申命記8:3他苦煉你,任你飢餓,將你和你列祖所不認識的嗎哪賜給你吃,使你知道,人活著不是單靠食物,乃是靠耶和華口裡所出的一切話.

真正的活著,不單是肉身維持

在滿足身體需要,生存要求之上,有更重要的價值

人之所以為人,不可單為了生存,而是尋找到生活的意義

焦點的試探:不單靠=需要靠,但不可是生命中的焦點

權力名利的試探

V5魔鬼又領他上了高山,霎時間把天下的萬國都指給他看,6對他說:「這一切權柄,榮華,我都要給你,因為這原是交付我的,我願意給誰就給誰.7你若在我面前下拜,這都要歸你.」

追求權柄,榮華是很多人人生的目標

一步到位更為人所羨慕

聖經沒有表示這追求是錯誤的,只是我們要反思追求的方法及途徑…

其實這是對耶穌主權的試探,魔鬼說:「因為這原是交付我的」…真的嗎?眼前的認知VS事實的真相

主禱文的教導:因為國度,權柄,榮耀…全是父的

V8耶穌說:「經上記著說:當拜主－－你的上帝,單要事奉他.」

申命記6:13你要敬畏耶和華－你的上帝,事奉他,指著他的名起誓.

申命記10:20你要敬畏耶和華－你的上帝,事奉他,專靠他,也要指著他的名起誓.

耶穌不是說權柄與榮耀不重要,不需要追求,而是表明真正的擁有者不是魔鬼,而是上帝,更知道祂將要得到的方法不是向魔鬼下拜(承認魔鬼的權柄),而是不完成上帝的心意,走上十字架.

濫用的試探,驕傲的試探

V9魔鬼又領他到耶路撒冷去,叫他站在殿頂上,對他說:「你若是上帝的兒子,可以從這裡跳下去； 10因為經上記著說:主要為你吩咐他的使者保護你；11他們要用手托著你,免得你的腳碰在石頭上.」

「有上帝的看顧保守,就可以從這裡跳下去」真的可以嗎?

上帝的看顧不等如我們要刻意製造險境

例子:有了藥水膠布就可以不斷弄傷自己?國內某燙火膏的示範:為表明藥膏有效,把雙手放在燒紅的鐵鏈後,馬上使用該藥膏以示其神奇效果…

人往往希望用刻意製造意外/險境的方式後,得到幫助及拯救,以「證明」上帝的真確,再讓自己放心地「信」

這其實才真正在試探上帝中

V12耶穌對他說:「經上說:『不可試探主－－你的上帝.』」

申命記6:16「你們不可試探耶和華－你們的上帝,像你們在瑪撒那樣試探他.

V13魔鬼用完了各樣的試探,就暫時離開耶穌.

聖經記載的三個試探,原來只是「各樣」試探的代表

三個試探都有的共通點

V3你若是上帝的兒子

V7你若在我面前下拜

V9你若是上帝的兒子

其實整段經文的重點,就是對耶穌身份的試探

耶穌受洗,聖靈充滿,在傳道之先,上帝透過這四十天,確認了耶穌基督的身份,祂就是上帝的兒子,完全的承繼,但不是以滿足人生存的需要為任務,不是以世人所期望的方法與途徑去獲得權柄榮耀,也不是來「表演」上帝的拯救,而是要以「基督」的身份,將會走上十架,成為為人犧牲生命的救主,以完成救恩,這才是耶穌真正要確立的使命.



\newpage

\section{今天明天後天我必須向前走}
\label{sec:741}
\textbf{林豪恩先生 - 今天明天後天我必須向前走}
\newline
\newline
\hyperref[sec:740]{< < < PREV SERMON < < <}
~
\hyperref[sec:toc]{[返目錄]}
~
\hyperref[sec:742]{> > > NEXT SERMON > > >}


講題:今天明天後天我必須向前走 Yet Today, Tomorrow, And The Next Day I Must Be On My Way

講員:林豪恩先生

場合:香港中文大學崇基學院禮拜堂主日崇拜

日期:2025年3月16日

一．遇到阻礙,向前還是退後?

有沒有試過被人阻礙你?例如阻你路?阻礙你做想做的事情?被人「阻頭阻勢」的確是令人感到煩擾.2013年11月7日,有工人到尖沙咀柯士甸路一大廈送貨,期間疑因貨物擺放問題與保安員發生爭執,工人指責保安「阻頭阻勢」,盛怒之下向保安員拳打腳踢.(東網)

今天的福音經課都好似有一件「阻人前進」的事.

路加福音13:31:正當那時,有幾個法利賽人來對耶穌說:「離開這裡去吧,因為希律想要殺你.」

「正當那時」,是甚麼時候呢?「離開這裡」,即是哪裡呢?上文有資料提供:13:22耶穌往耶路撒冷去

原來當時耶穌正要往耶路撒冷去,在路上,有「消息人士」來勸退他,「我們收到消息,希律要lor 你命!你都係返鄉下啦,耶路撒冷好危險嫁!」於勸退者來說,可能自覺是善意,「都係為你好即」；但在被勸退者來說,可能感到被威嚇逼使自己放棄想做的事情.遇到這些阻礙,人們通常有甚麼反應的呢?

其實,每個人都希望以自己的方法做到自己想做的事情.心想事成就最暢快的了.如果被人「阻頭阻勢」,難免感到挫折及憤怒.做不到自己想做的事情的挫折及憤怒,連小孩子也感覺得到.

有一個兩,三歲的小孩拿起玩具扔向麻麻,同時大哭大鬧.

媽媽:「BB,你為甚麼向麻麻扔玩具呀?」

小孩:「我要吃糖,麻麻不讓我吃!」

媽媽:「原來你想吃糖吃不到.」

女兒:「麻麻不給我.」

媽媽:「你想吃糖吃不到,所以你感到憤怒?」

女兒:「麻麻不給我糖吃.」

如果這情況出現在你家中,劇情將會如何發展呢?會是動作片?恐怖片?喜劇?還是教育電視?那就視乎不同家庭的狀況了.無論如何,這位小孩的脾氣反映出我們自小就對於「阻頭阻勢」反應很大.不能夠做到自己想做的事情的狀況並沒有隨著長大而減少,反而可能越來越多.除了自己的能力不足之外,更常常感到其他人及環境的阻力,令人不能「想做就去做」.因此,這種挫敗及憤怒,可能隨著年日的增多而增加.從前有一位歌星叫譚詠麟,他有一首歌叫「傲骨」,正唱出這種鬱結及憤怒.

誰知我內心多苦悶 一切沒法如願

誰知我內心多苦悶 因我甚覺疲倦

但到底 這鬥爭沒有完 想逼我改變

有誰會知我心內 有極強信念

從開始至今多考驗 使我步法常亂

從開始至今多考驗 手裡利劍常斷

話到底 我了解好清楚 不肯趁風轉

有時我想去屈服 想起也心酸

遇到阻礙,應該迎難而上,還是回頭是岸呢?這真是千古難題.就算用信仰的角度去問屬靈長輩,得到的回應也可以是兩極:

牧師,我喜歡的女孩子對我「已讀不回」,我應該如何理解上帝的心意呢?

牧師A可能回答:事情既然不順,可能是上帝提示你不要再糾纏下去.

牧師B可能回答:事情不順可能是上帝對你的考驗,不應放棄追求.

回到福音經課,法利賽人警告耶穌「離開這裡去吧,因為希律想要殺你.」耶穌如何回應呢?

二．為別人的福祉而前行

13:32耶穌說:「你們去告訴那個狐狸說:『今天,明天我趕鬼治病,第三天我的事就成全了.』

耶穌的回應爽快直接,不向恐嚇屈服,堅持走自己的路.如果你剛好有些事情猶豫不決,聽完耶穌的回應,你會不會因此而被激勵勇往直前呢?在採取行動之前,不妨思考一下我們能否如此簡單直接地應用聖經.雖然堅持是一種社會讚許的態度,然而,我們會否也需要考慮一下堅持做的是甚麼事情呢?有沒有聽過這句激勵人心的說話:「不是因為看見希望而堅持,而是因為堅持才看見希望」.這句說話雖然勵志,但其實沒有內容.堅持要做的是甚麼事情呢?這種事情是否值得堅持呢?會不會我看自己為堅持,其他人看我為固執呢?

在婚前輔導中,我邀請兩位準婚人士分享對方的優點.其中一位說:「我欣賞你的堅持,你想做的事情,總是不怕艱辛,排除萬難都要做.例如...」另一位一面聽,一面流眼淚.我問她有甚麼感受,她說:「聽到他說我堅持,我很感動,因為只有他說我堅持,其他人都說我固執.」其實,如何分辨堅持和固執呢?

還記得上個月鄧瑞強博士講約拿書嗎?約拿這位先知夠堅持了吧!為了不向尼尼微城的人傳遞上帝的警告,上帝叫他去東,他就去西.你還記得約拿書結束時的場面嗎?

4:9上帝對約拿說:「你因這棵蓖麻這樣發怒,對嗎?」他說:「我發怒以至於死,都是對的!」

你如何看約拿的堅持呢?或者有人會說,這不是堅持,而是固執.如此,又回到那個問題:如何分辨堅持和固執呢?這個問題逼使我們思考堅持的是甚麼事情,會不會有些事情值得勇往直前,有些事情其實回頭是岸較為適合呢?

讓我們回到耶穌的故事,耶穌所堅持的是甚麼事情呢?耶穌為了誰而繼續向前不退後呢?第一個問題耶穌直接回答:「今天,明天我趕鬼治病」,耶穌堅持做的事就是趕鬼治病.至於第二個問題,為了誰而堅持趕鬼治病呢?為豐厚的收入而堅持嗎?為了揚名立萬,威振天下而堅持嗎?福音書經常描述耶穌看見受苦的人而動憐憫之心,看來,耶穌堅持趕鬼治病很可能是為了解除他碰見的人的痛苦.

當我們說要學效耶穌一樣堅持到底的時候,我們不妨也問一問自己,堅持要做的是甚麼事情呢?為了誰的福祉而堅持呢?人類是很聰明,懂得把事情合理化,甚至自以為義.例如我們會說,我堅持這樣做都是為了你的需要,解決你的問題.然而,這到底是我的需要還是你的需要呢?是我的問題還是你的問題呢?有一個開罪媽媽的話是這樣說的:「世界上有一種凍,叫阿媽覺得你凍」.

有人認為在市場裡,有些商品及服務需要其實是製造出來的.人們本來沒有這種需要,市場推廣就是要為消費者製造需要,令他們去購買那些商品及服務.耶穌所做的並非如此,他沒有對健康的人說你有病,然後施展能力醫治他們,他只是醫治真正有需要的人；他也沒有為宮殿裡吃得腦滿腸肥的人施行五餅二魚神蹟,只是為在荒野沒有食物的群眾提供食物.世界上真正有需要的人多得是,耶穌堅持向前而不退後,就是為了這些有需要的人.

講了那麼多道理,如何應用呢?我們要學效耶穌去趕鬼醫病嗎?如果你有趕鬼醫病的恩賜,就去做吧!沒有這種能力的人又怎樣呢?可能上帝交給耶穌的任務未必是我的任務；上帝賜給耶穌的恩賜未必是我的恩賜；上帝讓耶穌碰見的人的需要未必是我碰見的人的需要.或者,我們可以問一問自己,上帝賜下甚麼恩賜給我呢?上帝讓我看見這個世界的甚麼需要呢?上帝感召我去做甚麼事情去回應這些需要呢?你的發現和你旁邊的人可能不同,那就好了,正如使徒保羅在羅馬12:6所說:「按我們所得的恩賜,各有不同.」這不同是有好處的,他在以弗所書4:16說:「百節各按各職,照著各體的功用彼此相助,使身體漸漸增長,在愛中建立自己.」當人能夠運用自己的才能,透過行動回應別人的需要,這樣既能「在愛中建立自己」,又能成就社群,即是保羅說的「使身體漸漸增長」.

三．為完成使命而前行

或者有人會對耶穌說:「你往別處去都可以繼續趕鬼醫病,為甚麼死都要去耶路撒冷呢?」耶穌可能會回答:「是的,我死都要去耶路撒冷,因為我去耶路撒冷就是要死.」以上只是設計對白,耶穌的回答是這樣的:

13:33雖然這樣,今天,明天,後天,我必須前行,因為先知在耶路撒冷之外喪命是不能的.

看來我們的設計對白也符合劇情,耶穌走向耶路撒冷,果然是走向死亡.這是他自願的,事先知情的.他曾經三次向門徒預言,自己如何被交在耶路撒冷的權貴手上,被殺害,然後第三天復活.換言之,耶穌每一天都活在死亡的陰影下.死亡的陰影是恐怕的,是嚇人的,可以令人跌入悽楚自憐,長嗟短嘆,或者不知所措之中.然而,這卻也是每個人不可逃避的道路.有人說,「我們從一出生就開始走向死亡」.這也是一種向前走的形式,而且是無法阻擋的向前走.分別在於是無意識地被推著向前走,是憂心忡忡地不得不向前走,還是為了完成使命有意識地向前走.

並不是說懷著使命有意識地向前走就必然是康莊大道.耶穌的狀況及情緒都是複雜及多樣的,他堅持向前走,竟然走到十字架上.在十字架上他竟然大聲呼喊:「我的上帝!我的上帝!為甚麼離棄我?」上帝的兒子走進人類苦難的深淵,被人性的邪惡與黑暗蹂躪的處境.相信當時跟隨耶穌的人,都感到大惑不解.如果耶穌不向耶路撒冷前進,豈不是不用受苦嗎?如果耶穌不走上十字架,豈不是可以為更多人趕鬼醫病嗎?這些想法都是有道理的.不過,耶穌的行動帶來另類可能性.耶穌說「第三天我的事就成全了」,會不會有一種使命是透過受苦去完成的呢?會不會有一種成全是透過到達終點成就的呢?可能世界不理解,也不接受,沒有建功立業,無福無祿無壽,都可以說是人生成全了?可以說生命完滿了?其實耶穌的門徒當時也不理解,不接受.後來,他們領悟到了,因此彼得在彼得前書說:

2:24他被掛在木頭上,親身擔當了我們的罪,使我們既然在罪上死,就得以在義上活.因他受的鞭傷,你們得了醫治.

事實上,的確有人是為了其他人而在某程度上獻出自己的.我有一位朋友結婚不久,正想像著以後甜蜜快樂的婚姻生活時,就傳來配偶患了小腦萎縮症的消息.從此以後,她就按著配偶的身體衰退而不斷改變她的工作及生活方式,以配合照顧他的需要,直至配偶安息主懷.在這十幾年間,他們合作寫了幾本書記錄病情起伏及生活點滴.在最後一本書中,他們表達類似保羅在提摩太後書4:7所說的話:「那美好的仗我已經打過了,當跑的路我已經跑盡了,所信的道我已經守住了.」世界可能有聲音說:「你還年輕,為甚麼被那位活不久的病人阻礙你的人生,撇下他去尋找你的人生豈不更好嗎?」但我的朋友感到她堅持照顧配偶,伴著他渡過患病的日子,正是成就自己的生命,走了自己感到完滿的路.

今天明天後天,耶穌說他必須向前走.我們呢?無論自覺還是不自覺,時間都推著我們向前走.當你走到終點的時候,你會對自己說些甚麼話呢?你會對其他人說甚麼呢?你會對迎接你的上主說甚麼呢?其實也不需要等到終點才說些甚麼話,很多人不時都會做階段性的小結,例如年度回顧,結束一份工作時回顧等.近期有兩位校長完成任命時都唱同一首歌表達心情.這首歌叫做My Way,其中幾段的意思如下:

如今,終點近了 And now, the end is near

在這最後一幕And so I face the final curtain

我的朋友,我要清楚地說 My friend, I'll say it clear

肯定地說 I'll state my case, of which I'm certain

我過了充實的一生I've lived a life that's full

我走過了我生命中的高山低谷I traveled each and every highway

最重要的是,我走完了自己的路And more, much more than this, I did it my way

後悔?我都有少許Regrets, I've had a few

但都是不足掛齒之事But then again, too few to mention

我做了我應該做的I did what I had to do

看破了紅塵,嚐盡甜酸苦辣And saw it through without exemption

我走過人生的大道I planned each charted course

也細味其中的偏僻小路Each careful step along the by way

最重要的是,這是我的人生路And more, much more than this, I did it my way

人生路上,有愛,有笑也有淚I've loved, I've laughed and cried

我的一生有圓也有缺I've had my fill, my share of losing

如今,抹去淚水後And now, as tears subside

赫然發現一路走來是如此美妙I find it all so amusing

回望一路以來的所為To think I did all that

請容我自豪地說And may I say, not in a shy way

這是我的人生路Oh, no, oh, no, not me, I did it my way

每位仍在旅途中每天向前走的人,當你完成一個階段的時候,不知道你會否唱這首歌呢?當我們走盡人生路的時候,保羅在提摩太後書的話會否就是我們的心聲呢?

4:7那美好的仗我已經打過了,當跑的路我已經跑盡了,所信的道我已經守住了,從此以後,有公義的冠冕為我存留.

這些話題對在座中的同學可能有點遙遠,我為你們改寫了保羅的說話:「那美好的知識我已經學過了,當交的功課已經交完了,所有的課外活動我已經參加了,從此以後,有大大張的畢業證書為我存留.」



\newpage

\section{當歡喜快樂}
\label{sec:742}
\textbf{高國雄牧師 - 當歡喜快樂}
\newline
\newline
\hyperref[sec:741]{< < < PREV SERMON < < <}
~
\hyperref[sec:toc]{[返目錄]}
~
\hyperref[sec:743]{> > > NEXT SERMON > > >}


路加福音 15:1-3及11b-32
\newline
\begin{longtable}{cl}
\hline
\hline
章節 & 經文 (和合本修訂版)\\
\hline
15:1 & \begin{tabularx}{0.7\textwidth}{X} 許多稅吏和罪人都挨近耶穌，要聽他講道。 \end{tabularx} \\ \\ \relax
15:2 & \begin{tabularx}{0.7\textwidth}{X} 法利賽人和文士私下議論說：「這個人接納罪人，又同他們吃飯。」 \end{tabularx} \\ \\ \relax
15:3 & \begin{tabularx}{0.7\textwidth}{X} 耶穌就用比喻對他們說： \end{tabularx} \\ \\ \relax
15:4 & \begin{tabularx}{0.7\textwidth}{X} 「你們中間誰有一百隻羊，失去其中的一隻，不把這九十九隻留在曠野，去找那失去的羊，直到找著呢？ \end{tabularx} \\ \\ \relax
15:5 & \begin{tabularx}{0.7\textwidth}{X} 找到了，他就歡歡喜喜地把羊扛在肩上。 \end{tabularx} \\ \\ \relax
15:6 & \begin{tabularx}{0.7\textwidth}{X} 他回到家裡，請朋友和鄰舍來，對他們說：『你們和我一同歡喜吧，我失去的羊已經找到了！』 \end{tabularx} \\ \\ \relax
15:7 & \begin{tabularx}{0.7\textwidth}{X} 我告訴你們，一個罪人悔改，在天上也要這樣為他歡喜，比為九十九個不用悔改的義人歡喜還大呢！」 \end{tabularx} \\ \\ \relax
15:8 & \begin{tabularx}{0.7\textwidth}{X} 「同樣，哪一個婦人有十塊錢，若失落一塊，不點上燈，打掃屋子，細細地找，直到找著呢？ \end{tabularx} \\ \\ \relax
15:9 & \begin{tabularx}{0.7\textwidth}{X} 找到了，她就請朋友和鄰舍來，對她們說：『你們和我一同歡喜吧，我失落的那塊錢已經找到了！』 \end{tabularx} \\ \\ \relax
15:10 & \begin{tabularx}{0.7\textwidth}{X} 我告訴你們，一個罪人悔改，神的使者也是這樣為他歡喜。」 \end{tabularx} \\ \\ \relax
15:11 & \begin{tabularx}{0.7\textwidth}{X} 耶穌又說：「一個人有兩個兒子。 \end{tabularx} \\ \\ \relax
15:12 & \begin{tabularx}{0.7\textwidth}{X} 小兒子對父親說：『父親，請你把我應得的家業分給我。』他父親就把財產分給他們。 \end{tabularx} \\ \\ \relax
15:13 & \begin{tabularx}{0.7\textwidth}{X} 過了不多幾天，小兒子把他一切所有的都收拾起來，往遠方去了。在那裡，他任意放蕩，浪費錢財。 \end{tabularx} \\ \\ \relax
15:14 & \begin{tabularx}{0.7\textwidth}{X} 他耗盡了一切所有的，又恰逢那地方有大饑荒，就窮困起來。 \end{tabularx} \\ \\ \relax
15:15 & \begin{tabularx}{0.7\textwidth}{X} 於是他去投靠當地的一個居民，那人打發他到田裡去放豬。 \end{tabularx} \\ \\ \relax
15:16 & \begin{tabularx}{0.7\textwidth}{X} 他恨不得拿豬所吃的豆莢充飢，也沒有人給他甚麼吃的。 \end{tabularx} \\ \\ \relax
15:17 & \begin{tabularx}{0.7\textwidth}{X} 他醒悟過來，就說：『我父親有多少雇工，糧食有餘，我倒在這裡餓死嗎？ \end{tabularx} \\ \\ \relax
15:18 & \begin{tabularx}{0.7\textwidth}{X} 我要起來，到我父親那裡去，對他說：父親！我得罪了天，又得罪了你， \end{tabularx} \\ \\ \relax
15:19 & \begin{tabularx}{0.7\textwidth}{X} 從今以後，我不配稱為你的兒子，把我當作一個雇工吧。』 \end{tabularx} \\ \\ \relax
15:20 & \begin{tabularx}{0.7\textwidth}{X} 於是他起來，往他父親那裡去。相離還遠，他父親看見，就動了慈心，跑去擁抱著他，連連親他。 \end{tabularx} \\ \\ \relax
15:21 & \begin{tabularx}{0.7\textwidth}{X} 兒子對他說：『父親！我得罪了天，又得罪了你，從今以後，我不配稱為你的兒子。』 \end{tabularx} \\ \\ \relax
15:22 & \begin{tabularx}{0.7\textwidth}{X} 父親卻吩咐僕人：『快把那上好的袍子拿出來給他穿，把戒指戴在他指頭上，把鞋穿在他腳上， \end{tabularx} \\ \\ \relax
15:23 & \begin{tabularx}{0.7\textwidth}{X} 把那肥牛犢牽來宰了，我們來吃喝慶祝； \end{tabularx} \\ \\ \relax
15:24 & \begin{tabularx}{0.7\textwidth}{X} 因為我這個兒子是死而復活，失而復得的。』他們就開始慶祝。 \end{tabularx} \\ \\ \relax
15:25 & \begin{tabularx}{0.7\textwidth}{X} 「那時，大兒子正在田裡。他回來，離家不遠時，聽見奏樂跳舞的聲音， \end{tabularx} \\ \\ \relax
15:26 & \begin{tabularx}{0.7\textwidth}{X} 就叫一個僮僕來，問是甚麼事。 \end{tabularx} \\ \\ \relax
15:27 & \begin{tabularx}{0.7\textwidth}{X} 僮僕對他說：『你弟弟回來了，你父親因為他無災無病地回來，把肥牛犢宰了。』 \end{tabularx} \\ \\ \relax
15:28 & \begin{tabularx}{0.7\textwidth}{X} 大兒子就生氣，不肯進去，他父親出來勸他。 \end{tabularx} \\ \\ \relax
15:29 & \begin{tabularx}{0.7\textwidth}{X} 他對父親說：『你看，我服侍你這麼多年，從來沒有違背過你的命令，而你從來沒有給我一隻小山羊，叫我和朋友們一同快樂。 \end{tabularx} \\ \\ \relax
15:30 & \begin{tabularx}{0.7\textwidth}{X} 但你這個兒子和娼妓吃光了你的財產，他一回來，你倒為他宰了肥牛犢。』 \end{tabularx} \\ \\ \relax
15:31 & \begin{tabularx}{0.7\textwidth}{X} 父親對他說：『兒啊！你常和我同在，我所有的一切都是你的； \end{tabularx} \\ \\ \relax
15:32 & \begin{tabularx}{0.7\textwidth}{X} 可是你這個弟弟是死而復活，失而復得的，所以我們理當歡喜慶祝。』」 \end{tabularx} \\ \\
[1ex]
\hline
\hline
\end{longtable}
講題:當歡喜快樂 Be Happy!

經文:路加福音15章1至3節及11b至32節

講員:高國雄牧師

場合:香港中文大學崇基學院禮拜堂主日崇拜

日期:2025年3月30日

我們熟悉的浪子回頭,相信是很多人被感動的故事,歷世歷代有極多藝術,音樂嘗試去記錄,分享這故事.

讓我們一同回顧路加福音的獨有記載:

路加福音15:11-24

11耶穌又說:「一個人有兩個兒子. 12小兒子對父親說:『父親,請你把我應得的家業分給我.』他父親就把產業分給他們. 13過了不多幾日,小兒子就把他一切所有的都收拾起來,往遠方去了.在那裡任意放蕩,浪費資財. 14既耗盡了一切所有的,又遇著那地方大遭饑荒,就窮苦起來. 15於是去投靠那地方的一個人；那人打發他到田裡去放豬. 16他恨不得拿豬所吃的豆莢充飢,也沒有人給他. 17他醒悟過來,就說:『我父親有多少的雇工,口糧有餘,我倒在這裡餓死嗎? 18我要起來,到我父親那裡去,向他說:父親!我得罪了天,又得罪了你； 19從今以後,我不配稱為你的兒子,把我當作一個雇工吧!』 20於是起來,往他父親那裡去.相離還遠,他父親看見,就動了慈心,跑去抱著他的頸項,連連與他親嘴. 21兒子說:『父親!我得罪了天,又得罪了你；從今以後,我不配稱為你的兒子.』 22父親卻吩咐僕人說:『把那上好的袍子快拿出來給他穿；把戒指戴在他指頭上；把鞋穿在他腳上； 23把那肥牛犢牽來宰了,我們可以吃喝快樂； 24因為我這個兒子是死而復活,失而又得的.』他們就快樂起來.

故事內容:兒子要分父親的家業,浪費耗盡,走投無路,回想起家中的父親,回轉,父親接納.

比喻的目的?

父親=全愛的天父

兒子=我們?

為何比喻中有小兒子及大兒子?

分別代表那些人?

我們是大兒子?還是小兒子呢?

V25-32的價值及意義?

路加福音15:25-32

25那時,大兒子正在田裡.他回來,離家不遠,聽見作樂跳舞的聲音, 26便叫過一個僕人來,問是甚麼事. 27僕人說:『你兄弟來了；你父親因為得他無災無病地回來,把肥牛犢宰了.』 28大兒子卻生氣,不肯進去；他父親就出來勸他. 29他對父親說:『我服事你這多年,從來沒有違背過你的命,你並沒有給我一隻山羊羔,叫我和朋友一同快樂. 30但你這個兒子和娼妓吞盡了你的產業,他一來了,你倒為他宰了肥牛犢.』 31父親對他說:『兒啊!你常和我同在,我一切所有的都是你的； 32只是你這個兄弟是死而復活,失而又得的,所以我們理當歡喜快樂.』」

大兒子生氣,合理嗎?我們也是如此嗎?

父親的回應合理嗎?

我們真的「理當」歡喜快樂嗎?

經文不是單獨的故事,路加福音記錄耶穌同時講了3個故事

路加福音15:4-7

4「你們中間誰有一百隻羊失去一隻,不把這九十九隻撇在曠野,去找那失去的羊,直到找著呢? 5找著了,就歡歡喜喜地扛在肩上,回到家裡, 6就請朋友鄰舍來,對他們說:『我失去的羊已經找著了,你們和我一同歡喜吧!』 7我告訴你們,一個罪人悔改,在天上也要這樣為他歡喜,較比為九十九個不用悔改的義人歡喜更大.」

路加福音15:8-10

8「或是一個婦人有十塊錢,若失落一塊,豈不點上燈,打掃屋子,細細地找,直到找著嗎? 9找著了,就請朋友鄰舍來,對他們說:『我失落的那塊錢已經找著了,你們和我一同歡喜吧!』 10我告訴你們,一個罪人悔改,在上帝的使者面前也是這樣為他歡喜.」

失羊的故事與及失錢的故事

3個故事=1個比喻,為要回應

路加福音15:1-3

1眾稅吏和罪人都挨近耶穌,要聽他講道. 2法利賽人和文士私下議論說:「這個人接待罪人,又同他們吃飯.」 3耶穌就用比喻說:

所以是「失羊」,「失錢」,「失人」…失而復得的比喻

是百分之一,十分之一,二分之一,愈發珍貴的三層推展式的比喻

為要突出失而復得的人的價值!

人的自我,自以為事

上帝的眼目,耶穌基督的心腸

羅馬書5:8

惟有基督在我們還作罪人的時候為我們死,上帝的愛就在此向我們顯明了.

守望你(City Beat)1988

詞:韋然. 曲:John Laudon

思念你 沒忘記

當日你 獨自離家

你充滿自信 披風雨獨闖

永不會後悔 你倔強的說

想念你 仍懷念你

希望你 能回家

我深愛著你 當你失意地哭

你可叫喚我 我定會親你 你莫痛悲

每天地球在轉 天光與天黑

關心你靜看你 如何笑或流淚

是誰負你 因此覺渺小

知否我極愛你 全意守望你

思念你 熱愛從未變

若有天 當你不得意

可叫喚我 當你失了目標

你不要害怕 我定會親你 你莫痛悲

每天地球在轉 天光與天黑

關心你靜看你 如何笑或流淚

是誰負你 因此覺渺小

知否我極愛你 全意守望你

熱愛從未變

希望你 能回家 喔



\newpage

\section{石頭也要呼叫}
\label{sec:743}
\textbf{李均熊牧師 - 石頭也要呼叫}
\newline
\newline
\hyperref[sec:742]{< < < PREV SERMON < < <}
~
\hyperref[sec:toc]{[返目錄]}
~
\hyperref[sec:744]{> > > NEXT SERMON > > >}


路加福音 19:28-40
\newline
\begin{longtable}{cl}
\hline
\hline
章節 & 經文 (和合本修訂版)\\
\hline
19:28 & \begin{tabularx}{0.7\textwidth}{X} 耶穌說完了這些話，就走在前面，上耶路撒冷去。 \end{tabularx} \\ \\ \relax
19:29 & \begin{tabularx}{0.7\textwidth}{X} 快到伯法其和伯大尼，在名叫橄欖山的地方，他打發兩個門徒， \end{tabularx} \\ \\ \relax
19:30 & \begin{tabularx}{0.7\textwidth}{X} 說：「你們往對面村子裡去，進去的時候會看見一匹驢駒拴在那裡，是從來沒有人騎過的，把牠解開，牽來。 \end{tabularx} \\ \\ \relax
19:31 & \begin{tabularx}{0.7\textwidth}{X} 若有人問為甚麼解開牠，你們就這樣說：『主要用牠。』」 \end{tabularx} \\ \\ \relax
19:32 & \begin{tabularx}{0.7\textwidth}{X} 被打發的人去了，所遇見的正如耶穌對他們所說的。 \end{tabularx} \\ \\ \relax
19:33 & \begin{tabularx}{0.7\textwidth}{X} 他們解開驢駒的時候，主人問他們：「為甚麼解開驢駒？」 \end{tabularx} \\ \\ \relax
19:34 & \begin{tabularx}{0.7\textwidth}{X} 他們說：「主要用牠。」 \end{tabularx} \\ \\ \relax
19:35 & \begin{tabularx}{0.7\textwidth}{X} 他們把驢駒牽到耶穌那裡，把自己的衣服搭在上面，扶耶穌騎上。 \end{tabularx} \\ \\ \relax
19:36 & \begin{tabularx}{0.7\textwidth}{X} 他前進的時候，眾人把衣服鋪在路上。 \end{tabularx} \\ \\ \relax
19:37 & \begin{tabularx}{0.7\textwidth}{X} 他將近耶路撒冷，正下橄欖山的時候，一大群門徒因所見過的一切異能，都歡呼起來，大聲讚美神， \end{tabularx} \\ \\ \relax
19:38 & \begin{tabularx}{0.7\textwidth}{X} 說：「奉主名來的王是應當稱頌的！在天上有和平；在至高之處有榮光。」 \end{tabularx} \\ \\ \relax
19:39 & \begin{tabularx}{0.7\textwidth}{X} 人群中有幾個法利賽人對耶穌說：「老師，責備你的門徒吧！」 \end{tabularx} \\ \\ \relax
19:40 & \begin{tabularx}{0.7\textwidth}{X} 耶穌回答：「我告訴你們，若是這些人閉口不說，石頭也要呼叫起來。」 \end{tabularx} \\ \\
[1ex]
\hline
\hline
\end{longtable}
講題:石頭也要呼叫 The Stones Will Cry Out

經文:路加福音19:28-40

講員:李均熊牧師

場合:香港中文大學崇基學院禮拜堂主日崇拜

日期:2025年4月13日

起:耶穌光榮地進入耶路撒冷?

負責兒童事工的弟兄姊妹安排棕枝主日的崇拜其實十分尷尬.

他們或會在今天專登買一些棕櫚葉,然後叫教會的兒童扮演迎接耶穌進入耶路撒冷的群眾,通常伴隨著棕櫚葉在空中尷尬地揮舞,(也許孩子們會用它們來打弟弟妹妹的頭),會眾配合著唱幾首有「和散拿」字句的讚美詩,然後趕走了喧鬧的孩子們,因為到了講道時間,我們的焦點要轉向十字架,因為這是耶穌受難週的開始,講道又要變得嚴肅又沉重,絕對兒童不兒.

明明耶穌光榮地進入耶路撒冷(聖壇布置也是用慶典的紅色),但因為聖經一早劇透了耶穌跟著是要受苦受死,於是乎一場漂亮的巡行就要尷尷尬尬又匆匆忘忘的作個過場,我們又要變得愁雲慘霧,好像預苦期是不能容納一絲的歡樂.

承:掃興的法利賽人

不同於其他福音書, 路加福音特別記載了一些掃興的法利賽人,當他們見到耶穌門徒因為「所見過的一切異能」,而興奮過頭,「歡樂起來,大聲讚美上帝,說:奉主名來的王是應當稱頌的!」就和耶穌說:「夫子,責備你的門徒吧!」

為甚麼法利賽人這麼掃興呢?很可能是那個「王」字特別刺耳吧?其實法利賽人的潛台詞可能是:「夫子啊,現在耶路撒冷正在過節時期,羅馬政府加派了大量的駐防軍, 以防在過節期間有甚麼風吹草動,就可以立刻平亂.但你的門徒這麼明目張膽叫你做「王」,你不怕羅馬軍把你捉拿,告你煽動顛覆嗎?夫子,你話o下你D 門徒啦!你被人拉jor 無所謂,攪到羅馬軍封jor 耶路撒冷城唔比我地過節就大件事!」

但耶穌怎樣回答這些法利賽人呢?

「我告訴你們,若是他們閉口不說,這些石頭必要呼叫起來.」

不要誤會,我想耶穌不是要我們周圍宣揚耶穌是要取代極權帝君的王.在福音書中,他是極力回避被人稱他作王,又或一律禁止別人擁戴他為王.在這樣關鍵時刻,他自然不想被羅馬兵丁帶走.你可以推想為甚麼耶穌要騎著驢駒進入耶路撒冷,就是要讓羅馬駐軍忽視他的存在,但同時讓熟識舊約像徵的猶太人明白他確實是奉主名而來的王.明就明,唔明就搵黎明…….

但是若我們閉口不說那一句出自詩篇118:26的那一句「奉主名來的是應當稱頌的」,若我們不去宣揚耶穌作為基督,就是那在天上的和平和至高之處的榮光,那麼就算路邊的石頭也要呼叫起來!

轉: 為耶路撒冷哀哭

「和平」與「榮光」doxa都像在路加福音哪裡出現過……不錯,在耶穌出生的時候,就有天使對著野外牧羊人高唱:在至高之處榮耀doxa歸與神!在地上平安歸與他所喜悅的人!」(路2:14)耶穌入耶路撒冷就如他在天上闖入人間去為我們帶來和平和救贖一樣.

至於詩118:26 「奉主名來的是應當稱頌的」也在路加福音出現過,就在路加福音13 章,當耶穌為耶路撒冷哀哭的時候出現.那時也有法利賽人出現:

31正當那時,有幾個法利賽人來對耶穌說:「離開這裡去吧,因為希律想要殺你.」32耶穌說:「你們去告訴那個狐狸說:『今天,明天我趕鬼治病,第三天我的事就成全了.』33雖然這樣,今天,明天,後天,我必須前行,因為先知在耶路撒冷之外喪命是不能的.34耶路撒冷啊!耶路撒冷啊!你常殺害先知,又用石頭打死那奉差遣到你這裡來的人.我多次願意聚集你的兒女,好像母雞把小雞聚集在翅膀底下；只是你們不願意.35看哪,你們的家成為荒場留給你們.我告訴你們,從今以後你們不得再見我,直等到你們說:『奉主名來的是應當稱頌的.』」

這些早前的耶穌說話,是以條件承諾的形式表達的,彷彿在說:「如果以色列悔改,並以稱頌迎接奉主名來的那一位,那麼以色列將會經歷救恩的來臨.」這段經文與經課這段經文之間的聯繫並不甚愉快.因為迎接耶穌的是耶穌的門徒和群眾,而不是耶路撒冷本身.——甚至法利賽人對門徒歡欣高唱「奉主名來的是應當稱頌的」的這句歌頌嗤之以鼻!

正如接下來的經文(41-44節是經課沒有的),耶穌再一次為耶路撒冷哀哭,同樣用上和平和石頭的字眼:

41耶穌快到耶路撒冷,看見那城,就為它哀哭,42說:「但願你在這日子知道有關你平安的事,不過這事現在是隱藏的,你的眼睛看不出來.43因為日子將到,你的仇敵要築起土壘包圍你,四面困住你,44並要消滅你和你裡頭的兒女,連一塊石頭也不留在另一塊石頭上,因為你不知道你蒙眷顧的時候.」

這裡路加福音用法利賽人的口,說出耶路撒冷的頑固,因為他們沒有認識到耶穌的來臨,就是救恩的來臨.

合:請不要拒絕基督

今日是棕枝主日,是紀念耶穌光榮進入耶路撒冷的聖日.不錯,預苦期進入最後一週,再過幾天,就是耶穌被針十字架的受難節.每年預苦期,當中你或會禁誡一些事情,例如禁食肉,好紀念基督的受苦,所以有時預苦期又稱為大齋期.但這個齋誡的舉動,即預苦期的四十日,是不計星期日在內的.因為星期日是我們紀念基督復活的日子.

同樣道理,棕枝主日是我們紀念耶穌受難之後所要帶來榮耀的日子,是我們悔改得蒙接納的日子,是我們崇拜歌頌上主救恩的日子,是我們應該歡喜快樂的日子.所以讓我們不再尷尷尬尬,而是能高聲歡唱「和散那,那奉主名而來的是應當稱頌的!」

請不要好像法利賽人一般,拒絕耶穌基督的救恩.但就算我們選擇閉口不說上主救恩的奇妙和基督帶來的平安榮耀,石頭也必要呼叫起來.



\newpage

\section{復活天人合一的故事}
\label{sec:744}
\textbf{鄧瑞強博士 - 復活:天人合一的故事}
\newline
\newline
\hyperref[sec:743]{< < < PREV SERMON < < <}
~
\hyperref[sec:toc]{[返目錄]}
~
\hyperref[sec:745]{> > > NEXT SERMON > > >}


約翰福音 20:11-18
\newline
\begin{longtable}{cl}
\hline
\hline
章節 & 經文 (和合本修訂版)\\
\hline
20:11 & \begin{tabularx}{0.7\textwidth}{X} 馬利亞卻站在墳墓外面哭。她哭的時候，低頭往墳墓裡看， \end{tabularx} \\ \\ \relax
20:12 & \begin{tabularx}{0.7\textwidth}{X} 看見兩個天使穿著白衣，在安放耶穌身體的地方坐著，一個在頭，一個在腳。 \end{tabularx} \\ \\ \relax
20:13 & \begin{tabularx}{0.7\textwidth}{X} 天使對她說：「婦人，你為甚麼哭？」她對他們說：「因為有人把我主移走了，我不知道他們把他放在哪裡。」 \end{tabularx} \\ \\ \relax
20:14 & \begin{tabularx}{0.7\textwidth}{X} 說了這些話，她轉過身來，看見耶穌站在那裡，卻不知道他是耶穌。 \end{tabularx} \\ \\ \relax
20:15 & \begin{tabularx}{0.7\textwidth}{X} 耶穌問她：「婦人，你為甚麼哭？你找誰？」馬利亞以為他是看園子的，就對他說：「先生，若是你把他移了去，請告訴我，你把他放在哪裡，我去把他移回來。」 \end{tabularx} \\ \\ \relax
20:16 & \begin{tabularx}{0.7\textwidth}{X} 耶穌對她說：「馬利亞。」馬利亞轉過身來，用希伯來話對他說：「拉波尼！」（「拉波尼」就是老師的意思。） \end{tabularx} \\ \\ \relax
20:17 & \begin{tabularx}{0.7\textwidth}{X} 耶穌對她說：「不要拉住我，因為我還沒有升上去見我的父。你到我弟兄那裡去告訴他們，我要升上去見我的父，也是你們的父，見我的神，也是你們的神。」 \end{tabularx} \\ \\ \relax
20:18 & \begin{tabularx}{0.7\textwidth}{X} 抹大拉的馬利亞就向門徒報信：「我已經看見了主。」她又把主對她說的話告訴他們。 \end{tabularx} \\ \\
[1ex]
\hline
\hline
\end{longtable}
講題:復活:天人合一的故事 Resurrection: The Unity of Heaven and Earth

經文:約翰福音20:11-18

講員:鄧瑞強博士

場合:香港中文大學崇基學院禮拜堂主日崇拜

日期:2025年4月20日

各位弟兄姊妹,早安.願神的恩典常與你們同在.

今日是復活節,是基督信仰最重要的節日之一,記念主耶穌從死裡復活,修復這個被罪惡和死亡操控的世界,讓我們看到,天父仍是天地的主,生命仍是天父終極要送贈給蒼生的禮物.世界的權勢用死亡去掌控他人,天父以生命去解救世界,這是基督信仰給我們看見的真相,這真相給我們希望,在殘酷的世界裡,仍堅持活出仁愛與真理.

我們看看今天的講道經文:約翰福音20:11-18

約20:11馬利亞卻站在墳墓外面哭.哭的時候,低頭往墳墓裡看,

約20:12就見兩個天使,穿著白衣,在安放耶穌身體的地方坐著,一個在頭,一個在腳.

約20:13天使對她說:「婦人,你為甚麼哭?」她說:「因為有人把我主挪了去,我不知道放在哪裡.」

約20:14說了這話,就轉過身來,看見耶穌站在那裡,卻不知道是耶穌.

約20:15耶穌問她說:「婦人,為甚麼哭?你找誰呢?」馬利亞以為是看園的,就對他說:「先生,若是你把他移了去,請告訴我,你把他放在哪裡,我便去取他.」

約20:16耶穌說:「馬利亞.」馬利亞就轉過來,用希伯來話對他說:「拉波尼!」(拉波尼就是夫子的意思.)

約20:17耶穌說:「不要摸我,因我還沒有升上去見我的父.你往我弟兄那裡去,告訴他們說,我要升上去見我的父,也是你們的父,見我的上帝,也是你們的上帝.」

約20:18抹大拉的馬利亞就去告訴門徒說:「我已經看見了主.」她又將主對她說的這話告訴他們.

(1)施恩座

耶穌死了,安躺在墳墓,馬利亞在哀悼.這馬利亞不是耶穌的母親,而是抹大拉的馬利亞,按其他福音書的記載,耶穌曾從這位馬利亞身上趕出七隻鬼.看來,這位女士曾被過往的痛苦經歷牢牢抓住,而耶穌解救了她.她領受耶穌的恩典很深,也愛耶穌愛得很深.她為耶穌的死深深哀傷.

約20:11馬利亞卻站在墳墓外面哭.哭的時候,低頭往墳墓裡看,

約20:12就見兩個天使,穿著白衣,在安放耶穌身體的地方坐著,一個在頭,一個在腳.

馬利亞發現,耶穌的遺體不見了,在原先安放耶穌遺體的地方,有兩個天使,一個在頭,一個在腳.我們對這幅畫面,沒有任何感覺.但猶太人一看,便會即刻意識到,這就是約櫃上的施恩座.

我們去以色列旅行,有時會買這約櫃做手信.施恩座其實是約櫃的頂蓋,代表神的臨在.神的臨在,守護著神與人的約的連繫.施恩座的兩邊,正是兩個天使相對.在贖罪日,祭牲的血會彈在施恩座上,仰望神恩典裡的寬恕.當然,若躺著耶穌遺體的石板成了施恩座,則耶穌的血也多少會沾在這施恩座上,表達出神對罪人的赦免.

在這復活日,天父進行其重造天地的偉大工程.祂將耶穌的死地,創造成新的施恩座.人類的罪,固然破壞神和人的關係.但耶穌因愛而流血的行動,表達出神以愛來回報人類的罪.神在安放耶穌遺體的地方,創造出新的施恩座,神在這裡,臨在人間.神和人,天和地,在耶穌的犧牲裡,重新復和.神不是不理人間死活的,祂在耶穌的死裡,臨在人間；在耶穌的復活裡,顯明祂勝過罪惡和死亡.這施恩座,是神在人間作王的記號.

安放耶穌遺體的地方,成了施恩座.在今日,或者可以說,十字架這受死記號,象徵著神的施恩座,神在受苦的記號中臨在.神的這種臨在的方式,其實不太符合人類對神明的構想.當我們認為某些人是敵人時,我們會肯定自己是正義的,對方是不義的,而神明的臨在,就是消滅對方.神明變成我們以暴力解決問題的方式.神拒絕這種臨在方式.以受死的十架為施恩座,為神臨在的地方,就是以愛還恨,唯有這樣,才能打破暴力的循環,天國才能圓滿實現.

(2)新的伊甸園

約20:14說了這話,就轉過身來,看見耶穌站在那裡,卻不知道是耶穌.

約20:15耶穌問她說:「婦人,為甚麼哭?你找誰呢?」馬利亞以為是看園的,就對他說:「先生,若是你把他移了去,請告訴我,你把他放在哪裡,我便去取他.」

這裡的場景,是一個園子,馬利亞以為耶穌是看園的.復活的主,在園中,猶太人一看,便會聯想到在創世記裡神在伊甸園中行走的故事.馬利亞以為耶穌是看園的,我們讀到這裡,多少都會認為馬利亞真是糊塗.其實,約翰福音有一種諷刺的書寫手法,表面上我們認為不可能是某東西的,事實上偏偏就是那東西.我們以為馬利亞是誤會了,耶穌怎會是園丁,但偏偏,耶穌就真是園丁.因著復活,這世界其實改變了.安放耶穌遺體的地方,成了施恩座,這是聖殿最核心的物件.墓園以至大地,因而成了聖殿.我們的世界,成了聖殿,成了新的伊甸園,神在其中行走.在太初,神在伊甸園行走,吩咐亞當和夏娃看守這伊甸園.如今,復活的耶穌,親自看守大地,也看守著我們.祂會一直看守,直到天國降臨人間.

(3)新的時代

約20:16耶穌說:「馬利亞.」馬利亞就轉過來,用希伯來話對他說:「拉波尼!」(拉波尼就是夫子的意思.)

約20:17耶穌說:「不要摸我〔和修版:「不要拉住我」〕,因我還沒有升上去見我的父.你往我弟兄那裡去,告訴他們說,我要升上去見我的父,也是你們的父,見我的上帝,也是你們的上帝.」

憑樣貌,馬利亞認不出耶穌,但她認出耶穌的呼喚.在創世記,亞當,夏娃犯罪後,躲藏起來,逃避神的呼喚.如今,在新的伊甸園裡,主耶穌呼喚馬利亞,馬利亞即時有反應,她回應說:「夫子」(或者譯「主」),這是人面對神時應有的反應.可以說,這是新伊甸園裡的新人性.馬利亞的反應,也代表著日後教會群體對主耶穌的正面回應:耶穌,你是主.

但明顯,馬利亞還不明白天地已經改變,時代已經轉換.馬利亞拉住耶穌,一般而言,這代表她想和耶穌返回以往的日子.像昔日,在加利利,一起吃飯,聽耶穌講道.她不明白,復活不是「翻生」.若耶穌只是翻生的話,則一切是可以回復從前的.但復活不是翻生,復活,是神的大能闖進人間,將腐朽的生命轉化成不朽的生命,這是被罪惡蠶食的人間在神的恩典下徹底被轉化.這不單是新的天地,也是新的時代.用新約的術語來說,就是:舊的世代過去了,新的世代開始了.在耶穌復活的那天,空間被改變了,這是新的天地；時間被改變了,這是新的世代.很多信徒可能意識不到,我們已活在新的時空裡.

馬利亞不能再以塵世的眼光看耶穌,她要理解耶穌的永恆性,這永恆性,以耶穌「升上去見他的父」來表達.馬利亞要看到耶穌這永恆生命來到人間所帶來的翻天覆地的轉化.因著神人復和,天人合一,與耶穌生命連結的信徒,生命也改變了.信徒不再是一般的人,耶穌稱我們為「兄弟」.天父不再是外太空的神明,而是我們的父,我們的神.我們以「神的兒女」的新生命,活在新的時空裡.希望弟兄姊妹,不要拉住,眷戀舊的世界,舊的價值,舊的渴求,舊的目標,舊的經歷,舊的人生.「不要效法這個世界,只要心意更新而變化,叫你們察驗何為神的善良,純全,可喜悅的旨意.」(羅12:2)「若有人在基督裡,他就是新造的人,舊事已過,都變成新的了.」(林後5:17)

信徒的問題,很多時就是不知道自己已進入了新的時代,新的世界,有了新的身分,仍像馬利亞那樣,想抓住舊的人生.這就像有些人移民去了外國,卻又學習不到新的語言,不能了解新的文化,完全無法進入一個新的世界裡,只是不斷依戀著舊的生活,舊有的一切.這樣的人生,格格不入,無法回到以往,也無法進入新的將來.

因著耶穌的復活,信徒已進入了新的世界,新的時代裡.舊世界的價值觀,即:對不同的被造物的追求,應該放下了.在這天人合一,如聖殿一樣的世界裡,應追求天國降臨,追求神的旨意行在地上,如同行在天上.想想,我們要活得多光明,世人才能在我們身上,看到神的真實,神的臨在?

在耶穌的復活裡,原先安放遺體的地方成了神臨在的施恩座,成了聖殿最核心的部分「至聖所」,這世界成了聖殿,我們的身分因而也成了祭司.正如彼得前書2:9的講法,信徒是「君尊的祭司」,是帶有王者氣度的祭司,像大祭司主耶穌一樣,在人間,溝通天地,促成神人復和.真正的復和,不是暴力壓制,而是以愛去聯繫,這就是十架的永恆信息.

(4)結語

復活節,就像二戰時的D-day.在這一天,盟軍反攻納粹德軍,破其防線,挫其精銳,奠定勝利的前景.雖然很多地方,表面看來,盟軍和納粹仍在周旋,但事實上,D-day之後,盟軍勝利已是仍在延續的小戰役的底層事實.一切已經改變了.

復活節,就是令天地改變的D-day,罪惡和死亡的權勢被擊敗了,神的生命大能已成為萬事萬物的底層真相.縱使表面上,罪惡和死亡的力量仍在張牙舞爪,但實際上,神的大能已然彰顯,天國的圓滿實現已可預期.

約20:18抹大拉的馬利亞就去告訴門徒說:「我已經看見了主.」她又將主對她說的這話告訴他們.

馬利亞很興奮,她說:「我已經看見了主.」她同時在復活主裡看見一切的改變,她對世界有信息可傳.她看見施恩座在人間,看見神在十架死亡裡的臨在,看見天地轉化成聖殿,看見天父成了我們的天父,這一切,意味著:她看見人生的使命.她不可能像舊時一樣生活,而會活得像帶有王者氣度的祭司,在人間,活出神聖.

在這復活節,期望各人,在人間,努力活出神聖.

願神祝福你們.

願榮耀歸予聖父,聖子,聖靈.



\newpage

\section{不要疑惑總要信}
\label{sec:745}
\textbf{高國雄牧師 - 不要疑惑,總要信}
\newline
\newline
\hyperref[sec:744]{< < < PREV SERMON < < <}
~
\hyperref[sec:toc]{[返目錄]}
~
\hyperref[sec:746]{> > > NEXT SERMON > > >}


約翰福音 20:19-31
\newline
\begin{longtable}{cl}
\hline
\hline
章節 & 經文 (和合本修訂版)\\
\hline
20:19 & \begin{tabularx}{0.7\textwidth}{X} 那日（就是七日的第一日）晚上，門徒因怕猶太人，所在的地方門都關了。耶穌來，站在當中，對他們說：「願你們平安！」 \end{tabularx} \\ \\ \relax
20:20 & \begin{tabularx}{0.7\textwidth}{X} 說了這話，他把手和肋旁給他們看。門徒一看見主就喜樂了。 \end{tabularx} \\ \\ \relax
20:21 & \begin{tabularx}{0.7\textwidth}{X} 於是耶穌又對他們說：「願你們平安！父怎樣差遣了我，我也照樣差遣你們。」 \end{tabularx} \\ \\ \relax
20:22 & \begin{tabularx}{0.7\textwidth}{X} 說了這話，他向他們吹一口氣，說：「領受聖靈吧！ \end{tabularx} \\ \\ \relax
20:23 & \begin{tabularx}{0.7\textwidth}{X} 你們赦免誰的罪，誰的罪就得赦免；你們不赦免誰的罪，誰的罪就不得赦免。」 \end{tabularx} \\ \\ \relax
20:24 & \begin{tabularx}{0.7\textwidth}{X} 那十二使徒中，有個叫低土馬的多馬，耶穌來的時候，他沒有和他們在一起。 \end{tabularx} \\ \\ \relax
20:25 & \begin{tabularx}{0.7\textwidth}{X} 其他的門徒就對他說：「我們已經看見主了。」多馬卻對他們說：「除非我看見他手上的釘痕，用我的指頭探入那釘痕，用我的手探入他的肋旁，我絕不信。」 \end{tabularx} \\ \\ \relax
20:26 & \begin{tabularx}{0.7\textwidth}{X} 過了八日，門徒又在屋裡，多馬也和他們在一起。門都關了，耶穌來，站在當中，說：「願你們平安！」 \end{tabularx} \\ \\ \relax
20:27 & \begin{tabularx}{0.7\textwidth}{X} 然後他對多馬說：「把你的指頭伸到這裡來，看看我的手；把你的手伸過來，探入我的肋旁。不要疑惑，總要信！」 \end{tabularx} \\ \\ \relax
20:28 & \begin{tabularx}{0.7\textwidth}{X} 多馬回答，對他說：「我的主！我的神！」 \end{tabularx} \\ \\ \relax
20:29 & \begin{tabularx}{0.7\textwidth}{X} 耶穌對他說：「你因為看見了我才信嗎？那沒有看見卻信的有福了。」 \end{tabularx} \\ \\ \relax
20:30 & \begin{tabularx}{0.7\textwidth}{X} 耶穌在他門徒面前另外行了許多神蹟，沒有記錄在這書上。 \end{tabularx} \\ \\ \relax
20:31 & \begin{tabularx}{0.7\textwidth}{X} 但記載這些事是要使你們信耶穌是基督，是神的兒子，並且使你們信他，好因著他的名得生命。 \end{tabularx} \\ \\
[1ex]
\hline
\hline
\end{longtable}
講題:不要疑惑,總要信 Do Not Doubt But Believe

經文:約翰福音20:19-31

講員:高國雄牧師

場合:香港中文大學崇基學院禮拜堂主日崇拜

日期:2025年4月27日

約翰福音 20:19-31

19那日(就是七日的第一日)晚上,門徒所在的地方,因怕猶太人,門都關了.耶穌來,站在當中,對他們說:「願你們平安!」

20說了這話,就把手和肋旁指給他們看.門徒看見主,就喜樂了.

21耶穌又對他們說:「願你們平安!父怎樣差遣了我,我也照樣差遣你們.」

22說了這話,就向他們吹一口氣,說:「你們受聖靈!

23你們赦免誰的罪,誰的罪就赦免了；你們留下誰的罪,誰的罪就留下了.」

24那十二個門徒中,有稱為低土馬的多馬；耶穌來的時候,他沒有和他們同在.

25那些門徒就對他說:「我們已經看見主了.」多馬卻說:「我非看見他手上的釘痕,用指頭探入那釘痕,又用手探入他的肋旁,我總不信.」

26過了八日,門徒又在屋裡,多馬也和他們同在,門都關了.耶穌來,站在當中說:「願你們平安!」

27就對多馬說:「伸過你的指頭來,摸我的手；伸出你的手來,探入我的肋旁.不要疑惑,總要信!」

28多馬說:「我的主!我的上帝!」

29耶穌對他說:「你因看見了我才信；那沒有看見就信的有福了.」

30耶穌在門徒面前另外行了許多神蹟,沒有記在這書上.

31但記這些事要叫你們信耶穌是基督,是上帝的兒子,並且叫你們信了他,就可以因他的名得生命.

耶穌基督復活後的顯現,帶來什麼變化?

l  路加福音24:13-49

往以馬忤斯的2個門徒,憂愁離開…

l  馬太福音28:1-10,馬可福音16:9-11,約翰福音20:11-18*(上主日的經課)

婦女因愛尋找耶穌遺體,尋見…

l  約翰福音 21:1-23*(下主日的經課)

打魚的門徒遇上復活顯現的耶穌,耶穌再次呼召彼得

約21:14耶穌從死裡復活以後,向門徒顯現,這是第三次.

多馬與耶穌

l  約翰福音 20:19-31

耶穌2次的顯現

一個顯現的故事:耶穌連續2個主日在同一地點顯現

約翰福音 20

19那日(就是七日的第一日)晚上,門徒所在的地方,因怕猶太人,門都關了.耶穌來,站在當中,對他們說:「願你們平安!」 20說了這話,就把手和肋旁指給他們看.門徒看見主,就喜樂了.

n  由恐懼,害怕,因主的復活而喜樂

24那十二個門徒中,有稱為低土馬的多馬；耶穌來的時候,他沒有和他們同在. 25那些門徒就對他說:「我們已經看見主了.」多馬卻說:「我非看見他手上的釘痕,用指頭探入那釘痕,又用手探入他的肋旁,我總不信.」

n  可憐的多馬…被理解為「多疑」的多馬,「理性」的多馬,其實他是「認真」的多馬,「與你我一樣」的多馬…

26過了八日,門徒又在屋裡,多馬也和他們同在,門都關了.

門為何關上?

喜樂未能驅走恐懼?因為多馬?

27就對多馬說:「伸過你的指頭來,摸我的手；伸出你的手來,探入我的肋旁.不要疑惑,總要信!」

n  耶穌主動各多馬伸手,耶穌極愛多馬

28多馬說:「我的主!我的上帝!」

耶穌的顯現為世界帶來突破

l  因極度失望,選擇往以馬忤斯要離開的2個門徒,經歷新的眼光,認出了耶穌,生命得改變…

路加福音24

30到了坐席的時候,耶穌拿起餅來,祝謝了,擘開,遞給他們. 31他們的眼睛明亮了,這才認出他來.忽然耶穌不見了. 32他們彼此說:「在路上,他和我們說話,給我們講解聖經的時候,我們的心豈不是火熱的嗎?」 33他們就立時起身,回耶路撒冷去

l  耶穌首次顯現的對象,就是被當時看為沒有地位,見證沒有公信力,連其他門徒都不信的,但一心愛主的婦女們.

馬可福音16

10她去告訴那向來跟隨耶穌的人；那時他們正哀慟哭泣. 11他們聽見耶穌活了,被馬利亞看見,卻是不信.

l  那個門徒中的班長,曾3次不認主的彼得,回轉成為教會中最重要的領袖…(下主日經課)

l  個人選擇,限制,情緒的突破

l  對復活身體認知的突破

l  社會中對婦女(弱勢社群)接納的突破

l  門徒身份角色的突破,承接耶穌的使命,牧養主的羊

多馬與顯現的主相遇的故事為我們帶來的突破

29耶穌對他說:「你因看見了我才信；那沒有看見就信的有福了.」

l  不憑眼見的信心,帶來福氣

30耶穌在門徒面前另外行了許多神蹟,沒有記在這書上. 31但記這些事要叫你們信耶穌是基督,是上帝的兒子,並且叫你們信了他,就可以因他的名得生命.

l  我們因他的見證,從而信耶穌是基督,也因此得了生命.

耶穌來,站在當中說:「願你們平安!」

不要疑惑,總要信



\newpage

\section{岸邊的陌生人}
\label{sec:746}
\textbf{李均熊牧師 - 岸邊的陌生人}
\newline
\newline
\hyperref[sec:745]{< < < PREV SERMON < < <}
~
\hyperref[sec:toc]{[返目錄]}
~
\hyperref[sec:747]{> > > NEXT SERMON > > >}


約翰福音 21:1-19
\newline
\begin{longtable}{cl}
\hline
\hline
章節 & 經文 (和合本修訂版)\\
\hline
21:1 & \begin{tabularx}{0.7\textwidth}{X} 這些事以後，耶穌在提比哩亞海邊又向門徒顯現。他怎樣顯現記在下面。 \end{tabularx} \\ \\ \relax
21:2 & \begin{tabularx}{0.7\textwidth}{X} 西門‧彼得、叫低土馬的多馬、加利利的迦拿人拿但業、西庇太的兩個兒子，和另外兩個門徒，都在一起。 \end{tabularx} \\ \\ \relax
21:3 & \begin{tabularx}{0.7\textwidth}{X} 西門‧彼得對他們說：「我打魚去。」他們對他說：「我們也和你一起去。」他們就出去，上了船；那一夜並沒有打著甚麼。 \end{tabularx} \\ \\ \relax
21:4 & \begin{tabularx}{0.7\textwidth}{X} 天剛亮的時候，耶穌站在岸上，門徒卻不知道他是耶穌。 \end{tabularx} \\ \\ \relax
21:5 & \begin{tabularx}{0.7\textwidth}{X} 耶穌就對他們說：「孩子們！你們有吃的沒有？」他們回答他：「沒有。」 \end{tabularx} \\ \\ \relax
21:6 & \begin{tabularx}{0.7\textwidth}{X} 耶穌對他們說：「你們把網撒在船的右邊，就會得到。」於是他們撒下網去，竟拉不上來了，因為魚很多。 \end{tabularx} \\ \\ \relax
21:7 & \begin{tabularx}{0.7\textwidth}{X} 耶穌所愛的那門徒對彼得說：「是主！」那時西門‧彼得赤著身子，一聽見是主，就束上外衣，跳進海裡。 \end{tabularx} \\ \\ \relax
21:8 & \begin{tabularx}{0.7\textwidth}{X} 其餘的門徒因離岸不遠，約有二百肘，就坐著小船把那網魚拉過來。 \end{tabularx} \\ \\ \relax
21:9 & \begin{tabularx}{0.7\textwidth}{X} 他們上了岸，看見那裡有炭火，上面有魚和餅。 \end{tabularx} \\ \\ \relax
21:10 & \begin{tabularx}{0.7\textwidth}{X} 耶穌對他們說：「把剛才打的魚拿幾條來。」 \end{tabularx} \\ \\ \relax
21:11 & \begin{tabularx}{0.7\textwidth}{X} 西門‧彼得就上船，把網拉到岸上，網裡滿了大魚，共一百五十三條；雖然魚這樣多，網卻沒有破。 \end{tabularx} \\ \\ \relax
21:12 & \begin{tabularx}{0.7\textwidth}{X} 耶穌對他們說：「你們來吃早飯。」門徒中沒有一個敢問他：「你是誰？」因為他們知道他是主。 \end{tabularx} \\ \\ \relax
21:13 & \begin{tabularx}{0.7\textwidth}{X} 耶穌走過來，拿餅給他們，也照樣拿魚給他們。 \end{tabularx} \\ \\ \relax
21:14 & \begin{tabularx}{0.7\textwidth}{X} 耶穌從死人中復活後向門徒顯現，這是第三次。 \end{tabularx} \\ \\ \relax
21:15 & \begin{tabularx}{0.7\textwidth}{X} 他們吃完了早飯，耶穌對西門‧彼得說：「約翰的兒子西門，你愛我比這些更深嗎？」彼得對他說：「主啊，是的，你知道我愛你。」耶穌說：「你餵養我的小羊。」 \end{tabularx} \\ \\ \relax
21:16 & \begin{tabularx}{0.7\textwidth}{X} 耶穌第二次又對他說：「約翰的兒子西門，你愛我嗎？」彼得對他說：「主啊，是的，你知道我愛你。」耶穌說：「你牧養我的羊。」 \end{tabularx} \\ \\ \relax
21:17 & \begin{tabularx}{0.7\textwidth}{X} 耶穌第三次對他說：「約翰的兒子西門，你愛我嗎？」彼得因為耶穌第三次對他說「你愛我嗎」，就憂愁，對耶穌說：「主啊，你無所不知，你知道我愛你。」耶穌說：「你餵養我的羊。 \end{tabularx} \\ \\ \relax
21:18 & \begin{tabularx}{0.7\textwidth}{X} 我實實在在地告訴你，你年輕的時候，自己束上帶子，隨意往來；但年老的時候，你要伸出手來，別人要把你束上，帶你到不願意去的地方。」 \end{tabularx} \\ \\ \relax
21:19 & \begin{tabularx}{0.7\textwidth}{X} 耶穌說這話，是指彼得會怎樣死來榮耀神。說了這話，耶穌對他說：「你跟從我吧！」 \end{tabularx} \\ \\
[1ex]
\hline
\hline
\end{longtable}
講題:岸邊的陌生人 Stranger on the Shore

經文:約翰福音21:1-19

講員:李均熊牧師

場合:香港中文大學崇基學院禮拜堂主日崇拜

日期:2025年5月4日

今日是復活節後第三主日,福音經課就是關於耶穌基督復活後在海邊向門徒顯現.一如約翰福音經常運用古希臘劇場的相認場景(recognition scenes),七個門徒——彼得,多馬,拿但業,西庇太的兩個兒子,和其餘兩個門徒,第一眼與耶穌相遇都只是以為他是一個岸邊的陌生人,但在門徒聽從這個陌生人向船的右邊下網,從而打了一網拉不上船的魚獲後,「耶穌所愛的門徒」就認出這個岸邊的陌生人不是別人,就是他們跟隨了三年的主!彼得二話不說,就跳下海游上岸.(經文細節說彼得束上外衣可能只是方便游泳)但上到岸後,整個場景大家都不發一言.反而耶穌——這個岸邊的陌生人,不單殷勤地弄上炭火烤魚和餅,叫大家來吃早餐,又自然地拿餅和魚給門徒吃.沒有一個敢問這個陌生人是誰,因為大家都知道他就是主.接著記載將焦點放在這個大家都知道是誰的陌生人,和尷尷尬尬的彼得之間的對話.吃完了早餐,耶穌問了彼得三次「你愛我嗎?」,彼得答了三次「我愛你」,最後耶穌就呼召彼得「你跟從我吧!」(經文接著是彼得詢問本段另一個主角——「耶穌所愛的門徒」的事情,經課則沒有選用)

返教會一段時間的第兄姊妹對這一段記載應該耳熟能詳.約翰福音的最後一章,也是學者十分關注的研究對象.一來,明明在約20:30-31「耶穌在門徒面前另外行了許多神蹟,沒有記在這書上.但記這些事要叫你們信耶穌是基督,是神的兒子,並且叫你們信了他,就可以因他的名得生命.」已經清楚的寫下一段結語,但不知甚麼原因約翰福音又加入第21章；二來,整晚捕魚毫無所獲,到因為耶穌的說話而獲得特別的漁獲,這種耶穌與彼得相遇的記載,在之前我們也在路加福音5章也讀過類似情節,究竟路加福音和約翰福音相似記載的關係是甚麼呢?這都用了很多學者的筆墨去討論.

不過,講道不是研究報告或教書,所以我就將我看過又認為合理的,關於這段經文的要點給大家講講.也因為講道不是研究報告或教書,我今次用一個比較另類對話的形式和大家分享,可能不會直接明白地講,但大家發揮一下想像力,可能大家的領受又會十分不同.

以下是一段關於年老的「耶穌所愛的門徒」和這老使徒的一位門徒,叫約翰(就叫他小約翰／約翰仔)之間的設計對白,用現代廣東話,模擬一段推斷出來的對話,關於約翰福音21章的寫作情形(大家可以當廣播劇黎聽):

老使徒:(對他的門徒)小約翰,你地前陣時幫我紀錄我地主基督事蹟ge 福音書我已經睇jor ,寫得好好……但係我想最後加返一段關於主復活之後另一次向我地顯現ge 事情.

小約翰:好ge ,老師.不過,主復活後除jor 向D 婦女顯現,同個次向多馬老師顯現之外,仲有另一次向你地顯現咩?

老使徒:係呀.你咁講,呢次算係第三次啦.你有無讀過保羅老師個宣教拍檔路加醫生寫ge ,關於主ge 福音書.佢記載過有一件事,就係彼得同我地一班做漁夫ge 門徒,有一次成晚打魚打唔到,後黎主上jor 彼得架船講,完道之後,叫佢再打魚,打jor 好多魚gwo 件事 ,你有無印象?

小約翰:當然有. 路加醫生仲話個次個網魚,多到連魚網都差D 裂開嘛!

老使徒: 其實個網魚雖然多,但個網一D 都冇穿冇爛,完好無缺. 我仲記得後黎數過gwo 網魚有153條魚咁多.

小約翰:老師你又會記得咁清楚?

老使徒:因為153 正好就係由1加到17 ge 一個三角形數.就係個D 希臘數學家提出ge 三角形數.你地唔係對於呢D 算術問題好有興趣ge 咩?

小約翰:係喎!又幾得意喎!

老使徒:我相信係有D 意思ge .你自己慢慢再參透啦.不過我想同你講ge 係,路加醫生應該係第二度聽到呢件事蹟,跟住佢就將呢段事蹟加jor 入佢ge 福音書裡面,等個D 操希臘語ge 猶太信徒,同個D 外邦信徒明白大宗徒彼得最初點樣蒙召.不過成件事其實係主復活之後先發生ge :當個日,我同彼得見到主ge 空墳墓之後,又聽到D 婦女話見到主復活,我地其實都半信半疑,我諗就係主以前成日話我地信得太遲.之後,我地呢班原本住係加利利ge 門徒就返jor 加利利,所以當主顯現比細孖(低土馬)ge 時候,即係多馬老師,我地其實唔係耶路撒冷.所以你地寫到20章,之後記載jor D 婦女同細孖見到復活ge 主,但無交待見完空墳墓ge 我同彼得有無遇到復活ge 主,就完jor 成卷書,我覺得應該要補充一下,想加返呢件事,順便等其他人知道下路加醫生記載ge 事情其實係有另一個版本.

小約翰:咁路加醫生係為佢教會群體記載主基督ge 事蹟,同我地為你創立ge 教會而寫呢卷書ge 目的係有D 唔同ge .

老門徒:無錯,係彼得殉主之後,我就要大家幫我寫呢卷書,為ge 就係等我都返去見主ge 時候,可以話所有同主生活過ge 使徒都唔係度,咁對未見過主,但因為我ge 見證而信主ge 人,就可以透過睇呢卷書而領受上主ge 道,又明白跟從主ge 意思係乜.

小約翰:我地明白ge ,因為我地呢D 第二代相信主ge 人,無親眼見過主,但又要相信佢,係比較困難ge ,所以我地係20章裡面加jor 主對多馬老師講ge 個句說話:「你因看見了我才信；那沒有看見就信的有福了.」

老門徒:你地仲用jor 好多篇幅寫主同我地最後晚餐講個D ye ,其實都係等大家雖然見唔到我地呢班同主生活過ge 門徒,仍然能夠彼此相愛,仍然能夠經歷大家見唔到ge 主,因為主留下平安比我地,其實亦都代表佢留下平安比你地.講開最後晚餐,令我諗返我同你講緊個件事,當我地打完魚之後,主都係好似係最後一餐同我地一齊咁一齊食ye 傾談.

小約翰:其實點解個日你地會走去打魚呢?你地唔係跟從jor 主之後就放低哂D 打魚ge 生意呢咩?

老門徒:D 打魚ge 船同器材係我老竇ge .講返轉頭,我地離開jor 耶路撒冷,返jor 加利利ge 個段時間,唔知係因為太無聊,還是太傷心,定乜ye 原因,有一日彼得突然話要去打魚,好以想放棄再跟從主ge 門徒咁.我地一班原本做漁夫ge 門徒就話跟埋去,除jor 彼得,我同我啊哥,仲有係你地寫個卷書一開始講ge 拿但業,同埋arm arm 趕返加利利,你頭先講ge 多馬老師,仲有邊兩個門徒呢……呀,我都唔係好記得係邊個啦.總之我記得係七個人分jor 兩艘漁船.點知打jor 成晚,連魚毛都網唔到.當時天蒙光,天色仲係有D 黑,仲有D 霧添,係岸上面有個人突然問我地「細路!你地有無ye 可以食架」我地初時以為係出黎搵早餐ge 鄉親,見到我地打緊魚以為可以即時lor 幾條新鮮烤黎食,點知我地其實成晚無收獲.我聽到唔知邊個嬲嬲地應佢話:「無呀!」點知呢個岸邊ge 陌生人好聲好氣咁講:「你地將網撒向船ge 右邊,就會有架啦.」於是我地見都打jor 成晚,算係最後一次落網,於是轉向船ge 右邊撒網——點知,個下我地個個都嚇到呆jor !!

小約翰:發生jor 咩事?

老門徒:因為我地點樣拉,都拉唔返個網上船!我開初以為因為我地打jor 成晚魚,手軟腳軟,所以拉唔返起個網,後黎發現網裡面擠滿哂都係魚,係我地打咁耐魚都未見過ge 情影!於是乎我驚覺呢單ye 唔簡單,再向岸上面gwo 個陌生人望一眼,就認到gwo 個岸上ge 陌生人就係我地ge 主,於是就大叫起黎:係主呀!彼得個時將件衫綁jor 係條腰度,佢聽到我咁大叫「係主」,佢就解jor 件衫著返上身,就跳落水度游返上岸.我地其實離岸唔係好遠,約有二三百尺左右,於是我地就唔將個滿jor 魚ge 網拉上船,而直接拉佢上岸.後黎發現個網一D 破損都無,我就話好似我地教會咁相愛合一,唔似路加醫生話就黎要破,可能佢咁講係想提醒佢教會D 操唔同語言ge 信徒唔好攪分裂掛.等佢地仲話得人如得魚.

小約翰:老師,你又黎啦,保羅老師同路加醫生向外邦人傳福音,你就多D 向猶太人傳福音,但D 猶太人咪仲煩過D 外邦人.

老門徒:我知道ge ,你地寫個卷書都講到D 猶太人好唔明白我地ge 主.哈,哈,唔同使徒係有唔同ge 召命ge .

小約翰:咁,彼得老師呢?雖然佢已經為主殉道,但路加醫生講到好似彼得老師因為見到咁誇張ge 魚獲,而話自己係罪人,而主當時就叫佢跟從佢,咁你講個次事情又係咪咁?

老門徒:嗯…..(陷入深深ge 沉思)

小約翰:老師,老師,你點呀?

老門徒:我諗返起你地寫個卷書.你地由第13章寫彼得到20章完結ge 時候,寫到彼得好似講ye 做事成日唔經大腦咁,佢份人其實唔係咁ge .

小約翰:老師你又有無發覺,我地每次寫彼得老師出場,你都會同場出現架!而且我地係好尊敬大宗徒彼得ge ,但係寫成彼得係咁,一來係我地成日聽你講彼得老師係一個真誠直爽ge 人物,另方面係想提醒我地第二代ge 信徒,無論係教會位置係高還是低,都可以將真性情係主面前敝開.

老門徒:無錯,彼得係我地ge 領袖,並唔係佢特別有手段,而係佢對主ge 真誠感動我地,反而有時講ye 做事隨心而發,更加令我地覺得佢係亳無機心,係真心能夠代表我地,成為我地ge 發言人.反而加略人猶大……唔好提佢啦.講返我同你講主向我地顯現ge 事,我覺得復活主係特別為彼得而係岸邊向我地顯現.我成日諗,主其實復活之後可以直接升天見返佢天上ge 父,我成日講佢道成肉身,所以復活ge 佢可以即時升天,唔再受肉身ge 困鎖,唔再會有傷痕痛楚,亦唔再需要肚餓食早餐.但我地ge 主唔單只選擇道成肉身,係肉身復活之後亦選擇仍然留係世上四十日先升天,為ge 就係幫助我地係主受苦時同時所經歷ge 創傷.先係向一班婦女,細孖多馬,然後向我地幾個,特別係大宗徒彼得,為ge 就係幫助我地由創傷中走出黎,而明白主復活對我地每一個相信佢ge 人ge 意義.我仲記得,個次食完早餐之後,主就問jor 彼得三次「你愛我嗎」,不是正好提醒彼得佢曾經三次唔認主嗎?後黎彼得同我講返,個日如果唔係主好似一個陌生人咁突然出現,可能會走jor 去無面見主.其實主每次問彼得「你愛我嗎?」彼得都亦清楚回應話他是愛主ge ,他的心一直沒有改變過.主跟住三次同佢講類似ge 說話,「你餵養我的小羊.」「你牧養我的羊.」「你餵養我的羊.」表明主不但沒有計交彼得三次唔認佢,更係我地面前肯定佢要接任主作為群羊大牧人ge 職責.我諗主咁樣係我地面前問佢三次,又委任佢三次,不單止係恢復大宗徒彼得ge 信心,更加係恢復我地對佢作為我地領導角色ge 信心.你諗下,彼得佢曾經三次不認主,總會有D 人對佢作為大宗徒大牧者有一定懷疑.佢在生時未必有人會講D 乜,而家佢殉主,唔係我地中間,可能會有D 多事之徒或貪圖教會權力ge 人會講佢壞話,破壞佢為主建立ge 教會.我要你係我在生時記低呢件事,就係為jor 避免心直口快ge 彼得被人中傷詆譭.

小約翰:你雖然咁講,但我地幫你寫完呢卷福音書之後,大家係咪真係會咁諗?

老門徒:唔緊要,小約翰.你要記住,我作為主所愛ge 門徒,為這些事作見證,係等相信主ge 信徒知道,主復活之後唔係就完成jor 在世之功,升天走jor 去.主ge 復活同顯現係要我地無論咩境況,都知道佢會挽回我地,幫助我地,無論係我地犯錯,又或係創傷ge 低谷,佢都唔會唔理會我地.

小約翰:老師,我會再細心想下你見證ge 主復活顯現ge 深意架啦.

老門徒:唔……講jor 咁耐,我都有D gui 啦,其他ge 細節,過幾日我再同你講清楚.到時你再諗下點樣編修入去你地卷福音書啦.你要記得,無論遇到咩ye 困難,主已經復活,能夠按我地ge 需要,顯現,幫助我地,就好似佢向彼得顯現一樣.



\newpage

\section{誰為你哭泣誰為你抹淚}
\label{sec:747}
\textbf{林豪恩先生 - 誰為你哭泣?誰為你抹淚?}
\newline
\newline
\hyperref[sec:746]{< < < PREV SERMON < < <}
~
\hyperref[sec:toc]{[返目錄]}
~
\hyperref[sec:748]{> > > NEXT SERMON > > >}


使徒行傳 9:36-43
\newline
\begin{longtable}{cl}
\hline
\hline
章節 & 經文 (和合本修訂版)\\
\hline
9:36 & \begin{tabularx}{0.7\textwidth}{X} 在約帕有一個女門徒，名叫大比大，翻出來的意思是多加；她廣行善事，多施賙濟。 \end{tabularx} \\ \\ \relax
9:37 & \begin{tabularx}{0.7\textwidth}{X} 當時，她患病死了，有人把她清洗後，停在樓上。 \end{tabularx} \\ \\ \relax
9:38 & \begin{tabularx}{0.7\textwidth}{X} 呂大原與約帕相近；門徒聽見彼得在那裡，就派兩個人去見他，央求他說：「請快到我們那裡去，不要耽延。」 \end{tabularx} \\ \\ \relax
9:39 & \begin{tabularx}{0.7\textwidth}{X} 彼得就起身和他們同去。他到了，就有人領他上樓。眾寡婦都站在彼得旁邊哭，拿多加與她們同在時所做的內衣外衣給他看。 \end{tabularx} \\ \\ \relax
9:40 & \begin{tabularx}{0.7\textwidth}{X} 彼得叫她們都出去，然後跪下禱告，轉身對著屍體說：「大比大，起來！」她就睜開眼睛，看見彼得，就坐了起來。 \end{tabularx} \\ \\ \relax
9:41 & \begin{tabularx}{0.7\textwidth}{X} 彼得伸手扶她起來，叫那些聖徒和寡婦都進來，把多加活活地交給他們。 \end{tabularx} \\ \\ \relax
9:42 & \begin{tabularx}{0.7\textwidth}{X} 這事傳遍了約帕，就有許多人信了主。 \end{tabularx} \\ \\ \relax
9:43 & \begin{tabularx}{0.7\textwidth}{X} 此後，彼得在約帕一個皮革匠西門的家裡住了好些日子。 \end{tabularx} \\ \\
[1ex]
\hline
\hline
\end{longtable}
詩篇 23:1-6
\newline
\begin{longtable}{cl}
\hline
\hline
章節 & 經文 (和合本修訂版)\\
\hline
23:1 & \begin{tabularx}{0.7\textwidth}{X} 耶和華是我的牧者，我必不致缺乏。 \end{tabularx} \\ \\ \relax
23:2 & \begin{tabularx}{0.7\textwidth}{X} 他使我躺臥在青草地上，領我在可安歇的水邊。 \end{tabularx} \\ \\ \relax
23:3 & \begin{tabularx}{0.7\textwidth}{X} 他使我的靈魂甦醒，為自己的名引導我走義路。 \end{tabularx} \\ \\ \relax
23:4 & \begin{tabularx}{0.7\textwidth}{X} 我雖然行過死蔭的幽谷，也不怕遭害，因為你與我同在；你的杖、你的竿，都安慰我。 \end{tabularx} \\ \\ \relax
23:5 & \begin{tabularx}{0.7\textwidth}{X} 在我敵人面前，你為我擺設筵席；你用油膏了我的頭，使我的福杯滿溢。 \end{tabularx} \\ \\ \relax
23:6 & \begin{tabularx}{0.7\textwidth}{X} 我一生一世必有恩惠慈愛隨著我；我且要住在耶和華的殿中，直到永遠。 \end{tabularx} \\ \\
[1ex]
\hline
\hline
\end{longtable}
啟示錄 7:9-17
\newline
\begin{longtable}{cl}
\hline
\hline
章節 & 經文 (和合本修訂版)\\
\hline
7:9 & \begin{tabularx}{0.7\textwidth}{X} 此後，我觀看，看見有許多人，沒有人能計算，是從各邦國、各支派、各民族、各語言來的，站在寶座和羔羊面前，身穿白衣，手拿棕樹枝， \end{tabularx} \\ \\ \relax
7:10 & \begin{tabularx}{0.7\textwidth}{X} 大聲喊著說：「願救恩歸於坐在寶座上我們的神，也歸於羔羊！」 \end{tabularx} \\ \\ \relax
7:11 & \begin{tabularx}{0.7\textwidth}{X} 眾天使都站在寶座和眾長老，以及四個活物的周圍，俯伏在寶座前，敬拜神， \end{tabularx} \\ \\ \relax
7:12 & \begin{tabularx}{0.7\textwidth}{X} 說：「阿們！頌讚、榮耀、智慧、感謝、尊貴、權能、力量都歸於我們的神，直到永永遠遠。阿們！」 \end{tabularx} \\ \\ \relax
7:13 & \begin{tabularx}{0.7\textwidth}{X} 長老中有一位回應我說：「這些穿白衣的是誰？是從哪裡來的？」 \end{tabularx} \\ \\ \relax
7:14 & \begin{tabularx}{0.7\textwidth}{X} 我對他說：「我主啊，你是知道的。」他向我說：「這些人是從大患難中出來的，他們曾用羔羊的血把衣裳洗得潔白。 \end{tabularx} \\ \\ \relax
7:15 & \begin{tabularx}{0.7\textwidth}{X} 所以，他們在神寶座前，晝夜在他殿中事奉他；那坐在寶座上的要用帳幕覆庇他們。 \end{tabularx} \\ \\ \relax
7:16 & \begin{tabularx}{0.7\textwidth}{X} 他們不再飢，不再渴；太陽必不傷害他們，任何炎熱也不傷害他們， \end{tabularx} \\ \\ \relax
7:17 & \begin{tabularx}{0.7\textwidth}{X} 因為寶座中的羔羊必牧養他們，領他們到生命水的泉源；神必擦去他們一切的眼淚。」 \end{tabularx} \\ \\
[1ex]
\hline
\hline
\end{longtable}
約翰福音 10:22-30
\newline
\begin{longtable}{cl}
\hline
\hline
章節 & 經文 (和合本修訂版)\\
\hline
10:22 & \begin{tabularx}{0.7\textwidth}{X} 那時正是冬天，在耶路撒冷有獻殿節。 \end{tabularx} \\ \\ \relax
10:23 & \begin{tabularx}{0.7\textwidth}{X} 耶穌在聖殿裡的所羅門廊下行走。 \end{tabularx} \\ \\ \relax
10:24 & \begin{tabularx}{0.7\textwidth}{X} 猶太人圍著他，對他說：「你讓我們猶豫不定到幾時呢？你若是基督，就明白地告訴我們。」 \end{tabularx} \\ \\ \relax
10:25 & \begin{tabularx}{0.7\textwidth}{X} 耶穌回答他們：「我已經告訴你們，你們卻不信。我奉我父的名所行的事可以為我作見證。 \end{tabularx} \\ \\ \relax
10:26 & \begin{tabularx}{0.7\textwidth}{X} 但是你們不信，因為你們不是我的羊。 \end{tabularx} \\ \\ \relax
10:27 & \begin{tabularx}{0.7\textwidth}{X} 我的羊聽我的聲音，我認識牠們，牠們也跟從我。 \end{tabularx} \\ \\ \relax
10:28 & \begin{tabularx}{0.7\textwidth}{X} 並且，我賜給他們永生；他們永不滅亡，誰也不能從我手裡把他們奪去。 \end{tabularx} \\ \\ \relax
10:29 & \begin{tabularx}{0.7\textwidth}{X} 我父所賜給我的比萬有都大，誰也不能從我父手裡把他們奪去。 \end{tabularx} \\ \\ \relax
10:30 & \begin{tabularx}{0.7\textwidth}{X} 我與父原為一。」 \end{tabularx} \\ \\
[1ex]
\hline
\hline
\end{longtable}
講題:誰為你哭泣?誰為你抹淚?  Who Weeps for You? Who Wipes Your Tears?

經文:使徒行傳9:36-43；詩篇23:1-6；啟示錄7:9-17；約翰福音10:22-30

講員:林豪恩先生

場合:香港中文大學崇基學院禮拜堂主日崇拜

日期:2025年5月11日

一. 兩個復活的故事

今天是母親節,願上主祝福天下母親,給你們上好的福分.祝願天下母親因手中的工而心裡喜悦,並深深知道自己在上主眼中的是美好的.按教會年曆,今日是復活節第四主日,剛才誦讀經課的時候,有沒有留意到與復活相關的兩種狀況呢?

第一種復活狀況記載在經課中的使徒行傳9:36-43.

9:36在約帕有一個女徒,名叫大比大,翻希臘話就是多加；她廣行善事,多施賙濟.

9:37當時,她患病而死,有人把她洗了,停在樓上.

9:38呂大原與約帕相近；門徒聽見彼得在那裡,就打發兩個人去見他,央求他說:「快到我們那裡去,不要耽延.」

9:39彼得就起身和他們同去；到了,便有人領他上樓.眾寡婦都站在彼得旁邊哭,拿多加與她們同在時所做的裡衣外衣給他看.

9:40彼得叫她們都出去,就跪下禱告,轉身對著死人說:「大比大,起來!」她就睜開眼睛,見了彼得,便坐起來.

9:41彼得伸手扶她起來,叫眾聖徒和寡婦進去,把多加活活地交給他們.

這事件中大比大的復活,用廣東話來說,我們稱為「返生」.好像大比大一樣,死了能夠再次復生,不知道有沒有燃起你的想望呢?如果自己萬一死了,想不想「返生」呢?有沒有期望自己的家人萬一死了,可以「返生」呢?不過,這種「返生」是暫時的,「返生」之後的人,都總有死亡的一日.就算如此,如果真是有「返生」的機會,相信仍然有人渴求的,正如某首舊歌所唱:「再坐一會,但願能再坐一會」,能夠留得住多一刻,總好過立即失去.然而,「返生」在現實中絕無僅有,就算在聖經中都是罕見的特殊個案而不是常態.如果可以有死後立即「返生」的機會,最感到興奮的會是甚麼人呢?相信是那些南征北討的王帝,例如希臘的亞力山大大帝,蒙古的成吉思汗等,他們可以繼續帶兵打仗,征服世界.然而,如果這些南征北討的王帝死了又「返生」,對世界也未必是好事,因為他們會為世界帶來更多殺戮.好像大比大一樣死了又「返生」的復活狀況,保留為特殊個案而不普及化,可能對世界來說並不是壞事.

曾經聽過有人說,上帝幫助大比大「返生」,因為她「廣行善事,多施賙濟」.聽來都有道理,因為讓大比大「返生」,總好過讓一個殺人狂魔「返生」.然而,可能還有其他人都「「廣行善事,多施賙濟」,為甚麼不讓他們「返生」呢?因此,對於誰有返生的機會,誰沒有返生的機會,或許我們只能夠說:不知道.

相對於類似大比大的暫時「返生」特殊個案,經課中的啟示錄揭示了另一種復活的狀況:

7:9此後,我觀看,見有許多的人,沒有人能數過來,是從各國,各族,各民,各方來的,站在寶座和羔羊面前,身穿白衣,手拿棕樹枝,

7:10大聲喊著說:「願救恩歸與坐在寶座上我們的上帝,也歸與羔羊!」

7:11眾天使都站在寶座和眾長老並四活物的周圍,在寶座前,面伏於地,敬拜上帝,

7:12說:阿們!頌讚,榮耀,智慧,感謝,尊貴,權柄,大力都歸與我們的上帝,直到永永遠遠.阿們!

7:13長老中有一位問我說:「這些穿白衣的是誰?是從哪裡來的?」

7:14我對他說:「我主,你知道.」他向我說:「這些人是從大患難中出來的,曾用羔羊的血把衣裳洗白淨了.

經文中那些站在寶座和羔羊面前,身穿白衣,手拿棕樹枝,頌讚上帝的人,他們這種復活狀況應該不是暫時的,不會如「返生」一樣再死一次.另外,這種復活狀況也應該不是特殊個案,而是普遍的.啟示錄20:12-13節如此說:

20:12我又看見死了的人,無論大小,都站在寶座前.案卷展開了,並且另有一卷展開,就是生命冊.死了的人都憑著這些案卷所記載的,照他們所行的受審判.

20:13於是海交出其中的死人；死亡和陰間也交出其中的死人；他們都照各人所行的受審判.

二. 一個哭泣的故事

兩個復活的故事講完了,但今日的題目是「誰為你哭泣?誰為你抹淚?」復活那麼美好,怎會哭泣呢?怎會有眼淚呢?有沒有聽過保羅在帖撒羅尼迦前書4章如此勸導信徒:

4:13論到睡了的人,我們不願意弟兄們不知道,恐怕你們憂傷,像那些沒有指望的人一樣.

4:14我們若信耶穌死而復活了,那已經在耶穌裡睡了的人,上帝也必將他們與耶穌一同帶來.

4:17以後我們這活著還存留的人必和他們一同被提到雲裡,在空中與主相遇.這樣,我們就要和主永遠同在.

有復活,有盼望,不用憂傷,因此保羅在腓立比書1章如此表達自己對死亡的心聲:

1:21因我活著就是基督,我死了就有益處.

1:23我正在兩難之間,情願離世與基督同在,因為這是好得無比的.

如果好像保羅所說如此理性和瀟灑,或許真的不用哭泣,不用流淚.不過,大比大的「返生」故事中卻是充滿眼淚的.

9:39彼得就起身和他們同去；到了,便有人領他上樓.眾寡婦都站在彼得旁邊哭,拿多加與她們同在時所做的裡衣外衣給他看.

其實,在聖經記載中,很多類似的狀況,我們都聽到哭泣聲,看見流淚:

福音書記載耶穌叫管會堂的睚魯的女兒復活的故事,有以下片段:

5:38他們來到管會堂的家裡；耶穌看見那裡亂嚷,並有人大大地哭泣哀號(馬可福音5:38)

在拉撒路復活事件中,約翰福音11章如此記載人們和耶穌的反應:

11:33耶穌看見她哭,並看見與她同來的猶太人也哭,就心裡悲歎,又甚憂愁,

11:34便說:「你們把他安放在哪裡?」他們回答說:「請主來看.」

11:35耶穌哭了.

11:36猶太人就說:「你看他愛這人是何等懇切.」

大比大身邊的眾寡婦為甚麼哭呢?馬利亞和鄉親為甚麼哭呢?耶穌為甚麼哭呢?當場的猶太人講出了事實:「你看,他愛這人是何等懇切!」哭甚麼?有甚麼好哭呢?「係愛呀,哈利!」有人說,愛有多深,痛有多深.當你愛的某個人被帶走時,您會感受到情感上的痛苦,損失越大,悲傷就越強烈.有時縱使相隔多年,想起的時候,仍然會默默流下眼淚.記得讀中學的時候,坐我旁邊的同學常常唱張學友「遙遠的她」這首歌,第二段如此唱:

但這天收到她爸爸的一封信 信裡面說血癌已帶走她

但覺得空虛的心彷彿已僵化 過去事像炮彈心中爆炸

在這半山 這天我悲痛悲痛不已在胡亂說話

夜雨中 似聽到她說

不要相約 縱使分隔 相愛不會害怕

人無覓處 心聲有否偏差

正是讓這愛試出 真與假

遙遠的她 彷彿借風聲跟我話

熱情若無變 那管它滄桑變化

遙遠的她 不可以再歸家

我在夢裡卻始終祇有她

遙遠的她 可知我心中的說話

熱情並無變 那管它滄桑變化

常言道:「男兒有淚不輕彈,只因未到傷心處」.大比大身邊的婦女哭了,馬利亞哭了,耶穌也哭了.上主是有感情的,上主按祂的形象創造的人類都是有感情的,會悲傷,會與其他人同感共情,哭泣正是人性的表達.有一次,我陪一位親人去醫院見醫生.醫生說,報告顯示你的癌細胞已經擴散,是第四期,根源統計,一般還能存活六至十一個月.當下,每一位在場的家人都流下眼淚.患病者的第一句說話是:「我走jor ,你地算點呀?」如果耶穌在場,可能耶穌都會哭泣.「耶穌哭了」,是聖經對耶穌最親切,最合乎人性的描述.這是你熟識的耶穌形象嗎?

或許有人質疑,耶穌知道不久後拉撒路就會復活,有甚麼好哭呢?為甚麼祂不當場教導眾人說:「各位別哭,現在不是哭泣的時候,我在這裡,記得五餅二魚嗎?記得瞎子得見嗎?記得平靜風浪嗎?三十秒內我要把哀哭變為歡笑!」幸好耶穌不是如此「不吃人間煙火」,不會對人們的感受無動於衷.耶穌哭泣是自然的事,因為祂也是人,哭泣是祂人性的表現.

縱使使徒保羅告訴我們離世與基督同在好得無比,而我們也相信我們自己及親友最終都會復活,但在這之前的年日,令我們流下眼淚的事情仍然很多.為失去而流淚,為自己的痛苦而流淚,為別人的痛苦而流淚,為世間的不公不義而流淚.唐代詩人陳子昂為天下蒼生而流淚,他在「登幽州台歌」這詩如此說:

前不見古人,

後不見來者.

念天地之悠悠,

獨愴然而涕下.

三．一個抹眼淚的教事

人世間如此多眼淚,耶穌基督除了陪我們一同流淚,還可以怎樣呢?啟示錄經課中有抹眼淚這一幕:

7:17因為寶座中的羔羊必牧養他們,領他們到生命水的泉源；上帝也必擦去他們一切的眼淚.

抹眼淚代表甚麼呢?我曾經在特殊學校教書,還記得有一位學生在課堂上被旁邊的同學投訴騷擾,於是我把他的座位移開,搬到一角獨坐.當我繼續上課時,聽到他說:「林Sir,我哭呀,你唔理我呀?我哭,你都唔理我呀?」他直率可愛的表達真令我哭笑不得.為甚麼那麼多年之後我仍然記得這位學生這句說話呢?他教導了我,原來哭是想人「理」,也是需要有人「理」的.最近,我在報章讀到以下這篇報導,標題是「出租淚療師暖男為你抹眼淚」 (香港01):

「小時候哭泣,總有人呵護你,替你抹眼淚.但其實傷心時得到別人愛錫,絕對不是小朋友的專利.日本就有一個網上平台,專門提供一項有趣的服務,就是由帥氣的男生幫女生抹眼淚.女生可以在網上選擇以及預約你們喜歡的「淚療師」,他們就會親身來到你的地方提供服務.淚療師會陪你看傷心電影,待你想哭想宣洩情緒的時候,他們就會拿出手帕,溫柔地為你抹眼淚.」

聖經說,為你抹眼淚的是誰呢?正如詩歌「有一位神」所唱,是那位「有權能創造宇宙萬物,也有溫柔雙手安慰受傷靈魂」的上帝.有沒有留意今天經段經課都描述上帝是一位照顧我們的牧者呢?

詩篇23:1耶和華是我的牧者,我必不致缺乏.

約翰福音10:27我的羊聽我的聲音,我也認識他們,他們也跟著我.

在啟示錄中,這位牧者雖然自己成為被宰的羔羊,但仍然沒有放下牧養我們:

7:17因為寶座中的羔羊必牧養他們,領他們到生命水的泉源；上帝也必擦去他們一切的眼淚.」

正如詩篇23所說,在人生的旅途中,有青草地,溪水旁,但也有死蔭的幽谷；有朋友,但也有敵人.令人哭泣的情況多的是,連耶穌基督也哭了.在我們仍然會哭泣的日子,詩人說上帝「有恩惠慈愛隨著我」；耶穌基督保證「誰也不能從我手裡把他們奪去」,直至有一天,「上帝也必擦去他們一切的眼淚.」



\newpage

\section{我們獻的神不喜悅神喜悅的我們願獻嗎}
\label{sec:748}
\textbf{鄧瑞強博士 - 我們獻的,神不喜悅；神喜悅的,我們願獻嗎?}
\newline
\newline
\hyperref[sec:747]{< < < PREV SERMON < < <}
~
\hyperref[sec:toc]{[返目錄]}
~
\hyperref[sec:749]{> > > NEXT SERMON > > >}


彌迦書 4:1-4
\newline
\begin{longtable}{cl}
\hline
\hline
章節 & 經文 (和合本修訂版)\\
\hline
4:1 & \begin{tabularx}{0.7\textwidth}{X} 末後的日子，耶和華殿的山必堅立，超乎諸山，高舉過於萬嶺；萬民都要流歸這山。 \end{tabularx} \\ \\ \relax
4:2 & \begin{tabularx}{0.7\textwidth}{X} 必有許多民族前往，說：「來吧，我們登耶和華的山，到雅各神的殿。他必將他的道指教我們，我們也要行他的路。」因為教誨必出於錫安，耶和華的言語必出於耶路撒冷。 \end{tabularx} \\ \\ \relax
4:3 & \begin{tabularx}{0.7\textwidth}{X} 他必在許多民族中施行審判，為遠方強盛的國斷定是非。他們要將刀打成犁頭，把槍打成鐮刀。這國不舉刀攻擊那國，他們也不再學習戰事。 \end{tabularx} \\ \\ \relax
4:4 & \begin{tabularx}{0.7\textwidth}{X} 人人都要坐在自己的葡萄樹和無花果樹下，無人使他們驚嚇；這是萬軍之耶和華親口說的。 \end{tabularx} \\ \\
[1ex]
\hline
\hline
\end{longtable}
講題:我們獻的,神不喜悅；神喜悅的,我們願獻嗎?  O Mortal, The Lord Has Told You What Is Good

經文:彌迦書4:1-4;6:1-13

講員:鄧瑞強博士

場合:香港中文大學崇基學院禮拜堂主日崇拜

日期:2025年5月18日

各位弟兄姊妹,早安.願神的恩典常與你們同在.

今日,我們一起看彌迦書.

先知彌迦,與先知以賽亞同一時代,他們看著北國以色列被亞述毀滅,也親身經歷亞述兵臨耶路撒冷的威脅.彌迦和以賽亞,都意識到南國猶大的大禍將臨,這是先知的憂患意識,他們呼喚人民從腐敗生活中改變過來.先知都是體會到神的內心世界的人物,他們呼喚人,回歸神,過一「天人合一」的人生.當然,這「天人合一」不是中國哲學後來漸漸發展的,將「天」納入「人」的「天人合一」,而是呼喚人謙虛下來,順服天命,是「人」去回應「天」的「天人合一」.

先知以賽亞是朝廷裡的先知,面對猶大國的危機,他較同情大衛王朝及耶路撒冷的命運.先知彌迦,來自一個小城鎮摩利沙,他體會到有權有勢者對升斗巿民的欺壓,他較不同情權力的核心耶路撒冷.他甚至預告,將來的新君王將來自不顯眼的小城鎮伯利恆(彌5:2),這算是對耶路撒冷的權勢的一種批判.

彌迦書全書交織著審判與盼望,這也看到先知複雜的心情.一方面,他看到權貴對弱小者的欺壓,他不得不作出警告與審判.另一方面,他體會到神嚴厲審判下人的悲慘命運,他怕這審判令人絕望,故他又不時言說神審判後的解救,叫人在絕望中懷著希望.像一個極想孩子好的父親,不得不責罰孩子的錯,但又不斷地作出安慰.先知如此,正因天父也如此.

頑劣的人很難改變,猶大國劫數難逃.總有一天,猶大會亡國.現在,我們看看先知對將亡的猶大留下什麼說話,日後,劫後餘生的猶大人,就憑藉這些應許,找到有盼望的明天.

讀經:彌迦書4:1-4

彌4:1末後的日子,耶和華殿的山必堅立,超乎諸山,高舉過於萬嶺；萬民都要流歸這山.

彌4:2必有許多國的民前往,說:來吧,我們登耶和華的山,奔雅各上帝的殿.主必將他的道教訓我們；我們也要行他的路.因為訓誨必出於錫安；耶和華的言語必出於耶路撒冷.

彌4:3他必在多國的民中施行審判,為遠方強盛的國斷定是非.他們要將刀打成犁頭,把槍打成鐮刀.這國不舉刀攻擊那國；他們也不再學習戰事.〔這經文也出現在:以賽亞書2:4〕

彌4:4人人都要坐在自己葡萄樹下和無花果樹下,無人驚嚇.這是萬軍之耶和華親口說的.

大劫過後,不是世界末日,不必絕望,而是開啟我們去想像一個由神掌控的前景.我們想像,終有一天,萬民在耶和華的殿裡,聆聽神的真理.人與人的相處,不再是按暴力原則,而是順服神的指引.這不單是猶大人的遠象,更是世人的遠象.

接著,我們看一場法庭辯論.

讀經:彌迦書6:1-8

彌6:1以色列人哪,當聽耶和華的話!要起來向山嶺爭辯,使岡陵聽你的話.

彌6:2山嶺和地永久的根基啊,要聽耶和華爭辯的話!因為耶和華要與他的百姓爭辯,與以色列爭論.

彌6:3我的百姓啊,我向你做了什麼呢?我在什麼事上使你厭煩?你可以對我證明.

彌6:4我曾將你從埃及地領出來,從作奴僕之家救贖你；我也差遣摩西,亞倫,和米利暗在你前面行.

彌6:5我的百姓啊,你們當追念摩押王巴勒所設的謀和比珥的兒子巴蘭回答他的話,並你們從什亭到吉甲所遇見的事,好使你們知道耶和華公義的

作為.

彌6:6我朝見耶和華,在至高上帝面前跪拜,當獻上什麼呢?豈可獻一歲的牛犢為燔祭嗎?

彌6:7耶和華豈喜悅千千的公羊,或是萬萬的油河嗎?我豈可為自己的罪過獻我的長子嗎?為心中的罪惡獻我身所生的嗎?

彌6:8世人哪,耶和華已指示你何為善.他向你所要的是什麼呢?只要你行公義,好憐憫,存謙卑的心,與你的上帝同行.

6:1:「以色列人哪,當聽耶和華的話!」先知傳召神的子民到法庭聽審.「要起來向山嶺爭辯,使岡陵聽你的話.」「爭辯」是法庭用語,神的子民要答辯.山嶺與岡陵,是人與神立盟約時的見證人.類似我們的「山盟海誓」,山與海,這些恆久的東西,見證著立約者在歷史的流逝中是否守約.

6:2:「山嶺和地永久的根基啊,要聽耶和華爭辯的話!因為耶和華要與他的百姓爭辯,與以色列爭論.」「爭辯」這個法庭用字,一而再,再而三的出現,神要祂的子民答辯.

6:3-5:神的發言,神的控訴.

彌6:3我的百姓啊,我向你做了什麼呢?我在什麼事上使你厭煩?你可以對我證明.

彌6:4我曾將你從埃及地領出來,從作奴僕之家救贖你；我也差遣摩西,亞倫,和米利暗在你前面行.

彌6:5我的百姓啊,你們當追念摩押王巴勒所設的謀和比珥的兒子巴蘭回答他的話,並你們從什亭到吉甲所遇見的事,好使你們知道耶和華公義的作為.

留意這裡的用字:「我的百姓啊」.神要起訴祂的子民違約,但祂仍然用一種互有關係的稱呼.在彌迦書第2-3章,神譴責權貴的貪得無厭和對弱者的欺壓,在這裡,神直指問題的核心,是人對上天的恩情毫無感覺,沒有半點感激之心.人對上天寡情,就會對人無義.毫無感恩之心,就不懂得餽贈.不懂分享,就只會掠奪.

神的控訴問了兩個問題:我向你做了什麼呢?我在什麼事上使你厭煩?

神在問:我做錯了什麼,你們這樣待我?神像是很委屈似的.

然後,神講述祂做過的事情,祂曾救贖祂的子民離開埃及,讓他們得享自由.「出埃及」是舊約常常提及的意象,這不是神的子民自己能做到的事情,也不是他們做了什麼而配得的賞賜,這完全是神對弱者的眷顧,是神的恩典,是愛的禮物,讓這群奴隸的悲慘命運得以改變,獲得他們想也不曾想過的自由.這一切,出自神主動的恩情.

神繼續講過去的故事:在神的子民出了埃及,邁向迦南地的途中,邪惡的摩押王買通不敬虔的先知巴蘭,去咒詛他們.然後,又設美人計,引誘神的子民在什亭犯姦淫,拜偶像.幸得一位敬畏神的人「非尼哈」出手阻止,他們才免於自我毀滅.

神訴說這些歷史,是讓神的子民看清楚,他們的人生其實出於神的恩情.但他們卻像什亭那些背棄上帝的人一樣,忘恩負義,為了當下的快樂,將神忘得一乾二淨.

在這法庭裡,神問:是我做錯了什麼嗎?這問題成了反問:你們真的不知自己做錯了什麼嗎?我們看看神的子民的答辯.

彌6:6我朝見耶和華,在至高上帝面前跪拜,當獻上什麼呢?豈可獻一歲的牛犢為燔祭嗎?

彌6:7耶和華豈喜悅千千的公羊,或是萬萬的油河嗎?我豈可為自己的罪過獻我的長子嗎?為心中的罪惡獻我身所生的嗎?

這是神的子民的答辯,聽起來,帶點脾氣似的(廣東話叫「發晦氣」).

神的子民看來是問:「那麼,我們當做什麼,去回報神的恩情?」

他們想到能做的事,便是獻些什麼給神.接著,他們說得很誇張.獻最好的牛嗎?獻千千的公羊,或是萬萬的油河嗎?看來是說:「神啊,你說你待我們這麼好,那麼,我們獻上這麼多,夠回報你了嗎?」

神的用字,其實是很親切的:「我的百姓」.但祂的子民的用字,卻強調兩者之間的距離:「至高上帝」.跟著的講法,不知是否刻意頂撞神.「神啊,你這麼崇高,在你看來,我只是罪人,為了贖罪,我是否需要獻我的長子呢?」

在我聽來,他們這樣的回覆,顯得不懂敬畏,不懂感恩,甚至將恩典的關係變成一場買賣.他們竟然講,用殺長子的方式獻祭,來贖罪,這根本完全違反神一直以來禁止獻人為祭的命令.神言說恩典,他們的回應卻完全是反恩典.

神喜悅什麼呢?

彌6:8世人哪,耶和華已指示你何為善.他向你所要的是什麼呢?只要你行公義,好憐憫,存謙卑的心,與你的上帝同行.

我們看到,神喜悅的,與神的子民提議的,相差十萬八千里.他們看來什麼也願意獻,但就是獻不出神喜悅的東西.他們表面敬虔,實質是褻瀆.不單不能親近神,反更招來神的審判.他們沒有感謝神的恩慈,反顯出他們內心不擇手段的殘忍.

神要的,是世人行公義.在舊約,有兩個希伯來文都會譯作「公義」,字義相關.第5節那裡提到「耶和華公義的作為」,那個「公義」的希伯來文是「tsedaqah」,傾向指「道德上的正直」.舊約講的「義人」,一般都是用這個字.第8節這裡用的字,是「mishpat」,傾向指「法律上的公道」.你買東西,人家找多了錢給你,你付少了,就應該補回給人.十元的東西,就應付十元,公公道道.有次,在機場,取行李,發現行李箱破損了.在場的工作人員問我,這行李箱原價多少錢?用了多久?我如實告之,沒有作大,沒有屈人,那人按我的報價折舊後賠償給我.沒有騙人,公公道道.

「行公義」之外,神也希望我們「好憐憫」.「憐憫」這個字(chesed),在舊約,很多時用於指稱神在約裡信實的,恆久的愛.故此,這字的意思不純粹指「仁慈」,更著重堅忍的委身.你要堅持公道,以愛待人,當人家懂得感恩時如此,當人家不懂得感恩,甚或令自己有所損失時,也堅持如此.

「存謙卑的心,與你的上帝同行」.什麼是「存謙卑的心」呢?「謙卑」這個希伯來文,在舊約只出現一次,其意思不能完全確定,有多種可能.其意思可能是:謙卑,明智,慎重,時刻準備好.不管哪個意思,都指向奴僕的態度.一個忠誠的奴僕,時刻準備好回應主人的吩咐.他放下了自我,完全與主人配合.一個人若完全聽命於神,就真是一「天人合一」的生命,上天的旨意在其生命中實現出來.

在我們的人生背後,是上天無盡的恩典,但我們好像不能領會,總是以交易的心態去對待神,以自利的方式去對待他人.世人離神好遠.

神要求的人生,已講得很清楚了.但,看來依然故我的人仍是居多.最後,我們看看神的警告,作為我們生命的提示.

彌6:9耶和華向這城呼叫,智慧人必敬畏他的名.你們當聽是誰派定刑杖的懲罰.

彌6:10惡人家中不仍有非義之財和可惡的小升斗嗎?

彌6:11我若用不公道的天平和囊中詭詐的法碼,豈可算為清潔呢?

彌6:12城裡的富戶滿行強暴；其中的居民也說謊言,口中的舌頭是詭詐的.

彌6:13因此,我擊打你,使你的傷痕甚重,使你因你的罪惡荒涼.

但願我們行事公道,就算受虧損仍堅持仁愛,時刻留心神的心意,並忠誠地實行.

願神祝福你們.

願榮耀歸予聖父,聖子,聖靈.



\newpage

\section{詩篇一百二十七篇的隱藏信息}
\label{sec:749}
\textbf{郭志榮牧師 - 《詩篇》一百二十七篇的隱藏信息}
\newline
\newline
\hyperref[sec:748]{< < < PREV SERMON < < <}
~
\hyperref[sec:toc]{[返目錄]}
~
\hyperref[sec:750]{> > > NEXT SERMON > > >}


詩篇 127:1-5
\newline
\begin{longtable}{cl}
\hline
\hline
章節 & 經文 (和合本修訂版)\\
\hline
127:1 & \begin{tabularx}{0.7\textwidth}{X} 若不是耶和華建造房屋，建造的人就枉然勞力；若不是耶和華看守城池，看守的人枉然警醒。 \end{tabularx} \\ \\ \relax
127:2 & \begin{tabularx}{0.7\textwidth}{X} 你們清晨早起，夜晚安歇，吃勞碌得來的飯，本是枉然；惟有耶和華所親愛的，必叫他安然睡覺。 \end{tabularx} \\ \\ \relax
127:3 & \begin{tabularx}{0.7\textwidth}{X} 看哪，兒女是耶和華所賜的產業，所懷的胎是他所給的賞賜。 \end{tabularx} \\ \\ \relax
127:4 & \begin{tabularx}{0.7\textwidth}{X} 人在年輕時生的兒女好像勇士手中的箭。 \end{tabularx} \\ \\ \relax
127:5 & \begin{tabularx}{0.7\textwidth}{X} 箭袋充滿的人有福了！他們在城門口和仇敵爭論時必不蒙羞。 \end{tabularx} \\ \\
[1ex]
\hline
\hline
\end{longtable}
講題:《詩篇》一百二十七篇的隱藏信息  Hidden Message of Psalm 127

經文:詩篇127:1-5

講員:郭志榮牧師

場合:香港中文大學崇基學院禮拜堂主日崇拜

日期:2025年5月25日

⼀．前⾔

弟兄姊妹,主內平安!

今天要跟⼤家宣講⼀篇⼤家很熟悉的詩篇,就是《詩篇》⼀百⼆⼗七篇.

這篇詩篇當中有很多⼤家常常獨立地引⽤的⾦句,卻是很少會把這篇詩篇當作⼀個整體來研讀.但只有把這篇詩篇視為⼀個整體,⼤家才可以把當中隱藏著的信息解讀出來,⽽這信息亦正是對⾝處在今天這個⼤時代中的信徒⼀個非常重要的提醒.

經文已經顯⽰在電幕上,讓我們⼀起誦讀.

【經文:詩⼀⼆七 1-5 】

1 若不是耶和華建造房屋,建造的⼈就枉然勞⼒；若不是耶和華看守城池,看守的⼈就枉然警醒.

2 你們清晨早起,夜晚安歇,吃勞碌得來的飯,本是枉然；惟有耶和華所親愛的,必叫他安然睡覺.

3  兒女是耶和華所賜的產業；所懷的胎是他所給的賞賜.

4  少年時所⽣的兒女好像勇⼠⼿中的箭.

5  箭袋充滿的⼈便為有福；他們在城門口和仇敵說話的時候,必不⾄於羞愧.

⼆．有關《詩篇》⼀百⼆⼗七篇的⼀些疑問

對這篇詩篇到底是要講甚麼,詩⼈正⾯對⼀個怎樣的處境,學者有很多的討論.這篇詩篇⼀開始提到「建造房屋」,「房屋」是指甚麼?⽽「建造房屋」和「看守城池」⼜有甚麼關係?經文跟著⼜提到「吃飯」和「睡覺」.詩⼈所關注的,是否就是住屋,吃喝和睡覺所代表的基本⽣活需要呢?就是這麼簡單?

這篇詩篇的下半部分(即 3-5 節)似乎在談論關於兒女的主題,這跟前⾯的經文⼜有甚麼關係?詩⼈是否希望解釋,住屋,吃喝和睡覺等基本生活需要以及⽣兒育女之間的關係呢?

如果 3-5 節單純是關於兒女的教導,這篇詩篇最後講到「在城⾨⼝和仇敵說話」(5節),究竟想表達甚麼意思呢?難道詩⼈想提醒⽗⺟要自⼩教導自⼰的兒女如何跟別⼈「嗌交」時不會「輸蝕」呢?當然不是這樣!有⼀點值得留意:第 5 節中的「在城門口」,正好與詩篇⼀開始時提到的「看守城池」(1 節)⾸尾呼應.這是⼀個重要的提⽰,說明這篇詩篇上下兩個部分有緊密的關係.

釋經學者對這篇詩篇應該如何解讀有很多的討論.有些學者斷⾔,這篇詩篇本來是兩篇獨立的詩篇,後來才被之後的⼀位編修者放在⼀起,並加以修訂成為現在的文本.1情形就好像在敬拜讚美的聚會中,主席領詩把兩⾸詩歌結合起來⼀起頌唱那樣.從上述的討論看,這並非不可能.但問題是:即使這篇詩篇本來包含了兩篇獨立的詩篇,為何後來的編修者要把它們放在⼀起?

《聖經》是神所默⽰.換⾔之,各個經卷不同時代,不同背景的作者在寫作時得到神的啟dik .不僅如此,神也會感動後來收集和編修經卷的⼈.《聖經》各經卷的「後期製作」本⾝就是神默⽰的⼀部分,經文的編修本⾝往往就有神的信息包含在其中.

所以,如果這篇詩篇真的由兩篇本來獨立的詩篇經過編輯⽽組合起來,⼤家就更加要追問和探究:負責「後期製作」的那位「詩⼈」正在⾯對甚麼⽣活處境?「詩⼈」想藉著這篇複合的詩篇回應甚麼信仰的挑戰呢?

三．釋經

A.    詩⼀⼆七 1-2

這篇《詩篇》頭兩節經文是:「若不是耶和華建造房屋,建造的⼈就枉然勞⼒；若不是耶和華看守城池,看守的⼈就枉然警醒.你們清晨早起,夜晚安歇,吃勞碌得來的飯,本是枉然；惟有耶和華所親愛的,必叫他安然睡覺.」(詩⼀⼆七 1-2)

1.     三個「枉然」

詩⼈⼀開始重覆⽤了三個「枉然」,即「空泛」之意,也就是「無實質」,「無價值」,「無作為」,2 來強烈表達⼀個重要的事實:若不是神作主導,即按照著神所定的⽬標和⽅向,依循神的策略和途徑,跟隨著神⼼⽬中的時間表,並倚靠著神的介入和幫助,就算⼈天未光就起來⼯作,⼀直做到夜深也不休息(「夜晚安歇」直譯

「延遲坐下來」),⼀切「建造房屋」和「看守城池」的努⼒都是徒然,都是毫無意義,只能換來「吃悲傷的麵包」(「bread of grief」3,「勞碌」即「悲傷」之意4)的結果,意即帶來⾝體上的痛苦和⼼靈裡的憂傷.5

2.     建造「房屋」?

「房屋」這個字可解作⼈的「居所」或「家庭」,⼜或者國王的「王宮」,但同⼀個字亦可以⽤來指稱神的「聖殿」.6  ⽽在這篇詩篇中最合理的解釋,當然就是指「聖殿」.所以,「建造房屋」所指的,其實是「建造聖殿」,⽽且是指「第⼆聖殿」的重建.7

補充少少背景資料:「第⼀聖殿」就是指那座由⼤衛王朝第⼆位國王所羅⾨所建成的聖殿,這聖殿在主前 587 年被入侵並消滅猶⼤國的巴比倫軍隊拆毀；⽽「第⼆聖殿」就是指那些被擄到巴比倫的猶⼤國⺠,後來返回耶路撒冷,並在主前 515 年完成重建的新聖殿.8

⼀位猶太聖經學者分析了《詩篇》⼀百⼆⼗七篇,並把當中的內容與《以斯拉記》和《哈該書》的描述作對比後,提出⼀個相當有說服⼒的解釋,就是把《詩篇》⼀百⼆⼗七的背景,設定於第⼀批被擄歸回耶路撒冷猶太⼈的時代,也就是《以斯拉記》和《哈該書》所講述的時期.9

循著這個解釋的⽅向,⼤家便會發現,《詩篇》⼀百⼆⼗七篇應該跟它的上⼀篇⼀起閱讀,就好像講述同⼀個故事的「上下集」那樣.

《詩篇》⼀百⼆⼗六篇形容,當這⼀批「錫安被擄的⼈」離開巴比倫⽽回到耶路撒冷的時候,他們就「好像做夢的⼈」,因為神「果然為他們⾏了⼤事」(詩⼀⼆六 1, 3).10 這些「被擄歸回」的猶太⼈經歷到⼀個偉⼤的神蹟:沒有⼈可以想像到,當時的「超級⼤國」巴比倫竟然會⼀下⼦就崩潰滅亡；更沒有⼈會想到,取⽽代之的波斯會推出新的政策,允許和⿎勵被擄到巴比倫的猶太族群歸回故⼟,重建聖殿這個猶太⼈先祖昔⽇敬拜耶和華的中⼼.

然⽽,這批充滿熱誠與盼望,希望要為神和為自⼰的⺠族幹⼀番⼤事的猶太⼈,很快就發覺要從耶路撒冷⼀片頹垣敗⽡中重建聖殿並非那麼輕易和簡單.這批猶太⼈雖然很快就「立定了聖殿的根基」(拉三 11),但卻因為居住在耶路撒冷⼀帶和周邊的⺠族(特別是撒⾺利亞⼈)的阻撓,重建聖殿的⼯程⼀直無法開展.這批猶太⼈最後幾經波折下,在⼗四年後才有機會重啟⼯程!下⾯⼤家會看到,《詩篇》⼀百⼆⼗七篇正好回應這批歸回耶路撒冷,在重建聖殿時遇到攔阻的猶太⼈的處境.

3.     詩⼀⼆七 1-2 節回應甚麼問題?

所以,這篇詩篇的 1-2 節,正是要提醒第⼀批被擄歸回耶路撒冷的猶太⼈:若非神作主導,⼈在重建聖殿和守衛聖城所付出的⼀切,都是「枉然」!

問題是:詩⼈為何要這麼說呢?要明⽩這⼀點,⼤家就要回到這批被擄歸回的猶太⼈所⾯對的處境去瞭解.這批猶太⼈在戰爭中遭到破壞,破落不堪的耶路撒冷⽣活,⽇⼦⼀點也不好過.不僅物質上匱乏,還要擔驚受怕,因為在耶路撒冷⼀帶和周邊地區的居⺠對這批猶太⼈充滿敵意,不時會襲擊他們,造成極⼤的威脅.

這批猶太⼈的使命是要重建聖殿,但⼯程的開展卻遙遙無期.可以想像,這批猶太⼈當中可能已有⼈因為意興闌珊,感到灰⼼失望⽽返回巴比達去.即使是留下來的猶太⼈,也可能充滿疑惑.

這些猶太⼈很可能會問:神為甚麼不施恩幫助他們?他們甚⾄可能會聯想到,是否當初有甚麼做錯了?是否不應該跟原來居住在那裡和周邊的居⺠交惡?《以斯拉記》第四章記載了,當猶太⼈開始要重建聖殿的時候,撒⾺利亞⼈說他們也想⼀起參與,但卻被猶太⼈領袖斷然拒絕.自此以後,撒⾺利亞⼈就與其他外邦⼈合⼒阻撓聖殿的重建⼯程.

這個決定是正確的,因為這些撒⾺利亞⼈根本就不懷好意,他們的信仰也不純正.⼤家或許可以作這樣的類比:當年神學院興建這個新聖堂時,資⾦不⾜,假若摩⾨教或耶和華⾒證⼈會提議不如⼀⿑合作,相信我們亦⼀定不會接受.

詩⼈在這裡正正就是要回應有這些疑惑的猶太⼈,「重建聖殿」和「看守聖城」時若不是按著神的⼼意⽽⾏,⽽跟這些撒⾺利亞⼈妥協,就算最終能把聖殿和聖城的⼯程加快完成,也是枉然,因為這根本不是神的⼼意!神的⼼意並非只是重建聖殿和聖城的硬件,更重要是重建⼀個「信仰純正」,是神「所親愛」的屬神群體!

讓神作主導,不⾛世界之⼦的捷徑,的確可能要付出更⼤的代價.然⽽,這是令事奉神的⼈不會枉然的唯⼀選擇.

B.    詩⼀⼆七 3-5:

經文的下半部分 3-5 節:「兒女是耶和華所賜的產業；所懷的胎是他所給的賞賜.少年時所⽣的兒女好像勇⼠⼿中的箭.箭袋充滿的⼈便為有福；他們在城⾨⼝和仇敵說話的時候,必不⾄於羞愧.」

⼤家可以設⾝處地,代入當時這班猶太⼈的處境,當聽到詩⼈這個的回應後,⼤家會有甚麼想法或感受?⼤家很可能也不能否認,詩⼈這麼說當然是很有道理的,但⼼中仍然會有疑問:那麼,當下是否只有等待,等待和等待呢?眼⾒⽇⼦⼀年⼜⼀年地過去,當下除了「躺平」,等候神出⼿以外,還有甚麼可以做或者應該做呢?

《詩篇》⼀百⼆⼗七的下半部分,即 3-5 節,就正是要回應這個問題.

「兒女是耶和華所賜的產業」中,原文應直譯為「兒女是耶和華的產業」.11 「兒女」是屬於神,不是屬於⽗⺟.⽽在舊約聖經中,「產業」這個字通常⽤來指「地⼟」,⽽並非指別的物質財產.12

《和合本》譯作「所懷的胎是他所給的賞賜」,是⼀個非常誤導的翻譯.「所懷的胎」直譯「肚腹的果⼦」,⽽「肚腹」原文就是指⼀個⼈的腹部或肚腩,13 在這裡是指⼀ 個⼈五臟六腑的總和,⽽並非⺟親的⼦宮.舊約聖經把⼈的五臟六腑視為掌管著⼀個

⼈內在⽣命的器管,故此「肚腹」正是指⼀個⼈的「inner self」,是這個⼈的「⼈格 我」,也就是這個⼈「知,情,意,良知」等等的總和.箴⾔廿⼆ 18「你若⼼中存記」中的「⼼」,原文就是「肚腹」這個字.14 因此,有拉比就解釋,「肚腹的果⼦」所指的,就是⼈內在⽣命中所結出的果⼦.15

把第 3 節上下兩半的經文⼀起看,就可以看到經文正要表達:神將兒女交託給⽗⺟,就是要⽗⺟好好開墾這塊「地⼟」,使這塊「地⼟」變得肥沃,使他們的內在⽣命成為為神多結果⼦的肥沃⼟壤.「肚腹」指的是兒女的「內在⽣命」,⽽不是指孕婦的腹部.

⼤家會問,這跟被擄歸回的猶太⼈的問題有甚麼關係呢?

答案很簡單:詩⼈就是要告訴當時那些感覺到自⼰好像是⼀籌莫展的猶太⼈,讓他們知道自⼰絕對不是沒有作為,絕對不只有「躺平」,被動地等候神出⼿.神要他們好好栽培他們的下⼀代,就是他們「少年時所⽣的兒女」,琢磨他們,使他們成為神這位勇⼠⼿中隨時可以動⽤的利箭,讓神在與仇敵的爭戰中使⽤.

四．引伸應⽤

A.    在逆境中的堅持

詩⼈所⾯對的處境,正是⼀個充滿挑戰的⼤時代.詩⼈和他同代的⼈所⾯對的難關甚⾄是困局,超過了他們的能⼒範圍以外.他們有⼈選擇離開,放棄了自⼰的信仰和使命.

留下來的⼈也要⾯對是否與敵⼈妥協的兩難抉擇.就好像當時的猶太⼈那樣,向撒⾺利亞⼈妥協,就會偏離信仰的真道.不妥協就要付出代價,要⾯對那些看似無法逾越的困難.從《以斯拉記》和《尼希米記》的記載可以看到,當這些留下來猶太⼈堅持他們的信念,他們每每就會經歷到神在歷史中的奇妙作為,震撼他們的⼼靈,改變他們的⽣命與⼈⽣.

我之前曾經跟我堂會的弟兄姊妹⼀起查考《以斯拉記》和《尼希米記》,⼤家都好有感受,看到這批猶太⼈怎樣在神的帶領和幫助下,「難關難過」卻是「關關過」.⼤家都可以看到,整個經歷根本就是⼀個重整和重塑以⾊列這個信仰群體的必要歷程.

今天⼤家或許也同樣處於⼀個充滿挑戰的⼤時代,社會以⾄整個世界的變遷,都衝擊著⼤家的信⼼.世界之⼦總是想向信徒提供⼀個跟隨神以外的捷徑,並作出應許,只要⼤家跟隨世界之⼦的⽅法,⼤家就不⽤折騰,不⽤多付代價,就可以開⼼快樂的過⽇⼦.詩⼈提醒⼤家,若不是跟從神所主導的那個計劃,活出與世界之⼦似是⽽非相反的⽣命樣式,⼀切都是枉然.

耶穌就是神成了⾁⾝的聖道,是神所差遣的彌賽亞,是萬眾期待的救贖主,要把神的國度建立起來.當時⼤部分的猶太⼈都期望耶穌會差遣祂的天軍天將,帶領猶太⼈起義,推翻羅⾺帝國的統治,復興以⾊列國.猶太⼈⼤概認為,這樣,天國便可以輕⽽易舉⾺上在地上完全實現.但這卻不是神的⼼意,不是神救贖⼈的計劃.地上的世俗權柄使⽤武⼒,通過殺戮別⼈來建立他們的國度,擴張他們的疆界.但耶穌沒有這樣做,祂反⽽親自⾛上⼗字架,成為被宰殺的羔⽺,粉碎了地上⼀切邪惡勢⼒背後那個那惡者的權勢.就是這樣,地上的國度以威嚇,以殺戮的⼿段來建立和維持,神的國度卻是以⼗字架,也就是神⾏動的慈愛來展開.

B.    現在能做甚麼?

活在⼀個充滿挑戰的⼤時代中,⼤家有時會感到非常迷惘,好像⼀籌莫展,在困局中好像只能等待,等待和等待.

詩⼈向他同代的⼈發出提醒和挑戰,不是這樣!他們的使命,就是要好好建立他們這個信仰群體,特別是要栽培好他們信仰群體的下⼀代.時機來到,這個信仰群體的下⼀代就會被神使⽤,參與在神所主導的那個屬靈爭戰當中.

詩⼈這個提醒和挑戰,是否也正正是這篇詩篇要向今天信徒所發出的信息?

最近曾經有⼀位老朋友跟我分享,覺得自⼰的事奉,甚⾄是整個堂會都好似失去了⽅向.這或許也是不少教牧或信徒領袖的⼼聲,⽽我好幾年前在⼼中也曾經有過類似的掙扎.我當時因為預備⼀個祈禱會,於是細⼼研讀這篇詩篇,⽽主聖靈就藉著當中的信息來啟發我和激勵我,提醒我教會的使命從來沒有改變.

教會的使命從來都是為了「造⼈」⽽不是「做事」!「造⼈」的「造」字,是「造 就」的「造」,也是「塑造」的「造」.意思就是:事⼯的發展從來不是焦點,最要緊的是建立⼈的⽣命.換⾔之,事⼯做得好是 bonus,是額外的獎賞.但即使事⼯做得不好,其實也沒有甚麼⼤不了.要建立⼈的⽣命,就是要使⼈真真正正認識三⼀真神這位⽣命的源頭,基礎和歸宿,因⽽知道和明⽩自⼰是從哪裡來,要往哪裡去,於是就能搞清楚自⼰⼈⽣在世的召命,要成為降服在神主權之下的僕⼈,要成為神⼿中隨時可⽤的利箭.

這⼀點正好與詩⼈對他同代猶太⼈的教導不謀⽽合.

1 Leslie C. Allen, Psalms 101–150 (Revised), vol. 21, Word Biblical Commentary (Dallas: Word, Incorporated, 2002), 237.

2 Victor P. Hamilton, “2338 א ,” in Theological Wordbook of the Old Testament, ed. R. Laird Harris, Gleason L. Archer Jr., and Bruce K. Waltke (Chicago: Moody Press, 1999), 908.

3 參 YLT 的翻譯.

4 Ronald B. Allen, “1666 ַָצבﬠ ,” in Theological Wordbook of the Old Testament, ed. R. Laird Harris, Gleason L. Archer Jr., and Bruce K. Waltke (Chicago: Moody, 1999), 687–688.

5 同上.

ַ

6 Louis Goldberg, “241 יתבִּ ,” in Theological Wordbook of the Old Testament, ed. R. Laird Harris, Gleason L. Archer Jr., and Bruce K. Waltke (Chicago: Moody, 1999), 105.

7 Eliyahu Assis, “Psalm 127 and the Polemic of the Rebuilding of the Temple in the Post Exilic Period,” Zeitschrift Für Die Alttestamentliche Wissenschaft 121, no. 2 (2009): 256-72.

8 Carol Meyers, “Temple, Jerusalem,” in The Anchor Yale Bible Dictionary, ed. David Noel Freedman (New York: Doubleday, 1992), 350–364.

9 Eliyahu Assis, “Psalm 127 and the Polemic of the Rebuilding of the Temple in the Post Exilic Period,” Zeitschrift Für Die Alttestamentliche Wissenschaft 121, no. 2 (2009): 256-72.

10 Tremper Longman III, Psalms: An Introduction and Commentary, ed. David G. Firth, vol. 15–16, Tyndale Old Testament Commentaries (Nottingham: Inter-Varsity, 2014), 423.

11 參 YLT 的翻譯.

12 J. Clinton McCann, JR. “The Book of Psalms,” The New Interpreter’s Bible, vol. IV, Leander E. Keck ed. (Nashville: Abingdon, 1996), 1198.

ֶֶ

13 David N. Freedman and J. Lundbom, “טןבּ ,” ed. G. Johannes Botterweck and Helmer Ringgren, trans. John T. Willis, Theological Dictionary of the Old Testament (Grand Rapids: William B. Eerdmans, 1977), 95.

ֶֶ

14 David N. Freedman and J. Lundbom, “טןבּ ,” ed. G. Johannes Botterweck and Helmer Ringgren, trans. John T. Willis, Theological Dictionary of the Old Testament (Grand Rapids: William B. Eerdmans, 1977), 97.

15  \url{https://www.chabad.org/library/bible_cdo/aid/16348/showrashi/true/jewish/Chapter-127.htm.}

書⽬

Allen, Leslie C. Psalms 101–150 (Revised), vol. 21, Word Biblical Commentary. Dallas: Word, Incorporated, 2002.

Assis, Eliyahu. “Psalm 127 and the Polemic of the Rebuilding of the Temple in the Post Exilic Period,” Zeitschrift Für Die Alttestamentliche Wissenschaft 121, no. 2 (2009): 256-72.

Brown, Francis, Samuel Rolles Driver, and Charles Augustus Briggs, Enhanced Brown-Driver-Briggs Hebrew and English Lexicon. Oxford: Clarendon Press, 1977.

Claissé-Walford, Nancy de, Rolf A. Jacobson, and Beth LaNeel Tanner, “The Songs of the Ascents: Psalms,” in The Book of Psalms, ed. E. J. Young, R. K. Harrison, and Robert L. Hubbard Jr., The New International Commentary on the Old Testament. Grand Rapids: William B. Eerdmans, 2014.

Hakham, Amos. Psalms 101-105: The Bible with Jerusalem Commentary. Jerusalem: Mosad Harav Kook, 2003.

Harris, R. Laird, Gleason L. Archer Jr., and Bruce K. Waltke, ed. Theological Wordbook of the Old Testament. Chicago: Moody Press, 1999.

Hossfeld, Frank-Lothar and Erich Zenger, Psalms 3: A Commentary on Psalms 101–150, ed. Klaus Baltzer, trans. Linda M. Maloney, Hermeneia—A Critical and Historical Commentary on the Bible. Minneapolis: Fortress, 2011.

Janzen, W. “Land,” ed. David Noel Freedman, The Anchor Yale Bible Dictionary (New York: Doubleday, 1992.

Kraus, Hans-Joachim. A Continental Commentary: Psalms 60–150 Minneapolis: Fortress, 1993.

Longman, Tremper III. Psalms: An Introduction and Commentary, ed. David G. Firth, vol. 15–16, Tyndale Old Testament Commentaries. Nottingham: Inter-Varsity, 2014.

McCann, J. Clinton, JR. “The Book of Psalms,” The New Interpreter’s Bible, vol. IV, Leander E. Keck ed. Nashville: Abingdon, 1996.

VanGemeren, Willem A. “Psalms,” in The Expositor’s Bible Commentary: Psalms, Proverbs, Ecclesiastes, Song of Songs, ed. Frank E. Gaebelein, vol. 5. Grand Rapids, MI: Zondervan Publishing House, 1991.

Weiser, Artur. The Psalms: A Commentary, ed. Peter Ackroyd et al, The Old Testament Library. Philadelphia: The Westminster, 1998.



\newpage

\section{一個關於宗教信仰與公共安全的故事}
\label{sec:750}
\textbf{林豪恩先生 - 一個關於宗教信仰與公共安全的故事}
\newline
\newline
\hyperref[sec:749]{< < < PREV SERMON < < <}
~
\hyperref[sec:toc]{[返目錄]}
~
\hyperref[sec:751]{> > > NEXT SERMON > > >}


使徒行傳 16:16-34
\newline
\begin{longtable}{cl}
\hline
\hline
章節 & 經文 (和合本修訂版)\\
\hline
16:16 & \begin{tabularx}{0.7\textwidth}{X} 後來，我們往那禱告的地方去時，有一個被占卜的靈附身的使女迎面走來，她使用法術使她的主人們發了大財。 \end{tabularx} \\ \\ \relax
16:17 & \begin{tabularx}{0.7\textwidth}{X} 她跟隨保羅和我們，喊著說：「這些人是至高神的僕人，對你們傳講救人的道路。」 \end{tabularx} \\ \\ \relax
16:18 & \begin{tabularx}{0.7\textwidth}{X} 她一連好幾天這樣喊叫，保羅就心中厭煩，轉身對那靈說：「我奉耶穌基督的名吩咐你從她身上出來！」那靈立刻出來了。 \end{tabularx} \\ \\ \relax
16:19 & \begin{tabularx}{0.7\textwidth}{X} 使女的主人們見發財的指望沒有了，就揪住保羅和西拉，拉他們到市上去見官； \end{tabularx} \\ \\ \relax
16:20 & \begin{tabularx}{0.7\textwidth}{X} 又帶他們到行政官長們面前，說：「這些騷擾我們城的，他們是猶太人， \end{tabularx} \\ \\ \relax
16:21 & \begin{tabularx}{0.7\textwidth}{X} 竟傳佈我們羅馬人所不可接受、不可遵守的規矩。」 \end{tabularx} \\ \\ \relax
16:22 & \begin{tabularx}{0.7\textwidth}{X} 群眾就一齊起來攻擊他們。官長們吩咐撕開他們的衣裳，用棍子打； \end{tabularx} \\ \\ \relax
16:23 & \begin{tabularx}{0.7\textwidth}{X} 打了許多棍，就把他們下在監裡，囑咐獄警嚴緊看守。 \end{tabularx} \\ \\ \relax
16:24 & \begin{tabularx}{0.7\textwidth}{X} 獄警領了這樣的命令，就把他們下在內監，兩腳拴在木架上。 \end{tabularx} \\ \\ \relax
16:25 & \begin{tabularx}{0.7\textwidth}{X} 約在半夜，保羅和西拉正在禱告，唱詩讚美神，眾囚犯也側耳聽著的時候， \end{tabularx} \\ \\ \relax
16:26 & \begin{tabularx}{0.7\textwidth}{X} 忽然，地大震動，甚至監牢的地基都搖動了，監門立刻全開，眾囚犯的鎖鏈也都解開了。 \end{tabularx} \\ \\ \relax
16:27 & \begin{tabularx}{0.7\textwidth}{X} 獄警一醒，看見監門全開，以為囚犯已經逃走，就拔刀要自殺。 \end{tabularx} \\ \\ \relax
16:28 & \begin{tabularx}{0.7\textwidth}{X} 保羅大聲呼叫：「不要傷害自己！我們都在這裡。」 \end{tabularx} \\ \\ \relax
16:29 & \begin{tabularx}{0.7\textwidth}{X} 獄警叫人拿燈來，就衝進去，戰戰兢兢地俯伏在保羅和西拉面前。 \end{tabularx} \\ \\ \relax
16:30 & \begin{tabularx}{0.7\textwidth}{X} 然後獄警領他們出來，說：「二位先生，我必須做甚麼才可以得救？」 \end{tabularx} \\ \\ \relax
16:31 & \begin{tabularx}{0.7\textwidth}{X} 他們說：「當信主耶穌，你和你一家都必得救。」 \end{tabularx} \\ \\ \relax
16:32 & \begin{tabularx}{0.7\textwidth}{X} 他們就把主的道講給他和他全家的人聽。 \end{tabularx} \\ \\ \relax
16:33 & \begin{tabularx}{0.7\textwidth}{X} 當夜，就在那時候，獄警把他們帶去，洗他們的傷；他和他所有的家人立刻都受了洗。 \end{tabularx} \\ \\ \relax
16:34 & \begin{tabularx}{0.7\textwidth}{X} 於是獄警領他們上自己的家裡去，給他們擺上飯。他和全家的人，因為信了神，都滿心喜樂。 \end{tabularx} \\ \\
[1ex]
\hline
\hline
\end{longtable}
講題:一個關於宗教信仰與公共安全的故事  A Story About Religious Faith and Public Security

經文:使徒行傳16章16至34節

講員:林豪恩先生

場合:香港中文大學崇基學院禮拜堂主日崇拜

日期:2025年6月1日

引言

被視為「家庭治療之母」,也是香港人比較熟識的家庭治療大師沙維雅(Virginia Satir)有一句名言:「問題本身不是問題；怎麼處理才是問題.」(Problems are not the problem; coping is the problem.)有些處理問題的手法能夠化繁為簡,藥到病除；有些處理問題的手法越搞越亂,小事變大.有人在網上分享說:「還記得在某公司做工程師的時候,每次開會都會看到PM拿著機關槍朝著RD開火,然後再看到RD躲回堡壘內,靠著窗戶再丟一顆手榴彈出去.結果呢?問題沒解決,還兩敗俱傷,連帶附近的人都受到牽連.」

再舉一個近期社會關注的例子,青少年吸食太空油現象.吸食太空油對吸食者做成傷害,是要處理的問題.新聞報導有一位父親發現兒子在家吸食太空油,於是報警,警方以懷疑藏毒罪名拘捕兒子.吸食太空油這個問題會不會是其他問題的轉移呢?明報週刊訪問了四位青少年,問他們為甚麼吸食太空油.第一位說她有急切的東西需要逃離,例如讀書,例如家人.她的英文與數學從未合格.在學校被老師罵,回家被媽媽罵.第二位說:「無人可以讓你放心傾偈,那一刻還要被家人罵.我壓力又大,上班又要俾人X.回到家又要被人罵,說甚麼既然搵不到錢,你返咩工?這些事情你又不能跟別人說,你就會一直在想自己是不是真的這麼差?為甚麼我賺不到錢?就會不停地吸太空油.」另外一位說小學時被同班同學欺凌,直到小學畢業,他還是一個朋友也沒有.但與不和他說話的同學相比,他更討厭駐校社工.「我寫字不美,又常寫錯字,他(她)很喜歡是非不分地罵我.小學六年,我便喊足六年.」即使回到家,亦不代表終於有人保護.家人說要懲罰他的爛成績,時常隨手拿起衣架就是一輪抽打.他每日用校服蓋好傷疤,來回兩個煉獄之間.

用如此長篇的引言,想說明一個現象,問題是會從一個層面轉移往另一個層面的.今天我們嘗試從問題轉移的現象來閱讀經課中使徒行傳的故事.我們會看見宗教信仰的問題如何被轉移為經濟問題,然後又如何被轉移為公共安全問題.以保羅為例子的宗教人士又如何回應呢?讓我們閱讀這個故事.

一．宗教信仰問題

16:16後來,我們往那禱告的地方去.有一個使女迎著面來,她被巫鬼所附,用法術,叫她主人們大得財利.

16:17她跟隨保羅和我們,喊著說:「這些人是至高上帝的僕人,對你們傳說救人的道.」

第一,事件的兩個主角都是宗教界人士:一方是保羅及同伴,他們拜的是「至高上帝」；另一方是一位使女,她拜的神明(如果都稱為神明的話)是巫鬼.

第二,事件的雙方的工作都在某程度上和宗教有關:保羅及同伴所做的是「傳說救人的道」,而使女所做的是「用法術叫她主人們大得財利」.

兩種宗教信仰相遇,其實也不一定需要比拼,不幸地比拼最終都發生了.

16:18她一連多日這樣喊叫,保羅就心中厭煩,轉身對那鬼說:「我奉耶穌基督的名,吩咐你從她身上出來!」那鬼當時就出來了.

比拼結果是耶穌基督勝出,巫鬼敗退.如果只是宗教界內部事情,故事就這樣結束了.

歷史自古至今都有宗教戰爭.宗教分歧為甚麼可以變成戰爭呢?如果宗教分歧的雙方不把問題向外延伸,也沒有公共權力介入,更沒有武器,最多都是「打口水戰」,或者「拳頭交」,應該不會演變成戰爭.變成戰爭的宗教問題,很可能是因為被轉移了往其他層面.讓我們繼續閱讀使徒行傳的故事,看看問題如何被轉移.

二．轉移為經濟問題

16:19使女的主人們見得利的指望沒有了,便揪住保羅和西拉,拉他們到市上去見首領；

第一,除了保羅和使女對壘的雙方之外,有第三者出場.第三者是誰呢?經文指出這些人是「使女的主人們」.那就奇怪了,如果純粹是宗教範疇,宗教信徒的主人不應該是他們的神明嗎?例如使女稱保羅和同伴是「至高上帝的僕人」(16:17),他們的主人就是「至高上帝」.按此道理,使女被巫鬼所附,使女豈不是「巫鬼的僕人」嗎?她的主人豈不是「巫鬼」嗎?現在,使女的主人竟然是人,而且多過一個人,那就複雜了.試想想,如果某宗教團體的主人不單是該宗教的神明,更是某些利益集團,那麼,該宗教團體應該向誰效忠呢?有事情的時候,這些利益團體介入,豈不是更複雜嗎?怪不得有人說,鬼和神都不是最可怕,最可怕的是人.

第二,原來對於某些人來說,宗教信仰並不是最終目的,只是手段而已.敬拜神也好,鬼也好,都是為了財利.使女的主人們為甚麼介入耶穌基督和巫鬼的對決呢?為甚麼拉保羅和西拉去見官呢?是不是要維護他們的宗教呢?經文指出原因是「見得利的指望沒有了」.(16:19)在這個案例中,使女的主人們所關心的根本不是宗教,對他們來說,宗教問題不是宗教問題,而是經濟問題.今時今日,可能有些宗教人士好像那位使女一樣,一個人要事奉多個主.也有些人好似使女的主人們一樣,視神明為財利的來源.例如每年某些日子都有很多人去向觀音借庫.這日子之外,不知道這些善信還會不會去拜觀音呢?

讓我們回到使徒行傳的故事,看看他們如何把宗教事件繼續轉移至其他範疇.

三．轉移為公共安全問題

16:20又帶到官長面前說:「這些人原是猶太人,竟騷擾我們的城,

16:21傳我們羅馬人所不可受不可行的規矩.」

16:22眾人就一同起來攻擊他們.官長吩咐剝了他們的衣裳,用棍打；

16:23打了許多棍,便將他們下在監裡,囑咐禁卒嚴緊看守.

本來是使女主人與保羅雙方的私人問題,現在轉移為公共問題；本來是個人經濟問題,現在轉移為城市安全問題.本來的宗教問題既然被移轉為公共安全問題,處理的人就按著公共安全的性質來辦理,「未見官先打八十」.公眾不要以為打幾個外來人沒有甚麼大不了,也不要以為惡待幾個非我族類的人不是甚麼一回事.其實這是對制度的嚴重破壞:第一,制度被人運用來進行個人報復的工具.第二,製造了一件冤假錯案,令制度蒙上污點.動輒把宗教分歧轉移為公共安全問題,結果所有持份者都受到損害,包括公共制度本身.使女主人奸計得逞,但不要開心得太早,當制度被破壞之後,將來他們也可能是受害者.

四．宗教人士以宗教方式回應

16:24禁卒領了這樣的命,就把他們下在內監裡,兩腳上了木狗.

連那位使女都肯定保羅和同伴「是至高上帝的僕人,對你們傳說救人的道.」保羅和同伴到底做錯了甚麼以至被官府棒打及監禁呢?被投進如此的冤獄,人們通常有甚麼反應呢?

在某地有一個人被判死刑,坐了16年冤獄之後被無罪釋放,他回憶說:「剛開始都在詛咒辦案的警察,那是每天醒來必做的功課...那時是非常悲憤的,我覺得世上不可能有這麼糟糕的事啊!我對司法,對人都很絕望,給爸爸的遺書裡還說我會到地獄裡向閻羅王告狀,也不是真的有什麼信仰,只是覺得在人世間已經沒有辦法.」

我們並不是批評這位冤假錯案的受害人,反而覺得他的反應是人之常情.被冤屈,感到憤怒；求助無門,感到絕望；今生沒有辦法懲治冤枉自己的人,唯望他們來世得到報應.

保羅和同伴的宣教行動被上綱上線轉移為破壞公共安全行為,他們如何回應呢?

16:25約在半夜,保羅和西拉禱告,唱詩讚美上帝,眾囚犯也側耳而聽.

保羅和西拉所做的事情,對於冤獄來說是出人意外的回應.雖然他們被屈破壞公共安全而受罰,但他們的回應不似是公共安全層次.他們以甚麼層次來回應呢?他們「禱告,唱詩讚美上帝」是宗教行為,他們以宗教層次來回應.無論對方把問題轉移為甚麼範疇,保羅和同伴始終是宗教人士,他們始終以宗教層次來回應.在混亂,甚至胡亂的世界中,他們保住了初心,守住了本意.

五．又變成公共安全事情

16:26忽然,地大震動,甚至監牢的地基都搖動了,監門立刻全開,眾囚犯的鎖鍊也都鬆開了.

16:27禁卒一醒,看見監門全開,以為囚犯已經逃走,就拔刀要自殺.

地震是甚麼問題呢?其本質是地理問題,但對這個監獄來說地震是公共安全問題.因為囚犯逃走,就變成了公共安全事故.既然是公共安全問題,禁卒就要承擔公共安全的責任.身為官府中的一員,禁卒當然知道自己將會如何被處置,與其被羞辱折磨至死,不如自行了斷.當事情被視為公共安全問題,還有情可以講嗎?執行的人自己都知道沒有了.幸好,不是所有人都以公共安全來看待事情,也不是所有人都沒有情講的.

六．變回宗教問題

16:28保羅大聲呼叫說:「不要傷害自己!我們都在這裡.」

16:29禁卒叫人拿燈來,就跳進去,戰戰兢兢地俯伏在保羅,西拉面前；

16:30又領他們出來,說:「二位先生,我當怎樣行才可以得救?」

16:31他們說:「當信主耶穌,你和你一家都必得救.」

16:32他們就把主的道講給他和他全家的人聽.

保羅作為宗教人士,他的首要關注是人本身.他不理會這個人是甚麼身分,也不計較對方是加害者還是朋友,救人要緊.惻隱之心,人皆有之,孟子說這是人性,宗教不會忽視人性,宗教應該是致力恢復上主賦予的人性.從禁卒和保羅接下來的對話來看,他們沒有討論公共安全問題,沒有討論司法問題,沒有討論經濟問題,他們討論的是宗教問題.保羅作為宗教人士,對於有機會成為公共安全問題的事件,都轉移到宗教層面.他在大街小巷進行宗教工作,在囚室也進行宗教工作.保羅沒有辦法阻止有人上綱上線指控他們破壞公共安全,也不能阻止有人誣衊他們經濟犯罪,但作為宗教人士由始至終,在任何狀況之下,都傳主耶穌拯救世人的道.如果保羅可以申辯,他可能會說:「我們是為了信仰」.

16:33當夜,就在那時候,禁卒把他們帶去,洗他們的傷；他和屬乎他的人立時都受了洗.

16:34於是禁卒領他們上自己家裡去,給他們擺上飯.他和全家,因為信了上帝,都很喜樂.

第一,保羅把事件轉移回宗教層面,但並沒有破壞制度.保羅救回禁卒一命,但禁卒卻沒有私下釋放保羅等人；保羅也沒有強禁卒所難要求禁卒讓他們離開.這並不是說保羅軟弱就範,我們知道他後來是根據當時的制度據理力爭,討回公道和尊嚴.然而,他知道據理力爭的對象不是他面前那位舉刀要自殺的禁卒,他沒有向弱者抽刀.

第二,在宗教的感染之下,禁卒開始把保羅他們視為人.既然是人,就有人的需要,禁卒回應他們的需要:替他們洗傷口,給他們擺上飯吃.人與人之間能夠善意相待,經文最後四個字是「都很喜樂」.為甚麼能夠「都很喜樂」呢?經文指出「因為信了上帝」,在上帝作為天父的國度中,無論大家是甚麼身分角色,本質上都是天父的兒女,都是兄弟姐妹,這是大家可以把對方視為人,以人性相待的基礎.願榮耀歸予聖父聖子聖靈三位一體上帝.



\newpage

\section{五旬節到了}
\label{sec:751}
\textbf{李均熊牧師 - 五旬節到了}
\newline
\newline
\hyperref[sec:750]{< < < PREV SERMON < < <}
~
\hyperref[sec:toc]{[返目錄]}
~
\hyperref[sec:752]{> > > NEXT SERMON > > >}


使徒行傳 2:1-21
\newline
\begin{longtable}{cl}
\hline
\hline
章節 & 經文 (和合本修訂版)\\
\hline
2:1 & \begin{tabularx}{0.7\textwidth}{X} 五旬節那日到了，他們全都聚集在一起。 \end{tabularx} \\ \\ \relax
2:2 & \begin{tabularx}{0.7\textwidth}{X} 忽然，有響聲從天上下來，好像一陣大風吹過，充滿了他們所坐的整座屋子； \end{tabularx} \\ \\ \relax
2:3 & \begin{tabularx}{0.7\textwidth}{X} 又有舌頭如火焰向他們顯現，分開落在他們每個人身上。 \end{tabularx} \\ \\ \relax
2:4 & \begin{tabularx}{0.7\textwidth}{X} 他們都被聖靈充滿，就按著聖靈所賜的口才說起別國的話來。 \end{tabularx} \\ \\ \relax
2:5 & \begin{tabularx}{0.7\textwidth}{X} 那時，有從天下各國來的虔誠的猶太人，住在耶路撒冷。 \end{tabularx} \\ \\ \relax
2:6 & \begin{tabularx}{0.7\textwidth}{X} 這聲音一響，許多人都來聚集，各人因為聽見門徒用他們各自的鄉談說話，就甚納悶， \end{tabularx} \\ \\ \relax
2:7 & \begin{tabularx}{0.7\textwidth}{X} 都詫異驚奇說：「看哪，這些說話的不都是加利利人嗎？ \end{tabularx} \\ \\ \relax
2:8 & \begin{tabularx}{0.7\textwidth}{X} 我們每個人怎麼聽見他們說我們生來所用的鄉談呢？ \end{tabularx} \\ \\ \relax
2:9 & \begin{tabularx}{0.7\textwidth}{X} 我們帕提亞人、瑪代人、以攔人，和住在美索不達米亞、猶太、加帕多家、本都、亞細亞、 \end{tabularx} \\ \\ \relax
2:10 & \begin{tabularx}{0.7\textwidth}{X} 弗呂家、旁非利亞、埃及的人，並靠近古利奈的利比亞一帶地方的人，僑居的羅馬人， \end{tabularx} \\ \\ \relax
2:11 & \begin{tabularx}{0.7\textwidth}{X} 包括猶太人和皈依猶太教的人， 克里特人和阿拉伯人，都聽見他們用我們的鄉談講論神的大作為。」 \end{tabularx} \\ \\ \relax
2:12 & \begin{tabularx}{0.7\textwidth}{X} 眾人就都驚奇困惑，彼此說：「這是甚麼意思呢？」 \end{tabularx} \\ \\ \relax
2:13 & \begin{tabularx}{0.7\textwidth}{X} 還有人譏誚，說：「他們是灌滿了新酒吧！」 \end{tabularx} \\ \\ \relax
2:14 & \begin{tabularx}{0.7\textwidth}{X} 彼得和十一個使徒站起來，他就高聲向眾人說：「猶太人和所有住在耶路撒冷的人哪，這件事你們要知道，要側耳聽我的話。 \end{tabularx} \\ \\ \relax
2:15 & \begin{tabularx}{0.7\textwidth}{X} 這些人並不像你們所想的喝醉了，因為現在才早晨九點鐘。 \end{tabularx} \\ \\ \relax
2:16 & \begin{tabularx}{0.7\textwidth}{X} 這正是藉著先知約珥所說的： \end{tabularx} \\ \\ \relax
2:17 & \begin{tabularx}{0.7\textwidth}{X} 『神說：在末後的日子，我要將我的靈澆灌凡血肉之軀的。你們的兒女要說預言；你們的少年要見異象；你們的老人要做異夢。 \end{tabularx} \\ \\ \relax
2:18 & \begin{tabularx}{0.7\textwidth}{X} 在那些日子，我要把我的靈澆灌，甚至給我的僕人和婢女，他們要說預言。 \end{tabularx} \\ \\ \relax
2:19 & \begin{tabularx}{0.7\textwidth}{X} 在天上，我要顯出奇事，在地下，我要顯出神蹟，有血，有火，有煙霧。 \end{tabularx} \\ \\ \relax
2:20 & \begin{tabularx}{0.7\textwidth}{X} 太陽要變為黑暗，月亮要變為血，這都在主大而光榮的日子未到以前。 \end{tabularx} \\ \\ \relax
2:21 & \begin{tabularx}{0.7\textwidth}{X} 那時，凡求告主名的都必得救。』 \end{tabularx} \\ \\
[1ex]
\hline
\hline
\end{longtable}
講題:五旬節到了 When the Day of Pentecost Had Fully Come

經文:使徒行傳2章1至21節

講員:李均熊牧師

場合:香港中文大學崇基學院禮拜堂主日崇拜

日期:2025年6月8日

起:帶來混亂的聖靈

五旬節(Pentecost)意即五十天,因為是由逾越節七周後舉行,故在舊約稱作七七節 (The Feast of Weeks),或音譯作Shavuoth [沙夫幼特](意即幾個星期weeks).所以使徒行傳2:1「五旬節到了」其實應譯作「五旬節的日子滿足了」.五旬節在古代以色列是一個慶祝農產豐收的日子,所以有時又稱作收割節(Festival of Reaping)或初熟節(Day of the First Fruits).但到了新約時期,猶太人將這節日的焦點轉為慶祝上主所賜的禮物,而猶太人認為上主所賜最大的禮物就是「妥拉」(Torah),會堂當日會選讀舊約的路得記,其一原因是表明外邦女子路得歸化入猶太教,接受了「妥拉」這份禮物.所以經文第一句「五旬節的日子滿足了」包含對天數和周數的計算被,可以理解為對上授與「妥拉」的期待和渴望.

當然我們基督徒進一步將五旬節理解為聖靈降臨節,又可以作基督教會誕生的日子.所以無論怎樣說,五旬節都是興祝快樂的日子.大家看到教會用上了婚禮或堂慶日子才用的紅色作為主色就可以感受到.

但使徒行傳記載的聖靈降臨,對當時在耶路撒冷過節的人來說,最初帶來的並不是歡喜快樂.仔細看經文,聖靈降臨帶來的首先是令人可畏的大風的響聲,和如火焰的舌頭降下.這些都可能是作者回響出埃及記所記載,在西乃山中上主顯現,頒佈妥拉時在山頂出現的雷轟閃電和烈火.然候門徒說起別國的話,效果是令當時在耶路撒冷過節的猶太「聽見門徒用眾人的鄉談說話,就甚納悶」,而且「都驚訝希奇」,大家因為「驚訝猜疑」就彼此說「這是甚麼意思呢!」

中文譯為「納悶」的字,有很多英譯本譯作bewilderment/bewildered,但查原文這個字,也是希臘文翻譯舊約創世記著名的「巴別塔」敘述中上主要「變亂」各人口音的那個希伯來文:「我們下去,在那裡變亂他們的口音,使他們的言語彼此不通.」(創11:7)接著的「驚訝希奇」的希臘文可以是正面也可以是負面的,但大多數英譯本都取比較正面的amazed 或marvelled來譯此字.但2:12 「驚訝猜疑」希臘文乃代表極大程度的費解,字典的解釋是to be thoroughly perplexed - 'to be very perplexed, to not know what to do, to be very confused.(徹底困惑 - 「非常困惑,不知道該怎麼辦,非常迷茫.」)

仔細看這一段的經文的用字,我認為使徒行傳記載聖靈降臨的那一個五旬節,散居於各處,甚至是出生於「天下各國」,又來到耶路撒冷過節的的猶太人,並不會因為在他過節的異地聽到鄉談而感到快慰或者高興.他們其實是「非常困惑,不知道該怎麼辦,非常迷茫.」聖靈降臨開首,竟如創世記記載的巴別塔事件一般,帶來的是混亂!

承:巴別塔和五旬節的共通處

有一類解經家以為五旬節聖靈降臨是要將巴別塔的事件倒轉過來,除消undo巴別塔事件帶來的語言不通.但若仔細看,與其說五旬節是倒轉巴別塔事件,不如說五旬切是重複巴別塔事件.將五旬節解說為巴別塔事件倒轉,可能是誤解了創世記中巴別塔事件的目的.不如讓我們用少少時間重溫一下巴塔事件的來龍去脈.創世記11:1-9實是這樣記載的

1. 那時,天下人的口音,言語都是一樣.2他們往東邊遷移的時候,在示拿地遇見一片平原,就住在那裡.3他們彼此商量說:「來吧!我們要做磚,把磚燒透了.」他們就拿磚當石頭,又拿石漆當灰泥.4他們說:「來吧!我們要建造一座城和一座塔,塔頂通天,為要傳揚我們的名,免得我們分散在全地上.」5耶和華降臨,要看看世人所建造的城和塔.6耶和華說:「看哪,他們成為一樣的人民,都是一樣的言語,如今既做起這事來,以後他們所要做的事就沒有不成就的了.7我們下去,在那裡變亂他們的口音,使他們的言語彼此不通.」8於是耶和華使他們從那裡分散在全地上；他們就停工,不造那城了.9因為耶和華在那裡變亂天下人的言語,使眾人分散在全地上,所以那城名叫巴別.

當時人要興建巴別塔其實有兩個目的(還是一個目的的兩個階段?):一是為了揚名,二是免得分散在全地上.先不論揚名的本意是好是壞,但免得分散顯然是上主變亂他們言語要抗衡的目的,因為創世記記載明確指出「於是耶和華使他們從那裡分散在全地上；他們就停工,不造那城了.因為耶和華在那裡變亂天下人的言語,使眾人分散在全地上,所以那城名叫巴別.」

轉:分散還是多元

若我向一間著重傳福音的福音派教會,可能講道到此就是挑戰弟兄姊妹們不要以為返教會就是為著聚集一起,圍爐自暖,卻應分散到天下,完成大使命,承擔使徒行傳開首的「但聖靈降臨在你們身上,你們就必得著能力,並要在耶路撒冷,猶太全地,和撒馬利亞,直到地極,作我的見證.」但崇基禮拜堂是設在大學校園的禮拜堂,主要事工並非傳福音,而是牧養校園裡面的師生(我想高牧師會同意吧).在座中有些來崇拜的弟兄姊妹,我想也不是因為想向校園裡的學生傳福音,而是感到這裡崇拜講道能牧養心靈,才會在這裡聚會吧(corret me if I am wrong)?所以我想大家繼續留意使徒行傳記載五旬節聖靈降臨,並非純粹重現巴別塔事件的目的,而是更進一步,將基督的福音,神的大作為,能讓不同語言,文化,身份的人,都能明白和認識.換句話說,分散不必是空間上的分散,也可以是文化身份的分散.聖靈降臨除了能使教會信徒能分散住普天下去,亦能建立教會能開放接納多元文化和身份的信徒.

若以使徒行傳的脈絡去看,聖靈降臨先是建立一所能承載多元文化,身份的教會,才讓信徒分散住普天下去.使徒行傳第6章,就是首先要處理說希臘語的語言／文化群體發怨言的問題,而選立聖靈充滿去管派寡婦飯食的人.由第8章起,教會才分散到各處傳揚福音,而到了第10章,彼得才首次向外邦人哥尼流傳福音和外邦人領受聖靈.(第11章?)

回到我們今日的經課,使徒行傳2:7-11仔細記載這些散居的猶太人是屬甚麼語言文化的人,來自亞,非,歐大陸總供十五個區域.當中包括來自耶路撒冷東面的地區:帕提亞人,米底亞人,以攔人,和住在美索不達米亞(現在的伊朗伊拉克的地方),又有南面的猶大,也有北面小亞細亞一帶:加帕多家,本都,亞細亞,弗呂家(今日的土耳其),更有西南面北非的埃及,利比亞,以及西面的羅馬,克里特和東南面的阿拉伯.不過大家要留意,這個名單並非複製自任何已知的名單；除了羅馬是當時主要的城市,一些猶太人主要的聚居地,如希臘或馬其頓,並不包括在這個名單內；名單的先後次序看來並沒有特定的方向和原則,像是隨意的包羅東南西北的語言文化的代表,而非根據一個一致的大原則去陳列.換句話說,雖然在這階段聖靈所賜的大能的顯現對像仍然是猶太人,但這些語言文化的羅致並非服務一個大原則,而是傾向多元分散.

最後,若我們看看彼得所引用的約珥書的經文,上主的靈要澆灌「凡有血氣的」,但進一步指明是「兒子」和「女兒」；「少年人」和「老年人」；「僕人」和「使女」,顯然凡有血氣的並不分男／女；老／幼；以及地位高低.使徒行傳記載五旬節到來,聖靈降臨就是要展現一個多元的教會群體,在任何語言文化,性別,年齡,階級,以及之後不同的種族,都能明白領受上帝的大能.

但不要忘記,聖靈降臨所開創的局面,首先帶來的卻是困惑與混亂,不安與猜疑.為甚麼會這樣呢?這是因為無論使徒行傳記載的帝國大背景,又或我們現在身處的高度現代化的社會,都想透過標準化(standardization)去增加管治/管理的效率.【標準化對某些失去身體某些能力的人士產生障礙】帝國亦因為政治的因素,要削弱局部差異,傳播大一統的意識形態.各地方必須學習和使用官方語言,在文化上只有中央文化才是優等,其他文化若不是野蠻的,就是次等的.今日進入後現代,後殖民的世界,我們才明白到世界不必要由一個文明(civilization)去同化,而必須透過對話去理解不同文化(culture)均有價值.若非五旬節聖靈的降臨和加力,我們作為普通人也只會願意和自己人交徒,活在同溫層的安舒區內,而不會與自己所不認識的異文化相處吧.[作為講道場合,我想不用在這裡多討論.大學的文學院,社會科學院,不同科目對此都有所研究,這也是大學裡的禮拜堂的好處吧】

合:合一非一律,開放與多元

五旬節的故事並非完在今日的經課,接著的經文是彼得對現場散居的猶太的講道,當日就有三千人受洗,建立了耶路撒冷教會.使徒行傳2章的結束,展現一個新的群體:「43眾人都懼怕；使徒又行了許多奇事神蹟.44信的人都在一處,凡物公用,45並且賣了田產,家業,照各人所需用的分給各人.46他們天天同心合意恆切地在殿裡,且在家中擘餅,存著歡喜,誠實的心用飯,47讚美神,得眾民的喜愛.主將得救的人天天加給他們.」

大家現在身處的乃是崇基學院神學院的聖堂,是師生在學期時間每週四一同敬拜眾集的地方.大家可以向上望,會見到一個環形結構.那是這聖堂興建時所設計的「合一環」.聖靈所賜的乃是合而為一的心,我們重視合一,正如以弗所書4章作者的勸勉:「凡事謙虛,溫柔,忍耐,用愛心互相寬容,用和平彼此聯絡,竭力保守聖靈所賜合而為一的心.」但我們亦會強調,合一非一律:合一並非每一個人都守著同一個規律,所有事情都要一致.在合一之前,我們也必須看見大家的不同,無論是語言上,文化上,性別／年齡／階級／身份都是不一樣,才是五旬節聖靈降臨要開創的信徒群體.

崇基學院禮拜堂,是所大學內的禮拜堂.我們有英語和普通話的傳譯,這也是我們的特色吧,願聖靈臨這些傳譯的生命中,成就上主的大能和美意.除了語言上的多元,也祈求聖靈臨在崇基校牧室,使我們能樣大學裡的多元群體認識上主的大能.



\newpage

\section{天地有正氣}
\label{sec:752}
\textbf{鄧瑞強博士 - 天地有正氣}
\newline
\newline
\hyperref[sec:751]{< < < PREV SERMON < < <}
~
\hyperref[sec:toc]{[返目錄]}
~
\hyperref[sec:753]{> > > NEXT SERMON > > >}


那鴻書 1:1-3
\newline
\begin{longtable}{cl}
\hline
\hline
章節 & 經文 (和合本修訂版)\\
\hline
1:1 & \begin{tabularx}{0.7\textwidth}{X} 論尼尼微的默示，伊勒歌斯人那鴻所見異象的書。 \end{tabularx} \\ \\ \relax
1:2 & \begin{tabularx}{0.7\textwidth}{X} 耶和華是忌邪、報應的神。耶和華施報應，大有憤怒；耶和華向他的敵人報應，向他的仇敵懷怒。 \end{tabularx} \\ \\ \relax
1:3 & \begin{tabularx}{0.7\textwidth}{X} 耶和華不輕易發怒，大有能力，但耶和華萬不以有罪的為無罪。他的道路在旋風和暴風之中，雲彩為他腳下的塵土。 \end{tabularx} \\ \\
[1ex]
\hline
\hline
\end{longtable}
講題:天地有正氣 God is Righteous

經文:那鴻書1:1-3;1:9-2:2;3:1-7,19

講員:鄧瑞強博士

場合:香港中文大學崇基學院禮拜堂主日崇拜

日期:2025年6月15日

各位弟兄姊妹,早安.願神的恩典常與你們同在.

崇基禮拜堂主日崇拜用的三年一循環的經課(lectionary),不是每卷書都包括在內的.今日我要講的那鴻書,就不包括在內.若我們嚴格按照經課的經文來講道的話,則大家返了一輩子崇基禮拜堂,都不會聽到那鴻書的講道.我今日的講道,就講講這卷被遺忘了的書卷.

主前722年,北國以色列被亞述所滅.自此以後,南國猶大便不斷受亞述欺凌,掠奪,直到亞述於主前612年被巴比倫所滅才停止.猶大國被亞述掐住頸項,前後有110年,超過一個世紀.

亞述非常強大,猶大十分微弱,他們無力反抗,只能啞忍.但,現實就只能是這樣的嗎?歷史的前景就只能由強權來書寫的嗎?天地還有沒有一種叫「正氣」的東西呢?

先知那鴻就在這歷史形勢下,宣告他的信息.

先看那鴻書1:1-3:

鴻1:1　論尼尼微的默示,就是伊勒歌斯人那鴻所得的默示.

鴻1:2　耶和華是忌邪施報的上帝.耶和華施報大有忿怒；向他的敵人施報,向他的仇敵懷怒.

鴻1:3　耶和華不輕易發怒,大有能力,萬不以有罪的為無罪.他乘旋風和暴風而來,雲彩為他腳下的塵土.

幾節經文,就為那鴻書定調.

「尼尼微」就是亞述的首都,這書宣告亞述的命運.以前講約拿書時,會給你一個印象,尼尼微城一聽到神的警誡,便全城悔改,聽起來他們很純良似的,而神也很眷顧他們.現實是尼尼微是一邪惡之城,神已在數算日子,消滅他們.

那鴻書1:2-3引用出埃及記34:6-7的經文,論述神的本性.

「忌邪」(qanno, jealous),是對神的首要描述.中文「忌邪」不太能表達這詞的含意.這詞的核心意思是:以激情去守護自己所愛的,以忠貞去維護約的關連,不容任何事物破壞這愛的連繫.在近東眾多的神明裡,只有耶和華有這特性.一般的神明都是反復無常,唯我獨專,不理他人死活,唯有耶和華,祂專一愛祂的子民,也希望祂的子民專一愛祂.在神人的專一關係裡,耶和華要求祂的子民維護正義,實踐仁愛,這要求讓信靠耶和華的人,看到世界的道德底蘊.當神的子民被不正義對待時,「忌邪」的神會大發熱心,為了祂所愛的子民,會維護正義,施行審判.神的這特性,讓神的子民在歷史的艱難中,保持希望.

為了祂所愛的,耶和華萬不以有罪的為無罪,祂有所行動.鴻1:2有一個字眼重復了三次:「施報」,重要的事情要說三次,神要守護祂所愛的,一定會報應加害祂子民的人.亞述過份殘暴,成了神的敵人,神要消滅他.

短短兩節經文,我們看到:這位神,對祂的子民充滿情感,祂看到他們受亞述的不正義對待,祂要行動,要保護祂的子民,他們是屬祂的,不是屬亞述的,神要以祂超越萬有的大能,守護正義,消滅不義的人.就因為這樣,神的子民在苦難中,仍可以堅持希望.畢竟,世界是有上帝的,歷史是有方向的,天地是有正氣的.

我們看第二段經文:那鴻書1:9-2:2.這段經文夾雜著審判亞述和解救猶大的宣告,一方面表達神對惡人發怒,另方面卻表達神對祂的子民的眷愛,正義與仁愛並行.神對亞述的審判是嚴厲的,徑直宣告他的滅亡.神對猶大的眷顧是溫柔的,最微弱的人民,不是孤獨地面對人間的苦難的,神與他們同在,解救他們.

神對亞述的審判之言(1:9-11)

鴻1:9　尼尼微人哪,設何謀攻擊耶和華呢?他必將你們滅絕淨盡；災難不再興起.

鴻1:10　你們像叢雜的荊棘,像喝醉了的人,又如枯乾的碎秸全然燒滅.

鴻1:11　有一人從你那裡出來,圖謀邪惡,設惡計攻擊耶和華.

神對猶大眷顧的話(1:12-13)

鴻1:12　耶和華如此說:尼尼微雖然勢力充足,人數繁多,也被剪除,歸於無有.猶大啊,我雖然使你受苦,卻不再使你受苦.

鴻1:13　現在我必從你頸項上折斷他的軛,扭開他的繩索.

神對亞述的審判之言(1:14)

鴻1:14　耶和華已經出令,指著尼尼微說:你名下的人必不留後；我必從你神的廟中除滅雕刻的偶像和鑄造的偶像；我必因你鄙陋,使你歸於墳墓.【新譯本:我必為你預備墳墓,因為你沒有存在的價值.】

神對猶大眷顧的話(1:15)

鴻1:15　看哪,有報好信傳平安之人的腳登山,說:猶大啊,可以守你的節期,還你所許的願吧!因為那惡人不再從你中間經過,他已滅絕淨盡了.

神對亞述的審判之言(2:1)

鴻2:1　尼尼微啊,那打碎邦國的上來攻擊你.你要看守保障,謹防道路,使腰強壯,大大勉力.

神對猶大眷顧的話(2:2)

鴻2:2　耶和華復興雅各的榮華,好像以色列的榮華一樣；因為使地空虛的,已經使雅各和以色列空虛,將他們的葡萄枝毀壞了.【將這經文的次序倒轉一下,較易理解,意譯如下:縱使毀壞者使地荒涼,也毀壞了葡萄樹,但耶和華要恢復雅各和以色列先前的榮美.】

耶和華好像一個壯男,走過來去保護祂那被調戲的女朋友一樣.祂很嚴厲的對那調戲者說:「你在作什麼?你今次死定了.」卻溫柔的對祂的愛人說:「不用怕,有我在,那人再不能傷害你了.」對不義者,耶和華是很嚴厲的.當祂的子民行不義時,耶和華對他們也是很嚴厲的.但神始終是信實的,縱使祂的子民失信,不守約,拜偶像,行不義,神自己依然守約,到了合適的時候,祂總會挽回他們.像何西阿先知所說,神是痴心的單戀者,無論祂的子民多麼負心,祂仍堅持愛他們.

神為何這麼嚴厲的對待亞述?我們看第三段經文,你們也評評理,看看他是否死有餘辜?我們看那鴻書3:1-7,19.

鴻3:1　禍哉!這流人血的城,充滿謊詐和強暴─搶奪的事總不止息.

以「流人血的城」來形容尼尼微,這城的生活內容是「謊詐」,「強暴」,「搶奪」,有這樣的城巿,是人類的悲劇.在今天,在戰場,仍有這種謊話連篇,強暴掠奪,視人命為草芥的城巿.

鴻3:2　鞭聲響亮,車輪轟轟,馬匹踢跳,車輛奔騰,【按:若是由杜甫翻譯,將是:車轔轔,馬蕭蕭】

鴻3:3　馬兵爭先,刀劍發光,槍矛閃爍,被殺的甚多,屍首成了大堆,屍骸無數,人碰著而跌倒,【按:正是刀光劍影,屍橫遍野】

鴻3:4　都因那美貌的妓女多有淫行,慣行邪術,藉淫行誘惑列國,用邪術誘惑〔原文是賣〕多族.【按:「極美的妓女」,代表任何吸引你而令你跌倒的東西.在啟示錄,也有一誘惑人的妓女,她代表人人追求的榮華富貴.】

鴻3:5　萬軍之耶和華說:我與你為敵；我必揭起你的衣襟,蒙在你臉上,使列國看見你的赤體,使列邦觀看你的醜陋.【按:讓世人看見你的裸露,看見你的羞恥.這其實指:讓世人明白你這種生命的醜陋.】

鴻3:6　我必將可憎污穢之物拋在你身上【按:Eugene Peterson的The Message Bible傳神地翻譯為:「用狗糞砸你」.可看到:神對暴力的憎惡】,辱沒你,為眾目所觀.

鴻3:7　凡看見你的,都必逃跑離開你,說:尼尼微荒涼了!有誰為你悲傷呢?我何處尋得安慰你的人呢?......

鴻3:19　你的損傷無法醫治；你的傷痕極其重大.凡聽你信息的,必都因此向你拍掌.你所行的惡,誰沒有時常遭遇呢?【按:那鴻書就在這裡結束.】

亞述的邪惡,就是為了人間的榮華富貴,採用極端的暴力,對別人的生命,沒有半點顧惜.大家說,亞述該不該死?所以當他滅亡時,人人拍掌,歡慶亞述這暴政的結束.

我由先知以賽亞講起,講到今日的先知那鴻,已是第11位先知了.大家應該看到,先知的信仰,這也就是基督信仰的根,其實是嵌在歷史裡面,回應時代困局的.一般信徒習慣的基督信仰,是「信耶穌,上天堂」,這種信仰是與歷史無關的.「信耶穌,上天堂」不理會當下歷史是在秦朝,抑或是在清朝,不在意當下任何社會問題,對現世歷史的前景也沒有什麼期望.先知一點也不認識這種信仰形態.

就以先知那鴻為例,他感受到當下猶大民族及其他民族在暴力的亞述底下的苦難,亞述予取予攜,任意掠奪.亞述的軍力是強大的,從現實政治來看,微弱者不可能有任何出路,只有絕望.先知的信仰,不是「信上帝,然後等上天堂」.而是在神的啟示裡,尋找激勵民族保持希望的信念.信仰不是可有可無的,不是不顧現世的,而是苦難的人保持希望的力量根源.先知看到,在這暴力橫行的亂世,背後仍有道德秩序的.先知看到,在拳頭鬥大的血腥世界背後,仍有一駕馭一切暴力的大能者,祂有能力主持正義.先知看到,微弱的人不是無助地在受苦,憐憫蒼生的神默默地看顧著他們.先知看到,世上看來無道,但事實上天地有正氣.先知看到,歷史不是在黑暗中無方向地前進的,正義會在歷史的將來顯明.先知看到,暴力地掌控今天的人,不能掌控將來,因為將來是在神的手中的.先知看到,暴力的人不是無敵的,因為神作為他們的敵人,消滅他們的力量.先知看到,在這個叫我們絕望的世界裡,我們仍有保持希望的理由.

深願各位弟兄姊妹,在亂世裡,相信上帝,保持希望.

願神祝福你們.

願榮耀歸予聖父,聖子,聖靈.



\newpage

\section{回到你的家}
\label{sec:753}
\textbf{高國雄牧師 - 回到你的家}
\newline
\newline
\hyperref[sec:752]{< < < PREV SERMON < < <}
~
\hyperref[sec:toc]{[返目錄]}
~
\hyperref[sec:754]{> > > NEXT SERMON > > >}


路加福音 8:26-39
\newline
\begin{longtable}{cl}
\hline
\hline
章節 & 經文 (和合本修訂版)\\
\hline
8:26 & \begin{tabularx}{0.7\textwidth}{X} 他們到了格拉森人的地區，就在加利利的對面。 \end{tabularx} \\ \\ \relax
8:27 & \begin{tabularx}{0.7\textwidth}{X} 耶穌上了岸，就有城裡一個被鬼附的人迎著他走來。這個人好久不穿衣服，不住在屋子裡，而住在墳墓裡。 \end{tabularx} \\ \\ \relax
8:28 & \begin{tabularx}{0.7\textwidth}{X} 他看見耶穌，就喊叫著俯伏在他面前，大聲說：「至高神的兒子耶穌，你為甚麼干擾我？我求你，不要叫我受苦！」 \end{tabularx} \\ \\ \relax
8:29 & \begin{tabularx}{0.7\textwidth}{X} 這是因耶穌曾吩咐污靈從這人身上出來。原來這污靈屢次抓住他；他常被人看守，又被鐵鏈和腳鐐捆鎖，他竟把鎖鏈掙斷，被鬼趕到曠野去。 \end{tabularx} \\ \\ \relax
8:30 & \begin{tabularx}{0.7\textwidth}{X} 耶穌問他：「你的名字叫甚麼？」他說：「群」；這是因為附著他的鬼多。 \end{tabularx} \\ \\ \relax
8:31 & \begin{tabularx}{0.7\textwidth}{X} 鬼就央求耶穌不要命令他們到無底坑裡去。 \end{tabularx} \\ \\ \relax
8:32 & \begin{tabularx}{0.7\textwidth}{X} 那裡有一大群豬正在山坡上吃食，鬼央求耶穌准他們進入豬裡；耶穌准了他們。 \end{tabularx} \\ \\ \relax
8:33 & \begin{tabularx}{0.7\textwidth}{X} 於是鬼從那人出來，進入豬裡，那群豬就闖下山崖，投進湖裡，淹死了。 \end{tabularx} \\ \\ \relax
8:34 & \begin{tabularx}{0.7\textwidth}{X} 放豬的看見這事就逃跑了，去告訴城裡和鄉下的人。 \end{tabularx} \\ \\ \relax
8:35 & \begin{tabularx}{0.7\textwidth}{X} 眾人出來，要看發生了甚麼事；到了耶穌那裡，發現那人坐在耶穌腳前，鬼已離開了他，穿著衣服，神智清醒，他們就害怕。 \end{tabularx} \\ \\ \relax
8:36 & \begin{tabularx}{0.7\textwidth}{X} 看見這事的人把被鬼附的人怎麼得醫治的事告訴他們。 \end{tabularx} \\ \\ \relax
8:37 & \begin{tabularx}{0.7\textwidth}{X} 格拉森周圍地區的人，因為害怕得很，都求耶穌離開他們；耶穌就上船回去了。 \end{tabularx} \\ \\ \relax
8:38 & \begin{tabularx}{0.7\textwidth}{X} 鬼已從身上出去的那人懇求要和耶穌在一起，耶穌卻打發他回去，說： \end{tabularx} \\ \\ \relax
8:39 & \begin{tabularx}{0.7\textwidth}{X} 「你回家去，傳講神為你做了多麼大的事。」他就走遍全城，傳揚耶穌為他做了多麼大的事。 \end{tabularx} \\ \\
[1ex]
\hline
\hline
\end{longtable}
講題:回到你的家 Return To Your Home

經課:路加福音8:26-39

講員:高國雄牧師

場合:香港中文大學崇基學院禮拜堂主日崇拜

日期:2025年6月22日

路加福音8:26-39

26他們到了格拉森人的地方,就是加利利的對面. 27耶穌上了岸,就有城裡一個被鬼附著的人迎面而來.這個人許久不穿衣服,不住房子,只住在墳塋裡. 28他見了耶穌,就俯伏在他面前,大聲喊叫,說:「至高上帝的兒子耶穌,我與你有甚麼相干?求你不要叫我受苦!」 29是因耶穌曾吩咐污鬼從那人身上出來.原來這鬼屢次抓住他；他常被人看守,又被鐵鍊和腳鐐捆鎖,他竟把鎖鍊掙斷,被鬼趕到曠野去. 30耶穌問他說:「你名叫甚麼?」他說:「我名叫『群』」；這是因為附著他的鬼多. 31鬼就央求耶穌,不要吩咐他們到無底坑裡去.32那裡有一大群豬在山上吃食.鬼央求耶穌,准他們進入豬裡去.耶穌准了他們, 33鬼就從那人出來,進入豬裡去.於是那群豬闖下山崖,投在湖裡淹死了. 34放豬的看見這事就逃跑了,去告訴城裡和鄉下的人. 35眾人出來要看是甚麼事；到了耶穌那裡,看見鬼所離開的那人,坐在耶穌腳前,穿著衣服,心裡明白過來,他們就害怕. 36看見這事的便將被鬼附著的人怎麼得救告訴他們. 37格拉森四圍的人,因為害怕得很,都求耶穌離開他們；耶穌就上船回去了. 38鬼所離開的那人懇求和耶穌同在；耶穌卻打發他回去,說: 39「你回家去,傳說上帝為你做了何等大的事.」他就去,滿城裡傳揚耶穌為他做了何等大的事.

耶穌趕鬼入豬群,一個我們很熟悉的故事…

我們因不同的原因,對此故事有不同的感受

我從少就很喜愛這神蹟,因為覺得很「精彩」

其震撼教導,讓我認識耶穌「很厲害」

但這福音書記載這神蹟的真正目的又是什麼呢?

讓我們重新一次,進入這個故事當中,一同體會耶穌基督的心意…

8:26他們到了格拉森人的地方,就是加利利的對面.

他們(耶穌與門徒們)為什麼要到格拉森呢?

8:22有一天,耶穌和門徒上了船,對門徒說:「我們可以渡到湖那邊去.」他們就開了船.

耶穌主動要求,吩咐門徒跟從而去

8:23-25 平靜風浪的故事

8:40-耶穌已離開,回到原來的地方

馬可福音5:21耶穌坐船又渡到那邊去

耶穌專程為格拉森的弟兄走這一趟

耶穌帶著門徒,冒生命危險,為要拯救他!

這個是一個怎麼樣的人呢?

馬可福音5:3那人常住在墳塋裡,沒有人能捆住他,就是用鐵鍊也不能； 4因為人屢次用腳鐐和鐵鍊捆鎖他,鐵鍊竟被他掙斷了,腳鐐也被他弄碎了；總沒有人能制伏他. 5他晝夜常在墳塋裡和山中喊叫,又用石頭砍自己.

你會如果形容他?

你願意接近他嗎?

8:28他見了耶穌,就俯伏在他面前,大聲喊叫,說:「至高上帝的兒子耶穌,我與你有甚麼相干?求你不要叫我受苦!」 29是因耶穌曾吩咐污鬼從那人身上出來.

耶穌第一次趕鬼,沒有即時「成功」

30耶穌問他說:「你名叫甚麼?」他說:「我名叫『群』」；這是因為附著他的鬼多.

耶穌問它名字,為了要讓我們明白當中的複雜性…

群=軍團,很多的意思

8:32那裡有一大群豬在山上吃食.鬼央求耶穌,准他們進入豬裡去.耶穌准了他們, 33鬼就從那人出來,進入豬裡去.於是那群豬闖下山崖,投在湖裡淹死了. 34放豬的看見這事就逃跑了,去告訴城裡和鄉下的人. 35眾人出來要看是甚麼事；到了耶穌那裡,看見鬼所離開的那人,坐在耶穌腳前,穿著衣服,心裡明白過來,他們就害怕.

8:37格拉森四圍的人,因為害怕得很,都求耶穌離開他們；耶穌就上船回去了.

經歷神蹟為我們帶來什麼?

我們期望神蹟/上帝的拯救在我們身上/身邊發生嗎?

「害怕」,「害怕得很」…為什麼?

馬可福音5:13耶穌准了他們,污鬼就出來,進入豬裡去.於是那群豬闖下山崖,投在海裡,淹死了.豬的數目約有二千.

二千頭豬的價值:一頭食用豬平均約200斤,在香港今天如每斤豬肉可賣\$100,每頭豬值\$20,000,二千頭豬=\$40,000,000

你願意花上四千萬去換取一個人的清醒/得救嗎?

8:38鬼所離開的那人懇求和耶穌同在；耶穌卻打發他回去…

耶穌拒絕這弟兄的跟從,為什麼?

難道耶穌不愛他嗎?絕對不是(要記得耶穌及門徒是冒生命危險而專程去拯救…)

不是只有人拒絕耶穌,耶穌接納所有人嗎?

耶穌清楚知道他的心,知道他真正的需要…

他需要被接納,不單是上主的接納.

所以耶穌邀請他

8:39「你回家去,傳說上帝為你做了何等大的事.」

「家」可以很溫馨,很有愛,很美好…

但也可以是我們最難面對的,因為可能太赤裸,距離太近,而讓我們懼怕.

這格拉森的弟兄,他的家,就正正是傷害仔最深,同時又是被他傷害最深的群體…

試想像當他回到家中,如何面對那些曾把他捆綁,鎖起,放在墳墓的家人親友…

又如何面對曾被他暴力傷害過的人…

更甚…他又如何面對因他而失去二千頭豬的那位豬農???

耶穌的邀請,有著極重的份量.

這亦可能是今天上主對我們的邀請

你願意接近這邀請嗎?

讓我分享一首詩歌

《回家》曲:李健達

詞:張祥志,李健達

步伐是否感覺疲倦了

現實令你感到無奈

只懂每日每夜忙碌裡

兜兜轉轉　似沒了沒完

但願讓你知道　神是愛

在路上每一秒同在

可否接受這份純真愛

一生歡欣　快樂精彩

回家　將傷痛放下

看透俗世謊話

迷失漆黑中　祂總把你念掛

回家　即使新世代

全然皆冰冷

在這屋裡　有主的愛　別再怕

(在這屋裡　有主的愛　不要害怕)

8:39      他就去,滿城裡傳揚耶穌為他做了何等大的事.

路加用簡單的一節交代了故事的結局,

我們今天又將會經歷一個怎樣「回家」的故事?



\newpage

\section{誰是我的鄰舍呢}
\label{sec:754}
\textbf{鄭躍培牧師 - 誰是我的鄰舍呢}
\newline
\newline
\hyperref[sec:753]{< < < PREV SERMON < < <}
~
\hyperref[sec:toc]{[返目錄]}
~
\hyperref[sec:755]{> > > NEXT SERMON > > >}


路加福音 10:25-37
\newline
\begin{longtable}{cl}
\hline
\hline
章節 & 經文 (和合本修訂版)\\
\hline
10:25 & \begin{tabularx}{0.7\textwidth}{X} 有一個律法師起來試探耶穌，說：「老師！我該做甚麼才可以承受永生？」 \end{tabularx} \\ \\ \relax
10:26 & \begin{tabularx}{0.7\textwidth}{X} 耶穌對他說：「律法上寫的是甚麼？你是怎樣念的呢？」 \end{tabularx} \\ \\ \relax
10:27 & \begin{tabularx}{0.7\textwidth}{X} 他回答說：「你要盡心、盡性、盡力、盡意愛主—你的神，又要愛鄰如己。」 \end{tabularx} \\ \\ \relax
10:28 & \begin{tabularx}{0.7\textwidth}{X} 耶穌對他說：「你回答得正確，你這樣做就會得永生。」 \end{tabularx} \\ \\ \relax
10:29 & \begin{tabularx}{0.7\textwidth}{X} 那人要證明自己有理，就對耶穌說：「誰是我的鄰舍呢？」 \end{tabularx} \\ \\ \relax
10:30 & \begin{tabularx}{0.7\textwidth}{X} 耶穌回答：「有一個人從耶路撒冷下耶利哥去，落在強盜手中。他們剝去他的衣裳，把他打個半死，丟下他走了。 \end{tabularx} \\ \\ \relax
10:31 & \begin{tabularx}{0.7\textwidth}{X} 偶然有一個祭司從那條路下來，看見他就從另一邊過去了。 \end{tabularx} \\ \\ \relax
10:32 & \begin{tabularx}{0.7\textwidth}{X} 又有一個利未人來到那裡，看見他，也照樣從另一邊過去了。 \end{tabularx} \\ \\ \relax
10:33 & \begin{tabularx}{0.7\textwidth}{X} 可是，有一個撒瑪利亞人路過那裡，看見他就動了慈心， \end{tabularx} \\ \\ \relax
10:34 & \begin{tabularx}{0.7\textwidth}{X} 上前用油和酒倒在他的傷處，包裹好了，扶他騎上自己的牲口，帶他到旅店裡去，照應他。 \end{tabularx} \\ \\ \relax
10:35 & \begin{tabularx}{0.7\textwidth}{X} 第二天，他拿出兩個銀幣來，交給店主，說：『請你照應他，額外的費用，我回來時會還你。』 \end{tabularx} \\ \\ \relax
10:36 & \begin{tabularx}{0.7\textwidth}{X} 你想，這三個人哪一個是落在強盜手中那人的鄰舍呢？」 \end{tabularx} \\ \\ \relax
10:37 & \begin{tabularx}{0.7\textwidth}{X} 他說：「是憐憫他的。」耶穌對他說：「你去，照樣做吧！」 \end{tabularx} \\ \\
[1ex]
\hline
\hline
\end{longtable}
雅各書 1:27
\newline
\begin{longtable}{cl}
\hline
\hline
章節 & 經文 (和合本修訂版)\\
\hline
1:27 & \begin{tabularx}{0.7\textwidth}{X} 在神—我們的父面前，清潔沒有玷污的虔誠就是看顧在患難中的孤兒寡婦，並且保守自己不沾染世俗。 \end{tabularx} \\ \\
[1ex]
\hline
\hline
\end{longtable}
講題:誰是我的鄰舍呢? Who Is My Neighbor?

經課:路加福音10:25-37;雅各書1:27

講員:鄭躍培牧師

場合:香港中文大學崇基學院禮拜堂主日崇拜

日期:2025年6月29日

大綱:

路10:25/37 有一個律法師起來試探耶穌,說:「夫子,我該做什麼才可以承受永生?」耶穌對他說:「律法上寫的是什麼?你念的是怎樣呢?」他回答說:「你要盡心,盡性,盡力,盡意愛主你的神,又要愛鄰舍如同自己.」耶穌說:「你回答的是.你這樣行,就必得永生.」那人要顯明自己有理,就對耶穌說:「誰是我的鄰舍呢?」耶穌回答說:「有一個人從耶路撒冷下耶利哥去,落在強盜手中.他們剝去他的衣裳,把他打個半死,就丟下他走了.偶然有一個祭司從這條路下來,看見他,就從那邊過去了.又有一個利未人來到這地方,看見他,也照樣從那邊過去了.唯有一個撒馬利亞人行路來到那裡,看見他,就動了慈心,上前用油和酒倒在他的傷處,包裹好了,扶他騎上自己的牲口,帶到店裡去照應他.第二天拿出二錢銀子來,交給店主說:『你且照應他,此外所費用的,我回來必還你.』你想,這三個人哪一個是落在強盜手中的鄰舍呢?」他說:「是憐憫他的.」耶穌說:「你去照樣行吧.」

雅1:27 在上帝我們的父面前,那清潔沒有玷污的虔誠,就是看顧在患難中的孤兒寡婦,並且保守自己不沾染世俗.

男女有別卻同樣寶貴

文化有別但平等機會

天父聽到女孩子的哭聲

女性被嚴重剝削

* 女性佔全球人口的一半,但竟佔了貧窮人口的七成

* 全球的難民中亦有八成是女性

* 在－百三十個國家的議會裡,只有百分之十二的議員是女性；

* 女性佔全球工作人口的三成半,卻只擁有全球工資收入的一成

* 全球的物業中僅有百分之一是屬於女性的

不能接受這不平等

教會要為被剝削的人申冤

路4:16/19 耶穌來到拿撒勒,就是他長大的地方,在安息日,照他平常的規矩進了會堂,站起來要念聖經.有人把先知以賽亞的書交給他,他就打開,找到一處寫著說:「主的靈在我身上,因為他用膏膏我,叫我傳福音給貧窮的人,差遣我報告被擄的得釋放,瞎眼的得看見,叫那受壓制的得自由,報告神悅納人的禧年.」



\newpage

\section{成為狼群中的羊}
\label{sec:755}
\textbf{高國雄牧師 - 成為狼群中的羊}
\newline
\newline
\hyperref[sec:754]{< < < PREV SERMON < < <}
~
\hyperref[sec:toc]{[返目錄]}
~
\hyperref[sec:756]{> > > NEXT SERMON > > >}


路加福音 10:1-20
\newline
\begin{longtable}{cl}
\hline
\hline
章節 & 經文 (和合本修訂版)\\
\hline
10:1 & \begin{tabularx}{0.7\textwidth}{X} 這些事以後，主另外指定七十二個人，差遣他們兩個兩個地在他前面，往自己所要到的各城各地去。 \end{tabularx} \\ \\ \relax
10:2 & \begin{tabularx}{0.7\textwidth}{X} 他對他們說：「要收的莊稼多，做工的人少。所以，你們要求莊稼的主差遣做工的人出去收他的莊稼。 \end{tabularx} \\ \\ \relax
10:3 & \begin{tabularx}{0.7\textwidth}{X} 你們去吧！看！我差你們出去，如同羔羊進入狼群。 \end{tabularx} \\ \\ \relax
10:4 & \begin{tabularx}{0.7\textwidth}{X} 不要帶錢囊，不要帶行囊，不要帶鞋子；在路上也不要向人問安。 \end{tabularx} \\ \\ \relax
10:5 & \begin{tabularx}{0.7\textwidth}{X} 無論進哪一家，先要說：『願這一家平安。』 \end{tabularx} \\ \\ \relax
10:6 & \begin{tabularx}{0.7\textwidth}{X} 那裡若有當得平安的人，你們所求的平安就必臨到那家，不然，將歸還你們。 \end{tabularx} \\ \\ \relax
10:7 & \begin{tabularx}{0.7\textwidth}{X} 你們要住在那家，吃喝他們所供給的，因為工人得工錢是應當的；不要從這家搬到那家。 \end{tabularx} \\ \\ \relax
10:8 & \begin{tabularx}{0.7\textwidth}{X} 無論進哪一城，人若接待你們，給你們擺上甚麼食物，你們就吃甚麼。 \end{tabularx} \\ \\ \relax
10:9 & \begin{tabularx}{0.7\textwidth}{X} 要醫治那城裡的病人，對他們說：『神的國臨近你們了。』 \end{tabularx} \\ \\ \relax
10:10 & \begin{tabularx}{0.7\textwidth}{X} 無論進哪一城，人若不接待你們，你們就到大街上去，說： \end{tabularx} \\ \\ \relax
10:11 & \begin{tabularx}{0.7\textwidth}{X} 『就是你們城裡的塵土黏在我們的腳上，我們也當著你們擦去。但是，你們該知道神的國臨近了。』 \end{tabularx} \\ \\ \relax
10:12 & \begin{tabularx}{0.7\textwidth}{X} 我告訴你們，在那日子，所多瑪所受的，比那城還容易受呢！」 \end{tabularx} \\ \\ \relax
10:13 & \begin{tabularx}{0.7\textwidth}{X} 「哥拉汛哪，你有禍了！伯賽大啊，你有禍了！因為在你們中間所行的異能，若行在推羅、西頓，他們早已披麻蒙灰，坐在地上悔改了。 \end{tabularx} \\ \\ \relax
10:14 & \begin{tabularx}{0.7\textwidth}{X} 在審判的時候，推羅和西頓所受的，比你們還容易受呢！ \end{tabularx} \\ \\ \relax
10:15 & \begin{tabularx}{0.7\textwidth}{X} 迦百農啊，你以為要被舉到天上嗎？你要被推下陰間！」 \end{tabularx} \\ \\ \relax
10:16 & \begin{tabularx}{0.7\textwidth}{X} 耶穌又對門徒說：「聽從你們的就是聽從我；棄絕你們的就是棄絕我；棄絕我的就是棄絕差遣我來的那位。」 \end{tabularx} \\ \\ \relax
10:17 & \begin{tabularx}{0.7\textwidth}{X} 那七十二個人歡歡喜喜地回來，說：「主啊，因你的名，就是鬼也服了我們。」 \end{tabularx} \\ \\ \relax
10:18 & \begin{tabularx}{0.7\textwidth}{X} 耶穌對他們說：「我看見撒但從天上墜落，像閃電一樣。 \end{tabularx} \\ \\ \relax
10:19 & \begin{tabularx}{0.7\textwidth}{X} 我已經給你們權柄可以踐踏蛇和蠍子，又勝過仇敵一切的能力，絕沒有甚麼能害你們。 \end{tabularx} \\ \\ \relax
10:20 & \begin{tabularx}{0.7\textwidth}{X} 然而，不要因靈服了你們就歡喜，而要因你們的名記錄在天上歡喜。」 \end{tabularx} \\ \\
[1ex]
\hline
\hline
\end{longtable}
講題:成為狼群中的羊 Be Sheep among Wolves

經課:路加福音10:1-20

講員:高國雄牧師

場合:香港中文大學崇基學院禮拜堂主日崇拜

日期:2025年7月6日

1這事以後,主又設立七十個人,差遣他們兩個兩個地在他前面,往自己所要到的各城各地方去,

2就對他們說:「要收的莊稼多,做工的人少.所以,你們當求莊稼的主打發工人出去收他的莊稼.

3你們去吧!我差你們出去,如同羊羔進入狼群. 4不要帶錢囊,不要帶口袋,不要帶鞋；在路上也不要問人的安.

5無論進哪一家,先要說:『願這一家平安.』 6那裡若有當得平安的人,你們所求的平安就必臨到那家；不然,就歸與你們了.

7你們要住在那家,吃喝他們所供給的,因為工人得工價是應當的；不要從這家搬到那家.

8無論進哪一城,人若接待你們,給你們擺上甚麼,你們就吃甚麼.

9要醫治那城裡的病人,對他們說:『上帝的國臨近你們了.』

10無論進哪一城,人若不接待你們,你們就到街上去,

11說:『就是你們城裡的塵土黏在我們的腳上,我們也當著你們擦去.雖然如此,你們該知道上帝的國臨近了.』

16又對門徒說:「聽從你們的就是聽從我；棄絕你們的就是棄絕我；棄絕我的就是棄絕那差我來的.」

17那七十個人歡歡喜喜地回來,說:「主啊!因你的名,就是鬼也服了我們.」

18耶穌對他們說:「我曾看見撒旦從天上墜落,像閃電一樣.

19我已經給你們權柄可以踐踏蛇和蠍子,又勝過仇敵一切的能力,斷沒有甚麼能害你們.

20然而,不要因鬼服了你們就歡喜,要因你們的名記錄在天上歡喜.」

我們經歷了信仰的美好,自然亦會很想與他人分享.

而傳福音更是耶穌所托付門徒的大使命,我們當然不可推辭…

但如何傳福音呢?為何我們覺我很難?好像怎樣努力也沒有大果效呢?

讓我們回想當日耶穌如何差遣門徒,期望從中得到一些啟示.

經文背景:這事以後

路加福音9:

門徒被差派,耶穌用五餅二魚讓五千人吃飽,彼得認出耶穌是基督,耶穌預言他即將受害與復活,耶穌山上變像,耶穌趕鬼,門徒爭大,面對不被接待,面對不願意馬上放下一切跟從主的人…

被差派的情況:

V1主又設立七十個人,差遣他們兩個兩個地在他前面,往自己所要到的各城各地方去…

有伴,有目標

V2要收的莊稼多,做工的人少.所以,你們當求莊稼的主打發工人出去收他的莊稼.

工作量多,人不夠,要渴望更多人加入

實際操作:

v 4不要帶錢囊,不要帶口袋,不要帶鞋；在路上也不要問人的安.

不要帶錢囊=不要擔心經濟需要

不要帶口袋=不要擔心沒有食物而餓

不要帶鞋=不要擔心生活所需,現有的已足夠了

在路上也不要問人的安=不要被其他人或事轉移焦點(sidetrack)

V5無論進哪一家,先要說:『願這一家平安.』 6那裡若有當得平安的人,你們所求的平安就必臨到那家；不然,就歸與你們了.

自主動發出祝福,不計較回報

V7你們要住在那家,吃喝他們所供給的,因為工人得工價是應當的；不要從這家搬到那家. 8無論進哪一城,人若接待你們,給你們擺上甚麼,你們就吃甚麼.

放心被接待,為主被欣待是應當的事

不要按自己的喜好選擇被誰接待

V9要醫治那城裡的病人,對他們說:『上帝的國臨近你們了.』 10無論進哪一城,人若不接待你們,你們就到街上去,

11說:『就是你們城裡的塵土黏在我們的腳上,我們也當著你們擦去.雖然如此,你們該知道上帝的國臨近了.』

要醫治

要宣告,報喜也喜悲

結果:

17那七十個人歡歡喜喜地回來,說:「主啊!因你的名,就是鬼也服了我們.」 18耶穌對他們說:「我曾看見撒旦從天上墜落,像閃電一樣. 19我已經給你們權柄可以踐踏蛇和蠍子,又勝過仇敵一切的能力,斷沒有甚麼能害你們. 20然而,不要因鬼服了你們就歡喜,要因你們的名記錄在天上歡喜.」

歡歡喜喜

看到主的名的權能

耶穌調正焦點:要因你們的名記錄在天上歡喜!

耶穌的差遣:

V3你們去吧!我差你們出去,如同羊羔進入狼群.

不要成為狼,不要忘記自己是羊的身份

這世界是狼群,這是不能/需回避的事實

帶著主的教導及信心前往

那一日,我們要歡歡喜喜地回報成果



\newpage

\section{好撒馬利亞人}
\label{sec:756}
\textbf{李均熊牧師 - 好撒馬利亞人}
\newline
\newline
\hyperref[sec:755]{< < < PREV SERMON < < <}
~
\hyperref[sec:toc]{[返目錄]}
~
\hyperref[sec:757]{> > > NEXT SERMON > > >}


路加福音 10:25-37
\newline
\begin{longtable}{cl}
\hline
\hline
章節 & 經文 (和合本修訂版)\\
\hline
10:25 & \begin{tabularx}{0.7\textwidth}{X} 有一個律法師起來試探耶穌，說：「老師！我該做甚麼才可以承受永生？」 \end{tabularx} \\ \\ \relax
10:26 & \begin{tabularx}{0.7\textwidth}{X} 耶穌對他說：「律法上寫的是甚麼？你是怎樣念的呢？」 \end{tabularx} \\ \\ \relax
10:27 & \begin{tabularx}{0.7\textwidth}{X} 他回答說：「你要盡心、盡性、盡力、盡意愛主—你的神，又要愛鄰如己。」 \end{tabularx} \\ \\ \relax
10:28 & \begin{tabularx}{0.7\textwidth}{X} 耶穌對他說：「你回答得正確，你這樣做就會得永生。」 \end{tabularx} \\ \\ \relax
10:29 & \begin{tabularx}{0.7\textwidth}{X} 那人要證明自己有理，就對耶穌說：「誰是我的鄰舍呢？」 \end{tabularx} \\ \\ \relax
10:30 & \begin{tabularx}{0.7\textwidth}{X} 耶穌回答：「有一個人從耶路撒冷下耶利哥去，落在強盜手中。他們剝去他的衣裳，把他打個半死，丟下他走了。 \end{tabularx} \\ \\ \relax
10:31 & \begin{tabularx}{0.7\textwidth}{X} 偶然有一個祭司從那條路下來，看見他就從另一邊過去了。 \end{tabularx} \\ \\ \relax
10:32 & \begin{tabularx}{0.7\textwidth}{X} 又有一個利未人來到那裡，看見他，也照樣從另一邊過去了。 \end{tabularx} \\ \\ \relax
10:33 & \begin{tabularx}{0.7\textwidth}{X} 可是，有一個撒瑪利亞人路過那裡，看見他就動了慈心， \end{tabularx} \\ \\ \relax
10:34 & \begin{tabularx}{0.7\textwidth}{X} 上前用油和酒倒在他的傷處，包裹好了，扶他騎上自己的牲口，帶他到旅店裡去，照應他。 \end{tabularx} \\ \\ \relax
10:35 & \begin{tabularx}{0.7\textwidth}{X} 第二天，他拿出兩個銀幣來，交給店主，說：『請你照應他，額外的費用，我回來時會還你。』 \end{tabularx} \\ \\ \relax
10:36 & \begin{tabularx}{0.7\textwidth}{X} 你想，這三個人哪一個是落在強盜手中那人的鄰舍呢？」 \end{tabularx} \\ \\ \relax
10:37 & \begin{tabularx}{0.7\textwidth}{X} 他說：「是憐憫他的。」耶穌對他說：「你去，照樣做吧！」 \end{tabularx} \\ \\
[1ex]
\hline
\hline
\end{longtable}
講題:好撒馬利亞人 Good Samaritan

經課:路加福音10:25-37

講員:李均熊牧師

場合:香港中文大學崇基學院禮拜堂主日崇拜

日期:2025年7月13日

起:「好撒瑪利亞人」只是教我們做好人?

我想就算不是基督徒都會聽過「好撒瑪利亞人」.「好撒瑪利亞人」就是不問回報地幫助人,甚至是無親無故的陌生人.做基督徒當然要做好人行好事,但其他宗教信仰不也是教我們做好人行好事的嗎?那「好撒瑪利亞人」這個故事又有何特別呢?大部份信徒也可能聽過這段經文的解釋,例如祭司與利未人因為怕沾染死人的不潔所以才從路上的另一邊行過,又或者猶太人和撒瑪利亞人是不和的兩個族群,所以撒瑪利亞人願意幫助一個受傷的猶太人更顯寶貴.但這些解釋其實都沒有經文根據(參下),猶太人和撒瑪利亞人的關係惡劣也未必符合歷史事實(*加入學者解釋?).更大一個問題就是,若我們只將「好撒瑪利亞人」的好人好事只理解為道德榜樣,那其實是我們將這個故事抽離了經文的上文下理,完全忽略了耶穌講這個故事的目的.

承:「好撒馬利亞人」的比喻是關於律法的解釋

耶穌為何會講這個「好撒瑪利亞人」的比喻呢?經文一開始是說到有一個律法師起來要試探耶穌,就問了一個問題:「夫子!我該做甚麼才可以承受永生?」我們聽到律法師「試探」耶穌,可能就會想這名「律法師」一定不壞好意,不安好心要找耶穌把抦.但這可先入為主的觀感.的確有些時候法利賽人或文士會想找耶穌把抦控告耶穌,但這裡你會見耶穌並不抗拒這個律法師提問,也沒有怎樣的與他針鋒相對.甚至從還認為這個律法師說得對.其實我們不用首先假設律法師一定是壞旦,「試探」一定不懷好意.律法師是精研律法的讀書人,所謂「試探」tempted,也可以譯作較中性的「測試」tested.(就好像我常說自己是在神學院教舊約的,就有對舊約好有興趣的會友來問我問題,想「測試」一下我是否熟識舊約一樣.)耶穌用當時拉比慣用的「以問題答問題」的教學手法去問他「律法上寫的是甚麼?你念的是怎樣呢?」這名律法師的答了申命記6章:「你要盡心,盡性,盡力,盡意愛主你的神」和利未記19章的「愛鄰舍如同自己」 兩段經文涵蓋了愛神和愛人的原則,這個答案連耶穌也認為答得好.但若然這個律法師一心求道,他這樣自問自答,並未能「測試」耶穌的水平,於是乎他便再進一步追問耶穌:「誰是我的鄰舍呢?」耶穌才說出這個家傳戶曉的「好撒馬利亞人」的比喻.可見這個「好撒馬利亞人」的比喻是為著討論和解釋律法而說出來的,而不是單純的教導我們做好人.

轉:行出律法在於憐憫

「好撒馬利亞人」的比喻其實內容很簡單,這也是它家傳戶曉的原因吧.但正因為它的簡單,其實很多細節是沒有交代的.例如為甚麼祭司和利未人見到一個被人剝去衣裳,打得半死,會從路的另一邊行過呢?耶穌的比喻裡其實並沒有交代祭司和利未人這樣做的原因是甚麼.之前我說大部人認為是因為他們害怕接觸到死人而不潔的這個理由,是因為我們對祭司和利未人這兩類人作為聖殿供職的「人設」(人物設定)想當然的解釋.有人就指出經文寫明祭司正在這條路「下來」,即是表明他已經完成在耶路撒冷聖殿的供職,沒有太大需要關注是否潔淨的問題.而且在猶太人的典籍裡,也有要求在祭司要埋葬沒人認領的屍體的,所以一些人認為我們不應好像祭司和利未人般,看重宗教潔淨條例的律法要求,先於愛鄰如己的律法要求,未必是比喻的核心.當然這樣理解也算是一種提醒,但比對引入這個比喻的經文兩次交代律法師問耶穌問題的動機,耶穌的比喻並沒有交代祭司和利未人的行事動機,反而交代那個撒馬利亞人,因為見到這個有被害受苦而有急切需要的人,就「動了慈心」,正因為如此他才會行出律法的要求,愛憐舍如同自己,反而更令人覺得耶穌比喻強調的並不是不行使律法你需要考慮的理由,而是強調你看見別人有需要,就應該去行出來.全段經文出現最多次的字眼,就是「行出來」(v. 25 「我該做甚麼」/我要如何行 ；v. 28「你這樣行」 v. 37 「行憐憫的」 v. 37 「你去照樣行吧」).一開始,律法師的問題是問v. 25 「我該做甚麼」/我要如何行,但律法師明知要如何做去滿足律法的要求,耶穌也直接提出他是對的,只要行出來就可以了；但律法師卻跌入律法細節的思考,還問誰是我的憐舍?耶穌就用了這個「好撒馬利亞人」的比喻,提醒這個律法師當見到人落在困苦需要下,其實是不用細想眼前的這個受苦有需要的人值不值得我去行出律法愛人如己的要求,而是按著當時需要行出來就可以了!所以比喻的大部份篇幅不是描寫比喻中角色想了甚麼,而是他們做了甚麼:

31偶然有一個祭司從這條路下來,看見他就從那邊過去了.32又有一個利未人來到這地方,看見他,也照樣從那邊過去了.33惟有一個撒馬利亞人行路來到那裡,看見他就動了慈心,34上前用油和酒倒在他的傷處,包裹好了,扶他騎上自己的牲口,帶到店裡去照應他.35第二天拿出二錢銀子來,交給店主,說:『你且照應他；此外所費用的,我回來必還你.』

而我們了解這個比喻,太多時只在意比喻沒有說的是甚麼,是否我們也要求關愛別人需要一些理由呢?(我兒子名字「意行」的意思)

合:受苦者視角

耶穌在最後的反問中,其實明顯是將思考的視點完全反轉過來:36你想,這三個人哪一個是落在強盜手中的鄰舍呢?」你若真要細想律法的細節,那你不是從「我」的角度出發,而是需要從「落在強盜手中的」那人的角度出發.或許我們可以反問,為何那撒馬利亞人會動了慈心呢?不就是他看到了受苦受害者看到甚麼嗎?(慈心是甚麼?)



\newpage

\section{被歷史烙印的家庭}
\label{sec:757}
\textbf{林豪恩先生 - 被歷史烙印的家庭}
\newline
\newline
\hyperref[sec:756]{< < < PREV SERMON < < <}
~
\hyperref[sec:toc]{[返目錄]}
~
\hyperref[sec:758]{> > > NEXT SERMON > > >}


講題:被歷史烙印的家庭 The Family Marked by History

經課:何西亞書1:2-10

講員:林豪恩先生

場合:香港中文大學崇基學院禮拜堂主日崇拜

日期:2025年7月27日

一．引言

大家有沒有想過自己的名字有甚麼意思呢?會不會反映出時代特色呢?

央視網有一段文章,名為「70年中國人起名大全,背後全是時代的記憶!」:「我們每個人的名字,既承載著長輩的美好期望,也留存著鮮明的時代烙印.在不同時期,很多人的名字都集中於某些詞彙,某種意向的表達.聽到這些名字,我們就能猜到他們是那個年代的人!」

猜猜這些名字是甚麼年代:建國,解放,國慶,國強.他們是1949—1950年出生的人.據統計,叫建國的有近30萬人.

以下這些名字呢?抗美,援朝,衛國.這些是在1951—1953年出生的人,有上萬人叫張衛國,李衛國,王衛國,也有幾千人叫王援朝,張援朝,孫援朝.

這些呢?建設,超英,衛星,圖強,抗洪.這些是50年代後期常見名字.

要武,衛東,永紅呢?這些是60年代的名字了.

剛才大家讀舊約經課的時候,有沒有讀到一個為孩子改名的故事呢?這幾位但孩子的名子有甚麼特別呢?其實,故事中兩位父母的名字也有意思.讓我們再仔細閱讀經文.

二．何西亞與歌篾

1:1當烏西雅,約坦,亞哈斯,希西家作猶大王,約阿施的兒子耶羅波安作以色列王的時候,耶和華的話臨到備利的兒子何西阿.

這個故事的主角叫何西阿,何西亞是甚麼意思呢?何西亞的意思是「拯救」.這個名叫「拯救」的人,要拯救誰呢?如何拯救呢?

1:2耶和華初次與何西阿說話,對他說:「你去娶淫婦為妻,也收那從淫亂所生的兒女；因為這地大行淫亂,離棄耶和華.」

1:3於是,何西阿去娶了滴拉音的女兒歌篾.

「淫婦」這個詞語令人聽來不舒服.不過,在何西亞書中,這是一個角色.是甚麼角色呢?代表甚麼呢?

3:1耶和華對我說:「你再去愛一個淫婦,就是她情人所愛的；好像以色列人,雖然偏向別神,喜愛葡萄餅,耶和華還是愛他們.」

這個角色代表以色列,經文指出「以色列人偏向別神」,即是對耶和華不忠.何西亞書用婚姻來比擬耶和華與以色列的關係,當耶和華取了丈夫的角色,以色列自然就是妻子的角色.另一個用家庭關係來比喻的例子是記載在福音書「浪子回頭」故事,上帝取了父親的角色,人自然就是兒子的角色.

何西亞故事中的這個角色其實是有名字的,她叫歌篾.歌篾是甚麼意思呢?歌篾的意思是「完全的」.那就奇怪了,既然是「完全的」,怎會需要被拯救呢?這個名字到底是嘲諷,還是要表達甚麼呢?

正如剛才說過,何西亞和歌篾是兩個角色,這兩個角色分別代表耶和華及以色列.以色列稱自己是「完全的」,是甚麼令以色列如此自信呢?假設你是面試主持人,問求職者「你有甚麼缺點呢?」如果求職者回答「我的缺點就是沒有缺點」.你覺得如何呢?那個時代,以色列何以如此自信認為自己沒有欠缺,是「完全的」呢?且看他們當時的狀況:

10:1以色列是茂盛的葡萄樹,結果繁多.

果子越多,就越增添祭壇；

地土越肥美,就越造美麗的柱像.

10:3他們必說:我們沒有王,

因為我們不敬畏耶和華.

王能為我們做甚麼呢?

厲害了,以色列!經濟繁榮,國力強盛.或許就是因為如此強盛,他們不需要耶和華了,也不需要耶和華作他們的王了.以色列不是與耶和華立了約的嗎?或者當時的以色列認為,他們與耶和華的盟約只是一份歷史文件,已經沒有現實意義了.以色列覺得自己足夠強大撕毀盟約,足夠智慧選擇自己喜歡的神明來敬拜.當以色列以為自己「完全了」的時候,先知看到是卻是另一面.

7:3他們行惡使君王歡喜,

說謊使首領喜樂.

5在我們王宴樂的日子,

首領因酒的烈性成病；

王與褻慢人拉手.

正所謂:上樑不正下樑歪,君王首領行為不正,下面的人也跟著行惡說謊,形成背叛公義和誠信的惡劣環境.君王首領熱衷於聲色犬馬,沉迷於宴樂,寵信奸佞,任用小人,其實岌岌可危,正如施洗約翰所說:「斧子已經放在樹根上」.先知心知以色列的確需要被拯救,然而,既然自視為「完全的」,怎會承認自己需要被拯救呢?一個飲醉了酒的人堅稱自己沒醉,還去駕車,接下來會發生甚麼事情呢?

三．第一個孩子「耶斯列」

1:3這婦人懷孕,給他生了一個兒子.

1:4耶和華對何西阿說:「給他起名叫耶斯列；因為再過片時,我必討耶戶家在耶斯列殺人流血的罪,也必使以色列家的國滅絕.

「耶斯列」這個名字是「上帝播種」的意思,而「耶斯列」的讀音跟「以色列」相似.從根源上來說,以色列的確是耶和華所播種的.同時,「耶斯列」又是一個地方名,當年以色列的亞哈王和耶洗別以耶斯列作為大本營,在全國鼓吹敬拜巴力.耶斯列又是一處殺戮的地方,耶戶就是在那裡殺死當時的以色列王及猶大王(王下9:27),並且對以色列王室大開殺戒(王下10:12-14),斬草除根.

何西亞的大兒子叫耶斯列這個語帶雙關的名字,是諷刺以色列嗎?還是審判以色列呢?耶斯列的意思是上帝播種.現實上,以色列在耶斯列這個地方播下甚麼種子呢?是不是背叛的種子?是不是殺戮的種子呢?何西亞的第一個孩子「耶斯列」是一個角色,這個名字讀出來就令人想到以色列,想到以色列在耶斯列這個地方所做的惡事.這些惡事,人間好像無可奈何,但耶和華不會視而不見,總會對這段歷史給出一個回答.

四．第二個孩子「羅‧路哈瑪」

1:6歌篾又懷孕生了一個女兒,耶和華對何西阿說:「給她起名叫羅‧路哈瑪；因為我必不再憐憫以色列家,決不赦免他們.

「路哈瑪」是「憐憫」的意思,「羅」是否定詞,意思是「不」,何西亞的女兒叫「羅‧路哈瑪」,就是「不憐憫」的意思.這也是一個角色,不代表何西亞的女兒不蒙憐憫,而是代表以色列不蒙憐憫.究竟以色列聽不聽得明白呢?以色列有甚麼反應呢?人在甚麼時候感到自己需要被憐憫的呢?如果自覺強大,怎會需要被憐憫呢?話說回到,就算耶和華不憐憫又有甚麼問題呢?此時的以色列已經精明地在宗教上「分散投資」,耶和華不憐憫,還有巴力,還有其他供奉的神明.

五．第三個孩子「羅‧阿米」

1:8歌篾給羅‧路哈瑪斷奶以後,又懷孕生了一個兒子.

1:9耶和華說:「給他起名叫羅‧阿米；因為你們不作我的子民,我也不作你們的上帝.」

「阿米」是「我的民」的意思,加上「羅」這個否定詞,變成「不是我的民」.「羅‧阿米」都是一個角色,不代表耶和華與這位孩子「割蓆」,而是警告以色列耶和華準備「割蓆」.以色列聽不聽得明白呢?以色列有甚麼反應呢?其實,在耶和華警告與以色列割蓆之前,以色列已經先行一步與耶和華割蓆了.他們一方面認為自己不需要耶和華作王,另一方面他們懂得「分散投資」,耶和華不再是以色列獨一的上帝,也不是唯一的依靠.

何西亞家庭的故事其實是以色列的故事,何西亞家庭的人名折射出以色列的歷史.我們曾經談及現代人的名字往往也有歷史的痕跡,其實就算名字沒有歷史的烙印,也不等於家庭不受歷史影響.當我們追溯自己的家庭最近三代的變遷,就足以發現歷史在家庭中的痕跡了.

六．歷史在我們的家庭的烙印

1945年二次大戰結束的時候,香港有多少人口呢?有資料說只有50—75萬.到1960年呢?增加到300萬,1970年400萬,1980年500萬.如此大量的人口從哪裡來呢?大家都知道是來自深圳河以北.香港不少家庭的祖父輩都是這個故事的一分子.被稱為香港四大才子的金庸,黃霑,倪匡,蔡瀾都不在香港出生,其中倪匡在1957年才來到香港.

倪匡的家庭變遷,正是時代的縮影.維基百科說倪匡出生於上海,有兄弟姐妹七人,他排行第四.大哥倪亦方留在遼寧,五弟倪亦平在香港,六妹倪亦舒即是小說作家亦舒,後來移居加拿大,七弟倪亦靖在新加坡.為甚麼一家人分散在不同地方生活呢?很可能是歷史的力量在推動.其他家庭就算沒有倪家庭那麼「國際化」,也可能有「兩頭家」,一個在香港,一個在內地.還記得那些年,人們在農曆新年前挑著大包小包回鄉探親的景象嗎?

70年代和80年代,香港很多工廠,後來漸漸減少,為甚麼呢?因為工廠北移,老闆北上開廠,師傅北上開工.這些北上的大多數是男人,於是又出現另一個時代的另一種香港和內地「兩頭家」的社會現象.那個時代的烙印深深地影響了牽涉其中的家庭,有些糾結至今仍未解開,有些眼淚而今仍未流乾.

倪匡在香港結婚生仔,寫作出版,相信在座不少人聽過「衛斯理」,或者「原振俠」.1992年,倪匡與妻子在移居美國.倪匡再次遷移,其實也是時代的縮影.這個歷史時空之下的「兩頭家」,變成了香港和美國,或香港和加拿大了.倪匡最後在哪裡安息主懷呢?在香港.他不是移居了美國嗎?回流,也曾經是一個時代的主題.

小學六年級的時候,老師給我們一題永恆的作文題目,名為「我的志願」,不知道你有沒有寫過呢?你還記得當年的志願嗎?當時被老師讀出來的志願非常特別,因此現在還記得.那位同學的志願是長大後找份工作,結婚生仔,安居樂業.當時其他同學聽到都捧腹大笑.同們學笑,是因為太年輕了.少年不知愁滋味,你以為成為立室,安居樂業是容易的事嗎?你以為一家人整整齊齊生活在一起是必然的嗎?你不知道有多少家庭在歷史的洪流中被沖散?有多少家庭在社會的動盪中被拋來拋去呢?

七．歷史無情,上主有愛

從前有一個人名叫拯救,他娶了一位名叫完全的人,生了三個兒女,第一個叫上帝播種,第二個叫不蒙憐憫,第三個叫不是我民.我們不知道這個可能被視為古怪的家庭後來有甚麼遭遇,我們知道的是這個家庭承載了上帝給以色列的信息.上帝要對以色列表達甚麼呢?期望自以為完全的以色列醒悟嗎?警告以色列上帝會拔出祂所栽種的嗎?宣告不憐憫以色列嗎?登報聲明和以色列脫離關係嗎?都可能是.然而,這個帶著歷史烙印的家庭所要告訴世人的不只於此,其實還有更重要的:

2:23我必將她種在這地.素不蒙憐憫的,我必憐憫；本非我民的,我必對他說:你是我的民；他必說:你是我的上帝.」

「我必將她種在這地」,回應了第一個孩子的名字,上帝的本意不是「拔出」,而是「栽種」,就算拔出,必再栽種.

「素不蒙憐憫的,我必憐憫」,回應了第二個孩子的名字,憐憫才是上帝本來的性情.

「本非我民的,我必對他說:你是我的民；他必說:你是我的上帝」,回應了第三個孩子的名字.正如詩人說:「他的怒氣不過是轉眼之間,他的恩典乃是一生之久」,上帝不會放棄與人的關係.

有時,人事在躁動,歷史在盲動,家庭在歷史的洪流中彷彿被無目的地拋擲,承受著歷史無情的惡果.然而,何西亞的故事告訴人們,盲動的歷史之上有一位鑑察世情的上帝,無情的歷史之上有一位始終充滿憐憫的上帝,流離失所的世間之上有一位尋找失散迷羊的牧者.

處身於歷史之中的家庭,或許沒有力量力挽狂瀾,也未必有能力扭轉乾坤,然而,如果我們知道上主的心意是栽種,是憐憫,是和好,就讓我們朝著上主的方向前行,在我們的家庭培養互相建立,彼此憐憫,以情相繫的氣氛.我們承認自己永遠不會是完全的,我們仰望那位完全的上主接納我們,拯救我們.



\newpage

\section{追啊追}
\label{sec:758}
\textbf{梁遠耀牧師 - 追啊,追}
\newline
\newline
\hyperref[sec:757]{< < < PREV SERMON < < <}
~
\hyperref[sec:toc]{[返目錄]}
~
\hyperref[sec:759]{> > > NEXT SERMON > > >}


路加福音 12:13-21
\newline
\begin{longtable}{cl}
\hline
\hline
章節 & 經文 (和合本修訂版)\\
\hline
12:13 & \begin{tabularx}{0.7\textwidth}{X} 人群中有一個人對耶穌說：「老師！請你吩咐我的兄弟和我分家產。」 \end{tabularx} \\ \\ \relax
12:14 & \begin{tabularx}{0.7\textwidth}{X} 耶穌對他說：「你這個人！誰立我作你們的判官，或給你們分家產的呢？」 \end{tabularx} \\ \\ \relax
12:15 & \begin{tabularx}{0.7\textwidth}{X} 於是他對他們說：「你們要謹慎自守，躲避一切的貪心，因為人的生命不在於家道豐富。」 \end{tabularx} \\ \\ \relax
12:16 & \begin{tabularx}{0.7\textwidth}{X} 然後他用比喻對他們說：「有一個財主，田地出產豐富。 \end{tabularx} \\ \\ \relax
12:17 & \begin{tabularx}{0.7\textwidth}{X} 他自己心裡想：『我的出產沒有地方儲藏，怎麼辦呢？』 \end{tabularx} \\ \\ \relax
12:18 & \begin{tabularx}{0.7\textwidth}{X} 就說：『我要這麼辦：要把我的倉庫拆了，另蓋更大的，在那裡好儲藏我一切的糧食和財物， \end{tabularx} \\ \\ \relax
12:19 & \begin{tabularx}{0.7\textwidth}{X} 然後要對我自己說：你這個人哪，你有許多財物積存，可供多年享用，只管安安逸逸吃喝快樂吧！』 \end{tabularx} \\ \\ \relax
12:20 & \begin{tabularx}{0.7\textwidth}{X} 神卻對他說：『無知的人哪！今夜就要你的性命，你所預備的要歸誰呢？』 \end{tabularx} \\ \\ \relax
12:21 & \begin{tabularx}{0.7\textwidth}{X} 凡為自己積財，在神面前卻不富足的，也是這樣。」 \end{tabularx} \\ \\
[1ex]
\hline
\hline
\end{longtable}
講題:追啊,追 Chasing, always chasing

經課:路加福音12:13-21

講員:梁遠耀牧師

場合:香港中文大學崇基學院禮拜堂主日崇拜

日期:2025年8月3日

一,問安與祈禱

各位弟兄姊妹及會眾,願主所賜的平安,常與你們同在!

在分享主道前,先請我們一同祈禱:「天父願我口中的語言,眾人心裡的意念,皆蒙你悅納.奉主耶穌基督名求,誠心所願!」

二,引言

1975年香港電影「遊戲人間三百年」(Enjoy Longevity-300Years)(導演:楊權,主演:譚炳文),以荒誕幽默的方式,講述男主角意外地救了一個神仙,神仙賜給男主角有取之不盡的財富,男主角更加要求可以長生不老,反映人的心思,就是不斷追求更多,不斷追啊,追.故事中的男主角因為財富與長生不老,結果為他帶來很多痛苦.這齣電影探討了一個嚴肅的哲學問題——物質與生命的有限與無限,帶出的訊息是知足.

想像一下,有一天你獲得一筆突然而來的遺產,足夠你買下一間豪宅,但接下來的問題是,這筆錢能保證你永遠幸福嗎?可能有些人會興奮地:「當有可以!」但今天,我們從耶穌的比喻中會看到一個截然不同的答案.這段經文提醒我們,生命中真正的富足,並不在於銀行帳戶的數字,而在於我們與上帝的.準備好了嗎?讓我們一起來探索耶穌所說的「天國中的富有」

三,經文的背景與解釋

這段經文的背景是耶穌向群眾講道時,一個人突然向祂請求,要求祂幫忙分家產(路12:13).耶穌並沒有直接解決他的問題,首先反問這個人:「你這個人!誰立我作你們的判官,或給你們分家產的呢?」(路12:14)從字面看,耶穌以反問的方式,表達自己的工作,不是要為人分家產.在前一段經文,即12章8至12節,耶穌教導人要在別人面前要認他,即肯定耶穌的身份,與及耶穌與信徒之間的關係.

故此,今日所選讀的福音經課中的財主比喻,藉此機會警告人們不要貪心,並透過比喻來教導真正的財富之道外,應該也是要肯定耶穌的身份及能力.

在這個比喻中,一個富足的人收成豐盛,他心想:「我要怎樣辦?因為我的倉房不夠放我的糧食.」(路12:17).於是,他決定拆掉舊有的倉庫,建造更大的倉房,並對自己說:「靈魂哪,你有許多財物積存,可供多年使用；只管安安逸逸地吃喝快樂吧.」(12:19).然而,上帝卻對他說:「無知的人哪,今夜必要你的靈魂,你所預備的要歸誰呢?」(12:20)「今夜必要你的靈魂」直譯是「他們今夜必要向你取回你的靈魂」.「他們」可指上帝或祂的使者,但由於這句話是上帝所說,所以更可能是指財主的「許多財物」,是財主的「許多財物」取去了財主的靈魂.耶穌最後總結:「凡為自己積財,在上帝面前卻不富足的,也是這樣.」(12:21).

正如剛才所說,耶穌說這個背喻的背景,是有一個人想請耶穌幫他的家庭分家產,做一個仲裁人.我認為這個人只是把耶穌看成一個德高望重,受人景仰及熟識律法的律法教師,卻沒有把耶穌視為彌賽亞或是生命的主宰.

我認為在這個比喻,耶穌並非全然否定人去追求物資上財富,耶穌提醒我們,追求世上的財富,但不要忽略了靈命上的富足,要認定上帝在他生命中的安排.比喻中的財主,他的計劃看似合理,卻缺少了一個關鍵因素——他未曾考慮上帝在他生命中的旨意,也未曾想到財富應如何用來祝福他人.因此,他的生命突然結束,他所積存的一切,都無法陪伴他進入永恆.

四,財主比喻中的財富觀

從在這個財主的比喻中,帶給我們對財富有以下的教導.

1. 財富的主權

財主以物質成為他的保障,卻沒有意識到生命的有限,更沒有認定生命的主權是在上帝手裡.我們要認識到世上的財富無法真正帶來安全感,因為它是短暫的,一旦生命終結,這些財物就不再屬於我們.

2. 屬靈財富才是永恆的

耶穌指出真正的富足是在上帝裡的富足,而非世俗的積累.真正的財富是信仰的堅定,愛的施予和上帝國度的事工.馬太福音6:19-20也提醒我們:「不要為自己積蓄財寶…要在天上積蓄財寶….」

3. 生命的優先次序

富足人只考慮如何讓自己生活舒適,卻忽略了靈性的需求.信徒要學習以上帝為中心,認識到生命的價值不在於物質,而是與上帝建立深厚的關係,並用所擁有的財富去祝福他人.

4. 服事與奉獻

財富的真正價值在於如何使用它.如果我們只把財富留給自己,而忽略了幫助有需要的人,那麼這些財富就只是虛空.我們應該運用上帝所賜的資源來做善工,使財富成為祝福而非束縛.

5. 預備生命的終結

富足人沒有想到自己會突然離世,而耶穌提醒我們要隨時預備迎見上帝.我們要問自己:「我是否只專注於世上的事,而忽略了永恆的價值?」我們應當時刻預備自己,活出有意義,合上帝心意的生命.

6. 以李清詞牧師為例

以故的中華基督教會前副總幹事李清詞牧師(Rev. Lee Ching Chee),她在生前曾向我們這些後輩分享,她小時候生長在一個富裕家庭,有一次,她的父親帶她進入家中的一個房間,房間lor 滿了各種古董,李牧師的父親向幼小的李牧師表示,你只要拿其中一樣,都足夠你一生過活.只是後來,因為戰亂,李牧師家道中落,沒有了那些古董,有些家人也失散了.但幸好李牧師認識耶穌,一生在心裡富足中渡過.

五,建議

讓我重申,這段比喻並非否定財富,而是要提醒我們,財富是上主的賜予.人要從中學校管理,包括累積財富和分享財富,多作善工,榮神益人.以下我有一些簡單的建議.

1. 學習理財

首先,我認為信徒都應該要學習理財,把天父所賜的財富好好管理.學習理財不單是學習理財知識,也是學習控制自己的慾望.每個人都要有適當的儲蓄或保障,讓自己和身體人在需要時,得到多一點的方法去面對.

2. 養成奉獻的習慣

第二,信徒除了學習理財和累積,也是在生活中,養成金錢奉獻的習慣.因為我們若能在平常日子,已經有養成金錢奉獻的習慣,我們就會較容易建立自己對財富的價值觀.

3. 遺贈／遺產奉獻

近年,我聽閒不少信徒在訂立遺囑／平安紙時,會考慮把部分遺產奉獻給教會或其他機構.在我們離開物質世界時,透過遺囑把財富捐獻出來,這是一種在生命終結時,透過遺囑,向世人宣告:「萬物都是從主而來,我們把從主要來的,獻給主.」

六,結語及祈禱

有一個故事,有一位媽媽在某一天的早上,喚醒她的兒子,叫他起床.但這位兒童不肯起床.媽媽問兒子:「請你給我一個理由,你今天為甚麼不起床上學.」兒子說:「因為學校裡同學,不跟我一起玩耍.」媽媽聽到後,便向這位兒子說:「原來是這樣,但我卻給你兩個理由,是你必須立即起床預備上學的.」兒子便說:「是哪兩個理由?」媽媽說:「第一,因為你已經四十多歲了；第二,因為你是校長.」

在人生不同階段,我們可能有不同的承擔和責任.按著上帝的心意,有時我們要學習積累,有時我們也要學習分享.耶穌透過這個比喻提醒我們,我們生命的中心,不是財富,而是上帝.不要只為自己積攢財富,而要在上帝面前富足.真正的財富不是銀行存款或物質擁有,而是我們在信仰上的成長,對他人的愛心,以及對上帝國度的奉獻.

財主最大的錯誤不是他的財富,而是他的態度.他忘記了財富的真正目的,也未曾思考生命的有限性.我們應當學習在生活中尋求上帝的智慧,將財富視為上帝的恩典,而非驕傲的憑藉.

願我們每位信徒都能在世上過一個忠心於上帝的生活,不只是為自己積財,而是成為祝福的管道,並且預備好迎接上帝的永恆應許.

請我們再次祈禱:「天父感謝你,因為你創造世界,也讓我們管理物質世界,但有時我們忘記了自己只是管家,沒有學校積累財富或分配財富,求你赦免我們.願我們能因著主耶穌你的教導,過以基督為中心的生活,以顯明你是萬有之主,一切的美善都是從你而來.奉主名求,誠心所願!」



\newpage

\section{警醒的僕人}
\label{sec:759}
\textbf{李均熊牧師 - 警醒的僕人}
\newline
\newline
\hyperref[sec:758]{< < < PREV SERMON < < <}
~
\hyperref[sec:toc]{[返目錄]}
~
\hyperref[sec:760]{> > > NEXT SERMON > > >}


路加福音 12:32-40
\newline
\begin{longtable}{cl}
\hline
\hline
章節 & 經文 (和合本修訂版)\\
\hline
12:32 & \begin{tabularx}{0.7\textwidth}{X} 你們這小群，不要懼怕，因為你們的父樂意把國賜給你們。 \end{tabularx} \\ \\ \relax
12:33 & \begin{tabularx}{0.7\textwidth}{X} 你們要變賣財產賙濟人，為自己預備永不壞的錢囊和用不盡的財寶在天上，就是賊不能近，蟲不能蛀的地方。 \end{tabularx} \\ \\ \relax
12:34 & \begin{tabularx}{0.7\textwidth}{X} 因為你們的財寶在哪裡，你們的心也在哪裡。」 \end{tabularx} \\ \\ \relax
12:35 & \begin{tabularx}{0.7\textwidth}{X} 「你們要束緊腰帶，燈也要點著， \end{tabularx} \\ \\ \relax
12:36 & \begin{tabularx}{0.7\textwidth}{X} 好像僕人等候自己的主人從婚宴上回來。他來叩門，就立刻給他開門。 \end{tabularx} \\ \\ \relax
12:37 & \begin{tabularx}{0.7\textwidth}{X} 主人來了，看見僕人警醒，那些僕人就有福了。我實在告訴你們，主人會叫他們坐席，自己束上腰帶，前來伺候他們。 \end{tabularx} \\ \\ \relax
12:38 & \begin{tabularx}{0.7\textwidth}{X} 他或是半夜來，或是天亮之前來，看見僕人這樣，那些僕人就有福了。 \end{tabularx} \\ \\ \relax
12:39 & \begin{tabularx}{0.7\textwidth}{X} 你們要知道，一家的主人若知道賊甚麼時候來，就不容賊挖穿房屋。 \end{tabularx} \\ \\ \relax
12:40 & \begin{tabularx}{0.7\textwidth}{X} 你們也要預備，因為在你們想不到的時候，人子就來了。」 \end{tabularx} \\ \\
[1ex]
\hline
\hline
\end{longtable}
講題:警醒的僕人 Watchful Servants

經課:路加福音12:32-40

講員:李均熊牧師

場合:香港中文大學崇基學院禮拜堂主日崇拜

日期:2025年8月10日

起:奇怪的分段與上文下理

今日的福音經課以路加福音12章經文為選段,開始於12:32-34 一小段關於主耶穌對財寶的忠告:「你們這些小群,不要懼怕,因為你們的父樂意把國賜給你們.你們要變賣財產賙濟人,為自己預備永不壞的錢囊和用不盡的財寶在天上,就是賊不能近,蟲不蛀的地方.因為你們的財寶在哪裡,你們的心也在哪裡.」這一小段其實是大家十分熟識的一段耶穌教導的結語:

22耶穌又對門徒說:「所以,我告訴你們,不要為生命憂慮吃甚麼,為身體憂慮穿甚麼.23因為生命勝於飲食,身體勝於衣裳.24你們想一想烏鴉:牠們既不種也不收,既沒有倉又沒有庫,神尚且養活牠們.你們比飛鳥要貴重得多呢!25你們哪一個能藉著憂慮使壽數多加一刻呢?26這最小的事你們尚且不能做,何必憂慮其餘的事呢?27你們想一想百合花是怎麼長起來的:它也不勞動,也不紡線.然而我告訴你們,就是所羅門極榮華的時候,他所穿戴的還不如這些花的一朵呢!28你們這小信的人哪!野地裡的草今天還在,明天就丟在爐裡,神還給它這樣的妝飾,何況你們呢?29你們不要求吃甚麼,喝甚麼,也不要掛慮.30這都是世上的外邦人所求的；你們需要這些東西,你們的父都知道.31你們只要求他的國,這些東西就必加給你們了.32你們這小群,不要懼怕,因為你們的父樂意把國賜給你們.33你們要變賣財產賙濟人,為自己預備永不壞的錢囊和用不盡的財寶在天上,就是賊不能近,蟲不能蛀的地方.34因為你們的財寶在哪裡,你們的心也在哪裡.」

但今日的經課選段,卻硬生生將耶穌教導的結語,連結一段耶穌提醒我們在主再來前我們要警醒的「警醒僕人」的比喻.這樣不單將兩段看似主題不同的經文夾硬連在一起,也將耶穌連續兩個僕人比喻(路12:35-48)抽離了原來的脈絡.再看看舊約,詩篇,和書信經課的選段,幾段的主題也沒有明顯的關係,看來這段經課選段並非互補性的選段.當然我可以跳過這個疑難,只集中講其中一段便可.不過有時這些困難,很大可能就是經文編輯者要我們留意的地方.

承:生活的恐懼與主的再來

時間關係,這裡不能做一個全面的分析.不過大家若將12章速讀一次(或只看和合本的標題),就會發覺耶穌一開始提醒門徒要防備法利賽人的假冒為善,然後提出不要懼怕殺身體而不能做甚麼的.可見當時的信徒落入一個危險的處境,很大機會是公元七十年的第一次羅馬與猶太人之戰,羅馬軍隊圍剿耶路撒冷和燒燬聖殿之後,猶太人並沒有得到法利賽人所代表的建制力量的保護,而信徒又正受羅馬多米田政府的逼害(當中細節大家可參考網上資源).信徒正是活在惶恐和焦慮當中.面對這一種懼怕,信徒的反應是甚麼呢?從經文推斷,很大機會他們轉為依賴錢財,以為財寶可以在此亂世中提供保障.但耶穌的教導正要指出錢財並不能為他們生活提供保護(還記得上星期路12:13-21那個無知財主的比喻嗎?).這也是為何主耶穌要對門徒說:「32你們這小群,不要懼怕,因為你們的父樂意把國賜給你們.33你們要變賣財產賙濟人,為自己預備永不壞的錢囊和用不盡的財寶在天上,就是賊不能近,蟲不能蛀的地方.34因為你們的財寶在哪裡,你們的心也在哪裡.」

另一方面,在面臨困難迫害的當下,信徒可能更加渴望主的再來,平定一切世上的不公和殘暴.但路加福音背後的讀者群體卻可能感到人子再來的日子遲遲沒有發生,且感覺這日子好像越來越遠,亦正因面對迫害困苦,當下失去等待的信心,不再持守主的道,以為主再來只是遙不可及的空想.於是乎信徒生活更加惶恐不安,反而更想依靠物質錢財的保障了(之後一段的路12:41-48更提到管家(即是信徒領袖?)變本加厲惡待同伴).於是乎耶穌就說了警醒僕人的比喻,鼓勵門徒人子必然再來,只是沒有人知道何時；但若主忽然來到,發現大家在警醒等候,「37主人來了,看見僕人警醒,那些僕人就有福了.我實在告訴你們,主人會叫他們坐席,自己束上腰帶,前來伺候他們.」

所以今日經課的主題,正是在因難失望的日子中不要懼怕,「因為你們的父樂意把國賜給你們」,因此也不用依靠金錢財寶去保障自己.「因為在你們想不的時候,人子就來了」.雖然今天面對壓迫與懼怕,我們也要警醒預備,因為當「主人來了,看見僕人警醒,那些僕人就有福了.」

轉:焦慮與乏味等待(Drudgery)

今日我們在香港做基督徒當然不會如第一,二世紀的信徒們,面對社會中種種的中傷和迫害,但這不表示我們今天當信徒就能免卻當時路加福音背後群體感受的焦慮與失落.無可否認,今天香港人大多生活都充滿焦慮和不安.而在商業消費社會中,人也容易以金錢作為自己的保障.最近我和我的保險經紀見面,其實主要討論我兒子的教育經費的問題,但她在會面時也不忘建議我再買一些儲蓄保險去為將來退休生活作保障.我在剛開始工作時其實已經有買保險,但我的保險經紀卻說當時的保障只到65歲,但現時大部份的保險產品都保障到80歲以後.人口老化,經濟下行,社會氣氛不明朗,這些都增加我們生活的焦慮和不安,若個身邊有個錢防身,的確可以有多些安全感.但這些是否真能令我們不再懼怕呢?我不是叫大家就不要在財務上作出計劃,在現代商業消費社會,如何管理財富是一個很重要的課題,希望大家也不要忽略做一個耶穌恩賜的好管家.但就像我以住的保險可以保到六十多歲,但在今天卻看來已經不管用,要買就買保到八十多歲的產品,看來金錢的價值能為人帶來安心的效果是越來越少；更可況人生多變,正如上主日經課的經文題醒我們:「『無知的人哪!今夜就要你的性命,你所預備的要歸誰呢?』21凡為自己積財,在神面前卻不富足的,也是這樣.」

在信徒面對極大的困難和焦慮中,主耶穌在這一章中,就不斷提醒他們:「五隻麻雀不是賣二銅錢嗎?但在神面前,一隻也不被忘記；就是你們的頭髮也都數過了.不要懼怕,你們比許多的麻雀還貴重!」(12:6-7);「你們要謹慎自守,躲避一切的貪心,因為人的生命不在於家道豐富.」(12:15)若你也感到生活在焦慮不安中,願你今天也在慈愛的主面前找到平安.

有時焦慮的源頭並非來自生活的不足和缺乏,而是在信仰上感到上主的遠離.難道我們理性上不知道慈愛的上主會看顧我們嗎?我們不知道我們的價值比麻雀貴重嗎?我們真的以為錢財在變幻莫測的世代裡真能保障我們嗎?就如12章內耶穌多番提醒門徒不要懼怕,但到後來要進一步以警醒僕人比喻再一次鼓勵他們「你們要束緊腰帶,燈也要點著,好像僕人等候自己的主人從婚宴上回來.他來叩門,就立刻給他開門.」

作為信仰基督的我們,真正令我們煩惱的,可能是感到耶穌的應許好像沒有實現過,說好要再來,要賜世人平安,但世界仍然充斥戰亂不公,謊言奸惡仍然當道.我們會否在世生活,慢慢因為平淡無期的等待,以至失去依靠的動力,盼望的信心呢?耶穌用上這個警醒僕人比喻,就是鼓勵你我,不要因長久等候而失去信心.無論遲早,主人都會從婚宴上回來.那預備好迎見上主的,主人必會叫他們坐席,自己束上腰帶,前來伺候他們.(例子:等開學?)

合:警醒的忠告

但警醒預備,其實是做甚麼.這個在經文中其實是十分明顯:

33你們要變賣財產賙濟人,為自己預備永不壞的錢囊和用不盡的財寶在天上,就是賊不能近,蟲不能蛀的地方.34因為你們的財寶在哪裡,你們的心也在哪.

主耶穌認為真正的財寶,並非有盡而又會壞掉的金錢,反而是將今世的財寶變賣了去幫助有需的人(賙濟人).能抗衡心中恐懼,在主內真正的平安,並不是透過擁有多少而得到,而是透過捨去和分享而產生的.主耶穌警醒僕人的比喻,正是叫我們預備好隨時放手棄去,讓我們所擁有的成為別人的幫助和祝福.



\newpage

\section{絕望中的希望}
\label{sec:760}
\textbf{鄧瑞強博士 - 絕望中的希望}
\newline
\newline
\hyperref[sec:759]{< < < PREV SERMON < < <}
~
\hyperref[sec:toc]{[返目錄]}
~
\hyperref[sec:761]{> > > NEXT SERMON > > >}


哈巴谷書 1:1-3:19
\newline
\begin{longtable}{cl}
\hline
\hline
章節 & 經文 (和合本修訂版)\\
\hline
1:1 & \begin{tabularx}{0.7\textwidth}{X} 哈巴谷先知所看見的默示。 \end{tabularx} \\ \\ \relax
1:2 & \begin{tabularx}{0.7\textwidth}{X} 耶和華啊，我呼求，你不應允，要到幾時呢？我向你呼喊「暴力！」你還不拯救？ \end{tabularx} \\ \\ \relax
1:3 & \begin{tabularx}{0.7\textwidth}{X} 你為何使我看見罪孽？你為何坐視奸惡呢？毀滅和兇暴在我面前，爭執與紛爭不斷發生。 \end{tabularx} \\ \\ \relax
1:4 & \begin{tabularx}{0.7\textwidth}{X} 因此律法無效，公理從未彰顯。惡人圍困義人，所以公理遭受扭曲。 \end{tabularx} \\ \\ \relax
1:5 & \begin{tabularx}{0.7\textwidth}{X} 你們要向列國觀看，注意看，要驚奇，再驚奇！因為在你們的日子，有一件事發生，儘管有人說了，你們還是不信。 \end{tabularx} \\ \\ \relax
1:6 & \begin{tabularx}{0.7\textwidth}{X} 看哪，我必興起迦勒底人，就是那殘忍暴躁之民，通行遍地，霸佔不屬自己的住處。 \end{tabularx} \\ \\ \relax
1:7 & \begin{tabularx}{0.7\textwidth}{X} 他威武可畏，審判與威權都由他而出。 \end{tabularx} \\ \\ \relax
1:8 & \begin{tabularx}{0.7\textwidth}{X} 他的馬比豹更快，比晚上的野狼更猛。他的戰馬跳躍，他的戰馬從遠方而來；他們飛跑，如鷹急速抓食， \end{tabularx} \\ \\ \relax
1:9 & \begin{tabularx}{0.7\textwidth}{X} 都為施行殘暴而來，他們的臉面向東，聚集俘虜，多如塵沙。 \end{tabularx} \\ \\ \relax
1:10 & \begin{tabularx}{0.7\textwidth}{X} 他譏誚列王，嘲諷領袖，嗤笑一切堡壘，堆土攻取它。 \end{tabularx} \\ \\ \relax
1:11 & \begin{tabularx}{0.7\textwidth}{X} 那時，他如風猛然掃過，他背叛，顯為有罪；他以自己的力量為神明。 \end{tabularx} \\ \\ \relax
1:12 & \begin{tabularx}{0.7\textwidth}{X} 耶和華－我的神，我的聖者啊，你不是從亙古就有嗎？我們必不致死。耶和華啊，你派他為要行審判；磐石啊，你立他為要懲治人。 \end{tabularx} \\ \\ \relax
1:13 & \begin{tabularx}{0.7\textwidth}{X} 你的眼目清潔，不看邪惡，也不看奸惡，為何你卻看著人行詭詐呢？惡人吞滅比自己公義的人，為何你保持沉默呢？ \end{tabularx} \\ \\ \relax
1:14 & \begin{tabularx}{0.7\textwidth}{X} 你為何使人如海中的魚，又如無人管轄的爬行動物呢？ \end{tabularx} \\ \\ \relax
1:15 & \begin{tabularx}{0.7\textwidth}{X} 他用鉤子把他們全拉上來，用羅網捕獲他們，拉漁網聚集他們。因此，他歡喜快樂， \end{tabularx} \\ \\ \relax
1:16 & \begin{tabularx}{0.7\textwidth}{X} 向羅網獻祭，向漁網燒香；因為他藉此得豐盛的收穫與肥美的食物。 \end{tabularx} \\ \\ \relax
1:17 & \begin{tabularx}{0.7\textwidth}{X} 但他豈可因此屢屢倒空羅網，時常殺戮列國的人，毫不顧惜呢？ \end{tabularx} \\ \\ \relax
2:1 & \begin{tabularx}{0.7\textwidth}{X} 我要站在我的瞭望臺，立在城樓上觀看，看耶和華要對我說甚麼，我可用甚麼話向他訴冤。 \end{tabularx} \\ \\ \relax
2:2 & \begin{tabularx}{0.7\textwidth}{X} 耶和華回答我，說：將這默示清楚地寫在看板上，使人容易朗讀。 \end{tabularx} \\ \\ \relax
2:3 & \begin{tabularx}{0.7\textwidth}{X} 因為這默示有一定的日期，論及終局，絕不落空。它雖然耽延，你要等候；因為它必臨到，不再遲延。 \end{tabularx} \\ \\ \relax
2:4 & \begin{tabularx}{0.7\textwidth}{X} 看哪，惡人自高自大，心不正直；惟義人必因他的信得生。 \end{tabularx} \\ \\ \relax
2:5 & \begin{tabularx}{0.7\textwidth}{X} 他因酒詭詐、狂傲、不安於位；他張開喉嚨，好像陰間，如死亡不能知足，他聚集萬國，招聚萬民全歸自己。 \end{tabularx} \\ \\ \relax
2:6 & \begin{tabularx}{0.7\textwidth}{X} 這些人豈不都要提起詩歌和俗語，嘲諷他說：禍哉！你增添不屬自己的財物，靠押金發財，要到幾時呢？ \end{tabularx} \\ \\ \relax
2:7 & \begin{tabularx}{0.7\textwidth}{X} 咬傷你的豈不忽然興起，擾害你的豈不突然崛起，你就成為他們的擄物嗎？ \end{tabularx} \\ \\ \relax
2:8 & \begin{tabularx}{0.7\textwidth}{X} 因你搶奪許多國家，流人的血，向土地、城鎮和全城的居民施行殘暴，各國殘存之民都必搶奪你。 \end{tabularx} \\ \\ \relax
2:9 & \begin{tabularx}{0.7\textwidth}{X} 禍哉！那為本家積蓄不義之財、在高處搭窩、指望得免災禍的人！ \end{tabularx} \\ \\ \relax
2:10 & \begin{tabularx}{0.7\textwidth}{X} 你圖謀剪除許多民族，犯了罪，使自己的家蒙羞，自害己命。 \end{tabularx} \\ \\ \relax
2:11 & \begin{tabularx}{0.7\textwidth}{X} 牆裡的石頭要呼叫，屋內的棟梁必應聲。 \end{tabularx} \\ \\ \relax
2:12 & \begin{tabularx}{0.7\textwidth}{X} 禍哉！那以鮮血建城、以罪孽造鎮的人！ \end{tabularx} \\ \\ \relax
2:13 & \begin{tabularx}{0.7\textwidth}{X} 看哪，這不都是出於萬軍之耶和華嗎？萬民勞碌得來的被火焚燒，萬族辛苦建造的，歸於虛空。 \end{tabularx} \\ \\ \relax
2:14 & \begin{tabularx}{0.7\textwidth}{X} 全地都必認識耶和華的榮耀，好像水充滿海洋一般。 \end{tabularx} \\ \\ \relax
2:15 & \begin{tabularx}{0.7\textwidth}{X} 禍哉！那給鄰舍酒喝，加上毒物，使人喝醉，為要看見他們下體的人！ \end{tabularx} \\ \\ \relax
2:16 & \begin{tabularx}{0.7\textwidth}{X} 你滿受羞辱，不得榮耀；你也喝吧，顯明你是未受割禮的！耶和華右手的杯必傳到你那裡，你的榮耀就變為羞辱。 \end{tabularx} \\ \\ \relax
2:17 & \begin{tabularx}{0.7\textwidth}{X} 黎巴嫩所受的殘暴必淹沒你，野獸所遭遇的毀滅使你驚嚇；因你流人的血，向土地、城鎮和全城的居民施行殘暴。 \end{tabularx} \\ \\ \relax
2:18 & \begin{tabularx}{0.7\textwidth}{X} 偶像有甚麼益處呢？製造者雕刻它，鑄成偶像，作虛假的教師；製造者倚靠的是自己所做的啞巴偶像。 \end{tabularx} \\ \\ \relax
2:19 & \begin{tabularx}{0.7\textwidth}{X} 禍哉！那對木頭說「醒起」，對啞巴石頭說「起來」的人！偶像豈能教導人呢？看哪，它以金銀包裹，其中並無氣息。 \end{tabularx} \\ \\ \relax
2:20 & \begin{tabularx}{0.7\textwidth}{X} 惟耶和華在他的聖殿中，全地都當在他面前肅靜。 \end{tabularx} \\ \\ \relax
3:1 & \begin{tabularx}{0.7\textwidth}{X} 哈巴谷先知的禱告，調用流離歌。 \end{tabularx} \\ \\ \relax
3:2 & \begin{tabularx}{0.7\textwidth}{X} 耶和華啊，我聽見你的名聲；耶和華啊，我懼怕你的作為。求你在這些年間復興你的作為，在這些年間將它顯明出來；在發怒的時候以憐憫為念。 \end{tabularx} \\ \\ \relax
3:3 & \begin{tabularx}{0.7\textwidth}{X} 神從提幔而來，聖者從巴蘭山臨到；（細拉）他的榮光遮蔽諸天，頌讚遍滿全地。 \end{tabularx} \\ \\ \relax
3:4 & \begin{tabularx}{0.7\textwidth}{X} 他的輝煌如同日光，從他手裡發出光芒，那裡隱藏他的能力。 \end{tabularx} \\ \\ \relax
3:5 & \begin{tabularx}{0.7\textwidth}{X} 在他前面有瘟疫流行，在他腳下有熱症發出。 \end{tabularx} \\ \\ \relax
3:6 & \begin{tabularx}{0.7\textwidth}{X} 他站立，震動大地，他觀看，震動列國。永久的山崩裂，長存的嶺塌陷，他的作為與古時一樣。 \end{tabularx} \\ \\ \relax
3:7 & \begin{tabularx}{0.7\textwidth}{X} 我見古珊的帳棚遭難，米甸地的幔子動搖。 \end{tabularx} \\ \\ \relax
3:8 & \begin{tabularx}{0.7\textwidth}{X} 耶和華啊，你豈是向江河發怒，向江河生氣，向海洋發烈怒嗎？你騎在馬上，坐在得勝的戰車上， \end{tabularx} \\ \\ \relax
3:9 & \begin{tabularx}{0.7\textwidth}{X} 你的弓全然顯露，箭是發誓的言語；（細拉）你以江河分開大地。 \end{tabularx} \\ \\ \relax
3:10 & \begin{tabularx}{0.7\textwidth}{X} 山嶺見你，無不戰抖；大水氾濫而過，深淵發聲，洶湧翻騰。 \end{tabularx} \\ \\ \relax
3:11 & \begin{tabularx}{0.7\textwidth}{X} 因你的箭射出光芒，你的槍閃出光耀，日月都停在原處。 \end{tabularx} \\ \\ \relax
3:12 & \begin{tabularx}{0.7\textwidth}{X} 你發怒遍行大地，以怒氣責打列國，如打穀一般。 \end{tabularx} \\ \\ \relax
3:13 & \begin{tabularx}{0.7\textwidth}{X} 你出來拯救你的百姓，拯救你的受膏者；你打破惡人之家的頭，暴露其根基，直到頸項。 （細拉） \end{tabularx} \\ \\ \relax
3:14 & \begin{tabularx}{0.7\textwidth}{X} 你以其戈矛刺透他戰士的頭；他們如旋風將我颳散，他們喜愛暗中吞吃困苦的人。 \end{tabularx} \\ \\ \relax
3:15 & \begin{tabularx}{0.7\textwidth}{X} 你騎馬踐踏海，踐踏洶湧的大水。 \end{tabularx} \\ \\ \relax
3:16 & \begin{tabularx}{0.7\textwidth}{X} 我聽見這聲音，身體戰兢，嘴唇發顫，骨中朽爛，在所立之處戰兢；但我安靜等候災難之日臨到那上來侵犯我們的民。 \end{tabularx} \\ \\ \relax
3:17 & \begin{tabularx}{0.7\textwidth}{X} 雖然無花果樹不發旺，葡萄樹不結果，橄欖樹也不收成，田地不出糧食，圈中絕了羊，棚內也沒有牛； \end{tabularx} \\ \\ \relax
3:18 & \begin{tabularx}{0.7\textwidth}{X} 然而，我要因耶和華歡欣，因救我的神喜樂。 \end{tabularx} \\ \\ \relax
3:19 & \begin{tabularx}{0.7\textwidth}{X} 主耶和華是我的力量，他使我的腳快如母鹿，又使我穩行在高處。這歌交給聖詠團長，用絲弦的樂器。 \end{tabularx} \\ \\
[1ex]
\hline
\hline
\end{longtable}
講題:絕望中的希望    Hope in Despair

經課:哈巴谷書1:1-3:19

講員:鄧瑞強博士

場合:香港中文大學崇基學院禮拜堂主日崇拜

日期:2025年8月17日

各位弟兄姊妹,早安.願神的恩典常與你們同在.

有沒有這種經驗:你求神某件事,當神應允後,你又後悔?有人求神賜他一位太太,神應允了,結婚兩個禮拜後,便後悔了.又或者,你牙痛,求神讓你永遠不再有牙痛.對神而言,這簡單不過,轉眼間,你的所有牙齒都消失了.即時,你便會後悔.這些故事教訓我們,不要亂祈禱.

先知哈巴谷,求神懲治人間的不義,神應允,但他很快便後悔了.

今日,我們一起讀哈巴谷書.

哈巴谷書1:1-4

哈1:1　先知哈巴谷所得的默示.

(哈巴谷的申訴)

哈1:2　他說:耶和華啊!我呼求你,你不應允,要到幾時呢?我因強暴哀求你,你還不拯救.【哀歌的語調】

哈1:3   你為何使我看見罪孽?你為何看著奸惡而不理呢?毀滅和強暴在我面前,又起了爭端和相鬥的事.

【先知看見社會的崩壞】

哈1:4　因此律法放鬆,公理也不顯明；惡人圍困義人,所以公理顯然顛倒.

第2節原文開頭的字,其實是「要到幾時」.先知叫救命,大聲呼喊「暴力」.他問:要叫多久,神才解救呢?我們看到,先知是一個渴望正義的人.

第3節原文開頭的字,是「為何?」「為什麼?」先知問:神為什麼讓他看見這麼多的罪孽,奸惡,毀滅,強暴,爭端,相鬥的事?這些事,光天化日之下,人人看到,神也必看到.但神,為何沒有作為?

律法鬆弛,公理不張,惡人圍著義人,全無道理可言.這是一個什麼樣的世界?在猶大國崩潰前夕,猶大社會就是這樣子.

神,你到底理不理這個世界?很多人常常問這類問題.

哈巴谷書1:5-6

(耶和華的回答)

哈1:5　耶和華說:你們要向列國中觀看,大大驚奇；因為在你們的時候,我行一件事,雖有人告訴你們,你們總是不信.【講你都唔信】

哈1:6　我必興起迦勒底人(即巴比倫人),就是那殘忍暴躁之民,通行遍地,佔據那不屬自己的住處.【神將興起巴比倫人懲罰猶大】

第5節原文開頭的字,是「觀」.之前,哈巴谷問個天,天有無眼?現在,神叫哈巴谷看.看什麼?看看:

一,神真的存在,有所作為,祂將懲治不義,你們用心看.

二,萬國在神手中,祂可隨意調動任何民族,包括猶大的敵人巴比倫人(巴比倫即現時的伊拉克),去成就祂的旨意.

三,如此,顯明耶和華不單是猶大國的神明,更是萬國的神明,是天地間的主宰,祂引導著歷史的變局.國際大變局未發生前,猶大人的信仰相對單純,他們相信耶和華保護著他們的民族與領土.耶和華和他們打交道,與其他民族關係不大.但在國際大變局中,猶大人的視野也開闊了,他們看到:耶和華原來不單是猶大的神,事實上也是巴比倫的神.祂是萬國之神.有這種信仰,是很重要的.當後來他們被擄去巴比倫時,他們才不會覺得那是沒有耶和華之異域.耶和華是萬國之主,處處都在.神不單是一時,一地,一形勢的主,神是萬有的主.祂賜福給猶大人,也降禍給犯罪的猶大人.

四,神興起「巴比倫人」,就是那些殘忍暴躁之民,去懲罰猶大國的不義.神不是沒有行動的,不過,祂有祂自己的行動方式.這些方式,在人看來,可能是難以理解的.

哈巴谷呼求神處理一下猶大社會的邪惡,怎知神使用殘暴不仁的大惡人巴比倫來滅猶大國.神聽了哈巴谷的祈禱,很可能,哈巴谷即時便後悔了.我叫你打一下那個頑皮仔,無叫你打死他.哈巴谷想都想不到,神一動手,猶大國便亡國了.他可能後悔,求神主持正義.所以,不要隨便求,有時,要得到你想要的東西,先要經過某些意想不到的過程.

哈巴谷書1:12-13

(哈巴谷的申訴)

哈1:12　耶和華—我的上帝,我的聖者啊,你不是從亙古而有嗎?我們必不致死.耶和華啊,你派定他為要刑罰人；磐石啊,你設立他為要懲治人.

哈1:13　你眼目清潔,不看邪僻,不看奸惡；行詭詐的,你為何看著不理呢?惡人吞滅比自己公義的,你為何靜默不語呢?

哈巴谷書1:12的翻譯及解釋,在學術上爭議非常大.我作出以下的理解:由於1:12的開頭及1:13整句是問句,1:12的下半句最好理解為問句.於是,1:12的意思是:「耶和華—我的上帝,我的聖者啊,你不是從亙古而有嗎?在你的旨意下,我們今次要死(滅亡)了嗎?耶和華啊,你真的派定他(巴比倫人)來刑罰我們嗎?磐石啊,你真的設立他(巴比倫人)來懲治我們嗎?」言下之意,是問神祂是否來真的,以巴比倫人懲罰猶大,這是神不再改變的旨意嗎?

然後,哈巴谷又投訴,他訴之於神的聖潔,神怎能用比猶大更不如的巴比倫去審判猶大國呢?打個比喻:神用公正的法官去懲治犯罪的猶大,是合理的；但神用像黑社會般的巴比倫人來懲治猶大,就不合理了.這違反神的聖潔.猶大人的生活又不義又暴力,對,但巴比倫入侵時,聖靈塗炭,其不義與暴力,何止百倍.這違反神的公義.面對猶大的罪惡,先知渴望「變」.但當這「變」看來變得更罪惡時,先知很疑惑.

我們像哈巴谷,在歷史的迷宮中,有很多不明白的地方.我們有時不明白神為何看來毫無反應,有時不明白歷史的變局為何是這樣,有時更不明白歷史好像變得越來越差.神還在掌權嗎?神是正義的嗎?歷史還有全盤意義的嗎?

哈巴谷書2:4

(耶和華的回答)

哈2:4　迦勒底人自高自大,心不正直；惟義人因信得生.

因保羅神學的緣故,「義人必因信得生」這金句,變成了基督教的核心信仰.經過簡化後,變成了「信耶穌,上天堂」.

但在哈巴谷書這裡,「義人因信得生」是神用來回答哈巴谷的歷史正義問題的.神用巴比倫去懲罰猶大,不表示巴比倫是正義的.這裡,對巴比倫的描繪是「自高自大,心不正直」.由於巴比倫在成為帝國的過程中,「自高自大,心不正直」,故此,神又會興起波斯(即現今的伊朗)來取代它.當波斯興起時,被擄去巴比倫的猶大人,便能回歸,重建家園,重整信仰,如此,墮落的民族才能重生.歷史的過程是很混亂的,誰義誰不義,很難說得明白.歷史是否向著終極美善的方向發展,一時間也看不清楚.就在這歷史的混亂局面裡,神說:「惟義人因信得生」.在歷史的黑洞裡,這是一條救生索,抓住它,沿著它走,就能走出迷霧.

「義人因信得生」講兩方面的事情.

一,這「信」,首要講的,不是人的信,而是神的信實(faithfulness).所謂「信實」,就是「一諾千金,講了便是」的永不失信.在歷史過程中,混亂非常,唯一恆常的,是神對祂的子民的承諾.神愛我們,甚至犧牲祂自己,也不會放棄我們.要死的話,祂自己去死,只望救我們得到生命.在迷茫中,我們能靠的,就是神的承諾,神的信實.

二,這「信」,當然也牽涉我們對神的信任.神有承諾,我們信任祂,這就構成「約」的關連.這當然不是現代社會的「契約」(contract),而是生死相繫的「盟約」(covenant).契約是很易解除的,盟約卻是以性命相許的.神將生命送贈給我們,我們將生命交付給祂,惟義人因這互信的盟約而得生,這就是我們走出歷史迷霧的救生索.

歷史過程中,有很多無法理解的地方,哈巴谷發出疑問,作出申訴.神的回答是:「惟義人因信得生.」神希望我們相信祂,相信祂對我們的愛,相信祂主宰著歷史的前進.在這信任中,我們要活出生命應有的「義」.倚靠神的應許,自己也努力活出美善,這才能令我們走出歷史的迷茫.

哈巴谷其實與神很親近,能與神對話,一問一答.這一問一答,在先知中,也是少有的.最後,他大徹大悟,悟到穿越歷史迷霧的信念.

哈巴谷書3:16-19

哈3:16　我聽見耶和華的聲音,身體戰兢,嘴唇發顫,骨中朽爛；我在所立之處戰兢(腳在戰抖).【聽到神定意審判猶大,想到家破人亡,先知整個人由頭到腳在顫動.】我只可安靜等候災難之日臨到,犯境之民上來.【大能之日既定,只能默然等這日的到來.】

哈3:17　雖然無花果樹不發旺,葡萄樹不結果,橄欖樹也不效力,田地不出糧食,圈中絕了羊,棚內也沒有牛；

哈3:18　然而,我要因耶和華歡欣,因救我的上帝喜樂.

哈3:19　主耶和華是我的力量；他使我的腳快如母鹿的蹄,又使我穩行在高處.這歌交與伶長,用絲弦的樂器.

先知不是阿Ｑ,不是說死到臨頭也算是勝利,他知道猶大民族要面臨的災劫時,他苦不堪言,戰慄不已.但他仍能講出一段充滿希望的話,在我看來,這段話展現出最高級的信仰人格,憑著對神的信任,在全無希望的絕境中,言說希望.希望不建基於環境的好壞,只建基於神的應許與恩典.哈巴谷,竟然在深不見底的黑暗中,相信光明的存在.這真是一個不簡單的人.

我好喜歡哈巴谷這段話,這段話充分表達什麼叫做信任神,和因著這信任而生起的盼望.我再讀這段話.哈巴谷說:

哈3:17　雖然無花果樹不發旺,葡萄樹不結果,橄欖樹也不效力,田地不出糧食,圈中絕了羊,棚內也沒有牛；

哈3:18　然而,我要因耶和華歡欣,因救我的上帝喜樂.

什麼也沒有了,真的一無所有了.用我們的講法是:錢沒有了,名譽沒有了,地位沒有了,前途沒有了,夢想沒有了,任何叫人生存下去的意義都沒有了.就在這一無所有的光景裡,哈巴谷說:

這現實的絕境不能成為我的牢籠；

這絕望的現實不能蒙蔽我的眼睛叫我看不到仁愛的上帝；

苦難並不是上帝,它不能在我身上有絕對的話事權；

無邊的苦楚並不能取消喜樂的理由；

不義者的勝利並不能決定我們的命運；

在看似封死了的前路仍有測不透的將來.

在這無盡的黑夜中,先知以信心,信任神救贖的承諾.他看到一個全新的世界,在看來一無所有的絕望世界中,看到神應許的希望,充滿人間.神的光明比一切黑暗的力量大,故此,我們仍有喜樂的理由.在絕望中,先知因信仰神而拒絕絕望.當人人都因困境而質問上帝時,先知因信任上帝而質問這困境.哈巴谷說:雖然現實的絕望是真的,然而,我仍因耶和華歡欣,因救我的上帝喜樂.

哈巴谷最後的話:

哈3:19　主耶和華是我的力量；他使我的腳快如母鹿的蹄,又使我穩行在高處.這歌交與伶長,用絲弦的樂器.【好明顯,哈巴谷的信息,已成為猶大信仰禮儀的一部分】

「他使我的腳快如母鹿的蹄,又使我穩行在高處」,這經文引自撒母耳記下22:34及詩篇18:33,看來是常用的意象,代表得到力量和救贖.

好在世上有哈巴谷這種人,懷著對神的信心,拒絕絕望.他激勵了後來亡國被擄的人,在歷史的黑暗長夜裡,仍然堅持那暫時未能看見的希望.

我們需要更多人,在黑夜中,相信光明的來臨.

深願各位弟兄姊妹,在絕望裡,相信上帝,保持希望.

願神祝福你們.

願榮耀歸予聖父,聖子,聖靈.阿們.



\newpage

\section{直起腰來}
\label{sec:761}
\textbf{葉菁華教授 - 直起腰來}
\newline
\newline
\hyperref[sec:760]{< < < PREV SERMON < < <}
~
\hyperref[sec:toc]{[返目錄]}
~
\hyperref[sec:762]{> > > NEXT SERMON > > >}


路加福音 13:10-17
\newline
\begin{longtable}{cl}
\hline
\hline
章節 & 經文 (和合本修訂版)\\
\hline
13:10 & \begin{tabularx}{0.7\textwidth}{X} 安息日，耶穌在一個會堂裡教導人。 \end{tabularx} \\ \\ \relax
13:11 & \begin{tabularx}{0.7\textwidth}{X} 有一個女人被靈附身，病了十八年，腰彎得一點都直不起來。 \end{tabularx} \\ \\ \relax
13:12 & \begin{tabularx}{0.7\textwidth}{X} 耶穌看見，就叫她過來，對她說：「婦人，你的病好了！」 \end{tabularx} \\ \\ \relax
13:13 & \begin{tabularx}{0.7\textwidth}{X} 於是用雙手按著她，她立刻直起腰來，就歸榮耀給神。 \end{tabularx} \\ \\ \relax
13:14 & \begin{tabularx}{0.7\textwidth}{X} 會堂的主管因為耶穌在安息日治病，就很生氣，對眾人說：「有六天應當做工，那六天之內可以來求醫，在安息日卻不可。」 \end{tabularx} \\ \\ \relax
13:15 & \begin{tabularx}{0.7\textwidth}{X} 主回答他：「假冒為善的人哪，難道你們各人在安息日不解開槽上的牛和驢，牽去喝水嗎？ \end{tabularx} \\ \\ \relax
13:16 & \begin{tabularx}{0.7\textwidth}{X} 何況她本是亞伯拉罕的後裔，被撒但捆綁了十八年，不該在安息日這天解開她的綁嗎？」 \end{tabularx} \\ \\ \relax
13:17 & \begin{tabularx}{0.7\textwidth}{X} 耶穌說這些話，他的敵人都慚愧了；所有的人因他所做一切榮耀的事都很歡喜。 \end{tabularx} \\ \\
[1ex]
\hline
\hline
\end{longtable}
講題:直起腰來 Stood Up Straight

經課:路加福音13:10-17

講員:葉菁華教授

場合:香港中文大學崇基學院禮拜堂主日崇拜

日期:2025年8月24日

第一頁:經文場景——十八年無法直起腰來

有一個女人被鬼附著,病了十八年,腰彎得一點直不起來.

在香港街頭,有時見到老婦人彎曲著腰,拖著一竹籃的廢紙皮……

十八年——足夠讓一個嬰孩長大進入中大.她每一步行走,看見的不是前方,而是腳下.這是疾病帶來的綑綁,她完全無可奈何.

身體的痛楚:脊椎的壓力,肌肉的僵硬,每一步的艱難

社交的隔離:在注重eye contact的文化中永遠處於劣勢

敬拜的諷刺:在讚美上帝的地方,卻無法抬頭仰望

即使無法「直起腰來」敬拜,她仍然堅持來到會堂.

像許多坐在這禮拜堂的人:即使生活壓彎了我們,我們仍然來尋求.

會堂裡的其他人已經習慣她的存在,視而不見(粵語:「當佢無到」).即使有人願意關心,但也愛莫能助.

第二頁:我們的處境——現代生活中的彎腰者

今天,我們身邊充滿了無法「直起腰來」的人.彎腰可能是物理姿勢,更可能是生命狀態.

照顧重擔下的彎腰

Caregivers of long-term patients

工作已疲於奔命,回家更要照顧中風,失智症等長期病患以至不能自理的家人

burnout

學術壓力下的彎腰

申請research grants 多次不成功

“Up or out” (升級或出局)制度下tenure/substantiation review 的焦慮

職場焦慮下的彎腰

常要為五斗米折腰,說些不想說的話,做些不想做的事情

劇烈競爭環境下擔心飯碗不保

拼搏精神下長期超時工作；996(中國內地)

情緒病下的彎腰

香港很多人患有抑鬱症等情緒病；中大學生也不少

有depression患者形容是整個人被invisible weight(無形重量)壓著

這些現代的「彎腰者」和經文中的女人一樣:不是不想「直起腰來」,是真的不能.他們需要的不是更多的建議「要堅強」,「要有信心」,「要正能量」,而是明白,關懷,醫治

第三頁:恩典介入——耶穌使人直起腰來

耶穌的行動

看見(seeing):耶穌看見

其他人視而不見,耶穌看到她,感受到她的困境

上帝永遠是subject,祂採取主動；女人沒有祈求,其他人沒有代求

恩典,意想不到

留意:這段經文之後,耶穌以芥菜種和麵酵比喻天國

呼喚(calling):叫她過來

耶穌主動邀請她從邊緣到中心

這個行動本身已能恢復尊嚴

宣告(declaring):對她說:「女人,你脫離這病了!」

不是未來的應許,而是現在的現實.釋放已經成就,只需要接受.

觸摸(touching):於是用兩隻手按著她

在一個視她為不潔的文化中,耶穌用身體接觸,表達接納,進一步醫治

行動的結果

她立刻直起腰來,就歸榮耀與上帝.

十八年的扭曲在瞬間被醫治後,她第一個行動就是敬拜.

會堂主管抗議(13:14):管會堂的因為耶穌在安息日治病,就氣忿忿地對眾人說:「有六日應當做工；那六日之內可以來求醫,在安息日卻不可.」

留意:他對群眾說話,不敢直接對耶穌

為何如此忿怒?因為耶穌犯規,犯法；違反既定程序

他關心什麼?醫治的適當時間表,宗教團體的秩序.

耶穌揭露他偽善,雙重標準:你們會為牲畜解綁喝水,卻不為人類解綁

主說:「假冒為善的人哪,難道你們各人在安息日不解開槽上的牛,驢,牽去飲嗎?況且這女人本是亞伯拉罕的後裔,被撒旦捆綁了這十八年,不當在安息日解開她的綁嗎?」

留意:牛驢牲畜的脊骨橫向,人的脊骨直向

安息日不工作,是為釋放人(暫停勞碌),不是為規管人

安息日是為人而設,人不是為安息日而設

「亞伯拉罕的後裔」:回復和肯定她在社會和歷史中的身份地位

第四頁:我們的回應——在現實處境中直起腰來

基督今天仍在我們中間使人「直起腰來」.祂深知我們的無能為力,帶來的不是建議,而是醫治的大能.

給照顧重擔下的彎腰者

對那些工作已筋疲力盡,回家還要照顧中風,失智症家人的caregivers,耶穌的「看見」(13:12)帶來深深的安慰.祂看見你犧牲所有personal time,看見你的眼淚和疲憊.

當burnout使你無法「直起腰來」,記得:這女人在會堂18年,神沒有忘記她.你的忍耐和愛心,天父都記念,並願意你脫離綑綁.

教會群體可作甚麼回應?

給學術壓力下的彎腰者

對那些research grants屢次被拒,面對 "up or out" 焦慮的學者,上帝看見你.你的價值不取決於funding success rate或tenure review的結果.

學生的愛戴,同事的欣賞,可能經歷上帝

甚至離開,可會是另類的釋放?

給職場焦慮者

對那些為五斗米折腰,在996文化中透不過氣的人,耶穌看見你,呼喚你過來,宣告你得釋放.

耶穌說:「凡勞苦擔重擔的人可以到我這裡來,我就使你們得安息.」

明白和接納你

耶穌也挑戰扭曲的系統.牛驢在安息日都能得釋放,何況按上帝形象所造的人?

教會群體如何與職場「打工仔」同行?

叫他們「返教會」?抑或提供另類聚會?

在工作中找尋意義?抑或與他們一起唱哀歌(lamentations)

給情緒病患者

不是講正能量,而是實際的同在,接觸

他的醫治未必是立刻,但繼續醫治

教會作為醫治的群體

群體的呼召

「眾人因他所行一切榮耀的事,就都歡喜了」(13:17)

當一個人「直起腰來」,整個群體都喜樂.我們是否如此?

教會群體有否無意中為人帶來綑綁?

教會群體(基督的身體)是否成為上帝的恩典的出口,使人直起腰來?

基督在那裡?

個人的呼召

開放自己,相信神跡(不一定是超自然,而是神的足跡)

上帝採取主動,祂看見你

承認無力

接受醫治,轉化

等待天國實現



\newpage

\section{破裂水池永不盈活水滋潤生命源}
\label{sec:762}
\textbf{林豪恩先生 - 破裂水池永不盈　活水滋潤生命源}
\newline
\newline
\hyperref[sec:761]{< < < PREV SERMON < < <}
~
\hyperref[sec:toc]{[返目錄]}
~
\hyperref[sec:763]{> > > NEXT SERMON > > >}


耶利米書 2:4-13
\newline
\begin{longtable}{cl}
\hline
\hline
章節 & 經文 (和合本修訂版)\\
\hline
2:4 & \begin{tabularx}{0.7\textwidth}{X} 雅各家，以色列家的各族啊，當聽耶和華的話， \end{tabularx} \\ \\ \relax
2:5 & \begin{tabularx}{0.7\textwidth}{X} 耶和華如此說：「你們的祖先看我有甚麼錯處，竟遠離我，隨從那虛無的神明，自己成為虛無呢？ \end{tabularx} \\ \\ \relax
2:6 & \begin{tabularx}{0.7\textwidth}{X} 他們並不問：『那領我們從埃及地上來，引導我們走過曠野、沙漠有坑洞之地，走過乾旱死蔭、無人經過、無人居住之地的耶和華在哪裡呢？』 \end{tabularx} \\ \\ \relax
2:7 & \begin{tabularx}{0.7\textwidth}{X} 我領你們進入肥沃之地，使你們得吃其中的果子和美物；你們進入時，卻使我的地玷污，使我的產業成為可憎的。 \end{tabularx} \\ \\ \relax
2:8 & \begin{tabularx}{0.7\textwidth}{X} 祭司從來不問：『耶和華在哪裡呢？』傳講律法的不認識我，官長違背我，先知藉巴力說預言，隨從無益的東西。」 \end{tabularx} \\ \\ \relax
2:9 & \begin{tabularx}{0.7\textwidth}{X} 「我因此必與你們爭辯，也與你們的子孫爭辯。這是耶和華說的。 \end{tabularx} \\ \\ \relax
2:10 & \begin{tabularx}{0.7\textwidth}{X} 你們且渡到基提海島察看，派人往基達去留心查考，看可曾有過這樣的事。 \end{tabularx} \\ \\ \relax
2:11 & \begin{tabularx}{0.7\textwidth}{X} 豈有一國換了它的神明嗎？其實那不是神明！但我的百姓將他們的榮耀換了那無益的東西。 \end{tabularx} \\ \\ \relax
2:12 & \begin{tabularx}{0.7\textwidth}{X} 諸天哪，要因此震驚，顫慄，極其淒涼！這是耶和華說的。 \end{tabularx} \\ \\ \relax
2:13 & \begin{tabularx}{0.7\textwidth}{X} 因為我的百姓做了兩件惡事：離棄我這活水的泉源；又為自己鑿出水池，卻是破裂不能儲水的池子。」 \end{tabularx} \\ \\
[1ex]
\hline
\hline
\end{longtable}
講題:破裂水池永不盈 活水滋潤生命源 Cracked Cisterns and the Fountain of Living Water

經課:耶利米書2:4-13

講員:林豪恩先生

場合:香港中文大學崇基學院禮拜堂主日崇拜

日期:2025年8月31日

引言:破池與活泉

「破裂水池」和「活水泉源」這兩種狀況令你想到甚麼呢?「破裂水池」會否想到一個有裂縫而漏水的水池?當你把水倒進那個破裂水池的時候,水就從裂縫流走,你繼續倒水進去,水又繼續流走.這種狀況對現代人來說可能不會有甚麼感受,因為扭開水喉就有水源源不絕地送到.不過,如果聯想到另一種香港人稱之為「水」的東西就有感覺了.破裂水池的狀況會否令我們想到稱為「月光族」的狀況呢?雖然每個月都有收入,但卻又每個月都花清光,銀行戶口個個月清零,都不能不感到無奈.至於不停有水湧出的噴泉又會想到甚麼狀況呢?會不會想到「收成期」的成功人士呢?他們累積下來的大量資產所產生的被動收入,源源不絕地送到,正所謂躺著也賺錢.

這個「破裂水池」以及那個「活水泉源」其實是出現在我們剛才讀的舊約經課耶利米書2章13節的,在經文中又代表甚麼狀況呢?

2:13因為我的百姓做了兩件惡事,

就是離棄我這活水的泉源,

為自己鑿出池子,

是破裂不能存水的池子.

相信大家都意會得到耶利米書2章13節提及的「破裂不能存水的池子」以及「活水的泉源」既不是按字面討論食水的問題,也不是比擬經濟上收入和支出的狀況.那麼,這兩種景象象徵甚麼呢?對於耶利米先知身處的時代有甚麼信息要傳達給他們呢?讓我們回到經課的開始2章4節一齊閱讀及思想.

一．離棄活水的泉源

首先,讓我們尋找第13節「離棄我這活水的泉源」所指的是甚麼.第5節如此寫道:

5你們的列祖見我有甚麼不義,

竟遠離我,隨從虛無的神,

自己成為虛妄的呢?

耶和華把自己比擬為「活水的泉源」,「離棄活水的泉源」即是遠離耶和華.誰遠離耶和華呢?第4節告訴了我們:

4雅各家,以色列家的各族啊,你們當聽耶和華的話.

原來遠離耶和華的是「雅各家,以色列家的各族」,那就弔詭了.他們是耶和華的選民,何以被指責遠離耶和華呢?他們做了甚麼事情或者沒有做甚麼事情而令他們被視為遠離耶和華呢?經文列出的狀況也值得我們這些跟隨上主的個人檢視自己:我們會否有意無意地遠離上主呢?同時也值得跟隨上主的教會反省:教會會不會有意無意地遠離上主呢?

第一,不記念耶和華對他們做過的好事

6他們也不說:

那領我們從埃及地上來,

引導我們經過曠野,

沙漠有深坑之地,

和乾旱死蔭,無人經過,無人居住之地的耶和華在哪裡呢?

耶和華曾經拯救以色列出埃及離開為奴之地,但耶利米時代的人不再提起耶和華,也不再記念耶和華對他們的恩惠.是不是因為歷史久遠而難以傳承呢?還是刻意忘記呢?還是在新時代之下不准記念呢?這些問題都有待研究.不過,從人類經驗來說,通常被善待的感覺容易過去,但意被惡待的記憶就穩固得多.耶穌也曾經遇過類似的狀況:

17:11耶穌往耶路撒冷去,經過撒馬利亞和加利利.

17:12進入一個村子,有十個長大痲瘋的,迎面而來,遠遠地站著,

17:13高聲說:「耶穌,夫子,可憐我們吧!」

17:14耶穌看見,就對他們說:「你們去把身體給祭司察看.」他們去的時候就潔淨了.

17:15內中有一個見自己已經好了,就回來大聲歸榮耀與上帝,

17:16又俯伏在耶穌腳前感謝他；這人是撒馬利亞人.

17:17耶穌說:「潔淨了的不是十個人嗎?那九個在哪裡呢?」

以耶穌的遭遇來推論,對恩惠有回應的是十分之一,當沒事發生過的是十分之九.由此可見,那邊才是大多數?那邊才是小數呢?不再提說從前恩待自己的人,好像有點忘恩,但卻是常態.或者不同歷史狀況有不同的考慮,不同時代有不同的需要.過去的都過去了,提說還是不提說又有甚麼關係呢?

古語有云:「施恩莫望報」,曾經拯救及帶領以色列的耶和華指責以色列不提說祂,會否顯出耶和華太過計較呢?耶和華是不是好像某些要「刷存在感」的人一樣,要以別人對自己的感戴為享受呢?如果不是,為甚麼那麼介意以色列提說不提說祂呢?這些問題,我們稍後再討論.

第二,把耶和華驅逐出生活之外

8祭司都不說,耶和華在哪裡呢?

傳講律法的都不認識我.

官長違背我；

不提說過去曾經拯救以色列的耶和華是一件事,更奇怪的是祭司繼續獻祭但不理會耶和華,律法師繼續傳講律法而不認識耶和華,以耶和華之名被膏立的官長更違背耶和華,即是說,他們都把耶和華驅逐出生活之外.這種狀況可以怎樣理解呢?杜斯妥也夫斯基的著名小說《卡拉馬助夫兄弟們》其中一個叫伊凡的兄弟,在小說裡面所說的幾句話或者能夠幫助我們思考這種現象.他說:如果沒有上帝,那麼為了滿足欲望,人類就可以「為所欲為」.那是甚麼意思呢?在伊凡的觀念中,上帝代表的就是社會道德本身.如果沒有上帝,就意味著所有道德都會變成相對,人就可以自己決定自己的道德準則,不再受限制.

不知道耶利米時代的以色列人會否覺得耶和華很煩厭以至要驅逐祂呢?大家還記得在周星馳西遊記電影中孫悟空為甚麼要唐三藏消失嗎?在其中一場戲中,孫悟空伸手去搶唐三藏手中的月光寶盒,當時唐三藏說:「你愛呀?悟空,你愛實要講啦!做咩唔出聲je ?你望住我做乜je ?雖然你咁有誠意噉望住我,我係好高興,你都要講架.lor 去啦!你唔係真係愛呀嘛?你真係愛咩?」

然後孫悟空一拳打倒唐三藏,帶點瘋狂自大地說:「大家見啦,條友講ye 長篇大論九唔搭八咿咿哦哦.就成日好似有一隻烏蠅摟住,嗡……對唔住,唔系一隻,系一紮烏蠅係你耳仔邊,一路飛入你ge 腦裡邊,救命啊!於是乎我捉住只烏蠅捻爆佢ge 肚,擳佢D 腸出黎,再將D 腸箍返佢條頸跟住用力一扯!各位,成條利根出jor 黎.我再手起刀落嘩－－!成個世界清淨曬.噉你而家明白點解我要殺jor 佢.」

當年尼采說「上帝已死」的時候,世人震驚.其實,在耶利米年代的以色列,祭司不提說耶和華,傳講律法的不認識耶和華,官長違背耶和華,不就是等於宣布「上帝已死」嗎?上帝已死不是上帝存在與否的意思,而是指人們對於上帝的信念已死.上帝的死除了標誌了人不信任上帝,亦同時顯示了對衪倡導的價值和世界觀不再理會.

大家對於耶和華的子民如此對待耶和華有甚麼感受呢?耶和華是感到震驚的:

12諸天哪,要因此驚奇,

極其恐慌,甚為淒涼!

為何天地都要「驚奇,極其恐慌,甚為淒涼」呢?因為耶和華被子民背叛而需要搏同情嗎?心靈受創需要被安撫嗎?如果人這樣猜測上帝,未免以小人之心度君子之腹了.我們稍後再討論這個問題.

二．為自己鑿出破裂不能存水的池子

我們剛才讀的經文指出「活水的泉源」就是耶和華,那麼13節下半節所說的「為自己鑿出池子,是破裂不能存水的池子」又是何事呢?

4「隨從虛無的神,自己成為虛妄」

8「先知藉巴力說預言,隨從無益的神」

人類以為宣布「上帝已死」就可以不用依靠上帝,也不用信靠自己以外的任何對象,以為只是靠自己就能夠頂天立地,靠自己就能夠打造美好的新世界,靠自己就能夠過有意義又滿足的生活.原來他們仍然有所依賴,只是以偶像來取代耶和華而已.我們聽到拜偶像這回事以為是無知者的迷信行為,其實拜偶像這現象從古而今從沒有斷絕,只是形式不同而已.在耶利米年代,那稱之為巴力,在我們這個年代又稱為甚麼呢?

人為了證明自己存在過,為了肯定自己,讓自己感到實在,總想䀆自己的能力抓住一些東西.社會告訴我們,當我們擁有金錢,權力,地位,學位,成就...等等,我們就會感到自己有價值,就能夠肯定自己,就會感到滿足.果真如此嗎?我們想深一層,可能會另有發現.舉學業績為例:假設有一個學生拿著考第一的成績表回家,卻被父母責罵:「為甚麼不是每一科都100分?」在這種情況之下,這份成績表本身能夠為這位學生帶來肯定和滿足嗎?舉才能為例:被稱為「才高八斗」的曹植才華出眾,但卻不被父親曹操肯定,後來更被兄長排擠,最後鬱悶而終.歷史上還多少才子好像曹植一樣落入懷才不遇的鬱結中呢?舉金錢為例,就算滿手鈔票,就是肯定和滿足的來源嗎? 2025年8月某天,在深水Po 福榮街某私人屋苑住客會所內,一群長者把鈔票堆滿在桌子上.他們在做甚麼呢?不是非法聚賭,而是慨嘆那堆鈔票那些鈔票可能永遠變成廢紙.因為那些是日佔時期的「軍票」,若沒有人承認這些鈔票的價值,這些鈔票只會帶損失和憤怒,而不會是肯定和滿足的來源.

上述這些例子給我們甚麼啟發呢?原來人類一直追求的東西,例如金錢,權力,地位,成就...等等,不是沒有用,只是單憑這些東西本身不能肯定我們,也不能使我們滿足.人需要的肯定和滿足來自哪裡呢?就是來自和自己一樣有思想,有感情,有意志,有反應的人.回頭看上述的例子,那位拿著考第一的成績表回家的學生,如果他看到父母愉快及滿足的笑容,聽到父母的欣賞,被父母擁抱入懷,這張考第一的成績表才會讓他感到肯定和滿足.其他情況也類似,對於勤奮工作,盡忠職守的員工,最感到委屈的是付出不被肯定.因此,管理學提示上司及機構,只要慷慨地對員工表示肯定,就算工作量多,員工的感受都仍然良好.

這是我們人類的經驗,就是我們所需要的接納,肯定和滿足只能源自活生生的人,而不可能來是死物.人類雖然能夠互相給予接納和肯定,然而,我們的經驗也告訴我們,人類互相給予的接納和肯定在時空上是有限制的,在品質上良莠不齊,也不穩定的,因此我們才有那麼多忐忑不安,甚至失望和傷痛,感到焦慮和憤怒.在天地之間,在永恆之中,誰能夠無限制地給予我們這種接納和肯定呢?誰能夠令我們的生存感到踏實和滿足呢?那必定是一位活生生,並且不受能力和時空限制的上主.以這種經驗回看耶利米時代以色列的舉動,他們以人手做的偶像取代一位有位格的耶和華上帝,難怪耶和華慨嘆天地都要「驚奇,極其恐慌,甚為淒涼」.他們離棄天地的主而去仰賴手做的偶像,正正是為自己鑿出「破裂不能存水的池子」.

總結:既要喝井水,也要活泉水

耶利米書提及的井水和活泉,有沒有令我們想起耶穌也講過井水和活水的故事呢?在約翰福音第四章,耶穌和撒馬利亞的婦人對話的時候如此說:

4:13「凡喝這水的還要再渴；

4:14人若喝我所賜的水就永遠不渴.我所賜的水要在他裡頭成為泉源,直湧到永生.」

人類作為生物之一的確需要「喝這水」；然而,人類又不止是一般生物,不是生存就算,還會尋找生存的意義,存在的價值等,因此還需要「耶穌所賜的活水」.寓意地說,人類的試探和糾結往往是以「這水」去取代「活水」,不自覺地把那些在時空上及品質上有限的事物變成了「直湧到永生的泉源」.不是那些事物本身不好,而是它們都有其限制,而且會改變而不穩定,因此不能成為我們穩定的安全和滿足的來源.然而,接納,肯定,安全,穩定,滿足這一切都是我們所需要,甚至賴以為生的,因此我們極其渴望可以得到；其實我們知道這些最終都是由人給予的,而不是由物給予的,因此我們才那麼介意別人怎樣看我們.我們來到這個世界就被放置在跑道上競爭,我們知道贏才會被看重,才會被肯定,因此,個人要贏,組織要贏,國家要贏,為了贏,人可以去得幾盡呢?其實,那個獎杯很重要嗎?大家也我們真正需要的不是那個獎杯,而是那份肯定.悲劇是,我們越追求那份肯定,我們竟然越失去那份肯定.對於如我一樣的凡夫俗子,就算贏不了,都不想輸,因為輸就等於沒有人看你,等於失去肯定.在追求接納和肯定的競爭中,我們彼此也在捱.

在夜深人靜,當我們撫心自問的時候,我們感受到自己已經很疲倦,我們要驅走的焦慮和不安其實如影隨形,午夜夢迴.或許,在那個時刻,我們在某程度上感受到奧古斯丁的體會:「上帝造人時,在人的靈魂裡留下一個像上帝的形狀一樣的破洞,因此只有當人找到上帝時,才能彌補心靈的空缺.」上主當天對以色列說,不要去鑿那破裂不能存水的池子了,回來我這活水的泉源吧!今天,我們也聽到上主對我們說這話嗎?上主那份接納和肯定,能夠在我們裡頭成為泉源,直湧到永生.



\newpage

\section{作門徒的代價}
\label{sec:763}
\textbf{李均熊牧師 - 作門徒的代價}
\newline
\newline
\hyperref[sec:762]{< < < PREV SERMON < < <}
~
\hyperref[sec:toc]{[返目錄]}
~
\hyperref[sec:764]{> > > NEXT SERMON > > >}


路加福音 14:25-33
\newline
\begin{longtable}{cl}
\hline
\hline
章節 & 經文 (和合本修訂版)\\
\hline
14:25 & \begin{tabularx}{0.7\textwidth}{X} 有一大群人和耶穌同行。他轉過來對他們說： \end{tabularx} \\ \\ \relax
14:26 & \begin{tabularx}{0.7\textwidth}{X} 「無論甚麼人到我這裡來，若不愛我勝過愛自己的父母、妻子、兒女、兄弟、姊妹，甚至自己的性命，就不能作我的門徒。 \end{tabularx} \\ \\ \relax
14:27 & \begin{tabularx}{0.7\textwidth}{X} 凡不背著自己的十字架來跟從我的，也不能作我的門徒。 \end{tabularx} \\ \\ \relax
14:28 & \begin{tabularx}{0.7\textwidth}{X} 你們哪一個要蓋一座樓，不先坐下來計算費用，看能不能蓋成？ \end{tabularx} \\ \\ \relax
14:29 & \begin{tabularx}{0.7\textwidth}{X} 免得安了地基，不能蓋成，看見的人都笑話他，說： \end{tabularx} \\ \\ \relax
14:30 & \begin{tabularx}{0.7\textwidth}{X} 『這個人開了工，卻不能完工。』 \end{tabularx} \\ \\ \relax
14:31 & \begin{tabularx}{0.7\textwidth}{X} 或是一個王出去和別的王打仗，豈不先坐下來酌量，他能不能用一萬兵去抵抗那領二萬兵來攻打他的嗎？ \end{tabularx} \\ \\ \relax
14:32 & \begin{tabularx}{0.7\textwidth}{X} 若是不能，他就趁敵人還遠的時候，派使者去談和平的條件。 \end{tabularx} \\ \\ \relax
14:33 & \begin{tabularx}{0.7\textwidth}{X} 這樣，你們無論甚麼人，若不撇下一切所有的，就不能作我的門徒。」 \end{tabularx} \\ \\
[1ex]
\hline
\hline
\end{longtable}
講題:作門徒的代價    The Cost of Discipleship

經課:路加福音14:25-33

講員:李均熊牧師

場合:香港中文大學崇基學院禮拜堂主日崇拜

日期:2025年9月7日

起:恨——點解要咁講?

若大家網上聖經apps或打開一本聯合聖經公會出版的和合本聖經,你會見到「愛我勝過愛」後面有一個註腳寫著「原文是恨」.若要原原本本將聖經原文希臘文譯出來,路14:26應該是「人到我這裡來,若不恨自己的父母,妻子,兒女,弟兄,姊妹,和自己的性命,就不能作我的門徒.」返教會一段時間的信徒,很多都知道路加福音這裡是直接用上一個「恨」字,和合本的翻譯者明顯覺得這樣直接用上一個「恨」字,一定令中文讀者誤解耶穌的意思,所以用上一「若不愛我勝過愛…」這樣比較性的句子去譯一個明顯解作恨字的希臘文.翻譯者的意圖不難理解,反而耶穌為何會用上一個「恨」字呢?若參考太10:37平行的經文,那裡也是用上比較性語言:「愛父母過於愛我的,不配作我的門徒；愛兒女過於愛我的,不配作我的門徒」看來馬太福音作者與和合本譯者,都用上讀者比較容易接受的字眼,去表達耶穌原本用上的一個「恨」字.學者都認為很大可能是路加在為自己服侍的說希臘語的群體重寫耶穌言行時,還原耶穌說話所用的「恨」字.

為何耶穌要這樣說呢?

耶穌當然不是要我們在心裡又或在行動上「仇恨」我們的家人甚或自己的生命.我們今日用「恨」字,多數是對恨的對像有極度的憎惡,討厭等強烈的情緒反應.但耶穌又怎會叫我們厭惡我們的家人又或自己?這並非平常耶穌的為人和教導.所以這個恨字,像馬太福音或和合本翻譯般理解為對比或者程度上「愛耶穌過於愛家人和自己」並沒有大問題.只不過,耶穌以誇張用語的修辭手法去對他的聽眾說話,目的就是迫使聽眾正視耶穌要求的嚴重性.所以一方面我們不能照字面解去「憎恨」我們家人甚或自己,卻又不能軟化耶穌跟從他的門徒的要求.

經課一開始說到「有極多的人和耶穌同行」,至令到耶穌用上非一般的用語去表達他意思的原因.當時一般的群眾為甚麼跟著耶穌呢?當然是覺得有著數.群眾只對耶穌有要求,卻少理耶穌對他們的要求!於是耶穌用上極端的修辭,要群眾反思跟從耶穌所需要放下和拾棄的關係(恨家人和自己性命),並將會承擔的風險(十字架).

承:耶穌所用的比喻

耶穌接著用兩個比喻,要求跟從他的人務必計算一下跟從耶穌的風險.這兩個比喻對當時人應該不難理解,就算今日我們也可以領會當中的意思.唯一要解釋一下的是第一個比喻所蓋的是一座怎樣的樓.這裡原文是解作望樓,或很多英譯本譯作tower.更準確一點應作watch-tower,通常是在葡萄園中作為守護的瞭望樓.當然我們今天也有爛尾樓,但起樓打杖並非我們平常會遇到的經驗.現在是大學的add-drop period,正好讓在坐的一些同學明白甚麼是計算風險.

但無論這兩個比喻,又或耶穌的要求,都叫我們了解到,信耶穌並不只是相信一些道理,耶穌要求跟從他的人,要成為門徒,而不是行行企企,純粹滿足自己的宗教需要.甚麼是成為門徒呢?這兩個比喻的共通點,就是要為著完成一個目標而計算代價.作門徒就是跟從耶穌,向著一個目標進發,成就上主的旨意.

轉:教會的任務——塑造門徒

今日如果我好似耶穌咁,一開聲就叫大家恨自己的家人,恨自己的性命,我想大部份在座ge 可能就會起身走人.就算你坐係到聲到最後,發覺原來我叫大家要捨棄一切,才可以成為門徒,也不知大家下星期會唔會再出現.無錯,我地以為教會ge 使命係叫更加多人信主,令教會坐無虛席.(好似我地神學院開ge 科目,我地都宣傳希望多D 中大學生黎讀.但耶穌見到大班人跟位佢,佢返而用D 好難入耳ge 說話要求人成為門徒,我諗趕走人多過留低人.但今日教會開門的目的係咩?係要一大班人坐係教堂,還是要令更多人成為主ge 門徒?

太28:19「所以,你們要去,使萬民作我的門徒,奉父,子,聖靈的名給他們施洗 .20凡我所吩咐你們的,都教訓他們遵守,我就常與你們同在,直到世界的末了.」有四個動詞:去,施洗,教導,使萬民作主的門徒,當中四個動詞唯一作為命令式的,是使萬民作我的門徒,其他三個都是分詞participle.所以教會的任務是透過聖禮的施行和聖言的教導,去使人作門徒.但當中的去,在耶穌要求我們跟從他的前題下,是有方向和目標的,而不是盲目胡亂的四處遊走,甚至只閒散地「做一日基督徒(和尚)返一日崇拜(敲一日鐘)」.你有想過返教會不是只在乎相信,更在乎跟從基督,承擔使命,為著完成上主所定的目標而計算,又付出其代價?

合:今日作們徒的代價

究竟今日,甚麼是作門徒的代價?我們要計算的又會是甚麼?

我相信每個人的召命都是不一樣的,作門徒就是全心全意的為主的國度擺上.你的召命是甚麼?你又願意全然擺上一切為主所用嗎?

作門徒的代價,不是一時衝動的決定,而是每天細水長流的選擇.耶穌邀請我們,不只是來聽祂的話,更是願意在生活中與祂同行.這條路不容易,但祂從不要求我們獨自走過.

你可能仍在思考自己的召命,也可能還未準備好放下一切.但主耶穌並不催逼,只是溫柔地呼喚:「你願意跟我走嗎?」祂知道我們的掙扎,也願意陪伴我們一步一步地學習信靠.

作耶穌基督門徒的代價不再只是物質上的捨棄,而是面對一個多元,世俗甚至敵對信仰的社會中,堅持跟隨基督的挑戰.今日福音經課提醒我們,作門徒不只是頭腦上的相信,甚至不是只有情感上的認同,而是全人生命的委身.「要恨你們的父母」——主耶穌呼召我們放下的,不只是外在的擁有,更是內心的主權——讓祂成為我們生命的中心.在這個充滿不同聲音與價值觀的城市裡,忠於基督可能意味著孤單,誤解,甚至犧牲.但正如耶穌所說,若不計算代價,就無法真正跟隨祂.這不是要我們懼怕,而是邀請我們進入一個更深的愛與信任之中.因為當我們願意放下世界所給的,我們就能擁抱祂所應許的——那永不動搖的國度與祂同在的喜樂.願我們都勇敢地回應這呼召,在香港這片土地上,活出門徒的真實與榮耀.願我們都願意在主的愛中慢慢學習,慢慢放下,慢慢走上這條門徒之路.不是靠自己的力量,而是靠祂的恩典與引導.



\newpage

\section{1比99更喜樂}
\label{sec:764}
\textbf{高國雄牧師 - 1比99更喜樂?}
\newline
\newline
\hyperref[sec:763]{< < < PREV SERMON < < <}
~
\hyperref[sec:toc]{[返目錄]}
~
\hyperref[sec:765]{> > > NEXT SERMON > > >}


路加福音 15:1-10
\newline
\begin{longtable}{cl}
\hline
\hline
章節 & 經文 (和合本修訂版)\\
\hline
15:1 & \begin{tabularx}{0.7\textwidth}{X} 許多稅吏和罪人都挨近耶穌，要聽他講道。 \end{tabularx} \\ \\ \relax
15:2 & \begin{tabularx}{0.7\textwidth}{X} 法利賽人和文士私下議論說：「這個人接納罪人，又同他們吃飯。」 \end{tabularx} \\ \\ \relax
15:3 & \begin{tabularx}{0.7\textwidth}{X} 耶穌就用比喻對他們說： \end{tabularx} \\ \\ \relax
15:4 & \begin{tabularx}{0.7\textwidth}{X} 「你們中間誰有一百隻羊，失去其中的一隻，不把這九十九隻留在曠野，去找那失去的羊，直到找著呢？ \end{tabularx} \\ \\ \relax
15:5 & \begin{tabularx}{0.7\textwidth}{X} 找到了，他就歡歡喜喜地把羊扛在肩上。 \end{tabularx} \\ \\ \relax
15:6 & \begin{tabularx}{0.7\textwidth}{X} 他回到家裡，請朋友和鄰舍來，對他們說：『你們和我一同歡喜吧，我失去的羊已經找到了！』 \end{tabularx} \\ \\ \relax
15:7 & \begin{tabularx}{0.7\textwidth}{X} 我告訴你們，一個罪人悔改，在天上也要這樣為他歡喜，比為九十九個不用悔改的義人歡喜還大呢！」 \end{tabularx} \\ \\ \relax
15:8 & \begin{tabularx}{0.7\textwidth}{X} 「同樣，哪一個婦人有十塊錢，若失落一塊，不點上燈，打掃屋子，細細地找，直到找著呢？ \end{tabularx} \\ \\ \relax
15:9 & \begin{tabularx}{0.7\textwidth}{X} 找到了，她就請朋友和鄰舍來，對她們說：『你們和我一同歡喜吧，我失落的那塊錢已經找到了！』 \end{tabularx} \\ \\ \relax
15:10 & \begin{tabularx}{0.7\textwidth}{X} 我告訴你們，一個罪人悔改，神的使者也是這樣為他歡喜。」 \end{tabularx} \\ \\
[1ex]
\hline
\hline
\end{longtable}
講題:1比99更喜樂?    1 to 99 Is More Joyful?

經課:路加福音15章1至10節

講員:高國雄牧師

場合:香港中文大學崇基學院禮拜堂主日崇拜

日期:2025年9月14日

1 許多稅吏和罪人都挨近耶穌,要聽他講道.

2 法利賽人和文士私下議論說:「這個人接納罪人,又同他們吃飯.」

3 耶穌就用比喻對他們說:

4 「你們中間誰有一百隻羊,失去其中的一隻,不把這九十九隻留在曠野,去找那失去的羊,直到找著呢?

5 找到了,他就歡歡喜喜地把羊扛在肩上.

6 他回到家裡,請朋友和鄰舍來,對他們說:『你們和我一同歡喜吧,我失去的羊已經找到了!』

7 我告訴你們,一個罪人悔改,在天上也要這樣為他歡喜,比為九十九個不用悔改的義人歡喜還大呢!」

8 「同樣,哪一個婦人有十塊錢,若失落一塊,不點上燈,打掃屋子,細細地找,直到找著呢?

9 找到了,她就請朋友和鄰舍來,對她們說:『你們和我一同歡喜吧,我失落的那塊錢已經找到了!』

10 我告訴你們,一個罪人悔改,上帝的使者也是這樣為他歡喜.」

路15:1-7   迷羊的比喻

請回想一次遺失物件的經驗…

心情:自責?失落?無奈?憤怒?

尋回時又怎樣呢?

一個曾經令我很困惑的比喻

為什麼上帝會選擇「放/留下」那99隻羊呢?他們也是主的羊,也應該要被照顧…??

常認為主題就是上帝的愛,但又無法接受,因擔心自己可能會是那99隻「被留在曠野」的羊之一…

而且經文中耶穌表達出這是「理所當然」的選出,真的嗎?

不怕那99隻羊有危險?不怕他們不喜歡/妒忌?

還需要為原本應該在的那隻走失的羊歡天喜地?不是應該責罰嗎?

我知道在上帝內應該很安全,但那走失的是自己的選擇,他應該要為此負責…

這比喻用「悔改」VS「尋回」

路15:8-10 失錢的比喻

8「同樣,哪一個婦人有十塊錢,若失落一塊,不點上燈,打掃屋子,細細地找,直到找著呢?

9找到了,她就請朋友和鄰舍來,對她們說:『你們和我一同歡喜吧,我失落的那塊錢已經找到了!』

10我告訴你們,一個罪人悔改,上帝的使者也是這樣為他歡喜.」

失落的一塊錢,在家中失落的…真的值得大費周章去尋找嗎?

同樣的結語:

尋回,與朋友鄰舍分享喜悅

失去被尋回=罪人悔改

路15:11-32浪子的比喻

11 耶穌又說:「一個人有兩個兒子.

12 小兒子對父親說:『父親,請你把我應得的家業分給我.』他父親就把財產分給他們.

13 過了不多幾天,小兒子把他一切所有的都收拾起來,往遠方去了.在那裡,他任意放蕩,浪費錢財.

14 他耗盡了一切所有的,又恰逢那地方有大饑荒,就窮困起來.

15 於是他去投靠當地的一個居民,那人打發他到田裡去放豬.

16 他恨不得拿豬所吃的豆莢充飢,也沒有人給他甚麼吃的.

17 他醒悟過來,就說:『我父親有多少雇工,糧食有餘,我倒在這裡餓死嗎?

18 我要起來,到我父親那裡去,對他說:父親!我得罪了天,又得罪了你,

19 從今以後,我不配稱為你的兒子,把我當作一個雇工吧.』

20 於是他起來,往他父親那裡去.相離還遠,他父親看見,就動了慈心,跑去擁抱著他,連連親他.

21 兒子對他說:『父親!我得罪了天,又得罪了你,從今以後,我不配稱為你的兒子.』

22 父親卻吩咐僕人:『快把那上好的袍子拿出來給他穿,把戒指戴在他指頭上,把鞋穿在他腳上,

23 把那肥牛犢牽來宰了,我們來吃喝慶祝；

24 因為我這個兒子是死而復活,失而復得的.』他們就開始慶祝.

25 「那時,大兒子正在田裡.他回來,離家不遠時,聽見奏樂跳舞的聲音,

26 就叫一個僮僕來,問是甚麼事.

27 僮僕對他說:『你弟弟回來了,你父親因為他無災無病地回來,把肥牛犢宰了.』

28 大兒子就生氣,不肯進去,他父親出來勸他.

29 他對父親說:『你看,我服侍你這麼多年,從來沒有違背過你的命令,而你從來沒有給我一隻小山羊,叫我和朋友們一同快樂.

30 但你這個兒子和娼妓吃光了你的財產,他一回來,你倒為他宰了肥牛犢.』

31 父親對他說:『兒啊!你常和我同在,我所有的一切都是你的；

32 可是你這個弟弟是死而復活,失而復得的,所以我們理當歡喜慶祝.』」

耶穌終於講出最到肉的比喻

我在2025年3月30日在崇拜中曾集中分享這比喻

主角:浪子?父親?大兒子?

失羊,失錢,失人的比喻

1:99(百分之一),1:9(十分之一),1:1(二分之一)

三個故事,一個比喻

一個愛的故事

一個不合理的愛的故事

提摩太前書 1:15

「基督耶穌降世,為要拯救罪人.」這話是可信的,是十分可佩服的.

上帝不只是要拯救,更是要拯救罪人

羅馬書5:8

惟有基督在我們還作罪人的時候為我們死,上帝的愛就在此向我們顯明了.

超越一般的價值觀

在這3個比喻的故事中,我們會不停地問,值得嗎?應該嗎?

耶穌所關注的與我們不一樣

祂看到人真正的需要

路15:1-3

1許多稅吏和罪人都挨近耶穌,要聽他講道. 2法利賽人和文士私下議論說:「這個人接納罪人,又同他們吃飯.」3耶穌就用比喻對他們說:

耶穌用比喻作教導,不是為要「教訓」人,不是要為自己的行為辯護…

而單單是要讓人明白什麼是上帝的愛?

是行動的愛,是分享的愛,是超越的愛

是不論二分之一,十分之一,百分之一,甚至歷世歷代千千萬萬眾人的之一的你及我,也是主等待著悔改,等待著被尋回的一個.



\newpage

\section{若果世人不是那麼邪惡}
\label{sec:765}
\textbf{鄧瑞強牧師 - 若果世人不是那麼邪惡}
\newline
\newline
\hyperref[sec:764]{< < < PREV SERMON < < <}
~
\hyperref[sec:toc]{[返目錄]}
~
\hyperref[sec:766]{> > > NEXT SERMON > > >}


西番雅書 1:1-12
\newline
\begin{longtable}{cl}
\hline
\hline
章節 & 經文 (和合本修訂版)\\
\hline
1:1 & \begin{tabularx}{0.7\textwidth}{X} 當亞們的兒子猶大王約西亞在位的時候，耶和華的話臨到希西家的玄孫，亞瑪利雅的曾孫，基大利的孫子，古示的兒子西番雅。 \end{tabularx} \\ \\ \relax
1:2 & \begin{tabularx}{0.7\textwidth}{X} 耶和華說：「我必從地面上徹底除滅萬物。 \end{tabularx} \\ \\ \relax
1:3 & \begin{tabularx}{0.7\textwidth}{X} 我必除滅人與牲畜，除滅空中的鳥、海裡的魚、絆腳石和惡人；我必把人從地面上剪除，這是耶和華說的。 \end{tabularx} \\ \\ \relax
1:4 & \begin{tabularx}{0.7\textwidth}{X} 我必伸手攻擊猶大和耶路撒冷所有的居民；從這地方剪除剩下的巴力、事奉偶像之祭司的名字與祭司； \end{tabularx} \\ \\ \relax
1:5 & \begin{tabularx}{0.7\textwidth}{X} 還有那些在屋頂拜天上萬象的，那些敬拜耶和華指著他起誓，卻又指著米勒公起誓的； \end{tabularx} \\ \\ \relax
1:6 & \begin{tabularx}{0.7\textwidth}{X} 並那些轉去不跟從耶和華，不尋求耶和華，也不求問他的。」 \end{tabularx} \\ \\ \relax
1:7 & \begin{tabularx}{0.7\textwidth}{X} 在主耶和華面前要靜默無聲，因為耶和華的日子快到了。耶和華已經預備祭物，將召來的人分別為聖。 \end{tabularx} \\ \\ \relax
1:8 & \begin{tabularx}{0.7\textwidth}{X} 「到了獻祭給耶和華的日子，我要懲罰領袖和王子，及所有穿外邦衣服的人。 \end{tabularx} \\ \\ \relax
1:9 & \begin{tabularx}{0.7\textwidth}{X} 到那日，我必懲罰所有跳過門檻，以殘暴和詭詐塞滿主人房屋的人。 \end{tabularx} \\ \\ \relax
1:10 & \begin{tabularx}{0.7\textwidth}{X} 「當那日，從魚門必發出悲哀的聲音，從第二城區發出哀號的聲音，從山間發出破裂的大響聲。這是耶和華說的。 \end{tabularx} \\ \\ \relax
1:11 & \begin{tabularx}{0.7\textwidth}{X} 瑪革提施的居民哪，你們要哀號，因為所有的商人都滅亡了，滿載銀子的人都被剪除。 \end{tabularx} \\ \\ \relax
1:12 & \begin{tabularx}{0.7\textwidth}{X} 那時，我必用燈巡查耶路撒冷，懲罰那些沉湎在酒渣上的人；他們心裡說：『耶和華必不降福，也不降禍。』 \end{tabularx} \\ \\
[1ex]
\hline
\hline
\end{longtable}
講題:若果世人不是那麼邪惡…… If The World Wasn't So Evil...

經課:西番雅書1:1-12；3:1-4；2:10-11；3:9-20

講員:鄧瑞強牧師

場合:香港中文大學崇基學院禮拜堂主日崇拜

日期:2025年9月21日

各位弟兄姊妹,早安.願神的恩典常與你們同在.

講少許題外話,多謝上個禮拜出席我的按牧典禮的弟兄姊妹,和沒有出席卻也默默祝福我的弟兄姊妹,願神祝福你們.最感動的,是收到你們寫的「感言」.通常,最常見到大量感言的場景,便是安息禮拜裡的記念冊.我認識一個人,搞了一個「生前喪禮」,人未死,卻辦一個快樂喪禮,像派對,當中,聽聽親友們親口講的對他的評價,讚賞和感謝,他也親自答謝親友們對他的看法.這做法其實幾好,因為一個人死了,便無法在此塵世聽到和回應各人對他的評價了.我真有幸,未死前,讀到你們寫的評價.這些類似記念冊上的評價,通常都是寫得較好的,謝謝你們.

一般牧者按立為牧師,是在其事奉生涯的開頭,像高牧師,在年輕力壯時,已按立為牧師,之後,他們要在漫長的事奉歲月裡,活出其牧職.而我,退休了,事奉的歲月將盡,才被按立.如此,我的按立和高牧師的按立,意義便很不同了.高牧師的按立,是賦予他事奉的前景.我的按立,算是肯定我過去的事奉.這有點似「墓誌銘」,述說這個人過去的一生.我都很希望,我的「墓誌銘」能寫上:一個牧師.

題外話講完了,言歸正傳,我今日講西番雅書.西番雅書雖然短,但學者認為,它是之前所有先知的審判信息的總結發言.按西番雅書在聖經裡的排位,也可看到西番雅書信息的總結性.題外話:在聖經研究裡,有一門學問,是研究經卷和經文如何透過其排列次序表達出其獨有的信息的.排列次序本身就表達意義,這一點中國人應很明白,在重大場合,行進的次序,誰在誰旁邊,都是大有學問的.聖經經卷的次序,也大有學問.西番雅書的位置,放在那鴻書及哈巴谷書之後.若大家還記得我上兩次的講道,那鴻書是審判亞述的,哈巴谷書是審判巴比倫的.今日講的西番雅書,更是審判全人類的罪,這可說是神對人的罪的總結發言.西番雅書之後,便是未來幾個月我會講的哈該書,撒迦利亞書,瑪拉基書,這些書卷講的,是審判到了極點後的重建.聖經經卷的次序,都是經過精心安排的.

(1)西番雅書1:1-3:回到太初的混沌

番1:1　當猶大王阿們的兒子約西亞在位的時候,耶和華的話臨到希西家的玄孫,亞瑪利雅的曾孫,基大利的孫子,古示的兒子西番雅.

書卷設定在約西亞的時代,那是國際風雲變幻的大時代.那時,傳統的帝國亞述逐漸衰落,新興的帝國巴比倫正在冒起,而另一傳統帝國埃及卻伺機而動.弱小的猶大國,位置剛好在這些列強爭霸必經之路,風雨飄搖,岌岌可危.那時,約西亞推行宗教改革,嘗試以純淨的耶和華信仰去一洗社會的頹敗,也讓人民在信仰裡找到危難中的希望.但是,縱使約西亞王有心,看來人民依舊活得紙醉金迷,對神的審判毫無意識.

番1:2　耶和華說:我必從地上除滅萬類.

神要除滅所有東西,這是非同小可的,這是神對罪最猛烈的審判.世界將返回創造之初,回到空虛混沌.

番1:3　我必除滅人和牲畜,與空中的鳥,海裡的魚,以及絆腳石和惡人；我必將人從地上剪除.這是耶和華說的.

在創世記,世人犯罪,滿地邪惡,神用洪水毀滅大地.如今,在西番雅書,這次審判比洪水的審判更徹底.在之前的洪水審判裡,「海裡的魚」是安然無恙的.今次,連「海裡的魚」也難逃劫數.若果世人不是那麼邪惡,神就不用審判全地,徹底地淨化這個世界.先知希望人嚴肅審視生活中的罪惡,但習慣了在罪惡中生活的人,根本沒有意識到生活中有什麼罪惡.

(2)西番雅書1:1-6,9,12; 3:1-4:猶大的罪

番1:4　我必伸手攻擊猶大和耶路撒冷的一切居民【神首要審判的對象,是號稱為神的子民的猶大人】,也必從這地方剪除所剩下的巴力,並基瑪林【異教祭司或舉行錯誤崇拜的耶和華祭司】的名和祭司,

番1:5　與那些在房頂上敬拜天上萬象的,並那些敬拜耶和華指著他起誓,又指著瑪勒堪起誓的,

番1:6　與那些轉去不跟從耶和華的,和不尋求耶和華也不訪問他的.

神審判祂的子民,因為他們的信仰不純正.他們又拜耶和華,又拜其他神明.這說明他們根本不當耶和華是獨一的真神,只當耶和華是給著數的神明,與其他神明一樣.給著數的神明,當然拜得越多越好.這裡提及「巴力」,「天上萬象」,「瑪勒堪」(這可能是亞捫人的神明米勒公).這些神明,很多變種,是自古以來人人都拜的神明.「巴力」,管雨水,管豐收,管風調雨順,其實即財神.世上哪有人不敬拜財富.「天上萬象」,預示運程,預告機會.哪個炒股的人不想預知事情的走勢.當時的人不是不敬拜耶和華,而是將耶和華和這些神明等同,多拜幾個無妨.

耶和華最討厭人拜偶像,拜偶像就是將被造物視為與創造主同級,將物質提升到神的地位.拜偶像的後果很嚴重,就是會為了物質的利益而放棄價值,道德,神聖.為了利益,甚至殺人.現在的世界看來就是這樣,不再講道理,不再講價值,而為了自身的利益,聲大夾惡,講拳頭而不講道義.這樣的世界,不可能不將大量資源放在發展武器之上.

番1:9　到那日,我必懲罰一切跳過門檻,將強暴和詭詐得來之物充滿主人房屋的.

「跳過門檻」可能指某種異教做法,或者指大搖大擺闖入民居掠奪,強搶別人的東西以自肥.這是拜偶像的人的生命的潛在傾向.

番1:12　那時,我必用燈巡查耶路撒冷；我必懲罰那些如酒在渣滓上澄清的【指那些醉臥在酒的污穢中的人】；他們心裡說:耶和華必不降福,也不降禍.

這裡用的意象,是西番雅書特有的:神像打更的人,拿著燈,在這黑暗的世界,看看誰醉臥街頭,誰醉生夢死.有太多的人,是當神死的.「他們心裡說:耶和華必不降福,也不降禍.」有神無神,世界都不會變.因此,不如盡量掠奪,盡情享受,那管明天如何.

番3:1　這悖逆,污穢,欺壓的城有禍了【「有禍了」其實在開頭】!

番3:2　她不聽從命令,不領受訓誨,不倚靠耶和華,不親近她的上帝.

番3:3　她中間的首領是咆哮的獅子；她的審判官是晚上的豺狼,一點食物也不留到早晨.

番3:4　她的先知是虛浮詭詐的人；她的祭司褻瀆聖所,強解律法.

3:1的開頭,是「禍哉」,也可譯「哀哉」,可能「哀哉」更貼切.城巿的悲哀景象.3:2一連四個「不」字.我想起「十誡」:「不可有別的神」,「不可拜偶像」,「不可妄稱耶和華的名」...等等；現在這四個「不」,明明違反十誡,以偶像為神:不聽從命令,不領受訓誨,不倚靠耶和華,不親近上帝.這樣的城巿,是掠奪的城巿,3:3說:首領如獅子,法官如豺狼,將別人的生命食到骨頭也沒有剩下來.宗教的情況也差勁,3:4說:講說神話語的先知沒有立場,為討好聽眾而拍馬屁.至於服事聖所的祭司,則根本分不清神聖抑或世俗,亂解聖經,褻瀆聖所.哀哉的是,世人都喜歡他們,因為他們講的話「中聽」,有鼓勵性,不單不需要人改變錯誤的生活,更祝福這種生活.看看這個城巿,不聽從神的真理,無道可言；首領剝削人,法官無公義；傳道人是浮誇的,服事神的人毫不聖潔.這個城巿,簡直無藥可救.無怪乎3:1這樣說:哀哉,這悖逆,污穢,欺壓的城.

從2:4開始,直到3:8,經文的主題從神審判猶大轉為神審判萬國.西番雅書是對全人類的罪的一個總審判,人人都因驕傲自大,目中無神而敗亡.由於講道時間有限,我們看這段落的其中兩節經文.

(3)西番雅書2:10-11:對列國的審判

番2:10　這事【毀滅】臨到他們是因他們驕傲,自誇自大,毀謗萬軍之耶和華的百姓.

番2:11　耶和華必向他們顯可畏之威,因他必叫世上的諸神瘦弱.列國海島的居民各在自己的地方敬拜他.

在2:11有一個意象,是西番雅書特有的:耶和華令所有偶像神明「瘦弱」,令它們「無啖好食」,餓壞它們.這個意象都頗幽默.藉這意象,先知想提示百姓:在這人人自肥,損人利己的世界裡,這種人的生命根本毫不豐盛,瘦弱不堪.

先知西番雅叫所有人,所有民族審視生活裡的罪惡元素,這些罪惡元素看來正是我們生活的內在本質.表面敬拜神,但當神是死的.以物質享受為最高目的,為了得到利益,毫無道德可言.不會聽從神的說話,當然也不會尊重任何神聖價值.這種生命,導致這個世界的淪亡.神希望人悔改.

(4)西番雅書3:9-20:天國的希望

神的審判是真的,像洪水毀滅一切.但洪水之上有方舟,神對人仍有恩典.西番雅書講一種劫後餘生的前景.

番3:9　那時,我必使萬民用清潔的言語好求告我─耶和華的名,同心合意地事奉我.

番3:10　祈禱我的,就是我所分散的民〔原文是女子；下同〕,必從古實河外來,給我獻供物.

要成為劫後餘生的人,整個人性要改變.要以「清潔的言語求告神」,「同心合意地事奉神」,「給神獻上供物」——這是宗教上的清心.

番3:11　當那日,你必不因你一切得罪我的事自覺羞愧；因為那時我必從你中間除掉矜誇高傲之輩,你也不再於我的聖山狂傲.

番3:12　我卻要在你中間留下困苦貧寒的民；他們必投靠我─耶和華的名.

番3:13　以色列所剩下的人必不作罪孽,不說謊言,口中也沒有詭詐的舌頭；而且吃喝躺臥,無人驚嚇.

經文提及五個「不」,3:11說:不再羞愧,不再狂傲；3:13說:不作罪孽,不說謊言,不再有詭詐的舌頭——這是行為上的聽從神的命令.

能在大劫中留下來的人,3:12說:是「困苦貧寒的民」——這是人性上的卑微.他們一無所有,一無所靠,並不自視為神,而專心的投靠神.天國的世界,由這群人承受.

罪人能在神的審判下存活,全因神的恩典.3:14-17將神的恩典表達得淋漓盡致.這4節經文是完美的平行.由於聖經的中文翻譯有時不跟原文詞序,以下我讀的時候,會作出少許修正.

番3:14　錫安的民哪,應當歌唱!以色列啊,應當歡呼!耶路撒冷的民哪,應當滿心歡喜快樂!

對應:番3:17下　...... 他... 必因你歡欣喜樂,默然愛你,且因你喜樂而歡呼.【留意:「默然愛你」在這裡是破壞了平行的】

番3:15上　耶和華已經除去你的刑罰,趕出你的仇敵.以色列的王─耶和華在你中間......

對應:番3:17上　耶和華─你的上帝〔在你中間〕,大有能力者〔戰士〕施行拯救...

番3:15下　...... 你必不再懼怕災禍.

對應:番3:16　當那日,必有話向耶路撒冷說:不要懼怕!錫安哪；不要手軟!

這完美的平行經文,頭尾講「歡呼喜樂」；再入一層,用了兩個意象:神是王及神是大能戰士.祂拯救祂的子民免於被消滅.中心位置,講「不要懼怕」.全段經文的中心,是叫神的子民不要懼怕.

為何能不懼怕呢?

關鍵在3:17.3:17對應著3:14,講「歡呼喜樂」,但3:17有一破格,忽然插入一短句「默然愛你」,這是完全沒有平行句的.有些學者為了句子夠平行工整,不惜移走這短句.但我覺得,這短句的不平行正是關鍵所在.「默然愛你」是什麼意思?學者的意見很分歧,解釋要多少有多少.學者Ehud Ben Zvi將這短句理解為:「神因愛你,祂對你過去做過的一切默然不語.」(He will be silent in his love)如此,就不是「默然愛你」,而是「愛你而默然」.這解釋一語道破罪人能喜樂,不用懼怕的理由,因為神在祂的愛裡,不再計較我們過去的一切.我們過去的錯事,祂一句都不再講了.我們吵架時,常常翻舊帳,一發不可收拾.但神因愛我們,對我們的過去,不再評論半句話了.這解釋了為何在一完美的平行經文中,突兀的插這短句進來.經文想指出:神這種因愛而對罪人的默然不語,在宇宙間,無一物能與之平行對應.在天地間,神的恩典在萬事萬物中突出來,無一物能與其相比.就因這無可比擬的恩典,崩壞的世界才有前景可言.

西番雅書的結尾,是殘缺無望的人,在神的救贖裡重新開始.

番3:18　那些屬你,為無大會愁煩,因你擔當羞辱的,我必聚集他們.

番3:19　那時,我必罰辦一切苦待你的人,又拯救你瘸腿的,聚集你被趕出的.那些在全地受羞辱的,我必使他們得稱讚,有名聲.

番3:20　那時,我必領你們進來,聚集你們；我使你們被擄之人歸回的時候,就必使你們在地上的萬民中有名聲,得稱讚.這是耶和華說的.

審判已過,預示亡國的猶大人的回歸.這結尾開啟下一卷書哈該書的內容,一無所有的人再次踏足故土,重建家園.我下次講哈該書.

結語:

西番雅書總結之前先知的審判信息,再沒有比他說得更嚴厲了.好好審視我們的罪:欺壓比我們弱少的人,不清心地敬畏神,以自己的利益為中心(這其實是終極的拜偶像).人的罪,引致天地的毀滅,回歸太初的混沌.在這漆黑一片的世界,唯有上天的惻隱,神的顧惜,我們才有希望.好好體會上天這份恩情,好好活出這份愛.

願神賜福你們.

願榮耀歸予聖父,聖子,聖靈.



\newpage

\section{在世界崩塌之際買一塊地}
\label{sec:766}
\textbf{林豪恩先生 - 在世界崩塌之際買一塊地}
\newline
\newline
\hyperref[sec:765]{< < < PREV SERMON < < <}
~
\hyperref[sec:toc]{[返目錄]}
~
\hyperref[sec:767]{> > > NEXT SERMON > > >}


耶利米書 32:1-15
\newline
\begin{longtable}{cl}
\hline
\hline
章節 & 經文 (和合本修訂版)\\
\hline
32:1 & \begin{tabularx}{0.7\textwidth}{X} 猶大王西底家第十年，就是尼布甲尼撒十八年，耶和華的話臨到耶利米。 \end{tabularx} \\ \\ \relax
32:2 & \begin{tabularx}{0.7\textwidth}{X} 那時巴比倫王的軍隊圍困耶路撒冷，耶利米先知被囚在猶大王宮中護衛兵的院內； \end{tabularx} \\ \\ \relax
32:3 & \begin{tabularx}{0.7\textwidth}{X} 因為猶大王西底家囚禁他，說：「你為甚麼預言耶和華如此說：『看哪，我要把這城交在巴比倫王的手中，他必攻下這城。 \end{tabularx} \\ \\ \relax
32:4 & \begin{tabularx}{0.7\textwidth}{X} 猶大王西底家必不能逃脫迦勒底人的手，定要交在巴比倫王手中，他要親眼看到巴比倫王，親口跟他說話。 \end{tabularx} \\ \\ \relax
32:5 & \begin{tabularx}{0.7\textwidth}{X} 巴比倫王要將西底家帶到巴比倫；西底家必住在那裡，直到我懲罰他的時候。你們雖與迦勒底人爭戰，卻不順利。這是耶和華說的。』」 \end{tabularx} \\ \\ \relax
32:6 & \begin{tabularx}{0.7\textwidth}{X} 耶利米說：「耶和華的話臨到我，說： \end{tabularx} \\ \\ \relax
32:7 & \begin{tabularx}{0.7\textwidth}{X} 『看哪，你叔父沙龍的兒子哈拿篾必到你這裡來，說：請你買我在亞拿突的那塊地，因為你有代贖的責任。』 \end{tabularx} \\ \\ \relax
32:8 & \begin{tabularx}{0.7\textwidth}{X} 我叔父的兒子哈拿篾果然照耶和華的話來到護衛兵的院內，對我說：『請你買我在便雅憫境內、亞拿突的那塊地；因為它應該由你來承受，而且你也有代贖的責任。請你買下它吧！』我就知道這確是耶和華的話。 \end{tabularx} \\ \\ \relax
32:9 & \begin{tabularx}{0.7\textwidth}{X} 「我便向我叔父的兒子哈拿篾買了亞拿突的那塊地，秤了十七舍客勒銀子給他。 \end{tabularx} \\ \\ \relax
32:10 & \begin{tabularx}{0.7\textwidth}{X} 我在契上簽字，將契封緘，又請證人來，用天平把銀子秤給他。 \end{tabularx} \\ \\ \relax
32:11 & \begin{tabularx}{0.7\textwidth}{X} 我又將按照法定條例所立的買契，就是封緘的那一張和敞開的那一張， \end{tabularx} \\ \\ \relax
32:12 & \begin{tabularx}{0.7\textwidth}{X} 在我叔父的兒子哈拿篾和簽字作證的人，並坐在護衛兵院內所有猶大人眼前，交給瑪西雅的孫子尼利亞的兒子巴錄。 \end{tabularx} \\ \\ \relax
32:13 & \begin{tabularx}{0.7\textwidth}{X} 我在眾人眼前囑咐巴錄說： \end{tabularx} \\ \\ \relax
32:14 & \begin{tabularx}{0.7\textwidth}{X} 『萬軍之耶和華－以色列的神如此說：你拿著這文件，就是封緘的和敞開的買契，把它們放在瓦器裡，以便長久保存。 \end{tabularx} \\ \\ \relax
32:15 & \begin{tabularx}{0.7\textwidth}{X} 因為萬軍之耶和華－以色列的神如此說：將來在這地必有人再購置房屋、田地和葡萄園。』」 \end{tabularx} \\ \\
[1ex]
\hline
\hline
\end{longtable}
講題:在世界崩塌之際買一塊地    Buy a Piece of Land on the Brink of Collapse

經課:耶利米書32:1-15

講員:林豪恩先生

場合:香港中文大學崇基學院禮拜堂主日崇拜

日期:2025年9月28日

一．引言:耶利米買地

有沒有想過在香港買一塊地呢?可能有人以為買地手續繁複,想也不敢想.地政總署說在香港買地流程相對簡單.第一,他們會在憲報刊登賣地章程及詳細資料,投標方法及截標日期.第二,有意競投人士入標競投.第三,若然標書出價達至或高於底價,會按價高者得原則批出土地,一般可在數天內公布結果.

可能有人會說,買地呀,怎會有那麼錢呢?反過來問,有錢是不是就會買呢?買還是不買,考慮甚麼因素呢?資料顯示過去幾年香港賣地狀況不理想,導致政府賣地收入創10年新低.有人指責地產商「合謀不投地」,也有人解釋並非「合謀不投地」,只是純粹「計唔掂數」.如果發展商預期樓市前景亮麗,地價貴也買；如果發展商預測樓市前景不佳,地價平也不會買.其實市民買樓也類似,越貴越追,越平越不敢買.買還是不買,原來取決於如何看待前景.

剛才讀舊約經課的時候,有沒有留意到這竟然是一段關於先知買地的記載呢?

32:9「我便向我叔父的兒子哈拿篾買了亞拿突的那塊地,秤了十七舍客勒銀子給他.

32:10我在契上簽字,將契封緘,又請證人來,用天平把銀子秤給他.

32:11我又將按照法定條例所立的買契,就是封緘的那一張和敞開的那一張,

32:12在我叔父的兒子哈拿篾和簽字作證的人,並坐在護衛兵院內所有猶大人眼前,交給瑪西雅的孫子尼利亞的兒子巴錄.

先知買地都算是特別事件了,除了耶利米,你還想到還有其他先知買地的記載嗎?耶利米先知在甚麼時候買地呢?當時是怎樣的狀況呢?猶大國的地產市道前景如何呢?

二．耶利米買地的時候:世界崩塌之際

32:1猶大王西底家第十年,就是尼布甲尼撒十八年,耶和華的話臨到耶利米.

32:2那時巴比倫王的軍隊圍困耶路撒冷,

32:3因為猶大王西底家囚禁他,說:「你為甚麼預言耶和華如此說:『看哪,我要把這城交在巴比倫王的手中,他必攻下這城.

耶利米買地的時候,原來是在「巴比倫王的軍隊圍困耶路撒冷」的危機之中.可能有人會以為,耶利米是先知,一定洞悉天機,知道耶路撒冷快將戰勝巴比倫,所以先人一步下注.事實剛好相反,他得到的天機是「這城交在巴比倫王的手中,他必攻下這城」.其實當時局勢非常凶險,外有強敵圍攻,內則朝政混亂,社會動盪,民不聊生.雖然耶利米買的那塊地在耶路撒冷城外的亞拿突,但亞拿突已經落在巴比倫軍隊手中.環境如此惡劣,精明的人都不會在這種狀況之下買地.就算要買,至少都等到前景明朗才再作決定.耶利米先知在猶大國崩塌之際買一塊地,所為何事呢?有甚麼意義呢?

三．在世界崩塌之際買一塊地:承擔責任

32:6耶利米說:「耶和華的話臨到我,說:

32:7『看哪,你叔父沙龍的兒子哈拿篾必到你這裡來,說:請你買我在亞拿突的那塊地,因為你有代贖的責任.』

32:8我叔父的兒子哈拿篾果然照耶和華的話來到護衛兵的院內,對我說:『請你買我在便雅憫境內,亞拿突的那塊地；因為它應該由你來承受,而且你也有代贖的責任.請你

買下它吧!』我就知道這確是耶和華的話.

從這些對話我們可以得知,耶利米並不是主動買地,也不是投資,而是在負起某種責任.

利未記25:25你的弟兄若漸漸窮乏,賣了幾分地業,他至近的親屬就要來把弟兄所賣的贖回.

親戚被逼出售產業,先知耶利米有責任贖回這塊地.聖經沒有說明哈拿篾為甚麼要出售這塊地.他可能是只是想逃避巴比倫的侵略,把固定資產兌現為流動現金,方便行動.也有可能是他債務纏身,不得不賣地還債.無論是甚麼原因令哈拿篾賣地,當時不是一個買地的好時機.在如此惡劣的環境之下買地,簡直如「倒錢落海」.然而,耶利米先知就算在惡劣的環境之中也照做,這不是出於利益的考慮,而是出於責任.原來,有些人在世界崩塌之際仍然負起責任,而不是逃之夭夭.責任是甚麼呢?從這個例子來看,責任是和身分有關的,也是關係的體現.賣地的哈拿篾是耶利米叔父的兒子,因此耶利米在這種關係中有這分責任.

在世界崩塌之際仍然負起責任不單是耶利米先知,我們的社會中也有這樣的人.2003年春天,一種不知名的病毒來到香港,全城陷入不安和驚恐之中,停課停工,人人自危,醫治病人的醫院成為了重災區.誰走入災區,誰就可能成為下一個倒下的人.在屯門醫院有一位三個月前才考獲呼吸系統科專科資格的醫生毫不猶疑地答允走到最前線,照顧男外科 A5沙士隔離病房的病人.朋友問她會否擔心受感染,她溫柔地說:「病人需要照顧.」,就是這樣,她成為了首位殉職的醫生.《愛是不能埋─謝婉雯傳》.醫生不單單是一份職業,更是一個身分,這個身分意味著和病人的關係.謝婉雯醫生殉職雖然激勵了全港市民,但在疫情中的公立醫院仍然每天在打仗,眼前的問題不斷湧現,病房內外每天都哭哭啼啼,處身其中,在威爾斯醫院的沈祖堯醫生說他當時沈重得連話都說不出來.在惡劣的環境中,正是對人性的考驗.他當時勉勵自己說:「不能消沉,仍有很多信任自己的病人在等著!」他知道很多堅守前線的同事整個月都沒有回家,身心飽受煎熬,於是決定在巡完病房後,邀請其他醫生到他的辦公室一起祈禱,彼此關懷互勉.

人們以為責任是冰冷和外加的壓力,我們可以用另外一個角度理解.其實責任是和身分相關的,而身分是在關係之中.因此,當關係得到充分實踐的時候,可以說體現了責任.在網上經常有人問:「跟另一半沒有愛了,只剩下責任,該不該離婚?」也有人感嘆說:「唉,現在到底是愛還是責任!」還記得在婚禮中所許下的諾言嗎?「我ＸＸＸ在上帝面前願意與你ＸＸＸ結為夫婦,無論疾病康健,富厚貧窮,我願意永遠敬愛你,安慰你,尊敬你,照顧你,與你一生共守,決無異心.」你認為這是愛還是責任呢?兩個人本來是沒有這種責任的,只是因為他們相愛而進入這種關係,進入了這種關係才帶來這種責任.用另一個角度來說,當這種關係被實踐之時,就體現了愛,也體現了責任.在風和日麗之時給予愛是沒有難度的,因此也不會想到責任問題；然而,在烏雲滿天的日子說我愛你就不容易了,那時就感到責任的存在.從這些經歷再回望耶利米在猶大國崩塌之際買一塊地的行動,是愛還是責任呢?無論是愛還是責任,在亂世之中,在困難之時仍然知道自己是誰,活出自己的身分,仍然沒有放棄和自己有關係的人,及仍然實踐這種關係,這表現是令人感動的,也令人感到縱使世界崩塌,人性仍然沒有被摧毀,人間仍有溫暖.

四．在世界崩塌之際買一塊地:散播盼望

32:13我在眾人眼前囑咐巴錄說:

32:14『萬軍之耶和華－以色列的神如此說:你拿著這文件,就是封緘的和敞開的買契,把它們放在瓦器裡,以便長久保存.

32:15因為萬軍之耶和華－以色列的神如此說:將來在這地必有人再購置房屋,田地和葡萄園.』」

有人說,先知的行動具有先知性質,耶利米先知在猶大國崩塌之際買一塊地也是一項先知性的行動.這個行動要傳遞甚麼信息呢?就是要告訴世界,雖然眼前盡是殺戮和毀壞,人民流離失所,田地荒廢不產,但要想像一個與現在不一樣的將來,「將來在這地必有人再購置房屋,田地和葡萄園」.這是甚麼呢?這是盼望,是耶和華賜給在世界崩塌之際的人的盼望.

在世界崩塌之際,人活在苦難之中,很自然會感到懷疑,甚至絕望,很自然會問:「活下去還有盼望嗎?」「堅持還有意義嗎?」有盼望還是沒有盼望,有意義還是沒有意義,足以影響生存的意志及生存的質素.

約伯在極度痛苦中,他的妻子對他說:「你仍然持守你的純正嗎?你棄掉上帝,死了吧!」(2:9)

是的,世界都崩塌了,還有甚麼需要持守呢?持守又有甚麼意義呢?然而,放棄純正,於棄上帝,放棄生命,甘心嗎?約伯就是不甘心,縱使他的世界崩塌了,他仍力排眾議,堅持尋找真實.盼望是甚麼呢?關注難民,人權,和社會公義的耶穌會會士馬可.雷珀如此說:「盼望並不是樂觀.樂觀是希望事情會好轉；盼望是忍受苦難的美德,衍生力量的恩典.盼望是心靈深處開出的應許,引領我們在茫茫前路中穩步邁進.」(Hope is not optimism. Optimism expects that things will get better. Hope is a virtue grounded in suffering. It is a grace that gives strength. Hope is a promise that takes root in the heart and is a guide to an unknown future—Mark Raper S.J.)

如果世界崩塌,如何維持盼望呢?我想起2022年9月的教職員團契午餐會龔立人教授的分享,當時他剛探望烏克蘭難民回來,講了這個小故事.

他與流離失所的烏克蘭人閒談時,聽到她們分享因為逃亡要在很短時間內匆忙收拾行李的經歷.其中一位說:「我將自己一生都放在這件行李裡.」教授好奇問這位女士帶了甚麼上路,那位中年女士回答說:「我帶了香水和唇膏.」教授更好奇地問下去.那位女士解釋說:「每早,我會噴香水和擦上唇膏.香水和唇膏讓我肯定自己不是社會媒體一直描繪的披頭散髮,衣衫襤褸的難民形象.香水和唇膏提醒我依然是我,逃亡或許影響了我,但我沒有因此失去了自己.」教授領悟到香水和唇膏不只有扮靚的意思,更提醒著逃亡的人不要因環境變化而看輕自己的尊嚴.

耶利米在猶大國崩塌之際買一塊地,這個先知性的行動是鼓勵將要被巴比倫打敗的人不要放棄,要活下去,並且要持守對耶和華的信仰.不但是猶大地的餘民要活下去,就算被擄到巴比倫的人民都要活下去,「因為萬軍之耶和華－以色列的神如此說:將來在這地必有人再購置房屋,田地和葡萄園.」

五．總結

最後,讓我們回到此時此地,生活在我們社會中的人活得有盼望嗎?尤其是社會中剛剛起步的青少年感覺到盼望嗎?

兩年前日本排放核廢水,市面傳說吃了核污染的水產會影響生育.記者在街頭訪問一對年青,問:「擔唔擔心核廢水?」得到的回答是:「唔擔心,因為我們不打算有下一代」.這句說話令人想起另一句類似的說話:「我們是最後的一代」.話說疫情時期,防疫人員雷厲風行,地對民眾說:「如果你們不這樣做,我們會記錄在案,影響你們的後代.」當時有市民回應說:「我們是最後一代了,謝謝!」這句話迅速成為網路流行語,被年輕一代用來表達對現實的無奈,對社會不公的抗議,以及「不反抗即反抗」的態度.如果年青人對社會感到沒有參與的機會,認為無論如何努力都無法改變現狀,就會喪失行動的動機和信心,心理學家稱為「習得性無助」(Learned Helplessness).到底如何能夠散播盼望呢?當年耶利米先知買在猶大國崩塌之際一塊地,此時此地又可以做甚麼呢?這個問題真是不簡單.我想起聖法蘭西斯的祈禱如此說:

在仇恨之處,播下愛；

在傷害之處,播下寬恕；

在懷疑之處,播下信心；

在絕望之處,播下盼望；

在幽暗之處,播下光明；

在憂愁之處,播下歡愉.

實踐這個祈禱的內容,不同人可能有不同的做法,因為不同人有不同的身分,相對於其他人有不同的關係,因此有同的責任.有一位中大校友,看見這個城市有不少人居住環境惡劣,但又無權無勢無財去在制度上整體地改善情況,唯有懷著「幫得一個得一個」的心態從小做起.他最初參與社企「要有光」的「光屋」項目,把閒置物業活化成為基層家庭可負擔的居所,同時提供家庭支援計劃,協助家庭成員訂立儲蓄及求職目標.後來進入大機構推動商界,社福界和學界合作,為有需要的人做一點事.這位校友接受訪問時說:「我們總有一天會離開這個工作崗位,或者不再在香港,甚至不在世,我們所做的工作會不會有些東西能夠留下來,影響到其他人,讓人們的生活變得更好?」(Cubic Zine)在這個城市中,如果多些在位的人如此思考及行動,或者生活在這個地方的人就能感受到多一點盼望.



\newpage

\section{請增加我們的信心}
\label{sec:767}
\textbf{姚志豪牧師 - 請增加我們的信心}
\newline
\newline
\hyperref[sec:766]{< < < PREV SERMON < < <}
~
\hyperref[sec:toc]{[返目錄]}
~
\hyperref[sec:768]{> > > NEXT SERMON > > >}


耶利米哀歌 1:1-6
\newline
\begin{longtable}{cl}
\hline
\hline
章節 & 經文 (和合本修訂版)\\
\hline
1:1 & \begin{tabularx}{0.7\textwidth}{X} 唉！先前人口稠密的城市，現在為何獨坐！先前在列國中為大的，現在竟如寡婦！先前在各省中為王后的，現在竟成為服苦役的人！ \end{tabularx} \\ \\ \relax
1:2 & \begin{tabularx}{0.7\textwidth}{X} 她夜間痛哭，淚流滿頰，在所有親愛的人中，找不到一個安慰她的。她的朋友都以詭詐待她，成為她的仇敵。 \end{tabularx} \\ \\ \relax
1:3 & \begin{tabularx}{0.7\textwidth}{X} 猶大被擄，遭遇苦難，多服勞役。她住在列國中，得不著安息；追逼她的在狹窄之地追上她。 \end{tabularx} \\ \\ \relax
1:4 & \begin{tabularx}{0.7\textwidth}{X} 錫安的道路因無人前來過節就哀傷，她的城門荒涼，祭司嘆息，少女悲傷；她自己充滿痛苦。 \end{tabularx} \\ \\ \relax
1:5 & \begin{tabularx}{0.7\textwidth}{X} 她的敵人作主，她的仇敵亨通；耶和華因她過犯多而使她受苦，她的孩童在敵人面前去作俘虜。 \end{tabularx} \\ \\ \relax
1:6 & \begin{tabularx}{0.7\textwidth}{X} 錫安的威榮全都失去。她的領袖如找不著草場的鹿，在追趕的人面前無力行走。 \end{tabularx} \\ \\
[1ex]
\hline
\hline
\end{longtable}
講題:請增加我們的信心    Increase Our Faith!

經課:耶利米哀歌1:1-6

講員:姚志豪牧師

場合:香港中文大學崇基學院禮拜堂主日崇拜

日期:2025年10月5日

一提起信心,我地好多時都會因著自己的困境向上主祈求.例如讀書ge 就關於學業,揀科,大學,拍拖；工作ge 就關於升職,轉職,人際關係；家庭就當然係家人ge 相處,同子女ge 關係,同另一伴ge 關係,同親友的關係….住屋,甚至是否移民.因此信心係比較著重個人化,今日ge 經課特別是舊約的經課,提醒我地信心唔單止只係著眼我地自己ge 事,有冇睇見上主所創造呢個ge 世界同我地的關係,上主擺我地係邊個ge 族群,甚至乎邊一個ge 地方,將視野由個人轉行至族群.

耶利米哀歌,同埋詩篇今日所講ge 正正就係呢一個咁樣ge 處境.

哀歌:耶路撒冷被巴比倫蹂躪之後,提到佢點樣從稠密地方變成獨坐(ALONE)係城門口(現:孤零零地躺在那裡),並且由過去的大國(有名聲,受尊重,D人爭住前來)成為jor 列國中ge 寡婦(被厭棄,甚至被歧視),由受歡迎變為一個不受歡迎的地方,從前是各省之重后,現在卻淪為一個苦役/奴婢.更加提到佢ge 朋友以詭詐對代佢(現:盟友出賣她).學者就認為其實就係鄰近身邊唔同ge 國族都離棄jor 佢.

哀歌,就是哀悼耶路撒冷在主前586年陷落,焚燬,被擄人民確實係發出呼喊,他們昔日覺得上主ge 愛,現在彷彿再睇唔見上主去愛佢地,反而係度問緊上主係咪忘記jor 佢地?

巴比倫對耶路撒冷城呢種ge 對待,其實直白地表達哀歌及詩篇,詩137:1「我們在巴比倫河邊,坐在那裡,追想錫安,就哭了.」果種嘆息,悲傷,充滿痛苦,見到敵人事事亨通,而且更成為自己既主人(現:哀1:5,她的敵人成為主人；他們的權力支配她.因為她作惡多端,),猶太人成為jor 俘虜,其實都冇乜野掩飾地去表達.

詩篇更加提到,佢地掠奪佢地ge ye ,成為jor 搶奪者ge 快樂,更要迫被擄者唱歌娛樂佢地.詩篇137篇6節更加提到「我若不記得你,不看你過於我最喜樂的,寧願我的舌頭貼於上膛.」咩意思呢?現代中文譯本表達「我要是不記得你,要是不以耶路撒冷為我最大ge 喜樂,願我ge 舌頭僵硬再也不能唱歌.」

非常的諷刺,敵人要求他們以唱歌去娛樂佢地,成為他們壓迫者的快樂.相反,被擄的人不想再唱歌,以舌頭僵硬黎代表gwo 種ge 苦.大家都係「唱歌」成為jor 一個好大的對比,十分諷刺.eg.集中營的琴聲(NOWTV)

因此詩137:7「耶路撒冷攻破的日子,以東人說:「拆毀!拆毀!直拆到根基!」求你(上主)記得.當我地面對苦難或者睇見世界不同地方被壓迫ge 時候,呢個會唔會都係我地ge 疑問呢?係咪覺得上主視若無賭呢?好想上主記住?

基督徒/活在沒有受戰亂,壓迫的人,有時習慣斬釘截鐵回應被壓迫者一句:「上帝冇放棄」,一句滿有信心的說話,是否太快的回應/太快下結論?有時缺乏對受苦者處境的思考,去感受他們的感受,去代入他們身處環境的煎熬.當我們跳過他們感受的時候,彷彿佢地所經驗的痛苦也需要被美化,甚至乎否定了他們無語問蒼天/AMANI.EG,BEYOND「麻木面對蒼生只懂不說話/祂怎一去不返,他可否會感到,峰煙掩蓋天空與未尺來.」.信徒好想同佢地講「上帝冇放棄」,但佢地確實覺得被上帝放棄.我地吾知係咪好驚上主被污名?好驚我吾咁講,上主不會工作?

Eg ,波蘭奧斯威辛二戰集中營.當我參觀空空如也的營地,卻有時房間堆滿孩子的鞋；有房間推滿男士用的行李箱；更有房屋示一塊大大的布(是用女性頭髮來織的),到今天納粹所作的行為,我也問上主?及後,我買一了本書”VOICES OF MEMORY”波蘭集中營聖餐.偷偷地ge 進行聖餐,卻散發出gwo 種無窮ge 力量.

成個行程去到最後我就問gwo 個導遊:其實每一次帶隊ge 時候你ge 感受如何呢?佢講其實係好悲傷,而每次ge 情緒都會被牽動.不過佢話,每次講都係好重要.我回應:「係!」,唔單止對波蘭人重要,對全世界都好重要,好感激你地ge 堅持.

福音書耶穌講芥菜種ge 信心ge 時候,佢講到若信心大,芥菜種的樹大到連根拔起種係海裡都仲得.芥菜種呢粒微小ge 種子我諗大家都聽過,由一粒小種子去到一棵樹,能夠乘涼,甚至為人遮風擋雨,佢ge 根部抓住ge 地下,抓得幾緊,抓得幾深,需要幾多時間會有今日,十年甚至過百年先能夠做到.雖然arm arm ge 颱風話比我地聽有D 樹木,都係會lum ge (不幸地有些是人為).樹木不倒呢個天然ge 定理依然是很有說服力的.芥菜種ge 信心,係一點一滴每日ge 累積每日ge 結果.就是呢種持久延續ge 生命力.好似gwo 個波蘭導遊一樣,即使係好傷感,即使內心遇到好似風吹雨打ge 經驗,佢仍然要繼續講繼續帶領,讓我於九千里外的香港人,也有機會明白到他們的苦,以至樣呢種ge 生命力能夠承傳落去.見證住戰爭結束,亦叫我們反思,要關注上主所做的世界,讓我們在人與人之間不要發起戰爭(不只國與國,人與人也是),佔領別人的地盤(如巴比倫一樣),讓自己不要成為壓迫人的人.芥菜種樹的生長,就如歷史長河,歷史叫人看見不爭的事實:巴比倫,波斯,羅馬,二戰納粹…….最終也要倒台,環看古今,歷史就是如此,上主在看,人要配合,要記下,要提醒(如奧斯威辛),不要視若無睹.

有時看似冇可能發生ge 事,信仰往往讓我地睇見另一個ge 不可能.耶穌就係用主人同僕人ge 關係,話緊比我地聽,當以奴隸社會為主年代,邊有僕人從田間/放羊回來後,主人會叫埋黎食飯同多謝佢呀,緊係要你服侍埋我啦.但耶穌似乎想講,信仰就係呢種ge 顛覆性,因為耶穌就是那個會為僕人洗腳,會叫佢埋位食飯果一位.只係我地冇辦法估到,或只選擇世界的慣性?愛隨波逐流,還是背道而行,逆流而上,事實係逆風才能飛.

上主往往要扭轉呢個ge 世界,是要讓我地一齊參與,歷世歷代每一個人做一D,gwo 份ge 堅持,讓呢一個世界慢慢地不經不覺被扭轉jor .剛剛過身ge 動物學家Jane Goodal(珍·古德爾博士)ge 一句說話:Think globally,Act locally,令我想起:雅各書講得岩:信心沒有為也是死的.要相信你所做,正是影響緊世界.

最後好想用,立陶宛的「十字架山」(Hill of Crosses),他們就是從第1個十字架插起,又被蘇維埃政權(前蘇聯)被推倒,再起,到今天估計約有10萬過十字架,當中更有人寫上禱文,去紀念/支持毗鄰的烏克蘭朋友.「十字架山」的紀念碑文上面寫著:感謝立陶宛人民對十字架山的貢獻,這座山向歐洲乃至全世界見證了這片土地人民的信仰.不論是十字架山/奧斯威辛集中營,代表著,人的眼目對世界關懷,人的信心,不只建基所相信,也要建基一點一滴行動.

今天是普世聖餐主日,讓我們從自己目光轉向世界/普世.



\newpage

\section{為這城求平安}
\label{sec:768}
\textbf{李均熊牧師 - 為這城求平安}
\newline
\newline
\hyperref[sec:767]{< < < PREV SERMON < < <}
~
\hyperref[sec:toc]{[返目錄]}
~
\hyperref[sec:769]{> > > NEXT SERMON > > >}


耶利米書 29:1-7
\newline
\begin{longtable}{cl}
\hline
\hline
章節 & 經文 (和合本修訂版)\\
\hline
29:1 & \begin{tabularx}{0.7\textwidth}{X} 耶利米先知從耶路撒冷送信給被擄倖存的長老，以及祭司、先知，和尼布甲尼撒從耶路撒冷擄到巴比倫去的眾百姓。 \end{tabularx} \\ \\ \relax
29:2 & \begin{tabularx}{0.7\textwidth}{X} 這是在耶哥尼雅王和太后、太監，並猶大和耶路撒冷的領袖，以及工匠、鐵匠都離開耶路撒冷之後。 \end{tabularx} \\ \\ \relax
29:3 & \begin{tabularx}{0.7\textwidth}{X} 他藉沙番的兒子以利亞薩和希勒家的兒子基瑪利的手送去；他們二人是猶大王西底家差往巴比倫去見巴比倫王尼布甲尼撒的。 \end{tabularx} \\ \\ \relax
29:4 & \begin{tabularx}{0.7\textwidth}{X} 信上說：「萬軍之耶和華－以色列的神對所有被擄的，就是我使他們從耶路撒冷被擄到巴比倫去的人如此說： \end{tabularx} \\ \\ \relax
29:5 & \begin{tabularx}{0.7\textwidth}{X} 你們要建造房屋，住在其中；要開墾田園，吃園中所出產的； \end{tabularx} \\ \\ \relax
29:6 & \begin{tabularx}{0.7\textwidth}{X} 要娶妻生兒養女，為你們的兒子娶妻，使你們的女兒嫁人，生兒養女。你們要在那裡生養眾多，不可減少。 \end{tabularx} \\ \\ \relax
29:7 & \begin{tabularx}{0.7\textwidth}{X} 我使你們被擄到的那城，你們要為那城求平安，為那城向耶和華祈求，因為那城得平安，你們也隨著得平安。 \end{tabularx} \\ \\
[1ex]
\hline
\hline
\end{longtable}
講題:為這城求平安    Seek the Peace of This City

經課:耶利米書29:1-7

講員:李均熊牧師

場合:香港中文大學崇基學院禮拜堂主日崇拜

日期:2025年10月12日

起:為甚麼要為被擄到的那城求平安?

耶利米書29 章記載的,是先知耶利米寫給被擄到到巴比倫的猶太人的幾封信,其中今日經課引用的4至7節,是寫給被擄到巴比倫的猶太人領袖的一封信.第1節中所謂「從耶路撒冷擄到巴比倫去的」,在2至3節(經課中沒有的)有進一步解釋.在耶利米書,總共記載了三次的「披擄」,這裡「論到這在耶哥尼雅王(即列王記所載的約雅斤王)和太后,太監,並猶大,耶路撒冷的首領,以及工匠,鐵匠都離了耶路撒冷以後.」可算是第一次被擄.在這之後,巴比倫立了耶哥尼雅／約雅斤王的叔父為猶大國的王,繼承了大衛王朝.換句話說,第一次被擄之後,猶大國仍然存在,大衛王朝還未傾覆.到西底家登基九年之後,西底家以為自己已有足夠力量,不甘再做巴比倫的傀儡王帝,於是起兵背叛巴比倫王,巴比倫王尼布甲尼撒領軍攻打耶路撒冷,圍城足有兩年,在西底家十一年攻破耶路撒冷,以色列人經歷第二次被擄,這時猶大國才算徹底消滅,而大衛王朝只剩之前被擄於巴比倫的約雅斤這一點血脈.(耶52)

若我們再看耶29的前文後理,當時明顯有一些「先知」傳神的話,認為在約雅斤王被擄之後,上主很快就會使這一批被擄的以色列民被釋放而回歸祖國(例如耶28:1-4:「1當年,就是猶大王西底家登基第四年五月,基遍人押朔的兒子,先知哈拿尼雅,在耶和華的殿中當著祭司和眾民對我說:2「萬軍之耶和華－以色列的神如此說:我已經折斷巴比倫王的軛.3二年之內,我要將巴比倫王尼布甲尼撒從這地掠到巴比倫的器皿,就是耶和華殿中的一切器皿都帶回此地.4我又要將猶大王約雅敬的兒子耶哥尼雅和被擄到巴比倫去的一切猶大人帶回此地,因為我要折斷巴比倫王的軛.這是耶和華說的.」),故此這一班被擄的以色列人有很大可能不願在巴比倫安定下來,甚至無時無刻想著將約雅斤王復辟,趕走西底家這個巴比倫所立的傀儡王帝.亦有可能這群被擄的以色列人一直找機會,想在巴比倫製造一些動亂,裡應外合,使擄走他們的敵國巴比倫添些麻煩.是這樣的處境下,耶利米呼籲被擄的巴比倫人要蓋房種田,生兒育女；並為「被擄到的那城」求平安.並且更進一步警告大家不要聽信那些假先知所謂傳講的神的話(耶29:8-9「8萬軍之耶和華－以色列的神如此說:不要被你們中間的先知和占卜的誘惑,也不要聽信自己所做的夢；9因為他們託我的名對你們說假預言,我並沒有差遣他們.這是耶和華說的.」)

承:不受歡迎的信息

可見,耶利米的信息並不受與他同時代的人歡迎,甚至對當時一些猶太人來說,耶利米叫人在被擄之地蓋房種田,生兒育女,簡直就是通番賣國的叛國言論!故此,今日我們讀到的經文,應該可能是一,二代之後,當他們(過了十多廿年之後)知道耶利米的說話比起那些認為被擄只是短暫的「先知」更加符合現實時,他們才認真看待耶利米的說話,並追憶,記載並傳遞,而成為今日的耶利米書.所以「你們要為那城求平安,為那城禱告耶和華；因為那城得平安,你們也隨著得平安.」這一句說話,並不是太平盛世之下,基督徒為了自身平安,而為自己身處的城巿添福祈禱的指引.這只會將這句耶利米苦口婆心的說話,變成了純粹為了讓自己心安理得的宗教術語；耶利米信中的這句說話,乃是在時勢凶惡,敵人仍然得呈的時候,某求安身立命的忠告.

與其和不可抗衡的帝國力量周旋,不如在惡劣的處境下休養生息.

這句說話背後其實是對上主說話的信靠,「10「耶和華如此說:為巴比倫所定的七十年滿了以後,我要眷顧你們,向你們成就我的恩言,使你們仍回此地.11耶和華說:我知道我向你們所懷的意念是賜平安的意念,不是降災禍的意念,要叫你們末後有指望.12你們要呼求我,禱告我,我就應允你們.13你們尋求我,若專心尋求我,就必尋見.14耶和華說:我必被你們尋見,我也必使你們被擄的人歸回,將你們從各國中和我所趕你們到的各處招聚了來,又將你們帶回我使你們被擄掠離開的地方.這是耶和華說的.」(29:10-14)就是在仍然被擄的一刻,相信上主有祂的時間表,也承認祂的旨意是美善的,雖然眼前壓迫處處,困難重重,但仍然深信祂的應許必要實現,而非以自己的滿腔熱血去自把自為,以為靠自己的手段可以成就神的旨意.

轉:在應許地外仍有生機

耶利米信中並不是只叫被擄的以色列人在敵邦保待宗教虔誠,只是在被擄的城巿求個福,希望可以保個平安便算.耶利米更是要求被擄的以色列人落手落腳參與被擄城巿的建設「蓋造房屋,住在其中；栽種田園,吃其中所產的」,自己不單娶妻生兒女,更要為「你們的兒子娶妻,使你們的女兒嫁人,生兒養女.在那裡生養眾多,不致減少.」不但是為要延續民族文化的傳承,更要保持自己人口不致減少.其實不單只香港,全球有很多地方的出生率都在下降,其中一個原因是人民對自己城巿的將來沒有把握和盼望,於是就不想有下一代.故此生兒養女其實是對將來有心信的投資.所以耶利米是在叫被擄的以色列人對被擄的城巿的將來有信心.被擄的以色列人會否相信在上主應許之地外,甚至是敵對邦國的城巿,仍有從主而來的生機?留意「生養眾多」在創世記第一章是造物主對人類在世的祝福.能夠生養眾多是深信上主的祝福,並不因以色列人被擄到外地而因此減少.

合:透過我們去普及萬民的恩典

「那城得平安,你們也隨著得平安」是否也可以反過來理解,「你們的平安,正好祝福了那城都有平安」,在被擄的逆景之中,你努力生活,展現上主賜下的生機和平安,也正是因為你在這個城巿中信靠上主,你的生機,平安和活力,不就祝福了你身邊的人嗎?但不要忘記,這裡仍要留意被擄的以色列人是身處消滅自己國家的敵對的環境.在因苦逆景中可能我們會著眼上主的拯救,甚至期望上主審判加害我們的人,但不要忘記我們所相信也是創造的主,祂的恩典普及眾生；神降雨給義人也給不義的人.正如耶利米要求被擄的以色列人要「生養眾多」去體現造物主的祝福,這個造物主的祝福又是否會透過你臨到其他的被造物?你又願否透過自己的得福去祝福別國的人呢?

我不知道你是否認為自己也是在被擄的處境:可能你在一個充滿敵視的環境工作；或你長期處於種種壓力下感到生活失去自由；甚至你感到世界或社會文化越來越彎曲悖謬,但正如耶利米提醒被擄到巴比倫的以色列人一樣,也在提醒著我們:神有祂復興萬事的時間表,但此時此刻,你會否本於上主的應許和他創世的祝福,去積極的建設生活,帶著對未來的盼望,展現生命的生機,甚至使我們生命領受的祝福,為這個我們身處的都巿帶來平安?



\newpage

\section{欲窮千里目更上一層樓}
\label{sec:769}
\textbf{鄧瑞強牧師 - 欲窮千里目,更上一層樓}
\newline
\newline
\hyperref[sec:768]{< < < PREV SERMON < < <}
~
\hyperref[sec:toc]{[返目錄]}
~
\hyperref[sec:770]{> > > NEXT SERMON > > >}


哈該書 1:1-2:23
\newline
\begin{longtable}{cl}
\hline
\hline
章節 & 經文 (和合本修訂版)\\
\hline
1:1 & \begin{tabularx}{0.7\textwidth}{X} 大流士王第二年六月初一，耶和華的話藉哈該先知向撒拉鐵的兒子猶大省長所羅巴伯和約撒答的兒子約書亞大祭司傳講，說： \end{tabularx} \\ \\ \relax
1:2 & \begin{tabularx}{0.7\textwidth}{X} 「萬軍之耶和華如此說，這百姓說，建造耶和華殿的時候還沒有到。」 \end{tabularx} \\ \\ \relax
1:3 & \begin{tabularx}{0.7\textwidth}{X} 耶和華的話藉哈該先知傳講，說： \end{tabularx} \\ \\ \relax
1:4 & \begin{tabularx}{0.7\textwidth}{X} 「這殿荒涼，你們自己還住天花板的房屋嗎？ \end{tabularx} \\ \\ \relax
1:5 & \begin{tabularx}{0.7\textwidth}{X} 現在，萬軍之耶和華如此說，你們要省察自己的行為。 \end{tabularx} \\ \\ \relax
1:6 & \begin{tabularx}{0.7\textwidth}{X} 你們撒的種多，收的卻少；你們吃，卻不得飽；喝，卻不得足；穿衣服，卻不得暖；領工錢的，領了工錢卻裝入有破洞的袋中。 \end{tabularx} \\ \\ \relax
1:7 & \begin{tabularx}{0.7\textwidth}{X} 「萬軍之耶和華如此說，你們要省察自己的行為。 \end{tabularx} \\ \\ \relax
1:8 & \begin{tabularx}{0.7\textwidth}{X} 你們要上山取木料，建造這殿，我就因此喜樂，且得榮耀。這是耶和華說的。 \end{tabularx} \\ \\ \relax
1:9 & \begin{tabularx}{0.7\textwidth}{X} 你們盼望多得，看哪，所得的卻少；你們收到家中，我就吹去。這是為甚麼呢？因為我的殿荒涼，你們各人卻只為自己的房屋奔走。這是萬軍之耶和華說的。 \end{tabularx} \\ \\ \relax
1:10 & \begin{tabularx}{0.7\textwidth}{X} 所以，因你們的緣故，天不降甘露，地也不出土產。 \end{tabularx} \\ \\ \relax
1:11 & \begin{tabularx}{0.7\textwidth}{X} 我命令乾旱臨到土地、山岡、五穀、新酒、新油和地上的出產，也臨到人和牲畜，以及一切人手勞碌得來的。」 \end{tabularx} \\ \\ \relax
1:12 & \begin{tabularx}{0.7\textwidth}{X} 那時，撒拉鐵的兒子所羅巴伯、約撒答的兒子約書亞大祭司，和所有倖存的百姓都聽從耶和華－他們神的話，就是哈該先知奉耶和華他們神差遣所說的話；百姓在耶和華面前存敬畏的心。 \end{tabularx} \\ \\ \relax
1:13 & \begin{tabularx}{0.7\textwidth}{X} 耶和華的使者哈該奉耶和華差遣對百姓說：「我與你們同在。這是耶和華說的。」 \end{tabularx} \\ \\ \relax
1:14 & \begin{tabularx}{0.7\textwidth}{X} 耶和華激發撒拉鐵的兒子猶大省長所羅巴伯、約撒答的兒子約書亞大祭司，和所有倖存百姓的心，他們就來為萬軍之耶和華－他們神的殿做工。 \end{tabularx} \\ \\ \relax
1:15 & \begin{tabularx}{0.7\textwidth}{X} 這是在大流士王第二年六月二十四日。 \end{tabularx} \\ \\ \relax
2:1 & \begin{tabularx}{0.7\textwidth}{X} 七月二十一日，耶和華的話藉哈該先知傳講，說： \end{tabularx} \\ \\ \relax
2:2 & \begin{tabularx}{0.7\textwidth}{X} 「你要曉諭撒拉鐵的兒子猶大省長所羅巴伯、約撒答的兒子約書亞大祭司，和所有倖存的百姓，說： \end{tabularx} \\ \\ \relax
2:3 & \begin{tabularx}{0.7\textwidth}{X} 『你們中間存留的，有誰見過這殿從前的榮耀呢？現在你們看如何？在你們眼中豈不是如同無有嗎？ \end{tabularx} \\ \\ \relax
2:4 & \begin{tabularx}{0.7\textwidth}{X} 所羅巴伯啊，現在，你當剛強！這是耶和華說的。約撒答的兒子約書亞大祭司啊，你當剛強！這是耶和華說的。這地的百姓啊，你們都當剛強做工，因為我與你們同在。這是萬軍之耶和華說的。 \end{tabularx} \\ \\ \relax
2:5 & \begin{tabularx}{0.7\textwidth}{X} 這是照著你們出埃及時我與你們立約的話。我的靈仍要住在你們中間，你們不必懼怕。 \end{tabularx} \\ \\ \relax
2:6 & \begin{tabularx}{0.7\textwidth}{X} 萬軍之耶和華如此說：過些時候，我必再一次震動天地、滄海與乾地。 \end{tabularx} \\ \\ \relax
2:7 & \begin{tabularx}{0.7\textwidth}{X} 我必震動萬國，萬國的珍寶都必運來，我就使這殿充滿榮耀。這是萬軍之耶和華說的。 \end{tabularx} \\ \\ \relax
2:8 & \begin{tabularx}{0.7\textwidth}{X} 銀子是我的，金子也是我的。這是萬軍之耶和華說的。 \end{tabularx} \\ \\ \relax
2:9 & \begin{tabularx}{0.7\textwidth}{X} 這後來的殿的榮耀必大過先前的榮耀。這是萬軍之耶和華說的。在這地方我必賜平安。這是萬軍之耶和華說的。』」 \end{tabularx} \\ \\ \relax
2:10 & \begin{tabularx}{0.7\textwidth}{X} 大流士王第二年九月二十四日，耶和華的話臨到哈該先知，說： \end{tabularx} \\ \\ \relax
2:11 & \begin{tabularx}{0.7\textwidth}{X} 「萬軍之耶和華如此說，你要向祭司請教律法，說： \end{tabularx} \\ \\ \relax
2:12 & \begin{tabularx}{0.7\textwidth}{X} 『看哪，若有人用衣服的邊兜聖肉，這衣服的邊接觸了餅，或湯，或酒，或油，或別的食物，這些是否成為聖呢？』」祭司回答說：「不。」 \end{tabularx} \\ \\ \relax
2:13 & \begin{tabularx}{0.7\textwidth}{X} 哈該又說：「若有人因摸屍體染了不潔淨，然後接觸任何東西，這東西就變為不潔淨嗎？」祭司回答說：「必不潔淨。」 \end{tabularx} \\ \\ \relax
2:14 & \begin{tabularx}{0.7\textwidth}{X} 於是哈該說：「耶和華說，在我面前這民如此，這國也是如此；他們手裡的各樣工作都是如此；他們在那裡所獻的都不潔淨。」 \end{tabularx} \\ \\ \relax
2:15 & \begin{tabularx}{0.7\textwidth}{X} 「現在，你們心裡要想一想，從今日起，耶和華的殿還沒有一塊石頭放在石頭上的情況。 \end{tabularx} \\ \\ \relax
2:16 & \begin{tabularx}{0.7\textwidth}{X} 那時你們怎麼了？有人來到二十斗的穀堆那裡，卻只得了十斗；有人來到酒池那裡要取五十桶，卻只得了二十桶。 \end{tabularx} \\ \\ \relax
2:17 & \begin{tabularx}{0.7\textwidth}{X} 我以焚風、霉爛、冰雹攻擊你們，和你們手上的各樣工作，你們仍不歸向我。這是耶和華說的。 \end{tabularx} \\ \\ \relax
2:18 & \begin{tabularx}{0.7\textwidth}{X} 你們心裡要想一想，從今日起，就是從這九月二十四日起，從立耶和華殿根基的日子起，你們心裡想一想： \end{tabularx} \\ \\ \relax
2:19 & \begin{tabularx}{0.7\textwidth}{X} 倉裡還有穀種嗎？葡萄樹、無花果樹、石榴樹、橄欖樹雖沒有結果子，從今日起，我必賜福。」 \end{tabularx} \\ \\ \relax
2:20 & \begin{tabularx}{0.7\textwidth}{X} 這月二十四日，耶和華的話再次臨到哈該，說： \end{tabularx} \\ \\ \relax
2:21 & \begin{tabularx}{0.7\textwidth}{X} 「你要告訴猶大省長所羅巴伯說，我必震動天地， \end{tabularx} \\ \\ \relax
2:22 & \begin{tabularx}{0.7\textwidth}{X} 傾覆列國的寶座，除滅列邦列國的勢力，並傾覆戰車和坐在其上的。馬和騎兵都必跌倒，各人被弟兄的刀所殺。 \end{tabularx} \\ \\ \relax
2:23 & \begin{tabularx}{0.7\textwidth}{X} 萬軍之耶和華說：撒拉鐵的兒子我僕人所羅巴伯啊，這是耶和華說的，到那日，我必以你為印，因我揀選了你。這是萬軍之耶和華說的。」 \end{tabularx} \\ \\
[1ex]
\hline
\hline
\end{longtable}
講題:欲窮千里目,更上一層樓 See Further, You and the World will be Transformed

經課:哈該書1:1-2:23

講員:鄧瑞強牧師

場合:香港中文大學崇基學院禮拜堂主日崇拜

日期:2025年10月19日

各位弟兄姊妹,早安.願神的恩典常與你們同在.

猶大亡國後,猶大人被擄去巴比倫的艱苦日子,終於過去了.差不多50年了,半個世紀的苦難,終於過去了.波斯興起,塞魯士大帝容讓猶大人回歸故土.王族後人所羅巴伯被任命為猶大省長,他帶領族人,浩浩蕩蕩,回歸耶路撒冷.在回歸旅途中,不知道他們會夢見什麼?或許,他們會夢見家鄉的故園,會夢見兒時捉魚的汲淪溪,會夢見京城的城牆,會夢見莊嚴的聖殿.不知道他們會否在夢中笑醒?畢竟,多少年來的夢想成真了.

終於到步了.噢,耶路撒冷,夢中的故鄉.眼前的景象卻是:城牆崩塌了,故園只剩下幾塊殘垣破瓦,聖殿更只是一片蕭瑟荒涼,雜草亂生.淚,便流下來了.

接著的日子很艱苦,比在巴比倫的時候更艱苦.要很努力蓋房子,起碼讓自己有個遮風擋雨的地方.要很努力種地,讓自己在吃完從巴比倫帶來的糧食後有新糧可食.胼手胝足,只求能活下去.日子真艱難.

回來後,他們曾經打算建聖殿,但由於種種可以想像的原因,工程擱置了.時光飛逝,連波斯的帝王都由塞魯士大帝改為大流士王了.猶大人的生活,也慢慢好了一點.這時,神藉先知哈該發言,鼓勵百姓建聖殿.

我們看哈該書.

該1:1　大流士王第二年六月初一日,耶和華的話藉先知哈該向猶大省長撒拉鐵的兒子所羅巴伯和約撒答的兒子大祭司約書亞說:

該1:2　萬軍之耶和華如此說:「這百姓說,建造耶和華殿的時候尚未來到.」

大流士王第二年六月初一日,這日子離開上次他們打算建聖殿的日子,轉眼已16年了.猶大人的生活,已逐漸上軌道了.這時,神在問:「百姓為何現在仍然認為建聖殿的時候尚未來到呢?」

當生活真的很艱難時,神沒有要求,只默默祝福他們.但已休養生息16年了,仍是對神的事情不聞不問嗎?一般香港大學生24歲畢業,16年後,40歲,已在管理層了.神在問:你們仍真的只顧自己的事業嗎?天國的事業,真的不重要嗎?經過歷史的大劫,現在能在故鄉逐漸重建人生,神好想祂的子民反省,生命的優先次序是什麼?在人生中,什麼是重要,什麼是次要?神的事情,你們覺得重要嗎?你們自己的人生,比神的事情還重要嗎?

該1:3　那時耶和華的話臨到先知哈該說:

該1:4　「這殿仍然荒涼,你們自己還住天花板的房屋嗎?

該1:5　現在萬軍之耶和華如此說:你們要省察自己的行為【「行為」原意:道路,方向,生活方式】.

該1:6　你們撒的種多,收的卻少；你們吃,卻不得飽；喝,卻不得足；穿衣服,卻不得暖；得工錢的,將工錢裝在破漏的囊中.」

神發出一個天問:「聖殿仍然荒涼,但你們自己卻住得安居樂業,這對嗎?」用今日的講法是:「你們各人住豪宅大屋,但神家的事情,卻無人理會.」各人努力去建立自己的幸福人生,這沒有什麼問題,問題只是:卻無人肯花心機去援助天國的工程.神說:「你們要省察你們到底是如何活的.」不將神放在優先次序的生活,到頭來,只是一種毫無所得的生活.現代人活得很不滿足,很多心理或精神的困擾,卻堅持只為自己打算,活在自我的世界中,從不相信神,從不善待他人.神說:這些人,吃,卻不得飽；喝,卻不得足；賺到的金錢,卻裝在破漏的銀包裡.這些人,一點都不滿足.

該1:7　萬軍之耶和華如此說:「你們要省察自己的行為【重複第5節,「行為」原意:道路,方向,生活方式】.

該1:8　你們要上山取木料,建造這殿,我就因此喜樂,且得榮耀.這是耶和華說的.

該1:9　你們盼望多得,所得的卻少；你們收到家中,我就吹去.這是為什麼呢?因為我的殿荒涼,你們各人卻顧【原文是奔】自己的房屋.這是萬軍之耶和華說的.

香港人說,「為口奔馳」,那些猶大人卻「為房屋奔馳」.簡單說,他們為生活營營役役,他們忘了人生.在此,神不厭其煩,再講一次,你們人生的優先次序,到底是什麼呢?你們真的將最重要的事,放在首位嗎?你們投注在自己的房屋上,資產上,但有沒有想過將生命投注在神的事業上呢?神問當時的猶大人,為何不去建聖殿呢?

在這裡,神教我們思考「生命的因果」.我們以為,賺得不夠,故無時間心力去事奉神了.神卻說,我們的因果搞錯了,正正因為我們不用心去事奉神,故我們永遠都賺得不夠.人生得多得少,這是神的恩典範圍.我們只顧自己,不理神,我們得到的,反而更少.我們尊重神,事奉神時,我們得到的,反而更豐足.

人的目光很短淺,只看到自己的需要,這令人的生命得不到益處.神教人有遠大的眼光,體會天心之所悅,渴求神聖之榮美,這廣闊的心靈讓人看到生命的豐盛所在.

主耶穌後來的講法也一樣:「不要憂慮說,吃什麼?喝什麼?穿什麼?這都是外邦人所求的.你們需用的這一切東西,你們的天父是知道的.你們要先求祂的國和祂的義,這些東西都要加給你們了.」(太6:31-33)

神這樣說,你們猜,百姓的反應如何呢?

該1:12　那時,撒拉鐵的兒子所羅巴伯和約撒答的兒子大祭司約書亞,並剩下的百姓,都聽從耶和華─他們上帝的話和先知哈該奉耶和華─他們上帝差來所說的話；百姓也在耶和華面前存敬畏的心.

該1:13　耶和華的使者哈該奉耶和華差遣對百姓說:「耶和華說:我與你們同在.」

該1:14　耶和華激動猶大省長撒拉鐵的兒子所羅巴伯和約撒答的兒子大祭司約書亞,並剩下之百姓的心,他們就來為萬軍之耶和華─他們上帝的殿做工.

該1:15　這是在大流士王第二年六月二十四日.

百姓的反應,有點出人意外.被擄前,先知講什麼,百姓都是不聽的.如今,百姓即時悔改,作出正面回應.有時,苦難的確是生命的好老師.

經文說,百姓敬畏神.神的回應:「我與你們同在」.一個人有神的同在,生命的需要和意義便有底,不用憂心了.

這時,是大流士王第二年六月二十四日,距離神初初發言的六月初一日,只是23日.嘩,這真是一快速的復興.短短23日,百姓便重整生命的優先次序,重整人生的方向.只要眼光能移離眼前的狹隘,望向廣闊的天地,望向悠長的歷史將來,我們會發現,有很多天國的事情比我們現時花盡心力去經營的事情重要得多.若我們都能有這樣的覺悟,重整人生,多好.

該2:1　七月二十一日,耶和華的話臨到先知哈該說:

該2:2　「你要曉諭猶大省長撒拉鐵的兒子所羅巴伯和約撒答的兒子大祭司約書亞,並剩下的百姓,說:

該2:3　你們中間存留的,有誰見過這殿從前的榮耀呢?現在你們看著如何?豈不在眼中看如無有嗎?」

該2:4　耶和華說:「所羅巴伯啊,雖然如此,你當剛強!約撒答的兒子大祭司約書亞啊,你也當剛強!這地的百姓,你們都當剛強做工,因為我與你們同在.這是萬軍之耶和華說的.」

七月二十一日,已是差不多1個月後了.來到工地,只見斷牆殘壁,荒涼至極.熱心,不會變走眼前的艱難.一腔熱誠,無法取消要背的十架.先知問那些超過50歲的百姓,你們誰見識過以前那聖殿的金雕玉砌,莊嚴宏偉?現在如何?現在如何?經文說,現在「看如無有」.一片廢墟,從哪處開始修建?越想,越是唏噓.

神說,要剛強,要剛強,要剛強,重要的事要說3次.領導要剛強,祭司要剛強,百姓也要剛強.在艱難時刻,任何一個環節軟弱放棄,都不可能成事.神再次承諾,與他們同在.不要垂頭看眼前的一無所有,當仰望神,神在,就有足夠的恩典了.不要懼怕困難,有神的恩典在,便夠了.像保羅所說:「我靠著那加給我力量的,凡事都能做.」(腓4:13)不要灰心,越是艱難,越顯得信心的寶貴,越體會神的恩典的真實.

該2:8　萬軍之耶和華說:「銀子是我的,金子也是我的.

該2:9　這殿後來的榮耀必大過先前的榮耀；在這地方我必賜平安.這是萬軍之耶和華說的.」

神叫這群資源缺乏的猶大人,想得再遠大一點.以前的聖殿,是所羅門王極榮華的時候建的,但所羅門王算得什麼?神說,金和銀都是我的,天地萬物都是我的,我讓你們建一座比所羅門王的聖殿更榮美的聖殿,有何難處呢?神教我們,在天國的事情上,在神的事情上,敢於想大一點.前人的成果是很美好的,但我們可以做得比他們更好,這不是不可能的.你的眼光若被當前的「一無所有」困住,失去信心,沒有勇氣,則將永遠都是一無所有.神說,「這殿後來的榮耀必大過先前的榮耀」,你們的目光要看到這將來的殿的榮美,才有勇氣將之建成.意識到你們對天國的貢獻,你們才會甘心委身.神說,在這新的聖殿裡,祂賜人平安.「平安」,shalom,指圓滿的生命.好好地建一座聖殿吧,好好地做些神聖的事情吧,藉這些東西,能令世人得到圓滿的生命的.

該2:18　你們要追想此日以前,就是從這九月二十四日起,追想到立耶和華殿根基的日子.

該2:19　倉裡有榖種嗎?葡萄樹,無花果樹,石榴樹,橄欖樹都沒有結果子.從今日起,我必賜福與你們.

9月24日,從神初初發言的6月1日到9月24日,只用了110多天,聖殿的根基便立好了.經文的意思是:從這日打後看,看看你們生活的變化.之前,榖種沒有,葡萄樹,無花果樹,石榴樹,橄欖樹也沒有果子,但從聖殿的根基立好的這日打後看,你們會明顯看到,神賜福給你們,倉裡有榖種,葡萄樹,無花果樹,石榴樹,橄欖樹都會結滿果子.你們可以看到,天地因你們立好聖殿的根基而即時改變.人在做,天地在變.人為神做的一件美事,會改變天地,造福多人.

我們看看哈該書裡最後的神諭.

該2:20　這月二十四日,耶和華的話二次臨到哈該說:

該2:21　「你要告訴猶大省長所羅巴伯說:我必震動天地.

該2:22　我必傾覆列國的寶座,除滅列邦的勢力,並傾覆戰車和坐在其上的.馬必跌倒,騎馬的敗落,各人被弟兄的刀所殺.

該2:23　萬軍之耶和華說:我僕人撒拉鐵的兒子所羅巴伯啊,到那日,我必以你為印,因我揀選了你.這是萬軍之耶和華說的.」

再提9月24日,這日真的很重要,資源短缺的人憑毅力立定聖殿的根基.神說,這日,就在這日,我必震動天地.用中國人的話說,這是驚天地,泣鬼神的日子.因人甘心建聖殿,不單天地因而改變,歷史也因而改道.2:22說的,是神消滅世上暴力之國,天國初現.2:23提到所羅巴伯,他的身分不再是波斯帝國的猶大省長,而是萬軍之耶和華的僕人.他以大衛王裔的身分,成了戴有皇印指環的君王,象徵著天國君王臨現人間.這是我們能望見的終極遠象:天國降臨,神的旨意行在地上.就因人願意建聖殿,我們主禱文祈求的內容開始實現.

哈該書很短,它將身處匱乏時代,前景未明的一群人的生命抉擇展示給我們看.這群人,大有理由,因困乏而自保,拒絕天命,圍爐取暖,只看眼下需要,只顧個人利益,無遠大目光,形役於物質.或許,他們能好好地活下去,但人生沒有半點永恆意義.

但我們看到這群人,超越了物質的當下需要,為了神的榮耀,為了神臨在人間,為了一種天命,他們努力張羅資源,建造聖殿.這群人敬畏神,明白何為信仰神,他們重整了人生的優先次序,放下了自己的渴求,先求神的國和神的義.就在他們立定聖殿根基的那日,聖經說,天地改變,大地開始結出豐盈的果實,歷史開展其天國的進程.只要你走上靈性的高樓,登高望遠,會望到神展示的宏大遠景.為這遠象,只要你為神做一點事情,那怕只是小小的工程,都能驚天動地,推動天國的進程.

你能否體會神的心意?你有沒有信仰的遠象?你是否願意為此付出一點心力?

願神賜福你們.

願榮耀歸予聖父,聖子,聖靈.



\newpage

\section{閉起雙眼前你最掛念誰}
\label{sec:770}
\textbf{林豪恩先生 - 閉起雙眼前你最掛念誰?}
\newline
\newline
\hyperref[sec:769]{< < < PREV SERMON < < <}
~
\hyperref[sec:toc]{[返目錄]}
~
\hyperref[sec:771]{> > > NEXT SERMON > > >}


提摩太後書 4:6-18
\newline
\begin{longtable}{cl}
\hline
\hline
章節 & 經文 (和合本修訂版)\\
\hline
4:6 & \begin{tabularx}{0.7\textwidth}{X} 至於我，我已經被澆獻，離世的時候到了。 \end{tabularx} \\ \\ \relax
4:7 & \begin{tabularx}{0.7\textwidth}{X} 那美好的仗我已經打過了，當跑的路我已經跑盡了，該信的道我已經守住了。 \end{tabularx} \\ \\ \relax
4:8 & \begin{tabularx}{0.7\textwidth}{X} 從此以後，有公義的冠冕為我存留，就是按著公義審判的主到了那日要賜給我的；不但賜給我，也賜給凡愛慕他顯現的人。 \end{tabularx} \\ \\ \relax
4:9 & \begin{tabularx}{0.7\textwidth}{X} 你要趕緊到我這裡來。 \end{tabularx} \\ \\ \relax
4:10 & \begin{tabularx}{0.7\textwidth}{X} 因為底馬貪愛現今的世界，已經離棄我，往帖撒羅尼迦去了；革勒士往加拉太去；提多往撻馬太去； \end{tabularx} \\ \\ \relax
4:11 & \begin{tabularx}{0.7\textwidth}{X} 只有路加在我這裡。你來的時候把馬可帶來，因為他在服事上於我有益。 \end{tabularx} \\ \\ \relax
4:12 & \begin{tabularx}{0.7\textwidth}{X} 我已經打發推基古往以弗所去。 \end{tabularx} \\ \\ \relax
4:13 & \begin{tabularx}{0.7\textwidth}{X} 我在特羅亞留給加布的那件外衣，你來的時候要帶來，那些書也帶來，特別是那幾卷羊皮的書。 \end{tabularx} \\ \\ \relax
4:14 & \begin{tabularx}{0.7\textwidth}{X} 銅匠亞歷山大多方害我；主必照他所行的報應他。 \end{tabularx} \\ \\ \relax
4:15 & \begin{tabularx}{0.7\textwidth}{X} 你也要防備他，因為他極力抗拒我們的話。 \end{tabularx} \\ \\ \relax
4:16 & \begin{tabularx}{0.7\textwidth}{X} 我初次上訴時，沒有人前來幫助，竟都離棄了我，但願這罪不歸在他們身上。 \end{tabularx} \\ \\ \relax
4:17 & \begin{tabularx}{0.7\textwidth}{X} 惟有主站在我身邊，加給我力量，使我能把福音完整地傳開，讓所有的外邦人都聽見；我也從獅子口裡被救出來。 \end{tabularx} \\ \\ \relax
4:18 & \begin{tabularx}{0.7\textwidth}{X} 主必救我脫離一切的兇惡，也必救我進他的天國。願榮耀歸給他，直到永永遠遠。阿們！ \end{tabularx} \\ \\
[1ex]
\hline
\hline
\end{longtable}
講題:閉起雙眼前你最掛念誰?Who Do You Miss the Most Before You Closed Your Eyes

經課:提摩太後書4章6至18節

講員:林豪恩先生

場合:香港中文大學崇基學院禮拜堂主日崇拜

日期:2025年10月26日

一．閉起雙眼前掛念甚麼?

聽說有一令天下男人苦惱的難題,你猜到是甚麼嗎?就是:「老婆同阿媽一齊跌落水,你會救邊個先?」網傳有位台灣男士成功破解了這條千古難題.他的破解方法,我們不在這裡劇透.不過,如果他的老婆或阿媽見到今天的講道題目,可以製造另一難題給他,就是:「閉起雙眼前,你最掛念老婆還是阿媽?」如果那個男人膽子夠大,他一句說話就可以破解:「兩個都不掛念!」這個回答有點無情,但在某程度上切合現實,因為市面上充斥著那些「死之前要做的五件事」,「死之前要做的十件事」多數都是與老婆或媽媽無關的.例如有網友列出死之前要做的十件事是:(一)環遊世界；(二)寫小說；(三)寫詩；(四)退休；(五)清書；(六)清戲；(七)挑戰極限運動；(八)海外生活；(九)開店做老闆；(十)擁有夢寐以求的居所.你認為這張清單如何呢?你也有死之前要做的多少件事的清單嗎?在今天的書信經課中,使徒保羅表示自己將要離世,他有沒有離世前要做的多少件事的清單呢?如果有,清單的內容是甚麼呢?

二．保羅如何看待死亡?

首先,讓我們看看保羅如何談論死亡.

4:6我現在被澆奠,我離世的時候到了.

保羅和大多數人一樣,沒有直接用「死亡」這個詞彙,他用的是「離世」這委婉的說法.不過,特別的地方是保羅用「被澆奠」來形容自己的離世.澆奠是在祭祀儀式中將酒澆在祭物上,使祭物更加馨香.保羅以此比喻他將自己作為馨香的祭物獻給主.保羅啟發了我們對死亡的想像.我們聽他說過把在生的日子當作活祭獻給上主,原來也可以視死亡為祭獻給上主.

人們有很多表達死亡的講法,不同講法隱含著對死亡的不同觀念,也帶來不同的感受.你慣常用甚麼講法表達死亡呢?粵語在這方面的詞彙非常豐富,例如:「瓜老襯」,「賣鹹鴨蛋」,「香jor 」,「收jor 皮」,「玩完」等等,這些詞彙隱含的看法也十分多元.

保羅用「被澆奠」來形容自己的死亡,作為馨香的祭物獻給上主.耶穌離世前所說的都類似:「父啊,我將我的靈交在你手裡!」(路23:46)有些人可能領受了類似保羅或耶穌的使命,以死為祭獻給上主,但我們大多數人的死亡並不如他們一樣轟轟烈烈,我們若能夠用「主懷安息」,「息勞歸主」,「榮歸天國」等詞彙已經感到不錯了.假如有一天你聽到有人說:「Walter走jor 」,你可能感到不安,自然會問:「佢去jor 邊?」但如果你聽到人說:「Walter榮歸天國」,你就會感覺安心,你知道他去了那裡,也感到這是光榮的,而不是羞辱.Walter當然不希望人們說他「去了賣鹹鴨蛋」,已經做了一世工作,死了還要做,聽到都覺得辛苦.如果說他「息勞歸主」,感覺就好得多,不用再操勞,安心「躺平」在上主的懷抱中.

三．保羅的清單有甚麼項目?

接著下來,讓我們看看保羅「死前要做的多少件事」的清單.

4:7那美好的仗我已經打過了,當跑的路我已經跑盡了,所信的道我已經守住了.

保羅的清單不是五件事或十件事,只是三項.這三個短句,原文都是用名詞開始的.

第一個名詞是「打美好的仗」的「仗」.「仗」這個字在軍事上指戰爭,在運動上是競賽.如果當作運動場上的競賽,保羅以美好來形容這場競賽.有時,我們也聽到運動員說,這場比場「好正」.「好正」的意思不一定拿了冠軍,可以比拿了冠軍更多意義.香港單車運動員李慧詩在2006 年一場訓練發生意外,被評定為「輕度殘障」,連教練也勸她退役.痛定思痛,她重新確定自己喜愛單車運動這份初心,堅持回到運動員的路軌上,比其他人努力百倍.就在那段日子,她信靠了上主.結果他在2007年復出,她說:「感謝上帝,靠著祂給我的力量,我再次堅強站起來,有機會重返單車賽場.」2021年東京奧運,李慧詩為香港增添一面銅牌,她在賽後接受訪問時說:「我知道很多人,包括自己都想贏得金牌,但人生就是這樣,即使很用力完成一件事,結果也有可能未如理想.但起碼自己盡過力,無悔就夠. 」

保羅說「那美好的仗我已經打過了」,不代表所有事情都完美,但他有落場參賽,全程投入,盡過力,流過血汗.回望時,他心滿意足地說:這場比賽「好正」.對基督徒來說,這場比賽或這場仗代表甚麼呢?我們可以參考耶穌曾經說過的話:

約翰福音17:15我不求你把他們從世上接走,只求你保全他們,使他們脫離那惡者.

基督徒生活在世界本身就是一場競賽,因為世界充滿引誘,隨時都有跌倒的可能.基督徒一方面感到聖靈的提示,另一方面又感受到那惡者的誘惑,到底要跟隨基督還是跟隨那惡者,要活出基督的樣式還是被那惡者擄去,這的確是一場不容易的仗.有些人覺得這場仗很辛苦,保羅沒有說不辛苦,但他說:「好正」,就正如喜歡踢足球的人落場踢波,雖然筋疲力盡,但卻感到「好正」.保羅啟發我們,可以用這種態度去投入人生,守住初心,盡力參與,到最後,我們向上主說:「這場競賽很多挑戰,真不容易,但「好正」」.

第二個名詞是「當跑的路我已經跑盡了」的「路」,原文只有「路程」和「完成」這兩個字.保羅的「路程」是甚麼呢?

使徒行傳20:24我卻不以性命為念,只要走完我的路程,完成我從主耶穌所領受的職分,為神恩典的福音作見證.

當保羅說「當跑的路我已經跑盡了」一方面表示完成了他領受的人生使命,另一方面,為了完成使命,他也真的跑了很多路.他走遍耶路撒冷,敘利亞,土耳其,希臘和羅馬.保羅知道自己「當跑的」是甚麼路,我們知道自己「當跑的」是甚麼路嗎?有人的路容易一點,有人的路比較艱難,其實可能都是「你睇我好,我睇你好」而已.有人跑得快,有人跑得慢,有人一面跑一面欣賞沿途風光,各適其式都同樣適宜.重要的不是快慢,而是沒有誤入歧途,也沒有放棄.保羅深信沿路有基督陪他跑,也深信在終點有人等著他分享完成的喜悅.

第三個名詞是「所信的道我已經守住了」的「信」,原文也是兩個字,分別是「信心」及「持守」.持守信心到最後原來並不是一件容易的事,保羅在下文提到一些持守不到信心的例子:

4:10因為底馬貪愛現今的世界,就離棄我往帖撒羅尼迦去了.

4:14-15銅匠亞歷山大……極力敵擋了我們的話.

正所謂「路遙知馬力,日久見人心」,到最後的日子仍能說「美好的仗我已經打過了,當跑的路我已經跑盡了,所信的道我已經守住了」其實是不容易的事.試看看聖經中著名的大衛王,他曾經是優秀的戰士,出色的君王,專心倚靠上帝.然而,後期的大衛卻弄得家破人亡,宮廷內亂,謀朝篡位,造成國家動盪,人民受苦.他的繼任人所羅門王又怎樣呢?他智慧超群,締造了以色列的全盛時期,並且建造聖殿.然而,晚年的所羅門王可以用白居易的兩句詩來形容:「春宵苦短日高起,從此君王不早朝」.列王記上說他有700位妻子和300位嬪妃,這些外邦公主和妃嬪帶來偶像崇拜,耶路撒冷附近充斥著祭祀偶像的丘壇.從大衛王和所羅門王可見,保住晚節其實是不容易的.最近讀到吳思源先生在基督教週報的短文,名為Finishing Well,有幾句如此說:

「曾立華牧師的安息禮拜,有逾千人在炎日下出席.有送花的店員問我,這位牧師是甚麼來頭的大人物?為甚麼有那麼多人來?我說他不是甚麼大教會的大牧師,也不是蜚聲國際的大佈道家,是他一生忠心的服事和崇高的品格,贏得無數人的愛戴.他的二兒子特別為父親創作了一首詩歌,歌名是“ Finishing Well”,在安息禮中獻唱.他說這是父親一生持守的信念,行事為人要貫徹始終,堅持到底,圓滿完成.真的,人生除了要活得好,同樣要完成得好.廣東話有句說話叫「衰收尾」,比「功虧一簣」更加糟糕.曾牧師語重心長勸喻我們做人要finishing well,而他自己更加是身體力行.」

四．牽掛仍然是人

雖然說「死之前要做的多少件事」多數都和個人想望有關,但有臨床觀察指出人在閉起雙眼前其實最掛念的仍然是人.在美國有一位87歲的老人家希望在他生命最後時刻被親人環繞,並且能再喝上一口自己最心愛的啤酒.在生命的最後時光,他在家人的簇擁下,有說有笑,一起回憶曾經度過的幸福時光.他的兒子用鏡頭記錄下這寶貴時刻,並把照片在自己家人朋友圈上分享.這張在生命最後時刻和家人一齊品嚐啤酒的照片傳開了,引起眾多陌生人的共鳴,紛紛表示得到安慰及啟發.在保羅的信中,我們看見他在離世之前也想見一些人,想起一些人.

第一,他想見提摩太和馬可:

4:9你要趕緊地到我這裡來.

4:11你來的時候,要把馬可帶來.

第二,他想到隊友:

4:11獨有路加在我這裡.

4:10革勒士往加拉太去；提多往撻馬太去；

4:12我已經打發推基古往以弗所去.

第三,他還想起令他難受的人:

4:10因為底馬貪愛現今的世界,就離棄我往帖撒羅尼迦去了；

4:14銅匠亞歷山大多多地害我；主必照他所行的報應他.

人生就是如此,有令你想見的人,也有令你難受的人.有人認為人生走到最後一程,應該甚麼也放下,無牽無掛地離開才美滿.如果仍有牽掛又如何呢?有人認為應該放下仇怨,無所執著地離開才是福氣.如何仍然不能忘記又如何呢?孫中山先生離世之時對局勢仍然憂心忡忡,留下遺言:「革命尚未成功,同志仍須努力.」保羅何嘗不是仍然掛心如樹苗初長的各地教會呢?他為甚麼仍然記住那些離經叛道的人呢?

4:15你也要防備他,因為他極力敵擋了我們的話.

他為的是要提醒提摩太,保護如幼樹苗般脆弱的初期教會.用孫中山先生遺言的格式重述保羅對提摩太的話,可以這樣說:「教會尚未穩健,提摩太仍須努力」.

五．最終把自己交托給誰?

「那美好的仗我已經打過了,當跑的路我已經跑盡了,所信的道我已經守住了」,這是令人肅然起敬的豪言壯語.然而,說出這番豪言壯語的保羅並沒有報喜不報憂,沒有美化自己的狀況:

4:16我初次申訴,沒有人前來幫助,竟都離棄我；但願這罪不歸與他們.

4:17惟有主站在我旁邊,加給我力量,使福音被我盡都傳明,叫外邦人都聽見；我也從獅子口裡被救出來.

4:18主必救我脫離諸般的凶惡,也必救我進他的天國.

生命的路途上有機會遇到離棄自己的人,有機會遇到諸般凶險.人生在世,或多或少,或輕或重,誰人不是如此的呢?

保羅如何面對呢?「有主站在我旁邊,加給我力量,使福音被我盡都傳明」.保羅懷著甚麼信念面對呢?「主必救我脫離諸般的凶惡,也必救我進他的天國」.在閉起雙眼前,無論我們尚有多少件死之前要做的事未完成,或者仍然掛念著甚麼人,保羅提示我們轉念思考一個問題:我們把自己交托給誰呢?保羅已經尋找到答案,他把自己交托給上主,他深信縱使他自己仍有牽掛,對於某些人及事仍然不能忘懷,他可以把這些人和事交托給上主,也把自己交托給上主,在上主裡面可以接納一切,融和一切,就算人間未臻圓滿,在上主裡面最終必然圓滿.讓我們以保羅這句滿懷信心和期盼的話作結:

4:8從此以後,有公義的冠冕為我存留,就是按著公義審判的主到了那日要賜給我的；不但賜給我,也賜給凡愛慕他顯現的人.



\newpage

\section{這要到幾時}
\label{sec:771}
\textbf{李均熊牧師 - 這要到幾時?}
\newline
\newline
\hyperref[sec:770]{< < < PREV SERMON < < <}
~
\hyperref[sec:toc]{[返目錄]}
~
\hyperref[sec:772]{> > > NEXT SERMON > > >}


哈巴谷書 1:1-4
\newline
\begin{longtable}{cl}
\hline
\hline
章節 & 經文 (和合本修訂版)\\
\hline
1:1 & \begin{tabularx}{0.7\textwidth}{X} 哈巴谷先知所看見的默示。 \end{tabularx} \\ \\ \relax
1:2 & \begin{tabularx}{0.7\textwidth}{X} 耶和華啊，我呼求，你不應允，要到幾時呢？我向你呼喊「暴力！」你還不拯救？ \end{tabularx} \\ \\ \relax
1:3 & \begin{tabularx}{0.7\textwidth}{X} 你為何使我看見罪孽？你為何坐視奸惡呢？毀滅和兇暴在我面前，爭執與紛爭不斷發生。 \end{tabularx} \\ \\ \relax
1:4 & \begin{tabularx}{0.7\textwidth}{X} 因此律法無效，公理從未彰顯。惡人圍困義人，所以公理遭受扭曲。 \end{tabularx} \\ \\
[1ex]
\hline
\hline
\end{longtable}
講題:這要到幾時?  How Long?

經課:哈巴谷書1:1-4;2:1-4

講員:李均熊牧師

場合:香港中文大學崇基學院禮拜堂主日崇拜

日期:2025年11月2日

起:先知的問題——這要到幾時?

今日舊約經課選用了哈巴谷書第一章和第二章的前四節.

但要瞭解哈巴谷書,關鍵在於一開始先知的提問:「要到幾時呢?」先知呼求上主,正因為他也知道上主也不喜悅強暴和罪孽；但最困擾先知的,是這種強暴,不公義要到何時?哈巴谷不是問「為甚麼會有這些事」；也不是問神為甚麼這樣不公平；甚至可以說,先知並非想求個明白,而只是希望事情過去.

當然,先知也知道在1:6-11中提到的迦勒底人,是上主興起來「為要刑罰人」和「懲治人」(「耶和華啊,你派定他為要刑罰人；磐石啊,你設立他為要懲治人.」1:12),但上主的審判,也令義人受苦,因為先知接著便說:「行詭詐的,你為何看著不理呢?惡人吞滅比自己公義的,你為何靜默不語呢?」(1:13)與其說先知不明白,不如說先知所經歷和看到的,是和他所了解和相信的,有太大出入.用今日的術語,這叫做認知失調(cognitive dissonance).

但信仰是為著要找答案嗎?哈巴谷書正是要告訴我們,信仰並不是來替苦罪不公找答案的.不是沒有答案,有時是答案太荒謬,荒謬到一個地步,是我們在理智上要拒絕的.在猶太人經歷納粹大屠殺之後,一些猶太人的思想家(他們有時拒絕被稱為神學家)就認為上主對當中猶太人的苦難和創傷不發一言,並不是不能解釋,而是不應解釋,因為這些邪惡並不應該在人間出現,你越解釋它,就好像要為大屠殺這種荒謬的苦難創傷在世上找一個合理的位置.面對如斯苦難,沉默反而是最皆的反應.

所以哈巴谷就算問why,都是為何要使我看見「惡人圍困義人,公理顯然顛倒」(1:4)這些荒謬的事情；而哈巴谷更關心的是上主不理會這些事情(1:13「你為何看著不理呢?惡人吞滅比自己公義的,你為何靜默不語呢?」),「要到幾時呢?」

承: 可以置身事外嗎?

哈巴谷所所見到的,是怎樣的惡呢?第一章看似是上帝派外國的軍隊去蹂躪自己的國家,但在第二章,哈巴谷所看到的,又不似一個軍事強國對國家的侵略.其實和合本／和修本都將經文沒有交代惡人身份的地方,都配上了迦勒底人(2:4, 5, 6) ,而令我們以為第二章的惡人仍然是1:6迦勒底人.但是1:4「因此律法放鬆,公理也不顯明；惡人圍困義人,所以公理顯然顛倒」當中用上律法(torah)和公理(mishphat),卻都是上主與子民立約的用語,明顯是針對猶大國內的罪行.

其實迦勒底人在哈巴谷全書只在1:6出現一次.若我們仔細看第二章的經文,則會看到當中的惡人是對城鎮經濟上的剝削者,和拜偶像的罪行.例如2:6 「這些國的民豈不都要提起詩歌並俗語譏刺他說:禍哉!迦勒底人,你增添不屬自己的財物,多多取人的當頭,要到幾時為止呢?」首先將中文翻譯加上去的迦勒底人拿走,就會見到那是一些過份的借貸,「增添不屬自己的財物,多多取人的當頭」；接著2:9「為本家積蓄不義之財,在高處搭窩,指望免災的」更似是發國難財,又可以住在安全地方的地方領導的罪行.其他如2:12「以人血建城」——這個惡人將城建在血上；2:15「給人酒喝,又加上毒物,使他喝醉,好看見他下體的有禍了!」——使鄰人醉倒以至羞辱他；2:8和2:17重複了「向國內的城並城中一切居民施行強暴.」這些都是針對自己國家人民的惡行,多於外國的侵略.至於2:18-19節「雕刻的偶像,人將它刻出來,有甚麼益處呢?鑄造的偶像就是虛謊的師傅.製造者倚靠這啞巴偶像有甚麼益處呢?對木偶說:醒起!對啞巴石像說:起來!那人有禍了!這個還能教訓人嗎?看哪,是包裹金銀的,其中毫無氣息.」明明就是針對猶大國中偶像的問題.

第二章經文的最終形式邀請讀者將惡人視為巴比倫人,但正如我剛提到,在某些地方(1.4和2.6-17的部分)對他們的描述暗示他們是猶大貴族.

很大可能,哈巴谷書成書時所針對的惡行,乃出現於在晚期猶大人仍有國君的王國時期,先知其實向耶和華抱怨,猶大上層階級的社會虐待.在第二章開首,他得到了神的回應,並被指示將其寫下來.這個回應是一系列的禍哉預言,宣告那些剝削窮人的人的厄運.但當耶路撒冷被毀後,這篇作品通過添加1:5-11和15-17的神諭進行了重新編寫,這些神諭描繪了巴比倫人進攻的暴力.這些插入部分將巴比倫人描繪為耶和華用來懲罰猶大罪行的工具,因為猶大貴族過度擴展了自己的權限,而招致了罪責.

對於後來的讀者,可能已經對猶大皇族的惡行沒有太多印像,但巴比倫毀滅國家的創傷仍瀝瀝在目,故此哈巴谷書後來世代所看到的是猶大貴族領袖的惡,混合著侵略者對國家的蹂躪.

//然而,同樣有可能的是,在後來的這個階段,流亡的聽眾會將1.2-4中的惡人直接理解為巴比倫人.這種重新編寫也可以在第二章的「禍哉」神諭中觀察到.這裡有一些添加將對社會不公的字面描述,轉變為國際事務中非法使用權力的隱喻.例如,2:6-8在預言債權人通過不正當金融交易積累財富的人的「禍哉」神諭(2.6b-7a)中,添加了一句話(2.7b-8),將這些債權人識別為“許多國家”和“剩餘的人民”.這一添加將2.6b中的貸款隱喻為2.8b中的流血和暴力.// (Hadjiev 2020)

為何要將這個編修成書過程告訴大家,其實要提醒大家哈巴谷書所呈現的並不單是外敵的入侵,更是如經文所記1:12,是上主透過侵略者的惡,去刑罰國家的惡.

這也提醒我們,不是只要我們做好人義人(或只是不參與社會上的惡),就能免去這些惡事——因為上主審判社會上流階層,其實會令社會各個階層都受苦.我們並不能置身苦罪之外,我們也是罪中之人,所以不能自以為義,以為別人受罰應該,自己受罰是不該.我們在社會的惡面前,就算不是施罪者,也會是被罪者,社會的惡所引發更大的惡,是我們不能逃避的.但我們也不能否認,我們已經盡量做好,也仍會受苦,那是因為有更大的邪惡在運作.我們是「被罪者」——即受害者……邪惡本質上就是邪惡,上帝不會造一些惡事來「塑造」我們,上帝的塑造只在乎我們對惡的回應.

哈巴谷正是不會逃避,亦不會因為面對不公義而放棄.哈巴谷就著社會的惡去質問上主,正是因為他不能接受「係咁架啦」,或只著眼自己的屬靈小天地,小問題,而免疫於社會的惡對人所帶來的創傷苦痛.

轉:信仰承載苦難人生

哈巴谷對神沒有信心嗎?不!哈巴谷正是相信神不會不理這種荒謬的邪惡,才會去面質主.先知並不是不理苦難, 他仍然勇於來到神前問:「這要到幾時?」 他要保持和上帝的關係和溝通,他就是要捉住神,不放過神.他並不是要找答案,但是他「就黎頂唔順」 所以先問, 到幾時呀!!!!!! 神的回答卻又不是給他一個時限, 只是強化了 「我最終會攪惦」 2:3: 因為這默示有一定的日期,快要應驗,並不虛謊.雖然遲延,還要等候；因為必然臨到,不再遲延.就像說 ‘hang on there!’ 所以,信仰帶來盼望,但是我們仍要面對爭扎的.

合:雖然…仍要….苦中承擔

社會的罪帶來社會的審判,無人能免,所以基督徒不能置身苦罪之外.卻要學習苦中承擔,作基督的見証.

哈巴谷書的結束,或許能作為我們對世間不公和苦罪的回應.

雖然哈三詩題標為哈巴谷的祈禱,但其內容實是一首敬拜的詩歌.全詩強調的是上帝的榮耀和得勝!

他的榮光遮蔽諸天；頌讚充滿大地.他的輝煌如同日光；從他手裡射出光線,在其中藏著他的能力.(3:3-4)

耶和華啊,你乘在馬上,坐在得勝的車上……你乘馬踐踏紅海,就是踐踏洶湧的大水.(3:8, 15)

詩人在敬拜的最後回應是3:17-19

3:17雖然無花果樹不發旺,葡萄樹不結果,橄欖樹也不效力,田地不出糧食,圈中絕了羊,棚內也沒有牛；3:18然而,我要因耶和華歡欣,因救我的神喜樂.3:19主耶和華是我的力量；他使我的腳快如母鹿的蹄,又使我穩行在高處.這歌交與伶長,用絲弦的樂器.

雖然… 仍要(not because of, but in spite of….)　——因為人並不能自救,只有信靠耶穌基督的救恩；而信徒完全的得勝,也只有末日基督再來才可以成就.其間只好憑信心堅持去問:「要到幾時呢?」

但如何得到這樣的信心?

哈三作為敬拜頌讚詩歌告訴我們:只有透過敬拜上帝,經歷上帝的榮耀,我們就有更多資源去面對苦困和不公.

信仰是承擔,是承載,不是用來解決問題的——當然,信仰的力量也能幫助我們解決問題.問題得到解決,也是信仰幫助我們承載生命的方法.

但那不是最主要的.

更重要的是在上主的榮耀中敬拜,使我們有堅持下去的力量和智慧.



\newpage

\section{來認識活人的上帝}
\label{sec:772}
\textbf{高國雄牧師 - 來認識活人的上帝}
\newline
\newline
\hyperref[sec:771]{< < < PREV SERMON < < <}
~
\hyperref[sec:toc]{[返目錄]}
~
\hyperref[sec:773]{> > > NEXT SERMON > > >}


路加福音 20:27-38
\newline
\begin{longtable}{cl}
\hline
\hline
章節 & 經文 (和合本修訂版)\\
\hline
20:27 & \begin{tabularx}{0.7\textwidth}{X} 有些撒都該人來見耶穌。他們說沒有復活這回事，於是問耶穌： \end{tabularx} \\ \\ \relax
20:28 & \begin{tabularx}{0.7\textwidth}{X} 「老師，摩西為我們寫下這話：『某人的哥哥若死了，有妻無子，他該娶哥哥的妻子，為哥哥生子立後。』 \end{tabularx} \\ \\ \relax
20:29 & \begin{tabularx}{0.7\textwidth}{X} 那麼，有兄弟七人，第一個娶了妻，沒有孩子死了。 \end{tabularx} \\ \\ \relax
20:30 & \begin{tabularx}{0.7\textwidth}{X} 第二個、 \end{tabularx} \\ \\ \relax
20:31 & \begin{tabularx}{0.7\textwidth}{X} 第三個也娶過她；同樣地，七個人都娶過她，沒有留下孩子就死了。 \end{tabularx} \\ \\ \relax
20:32 & \begin{tabularx}{0.7\textwidth}{X} 後來，那婦人也死了。 \end{tabularx} \\ \\ \relax
20:33 & \begin{tabularx}{0.7\textwidth}{X} 那麼，在復活的時候，那婦人是哪一個的妻子呢？因為他們七個人都娶過她。」 \end{tabularx} \\ \\ \relax
20:34 & \begin{tabularx}{0.7\textwidth}{X} 耶穌對他們說：「這世代的人有娶有嫁， \end{tabularx} \\ \\ \relax
20:35 & \begin{tabularx}{0.7\textwidth}{X} 惟有配得那要來的世代和從死人中復活的人不娶也不嫁。 \end{tabularx} \\ \\ \relax
20:36 & \begin{tabularx}{0.7\textwidth}{X} 因為他們不能再死，和天使一樣；既然是復活的人，他們就是神的兒子。 \end{tabularx} \\ \\ \relax
20:37 & \begin{tabularx}{0.7\textwidth}{X} 至於死人復活，摩西在《荊棘篇》上就指明了，他稱主是亞伯拉罕的神，以撒的神，雅各的神。 \end{tabularx} \\ \\ \relax
20:38 & \begin{tabularx}{0.7\textwidth}{X} 神不是死人的神，而是活人的神，因為對他來說，人都是活的。」 \end{tabularx} \\ \\
[1ex]
\hline
\hline
\end{longtable}
講題:來認識活人的上帝  Come to Know the God of the Living

經課:路加福音20:27-38

講員:高國雄牧師

場合:香港中文大學崇基學院禮拜堂主日崇拜

日期:2025年11月9日

27撒都該人常說沒有復活的事.有幾個來問耶穌說: 28「夫子!摩西為我們寫著說:『人若有妻無子就死了,他兄弟當娶他的妻,為哥哥生子立後.』 29有弟兄七人,第一個娶了妻,沒有孩子死了； 30第二個,第三個也娶過她； 31那七個人都娶過她,沒有留下孩子就死了. 32後來婦人也死了. 33這樣,當復活的時候,她是哪一個的妻子呢?因為他們七個人都娶過她.」 34耶穌說:「這世界的人有娶有嫁； 35惟有算為配得那世界,與從死裡復活的人也不娶也不嫁； 36因為他們不能再死,和天使一樣；既是復活的人,就為上帝的兒子. 37至於死人復活,摩西在荊棘篇上,稱主是亞伯拉罕的上帝,以撒的上帝,雅各的上帝,就指示明白了. 38上帝原不是死人的上帝,乃是活人的上帝；因為在他那裡,人都是活的.」

今日的福音書經課.

我們看到耶穌與當時的一些領袖,包括一群稱為撒都該人的人,討論了關於復活的事宜.

讓我們先再一次回顧這個故事.

首先,有人來問耶穌一個情況,就是說,

V 28夫子!摩西為我們寫著說:『人若有妻無子就死了,他兄弟當娶他的妻,為哥哥生子立後.』

這個其實有一個背景,我們首先要去認識這個描述.這個描述其實在聖經舊約的《申命記》第25章第5-10節.

摩西頒布律法時,其中一個內容就是講到:

申25:5-10

5「弟兄同居,若死了一個,沒有兒子,死人的妻不可出嫁外人,她丈夫的兄弟當盡弟兄的本分,娶她為妻,與她同房. 6婦人生的長子必歸死兄的名下,免得他的名在以色列中塗抹了. 7那人若不願意娶他哥哥的妻,他哥哥的妻就要到城門長老那裡,說:『我丈夫的兄弟不肯在以色列中興起他哥哥的名字,不給我盡弟兄的本分.』 8本城的長老就要召那人來問他,他若執意說:『我不願意娶她』, 9他哥哥的妻就要當著長老到那人的跟前,脫了他的鞋,吐唾沫在他臉上,說:『凡不為哥哥建立家室的都要這樣待他.』 10在以色列中,他的名必稱為脫鞋之家.」

弟兄沒有兒子死了,死人的妻不可以出嫁外人,他的兄弟當盡弟兄的本分,取其妻,與他同房.婦人生的長子必立在死兄的名下,免得他的名在以色列中被塗抹.

這是一個特別的描述,其實當時回應著一個社會的狀況.社會的狀況就是,當一個丈夫過世,而她沒有兒子,留下的寡婦,當時摩西的教導或者他們的傳統,就是要他的弟弟去娶哥哥的妻子.

也就是所謂的「叔娶寡嫂」.

這樣,其實主要原因是想讓當時的人的生活狀況得以完整.

為什麼這樣說呢?因為當時的狀況是,當一個男士離世,沒有兒子的話,他所留下的妻子就沒有社會地位,沒有角色,因為當時古代社會,女性是要依靠男性生活.如果她失去了丈夫,又沒有兒子,她就要改嫁,或者要找其他辦法依賴生活.如果她改嫁外人,她與及她之後生的孩子,就是外人,不再是亡夫家族的人,這家族的血脈就會因此而終斷了.

另外,離世了的哥哥,因沒有後裔,家業沒有後人承繼,所以為了他們承繼家業,家族裡的人就有本分,有責任去保護,去處理.

而當時所提供的處理方法,就是做弟弟的要有責任,要娶離世的哥哥的妻,並為他生兒承繼.

所以叔娶寡嫂這舉動,在於要為他立後而生的長子,就要當作是哥哥的兒子,成為哥哥血脈的一部分,這樣哥哥的血脈就可以延續.這可以令到整個家族脈絡得以成全,也可以因此而減低了很多紛爭,並且寡婦會得到照顧等等.所以有這樣的律例,這應該也是當時一直進行的運作.

但是當到了耶穌時代,這班人來問耶穌,他們的關注點,並不是究竟這件事合不合適,是不是應該繼續運作及延伸.究竟叔娶寡嫂的特色,方法是不是好事.

他們不是在討論這個內容,甚至裡面的道德等等也沒有討論過.

他們反而提出一個問題,你看《路加福音》第20章第29至33節這樣講:

29有弟兄七人,第一個娶了妻,沒有孩子死了； 30第二個,第三個也娶過她； 31那七個人都娶過她,沒有留下孩子就死了. 32後來婦人也死了. 33這樣,當復活的時候,她是哪一個的妻子呢?因為他們七個人都娶過她.」

這是一個特別的問題,我們可能沒有想過.

當然,今時今日社會更加複雜,甚至離婚再婚很多狀況,可能比這個更複雜,但我們的關注點都應該是,他們的生活是怎樣,或者上帝怎樣看這件事.

但這些撒都該人對耶穌所提問,背後所挑戰的重點及邏輯,就是因為解決不了人覺得難處理的複雜生活問題.因為這故事例子中,七個人都娶過同一個妻,又沒有留下兒子,因此人就無法處理這個妻屬誰,當然也不能同時屬於七兄弟,而這難題是因為如果人會復活就要面對,只要沒有復活的事發生,這難題就不需要面對,因此就不應該有復活的事.

其實這是一個好奇怪的邏輯.如果復活真的會發生,就算多位兄弟因這叔娶寡嫂的制度下而娶過同一人,到天家誰是誰的妻上帝絕對有能力處理及解決,沒有可能因為我們用我們現在的邏輯理解不到這件事,搞不清楚這麼複雜的問題,所以就不應該有復活.我們人生當中一定有大大小小我們無法解決處理,不能合我們自己的邏輯的事情,難道就可以因此而拒絕接受那超越我們所能想像的事情嗎?

而當時他們問這個問題,就是在挑戰耶穌:你講復活嘛?怎麼復活,到時這樣混亂.

很明顯,他們把焦點放錯了.

耶穌在這裡,將焦點擺得清楚.

其實我們的信仰也常常出現了焦點的問題.

我們按我們的邏輯去理解事情,當我們搞不清楚的時候,我們就把責任推在上帝身上.

我們搞不清楚,就會對上帝說:「你為什麼搞到我搞不清楚?」明明是我們的能力,或者我們不願意,或者我們不盡力,但當我們覺得我們沒有辦法處理的事情,我們就要歸咎於上帝,就如經典的問題:上帝,為什麼令我生在這個時代?生在這個屋企,都不是我決定的.同樣,很多人埋怨父母:為什麼生我,你都沒有問過我.當然這是邏輯問題,怎麼問你生你嘛?

事實上,很多時候我們就會將我們解決不了的問題,或者我們覺得在我們的理解上面,邏輯上面比較混亂,難處理的問題,我們就會歸咎於上帝,心底那一句就是「上帝你不對」.

就像這段經文,撒都該人所表達不信復活,復活不應該有的原因,就是因為復活就會搞到人生活混亂.沒有復活不就有這些混亂.

很有趣對不對?如果我沒有在這個混亂的世界出生,那我就不用受這些苦.

我們很容易用我們的邏輯,用我們的想法自圓其說,其實是在逃避,甚至否認一個上帝的話語,他的心意.拒絕相信上帝,上帝的作為不應該是這樣的,上帝不可能是這樣的,因為我「攪唔掂」我的事情.

耶穌怎麼回答呢?

34耶穌說:「這世界的人有娶有嫁； 35惟有算為配得那世界,與從死裡復活的人也不娶也不嫁；

耶穌的第一個重點,就是這世代,我們會有嫁娶,去到天家就沒有嫁娶的需要.當然,這重點其實不是在講天家有什麼什麼,而是指出,就是有些東西在我們生命裡面覺得重要,甚至是困擾的,去到天家,可能不再重要,因為在天國裡面,不再需要嫁娶,因為在世,人有需要有自己的後代,或者需要有個民族,要有那種關愛連結.去到天家,每一個人與主將有最深厚的連結,所以不再需要嫁娶,所以根本不再需要處理誰是誰的妻這問題.因為人之間的婚盟已經完結.

36因為他們不能再死,和天使一樣；既是復活的人,就為上帝的兒子.

耶穌繼續指出,因為人死後,在天國裡,就有永生,而不能再死,在天國裡的「生命」本質上與現世的生命完全不同,所以現世所面對短暫的問題,到那日就不再需要面對.

37至於死人復活,摩西在荊棘篇上,稱主是亞伯拉罕的上帝,以撒的上帝,雅各的上帝,就指示明白了. 38上帝原不是死人的上帝,乃是活人的上帝；因為在他那裡,人都是活的.」

所以,這裡耶穌解說為什麼人需要復活,因為人本來就是活的,因為我們有上帝活的.

事實上,有一天我們都會離開這肉身的世界,無論如何,我們都會有這一天.

但死後的世界是不會再死.因為在主裡面,這個復活的生命是永恆生命,這是一個很重要的認知.

當然,我們在天國裡面那種生活的形態,不是我們所能完全想像得到的.我們常常以為生命的事,我們是可以明白得到,想像得到,可以用我們現在這一套去理解.

例子:在慕道班,當教授到使徒信經時,就會討論到「身體復活」這問題.

必會有人問,復活的個身體是那個身體啊?是10歲的,是50歲的,還是90歲?我現在年紀大才信主,我不喜歡年老的身體…

結論當然是我不知道.

但我知道一定是最好的,是大家也能認出來的,更是最榮耀的身體.

因為這就是奧秘.

在信仰當中一定有不少我們無法理解(透徹理解)的事情.

但我常常分享,只要我們願意盡力地去回應那些我們能理解的就已經很足夠了…

我們常常用著我們的智慧,想理解上帝的所有,當然用智慧用知識去追求以至去認識上帝更多,這是重要的,所以我們要有學識,我們要看聖經,我們要研讀聖經,我們要讀神學.

但更重要的是我們要明白,人生裡面有好多個奧秘,在上帝裡面有好多個奧秘,是我們沒辦法用我們現有的一切邏輯理解到的.如果我們困在我們的邏輯度,想像一些古怪的場景,狀況,就如這些希望為難耶穌的撒都該人,七個兄弟跟一個妻子的故事…

這根本不是在實際生活上要應對,要被牧養的事情.

我們身邊不是有更多需要我們關心的人和事嗎?

要學懂接受生命中的限制.

我們在限制當中活,就是在一個盲人摸象,或者是隔著望遠鏡看東西的狀況裡面.

但要謹記:我們的上帝是活人的上帝,所以我們不用擔心,不用掛心,我們必會復活,每一個人都要復活,每一個人去到主面前,都有新的身體,新的生活,特別那些被上帝所揀選,能夠在主裡面得救的,就為天國裡面享受永恆的美好.  我們會同樣經歷復活的大能,這是上帝的心意及應許,請憑信心去接受的,這是有一天我們期盼的,這是我們信仰裡面好重要的目標.

接著留下的問題就是我們而今在世要如何的「活」,

今日經課去到38節,請留意39節之後的內容.

39有幾個文士說:「夫子!你說得好.」 40以後他們不敢再問他甚麼.

當這班撒都該人前來挑戰耶穌的時候,有另一班人,有幾個文士,就是那些熟悉聖經的人,他們就評價耶穌說:「老師講得好.」

一個評價之後,大家又不敢出聲,一個專家,一個聖經專家走出來,大家就不出聲.

這樣,是不是因為這個評價(好評)就令到事情平息了?

耶穌是不是真的認為,這群文士真好,讚賞我,支持我?

45眾百姓聽的時候,耶穌對門徒說: 46「你們要防備文士.他們好穿長衣遊行,喜愛人在街市上問他們安,又喜愛會堂裡的高位,筵席上的首座； 47他們侵吞寡婦的家產,假意作很長的禱告.這些人要受更重的刑罰!」

不,耶穌轉頭馬上向門徒教導.

當撒都該人緊張一些邏輯,討論到死後有沒有復活的事?

而文士就去同他回應:耶穌講得好.

耶穌真正所關注的是什麼?

一方面,耶穌告訴我們,我們的主是活人的主,所以我們不擔心,一定會有復活的事情.

但同樣重要,上帝是活人的上帝,祂更關注我們這些活人的需要.

所以當文士以為自己學識高,信仰純正,以為自己宗教程度高的時候.

但如果他們扮得好好在人面前去請安,坐高位,就算坐筵席的首座,結果他們行為卻去侵吞寡婦財產,做好多好長的假意禱告,是必要受重刑的.

活人的上帝的意思,不是在分析與討論,什麼只理解為我們死後仍然可以同活的上帝相處(上天堂),而是我們今日活著時好好按上帝的心意而活(活出永恆豐富的生命),就才是上帝最關注的事情.

信仰需要思考,同時更需要行動.

好多時候更大的問題就是我們今日雖然是活著,但我們同死了沒有分別,我們仍然是有氣息,但我們生命是死的,沒有生活,沒有愛,眼中沒有他人,同死的人是沒有分別的,如是這樣,復活與你也無關.



\newpage

\section{甚麼時候有這些事呢}
\label{sec:773}
\textbf{關瑞文教授 - 甚麼時候有這些事呢?}
\newline
\newline
\hyperref[sec:772]{< < < PREV SERMON < < <}
~
\hyperref[sec:toc]{[返目錄]}
~
\hyperref[sec:774]{> > > NEXT SERMON > > >}


路加福音 21:5-19
\newline
\begin{longtable}{cl}
\hline
\hline
章節 & 經文 (和合本修訂版)\\
\hline
21:5 & \begin{tabularx}{0.7\textwidth}{X} 有人談論聖殿是用美石和供物裝飾的，耶穌就說： \end{tabularx} \\ \\ \relax
21:6 & \begin{tabularx}{0.7\textwidth}{X} 「你們所看見的這一切，日子將到，沒有一塊石頭會留在另一塊石頭上而不被拆毀的。」 \end{tabularx} \\ \\ \relax
21:7 & \begin{tabularx}{0.7\textwidth}{X} 他們問他：「老師，甚麼時候有這些事呢？這些事將臨到的時候有甚麼預兆呢？」 \end{tabularx} \\ \\ \relax
21:8 & \begin{tabularx}{0.7\textwidth}{X} 耶穌說：「你們要謹慎，不要受迷惑，因為將有好些人冒我的名來，說『我是基督』，又說『時候近了』，你們不要跟從他們！ \end{tabularx} \\ \\ \relax
21:9 & \begin{tabularx}{0.7\textwidth}{X} 當你們聽見打仗和動亂的事，不要驚惶；因為這些事必須先發生，但終結不會立刻就到。」 \end{tabularx} \\ \\ \relax
21:10 & \begin{tabularx}{0.7\textwidth}{X} 於是耶穌對他們說：「民要攻打民，國要攻打國， \end{tabularx} \\ \\ \relax
21:11 & \begin{tabularx}{0.7\textwidth}{X} 將有大地震，多處必有饑荒、瘟疫，又有可怕的異象和大神蹟從天上顯現。 \end{tabularx} \\ \\ \relax
21:12 & \begin{tabularx}{0.7\textwidth}{X} 但這一切的事以前，有人要下手拿你們，迫害你們，把你們交給會堂，並且關在監裡，又為我名的緣故拉你們到君王和統治者面前。 \end{tabularx} \\ \\ \relax
21:13 & \begin{tabularx}{0.7\textwidth}{X} 但這些事終必成為你們作見證的機會。 \end{tabularx} \\ \\ \relax
21:14 & \begin{tabularx}{0.7\textwidth}{X} 所以，你們要立定心意，不要預先考慮怎樣申辯； \end{tabularx} \\ \\ \relax
21:15 & \begin{tabularx}{0.7\textwidth}{X} 因為我必賜你們口才和智慧，是你們一切敵人所敵不住、駁不倒的。 \end{tabularx} \\ \\ \relax
21:16 & \begin{tabularx}{0.7\textwidth}{X} 連你們的父母、兄弟、親族、朋友也要把你們交給官府；你們中間也將有被他們害死的。 \end{tabularx} \\ \\ \relax
21:17 & \begin{tabularx}{0.7\textwidth}{X} 你們要為我的名被眾人憎恨。 \end{tabularx} \\ \\ \relax
21:18 & \begin{tabularx}{0.7\textwidth}{X} 然而，你們連一根頭髮也不會損失。 \end{tabularx} \\ \\ \relax
21:19 & \begin{tabularx}{0.7\textwidth}{X} 你們憑著堅忍，就必保全性命。」 \end{tabularx} \\ \\
[1ex]
\hline
\hline
\end{longtable}
講題:甚麼時候有這些事呢?When Will This Be?

經課:路加福音21:5-19

講員:關瑞文教授

場合:香港中文大學崇基學院禮拜堂主日崇拜

日期:2025年11月16日

引言:當我們所倚仗的聖殿崩塌時

各位弟兄姊妹,主內平安.

今天在這個優美,莊嚴而充滿歷史感的禮拜堂,與各位一同崇拜,我心裡實在充滿感恩.

我們身處的這座崇基禮拜堂,它矗立在這裡超過半個世紀,安靜地承載著一代又一代人的禱告,頌讚,婚禮和安息禮.它就像一座穩固的聖殿,是我們心靈的座標,是這所大學的精神象徵.

我想,兩千年前,當耶穌的門徒走在耶路撒冷,他們的心情大概也和我們此刻相似.路加福音21章第5節記載:「有人談論聖殿是用美石和供物妝飾的.」他們看著那座宏偉壯麗,金碧輝煌的聖殿,心中充滿了宗教的虔誠.這座聖殿,是是猶太人信仰的中心,是上帝臨在的記號,是他們眼中永恆不變的象徵.他們相信,只要聖殿還在,上帝的約就在,就穩固.

然而,就在門徒們滿心讚嘆之際,耶穌卻說出了一句石破天驚的話.祂說:「論到你們所看見的這一切,將來日子到了,在這裡沒有一塊石頭留在石頭上,不被拆毀了.」(路21:6)

試想像一下當時的震撼.這不只是一句建築學上的預言,這是在宣告一個時代的終結.這等於是說,他們信仰的根基,身份的認同,都將會灰飛煙滅.耶穌的話,直接衝擊了他們心中最穩固,最賴以安身立命的「聖殿」.

弟兄姊妹,今天我們或許沒有一座實體的聖殿可以倚靠.但我們每一個人的生命中,其實都有著各式各樣,大小不同的「聖殿」——那些我們視為理所當然,賴以為生的穩固基石.可能是我們健康的身體,一份安穩的事業,一段美滿的關係,我們的學術成就,甚至是我們所熟悉的香港這個家.我們慣性地將盼望和安全感,建立在這些「聖殿」之上.

但耶穌的這番話,像暮鼓晨鐘,跨越兩千年,今天依然在對我們說話.祂提醒我們,所有眼所能見的「聖殿」,無論多麼華美,多麼堅固,都有崩塌的一天.那麼,當我們生命中的聖殿開始動搖,當我們所倚仗的不再穩固,當「沒有一塊石頭留在石頭上」的混亂臨到我們時,我們該如何自處?是恐懼,是絕望,還是有另一種可能?

今天,我想邀請大家一同回到路加的文本中,去聆聽耶穌基督在混亂和終末的預言中,賜給門徒,也賜給今天我們的,那個充滿盼望與使命的核心信息.

(一)從末日的現象,到歷史的現場

當門徒聽到聖殿將被拆毀,他們的第一個反應,自然是問:「夫子!甚麼時候有這事呢?這事將到的時候有甚麼預兆呢?」(路21:7)這是一個非常人性的問題.面對不可知的未來,我們總想尋找一些確據,一些可以掌控的線索,好讓我們能有所準備.

而耶穌接下來的回答,就是我們所熟悉的「小啟示錄」(Little Apocalypse).在符類福音的馬太,馬可福音中,都有類似的記載.經文充滿了戰爭,饑荒,地震,天象變化等令人不安的末日景象.長久以來,很多人讀這段經文,都把它當成一份解讀世界末日的密碼圖,熱衷於把經文中的預兆與國際新聞對號入座,試圖計算出主再來的準確時間.

然而,當我們仔細研讀路加的記載,並對比馬可福音的平行經文時,我們會發現路加福音的作者,似乎刻意在引導他的讀者,將焦點從「猜測世界何時終結」,轉移到「如何在身處的遭遇中活出信仰」.

第一,路加福音的寫作日期,學界普遍認為是在主後70年耶路撒冷被羅馬軍隊摧毀之後.他的讀者——很可能是一個以外邦基督徒為主體的教會——對於這場慘烈的戰爭記憶猶新.路加直接指出,耶穌所預言的,並不是某個遙遠未來的神秘災難,而是已經發生在他們身邊,具體的歷史事件.他彷彿在對他的讀者說:你們所經歷的這場浩劫,不是上帝失控了,這一切,早已在主的預言之中.

第二,路加回憶,『10當時,耶穌對他們說:「民要攻打民,國要攻打國；11地要大大震動,多處必有饑荒,瘟疫,又有可怕的異象和大神蹟從天上顯現.12但這一切的事以先,人要下手拿住你們,逼迫你們,把你們交給會堂,並且收在監裡,又為我的名拉你們到君王諸侯面前.』路加要告訴他的教會,他們所面對的逼迫和苦難,並不是世界末日來臨前的最後一場大災難,而是他們此時此刻,每日每夜都要面對的「日常」.苦難不是將來式,而是現在進行式.

路加的編排,把讀者的眼光從天上拉回到地上,從未來拉回到現在.他要處理的,不是一個關於末世時間表的神學問題,而是一個極其迫切的牧養關懷:當人正活在一個不穩定,憂慮,煎熬的日子時,應該如何理解自己的處境?應該如何活下去?

(二)轉危為機:苦難成為見證的契機

這就引導我們進入整段經文的核心,——第13節.耶穌在預告了種種煎熬之後,祂沒有說:「你們要忍耐,因為天堂的獎賞是大的.」祂也沒有說:「你們要趕快逃跑.」祂說了一句顛覆我們常理的話(《和修本》):「13但這些事終必成為你們作見證的機會.」

在路加的寫作中,「作見證」是一個核心主題,貫穿了福音書和使徒行傳.見證,不單是用口宣告福音,更是用整個生命去活出所信的道.耶穌在這裡提出一個革命性的觀點:苦難,不是上帝的懲罰,也不是上帝失蹤的記號.相反,它是一個獨特的「機會」,一個讓門徒的生命,能夠成為福音最有力的「見證」的場景.

弟兄姊妹,這是一個何等巨大的範式轉移!我們通常如何看待苦難?我們視之為需要避開的麻煩,需要解決的問題,需要挪去的重擔.我們禱告,是求主叫我們「免於」苦難.但耶穌卻說,苦難本身,可以「成為」一個機會.

那麼,如何「成為」?耶穌接著提供了兩個應許和一個命令.

第一,是「智慧言語的應許」.祂說:「因為我必賜你們口才,智慧,是你們一切敵人所敵不住,駁不倒的.」(路21:15)祂吩咐門徒,到時候「不要預先思想說什麼」.這不是叫我們懶惰,放棄思考.這是要我們在最關鍵的時刻,學習全然的倚靠.我們在使徒行傳中,看到司提反,看到彼得,看到保羅,他們在遭逼迫時,都印證了這個應許.

第二,是「終極保守的應許」.耶穌說:「然而,你們連一根頭髮也必不損壞.」(路21:18)這句話聽起來非常矛盾.因為就在兩節之前,祂才剛剛說過:「連你們的父母,弟兄,親族,朋友也要把你們交官；你們也有被他們害死的.」(路21:16)既然會被害死,又怎能說連一根頭髮也不損壞呢?

這正是福音的奧妙所在.耶穌所應許的,不是肉身性命的安然無恙,而是我們整個存有(our whole being)在上帝裡的終極安全.人世間的厄運,最多只能殺害我們的身體,卻無法觸及那在基督裡的新生命.我們的真實生命,我們的永恆價值,是安全地存放在上帝的手中.從永恆的視角看,沒有任何事物,能叫我們與神的愛隔絕.這是一個關乎「得著生命」的應許.

在這兩個應許之上,耶穌給出了一個命令:「你們常存忍耐,就必保全靈魂.」(路21:19)「忍耐」這個詞,希臘文是 ὑπομονή (hypomonē),它不是指一種消極被動的忍受,而是一種積極,主動,昂然站立的「恆忍」.它描畫的,是一幅士兵在戰場上堅守陣地的圖畫.無論周圍的炮火多麼猛烈,他依然屹立不倒.

所以,耶穌的教導是:在人生跌入混亂的時刻,門徒的使命不是去猜測終局,而是要在自己的崗位上「站穩」.當人際關係因信仰而破裂,當世界充滿問題,當未來充滿不確定性時,基督徒的見證,就在於那份不輕易動搖的「忍耐」.這份忍耐本身,就是一篇最有力的道.它無聲地宣告:我們所信靠的,比眼前的風暴更大.

(三)我的頸椎

弟兄姊妹,以上的經文解釋,可能聽起來還是有些遙遠.我們今天的世界雖然真的有「民要攻打民,國要攻打國 … 大地震… 饑荒 … 瘟疫」(v10), 但身處香港的我們,大概都是安全.在座的,大概也不會如路加教會一樣 (v12) 會遭遇到「人要下手拿住你們,逼迫你們,把你們交給會堂,並且收在監裡,又 … 拉你們到君王諸侯面前.」

不過,我想和大家分享一段我個人非常切身的經歷,來闡釋這段經文如何在我生命中,成為真實而震撼的提醒.

其實我的身體一直都有些「舊患」.從很多年前開始,我就有腰椎間盤突出的問題.而頸椎的不適,也斷斷續續地困擾著我,引起頸部疼痛,僵硬.但多年來我已經習慣了,所以一直不以為意.

直到今年三月,我的二姐大概是看不過眼,話出錢請我去做磁力共振掃描(MRI),催促我去取個最新的診斷.盛情難卻之下,我便去了.而掃描的結果,就是今天早上我帶來的這張影像圖片.

這是我頸椎的側面圖.醫生指給我看,我的問題相當複雜.首先,我有好幾節頸椎的椎間盤都出現了突出.但更關鍵的是,我的脊髓管道內的「腦脊髓液」(Cerebrospinal Fluid)已經大量流失.這腦脊髓液其中一個極其重要的功能,就是作為天然的緩衝墊(cushion),無時無刻保護著我們脆弱的脊髓神經,不怕一般的撞擊.現在,因為失去了這層保護液,加上突出的椎間盤的擠壓,從影像中可以看到,我的中樞神經已經被擠壓得非常要命.

接著,他指著我的影像,用一種非常平靜但字字千鈞的語氣說:「你現在的脊髓神經,就處於一個極度脆弱,不受保護的狀態.我的情況,就好像頸項上懸著一把刀.任何一次意外,比如在巴士上一個急剎車,或是不小心摔倒,只要頸部受到不當的撞擊,都可能導致你的脊髓神經中樞受到永久性損傷,結果就是——下半身癱瘓.」

醫生還怕我不明白,就引用了一個近年的例子.他說:「我這樣說,你可能更明白.你記得Mirror演唱會中那位受重傷的舞蹈員阿Mo嗎?」

那一刻,我整個人都驚愕了(shocked).這個潛在的風險,是我完全始料未及的.我原以為不過是普通的退化問題,最多是痛一點,麻痺一點,從沒想過會和「癱瘓」這個詞扯上關係.雖然醫生強調這只是一個「潛在風險」——意思是,如果沒有意外發生,我的活動能力將一如往常.但這個「如果」,像一把無形的刀,懸在了我的頸上.為了確認,我把百多張的MRI片子拿給了幾位非常資深,也是我好友的物理治療師和醫生的朋友看.他們都以非常嚴肅的態度,證實了這個風險的真實性和嚴重性.

弟兄姊妹,那一刻,我感覺自己生命中那座精心搭建的「聖殿」,突然間劇烈地動搖起來.我本該在去年就退休了.但為了我所事奉的崇基學院神學院,因為正面臨著教授團隊青黃不接的時期,我申請了並獲大學批准了延任兩年,明年才正式退下來.我心裡早就為自己規劃好了一幅美好的退休藍圖,那就是要享受一個「自由自在的人生」,去旅遊,更多時間在教會事奉,去做自己喜歡的研究.這座「退休聖殿」,是我近年來對未來美好的期盼.

但是,這個診斷,無情地在這座聖殿上鑿開了一道裂縫.一個微小的意外,就可能使我所有「自由自在」的計劃,變成在輪椅上渡過餘生.我心中充滿了疑惑:我所倚仗的健康身體,原來如此脆弱.我所期盼的美麗退休人生,原來可以不堪一擊.

就在那段充滿憂慮,反覆消化這個消息的日子裡,我重新回到路加福音21章的教導:第13節,「但這些事終必成為你們作見證的機會」,這句話,像一道銳利的光,刺穿了我疑惑的黑霧.

我開始問自己:我這個脆弱的頸椎,我這個被由風險籠罩的「下半場人生」,如何能成為一個「作見證的機會」?

在禱告和反思中,我的心境慢慢起了變化.我做了幾件事.首先,我坦誠地告訴了我的家人,讓他們對我的狀況有心理準備.這不是在散播恐懼,而是在學習一種信靠中的坦誠.

接著,我開始嘗試用聖經的教導,來重構我對這件事的反應.

第一,這不是要問「老師,甚麼時候有這些事呢?這些事將臨到的時候有甚麼預兆呢?」,而是要問,「老師,遭遇這件事,你要我作什麼?」

第二,這是我作一個關於「倚靠」的見證.這教導提醒我,我不能靠自己的憂慮去掌控未來.我過去或許倚仗自己的計劃和健康的身體.主透過這個脆弱的頸椎,邀請我學習全然的交託.我的見證,不再是向人展示一個完美的「自由自在」的退休計劃,而是向人展示,即使在無法掌控的風險中,我依然可以倚靠那位掌管生命的主,活好每一天.我已經在不同場合與不同教會的弟兄姊妹直接分享過這個體會.

第三,這是我作一個關於「真生命」的見證.耶穌應許「你們連一根頭髮也必不損壞」,又命令我們要憑忍耐「得著生命」.這個矛盾的應許,在此刻對我尤其真實.我的肉身生命,隨時可能因意外而「損壞」.但我真正的,那在基督裡的永恆生命,是世界任何意外都不能奪去的.這個危機,迫使我更深刻地去分辨:我所追求的,究竟是「自由自在」的肉體生活,還是那份無論在何種境況中,都能在主裡站立得穩的「靈性」?我的見證,就是要在這份張力中,活出對「一切都在主裡的生命」更深的渴求與確信.

最後,這是我作一個關於「恆忍」(hypomonē)的見證.恆忍,不是消極地等待厄運降臨,而是積極地,堅定地在自己的崗位上站穩.我們在負面的遭遇中,不會被不確定的未來所定義,而是被呼召與使命所引導.

我的頸椎,依然有它的問題.那把懸在頸上的刀,也依然存在.但我感謝主,祂透過這座搖搖欲墜的「聖殿」,讓我看見了信仰最真實的風景.原來,那些生命中崩塌的聖殿,真的可以成為我們為祂作見證的基石.

結論:在崩壞的世界中,屹立作見證

弟兄姊妹,我們人生在世,又何嘗不是可以一霎那間發現自己生命中的「聖殿」原來在劇烈搖動?可能是學業的壓力,工作的瓶頸,家庭的危機,關係的破裂,又或者,就像我一樣,是身體突然響起的警號.

在這樣一個時代,耶穌基督的話語,給了我們一條出路.這條路,不是逃避現實,也不是坐以待斃.這條路,是「站穩」的路,是「忍耐」的路,是「見證」的路.

祂沒有應許我們天色常藍,花香常漫.但祂卻給了我們一個更偉大的應許:在一切的混亂中,祂與我們同在；在一切的遭遇中,祂要將之轉化為我們生命最榮耀的見證.

所以,親愛的弟兄姊妹,今天你生命中正在動搖的「聖殿」是甚麼?是甚麼讓你感到疑惑和不安?

讓我們一起,把目光從崩塌的石頭上,轉向那位預言這一切,也掌管這一切的基督.讓我們不再問「甚麼時候」,「有甚麼預兆」,而是問:「主啊,在此時此刻,在這個處境中,我如何為你作見證?」

或許,你的見證,就是在充滿戾氣的辦公室裡,活出溫柔與和平.

或許,你的見證,就是在學術的激烈競爭中,持守誠實與謙卑.

或許,你的見證,就是在家庭的衝突撕裂中,選擇饒恕與復和.

或許,你的見證,就是在病痛的折磨和對未來的憂慮中,展現出那份從主而來的,不動搖的盼望和忍耐.

我們每一個人的生命,都可以成為一篇道.我們在軟弱中的忍耐,我們在絕望中的信靠,將會成為這個崩壞世界中最有力的見證,告訴世人:有一位上帝,祂的愛,比死亡更堅強；祂的國度,永不倒塌.

願我們都能靠著主的恩典,在各自的風暴中,屹立不倒,活出這份獨一無二的見證.

阿們.



\newpage

\section{我們的上帝要施行救贖}
\label{sec:774}
\textbf{高國雄牧師 - 我們的上帝要施行救贖}
\newline
\newline
\hyperref[sec:773]{< < < PREV SERMON < < <}
~
\hyperref[sec:toc]{[返目錄]}
~
\hyperref[sec:775]{> > > NEXT SERMON > > >}


路加福音 23:33-43
\newline
\begin{longtable}{cl}
\hline
\hline
章節 & 經文 (和合本修訂版)\\
\hline
23:33 & \begin{tabularx}{0.7\textwidth}{X} 到了一個地方，名叫髑髏地，他們就在那裡把耶穌釘在十字架上，又釘了兩個犯人：一個在右邊，一個在左邊。 \end{tabularx} \\ \\ \relax
23:34 & \begin{tabularx}{0.7\textwidth}{X} 〔 這時，耶穌說：「父啊！赦免他們，因為他們所做的，他們不知道。」〕士兵就抽籤分他的衣服。 \end{tabularx} \\ \\ \relax
23:35 & \begin{tabularx}{0.7\textwidth}{X} 百姓站在那裡觀看。官長也嘲笑他，說：「他救了別人，他若是基督，是神所揀選的，救救他自己吧！」 \end{tabularx} \\ \\ \relax
23:36 & \begin{tabularx}{0.7\textwidth}{X} 士兵也戲弄他，上前拿醋送給他喝， \end{tabularx} \\ \\ \relax
23:37 & \begin{tabularx}{0.7\textwidth}{X} 說：「你若是猶太人的王，救救你自己吧！」 \end{tabularx} \\ \\ \relax
23:38 & \begin{tabularx}{0.7\textwidth}{X} 在耶穌上方有一個牌子寫著：「這是猶太人的王。」 \end{tabularx} \\ \\ \relax
23:39 & \begin{tabularx}{0.7\textwidth}{X} 同釘的犯人中有一個譏笑他，說：「你不是基督嗎？救救你自己和我們吧！」 \end{tabularx} \\ \\ \relax
23:40 & \begin{tabularx}{0.7\textwidth}{X} 另一個就應聲責備他，說：「你是一樣受刑的，還不怕神嗎？ \end{tabularx} \\ \\ \relax
23:41 & \begin{tabularx}{0.7\textwidth}{X} 我們是應得的，因為我們是自作自受，但這個人沒有做過一件不對的事。」 \end{tabularx} \\ \\ \relax
23:42 & \begin{tabularx}{0.7\textwidth}{X} 他對耶穌說：「耶穌啊，你進入你國的時候，求你記念我。」 \end{tabularx} \\ \\ \relax
23:43 & \begin{tabularx}{0.7\textwidth}{X} 耶穌對他說：「我實在告訴你，今日你要同我在樂園裡了。」 \end{tabularx} \\ \\
[1ex]
\hline
\hline
\end{longtable}
講題:我們的上帝要施行救贖 Our God is a God who Saves

經課:路加福音23:33-43

講員:高國雄牧師

場合:香港中文大學崇基學院禮拜堂主日崇拜

日期:2025年11月23日

我們的上帝要施行拯救

耶利米書 23:5-6

5 耶和華說:「日子將到,我要給大衛興起一個公義的苗裔；他必掌王權,行事有智慧,在地上施行公平和公義. 6 在他的日子, 猶大必得救, 以色列也安然居住.他的名必稱為『耶和華－－我們的義』.

拯救:

興起公義的苗裔,行事有智慧,施行公平和公義.

路加福音 1:71-77

71 拯救我們脫離仇敵和一切恨我們之人的手,

74 叫我們既從仇敵手中被救出來,

75 就可以終身在他面前,坦然無懼地用聖潔,公義事奉他.

77 叫他的百姓因罪得赦,就知道救恩.

拯救:

脫離仇敵之手,坦然無懼地用聖潔,公義事奉主,罪得赦.

歌羅西書 1:11-20

13 他救了我們脫離黑暗的權勢,把我們遷到他愛子的國裡； 14 我們在愛子裡得蒙救贖,罪過得以赦免.

20 既然藉著他在十字架上所流的血成就了和平,便藉著他叫萬有－ －無論是地上的,天上的－ －都與自己和好了.

拯救:

脫離黑暗的權勢,成就了和平,與自己和好.

路加福音23:33-43

33到了一個地方,名叫「髑髏地」,就在那裡把耶穌釘在十字架上,又釘了兩個犯人:一個在左邊,一個在右邊. 34當下耶穌說:「父啊!赦免他們；因為他們所做的,他們不曉得.」兵丁就拈鬮分他的衣服. 35百姓站在那裡觀看.官府也嗤笑他,說:「他救了別人；他若是基督,上帝所揀選的,可以救自己吧!」 36兵丁也戲弄他,上前拿醋送給他喝, 37說:「你若是猶太人的王,可以救自己吧!」 38在耶穌以上有一個牌子寫著:「這是猶太人的王.」

39那同釘的兩個犯人有一個譏笑他,說:「你不是基督嗎?可以救自己和我們吧!」 40那一個就應聲責備他,說:「你既是一樣受刑的,還不怕上帝嗎? 41我們是應該的,因我們所受的與我們所做的相稱,但這個人沒有做過一件不好的事.」 42就說:「耶穌啊,你得國降臨的時候,求你記念我!」 43耶穌對他說:「我實在告訴你,今日你要同我在樂園裡了.」

向「拯救」回應:

明白上主的愛,承認自己的罪,向主求憐憫,得著盼望!



\newpage

\section{榮耀的王將要進來}
\label{sec:775}
\textbf{郭志榮牧師 - 榮耀的王將要進來}
\newline
\newline
\hyperref[sec:774]{< < < PREV SERMON < < <}
~
\hyperref[sec:toc]{[返目錄]}
~
\hyperref[sec:776]{> > > NEXT SERMON > > >}


詩篇 24:1-10
\newline
\begin{longtable}{cl}
\hline
\hline
章節 & 經文 (和合本修訂版)\\
\hline
24:1 & \begin{tabularx}{0.7\textwidth}{X} 地和其中所充滿的，世界和住在其中的，都屬耶和華。 \end{tabularx} \\ \\ \relax
24:2 & \begin{tabularx}{0.7\textwidth}{X} 他把地建立在海上，安定在江河之上。 \end{tabularx} \\ \\ \relax
24:3 & \begin{tabularx}{0.7\textwidth}{X} 誰能登耶和華的山？誰能站在他的聖所？ \end{tabularx} \\ \\ \relax
24:4 & \begin{tabularx}{0.7\textwidth}{X} 就是手潔心清，意念不向虛妄，起誓不懷詭詐的人。 \end{tabularx} \\ \\ \relax
24:5 & \begin{tabularx}{0.7\textwidth}{X} 他必蒙耶和華賜福，又蒙救他的神使他成義。 \end{tabularx} \\ \\ \relax
24:6 & \begin{tabularx}{0.7\textwidth}{X} 這是尋求耶和華的族類，是尋求你面的雅各。（細拉） \end{tabularx} \\ \\ \relax
24:7 & \begin{tabularx}{0.7\textwidth}{X} 眾城門哪，要抬起頭來！永久的門戶啊，你們要被舉起！榮耀的王將要進來！ \end{tabularx} \\ \\ \relax
24:8 & \begin{tabularx}{0.7\textwidth}{X} 這榮耀的王是誰呢？就是有力有能的耶和華，在戰場上大有能力的耶和華！ \end{tabularx} \\ \\ \relax
24:9 & \begin{tabularx}{0.7\textwidth}{X} 眾城門哪，要抬起頭來！永久的門戶啊，你們要高舉！榮耀的王將要進來！ \end{tabularx} \\ \\ \relax
24:10 & \begin{tabularx}{0.7\textwidth}{X} 這榮耀的王是誰呢？萬軍之耶和華是榮耀的王！（細拉） \end{tabularx} \\ \\
[1ex]
\hline
\hline
\end{longtable}
講題:榮耀的王將要進來 The King of Glory May Come In

經課:詩篇24:1-10

講員:郭志榮牧師

場合:香港中文大學崇基學院禮拜堂主日崇拜

日期:2025年11月30日

⼀．經文:詩廿四

1 地和其中所充滿的,世界和住在其間的,都屬耶和華.

2 他把地建立在海上,安定在⼤⽔之上.

3 誰能登耶和華的⼭?誰能站在他的聖所?

4 就是⼿潔⼼清,不向虛妄,起誓不懷詭詐的⼈.

5 他必蒙耶和華賜福,⼜蒙救他的神使他稱義.

6 這是尋求耶和華的族類,是尋求你⾯的雅各.

7 眾城⾨哪,你們要抬起頭來!永久的⾨⼾,你們要被舉起!那榮耀的王將要進來!

8 榮耀的王是誰呢?就是有⼒有能的耶和華,在戰場上有能的耶和華!

9 眾城⾨哪,你們要抬起頭來!永久的⾨⼾,你們要把頭抬起!那榮耀的王將要進 來!

10 榮耀的王是誰呢?萬軍之耶和華,他是榮耀的王!(細拉)

⼆．前⾔:解讀《詩篇》廿四篇的問題

弟兄姊妹,主內平安!

今天跟⼤家分享的經文來自《詩篇》廿四篇.經文已經在電幕上顯⽰,讓我們⼀同誦讀.

這是⼀篇⼤家很熟悉的詩篇,⼤家很可能在讚美祈禱會或主⽇崇拜中曾經聽過或誦讀 過,特別是3-4節和7-10節.這是⼀篇引起不少討論的詩篇.

曾經有⼀個時期的學者,就認為這是⼀篇在特定節慶或典禮中使⽤的詩篇.他們當中 有部分⼈主張,這是⼤衛把約櫃接回耶路撒冷時頌唱.亦有⼈推斷,這詩篇是寫給到 錫安聖殿的朝聖者在進殿的禮儀中使⽤.但是否真的是這樣嗎?

剛才誦讀這篇詩篇時,⼤家應該不難發現這篇詩篇⼤致可分為三個部分,分 別是:1-2節,3-6節和7-10節.⽽每⼀個部分表⾯上都好像有各自的主題和文 體風格, 所以要斷定這篇詩篇的背景,⽤途和整篇詩篇統⼀的主題並不容易.那麼, 這篇詩篇到底要表達甚麼信仰的真理呢?

三．釋經

A.《詩篇》廿四篇三個部分的釋義

先把這篇詩篇的㆔個部分逐㆒解釋.

1)「㆞和其㆗所充滿的,世界和住在其間的,都屬耶和華.他把㆞建立在海㆖,安定在大水之㆖.」(1-2節)

《詩篇》廿四篇第 1 節提到,神對「地和其中所充滿的」以及「世界和住在其間的」擁有主權,當然亦包括對這⼀切的管治權.⽽第 2 節則提供了背後的理由,正是神在洶湧的⼤海和急流中使「地」的根基建立起來.

最近看到⼀段新聞報導,提及地球核⼼的轉速和轉向正在改變,⽽科學家指出這些改變或會影響到⼀直保護著地球,擋開⼤部分宇宙輔射的地球磁場. ⼈類和各類的⽣物能在地球上⽣存,其實需要很多條件配合才能發⽣,⽽這當然是神精密的設計.

古代近東的⼈⺠當然沒有這些科學知識,不知道地殼是浮在溶岩上⽽不是在⽔⾯上.當詩⼈在第 2 節中提到,「地」是建立在海和⼤⽔之上時,實際上是表達詩⼈也體會到「地」的脆弱.所以,若不是神把「地」的根基立定,⼀個適合⼈居住的「地」根本不會存在.若不是神持續「拱托」和「維護」著,這「地」甚⾄可能會被「海」和「⼤⽔」所代表的混沌⼒量吞噬⽽完全消失.

第 1 節中的「地」和「世界」是⼀對平⾏的近義詞.但兩者有甚麼分別呢?有猶太拉比解釋這兩者的分別在於:「地」是指⼀切的地⼟,⽽「世界」則是指有⼈聚居的地⽅.所以,「世界」包括在「地」之中,是「地」的⼀部分.

既然如此,原則上有前半節的詩句便⾜夠,因為下半節詩句的內容已包含在上半節經文裡.但詩⼈這重覆的做法並非多此⼀舉,⽽是有刻意加強語氣和指出重點的作⽤. 換⾔之,詩⼈要特別強調和提醒⼤家,「世界」這個「⼈聚居和⽣活的領域」是屬於神的.

⼤家⼀定會問:為何詩⼈要這樣刻意強調?從逆向的⾓度看,詩⼈要如此刻意強調,就正好說明,儘管⼤家知道「地」和其中的⼀切都屬於神,在神的管轄之下,但許多⼈的⼼底深處並不真正願意接受,神擁有「世界」這個「⼈聚居和⽣活的領域」和「住在其間的自己」的主權.在骨子裡,人總是希望自己「話事」,按自己的喜好, 運用自己有限的智慧,去選擇自己要怎樣生活.

這些人骨子裡的潛台詞是:「世界」以外的大地,才是交給神掌管的.他們會認為神有責任要維持大地的運作井然有序,使大地綠草如茵,空氣清新,天氣晴朗,沒有天災地震,適合人居住,否則就會埋怨神,質疑神為何讓災禍降臨.然而, 當涉及「人的領域」時,這些人就不想讓神參與和過問,希望由自己來主宰.

結果如何?大家看看這個「世界」便知道!十九世紀末,二十世紀初,因為科學的發展,人對自己的理性,對世界的發展非常樂觀,到了一個程度,甚至認為人類已經足夠成熟,不再需要神,都能夠倚靠人的能力在地上建立「天堂」般的世界.接著,就是兩次「世界大戰」的爆發.

大家環顧身邊的人的表現,甚至認真省察一下自己的內心,就會明白人很容易有這樣的毛病.詩人強調「世界和住在其間的」都屬於神,正是要指出人這錯誤的傾向.

為何這篇詩篇要有這兩節經文作為開場白呢?「世界」所指的,不只是以色列或猶大國的疆土,而是指「地」上一切有人聚居棲息的地方,而「住在其間的」當然指一切的民族,而並非只是指以色列或猶大國的人民.

因此,在第三部分提到「要進來」的「榮耀的王」,也是列國萬民的君王.祂統治的 國度,超越了以色列甚至任何一個地上霸權的範圍.

2)「誰能登耶和華的山?誰能站在他的聖所?就是手潔心清,不向虛妄,起誓不懷詭詐的人.他必蒙耶和華賜福,又蒙救他的神使他稱義.這是尋求耶和華的族類, 是尋求你面的雅各.」(3-6節)

「登神的山」和「站在神的聖所」,代表人要親近神,進入神的神聖靈域之中.這篇詩篇的第二部分,就是要回應「誰能」亦即「誰有資格」的問題.

這段經文看似是一段在進殿的禮儀中朝聖者與守在聖殿入口的祭司之間對答的禮文. 然而,詩人在這裡列出的四個條件,包括「手潔」'「心清」'「不向虛妄」和「起 誓不懷詭詐」(4節),都是倫理道德和信仰上的要求,而不是指某些宗教儀式的規定.

「手潔」指一個人「外在的行為」清白無罪,而「心清」就是指這人「內心的思想」 純潔,沒有邪惡的念頭,沒有不純良的動機和謀算.

「不向虛妄」就是指不把生命「抬起」交給那「虛妄的」,代表這人敬拜神而不是那沒有實質的偶像,亦即與神有正確的關係.「起誓不懷詭詐」代表這是一個誠實的人, 不會藉發假誓來欺騙別人,亦即與人建立起真誠'得到別人信任的關係.

這樣的⼈⾏為和思想正直,與神,與⼈都有合宜的關係,當然能來到神⾯前,親近神, 並因⽣命結連於神⽽得著⽣命的⼒量,特別是⼼靈的⼒量(mental strength).這樣的⼈無論⾯對怎樣的困難,壓⼒和挫敗,都能以正⾯的態度⾯對.跌倒了,能再站起來；失敗了,能總結經歷並再作不同的嘗試；受傷了,願意讓神為自⼰包紮和治療.這樣的⼈當然有更多機會經歷神的醫治與更新,當然得到神更多,更⼤的賜福.他們當然也會得到神的嘉許,被神宣判為義⼈,因為他⾏神眼中看為正確的事.

但問題是:有誰能完全做到這四點的要求?假若要徹底「⼿潔⼼清」,只單單專注仰 望神⽽⼀點也「不向虛妄」,還要⼝裡半點「詭詐」也沒有,⼀個⼈才能親近神,得 蒙神的賜福.這樣,⼈還有盼望嗎?⼈⼀⽅⾯因為罪,無法拉近與神的距離,但另⼀ ⽅⾯,⼈遠離了神,就沒有離開罪惡,⽣命更新的動⼒.

這是⼀個兩難的困局!

「尋求你⾯的雅各」可譯作「尋求雅各的神的⾯」.「這是尋求耶和華的族類」和 「尋求雅各的神的⾯」,正是⼀群尋求神的⼈,在看似無望的,兩難的困局之中,向神所發出的呼喊.

3)「眾城門哪,你們要抬起頭來!永久的門戶,你們要被舉起!那榮耀的王將要進來!榮耀的王是誰呢?就是有力有能的耶和華,在戰場㆖有能的耶和華!眾城門哪,你們要抬起頭來!永久的門戶,你們要把頭抬起!那榮耀的王將要進來!榮耀的王是誰呢?萬軍之耶和華,他是榮耀的王!」(7-10節)

上⼀段的經文提到,尋求神的⼈要符合甚麼條件,才能進入神的神聖靈域.⽽來到這段詩篇,詩⼈卻好像突然作出了⼀百八⼗度轉變,宣告「榮耀的王將要進來」.換⾔之,就是把焦點由⼈進入神的⼭和神的聖所,轉到神將要「進城」的盛況.

有學者推斷,這可能是⼤衛把約櫃接回耶路撒冷時,跟約櫃⼀起巡遊的隊伍與看守城⾨的⼈互相對唱的詩歌.亦有學者提出,這段經文也可能是以⾊列或猶⼤國軍隊帶著約櫃出戰並打勝仗後,在回歸時所頌唱的凱旋之歌.是否如此,下⾯會進⼀步討論.

但無論是哪⼀種可能性,這段詩篇所營造出來的氣氛都是熱烈,⾼漲和澎湃的,是⼠氣⾼昂的,並充滿了歡欣喜樂的⼼情.

「眾城⾨」在這裡並不是指城牆出入⼝的那扉「⾨」,⽽是指⼀座城出入⼝的範圍, 包括⾨外和⾨內的廣場,也就是城內和周邊⼀帶居⺠聚集的地⽅. 所以,向「眾城⾨」的宣告,也就是對著這些在城⾨⼝聚集的⼈發⾔.

詩⼈叫這些⼈要「抬起頭來」,⽽「抬起頭來」是⼀個近東流⾏的成語,⽤來呼籲⼤家要「歡欣」和「⼼懷盼望」, 也有激勵⼤家「⿎起勇氣」的含義.

以逆向的思維來思考,詩⼈要這樣呼籲,表⽰⼤家本來並不歡欣,並不喜樂.這反映他們很可能對所⾯對的處境和形勢並不樂觀,因此感到憂⼼忡忡甚⾄⼼中充滿恐懼. 詩⼈作出「抬起頭來」的呼籲,正是要激勵⼤家的⼠氣,不⽤意志消沉.

甚麼是詩⼈叫⼤家「抬起頭來」重新振作的根據呢?詩⼈解釋,這是因為「榮耀的王將要進來」,⽽詩⼈反覆強調這「榮耀的王」不是地上的君王或領袖,⽽是神自⼰!

在這個段落中,詩⼈⽤了「榮耀的王」,「耶和華」和「萬軍之耶和華」來稱呼神. 這三個名字代表甚麼意思呢?

「榮耀」最基本的意思,就是「有份量」和「有重要性」,所以「榮耀的王」強調神是⼀位真正有份量和有重要性的王者,表明祂的主權沒有⼈能輕看或忽略,祂的統治要帶來有重要性的影響.神以「榮耀的王」的⾝份「進來」,意味著神要親自帶領祂的⼦⺠,不會讓祂的⼦⺠因為沒有可跟隨的領袖⽽潰散.這位「榮耀的王者」要藉此從這班⼈當中建立起⼀個彰顯祂份量和重要性的國度.

接下來,再解釋「耶和華」這個名字.神差遣摩⻄帶領以⾊列出埃及時,⽤了「耶和華」這個名字.所以「耶和華」這個名字讓⼈聯想到神對祂⼦⺠的救贖之恩,也讓⼈聯想到神與⼈有⽣命之約的關係.這表明神的「進來」,就是要履⾏祂在⽣命約中的應許,施恩拯救屬祂的⼦⺠.

「萬軍之耶和華」這個名字⼀般在有關戰爭的處境中使⽤, 配合「有⼒有能」和「在戰場上有能」的描述,說明神的「進來」就是要帶領祂的⼦⺠,也是為祂的⼦⺠⽽爭戰.

總結來說,在詩⼈⼼⽬中,正是因為神「進來」與⼈同在和同⾏,讓⼈可以在看似不樂觀的處境中「抬起頭來」.

另⼀點要注意:跟上⼀個段落(3-6節)不同,這裡沒有提到⼈要達到怎樣的條件神 才會「進來」.唯⼀的要求,就是要把「永久的⾨⼾」——即屬於「永恆的神」的那 扉⾨——⼤⼤的打開.

B.《詩篇》廿四篇整體的解讀

經過上述對這篇詩篇三個部分的解釋後,相信⼤家都會發覺,這篇詩篇的三個部分看似是三個獨立的單元.⼤家都明⽩每個部分的內容,但三個部分放在⼀起的時候,就好像不知道詩⼈想表達甚麼.

不少的學者提出不同的理論,嘗試作出⼀些「假想」的慶典或節期,以「解釋」這篇詩篇的⽤途.但總的來說,這些解釋都未能完全符合詩篇的內容,特別是第⼆和第三部分的描述難以說是在同⼀個禮儀中發⽣:在第⼆部分提到⼈要登神的⼭,進神的殿； 但第三部分卻好像說,要打開城⾨讓神進入,亦即神還未在聖城之中.

有學者推斷,這是因為這篇詩篇其實由不同時期創作的詩篇組合⽽成——7-10 節或許在⼤衛王朝的時代已出現,⽽ 1-6節則是被擄歸回之後的作品. 有學者甚⾄認為,這篇詩篇由三篇不同的詩篇組成. 但即使是這樣,問題仍然存在,⼤家仍會問:為何《詩篇》的編修者要把這幾篇詩篇放在⼀起?他們不會「無厘頭」地把完全不相關的材料編在同⼀篇詩篇裡!那麼,這些編修者到底想表達甚麼信息?

當中的討論繁多,非常複雜,但總結起來,答案就是⼀篇⽤作教導的詩篇,⽽不是為了某個慶典或節期⽽創作. 當中⼀個非常重要的根據就是,《詩篇》廿四篇並非獨立存在,⽽是跟前⾯的九篇詩篇(即《詩篇》⼗五⾄廿四篇)形成⼀個「專輯」.這個「專輯」中前⾯和後⾯的詩篇彼此呼應,詩篇⼗五篇與廿四篇平⾏,⼗六篇與 詩廿三篇平⾏,⼗七篇與廿⼆篇平⾏,如此類推.相應的詩篇有相同的主題,但後者會對前者作出進⼀步的闡釋,為前者所留下來的問題提供解答.

所以,《詩篇》廿四篇正呼應和補充《詩篇》⼗五篇的內容.那麼,詩篇⼗五篇說了甚麼呢?跟今天詩篇的第⼆部分⼀樣,《詩篇》⼗五篇指出了⼈進入神的神聖靈域的必要條件,包括要「⾏為正直,做事公義,⼼裡說實話」(詩⼗五 2).所以,今天的詩篇就是要補充⼈進入神聖靈域的問題.

這問題就是先前所提到的「兩難」:⼈尋求神,要進入神的領域,就必須要達到今天詩篇第⼆部分或《詩篇》⼗五篇中所提到那些條件,但根本沒有⼈能靠自⼰滿⾜這些要求.試問有誰能完全做到真正的「⼿潔⼼清」呢?今天詩篇的第三部分,正是指出這個「兩難」困局的出路:就在⼈因「⼿潔⼼清」的要求感到無望,不知如何是好的時候,是神反過來作出主動,先進入這些尋求神的⼈的⽣命之中,在他們的內⼼之加⼒給他們,使他們得以逐步成聖和稱義——「成聖」即逐步擺脫世俗對他們的羈絆(即「⼼清」),⽽「稱義」就是愈來愈能實踐神眼中看為正的事(即「⼿潔」)——這樣,「尋求神的⼈」就能逐步⾛近和進入神臨格的神聖靈域之內.

有學者認為,《詩篇》廿四篇的第⼆和第三個部分,說明了這是神與⼈互相進入對⽅的領域之中,是⼀個相互的「mutual advent」. 但當中的次序不能搞錯,是神的「進來」在先,因為⼈是蒙恩得救,⽽不是靠自⼰的⾏為和自潔⽽得救.按著今天詩篇的講法,尋求神的⼈要做的,就是把本來緊閉的⼼⾨打開.

今天是今年「將臨期」的第⼀個主⽇,⽽「將臨期」就是要記念主耶穌基督的降⽣, 以及盼望祂的再來.今天的經文正要強調,是三⼀上主作主動,讓祂的聖道打破神聖和凡俗的界限,成為有⾎有⾁的拿撒勒⼈耶穌,進到這個充滿罪惡,也充滿著苦難,傷痛,眼淚,焦慮,壓抑,憤怒和死亡的世界之中.祂要打開⼀個新的局⾯,就是讓人可以重新回歸到神的神聖靈域之中,使他們的生命本性和人生路徑都得以逆轉,人類歷史的軌跡也因此得以改變.

四•引伸應用

近年坊間流行「正念」或「靜觀」(mindfulness )的操練,甚至政府衛生署的網站都有介紹這操練在臨床醫療和精神健康中的作用.不過,這跟信徒的靈修有極大的分別.彼此的分別在於,無論是動態還是靜態的靈修,都是要把自己的心門打開——而且不只是最外面的那扉門,也包括自己內心中各個「暗室」的門——讓神進入自己生命的每一個領域之中.

這樣說有點抽象.大家有沒有看過〈反轉腦朋友2〉這齣卡通片?這套卡通片把人內心的運作非常形像化表達出來,指出人內心中有很多不同的「地方」:不想要的記憶就丟到遠遠的後山；有些「黑歷史」要鎖起,放在保險庫中；有大房專門負責人想像力的運作；還有主控制室.神要進入這些地方,幫助一個人認識和面對一些可能連他自己都不知道或遺忘了的「真正的自己」,並且加力給他作出改變.

所以,如果信徒只是一味想追求進入神聖的領域中,希望得到異象,又或者希望得到特殊的經驗或感動,卻沒有懷著謙卑順服的心來讓神更新改變自己,嚴格來說,這樣的靈修仍未達到應有的果效.

在去年的某次退修中曾做過一次「觀畫靈修」,而其中一副畫就是基督站在門外叩門的圖畫.那扉門已經長滿了野草,說明這門緊閉了很久,沒有讓基督(即基督的靈) 進入.大家可以想一想,這意味著基督只能隔著門跟這人說話.

弟兄姊妹,這是否也正是自己今天生命的光景?

當今天這篇詩篇呼籲大家要把「城門」「抬起」時,正是提醒大家要好好思想,自己是甚麼原因把曾經打開的「城門」緊密起來?還是從來也沒有打開過?

基督現在就在大家的心門之外,請把「城門」打開!



\newpage

\section{耶西的根}
\label{sec:776}
\textbf{李均熊牧師 - 耶西的根}
\newline
\newline
\hyperref[sec:775]{< < < PREV SERMON < < <}
~
\hyperref[sec:toc]{[返目錄]}
~
\hyperref[sec:777]{> > > NEXT SERMON > > >}


以賽亞書 11:1-10
\newline
\begin{longtable}{cl}
\hline
\hline
章節 & 經文 (和合本修訂版)\\
\hline
11:1 & \begin{tabularx}{0.7\textwidth}{X} 從耶西的殘幹必長出嫩枝，他的根所抽的枝子必結果實。 \end{tabularx} \\ \\ \relax
11:2 & \begin{tabularx}{0.7\textwidth}{X} 耶和華的靈必住在他身上，就是智慧和聰明的靈，謀略和能力的靈，知識和敬畏耶和華的靈。 \end{tabularx} \\ \\ \relax
11:3 & \begin{tabularx}{0.7\textwidth}{X} 他必以敬畏耶和華為樂；行審判不憑眼見，斷是非也不憑耳聞； \end{tabularx} \\ \\ \relax
11:4 & \begin{tabularx}{0.7\textwidth}{X} 卻要以公義審判貧寒人，以正直判斷地上的困苦人，以口中的棍擊打全地，以嘴裡的氣殺戮惡人。 \end{tabularx} \\ \\ \relax
11:5 & \begin{tabularx}{0.7\textwidth}{X} 公義必當他的腰帶，信實必作他脅下的帶子。 \end{tabularx} \\ \\ \relax
11:6 & \begin{tabularx}{0.7\textwidth}{X} 野狼必與小綿羊同住，豹子與小山羊同臥；少壯獅子、牛犢和肥畜同群；孩童要牽引牠們。 \end{tabularx} \\ \\ \relax
11:7 & \begin{tabularx}{0.7\textwidth}{X} 牛必與熊同食，牛犢與小熊同臥；獅子與牛一樣吃草。 \end{tabularx} \\ \\ \relax
11:8 & \begin{tabularx}{0.7\textwidth}{X} 吃奶的嬰孩在虺蛇的洞口玩耍，斷奶的幼兒必按手在毒蛇的穴上。 \end{tabularx} \\ \\ \relax
11:9 & \begin{tabularx}{0.7\textwidth}{X} 在我聖山各處，牠們都不傷人，不害物；因為認識耶和華的知識要遍滿全地，好像水充滿海洋一般。 \end{tabularx} \\ \\ \relax
11:10 & \begin{tabularx}{0.7\textwidth}{X} 到那日，耶西的根立作萬民的大旗；列國的人必尋求他，他安歇之所大有榮耀。 \end{tabularx} \\ \\
[1ex]
\hline
\hline
\end{longtable}
講題:耶西的根 A Root of Jesse

經課:以賽亞書11:1-10

講員:李均熊牧師

場合:香港中文大學崇基學院禮拜堂主日崇拜

日期:2025年12月7日

起:理想君王與理想國度

今日選讀的舊約經課是以賽亞書11章1至10節.在經文6-9節展現的,正好與今日啟應詩篇回響.詩篇72篇明確地祈求君王的特質,所以我們可以推斷賽11章2-5節是在寫出理想的君王,而緊接的6-9節則是寫出一個理想國度.詩篇72篇的祈求,在以賽亞書成了上主的應許.

今日不能詳細討論6-9節的理想國度／社會的意象,只可以說它應該不會照字面般應驗.因為今日我們從生態研究知道,屬於食物鏈上層的狩獵類動物,對整個生態系統是十分重要的,若然牠們好像先知這裡所描繪的不獵殺,而且吃草作為糧食,那可能會是一場生態災難.不過上主有奇妙的作為,將來害人的,兇殘的,無論是人是獸,如何弱小如嬰孩的人可以和平共處,我們有限的目光真是不能看透.但無論如何,和平的社會都是今日我們極度祈盼的.

但詩篇72篇,又或賽11:2-5所描寫的理想國君,在香港人的經驗中,可以說是十分容易明白,也是我們十分渴望出現的社會領袖.他要以公平公義治理,並且以信實公義審判,絕不偏袒.行審判不憑眼見,斷是非也不憑耳聞,可以想像為香港終審法院建築頂的帶著眼罩的正義女神像.現代的解釋通常不會認為她「無眼睇」,而是公正不偏袒任何人,特別不會因為你是權貴富豪就對你特別優待.

「公平」與「公義」在古代近東/西亞文獻常作為重言用法(hendiadys),即兩個詞語同時使用時,「公平和公義」(和合本多譯作「秉公行義」)多是指到社會公義的情況(social justice),而不僅是司法管轄的概念(jurisdiction concept),它與制定法律並負責執行法律的社會政治領導者更為相關(much more relevant for the social-political leaders who create the laws and are reponsible for their execution, M. Weinfeld 1995, p.44).在古代文獻的記載中,建立公平公義的社會,乃是君王——就是古時制定法律並負責執行法律的社會政治領導者——的任務,就是要照顧社會中無權無勢,以孤兒寡婦為代表的弱勢社群,並且審判不能偏私權貴.所以先知提到上主應許的理想君王,並非內聖外王德智具全的聖賢人物(sages who possess both inner virtue and outer wisdom),而是關心社會,建立公平公義制度,不讓無權無勢的弱小巿民,被自以為高人一等,以為有權有勢就可以隻手遮天,屈枉正直的權貴所欺負.這不也是我們今日所祈望理想的「制定法律並負責執行法律的社會政治領導者」?

承:回歸耶西的根

但由經課開始第1節及結束的第10節,都提到耶西的根.耶西是誰?熟識舊約的朋友可能會知道他就是大衛的父親.所以耶西的根所生的枝子,不就是大衛了嗎?那麼2-5節的各種理想君王的特點,不就是指大衛了嗎?但經文用上耶西的根,而不是直指大衛,其實是有其深意的.只靠這裡一段經文,背後的意思可能不太明顯；但若我們將整卷以賽亞書君王的主題貫串來看,回歸耶西的根,正是要將當時人們把理想君王的祈盼,全放在人間的大衛王朝之上的這種想法,扭轉過來.

在以色列人被擄流散前,先知可能仍然相信,上主其實揀選了大衛和他的家族王朝,去成為實踐公平公義的理想君王(賽16:5「必有寶座因慈愛堅立；必有一位誠誠實實坐在其上,在大衛帳幕中施行審判,尋求公平,速行公義.」),但隨著帝國的侵略,以色列國和猶大國因上主審判而相繼滅亡,先知便明白到人間的王朝並不能實現上主心目中的理想.在以賽亞書後段呼應今天經課的經文中,賽32:1-5:

32:1 看哪,必有一王憑公義行政；必有首領藉公平掌權.

32:2 必有一人像避風所和避暴雨的隱密處,又像河流在乾旱之地,像大磐石的影子在疲乏之地.

32:3 那能看的人,眼不再昏迷；能聽的人,耳必得聽聞.

32:4 冒失人的心必明白知識；結吧人的舌必說話通快.

32:5 愚頑人不再稱為高明；吝嗇人不再稱為大方.

(這段相信是亞述消滅北國以色列期間寫成的),先知已經不再將這個理想君王的角色,與大衛王朝捆綁在一起；甚至在我們熟識的「以馬內利」預言中(賽7:13-14),也不明言將要出生的「以馬內利」是否大衛家族中的子孫(亞哈斯兒子希西家王?),在整段敘事中(先知和亞哈斯王的會面),這個留白更顯突兀,也浪費了很多解經家的筆墨.

其實早在以賽亞書6章我們就了解到,先知在經驗人間國君的過去,就看見上主才是在高高的寶座上作王的那一位君主(賽6:1「當烏西雅王崩的那年,我見主坐在高高的寶座上.」).但要到賽40-55章中,從內容上看是屬於被據流散時期寫作的經文,就更明顯不再把盼望放在大衛王朝,而是要等候上主親自的指派和揀選.這裡只抽兩段經文去展示一下:

41:2 誰從東方興起一人,憑公義召他來到腳前呢?耶和華將列國交給他,使他管轄君王,把他們如灰塵交與他的刀,如風吹的碎稭交與他的弓.

41:3 他追趕他們,走他所未走的道,坦然前行.

41:4 誰行做成就這事,從起初宣召歷代呢?就是我－耶和華!我是首先的,也與末後的同在.

大家更為熟識的賽42章,更在某程度上回應著今日的舊約經課:

42:1 看哪,我的僕人－我所扶持所揀選,心裡所喜悅的!我已將我的靈賜給他；他必將公理傳給外邦.

42:2 他不喧嚷,不揚聲,也不使街上聽見他的聲音.

42:3 壓傷的蘆葦,他不折斷；將殘的燈火,他不吹滅.他憑真實將公理傳開.

42:4 他不灰心,也不喪膽,直到他在地上設立公理；海島都等候他的訓誨.

42:5 創造諸天,鋪張穹蒼,將地和地所出的一併鋪開,賜氣息給地上的眾人,又賜靈性給行在其上之人的神耶和華,他如此說:

42:6 我－耶和華憑公義召你,必攙扶你的手,保守你,使你作眾民的中保,作外邦人的光,

42:7 開瞎子的眼,領被囚的出牢獄,領坐黑暗的出監牢.

42:8 我是耶和華,這是我的名；我必不將我的榮耀歸給假神,也不將我的稱讚歸給雕刻的偶像.

42:9 看哪,先前的事已經成就,現在我將新事說明,這事未發以先,我就說給你們聽.

轉:失效人君

換句話說,以賽亞書此處提到耶西的根,並不是期望大衛王朝的復興.其實這裡提到耶西的根,乃延續以賽亞書10章關於「樹」的意象.你若有聖經,可以退後幾節,就可以讀到:

10:33 看哪,主－萬軍之耶和華以驚嚇削去樹枝；長高的必被砍下,高大的必被伐倒.

10:34 稠密的樹林,他要用鐵器砍下；黎巴嫩的樹木必被大能者伐倒.

若我們再往前看,這其實是先知對當時「制定法律並負責執行法律的社會政治領導者」(君王?)的審判:

10:1 禍哉!那些設立不義之律例的和記錄奸詐之判語的,

10:2 為要屈枉窮乏人,奪去我民中困苦人的理,以寡婦當作擄物,以孤兒當作掠物.

10:3 到降罰的日子,有災禍從遠方臨到,那時,你們怎樣行呢?你們向誰逃奔求救呢?你們的榮耀存留何處呢?

10:4 他們只得屈身在被擄的人以下,仆倒在被殺的人以下.雖然如此,耶和華的怒氣還未轉消；他的手仍伸不縮.

10:5 亞述是我怒氣的棍,手中拿我惱恨的杖.

10:6 我要打發他攻擊褻瀆的國民,吩咐他攻擊我所惱怒的百姓,搶財為擄物,奪貨為掠物,將他們踐踏,像街上的泥土一樣.

亞述實為上主審判猶大國的大衛王朝——當時的「制定法律並負責執行法律的社會政治領導者」——的工具.

今日舊約經課的第一節:從耶西的本必發一條；從他更新的枝子必結果實.實乃延續由第七章開始,先知對大衛家的審判:以賽亞說:「大衛家啊,你們當聽!你們使人厭煩豈算小事,還要使我的神厭煩嗎?(7:13)與其說審判,不如說是當時王室並無履行他們應有的「秉公行義」的責任,反而拒絕信靠上主,並且自以為了不起,急於尋找政治同盟去抵抗亞述,最終只會自食其果:看哪,主－萬軍之耶和華以驚嚇削去樹枝；長高的必被砍下,高大的必被伐倒.(賽10:33).就算之前上主應許大衛永遠的約,但你自高自大,不實踐上主對君王(「制定法律並負責執行法律的社會政治領導者」)的「秉公行義」的要求,上主也可以推倒重來,由耶西之根再選一次一個真正符合上主理想的君王.看來福音經課中施洗約翰對當時對法利賽人和撒都該人的挑戰也可以如此理解:

9 不要自己心裡說:『有亞伯拉罕為我們的祖宗.』我告訴你們,上帝能從這些石頭中給亞伯拉罕興起子孫來.10 現在斧子已經放在樹根上,凡不結好果子的樹就砍下來,丟在火裡.(太3:9-10)

合:將臨期的等候

所以今日經課所選以賽亞的經文,承載著對人間君王的失望,就是對失效於「秉公行義」的「制定法律並負責執行法律的社會政治領導者」失望,但也同時承載著深信只有上主的靈,才能夠建立真正理想的君王與和平國度的盼望.

今日我們對失效的君王有多麼的失望；對未能「秉公行義」的「制定法律並負責執行法律的社會政治領導者」多麼的忿怒；對以賽亞書10章開始所講的「那些設立不義之律例的和記錄奸詐之判語的,為要屈枉窮乏人,奪去我民中困苦人的理,以寡婦當作擄物,以孤兒當作掠物.」(賽10:1-2)的社會情況多麼的痛心——這些經文,其實都提醒我們,人間的君王是多麼不可靠,不可信.正是因為這樣,我們更加要仰望那能「使人有智慧和聰明的靈,謀略和能力的靈,知識和敬畏耶和華的靈」的上主之靈.等候那能帶來和平國度的耶西的根.

等等.耶西的根,不是已經來了?我們將臨期等待的,不就那來了一次,又要再來的耶西的根嗎?

羅15:12-13

12 又有以賽亞說:將來有耶西的根,就是那興起來要治理外邦的；外邦人要仰望他.

13 但願使人有盼望的上帝,因信將諸般的喜樂,平安充滿你們的心,使你們藉著聖靈的能力大有盼望!



\newpage

\section{有甚麼值得喜樂}
\label{sec:777}
\textbf{林豪恩先生 - 有甚麼值得喜樂?}
\newline
\newline
\hyperref[sec:776]{< < < PREV SERMON < < <}
~
\hyperref[sec:toc]{[返目錄]}
~
\hyperref[sec:778]{> > > NEXT SERMON > > >}


以賽亞書 35:1-10
\newline
\begin{longtable}{cl}
\hline
\hline
章節 & 經文 (和合本修訂版)\\
\hline
35:1 & \begin{tabularx}{0.7\textwidth}{X} 曠野和乾旱之地必然歡喜，沙漠也必快樂；又如玫瑰綻放， \end{tabularx} \\ \\ \relax
35:2 & \begin{tabularx}{0.7\textwidth}{X} 朵朵繁茂，其樂融融，而且歡呼。黎巴嫩的榮耀，並迦密與沙崙的華美，必賜給它。人要看見耶和華的榮耀，看見我們神的榮美。 \end{tabularx} \\ \\ \relax
35:3 & \begin{tabularx}{0.7\textwidth}{X} 你們要使軟弱的手強壯，使無力的膝蓋穩固； \end{tabularx} \\ \\ \relax
35:4 & \begin{tabularx}{0.7\textwidth}{X} 對心裡焦急的人說：「要剛強，不要懼怕。看哪，你們的神要來施報，要施行極大的報應，他必來拯救你們。」 \end{tabularx} \\ \\ \relax
35:5 & \begin{tabularx}{0.7\textwidth}{X} 那時，盲人的眼必睜開，聾子的耳必開通。 \end{tabularx} \\ \\ \relax
35:6 & \begin{tabularx}{0.7\textwidth}{X} 那時，瘸子必跳躍如鹿，啞巴的舌頭必歡呼。在曠野有水噴出，在沙漠有江河湧流。 \end{tabularx} \\ \\ \relax
35:7 & \begin{tabularx}{0.7\textwidth}{X} 火熱之地要變為水池，乾渴之地要變為泉源。野狗躺臥休息之處必長出青草、蘆葦和蒲草。 \end{tabularx} \\ \\ \relax
35:8 & \begin{tabularx}{0.7\textwidth}{X} 在那裡必有一條大道，就是一條路，稱為聖路。污穢的人不得經過，是專為走路的人預備的，愚昧的人也不會迷路。 \end{tabularx} \\ \\ \relax
35:9 & \begin{tabularx}{0.7\textwidth}{X} 在那裡沒有獅子，猛獸也不經過；在那裡牠們未現蹤跡，只有救贖的民在那裡行走。 \end{tabularx} \\ \\ \relax
35:10 & \begin{tabularx}{0.7\textwidth}{X} 耶和華救贖的民必歸回，歌唱來到錫安；永遠的快樂必歸到他們頭上，他們必得著歡喜快樂，憂傷嘆息盡都逃避。 \end{tabularx} \\ \\
[1ex]
\hline
\hline
\end{longtable}
馬太福音 11:2-11
\newline
\begin{longtable}{cl}
\hline
\hline
章節 & 經文 (和合本修訂版)\\
\hline
11:2 & \begin{tabularx}{0.7\textwidth}{X} 約翰在監獄裡聽見基督所做的事，就派他的門徒去， \end{tabularx} \\ \\ \relax
11:3 & \begin{tabularx}{0.7\textwidth}{X} 問耶穌：「將要來的那位就是你嗎？還是我們要等候另一位呢？」 \end{tabularx} \\ \\ \relax
11:4 & \begin{tabularx}{0.7\textwidth}{X} 耶穌回答他們：「你們去，把所聽見、所看見的告訴約翰： \end{tabularx} \\ \\ \relax
11:5 & \begin{tabularx}{0.7\textwidth}{X} 就是盲人看見，瘸子行走，痲瘋病人得潔淨，聾子聽見，死人復活，窮人聽到福音。 \end{tabularx} \\ \\ \relax
11:6 & \begin{tabularx}{0.7\textwidth}{X} 凡不因我跌倒的有福了！」 \end{tabularx} \\ \\ \relax
11:7 & \begin{tabularx}{0.7\textwidth}{X} 他們一走，耶穌就對眾人談到約翰，說：「你們從前到曠野去，是要看甚麼呢？看風吹動的蘆葦嗎？ \end{tabularx} \\ \\ \relax
11:8 & \begin{tabularx}{0.7\textwidth}{X} 你們出去到底是要看甚麼？看穿細軟衣服的人嗎？那穿細軟衣服的人是在王宮裡。 \end{tabularx} \\ \\ \relax
11:9 & \begin{tabularx}{0.7\textwidth}{X} 你們出去究竟是要看甚麼？是先知嗎？是的，我告訴你們，他比先知大多了。 \end{tabularx} \\ \\ \relax
11:10 & \begin{tabularx}{0.7\textwidth}{X} 這個人就是經上所說的：『看哪，我要差遣我的使者在你面前，他要在你前面為你預備道路。』 \end{tabularx} \\ \\ \relax
11:11 & \begin{tabularx}{0.7\textwidth}{X} 我實在告訴你們，凡女子所生的，沒有一個比施洗約翰大；但在天國裡，最小的比他還大。 \end{tabularx} \\ \\
[1ex]
\hline
\hline
\end{longtable}
講題:有甚麼值得喜樂? What is Worthy of Joy?

經課:以賽亞書35:1-10;馬太福音11:2-11

講員:林豪恩先生

場合:香港中文大學崇基學院禮拜堂主日崇拜

日期:2025年12月14日

引言:你快樂嗎?

大家有沒有留意今天燃燭禮所燃起的不是紫色蠟燭,而是粉紅色蠟燭呢?因為今天是將臨期第三主日,稱為「喜樂主日」,燃點的是代表喜樂的粉紅色蠟燭.如果問這個城市的人,你快樂嗎?會得到甚麼回答呢?或許有人很清醒地指出,這是喜樂主日,你問人是否快樂,偷換概念呢!其實,在聖經中關於喜樂或快樂的詞彙的確不止一個,不過,它們多是相通的,有時翻譯為喜樂,有時翻譯為快樂,有時翻譯為歡喜快樂,就好像英文Happy,是開心,快樂,還是幸福呢?每年一度的World Happiness Report,很少譯為世界開心報告,通常譯為世界快樂報告,或世界幸福報告.我們暫且把快樂和喜樂視為類似的狀況吧.

講起世界快樂報告,香港人和其他地方的人相比,快樂程度如何呢?2025年的報告顯示香港人的快樂指數排第88位,有新聞標題說「香港人的快樂程度低過菲律賓和印尼」(香港01).如果我們問香港人為甚麼不快樂,他們可能會反問:有甚麼值得快樂呢?今天又是每年一度的喜樂主日,如果我們問自己這個問題:我喜樂嗎?或許我們也如大眾市民一樣心中泛起疑問:有甚麼值得喜樂嗎?如果我們帶著這個問題去閱讀今天喜樂主日的經課,經課中的人物又會怎樣回答這個問題呢?

一．疑惑,而不是喜樂

首先,讓我們問一問施洗約翰:約翰,你喜樂嗎?施洗約翰可能請我們從馬太福音的記述去尋找答案.

11:2約翰在監裡聽見基督所做的事,就打發兩個門徒去,

11:3問他說:「那將要來的是你嗎?還是我們等候別人呢?」

讀到這兩節經文,可能大家都領悟到了.約翰的處境以及他的心情,相信不是喜樂,而是疑惑.是甚麼令他感到疑惑呢?就是他所相信及所宣揚的耶穌基督.竟然是耶穌令他疑惑?就是耶穌令他疑惑!你心裡曾經湧現過對耶穌的疑惑嗎?你有否問過類似的問題:「那將要來的是你嗎?還是我們等候別人呢?」約翰為甚麼對耶穌產生疑惑呢?因為他曾經對耶穌投以信心,還記得他如何向群眾介紹耶穌嗎?

3:11我是用水給你們施洗,叫你們悔改.但那在我以後來的,能力比我更大,我就是給他提鞋也不配.他要用聖靈與火給你們施洗.

3:12他手裡拿著簸箕,要揚淨他的場,把麥子收在倉裡,把糠用不滅的火燒盡了.」

約翰曾經宣認的那位大能的耶穌,越看越不似是振臂一呼,扭轉乾坤,拯救人民於水火之的英雄.如果約翰是旁觀者,只是「跌了眼鏡」,無傷大雅.然而,約翰已經為這位耶穌押上了一生,現實的落差卻如此巨大,怎能不產生懷疑呢?耶穌來是來了,但世界竟然一切依舊.約翰呼籲悔改,但世界不為所動,自己反而被囚牢獄,怎可能不疑惑呢?本來令約翰喜樂的,現在變成了令約翰疑惑.如此的狀況有甚麼值得喜樂呢?各位朋友,在你的信仰經歷中,有過這種狀況嗎?有人說,越認真的基督徒,就越有可能經驗這樣的疑惑.

有一位牧師,家中3名親人被困火場.當時家人們對牧師說,「你係我地屋企唯一ge 牧師,你可唔可以祈禱搵番佢地出黎?」在馮梁結中學等候期間,牧師與信主的家人圍在一起禱告.當刻,身邊許多不認識的人也走到他們身旁陪伴祈禱.「gwo 一刻我覺得有盼望,應該得ge .」然而,他們最後還是收到家人離世的消息.牧師說,那一刻「感覺好苦」,眼淚潰堤,「因為作為一個牧師,我係度祈禱,我發現我做唔到,救唔到屋企人.」而其親人亦拉著她說:「大表姐你係牧師黎,你快D 祈禱等佢地復活,等佢地翻生啦!」(時代論壇2025/12/6)

「那將要來的是你嗎?還是我們等候別人呢?」德蘭修女也問過類似的問題.她曾經向一位神父寫一小紙條說:「神父,請為我祈禱.耶穌在哪裡呢?」又有一次,她體力不支昏倒,醒來發現自己躺在床上,她流著眼淚說:「我辜負了他們,我辜負了那些貧民窟的人.」德蘭修女對自己的選擇產生了懷疑:「這真是天主要我做的嗎?我真的要在這地獄般的貧民窟裡度過我的一生嗎?世界上有那麼多的窮人,而我個人所能做的又這麼有限,就好像在黑暗的地獄裡,裝飾出一個明亮的小角落,僅此而已,我的工作究竟有什麼意義呢?」(德蘭修女傳節錄)

二．怎會值得喜樂呢?

到底有甚麼值得喜樂呢?有沒有事情值得喜樂,答案或許是和人們認為甚麼狀況值得喜樂有關.耶穌在回答約翰的問題時說了一句耐人尋味的話:

11:6凡不因我跌倒的就有福了!

耶穌是基督,並不是要來迷惑人的魔鬼,人們怎可能因祂跌倒呢?怎麼不可能呢?試試看耶穌向約翰提出證明自己就是那位「將要來的」的「證據」是甚麼:

11:5就是瞎子看見,瘸子行走,長大痲瘋的潔淨,聾子聽見,死人復活,窮人有福音傳給他們.

聽到如此的「證明」,人們可能反問:和我有甚麼關係?我不是瞎子；又不是瘸子,也沒有大痲瘋,還未死去,也不是窮人.如果彌賽亞就只是做這些低調又個別的事情,怎麼不會令人跌倒呢?

如果你問:「瞎子看見,瘸子行走,長大痲瘋的潔淨,聾子聽見,死人復活,窮人有福音傳給他們」有甚麼不好?人們可能會回答,沒有甚麼不好,只是和我沒有關係.有專欄作家指出,香港人的「獅子山精神」其實是拚搏賺錢的精神.拚搏是為甚麼?是為了「錢」!捱幾份工為甚麼?為多個錢,令到自己早日唔使做,可以去嘆世界.因此,香港人最喜歡看有錢佬發達的故事.但是,很多人等足幾十年,等足一世都沒有發達機會,於是心中對發了達的人又忌又恨.香港人「獅子山精神」的反面變成了「憎人富貴厭人窮」的心態,其實是說不出多麼的苦澀和失落.(劉天賜:憎人富貴厭人窮)

三．懂的,就會喜樂

施洗約翰聽到耶穌的回答,能夠釋除他的疑惑嗎?他會因這答案而感到喜樂嗎?馬太福音沒有跟進約翰的回應.不過,熟識舊約聖經的約翰會從耶穌的回答聽得出,那正正是古時先知的呼聲,是歷世歷代盼望的所在.這呼聲在哪裡發出的呢,就是今天的舊約經課以賽亞書35章:

35:5那時,瞎子的眼必睜開；聾子的耳必開通.

35:6那時,瘸子必跳躍像鹿；啞巴的舌頭必能歌唱.

這些現象代表甚麼呢?原來這些現象不但和那些有缺欠的人相關,也和整個群體有關.以賽亞書35章描繪的景況是這樣的:

35:1曠野和乾旱之地必然歡喜；沙漠也必快樂；又像玫瑰開花,

35:2必開花繁盛,樂上加樂,而且歡呼.黎巴嫩的榮耀,並迦密與沙崙的華美,必賜給它.人必看見耶和華的榮耀,我們上帝的華美.

為甚麼和整個群體有關呢?因為那是被擄回歸的時刻.

35:8在那裡必有一條大道,稱為聖路.污穢人不得經過,必專為贖民行走；行路的人雖愚昧,也不致失迷.

35:9在那裡必沒有獅子,猛獸也不登這路；在那裡都遇不見,只有贖民在那裡行走.

35:10並且耶和華救贖的民必歸回,歌唱來到錫安；永樂必歸到他們的頭上；他們必得著歡喜快樂,憂愁歎息盡都逃避.

「瞎子的眼睜開；聾子的耳開通.瘸子跳躍；啞巴歌唱」這現象代表耶和華拯救來到的時刻.當耶穌如此說時,約翰領悟到了,耶穌就是那位拯救者,那位歷世歷代所盼望的.這份確認沒有令約翰處身的環境改變,然而,他確認耶穌就是那位要來的,正如在聖殿的西面看見嬰兒耶穌就說:

2:29主啊!如今可以照你的話,釋放僕人安然去世；

2:30因為我的眼睛已經看見你的救恩－

2:31就是你在萬民面前所預備的:

四．一個肢體受苦,所有的肢體就一同受苦

有甚麼值得歡喜快樂呢?答案視乎人們看甚麼為重要.正如上述那位專欄作家一樣,人人以為香港人為了發達而只顧日以繼夜地拼搏,不理會其他人死活.然而,最近發生的災難令每一個人都感受到保羅所講的話:「一個肢體受苦,所有的肢體就一同受苦」.不但災難受害者的親友感到擔心和傷痛,很多和受災者不相識的市民都感都到擔憂和傷痛；不但受災者的親友流淚,連平日碰面都不會打招呼的陌生人都流淚；在災後的那個星期六和星期日,那條往現場悼念的人龍,從太和站排到大埔墟站延線不絕.原來,人和人之間是休戚與共的,其他人的福祉和自己是有關係的.如果社會中的瞎子沒有人理會,瘸子沒有人照顧,長大痲瘋的沒有人醫治,聾子被邊境化,死人沒有人哀悼,窮人被忽視,生活於如此的社會中,人們是會感到難過的.

有人說:「一個文明的衡量標準,在於它如何對待它最弱勢的成員」.「瞎子看見,瘸子行走,長大痲瘋的潔淨,聾子聽見,死人復活,窮人有福音傳給他們」,如此的社會,不正是一個文明的社會嗎?不正是值得歡喜快樂的情景嗎?

五．總結:若一個肢體得榮耀,所有的肢體就一同快樂

現實的落差總是令人失望,令人產生疑惑.尤其是看見人類的自私和貪婪如此根深柢固,所製造出來的災難帶來如此巨大的傷害,令人感到痛心疾首.連跟隨耶穌的人都要問:「那將要來的是你嗎?還是我們等候別人呢?」我們多多少少都感受到施洗約翰的疑惑,耶穌來了,但世界似乎沒有改變,人類的自私和貪婪似乎沒有收歛.這是令人們氣餒的現實,然而,這也不是現實的全部.請看看,在烈火焚城之時,有多少無私付出的畫面:消防員冒著生命危險進入火場救人；有被困火場的住客在生死存亡之際仍然照顧鄰居的長者,在死亡的威協之下互相守望；醫院全體總動員日以計夜救治傷者；港九新界不同的餐廳,超市立即出現提供協助,貨運公司一車車的物質趕入新界.地區教會,社區會堂,藝術中心開放之外,有party room取消當晚的預訂留給有需要的災民.現場甚至有寵物氧氣箱,有寵物醫院.就連便利店充電器都出現全大埔免費充電,免費借用的服務.火災發生的第二日,中大有同學們聚集在禮拜堂地下祈禱.有一位同學說,他在宿舍望到遠處的煙火,又看到收集物資的帖文,想把自己的棉被送出去,但又怕尷尬.掙扎了一會,他和房友決定抱著自己的棉被送去收集站.就在往物資收集站的路上,他看見有同學帶著衣物,有同學拿著杯麵,有同學拿著充電器,一齊行去物資收集站.他說,有些平時看似冷漠的同學,竟然都行出來了.原來,人是會因其他人受苦而感到苦,當其他人的苦得以舒緩,才會感到快樂.

大家還記得那位家人在火災中離世的牧師嗎?她後來怎樣呢?牧師提到在她最難過,不懂如何向家人解釋所發生的事的時候,有人向她遞上兩粒糖,讓她在苦中得到一點甜.她說「呢個係我地屋企ge 苦難,亦都係香港人共同面對ge 苦難.但我好相信,係苦難當中,仲有愛.係苦難當中,睇到基督在人間.」

德蘭修女呢?她怎樣在無論幫助了多少貧弱者,都仍有更多貧弱者等待幫助的困境中肯定自己的呼召和服務呢?「一九八七年復活節前棕樹主日那天,羅馬城內一位天主教神父突然聽到奇異聲音,吩咐他告訴德蘭修女一句話:「我渴了」(I thirst).這位神父根本不認識德蘭修女,也從沒經歷過神蹟奇事,是非常理性的知識份子,他以為自己聽錯了.誰知,同樣的信息接連出現三次,他無法抗拒,拿出信紙把經歷寫下來,寄去印度給德蘭修女.德蘭修女看到這封信仿似晴天霹靂,原來一九四六年她在開往大吉嶺的火車上,聽到耶穌基督呼召她,關鍵的一句就是基督說:「我渴了」(這是福音書記載耶穌在十字架上說的七句話中的一句),這句話她沒有告訴任何人.上帝藉著這封信讓德蘭修女知道,雖然她感覺不到親密甘甜的屬靈安慰,但上帝的同在仍然是千真萬確的.」(時代論壇2017年9月23日)

使徒保羅在哥林多前書說:

12:23身上肢體,我們看為不體面的,越發給它加上體面；不俊美的,越發得著俊美.

12:24我們俊美的肢體,自然用不著裝飾；但上帝配搭這身子,把加倍的體面給那有缺欠的肢體,

12:25免得身上分門別類,總要肢體彼此相顧.

12:26若一個肢體受苦,所有的肢體就一同受苦；若一個肢體得榮耀,所有的肢體就一同快樂.



\newpage

\section{不是倚靠勢力不是倚靠才能乃是倚靠神的靈}
\label{sec:778}
\textbf{鄧瑞強牧師 - 不是倚靠勢力,不是倚靠才能,乃是倚靠神的靈}
\newline
\newline
\hyperref[sec:777]{< < < PREV SERMON < < <}
~
\hyperref[sec:toc]{[返目錄]}
~
\hyperref[sec:779]{> > > NEXT SERMON > > >}


撒迦利亞書 4:1-14
\newline
\begin{longtable}{cl}
\hline
\hline
章節 & 經文 (和合本修訂版)\\
\hline
4:1 & \begin{tabularx}{0.7\textwidth}{X} 那與我說話的天使又來叫醒我，好像人睡覺時被喚醒一樣。 \end{tabularx} \\ \\ \relax
4:2 & \begin{tabularx}{0.7\textwidth}{X} 他問我：「你看見甚麼？」我說：「我看見了，看哪，有一個純金的燈臺，頂上有燈座，其上有七盞燈，每盞燈的上頭有七根管子； \end{tabularx} \\ \\ \relax
4:3 & \begin{tabularx}{0.7\textwidth}{X} 旁邊有兩棵橄欖樹，一棵在燈座的右邊，一棵在燈座的左邊。」 \end{tabularx} \\ \\ \relax
4:4 & \begin{tabularx}{0.7\textwidth}{X} 我問與我說話的天使說：「主啊，這是甚麼意思？」 \end{tabularx} \\ \\ \relax
4:5 & \begin{tabularx}{0.7\textwidth}{X} 與我說話的天使回答，對我說：「你不知道這是甚麼意思嗎？」我說：「主啊，我不知道。」 \end{tabularx} \\ \\ \relax
4:6 & \begin{tabularx}{0.7\textwidth}{X} 他回答我說：「這是耶和華指示所羅巴伯的話。萬軍之耶和華說：不是倚靠勢力，不是倚靠才能，乃是倚靠我的靈方能成事。 \end{tabularx} \\ \\ \relax
4:7 & \begin{tabularx}{0.7\textwidth}{X} 大山哪，你算甚麼呢？在所羅巴伯面前，你必夷為平地。他安放頂上的那塊石頭，人就歡呼：『願恩惠、恩惠歸與這殿！』」 \end{tabularx} \\ \\ \relax
4:8 & \begin{tabularx}{0.7\textwidth}{X} 耶和華的話臨到我，說： \end{tabularx} \\ \\ \relax
4:9 & \begin{tabularx}{0.7\textwidth}{X} 「所羅巴伯的手立了這殿的根基，他的手也必完成這工，你就知道萬軍之耶和華差遣我到你們這裡。 \end{tabularx} \\ \\ \relax
4:10 & \begin{tabularx}{0.7\textwidth}{X} 誰藐視這日的事為小呢？他們見所羅巴伯手拿石垂線就歡喜。這七盞燈是耶和華的眼睛，遍察全地。」 \end{tabularx} \\ \\ \relax
4:11 & \begin{tabularx}{0.7\textwidth}{X} 我問天使說：「那麼在燈臺左右的這兩棵橄欖樹是甚麼意思呢？」 \end{tabularx} \\ \\ \relax
4:12 & \begin{tabularx}{0.7\textwidth}{X} 我再次問他：「這兩根橄欖樹枝在兩根流出金色油的金嘴旁邊，是甚麼意思呢？」 \end{tabularx} \\ \\ \relax
4:13 & \begin{tabularx}{0.7\textwidth}{X} 他對我說：「你不知道這是甚麼意思嗎？」我說：「主啊，我不知道。」 \end{tabularx} \\ \\ \relax
4:14 & \begin{tabularx}{0.7\textwidth}{X} 他說：「這是兩位受膏者，侍立在全地之主的旁邊。」 \end{tabularx} \\ \\
[1ex]
\hline
\hline
\end{longtable}
講題:不是倚靠勢力,不是倚靠才能,乃是倚靠神的靈 Not by Might, nor by Power, but by God’s Spirit

經課:撒迦利亞書4:1-14

講員:鄧瑞強牧師

場合:香港中文大學崇基學院禮拜堂主日崇拜

日期:2025年12月21日

各位弟兄姊妹,早安.願神的恩典常與你們同在.

現在是「將臨期」,按道理,我應講些與將臨期相關的經文.但大家都知道,我按著聖經經卷一卷一卷的講,剛講到撒迦利亞書,我想了想,還是照講撒迦利亞書吧,就借此機會看看撒迦利亞書與將臨期的連繫.

先交代一下撒迦利亞書與耶穌的關連.關於耶穌的生平,若要在舊約找根據,大家最常想到的,是詩篇和以賽亞書,很少想到撒迦利亞書.其實,撒迦利亞書的內容,不少與耶穌的生平有關.

撒迦利亞書9:9寫:「看哪,你的王來到你這裡!...... 謙謙和和地騎著驢,就是騎著驢的駒子.」馬太福音引用這經文,來描繪耶穌騎驢進入耶路撒冷的情景.

撒迦利亞書11:13講到:將眾人估定的三十塊錢,丟給窯戶.這正是出賣耶穌的猶大後來所做的.

撒迦利亞書12:10寫:「他們必仰望我,就是他們所扎的」.約翰福音19:37引用了這經文,去描繪十架上的耶穌.

撒迦利亞書13:7寫:「擊打牧人,羊就分散」.這正是耶穌被捕時,門徒四散的光景.

按傳統講法,耶穌再來時,會站在橄欖山上,這講法出自撒迦利亞書14:4.

福音書的作者寫作福音書時,明顯意識到:先知撒迦利亞的信息應驗在耶穌身上.基督教聖經將撒迦利亞書放在舊約的結尾部分,是作為一橋樑,連接舊約和新約.它的作用,是讓那些在舊約結束時期百廢待興的神的子民,向前盼望,盼望新約時代開始時由耶穌降生所帶來的全新世界.換句話說,撒迦利亞書的信息,是站在舊約的結束時代,期盼新約的全新開始.在荒涼中,期待豐盛.在生活的失望中,盼望上帝創造新的意義.而這在黑暗中對光明的期待,正是「將臨期」的真意.

讓我們回到主前537年,王族後人所羅巴伯及大祭司約書亞從流放之地巴比倫回到耶路撒冷,眼前滿目瘡痍,百孔千瘡.要重建廢墟,他們選擇從重建一精神支柱開始.他們著眼於重建聖殿.大家會說,這麼多事情要幹,這麼多建築物要重建,幹嗎要耗費如此重要的有限資源,去搞這看來無迫切性的聖殿呢?

在絕望的人間,聖殿是重要的.聖殿象徵著上帝的臨在,代表著上帝生活在祂的子民中間,這其實就是聖誕節的意義.從約翰福音的神學來講,耶穌降生,他本身就是新的聖殿,是新的聖殿在人間,在他身上,我們看見上帝.約1:14:「道成了肉身,住在我們中間,充充滿滿地有恩典有真理」.「住」這個字,意思是搭起帳幕,有些英文聖經譯作「tabernacle」,直接點出耶穌降世,他本身就是以色列人出埃及時在曠野的會幕,這會幕是流動聖殿,象徵著上帝住在祂的子民中間,與他們同行.建聖殿,體會上帝的臨在是重要的,就像我們在耶穌這聖殿裡體會上帝的臨在一樣重要.對於一無所有的猶大人而言,聖殿象徵著上帝沒有撇棄他們,上帝仍愛他們,他們在上帝面前仍是上帝寶貴的兒女,他們在上帝那裡能迎向上帝應許的明天,在看來無望的黑暗歲月裡仍能盼望光明.重建聖殿,其實就是重建人生.

但當時的生活實在艱難,人心也渙散,也有很多客觀的阻礙,16年過去了,一切依舊,建殿的工程沒有寸進.

上帝先後差派先知哈該和撒迦利亞,去鼓勵群眾建聖殿.

上次,我講哈該書,先知哈該說:「這殿後來的榮耀必大過先前的榮耀；在這地方我必賜平安.這是萬軍之耶和華說的.」(該2:9)在殘垣敗瓦之上,以卑微的資源建起的聖殿,會比極榮華的所羅門王建的聖殿更美.這是先知給人的遠象.人要看到前景,才能有走下去的力量.

先知撒迦利亞和先知哈該同期,他們同時在耶路撒冷,鼓勵百姓建聖殿.建聖殿,意味著盼望上帝臨在人間,猶如我們盼望耶穌降生人間一樣,這正是「將臨期」所期盼的.在絕望的歲月,不能失去盼望.

撒迦利亞書上半部有8個異象,我們看看其中的第5個異象.

亞4:1　那與我說話的天使又來叫醒我,好像人睡覺被喚醒一樣.

亞4:2　他問我說:「你看見了什麼?」我說:「我看見了一個純金的燈臺,頂上有燈盞,燈臺上有七盞燈,每盞有七個管子.

亞4:3　旁邊有兩棵橄欖樹,一棵在燈盞的右邊,一棵在燈盞的左邊.」

亞4:4　我問與我說話的天使說:「主啊,這是什麼意思?」

亞4:5　與我說話的天使回答我說:「你不知道這是什麼意思嗎?」我說:「主啊,我不知道.」

眾人皆睡,唯獨先知被喚醒.他看見一個異象,大致是這樣:有一金燈臺,燈臺頂有一油盆,油盆有管子,將油注入燈臺上的七個燈杯上.頂上的油盆左右,各有一橄欖樹.

先知問天使:這是什麼意思呢?

天使反問:你不知道這是什麼意思嗎?

先知說:「我不知道」.若我在場,我會說「阿們」,我完全認同先知的講法,一點也不知道.

我們看看天使對這異象的解釋.

亞4:6　他對我說:「這是耶和華指示所羅巴伯的.萬軍之耶和華說:不是倚靠勢力(might),不是倚靠才能(power),乃是倚靠我的靈方能成事(原文沒有).

亞4:7　大山哪,你算什麼呢?在所羅巴伯面前,你必成為平地.他必搬出一塊石頭,安在殿頂上.人且大聲歡呼說:『願恩惠,恩惠歸與這殿!』」

亞4:8　耶和華的話又臨到我說:

亞4:9　「所羅巴伯的手立了這殿的根基,他的手也必完成這工,你就知道萬軍之耶和華差遣我到你們這裡來了.

亞4:10　誰藐視這日的事為小呢?這七眼(原文沒有)乃是耶和華的眼睛,遍察全地,見所羅巴伯手拿線鉈就歡喜.」

聽完了天使的講解,你明白他講什麼嗎?

我們可能跟先知再講一次:「我不知道」.

這些圖象:金燈臺,兩棵橄欖樹,七眼,都在啟示錄中重現.在文學體裁上,兩者都屬「天啟文學」,背後應有相近的象徵傳統.幸好,在啟示錄裡,這些象徵都解得明白.

「金燈臺」:在啟示錄象徵「教會」,即象徵神的子民群體.(參啟1:20)

「兩棵橄欖樹」:在啟示錄象徵「忠貞的僕人」.(參啟11:4)

「七眼」:在啟示錄象徵「七靈」,即聖靈.(啟5:6)

若以上的講法有效,則天使對撒迦利亞的講解便明白了.「金燈臺」是上帝的子民,若要發光,便需要倚靠上帝的靈.當時的處境是:所羅巴伯立了聖殿的根基,但不知何時完工.百姓根本輕忽建殿這件事情.「誰藐視這日的事為小呢?」(亞4:10)在《新譯本》譯作:「誰藐視這日子,以為所作的是小事呢?」百姓根本不放建殿在眼內.或許,他們以為:自己沒有勢力(might),也沒有能力(power),無權無勢,也沒有資源,怎能蓋得成呢?這個異象是提醒百姓,重點不在權勢,不必有權力迷思,不要以為擁有世上的力量才能成事,重點在上帝的靈.

有了上帝的靈,「大山」算什麼呢?要做好上帝的事工,總有很多「大山」的.天不時,地不利,人不和,寸步難行,萬念俱灰.天使說:只要上帝的靈在,「大山」也必成為平地.

天使說:只要人倚靠聖靈,排除萬難,聖殿總有平頂的一天,總有完工的一天.4:10的文字不太明晰,但意思是清楚的:那些藐視建殿,以為這只是不重要的小事的人,有一天,當看見聖殿完工,所羅巴伯拿著工具作最後測量時,他們必喜出望外.他們以為不會成就的事,最後竟然能完成,經文說,因為「耶和華的眼睛,遍察全地」,或者說,上帝的看顧,遍在人間.

在這裡,我花時間多講一點4:10這經文.經文中提及的「七」字(原文沒有「眼」字),應指上文金燈臺的七盞燈.金燈臺指上帝的子民,其能發亮,是因那七盞燈.4:10指這「七」(盞燈),就是「耶和華的眼睛」,或者簡單說,就是耶和華的看顧.後來在啟示錄,沿用這象徵,並直接用「七眼」這個詞,並點明這就是聖靈.將以上資料合起來,便明白天使講什麼了.天使說:上帝的子民(金燈臺)能發亮(七盞燈),是因聖靈的大能.建聖殿,是在人間點一枝燭光.能發光,全因上帝的看顧,聖靈的大能.

因上帝的看顧,靠著上帝的靈,人以為不可能建成的殿,最終會建成,標示著上帝臨在人間,廣施恩典.用聖誕節的情節來講,便是:因上帝的看顧,靠著上帝的靈感孕瑪利亞,人以為不可能的「上帝成為人」的事件,最終竟然發生,這標示著上帝臨到人間,廣施恩典.

我們來到這異象的最後部分:

亞4:11　我又問天使說:「這燈臺左右的兩棵橄欖樹是什麼意思?」

亞4:12　我二次問他說:「這兩根橄欖枝在兩個流出金色油(原文沒有)的金嘴旁邊是什麼意思?」

亞4:13　他對我說:「你不知道這是什麼意思嗎?」我說:「主啊,我不知道.」

亞4:14　他說:「這是兩個受膏者站在普天下主的旁邊.」

先知問:「這燈臺左右的兩棵橄欖樹是什麼意思?」

天使答:「這是兩個受膏者站在普天下主的旁邊.」

啟示錄承接著這象徵傳統,啟示錄裡的「兩棵橄欖樹」象徵「忠貞的僕人」.在舊約,「受膏者」一般的希伯來用字是「ma.shi.ach」,即「彌賽亞」這個字,但這裡卻刻意用了其他用字,直譯是「油之子」(sons of oil).這兩個人物,供應油給燈臺,即以聖靈之力支援上帝的子民.一般認為,這兩個人物就是王族的所羅巴伯及大祭司約書亞,他們在風雨飄搖的日子,以聖靈之力,鼓舞著他們的同胞.他們如此重要,百姓可能會誤以為他們便是民族等待已久的「彌賽亞」,救世主,經文刻意避開「彌賽亞」這用字,就是讓我們知道,「彌賽亞」另有其人.

來到這裡,經文為讀者留下一個問題:誰是我們真正等待的彌賽亞呢?他要兼備王者和祭司的身分,他的存在要發揮出聖殿的本意,讓人在他身上看到上帝在人間.他來到,要將這破敗的世界,重建為上帝居住的天地.要讓對前景無望的人,看到上帝賜予的光明前景.要讓罪孽深重的世人,看到新生命的可能.先知不知道這是誰,他也在等待.

先知撒迦利亞,站在舊約與新約的夾縫中,呼喚人建聖殿,也即呼喚人重視上帝的臨在,這是重建生命的基礎.他的信息隱約透露,他在等待一位彌賽亞,一位救主,這人將以聖靈之大能,讓苦難的人與上帝相遇,從而體會那全新的世界.先知在等待,那種「將臨期」式的等待.

生命無常,擁有的東西會轉眼成空,但請不要像那些在廢墟中的猶大人那樣說:我們無權無勢,而聖殿,信仰,上帝等東西虛無飄渺,毫無意義,我們注定要在這破敗的人間虛度一生.先知鼓勵我們,不要這樣說,不要這樣想.請靠賴聖靈,懷著希望,你會看到上帝在人間,向你施予恩典的.這是先知撒迦利亞給我們的信息,也是聖誕節的信息.

願神賜福你們.

願榮耀歸予聖父,聖子,聖靈.



\newpage

\section{你聽過嗎}
\label{sec:779}
\textbf{高國雄牧師 - 你聽過嗎?}
\newline
\newline
\hyperref[sec:778]{< < < PREV SERMON < < <}
~
\hyperref[sec:toc]{[返目錄]}
~
\hyperref[sec:780]{> > > NEXT SERMON > > >}


路加福音 2:1-14
\newline
\begin{longtable}{cl}
\hline
\hline
章節 & 經文 (和合本修訂版)\\
\hline
2:1 & \begin{tabularx}{0.7\textwidth}{X} 在那些日子，凱撒奧古斯都降旨，叫全國人民都登記戶籍。 \end{tabularx} \\ \\ \relax
2:2 & \begin{tabularx}{0.7\textwidth}{X} 這第一次登記戶籍是在居里扭作敘利亞總督的時候行的。 \end{tabularx} \\ \\ \relax
2:3 & \begin{tabularx}{0.7\textwidth}{X} 眾人各歸各城，辦理登記。 \end{tabularx} \\ \\ \relax
2:4 & \begin{tabularx}{0.7\textwidth}{X} 約瑟也從加利利的拿撒勒城上猶太去，到了大衛的城名叫伯利恆，因為他是大衛家族的人， \end{tabularx} \\ \\ \relax
2:5 & \begin{tabularx}{0.7\textwidth}{X} 要和他所聘之妻馬利亞一同登記戶籍。那時馬利亞已經懷孕。 \end{tabularx} \\ \\ \relax
2:6 & \begin{tabularx}{0.7\textwidth}{X} 他們在那裡的時候，馬利亞的產期到了， \end{tabularx} \\ \\ \relax
2:7 & \begin{tabularx}{0.7\textwidth}{X} 就生了頭胎的兒子，用布包起來，放在馬槽裡，因為客店裡沒有地方。 \end{tabularx} \\ \\ \relax
2:8 & \begin{tabularx}{0.7\textwidth}{X} 在伯利恆的野外有牧羊人，夜間值班看守羊群。 \end{tabularx} \\ \\ \relax
2:9 & \begin{tabularx}{0.7\textwidth}{X} 有主的一個使者站在他們旁邊，主的榮光四面照著他們，牧羊人就很懼怕。 \end{tabularx} \\ \\ \relax
2:10 & \begin{tabularx}{0.7\textwidth}{X} 那天使對他們說：「不要懼怕！看哪！因為我報給你們大喜的信息，是關乎萬民的： \end{tabularx} \\ \\ \relax
2:11 & \begin{tabularx}{0.7\textwidth}{X} 因今天在大衛的城裡，為你們生了救主，就是主基督。 \end{tabularx} \\ \\ \relax
2:12 & \begin{tabularx}{0.7\textwidth}{X} 你們要看見一個嬰孩，包著布，臥在馬槽裡，那就是給你們的記號。」 \end{tabularx} \\ \\ \relax
2:13 & \begin{tabularx}{0.7\textwidth}{X} 忽然，有一大隊天兵同那天使讚美神說： \end{tabularx} \\ \\ \relax
2:14 & \begin{tabularx}{0.7\textwidth}{X} 「在至高之處榮耀歸與神！在地上平安歸與他所喜悅的人！」 \end{tabularx} \\ \\
[1ex]
\hline
\hline
\end{longtable}
講題:你聽過嗎?Have You Heard?

經課:路加福音2:1-14

講員:高國雄牧師

場合:香港中文大學崇基學院禮拜堂平安夜燭光崇拜

日期:2025年12月24日

親愛的弟兄姊妹,平安夜平安!

在這個聖誕前夕,我們聚集在一起,慶祝主耶穌基督的降生.

今天,我想以一個簡單的問題開始:「你聽過嗎?Have you heard?」,這亦是今晚詩班所分享的主題.

你聽過那大喜的信息嗎?那關乎萬民的好消息——救主已經降生!

路加福音2章至14節:

「1在那些日子,凱撒奧古斯都降旨,叫全國人民都登記戶籍. 2這第一次登記戶籍是在居里扭作敘利亞總督的時候行的. 3眾人各歸各城,辦理登記. 4約瑟也從加利利的拿撒勒城上猶太去,到了大衛的城名叫伯利恆,因為他是大衛家族的人, 5要和他所聘之妻馬利亞一同登記戶籍.那時馬利亞已經懷孕. 6他們在那裡的時候,馬利亞的產期到了, 7就生了頭胎的兒子,用布包起來,放在馬槽裡,因為客店裡沒有地方.8在伯利恆的野外有牧羊人,夜間值班看守羊群. 9有主的一個使者站在他們旁邊,主的榮光四面照著他們,牧羊人就很懼怕. 10那天使對他們說:「不要懼怕!看哪!因為我報給你們大喜的信息,是關乎萬民的: 11因今天在大衛的城裡,為你們生了救主,就是主基督. 12你們要看見一個嬰孩,包著布,臥在馬槽裡,那就是給你們的記號.」 13忽然,有一大隊天兵同那天使讚美上帝說:14「在至高之處榮耀歸與上帝!在地上平安歸與他所喜悅的人!」

這段經文描述了耶穌降生時,牧羊人的經歷.他們是平凡人,在黑夜中按著吩咐與習慣看守羊群.忽然,天使向他們顯現,帶來大喜的信息:救主降生了!不是在宮殿裡,而是在馬槽中.

天使說:「不要懼怕!」

牧羊人聽見後,沒有猶豫,他們急忙去尋找,找到了嬰孩這記號,就把這好消息傳開.

你聽過嗎?這就是聖誕的核心:上帝差遣他的獨生子來到世上,為要拯救我們.

平安夜真的必然平安嗎?

想起1941年的香港平安夜,正是香港保衛戰最激烈的時刻.那年12月8日,日軍入侵香港,戰火迅速蔓延.到12月18日,日軍登陸香港島,激烈的戰鬥在淺水灣,黃泥涌等地展開.許多人失去家園,面對轟炸,饑餓和死亡.那個聖誕節,被稱為「黑色聖誕」,香港於12月25日淪陷,進入三年零八個月的艱難時期.

在戰爭的黑暗中,人們最需要什麼?

不是武器,不是權力,而是平安,希望和救贖.

就像當年的香港人一樣,他們在恐懼中盼望光明天使帶來的信息.

今年上月底大埔宏福苑的世紀火災,奪去了至少161條寶貴生命,包括一名殉職消防員,數十人受傷,數百家庭失去家園.火勢迅速蔓延七座大廈,許多長者,弱勢居民無法逃生.災難暴露了建築安全,維修工程的問題,也讓我們看見人的脆弱:在突如其來的災難中,我們需要救援,需要安慰,需要重建的希望.

在危難中人們最深刻的需要是什麼?

是安全,是被愛,是有人伸出援手,是在絕望中聽見一聲「不要懼怕」.

各位弟兄姊妹,各位朋友,你聽過嗎?

在這樣的黑暗時刻,上帝的信息依然響起:救主已經降生!耶穌來到世上,不是成為一般人認為的強者,而是軟弱的人；不是成為富足的人,而是成為貧窮,恐懼,痛苦的人.

牧羊人是社會底層,上帝先向他們宣告好消息,在平凡的牧羊日子中,上帝就向他們顯現,在恐懼驚惶中,大喜的信息就被天使宣告,耶穌就是那位好牧人,為羊捨命的大牧者.

在1941年的戰火中,上帝聽見了香港人的呼求；在今天,上帝也看見我們的眼淚.

耶穌降生,就是上帝親自來到我們中間,道成了肉身,為我們帶來真正的平安…不是世界暫時的和平,而是心靈的平安,與上帝,與自己和好的平安.

今天,在這個平安夜,你聽過嗎?

那大喜的信息依然為你而來:

不要懼怕!救主基督已經降生,他要拯救你脫離罪惡,恐懼和死亡.他要給你勇氣,力量及盼望.

親愛的弟兄姊妹,你聽過嗎?



\newpage

\section{天國大歷史之嬰孩耶穌逃命記}
\label{sec:780}
\textbf{林豪恩先生 - 天國大歷史之嬰孩耶穌逃命記}
\newline
\newline
\hyperref[sec:779]{< < < PREV SERMON < < <}
~
\hyperref[sec:toc]{[返目錄]}
~
\hyperref[sec:781]{> > > NEXT SERMON > > >}


馬太福音 2:13-23
\newline
\begin{longtable}{cl}
\hline
\hline
章節 & 經文 (和合本修訂版)\\
\hline
2:13 & \begin{tabularx}{0.7\textwidth}{X} 他們走後，忽然主的使者在約瑟夢中向他顯現，說：「起來！帶著小孩子和他母親逃往埃及，住在那裡，等我的指示；因為希律要搜尋那小孩子來殺害他。」 \end{tabularx} \\ \\ \relax
2:14 & \begin{tabularx}{0.7\textwidth}{X} 約瑟就起來，連夜帶著小孩子和他母親往埃及去， \end{tabularx} \\ \\ \relax
2:15 & \begin{tabularx}{0.7\textwidth}{X} 住在那裡，直到希律死了。這是要應驗主藉先知所說的話：「我從埃及召我的兒子出來。」 \end{tabularx} \\ \\ \relax
2:16 & \begin{tabularx}{0.7\textwidth}{X} 希律見自己被博學之士愚弄，極其憤怒，差人將伯利恆城裡和四境所有的男孩，根據他向博學之士仔細查問到的時間，凡兩歲以內的，都殺盡了。 \end{tabularx} \\ \\ \relax
2:17 & \begin{tabularx}{0.7\textwidth}{X} 這就應驗了耶利米先知所說的話： \end{tabularx} \\ \\ \relax
2:18 & \begin{tabularx}{0.7\textwidth}{X} 「在拉瑪聽見號咷大哭的聲音，是拉結哭她兒女；她不肯受安慰，因為他們都不在了。」 \end{tabularx} \\ \\ \relax
2:19 & \begin{tabularx}{0.7\textwidth}{X} 希律死了以後，在埃及，忽然主的使者在約瑟夢中向他顯現， \end{tabularx} \\ \\ \relax
2:20 & \begin{tabularx}{0.7\textwidth}{X} 說：「起來，帶著小孩子和他母親回以色列地去！因為要殺害這小孩子的人已經死了。」 \end{tabularx} \\ \\ \relax
2:21 & \begin{tabularx}{0.7\textwidth}{X} 約瑟就起來，帶著小孩子和他母親進入以色列地去。 \end{tabularx} \\ \\ \relax
2:22 & \begin{tabularx}{0.7\textwidth}{X} 但是他因聽見亞基老繼承他父親希律作了猶太王，怕到那裡去；又在夢中得到主的指示，就往加利利境內去了。 \end{tabularx} \\ \\ \relax
2:23 & \begin{tabularx}{0.7\textwidth}{X} 他們到了一座城，名叫拿撒勒，就住在那裡。這是要應驗先知所說的話：「他將稱為拿撒勒人。」 \end{tabularx} \\ \\
[1ex]
\hline
\hline
\end{longtable}
講題:天國大歷史之嬰孩耶穌逃命記 Baby Jesus on the Run

經課:馬太福音2:13-23

講員:高國雄牧師

場合:香港中文大學崇基學院禮拜堂主日崇拜

日期:2025年12月28日

引言

如果有選擇,你會選擇在甚麼時代出生呢?在怎樣的家庭出生呢?「小紅書上的一個帖子,掀起了網友極其熱烈的討論.帖子是一名新手媽媽寫的.她和丈夫經旅居國外多年,丈夫是美國加拿大華裔,最近,他們的女兒剛剛出生.他們給女兒準備了兩件禮物:

第一件是三個國家的護照.夫妻雙方都擁有幾個國籍,所以,女兒一出生就擁有了當「世界公民」的資格.第二件是他們將在這獨生女兒身上投注的巨大金錢和關愛.更令人震驚的是,她還給女兒安排了一套北京東城區的學區房.如果她長大了,不想留在美國,一回國就能在一線城市輕鬆擁有一套房.

而在帖子裡,最讓人心有戚戚的是她寫給女兒那一句話:「媽咪已經幫你體驗過了,確認了人世間是幸福的,快樂的,所以帶你來體驗一回.」

網友說:看到這些,我終於明白:這個世界的參差和溝壑,或許用一輩子的時間都難以填平.人和人之間,並不是生來就平等.有人一出生就在羅馬,而有的人從一出生開始,就需要朝著羅馬的方向前行.(\url{https://www.facebook.com/Wanderinghorizon/posts/})

而我們今天讀的福音經課介紹的那個人,他一出生就被羅馬的爪牙追殺.這個人是誰?就是耶穌.如果我們問耶穌,你為甚麼不出生在太平盛世,而出生在羅馬帝國高壓管治的年代呢?你為甚麼不出生在羅馬富貴家庭呢?為甚麼有些人一出世就要逃命?為甚麼有些人一出生就死在刀下?如果我們碰見馬太福音的作者,我們會問為甚麼聖誕不是一個溫馨的故事,而是一個逃命和眼淚的故事呢?

其實,每個時代的每個人都有很多「為甚麼」的問題,例如在鄭國江先生填詞的「為甚麼」這首歌,就一連問了幾個為甚麼:「為甚麼生世間上,此間許多哀與傷」／「為甚麼雙鬢斑白,光彩消失面容上」／「為甚麼竟會生病,輾轉反側在床上」.遺憾地,很多「為甚麼」都得不到滿意的答案.Steve Jobs曾經說過:「你無法預先把點點滴滴串連起來,只有在未來回顧時,你才會明白那些點點滴滴是如何串在一起的,所以你得相信,眼前你經歷的種種,將來多少會連結在一起.」.當我們讀到嬰孩耶穌的故事時,對眾多點點滴滴或許有很多問題,為甚麼這樣?為甚麼那麼?馬太福音講述這個故事的時候,建議讀者從更大的背景及更長遠的目光來看待嬰孩耶穌逃命事件.讓我們看看馬太福音如何把那些點點滴滴串連在一起,以及所產生的意義,和對我們的啟發.

一．「我從埃及召出我的兒子來.」

故事的第一部分是2:13-15節:

2:13他們去後,有主的使者向約瑟夢中顯現,說:「起來!帶著小孩子同他母親逃往埃及,住在那裡,等我吩咐你；因為希律必尋找小孩子,要除滅他.」

2:14約瑟就起來,夜間帶著小孩子和他母親往埃及去,

2:15住在那裡,直到希律死了.這是要應驗主藉先知所說的話,說:「我從埃及召出我的兒子來.」

當危險出現,人們的第一個反應是甚麼呢?相信不是問為甚麼,而是立即逃命,這是本能反應.逃難到何時呢?當然是等到危險消失之後.約瑟家庭面對的危險不是天災,也不是戰亂,而是一個人——希律王,因此,他們要和希律王「鬥長命」.俗語說,「你長命過佢,你就會贏佢」.如果有人問為甚麼要逃往埃及呢?這個問題可以用處境回答:第一,約瑟在夢中得到使者的指示；第二,從地緣政治來說,埃及是當時最近而又是希律王的法外之地.這樣的解釋合理,但卻不能滿足人心,因為沒有多少意義.正如人們問「為甚麼我會患癌病」,符合事實的回答是:「細胞異常增生,形成腫瘤」,這樣的答案未必能滿足人心.人類是一種要求意義的動物,而意義更是支撐著人類存活和文明的因素.嬰孩耶穌逃難這件事情,有甚麼意義呢?馬太福音的作者指出,其意義非常重大,影響非常深遠.作者引述何西亞書來解釋:

11:1以色列年幼的時候,我愛他,就從埃及召出我的兒子來.

馬太福音引導讀者從以色列出埃及事件來思考耶穌逃難往埃及的事件.當年以色列人在摩西的領導下,脫離埃及奴役,進入應許之地.這是以色列賴以為生的拯救故事,是以色列與耶和華建立關係的事件.如何看待嬰孩耶穌逃難往埃及的事件呢?馬太福音指出這是拯救的故事,而拯救的動機是愛.從歷史狀況來看,出生在高壓管治的希律王年代絕對不是甚麼好環境,遇上心狠手辣的希律王追殺,更加是苦難.這些都是事實,但不是意義本身,有些人是能夠從苦難中領悟意義,從而得到更新.最近,有一位在座的肢體傳來他的分享,我徵得他的同意在這裡讀出來.他吩咐我不要透露他的名字,但可以用「一位高大衛猛的弟兄」相稱.

「2025年整個三月奔波於各醫院診所,進行檢查,抽取組織,最終面對確診鼻咽癌第三期的事實.我內心充滿徬徨及無助,死亡突然如此接近.

緊接著的是長達半年,如同一場艱苦戰役的療程:先接受九週六次化療,再承受連續33次令我人生第一次食嗎啡止痛的電療.化療電療帶來的極度疲憊,白血球驟降,食慾大減,讓我大部分時間只能臥床,與世界隔絕.

然而,正是在這看似無盡的幽谷中,我經歷了信仰最真實的淬煉與照亮.聖經信息成了我屬靈的嗎哪和力量泉源.

我首先翻開《約伯記》.約伯的掙扎與追問——為何義人遭難?上帝在旋風中的回應,並非解答苦難的原因,上帝展現祂創造的奧秘和主權,我們無知暗昧根本完全無能力去明白上帝的創造和旨意.那段經文讓我深刻體會到自己的有限與渺小,正如我無法明白苦難為何臨到,但我可以確信上帝的同在,祂就在我身邊,清楚知道我的遭遇.「我從前風聞有你,現在親眼看見你.」(約伯記42:5)我選擇信靠祂與我同行,而不是追問為何有患難...

當我讀到《彼得前書》五章8節,感到如警鐘響起:「務要謹守,警醒.因為你們的仇敵魔鬼,如同吼叫的獅子,遍地遊行,尋找可吞吃的人.」撒旦不是挑戰約伯的信心從而令他遠離上帝嗎?它也想藉疾病吞噬我的信心與盼望.這反而激勵我更要振作,緊緊倚靠聖靈,不讓自己成為撒旦的掠物.相反,我要透過這經歷,更認識神,更敬畏祂,讓生命淬煉出更深的意義.

如今,我滿心感恩.最近一次的鼻咽組織抽取及頸部MRI檢查,結果顯示沒有癌細胞殘留和進展良好.現階段,我正學習適應治療後的各種辛苦的後遺症,並每天服用為期一年的輔助化療藥丸,以降低復發風險.

我不知我生命還有多久,但想花一點時間以陪伴者,同行者的身份,將這份「化妝的祝福」傳遞下去,報答上帝深不可測的恩情.我永遠不會忘記祂在患難中與我同行,醫治與更新.從恐懼走向平安,從絕望走向盼望,從疑問走向更深的信靠.這一切,只因為在患難中憑著單一的信心遇見祂.」

為甚麼是我遇到這些事情?為甚麼耶穌一出生就要逃難?馬太福音啟發我們不要只把它視為一個苦難的故事,這其實是一個拯救的故事.

二．「在拉瑪聽見號咷大哭的聲音」

故事的第一節結束在2:15「住在那裡,直到希律死了.」順理成章接下去應該就是2:19「希律死了以後,有主的使者在埃及向約瑟夢中顯現」.然而,中間卻有16-18節.這一段有甚麼意思呢?作者可能想告訴我們,事件就算有意義,也不能抹殺當事人痛苦的經驗.

2:16希律見自己被博士愚弄,就大大發怒,差人將伯利恆城裡並四境所有的男孩,照著他向博士仔細查問的時候,凡兩歲以裡的,都殺盡了.

馬太福音如何表達那種慘絕人寰的悲痛呢?2:18節引述耶利米書:

31:15在拉瑪聽見號咷痛哭的聲音,是拉結哭她兒女,不肯受安慰,因為他們都不在了.

拉結是雅各的愛妻,象徵以色列十二支派的母親,她在這裡被描繪為為失去的兒女而悲傷.失去的兒女是誰呢?就是被巴比倫滅國,聖靈塗炭,哀鴻遍野的慘況.希律王屠殺男孩帶來的慘況,可以怎樣比擬呢?就如耶路撒冷被攻破,人民被殺被擄一樣的悲慘和痛苦.

面對如此悲慘的苦難,有甚麼可以說呢?有甚麼說話能夠帶來安慰呢?馬太福音在這裡的回答是:沒有.讓號咷的號咷吧!讓痛哭的痛哭吧!不要試圖解釋掉別人的痛苦,不要試圖阻止別人哀悼,更不要粉飾慘案.

面對傷痛,很多人都想說一些安慰的話.有善終服務人士列出「十句最廢的安慰說話」(侯偉新,伍桂麟,梁梓敦)參人們參考:

1. 節哀順變.

2. 時間可以沖淡一切.

3. 時間都過了這麼久,還想來做什麼.

4. 這都是神的安排.

5. 你看看那個人的情況比你更慘,其實你已經不太差.

6. 其實他離開了對你來說是好事,你不用再這麼辛苦照顧他,對你也是解脫.

7. 你還年輕,再找個更好的吧.

8. 這樣更好啦,他對你們這麼差,死了便一了百了.

9. 我很明白你的感受.

10. 我都試過你的經歷啦,好快無事!

他們建議,如果真的無話可說,不如安靜的陪伴.

「苦難是化妝的祝福」這句話,只能夠由當事人經過深刻的體會後心悅誠服地自己說出來,而不能夠由其他人代替受苦者說,否則就是二次傷害.痛是痛在受傷者身上的,苦是苦在受苦者心裡的,這是真實的.嬰孩耶穌逃難事件的記述,為甚麼插入2:16-18節這小段呢?會不會是馬太福音認為縱使耶穌逃亡往埃及有重大意義,但也不能抹殺希律王的暴行帶來給人間的悲痛呢?

三．「他將稱為拿撒勒人」

記下了希律王的暴行,哀悼過被屠殺的兒童之後,回到嬰兒耶穌逃難埃及的故事.

有一個好消息和一個壞消息,你想先聽哪一個?好消息是:暴君必有死去的一日.

2:19希律死了以後,有主的使者在埃及向約瑟夢中顯現,說:

2:20「起來!帶著小孩子和他母親往以色列地去,因為要害小孩子性命的人已經死了.」

2:21約瑟就起來,把小孩子和他母親帶到以色列地去；

壞消息是:接任的人比上一任的人更差.

2:22只因聽見亞基老接著他父親希律作了猶太王,就怕往那裡去,又在夢中被主指示,便往加利利境內去了,

但壞消息裡仍有好消息,就是接任的亞基老能力不及他的父王,未能承繼大希律的全部領土,因此,人們只要「行遠一些」就可以避開他的暴政.後來,耶穌長大成人往耶路撒冷傳道的時候,再沒有提亞基老了.為甚麼呢?因為這個人實在太過瘋狂而且無能,連羅馬都「頂佢唔順」,流放了他,直接派巡撫彼拉多來管治.

一如既往,馬太福音不會認為事情完全是處境使然,耶穌事件也不是純粹處境的產物,而是有貫穿歷史繼往開來的意義:

2:23到了一座城,名叫拿撒勒,就住在那裡.這是要應驗先知所說,他將稱為拿撒勒人的話了.

聖經學者翻查整本舊約,都找不到有那一位先知說過這句話.最接近的可能是將臨期第二主日讀過的以賽亞書

11:1從耶西的殘幹必長出嫩枝(和修)

有學者認為嫩枝的讀音拿撒勒的讀音相似,因此而產生彌賽亞的連繫.不過,在耶穌的年代被稱為拿撒勒人有甚麼好處嗎?沒有.還記得約翰福音記載拿但業如何看待拿撒勒這個地方嗎?

1:46拿但業對他說:「拿撒勒還能出甚麼好的嗎?」

拿撒勒是一個偏僻,貧窮,不起眼的小村莊,因名不見經傳而受到輕視.如果耶路撒冷人等於城市人,拿撒勒人即是「鄉下仔」.就算拿撒勒的讀音令人聯想到彌賽亞,但為甚麼彌賽亞要做被看扁的「鄉下仔」呢?今天書信經課的希伯來書如此提示我們:

2:17所以,他凡事該與他的弟兄相同,為要在上帝的事上成為慈悲忠信的大祭司,為百姓的罪獻上挽回祭.

我們曾經問過:耶穌為甚麼出生在那個時代?為甚麼遭遇那些事情?為甚麼在那種環境成長?其實,這些不單是讀耶穌故事時的問題,也往往是我們對著自己的際遇所問的問題.對於我們來說,出生在甚麼年代,遭遇甚麼事情,在甚麼環境成長都未必是我們可以控制.耶穌呢?如果祂有選擇,為甚麼仍如此呢?希伯來書告訴我們耶穌凡事與他的弟兄相同,「他自己既然被試探而受苦,就能搭救被試探的人」(2:18).



\newpage

\section{因為我要使他們的悲哀變為歡喜}
\label{sec:781}
\textbf{李均熊牧師 - 因為我要使他們的悲哀變為歡喜}
\newline
\newline
\hyperref[sec:780]{< < < PREV SERMON < < <}
~
\hyperref[sec:toc]{[返目錄]}
~
\hyperref[sec:782]{> > > NEXT SERMON > > >}


耶利米書 31:7-14
\newline
\begin{longtable}{cl}
\hline
\hline
章節 & 經文 (和合本修訂版)\\
\hline
31:7 & \begin{tabularx}{0.7\textwidth}{X} 耶和華如此說：「你們當為雅各歡樂歌唱，為萬國中為首的歡呼。當傳揚，頌讚說：『耶和華啊，求你拯救你的百姓，拯救以色列的餘民。』 \end{tabularx} \\ \\ \relax
31:8 & \begin{tabularx}{0.7\textwidth}{X} 看哪，我必將他們從北方之地領來，從地極召集而來；同他們來的有盲人、瘸子、孕婦、產婦；他們必成群結隊回到這裡。 \end{tabularx} \\ \\ \relax
31:9 & \begin{tabularx}{0.7\textwidth}{X} 他們要哭泣而來。我要照他們懇求的引導他們，使他們在河水旁行走正直的路，他們在其上必不致絆跌；因為我是以色列的父，以法蓮是我的長子。 \end{tabularx} \\ \\ \relax
31:10 & \begin{tabularx}{0.7\textwidth}{X} 列國啊，要聽耶和華的話，要在遠方的海島傳揚，說：「趕散以色列的必召集他，看守他，如牧人看守羊群。」 \end{tabularx} \\ \\ \relax
31:11 & \begin{tabularx}{0.7\textwidth}{X} 因為耶和華救贖了雅各，救贖他脫離比他更強之人的手。 \end{tabularx} \\ \\ \relax
31:12 & \begin{tabularx}{0.7\textwidth}{X} 他們來到錫安的高處歌唱，因耶和華的宏恩而喜樂洋溢，就是五穀、新酒和新的油，並羔羊和牛犢。他們必像有水澆灌的園子，一點也不再有愁煩。 \end{tabularx} \\ \\ \relax
31:13 & \begin{tabularx}{0.7\textwidth}{X} 那時，少女必歡樂跳舞；年輕的、年老的，都一同歡樂；因為我要使他們的悲哀變為歡喜，並要安慰他們，使他們的愁煩轉為喜樂。 \end{tabularx} \\ \\ \relax
31:14 & \begin{tabularx}{0.7\textwidth}{X} 我必以肥油使祭司的心滿足，我的百姓也要因我的恩惠知足。這是耶和華說的。 \end{tabularx} \\ \\
[1ex]
\hline
\hline
\end{longtable}
講題:因為我要使他們的悲哀變為歡喜 For I Will Turn Their Mourning into Joy

經課:耶利米書31:7-14

講員:李均熊牧師

場合:香港中文大學崇基學院禮拜堂主日崇拜

日期:2026年1月4日

起:長夜仍是漆黑

耶利米書這一段說話,並非在以色列人看到曙光,勝利在望,才會叫人「當為雅各歡樂歌唱」,又或「要來到錫安的高處歌唱.」有一些學者會將耶利米書的這一段經文與以賽亞書40-55章作比較,因為當中經文也叫人要歡呼歌唱,例如賽52:9-10「發起歡聲,一同歌唱；因為耶和華安慰了他的百姓,救贖了耶路撒冷.耶和華在萬國眼前露出聖臂；地極的人都看見我們神的救恩了.」但細看以賽亞書,當中賽45:1提到波斯大帝古列／塞魯士,整個段落一開始(賽40:2),便說到要向耶路撒冷宣告她爭戰的日子滿了.明顯神子民的歷史正進入一個新階段,我們可以推斷,以賽亞書40-55章的內容是波斯帝國正大國崛起,巴比倫國勢卻已江河日下.所以宣告神的拯救和釋放已近於咫尺,人心應該從灰心失望中振奮起來,是可以理解的.

但今日經課的耶利米書,我們若仔細查考,就會發覺這一段經文很大可能是南國猶大剛經歷巴比倫大軍的攻擊,甚至可能在耶路撒冷陷落,以色列人成為戰俘,在被擄的日子當中宣告的.所以今日我們的舊約經課與以賽亞書40-55章的歡樂歌唱是截然不同的處境.以賽亞書是黎明漸現,耶利米書則正進入漆黑長夜.

怎樣可以在黑暗的時候,仍去宣告歡呼頌讚,歡樂跳舞呢?其實經文並沒有無視以色列人被擄受苦的事實.大家可以留意一下,當中可以回歸的「剩下的人」是誰.瞎子,瘸子,孕婦,產婦,他們可能正因為身體的狀況,反而能活命過來；處女,年少的和年老的,不是反映社會已經沒有了精壯男丁了嗎?31:9說到他們要哭泣而來,那究竟是歡喜的眼淚,還是困苦悲哀的眼淚呢?

承:信心不是盲目樂觀

經文沒有否認現實的痛苦.耶利米清楚知道百姓正面對失去國家,失去家園,失去盼望的處境.31:8提到「瞎子,瘸子,孕婦,產婦」——這些人是社會中最脆弱的群體,卻成為「剩下的人」.這反映出國家已經失去了精壯的男丁,社會結構崩潰,前景黯淡.31:9說「他們要哭泣而來」,這不是勝利的歡呼,而是帶著眼淚的歸回.

然而,神在這裡呼召百姓「要為雅各歡樂歌唱」,不是因為現實已經改變,而是因為神的應許是真實的.信心不是否認黑暗,而是在黑暗中仍然選擇相信神的話.

轉:神的身份與作為

為什麼可以在絕望中仍有盼望?因為神不僅是「趕散以色列的主」,祂也是「招聚他們的主」(31:10).祂是以色列的父,也是看守祂的羊,疼愛羊群的牧羊人.這是一個極大的安慰:即使百姓因罪受罰,被擄流離,神仍然認他們為子民,為兒子,仍然看顧他們.31:9說「我要照他們懇求的引導他們,使他們在河水旁走正直的路,在其上不致絆跌」；上主「又看守他,好像牧人看守羊群.」牧人不會丟棄羊群,反而在危險中更加警醒保護.神的身份和作為,愛與信實,成為我們在黑暗中仍能盼望的根基.弟兄姊妹,今天我們也可能面對失望,困境,甚至看不到出路.但今日經課的信息告訴我們:信心不是等到看到光明才開始,而是在黑暗中,仍然相信神的應許.在黑暗中,耶利米提醒我們要記得上主是我們的父,是牧人,是看守我們的主.

合:悲哀變為歡喜

經文最後描繪了一個美麗的圖畫:百姓要回到錫安的高處歌唱,神要供應他們的所需,甚至「他們的心必像澆灌的園子」(31:12).神要使用不能自救的瞎子,瘸子,孕婦,產婦,那些剩下的人,回歸上主的應許地,透過軟弱無力的處女,年少的,年老的,重新開闢上主施恩之地,使眼前被戰爭蹂躪的地土,重新充滿生機,成為有五穀,新酒,和油,並羊羔,牛犢之地.這不是空洞的安慰,而是神親自應許:「我要使他們的悲哀變為歡喜,使他們的愁煩變為安慰.」(31:13)

當我預備講章時,很長一段時間都找不到要傳講的信息,因為經文提到的悲哀,都令我聯想到我們新近經歷的悲哀,就是宏福苑的大火災.上個月講道,與災情更加接近的時候,我心仍對神充滿疑惑,對世事充滿悲哀和忿怒,所以大家可能會覺得我上一篇道想回應一下,卻又不知所云.

但當我不斷細讀今天的講道經文,神好像要我留意的是那些充滿著需要的「瞎子,瘸子,孕婦,產婦」,和那些軟弱和缺乏能力的「處女,少年的和年老的」.在悲哀中我發覺我關心了錯的對象.我期望上主施展大能的手,去使到令世界發生災禍的不公不義的事情消失,以為只要有強而有力的「壯丁」,世界才會復興改變(就像以色列人期望一位強人彌賽亞去解放他們)；但上主其實要我們注目的是身處悲哀,充滿需要又軟弱無能的人身上.

在大火發生之後,是甚麼令我最受觸動?是市民無條件,快速又有效的對災民的協助.

今日的經文還要我們注目留意那些需要人幫助的和軟弱的人,因為神的施恩就是從他們開始,甚至可以說,神的應許也是由他們去實現!

這提醒我們,神的拯救雖然未必立刻臨到,但祂的應許必不落空.今天我們可能仍在黑暗中,仍在流淚,但我們可以選擇在信心中讚美,因為神必成就祂的話:祂必使悲哀變為歡喜,愁煩變為安慰.



\newpage

\section{盡諸般的義}
\label{sec:782}
\textbf{高國雄牧師-盡諸般的義}
\newline
\newline
\hyperref[sec:781]{< < < PREV SERMON < < <}
~
\hyperref[sec:toc]{[返目錄]}
~
\hyperref[sec:783]{> > > NEXT SERMON > > >}


馬太福音 3:13-17
\newline
\begin{longtable}{cl}
\hline
\hline
章節 & 經文 (和合本修訂版)\\
\hline
3:13 & \begin{tabularx}{0.7\textwidth}{X} 當時，耶穌從加利利來到約旦河，到了約翰那裡，請約翰為他施洗。 \end{tabularx} \\ \\ \relax
3:14 & \begin{tabularx}{0.7\textwidth}{X} 約翰想要阻止他，說：「我應該受你的洗，你怎麼到我這裡來呢？」 \end{tabularx} \\ \\ \relax
3:15 & \begin{tabularx}{0.7\textwidth}{X} 耶穌回答他：「暫且這樣做吧，因為我們理當這樣履行全部的義。」於是約翰就依了他。 \end{tabularx} \\ \\ \relax
3:16 & \begin{tabularx}{0.7\textwidth}{X} 耶穌受了洗，隨即從水裡上來。天忽然為他開了，他看見神的靈降下，彷彿鴿子落在他身上。 \end{tabularx} \\ \\ \relax
3:17 & \begin{tabularx}{0.7\textwidth}{X} 這時，天上有聲音說：「這是我的愛子，我所喜愛的。」 \end{tabularx} \\ \\
[1ex]
\hline
\hline
\end{longtable}
講題:盡諸般的義ToFulfillAllRighteousness

經課:馬太福音3:13-17

講員:高國雄牧師

場合:香港中文大學崇基學院禮拜堂主日崇拜

日期:2026年1月11日

弟兄姊妹,今天我們的思想焦點,是主耶穌在約旦河邊對約翰說的一句話:

「你暫且許我,因為我們理當這樣盡諸般的義.」(馬太福音3:15)

這「理當」不是可有可無的選擇,而是「必須如此」——完全合乎上帝心意的義.耶穌一生都在實踐這句話,從馬槽到十字架,祂的每一步都是對天父完全的順服.

對於我們,「盡諸般的義」如何實踐呢?是否就是按一般人的期望而完成某些該做的事?

一,聖誕慶祝:不止於節日,更要記念道成肉身的順服

每年12月,世界熱鬧慶祝聖誕:燈飾,禮物,派對,聖誕歌,教會也舉辦盛大活動,甚至去年出現中國版的聖誕老人…

但若我們只停留在「慶祝」的層面,是否就盡了諸般的義?

聖誕的核心,是記念道成肉身的耶穌降生,為要拯救罪人.

約翰福音1:14說:「道成了肉身,住在我們中間,充充滿滿地有恩典有真理.」

腓立比書2:6-8耶穌「虛己,取了奴僕的形像……存心順服,以至於死,且死在十字架上」.

耶穌降生本身,就是順服的起點.若我們的聖誕只帶來年度情感的高峰,卻沒有提醒我們讓基督在我們裡面繼續「虛己」——在工作,家庭,教會決策中活出順服——那麼我們的慶祝仍欠缺深度.

二,參與教會:不止於出席,更要與罪人認同

許多人想:「我每主日聚會,奉獻,服事,這就夠了吧?」參與教會當然寶貴,但若只是「參與」,是否就盡了諸般的義?

約翰的洗是「悔改的洗」(馬太福音3:11),本為罪人預備.主耶穌明明無罪,卻站在罪人的隊伍裡,受這洗禮——不是因為祂需要潔淨,而是要與罪人認同,完全順服天父的救贖計劃.

教會正是這認同的延續.我們來教會,不是來一個機構消費屬靈資源,而是進入基督與罪人認同的聖約團體:在這裡一同承認自己是罪人,一同領聖餐記念祂的死,一同在軟弱肢體身上看見基督的臉.

若我們來教會多年,卻仍帶著「我來貢獻,我來得造就」的消費者心態,而非讓基督在我們裡面繼續祂的順服,那麼即使忠心服事,我們仍未摸著「盡諸般的義」的深度.真正的教會參與,是讓我們不斷被「洗」,與基督同死同復活的過程.當我們如此順服,上帝的同在與喜悅必在我們中間彰顯.

三,領受洗禮:不止於儀式,更是順服的起點與終生錨點

耶穌基督的洗禮,很自然地就會想起我們的洗禮.

還記得你當日如何接受洗禮嗎?

對已受洗的信徒而言,洗禮絕非一次過的儀式,而是我們一生不斷回歸的錨點.每當我們軟弱,犯罪,偏離,都要回到洗禮的恩典:我們已在基督裡與祂同死,同埋葬,同復活,祂的義已完全歸給我們.

若受洗多年後,仍靠行為,經驗,服事年資建立自己的義,那就是保羅在腓立比書3:9所拒絕的「自己因律法而得的義」.真正的義,是「本於信,在基督耶穌裡的義」.洗禮不是終點,而是公開宣告:我願一生跟隨主,讓聖靈掌管,結出聖靈的果子.

四,做個好人:不止於道德,更是讓基督在我們裡面活出祂的義

最後,有人想:「我做個好人,不做壞事,幫助別人,守規矩,這就夠了吧?」在世人眼中,這值得稱讚；但聖經說:「我們都像不潔淨的人；所有的義都像污穢的衣服.」(以賽亞書64:6)

主耶穌是史上「最好的人」——行神蹟,教導真理,醫治病人,從不犯罪——卻仍來受洗,因為單靠「做好人」無法盡諸般的義.真正的義,只有從基督而來:

哥林多後書5:21「上帝使那無罪的,替我們成為罪,好叫我們在他裡面成為上帝的義.」

成聖不是「努力做好人」,乃是越來越讓基督在我們裡面活(加拉太書2:20).對資深信徒而言,最大的危險不是犯罪,而是自義——以為自己的好行為能讓上帝更喜悅.真正的深度,是越來越看見自己的無有,越來越單純倚靠基督已成就的義.

結語:從「我理當」到「這是我所喜悅的」

弟兄姊妹,慶祝聖誕,參與教會,領受洗禮,做個好人——這些都不是錯的,但若停留在表面,我們仍未盡諸般的義.

主耶穌在約旦河給我們立下榜樣:祂無罪,卻完全順服；祂不需要,卻仍成就上帝的一切旨意.受洗後,天開了,聖靈降下,天父說:「這是我的愛子,我所喜悅的.」

今天,上帝也渴望對我們說同樣的話.當我們在生活每一件事上——家庭,工作,教會,日常決定——都像耶穌一樣說:「我們理當這樣盡諸般的義」,聖靈必運行,天父必喜悅.

我邀請你現在默想:你的生命中,哪一部分還在抗拒上帝的引導?你是否願意說:「主啊,我理當這樣行」?

若你仍在信仰探索當中,今天就是接受祂的時候；若你已信主多年,今天就是重新委身,讓基督在你裡面繼續祂順服的時候.

在新的一年開始之際,讓我們再次來到約旦河邊,聽主對我們說:「我們理當這樣盡諸般的義.」這「我們」包括祂,也包括你和我.祂已為我們盡了,我們如今在祂裡面,一同活出這義——不是靠自己,乃是讓祂在我們裡面繼續祂完全的順服.

願天為我們開啟,願我們常常聽見天父的聲音:「這是我的愛子,我所喜悅的.」



\newpage

\section{你們用話頂撞我}
\label{sec:783}
\textbf{鄧瑞強牧師 - 你們用話頂撞我}
\newline
\newline
\hyperref[sec:782]{< < < PREV SERMON < < <}
~
\hyperref[sec:toc]{[返目錄]}
~
\hyperref[sec:784]{> > > NEXT SERMON > > >}


瑪拉基書 3:13-4:6
\newline
\begin{longtable}{cl}
\hline
\hline
章節 & 經文 (和合本修訂版)\\
\hline
3:13 & \begin{tabularx}{0.7\textwidth}{X} 耶和華說：「你們用話頂撞我。」你們卻說：「我們說了甚麼話頂撞你呢？」 \end{tabularx} \\ \\ \relax
3:14 & \begin{tabularx}{0.7\textwidth}{X} 你們說：「事奉神是枉然，我們遵守神所吩咐的，在萬軍之耶和華面前哀痛而行，有甚麼益處呢？ \end{tabularx} \\ \\ \relax
3:15 & \begin{tabularx}{0.7\textwidth}{X} 現在，我們稱狂傲的人為有福，並且行惡的人得以建立；他們雖然試探神，卻得以逃脫。」 \end{tabularx} \\ \\ \relax
3:16 & \begin{tabularx}{0.7\textwidth}{X} 那時，敬畏耶和華的人彼此談論，耶和華側耳而聽，且有紀念冊在他面前，記錄那敬畏耶和華、思念他名的人。 \end{tabularx} \\ \\ \relax
3:17 & \begin{tabularx}{0.7\textwidth}{X} 萬軍之耶和華說：「在我所定的日子，他們必屬我，是我寶貴的產業。我必憐憫他們，如同人憐憫那服侍他的兒子。 \end{tabularx} \\ \\ \relax
3:18 & \begin{tabularx}{0.7\textwidth}{X} 那時你們必再一次看出義人和惡人，事奉神和不事奉神的人有何差別。」 \end{tabularx} \\ \\ \relax
4:1 & \begin{tabularx}{0.7\textwidth}{X} 萬軍之耶和華說：「看哪，那日臨近，勢如燒著的火爐，凡狂傲的和行惡的都如碎秸，在那日被燒盡，根與枝條無一存留。 \end{tabularx} \\ \\ \relax
4:2 & \begin{tabularx}{0.7\textwidth}{X} 但是，對你們敬畏我名的人，必有公義的太陽出現，其光線有醫治的能力。你們必出來跳躍如圈裡的牛犢。 \end{tabularx} \\ \\ \relax
4:3 & \begin{tabularx}{0.7\textwidth}{X} 你們必踐踏惡人；在我所定的日子，他們必成為你們腳掌下的灰塵。這是萬軍之耶和華說的。 \end{tabularx} \\ \\ \relax
4:4 & \begin{tabularx}{0.7\textwidth}{X} 「你們當記念我僕人摩西的律法，就是我在何烈山為以色列眾人所吩咐他的律例典章。 \end{tabularx} \\ \\ \relax
4:5 & \begin{tabularx}{0.7\textwidth}{X} 「看哪，耶和華大而可畏之日未到以前，我要差遣以利亞先知到你們那裡去。 \end{tabularx} \\ \\ \relax
4:6 & \begin{tabularx}{0.7\textwidth}{X} 他必使父親的心轉向兒女，兒女的心轉向父親，免得我來詛咒這地。」 \end{tabularx} \\ \\
[1ex]
\hline
\hline
\end{longtable}
講題:你們用話頂撞我 Your Words have been Harsh against Me

經課:瑪拉基書3:13-4:6

講員:鄧瑞強牧師

場合:香港中文大學崇基學院禮拜堂主日崇拜

日期:2026年1月18日

各位弟兄姊妹,早安.願神的恩典常與你們同在.

今日講道的經文出自瑪拉基書,若大家用實體聖經的話,便能直接感知這卷書位於舊約的最尾部分,下一卷書便是新約的馬太福音了.但我們現在很少觸摸聖經,漸漸失去對聖經的直接感覺.我們現在習慣用手機版的聖經,或看禮堂熒幕上powerpoint打出來的經文,再過多一些日子,新一代的信徒或許不知道聖經有多厚,也對哪卷書排在哪卷書之前,哪卷書跟在哪卷書之後再沒有任何直覺意識,這是閱讀聖經經驗的碎片化.新一代年輕信徒接觸的經文,多是科技切割出來的一小段經文,脫離了上下文,這使人更難對聖經有一氣呵成的理解；對於神與祂的子民同行了過千年的歷史記載,也難有宏觀的認識.這碎片化的經驗,對新一代信徒的靈性塑造,或會有深遠影響.

從瑪拉基書到馬太福音,是從舊約時代過渡到新約時代,這段時期一般稱為「兩約中間時期」,為時約400年.先知瑪拉基之後,便再沒有列入正典的先知了.在這沉寂的400年裡,神再沒有藉先知發言,神的子民在漫長地等待,直到黎明初現,基督降臨.

瑪拉基書沒有標示寫作日期,但大致與尼希米同期,約於主前5世紀中期.猶大亡國,已是130多年前的往事了.耶路撒冷的聖殿重建完成,到瑪拉基的時代已有60多年了.至於瑪拉基先知是早於抑或晚於尼希米建造耶路撒冷的城牆,就難以確定了.按經文所反映的情況,當時經濟低迷,糧食失收,人心虛怯.神的子民在各方面都對神不滿,他們看不到神有任何眷顧.曾經,他們可能響應尼希米的號召,從巴比倫返回猶大,心想會有好日子過,怎知大失所望,返到故鄉,原來百廢待興,生活比在巴比倫時更艱難.

瑪拉基書全書都是人心底的牢騷,神聆聽,神一一回應.不計開頭和結語,瑪拉基書大分六段,如此,就有六個大牢騷.今日,我們看最後一個牢騷.這牢騷寫在瑪拉基書的結尾,也就是在基督教聖經舊約部分的結尾,這牢騷可以看成是神的子民對千年以來神與他們同行的最後結局的一個大悲歎.暫時不論人當負什麼責任,對當時的神的子民來看,從先祖亞伯拉罕到當下,如此悠長歲月,他們相信神,當然中間有種種曲折,但他們畢竟仍算是相信神的,弄到最後,竟落得如此下場.他們發出千古一嘆.

我們看看在舊約聖經裡,人對神發出的最後投訴.

瑪3:13　耶和華說:「你們用話頂撞我,你們還說:『我們用什麼話頂撞了你呢?』

瑪3:14　你們說:『事奉上帝是徒然的,遵守上帝所吩咐的,在萬軍之耶和華面前苦苦齋戒,有什麼益處呢?

瑪3:15　如今我們稱狂傲的人為有福,並且行惡的人得建立；他們雖然試探上帝,卻得脫離災難.』」

神講出那些人的心裡話,他們用好「硬淨」的話來頂撞神.他們心裡說什麼呢?用香港口語來說,就是:「神都好廢,服事祂的,得個桔,你看,我們如今淪落到什麼田地!反而那些大壞蛋,無法無天的人,他們撈到風生水起.」那些生活艱難的猶大人,他們藐視神,卻讚賞那些心中無神的大壞蛋,藉此表達對神的最大不滿.其實,他們的心情也不難理解,若我們信耶穌信了一輩子,又敬虔,又努力事奉,到晚年,卻過著像難民一樣的生活,我們可能也會「發爛渣」,頂撞神.

神如何回應呢?

瑪3:16　那時,敬畏耶和華的彼此談論,耶和華側耳而聽,且有紀念冊在他面前,記錄那敬畏耶和華,思念他名的人.

瑪3:17　萬軍之耶和華說:「在我所定的日子,他們必屬我,特特歸我.我必憐恤他們,如同人憐恤服事自己的兒子.

瑪3:18　那時你們必歸回,將善人和惡人,事奉上帝的和不事奉上帝的,分別出來.」

在百廢待興的耶路撒冷裡,也不是沒有敬虔的人的.這裡,先知以神對待這些敬虔人的態度,側面回答不信者的疑惑.先知說,敬虔者的談論,神是用心傾聽的.神用記事簿記下他們的言談,言之之意,神真心留意敬虔者的生活細節.神記下一切,不是用來提醒自己,而是讓敬虔者明白,神多麼看重他們.現代的父母都為孩童時代的兒女拍很多相片,去過哪裡,吃過什麼,都留下記錄.這是給兒女們長大後看的,讓他們明白,父母多麼愛惜他們.神的記錄,也是留給敬虔者將來看的.

瑪3:17提到一個神設定的日子,很多英文譯本將之理解成神整理自己產業的日子.在這日,神將敬虔人納入為自己的產業,視他們為珍寶.經文採用父子的意象,神像父親,珍愛那些敬重祂的兒女.總有一天,人的本相會被揭露出來,敬畏神的和藐視神的,事奉神的和不事奉神的,是有不同結果的.神只收納敬畏祂的人,為永恆的產業.

神其實在回答頂撞祂的人,不要以為敬畏神和不敬畏神是沒有分別,沒有後果的.神用心留意著敬畏神的人的人生,至終,敬畏神的人,能成為神的至愛,分享神的豐盛.神提醒那些頂撞祂的人,不用因當下的艱難而「發爛渣」,應當懷著信心向前看,看到敬虔者在神那裡得到的豐盛人生.

接著,神更清楚地表述敬虔者和不敬虔者的相反下場.

瑪4:1　萬軍之耶和華說:「那日臨近,勢如燒著的火爐,凡狂傲的和行惡的必如碎稭,在那日必被燒盡,根本枝條一無存留.

瑪4:2　但向你們敬畏我名的人必有公義的日頭出現,其光線有醫治之能.你們必出來跳躍如圈裡的肥犢.

瑪4:3　你們必踐踏惡人；在我所定的日子,他們必如灰塵在你們腳掌之下.」這是萬軍之耶和華說的.

對不敬虔的人而言,神定的日子,卻是審判的日子.他們的人生,會灰飛煙滅,歸於無有.對敬虔的人而言,這日子是救贖的日子.有公義的日頭出現,照耀他們.在這裡,「公義的日頭」這意象,指向救主.救主來到,會醫好敬虔者一生的哀傷.到那時,他們便會如出牛棚的小牛,又開心,又自在.神告訴那些頂撞祂的人,請不要說惡人的結局比敬虔人的結局更好,事實不是這樣的.惡人的一切,歸於無有.敬虔者卻領受豐盛而圓滿的人生.

最後,是瑪拉基書的結語,也就是基督教聖經舊約部分的結語.

瑪4:4　「你們當記念我僕人摩西的律法,就是我在何烈山為以色列眾人所吩咐他的律例典章.

瑪4:5　「看哪,耶和華大而可畏之日未到以前,我必差遣先知以利亞到你們那裡去.

瑪4:6　他必使父親的心轉向兒女,兒女的心轉向父親,免得我來咒詛遍地.」

在歷史的大變局中,人心不安,很多神的兒女控訴神,比不信的人的控訴更甚.神最後的回應及勸勉,是叫人回歸最基本的東西,就是遵守摩西的律法.摩西的律法,是神的兒女最根本的法則.活出這些法則,就是活出神的兒女的獨特身分.這身分連接天地,將神的心意實踐在人間.唯這生命,能在動盪的大地上站穩.

摩西之後,經文提及以利亞.摩西代表律法,以利亞代表先知.律法塑造身分,先知展望將來.摩西加以利亞,律法加先知,就是鼓勵人,以神兒女的身分,開展神救贖的歷史.現實可能是殘酷的,但要有先知的視野,向前展望,看到神終極救贖的日子.在這裡,經文說,神會差遣以利亞作先頭部隊,開展神更偉大的救贖工程.我們知道,400年後,有一個像以利亞的人物出現,那就是施洗約翰,他作為耶穌的開路先鋒,開展一個全新的救贖時代.這就是瑪拉基書下一卷書馬太福音開頭的情節了.後來,當耶穌登山變像時,同時出現的神聖人物便是摩西和以利亞了,他們一起商討,改變天地的救贖工程.

瑪4:6提到,這位以利亞般的人物要「使父親的心轉向兒女,兒女的心轉向父親」.「父親」與「兒女」的原文,都是男性眾數,故此,是「父親們」和「兒子們」.按上文,經文都是關乎勸勉人不要藐視神,要敬畏神而論的,故此經文的意思不應指向一般的家庭關係的復和.有不少家庭事工,用這節經文來作為事工的基礎,鼓勵父親與兒女們復和.這是善意的引用,但應不是經文原本的意思.看來,這裡所謂「父親」,應指信仰傳統裡的列祖,先知以利亞的事奉,是讓列祖的信仰傳統能傳遞到後來一代又一代的子孫們,而後來的子孫們又願意承接先輩的信仰遺產.先知期望的,是源遠流長的道統不致斷絕.

道統斷了,人不再敬畏神,大地就崩壞,歷史也暗淡無光,這引來神對大地的審判.然而,讓我們再看看舊約最後一節經文:「免得我來咒詛遍地」.神特別派以利亞來,就是「免得」大地真正面臨崩壞.托著這個世界的,是慈愛的天父.祂不斷差派祂的僕人來,勸勉世人回歸祂的正道.當然,我們知道,最後神親自來到人間,「免得」世人走向滅亡.這都是新約講的故事了.我下一個月開始講.

瑪拉基書雖然簡單,但可看作是以色列民千多年來的信仰歷史的回顧.人生的艱難依舊,有人頂撞神,有人敬畏神,一直以來,都是如此.神卻像慈父,叮囑人走正路,鼓勵人有盼望,希望人對正義的天國不要失去渴望.在這苦難人間,咒罵神的人其實居多,但孤寂的神仍默默地背起人類罪惡造出來的十架,就因這份愛,這世界仍能存留至今.人,到底幾時才能醒悟?

在這裡,我講了39個禮拜的道,每次講一卷舊約的經卷,終於講完了.我作一簡單的總結.基督教的舊約,大致按「經史子集」來結集.「經」是指摩西五經,五經講明神創造萬物的美意,也點明神在歷史中的救贖.祂塑造一群屬祂的子民,透過他們,向人間啟明救贖的正道.「史」是指從約書亞記到以斯帖記一系列的歷史書卷,這些書卷挑明歷史變化背後的信仰原則.歷史的興亡,不在乎武力,而在乎信仰神,行事公義,實踐仁愛.「集」是指從約伯記到雅歌這幾卷詩歌智慧文學,智慧文學不強調救贖神學,強調的是創造神學.生有時,死有時；得有時,失有時,這不關乎神救不救你,而是有限人生的現實寫照.要接受現實,你不接受,現實仍是這樣.智慧的人,就是在這限制中,以敬畏的心,過感恩的生活.舊約最後的,是「子」,是從以賽亞書到瑪拉基書一系列的先知書.就著歷史的不同處境,先知發出適時的指導.面對腐敗的世情,先知呼喚人悔改.面對歷史的急劇轉變,先知提示人信靠神.面對絕望,先知鼓舞人向前盼望.無論神的子民去到什麼光景,先知總在叮嚀.先知的呼喚,代表著神的臨在,代表著神與祂的子民同行.在這苦難人間,人並不寂寞,在寂寞深處,總暗藏著那位消磨著自己生命去支撐我們生命的神.最後,這位神來到人間,背起了十架,犧牲了生命,用愛承托著蒼生.這就是新約的故事了.下次,我們一起看這故事.

願榮耀歸予聖父,聖子,聖靈.



\newpage

\section{十字架的道理}
\label{sec:784}
\textbf{王家輝牧師-十字架的道理}
\newline
\newline
\hyperref[sec:783]{< < < PREV SERMON < < <}
~
\hyperref[sec:toc]{[返目錄]}
~
\hyperref[sec:785]{> > > NEXT SERMON > > >}


哥林多前書 1:10-18
\newline
\begin{longtable}{cl}
\hline
\hline
章節 & 經文 (和合本修訂版)\\
\hline
1:10 & \begin{tabularx}{0.7\textwidth}{X} 弟兄們，我藉我們主耶穌基督的名勸你們說話要一致。你們中間不可分裂，只要一心一意彼此團結。 \end{tabularx} \\ \\ \relax
1:11 & \begin{tabularx}{0.7\textwidth}{X} 我的弟兄們，革來氏家裡的人曾對我提起你們，說你們中間有紛爭。 \end{tabularx} \\ \\ \relax
1:12 & \begin{tabularx}{0.7\textwidth}{X} 我的意思是，你們各人說：「我是屬保羅的」；「我是屬亞波羅的」；「我是屬磯法的」；「我是屬基督的。」 \end{tabularx} \\ \\ \relax
1:13 & \begin{tabularx}{0.7\textwidth}{X} 基督是分裂的嗎？保羅為你們釘了十字架嗎？你們是奉保羅的名受了洗嗎？ \end{tabularx} \\ \\ \relax
1:14 & \begin{tabularx}{0.7\textwidth}{X} 我感謝神，除了基利司布和該猶以外，我沒有給你們中的任何一個人施洗， \end{tabularx} \\ \\ \relax
1:15 & \begin{tabularx}{0.7\textwidth}{X} 免得有人說你們是奉我的名受洗的。 \end{tabularx} \\ \\ \relax
1:16 & \begin{tabularx}{0.7\textwidth}{X} 我曾為司提法那家施過洗；此外我已記不清有沒有給別人施過洗。 \end{tabularx} \\ \\ \relax
1:17 & \begin{tabularx}{0.7\textwidth}{X} 因為基督差遣我不是為施洗，而是為傳福音；並不是用智慧的言論，免得基督的十字架落了空。 \end{tabularx} \\ \\ \relax
1:18 & \begin{tabularx}{0.7\textwidth}{X} 因為十字架的道理，在那滅亡的人是愚拙，在我們得救的人卻是神的大能。 \end{tabularx} \\ \\
[1ex]
\hline
\hline
\end{longtable}
講題:十字架的道理TheMessageofTheCross

經課:哥林多前書1:10-18

講員:王家輝牧師

場合:香港中文大學崇基學院禮拜堂主日崇拜

日期:2026年1月25日

經文(新標點和合本)

10弟兄們,我藉我們主耶穌基督的名勸你們都說一樣的話.你們中間也不可分黨,只要一心一意,彼此相合.11因為革來氏家裡的人曾對我提起弟兄們來,說你們中間有紛爭.12我的意思就是你們各人說:「我是屬保羅的」；「我是屬亞波羅的」；「我是屬磯法的」；「我是屬基督的」.13基督是分開的嗎?保羅為你們釘了十字架嗎?你們是奉保羅的名受了洗嗎?14我感謝上帝,除了基利司布並該猶以外,我沒有給你們一個人施洗,15免得有人說,你們是奉我的名受洗.16我也給司提法那家施過洗,此外給別人施洗沒有,我卻記不清.17基督差遣我,原不是為施洗,乃是為傳福音,並不用智慧的言語,免得基督的十字架落了空.

18因為十字架的道理,在那滅亡的人為愚拙；在我們得救的人,卻為上帝的大能.

我們都活在不能認輸的世界

從小到大,我們都在跟別人比較,或被別人跟自己作比較.父母在有意,無意中比較子女的各方面表現,亦因此令兄弟姊妹間作比較.在學校,我們要跟同學作比較.甚至學校之間也要作比較,分為不同等級(BandOne,TwoThree).當然大學也是一樣,是全世界作比較,大家會關注大學今年的QS排名是升,還是降?

非常不幸,在同一教會之內也會有比較.你是初信,我是資深.既然是資深信徒,那麼就被想像成是比較「屬靈」.而教會與教會之間,同樣也會被比較.那些是超大型教會,那些是小型教會,大家都羨慕成為超大型教會.

除了要互相比較之外,我們也習慣跟「同聲同氣」的人走在一起,而社交媒體的大流行,令這種稱為「同質化」的現象更加強烈.我們今日對事情的了解很容易受「意見領袖」(KOL)影響,對於某些人對時事分析,我們幾乎奉為真理.當然有好些KOL是做了大量的資料分析,然後才發表他們的意見,可是亦有不少KOL為了增加點擊率而鋪陳嘩眾取寵的內容.結果因為社交媒體的「過濾泡泡效應」,令我們的偏好被不斷的強化,結果令我們跟不同想法的人,造成更大的對立面.這樣的情況,對於政治立場的塑造更是明顯的.最令人擔憂是令不同立場的人趨向極端,以至產生極端行為來排斥異己.在這樣的氛圍下,我們不自覺地被主流價值觀所影響,並且走進分黨結派的旋渦之中.

剛才提到我們從小到大都在比較,而大部份人也不希望被人比下去,故此我們就要力爭上游,只能「認叻」,而不能認輸!為了要勝過別人,其中一個重要方法就是「壯大自己勢力」,勢力越大就代表越成功.由於勢力是抽象的,因此要把它具體化,最實際的方式就是以財富和市場份額作衡量.結果今時今日馬斯克(不是馬克斯!),黃仁勳等就被視為新一代的超級成功人士.他們的影響力可說是全球性!

那麼基督徒又如何?我們追隨的是什麼呢?我們是跟隨世界的價值觀,還是跟隨基督呢?當然我們在頭腦上知道要跟隨基督,可是在現實生活上,我們卻不是這樣的實踐.

我們喜歡主導一切,我們喜歡有跟隨者,成為強者才會成功.若然在工作中,沒有地位,那就在家中確立自己的權威.若然在家庭,在工作也沒有地位,那就在教會確立自己的權威.

當這樣的想法不斷蘊釀時,結果令教會發生不同的衝突.情況猶如今天的經文哥林多前書一章10至18節所講述的相似.讓我們看一看.

哥林多教會反映一個不能認輸的世界

保羅時代的哥林多跟今天的都市生活有點相似,當中是充滿階級差異和領導風格衝突.

在這段經文中,我們看到保羅面對一個嚴重的問題:哥林多教會的分裂.想像一下,一個小小的教會,成員可能只有幾十人,卻分成好幾派.有些人說「我是屬保羅的」,因為保羅是開拓者,講道風格直率,像個哲學家；有些人說「我是屬亞波羅的」,亞波羅是個能言善道的知識份子,從亞歷山太來,講道充滿修辭學的魅力；還有「屬磯法的」,磯法就是彼得,代表傳統猶太根基；甚至有自以為最有見地的說「我是屬基督的」,聽起來很屬靈,但其實亦是另一種派系鬥爭.

哥林多教會為何有這樣的表現?其實背後也是反映著一個不能認輸的世界.我們清楚地看見,保羅藉書信呼籲弟兄姊妹要合一.他強調分裂不是來自深奧的信仰思想的對立,而是日常生活事情的爭執,如富人輕視窮人(見11:18-22).保羅沒有嘗試找出假想敵,然後想方設法把打倒對手,而是批評社會結構中的競爭和傲慢,這是源自羅馬帝國的階級和權力遊戲.保羅聽到從革來氏家裡傳來的報告(這個革來氏家裡的人可能是奴隸或一群婦女),他們認為這種分裂不妥當,因此要讓保羅知道哥林多教會究竟發生什麼事.這樣的狀況令保羅為教會憂心,因為這反映出人類的本性:我們總愛壯大自己的勢力,追隨偶像,而不是在基督裡合一.這彷彿是古今中外,有人的地方就會出現的問題.

在這樣的狀況下,教會面對的不只是個人紛爭,而是核心信仰危機.保羅質問:「基督是分開的嗎?」這句話像一把利劍,直刺人心.十字架本該是合一的記號,但現在卻帶來了分裂.保羅甚至自嘲自己記不清給誰施洗,這是幽默的筆觸,強調施洗不是他的重點工作,傳揚基督在十字架上犧牲然後復活的福音才是他的工作重點.但在哥林多人眼中,他們的「智慧」——希臘哲學的辯論,羅馬的權力遊戲——讓他們看不見十字架的真諦.十字架在他們看來,是愚拙的,軟弱的標記.耶穌沒有推翻羅馬帝國,沒有反抗不公審判,甚至在極大折磨下默然如羔羊.這在主流思想中,是徹頭徹尾的失敗.這種「智慧」是世界結構的一部分,包括教育,競爭和誇耀,阻礙了在基督十字架前的合一.

這是嚴肅的神學問題,是我們對人性的了解,那就是罪的表現——驕傲,分裂,追隨人,甚至妄想取代上帝的位置.

在人的眼中,十字架是失敗,因為耶穌沒有發動革命,沒有反抗鞭打.這在主流思想中,是軟弱.主流思想鼓吹「各人看自己比別人強」,所以有部份哥林多教會的信徒稱自己「是屬基督的」,其實心裡就是想著成為最強的派系,但這卻帶來孤立和衝突.保羅的憂慮是哥林多教會眾人仍依戀舊的智慧,而舊的智慧跟人的限制是有關的.

十字架的道理

人的限制是無力自救.在罪性中,我們無法靠自己的力量跳出來,我們需要一個比我們更強的力量,把我們從罪的深淵拉上來,這就是十字架上的拯救.保羅說:「基督差遣我,原不是為施洗,乃是為傳福音,並不用智慧的言語,免得基督的十字架落了空.因為十字架的道理,在那滅亡的人為愚拙；在我們得救的人,卻為神的大能.」

這裡,恩典是閃耀著的:十字架不是失敗,而是上帝的大能.十字架是超越自然界的力量和智慧.保羅呼籲信徒要有「要一心一意,彼此相合」,這是跟哥林多教會的狀況形成反差.若然他們想做到彼此相合,他們需要領略到十字架的道理——倒空自己,為更大的事犧牲.我喜歡用一個動作來示範十字架可以怎樣把要彼此排斥的人連在一起,那就是當我們看見在十字架上張開雙手的基督,他每一邊的手拉著一位不同陣營,不同想法的人.因著他這樣犧牲,我們這些原本不相干的人就可以連在一起了.當我們連在一起時,我們再不是屬於某一個派系,而是屬於在十字架上的基督.

我們知道在羅馬帝國,十字架是刑具,是抵觸羅馬法律的人才會被釘在其上,而按照聖經的記載,耶穌被治死的罪名是一項政治罪名,因為在羅馬帝國內,耶穌竟然被宣稱為猶太人的王,這是帝國不能接受的事.如果耶穌就此徹底死了,當然十字架就成為失敗的記號,但因著基督的復活,突破死亡的限制,十字架就成為一個顛覆想像,戰勝一切的記號.

這就是經文所言「十字架的道理,在那滅亡的人為愚拙；在我們得救的人,卻為上帝的大能.」的意思.耶穌基督的十字架看似軟弱——沒有打倒羅馬,沒有反抗——但正是這順服,卻戰勝了死亡.

保羅想指出,這個「世界的智慧」包括了競爭,壓制人的權力結構,甚至像妖魔般的統治者,是屬地上的,也是無形的.不過按以賽亞書29:14和耶利米書9:23-24,強調上帝使人的智慧歸於無有.

十字架成為拯救的記號,因為基督的謙卑——看別人比自己強,甚至為罪人死——帶來合一.

我們終於要認輸——認輸給十字架的大能

保羅對哥林多教會的提醒,對我們今日而言有什麼意思呢?首先,我們要好好內省,問問自己,我是跟隨基督,還是隨波逐流.若主流思想是要壯大自己的勢力,但基督卻給我們謙卑的榜樣,我們能夠接受別人比自己強的事實嗎?保羅呼籲足信徒要彼此相合,因為我們共同連於十字架的基督,這意味著我們是屬於彼此,超越競爭對手的關係.

我們看看自己一雙手,十隻手指有長短.有長短應該是正常的,如果每隻手指一樣長,相信對我們抓取物件會帶來許多不便.或許人都是喜歡成為領導者,話事人,這是古今相同的,不過保羅卻提醒我們,最重要是忠於基督.今天,我們面對社交媒體影響而令人與人之間造成更多的對立,我們需要認真明白十字架的道理:不以辱罵作回擊,而是以愛心回應.

由我們家裡開始做起,因為我們最易對家人缺乏耐性與愛心.如果你對家人也做得到,你向外人就應該有更大的能力做得到.

基督的犧牲是在軟弱中顯大能.讓我們實踐:在家裡,看配偶比自己強；在職場,看見同事的貢獻；在教會,樂意合一服事.十字架的道理總歸就是抗拒被主流思想所影響,活出跟隨基督的新生命.

願我們不再追逐世界的智慧與勢力,乃是單單傳揚十字架的福音.願上帝的大能透過我們這些真誠向基督「認輸」的人,顯明祂的智慧與能力,讓主的教會,讓這個地方,讓世界每一處也這樣行.

願榮耀歸於那被釘十字架,卻最終復活的主耶穌基督!阿們!



\newpage

\section{信仰的要求}
\label{sec:785}
\textbf{李均熊牧師-信仰的要求}
\newline
\newline
\hyperref[sec:784]{< < < PREV SERMON < < <}
~
\hyperref[sec:toc]{[返目錄]}
~
\hyperref[sec:786]{> > > NEXT SERMON > > >}


彌迦書 6:1-8
\newline
\begin{longtable}{cl}
\hline
\hline
章節 & 經文 (和合本修訂版)\\
\hline
6:1 & \begin{tabularx}{0.7\textwidth}{X} 當聽耶和華說的話：起來，向山嶺爭辯，使岡陵聽見你的聲音。 \end{tabularx} \\ \\ \relax
6:2 & \begin{tabularx}{0.7\textwidth}{X} 山嶺啊，要聽耶和華的指控！大地永久的根基啊，要聽！因耶和華控告他的百姓，與以色列爭辯。 \end{tabularx} \\ \\ \relax
6:3 & \begin{tabularx}{0.7\textwidth}{X} 「我的百姓啊，我向你做了甚麼呢？我在甚麼事上使你厭煩？你回答我吧！ \end{tabularx} \\ \\ \relax
6:4 & \begin{tabularx}{0.7\textwidth}{X} 我曾將你從埃及地領出來，從為奴之家救贖你，我差遣摩西、亞倫和米利暗在你前面帶領。 \end{tabularx} \\ \\ \relax
6:5 & \begin{tabularx}{0.7\textwidth}{X} 我的百姓啊，當記念從前摩押王巴勒如何籌算，比珥的兒子巴蘭如何回應他，當記念從什亭到吉甲所發生的事，好使你們明白耶和華公義的作為。」 \end{tabularx} \\ \\ \relax
6:6 & \begin{tabularx}{0.7\textwidth}{X} 「我朝見耶和華，在至高神面前跪拜，當獻上甚麼呢？難道獻一歲的牛犢為燔祭來朝見他嗎？ \end{tabularx} \\ \\ \relax
6:7 & \begin{tabularx}{0.7\textwidth}{X} 耶和華豈喜悅千千的公羊，或是萬萬的油河嗎？我豈可為自己的過犯獻我的長子，為自己的罪惡獻我所親生的嗎？」 \end{tabularx} \\ \\ \relax
6:8 & \begin{tabularx}{0.7\textwidth}{X} 世人哪，耶和華已指示你何為善。他向你所要的是甚麼呢？只要你行公義，好憐憫，存謙卑的心與你的神同行。 \end{tabularx} \\ \\
[1ex]
\hline
\hline
\end{longtable}
講題:信仰的要求TheDemandofFaith

經課:彌迦書6:1-8

講員:李均熊牧師

場合:香港中文大學崇基學院禮拜堂主日崇拜

日期:2026年2月1日

起:熟識等於認識?

「世人哪,耶和華已指示你何為善.他向你所要的是甚麼呢?只要你行公義,好憐憫,存謙卑的心,與你的神同行.」(彌6:8)

這一句經文相信是大家十分熟識的經文.但熟識不等於認識,彌迦書這裡真是這麼直接簡單告訴我們信仰的要求嗎?這樣只是抽出一句,然後探討這一句的意思,再將對只是一句經文的理解加諸於我們的信仰中,這樣其實是非常危險的.不僅聖經的某些用語其實是和現代的用語十分之不同,就如古代的公義和現代的公義概念相差何止萬里,我們僅僅抽出一節,沒有考慮到經文的上文下理,這樣的理解就會變得十分片面.試想想有一天突然一個人走過來對你說:「小心一點」,沒有前因後果,上文下理,他的意思可能就不是那麼容易確定了.

承:陌生化處理

為了進一步去解釋究竟彌迦書這一節經文的意思,容許我先以原文／聖經希伯來文重新翻譯一次這一句說話,目的不是讓它更通順易明,而是令你對此節再一次陌生起來,迫使我們細想經文遣詞用字的特色.

這一節直接由原文翻出來,可以是這樣的

「他已告訴了你,亞當(Adam)!甚麼是「好」(good).上主要求你的是甚麼呢?無非是要你實施misphat(do/excutemisphat)和愛hased(lovehased),並謙卑謹慎(希臘文:如僕婢般預備好服侍)地與你的神同行.」

謙卑:原文意思不太容易確定,但古代希臘文譯者選了一個僕人預備侍候主人意思的詞來翻譯,所以很多新的中文譯本都譯為謙卑謹慎

世人啊:原文是用上Adam亞當

行公義:原文可以是實施公義.但「公義」原文misphat,在和合本多數譯作「公平」；而「公義」則多數翻譯另一個字zedaq,然而這2個字大部分時間都是一齊出現的.公平公義這兩個字一齊出現,其實是古代文獻一個獨有的專用名詞,其上文下理通常是當立約的神明/宗主國立一國君,便要求君王照顧國中孤寡無倚的人,和實施公平的法律,絕對不可以屈枉正直,只要我們看看彌迦書其他章節用上misphat的上文下理就可以理解(不用擔心,主要就是第三章):

1於是我說:雅各的領袖,以色列家的官長啊,你們要聽!你們豈不知道公平嗎?2你們惡善好惡,剝我百姓身上的皮,從他們的骨頭上剔肉,3你們吃我百姓的肉,剝他們的皮,打斷他們的骨頭,如切塊下鍋,如釜中的肉……9當聽這話,雅各家的領袖,以色列家的官長啊!你們厭棄公平,在一切事上屈枉正直；10以血建立錫安,以罪孽建造耶路撒冷.11城裡的領袖為賄賂行審判……

好憐憫:原文是愛hesad.Hesad也是一個專用名詞,通常是指對盟約忠誠.所以應該用loyalty忠誠來翻譯更好(如果用廣東話,甚至應該用「義氣」來譯)找回彌迦書內的使用,可以看7:18-20

18有哪一個神明像你,赦免罪孽,饒恕他產業中餘民的罪過?他不永遠懷怒,喜愛施恩.19他必轉回憐憫我們,把我們的罪孽踏在腳下.你必將他們一切的罪投於深海.20你必按古時向我們列祖起誓的話,以信實待雅各,向亞伯拉罕施慈愛.

這一段表明,上主是首先忠於祂對子民的約.所以行公義,好憐憫並不能望文生義地去理解為今日的社會公義或好心幫助人,而是強調對忠誠的上主約的回應,強調當中對約的忠誠和領袖在立約子民中需踐行的要求.

轉:上主與子民爭辯(嗌交)

但真要了解這一句經文,我們一定不可脫離經文上文下理的脈絡.

這一節經文其實出現於上主要和祂的子民的爭辯中(1-2節出現了三次)!初看,這是一場法庭對決.控方／原告將要就著辯方是否違背西乃之約而提出訴訟；辯方則要為其一方確實嚴守盟約而據理力爭.這種「爭辯」,或說「法庭訴訟」的文體,在先知書中十分普遍.訴訟的程序首先是傳召見證人作仲裁者(1-2節),然後控方提出指控(3-5節),接著由辯方回應指控(6-7節).最後的裁決,往往會指出受控訴的一方該如何行動,補贖遭破壞的盟約(8節).不過,學者開始質疑這種「爭辯」是否真的可以稱為「法庭訴訟」.法庭的場景是否切合爭辯的方式和內容?這種對文體定位的疑問,十分值得我們繼續探討,免得我們將盟友,家庭,甚至情人之間的「強烈溝通」(「嗌交?!」【意即吵鬧】)理解為對簿公堂,你死我活的法律程序.無論如何,「爭辯」的出現,皆因盟約雙方已經出現裂痕,其中一方的行為已經對盟約造成危害,必須透過這樣的爭辯來正視和提出修正.

是誰危害盟約關係呢?上主接著的問題和說話(3-5節),似將矛頭直指自己!「我的百姓啊,我向你做了甚麼呢?我在甚麼事上使你厭煩?」(3節)究竟上主說這話時,用的是怎樣的語氣?祂是在諷刺以色列人,挖苦他們仍不自知盟約的破壞全是他們的責任嗎?還是上主深受傷害,心感疑惑:究竟自己在何事上出錯呢?為甚麼自己的子民這樣討厭自己?上主道出祂與子民過往的故事,究竟是控告以色列人不知感恩圖報,還是上主自省:我不是已經顯明我對你的愛嗎?難道那還不足夠?

其後的兩節經文(6-7節)似是以色列人的回答:「我知是我的錯,不過,不是我不想到你面前贖罪,而是你的標準是甚麼呢?是初生的牛犢嗎?難道你喜歡千千的公羊,萬萬的油河嗎?你不是要我們奉獻自己的長子吧?」那是怎樣的問題?他們真心感到疑惑,不曉得上主的要求；還是他們還在強嘴:誰危害盟約關係呢?是你!盟約的要求這麼高,叫我們怎樣達到?

若說這是一段法庭訴訟的對話,仲裁的法官(或是撮合這場爭辯,站在一旁卻再也按捺不住的先知)必然速速叫停:肅靜(收聲)!這算甚麼法庭辯論(你們在鬧甚麼)?你們不清清楚楚陳明理據,卻將心事含混在嘴邊?「【不是】已經指示你──世人哪!甚麼是善?上主向你所要的又是甚麼呢?無非是要你行公義,好憐憫,謙虛謹慎與你的神同行.【那不是可以了嗎?】」

彌迦書6:8不是一條抽離脈絡,放諸四海皆準的「通用金句」,而是上主與子民爭辯(訴訟/強烈溝通)中的裁決總結.在6:1–7描述的張力裡,百姓或以誇張祭物為藉口,或以無知為推搪,試圖將盟約的破裂合理化.先知於是以第8節正面,清楚地指出(或許還帶著幾分激動):上主的要求早已明明啟示,毫不含糊——就是在立約關係中踐行「公平／審判」(misphat)與「忠誠／慈愛」(hesad),並以僕人式的謙卑與神同行:就是時刻留心主人的心意,先問「主今日要我怎樣做?」然後𧗊所能去實完成.

因此,「行公義,好憐憫」並非現代社會倫理的空泛口號,而是對西乃之約的實際履行:特別臨在於領袖與百姓如何保障弱小,拒絕貪婪,秉行不枉不縱,以及在群體中持守對上主的忠誠之愛.結論很直接:面對上主的質問與指證,神的子民不能再以「不知道標準」或「加倍獻祭」以堆疊宗教熱情來逃避；唯一合宜的回應,就是回到與上主立約的核心,把misphat公平與hesad忠誠具體活出,並以謙卑的腳步與神並行.

弟兄姊妹,熟識不等於認識.我們也許早把「行公義,好憐憫」背得滾瓜爛熟,卻仍可能把它化作抽離,籠統的空泛倫理或口號.我認為,先知是帶著激情來提醒我們:神要的不是更多宗教熱心,而是正確地履行祂已啟示的要求.讓我們知而去行,不再找藉口；以公平,忠誠與謙卑,與神同行,使我們的生命與群體成為這城的祝福.阿們.



\newpage

\section{不切實際的耶穌}
\label{sec:786}
\textbf{鄧瑞強牧師 - 不切實際的耶穌}
\newline
\newline
\hyperref[sec:785]{< < < PREV SERMON < < <}
~
\hyperref[sec:toc]{[返目錄]}
~
\hyperref[sec:787]{> > > NEXT SERMON > > >}


馬太福音 5:38-48
\newline
\begin{longtable}{cl}
\hline
\hline
章節 & 經文 (和合本修訂版)\\
\hline
5:38 & \begin{tabularx}{0.7\textwidth}{X} 「你們聽過有話說：『以眼還眼，以牙還牙。』 \end{tabularx} \\ \\ \relax
5:39 & \begin{tabularx}{0.7\textwidth}{X} 但是我告訴你們：不要與惡人作對。有人打你的右臉，連另一邊也轉過去由他打。 \end{tabularx} \\ \\ \relax
5:40 & \begin{tabularx}{0.7\textwidth}{X} 有人想要告你，要拿你的裡衣，連外衣也由他拿去。 \end{tabularx} \\ \\ \relax
5:41 & \begin{tabularx}{0.7\textwidth}{X} 有人強迫你走一里路，你就跟他走二里。 \end{tabularx} \\ \\ \relax
5:42 & \begin{tabularx}{0.7\textwidth}{X} 有求你的，就給他；有向你借貸的，不可推辭。」 \end{tabularx} \\ \\ \relax
5:43 & \begin{tabularx}{0.7\textwidth}{X} 「你們聽過有話說：『要愛你的鄰舍，恨你的仇敵。』 \end{tabularx} \\ \\ \relax
5:44 & \begin{tabularx}{0.7\textwidth}{X} 但是我告訴你們：要愛你們的仇敵，為那迫害你們的禱告。 \end{tabularx} \\ \\ \relax
5:45 & \begin{tabularx}{0.7\textwidth}{X} 這樣，你們就可以作天父的兒女了。因為他叫太陽照好人，也照壞人；降雨給義人，也給不義的人。 \end{tabularx} \\ \\ \relax
5:46 & \begin{tabularx}{0.7\textwidth}{X} 你們若只愛那愛你們的人，有甚麼賞賜呢？就是稅吏不也是這樣做嗎？ \end{tabularx} \\ \\ \relax
5:47 & \begin{tabularx}{0.7\textwidth}{X} 你們若只請你弟兄的安，有甚麼比別人強呢？就是外邦人不也是這樣做嗎？ \end{tabularx} \\ \\ \relax
5:48 & \begin{tabularx}{0.7\textwidth}{X} 所以，你們要完全，如同你們的天父是完全的。」 \end{tabularx} \\ \\
[1ex]
\hline
\hline
\end{longtable}
講題:不切實際的耶穌 Jesus, so unrealistic

經課:馬太福音5:38-48

講員:鄧瑞強牧師

場合:香港中文大學崇基學院禮拜堂主日崇拜

日期:2026年2月15日

各位弟兄姊妹,早安.願神的恩典常與你們同在.

舊約講完了,我的講道終於來到新約.天父若許可,我未來27次的講道,會一次講一卷新約書卷,直到講完新約全部27卷書卷為止.我其實有一個卑微的渴望,是渴望弟兄姊妹能對整本聖經有一較整全的理解.聖經其實不只是幫助我們信耶穌,得永生,聖經更是承托著我們整個人生,讓我們無論面對什麼光景,都看到神的真實與恩典.以舊約為例,39卷書卷涵蓋人生的方方面面,無論在什麼光景,神都在.在遠古荒涼的時代,那個不知人生意義為何的亞伯拉罕,他聽到神的呼喚,讓他明白自己的天命.當亞伯拉罕仍只有一個小家庭時,神在,讓他開展一個民族.當這民族在埃及做奴隸時,神在,讓他們得釋放,並在迦南地建立家園.當他們在迦南地群龍無首時,神在,為他們揀選王者,建立君主制度.當君王腐敗時,神在,差派先知提點君王.當社會禮崩樂壞時,神在,興起列國來審判他們.當國破家亡時,神在,給他們盼望.當神的子民從被擄之地返回故鄉,面對頹垣敗瓦而失望時,神在,教他們建立生命,以生命為殿.無論是喜是悲,無論是得是失,聖經教我們看到:神在,祂以無限仁慈支撐著我們.聖經的不同經卷,對應不同的人生處境,讓我們看到,無論我們面對什麼人生大問題,神都在,真切地指引我們.

舊約的所有經卷,透露神是那位與我們同在的神,祂不會放棄我們,這就是神與我們之間的約的連繫.就因神人之間的這種緊密連繫,最後,神來到人間,成為人.馬太福音一開首便表明耶穌降生的真義便是「以馬內利」,意思是:神與我們同在.三一神透過耶穌啟示這個神人聯繫的世界的真相.來到新約,是來到神啟示的新階段,透過耶穌,我們更清楚地看到這個世界的真正面目.今日,我們一起看新約的第一卷書馬太福音,跟著耶穌腳步,我們將看到一個奇異的新世界.

在馬太福音裡,耶穌有多重角色.他是神的兒子,他是天國的王,他是來到人間的救主,他也是一個新摩西,教導我們,如何做天國的子民.馬太福音最著名的段落之一,就是第5至7章的「登山寶訓」.耶穌像摩西,走上山,向我們頒佈天國的誡命.今日,我們看其中一段.

馬太福音5:38-48

太5:38　「你們聽見有話說:『以眼還眼,以牙還牙.』

太5:39　只是我告訴你們,不要與惡人作對.有人打你的右臉,連左臉也轉過來由他打；

太5:40　有人想要告你,要拿你的裡衣,連外衣也由他拿去；

太5:41　有人強逼你走一里路,你就同他走二里；

太5:42　有求你的,就給他；有向你借貸的,不可推辭.」

太5:43　「你們聽見有話說:『當愛你的鄰舍,恨你的仇敵.』

太5:44　只是我告訴你們,要愛你們的仇敵,為那逼迫你們的禱告.

太5:45　這樣就可以作你們天父的兒子；因為他叫日頭照好人,也照歹人；降雨給義人,也給不義的人.

太5:46　你們若單愛那愛你們的人,有甚麼賞賜呢?就是稅吏不也是這樣行嗎?

太5:47　你們若單請你弟兄的安,比人有甚麼長處呢?就是外邦人不也是這樣行嗎?

太5:48　所以,你們要完全,像你們的天父完全一樣.」

主耶穌一開頭便說:「你們聽見有話說...... 只是我告訴你們......」,便提示我們面對一個很大的對比,這對比大得後者否定前者.主耶穌來世,不是要肯定世人習以為常的做法,而是否定這做法,展示另一種做法,一種看起來不合常理的做法.主耶穌來世,不是教我們如何在世人建構的世界裡獲得成功,而是教我們想像另一個世界,一個神聖的世界.

登山寶訓的誡命,在釋經上,倫理上,神學上都引起極大爭論.這些誡命的性質,到底是什麼?「有人打你的右臉,連左臉也轉過來由他打；有人要拿你的裡衣,連外衣也由他拿去」,若真的照字面去做的話,有人認為,這只會令社會亂作一團.這種教導看來不單沒有制止惡人作惡,反鼓勵他們變本加厲.若是這樣,這個不義的世界,還能管治嗎?

有人認為,耶穌的教導,是一種明智的策略.你對敵人和善,會令敵人變成你的朋友.你為敵人多走一步,會激發他的同理心.害你「炒魷魚」的那位同事,你送禮物給他,他會變得良善,回頭會介紹份更好的工作給你.若你真的這樣想,你真是天真得可憐.從現實策略來說,登山寶訓並不是好的策略.

有人認為,耶穌講的,是一種正向思考,教我們積極地面對逆境.正向思考並不質疑世人建構的世界,和世人追求的目標,但耶穌明顯要和現實世界的一切作出分別,作出對比,他根本不認同現實的一切.若他只是一個正向的人,教人在現實世界中獲取成功,他就根本不會被人釘在十架上,反而會成為社會的成功典範.

有人認為,由於登山寶訓的教導並不現實,不能真的作為信仰社群長遠運作的指引,故此,這應該只是一種特殊時期的特殊做法.這種觀點認為,當時的信徒以為主耶穌很快會回來,故採取極端的做法.若你肯定主耶穌下禮拜會回來,你應該願意變賣一切分給窮人.一個信徒若認為主耶穌很快會回來,有人要拿他的裡衣,他連外衣也由這人拿去,是一點不出奇的.但這種觀點,會令登山寶訓變得和我們悠長的歷史生活無關.

有人認為,由於登山寶訓難以實踐,故它只能是一種遙不可及的理想.由於是理想,做不到,便情有可原.又由於人人現實上都做不到,有人認為,登山寶訓的重點其實是凸顯我們多麼需要恩典.若真是這樣,登山寶訓就不單不是信徒努力實踐的倫理,而是誘導信徒不去實踐,這些信徒可能心想:硬去實踐這麼難的行為,不就是靠行為得救麼,這是遠離恩典了.假恩典之名,這些信徒就可以安心放下登山寶訓的教導,而依然故我.

不同學者千方百計為登山寶訓這些教導找理由,就因為這些教導看起來不切實際,助紂為虐,無法實行.想想:若有人打你的右臉,你連左臉也轉過來由他打,他會收手嗎?若有人要拿你的裡衣,你連外衣也由他拿去,你還剩下什麼呢?愛你們的仇敵,為那逼迫你們的禱告,這不是縱容壞人作惡嗎?面對這些教導,不同學者找不同理由,但想真點,這些理由都沒有什麼理由.

問題出在哪裡?為什麼這些教導看來不切實際,但主耶穌卻又語重心長地教導我們?問題出在哪裡?

問題出在:我們將這些教導,從主耶穌所建造的天國那裡抽了出來,而硬套進我們習慣的世俗領域裡去,結果當然是格格不入.想像一下,有一個外星人來到地球,這個外星人與地球人最不同的地方,是他不用吃東西,他能像植物那樣進行光合作用便已獲得足夠養分.他告訴地球人,不用為了爭奪生存的資源而開戰,只需到戶外走走,曬一下太陽,吸一口碳氣,便已有足夠的生存養分了.我們會覺得他講的東西,完全不切實際.又設想這外星人能夠在地球上建構一平行宇宙,地球人若進入他建構的世界裡,生命會轉化,而真能自行光合作用,明白他講的,都是事實.他講的,不是完全無法實踐的道理,而是呈現一個平行宇宙的真相.問題出現了,大部分人都沒有進入他的平行宇宙,而只是站在我們習慣的這個世界上,並批評這外星人講的東西,不符現實,不切實際.

耶穌當然不是外星人,他比外星人強得多,他是神,他見到的真相比我們習慣的真相,差異更大,對比更強烈.耶穌來到人間,以獨特的生命,呈現一個奇異的,神聖的世界.耶穌以他的一生,呈現他眼中的真相.真相是:

相信天父,比相信人類的各種想法更好:若天父的旨意叫他背負十架,他就全力以赴,並不認為千方百計去增加自己的利益更化算.人生追求的目標,不應該交由世俗去設定,而只能由天父設定.

相信天父是歷史的主宰,相信天父的仁愛,並不同意人類的各種殘酷手段更能推進歷史的文明進步.

當被人捉拿時,耶穌叫彼得放下武器,他看到真正的保護,只能來自天父.人心真正的平安,不能靠以暴易暴來保證.

在十架上時,耶穌第一句話,是求天父赦免釘他上十架的人.他看到,人間的痛苦與暴力,並不是終極的真相,並不需要讓痛苦與暴力主宰我們的理性與情感.真正的終極真相,是仁慈的天父,耶穌信任祂的仁慈,而不是世人的暴力.

若果你看到耶穌所看到的真相,你會明白,登山寶訓的教導,並不是匪夷所思,超乎常理的.登山寶訓的教導,放在耶穌所呈現的真相裡,是完全合理的.但當抽離了那終極真相,而硬塞入我們世俗所建構的世界裡,才顯得不切實際,格格不入.

若我們追求以最有效的策略,帶來最大的果效,則有人強逼你走一里路,你就同他走夠二里,只能是下下之策.若我們看個人的成敗,大過順從天父的旨意,則有人打你的右臉,連左臉也轉過來由他打,就根本背離正義.若我們認為「誰大誰惡誰正確」就是真相,不認為仁慈的天父才是終極真相,則「愛你們的仇敵,為那逼迫你們的禱告」只能是荒謬至極的歪理.

登山寶訓不獨是講一種古怪的倫理,它其實是呈現一奇異的世界.天父透過耶穌的一生,呈示這世界的真相.在這真實的世界裡,放下世俗的一切跟從主,是好的.在這真實的世界裡,施比受更為有福.在這真實的世界裡,寬恕比記仇好,不還手比還手好,十架比利劍好.這才是人生的真相嗎?是,因為耶穌就是這樣活的.耶穌是神,神也是這樣活的.耶穌順從的天父,祂也是這樣活的.

太5:44　只是我告訴你們,要愛你們的仇敵,為那逼迫你們的禱告.

太5:45　這樣就可以作你們天父的兒子；因為祂叫日頭照好人,也照歹人；降雨給義人,也給不義的人.

真正的信耶穌,就是跟著耶穌行.在這跟隨中,我們慢慢離開世俗的軌跡,進入神的軌跡.我們的視界,隨著這軌跡而開闊,而慢慢看到耶穌所看到的世界.真的,耶穌來了,他開啟了一個新世界,呈現了一個新世界,活現出一個新世界.他呼召我們,期望我們,在生活中,活出這個世界來.耶穌以一生,呈現過這個世界,也期望我們,向世人展示這世界.

今日的講道經文,最後一節是這樣的:

太5:48　「所以,你們要完全,像你們的天父完全一樣.」

這裡講的「完全」,當然不是指完美無瑕,而是指一種屬乎天父的印記.天父是完全的,屬乎祂的人,打下屬乎祂的印記.有這印記的人,應該順從天父的旨意,以天父的仁愛為目標.這種人放下了種種暴力,平安自在地在人間創造和平.就是這種人,呈現天國的真實.耶穌在登山寶訓裡,教我們以主禱文祈禱:「天父啊,願你的國降臨,願你的旨意行在地上,如同行在天上.」當我們看到耶穌生活背後的那個神聖世界,又願意跟著耶穌而行時,我們其實就是將天父的旨意行在地上,在這實踐的過程中,同時是讓天國的真相呈現人間.

弟兄姊妹,你們要完全,像你們的天父完全一樣.

願神賜福你們.

願榮耀歸予聖父,聖子,聖靈.阿們.



\newpage

\section{三項試探一個中心}
\label{sec:787}
\textbf{林豪恩先生-三項試探,一個中心}
\newline
\newline
\hyperref[sec:786]{< < < PREV SERMON < < <}
~
\hyperref[sec:toc]{[返目錄]}
~
\hyperref[sec:788]{> > > NEXT SERMON > > >}


詩篇 32:1-11
\newline
\begin{longtable}{cl}
\hline
\hline
章節 & 經文 (和合本修訂版)\\
\hline
32:1 & \begin{tabularx}{0.7\textwidth}{X} 過犯得赦免，罪惡蒙遮蓋的人有福了！ \end{tabularx} \\ \\ \relax
32:2 & \begin{tabularx}{0.7\textwidth}{X} 耶和華不算為有罪，內心沒有詭詐的人有福了！ \end{tabularx} \\ \\ \relax
32:3 & \begin{tabularx}{0.7\textwidth}{X} 我閉口不認罪的時候，因終日呻吟而骨頭枯乾。 \end{tabularx} \\ \\ \relax
32:4 & \begin{tabularx}{0.7\textwidth}{X} 黑夜白日，你的手壓在我身上沉重；我的精力耗盡，如同夏天的乾旱。（細拉） \end{tabularx} \\ \\ \relax
32:5 & \begin{tabularx}{0.7\textwidth}{X} 我向你陳明我的罪，不隱瞞我的惡。我說：「我要向耶和華承認我的過犯」；你就赦免我的罪惡。（細拉） \end{tabularx} \\ \\ \relax
32:6 & \begin{tabularx}{0.7\textwidth}{X} 為此，凡虔誠人都當趁你可尋找的時候向你禱告；大水氾濫的時候，必不臨到他。 \end{tabularx} \\ \\ \relax
32:7 & \begin{tabularx}{0.7\textwidth}{X} 你是我藏身之處，你必保佑我脫離苦難，以得救的歡呼四面環繞我。（細拉） \end{tabularx} \\ \\ \relax
32:8 & \begin{tabularx}{0.7\textwidth}{X} 我要教導你，指示你當行的路，我要定睛在你身上勸戒你。 \end{tabularx} \\ \\ \relax
32:9 & \begin{tabularx}{0.7\textwidth}{X} 你不可像那無知的騾馬，須用嚼環韁繩勒住，不然，牠就不會靠近你。 \end{tabularx} \\ \\ \relax
32:10 & \begin{tabularx}{0.7\textwidth}{X} 惡人必多受苦楚；惟獨倚靠耶和華的，必有慈愛四面環繞他。 \end{tabularx} \\ \\ \relax
32:11 & \begin{tabularx}{0.7\textwidth}{X} 義人哪，你們應當靠耶和華歡喜快樂，心裡正直的人哪，你們都當歡呼。 \end{tabularx} \\ \\
[1ex]
\hline
\hline
\end{longtable}
羅馬書 5:12-19
\newline
\begin{longtable}{cl}
\hline
\hline
章節 & 經文 (和合本修訂版)\\
\hline
5:12 & \begin{tabularx}{0.7\textwidth}{X} 為此，正如罪是從一人進入世界，死又從罪而來，於是死就臨到所有的人，因為人人都犯了罪。 \end{tabularx} \\ \\ \relax
5:13 & \begin{tabularx}{0.7\textwidth}{X} 沒有律法之前，罪已經在世上，但沒有律法，罪也不算罪。 \end{tabularx} \\ \\ \relax
5:14 & \begin{tabularx}{0.7\textwidth}{X} 然而，從亞當到摩西，死就掌了權，連那些不與亞當犯一樣罪過的，也在死的權下。亞當是那以後要來之人的預像。 \end{tabularx} \\ \\ \relax
5:15 & \begin{tabularx}{0.7\textwidth}{X} 但是過犯不如恩賜，若因一人的過犯，眾人都死了，那麼，神的恩典，與那因耶穌基督一人而來的恩典中的賞賜，豈不加倍地臨到眾人嗎？ \end{tabularx} \\ \\ \relax
5:16 & \begin{tabularx}{0.7\textwidth}{X} 因一人犯罪而來的後果，也不如賞賜，原來審判是由一人而定罪，恩賜乃是由許多過犯而稱義。 \end{tabularx} \\ \\ \relax
5:17 & \begin{tabularx}{0.7\textwidth}{X} 若因一人的過犯，死就因這一人掌權，那些受洪恩又蒙所賜之義的，豈不更要因耶穌基督一人在他們生命中掌權嗎？ \end{tabularx} \\ \\ \relax
5:18 & \begin{tabularx}{0.7\textwidth}{X} 這樣看來，因一次的過犯，所有的人都被定罪；照樣，因一次的義行，所有的人也就被稱義而得生命了。 \end{tabularx} \\ \\ \relax
5:19 & \begin{tabularx}{0.7\textwidth}{X} 因一人的悖逆，眾人成為罪人；照樣，因一人的順從，眾人也成為義了。 \end{tabularx} \\ \\
[1ex]
\hline
\hline
\end{longtable}
馬太福音 4:1-11
\newline
\begin{longtable}{cl}
\hline
\hline
章節 & 經文 (和合本修訂版)\\
\hline
4:1 & \begin{tabularx}{0.7\textwidth}{X} 當時，耶穌被聖靈引到曠野，受魔鬼的試探。 \end{tabularx} \\ \\ \relax
4:2 & \begin{tabularx}{0.7\textwidth}{X} 他禁食四十晝夜，後來就餓了。 \end{tabularx} \\ \\ \relax
4:3 & \begin{tabularx}{0.7\textwidth}{X} 那試探者進前來對他說：「你若是神的兒子，叫這些石頭變成食物吧。」 \end{tabularx} \\ \\ \relax
4:4 & \begin{tabularx}{0.7\textwidth}{X} 耶穌卻回答說：「經上記著：『人活著，不是單靠食物，乃是靠神口裡所出的一切話。』」 \end{tabularx} \\ \\ \relax
4:5 & \begin{tabularx}{0.7\textwidth}{X} 魔鬼就帶他進了聖城，叫他站在聖殿頂上， \end{tabularx} \\ \\ \relax
4:6 & \begin{tabularx}{0.7\textwidth}{X} 對他說：「你若是神的兒子，就跳下去！因為經上記著：『主要為你命令他的使者，用手托住你，免得你的腳碰在石頭上。』」 \end{tabularx} \\ \\ \relax
4:7 & \begin{tabularx}{0.7\textwidth}{X} 耶穌對他說：「經上又記著：『不可試探主—你的神。』」 \end{tabularx} \\ \\ \relax
4:8 & \begin{tabularx}{0.7\textwidth}{X} 魔鬼又帶他上了一座很高的山，將世上的萬國和萬國的榮華都指給他看， \end{tabularx} \\ \\ \relax
4:9 & \begin{tabularx}{0.7\textwidth}{X} 對他說：「你若俯伏拜我，我就把這一切賜給你。」 \end{tabularx} \\ \\ \relax
4:10 & \begin{tabularx}{0.7\textwidth}{X} 耶穌說：「撒但，退去！因為經上記著：『要拜主—你的神，惟獨事奉他。』」 \end{tabularx} \\ \\ \relax
4:11 & \begin{tabularx}{0.7\textwidth}{X} 於是，魔鬼離開了耶穌，立刻有天使來伺候他。 \end{tabularx} \\ \\
[1ex]
\hline
\hline
\end{longtable}
講題:三項試探,一個中心ThreeTrials,OneFocus

經課:創世記2:15-17,3:1-7；詩篇32:1-11；羅馬書5:12-19；馬太福音4:1-11

講員:林豪恩先生

場合:香港中文大學崇基學院禮拜堂主日崇拜

日期:2026年2月22日

一.三項試探和我們有沒有關係?

從今天的福音經課,我們又讀到耶穌的三個試探.學生團契討論耶穌的三個試探時,有同學認為和自己沒有多少關係.為甚麼呢?第一,耶穌有能力吩咐石頭變成食物,但他沒有這種能力；第二,他有畏高症,要站在殿頂跳下去,未跌死已經嚇死；第三,他覺得自己連自己都管理不好,例如要上的早堂去不到上,要交的功課未能準時交,如果把萬國交給他只會令他很苦惱.你覺得這位同學說得有道理嗎?

如果是這樣,耶穌的三項試探有甚麼意義呢?我把這個問題問AI,AI真厲害,它懂得以經解經,引用《約翰一書》2:16「肉體的情慾」,「眼目的情慾」,「今生的驕傲」來概括耶穌的三項試探.

吩咐石頭變成食物就是「肉體的情慾」

萬國和萬國的榮華就是「眼目的情慾」

從殿頂跳下就是「今生的驕傲」

聽起來似乎有道理.不過,我們也可以與AI商榷一下.第一,肚餓要食物乃人之常情,未必可以貶稱為「肉體的情慾」吧?第二,天地萬物本來就是上帝的,看一看自己的東西就等於「眼目的情慾」?第三,如果耶穌有能力從殿頂跳下去而不受傷,何以被說成是「今生的驕傲」呢?例如今年1月25日,有一位徒手攀岩運動員赤手攀登台北101,有沒有人評論他「驕傲」呢?最多都是說他魯莽,稱他為「狂人」.有媒體按這位攀岩高手的能力評估,指出對他來說攀登台北101只是「小菜一碟」.對於耶穌來說,從殿頂跳下去又算得是甚麼呢?

二.三項試探的中心是甚麼?

那麼,耶穌面對的三項試探到底在試探甚麼呢?讓我們參考經課,看看經課給我們的啟發.首先是創世記:

3:1耶和華上帝所造的,惟有蛇比田野一切的活物更狡猾.蛇對女人說:「上帝豈是真說不許你們吃園中所有樹上的果子嗎?」

3:2女人對蛇說:「園中樹上的果子,我們可以吃,

3:3惟有園當中那棵樹上的果子,上帝曾說:『你們不可吃,也不可摸,免得你們死.』」

3:4蛇對女人說:「你們不一定死；

3:5因為上帝知道,你們吃的日子眼睛就明亮了,你們便如上帝能知道善惡.」

3:6於是女人見那棵樹的果子好作食物,也悅人的眼目,且是可喜愛的,能使人有智慧,就摘下果子來吃了,又給她丈夫,她丈夫也吃了.

3:7他們二人的眼睛就明亮了,才知道自己是赤身露體,便拿無花果樹的葉子為自己編做裙子.

現代有KK園騙案,原來古代就已經有伊甸園騙案.KK園騙徒要騙的是金錢,伊甸園騙徒要騙甚麼呢?是騙情.騙甚麼情呢?是騙天父和祂兒女之間的感情.亞當和夏娃作為人類的始祖,他們是誰呢?大家還記得路加福音的族譜嗎?族譜的最後一句是甚麼呢?就是「亞當是上帝的兒子」(3:38).大家有沒有留意「上帝的兒子」這個身分在耶穌的三項試探中又再出現呢?看來魔鬼瞄準「上帝的兒子」,以「上帝的兒子」為攻擊對象.

4:3「你若是上帝的兒子,叫這些石頭變成食物吧.」

4:6「你若是上帝的兒子,就跳下去!因為經上記著:『主要為你命令他的使者,用手托住你,免得你的腳碰在石頭上.』」

4:9對他說:「你若俯伏拜我,我就把這一切賜給你.」

其實,需要食物,想玩「笨豬跳」,甚至收回天下萬國,本身不一定有問題,問題是要付出甚麼代價.DSE中文範文其中一篇是「魚我所欲也」,孟子說:「魚,我所欲也,熊掌,亦我所欲也;二者不可得兼,舍魚而取熊掌者也.生亦我所欲也,義亦我所欲也;二者不可得兼,舍生而取義者也.」耶穌的試探中有一個和食物相關,孟子也提出一個類似的例子:「一筐飯,一碗湯,得到便可以活下去,失去它就要死亡.呼喝著把食物給人吃,流浪者不會接受；踐踏後給人吃,乞丐也不屑接受.」

道成肉身的耶穌會肚餓,如常人一樣有各種生理和心理的需要.作為上帝的兒子,在滿足這些需要的時候有沒有底線呢?上帝的兒子會如何對待上帝呢?在伊甸園中,亞當和夏娃作為上帝的兒女怎樣對待上帝呢?他們被魔鬼成功地離間了和上帝的關係,使原本的信任變成猜忌,於是做出背叛上帝的行為,使原來的親密變成疏離,最後要躲避上帝.

現在,魔鬼又向上帝的兒子耶穌重施故技,魔鬼能夠離間耶穌對上帝的信任嗎?

4:4耶穌卻回答說:「經上記著說:人活著,不是單靠食物,乃是靠上帝口裡所出的一切話.」

魔鬼能夠令耶穌懷疑上帝而試探上帝嗎?

4:7耶穌對他說:「經上又記著說:『不可試探主－你的上帝.』」

魔鬼能夠令耶穌做出背叛上帝的行動嗎?

4:10耶穌說:「撒旦,退去吧!因為經上記著說:當拜主－你的上帝,單要事奉他.」

亞當和夏娃面對魔鬼試探時,任由魔鬼帶著走,似乎看不見他們的底線.耶穌面對魔鬼試探時,你看見耶穌的底線是甚麼嗎?從古至今,人們都傾向以行為守則來定立底線,於是創作出一大堆要遵守的規條,以及另一大堆不可做的事情.如此,就變成了律法主義.人們好似小朋友一樣經常問:「為甚麼他可以吃雪糕,而我不可以?」判斷好壞,有時不只是行為本身,而是和身分相關的.不知道你有沒有聽過這個故事:

一大清早,老媽衝進房間叫兒子起床:「快點起床!不要賴床!你今天還要上學呢!」

兒子死賴在被窩裡,蒙著頭說:「我不要!我不想上學!」

老媽氣得大吼:「不行!你一定要去!」

兒子哭喪著臉說:「為甚麼隔壁的小明可以不上學,而我一定要去!」

老媽一把掀開被子罵道:「再不起來就來不及了!你難道忘了你是校長嗎?!」

上帝的兒子是甚麼呢?其實,上帝的兒子不是獨立存在的,有上帝的兒子,就有上帝.上帝的兒子應該是怎樣的?這個問題也不是獨立的問題,而必須放在上帝的兒子與上帝兩者之間來思考的.於是,問題就變成:上帝的兒子和上帝的關係應該是怎樣的?關係不是我們平時所說的「感覺好不好」的情緒性描述.從魔鬼對亞當和夏娃,以及對耶穌的試探可見,這是關乎誠信.即是說,雙方是否以信任,誠實,不欺騙,不背叛來互相對待.申命記7:9用信實來描述上帝:「你要知道耶和華－你的上帝,他是上帝,是信實的上帝」.「信實」包含有可靠,穩固,誠實,忠誠不變及始終如一的意思.上帝以信實對待兒女,上帝的兒女會以信實對待上帝嗎?其實,不但是上帝和上帝的兒子之間期望以信實相待,難道人類不也都渴望任何關係都是誠實的嗎?難道人們不都是討厭欺騙和背叛的嗎?經課中的詩篇32篇就帶領我們向這個層面思考.

三．詩篇32篇大衛王的反例:王帝與子民

詩篇32篇相傳是大衛王認罪得赦後寫的詩歌.這首詩歌和今天的福音經課有甚麼關係呢?首先,也讓我們問一個身分的問題,大衛王是誰呢?是王.王是怎樣的呢?其實,王是不能獨立存在的,有王就必然有子民,因此,王是和子民存在著相互關係.如果要問王應該是怎樣的?就要放在王與子民之間的關係來思考.王與子民的關係不需要用「愛」與「不愛」這些感性的詞彙去描述,更重要的是雙方之間的誠信,是否以誠實,守信,不欺騙,不背叛來互相對待.

大衛王為甚麼要認罪呢?他犯了甚麼罪呢?撒母耳記下11-12章記載大衛王和一個女人同房,然後殺了一個男人.古代君王後宮三千,和一個女人同房有甚麼罪呢?古代君王操生殺大權,殺一個人又算得是甚麼呢?聖經指出,那個女人是烏利亞的妻子拔示巴,那個另人是拔示巴的丈夫烏利亞.烏利亞是誰呢?烏利亞如何對待大衛王呢?

11:8大衛(召烏利亞從戰場上回來)說:「你回家去,洗洗腳吧!」

11:11烏利亞對大衛說:「約櫃和以色列與猶大兵都住在棚裡,我主約押和我主的僕人都在田野安營,我豈可回家吃喝,與妻子同寢呢?我敢在王面前起誓:我決不行這事!」

烏利亞以誠實,守信,不欺騙,不背叛來對待大衛王,國家,以及同伴,大衛王如何對待對這位忠心的戰士呢?

11:14次日早晨,大衛寫信與約押,交烏利亞隨手帶去.

11:15信內寫著說:「要派烏利亞前進,到陣勢極險之處,你們便退後,使他被殺.」

沒有人否定大衛王有性需要,但有需要就可以召別人的妻子來洩慾的嗎?烏利亞是對王,對國家忠誠的戰士,大衛王竟然暗殺他?君王和子民之間的關係是這樣的嗎?大衛王犯的是甚麼罪呢?謀殺只是罪名,更令天地不容的是這位君王失去誠信,出賣和背叛子民.不知道你讀這個故事時有沒有感到憤怒和悲傷呢?信實,不騙欺,不出賣是人類社會的共識,人間多少災難,痛苦和悲劇,就是因為放棄誠信,欺騙和背叛而產生的.

大廈維修公司應該是怎樣的呢?工程顧問公司應該是怎樣的呢?安全督導應該是怎樣的呢?物料供應商應該是怎樣的呢?如果你把你居住的大廈交給他們維修,上述的問題,都不是獨立的問題,而是相對於居住在大廈裡的居民而言的.大廈的居民不但把金錢交給他們,甚至把性命也交給他們,他們應該如何看待把金錢和性命都交付他們的居民呢?如果這些公司和專業人士破壞與居民之間的誠信,欺騙居民,背叛居民,所產生的災難和悲劇我們都看見了；所產生的痛苦,很多居民現在仍在承受.

一個人應該是怎樣的,需要放在關係之中思考,因為一個人和其他有關係的人是存在著誠信的基礎.古人說,「食色性也」,親密是人的需要,性滿足也是人的需要,問題是:你是誰?對方又是誰?

著名基督教作家PhilipYancey(楊腓力),最近發表聲明,承認與一名已婚女子保持了長達八年的婚外情.他的聲明部分內容如下:

「我的行為違背了我對婚姻的一切信念,也與我的信仰和著述完全背道而馳,給她的丈夫和我們雙方的家庭都帶來了深深的痛苦.我在道德和靈性上都犯了錯,我為自己造成的傷害感到無比悲痛.我意識到我的行為會讓那些曾經信任我作品的讀者感到失望.更糟的是,我的罪羞辱了上帝.我充滿了悔恨和懺悔,除了上帝的憐憫和恩典,我別無依靠.」(ChristianityToday　6/1/2026)

一個男人和一個女人發展感情關係本身不是問題,如果把這段感情放回那個男人和那個女人的關係網中,就出現PhilipYancey所承認的「給她的丈夫和我們雙方的家庭都帶來了深深的痛苦」；如果把這段感情放進那個男人和信任他的讀者的關係中,就是「讓那些曾經信任我作品的讀者感到失望」；如果把那段感情放進那個男人和上帝的關係中,就是「我的罪羞辱了上帝」.

四．有沒有人能夠在關係中活出信實呢?

知道自己是誰,誠信地活出自己,是非常難得的狀況,也是帶有拯救效果的狀況,世間上有人能夠如此活著嗎?使徒保羅認為非常罕有,他在今天的書信經課羅馬書第5章中指出,有一位:

5:18如此說來,因一次的過犯,眾人都被定罪；照樣,因一次的義行,眾人也就被稱義得生命了.

5:19因一人的悖逆,眾人成為罪人；照樣,因一人的順從,眾人也成為義了.

這位完全知道自己是誰,活出上帝的兒子的耶穌,其拯救的效果是普及全人類.我們不可能是百分之一百知道自己是誰,也未能百分之一百活出自己,然而,若能達到某個程度,就已經帶來某個程度的拯救效果了.

例如在火災之中,縱使所有人都逃離火場,有些人知道自己不可以逃離,反而要衝進火場,他們是誰呢?就是消防員.在火災被困的居民在等候誰呢?就是消防員.由此可見,在那等別的時刻,消防員和災民之間是存在著這種特別關係的.當消防員活出他的身分的時候,火場中的居民得到拯救.

我們不是消防員,但每個人都活在與其他人有關係的身分之中,我們如何信實地活出自己,不背叛,不出賣呢?

天主教於去年12月冊封一位阿根廷平信徒商人為地位僅次於聖人的「真福」.這位商人何以獲得冊封呢?1961年,他的其中一家公司被售予美國的信託基金,信託基金其後決定裁掉1,200名員工.當時已罹患末期癌症的商人強烈反對裁員,並提出保留所有員工的改革計劃.他的友人表示,在他接受癌症治療,需要輸血的時候,在其工廠打工的人曾在醫院外排隊輸血,人數之多令醫護大感驚訝,一度以為他是工會成員.(時代論壇/2025年12月30日)

老闆是誰呢?老闆是怎樣的呢?老闆是存在於和員工的關係之中,當雙方以信實相待,雙方都得到拯救.



\newpage

\section{上帝愛世界}
\label{sec:788}
\textbf{高國雄牧師-上帝愛世界}
\newline
\newline
\hyperref[sec:787]{< < < PREV SERMON < < <}
~
\hyperref[sec:toc]{[返目錄]}
~
\hyperref[sec:789]{> > > NEXT SERMON > > >}


約翰福音 3:1-17
\newline
\begin{longtable}{cl}
\hline
\hline
章節 & 經文 (和合本修訂版)\\
\hline
3:1 & \begin{tabularx}{0.7\textwidth}{X} 有一個法利賽人，名叫尼哥德慕，是猶太人的官。 \end{tabularx} \\ \\ \relax
3:2 & \begin{tabularx}{0.7\textwidth}{X} 這人夜裡來見耶穌，對他說：「拉比，我們知道你是由神那裡來作老師的；因為你所行的神蹟，若沒有神同在，無人能行。」 \end{tabularx} \\ \\ \relax
3:3 & \begin{tabularx}{0.7\textwidth}{X} 耶穌回答他說：「我實實在在地告訴你，人若不重生，就不能見神的國。」 \end{tabularx} \\ \\ \relax
3:4 & \begin{tabularx}{0.7\textwidth}{X} 尼哥德慕對他說：「人已經老了，如何能重生呢？豈能再進母腹生出來嗎？」 \end{tabularx} \\ \\ \relax
3:5 & \begin{tabularx}{0.7\textwidth}{X} 耶穌回答：「我實實在在地告訴你，人若不是從水和聖靈生的，就不能進神的國。 \end{tabularx} \\ \\ \relax
3:6 & \begin{tabularx}{0.7\textwidth}{X} 從肉身生的就是肉身；從靈生的就是靈。 \end{tabularx} \\ \\ \relax
3:7 & \begin{tabularx}{0.7\textwidth}{X} 我說『你們必須重生』，你不要驚訝。 \end{tabularx} \\ \\ \relax
3:8 & \begin{tabularx}{0.7\textwidth}{X} 風隨著意思吹，你聽見風的聲音，卻不知道是從哪裡來，往哪裡去；凡從聖靈生的也是如此。」 \end{tabularx} \\ \\ \relax
3:9 & \begin{tabularx}{0.7\textwidth}{X} 尼哥德慕問他：「怎麼能有這些事呢？」 \end{tabularx} \\ \\ \relax
3:10 & \begin{tabularx}{0.7\textwidth}{X} 耶穌回答，對他說：「你是以色列人的老師，還不明白這些事嗎？ \end{tabularx} \\ \\ \relax
3:11 & \begin{tabularx}{0.7\textwidth}{X} 我實實在在地告訴你，我們所說的是我們知道的，我們所見證的是我們見過的，你們卻不領受我們的見證。 \end{tabularx} \\ \\ \relax
3:12 & \begin{tabularx}{0.7\textwidth}{X} 我對你們說地上的事，你們尚且不信，若對你們說天上的事，如何能信呢？ \end{tabularx} \\ \\ \relax
3:13 & \begin{tabularx}{0.7\textwidth}{X} 除了從天降下的人子，沒有人升過天。 \end{tabularx} \\ \\ \relax
3:14 & \begin{tabularx}{0.7\textwidth}{X} 摩西在曠野怎樣舉蛇，人子也必須照樣被舉起來， \end{tabularx} \\ \\ \relax
3:15 & \begin{tabularx}{0.7\textwidth}{X} 要使一切信他的人都得永生。 \end{tabularx} \\ \\ \relax
3:16 & \begin{tabularx}{0.7\textwidth}{X} 「神愛世人，甚至將他獨一的兒子賜給他們，叫一切信他的人不致滅亡，反得永生。 \end{tabularx} \\ \\ \relax
3:17 & \begin{tabularx}{0.7\textwidth}{X} 因為神差他的兒子到世上來，不是要定世人的罪，而是要使世人因他得救。 \end{tabularx} \\ \\
[1ex]
\hline
\hline
\end{longtable}
講題:上帝愛世界GodsoLovedtheWorld

經課:約翰福音3章1至17節

講員:高國雄牧師

場合:香港中文大學崇基學院禮拜堂主日崇拜

日期:2026年3月1日

約翰福音3:1-17

1有一個法利賽人,名叫尼哥德慕,是猶太人的官.2這人夜裡來見耶穌,說:「拉比,我們知道你是由上帝那裡來作師傅的；因為你所行的神蹟,若沒有上帝同在,無人能行.」3耶穌回答說:「我實實在在地告訴你,人若不重生,就不能見上帝的國.」4尼哥德慕說:「人已經老了,如何能重生呢?豈能再進母腹生出來嗎?」5耶穌說:「我實實在在地告訴你,人若不是從水和聖靈生的,就不能進上帝的國.6從肉身生的就是肉身；從靈生的就是靈.7我說:『你們必須重生』,你不要以為希奇.8風隨著意思吹,你聽見風的響聲,卻不曉得從哪裡來,往哪裡去；凡從聖靈生的,也是如此.」9尼哥德慕問他說:「怎能有這事呢?」10耶穌回答說:「你是以色列人的先生,還不明白這事嗎?11我實實在在地告訴你,我們所說的是我們知道的；我們所見證的是我們見過的；你們卻不領受我們的見證.12我對你們說地上的事,你們尚且不信,若說天上的事,如何能信呢?13除了從天降下,仍舊在天的人子,沒有人升過天.14摩西在曠野怎樣舉蛇,人子也必照樣被舉起來,15叫一切信他的都得永生.

16「上帝愛世人,甚至將他的獨生子賜給他們,叫一切信他的,不致滅亡,反得永生.17因為上帝差他的兒子降世,不是要定世人的罪,乃是要叫世人因他得救.

在當代語境中,「上帝愛世人」或許是基督教流傳最廣,卻也最被誤解的一句話.對於習慣批判性思考的我們來說,「愛」這個詞往往顯得過於感性,甚至在充滿苦難,荒謬與不公的世界面前,顯得有些軟弱無力.

又或已經變得十分空泛的宗教用語…

大齋期(Lent)是一個曠野的季節.它要求我們剝離表面的繁華,直視生命的幽暗與脆弱.

在第二主日,我們跟隨耶穌的腳步走向耶路撒冷,這是一個走向死亡的旅程.

當我們在這個背景下談論「上帝愛世界」,我們談論的不是一種溫馨的情感安慰,而是一場神學上的冒險與更新.

深夜裡的知識分子

V1有一個法利賽人,名叫尼哥德慕,是猶太人的官.

今日的經文始於一個深夜.尼哥德慕,當時的知識分子代表——他是法利賽人,是拉比,擁有嚴謹的邏輯與神學架構,一位擁有極高學術地位與宗教權力的知識分子,在黑暗中走向耶穌.

在當下高速生活中,許多人也正處於某種「深夜」:對大時代的無力感,對前景的迷惘,以及對既有知識體系崩潰的焦慮.

我們帶著「我們知道」的自信而來,卻發現自己陷入了更深的生命困境.

理性的邊界與「由上而生」

知識的困境

V2這人夜裡來見耶穌,說:「拉比,我們知道你是由上帝那裡來作師傅的；因為你所行的神蹟,若沒有上帝同在,無人能行.」

尼哥德慕開場便說:「拉比,我們知道……」.

這是他的困境:他試圖用既有的知識框架(「我們知道你是由上帝那裡來作老師的」)來定義耶穌.他追求的是一套解釋世界的理論,試圖用嚴謹的因果邏輯來定義耶穌.

對知識分子而言,「理解」是掌控世界的方式.

V3耶穌回答說:「我實實在在地告訴你,人若不重生,就不能見上帝的國.」

但耶穌隨即打破了他的框架:「人若不重生(或譯:由上而生),就不能見上帝的國.」重生不是知識的增長,而是轉向.耶穌指出,認識上帝不是靠觀察與分析(學術研究),而是靠生命的徹底重整.

更要讓我們明白,生命中追求的不是要「做師傅」,「行大能」,做偉大的事.

而是要見(進入)上帝的國,這是一切的首要!

太6:33你們要先求他的國和他的義,這些東西都要加給你們了.

風的靈性

V4尼哥德慕說:「人已經老了,如何能重生呢?豈能再進母腹生出來嗎?」5耶穌說:「我實實在在地告訴你,人若不是從水和聖靈生的,就不能進上帝的國.6從肉身生的就是肉身；從靈生的就是靈.7我說:『你們必須重生』,你不要以為希奇.8風隨著意思吹,你聽見風的響聲,卻不曉得從哪裡來,往哪裡去；凡從聖靈生的,也是如此.」

尼哥德慕繼續用他熟練的邏輯追問…

耶穌的回應,重生的途徑,只有透過水及聖靈,這唯一的途徑,超越尼哥德慕的想像,也超過我們的想像.

耶穌又用「風」來比喻聖靈.風不可見,卻有其影響力.在凡事追求精準計畫,數據指標(KPI)的香港社會,耶穌提醒我們:上帝的愛與工作,往往在我們專業邏輯的「不可測」之處.

愛世界,首先要承認上帝的主權高於我們的理性預測.

世界與人的掙扎

地上的事與天上的事

V11我實實在在地告訴你,我們所說的是我們知道的；我們所見證的是我們見過的；你們卻不領受我們的見證.12我對你們說地上的事,你們尚且不信,若說天上的事,如何能信呢?13除了從天降下,仍舊在天的人子,沒有人升過天.

耶穌責備尼哥德慕:「我對你們說地上的事,你們尚且不信……」這反映了知識我們的通病:我們過於沉溺於地上的結構,政治與邏輯,以至於失去了「屬靈的想像力」.

銅蛇的預表

V14摩西在曠野怎樣舉蛇,人子也必照樣被舉起來,15叫一切信他的都得永生.

耶穌引用民數記銅蛇的典故.當時以色列人在曠野因怨毒而死,仰望銅蛇便得活.

我們今天的生活,可能也讓我們感到如處身曠野,充滿了怨恨與毒素.

耶穌強調,上帝拯救世界的方式,不是透過一套完美的社會理論,而是透過一個「被舉起來」的受苦者.

這對崇尚權力與效率的知識分子來說,是極大的冒險與諷刺.

上帝愛「世界」

V16「上帝愛世人,甚至將他的獨生子賜給他們,叫一切信他的,不致滅亡,反得永生.

這是金句中的金句.

「上帝愛世人」中的「世人/世界」,在約翰福音中通常具有負面含義,指涉那個不認識上帝,甚至是敵對上帝的墮落體系.

上帝愛的不是一個「完美的世界」,而是一個「反叛的世界」.

愛不是一個可以被論證的命題,而是一個主動介入的事件.對我們而言,承認「我不知道」才是重生的開始.

上帝愛世界,是愛那個「真實卻殘缺」的世界,而不是我們頭腦中那個「理應完美」的理想國.

祂所愛的不是抽象的「人類」,而是每個具體,有罪,在體系中掙扎的(可能不認識,拒絕,甚至敵擋上帝的「世人」.

上帝的愛不是溺愛,而是透過「賜下獨生子」這場代價極重的介入,將世界從混亂中拯救出來.

不是要審判,而是要得救

V17因為上帝差他的兒子降世,不是要定世人的罪,乃是要叫世人因他得救.

這節經文對當下至關重要.

在充滿指責與審判氛圍的社會中,經文宣告:

上帝差子降世,原意不是要定世界的罪.

無錯,這份愛本身就應該帶有審判性.

當光進入世界,黑暗若拒絕光,這拒絕本身就是審判.

但上帝選擇了愛.

耶穌基督道成肉身的愛.

若上帝對這破碎的世界仍存憐憫而非單純的審判,我們這群自詡擁有知識與洞察力的人,是否也應放下那種「旁觀者的審判感」,轉而參與上帝那救贖性的,修復性的愛?

結論:從夜間走向白日

尼哥德慕的故事沒有結束在第三章.

約19:39-40又有尼哥德慕,就是先前夜裡去見耶穌的,帶著沒藥和沉香約有一百斤前來.他們就照猶太人殯葬的規矩,把耶穌的身體用細麻布加上香料裹好了.

他在約翰福音結尾時,於光天化日之下站出來安葬耶穌.

尼哥德慕開始在深夜偷偷而來,最後卻在白日為耶穌安葬.

這是一個知識分子從「旁觀」到「委身」的過程.

上帝愛世界,這份愛邀請我們用行動回應.

承認我們的專業與理性無法解決生命的荒涼,唯有「由上而生」.

在各自的崗位上,效法基督「道成肉身」的介入,守護真理,守護每一個具體的「世人」.

在大齋期的荒涼中,不再只看社會的崩壞,而是仰望那被舉起的基督,相信祂的愛正在聖別這座城市.

願我們不再只在深夜徘徊,而是在基督的愛中,得著重生的勇氣.

因為上帝愛世界,所以我們沒有理由絕望.

因為上帝愛世界,所以我們有責任去愛.

願主賜我們力量,勇氣,盼望與及愛,阿們.



\newpage

\section{活水之道}
\label{sec:789}
\textbf{李均熊牧師 - 活水之道}
\newline
\newline
\hyperref[sec:788]{< < < PREV SERMON < < <}
~
\hyperref[sec:toc]{[返目錄]}
~
\hyperref[sec:790]{> > > NEXT SERMON > > >}


約翰福音 4:5-42
\newline
\begin{longtable}{cl}
\hline
\hline
章節 & 經文 (和合本修訂版)\\
\hline
4:5 & \begin{tabularx}{0.7\textwidth}{X} 於是到了撒瑪利亞的一座城，名叫敘加，靠近雅各給他兒子約瑟的那塊地。 \end{tabularx} \\ \\ \relax
4:6 & \begin{tabularx}{0.7\textwidth}{X} 雅各井就在那裡；耶穌因旅途疲乏，坐在井旁。那時約是正午。 \end{tabularx} \\ \\ \relax
4:7 & \begin{tabularx}{0.7\textwidth}{X} 有一個撒瑪利亞婦人來打水。耶穌對她說：「請給我水喝。」 \end{tabularx} \\ \\ \relax
4:8 & \begin{tabularx}{0.7\textwidth}{X} 因為那時門徒進城買食物去了。 \end{tabularx} \\ \\ \relax
4:9 & \begin{tabularx}{0.7\textwidth}{X} 撒瑪利亞婦人對他說：「你是猶太人，怎麼向我一個撒瑪利亞女人要水喝呢？」因為猶太人和撒瑪利亞人沒有來往。 \end{tabularx} \\ \\ \relax
4:10 & \begin{tabularx}{0.7\textwidth}{X} 耶穌回答她說：「你若知道神的恩賜，和對你說『請給我水喝』的是誰，你早就會求他，他也早就會給了你活水。」 \end{tabularx} \\ \\ \relax
4:11 & \begin{tabularx}{0.7\textwidth}{X} 婦人對耶穌說：「先生，你沒有打水的器具，井又深，哪裡去取活水呢？ \end{tabularx} \\ \\ \relax
4:12 & \begin{tabularx}{0.7\textwidth}{X} 我們的祖宗雅各把這井留給我們，他自己和兒女以及牲畜都喝這井裡的水，難道你比他還大嗎？」 \end{tabularx} \\ \\ \relax
4:13 & \begin{tabularx}{0.7\textwidth}{X} 耶穌回答，對她說：「凡喝這水的，還要再渴； \end{tabularx} \\ \\ \relax
4:14 & \begin{tabularx}{0.7\textwidth}{X} 誰喝我所賜的水，就永遠不渴。我所賜的水要在他裡面成為泉源，直湧到永生。」 \end{tabularx} \\ \\ \relax
4:15 & \begin{tabularx}{0.7\textwidth}{X} 婦人對他說：「先生，請把這水賜給我，使我不渴，也不用到這裡來打水。」 \end{tabularx} \\ \\ \relax
4:16 & \begin{tabularx}{0.7\textwidth}{X} 耶穌對她說：「你去，叫你的丈夫，再到這裡來。」 \end{tabularx} \\ \\ \relax
4:17 & \begin{tabularx}{0.7\textwidth}{X} 婦人回答，對耶穌說：「我沒有丈夫。」耶穌說：「你說沒有丈夫是對的。 \end{tabularx} \\ \\ \relax
4:18 & \begin{tabularx}{0.7\textwidth}{X} 你已經有過五個丈夫，你現在有的並不是你的丈夫。你這話是真的。」 \end{tabularx} \\ \\ \relax
4:19 & \begin{tabularx}{0.7\textwidth}{X} 婦人對他說：「先生，我看你是一位先知。 \end{tabularx} \\ \\ \relax
4:20 & \begin{tabularx}{0.7\textwidth}{X} 我們的祖宗在這山上敬拜神，你們倒說，應當敬拜的地方是在耶路撒冷。」 \end{tabularx} \\ \\ \relax
4:21 & \begin{tabularx}{0.7\textwidth}{X} 耶穌對她說：「婦人，你要信我。時候將到，你們敬拜父，既不在這山上，也不在耶路撒冷。 \end{tabularx} \\ \\ \relax
4:22 & \begin{tabularx}{0.7\textwidth}{X} 你們所敬拜的，你們不知道；我們所敬拜的，我們知道，因為救恩是從猶太人出來的。 \end{tabularx} \\ \\ \relax
4:23 & \begin{tabularx}{0.7\textwidth}{X} 時候將到，現在就是了，那真正敬拜父的，要用心靈和誠實敬拜他，因為父要這樣的人敬拜他。 \end{tabularx} \\ \\ \relax
4:24 & \begin{tabularx}{0.7\textwidth}{X} 神是靈，所以敬拜他的必須用心靈和誠實敬拜他。」 \end{tabularx} \\ \\ \relax
4:25 & \begin{tabularx}{0.7\textwidth}{X} 婦人對他說：「我知道彌賽亞—就是那稱為基督的—要來；他來了，會把一切的事都告訴我們。」 \end{tabularx} \\ \\ \relax
4:26 & \begin{tabularx}{0.7\textwidth}{X} 耶穌對她說：「我就是，正在跟你說話呢！」 \end{tabularx} \\ \\ \relax
4:27 & \begin{tabularx}{0.7\textwidth}{X} 正在這時，門徒回來了。他們對耶穌正在和一個婦人說話感到驚訝，可是沒有人說：「你要甚麼？」或說：「你為甚麼和她說話？」 \end{tabularx} \\ \\ \relax
4:28 & \begin{tabularx}{0.7\textwidth}{X} 那婦人留下水罐，往城裡去，對眾人說： \end{tabularx} \\ \\ \relax
4:29 & \begin{tabularx}{0.7\textwidth}{X} 「你們來看！有一個人把我素來所做的一切事都說了出來，難道這個人就是基督嗎？」 \end{tabularx} \\ \\ \relax
4:30 & \begin{tabularx}{0.7\textwidth}{X} 他們就出城，來到耶穌那裡。 \end{tabularx} \\ \\ \relax
4:31 & \begin{tabularx}{0.7\textwidth}{X} 就在這個時候，門徒求耶穌說：「拉比，請吃吧。」 \end{tabularx} \\ \\ \relax
4:32 & \begin{tabularx}{0.7\textwidth}{X} 耶穌對他們說：「我有食物吃，是你們不知道的。」 \end{tabularx} \\ \\ \relax
4:33 & \begin{tabularx}{0.7\textwidth}{X} 門徒就彼此說：「難道有人拿甚麼給他吃了嗎？」 \end{tabularx} \\ \\ \relax
4:34 & \begin{tabularx}{0.7\textwidth}{X} 耶穌對他們說：「我的食物就是要遵行差我來那位的旨意，完成他的工作。 \end{tabularx} \\ \\ \relax
4:35 & \begin{tabularx}{0.7\textwidth}{X} 你們不是說『到收割的時候還有四個月』嗎？我告訴你們，舉目向田觀看，莊稼熟了，可以收割了。 \end{tabularx} \\ \\ \relax
4:36 & \begin{tabularx}{0.7\textwidth}{X} 收割的人已經得工錢，為永生儲存五穀，使撒種的和收割的一同快樂。 \end{tabularx} \\ \\ \relax
4:37 & \begin{tabularx}{0.7\textwidth}{X} 『那人撒種，這人收割』，這話可見是真的。 \end{tabularx} \\ \\ \relax
4:38 & \begin{tabularx}{0.7\textwidth}{X} 我差你們去收你們所沒有辛勞的；別人辛勞，你們享受他們辛勞的成果。」 \end{tabularx} \\ \\ \relax
4:39 & \begin{tabularx}{0.7\textwidth}{X} 那城裡有好些撒瑪利亞人信了耶穌，因為那婦人作見證，說：「他把我素來所做的一切事都說了出來。」 \end{tabularx} \\ \\ \relax
4:40 & \begin{tabularx}{0.7\textwidth}{X} 於是撒瑪利亞人來見耶穌，求他在他們那裡住下，他就在那裡住了兩天。 \end{tabularx} \\ \\ \relax
4:41 & \begin{tabularx}{0.7\textwidth}{X} 因為耶穌的話，信的人就更多了。 \end{tabularx} \\ \\ \relax
4:42 & \begin{tabularx}{0.7\textwidth}{X} 他們對那婦人說：「現在我們信，不再是因為你的話，而是我們親自聽見了，知道這人真是世界的救主。」 \end{tabularx} \\ \\
[1ex]
\hline
\hline
\end{longtable}
講題:活水之道 The Way of Living Water

經課:約翰福音4:5-42

講員:李均熊牧師

場合:香港中文大學崇基學院禮拜堂主日崇拜

日期:2026年3月8日

耶穌同撒瑪利亞婦人談道,我相信是一個大家都好熟識ge 故事.正因為這個原因,我今天又再嘗試以另外的一些角度去重述這個在約翰福音的記載.而我今日選擇的視角,是和這個婦人住在同一城巿,亦即撒馬利亞地區中,敘加城的居民,希望透過他們的理解.去重新給我們新鮮的角度去思考耶穌和這一個撒瑪尼亞婦人在井旁談道,對我們今日的意思.不過開始之前,讓我交待一下,在三代經課中的甲年,福音經課都主要都是選馬太福音,所以有我們又會叫馬太年,但由大齋期第二主日的,到大齋期第五主日的,福音經課均是選用約翰福音的經文,而且這些經文都是記載耶穌與某人的相遇.上星期大齋期第二主日,在夜間,於耶穌在耶路撒冷過節的時候,一位德高望重的法利賽人去探問耶穌,今天大齋期第三主日則是耶穌在正午,在井旁,與一個曾有四個丈夫,但現在和她一起的卻不是她丈夫的撒瑪利亞婦人相遇.下主日第四主日,則是耶穌在西羅亞池附近遇見一個生來瞎眼的人,使他重見光明；最後在大齋期第五主日,則耶穌在墳地中與拉撒路和他的家姐相遇.這個設計其實很有心思,不過今天沒有幾會詳談.有些在復活節與行洗禮儀節的教會,有時會用福音經課的篇排,來顯示受洗者由疑惑到相信到復活的靈性旅程.也可以值得我們參考下,尤其由上星期尼歌德慕那份黑夜查問耶穌,到完結時仍未知他的信心狀況,到今日係正午耶穌主動引導井旁婦人,至使敘加城的居民承認耶穌就是全人類的救主,顯出信仰歷程的多種可能性.

以下是以撒瑪利亞敘加城居民視角重述約翰福音 4:5–42的記載:

朋友,

我唔係耶穌,唔係門徒,唔係當日同耶穌談話ge 撒瑪利亞婦人.

我係敘加城裡ge 其中一個人——

就係個D 經常被視為「背景人物」gwo 種人.可能你有時都會有種咁ge 感覺,就係你唔係主角,係世界呢盤大棋中只係一個普普通通唔起眼ge 普通路人甲乙丙.

但我想話俾你知——

gwo 日,井邊發生ge 事,

當初改變ge 唔止係gwo 個女人,就係連我地全城ge 人都被卷入當中ge 改變.

而家,等我由我ge 角度,

講返gwo 件令我地全城震撼ge 事出黎.

我地城入面每個人都識得佢.

你地叫佢「撒瑪利亞婦人」又或者「井旁的女人」(都唔知點解福音書無記載佢個名)

但我地叫佢「gwo 個女人」,甚至更唔好聽ge 稱呼……

講真,

佢係我地城裡最唔想遇見ge 人.

感情史複雜,名聲唔好,

講佢兩句都唔會有人幫佢講好說話.但其實我地都知,佢係一個苦命ge 女人,係我地生活ge 日子,一個女人無jor 男人的依靠,根本就生存唔到.佢年輕時應該都係一般女仔咁,好早結婚,但係佢同時間好早做jor 寡婦.佢無一般女士咁樣守寡,好快又結婚,可惜個男人對佢唔好,成日話佢腳頭唔好,克死以前個丈夫,最後仲因為小小事就休jor 佢.之後我諗佢可能都試jor 幾次婚姻,我都唔知佢結過幾次婚,最後佢應該對男人無乜希望,依家跟jor 個肯包養佢ge 有婦之夫,叫做生活有D 保障,唔使去做「個行」.但係佢就成為jor 敘加城居民茶餘飯後ge 話柄,而家我地諗返轉頭,果陣時冇憐憫佢ge 不幸,反而仲係佢ge 傷口里面加鹽,真係有D 唔arm ……

咁呢個敘加城居民對佢ge 睇法,佢自己都知,所以佢總係揀正午去打水——

gwo 個最熱,最冇人會出現ge 時候.

但gwo 日完全唔同.

佢冇再避開人,

佢跑返黎搵我地!

仲大聲喊住話:「你地黎睇o下——我遇見jor 一個人!」

我地當時ge 反應係:

「咩事?佢又攪jor D 咩ye ?」

「跑成咁,即係點?」

但,佢雙眼發光……

gwo 份興奮,gwo 種自由,gwo 份毫無隱藏ge 勇氣——

我地從未見過.

佢以前係避世ge ,

但今日佢係衝住世間黎去做見證.

朋友,

你有冇見過一個人突然變得「好自由」?

自由到你都唔認得佢?又好似好熟識,變到好陌生?

gwo 一刻,我地就係咁.

佢氣喘喘咁同我地講:

「有一個人……

將我以前ge 事全部講晒出黎……

但佢唔係嫌棄我,

佢係接納我,理解我……

難道佢就係基督?」

你知唔知,

佢講「我以前所做ge 事」,

我地全城都知佢講緊乜.而且口中講ge 基督,係唔單止猶太人,亦係我地撒瑪亞人祈盼上主要派黎係我地而家不堪境況拯救我地ge 救主.係方面,呢個女人特別虔誠,我地基利心山ge 聖殿無耶路撒冷gwo 個咁多規限,所以好似佢咁ge 女人,都可以係我地ge 聖殿供獻.佢突然話有個知哂佢一齊ge 人,仲要係基督,我地都好疑惑.

但佢講ge 時候,

佢唔再羞愧.

佢冇再閃避.

佢面容甚至帶住D 光.

以前ge 佢,怕人講.

今日ge 佢,生命好似被打開jor 咁.

有D 人開始竊竊私語:

「……究竟發生咩事?」

「佢遇到邊個?」

「邊個可以令一個人變成咁?」

而我心裡面ge 問題係:

「究竟佢遇到個基督點能夠使佢有咁大改變?」

於是,我地決定親自去睇——到底佢講ge 係真定假」

我地其實唔係好完全信佢.

佢人生咁多問題,

突然咁熱情,聽日又唔知會點.

但佢gwo 份「被觸動」ge 真實,

係假唔到.

於是我地話:

「不如自己走去睇o下.」

然後,我就見到佢——

坐係井邊ge 耶穌.

佢睇住我地,

唔係用猶太人平時睇撒瑪利亞人ge 厭惡眼神；

而係用一種好似對住朋友ge 目光咁望向我地——

我後黎都有問返gwo 個女人,問佢你見到個猶太男人同你講ye 你唔驚架咩?佢同我講,佢行到井邊發現個男人時,其實想即刻拎起個水罐走開,但個男人突然問佢lor 水飲,於是佢先至停落黎,一來出於好奇,二來呢個男人有求於佢,佢比人拒絕得多,所以都想幫o下佢,於是問jor 佢一句,「你既是猶太人,怎麼向我一個撒馬利亞婦人要水喝呢?」點知佢同呢個佢後黎知道叫耶穌ge 猶太人ge 對話就越黎越離奇.開頭耶穌同佢講會可比佢活水,佢以為佢講緊佢知道係邊度有個水泉,D 水會流動ge 「活」水,咁若果耶穌話比佢聽邊度有個咁ge 活泉,咁佢就唔使咁辛苦日日為jor 避開D 人正午過黎打水.但耶穌突然叫佢叫自己個丈夫過埋黎,佢就打打jor 個突,無細想之下回jor 句「我無丈夫」,點知耶穌就將佢以住有五個丈夫,而家同佢係一齊ge 唔係佢丈夫ge 呢D 事講哂出黎.聽到呢度,我見佢唔再無掩飾自己ge 過去,羞恥被翻轉成見證.故事無變,但說故事的人變jor .就好似佢將個水罐留jor 係井旁咁,同時將過去ge 羞愧留jor 係井旁,留jor 係耶穌身邊.而佢裡面就有活水流出黎.點解我咁清楚佢留jor 個水罐係井旁?因為耶穌離開jor 敘加之後,佢有一次黎問我借水罐用,話佢個日之後搵唔返個水罐.我有問佢點解會留底jor 個水罐係井旁,佢就話其實有兩個答案,一係其實佢真係唔記得,因為佢gwo 刻個心好亂,好興奮,好似有人心入面打開jor 扇窗,好似有陣風湧jor 入黎,就一時唔記得jor 個水罐,後黎返到城裡面,見到大家一齊出去搵耶穌ge 時候,佢就覺得其實留低水罐好似係一個記號—我黎井旁係為jor 裝水；離開時,生命入面有耶穌所賜ge 活水.我以前靠個罐,靠自己揹水返屋企；而家—我可以唔再靠gwo D 去解渴.所以,忘記同釋放,其實一齊發生；耶穌ge 釋放,就發生係人最自然ge 反應入面.你以為係唔記得,其實係放低jor .咁我就問,咁你又問我借水罐用?佢話係因為活水唔會取代責任.我仍然打水,煮飯,做人；但我而家唔係靠水罐維生,係靠基督賜ge 活水活著.— 我會還番個罐,亦會繼續講我點樣遇見基督.佢咁講,我好似明白多jor :屬靈經歷唔係逃離生活,而係更新生活.活水＝心裡的泉源；水罐＝日常的責任.兩者係並行ge .

後黎耶穌同我地講jor 好多,

但有一句我一世唔會忘記——

「救恩係比全世界ge .」

唔再分猶太,撒瑪利亞,

唔再分男人,女人,好人,壞人,正常人,失敗人.

祂看穿我地,

但冇一句責備,

只有接納,只有邀請.

我終於明白——

gwo 個女人點解改變得咁快.

佢見到ge 耶穌,

原來我地都需要.

佢gwo 份自由,gwo 份光亮,gwo 份無需再掩飾ge 勇氣,

原來就係「活水」ge 果效

耶穌講ge 「活水」,

唔係淨係比佢——

係比我地每一個城裡ge 人.

我地城裡ge 好多時自以為比呢個女人好,

比呢個女人乾淨,

比呢個女人幸福安定.

但其實,

呢個女人比我地更早「嚐到」真正ge 恩典.

我地擁有禮儀,擁有傳統,擁有對神ge 知識,

但佢擁有一樣ye ——

她真誠地承認自己的口渴.

而耶穌最先回應ge ,

就係真心承認自己心靈「乾渴」ge 人.

起初,我地因為佢ge 見證而去；

後來,我地親自聽見耶穌.

最後,我地只能講一句:

「祂唔單止係妳ge 救主,

祂係我地全城ge 救主.」

佢生命被翻轉,

原來係為jor 讓我地一齊被翻轉.

佢留下水罐,

原來係要提醒我地:

我地城裡每個人都有自己揹住ge 水罐.

耶穌能解佢ge 渴,

就能解我地ge 渴.

而gwo 日我明白ge 係:

原來被改變ge ,不止佢一個.

耶穌係黎改變我地全部.

朋友,你ge 「視角」係咩?

今日,

係你生命裡,

你可能以為自己係「旁觀者」,

係「故事以外ge 人」,

係「只係聽別人見證ge 人」.

但我想對你講:

城裡人都可以遇見耶穌.

你都可以被活水觸動.

你都可以講出你ge 一句——

『祂真係我ge 救主.』



\newpage

\section{從深處的覺醒}
\label{sec:790}
\textbf{郭志榮牧師 - 從深處的覺醒}
\newline
\newline
\hyperref[sec:789]{< < < PREV SERMON < < <}
~
\hyperref[sec:toc]{[返目錄]}
~
\hyperref[sec:791]{> > > NEXT SERMON > > >}


詩篇 130:1-8
\newline
\begin{longtable}{cl}
\hline
\hline
章節 & 經文 (和合本修訂版)\\
\hline
130:1 & \begin{tabularx}{0.7\textwidth}{X} 耶和華啊，我從深處求告你！ \end{tabularx} \\ \\ \relax
130:2 & \begin{tabularx}{0.7\textwidth}{X} 主啊，求你聽我的聲音！求你側耳聽我懇求的聲音！ \end{tabularx} \\ \\ \relax
130:3 & \begin{tabularx}{0.7\textwidth}{X} 耶和華啊，你若究察罪孽，主啊，誰能站得住呢？ \end{tabularx} \\ \\ \relax
130:4 & \begin{tabularx}{0.7\textwidth}{X} 但在你有赦免之恩，要叫人敬畏你。 \end{tabularx} \\ \\ \relax
130:5 & \begin{tabularx}{0.7\textwidth}{X} 我等候耶和華，我的心等候；我也仰望他的話。 \end{tabularx} \\ \\ \relax
130:6 & \begin{tabularx}{0.7\textwidth}{X} 我的心等候主，勝於守夜的等候天亮，勝於守夜的等候天亮。 \end{tabularx} \\ \\ \relax
130:7 & \begin{tabularx}{0.7\textwidth}{X} 以色列啊，你當仰望耶和華，因耶和華有慈愛，有豐盛的救恩。 \end{tabularx} \\ \\ \relax
130:8 & \begin{tabularx}{0.7\textwidth}{X} 他必救贖以色列脫離一切的罪孽。 \end{tabularx} \\ \\
[1ex]
\hline
\hline
\end{longtable}
講題:從深處的覺醒 Awakened Out of the Depths

經課:詩篇130:1-8

講員:郭志榮牧師

場合:香港中文大學崇基學院禮拜堂主日崇拜

日期:2026年3月22日

一．經文:詩一百三十1-8

1     耶和華啊,我從深處求告你!

2     主啊,求你聽我的聲音!求你側耳聽我懇求的聲音!

3     耶和華啊,你若究察罪孽,主啊,誰能站得住呢?

4     但在你有赦免之恩,要叫人敬畏你.

5     我等候耶和華,我的心等候；我也仰望他的話.

6     我的心等候主,勝於守夜的等候天亮,勝於守夜的等候天亮.

7     以色列啊,你當仰望耶和華,因耶和華有慈愛,有豐盛的救恩.

6     他必救贖以色列脫離一切的罪孽.

二．前言

弟兄姊妹,主內平安!

來到大齋期的最後兩週.今日宣講的經文就是來自本週經課的詩篇——《詩篇》第十百三十篇.

這篇詩篇是十五篇〈上行之詩〉的第十一篇,[1] 而現在學者普遍同意,〈上行之詩〉這個「合輯」,是朝聖者在每年「上行」到耶路撒冷過節的旅途上或某些禮儀中頌唱的詩歌.[2] 而有人形容,大齋期就是一個與主耶穌一同走向十字架的屬靈旅程.[3]

詩句中用到「罪孽」,「赦免」,「救恩／救贖」等字眼,在四世紀時,教會早期教父奧古斯丁已把它跟其餘六篇詩篇放在一起,[4] 後來稱之為〈懺悔詩〉.[5] 而大齋期也是一個特別強調認罪懺悔的節期.[6]

剛才讀完這篇詩篇之後,不知道大家會否覺得,這篇詩篇好像沒有甚麼特別之處,只是在重覆《詩篇》中一些經常出現的主題?

然而,宗教改革者馬丁路德就對這篇詩篇推崇備至,稱它為《聖經》中「真正的大師和教師」(a proper master and doctor of Scripture),指它把「福音的基本真理」(the basic truth of the gospel)勾劃出來.[7] 約翰衛斯理(循道衛理會的創始者),在靈性陷入低潮的一個晚上得著「得救確據」,心中感到「異常溫暖」.他在日記中就特別提及,在當天下午的祈禱會中,他正好頌唱了這篇詩篇.[8]

到底這篇詩篇為甚麼能得到如此的讚譽,以及能感動和激勵人的心靈?

三．釋經

按著說話的對象來劃分,《詩篇》一百三十篇可以簡單分為1-4節和5-8節兩個大段落,前半部分是詩人對神說話,後半部分則不然.而上下兩個大段落又各自可分為兩個小段,每小段包括兩節經文.[9] 這四個小段也就形成了這篇詩篇的四個篇章.

以下將會把這四個篇章逐一拆解,以理解這篇詩篇的信息,並從而看到詩人屬靈覺醒的進程.

A. 從身處「深處」到迫切「求告神」的覺醒

第一個篇章所表達的內容是詩人向神的求告.詩人這樣說:「耶和華啊,我從深處求告你!」( 1節)詩人在這裡特別強調,他的求告是「從深處」向神發出,表明詩人意識到自己正陷於一個「深處」的處境之中.

「深處」即「海之深處」,[10] 比喻一個人正處於一個極之嚴峻的惡劣處境裡,好像遇溺一樣,而且並非在水面上掙扎,而是已下沉到深海之中.[11]「海」在《舊約聖經》中常代表那股為生命帶來「破壞」,「毀滅」和「死亡」的「混沌的力量」.[12] 而沉沒在「海之深處」,意味著人被大水所包圍和覆蓋,因此而與神隔絕.[13] 換言之,詩人意識到自己的活力,動力甚至生命氣息,都正在逐漸流逝和消失.詩人更發覺單靠自己的力量掙扎完全於事無補,只會徒然增加他的無力感.

面對「死亡」的危機,人常會有「否認」,「憤怒」,「討價還價」和「沮喪」的反應.[14] 但詩人沒有因此而被癱瘓.他沒有陷於對一切事情皆加以否定的狀態,他拒絕面對現實；他沒有被憤怒的情緒所困,終日怨天尤人,質問不幸的事情為何臨到自己身上；他沒有停留在討價還價的空想之中,希望惡事禍事會突然消失,甚至從未發生；他沒有陷溺在沮喪的心境裡,被焦慮和自憐自憫困住.

詩人以對神懇切的求告作為應對:他求神聽他的聲音!側耳留心他懇求的聲音!(參2節)「聽」這個字包含了「回應」的意思.[15] 而「懇求的聲音」中,「懇求」特別指祈求神的「憐憫與施恩」.[16] 所以,詩人不僅希望神成為他的聆聽者,[17] 他更是需要神對他的呼求作出發自憐憫之情的回應行動.

詩人如此應對他陷於「深處」的困局,表示他認識到自己無力自救的事實,而他更知道神的回應和幫助才是他唯一的出路.

B. 從「罪」到對神「赦免之恩」的覺醒

詩人因為覺醒到自己在「深處」中無力自救而向神發出求告後,他突然把話題轉換到神的「赦免之恩」之上.詩人對著神「自問自答」:「耶和華啊!你若究察罪孽,主啊,誰能站得住呢?但在你有赦免之恩,要叫人敬畏你.」(3-4節)

「究察」不僅指「細心留意」,[18] 更有「記錄下來」(keep record)的含義.[19] 詩人用「反問」的方式表達,神當然看得到人所作一切的罪孽,也就是一切偏離神,[20] 亦即違反神話語教導的行為,言語甚至是最隱密的思想,神不會忽視這個世界的罪惡,[21] 但神不把人犯罪的記錄鉅細無遺地收集起來,作為審判人的根據,否則,根本沒有人能在神的審判台前「站得住」而存活下來.[22]

這是詩人「信心的宣告」(a statement of confidence),[23] 但大家一定會問,這位正在「深處中求告神」的詩人心裡到底想起了甚麼事情,因而產生這個思想上的「跳躍」呢?

答案就是:詩人意識到自己陷於「深處」之中並非無因,他自己其實也要負上相當的責任.[24] 這是詩人極之深刻的反省和生命的洞見.「深處」的覺醒使他認識到自己只能尋求神的幫助,而這樣的「覺醒」和「認識」使他「甦醒」過來,以一個新的角度回顧自己的人生,重新審視自己生命的真相.

「罪孽」有「扭曲」,「偏離」的含義.[25].他知道問題的根源是自己離開了神,也就是他對神的不信——他不信把自己的生命主權交在神手中會比由自己掌握更穩妥；不信神的旨意(即《聖經》的教導)比自己所謂「理性」的判斷更正確；不信神對他的愛會如此深和徹底,比他對自己的愛更多更深.

「深處」之所以為「深處」,重點不在於外在的環境如何,而是有沒有神的同在.受造的人離開了賜生命的創造之主,自然會陷於「深處」而步向滅亡!這是必然之事,問題是人有沒有這個「覺醒」.

詩人提到神有「赦免之恩」,原文是「神有赦免」(forgiveness).莎士比亞曾經說過一句名言:「沒有甚麼比憐憫更能助長罪惡的膽量.」(Nothing emboldens sin so much as mercy.)[26] 有猶太拉比問過這樣的問題:假若神總是赦免人的罪,人會不會因此而變得更不用畏懼犯罪而得罪神呢?[27] 答案當然是不會!

有拉比作了一個比喻來解釋「赦免」的作用:假若一個債仔發覺他所欠的債務,是他怎樣努力也不可能如期償還.如果他的債主堅持這個債仔必須按時清還所有欠債,否則便會即時提出申請破產的法律行動,這個債仔很可能會選擇放棄繼續還債,任由債主申請他破產,因為反正結果都是一樣.但假若債主延長債仔的還款期限,甚至減免債仔的部分債務,債仔反而覺得仍有希望,於是便繼續還款.[28]

這位拉比的講法當然有他的道理,但卻仍未點出最重要的重點.最重要的是:神的「赦免」除了指抹去一個人罪孽的記錄外,[29] 也包括「除罪」(removing sins)的意思,[30] 換言之也是一個改變和轉化罪人生命的過程.而詩人正好強調,神的「赦免」是使人會因感恩而「敬畏」神.所以,「赦免」絕對不是縱容人放肆繼續犯罪,而是使人感恩以致生命轉變,並活出「敬畏」神的人生,這才是神「赦免」的目的.[31]

C. 三重「等候」的覺醒

接下來,詩人便作出堅定的信仰宣言,而這也是他的立志:「我等候耶和華,我的心等候；我也仰望他的話.我的心等候主,勝於守夜的等候天亮,勝於守夜的等候天亮.」(5-6節)

在這短短的兩節經文中,詩人反覆使用了三個不同,卻同樣可譯作「等候」的希伯來文詞語.有些學者認為,詩人的說法表明,他雖然仍身處於「深處」那看似沒有盡頭的黑暗之中,但卻懷著信心與盼望,耐心等待神的介入與拯救所帶來的光明.[32] 但這樣的理解過於消極,未能把這三重「等候」的深層意義表達得清楚.

1. 第一重「等候」:「我等候耶和華,我的心等候」

第一個「等候」所用的那個字(qwh),字根跟繩索有關.[33] 大家一定會好奇:繩索和「等候」有甚麼關係呢?大家可以從兩個方面去理解.首先,一條繩索是由多條較幼的繩子「交織」而成. 引伸出來,「等候」就是把「等候者」和「等候的對象」纏繞在一起的模樣.[34] 另一個理解的方式就更簡單和直接,「等候」的功能就好像一條拉緊了的繩索,[35] 一端連在「等候的對象」身上,另一端就在「等候者」手中緊緊拉著,讓「等候者」能夠保持甚至拉近彼此之間的距離.[36]

無論是哪一種解讀,第一重「等候」所強調的,是「等候者」與「等候的對象」的緊密關係,而這是生命結連的關係!因此,「等候者」絕對不是因為無可奈何,只能聽天由命,被動地「等運到」,而是主動地,竭力地與「等候的對象」建立愈來愈緊密的生命結連.他會積極地作好充分的準備,迎接「等候的對象」的臨到.

詩人說:「我等候耶和華,我的心等候.」前者表明詩人「等候的對象」是誰,就是耶和華這位救贖以色列民的,與他們有生命之約的獨一真神.

「我的心」中的「心」原文是「生命」的意思,[37] 所以,「我的心」指詩人的整個生命,也指他整個人.詩人想表達的是:他整個人——包括他的身體,他的思想,他的情感,他的意志——都要結連於神,都要交在神手中,成為神合用的器皿.而從此以後,詩人生命的主權屬於神,他人生的目標和生命方式的選取,皆離不開神的心意,亦即神話語的教導.

2. 第二重「等候」:「我也仰望他的話」

詩人接著說:「我也仰望他的話.」「仰望」一詞也包含了「等候」的意思,[38] NASB和JPS Tanakh英文譯本把這個字分別譯作「wait for」和「await」,就是這個原因.

第二重「等候」有「信心」,「信靠」和「盼望」的含義.[39] 而這個字的字根與「陷於泥沼」或「陷於困境」有關,[40] 所以第二重的「等候」常用來指「在困難處境中等候」,就正如詩人正處於「深處」之中那樣.

因此,《聖經》幾個不同的中文譯本把這個字譯作「仰望」,[41] 其實很有意思,表達出「等候者」是抬頭注目於天上那位滿有大能,值得信靠的信實之主,而不是單單低頭望著地上那些面前的難處.

詩人在困難中懷著信心和盼望所「等候」的,不是別的事物,而是「神的話」.詩人向上「仰望」,期盼「神的話」臨到.「神的話」包括了神的安慰,讓詩人的心得到堅固和激勵.「神的話」也包含了神的指引,讓詩人知道在當前的困境中要怎樣做.「神的話」更充滿了神的應許,詩人「等待」信實的神以祂的大能實現祂的應許.

這這是第二重「等候」所表達的重點.

3. 第三重「等候」:「我的心等候主,勝於守夜的等候天亮,勝於守夜的等候天亮.」

第6節經文:「我的心等候主,勝於守夜的等候天亮,勝於守夜的等候天亮.」原文應直譯為:「我的生命向主,勝於那等候天亮的,勝於那等候天亮的.」[42]

第三重的「等候」把詩人和那些「等候天亮」的人作對比.這些「等候天亮」的人可以指那些在晚上看守城門的人,[43] 但更可能是指在聖殿供職,負責守夜的利未祭司,他們等待晨光出現出現的那刻,便獻上當天的「每日燔祭」(daily sacrifice).[44] 這兩種「等候」的共通點,就是不能鬆懈,要專注,謹慎和警醒.

「勝於」原文是一個介詞(類似英文中的preposition),既可解作「勝於」,但也用來表達「在他們當中」的意思.[45] JPS Tanakh英文譯本跟大部分的中英文《經譯》本一樣,採用了前者的意思,表示詩人比那些守城或在聖殿守夜的人更加專注,謹慎和警醒.而Judaica Press的翻譯側使用了後者的解釋,[46] 強調詩人是與其他人一起「等候」那「救贖的黎明」到來,[47] 顯明詩人是滿有群體意識的信仰者.等一會大家將會看到,後者的解釋有優勝之處,就是與第四段篇章(即7-8節)銜接得更好.

4. 小結

總結而言,這三重「等候」都絕對不是被動和消極地「捱日子」,不會認為當下是沒有意義,是毫無作為的.

這樣的三重「等候」,正呼應著前一個篇章結尾時所強調的重點,就是生命要更新轉化,活出「敬畏神」的人生.

D. 「群體意識」的覺醒

這篇詩篇沒有在詩人「等候」的宣告中完結,詩人最後向「以色列」發出呼籲:「以色列啊,你當仰望耶和華,因耶和華有慈愛,有豐盛的救恩.他必救贖以色列脫離一切的罪孽.」(7-8節)

這表明詩人有更清晰的「覺醒」,知道他所持守的是一個「群體的信仰」,而不是單單追求自己個人的救贖.詩人知道他不是孤單地在「深處」之中受困和受苦,所以他要與他的同伴一起求告神,一起倚靠神的赦免之恩以致生命轉變,活出敬畏神的人生.他與群體中的同伴一起在「等候」,同把生命結連於神,一起在「等候」中專注仰望神的話語,並一起過謹慎和警醒「等候」的生活.

當詩人重提神「有慈愛,有豐盛的救恩」,以及神「必救贖以色列脫離一切的罪孽」的時候,在熟悉《舊約聖經》的詩人心中,好可能正是聯想到《出埃及記》中的那段金句:「耶和華,耶和華,有憐憫,有恩惠的神,不輕易發怒,且有豐盛的慈愛和信實,為千代的人存留慈愛,赦免罪孽,過犯和罪惡,萬不以有罪的為無罪,必懲罰人的罪,自父及子,直到三,四代.」(出卅四6-7,和修版)[48]

這是神在「金牛犢事件」後,跟以色列民重新立約前向以色列民自我啟示的說話.詩人的呼籲提醒以色列的同胞,他們是與神有生命之約的屬神群體.

在這個群體中,大家更易得到勉勵和餵養,更能實踐神要他們承擔的使命,就是要按著神話語而生活,彼此激勵和扶持,抗衡罪惡,向群體以外的人彰顯神的榮耀,讓人認識到神在人生命中的份量和重要性.

四．引伸應用

A. 詩人「屬靈覺醒」的進程

這篇詩篇的四個篇章,一環緊扣著一環,正好示範了一個「屬靈覺醒」進程的四步曲.

首先是對自身正陷於「深處」的覺醒,也就是意識到自己的生命正在逐步流逝.這是「屬靈覺醒」的第一步,但人有這樣的覺醒並非必然,陷在罪中,步向滅亡卻不自知的人,實在多的是!他們有時甚至為了合理化自己的惡行,好讓自己能繼續享受罪中之樂,於是把判斷「是非對錯善惡好壞」的真理和價值觀全然扭曲.

即使是信徒,也要時常省察自己生命的景況如何,是否在某些方面也有逐漸走近「深處」的情況,甚至是有墮進「深處」的危機!

第二步是「覺醒」到自己之所以陷於「深處」的終極原因,是自己因為不信的惡心而與創造人並賜下生命氣息的萬有主宰隔絕了.在「深處」掙扎的人,若看不清這才是淪落到如斯田地的終極原因,就有可能會被困在憤怒和自憐的情緒之中,只會怨天尤人或去搞破壞.

作為信徒也要省察,神在自己的生命中是遠是近?例如,相比起看手機,自己有多少時間會安靜在神面前,親近神,進入深度的禱告,聆聽聖靈藉著神話語的提醒?

第三步是要有那三重「等候」的覺醒,知道要以積極的「等候」,把自己的整個生命全方位地與神建立起生命的結連,要注目於神話語的指引和教導,並要時刻謹慎和警醒,以致對神的作為和提示有敏銳的觸覺.這個漫長的「等候」過程,實際上也是信徒生命成長的進程.

最後是「群體意識」的覺醒,而這也是信徒的「屬靈意識」由「個人」走向「群體」的提升進程.一個永遠停留在「個人意識」,不願為群體付出而放下自我的信徒,就難以明白何謂基督大愛的真諦.

在今年的大齋期中,讓我們也從這四個方面省察我們屬靈生命的實況.

B. 怎樣的「深處」?——傳福音的起始點

以上的框架除了勾劃出「屬靈覺醒」進程的四步曲外,它也給我們一個重要的提示,就是現在向人傳福音時要從哪裡入手,而答案就是要問:現在的福音對象到底處於一個怎樣的「深處」?

是有家人甚至有很多朋友,但卻無人能傾訴,內心感到異常孤單?是壓力無處不在,極其沉重,令人喘不過氣來?是工作忙碌,營營役役,停不下來,卻找不到真正的人生意義,不知這一切為了甚麼?是活在焦慮之中,缺乏安全感,卻因找錯了生命的保障而心靈無法安頓下來?是被社會扭曲的價值所牽引而陷在罪的折騰之中?還是別的「處境」?

明白福音對象的處境,大家便可以循那個方向作為分享福音的起始點,觸動福音對象的心靈,認識自己生命的真相,引導他們尋求神,讓基督的福音改變他們的生命.

書目

Allen, Leslie C. Psalms 101–150 (Revised), vol. 21, Word Biblical Commentary. Dallas: Word, Incorporated, 2002.

Botterweck, G. Johannes and Helmer Ringgren. Theological Dictionary of the Old Testament, ed., trans. David E. Green. Grand Rapids: William B. Eerdmans Publishing Company, 1990.

Dahood, Mitchell S.J. Psalms III: 101-150: Introduction, Translation, and Notes with an Appendix: The Grammar of the Psalter, vol. 17A, Anchor Yale Bible. New Haven: Yale University Press, 2008.

Hakham, Amos. The Bible: Psalms with the Jerusalem Commentary, vol 3. Jerusalem, Mosad Harav Kook, 2003.

Harris, R. Laird, Gleason L. Archer Jr., and Bruce K. Waltke. Theological Wordbook of the Old Testament. Chicago: Moody Press, 1999.

Fruchtenbaum, Arnold G. The Messianic Bible Study Collection, vol. 152. Tustin: Ariel Ministries, 1983.

Gesenius, Wilhelm and Samuel Prideaux Tregelles. Gesenius’ Hebrew and Chaldee Lexicon to the Old Testament Scriptures. Bellingham: Logos Bible Software, 2003.

Goldingay, John. Baker Commentary on the Old Testament: Psalms 90–150, ed. Tremper Longman III, vol. 3. Grand Rapids: Baker Academic, 2006.

Grogan, Geoffrey W. Psalms, The Two Horizons Old Testament Commentary. Grand Rapids: William B. Eerdmans, 2008.

Jenni, Ernst and Claus Westermann. Theological Lexicon of the Old Testament. Peabody: Hendrickson Publishers, 1997.

Kirkpatrick, A. F. The Book of Psalms, The Cambridge Bible for Schools and Colleges. Cambridge: Cambridge University Press, 1906.

Marcus, Yosef B. Tehillim: Book of Psalms with Commentary from the Talmud, Midrash, Kabbalah, Classic Commentators, and the Chasidic Masters. New York: Kehot Publication Society, 2016.

Mays, James Luther. Psalms, Interpretation, a Bible Commentary for Teaching and Preaching. Louisville: John Knox Press, 1994.

McCann, J. Clinton Jr. “The Book of Psalms,” in New Interpreter’s Bible, ed. Leander E. Keck, vol. 4. Nashville: Abingdon Press, 1994–2004.

Melina, Livio. The Epiphany of Love: Toward a Theological Understanding of Christian Action. Grand Rapids: William B. Eerdmans, 2010.

Kidner, Derek. Psalms 73–150: An Introduction and Commentary, vol. 16, Tyndale Old Testament Commentaries. Downers Grove: InterVarsity Press, 1975.

Ross, Allen P. A. Commentary on the Psalms (90–150): Commentary, vol. 3, Kregel Exegetical Library. Grand Rapids: Kregel Academic, 2016.

Waltke, Bruce K. James M. Houston, and Erika Moore. The Psalms as Christian Lament: A Historical Commentary. Grand Rapids: William B. Eerdmans, 2014.

[1] James Luther Mays, Psalms, Interpretation, a Bible Commentary for Teaching and Preaching (Louisville: John Knox Press, 1994), 405.

[2] J. Clinton Mccann Jr., “The Book of Psalms,” in New Interpreter’s Bible, ed. Leander E. Keck, vol. 4 (Nashville: Abingdon Press, 1994–2004), 1176.

[3] Sherwood E. Wirt, “Let’s Lengthen Lent,” Christianity Today (Washington, D.C.: Christianity Today, 1974), 546.

[4] Bruce K. Waltke, James M. Houston, and Erika Moore, The Psalms as Christian Lament: A Historical Commentary (Grand Rapids: William B. Eerdmans, 2014), 14.

[5] Mitchell Dahood S.J., Psalms III: 101-150: Introduction, Translation, and Notes with an Appendix: The Grammar of the Psalter, vol. 17A, Anchor Yale Bible (New Haven: Yale University Press, 2008), 234.

[6] Andrea Wigodsky, “Pastoral Perspective on Psalm 32,” in Feasting on the Word: Preaching the Revised Common Lectionary: Year A, ed. David L. Bartlett and Barbara Brown Taylor, vol. 2 (Louisville: Westminster John Knox Press, 2010), 32.

[7] James Luther Mays, Psalms, Interpretation, a Bible Commentary for Teaching and Preaching (Louisville,: John Knox Press, 1994), 405.

[8] \url{https://ap.org.au/2022/05/24/i-remember-24-may-1738-john-wesleys-happy-day}\%EF\%BF\%BC/

[9] Amos Hakham, The Bible: Psalms with the Jerusalem Commentary, vol 3 (Jerusalem, Mosad Harav Kook, 2003), 337.

[10] Ronald B. Allen, “1644 עָמֹק,” in Theological Wordbook of the Old Testament, ed. R. Laird Harris, Gleason L. Archer Jr., and Bruce K. Waltke (Chicago: Moody Press, 1999), 679.

[11] Amos Hakham, The Bible: Psalms with the Jerusalem Commentary, vol 3 (Jerusalem, Mosad Harav Kook, 2003), 337.

[12] J. Clinton Mccann Jr., “The Book of Psalms,” in New Interpreter’s Bible, ed. Leander E. Keck, vol. 4 (Nashville: Abingdon Press, 1994–2004), 1205.

[13] Leslie C. Allen, Psalms 101–150 (Revised), vol. 21, Word Biblical Commentary (Dallas: Word, Incorporated, 2002), 255.

[14] \url{https://www.change-management-coach.com/kubler-ross.html}

[15] Allen P. Ross, A Commentary on the Psalms (90–150): Commentary, vol. 3, Kregel Exegetical Library (Grand Rapids: Kregel Academic, 2016), 712.

[16] Edwin Yamauchi, “694 חָנַן,” in Theological Wordbook of the Old Testament, ed. R. Laird Harris, Gleason L. Archer Jr., and Bruce K. Waltke (Chicago: Moody Press, 1999), 302.

[17] Allen P. Ross, A Commentary on the Psalms (90–150): Commentary, vol. 3, Kregel Exegetical Library (Grand Rapids: Kregel Academic, 2016), 712.

[18] R. Laird Harris, Gleason L. Archer Jr., and Bruce K. Waltke, eds., Theological Wordbook of the Old Testament (Chicago: Moody Press, 1999), 939.

[19] Mitchell Dahood S.J., Psalms III: 101-150: Introduction, Translation, and Notes with an Appendix: The Grammar of the Psalter, vol. 17A, Anchor Yale Bible (New Haven: Yale University Press, 2008), 235.

[20] Carl Schultz, “1577 עָוָה,” in Theological Wordbook of the Old Testament, ed. R. Laird Harris, Gleason L. Archer Jr., and Bruce K. Waltke (Chicago: Moody Press, 1999), 650–651.

[21] Leslie C. Allen, Psalms 101–150 (Revised), vol. 21, Word Biblical Commentary (Dallas: Word, Incorporated, 2002), 256.

[22] Amos Hakham, The Bible: Psalms with the Jerusalem Commentary, vol 3 (Jerusalem, Mosad Harav Kook, 2003), 338.

[23] John Goldingay, Baker Commentary on the Old Testament: Psalms 90–150, ed. Tremper Longman III, vol. 3 (Grand Rapids: Baker Academic, 2006), 526.

[24] 同上.

[25] Carl Schultz, “1577 עָוָה,” in Theological Wordbook of the Old Testament, ed. R. Laird Harris, Gleason L. Archer Jr., and Bruce K. Waltke (Chicago: Moody Press, 1999), 650–651.

[26] Leslie C. Allen, Psalms 101–150 (Revised), vol. 21, Word Biblical Commentary (Dallas: Word, Incorporated, 2002), 256.

[27] Yosef B. Marcus, Tehillim: Book of Psalms with Commentary from the Talmud, Midrash, Kabbalah, Classic Commentators, and the Chasidic Masters (New York: Kehot Publication Society, 2016), 272a.

[28] 同上.

[29] Geoffrey W. Grogan, Psalms, The Two Horizons Old Testament Commentary (Grand Rapids; Cambridge, U.K.: William B. Eerdmans Publishing Company, 2008), 207.

[30] Allen P. Ross, A Commentary on the Psalms (90–150): Commentary, vol. 3, Kregel Exegetical Library (Grand Rapids: Kregel Academic, 2016), 714.

[31] 同上.

[32] Derek Kidner, Psalms 73–150: An Introduction and Commentary, vol. 16, Tyndale Old Testament Commentaries (Downers Grove: InterVarsity Press, 1975), 483.

[33] Ernst Jenni and Claus Westermann, Theological Lexicon of the Old Testament (Peabody: Hendrickson Publishers, 1997), 1126.

[34] Wilhelm Gesenius and Samuel Prideaux Tregelles, Gesenius’ Hebrew and Chaldee Lexicon to the Old Testament Scriptures (Bellingham: Logos Bible Software, 2003), 726–727.

[35] G. Waschke, “קוה,” in Theological Dictionary of the Old Testament, ed. G. Johannes Botterweck, Helmer Ringgren, and Heinz-Josef Fabry, trans. Douglas W. Stott (Grand Rapids: William B. Eerdmans Publishing Company, 2003), 565.

[36] Livio Melina, The Epiphany of Love: Toward a Theological Understanding of Christian Action (Grand Rapids: William B. Eerdmans, 2010), 20.

[37] Bruce K. Waltke, “1395 נָפַשׁ,” in Theological Wordbook of the Old Testament, ed. R. Laird Harris, Gleason L. Archer Jr., and Bruce K. Waltke (Chicago: Moody Press, 1999), 587–591.

[38] Paul R. Gilchrist, “859 יָחַל,” in Theological Wordbook of the Old Testament, ed. R. Laird Harris, Gleason L. Archer Jr., and Bruce K. Waltke (Chicago: Moody Press, 1999), 373.

[39] Paul R. Gilchrist, “859 יָחַל,” in Theological Wordbook of the Old Testament, ed. R. Laird Harris, Gleason L. Archer Jr., and Bruce K. Waltke (Chicago: Moody Press, 1999), 373.

[40] Ch. Barth, “יָחַל,” in Theological Dictionary of the Old Testament, ed. G. Johannes Botterweck and Helmer Ringgren, trans. David E. Green (Grand Rapids, MI; Cambridge, U.K.: William B. Eerdmans Publishing Company, 1990), 49.

[41] 包括《和合本》,《和修版》,《新譯本》

[42] 參考Young Literal Translation的翻譯.

[43] Arnold G. Fruchtenbaum, The Messianic Bible Study Collection, vol. 152 (Tustin: Ariel Ministries, 1983), 22.

[44] A. F. Kirkpatrick, The Book of Psalms, The Cambridge Bible for Schools and Colleges (Cambridge: Cambridge University Press, 1906), 760.

[45] R. Laird Harris, “1212 מִן,” in Theological Wordbook of the Old Testament, ed. R. Laird Harris, Gleason L. Archer Jr., and Bruce K. Waltke (Chicago: Moody Press, 1999), 512.

[46] \url{https://www.chabad.org/library/bible_cdo/aid/16351/showrashi/true}

[47] Yosef B. Marcus, Tehillim: Book of Psalms with Commentary from the Talmud, Midrash, Kabbalah, Classic Commentators, and the Chasidic Masters (New York: Kehot Publication Society, 2016), 272b.

[48] J. Clinton Mccann Jr., “The Book of Psalms,” in New Interpreter’s Bible, ed. Leander E. Keck, vol. 4 (Nashville: Abingdon Press, 1994–2004), 1205.



\newpage

\section{路不能亂走車不可亂坐}
\label{sec:791}
\textbf{林豪恩先生 - 路不能亂走,車不可亂坐}
\newline
\newline
\hyperref[sec:790]{< < < PREV SERMON < < <}
~
\hyperref[sec:toc]{[返目錄]}
~
\hyperref[sec:792]{> > > NEXT SERMON > > >}


馬太福音 21:1-11
\newline
\begin{longtable}{cl}
\hline
\hline
章節 & 經文 (和合本修訂版)\\
\hline
21:1 & \begin{tabularx}{0.7\textwidth}{X} 耶穌和門徒快到耶路撒冷，進了橄欖山的伯法其時，打發兩個門徒， \end{tabularx} \\ \\ \relax
21:2 & \begin{tabularx}{0.7\textwidth}{X} 對他們說：「你們往對面村子裡去，會立刻看見一匹驢拴在那裡，還有驢駒同在一處，解開牠們，牽到我這裡來。 \end{tabularx} \\ \\ \relax
21:3 & \begin{tabularx}{0.7\textwidth}{X} 若有人對你們說甚麼，你們就說：『主要用牠們。』那人會立刻讓你們牽來。」 \end{tabularx} \\ \\ \relax
21:4 & \begin{tabularx}{0.7\textwidth}{X} 這事發生是要應驗先知所說的話： \end{tabularx} \\ \\ \relax
21:5 & \begin{tabularx}{0.7\textwidth}{X} 「要對錫安的兒女說：看哪，你的王來到你這裡，謙和地騎著驢，騎著小驢—驢的駒子。」 \end{tabularx} \\ \\ \relax
21:6 & \begin{tabularx}{0.7\textwidth}{X} 門徒就照耶穌所吩咐的去做， \end{tabularx} \\ \\ \relax
21:7 & \begin{tabularx}{0.7\textwidth}{X} 牽了驢和驢駒來，把他們的衣服搭在上面，耶穌就騎上。 \end{tabularx} \\ \\ \relax
21:8 & \begin{tabularx}{0.7\textwidth}{X} 許許多多的人把自己的衣服鋪在路上，還有人砍下樹枝來鋪在路上。 \end{tabularx} \\ \\ \relax
21:9 & \begin{tabularx}{0.7\textwidth}{X} 前呼後擁的人群喊著說：「和散那歸於大衛之子！奉主名來的是應當稱頌的！至高無上的，和散那！」 \end{tabularx} \\ \\ \relax
21:10 & \begin{tabularx}{0.7\textwidth}{X} 耶穌進了耶路撒冷，全城都驚動了，說：「這是誰？」 \end{tabularx} \\ \\ \relax
21:11 & \begin{tabularx}{0.7\textwidth}{X} 眾人說：「這是從加利利的拿撒勒來的先知耶穌。」 \end{tabularx} \\ \\
[1ex]
\hline
\hline
\end{longtable}
講題:路不能亂走,車不可亂坐 Cautious Steps, Thoughtful Rides

經課:馬太福音21:1-11

講員:林豪恩先生

場合:香港中文大學崇基學院禮拜堂主日崇拜

日期:2026年3月29日

引言

你乘搭飛機的時候,有沒有試過驚喜地被升級上頭等艙呢?有網友分享說:「我拿著經濟艙機票要登機的時候,空姐把我的票換了一張,說:「女士,你的座位被升級了!」我一時反應不過來拿著機票向前走,看著機票寫著座位01A,心想,這可能是經濟艙比較寬大的位置.沒想到我是被升級成商務艙了!於是我蹦蹦跳跳地走進寬敞的商務艙走道,坐下就馬上開始玩各種座椅的按鈕!」這位網友總結說:「好幸運喔!人生的第一次商務艙成就達成!超爽!」不知道你有沒有這樣的經驗呢?或者有沒有期望這樣的待遇呢?如果突如其來被升級商務艙,的確是意外驚喜.不過,常言道,水能載舟,亦能覆舟,驚喜之後有可能變成驚嚇.

(香港01 2026-1-22)香港駐天津聯絡處主任在社交平台分享,獲噴射飛航高級管理人員安排港澳航程坐私人艙室由香港去澳門.政制及內地事務局獲悉事件後表示,已決定即時終止該職員的合約,經進一步調查,認為該在早前乘船前往澳門時座位獲「升級」的事件,有違反《防止賄賂條例》下有關公務員未經許可接受利益的規定之嫌,已把個案轉介執法部門跟進.

自古至今,代步工具都可能帶有某種含意.例如香港開埠初期最常見的交通工具有轎子,人力車,馬車等.然而,這些代步工具都是殖民地官員或富有的洋行大班用的.一般平民百姓的交通工具是甚麼呢?就是自己的雙腳.

更久遠的耶穌年代呢?當時的一般人從一個地方到另外一個地方通常都是步行,聖經也經常提及耶穌和門徒走路.然而,今天我們讀到的福音經課,卻看到耶穌煞有介事地安排代步工具,到底所為何事呢?

一．為甚麼耶穌要騎驢進城?

馬太福音

21:1耶穌和門徒將近耶路撒冷,到了伯法其,在橄欖山那裡.

21:2耶穌就打發兩個門徒,對他們說:「你們往對面村子裡去,必看見一匹驢拴在那裡,還有驢駒同在一處；你們解開,牽到我這裡來.

21:3若有人對你們說甚麼,你們就說:『主要用牠.』那人必立時讓你們牽來.

21:6門徒就照耶穌所吩咐的去行,

21:7牽了驢和驢駒來,把自己的衣服搭在上面,耶穌就騎上.

當時除了耶穌和門徒之外,其實往耶路撒冷路上的行人絡繹不絕,因為猶太人的重要節期快到,人們都去耶路撒冷過節.一般平民百姓大多數都如耶穌和門徒一樣步行前往耶路撒冷.現在,他們已經來到位於橄欖山的伯法其,這伯法其在耶路撒冷的東面,距離耶路撒冷只有1.6公里,只要多行一會就到達了.耶穌為甚麼要在這個時候找一隻驢去騎呢?是耶穌體力不支而需要坐騎嗎?還是耶穌趕時間而需要坐騎呢?看來兩者都不是.那麼,特意安排坐驢進城是甚麼事情呢?這可能是一場進城儀式.

進城儀式是甚麼呢?在古代,進城儀式是一種具有強烈象徵意義的隆重典禮.古羅馬進城儀式中最著名的是「凱旋式」,授予戰勝將領或皇帝的最高軍事與宗教榮譽.獲勝者身穿神聖服飾,乘坐金色戰車,率軍隊,戰利品及俘虜穿過城門至神廟,藉此展示權力,慶祝勝利.

與古羅馬「凱旋式」的進城儀式相比,耶穌坐的不是金色戰車,而是一匹民間運輸用的驢；耶穌率領的不是軍隊,而是門徒.看在路人眼中,這算是甚麼樣的進城儀式呢?有趣的是有些路人樂意配合,演好這一場戲.

21:8眾人多半把衣服鋪在路上；還有人砍下樹枝來鋪在路上.

21:9前行後隨的眾人喊著說:和散那歸於大衛的子孫!奉主名來的是應當稱頌的!高高在上和散那!

馬太福音記述這事件要表達甚麼意思呢?相信大家都留意得到馬太福音從預告耶穌出生就經常用舊約聖經來解釋事件的意義,這次也如是.

21:4這事成就是要應驗先知的話,說:

21:5要對錫安的居民說:看哪,你的王來到你這裡,是溫柔的,又騎著驢,就是騎著驢駒子.」

講這幾句話的先知是誰呢?他就是撒迦利亞,這段引文出自撒迦利亞書:

9:9錫安的民哪,應當大大喜樂；耶路撒冷的民哪,應當歡呼.看哪,你的王來到你這裡!他是公義的,並且施行拯救,謙謙和和地騎著驢,就是騎著驢的駒子.

馬太福音指出,雖然耶穌的進城儀式和古羅馬「凱旋式」的進城儀式相距甚遠,但卻不應該被視為「兒戲」,因為古時有先知預告,錫安的王是騎驢進城,而不是驢戰馬進城的.然而,眼前看見的狀況可能令當時認真思考的人非常困惑,因為眼前看見的狀況和現實的需要差距實在太大了.現實的需要是甚麼呢?在羅馬帝國的鐵蹄之下,在官府和地方勢力的高壓管治之下,騎驢駒進城的王有甚麼用呢?騎驢的人,無論多麼謙和,多麼公義,只會是「遇著兵」的「秀才」,能有甚麼作為呢?根據「華人基督徒聖地和平網絡」2026年3月23日報導,以色列總理內塔尼亞胡在記者會回應問題時,暗諷有些人「天真」,「自欺欺人」,以為光靠道德和正義就夠打敗敵人,他援引歷史學家杜蘭特(Will Durant)在《歷史的教訓》(The Lessons of History)所寫,指「歷史證明,遺憾且不幸的是,耶穌基督並未比成吉思汗強.因為若你夠強大,夠冷酷,夠強勢,邪惡終將戰勝善良,侵略終將戰勝克制」.或許2000年前那個地方的人都是懷著類似想法,當他們知道耶穌只是一位從加利利拿撒勒來手無寸鐵的先知,或許都希望他做完這場戲就乖乖地做回來過節的朝聖者,免得節外生枝.畢竟這是猶太人的大節日,人們是來慶祝的,不是來受難的；是來向耶路撒冷讚美的,不是來哀悼的.

二．謙和的王何以受難?

現實卻是,就算耶穌是手無寸鐵的先知,也遭到嚴刑峻罰,被有殺錯無放過地對待.為何這位驅驢進城的謙和之士也在被殺之列呢?他沒有戰馬,也沒有軍隊,能夠如何威脅當權者呢?會不會只是因為他在暴力當道的世界批判暴力,呼籲謙和才是拯救之道呢?會不會只是因為他在聖殿被異化的宗教狀況之下,拆除捆綁心靈的律法主義,帶領人歸向上主呢?會不會只因為他應該的拯救及自由促成了宗教和帝國的合作而被判為罪人呢?

這個人做錯了甚麼以至被本族人出賣,被羅馬軍兵釘上十字架呢?如果沒有做錯甚麼,好人為甚麼遭遇苦難呢?好人遭遇苦難確是令人感到傷感和不解的千古疑問.耶穌就是其中一個義人受苦的例子.今天是大齋期最後一個主日,也是聖週的開始.聖週,又稱受難週,在基督教傳統中,是復活節之前的一週,用來紀念耶穌受難.原來教會從來沒有輕視好人遭遇苦難這反常狀況,也沒有忘記那些為義受苦的無辜者.

今天經課中的以賽亞書指出,受苦不一定因為有罪,受苦可能只因為他忠心地活出作為聽從上主的僕人:

50:5主耶和華開通我的耳朵；我並沒有違背,也沒有退後.

50:6人打我的背,我任他打；人拔我腮頰的鬍鬚,我由他拔；人辱我,吐我,我並不掩面.

50:7主耶和華必幫助我,所以我不抱愧.我硬著臉面好像堅石；我也知道我必不致蒙羞.

50:8稱我為義的與我相近；誰與我爭論,可以與我一同站立；誰與我作對,可以就近我來.

50:9主耶和華要幫助我；誰能定我有罪呢?

三．真正有實力的才能謙和

如何面對好人都要受苦這種逆境呢?怎樣的人能夠面對辱罵而不抱愧呢?怎樣的人能夠面對攻擊而仍然站立得穩呢?怎樣的人才能夠面對邪惡而自己不變為邪惡呢?這樣的人雖然看似軟弱,但其實不可能是軟弱的人.其實,堅強的人才能不被逆境摧毀,有實力的人才能夠謙和,「行得正企得正」的人才能夠站立得穩.今天的書信經課腓立比書,保羅指出,手無寸鐵的耶穌其實是如此堅實的人:

2:6他本有上帝的形像,不以自己與上帝同等為強奪的；

2:7反倒虛己,取了奴僕的形像,成為人的樣式；

2:8既有人的樣子,就自己卑微,存心順服,以至於死,且死在十字架上.

2:9所以,上帝將他升為至高,又賜給他那超乎萬名之上的名,

2:10叫一切在天上的,地上的,和地底下的,因耶穌的名無不屈膝,

2:11無不口稱「耶穌基督為主」,使榮耀歸與父上帝.

耶穌騎驢駒進城而不騎戰馬,不是因為他沒有能力召來戰馬,只是因為他「自己卑微」；耶穌被釘死在十字架上,不是因為他沒有能力救自己,只是因為他「不以自己與上帝同等為強奪」.有的,才有可以放下的東西,沒有的,不可能有東西放下.耶穌自我降卑,顯出他是真正的強者,是真正有實力的人.

四．以基督耶穌的心為心

耶穌為甚麼要選取騎驢進城,以柔和謙卑的態度出現,甚至被釘十字架而不反擊呢?馬太福音引述撒迦利亞書那句話的下文是這樣說的:

9:10我必除滅以法蓮的戰車和耶路撒冷的戰馬；爭戰的弓也必除滅.他必向列國講和平；他的權柄必從這海管到那海,從大河管到地極.

這位騎驢的王宣揚的是和平而不是戰爭,他認為和平不是來自軍備競賽,相反是透過裁軍帶來的；平安不是透過武力得到的,而是放棄武力才有的.如果有一位武林高手自廢武功,會否被欺負呢?如果與上帝同等的聖子「取了奴僕的形像」,會否被苦待呢?在這個世界中,如此的遭遇非常大機會,而且難免流血.撒迦利亞先知繼續說:

9:11錫安哪,我因與你立約的血,將你中間被擄而囚的人從無水的坑中釋放出來.

縱使流血,也唯有如此,才能帶來真正的釋放.相反來說,當每個人都武裝起來,帶來的只會是更嚴重的擔憂和懼怕,距離釋放和自由起走越遠.

還記得保羅對腓立比信徒的勸勉嗎?

2:5你們當以基督耶穌的心為心

從今天棕枝主日到下個主日復活主日,讓我們默想耶穌基督最後一週的受難,死亡與復活.透過「為門徒洗腳」的愛,「客西馬尼園」的順服,「十架」的痛苦,進入「墳墓」, 耶穌一步一步downgrade.對於永遠追求upgrade的現代人,耶穌的路途是奇怪的,但卻是人類福祉至關重要的.為甚麼downgrade之路重要及值得學習呢?讓我以這個例子作結.

還記得十多年前真人騷節目「窮富翁大作戰」嗎?其中一位參加者是田北辰,他曾堅信:「如果你有鬥志,即使是弱者,亦可以變成強者.」在節目中,他化身為清潔工人,靠每天的工作賺取生活費,並住在「豪華籠屋」,徹底體驗「普通香港清潔工」的人生.他做了兩天就捱不下去,並且說:「很淒慘,沒有意思,有點像等日子過,很無奈,沒有盼望.」田北辰最後總結說,他這兩天只想著吃東西,沒有其他盼望,一心只想著吃飯,他終於明白貧窮人口的心聲,「怎麼會想計畫?最重要的是解決下一餐.」他不禁問:「當市場經濟到了盡頭,還有沒有公義?」

有一位參與者在體會過清潔工的勞動生活,每晚蜷曲在不及20呎的床位後,毅然辭去珠寶公司行政總裁的職位,與社福機構合作開展了扶貧的工作.

人們以為upgrade才值得追求,原來耶穌自我downgrade能夠拯救世界,上述的例子也顯示有能力的人願意downgrade,也能成為別人的幫助,使世界更美好.



\newpage

\section{基督今復活}
\label{sec:792}
\textbf{高國雄牧師 - 基督今復活}
\newline
\newline
\hyperref[sec:791]{< < < PREV SERMON < < <}
~
\hyperref[sec:toc]{[返目錄]}
~
\hyperref[sec:793]{> > > NEXT SERMON > > >}


約翰福音 20:1-18
\newline
\begin{longtable}{cl}
\hline
\hline
章節 & 經文 (和合本修訂版)\\
\hline
20:1 & \begin{tabularx}{0.7\textwidth}{X} 七日的第一日清早，天還黑的時候，抹大拉的馬利亞來到墳墓，看見石頭已從墳墓挪開了， \end{tabularx} \\ \\ \relax
20:2 & \begin{tabularx}{0.7\textwidth}{X} 就跑來見西門‧彼得和耶穌所愛的那個門徒，對他們說：「有人從墳墓裡把主移走了，我們不知道他們把他放在哪裡。」 \end{tabularx} \\ \\ \relax
20:3 & \begin{tabularx}{0.7\textwidth}{X} 彼得和那門徒就出來，往墳墓去。 \end{tabularx} \\ \\ \relax
20:4 & \begin{tabularx}{0.7\textwidth}{X} 兩個人同跑，那門徒比彼得跑得快，先到了墳墓， \end{tabularx} \\ \\ \relax
20:5 & \begin{tabularx}{0.7\textwidth}{X} 低頭往裡看，看見細麻布還放在那裡，只是沒有進去。 \end{tabularx} \\ \\ \relax
20:6 & \begin{tabularx}{0.7\textwidth}{X} 西門‧彼得隨後也到了，進了墳墓，看見細麻布放在那裡， \end{tabularx} \\ \\ \relax
20:7 & \begin{tabularx}{0.7\textwidth}{X} 又看見耶穌的裹頭巾沒有和細麻布放在一起，是另在一處捲著。 \end{tabularx} \\ \\ \relax
20:8 & \begin{tabularx}{0.7\textwidth}{X} 然後先到墳墓的那門徒也進去，他看見就信了。 \end{tabularx} \\ \\ \relax
20:9 & \begin{tabularx}{0.7\textwidth}{X} 他們還不明白聖經所說耶穌必須從死人中復活的意思。 \end{tabularx} \\ \\ \relax
20:10 & \begin{tabularx}{0.7\textwidth}{X} 於是兩個門徒回自己的住處去了。 \end{tabularx} \\ \\ \relax
20:11 & \begin{tabularx}{0.7\textwidth}{X} 馬利亞卻站在墳墓外面哭。她哭的時候，低頭往墳墓裡看， \end{tabularx} \\ \\ \relax
20:12 & \begin{tabularx}{0.7\textwidth}{X} 看見兩個天使穿著白衣，在安放耶穌身體的地方坐著，一個在頭，一個在腳。 \end{tabularx} \\ \\ \relax
20:13 & \begin{tabularx}{0.7\textwidth}{X} 天使對她說：「婦人，你為甚麼哭？」她對他們說：「因為有人把我主移走了，我不知道他們把他放在哪裡。」 \end{tabularx} \\ \\ \relax
20:14 & \begin{tabularx}{0.7\textwidth}{X} 說了這些話，她轉過身來，看見耶穌站在那裡，卻不知道他是耶穌。 \end{tabularx} \\ \\ \relax
20:15 & \begin{tabularx}{0.7\textwidth}{X} 耶穌問她：「婦人，你為甚麼哭？你找誰？」馬利亞以為他是看園子的，就對他說：「先生，若是你把他移了去，請告訴我，你把他放在哪裡，我去把他移回來。」 \end{tabularx} \\ \\ \relax
20:16 & \begin{tabularx}{0.7\textwidth}{X} 耶穌對她說：「馬利亞。」馬利亞轉過身來，用希伯來話對他說：「拉波尼！」（「拉波尼」就是老師的意思。） \end{tabularx} \\ \\ \relax
20:17 & \begin{tabularx}{0.7\textwidth}{X} 耶穌對她說：「不要拉住我，因為我還沒有升上去見我的父。你到我弟兄那裡去告訴他們，我要升上去見我的父，也是你們的父，見我的神，也是你們的神。」 \end{tabularx} \\ \\ \relax
20:18 & \begin{tabularx}{0.7\textwidth}{X} 抹大拉的馬利亞就向門徒報信：「我已經看見了主。」她又把主對她說的話告訴他們。 \end{tabularx} \\ \\
[1ex]
\hline
\hline
\end{longtable}
講題:基督今復活 Christ the Lord is Risen Today

經課:約翰福音20:1-18

講員:高國雄牧師

場合:香港中文大學崇基學院禮拜堂主日崇拜

日期:2026年4月5日

約翰福音 20:1-18

1七日的第一日清早,天還黑的時候,抹大拉的馬利亞來到墳墓那裡,看見石頭從墳墓挪開了, 2就跑來見西門‧彼得和耶穌所愛的那個門徒,對他們說:「有人把主從墳墓裡挪了去,我們不知道放在哪裡.」 3彼得和那門徒就出來,往墳墓那裡去. 4兩個人同跑,那門徒比彼得跑得更快,先到了墳墓, 5低頭往裡看,就見細麻布還放在那裡,只是沒有進去. 6西門‧彼得隨後也到了,進墳墓裡去,就看見細麻布還放在那裡, 7又看見耶穌的裹頭巾沒有和細麻布放在一處,是另在一處捲著. 8先到墳墓的那門徒也進去,看見就信了.( 9因為他們還不明白聖經的意思,就是耶穌必要從死裡復活.) 10於是兩個門徒回自己的住處去了.11 馬利亞卻站在墳墓外面哭.哭的時候,低頭往墳墓裡看, 12就見兩個天使,穿著白衣,在安放耶穌身體的地方坐著,一個在頭,一個在腳. 13天使對她說:「婦人,你為甚麼哭?」她說:「因為有人把我主挪了去,我不知道放在哪裡.」 14說了這話,就轉過身來,看見耶穌站在那裡,卻不知道是耶穌. 15耶穌問她說:「婦人,為甚麼哭?你找誰呢?」馬利亞以為是看園的,就對他說:「先生,若是你把他移了去,請告訴我,你把他放在哪裡,我便去取他.」 16耶穌說:「馬利亞.」馬利亞就轉過來,用希伯來話對他說:「拉波尼!」(拉波尼就是夫子的意思.) 17耶穌說:「不要摸我,因我還沒有升上去見我的父.你往我弟兄那裡去,告訴他們說,我要升上去見我的父,也是你們的父,見我的上帝,也是你們的上帝.」 18抹大拉的馬利亞就去告訴門徒說:「我已經看見了主.」她又將主對她說的這話告訴他們.

黑暗中的「看見」與誤解

各位弟兄姊妹,主基督復活了!

約翰福音的復活敘事始於「七日的第一日清晨,天還黑的時候」.抹大拉的馬利亞來到墳墓.

她看見的是什麼?

她看見的是「不在」(Absence)——石頭挪開了,身體不見了.

我們的人生也常經歷這種「黑暗中的空洞」.

眼中所看到的就是「失去」.

無論面前是什麼情況,也只會看到「失出了的東西」.

當我們失去工作,失去至親,或是對社會現狀感到絕望時,我們看到的往往只是「空墳墓」.

當只有「失去」時,我們又會有什麼反應及行動呢?

馬利亞的第一反應是典型的理性邏輯:「有人把主從墳墓裡挪走了.」

我們是否也常把上帝局限在「墳墓」裡?

我們以為上帝只存在於過去,存在於儀式中,一旦現實發生震盪,我們就以為上帝「被挪走了」.

奔跑的群體:競爭還是共鳴?

接著是彼得與「耶穌所愛的那門徒」的奔跑競賽.

約翰細膩地描寫了誰先到,誰先看見,誰先進入.

細麻布與裹頭巾: 約翰特別提到裹頭巾是「捲著,在一處獨放」.這不是盜墓者的行徑,盜墓者會匆忙亂丟.耶穌是在秩序中復活的.

看見就信了: 那個門徒「看見就信了」,但經文隨即說「他們還不明白」.

很多時候,我們的「相信」要走在「理解」之前.

信仰不是與他們競賽,而是要以信心及行動投入.

復活的行動力不需要你完全搞懂神學體系,而是在於你願意跑向現場,願意在混亂中看見上帝留下的秩序.

轉身後的相遇:從「先生」到「拉波尼」

門徒回去了,唯有馬利亞站在墳墓外面哭.

馬利亞為什麼哭?

她甚至對天使和耶穌(她以為是園丁)說:「若是你移了他,請告訴我……」.她依然在尋找那具「屍體」.

我們往往被悲哀蒙蔽了雙眼,以至於復活的救主站在面前,我們卻以為祂只是個路人.

耶穌只說了一個詞:「馬利亞.」

這一聲呼喚,擊碎了死亡的冰冷.

約翰福音 10:3下「他按著名叫自己的羊,把羊領出來.」

復活不單只是一個歷史事實,更是一個個人關係的恢復.

當你聽見主呼喚你的名字,你就不再是受害者,而是見證者.

耶穌叫她不要拉住祂,因為祂要升天.

這是在告訴馬利亞(和我們):不能再用舊的方式去擁有耶穌.

祂不再局限於肉身,而是要透過聖靈,與每一位信徒同在.

第一位使徒:最具動力的宣告

馬利亞從一個「哭泣的婦女」轉變為「使徒中的使徒」(Apostle to the Apostles).

她的行動: 她跑去告訴門徒:「我已經看見了主!」

復活節的行動力不是去辯論空墳墓的真偽,而是去分享你的「看見」.

你是否在生命中看見了主?在破碎中看見了醫治?在絕望中看見了光?

清晨的派遣

今天,天已經亮了.

回轉(Turn around): 不要盯著你生命中的「空墳墓」或遺憾不放,轉過身來,聽聽那呼喚你名字的聲音.

整理你的「裹頭巾」: 復活代表秩序勝過混亂.請在你的混亂生活中(可能是人際關係或情緒),靠著主重新建立秩序.

大膽宣告: 找一個機會,對一個身處黑暗的人說:「我已經看見了主.」不需要神學辯論,只需要分享生命.

基督已經復活了!

祂不在此處,祂已在我們心裡,差遣我們走向世界.



\newpage

\section{我信身體復活}
\label{sec:793}
\textbf{李均熊牧師 - 我信身體復活}
\newline
\newline
\hyperref[sec:792]{< < < PREV SERMON < < <}
~
\hyperref[sec:toc]{[返目錄]}
~
\hyperref[sec:794]{> > > NEXT SERMON > > >}


使徒行傳 2:14,22-32
\newline
\begin{longtable}{cl}
\hline
\hline
章節 & 經文 (和合本修訂版)\\
\hline
2:14 & \begin{tabularx}{0.7\textwidth}{X} 彼得和十一個使徒站起來，他就高聲向眾人說：「猶太人和所有住在耶路撒冷的人哪，這件事你們要知道，要側耳聽我的話。 \end{tabularx} \\ \\ \relax
2:15 & \begin{tabularx}{0.7\textwidth}{X} 這些人並不像你們所想的喝醉了，因為現在才早晨九點鐘。 \end{tabularx} \\ \\ \relax
2:16 & \begin{tabularx}{0.7\textwidth}{X} 這正是藉著先知約珥所說的： \end{tabularx} \\ \\ \relax
2:17 & \begin{tabularx}{0.7\textwidth}{X} 『神說：在末後的日子，我要將我的靈澆灌凡血肉之軀的。你們的兒女要說預言；你們的少年要見異象；你們的老人要做異夢。 \end{tabularx} \\ \\ \relax
2:18 & \begin{tabularx}{0.7\textwidth}{X} 在那些日子，我要把我的靈澆灌，甚至給我的僕人和婢女，他們要說預言。 \end{tabularx} \\ \\ \relax
2:19 & \begin{tabularx}{0.7\textwidth}{X} 在天上，我要顯出奇事，在地下，我要顯出神蹟，有血，有火，有煙霧。 \end{tabularx} \\ \\ \relax
2:20 & \begin{tabularx}{0.7\textwidth}{X} 太陽要變為黑暗，月亮要變為血，這都在主大而光榮的日子未到以前。 \end{tabularx} \\ \\ \relax
2:21 & \begin{tabularx}{0.7\textwidth}{X} 那時，凡求告主名的都必得救。』 \end{tabularx} \\ \\ \relax
2:22 & \begin{tabularx}{0.7\textwidth}{X} 「以色列人哪，你們要聽我這些話：拿撒勒人耶穌就是神以異能、奇事、神蹟向你們證明出來的人，這些事是神藉著他在你們中間施行，正如你們自己知道的。 \end{tabularx} \\ \\ \relax
2:23 & \begin{tabularx}{0.7\textwidth}{X} 他既按著神確定的旨意和預知被交與人，你們就藉著不法之人的手把他釘在十字架上，殺了。 \end{tabularx} \\ \\ \relax
2:24 & \begin{tabularx}{0.7\textwidth}{X} 神卻將死的痛苦解除，使他復活了，因為他原不能被死拘禁。 \end{tabularx} \\ \\ \relax
2:25 & \begin{tabularx}{0.7\textwidth}{X} 大衛指著他說：『我看見主常在我眼前，他在我右邊，使我不至於動搖。 \end{tabularx} \\ \\ \relax
2:26 & \begin{tabularx}{0.7\textwidth}{X} 所以我心裡歡喜，我的舌頭快樂，而且我的肉身要安居在指望中。 \end{tabularx} \\ \\ \relax
2:27 & \begin{tabularx}{0.7\textwidth}{X} 因你必不將我的靈魂撇在陰間，也不讓你的聖者見朽壞。 \end{tabularx} \\ \\ \relax
2:28 & \begin{tabularx}{0.7\textwidth}{X} 你已將生命的道路指示我，必使我在你面前充滿快樂。』 \end{tabularx} \\ \\ \relax
2:29 & \begin{tabularx}{0.7\textwidth}{X} 「諸位弟兄，先祖大衛的事，我可以坦然地對你們說：他死了，也埋葬了，而且他的墳墓直到今日還在我們這裡。 \end{tabularx} \\ \\ \relax
2:30 & \begin{tabularx}{0.7\textwidth}{X} 既然大衛是先知，他知道神曾向他起誓，要從他的後裔中立一位坐在他的寶座上。 \end{tabularx} \\ \\ \relax
2:31 & \begin{tabularx}{0.7\textwidth}{X} 他預先看見了，就講論基督的復活，說：『他不被撇在陰間；他的肉身也不見朽壞。』 \end{tabularx} \\ \\ \relax
2:32 & \begin{tabularx}{0.7\textwidth}{X} 這耶穌，神已經使他復活了，我們都是這事的見證人。 \end{tabularx} \\ \\
[1ex]
\hline
\hline
\end{longtable}
詩篇 16:1-11
\newline
\begin{longtable}{cl}
\hline
\hline
章節 & 經文 (和合本修訂版)\\
\hline
16:1 & \begin{tabularx}{0.7\textwidth}{X} 神啊，求你保佑我，因為我投靠你。 \end{tabularx} \\ \\ \relax
16:2 & \begin{tabularx}{0.7\textwidth}{X} 我曾對耶和華說：「你是我的主，我的福氣惟獨從你而來。」 \end{tabularx} \\ \\ \relax
16:3 & \begin{tabularx}{0.7\textwidth}{X} 論到世上的聖民，他們是尊貴的人，是我最喜悅的。 \end{tabularx} \\ \\ \relax
16:4 & \begin{tabularx}{0.7\textwidth}{X} 追逐別神的，他們的愁苦必增加；他們所澆奠的血我不獻上，我嘴唇也不提別神的名號。 \end{tabularx} \\ \\ \relax
16:5 & \begin{tabularx}{0.7\textwidth}{X} 耶和華是我的產業，是我杯中的福分；我所得的，你為我持守。 \end{tabularx} \\ \\ \relax
16:6 & \begin{tabularx}{0.7\textwidth}{X} 用繩量給我的地界，坐落在佳美之處；我的產業實在美好。 \end{tabularx} \\ \\ \relax
16:7 & \begin{tabularx}{0.7\textwidth}{X} 我要稱頌那指引我的耶和華，在夜間我的心腸也指教我。 \end{tabularx} \\ \\ \relax
16:8 & \begin{tabularx}{0.7\textwidth}{X} 我讓耶和華常在我面前，因他在我右邊，我就不致動搖。 \end{tabularx} \\ \\ \relax
16:9 & \begin{tabularx}{0.7\textwidth}{X} 因此，我的心歡喜，我的靈快樂；我的肉身也要安然居住。 \end{tabularx} \\ \\ \relax
16:10 & \begin{tabularx}{0.7\textwidth}{X} 因為你必不將我的靈魂撇在陰間，也不讓你的聖者見地府。 \end{tabularx} \\ \\ \relax
16:11 & \begin{tabularx}{0.7\textwidth}{X} 你必將生命的道路指示我。在你面前有滿足的喜樂，在你右手中有永遠的福樂。 \end{tabularx} \\ \\
[1ex]
\hline
\hline
\end{longtable}
彼得前書 1:3-9
\newline
\begin{longtable}{cl}
\hline
\hline
章節 & 經文 (和合本修訂版)\\
\hline
1:3 & \begin{tabularx}{0.7\textwidth}{X} 願頌讚歸於我們主耶穌基督的父神！他曾照自己的大憐憫，藉著耶穌基督從死人中復活，重生了我們，使我們有活的盼望， \end{tabularx} \\ \\ \relax
1:4 & \begin{tabularx}{0.7\textwidth}{X} 好得到不朽壞、不玷污、不衰殘、為你們存留在天上的基業， \end{tabularx} \\ \\ \relax
1:5 & \begin{tabularx}{0.7\textwidth}{X} 就是為你們這些藉著信、蒙神大能保守的人，能獲得他所預備、到末世要顯現的救恩。 \end{tabularx} \\ \\ \relax
1:6 & \begin{tabularx}{0.7\textwidth}{X} 雖然你們必須在百般試煉中暫時憂愁，你們要為此喜樂 ， \end{tabularx} \\ \\ \relax
1:7 & \begin{tabularx}{0.7\textwidth}{X} 使你們的信心既被考驗，就比那被火試煉仍然能壞的金子更顯寶貴，可以在耶穌基督顯現的時候得著稱讚、榮耀、尊貴。 \end{tabularx} \\ \\ \relax
1:8 & \begin{tabularx}{0.7\textwidth}{X} 雖然你們沒有見過他，卻是愛他；如今雖看不見，你們卻因信他而有說不出來、滿有榮光的喜樂， \end{tabularx} \\ \\ \relax
1:9 & \begin{tabularx}{0.7\textwidth}{X} 因為你們得到信心的效果，就是靈魂的得救。 \end{tabularx} \\ \\
[1ex]
\hline
\hline
\end{longtable}
約翰福音 20:19-31
\newline
\begin{longtable}{cl}
\hline
\hline
章節 & 經文 (和合本修訂版)\\
\hline
20:19 & \begin{tabularx}{0.7\textwidth}{X} 那日（就是七日的第一日）晚上，門徒因怕猶太人，所在的地方門都關了。耶穌來，站在當中，對他們說：「願你們平安！」 \end{tabularx} \\ \\ \relax
20:20 & \begin{tabularx}{0.7\textwidth}{X} 說了這話，他把手和肋旁給他們看。門徒一看見主就喜樂了。 \end{tabularx} \\ \\ \relax
20:21 & \begin{tabularx}{0.7\textwidth}{X} 於是耶穌又對他們說：「願你們平安！父怎樣差遣了我，我也照樣差遣你們。」 \end{tabularx} \\ \\ \relax
20:22 & \begin{tabularx}{0.7\textwidth}{X} 說了這話，他向他們吹一口氣，說：「領受聖靈吧！ \end{tabularx} \\ \\ \relax
20:23 & \begin{tabularx}{0.7\textwidth}{X} 你們赦免誰的罪，誰的罪就得赦免；你們不赦免誰的罪，誰的罪就不得赦免。」 \end{tabularx} \\ \\ \relax
20:24 & \begin{tabularx}{0.7\textwidth}{X} 那十二使徒中，有個叫低土馬的多馬，耶穌來的時候，他沒有和他們在一起。 \end{tabularx} \\ \\ \relax
20:25 & \begin{tabularx}{0.7\textwidth}{X} 其他的門徒就對他說：「我們已經看見主了。」多馬卻對他們說：「除非我看見他手上的釘痕，用我的指頭探入那釘痕，用我的手探入他的肋旁，我絕不信。」 \end{tabularx} \\ \\ \relax
20:26 & \begin{tabularx}{0.7\textwidth}{X} 過了八日，門徒又在屋裡，多馬也和他們在一起。門都關了，耶穌來，站在當中，說：「願你們平安！」 \end{tabularx} \\ \\ \relax
20:27 & \begin{tabularx}{0.7\textwidth}{X} 然後他對多馬說：「把你的指頭伸到這裡來，看看我的手；把你的手伸過來，探入我的肋旁。不要疑惑，總要信！」 \end{tabularx} \\ \\ \relax
20:28 & \begin{tabularx}{0.7\textwidth}{X} 多馬回答，對他說：「我的主！我的神！」 \end{tabularx} \\ \\ \relax
20:29 & \begin{tabularx}{0.7\textwidth}{X} 耶穌對他說：「你因為看見了我才信嗎？那沒有看見卻信的有福了。」 \end{tabularx} \\ \\ \relax
20:30 & \begin{tabularx}{0.7\textwidth}{X} 耶穌在他門徒面前另外行了許多神蹟，沒有記錄在這書上。 \end{tabularx} \\ \\ \relax
20:31 & \begin{tabularx}{0.7\textwidth}{X} 但記載這些事是要使你們信耶穌是基督，是神的兒子，並且使你們信他，好因著他的名得生命。 \end{tabularx} \\ \\
[1ex]
\hline
\hline
\end{longtable}
講題:我信身體復活 I Believe in the Resurrection of the Body

經課:使徒行傳2:14,22-32;詩篇16:1-11;彼得前書1:3-9;約翰福音20:19-31

講員:李均熊牧師

場合:香港中文大學崇基學院禮拜堂主日崇拜

日期:2026年4月12日

起:基督徒信甚麼?

其實我懷疑很多基督徒,其實並不相信身體復活.

我們每星期 都會念信經哪一句:我信身體復活,但大家如果誠實一點問自己:

我們心裡真的相信這一句嗎?

但通常基督徒其實相信甚麼呢?我懷疑一般基督徒相信的可能是另一套東西:靈魂不滅,死後上天堂.

因為很多時候,我們以為基督教的盼望是這樣的:

人死了,身體留在地上, 最後 腐化了甚麼也沒有了；

真正的「我」——那個「靈魂」——就去到天堂,

遠離這個會老,會痛,會壞的身體.

這套想法聽起來很熟悉,

但其實它不是聖經的語言,

而是希臘的思想.

希臘人相信:

身體是次等的,暫時的；

靈魂才是不朽,才值得拯救.

其實,弟兄姊妹,

這種把「靈魂」和「身體」徹底分開的想法,

不是從聖經來的.

在今日的詩篇經課,詩篇第 16 篇,

詩人這樣說:

「因為你必不將我的靈魂撇在陰間,也不叫你的聖者見朽壞.」(詩16:10)

我們今天讀「靈魂」,

很容易就以為是:

一個可以離開身體,單獨存在的部分.

但在舊約裡,

「靈魂」這個字,

指的不是一個脫離肉身的東西,

而是指整個活著的人.

換句話說,

在舊約的信仰裡,

人不是「有一個靈魂」,

人本身就是一個有身體的生命.

死亡不是靈魂離開身體,

而是整個人落入死亡；

而神的拯救,

也不是只拯救一部分,

而是使整個人不被死亡拘禁.

後來,當這些經文被翻譯成希臘文,

「靈魂」被翻成 psyche(psychology 的字根),

慢慢就被希臘哲學的思想侵佔——

把靈魂想成可以逃離肉身,不朽的本體.

但彼得在《使徒行傳》引用詩篇 16 篇時,

他要證明的,

不是靈魂不死,

而是——

神必使那被埋葬的耶穌,從死裡復活.

這從一開始,

就不是一個「靈魂上天堂」的信仰,

而是一個身體從死亡中被拯救的盼望.

不知不覺,連教會也被這種思想影響.

我們口裡說「身體復活」,

心裡卻可能只期待「靈魂上天堂」.

承:多馬要求見到的「復活」

然而,今日福音經課中,我們聽到我們熟識的多馬,所要求要見到的,卻是身體復活的證據!「我非看見他手上的釘痕,用指頭探入那釘痕,又用手探入他的肋旁,我總不信.」

或者大家真心想一想,如果我問你:「你死後希望發生什麼事?」

很多人或者會答:

「希望靈魂去天堂.」

但我們唸信經時,唸的是什麼?

「我信身體復活.」

或者我們所熟識的多馬,是個多疑,看不到摸不到就不會相信的多馬,但我們卻沒有留意到,多馬的信心可能比我們大部份的信徒更準確,更符合信經的信仰.多馬拒絕一個沒有身體的復活；

而我們,卻可能滿足於這樣的信仰.這也可能最後耶穌不單沒有怪責他的多疑,反而更滿足他的要求.

因為經課並沒有停在多馬的要求那裡.

約翰福音告訴我們,

耶穌再次來到門徒中間,

並沒有對多馬說一句責備的話.

耶穌反而對他說:

「伸過你的指頭來,摸我的手；

伸出你的手來,探入我的肋旁.」

請大家留意:

耶穌沒有說——

「你不應該這樣要求.」

「你應該學會屬靈一點.」

耶穌給多馬的,

正正就是一個帶著釘痕,帶著傷口的身體.

換句話說,

耶穌沒有否定多馬對「身體復活」的期待,

反而用自己的身體,去肯定這正是復活的內容.

那一刻,多馬所見到的,

不是靈魂顯現,

不是天堂異象,

而是——

被釘十字架的身體,卻站在他面前.

轉:原來神沒有放棄身體,也沒有放棄這個佢創造ge 物質世界

弟兄姊妹,

這就是為甚麼教會一直誦讀那句聽起來「奇怪」的信經:

「我信身體復活.」

這句話不是多餘的,

也不是象徵性的,

而是對世界的一個宣告——

神沒有放棄身體,

也沒有放棄這個受傷的世界.

復活的耶穌仍然有傷痕,

意味著:

受苦的經歷不會當成沒存在過,好似粉筆字被抹掉,而係復活ge 時候 被 賦予新的意義.

因為死亡沒有最後發言權.

所以基督教的盼望,

不是靈魂逃離身體,

而是神親自進入死亡,

並且把死亡裡的身體,帶進新生命.

今天我們再唸信經的時候,

我們唸的不只是一句傳統的說話,

而是在宣告一個更大的盼望:

不是「我將來會去哪裡」,

而是——

神將要把我,連同這個世界,一同更新,再創造.

合:對我們今天的意義?

那麼,弟兄姊妹,

「我信身體復活」對我們今天究竟有甚麼意思?

首先,這句話提醒我們:

我們今天的身體,不是多餘的,

不是暫時用完就可以丟棄的外殼.

這個會疲倦,會生病,會衰老的身體,

正正就是神所愛,所拯救,將來要更新的身體.(當然神怎樣去更新要到將來才知道)

所以,當我們面對病痛,軟弱,年老,

甚至面對死亡的時候,

我們不是說:「不要緊,反正這個身體沒甚麼用.」

而是敢於相信:

神仍然在乎,神仍然記得,神仍然會拯救.

第二,「我信身體復活」也改變我們今天怎樣看待別人的身體.

如果神沒有放棄人的身體,

那麼貧窮的,受傷的,被忽略的,被壓迫的身體,

就不可能是無關重要的.

我們的信仰,

不只是關乎將來靈魂去到哪裡,

而是關乎今天如何對待眼前的生命.

最後,「我信身體復活」,

讓我們在這個仍然充滿死亡與破碎的世界裡,

活得不一樣.

因為我們知道:

這個世界不是注定要被丟棄,

而是正在等候將來神的更新.

而且這個將來,不是模糊的,也不是不確定的.

正如今日書信經課彼得前書的說話,

藉著耶穌基督從死裡復活,

神已經重生了我們,

叫我們有一個活潑的盼望,

可以得著那

不能朽壞,不能玷污,不能衰殘,

為我們存留在天上的基業.

這個「存留在天上」的基業,

不是叫我們逃離地上的肉身,

而是確定一件事:

今日的我, 種種與我身體關連的事情:無論是我在世上的學習,經歷,成長,甚至因肉身而有同所經營的關係,

已經在神手中被保守,不是靠我們抓緊,所以彼得說,

我們這些因信蒙神能力保守的人,

必能得著

那到末世要顯現的救恩.

弟兄姊妹,

這就是為甚麼我們今天可以在試煉中仍然站立,

在破碎中仍然盼望,

因為我們的將來,

不是一個問號,

而是一個已經因為主的復活而確立的應許.

係靠上帝的能力托住

不會被死亡奪走,

不會被時間消磨.甚至更新以後就不能朽壞,不能玷污,不能衰殘.

所以今天,

當我們再次誦讀信經,

我們不是只在背誦一段傳統的文字,

而是在用信心宣告:

我信身體復活,

我信神沒有放棄我,

也沒有放棄這個世界.

我信,

今天一切會朽壞的,

不會是結局；

一切正在破碎的,

不會是終結；

因為在復活的主裡,

一切都要被復原,被更生,成為新的了.

其實我懷疑很多基督徒,其實並不相信身體復活.

我們每星期 都會念信經哪一句:我信身體復活,但大家如果誠實一點問自己:

我們心裡真的相信這一句嗎?

但通常基督徒其實相信甚麼呢?我懷疑一般基督徒相信的可能是另一套東西:靈魂不滅,死後上天堂.

因為很多時候,我們以為基督教的盼望是這樣的:

人死了,身體留在地上, 最後 腐化了甚麼也沒有了；

真正的「我」——那個「靈魂」——就去到天堂,

遠離這個會老,會痛,會壞的身體.

這套想法聽起來很熟悉,

但其實它不是聖經的語言,

而是希臘的思想.

希臘人相信:

身體是次等的,暫時的；

靈魂才是不朽,才值得拯救.

其實,弟兄姊妹,

這種把「靈魂」和「身體」徹底分開的想法,

不是從聖經來的.

在今日的詩篇經課,詩篇第 16 篇,

詩人這樣說:

「因為你必不將我的靈魂撇在陰間,也不叫你的聖者見朽壞.」(詩16:10)

我們今天讀「靈魂」,

很容易就以為是:

一個可以離開身體,單獨存在的部分.

但在舊約裡,

「靈魂」這個字,

指的不是一個脫離肉身的東西,

而是指整個活著的人.

換句話說,

在舊約的信仰裡,

人不是「有一個靈魂」,

人本身就是一個有身體的生命.

死亡不是靈魂離開身體,

而是整個人落入死亡；

而神的拯救,

也不是只拯救一部分,

而是使整個人不被死亡拘禁.

後來,當這些經文被翻譯成希臘文,

「靈魂」被翻成 psyche(psychology 的字根),

慢慢就被希臘哲學的思想侵佔——

把靈魂想成可以逃離肉身,不朽的本體.

但彼得在《使徒行傳》引用詩篇 16 篇時,

他要證明的,

不是靈魂不死,

而是——

神必使那被埋葬的耶穌,從死裡復活.

這從一開始,

就不是一個「靈魂上天堂」的信仰,

而是一個身體從死亡中被拯救的盼望.

不知不覺,連教會也被這種思想影響.

我們口裡說「身體復活」,

心裡卻可能只期待「靈魂上天堂」.

承:多馬要求見到的「復活」

然而,今日福音經課中,我們聽到我們熟識的多馬,所要求要見到的,卻是身體復活的證據!「我非看見他手上的釘痕,用指頭探入那釘痕,又用手探入他的肋旁,我總不信.」

或者大家真心想一想,如果我問你:「你死後希望發生什麼事?」

很多人或者會答:

「希望靈魂去天堂.」

但我們唸信經時,唸的是什麼?

「我信身體復活.」

或者我們所熟識的多馬,是個多疑,看不到摸不到就不會相信的多馬,但我們卻沒有留意到,多馬的信心可能比我們大部份的信徒更準確,更符合信經的信仰.多馬拒絕一個沒有身體的復活；

而我們,卻可能滿足於這樣的信仰.這也可能最後耶穌不單沒有怪責他的多疑,反而更滿足他的要求.

因為經課並沒有停在多馬的要求那裡.

約翰福音告訴我們,

耶穌再次來到門徒中間,

並沒有對多馬說一句責備的話.

耶穌反而對他說:

「伸過你的指頭來,摸我的手；

伸出你的手來,探入我的肋旁.」

請大家留意:

耶穌沒有說——

「你不應該這樣要求.」

「你應該學會屬靈一點.」

耶穌給多馬的,

正正就是一個帶著釘痕,帶著傷口的身體.

換句話說,

耶穌沒有否定多馬對「身體復活」的期待,

反而用自己的身體,去肯定這正是復活的內容.

那一刻,多馬所見到的,

不是靈魂顯現,

不是天堂異象,

而是——

被釘十字架的身體,卻站在他面前.

轉:原來神沒有放棄身體,也沒有放棄這個佢創造ge 物質世界

弟兄姊妹,

這就是為甚麼教會一直誦讀那句聽起來「奇怪」的信經:

「我信身體復活.」

這句話不是多餘的,

也不是象徵性的,

而是對世界的一個宣告——

神沒有放棄身體,

也沒有放棄這個受傷的世界.

復活的耶穌仍然有傷痕,

意味著:

受苦的經歷不會當成沒存在過,好似粉筆字被抹掉,而係復活ge 時候 被 賦予新的意義.

因為死亡沒有最後發言權.

所以基督教的盼望,

不是靈魂逃離身體,

而是神親自進入死亡,

並且把死亡裡的身體,帶進新生命.

今天我們再唸信經的時候,

我們唸的不只是一句傳統的說話,

而是在宣告一個更大的盼望:

不是「我將來會去哪裡」,

而是——

神將要把我,連同這個世界,一同更新,再創造.

合:對我們今天的意義?

那麼,弟兄姊妹,

「我信身體復活」對我們今天究竟有甚麼意思?

首先,這句話提醒我們:

我們今天的身體,不是多餘的,

不是暫時用完就可以丟棄的外殼.

這個會疲倦,會生病,會衰老的身體,

正正就是神所愛,所拯救,將來要更新的身體.(當然神怎樣去更新要到將來才知道)

所以,當我們面對病痛,軟弱,年老,

甚至面對死亡的時候,

我們不是說:「不要緊,反正這個身體沒甚麼用.」

而是敢於相信:

神仍然在乎,神仍然記得,神仍然會拯救.

第二,「我信身體復活」也改變我們今天怎樣看待別人的身體.

如果神沒有放棄人的身體,

那麼貧窮的,受傷的,被忽略的,被壓迫的身體,

就不可能是無關重要的.

我們的信仰,

不只是關乎將來靈魂去到哪裡,

而是關乎今天如何對待眼前的生命.

最後,「我信身體復活」,

讓我們在這個仍然充滿死亡與破碎的世界裡,

活得不一樣.

因為我們知道:

這個世界不是注定要被丟棄,

而是正在等候將來神的更新.

而且這個將來,不是模糊的,也不是不確定的.

正如今日書信經課彼得前書的說話,

藉著耶穌基督從死裡復活,

神已經重生了我們,

叫我們有一個活潑的盼望,

可以得著那

不能朽壞,不能玷污,不能衰殘,

為我們存留在天上的基業.

這個「存留在天上」的基業,

不是叫我們逃離地上的肉身,

而是確定一件事:

今日的我, 種種與我身體關連的事情:無論是我在世上的學習,經歷,成長,甚至因肉身而有同所經營的關係,

已經在神手中被保守,不是靠我們抓緊,所以彼得說,

我們這些因信蒙神能力保守的人,

必能得著

那到末世要顯現的救恩.

弟兄姊妹,

這就是為甚麼我們今天可以在試煉中仍然站立,

在破碎中仍然盼望,

因為我們的將來,

不是一個問號,

而是一個已經因為主的復活而確立的應許.

係靠上帝的能力托住

不會被死亡奪走,

不會被時間消磨.甚至更新以後就不能朽壞,不能玷污,不能衰殘.

所以今天,

當我們再次誦讀信經,

我們不是只在背誦一段傳統的文字,

而是在用信心宣告:

我信身體復活,

我信神沒有放棄我,

也沒有放棄這個世界.

我信,

今天一切會朽壞的,

不會是結局；

一切正在破碎的,

不會是終結；

因為在復活的主裡,

一切都要被復原,被更生,成為新的了.



\newpage

\section{只有一位受苦的上帝能救我們}
\label{sec:794}
\textbf{鄧瑞強牧師 - 只有一位受苦的上帝能救我們}
\newline
\newline
\hyperref[sec:793]{< < < PREV SERMON < < <}
~
\hyperref[sec:toc]{[返目錄]}
~
\hyperref[sec:795]{> > > NEXT SERMON > > >}


馬可福音 8:29-9:1
\newline
\begin{longtable}{cl}
\hline
\hline
章節 & 經文 (和合本修訂版)\\
\hline
8:29 & \begin{tabularx}{0.7\textwidth}{X} 他又問他們：「你們說我是誰？」彼得回答他：「你是基督。」 \end{tabularx} \\ \\ \relax
8:30 & \begin{tabularx}{0.7\textwidth}{X} 於是耶穌切切地囑咐他們不可對任何人說起他。 \end{tabularx} \\ \\ \relax
8:31 & \begin{tabularx}{0.7\textwidth}{X} 從此，他教導他們說：「人子必須受許多的苦，被長老、祭司長和文士棄絕，並且被殺，三天後復活。」 \end{tabularx} \\ \\ \relax
8:32 & \begin{tabularx}{0.7\textwidth}{X} 耶穌明白地說了這話，彼得就拉著他，責備他。 \end{tabularx} \\ \\ \relax
8:33 & \begin{tabularx}{0.7\textwidth}{X} 耶穌轉過來看著門徒，斥責彼得說：「撒但，退到我後邊去！因為你不體會神的心意，而是體會人的意思。」 \end{tabularx} \\ \\ \relax
8:34 & \begin{tabularx}{0.7\textwidth}{X} 於是他叫眾人和門徒來，對他們說：「若有人要跟從我，就當捨己，背起自己的十字架來跟從我。 \end{tabularx} \\ \\ \relax
8:35 & \begin{tabularx}{0.7\textwidth}{X} 因為凡要救自己生命的，必喪失生命；凡為我和福音喪失生命的，必救自己的生命。 \end{tabularx} \\ \\ \relax
8:36 & \begin{tabularx}{0.7\textwidth}{X} 人就是賺得全世界，賠上自己的生命，有甚麼益處呢？ \end{tabularx} \\ \\ \relax
8:37 & \begin{tabularx}{0.7\textwidth}{X} 人還能拿甚麼換生命呢？ \end{tabularx} \\ \\ \relax
8:38 & \begin{tabularx}{0.7\textwidth}{X} 凡在這淫亂罪惡的世代，把我和我的道當作可恥的，人子在他父的榮耀裡與聖天使一同來臨的時候，也要把那人當作可恥的。」 \end{tabularx} \\ \\ \relax
9:1 & \begin{tabularx}{0.7\textwidth}{X} 耶穌又對他們說：「我實在告訴你們，站在這裡的，有人在沒經歷死亡以前，必定看見神的國帶著能力臨到。」 \end{tabularx} \\ \\
[1ex]
\hline
\hline
\end{longtable}
講題:只有一位受苦的上帝能救我們  Only a Suffering God can Help

經課:馬可福音8:29-9:1

講員:鄧瑞強牧師

場合:香港中文大學崇基學院禮拜堂主日崇拜

日期:2026年4月19日

各位弟兄姊妹,早安.願神的恩典常與你們同在.

在這蒼茫大地,群雄逐鹿中原,拜的,都是戰神.當大家不再講規矩,返回叢林法則,就看看誰能在廝殺中取勝了.一生追求「溝通理性」(Communicative Rationality)的哲人哈伯瑪斯(Jürgen Habermas)近日離世,屍骨未寒,各國行動的重點卻是儲備軍火.大家都在想,只有戰神,才能救自己.

第一世紀,羅馬帝國的鐵蹄,踐踏在巴勒斯坦的土地上.猶太人苦不堪言,他們祈求,救世主像無敵戰神那樣來臨,驅散敵人,救助他們.他們望穿秋水,終於來了耶穌.他們看見他行神蹟,對他滿有期望.他們想:他就是那戰神嗎?

我們聽聽耶穌的話.

可8:29〔耶穌〕問他們說:「你們說我是誰?」彼得回答說:「你是基督.」

可8:30耶穌就禁戒他們,不要告訴人.

可8:31從此,他教訓他們說:「人子必須受許多的苦,被長老,祭司長,和文士棄絕,並且被殺,過三天復活.」

可8:32耶穌明明地說這話,彼得就拉著他,勸他【按:「勸」這個字,與下一節「責備」這個字,是同一個字.耶穌「斥責」污鬼時,也是用這個字.在這裡,彼得直頭是在斥責耶穌.】.

可8:33耶穌轉過來,看著門徒,就責備彼得說:「撒旦,退我後邊去吧!因為你不體貼上帝的意思,只體貼人的意思.」【按:非常罕有的場面,使徒互相斥責,耶穌責備彼得站在撒旦那一邊.】

【耶穌面對更大的群眾】

可8:34於是叫眾人和門徒來,對他們說:「若有人要跟從我,就當捨己,背起他的十字架來跟從我.

可8:35因為,凡要救自己生命的,必喪掉生命；凡為我和福音喪掉生命的,必救了生命.

可8:36人就是賺得全世界,賠上了自己的生命,有什麼益處呢?

可8:37人還能拿什麼換生命呢?

可8:38凡在這淫亂罪惡的世代,把我和我的道當做可恥的,人子在他父的榮耀裡,同聖天使降臨的時候,也要把那人當做可恥的.」

可9:1耶穌又對他們說:「我實在告訴你們,站在這裡的,有人在沒嘗死味以前,必要看見上帝的國大有能力臨到.」【按:這是指他的復活,他的復活更新了天地,讓這被死亡氣息濃罩的世界,窺見生命的光明.】

耶穌講的,與當日門徒的期望,與今天信徒的期望,落差很大.

彼得很激動,對耶穌說:我們跟你,就是因為你勁,你強,你好打.像黑幫大佬,在這講拳頭的世界,能殺出一條血路.怎知你講受苦,被棄絕,被殺.若真有能力復活,就不要在死後才顯這大能.在死前顯,效果不是更大嗎?在十架上,大搖大擺地走下來,殺盡釘人十架的羅馬兵,宣傳效果不更震撼嗎?跟你,就是不想受苦.

今天,信徒對耶穌的想像,仍是希望他是戰神,無堅不摧,無往不利.病的時候,是希望耶穌醫治,而不是希望他叫我們承受病的痛苦,體會受苦的意義,然後以這體會去鼓勵其他病人.投資的時候,是祈求耶穌祝福,一切順風順水,一本萬利,而不希望耶穌提醒我們要知足,要親手勞碌做些有意義的事.不少學生現時做功課,是叫人工智能將文章或書本的精要抽出來,並按相關觀點,輸出成論文.學生的祈禱,是這論文能通過大學的剽竊檢測系統,能拿高分.從來不希望耶穌吩咐我們,就算時間多麼短促,也要一篇文章一篇文章的讀；不計較分數的高低,也要憑自己的努力寫一篇論文來.凡走捷徑救了分數的,必喪掉學術生命；凡放下取分數的旁門左道的,必救了學術生命.有些神學生,畢業時,找教會,入職作傳道.他們的祈求,是人工高,福利好,責任少.他們想的,是他們心中渴望的教會,而不是思考,哪裡才是最需要他們捨己服事的教會.不少人的人生規劃,都渴想蘇軾所講的:「無災無難到公卿」.耶穌講的,卻是:要有天國的遠象,為了讓上帝的恩典改變這個世界,要背起十架,要甘心多背負一點仁義的責任,要比別人多作一點犧牲.耶穌講的人生,是一門虧本生意.講這些東西,在今天,很多信徒都會斥責他.… 或者,你不敢斥責耶穌,但肯定會斥責我.

我們見到,耶穌背著十架前行.他的苦難,來自世人的反對,因他不滿足世人及一般信徒的期望,反而想改變這個令不少人活得舒舒服服的世界.

俄國大文豪陀思妥耶夫斯基(Fyodor Dostoevsky)的巨著《卡拉馬佐夫兄弟》(The Brothers Karamazov)裡有一段出名的情節,話說,在16世紀,耶穌再次來到這個世界,他想給人「自由」,或者說,他渴望帶領那些迷失在世界種種虛假價值裡的人,重新找回生命的永恆意義.怎知,宗教大法官(一個大祭司般的人物)對他講,你走吧,這裡的信徒其實要的是麵包,不是自由.你講的生命意義,反以令習慣了虛假價值的信徒混亂,搞亂了這個世界.你走吧,你若不走,我會處死你.

只要彼得或信徒,要的,是世界的價值,就會與耶穌爭論,衝突.只要世界敬拜的是戰神,而耶穌在努力的背十架,世界和耶穌,就會衝突不休.

20世紀,二次大戰期間,在納粹的集中營,住著一個囚犯,一個反納粹的德國人,一個德國的神學家,他叫潘霍華(Dietrich Bonhoeffer).在監倉裡,他日夜思考幾個關乎歷史的將來的問題:戰後德國的重建問題,基督信仰在戰後的歷史前景問題,主耶穌對這世界意味著什麼這問題.

在他的《監獄書信》裡,他寫道:「在世上,上帝是軟弱的,無能的,正正就是這樣,也唯獨是這樣,祂在我們身邊,幫助我們.…基督幫助我們,不是靠他的全能,而是靠他的軟弱和受苦.…在世上,人類一般的宗教渴望指引人追求上帝的大能,一個懂得在困境中變法術的上帝.聖經卻指引人面向上帝的無能和受苦,只有這位受苦的上帝能救我們.」(注:1944年7月16日致Eberhard Bethge信)

潘霍華想的,是歷史前景的問題.一個遲早會戰敗的德國,將以什麼來重建呢?支撐著歐洲文明的基督教,前途從何說起?對這殘垣敗瓦的世界,基督還能意味著什麼嗎?

戰火仍在燃燒,潘霍華想到,人類歷史要有前途,就不要再敬拜戰神,而要跟從那位背著十架前行的耶穌.人人都拜戰神,人人都要勝過別人,人人都要取得最大的利益,則無可避免地在叢林中混戰,以力取勝.以暴易暴,將會在暴力中無限輪迴.他殺了我的父親,我要上山學武,有朝一日,下山為父報仇.終於,你殺了他.然後,他的兒子又上山學武,為的是有一朝,他能報殺父之仇.歷史就是無盡的後人,不斷練武,練的武功越來越利害,殺人的能力越來越強.未來歲月,核彈會越來越多,射程會越來越遠,殺傷力會越來越大.來到如此光景,是因為無論國家,世人,信徒,都敬拜戰神,對耶穌的排斥,力度越來越大.

潘霍華想到,人類要有前途,只能靠那背著十架前行的上帝.這位上帝沒有以核彈,炸毀祂看不順眼的世界.祂放下身段,捨棄自己的尊嚴,與無能的人同在,與受苦的人同行.這位上帝,以脆弱與無能,去反抗這個暴力的世界,去揭露神聖的天國.就因為上帝甘願如此受苦,祂真的明白軟弱而受苦的人.看似離棄世人的上帝,祂原來不單沒有離棄我們,還在苦難中深深與我們結連.正因祂在苦難中的同在,我們才能在苦難中找到來自上帝的希望.耶穌對眾人說,經過十架的黑夜之後,會看到天國的黎明.

若祂是戰神,則祂固然威風凜凜,大有能力,但祂只能高高在上,無法親自體驗我們的苦難,無法在苦難的深淵裡與我們同在.一顆子彈,能解決惡人,但這顆子彈不會為我們的痛苦流淚,不能給苦難的我們一點同情安慰,也不可以令我們對天地間的仁愛有多一點信任.唯有那位為我們擋子彈的耶穌,那位承受苦難而與苦難的我們不離不棄的耶穌,才能在我們心靈深處,生發出對仁愛的信任,把捉到來自上天的救贖,燃點起一丁點照亮前路的火光.那怕是一丁點火光,那怕是柔弱不堪的仁愛幼苗,就能給人勇氣,讓人放下武器,以愛前行.耶穌背著十架,緩慢前行,跟著他,反能看到歷史的出路.

這些講法,必然引來世人的譏笑和漫罵.在這叢林世界,講的是現實政治.你不拿起武器,下一刻,你便不在了.現實真的就是這樣.畢竟,耶穌真的死了.但,這證明耶穌不對嗎?或者,應反過來想想,若沒有耶穌,這個叢林世界還有希望嗎?

或遲或早,人總要抉擇:在叢林世界裡,以力決戰,看看誰能留到最後,抑或,甘願吃一點虧,放下一點尊嚴,多承受一點責任,好能終止暴力的循環.神學家潘霍華說:只有這位受苦的上帝能救我們.

耶穌毫不含糊,也絕無妥協地要求門徒捨己,背起十字架來跟從他.或許,是時候檢視一下我們的整個人生.十架,與世人追求的價值和目標相反.世人的目標,是權力和名利.採取的手段,是暴力.十架卻是放下身段,與卑下的人同在,與受苦的人同行.堅拒以世界的價值,作為人生的目標.當時的彼得,拒絕耶穌這種想法.但耶穌堅持,這才是正路.

最後,讓我們靜下來,再次聽聽耶穌的呼喚.

可8:34〔耶穌〕說:「若有人要跟從我,就當捨己,背起他的十字架來跟從我.

可8:35因為,凡要救自己生命的,必喪掉生命；凡為我和福音喪掉生命的,必救了生命.」

願神賜福你們.

願榮耀歸予聖父,聖子,聖靈.



\newpage

\section{活石}
\label{sec:795}
\textbf{李均熊牧師 - 活石}
\newline
\newline
\hyperref[sec:794]{< < < PREV SERMON < < <}
~
\hyperref[sec:toc]{[返目錄]}
~
\hyperref[sec:796]{> > > NEXT SERMON > > >}


彼得前書 2:2-10
\newline
\begin{longtable}{cl}
\hline
\hline
章節 & 經文 (和合本修訂版)\\
\hline
2:2 & \begin{tabularx}{0.7\textwidth}{X} 要愛慕那純淨的靈奶，像初生的嬰孩愛慕奶一樣，好使你們藉著它成長，以致得救， \end{tabularx} \\ \\ \relax
2:3 & \begin{tabularx}{0.7\textwidth}{X} 因為你們已經嘗過主恩的滋味。 \end{tabularx} \\ \\ \relax
2:4 & \begin{tabularx}{0.7\textwidth}{X} 要親近主，他是活石，雖然被人所丟棄，卻是神所揀選、所珍貴的。 \end{tabularx} \\ \\ \relax
2:5 & \begin{tabularx}{0.7\textwidth}{X} 你們作為活石，要被建造成屬靈的殿，成為聖潔的祭司，藉著耶穌基督獻上蒙神悅納的屬靈祭物。 \end{tabularx} \\ \\ \relax
2:6 & \begin{tabularx}{0.7\textwidth}{X} 因為經上說：「看哪，我把一塊石頭放在錫安—一塊蒙揀選、珍貴的房角石；信靠他的人必不蒙羞。」 \end{tabularx} \\ \\ \relax
2:7 & \begin{tabularx}{0.7\textwidth}{X} 所以，這石頭在你們信的人是珍貴的；在那不信的人卻有話說：「匠人所丟棄的石頭已作了房角的頭塊石頭。」 \end{tabularx} \\ \\ \relax
2:8 & \begin{tabularx}{0.7\textwidth}{X} 又說：「作了絆腳的石頭，使人跌倒的磐石。」他們絆跌，因為不順從這道，這也是預定的。 \end{tabularx} \\ \\ \relax
2:9 & \begin{tabularx}{0.7\textwidth}{X} 不過，你們是被揀選的一族，是君尊的祭司，是神聖的國度，是屬神的子民，要使你們宣揚那召你們出黑暗入奇妙光明者的美德。 \end{tabularx} \\ \\ \relax
2:10 & \begin{tabularx}{0.7\textwidth}{X} 「你們從前不是子民，現在卻成了神的子民；從前未曾蒙憐憫，現在卻蒙了憐憫。」 \end{tabularx} \\ \\
[1ex]
\hline
\hline
\end{longtable}
講題:活石  Living Stones

經課:彼得前書2:2-10

講員:李均熊牧師

場合:香港中文大學崇基學院禮拜堂主日崇拜

日期:2026年5月3日

起:石頭——不是活石,而是殺人的石頭

復活節過後,我們往往以為,教會會慢慢帶領我們離開黑暗與死亡,走向較為光明,確定的敘事.經文應該一步一步把我們推向喜樂與盼望,好像一切都已經翻進新一頁.然而,教會卻在復活後第五主日,把我們重新帶回一個極度不安的場景──史提反倒在石頭之下.

這個故事被放在復活期中段,本身就構成了一個提問:如果耶穌真的復活了,為甚麼世界仍然會用石頭來回應那些見證復活的人?史提反所看見的,是站立在上帝右邊的主；但他所承受的,卻不是理解或對話,而是用來處決他的石頭.復活並沒有令世界立刻變得安全可靠.

耶穌在今日的福音經課中對門徒說:「你們心裡不要憂愁.」但這不是在問題解決之後才說的話,而是在門徒仍然困惑,仍然害怕,仍然不知道前路何在之前說的.復活信仰從來不是一條沒有問題的道路,而是在問題仍然真實存在的世界中繼續前行.

承:主復活後的教會

從史提反開始,教會已經知道,見證復活並不保證免於受傷.到了彼得前書成書的時候,傷害的形式不再只是公開迫害,而是轉化為日常而地方性的排擠.基督徒沒有天天面對刀劍,卻在日復一日的生活裡被慢慢推離社會核心.

這種排擠有著十分具體的歷史原因……(尼祿當政對基督的「迫害」)……在古代羅馬世界,宗教不是私人信念,而是社會忠誠的一部分.行業行會的聚會,往往以祭祀與共膳開始,既是宗教行為,也是經濟與社會網絡的關鍵.當基督徒因信仰拒絕參與這些行為時,並非只是『唔拜』,而是被理解為拒絕一套被視為理所當然的公共責任.

因此,基督徒不一定犯法,卻逐漸被視為不合群,不可靠,甚至對城市秩序構成潛在威脅.再加上他們獨特的生活方式與倫理選擇,被謠言與誤解包圍,被指為道德可疑的一群.這些指控未必立即帶來官方行動,卻足以讓一個群體在鄰里,家庭與經濟網絡中慢慢失去位置.

於是,彼得前書的信徒活在一種極具磨蝕性的張力之中:他們仍然留在羅馬帝國的城巿中,好像生活如常,但其實愈來愈不像身處城市的必要部分.這不是劇烈的衝突,而是長期的耗損,迫使人反覆自問:我是否仍然在能見證上帝,在建立天國的工程之中?

轉:主乃活石

正是在這樣的處境之中,彼得選擇用「石頭」這個意象,並刻意不把焦點放在信徒身上.他沒有說教會要成為甚麼樣的建築,也沒有鼓勵信徒要把自己磨成有用的材料；他首先說的是基督——那一位「雖被人所棄,卻是神所揀選,所寶貴的」活石.在古代建築的語境中,被棄的石頭意味著它曾被檢驗,被判定,被拒絕,失去了任何進入工程的可能.彼得並沒有否定這個判定的真實性,反而將之視為理解基督道路的關鍵.反而指出:耶穌正是被世界判定為無用的那一塊.然而,這塊被棄的石頭,卻成了房角石,決定整座建築的方向與穩定.這並不是成功的保證,而是一個顛覆性的重新理解:被棄,未必等於走錯路；不被需要,也未必代表不被揀選.

更為關鍵的,並不在於那塊石頭本身是否帶來改變,而在於整個建築的取向轉移至誰的手中.那一塊被棄的石頭,並沒有憑自身的價值翻轉處境,也沒有為自己爭取位置；它只是被動地被安放於另一個位置,成為房角石.經文以一連串被動的動詞描寫這個過程——被棄,被揀選,被立——彷彿刻意淡化石頭本身的能動性,將注意力引向那真正推動建造的主體.

基督之所以成為房角石,不是因為世界重新認可了祂,而是因為神在自身的工程中,重新確立了建築的中心與方向.那曾被擱置的石頭,正是在這樣的轉移之中,被放置為決定整體結構穩定的關鍵所在.於是,房角石的意義不在於證明成功或翻身,而在於揭示一種次序的轉換:建築的根基與指向,已不再由人的判斷決定,而是交付於那位在被棄之處行使主權的神.

也正是在這個神學邏輯下,彼得才進一步稱信徒為「活石」,而且不是會自己活／自己移動自行砌合的石頭,而是「被建造成為靈宮」的石頭.教會不是一個靠群體共識慢慢搭建的建築,而是一個正在被建造的工程；而信徒的身份,也不是靠自我定位來獲得,而是在基督作為房角石的前提下,被放進去,被承托,被安排.

合:等候被放置的活石,被建造的教會

當我們把這段經文放回今日香港讀,可以類比的可能是那種生活節奏.我們仍然可以敬拜,教會仍然存在,信仰仍然合法；但不少信徒都感到,信仰的公共位置正在收縮,有些價值不再容易被理解,有些信念被歸類為純粹私人的選擇.

這種轉變未必帶來恐懼,卻慢慢塑造出一種退後與收斂的生活姿態.在這樣的處境中,彼得前書所呈現的,並不是一條指向重新站回中心的位置的道路.信徒也未被催促去證明自己仍然具有某種社會上的價值.經文反而把目光引向另一個方向——回到與主的關係之中,留在那位活石附近,在尚未完成的工程裡,等候被放置.

這種等候並非消極,而是具體的生活操練:在小而實在的信仰群體中彼此守望,在不確定中持續禱告,在仍然可以的空間裡誠實而不張揚地活出信念,使自己不至於成為一塊孤立的石頭.

或許今日的教會,更像一個仍在施工的工地,而不是一座已完成的殿宇.結構未明,方向未清,但正是在這樣的狀態之中,我們被召繼續靠近那位被棄,卻復活的房角石,在等候之中,被建造.

彼得前書所描繪的,其實對今日信徒可能並不陌生的感受.當人仍然生活在城中,仍然如常工作與往來,卻發現自己愈來愈不像是其中不可或缺的一部分,孤立感便在不知不覺中形成.原本理所當然的連結開始鬆動,一些關係不再穩固,信徒也很容易在這樣的過程中,把自己理解為一塊孤零零的石頭——需要自行證明價值,自行承受壓力,甚至自行決定是否還有留下來的理由.

這正是彼得前書讀者所處的境況.他們並非突然被趕離城市,而是在日常生活中,慢慢被推向邊緣；而這種變化,往往令人把結構性的排擠內化為個人的責任.當支持的網絡一點一點鬆開,人便不自覺地以為,撐得住是責任,撐不住則是失敗.於是,本應由群體共同承當的重量,逐漸化成個人默默硬撐的重負.

焦慮正是在這裡悄悄生成.因為一旦把自己理解為孤立的石頭,信徒便會以為:若我撐不住,站不穩,派不上用場,最終便只剩下被丟棄的命運.於是,那些原本由眾人一同承受的壓力,逐漸被誤會成必須各自面對的責任；那些本應彼此分擔的重量,也在不知不覺中,被推回到個人的肩上.人在這樣的處境裡,很容易以為,撐得住是本分,撐不住便成了自己的問題.然而,彼得前書中所描繪的信徒,從來不是被召去單獨完成甚麼,更不是被留下來處理自身去留的命運.他們所面對的,乃是一個早已開始,且並不由他們掌控的建造過程；他們被召進入其中,不是獨力承擔,而是在既有的結構裡,被放置,被承托.

正因如此,「活石」這個稱呼本身,已經否定了孤立生存的可能性.石頭之所以能夠承受重量,從來不是因為它本身足夠堅硬,而是因為它被放置在一個結構之內,與其他石頭彼此貼合,彼此支撐,彼此分擔承重.離開了這種彼此關係,石頭在工程中便失去了存在的意義.同樣,彼得所描繪的教會,也不是一群各自站立,各自硬撐的信徒,而是一個在基督內逐步被建造,彼此承托的整體.

而更關鍵的,是這整個結構的承重焦點,始終不在任何一塊活石身上.彼得不斷把視線拉回那一位已被立為房角石的基督.房角石的功能,並不是取代其他石頭,而是承托,定位,校準整座建築的方向與穩定.在古代的建築世界裡,人們並不像我們今天這樣,先在地下建立一個看不見,卻可以完全承托整座建築的底層結構,然後才開始向上建造.建築的穩定,更多依賴一塊被準確放置,分量極重的房角石.這塊石頭不單承受壓力,更重要的是,它決定了整座建築其他石頭如何排列,如何彼此貼合,如何共同分擔重量.只要房角石放得正,其他石頭才能站得穩.

彼得稱基督為房角石,所描繪的,並不是一個所有重量聚集的中心,而是一個使整座建築得以成形的起點.其餘的石頭,正是在與祂的對齊之中,找到位置,也分擔重量.

沒有房角石,各自站立的石頭只會變成一堆不相干,無法承重的碎石；有了房角石,它們才開始成為一座能夠承托壓力的建築.正因為如此,活石才能承重,卻從來不是獨自承重；我們彼此支撐,卻不是彼此救贖；我們共同分擔壓力,卻不是共同決定根基.真正承擔重量的,始終是那一位先一步承擔了「被棄」,並仍然托住整個工程的基督.

或許,彼得前書最後留給我們的,並不是一個清楚可控的答案,而是一個需要被持續重讀的圖像.我們不是被留下來各自站立的石頭,也不是被要求證明自己是否足夠堅強的一群.我們只是被放在一座仍在建造中的結構之內,而那結構是否能夠站立,從來不取決於任何一塊活石本身.房角石已經被放定,其餘的石頭,正在被安排,被承托,被彼此結連.

在這樣的建築之中,我們未必知道自己的位置有多重要,也未必看得見完成的形態；但只要我們仍然留在那位活石所設定的關係裡,建造便沒有停止.

或許,信心並不是確定自己撐得住,而是容許自己不再單獨承重；不是急於看見結果,而是在仍然未完成的狀態中,繼續被放置.

因此,在講章將要結束的時候,我們並不是被帶到一個結論之前,而是被帶到一個動作之中.彼得前書所描繪的建造,並不是一項已完成的工程,而是一個仍然進行中的現實；而今天,我們所要做的,也不是為這個建造加上最後一句話,而是走向桌前,領受.

在接著的聖餐神中,讓我們再次被提醒,我們並不是憑著自己的站立而成為教會,也不是靠著各自撐住而連結在一起.我們所倚靠的,始終是那一位已經為我們被棄,被立的基督.餅被擘開,杯被遞上,並不是因為我們已經找到了自己的位置,而是因為我們願意再一次,把自己交回那位房角石的次序之中.

讓我們來到聖桌前,讓我們留意站在身邊弟兄姊妹彼此的存在,記得我們不是以孤立的石頭前來,而是以一群被呼召,被連結,仍在被建造的活石前來.在同一張桌子前領受,同一個身體被分給眾人,我們再次被置於彼此之間,也被置於基督之內.或許建造仍未完成,但在這份領受之中,我們被提醒:真正承托我們的,並不是自己,而是那位為我們成為生命糧的主.



\newpage

\section{耶穌不在的時候}
\label{sec:796}
\textbf{林豪恩先生 - 耶穌不在的時候}
\newline
\newline
\hyperref[sec:795]{< < < PREV SERMON < < <}
~
\hyperref[sec:toc]{[返目錄]}
~
\hyperref[sec:797]{> > > NEXT SERMON > > >}


彼得前書 3:13-22
\newline
\begin{longtable}{cl}
\hline
\hline
章節 & 經文 (和合本修訂版)\\
\hline
3:13 & \begin{tabularx}{0.7\textwidth}{X} 你們若熱心行善，有誰會害你們呢？ \end{tabularx} \\ \\ \relax
3:14 & \begin{tabularx}{0.7\textwidth}{X} 即使你們為義受苦，也是有福的。不要怕人的威嚇，也不要驚慌； \end{tabularx} \\ \\ \relax
3:15 & \begin{tabularx}{0.7\textwidth}{X} 只要心裡奉主基督為聖，尊他為主。有人問你們心中盼望的理由，要隨時準備答覆； \end{tabularx} \\ \\ \relax
3:16 & \begin{tabularx}{0.7\textwidth}{X} 不過，要以溫柔、敬畏的心回答。要存無虧的良心，使你們在何事上被毀謗，就在何事上使那些凌辱你們在基督裡有好品行的人自覺羞愧。 \end{tabularx} \\ \\ \relax
3:17 & \begin{tabularx}{0.7\textwidth}{X} 神的旨意若是要你們因行善受苦，這總比因行惡受苦好。 \end{tabularx} \\ \\ \relax
3:18 & \begin{tabularx}{0.7\textwidth}{X} 因為基督也曾一次為罪受苦 ，就是義的代替不義的，為要引領你們到神面前。在肉體裡，他被治死；但在靈裡，他復活了。 \end{tabularx} \\ \\ \relax
3:19 & \begin{tabularx}{0.7\textwidth}{X} 他藉這靈也曾去向那些在監獄裡的靈傳道， \end{tabularx} \\ \\ \relax
3:20 & \begin{tabularx}{0.7\textwidth}{X} 就是那些從前在挪亞預備方舟、神容忍等待的時候不信從的人。當時進入方舟，藉著水得救的不多，只有八個人。 \end{tabularx} \\ \\ \relax
3:21 & \begin{tabularx}{0.7\textwidth}{X} 這水所預表的洗禮，現在藉著耶穌基督的復活拯救你們，不是除掉肉體的污穢，而是向神懇求有無虧的良心。 \end{tabularx} \\ \\ \relax
3:22 & \begin{tabularx}{0.7\textwidth}{X} 耶穌已經到天上去，在神的右邊，眾天使、有權柄的、有權能的都服從了他。 \end{tabularx} \\ \\
[1ex]
\hline
\hline
\end{longtable}
約翰福音 14:15-21
\newline
\begin{longtable}{cl}
\hline
\hline
章節 & 經文 (和合本修訂版)\\
\hline
14:15 & \begin{tabularx}{0.7\textwidth}{X} 「你們若愛我，就會遵守我的命令。 \end{tabularx} \\ \\ \relax
14:16 & \begin{tabularx}{0.7\textwidth}{X} 我要求父，父就賜給你們另外一位保惠師，使他永遠與你們同在。 \end{tabularx} \\ \\ \relax
14:17 & \begin{tabularx}{0.7\textwidth}{X} 他就是真理的靈，是世人不能接受的。因為他們既看不見他，也不認識他；你們卻認識他，因他常與你們同在，也要在你們裡面。 \end{tabularx} \\ \\ \relax
14:18 & \begin{tabularx}{0.7\textwidth}{X} 我不會撇下你們為孤兒，我必到你們這裡來。 \end{tabularx} \\ \\ \relax
14:19 & \begin{tabularx}{0.7\textwidth}{X} 再過不久，世人不再看見我，你們卻會看見我，因為我活著，你們也要活著。 \end{tabularx} \\ \\ \relax
14:20 & \begin{tabularx}{0.7\textwidth}{X} 到那日，你們就會知道我在父裡面，你們在我裡面，我也在你們裡面。 \end{tabularx} \\ \\ \relax
14:21 & \begin{tabularx}{0.7\textwidth}{X} 有了我的命令而又遵守的人，就是愛我的；愛我的人，我父要愛他，我也要愛他，並且要親自向他顯現。」 \end{tabularx} \\ \\
[1ex]
\hline
\hline
\end{longtable}
講題:耶穌不在的時候  When Jesus Is Not Here

經課:彼得前書3:13-22;約翰福音14:15-21

講員:林豪恩先生

場合:香港中文大學崇基學院禮拜堂主日崇拜

日期:2026年5月10日

一. 約翰福音:耶穌要離開了

今天是母親節,祝願天下母親身心靈健康,願上主的恩典常與你們同在.談到母親,想起一個和母親有關的心理學實驗.實驗是這樣進行的:母親和幼兒被帶到一個不熟悉的房間,房間裡擺滿了玩具.當母親在場時,幼兒會被鼓勵去探索周遭環境；但在幾分鐘後,一位陌生人進來,然後母親離開.在經過這次短暫分離之後,母親又會回到現場.研究人員要觀察的是,母親離開房間以及重回房間時,幼兒的反應如何?這個心理學實驗和今天的經課有甚麼關係呢?剛才我們聆聽約翰福音的時候,聽到耶穌預告自己要離開門徒的片段:

14:18我不撇下你們為孤兒,我必到你們這裡來.

14:19還有不多的時候,世人不再看見我,你們卻看見我；因為我活著,你們也要活著.

耶穌要離開,當耶穌不在的時候,跟隨耶穌的人會有甚麼反應呢?然後,有一天耶穌重來的時候,跟隨者又有甚麼反應呢?

回到那個心理學實驗,研究人員把幼兒對母親離開後及重回時的反應最初分成三種,第一種稱為安全型依附(Secure Attachment):安全型依附的幼兒在母親在場時,將母親視為探索世界的安全基地(secure base)；當母親離去之後,幼兒會變得很沮喪,但是當母親返回時,幼兒會去接近母親並感到很快樂.第二種反應稱為迴避型不安全依附(Anxious-Avoidant Insecure Attachment):迴避型不安全依附則會和母親顯得疏遠,母親在場時無視母親,母親離開時也不會產生明顯失望的情緒,母親再次回來時也不會去尋求母親支持.第三種稱為焦慮/矛盾型不安全依附(Anxious-Resistant Insecure Attachment):焦慮/矛盾型不安全依附的小孩,則會在母親離去時顯得相當不安與焦慮,但在母親回來時,也不會變得特別開心,他們無法輕易的被安撫,同時對母親抱持著憤怒及抗拒；但在此同時,他們又很渴望和母親接觸.這個心理學實驗的發現後來被擴充及廣泛應用於親子關係及其他類型的人際關係.今天,我們借用這個角度來閱讀剛才讀過的幾段經課.

我們要問的是,耶穌離開之後,他的跟隨者有甚麼反應呢?他們和耶穌維持著怎樣的關係呢?雖然當年耶穌的跟隨者的處境和我們現在的處境不一樣,然而,類似的問題也值得我們問一下自己,我與那位不在的耶穌的關係是如何的呢?感到安全?想迴避他?還是矛盾地對徘徊於迎拒之間呢?

其實耶穌離開的情況和那個心理學實驗並不相同.在實驗中,母親離開之前沒有預告,但耶穌離開之前卻充分向跟隨著說明,是「有交帶」的.

第一,說明會回來

14:18我不會撇下你們為孤兒,我必到你們這裡來.

第二,賜下聖靈作為保惠師

14:16我要求父,父就另外賜給你們一位保惠師,叫他永遠與你們同在,

第三,留下指引及應許

14:21有了我的命令又遵守的,這人就是愛我的；愛我的必蒙我父愛他,我也要愛他,並且要向他顯現.」

不過,就算耶穌有「交帶」,提供應許,並賜下聖靈作為保惠師,是否就等於他沒有離開呢?耶穌在與不在,跟隨者的經驗是不同的.耶穌在的時候是看得見,摸得到的；耶穌離開後,是看不見,摸不到的.從彼得前書可見,作者是需要因為耶穌不在而向跟隨著解釋.

二. 耶穌不在的時候:彼得前書受眾的困苦

3:22耶穌已經到天上去,在上帝的右邊,眾天使,有權柄的,有權能的都服從了他.

跟隨者被告知耶穌在天上,有權柄；然而,對於跟隨者來說,耶穌在天上的權柄卻沒有被他們在地上體會得到,因為字裡行間顯示他們正承受著困苦.

3:14你們就是為義受苦,也是有福的.不要怕人的威嚇,也不要驚慌；

3:16存著無虧的良心,叫你們在何事上被毀謗,就在何事上可以叫那誣賴你們在基督裡有好品行的人自覺羞愧.

書信中出現「受苦」,「威嚇」,「驚慌」,「被毀謗」這些字眼,顯示出耶穌不在的時候,跟隨者直接面對攻擊和承受困苦.耶穌在的時候,狀況不一樣嗎?是不一樣的,因為有耶穌為他們遮風擋雨,讓我們看看約翰福音這個片段:

18:3猶大領了一隊兵,和祭司長並法利賽人的差役,拿著燈籠,火把,兵器,就來到園裡.

18:4耶穌知道將要臨到自己的一切事,就出來對他們說:「你們找誰?」

18:5他們回答說:「找拿撒勒人耶穌.」耶穌說:「我就是!」賣他的猶大也同他們站在那裡.

18:6耶穌一說「我就是」,他們就退後倒在地上.

18:7他又問他們說:「你們找誰?」他們說:「找拿撒勒人耶穌.」

18:8耶穌說:「我已經告訴你們,我就是.你們若找我,就讓這些人去吧.」

由此可見,耶穌在的時候,跟隨者有他保護.耶穌不在的時候,跟隨者直接受到攻擊.因此,對於耶穌不在的時候跟隨著所受的苦,彼得是需要向他們解說的:

3:17上帝的旨意若是叫你們因行善受苦,總強如因行惡受苦.

3:18因基督也曾一次為罪受苦,就是義的代替不義的,為要引我們到上帝面前.按著肉體說,他被治死；按著靈性說,他復活了.

彼得勸勉跟隨者學習耶穌的榜樣「為義受苦」,然而,無論有甚麼應許,將來有甚麼獎賞,「苦」是當下的狀況,是此刻的煎熬.曾經在痛苦中的人必定感受至深——無論那是身體上的痛苦,心靈上的痛苦,或者關係上的痛苦.雖然有人告訴你耶穌在你心中,聖靈在你心中,但在心中反而是看不到,捉不實,也不能好像當年耶穌在的時候一樣出來遮風擋雨.耶穌不在,會否令人疑惑:「聖靈在心中／耶穌在心中」這些說法只是一雙情願的想像呢?對正在受苦的人有甚麼幫助呢?甚至連「聖靈在心中／耶穌在心中」都不能肯定又怎樣呢?如果有人出現這些疑問,是不是等於他們不屬靈呢?是不是等於他們放棄神呢?不一定,因為在歷史中的屬靈先賢也有類似的經驗.

德蘭修女就經歷了四十年的心靈黑夜.在靈魂深處,她感受到耶穌經常缺席.心靈的黑夜不是德蘭修女的專利.聖十字約翰 (St John of the Cross) 談及黑夜「是人類和基督徒的典型經驗」.耶穌不在的時候,跟隨耶穌的人也有可能經常掙扎在黑夜中.改編自遠藤周作的同名小說的電影《沉默》(Silence),探討信仰,苦難與神的沉默,其中有句對白說:「你沉重的沉默如此可怕.我祈禱,但我迷失了.或者我只是在向虛無祈禱嗎?」

三．耶穌不在,就可以宣告上帝已死嗎?

古時有人如此挑戰受苦的約伯:「你仍然持守你的純正嗎?你背棄上帝,死了吧!」(約伯記2:9)現代有人宣告「上帝已死」.可能這些挑戰及口號都反映著人們的處境及經驗,然而,人類真的願意否定神嗎?人類果真可以「當神死了」嗎?使徒行傳記載了保羅在雅典遇見的狀況給我們認識人類的另一個面向.

17:22保羅站在亞略‧巴古當中,說:「眾位雅典人哪,我看你們凡事很敬畏鬼神.

17:23我遊行的時候,觀看你們所敬拜的,遇見一座壇,上面寫著『未識之神』.

雅典在希臘世界中是文化之都,是學術重鎮,然而,雅典人的心中仍然有神,只是因為對神的認識有限,於是稱之為「未識之神」.耶穌離開了,人類不能在物質世界中見到耶穌,是否等於人類能夠否定耶穌呢?或者有人認為可以,他們認為隨著人類的發展和科技的進步,宗教就會消失,神也會被消滅.人類世界的確在發展,科技的確在進步,然而宗教消失了嗎?神被消滅了嗎?以下三種發生在我們身邊的狀況顯示出人類的需要:

第一種狀況發生在五,六十年前的某個地方,當時宗教遭到毀滅性打擊,寺廟,教堂,道觀等被破壞,經典被燒毀,僧尼和神職人員被批鬥及強迫還俗.縱使如此暴烈,人們的宗教信仰被根除了嗎?暴風雨過後,鄉間的小寺廟被重建起來,香火重新燃點；在大城及名山,那些被封閉的寺廟,教堂,道觀都重新開放,善男信女不斷增加.他們不只是升斗小民,還有知識分子,透過求神問卜,酬神進香,冥想,讀經,參加宗教聚會等,安頓一己的身心靈,更有的致力於社會倫理的昇華改造,重建人文世界.

第二個例子發生在我們附近:曾經有一間「應用科技研究院」被揭發豪花十八萬元公帑「睇風水」,引起政府及社會的強烈不滿.有議員批評應科院作為科研機構,不應該豪花十八萬元公帑看風水,反問「風水仲緊要過科研?」

第三個例子是喪禮:如果大家去過殯儀館,就會發現原來很多喪禮是以宗教儀式舉行的,最常見有道教(如破地獄,燒紙紮)與佛教(如誦經,超渡),以及基督教宗教的安息禮拜.這種現象所反映的人類的觀念及需要,相信都不用多解釋了.

四．耶穌不在的時候分享耶穌的在

雖然耶穌不在,但人類也不輕易就此而「當神無到」.如何在耶穌不在的時候言說耶穌呢?如何可能在拿不出個神來的狀況之下言說神呢?使徒保羅在使徒行傳中作了一個嘗試.

17:24創造宇宙和其中萬物的上帝,既是天地的主,就不住人手所造的殿,

17:25也不用人手服事,好像缺少甚麼；自己倒將生命,氣息,萬物,賜給萬人.

17:26他從一本造出萬族的人,住在全地上,並且預先定準他們的年限和所住的疆界,

17:27要叫他們尋求上帝,或者可以揣摩而得,其實他離我們各人不遠；

17:28我們生活,動作,存留,都在乎他.就如你們作詩的,有人說:『我們也是他所生的.』

17:29我們既是上帝所生的,就不當以為上帝的神性像人用手藝,心思所雕刻的金,銀,石.

17:30世人蒙昧無知的時候,上帝並不監察,如今卻吩咐各處的人都要悔改.

17:31因為他已經定了日子,要藉著他所設立的人按公義審判天下,並且叫他從死裡復活,給萬人作可信的憑據.」

古希臘哲學家認為人類有三大人生問題,其中兩題是「我從哪裡來」,以及「我往哪裡去」.大家可能已經看出,保羅向雅典人介紹神大致上就是請他們思考「我從哪裡來」,以及「我往哪裡去」這兩個人生必問的問題.歷世歷代不少先賢先哲就是在思索這兩個人生問題的時候認識神,並且與神建立關係.

除了使徒保羅的示範,詩篇也提供了另一種言說那位看不見的神的可能性:

66:10上帝啊,你曾試驗我們,熬煉我們,如熬煉銀子一樣.

66:11你使我們進入網羅,把重擔放在我們的身上.

66:12你使人坐車軋我們的頭；我們經過水火,你卻使我們到豐富之地.

66:16凡敬畏上帝的人,你們都來聽!我要述說他為我所行的事.

66:17我曾用口求告他；我的舌頭也稱他為高.

66:18我若心裡注重罪孽,主必不聽.

66:19但上帝實在聽見了；他側耳聽了我禱告的聲音.

66:20上帝是應當稱頌的!他並沒有推卻我的禱告,也沒有叫他的慈愛離開我.

詩人如何述說那位摸不到,看不見的神呢?他用的方法是分享與神交往的經驗,類似我們現在的「講見證」.希伯來書也是如此勸勉信徒:

12:1所以,既然我們有這許多見證人如同雲彩圍繞著我們,就該卸下各樣重擔和緊緊纏累的罪,以堅忍的心奔那擺在我們前頭的路程,

12:2仰望我們信心的創始成終者耶穌.

有人認為信仰是對真理一個本質性的,肯定的贊同,卻不是直接的抓實.回想那個心理學實驗,可以啟發我們檢視和神的關係嗎?例如我們和神之間是「安全型依附」關係,還是「迴避型不安全依附」關係呢?會不會是「焦慮/矛盾型不安全依附」關係呢?無論是哪一種關係,其實都不太重要,因為我們和神的關係並不決定於我們,而是建基於神的恩典及憐憫,以及信賴祂的信實.



\newpage

\section{夕陽西下斷腸人在天涯}
\label{sec:797}
\textbf{鄧瑞強牧師 - 夕陽西下,斷腸人在天涯}
\newline
\newline
\hyperref[sec:796]{< < < PREV SERMON < < <}
~
\hyperref[sec:toc]{[返目錄]}
~
\hyperref[sec:798]{> > > NEXT SERMON > > >}


路加福音 7:11-17
\newline
\begin{longtable}{cl}
\hline
\hline
章節 & 經文 (和合本修訂版)\\
\hline
7:11 & \begin{tabularx}{0.7\textwidth}{X} 過了不久，耶穌往一座城去，這城名叫拿因，他的門徒和一大群人與他同行。 \end{tabularx} \\ \\ \relax
7:12 & \begin{tabularx}{0.7\textwidth}{X} 當他走近城門時，有一個死人被抬出來。這人是他母親獨生的兒子，而他母親又是寡婦。城裡的許多人與她一同送殯。 \end{tabularx} \\ \\ \relax
7:13 & \begin{tabularx}{0.7\textwidth}{X} 主看見那寡婦就憐憫她，對她說：「不要哭。」 \end{tabularx} \\ \\ \relax
7:14 & \begin{tabularx}{0.7\textwidth}{X} 於是耶穌進前來，按著槓，抬的人就站住了。耶穌說：「年輕人，我吩咐你，起來！」 \end{tabularx} \\ \\ \relax
7:15 & \begin{tabularx}{0.7\textwidth}{X} 那死人就坐了起來，開始說話，耶穌就把他交給他的母親。 \end{tabularx} \\ \\ \relax
7:16 & \begin{tabularx}{0.7\textwidth}{X} 眾人都驚奇，歸榮耀給神，說：「有大先知在我們當中興起了！」又說：「神眷顧了他的百姓！」 \end{tabularx} \\ \\ \relax
7:17 & \begin{tabularx}{0.7\textwidth}{X} 關於耶穌的這事就傳遍了猶太和周圍地區。 \end{tabularx} \\ \\
[1ex]
\hline
\hline
\end{longtable}
講題:夕陽西下,斷腸人在天涯  A Lonely Soul Strays at the Horizon’s End

經課:路加福音7:11-17

講員:鄧瑞強牧師

場合:香港中文大學崇基學院禮拜堂主日崇拜

日期:2026年5月17日

各位弟兄姊妹,早安.願神的恩典常與你們同在.

大家有沒有試過絕望?

親手經營幾十年,以為可以養老的家園,轉眼間,便灰飛煙滅了.一起成長的玩伴,幾十年來互相支持的老街坊,忽然間,告別人間.至愛的家人,相依的老伴,一下子,便永別了.當中的淒慘,真無法言傳.個中人,連絕望的淚,都流不出來.

曾經,有一個晚上,這一刻我還和太太說話,下一刻她就離開了這個塵世.整個世界,一剎那間,就只是一剎那,就返回太古之前的蒼白與虛無.絕望臨到的時候,只是空白一片,空白到連眼淚都不懂得流.

路加福音筆下的耶穌,就是神走入人間,擁抱絕望的人.

不同的福音書,描繪出耶穌的不同畫像.馬太筆下的耶穌,像摩西,走上山,頒佈律法,宣講登山寶訓,塑造神的新子民.馬可筆下的耶穌,背著沉重的十架,呼喚信徒跟著他.而路加筆下的耶穌,最富人情味,最喜歡和人吃飯.喜歡飯局的弟兄姊妹,應多讀路加福音.這福音書有不少獨有的故事,都是很有人性溫度的,如:好撒馬利亞人的故事,浪子回頭的故事,稅吏撒該的故事.稅吏撒該,因耶穌和他吃了一餐飯,就將一半財富分給窮人.這頓飯好貴,但他認為值得.

今日講道,講路加福音的一個小故事,也是路加獨有的.故事以悲劇為背景,在夕陽西下的歲月,一個斷腸人,在天涯.絕望至極,哀傷至極.耶穌,以神聖的仁慈,與這絕望的人相遇.人性的溫暖,讓冰冷的生命,重拾希望.

經文:路加福音7:11-17

路7:11過了不多時,耶穌往一座城去,這城名叫拿因,他的門徒和極多的人與他同行.

路7:12將近城門,有一個死人被抬出來.這人是他母親獨生的兒子；他母親又是寡婦.有城裡的許多人同著寡婦送殯.

路7:13主看見那寡婦,就憐憫她,對她說:「不要哭!」

路7:14於是進前按著槓,抬的人就站住了.耶穌說:「少年人,我吩咐你,起來!」

路7:15那死人就坐起,並且說話.耶穌便把他交給他母親.

路7:16眾人都驚奇,歸榮耀與上帝,說:「有大先知在我們中間興起來了!」又說:「上帝眷顧了他的百姓!」

路7:17他這事的風聲就傳遍了猶太和周圍地方.

這個故事很簡單,只有七節經文,沒有冗長情節,沒有長篇大論,卻包含了福音的多項元素:死亡與生命,失去與重獲,絕望與希望,哭泣與安慰,「天無眼」與「天有眼」,人間的苦難與上帝的憐憫,「獨生子」的死亡與復活,神話語的創造權能…等等,有人說,這是福音信息最濃縮的版本,是福音的「微型」.福音,並不是抽象的大道理.福音,是神與人相遇的故事,是神擁抱絕望的人.

(1)生與死的相遇

這個小故事,始於生與死的相遇.

路7:11…他(耶穌)的門徒和極多的人與他同行.

耶穌是生命的主,這是一散發生命的行列.

另一邊廂,是死亡的行列,是絕望的行列,是世人的行列.

路7:12將近城門,有一個死人被抬出來.這人是他母親獨生的兒子；他母親又是寡婦.有城裡的許多人同著寡婦送殯.

寥寥數語,已道盡當事人的淒涼.

這是一個寡婦,她過去所倚靠的丈夫,已經死了.她有一個獨生子,她的將來,她的希望,全靠他了.如今,這兒子也死了.沒有了過去,失去了將來,這寡婦的當下,只是一片空白.在這故事裡,她的無言,是震天動地的哀號.

在這生與死的相遇中,耶穌主動中止死亡行列的前行.

(2)主看見,主憐憫

路7:13主看見那寡婦,就憐憫她……

耶穌看見這位寡婦,耶穌看見在痛苦深淵裡絕望的人,耶穌看見我們.

耶穌憐憫她.憐憫,就是神的本性.人的憐憫,常只是一種同情的情感.神的憐憫,卻是行動,一種堅決的行動.因為憐憫,神進入人間,將人的苦難扛在自己的肩頭上.這個故事裡的「獨生子」的死亡,預示耶穌的死.對神來說,憐憫一個人,就不只是同情他,而是為他付出自己的一切.神的憐憫,是以救贖的行動,對人類的苦難大聲說「不」.憐憫的神,進入人間,與絕望的我們相遇,要解救我們.

(3)不要哭

路7:13主……對她說:「不要哭!」

在絕望的處境裡,叫人「不要哭」,是講風涼話嗎!能夠不哭,就不是真正的絕望.很多時,叫人「不要哭」,或者說「節哀順變」,只是些廉價的安慰,不單無法安慰人,更顯出你根本不明白別人的哀傷.在這裡,耶穌對這寡婦說:「不要哭」,這不是作些廉價的安慰,而是讓這寡婦體會神的介入.她將要面對的,不是苦難,而是闖入她生命的神.在這苦難人間,我們有千百個哭泣的理由.但在神的恩典下,我們有千百個不去哭泣的理由.生命充滿悲劇,要哭,總有理由.要不哭,就必須有一位克服悲劇的神.這位神,與這寡婦相遇.

(4)耶穌出手了,發言了

路7:14於是進前按著槓,抬的人就站住了.耶穌說:「少年人,我吩咐你,起來!」

耶穌主動進前,用手按著棺材架,抬的人就站住了.棺材架是不潔的,帶有死亡的氣息,劃出死亡的界線.一般人不應觸碰這界線,何況是尊貴的拉比耶穌,觸碰的人會沾染不潔.耶穌,卻是刻意的,進前來,按著這棺材架,接觸這不應觸碰的界線.真正的神聖,不是遠離這些所謂不潔的領域,而是觸碰它,瓦解它背後導致絕望的根源.耶穌正面攔住死亡行列的前進,耶穌就在死亡界線前面,攔住它,阻止它前進.耶穌闖入這死亡的界限,要在絕望的深淵裡,燃起生命的火.

然後,耶穌發言了,發出神聖主權的命令.這是一道沒有任何東西能違抗的命令.耶穌說:「少年人,我吩咐你,起來!」耶穌不需要與死神討價還價,耶穌不需要與人討論死人復生在科學上是否可能,耶穌不需要考慮這是否合乎人類理性的框架,耶穌直接發出命令.「少年人,我吩咐你,起來!」

這是創造主在創造萬有時發出的命令.祂說「要有光」,便有了光.祂說有,就有.命立,就立.祂從虛無創造出萬有,祂命令光明從黑暗裡照出來,祂從死亡中呼喚生命.神的說話帶著權能,神的命令一出,事就這樣成了.耶穌以神的命令,創造出一個新的世界.神的話語一出,對那寡婦來說,舊事已過,都變成新的了.

路7:15那死人就坐起,並且說話.耶穌便把他交給他母親.

就算死人,也得聽從神的命令.那死人就坐起,並且說話.原來,死亡並沒有話事權,我們全都高估了它.在神的話語面前,死亡也得俯首聽命.生命雖然有限,但死亡並沒有終極的權柄,它也要降服於神.耶穌讓我們看到,在死亡的界限以外,有更廣闊的生命世界,有神的話語創造的明天,有神開啟的希望.

耶穌將那復生的人,交給他母親.不知這母親如何想.這顛覆了她的一切想法.明明無可能,怎會變成有可能!明明是死亡,黑暗,絕望,怎麼可能變成生命,光明,希望.她一生的生活邏輯,是死去的,就是死去了,無法挽回.殊不知,在她的邏輯外,有一永活的神,祂闖入人間,改變了她的一生.

(5)恩典

由頭至尾,這寡婦都沒有出過一句聲.她沒有求過什麼,全部事情都是由耶穌主動引發的,這就是神白白的恩典.

各位弟兄姊妹,無論你在什麼光景,或憂慮,或絕望,施恩的神都會介入你的生命,祂看見你,憐憫你,祂正面對抗虛無的力量,以權能的命令創造光明.這一切都是神主動送贈的禮物,都是恩典,這就是福音.

恩典在,福音在,神在,故此,弟兄姊妹,請「不要哭」.

願神賜福你們.

願榮耀歸予聖父,聖子,聖靈.



\newpage

\section{我們今天一同得這能力吧}
\label{sec:798}
\textbf{高國雄牧師 - 我們今天一同得這能力吧!}
\newline
\newline
\hyperref[sec:797]{< < < PREV SERMON < < <}
~
\hyperref[sec:toc]{[返目錄]}
~
\hyperref[sec:799]{> > > NEXT SERMON > > >}


使徒行傳 2:1-21
\newline
\begin{longtable}{cl}
\hline
\hline
章節 & 經文 (和合本修訂版)\\
\hline
2:1 & \begin{tabularx}{0.7\textwidth}{X} 五旬節那日到了，他們全都聚集在一起。 \end{tabularx} \\ \\ \relax
2:2 & \begin{tabularx}{0.7\textwidth}{X} 忽然，有響聲從天上下來，好像一陣大風吹過，充滿了他們所坐的整座屋子； \end{tabularx} \\ \\ \relax
2:3 & \begin{tabularx}{0.7\textwidth}{X} 又有舌頭如火焰向他們顯現，分開落在他們每個人身上。 \end{tabularx} \\ \\ \relax
2:4 & \begin{tabularx}{0.7\textwidth}{X} 他們都被聖靈充滿，就按著聖靈所賜的口才說起別國的話來。 \end{tabularx} \\ \\ \relax
2:5 & \begin{tabularx}{0.7\textwidth}{X} 那時，有從天下各國來的虔誠的猶太人，住在耶路撒冷。 \end{tabularx} \\ \\ \relax
2:6 & \begin{tabularx}{0.7\textwidth}{X} 這聲音一響，許多人都來聚集，各人因為聽見門徒用他們各自的鄉談說話，就甚納悶， \end{tabularx} \\ \\ \relax
2:7 & \begin{tabularx}{0.7\textwidth}{X} 都詫異驚奇說：「看哪，這些說話的不都是加利利人嗎？ \end{tabularx} \\ \\ \relax
2:8 & \begin{tabularx}{0.7\textwidth}{X} 我們每個人怎麼聽見他們說我們生來所用的鄉談呢？ \end{tabularx} \\ \\ \relax
2:9 & \begin{tabularx}{0.7\textwidth}{X} 我們帕提亞人、瑪代人、以攔人，和住在美索不達米亞、猶太、加帕多家、本都、亞細亞、 \end{tabularx} \\ \\ \relax
2:10 & \begin{tabularx}{0.7\textwidth}{X} 弗呂家、旁非利亞、埃及的人，並靠近古利奈的利比亞一帶地方的人，僑居的羅馬人， \end{tabularx} \\ \\ \relax
2:11 & \begin{tabularx}{0.7\textwidth}{X} 包括猶太人和皈依猶太教的人， 克里特人和阿拉伯人，都聽見他們用我們的鄉談講論神的大作為。」 \end{tabularx} \\ \\ \relax
2:12 & \begin{tabularx}{0.7\textwidth}{X} 眾人就都驚奇困惑，彼此說：「這是甚麼意思呢？」 \end{tabularx} \\ \\ \relax
2:13 & \begin{tabularx}{0.7\textwidth}{X} 還有人譏誚，說：「他們是灌滿了新酒吧！」 \end{tabularx} \\ \\ \relax
2:14 & \begin{tabularx}{0.7\textwidth}{X} 彼得和十一個使徒站起來，他就高聲向眾人說：「猶太人和所有住在耶路撒冷的人哪，這件事你們要知道，要側耳聽我的話。 \end{tabularx} \\ \\ \relax
2:15 & \begin{tabularx}{0.7\textwidth}{X} 這些人並不像你們所想的喝醉了，因為現在才早晨九點鐘。 \end{tabularx} \\ \\ \relax
2:16 & \begin{tabularx}{0.7\textwidth}{X} 這正是藉著先知約珥所說的： \end{tabularx} \\ \\ \relax
2:17 & \begin{tabularx}{0.7\textwidth}{X} 『神說：在末後的日子，我要將我的靈澆灌凡血肉之軀的。你們的兒女要說預言；你們的少年要見異象；你們的老人要做異夢。 \end{tabularx} \\ \\ \relax
2:18 & \begin{tabularx}{0.7\textwidth}{X} 在那些日子，我要把我的靈澆灌，甚至給我的僕人和婢女，他們要說預言。 \end{tabularx} \\ \\ \relax
2:19 & \begin{tabularx}{0.7\textwidth}{X} 在天上，我要顯出奇事，在地下，我要顯出神蹟，有血，有火，有煙霧。 \end{tabularx} \\ \\ \relax
2:20 & \begin{tabularx}{0.7\textwidth}{X} 太陽要變為黑暗，月亮要變為血，這都在主大而光榮的日子未到以前。 \end{tabularx} \\ \\ \relax
2:21 & \begin{tabularx}{0.7\textwidth}{X} 那時，凡求告主名的都必得救。』 \end{tabularx} \\ \\
[1ex]
\hline
\hline
\end{longtable}
詩篇 104:24-34
\newline
\begin{longtable}{cl}
\hline
\hline
章節 & 經文 (和合本修訂版)\\
\hline
104:24 & \begin{tabularx}{0.7\textwidth}{X} 耶和華啊，你所造的何其多！都是你用智慧造成的，全地遍滿了你所造之物。 \end{tabularx} \\ \\ \relax
104:25 & \begin{tabularx}{0.7\textwidth}{X} 那裡有海，又大又廣，其中有無數的動物，大小活物都有。 \end{tabularx} \\ \\ \relax
104:26 & \begin{tabularx}{0.7\textwidth}{X} 那裡有船行走，有你所造的力威亞探悠游在其中。 \end{tabularx} \\ \\ \relax
104:27 & \begin{tabularx}{0.7\textwidth}{X} 這些都仰望你按時給牠們食物。 \end{tabularx} \\ \\ \relax
104:28 & \begin{tabularx}{0.7\textwidth}{X} 你給牠們，牠們就拾起來；你張手，牠們就飽得美食。 \end{tabularx} \\ \\ \relax
104:29 & \begin{tabularx}{0.7\textwidth}{X} 你轉臉，牠們就驚惶；你收回牠們的氣，牠們就死亡，歸於塵土。 \end{tabularx} \\ \\ \relax
104:30 & \begin{tabularx}{0.7\textwidth}{X} 你差遣你的靈，牠們就受造；你使地面更換為新。 \end{tabularx} \\ \\ \relax
104:31 & \begin{tabularx}{0.7\textwidth}{X} 願耶和華的榮耀存到永遠！願耶和華喜愛自己所造的！ \end{tabularx} \\ \\ \relax
104:32 & \begin{tabularx}{0.7\textwidth}{X} 他看地，地便震動；他摸山，山就冒煙。 \end{tabularx} \\ \\ \relax
104:33 & \begin{tabularx}{0.7\textwidth}{X} 我一生要向耶和華唱詩！我還活的時候，要向我的神歌頌！ \end{tabularx} \\ \\ \relax
104:34 & \begin{tabularx}{0.7\textwidth}{X} 願他悅納我的默念！我要因耶和華歡喜！ \end{tabularx} \\ \\
[1ex]
\hline
\hline
\end{longtable}
哥林多前書 12:3-13
\newline
\begin{longtable}{cl}
\hline
\hline
章節 & 經文 (和合本修訂版)\\
\hline
12:3 & \begin{tabularx}{0.7\textwidth}{X} 所以，我要你們知道，被神的靈感動的，沒有人會說「耶穌該受詛咒」；若不是被聖靈感動的，也沒有人能說「耶穌是主」。 \end{tabularx} \\ \\ \relax
12:4 & \begin{tabularx}{0.7\textwidth}{X} 恩賜有許多種，卻是同一位聖靈所賜。 \end{tabularx} \\ \\ \relax
12:5 & \begin{tabularx}{0.7\textwidth}{X} 事奉有許多種，卻是事奉同一位主。 \end{tabularx} \\ \\ \relax
12:6 & \begin{tabularx}{0.7\textwidth}{X} 工作有許多種，卻是同一位神在萬人中運行萬事。 \end{tabularx} \\ \\ \relax
12:7 & \begin{tabularx}{0.7\textwidth}{X} 聖靈彰顯在各人身上，是要使人得益處。 \end{tabularx} \\ \\ \relax
12:8 & \begin{tabularx}{0.7\textwidth}{X} 有人藉著聖靈領受智慧的言語；有人也靠著同一位聖靈領受知識的言語； \end{tabularx} \\ \\ \relax
12:9 & \begin{tabularx}{0.7\textwidth}{X} 又有人由同一位聖靈領受信心；還有人由同一位聖靈領受醫病的恩賜； \end{tabularx} \\ \\ \relax
12:10 & \begin{tabularx}{0.7\textwidth}{X} 又有人能行異能，又有人能作先知，又有人能辨別諸靈，又有人能說方言，又有人能翻方言。 \end{tabularx} \\ \\ \relax
12:11 & \begin{tabularx}{0.7\textwidth}{X} 這一切都是由惟一的、同一位聖靈所運行，隨著自己的旨意分給各人的。 \end{tabularx} \\ \\ \relax
12:12 & \begin{tabularx}{0.7\textwidth}{X} 就如身體是一個，卻有許多肢體，身體的肢體雖多，仍是一個身體；基督也是這樣。 \end{tabularx} \\ \\ \relax
12:13 & \begin{tabularx}{0.7\textwidth}{X} 我們無論是猶太人是希臘人，是為奴的是自主的，都從一位聖靈受洗成了一個身體，並且共享這位聖靈。 \end{tabularx} \\ \\
[1ex]
\hline
\hline
\end{longtable}
約翰福音 20:19-23
\newline
\begin{longtable}{cl}
\hline
\hline
章節 & 經文 (和合本修訂版)\\
\hline
20:19 & \begin{tabularx}{0.7\textwidth}{X} 那日（就是七日的第一日）晚上，門徒因怕猶太人，所在的地方門都關了。耶穌來，站在當中，對他們說：「願你們平安！」 \end{tabularx} \\ \\ \relax
20:20 & \begin{tabularx}{0.7\textwidth}{X} 說了這話，他把手和肋旁給他們看。門徒一看見主就喜樂了。 \end{tabularx} \\ \\ \relax
20:21 & \begin{tabularx}{0.7\textwidth}{X} 於是耶穌又對他們說：「願你們平安！父怎樣差遣了我，我也照樣差遣你們。」 \end{tabularx} \\ \\ \relax
20:22 & \begin{tabularx}{0.7\textwidth}{X} 說了這話，他向他們吹一口氣，說：「領受聖靈吧！ \end{tabularx} \\ \\ \relax
20:23 & \begin{tabularx}{0.7\textwidth}{X} 你們赦免誰的罪，誰的罪就得赦免；你們不赦免誰的罪，誰的罪就不得赦免。」 \end{tabularx} \\ \\
[1ex]
\hline
\hline
\end{longtable}
講題:我們今天一同得這能力吧!  Let’s Receive the Power of Holy Spirit Today!

經課:使徒行傳2:1-21;詩篇104:24-34;哥林多前書12:3-13;約翰福音20:19-23

講員:高國雄牧師

場合:香港中文大學崇基學院禮拜堂主日崇拜

日期:2026年5月24日

引言:我們渴求能力嗎?

各位弟兄姊妹,平安.

在今天的社會,每個人都在追求「能力」.學生追求學習的能力,職場人士追求競爭的能力,國家追求經濟與軍事的能力.然而,當我們環顧世界——面對戰爭,氣候變遷,或是我們個人生命中的焦慮,破碎與無力感時,我們不禁要問:我們真正需要的是什麼樣的能力?

今天是聖靈降臨節(Pentecost).這不只是一個紀念歷史事件的宗教節日,這是一個關於「Empower 充權」的慶典.今天的主題是:「我們今天一同得這能力吧!」 讓我們透過今天的經課,看聖靈如何將一群躲在關上門的屋子裡的門徒,轉化為翻天覆地改變世界的群體.

能力的來源:從懼怕到平安(約 20:19-23)

讓我們先回到最原始的起點.在約翰福音中,復活後的耶穌出現在門徒中間.當時門徒的狀態是什麼?

約翰福音20:19

19那日(就是七日的第一日)晚上,門徒所在的地方,因怕猶太人,門都關了.耶穌來,站在當中,對他們說:「願你們平安!」

恐懼: 他們因為怕猶太人,門都關上了.

無力: 他們失去了重心,前路茫茫.

這就是我們許多人的寫照.我們關上心門,保護自己不受外界傷害,卻也把自己困在恐懼中.耶穌對他們說的第一句話是:「願你們平安!」 隨後,祂做了一個非常關鍵的動作:

約翰福音20:22

22說了這話,就向他們吹一口氣,說:「你們受聖靈!

重點:聖靈的能力,不一定是外在的轟轟烈烈,而是由內在的生命更新開始.

希伯來文的「靈」與「氣息」是同一個字.就像上帝當初吹氣在亞當鼻孔裡使他成為有靈的活人,耶穌此時正對著門徒進行「新創造」.這份能力首先帶給我們的是超越環境的平安.

能力的彰顯:從沈默到宣告(徒 2:1-21)

五十天後,這份平安轉化成了巨大的動能.使徒行傳記載,聖靈如大風吹過,舌頭的火焰落在他們頭上.

突破語言的隔閡

使徒行傳2:1-4

1五旬節到了,門徒都聚集在一處. 2忽然,從天上有響聲下來,好像一陣大風吹過,充滿了他們所坐的屋子, 3又有舌頭如火焰顯現出來,分開落在他們各人頭上. 4他們就都被聖靈充滿,按著聖靈所賜的口才說起別國的話來.

聖靈降臨時,最明顯的神蹟是門徒開始說起別國的話.這代表了聖靈的能力具有「溝通與復和」的特質.巴別塔事件使人類語言不通,分散各地；聖靈降臨則是讓不同背景,不同膚色,不同方言的人,能聽見同一件大事——上帝的大作為.

凡有血氣的都要得著

使徒行傳2:17-21

17上帝說:在末後的日子,我要將我的靈澆灌凡有血氣的.你們的兒女要說預言；你們的少年人要見異象；老年人要做異夢.18在那些日子,我要將我的靈澆灌我的僕人和使女,他們就要說預言.19在天上,我要顯出奇事；在地下, 我要顯出神蹟；有血,有火,有煙霧.20日頭要變為黑暗,月亮要變為血；這都在主大而明顯的日子未到以前.21到那時候,凡求告主名的,就必得救.

彼得引用約珥書的預言:「我要將我的靈澆灌凡有血氣的.」

少年人要見異象(看見未來的希望).

老年人要做異夢(重燃對生命的熱情).

僕人和使女也要受靈(打破階級的枷鎖).

這能力不是給少數精英的,是給每一個渴慕上帝的人. 聖靈的能力不是讓我們自我膨脹,而是讓我們有勇氣跨出舒適圈,去對外宣告耶穌基督的救恩.

能力的運作:多元中的合一(林前 12:3b-13)

如果說使徒行傳讓我們看到聖靈的「爆發力」,那麼哥林多前書就讓我們看到聖靈的「持續力」.

哥林多前書12:4-7

4恩賜原有分別,聖靈卻是一位. 5職事也有分別,主卻是一位. 6功用也有分別,上帝卻是一位,在眾人裡面運行一切的事. 7聖靈顯在各人身上,是叫人得益處.

保羅提醒我們,恩賜有許多種,但聖靈卻是一位.

恩賜不同,職事不同,功用不同.

目的只有一個: 是叫人得益處(共同的福祉).

我們有時會誤以為,有能力就是「我能做別人不能做的事」.但聖靈的能力在教會中運作時,是讓我們看見自己的不足與他人的價值.「我們不拘是猶太人,是希臘人,是為奴的,是自主的,都從一位聖靈受洗,成了一個身體.」

這就是今天的主題「一同」得能力.能力不是為了競爭,而是為了補足.在聖靈裡,你的強項是為了服事我的弱項,而我的恩賜是為了祝福你的生命.

結語:在萬物中看見上帝的靈(詩 104)

詩篇 104 篇帶領我們將眼光放大到整個宇宙.

詩篇104:28b-29

你張手,牠們飽得美食.

29你掩面,牠們便驚惶；

你收回牠們的氣,牠們就死亡,歸於塵土.

上帝張手,萬物就得飽足；上帝收回氣息,萬物就歸回塵土.

聖靈不僅在教會中,聖靈更運行在全地.

親愛的弟兄姊妹,我們今天如何「一同得這能力」?

承認軟弱: 像門徒一樣,來到主面前,承認我們靠自己無法面對恐懼.

呼求氣息: 向主呼求:「主啊,請向我吹氣!讓聖靈更新我的生命.」

實踐合一: 在教會中,在家庭裡,放下成見,用聖靈給予的恩賜彼此建立.

「我們今天一同得這能力吧!」

這能力不是要我們去統治別人,而是要我們去愛；

這能力不是要我們避開苦難,而是給我們盼望；

這能力不是要我們獨自成功,而是要我們一同見證上帝的榮耀.

願聖靈的火,今天重新點燃我們每一個人的心.



\newpage

\section{放心你的信治好你了}
\label{sec:799}
\textbf{李均熊牧師 - 放心……你的信治好你了}
\newline
\newline
\hyperref[sec:798]{< < < PREV SERMON < < <}
~
\hyperref[sec:toc]{[返目錄]}
~
\hyperref[sec:800]{> > > NEXT SERMON > > >}


創世記 12:1-9
\newline
\begin{longtable}{cl}
\hline
\hline
章節 & 經文 (和合本修訂版)\\
\hline
12:1 & \begin{tabularx}{0.7\textwidth}{X} 耶和華對亞伯蘭說：「你要離開本地、本族、父家，往我所要指示你的地去。 \end{tabularx} \\ \\ \relax
12:2 & \begin{tabularx}{0.7\textwidth}{X} 我必使你成為大國，我必賜福給你，使你的名為大；你要使別人得福。 \end{tabularx} \\ \\ \relax
12:3 & \begin{tabularx}{0.7\textwidth}{X} 為你祝福的，我必賜福給他；詛咒你的，我必詛咒他。地上的萬族都必因你得福。」 \end{tabularx} \\ \\ \relax
12:4 & \begin{tabularx}{0.7\textwidth}{X} 亞伯蘭就遵照耶和華的吩咐去了；羅得也和他同去。亞伯蘭離開哈蘭的時候年七十五歲。 \end{tabularx} \\ \\ \relax
12:5 & \begin{tabularx}{0.7\textwidth}{X} 亞伯蘭帶著他妻子撒萊和姪兒羅得，以及他們在哈蘭積蓄的財物、獲得的人口，往迦南地去。他們就來到了迦南地。 \end{tabularx} \\ \\ \relax
12:6 & \begin{tabularx}{0.7\textwidth}{X} 亞伯蘭經過那地，直到示劍地方，摩利橡樹那裡；當時迦南人住在那地。 \end{tabularx} \\ \\ \relax
12:7 & \begin{tabularx}{0.7\textwidth}{X} 耶和華向亞伯蘭顯現，說：「我要把這地賜給你的後裔。」亞伯蘭就在那裡為向他顯現的耶和華築了一座壇。 \end{tabularx} \\ \\ \relax
12:8 & \begin{tabularx}{0.7\textwidth}{X} 從那裡他又遷到伯特利東邊的山，支搭帳棚；西邊是伯特利，東邊是艾。他在那裡又為耶和華築了一座壇，求告耶和華的名。 \end{tabularx} \\ \\ \relax
12:9 & \begin{tabularx}{0.7\textwidth}{X} 後來亞伯蘭漸漸遷往尼革夫去。 \end{tabularx} \\ \\
[1ex]
\hline
\hline
\end{longtable}
詩篇 33:1-12
\newline
\begin{longtable}{cl}
\hline
\hline
章節 & 經文 (和合本修訂版)\\
\hline
33:1 & \begin{tabularx}{0.7\textwidth}{X} 義人哪，你們當因耶和華歡呼，正直人理當讚美耶和華。 \end{tabularx} \\ \\ \relax
33:2 & \begin{tabularx}{0.7\textwidth}{X} 你們要彈琴稱謝耶和華，用十弦瑟歌頌他。 \end{tabularx} \\ \\ \relax
33:3 & \begin{tabularx}{0.7\textwidth}{X} 應當向他唱新歌，彈得巧妙，聲音洪亮。 \end{tabularx} \\ \\ \relax
33:4 & \begin{tabularx}{0.7\textwidth}{X} 因為耶和華的言語正直，他的作為盡都信實。 \end{tabularx} \\ \\ \relax
33:5 & \begin{tabularx}{0.7\textwidth}{X} 他喜愛公義和公平，遍地滿了耶和華的慈愛。 \end{tabularx} \\ \\ \relax
33:6 & \begin{tabularx}{0.7\textwidth}{X} 諸天藉耶和華的話而造，萬象藉他口中的氣而成。 \end{tabularx} \\ \\ \relax
33:7 & \begin{tabularx}{0.7\textwidth}{X} 他聚集海水如壘，收藏深洋在倉庫。 \end{tabularx} \\ \\ \relax
33:8 & \begin{tabularx}{0.7\textwidth}{X} 願全地都敬畏耶和華！願世上的居民都懼怕他！ \end{tabularx} \\ \\ \relax
33:9 & \begin{tabularx}{0.7\textwidth}{X} 因為他說有，就有，命立，就立。 \end{tabularx} \\ \\ \relax
33:10 & \begin{tabularx}{0.7\textwidth}{X} 耶和華使列國的籌算歸於無有，使萬民的計謀全無功效。 \end{tabularx} \\ \\ \relax
33:11 & \begin{tabularx}{0.7\textwidth}{X} 耶和華的籌算永遠立定，他心中的計劃萬代長存。 \end{tabularx} \\ \\ \relax
33:12 & \begin{tabularx}{0.7\textwidth}{X} 以耶和華為神的，那國有福了！耶和華揀選為自己產業的，那民有福了！ \end{tabularx} \\ \\
[1ex]
\hline
\hline
\end{longtable}
羅馬書 4:13-25
\newline
\begin{longtable}{cl}
\hline
\hline
章節 & 經文 (和合本修訂版)\\
\hline
4:13 & \begin{tabularx}{0.7\textwidth}{X} 因為神給亞伯拉罕和他後裔承受世界的應許不是藉著律法，而是藉著信而得的義。 \end{tabularx} \\ \\ \relax
4:14 & \begin{tabularx}{0.7\textwidth}{X} 若是屬於律法的人才是後嗣，信就落空了，應許也就失效了。 \end{tabularx} \\ \\ \relax
4:15 & \begin{tabularx}{0.7\textwidth}{X} 因為律法是惹動憤怒的，哪裡沒有律法，哪裡就沒有過犯。 \end{tabularx} \\ \\ \relax
4:16 & \begin{tabularx}{0.7\textwidth}{X} 所以，人作後嗣是出於信，因此就屬乎恩，以致應許保證歸給所有的後裔，不但歸給那屬於律法的，也歸給那效法亞伯拉罕之信的人。 \end{tabularx} \\ \\ \relax
4:17 & \begin{tabularx}{0.7\textwidth}{X} 亞伯拉罕所信的是那叫死人復活、使無變為有的神，在這位神面前亞伯拉罕成為我們眾人的父，如經上所記：「我已經立你作多國之父。」 \end{tabularx} \\ \\ \relax
4:18 & \begin{tabularx}{0.7\textwidth}{X} 他在沒有盼望的時候，仍存著盼望來相信，就得以作多國之父，正如先前所說：「你的後裔將要如此。」 \end{tabularx} \\ \\ \relax
4:19 & \begin{tabularx}{0.7\textwidth}{X} 他將近百歲的時候，雖然想到自己的身體如同已死，撒拉也不可能生育，他的信心還是不軟弱， \end{tabularx} \\ \\ \relax
4:20 & \begin{tabularx}{0.7\textwidth}{X} 仍仰望神的應許，總沒有因不信而起疑惑，反倒因信而剛強，將榮耀歸給神， \end{tabularx} \\ \\ \relax
4:21 & \begin{tabularx}{0.7\textwidth}{X} 且滿心相信神所應許的必能成就。 \end{tabularx} \\ \\ \relax
4:22 & \begin{tabularx}{0.7\textwidth}{X} 所以這也就算他為義。 \end{tabularx} \\ \\ \relax
4:23 & \begin{tabularx}{0.7\textwidth}{X} 「算他為義」這句話不是單為他寫的， \end{tabularx} \\ \\ \relax
4:24 & \begin{tabularx}{0.7\textwidth}{X} 也是為我們將來得算為義的人寫的，就是為我們這些信神使我們的主耶穌從死人中復活的人寫的。 \end{tabularx} \\ \\ \relax
4:25 & \begin{tabularx}{0.7\textwidth}{X} 耶穌被出賣，是為我們的過犯；他復活，是為使我們稱義。 \end{tabularx} \\ \\
[1ex]
\hline
\hline
\end{longtable}
馬太福音 9:9-13,18-26
\newline
\begin{longtable}{cl}
\hline
\hline
章節 & 經文 (和合本修訂版)\\
\hline
9:9 & \begin{tabularx}{0.7\textwidth}{X} 耶穌從那裡往前走，看見一個人名叫馬太，在稅關坐著，就對他說：「來跟從我！」他就起來跟從耶穌。 \end{tabularx} \\ \\ \relax
9:10 & \begin{tabularx}{0.7\textwidth}{X} 耶穌在屋裡坐席的時候，有好些稅吏和罪人來，與耶穌和他的門徒一同坐席。 \end{tabularx} \\ \\ \relax
9:11 & \begin{tabularx}{0.7\textwidth}{X} 法利賽人看見，就對耶穌的門徒說：「你們的老師為甚麼與稅吏和罪人一同吃飯呢？」 \end{tabularx} \\ \\ \relax
9:12 & \begin{tabularx}{0.7\textwidth}{X} 耶穌聽見，就說：「健康的人用不著醫生；有病的人才用得著。 \end{tabularx} \\ \\ \relax
9:13 & \begin{tabularx}{0.7\textwidth}{X} 經上說：『我喜愛憐憫，不喜愛祭祀。』這句話的意思，你們去揣摩。我不是來召義人，而是召罪人。」 \end{tabularx} \\ \\ \relax
9:14 & \begin{tabularx}{0.7\textwidth}{X} 那時，約翰的門徒來見耶穌，說：「我們和法利賽人常常禁食，你的門徒卻不禁食，這是為甚麼呢？」 \end{tabularx} \\ \\ \relax
9:15 & \begin{tabularx}{0.7\textwidth}{X} 耶穌對他們說：「新郎和賓客在一起的時候，賓客怎麼能哀慟呢？但日子將到，新郎要被帶走，那時候他們就要禁食了。 \end{tabularx} \\ \\ \relax
9:16 & \begin{tabularx}{0.7\textwidth}{X} 沒有人把新布補在舊衣服上；因為所補上的會撕破那衣服，裂口就更大了。 \end{tabularx} \\ \\ \relax
9:17 & \begin{tabularx}{0.7\textwidth}{X} 也沒有人把新酒裝在舊皮袋裡，若是這樣，皮袋會脹破，酒就漏出來，皮袋也糟蹋了。相反地，把新酒裝在新皮袋裡，兩樣就都保全了。」 \end{tabularx} \\ \\ \relax
9:18 & \begin{tabularx}{0.7\textwidth}{X} 耶穌說這些話的時候，有一個會堂主管來，向他下跪，說：「我女兒剛死了，求你去按手在她身上，她就會活過來。」 \end{tabularx} \\ \\ \relax
9:19 & \begin{tabularx}{0.7\textwidth}{X} 耶穌就起來跟他去；門徒也跟了去。 \end{tabularx} \\ \\ \relax
9:20 & \begin{tabularx}{0.7\textwidth}{X} 這時，有一個女人，患了經血不止的病有十二年，來到耶穌背後，摸他的衣裳繸子； \end{tabularx} \\ \\ \relax
9:21 & \begin{tabularx}{0.7\textwidth}{X} 因為她心裡說：「我只要摸他的衣裳，就會痊癒。」 \end{tabularx} \\ \\ \relax
9:22 & \begin{tabularx}{0.7\textwidth}{X} 耶穌轉過來，看見她，就說：「女兒，放心！你的信救了你。」從那時起，這女人就痊癒了。 \end{tabularx} \\ \\ \relax
9:23 & \begin{tabularx}{0.7\textwidth}{X} 耶穌到了會堂主管的家裡，看見吹鼓手和亂哄哄的一群人， \end{tabularx} \\ \\ \relax
9:24 & \begin{tabularx}{0.7\textwidth}{X} 就說：「退去吧！這女孩不是死了，而是睡著了。」他們就嘲笑他。 \end{tabularx} \\ \\ \relax
9:25 & \begin{tabularx}{0.7\textwidth}{X} 眾人被趕出後，耶穌就進去，拉著她的手，女孩就起來了。 \end{tabularx} \\ \\ \relax
9:26 & \begin{tabularx}{0.7\textwidth}{X} 於是這消息傳遍了那地方。 \end{tabularx} \\ \\
[1ex]
\hline
\hline
\end{longtable}
講題:放心……你的信治好你了 Take Heart...Your Faith Has Made You Well

經課:創世記12:1-9；詩篇33:1-12；羅馬書4:13-25；馬太福音9:9-13,18-26

講員:李均熊牧師

場合:香港中文大學崇基學院禮拜堂主日崇拜

日期:2026年6月7日

起:為甚麼婦人會偷偷摸摸呢?

為甚麼婦人會偷偷摸摸接近耶穌?這段經文我們都很熟,所以我們很容易很快就會說,這是一個有信心的故事.不過,每次再讀的時候,就會發現,她的行動本身其實十分不解.她沒有走到耶穌面前,她沒有開口求,她沒有讓人注意到自己.她只是在一群人之中,從後面靠近,然後輕輕摸一下衣角,好像事情完成之後,她就可以離開.這個動作本身,就已經令我們要問:為甚麼她要這樣做?

傳統的理解會說,因為血漏在律法中屬於不潔,所以她不能公開接觸別人.這個背景當然重要,因為它幫助我們看見她在群體中所處的位置.不過,如果我們只停在這裡,可能還未完全讀到這段故事的張力.因為她不只是被規條限制,她也可能已經習慣了一種方式,就是不太期望自己被看見,也不太相信自己可以被正面地接納.她不是完全拒絕幫助,但她好像只容許自己用一種不驚動任何人的方式去得到幫助.

如果放在今天的處境…… 【宿生失眠的例子】

不過我要提醒一下大家,當我們嘗試用今日的處境去理解經文時,很容易不知不覺將這段故事讀成,它只是關於某一類特殊的人,或者某些極端的生命處境.如果這樣讀,這段經文就會慢慢離我們很遠,變成只是關於「他們」,而不是讓我們看見一種更普遍的人性(現代人?後covid?)經驗.也許,這段經文不是要我們找出誰最像這個婦人,而是意識到,原來人在很多不同的處境裡,都可能有一種相似的傾向,就是願意接近幫助,但同時又避免被看見.

所以,當我們再回到這位婦人,我們看到的就不只是她有沒有信心,也包括她怎樣接近耶穌.她想改變自己的處境,但不一定想被帶出來.她願意伸手,但未必準備好被認出.她的動作很細小,但正正是這個細小的動作,反而將她的心態放大了.

承:她想要的是醫治,還是只想不被看見地得到幫助?

而如果將這段經文放回整個敘事,就會發現它其實夾在另一個故事之中,就是會堂主管(其他福音書稱他為睚魯)的故事；一個人公開走到耶穌面前,求祂幫助,甚至邀請祂進入自己的家,於是兩個畫面同時存在,顯出特別的對比:一個是公開呼求,一個是隱藏接近,一個是邀請耶穌進入,一個是只想碰一下就離開……這種對比並沒有直接說明,但讀者聽眾也應感受到當中的張力.

回到那沒名字的婦人,經文說她心裡想:「我只要摸他的衣裳,就必痊癒.她的確有信心,因為她相信耶穌有能力,甚至她相信只要摸一摸耶穌的衣裳,就能得著這種能力.這樣說當然沒有錯.但她的信心卻只是停在一個很特別的位置.她相信耶穌可以幫助她,但她選擇一種不需要進入關係的方式去接近祂.她可以「觸摸」耶穌,但不一定願意讓耶穌「觸摸」她,看見她.

整個敘事的結構本身也有一種「夾敘」的效果,婦人的故事被放在會堂主管的故事中間,一邊是女兒即將死亡甚至已經死去的敘事,另一邊是一個流血十二年的身體,一種是急性的死亡,一種是慢性的流失.

於是問題好像不只是誰比較有信心,而是兩種不同的接近耶穌的方式,她不是隨便碰碰運氣,而是帶著一種很強的期待來到耶穌面前/後面?.可是,她雖然期待改變,卻仍然停留在一個很安全的位置.她不開口,她不對話,她不需要任何人知道,她也不需要耶穌在這一刻對她說甚麼.她只想伸手,只想得到,只想離開.這就令我們開始思考,她所追求的「痊癒」,究竟是甚麼.而真正的痊癒又是甚麼.

這樣的狀態,其實在今天香港信徒身上,也不是完全陌生.我們很熟悉信仰的語言,也知道甚麼時候應該禱告,知道神可以幫助,也知道在難處裡要來到神面前.但與此同時,我們未必習慣讓信仰真正進入那些最深,最難講,最怕被人知道的地方.我們可以維持一種正常運作的生活,也可以在教會裡參與,事奉,出現,但這並不等於我們真的願意被看見.有些人不是完全離開信仰,而是慢慢退到一個較遠的位置,仍然保留某種接觸,但不再願意讓自己太深地進入(後covid?離堂?網上崇拜?).於是,信仰有時便慢慢變成一種可以在需要時使用的資源,而不再是一段要投入,要被觸碰,甚至要被改變的關係.

如果這樣看,那麼這位婦人的故事就不只是古代的一個醫治故事.它也讓我們看見一種很熟悉的信仰姿態,就是人可以一邊相信耶穌有能力,一邊又避免讓自己真正暴露在祂面前.這種姿態並不等於完全沒有信心,但它也顯示,人的信心可以仍然帶著距離.

轉:耶穌不讓她匿名離開

耶穌沒有讓事情就這樣過去.祂轉過來,看見那個本來想離開的婦人,然後對她說:「女兒,放心!你的信救了你.」這一句話表面上看來,好像只是為醫治畫上一個句號,但如果我們留意當中的用字,就會發現耶穌所說的,其實比一個醫療結果更深.

「救了你」這裡的原文,是希臘文 σῴζω[sozo]就是婦人心裡想「我只要摸他的衣裳,就必痊癒」的痊癒是同一個字根的.這個字既可以指醫治,也可以指拯救,是帶有salvation 的神學意思,所以在一些英文譯本中,這句話會被翻譯為:Take heart, your faith has made you well.這裡的 made you well,究竟只是指身體恢復正常,還是指一個人被帶進一種更深,更完整的狀態?

如果再看耶穌在故事中的行動,這一點就更加明顯.耶穌不只是確認她已經得醫治,然後容許她離開.祂停下來,祂看見她,祂對她說話,祂稱她為「女兒」.換句話說,祂不容許她繼續停留在不被辨認的位置.原本她只是想伸手,然後退後；原本她只是想得到幫助,而不需要進入關係；但在這一刻,耶穌把她帶到另一個位置,讓她不只是「拿了就走」,而是被看見,被呼喚,被重新命名.

如果這樣理解,這裡的 σῴζω 就不只是「好了」那麼簡單.這裡發生的,不只是身體的改變,也包括一個人從原本孤立,隱藏,不能被看見的位置,被帶回一段關係之中.她原本認為的痊癒,只求一個結果,但耶穌給她的,似乎是一種更深的痊癒.祂不只是讓她身體恢復正常,而是讓她重新以一個可以被稱呼,可以被接納的身份站在祂面前.

合:不只是信心,而是信心停在哪裡

這段經文留下來的,好像就不只是一個關於信心大小的問題,而是一個關於信心位置的問題.我們有沒有信心,當然重要；但同樣重要的是,我們的信心停在哪裡.是在一個只想接觸,卻不想進入的地方,還是在一個願意被看見,願意被帶進關係的地方.

很多時候,我們都很習慣在需要的時候來到神面前,也可能真的經歷過祂的幫助.不過同時,我們也可能保留一些距離.我們不太想被看見,不太想讓某些部分被觸碰.我們可以有宗教行動,可以有屬靈語言,也可以有某種與神的接觸,但這一切未必已經成為一段真正讓自己進入的關係.

到最後我們要問的,不只是我們有沒有來到耶穌面前,而是當耶穌轉過來,看見我們的時候,我們是否仍然想停在原來的位置.又或者,我們是否願意讓祂對我們說話,讓祂重新定義我們,讓祂帶我們進入那種不只是被醫治,而是真正的痊癒.



\newpage

\section{在主前蒙恩}
\label{sec:800}
\textbf{高國雄牧師 - 在主前蒙恩}
\newline
\newline
\hyperref[sec:799]{< < < PREV SERMON < < <}
~
\hyperref[sec:toc]{[返目錄]}
~
\hyperref[sec:801]{> > > NEXT SERMON > > >}


創世記 18:1-15
\newline
\begin{longtable}{cl}
\hline
\hline
章節 & 經文 (和合本修訂版)\\
\hline
18:1 & \begin{tabularx}{0.7\textwidth}{X} 耶和華在幔利橡樹那裡向亞伯拉罕顯現。天正熱的時候，亞伯拉罕坐在帳棚門口。 \end{tabularx} \\ \\ \relax
18:2 & \begin{tabularx}{0.7\textwidth}{X} 他舉目觀看，看哪，有三個人站在他附近。他一看見，就從帳棚門口跑去迎接他們，俯伏在地， \end{tabularx} \\ \\ \relax
18:3 & \begin{tabularx}{0.7\textwidth}{X} 說：「我主，我若在你眼前蒙恩，請不要離開你的僕人走過去。 \end{tabularx} \\ \\ \relax
18:4 & \begin{tabularx}{0.7\textwidth}{X} 容我拿點水來，請你們洗腳，在樹下休息。 \end{tabularx} \\ \\ \relax
18:5 & \begin{tabularx}{0.7\textwidth}{X} 既然你們來到僕人這裡了，我再拿點餅來，讓你們恢復心力，然後再走。」他們說：「就照你說的去做吧。」 \end{tabularx} \\ \\ \relax
18:6 & \begin{tabularx}{0.7\textwidth}{X} 亞伯拉罕急忙進帳棚到撒拉那裡，說：「你趕快拿三細亞細麵，揉麵做餅。」 \end{tabularx} \\ \\ \relax
18:7 & \begin{tabularx}{0.7\textwidth}{X} 亞伯拉罕又跑到牛群裡，牽了一頭又嫩又好的牛犢來，交給僕人，僕人就急忙去預備。 \end{tabularx} \\ \\ \relax
18:8 & \begin{tabularx}{0.7\textwidth}{X} 亞伯拉罕取了乳酪和奶，以及預備好了的牛犢來，擺在他們面前，自己在樹下站在旁邊，他們就吃了。 \end{tabularx} \\ \\ \relax
18:9 & \begin{tabularx}{0.7\textwidth}{X} 他們對亞伯拉罕說：「你妻子撒拉在哪裡？」他說：「看哪，在帳棚裡。」 \end{tabularx} \\ \\ \relax
18:10 & \begin{tabularx}{0.7\textwidth}{X} 有一位說：「明年這時候 ，我一定會回到你這裡。看哪，你的妻子撒拉會生一個兒子。」撒拉在那人後面的帳棚門口也聽見了。 \end{tabularx} \\ \\ \relax
18:11 & \begin{tabularx}{0.7\textwidth}{X} 亞伯拉罕和撒拉都年紀老邁，撒拉的月經已停了。 \end{tabularx} \\ \\ \relax
18:12 & \begin{tabularx}{0.7\textwidth}{X} 撒拉心裡竊笑，說：「我已衰老，我的主也老了，怎能有這喜事呢？」 \end{tabularx} \\ \\ \relax
18:13 & \begin{tabularx}{0.7\textwidth}{X} 耶和華對亞伯拉罕說：「撒拉為甚麼竊笑，說：『我已年老，果真能生育嗎？』 \end{tabularx} \\ \\ \relax
18:14 & \begin{tabularx}{0.7\textwidth}{X} 耶和華豈有難成的事嗎？到了所定的時候，我必回到你這裡。明年這時候，撒拉會生一個兒子。」 \end{tabularx} \\ \\ \relax
18:15 & \begin{tabularx}{0.7\textwidth}{X} 撒拉因為害怕，就不承認，說：「我沒有笑。」那人說：「不，你的確笑了。」 \end{tabularx} \\ \\
[1ex]
\hline
\hline
\end{longtable}
羅馬書 5:1-8
\newline
\begin{longtable}{cl}
\hline
\hline
章節 & 經文 (和合本修訂版)\\
\hline
5:1 & \begin{tabularx}{0.7\textwidth}{X} 所以，我們既因信稱義，就藉著我們的主耶穌基督得以與神和好。 \end{tabularx} \\ \\ \relax
5:2 & \begin{tabularx}{0.7\textwidth}{X} 我們又藉著他，因信得以進入現在所站立的這恩典中，並且歡歡喜喜盼望神的榮耀。 \end{tabularx} \\ \\ \relax
5:3 & \begin{tabularx}{0.7\textwidth}{X} 不但如此，就是在患難中也是歡歡喜喜的，因為知道患難生忍耐， \end{tabularx} \\ \\ \relax
5:4 & \begin{tabularx}{0.7\textwidth}{X} 忍耐生老練，老練生盼望， \end{tabularx} \\ \\ \relax
5:5 & \begin{tabularx}{0.7\textwidth}{X} 盼望不至於落空，因為神的愛，已藉著所賜給我們的聖靈，澆灌在我們心裡。 \end{tabularx} \\ \\ \relax
5:6 & \begin{tabularx}{0.7\textwidth}{X} 我們還軟弱的時候，基督就在特定的時刻為不敬虔之人死。 \end{tabularx} \\ \\ \relax
5:7 & \begin{tabularx}{0.7\textwidth}{X} 為義人死，是少有的；為仁人死，或者有敢做的。 \end{tabularx} \\ \\ \relax
5:8 & \begin{tabularx}{0.7\textwidth}{X} 惟有基督在我們還作罪人的時候為我們死，神的愛就在此向我們顯明了。 \end{tabularx} \\ \\
[1ex]
\hline
\hline
\end{longtable}
馬太福音 9:35-10:8
\newline
\begin{longtable}{cl}
\hline
\hline
章節 & 經文 (和合本修訂版)\\
\hline
9:35 & \begin{tabularx}{0.7\textwidth}{X} 耶穌走遍各城各鄉，在他們的會堂裡教導人，宣講天國的福音，又醫治各樣的病症。 \end{tabularx} \\ \\ \relax
9:36 & \begin{tabularx}{0.7\textwidth}{X} 他看見一大群人，就憐憫他們；因為他們困苦無助，如同羊沒有牧人一樣。 \end{tabularx} \\ \\ \relax
9:37 & \begin{tabularx}{0.7\textwidth}{X} 於是他對門徒說：「要收的莊稼多，做工的人少。 \end{tabularx} \\ \\ \relax
9:38 & \begin{tabularx}{0.7\textwidth}{X} 所以，你們要求莊稼的主差遣做工的人出去收他的莊稼。」 \end{tabularx} \\ \\ \relax
10:1 & \begin{tabularx}{0.7\textwidth}{X} 耶穌叫了十二個門徒來，給他們權柄，能驅趕污靈和醫治各樣的疾病。 \end{tabularx} \\ \\ \relax
10:2 & \begin{tabularx}{0.7\textwidth}{X} 這十二使徒的名字如下：頭一個叫西門（又稱彼得），還有他弟弟安得烈，西庇太的兒子雅各和雅各的弟弟約翰， \end{tabularx} \\ \\ \relax
10:3 & \begin{tabularx}{0.7\textwidth}{X} 腓力和巴多羅買，多馬和稅吏馬太，亞勒腓的兒子雅各，和達太， \end{tabularx} \\ \\ \relax
10:4 & \begin{tabularx}{0.7\textwidth}{X} 激進黨的西門，還有出賣耶穌的加略人猶大。 \end{tabularx} \\ \\ \relax
10:5 & \begin{tabularx}{0.7\textwidth}{X} 耶穌差遣這十二個人出去，吩咐他們說：「外邦人的路，你們不要走；撒瑪利亞人的城，你們不要進； \end{tabularx} \\ \\ \relax
10:6 & \begin{tabularx}{0.7\textwidth}{X} 寧可往以色列家迷失的羊那裡去。 \end{tabularx} \\ \\ \relax
10:7 & \begin{tabularx}{0.7\textwidth}{X} 要邊走邊傳，說『天國近了』。 \end{tabularx} \\ \\ \relax
10:8 & \begin{tabularx}{0.7\textwidth}{X} 要醫治病人，使死人復活，使痲瘋病人潔淨，把鬼趕出去。你們白白地得來，也要白白地給人。 \end{tabularx} \\ \\
[1ex]
\hline
\hline
\end{longtable}
講題:在主前蒙恩 Blessed by the Lord

經課:創世記18章1至15節；詩篇116篇1至2節及12至19節；羅馬書5章1至8節；馬太福音9章35至10章8節

講員:高國雄牧師

場合:香港中文大學崇基學院禮拜堂主日崇拜

日期:2026年6月14日

「蒙恩」是我們基督徒非常熟悉,也常常掛在嘴邊的詞.我們看到別人生活順遂,家庭幸福,會說「你真蒙恩」；我們自己病得醫治,工作升職,也會感謝主說「我蒙了恩典」.

但是,聖經中所說的「在主前蒙恩」,僅僅是指我們日子過得順風順水嗎?當我們面對生命中的荒涼,身體的軟弱,甚至罪惡的捆綁時,我們還是一個「蒙恩」的人嗎?

今天,我們要透過四段看似截然不同,卻被聖靈奇妙串聯在一起的經文,從舊約到新約,從應許到踐行,一起來思想:一個在主前蒙恩的人,究竟擁有怎樣的生命確據?

蒙恩,是因上帝垂聽軟弱者的呼求

首先,在主前蒙恩,意味著我們的聲音被上帝聽見了.

詩篇 116:1-2

「我愛耶和華,因為他聽了我的聲音和我的懇求.他既向我側耳,我一生要求告他.」

詩篇的作者在這裡表達了一種極其親密的關係.他遭遇了什麼?如果我們往後看,會發現他當時被「死亡的繩索纏繞」,「遭受痛苦與愁悶」.在最絕望的時候,他只有呼求.

請大家注意中文翻譯用這個生動的動詞:「側耳」.

創造宇宙萬物的主宰,竟然為了我們這如塵土般微小的人,把耳朵貼近我們.就像一個母親,在深夜裡敏銳地捕捉到嬰兒最微弱的哭聲,立刻俯身過去一樣.

​恩典的起點: 不是我們有多偉大,而是上帝願意垂聽.

​我們的回應: 既然上帝向我們側耳,我們一生的回應就應該是「求告祂」.在主前蒙恩的人,是一個深知自己軟弱,卻隨時隨地敢向天父呼求的人.

​

蒙恩,是看見上帝在不可能中成就應許

​接著,我們把目光投向創世記中那位「蒙恩」的信心之父——亞伯拉罕.

​創世記 18:1-3, 10, 14

亞伯拉罕在幔利橡樹旁接待了三位神秘的客人.他對那領頭的說:「我主,我若在你眼前蒙恩,求你不要離開僕人往前去.」

耶和華說:「到明年這時候,我必要回到你這裡；你的妻子撒拉必生一個兒子.」……「耶和華豈有難成的事嗎?」

亞伯拉罕求在主眼前「蒙恩」,而上帝給他的恩典,完全打破了人的理性限制.當時亞伯拉罕和撒拉都已經年老髮白,撒拉的生育斷絕了.當撒拉在帳棚門口聽見這話時,她「心裡暗笑」——這是一種類似於「別開玩笑了,怎麼可能」的苦笑與無奈.

我們是不是也常常像撒拉一樣「暗笑」?

「我的壞習慣改不掉了,怎麼可能改變?」

「我和家人的關係凍結這麼多年了,怎麼可能和好?」

「世界這麼混亂,上帝真的還在掌權嗎?」

但上帝對撒拉的暗笑發出了靈魂之問:「耶和華豈有難成的事嗎?」

在主前蒙恩,不是因為我們「能」,而是因為上帝「能」.恩典,就是在我們已經絕望,毫無辦法的地方,上帝的應許依然能夠成就.

蒙恩,是在我們還作罪人時就被愛

如果說舊約的恩典是一幅預表的圖畫,那麼新約羅馬書則把「蒙恩」的核心,最清晰,最震撼地展現了出來.

羅馬書 5:6, 8

「因我們還軟弱的時候,基督就按時候為罪人死.……惟有基督在我們還作罪人的時候為我們死,神的愛就在此向我們顯明了.」

保羅在這裡用了三個詞來形容我們未蒙恩前的狀態:「軟弱」,「罪人」,「仇敵」.

世俗的邏輯是:你表現得好,我才愛你；你對我有價值,我才施恩給你.這叫「交換」,不叫恩典.

但神的恩典是:在我們還是一團糟,還在背叛祂,還在扮演祂的「仇敵」時,耶穌基督就已經為我們釘在了十字架上.

白白的恩典: 我們得以與神相通,有著所指望的榮耀,完全不是因為我們的功勞.

患難中的歡歡喜喜: 因為這恩典的根基是基督的死與復活,無比堅固,所以保羅說,我們甚至在患難中也是歡歡喜喜的.因為知道神的愛已經藉著聖靈澆灌在我們心裡.在主前蒙恩,是給了我們在苦難中依然大有盼望的底氣.

蒙恩,是白白得來,也要白白捨去

最後,蒙恩的人不能只把恩典據為己有.馬太福音向我們展示了,一個在主前蒙恩的群體,應該如何活在世界上.

馬太福音 9:35-36, 10:1, 8

耶穌走遍各城各鄉……看見許多的人,就憐憫他們,因為他們困苦流離,如同羊沒有牧人一般.

耶穌叫了十二個門徒來,給他們權柄……耶穌差遣這十二個人去,吩咐他們說:「……醫治病人,叫死人復活,叫長大痲瘋的乾淨,把鬼趕出去.你們白白地得來,也要白白地捨去.」

這段經文是主耶穌差遣門徒的宣教大委任.這一切的起點,是耶穌的「憐憫」.看見世人的困苦流離,耶穌動了慈心.而這十二個門徒,本身也是跟隨耶穌,白白得著醫治與生命翻轉的蒙恩者.

耶穌對他們,也是對今天的我們說了一句至關重要的話:

「你們白白地得來,也要白白地捨去.」

我們在主前蒙了赦免的恩,醫治的恩,安慰的恩,不是為了讓我們築起一道高牆,自己躲在裡面享受屬靈的安逸.

恩典的管道: 蒙恩的人,必須成為施恩的人.

教會的使命: 面對這個世代無數「困苦流離」的靈魂,我們不能冷眼旁觀.神給了我們福音的權柄,我們要走出去,把這白白的恩典分享給身邊的人.

活出蒙恩者的生命

各位弟兄姊妹,今天這四段經文,為我們拼湊出一幅完整的「在主前蒙恩」的圖畫:

因為蒙恩,我們的生命不再孤單,有一位向我們側耳聽呼求的神.

因為蒙恩,我們的生命不再絕望,因為在神沒有難成的事.

因為蒙恩,我們的生命滿有安全感,因為在我們還作罪人時祂就愛了我們.

因為蒙恩,我們的生命有了宣教的使命——白白得來,白白捨去.

願我們每一個人,每天都深深地浸泡在主前的恩典中,並且挺起胸膛,走入人群,成為這個時代傳遞天國恩典的管道.



\newpage

\section{從一扇門到心深處}
\label{sec:801}
\textbf{鄧瑞強牧師 - 從一扇門,到心深處}
\newline
\newline
\hyperref[sec:800]{< < < PREV SERMON < < <}
~
\hyperref[sec:toc]{[返目錄]}
~
\hyperref[sec:802]{> > > NEXT SERMON > > >}


約翰福音 10:1-18
\newline
\begin{longtable}{cl}
\hline
\hline
章節 & 經文 (和合本修訂版)\\
\hline
10:1 & \begin{tabularx}{0.7\textwidth}{X} 「我實實在在地告訴你們，那不從門進羊圈，倒從別處爬進去的，就是賊，就是強盜。 \end{tabularx} \\ \\ \relax
10:2 & \begin{tabularx}{0.7\textwidth}{X} 那從門進去的才是羊的牧人。 \end{tabularx} \\ \\ \relax
10:3 & \begin{tabularx}{0.7\textwidth}{X} 看門的給他開門，羊也聽他的聲音。他按著名叫自己的羊，把羊領出來。 \end{tabularx} \\ \\ \relax
10:4 & \begin{tabularx}{0.7\textwidth}{X} 當他把自己的羊都放出來，就走在前面，羊也跟著他，因為牠們認得他的聲音。 \end{tabularx} \\ \\ \relax
10:5 & \begin{tabularx}{0.7\textwidth}{X} 羊絕不跟陌生人，反而會逃走，因為不認得陌生人的聲音。」 \end{tabularx} \\ \\ \relax
10:6 & \begin{tabularx}{0.7\textwidth}{X} 耶穌把這比方告訴他們，但他們不明白他所說的是甚麼。 \end{tabularx} \\ \\ \relax
10:7 & \begin{tabularx}{0.7\textwidth}{X} 所以，耶穌又對他們說：「我實實在在地告訴你們，我就是羊的門。 \end{tabularx} \\ \\ \relax
10:8 & \begin{tabularx}{0.7\textwidth}{X} 凡在我以前來的都是賊，是強盜；羊沒有聽從他們。 \end{tabularx} \\ \\ \relax
10:9 & \begin{tabularx}{0.7\textwidth}{X} 我就是門，凡從我進來的，必得安全，並且可進出，找到草吃。 \end{tabularx} \\ \\ \relax
10:10 & \begin{tabularx}{0.7\textwidth}{X} 盜賊來，無非要偷竊、殺害、毀壞；我來了，是要羊得生命，並且得的更豐盛。 \end{tabularx} \\ \\ \relax
10:11 & \begin{tabularx}{0.7\textwidth}{X} 「我是好牧人，好牧人為羊捨命。 \end{tabularx} \\ \\ \relax
10:12 & \begin{tabularx}{0.7\textwidth}{X} 雇工不是牧人，羊不是他自己的，他一看見狼來，就撇下羊群逃跑；狼抓住羊，把牠們趕散。 \end{tabularx} \\ \\ \relax
10:13 & \begin{tabularx}{0.7\textwidth}{X} 雇工逃走，因為他是雇工，對羊毫不關心。 \end{tabularx} \\ \\ \relax
10:14 & \begin{tabularx}{0.7\textwidth}{X} 我是好牧人；我認識我的羊，我的羊也認識我， \end{tabularx} \\ \\ \relax
10:15 & \begin{tabularx}{0.7\textwidth}{X} 正如父認識我，我也認識父一樣；並且我為羊捨命。 \end{tabularx} \\ \\ \relax
10:16 & \begin{tabularx}{0.7\textwidth}{X} 我另外有羊，不屬這圈裡的，我必須領牠們來，牠們也要聽我的聲音，並且要合成一群，歸一個牧人。 \end{tabularx} \\ \\ \relax
10:17 & \begin{tabularx}{0.7\textwidth}{X} 為此，我父愛我，因為我把命捨去，好再取回來。 \end{tabularx} \\ \\ \relax
10:18 & \begin{tabularx}{0.7\textwidth}{X} 沒有人奪去我的命，是我自己捨的；我有權捨棄，也有權再取回。這是我從我父所受的命令。」 \end{tabularx} \\ \\
[1ex]
\hline
\hline
\end{longtable}
講題:從一扇門,到心深處 From the Gate to the Inner Self

經課:約翰福音10:1-18

講員:鄧瑞強牧師

場合:香港中文大學崇基學院禮拜堂主日崇拜

日期:2026年6月21日

各位弟兄姊妹,早安.願神的恩惠常與你們同在.

今日是父親節,願神祝福每位父親.每位父親,都是天父在人間的代表,付出自己,守護家庭,成就兒女的生命.保羅在哥林多前書4:15裡說:「…師傅雖有一萬,為父的卻是不多…」,所謂「師傅」,類乎教練,重點在訓練孩子成功；保羅口中的「父親」,卻不同,重點在於以高貴的人格,引領孩子歸向神,提升孩子生命的神聖內涵.近來有一則熱門新聞,一位出名的教練以公開羞辱的方式訓練球員,效果可能很成功,令球隊的成績進步,但這肯定不是保羅口中的「父親」.為父的,教孩子尊重他人,教孩子領悟生命的尊嚴比成功更重要.多一點這樣的父親,將來的世界才不會變成格鬥的場所.願天父引領每位父親,成為真正的父親.

之前,我已講過頭三卷福音書.馬太福音用心表明:耶穌在什麼意義上超越猶太教,以激發耶穌的同胞用心在耶穌身上領悟神的新作為.馬可福音表明耶穌的十架人生的沉重,讓我們以嚴肅的態度視之.路加福音卻藉耶穌顯明天父的大愛在人間,人生被天父的大愛承托著.約翰福音卻表明耶穌是宇宙的大道,這大道貫注人心,讓人成為流露大道的器皿.用約翰福音的話來說,這就是永生——真正豐盛的生命.

今日,我們看約翰福音一段熟悉的經文,那段「耶穌是好牧人」的經文,在耶穌這「大道」裡,細味我們應有的信仰生命.這段經文開頭時講「門」,最後卻講內心的最深領悟,像是引領我們,從一扇門進入,然後,一步一步,進到內心深處.

我們看約翰福音10章1-18節.經文頗長,我將之分成五個段落.

(1)約10: 1-6

約10:1 「我實實在在地告訴你們,人進羊圈,不從門進去,倒從別處爬進去,那人就是賊,就是強盜.

約10:2 從門進去的,才是羊的牧人.

約10:3 看門的就給他開門；羊也聽他的聲音.他按著名叫自己的羊,把羊領出來.

約10:4 既放出自己的羊來,就在前頭走,羊也跟著他,因為認得他的聲音.

約10:5 羊不跟著生人；因為不認得他的聲音.必要逃跑.」

約10:6 耶穌將這比喻告訴他們,但他們不明白所說的是什麼意思.

約10:1:「我實實在在地告訴你們」,耶穌認真的講,希望我們認真的聽.這裡用了舊約「牧人與羊」的意念,這是舊約描繪「神與祂的子民」的關係的經典圖畫,如:我們熟悉的詩篇23篇,「耶和華是我的牧者」.耶穌在這裡認真地講,神與祂的子民的關係的重要道理.

這段經文的主要動作,是「進羊圈」.羊,就是神的子民,就是我們.「進羊圈」,像是帶領我們,一步一步了解信仰生命的真正含意.這裡分別兩種人,兩種進法.「不從門進去」的,是賊,是強盜,是屬於黑暗的.「從門進去的」,是牧羊人,是屬於光明的.黑暗與光明的對立,這是約翰福音一貫的主題.

唯有從正門進羊圈的,才是真牧人.羊認得牧羊人的聲音,跟著他走.我們這裡應該無人帶過羊,但很多人帶過狗.你叫牠的名,Bobby,牠老遠,都即刻跑過來,搖下條尾.你會摸下牠,對牠說:good boy.若你叫牠,牠不應,卻跟著一個給牠薯條吃的陌生人,還越走越遠,這條狗一定是bad boy.

這是信仰生命的起點,認定耶穌,認定他是我們的牧人,就緊緊跟著他,不再跟別人了.縱或這邊有人應許成功,那邊有人承諾富貴,都不跟了,唯獨跟從耶穌.

(2)約10:7-10

約10:7 所以,耶穌又對他們說:「我實實在在地告訴你們,我就是羊的門.

約10:8 凡在我以先來的都是賊,是強盜；羊卻不聽他們.

約10:9 我就是門；凡從我進來的,必然得救,並且出入得草吃.

約10:10 盜賊來,無非要偷竊,殺害,毀壞；我來了,是要叫羊〔或譯:人〕得生命,並且得的更豐盛.」

耶穌再一次講:「我實實在在地告訴你們」,他真是很認真的.剛才提到「門」,現在開展關於「門」的討論.門在哪裡?耶穌說:「我就是門」.明顯,這「門」與豐盛人生有關.這是「生命的門」.

約10:8:「凡在我以先來的都是賊,是強盜」,所謂「以先」,有兩種可能解釋.

一,「時間」上的「以先」,按上文下理,是指第9章那些在耶穌出道以先,已在教導猶太人的法利賽人.

二,「以先」也可指「優次上的排在前頭」,即:那些將自己放在比耶穌更優先次序的人.我覺得這解釋更有意思.按這解釋,法利賽人或我們任何一個人,都可以是「放自己在耶穌以先」的人.這些人,是賊,是強盜.

這種人,對別人的生命,造成很多破壞.經文說:「盜賊來,無非要偷竊,殺害,毀壞」.將自己放在耶穌之上的人,只會利用耶穌,來滿足自己.在群體裡,他們會引導別人注意自己,而不是教人仰望耶穌.只願自己的旨意成就,而不是願神的旨意成就.這類人,充滿私意,令人看不到神.

耶穌是「門」.從這「門」進入的,經文說:

一,「必然得救,並且出入得草吃」,這就是生命得到豐盛的供應,得到真正的養分.

二,「叫羊得生命,並且得的更豐盛」,這不是指「物質財富」,而是指生命意義的圓滿.孔子說:「朝聞道,夕死可矣」.朝早悟到大道,晚上死去,也無憾了.孔子很明白,了悟生命的真正意義,比生命的長短更重要.

這是信仰生命的第二個層次.放下自己,將耶穌放在最高的位置上.不單聽從耶穌的聲音,更以耶穌為進入圓滿生命的大門.

(3)約10:11-15

約10:11 「我是好牧人；好牧人為羊捨命.

約10:12 若是雇工,不是牧人,羊也不是他自己的,他看見狼來,就撇下羊逃走；狼抓住羊,趕散了羊群.

約10:13 雇工逃走,因他是雇工,並不顧念羊.

約10:14 我是好牧人；我認識我的羊,我的羊也認識我,

約10:15 正如父認識我,我也認識父一樣；並且我為羊捨命.」

耶穌進一步以牧羊人與羊的關係,來表明合乎正道的生活態度.

牧人與雇工,兩種不同的存活方式.先講雇工.雇工的重點是「賺錢」,與人相處,都只是找利益.身邊的人,對他而言,只是討生活的「工具」.無事時,工作看來正常.狼來了,危險來了,要找人背鑊了,他便不理對方,一走了之.經文說:「雇工逃走,因他是雇工,並不顧念羊.」他根本不愛身邊的人.他只是討生活,虧本的事,他不去做的.

好牧人,則不同,是另一種存活方式.就正正因為狼來了,危險來了,他更守住崗位.他知道狼會來,而他明白,他的存在就是為了這個時刻.他努力守住崗位,因他想好好守護羊,他知道,這是他的天職.狼來了,不是他逃走的理由,反而是他站立在那裡的理由.他認識他的羊,他的羊也認識他.一種生命的連繫,愛的連繫.這愛的連繫,深刻到像聖父與聖子的連繫一樣.就因這愛的連繫,好牧人不會放棄他的羊,他甚至願意,為他們捨命.

信仰生命來到這個層次,要求我們對自己的存活方式,有一深刻反省.我們是雇工,抑或是好牧人?我們如何對待他人?我們的生活追求什麼?我們在人類悠長的歷史中出現在這個特定的時間點,我們的天命何在?時代的艱難是我們逃避責任的藉口,抑或是我們背起十架的機遇?人間有無數「討生活的雇工」,但若能多一點好牧人,人間將多麼不同.

(4)約10:16

約10:16 「我另外有羊,不是這圈裡的；我必須領他們來,他們也要聽我的聲音,並且要合成一群,歸一個牧人了.」

耶穌講到,他不單有眼前的羊,即眼前的猶太人,他還有另外的羊,即全世界的非猶太人,他們都是神所愛的羊.這觀念突破了當時傳統猶太教限定的「神的子民」的範圍.傳統猶太教規定,要成為神的子民,必先成為猶太人,守猶太教的一切規矩.這其實是一種排外的思想.

耶穌講的,是心靈的擴闊,突破人為的界線,擁抱界外之人.現代世界為了守護自己的利益,越來越排外,合作越來越難.耶穌想的,卻是要打破藩籬,促成不同社群間的和解,合一.

我們每個禮拜認信時宣讀的「使徒信經」,有一句是「我信聖而公(catholic)之教會」.所謂「公」,即「大公精神」,一種「有容乃大」的包融精神.這包括「民族,文化,經濟,地域」的「大公精神」,即擁抱不同民族,文化,經濟背景,居住地域的人.也包括「時間」上的「大公精神」,即思念過去的人,關心當下的人,顧及將來的人.

耶穌的心靈,是一「擴而充之」的心靈,突破人為的界限,使仁愛及於未及之處.信仰不應使人越來越封閉,應使人打開心門,向他人開放.越多這樣的信徒,基督信仰才有將來.

(5)約10:17-18

約10:17 「我父愛我；因我將命捨去,好再取回來.

約10:18 沒有人奪我的命去,是我自己捨的.我有權柄捨了,也有權柄取回來.這是我從我父所受的命令.」

耶穌這裡表達他對生命的最高領悟.面對生死,他不訴之命運,而訴之自由.面對十字架,這不是遇著惡運,而是自己權能下的抉擇.耶穌一生,浸淫在天父的愛中.若活著能流露天父對人的愛,他就選擇活著.若死去能讓人驚訝於天父的愛的闖入,他就選擇死去.生的欲望,死的恐懼,並不主宰耶穌.主宰耶穌的,是天父的心,是他從天父那裡領受的天命.他面向天父,展現真正的自由,自由地出入生死,自由地去愛,自由地活出生命的豐盛內容.

這是信仰生命的最高體會了.活著,就是實踐天命,流露天心.不是有人強迫他捨命,不是因為要履行某些律法而捨命,捨命,是他出於自由的行動.他不會問:這公平嗎?他不會想:如何可逃避呢?他捨命,是為了愛.他真心去愛,因為他體會這是天父內心的想法.他的想法,與天父的想法,完全一致.活著,就是天父的心的流露.

2010年,有齣法語電影上映,法語片名直譯是《人與神》,英文片名是《Of Gods and Men》,香港片名是《人神之間》.根據真人真事改編.我在神學院教「教牧神學」或「教牧倫理」時,這電影是必看的,是堂上討論的內容.

話說,在北非阿爾及利亞的一條回教村落裡,有一間天主教修道院,住著9名法國修士,他們在那裡修道,宣教,醫病,和那裡的回教村民和諧共處.忽然,一群極端的回教武裝分子出現,殺害外國人.之後,整齣電影的內容就是:在死亡的威脅下,這9名修士如何不斷討論去留問題.起初,不少修士恐懼死亡,想返回法國.修士對村民說,我們是鳥,有天會飛走的.但村民說,我們才是鳥,你們是樹,你們堅立在這裡,我們才能有安身的地方.

電影一幕一幕地呈現修院的生活,同時也逐步展示基督信仰對人生命的轉化.聖誕節的崇拜,讓修士思考基督卑微的臨在,與真正的和平.每天的聖餐環節,讓修士明白基督的犧牲,及基督在受苦群體中的臨在.在受難節的最後晚餐中,體會基督的受難,並在死人中復活所帶來的希望.

最後,所有修士都與基督的生命融成一體,以謙和的生命,與受苦者同在,以自由,選擇置生死於度外,繼續在那貧窮社區,靜靜地流露天父的仁愛.結束時,恐怖分子襲擊修院,擄走7位修士,有2位躲藏了而避過一刼.這7位修士最後被殺害.他們陪伴基督,活出了豐盛的人生.

保羅曾說:「我們若活著,是為主而活；若死了,是為主而死.所以,我們或活或死總是主的人.」(羅14:8)又說:「我已經與基督同釘十字架,現在活著的不再是我,乃是基督在我裡面活著.」(加2:20)一個人,真的可以,或活或死,都是神的天心的流露.我希望,你能體會這種生命,活出這種生命,這是最豐盛的生命,這是永生.

願神賜福你們.

願榮耀歸予聖父,聖子,聖靈.



\newpage

\section{要到幾時呢}
\label{sec:802}
\textbf{林豪恩先生 - 要到幾時呢?}
\newline
\newline
\hyperref[sec:801]{< < < PREV SERMON < < <}
~
\hyperref[sec:toc]{[返目錄]}
~
\hyperref[sec:803]{> > > NEXT SERMON > > >}


創世記 22:1-14
\newline
\begin{longtable}{cl}
\hline
\hline
章節 & 經文 (和合本修訂版)\\
\hline
22:1 & \begin{tabularx}{0.7\textwidth}{X} 這些事以後，神考驗亞伯拉罕，對他說：「亞伯拉罕！」他說：「我在這裡。」 \end{tabularx} \\ \\ \relax
22:2 & \begin{tabularx}{0.7\textwidth}{X} 神說：「你要帶你的兒子，就是你所愛的獨子以撒，往摩利亞地去，在我指示你的一座山上，把他獻為燔祭。」 \end{tabularx} \\ \\ \relax
22:3 & \begin{tabularx}{0.7\textwidth}{X} 亞伯拉罕清早起來，預備了驢，帶著跟他一起的兩個僕人和他兒子以撒，劈好了燔祭的柴，就起身往神指示他的地方去了。 \end{tabularx} \\ \\ \relax
22:4 & \begin{tabularx}{0.7\textwidth}{X} 到了第三日，亞伯拉罕舉目遙望那地方。 \end{tabularx} \\ \\ \relax
22:5 & \begin{tabularx}{0.7\textwidth}{X} 亞伯拉罕對他的僕人說：「你們和驢留在這裡，我和孩子要去那裡敬拜，然後回到你們這裡來。」 \end{tabularx} \\ \\ \relax
22:6 & \begin{tabularx}{0.7\textwidth}{X} 亞伯拉罕把燔祭的柴放在他兒子以撒身上，自己手裡拿著火與刀；於是二人同行。 \end{tabularx} \\ \\ \relax
22:7 & \begin{tabularx}{0.7\textwidth}{X} 以撒對他父親亞伯拉罕說：「我父啊！」亞伯拉罕說：「我兒，我在這裡。」以撒說：「看哪，火與柴都有了，但燔祭的羔羊在哪裡呢？」 \end{tabularx} \\ \\ \relax
22:8 & \begin{tabularx}{0.7\textwidth}{X} 亞伯拉罕說：「我兒，神必自己預備燔祭的羔羊。」於是二人同行。 \end{tabularx} \\ \\ \relax
22:9 & \begin{tabularx}{0.7\textwidth}{X} 他們到了神指示他的地方，亞伯拉罕在那裡築壇，把柴擺好，綁了他兒子以撒，放在壇的柴上。 \end{tabularx} \\ \\ \relax
22:10 & \begin{tabularx}{0.7\textwidth}{X} 亞伯拉罕就伸手拿刀，要殺他的兒子。 \end{tabularx} \\ \\ \relax
22:11 & \begin{tabularx}{0.7\textwidth}{X} 耶和華的使者從天上呼喚他說：「亞伯拉罕！亞伯拉罕！」他說：「我在這裡。」 \end{tabularx} \\ \\ \relax
22:12 & \begin{tabularx}{0.7\textwidth}{X} 天使說：「不可在這孩子身上下手！一點也不可傷害他！現在我知道你是敬畏神的人了，因為你沒有把你的兒子，就是你的獨子，留下不給我。」 \end{tabularx} \\ \\ \relax
22:13 & \begin{tabularx}{0.7\textwidth}{X} 亞伯拉罕舉目觀看，看哪，一隻公綿羊兩角纏在灌木叢中。亞伯拉罕就去牽了那隻公綿羊，獻為燔祭，代替他的兒子。 \end{tabularx} \\ \\ \relax
22:14 & \begin{tabularx}{0.7\textwidth}{X} 亞伯拉罕給那地方起名叫「耶和華以勒」 。直到今日人還說：「在耶和華的山上必有預備。」 \end{tabularx} \\ \\
[1ex]
\hline
\hline
\end{longtable}
羅馬書 6:12-23
\newline
\begin{longtable}{cl}
\hline
\hline
章節 & 經文 (和合本修訂版)\\
\hline
6:12 & \begin{tabularx}{0.7\textwidth}{X} 所以，不要讓罪在你們必死的身上掌權，使你們順從身體的私慾。 \end{tabularx} \\ \\ \relax
6:13 & \begin{tabularx}{0.7\textwidth}{X} 也不要把你們的肢體獻給罪作不義的工具，倒要像從死人中活著的人，把自己獻給神，並把你們的肢體獻給神作義的工具。 \end{tabularx} \\ \\ \relax
6:14 & \begin{tabularx}{0.7\textwidth}{X} 罪必不能作你們的主，因你們不在律法之下，而是在恩典之下。 \end{tabularx} \\ \\ \relax
6:15 & \begin{tabularx}{0.7\textwidth}{X} 那又怎麼樣呢？我們在恩典之下，不在律法之下，就可以犯罪嗎？絕對不可！ \end{tabularx} \\ \\ \relax
6:16 & \begin{tabularx}{0.7\textwidth}{X} 難道你們不知道，你們獻自己作奴僕，順從誰就作誰的奴僕嗎？或作罪的奴隸，以至於死；或作順服的奴僕，以至於成義。 \end{tabularx} \\ \\ \relax
6:17 & \begin{tabularx}{0.7\textwidth}{X} 感謝神！因為你們從前雖然作罪的奴隸，現在卻從心裡順服了所傳給你們教導的典範。 \end{tabularx} \\ \\ \relax
6:18 & \begin{tabularx}{0.7\textwidth}{X} 你們既從罪裡得了釋放，就作了義的奴僕。 \end{tabularx} \\ \\ \relax
6:19 & \begin{tabularx}{0.7\textwidth}{X} 我因你們肉體的軟弱，就以人的觀點來說。你們從前怎樣把肢體獻給不潔不法作奴隸，以至於不法；現在也要照樣將肢體獻給義作奴僕，以至於成聖。 \end{tabularx} \\ \\ \relax
6:20 & \begin{tabularx}{0.7\textwidth}{X} 因為你們作罪的奴隸時，不被義所約束。 \end{tabularx} \\ \\ \relax
6:21 & \begin{tabularx}{0.7\textwidth}{X} 那麼，你們現在所看為羞恥的事，當時有甚麼果子呢？那些事的結局就是死。 \end{tabularx} \\ \\ \relax
6:22 & \begin{tabularx}{0.7\textwidth}{X} 但如今，你們既從罪裡得了釋放，作了神的奴僕，就結出果子，以至於成聖，那結局就是永生。 \end{tabularx} \\ \\ \relax
6:23 & \begin{tabularx}{0.7\textwidth}{X} 因為罪的工價乃是死；惟有神的恩賜，在我們的主基督耶穌裡，乃是永生。 \end{tabularx} \\ \\
[1ex]
\hline
\hline
\end{longtable}
馬太福音 10:40-42
\newline
\begin{longtable}{cl}
\hline
\hline
章節 & 經文 (和合本修訂版)\\
\hline
10:40 & \begin{tabularx}{0.7\textwidth}{X} 「接納你們的就是接納我；接納我的就是接納差遣我來的那位。 \end{tabularx} \\ \\ \relax
10:41 & \begin{tabularx}{0.7\textwidth}{X} 把先知當作先知接納的，必得先知的賞賜；把義人當作義人接納的，必得義人的賞賜。 \end{tabularx} \\ \\ \relax
10:42 & \begin{tabularx}{0.7\textwidth}{X} 無論誰，只因門徒的名，就算把一杯涼水給這些小子中的一個喝，我實在告訴你們，他一定會得到賞賜。」 \end{tabularx} \\ \\
[1ex]
\hline
\hline
\end{longtable}
講題:要到幾時呢? How Long, O LORD?

經課:創世記22章1至14節；詩篇13篇1至6節；羅馬書6章12至23節；馬太福音10章40至42節

講員:林豪恩先生

場合:香港中文大學崇基學院禮拜堂主日崇拜

日期:2026年6月28日

引言:

有人覺得時間過得很快,又有人覺得時間過得很慢.你覺得時間過得快還是慢呢?時間不是客觀地而獨立的嗎?一日二十四小時就是二十四小時,不可能有人一日只有十二小時,而另一人一日卻有三十小時.不過,人們感覺時間的長短並不決定於客觀時間,而是受到「主觀時間知覺」影響,「主觀時間知覺」又與人的狀況和情緒相關.很多人都說「快樂不知時日過」,所以感到放假的日子過得特別快.也有人感到工作的日子過得特別慢,經常聽到人說:「要等都幾時才有假放呢?」原來在某些狀況之下,人們感到時間被拉長,正所謂「度日如年」.剛才我們讀詩篇時候,聽到詩人呼喊:「要到幾時呢?」是甚麼狀況令詩人發出「要到幾時呢」這呼喚呢?除了詩人之外,經課還提及甚麼類似的狀況呢?他們如何渡過「要到幾時呢」這種漫長等待呢?讓我們一齊閱讀他們的經驗和出路,思考如何成為我們的啟發.

制度的缺失要都幾時呢(詩篇)

13:1耶和華啊,你忘記我要到幾時呢?要到永遠嗎?你掩面不顧我要到幾時呢?

13:2我心裡籌算,終日愁苦,要到幾時呢?我的仇敵升高壓制我,要到幾時呢?

13:3耶和華－我的上帝啊,求你看顧我,應允我!使我眼目光明,免得我沉睡至死；

13:4免得我的仇敵說:我勝了他；免得我的敵人在我搖動的時候喜樂.

從詩句可見,詩人「被仇敵壓制」,因而感到「終日愁苦」,苦無出路看不見曙光,感到連上帝也忘記了自己,甚至刻意「掩面不顧他」.如此困苦的日子,要到幾時呢?其實具體是甚麼事情況呢?詩歌沒有說明.相傳這是大衛的詩,於是人們估計可能和下列兩件事相關:其一是兒子押沙龍發動軍事政變奪取王位所引發的逃亡及戰亂,事件令大衛經歷家破人亡的震驚及哀傷.另一個可能是大衛被掃羅追殺,因而四處逃亡躲避.

詩人的痛苦看似是個人的,而的確也是個人的；然而,卻又不單單是個人的,詩人的痛苦其實反映某種歷史時空的社會狀況,他的困境是因為被捲入這種歷史狀況而產生的.讓我們回頭問:掃羅為甚麼要追殺大衛呢?會不會是因為當時沒有能夠和平處理權力鬥爭的制度呢?當一山不能藏二虎的時候,只有透過消滅對方來解除威脅.大衛的兒子押沙龍為甚麼要發動軍事政變奪取王位呢?會不會是當時沒有能夠和平處理權力交接的制度呢?沒有和平的權力交接制度,眾王子爭位就變成可預期的兄弟相殺,引發宮廷內亂,這些悲劇在古代王朝比比皆是.大衛感到痛苦當然可以呼喊「要到幾時呢」,然而,以大衛所在的歷史空間及權力位置,例如在掃羅作王時被視為權力競爭者,在作王晚年時眾子各懷鬼胎爭奪王位,這些狀況引至大衛痛苦是可以想像的.類似大衛的困擾要到幾時呢?或許要到人們已經建立了和平非暴力的權力交替制度.因此,「要到幾時呢?」這呼喊,同時也是對壞制度的投訴,以及對良好制度的盼望.就算在我們身處的時空,有些痛苦的狀況都未必是受苦者自己引起的,有可能只是因為當事人剛好處身於不良制度之下,無可避免地受到傷害.

在宏福苑大火聽證會場外,有居民接受訪問時含著眼淚說:「我地宏福苑呢單,只不過係今日香港所發生ge 其中一件事,再咁樣落去,仲有其他ge 事情,一路會發生......宏福苑只係話比你知,香港係有問題」.

傳媒報導指出,「大維修」計劃初推行時,便已有質疑它會助長管理方貪污及操縱投標 ,以及承建商誇大費用.當年廉政公署審查貪污舉報諮詢委員會主席施祖祥也認為此一計劃因由政府注入大量資金,容易成為不少投機分子的目標,有必要更慎重監管.人們一直期望健全的監管制度,但「要到幾時呢」?多年來宏福苑居民多次多方向政府部門投訴大維修圍標,棚網違規及工人吸煙等問題,亦曾投訴有人以發泡膠板封窗阻礙救援,這些狀況要到幾時才得到認真回應呢?要到幾時才肯處理這些危機呢?要到幾時才肯改善管理呢?

二．心愛的受苦要都幾時呢(創世記)

有些「要到幾時呢」的呼喊是以年月計算的,但有些「要到幾時呢」的呼喊可能是轉瞬之間的.可能有人感到奇怪,轉瞬之間的短時間也忍受不住嗎?可以的,因為我們剛才已經說過,時間感覺是和狀況相關的.是甚麼狀況令人感到一刻也難捱呢?

22:1這些事以後,上帝要試驗亞伯拉罕,就呼叫他說:「亞伯拉罕!」他說:「我在這裡.」

22:2上帝說:「你帶著你的兒子,就是你獨生的兒子,你所愛的以撒,往摩利亞地去,在我所要指示你的山上,把他獻為燔祭.」

22:3亞伯拉罕清早起來,備上驢,帶著兩個僕人和他兒子以撒,也劈好了燔祭的柴,就起身往上帝所指示他的地方去了.

22:4到了第三日,亞伯拉罕舉目遠遠地看見那地方.

22:5亞伯拉罕對他的僕人說:「你們和驢在此等候,我與童子往那裡去拜一拜,就回到你們這裡來.」

22:6亞伯拉罕把燔祭的柴放在他兒子以撒身上,自己手裡拿著火與刀；於是二人同行.

22:7以撒對他父親亞伯拉罕說:「父親哪!」亞伯拉罕說:「我兒,我在這裡.」以撒說:「請看,火與柴都有了,但燔祭的羊羔在哪裡呢?」

22:8亞伯拉罕說:「我兒,上帝必自己預備作燔祭的羊羔.」於是二人同行.

22:9他們到了上帝所指示的地方,亞伯拉罕在那裡築壇,把柴擺好,捆綁他的兒子以撒,放在壇的柴上.

壇已築好,柴已lor 好,刀已在手,亞伯拉罕等待那「上帝必自己預備作燔祭的羊羔」的時候有多長呢?客觀上來說不長,因為不久之後天上就有聲音出現,然後亞伯拉罕就見到祭物就出現了.然而,對亞伯拉罕來說,在那生死交關的瞬間,焦急等待的煎熬,一秒鐘彷如是一個世紀.經文沒有描述亞伯拉罕的內心感受,然而,如果亞伯拉罕是一個正常的父親,怎會沒有掙扎呢?怎會看見自己所愛的人受傷害而沒有感受呢?要到幾時呢?是因為心愛的人處於危難之中,有人分享了在手術室外等候的心情:「手術室的門緊閉著,那扇冰冷的門後是我最深切的擔憂.我呆滯地盯著手術室的大門,心中默默祈禱.時間仿佛在這一刻停滯了,每一分鐘都像是被拉長了一個世紀.11點,12點,1點……螢幕上「手術中」三個字像是一把懸在我心頭的劍,讓我無法平靜.直到下午3點,手術室的門依舊緊閉,螢幕上「手術中」的字樣刺眼而沉重.到了下午4點,我再也坐不住了.那種無助和焦慮幾乎讓我窒息,我再也無法忍受這種等待的煎熬,焦急地向工作人員詢問.5點,我在手術室外焦急地踱步,生怕錯過從手術室推出來的每一個瞬間.每一次手術室的門打開,我的心都會跟著緊張地跳動,仿佛要跳出胸膛.」

三．罪在你們必身上作王要到幾時呢?(羅馬書)

「要到幾時呢?」常常是人們的呼喊.然而,我們有沒有想為,除了人之外,還有誰都可能發出類似的呼喊呢?

6:12所以,不要容罪在你們必死的身上作王,使你們順從身子的私慾.

6:13也不要將你們的肢體獻給罪作不義的器具；倒要像從死裡復活的人,將自己獻給上帝,並將肢體作義的器具獻給上帝.

當人類呼喊「要到幾時呢?」的時候,會不會造物主也在呼喊:你們人類把自己「獻給罪作不義的器具」要到幾時呢?使徒保羅啟發我們思考:真正需要被終止的不只是苦況本身,也包括罪在人生命中的統治.人若繼續把自己交給私慾與不義,生命就仍被罪作主.當人類呼喊痛苦要到幾時才終止的時候,或許人類也要問自己,我們到幾時才肯悔改呢?到幾時才肯離開罪惡,行善行義呢?

6:21你們現今所看為羞恥的事,當日有甚麼果子呢?那些事的結局就是死.

6:22但現今,你們既從罪裡得了釋放,作了上帝的奴僕,就有成聖的果子,那結局就是永生.

6:23因為罪的工價乃是死；惟有上帝的恩賜,在我們的主基督耶穌裡,乃是永生.

人類在等待上帝介入,主持正義,撥亂反正,賞善罰惡；同時間,上帝會否也在等待人類悔改,離開罪惡,活出上帝創造的美好呢?

四．生活在「要到幾時呢」的呼喊中

聖經的確應許上帝終必會主持正義,撥亂反正,賞善罰惡；聖經也呼籲人們要悔改,離棄罪惡,歸向上帝.在這個尚未完全實現的過程中,我們可以怎樣活在「要到幾時呢」這種過渡期呢?讓我們看看福音經課給我們的啟發:

馬太福音

10:40「人接待你們就是接待我；接待我就是接待那差我來的.

10:41人因為先知的名接待先知,必得先知所得的賞賜；人因為義人的名接待義人,必得義人所得的賞賜.

10:42無論何人,因為門徒的名,只把一杯涼水給這小子裡的一個喝,我實在告訴你們,這人不能不得賞賜.」

「一杯涼水」代表盼望不一定只以巨大奇蹟的方式出現,也可能透過極小,極日常的接待與善意被看見.上文指出門徒受逼迫,但「一杯涼水」出現顯示門徒不是完全孤立無援；上帝的國度也透過人對門徒,先知,義人的接納而具體顯明.盼望從抽象應許落到具體行動上,在艱難中,哪怕是一點點善意,也帶有屬靈重量.

2026年6月7日,明報刊出了一篇宏福居民的訪問,標題是:有愛有家在— 「感謝所有幫過我們的人」Tina很感恩在危難關頭捐款,捐衣服給宏福居民的香港人.空無一物時,她收到新的雪櫃,微波爐,風筒,電飯煲,iPad……連她平時沒有用開的電動牙刷都有.她說:「無緣無故這麼多善款給我們,都足夠我們買新的東西,足夠我們暫時住的地方,我已經滿足的了,我很感謝所有幫過我們的人」.

失去了間屋,不代表失去了家.Tina說,有愛有家在,「屋只不過是住的地方,它不會產生感情,有感情的是人.我的兒子沒有死,我沒有死,我們在一起,我們就有個家.就算我們不是住在一起,家也在,心也在」.

Tina沒有為自己流過淚,但她聽到有鄰居在火災中失去生命,失去家人時,總會為他們流淚.她有留意聽證會,也感到憤怒,為什麼有些人不做好自己本分,「我為那些死的人惋惜,我覺得他們不應該死,但是沒辦法,我又不可以代替他們.我希望這個血的教訓,將來的人不會再重蹈覆轍,只可以這樣想.」

面對巨變和令人心傷的消息,她更頻繁地念經,「有時睡不著就念經,有時我都會為所有死了的人念經,還有為這件事受痛苦的那些人,我為他們念經」.「念經的時候,有一種對我的抑鬱焦慮很有幫助的感覺,因為全心全意集中精神在念經.」念經時也要懺悔,為自己前生做的壞事懺悔,也為所有的眾生懺悔,放下了一些鬱結和心情.

「要到幾時呢?」這呼喊表達了對於苦難,沉默,試煉,逼迫何時過去的詰問.然而,同時也在問我們,在等待中,我們會成為甚麼樣的人?我們的社會成為怎樣的社會?願上主幫助我們.



\newpage

\section{好人難做}
\label{sec:803}
\textbf{高國雄牧師 - 好人難做}
\newline
\newline
\hyperref[sec:802]{< < < PREV SERMON < < <}
~
\hyperref[sec:toc]{[返目錄]}
~
\hyperref[sec:804]{> > > NEXT SERMON > > >}


羅馬書 7:15-25
\newline
\begin{longtable}{cl}
\hline
\hline
章節 & 經文 (和合本修訂版)\\
\hline
7:15 & \begin{tabularx}{0.7\textwidth}{X} 因為我所做的，我自己不明白。我所願意的，我並不做；我所恨惡的，我反而去做。 \end{tabularx} \\ \\ \relax
7:16 & \begin{tabularx}{0.7\textwidth}{X} 如果我所做的是我所不願意的，我得承認律法是善的。 \end{tabularx} \\ \\ \relax
7:17 & \begin{tabularx}{0.7\textwidth}{X} 事實上，這不是我做的，而是住在我裡面的罪做的。 \end{tabularx} \\ \\ \relax
7:18 & \begin{tabularx}{0.7\textwidth}{X} 我也知道，住在我裡面的，就是我肉體之中，沒有善。因為立志為善由得我，只是行出來由不得我。 \end{tabularx} \\ \\ \relax
7:19 & \begin{tabularx}{0.7\textwidth}{X} 我所願意的善，我不去做；我所不願意的惡，我反而去做。 \end{tabularx} \\ \\ \relax
7:20 & \begin{tabularx}{0.7\textwidth}{X} 如果我去做我不願意做的，就不是我做的，而是住在我裡面的罪做的。 \end{tabularx} \\ \\ \relax
7:21 & \begin{tabularx}{0.7\textwidth}{X} 我覺得有個律，就是我願意行善的時候，就有惡纏著我。 \end{tabularx} \\ \\ \relax
7:22 & \begin{tabularx}{0.7\textwidth}{X} 因為，按著我裡面的人，我喜歡神的律， \end{tabularx} \\ \\ \relax
7:23 & \begin{tabularx}{0.7\textwidth}{X} 但我看出肢體中另有個律和我內心的律交戰，把我擄去，使我附從那肢體中罪的律。 \end{tabularx} \\ \\ \relax
7:24 & \begin{tabularx}{0.7\textwidth}{X} 我真苦啊！誰能救我脫離這必死的身體呢？ \end{tabularx} \\ \\ \relax
7:25 & \begin{tabularx}{0.7\textwidth}{X} 感謝神，靠著我們的主耶穌基督就能！這樣看來，一方面，我內心順服神的律，另一方面，肉體卻順服罪的律了。 \end{tabularx} \\ \\
[1ex]
\hline
\hline
\end{longtable}
馬太福音 11:16-19
\newline
\begin{longtable}{cl}
\hline
\hline
章節 & 經文 (和合本修訂版)\\
\hline
11:16 & \begin{tabularx}{0.7\textwidth}{X} 「我該用甚麼來比這世代呢？這正像孩童坐在街市上向同伴呼喊： \end{tabularx} \\ \\ \relax
11:17 & \begin{tabularx}{0.7\textwidth}{X} 『我們為你們吹笛，你們不跳舞；我們唱哀歌，你們不捶胸。』 \end{tabularx} \\ \\ \relax
11:18 & \begin{tabularx}{0.7\textwidth}{X} 約翰來了，既不吃也不喝，人們就說他是被鬼附的； \end{tabularx} \\ \\ \relax
11:19 & \begin{tabularx}{0.7\textwidth}{X} 人子來了，也吃也喝，他們又說這人貪食好酒，是稅吏和罪人的朋友。而智慧是由它的果子來證實的。」 \end{tabularx} \\ \\
[1ex]
\hline
\hline
\end{longtable}
講題:好人難做 It Is Hard To Be A Good Man

經課:羅馬書7:15-25；馬太福音11:16-19

講員:高國雄牧師

場合:香港中文大學崇基學院神學院容啟東校長紀念樓G/F聖堂

日期:2026年7月5日

現代人的「好人困境」

生活共鳴: 在日常生活中,我們常聽人感嘆「好人難做」.好心幫人可能被當成理所當然,甚至被反咬一口；想面面俱圓,結果卻是兩頭落空.

但聖經將這個問題帶得更深——「好人難做」不只是人際關係的無奈,更是屬靈現實的掙扎.今天這兩段經文,分別從「外在世界」和「內在自我」解構了這個困境.

外在的困境:順得哥情失嫂意(太 11:16-19)

馬太福音中,耶穌用了一個生動的比喻:孩童在街市上玩耍,吹笛時別人不跳舞,舉哀時別人不捶胸.這正反映了世人對「好」的標準是多變,主觀且矛盾的.

世界永遠有挑剔的理由:

施洗約翰來了,禁食奔走,世人說他是「被鬼附的」(太過極端,嚴厲).

耶穌來了,與罪人一同吃喝,世人又說祂是「貪食好酒的人,是稅吏和罪人的朋友」(太過隨性,沒原則).

反思: 如果你想單憑「滿足世界的標準」來當一個好人,你註定會失敗.

因為世界總是用雙重標準在審視屬靈的人.當你試圖討好所有人時,你往往誰也討好不了.

智慧的答案:

「但智慧之子總以智慧為是.」(太 11:19).

好人的標準不在於世界的掌聲,而在於是否行在神的旨意(智慧)之中.

內在的絕望:立志行善由得我,行出來由不得我(羅 7:15-25a)

如果外在環境的刁難還能靠意志力撐過去,那內心的「背叛」才是最讓人絕望的.保羅在羅馬書中赤裸裸地剖白了人性最深處的無能為力.

理智與行為的分裂:

「我所願意的善,我反不做；我所不願意的惡,我倒去做.」(羅 7:19)

我們想當個好丈夫/好妻子,想當個誠實的基督徒,想勒住舌頭不說八卦……我們理智上知道那是「好」的,但一到關鍵時刻,內心的另一條「律」就把我們擄去.

「好人」的真相:

保羅並非不想做好人

「按著裡面的意思,我是喜歡上帝的律」(羅 7:22)

但殘酷的現實是,光靠「立志」是無法勝過罪的權勢的.

靈魂的吶喊:

「我真是苦啊!誰能救我脫離這取死的身體呢?」(羅 7:24)

這是所有渴望做好人,做義人,卻在罪中不斷跌倒之人的真實寫照.

唯一的出路:不是靠自己做,而是「到祂那裡去」

既然外有世界的刁難,內有罪性的捆綁,「好人」是不是就真的做不下去了?感謝主,經文在這裡給了我們驚天動地的轉折.

轉向基督的恩典

保羅在悲痛吶喊後,立刻宣告

「感謝上帝,靠著我們的主耶穌基督就能脫離了.」(羅 7:25a)

我們不需要再靠自己的殘破的力量去證明自己是好人,而是承認自己的無能,接受耶穌在十字架上為我們成就的義.

卸下世界的重擔,換上耶穌的軛(太 11:25-30)

耶穌向那些在「好人難做」的律法與世俗標準下,精疲力竭的人發出邀請:

「凡勞苦擔重擔的人可以到我這裡來,我就使你們得安息.」(太11:28)

世界要你面面俱圓,耶穌只要你「心裡柔和謙卑」.

耶穌的「軛」是容易的,祂的「擔子」是輕省的.因為祂不是在旁邊挑剔你的法官,而是與你同負一軛,替你承擔重量的恩主.

結語與回應:從「努力做好人」到「活出基督」

觀念的更新: 基督信仰的核心從來不是教導我們「如何靠自己努力做個好人」,因為羅馬書已經證明了這條路行不通.基督信仰是關於「一個承認自己做不到的好人,如何每天來到主面前支取生命」.

呼召與應用:

放下討好世界的包袱: 不用再為世界矛盾的標準而焦慮(太 11:16-19).

接納自己的軟弱,停止自我定罪: 面對內心的掙扎(羅 7),來到主面前坦然承認「我真是苦」,並接受祂的恩典.

每天與主同負一軛: 將生命的掌控權交給耶穌,讓祂的生命從我們裡面活出來.這時,行善不再是沉重的律法,而是生命自然流露的果子.

【生命的歌】曲,詞:伍偉基

琴鍵奏起,是你的歌,是每日愛惜我的笑臉,

常勵勉我,又提醒我,要努力奏生命的歌.

眼淚有你,歡呼聲有你,未放棄是因你的信念,

未能看見,但仍相信,神保守我美麗明天.

惟有你的愛感動我,惟有你擁抱改寫我,

你教我懂得去生活,今我知活著要為何,

我願意將愛給予你,回應你對我不捨棄,

琴聲歌聲一天會終結,在空氣中環繞我是你.



\newpage

\section{我都就死這於我有甚麼益處}
\label{sec:804}
\textbf{李均熊牧師 - 我都就死,這於我有甚麼益處?}
\newline
\newline
\hyperref[sec:803]{< < < PREV SERMON < < <}
~
\hyperref[sec:toc]{[返目錄]}
~
\hyperref[sec:805]{> > > NEXT SERMON > > >}


創世記 25:19-34
\newline
\begin{longtable}{cl}
\hline
\hline
章節 & 經文 (和合本修訂版)\\
\hline
25:19 & \begin{tabularx}{0.7\textwidth}{X} 這是亞伯拉罕的兒子以撒的後代。亞伯拉罕生以撒。 \end{tabularx} \\ \\ \relax
25:20 & \begin{tabularx}{0.7\textwidth}{X} 以撒四十歲時娶利百加為妻。利百加是巴旦‧亞蘭地的亞蘭人彼土利的女兒，是亞蘭人拉班的妹妹。 \end{tabularx} \\ \\ \relax
25:21 & \begin{tabularx}{0.7\textwidth}{X} 以撒因他妻子不生育，就為她祈求耶和華。耶和華應允他的祈求，他的妻子利百加就懷了孕。 \end{tabularx} \\ \\ \relax
25:22 & \begin{tabularx}{0.7\textwidth}{X} 胎兒們在她腹中彼此相爭，她就說：「若是如此，我為甚麼會這樣呢？」她就去求問耶和華。 \end{tabularx} \\ \\ \relax
25:23 & \begin{tabularx}{0.7\textwidth}{X} 耶和華對她說：兩國在你腹中；兩族要從你身上分立。這族必強於那族；將來大的要服侍小的。 \end{tabularx} \\ \\ \relax
25:24 & \begin{tabularx}{0.7\textwidth}{X} 到了生產的日期，看哪，腹中是對雙胞胎。 \end{tabularx} \\ \\ \relax
25:25 & \begin{tabularx}{0.7\textwidth}{X} 先出生的身體帶紅，渾身有毛，好像皮衣；他們就給他起名叫以掃。 \end{tabularx} \\ \\ \relax
25:26 & \begin{tabularx}{0.7\textwidth}{X} 隨後，以掃的弟弟也出生，他的手抓住以掃的腳跟，因此給他起名叫雅各。兩個兒子出生時，以撒六十歲。 \end{tabularx} \\ \\ \relax
25:27 & \begin{tabularx}{0.7\textwidth}{X} 兩個孩子漸漸長大，以掃善於打獵，常在田野；雅各為人安靜，常住在帳棚裡。 \end{tabularx} \\ \\ \relax
25:28 & \begin{tabularx}{0.7\textwidth}{X} 以撒愛以掃，因為常吃他的野味；利百加卻愛雅各。 \end{tabularx} \\ \\ \relax
25:29 & \begin{tabularx}{0.7\textwidth}{X} 有一天，雅各熬了湯，以掃從田野回來，疲憊不堪。 \end{tabularx} \\ \\ \relax
25:30 & \begin{tabularx}{0.7\textwidth}{X} 以掃對雅各說：「我累死了，請你讓我吃這紅的，這紅的湯吧！」因此以掃又叫以東。 \end{tabularx} \\ \\ \relax
25:31 & \begin{tabularx}{0.7\textwidth}{X} 雅各說：「你今日把長子的名分賣給我吧。」 \end{tabularx} \\ \\ \relax
25:32 & \begin{tabularx}{0.7\textwidth}{X} 以掃說：「看哪，我快要死了，這長子的名分對我有甚麼用呢？」 \end{tabularx} \\ \\ \relax
25:33 & \begin{tabularx}{0.7\textwidth}{X} 雅各說：「你今日對我起誓吧。」以掃就向他起誓，把長子的名分賣給了雅各。 \end{tabularx} \\ \\ \relax
25:34 & \begin{tabularx}{0.7\textwidth}{X} 於是雅各把餅和豆湯給了以掃，以掃吃喝以後，起來走了。這樣，以掃輕看他長子的名分。 \end{tabularx} \\ \\
[1ex]
\hline
\hline
\end{longtable}
講題:我都就死,這於我有甚麼益處?I Am Going To Die, What Profit Shall This Do To Me?

經課:創世記25:19-34

講員:李均熊牧師

場合:香港中文大學崇基學院神學院容啟東校長紀念樓G/F聖堂

日期:2026年7月12日

起:信耶穌唔通有飯食?

今日舊約經課當中以掃講ge gwo 一句說話:「我都就死,這於我有甚麼益處?」就好似上一輩成日講ge 一句說話「信耶穌唔通有飯食咩?」這句說話,表面聽落好似好世俗,好功利,甚至有點無賴.但有時,正正是這種說話,最能夠講出人心底深處的關注.信仰講應許,講將來,講上主的帶領；但人活在眼前,活在當下,會餓,會累,會驚,會覺得撐唔住.當眼前的需要太近,將來的價值就會變得很遠.

所以如果我們太快把以掃那句「我都就死,這於我有甚麼益處?」否定,讀成一個反面教材,說他只是貪食,短視,輕看屬靈的事,可能這段舊約經課對我們意義不大.創世記25章經課選讀不是一開始就叫我們批評以掃；它先讓我們看見,雅各與以掃還未出世,鬥爭已經開始.兩個兄弟在母腹之中相爭,出生時一個先出來,另一個抓住哥哥的腳跟.顯然這不是一場普通兄弟爭執,我們直看下去,可以睇到一段關於身份,位置,承繼和家庭糾紛的故事.

學者提醒我們,雅各與以掃的故事從來不只是兩個人的故事.以掃後來與以東民族的起源重疊,雅各後來更加係以色列十二支派先祖ge 父親；兄弟之間的緊張,成為兩大民族之間又似兄弟又似敵人的關係.創世記一路記載的段反鬧劇,由母腹中的掙扎,出生時的抓緊,長子名分的交易,祝福的爭奪,都像在預示一段關係會怎樣被拉長入漫長的歷史裡面.

但以掃不是一個可以輕輕放在反面教材位置的人.從經文中我們看見他有體毛,有顏色,有田野的生活.他從外面回來,疲倦到一個地步,眼前一碗紅豆湯,比任何長遠的身份都更真實.雅各也不是單純的英雄.他懂得安靜的等,懂得看時機,懂得在有限資源中計算下一步；這種生存方式,其實很多香港人都明白,我地有時就會叫呢D 人做香港仔.經文不是叫我們立刻分出誰比較屬靈,而是讓我們看見兩種生活方式,當兩個咁唔同ge 人,放在同一個家庭裡,遲早會遇上一個不能避免的衝突時刻(正如而家好多人講ge 一句說話,所謂冇比較無傷害).

所以當以掃說:「我都就死,這長子名分於我有甚麼益處?」我們聽到的不只是一句 低俗無賴又愚蠢的話,而是一個人被疲乏壓到至只剩眼前的需要.長子名分不是不重要,應許也不是不重要,只是那一刻,它們不容易滿足他即時的需要,不容易回答他的過度的疲倦.

承:長子名分到底有幾重要?

不過我們要明白,長子名分在古代世界不是一個小問題.今日我們多數關心個人前途,職涯選擇,自由戀愛關係,甚至用呢D 概念來理解人生；但古代氏族社會不是這樣.一個人承受長子名分,不只是多分一點家產,而是承擔整個家族的名號,責任,記憶和家族ge 遺傳(legancy).這個身份把上一代與下一代連結起來,牽涉一個家族如何繼續存在於古代艱難ge 世界之中.

如果用今日比較容易明白的類比,好似係睇電視見到ge 一些王朝故事,《權力遊戲》或《龍之家族》那種世界甚至,長子名分不是單純「邊個分多D 」.它關乎誰有資格代表家族,誰承繼家族的歷史,家族ge 遺傳(legancy),誰可能站在家族最關鍵的位置上,有權為家族ge 將來作出決定.當然,創世記不是叫我們投入那種權力遊戲；但這個類比可以幫助我們明白,長子名分不只是私人喜好,而是整個家庭故事的重心.

在亞伯拉罕,以撒,雅各的敘事裡,這份名分更不只是制度,而是與上主的應許連在一起.上主對亞伯拉罕所說的後裔,土地,祝福萬民,不是抽象地留在天上,而是具體地在一個被呼召的家族中被承載和延續.因此,長子名分唔單止關乎一個人,更加係超越個人,而係關乎成個家族點樣延續落去.對亞伯拉罕ge 家族黎講,呢個位置仲承載住上主ge 應許.

所以,那碗紅豆湯與長子名分的交換,比我們想像中複雜.一邊是馬上入口的食物,一邊是一個要用很多年,甚至要用後來歷史才看得見後果的身份.以掃問「有甚麼益處」,不是因為長子名分真的沒有益處,而是因為它的益處太長,太遠,太依賴未來；而他當時面對的,係眼前gwo 一刻最迫切ge 需要.

重要的東西仍然重要,但重要不等於在每一刻都可以被人感覺得到.信仰也是這樣.應許,呼召,身份,上主的同在,這些不是不重要；只是有時人在疲累裡,在壓力裡,在生活催迫之下,未必感受到它們有幾真實.亦真的未必感受到它們有甚麼即時益處.以前gwo 一代講「信耶穌真係可以當飯食咩」——正正是因為他們的生活太窮困,太艱苦,基督信仰未必能帶來即時幫助和滿足,所以我們也不可以隨便當這句說話只是說說笑.

轉:誰對誰錯?

故事去到這裡,我們很自然想問:究竟誰對誰錯?雅各計算,以掃衝動；雅各懂得等待,以掃被眼前需要牽引；雅各似乎走上應許的路線,以掃似乎把自己放在外面.如果只用結果倒推過程,雅各很容易變成正確的一方,以掃很容易變成失敗的一方.

但創世記似乎不容許我們這麼快落結論.經文沒有把雅各寫成因為義而被揀選的人,也沒有把以掃寫成因為惡而被拒絕的人.我們要小心,因為我們太容易把揀選變成德性證明,把失去變成道德懲罰.可是經文沒有這樣整理人物.雅各的行動帶著抓取,甚至有種乘虛而入ge 味道；以掃的行動帶著輕率,卻也帶著一種受困於當下的真實.經文最後說以掃輕看了長子名分,但它沒有因此替雅各的手段洗白.

好多決定其實都係咁樣作出來.眼前感受到ge ,往往比記憶裡面知道重要ge ye 更加有份量.身體頂唔順,情緒撐唔住,壓力太大ge 時候,人自然會先處理眼前gwo 一關.本來重要ge ye 冇消失,只係暫時退到後面.長子名分仲係度,上主ge 應許都仲係度,只係對gwo 刻ge 以掃黎講,一碗湯比呢D 更加真實.所以以掃的「有咩益處」不單是一個計較的問題,也是一個人在疲倦裡面會問出來的問題.有時唔係人刻意放棄重要ge ye ,而係眼前ge 事,實在太佔據整個人,佔據到再冇空間去想其他ye .

這樣,以掃就不只是一個遠古人物.讀到以掃,我們未必只見到一個衝動,世俗的人輕看屬靈福份的失敗哥哥,反而可能反映我們人的景況.有時我們都知道有些東西重要:信仰重要,關係重要,承諾重要,身份重要.但當眼前的疲累,壓力,焦慮實在太埋身,重要的東西就好像突然變得很遠.不是消失了,而是那一刻,我們承擔不到.

合:應許在不完整的家庭裡前行

故事沒有停在紅豆湯,也沒有停在一句「輕看長子名分」.若果你讀過創世記,就會知道創世記繼續往前走.後來雅各要面對逃亡,被欺騙和被算計；以掃亦沒有從故事裡消失,兄弟二人最終再次相遇,一同埋葬父親.

創世記讓我們看見,應許的前行從來不是一條直線.人的選擇有後果,以掃的交易不是無關痛癢,雅各的抓取也不是沒有代價；上主卻沒有因此停下祂的工作.

好多家庭未必有激烈衝突,但總有D 說話唔想提,總有D 事情放jor 好多年.大家照樣見面,照樣食飯,但關係有時始終隔住一段距離.

邊個付出多D ,邊個被了解多D ,邊個一直覺得自己被父母偏心；呢D 事未必講出口,但會一直留係心裡.於是有時人都會忍不住問一句:我咁多年為呢個家庭付出,究竟有乜益處?

其實雅各同以掃離我地唔遠.

創世記冇將呢個家庭寫聖所有問題都解決晒,然後上主先開始工作.相反,上主ge 應許正正係仍然有距離,仍然有舊傷,仍然未完全講得清楚ge 關係裡面繼續前行.

所以講到最後,我想返去以掃問gwo 個問題:「我都就死,這於我有甚麼益處?」

因為有時我地都會用唔同方式問類似ge 問題:祈禱又解決唔到而家ge 問題,咁有咩用?」

「讀經又唔會令份工容易D ,咁有咩益處?」

「返教會又唔能夠即刻改善家庭關係,咁有咩意思?」

「我已經咁辛苦,上主究竟係邊度?」

當眼前ge 壓力太大,太迫切ge 時候,我地未必會用以掃ge 說話,但其實問緊ge ,往往係同一個問題.

「信耶穌唔通有飯食咩?」,然後帶住呢個問題繼續生活,繼續返屋企,繼續面對家人,繼續面對gwo D 一直未真正解開ge 心結.

也許信仰唔係要我地即刻明白每一件事點解發生,亦唔係要我地太快為一段關係下判語.因為好多事情仲未講完,好多關係亦未走到最後.

而創世記提醒我地:故事未完,上主亦未離開.今日我地未必睇得明白,未必感受得到,但上主ge 應許仍然係人ge 故事裡面,係家庭ge 故事裡面,繼續實現.若果你讀過創世記,就會知道創世記繼續往前走.雅各還要面對欺騙帶來的逃亡,在舅父拉班家裡經歷另一種被算計和出賣,也要在夜裡與神,與自己,甚至可以話同兄弟的影子摔跤.以掃也沒有即時從故事裡消失；他會再出現,會與雅各相遇,會接納雅各ge 禮物,也會與雅各一同埋葬父親.

這意味著,應許的前行不是一條直線,不是今天做對,明天就清楚；今天做錯,明天就出局以後都冇得翻身.應許是在一個不整齊的家庭裡前行,在一段充滿偏愛,抓取,疲乏,誤判的歷史裡前行.人的選擇有後果,這一點不能取消；以掃的交易不是無關痛癢,雅各的抓取也不是沒有代價.但應許的故人會失手,但上主唔會因為人一次失手就停低祂ge 工作.

講到呢度,或者我地都會諗起自己屋企.因為一講到兄弟姊妹,一講到比較,一講到邊個得到多D ,邊個好似蝕底D ,我地其實都唔難明白雅各同以掃之間gwo 種張力.

好多家庭未必有激烈衝突,但總有D 說話唔想提,總有D 事情放jor 好多年.大家照樣見面,照樣食飯,但有時關係始終隔住一段距離.

邊個付出多D ,邊個被明白多D ,邊個一直覺得自己被比較；呢D 事未必講出口,但會一直留係心裡.於是有時人都會忍不住問一句:我咁多年為呢個家庭付出,究竟有乜益處?

以掃和雅各,其實離我地唔會太遠.家人之間有時未必會完全反目成仇,但也未必真的和好；只是關係仍在,距離也仍在.創世記冇將呢個家庭寫成後來大家和和氣氣,所有問題都解決晒,然後上主先開始工作.相反,上主ge 應許正正係係呢D 仍然有距離,仍然有舊傷,仍然未完全講得清楚ge 關係裡面繼續前行.呢個唔係叫我地當冇事發生過,亦唔係一句「神有祂ge 計劃」就算數,好似連人ge 傷痛都唔使面對.相反,它提醒我們:人的家庭故事可能很亂,但亂不等於上主不在其中；關係未完全修復,不等於應許已經停止；今日我們感受不到益處,不等於那份呼召,那份身份,那份上主的同在已經沒有價值.

所以講到最後,我想返去以掃問gwo 個問題:「我都就死,這於我有甚麼益處?」

因為有時我地都會唔同ge 方式問類似ge 問題.然後帶住呢個問題繼續生活,繼續返屋企,繼續面對家人,繼續面對gwo D 一直未真正解開ge 心結.

也許信仰唔係要我地即刻明白每一件事點解發生,亦唔係要我地太快為一段關係下個判語.好多事情仲未講完,好多關係亦仲未走到最後.

雅各同以掃ge 故事係咁,我地自己屋企ge 故事好多時亦都係咁.

而創世記提醒我地,故事未完,上主亦未離開.就算今日我地未必睇得明白,未必感受得到,但上主ge 應許仍然係人ge 故事裡面,係家庭ge 故事裡面,繼續實現.



\newpage

\section{直搗羅馬}
\label{sec:805}
\textbf{鄧瑞強牧師 - 直搗羅馬}
\newline
\newline
\hyperref[sec:804]{< < < PREV SERMON < < <}
~
\hyperref[sec:toc]{[返目錄]}
~
\hyperref[sec:806]{> > > NEXT SERMON > > >}


使徒行傳 16:6-10
\newline
\begin{longtable}{cl}
\hline
\hline
章節 & 經文 (和合本修訂版)\\
\hline
16:6 & \begin{tabularx}{0.7\textwidth}{X} 因為聖靈禁止他們在亞細亞講道，他們就經過弗呂家、加拉太一帶地方。 \end{tabularx} \\ \\ \relax
16:7 & \begin{tabularx}{0.7\textwidth}{X} 到了每西亞的邊界，他們想要往庇推尼去，耶穌的靈卻不許。 \end{tabularx} \\ \\ \relax
16:8 & \begin{tabularx}{0.7\textwidth}{X} 他們就越過每西亞，下特羅亞去。 \end{tabularx} \\ \\ \relax
16:9 & \begin{tabularx}{0.7\textwidth}{X} 夜間，有異象向保羅顯現。有一個馬其頓人站著求他說：「請你過來，到馬其頓來幫助我們！」 \end{tabularx} \\ \\ \relax
16:10 & \begin{tabularx}{0.7\textwidth}{X} 保羅既看見這異象，我們就立即設法往馬其頓去，認為神呼召我們傳福音給那裡的人。 \end{tabularx} \\ \\
[1ex]
\hline
\hline
\end{longtable}
講題:直搗羅馬  March Straight on Rome

經課:使徒行傳16:6-10

講員:鄧瑞強牧師

場合:香港中文大學崇基學院神學院容啟東校長紀念樓G/F聖堂

日期:2026年7月19日

各位弟兄姊妹,早安.願神的恩典常與你們同在.

講完四福音,今日講《使徒行傳》,看看耶穌升天後,教會如何發展.教會的發展不是一帆風順的,很多波折,但在神的介入下,也成了很多機遇.今日,看看教會發展的一個關鍵時刻.

先看看經文:使徒行傳16:6-10

16:6 聖靈既然禁止他們在亞細亞講道,他們就經過弗呂家,加拉太一帶地方.

16:7 到了每西亞的邊界,他們想要往庇推尼去,耶穌的靈卻不許.

16:8 他們就越過每西亞,下到特羅亞去.

16:9 在夜間有異象現與保羅.有一個馬其頓人站著求他說:「請你過到馬其頓來幫助我們.」

16:10 保羅既看見這異象,我們隨即想要往馬其頓去,以為上帝召我們傳福音給那裡的人聽.

一般來說,這故事被稱為「馬其頓異象」,很多機構和事工都以此為名.這故事耳熟能詳,今日,我講一下這故事背後涉及的歷史典故及神學意義.

(1)「特羅亞」(Troas)

這事件發生在「特羅亞」.一個關鍵問題:「特羅亞」(Troas)是一個什麼地方?

從字源來說,「特羅亞」是指「特洛伊人的地區」,這地區最重要的城巿,當然是「特洛伊」(Troy).在文學上,很多時「特羅亞」和「特洛伊」是同義的,互用的.在我這講道裡,這兩詞也是互用的.

特羅亞的重要性,在於其地理位置,它位於現在的土耳其.從歐洲的希臘本土出發,要通往亞洲,最直接簡便的方法,便是從水路,越過愛琴海,直達亞洲的特羅亞.這是歐亞貿易的主要通道,因而也成了兵家必爭之地.有點似現時的霍爾木茲海峽,扼著能源貿易的咽喉,成了兵家必爭之地.「特洛伊」成了著名的古戰場,亞洲帝國想征服歐洲,或者歐洲帝國想征服亞洲,也要在「特洛伊」一決雌雄.

當使徒行傳提到特羅亞時,當地人一聽,便會直接聯想到「特洛伊」這個古戰場.每個文化都有些出名的古戰場,一提到那些戰場的名字,不用多講,那古戰場相關的歷史,便會自然浮現在那文化裡的人的腦海裡.例如:在中國文化圈裡,一提到「赤壁」,中國人就會聯想到孫權連同劉備與曹操的大戰；聯想到孔明草船借箭,借東風,火燒連環船；聯想到梟雄曹操敗走華容道,義氣大過天的關羽放走他；當然也會想到蘇東坡的《赤壁懷古》.

一講到特羅亞,想到特洛伊,那裡的人便會想起「特洛伊戰爭」,即木馬屠城記的故事.那是一場神與人共同參與的戰爭,有點似中國《封神榜》武王伐紂的故事,神與人共同參戰.在木馬屠城記的故事裡,有個神明人物,不得不講,他叫「阿基里斯」,是希臘第一勇士.他刀槍不入,但有一死門,就是腳跟,後人稱腳跟上的腱為「阿基里斯腱」.早前,我就斷了「阿基里斯腱」,故我一生都記得他.

話說回來,「特洛伊戰爭」的本質,當然不是為了爭奪美女Helen,而是歐洲的希臘與亞洲的特洛伊為爭奪貿易的黃金航道而爆發的地緣戰爭.結果是:歐洲的希臘贏了.

作為古戰場,「特羅亞」或「特洛伊」令人想起的,是戰爭,是王者之戰.戰勝的,便能一統歐亞,一統當時的天下,成為王者中的王者.「馬其頓異象」這段經文,背後講的,是王者之戰.

(2)那位「馬其頓人」

在講希臘話的文化圈裡,一講到「一位馬其頓人」,他們即刻會聯想到亞歷山大大帝.正如我們一聽到「烏江」,便想到西楚霸王.一講起「曲阜」,便想到孔子.事實上,經文提及的「特羅亞」,是指特羅亞地區的一座港口城巿,其全名正是「亞歷山大特羅亞」(Alexandria Troas),離開特洛伊古城約25公里.

古代的帝國,除了埃及,便是亞洲的帝國,不計中國,便是亞述,巴比倫,波斯等,西方帝國入侵亞洲帝國,亞歷山大首開先河,他滅了波斯,一路殺到去印度的邊界.正是這一馬其頓人,一統歐亞,一統當時的天下.

來到亞洲,進軍波斯之前,亞歷山大刻意去「特洛伊」獻祭.他聲稱自己是希臘那位刀槍不入的勇士「阿基里斯」的後裔,將這場即將爆發的歐亞大戰宣傳為「新的特洛伊戰爭」,這大大激勵了士氣,最後,這歐洲力量真的打敗了亞洲的霸主波斯.

在第一世紀猶太歷史學家約瑟夫(Flavius Josephus)的《猶太古史》(The Antiquities of the Jews)裡,記載了亞歷山大大帝的一件軼事.話說,亞歷山大大帝進軍亞洲時,順道進軍耶路撒冷,人人以為,這將會屠城.怎知,入城後,亞歷山大見到大祭司Jaddua,竟向他致敬,並讚頌大祭司冠冕上寫著的神的名字耶和華.眾人很驚訝,亞歷山大大帝為何如此做呢?原來,當亞歷山大還在馬其頓,正猶疑如何進軍亞洲時,在夢中,有一人向他顯現,叫他勇敢渡海,這人還承諾會帶領他的軍隊,奪取波斯.這人穿的服飾,正是這耶路撒冷大祭司的服飾,亞歷山大一眼便看出來.他認為,正是這大祭司敬拜的神,幫助他得勝.在這故事裡,亞歷山大這個馬其頓人感謝耶和華,認為是耶和華幫助了他.(參:約瑟夫:《猶太古史》11.8.5)

對比一下歷史軼事與聖經經文:這位馬其頓人,在夢中,看見一個人指示他,進軍亞洲,幫助他一統天下.如今,保羅,在夢中,看見一個馬其頓人,求他往歐洲的馬其頓.我相信,經文的潛台詞是:保羅的福音,將橫跨亞歐,一統天下.

關於經文,有兩點值得一提.一,在《使徒行傳》裡,一般而言,是用「聖靈」來稱呼聖靈的,但在這裡,卻是刻意用「耶穌的靈」來稱呼聖靈.或許,這是因為要對應亞歷山大的夢境,在那夢境中,那人是耶和華的僕人,能被亞歷山大認出來.如今,要讓聖靈被認出,需要一可辨認的「面貌」,這是「耶穌」的靈.二,經文裡,那位馬其頓人求保羅去幫助他們,一反亞歷山大大帝征服他人的霸氣.亞歷山大一統歐亞是靠武力,如今,一統亞歐的,是福音.真正讓天下歸一的,是神恩典的福音.

(3)特洛伊與羅馬

話說,「特洛伊」被毀後,王族後人「艾尼亞斯」(Aeneas)往西走,逃到歐洲,成了被狼養大的一對孖生兄弟的祖先,這對兄弟後來建立了羅馬城.在羅馬的敘事裡,「羅馬」就是「特洛伊」的延續.你可能覺得奇怪,羅馬為何認祖歸宗於「特洛伊戰爭」的失敗者.這是由於他們要宣稱自己源於一不同於希臘文明,卻又與希臘文明平起平坐的血統.而後來當羅馬征服希臘時,這被宣傳成「特洛伊」後人對希臘人的復仇之戰,也給予羅馬對希臘的統治以合法性.說了這麼多,簡單來說,談論「特洛伊」,終極來說,就是談論「羅馬」.從「特洛伊」往西行,越過馬其頓後,就直指羅馬.

講起羅馬,不得不講凱撒大帝.話說當年,他打遍高盧地區,兵威極盛,羅馬的元老們怕他威脅到自己的利益,於是,下令解除他的兵權.主前49年,凱撒連同他的大軍,來到「盧比孔河」(Rubicon River)河邊.這河被視為一條界線,任何軍隊未獲授權而越過這界線,被視為背叛羅馬.放下軍權過河,返回羅馬城,凱撒預計元老院會處死他.不放下軍權而過河,則會與羅馬帝國的主力正面交鋒.何去何從?

在「盧比孔河」河邊,凱撒到底想什麼,是很多文學家發揮想像的題材.有一個傳統這樣說:當凱撒猶疑之際,忽見異象,神明指示他帶著大軍過河.於是,凱撒講出那句名垂千古的話:「the die is cast」(孤注一擲),一切交由命運決定了.之後,他越過界線,直指羅馬.

(4)《使徒行傳》這段經文,到底講什麼?

在文學上,「特羅亞」即「特洛伊」,這是古戰場,是王者爭雄的地方.在這裡,有很多祭壇,作戰前,任何一方都要求問神明,求神指引,求神幫助.

在這裡,最出名的戰爭,是古代的「特洛伊戰爭」,西方的馬其頓戰勝了亞洲的特洛伊.

馬其頓的亞歷山大大帝曾在這裡獻祭,渴求神明指引.但真正在夢中指引他的,是他在馬其頓時夢見的耶和華的僕人.這僕人預告,他的軍隊必能戰勝亞洲的波斯.真正掌管歷史的,是耶和華.

馬其頓人亞歷山大以武力征服天下,但真正統一天下的,《使徒行傳》說,是福音,是耶穌的福音.馬其頓人真正需要的,不是暴力,而是和平的福音.

「特洛伊」是「羅馬」的根,談論「特洛伊」,就不能不談「羅馬」.凱撒大帝在夢中,得神明指引,越過不應超越的界線,直指羅馬.如今,在《使徒行傳》,保羅得耶穌的靈的指引,越過亞歐的界線,到達福音前所未到的領域.從「特洛伊」往西行,便直指羅馬了.

在保羅時代,羅馬是一統地中海的天下的帝國,它以武力主宰著歷史的軌跡.《使徒行傳》卻說,在萬有之上,其實是耶穌在掌權,他才是萬王之王.他指引著他的僕人,突破人間的界限,以仁愛的福音,直取羅馬.真正能令天下一家的,不能靠武力,只能靠福音.到《使徒行傳》結尾時,保羅就在羅馬,傳講天國——人間真正美善的國度.

願神賜福你們.願榮耀歸予聖父,聖子,聖靈.

\allsectionsfont{\centering}

\setlength\parindent{0pt}
\setlength{\columnsep}{1.25em}
\setlength{\parfillskip}{0pt}
\setlength{\tabcolsep}{1em}
\raggedbottom

\pagenumbering{gobble}


\newfontfamily\leftfont[Path=../fonts/fell_french_canon/, Ligatures=TeX, ItalicFont=IMFeFCit29C.otf, BoldFont=AveriaLibre-Bold.ttf]{IMFeFCrm29C.otf}
\newfontfamily\leftcitationfont[Path=../fonts/frankruehl/]{FrankRuehlCLM-Medium.ttf}
\newfontfamily\centerfont[Path=../fonts/garamond/, Ligatures=TeX, ItalicFont=EBGaramond-SemiBoldItalic.ttf]{EBGaramond-SemiBold.ttf}
\newfontfamily\rightfont[Path=../fonts/averia/, Ligatures=TeX, ItalicFont=AveriaLibre-RegularItalic.ttf, BoldFont=AveriaLibre-Bold.ttf, BoldItalicFont=AveriaLibre-BoldItalic.ttf]{AveriaLibre-Light.ttf}
\newfontfamily\rightcitationfont[Path=../fonts/rashi/]{Mekorot-Rashi.ttf}
\definecolor{hcolor}{HTML}{D3230C}
\definecolor{rcolor}{HTML}{D36F0C}
\newcommand{\chfont}[1]{\centerfont{\huge\textcolor{hcolor}{#1}}}
\newcommand{\leftcitation}[1]{\leftcitationfont{\Large\textcolor{hcolor}{#1}}}
\newcommand{\rightcitation}[1]{\rightcitationfont{\normalsize\textcolor{rcolor}{#1}}}
\newfontfamily\flowerfont[Path=../fonts/fell_flowers/]{IMFeFlow2.otf}

\begin{sloppypar}

\chapter*{\chfont{編按結語}}

\columnratio{0.5,0.5}\begin{paracol}{2}

\fontsize{11}{13}\leftfont \Large \leftcitation{א} \leftfont 余少好文．宏志博覽群書而不忘．善存藏經籍文獻備後時之用。\leftcitation{ב} \leftfont 歸主年時．受友所薦．聞道網海．\switchcolumn\fontsize{11}{13}\rightfont \Large \leftcitation{ח} \rightfont 有見粵道之危．國之封講道千言亦將就至．急之何則為？\leftcitation{ט} \rightfont 嘗聞猶太者之傳承．在其力守口述之

\end{paracol}


\columnratio{0.32,0.32,0.32}\begin{paracol}{3}

\fontsize{11}{13}\leftfont \Large 尤以吳約翰遜者 \switchcolumn[2]\fontsize{11}{13}\rightfont \Large 統．以煉千載不

\end{paracol}

\columnratio{0.32,0.32,0.32}
\begin{paracol}{3}\fontsize{11}{13}\leftfont \Large 為重．其載上之粵語講道緩緩入耳．收之藏其音頻．善妥整存．反復而嚼．受益無窮。\leftcitation{ג} \leftfont 我城我國既限．歷一四一九之不測．肺疫延年．信徒靈長屢受圍創．神州燈臺數盡指日可待．粵道之求與日俱增。\leftcitation{ד} \leftfont 觀乎社、經、法、媒、言、信、網之地．愈趨受鋤．自翔不果．授受壓力．粵道聖言亦愈漸艱難。\leftcitation{ה} \leftfont 況崇基例乎．學苑講道屢逆權勢者．其言末強受壓．舊章盡刪以存其身。講道釋數失傳．徒嘆奈何。

\switchcolumn

\fontsize{11}{13}\centerfont 
\begin{tikzpicture}
    \node (0,0) [xshift=-0.10cm, yshift=-1.0cm, opacity=0.10]{\includegraphics[width=0.30\textwidth]{../ot_frontcover.png}} ;
    \node (0,0) [xshift=+0.20cm, yshift=+2.0cm, opacity=0.10]{\includegraphics[width=0.20\textwidth]{../christ_on_cross.png}} ;
\end{tikzpicture}
\Large 

\leftcitation{ס} \centerfont 詩百又廿七載：
\leftcitation{ע} \centerfont 非耶和華建屋宇.則匠人之經營徒.
\leftcitation{פ} \centerfont 非耶和華衛城邑.則守者之儆醒徒.
\leftcitation{צ} \centerfont 余獻是卷予華人社區．願為福音流通之器．願獻斯微材為祭榮耀上帝．
\leftcitation{ק} \centerfont 阿門

\switchcolumn

\fontsize{11}{13}\rightfont \Large 滅．時越次聖殿期及當今。\leftcitation{י} \rightfont 猶太者力廣納之．筆錄以卷軸．便以傳、閱、頌、攜、守、鎖、抄、譯、釋、編，得書塔木德、密示拿等經傳．家喻戶曉．傳流若芳。\leftcitation{כ} \rightfont 猶太者文以載道．傳其口述．今我輩粵道之傳應當作如是．遂力行粵音識辨之法．載言載道．以盡忠傳粵道以待興。\leftcitation{ל} \rightfont 蒙下賜恩惠．無畏海量字音文書．既馭上帝之道．今廣及粵語講道．重駛編程之技．匯導粵音遂字稿．重塑講道現場．以傚猶太卷軸之舉便以傳流。\leftcitation{מ} \rightfont 是卷乃粵音口述傳之屬．莫通華文白話之語．

\end{paracol}

\columnratio{0.5,0.5}
\begin{paracol}{2}\fontsize{11}{13}\leftfont \Large \leftcitation{ו} \leftfont 斯殺一違儆百逆．既禁壓之．我輩聞風無奈．在所難免。\leftcitation{ז} \leftfont 另有異人例乎．以版權之名．脅網絡頻道之舉．同授礙予粵道之存流。

\switchcolumn

\fontsize{11}{13}\rightfont \Large 惟待後繼來者之傚．以譯釋傳之於神州華文地。\leftcitation{נ} \rightfont 今能排程驅馭圖靈以編彙文檔，其碼長共數千千亦無逢大礙．全蒙上帝保守。

\end{paracol}



\columnratio{1}\begin{paracol}{1}

\fontsize{11}{13}\rightfont \Large
~~~~~~~~~~~~~~~~~~~~~~~~~~~~~~~~~~~~~~~~~~~~~~~~~~~~~~~~~~~~~~~~~~~~~~~~~~~~~~~\leftcitation{ר} \rightfont 二零二三年二月一日

~~~~~~~~~~~~~~~~~~~~~~~~~~~~~~~~~~~~~~~~~~~~~~~~~~~~~~~~~~~~~~~~~~~~~~~~~~~~~~~\leftcitation{ש} \rightfont 米迦勒

~~~~~~~~~~~~~~~~~~~~~~~~~~~~~~~~~~~~~~~~~~~~~~~~~~~~~~~~~~~~~~~~~~~~~~~~~~~~~~~\leftcitation{ת} \rightfont 書於香港

\end{paracol}

\end{sloppypar}
\end{document}
